% Seminar in Algebraic Geometry - complete cumulative English LaTeX.
% Nine complete native documents; use split_cumulative_tex.py before compilation.
% Dependencies, native modular projects and PDF reconstruction are in the accompanying source ZIP.
% Correction integration and assembly: OpenAI Codex - GPT-6 Astra, Ultra effort.

% SGA-CUMULATIVE-BEGIN 01_SGA1_EN_COMPLETE_PUBLISHED_SCOPE_READER.tex
% Complete English SGA 1 editable source. Every native text input is embedded.
% Compile in a fresh directory with a standard complete TeX distribution.
% The companion source ZIP preserves the original modular files and build scripts.
% Original input: drafts/SGA1_III_1_4_THROUGH_1_9_English.texfrag
\begin{filecontents*}[overwrite,noheader]{sga1-cumulative-input-001.tex}
% SGA 1, Exposé III, French TeX lines 4648--4810
% (Proposition III.1.4 through Proposition III.1.9 and its proof).

\begin{proposition}
\label{III.1.4}
Let \(A\to B\) be a local homomorphism of local rings with finite residue
extension, and let \(A'\) be a finite local \(A\)-algebra. Then
\(B'=B\otimes_AA'\) is finite over~\(B\), and is therefore a Noetherian
semilocal ring.
\begin{enumerate}
\item[\textup{(i)}] If \(B\) is formally smooth over~\(A\), then the
  localizations of~\(B'\) at its maximal ideals are formally smooth
  over~\(A'\).
\item[\textup{(ii)}] The converse holds if \(A'\) is free over~\(A\).
\end{enumerate}
\end{proposition}

We reduce at once to the case in which \(A\) and~\(B\) are complete.

For \textup{(i)}, let \(A''\) be a finite free local extension of~\(A\)
such that the local components of
\(B''=B\otimes_AA''\) are formal power-series algebras over~\(A''\).
Make the scalar extension
\[
  A''\longrightarrow A''\otimes_AA'\longrightarrow A''',
\]
where \(A'''\) is a local component of~\(A''\otimes_AA'\). The local
components of
\[
  B''\otimes_{A''}A'''=B\otimes_AA'''
\]
are then formal power-series algebras over~\(A'''\). We also have
\[
  B\otimes_AA'''
  =(B\otimes_AA')\otimes_{A'}A'''
  =B'\otimes_{A'}A'''.
\]
Moreover, since \(A''\) is free over~\(A\), the algebra
\(A''\otimes_AA'\) is free over~\(A'\); its direct factor \(A'''\) is
therefore free over the local ring~\(A'\). This proves that the local
components of~\(B'\) are formally smooth over~\(A'\).

For \textup{(ii)}, let \(A''\) be a finite free local
\(A'\)-algebra such that the local components of
\[
  B'\otimes_{A'}A''=B\otimes_AA''
\]
are formal power-series algebras over~\(A''\). Since \(A'\) is free
over~\(A\), the algebra \(A''\) is also free over~\(A\), and hence \(B\)
is formally smooth over~\(A\).

\begin{proposition}
\label{III.1.5}
Let \(A\to B\) be a local homomorphism of local rings with
\emph{trivial} residue extension. Then \(B\) is formally smooth over~\(A\)
if and only if \(\overline B\) is isomorphic to a formal power-series
algebra over~\(\overline A\).
\end{proposition}

Only necessity remains to be proved, and we may suppose that \(A\) and
\(B\) are complete. Let \(\mathfrak m\) and \(\mathfrak n\) be the
maximal ideals of \(A\) and~\(B\), respectively, and choose elements
\(t_1,\dots,t_n\in\mathfrak n\) whose classes form a basis of the vector
space
\[
  \bigl(\mathfrak n/\mathfrak n^2\bigr)
  \big/\operatorname{Im}\bigl(\mathfrak m/\mathfrak m^2\bigr)
  =\mathfrak n/(\mathfrak n^2+\mathfrak mB).
\]
These elements define a local homomorphism of \(A\)-algebras
\[
  B_1=A[[t_1,\dots,t_n]]\longrightarrow B.
\]
\marginpar{\raggedright\footnotesize French source p.~61}%
We shall prove that it is an isomorphism. It is enough to prove, for every
power \(\mathfrak m^q\), that reduction modulo~\(\mathfrak m^q\) gives an
isomorphism
\[
  B_1/\mathfrak m^qB_1\longrightarrow B/\mathfrak m^qB,
\]
since \(B_1\) and~\(B\) are the projective limits of the corresponding
rings modulo~\(\mathfrak m^q\), as \(q\) varies. The \(A\)-modules \(B\)
and~\(B_1\) are flat. Consequently, their associated graded rings for the
\(\mathfrak m\)-adic filtration are obtained by tensoring
\(B_1/\mathfrak mB_1\), respectively \(B/\mathfrak mB\), over
\(k=A/\mathfrak m\) with \(\operatorname{gr}(A)\). We are thus reduced to
showing that
\[
  B_1/\mathfrak mB_1\longrightarrow B/\mathfrak mB
\]
is an isomorphism. In view of Lemma~\ref{III.1.3}, this reduces us to the
case in which \(A\) is a field~\(k\).

On the other hand, let \(A'\) be a finite free local \(A\)-algebra such
that \(B\otimes_AA'\) is a formal power-series algebra over~\(A'\). (This
algebra is local, since the residue extension of~\(B\) over~\(A\) is
trivial.) To prove that \(B_1\to B\) is an isomorphism, it is enough to
prove that
\[
  B_1\otimes_AA'\longrightarrow B\otimes_AA'
\]
is an isomorphism. This reduces us to the case in which \(B\) is already a
formal power-series algebra. This reduction should in fact have been made
before reducing to the case of a base field. But then \(B\) is a regular
local ring with coefficient field~\(k\), and it is well known---and follows
immediately by considering the associated graded rings for the
\(\mathfrak n_1\)-adic and \(\mathfrak n\)-adic filtrations on \(B_1\)
and~\(B\), where \(\mathfrak n_1\) is the maximal ideal of~\(B_1\)---that
\(B_1\to B\) is an isomorphism. This completes the proof.

\begin{corollary}
\label{III.1.6}
If \(B\) is formally smooth over~\(A\), then there exists a finite local
\(A\)-algebra \(A'\) such that the local components of
\[
  \overline B\otimes_{\overline A}\overline{A'}
  =\overline{B\otimes_AA'}
\]
are isomorphic to formal power-series algebras over~\(\overline{A'}\).
\end{corollary}

Indeed, let \(L/k\) be the residue extension of~\(B/A\), and choose a
finite extension \(k'/k\) such that the residue extensions of the local
components of the \(k'\)-algebra \(L\otimes_kk'\) are trivial. Let \(A'\)
be a finite free \(A\)-algebra such that
\(A'/\mathfrak mA'=k'\). Such an algebra exists: one may reduce step by
step to the case in which \(k'/k\) is generated by one element, and then
lift to~\(A\) the coefficients of the minimal polynomial of a generator
of~\(k'\) over~\(k\). The algebra \(A'\) is local. At each maximal ideal
of \(B\otimes_AA'\), the residue extension over the residue field \(k'\)
of~\(A'\) is trivial, and Proposition~\ref{III.1.5} completes the proof.

\begin{corollary}
\label{III.1.7}
Let \(A\to B\) be a local homomorphism of local rings. Then \(B\) is
formally smooth over~\(A\) if and only if \(B\) is flat over~\(A\) and
\(B/\mathfrak mB\) is formally smooth over \(k=A/\mathfrak m\).
\end{corollary}

After making a suitable finite free local extension \(A'\) of~\(A\) and
using Proposition~\ref{III.1.4}\textup{(ii)}, we reduce to the case in
which the residue extension of~\(B/A\) is trivial.
\marginpar{\raggedright\footnotesize French source p.~62}%
Proposition~\ref{III.1.4}\textup{(i)} and
Lemma~\ref{III.1.3} show that the stated conditions are necessary. For
sufficiency, it is enough to observe that, under the present hypotheses,
the proof of
Proposition~\ref{III.1.5} shows that \(B\) is a formal power-series algebra
over~\(A\), after replacing \(A\) and~\(B\) by their completions, as we
may.

\refstepcounter{proposition}
\paragraph{Remark \theproposition.}\label{III.1.8}
It would not be difficult to develop, for formally smooth homomorphisms,
analogues of all the properties of smooth morphisms studied in
\hyperref[II]{Exposé~II}.
For differential properties, however, this requires modifying the usual
definition of Kähler differentials (cf.\ \ifthenelse{\boolean{orig}}{I~}{}\ref{I.1}), with completed
tensor products replacing ordinary tensor products. We shall merely
gesture toward these abysses here; the preceding results suffice for our
purposes.

It remains to connect formal smoothness with the notion of smoothness
developed in \hyperref[II]{Exposé~II}, which we have not yet used at all:

\begin{proposition}
\label{III.1.9}
Let \(A\to B\) be a local homomorphism, where \(B\) is a localization of
an \(A\)-algebra of finite type. Then \(B\) is smooth over~\(A\) if and
only if it is formally smooth over~\(A\).
\end{proposition}

Using Corollary~\ref{III.1.7} and Theorem~\ref{II.2.1}, we reduce to the
case in which \(A\) is a field~\(k\).

Using Proposition~\ref{III.1.4}\textup{(ii)} and
Corollary~\ref{II.4.13}, a suitable extension \(k'/k\) reduces us to the
case in which the residue extension of~\(B/k\) is trivial. By
Corollary~\ref{II.5.2}, respectively Proposition~\ref{III.1.5}, the
algebra \(B\) is smooth over~\(k\), respectively formally smooth over~\(k\),
if and only if \(B\) is a regular local ring, respectively its completion
is a formal power-series algebra over~\(k\). These two conditions are well
known to be equivalent when the residue extension is trivial.
\end{filecontents*}
% Original input: drafts/SGA1_III_1_OPENING_THROUGH_1_3_English.texfrag
\begin{filecontents*}[overwrite,noheader]{sga1-cumulative-input-002.tex}
\section{Formally smooth homomorphisms}
\label{III.1}

In \hyperref[II]{Exposé~II} we confined ourselves to homomorphisms of finite type and,
consequently, in the case of local homomorphisms \(A\to B\) of local
rings, to the case where \(B\) is isomorphic to a localization of a
finitely generated \(A\)-algebra. This case is insufficient for various
applications, especially in formal or analytic geometry. For example,
the formal power-series ring
\(B=A[[t_1,\dots,t_n]]\) has, from the viewpoint of formal geometry,
the properties of a smooth algebra over~\(A\). Likewise, in analytic
geometry, the local ring at a point \((y,z)\) of a product
\(Y\times\mathbb C^n\), regarded as an algebra over the local ring
at~\(y\), has these properties. Indeed, its completion is isomorphic to
the algebra of formal power series in \(n\) indeterminates over the
completion of the base ring \(\mathcal O_y\). This leads to the
following definition.

\begin{definition}
\label{III.1.1}
Let \(u\colon A\to B\) be a local homomorphism of local rings, which,
as always, are assumed Noetherian. Suppose that the residue-field
extension \(\kappa(B)/\kappa(A)\) is finite. We say that \(u\) is a
\emph{formally smooth homomorphism}, or that the algebra \(B\) is
\emph{formally smooth over}~\(A\), if there exists a finite local
\(\overline A\)-algebra \(A'\), free over~\(\overline A\), such that
the local components of the semilocal ring
\[
  B'=\overline B\otimes_{\overline A}A'
\]
are \(A'\)-isomorphic to formal power-series algebras over~\(A'\).
\footnote{For a more general and more conceptual definition, motivated
by Theorem~\ref{III.2.1} below, see
EGA~\(\mathrm{0}_{\mathrm{IV}}\), 19.3.1.}
\end{definition}

Here \(\overline A\) and \(\overline B\) denote the completions of
\(A\) and~\(B\). Since \(B'\) is finite and free over~\(\overline B\),
it is indeed a semilocal ring, a direct product of complete local rings.
Each of these is again a free \(\overline B\)-module, and therefore has
the same dimension as \(\overline B\), hence as~\(B\). It follows that
the number of variables \(t_i\) in the formal power-series rings of
Definition~\ref{III.1.1} is
\[
  \dim\overline B-\dim\overline A=\dim B-\dim A,
\]
and in particular is independent of the chosen local component.
\marginpar{\raggedright\footnotesize French source p.~59}%
It is immediately seen to be also the dimension of the ring
\(B\otimes_A k=B/\mathfrak mB\), where
\(k=A/\mathfrak m\) is the residue field of~\(A\). We call it the
\emph{relative dimension of \(B\) over \(A\)}.

\begin{remarques}
\label{III.1.2}
Definition~\ref{III.1.1} plainly depends only on the homomorphism of
completions \(\overline A\to\overline B\) induced by \(A\to B\), which
to some extent justifies the terminology. We here retract the
terminology of \ifthenelse{\boolean{orig}}{I~}{}\ref{I.3.2}\textup{(b)} and
\ifthenelse{\boolean{orig}}{I~}{}\ref{I.4.1}\textup{(b)}, which risks being misleading, and prefer to
say \lq\lq formally unramified\rq\rq{} and
\lq\lq formally étale\rq\rq{} in the cases treated by those
definitions. We reserve \lq\lq unramified\rq\rq{} and
\lq\lq étale\rq\rq{} for the case where \(B\) is a localization of a
finitely generated \(A\)-algebra.
\footnote{Or, better, \lq\lq essentially unramified\rq\rq{} and
\lq\lq essentially étale\rq\rq{}; compare EGA~IV, 18.6.1.}

The reader will verify at once that \lq\lq formally étale\rq\rq{} is
equivalent to \lq\lq formally smooth and quasi-finite\rq\rq{}.
Finally, there is a reasonable definition of
\lq\lq formally smooth\rq\rq{} without any preliminary hypothesis on
the residue extension \(\kappa(B)/\kappa(A)\), assumed finite here.
Among other cases, it includes local homomorphisms \(A\to B\) for which
\(B\) is flat over~\(A\) and \(B/\mathfrak mB\) is a separable
extension of \(A/\mathfrak m=k\), not necessarily of finite type. For
example, a Cohen \(p\)-ring is formally smooth over the ring of
\(p\)-adic integers. In this general setting, the lifting property for
homomorphisms (compare Theorem~\ref{III.2.1}) should be taken as the
definition. For the applications we have in mind, the case treated in
Definition~\ref{III.1.1} will suffice. Henceforth,
\lq\lq formally smooth\rq\rq{} will tacitly mean
\lq\lq with finite residue extension\rq\rq{}.
\end{remarques}

\begin{lemma}
\label{III.1.3}
If \(B\) is formally smooth over~\(A\), then \(B\) is flat over~\(A\).
\end{lemma}

Since flatness is invariant under completion, we may suppose that
\(A\) and \(B\) are complete. Since flatness is also invariant under a
flat local, hence faithfully flat, extension of the base ring,
Definition~\ref{III.1.1} reduces us to the case where \(B\) is a formal
power-series algebra over~\(A\). As an \(A\)-module, \(B\) is then
isomorphic to a product of \(A\)-modules isomorphic to~\(A\), and is
therefore \(A\)-flat, since the base ring \(A\) is Noetherian.

We now place ourselves under the hypotheses of
Definition~\ref{III.1.1}. Since the residue extensions of the local
components of \(B'\) over~\(A'\) are trivial,
\marginpar{\raggedright\footnotesize French source p.~60}%
\(L\otimes_k k'\) is an Artinian \(k'\)-algebra whose local components
have trivial residue extensions, where \(L\), \(k\), and \(k'\) are the
residue fields of \(B\), \(A\), and \(A'\), respectively. This
condition, necessary for the finite free extension \(A'\) to satisfy
the requirement in Definition~\ref{III.1.1}, is also sufficient, as
follows immediately from Proposition~\ref{III.1.4}\textup{(i)} and
Proposition~\ref{III.1.5} below.
\end{filecontents*}
% Original input: drafts/SGA1_III_2_COMPLETE_English.texfrag
\begin{filecontents*}[overwrite,noheader]{sga1-cumulative-input-003.tex}
% SGA 1, Exposé III, French TeX lines 4811--5070.
% Complete Section III.2 (Theorem III.2.1 and Corollaries III.2.2--III.2.3).

\input{sga1-cumulative-input-004.tex}

\medskip
\noindent\textup{(iii)}\(\Rightarrow\)\textup{(ii)}.
\marginpar{\raggedright\footnotesize French source p.~65}%
Let \(\mathfrak n\) be the maximal ideal of~\(C\), and, for each
integer \(q>0\), set
\[
  C_q=C/\mathfrak n^q,
  \qquad
  J_q=(J+\mathfrak n^q)/\mathfrak n^q.
\]
Thus \(C_q/J_q\) is naturally a quotient algebra of~\(C/J\). Moreover,
the composite homomorphism
\[
  u_q\colon B\longrightarrow C/J\longrightarrow C_q/J_q
\]
is continuous when \(C_q/J_q\) is given the discrete topology, and
\(J_q\) is a nilpotent ideal of~\(C_q\). We construct successively
\(A\)-homomorphisms
\[
  v_q\colon B\longrightarrow C_q
\]
such that \textup{(a)} \(v_q\) lifts \(u_q\), and \textup{(b)} \(v_q\)
lifts \(v_{q-1}\). The induction can be carried out as follows. The maps
\[
  u_q\colon B\longrightarrow C/(J+\mathfrak n^q),
  \qquad
  v_{q-1}\colon B\longrightarrow C/\mathfrak n^{q-1}
\]
induce the same homomorphism
\[
  B\longrightarrow
  C/\bigl((J+\mathfrak n^q)+\mathfrak n^{q-1}\bigr)
  =C/(J+\mathfrak n^{q-1})=C_{q-1}/J_{q-1},
\]
namely \(u_{q-1}\). They therefore determine a homomorphism
\[
  B\longrightarrow C/J'_q,
  \qquad
  J'_q=(J+\mathfrak n^q)\cap\mathfrak n^{q-1}
  \supseteq\mathfrak n^q,
\]
from which both are obtained by reduction. It remains only to lift this
homomorphism from~\(B\) to the quotient of~\(C_q\) by
\(J'_q/\mathfrak n^q\). This ideal is contained in~\(J_q\), hence is
nilpotent, so the required lift exists by hypothesis~\textup{(iii)}.

The maps \(v_q\) now determine a homomorphism from~\(B\) to the inverse
limit \(C\) of the rings~\(C_q\). Since \(J\) is closed, it is the
inverse limit of the ideals~\(J_q\); consequently this homomorphism
lifts~\(u\).

\medskip
\noindent\textup{(iv)}\(\Rightarrow\)\textup{(i)}.
First observe that if condition~\textup{(iv)} holds, then it also holds
over~\(A'\) for every local component of \(B\otimes_AA'\), whenever
\(A'\) is a finite local \(A\)-algebra. Choose \(A'\) free over~\(A\)
and such that the residue extensions of the local components of
\(B'=B\otimes_AA'\) over~\(A'\) are trivial. Proposition~\ref{III.1.4}
\textup{(ii)} then reduces us to the case in which the residue extension
of~\(B\) over~\(A\) is trivial. We shall prove the following slightly
more precise result.

\begin{corollary}
\label{III.2.2}
Under the hypotheses of Theorem~\ref{III.2.1}, suppose in addition that
the residue extension of~\(B\) over~\(A\) is trivial. Then the equivalent
conditions of Theorem~\ref{III.2.1} are also equivalent to each of the
following two conditions (with \(A\) and~\(B\) assumed complete in
\textup{(v)}):
\marginpar{\raggedright\footnotesize French source p.~66}%
\begin{enumerate}
\item[(iv ter)] The same assertion as in \textup{(iv)}, but with the
  finite local Artinian \(A\)-algebra \(C\) required to have trivial
  residue extension (and, if desired, with \(J^2=0\)).
\item[(v)] There is a local \(A\)-homomorphism, where
  \[
    n=\dim\bigl(\mathfrak n/(\mathfrak n^2+\mathfrak mB)\bigr),
  \]
  \[
    u\colon B\longrightarrow B_1=A[[t_1,\dots,t_n]],
  \]
  which induces an isomorphism
  \[
    \mathfrak n/(\mathfrak n^2+\mathfrak mB)
    \xrightarrow{\sim}
    \mathfrak n_1/(\mathfrak n_1^2+\mathfrak mB_1),
  \]
  where \(\mathfrak n\) and \(\mathfrak n_1\) are the maximal ideals
  of \(B\) and \(B_1\), respectively, and \(\mathfrak m\) is the
  maximal ideal of~\(A\).
\end{enumerate}
\end{corollary}

\subsubsection*{Proof}
Since \textup{(iv bis)} plainly implies \textup{(iv ter)}---even without
using the square-zero proviso---it is enough to prove
\[
  \textup{(iv ter)}\Rightarrow\textup{(v)}\Rightarrow\textup{(i)}.
\]

\noindent\textup{(iv ter)}\(\Rightarrow\)\textup{(v)}.
Choose a basis \(a_1,\dots,a_n\) of
\(\mathfrak n/(\mathfrak n^2+\mathfrak mB)\). This determines a local
homomorphism of \(A\)-algebras
\[
  B\longrightarrow
  B_1/(\mathfrak n_1^2+\mathfrak mB_1)
  =k[t_1,\dots,t_n]/(t_1,\dots,t_n)^2.
\]
By \textup{(iv ter)} it can be lifted successively to homomorphisms of
\(A\)-algebras from~\(B\) to
\(B_1/\mathfrak n_1^2,B_1/\mathfrak n_1^3,\dots\). Passing to the
inverse limit gives a homomorphism \(B\to B_1\) having the required
property.

\noindent\textup{(v)}\(\Rightarrow\)\textup{(i)}.
In the commutative diagram
\[
\begin{array}{ccccccc}
\mathfrak m/\mathfrak m^2 & \longrightarrow &
\mathfrak n/\mathfrak n^2 & \longrightarrow &
\mathfrak n/(\mathfrak n^2+\mathfrak mB) & \longrightarrow & 0 \\
\big\downarrow && \big\downarrow && \big\downarrow && \\
\mathfrak m/\mathfrak m^2 & \longrightarrow &
\mathfrak n_1/\mathfrak n_1^2 & \longrightarrow &
\mathfrak n_1/(\mathfrak n_1^2+\mathfrak mB_1) & \longrightarrow & 0.
\end{array}
\]
the rows are exact and the two outer vertical arrows are surjective.
The middle vertical arrow is therefore surjective as well; since \(B\)
is complete, it follows that \(u\colon B\to B_1\) is surjective. Choose
elements \(x_i\in B\), \(1\leq i\leq n\), lifting the~\(t_i\). They
define an \(A\)-algebra homomorphism \(B_1\to B\), which is surjective
for the same reason as~\(u\), and whose composite with~\(u\) is the
identity by construction. Thus \(B_1\to B\) is also injective and hence
is an isomorphism. We have therefore also proved the following.

\begin{corollary}
\label{III.2.3}
\marginpar{\raggedright\footnotesize French source p.~67}%
Under the hypotheses of Corollary~\ref{III.2.2}\textup{(v)}, the map
\(u\) is necessarily an isomorphism.
\end{corollary}

This completes the proof that \(B\) is formally smooth over~\(A\). At
the same time we have recovered Proposition~\ref{III.1.5}, although
there is little merit in doing so.
\end{filecontents*}
% Original input: drafts/SGA1_III_2_OPENING_THROUGH_PROOF_I_TO_III_English.texfrag
\begin{filecontents*}[overwrite,noheader]{sga1-cumulative-input-004.tex}
\section[Lifting property]{The characteristic lifting property of
formally smooth homomorphisms}
\label{III.2}

\begin{theorem}
\label{III.2.1}
Let \(A\to B\) be a local homomorphism of local rings inducing a finite
residue-field extension. The following conditions are equivalent:
\begin{enumerate}
\item[(i)]
\marginpar{\raggedright\footnotesize French source p.~63}%
\(B\) is formally smooth over~\(A\).
\item[(ii)] For every local homomorphism \(A\to C\), where \(C\) is a
complete local ring, every ideal \(J\subseteq\mathfrak r(C)\), and
every local \(A\)-homomorphism \(B\to C/J\), there exists an
\(A\)-homomorphism \(B\to C\), necessarily local, which lifts it.
\item[(iii)] For every \(A\)-algebra \(C\), not necessarily
Noetherian, every nilpotent ideal \(J\) of~\(C\), and every continuous
\(A\)-homomorphism \(B\to C/J\), that is, one which vanishes on some
power of \(\mathfrak r(B)\), there exists an \(A\)-homomorphism
\(B\to C\) lifting it; this lift is necessarily continuous as well.
\item[(iv)] The same lifting assertion as in \textup{(ii)} and
\textup{(iii)}, with \(C\) restricted to be a local Artinian ring
finite over~\(A\).
\item[(iv bis)] The same as \textup{(iv)}, with the additional
requirement \(J^2=0\).
\end{enumerate}
\end{theorem}

\paragraph{Remark.}
In the remainder of this Exposé we shall use only the implication
\textup{(iv)}\(\Rightarrow\)\textup{(i)}, or
\textup{(iv bis)}\(\Rightarrow\)\textup{(i)}. The direct implication
\textup{(i)}\(\Rightarrow\)\textup{(ii)} will be proved by another
method in the next section when \(B\) is a localization of a finitely
generated \(A\)-algebra. Recall that in the
\lq\lq good\rq\rq{} theory of Cohen
theorems,\footnote{See EGA~\(\mathrm{0}_{\mathrm{IV}}\), 19.3 and
19.8.} property \textup{(ii)} or \textup{(iii)} becomes the definition
of formally smooth homomorphisms, whereas Definition~\ref{III.1.1}
becomes a characteristic property valid only when the residue
extension is finite.

Notice that neither property \textup{(ii)} nor property \textup{(iii)}
is more general than the other. One could give an equivalent property
encompassing both by taking a linearly topologized ring \(C\), separated
and complete, a closed topologically nilpotent ideal of~\(C\), and a
continuous homomorphism \(A\to C\), thereby making \(C\) a topological
\(A\)-algebra. We leave this modification to the reader.

\subsubsection*{Proof of Theorem~\ref{III.2.1}}

We shall prove
\[
  \textup{(i)}\Longrightarrow\textup{(iii)}
  \Longrightarrow\textup{(ii)}
\]
and then \(\textup{(iv)}\Rightarrow\textup{(i)}\). Since
\(\textup{(ii)}\Rightarrow\textup{(iv)}\) is trivial, and the
equivalence of \textup{(iv)} and \textup{(iv bis)} follows by immediate
induction on an integer \(n\) such that \(J^n=0\), this will complete
the proof.

\paragraph{\textup{(i)}\(\Rightarrow\)\textup{(iii)}.}
An immediate induction reduces us to the case \(J^2=0\). Continuity of
the given map, together with \(J^2=0\), implies that some power of the
maximal ideal \(\mathfrak m\) of~\(A\) acts trivially on~\(C\).
After quotienting by that power of \(\mathfrak m\), and observing from
Proposition~\ref{III.1.4}\textup{(i)} that
\(B/\mathfrak m^qB\) remains formally smooth over
\(A/\mathfrak m^q\), we may suppose that \(A\) is Artinian.
\marginpar{\raggedright\footnotesize French source p.~64}%
By Lemma~\ref{III.1.3}, \(B\) is flat over~\(A\), and hence is free as
an \(A\)-module because \(A\) is Artinian. There is therefore an
\(A\)-linear map
\[
  w\colon B\longrightarrow C
\]
lifting the given homomorphism \(u\colon B\to C/J\). Put
\[
  f(x,y)=w(xy)-w(x)w(y)\qquad(x,y\in B).
\]
Then \(f(x,y)\in J\), so \(f\) is an \(A\)-bilinear map from
\(B\times B\) to~\(J\). A lift \(v\colon B\to C\) of~\(u\) is an
algebra homomorphism if and only if there is an \(A\)-linear map
\(g\colon B\to J\) such that \(v=w+g\) and
\begin{align*}
  g(1)&=1-w(1),\\
  g(xy)-u(x)g(y)-u(y)g(x)&=-f(x,y)
  \qquad(x,y\in B).
\end{align*}
Here the action of \(C/J\) on~\(J\) is well defined because \(J^2=0\).

This is a system of linear equations in \(\operatorname{Hom}_A(B,J)\),
with right-hand sides in~\(J\). It has a solution if and only if the
corresponding system in
\[
  \operatorname{Hom}_A(B,J)\otimes_AA',
  \qquad J'=J\otimes_AA',
\]
has a solution, where \(A'\) is any faithfully flat
\(A\)-algebra. Choose a finite free local \(A\)-algebra \(A'\) such
that \(B'=B\otimes_AA'\) is a formal power-series algebra over~\(A'\);
we may suppose \(A\) and \(B\) complete in this argument. Since \(A'\)
is finite free over~\(A\),
\[
  \operatorname{Hom}_A(B,J)\otimes_AA'
  =\operatorname{Hom}_{A'}(B',J').
\]
The resulting system in \(\operatorname{Hom}_{A'}(B',J')\) is exactly
the one determining the \(A'\)-algebra homomorphisms
\[
  B'\longrightarrow C'=C\otimes_AA'
\]
which lift the scalar extension
\(u'\colon B'\to C'/J'\) of~\(u\), by correcting the scalar extension
\(w'\colon B'\to C'\) of~\(w\) with an \(A'\)-linear map
\(g'\colon B'\to J'\). Notice that \(B\) generates \(B'\) as an
\(A'\)-module.

We are thus reduced to proving \textup{(iii)} when
\[
  B=A[[t_1,\dots,t_n]]
\]
is a formal power-series algebra over~\(A\). Lift the images of the
\(t_i\) in \(C/J\) arbitrarily to elements \(z_i\in C\). Since the
\(z_i\bmod J\) are nilpotent, by continuity of
\(u\colon B\to C/J\), and \(J\) itself is nilpotent, the \(z_i\) are
nilpotent. They therefore define a continuous homomorphism of
topological \(A\)-algebras from \(B\) to the discrete ring~\(C\), which
plainly lifts~\(u\).

\end{filecontents*}
% Original input: drafts/SGA1_III_3_English.texfrag
\begin{filecontents*}[overwrite,noheader]{sga1-cumulative-input-005.tex}
\section[Local infinitesimal lifting of morphisms]{Local infinitesimal
lifting of morphisms into a smooth S-scheme}
\label{III.3}

\begin{theorem}
\label{III.3.1}
Let \(f\colon X\to Y\) be a morphism locally of finite type. The
following conditions are equivalent:
\begin{enumerate}
\item[(i)] \(f\) is smooth.
\item[(ii)] For every prescheme \(Y'\) over~\(Y\), every closed
subprescheme \(Y'_0\subseteq Y'\) having the same underlying space as
\(Y'\), every \(Y\)-morphism \(g_0\colon Y'_0\to X\), and every
\(z\in Y'_0\), there exist an open neighborhood \(U\) of~\(z\) in
\(Y'\) and a \(Y\)-morphism \(g\colon U\to X\) extending
\(g_0|_{Y'_0\cap U}\).
\item[(ii bis)] For \(Y'\), \(Y'_0\), and \(z\) as in \textup{(ii)},
put
\[
  X'=X\times_Y Y',
  \qquad
  X'_0=X\times_Y Y'_0.
\]
Every section of \(X'_0\) over~\(Y'_0\) extends to a section of~\(X'\)
over an open neighborhood \(U\) of~\(z\).
\item[(iii)] For every \(Y\)-scheme \(Y'\) which is the spectrum of a
local Artinian ring finite over some \(\mathcal O_y\), \(y\in Y\),
every nonempty closed subprescheme \(Y'_0\subseteq Y'\), and every
\(Y\)-morphism \(g_0\colon Y'_0\to X\), there exists a \(Y\)-morphism
\(g\colon Y'\to X\) extending~\(g_0\).
\item[(iii bis)] For \(Y'\) and \(Y'_0\) as in \textup{(iii)}, put
\[
  X'=X\times_Y Y',
  \qquad
  X'_0=X\times_Y Y'_0.
\]
Every section of \(X'_0\) over~\(Y'_0\) extends to a section of~\(X'\)
over~\(Y'\).
\end{enumerate}
\end{theorem}

\subsubsection*{Proof}

The equivalence of \textup{(ii)} and \textup{(ii bis)}, and that of
\textup{(iii)} and \textup{(iii bis)}, are immediate, as is the
implication \(\textup{(ii)}\Rightarrow\textup{(iii)}\). It remains to
prove
\[
  \textup{(i)}\Rightarrow\textup{(ii)}
  \qquad\text{and}\qquad
  \textup{(iii)}\Rightarrow\textup{(i)}.
\]

\paragraph{\textup{(i)}\(\Rightarrow\)\textup{(ii)}.}
Put \(x=g_0(z)\). By replacing \(X\) with a suitable open neighborhood
of~\(x\), and \(Y'\) with the prescheme induced on the inverse image of
that open set under~\(g_0\), we may suppose that \(X\) is étale over
\(Y[t_1,\dots,t_n]\). The composite \(Y\)-morphism
\[
  Y'_0\longrightarrow X\longrightarrow Y[t_1,\dots,t_n]
\]
is defined by \(n\) sections of \(\mathcal O_{Y'_0}\). Near~\(z\),
these extend to sections of \(\mathcal O_{Y'}\), so after shrinking we
may suppose that the displayed morphism extends to a \(Y\)-morphism
\[
  Y'\longrightarrow Y[t_1,\dots,t_n].
\]
\marginpar{\raggedright\footnotesize French source p.~68}%
By \ifthenelse{\boolean{orig}}{I~}{}\ref{I.5.6}, there is then a unique \(Y\)-morphism
\(g\colon Y'\to X\) lifting this morphism and at the same time
extending~\(g_0\).

\paragraph{\textup{(iii)}\(\Rightarrow\)\textup{(i)}.}
Since the locus where \(f\) is smooth is open, it is enough to show that
it contains every \(x\in X\) which is closed in its fiber. Put
\(y=f(x)\). Then \(\mathcal O_x\) is an \(\mathcal O_y\)-algebra which
is a localization of a finitely generated algebra and has finite
residue extension. Hypothesis \textup{(iii)} implies that every
homomorphism from \(\mathcal O_x\) to \(A/J\), where \(A\) is a local
Artinian algebra finite over~\(\mathcal O_y\) and \(J\) is an ideal
contained in its radical, lifts to a homomorphism
\(\mathcal O_x\to A\). Here we use the fact that a morphism
\(\operatorname{Spec}(B)\to X\), with \(B\) local and the closed point
mapping to \(x\), is determined bijectively by a local homomorphism
\(\mathcal O_x\to B\). Theorem~\ref{III.2.1} therefore shows that
\(\mathcal O_x\) is formally smooth over~\(\mathcal O_y\), and hence
smooth over~\(\mathcal O_y\) by Proposition~\ref{III.1.9}.

\begin{corollary}
\label{III.3.2}
Let \(f\colon X\to Y\) be as in Theorem~\ref{III.3.1}. The following
conditions are equivalent:
\begin{enumerate}
\item[(i)] \(f\) is étale.
\item[(ii)] Condition \textup{(ii)} of Theorem~\ref{III.3.1} holds,
with the extension \(g\) of~\(g_0\) to~\(U\) unique.
\item[(iii)] Condition \textup{(iii)} of Theorem~\ref{III.3.1} holds,
with \(g\) unique.
\end{enumerate}
\end{corollary}

It is enough to observe in the proof of
\(\textup{(i)}\Rightarrow\textup{(ii)}\) above that uniqueness, when
\(Y'_0\) is not identical to~\(Y'\) near~\(z\), is possible only if
\(n=0\), a condition already known to be sufficient.

\begin{corollary}
\label{III.3.3}
Let \(X\) be a prescheme locally of finite type over a complete local
ring~\(A\), let \(y\) be the closed point of
\(Y=\operatorname{Spec}(A)\), let \(f\colon X\to Y\) be the structural
morphism, and let \(x\in f^{-1}(y)\) be rational
over~\(\kappa(y)\). If \(X\) is smooth over~\(A\) at~\(x\), then there
exists a section \(s\) of \(X\) over~\(Y\) passing through~\(x\), that
is, satisfying \(s(y)=x\).
\end{corollary}

In particular, if \(X\) is smooth over~\(A\), then the natural map
\[
  \Gamma(X/Y)\longrightarrow\Gamma(X\otimes_Ak/k)
\]
from the sections of \(X\) over~\(Y\) to the set of \(k\)-rational
points of the fiber \(f^{-1}(y)=X\otimes_Ak\) is surjective. This fact
was especially well known and used when \(A\) is a discrete valuation
ring and \(X\) is proper over~\(A\), in fact projective. In that case
the sections of \(X\) over~\(Y\), that is, the \(A\)-valued points
of~\(X\), also identify with the rational sections, namely with the
\(K\)-valued points of
\[
  X\otimes_AK=X_K,
\]
which is a proper smooth scheme over the fraction field \(K\) of~\(A\).
\marginpar{\raggedright\footnotesize French source p.~69}%
\end{filecontents*}
% Original input: drafts/SGA1_III_4_English.texfrag
\begin{filecontents*}[overwrite,noheader]{sga1-cumulative-input-006.tex}
\section{Local infinitesimal lifting of smooth S-schemes}
\label{III.4}

\begin{theorem}
\label{III.4.1}
Let \(Y\) be a locally Noetherian prescheme, let \(Y_0\) be a closed
subprescheme with the same underlying space, let \(X_0\) be a smooth
\(Y_0\)-prescheme, and let \(x\in X_0\). There exist an open
neighborhood \(U_0\) of~\(x\), a prescheme \(X\) smooth over~\(Y\), and
a \(Y_0\)-isomorphism
\[
  h\colon U_0\xrightarrow{\sim}X\times_Y Y_0.
\]
Moreover, if \((U'_0,X',h')\) is another solution of this problem, then
it is isomorphic to the first in a neighborhood of~\(x\).
\end{theorem}

We leave it to the reader to make the last assertion precise. For a
fixed \(U_0\), a solution of the problem amounts to giving on~\(U_0\)
a sheaf of algebras \(\mathcal B\) over
\(f_0^{-1}(\mathcal O_Y)\), where \(f_0\) is the continuous map
underlying the structural morphism \(U_0\to Y_0\), together with a ring
homomorphism
\[
  \mathcal B\longrightarrow\mathcal O_{U_0}
\]
compatible with
\(f_0^{-1}(\mathcal O_Y)\to f_0^{-1}(\mathcal O_{Y_0})\), such that:
\begin{enumerate}
\item[(a)] this homomorphism induces an isomorphism
\[
  \mathcal B\otimes_{f_0^{-1}(\mathcal O_Y)}
  f_0^{-1}(\mathcal O_{Y_0})
  \xrightarrow{\sim}\mathcal O_{U_0};
\]
\item[(b)] \(U_0\), equipped with \(\mathcal B\), becomes a smooth
\(Y\)-prescheme.
\end{enumerate}
In this formulation, the precise meaning of the local uniqueness
assertion becomes particularly clear.

\subsubsection*{Proof}

We may first suppose that \(X_0\) is étale over
\[
  Y_0[t_1,\dots,t_n]=Y'_0.
\]
The latter may be regarded as a closed subprescheme, with the same
underlying space, of \(Y'=Y[t_1,\dots,t_n]\). By
\ifthenelse{\boolean{orig}}{I~}{}\ref{I.8.3}, there exist a prescheme \(X\) étale over~\(Y'\) and a
\(Y'_0\)-isomorphism
\[
  X\times_{Y'}Y'_0\xrightarrow{\sim}X_0.
\]
This proves existence. For uniqueness, use the lifting property
in Theorem~\ref{III.3.1}\textup{(ii)} for smooth morphisms, together
with the following lemma.
\marginpar{\raggedright\footnotesize French source p.~70}%

\begin{lemma}
\label{III.4.2}
Let \(Y\) be a prescheme, let \(Y_0\) be a closed subprescheme defined
by a locally nilpotent sheaf of ideals \(\mathcal J\), let \(X\) and
\(X'\) be \(Y\)-preschemes, and let \(u\colon X\to X'\) be a
\(Y\)-morphism. Suppose that \(X\) is flat over~\(Y\). Then \(u\) is an
isomorphism if and only if
\[
  u_0\colon X\times_Y Y_0\longrightarrow X'\times_Y Y_0
\]
is an isomorphism.
\end{lemma}

This follows easily by passing to the affine case and considering the
associated graded objects. Notice also that the analogous statement
obtained by replacing \lq\lq isomorphism\rq\rq{} with
\lq\lq closed immersion\rq\rq{} is valid without the flatness
hypothesis.

\refstepcounter{proposition}
\paragraph{Remark \theproposition.}\label{III.4.3}
It is essential to observe that the local lifting \(X\) obtained in
Theorem~\ref{III.4.1} is \emph{not canonical}. In other words, the
local isomorphism between two solutions is not unique: in general there
are nontrivial \(Y\)-automorphisms of~\(X\) inducing the identity on the
closed subprescheme \(X_0=X\times_Y Y_0\). Consequently, in constructing
global infinitesimal liftings of smooth preschemes one should expect a
cohomological obstruction, which will be made precise below in
Section~\ref{III.6}.
\end{filecontents*}
% Original input: drafts/SGA1_III_5_3_THROUGH_END_English.texfrag
\begin{filecontents*}[overwrite,noheader]{sga1-cumulative-input-007.tex}
\begin{corollary}
\label{III.5.3}
Let \(S,X,Y,\mathcal I,g_n\) be as above. Let
\(\mathcal P(g_n)\) be the sheaf on~\(Y\) whose sections over an open
set \(U\) are the extensions \(g_{n+1}\) of~\(g_n|_{U_n}\) to an
\(S_{n+1}\)-morphism \(U_{n+1}\to X_{n+1}\), where
\(U_m=U\times_S S_m\). Then
\(\mathcal P(g_n)\) is a pseudo-torsor under the sheaf of groups
\[
  \mathcal G=
  \mathcal{H}\!om_{\mathcal O_{Y_0}}
  \left(
    g_0^*\Omega^1_{X_0/S_0},
    \operatorname{gr}^{n+1}_{\mathcal I\mathcal O_Y}(\mathcal O_Y)
  \right).
\]
\end{corollary}

In particular:

\begin{corollary}
\label{III.5.4}
If, in addition, \(X\) is smooth over~\(S\), at least at the points of
\(g_0(Y_0)\), then \(\mathcal P(g_n)\) is a torsor. It therefore
defines an obstruction class in \(H^1(Y_0,\mathcal G)\), whose
vanishing is necessary and sufficient for the existence of a global
extension \(g_{n+1}\) of~\(g_n\). If such an extension exists, the set
of all global extensions is a principal homogeneous space under
\(H^0(Y_0,\mathcal G)\). In this case the sheaf \(\mathcal G\) may
also be written
\[
  \mathcal G=
  g_0^*(\mathfrak g_{X_0/S_0})
  \otimes_{\mathcal O_{Y_0}}
  \operatorname{gr}^{n+1}_{\mathcal I\mathcal O_Y}(\mathcal O_Y).
\]
\end{corollary}

Proceeding step by step, we see that if all the
\(H^1(Y_0,\mathcal G_n)\) vanish, where
\[
  \mathcal G_n=
  g_0^*(\mathfrak g_{X_0/S_0})
  \otimes
  \operatorname{gr}^{n}_{\mathcal I\mathcal O_Y}(\mathcal O_Y),
\]
then, starting with any \(g_k\), one may successively extend it to
\(g_{k+1},\dots\). In particular, if \(\mathcal I\) is nilpotent, an
extension of~\(g_k\) to~\(Y\) exists. The required \(H^1\)-vanishing
holds in particular when \(Y_0\) is affine. Thus:

\begin{corollary}
\label{III.5.5}
\marginpar{\raggedright\footnotesize French source p.~75}%
In Theorem~\ref{III.3.1}, one obtains another equivalent necessary and
sufficient condition by restricting the \(Y'\) occurring in
\textup{(ii)}, or in \textup{(ii bis)}, to be affine and requiring a
global extension \(g\) of~\(g_0\) over all of~\(Y'\).
\end{corollary}

Notice that the proof of Theorem~\ref{III.3.1} could not have given this
result directly.

An important case is that in which \(Y\) is flat over~\(S\). Then
\[
  \operatorname{gr}^n(\mathcal O_Y)
  =
  \operatorname{gr}^n(\mathcal O_S)
  \otimes_{\mathcal O_{S_0}}\mathcal O_{Y_0}.
\]
If, in addition, the \(\operatorname{gr}^n(\mathcal O_S)\) are locally
free over~\(S_0\), then
\[
  \mathcal G_n=
  \mathcal{H}\!om_{\mathcal O_{Y_0}}
  \bigl(g_0^*\Omega^1_{X_0/S_0},\mathcal O_{Y_0}\bigr)
  \otimes_{\mathcal O_{S_0}}\operatorname{gr}^n(\mathcal O_S).
\]
Equivalently, if \(\Omega^1_{X_0/S_0}\) is also locally free, for
example if \(X\) is smooth over~\(S\),
\[
  \mathcal G_n=
  g_0^*(\mathfrak g_{X_0/S_0})
  \otimes_{\mathcal O_{S_0}}\operatorname{gr}^n(\mathcal O_S).
\]
For example, suppose that \(S\) is affine with coordinate ring~\(A\),
and that \(\mathcal I\) is defined by an ideal \(I\subseteq A\). Then,
for every \(i\),
\[
  H^i(Y_0,\mathcal G_n)
  =
  H^i(Y_0,\mathcal G_0)\otimes_A\operatorname{gr}_I^n(A).
\]
Indeed, the question is local on~\(S_0\), and reduces to tensoring with
a free module. In this case, the vanishing of
\(H^1(Y_0,\mathcal G_0)\) implies that every obstruction to the
successive extensions of~\(g_n\) vanishes. We obtain:

\begin{corollary}
\label{III.5.6}
Let \((S,X,Y,\mathcal I,g_n)\) be as above. Suppose in addition that
\(X\) is smooth over~\(S\), that \(Y\) is flat over~\(S\), that \(S\)
is affine, and that
\[
  \operatorname{gr}^n(\mathcal O_S)
  =\mathcal I^n/\mathcal I^{n+1}
\]
is locally free for every~\(n\). Then the obstruction to constructing
\(g_{n+1}\) lies in
\[
  H^1(Y_0,\mathcal G_0)
  \otimes_A\operatorname{gr}_I^{n+1}(A),
\]
where \(A\) is the ring of~\(S\), \(I\) is the ideal defining
\(\mathcal I\), and
\[
  \mathcal G_0=g_0^*(\mathfrak g_{X_0/S_0}).
\]
If \(H^1(Y_0,\mathcal G_0)=0\), then \(g_n\) extends to an
\(\widehat S\)-morphism
\[
  \widehat g\colon\widehat Y\longrightarrow\widehat X.
\]
\end{corollary}

The same result plainly remains valid if, instead of ordinary
\(S\)-preschemes \(X\) and~\(Y\), one starts with adic formal
\(\widehat S\)-preschemes \(\mathfrak X\) and~\(\mathfrak Y\).
\marginpar{\raggedright\footnotesize French source p.~76}%
For example, it can be used to show that certain formal schemes proper
over a complete local ring are in fact algebraic. Proceeding as in
Lemma~\ref{III.4.2}, one obtains:

\begin{corollary}
\label{III.5.7}
Under the hypotheses of Corollary~\ref{III.5.6}, if \(g_0\) is an
isomorphism, then so is~\(\widehat g\).
\end{corollary}

The same statement holds for closed immersions.

\begin{proposition}
\label{III.5.8}
Let \(A\) be a complete local ring with maximal ideal
\(\mathfrak m\) and residue field~\(k\). Let \(\mathfrak X\) and
\(\mathfrak Y\) be two \(\mathfrak m\)-adic formal preschemes over~\(A\),
flat over~\(A\); that is, for every \(n\), \(X_n\) and \(Y_n\) are flat
over \(A_n=A/\mathfrak m^{n+1}\). Suppose that
\[
  X_0=\mathfrak X\otimes_Ak
\]
is smooth over~\(k\) and that
\[
  H^1(X_0,\mathfrak g_{X_0/k})=0.
\]
Then every \(k\)-isomorphism \(Y_0\xrightarrow{\sim}X_0\) extends to an
\(A\)-isomorphism \(\mathfrak Y\xrightarrow{\sim}\mathfrak X\). This
extension is unique if, in addition,
\[
  H^0(X_0,\mathfrak g_{X_0/k})=0.
\]
\end{proposition}

This gives, in particular, a uniqueness result for a smooth formal
prescheme over~\(A\) reducing to a given prescheme \(X_0\), provided
\(H^1(X_0,\mathfrak g_{X_0/k})=0\). Furthermore, if
\(\mathfrak X\) and \(\mathfrak Y\) arise as the formal completions of
ordinary schemes \(X\) and~\(Y\) proper over~\(A\), then the existence
theorem for sheaves in formal geometry
(see Bourbaki Seminar, no.~182, and EGA~III, 5.4.1) gives a bijection
between the \(A\)-isomorphisms \(Y\xrightarrow{\sim}X\) and the
\(A\)-isomorphisms of their formal completions. Hence:

\begin{corollary}
\label{III.5.9}
Proposition~\ref{III.5.8} remains valid when \(\mathfrak X\) and
\(\mathfrak Y\) are replaced by ordinary \(A\)-schemes \(X\) and~\(Y\)
proper over~\(A\).
\end{corollary}

Finally, suppose that \(\mathfrak X\) is a formal scheme proper
over~\(A\), while \(\mathfrak Y=\widehat Y\) for an ordinary scheme
\(Y\) proper over~\(A\). Proposition~\ref{III.5.8} gives sufficient
conditions for an isomorphism
\(\mathfrak X\xrightarrow{\sim}\widehat Y\) to exist, and hence for
\(\mathfrak X\) to be algebraic, that is, isomorphic to
\(\widehat X\) for an ordinary scheme \(X\) proper over~\(A\); this
\(X\) is then canonically determined. This applies in particular when
\(X_0=\mathbb P_k^r\), and more generally when \(X_0\) is a
Severi--Brauer scheme, that is, becomes isomorphic to projective space
over the algebraic closure of~\(k\).
\marginpar{\raggedright\footnotesize French source p.~77}%
Every formal scheme proper and flat over~\(A\) with special fiber
\(\mathbb P_k^r\) is algebraizable; more precisely, it is isomorphic to
the \(\mathfrak m\)-adic formal completion of~\(\mathbb P_A^r\).
Consequently, by the existence theorem, every ordinary scheme proper
over a complete local ring~\(A\), with special fiber
\(\mathbb P_k^r\), is isomorphic to~\(\mathbb P_A^r\).

Using descent theory, one can prove that when \(A\) is not complete,
\(X\) becomes isomorphic to \(\mathbb P^r\) after a finite étale base
extension \(A\to A'\). In this form the result remains valid when the
special fiber is a Severi--Brauer scheme.

\end{filecontents*}
% Original input: drafts/SGA1_III_5_OPENING_THROUGH_5_2_English.texfrag
\begin{filecontents*}[overwrite,noheader]{sga1-cumulative-input-008.tex}
% SGA 1, Exposé III, French TeX lines 5275--5520
% (section III.5 opening through Corollary III.5.2 and the setup for III.5.3).

\section{Global infinitesimal lifting of morphisms}
\label{III.5}

Let \(T\) be a topological space, let \(\mathcal G\) be a sheaf of groups
on~\(T\), and let \(\mathcal P\) be a sheaf of sets on~\(T\) on which
\(\mathcal G\) acts, on the right for definiteness. We say that
\(\mathcal P\) is \emph{formally principal homogeneous} under
\(\mathcal G\) if the familiar homomorphism
\[
  \mathcal G\times\mathcal P\longrightarrow
  \mathcal P\times\mathcal P
\]
of sheaves of sets, induced by the action of~\(\mathcal G\) on
\(\mathcal P\), is an \emph{isomorphism}. Equivalently, for every
\(x\in T\), the stalk \(\mathcal P_x\) is either empty or a principal
homogeneous space under the ordinary group \(\mathcal G_x\). Equivalently
again, for every open subset \(U\subseteq T\), the set \(\mathcal P(U)\)
is either empty or a principal homogeneous space under the ordinary group
\(\mathcal G(U)\).

We say that \(\mathcal P\) is a \emph{principal homogeneous sheaf} under
\(\mathcal G\) if it is formally principal homogeneous and all its stalks
\(\mathcal P_x\) are nonempty; in other words, if every \(\mathcal P_x\)
is a nonempty principal homogeneous space under~\(\mathcal G_x\).
\marginpar{\raggedright\footnotesize French source p.~71}%
\footnote{It seems preferable to adopt the shorter and more expressive
term \lq\lq\emph{torsor under}~\(\mathcal G\)\rq\rq{}, introduced in the
thesis of J.\ Giraud.}
Recall that the set of isomorphism classes of principal homogeneous
sheaves under~\(\mathcal G\) is identified with the cohomology set
\(H^1(T,\mathcal G)\). When \(\mathcal G\) is commutative, this is the
usual cohomology group of~\(T\) with coefficients in~\(\mathcal G\).
Thus every principal homogeneous sheaf \(\mathcal P\) has a characteristic
class
\[
  c(\mathcal P)\in H^1(T,\mathcal G).
\]
This class vanishes if and only if \(\mathcal P\) is trivial---that is,
isomorphic to~\(\mathcal G\), acting on itself by right translations---or,
equivalently, if and only if \(\mathcal P\) has a section.

\begin{proposition}
\label{III.5.1}
Let \(S\) be a prescheme, let \(X\) and~\(Y\) be preschemes over~\(S\),
and let \(Y_0\) be a closed subprescheme of~\(Y\), defined by an ideal
sheaf \(\mathcal J\) on~\(Y\) with \(\mathcal J^2=0\). Let
\(g_0\colon Y_0\to X\) be an \(S\)-morphism, and let
\(\mathcal P(g_0)\) be the sheaf on~\(Y\) whose sections on an open subset
\(U\) are the \(S\)-morphisms \(g\colon U\to X\) extending
\(g_0|_{U\cap Y_0}\). Then \(\mathcal P(g_0)\) is naturally a
\emph{formally principal homogeneous sheaf} under the commutative sheaf of
groups
\[
  \mathcal G
  =\mathcal{H}\!om_{\mathcal O_{Y_0}}
    \bigl(g_0^*\Omega^1_{X/S},\mathcal J\bigr).
\]
\end{proposition}

Put \(\mathcal P=\mathcal P(g_0)\). For every open subset \(U\) of~\(Y\),
we must define a map
\[
  \mathcal P(U)\times\mathcal G(U)\longrightarrow\mathcal P(U)
\]
such that: \textup{(a)} for fixed \(g\in\mathcal P(U)\), the map
\(s\mapsto gs\) from \(\mathcal G(U)\) to \(\mathcal P(U)\) is bijective;
\textup{(b)} \(\mathcal P(U)\) becomes a right \(\mathcal G(U)\)-set; and
\textup{(c)} these maps are compatible with restriction to every open
subset \(V\subseteq U\). Condition~\textup{(c)} will be immediate, so for
simplicity we may take \(U=Y\). We leave condition~\textup{(b)}, which is
local in nature, to the reader. It remains, for a given
\(g\in\mathcal P(Y)\), to define a natural bijection from
\(\mathcal G(Y)\) onto \(\mathcal P(Y)\).

We are therefore given an \(S\)-morphism \(g\colon Y\to X\), and seek a
canonical bijection
\begin{equation*}
\tag{*}\label{eq:III.5.1.*}
  \operatorname{Hom}_{\mathcal O_{Y_0}}
  \bigl(g_0^*\Omega^1_{X/S},\mathcal J\bigr)
  \xrightarrow{\ \sim\ }\mathcal P(Y),
\end{equation*}
where \(\mathcal P(Y)\) is the set of \(S\)-morphisms \(g'\colon Y\to X\)
that induce on~\(Y_0\) the same morphism \(g_0\colon Y_0\to X\) as~\(g\).
Giving such a \(g'\) is equivalent to giving an \(S\)-morphism
\(h\colon Y\to X\times_SX\) such that
\[
  \operatorname{pr}_1\circ h=g,
  \qquad
  h\circ i=(g_0,g_0),
\]
where \(\operatorname{pr}_1\colon X\times_SX\to X\) is the first
projection, \(i\colon Y_0\to Y\) is the canonical immersion, and
\((g_0,g_0)=\Delta_{X/S}g_0\colon Y_0\to X\times_SX\) has components
\(g_0,g_0\).
\marginpar{\raggedright\footnotesize French source p.~72}%
This is summarized by the diagram
\[
\begin{array}{ccc}
Y_0 & \xrightarrow{\ h_0=(g_0,g_0)\ } & X\times_SX
\\[2pt]
{\scriptstyle i}\Big\downarrow & \nearrow_{\scriptstyle h} &
  \Big\downarrow{\scriptstyle\operatorname{pr}_1}
\\[2pt]
Y & \xrightarrow{\ g\ } & X.
\end{array}
\]

The morphism \(h_0\) factors through the diagonal immersion
\(\Delta_{X/S}\). Since \(Y\) lies in the first infinitesimal
neighborhood of~\(Y_0\), that is, since \(\mathcal J^2=0\), every such
\(h\) necessarily factors, uniquely, through the first infinitesimal
neighborhood of the diagonal. As an \(X\)-prescheme via
\(\operatorname{pr}_1\), this neighborhood is identified with the
relative spectrum \(X'\) of the sheaf of algebras
\[
  \mathcal O_X\oplus\Omega^1_{X/S},
\]
where the second summand is a square-zero ideal; the diagonal morphism
\(X\to X'\) corresponds to the canonical augmentation of this algebra
sheaf.

Put
\[
  Y'=X'\times_XY,
  \qquad
  Y'_0=Y'\times_YY_0=X'\times_XY_0.
\]
The desired morphisms \(h\) are then in bijection with the sections
\(u\) of~\(Y'\) over~\(Y\) that extend a fixed section \(u_0\) of
\(Y'_0\) over~\(Y_0\). More explicitly, \(Y'\) is the relative spectrum
over~\(Y\) of
\[
  \mathcal A
  =g^*\bigl(\mathcal O_X\oplus\Omega^1_{X/S}\bigr)
  =\mathcal O_Y\oplus g^*\Omega^1_{X/S},
\]
and \(Y'_0\) is the relative spectrum over~\(Y_0\) of
\[
  \mathcal A_0
  =\mathcal A\otimes_{\mathcal O_Y}\mathcal O_{Y_0}
  =\mathcal O_{Y_0}\oplus g_0^*\Omega^1_{X/S}.
\]
The section \(u_0\) is defined by the canonical augmentation
\(\mathcal A_0\to\mathcal O_{Y_0}\). Hence \(\mathcal P(Y)\) is
identified with the algebra homomorphisms
\(\mathcal A\to\mathcal O_Y\) inducing that augmentation on
\(\mathcal A_0\).

For brevity put \(\mathcal M=g^*\Omega^1_{X/S}\) and
\(\mathcal M_0=\mathcal M\otimes_{\mathcal O_Y}\mathcal O_{Y_0}\).
Such an algebra homomorphism is determined by an
\(\mathcal O_Y\)-linear homomorphism \(\mathcal M\to\mathcal O_Y\). The
condition on \(\mathcal A_0\) says precisely that its image lies in
\(\mathcal J\); conversely, because \(\mathcal J^2=0\), every
\(\mathcal O_Y\)-linear map \(\mathcal M\to\mathcal J\) defines such an
algebra homomorphism. We therefore obtain
\[
  \operatorname{Hom}_{\mathcal O_Y}(\mathcal M,\mathcal J)
  =\operatorname{Hom}_{\mathcal O_{Y_0}}(\mathcal M_0,\mathcal J),
\]
because the action of \(\mathcal O_Y\) on~\(\mathcal J\) factors through
\(\mathcal O_{Y_0}\). This is precisely the canonical bijection
\eqref{eq:III.5.1.*}.

Taking into account the implication
\(\textup{(i)}\Rightarrow\textup{(iii)}\) in
Theorem~\ref{III.3.1}, we obtain:

\begin{corollary}
\label{III.5.2}
\marginpar{\raggedright\footnotesize French source p.~73}%
If \(X\) is smooth over~\(S\), at least at the points of~\(g_0(Y_0)\),
then \(\mathcal P\) is in fact a \emph{principal homogeneous sheaf}
under the commutative sheaf of groups~\(\mathcal G\). In this case it may
also be written
\[
  \mathcal G
  =g_0^*(\mathfrak g_{X/S})
   \otimes_{\mathcal O_{Y_0}}\mathcal J,
\]
where \(\mathfrak g_{X/S}=(\Omega^1_{X/S})^\vee\)
\label{indnot:cb}
is the sheaf on~\(X\)
dual to \(\Omega^1_{X/S}\), that is, the \emph{tangent sheaf}, or
\emph{sheaf of derivations}, of~\(X\) relative to~\(S\). The displayed
formula follows from the fact that \(\Omega^1_{X/S}\) is then locally free
of finite rank.
\end{corollary}

In particular, this principal homogeneous sheaf determines a cohomology
class in \(H^1(Y_0,\mathcal G)\). Its vanishing is necessary and sufficient
for the existence of an \(S\)-morphism \(g\) extending~\(g_0\). If such an
extension exists, the set of all possible extensions is a homogeneous
space under the group \(H^0(Y_0,\mathcal G)\).

In applications of the methods of formal geometry, the situation is most
often the following. Let \(X\) and~\(Y\) be \(S\)-preschemes, let
\(\mathcal I\) be a coherent ideal sheaf on~\(S\), let \(S_n\) denote the
closed subprescheme of~\(S\) defined by \(\mathcal I^{n+1}\), and put
\[
  X_n=X\times_SS_n,
  \qquad
  Y_n=Y\times_SS_n.
\]
Suppose that an \(S_n\)-morphism
\[
  g_n\colon Y_n\longrightarrow X_n
\]
is given. Equivalently, this is an \(S\)-morphism \(Y_n\to X\), or an
\(S_{n+1}\)-morphism \(Y_n\to X_{n+1}\), since any such morphism
necessarily induces \(Y_n\to X_n\). We seek to extend it to an
\(S_{n+1}\)-morphism
\[
  g_{n+1}\colon Y_{n+1}\longrightarrow X_{n+1}.
\]
If this process can be continued indefinitely, it yields a morphism
\(\widehat Y\to\widehat X\) between the formal preschemes obtained by
completing \(Y\) and~\(X\) with respect to the ideals
\(\mathcal I\mathcal O_Y\) and~\(\mathcal I\mathcal O_X\), respectively.

Apply Proposition~\ref{III.5.1} with
\[
  (S,X,Y,Y_0,g_0)
  =(S_{n+1},X_{n+1},Y_{n+1},Y_n,g_n).
\]
The sheaf \(\mathcal G\) is now the sheaf of homomorphisms from
\[
  g_n^*\Omega^1_{X_{n+1}/S_{n+1}}
\]
to
\[
  \mathcal J
  =\mathcal I^{n+1}\mathcal O_Y/
   \mathcal I^{n+2}\mathcal O_Y.
\]
Since \(\mathcal J\) is annihilated by
\(\mathcal I\mathcal O_Y\), the first sheaf may be replaced by its
restriction to \(Y_0=Y\times_SS_0\).
\marginpar{\raggedright\footnotesize French source p.~74}%
Let \(h_0\colon Y_0\to X_{n+1}\) be the composite
\(Y_0\to Y_n\xrightarrow{g_n}X_n\to X_{n+1}\), and let
\(g_0\colon Y_0\to X_0\) be the morphism induced by~\(g_n\). Equivalently,
\(h_0\) is the composite \(Y_0\xrightarrow{g_0}X_0\to X_{n+1}\). The
restricted sheaf is
\(h_0^*\Omega^1_{X_{n+1}/S_{n+1}}\). Since the inverse image of
\(\Omega^1_{X_{n+1}/S_{n+1}}\) on
\[
  X_0=X_{n+1}\times_{S_{n+1}}S_0
\]
is \(\Omega^1_{X_0/S_0}\), we may write
\[
  \mathcal G
  =\mathcal{H}\!om_{\mathcal O_{Y_0}}
    \left(
      g_0^*\Omega^1_{X_0/S_0},
      \mathcal I^{n+1}\mathcal O_Y/
      \mathcal I^{n+2}\mathcal O_Y
    \right).
\]
Thus we obtain the following corollary.
\end{filecontents*}
% Original input: drafts/SGA1_III_6_OPENING_THROUGH_OBSTRUCTION_DERIVATION_English.texfrag
\begin{filecontents*}[overwrite,noheader]{sga1-cumulative-input-009.tex}
\section{Global infinitesimal lifting of smooth S-schemes}
\label{III.6}

Under the hypotheses of Theorem~\ref{III.4.1}, we ask whether there is
a prescheme \(X\) smooth over~\(Y\) such that
\(X\times_Y Y_0\) is \(Y_0\)-isomorphic to~\(X_0\), knowing that such
a scheme exists locally on~\(X_0\). Returning to the step-by-step
construction, replace \(Y\) by~\(S\). Let \(S_0\) be a closed
subprescheme of~\(S\) defined by a sheaf of ideals \(\mathcal I\), no
longer necessarily locally nilpotent, and let \(S_n\) be the closed
subprescheme of~\(S\) defined by \(\mathcal I^{n+1}\). Suppose that a
prescheme \(X_n\), smooth over~\(S_n\), has been given. We seek an
\(S_{n+1}\)-prescheme \(X_{n+1}\) reducing to~\(X_n\), that is,
equipped with an isomorphism
\[
  X_n\xrightarrow{\sim}
  X_{n+1}\times_{S_{n+1}}S_n,
\]
and smooth over~\(S_{n+1}\), or equivalently flat over~\(S_{n+1}\) by
\ifthenelse{\boolean{orig}}{II~}{}\ref{II.2.1}.

As observed in Section~\ref{III.4}, such data amount to a sheaf of
algebras \(\mathcal B\) over
\(f^{-1}(\mathcal O_{S_{n+1}})\), where \(f\) is the continuous map
underlying the structural morphism \(X_n\to S_n\), together with an
augmentation \(\mathcal B\to\mathcal O_{X_n}\) compatible with
\[
  f^{-1}(\mathcal O_{S_{n+1}})
  \longrightarrow f^{-1}(\mathcal O_{S_n}),
\]
and satisfying the two conditions \textup{(a)} and \textup{(b)} already
given there. We need only recall that they are local on the underlying
topological space of~\(X_n\). Theorem~\ref{III.4.1} gives a solution
locally, unique locally up to a nonunique isomorphism. We first make
this last point precise.

\begin{proposition}
\label{III.6.1}
\marginpar{\raggedright\footnotesize French source p.~78}%
Let \(X_{n+1}\) over~\(S_{n+1}\) reduce to \(X_n\) over~\(S_n\). The
sheaf, on the underlying topological space of \(X_n\), equivalently of
\(X_0\), of \(S_{n+1}\)-automorphisms of~\(X_{n+1}\) which induce the
identity on~\(X_n\) is canonically isomorphic, as a sheaf of groups, to
\[
  \mathcal G=
  \mathfrak g_{X_0/S_0}
  \otimes_{\mathcal O_{S_0}}
  \operatorname{gr}_{\mathcal I}^{n+1}(\mathcal O_S).
\]
\end{proposition}

Indeed, by Corollary~\ref{III.5.4} and Lemma~\ref{III.4.2}, this
automorphism sheaf is a torsor under~\(\mathcal G\). It has the
distinguished section given by the identity automorphism of
\(X_{n+1}\), and hence identifies with \(\mathcal G\) as a sheaf of
sets. Compatibility with the group structures is immediate. It is also
a special case of a more general compatibility, omitted here, between
the torsor structures of Proposition~\ref{III.5.1} and
Corollary~\ref{III.5.3} and the composition of morphisms.

In particular, the sheaf on~\(X_0\) of germs of automorphisms of
\(X_{n+1}\), with the structures just described, is commutative. Thus,
if \(X'_{n+1}\) is another solution and is isomorphic to~\(X_{n+1}\)
over an open set \(U\subseteq X_0\), the isomorphism
\[
  \mathcal{A}ut(X_{n+1})|_U
  \xrightarrow{\sim}
  \mathcal{A}ut(X'_{n+1})|_U
\]
obtained by transport of structure through an isomorphism
\(X_{n+1}|_U\xrightarrow{\sim}X'_{n+1}|_U\) does not depend on the
chosen isomorphism. Under the identifications of
Proposition~\ref{III.6.1}, it is simply the identity of~\(\mathcal G\).

We deduce:

\begin{corollary}
\label{III.6.2}
Let \(X_{n+1}\) and \(X'_{n+1}\) be smooth over~\(S_{n+1}\) and reduce
to~\(X_n\). Then the sheaf, on the underlying space of~\(X_0\), of
\(S_{n+1}\)-isomorphisms
\[
  X_{n+1}\xrightarrow{\sim}X'_{n+1}
\]
which induce the identity on~\(X_n\) is naturally a torsor
under~\(\mathcal G\).
\end{corollary}

This says precisely that \(X_{n+1}\) and \(X'_{n+1}\) are locally
isomorphic and that the sheaf of germs of automorphisms of either one
is~\(\mathcal G\).

By Theorem~\ref{III.4.1}, one can choose an open cover \((U_i)\) of
\(X_n\), which may be assumed affine, and for each~\(i\) a smooth
scheme \(X^i\) over~\(S_{n+1}\) reducing to~\(U_i\).
\marginpar{\raggedright\footnotesize French source p.~79}%
For simplicity suppose that \(X_n\) is separated. The intersections
\(U_{ij}=U_i\cap U_j\) are then affine. Since
\(H^1(U_{ij},\mathcal G)=0\) for the quasi-coherent sheaf
\(\mathcal G\), Corollary~\ref{III.6.2} shows that
\(X^i|_{U_{ij}}\) and \(X^j|_{U_{ij}}\) are isomorphic. Choose an
isomorphism
\[
  f_{ji}\colon X^i|_{U_{ij}}\xrightarrow{\sim}X^j|_{U_{ij}}.
\]
It is determined up to a section of \(\mathcal G\) over~\(U_{ij}\).
For every triple of indices put
\[
  f_{ji}^{(k)}=f_{ji}|_{U_{ijk}},
  \qquad
  U_{ijk}=U_i\cap U_j\cap U_k.
\]
If
\begin{equation}
\label{eq:III.6.1}
  f_{kj}^{(i)}f_{ji}^{(k)}=f_{ki}^{(j)},
\tag{1}
\end{equation}
then the \(X^i\) glue through the \(f_{ji}\), and hence define the
desired solution \(X=X_{n+1}\). More generally, a solution exists if
the \(f_{ji}\) can be modified to
\begin{equation}
\label{eq:III.6.2}
  f'_{ji}=f_{ji}g_{ji},
  \qquad
  g_{ji}\in\Gamma(U_{ij},\mathcal G),
\tag{2}
\end{equation}
so that the \(f'_{ji}\) satisfy the preceding cocycle condition. This
sufficient condition is also necessary. Indeed, any solution \(X\) is
locally isomorphic to \(X^i\); choosing isomorphisms
\[
  f_i\colon X|_{U_i}\xrightarrow{\sim}X^i
\]
gives
\[
  f'_{ji}
  =(f_j|_{U_{ij}})(f_i|_{U_{ij}})^{-1}
  \colon X^i|_{U_{ij}}\xrightarrow{\sim}X^j|_{U_{ij}},
\]
which satisfy the gluing condition.

Now set
\begin{equation}
\label{eq:III.6.3}
  f_{ijk}
  =(f_{ki}^{(j)})^{-1}f_{kj}^{(i)}f_{ji}^{(k)}.
\tag{3}
\end{equation}
This is an automorphism of \(X^i|_{U_{ijk}}\), which we identify with a
section of \(\mathcal G\) by Proposition~\ref{III.6.1}. A direct formal
calculation, left to the reader, shows that
\(f=(f_{ijk})\) is a Čech \(2\)-cocycle of the open cover
\(\mathcal U=(U_i)\) with coefficients in~\(\mathcal G\).
\marginpar{\raggedright\footnotesize French source p.~80}%
The same calculation shows that, after the modification
\eqref{eq:III.6.2}, the gluing condition \eqref{eq:III.6.1} for the
\(f'_{ij}\) is equivalent to
\begin{equation}
\label{eq:III.6.4}
  f=dg,
\tag{4}
\end{equation}
where \(g=(g_{ij})\) is regarded as a Čech \(1\)-cochain of
\(\mathcal U\) with coefficients in~\(\mathcal G\). Thus the necessary
and sufficient condition for a solution to exist is the vanishing of
the cohomology class in \(H^2(\mathcal U,\mathcal G)\) defined by the
cocycle \eqref{eq:III.6.3}.

Since \(\mathcal U\) is an affine open cover of the scheme~\(X_0\),
\(H^2(\mathcal U,\mathcal G)\) identifies with
\(H^2(X_0,\mathcal G)\). The resulting class in
\(H^2(X_0,\mathcal G)\) is plainly independent of the chosen affine
cover. We call it the obstruction class to extending~\(X_n\) to a
scheme \(X_{n+1}\) smooth over~\(S_{n+1}\).

Suppose this obstruction vanishes. The argument just given shows that
every solution \(X=X_{n+1}\) is isomorphic to one obtained by gluing
with isomorphisms \(f'_{ji}\), which may be assumed of the form
\eqref{eq:III.6.2}. The gluing condition is then
\eqref{eq:III.6.4}. The set of admissible \(g\) is a principal
homogeneous space under the group
\(Z^1(\mathcal U,\mathcal G)\) of Čech \(1\)-cocycles. Moreover, two
cochains \(g\) and \(g'\), with \(dg=dg'=f\), define isomorphic
solutions if and only if the cocycle \(g-g'\) is of the form \(dh\),
where \(h=(h_i)\in C^0(\mathcal U,\mathcal G)\). We therefore obtain:

\end{filecontents*}
% Original input: drafts/SGA1_III_6_THEOREM_6_3_THROUGH_END_English.texfrag
\begin{filecontents*}[overwrite,noheader]{sga1-cumulative-input-010.tex}
\begin{theorem}
\label{III.6.3}
Let \((S,\mathcal I,X_n)\) be as above, and suppose that \(X_n\) is
separated.\footnote{This hypothesis is in fact unnecessary, and the
cocycle calculations above can be avoided. See J.~Giraud,
\emph{Cohomologie non abélienne}, forthcoming from Springer-Verlag in
1971, and compare Remarks~\ref{III.6.4}.}
There is a canonically defined obstruction class in
\(H^2(X_0,\mathcal G)\), where \(\mathcal G\) is defined in
Proposition~\ref{III.6.1}. Its vanishing is necessary and sufficient
for the existence of a scheme \(X_{n+1}\), smooth over~\(S_{n+1}\),
which reduces to~\(X_n\). If the obstruction vanishes, the set of
isomorphism classes, with the isomorphisms inducing the identity
on~\(X_n\), of \(S_{n+1}\)-preschemes \(X_{n+1}\) reducing to~\(X_n\)
is naturally a principal homogeneous space under
\(H^1(X_0,\mathcal G)\).
\end{theorem}

\begin{remarques}
\label{III.6.4}
From Proposition~\ref{III.6.1} onward, the arguments are purely formal
and are better expressed in the language of local categories, or even
general fibered categories.
\marginpar{\raggedright\footnotesize French source p.~81}%
In that general setting, one obtains an obstruction class to the
existence of a global object of a category in which objects exist
locally, any two objects are locally isomorphic, and the automorphism
group of every object is commutative. As a special case it contains the
second connecting homomorphism in an exact sequence of sheaves of
groups which need not be commutative, studied for example by
Grothendieck in Kansas or Tôhoku. The elementary cocycle calculation
given here should therefore be regarded as a makeshift necessitated by
the absence of a satisfactory reference.
\end{remarques}

\refstepcounter{proposition}
\paragraph{\theproposition.}\label{III.6.5}
In Theorem~\ref{III.6.3}, the principal homogeneous space under
\(H^1(X_0,\mathcal G)\) generally has no distinguished element. After
localizing on~\(S\), this is reflected by a torsor on~\(S_0\) with
structural group \(R^1f_*(\mathcal G)\), which need not be trivial and
thus defines a cohomology class in
\[
  H^1\bigl(S_0,R^1f_*(\mathcal G)\bigr)
\]
which need not vanish. This is the situation when the class
\(d\in H^2(X_0,\mathcal G)\), although not zero, vanishes locally
over~\(S\); equivalently, it defines the zero section of
\(R^2f_*(\mathcal G)\), or the zero element of
\(H^0(S_0,R^2f_*(\mathcal G))\).

\refstepcounter{proposition}
\paragraph{\theproposition.}\label{III.6.6}
At present almost nothing is known about the general algebraic
mechanism of the cohomology classes introduced in this section and the
preceding one, and nothing precise can be said even in the simplest
special cases, such as abelian schemes over Artinian rings.
\footnote{It is now known that the obstruction always vanishes in this
case. [Added in 2003 (MR): see F.~Oort,
\lq\lq Finite group schemes, local moduli for abelian varieties and
lifting problems,\rq\rq{} \emph{Algebraic Geometry, Oslo 1970},
Wolters-Noordhoff, 1972, pp.~223--254.]}
It is to be hoped that someone will work out this particularly
interesting question in detail. It is closely connected, in particular,
with the moduli theory of algebraic structures.

\begin{corollary}
\label{III.6.7}
If \(H^2(X_0,\mathcal G)=0\), then an \(X_{n+1}\) exists; it is unique
up to isomorphism if, in addition, \(H^1(X_0,\mathcal G)=0\).
\end{corollary}

Proceeding step by step, and recalling that an affine scheme is
acyclic for a quasi-coherent sheaf, we obtain:

\begin{corollary}
\label{III.6.8}
\marginpar{\raggedright\footnotesize French source p.~82}%
Under the hypotheses of Theorem~\ref{III.4.1}, if \(X_0\) is affine,
then there exists a prescheme \(X\), smooth over~\(Y\), reducing
to~\(X_0\); it is unique up to a nonunique isomorphism.
\end{corollary}

Notice that the direct proof of Theorem~\ref{III.4.1} could not have
given this result.

\begin{corollary}
\label{III.6.9}
Under the hypotheses of Theorem~\ref{III.6.3}, suppose that \(S\) is
affine with ring~\(A\), that \(\mathcal I\) is defined by an ideal
\(I\subseteq A\), and that
\[
  \operatorname{gr}_{\mathcal I}^n(\mathcal O_S)
  =\mathcal I^n/\mathcal I^{n+1}
\]
is locally free for every~\(n\). Then
\[
  H^i(X_0,\mathcal G)
  =
  H^i(X_0,\mathcal G_0)
  \otimes_A\operatorname{gr}_I^{n+1}(A),
  \qquad
  \mathcal G_0=\mathfrak g_{X_0/S_0}.
\]
Thus the obstruction class to extending~\(X_n\) lies in
\[
  H^2(X_0,\mathcal G_0)
  \otimes_A\operatorname{gr}_I^{n+1}(A).
\]
If it vanishes, the set of isomorphism classes of solutions is a
principal homogeneous space under
\[
  H^1(X_0,\mathcal G_0)
  \otimes_A\operatorname{gr}_I^{n+1}(A).
\]
\end{corollary}

In particular:

\begin{corollary}
\label{III.6.10}
Under the hypotheses of Corollary~\ref{III.6.9}, suppose that
\[
  H^2(X_0,\mathfrak g_{X_0/S_0})=0.
\]
Then there exists an \(I\)-adic formal scheme \(\mathfrak X\) over the
\(\mathcal I\)-adic formal completion \(\widehat S\) of~\(S\), smooth
over~\(\widehat S\) in the sense that each \(\mathfrak X_p\) is smooth
over~\(S_p\), and reducing to~\(X_n\); that is, equipped with an
isomorphism
\[
  X_n\xrightarrow{\sim}\mathfrak X\times_S S_n.
\]
If, in addition,
\[
  H^1(X_0,\mathfrak g_{X_0/S_0})=0,
\]
then such a formal scheme \(\mathfrak X\) is unique up to isomorphism.
\end{corollary}

Indeed, one constructs \(X_{n+1},X_{n+2},\dots\) successively and
passes to the limit to obtain~\(\mathfrak X\). The uniqueness assertion
was already proved in the preceding paragraph.

\end{filecontents*}
% Original input: drafts/SGA1_III_7_English.texfrag
\begin{filecontents*}[overwrite,noheader]{sga1-cumulative-input-011.tex}
\section[Application to formal schemes]{Application to the construction
of formal and ordinary schemes smooth over a complete local ring \(A\)}
\label{III.7}

The results of the preceding section sometimes prove the existence of
an \(\mathfrak m\)-adic formal scheme over such a ring, reducing to a
given smooth scheme \(X_0\) over~\(k\).
\marginpar{\raggedright\footnotesize French source p.~83}%
We distinguish two cases.

\paragraph{\textup{(a)} Equal characteristics.}
This includes, in particular, the case where \(k\) has characteristic
zero. There is then a coefficient field in~\(A\), that is, a subfield
\(k'\subseteq A\) such that \(A\to k\) induces an isomorphism
\(k'\xrightarrow{\sim}k\). In fact, there is then an ordinary scheme
smooth over~\(A\) reducing to~\(X_0\), namely
\[
  X=X_0\otimes_kA,
\]
where \(A\) is regarded as a \(k\)-algebra through the section
\(k\to k'\to A\) defined by~\(k'\).

This construction is not natural. Already for the dual-number algebra
\(A=k[t]/(t^2)\), a different lifting homomorphism \(k\to A\), defined
in this case by an absolute derivation of~\(k\) into itself, gives an
\(A\)-scheme \(X'\) which in general is not isomorphic to~\(X\) when
\[
  H^1(X_0,\mathfrak g_{X_0/k})\ne0.
\]
It would be interesting, when \(k\) has characteristic zero, or is
imperfect of characteristic \(p>0\), to determine which smooth
\(A\)-schemes arise in this way and when two homomorphisms \(k\to A\)
give isomorphic \(A\)-schemes.

Nevertheless, the existence of~\(k'\) implies that the first
obstruction to lifting~\(X_0\), which lies in
\[
  H^2(X_0,\mathfrak g_{X_0/k})
  \otimes_k\mathfrak m/\mathfrak m^2,
\]
must vanish. After \(X_0\) has been lifted to an \(X_1\) smooth over
\(A/\mathfrak m^2\), however, the new obstruction to constructing
\(X_2\) need not vanish. It depends on a variable element in a certain
principal homogeneous space under
\[
  H^1(X_0,\mathcal G_0)\otimes\mathfrak m/\mathfrak m^2
\]
and lies in
\[
  H^2(X_0,\mathcal G_0)
  \otimes\mathfrak m^2/\mathfrak m^3.
\]
The situation deserves a detailed study.
\footnote{It is presumably described by the Kodaira--Spencer bracket;
see the Cartan Seminar 1960/61, Exposé~4.}

\paragraph{\textup{(b)} Unequal characteristics.}
In this case essentially nothing is known, except when by good fortune
\(H^2(X_0,\mathfrak g_{X_0/k})=0\), in which case an
\(\mathfrak m\)-adic formal scheme smooth over~\(A\) and reducing
to~\(X_0\) can be constructed. Even for
\(A=\mathbb Z/p^2\mathbb Z\) and \(X_0\) an abelian scheme of
dimension~\(2\), it was not known whether \(X_0\) could be lifted to an
\(X=X_1\) smooth over~\(A\).
\footnote{This has now been proved; see the note to
Section~\ref{III.6.6}.}
Nor was there an example proved not to arise from an ordinary scheme
smooth over~\(A\); one expected such examples among projective
surfaces.
\footnote{\label{nIII.3}Such an example was subsequently constructed
by J.-P.~Serre, \emph{Proc. Nat. Acad. Sci. USA} \textbf{47} (1961),
no.~1, pp.~108--109, at least in certain dimensions. D.~Mumford gave an
unpublished example with an algebraic surface.}

By Cohen's theorem there is a Cohen \(p\)-ring \(B\) with residue
field~\(k\) and a homomorphism \(B\to A\) inducing the identity on
residue fields. The strongest lifting result is therefore obtained by
taking \(A\) itself to be a Cohen \(p\)-ring: if an ordinary or formal
solution exists over such a ring, one exists over every complete local
ring with residue field~\(k\).
\marginpar{\raggedright\footnotesize French source p.~84}%
For a Cohen \(p\)-ring, \(\mathfrak m/\mathfrak m^2\) identifies
canonically with~\(k\). Hence, for every smooth scheme \(X_0\) over a
field \(k\) of characteristic \(p>0\), there is a cohomology class in
\[
  H^2(X_0,\mathfrak g_{X_0/k})
\]
which is the first obstruction to lifting \(X_0\) to a smooth scheme
over a Cohen \(p\)-ring. It was not known whether this class could be
nonzero.\footnote{It can be nonzero, as noted in
Footnote~\ref{nIII.3}.}

Even if the \(X_n\) reducing to~\(X_0\) can be constructed
successively, this generally yields only a formal scheme
\(\mathfrak X\), smooth over~\(A\), reducing to~\(X_0\). When \(X_0\)
is proper over~\(k\), one must still decide whether \(\mathfrak X\) is
algebraizable, so as to obtain an ordinary scheme proper and smooth
over~\(A\) reducing to~\(X_0\). The only known criterion, stated in the
Bourbaki Seminar and in EGA~III, 4.7.1, is the following. Suppose that
\(\mathfrak X\) is proper over~\(A\), and that \(\mathcal L\) is an
invertible sheaf on~\(\mathfrak X\) whose restriction
\(\mathcal L_0\) to~\(X_0\) is ample, meaning that some positive tensor
power \(\mathcal L_0^{\otimes n}\) comes from a projective immersion
of~\(X_0\). Then there exist a scheme \(X\) projective over~\(A\) and
an ample invertible sheaf on~\(X\) whose \(\mathfrak m\)-adic
completion is \((\mathfrak X,\mathcal L)\).

This reduces us to extending a locally free sheaf \(\mathcal E_0\)
on~\(X_0\), chosen invertible and ample for the present purpose, to a
locally free sheaf \(\mathcal E\) on~\(\mathfrak X\). One constructs
locally free sheaves \(\mathcal E_n\) on the \(X_n\) successively. The
discussion is entirely analogous to Section~\ref{III.6}; compare
Remarks~\ref{III.6.4}. The essential role is played by the sheaf of
automorphisms of an \(\mathcal E_{n+1}\) which induce the identity on
\(\mathcal E_n\). This sheaf identifies with
\begin{equation}
\label{eq:III.7.*}
  \mathcal G
  =
  \mathcal{H}\!om_{\mathcal O_{X_0}}
    \left(\mathcal E_0,
      \mathcal E_0\otimes
      \operatorname{gr}_{\mathcal I}^{n+1}(\mathcal O_X)\right)
  =
  \mathcal{E}nd_{\mathcal O_{X_0}}(\mathcal E_0)
  \otimes\operatorname{gr}_{\mathcal I}^{n+1}(\mathcal O_X),
\tag{*}
\end{equation}
which is again a sheaf of commutative groups. We obtain:

\begin{proposition}
\label{III.7.1}
Let \(S\) be a prescheme equipped with a quasi-coherent sheaf of ideals
\(\mathcal I\), let \(X\) be a prescheme over~\(S\), let \(S_n\) be
the closed subprescheme of~\(S\) defined by \(\mathcal I^{n+1}\), and
put \(X_n=X\times_SS_n\) for every~\(n\). Let \(\mathcal E_n\) be a
locally free sheaf on~\(X_n\). The obstruction to extending it to a
locally free sheaf \(\mathcal E_{n+1}\) on~\(X_{n+1}\) is a canonical
class in \(H^2(X_0,\mathcal G)\), where \(\mathcal G\) is the
quasi-coherent sheaf in \eqref{eq:III.7.*}. Its vanishing is necessary
and sufficient for such an extension to exist. If it vanishes, the set
of isomorphism classes, with isomorphisms inducing the identity on
\(\mathcal E_n\), of extensions \(\mathcal E_{n+1}\) is a principal
homogeneous space under \(H^1(X_0,\mathcal G)\).
\marginpar{\raggedright\footnotesize French source p.~85}%
\end{proposition}

This proposition has the usual corollaries. We record only that if
\(X\) is flat over~\(S\), then
\[
  \mathcal G=
  \mathcal{E}nd_{\mathcal O_{X_0}}(\mathcal E_0)
  \otimes_{\mathcal O_{S_0}}
  \operatorname{gr}_{\mathcal I}^{n+1}(\mathcal O_S).
\]
Consequently, if \(S\) is affine with ring~\(A\) and the
\(\mathcal I^n/\mathcal I^{n+1}\) are locally free, the condition
\[
  H^2(X_0,\mathcal G_0)=0,
  \qquad
  \mathcal G_0=
  \mathcal{E}nd_{\mathcal O_{X_0}}(\mathcal E_0),
\]
is sufficient for \(\mathcal E_{n+1}\) to exist, and hence, step by
step, for all successive extensions \(\mathcal E_m\),
\(m=n,n+1,\dots\), to exist.

Returning to the original situation, we obtain:

\begin{proposition}
\label{III.7.2}
Let \(A\) be a complete local ring, and let \(\mathfrak X\) be a
formal scheme proper and flat over~\(A\). Suppose that \(X_0\) is
projective and
\[
  H^2(X_0,\mathcal O_{X_0})=0.
\]
Then there exists a scheme \(X\) projective over~\(A\) whose
\(\mathfrak m\)-adic formal completion is isomorphic to~\(\mathfrak X\).
\end{proposition}

Combining this with Corollary~\ref{III.6.10}, we obtain:

\begin{theorem}
\label{III.7.3}
Let \(A\) be a complete local ring with residue field~\(k\), and let
\(X_0\) be a smooth projective scheme over~\(k\) such that
\[
  H^2(X_0,\mathfrak g_{X_0/k})
  =
  H^2(X_0,\mathcal O_{X_0})
  =0.
\]
Then there exists a smooth projective scheme \(X\) over~\(A\) reducing
to~\(X_0\).
\end{theorem}

More generally, if \(X_n\) is smooth over
\(A_n=A/\mathfrak m^{n+1}\) and reduces to~\(X_0\), then there exist a
scheme \(X\) smooth and proper over~\(A\) and an isomorphism
\[
  X\otimes_AA_n\xrightarrow{\sim}X_n.
\]

\begin{corollary}
\label{III.7.4}
Every smooth proper curve over~\(k\) arises by reduction from a smooth
proper curve over~\(A\).
\end{corollary}

This result, together with the existence theorem for sheaves in formal
geometry, will be the essential tool for studying the fundamental
group of~\(X_0\) by transcendental methods.
\marginpar{\raggedright\footnotesize French source p.~86}%

\end{filecontents*}
% Original input: drafts/SGA1_II_1_English.texfrag
\begin{filecontents*}[overwrite,noheader]{sga1-cumulative-input-012.tex}
% SGA 1, Exposé II, French TeX lines 2519--2667 (opening and section II.1).

\chapter{Smooth morphisms: generalities and differential properties}
\label{II}
\marginpar{\raggedright\footnotesize French source p.~29}%

References to Exposé~\ref{I} are indicated by~I. Recall that all rings are
noetherian and all preschemes locally noetherian.

\section{Generalities}
\label{II.1}

Let \(Y\) be a prescheme, and let \(t_1,\dots,t_n\) be indeterminates. We
write
\begin{equation}
\label{eq:II.1.1}
  Y[t_1,\dots,t_n]
  =Y\otimes_{\mathbb Z}\mathbb Z[t_1,\dots,t_n].
\end{equation}
Thus \(Y[t_1,\dots,t_n]\) is a \(Y\)-scheme, affine over~\(Y\), defined by
the quasi-coherent sheaf of algebras
\(\mathcal O_Y[t_1,\dots,t_n]\). Giving a section of this prescheme over~\(Y\)
is therefore equivalent to giving \(n\) sections of~\(\mathcal O_Y\), which
correspond to the images of the \(t_i\) under the associated homomorphism. If
\(Y'\) lies over~\(Y\), then
\begin{equation}
\label{eq:II.1.2}
  Y[t_1,\dots,t_n]\times_Y Y'=Y'[t_1,\dots,t_n].
\end{equation}
Consequently, giving a \(Y\)-morphism from~\(Y'\) to
\(Y[t_1,\dots,t_n]\) is equivalent to giving \(n\) sections of
\(\mathcal O_{Y'}\). We also have
\begin{equation}
\label{eq:II.1.3}
  \bigl(Y[t_1,\dots,t_n]\bigr)[t_{n+1},\dots,t_m]
  =Y[t_1,\dots,t_m],
\end{equation}
by the analogous formula for polynomial rings over~\(\mathbb Z\).
Formula~\eqref{eq:II.1.2} shows that \(Y[t_1,\dots,t_n]\) varies
functorially with~\(Y\).

The \(Y\)-scheme \(Y[t_1,\dots,t_n]\) is of finite type and flat over~\(Y\).

\begin{sectiondefinition}
\label{II.1.1}
Let \(f\colon X\to Y\) be a morphism, making \(X\) a \(Y\)-prescheme. We
say that \(f\) is
\emph{smooth}\footnote{Former terminology: \(f\) is \emph{simple} at~\(x\),
or \(x\) is a \emph{simple} point for~\(f\). This terminology caused
confusion in various contexts (simple algebras, simple groups) and had to be
abandoned.}
at \(x\in X\), or that \(X\) is \emph{smooth over~\(Y\) at~\(x\)}, if there
exist an integer \(n\geq0\), an open neighborhood \(U\) of~\(x\), and an
étale \(Y\)-morphism from~\(U\) to \(Y[t_1,\dots,t_n]\). We say that \(f\)
is \emph{smooth} (respectively, that \(X\) is \emph{smooth over~\(Y\)}) if
it is smooth at every point of~\(X\).

An algebra~\(B\) over a ring~\(A\) is said to be smooth at a prime ideal
\(\mathfrak p\) of~\(B\) if \(\operatorname{Spec}(B)\) is smooth over
\(\operatorname{Spec}(A)\) at~\(\mathfrak p\); \(B\) is said to be smooth
over~\(A\) if it is smooth over~\(A\) at every prime ideal
\(\mathfrak p\) of~\(B\). Finally, a local homomorphism \(A\to B\) of local
rings is said to be smooth (or \(B\) is said to be smooth over~\(A\))
\footnote{It is preferable in this case, as in EGA~IV, 18.6.1, to say that
\(B\) is \emph{essentially smooth} over~\(A\).}
if \(B\) is a localization of an algebra \(B_1\) of finite type and smooth
over~\(A\).
\end{sectiondefinition}

The notion that \(X\) is smooth over~\(Y\) is local on~\(X\) and on~\(Y\).
If \(X\) is smooth over~\(Y\), then it is locally of finite type over~\(Y\).

\begin{proposition}
\label{prop:II.1.1}
The set of points \(x\) of~\(X\) at which \(f\) is smooth is open.
\end{proposition}

This is immediate from the definition.

\begin{corollary}
\label{II.1.2}
If \(B\) is smooth over~\(A\) at~\(\mathfrak p\), then it is smooth over
\(A\) at every prime ideal \(\mathfrak q\) of~\(B\) contained in
\(\mathfrak p\).
\end{corollary}

Proposition~\ref{prop:II.1.1} also implies that the last two definitions
in Definition~\ref{II.1.1} agree wherever both apply.

\begin{proposition}
\label{II.1.3}
\textup{(i)} An étale morphism, and in particular an open immersion or an
identity morphism, is smooth.

\textup{(ii)} A base extension of a smooth morphism is smooth.

\textup{(iii)} The composite of two smooth morphisms is smooth.
\end{proposition}

Assertion~\textup{(i)} is immediate from the definition. More precisely, we
have the following.

\begin{corollary}
\label{II.1.4}
\[
  \text{étale}=\text{quasi-finite}+\text{smooth}.
\]
\end{corollary}

Assertion~\textup{(ii)} follows at once from the analogous fact for étale
morphisms (\ifthenelse{\boolean{orig}}{I~}{}\ref{I.4.6}) and for the projections
\(Y[t_1,\dots,t_n]\to Y\); see~\eqref{eq:II.1.2}. Assertion~\textup{(iii)}
follows formally because it holds separately for étale morphisms
(\ifthenelse{\boolean{orig}}{I~}{}\ref{I.4.6}) and for projections of the form
\(Y[t_1,\dots,t_n]\to Y\), by~\eqref{eq:II.1.3}, together with the two
facts just cited for~\textup{(ii)}. Indeed, suppose that \(Y\) is smooth
over~\(Z\) and that \(X\) is smooth over~\(Y\); we prove that \(X\) is smooth
over~\(Z\). We may suppose that \(Y\) is étale over
\(Z[t_1,\dots,t_n]\), and that \(X\) is étale over
\(Y[s_1,\dots,s_m]\). The first hypothesis implies that
\(Y[s_1,\dots,s_m]\) is étale over
\[
  \bigl(Z[t_1,\dots,t_n]\bigr)[s_1,\dots,s_m]
  =Z[t_1,\dots,t_n,s_1,\dots,s_m].
\]
It follows that \(X\) is étale over
\(Z[t_1,\dots,t_n,s_1,\dots,s_m]\), as required.

\refstepcounter{proposition}
\paragraph{Remark \theproposition.}\label{II.1.5}
The integer \(n\) occurring in Definition~\ref{II.1.1} is indeed determined:
it is immediately seen to be the dimension of the local ring of~\(x\) in
its fiber \(f^{-1}\bigl(f(x)\bigr)\).
\marginpar{\raggedright\footnotesize French source p.~31}%
It is called the \emph{relative dimension} of~\(X\) over~\(Y\). Relative
dimension is additive under composition of morphisms.
\end{filecontents*}
% Original input: drafts/SGA1_II_2_English.texfrag
\begin{filecontents*}[overwrite,noheader]{sga1-cumulative-input-013.tex}
% SGA 1, Exposé II, French TeX lines 2669--2877 (section II.2).

\section{Some criteria for smoothness of a morphism}
\label{II.2}

\begin{theorem}
\label{II.2.1}
Let \(f\colon X\to Y\) be a morphism locally of finite type, let \(x\in X\),
and put \(y=f(x)\). Then \(f\) is smooth at~\(x\) if and only if
\textup{(a)} \(f\) is flat at~\(x\), and \textup{(b)} the fiber
\(f^{-1}(y)\) is smooth over~\(\kappa(y)\) at~\(x\).
\end{theorem}

Since the composite of two flat morphisms is flat, and
\(Y[t_1,\dots,t_n]\to Y\) is flat, every smooth morphism is flat.
Together with Proposition~\ref{II.1.3}\textup{(ii)}, this proves the
necessity of the two conditions. Conversely, suppose that \textup{(a)} and
\textup{(b)} hold. Let \(V\) be an affine neighborhood of~\(y\), with
ring~\(A\), and let \(U\) be an affine neighborhood of~\(x\) above~\(V\),
with ring~\(B\). After shrinking \(U\), condition~\textup{(b)} allows us to
suppose that there is an étale \(\kappa(y)\)-morphism
\[
  g\colon U|_{f^{-1}(y)}
  \longrightarrow \operatorname{Spec} k[t_1,\dots,t_n],
  \qquad k=\kappa(y),
\]
defined by \(n\) sections \(g_i\) of the structure sheaf of
\(U|_{f^{-1}(y)}\). We may plainly suppose that the \(g_i\), which a priori
are elements of
\[
  B\otimes_A k=BS^{-1},
  \qquad S=A-\mathfrak p,
\]
where \(\mathfrak p\) is the prime ideal of~\(A\) corresponding to~\(y\),
come from sections of the structure sheaf of~\(U\). Thus \(g\) is induced
by a morphism, still denoted by~\(g\),
\[
  g\colon U\longrightarrow Y[t_1,\dots,t_n],
\]
after multiplying all the \(g_i\), if necessary, by the same nonzero
element of~\(k\). Now \(U\) is flat over~\(Y\) by~\textup{(a)}, as is
\(Y[t_1,\dots,t_n]\); moreover, \(g\) induces an étale morphism between the
fibers over~\(y\). Therefore \(g\) is étale at~\(x\), by
\ifthenelse{\boolean{orig}}{I~}{}\ref{I.5.8}, as required.

\begin{corollary}
\label{II.2.2}
% SGA1-CORR-0017: smoothness of X/S is not implied under only flatness of Y/S.
Let \(S\) be a prescheme, and let \(f\colon X\to Y\) be an \(S\)-morphism
of finite type, where \(Y\) is of finite type and flat over~\(S\). Let
\(x\in X\), and let \(s\) be its image in~\(S\). Then \(f\) is smooth at
\(x\) if and only if \(X\) is flat over~\(S\) at~\(x\), and
the induced morphism on the fibers over~\(s\),
\[
  f_s\colon X_s\longrightarrow Y_s,
\]
is smooth at~\(x\).
\end{corollary}

Only sufficiency requires proof; it follows from the criterion in
Theorem~\ref{II.2.1}, together with the flatness criterion
\ifthenelse{\boolean{orig}}{I~}{}\ref{I.5.9}.

\marginpar{\raggedright\footnotesize French source p.~32}%
To state the next result, \lq\lq recall\rq\rq{} that a morphism
\(f\colon X\to Y\) locally of finite type is said to be
\emph{equidimensional} at \(x\in X\) if, on writing \(y=f(x)\), there is an
open neighborhood \(U\) of~\(x\), every component of which dominates a
component of~\(Y\), such that, for every \(y'\in Y\), all the irreducible
components of \(f^{-1}(y')\cap U\) have one and the same dimension,
independent of~\(y'\). In this condition it is enough to take for \(y'\)
the generic points of the irreducible components of~\(Y\) passing
through~\(y\), \emph{and} the point~\(y\) itself. If, for example, \(X\)
and~\(Y\) are integral and \(f\) is dominant, the condition says that the
components of \(f^{-1}(y)\) passing through~\(x\) have the
\lq\lq correct\rq\rq{} dimension, that is, the dimension of the generic
fiber. Recall that their dimension is always at least that of the generic
fiber.

Suppose that \(f\) is equidimensional at~\(x\), that the dimension of its
fiber at~\(x\) is~\(n\), and that
\[
  g\colon U\longrightarrow Y'=Y[t_1,\dots,t_n]
\]
is a \(Y\)-morphism from a neighborhood \(U\) of~\(x\), inducing on the
fibers over~\(y\) a morphism that is quasi-finite at~\(x\). Equivalently,
suppose that \(g\) itself is quasi-finite at~\(x\). One then shows that
every irreducible component of~\(U\) passing through~\(x\) dominates an
irreducible component of~\(Y'\). Moreover, by the
\lq\lq normalization lemma\rq\rq{}, such a \(g\) always exists. Conversely,
if there is a quasi-finite \(Y\)-morphism \(g\) from an open neighborhood
\(U\) of~\(x\) to a \(Y\)-scheme of the form
\(Y'=Y[t_1,\dots,t_n]\), such that every component of~\(U\) passing
through~\(x\) dominates a component of~\(Y'\), then \(f\) is
equidimensional at~\(x\). With this understood, we have:

\begin{proposition}
\label{II.2.3}
Let \(f\colon X\to Y\) be a morphism locally of finite type, let \(x\in X\),
and put \(y=f(x)\). Suppose that \(\mathcal O_y\) is normal. Then \(f\) is
smooth at~\(x\) if and only if \(f\) is equidimensional at~\(x\) and the
fiber \(f^{-1}(y)\) is smooth over~\(\kappa(y)\) at~\(x\).
\end{proposition}

It follows at once from the definition that a smooth morphism is
equidimensional. Notice, however, that a flat morphism of finite type need
not be equidimensional at~\(x\), even if its fiber at~\(x\) is irreducible.
Let us prove the converse. Since \(f^{-1}(y)\) is smooth over~\(\kappa(y)\)
at~\(x\), we may suppose, after replacing \(X\) by a suitable neighborhood
of~\(x\) if necessary, that there is a \(Y\)-morphism
\[
  g\colon X\longrightarrow Y[t_1,\dots,t_n]=Y'
\]
which induces an étale morphism on the fibers over~\(y\), and is therefore
quasi-finite at~\(x\). Thus \(g\) is unramified, and, since \(f\) is
equidimensional at~\(x\), the irreducible components of~\(X\) passing
through~\(x\) each dominate an irreducible component of~\(Y'\).
\marginpar{\raggedright\footnotesize French source p.~33}%
Consequently, the homomorphism
\[
  \mathcal O_{y'}\longrightarrow\mathcal O_x,
  \qquad y'=g(x),
\]
induced by~\(g\) is injective. This homomorphism is also unramified, and
\(\mathcal O_{y'}\) is normal because it is a localization of
\(\mathcal O_y[t_1,\dots,t_n]\), which is normal since \(\mathcal O_y\) is
normal. Hence \(\mathcal O_{y'}\to\mathcal O_x\) is étale, by
\ifthenelse{\boolean{orig}}{I~}{}\ref{I.9.5}\textup{(ii)}.

\refstepcounter{proposition}
\paragraph{Remarks \theproposition.}\label{II.2.4}
The preceding statement remains valid if the hypothesis that
\(\mathcal O_y\) is normal is replaced by the weaker hypothesis that \(Y\)
is \emph{geometrically unibranch} at~\(y\); see \ifthenelse{\boolean{orig}}{I~}{}\ref{I.11}. Indeed,
\ifthenelse{\boolean{orig}}{I~}{}\ref{I.9.5} remains valid under this hypothesis.

Let us also take this opportunity to point out that, if the residue field of
a local integral domain~\(A\) is algebraically closed, then analytic
irreducibility---that is, the condition that \(\widehat A\) be an integral
domain---implies geometric unibranchedness. The converse also holds in every
category of \lq\lq good rings\rq\rq{}, more precisely in a category of rings
stable under the usual operations and in which the completion of a normal
local ring is normal. By Zariski's \lq\lq theorem of analytic
normality\rq\rq{}, this condition is satisfied in the category of affine
algebras and their localizations.\footnote{Cf.\ EGA~IV, 7.8.}

Finally, \lq\lq recall\rq\rq{} in the present context the following result
of Hironaka,\footnote{Cf.\ EGA~IV, 5.12.10.} which can sometimes be used to
ensure that \(f^{-1}(y)\) is a reduced scheme, that is, that it coincides
with what many algebraic geometers would loosely regard as the fiber of~\(f\)
above~\(y\) without multiplicities, namely \(f^{-1}(y)_{\mathrm{red}}\).

\begin{proposition}
\label{II.2.5}
Let \(f\colon X\to Y\) be a dominant morphism of finite type between reduced
preschemes, and let \(y\in Y\) be a point such that \(\mathcal O_y\) is
regular. Suppose that every component of \(f^{-1}(y)\) has multiplicity~\(1\)
(in the sense defined below), and that \(f^{-1}(y)_{\mathrm{red}}\) is
normal. Then \(f^{-1}(y)\) is reduced, hence normal; \(X\) is normal at every
point of \(f^{-1}(y)\); and \(X\) is flat over~\(Y\) at every point of
\(f^{-1}(y)\).
\end{proposition}

\marginpar{\raggedright\footnotesize French source p.~34}%
A component \(Z\) of \(f^{-1}(y)\) is said to have
\emph{multiplicity~\(1\)} if, with \(x\) denoting the generic point of~\(Z\),
the following conditions hold:

\textup{(i)} \(\dim\mathcal O_x=\dim\mathcal O_y\), that is, \(Z\) is not an
\lq\lq excess component\rq\rq{} or, equivalently, is not of
\lq\lq excessively large dimension\rq\rq{};

\textup{(ii)} the maximal ideal of~\(\mathcal O_x\) is generated by the
maximal ideal of~\(\mathcal O_y\), which a priori, by the choice of~\(x\),
generates an ideal of definition of~\(\mathcal O_x\).

Taking Proposition~\ref{II.2.3}, or Theorem~\ref{II.2.1}, into account, we
obtain:

\begin{corollary}
\label{II.2.6}
Let \(f\colon X\to Y\) be a dominant morphism of finite type between reduced
preschemes, and let \(y\in Y\) be a point such that \(\mathcal O_y\) is
regular. Then \(f\) is smooth at all points of~\(X\) above~\(y\) if and only
if the components of \(f^{-1}(y)\) have multiplicity~\(1\) and
\(f^{-1}(y)_{\mathrm{red}}\) is smooth over~\(\kappa(y)\).
\end{corollary}

In the past, this situation was considered especially when \(Y\) was the
spectrum of a discrete valuation ring~\(A\), and was commonly described in
such terms as: \lq\lq if the reduction of~\(X\) with respect to the given
valuation is nice\rq\rq{}\dots. Moreover, \(X\) then denoted a closed
subscheme (if one may say so) of some \(\mathbb P_K^n\), where \(K\) is the
fraction field of~\(A\); and, for want of adequate language, the more
intrinsic role of an object \lq\lq defined over~\(A\)\rq\rq{} rather than
merely over~\(K\) was scarcely apparent.
\end{filecontents*}
% Original input: drafts/SGA1_II_3_English.texfrag
\begin{filecontents*}[overwrite,noheader]{sga1-cumulative-input-014.tex}
% SGA 1, Exposé II, French TeX lines 2878--2927 (section II.3).

\section{Permanence properties}
\label{II.3}

\begin{proposition}
\label{II.3.1}
Let \(f\colon X\to Y\) be a morphism, let \(x\in X\), and put \(y=f(x)\).
Suppose that \(f\) is smooth at~\(x\). Then \(\mathcal O_x\) is reduced
(respectively, regular; respectively, normal) if and only if
\(\mathcal O_y\) has the corresponding property.
\end{proposition}

This statement is known when \(X\) is of the form
\(Y[t_1,\dots,t_n]\), and it was proved in I, no.~\ref{I.9}, for an étale
morphism. The general case follows at once from
Definition~\ref{II.1.1}.

We shall not detail here the other permanence properties, which already
follow either from flatness alone or from the fact that \(X\) is locally
quasi-finite and flat over a \(Y\)-prescheme of the form
\(Y[t_1,\dots,t_n]\), that is, as we shall say, from the fact that \(X\) is
Cohen--Macaulay over~\(Y\).
\marginpar{\raggedright\footnotesize French source p.~35}%
We point out only that this latter fact implies
\begin{equation}
\label{eq:II.3.1}
\begin{aligned}
  \dim\mathcal O_x
    &=\dim\mathcal O_y+n-d,\\
  \operatorname{depth}\mathcal O_x
    &=\operatorname{depth}\mathcal O_y+n-d,
\end{aligned}
\end{equation}
where \(n\) is the dimension of the fiber of~\(f\) at~\(x\), and \(d\) is
the transcendence degree of~\(\kappa(x)\) over~\(\kappa(y)\). Thus, on
putting
\[
  \operatorname{codepth}=\dim-\operatorname{depth},
\]
we obtain\footnote{For these formulas, cf.\ EGA~IV, 6.1 and~6.3.}
\begin{equation}
\label{eq:II.3.2}
  \operatorname{codepth}\mathcal O_x
  =\operatorname{codepth}\mathcal O_y.
\end{equation}
It follows, for example, that \(\mathcal O_x\) is Cohen--Macaulay
(respectively, has no embedded components) if and only if the same is true
of~\(\mathcal O_y\).
\end{filecontents*}
% Original input: drafts/SGA1_II_4_15_THROUGH_4_19_English.texfrag
\begin{filecontents*}[overwrite,noheader]{sga1-cumulative-input-015.tex}
% SGA 1, Exposé II, French TeX lines 3728--4046
% (Theorem II.4.15 through Remark II.4.19; end of section II.4).

\begin{theorem}
\label{II.4.15}
Let \(X\) be an \(S\)-prescheme locally of finite type, let \(Y\) be a
closed subprescheme of~\(X\) defined by a coherent sheaf of
ideals~\(\mathcal J\), and let \(x\in X\). Suppose now that \(Y\) is
\emph{smooth over~\(S\) at~\(x\)}, with no assumption on~\(X\). Then the
following conditions are equivalent:
\begin{enumerate}
\item[\textup{(i)}]
\(X\) is smooth over~\(S\) at~\(x\);
\item[\textup{(ii)}]
the immersion \(i\colon Y\to X\) is regular at~\(x\), that is,
\(\mathcal J_x\) is a regular ideal of~\(\mathcal O_x\).
\end{enumerate}
\end{theorem}

\begin{corollary}
\label{II.4.16}
Suppose that \(Y\) is smooth over~\(S\). Then \(X\) is smooth over~\(S\) in
a neighborhood of~\(Y\), that is, at every point of~\(Y\), if and only if
\(Y\) is regularly embedded in~\(X\), or equivalently the immersion
\(i\colon Y\to X\) is regular.
\end{corollary}

\paragraph{Proof of Theorem~\ref{II.4.15}.}
Suppose first that \textup{(i)} holds. Apply
criterion~\textup{(ii)} of Theorem~\ref{II.4.10}. Since the resulting
morphism \(g\colon X_1\to X'\) is flat, regular immersions are preserved
under its inverse image. It is therefore enough to prove that
\[
  Y'=S[t_{p+1},\dots,t_n]
\]
is regularly embedded in \(X'=S[t_1,\dots,t_n]\), which is immediate:
\(t_1,\dots,t_p\) form a regular system of generators of the ideal defining
\(Y'\) in~\(X'\).

Conversely, suppose that \textup{(ii)} holds. Let
\(g_i\), \(1\leq i\leq p\), be a regular system of generators of
\(\mathcal J_x\), and choose elements \(g_i\),
\(p+1\leq i\leq n\), of~\(\mathcal O_{X,x}\) whose images \(g'_i\) in
\(\mathcal O_{Y,x}\) define an étale morphism
\[
  Y_1\longrightarrow Y'=S[t_{p+1},\dots,t_n]
\]
from a neighborhood \(Y_1\) of~\(x\) in~\(Y\). The \(g_i\),
\(1\leq i\leq n\), come from sections of~\(\mathcal O_X\) on a
neighborhood \(X_1\) of~\(x\). We may suppose that \(X_1=X\) and
\(Y_1=Y\). They define a morphism
\[
  g\colon X\longrightarrow X'=S[t_1,\dots,t_n],
\]
and it remains only to prove that \(g\) is étale at~\(x\).

After shrinking \(X\), we may suppose that
\(g_1,\dots,g_p\) form a regular system of generators of~\(\mathcal J\)
throughout~\(X\).
\marginpar{\raggedright\footnotesize French source p.~48}%
In particular, they generate~\(\mathcal J\), so the closed subprescheme
\(Y\) of~\(X\) is the inverse image under~\(g\) of the closed
subprescheme~\(Y'\) of~\(X'\). Put \(x'=g(x)\). The fiber of
\(X\to X'\) over~\(x'\) is then the same as the fiber of
\(Y\to Y'\) over~\(x'\), and is therefore étale over~\(\kappa(x')\).
Thus \(g\) is unramified at~\(x\); it remains to prove that \(g\) is flat
there.

Let \(\mathcal J'_{x'}\) be the ideal of~\(Y'\) in
\(\mathcal O_{X',x'}\). The associated graded ring of
\(\mathcal O_{X',x'}\), filtered by the powers of
\(\mathcal J'_{x'}\), is free over~\(\mathcal O_{Y',x'}\) in every
degree. The associated graded ring of~\(\mathcal O_{X,x}\), filtered by
the powers of
\(\mathcal J_x=\mathcal J'_{x'}\mathcal O_{X,x}\), is canonically
isomorphic to the tensor product of the former graded ring with
\(\mathcal O_{Y,x}\); indeed, the two graded rings are polynomial rings in
\(p\) indeterminates over \(\mathcal O_{Y',x'}\) and
\(\mathcal O_{Y,x}\), respectively. Finally,
\[
  \mathcal O_{X,x}/\mathcal J_x=\mathcal O_{Y,x}
\]
is flat over
\[
  \mathcal O_{X',x'}/\mathcal J'_{x'}=\mathcal O_{Y',x'}.
\]
A general flatness criterion now applies: for a local homomorphism of
noetherian local rings \(A'\to A\), with an ideal \(J'\neq A'\) such that
the associated graded ring is free over~\(A'/J'\) in every degree, these
conditions imply flatness. Hence \(X\) is flat over~\(X'\) at~\(x\), as
required.

\begin{corollary}
\label{II.4.17}
Let \(X\) be a prescheme locally of finite type over~\(Y\), let \(i\) be a
section of~\(X\) over~\(Y\), let \(y\in Y\), and put \(x=i(y)\). Let
\(\mathcal J\) be the sheaf of ideals on~\(X\) defined by the
subprescheme \(i(Y)\), which we assume closed in order to simplify the
statement, a condition automatically satisfied when \(X\) is a scheme.
The following conditions are equivalent:
\begin{enumerate}
\item[\textup{(i)}]
\(X\) is smooth over~\(Y\) at~\(x\);
\item[\textup{(ii)}]
\(i\) is a regular immersion at~\(y\);
\item[\textup{(iii)}]
the completion of the \(\mathcal O_y\)-algebra \(\mathcal O_x\) for the
topology defined by the powers of~\(\mathcal J_x\) is isomorphic to a
formal power-series algebra
\[
  \mathcal O_y[[t_1,\dots,t_n]];
\]
\item[\textup{(iii bis)}]
there is an open neighborhood \(U\) of~\(y\) such that the sheaf of
\(\mathcal O_Y\)-algebras
\[
  \varprojlim_m i^*\bigl(\mathcal O_X/\mathcal J^{m+1}\bigr)
\]
over~\(U\) is isomorphic to a sheaf of the form
\(\mathcal O_Y[[t_1,\dots,t_n]]\);
\item[\textup{(iv)}]
there exist an open neighborhood \(U\) of~\(y\), an open neighborhood
\(V\) of~\(x\), and a \(Y\)-morphism
\[
  g\colon V\longrightarrow U[t_1,\dots,t_n]
\]
such that \(g\) is étale, \(i\) induces a section of~\(V\) over~\(U\), and
\(g\) carries this section to the zero section of
\(U[t_1,\dots,t_n]\) over~\(U\).
\end{enumerate}
\end{corollary}

The equivalence of \textup{(i)} and \textup{(ii)} is the special case
\(Y=S\) of Theorem~\ref{II.4.15}.
\marginpar{\raggedright\footnotesize French source p.~49}%
The equivalence of \textup{(ii)} and \textup{(iii)}, and morally that of
\textup{(ii)} and \textup{(iii bis)}, was noted in the preceding
\lq\lq reminders\rq\rq{}. The equivalence of \textup{(i)} and
\textup{(iv)} follows readily from the equivalence of
\textup{(i)} and \textup{(ii)} in Theorem~\ref{II.4.10}.

\begin{corollary}
\label{II.4.18}
Let \(X\) be a prescheme smooth over~\(S\). Then the diagonal morphism
\[
  \Delta_{X/S}\colon X\longrightarrow X\times_SX
\]
is a \emph{regular immersion}; equivalently, \(X\) is
\emph{differentially smooth} over~\(S\).
\end{corollary}

Indeed, this is a special case of Corollary~\ref{II.4.16}, since both
\(X\) and \(X\times_SX\) are smooth over~\(S\).

\paragraph{Remarks II.4.18.}
Recall from \ifthenelse{\boolean{orig}}{I~}{}\ref{I.1} that, if \(X\) is a prescheme over~\(S\), one
introduces on~\(X\) the quasi-coherent sheaves of algebras
\[
  \mathcal P^n_{X/S}
  =\mathcal O_{X\times_SX}/\mathcal I_X^{n+1},
\]
where \(\mathcal I_X\) is the sheaf of ideals defining the diagonal in
\(X\times_SX\). It is regarded as a sheaf of \(\mathcal O_X\)-algebras by
means of the first projection
\(\operatorname{pr}_1\colon X\times_SX\to X\). The
\(\mathcal P^n_{X/S}\) form a projective system of algebras on~\(X\). Its
projective limit, denoted by \(\mathcal P^\infty_{X/S}\), is the structure
sheaf of the formal completion of \(X\times_SX\) along the diagonal, where
we now suppose that \(X\) is locally of finite type over~\(S\), so that the
\(\mathcal P^n_{X/S}\) are coherent.

To say that \(X\) is differentially smooth over~\(S\), that is, that the
diagonal \(\Delta_{X/S}\) is a regular immersion, is also to say that
\(\mathcal P^\infty_{X/S}\) is regular as a sheaf of algebras augmented
toward~\(\mathcal O_X\). Equivalently, \(\Omega^1_{X/S}\) is locally free
and the canonical surjective homomorphism
\[
  \operatorname{Sym}_{\mathcal O_X}\bigl(\Omega^1_{X/S}\bigr)
  \longrightarrow
  \operatorname{gr}_*\bigl(\mathcal P^\infty_{X/S}\bigr)
\]
is an isomorphism. Equivalently again, every point of~\(X\) has an open
neighborhood on which the augmented sheaf of algebras
\(\mathcal P^\infty_{X/S}\) is isomorphic to a sheaf
\(\mathcal O_X[[t_1,\dots,t_n]]\).

Let \(s\) be a section of~\(X\) over~\(S\), and let \(\mathcal J\) be the
sheaf of ideals on~\(X\) that it defines; for simplicity, suppose that
\(s(S)\) is closed.
\marginpar{\raggedright\footnotesize French source p.~50}%
There are then canonical isomorphisms of augmented
\(\mathcal O_X\)-algebras
\begin{equation}
\label{eq:II.4.4}
  s^*\bigl(\mathcal P^n_{X/S}\bigr)
  =\mathcal O_X/\mathcal J^{n+1},
  \qquad
  s^*\bigl(\mathcal P^\infty_{X/S}\bigr)
  =\varprojlim_n\mathcal O_X/\mathcal J^{n+1}.
\end{equation}
These isomorphisms are functorial under base change in the evident sense
and, taking this fact into account, give another characterization of the
sheaves of algebras \(\mathcal P^n_{X/S}\) on~\(S\). If, for example,
\(S=\operatorname{Spec}(k)\), where \(k\) is a field, then giving a section
\(s\) of~\(X\) over~\(S\) is equivalent to giving a \(k\)-rational point
\(x\in X\), and the preceding formulas say that there is an isomorphism of
\(k\)-algebras
\begin{equation}
\label{eq:II.4.5}
  \mathcal P^n_{X/S}(x)
  =\mathcal O_x/\mathfrak m_x^{n+1}.
\end{equation}
This justifies the name \emph{sheaf of principal parts of order~\(n\) on
\(X\) relative to~\(S\)} for \(\mathcal P^n_{X/S}\).
\label{indnot:ad1}

Formula~\eqref{eq:II.4.4} also shows that \emph{if \(X\) is differentially
smooth over~\(S\) at every point of~\(s(S)\), then \(X\) is smooth
over~\(S\) at every point of~\(s(S)\)}, by
Corollary~\ref{II.4.17}; the converse is also true, by
Corollary~\ref{II.4.18}. Together with Corollary~\ref{II.4.13}, this readily
implies that, if \(X\) is an \(S\)-prescheme locally of finite type, then
\emph{\(X\) is smooth over~\(S\) if and only if it is flat and
differentially smooth over~\(S\)}. Notice that the flatness hypothesis is
essential, as one sees by taking \(X\) to be a closed subprescheme of~\(S\).

Recall also that \(\mathcal P^n_{X/S}\) has a \emph{second algebra
structure}, obtained from the second projection
\(\operatorname{pr}_2\colon X\times_SX\to X\). This structure is obtained
from the first by means of the \emph{canonical involution} of the sheaf of
rings \(\mathcal P^n_{X/S}\) induced by the symmetry automorphism of
\(X\times_SX\). We denote by \(d^n_{X/S}\), or simply by~\(d^n\), the
\label{indnot:bb}
homomorphism of sheaves of rings
\begin{equation}
\label{eq:II.4.6}
  d^n_{X/S}\colon\mathcal O_X\longrightarrow\mathcal P^n_{X/S}
\end{equation}
corresponding to this second algebra structure. Under the isomorphism
\eqref{eq:II.4.4}, this homomorphism carries a section \(f\) of
\(\mathcal O_X\) to a section \(d^n(f)\) of \(\mathcal P^n_{X/S}\) whose
inverse image under a section \(s\) of~\(X\) over~\(S\) is identified with
the canonical image of~\(f\) in
\(\Gamma(X,\mathcal O_X/\mathcal J^{n+1})\). This justifies the name
\emph{system of principal parts of order~\(n\) of~\(f\)} for~\(d^nf\),
especially when \(S=\operatorname{Spec}(k)\), as in
\eqref{eq:II.4.5}.
\marginpar{\raggedright\footnotesize French source p.~51}%

Finally, the homomorphism \eqref{eq:II.4.6} may be regarded as the
\emph{universal differential operator of order at most~\(n\)}
on~\(\mathcal O_X\), relative to the constant prescheme~\(S\).
\footnote{For everything concerning this paragraph, see
EGA~IV, 16.8--16.12.}
Here a differential operator of order at most~\(n\) from
\(\mathcal O_X\) to a module~\(F\) means a homomorphism of sheaves \(D\)
that factors as
\[
  D\colon
  \mathcal O_X\xrightarrow{\ d^n\ }\mathcal P^n_{X/S}
  \xrightarrow{\ u\ }F,
\]
where \(u\) is a homomorphism of \(\mathcal O_X\)-modules, uniquely
determined by~\(D\). This agrees with the intuitive recursive definition:
\(D\) has order at most~\(n\) if, for every section \(g\) of
\(\mathcal O_X\) on an open set \(U\subseteq X\), the map
\[
  f\longmapsto D(fg)-gD(f)
\]
is a differential operator of order at most \(n-1\) on~\(U\).

It follows that, \emph{if \(X\) is differentially smooth over~\(S\), the
sheaf of rings of differential operators of all orders has the familiar
simple structure} from differential calculus on differentiable manifolds.
In particular, it has locally a basis as an \(\mathcal O_X\)-module formed
by the \emph{divided powers} in mutually commuting operators
\(\delta/\delta x_i\), \(1\leq i\leq n\). If \(\mathcal O_S\) is a sheaf of
\(\mathbb Q\)-algebras, where \(\mathbb Q\) is the field of rational
numbers, it is enough to take ordinary polynomials in the
\(\delta/\delta x_i\). In this case, exceptionally, local freeness of
\(\Omega^1_{X/S}\) already suffices for \(X\) to be differentially smooth
over~\(S\).

\paragraph{Remark II.4.19.}
The terminology \lq\lq regular immersion\rq\rq{},
\lq\lq regular ideal\rq\rq{}, and so forth, introduced in this section met
with fairly strong and general opposition from Chevalley and Serre. The
terms \lq\lq Cohen--Macaulay ideal\rq\rq{},
\lq\lq Macaulay ideal\rq\rq{}, and
\lq\lq Macaulayan ideal\rq\rq{} were proposed; morally, this would also
require the adoption of \lq\lq Cohen--Macaulay immersion\rq\rq{} or
\lq\lq Macaulay immersion\rq\rq{}. That terminology, however, conflicts
with another already used in forthcoming versions of the
\emph{multiplodoque}, where a morphism of finite type is said to be
\lq\lq Cohen--Macaulay\rq\rq{} at a point if it is flat there and the
local ring of the fiber through that point is Cohen--Macaulay. Until a
satisfactory solution is found, we shall retain, with every reservation,
the terminology introduced in this section.\footnote{This is the
terminology adopted in EGA~\(\mathrm{0}_{\mathrm{IV}}\), 15.1.7.}
\end{filecontents*}
% Original input: drafts/SGA1_II_4_9_THROUGH_4_13_English.texfrag
\begin{filecontents*}[overwrite,noheader]{sga1-cumulative-input-016.tex}
% SGA 1, Exposé II, French TeX lines 3235--3499
% (Proposition II.4.9 through Corollary II.4.13 and its proof).

\begin{proposition}
\label{II.4.9}
Let \(X\) be an \(S\)-prescheme, let \(Y\) be a closed subprescheme of~\(X\)
defined by a coherent sheaf of ideals~\(\mathcal J\), and let \(x\in X\).
Let \(g_i\), \(1\leq i\leq n\), be sections of~\(\mathcal O_X\) defining an
\(S\)-morphism
\[
  g\colon X\longrightarrow S[t_1,\dots,t_n]=X',
\]
and let \(p\) be an integer with \(0\leq p\leq n\). Suppose that \(X\) is
\emph{smooth over~\(S\) at~\(x\)}. Then the following conditions are
equivalent:
\begin{enumerate}
\item[\textup{(i)}]
There is an open neighborhood \(X_1\) of~\(x\) in~\(X\) such that
\(g|_{X_1}\) is \emph{étale} and
\(Y_1=Y\cap X_1\), the trace of~\(Y\) on~\(X_1\), is the \emph{inverse
image} of the closed subprescheme
\[
  Y'=S[t_{p+1},\dots,t_n]
  \quad\text{of}\quad
  X'=S[t_1,\dots,t_n].
\]
Equivalently, the \(g_i\), \(1\leq i\leq p\), generate
\(\mathcal J|_{X_1}\). Thus the diagram
\[
\begin{array}{ccc}
Y_1 & \longrightarrow & X_1\\
\big\downarrow && \big\downarrow{\scriptstyle\text{ étale}}\\
Y'=S[t_{p+1},\dots,t_n] & \longrightarrow &
X'=S[t_1,\dots,t_n]
\end{array}
\]
is cartesian.
\item[\textup{(ii)}]
\(Y\) is \emph{smooth over~\(S\) at~\(x\)}; the germs of the \(g_i\),
\(1\leq i\leq p\), lie in~\(\mathcal J_x\); the elements
\(\mathrm dg_i(x)\), \(1\leq i\leq n\), form a \emph{basis} of the
\(\kappa(x)\)-vector space \(\Omega^1_{X/S}(x)\); and the
\(\mathrm dg'_i(x)\), \(p+1\leq i\leq n\), form a \emph{basis} of
\(\Omega^1_{Y/S}(x)\), where \(g'_i\) denotes the restriction of~\(g_i\)
to~\(Y\) and all differentials are relative to~\(S\).
\marginpar{\raggedright\footnotesize French source p.~40}%
\item[\textup{(iii)}]
The \(g_i\), \(1\leq i\leq p\), form a \emph{system of generators} of
\(\mathcal J_x\), and the \(\mathrm dg_i(x)\), \(1\leq i\leq n\), form a
\emph{basis} of \(\Omega^1_{X/S}(x)\) over~\(\kappa(x)\).
\item[\textup{(iv)}]
\(Y\) is \emph{smooth over~\(S\) at~\(x\)}; the \(g_i\),
\(1\leq i\leq p\), form a \emph{minimal system of generators} of
\(\mathcal J_x\); and the \(\mathrm dg'_i(x)\),
\(p+1\leq i\leq n\), form a \emph{basis} of
\(\Omega^1_{Y/S}(x)\) over~\(\kappa(x)\).
\end{enumerate}

Moreover, under these conditions, \(\mathcal J/\mathcal J^2\) is free
over~\(Y\) at~\(x\), with the classes of the \(g_i\),
\(1\leq i\leq p\), as a basis at~\(x\), and the canonical homomorphism
\[
  \mathcal J/\mathcal J^2
  \longrightarrow
  \Omega^1_{X/S}\otimes_{\mathcal O_X}\mathcal O_Y
\]
is \emph{universally injective at~\(x\)}.
\end{proposition}

\paragraph{Remark \textup{*}.}
It follows that \(p\) is determined by the other conditions: it is the
\emph{rank} of the free \(Y\)-module \(\mathcal J/\mathcal J^2\) at~\(x\),
the \emph{minimum number of generators} of~\(\mathcal J_x\) on~\(X\), and
also the integer characterized by the fact that the relative dimension of
\(Y\) over~\(S\) at~\(x\) is \(n-p\).

\paragraph{Proof of Proposition~\ref{II.4.9}.}
Suppose first that \textup{(i)} holds. By
\ifthenelse{\boolean{orig}}{I~}{}\ref{I.4.6}\textup{(iii)}, \(Y_1\) is étale over~\(Y'\); by definition it
is therefore smooth over~\(S\) at~\(x\), of relative dimension \(n-p\), and
the same is true of~\(Y\). Corollary~\ref{II.4.8}, applied to \(g\) and to
its restriction, shows that the \(\mathrm dg_i\),
\(1\leq i\leq n\), form a basis of \(\Omega^1_{X/S}\) at~\(x\), and that
the \(\mathrm dg'_i\), \(p+1\leq i\leq n\), form a basis of
\(\Omega^1_{Y/S}\) at~\(x\). The exact sequence \eqref{eq:II.4.3} then shows
that the classes of the \(g_i\), \(1\leq i\leq p\), are linearly independent
in the \(Y\)-module \(\mathcal J/\mathcal J^2\) at~\(x\). Since these
\(g_i\) generate~\(\mathcal J_x\), their classes modulo
\(\mathcal J_x^2\) form a basis of \(\mathcal J/\mathcal J^2\) at~\(x\).
Consequently, the \(g_i\), \(1\leq i\leq p\), form a minimal system of
generators of~\(\mathcal J_x\), and the homomorphism
\[
  \mathcal J/\mathcal J^2
  \longrightarrow
  \Omega^1_{X/S}\otimes_{\mathcal O_X}\mathcal O_Y
\]
in \eqref{eq:II.4.3} is universally injective at~\(x\): it carries a basis
of a module free at~\(x\) to part of a basis of another such module, both
being \(Y\)-modules. This proves that \textup{(i)} implies
\textup{(ii)}, \textup{(iii)}, and \textup{(iv)}, as well as the final
assertions of Proposition~\ref{II.4.9}.

Condition~\textup{(iii)} implies \textup{(i)} by
Corollary~\ref{II.4.8}.

Condition~\textup{(ii)} also implies \textup{(i)}.
\marginpar{\raggedright\footnotesize French source p.~41}%
Indeed, the requirement that the germs of \(g_1,\dots,g_p\) lie in
\(\mathcal J_x\) means that, after replacing \(X\) by an open neighborhood
of~\(x\), the map \(g\) induces a morphism \(h\colon Y\to Y'\). By
Corollary~\ref{II.4.8}, the two basis requirements mean that \(g\) and~\(h\)
are étale at~\(x\). Let \(Y''\) be the inverse
image of~\(Y'\) under~\(g\). Then \(Y\) is a closed subprescheme of~\(Y''\),
and \(Y''\) is étale over~\(Y'\) at~\(x\), by
\ifthenelse{\boolean{orig}}{I~}{}\ref{I.4.6}\textup{(iii)}, because \(g\) is étale at~\(x\). The immersion
\(Y\to Y''\) is therefore itself étale, by \ifthenelse{\boolean{orig}}{I~}{}\ref{I.4.8}, and hence is an
open immersion, by \ifthenelse{\boolean{orig}}{I~}{}\ref{I.5.8} or \ifthenelse{\boolean{orig}}{I~}{}\ref{I.5.2}. Replacing \(X\) once
more by a suitable open neighborhood \(X_1\) of~\(x\), we obtain
\textup{(i)}.

The preceding argument establishes the equivalence of
\textup{(i)}, \textup{(ii)}, and \textup{(iii)}, and shows that they imply
\textup{(iv)}. It remains to prove that \textup{(iv)} implies
\textup{(ii)}. Since \(\Omega^1_{X/S}\) is free at~\(x\), this will be
immediate once we know that the smoothness of~\(Y\) over~\(S\) at~\(x\)
implies that \(\mathcal J/\mathcal J^2\) is free over~\(Y\) at~\(x\), and
that
\[
  \mathcal J/\mathcal J^2
  \longrightarrow
  \Omega^1_{X/S}\otimes_{\mathcal O_X}\mathcal O_Y
\]
is universally injective at~\(x\). This assertion is included in the
following theorem.

\begin{theorem}
\label{II.4.10}
Let \(X\) be an \(S\)-prescheme smooth over~\(S\), let \(Y\) be a closed
subprescheme of~\(X\) defined by a coherent sheaf of ideals~\(\mathcal J\),
and let \(x\in X\). The following conditions are equivalent:
\begin{enumerate}
\item[\textup{(i)}]
\(Y\) is \emph{smooth over~\(S\) at~\(x\)}.
\item[\textup{(ii)}]
There exist an open neighborhood \(X_1\) of~\(x\) in~\(X\), an étale
\(S\)-morphism
\[
  g\colon X_1\longrightarrow X'=S[t_1,\dots,t_n],
\]
and an integer \(p\), such that \(Y_1=Y\cap X_1\) is the
\emph{inverse image} under~\(g\) of the closed subprescheme
\[
  Y'=S[t_{p+1},\dots,t_n]
  \quad\text{of}\quad
  X'=S[t_1,\dots,t_n].
\]
\item[\textup{(iii)}]
There exist generators \(g_i\), \(1\leq i\leq p\), of~\(\mathcal J_x\)
such that the \(\mathrm dg_i\) form part of a basis of
\(\Omega^1_{X/S}\) at~\(x\), or equivalently such that the
\(\mathrm dg_i(x)\) are linearly independent over~\(\kappa(x)\) in
\(\Omega^1_{X/S}(x)\).
\item[\textup{(iv)}]
The sheaf \(\mathcal J/\mathcal J^2\) is free over~\(Y\) at~\(x\), and the
canonical homomorphism
\[
  \mathrm d\colon
  \mathcal J/\mathcal J^2
  \longrightarrow
  \Omega^1_{X/S}\otimes_{\mathcal O_X}\mathcal O_Y
\]
is universally injective at~\(x\). Equivalently, the sequence of canonical
homomorphisms
\[
  0\longrightarrow\mathcal J/\mathcal J^2
  \longrightarrow
  \Omega^1_{X/S}\otimes_{\mathcal O_X}\mathcal O_Y
  \longrightarrow
  \Omega^1_{Y/S}
  \longrightarrow0
\]
is exact at~\(x\), and \(\Omega^1_{Y/S}\) is locally free at~\(x\).
\end{enumerate}
\end{theorem}

\paragraph{Proof.}
We already know from the preceding argument that \textup{(ii)} implies
\textup{(i)}, \textup{(iii)}, and \textup{(iv)}. Let us prove that
\textup{(i)} implies \textup{(ii)}, which will at the same time complete
the proof of Proposition~\ref{II.4.9}.
\marginpar{\raggedright\footnotesize French source p.~42}%
By Theorem~\ref{II.4.3}\textup{(ii)}, the last two terms of the exact
sequence \eqref{eq:II.4.3} are locally free \(Y\)-modules at~\(x\). The
images of the elements \(\mathrm dg\), \(g\in\mathcal O_X\), generate
\(\Omega^1_{X/S}\otimes_{\mathcal O_X}\mathcal O_Y\) at~\(x\), and hence
their images generate \(\Omega^1_{Y/S}\) there. We can therefore choose
\(g_i\), \(p+1\leq i\leq n\), in~\(\mathcal O_X\) such that the
\(\mathrm dg'_i\) form a basis of~\(\Omega^1_{Y/S}\). By the exactness of
\eqref{eq:II.4.3}, we can then complete the
\(\mathrm dg_i\), \(p+1\leq i\leq n\), to a basis of the middle term by
elements of the form \(\mathrm dg_i\), \(1\leq i\leq p\), with
\(g_i\in\mathcal J_x\). The \(g_i\) come from sections of
\(\mathcal O_X\) on a neighborhood of~\(x\), which we may take to be~\(X\).
We are now in condition~\textup{(ii)} of Proposition~\ref{II.4.9}, and we
have proved that it implies condition~\textup{(i)} there. This gives
condition~\textup{(ii)} of the present theorem.

The implication \textup{(iii)} \(\Rightarrow\) \textup{(ii)} in the present
theorem follows immediately from the implication
\textup{(iii)} \(\Rightarrow\) \textup{(i)} in
Proposition~\ref{II.4.9}. Thus \textup{(i)}, \textup{(ii)}, and
\textup{(iii)} are equivalent and imply \textup{(iv)}. Finally,
\textup{(iv)} implies \textup{(iii)}: choose elements
\(g_i\in\mathcal J_x\) whose classes form a basis of
\(\mathcal J_x/\mathcal J_x^2\). They generate \(\mathcal J_x\) by
Nakayama's lemma, and universal injectivity says that their differentials
form part of a basis.

The preceding proof also gives the following result.

\begin{corollary}
\label{II.4.11}
Let \(X\) be an \(S\)-prescheme, let \(Y\) be a closed subprescheme defined
by a coherent sheaf of ideals~\(\mathcal J\), and let \(x\in Y\). Suppose
that both \(X\) and~\(Y\) are \emph{smooth over~\(S\) at~\(x\)}. Let
\(g_i\), \(1\leq i\leq p\), be sections of~\(\mathcal J\). The following
conditions are equivalent:
\begin{enumerate}
\item[\textup{(i)}]
The \(g_i\) \emph{generate} \(\mathcal J_x\), and the
\(\mathrm dg_i(x)\) are \emph{linearly independent} over~\(\kappa(x)\) in
\(\Omega^1_{X/S}(x)\).
\item[\textup{(ii)}]
The classes of the \(g_i\) modulo~\(\mathcal J^2\) form a basis of
\(\mathcal J/\mathcal J^2\) at~\(x\).
\item[\textup{(iii)}]
The \(g_i\) form a minimal system of generators of~\(\mathcal J_x\).
\item[\textup{(iv)}]
There exist further sections \(g_i\), \(p+1\leq i\leq n\), of
\(\mathcal O_X\) on a neighborhood \(X_1\) of~\(x\) in~\(X\), such that,
together with the preceding sections, they define an étale morphism
\[
  g\colon X_1\longrightarrow X'=S[t_1,\dots,t_n]
\]
for which \(Y_1=Y\cap X_1\) is the \emph{inverse image} of the closed
subprescheme
\[
  Y'=S[t_{p+1},\dots,t_n]
  \quad\text{of}\quad
  X'=S[t_1,\dots,t_n].
\]
\end{enumerate}
\end{corollary}

In particular:
\marginpar{\raggedright\footnotesize French source p.~43}%

\begin{corollary}
\label{II.4.12}
Let \(X\) be an \(S\)-prescheme, let \(F\) be a section of~\(\mathcal O_X\),
let \(Y\) be its zero subprescheme, defined by the coherent ideal
\(F\mathcal O_X\), and let \(x\in Y\). Suppose that \(X\) is smooth over
\(S\) at~\(x\). Then \(Y\) is smooth over~\(S\) at~\(x\) if and only if
either \(F\) vanishes in a neighborhood of~\(x\), or
\(\mathrm dF(x)\neq0\), where \(\mathrm dF(x)\) denotes the image of
\(\mathrm dF\) in the \(\kappa(x)\)-vector space
\(\Omega^1_{X/S}(x)\).
\end{corollary}

The condition is sufficient by criterion~\textup{(iii)} of
Theorem~\ref{II.4.10}. It is necessary because, as \(\mathcal J\) is
generated by one element, \(\mathcal J/\mathcal J^2\) at~\(x\) must first
be free of rank at most~\(1\). If its rank is~\(0\), that is, if
\(\mathcal J/\mathcal J^2=0\) at~\(x\), then
\(\mathcal J=0\) at~\(x\) by Nakayama's lemma, so \(F\) vanishes in a
neighborhood of~\(x\). If its rank is~\(1\), then \(F\) is a minimal system
of generators of~\(\mathcal J\) at~\(x\), and the conclusion follows from
the equivalence of \textup{(i)} and \textup{(iii)} in
Corollary~\ref{II.4.11}.

\begin{corollary}
\label{II.4.13}
Let \(Y\) be an \(S\)-prescheme locally of finite type, let \(S'\) be a
flat \(S\)-prescheme, put \(Y'=Y\times_SS'\), let \(x'\in Y'\), and let
\(x\in Y\) be its canonical image. Then \(Y\) is smooth over~\(S\) at~\(x\)
if and only if \(Y'\) is smooth over~\(S'\) at~\(x'\). In particular, if
\(S'\to S\) is flat and surjective, then \(Y\) is smooth over~\(S\) if and
only if \(Y'\) is smooth over~\(S'\).
\end{corollary}

Only sufficiency remains to be proved; necessity was established in
Proposition~\ref{II.1.3}\textup{(ii)}. After replacing \(Y\) by a suitable
neighborhood of~\(x\), and \(Y'\) by its inverse image, we may suppose that
\(Y\) is affine and of finite type over the affine prescheme~\(S\). Thus
\(Y\) is isomorphic to a closed subprescheme of
\[
  X=S[t_1,\dots,t_n].
\]
It follows that \(Y'\) is a closed subprescheme of
\(X'=X\times_SS'\). Since \(X\) is smooth over~\(S\), the prescheme
\(X'\) is smooth over~\(S'\), and the smoothness criteria in
Theorem~\ref{II.4.10} apply. Criterion~\textup{(iv)} gives the result
immediately.
\end{filecontents*}
% Original input: drafts/SGA1_II_4_LINES_3500_3727_English.texfrag
\begin{filecontents*}[overwrite,noheader]{sga1-cumulative-input-017.tex}
% Source: smf_doc-math_3_01.tex, lines 3500--3727 inclusive.

\begin{remarques}
\label{II.4.14}
Criterion (iii) of Theorem~\Ref{II.4.10} deserves to be called the
\emph{Jacobian criterion for smoothness}.
%
\index{Jacobian criterion for smoothness|hyperpage}%
\index{Jacobian (criterion for smoothness)|hyperpage}%
\index{smoothness (Jacobian criterion for)|hyperpage}%
%
It makes it possible, in theory, to determine whether a given
$S$-prescheme $Y$ is smooth over~$S$ at a point $x$ of~$Y$, since there
always exists a neighborhood of~$x$ in~$Y$ isomorphic to a subprescheme of a
smooth $S$-prescheme $X$, for example $X=S[t_1,\dots,t_n]$. Moreover, it is
for $X=S[t_1,\dots,t_n]$, $S=\Spec(A)$, that the Jacobian criterion is
usually stated (of course,
% page: 44
\marginpar{44}
%
in the classical case considered by Zariski, $A$ was a field). We leave it
to the reader to give the statement for an ideal $J$ of
~$A[t_1,\dots,t_n]$ and a prime ideal containing it, to which one is thus
led. Let us note that it now seems clear (especially since Nagata succeeded,
by nondifferential methods, in generalizing Zariski's theorem that the set
of regular points of an algebraic scheme is open) that the Jacobian
criterion is of scarcely any interest except in the form in which we give
it here (\ie using exclusively \emph{relative} differentials and not
\emph{absolute} differentials, \ie differentials relative to the ring of
absolute constants~$\ZZ$). As is so often the case, the use of differentials
is more convenient here than that of derivations. Finally, let us note that
% SGA1-CORR-0022: the relative dimension is m; n=m+1 is the ambient dimension.
if $Y$ is smooth over~$S$ at $x$, of relative dimension $m$, then there
exists an open neighborhood of~$x$ in~$Y$ isomorphic to a subprescheme of
$X=S[t_1,\dots,t_n]$ with $n=m+1$, as follows from the definition and
\ifthenelse{\boolean{orig}}{from~I~\Ref{I.7.6}).}{from~(\Ref{I.7.6}).}
% correction: the opening parenthesis is missing
% (I~\Ref{I.7.6})

Let $A$ be a noetherian ring, let $x_i$ ($1 \leq i \leq n$) be elements
of~$A$, and let $J$ be the ideal generated by the $x_i$. The $x_i$ are said
to form a \emph{regular system of generators}
%
\index{regular (system of generators)@regular (system of generators)|hyperpage}%
\index{regular system of generators)@regular system of generators|hyperpage}%
%
of~$J$ if the canonical surjective homomorphism
$$
(A/J)[t_1,\dots,t_n] \to \gr^J(A)
$$
defined by the $x_i$ (where the second member denotes the graded ring
associated with $A$ filtered by the powers of~$J$) is an
\emph{isomorphism}. This condition also means that
\begin{enumerate}
\item[(i)] The canonical surjective homomorphism
$$
S_{A/J} (J/J^2) \to \gr^J(A)
$$
(where the first member denotes the symmetric algebra of the $A/J$-module
$J/J^2$) is an isomorphism, and
\item[(ii)] $J/J^2$ is free and has as a basis the classes of the $x_i
\mod J^2$.
\end{enumerate}
In this form, we see that if $J\neq A$, the $x_i$ form a \emph{minimal
system of generators} of~$J$, and that \emph{every other minimal system of
generators} of~$J$ \emph{is a regular system of generators} (N.B.
\og minimal\fg is taken in the strict sense: minimum number of elements,
which is equivalent to the sense: minimal under inclusion only if $A$ is
local); on the other hand, if $J=A$, every system of generators of~$J$ is
regular.

The
% page: 45
\marginpar{45}
%
regularity condition for a system of generators of an ideal is stable under
localization by any multiplicatively closed set; on the other hand, it is
immediately seen that in order for $(x_i)$ to be a minimal system of
generators of~$J$, it already suffices that for every \emph{maximal ideal
$\goth m$ containing $J$}, the $x_i$ define a regular system of generators
of~$JA_{\goth m}$ in $A_{\goth m}$. We are thus reduced to the case where
$A$ is a local ring with maximal ideal $\goth m$, and where the $x_i$ lie
in $\goth m$. \emph{Then the $x_i$ form a regular system of generators
of~$J$ if and only if they form an $A$-sequence in the sense of
Serre}\footnote{We shall now rather say \og $A$-regular sequence\fg, \cf EGA
$0_{\textup{IV}}$~15.1.7 and 15.1.11.}, \ie if for every $i$ such that $1
\leq i \leq n$, $x_i$ is a non-zero-divisor of~$0$ in
\ifthenelse{\boolean{orig}}{$A/(x_1,\dots,x_i)A$.}
{$A/(x_1,\dots,x_{i-1})A$.}
% correction: $A/(x_1,\dots,x_{i-1})A$.

Finally, in the case where $A$ is an algebra over a ring $B$, and where
$A/J$ is isomorphic as a $B$-algebra to $B$ (so that $J$ is the kernel of a
homomorphism of $B$-algebras $A\to B$), the $x_i$ form a regular system of
generators of~$J$ if and only if the canonical homomorphism
$$
B[[ t_1,\dots,t_n]]\to \hat{A}
$$
defined by the $x_i$ (where the second member denotes the separated
completion $\varprojlim A/J^{n+1}$ of~$A$ for the topology defined by the
powers of~$J$) is an \emph{isomorphism} (it is in any case
\emph{surjective}).

All these facts are well known, and doubtless appear, up to minor
differences, in Serre's course on commutative algebra written up by Gabriel.
(Where one finds $N$ other characterizations of $A$-sequences, in the case
where $A$ is a local ring).

Let $J$ be an ideal in a noetherian ring $A$. We shall say that $J$ is a
\emph{regular ideal}
%
\index{regular ideal|hyperpage}%
\index{regular (ideal)@regular (ideal)|hyperpage}%
%
if for every prime ideal $\goth p$ of~$A$, $JA_{\goth p}$ admits a regular
system of generators. It is plainly enough to verify this for
$\goth p \supset J$, and one may moreover restrict oneself to maximal
$\goth p$. More generally, let $\cal{J}$ be an Ideal on a locally
noetherian prescheme $X$; one says that $\cal{J}$ is a \emph{regular Ideal}
if for every $x\in X$, $\cal{J}_x$ is an ideal of $\cal{O}_x$ which admits a
regular system of generators. This is equivalent to the conjunction of the
following two conditions:
\begin{enumerate}
\item[(a)] The canonical surjective homomorphism
$$
\ifthenelse{\boolean{orig}}
{S_{\cal{O}/\cal{J}}(\cal{J}/\cal{J}^2) \to \gr^{\cal{J}}(\cal{O}_X)}
{S_{\cal{O}_X/\cal{J}}(\cal{J}/\cal{J}^2) \to \gr^{\cal{J}}(\cal{O}_X)}
$$
% correction:
% S_{\cal{O}_X/\cal{J}}(\cal{J}/\cal{J}^2) \to \gr^{\cal{J}}(\cal{O}_X)
is an isomorphism, and
\item[(b)] The sheaf of $\cal{O}_X/\cal{J}$-Modules
$\cal{J}/\cal{J}^2$ is locally free.
\end{enumerate}

One
% page: 46
\marginpar{46}
%
also says then that the subprescheme $Y$ of~$X$ defined by $\cal{J}$ (thus
such that $\cal{O}_Y$, extended by $0$, is isomorphic to
$\cal{O}_X/\cal{J}$) is \emph{regularly immersed}
%
\index{regularly immersed (prescheme)@regularly immersed (prescheme)|hyperpage}%
%
in~$X$, and one defines in the same way (in the evident manner) the notion
of a morphism that is a \emph{regular immersion},
%
\index{regular immersion|hyperpage}%
\index{regular (immersion)@regular (immersion)|hyperpage}%
%
(\resp \emph{regular at a point $x$}), an immersion morphism $Y\to X$
identifying $Y$ (\resp a suitable neighborhood of $x$) with a regularly
immersed closed subprescheme of an open subset of~$X$. (One must not say:
regular subprescheme, for that would mean that the local rings of~$Y$ are
regular). Finally, sections $x_i$ of~$\cal{J}$ are called a \emph{regular
system of generators} if for every $x\in X$, the corresponding elements
of~$\cal{O}_x$ form a regular system of generators of~$\cal{J}_x$
(terminology compatible with that introduced for generators of an ideal of
a ring). This also means that the canonical surjective homomorphism
$$
\cal{O}_Y[t_1,\dots,t_n] \to \gr^{\cal{J}} (\cal{O}_X)
$$
defined by the $x_i$ is an isomorphism. If it is known in advance that the
Ideal $\cal{J}$ is regular, this means simply that at every point $x$
\emph{of~$Y$}, the $x_i$ define a \emph{basis} of~$\cal{J}/\cal{J}^2$
over~$\cal{O}_{Y,x}$. (N.B. this condition is empty if $Y$ is empty). Thus,
in order that $\cal{J}$ admit a regular system of generators, it is
necessary and sufficient that $\cal{J}$ be regular, and that the
$\cal{O}_Y$-Module $\cal{J} /\cal{J}^2$ be globally free (and not merely
locally free), \ie that the canonical homomorphism
$S_{\cal{O}_Y}(\cal{J}/\cal{J}^2) \to
\gr^{\cal{J}} (\cal{O}_X)$ be surjective, and that the
$\cal{O}_Y$-Module $\cal{J}/\cal{J}^2$ be globally free.

An \emph{augmented ring is said to be regular}
%
\index{regular augmented ring|hyperpage}%
\index{regular (augmented ring@regular (augmented ring)|hyperpage}%
%
if the augmentation ideal is regular. Thus, if $A$ is a local ring,
regarded as augmented over its residue field $k$, then $A$ is a regular
local ring if and only if it is a regular augmented ring.

(To tell the truth, it seems that it was unnecessary to begin by giving the
preliminary chain of reasoning for rings; it is advantageous to begin with
sheaves at once. If one wants something in the noetherian case, it is the
definition adopted here --- a priori less restrictive than that by Serre's
$A$-sequences --- which seems preferable for the needs of differential
calculus. Of course, to do things properly, one would also have to develop
at least part of the theory of smooth morphisms in the nonnoetherian
setting\footnote{As is said in the foreword, this has now been done, \cf EGA
IV 17, 18}, probably by
% page: 47
\marginpar{47}
%
starting from the Jacobian criterion, so as to obtain, if possible, all the
essential formal properties of smooth morphisms and of \`etale morphisms,
\ie smooth and quasi-finite morphisms; only the converses making use of
noetherian hypotheses).
\end{remarques}

After these long terminological preliminaries, a little theorem:
\end{filecontents*}
% Original input: drafts/SGA1_II_4_OPENING_THROUGH_4_8_English.texfrag
\begin{filecontents*}[overwrite,noheader]{sga1-cumulative-input-018.tex}
% SGA 1, Exposé II, French TeX lines 2928--3234
% (section II.4 opening through Corollary II.4.8 and (4.3 bis)).

\section{Differential properties of smooth morphisms}
\label{II.4}

For simplicity, we shall for the most part restrict ourselves to
first-order differential calculus, giving only brief indications for higher
orders, where the results are just as simple.

For the definition of the sheaf \(\Omega^1_{X/Y}\) of first differentials of
a \(Y\)-prescheme~\(X\), see I, no.~\ref{I.1}. Suppose that \(X\) and~\(Y\)
are \(S\)-preschemes and that the morphism \(f\colon X\to Y\) is an
\(S\)-morphism. Then \(f\) defines a homomorphism of sheaves of modules,
compatible with~\(f\),
\begin{equation}
\label{eq:II.4.1}
  f^*\colon\Omega^1_{Y/S}\longrightarrow\Omega^1_{X/S}.
\end{equation}
In other words, \(\Omega^1_{X/S}\) is \emph{contravariant} in the
\(S\)-prescheme~\(X\). Equivalently, \eqref{eq:II.4.1} gives a homomorphism
of modules on~\(X\),
\begin{equation*}
\tag{4.1\(\mathrm{bis}\)}
\label{eq:II.4.1bis}
  f^*\bigl(\Omega^1_{Y/S}\bigr)\longrightarrow\Omega^1_{X/S},
\end{equation*}
also denoted by \(f^*\), for lack of a better notation. It belongs to the
canonical exact sequence
\begin{equation}
\label{eq:II.4.2}
  f^*\bigl(\Omega^1_{Y/S}\bigr)
  \longrightarrow\Omega^1_{X/S}
  \longrightarrow\Omega^1_{X/Y}
  \longrightarrow0.
\end{equation}
All these homomorphisms are defined by the requirements that they be local,
which reduces their definition to the affine case, and that they commute
with the differential operators \(\mathrm d\). The exactness of
\eqref{eq:II.4.2} is classical and immediate. In the affine case, for a
homomorphism \(B\to C\) of \(A\)-algebras, it becomes the exact sequence
\begin{equation*}
\tag{4.2\(\mathrm{bis}\)}
\label{eq:II.4.2bis}
  \Omega^1_{B/A}\otimes_B C
  \longrightarrow\Omega^1_{C/A}
  \longrightarrow\Omega^1_{C/B}
  \longrightarrow0.
\end{equation*}

\begin{lemma}
\label{II.4.1}
Let \(f\colon X\to Y\) be a morphism of \(S\)-preschemes.
\marginpar{\raggedright\footnotesize French source p.~36}%
If \(f\) is unramified (respectively, étale), then
\[
  f^*\bigl(\Omega^1_{Y/S}\bigr)\longrightarrow\Omega^1_{X/S}
\]
is surjective (respectively, an isomorphism). The converse in the
unramified case is true if \(f\) is assumed locally of finite type.
\end{lemma}

The unramified assertion follows from the exact sequence
\eqref{eq:II.4.2} and \ifthenelse{\boolean{orig}}{I~}{}\ref{I.3.1}; it may also be seen directly as part
of the étale case. For a direct proof of the latter, consider the diagram
\[
\begin{array}{ccccc}
X & \xrightarrow{\ \Delta_{X/Y}\ } & X\times_Y X
  & \longrightarrow & X\times_S X\\
&& \big\downarrow && \big\downarrow\\
&& Y & \xrightarrow{\ \Delta_{Y/S}\ } & Y\times_S Y .
\end{array}
\]
Here \(X\times_YX\) is the fiber product of~\(Y\) and \(X\times_SX\) over
\(Y\times_SY\). Since \(f\) is unramified,
\(X\to X\times_YX\) is an open immersion. Hence the conormal sheaf of the
composite immersion \(\Delta_{X/S}\) is isomorphic to the inverse image on
\(X\) of the conormal sheaf of
\(X\times_YX\to X\times_SX\). In the étale case, \(X\to Y\) is flat, so
\(X\times_SX\to Y\times_SY\) is flat. The conormal sheaf of
\(X\times_YX\to X\times_SX\) is therefore the inverse image of the
conormal sheaf of \(Y\to Y\times_SY\), that is, the inverse image of
\(\Omega^1_{Y/S}\). The assertion follows.

\begin{lemma}
\label{II.4.2}
Let \(X=Y[t_1,\dots,t_n]\), where \(Y\) is an \(S\)-prescheme. Then the
sequence of canonical homomorphisms
\[
  0\longrightarrow f^*\bigl(\Omega^1_{Y/S}\bigr)
  \longrightarrow\Omega^1_{X/S}
  \longrightarrow\Omega^1_{X/Y}
  \longrightarrow0
\]
is exact, and \(\Omega^1_{X/Y}\) is free with basis
\(\mathrm d_{X/Y}t_1,\dots,\mathrm d_{X/Y}t_n\).
\end{lemma}

The verification, which is purely affine, is immediate. Notice that the
exactness at the right was already known from \eqref{eq:II.4.2}.

Combining these two statements with Definition~\ref{II.1.1}, we obtain:

\begin{theorem}
\label{II.4.3}
Let \(f\colon X\to Y\) be a smooth morphism of \(S\)-preschemes. Then:
\begin{enumerate}
\item[\textup{(i)}]
the sequence of canonical homomorphisms
\[
  0\longrightarrow f^*\bigl(\Omega^1_{Y/S}\bigr)
  \longrightarrow\Omega^1_{X/S}
  \longrightarrow\Omega^1_{X/Y}
  \longrightarrow0
\]
is exact;
\item[\textup{(ii)}]
\(\Omega^1_{X/Y}\) is locally free, and its rank \(n\) at~\(x\) is equal to
the relative dimension of~\(f\) at~\(x\).
\end{enumerate}
\end{theorem}

\begin{corollary}
\label{II.4.4}
The homomorphism
\[
  f^*\bigl(\Omega^1_{Y/S}\bigr)\longrightarrow\Omega^1_{X/S}
\]
is injective, and its image in \(\Omega^1_{X/S}\) is locally a direct
summand.
\marginpar{\raggedright\footnotesize French source p.~37}%
\end{corollary}

Let \(u\colon F\to G\) be a homomorphism of modules on the prescheme~\(X\).
It is said to be \emph{universally injective} at~\(x\in X\) if the
homomorphism \(F_x\to G_x\) of \(\mathcal O_x\)-modules is injective and
remains so after tensoring with every \(\mathcal O_x\)-algebra, or,
equivalently, with every \(\mathcal O_x\)-module. It is enough, for example,
that there exist an open neighborhood \(U\) of~\(x\) on which \(u\)
identifies \(F|_U\) with a direct summand of~\(G|_U\). This condition is also
necessary when \(F\) and \(G\) are free of finite type in a neighborhood
of~\(x\). More precisely, in that case the following conditions are
equivalent:
\begin{enumerate}
\item[\textup{(i)}]
\(u\) is injective at~\(x\), and \(\operatorname{Coker}(u)\) is free
at~\(x\);
\item[\textup{(ii)}]
there exists an open neighborhood \(U\) of~\(x\) on which \(u\) identifies
\(F|_U\) with a direct summand of~\(G|_U\);
\item[\textup{(iii)}]
\(u\) is universally injective at~\(x\);
\item[\textup{(iv)}]
the homomorphism on the fibers induced by~\(u\),
\[
  F_x\otimes_{\mathcal O_x}\kappa(x)
  \longrightarrow
  G_x\otimes_{\mathcal O_x}\kappa(x),
\]
is injective;
\item[\textup{(v)}]
the transpose homomorphism \(G^\vee\to F^\vee\) is surjective at~\(x\), or,
equivalently, in a neighborhood of~\(x\).
\end{enumerate}
The implications form a circle: \textup{(iv)} implies \textup{(v)} by
Nakayama's lemma, while \textup{(v)} implies \textup{(i)} because a locally
free quotient sheaf is necessarily a direct summand. Geometrically, the
situation means that \(u\) corresponds to an isomorphism of the vector
bundle whose sheaf of sections is~\(F\) onto a vector subbundle of the
analogous vector bundle defined by~\(G\). It is, of course, not enough for
this that \(F\to G\) merely be injective.

\begin{corollary}
\label{II.4.5}
Let \(f\colon X\to Y\) be a morphism of \(S\)-preschemes, locally of finite
type. Let \(x\in X\), put \(y=f(x)\), and let \(s\) be the common image of
\(x\) and~\(y\) in~\(S\). Suppose that \(Y\) is smooth over~\(S\) at~\(y\).
Then the following conditions are equivalent:
\begin{enumerate}
\item[\textup{(i)}]
\(f\) is smooth at~\(x\);
\item[\textup{(ii)}]
\(X\) is smooth over~\(S\) at~\(x\), and
\[
  f^*\bigl(\Omega^1_{Y/S}\bigr)\longrightarrow\Omega^1_{X/S}
\]
is universally injective at~\(x\); that is, it is injective at~\(x\), and
its cokernel \(\Omega^1_{X/Y}\) is free at~\(x\).
\end{enumerate}
\end{corollary}

Necessity follows from Proposition~\ref{II.1.3}\textup{(iii)} and
Theorem~\ref{II.4.3}\textup{(i)}, \textup{(ii)}. Let us prove sufficiency.
Since the elements \(\mathrm dg\), \(g\in\mathcal O_x\), generate
\((\Omega^1_{X/Y})_x\), we can choose \(g_i\), \(1\leq i\leq n\), such that
the images of the \(\mathrm dg_i\) form a basis of this module.
\marginpar{\raggedright\footnotesize French source p.~38}%
After shrinking \(X\), we may suppose that the \(g_i\) come from sections of
\(\mathcal O_X\), and hence define a \(Y\)-morphism
\[
  g\colon X\longrightarrow Y'=Y[t_1,\dots,t_n].
\]
Using the hypothesis and Lemma~\ref{II.4.2}, one sees immediately that the
corresponding homomorphism
\[
  g^*\bigl(\Omega^1_{Y'/S}\bigr)\longrightarrow\Omega^1_{X/S}
\]
is an isomorphism at~\(x\). We are thus reduced to proving the following
corollary.

\begin{corollary}
\label{II.4.6}
Let \(f\colon X\to Y\) be a morphism between \(S\)-preschemes that are
smooth over~\(S\). Then \(f\) is étale at~\(x\in X\) if and only if
\[
  f^*\bigl(\Omega^1_{Y/S}\bigr)\longrightarrow\Omega^1_{X/S}
\]
is an isomorphism at~\(x\).
\end{corollary}

Necessity is known from Lemma~\ref{II.4.1}, and the same lemma shows that
the stated condition implies that \(f\) is unramified at~\(x\). By
Corollary~\ref{II.2.2}, we are reduced to the case
\(S=\operatorname{Spec}(k)\). Since \(Y\) is smooth over~\(k\), it is
regular, hence in particular normal. By \ifthenelse{\boolean{orig}}{I~}{}\ref{I.9.5}\textup{(ii)}, it
remains only to prove that \(\mathcal O_y\to\mathcal O_x\) is injective, or,
equivalently, that \(\mathcal O_y\) and \(\mathcal O_x\) have the same
dimension. These dimensions are respectively the ranks of
\(\Omega^1_{Y/k}\) and \(\Omega^1_{X/k}\) at \(y\) and~\(x\), and are
therefore equal by hypothesis.

\refstepcounter{proposition}
\paragraph{Remarks \theproposition.}\label{II.4.7}
When \(X\) and \(Y\) are smooth over~\(S\), the smoothness criterion
in Corollary~\ref{II.4.5}\textup{(ii)} for \(f\colon X\to Y\) can also be
stated by saying that, for every \(x\in X\), the \emph{tangent} map of~\(f\)
at~\(x\), relative to the base~\(S\), is surjective. This tangent map is the
transpose of the homomorphism between the finite-dimensional
\(\kappa(x)\)-vector spaces obtained as the fibers at~\(x\) of
\(f^*(\Omega^1_{Y/S})\) and \(\Omega^1_{X/S}\). This is a familiar
hypothesis, especially among those who work with analytic spaces. Their
usual nonsingularity hypothesis, which says that \(X\) and \(Y\) are
\lq\lq smooth over~\(\mathbb C\)\rq\rq{}; see
\hyperref[II.5]{no.~II.5}, seems to be
due only to the fear that singular points of algebraic varieties or analytic
spaces still inspire in many geometers.

We record the following special case of Corollary~\ref{II.4.6}.

\begin{corollary}
\label{II.4.8}
Let \(X\) be an \(S\)-prescheme, and let
\[
  g\colon X\longrightarrow S[t_1,\dots,t_n]
\]
be an \(S\)-morphism defined by sections \(g_i\), \(1\leq i\leq n\), of
\(\mathcal O_X\). Let \(x\in X\) be a point at which \(X\) is smooth
over~\(S\). Then \(g\) is étale at~\(x\) if and only if the
\(\mathrm dg_i\), \(1\leq i\leq n\), form a basis of
\(\Omega^1_{X/S}\) at~\(x\). Equivalently, their images form a basis of the
\(\kappa(x)\)-vector space
\[
  \Omega^1_{X/S}(x)
  =(\Omega^1_{X/S})_x\otimes_{\mathcal O_x}\kappa(x).
\]
\marginpar{\raggedright\footnotesize French source p.~39}%
\end{corollary}

Let \(X\) be a prescheme, and let \(Y\) be a closed subprescheme of~\(X\)
defined by a coherent sheaf of ideals~\(\mathcal J\). Then
\(\mathcal J/\mathcal J^2\) may be regarded as a coherent sheaf on~\(Y\);
it is the \emph{conormal sheaf} of~\(Y\) in~\(X\). If \(X\) is an
\(S\)-prescheme, there is a canonical exact sequence of quasi-coherent
sheaves on~\(Y\):
\begin{equation}
\label{eq:II.4.3}
  \mathcal J/\mathcal J^2
  \xrightarrow{\ \mathrm d\ }
  \Omega^1_{X/S}\otimes_{\mathcal O_X}\mathcal O_Y
  \longrightarrow\Omega^1_{Y/S}
  \longrightarrow0.
\end{equation}
The right-hand part is simply \eqref{eq:II.4.2}, with the roles of \(X\)
and~\(Y\) interchanged and with \(\Omega^1_{Y/X}=0\). The homomorphism
\[
  \mathcal J/\mathcal J^2
  \longrightarrow
  \Omega^1_{X/S}\otimes_{\mathcal O_X}\mathcal O_Y
\]
is induced, on passing to quotients, by the generally nonlinear map
\(g\mapsto\mathrm dg\). The exactness of \eqref{eq:II.4.3} is classical and
immediate. In the affine case, for a surjective homomorphism \(B\to C\) of
\(A\)-algebras with kernel~\(J\), it becomes
\begin{equation*}
\tag{4.3\(\mathrm{bis}\)}
\label{eq:II.4.3bis}
  J/J^2
  \longrightarrow\Omega^1_{B/A}\otimes_B C
  \longrightarrow\Omega^1_{C/A}
  \longrightarrow0,
  \qquad C=B/J.
\end{equation*}
This is the exact sequence already used implicitly in the proof of
\ifthenelse{\boolean{orig}}{I~}{}\ref{I.7.2}.
\end{filecontents*}
% Original input: drafts/SGA1_II_5_6_THROUGH_5_7_English.texfrag
\begin{filecontents*}[overwrite,noheader]{sga1-cumulative-input-019.tex}
\begin{corollary}
\label{II.5.6}
Let \(K\) be a local Artinian ring which is a localization of a
finitely generated algebra over the field~\(k\) (for example, \(K\) may
be a finitely generated extension of~\(k\)), and let \(n\) be the
transcendence degree over~\(k\) of its residue field. The following
conditions are equivalent:
\begin{enumerate}
\item[(i)] \(K\) is a finite separable extension of a purely
transcendental extension \(k(t_1,\dots,t_n)\) of~\(k\).
\item[(ii)]
\marginpar{\raggedright\footnotesize French source p.~54}%
\(\Omega^1_{K/k}\) is a free \(K\)-module of rank~\(n\).
\item[(ii bis)] \(\Omega^1_{K/k}\) is a \(K\)-module generated by
\(n\) elements.
\item[(iii)] The completion \(O'\) of \(K\otimes_k K\), for the
topology defined by the powers of the augmentation ideal of
\(K\otimes_k K\to K\), is a \lq\lq regular\rq\rq{} augmented
\(K\)-algebra, that is, isomorphic to a formal power-series algebra
over~\(K\). (If \(K\) is a field, this is equivalent to saying that
\(O'\) is a regular local ring.)
\item[(iv)] \(K\) is a separable extension of~\(k\).
\end{enumerate}
\end{corollary}

Indeed, \(K\) may always be regarded as the local ring at the generic
point of an irreducible scheme \(X\) of finite type over~\(k\), and the
conditions in question are the conditions bearing the same names in
Theorem~\ref{II.5.5}, taking for \(k'\) in~\textup{(iv)} an
algebraically closed extension of~\(k\). Only the implication
\[
  K\text{ separable over }k
  \quad\Longrightarrow\quad
  X\text{ smooth over }k\text{ at }x
\]
requires proof. By Corollary~\ref{II.4.13}, we are immediately reduced
to the case where the base field is \(k'\), hence algebraically closed;
there is then a \(k\)-rational point \(a\) of~\(X\). By
Proposition~\ref{II.5.4}, \(X\) is smooth over~\(k\) at~\(a\), and is
therefore all the more smooth over~\(k\) at the generic point~\(x\).
\footnote{See the Errata at the end of the present
\hyperref[II.fin.errata]{Exposé~II}
(p.~\pageref{II.fin.errata}).}

When \(K\) is a finitely generated extension of~\(k\), the equivalence
of \textup{(i)}, \textup{(ii)}, \textup{(ii bis)}, and \textup{(iv)}
is well known. Notice, however, that we have used none of these already
known equivalences. Of course, Proposition~\ref{II.5.1} contains as a
special case the familiar fact that a sequence of elements \(x_i\),
\(1\leq i\leq n\), is a separating transcendence basis of \(K\)
over~\(k\) if and only if the \(\mathrm d x_i\) form a basis of the
\(K\)-module \(\Omega^1_{K/k}\).

\begin{corollary}
\label{II.5.7}
Let \(X\) be a prescheme of finite type over a field~\(k\). Then \(X\)
is smooth over~\(k\) if and only if \(\Omega^1_{X/k}\) is locally free
and the local rings at the generic points of the irreducible components
of~\(X\) are separable extensions of~\(k\). (The latter condition is
automatically satisfied if \(k\) is perfect and \(X\) is reduced.)
\end{corollary}

We may assume that \(X\) is connected. Let \(n\) be the rank of
\(\Omega^1_{X/k}\), assumed locally free. By the hypothesis and
Corollary~\ref{II.5.6}, \(n\) is also the transcendence degree of the
extensions of~\(k\) defined by the local rings at the generic points
of~\(X\).
\marginpar{\raggedright\footnotesize French source p.~55}%
Thus all the irreducible components of~\(X\) have dimension~\(n\), and
the conclusion follows from Theorem~\ref{II.5.5}.

Care is needed here. If \(K\) is a finite extension of~\(k\), not
necessarily separable, then \(\Omega^1_{K/k}\) is a free \(K\)-module.
Thus, putting \(X=\operatorname{Spec}(K)\), the sheaf
\(\Omega^1_{X/k}\) is locally free and \(X\) is reduced, although \(X\)
need not be smooth over~\(k\). Extending scalars to the algebraic
closure of~\(k\) then gives a similar example in which \(k\) is
algebraically closed, but \(X\) is no longer reduced.
\end{filecontents*}
% Original input: drafts/SGA1_II_5_8_THROUGH_5_10_English.texfrag
\begin{filecontents*}[overwrite,noheader]{sga1-cumulative-input-020.tex}
\begin{corollary}
\label{II.5.8}
Let \(X\) be a prescheme of finite type over the field~\(k\), let
\(x\in X\), let \(n\) be the dimension of~\(X\) at~\(x\), and let
\(p=\dim\mathcal O_x\), that is, the codimension in~\(X\) of the
closure \(Y\) of~\(x\). Thus \(n-p\) is the transcendence degree of
\(\kappa(x)\) over~\(k\). Let \(f_i\), \(1\leq i\leq n\), be elements
of~\(\mathcal O_x\), with \(f_i\in\mathfrak m_x\) for
\(1\leq i\leq p\). The following conditions are equivalent:
\begin{enumerate}
\item[(i)] The germ at~\(x\) of the morphism
\[
  X\longrightarrow\operatorname{Spec}\bigl(k[t_1,\dots,t_n]\bigr)
\]
defined by the \(f_i\) is étale at~\(x\).
\item[(ii)] The elements \(f_i\), \(1\leq i\leq p\), generate
\(\mathfrak m_x\), that is, they form a regular system of parameters
of~\(\mathcal O_x\), and the classes in \(\kappa(x)\) of the \(f_j\),
\(p+1\leq j\leq n\), form a separating transcendence basis. In other
words, the \(\mathrm d\overline f_j\), \(p+1\leq j\leq n\), form a
basis of \(\Omega^1_{\kappa(x)/k}\), or equivalently generate that
module.
\end{enumerate}
\end{corollary}

Suppose first that \textup{(i)} holds. By
Corollary~\ref{II.4.8}, the \(\mathrm d f_i(x)\) form a basis of
\(\Omega^1_{X/k}(x)\); hence their images
\(\mathrm d\overline f_i(x)\) generate \(\Omega^1_{\kappa(x)/k}\)
over~\(\kappa(x)\). Since \(\overline f_i=0\) for
\(1\leq i\leq p\), it is enough to retain the
\(\mathrm d\overline f_i(x)\) with \(p+1\leq i\leq n\). The
transcendence degree of \(\kappa(x)\) over~\(k\) is \(n-p\), so
Corollary~\ref{II.5.6}, applied to \(K=\kappa(x)\), shows that \(Y\) is
smooth over~\(k\) at its generic point~\(x\), and that the
\(\mathrm d\overline f_i(x)\), \(p+1\leq i\leq n\), form a basis of
\(\Omega^1_{\kappa(x)/k}\) over~\(\kappa(x)\). Condition
\textup{(ii)} of Proposition~\ref{II.4.9} is therefore satisfied, and
so is condition \textup{(iii)}; in particular the \(f_i\),
\(1\leq i\leq p\), generate \(\mathfrak m_x\). Since
\marginpar{\raggedright\footnotesize French source p.~56}%
\(\mathcal O_x\) has dimension~\(p\), they form a regular system of
parameters at~\(x\). This proves \textup{(ii)}.

Conversely, suppose that \textup{(ii)} holds. The exact sequence
\eqref{eq:II.5.1} shows that the \(\mathrm d f_i(x)\) generate
\(\Omega^1_{X/k}(x)\), and \textup{(i)} follows from
Proposition~\ref{II.5.1}.

\begin{corollary}
\label{II.5.9}
Let \(X\) be a prescheme of finite type over the field~\(k\), let
\(x\in X\), let \(n\) be the dimension of~\(X\) at~\(x\), and let
\(p=\dim\mathcal O_x\), that is, the codimension of the closure \(Y\)
of~\(x\) in~\(X\). Thus \(n-p\) is the transcendence degree of
\(\kappa(x)\) over~\(k\). The following conditions are equivalent:
\begin{enumerate}
\item[(i)] \(\mathcal O_x\) is regular and \(\kappa(x)\) is a
separable extension of~\(k\).
\item[(ii)] \(X\) is smooth over~\(k\) at~\(x\), and the canonical
homomorphism
\[
  \mathfrak m_x/\mathfrak m_x^2
  \longrightarrow
  \Omega^1_{\mathcal O_x/k}\otimes_{\mathcal O_x}\kappa(x)
  =\Omega^1_{X/k}(x)
\]
is injective.
\item[(iii)] There exist \(f_i\in\mathcal O_x\),
\(1\leq i\leq n\), with \(f_i\in\mathfrak m_x\) for
\(1\leq i\leq p\), such that the germ at~\(x\) of the morphism from
\(X\) to \(\operatorname{Spec}\bigl(k[t_1,\dots,t_n]\bigr)\) defined
by the \(f_i\) is étale at~\(x\). Equivalently, by
Proposition~\ref{II.5.1}, the \(\mathrm d f_i(x)\) generate
\(\Omega^1_{X/k}(x)\).
\item[(iv)] There exist \(f_i\in\mathcal O_x\),
\(1\leq i\leq n\), such that the \(f_i\),
\(1\leq i\leq p\), generate \(\mathfrak m_x\), and the
\(\mathrm d f_j(x)\), \(p+1\leq j\leq n\), generate
\(\Omega^1_{\kappa(x)/k}\) over~\(\kappa(x)\).
\end{enumerate}
\end{corollary}

The equivalence of \textup{(iii)} and \textup{(iv)} follows from
Corollary~\ref{II.5.8}. By Proposition~\ref{II.4.9}, these conditions
are also equivalent to \(X\) being smooth over~\(k\) at~\(x\) together
with condition \textup{(ii)} of Theorem~\ref{II.4.10}. They are
therefore equivalent to \(X\) being smooth over~\(k\) at~\(x\)
together with condition \textup{(iv)} of that theorem, which is
condition \textup{(ii)} above. Equally, they are equivalent to \(X\)
being smooth over~\(k\) at~\(x\) together with condition
\textup{(i)} of Theorem~\ref{II.4.10}; here the latter simply means
that \(\kappa(x)\) is separable over~\(k\). This implies
condition \textup{(i)} above. To prove the converse it remains only to
establish the following corollary.

\begin{corollary}
\label{II.5.10}
Let \(x\) be a point of a prescheme \(X\) of finite type over the
field~\(k\), and suppose that \(\kappa(x)\) is separable over~\(k\).
Then \(X\) is smooth over~\(k\) at~\(x\) if and only if it is regular
at~\(x\), that is, if and only if the local ring \(\mathcal O_x\) is
regular.
\end{corollary}

Indeed, if \(X\) is regular at~\(x\), one can plainly find elements
\(f_i\in\mathcal O_x\), \(1\leq i\leq n\), satisfying
condition~\ref{II.5.9}\textup{(iv)}.

\end{filecontents*}
% Original input: drafts/SGA1_II_5_OPENING_THROUGH_5_5_English.texfrag
\begin{filecontents*}[overwrite,noheader]{sga1-cumulative-input-021.tex}
% SGA 1, Exposé II, French TeX lines 4048--4189
% (section II.5 opening through Theorem II.5.5 and its proof).

\section{Case of a base field}
\label{II.5}

\marginpar{\raggedright\footnotesize French source p.~52}%
\begin{proposition}
\label{II.5.1}
Let \(k\) be a field, let \(X\) be a prescheme of finite type over~\(k\),
let \(x\in X\), and let \(n\) be the dimension of~\(X\) at~\(x\). Let
\[
  f\colon X\longrightarrow
  \operatorname{Spec}\bigl(k[t_1,\dots,t_n]\bigr)=Y
\]
be the morphism defined by elements
\(f_i\in\Gamma(X,\mathcal O_X)\). The following conditions are equivalent
(and imply that \(X\) is smooth over~\(k\) at~\(x\), and hence, by
Proposition~\ref{II.3.1}, regular at~\(x\)):
\begin{enumerate}
\item[\textup{(i)}] \(f\) is étale at~\(x\);
\item[\textup{(ii)}] the \(\mathrm df_i\) form a basis of
  \(\Omega^1_{X/k}\) at~\(x\);
\item[\textup{(iii)}] the \(\mathrm df_i\) generate
  \(\Omega^1_{X/k}\) at~\(x\).
\end{enumerate}
\end{proposition}

Since \textup{(i)} implies that \(X\) is smooth over~\(k\) at~\(x\), the
implication \textup{(i)}\(\Rightarrow\)\textup{(ii)} is a special case of
Corollary~\ref{II.4.8}, and
\textup{(ii)}\(\Rightarrow\)\textup{(iii)} is immediate. It remains to
prove \textup{(iii)}\(\Rightarrow\)\textup{(i)}. If \textup{(iii)} holds,
Lemma~\ref{II.4.1} shows that \(f\) is unramified at~\(x\). After replacing
\(X\) by an open neighborhood of~\(x\), it is therefore quasi-finite, and
hence dominant by dimension considerations. Since \(Y\) is regular, it
follows from \ifthenelse{\boolean{orig}}{I~}{}\ref{I.9.5}\textup{(ii)}, or from \ifthenelse{\boolean{orig}}{I~}{}\ref{I.9.11}, that \(f\)
is étale.

\begin{corollary}
\label{II.5.2}
Under the hypotheses preceding the equivalent conditions in
Proposition~\ref{II.5.1}, suppose that \(\kappa(x)\) is a
\emph{finite separable} extension of~\(k\), and that
\(f_i\in\mathfrak m_x\) for \(1\leq i\leq n\). Then the preceding
conditions are also equivalent to:
\begin{enumerate}
\item[\textup{(iv)}] the \(f_i\) form a system of generators of
  \(\mathfrak m_x\); equivalently, their classes modulo
  \(\mathfrak m_x^2\) form a basis of
  \(\mathfrak m_x/\mathfrak m_x^2\) over~\(\kappa(x)\).
\end{enumerate}
\end{corollary}

Indeed, \textup{(iv)} implies \textup{(iii)} by the exact sequence
\begin{equation}
\label{eq:II.5.1}
  \mathfrak m_x/\mathfrak m_x^2
  \longrightarrow \Omega^1_{\mathcal O_x/k}
  \longrightarrow \Omega^1_{\kappa(x)/k}
  \longrightarrow 0,
\end{equation}
since \(\Omega^1_{\kappa(x)/k}=0\), because \(\kappa(x)\) is étale
over~\(k\). Conversely, \textup{(ii)} implies \textup{(iv)}: both \(X\)
and \(\operatorname{Spec}(\kappa(x))\) are smooth over~\(k\) at~\(x\), so
by Theorem~\ref{II.4.10}\textup{(iv)} the preceding exact sequence
may be preceded by~\(0\).

\begin{corollary}
\label{II.5.3}
Let \(x\) be a point of a prescheme \(X\) of finite type over~\(k\). If
\(X\) is smooth over~\(k\) at~\(x\), then \(\mathcal O_x\) is regular.
Conversely, if \(\kappa(x)\) is a finite separable extension of~\(k\) and
\(\mathcal O_x\) is regular, then \(X\) is smooth over~\(k\) at~\(x\).
\end{corollary}

\marginpar{\raggedright\footnotesize French source p.~53}%
Indeed, the converse follows from Corollary~\ref{II.5.2}, by taking a
regular system \((f_i)\) of generators of~\(\mathfrak m_x\). (Instead of
Corollary~\ref{II.5.2}, one may also invoke Theorem~\ref{II.4.15}.) We
conclude:

\begin{proposition}
\label{II.5.4}
Let \(X\) be a prescheme of finite type over~\(k\). If \(X\) is smooth
over~\(k\), then it is regular. Conversely, if \(k\) is perfect and \(X\)
is regular, then \(X\) is smooth over~\(k\).
\end{proposition}

For the converse, Corollary~\ref{II.5.3} shows that \(X\) is smooth
over~\(k\) at every closed point, hence everywhere, since the locus on
which it is smooth is open.

\begin{theorem}
\label{II.5.5}
Let \(X\) be a prescheme of finite type over~\(k\), let \(x\in X\), let
\(n\) be the dimension of~\(X\) at~\(x\), and let \(k'\) be a perfect
field extending~\(k\). The following conditions are equivalent:
\begin{enumerate}
\item[\textup{(i)}] \(X\) is smooth over~\(k\) at~\(x\);
\item[\textup{(ii)}] \(\Omega^1_{X/k}\) is locally free of rank \(n\)
  at~\(x\);
\item[\textup{(ii bis)}] \(\Omega^1_{X/k}\) is generated by \(n\)
  elements at~\(x\);
\item[\textup{(iii)}] \(X\) is differentially smooth over~\(k\) at~\(x\);
\item[\textup{(iv)}] there exists an open neighborhood \(U\) of~\(x\)
  such that \(U\otimes_k k'\) is regular; that is, all its local rings are
  regular.
\end{enumerate}
\end{theorem}

Theorem~\ref{II.4.3}\textup{(ii)} gives
\textup{(i)}\(\Rightarrow\)\textup{(ii)}, the implication
\textup{(ii)}\(\Rightarrow\)\textup{(ii bis)} is immediate, and
Proposition~\ref{II.5.1} gives
\textup{(ii bis)}\(\Rightarrow\)\textup{(i)}. Since \(X\) is flat
over~\(k\), Corollary~\ref{II.4.18} gives
\textup{(i)}\(\Leftrightarrow\)\textup{(iii)}. Condition \textup{(i)}
implies \textup{(iv)}, because smoothness is preserved under base extension
and implies regularity. Conversely, if \textup{(iv)} holds, then
Proposition~\ref{II.5.4} shows that \(U\otimes_k k'\) is smooth over~\(k'\),
and hence \(U\) is smooth over~\(k\) by Corollary~\ref{II.4.13}.

Taking \(x\) to be the generic point of~\(X\), with \(X\) assumed
irreducible, we obtain:
\end{filecontents*}
% Original input: drafts/SGA1_II_ERRATA_AND_III_HEADING_English.texfrag
\begin{filecontents*}[overwrite,noheader]{sga1-cumulative-input-022.tex}
\subsection*{Errata}
\label{II.fin.errata}

\marginpar{\raggedright\footnotesize French source p.~57}%
In the proof of Corollary~\ref{II.5.6} above, we used the fact that a
nonempty reduced scheme of finite type over an algebraically closed
field has at least one regular, hence smooth, point. This fact is
usually proved by differential methods, using Zariski's theorem that
the set of regular points of~\(X\) is open. To avoid a circular
argument, one must prove the following directly. Let \(K/k\) be a
finitely generated separable extension, and suppose that
\(f_i\in K\), \(1\leq i\leq n\), are such that the
\(\mathrm d_{K/k}f_i\) form a basis of \(\Omega^1_{K/k}\). Then \(n\)
is the transcendence degree of \(K\) over~\(k\); equivalently, the
\(f_i\) are algebraically independent.

The usual proof uses Mac~Lane's criterion; see Bourbaki,
\emph{Algèbre}, Chapter~V, §9, Theorem~2. Choose a nonzero polynomial
\(g\in k[t_1,\dots,t_n]\) of minimal degree such that
\(g(f_1,\dots,f_n)=0\). Then
\[
  \sum_i
  \frac{\partial g}{\partial t_i}(f_1,\dots,f_n)\,
  \mathrm d f_i=0.
\]
Since the \(\mathrm d f_i\) form a basis of \(\Omega^1_{K/k}\), each
\(\partial g/\partial t_i\) vanishes at
\((f_1,\dots,f_n)\), and hence is the zero polynomial by the minimality
of~\(g\). In characteristic zero this gives \(g=0\). In
characteristic \(p\ne0\), one has
\(g=h(t_1^p,\dots,t_n^p)\). Mac~Lane's criterion then shows that
\(h\in k[t_1,\dots,t_n]\) also vanishes at
\((f_1,\dots,f_n)\), and the minimality of~\(g\) again gives \(g=0\).

\chapter{Smooth morphisms: lifting properties}
\label{III}
\marginpar{\raggedright\footnotesize French source p.~58}%

\end{filecontents*}
% Original input: drafts/SGA1_IV_1_English.texfrag
\begin{filecontents*}[overwrite,noheader]{sga1-cumulative-input-023.tex}
\section{Elementary facts about flat modules}
\label{IV.1}

An \(A\)-module \(M\) is called \emph{flat}, or \(A\)-flat when the
ring must be specified, if the functor
\[
  T_M\colon N\longmapsto M\otimes_A N,
\]
which is always right exact, is exact; equivalently, if it takes
monomorphisms to monomorphisms.
\marginpar{\raggedright\footnotesize French source p.~88}%
This is equivalent to the vanishing of the first left-derived functor,
or of all the left-derived functors:
\[
  \operatorname{Tor}_1^A(M,N)=0
  \qquad\text{for every }N,
\]
or, respectively,
\[
  \operatorname{Tor}_i^A(M,N)=0
  \qquad\text{for every }i>0\text{ and every }N.
\]
Since the \(\operatorname{Tor}_i\) commute with filtered direct limits,
it is enough to verify these conditions for finitely generated~\(N\).
Using a filtration of~\(N\) with cyclic quotients, it is enough even to
require
\[
  \operatorname{Tor}_1^A(M,N)=0
\]
when \(N\) is cyclic, hence of the form \(A/I\) for an ideal
\(I\subseteq A\). Moreover,
\[
  \operatorname{Tor}_1^A(M,A/I)=0
  \quad\Longleftrightarrow\quad
  I\otimes_A M\longrightarrow M=A\otimes_A M
  \text{ is injective}.
\]
This follows from the exact \(\operatorname{Tor}\)-sequence, since
\(\operatorname{Tor}_1^A(M,A)=0\). Thus \(M\) is flat if and only if,
for every ideal~\(I\), the natural homomorphism
\[
  I\otimes_A M\longrightarrow IM
\]
is an isomorphism. It is enough to verify this for finitely generated
ideals~\(I\); a fortiori it is enough to verify that
\(M\otimes_A-\) is exact on finitely generated modules.

As for every exact functor \(T\), if \(T(N')\) is identified with a
subobject of \(T(N)\) whenever \(N'\subseteq N\), then
\[
  T(N'\cap N'')=T(N')\cap T(N''),
  \qquad
  T(N'+N'')=T(N')+T(N'')
\]
for any two subobjects \(N'\) and \(N''\) of~\(N\).

A direct sum of flat modules is flat, and a direct summand of a flat
module is flat. Since \(A\) is flat over itself, every free module, and
hence every projective module, is flat. The tensor product of two flat
modules is flat. If \(M\) is flat over~\(A\), then
\(M\otimes_A B\) is flat over~\(B\) after every base change
\(A\to B\). These assertions follow from associativity of the tensor
product and from the fact that a composite of exact functors is exact.
\marginpar{\raggedright\footnotesize French source p.~89}%
Likewise, if \(M\) is flat over~\(B\) and \(B\) is flat over~\(A\),
then \(M\) is flat over~\(A\).

The exact \(\operatorname{Tor}\)-sequence, together with the symmetry
of \(\operatorname{Tor}\), gives:

\begin{proposition}
\label{IV.1.1}
Let
\[
  0\longrightarrow M'\longrightarrow M\longrightarrow M''
  \longrightarrow0
\]
be an exact sequence of \(A\)-modules, with \(M''\) flat. Then:
\begin{enumerate}
\item[(i)] the sequence remains exact after tensoring with any
\(A\)-module \(N\);
\item[(ii)] \(M\) is flat if and only if \(M'\) is flat.
\end{enumerate}
\end{proposition}

Thus, from the viewpoint of tensor products, flat modules are as well
behaved as free or projective modules. In particular, the exact
sequence in Proposition~\ref{IV.1.1} behaves as well as if it split.

Let \(S\) be a multiplicatively closed subset of~\(A\). Then
\(S^{-1}A\) is flat over~\(A\), since
\[
  S^{-1}A\otimes_A N=S^{-1}N
\]
is an exact functor of~\(N\). If \(M\) is \(A\)-flat, then
\(S^{-1}M=S^{-1}A\otimes_AM\) is \(S^{-1}A\)-flat. Conversely, if
\(M\to S^{-1}M\) is an isomorphism, equivalently if every \(s\in S\)
acts bijectively on~\(M\), then \(S^{-1}A\)-flatness of \(M\) implies
\(A\)-flatness by transitivity.

More generally, the study of the flatness over~\(Y\) of a
quasi-coherent sheaf \(\mathcal F\) on a prescheme \(X\to Y\) leads to
the following two-ring situation.

\begin{proposition}
\label{IV.1.2}
Let \(A\to B\) be a ring homomorphism, let \(M\) be a \(B\)-module, and
let \(T\) be a multiplicatively closed subset of~\(B\).
\begin{enumerate}
\item[(i)] If \(M\) is \(A\)-flat, then \(T^{-1}M\) is \(A\)-flat.
Consequently, it is also \(S^{-1}A\)-flat for every multiplicatively
closed subset \(S\subseteq A\) whose image lies in~\(T\).
\item[(ii)] Conversely, suppose that for every maximal ideal
\(\mathfrak n\) of~\(B\), \(M_{\mathfrak n}\) is flat over~\(A\), or
equivalently over \(A_{\mathfrak m}\), where
\(\mathfrak m\) is the inverse image of~\(\mathfrak n\) in~\(A\).
Then \(M\) is \(A\)-flat.
\end{enumerate}
\end{proposition}

Indeed, functorially in the \(A\)-module~\(N\),
\[
  T^{-1}M\otimes_A N
  =
  T^{-1}(M\otimes_A N).
\]
Both sides are naturally isomorphic, by associativity, to
\[
  T^{-1}B\otimes_B M\otimes_B(N\otimes_A B).
\]
Thus exactness of \(M\otimes_A N\) in~\(N\) implies that of
\(T^{-1}M\otimes_A N\), proving \textup{(i)}. It also proves
\textup{(ii)}, since exactness of a sequence of \(B\)-modules may be
checked after localization at every maximal ideal of~\(B\).

\begin{proposition}
\label{IV.1.3}
\begin{enumerate}
\item[(i)]
\marginpar{\raggedright\footnotesize French source p.~90}%
Let \(M\) be a flat \(A\)-module. If \(x\in A\) is not a zero divisor
in~\(A\), then it is not a zero divisor on~\(M\). In particular, if
\(A\) is an integral domain, then \(M\) is torsion-free.
\item[(ii)] Suppose that \(A\) is an integral domain and that
\(A_{\mathfrak m}\) is a principal ideal domain for every maximal
ideal \(\mathfrak m\) of~\(A\), for example when \(A\) is a Dedekind
domain or a principal ideal domain. An \(A\)-module \(M\) is flat if
and only if it is torsion-free.
\end{enumerate}
\end{proposition}

For \textup{(i)}, multiplication by~\(x\) on~\(M\) is obtained by
tensoring multiplication by~\(x\) on~\(A\) with~\(M\). For
\textup{(ii)}, Proposition~\ref{IV.1.2}\textup{(ii)} reduces us to the
case where \(A\) is a principal ideal domain. If \(M\) is torsion-free,
then for every ideal \(I=(x)\), tensoring the injection \(I\to A\) with
\(M\) remains injective, since \(x\) is not a zero divisor on~\(M\).
This is precisely the flatness criterion above.

\end{filecontents*}
% Original input: drafts/SGA1_IV_2_English.texfrag
\begin{filecontents*}[overwrite,noheader]{sga1-cumulative-input-024.tex}
\section{Faithfully flat modules}
\label{IV.2}

A functor \(F\) from one category to another is called \emph{faithful}
if, for every pair of objects \(X,Y\), the map
\[
  \operatorname{Hom}(X,Y)
  \longrightarrow
  \operatorname{Hom}(F(X),F(Y))
\]
is injective. For an additive functor between additive categories, this
is equivalent to saying that \(F(u)=0\) implies \(u=0\), and it implies
that \(F(X)=0\) only if \(X=0\).

For a functor between abelian categories to be faithful and exact, it
is necessary and sufficient that a sequence
\[
  M'\longrightarrow M\longrightarrow M''
\]
be exact if and only if its image
\[
  F(M')\longrightarrow F(M)\longrightarrow F(M'')
\]
is exact. Equivalently, \(F\) must be exact and \(F(X)=0\) must imply
\(X=0\).

Suppose that a family \((M_i)\) of nonzero objects has the property
that every nonzero object has a subobject admitting a quotient
isomorphic to some~\(M_i\). Then \(F\) is faithful and exact if and
only if \(F\) is exact and \(F(M_i)\ne0\) for every~\(i\). In the
category of modules over a ring~\(A\), one may take the family of
modules \(A/\mathfrak m\), where \(\mathfrak m\) runs through the
maximal ideals of~\(A\). Indeed, every nonzero module has a nonzero
cyclic submodule, isomorphic to \(A/I\) for some proper ideal~\(I\),
and by Krull's theorem this has a quotient \(A/\mathfrak m\). Hence:

\begin{proposition}
\label{IV.2.1}
For an \(A\)-module \(M\), the following conditions are equivalent:
\begin{enumerate}
\item[(i)] The functor \(M\otimes_A-\) is faithful and exact.
\item[(i bis)]
\marginpar{\raggedright\footnotesize French source p.~91}%
\(M\) is flat, and \(M\otimes_A N=0\) implies \(N=0\).
\item[(i ter)] \(M\) is flat, and
\[
  M\otimes_AA/\mathfrak m\ne0
\]
for every maximal ideal \(\mathfrak m\) of~\(A\).
\item[(ii)] For every sequence of homomorphisms
\(N'\to N\to N''\), the sequence obtained by tensoring with~\(M\) is
exact if and only if the original sequence is exact.
\end{enumerate}
\end{proposition}

In this case \(M\) is called a \emph{faithfully flat \(A\)-module}. In
particular, if \(M\) is faithfully flat, a homomorphism
\(N\to N'\) is a monomorphism, epimorphism, or isomorphism if and only
if the homomorphism obtained by tensoring with~\(M\) has the same
property. A faithfully flat module is faithful, since multiplication
by \(f\) on~\(M\) is obtained by tensoring multiplication by \(f\)
on~\(A\) with~\(M\).

The usual transitivity properties follow as in Section~\ref{IV.1}. The
tensor product of two faithfully flat modules is faithfully flat. If
\(M\) is faithfully flat over~\(A\), then \(M\otimes_AB\) is
faithfully flat over~\(B\) after every base extension \(A\to B\). If
\(B\) is a faithfully flat \(A\)-algebra and \(M\) is faithfully flat
as a \(B\)-module, then \(M\) is faithfully flat over~\(A\).

\begin{corollary}
\label{IV.2.2}
Let \(A\to B\) be a local homomorphism of local rings, and let \(M\) be
a finitely generated \(B\)-module. Then \(M\) is faithfully flat
over~\(A\) if and only if it is flat over~\(A\) and nonzero.
\end{corollary}

This follows from criterion \textup{(i ter)} and Nakayama's lemma. In
particular, \(B\) is flat over~\(A\) if and only if it is faithfully
flat over~\(A\).

\begin{proposition}
\label{IV.2.3}
Let \(A\to B\) be a ring homomorphism, and let \(M\) be a \(B\)-module
which is faithfully flat over~\(A\). For every prime ideal
\(\mathfrak p\) of~\(A\), there exists a prime ideal \(\mathfrak q\)
of~\(B\) whose contraction to~\(A\) is~\(\mathfrak p\).
\end{proposition}

After quotienting by~\(\mathfrak p\), we reduce to
\(\mathfrak p=0\). Localizing at the prime ideal~\(0\), we reduce to
the case where \(A\) is a field. Since \(M\) is faithfully flat
over~\(A\), it is nonzero, and hence \(B\ne0\). The ring \(B\) has a
prime ideal, which necessarily contracts to the unique prime ideal
of~\(A\). Geometrically, the existence on
\(X=\operatorname{Spec}(B)\) of a quasi-coherent sheaf which is
faithfully flat relative to~\(A\) implies that
\[
  X\longrightarrow Y=\operatorname{Spec}(A)
\]
is surjective.

\begin{corollary}
\label{IV.2.4}
Suppose that \(M\) is flat over~\(A\), finitely generated over~\(B\),
and has full support:
\[
  \operatorname{Supp}M=\operatorname{Spec}(B).
\]
For every prime ideal \(\mathfrak p\) of~\(A\), each prime ideal
\(\mathfrak q\) of~\(B\) minimal among those containing
\(\mathfrak pB\) contracts to~\(\mathfrak p\).
\end{corollary}

\marginpar{\raggedright\footnotesize French source p.~92}%
All the hypotheses are preserved after quotienting, so once again we
reduce to \(\mathfrak p=0\), with \(A\) an integral domain. We are
reduced to:

\begin{corollary}
\label{IV.2.5}
Let \(M\) be as above. Every minimal prime \(\mathfrak q\) of~\(B\)
contracts to a minimal prime \(\mathfrak p\) of~\(A\).
\end{corollary}

Indeed, after localizing at \(\mathfrak p\) and~\(\mathfrak q\), it is
enough to prove the following. Let \(A\) and \(B\) be local rings, let
\(A\to B\) be local, let \(M\) be a nonzero \(B\)-module flat
over~\(A\), and suppose that \(B\) has dimension zero. Then \(A\) has
dimension zero. By Corollary~\ref{IV.2.2} and
Proposition~\ref{IV.2.3}, every prime ideal of~\(A\) is induced by a
prime ideal of~\(B\), hence by the maximal ideal of~\(B\), and is
therefore the maximal ideal of~\(A\).

Geometrically, Corollary~\ref{IV.2.5} says that every irreducible
component of \(X=\operatorname{Spec}(B)\) dominates an irreducible
component of \(Y=\operatorname{Spec}(A)\), provided there is a
finitely generated quasi-coherent sheaf on~\(X\), with support~\(X\),
which is flat relative to~\(Y\).

Notice that Corollary~\ref{IV.2.4} did not require \(M\) to be
faithfully flat over~\(A\). In that case, however, there is no
guarantee that a prime ideal containing \(\mathfrak pB\) exists.

\begin{proposition}
\label{IV.2.6}
Let \(i\colon A\to B\) be a ring homomorphism. The following conditions
are equivalent:
\begin{enumerate}
\item[(i)] \(B\) is faithfully flat as an \(A\)-module.
\item[(ii)] \(B\) is flat over~\(A\), and
\(\operatorname{Spec}(B)\to\operatorname{Spec}(A)\) is surjective.
\item[(ii bis)] \(B\) is flat over~\(A\), and every maximal ideal
of~\(A\) is the contraction of a prime ideal of~\(B\).
\item[(iii)] \(i\) is injective and \(\operatorname{Coker}i\) is a
flat \(A\)-module.
\item[(iv)] The functor
\[
  M\longmapsto M_{(B)}=M\otimes_AB
\]
is exact, and the canonical natural homomorphism
\(M\to M_{(B)}\) is injective.
\item[(iv bis)] For every ideal \(I\subseteq A\), the map
\[
  I\otimes_AB\longrightarrow IB
\]
is an isomorphism, and the inverse image of \(IB\) in~\(A\) is~\(I\).
\end{enumerate}
\end{proposition}

The implication \(\textup{(i)}\Rightarrow\textup{(ii)}\) follows from
Proposition~\ref{IV.2.3};
\(\textup{(ii)}\Rightarrow\textup{(ii bis)}\) is immediate; and
\(\textup{(ii bis)}\Rightarrow\textup{(i)}\) follows from criterion
\textup{(i ter)} of Proposition~\ref{IV.2.1}. Proposition~\ref{IV.1.1}
gives \(\textup{(iii)}\Rightarrow\textup{(iv)}\), and
\(\textup{(iv)}\Rightarrow\textup{(iv bis)}\) is immediate.
Condition \textup{(iv bis)} implies \textup{(i)} by the ideal
criterion for flatness in Section~\ref{IV.1} and criterion
\textup{(i ter)} of Proposition~\ref{IV.2.1}. The implication
\(\textup{(iv)}\Rightarrow\textup{(iii)}\) follows from the converse
form of Proposition~\ref{IV.1.1}.

Finally, \textup{(i)} implies \textup{(iv)} as follows. If \(N\) is
the kernel of
\[
  M\longrightarrow M\otimes_AB=T(M),
\]
then exactness of~\(T\) makes the map \(N\to T(N)\) zero. Hence
\(T(N)=N\otimes_AB=0\), and faithful flatness gives \(N=0\).
\marginpar{\raggedright\footnotesize French source p.~93}%

\end{filecontents*}
% Original input: drafts/SGA1_IV_3_AND_4_English.texfrag
\begin{filecontents*}[overwrite,noheader]{sga1-cumulative-input-025.tex}
\section{Relations with completion}
\label{IV.3}

Let \(A\) be a Noetherian ring, let \(I\) be an ideal of~\(A\), and
let \(\widehat A\) be the separated completion of~\(A\) for the
\(I\)-adic topology. For every \(A\)-module \(M\), let
\(\widehat M\) be its \(I\)-adic completion. It is an
\(\widehat A\)-module, and hence there is a canonical homomorphism
\[
  M\otimes_A\widehat A\longrightarrow\widehat M.
\]

On finitely generated modules, the functor \(M\mapsto\widehat M\) is
exact. This follows readily from Krull's theorem: if \(N\subseteq M\),
the \(I\)-adic topology of~\(N\) is the topology induced from~\(M\).
Since \(M\mapsto M\otimes_A\widehat A\) is right exact, a finite free
presentation \(L\to L'\to M\) shows that the displayed natural
homomorphism is an isomorphism, completion also being right exact.
Consequently \(M\mapsto M\otimes_A\widehat A\) is exact on finitely
generated modules. Hence:

\begin{proposition}
\label{IV.3.1}
Let \(A\) be a Noetherian ring and \(I\) an ideal of~\(A\). The
separated \(I\)-adic completion \(\widehat A\) of~\(A\) is flat
over~\(A\).
\end{proposition}

\begin{corollary}
\label{IV.3.2}
The ring \(\widehat A\) is faithfully flat over~\(A\) if and only if
\(I\) is contained in the Jacobson radical of~\(A\).
\end{corollary}

Apply criterion \textup{(i ter)} of Proposition~\ref{IV.2.1}.

These statements summarize everything that linear algebra alone says
about the relation between \(A\) and~\(\widehat A\).
Corollary~\ref{IV.3.2} is used especially when \(A\) is a Noetherian
local ring and \(I\) is contained in its maximal ideal
\(\mathfrak m\), most often with \(I=\mathfrak m\).

\section{Relations with free modules}
\label{IV.4}

\begin{proposition}
\label{IV.4.1}
Let \(A\) be a ring, \(I\) an ideal of~\(A\), and \(M\) an
\(A\)-module. Suppose that one of the following hypotheses holds:
\begin{enumerate}
\item[(a)] \(I\) is nilpotent;
\item[(b)] \(A\) is Noetherian, \(I\) is contained in the Jacobson
radical of~\(A\), and \(M\) is finitely generated.
\end{enumerate}
Then \(M\) is free over~\(A\) if and only if \(M\otimes_AA/I\) is free
over~\(A/I\) and
\[
  \operatorname{Tor}_1^A(M,A/I)=0.
\]
\end{proposition}

Necessity is clear; we prove sufficiency.
\marginpar{\raggedright\footnotesize French source p.~94}%
Choose elements \((e_i)\) of~\(M\) whose images form a basis of
\[
  M\otimes_AA/I=M/IM
\]
over~\(A/I\); the family is finite in case~\textup{(b)}. Let \(L\) be
the free \(A\)-module on the same index set. The resulting homomorphism
\(L\to M\) becomes an isomorphism after tensoring with~\(A/I\). If
\(Q\) is its cokernel, then \(Q/IQ=0\), and Nakayama's lemma, valid
under either hypothesis, gives \(Q=0\). Thus \(L\to M\) is surjective.
Let \(R\) be its kernel:
\[
  0\longrightarrow R\longrightarrow L\longrightarrow M
  \longrightarrow0.
\]
Since \(\operatorname{Tor}_1^A(M,A/I)=0\), tensoring with~\(A/I\)
gives an exact sequence
\[
  0\longrightarrow R/IR\longrightarrow L/IL\longrightarrow M/IM
  \longrightarrow0.
\]
Hence \(R/IR=0\), and Nakayama again gives \(R=0\). In case
\textup{(b)}, note that \(R\) is finitely generated because \(A\) is
Noetherian.

\begin{corollary}
\label{IV.4.2}
The condition \(\operatorname{Tor}_1^A(M,A/I)=0\) may be replaced by
the condition that the canonical surjection
\begin{equation}
\label{eq:IV.4.*}
  \operatorname{gr}_I^0(M)
  \otimes_{A/I}\operatorname{gr}_I(A)
  \longrightarrow\operatorname{gr}_I(M)
\tag{*}
\end{equation}
is an isomorphism.
\end{corollary}

The condition plainly holds when \(M\) is free. Conversely, suppose
that \(M\otimes_AA/I\) is free over~\(A/I\) and that
\eqref{eq:IV.4.*} is an isomorphism. Construct \(L\to M\) as in the
proof above. The hypothesis shows that this homomorphism induces an
isomorphism on associated graded modules. Its kernel is therefore
contained in \(\bigcap_n I^nL\), which is zero: trivially in case
\textup{(a)}, and by the standard separation theorem in case
\textup{(b)}.

\begin{corollary}
\label{IV.4.3}
Suppose that \(A/I\) is a field. The following conditions on~\(M\) are
equivalent:
\begin{enumerate}
\item[(i)] \(M\) is free;
\item[(ii)] \(M\) is projective;
\item[(iii)] \(M\) is flat;
\item[(iv)] \(\operatorname{Tor}_1^A(M,A/I)=0\);
\item[(v)] the canonical homomorphism \eqref{eq:IV.4.*} is an
isomorphism.
\end{enumerate}
\end{corollary}

Indeed, \(M\otimes_AA/I\) is automatically free in this case. The
corollary applies in either of the following situations:
\begin{enumerate}
\item[(a)] \(M\) is an arbitrary module over a local ring~\(A\) whose
maximal ideal \(I\) is nilpotent, for example an Artinian local ring;
\item[(b)] \(M\) is a finitely generated module over a Noetherian local
ring~\(A\).
\end{enumerate}

For reference, we recall:

\begin{corollary}
\label{IV.4.4}
\marginpar{\raggedright\footnotesize French source p.~95}%
Let \(A\) be a Noetherian local integral domain with maximal ideal
\(\mathfrak m=I\), residue field \(k=A/I\), and fraction field~\(K\).
Let \(M\) be a finitely generated \(A\)-module. Conditions
\textup{(i)}--\textup{(v)} above are also equivalent to:
\begin{enumerate}
\item[(vi)] The vector spaces \(M\otimes_AK\) and \(M\otimes_Ak\)
have the same dimension; equivalently, the rank of~\(M\) over~\(A\)
equals the minimum number of generators of the \(A\)-module~\(M\).
\end{enumerate}
\end{corollary}

The proof is immediate. The reader may generalize it to the case where
\(A\) is only assumed reduced: one must then require the ranks of~\(M\)
at all minimal prime ideals of~\(A\) to equal the dimension of the
vector space \(M\otimes_Ak\).

\end{filecontents*}
% Original input: drafts/SGA1_IV_5_English.texfrag
\begin{filecontents*}[overwrite,noheader]{sga1-cumulative-input-026.tex}
% SGA 1, Exposé IV, French TeX lines 6829--7127 (section IV.5).

\section{Local criteria for flatness}
\label{IV.5}

\begin{proposition}
\label{IV.5.1}
Let \(A\) be a ring equipped with an ideal~\(I\), and let \(M\) be an
\(A\)-module. Suppose that
\[
  \operatorname{Tor}_1^A(M,A/I^n)=0
  \qquad\text{for every \(n>0\)}.
\]
Then the canonical surjective homomorphism
\begin{equation}
\label{eq:IV.5.1.*}
  \operatorname{gr}_I^0(M)
  \otimes_{A/I}\operatorname{gr}_I(A)
  \longrightarrow\operatorname{gr}_I(M)
\tag{*}
\end{equation}
is an isomorphism. The converse holds if \(I\) is nilpotent.
\end{proposition}

The hypothesis means that the homomorphisms
\[
  I^n\otimes_AM\longrightarrow I^nM
\]
are isomorphisms. It follows immediately that the homomorphisms
\[
  (I^n/I^{n+1})\otimes_AM
  \longrightarrow I^nM/I^{n+1}M
\]
are isomorphisms. Conversely, suppose this last condition holds and that
\(I\) is nilpotent. We prove that
\(\operatorname{Tor}_1^A(M,A/I^n)=0\) for every~\(n\). This is true for
large~\(n\); proceed by descending induction, assuming it for~\(n+1\).
There is a commutative diagram
\[
\begin{array}{ccccccccc}
 && M\otimes_AI^{n+1} &\longrightarrow& M\otimes_AI^n
 &\longrightarrow& M\otimes_A(I^n/I^{n+1}) &\longrightarrow& 0
\\
 && \Big\downarrow && \Big\downarrow && \Big\downarrow &&
\\
0 &\longrightarrow& MI^{n+1} &\longrightarrow& MI^n
 &\longrightarrow& MI^n/MI^{n+1} &\longrightarrow& 0,
\end{array}
\]
whose rows are exact.
\marginpar{\raggedright\footnotesize French source p.~96}%
By hypothesis, the last vertical arrow is an isomorphism, while the
induction hypothesis says that the first vertical arrow is an isomorphism.
The middle vertical arrow is therefore also an isomorphism, completing the
proof.

The following proposition was isolated by Serre during the seminar; it
allows substantial simplifications in this section.

\begin{proposition}
\label{IV.5.2}
Let \(A\to B\) be a ring homomorphism and let \(M\) be an \(A\)-module.
The following conditions are equivalent:
\begin{enumerate}
\item[\textup{(i)}] for every \(B\)-module \(N\),
  \(\operatorname{Tor}_1^A(M,N)=0\);
\item[\textup{(ii)}] \(\operatorname{Tor}_1^A(M,B)=0\), and
  \(M_{(B)}=M\otimes_AB\) is flat over~\(B\).
\end{enumerate}
\end{proposition}

There is a functorial isomorphism
\[
  M\otimes_AN=(M\otimes_AB)\otimes_BN.
\]
Regarding the left-hand side as a functor of~\(M\), this expresses it as
the composite of the two functors
\[
  M\longmapsto M\otimes_AB,
  \qquad
  P\longmapsto P\otimes_BN.
\]
The first carries free \(A\)-modules to free \(B\)-modules, and hence
projective modules to projective modules. The spectral sequence for
composite functors therefore gives
\[
  E^2_{p,q}
  =\operatorname{Tor}_p^B
    \bigl(\operatorname{Tor}_q^A(M,B),N\bigr)
  \Longrightarrow
  \operatorname{Tor}_{p+q}^A(M,N),
\]
and its low-degree terms give an exact sequence
\[
  \operatorname{Tor}_1^A(M,B)\otimes_BN
  \longrightarrow \operatorname{Tor}_1^A(M,N)
  \longrightarrow \operatorname{Tor}_1^B(M\otimes_AB,N)
  \longrightarrow 0.
\]
If \textup{(i)} holds, then
\(\operatorname{Tor}_1^A(M,B)=0\), and the exact sequence shows that
\(\operatorname{Tor}_1^B(M\otimes_AB,N)=0\) for every~\(N\). Thus
\(M\otimes_AB\) is flat over~\(B\), proving \textup{(ii)}. Conversely,
if \textup{(ii)} holds, the two terms surrounding
\(\operatorname{Tor}_1^A(M,N)\) vanish, proving \textup{(i)}.

\begin{corollary}
\label{IV.5.3}
Suppose that \(B=A/I\). Then the preceding conditions are also equivalent
to:
\begin{enumerate}
\item[\textup{(iii)}] \(\operatorname{Tor}_1^A(M,N)=0\) for every
  \(A\)-module \(N\) annihilated by a power of~\(I\).
\end{enumerate}
\end{corollary}

Indeed, condition~\textup{(i)} says this when \(N\) is annihilated by
\(I\). Condition~\textup{(iii)} follows by applying this assertion
successively to the modules \(I^nN/I^{n+1}N\).

\begin{corollary}
\label{IV.5.4}
Under the hypotheses of Corollary~\ref{IV.5.3}, the conditions under
consideration imply that the functorial homomorphism
\marginpar{\raggedright\footnotesize French source p.~97}%
\begin{equation}
\label{eq:IV.5.*}
  \operatorname{gr}_I^0(M)
  \otimes_{A/I}\operatorname{gr}_I(A)
  \longrightarrow\operatorname{gr}_I(M)
\tag{*}
\end{equation}
is an isomorphism, and that \(M\otimes_AA/I\) is flat over~\(A/I\).
\end{corollary}

Apply condition~\textup{(iii)} and Proposition~\ref{IV.5.1}. Using the
converse in Proposition~\ref{IV.5.1} when \(I\) is nilpotent gives:

\begin{corollary}
\label{IV.5.5}
Let \(A\) be a ring equipped with a nilpotent ideal~\(I\), and let \(M\)
be an \(A\)-module. The following conditions are equivalent:
\begin{enumerate}
\item[\textup{(i)}] \(M\) is flat over~\(A\);
\item[\textup{(ii)}] \(M\otimes_AA/I\) is flat over~\(A/I\), and
  \(\operatorname{Tor}_1^A(M,A/I)=0\);
\item[\textup{(iii)}] \(M\otimes_AA/I\) is flat over~\(A/I\), and the
  canonical homomorphism \eqref{eq:IV.5.*} on associated graded modules
  is an isomorphism.
\end{enumerate}
\end{corollary}

These are, respectively, condition~\textup{(iii)} and
condition~\textup{(ii)} above, and the conditions of
Corollary~\ref{IV.5.4}.

Now drop the assumption that \(I\) is nilpotent. A priori, in
Corollary~\ref{IV.5.5} we retain only the implications
\[
  \textup{(i)}\Rightarrow\textup{(ii)}\Rightarrow\textup{(iii)}.
\]
Since condition~\textup{(iii)} remains stable after quotienting by a power
of~\(I\), Corollary~\ref{IV.5.5} shows that it implies:
\begin{enumerate}
\item[\textup{(iv)}] \emph{for every \(n>0\),
  \(M\otimes_AA/I^n\) is flat over~\(A/I^n\).}
\end{enumerate}

We seek conditions under which \textup{(iv)} implies \textup{(i)}, that
is, under which \(M\) is flat over~\(A\). It is enough that \(A\) be
Noetherian and that \(M\) satisfy the following finiteness condition:
\emph{for every finitely generated \(A\)-module \(N\), the module
\(M\otimes_AN\) is separated for the \(I\)-adic topology}. It would be
enough to check this when \(N\) is a finitely generated ideal of~\(A\).

Indeed, suppose these conditions hold, and let \(N'\to N\) be an
injective homomorphism of finitely generated modules. We prove that
\[
  M\otimes_AN'\longrightarrow M\otimes_AN
\]
is injective. It is enough to show that its kernel is contained in every
\[
  I^n(M\otimes_AN')
  =\operatorname{Im}
   \bigl(M\otimes_AI^nN'\longrightarrow M\otimes_AN'\bigr),
\]
or, equivalently, in every
\[
  \operatorname{Im}
   \bigl(M\otimes_AV'_n\longrightarrow M\otimes_AN'\bigr)
  =\operatorname{Ker}
   \bigl(M\otimes_AN'\longrightarrow M\otimes_A(N'/V'_n)\bigr),
\]
where \((V'_n)\) runs through a countable fundamental system of
neighborhoods of zero in~\(N'\), equipped with its \(I\)-adic topology.
By Krull's theorem, the \(I\)-adic topology of~\(N'\) is induced by that
of~\(N\), so we may take
\[
  V'_n=N'\cap I^nN.
\]
Consider the commutative diagram
\[
\begin{array}{ccc}
M\otimes_AN' &\longrightarrow& M\otimes_A(N'/V'_n)
\\
\Big\downarrow && \Big\downarrow
\\
M\otimes_AN &\longrightarrow& M\otimes_A(N/I^nN).
\end{array}
\]
\marginpar{\raggedright\footnotesize French source p.~98}%
Since \(N'/V'_n\) and \(N/I^nN\) are annihilated by~\(I^n\), the
right-hand vertical homomorphism is obtained from the \emph{injective}
homomorphism
\[
  N'/V'_n\longrightarrow N/I^nN
\]
by tensoring over~\(A/I^n\) with the \emph{flat}
\((A/I^n)\)-module \(M\otimes_AA/I^n\). It is therefore injective. Hence
the kernel of \(M\otimes_AN'\to M\otimes_AN\) is contained in the kernel
of
\[
  M\otimes_AN'\longrightarrow M\otimes_A(N'/V'_n),
\]
as required.

The finiteness condition just imposed on~\(M\) holds in particular if
\(M\) is a finitely generated module over a Noetherian \(A\)-algebra
\(B\) such that \(IB\) is contained in the Jacobson radical of~\(B\).
Indeed, for every finitely generated \(A\)-module \(N\), the module
\(M\otimes_AN\) is finitely generated over~\(B\), and is therefore
separated by Krull's theorem for the \(I\)-adic topology, which is its
\((IB)\)-adic topology. We have proved:

\begin{theorem}
\label{IV.5.6}
Let \(A\to B\) be a homomorphism of Noetherian rings, let \(I\) be an
ideal of~\(A\) such that \(IB\) is contained in the Jacobson radical
of~\(B\), and let \(M\) be a finitely generated \(B\)-module. The
following conditions are equivalent:
\begin{enumerate}
\item[\textup{(i)}] \(M\) is flat over~\(A\);
\item[\textup{(ii)}] \(M\otimes_AA/I\) is flat over~\(A/I\), and
  \(\operatorname{Tor}_1^A(M,A/I)=0\);
\item[\textup{(iii)}] \(M\otimes_AA/I\) is flat over~\(A/I\), and the
  canonical homomorphism
  \[
    \operatorname{gr}_I^0(M)
    \otimes_{A/I}\operatorname{gr}_I(A)
    \longrightarrow\operatorname{gr}_I(M)
  \]
  is an isomorphism;
\item[\textup{(iv)}] for every \(n>0\), \(M\otimes_AA/I^n\) is flat
  over~\(A/I^n\).
\end{enumerate}
\end{theorem}

This result is used chiefly when \(A\) and~\(B\) are Noetherian
\emph{local} rings, \(A\to B\) is a local homomorphism, and \(I\) is an
ideal of~\(A\) contained in its maximal ideal; Theorem~\ref{IV.5.6} may
in fact be reduced immediately to this case. An interesting special case
is that in which \(A/I\) is a field, that is, \(I\) is maximal. The
condition that \(M\otimes_AA/I\) be flat over~\(A/I\) is then automatic.
Moreover, the rings \(A/I^n\) are Artinian local rings, so
condition~\textup{(iv)} says that the modules \(M\otimes_AA/I^n\) are
\emph{free} over~\(A/I^n\).

\begin{corollary}
\label{IV.5.7}
\marginpar{\raggedright\footnotesize French source p.~99}%
Let \(A\to B\) be a local homomorphism of Noetherian local rings, and let
\(u\colon M'\to M\) be a homomorphism of finitely generated
\(B\)-modules. Suppose that \(M\) is flat over~\(A\). The following
conditions are equivalent:
\begin{enumerate}
\item[\textup{(i)}] \(u\) is injective and \(\operatorname{Coker}(u)\)
  is flat over~\(A\);
\item[\textup{(ii)}] the homomorphism
  \[
    u\otimes_Ak\colon M'\otimes_Ak\longrightarrow M\otimes_Ak
  \]
  is injective, where \(k\) is the residue field of~\(A\).
\end{enumerate}
\end{corollary}

The implication \textup{(i)}\(\Rightarrow\)\textup{(ii)} follows from
Proposition~\ref{IV.1.1}. Conversely, \(u\) is first of all injective: it
is enough to check this on associated graded modules, where it follows
from a commutative square that we leave to the reader. Put
\(M''=\operatorname{Coker}(u)\). We then have an exact sequence
\[
  0\longrightarrow M'\longrightarrow M\longrightarrow M''
  \longrightarrow0.
\]
The long exact sequence of \(\operatorname{Tor}\), together with
hypothesis~\textup{(ii)} and
\(\operatorname{Tor}_1^A(M,k)=0\), gives
\(\operatorname{Tor}_1^A(M'',k)=0\). Theorem~\ref{IV.5.6} therefore shows
that \(M''\) is flat over~\(A\).

\begin{corollary}
\label{IV.5.8}
Under the hypotheses of Theorem~\ref{IV.5.6}, let \(J\) be an ideal
of~\(B\) containing \(IB\) and contained in the Jacobson radical. Let
\(\widehat A\) be the \(I\)-adic completion of~\(A\), and let
\(\widehat B\) and~\(\widehat M\) be the \(J\)-adic completions of
\(B\) and~\(M\), respectively. Then \(M\) is flat over~\(A\) if and only
if \(\widehat M\) is flat over~\(\widehat A\).
\end{corollary}

The sufficiency already follows readily from Corollary~\ref{IV.3.2}.
Apply criterion~\textup{(iii)} of Theorem~\ref{IV.5.6} to the situations
\[
  (A,B,I,M)
  \qquad\text{and}\qquad
  (\widehat A,\widehat B,I\widehat A,\widehat M).
\]
The resulting conditions in the two cases are equivalent by
Corollary~\ref{IV.3.2}.

\begin{corollary}
\label{IV.5.9}
Let \(A\to B\to C\) be local homomorphisms of Noetherian local rings,
and let \(M\) be a finitely generated \(C\)-module. (The ring \(C\)
occurs only so that a finiteness condition can be imposed on~\(M\).)
Suppose that \(B\) is flat over~\(A\), and let \(k\) be the residue field
of~\(A\). The following conditions are equivalent:
\begin{enumerate}
\item[\textup{(i)}] \(M\) is flat over~\(B\);
\item[\textup{(ii)}] \(M\) is flat over~\(A\), and \(M\otimes_Ak\) is
  flat over~\(B\otimes_Ak\).
\end{enumerate}
\end{corollary}

The implication \textup{(i)}\(\Rightarrow\)\textup{(ii)} is immediate.
For the converse, apply criterion~\textup{(iii)} of
Theorem~\ref{IV.5.6} to
\[
  (B,C,I=\mathfrak mB,M),
\]
where \(\mathfrak m\) is the maximal ideal of~\(A\). Since
\[
  M\otimes_B(B/I)
  =M\otimes_B(B\otimes_Ak)
  =M\otimes_Ak,
\]
the first condition of that criterion says precisely that
\(M\otimes_Ak\) is flat over~\(B\otimes_Ak\), as required.
\marginpar{\raggedright\footnotesize French source p.~100}%
The second condition is satisfied because \(M\) and~\(B\) are flat
over~\(A\), by associativity of the tensor product.

Of course, using Corollary~\ref{IV.5.5} instead of
Theorem~\ref{IV.5.6} gives an analogous statement without Noetherian or
finiteness hypotheses, provided instead that the ideal~\(\mathfrak m\)
of~\(A\) is nilpotent. The fact that \(\mathfrak m\) was taken maximal
did not enter the argument; nevertheless, in a sense, the case of maximal
\(\mathfrak m\) is the best possible one.

\end{filecontents*}
% Original input: drafts/SGA1_IV_6_English.texfrag
\begin{filecontents*}[overwrite,noheader]{sga1-cumulative-input-027.tex}
\section{Flat morphisms and open subsets}
\label{IV.6}

We first recall some results on constructible sets, also proved in
circulating notes from the Dieudonné--Rosenlicht Seminar on schemes.
\footnote{See EGA~\(\mathrm{0}_{\mathrm{III}}\), 9, and EGA~IV, 1.8
and 1.10.}

Following Chevalley, a subset of a topological space is called
\emph{constructible} if it is a finite union of locally closed subsets.

\begin{lemma}
\label{IV.6.1}
Let \(X\) be a Noetherian topological space and \(Z\subseteq X\). Then
\(Z\) is constructible if and only if, for every irreducible closed
subset \(Y\subseteq X\), the intersection \(Z\cap Y\) is either not
dense in~\(Y\) or contains a nonempty open subset of~\(Y\).
\end{lemma}

Using a familiar lemma from commutative algebra, one deduces:

\begin{lemma}[Chevalley]
\label{IV.6.2}
Let \(f\colon X\to Y\) be a morphism of finite type between
preschemes, with \(Y\) Noetherian. Then \(f(X)\) is constructible.
\end{lemma}

\begin{lemma}
\label{IV.6.3}
Let \(X\) be a Noetherian topological space in which every irreducible
closed subset has a generic point, let \(U\subseteq X\) be
constructible, and let \(x\in X\). Then \(U\) is a neighborhood of~\(x\)
if and only if every generalization \(y\) of~\(x\), that is, every
\(y\) such that \(x\in\overline{\{y\}}\), belongs to~\(U\).
\end{lemma}

In particular:

\begin{corollary}
\label{IV.6.4}
Let \(X\) be a Noetherian topological space in which every irreducible
closed subset has a generic point. A subset \(U\subseteq X\) is open if
and only if it satisfies both of the following conditions:
\begin{enumerate}
\item[(a)] \(U\) contains every generalization of each of its points;
\item[(b)] if \(x\in U\), then \(U\cap\overline{\{x\}}\) contains a
nonempty open subset of \(\overline{\{x\}}\).
\end{enumerate}
\marginpar{\raggedright\footnotesize French source p.~101}%
\end{corollary}

Condition \textup{(b)} and Lemma~\ref{IV.6.1} show that \(U\) is
constructible, and Lemma~\ref{IV.6.3} then shows that \(U\) is a
neighborhood of each of its points.

\begin{corollary}
\label{IV.6.5}
Let \(f\colon X\to Y\) be a morphism of finite type between
preschemes, with \(Y\) locally Noetherian. Let \(x\in X\) and
\(y=f(x)\). Then \(f\) maps every neighborhood of~\(x\) onto a
neighborhood of~\(y\) if and only if, for every generalization \(y'\)
of~\(y\), there is a generalization \(x'\) of~\(x\) such that
\(f(x')=y'\).
\end{corollary}

We may suppose that \(X\) and \(Y\) are affine, hence Noetherian. The
condition is sufficient because it is enough to show that \(f(X)\) is
a neighborhood of~\(y\). The image \(f(X)\) is constructible by
Lemma~\ref{IV.6.2}, and Lemma~\ref{IV.6.3} applies.

Conversely, let \(Y'=\overline{\{y'\}}\), and let \(F\) be the union of
the irreducible components of \(f^{-1}(Y')\) which do not contain~\(x\).
Then \(X-F\) is an open neighborhood of~\(x\), so its image is a
neighborhood of~\(y\), and in particular contains~\(y'\). Choose
\(x'_1\in X-F\) with \(f(x'_1)=y'\). An irreducible component of
\(f^{-1}(Y')\) containing~\(x'_1\) must also contain~\(x\), for
otherwise it would lie in~\(F\). Let \(x'\) be its generic point.
Then \(x'\) is a generalization of~\(x\), and \(f(x')\) is a
generalization of \(f(x'_1)=y'\) lying in~\(Y'\); hence \(f(x')=y'\).

\begin{theorem}
\label{IV.6.6}
Let \(f\colon X\to Y\) be a morphism locally of finite type, with \(Y\)
locally Noetherian, and let \(F\) be a coherent sheaf on~\(X\), with
support~\(X\), flat relative to~\(Y\). Then \(f\) is an open morphism;
that is, it maps open sets to open sets.
\end{theorem}

It is enough to verify the criterion of Corollary~\ref{IV.6.5} at each
\(x\in X\). Generalizations of~\(x\) correspond to prime ideals of
\(\mathcal O_x\), and generalizations of \(y=f(x)\) correspond to
prime ideals of~\(\mathcal O_y\). We must therefore show that every
prime ideal of~\(\mathcal O_y\) is the contraction of a prime ideal of
\(\mathcal O_x\). The stalk \(F_x\) is a nonzero finitely generated
\(\mathcal O_x\)-module flat over~\(\mathcal O_y\), hence faithfully
flat over~\(\mathcal O_y\) by Corollary~\ref{IV.2.2}.
Proposition~\ref{IV.2.3} now finishes the proof.

\paragraph{Remarks.}
Since flatness is preserved under base change, the preceding theorem
also shows that \(f\) is universally open. If \(Y\) is integral and
\(X\) is of finite type over~\(Y\), it was not known whether the
restriction of~\(f\) to each irreducible component \(X_i\) of~\(X\)
must be open, or even merely equidimensional.
\footnote{The answer to the second question is affirmative and to the
first negative, even when \(f\) is étale; see EGA~IV, 12.1.1.5, and
EGA Errata~IV, 33.}
Here equidimensional means that all components of all fibers have the
same dimension; one knows only that \(X_i\) dominates~\(Y\).
\marginpar{\raggedright\footnotesize French source p.~102}%
The question is related to the following one. Let \(A\to B\) be a
local homomorphism of Noetherian local rings, with \(B\) flat over~\(A\)
and \(\mathfrak mB\) an ideal of definition of~\(B\), which in
particular implies \(\dim B=\dim A\). Is it true that
\[
  \dim B/\mathfrak p_i=\dim B
\]
for every minimal prime ideal \(\mathfrak p_i\) of~\(B\)?

We merely note that the answer to the first question is negative if
the flatness hypothesis in Theorem~\ref{IV.6.6} is replaced by the
assumption that \(f\) is universally open.

\begin{lemma}
\label{IV.6.7}
Let \(A\) be a Noetherian integral domain, let \(B\) be a finitely
generated \(A\)-algebra, and let \(M\) be a finitely generated
\(B\)-module. There exists a nonzero \(f\in A\) such that \(M_f\) is
free, and hence flat, over~\(A_f\).
\end{lemma}

Let \(K\) be the fraction field of~\(A\). Then \(B\otimes_AK\) is a
finitely generated \(K\)-algebra and \(M\otimes_AK\) is a finitely
generated module over it. Let \(n\) be the dimension of the support of
this module; we argue by induction on~\(n\). If \(n<0\), that is, if
\(M\otimes_AK=0\), choose finitely many generators of~\(M\) over~\(B\).
Some nonzero \(f\in A\) annihilates all of them, hence annihilates~\(M\),
so \(M_f=0\).

Suppose \(n\ge0\). The \(B\)-module \(M\) has a prime filtration whose
successive quotients are \(B/\mathfrak p_i\). Since an extension of
free modules is free, we reduce to \(M=B/\mathfrak p\), or,
equivalently, to \(M=B\) with \(B\) an integral \(A\)-algebra. Applying
Noether normalization to \(B\otimes_AK\), after inverting a suitable
nonzero element of~\(A\) we may suppose that \(B\) is integral over
\[
  C=A[t_1,\dots,t_n].
\]
It is then a finitely generated torsion-free \(C\)-module. If \(m\) is
its rank, there is an exact sequence of \(C\)-modules
\[
  0\longrightarrow C^m\longrightarrow B\longrightarrow M'
  \longrightarrow0
\]
with \(M'\) a torsion \(C\)-module. The Krull dimension of the
\(C\otimes_AK\)-module \(M'\otimes_AK\) is strictly smaller than the
dimension \(n\) of~\(C\otimes_AK\). By induction, after inverting
another suitable nonzero \(f\in A\), \(M'\) is free over~\(A\).
\marginpar{\raggedright\footnotesize French source p.~103}%
The module \(C^m\) is also free over~\(A\), so the sequence splits over
\(A\) and \(B\) is free over~\(A\).

\begin{lemma}
\label{IV.6.8}
Let \(A\) be Noetherian, let \(B\) be a finitely generated
\(A\)-algebra, let \(M\) be a finitely generated \(B\)-module, and let
\(\mathfrak p\) be a prime ideal of~\(B\), with contraction
\(\mathfrak q\) in~\(A\). Suppose that \(M_{\mathfrak p}\) is flat
over \(A_{\mathfrak q}\), equivalently over~\(A\). Then there exists
\(g\in B-\mathfrak p\) such that:
\begin{enumerate}
\item[(a)] \((M/\mathfrak qM)_g\) is flat over~\(A/\mathfrak q\);
\item[(b)] \(\operatorname{Tor}_1^A(M,A/\mathfrak q)_g=0\).
\end{enumerate}
\end{lemma}

Apply Lemma~\ref{IV.6.7} to
\[
  (A/\mathfrak q,\ B/\mathfrak qB,\ M/\mathfrak qM).
\]
There is first an \(f\in A-\mathfrak q\) such that
\((M/\mathfrak qM)_f\) is flat over~\(A/\mathfrak q\). Since
\(M_{\mathfrak p}\) is flat over~\(A\),
\[
  \operatorname{Tor}_1^A(M,A/\mathfrak q)_{\mathfrak p}
  =
  \operatorname{Tor}_1^A(M_{\mathfrak p},A/\mathfrak q)
  =0.
\]
The left-hand module is finitely generated over~\(B\), so there exists
\(g\in B-\mathfrak p\) for which \textup{(b)} holds. Replacing \(g\)
by \(gf\), we may arrange that \textup{(a)} holds as well.

\begin{corollary}
\label{IV.6.9}
With the notation of Lemma~\ref{IV.6.8}, let \(g\) be as there. For
every prime ideal \(\mathfrak p'\) of~\(B\) containing
\(\mathfrak p\) and not containing~\(g\), the module
\(M_{\mathfrak p'}\) is flat over~\(A\), equivalently over
\(A_{\mathfrak q'}\), where \(\mathfrak q'\) is the contraction of
\(\mathfrak p'\) to~\(A\).
\end{corollary}

Apply criterion \textup{(ii)} of Theorem~\ref{IV.5.6} to
\[
  (A,\ B_{\mathfrak p'},\ \mathfrak q,\ M_{\mathfrak p'}),
\]
using localization of \(\operatorname{Tor}\).

\begin{theorem}
\label{IV.6.10}
Let \(f\colon X\to Y\) be a morphism of finite type, with \(Y\) locally
Noetherian, and let \(F\) be a coherent sheaf on~\(X\). The set
\[
  U=\{x\in X\mid F_x\text{ is flat over }\mathcal O_{f(x)}\}
\]
is open.
\end{theorem}

\subsubsection*{Proof}

We may suppose that \(X\) and \(Y\) are affine, with rings \(B\) and
\(A\), and that \(F\) is defined by a finitely generated \(B\)-module
\(M\). Apply Corollary~\ref{IV.6.4}. Condition \textup{(a)} follows
immediately from Proposition~\ref{IV.1.2}\textup{(i)}, and condition
\textup{(b)} is exactly Lemma~\ref{IV.6.8} and
Corollary~\ref{IV.6.9}.

For many purposes the following weaker form of
Theorem~\ref{IV.6.10} suffices. It already follows from
Lemma~\ref{IV.6.7}, and therefore requires neither constructible sets
nor Theorem~\ref{IV.5.6}.

\begin{corollary}
\label{IV.6.11}
\marginpar{\raggedright\footnotesize French source p.~104}%
Under the hypotheses of Theorem~\ref{IV.6.10}, suppose that \(Y\) is
integral. There exists a nonempty open subset \(V\subseteq Y\) such
that \(F\) is flat relative to~\(Y\) at every point of
\(f^{-1}(V)\).
\end{corollary}

Indeed, the open set \(U\) contains the fiber over the generic point
of~\(Y\), since its local ring is a field. As \(X\) is of finite type
over~\(Y\), \(U\) therefore contains an open set of the form
\(f^{-1}(V)\).

Corollary~\ref{IV.6.11} also readily implies the following statement
when \(Y\) is Noetherian but not necessarily integral. There exists a
partition of~\(Y\) into locally closed subsets \(Y_i\) such that,
after giving \(Y_i\) its induced reduced structure, the restriction of
\(F\) to \(X_i=X\times_YY_i\) is flat relative to~\(Y_i\).

\end{filecontents*}
% Original input: drafts/SGA1_IV_INTRODUCTION_English.texfrag
\begin{filecontents*}[overwrite,noheader]{sga1-cumulative-input-028.tex}
% SGA 1, Exposé IV introduction, French TeX lines 6269--6311.

\chapter{Flat morphisms}
\label{IV}
\marginpar{\raggedright\footnotesize French source p.~87}%

We give here chiefly the properties of flatness used in the preceding
exposés. A more detailed study will appear in Chapter~IV of the
forthcoming \emph{Éléments de Géométrie Algébrique}\footnote{See EGA
IV, 11 and 12.}, where the following situation is treated
systematically: \(X\) is locally of finite type over a locally
Noetherian prescheme~\(Y\), and \(\mathcal F\) is a coherent sheaf on
\(X\), flat over~\(Y\); one seeks relations among properties of~\(Y\),
properties of~\(\mathcal F\), and properties of the coherent sheaves
induced by~\(\mathcal F\) on the fibers of \(X\to Y\), especially with
respect to dimension, cohomological dimension, and depth. In
particular, this yields a systematic way of obtaining theorems of
Seidenberg or Bertini type for hyperplane sections.

The essential result for applying flatness methods in this setting is
the following, to be proved below: if \(Y\) is integral, \(X\) is of
finite type over~\(Y\), and \(\mathcal F\) is coherent on~\(X\), then
there exists a nonempty open subset \(U\subseteq Y\) such that
\(\mathcal F\) is flat over~\(Y\) at every point of~\(X\) lying
over~\(U\).

A second, and perhaps still more important, way in which flatness enters
algebraic geometry is through descent theory. See, for example,
Grothendieck's two Bourbaki Seminar talks on the subject, and, for a more
detailed account, Exposés~VIII and~IX below. Flatness thus appears to be
one of the central technical notions of algebraic geometry.

Recall that Serre introduced the notions of flatness and faithful
flatness in GAGA. The material of Sections~\ref{IV.1} and~\ref{IV.2}
below may also be found in Bourbaki's \emph{Algèbre commutative}, which,
as its title indicates, does not restrict itself to commutative base
rings.\footnote{N.~Bourbaki, \emph{Algèbre commutative}, Chapter~I
(Modules plats), Act. Sci. Ind. 1290, Hermann, Paris, 1961.}

Unlike the preceding exposés, we do not assume that the rings under
consideration are necessarily Noetherian.
\end{filecontents*}
% Original input: drafts/SGA1_IX_2_English.texfrag
\begin{filecontents*}[overwrite,noheader]{sga1-cumulative-input-029.tex}
\section{Submersive and universally submersive morphisms}
\label{IX.2}

\begin{definition}
\label{IX.2.1}
A morphism \(g\colon S'\to S\) of preschemes is called
\emph{submersive}
%
\index{submersive (morphism)|hyperpage}%
%
if it is surjective and makes \(S\) a quotient topological space
of~\(S'\) (\ie every subset \(U\) of~\(S\) for which \(g^{-1}(U)\)
is open is itself open). We say that \(g\) is \emph{universally
submersive}
%
\index{universally submersive (morphism)|hyperpage}%
%
if, for every morphism \(T\to S\), the morphism
\(g'\colon T'=S'\times_ST\to T\) deduced from~\(g\) by base change
is submersive.
\end{definition}

It is immediate that the composite of two submersive (respectively
universally submersive) morphisms is submersive (respectively
universally submersive), and that a base change of a universally
submersive morphism is universally submersive, as the definition was
made precisely for this purpose. If \(fg\) is submersive
(respectively universally submersive), then \(f\) is so.

\begin{exemples}
\label{IX.2.2}
a) A surjective morphism which is open or closed is submersive.
Consequently, a surjective universally closed or universally open
morphism is universally submersive. For example, \emph{a proper
surjective morphism is universally submersive}. Moreover, \emph{a
faithfully flat and quasi-compact morphism is universally
submersive} (\ifthenelse{\boolean{orig}}{VIII~}{}\Ref{VIII.4.3}). These will be the two most
important cases for us.
\end{exemples}

The arguments of \ifthenelse{\boolean{orig}}{VIII~}{}\Ref{VIII.4.3} may
% page: 232
\marginpar{232}
%
be applied to a submersive or universally submersive morphism
\(g\colon S'\to S\). In particular, one obtains:

\begin{proposition}
\label{IX.2.3}
Suppose that \(g\colon S'\to S\) is submersive. Then the following
diagram of maps is exact:
\[
  \xymatrix@C=.5cm{
  \Ouv(S)\ar[r]&
  \Ouv(S')\ar@<2pt>[r]\ar@<-2pt>[r]&
  \Ouv(S''),
  }
\]
where \(S''=S'\times_SS'\), and where \(\Ouv(X)\) denotes the set of
open subsets of the prescheme~\(X\).
\end{proposition}

\begin{proposition}
\label{IX.2.4}
Let \(g\colon S'\to S\) be universally submersive, let
\(f\colon X\to Y\) be an \(S\)-morphism, and let
\(f'\colon X'\to Y'\) be the \(S'\)-morphism deduced from~\(f\) by
base change. For \(f\) to be open (respectively closed), it is
sufficient that \(f'\) be so. The morphism \(f\) is universally open,
respectively universally closed, respectively separated, if and only
if \(f'\) is so. If, in addition, \(g\) is quasi-compact and \(f\) is
locally of finite type, then \(f\) is proper if and only if \(f'\) is
proper.
\end{proposition}

For the last assertion, observe that if \(f'\) is proper, and hence
quasi-compact, then \(f\) is quasi-compact
(\ifthenelse{\boolean{orig}}{VIII~}{}\Ref{VIII.3.3}), and therefore of finite type since it is
locally of finite type. Moreover, by what precedes, it is separated
and universally closed, and hence proper.

\begin{proposition}
\label{IX.2.5}
Let \(S'\) be a prescheme of finite type over the spectrum \(S\) of
a complete noetherian local ring. Suppose that the fiber over the
closed point \(s\) of~\(S\) is finite, so that the local rings
of~\(S'\) at the points \(s'\) of this fiber are finite over
\(A=\cal O_s\). Let \(S''\) be the scheme which is the disjoint sum
of the spectra of the rings \(\cal O_{S',s'}\) in question, regarded
as a finite \(S\)-scheme. Then \(g\colon S'\to S\) is universally
submersive if and only if the structural morphism \(S''\to S'\) is
surjective.
\end{proposition}

Since there is a natural \(S\)-morphism \(S''\to S'\), and a finite
surjective morphism is universally submersive by~\Ref{IX.2.2}, the
stated condition is sufficient. We therefore show that if
\(S''\to S'\) is not surjective, then \(g\) is not universally
submersive. Let \(t\) be a point of~\(S\) not in the image of~\(S''\).
There is then an \(S\)-scheme \(T\), the spectrum of a discrete
valuation ring, whose image in~\(S\) is
% page: 233
\marginpar{233}
%
\(\{s,t\}\). The image of~\(S''\) in~\(S'\) is open, since the
morphism \(S''\to S'\) is a local isomorphism; moreover, this image
contains \(S'_s\) and does not meet \(S'_t\). Its inverse image in
\(T'=S'\times_ST\) is therefore \emph{open}, and is precisely the
inverse image of the closed point of~\(T\). This shows that
\(T'\to T\) is not submersive, and hence that \(S'\to S\) is not
universally submersive.

\begin{remarque}
\label{IX.2.6}
Using the criterion \ifthenelse{\boolean{orig}}{IV~}{}\Ref{IV.6.3} for a constructible subset of a
noetherian space to be open, one readily obtains the following
valuative criterion for a morphism \(g\colon S'\to S\) \emph{of
finite type}, with \(S\) locally noetherian, to be universally
submersive: it is necessary and sufficient that, for every
\(S\)-scheme \(T\) which is the spectrum of a discrete valuation
ring, on setting \(T'=S'\times_ST\), the inverse image in~\(T'\) of
the closed point of~\(T\) is not open.
\end{remarque}

\end{filecontents*}
% Original input: drafts/SGA1_IX_3_English.texfrag
\begin{filecontents*}[overwrite,noheader]{sga1-cumulative-input-030.tex}
\section{Descent of \'etale morphisms of preschemes}
\label{IX.3}

\begin{proposition}
\label{IX.3.1}
Let \(g\colon S'\to S\) be a \emph{surjective} morphism of
preschemes, let \(X\) and \(Y\) be two preschemes over~\(S\), and let
\(X'\) and \(Y'\) be their inverse images over~\(S'\). If \(Y\) is
unramified over~\(S\), then the canonical map
\[
  \Hom_S(X,Y)\longrightarrow\Hom_{S'}(X',Y')
\]
is injective.
\end{proposition}

Indeed, by~\Ref{IX.1.6}, an \(S\)-morphism \(f\colon X\to Y\) is
determined by the underlying set of its graph \(\Gamma\), which is a
subset of \(Z=X\times_SY\). Since
\[
  Z'=Z\times_SS'=X'\times_{S'}Y'\longrightarrow Z
\]
is surjective, because \(S'\to S\) is surjective, the subset
\(\Gamma\) is determined by its inverse image in
\(X'\times_{S'}Y'\), which is precisely the underlying set of the
graph of~\(f'\). This proves the assertion.

A subset \(\Gamma\) of~\(Z\) is the graph of an \(S\)-morphism
\(f\colon X\to Y\) if and only if it is open and the morphism
\(\Gamma\to X\) induced by~\(\pr_1\) is radicial
% page: 234
\marginpar{234}
%
and surjective; see~\Ref{IX.1}. When the first property holds, the
second holds if and only if the inverse image \(\Gamma'\) of
\(\Gamma\) in~\(Z'\) satisfies the same condition
(\ifthenelse{\boolean{orig}}{VIII~}{}\Ref{VIII.3.1}). Finally, if \(Z'\to Z\) is known to be
submersive, as is in particular the case when \(S'\to S\) is
universally submersive, then \(\Gamma\) is open if and only if
\(\Gamma'\) is open.

Thus \(\Hom_S(X,Y)\) is in bijection with the set of open subsets
\(\Gamma'\) of~\(Z'\) for which the projection
\(\pr_1\colon\Gamma'\to X'\) is radicial and surjective (\ie those
corresponding to an \(S'\)-morphism \(f'\colon X'\to Y'\)), and which
are saturated for the equivalence relation defined by \(Z'\to Z\).
The latter condition says that their two inverse images in
\[
  Z''=Z'\times_ZZ'=Z\times_SS'',
  \qquad S''=S'\times_SS',
\]
under the two projections are equal. These inverse images are the
graphs of the two \(S''\)-morphisms \(X''\to Y''\) deduced from~\(f'\)
by base change along the two projections \(S''\to S'\). We have
therefore proved:

\begin{proposition}
\label{IX.3.2}
Let \(g\colon S'\to S\) be a \emph{universally submersive} morphism
of preschemes, let \(S''=S'\times_SS'\), let \(X\) and \(Y\) be two
\(S\)-preschemes, let \(X'\) and \(Y'\) be their inverse images
over~\(S'\), and let \(X''\) and \(Y''\) be their inverse images
over~\(S''\). If \(Y\) is \'etale over~\(S\), then the following
canonical diagram of maps is exact:
\[
  \xymatrix@C=.5cm{
  {\Hom_S(X,Y)}\ar[r]&
  {\Hom_{S'}(X',Y')}\ar@<2pt>[r]\ar@<-2pt>[r]&
  {\Hom_{S''}(X'',Y'')}.
  }
\]
\end{proposition}

Taking \(X\) and \(Y\) \'etale over~\(S\), one obtains the following
statement, which moreover recovers~\Ref{IX.3.2}, even after restricting
to \(X=S\), a case to which one may always reduce
in~\Ref{IX.3.2} by the base change \(X\to S\):

\begin{corollary}
\label{IX.3.3}
A universally submersive morphism of preschemes is a descent
morphism for the fibered category of preschemes \'etale over other
preschemes.
\end{corollary}

I do not know whether it is necessarily an \emph{effective} descent
morphism for the fibered category in question, even under the
additional hypotheses that \(S\) is noetherian and \(g\) is of finite
type, and even if one restricts to \'etale coverings. In the next
section we shall nevertheless give useful effectivity criteria.

\begin{corollary}
\label{IX.3.4}
Let
% page: 235
\marginpar{235}
%
\(g\colon S'\to S\) be a universally submersive morphism whose
fibers \(g^{-1}(s)\) are \emph{geometrically connected},
%
\index{geometrically connected|hyperpage}%
%
\ie such that for every extension \(K/\kres(s)\), the scheme
\(g^{-1}(s)\otimes_{\kres(s)}K\) is connected. Then \(S'\) is
connected whenever \(S\) is connected. The functor, defined by~\(g\),
from the category of preschemes \'etale over~\(S\) to the
\ifthenelse{\boolean{orig}}
{categories}
{category}
% correction: category.
of preschemes \'etale over~\(S'\) is fully faithful.
\end{corollary}

A subset of~\(S'\) which is both open and closed is saturated for the
set-theoretic equivalence relation defined by~\(g\), since the fibers
are connected. It is therefore the inverse image of a subset of~\(S\),
which is necessarily open and closed because \(g\) is submersive. Thus
if \(S\) is connected, so is \(S'\).

This also implies the following fact: if \(f\) and \(g\) are two
morphisms with universally connected fibers,
%
\index{universally connected: synonym of geometrically connected|hyperpage}%
%
and \(f\) is universally submersive, then the composite \(fg\) has
universally connected fibers. If \(S_1'\) and \(S_2'\) over~\(S\)
have universally connected fibers, then so does
\(S_1'\times_SS_2'\). In particular, under the hypotheses of
\Ref{IX.3.4}, \(S''\) has universally connected fibers over~\(S\).

Now let \(X\) and \(Y\) be \'etale over~\(S\), and let \(u'\) be an
\(S'\)-morphism from~\(X'\) to~\(Y'\). We prove that it is compatible
with the descent data, which gives the desired conclusion by
\Ref{IX.3.3}. Let \(u_1''\) and \(u_2''\) be the two
\(S''\)-morphisms \(X''\to Y''\) deduced from~\(u'\). The
subprescheme of~\(S''\) on which \(u_1''\) and \(u_2''\) coincide is
an induced open subprescheme, closed on each fiber, since it is the
inverse image of the diagonal prescheme of~\(Y''\) over~\(S''\)
\footnote{Observe that the fibers of~\(\ifthenelse{\boolean{orig}}{S'}{Y}\) over~\(S\) are
separated!}. It is therefore the inverse image of a subset of~\(S\).
Since it contains the diagonal in~\(S''\), it is all of~\(S''\).
Thus \(u_1''=u_2''\), as required.

\end{filecontents*}
% Original input: drafts/SGA1_IX_4_English.texfrag
\begin{filecontents*}[overwrite,noheader]{sga1-cumulative-input-031.tex}
\section{Descent of \'etale preschemes: effectivity
criteria}
\label{IX.4}

\ifthenelse{\boolean{orig}}{}
{\enlargethispage{.7cm}}%
\begin{proposition}
\label{IX.4.1}
Let $g\colon S'\to S$ be a faithfully flat and
quasi-compact morphism. Then~$g$ is an effective descent morphism for
the fibered category of preschemes that are \'etale,
separated, and of finite type over other preschemes.
\end{proposition}

Indeed, it is a descent morphism for the
fibered category in question, by~\Ref{IX.3.3} or
\ifthenelse{\boolean{orig}}{VIII~}{}\Ref{VIII.5.2}, as one prefers. It remains to show that if~$X'$ is
\'etale, separated, and of finite type over~$S'$, and endowed with a
descent datum relative to $g\colon S'\to S$, the
latter is effective in the fibered category in
question. But one sees easily that if~$X$ is a prescheme
over~$S$, then
% page: 236
\marginpar{236}
%
it is \'etale over~$S$ if and only if it is \'etale over~$S'$
(by Definition~\Ref{IX.1.1} and \loccit \Ref{VIII.3.6}). Thus it is \'etale, separated, and of
finite type over~$S$ if and only if~$X'$ is so over~$S'$, \cf for
example~\Ref{IX.2.4}. It therefore suffices to ensure the
effectivity of the descent datum on~$X\ifthenelse{\boolean{orig}}{}{'}$ for the
fibered category of arrows in~$\Sch$. But~$X'$ is
quasi-affine over~$S'$ by \ifthenelse{\boolean{orig}}{VIII~}{}\Ref{VIII.6.2}
and~\Ref{VIII.6.6}. We may then conclude by using
\ifthenelse{\boolean{orig}}{VIII~}{}\Ref{VIII.7.9}. The reader will moreover note that the
proof requires less if one restricts oneself to preschemes that are
\'etale and \emph{finite} over other preschemes, since one may then invoke
\ifthenelse{\boolean{orig}}{VIII~}{}\Ref{VIII.2.1} directly.

\begin{corollaire}
\label{IX.4.2}
Let $g\colon S'\to S$ be a universally submersive morphism, let $X'$ be
an
\ifthenelse{\boolean{orig}}
{$S'$-preschemes}
{$S'$-prescheme}
% correction: $S'$-pr\'esch\'ema
that is \'etale, separated, and of finite type, endowed with a
descent datum relative to~$g$, let $S_1\to S$ be a
faithfully flat and quasi-compact morphism, and let $S_1'$ and~$X_1'$ be deduced
from~$S'$ and~$X'$ by base change, so that $S_1'\to S_1$
is universally submersive, while $X_1'$ is \'etale, separated, and
of finite type over~$S_1'$, and is endowed with a descent datum
relative to $g_1\colon S_1'\to S_1$. The descent datum
on~$X'$ is effective if and only if the
descent datum on~$X_1'$ is effective.
\end{corollaire}

This follows from descent theory in
categories~\cite{IX.D}, in view of~\Ref{IX.4.1} and~\Ref{IX.3.3}.

One proves similarly:

\begin{corollaire}
\label{IX.4.3}
Let $g\colon S'\to S$ be a universally submersive morphism, let $X'$ be
an~$S'$-prescheme that is \'etale and endowed with a descent datum
relative to~$g$, and let $(S_i)$ be a covering of~$S$ by
open subsets. The descent datum is effective if and only
if, for every~$i$, the corresponding descent datum
on~$X_i'=X\ifthenelse{\boolean{orig}}{}{'}\times_S S_i$, relative to the morphism
$g_i\colon S_i'=S'\times_S S_i\to S_i$, is effective.
\end{corollaire}

This last result leads to the following local
effectivity criterion:

\enlargethispage{.3cm}
\begin{proposition}
\label{IX.4.4}
Let $g\colon S'\to S$ be a morphism of finite presentation
\textup{VIII}~\Ref{VIII.3.6}
% page: 237
\marginpar{237}
%
and universally submersive, let $X'$ be a prescheme that is \'etale and
of finite presentation over~$S'$, endowed with a descent datum
relative to~$g$, and finally let~$a$ be a point of~$S$. In order that there exist
an open neighborhood~$U$ of~$a$ such that the corresponding descent datum
on~$X_U'=X'\times_S U$ relative to the morphism
$$g_U\colon S_U'=S'\times_S U\to S_U=U$$ is effective, it is necessary and it
is sufficient that the corresponding descent datum on $X_a'=X'\times_S\Spec(\cal{O}_a)$, relative to the morphism
$$g_a\colon S_a'=S'\times_S\Spec(\cal{O}_a)\to S_a=\Spec(\cal{O}_a),$$
be effective.
\end{proposition}

Necessity being trivial, let us prove sufficiency. We
thus have an \'etale prescheme of finite type~$X_a$
over~$S_a$, and an isomorphism
\begin{equation*}
\label{eq:IX.4.*}
\tag{$*$} {X_a'\isomto X_a\times_{S_a}S_a'}
\end{equation*}
compatible with the descent data. In accordance with
a general elementary limit argument for preschemes
defined over a direct limit of rings (here the rings~$A_f$,
where~$A$ is the ring of an affine open neighborhood of~$a$, and
where~$f$ ranges over the elements of~$A$ that do not belong to
the prime ideal corresponding to~$a$), one can find an
open neighborhood~$U$ of~$a$, an \'etale prescheme of finite
type~$X_U$ over~$U=S_U$, and an $S_a$-isomorphism $X_a\isomto
X_U\times_{S_U}S_a$. Moreover, on taking~$U$ sufficiently small, one may then
suppose that the isomorphism~\eqref{eq:IX.4.*} comes from an
isomorphism:
$$
X_U'\isomto X_U\times_{S_U} S_U';
$$
the latter might not be compatible with the descent
data; nevertheless, after shrinking~$U$, it will be
compatible with the descent data. This completes the
proof.

\begin{corollaire}
\label{IX.4.5}
Under the conditions of~\Ref{IX.4.4}, the descent datum
on~$X'$ is effective if and only if, for every
$a\in S$, the corresponding descent datum on~$X_a'$,
relative to the morphism
$S_a'=S'\times_S\Spec(\cal{O}_a)\to\Spec(\cal{O}_a)$, is effective.
When~$S$ is locally Noetherian and~$X'$ is separated
over~$S'$, one may in the preceding criterion also
replace~$\cal{O}_a$ by its completion.
\end{corollaire}

The
% page: 238
\marginpar{238}
%
first assertion follows from~\Ref{IX.4.4} and~\Ref{IX.4.3}; the
second is then a consequence of~\Ref{IX.4.2}. Using~\Ref{IX.4.2}
once more and the fact that for every Noetherian local
ring~$A$, one can find a complete Noetherian local
ring~$B$, and a local homomorphism $A\to B$, such that~$B$ is flat
over~$A$ and $B/\goth{m}B$ is a prescribed extension of the residue
field $k=A/\goth{m}$ of~$A$, one obtains:

\begin{corollaire}
\label{IX.4.6}
Under the conditions of~\Ref{IX.4.4}, suppose in addition that~$X'$ is
separated over~$S'$, and that~$S$ is locally Noetherian. The
descent datum on~$X'$ is effective if and only
if, for every prescheme~$S_1$ over~$S$ that is the spectrum of a
complete local ring with algebraically closed residue field, the
corresponding descent datum on~$X_1'=X'\times_S S_1$,
relative to the morphism $g_1\colon S_1'\to S_1$, is effective.
\end{corollaire}

\begin{theoreme}
\label{IX.4.7}
Let $g\colon S'\to S$ be a finite and surjective morphism of finite
presentation (the latter hypothesis being a
consequence of the others if~$S$ is locally
Noetherian)\footnote{One can show that it in fact suffices for~$g$
to be an \emph{integral} morphism, by reducing to the case in the text through a
limit process in the style of EGA IV~8.}.
Then~$g$ is an effective descent morphism for the fibered
category of preschemes that are \'etale, separated, and of
finite type over other preschemes.
\end{theoreme}

We must show that if~$X'$ is \'etale, separated, and of
finite type over~$S'$, and endowed with a descent datum relative
to~$g$, then this datum is effective.
Using~\Ref{IX.4.3}, one reduces easily to the case where~$S$
is Noetherian. By~\Ref{IX.4.5}, we may therefore
suppose that~$S$ is the spectrum of a Noetherian local ring, and
hence in particular that
$$
\dim S=n<+\infty.
$$
We then argue by induction on~$\dim S=n$, the assertion
being trivial for $n<0$. Suppose therefore that $n\geq 0$ and that the
theorem has been proved for dimensions $n'<n$. By~\Ref{IX.4.6},
we are reduced to the case where~$S$ is the
spectrum of a complete local ring, so~$S'$ is a finite union
of spectra of complete local rings. We therefore have
$$
X'=X_1'\sqcup X_2'
$$
where~$X_1'$ is \emph{finite} over~$S'$, and where~$X_2'$ has no
point above any of the closed points
% page: 239
\marginpar{239}
%
of~$S'$. Consider the morphisms
$$
\xymatrix@C=.5cm{ q_1,q_2\colon X'' \ar@<2pt>[r] \ar@<-2pt>[r] & X' }
$$
corresponding to the descent datum, compatible with
$\xymatrix@C=.5cm{p_1,p_2\colon S'' \ar@<2pt>[r] \ar@<-2pt>[r] & S'}$. One
sees at once that
$$
X''=q_i^{-1}(X_1')\sqcup q_i^{-1}(X_2')\quad i=1,2
$$
is the analogous canonical decomposition of~$X''$ over~$S''$, which
implies $q_1^{-1}(X_1')=\allowbreak q_2^{-1}(X_1')$, and consequently~$X_1'$ and~$X_2'$
are endowed with induced descent data. Now let~$T$ be the open subset
of~$S$ complementary to its closed point; then $T'=S'\times_S
T$ is the part of~$S'$ complementary to the set of closed
points, and~$X_2'$, which lies entirely over~$T'$,
is endowed with a descent datum relative to the morphism $T'\to
T$ induced by~$g$. Since the latter is finite and surjective, and $\dim
T<\dim S=n$, this descent datum is effective by
the induction hypothesis. Thus one sees that it suffices to
prove that the descent datum on~$X_1'$ is effective, so
we may now suppose that~$X'$ is \'etale and \emph{finite} over~$S'$.
(N.B. the inductive argument is unnecessary if one restricts
to \'etale coverings in the statement~\Ref{IX.4.7}).
Let~$S_0$ then be the spectrum of the residue field of~$A$, let
$S_0'=S'\times_S S_0$, and define~$S_0''$, $S_0'''$ similarly
from the fibered squares and cubes~$S''$ and~$S'''$
of~$S'$ over~$S$. By~\Ref{IX.1.8}, the morphisms $S_0\to S$,
$S_0'\to S'$, etc. induce equivalences for the
categories of \'etale coverings of~$S$ and~$S_0$, on the one
hand, and of~$S'$ and~$S_0'$ on the other hand, \ifthenelse{\boolean{orig}}{etc...}{etc.} By the standard consequences of
descent theory in categories~\cite{IX.D}, it follows
that $g\colon S'\to S$ is an effective descent morphism
for the fibered category of \'etale coverings if and
only if the same is true of~$g_0\colon S_0'\to S_0$.
But this is indeed the case, for example as a special case of~\Ref{IX.4.1}.
This completes the proof.

\begin{corollaire}
\label{IX.4.8}
The conclusion of~\Ref{IX.4.7} remains valid if one assumes only that
$S'\to S$ is universally submersive, of finite type, and quasi-finite,
provided~$S$ is assumed locally Noetherian.
\end{corollaire}

Indeed,
% page: 240
\marginpar{240}
%
by~\Ref{IX.4.6}, we may suppose that~$S$ is the
spectrum of a complete Noetherian local ring. Then, by~\Ref{IX.2.5},
there exists a finite and surjective morphism $S_1\to S$,
and an $S$-morphism $S_1\to S'$. Since $S_1\to S$ is a universal
strict descent morphism for the fibered category under
consideration, by~\Ref{IX.4.7}, and $S'\to S$ is a
universal descent morphism for that category, \Ref{IX.4.8} follows
from the general formal consequences~\cite{IX.D}.

\begin{corollaire}
\label{IX.4.9}
Let $g\colon S'\to S$ be a morphism of finite type, surjective and
universally open, with~$S$ locally Noetherian. Then~$g$
is an effective descent morphism for the fibered
category of preschemes that are \'etale, separated, and of
finite type over other preschemes.
\end{corollaire}

Proceeding as in~\Ref{IX.4.7}, we reduce to the case
where~$S$ is the spectrum of a complete Noetherian local
ring~$A$. Let~$A_1$ be a finite algebra over~$A$, with
spectrum~$S_1$, such that $S_1\to S$ is finite and \emph{surjective},
and hence a universal effective descent morphism for the
fibered category under consideration, by~\Ref{IX.4.7}.
It then follows from the general theorems~\cite{IX.D}
that~$g$ is an effective descent morphism for the fibered
category under consideration if and only if the corresponding morphism
$g_1\colon S_1'=S'\times_S S_1\to S_1$ is so. Since the
latter satisfies the same hypotheses as~$g$, we are
reduced to proving~\Ref{IX.4.9} with~$S_1$ in place of~$S$.
Taking first for~$A_1$ the direct product of the $A/\goth{p}_i$,
for the minimal prime ideals~$\goth{p}_i$ of~$A$, we are
reduced to the case where~$A$ is an \emph{integral domain}. One then
shows\footnote{\Cf EGA IV~14.3.13 and~14.5.4.} that there exists an
integral subscheme~$S_1$ of~$S'$, quasi-finite over~$S$ and
dominating~$S$, passing through a point of the fiber of~$S'$ over the closed
point~$y$ of~$S$ (using the fact that~$S'$ is universally
open and of finite type over the integral Noetherian local scheme~$S$, and
$S_y'\neq\emptyset$). Since~$A$ is complete, $S_1$ is finite over~$S$,
and since it dominates~$S$, the morphism $S_1\to S$ is surjective.
Replacing~$S$ once more by~$S_1$, we are reduced to the
case where~$S'$ has a section over~$S$, in which case the statement is
trivial.

\begin{theoreme}
\label{IX.4.10}
Let $g\colon S'\to S$ be a finite, radicial, and surjective morphism of
finite presentation (the latter condition being
superfluous if~$S$ is locally Noetherian\footnote{It even suffices
for~$g$ to be integral, radicial, and surjective, as follows from
an easy reduction to the case in the text, in the style of EGA~IV~8, \cf SGA~4
VIII~1.1.}).
% page: 241
\marginpar{241}
%
Then the inverse-image functor induces an equivalence from the
category of preschemes that are \'etale over~$S$ to the
category of preschemes that are \'etale over~$S'$.
\end{theoreme}

Since the diagonal morphisms from~$S'$ into $S'\times_S S'$ and
$S'\times_S S'\times_S S'$ are surjective immersions, and hence
by~\Ref{IX.1.9} induce equivalences from the
categories of preschemes that are \'etale over~$S'\times_S S'$
\resp $S'\times_S S'\times_S S'$ to the category of
preschemes that are \'etale over~$S'$, it follows from the standard consequences of
descent~\cite{IX.D} that every~$X'$ \'etale over~$S'$ is endowed with
one and only one descent datum relative to $g\colon
S'\to S$. Thus~\Ref{IX.3.3} implies that the inverse-image functor
under~$g$, from the category of preschemes that are \'etale over~$S$
to the category of preschemes that are \'etale over~$S'$, is
\emph{fully faithful}. It remains to show that it is
essentially surjective, \ie that every~$X'$ \'etale over~$S'$ is
isomorphic to the inverse image of an~$X$ \'etale over~$S$. Since the
question is plainly local on~$S$ \emph{and on}~$X'$,
we may suppose that $S$,~$S'$,~$X'$ are affine. But then~$X'$ is
separated and of finite type over~$S'$, and we may apply the
effectivity criterion~\Ref{IX.4.7}.

\begin{corollaire}
\label{IX.4.11}
The conclusion of~\ifthenelse{\boolean{orig}}{\Ref{IX.4.9}}{\Ref{IX.4.10}} remains valid if the hypothesis
on~$g$ is replaced by: $g$~is faithfully flat, quasi-compact, and radicial.
\end{corollaire}

The same proof applies, invoking~\Ref{IX.4.1} instead
of~\Ref{IX.4.7}.

Note that the proof of~\Ref{IX.4.7} is
\og elementary\fg in that it does not use the theorems
\ifthenelse{\boolean{orig}}
{}
{on}
% correction: insert `de'.
finiteness and comparison for proper morphisms (EGA~III~3,~4,~5). The same is not true of the following result:

\begin{theoreme}
\label{IX.4.12}
Let $g\colon S'\to S$ be a proper and surjective morphism of finite
presentation (the latter hypothesis being a
consequence of the first if~$S$ is locally
Noetherian). Then~$g$ is an effective descent morphism
for the fibered category of \'etale coverings of
preschemes.
\end{theoreme}

By~\Ref{IX.3.3} and~\Ref{IX.2.2}, we are reduced to
proving that for every \'etale covering~$X'$ over~$S'$, endowed
with a descent datum relative to $g\colon S'\to S$,
this descent datum is effective. Using~\Ref{IX.4.3},
we reduce easily
% page: 242
\marginpar{242}
%
to the case where~$S$ is Noetherian, and using~\Ref{IX.4.6}, we
may therefore suppose that~$S$ is the spectrum of a
\emph{complete} Noetherian local ring~$A$. Introduce~$S''$ and~$S'''$ as
usual, let~$S_0$ be the spectrum of the residue field of~$A$, and
let~$S_0'$,~$S_0''$,~$S_0'''$ be deduced from~$S'$,~$S''$,~$S'''$
by the base change $S_0\to S$, \ie the fibers
of~$S'$,~$S''$,~$S'''$ at the closed point of~$S$.
By~\Ref{IX.1.10}, the morphisms $S_0\to S$, $S_0'\to S'$,
etc. induce equivalences from the category of
\'etale coverings over the target scheme to the category
of \'etale coverings over the source scheme. Consequently,
$g\colon S'\to S$ is a strict descent morphism for the
fibered category of \'etale coverings of
preschemes if and only if $g_0\colon S_0'\to S_0$ is so,
which is indeed the case by~\Ref{IX.4.1}. This completes the
proof of~\Ref{IX.4.12}. (In this argument, from~\Ref{IX.1.10} we
needed only the fact that the functor considered
in~\Ref{IX.1.10} is \emph{fully faithful}, which does
\emph{not} use the existence theorem for coherent sheaves
in algebraic geometry.)
\end{filecontents*}
% Original input: drafts/SGA1_IX_5_English.texfrag
\begin{filecontents*}[overwrite,noheader]{sga1-cumulative-input-032.tex}
\section{Translation in terms of the fundamental group}
\label{IX.5}

Let
$$
g\colon S'\to S
$$
be an \emph{effective descent morphism} for the fibered category of
\emph{\'etale coverings} of preschemes, for example a proper,
surjective morphism of finite presentation \eqref{IX.4.12}, or a
faithfully flat and quasi-compact morphism. Introducing~$S''$,~$S'''$
as usual, and denoting by~$\cal{C}$,~$\cal{C}'$,~$\cal{C}''$,
and~$\cal{C}'''$ the categories of \'etale coverings
of~$S$,~$S'$,~$S''$, and~$S'''$, respectively, we thus have a
$2$-exact diagram of categories
\begin{equation*}
\label{eq:IX.5.*}
\tag{$*$} 
\xymatrix@C=.5cm{
{\cal{C}} \ar[r]^-{~p^*} & {\cal{C}'}
\ar@<2pt>[rr]^{p_1^*,p_2^*} \ar@<-2pt>[rr] && {\cal{C}''}
\ar@<4pt>[rrr]^{p_{21}^*,p_{32}^*,p_{31}^*} \ar[rrr] \ar@<-4pt>[rrr] &&&
{\cal{C}'''}
}
\end{equation*}
corresponding to the diagram
$$
\xymatrix@C=.5cm{ S & \ar[l]_-{~p} S' && \ar@<2pt>[ll] \ar@<-2pt>[ll]_-{p_1,p_2}
S'' &&& \ar@<4pt>[lll] \ar[lll] \ar@<-4pt>[lll]_{p_{21},p_{32},p_{31}} S'''.
}
$$
Suppose
% page: 243
\marginpar{243}
%
that the preschemes~$S$,~$S'$,~$S''$, and~$S'''$ are disjoint
unions of connected preschemes, as will be the case in particular if
they are locally connected preschemes, and a fortiori if they are
locally noetherian (for example, if~$S'$ is of finite type over the
locally noetherian prescheme~$S$). The
categories~$\cal{C}$,~$\cal{C}'$ ... in~\eqref{eq:IX.5.*} are then
multigalois categories \ifthenelse{\boolean{orig}}{V~}{}\Ref{V.9}, each described by a collection of
compact totally disconnected topological groups, namely the
fundamental groups of the connected components of the
preschemes~$S$,~$S'$,~$S''$, and~$S'''$. For simplicity, we suppose
that~$S$ is connected, and we shall then give a procedure for computing
its fundamental group in terms of the fibered category formed
from~$\cal{C}'$,~$\cal{C}''$, and~$\cal{C}'''$, suitably made
explicit by means of the fundamental groups describing these
categories. The reader will note that the procedure outlined is in
fact valid in the general framework of multigalois categories (which
need not arise from given preschemes~$S$,~$S'$,~$S''$, and~$S'''$).
It is moreover the analogue of the well-known procedure for computing
the fundamental group of a topological space~$S$ that is a locally
finite union of closed subspaces~$S_i$ (or an arbitrary union of open
subspaces~$S_i$), by means of the fundamental groups of the components
\ifthenelse{\boolean{orig}}
{of the~$X_i$}
{of the~$S_i$}
%correction: X_i -> S_i
and of the components of the $S_i\cap S_j$. Of course, the analogous
situation for preschemes does fall within the general framework of
descent, by introducing the prescheme~$S'$ that is the disjoint union
of the~$S_i$ and the canonical morphism
$g\colon S'\to S$.

Set
$$
E'=\pi_0(S'), \quad E''=\pi_0(S''), \quad E'''=\pi_0(S'''),
$$
where~$\pi_0$ denotes the \og set of connected components\fg\
functor. Since the fiber products of~$S'$ over~$S$ form a simplicial
object of~$\Sch$, it is transformed by the functor~$\pi_0$ into a
simplicial set of which~$E'$,~$E''$, and~$E'''$ are the components
in dimensions~$0$,~$1$, and~$2$. We shall have to use the simplicial
maps
$$
q_i=\pi_0(p_i),\quad(i=1,2)\quad\text{and}
\quad q_{ij}=\pi_0(p_{ij}),\quad (i,j)=(2,1),(3,2),(3,1),
$$
displayed in the diagram
\begin{equation*}
\label{eq:IX.5.1}
\tag{1} {\xymatrix@C=.5cm{ E' && \ar@<2pt>[ll] \ar@<-2pt>[ll]_{q_1,q_2} E'' &&&
\ar@<4pt>[lll] \ar[lll] \ar@<-4pt>[lll]_{q_{21},q_{32},q_{31}} E'''. }}
\end{equation*}
The elements
% page: 244
\marginpar{244}
%
of~$E'$ will be denoted with one prime, as in~$s'$, while those
of~$E''$ \resp $E'''$ will be denoted with~$''$ \resp with~$'''$.
The fact that~$S$ is connected is expressed by $\pi_0(K)=0$, where~$K$
is the simplicial set defined by $g\colon S'\to S$, or equivalently by
the fact that the equivalence relation on~$E'$ generated by the pair
of maps $(q_1,q_2)$ is transitive.

We shall choose once and for all an element~$s_0'$ in~$E'$, and, for
every~$s'$ in~$E'$, an element~$\overline{s'}\in E''$ such
that\footnote{Care must be taken here: the element~$\overline{s'}$,
whose existence is implicitly assumed, does not exist in every case.
Consequently, Theorem~\Ref{IX.5.1}, as stated, does not apply in every
case. It is not difficult, however, by following the ideas of the
written text, to modify the statement of this theorem so that it gives
a computational method applicable in every case. In particular, the
corollaries of the said theorem remain valid as stated.}
$$
q_1(\overline{s'})=s_0', \quad q_2(\overline{s'})=s',
$$
thus exhibiting the connectedness of~$S$. For every~$s'\in E'$,
choose a geometric point~$\underline{s}'$ in the connected
component~$s'$ of~$S'$; this point will in fact enter through the
corresponding fiber functor~$F_{s'}'$ on the multigalois
category~$\cal{C}'$. The automorphism group of this functor, \ie the
fundamental group of~$S'$ at~$\underline{s}'$, will be
denoted~$\pi_{s'}$. We similarly choose~$\underline{s}''$ and
$\underline{s}'''$, hence functors~$F_{s''}''$ and~$F_{s'''}'''$,
and thus fundamental groups~$\pi_{s''}$ and~$\pi_{s'''}$. Hence
$$
\pi_{s'}=\pi_1(S',\underline{s}'), \quad
\pi_{s''}=\pi_1(S'',\underline{s}''), \quad
\pi_{s'''}=\pi_1(S''',\underline{s}''').
$$
For every $s''\in E''$, $p_1(\underline{s}'')$ lies in the same
connected component as $\underline{q_1(s'')}$, so there exists an
isomorphism of functors
\ifthenelse{\boolean{orig}}
{$F''_{s''}\circ p_1^* F'_{s'}$}
{$F''_{s''}\circ p_1^*\isomto F'_{q_1(s'')}$}
% correction: $F''_{s''}\circ p_1^*\isomto F'_{s'}$
(\ie a \og path class\fg\ from~$p_1(\underline{s}'')$ to
$\underline{q_1(s'')}$). This observation applies likewise to~$q_2$
and to the~$q_{ij}$. Choose all these path classes:
$$
F''_{s''}\circ p_i^*\isomto F'_{q_i(s'')},\quad F_{s'''}'''\circ
p_{ij}^*\isomto \ifthenelse{\boolean{orig}}{F'''_{q_{ij}(s''')}}{F''_{q_{ij}(s''')}},
$$
(for
% page: 245
\marginpar{245}
%
$i=1,2$ and $(i,j)=(2,1),(3,2),(3,1)$). In particular, this yields
group homomorphisms:
\begin{equation*}
\label{eq:IX.5.2}
\tag{2} {q_i^{s''}\colon \pi_{s''}\to\pi_{q_i(s'')},\quad
q_{ij}^{s'''}\colon \pi_{s'''}\to\pi_{q_{ij}(s''')},}
\end{equation*}
(with the same values of~$i$ and $(i,j)$). Finally, recall that the
structure of a cloven category with
fibers~$\cal{C}'$,~$\cal{C''}$, and~$\cal{C'''}$ also includes
isomorphisms of functors:
$$
p_{21}^*p_1^*\isomto p_{31}^*p_1^*,\quad p_{21}^*p_2^*\isomto
p_{32}^*p_1^*,\quad p_{31}^*p_2^*\isomto p_{32}^*p_2^*,
$$
deduced from isomorphisms of the two sides, respectively, with the
$u_i^*$~($i=1,2,3$), where the~$u_i$ are the three projections
from~$S'''$ to~$S'$. When these data are made explicit, one obtains,
for every~$s'''$, a well-determined element
\begin{equation*}
\label{eq:IX.5.3}
\tag{3} {a_i^{s'''}\in\pi_{v_i(s''')},}
\end{equation*}
(where the~$v_i$,~$i=1,2,3$, are the three maps $E'''\to E'$ defined
by $v_i=\pi_0(u_i)$), which moreover satisfies the conditions:
$$
q_1^{s_1''}q_{21}^{s'''}=\mathrm{int}(a_1^{s'''})
q_1^{s_2''}q_{31}^{s'''}\qquad (s_1''=q_{21}(s'''),\
s_2''=q_{31}(s''')),
$$
and the two analogous conditions involving~$a_2$ and~$a_3$. The
reader will moreover note that the
data~\eqref{eq:IX.5.1},~\eqref{eq:IX.5.2}, and~\eqref{eq:IX.5.3}
allow one to reconstruct, up to an equivalence of fibered categories,
the fibered category under consideration with
fibers~$\cal{C}'$,~$\cal{C}''$, and~$\cal{C}'''$. They must therefore,
in principle, allow one to reconstruct~$\cal{C}$ up to equivalence,
and hence its fundamental group up to isomorphism. In fact, we shall
determine the fundamental group at the geometric point
$p(\underline{s}_0')$ of~$S$, \ie the automorphism group
of~$F'_{s_0'}\circ p^*$.

Note that the datum of an object~$X'$ of~$\cal{C}'$ is
essentially equivalent to the datum of finite
sets~$X'_{s'}$ ($s'\in E'$) on which the~$\pi_{s'}$ act continuously. A
% page: 246
\marginpar{246}
%
gluing datum on such an object then amounts to the datum, for every
$s''\in E''$, of a bijection
$$
\varphi_{s''}\colon X'_{q_1(s'')}\isomto X'_{q_2(s'')}
$$
compatible with the actions of~$\pi_{s''}$, acting on the two sides
through the homomorphisms $q_i^{s''}\colon
\pi_{s''}\to\pi_{q_i(s'')}$. Taking first the~$s''$ of the
form~$\overline{s'}$, one sees that such a datum defines bijections
$$
\psi_{s'}\colon X'_{s_0'}=F_0'(X')\to X_{s'}'
$$
which allow the~$X'_{s'}$ to be identified with the same set
$F_0'(X')=X'_{s_0'}$, on which all the groups~$\pi_{s'}$ will
henceforth act. With this understood, the bijections~$\varphi_{s''}$
correspond to bijections
$$
g_{s''}\colon F_0'(X')\isomto F_0'(X'),
$$
subject, on the one hand, to the commutation relations
with~$\pi_{s''}$:
\begin{equation*}
\label{eq:IX.5.a}
\tag*{a)} {g_{s''}q_1^{s''}(g'')=q_2^{s''}(g'')g_{s''}\qquad (s''\in E'',\ g''\in\pi_{s''}),}
\end{equation*}
and, on the other hand, to the relations
\begin{equation*}
\label{eq:IX.5.b}
\tag*{b)} {g_{\overline{s'}}=g_{\overline{s_0'}}\qquad (s'\in E'),}
\end{equation*}
which express the manner in which we identified the~$X\ifthenelse{\boolean{orig}}{}{'}_{s'}$ with
one another. On making explicit the condition for such a gluing datum
to be a descent datum, one obtains the relations
\begin{equation*}
\label{eq:IX.5.c}
\tag*{c)} {a_3^{s'''}g_{q_{31}(s''')}a_1^{s'''}=
g_{q_{32}(s''')}a_2^{s'''}g_{q_{21}(s''')}\qquad (s'''\in E''').}
\end{equation*}

This gives an equivalence between the category of objects
of~$\cal{C}'$ equipped with a
% page: 247
\marginpar{247}
%
descent datum and the category of finite sets on which the
groups~$\pi_{s'}$ act continuously, equipped in addition with
bijections~$g_{s''}$ satisfying the
relations~\Ref{eq:IX.5.a},~\Ref{eq:IX.5.b}, and~\Ref{eq:IX.5.c}.
Let~$G$ be the group generated by the groups~$\pi_{s'}$ and the new
generators~$g_{s''}$, subject to the
relations~\Ref{eq:IX.5.a},~\Ref{eq:IX.5.b}, and~\Ref{eq:IX.5.c}, and
let~$\pi$ be the inverse-limit group of the quotients of~$G$ by the
subgroups of finite index whose inverse images in the
groups~$\pi_{s'}$ are open subgroups. One also says that~$\pi$ is the
\emph{group of Galois type generated by the~$\pi_{s'}$ and
the~$g_{s''}$, subject to the
relations~\Ref{eq:IX.5.a},~\Ref{eq:IX.5.b}, and~\Ref{eq:IX.5.c}}.
It is immediately seen that the category under consideration is also
equivalent to the category of finite sets on which the topological
group~$\pi$ acts continuously. This proves the following statement:

\begin{theorem}
\label{IX.5.1}
Let $g\colon S'\to S$ be a morphism of preschemes which is an
effective descent morphism for the fibered category of \'etale
coverings of preschemes (\cf \Ref{IX.4.9} and~\Ref{IX.4.12}). Suppose
that~$S$ is connected, and that~$S'$, its fiber square~$S''$, and its
fiber cube~$S'''$ are disjoint unions of connected preschemes (as is
the case, for example, if~$S'$ is of finite type over a locally
noetherian connected~$S$). Choose as above: a geometric point in every
connected component of~$S'$,~$S''$, and~$S'''$, certain path classes,
an $s_0'\in E'$, and, for every $s'\in E'$, an
\ifthenelse{\boolean{orig}}
{$s'\in E''$}
{$s''\in E''$}
% correction: $s''\in E''$
whose two images in~$E'$ are~$s_0'$ and~$s'$.
($E'$,~$E''$, and~$E'''$ denote respectively the sets of connected
components of~$S'$,~$S''$, and~$S'''$.) Then the fundamental group
of~$S$ at the geometric point that is the image of~$s_0'$ is
canonically isomorphic to the group of Galois type generated by the
$\pi_{s'}=\pi_1(S',\underline{s}')$ ($s'\in E'$) and
generators~$g_{s''}$ ($s''\in E''$), subject to the above
relations~\Ref{eq:IX.5.a},~\Ref{eq:IX.5.b}, and~\Ref{eq:IX.5.c},
which involve the elements of the groups
$\pi_{s''}=\pi_1(S'',\underline{s}'')$, and the
elements~$a_i^{s'''}$ ($i=1,2,3$, $s'''\in E'''$) introduced above.
\end{theorem}

\begin{corollary}
\label{IX.5.2}
Suppose that~$S'$ and~$S''$ have only finitely many connected
components, and that the fundamental groups of the connected
components of~$S'$ are topologically finitely generated. Then the
fundamental group of~$S$ is topologically finitely generated.
\end{corollary}

Thus,
% page: 248
\marginpar{248}
%
we shall prove later that the fundamental group of a normal projective
scheme over an algebraically closed field is topologically finitely
generated. Using Chow's lemma and the normalization of algebraic
schemes, it will follow that the same result holds for every proper
scheme over an algebraically closed field.

\begin{corollary}
\label{IX.5.3}
Suppose that~$S'$,~$S''$, and~$S'''$ have only finitely many connected
components, that the fundamental groups of the connected components
of~$S'$ are topologically \emph{finitely presented}, and that the
fundamental groups of the connected components of~$S''$ are
topologically \emph{finitely generated}. Then the fundamental group
of~$S$ is topologically \emph{finitely presented}.
\end{corollary}

Note that~\ifthenelse{\boolean{orig}}{\Ref{IX.4.9}}{\Ref{IX.4.10}} (restricted to \'etale \emph{coverings}) may
be expressed by saying that \emph{a finite radicial surjective
morphism of noetherian preschemes induces an isomorphism of fundamental
groups}; figuratively speaking, one may therefore say that the
fundamental group is a \emph{topological invariant} for preschemes.
More generally, using~\Ref{IX.5.1}, one can make explicit the effect on
the fundamental group of operations on preschemes, such as the
\og pinching\fg of a prescheme along a finite set of points, which has
a simple topological meaning. For example, one obtains:

\begin{corollary}
\label{IX.5.4}
Let $g\colon S'\to S$ be a finite morphism of finite presentation,
and let $T$ be a discrete subset of~$S$. For every $s\in S$, let
$n(s)$ be the \og geometric number of points\fg in the fiber
$g^{-1}(s)$ (which may also be described as the separable degree of
$g^{-1}(s)$ over~$\kres(s)$, the sum of the separable degrees of its
residue extensions). Suppose that $n(s)=1$ for $s\in S-T$. For every
$s\in T$, let $K_s$ be an algebraically closed extension of
$\kres(s)$, let $I_s$ be the set of geometric points of~$S'$ with
values in~$K_s$ (a set with $n(s)$ elements), let $I_s'$ be the
complement of a chosen point of~$I_s$, and finally let $I'$ be the
union of the~$I_s'$. Suppose that $S'$ is connected. Then the
fundamental group of~$S$ is isomorphic to the group of Galois type
generated by the fundamental group of~$S'$ and generators $g_i$
($i\in I'$), subject to no additional relations.
\end{corollary}

The
% page: 249
\marginpar{249}
%
details of the proof are left to the reader; the resulting statement
is merely the translation into the language of group theory of the
fact that the category $C$ of \'etale coverings of~$S$ is equivalent
to the category of \'etale coverings $X'$ of~$S'$, \emph{equipped},
for every $s\in T$, with a transitive system of bijections between the
$n(s)$ fibers of~$X'$ at the points of~$g^{-1}(s)$ with values
in~$K_s$. (In this intrinsic form, of course, it is no longer
necessary to suppose that $S'$ is connected.)

\begin{example}
\label{IX.5.5}
It is easy to prove that the rational curve $\PP^1_k$ over an
algebraically closed field $k$ is simply connected\footnote{\Cf
Exp\ptbl \ifthenelse{\boolean{orig}}{XI~}{}\Ref{XI.1.1}.}. Hence the fundamental group of a complete
rational curve having exactly one double point, with $n$ analytic
branches, is the free group of Galois type generated by $n-1$
generators. For example, in the case of an ordinary double point, one
obtains the fundamental group $\hat\ZZ$, as announced in
(\ifthenelse{\boolean{orig}}{I~}{}\Ref{I.11}~a)). By contrast, the presence of a cusp (which is a
\og geometrically unibranch\fg point) has no effect on the fundamental
group.
\end{example}

\begin{corollaire}
\label{IX.5.6}
Let $g\colon S'\to S$ be a universally submersive morphism of
preschemes with geometrically connected fibers, where $S$ is connected.
Then $S'$ is connected, and, on choosing a geometric point $s'$ of~$S'$
and denoting its image in~$S$ by~$s$, the homomorphism
$$
\pi_1(S',s')\to \pi_1(S,s)
$$
is \emph{surjective}. If $g$ is an effective descent morphism for the
fibered category of \'etale coverings of preschemes
(\cf \emph{\Ref{IX.4.12}}), introduce the geometric point
$s^{\prime\prime}=\diag(s')$ of~$S^{\prime\prime}=S'\times_S S'$ and
the two homomorphisms
$$
p_{1*},p_{2*}\colon \pi_1(S^{\prime\prime},s^{\prime\prime}) \to
\pi_1(S',s')
$$
induced by the two projections. Then $\pi_1(S,s)$ is isomorphic to the
cokernel of this pair of homomorphisms in the category of groups of
Galois type, \ie to the quotient of~$\pi_1(S',s')$ by the closed normal
subgroup generated by the
% page: 250
\marginpar{250}
%
elements of the form
$p_{1*}(g^{\prime\prime})p_{2*}(g^{\prime\prime})^{-1}$, with
$g^{\prime\prime}\in \pi_1(S^{\prime\prime},s^{\prime\prime})$.
\end{corollaire}

Indeed, by~\Ref{IX.3.4}, the functor $X\mto X\times_SS'$ from \'etale
coverings of~$S$ to \'etale coverings of~$S'$ is fully faithful, which is
equivalent to the homomorphism on fundamental groups being an
epimorphism (\ifthenelse{\boolean{orig}}{V~}{}\Ref{V.6.9}). The last assertion follows immediately
from the description in~\Ref{IX.5.1}.

\begin{remarque}
\label{IX.5.7}
It is not known at present whether the fundamental group of a proper
scheme over an algebraically closed field $k$ is topologically finitely
presented\footnote{This seems very unlikely in the case of smooth curves
of genus $g\geq 2$ in characteristic $p>0$. On the other hand, when
$\pi_1$ is replaced by its largest prime-to-$p$ quotient, it appears
that well-known techniques give an affirmative answer, even without
the properness hypothesis. \Cf forthcoming work of J.P\ptbl Murre.}.
Using~\Ref{IX.5.3}, a well-known hyperplane-section technique, and the
resolution of singularities of normal surfaces, one is reduced to the
case of a \emph{smooth surface over}~$k$. This at least makes it
possible to show, by transcendental methods, that the answer is
affirmative in characteristic~$0$ (without having to assume the
triangulability of singular algebraic varieties). In characteristic
$p>0$, the principal difficulty seems to lie in the case of curves,
whose fundamental group is known only to be a quotient of the group
arising in the classical case (\cf the following expos\'e), while the
kernel by which one takes the quotient is very poorly understood.
\end{remarque}

\begin{remarque}
\label{IX.5.8}
One could make explicit other special cases besides~\Ref{IX.5.4}
and~\Ref{IX.5.6} in which~\Ref{IX.5.1} takes a particularly simple
form. An interesting case is that in which $S$ is the quotient of~$S'$
by a finite group of automorphisms~$\Gamma$. \emph{Then the category
of \'etale coverings of~$S\ifthenelse{\boolean{orig}}{'}{}$ is equivalent to the category of \'etale
coverings $X'$ of~$S'$ on which the group $\Gamma$ acts compatibly
with its action on~$S'$, in such a way that for every $s'\in S'$ and
every $g\in \Gamma_{s'}$} (where $\Gamma_{s'}$ denotes \emph{the
inertia group} of~$s'$ in~$\Gamma$), \emph{$g$ acts trivially on the
fiber}~$X_{s'}^\prime$. If $S'$ is connected, this statement is
interpreted as follows. Let $C'_0$ be the category of \'etale coverings
of~$S'$ on which $\Gamma$ acts compatibly with its action on~$S'$
(but without necessarily satisfying the preceding condition on the
inertia groups of the points of~$S'$). One sees easily that this is a
Galois category (\ifthenelse{\boolean{orig}}{V~}{}\Ref{V.5}), and that for
% page: 251
\marginpar{251}
%
every geometric point $a'$ of~$S'$, the fiber functor
$X'\mto X_{a'}'$ on~$C'_0$ is a fundamental functor. Let
$\pi_1(S',\Gamma;a')=G$ be the group of automorphisms of this functor,
endowed with its usual topology. There is then a canonical exact
sequence
$$
e\to \pi_1(S',a')\to G\to \Gamma\to e
$$
(a special case of (\ifthenelse{\boolean{orig}}{V~}{}\Ref{V.6.13}), where one takes for $S$ the
trivial covering $S'\times \Gamma$ of~$S'$ defined by~$\Gamma$, on
which $\Gamma$ is made to act in the evident way). Moreover, for every
geometric point $b'$ of~$S'$, there is an isomorphism
$\pi_1(S', \Gamma ; b' )\to G = \pi_1(S', \Gamma; a')$, defined up to
an inner automorphism arising from~$\pi_1(S',a')$; and since
$\Gamma_{b'}$ maps in the evident way into the first group, one obtains
a homomorphism
$$
u_{b'}\colon \Gamma_{b'}\to G \quoi,
$$
defined up to an inner automorphism (arising
\ifthenelse{\boolean{orig}}
{from~$\pi_1(S',a')$,}
{from~$\pi_1(S',a')$),}
% correction: from~$\pi_1(S',a')$),
whose composite with the canonical homomorphism $G\to\Gamma$ is the
canonical inclusion $\Gamma_{b'}\to\Gamma$. With this understood,
\emph{the fundamental group $\pi_1(S,a)$ \emph{is} canonically
isomorphic to the quotient of~$G=\pi_1(S',\Gamma;a')$ by the closed
normal subgroup generated by the images of the homomorphisms
$\Gamma_{b'}\to G$}. In particular, the image of~$\pi_1(S',a')$ in
$\pi_1(S,a)$ is a normal subgroup, and the corresponding quotient is
isomorphic to a quotient of~$\Gamma$. One can moreover reduce the
number of \og relations\fg introduced by considering, for every
$g\in\Gamma$, $g\neq e$, the coincidence subprescheme $S'_g$ of the
automorphisms \ifthenelse{\boolean{orig}}{$\id_S$ and $g$ of~$S$}{$\id_{S'}$ and $g$ of~$S'$}, choosing a geometric point
$b'_{g,i}$ in each connected component of~$S'_g$, and then choosing
one of the corresponding homomorphisms
$\pi_1(S',\Gamma;b'_{g,i})\to G$, thereby obtaining lifts
$\bar{g}_i$ of~$g$ in~$G$. It then suffices to take the quotient of~$G$
by the closed normal subgroup of~$G$ generated by the~$\bar{g}_i$.

When $a'$ is fixed by $\Gamma$, one sees easily that $\Gamma$ acts
naturally on~$\pi_1(S',a')$, and $G$ identifies with the corresponding
semidirect product. Identifying~$\Gamma$
\ifthenelse{\boolean{orig}}
{to}
{with}
% correction: \`a
a subgroup of~$G$, one sees that among the relations
% page: 252
\marginpar{252}
%
introduced above, on taking $b'=a'$, one finds \og $g=e$\fg for
$g\in \ifthenelse{\boolean{orig}}{I}{\Gamma}$. Thus \emph{if $S'$ has a geometric point $a'$ fixed by}~$\ifthenelse{\boolean{orig}}{I}{\Gamma}$
(\ie a point $s'$ whose inertia group is $\Gamma$), \emph{then
$\pi_1(S,a)$ is a quotient group of the Galois-type quotient group
of~$\pi_1(S',a')$ obtained by \emph{\og making trivial\fg} the actions
of~$\Gamma$ on~$\pi_1(S',a')$; and it is even isomorphic to this latter
group if one assumes that for every $g\in \ifthenelse{\boolean{orig}}{G}{\Gamma}$, the inertia locus $S_g'$
is connected, and therefore contains the underlying point of~$a'$}.
Indeed, this last assertion is contained in the second description
given above of the relations to be imposed on~$G$.

This last result applies in particular if one takes for $S'$ the
Cartesian power $X^n$ of a connected prescheme over an algebraically
closed field, and for $\Gamma$ the symmetric group
$\Gamma={\goth {S}}_n$, acting in the usual way, so that $S$ is the
symmetric power of degree~$n$ of~$X$. Taking for $a'$ a geometric point whose
underlying point lies on the diagonal, one is under the preceding
conditions, since all the inertia loci $S_g'$ contain the diagonal.
Using the fact, proved in the following expos\'e, that if $X$ is proper
and connected over~$k$, the fundamental group of~$X^n$ identifies
with $\pi_1(X)^n$, one obtains the following amusing result:
\emph{If $X$ is proper and connected over an algebraically closed
field~$k$, the fundamental group of its symmetric power of degree~$n$,
$n\geq 2$, is isomorphic to the fundamental group of~$X$ made
abelian}. (I do not know whether the analogous fact in algebraic
topology is known; it should be possible to establish it by the same
descent method.) Take, for example, the rational curve
$X=\PP^1_k$; one obtains yet another proof of the fact that $\PP^r_k$ is
simply connected, using the fact that $\PP^1_k$ is so. Now take for
$X$ a nonsingular curve over~$k$, and let $n\geq 2g-1$, so that
$\Symm^n(X)$ is fibered over the Jacobian~$J$, with projective spaces
as fibers, and hence (as will be seen using the results of the next
two expos\'es) has the same fundamental group as~$J$. One thereby
recovers, without d\'evissage, the well-known fact that \emph{the
fundamental group of the Jacobian of~$X$ is isomorphic to the
fundamental group of~$X$ made abelian}.
\end{remarque}
\end{filecontents*}
% Original input: drafts/SGA1_IX_6_A_English.texfrag
\begin{filecontents*}[overwrite,noheader]{sga1-cumulative-input-033.tex}
\section[A fundamental exact sequence]{A fundamental exact sequence. Descent by morphisms with relatively connected fibers}
\label{IX.6}
% page: 253
\marginpar{253}
%

\begin{theorem}
\label{IX.6.1}
Let \(S\) be the spectrum of an Artinian ring \(A\) with residue field
\(k\), let \(\overline k\) be an algebraic closure of~\(k\), let
\(X\) be an \(S\)-prescheme, set
\[
  X_0=X\otimes_Ak,
  \qquad
  \overline X_0=X\otimes_A\overline k,
\]
and let \(\overline a\) be a geometric point of~\(\overline X_0\),
with image \(a\) in~\(X\) and image \(b\) in~\(S\). Suppose that
\(X_0\) is quasi-compact and \emph{geometrically connected} over~\(k\).
(If \(X\) is proper over~\(S\), this means that
\(\H^0(X_0,\cal O_{X_0})\) is a \emph{local} Artinian ring whose
residue field is \emph{radicial} over~\(k\).) Then the sequence of
canonical homomorphisms
\[
  e\mto\pi_1(\overline X_0,\overline a)
  \longrightarrow\pi_1(X,a)
  \longrightarrow\pi_1(S,b)\longrightarrow e
\]
is exact, and
\[
  \pi_1(S,b)\isomfrom\pi_1(k,\overline k)
  =\text{the Galois group of \(\overline k\) over~\(k\)}.
\]
\end{theorem}

Since fundamental groups are unchanged on killing nilpotent elements,
we may suppose that \(A=k\), which already makes the last isomorphism
evident. Let \(k'\) be the separable closure of~\(k\) in
\(\overline k\), consider \(X'=X\otimes_kk'\), and let \(a'\) be the
image of~\(\overline a\) in~\(X'\). There is a sequence of canonical
homomorphisms
\[
  e\mto\pi_1(\overline X_0,\overline a)
  \longrightarrow\pi_1(X',a')
  \longrightarrow\pi_1(S',b')\longrightarrow e,
\]
where \(S'=\Spec(k')\). The diagram
\[
  \xymatrix{
  S&X\ar[l]&\overline X_0\ar[l]\\
  S'\ar[u]&X'\ar[u]\ar[l]&\overline X_0\ar[l]\ar@{=}[u]
  }
\]
gives a canonical homomorphism from this sequence to the corresponding
sequence for~\(X/k\). By~\Ref{IX.4.11}, this homomorphism of sequences
of groups is an isomorphism. We are therefore reduced to proving that
the second sequence is exact, \ie we may assume that \(k\) is
\emph{perfect}.

Let \(k_i\) run through the finite Galois subextensions of~\(k\) in
\(\overline k\), set \(X_i=X\otimes_kk_i\), and let \(a_i\) be the
image of~\(\overline a\) in~\(X_i\). We leave it to the reader
\ifthenelse{\boolean{orig}}
{\ignorespaces}
{to verify}
% correction: words were missing in the source.
that the natural homomorphism
\[
  \pi_1(\overline X_0,\overline a)
  \isomto\varprojlim_i\pi_1(X_i,a_i)
\]
% page: 254
\marginpar{254}
%
is an isomorphism. This simply means that an \'etale covering of
\(\overline X_0\) comes from an \'etale covering of some~\(X_i\), and
that the latter is essentially unique after passing to an~\(X_j\)
with \(j\geq i\).

Let \(\pi_i\) be the Galois group of~\(k_i\) over~\(k\), \ie the
group opposite to the group of \(S\)-automorphisms of
\(S_i=\Spec(k_i)\). Since the functor
\[
  S'\mto X\times_SS'
\]
from \'etale coverings of~\(S\) to \'etale coverings of~\(X\) is
fully faithful by~\eqref{IX.3.4}, it follows that \(\pi_i\) is also
isomorphic to the group opposite to the group of \(X\)-automorphisms
of the connected principal covering \(X_i\) of~\(X\).
Consequently, \ifthenelse{\boolean{orig}}{V~}{}\Ref{V.6.13} gives an exact sequence
\[
  e\mto\pi_1(X_i,a_i)\longrightarrow\pi_1(X,a)
  \longrightarrow\pi_i\longrightarrow e.
\]
Passing to the inverse limit over~\(i\) in these exact sequences gives
an exact sequence, since we are in the category of groups of Galois
type. This is precisely the sequence considered in~\Ref{IX.6.1}, and
the proof is complete.

The geometric translation of exactness on the right
in~\Ref{IX.6.1} is the following:

\begin{corollary}
\label{IX.6.2}
With the preceding notation, let \(X'\) be an \'etale covering
of~\(X\), and let \(\overline{X'}_0\) be the corresponding \'etale
covering of~\(\overline X_0\). The following conditions are
equivalent:

\begin{enumerate}
\item[(i)] There exist an \(S'\) \'etale over~\(S\) and an
\(X\)-isomorphism \(X'\isomto X\times_SS'\). In that case \(S'\) is
determined up to unique isomorphism by~\Ref{IX.3.4}.

\item[(ii)] The covering \(\overline{X'}_0\) is completely split
over~\(\overline X_0\).

If \(X'\) is connected, these conditions are also equivalent to:

\item[(ii bis)] The covering \(\overline{X'}_0\) has a section
over~\(\overline X_0\).
\end{enumerate}
\end{corollary}

The last qualification is essential. The equivalence of~(i) and~(ii)
means only that \(\pi_1(S,b)\) is the quotient of \(\pi_1(X,a)\) by
the closed normal subgroup
% page: 255
\marginpar{255}
%
generated by the image of
\(\pi_1(\overline X_0,\overline a)\), and not by that image itself.
Under the preceding hypotheses, we shall say that \(X'\) is a
\emph{geometrically trivial} covering
%
\index{geometrically trivial (covering)|hyperpage}%
%
of~\(X\).

\begin{remarque}
\label{IX.6.3}
In~\Ref{IX.6.1}, one cannot replace \(\overline k\) by an arbitrary
algebraically closed extension of~\(k\), even if \(k\) is already
assumed algebraically closed. In other words, it is not true in
general that, if \(X\) is a connected algebraic scheme over an
algebraically closed field~\(k\), its fundamental group is unchanged
when \(k\) is replaced by an algebraically closed extension. This
already fails, for example, in characteristic \(p>0\) for the affine
line over~\(k\), because of phenomena of \emph{higher ramification}
at the point at infinity, which give the fundamental group a
\emph{continuous} structure. In the next expos\'e we shall see,
however, that such phenomena cannot occur if \(X\) is \emph{proper}
over~\(k\). We shall also prove by transcendental methods that the
same holds when \(k\) has characteristic zero.
\end{remarque}

\begin{corollary}
\label{IX.6.4}
Suppose that \(a\) is localized at a point \(x\in X\) rational
over~\(k\), or more generally with residue field radicial over~\(k\).
Then the exact sequence~\Ref{IX.6.1} splits.
\end{corollary}

We may suppose that \(S=\Spec(k)\). If \(x\) is rational over~\(k\),
it corresponds to a section \(S\to X\) of~\(X\) over~\(S\) taking
\(b\) to~\(a\). This section defines a homomorphism
\(\pi_1(S,b)\to\pi_1(X,a)\), which is the required splitting. If
\(\kres(x)\) is radicial over~\(k\), we reduce to the preceding case
by the base extension \(\Spec(\kres(x))\to\Spec(k)\).

\begin{theorem}
\label{IX.6.5}
Let \(f\colon X\to S\) be a proper surjective morphism of finite
presentation with geometrically connected fibers, let \(X'\) be a
prescheme of finite presentation and proper over~\(X\), let \(s\) be
a point of~\(S\), let \(F=X_s\) be the fiber of~\(X\) at~\(s\), and
let \(F_1'\) be a connected component of the fiber \(F'=X'_s\)
of~\(X'\) at~\(s\). There exist an open neighborhood \(X_1'\) of
\(F_1'\) in~\(X'\), an \(S\)-scheme \(S_1'\) \'etale over~\(S\),
and an \(X\)-isomorphism
\[
  X_1'\isomto S_1'\times_SX
\]
if and only if \(X'\) is \'etale over~\(X\) at every point of
\(F_1'\), and \(F_1'\) is a geometrically trivial covering of~\(F\).
\end{theorem}

The
% page: 256
\marginpar{256}
%
necessity of the condition is trivial, so it remains to prove
sufficiency. We reduce readily to the case where \(S\) is noetherian.
Consider the Stein factorization \(X\to T\to S\) of~\(f\), where
\(T\) is the spectrum of the algebra \(f_*(\cal O_X)\) on~\(S\).
Since the fibers of~\(X\) over~\(S\) are geometrically connected and
\(f\) is surjective, the morphism \(T\to S\) is finite, surjective,
and radicial. Hence, by~\eqref{IX.4.10}, every \(T'\) \'etale
over~\(T\) is the inverse image of an \(S'\) \'etale over~\(S\).
This reduces us to
\ifthenelse{\boolean{orig}}
{proving~\Ref{IX.6.5}.}
{proving~\Ref{IX.6.5}}
% correction: remove the period.
after replacing \(S\) by~\(T\), \ie to the case where
\(f_*(\cal O_X)=\cal O_S\). Now consider the Stein factorization
\[
  X'\longrightarrow S'\longrightarrow S
\]
of the proper morphism \(h\colon X'\to S\), where \(S'\) is the
spectrum of the algebra \(h_*(\cal O_{X'})\). The morphisms
\(X'\to X\) and \(X'\to S'\) define a canonical morphism
\[
  X'\longrightarrow X\times_SS',
\]
and our assertion is contained in the following result.

\end{filecontents*}
% Original input: drafts/SGA1_IX_6_B_English.texfrag
\begin{filecontents*}[overwrite,noheader]{sga1-cumulative-input-034.tex}
\begin{corollary}
\label{IX.6.6}
Let \(f\colon X\to S\) be a proper morphism of locally noetherian
preschemes such that \(f_*(\cal O_X)=\cal O_S\), and let \(X'\) be a
prescheme proper over~\(X\). Consider the \emph{Stein} factorization
\[
  X'\longrightarrow S'\longrightarrow S
\]
of \(X'\to S\), and the canonical morphism
\(X'\to X\times_SS'\). Let \(s\) be a point of~\(S\), and let \(s'\)
be a point of~\(S'\) above~\(s\), corresponding to a connected
component \(F_1'\) of the fiber \(X'_s\). The morphism
\(X'\to X\times_SS'\) is an isomorphism over an open neighborhood
\(U'\) of~\(s'\) which is \'etale over~\(S\) if and only if \(X'\)
is \'etale over~\(X\) at the points of~\(F_1'\), and \(F_1'\) is a
geometrically trivial covering of the fiber \(F=X_s\).
\end{corollary}

Necessity is again trivial, so it remains to prove sufficiency. The
conclusion also says that
\begin{enumerate}
\item[a)] the morphism deduced from \(X'\to X\times_SS'\) by the
base change \(\Spec(\widehat{\cal O}_{s'})\to S'\) is an isomorphism;
\item[b)] \(S'\) is \'etale over~\(S\) at~\(s'\), \ie
\(\widehat{\cal O}_{s'}\) is \'etale over
\(\widehat{\cal O}_s\).
\end{enumerate}
In this form, the conclusion is invariant under the base change
\(\Spec(\widehat{\cal O}_s)\to S\). Since the hypotheses are also
stable under this base change, we may suppose that \(S\) is the
spectrum of a complete noetherian local ring. We may moreover plainly
suppose that \(X'\) is connected, which here implies that
\[
  S'=\Spec(\cal O_{s'}),
  \qquad F'=F_1'.
\]

The set
% page: 257
\marginpar{257}
%
of points of~\(X'\) at which \(X'\) is \'etale over~\(X\) is open
and contains the fiber \(X'_{s'}=F'\). Since \(X'\) is proper
over~\(S\), it follows that \(X'\) is \'etale over~\(X\). On
\(F=X_s\), it induces an \'etale covering isomorphic to
\[
  F\otimes_{\kres(s)}L,
\]
where \(L\) is \'etale over~\(\kres(s)\). It follows from
\Ref{IX.1.10} that this covering is isomorphic to one of the form
\(X\times_ST\), with \(T\) \'etale over~\(S\). Here again, it is
enough to use the full faithfulness of the functor
in~\Ref{IX.1.10}; this follows from the fact that a formal
isomorphism of coherent sheaves on~\(X\) comes from an isomorphism
of those sheaves.

Thus, if \(T\) is defined by the finite \(A\)-algebra~\(B\), then
\(X'\) is identified with the spectrum of the algebra
\(\cal O_X\otimes_AB\) on~\(X\). Since
\(f_*(\cal O_X)=\cal O_S\), it follows at once that
\(h_*(\cal O_{X'})\) is defined by~\(B\). Consequently, the canonical
homomorphism \(X'\to X\times_SS'\) is precisely the isomorphism
\(X'\isomto X\times_ST\) under consideration. This completes the
proof.

\begin{corollary}
\label{IX.6.7}
Under the hypotheses of~\Ref{IX.6.5}, there exist a prescheme \(S'\)
\'etale over~\(S\) and an \(X\)-isomorphism
\[
  X'\isomto X\times_SS'
\]
if and only if \(X'\) is \'etale over~\(X\) and, for every
\(s\in S\), \(X'_s\) is a geometrically trivial covering of~\(X_s\).
\end{corollary}

Indeed, if this is so, then \(X'\) is the union of open subsets
\(X_i'\) isomorphic to inverse images of \(S_i'\) \'etale over~\(S\).
It is readily seen that these \(S_i'\) glue to an \(S'\) \'etale
over~\(S\), yielding an isomorphism \(X'\isomto X\times_SS'\).
For example, one may say that the \(X_i'\) carry descent data relative
to \(X\to S\), and that these necessarily glue to a descent datum on
all of~\(X'\) relative to \(X\to S\). Since the latter is effective
on the \(X_i'\), it follows readily, by a sorites omitted from
Section~\Ref{IX.4}, that it is effective. We may also state
\Ref{IX.6.7} in the following form:

\begin{corollary}
\label{IX.6.8}
Let \(f\colon X\to S\) be a proper surjective morphism of finite
presentation with geometrically connected fibers. Then \(f\) is an
effective descent morphism for the fibered category of preschemes
finite and \'etale over other preschemes. The functor
\[
  S'\mto X\times_SS'
\]
induces an equivalence from the category of preschemes finite and
\'etale over~\(S\) to the category of preschemes
% page: 258
\marginpar{258}
%
finite and \'etale over~\(X\) which induce a geometrically trivial
covering on every fiber \(X_s\).
\end{corollary}

\begin{remarque}
\label{IX.6.9}
Let \(f\colon X\to S\) be proper and surjective, with \(S\) locally
noetherian. Then \(f\) factors as a morphism \(X\to S'\) satisfying
the hypothesis of~\Ref{IX.6.8}, followed by a finite surjective
morphism \(S'\to S\) to which~\Ref{IX.4.7} applies. Thus \(f\) is the
composite of two morphisms which are \emph{universal effective descent
morphisms} for the fibered category of preschemes finite and \'etale
over other preschemes. It follows that \(f\) itself is a universal
effective descent morphism for the fibered category in question. This
recovers~\Ref{IX.4.12} by a different method.
\end{remarque}

\begin{remarque}
\label{IX.6.10}
The conclusion of~\Ref{IX.6.7} no longer holds if the hypothesis that
\(f\) is proper is replaced by the following: \(X\) is of finite
type over~\(S\) and has a section over~\(S\). Under this replacement,
\(f\) is universally submersive and is a descent morphism for the
fibered category of preschemes \'etale over other preschemes, but the
conclusion fails even when \(S\) is the spectrum of a discrete
valuation ring and \(X'\) is an \'etale covering of~\(X\).

To see this, start with a \(Z\) proper over~\(S\) whose generic
% remark: or general?
fiber is a nonsingular rational curve and whose special fiber \(Z_0\)
consists of two intersecting lines. For example, if \(t\) is a
uniformizer of the valuation ring~\(A\), take the closed subscheme
\(Z\) of~\(\PP_A^2\) defined by the homogeneous equation
\[
  x^2+y^2+tz^2=0
\]
in homogeneous coordinates \(x,y,z\). Let \(X\) be the complement
of the singular point \(a\) of~\(Z_0\) in the union
\(Z\cup\PP_k^2\). The fibers of~\(X\) are \(\PP_k^1\) and
\(\PP_k^2-a\), and are therefore geometrically simply connected
(\ie every \'etale covering of such a fiber is geometrically
trivial).

Nevertheless,
% page: 259
\marginpar{259}
%
arguing as in Section~\Ref{IX.4}, one readily constructs \'etale
coverings of~\(X\) which do not come from \'etale coverings of~\(S\),
by gluing trivial coverings of \(Z-a\) and \(\PP_k^2-a\). It is
possible, on the other hand, that the conclusion of~\Ref{IX.6.7}
remains valid if properness is replaced by the assumption that \(X\)
is \emph{universally open} and of finite presentation over~\(S\)
\footnote{This has now been proved with \(\ifthenelse{\boolean{orig}}{g}{f}\) merely universally open
and surjective; \cf SGA~4, XV~1.15.}. This is at least true if the
fibers of~\(X\) over~\(S\) are assumed geometrically irreducible, and
not merely geometrically connected. We mention only that this question
may be reduced to the case where \(S\) is the spectrum of a complete
discrete valuation ring with algebraically closed residue field.
\end{remarque}

The interpretation of~\Ref{IX.6.7} in terms of the fundamental group
is the following:

\begin{corollary}
\label{IX.6.11}
Let \(f\colon X\to S\) be a proper surjective morphism of finite
presentation with geometrically connected fibers. Suppose that \(X\),
and hence \(S\), is connected. Let \(a\) be a geometric point of~\(X\)
and \(b\) its image in~\(S\). For every \(s\in S\), choose an
algebraic closure \(\overline{\kres(s)}\) of~\(\kres(s)\), a
geometric point \(a_s\) of~\(X_s\) with values in this extension, and
a path class from~\(a_s\) to~\(a\). This gives a homomorphism
\[
  \pi_1(\overline X_s,a_s)\longrightarrow\pi_1(X,a),
  \qquad
  \overline X_s=X_s\otimes_{\kres(s)}\overline{\kres(s)}.
\]
Then the homomorphism \(\pi_1(X,a)\to\pi_1(S,b)\) is surjective, and
its kernel is the closed normal subgroup of~\(\pi_1(X,a)\) generated
by the images of the groups \(\pi_1(\overline X_s,a_s)\).
\end{corollary}

\begin{remarque}
\label{IX.6.12}
Under the hypotheses of~\Ref{IX.6.7}, with \(S\) noetherian, it is
readily seen that the set of points \(s\in S\) for which \(X'_s\) is
geometrically trivial over~\(X_s\) is constructible. If \(X'\) is
proper over~\(X\), it is even open, as follows from~\Ref{IX.6.6}.
Consequently, if \(S\) is a Jacobson prescheme (for example, of finite
type over a field), or
% page: 260
\marginpar{260}
%
if \(X'\) is proper over~\(X\), it is enough, in verifying the
hypotheses of~\Ref{IX.6.7}, to consider the \emph{closed} points
\(s\) of~\(S\). Likewise, in~\Ref{IX.6.11}, it is then enough to use
the groups \(\pi_1(\overline X_s,a_s)\) for closed points \(s\)
of~\(S\).
\end{remarque}

%\section{Bibliography}
\begin{thebibliography}{D}

\bibitem[D]{IX.D} J\ptbl \textsc{Giraud}, M\'ethode de la
\ifthenelse{\boolean{orig}}
{d\'escente,}
{descente,}
% correction: descente.
M\'emoire \no 2 de la Soc.\ Math.\ Fr., 1964.

\bibitem[1]{IX.1} A\ptbl \textsc{Grothendieck}, G\'eom\'etrie
Alg\'ebrique et G\'eom\'etrie Formelle, S\'eminaire Bourbaki
t\ptbl 11, 1959, \No 182.
\end{thebibliography}
\end{filecontents*}
% Original input: drafts/SGA1_IX_OPENING_AND_1_English.texfrag
\begin{filecontents*}[overwrite,noheader]{sga1-cumulative-input-035.tex}
\chapter[Descent of \'etale morphisms]{Descent of \'etale morphisms.\\ Application to the fundamental group}
\label{IX}
% page: 228
\marginpar{228}
%

\section{Review of \'etale morphisms}
\label{IX.1}

We shall review here the properties of \'etale morphisms (developed
in~\Ref{I}) that we shall use, taking this opportunity to remove the
superfluous noetherian hypotheses from the theory. The reader should
note that even if one is interested only in noetherian schemes, the
technique of descent leads one to introduce nonnoetherian schemes
(such as \(\Spec(\widehat A\otimes_A\widehat A)\), where \(A\) is a
noetherian local ring). In order to use the language of fibered
categories, it is therefore important to define notions such as
\'etale morphisms, \ifthenelse{\boolean{orig}}{etc...}{etc.}, without
introducing noetherian restrictions. A reader reluctant to verify or
accept that the statements below hold without noetherian hypotheses
may simply take them under the noetherian hypotheses of~\Ref{I},
provided that the same noetherian hypotheses are inserted into the
statements of the following sections, and that Definition~\Ref{IX.1.1}
below is used for the nonnoetherian schemes that arise in the
arguments.

\begin{definition}
\label{IX.1.1}
Let \(f\colon X\to S\) be a morphism of preschemes and \(x\) a point
of~\(X\). We say that \(f\) is \emph{\'etale at}~\(x\),
%
\index{etale morphism@\'etale (morphism, algebra)|hyperpage}%
%
or that \(X\) is \'etale over~\(S\) at~\(x\), if there exist an
affine open neighborhood \(U\) of \(s=f(x)\), an affine open
neighborhood \(V\) of~\(x\) above~\(U\), an affine
\emph{noetherian} scheme \(U_0\), an affine \(U_0\)-scheme \(V_0\)
which is \emph{\'etale} \eqref{I}, a morphism \(U\to U_0\), and a
\(U\)-isomorphism
\[
  V\isomto V_0\times_{U_0}U.
\]
\end{definition}

When \(S\)
% page: 229
\marginpar{229}
%
is locally noetherian, this terminology agrees with that of \loccit
Likewise, we say that \(f\) \emph{is \'etale}, or that \(X\)
\emph{is \'etale over}~\(S\), if \(f\) is \'etale at every point
\(x\) of~\(X\). With these definitions, the propositions below
reduce without difficulty to the noetherian case, where they are
proved in I, \No \Ref{I.4}, \Ref{I.5}, and~\Ref{I.7}. For details,
the reader may consult EGA~IV\footnote{More precisely, EGA~IV,
17 and~18.}.

\begin{remarques}
\label{IX.1.2}
If \(f\) is \'etale at~\(x\), then \(f\) is \emph{of finite
presentation at}~\(x\) (\ifthenelse{\boolean{orig}}{VIII~}{}\Ref{VIII.3.5}), the local ring of~\(x\)
in the fiber \(f^{-1}(s)\) is a \emph{finite separable extension} of
\(\kres(s)\), and \(f\) is \emph{flat at}~\(x\). One can prove the
converse, so Definition~\Ref{IX.1.1} is the same as in the case where
\(S\) is locally noetherian, except that the condition \emph{of finite
type at \(x\)} must be replaced by \emph{of finite presentation at
\(x\)}. Since this result is delicate to prove, we have not chosen it
here as the definition of an \'etale morphism; it does not lend
itself directly to proving the properties that follow.
\end{remarques}

First, we have trivially:

\begin{proposition}
\label{IX.1.3}
If \(f\colon X\to S\) is \'etale, then every morphism
\(f'\colon X'\to S'\) deduced from it by a base change \(S'\to S\)
is likewise \'etale.
\end{proposition}

Thus one may say that \'etale morphisms form a \emph{fibered
subcategory} of the category of arrows in~\(\Sch\)
\ifthenelse{\boolean{orig}}
{(\cf VI~\Ref{VI.11}~a)}
{(\cf \ifthenelse{\boolean{orig}}{VI~}{}\Ref{VI.11}~a)}. The purpose of the present expos\'e is to
study the exactness properties of this category fibered over~\(\Sch\).

\begin{proposition}
\label{IX.1.4}
Let \(f\colon X\to S\) be a morphism of preschemes. It is an open
immersion if and only if it is \'etale and radicial.
\end{proposition}

\Cf \ifthenelse{\boolean{orig}}{I~}{}\Ref{I.5.1}. It follows that if \(X\) is \'etale over~\(S\),
every section of~\(X\) over~\(S\) is an open immersion. Applying
\Ref{IX.1.4} once more gives:

\begin{corollary}
\label{IX.1.5}
Let \(X\) be an \'etale \(S\)-prescheme. There is a bijection between
the set of sections of~\(X\) over~\(S\) and the set of open
% page: 230
\marginpar{230}
%
subsets \(\Gamma\) of~\(X\) for which the morphism
\(\Gamma\to S\) induced by the structural morphism is
\emph{radicial} and \emph{surjective}.
\end{corollary}

If \(X\) is separated over~\(S\), then \(\Gamma\) is both open and
closed in~\(X\), but this is immaterial. After an evident base
change, \Ref{IX.1.5} may be put in the following apparently more
general form:

\begin{corollary}
\label{IX.1.6}
Let \(X\) and \(Y\) be two \(S\)-preschemes, with \(Y\) \'etale
over~\(S\). The map \(f\mto\Gamma_f\), which assigns to every
\(S\)-morphism \(f\colon X\to Y\) the underlying subset of the graph
of~\(f\) in~\(X\times_SY\), is a bijection from \(\Hom_S(X,Y)\) onto
the set of open subsets \(\Gamma\) of~\(X\times_SY\) such that the
morphism \(\Gamma\to X\) induced by \(\pr_1\) is \emph{radicial} and
\emph{surjective}.
\end{corollary}

\begin{proposition}
\label{IX.1.7}
Let \(S_0\) be the subprescheme of~\(S\) defined by a quasi-coherent
nil ideal, \ie such that \(S_0\) has the same underlying set as~\(S\).
Then the functor
\[
  X\mto X\times_SS_0
\]
from the category of preschemes \'etale over~\(S\) to the category
of preschemes \'etale over~\(S_0\) is an equivalence of categories.
\end{proposition}

The fact that this functor is fully faithful is an immediate
consequence of~\Ref{IX.1.6}. Its essential surjectivity is contained
in \ifthenelse{\boolean{orig}}{I~}{}\Ref{I.8.3}. Observe that under the preceding equivalence,
\(X\) is of finite type, \ie quasi-finite, over~\(S\) (respectively
finite, \ie an \'etale covering of~\(S\)) if and only if \(X_0\)
satisfies the analogous condition over~\(S_0\). The same observation
applies to separatedness. These facts are immediate, and are also
contained in~\Ref{IX.2.4} below.

\begin{corollary}
\label{IX.1.8}
Let \(A\) be a complete noetherian local ring with residue field
\(k\). Then the functor
\[
  B\mto B\otimes_Ak
\]
is an equivalence from the category of finite \'etale algebras
over~\(A\) to the category of finite \'etale algebras over~\(k\)
(\ie finite products of finite separable extensions of~\(k\)).
\end{corollary}

\begin{proposition}
\label{IX.1.9}
A prescheme \(X\) is an \emph{\'etale covering} of~\(S\), \ie finite
and \'etale over~\(S\), if and only if \(X\) is \(S\)-isomorphic to
the spectrum of an algebra \(\cal A\) on~\(S\) which is finite locally
free as a module and such that, for every
% page: 231
\marginpar{231}
%
\(s\in S\), the algebra
\(\cal A_s\otimes_{\cal O_s}\kres(s)\) is separable over~\(\kres(s)\),
and hence in this case is a finite product of finite separable
extensions of~\(\kres(s)\).
\end{proposition}

Finally, the following result is less elementary, being the
conjunction of \ifthenelse{\boolean{orig}}{I~}{}\Ref{I.8.4} and the \emph{existence theorem for
sheaves in algebraic geometry} (EGA~III, 5; \cf also
\cite{IX.1}, Th.~3).

\begin{theorem}
\label{IX.1.10}
Let \(S\) be the spectrum of a complete noetherian local ring, let
\(X\) be a proper \(S\)-scheme, and let \(X_0\) be the fiber of~\(X\)
over the closed point of~\(S\), so that \(X_0\) is a closed subscheme
of~\(X\). Then the restriction functor
\[
  X'\mto X'\times_XX_0
\]
is an equivalence from the category of \'etale coverings of~\(X\) to
the category of \'etale coverings of~\(X_0\).
\end{theorem}

\end{filecontents*}
% Original input: drafts/SGA1_I_10_1_English_source_draft.texfrag
\begin{filecontents*}[overwrite,noheader]{sga1-cumulative-input-036.tex}
\section{Étale coverings of a normal scheme}
\label{I.10}

\begin{proposition}
\label{I.10.1}
% SGA1-CORR-0009: propagate the accepted corrected-French base scheme.
Let \(X\) be a prescheme that is étale and separated over a connected normal
prescheme~\(Y\) with function field~\(K\). Then the connected components
\(X_i\) of~\(X\) are integral, their function fields \(K_i\) are finite
separable extensions of~\(K\), and each \(X_i\) is identified with a nonempty
open subset of the normalization of~\(Y\) in~\(K_i\) (and hence \(X\) is
identified with a dense open subset of the normalization
of~\(Y\) in \(R(X)=L=\prod K_i\)).
\end{proposition}
\end{filecontents*}
% Original input: drafts/SGA1_I_10_1_PROOF_English_source_draft.texfrag
\begin{filecontents*}[overwrite,noheader]{sga1-cumulative-input-037.tex}
% SGA 1, Expose I, proof of Proposition I.10.1.
% French authority semantic boundary: lines 2044--2049,
% 6 LF lines / 367 bytes; LF-normalized SHA-256:
% BD95FBC222429DF8B699835A75B8C84A2D450E9C65334F51CE8B82673EE9BB46.
% The blank line 2050 closes the source envelope; Corollary I.10.2 begins
% at line 2051 and is excluded.
% SGA1-CORR-0009: the corrected French restores quasi-finiteness, Y as the
% normalization base, and X_i as the scheme with function field K_i.

By Corollary~\ref{I.9.10}, \(X\) is normal; a fortiori its local rings are
integral domains, so the connected components of~\(X\) are irreducible.
Since \(X_i\) is normal, quasi-finite, and dominant over~\(Y\), it follows from a
special case (an almost trivial one, moreover) of the Main Theorem that
\(X_i\) is an open subset of the normalization of~\(Y\) in the field \(K_i\)
of~\(X_i\).
\end{filecontents*}
% Original input: drafts/SGA1_I_10_2_English_source_draft.texfrag
\begin{filecontents*}[overwrite,noheader]{sga1-cumulative-input-038.tex}
% SGA 1, Expose I, Corollary I.10.2 statement only.
% French authority semantic boundary: lines 2051--2057,
% 7 LF lines / 266 bytes; LF-normalized SHA-256:
% 09E280D81EFDC9F8F43E6934AD8CA53A3802323E4EA72ED716A8910C4FC02F2B.
% The blank line 2058 closes the source envelope; the unheaded proof begins
% at line 2059 and is excluded.

\begin{corollary}
\label{I.10.2}
Under the hypotheses of Proposition~\ref{I.10.1}, \(X\) is finite over~\(Y\)
(i.e. an étale covering of~\(Y\)) if and only if \(X\) is isomorphic to the
normalization \(Y'\) of~\(Y\) in \(L=R(X)\) (the ring of rational functions
on~\(X\)).
\end{corollary}
\end{filecontents*}
% Original input: drafts/SGA1_I_10_2_PROOF_English_source_draft.texfrag
\begin{filecontents*}[overwrite,noheader]{sga1-cumulative-input-039.tex}
% SGA 1, Expose I, proof of Corollary I.10.2.
% French authority semantic boundary: lines 2059--2062,
% 4 LF lines / 220 bytes; LF-normalized SHA-256:
% 7F492A882C0C6B690AAEAF417A4A1DF104E6A74CCDC945BE0A73BA1C0EDF92B2.
% The blank line 2063 closes the source envelope; the next definition begins
% at line 2064 and is excluded.
% SGA1-CORR-0010: propagate the accepted corrected-French extension L/K.

Indeed, we know that this normalization is finite over~\(Y\) (\(Y\) being
normal and \(L/K\) separable); conversely, if \(X\) is finite over~\(Y\), then
it is finite over~\(Y'\), so its image in~\(Y'\) is closed; on the other hand,
that image is dense.
\end{filecontents*}
% Original input: drafts/SGA1_I_10_3_DEFINITION_English_source_draft.texfrag
\begin{filecontents*}[overwrite,noheader]{sga1-cumulative-input-040.tex}
% SGA 1, Expose I, definition preceding Corollary I.10.3.
% French authority semantic boundary: lines 2064--2070,
% 7 LF lines / 408 bytes; LF-normalized SHA-256:
% CD605B35B1A40FE92DCF1B4FDF6A8A774956937B71B8D911E24E3647E7A37C2D.
% The blank line 2071 closes the source envelope; Corollary I.10.3 begins at
% line 2072 and is excluded.
% SGA1-CORR-0011: propagate the accepted corrected-French base scheme Y.

An algebra \(L\) of finite rank over~\(K\) will be said to be
\emph{unramified over~\(Y\)} (or simply unramified over~\(K\), if \(Y\) is
understood) if \(L\) is a separable algebra over~\(K\), that is, a direct
product of separable extensions \(K_i\), and if the normalization \(Y'\)
of~\(Y\) in~\(L\) (the disjoint union of the normalizations of~\(Y\) in
the \(K_i\)) is unramified (that is, étale by
Corollary~\ref{I.9.11}) over~\(Y\). Thus:
\end{filecontents*}
% Original input: drafts/SGA1_I_10_3_English_source_draft.texfrag
\begin{filecontents*}[overwrite,noheader]{sga1-cumulative-input-041.tex}
% SGA 1, Expose I, Corollary I.10.3 statement only.
% French authority semantic boundary: lines 2072--2086,
% 15 LF lines / 582 bytes; LF-normalized SHA-256:
% 590043E258138366395EEB9C9CCEA5E8A81D9AD8B7C9339B255F320C0EA6725E.
% The blank line 2087 closes the source envelope; the unheaded proof begins
% at line 2088 and is excluded.

\begin{corollary}
\label{I.10.3}
For every \(X\) finite over~\(Y\) such that every irreducible component of
\(X\) dominates~\(Y\), let \(R(X)\) denote the ring of rational functions
on~\(X\) (the product of the
\marginpar{\raggedright\footnotesize French source p.~22}%
local rings at the
generic points of the irreducible components of~\(X\)), so that
\(X\longmapsto R(X)\) is a functor with values in algebras of finite rank
over \(K=R(Y)\). This functor establishes an equivalence between the category
of connected étale coverings of~\(Y\) and the category of extensions \(L\)
of~\(K\) unramified over~\(Y\).
\end{corollary}
\end{filecontents*}
% Original input: drafts/SGA1_I_10_3_PROOF_AFFINE_REDUCTION_English_source_draft.texfrag
\begin{filecontents*}[overwrite,noheader]{sga1-cumulative-input-042.tex}
% SGA 1, Expose I, Corollary I.10.3 proof: affine reduction and discriminant criterion.
% French authority semantic boundary: lines 2090--2121,
% 32 LF lines / 1707 bytes; LF-normalized SHA-256:
% FA596E3A812E7D4821098F15509E9270C26B0137543CD76A98320D828BAABDB4.
% The blank line 2122 closes the source envelope; line 2123 begins the
% excluded sorites bridge. Envelope SHA-256:
% D027231349D8053067BB2A289A5D72A6BFABC735E22D244C4B86071D3921BE3A.

Suppose that \(Y\) is affine, say \(Y=\operatorname{Spec}(A)\), where
\(A\) is normal with field of fractions~\(K\). Let \(L\) be a finite
\(K\)-algebra that is a product of fields. By definition, the normalization
\(Y'\) of~\(Y\) in~\(L\) is then isomorphic to
\(\operatorname{Spec}(A')\), where \(A'\) is the normalization of~\(A\)
in~\(L\). To say that \(L\) is unramified over~\(Y\) means that \(A'\) is
unramified over~\(A\), or equivalently that it is étale over~\(A\). If
\(A\) is local, this is equivalent to saying that the local rings
\(A'_{\mathfrak n}\), where \(\mathfrak n\) ranges over the finite set of
maximal ideals of~\(A'\), that is, its prime ideals lying over the maximal
ideal \(\mathfrak m\) of~\(A\), are unramified (that is, étale) over the
local ring~\(A\).
Finally, let us also note that the discriminant criterion
(Proposition~\ref{I.4.10}) can be applied in this situation. More generally,
a variant of that criterion should be stated as follows, without a
preliminary flatness condition when \(X\) dominates \(Y\), with \(Y\)
nevertheless assumed locally integral: the maps
\(A\longrightarrow B\) and \(B\longrightarrow B\otimes_A K=L\) are
injective---then \(\operatorname{Tr}_{L/K}\) is defined---and
\(\operatorname{Tr}_{L/K}(xy)\) induces a \emph{fundamental bilinear form}
\(B\times B\longrightarrow A\); that is, there exist elements \(x_i\in B\)
(\(1\leq i\leq n\), where \(n=\operatorname{rank}_K L\)) such that
\(\operatorname{Tr}_{L/K}(x_i x_j)\in A\) for all \(i,j\), and
\(\det\!\left(\bigl(\operatorname{Tr}_{L/K}(x_i x_j)\bigr)_
{1\leq i,j\leq n}\right)\) is invertible in~\(A\).
\end{filecontents*}
% Original input: drafts/SGA1_I_10_3_PROOF_OPENING_English_source_draft.texfrag
\begin{filecontents*}[overwrite,noheader]{sga1-cumulative-input-043.tex}
% SGA 1, Expose I, Corollary I.10.3 proof opening only.
% French authority semantic boundary: line 2088,
% 1 LF line / 51 bytes; LF-normalized SHA-256:
% 1C376ABB7DD5E5FE3AA7AFBA41ECE80D4A75B258B06DEF8921EEF7556CF68421.
% The blank line 2089 closes the source envelope; the affine reduction begins
% at line 2090 and is excluded.

The inverse functor is the normalization functor.
\end{filecontents*}
% Original input: drafts/SGA1_I_10_3_SORITES_BRIDGE_English_source_draft.texfrag
\begin{filecontents*}[overwrite,noheader]{sga1-cumulative-input-044.tex}
% SGA 1, Expose I, Corollary I.10.3 proof: sorites bridge only.
% French authority semantic boundary: lines 2123--2124,
% 2 LF lines / 94 bytes; LF-normalized SHA-256:
% 5260A725EB492499518BDCA386A613A079D04EE2C34F49C9C69077C13B29C28B.
% The blank line 2125 closes the source envelope; Proposition I.10.4 begins
% at line 2126 and is excluded. Envelope SHA-256:
% 5C03F6C1270A1EF8A11C9AC7336285C2B036F6712FD1287D1166728192F9455A.

The sorites (\ref{I.4.6}) immediately implies the sorites of
unramifiedness in the classical setting:
\end{filecontents*}
% Original input: drafts/SGA1_I_10_4_English_source_draft.texfrag
\begin{filecontents*}[overwrite,noheader]{sga1-cumulative-input-045.tex}
% SGA 1, Expose I, Proposition I.10.4.
% French authority semantic boundary: lines 2126--2143,
% 18 LF lines / 885 bytes; LF-normalized SHA-256:
% 15164185937F81EA3C39667EDAA74A15AF9D150FCDCCA83E456CBE39E7A4FBC3.
% The blank line 2144 closes the source envelope; the split phrase
% `De ... plus:' and its page apparatus begin at line 2145 and are excluded.
% Envelope SHA-256:
% F66A9ABED4E15F584BD27B7840253D17080653E07B448929E16B0ABC298F56F3.

\begin{proposition}
\label{I.10.4}
% SGA1-CORR-0012: propagate the accepted corrected-French transitivity base.
Let \(Y\) be a normal integral prescheme with function field~\(K\).
\textup{(i)} \(K\) is unramified over~\(Y\).
\textup{(ii)} If \(L\) is an extension of~\(K\) unramified over~\(Y\), if
\(Y'\) is a normal prescheme with function field~\(L\) that dominates~\(Y\)
(for example, the normalization of~\(Y\) in~\(L\)), and if \(M\) is an
extension of~\(L\) unramified over~\(Y'\), then \(M/K\) is unramified
over~\(Y\) (\emph{transitivity} of unramifiedness).
\textup{(iii)} Let \(Y'\) be a normal integral prescheme dominating~\(Y\),
with function field extension~\(K'/K\); if \(L\) is an extension of~\(K\)
unramified over~\(Y\), then \(L\otimes_K K'\) is an extension of~\(K'\)
unramified over~\(Y'\) (the \emph{translation} property).
\end{proposition}
\end{filecontents*}
% Original input: drafts/SGA1_I_10_5_English_source_draft.texfrag
\begin{filecontents*}[overwrite,noheader]{sga1-cumulative-input-046.tex}
% SGA 1, Expose I, page-23 connective and Corollary I.10.5.
% French authority semantic boundary: lines 2145--2161,
% 17 LF lines / 352 bytes; LF-normalized SHA-256:
% 553FF11C1B02FA422759AB738B7D0C178E0E905B445ED51115CDBE834064E500.
% The blank line 2162 closes the source envelope; the prose introducing
% Corollary I.10.6 begins at line 2163 and is excluded. Envelope SHA-256:
% 1BFBC0BFBA3ACBD330A3FB96CDC1104DC2CA119D435B5D5E8FF8E0EC8B98F8C0.

\marginpar{\raggedright\footnotesize French source p.~23}%
Moreover:

\begin{corollary}
\label{I.10.5}
Under the hypotheses of Proposition~\ref{I.10.4}\textup{(iii)}, if
\(Y=\operatorname{Spec}(A)\) and \(Y'=\operatorname{Spec}(A')\), then
\(\bar{A'}\),\footnote{The original printing has the unbarred \(A'\) at
this point; the corrected French \TeX{} replaces it with \(\bar{A'}\),
and the target follows that corrected branch.} the coordinate ring of the
normalization of~\(Y'\) in \(L'=L\otimes_K K'\), is identified with
\(\bar{A}\otimes_A A'\), where \(\bar{A}\) is the normalization of~\(A\)
in~\(L\).
\end{corollary}
\end{filecontents*}
% Original input: drafts/SGA1_I_10_6_THROUGH_I_10_12_English.texfrag
\begin{filecontents*}[overwrite,noheader]{sga1-cumulative-input-047.tex}
% SGA 1, Expose I, French TeX lines 2163--2353 (end of section I.10).

Usually, people (who are reluctant to consider nonintegral rings, even when
they are direct products of fields) state the translation property in the
following (weaker) form:

\begin{corollary}
\label{I.10.6}
Under the hypotheses of Proposition~\ref{I.10.4}\textup{(iii)}, let \(L_1\)
be a \emph{composite extension} of \(L/K\) (unramified over~\(Y\)) and of
\(K'/K\). Then \(L_1/K'\) is unramified over~\(Y'\). If
\(Y=\operatorname{Spec}(A)\) and \(Y'=\operatorname{Spec}(A')\), one also has
\[
  \overline{A'}=A[\overline{A},A'].
\]
That is, the ring \(\overline{A'}\), the normalization of~\(A'\) in~\(L_1\),
is the \(A\)-algebra generated by~\(A'\) and by the normalization
\(\overline A\) of~\(A\) in~\(L\).
\end{corollary}

This last assertion is false without the unramifiedness hypothesis, even for
composite extensions of number fields~\dots

To conclude this section, we shall give the interpretation of the notion of
an étale covering that corresponds to its intuitive picture: there should be
the \lq\lq maximum number\rq\rq{} of points above the point \(y\in Y\) under
consideration, and in particular there should not be \lq\lq several points
merged together\rq\rq{} above~\(y\). To prove the results in this direction
in the desired generality, we shall admit Proposition~\ref{I.10.7} below
(whose proof will appear in the \emph{multiplodoque}, Chapter~IV, no.~15,
and uses Chevalley's technique of constructible sets and a little descent
theory~\dots).

A morphism of finite type \(f\colon X\to Y\) is called
\emph{universally open} if, for every base extension \(Y'\to Y\) (with
\(Y'\) locally noetherian), the morphism
\[
  f'\colon X'=X\times_Y Y'\longrightarrow Y'
\]
is open, that is, carries open sets to open sets. It is enough to consider
the case in which \(Y'\) is of finite type over~\(Y\) (and even the case in
which \(Y'\) is of the form \(Y[t_1,\dots,t_r]\), where the \(t_i\) are
indeterminates).

A universally open morphism is \emph{a fortiori} open (the converse is
false). On the other hand, if \(f\) is open and \(X\) and \(Y\) are
irreducible, then all the components of all the fibers of~\(f\) have the same
dimension, namely the dimension of the generic fiber
\marginpar{\raggedright\footnotesize French source p.~24}%
\(f^{-1}(z)\), where \(z\) is the generic point of~\(Y\). Finally, if \(Y\)
is normal, this last condition already implies that \(f\) is universally
open (Chevalley's theorem). It follows, for example, that if
\(f\colon X\to Y\) is quasi-finite, with \(Y\) normal and irreducible, then
\(f\) is universally open (equivalently, open) if and only if every
irreducible component of~\(X\) dominates~\(Y\). Recall also that a flat
morphism of finite type is open, and hence universally open. With these
preliminaries established, let us \lq\lq recall\rq\rq{} the following
proposition.

\begin{proposition}
\label{I.10.7}
Let \(f\colon X\to Y\) be a quasi-finite, separated, universally open
morphism. For every \(y\in Y\), let \(n(y)\) be the \lq\lq geometric number
of points of the fiber \(f^{-1}(y)\)\rq\rq, equal to the sum of the separable
degrees of the residue-field extensions \(\kappa(x)/\kappa(y)\), for
\(x\in f^{-1}(y)\). Then the function \(y\longmapsto n(y)\) on~\(Y\) is
upper semicontinuous. If it is constant in a neighborhood of~\(y\) (that is,
if \(n(y)=n(z_i)\), where the \(z_i\) are the generic points of the
irreducible components of~\(Y\) that contain \(y\)), then there must exist a
neighborhood \(U\) of~\(y\) such that \(X|_U\) is finite over~\(U\).
\footnote{Cf.\ EGA~IV, 15.5.1.}
\end{proposition}

\begin{corollary}
\label{I.10.8}
If \(y\longmapsto n(y)\) is constant and \(Y\) is geometrically
unibranch,\footnote{For the definition, see no.~\ref{I.11} below.} then the
irreducible components of~\(X\) are disjoint.
\end{corollary}

\begin{proposition}
\label{I.10.9}
Let \(f\colon X\to Y\) be an \emph{étale}, separated morphism. With the
notation of Proposition~\ref{I.10.7}, the function \(y\longmapsto n(y)\) is
upper semicontinuous. For it to be constant in a neighborhood of~\(y\)
(that is, for \(n(y)=n(z_i)\), where the \(z_i\) are the generic points of
the irreducible components of~\(Y\) that contain \(y\)), it is necessary and
sufficient that there exist an open neighborhood \(U\) of~\(y\) such that
\(X|_U\) is finite over~\(U\), that is, is an \emph{étale covering} of~\(U\).
\end{proposition}

\begin{corollary}
\label{I.10.10}
For a separated étale morphism \(f\colon X\to Y\), with \(Y\) connected, to
be finite (that is, to make \(X\) an \emph{étale covering} of~\(Y\)), it is
necessary and sufficient that all the fibers of~\(f\) have the same
geometric number of points.
\end{corollary}

In Proposition~\ref{I.10.7} and its corollary there was no normality
hypothesis on~\(Y\). If such a hypothesis is imposed, one obtains the
stronger statement below (which is most often taken as the definition of the
unramifiedness of a covering).

\begin{theorem}
\label{I.10.11}
\marginpar{\raggedright\footnotesize French source p.~25}%
Let \(f\colon X\to Y\) be a \emph{quasi-finite}, separated morphism. Suppose
that \(Y\) is \emph{irreducible}, that every component of~\(X\) dominates
\(Y\), and that \(X\) is reduced (that is, \(\mathcal O_X\) has no nilpotent
elements). Let \(n\) be the degree of~\(X\) over~\(Y\) (the sum of the
degrees, over the field \(K\) of~\(Y\), of the fields \(K_i\) of the
irreducible components \(X_i\) of~\(X\)). Let \(y\) be a normal point of
\(Y\). Then the geometric number \(n(y)\) of points of~\(X\) above~\(y\) is
at most~\(n\), and equality holds if and only if there exists an open
neighborhood \(U\) of~\(y\) such that \(X|_U\) is an \emph{étale covering}
of~\(U\).
\end{theorem}

The implication from the existence of such an étale neighborhood to
\(n(y)=n\) is trivial. Let us prove the converse.\footnote{The French text
names these two directions in the opposite order; the argument itself starts
from \(n(y)=n\) and proves the existence of the étale neighborhood, as
translated here.} Let \(z\) be the generic point of~\(Y\). We have
\[
  n(z)=\text{the sum of the separable degrees of the }K_i/K\leq n,
\]
and Proposition~\ref{I.10.7} gives \(n(y)\leq n(z)\leq n\). Equality implies
that \(X|_U\) is finite over~\(U\) for a suitable neighborhood \(U\) of~\(y\).
We may therefore suppose that \(X\) is finite over~\(Y\) and that the
function \(n(y')\) on~\(Y\) is constant. By Corollary~\ref{I.10.8}, \(X\) is
then the disjoint union of its irreducible components; to prove that it is
unramified at~\(y\), we are reduced to the case in which \(X\) is
irreducible, hence integral. Finally, we may suppose that
\(Y=\operatorname{Spec}(\mathcal O_y)\). The theorem then reduces to the
following classical statement.

\begin{corollary}
\label{I.10.12}
Let \(A\) be a normal local ring (noetherian, as always) with fraction field
\(K\); let \(L\) be a finite extension of~\(K\), of degree \(n\) and
separable degree \(n_s\); and let \(B\) be a subring of~\(L\), finite over
\(A\), with fraction field~\(L\). Let \(\mathfrak m\) be the maximal ideal
of~\(A\), and let \(n'\) be the separable degree of
\(B/\mathfrak mB\) over \(A/\mathfrak mA=k\) (that is, the sum of the
separable degrees of the residue-field extensions of this ring). Then
\(n'\leq n_s\), and \emph{a fortiori} \(n'\leq n\). Equality in the latter
inequality holds if and only if \(B\) is unramified (equivalently, étale)
over~\(A\).
\end{corollary}

It remains only to show that \(n'=n\) implies that \(B\) is étale over~\(A\).
Recall the proof when \(k\) is infinite. It is enough to show that
\(R=B/\mathfrak mB\) is separable over~\(k\). If it were not, a familiar
lemma would imply that there is an element \(a\in R\) whose minimal
polynomial over~\(k\) has degree \(>n'\). This element is the image of an
element \(x\in B\), whose minimal polynomial over~\(K\) (as an element of
\(L\)) has degree at most~\(n\). Moreover, the coefficients of this last
polynomial lie in~\(A\), since \(A\) is normal. Reduction modulo
\(\mathfrak m\) therefore gives a monic polynomial \(F\in k[t]\), of degree
at most \(n=n'\), such that \(F(a)=0\), a contradiction.

\marginpar{\raggedright\footnotesize French source p.~26}%
In the general case (where \(k\) may be finite), return to geometric
language and put \(Y=\operatorname{Spec}(A)\), \(X=\operatorname{Spec}(B)\).
Consider \(Y'=\operatorname{Spec}(A[t])\), which is faithfully flat over
\(Y\), and the generic point \(y'\) of the fiber
\(\operatorname{Spec}(k[t])\) of~\(Y'\) over~\(y\). Then \(X\) is unramified
over~\(Y\) at~\(y\) if and only if
\[
  X'=X\times_Y Y'=\operatorname{Spec}(B[t])
\]
is unramified at~\(y'\) over~\(Y'\), as one sees immediately. By the choice
of~\(y'\), its residue field is \(k(t)\), hence infinite. Since \(y'\) is a
normal point of~\(Y'\), this reduces the assertion to the preceding case.
\end{filecontents*}
% Original input: drafts/SGA1_I_11_English.texfrag
\begin{filecontents*}[overwrite,noheader]{sga1-cumulative-input-048.tex}
% SGA 1, Exposé I, French TeX lines 2379--2517 (section I.11).

\section{Some complements}
\label{I.11}

We have already said that a connected étale covering of an integral scheme
need not be integral. Here are two examples of this fact.

\noindent
a) Let \(C\) be an algebraic curve with an ordinary double point~\(x\), let
\(C'\) be its normalization, and let \(a\) and~\(b\) be the two points of
\(C'\) above~\(x\). Let \(C_i'\) (\(i=1,2\)) be two copies of~\(C'\), and
let \(a_i\) and~\(b_i\) be the points of \(C_i'\) corresponding respectively
to \(a\) and~\(b\). In the disjoint union \(C_1'\amalg C_2'\), identify
\(a_1\) with~\(b_2\), and \(a_2\) with~\(b_1\). (The reader is left to make
the identification process precise; it will be explained in Chapter~VI of
the \emph{multiplodoque}, but, for curves over an algebraically closed field,
it is treated in Serre's book on algebraic curves.) The result is a
\emph{connected} and \emph{reducible} curve \(C''\), which is an étale
covering of~\(C\) of degree~\(2\). More generally, the reader will verify
that the connected \lq\lq Galois\rq\rq{} étale coverings \(C''\) of~\(C\)
for which the inverse image \(C''\times_C C'\) is a \emph{trivial} covering
of~\(C'\) (that is, isomorphic to a disjoint union of some number of copies
of~\(C'\)) are \lq\lq cyclic\rq\rq{} of degree~\(n\), and that, for every
integer \(n>0\), one can construct a connected cyclic étale covering of
degree~\(n\). In the language of the fundamental group, which will be
developed later, this means that the quotient of~\(\pi_1(C)\) by the closed
normal subgroup generated by the image of
\[
  \pi_1(C')\longrightarrow\pi_1(C)
\]
(the homomorphism induced by the projection) is isomorphic to the profinite
completion of~\(\mathbb Z\). More precisely, one should be able to show that the
fundamental group of~\(C\) is isomorphic to the topological free product of
the fundamental group of~\(C'\) with the profinite completion of~\(\mathbb Z\).
Notice that questions of this kind gave rise to \lq\lq descent
theory\rq\rq{} for schemes.

\noindent
b) Let \(A\) be a complete local integral domain. Its normalization \(A'\) is
known to be finite over~\(A\) (Nagata); it is therefore a complete semilocal
ring, and hence local because it is an integral domain.
\marginpar{\raggedright\footnotesize French source p.~27}%
Suppose that the residue-field extension \(L/k\) thereby defined is not
radicial (in the contrary case, \(A\) will be called geometrically
unibranch; see below). This is the case, for example, for the ring
\[
  \mathbb R[[s,t]]/(s^2+t^2)\mathbb R[[s,t]],
\]
where \(\mathbb R\) is the field of real numbers. Let \(k'\) be a finite Galois
extension of~\(k\) such that \(L\otimes_k k'\) splits, and let \(B\) be
a finite étale algebra over~\(A\) corresponding to the residue-field
extension~\(k'\) (recall that \(B\) is essentially unique). Then
\[
  B'=A'\otimes_A B
\]
over~\(B\) has residue algebra \(L\otimes_k k'\), which is not local. Thus
\(B'\) is not a local ring and, being complete, has zero divisors. But \(B'\)
is contained in the total ring of fractions of~\(B\): indeed, it is free
over~\(A'\), hence torsion-free over~\(A'\), hence torsion-free over~\(A\),
and is consequently contained in
\[
  B'\otimes_A K
  =B'_{(K)}
  =A'_{(K)}\otimes_K B_{(K)}
  =B_{(K)},
\]
since \(A'_{(K)}=K\). It follows that \(B\) is not an integral domain. In the case of
the ring
\[
  \mathbb R[[s,t]]/(s^2+t^2)\mathbb R[[s,t]],
\]
taking \(k'/k=\mathbb C/\mathbb R\), the ring \(B\) obtained is the local ring, at their
point of intersection, of two intersecting lines in the plane.

Notice also that, if there exists a connected étale covering \(X\) of an
integral~\(Y\) that is not irreducible, then every irreducible component
of~\(X\) gives an example of an unramified covering \(X'\) of~\(Y\),
dominating~\(Y\), that is not étale over~\(Y\). Thus, in example~a), \(C'\)
is unramified over~\(C\) without being étale at the two points \(a\) and~\(b\).
This can also be seen directly by inspecting the completions of the local
rings at \(x\) and~\(a\): from the \lq\lq formal\rq\rq{} point of view,
\(C'\) at~\(a\) is identified with a closed subscheme of~\(C\) at~\(x\),
namely one of the two \lq\lq branches\rq\rq{} of~\(C\) passing through~\(x\).

In a) and~b), the failure of the conclusions of
Proposition~\ref{I.9.5}\textup{(i)} and~\textup{(ii)} is seen to be directly
connected with the fact that a point of~\(Y\) \lq\lq splits\rq\rq{} into
\emph{distinct} points of the normalization. In b), the fact that the
residue-field extension is not radicial is to be interpreted geometrically
in this way. More precisely, we shall say that a local integral domain~\(A\)
is \emph{geometrically unibranch} if its normalization has only one maximal
ideal and the corresponding residue-field extension is radicial. A point
\(y\) of an integral prescheme is said to be geometrically unibranch if its
local ring is so. Examples are a normal point, an ordinary cusp of a curve,
and so on. It seems that, if \(Y\) has a point that is not unibranch, there
always exists a connected, nonirreducible étale covering of~\(Y\); this is
at least what we have shown in case~b), when \(Y\) is the spectrum of a
complete local ring.
\marginpar{\raggedright\footnotesize French source p.~28}%
On the other hand, one can show that \emph{if every point of~\(Y\) is
geometrically unibranch, then every connected unramified \(Y\)-prescheme
dominating~\(Y\) is étale} and irreducible. The proof repeats that of
Proposition~\ref{I.9.5}, using the following generalization of
Theorem~\ref{I.8.3}, which will be proved later by means of descent
theory:\footnote{Cf.\ \hyperref[IX.4.10]{IX~4.10}. For a more direct proof, cf.\
EGA~IV, 18.10.3, using a variant of Proposition~\ref{I.9.5} for
geometrically unibranch local rings.}

\emph{Let \(Y'\to Y\) be a finite, radicial, surjective morphism (that is,
what might be called a \lq\lq universal homeomorphism\rq\rq). Consider the
functor
\[
  X\longmapsto X\times_Y Y'=X'
\]
from \(Y\)-preschemes to \(Y'\)-preschemes. This functor induces an
equivalence between the category of étale \(Y\)-schemes and the category of
étale \(Y'\)-schemes.} For example, this result may be applied when \(Y'\)
is the normalization of~\(Y\), with \(Y\) assumed unibranch (and \(Y'\)
finite over~\(Y\), as is true in all cases encountered in practice), or
when a \(Y''\) is \lq\lq sandwiched\rq\rq{} between~\(Y\) and its
normalization; in the latter case the normalization need no longer be finite
over~\(Y\).
\end{filecontents*}
% Original input: drafts/SGA1_I_5_English_source_draft.texfrag
\begin{filecontents*}[overwrite,noheader]{sga1-cumulative-input-049.tex}
% SGA 1, Expose I, section I.5.
% French authority lines 952--1167; source-first reviewed working draft.
% Raw-slice SHA-256:
% 4E0A4AA4FBB2C5F74E83C626A1DCCF0D20456E28922F65F08E5AD7016824F285.

\section{The fundamental property of étale morphisms}
\label{I.5}

\begin{theorem}
\label{I.5.1}
Let \(f\colon X\to Y\) be a morphism of finite type.  For \(f\) to be an
open immersion, it is necessary and sufficient that it be
\emph{étale and radicial}.
\end{theorem}

Recall that radicial means injective, with radicial residue-field
extensions; one may also recall that this means that the morphism remains
injective after every base change.  Necessity is trivial; it remains to
prove sufficiency.  We shall give two different proofs, the first shorter
and the second more elementary.

\begin{enumerate}
\item[1)] A flat morphism is open, so (replacing \(Y\) by \(f(X)\)) we may
  suppose that \(f\) is a \emph{homeomorphism} onto~\(Y\).  After every
  base change, \(f\) remains flat, radicial, and surjective, hence a
  homeomorphism and therefore closed.  Thus \(f\) is \emph{proper}.  Hence
  \(f\) is \emph{finite} (reference: Chevalley's theorem) and is defined by a coherent
  sheaf \(\mathcal B\) of algebras.  The sheaf \(\mathcal B\) is locally
  free and, by the hypothesis, has rank~\(1\) everywhere.  Thus \(X=Y\),
  as required.

\item[2)] One may suppose \(Y\) and \(X\) \emph{affine}.  Moreover, one easily
  reduces to proving the following: if
  \(Y=\operatorname{Spec}(A)\), with \(A\) local, and if \(f^{-1}(y)\) is
  nonempty, where \(y\) is the closed point of~\(Y\), then \(X=Y\).
  Indeed, this will imply that every \(y\in f(X)\) has an open
  neighborhood \(U\) such that \(X|U=U\).  We shall have
  \(X=\operatorname{Spec}(B)\), and we want to prove \(A=B\).  For this,
  it is enough to prove the analogous assertion after replacing \(A\)
  by~\(\widehat A\) and \(B\) by \(B\otimes_A\widehat A\), taking into
  account that \(\widehat A\) is faithfully flat over~\(A\).  We may
  therefore suppose that \(A\) is \emph{complete}.  Let \(x\) be the point
  above~\(y\).  By Corollary~\ref{I.2.2}, \(\mathcal O_x\) is finite
  over~\(A\), and hence, being flat and radicial over~\(A\), is \(A\)
  itself.  Thus
  \[
    X=Y\amalg X'
  \]
  as a disjoint union.  Since \(X\) is radicial over~\(Y\), \(X'\) is
  empty.  This completes the proof.
\end{enumerate}

\begin{corollary}
\label{I.5.2}
Let \(f\colon X\to Y\) be both a \emph{closed immersion} and \emph{étale}.  If \(X\) is
connected, \(f\) is an isomorphism from~\(X\) onto a connected component
of~\(Y\).
\end{corollary}

Indeed, \(f\) is also an open immersion.  It follows that:

\marginpar{\raggedright\footnotesize French source p.~7}%
\begin{corollary}
\label{I.5.3}
Let \(X\) be an unramified \(Y\)-scheme, with \(Y\) connected.  Then every
section of~\(X\) over~\(Y\) is an isomorphism from~\(Y\) onto a connected
component of~\(X\).  There is therefore a one-to-one correspondence
between the set of such sections and the set of connected components
\(X_i\) of~\(X\) for which the projection \(X_i\to Y\) is an isomorphism
(or, equivalently by Theorem~\ref{I.5.1}, is surjective, étale, and
radicial).\footnote{The French TeX and the original printing give only
``surjective and radicial'' and cite Theorem~\ref{I.5.1}, which also requires
étaleness.  Without it the claim is false: for example,
\(\operatorname{Spec}(k)\to
\operatorname{Spec}(k[\varepsilon]/(\varepsilon^2))\) is surjective on
points, radicial, unramified, separated, and of finite type, but is not an
isomorphism.  The logically correct criterion is used here with the source defect
disclosed.}
In particular, a section is determined by its value at one point.
\end{corollary}

Only the first assertion requires proof.  By Corollary~\ref{I.5.2}, it is
enough to observe that a section is a closed immersion (because \(X\) is
separated over~\(Y\)) and is étale by Corollary~\ref{I.4.8}.

\begin{corollary}
\label{I.5.4}
Let \(X\) and \(Y\) be preschemes over~\(S\), with \(X\) unramified and
separated over~\(S\) and \(Y\) connected.  Let \(f,g\colon Y\to X\) be
two \(S\)-morphisms, and let \(y\in Y\).  Suppose that
\(f(y)=g(y)=x\) and that the homomorphisms of residue fields
\(\kappa(x)\to\kappa(y)\) induced by \(f\) and \(g\) are identical
(``\(f\) and \(g\) coincide geometrically at~\(y\)'').  Then \(f\) and
\(g\) are identical.
\end{corollary}

This follows from Corollary~\ref{I.5.3} by reducing to the case \(Y=S\),
replacing \(X\) by \(X\times_S Y\).

Here is a particularly important variant of Corollary~\ref{I.5.3}.

\begin{theorem}
\label{I.5.5}
Let \(S\) be a prescheme, let \(X\) and \(Y\) be two \(S\)-preschemes,
and let \(S_0\) be a closed subprescheme of~\(S\) having the same
underlying space as~\(S\).  Put
\[
  X_0=X\times_S S_0
  \qquad\text{and}\qquad
  Y_0=Y\times_S S_0,
\]
the ``restrictions'' of \(X\) and \(Y\) to~\(S_0\).  Suppose that \(X\)
is étale over~\(S\).  Then the natural map
\[
  \operatorname{Hom}_S(Y,X)
  \longrightarrow
  \operatorname{Hom}_{S_0}(Y_0,X_0)
\]
is \emph{bijective}.
\end{theorem}

We are again reduced to the case \(Y=S\); the assertion then follows
from the ``topological'' description of the sections of \(X/Y\) given
in Corollary~\ref{I.5.3}.

\paragraph{Scholium.}
This result includes both a \emph{uniqueness} assertion and an
\emph{existence} assertion for \emph{morphisms}.  It may also be
expressed, when \(X\) and \(Y\) are both taken to be étale over~\(S\),
by saying that the functor
\[
  X\longmapsto X_0
\]
from the category of étale \(S\)-schemes to the category of étale
\(S_0\)-schemes is \emph{fully faithful}; that is, it establishes an
\emph{equivalence} of the first category with a \emph{full subcategory}
of the second.  We shall see below that it is actually an equivalence
between the first and the second (which will be an \emph{existence
theorem for étale \(S\)-schemes}).

\marginpar{\raggedright\footnotesize French source p.~8}%
The following form of Theorem~\ref{I.5.5}, apparently more general, is
often convenient.

\begin{corollary}[``Extension theorem for liftings'']
\label{I.5.6}
Consider a commutative diagram of morphisms
\[
  \begin{array}{ccc}
    X & \longleftarrow & Y_0\\
    \downarrow && \downarrow\\
    S & \longleftarrow & Y
  \end{array}
\]
in which \(X\to S\) is étale and \(Y_0\to Y\) is a bijective closed
immersion.  Then there is a unique morphism \(Y\to X\) that makes the
two corresponding triangles commute.
\end{corollary}

Indeed, replacing \(S\) by~\(Y\) and \(X\) by \(X\times_S Y\), we are
reduced to the case \(Y=S\); this is then the special case of
Theorem~\ref{I.5.5} with \(Y=S\).

We should also note the following immediate consequence of
Theorem~\ref{I.5.1} (which we did not state as Corollary~1, so as not to
interrupt the line of ideas developed after Theorem~\ref{I.5.1}):

\begin{proposition}
\label{I.5.7}
Let \(X\) and \(X'\) be two preschemes of finite type and flat over~\(Y\),
and let \(g\colon X\to X'\) be a \(Y\)-morphism.  For \(g\) to be an open
immersion (respectively, an isomorphism), it is necessary and sufficient
that, for every \(y\in Y\), the induced morphism on the fibers
\[
  g\otimes_Y\kappa(y)\colon
  X\otimes_Y\kappa(y)
  \longrightarrow
  X'\otimes_Y\kappa(y)
\]
be so.
\end{proposition}

It is enough to prove sufficiency.  Since the assertion holds for
surjectivity, we are reduced to the case of an open immersion.  By
Theorem~\ref{I.5.1}, one must verify that \(g\) is \emph{radicial}, which
is immediate, and that it is \emph{étale}, which follows from
Corollary~\ref{I.5.9} below.

\begin{corollary}
\label{I.5.8}
(This should be moved to no.~\ref{I.3}.)  Let \(X\) and \(X'\) be
two preschemes over~\(Y\), let \(g\colon X\to X'\) be a \(Y\)-morphism, let
\(x\in X\), and let \(y\) be its image in~\(Y\).  For \(g\) to be
quasi-finite (respectively, unramified) at~\(x\), it is necessary and
sufficient that the same be true of \(g\otimes_Y\kappa(y)\).
\end{corollary}

Indeed, the two algebras over \(\kappa(g(x))\) that one must examine to
verify that the morphism is quasi-finite (respectively, unramified) at~\(x\)
are the same for \(g\) and for \(g\otimes_Y\kappa(y)\).

\marginpar{\raggedright\footnotesize French source p.~9}%
\begin{corollary}
\label{I.5.9}
With the notation of Corollary~\ref{I.5.8}, suppose that \(X\) and \(X'\)
are flat and of finite type over~\(Y\).  For \(g\) to be flat
(respectively, étale) at~\(x\), it is necessary and sufficient that
\(g\otimes_Y\kappa(y)\) be so.
\end{corollary}

For ``flat,'' the statement is included only as a reminder: it is one of
the fundamental criteria for flatness.\footnote{Cf.\
\hyperref[IV]{Exposé~IV},
\hyperref[IV.5.9]{no.~5.9}.}  For ``étale,'' the assertion follows from this, taking
Corollary~\ref{I.5.8} into account.
\end{filecontents*}
% Original input: drafts/SGA1_I_6_English_source_draft.texfrag
\begin{filecontents*}[overwrite,noheader]{sga1-cumulative-input-050.tex}
% SGA 1, Expose I, section I.6.
% French authority lines 1168--1216; bounded source-first working draft.
% LF-normalized source-slice SHA-256:
% 7F9831D26582DB33861D2C9D48F6DD09C6956F639C566A92FAFA0514F08DDFCD.

\section{Application to étale extensions of complete local rings}
\label{I.6}

This number is a special case of results on formal preschemes that will have
to appear in the \emph{multiplodoque}.  Nevertheless, here one can get by
more cheaply, that is, without the explicit local determination of étale
morphisms in \hyperref[I.7]{no.~I.7} (using the Main Theorem).  This is perhaps a
sufficient reason to keep the present number (even in the
\emph{multiplodoque}) in this place.

\begin{theorem}
\label{I.6.1}
Let \(A\) be a complete local ring (noetherian, of course), with residue
field~\(k\).  For every \(A\)-algebra~\(B\), put

\[
  R(B)=B\otimes_A k,
\]

regarded as a \(k\)-algebra; it thus depends functorially on~\(B\).  Then \(R\)
defines an \emph{equivalence} between the category of \(A\)-algebras
\emph{finite and étale over~\(A\)} and the category of algebras
\emph{of finite rank and separable} over~\(k\).
\end{theorem}

First, the functor in question is fully faithful, as follows from the more
general fact:

\begin{corollary}
\label{I.6.2}
Let \(B\) and \(B'\) be two \(A\)-algebras finite over~\(A\).  If \(B\) is
étale over~\(A\), then the canonical map
\[
  \operatorname{Hom}_{A\text{-alg}}(B,B')
  \longrightarrow
  \operatorname{Hom}_{k\text{-alg}}\bigl(R(B),R(B')\bigr)
\]
is bijective.
\end{corollary}

We are reduced to the case in which \(A\) is artinian (by replacing \(A\) with
\(A/\mathfrak m^n\)); this is then a special case of
Theorem~\ref{I.5.5}.

It remains to prove that for every finite separable \(k\)-algebra~\(L\)
(why not say ``étale''? It is shorter), there is an algebra \(B\), étale
over~\(A\), such that \(R(B)\) is isomorphic to~\(L\).  We may suppose that
\(L\) is a separable extension of~\(k\); as such, it has a generator~\(x\),
that is, it is isomorphic to an algebra \(k[t]/F k[t]\), where
\(F\in k[t]\) is a monic polynomial.  Lift \(F\) to a monic polynomial
\(F_1\in A[t]\), and take
\[
  B=A[t]/F_1A[t].
\]
\end{filecontents*}
% Original input: drafts/SGA1_I_7_English_source_draft.texfrag
\begin{filecontents*}[overwrite,noheader]{sga1-cumulative-input-051.tex}
% SGA 1, Expose I, section I.7.
% French authority lines 1217--1492; bounded source-first working draft.
% LF-normalized source-slice SHA-256:
% 3D058327D83BC14D60A40BC46AA3B468BAD9D31AB0A749598EE70CEF7F2112F4.

\marginpar{\raggedright\footnotesize French source p.~10}%
\section{Local construction of unramified and étale morphisms}
\label{I.7}

\begin{proposition}
\label{I.7.1}
Let \(A\) be a noetherian ring, \(B\) a finite \(A\)-algebra, \(u\) a
generator of \(B\) over~\(A\), and \(F\in A[t]\) such that \(F(u)=0\)
(we do not assume that \(F\) is monic).  Put \(u'=F'(u)\), where \(F'\)
is the derivative of~\(F\).  Let \(\mathfrak q\) be a prime ideal of~\(B\)
that does not contain~\(u'\), and let \(\mathfrak p\) be its contraction
to~\(A\).  Then \(B_{\mathfrak q}\) is unramified over
\(A_{\mathfrak p}\).
\end{proposition}

In other words, put \(Y=\operatorname{Spec}(A)\),
\(X=\operatorname{Spec}(B)\), and
\(X_{u'}=\operatorname{Spec}(B_{u'})\).  Then \(X_{u'}\) is unramified
over~\(Y\).  The assertion follows from the following, more precise one.

\begin{corollary}
\label{I.7.2}
The different ideal of \(B/A\) contains \(u'B\), and is equal to it if the
natural homomorphism
\[
  A[t]/FA[t]\longrightarrow B,
\]
which sends \(t\) to~\(u\), is an isomorphism.
\end{corollary}

Let \(J\) be the kernel of the homomorphism \(C=A[t]\to B\).  This kernel
contains \(FA[t]\), and is equal to it in the second case considered in
Corollary~\ref{I.7.2}.  Since the homomorphism is surjective,
\(\Omega^1_{B/A}\) is identified with the quotient of
\(\Omega^1_{C/A}\) by the submodule generated by
\(J\Omega^1_{C/A}\) and \(d(J)\) (the definition of the homomorphism~\(d\),
and the computation of~\(\Omega^1\) for a polynomial algebra, should have
been made explicit in no.~\ref{I.1}).  Identifying
\(\Omega^1_{C/A}\) with~\(C\) by means of the basis~\(dt\), one obtains
\[
  B/(B\cdot J'),
\]
where \(J'\) is the set of images in~\(B\) of the derivatives of the
polynomials \(G\in J\); thus the different ideal is generated by~\(J'\)
(and it is enough to take polynomials~\(G\) that generate~\(J\)).  Since
\(F\in J\), respectively since \(F\) is a generator of~\(J\), the proof is
complete.  (N.B. Corollary~\ref{I.7.2} should be a proposition and
Proposition~\ref{I.7.1} a corollary.)  We obtain:

\begin{corollary}
\label{I.7.3}
Under the hypotheses of Proposition~\ref{I.7.1}, suppose that \(F\) is
monic and that \(A[t]/FA[t]\to B\) is an isomorphism.  Then
\(B_{\mathfrak q}\) is étale over \(A_{\mathfrak p}\) if and only if
\(\mathfrak q\) does not contain~\(u'\).
\end{corollary}

Indeed, since \(B\) is flat over~\(A\), being étale is equivalent to being
unramified, and Corollary~\ref{I.7.2} applies.

\begin{corollary}
\label{I.7.4}
Under the hypotheses of Corollary~\ref{I.7.3}, \(B\) is étale over~\(A\)
if and only if \(u'\) is invertible, or, equivalently, if and only if the
ideal generated by \(F\) and \(F'\) in \(A[t]\) is the unit ideal.
\end{corollary}

The last criterion follows from the first and Nakayama's lemma (in~\(B\)).

A monic polynomial \(F\in A[t]\) having the property stated in
\marginpar{\raggedright\footnotesize French source p.~11}%
Corollary~\ref{I.7.4}
is called a \emph{separable polynomial}.  If \(F\)
is not monic, one should at least require the coefficient of its leading
term to be invertible; when \(A\) is a field, this recovers the usual
definition.
% SOURCE INDEX DEBT: separable (polynomial), French lines 1291--1294.

\begin{corollary}
\label{I.7.5}
Let \(B\) be a \emph{finite} algebra over the \emph{local} ring~\(A\).
Suppose that \(K(A)\) is infinite or that \(B\) is local.  Let \(n\) be the
rank of \(L=B\otimes_A K(A)\) over \(K(A)=k\).  Then \(B\) is unramified
(respectively, étale) over~\(A\) if and only if \(B\) is isomorphic to a
quotient of (respectively, isomorphic to)
\[
  A[t]/FA[t],
\]
where \(F\) is a monic separable polynomial that may be assumed
(respectively, is necessarily) to have degree~\(n\).
\end{corollary}

Only necessity has to be proved.  Suppose that \(B\) is unramified over~\(A\),
so that \(L\) is separable over~\(k\).  The hypothesis then implies that
\(L/k\) has a generator~\(\xi\), and hence the \(\xi^i\),
\(0\leq i<n\), form a basis of \(L\) over~\(k\).  Let \(u\in B\) lift
\(\xi\).  By Nakayama's lemma, the \(u^i\), \(0\leq i<n\), generate the
\(A\)-module~\(B\) (respectively, form a basis of it).  In particular, one
can find a monic polynomial \(F\in A[t]\) such that \(F(u)=0\), and \(B\)
is isomorphic to a quotient of (respectively, isomorphic to)
\(A[t]/FA[t]\).  Finally, by Corollary~\ref{I.7.4} applied to \(L/k\),
\(F\) and \(F'\) generate \(A[t]\) modulo \(\mathfrak mA[t]\).  Hence, by
Nakayama's lemma in \(A[t]/FA[t]\), \(F\) and \(F'\) generate \(A[t]\).

\begin{theorem}
\label{I.7.6}
Let \(A\) be a local ring, and let \(A\to\mathcal O\) be a local
homomorphism such that \(\mathcal O\) is isomorphic to a localization of an
algebra of finite type over~\(A\).  Suppose that \(\mathcal O\) is
unramified over~\(A\).  Then there exist an \(A\)-algebra~\(B\) integral
over~\(A\), a maximal ideal \(\mathfrak n\) of~\(B\), a generator~\(u\) of
\(B\) over~\(A\), and a monic polynomial \(F\in A[t]\), such that
\[
  F'(u)\notin\mathfrak n
  \qquad\text{and}\qquad
  \mathcal O\simeq B_{\mathfrak n}
\]
as \(A\)-algebras.  If \(\mathcal O\) is étale over~\(A\), one may take
\(B=A[t]/FA[t]\).
\end{theorem}

(Of course, the conditions displayed here are also sufficient\ldots)

Let us first point out the pleasant corollaries.

\begin{corollary}
\label{I.7.7}
The algebra \(\mathcal O\) is unramified over~\(A\) if and only if it is
isomorphic to a quotient of an algebra of the same form that is étale
over~\(A\).
\end{corollary}

Indeed, take \(\mathcal O'=B'_{\mathfrak n'}\), where
\(B'=A[t]/FA[t]\) and \(\mathfrak n'\) is the inverse image of
\(\mathfrak n\) in~\(B'\).

\marginpar{\raggedright\footnotesize French source p.~12}%
\begin{corollary}
\label{I.7.8}
Let \(f\colon X\to Y\) be a morphism of finite type and let \(x\in X\).
Then \(f\) is unramified at~\(x\) if and only if there is an open
neighborhood~\(U\) of~\(x\) such that \(f|_U\) factors as
\[
  U\longrightarrow X'\longrightarrow Y,
\]
where the first arrow is a closed immersion and the second is an étale
morphism.
\end{corollary}

This is a straightforward translation of Corollary~\ref{I.7.7}.

Let us show how the statement of Theorem~\ref{I.7.6} follows from the main
assertion.  By Corollary~\ref{I.7.7}, there is a surjective homomorphism
\(\mathcal O'\to\mathcal O\), where \(\mathcal O'\) has the desired
form.\footnote{The maintained French TeX says literally
``the jargon of Theorem~\ref{I.7.6},'' while the original printing has the typo ``jigon.'' Both write the unprimed \(\mathcal O\)
after ``where,'' and move directly to ``since \(\mathcal O'\) and
\(\mathcal O\) are étale.''  Corollary~\ref{I.7.7} constructs the primed
algebra with the desired étale presentation, while the latter sentence
concerns the theorem's final case, in which \(\mathcal O\) is additionally
assumed étale.  Both points are made explicit here.}  If \(\mathcal O\) is étale over~\(A\), then both
\(\mathcal O'\) and \(\mathcal O\) are étale over~\(A\); the morphism
\(\mathcal O'\to\mathcal O\) is therefore étale by
Corollary~\ref{I.4.8}, and hence is an isomorphism.

\subsubsection*{Proof of Theorem~\ref{I.7.6}}

This repeats a proof from the Chevalley seminar.  By the \emph{Main
Theorem}, one has \(\mathcal O=B_{\mathfrak n}\), where \(B\) is a finite
algebra over~\(A\) and \(\mathfrak n\) is a maximal ideal of~\(B\).  Then
\(B/\mathfrak n=K(\mathcal O)\) is a separable, hence monogenic, extension
of~\(k\).  If \(\mathfrak n_i\), \(1\leq i\leq r\), are the maximal ideals
of~\(B\) distinct from \(\mathfrak n\), there is therefore an element
\(u\in B\) that belongs to every \(\mathfrak n_i\) and whose image in
\(B/\mathfrak n\) generates that algebra.  Moreover,
\[
  B/\mathfrak n
  =B_{\mathfrak n}/\mathfrak nB_{\mathfrak n}
  =B_{\mathfrak n}/\mathfrak mB_{\mathfrak n},
\]
where \(\mathfrak m\) is the maximal ideal of~\(A\).  Let us admit for the
moment the following lemma.

\begin{lemma}
\label{I.7.9}
Let \(A\) be a local ring, \(B\) a finite \(A\)-algebra, and
\(\mathfrak n\) a maximal ideal of~\(B\).  Let \(u\in B\) be an element
whose image in \(B_{\mathfrak n}/\mathfrak mB_{\mathfrak n}\) generates
that algebra over \(k=A/\mathfrak m\), and which belongs to every maximal
ideal of~\(B\) distinct from \(\mathfrak n\).  Put
\[
  B'=A[u],\qquad \mathfrak n'=\mathfrak n\cap B'.
\]
Then the canonical homomorphism
\[
  B'_{\mathfrak n'}\longrightarrow B_{\mathfrak n}
\]
is an isomorphism.\footnote{The French TeX and the original printing give
\(B'=B[u]\) and \(\mathfrak n'=\mathfrak nB'\).  Since \(u\in B\), the
first formula would make \(B'=B\), while the proof uses \(B'\) as the
finite \(A\)-subalgebra generated by~\(u\) and uses \(\mathfrak n'\) as the
contraction of \(\mathfrak n\).  The type-correct formulas required by the
proof are used here with the source reading disclosed.}
\end{lemma}

\begin{lemma}
\label{I.7.10}
(This should have appeared as a corollary to Proposition~\ref{I.7.1},
before Corollary~\ref{I.7.5}, which it implies.)  Let \(B\) be a finite
\(A\)-algebra generated by an element~\(u\), and let \(\mathfrak n\) be a
maximal ideal of~\(B\) such that \(B_{\mathfrak n}\) is unramified over~\(A\).
Then there is a monic polynomial \(F\in A[t]\) such that
\[
  F(u)=0
  \qquad\text{and}\qquad
  F'(u)\notin\mathfrak n.
\]
\end{lemma}

Indeed, let \(n\) be the rank of the \(k\)-algebra
\(L=B\otimes_A k\).  By Nakayama's lemma there is a monic polynomial
\(F\in A[t]\) of degree~\(n\) such that \(F(u)=0\).  Let \(f\) be the
reduction of~\(F\) modulo \(\mathfrak m\).  Then \(L\) is isomorphic, as a
\(k\)-algebra, to \(k[t]/fk[t]\).  Hence, by Corollary~\ref{I.7.3},
\(f'(\xi)\) does not belong to the maximal ideal of~\(L\) corresponding to
\(\mathfrak n\), where \(\xi\) denotes the image of~\(t\) in~\(L\), that
is, the image of~\(u\) in~\(L\).  Since \(f'(\xi)\) is the image of
\(F'(u)\), the proof is complete.

\marginpar{\raggedright\footnotesize French source p.~13}%
Theorem~\ref{I.7.6} now follows by combining Lemmas~\ref{I.7.9} and
\ref{I.7.10}.  It remains to prove Lemma~\ref{I.7.9}.  Put
\(S'=B'\setminus\mathfrak n'\), so that
\(B'{S'}^{-1}=B'_{\mathfrak n'}\).  We have the diagram
\[
\begin{array}{ccccc}
B&\longrightarrow&B{S'}^{-1}&\longrightarrow&BS^{-1}=B_{\mathfrak n}\\
\Big\uparrow&&\Big\uparrow&&\\
B'&\longrightarrow&B'{S'}^{-1}
\rlap{\(\displaystyle =B'_{\mathfrak n'}\)}&&
\end{array}
\]
where, similarly, \(S=B\setminus\mathfrak n\), so that
\(BS^{-1}=B_{\mathfrak n}\).  There is therefore a natural homomorphism
\[
  B{S'}^{-1}\longrightarrow BS^{-1}=B_{\mathfrak n}.
\]
We prove that it is an isomorphism.  Equivalently, the elements of~\(S\)
must be invertible in \(B{S'}^{-1}\); equivalently, every maximal ideal
\(\mathfrak p\) of the latter ring must be disjoint from~\(S\), that is,
must induce \(\mathfrak n\) on~\(B\).  Since \(B{S'}^{-1}\) is finite over
\(B'{S'}^{-1}=B'_{\mathfrak n'}\), the ideal \(\mathfrak p\) contracts to
the unique maximal ideal
\(\mathfrak n'B'_{\mathfrak n'}\) of
\(B'_{\mathfrak n'}\),\footnote{The French TeX and the original printing
write \(\mathfrak n'B_{\mathfrak n'}\) here.  The ring in the preceding
clause is \(B'_{\mathfrak n'}\), so the omitted prime on~\(B\) is
restored explicitly.} and hence to the maximal ideal \(\mathfrak n'\) of
\(B'\).  Since \(B\) is finite over~\(B'\), the ideal \(\mathfrak q\) of
\(B\) induced by \(\mathfrak p\), lying over \(\mathfrak n'\), is
necessarily maximal.  It does not contain~\(u\), and is therefore equal to
\(\mathfrak n\).  Here we used the fact that \(u\) belongs to every
maximal ideal of~\(B\) distinct from \(\mathfrak n\).

It remains to prove that \(B{S'}^{-1}=B'{S'}^{-1}\).  Since the former is
finite over the latter, Nakayama's lemma reduces us to proving equality
modulo \(\mathfrak n'B{S'}^{-1}\), and it is enough, a fortiori, to prove
equality modulo \(\mathfrak mB{S'}^{-1}\).  But
\[
  B{S'}^{-1}/\mathfrak mB{S'}^{-1}
  =B_{\mathfrak n}/\mathfrak mB_{\mathfrak n}
\]
is generated over~\(k\) by~\(u\) (here we use the other property of~\(u\)).
Thus the image of~\(B'\), and a fortiori of \(B'{S'}^{-1}\), in this
quotient is the whole ring, since it is a subring containing~\(k\) and the
image of~\(u\).

\paragraph{Remark \textup{*}.}
It should be possible to state Theorem~\ref{I.7.6} for an algebra
\(\mathcal O\) that is only semilocal, so as to subsume
Corollary~\ref{I.7.5} as well.  One would assume that
\(\mathcal O/\mathfrak m\mathcal O\) is a \emph{monogenic} \(k\)-algebra.
One could then find \(u\in B\) whose image in
\(B/\mathfrak mB\) is a generator and which belongs to every maximal ideal
of~\(B\) not arising from \(\mathcal O\).  Lemmas~\ref{I.7.9} and
\ref{I.7.10} should adapt without difficulty.  More generally, \ldots
\end{filecontents*}
% Original input: drafts/SGA1_I_8_English_source_draft.texfrag
\begin{filecontents*}[overwrite,noheader]{sga1-cumulative-input-052.tex}
% SGA 1, Expose I, section I.8.
% French authority lines 1493--1653; bounded source-first working draft.
% LF-normalized source-slice SHA-256:
% A8885501DFCE9760EA0EC20EB6894F8E45B53E71C131132238F24C7579D0A86D.

\section[Infinitesimal lifting of étale schemes]{Infinitesimal lifting of
étale schemes. Application to formal schemes}
\label{I.8}

\begin{proposition}
\label{I.8.1}
Let \(Y\) be a prescheme, let \(Y_0\) be a subprescheme, let \(X_0\) be
an étale \(Y_0\)-scheme, and let \(x\) be a point of~\(X_0\). Then there
exist an étale \(Y\)-scheme~\(X\), a neighborhood \(U_0\) of~\(x\) in
\(X_0\), and a \(Y_0\)-isomorphism
\[
  U_0\xrightarrow{\sim}X\times_Y Y_0.
\]
\end{proposition}

Indeed, let \(y\) be the image of~\(x\) in~\(Y_0\). Applying
Theorem~\ref{I.7.6} to the étale local homomorphism \(A_0\to B_0\) of the
local rings at \(y\) and~\(x\) in \(Y_0\) and~\(X_0\), respectively, one
obtains an isomorphism
\marginpar{\raggedright\footnotesize French source p.~14}%
\[
  B_0=(C_0)_{\mathfrak n_0},
  \qquad
  C_0=A_0[t]/F_0A_0[t],
\]
where \(F_0\) is monic and \(\mathfrak n_0\) is a maximal ideal of~\(C_0\)
that does not contain the class of \(F'_0(t)\) in~\(C_0\). Let \(A\) be
the local ring at~\(y\) in~\(Y\), and let \(F\in A[t]\) be a monic
polynomial whose image under the surjective homomorphism \(A\to A_0\) is
\(F_0\) (one lifts the coefficients of~\(F_0\)). Finally, put
\[
  C=A[t]/FA[t],
\]
and let \(\mathfrak n\) be the inverse image of~\(\mathfrak n_0\) under
the natural epimorphism
\[
  C\longrightarrow C\otimes_A A_0=C_0.
\]
Put
\[
  B=C_{\mathfrak n}.
\]
It follows immediately from the construction and
Proposition~\ref{I.7.1} that \(B\) is étale over~\(A\), and there is an
isomorphism
\[
  B\otimes_A A_0\xrightarrow{\sim}B_0.
\]
\footnote{At this point both the French TeX and the original printing
give \(B\otimes_A A_0=A_0\). Since \(B=C_{\mathfrak n}\) and
\(C\otimes_A A_0=C_0\), base change gives
\((C_0)_{\mathfrak n_0}=B_0\). The type-correct target \(B_0\) is used
here with the shared source defect disclosed.}
We know from EGA, Chapter~I, that there exist a \(Y\)-scheme~\(X\) of
finite type and a point \(z\in X\) above~\(y\) such that
\(\mathcal O_z\) is \(A\)-isomorphic to~\(B\).\footnote{The French TeX
and the original printing have \(C\) here. The following equality with
\(B\otimes_A A_0\), and the appeal to the étaleness of the localized
algebra, require \(B=C_{\mathfrak n}\); \(B\) is therefore restored
explicitly.} Since the latter is étale over \(A=\mathcal O_y\), by taking
\(X\) small enough we may assume that \(X\) is étale over~\(Y\). Put
\(X'_0=X\times_Y Y_0\). The local ring at~\(z\) in~\(X'_0\) is then
identified with
\[
  \mathcal O_z\otimes_A A_0=B\otimes_A A_0,
\]
and is consequently isomorphic to~\(B_0\). This isomorphism is induced by
an isomorphism from a neighborhood \(U_0\) of~\(x\) in~\(X_0\) onto a
neighborhood of~\(z\) in~\(X'_0\) (loc.~cit.); by taking \(X\) small
enough, we may assume that the latter neighborhood is all of~\(X'_0\).
This proves the proposition.

\begin{corollary}
\label{I.8.2}
There is an analogous statement for étale \emph{coverings}, provided the
residue field \(\kappa(y)\) is infinite.
\end{corollary}

The proof is the same, with Corollary~\ref{I.7.5} in place of
Theorem~\ref{I.7.6}.

\begin{theorem}
\label{I.8.3}
The functor considered in Theorem~\ref{I.5.5} is an \emph{equivalence} of
\emph{categories}.
\end{theorem}

By Theorem~\ref{I.5.5}, it remains to show that every étale
\(S_0\)-scheme~\(X_0\) is isomorphic to an \(S_0\)-scheme
\(X\times_S S_0\), where \(X\) is an étale \(S\)-scheme. The underlying
topological space of~\(X\) must necessarily be identical with that of
\(X_0\), with \(X_0\) moreover identified with a closed subprescheme
of~\(X\). The problem is therefore equivalent to the following one: on
the topological space \(\lvert X_0\rvert\) underlying~\(X_0\), find a
sheaf of algebras \(\mathcal O_X\) over \(f_0^*(\mathcal O_S)\) (where
\(f_0\colon X_0\to S_0\) is here regarded as a continuous map of the
underlying spaces) that makes \(\lvert X_0\rvert\) into an étale
\(S\)-prescheme~\(X\), together with a homomorphism of algebras
\[
  \mathcal O_X\longrightarrow\mathcal O_{X_0}
\]
compatible with the homomorphism
\[
  f_0^*(\mathcal O_S)\longrightarrow f_0^*(\mathcal O_{S_0})
\]
on the scalar sheaves, and inducing an isomorphism
\[
  \mathcal O_X\otimes_{f_0^*(\mathcal O_S)}f_0^*(\mathcal O_{S_0})
  \xrightarrow{\sim}\mathcal O_{X_0}.
\]
Then \(X\) will be an étale \(S\)-prescheme whose reduction is~\(X_0\),
\marginpar{\raggedright\footnotesize French source p.~15}%
and hence will be separated over~\(S\), since \(X_0\) is separated
over~\(S_0\); thus \(X\) answers the question. Moreover, if \((U_i)\) is
an open covering of~\(X_0\), and if a solution to the problem has been
found on each~\(U_i\), it follows from the uniqueness theorem
\ref{I.5.5} that these solutions glue (that is, the sheaves of algebras
that define them, equipped with their augmentation homomorphisms, glue).
One checks at once that the ringed space over~\(S\) constructed in this
way is an étale \(S\)-prescheme~\(X\), equipped with an isomorphism
\[
  X_0\xrightarrow{\sim}X\times_S S_0.
\]
It is therefore enough to find a solution locally, and this is provided
by Proposition~\ref{I.8.1}.

\begin{corollary}
\label{I.8.4}
Let \(S\) be a locally noetherian formal prescheme, equipped with an ideal
of definition~\(\mathcal J\),\footnote{The French TeX and the original
printing alternate \(J\) and \(\mathcal J\) for this ideal of definition
and its powers. Since the quotient sheaves all use the same ideal sheaf,
\(\mathcal J\) is used consistently here.} and let
\[
  S_0=\bigl(\lvert S\rvert,\mathcal O_S/\mathcal J\bigr)
\]
be the corresponding ordinary prescheme. Then the functor
\[
  \mathfrak X\longmapsto\mathfrak X\times_S S_0
\]
from the category of étale coverings of~\(S\) to the category of étale
coverings of~\(S_0\) is an equivalence of categories.
\end{corollary}

Of course, an étale covering of a \emph{formal} prescheme~\(S\) means a
covering of~\(S\), that is, a formal prescheme over~\(S\) defined by a
coherent sheaf of algebras \(\mathcal B\), such that \(\mathcal B\) is
\emph{locally free} and the residual fibers
\[
  \mathcal B_s\otimes_{\mathcal O_s}\kappa(s)
\]
of~\(\mathcal B\) are \emph{separable} algebras over~\(\kappa(s)\). If
\(S_n\) denotes the ordinary prescheme
\[
  S_n=\bigl(\lvert S\rvert,\mathcal O_S/\mathcal J^{n+1}\bigr),
\]
then giving a coherent sheaf of algebras \(\mathcal B\) on~\(S\) is
equivalent to giving a sequence of coherent sheaves of algebras
\(\mathcal B_n\) on the \(S_n\), equipped with a transitive system of
homomorphisms
\[
  \mathcal B_m\longrightarrow\mathcal B_n
  \qquad(m\geq n)
\]
that define isomorphisms
\[
  \mathcal B_m\otimes_{\mathcal O_{S_m}}\mathcal O_{S_n}
  \xrightarrow{\sim}\mathcal B_n.
\]
It is immediate that \(\mathcal B\) is locally free if and only if all the
\(\mathcal B_n\) on the \(S_n\) are locally free, and that the separability
condition holds if and only if it holds for~\(\mathcal B_0\), or,
equivalently, for all the \(\mathcal B_n\). Thus \(\mathcal B\) is étale
over~\(S\) if and only if the \(\mathcal B_n\) are étale over the
\(S_n\). In view of this, Corollary~\ref{I.8.4} follows at once from
Theorem~\ref{I.8.3}.

\paragraph{Remark \textup{*}.}
It was not necessary in Corollary~\ref{I.8.4} to restrict oneself to the
case of \emph{coverings}. This is, however, the only case used for the
moment.
\end{filecontents*}
% Original input: drafts/SGA1_I_9_10_English_source_draft.texfrag
\begin{filecontents*}[overwrite,noheader]{sga1-cumulative-input-053.tex}
\begin{corollary}
\label{I.9.10}
Let \(f\colon X\to Y\) be an étale morphism.  If \(Y\) is normal,
then \(X\) is normal; the converse holds if \(f\) is surjective.
\end{corollary}
\end{filecontents*}
% Original input: drafts/SGA1_I_9_11_English_source_draft.texfrag
\begin{filecontents*}[overwrite,noheader]{sga1-cumulative-input-054.tex}
\begin{corollary}
\label{I.9.11}
Let \(f\colon X\to Y\) be a dominant morphism, with \(Y\) normal and
\(X\) connected.  If \(f\) is unramified, then \(f\) is étale; hence
\(X\) is normal and therefore (being connected) irreducible.
\end{corollary}
\end{filecontents*}
% Original input: drafts/SGA1_I_9_11_PROOF_English_source_draft.texfrag
\begin{filecontents*}[overwrite,noheader]{sga1-cumulative-input-055.tex}
% SGA 1, Expose I, proof of Corollary I.9.11.
% French authority envelope: lines 1988--2003,
% 16 LF lines / 897 bytes; LF-normalized SHA-256:
% 1CC6002323446CB167F636BE913569BF360AA2779954CF87088200C5EC8F23F7.
% Corollary I.9.12 begins at line 2004 and is excluded.
% Source correction: line 1995 and the original printing have
% X_i\overset{v}{=}\bar{U}_i; the corrected/current reading uses
% X_i\overset{\mathrm{def}}{=}\bar{U}_i as the definition marker.

Let \(U\) be the set of points at which \(f\) is étale.  It is open,
and it suffices to show that it is also closed and nonempty.  The set
\(U\) contains the inverse image of the generic point of~\(Y\) (for an
algebra over a field, unramified is equivalent to étale); this inverse
image is nonempty because \(X\) dominates~\(Y\).  If \(x\) lies in the
closure of~\(U\), then it lies in the closure of an irreducible component
\(U_i\) of~\(U\), hence on an irreducible component
\(X_i\overset{\mathrm{def}}{=}\overline{U_i}\)\footnote{The French print
and diplomatic reading place the unbound superscript \(v\) above the
equality sign. Since the sentence defines \(X_i\) as the closure of
\(U_i\), the corrected/current branch uses the definition marker
\(\mathrm{def}\).}
of~\(X\) that meets~\(U\), and consequently dominates~\(Y\) (for every
irreducible component of~\(U\), being flat over~\(Y\), dominates~\(Y\)).
Let \(y=f(x)\) be the image of~\(x\) in~\(Y\).  Then the homomorphism
\(\mathcal O_y\to\mathcal O_x\) is \emph{injective} (since \(\mathcal O_y\)
is an integral domain).  Since \(\mathcal O_y\) is normal and
\(\mathcal O_y\to\mathcal O_x\) is unramified, Theorem~\ref{I.9.5}(ii)
shows that this local homomorphism is étale.  Thus \(x\in U\), so \(U\) is
closed.  Since \(U\) is nonempty and \(X\) is connected, \(U=X\), which
proves the corollary.
\end{filecontents*}
% Original input: drafts/SGA1_I_9_12_English_source_draft.texfrag
\begin{filecontents*}[overwrite,noheader]{sga1-cumulative-input-056.tex}
% SGA 1, Expose I, Corollary I.9.12 statement.
% French authority envelope: lines 2004--2013,
% 10 LF lines / 371 bytes; LF-normalized SHA-256:
% E4DF7B867508B98A3F1CB594EE426FD48C4A1A5139C5D1F63CFDC28A3057BA9E.
% The associated parenthetical note begins at line 2014 and is excluded.

\begin{corollary}
\label{I.9.12}
Let \(f\colon X\to Y\) be a dominant morphism of finite type, with \(Y\)
normal and \(X\) irreducible.  Then the set of points at which \(f\) is
étale coincides with the complement of the support of
\(\Omega^1_{X/Y}\), that is, with the complement of the closed
subprescheme of~\(X\) defined by the different ideal
\(\mathfrak d_{X/Y}\).
\end{corollary}
\end{filecontents*}
% Original input: drafts/SGA1_I_9_12_NOTE_English_source_draft.texfrag
\begin{filecontents*}[overwrite,noheader]{sga1-cumulative-input-057.tex}
\marginpar{\raggedright\footnotesize French source p.~21}%
(This is the ``less trivial'' statement alluded to in the remark in
no.~\ref{I.4}.)
\end{filecontents*}
% Original input: drafts/SGA1_I_9_12_STARRED_REMARK_English_source_draft.texfrag
\begin{filecontents*}[overwrite,noheader]{sga1-cumulative-input-058.tex}
\paragraph{Remark \textup{*}.}
One should beware of assuming that a connected étale covering of an
irreducible scheme is itself irreducible when the base is not assumed to be
normal. This question will be studied in \hyperref[I.11]{no.~I.11}.
\end{filecontents*}
% Original input: drafts/SGA1_I_9_1_English_source_draft.texfrag
\begin{filecontents*}[overwrite,noheader]{sga1-cumulative-input-059.tex}
% SGA 1, Expose I, section I.9 opening through Proposition I.9.1.
% Semantic French authority body: lines 1654--1702, 49 lines / 2,276 bytes;
% LF-normalized SHA-256:
% 17A095936F0E25B8F86A22134EE44398DCD7E6F5712520BA7758FAEB4AC81A08.
% The bounded structural envelope also retains blank separator line 1703:
% lines 1654--1703, 50 lines / 2,277 bytes; LF-normalized SHA-256
% 0AF0545B813A2676BC7B94A917B5C5BDE113D7391D26A783C8A6460E2D7FAA2C.

\clearpage
\section{Permanence properties}
\label{I.9}

Let \(A\to B\) be a local étale homomorphism. We examine here several
cases in which a given property of~\(A\) implies the same property
for~\(B\), or conversely.
\marginpar{\raggedright\footnotesize French source p.~16}%
A number of such assertions already follow from the simple fact that \(B\)
is \emph{quasi-finite} and \emph{flat} over~\(A\), and we shall confine
ourselves to ``recalling'' a few of them: \emph{\(A\) and~\(B\) have the
same Krull dimension and the same depth} (Serre's ``cohomological
codimension'', in the terminology still in use). It follows, for example,
that \emph{\(A\) is Cohen--Macaulay if and only if \(B\) is}.
Moreover, for every prime ideal~\(\mathfrak q\) of~\(B\) whose contraction
to~\(A\) is~\(\mathfrak p\), the ring \(B_{\mathfrak q}\) is again
quasi-finite and flat over~\(A_{\mathfrak p}\), provided that \(B\) is a
localization of an \(A\)-algebra of finite type. This follows from the fact
that, for a morphism of finite type, the locus where it is quasi-finite,
respectively flat, is open. Furthermore, \emph{every} prime ideal~\(\mathfrak p\)
of~\(A\) is the contraction of a prime ideal~\(\mathfrak q\) of~\(B\), since \(B\)
is \emph{faithfully} flat over~\(A\). It follows, for example, that
\emph{\(\mathfrak p\) and~\(\mathfrak q\) have the same height}, and also
that \emph{\(A\) has no embedded prime ideal if and only if \(B\) has none}.

We shall therefore restrict ourselves to assertions that are more specific
to étale morphisms.

\begin{proposition}
\label{I.9.1}
Let \(A\to B\) be a local étale homomorphism. Then \(A\) is regular if and
only if \(B\) is regular.
\end{proposition}

Indeed, let \(k\) be the residue field of~\(A\), and let \(L\) be that
of~\(B\). Since \(B\) is flat over~\(A\) and
\(L=B\otimes_A k\), that is, \(\mathfrak n=\mathfrak mB\), where \(\mathfrak m\) and
\(\mathfrak n\) are the maximal ideals of \(A\) and~\(B\), respectively,
the \(\mathfrak m\)-adic filtration on~\(B\) is identical to its
\(\mathfrak n\)-adic filtration, and hence
\[
  \operatorname{gr}^{*}(B)
  =\operatorname{gr}^{*}(A)\otimes_k L.
\]
It follows that \(\operatorname{gr}^{*}(B)\) is a polynomial algebra
over~\(L\) if and only if \(\operatorname{gr}^{*}(A)\) is a polynomial
algebra over~\(k\). This proves the proposition. (N.B. We did not use
the fact that \(L/k\) is separable.)
\end{filecontents*}
% Original input: drafts/SGA1_I_9_2_English_source_draft.texfrag
\begin{filecontents*}[overwrite,noheader]{sga1-cumulative-input-060.tex}
% SGA 1, Expose I, Corollary I.9.2 and Proposition I.9.2.
% French authority body and numbering apparatus: lines 1704--1720,
% 17 lines / 538 bytes; LF-normalized SHA-256:
% A561F19005B1040B3B2C12C7FBD9EC20B7EB761D8AEBC331647E5B1170ED9562.
% The bounded structural envelope also retains blank separator line 1721:
% lines 1704--1721, 18 lines / 539 bytes; LF-normalized SHA-256:
% 55F044CF7103997F3B2705898EAAF4CBDD40B4551C44A1A10C3C5509AE1D5980.

\begin{corollary}
\label{I.9.2}
Let \(f\colon X\to Y\) be an étale morphism. If \(Y\) is regular, then
\(X\) is regular; the converse holds if \(f\) is surjective.
\end{corollary}

% The corrected French TeX records and retains the original duplicate I.9.2.
\setcounter{proposition}{1}

% Preserve the visible duplicate source number while giving the second
% statement a distinct internal Hyperref destination.
\begingroup
\renewcommand{\theHproposition}{\theHsection.\arabic{proposition}.second}
\begin{proposition}
\label{prop:I.9.2}
Let \(f\colon X\to Y\) be an étale morphism. If \(Y\) is reduced, then
so is \(X\); the converse holds if \(f\) is surjective.
\end{proposition}
\endgroup
\end{filecontents*}
% Original input: drafts/SGA1_I_9_3_English_source_draft.texfrag
\begin{filecontents*}[overwrite,noheader]{sga1-cumulative-input-061.tex}
% SGA 1, Expose I, bridge, Corollary I.9.3, and proof.
% French authority body: lines 1722--1759,
% 38 LF lines / 1,881 bytes; LF-normalized SHA-256:
% A51833154B4AB4DB10E305F2D6B04776E069B42D777B79DC8116F3D6F577DCCD.
% The bounded structural envelope also retains blank separator line 1760:
% 39 LF lines / 1,882 bytes; LF-normalized SHA-256:
% 90DDCDB85D7271081E2E1C3964F66E0593E5881FB157D9FCA9ACEC7619674D0D.
% Page-mapping note: the corrected French TeX places its p. 17 marker after
% "isomorphe \`a une", whereas the original print breaks before "\`a une".

This is equivalent to the following corollary.

\begin{corollary}
\label{I.9.3}
Let \(f\colon A\to B\) be a local étale homomorphism, with \(B\) isomorphic
\marginpar{\raggedright\footnotesize French source p.~17}%
to a localization of an \(A\)-algebra of finite type. Then \(A\) is reduced
if and only if \(B\) is reduced.
\end{corollary}

If \(B\) is reduced, then \(A\) is reduced, since \(A\to B\) is injective
(because \(B\) is faithfully flat over~\(A\)). Conversely, suppose that
\(A\) is reduced, and let \(\mathfrak p_i\) be the minimal prime ideals
of~\(A\). By hypothesis, the natural map
\(A\to\prod_i A/\mathfrak p_i\) is injective. Tensoring with the flat
\(A\)-module~\(B\) shows that
\(B\to\prod_i B/\mathfrak p_iB\) is injective, so it is enough to prove
that the rings \(B/\mathfrak p_iB\) are reduced. Since
\(B/\mathfrak p_iB\) is étale over \(A/\mathfrak p_i\), we are reduced to
the case where \(A\) is an integral domain.

Let \(K\) be its field of fractions. Since \(A\to K\) is injective, so is
\(B\to B\otimes_A K\), because \(B\) is flat over~\(A\). It therefore
remains to prove that the latter ring is reduced. Now \(B\), being a
localization of an \(A\)-algebra of finite type, is the local ring at a
point~\(x\) of a scheme \(X=\operatorname{Spec}(C)\) of finite type and
étale over \(Y=\operatorname{Spec}(A)\). Consequently
\(B\otimes_A K\) is a localization, with respect to a suitable
multiplicatively closed set, of the ring \(C\otimes_A K\) of
\(X\otimes_A K\). Since \(X\otimes_A K\) is étale over~\(K\), its ring
is a finite product of fields that are finite separable extensions of~\(K\);
the same is therefore true of \(B\otimes_A K\). This proves the corollary.
\end{filecontents*}
% Original input: drafts/SGA1_I_9_4_English_source_draft.texfrag
\begin{filecontents*}[overwrite,noheader]{sga1-cumulative-input-062.tex}
% SGA 1, Expose I, Corollary I.9.4 and proof.
% French authority envelope: lines 1761--1775,
% 15 LF lines / 486 bytes; LF-normalized SHA-256:
% 6F036BA67752677BC70C7C4F53D7714C05498AF26878519EF15B15B88CC7F3A2.
% Theorem I.9.5 begins at line 1776 and is excluded.
% Corrected-source note: orig=false selects \hat{B} at line 1772; the
% original-print branch has B and is retained only as an adverse witness.

\begin{corollary}
\label{I.9.4}
Let \(f\colon A\to B\) be a local étale homomorphism, and suppose that
\(A\) is analytically reduced (that is, the completion \(\widehat A\)
of~\(A\) has no nilpotent elements). Then \(B\) is analytically reduced,
and a fortiori reduced.
\end{corollary}

Indeed, \(\widehat B\) is \emph{finite} and étale over
\(\widehat A\), so Corollary~\ref{I.9.3} applies.
\end{filecontents*}
% Original input: drafts/SGA1_I_9_5_English_source_draft.texfrag
\begin{filecontents*}[overwrite,noheader]{sga1-cumulative-input-063.tex}
% SGA 1, Expose I, Theorem I.9.5 statement only.
% French authority body: lines 1776--1787,
% 12 LF lines / 429 bytes; LF-normalized SHA-256:
% 1BD713844A82B14BE64A5E711030856D7A6B3B42DC4E8A95881DBF7E38FCA9AC.
% The bounded structural envelope retains blank separator line 1788:
% lines 1776--1788, 13 LF lines / 430 bytes; LF-normalized SHA-256:
% 3399E7A36B652770D1C166309F34D8294196B07D28FFBA5E4037C30EBF41C179.
% The explanatory proof preface begins at line 1789 and is excluded.

\begin{theorem}
\label{I.9.5}
Let \(f\colon A\to B\) be a local homomorphism, with \(B\) isomorphic
to a localization of an \(A\)-algebra of finite type. Then:
\begin{enumerate}
\item[(i)] If \(f\) is étale, then \(A\) is normal if and only if \(B\)
is normal.
\item[(ii)] If \(A\) is normal, then \(f\) is étale if and only if
\(f\) is injective and unramified (and in that case \(B\) is normal by
\textup{(i)}).
\end{enumerate}
\end{theorem}

\end{filecontents*}
% Original input: drafts/SGA1_I_9_5_P1_English_source_draft.texfrag
\begin{filecontents*}[overwrite,noheader]{sga1-cumulative-input-064.tex}
% SGA 1, Expose I, Theorem I.9.5 proof roadmap and first proof of (i).
% French authority semantic body: lines 1789--1828,
% 40 LF lines / 1,781 bytes; LF-normalized SHA-256:
% A22F9904545AE340B9A4F60251CCCDDD23E25D69682CB37A3B7F455954CE8AB0.
% The bounded structural envelope retains blank separator line 1829:
% lines 1789--1829, 41 LF lines / 1,782 bytes; LF-normalized SHA-256:
% CD924FE11494EAA408AEC7FB8312BAC88A77254C15A646CE37BEB9224329280A.
% The second-proof heading begins at line 1830 and is excluded.

We shall give two different proofs of~\textup{(i)}.  The first uses certain
properties of quasi-finite flat morphisms (recalled at the beginning of this
number) without using Theorem~\ref{I.7.6} (and hence the Main Theorem); in the
second proof the roles are reversed.  Finally, for~\textup{(ii)}, it seems
that the Main Theorem is needed in any case.

\subsubsection*{First proof}
We use the following necessary and sufficient condition
\marginpar{\raggedright\footnotesize French source p.~18}%
for a noetherian local ring~\(A\) of nonzero dimension to be normal.

\begin{serrecriterion}
\begin{enumerate}
\item[(i)] For every prime ideal~\(\mathfrak p\) of~\(A\) of height~\(1\),
\(A_{\mathfrak p}\) is normal (or, equivalently, regular).
\item[(ii)] For every prime ideal~\(\mathfrak p\) of~\(A\) of
height~\(\geq 2\), one has
\(\operatorname{depth} A_{\mathfrak p}\geq 2\).\footnote{Cf.\ EGA~IV~5.8.6.}
\end{enumerate}
\end{serrecriterion}

We shall admit this criterion here; it is supposed to appear in the section
on flat morphisms.  Its principal advantage is that it does not assume
a priori that~\(A\) is reduced, nor a fortiori that it is an integral domain.
Here we may already assume
\(\dim A=\dim B\neq 0\).

By the facts recalled at the beginning of this number, the prime
ideals~\(\mathfrak p\) of~\(A\) of height~\(1\) (respectively,
height~\(\geq 2\)) are exactly the contractions to~\(A\) of the prime
ideals~\(\mathfrak q\) of~\(B\) of height~\(1\) (respectively,
height~\(\geq 2\)).  Finally, if~\(\mathfrak p\) and~\(\mathfrak q\)
correspond, then \(B_{\mathfrak q}\) is étale over~\(A_{\mathfrak p}\),
so it has the same depth as~\(A_{\mathfrak p}\), and it is regular if and
only if~\(A_{\mathfrak p}\) is regular (Proposition~\ref{I.9.1}).  Applying
Serre's criterion, we find that~\(A\) is normal if and only if~\(B\) is
normal.

\end{filecontents*}
% Original input: drafts/SGA1_I_9_5_P2_English_source_draft.texfrag
\begin{filecontents*}[overwrite,noheader]{sga1-cumulative-input-065.tex}
% SGA 1, Expose I, Theorem I.9.5(i), complete second proof.
% French authority semantic body: lines 1830--1951,
% 122 LF lines / 5,185 bytes; LF-normalized SHA-256:
% 224F6855592EF1D4CC1BEAD7725D05D60FE20D249A082BDFA60DAD18EDD3DB51.
% The bounded structural envelope retains blank separator line 1952:
% lines 1830--1952, 123 LF lines / 5,186 bytes; LF-normalized SHA-256:
% 1D670DBFC69B9D9BDF1AFD653431708C748E71F9CB3FDB3A7B753D86D25E2EEC.
% The proof of item (ii) begins at line 1953 and is excluded.

\subsubsection*{Second proof}
Suppose that~\(B\) is normal, and let~\(L\) be its field of fractions
and~\(K\) that of~\(A\) (the ring~\(A\) is an integral domain because
\(B\) is).  We saw in the proof of Corollary~\ref{I.9.3} that
\(B\otimes_A K\) is a finite product of fields.  Since it is contained
in~\(L\), it is a field, and since it contains~\(B\), it is~\(L\).
An element of~\(K\) that is integral over~\(A\) is integral over~\(B\),
so it belongs to~\(B\), since~\(B\) is normal.  It therefore belongs
to~\(A\), because \(B\cap K=A\) (as follows from the faithful flatness
of~\(B\) over~\(A\)).

Now suppose that~\(A\) is normal; we prove that~\(B\) is normal.  By
Theorem~\ref{I.7.6}, we have \(B=B'_{\mathfrak n}\), where
\[
  B'=A[t]/F A[t],
\]
with~\(F\) and~\(\mathfrak n\) as in Theorem~\ref{I.7.6}.  Thus
\(L=B\otimes_A K\) is a localization of
\[
  B'\otimes_A K=K[t]/F K[t]
\]
and is a product of fields, each a finite separable extension of~\(K\).
This product, as happens whenever an artinian ring---here
\(B'_K=B'\otimes_A K\)---is localized at a multiplicatively closed set,
is a direct factor of~\(B'_K\).  It therefore corresponds to a factorization
\(F=F_1F_2\) in~\(K[t]\), and the generator of~\(L\) corresponding to~\(t\)
is already annihilated by~\(F_1\).

But \emph{because~\(A\) is normal, the polynomials~\(F_i\) belong
to~\(A[t]\)} (when they are taken to be monic).  The map
\(B\to L=B\otimes_A K\) is injective, since \(A\to K\) is injective
and~\(B\) is flat over~\(A\).  Consequently, we already have
\(F_1(u)=0\), where~\(u\) is the class of~\(t\) in~\(L\).  If~\(F\)
has been chosen of minimal degree, it follows that \(F_2=1\).  (N.B. one
has
\[
  F'(u)=F'_1(u)F_2(u)+F_1(u)F'_2(u)=F'_1(u)F_2(u),
\]
because \(F_1(u)=0\); hence \(F'_1(u)\neq0\), since \(F'(u)\neq0\).)

Thus
\marginpar{\raggedright\footnotesize French source p.~19}%
\begin{equation*}
\label{eq:I.9.5.*}
\tag{*}
  L=B\otimes_A K=K[t]/F K[t].
\end{equation*}
Consequently, \(F\) is separable over~\(K\), though of course not
necessarily separable over~\(A\).\footnote{The machine-readable French preimage
and the original printing both read \emph{par nécessairement}; the corrected/current
reading is \emph{pas nécessairement}. Since \(F\in A[t]\) was already fixed,
the contrast concerns separability over~\(K\) and over~\(A\), not membership
in the two polynomial rings.}  (N.B. for the moment, we have
essentially shown only that in Theorem~\ref{I.7.6} one may choose~\(F\)
and~\(\mathfrak n\) in such a way that, with the notation used here, the map
\[
  B'\longrightarrow B'_{\mathfrak n}=B
\]
is \emph{injective}.  We used the normality of~\(A\) for this; I do not know
whether the assertion remains true without the normality hypothesis.)

We now recall the well-known lemma extracted from Serre's course of last
year.

\begin{lemma}
\label{I.9.6}
Let~\(K\) be a ring, let \(F\in K[t]\) be a monic separable polynomial,
let \(L=K[t]/F K[t]\), and let~\(u\) be the class of~\(t\) in~\(L\) (so
that \(F'(u)\) is a unit of~\(L\)).  Then, with \(n=\deg F\),
\begin{align*}
  \operatorname{Tr}_{L/K}\!\left(\frac{u^i}{F'(u)}\right)&=0
    &&\text{if }0\leq i<n-1,\\
  \operatorname{Tr}_{L/K}\!\left(\frac{u^{n-1}}{F'(u)}\right)&=1.
\end{align*}
\end{lemma}

\begin{corollary}
\label{I.9.7}
The determinant of the matrix
\[
  \left(
    \operatorname{Tr}_{L/K}\!\left(\frac{u^j u^i}{F'(u)}\right)
  \right)_{0\leq i,j\leq n-1}
\]
is equal to \((-1)^{n(n-1)/2}\), and hence is a unit in every subring~\(A\)
of~\(K\).\footnote{The French \TeX{} and the original printing omit
\(\operatorname{Tr}_{L/K}\) from these matrix entries.  Read literally,
the printed matrix is an outer-product matrix of rank at most one, so for
\(n>1\) its determinant is zero, contrary to the stated unit.  Lemma~\ref{I.9.6}
and Corollary~\ref{I.9.8} determine the intended trace-pairing matrix; the
trace operator is restored here with the source defect disclosed.  The
original printing also gives the exponent~\(n\); the active French source
uses \(n(n-1)/2\), and its comment records Serre's correction of
27~November~2002.}
\end{corollary}

\begin{corollary}
\label{I.9.8}
Let~\(A\) be a subring of~\(K\), and suppose that \(F\in A[t]\).\footnote{The
French \TeX{} and the original printing omit the hypothesis \(F\in A[t]\)
from this corollary.  It is needed for the asserted basis.  For example, take
\(A=\mathbb Z\), \(K=\mathbb Q\), and \(F=t(t-\tfrac12)\).  Then
\(L\simeq\mathbb Q\times\mathbb Q\) and \(u/F'(u)=(0,1)\), but pairing with
\(u\) has trace \(1/2\notin\mathbb Z\).  Corollary~\ref{I.9.9} imposes the
missing coefficient hypothesis, so the application below is unaffected.}
Let~\(V\) be the \(A\)-submodule of~\(L\) generated by the elements \(u^i\),
\(0\leq i\leq n-1\), and let~\(V'\)
be the \(A\)-submodule of~\(L\) consisting of the elements \(x\in L\)
such that
\[
  \operatorname{Tr}_{L/K}(xy)\in A
\]
for every \(y\in V\) (that is, for \(y=u^i\), \(0\leq i\leq n-1\)).
Then~\(V'\) has as an \(A\)-basis the elements
\[
  \frac{u^i}{F'(u)},\qquad 0\leq i\leq n-1.
\]
\end{corollary}

\begin{corollary}
\label{I.9.9}
Suppose that~\(K\) is the field of fractions of a normal integral
domain~\(A\), and that the coefficients of~\(F\) belong to~\(A\).  With
the notation of Corollary~\ref{I.9.8}, the module~\(V'\) then contains the
integral closure~\(A'\) of~\(A\) in~\(L\).  Consequently,
\[
  A'\subseteq \frac{A[u]}{F'(u)}=F'(u)^{-1}A[u]
  \subseteq A[u][F'(u)^{-1}].
\]
\end{corollary}

Apply the last corollary to the situation obtained in the proof.  Since
\(F'(u)\) is a unit in~\(B\), and~\(B\) contains~\(A[u]\), the ring~\(B\)
contains~\(A'\).  By the Main Theorem (or directly from
\(B=A[u]_{\mathfrak n}\)), \(B\) is a localization of~\(A'\).  Since
\(A'\) is normal, so is~\(B\).
\end{filecontents*}
% Original input: drafts/SGA1_I_9_5_P3_English_source_draft.texfrag
\begin{filecontents*}[overwrite,noheader]{sga1-cumulative-input-066.tex}
\subsubsection*{Proof of \textup{(ii)}}
We
\marginpar{\raggedright\footnotesize French source p.~20}%
proceed as in the preceding proof to show that, in
Theorem~\ref{I.7.6}, \(F\) can be chosen so that
\eqref{eq:I.9.5.*} still holds.  The only a priori obstacle is that,
since \(B\) is no longer assumed flat over~\(A\), we can no longer
assert that \(B\to L\) is injective; hence, a priori, the argument
applies only to the image \(B_1\) of~\(B\) under that homomorphism.
It follows at once that \(B_1\) is flat over~\(A\) (being a localization
of an \(A\)-algebra that is free as an \(A\)-module).  By
Corollary~\ref{I.4.8}, the morphism \(B\to B_1\) is étale, hence an
isomorphism, which completes the proof.

(From the editorial point of view, the last two proofs should be
interchanged, and the formal calculations of the lemma and its
corollaries should be put in a separate number.)
\end{filecontents*}
% Original input: drafts/SGA1_VIII_1_A_English.texfrag
\begin{filecontents*}[overwrite,noheader]{sga1-cumulative-input-067.tex}
\chapter{Faithfully flat descent}
\label{VIII}
% page: 195
\marginpar{195}
%

\section{Descent of quasi-coherent modules}
\label{VIII.1}

Let $\Sch$\label{indnot:hb} denote the category of preschemes.  Proceeding as
in~\ifthenelse{\boolean{orig}}{VI~}{}\Ref{VI.11}.b, one finds that the category of pairs $(X,F)$
consisting of a prescheme~$X$ and a module~$F$ on~$X$ (with morphisms
defined as in \loccit by means of the direct image of a module under a
morphism of ringed spaces) may be regarded as a fibered category
over~$\Sch$.  The base-change functor associated with a morphism
$f\colon X\to Y$ in~$\Sch$ is the inverse-image functor for modules
under~$f$.  Notice that the fiber category at
$X\in\Ob(\Sch)$ of the preceding fibered category is the category
\emph{opposite} to the category of modules on~$X$.

Since the inverse image of a quasi-coherent module is quasi-coherent,
the full subcategory of the category of pairs $(X,F)$ formed by those
pairs for which $F$ is quasi-coherent is a fibered subcategory of the
preceding fibered category.  (By contrast, without hypotheses on~$f$,
the direct image of a quasi-coherent module is not in general a
quasi-coherent module.)  We shall simply call this the \emph{fibered
category of quasi-coherent modules on preschemes}.

Recall also that a morphism $f\colon X\to Y$ of ringed spaces is called
\emph{faithfully flat} if it is \emph{flat} (that is, for every
$x\in X$, $\mathcal{O}_{X,x}$ is a flat module over
$\mathcal{O}_{Y,f(x)}$, \eqref{IV}) and \emph{surjective}.  The
morphism~$f$ is called \emph{quasi-compact} if the inverse image under
$f$ of every quasi-compact subset is quasi-compact.  When $f$ is a
morphism of preschemes, this also means that the inverse image under~$f$
of every affine open subset of~$Y$ is a \emph{finite} union of affine
open subsets of~$X$.

\begin{theorem}
\label{VIII.1.1}
Let
% page: 196
\marginpar{196}
%
$\cal{F}$ be the fibered category of quasi-coherent modules on
preschemes.  Let $g\colon S'\to S$ be a faithfully flat and
quasi-compact morphism of preschemes.  Then $g$ is a morphism of
effective $\cal{F}$-descent.
\end{theorem}

Recall\footnote{We shall admit here the general theory of descent set
out in detail in the article of J.~\textsc{Giraud} cited in the
footnote to the Warning, a work that we shall cite below as
\cite{VIII.D}.  Cf.~also \cite{VIII.2} for a succinct account.} that
this means two things.

\begin{corollary}[Descent of homomorphisms of modules]
\label{VIII.1.2}
Let $g\colon S'\to S$ be a faithfully flat and quasi-compact morphism
of preschemes, let $F$ and $G$ be two quasi-coherent modules on~$S$,
let $F'$ and $G'$ be their inverse images on~$S'$, and finally let
$F''$ and $G''$ be their inverse images on
$S''=S'\times_SS'$.  Consider the diagram of maps of sets defined by
the base-change functors for $g$, $p_1$, and $p_2$, where
$\xymatrix@C=.5cm{p_1,p_2\colon
S'\times_SS'\ar@<2pt>[r]\ar@<-2pt>[r]&S'}$ are the two projections:
$$
\xymatrix@C=.5cm{
\Hom_S(F,G)\ar[r]&
\Hom_{S'}(F',G')\ar@<2pt>[r]\ar@<-2pt>[r]&
\Hom_{S''}(F'',G'').}
$$
This diagram is exact; that is, it identifies the first set
bijectively with the equalizer of the two displayed maps from the
second set to the third.
\end{corollary}

In other words, the base-change functor for~$g$,
$F\mto F'$, defines a \emph{fully faithful} functor from the category
of quasi-coherent modules on~$S$ to the category of quasi-coherent
modules on~$S'$ endowed with descent data relative to~$g$.  Moreover:

\begin{corollary}[Descent of modules]
\label{VIII.1.3}
For every quasi-coherent module~$F'$ on~$S'$, every descent datum on
$F'$ relative to~$g$ is \emph{effective}; that is, $F'$, together with
its descent datum, is isomorphic to the inverse image under~$g$ of a
quasi-coherent module on~$S$ (determined up to unique isomorphism
by~\Ref{VIII.1.2}).
\end{corollary}

In other words, the preceding fully faithful functor is in fact an
\emph{equivalence}.  In practical terms, giving a quasi-coherent module
on~$S$ amounts to the same thing as giving a quasi-coherent module
on~$S'$ endowed with descent data relative to~$g$.

\subsubsection*{Proof of~\Ref{VIII.1.1}}
Let
% page: 197
\marginpar{197}
%
$T$ first be an $S$-prescheme that is $S$-isomorphic to the disjoint
union of a family of open subpreschemes $S_i$ induced from~$S$ and
covering~$S$.  It is then evident that the structural morphism
$T\to S$ is a morphism of effective $\cal{F}$-descent.  This means
precisely that giving a quasi-coherent module~$F$ on~$S$ is equivalent
to giving quasi-coherent modules $F_i$ on the~$S_i$, together with
gluing isomorphisms
$$
\varphi_{ji}\colon
F_i|S_i\cap S_j\longrightarrow F_j|S_i\cap S_j
$$
satisfying the familiar cocycle condition.  By VII,~8, it follows that,
to verify that $g\colon S'\to S$ is a morphism of effective
$\cal{F}$-descent, it suffices to verify this for the morphism
$$
g_T\colon T'=T\times_SS'\longrightarrow T
$$
obtained from~$g$ by the base change $T\to S$.  Notice that the
hypothesis on $T\to S$ is stable under arbitrary base change, so that
$T\to S$ is in fact a \emph{universal} morphism of effective
$\cal{F}$-descent.  Taking the~$S_i$ to be affine open subsets
covering~$S$, we are therefore reduced to the case where~$S$ is affine.

Then $S'$ is a finite union of affine open subsets.  Taking their
disjoint union as an $S$-scheme gives an affine $S$-scheme~$S_1$ and a
flat surjective $S$-morphism $S_1\to S'$.  Thus $S_1$ is also faithfully
flat over~$S$.  If, therefore, one proves that a faithfully flat affine
morphism is a morphism of effective $\cal{F}$-descent, hence a strict
universal morphism of $\cal{F}$-descent (the hypothesis being stable
under base change), it follows in particular that the structural
morphism $S_1\to S$ is a strict universal morphism of
$\cal{F}$-descent.  Since there is an $S$-morphism $S_1\to S'$, it then
follows from~\cite{VIII.D} that $g\colon S'\to S$ is indeed a strict
morphism of $\cal{F}$-descent.

We are thus reduced to the case in which~$g$ is affine.  As we have
seen, we may then also suppose that~$S$ is affine; hence \emph{we may
suppose that $S$ and $S'$ are affine}.  In this case,
\Ref{VIII.1.2} is equivalent to the following lemma.

\begin{lemma}
\label{VIII.1.4}
Let $A$ be a ring, let $A'$ be a faithfully flat $A$-algebra, let $M$
and $N$ be two $A$-modules, let $M'$ and $N'$ be the $A'$-modules
obtained by the change of rings $A\to A'$, and let $M''$ and $N''$ be
the $A''=A'\otimes_AA'$-modules obtained by the change of rings
$A\to A''$.  Then the sequence of maps of sets
$$
\xymatrix@C=.5cm{
\Hom_A(M,N)\ar[r]&
\Hom_{A'}(M',N')\ar@<2pt>[r]\ar@<-2pt>[r]&
\Hom_{A''}(M'',N'')}
$$
is exact.
\end{lemma}

Since the homomorphism $N\to N'$ is injective (because $A'$ is
faithfully flat over~$A$), the first arrow is injective.  It remains to
prove that if an $A'$-homomorphism $u'\colon M'\to N'$ is compatible
with the descent data, then it comes from an $A$-homomorphism
$u\colon M\to N$.  This simply means that $u'$ maps the subset~$M$
of~$M'$ into the subset~$N$ of~$N'$.  The induced map
$u\colon M\to N$ is then automatically $A$-linear because $u'$ is
$A'$-linear, and one sees similarly that $u'$ is necessarily equal to
$u\otimes_AA'$.

Now, if $x\in M$, then $u'(x)$ belongs to the equalizer of the pair of
maps
$\xymatrix@C=.5cm{N'\ar@<2pt>[r]\ar@<-2pt>[r]&N''}$.
Thus, to prove~\Ref{VIII.1.4}, we are reduced to the following special
case (corresponding to $M=A$):
\end{filecontents*}
% Original input: drafts/SGA1_VIII_1_B_English.texfrag
\begin{filecontents*}[overwrite,noheader]{sga1-cumulative-input-068.tex}
\begin{corollary}
\label{VIII.1.5}
Let \(N\) be an \(A\)-module. Then the sequence of maps of sets
\[
  \xymatrix@C=.5cm{
  N\ar[r]&N'\ar@<2pt>[r]\ar@<-2pt>[r]&N''
  }
\]
is exact.
\end{corollary}

Let \(A_1\) be a faithfully flat \(A\)-algebra. To prove exactness of
the displayed sequence, it is enough to prove exactness after the base
change \(A\to A_1\). The resulting sequence is plainly the sequence
associated with the
\ifthenelse{\boolean{orig}}
{\(A\)-module}
{\(A_1\)-module}
\[
  N_1=N\otimes_AA_1
\]
and the \(A_1\)-algebra \(A_1'=A_1\otimes_AA'\). It is therefore
enough to find a faithfully flat \(A\)-algebra \(A_1\) for which
\[
  \Spec(A_1')\longrightarrow\Spec(A_1)
\]
is a strict \(\cal F\)-descent morphism. We may take \(A_1=A'\):
the preceding morphism then has a right inverse, and hence, by
\cite{VIII.D}, is an effective descent morphism for every category
fibered over~\(\Sch\).

It remains to show that if \(N'\) is an \(A'\)-module equipped with a
descent datum relative to \(A\to A'\), \ie with an isomorphism
\[
  \varphi\colon N_1'\isomto N_2'
\]
between the two modules obtained from \(N'\) by the two base changes
\[
  \xymatrix@C=.5cm{
  A'\ar@<2pt>[r]\ar@<-2pt>[r]&A'\otimes_AA',
  }
\]
then \(N'\), together with its descent datum,
% page: 199
\marginpar{199}
%
is isomorphic to a module of the form \(N\otimes_AA'\). In view
of~\Ref{VIII.1.5}, this statement is readily seen to be equivalent to
the following:

\begin{lemma}
\label{VIII.1.6}
Let \(N'\) be an \(A'\)-module equipped with a descent datum relative
to \(A\to A'\), where \(A'\) is an \(A\)-algebra. Let \(N\) be the
\(A\)-submodule of~\(N'\) formed by the elements \(x\) such that
\[
  \varphi(x\otimes_{A'}1_{A'})
  =
  1_{A'}\otimes_{A'}x,
\]
and consider the canonical homomorphism
\[
  N\otimes_AA'\longrightarrow N',
\]
which is compatible with the descent data. If \(A'\) is faithfully
flat over~\(A\), this homomorphism is an isomorphism.
\end{lemma}

Let us prove the lemma. Again let \(A_1\) be a faithfully flat
\(A\)-algebra. To show that the homomorphism in question is an
isomorphism, it is enough to show that it becomes one after the base
change
\ifthenelse{\boolean{orig}}
{\(A_1\to A\).}
{\(A\to A_1\).}
Using the flatness of~\(A_1\) over~\(A\), one sees that the resulting
homomorphism is precisely the one obtained directly from the
\(A_1'=A'\otimes_AA_1\)-module \(N'\otimes_AA_1\), equipped with the
descent datum relative to \(A_1\to A_1'\) obtained canonically by base
change from the descent datum on~\(N'\). Thus it is enough to choose a
faithfully flat \(A\)-algebra \(A_1\) such that
\(\Spec(A_1')\to\Spec(A_1)\) is an effective
\(\cal F\)-descent morphism. As above, take \(A_1=A'\). This proves
\Ref{VIII.1.6}, and hence also~\Ref{VIII.1.1}.

\begin{corollary}[Descent of sections of modules]
\label{VIII.1.7}
%
\index{descent of sections of modules|hyperpage}%
%
Let \(g\colon S'\to S\) be a faithfully flat and quasi-compact
morphism of preschemes. For every quasi-coherent module \(G\) on~\(S\),
let \(G'\) and \(G''\) be its inverse images on~\(S'\) and
\(S''=S'\times_SS'\), respectively. Then the diagram of homomorphisms
of modules on~\(S\)
\[
  \xymatrix@C=.5cm{
  G\ar[r]&g_*(G')\ar@<2pt>[r]\ar@<-2pt>[r]&h_*(G'')
  }
\]
where \(h\colon S''\to S\) is the structural morphism, is
\emph{exact}.
\end{corollary}

Indeed,
% page: 200
\marginpar{200}
%
this means that, for every open subset \(U\) of~\(S\), the
corresponding diagram of sections over~\(U\) is exact. We may plainly
assume \(U=S\), and the asserted exactness is then the special case
of~\Ref{VIII.1.2} obtained by taking \(F=\mathcal O_S\).

Since the inverse-image functor on modules is right exact,
\Ref{VIII.1.1} formally implies:

\begin{corollary}[Descent of quotient modules]
\label{VIII.1.8}
%
\index{descent of quotient modules|hyperpage}%
%
With the notation of~\Ref{VIII.1.7}, let
\(\operatorname{Quot}(F)\), for every quasi-coherent module \(F\) on a
prescheme, denote the set of quasi-coherent quotient modules of~\(F\).
Then the diagram of maps of sets
\[
  \xymatrix@C=.5cm{
  \operatorname{Quot}(G)\ar[r]&
  \operatorname{Quot}(G')\ar@<2pt>[r]\ar@<-2pt>[r]&
  \operatorname{Quot}(G'')
  }
\]
is exact.
\end{corollary}

(The same statement plainly holds with submodules in place of quotient
modules, since the two are in bijection.) Taking
\(G=\mathcal O_S\), we obtain:

\begin{corollary}[Descent of closed subpreschemes]
\label{VIII.1.9}
%
\index{descent of closed subpreschemes|hyperpage}%
%
For every prescheme \(X\), let \(H(X)\) denote the set of closed
subpreschemes of~\(X\). With this notation, and under the hypotheses
of~\Ref{VIII.1.7}, the diagram of maps of sets
\[
  \xymatrix@C=.5cm{
  H(S)\ar[r]&H(S')\ar@<2pt>[r]\ar@<-2pt>[r]&H(S'')
  }
\]
\emph{is exact}.
\end{corollary}

Theorem~\Ref{VIII.1.1} should be supplemented by the following result:

\begin{proposition}[Descent of properties of modules]
\label{VIII.1.10}
%
\index{descent of properties of modules|hyperpage}%
%
Let \(g\colon S'\to S\) be faithfully flat and quasi-compact, and let
\(F\) be a quasi-coherent module on~\(S\). Then \(F\) is of finite
type, respectively of finite presentation, respectively finite
locally free, if and only if its inverse image \(F'\) on~\(S'\) has
the corresponding property.
\end{proposition}

Only sufficiency requires proof. We may plainly suppose that \(S\) is
affine.
% page: 201
\marginpar{201}
%
Replacing \(S'\) by
\ifthenelse{\boolean{orig}}
{the}
{a}
sum of affine open subsets covering~\(S'\), we reduce to the case in
which \(S'\) is also affine. The assertion is then equivalent to the
following:

\begin{corollary}
\label{VIII.1.11}
Let \(A\) be a ring, let \(A'\) be a faithfully flat \(A\)-algebra,
let \(M\) be an \(A\)-module, and let
\(M'=M\otimes_AA'\). Then \(M\) is of finite type, respectively of
finite presentation, respectively finite locally free, if and only if
\(M'\) has the corresponding property.
\end{corollary}

Indeed, write \(M=\varinjlim_iM_i\), where the \(M_i\) are the
finite-type submodules of~\(M\). Then
\(M'=\varinjlim_iM_i'\). If \(M'\) is of finite type, it equals one
of the \(M_i'\); faithful flatness then gives \(M=M_i\), so \(M\) is
of finite type. Consequently, there is an exact sequence
\[
  0\longrightarrow R\longrightarrow L\longrightarrow M
  \longrightarrow0
\]
with \(L\) finite free, and hence an exact sequence
\[
  0\longrightarrow R'\longrightarrow L'\longrightarrow M'
  \longrightarrow0
\]
with \(L'\) finite free. If \(M'\) is finitely presented, then \(R'\)
is of finite type; by what precedes, \(R\) is of finite type, and
therefore \(M\) is finitely presented. Finally, \(M\) is finite
locally free precisely when it is finitely presented and flat
(\cf \Ref{IV} in the noetherian case; the general case is left to the
reader). Each of these properties descends, and so does their
conjunction. This completes the proof.

\begin{remarque}
\label{VIII.1.12}
Combining~\Ref{VIII.1.1} and~\Ref{VIII.1.10} shows that
\Ref{VIII.1.1} remains valid if the fibered category \(\cal F\) is
replaced by the fibered subcategory consisting of quasi-coherent
modules of finite type, respectively of finite presentation,
respectively finite locally free, respectively locally free of a
specified rank~\(n\).
\end{remarque}

\end{filecontents*}
% Original input: drafts/SGA1_VIII_2_English.texfrag
\begin{filecontents*}[overwrite,noheader]{sga1-cumulative-input-069.tex}
\section{Descent of affine preschemes over another}
\label{VIII.2}
% page: 202
\marginpar{202}
%

Since the inverse-image functor for modules is compatible with tensor
products and other tensor operations, Theorem~\Ref{VIII.1.1} implies
various variants obtained by considering, in place of a single
quasi-coherent module, a quasi-coherent module or a system of
quasi-coherent modules endowed with various additional structures
expressed by means of tensor operations.  For example, giving three
quasi-coherent modules $F$, $G$, and $H$ on~$S$ and a pairing
$$
F\otimes G\longrightarrow H
$$
is equivalent to giving three quasi-coherent modules $F'$, $G'$, and
$H'$ on~$S'$, endowed with descent data relative to
$g\colon S'\to S$, together with a pairing
$$
F'\otimes G'\longrightarrow H'
$$
``compatible'' with those descent data, in the evident sense of the
term.  For example, if $F=G=H$, one sees that giving a quasi-coherent
module~$F$ on~$S$ endowed with an algebra law (which for the moment we
do not require to satisfy any axiom of associativity, commutativity, or
the existence of a unit section) is equivalent to giving the same data
on~$S'$, together with descent data.  Using the results of the preceding
section, one sees at once that $F$ satisfies any one of the usual axioms
just mentioned if and only if $F'$ does.

For example, giving a quasi-coherent algebra $\mathcal{A}$ on~$S$ (by
which from now on we mean associative, commutative, and equipped with a
unit section) is equivalent to giving a quasi-coherent algebra
$\mathcal{A}'$ on~$S'$ endowed with descent data relative to
$g\colon S'\to S$.  Recalling the equivalence between the opposite
category of quasi-coherent algebras on~$S$ and the category of affine
$S$-preschemes over~$S$ (EGA~II, \S1), one obtains at once:

\begin{theorem}
\label{VIII.2.1}
Let $\cal{F}'$ be the fibered category of affine morphisms of
preschemes $f\colon X\to S$, regarded as a fibered subcategory of the
fibered category
% page: 203
\marginpar{203}
%
of arrows in the category $\Sch$ of preschemes
\textup{(\ifthenelse{\boolean{orig}}{VI~}{}\Ref{VI.11}.a)}.  Let $g\colon S'\to S$ be a faithfully
flat and quasi-compact morphism of preschemes.  Then $g$ is a morphism
of effective $\cal{F}'$-descent.
\end{theorem}

\end{filecontents*}
% Original input: drafts/SGA1_VIII_3_English.texfrag
\begin{filecontents*}[overwrite,noheader]{sga1-cumulative-input-070.tex}
\section[Descent of set-theoretic properties]{Descent of set-theoretic
properties and finiteness properties of morphisms\protect\footnotemark}
\label{VIII.3}
\footnotetext{\hypertarget{Hfootnote.2}{}For other results of the kind treated in Sections~3
and~4, cf.~EGA~IV~2.3, 2.6, and~2.7.}

\begin{proposition}
\label{VIII.3.1}
Let $f\colon X\to Y$ be an $S$-morphism, let $g\colon S'\to S$ be a
surjective morphism, and let
$$
f'\colon X'=X\times_SS'\longrightarrow Y'=Y\times_SS'
$$
be the morphism obtained from~$f$ by base change along
$g\colon S'\to S$.  The morphism $f$ is surjective (\resp radicial) if
and only if $f'$ is so.
\end{proposition}

Notice that $f'$ can also be obtained by the base change $Y'\to Y$,
which is itself surjective because it is obtained from the surjective
morphism $g\colon S'\to S$.  On the other hand, for every $y\in Y$ and
every $y'\in Y'$ lying over~$y$, there is an isomorphism
$$
X'_{y'}\isomfrom X_y\otimes_{\kres(y)}\kres(y'),
$$
where $X_y$ denotes the fiber of~$X$ at~$y$, and $X'_{y'}$ that of
$X'$ at~$y'$.  It follows that $X_y$ is nonempty (\resp has at most one
point, and that point corresponds to a radicial residue-field extension)
if and only if $X'_{y'}$ has the same property.  This proves
\Ref{VIII.3.1}.

\begin{corollary}
\label{VIII.3.2}
Under the hypotheses of~\Ref{VIII.3.1}, if $f'$ is injective
(\resp bijective), then $f$ is likewise.
\end{corollary}

Indeed, if $X'_{y'}$ has at most one point (\resp exactly one point),
the same is true of~$X_y$.  This follows because the morphism
$X'_{y'}\to X_y$ is surjective, being obtained from the surjective
morphism $\Spec(\kres(y'))\to\Spec(\kres(y))$.

\begin{proposition}
\label{VIII.3.3}
With the notation of~\Ref{VIII.3.1}, suppose that
$g\colon S'\to S$ is surjective and quasi-compact (\resp faithfully
flat and quasi-compact).  The morphism $f$ is quasi-compact
(\resp of finite type) if and only if $f'$ is so.
\end{proposition}

Only the sufficiency needs to be proved.  We may clearly suppose
$S=Y$, since the hypothesis on $g\colon S'\to S$ is preserved for
$Y'\to Y$.  We may also suppose that~$Y$ is affine.  Then $Y'$ is
quasi-compact, and therefore $X'$ is quasi-compact, since $f'$ is so by
hypothesis.  Let $(X_i)_{i\in I}$ be a family of affine open subsets of
$X$ covering~$X$.  The $X'_i$ are open subsets covering~$X'$, so a
finite subfamily covers~$X'$.  Since $X'\to X$ is surjective, the
corresponding $X_i$ already cover~$X$.  Thus $X$ is quasi-compact, that
is, $f$ is quasi-compact.

Now suppose that $f'$ is of finite type, and let us prove that $f$ is
of finite type, assuming $g$ faithfully flat.  Replacing $Y'$ by the
disjoint union of a family of affine open subsets covering it, we may
suppose that $Y'$ is affine.  Finally, by the preceding argument $X$ is
covered by finitely many affine open subsets~$X_i$; it remains to show
that they are of finite type over~$Y$, knowing that $X'_i$ is of finite
type over~$Y'$.  This reduces us to the following result.

\begin{corollary}
\label{VIII.3.4}
Let $B$ be an $A$-algebra, let $A'$ be a faithfully flat $A$-algebra,
and let $B'=B\otimes_AA'$ be the $A'$-algebra obtained from~$B$ by
change of rings.  Then $B$ is of finite type if and only if $B'$ is of
finite type.
\end{corollary}

Again, only the sufficiency needs proof.  We have
$B=\varinjlim_i B_i$, where the $B_i$ are the finite-type subalgebras
of~$B$.  Hence $B'=\varinjlim_i B'_i$, and if $B'$ is of finite type
over~$A'$, then $B'$ is equal to one of the~$B'_i$.  Faithful flatness
then gives $B=B_i$, so $B$ is of finite type.

\begin{corollary}
\label{VIII.3.5}
Suppose again that the base-change morphism $g\colon S'\to S$ is
faithfully flat and quasi-compact.  The morphism $f$ is quasi-finite if
and only if $f'$ is quasi-finite.
\end{corollary}

Indeed, the property of being quasi-finite is, by definition, the
conjunction of being of finite type and having finite fibers.  Each of
these properties descends under~$g$: the first by~\Ref{VIII.3.3}, and
the second by the argument of~\Ref{VIII.3.1}, which uses only the
surjectivity of~$g$.

\begin{remarques}
\label{VIII.3.6}
Let $A$ be a ring and $X$ an $A$-prescheme.  One sees easily that
% page: 205
\marginpar{205}
%
the following conditions are equivalent:

\begin{enumerate}
\item[(i)] There exist a noetherian ring~$A_0$ (which, if desired, may
be assumed to be a finite-type subring of~$A$), an $A_0$-prescheme
$X_0$ of finite type, a homomorphism $A_0\to A$, and an $A$-isomorphism
$$
X\isomto X_0\times_{A_0}A.
$$

\item[(ii)] The diagonal morphism
$X\to X\times_{\Spec(A)}X$ is quasi-compact (a vacuous condition if
$X$ is separated over~$A$), and $X$ is the union of finitely many
affine open subsets~$X_i$ whose rings~$B_i$ are algebras of finite
presentation over~$A$; that is, quotients of polynomial algebras in
finitely many indeterminates by ideals of finite type.
\end{enumerate}

If $X$ is itself affine, with ring~$B$, these conditions simply mean
that $B$ is an algebra of finite presentation over~$A$.

A morphism $f\colon X\to Y$ is called a \emph{morphism of finite
presentation}, and $X$ is also said to be of finite presentation
over~$Y$, if $Y$ is the union of affine open subsets~$Y_i$ such that
$X|Y_i$, as a $Y_i$-prescheme, satisfies the preceding equivalent
conditions.  The same is then true of $X|Y'$ for \emph{every} affine
open subset~$Y'$ of~$Y$.  This property is stable under base change,
and the composite of two morphisms of finite presentation is again of
finite presentation.

With these notions in place, condition~(ii), by the same argument as
in~\Ref{VIII.1.10}, shows that Proposition~\Ref{VIII.3.3} remains valid
when ``of finite type'' is replaced by ``of finite presentation.''
\end{remarques}
\end{filecontents*}
% Original input: drafts/SGA1_VIII_4_English.texfrag
\begin{filecontents*}[overwrite,noheader]{sga1-cumulative-input-071.tex}
\section{Descent of topological properties}
\label{VIII.4}

\begin{theorem}
\label{VIII.4.1}
Let $g\colon Y'\to Y$ be a morphism, and let $Z$ be a subset of~$Y$.
Suppose that~$g$ is flat and that there exists a quasi-compact
morphism~$f\colon X\to Y$ such that~$Z=f(X)$ (N.B. if $Y$ is
noetherian, this last condition is implied by
\ifthenelse{\boolean{orig}}
{$Z$ being constructible}
{\og $Z$ being constructible\fg}). Then
$$
g^{-1}(\overline{Z})=\overline{g^{-1}(Z)}\quoi.
$$
\end{theorem}

\ifthenelse{\boolean{orig}}
% this is a long one!
{We may assume~$Y$ affine, and then $Y'$ affine. Since $Y$ is affine,
$X$ is a \emph{finite
% page: 206
\marginpar{206}
%
union of affine open subsets~$X_{i}$, and, replacing~$X$ by the
disjoint union of the~$X_{i}$, we may also assume~$X$ affine.
Let~$A,A',B$ be the rings of~$Y,Y',X$, let $B'=B\otimes_{A}A'$ be that
of~$X'=X\times_{Y}Y'$, let $I$ be the kernel of~$A\to B$, and let
$I'$ be the kernel of~$A'\to B'$. Thus the closed subsets of~$Y$
and~$Y'$ defined by these ideals are respectively the closure
of~$Z=f(X)$ and the closure of~$Z'=f'(X')=g^{-1}(Z)$. We wish to show
that the latter equals $g^{-1}(\overline{Z})$; this follows
from~$I'=IA'$, which in turn follows from the flatness of~$A'$ over~$A$}.}
%
{We may assume~$Y$ affine, and then $Y'$ affine. Since $Y$ is affine,
$X$ is a finite
% page: 206
\marginpar{206}
%
union of affine open subsets~$X_{i}$, and, replacing~$X$ by the
disjoint union of the~$X_{i}$, we may also assume~$X$ affine.
Let~$A,A',B$ be the rings of~$Y,Y',X$, let $B'=B\otimes_{A}A'$ be that
of~$X'=X\times_{Y}Y'$, let $I$ be the kernel of~$A\to B$, and let
$I'$ be the kernel of~$A'\to B'$. Thus the closed subsets of~$Y$
and~$Y'$ defined by these ideals are respectively the closure
of~$Z=f(X)$ and the closure of~$Z'=f'(X')=g^{-1}(Z)$. We wish to show
that the latter equals $g^{-1}(\overline{Z})$; this follows
from~$I'=IA'$, which in turn follows from the flatness of~$A'$ over~$A$.}
% correction: this part should not be underlined (note also that the
% theorem was not underlined; they underlined the wrong part).

\begin{corollary}
\label{VIII.4.2}
Let $g\colon Y'\to Y$ be a flat and quasi-compact morphism, and let
$Z'$ be a closed subset of~$Y'$ saturated for the set-theoretic
equivalence relation defined by~$g$. Then
$Z'=g^{-1}(\overline{g(Z')})$.
\end{corollary}

\ifthenelse{\boolean{orig}}{}
{\enlargethispage{.5cm}}%
Indeed, $Z'=g^{-1}(Z)$, with~$Z=g(Z')$. We may then
apply~\Ref{VIII.4.1}, noting that the condition imposed on~$Z$
in~\Ref{VIII.4.1} is indeed satisfied by taking for~$X$ the
prescheme~$Z'$ equipped with the reduced structure induced
by~$Y'$. (The fact that~$g$ is quasi-compact then ensures that the
induced morphism~$f\colon Z'\to Y$ is quasi-compact.)

Statement~\Ref{VIII.4.2} also means that \emph{the topology
on~$g(Y')$ induced by~$Y$ is the quotient of that on~$Y'$}. In
particular:

\begin{corollary}
\label{VIII.4.3}
Let $g\colon Y'\to Y$ be a faithfully flat and quasi-compact
morphism. Then $g$ makes~$Y$ a quotient topological space of~$Y'$,
\ie for a subset~$Z$ of~$Y$, $Z$ is closed (\resp open) if and only
if $Z'=g^{-1}(Z)$ is so.
\end{corollary}

Recall now that two elements $a,b$ of~$Y'$ have the same image in
$Y$ if and only if they are of the form $p_{1}(c),p_{2}(c)$ for a
suitable element $c$ of~$Y''=Y'\times_{Y}Y'$. It follows that if $g$
is surjective, there is an \emph{exact} diagram of sets
$$
\xymatrix@C=.5cm{\cal{P}(Y)\ar[r] & \cal{P}(Y')\ar@<2pt>[r]\ar@<-2pt>[r] &
\cal{P}(Y'')}\quoi,
$$
where, for every set~$E$, $\cal{P}(E)$ denotes its power set. With
this notation, \Ref{VIII.4.3} may also be interpreted as follows:

\begin{corollary}[Descent of open \resp closed subsets]
\label{VIII.4.4}
Let
% page: 207
\marginpar{207}%
%
%
\index{descent of open subsets (\resp closed subsets)|hyperpage}%
%
$g\colon Y'{\to} Y$ be as in~\Ref{VIII.4.3}. For every
prescheme~$X$, let $\Ouv(X)$ \resp $\Fer(X)$ be the set of its open
subsets \resp the set of its closed subsets. Then the following are
\emph{exact} diagrams of maps of sets (deduced from~$g$ and the two
projections of $Y''=Y'\times_{Y}Y'$):
\vspace*{-5pt}
\begin{gather*}
\xymatrix@C=.5cm{\Ouv(Y)\ar[r] & \Ouv(Y')\ar@<2pt>[r]\ar@<-2pt>[r] &
\Ouv(Y'')}\\[-3pt]
\xymatrix@C=.5cm{\Fer(Y)\ar[r] & \Fer(Y')\ar@<2pt>[r]\ar@<-2pt>[r] &
\Fer(Y'')}
\end{gather*}
\end{corollary}

The following supplements~\Ref{VIII.4.3}:

\begin{corollary}
\label{VIII.4.5}
Let $g\colon Y'\to Y$ be as in~\Ref{VIII.4.3}, and let~$Z$ be a
subset of~$Y$ for which there exists a quasi-compact morphism~$f\colon
X\to Y$ with image~$Z$ (for example, $Z$ constructible and $Y$
noetherian). Then $Z$ is a locally closed subset of~$Y$ if and only
if $Z'=g^{-1}(Z)$ is a locally closed subset of~$Y'$.
\end{corollary}

It suffices to prove the \og if\fg direction. Let $Y_{1}$ be the
closed subprescheme of~$Y$ given by the closure of~$Z$ equipped with
the induced reduced structure, and let~$Y_{1}'=Y_{1}\times_{Y}Y'$ be
the closed subprescheme of~$Y'$ that is the inverse image of~$Y_{1}$.
Its underlying set is the inverse image
$g^{-1}(Y_{1})=g^{-1}(\overline{Z})$, and hence equals
$\overline{Z'}$ by~\Ref{VIII.4.1}. Since $Z'$ is locally closed
in~$Y'$, it is open in~$\overline{Z'}$, and hence open in~$Y_{1}'$.
The latter is faithfully flat and quasi-compact over~$Y_{1}$, so
by~\Ref{VIII.4.3} we conclude that~$Z$ is open in~$Y_{1}$,
\ie in~$\overline{Z}$, which means that $Z$ is locally closed.

\begin{corollary}
\label{VIII.4.6}
Let $g\colon S'\to S$ be a faithfully flat and quasi-compact
morphism, let $f\colon X\to Y$ be an $S$-morphism, and let
$f'\colon X'\to Y'$ be the $S'$-morphism obtained from it by base
change. Suppose that~$f'$ is an open map (\resp a closed map,
\resp \ifthenelse{\boolean{orig}}
{quasi-compact}
{quasi-compact}
% correction: `quasi-compact' should be `quasi-compacte'.
and a homeomorphism into, \resp a homeomorphism onto); then $f$ has
the same property.
\end{corollary}

Since $Y'$ is faithfully flat and quasi-compact over~$Y$, we may
assume~$Y=S$. Let~$Q$
% page: 208
\marginpar{208}
%
be a subset of~$X$. Then (denoting by $h$ the projection morphism
$X'\to X$) one has:
$$
g^{-1}(f(Q))=f'(h^{-1}(Q))\quoi.
$$

If $Q$ is open (\resp closed), the same is true of~$h^{-1}(Q)$, and
hence also of~$f'(h^{-1}(Q))$ if $f'$ is assumed to be an open
(\resp closed) map; the same is therefore true of~$f(Q)$ by the
preceding formula and~\Ref{VIII.4.3}. This proves the first two
assertions in~\Ref{VIII.4.6}. It remains to consider the case in
which~$f'$ is a homeomorphism into, and then to prove that~$f$ is a
homeomorphism into. (The case of a homeomorphism onto will then follow
from~\Ref{VIII.3.1}.) By
\ifthenelse{\boolean{orig}} {\Ref{VIII.3.1}} {\Ref{VIII.3.2}}
% correction: \Ref{VIII.3.2}.
$f$ is injective; it remains to prove that the map $X\to f(X)$ is
open. We already know that~$f$ is quasi-compact by~\Ref{VIII.3.3}.
It suffices to prove that, for every closed subset~$Z$ of~$X$, one
has $Z=f^{-1}(\overline{f(Z)})$. This is equivalent to the analogous
formula for the inverse images under the surjective
map~$h\colon X'\to X$, \ie to
$$
Z'=f'^{-1}(g^{-1}(\overline{f(Z)}))\quoi,
$$
where $Z'=h^{-1}(Z)$. But by~\Ref{VIII.4.1}, applied to the
subset~$f(Z)$ of~$Y$, one has
$g^{-1}(\overline{f(Z)})=\overline{g^{-1}(f(Z))}$, and the formula to
be proved is equivalent to
$$
Z'=f'^{-1}(\overline{f'(Z')})\quoi,
$$
which follows from the hypothesis that $f'$ is a homeomorphism into.

N.B. In this last argument, already supposing that~$f$ is
quasi-compact, we did not use the quasi-compactness of~$g$, but only
that
\ifthenelse{\boolean{orig}}
{$g'$}
{$g$}
% correction: `$g'$' should be `$g$'.
is faithfully flat. Thus it is under this hypothesis that one may
descend the property \og homeomorphism into\fg, or \og homeomorphism
onto\fg, or, by the preceding argument, the property \og $f'$ is
quasi-compact and makes~$f'(X')$ a quotient topological space
of~$X'$\fg.

We shall say that a morphism~$f\colon X\to Y$ of preschemes is
\emph{universally open} (\resp \emph{universally closed},
\resp \emph{universally bicontinuous},
\ifthenelse{\boolean{orig}}{etc...}{etc.})
%
\index{universally open (\resp universally closed, \resp
universally bicontinuous, etc.)|hyperpage}%
%
if for every base change~$Y'\to Y$, $f'\colon X'\to Y'$ is an open
morphism (\resp closed, \resp a
\ifthenelse{\boolean{orig}}
{homeom.}
{homeomorphism}
% correction: homeomorphism
onto its image space). We then deduce from~\Ref{VIII.4.6}:

\begin{corollary}
\label{VIII.4.7}
Under
% page: 209
\marginpar{209}
%
the conditions of~\Ref{VIII.4.6}, $f$ is universally open
(\resp universally closed, \resp a universal homeomorphism into,
\resp a universal homeomorphism) if and only if $f'$ is so.
\end{corollary}

\begin{corollary}
\label{VIII.4.8}
Under the conditions of~\Ref{VIII.4.6}, $f$ is separated
(\resp proper) if and only if $f'$ is so.
\end{corollary}

To say that~$f$ is separated means that the diagonal
morphism~$X\to X\times_{Y}X$ is closed, or equivalently universally closed,
and the first assertion
\ifthenelse{\boolean{orig}}
{}
{of}
% correction: insert `de'.
\Ref{VIII.4.8} therefore follows from~\Ref{VIII.4.7}. To say that~$f$
is proper means that $f$ satisfies the conditions a) $f$ is of finite
type, b) $f$ is separated, c) $f$ is universally closed. Condition
a) does descend by~\Ref{VIII.3.3}, b) does as well by what we have
just seen, and finally c) does by~\Ref{VIII.4.7}.

\begin{remarques}
\label{VIII.4.9}
Recall that when $g\colon Y'\to Y$ is a flat morphism of finite type,
with~$Y$ locally noetherian, then $g$ is an open morphism
(\ifthenelse{\boolean{orig}}{VI}{IV}~\Ref{IV.6.6}), which is a more precise result
than~\Ref{VIII.4.3}. Note, however, that if~$f$ is a faithfully flat
and quasi-compact morphism of noetherian preschemes, then $f$ is not
in general an open morphism. For example, let $Y$ be an irreducible
scheme whose generic point~$y$ is not open (for example, an algebraic
curve), and take for~$Y'$ the disjoint-union
scheme~$Y\amalg \Spec(\kres(y))$. Then the image of the open
subset~$\Spec(\kres(y))$ under the structure morphism~$Y'\to Y$ is not an
open subset of~$Y$. The reader will also note that various statements
of the present expos\'e become false if the hypothesis that the
faithfully flat morphism under consideration is also quasi-compact is
dropped. The typical counterexample is obtained by taking for~$Y'$ the
disjoint-union scheme of the spectra of the local rings at the points
of~$Y$. For example, again taking~$Y$ to be an irreducible algebraic
curve and~$Z$ to be the subset of~$Y$ consisting of the generic point,
its inverse image in~$Y'$ is open although~$Z$ is not open.
\end{remarques}

\subsection{}
\label{VIII.4.10}
Various
% page: 210
\marginpar{210}
%
statements of the present expos\'e remain valid if the hypothesis
that~$Y'$ is flat over~$Y$ is replaced by the following: there exists a
module of finite type~$F$ on~$Y'$, with support~$Y'$, that is flat
relative to~$Y$; the hypothesis of faithful flatness is then replaced
by the preceding one together with the hypothesis that $Y'\to Y$ is
surjective. This applies to the first two assertions
in~\Ref{VIII.1.10}, to~\Ref{VIII.3.3}, \Ref{VIII.3.5},
\Ref{VIII.4.1}, and consequently to all the results of the present
section.
\end{filecontents*}
% Original input: drafts/SGA1_VIII_5_English.texfrag
\begin{filecontents*}[overwrite,noheader]{sga1-cumulative-input-072.tex}
\section{Descent of morphisms of preschemes}
\label{VIII.5}

\begin{proposition}
\label{VIII.5.1}
Let $g\colon S'\to S$ be a morphism of preschemes.
\begin{enumerate}
\item[a)] Suppose that~$g$ is surjective, and that the homomorphism
$$
g^*\colon \mathcal{O}_{S}\to g_*(\mathcal{O}_{S'})
$$
is injective. Then~$g$ is an epimorphism in the category of
preschemes, and even in the category of ringed spaces.
\item[b)] Suppose that~$g$ is surjective and makes~$S$ a quotient
topological space of~$S'$. Let~$S''=S'\times_{S}S'$ and
let~$h\colon S''\to S$ be the structural morphism; consider the
canonical diagram of homomorphisms:
$$
\xymatrix@C=.5cm{\mathcal{O}_{S} \ar[r] &
g_*(\mathcal{O}_{S'})\ar@<2pt>[r]\ar@<-2pt>[r] &
h_*(\mathcal{O}_{S''})}\quoi,
$$
and suppose that this diagram is \emph{exact}. Then $g$ is an
effective epimorphism
%
\index{effective epimorphism@effective epimorphism|hyperpage}%
%
in the category of preschemes (and also in the category of ringed
spaces), \ie the diagram
$$
\xymatrix@C=.5cm{S & S'\ar[l] & S''\ar@<2pt>[l]\ar@<-2pt>[l]}
$$
is exact.
\end{enumerate}
\end{proposition}

\subsubsection*{Proof} a) We must show that a morphism of ringed
spaces~$f\colon S\to Z$ is determined once~$fg$ is known. Since~$g$
is surjective, the underlying set map~$f_{0}$ of~$f$ is known; it
remains to determine the homomorphism of sheaves of
rings~$\mathcal{O}_{Z}\to \mathcal{O}_{S}$, or, what amounts to the
same thing, the homomorphism
% page: 211
\marginpar{211}
%
of sheaves of rings
$$
u\colon f_{0}^{-1}(\mathcal{O}_{Z})\to \mathcal{O}_{S}
$$
defined by~$f$. We already know the homomorphism
$$
(fg)_{0}^{-1}(\mathcal{O}_{Z})=
g_{0}^{-1}(f_{0}^{-1}(\mathcal{O}_{Z}))\to \mathcal{O}_{S'}
$$
defined by~$fg$, or, what amounts to the same thing, we have a
homomorphism
$$
f_{0}^{-1}(\mathcal{O}_{Z})\to
g_{0*}(\mathcal{O}_{S'})=g_*(\mathcal{O}_{S'})\quoi.
$$
One immediately sees that the latter is none other than the composite
of~$g^*\colon \mathcal{O}_{S}\to g_*(\mathcal{O}_{S'})$ with~$u$,
and since $g^*$ is injective, $u$ is known when~$g^*u$ is known.
[N.B. we have obviously not used the fact that~$g\colon S'\to S$ is
a morphism of preschemes; the statement would hold for an arbitrary
morphism of ringed spaces; the same remark applies to b), both in the
category of ringed spaces and in the category of locally ringed
spaces. Note also that if~$g$ is a morphism of preschemes, not
necessarily surjective, but such
that~$g^*\colon \mathcal{O}_{S}\to g_*(\mathcal{O}_{S'})$ is
injective, then for two morphisms~$f_{1},f_{2}$ from~$S$ to a
\emph{scheme}~$Z$ such that~$f_{1}g=f_{2}g$, one has $f_{1}=f_{2}$;
indeed, if~$I$ is the Ideal on~$S$ that defines the subprescheme
of~$S$ where~$f_{1},f_{2}$ coincide (the inverse image of the diagonal
subprescheme of~$Z\times Z$ under~$(f_{1},f_{2})$), one sees that~$I$
is contained in~$\Ker(g^*)$].

b) We must show that for every ringed space~$Z$, the following diagram
of maps
$$
\xymatrix@C=.5cm{\Hom(S,Z)\ar[r] & \Hom(S',Z)\ar@<2pt>[r]\ar@<-2pt>[r] &
\Hom(S'',Z)}
$$
is exact, and that the same holds when~$Z$ is a locally ringed space
and one restricts to morphisms of locally ringed spaces. Since we
already know from a) that the first map is injective, it remains to see
that if~$f'\colon S'\to Z$ is a morphism of ringed spaces such
that~$f'p_{1}=f'p_{2}$,
% page: 212
\marginpar{212}
%
then~$f'$ is of the form~$fg$, where~$f\colon S\to Z$ is a morphism
of ringed spaces. Since~$g$ is surjective, it is then evident that
if~$f'$ is a morphism of locally ringed spaces, the same will hold
for~$f$.

It follows from the hypothesis on~$f'$ that the underlying set
map~$f_{0}'$ is constant on the fibers of the map~$g_{0}$; hence,
since the latter is surjective, $f_{0}'$ factors uniquely
as~$f_{0}'=f_{0}g_{0}$, where~$f_{0}\colon S\to Z$ is a map,
necessarily continuous since~$g_{0}$ identifies~$S$ with a quotient
topological space of~$S'$. Now consider the homomorphism
$$
f_{0}^{-1}(\mathcal{O}_{Z})\to g_*(\mathcal{O}_{S'})
$$
deduced from the homomorphism~$(f_{0}g_{0})^{-1}(\mathcal{O}_{Z})\to
\mathcal{O}_{S'}$ corresponding to~$f'$. The
hypothesis~$f'p_{1}=f'p_{2}$ is then interpreted as saying that the
composites of the preceding homomorphism with the two homomorphisms
$\xymatrix@C=.5cm{g_*(\mathcal{O}_{S'})\ar@<2pt>[r]\ar@<-2pt>[r] &
h_*(\mathcal{O}_{S''})}$ are the same; hence, by hypothesis b), it
factors through a morphism
$$
f_{0}^{-1}(\mathcal{O}_{Z})\to \mathcal{O}_{S}\quoi.
$$
The latter defines a morphism of ringed spaces~$f\colon S\to Z$,
which is the desired morphism.

\begin{theoreme}
\label{VIII.5.2}
Let $\cal{F}$ be the fibered category of arrows in the category~$\Sch$
of preschemes \textup{(\ifthenelse{\boolean{orig}}{VI~}{}\Ref{VI.11}.a)}. Then every faithfully flat
and quasi-compact morphism~$g\colon S'\to S$ is a morphism
of~$\cal{F}$-descent (or, as one also says, a descent morphism
in~$\Sch$).
%
\index{descent (morphism of)|hyperpage}%
%
\end{theoreme}

This means the following: let~$S''=S'\times_{S}S'$, and for two
preschemes~$X,Y$ over~$S$, consider their inverse images~$X',Y'$
over~$S'$ and their inverse images~$X'',Y''$ over~$S''$, giving a
diagram of maps
$$
\xymatrix@C=.5cm{\Hom_{S}(X,Y)\ar[r] &
\Hom_{S'}(X',Y')\ar@<2pt>[r]\ar@<-2pt>[r] & \Hom_{S''}(X'',Y'')}\quoi;
$$
with these notations, \Ref{VIII.5.2} means that this diagram is exact.
It should be noted that it
% page: 213
\marginpar{213}
%
is not true in general that~$g$ is a morphism of effective descent,
\ie that for every prescheme~$X'$ over~$S'$, every descent datum
on~$X'$ relative to~$g\colon S'\to S$ is effective. The question of
effectivity, often delicate, will be examined in~\No \Ref{VIII.7}.

We have seen in \cite{VIII.D} (taking into account that fiber products
exist in~$\Sch$) that statement~\Ref{VIII.5.2} is equivalent to the
following:

\begin{corollaire}
\label{VIII.5.3}
A faithfully flat and quasi-compact morphism of preschemes is a
universal effective epimorphism.
\end{corollaire}

Since a faithfully flat and quasi-compact morphism remains so under
every base change, we are reduced to proving that it is an effective
epimorphism. We then apply the criterion~\Ref{VIII.5.1} b), which gives
the desired result, taking~\Ref{VIII.4.3} and~\Ref{VIII.1.7} into
account.

\begin{corollaire}
\label{VIII.5.4}
Let $g\colon S'\to S$ be a faithfully flat and quasi-compact morphism,
let $f\colon X\to Y$ be an~$S$-morphism, and let $f'\colon X'\to Y'$
be the~$S'$-morphism deduced from it by the base change~$S'\to S$.
Then~$f$ is an isomorphism if and only if~$f'$ is.
\end{corollaire}

Indeed, if $f'$ is an isomorphism, it is also an isomorphism for the
natural descent structures on~$X',Y'$, and since the functor~$X\mto X'$
from~$\Sch_{/S}$ to the category of preschemes over~$S'$ endowed with
a descent datum relative to~$g$ is fully faithful by~\Ref{VIII.5.2},
the desired conclusion follows.

\begin{corollaire}
\label{VIII.5.5}
Under the conditions of~\Ref{VIII.5.4}, $f$ is a closed immersion
(\resp an open immersion, \resp a quasi-compact immersion) if and only
if~$f'$ is.
\end{corollaire}

As usual, we may suppose~$Y=S$, and only the \og if\fg{} direction has
to be proved. Note that the fact that~$X'/Y'$ is endowed with a descent
datum relative to~$g\colon Y'\to Y$, and that the structural
morphism~$f'\colon X'\to Y'$ is an immersion, hence a monomorphism,
implies that the two subobjects of~$Y''$ obtained as
% page: 214
\marginpar{214}
%
inverse images of~$X'/Y'$ under either projection from~$S''$ to~$S'$
are the same. If~$f'$ is a closed immersion, it follows by
virtue of~\Ref{VIII.1.9} that there exists a closed subprescheme $X_{1}$
of~$Y$ whose inverse image under~$g\colon Y'\to Y$ is~$X'$. Thus,
by the uniqueness of the solution to a descent problem relative to a
morphism of $\cal{F}$-descent,
\ifthenelse{\boolean{orig}}
{}
{it}
% correction: insert `it'.
follows that~$X_{1}$ is $Y$-isomorphic to~$X$, so $f\colon X\to Y$
is a closed immersion. We proceed in the same way for an open
immersion, using~\Ref{VIII.4.4}. Finally, if $f'$ is a quasi-compact
immersion, $f$ is quasi-compact by virtue of~\Ref{VIII.3.3}, so we may
apply the criterion~\Ref{VIII.4.5} to the subset~$f(X)$ of~$Y$, which
proves that~$f(X)$ is locally closed since its inverse image~$f'(X')$
in~$Y'$ is so. Replacing~$Y$ by an open subset in which~$f(X)$ is
closed, we are reduced to the case where~$f'$ is a closed immersion,
so $f$ is one by what precedes.

\begin{corollaire}
\label{VIII.5.6}
Under the conditions of~\Ref{VIII.5.4}, $f$ is affine if and only
if~$f'$ is.
\end{corollaire}

We proceed as in~\Ref{VIII.5.5}, using~\Ref{VIII.2.1}
(One may also use Serre's cohomological criterion [EGA II 5.2], which
proves~\Ref{VIII.5.6} without using descent techniques).

\begin{corollaire}
\label{VIII.5.7}
Under the conditions of~\Ref{VIII.5.4}, $f$ is integral
(\resp finite, \resp finite and locally free) if and only if~$f'$ is.
\end{corollaire}

Only the \og if\fg{} direction has to be proved, and as usual we may
suppose $Y=S$, with $Y$ affine and $Y'$ affine. Since the hypothesis
implies that~$f'$ is affine, the same is true of~$f$
by~\Ref{VIII.5.6}, so~$X$ and consequently~$X'$ are affine.
Let~$A,A',B,B'=B\otimes_{A}A'$ be the rings of~$Y,Y',X,X'$. We have
$B=\varinjlim_i B_{i}$, where~$B_{i}$ ranges over the
sub-$A$-algebras of~$B$ that are of finite type over~$A$, whence
$B'=\varinjlim_i B_{i}'$, where the~$B_{i}'$ are finite-type
subalgebras of the $A'$-algebra~$B'$. If $B'$ is integral
\ifthenelse{\boolean{orig}}
{over~$A$,}
{over~$A'$,}
% correction: `$A$' should be `$A'$'.
the~$B_{i}'$ are modules of finite type over~$A'$; since $A'$ is
faithfully flat over~$A$, the~$B_{i}$ are modules of finite type
over~$A$, \ie $B$ is integral over~$A$. One sees in the same way that
if~$B'$ is finite over~$A'$, $B$ is finite over~$A$. The same
conclusion holds for \og locally free of finite type\fg, \cf
\Ref{VIII.1.11}.

\begin{corollaire}
\label{VIII.5.8}
Under
% page: 215
\marginpar{215}
%
the conditions of~\Ref{VIII.5.4}, suppose~$f$ quasi-compact and
let~$\mathcal{L}$ be an invertible Module on~$X$, and~$\mathcal{L}'$
its inverse image on~$X'$. Then~$\mathcal{L}$ is ample
(\resp very ample) relative to~$f$ if and only if
\ifthenelse{\boolean{orig}}
{~$f'$}
{~$\cal{L}'$}
% correction: `f'' should be `\cal{L}''.
is ample (\resp very ample) relative to~$f'$.
\end{corollaire}

Only the \og if\fg{} direction has to be proved. The hypothesis
on~$\mathcal{L}\ifthenelse{\boolean{orig}}{}{'}$ implies in any case that~$f'$ is separated, hence
$f$ is separated by~\Ref{VIII.4.8}; and since~$f$ is quasi-compact
and~$g\colon Y'\to Y$ is flat, the computation of direct images by
affine coverings shows that there are isomorphisms
$$
g^*(f_*(\mathcal{L}^{\otimes n}))\isomto
f_*'(\mathcal{L}'^{\otimes n})
$$
for every integer~$n$, and consequently there is an isomorphism
$$
g^*(\mathcal{S})\isomto \mathcal{S}'\quoi,
$$
where~$\mathcal{S}$ (\resp $\mathcal{S}'$) denotes the quasi-coherent
graded Algebra on~$Y$ (\resp on~$Y'$) that is the direct sum of
the~$f_*(\mathcal{L}^{\otimes n})$
(\resp the~$f_*'(\mathcal{L}'^{\otimes n})$) for~$n\geq 0$. Note that
for every~$n\geq 0$, the cokernel of the canonical
homomorphism~$f'^*(\mathcal{S}_{n}')\to \mathcal{L}'^{\otimes n}$ is
the inverse image under~$X'\to X$ of the cokernel
of~$f^*(\mathcal{S}_{n})\to \mathcal{L}^{\otimes n}$, and hence its
support~$Z_{n}'$ is the inverse image of the support~$Z_{n}$. If
\ifthenelse{\boolean{orig}}
{$f$}
{$\cal{L}'$}
% correction: $f$ should be $\cal{L}'$
is ample, the intersection of the~$Z_{n}'$ is empty; hence, since
$X'\to X$ is surjective, the intersection of the~$Z_{n}$ is empty,
\ie there is a canonical morphism
$$
j\colon X\to \Proj(\mathcal{S})
$$
(EGA~II 3). Moreover, the analogous morphism
$$
j'\colon X'\to \Proj(\mathcal{S}')
$$
is none other than the one deduced from the preceding morphism by the
base change~$Y'\to Y$ (\loccit). This being so, to say that
$\mathcal{L}'$ is ample relative to~$f'$ means that~$j'$ is an
immersion, necessarily quasi-compact since~$f'$ is quasi-compact.
Thus, by virtue of~\Ref{VIII.5.5}, $j$ is an immersion, \ie
$\mathcal{L}$ is ample relative to~$f$. --- We proceed in an entirely
analogous manner in the case of \og very
% page: 216
\marginpar{216}
%
ample\fg, restricting the above to~$n=1$, and replacing the
consideration of~$\Proj(\mathcal{S})$ by that of the projective
bundle~$\mathcal{P}(\mathcal{S}_{1})$ associated
with~$\mathcal{S}_{1}$.

Recall (EGA~II 5.1.1) that a quasi-compact morphism~$f$ is called
\emph{quasi-affine} if for every affine open subset~$U$ of~$Y$,
$f^{-1}(U)$ is a prescheme isomorphic to an open subscheme of an
affine scheme. It is shown (\loccit) that this is equivalent to saying
that~$\mathcal{O}_{X}$ is ample (or also: very ample) relative to~$f$.
Thus~\Ref{VIII.5.8} implies:

\begin{corollaire}
\label{VIII.5.9}
Under the conditions of~\Ref{VIII.5.4}, and assuming~$f$
quasi-compact, $f$ is quasi-affine if and only if~$f'$ is.
\end{corollaire}

\begin{remarques}
\label{VIII.5.10}
Hironaka's example of a nonprojective variety shows that one can have
a proper morphism~$f\colon X\to Y$ of nonsingular algebraic varieties
(with~$Y$ projective), such that~$Y$ is the union of two open
subsets~$Y_{i}$ for which~$X_{i}=X\times_{Y}Y_{i}$ is projective
over~$Y_{i}$, while~$f$ is not projective. Thus, putting
$Y'=Y_{1}\amalg Y_{2}$, $Y'$ is faithfully flat and quasi-compact
(and even quasi-finite) over~$Y$, and $f'\colon X'\to Y'$ is
projective, but~$f$ is not projective. One must therefore take care:
in order to apply~\Ref{VIII.5.8}, and to deduce from the fact that~$f'$
is projective the same conclusion for~$f$, one must already have
on~$X'$ an invertible Module~$\mathcal{L}'$ ample for~$f'$,
\emph{endowed with a descent datum relative to}~$X'\to X$ (which
allows~$\mathcal{L}'$ to be regarded as the inverse image of an
invertible Module~$\mathcal{L}$ on~$X$, which will then be ample
for~$f$ by virtue of~\Ref{VIII.5.8}). When~$g\colon S'\to S$ is
finite and locally free, however, see~\Ref{VIII.7.7}.
\end{remarques}
\end{filecontents*}
% Original input: drafts/SGA1_VIII_6_English.texfrag
\begin{filecontents*}[overwrite,noheader]{sga1-cumulative-input-073.tex}
\section[Application to finite and quasi-finite morphisms]{Application to finite and quasi-finite morphisms\protect\footnotemark}
\label{VIII.6}
\footnotetext{\hypertarget{Hfootnote.3}{}\Cf EGA IV 18.12 for generalizations to preschemes that are not necessarily locally Noetherian}
% the footnote was in the title of the section.

We shall prove the following two theorems:

\begin{theoreme}
\label{VIII.6.1}
Let~$f\colon X\to Y$ be a morphism that is \emph{proper with finite fibers},
with~$Y$ locally Noetherian. Then~$f$ is finite.
\end{theoreme}

\begin{theoreme}
\label{VIII.6.2}
Let~$f\colon X\to Y$ be a \emph{quasi-finite} and
\emph{separated} morphism, with~$Y$ locally Noetherian. Then~$f$
is quasi-affine, and a fortiori quasi-projective.
\end{theoreme}

\begin{remarques}
\label{VIII.6.3}
The
% page: 217
\marginpar{217}
%
theorem~\Ref{VIII.6.1} is well known, and is due to
Chevalley in the case of algebraic varieties; a simple proof may also
be found in [EGA III 4], using the \og theorem on holomorphic functions\fg. The
proof given here does not use the latter
theorem, but instead uses descent theory; we
give it as a bonus to the reader, since one obtains it \og for free\fg at the
same time as the proof of~\Ref{VIII.6.2}. Recall also ([EGA III
4] or \cite{VIII.1}) that the global form of Zariski's \og Main Theorem\fg,
deduced from the \og theorem on holomorphic functions\fg, asserts
that if~$f\colon X\to Y$ is quasi-finite and \emph{quasi-projective}, with~$Y$
Noetherian, then~$X$ is $Y$-isomorphic to an
open subprescheme of a~$Y$-prescheme
\emph{finite}~$Z$. The conjunction of the \og Main Theorem\fg and
of~\Ref{VIII.6.2} may therefore be stated as follows:
\end{remarques}

\begin{corollaire}
\label{VIII.6.4}
Let~$f\colon X\to Y$ be a quasi-finite and separated morphism,
with~$Y$ Noetherian. Then~$X$ is $Y$-isomorphic to an
open subprescheme of a finite~$Y$-prescheme~$Z$.
\end{corollaire}

Another interesting consequence of~\Ref{VIII.6.2} for
descent theory will be given along with~\Ref{VIII.7.9}.

\subsubsection*{Proof of~\Ref{VIII.6.1} and~\Ref{VIII.6.2}}
We shall use the following fact, whose proof is
easy\footnote{\Cf EGA IV 8}:

\begin{lemme}
\label{VIII.6.5}
Let~$X$ be a prescheme of finite type over a locally
Noetherian prescheme~$Y$, and let~$y\in Y$. In order that there exist an open
neighborhood~$U$ of~$y$ such that~$X|U$ is finite (\resp quasi-affine,
\resp\ ...) over~$U$, it is necessary and sufficient
that~$X\times_{Y}\Spec(\mathcal{O}_{y})$ be finite (\resp quasi-affine,
\resp\ ...) over~$\Spec(\mathcal{O}_{y})$.
\end{lemme}

Since, on the other hand, the property for~$f\colon X\to Y$
of being finite, \resp quasi-affine, is local on~$Y$, we are
reduced, in order to prove~\Ref{VIII.6.1} and~\Ref{VIII.6.2}, to the case
where~$Y$ is the spectrum of a local ring, and hence in particular is of
finite dimension. (N.B. by the \emph{dimension of a
prescheme}~$Y$ we mean the supremum of the Krull dimensions of its local
rings). We proceed by induction on~$$
n=\dim(Y)\quoi,
$$
the assertion
% page: 218
\marginpar{218}
%
being trivial for~$n<0$. We may therefore suppose~$n\geq 0$,
and the assertion proved for dimensions~$n'<n$. We may
again suppose that~$Y$ is the spectrum of a Noetherian local
ring~$A$, of dimension~$n$. Note that the hypotheses
made in~\Ref{VIII.6.1} and~\Ref{VIII.6.2} are stable under
base change (this was already used in the initial
reduction), so they will remain true after the base
change~$\Spec(\widehat{A})\to \Spec(A)$. Since the latter is
faithfully flat and quasi-compact, the
statements~\Ref{VIII.5.7} and~\Ref{VIII.5.9} reduce us to the
case where~$A$ is moreover complete. Using then the fact that every
Noetherian local ring~$B$ over~$A$ that is quasi-finite over~$A$ is finite
over~$A$, and the fact that~$X$ is separated over~$Y$ and the fiber
at~$y$ consists of isolated points, we obtain a
decomposition
$$
X=X'\amalg X''
$$
where~$X'$ is \emph{finite} over~$Y$, and where the fiber of~$X''$
at~$y$ is empty. If~$X$ is proper over~$Y$, then so is~$X''$,
hence its image in~$Y$ is closed, and since it does not
contain~$y$, it is empty; thus~$X''=\emptyset$ and hence~$X=X'$,
which shows that~$X$ is finite over~$Y$ and proves~\Ref{VIII.6.1}
(N.B. the induction hypothesis is unnecessary here). If~$X$ is
quasi-finite over~$Y$, so is~$X''$, but~$X''$ in fact lies over
\ifthenelse{\boolean{orig}}
{the open subset~$Y-(y)$}
{the open subset~$Y-\{ y\}$}
% correction: replace `$Y-(y)$' by `$Y-\{ y\}$'.
of~$Y$, \emph{which has dimension} $<n$. By the induction
hypothesis, $X''$ is quasi-affine
\ifthenelse{\boolean{orig}}
{over~$Y-(y)$,}
{over~$Y-\{ y\}$,}
% correction: replace `$Y-(y)$' by `$Y-\{ y\}$'.
hence also over~$Y$; the same is plainly true of~$X'$, and hence
also of their sum~$X$, which proves~\Ref{VIII.6.2}.

\begin{remarque}
\label{VIII.6.6}
Theorems~\Ref{VIII.6.1} and~\Ref{VIII.6.2} remain
valid if~$Y$ is no longer assumed locally Noetherian, provided
one specifies that~$f$ is assumed to be of finite presentation
\ifthenelse{\boolean{orig}}
{(\cf \Ref{VIII.3.6}).}
{(\cf \Ref{VIII.3.6}).}
% correction: \cf
Indeed, we may still suppose~$Y$ affine, and then one verifies
without difficulty that the situation~$f\colon X\to Y$ is
\ifthenelse{\boolean{orig}}
{deduced,}
{deduced,}
% correction: `d\'eduit' should be `d\'eduite'.
by a base change~$Y\to Y_{0}$, from a situation~$f_{0}\colon
X_{0}\to Y_{0}$ satisfying the same hypotheses as~$f$,
with~$Y_{0}$ \emph{Noetherian}. Therefore, by~\Ref{VIII.6.1}
\resp \Ref{VIII.6.2}, $f_{0}$ is finite
\resp quasi-affine, and hence so is~$f$. This kind of
argument is often useful for dispensing with
Noetherian hypotheses,
% page: 219
\marginpar{219}
%
(which always end up being troublesome in
applications).
\end{remarque}
\end{filecontents*}
% Original input: drafts/SGA1_VIII_7_A_English.texfrag
\begin{filecontents*}[overwrite,noheader]{sga1-cumulative-input-074.tex}
\section{Criteria for the effectivity of a descent datum}
\label{VIII.7}

As usual, consider a morphism of preschemes
\[
  g\colon S'\longrightarrow S
\]
and an \(S'\)-prescheme~\(X'\). In accordance with the general
facts in
\ifthenelse{\boolean{orig}}
{VII,~9,}
{VII,~9\footnote{This section does not exist; see Expos\'e 212 of
the S\'eminaire Bourbaki (1962).},}
% correction: the cited section does not exist.
giving a descent datum on~\(X'\) relative to~\(g\) is equivalent to
giving an equivalence pair \cite{VIII.3}
\[
  \xymatrix@C=.5cm{
  q_1,q_2\colon X''\ar@<2pt>[r]\ar@<-2pt>[r]&X'
  }
\]
such that the structural morphism \(X'\to S'\) is compatible with
this pair and the equivalence pair
\[
  \xymatrix@C=.5cm{
  p_1,p_2\colon S''=S'\times_SS'
  \ar@<2pt>[r]\ar@<-2pt>[r]&S'
  }
\]
defined by~\(g\), and such that the two squares (or either one of
them, which amounts to the same thing by symmetry) extracted from the
corresponding diagram
\[
  \xymatrix{
  X'\ar[d]&X''\ar@<2pt>[l]\ar@<-2pt>[l]\ar[d]\\
  S'&S''\ar@<2pt>[l]\ar@<-2pt>[l]
  }
\]
(using either \(p_1,q_1\) or \(p_2,q_2\)) are \emph{cartesian}. A
solution to the descent problem posed by this descent datum, \ie an
object \(X\) over~\(S\) endowed with an isomorphism
\(X\times_SS'\isomfrom X'\) compatible with the descent data, is
equivalent to a \emph{cartesian} square
\[
  \xymatrix{
  X\ar[d]&X'\ar[l]_{h}\ar[d]\\
  S&S'\ar[l]_{g}
  }
\]
satisfying \(hq_1=hq_2\).

Since the class
% page: 220
\marginpar{220}
%
of faithfully flat and quasi-compact morphisms is stable under base
change, and since a faithfully flat and quasi-compact morphism is an
effective epimorphism by~\Ref{VIII.5.3}, the general theory
\cite{VIII.D} gives:

\begin{proposition}
\label{VIII.7.1}
Suppose that \(g\colon S'\to S\) is faithfully flat and
quasi-compact. A descent datum on~\(X'\) relative to~\(g\) is
effective if and only if the equivalence relation
\(R=(q_1,q_2)\) that it defines is effective (\ie the quotient
\(X'/R\) exists and \(X''\) becomes the fiber square of~\(X'\) over
\(X'/R\)), and the canonical morphism \(X'\to X'/R\) is faithfully
flat and quasi-compact.
\end{proposition}

Thus the question of the effectivity of a descent datum is a special
case of the question of the effectivity of an equivalence graph, and
various effectivity criteria given in this section may be obtained in
this way. Nevertheless, in the context of descent one has
Theorem~\Ref{VIII.2.1}, which implies that \emph{if \(X'\) is affine
over~\(S'\), every descent datum on~\(X'\) relative to~\(g\) is
effective}. This statement has no analogue for passage to the
quotient by a general flat equivalence graph. All the effectivity
criteria given here may also be regarded as consequences of that
theorem.

Let \(U'\) be a subprescheme of~\(X'\) (or, more generally, a
subobject of~\(X'\) in the category~\(\Sch\)). We say that \(U'\) is
\emph{stable under the descent datum} on~\(X'\) if \(U'\) can be
endowed with a descent datum relative to~\(g\) for which the immersion
\(U'\to X'\) is compatible with the descent data. This also means
that the inverse images of~\(U'\) in~\(X''\) under \(q_1\) and
\(q_2\) are the same (or, as one also says, that \(U'\) is
\emph{stable under the equivalence relation}~\(R\)). The descent
datum in question on~\(U'\) is then plainly unique and is called the
\emph{induced descent datum}
%
\index{induced descent datum|hyperpage}%
%
from that on~\(X'\). With this terminology:

\begin{proposition}
\label{VIII.7.2}
Let \((X_i')\) be a covering of~\(X'\) by open subsets \(X_i'\)
stable
% page: 221
\marginpar{221}
%
% correction: the following source lines were not underlined.
under the descent datum. The descent datum on~\(X'\) is effective if
and only if the induced descent data on all the \(X_i'\) are
effective.
\end{proposition}

This is, for example, an immediate consequence of~\Ref{VIII.7.1};
the details of the proof are left to the reader.

\begin{corollary}
\label{VIII.7.3}
Let \((S_i)\) be an open covering of~\(S\), and for every~\(i\) let
\(S_i'\) and \(X_i'\) be deduced from \(S'\)
\ifthenelse{\boolean{orig}}
{from}
{and}
% correction: replace from by and.
\(X'\) by the base change \(S_i\to S\). The descent datum on~\(X'\)
is effective if and only if, for every~\(i\), the descent datum on
\(X_i'\) relative to \(g_i\colon S_i'\to S_i\) is effective.
\end{corollary}

In practice this criterion always reduces us to the case in which
\(S\) is affine. When \(S'\) is affine as well, as is most often the
case in applications, we have:

\begin{corollary}
\label{VIII.7.4}
Suppose that \(S\) and \(S'\) are affine. The descent datum on~\(X'\)
is effective if and only if \(X'\) is the union of affine open subsets
\(X_i'\) stable under the descent datum.
\end{corollary}

Sufficiency follows from~\Ref{VIII.7.2} and the fact that, if \(X_i'\)
is affine, then it is affine over~\(S'\), so that
\Ref{VIII.2.1} applies. For necessity, observe that if \(X'\) comes
from~\(X\), and if \(X\) is covered by affine open subsets \(X_i\),
then the \(X_i'=X_i\times_SS'\) are affine open subsets stable under
the descent datum and covering~\(X'\).

\begin{corollary}
\label{VIII.7.5}
Let \(g\colon S'\to S\) be a faithfully flat, quasi-compact, and
\emph{radicial} morphism. Then \(g\) is an \emph{effective descent}
morphism, \ie for every \(X'\) over~\(S'\), every descent datum on
\(X'\) relative to \(g\colon S'\to S\) is effective.
\end{corollary}

Indeed, by~\Ref{VIII.7.3} we may suppose that \(S\) is affine. Since
\(S'\) is radicial over~\(S\), and hence separated over~\(S\), it is
separated. Moreover, for every \(x'\in X'\), the fiber
\[
  R(x')=q_2(q_1^{-1}(x'))
\]
of the set-theoretic equivalence relation defined
% page: 222
\marginpar{222}
%
by~\(R\) consists of a single point. Indeed, since \(g\) is
radicial, the same holds for \(p_1,p_2\), which are deduced from~\(g\)
by the base change \(S'\to S\), and hence also for \(q_1,q_2\), which
are deduced from the former by the base change \(X'\to S''\).
Therefore \emph{every open subset of~\(X'\) is stable} under the
descent datum. Cover \(X'\) by affine open subsets \(X_i'\). They are
affine over~\(S'\), since \(S'\) is separated, so the induced descent
datum is effective by~\Ref{VIII.2.1}. We conclude by~\Ref{VIII.7.2}.

Observe that~\Ref{VIII.7.5} gives the only known instance of an
effective descent morphism in the category of preschemes, and it is
probably the only instance, even if one restricts to noetherian
schemes or to schemes of finite type over a field.

When \(S\) is assumed locally noetherian and \(S'\) of finite type
over~\(S\), statement~\Ref{VIII.7.5} is also a special case of the
following result (which generalizes Weil's Galois descent and
Cartier's inseparable descent):

\begin{corollary}
\label{VIII.7.6}
Let \(g\colon S'\to S\) be a finite locally free morphism (\ie
defined by an algebra on~\(S\) which, as a module, is finite locally
free) and suppose that \(g\) is surjective (so \(g\) is faithfully
flat and quasi-compact, and hence a descent morphism). Let \(X'\) be
an \(S'\)-prescheme endowed with a descent datum. This datum is
effective if and only if, for every \(x'\in X'\), the fiber
\(R(x')=q_2(q_1^{-1}(x'))\) is contained in an affine open subset
(a condition automatically satisfied if \(X'\) is quasi-projective
over~\(S'\)).
\end{corollary}

The parenthetical observation follows from the fact that, if \(s\)
is the point of~\(S\) below~\(x'\), then \(R(x')\) is finite and
contained in the
\ifthenelse{\boolean{orig}}
{fiber \(X_s\),}
{fiber \(X'_s\),}
% correction: the fiber is that of X prime.
while, since \(X'\) is quasi-projective over~\(S'\) and \(S'\) is
finite over~\(S\), \(X'\) is quasi-projective over~\(S\); this implies
that a fiber
\ifthenelse{\boolean{orig}}
{of \(X/S\)}
{of \(X'\) over~\(S\)}
% correction: the fiber is that of X prime over S.
is contained in an affine open subset.

Since every finite subset of an affine scheme has a fundamental
system of affine neighborhoods, the hypothesis is preserved when we
restrict over an affine open subset of~\(S\). By~\Ref{VIII.7.3}, this
reduces us to the case in which \(S\) is affine. By~\Ref{VIII.7.4},
it is enough to show
\ifthenelse{\boolean{orig}}
{that \(X'\)}
{that \(x'\)}
% correction: X prime should be x prime.
is contained
% page: 223
\marginpar{223}
%
in an affine open subset \emph{stable} under the descent datum.
Indeed, let \(U\) be an affine open subset containing \(R(x')\).
Then the saturation
\[
\ifthenelse{\boolean{orig}}
{R(X-U)=q_2(q_1^{-1}(X-U))}
{R(X'-U)=q_2(q_1^{-1}(X'-U))}
% correction: X should be X prime in both occurrences.
\]
does not meet \(R(x')\). Moreover, since \(q_2\) is finite (because
\(g\), and hence \(p_2\), is finite), and therefore closed, the
right-hand side is a closed subset of~\(X'\). Let \(U'\) be its
complement in~\(X'\). Thus \(U'\) is a \emph{saturated} open subset,
and
\[
  R(x')\subset U'\subset U,
\]
where \(U\) is affine, although \(U'\) need not be affine a priori.
Since the finite subset \(R(x')\) of the affine scheme~\(U\) has a
fundamental system of affine neighborhoods of the form~\(U_f\), one
sees, by replacing \(f\) by its restriction to~\(U'\), that there is
a section \(f\) of~\(\mathcal O_U\) such that
\[
  R(x')\subset U_f',
  \qquad U_f'\ \text{is affine}.
\]
Set \(U''=q_1^{-1}(U')=q_2^{-1}(U')\), continue to write \(q_1,q_2\)
for the induced morphisms \(U''\to U'\), and consider
\[
  f'=\Norm_{q_2}(q_1^*(f)),
\]
where \(\Norm_{q_2}\) denotes the \emph{norm} relative to the finite
locally free morphism \(q_2\colon U''\to U'\). Compatibility of the
formation of the norm with base change readily implies that \(f'\)
is an \emph{invariant} section:
\[
  q_1^*(f')=q_2^*(f').
\]
It follows that \(U_{f'}'\) is a saturated open subset of~\(U'\).
More precisely, writing \(Z(f')\) for the zero set of a section
\(f'\), the properties of norms give
\[
  Z(f')=q_2(Z(q_1^*(f)))
       =q_2(q_1^*(Z(f)))
       =R(U'-U_f').
\]
Consequently,
% page: 224
\marginpar{224}
%
\(U_{f'}'=U'-Z(f')\) is saturated, contains \(R(x')\), and is
contained in~\(U_f'\). Since the latter is affine, \(U_{f'}'\) is
affine as well (it is equal to \((U_f')_{f''}\), with
\(f''=f'|_{U_f'}\)). It is therefore a saturated affine open subset
containing \(R(x')\), and hence \(x'\). This completes the proof.

\end{filecontents*}
% Original input: drafts/SGA1_VIII_7_B_English.texfrag
\begin{filecontents*}[overwrite,noheader]{sga1-cumulative-input-075.tex}
The same argument applies whenever one has an equivalence relation
(or even merely a pre-equivalence relation, \cf \cite{VIII.3}) on a
prescheme~\(X'\) which is finite and locally free. Moreover,
\Ref{VIII.7.6} is itself a special case of the analogous result for
finite locally free pre-equivalence relations; see \loccit The same
observation applies to~\Ref{VIII.7.7} below.

Alternatively, once a saturated quasi-affine open subset \(U'\)
containing~\(x'\) has been obtained, one may invoke
\Ref{VIII.7.9} and~\Ref{VIII.7.2}, thereby avoiding the use of norms.

Observe also that under the hypotheses of~\Ref{VIII.7.6}, if the
descent datum on~\(X'\) is effective and \(X'\) comes from~\(X\)
over~\(S\), then the morphism \(X'\to X\) is finite, locally free,
and surjective, since it is deduced from~\(g\) by the base change
\(X\to S\). It follows (EGA~II, 6.6.4) that if \(X'\) is
quasi-projective over~\(S'\), and hence over~\(S\), then \(X\) is
quasi-projective over~\(S\). A relatively ample invertible sheaf on
\(X\) is obtained by taking the \emph{norm} of an invertible sheaf
on~\(X'\) that is relatively ample over~\(S\), or equivalently
over~\(S'\). Thus we obtain:

\begin{corollary}
\label{VIII.7.7}
A finite locally free and surjective morphism
\(g\colon S'\to S\) is an effective descent morphism for the fibered
category of preschemes quasi-projective over other preschemes. In
other words, for every \(X'\) quasi-projective over~\(S'\), every
descent datum on~\(X'\) relative to~\(g\) is effective, and the
descended \(S\)-prescheme~\(X\) is quasi-projective over~\(S\).
\end{corollary}

\begin{proposition}
\label{VIII.7.8}
Let \(g\colon S'\to S\) be faithfully flat and quasi-compact. Then
\(g\) is an effective descent morphism for the fibered category of
preschemes \(Z\) quasi-compact over a prescheme~\(T\), equipped with
an invertible sheaf ample relative
% page: 225
\marginpar{225}
%
to~\(T\). In particular, let \(X'\) be an \(S'\)-prescheme endowed
with a descent datum relative to \(g\colon S'\to S\), and let
\(\mathcal L'\) be an invertible sheaf on~\(X'\), ample relative
to~\(S'\), likewise endowed with a descent datum relative to the
given datum on~\(X'\) (\ie endowed with an isomorphism between
\(q_1^*(\mathcal L')\) and \(q_2^*(\mathcal L')\) satisfying the
usual transitivity condition). Then the descent datum on~\(X'\) is
effective, and the invertible sheaf \(\mathcal L\) on the descended
prescheme~\(X\), obtained from~\(\mathcal L'\) by descent, is ample
relative to~\(S\).
\end{proposition}

The proof is entirely analogous to that of~\Ref{VIII.5.8}. On the
quasi-coherent graded algebra \(\mathcal S'\) over~\(S'\) defined by
\(\mathcal L'\), there is a descent datum. By~\Ref{VIII.1.1}, this
allows one to construct a quasi-coherent graded algebra
\(\mathcal S\) over~\(S\), and hence an
\(S\)-prescheme \(P=\Proj(\mathcal S)\), such that
\(P'=\Proj(\mathcal S')\), together with its descent datum, is
identified with \(P\times_SS'\). By hypothesis, \(X'\) is identified
with an open subset of~\(P'\), necessarily stable under the descent
datum on~\(P'\). The descent datum on~\(X'\) is therefore effective
as well, and the descended prescheme is obtained as an open subset
of~\(P\). The details are left to the reader. In particular, taking
\(\mathcal L'=\mathcal O_{X'}\), we find:

\begin{corollary}
\label{VIII.7.9}
Let \(g\colon S'\to S\) be a faithfully flat and quasi-compact
morphism, and let \(X'\) be a prescheme \emph{quasi-affine}
over~\(S'\). Then every descent datum on~\(X'\) relative to~\(g\) is
effective, and the descended prescheme~\(X\) is quasi-affine
over~\(S\).
\end{corollary}

By~\Ref{VIII.6.2}, this result applies in particular when \(S'\) is
locally noetherian and \(X'\) is quasi-finite and separated
over~\(S'\), and more generally when \(S'\) is arbitrary and \(X'\)
is of finite presentation, quasi-finite, and separated over~\(S'\)
(\cf \Ref{VIII.6.6}).

\begin{remarques}
\label{VIII.7.10}
The results given in this section exhaust the currently known
effectivity criteria, and probably even all the useful criteria that
exist\footnote{This opinion proved partly mistaken; see, for example,
J.P\ptbl \textsc{Murre}, S\'em. Bourbaki 294 (Appendix), May 1965,
and special results, notably those of \textsc{N\'eron} and
\textsc{Raynaud}. For the descent of group schemes, \cf
M\ptbl Raynaud, thesis (1968).}. The following counterexamples
support this assertion:

\begin{enumerate}
\item[(i)]
% page: 226
\marginpar{226}
%
If \(S\) is the spectrum of a field and \(S'\) is the spectrum of a
quadratic Galois extension, there exists an \(X'\) over~\(S'\), proper
and smooth over~\(S'\), of dimension~3, endowed with a descent datum
which is not effective (Serre).

\item[(ii)] There exist an \(S\) which is the spectrum of a regular
local ring of dimension~3 (if desired, the local ring of an algebraic
scheme over a field of prescribed characteristic), a principal
covering \(T\) of~\(S\) with group \(\ZZ/2\ZZ\), and the following
property. If \(t\) denotes one of the points of~\(T\) above the closed
point \(s\) of~\(S\), and \(S'=T-\{t\}\), there exists a regular
\(X'\), \emph{projective} over~\(S'\), endowed with a descent datum
relative to \(g\colon S'\to S\), which is not effective.
\end{enumerate}

These constructions use Hironaka's example of nonprojective
varieties. For~(i), it is enough to use the fact that there exists
over~\(k\) a proper smooth scheme \(X_0\) of dimension~3 on which
\(G=\ZZ/2\ZZ\) acts without inertia, and in which there are two
rational points \(a,b\),
\ifthenelse{\boolean{orig}}
{over \(P\),}
{over \(S\),}
% correction: P should be S.
conjugate under~\(G\), which are not contained in an affine open
subset. Set \(X'=X_0\times_k k'\), and let \(G\) act on~\(X'\) by
its actions on both factors. This gives a descent datum on~\(X'\)
relative to \(g\colon\Spec(k')\to\Spec(k)\). Above \(a\),
respectively \(b\), there is exactly one point \(a'\), respectively
\(b'\), with quadratic residue-field extension, and \(a'\) and
\(b'\) are conjugate under~\(G\), since \(X'\to X_0\) is compatible
with the actions of~\(G\). The points \(a'\) and \(b'\) cannot be
contained in an affine open subset \(U'\). Indeed, otherwise
\[
  U=X_0-\operatorname{Im}(X'-U')
\]
would be an open subset of~\(X_0\) containing \(a,b\), whose inverse
image in~\(X'\) would be contained in~\(U'\), and hence would be
quasi-affine. It would follow that \(U\) is quasi-affine, and thus
that \(a,b\) have an affine neighborhood in~\(U\).

For~(ii), one uses the fact that in Hironaka's example \(X_0\) is
obtained as a prescheme proper over a projective \(k\)-scheme~\(Y\),
smooth over~\(k\). The morphism \(f\colon X_0\to Y\) is in fact
birational, but this is irrelevant. The group \(G\) also acts on~\(Y\)
compatibly with its action on~\(X_0\). Finally, on setting
\(S'=Y-\ifthenelse{\boolean{orig}}{f(b)}{\{f(b)\}}\) and \(X'=X_0|S'\), the prescheme \(X'\) is projective
over~\(S'\). The prescheme \(X_0\) is then endowed with a natural
descent datum relative to the canonical morphism
% page: 227
\marginpar{227}
%
\(Y\to S=Y/G\), by virtue of the action of~\(G\) on~\(X_0\)
compatible with its action on~\(Y\). This descent datum is not
effective, since \(a,b\) are not contained in an affine open subset.
The induced descent datum on~\(X'\) relative to
\(g\colon S'\to S\) is then not effective either, as is readily
verified.
\end{remarques}

%\section{Bibliography}
\begin{thebibliography}{D}

\bibitem[D]{VIII.D} J\ptbl \textsc{Giraud}, M\'ethode de la descente,
M\'emoire \no 2 de la Soc.\ Math.\ Fr., 1964.

\bibitem[1]{VIII.1} A\ptbl \textsc{Grothendieck}, S\'eminaire
Bourbaki: G\'eom\'etrie formelle et G\'eom\'etrie alg\'ebrique,
May 1959, \No 182.

\bibitem[2]{VIII.2} A\ptbl \textsc{Grothendieck}, S\'eminaire
Bourbaki: Technique de descente et Th\'eor\`emes d'existence~I,
December 1959, \No 190.

\bibitem[3]{VIII.3} A\ptbl \textsc{Grothendieck}, S\'eminaire
Bourbaki: Technique de descente et Th\'eor\`emes d'existence~III,
February 1961, \No 212.
\end{thebibliography}
\end{filecontents*}
% Original input: drafts/SGA1_VI_10_English.texfrag
\begin{filecontents*}[overwrite,noheader]{sga1-cumulative-input-076.tex}
\section{Cofibered categories and bifibered categories}
\label{VI.10}

Consider a category \(\cal F\) over~\(\cal E\), with projection
functor
\[
  p\colon\cal F\longrightarrow\cal E.
\]
It defines a category \(\cal F^\circ\) over~\(\cal E^\circ\), with
projection functor
\[
  p^\circ\colon\cal F^\circ\longrightarrow\cal E^\circ.
\]
A morphism \(\alpha\colon\eta\to\xi\) in~\(\cal F\) is called
\emph{cocartesian} if it is Cartesian for \(\cal F^\circ\) over
\(\cal E^\circ\). Explicitly, this means that, for every object
\(\xi'\) of~\(\cal F_S\), the map \(u\mto u\circ\alpha\),
\[
  \Hom_S(\xi,\xi')
  \longrightarrow
  \Hom_f(\eta,\xi'),
\]
is bijective. The pair \((\xi,\alpha)\) is then also called a
\emph{direct image}
%
\index{direct image|hyperpage}%
%
of~\(\eta\) under~\(f\), in the category \(\cal F\) over~\(\cal E\).
If such a direct image exists for every \(\eta\) in~\(\cal F_T\), the
direct-image functor under~\(f\) is said to exist; once chosen, it is
denoted by
% page: 182
\marginpar{182}
%
\(f_*^{\cal F}\), or simply \(f_*\).
\label{indnot:ft}\oldindexnot{\(f_*^{\cal F}\) or \(f_*\)|hyperpage}%
It is characterized by an isomorphism of bifunctors on
\(\cal F_T^\circ\times\cal F_S\):
\[
  \Hom_S(f_*(\eta),\xi)
  \isomto
  \Hom_f(\eta,\xi).
\]

If \(f_*\) exists, then \(f^*\) exists if and only if
\ifthenelse{\boolean{orig}}{\(f\)}{\(f_*\)}
admits a right adjoint, \ie if and only if there exist a functor
\[
  f^*\colon\cal F_S\longrightarrow\cal F_T
\]
and an isomorphism of bifunctors
\[
\ifthenelse{\boolean{orig}}
{\Hom_E(f_*(\eta),\xi)\isomto\Hom_T(\eta,f^*(\xi)).}
{\Hom_S(f_*(\eta),\xi)\isomto\Hom_T(\eta,f^*(\xi)).}
\]

Let \(g\colon U\to T\) be another morphism in~\(\cal E\), and suppose
that the inverse and direct images under \(f,g,\) and \(fg\) exist.
Consider the functorial homomorphisms
\[
\begin{array}{rcl}
c^{f,g}\colon f_*g_*&\longleftarrow&(fg)_*,\\
c_{f,g}\colon g^*f^*&\longrightarrow&(fg)^*.
\end{array}
\]
If \(f_*g_*\) and \(g^*f^*\) are regarded as a pair of adjoint
functors, as are \((fg)_*\) and \((fg)^*\), then these two
homomorphisms are adjoint to one another. Hence one is an isomorphism
if and only if the other is. In particular:

\begin{proposition}
\label{VI.10.1}
Suppose that the category \(\cal F\) over~\(\cal E\) is both
prefibered and coprefibered. Then it is fibered if and only if it is
\ifthenelse{\boolean{orig}}{cofibered}{cofibered}.
\end{proposition}

Of course, \(\cal F\) is called coprefibered, respectively cofibered,
%
\index{cofibered category|hyperpage}%
%
over~\(\cal E\) if \(\cal F^\circ\) is prefibered, respectively
fibered, over~\(\cal E^\circ\). We call \(\cal F\)
\emph{bifibered over~\(\cal E\)}
%
\index{bifibered category|hyperpage}%
%
if it is both fibered and cofibered over~\(\cal E\).

\end{filecontents*}
% Original input: drafts/SGA1_VI_11_English.texfrag
\begin{filecontents*}[overwrite,noheader]{sga1-cumulative-input-077.tex}
\section{Miscellaneous examples}
\label{VI.11}

\begin{enumerate}

\item[a)] \textbf{Arrow categories of}~$\cal{E}$. Let
$\cal{E}$ be a category. Denote by $\mathbf{\Delta^1}$ the category
associated with the totally ordered two-element set $[0,1]$;
% page: 183
\marginpar{183}
%
it thus has two objects, 0 and 1, and, in addition to the two identity
morphisms, an arrow $(0,1)$ with source 0 and target 1. Let
$$
\CatFl(\cal{E})=\SheafHom(\mathbf{\Delta^1},\cal{E})
$$
which is called the \emph{arrow category of} $\cal{E}$. The object 1
of~$\mathbf{\Delta^1}$ defines a canonical functor, called the
\emph{target functor},
$$
\CatFl(\cal{E})\to \cal{E}
$$
(the functor defined by the object 0 of~$\mathbf{\Delta^1}$ is called
the \emph{source functor}). For every object $S$ of~$\cal{E}$, the
fiber category $\CatFl(\cal{E})_S$ is canonically isomorphic to the
category $\cal{E}_{/S}$ of objects of~$\cal{E}$ over~$S$.

Consider a morphism $f\colon T\to S$ in $\cal{E}$. It corresponds to
a canonical functor
$$
f_*\colon \cal{E}_{/T}=\cal{F}_T\to \cal{E}_{/S}=\cal{F}_S
$$
and a functorial isomorphism
$$
\Hom_S(f_*(\eta),\xi)\isomto\Hom_f(\eta,\xi)
$$
which thus makes~$f_*$ a direct-image functor for~$f$ in~$\cal{F}$.
Moreover, here one has
$$
(\id_S)_*=\id_{\cal{F}_S},\quad (fg)_*=f_*g_*,\quad
\ifthenelse{\boolean{orig}}{c^{f,g}=\id_{(fg)},}{c^{f,g}=\id_{(fg)_*},}
$$
\ie $\cal{F}$ is equipped with a cosplitting over~$\cal{E}$. A
fortiori, $\cal{F}$ is cofibered over~$\cal{E}$. Note now that the set
of morphisms in $\cal{F}$ is in bijection with the set of commutative
square diagrams in $\cal{E}$:
$$
\xymatrix{
X\ar[d]_-u&\ar[l]_-{f'}Y\ar[d]^-v\\
S&\ar[l]_-{f}T
}
$$
By
% page: 184
\marginpar{184}
%
definition, the morphism in question is Cartesian if the square is
Cartesian in $\cal{E}$, \ie if it makes~$Y$ a fiber product of~$X$ and
$T$ over~$S$. The inverse-image functor
\ifthenelse{\boolean{orig}}
{$f$}
{$f^*$}
% correction: $f^*$
therefore exists if and only if, for every object $X$ over~$S$, the
fiber product $X\times_ST$ exists. It follows from~\Ref{VI.10.1} that
if the product of two objects over a third always exists in~$\cal{E}$,
\ie if $\cal{F}$ is prefibered over~$\cal{E}$, then $\cal{F}$ is in
fact fibered over~$\cal{E}$.

\medskip

\item[b)] \textbf{Category of presheaves or sheaves on varying
spaces.}

Let $\cal{E}=\mathbf{Top}$ be the category of topological spaces.  If
$T$ is a topological space, we denote by $\cal{U}(T)$ the category of
open subsets of~$T$, whose morphisms are the inclusion maps.  If
$\cal{C}$ is a category, a functor
$\cal{U}(T)^\circ\to\cal{C}$ is called a \emph{presheaf} on~$T$ with
values in~$\cal{C}$, and a \emph{sheaf} if it satisfies a left-exactness
condition that we shall not repeat here.  The \emph{category
$\cal{P}(T)$ of presheaves on~$T$ with values in~$\cal{C}$} is, by
definition, the category
$\SheafHom(\cal{U}(T)^\circ,\cal{C})$, and the category $\cal{F}(T)$
of sheaves on~$T$ with values in~$\cal{C}$ is the full subcategory whose
objects are those objects of $\SheafHom(\cal{U}(T)^\circ,\cal{C})$
that are sheaves.  If $f\colon T\to S$ is a morphism in~$\cal{E}$,
\ie a continuous map of topological spaces, the order-preserving map
$U\mto f^{-1}(U)$ gives a functor
$\cal{U}(S)\to\cal{U}(T)$, hence a functor
$$
f_*\colon\SheafHom(\cal{U}(T)^\circ,\cal{C})\to
\SheafHom(\cal{U}(S)^\circ,\cal{C}),
$$
called the \emph{direct-image functor for presheaves under}~$f$.
It is immediate that the direct image of a sheaf is a sheaf.  Thus the
functor
$$
f_*\colon\cal{P}(T)\to\cal{P}(S)
$$
induces a functor, also denoted
$$
f_*\colon\cal{F}(T)\to\cal{F}(S).
$$
Moreover, one checks trivially, by associativity of the composition of
functors, that for a second continuous map $g\colon U\to T$ one has
$$
\ifthenelse{\boolean{orig}}{(gf)_*=g_*f_*,}{(fg)_*=f_*g_*,}\qquad
\text{and likewise}\qquad (\id_S)_*=\id_{\cal{P}(S)}.
$$
In this way we have obtained a functor
$$
S\mto\cal{P}(S),
$$
\resp
$$
S\mto\cal{F}(S),
$$
from
% page: 185
\marginpar{185}
%
$\cal{E}$ to~$\Cat$.  In fact, we are interested in the corresponding
functor
$$
S\mto\cal{P}(S)^\circ,
\qquad\text{\resp}\qquad
S\mto\cal{F}(S)^\circ.
$$
It defines a cofibered, and indeed cosplit, category over the category
of topological spaces, called the \emph{cofibered category of
presheaves} (\resp \emph{sheaves}) \emph{with values in}~$\cal{C}$
(understood: on varying spaces).  Spelling out the construction
of~\No\Ref{VI.8}, one sees that a morphism from a presheaf $B$ on~$T$
to a presheaf $A$ on~$S$ is a pair $(f,u)$ consisting of a continuous
map from~$T$ to~$S$ and a morphism
$u\colon A\to f_*(B)$ in the category $\cal{P}(S)$.  This description
is equally valid for morphisms of sheaves, since $\cal{F}$ is a full
subcategory of~$\cal{P}$.

In the most important cases, the categories $\cal{P}$ and $\cal{F}$
over~$\cal{E}$ are also fibered categories; that is, for every
continuous map, the direct-image functors
$\cal{P}(T)\to\cal{P}(S)$ and $\cal{F}(T)\to\cal{F}(S)$ have an
adjoint functor, then denoted $f^*$ and called the inverse-image functor
for presheaves, \resp the inverse-image functor for sheaves, under the
continuous map~$f$.  This functor exists, for example, if
$\cal{C}=\Ens$.  One can show that
$f^*\colon\cal{P}(S)\to\cal{P}(T)$ exists whenever inductive limits
(relative to diagrams in the Universe under consideration) exist
in~$\cal{C}$.  The question is less easy for~$\cal{F}$.  Indeed, even
when $\cal{C}=\Ens$, the inverse image of a presheaf that is a sheaf is
not in general a sheaf; in other words, the inverse-image functor for
sheaves is not isomorphic to the functor induced by the inverse-image
functor for presheaves (despite the common notation $f^*$).  Thus
$\cal{F}$ is a cofibered subcategory of~$\cal{P}$, but not a fibered
subcategory; that is, \emph{the inclusion functor
$\cal{F}\to\cal{P}$ is not fibrant}.

The cofibered category $\cal{P}$ can be deduced from a more general
cofibered category---or rather, fibered category---constructed as
follows.  For every category $\cal{U}$ (in the fixed Universe), put
$$
\cal{P}(\cal{U})=\SheafHom(\cal{U},\cal{C}),
$$
and
% page: 186
\marginpar{186}
%
note that $\cal{U}\mto\cal{P}(\cal{U})$ is naturally a contravariant
functor in~$\cal{U}$, from $\Cat$ to~$\Cat$.  It therefore defines a
split category over $\cal{E}=\Cat$, which we denote
$\Cat_{/\!/\cal{C}}$.  The objects of this category are pairs
$(\cal{U},p)$ consisting of a category $\cal{U}$ and a functor
$p\colon\cal{U}\to\cal{C}$, and a morphism from $(\cal{U},p)$ to
$(\cal{V},q)$ is essentially a pair $(f,u)$, where $f$ is a functor
$\cal{U}\to\cal{V}$ and $u$ is a homomorphism of functors
$u\colon p\to qf$.  We leave it to the reader to spell out the
composition of morphisms in $\Cat_{/\!/\cal{C}}$.  The projection
functor
$$
\cal{F}=\Cat_{/\!/\cal{C}}\to\cal{E}=\Cat
$$
assigns to the pair $(\cal{U},p)$ the object~$\cal{U}$; the fiber
category at~$\cal{U}$ is the category
$\SheafHom(\cal{U},\cal{C})$, up to isomorphism.  When inductive
limits exist in~$\cal{C}$, one easily shows that the fibered category
$\Cat_{/\!/\cal{C}}$ over~$\Cat$ is also cofibered over~$\Cat$;
that is, one can define the notion of the \emph{direct image of a
functor} $p\colon\cal{U}\to\cal{C}$ under a functor
$f\colon\cal{U}\to\cal{V}$.  The category of presheaves is obtained
from the preceding fibered category by the base change
$$
\mathbf{Top}^\circ\to\Cat
$$
(the functor $S\mto\cal{U}(S)$ defined above).  This gives a fibered
category over $\mathbf{Top}^\circ$, and on passing to the opposite
category one obtains the cofibered category $\cal{P}$ of presheaves
over~$\mathbf{Top}$.  The notion of inverse image of a functor
corresponds to that of direct image of a presheaf, and the notion of
direct image of a functor to that of inverse image of a presheaf.

\medskip
\item[c)] \textbf{Objects with operators over an object with
operators.}

Let $\cal{F}$ be a category over~$\cal{E}$, and let $S$ be an object
of~$\cal{E}$ on which a group $G$ acts, on the left to fix ideas.
This object with operators can be interpreted as corresponding to a
functor $\lambda\colon\cal{E}'\to\cal{E}$ from the category
$\cal{E}'$ defined by~$G$ (having a single object and $G$ as its group
of endomorphisms) to the category~$\cal{E}$.  By base change it
therefore defines a category $\cal{F}'$ over~$\cal{E}'$, which is
fibered, \resp cofibered, when $\cal{F}$ is so over~$\cal{E}$.  A
% page: 187
\marginpar{187}
%
section of \ifthenelse{\boolean{orig}}{$\cal{E}'$ over~$\cal{F}'$}{$\cal{F}'$ over~$\cal{E}'$} (necessarily Cartesian, since
$\cal{E}'$ is a groupoid and every isomorphism in~$\cal{F}'$ is
Cartesian by~\Ref{VI.6.12}) can also be interpreted as an
$\cal{E}$-functor $\cal{E}'\to\cal{F}$ over~$\lambda$, or as an object
with operators~$\xi$ in~$\cal{F}$ ``over'' the object with
operators~$S$.

\medskip
\item[d)] \textbf{Pairs of quasi-inverse adjoint functors;
autodualities.}

When the base category $\cal{E}$ is reduced to two objects $a,b$ and,
besides the identity arrows, two mutually inverse isomorphisms
$f\colon a\to b$ and $g\colon b\to a$ (that is, $\cal{E}$ is a rigid
connected groupoid with two objects), a normalized cloven category
over~$\cal{E}$ is essentially the same thing as a system consisting of
two categories $\cal{F}_a$ and $\cal{F}_b$ and a \emph{pair of adjoint
functors}
$$
G\colon\cal{F}_a\to\cal{F}_b,
\qquad
F\colon\cal{F}_b\to\cal{F}_a,
$$
which are equivalences of categories (and hence quasi-inverse to one
another).  One takes for $\cal{F}_a$ and $\cal{F}_b$ the fiber
categories of~$\cal{F}$, for $F$ and $G$ the functors $f^*$ and $g^*$,
and the two isomorphisms
$$
u\colon FG\isomto\id_{\cal{F}_a},
\qquad
v\colon GF\isomto\id_{\cal{F}_b}
$$
are $c_{g,f}$ and $c_{f,g}$.  The two usual compatibility conditions
between $u$ and $v$ are nothing other than condition~\Ref{VI.7.4}~B)
for the composites $fgf$ and $gfg$.  It is easy to show that these
conditions suffice to imply that one indeed has a pseudofunctor
$\cal{E}^\circ\to\Cat$.

An interesting case is that in which
$$
\cal{F}_b=\cal{F}_a^\circ,
\qquad G=F^\circ,
\qquad v=u^\circ.
$$
An \emph{autoduality} in a category $\cal{C}$ means the data of a
functor $D\colon\cal{C}\to\cal{C}^\circ$ and an isomorphism
\ifthenelse{\boolean{orig}}{$u\colon DD^\circ\isomto\id_{\cal{C}}$}{$u\colon D^\circ D\isomto\id_{\cal{C}}$} such that $u$ and the
isomorphism
\ifthenelse{\boolean{orig}}{$u^\circ\colon D^\circ D\isomto\id_{\cal{C}^\circ}$}{$u^\circ\colon DD^\circ\isomto\id_{\cal{C}^\circ}$} make
$(D,D^\circ)$ a pair of adjoint functors (necessarily quasi-inverse to
one another).  This condition is written
$$
D(u(x))=u(D(x))\ifthenelse{\boolean{orig}}{qquad}{\qquad}\text{for every $x\in\Ob(\cal{C})$.}
$$

\item[e)]
% page: 188
\marginpar{188}
%
\textbf{Categories over a discrete category $\cal{E}$.}
One says that $\cal{E}$ is a \emph{discrete category} if every arrow in
it is an identity arrow, so that $\cal{E}$ is determined up to unique
isomorphism by the set $I=\Ob(\cal{E})$.  The data of a category
$\cal{F}$ over~$\cal{E}$ is therefore equivalent (up to unique
isomorphism) to the data of a family of categories $\cal{F}_i$
($i\in I$), the fiber categories.  Every category $\cal{F}$ over
$\cal{E}$ is fibered, every $\cal{E}$-functor
$\cal{F}\to\cal{G}$ is Cartesian, and there is a canonical
isomorphism
$$
\SheafHom_{\cal{E}/-}(\cal{F},\cal{G})
\isomto
\underset{i}{\prod}\SheafHom(\cal{F}_i,\cal{G}_i).
$$
In particular, one obtains
$$
\boldsymbol{\Gamma}(\cal{F}/\cal{E})
=\varprojLim\cal{F}/\cal{E}
\isomto
\underset{i}{\prod}\cal{F}_i.
$$

\item[f)] \textbf{Suppose that $\cal{E}$ has exactly two objects
$S$ and $T$,} and, besides the \textbf{identity morphisms, one
morphism} $f\colon T\to S$.  Then a category $\cal{F}$ over
$\cal{E}$ is determined, up to unique $\cal{E}$-isomorphism, by two
categories $\cal{F}_S$ and $\cal{F}_T$ and a bifunctor
$H(\eta,\xi)$ on $\cal{F}_T^\circ\times\cal{F}_S$ with values
in~$\Ens$.  Indeed, if $\cal{F}$ is a category over~$\cal{E}$, one
associates with it the two fiber categories $\cal{F}_S$ and
$\cal{F}_T$ and the bifunctor
$H(\eta,\xi)=\Hom_f(\eta,\xi)$.  We leave it to the reader to spell
out the construction in the opposite direction.  For the category in
question to be fibered (or prefibered, which amounts to the same
thing), it is necessary and sufficient that the functor $H$ be
representable with respect to the argument~$\xi$.  For it to be
cofibered, it is necessary and sufficient that $H$ be representable
with respect to the argument~$\eta$.

\item[g)] Let $\cal{F}=\cal{C}\times\cal{E}$, regarded as a category
over~$\cal{E}$ by means of~$\pr_2$.  Then $\cal{F}$ is fibered and
cofibered over~$\cal{E}$, and is even endowed with a canonical splitting
and cosplitting, corresponding to the constant functor on~$\cal{E}$,
\resp on~$\cal{E}^\circ$, with values in~$\Cat$ and value~$\cal{C}$.
One has
$$
\boldsymbol{\Gamma}(\cal{F}/\cal{E})
\simeq\SheafHom(\cal{E},\cal{C}),
$$
and
% page: 189
\marginpar{189}
%
$\varprojLim\cal{F}/\cal{E}$ corresponds to the full subcategory
formed by the functors $F\colon\cal{E}\to\cal{C}$ that send arbitrary
morphisms to isomorphisms.

\end{enumerate}
\end{filecontents*}
% Original input: drafts/SGA1_VI_12_English.texfrag
\begin{filecontents*}[overwrite,noheader]{sga1-cumulative-input-078.tex}
\section{Functors on a cloven category}
\label{VI.12}

Let \(\cal F\) be a normalized cloven category over~\(\cal E\). For
every object \(S\) of~\(\cal E\), denote by
\[
  i_S\colon\cal F_S\longrightarrow\cal F
\]
the inclusion functor. Thus, for every morphism
\(f\colon T\to S\) in~\(\cal E\), there is a functorial homomorphism
\[
  \alpha_f\colon i_Tf^*\longrightarrow i_S,
\]
where \(f^*\colon\cal F_S\to\cal F_T\) is the base-change functor
defined by the cleavage. Now let
\[
  F\colon\cal F\longrightarrow\cal C
\]
be a functor from~\(\cal F\) to a category~\(\cal C\). For every
\(S\in\Ob(\cal E)\), set
\[
  F_S=F\circ i_S\colon\cal F_S\longrightarrow\cal C,
\]
and, for every \(f\colon T\to S\) in~\(\cal E\), set
\[
  \varphi_f=F*\alpha_f\colon F_Tf^*\longrightarrow F_S.
\]
Thus every functor \(F\colon\cal F\to\cal C\) determines a family
\((F_S)\) of functors \(\cal F_S\to\cal C\) and a family
\((\varphi_f)\) of functorial homomorphisms
\(F_Tf^*\to F_S\). These families satisfy the following conditions:

\[
\text{a)}\qquad
\varphi_{\id_S}=\id_{F_S}.
\]

For morphisms \(f\colon T\to S\) and \(g\colon U\to T\) in~\(\cal E\),
the following square of functorial homomorphisms is commutative:
\[
\xymatrix@C=2cm{
F_Ug^*f^*
  \ar[r]^-{F_U*c_{f,g}}
  \ar[d]_{\varphi_g*f^*}
&
F_U(fg)^*
  \ar[d]^{\varphi_{fg}}
\\
F_Tf^*
  \ar[r]^-{\varphi_f}
&
F_S.
}
\tag{b)}
\]

% page: 190
\marginpar{190}%
%
The first relation is immediate, and the second follows by applying
the functor~\(F\) to the commutative diagram
\[
\xymatrix@C=2cm{
g^*f^*(\xi)
  \ar[r]^-{c_{f,g}(\xi)}
  \ar[d]_{\alpha_g(f^*(\xi))}
&
(fg)^*(\xi)
  \ar[d]^{\alpha_{fg}(\xi)}
\\
f^*(\xi)
  \ar[r]^-{\alpha_f(\xi)}
&
\xi
}
\]
for a varying object \(\xi\) in~\(\cal F_S\).

If \(G\colon\cal F\to\cal C\) is a second functor, giving rise to
functors \(G_S\colon\cal F_S\to\cal C\) and functorial homomorphisms
\(\psi_f\colon G_Tf^*\to G_S\), and if \(u\colon F\to G\) is a
functorial homomorphism, then the homomorphisms \(u*i_S\) give
\[
  u_S\colon F_S\longrightarrow G_S.
\]
For every morphism \(f\colon T\to S\) in~\(\cal E\), the squares
\[
\begin{array}{c}
\xymatrix@C=1.5cm{
F_Tf^*
  \ar[r]^-{\varphi_f}
  \ar[d]_{u_T*f^*}
&
F_S
  \ar[d]^{u_S}
\\
G_Tf^*
  \ar[r]^-{\psi_f}
&
G_S
}
\end{array}
\tag{c)}
\]
are commutative.

\begin{proposition}
\label{VI.12.1}
Let \(\cal H(\cal F,\cal C)\) be the category whose objects are pairs
consisting of a family \((F_S)\), indexed by \(S\in\Ob(\cal E)\), of
functors \(\cal F_S\to\cal C\), and a family \((\varphi_f)\), indexed
by \(f\in\Fl(\cal E)\), of functorial homomorphisms
\(F_Tf^*\to F_S\), satisfying conditions \emph{a) and b)}. Its
morphisms are the families \((u_S)\), indexed by
\(S\in\Ob(\cal E)\), of homomorphisms
% page: 191
\marginpar{191}
%
\(F_S\to G_S\) satisfying commutativity condition \emph{c)} above;
composition is defined by composing the homomorphisms of functors
\(\cal F_S\to\cal C\). Then the two constructions described above
define an \emph{isomorphism}
\[
  K\colon\SheafHom(\cal F,\cal C)
  \isomto
  \cal H(\cal F,\cal C).
\]
\end{proposition}

It is immediate that these constructions define a functor from the
first category to the second. This functor is fully faithful. For
given \(F,G\), the map
\[
  \Hom(F,G)\longrightarrow\Hom(K(F),K(G))
\]
is plainly injective. To prove surjectivity, note that condition c)
expresses the functoriality of the maps
\[
  u(\xi)=u_S(\xi)\colon F(\xi)=F_S(\xi)\longrightarrow
  G(\xi)=G_S(\xi)
\]
for homomorphisms of the form \(\alpha_f(\xi)\) in~\(\cal F\).
Functoriality also holds on every fiber category, \ie for the
morphisms in~\(\cal F\) that are \(T\)-morphisms, with
\(T\in\Ob(\cal E)\). It therefore holds for every morphism in
\(\cal F\), since every \(f\)-morphism, with
\(f\colon T\to S\) in
\ifthenelse{\boolean{orig}}
{\((\cal E)\)}
{\(\cal E\),}
is uniquely the composite of a morphism \(\alpha_f(\xi)\) and a
\(T\)-morphism.

It remains to show that \(K\) is bijective on objects. The preceding
argument already shows that \(K\) is injective on objects. For
surjectivity, start with a system \((F_S),(\varphi_f)\) satisfying
a) and b), and define a map \(\Ob(\cal F)\to\Ob(\cal C)\) by
\[
  F(\xi)=F_S(\xi)
  \qquad
  \text{for }\xi\in\Ob(\cal F_S)\subset\Ob(\cal F),
\]
and a map \(\Fl(\cal F)\to\Fl(\cal C)\) by
\[
  F(\alpha_f(\xi)u')
  =
  \varphi_f(\xi)\,F_T(u')
\]
for every morphism \(f\colon T\to S\) in~\(\cal E\), every object
\(\xi\) of~\(\cal F_S\), and every \(T\)-morphism \(u'\) with target
\(f^*(\xi)\). These maps define a \emph{functor}
\(F\colon\cal F\to\cal C\). Indeed,
\(F(\id_\xi)=\id_{F(\xi)}\) is immediate. It remains to prove
multiplicativity \(F(uv)=F(u)F(v)\) for an \(f\)-morphism
\(u\colon\eta\to\xi\) and a \(g\)-morphism
\(v\colon\zeta\to\eta\), where \(f\colon T\to S\) and
\(g\colon U\to T\) are morphisms of~\(\cal E\). Setting \(w=uv\),
we have
\[
  u=\alpha_f(\xi)u',
  \qquad
  v=\alpha_g(\eta)v',
  \qquad
  w=\alpha_{fg}(\xi)w',
\]

% page: 192
\marginpar{192}%
%
with
\[
  w'=c_{f,g}(\xi)g^*(u')v'
  \qquad(\cf \No\Ref{VI.8}).
\]
With this notation, we must prove commutativity along the outer
boundary of the following diagram:
\[
\xymatrix{
F_U(\zeta)
  \ar[rr]^{F_U(v')}
  \ar[drr]_{F(v)}
  \ar@/^1.5pc/[rrrrrr]^{F_U(w')}
&&
F_Ug^*(\eta)
  \ar[rr]^{F_Ug^*(u')}
  \ar[d]^{\varphi_g(\eta)}
&&
F_Ug^*f^*(\xi)
  \ar[rr]^{F_U(c_{f,g}(\xi))}
  \ar[d]^{\varphi_g(f^*(\xi))}
&&
F_U(fg)^*(\xi)
  \ar[d]^{\varphi_{fg}(\xi)}
\\
&&
F_T(\eta)
  \ar[rr]_{F_T(u')}
  \ar@/_1.2pc/[rrrr]_{F(u)}
&&
F_Tf^*(\xi)
  \ar[rr]_{\varphi_f(\xi)}
&&
F_S(\xi).
}
\]
The left triangle is commutative by the definition of~\(F(v)\). The
middle square is commutative because it is obtained from the
homomorphism \(u'\colon\eta\to f^*(\xi)\) by the functorial
homomorphism~\(\varphi_g\), and the square on the right is commutative
by condition b). The desired conclusion follows.

Now suppose that \(\cal C\) is itself a normalized cloven category
over~\(\cal E\), which we shall henceforth call~\(\cal G\), and that
we are interested in the \(\cal E\)-functors from~\(\cal F\)
to~\(\cal G\). Every such functor \(F\) induces functors
\[
  F_S\colon\cal F_S\longrightarrow\cal G_S
\]
on the fiber categories. Moreover, for every morphism
\(f\colon T\to S\) in~\(\cal E\) and every object \(\xi\) of
\(\cal F_S\), the \(f\)-morphism \(F(\alpha_f(\xi))\) factors uniquely
through a \(T\)-morphism
\[
  \varphi_f(\xi)\colon
  F_T(f_{\cal F}^*(\xi))
  \longrightarrow
  f_{\cal G}^*(F_S(\xi)).
\]
Here the subscript \(\cal F\) or \(\cal G\) indicates the cloven
category whose inverse-image functor is being used. We thereby obtain
a functorial homomorphism between functors from~\(\cal F_S\)
to~\(\cal G_T\):
\[
  \varphi_f\colon
  F_Tf_{\cal F}^*
  \longrightarrow
  f_{\cal G}^*F_S.
\]

% page: 193
\marginpar{193}%
%
The two systems \((F_S)\) and \((\varphi_f)\) satisfy the following
conditions:
\[
\text{a')}\qquad
\varphi_{\id_S}=\id_{F_S}.
\]
For morphisms \(f\colon T\to S\) and \(g\colon U\to T\) in~\(\cal E\),
the following diagram of functorial homomorphisms is commutative:
\[
\xymatrix{
F_Ug_{\cal F}^*f_{\cal F}^*
  \ar[rr]^{F_U*c_{f,g}^{\cal F}}
  \ar[d]_{\varphi_g*f_{\cal F}^*}
&&
F_U(fg)_{\cal F}^*
  \ar[dd]^{\varphi_{fg}}
\\
g_{\cal G}^*F_Tf_{\cal F}^*
  \ar[d]_{g_{\cal G}^**\varphi_f}
&&
\\
g_{\cal G}^*f_{\cal G}^*F_S
  \ar[rr]^{c_{f,g}^{\cal G}*F_S}
&&
(fg)_{\cal G}^*F_S.
}
\tag{b')}
\]
We leave to the reader the verification, as well as the statement and
proof of the analogue of Proposition~\Ref{VI.12.1}. This analogue
shows that the set of \(\cal E\)-functors from~\(\cal F\) to~\(\cal G\)
is in bijection with the set of systems \((F_S),(\varphi_f)\)
satisfying conditions a') and b') above. Under this correspondence,
the Cartesian functors are precisely those for which the
homomorphisms \(\varphi_f\) are isomorphisms.

\begin{remarquestar}
It is usually preferable, of course, to reason directly with fibered
categories without using explicit cleavages. In particular, this
avoids having to use a cumbersome interpretation such as the one
above for the simple notion of an \(\cal E\)-functor or a Cartesian
\(\cal E\)-functor. It was in order to avoid unbearable heaviness and
to obtain more intrinsic statements
% page: 194
\marginpar{194}
%
that we had to abandon the approach taken in~\cite{VI.2}, starting
from the notion of a cloven category, called there a \lq\lq fibered
category\rq\rq{}, which now takes second place behind the notion of a
fibered category. It is moreover likely that, contrary to the still
predominant usage inherited from older habits of thought, universal
problems will eventually prove more conveniently handled by not
emphasizing \emph{one} solution chosen once and for all, but by placing
all solutions on an equal footing.
\end{remarquestar}

\begin{thebibliography}{9}
\bibitem[1]{VI.1}
A\ptbl\ \textsc{Grothendieck}, \emph{Sur quelques points d'algèbre
homologique},
\ifthenelse{\boolean{orig}}
{Tohoku Math.\ Journal, pp\ptbl\ 119--221, Vol\ptbl\ 9, 1957.}
{T\^ohoku Math.\ J. \textbf{9} (1957), pp\ptbl\ 119--221.}

\bibitem[2]{VI.2}
A\ptbl\ \textsc{Grothendieck}, \emph{Technique de descente et
Théorèmes d'existence, I}, Séminaire Bourbaki~190,
December~1959.
\end{thebibliography}

%% Exposé VII does not exist.

\refstepcounter{chapter}
\addtocontents{toc}{\protect\medskip\protect\par\protect\noindent
\protect\textbf{VII: does not exist}}

\end{filecontents*}
% Original input: drafts/SGA1_VI_2_English.texfrag
\begin{filecontents*}[overwrite,noheader]{sga1-cumulative-input-079.tex}
% SGA 1, Exposé VI, French TeX lines 10121--10350.

\section{Categories over another category}
\label{VI.2}

Let \(\mathcal E\) be a category belonging to \(\mathfrak U\), hence an
object of \(\mathbf{Cat}\). We may therefore consider the category
\[
  \mathbf{Cat}_{/\mathcal E}
\]
\label{indnot:fh}
of ``objects of \(\mathbf{Cat}\) over \(\mathcal E\).'' An object of this
category is a functor
\[
  p\colon \mathcal F\longrightarrow\mathcal E.
\]
The category \(\mathcal F\), equipped with such a functor, is also called a
\emph{category over \(\mathcal E\)}, or an \(\mathcal E\)-\emph{category}.
Accordingly, an \(\mathcal E\)-functor from a category \(\mathcal F\) over
\(\mathcal E\) to a category \(\mathcal G\) over \(\mathcal E\) is a functor
\[
  f\colon\mathcal F\longrightarrow\mathcal G
\]
such that
\[
  qf=p,
\]
where \(p\) and \(q\) are the projection functors for \(\mathcal F\) and
\(\mathcal G\), respectively. The set of \(\mathcal E\)-functors from
\(\mathcal F\) to \(\mathcal G\) is thus in bijection with the set of arrows
from \(\mathcal F\) to \(\mathcal G\) in \(\mathbf{Cat}_{/\mathcal E}\).
This is not literally an identity, since the datum of \(f\) alone does not
determine \(\mathcal F\) and \(\mathcal G\) as categories over
\(\mathcal E\); nevertheless, as in every category \(\mathcal C_{/S}\), we
shall routinely make the harmless abuse of identifying
\(\mathcal E\)-functors in the above sense with arrows in
\(\mathbf{Cat}_{/\mathcal E}\).
\marginpar{\raggedright\footnotesize French source p.~149}%

We denote by
\[
  \operatorname{Hom}_{\mathcal E}(\mathcal F,\mathcal G)
\]
the set of \(\mathcal E\)-functors from \(\mathcal F\) to \(\mathcal G\).
Of course, a composite of \(\mathcal E\)-functors is an
\(\mathcal E\)-functor; by definition, the composition in question
corresponds to composition of arrows in \(\mathbf{Cat}_{/\mathcal E}\).

Now consider two \(\mathcal E\)-functors
\[
  f,g\colon\mathcal F\longrightarrow\mathcal G
\]
and a natural transformation
\[
  u\colon f\longrightarrow g.
\]
We call \(u\) an \(\mathcal E\)-\emph{homomorphism}, or a
\emph{homomorphism of \(\mathcal E\)-functors}, if, for every
\(\xi\in\operatorname{Ob}(\mathcal F)\),
\[
  q\bigl(u(\xi)\bigr)=\operatorname{id}_{p(\xi)}.
\]
In words, put
\[
  S=p(\xi)=qf(\xi)=qg(\xi)\in\operatorname{Ob}(\mathcal E).
\]
Then the morphism
\[
  u(\xi)\colon f(\xi)\longrightarrow g(\xi)
\]
in \(\mathcal G\) is an \(\operatorname{id}_S\)-morphism. More generally,
for every morphism \(\alpha\colon T\to S\) in \(\mathcal E\), and every
category \(\mathcal G\) over \(\mathcal E\), a morphism \(v\) in
\(\mathcal G\) is called an \(\alpha\)-\emph{morphism} if
\(q(v)=\alpha\), where \(q\colon\mathcal G\to\mathcal E\) is the
projection functor.

If \(h\colon\mathcal F\to\mathcal G\) is a third
\(\mathcal E\)-functor and \(v\colon g\to h\) is an
\(\mathcal E\)-homomorphism, then \(vu\) is again an
\(\mathcal E\)-homomorphism. Thus the \(\mathcal E\)-functors from
\(\mathcal F\) to \(\mathcal G\), together with their
\(\mathcal E\)-homomorphisms, form a subcategory of the category
\(\operatorname{Fun}(\mathcal F,\mathcal G)\) of all functors from
\(\mathcal F\) to \(\mathcal G\). We call it the \emph{category of
\(\mathcal E\)-functors from \(\mathcal F\) to \(\mathcal G\)} and denote
it by
\[
  \operatorname{Fun}_{\mathcal E/-}(\mathcal F,\mathcal G).
\]
\label{indnot:fi}
\marginpar{\raggedright\footnotesize French source p.~150}%
It is also the kernel subcategory of the pair of functors
\[
  \xymatrix@C=1.1cm{
    \operatorname{Fun}(\mathcal F,\mathcal G)
      \ar@<2pt>[r]^-R\ar@<-2pt>[r]_-S &
    \operatorname{Fun}(\mathcal F,\mathcal E),
  }
\]
where \(R\) is the constant functor defined by the object
\(p\in\operatorname{Fun}(\mathcal F,\mathcal E)\), while \(S\) is the
functor \(f\mapsto q\circ f\) defined by
\(q\colon\mathcal G\to\mathcal E\).

To finish these generalities, it remains to define the natural pairings
among the categories
\(\operatorname{Fun}_{\mathcal E/-}(\mathcal F,\mathcal G)\) by
composition of \(\mathcal E\)-functors. In other words, for three
categories \(\mathcal F,\mathcal G,\mathcal H\) over \(\mathcal E\), we
want to define a ``composition functor''
\begin{equation*}
  \label{eq:VI.2.i}
  \tag{i}
  \operatorname{Fun}_{\mathcal E/-}(\mathcal F,\mathcal G)
  \times\operatorname{Fun}_{\mathcal E/-}(\mathcal G,\mathcal H)
  \longrightarrow
  \operatorname{Fun}_{\mathcal E/-}(\mathcal F,\mathcal H)
\end{equation*}
which, on objects, induces the composition map
\((f,g)\mapsto gf\) for \(\mathcal E\)-functors
\(f\colon\mathcal F\to\mathcal G\) and
\(g\colon\mathcal G\to\mathcal H\).

For this purpose, recall the canonical functor
\begin{equation*}
  \label{eq:VI.2.ii}
  \tag{ii}
  \operatorname{Fun}(\mathcal F,\mathcal G)
  \times\operatorname{Fun}(\mathcal G,\mathcal H)
  \longrightarrow\operatorname{Fun}(\mathcal F,\mathcal H),
\end{equation*}
which, on objects, is just the composition map \((f,g)\mapsto gf\).
It sends an arrow \((u,v)\), where
\[
  u\colon f\to f',\qquad v\colon g\to g'
\]
are arrows in \(\operatorname{Fun}(\mathcal F,\mathcal G)\) and
\(\operatorname{Fun}(\mathcal G,\mathcal H)\), respectively, to the arrow
\[
  v*u\colon gf\longrightarrow g'f'
\]
\label{indnot:fj}
defined by
\[
  (v*u)(\xi)
    =v\bigl(f'(\xi)\bigr)\circ g\bigl(u(\xi)\bigr)
    =g'\bigl(u(\xi)\bigr)\circ v\bigl(f(\xi)\bigr).
\]
\marginpar{\raggedright\footnotesize French source p.~151}%
It is well known that this does define a natural transformation from
\(gf\) to \(g'f'\), and that, as \(f,g,u,v\) vary, it defines the functor
in~\eqref{eq:VI.2.ii}; that is,
\begin{equation*}
  \label{eq:VI.2.I}
  \tag{I}
  \operatorname{id}_g*\operatorname{id}_f=\operatorname{id}_{gf},
\end{equation*}
and
\begin{equation*}
  \label{eq:VI.2.II}
  \tag{II}
  (v'*u')\circ(v*u)=(v'\circ v)*(u'\circ u).
\end{equation*}

Recall also the associativity formula for the canonical pairings
in~\eqref{eq:VI.2.ii}. It is expressed both by the associativity
\((hg)f=h(gf)\) of composition of functors and by
\begin{equation*}
  \label{eq:VI.2.III}
  \tag{III}
  (w*v)*u=w*(v*u)
\end{equation*}
for composition products of natural transformations. Here
\(u\colon f\to f'\) and \(v\colon g\to g'\) are as above, and one is
also given a natural transformation \(w\colon h\to h'\) between functors
\(h,h'\colon\mathcal H\to\mathcal K\).

We now claim that, when \(\mathcal F\) and \(\mathcal G\) are
\(\mathcal E\)-categories, the canonical composition functor
in~\eqref{eq:VI.2.ii} induces the functor in~\eqref{eq:VI.2.i}. Since we
already know that the composite of two \(\mathcal E\)-functors is an
\(\mathcal E\)-functor, this amounts to saying that if
\(u\colon f\to f'\) and \(v\colon g\to g'\) are homomorphisms of
\(\mathcal E\)-functors, then
\(v*u\colon gf\to g'f'\) is again a homomorphism of
\(\mathcal E\)-functors. This follows immediately from the definitions.
Because the pairings in~\eqref{eq:VI.2.i} are induced by those
in~\eqref{eq:VI.2.ii}, they satisfy the same associativity property,
expressed by the formulas \((hg)f=h(gf)\) and
\((w*v)*u=w*(v*u)\) for \(\mathcal E\)-functors and homomorphisms of
\(\mathcal E\)-functors.

To complete the collection of formulas
\eqref{eq:VI.2.I}--\eqref{eq:VI.2.III}, recall also
\begin{equation*}
  \label{eq:VI.2.IV}
  \tag{IV}
  v*\operatorname{id}_{\mathcal F}=v,
  \qquad
  \operatorname{id}_{\mathcal G}*u=u.
\end{equation*}
For brevity, one writes \(v*f\) or \ifthenelse{\boolean{orig}}{\(u*g\)}{\(g*u\)} instead of \(v*u\) when
\(u\), respectively \(v\), is the identity automorphism of \(f\),
respectively \(g\).
\marginpar{\raggedright\footnotesize French source p.~152}%

It follows from the definition of the pairings in~\eqref{eq:VI.2.i} that
\(\operatorname{Fun}_{\mathcal E/-}(\mathcal F,\mathcal G)\) is a functor
in \(\mathcal F\) and \(\mathcal G\), from the product category
\[
  \mathbf{Cat}_{/\mathcal E}^{\mathrm{op}}
  \times\mathbf{Cat}_{/\mathcal E}
  \quad\text{to}\quad\mathbf{Cat}.
\]
Indeed, if \(g\colon\mathcal G\to\mathcal G_1\) is an
\(\mathcal E\)-functor, hence an object of
\(\operatorname{Fun}_{\mathcal E/-}(\mathcal G,\mathcal G_1)\), then,
on putting \(\mathcal H=\mathcal G_1\) in~\eqref{eq:VI.2.i}, it gives a
functor
\[
  g_*\colon
  \operatorname{Fun}_{\mathcal E/-}(\mathcal F,\mathcal G)
  \longrightarrow
  \operatorname{Fun}_{\mathcal E/-}(\mathcal F,\mathcal G_1).
\]
Similarly, an \(\mathcal E\)-functor
\(f\colon\mathcal F_1\to\mathcal F\) defines a functor
\[
  f^*\colon
  \operatorname{Fun}_{\mathcal E/-}(\mathcal F,\mathcal G)
  \longrightarrow
  \operatorname{Fun}_{\mathcal E/-}(\mathcal F_1,\mathcal G).
\]
For short, these functors are also denoted by the symbols
\(f\mapsto g\circ f\) and \(g\mapsto g\circ f\), respectively; strictly
speaking, those symbols describe only the corresponding maps on sets of
objects. The associativity property established above shows that this does
indeed give the announced functor
\[
  \mathbf{Cat}_{/\mathcal E}^{\mathrm{op}}
  \times\mathbf{Cat}_{/\mathcal E}\longrightarrow\mathbf{Cat}.
\]
\end{filecontents*}
% Original input: drafts/SGA1_VI_3_English.texfrag
\begin{filecontents*}[overwrite,noheader]{sga1-cumulative-input-080.tex}
\section{Base change in categories over~$\mathcal{E}$}
\label{VI.3}

Since inverse limits (with respect to categories $\cal{I}$ that are
elements of~$\Univ$) exist in~$\Cat$, the same is true in
$\Cat_{/\cal{E}}$; in particular, Cartesian products exist there, and
are interpreted as fiber products in~$\Cat$. In accordance with the
general notation, if $\cal{F}$ and $\cal{G}$ are categories over~$\cal{E}$,
we denote by
$$
\cal{F}\times_{\cal{E}}\cal{G}
$$
\label{indnot:fk}\oldindexnot{$\cal{F}\times_{\cal{E}}\cal{G}$|hyperpage}%
%
their product in~$\Cat_{/\cal{E}}$, that is, their fiber product over~$\cal{E}$
in~$\Cat$,
% page: 153
\marginpar{153}
%
regarded as a category over~$\cal{E}$. Thus
$\cal{F}\times_{\cal{E}} \cal{G}$ is equipped with two
$\cal{E}$-functors $\pr_1$ and $\pr_2$, which define, for every
category $\cal{H}$ over~$\cal{E}$, a bijection
$$
\Hom_{\cal{E}}(\cal{H},\cal{F}\times_{\cal{E}}\cal{G})\isomto
\Hom_{\cal{E}}(\cal{H},\cal{F})\times\Hom_{\cal{E}}(\cal{H},\cal{G}).
$$
Moreover, this bijection is obtained from an isomorphism of categories
$$
\SheafHom_{\cal{E}/-}(\cal{H},\cal{F}\times_{\cal{E}}
\cal{G})\isomto\SheafHom_{\cal{E}/-}(\cal{H},\cal{F})\times
\SheafHom_{\cal{E}/-}(\cal{H},\cal{G})
$$
by taking the sets of objects on both sides, where the displayed functor
has as its components the functors $h\mto \pr_1\circ h$ and
$h\mto \pr_2\circ h$ from the left-hand side to the two factors on the
right. We leave it to the reader to verify that this indeed gives an
isomorphism (the analogous fact holds, more generally, whenever one has
an inverse limit of categories---and not only in the case of a fiber
product).

Recall also that (as was stated in~\No \Ref{VI.1}):
\begin{eqnarray*}
\Ob(\cal{F}\times_\cal{E}\cal{G})&=&
\Ob(\cal{F})\times_{\Ob(\cal{E})}\Ob(\cal{G})\\
\Fl(\cal{F}\times_\cal{E}\cal{G})&=&
\Fl(\cal{F})\times_{\Fl(\cal{E})}\Fl(\cal{G})\quoi,
\end{eqnarray*}
the composition of arrows being performed componentwise.

In what follows, we consider a functor
$$
\lambda\colon\cal{E}'\to\cal{E}
$$
and, for every category $\cal{F}$ over~$\cal{E}$, regard
$\cal{F}\times_\cal{E}\cal{E}'$ as a category over~$\cal{E}'$ by
means of~$\pr_2$; in other words, we interpret the \og fiber product\fg
operation as a \emph{\og base change\fg} operation, and call the functor
$\lambda\colon\cal{E}'\to\cal{E}$ a \emph{\og base-change functor\fg.}
%
\index{changement de base (foncteur)|hyperpage}%
\index{foncteur changement de base|hyperpage}%
%
In accordance with the familiar general facts, this yields a functor,
called
% page: 154
\marginpar{154}
%
the \emph{base-change functor} for~$\lambda$:
$$
\lambda^*\colon\Cat_{/\cal{E}}\to\Cat_{/\cal{E}'}\quoi,
$$
\label{indnot:fl}\oldindexnot{$\lambda^*\colon\Cat_{/\cal{E}}\to\Cat_{/\cal{E}'}$|hyperpage}%
%
(adjoint to the \og restriction of base\fg functor which, to every
category $\cal{F}'$ over~$\cal{E}'$ with projection functor~$p'$,
assigns~$\cal{F}'$, regarded as a category over~$\cal{E}$ through the
functor~$p=\lambda p'$). As is well known in the general case of a base-change
functor in a category, the base-change functor \og commutes with inverse
limits\fg, and in particular \og transforms\fg fiber products over~$\cal{E}$
into fiber products over~$\cal{E}'$.

Let $\cal{F}$ and $\cal{G}$ be two categories over~$\cal{E}$. We shall
define a \emph{canonical isomorphism}:
\begin{multline*}
\label{eq:VI.3.i}
\tag{i}
\SheafHom_{\cal{E}'/-}(\cal{F}',\cal{G}')\isomto
\SheafHom_{\cal{E}/-}(\cal{F}\times_{\cal{E}}\cal{E}',\cal{G})\\
\text{where }\cal{F}'=\cal{F}\times_{\cal{E}}\cal{E}'\text{, }
\cal{G}'=\cal{G}\times_{\cal{E}}\cal{E}'.
\end{multline*}
For this, consider the functor
$$
\pr_1\colon\cal{G}'=\cal{G}\times_{\cal{E}}\cal{E}'\to \cal{G}\quoi,
$$
and define~\eqref{eq:VI.3.i} by
$$
F\mto \pr_1\circ F\quoi,
$$
which a priori denotes a functor
\begin{equation*}
\label{eq:VI.3.ii}
\tag{ii} {\SheafHom(\cal{F}',\cal{G}')\to\SheafHom(\cal{F}',\cal{G})}
\end{equation*}
It therefore remains only to verify that the latter induces a functor
on the subcategories in~\eqref{eq:VI.3.i}, and that this functor is an
isomorphism. That~\eqref{eq:VI.3.ii} induces a bijection
$$
\Hom_{\cal{E}'/-}(\cal{F}',\cal{G}')\isomto
\Hom_{\cal{E}/-}(\cal{F}\times_{\cal{E}}\cal{E}',\cal{G})
$$
is the characteristic property of the base-change functor. It therefore
remains to
% page: 155
\marginpar{155}
%
prove that if $F$, $G$ are $\cal{E}'$-functors
$\cal{F}'\to\cal{G}'$, then \emph{the map}
$$
u\mto \pr_1\circ u
$$
\emph{induces a bijection}
$$
\Hom_{\cal{E}'}(F,G)\isomto\Hom_{\cal{E}}(\pr_1\circ F,\pr_1\circ
G)\quoi.
$$
The verification of this fact is immediate and is left to the
reader.

It follows from this isomorphism~\eqref{eq:VI.3.i}, and from the end of
the preceding subsection, that
$$
\SheafHom_{\cal{E}'/-}(\cal{F}\times_{\cal{E}}\cal{E}',\cal{G}
\times_{\cal{E}}\cal{E}')
$$
\emph{may be regarded as a functor in}
$\cal{E}',\cal{F},\cal{G}$, \emph{from the category}
$\Cat_{/\cal{E}}^\circ\times\Cat_{/\cal{E}}^\circ\allowbreak\times\Cat_{/\cal{E}}$
\emph{to the category}~$\Cat$, isomorphic to the functor defined
by the expression
$\SheafHom_{\cal{E}/-}(\cal{F}\times_{\cal{E}}\cal{E}',\cal{G})$. In
particular, for fixed $\cal{F}$,~$\cal{G}$, one obtains a functor in
$\cal{E}'$, $\ifthenelse{\boolean{orig}}{\cal{E}'\to}{\cal{E}'\mto}
\ifthenelse{\boolean{orig}}{\SheafHom_{\cal{E}''/-}(\cal{F}',\cal{G}')=}{\SheafHom_{\cal{E}'/-}(\cal{F}',\cal{G}')=}
\SheafHom_{\cal{E}'/-}(\cal{F}\times_{\cal{E}}
\cal{E}',\cal{G}\times_{\cal{E}}\cal{E}')$, and in particular the
projection $\cal{E}$-functor $\lambda\colon \cal{E}'\to\cal{E}$
defines a morphism \ie a functor
$$
\lambda^*_{\cal{F},\cal{G}}\colon
\SheafHom_{\cal{E}/-}(\cal{F},\cal{G})\to
\SheafHom_{\cal{E}'/-}(\cal{F}',\cal{G}')
$$
which we shall make explicit. On the sets of objects of the two
sides, this is the map
$$
f\mto f'=f\times_{\cal{E}}\cal{E}'
$$
which expresses the functorial dependence of
$\cal{F}\times_{\cal{E}}\cal{E}'$ on the object $\cal{F}$
over~$\cal{E}$. On the other hand, consider two $\cal{E}$-functors
$$
f,g\colon\cal{F}\to \cal{G}
$$
and a homomorphism of~$\cal{E}$-functors
$$
u\colon f\to g\quoi,
$$
we
% page: 156
\marginpar{156}
%
shall make explicit the corresponding homomorphism of~$\cal{E}'$-functors:
$$
u'\colon f'\to g'\quoi.
$$
For every
$$
\xi'=(\xi,S')\in\Ob(\cal{F}')
$$
with
$$
\xi\in\Ob(\cal{F}),\quad S'\in\Ob(\cal{E}'),\quad p(\xi)=\lambda(S')=S
$$
the morphism
$$
u'(\xi')\colon f'(\xi')=(f(\xi),S')\to g'(\xi')=(g(\xi),S')\quad
\text{in } \cal{G}'
$$
is defined by the formula
$$
u'(\xi')=(u(\xi),\id_{S'})
$$
(which is indeed an $S'$-morphism in~$\cal{G}'$, since
$q(u(\xi))=\lambda(\id_{S'})=\id_S$).

Now consider an arbitrary $\cal{E}$-functor
$$
\lambda'\colon\cal{E}''\to\cal{E}'
$$
and the corresponding functor
$$
\SheafHom_{\cal{E}'/-}(\cal{F}\times_{\cal{E}}\cal{E}',
\cal{G}\times_{\cal{E}}\cal{E}')\to
\SheafHom_{\cal{E}''/-}(\cal{F}\times_{\cal{E}}\cal{E}'',
\cal{G}\times_{\cal{E}}\cal{E}'')\quoi,
$$
I claim that this functor is none other than the functor obtained by
the preceding procedure, starting from~$\cal{F}'$ and
$\cal{G}'$ over~$\cal{E}'$ and regarding $\cal{E}''$ as a
category over~$\cal{E}'$, taking into account the isomorphisms of
\emph{\og transitivity of base change\fg}:
$$
\ifthenelse{\boolean{orig}}
{\cal{F}'\times_{\cal{E}'}\cal{E}''\isomto \cal{F}''=
\cal{F}\times_{\cal{E}}\cal{E}''\quad
\text{and}\quad\cal{G}'\times_{\cal{E}}\cal{E}''\isomto\cal{G}''
=\cal{G}\times_{\cal{E}}\cal{E}''\quoi,}
{\cal{F}'\times_{\cal{E}'}\cal{E}''\isomto
\cal{F}''=\cal{F}\times_{\cal{E}}\cal{E}''\quad
\text{and}\quad\cal{G}'\times_{\cal{E}'}\cal{E}''\isomto\cal{G}''
=\cal G\times_{\cal E}\cal E''\quoi,}
$$
% correction: a prime is missing; the correct formula should be:
% \cal{F}'\times_{\cal{E}'}\cal{E}''\isomto
% \cal{F}''=\cal{F}\times_{\cal{E}}\cal{E}''\quad
% \text{and}\quad\cal{G}'\times_{\cal{E}'}\cal{E}''\isomto\cal{G}''
% =\cal G\times_{\cal E}\cal E''\quoi,
which induce a canonical isomorphism
$$
\SheafHom_{\cal{E}''/-}(\cal{F}'\times_{\cal{E}'}\cal{E}'',
\cal{G}'\times_{\cal{E}'}\cal{E}'')\isomto
\SheafHom_{\cal{E}''/-}(\cal{F}\times_{\cal{E}}\cal{E}'',
\cal{G}\times_{\cal{E}}\cal{E}'')
$$
% remark: a prime is oddly placed in the original text
% original
The
% page: 157
\marginpar{157}
%
verification of this compatibility is immediate and is left to the
reader.

The functors just defined are compatible with the pairings defined in
the preceding subsection, in the following precise sense: if $\cal{F}$,
$\cal{G}$, and~$\cal{H}$ are categories over~$\cal{E}$ and if we set
$$
\cal{F}'=\cal{F}\times_{\cal{E}}\cal{E}'\quoi,\quad \cal{G}'=
\cal{G}\times_{\cal{E}}\cal{E}'\quoi,\quad \cal{H}'=
\cal{H}\times_{\cal{E}}\cal{E}'\quoi,
$$
then the following diagram of functors is commutative:
$$
\xymatrix{ \SheafHom_{\cal{E}/-}(\cal{F},\cal{G})\times
\SheafHom_{\cal{E}/-}(\cal{G},\cal{H}) \ar[r]
\ar[d]_{\lambda^*_{\cal{F},\cal{G}}\times \lambda^*_{\cal{G},\cal{H}}}
& \SheafHom_{\cal{E}/-}(\cal{F},\cal{H})
\ar[d]^{\lambda^*_{\cal{F},\cal{H}}} \\
\SheafHom_{\cal{E}'/-}(\cal{F}',\cal{G}')\times
\SheafHom_{\cal{E}'/-}(\cal{G}',\cal{H}')\ar[r] &
\SheafHom_{\cal{E}'/-}(\cal{F}',\cal{H}')}
$$
% correction: In the first factor of the lower-left entry of this
% square, a prime is missing on \cal{E} in the \SheafHom_{\cal
% E/-}, so that the arrow to the right is well-defined.
where the horizontal arrows are the composition functors defined in
the preceding subsection. This commutativity is expressed by the
formulas
$$
(gf)'=g'f'
$$
for $f\in\Hom_{\cal{E}}(\cal{F},\cal{G})$,
$g\in\Hom_{\cal{E}}(\cal{G},\cal{H})$ (a formula that simply expresses
the functoriality of base change), and by the formula
$$
(v*u)'=v'*u'
$$
when $u\colon f\to f_1$ is an arrow of
$\SheafHom_{\cal{E}/-}(\cal{F},\cal{G})$ and $v\colon g\to g_1$ is an
arrow of~$\SheafHom_{\cal{E}/-}(\cal{G},\cal{H})$. The verification of
this formula follows readily from the definitions.

In what follows, we shall be principally interested in
$\SheafHom_{\cal{E}}(\cal{F},\cal{G})$ (and certain noteworthy
subcategories thereof) when $\cal{F}=\cal{E}$, and for this reason we
introduce special notation:
$$
\bf{\Gamma}(\cal{G}/\cal{E})=
\SheafHom_{\cal{E}}(\cal{E},\cal{G})\quoi,\quad
\Gamma(\cal{G}/\cal{E})=\Ob(\bf{\Gamma}(\cal{G}/\cal{E}))=
\ifthenelse{\boolean{orig}}
{\Hom_{\cal{E}}(\cal{G},\cal{E})\quoi.}
{\Hom_{\cal{E}}(\cal{E},\cal{G})\quoi.}
$$
% correction: replace `\Hom_{\cal{E}}(\cal{G},\cal{E})\, .' by
% `\Hom_{\cal{E}}(\cal{E},\cal{G})\, .'%
\label{indnot:fm}\oldindexnot{$\bf{\Gamma}(\cal{G}/\cal{E})$ and $\Gamma(\cal{G}/\cal{E})$|hyperpage}%
%

\begin{remarquesstar}
When
% page: 158
\marginpar{158}
%
$\cal{E}$ is a one-point category,
\ie when $\Ob(\cal{E})$ and $\Fl(\cal{E})$ are each reduced to a single
element, which also means that $\cal{E}$ is a final object of the
category~$\Cat$, the datum of a category over~$\cal{E}$ is equivalent
to the datum of a category, plain and simple (since there is a unique
functor from $\cal{F}$ to~$\cal{E}$). More precisely,
$\Cat_{/\cal{E}}$ is then isomorphic to~$\Cat$. Moreover, the
categories $\SheafHom_{\cal{E}/-}(\cal{F},\cal{G})$ are then precisely
the categories $\SheafHom(\cal{F},\cal{G})$. Recall that the
fundamental formula
$$
\Hom(\cal{H},\SheafHom(\cal{F},\cal{G}))\isomto
\Hom(\cal{F}\times\cal{H},\cal{G})
$$
(an isomorphism functorial in the three arguments occurring therein)
allows $\SheafHom(\cal{F},\cal{G})$ to be interpreted axiomatically,
in terms internal to the category~$\Cat$, so that the familiar
formalism of the categories $\SheafHom$ appears as a special case of
a formalism valid in categories such as~$\Cat$, in which \og
$\SheafHom$ objects\fg\ (defined by the preceding formula) exist.
There is an analogous interpretation of
$\SheafHom_{\cal{E}/-}(\cal{F},\cal{G})$ when $\cal{E}$ is again
arbitrary, by means of the formula
$$
\Hom(\cal{H},\SheafHom_{\cal{E}/-}(\cal{F},\cal{G}))\isomto
\Hom_{\cal{E}}(\cal{F}\times\cal{H},\cal{G})
$$
(an isomorphism functorial in the three arguments). Thus the formal
properties set out in \No \Ref{VI.2},~\Ref{VI.3} are special cases of
more general results, valid in categories in which the objects
$\SheafHom_{\cal{E}/-}(\cal{F},\cal{G})$ (when $\cal{F}$ and~$\cal{G}$
are two objects of the category over a third object~$\cal{E}$) exist.
\end{remarquesstar}
\end{filecontents*}
% Original input: drafts/SGA1_VI_4_English.texfrag
\begin{filecontents*}[overwrite,noheader]{sga1-cumulative-input-081.tex}
\section[Fiber categories]{Fiber categories; equivalence of
\(\mathcal E\)-categories}
\label{VI.4}

Let \(\cal F\) be a category over~\(\cal E\), and let
\(S\in\Ob(\cal E)\). The \emph{fiber category of \(\cal F\) at \(S\)}
%
\index{fiber category|hyperpage}%
%
is the subcategory \(\cal F_S\)
\label{indnot:fn}\oldindexnot{\(\cal F_S\)|hyperpage}%
%
of~\(\cal F\) that is the inverse image of the one-point subcategory
of~\(\cal E\) defined by~\(S\). Thus the objects of~\(\cal F_S\) are
the objects \(\xi\) of~\(\cal F\) such that \(p(\xi)=S\), and its
morphisms are the morphisms \(u\) of~\(\cal F\) such that
\(p(u)=\id_S\), \ie the \(S\)-morphisms in~\(\cal F\).

% page: 159
\marginpar{159}
%
Of course, \(\cal F_S\) is canonically isomorphic to the fiber product
\(\cal F\times_{\cal E}\{S\}\), where \(\{S\}\) denotes the one-point
subcategory of~\(\cal E\) defined by~\(S\), equipped with its inclusion
functor into~\(\cal E\). It follows, in view of the transitivity of base
change, that if one makes the base change
\(\lambda\colon\cal E'\to\cal E\), then for every
\(S'\in\Ob(\cal E')\), \emph{the projection}
\[
  \pr_1\colon
  \cal F'=\cal F\times_{\cal E}\cal E'\longrightarrow\cal F
\]
\emph{induces an isomorphism}
\[
  \cal F'_{S'}\longrightarrow\cal F_S
  \qquad\text{(where \(S=\lambda(S')\))}
\ifthenelse{\boolean{orig}}{}{.}
\]

\begin{proposition}
\label{VI.4.1}
Let \(f\colon\cal F\to\cal G\) be an \(\cal E\)-functor. If \(f\) is
fully faithful, then for every base change \(\cal E'\to\cal E\), the
corresponding functor
\[
  f'\colon
  \cal F'=\cal F\times_{\cal E}\cal E'
  \longrightarrow
  \cal G'=\cal G\times_{\cal E}\cal E'
\]
is fully faithful.
\end{proposition}

The verification is immediate. More generally, one can show that every
projective limit of
\ifthenelse{\boolean{orig}}{a fully faithful functor}
{fully faithful functors}
(here \(f\) and the identity functors in~\(\cal E,\cal E'\)) is a
fully faithful functor.

The analogue of Proposition~\Ref{VI.4.1} in which
\ifthenelse{\boolean{orig}}{where}{where}
\lq\lq fully faithful\rq\rq{} is replaced by
\lq\lq equivalence of categories\rq\rq{} is false, already when
\(\cal G=\cal E\). Nevertheless:

\begin{proposition}
\label{VI.4.2}
Let \(f\colon\cal F\to\cal G\) be an \(\cal E\)-functor. The following
conditions are equivalent:
\begin{enumerate}
\item[(i)] There exist an \(\cal E\)-functor
\(g\colon\cal G\to\cal F\) and \(\cal E\)-isomorphisms
\[
  gf\isomto\id_{\cal F},
  \qquad
  fg\isomto\id_{\cal G}
\ifthenelse{\boolean{orig}}{}{.}
\]
\item[(ii)] For every category \(\cal E'\) over~\(\cal E\), the
functor
\[
  f'=f\times_{\cal E}\cal E'\colon
  \cal F'=\cal F\times_{\cal E}\cal E'
  \longrightarrow
  \cal G'=\cal G\times_{\cal E}\cal E'
\]
is an equivalence of categories.
\item[(iii)]
% page: 160
\marginpar{160}
%
\(f\) is an equivalence of categories, and for every
\(S\in\Ob(\cal E)\), the functor
\[
  f_S\colon\cal F_S\longrightarrow\cal G_S
\]
induced by~\(f\) is an equivalence of categories.
\item[(iii~bis)] \(f\) is fully faithful, and for every
\(S\in\Ob(\cal E)\) and every
\ifthenelse{\boolean{orig}}
{\(\eta\in\Ob\cal G_S\),}
{\(\eta\in\Ob(\cal G_S)\),}
there exist \(\xi\in\Ob(\cal F_S)\) and an \(S\)-isomorphism
\ifthenelse{\boolean{orig}}
{\(u\colon f((\xi)\to\eta\).}
{\(u\colon f(\xi)\to\eta\).}
\end{enumerate}
\end{proposition}

\subsubsection*{Proof}
Condition (i) plainly implies that \(f\) is an equivalence of
categories, a notion defined by the same condition except that the
functor isomorphisms are not required to be \(\cal E\)-morphisms.
Moreover, the functoriality established in the preceding section shows
that condition (i) is preserved after a base change
\(\cal E'\to\cal E\). Thus (i)\(\To\)(ii). Clearly
(ii)\(\To\)(iii), since it is enough to take \(\cal E'=\cal E\) and
\(\cal E'=\{S\}\). The implication (iii)\(\To\)(iii~bis) is still more
immediate, so it remains to prove (iii~bis)\(\To\)(i).

For every \(\eta\in\Ob(\cal G)\), choose
\(g(\eta)\in\Ob(\cal F)\) and an isomorphism
\[
  u(\eta)\colon f(g(\eta))\longrightarrow\eta
\]
such that \(q(u(\eta))=\id_S\), where \(S=q(\eta)\). This is possible
by the second condition in (iii~bis). As is well known and immediate,
the full faithfulness of~\(f\) implies that \(g\) can be made, in a
unique way, into a functor from~\(\cal G\) to~\(\cal F\) so that the
\(u(\eta)\) define a functorial homomorphism, hence an isomorphism,
\[
  u\colon fg\isomto\id_{\cal G}.
\]
Moreover, by construction, \(g\) is an \(\cal E\)-functor and \(u\)
an \(\cal E\)-homomorphism. The preceding data then determine a
functorial isomorphism
\[
  v\colon gf\longrightarrow\id_{\cal F},
\]
defined by the condition \(f*v=u*f\), and one sees at once that this,
too, is an \(\cal E\)-morphism. This proves the proposition.

\begin{definition}
\label{VI.4.3}
When the preceding conditions hold, \(f\) is called an
\emph{equivalence of categories over \(\cal E\)}, or an
\emph{\(\cal E\)-equivalence}.
\end{definition}

\begin{corollary}
\label{VI.4.4}
Suppose that the projection functor
\(p\colon\cal F\to\cal E\) is transportable, \ie that for every
isomorphism \(\alpha\colon T\to S\) in~\(\cal E\) and every object
\(\xi\) in~\(\cal F_T\), there exists an isomorphism \(u\) in~\(\cal F\)
with source \(\xi\) such that \(p(u)=\alpha\). Then every
\(\cal E\)-functor \(f\colon\cal F\to\cal G\) that is an equivalence
of categories is an \(\cal E\)-equivalence.
\end{corollary}

This follows
% page: 161
\marginpar{161}
%
from criterion (iii~bis).

\begin{corollary}
\label{VI.4.5}
Let \(f\colon\cal F\to\cal G\) be an \(\cal E\)-equivalence. Then, for
every category \(\cal H\) over~\(\cal E\), the corresponding functors
\begin{align*}
\SheafHom_{\cal E/-}(\cal G,\cal H)
  &\longrightarrow\SheafHom_{\cal E/-}(\cal F,\cal H),\\
\SheafHom_{\cal E/-}(\cal H,\cal F)
  &\longrightarrow\SheafHom_{\cal E/-}(\cal H,\cal G)
\end{align*}
(\cf \No\Ref{VI.2}) are equivalences of categories.
\end{corollary}

This follows from criterion (i) by the usual argument.

\end{filecontents*}
% Original input: drafts/SGA1_VI_5_English.texfrag
\begin{filecontents*}[overwrite,noheader]{sga1-cumulative-input-082.tex}
\section{Cartesian morphisms, inverse images, Cartesian functors}
\label{VI.5}

Let \(\cal F\) be a category over~\(\cal E\), with projection
functor~\(p\).
\ifthenelse{\boolean{orig}}
{}
{\enlargethispage{.5cm}}

\begin{definition}
\label{VI.5.1}
Consider a morphism
\[
  \alpha\colon\eta\longrightarrow\xi
\]
in~\(\cal F\), and let
\[
  S=p(\xi),\qquad T=p(\eta),\qquad f=p(\alpha)
\ifthenelse{\boolean{orig}}{}{.}
\]
The morphism \(\alpha\) is called a \emph{Cartesian morphism}
%
\index{Cartesian morphism|hyperpage}%
%
if, for every \(\eta'\in\Ob(\cal F_T)\) and every \(f\)-morphism
\(u\colon\eta'\to\xi\), there exists a unique \(T\)-morphism
\(\overline u\colon\eta'\to\eta\) such that
\(u=\alpha\circ\overline u\).
\end{definition}

Thus, for every \(\eta'\in\Ob(\cal F_T)\), the map
\(v\mto\alpha\circ v\),
\begin{equation}
\tag{i}
  \Hom_T(\eta',\eta)\longrightarrow\Hom_f(\eta',\xi),
\end{equation}
is bijective. Equivalently, the pair \((\eta,\alpha)\)
\emph{represents the functor in \(\eta'\)}
\(\cal F_T^\circ\to\Ens\) appearing on the right.

If, for a given morphism \(f\colon T\to S\) in~\(\cal E\) and a given
\(\xi\in\Ob(\cal F_S)\), such a pair \((\eta,\alpha)\) exists, \ie if
there is a Cartesian morphism \(\alpha\) in~\(\cal F\) with target
\(\xi\) and \(p(\alpha)=f\), then
% page: 162
\marginpar{162}
%
\(\eta\) is determined in~\(\cal F_T\) up to unique isomorphism. One
then says that \emph{the inverse image of \(\xi\) under \(f\) exists},
and an object \(\eta\) of~\(\cal F_T\), together with a Cartesian
\(f\)-morphism \(\alpha\colon\eta\to\xi\), is called \emph{an inverse
image of \(\xi\) under \(f\)}.
%
\index{inverse image|hyperpage}%
%
Often such an inverse image is chosen whenever it exists, with
\(\cal F\) fixed. It is then denoted by
\(f^*_{\cal F}(\xi)\), or simply \(f^*(\xi)\),
\label{indnot:fo}\oldindexnot{\(f^*_{\cal F}(\xi)\) or \(f^*(\xi)\)|hyperpage}%
%
\ifthenelse{\boolean{orig}}{where \(\xi\times_ST\)}
{or \(\xi\times_ST\)}
when these notations cause no confusion. The canonical morphism
\(\alpha\colon\eta\to\xi\) will be denoted below by~\(\alpha_f(\xi)\).
\label{indnot:fp}\oldindexnot{\(\alpha_f(\xi)\)|hyperpage}%

If the inverse image of every \(\xi\in\Ob(\cal F_S)\) under~\(f\)
exists, one also says that \emph{the inverse-image functor under \(f\)
in~\(\cal F\) exists}; then \(f^*(\xi)\) becomes a covariant functor
in~\(\xi\), from~\(\cal F_S\) to~\(\cal F_T\). This follows because
the right-hand side of (i) depends covariantly on~\(\xi\), or more
precisely defines a functor from
\(\cal F_T^\circ\times\cal F_S\) to~\(\Ens\).

This functorial dependence of \(f^*(\xi)\) is explicit as follows.
Consider Cartesian \(f\)-morphisms
\[
  \alpha\colon\eta\longrightarrow\xi,
  \qquad
  \alpha'\colon\eta'\longrightarrow\xi',
\]
and an \(S\)-morphism \(\lambda\colon\xi\to\xi'\). There exists a
unique \(T\)-morphism \(\mu\colon\eta\to\eta'\) such that
\[
  \alpha'\mu=\lambda\alpha,
\]
as follows from the fact that \(\alpha'\) is Cartesian.

We also record the following immediate fact. Consider a commutative
diagram
\[
\xymatrix@C=1.2cm{
\xi\ar[d]_-{\lambda}& \ar[l]_-{\alpha}\eta \ar[d]^-{\mu}\\
\xi' &\ar[l]_-{\alpha'}\eta'
}
\]
in~\(\cal F\),
% page: 163
\marginpar{163}
%
where \(\alpha\) and \(\alpha'\) are \(f\)-morphisms, \(\lambda\) is
an \(S\)-isomorphism, and \(\mu\) a \(T\)-isomorphism. Then
\emph{\(\alpha\) is Cartesian if and only if \(\alpha'\) is
Cartesian}.

\begin{definition}
\label{VI.5.2}
An \(\cal E\)-functor \(F\colon\cal F\to\cal G\) is called a
\emph{Cartesian functor}
%
\index{Cartesian functor|hyperpage}%
%
if it carries Cartesian morphisms to Cartesian morphisms. Denote by
\(\SheafHom_\cart(\cal F,\cal G)\)
\label{indnot:fq}\oldindexnot{\(\SheafHom_\cart(\cal F,\cal G)\)|hyperpage}%
%
the full subcategory of
\(\SheafHom_{\cal E/-}(\cal F,\cal G)\) formed by the Cartesian
functors.

For example, regard~\(\cal E\) as a category over itself by means of
the identity functor. Every morphism of~\(\cal E\) is Cartesian, so a
Cartesian functor from~\(\cal E\) to~\(\cal F\) is a section functor
\(F\colon\cal E\to\cal F\) that carries every morphism of~\(\cal E\)
to a Cartesian morphism. Such a functor is called a
\emph{Cartesian section} of~\(\cal F\) over~\(\cal E\).
\end{definition}

\begin{proposition}
\label{VI.5.3}
\begin{enumerate}
\item[(i)] A functor \(F\colon\cal F\to\cal G\) that is an
\(\cal E\)-equivalence is Cartesian.
\item[(ii)] Let \(F,G\colon\cal F\to\cal G\) be isomorphic
\(\cal E\)-functors. If one is Cartesian, then so is the other.
\item[(iii)] The composite of two Cartesian functors
\(\cal F\to\cal G\) and \(\cal G\to\cal H\) is Cartesian.
\end{enumerate}
\end{proposition}

Assertion (iii) is immediate from the definition, (ii) follows from
the remark preceding~\Ref{VI.5.2}, and (i) follows readily from the
definition and criterion~\Ref{VI.4.2}~(iii). More precisely, a
morphism \(\alpha\) in~\(\cal F\) is Cartesian if and only if
\(F(\alpha)\) is Cartesian.

\begin{corollary}
\label{VI.5.4}
Let \(F\colon\cal F\to\cal G\) be an \(\cal E\)-equivalence. Then,
for every category \(\cal H\) over~\(\cal E\), the corresponding
functors \(G\mto G\circ F\) and \(G\mto F\circ G\) induce equivalences
of categories
\begin{eqnarray*}
\SheafHom_\cart(\cal G,\cal H)&\lto{\approx}&
\SheafHom_\cart(\cal F,\cal H),\\
\SheafHom_\cart(\cal H,\cal F)&\lto{\approx}&
\SheafHom_\cart(\cal H,\cal G).
\end{eqnarray*}
\end{corollary}

This follows in the usual way from criterion~\Ref{VI.4.2}~(i) and
\Ref{VI.5.3}~(i), (ii), and~(iii).

% page: 164
\marginpar{164}
%
More precisely, \emph{an \(\cal E\)-functor
\(G\colon\cal G\to\cal H\) is Cartesian if and only if
\(G\circ F\) is Cartesian}; likewise, \emph{an \(\cal E\)-functor
\(G\colon\cal H\to\cal F\) is Cartesian if and only if
\(F\circ G\) is Cartesian}.

It follows from~\Ref{VI.5.3}~(iii) that, if one considers the
subcategory \(\Cat^\cart_{/\cal E}\)
\label{indnot:fr}\oldindexnot{\(\Cat^\cart_{/\cal E}\)|hyperpage}%
%
of~\(\Cat_{/\cal E}\) having the same objects as~\(\Cat_{/\cal E}\)
and Cartesian functors as its morphisms, then, as in~\no\Ref{VI.2},
there are composition pairings
\[
  \SheafHom_\cart(\cal F,\cal G)\times
  \SheafHom_\cart(\cal G,\cal H)
  \longrightarrow
  \SheafHom_\cart(\cal F,\cal H)
\]
\ifthenelse{\boolean{orig}}{induced}{induced}
by those of~\no\Ref{VI.2}. These allow
\(\SheafHom_\cart(\cal F,\cal G)\) to be regarded as a functor in
\(\cal F,\cal G\), from
\(\bigl(\Cat^\cart_{/\cal E}\bigr)^\circ
\times\Cat^\cart_{/\cal E}\)
\ifthenelse{\boolean{orig}}{to~\(\Cat\)).}{to~\(\Cat\).}
We shall chiefly need this remark when \(\cal F=\cal G\).

\begin{definition}
\label{VI.5.5}
Let \(\cal F\) be a category over~\(\cal E\). Denote by
\[
  \varprojLim\cal F/\cal E
\]
\label{indnot:fs}\oldindexnot{\(\varprojLim\cal F/\cal E\)|hyperpage}%
%
the category of Cartesian \(\cal E\)-functors
\(\cal E\to\cal F\), \ie of Cartesian sections of~\(\cal F\)
over~\(\cal E\).
\end{definition}

By what has just been said, \(\varprojLim\cal F/\cal E\) is a functor
in~\(\cal F\), from the category \(\Cat^\cart_{/\cal E}\) to the
category~\(\Cat\).

We shall see below the relations between this operation
\(\varprojLim\) and the notion of a projective limit of categories, as
well as many examples.

\end{filecontents*}
% Original input: drafts/SGA1_VI_6_OPENING_English.texfrag
\begin{filecontents*}[overwrite,noheader]{sga1-cumulative-input-083.tex}
\section[Fibered categories and prefibered categories]{Fibered
categories and prefibered categories. Products and base change therein}
\label{VI.6}

\begin{definition}
\label{VI.6.1}
A category $\cal{F}$ over~$\cal{E}$ is called a \emph{fibered category}
%
\index{fibered (category)|hyperpage}%
%
(and the functor $\cal{F}\to\cal{E}$ is then said to be
\emph{fibered})
%
\index{fibered (functor)|hyperpage}%
%
if it satisfies the following two axioms:
\begin{itemize}
\item{$\Fib_{\mathrm{I}}$}
% page: 165
\marginpar{165}
%
For every morphism $f\colon T\to S$ in~$\cal{E}$, the inverse-image
functor by~$f$ in~$\cal{F}$ exists.
\item{$\Fib_{\mathrm{II}}$} The composite of two Cartesian morphisms is
Cartesian.
\end{itemize}
A category $\cal{F}$ over~$\cal{E}$ satisfying condition
$\Fib_{\mathrm{I}}$ is called a \emph{prefibered category over~$\cal{E}$}.
%
\index{prefibered (category)|hyperpage}%
%
\end{definition}

If~$\cal{F}$ is a fibered category
(\resp prefibered category) over~$\cal{E}$, a subcategory
$\cal{G}$ of~$\cal{F}$ is called a \emph{fibered subcategory}
%
\index{fibered subcategory (\resp prefibered)|hyperpage}%
%
(\resp a \emph{prefibered subcategory}) if it is a fibered category
(\resp prefibered category) over~$\cal{E}$ and, moreover, the inclusion
functor is Cartesian. If, for example, $\cal{G}$ is a \emph{full}
subcategory of~$\cal{F}$, one sees that this means that, for every
morphism $f\colon T\to S$ in $\cal{E}$ and every
$\xi \in \Ob(\cal{G}_S)$, $f^*_{\cal{F}}(\xi)$ is $T$-isomorphic to
an object of~$\cal{G}_T$. Another interesting case is the following:
with $\cal{F}$ a fibered category over~$\cal{E}$, consider the
subcategory $\cal{G}$ of~$\cal{F}$ having the same objects and whose
morphisms are the \emph{Cartesian} morphisms of~$\cal{F}$; in
particular, the morphisms of~$\cal{G}_S$ are the isomorphisms of
$\cal{F}_S$. One sees at once that this is indeed a fibered subcategory
of~$\cal{F}$, since in the bijection
$$
\Hom_T(\eta',\eta) \isomto \Hom_f(\eta',\xi)
$$
associated with a Cartesian $f$-morphism $\alpha$ in~$\cal{F}$, the
$T$-isomorphisms on the left correspond to the Cartesian morphisms on
the right. By definition, the Cartesian sections
$\cal{E}\to\cal{F}$
\ifthenelse{\boolean{orig}}{(correspond}{correspond}
% correction: `(correspond' should be `correspond'.
bijectively to arbitrary $\cal{E}$-functors
$\cal{E}\to\cal{G}$ (but note that the natural functor
$$
\SheafHom_{\cal{E}/-}(\cal{E},\cal{G})\to
\SheafHom_\cart(\cal{E},\cal{F})=
\varprojLim(\cal{F}/\cal{E})
$$
is faithful, but in general is not fully faithful, \ie is not an
isomorphism).

\begin{remarquesstar}
Let $\cal{F}$ be a category over~$\cal{E}$. The following conditions
are equivalent:~(i)~Every morphism of~$\cal{F}$ is Cartesian;
(ii)~$\cal{F}$ is a fibered category over~$\cal{E}$, and the
$\cal{F}_S$ are groupoids%
\ifthenelse{\boolean{orig}}{,}{}
% correction: `groupoids,' should be `groupoids'
% (no comma before a parenthesis).
(\ie every morphism in $\cal{F}_S$ is an isomorphism). In this case
$\cal{F}$ is said to be a category \emph{fibered in groupoids}
%
\index{fibered in groupoids|hyperpage}%
%
% page: 166
\marginpar{166}
%
over~$\cal{E}$. These are the categories encountered especially in
\og moduli theory\fg. If $\cal{E}$ is a groupoid, one shows that
conditions (i) and (ii) are also equivalent to the following:
(iii)~$\cal{F}$ is a groupoid, and the projection functor
$p\colon \cal{F}\to\cal{E}$ is transportable (\cf \Ref{VI.4.4}). For
example, if $\cal{E}$ and $\cal{F}$ are groupoids such that
\ifthenelse{\boolean{orig}}
{$\Ob \cal{E}$ and $\Ob \cal{F}$}
{$\Ob(\cal{E})$ and $\Ob(\cal{F})$}
% correction: `\Ob \cal{E}' should be `\Ob(\cal{E})'.
% correction: `\Ob \cal{F}' should be `\Ob(\cal{F})'.
are each reduced to one point, so that $\cal{E}$ and $\cal{F}$ are
defined, up to isomorphism, by groups $E$ and $F$, and the functor
$p\colon \cal{F}\to\cal{E}$ is defined by a group homomorphism
$p\colon F\to E$, then $\cal{F}$ is fibered over~$\cal{E}$ if and
only if $p$ is surjective, \ie if $p$ defines an \emph{extension} of
the group $E$ by the group $G = \Ker p$.
\end{remarquesstar}
\end{filecontents*}
% Original input: drafts/SGA1_VI_6_PROPOSITIONS_English.texfrag
\begin{filecontents*}[overwrite,noheader]{sga1-cumulative-input-084.tex}
\begin{proposition}
\label{VI.6.2}
Let \(F\colon\cal F\to\cal G\) be an \(\cal E\)-equivalence. Then
\(\cal F\) is a fibered (respectively, prefibered) category over
\(\cal E\) if and only if \(\cal G\) is.
\end{proposition}

This follows readily from the definitions and from the observation
made above that a morphism
\ifthenelse{\boolean{orig}}{\(\alpha x\)}{\(\alpha\)}
in~\(\cal F\) is Cartesian if and only if \(F(\alpha)\) is Cartesian.

\begin{proposition}
\label{VI.6.3}
Let \(\cal F_1,\cal F_2\) be two categories over~\(\cal E\), and let
\(\alpha=(\alpha_1,\alpha_2)\) be a morphism in
\(\cal F=\cal F_1\times_{\cal E}\cal F_2\). Then \(\alpha\) is
Cartesian if and only if both its components are Cartesian.
\end{proposition}

Let \(\xi_i\) be the target and \(\eta_i\) the source of~\(\alpha_i\),
and let \(f\colon T\to S\) be the morphism of~\(\cal E\) for which
\(\alpha_1\) and \(\alpha_2\) are \(f\)-morphisms. For every
\(\eta'=(\eta'_1,\eta'_2)\) in~\(\cal F_T\), there is a commutative
diagram
\[
\xymatrix{
\Hom_T(\eta',\eta) \ar[r] \ar[d] &
\Hom_f(\eta',\xi) \ar[d]\\
\Hom_T(\eta'_1,\eta_1)\times\Hom_T(\eta'_2,\eta_2) \ar[r] &
\Hom_f(\eta'_1,\xi_1)\times\Hom_f(\eta'_2,\xi_2)
}
\]
whose vertical arrows are bijections. Thus, if either horizontal arrow
is a bijection, then so is the other. This already shows that if
\(\alpha_1\) and \(\alpha_2\) are Cartesian, so that the lower
horizontal arrow is bijective, then \(\alpha\) is Cartesian.
Conversely, in the preceding diagram first put
\(\eta'_2=\eta_2\), so that
\(\Hom_T(\eta'_2,\eta_2)\ne\emptyset\); this proves that \(\alpha_1\)
is Cartesian. Then put \(\eta'_1=\eta_1\), proving that \(\alpha_2\)
is Cartesian.

\begin{corollary}
\label{VI.6.4}
% page: 167
\marginpar{167}
%
Let \(\cal F=\cal F_1\times_{\cal E}\cal F_2\), and let
\(F=(F_1,F_2)\) be an \(\cal E\)-functor \(\cal G\to\cal F\). Then
\ifthenelse{\boolean{orig}}{\(\cal F\)}{\(F\)}
is Cartesian if and only if \(F_1\) and \(F_2\) are Cartesian. This
gives an isomorphism of categories
\[
  \SheafHom_\cart(\cal G,\cal F_1\times_{\cal E}\cal F_2)
  \isomto
  \SheafHom_\cart(\cal G,\cal F_1)
  \times
  \SheafHom_\cart(\cal G,\cal F_2),
\]
and, in particular, by taking \(\cal G=\cal E\), an isomorphism of
categories
\[
  \varprojLim(\cal F_1\times_{\cal E}\cal F_2/\cal E)
  \isomto
  \varprojLim(\cal F_1/\cal E)
  \times
  \varprojLim(\cal F_2/\cal E).
\]
\end{corollary}

\begin{corollary}
\label{VI.6.5}
Let \(\cal F_1\) and \(\cal F_2\) be two fibered (respectively,
prefibered) categories over~\(\cal E\). Then their fiber product
\(\cal F=\cal F_1\times_{\cal E}\cal F_2\) is a fibered
(respectively, prefibered) category over~\(\cal E\).
\end{corollary}

These results extend to the fiber product of an arbitrary family of
categories over~\(\cal E\).

\begin{proposition}
\label{VI.6.6}
Let \(\cal F\) be a category over~\(\cal E\), with projection functor
\(p\), and let \(\lambda\colon\cal E'\to\cal E\) be a functor. Regard
\[
  \cal F'=\cal F\times_{\cal E}\cal E'
\]
as a category over~\(\cal E'\) by means of the projection functor
\(p'=p\times_{\cal E}\id_{\cal E'}\). Let \(\alpha'\) be a morphism
of~\(\cal F'\). Then \(\alpha'\) is Cartesian if and only if its image
\(\alpha\) in~\(\cal F\) is Cartesian.
\end{proposition}

The proof is immediate and is left to the reader.

\begin{corollary}
\label{VI.6.7}
For every Cartesian functor \(F\colon\cal F\to\cal G\) between
categories over~\(\cal E\), the functor
\[
  F'=F\times_{\cal E}\cal E'\colon
  \cal F'=\cal F\times_{\cal E}\cal E'
  \longrightarrow
  \cal G'=\cal G\times_{\cal E}\cal E'
\]
is Cartesian.
\end{corollary}

Consequently, the functor
\[
  \SheafHom_{\cal E}(\cal F,\cal G)
  \longrightarrow
  \SheafHom_{\cal E'}(\cal F',\cal G')
\]
considered in
\ifthenelse{\boolean{orig}}{\ignorespaces}{\no}
\Ref{VI.3} induces a functor
\[
  \SheafHom_\cart(\cal F,\cal G)
  \longrightarrow
  \SheafHom_\cart(\cal F',\cal G').
\]
In other words, for fixed \(\cal F,\cal G\), one may regard
\[
  \SheafHom_\cart
  (\cal F\times_{\cal E}\cal E',
   \cal G\times_{\cal E}\cal E')
\]
\emph{as
% page: 168
\marginpar{168}
%
a functor in~\(\cal E'\), from the category
\({\Cat_{/\cal E}}^\circ\) to~\(\Cat\)}. If \(\cal F,\cal G\) are
also allowed to vary, one obtains
\ifthenelse{\boolean{orig}}{,}{}
a functor from
\[
  {\Cat_{/\cal E}}^\circ
  \times
  \bigl(\Cat^\cart_{/\cal E}\bigr)^\circ
  \times
  \Cat^\cart_{/\cal E}
\]
to~\(\Cat\).

Taking into account the isomorphism
\[
  \Hom_{\cal E'}(\cal F',\cal G')
  \isomto
  \Hom_{\cal E}(\cal F\times_{\cal E}\cal E',\cal G)
\]
considered in~\no\Ref{VI.3}, the Cartesian \(\cal E'\)-functors from
\(\cal F'\) to~\(\cal G'\) correspond to the \(\cal E\)-functors
\[
  \cal F\times_{\cal E}\cal E'\longrightarrow\cal G
\]
that carry every morphism whose first projection is a Cartesian
morphism of~\(\cal F\) to a Cartesian morphism of~\(\cal G\). Taking
\(\cal F=\cal E\) and changing notation gives:

\begin{corollary}
\label{VI.6.8}
\(\varprojLim(\cal F'/\cal E')\) is isomorphic to the full subcategory
of \(\SheafHom_{\cal E/-}(\cal E',\cal F)\) formed by those
\(\cal E\)-functors \(\cal E'\to\cal F\) that carry arbitrary
morphisms to Cartesian morphisms. In particular, if \(\cal F\) is a
fibered category and \(\widetilde{\cal F}\) is the subcategory
of~\(\cal F\) whose morphisms are the Cartesian morphisms of~\(\cal F\),
then there is a bijection
\[
  \Ob\varprojLim(\cal F'/\cal E')
  \isomto
  \Hom_{\cal E/-}(\cal E',\widetilde{\cal F}).
\]
\end{corollary}

This specifies how the expression
\[
  \varprojLim(\cal F\times_{\cal E}\cal E'/\cal E')
\]
is to be regarded as a functor in~\(\cal E'\) and~\(\cal F\), from
\({\Cat_{/\cal E}}^\circ\times\Cat^\cart_{/\cal E}\) to~\(\Cat\).
We shall later obtain a fuller functorial dependence on~\(\cal E'\)
when \(\cal F\) is required to be a fibered category.

\begin{corollary}
\label{VI.6.9}
Let \(\cal F\) be a fibered (respectively, prefibered) category
over~\(\cal E\). Then
\(\cal F'=\cal F\times_{\cal E}\cal E'\) is a fibered
(respectively, prefibered) category over~\(\cal E'\).
\end{corollary}

\begin{proposition}
\label{VI.6.10}
Let \(\cal F,\cal G\) be prefibered categories over~\(\cal E\), and
let \(F\) be a Cartesian \(\cal E\)-functor from~\(\cal F\)
to~\(\cal G\). Then \(F\) is faithful
% page: 169
\marginpar{169}
%
\ifthenelse{\boolean{orig}}
{respectively fully faithful, (respectively an \(\cal E\)-equivalence)}
{(respectively fully faithful, respectively an \(\cal E\)-equivalence)}
if and only if, for every
\ifthenelse{\boolean{orig}}
{\(S\in\Ob\cal E\),}
{\(S\in\Ob(\cal E)\),}
the induced functor
\[
  F_S\colon\cal F_S\longrightarrow\cal G_S
\]
is faithful (respectively fully faithful, respectively an
equivalence).
\end{proposition}

The proof is immediate from the definitions.

To conclude this section, we give some properties of fibered categories
that use axiom \(\Fib_{\mathrm{II}}\).

\begin{proposition}
\label{VI.6.11}
Let \(\cal F\) be a prefibered category over~\(\cal E\). Then
\(\cal F\) is fibered if and only if it satisfies the following
condition:

\(\Fib_{\mathrm{II}}'\): Let \(\alpha\colon\eta\to\xi\) be a
Cartesian morphism in~\(\cal F\) lying over a morphism
\(f\colon T\to S\) of~\(\cal E\). For every morphism
\(g\colon U\to T\) in~\(\cal E\), and every
\ifthenelse{\boolean{orig}}
{\(\zeta\in\Ob\cal F_U\),}
{\(\zeta\in\Ob(\cal F_U)\),}
the map \(u\mto\alpha\circ u\),
\[
  \Hom_g(\zeta,\eta)\longrightarrow\Hom_{fg}(\zeta,\xi),
\]
is bijective.
\end{proposition}

In other words, in a category \emph{fibered} over~\(\cal E\),
Cartesian diagrams are characterized by a property that is a priori
stronger than the one in the
\ifthenelse{\boolean{orig}}
{definition, (obtained by taking \(g=\id_T\) in}
{definition (obtained by taking \(g=\id_T\) in}
the preceding statement).

\begin{corollary}
\label{VI.6.12}
Let \(\cal F\) be a category over~\(\cal E\), and let \(\alpha\) be a
morphism in~\(\cal F\). If \(\alpha\) is an isomorphism, then
\(p(\alpha)=f\) is an isomorphism and \(\alpha\) is Cartesian. The
converse holds if \(\cal F\) is fibered over~\(\cal E\).
\end{corollary}

Indeed, if \(\alpha\) is an isomorphism, then
\(f=p(\alpha)\) plainly is as well. For every
\ifthenelse{\boolean{orig}}
{\(\eta'\in\Ob\cal F_T\),}
{\(\eta'\in\Ob(\cal F_T)\),}
the map \(u\mto\alpha\circ u\),
\[
  \Hom(\eta',\eta)\longrightarrow\Hom(\eta',\xi),
\]
is bijective. Since \(f\) is an isomorphism, an element on the left is
a \(T\)-morphism if and only if its image on the right is an
\(f\)-morphism. We therefore obtain a bijection
\[
  \Hom_T(\eta',\eta)\longrightarrow\Hom_f(\eta',\xi),
\]
which
% page: 170
\marginpar{170}
%
proves the first assertion.

Conversely, suppose that \(f\) is an isomorphism and that \(\alpha\)
satisfies condition \(\Fib_{\mathrm{II}}'\), which, when \(\cal F\)
is
\ifthenelse{\boolean{orig}}
{fibered}
{fibered}
over~\(\cal E\), means that \(\alpha\) is Cartesian. Then one sees at
once that, for every \(\zeta\in\Ob(\cal F)\), the map
\[
  \Hom(\zeta,\eta)\longrightarrow\Hom(\zeta,\xi),
  \qquad u\longmapsto\alpha\circ u,
\]
is bijective. Thus \(\alpha\) is an isomorphism.

\begin{corollary}
\label{VI.6.13}
Let \(\alpha\colon\eta\to\xi\) and
\(\beta\colon\zeta\to\eta\) be two composable morphisms in the
category \(\cal F\) fibered over~\(\cal E\). If \(\alpha\) is
Cartesian, then \(\beta\) is Cartesian if and only if
\(\alpha\beta\) is Cartesian.
\end{corollary}

Use the strengthened form of the definition of Cartesian morphisms
given in~\Ref{VI.6.11}.

\end{filecontents*}
% Original input: drafts/SGA1_VI_7_English.texfrag
\begin{filecontents*}[overwrite,noheader]{sga1-cumulative-input-085.tex}
% SGA 1, Exposé VI, French TeX lines 11484--11746.

\section{Cloven categories over \(\mathcal E\)}
\label{VI.7}

\begin{definition}
\label{VI.7.1}
Let \(\mathcal F\) be a category over \(\mathcal E\). A
\emph{cleavage} of \(\mathcal F\) over \(\mathcal E\) is a function
that assigns to every \(f\in\operatorname{Fl}(\mathcal E)\) an
inverse-image functor \(f^*\) for \(f\) in \(\mathcal F\). The
cleavage is called \emph{normalized} if
\(f=\operatorname{id}_S\) implies
\(f^*=\operatorname{id}_{\mathcal F_S}\). A \emph{cloven category}
(respectively a \emph{normalized cloven category}) is a category
\(\mathcal F\) over \(\mathcal E\) equipped with a cleavage
(respectively a normalized cleavage).
\end{definition}

It is clear that \(\mathcal F\) admits a cleavage if and only if it is
prefibered over \(\mathcal E\), and in that case it admits a normalized
cleavage. The set of cleavages on \(\mathcal F\) is in bijection with
the set of subsets \(K\subset\operatorname{Fl}(\mathcal F)\) satisfying
the following conditions:
\begin{enumerate}
\item[(a)] Every \(\alpha\in K\) is cartesian.
\item[(b)] For every morphism \(f\colon T\to S\) in \(\mathcal E\) and
  every \(\xi\in\operatorname{Ob}(\mathcal F_S)\), there is a unique
  \(f\)-morphism in \(K\) with target \(\xi\).
\end{enumerate}
The cleavage defined by \(K\) is normalized if and only if \(K\) also
satisfies
\begin{enumerate}
\item[(c)] The identity morphisms of \(\mathcal F\) belong to \(K\).
\end{enumerate}
The elements of \(K\) will be called the \emph{transport morphisms} for
the cleavage in question.
\marginpar{\raggedright\footnotesize French source p.~171}%

The notion of an isomorphism of cloven categories over \(\mathcal E\)
is clear. More generally, one may define morphisms of cloven
\(\mathcal E\)-categories to be the \(\mathcal E\)-functors
\(\mathcal F\to\mathcal G\) that send transport morphisms to transport
morphisms; in particular, these are cartesian functors. In this way the
cloven categories over \(\mathcal E\) are the objects of a category,
the \emph{category of cloven categories over \(\mathcal E\)}. The
reader may spell out the existence of products: if a category over
\(\mathcal E\) is the product of categories \(\mathcal F_i\) over
\(\mathcal E\), each supplied with a cleavage, then their product has
the corresponding natural cleavage. We likewise leave the notion of
base change for cloven categories to the reader.

We denote by \(\alpha_f(\xi)\) the canonical morphism
\[
  \alpha_f(\xi)\colon f^*(\xi)\longrightarrow\xi.
\]
As already noted, it is functorial in \(\xi\); that is, there is a
natural transformation
\[
  \alpha_f\colon i_Tf^*\longrightarrow i_S,
\]
where, for every \(S\in\operatorname{Ob}(\mathcal E)\),
\(i_S\colon\mathcal F_S\to\mathcal F\) denotes the inclusion functor.

Now consider morphisms
\[
  f\colon T\longrightarrow S,
  \qquad g\colon U\longrightarrow T
\]
in \(\mathcal E\), and let
\(\xi\in\operatorname{Ob}(\mathcal F_S)\). There is a unique
\(U\)-morphism
\[
  c_{f,g}(\xi)\colon g^*f^*(\xi)\longrightarrow(fg)^*(\xi)
\]
making the diagram
\[
\xymatrix{
 f^*(\xi) \ar[d]_{\alpha_f(\xi)} &&
 g^*f^*(\xi) \ar[ll]_{\alpha_g(f^*(\xi))}
              \ar[d]^{c_{f,g}(\xi)}\\
 \xi && (fg)^*(\xi) \ar[ll]^{\alpha_{fg}(\xi)}
}
\]
commutative, by the definition of \((fg)^*(\xi)\).
\marginpar{\raggedright\footnotesize French source p.~172}%
As \(\xi\) varies, this morphism is functorial; in other words, there
is a natural transformation
\[
  c_{f,g}\colon g^*f^*\longrightarrow(fg)^*
\]
between functors \(\mathcal F_S\to\mathcal F_U\).

\begin{proposition}
\label{VI.7.2}
A cloven category \(\mathcal F\) over \(\mathcal E\) is fibered if
and only if all the transformations \(c_{f,g}\) are isomorphisms.
\end{proposition}

Taking \(f\) to be an isomorphism and \(g\) its inverse, and considering
\(c_{f,g}\) and \(c_{g,f}\), gives the following.

\begin{corollary}
\label{VI.7.3}
If \(\mathcal F\) is a cloven category fibered over \(\mathcal E\),
then, for every isomorphism \(f\colon T\to S\) in \(\mathcal E\), the
functor
\[
  f^*\colon\mathcal F_S\longrightarrow\mathcal F_T
\]
is an equivalence of categories.
\end{corollary}

\begin{proposition}
\label{VI.7.4}
Let \(\mathcal F\) be a cloven category over \(\mathcal E\). Then
\[
\textup{A)}\qquad
\begin{cases}
 c_{f,\operatorname{id}_T}(\xi)
   =\alpha_{\operatorname{id}_T}(f^*(\xi)),\\
 c_{\operatorname{id}_S,f}(\xi)
   =f^*(\alpha_{\operatorname{id}_S}(\xi)),
\end{cases}
\]
and
\[
\textup{B)}\qquad
 c_{f,gh}(\xi)\circ c_{g,h}(f^*(\xi))
 =c_{fg,h}(\xi)\circ h^*(c_{f,g}(\xi)).
\]
Here \(f,g,h\) are successive morphisms
\[
  V\longrightarrow U\longrightarrow T\longrightarrow S
\]
and \(\xi\) is an object of \(\mathcal F_S\).
\end{proposition}
\marginpar{\raggedright\footnotesize French source p.~173}%

For a normalized cleavage, the first two relations take the simpler
form
\begin{equation*}
\label{eq:VI.7.A'}
\tag{A'}
 c_{f,\operatorname{id}_T}=\operatorname{id}_{f^*},
 \qquad
 c_{\operatorname{id}_S,f}=\operatorname{id}_{f^*}.
\end{equation*}
The third relation is visualized by the commutative diagram
\begin{equation*}
\label{eq:VI.7.D}
\tag{D}
\xymatrix@C=2.2cm{
 h^*g^*f^*(\xi)
   \ar[r]^-{c_{g,h}(f^*(\xi))}
   \ar[d]_{h^*(c_{f,g}(\xi))} &
 (gh)^*f^*(\xi) \ar[d]^{c_{f,gh}(\xi)}\\
 h^*(fg)^*(\xi) \ar[r]_-{c_{fg,h}(\xi)} & (fgh)^*(\xi).
}
\end{equation*}
For fibered categories, where the \(c_{f,g}\) are isomorphisms, this
commutativity says intuitively that successive use of isomorphisms of
the form \(c_{f,g}\) does not lead to ``contradictory identifications.''
The same formula, written without the argument \(\xi\) and using the
convolution product of natural transformations, is
\[
 c_{fg,h}\circ(h^**c_{f,g})
 =c_{f,gh}\circ(c_{g,h}*f^*).
\]

The first two formulas in \hyperref[VI.7.4]{VI.7.4} are immediate; we sketch the proof of
the third. In addition to square (D), consider the square of natural
transformations
\begin{equation*}
\label{eq:VI.7.D'}
\tag{D'}
\xymatrix@C=2cm{
 g^*f^*(\xi) \ar[r]^-{\alpha_g(f^*(\xi))}
                 \ar[d]_{c_{f,g}(\xi)} &
 f^*(\xi) \ar[d]^{\alpha_f(\xi)}\\
 (fg)^*(\xi) \ar[r]_-{\alpha_{fg}(\xi)} & \xi,
}
\end{equation*}
which is commutative by the definition of \(c_{f,g}(\xi)\).
\marginpar{\raggedright\footnotesize French source p.~174}%
Join the vertices of (D) to the corresponding vertices of (D') by the
natural transformations
\[
\begin{array}{ll}
 \alpha_h(g^*f^*(\xi)),&\alpha_{gh}(f^*(\xi)),\\
 \alpha_h((fg)^*(\xi)),&\alpha_{fgh}(\xi).
\end{array}
\]
The four lateral faces of the resulting cube are also commutative. For
the left face this follows because the left column of (D) is obtained
from the left column of (D') by applying \(h^*\), and because
\(\alpha_h\) is natural. For the other three faces it is simply the
definition of the \(c\)-operations on the remaining three sides of
(D). Thus the five faces other than the top face commute. It follows
that the two \((fgh)\)-morphisms
\[
  h^*g^*f^*(\xi)\longrightarrow(fgh)^*(\xi)
\]
defined by (D) have the same composite with
\(\alpha_{fgh}(\xi)\colon(fgh)^*(\xi)\to\xi\); hence they are equal by
the definition of \((fgh)^*\).

For what follows, restrict attention to \emph{normalized} cloven
categories. Such a category gives rise to the following data:
\begin{enumerate}
\item[(a)] A map \(S\mapsto\mathcal F_S\) from
  \(\operatorname{Ob}(\mathcal E)\) to \(\mathbf{Cat}\).
\item[(b)] A map \(f\mapsto f^*\) assigning to every
  \(f\in\operatorname{Fl}(\mathcal E)\), with source \(T\) and target
  \(S\), a functor
  \(f^*\colon\mathcal F_S\to\mathcal F_T\).
\item[(c)] A map \((f,g)\mapsto c_{f,g}\) assigning to every composable
  pair of arrows \((f,g)\) of \(\mathcal E\) a natural transformation
  \(c_{f,g}\colon g^*f^*\to(fg)^*\).
\end{enumerate}
These data satisfy conditions (A') and (B) above. If we had not
restricted to normalized cleavages, one further datum would be needed:
a map \(S\mapsto\alpha_S\) assigning to every object \(S\) of
\(\mathcal E\) a natural transformation
\[
  \alpha_S\colon(\operatorname{id}_S)^*
  \longrightarrow\operatorname{id}_{\mathcal F_S};
\]
condition (A') would then be replaced by condition (A).

We shall now show how the preceding data reconstruct, up to unique
isomorphism, the normalized cloven category \(\mathcal F\) over
\(\mathcal E\).
\marginpar{\raggedright\footnotesize French source p.~175}%
\end{filecontents*}
% Original input: drafts/SGA1_VI_8_English.texfrag
\begin{filecontents*}[overwrite,noheader]{sga1-cumulative-input-086.tex}
\section[Cloven category defined by a pseudofunctor]{Cloven category
defined by a pseudofunctor \(\mathcal E^\circ\to\Cat\)}
\label{VI.8}

For brevity, call a \emph{pseudofunctor} from~\(\cal E^\circ\)
to~\(\Cat\)---more precisely, a normalized pseudofunctor---a set of
data a), b), c) as above satisfying conditions \(A')\) and~\(B)\).
In the preceding section we associated a pseudofunctor
\(\cal E^\circ\to\Cat\) with a normalized cloven category over~\(\cal E\);
we now describe the inverse construction. We leave to the reader the
verification of most of the details, as well as the fact that these
constructions are indeed inverse to one another. More precisely, one
should regard the pseudofunctors \(\cal E^\circ\to\Cat\) as the objects
of a new category and show that our constructions give mutually
quasi-inverse equivalences between this category and the category of
cloven categories over~\(\cal E\) defined in the preceding section.

Set
\[
  \cal F_\circ
  =
  \underset{S\in\Ob(\cal E)}{\coprod}
  \ifthenelse{\boolean{orig}}
  {\Ob\cal F(S),}
  {\Ob(\cal F(S)),}
\]
the disjoint union of the sets
\ifthenelse{\boolean{orig}}
{\(\Ob\cal F(S)\)}
{\(\Ob(\cal F(S))\)}.
(N.B. Here we write \(\cal F(S)\), rather than \(\cal F_S\), for the
value of the given pseudofunctor at the object \(S\) of~\(\cal E\), in
order to avoid notational confusion below.) There is an evident map
\[
  p_\circ\colon\cal F_\circ\longrightarrow\Ob(\cal E).
\]
Let
\[
  \ifthenelse{\boolean{orig}}
  {\xi}
  {\overline\xi}
  =(S,\xi),
  \qquad
  \overline\eta=(T,\eta)
  \qquad
  \text{with \(\xi\in\Ob(\cal F(S))\),
  \(\eta\in\Ob(\cal F(T))\)}
\]
be two elements of~\(\cal F_\circ\), and let
\(f\in\Hom(T,S)\). Define
\[
  h_f(\overline\eta,\overline\xi)
  =
  \Hom_{\cal F(T)}(\eta,f^*(\xi)).
\]

% page: 176
\marginpar{176}%
%
If, in addition, \(g\colon U\to T\) is a morphism in~\(\cal E\) and
\(\zeta\in\Ob(\cal F(U))\), define a map, denoted
\((u,v)\mto u\circ v\),
\[
  h_f(\overline\eta,\overline\xi)
  \times
  h_g(\overline\zeta,\overline\eta)
  \longrightarrow
  h_{fg}(\overline\zeta,\overline\xi),
\]
\ie a map
\[
\begin{split}
  \Hom_{\cal F(T)}(\eta,f^*(\xi))
  \times
  \Hom_{\cal F(U)}(\zeta,g^*(\eta))
  \longrightarrow
  \Hom_{\cal F(U)}(\zeta,(fg)^*(\xi)),
\end{split}
\]
by the formula
\[
  u\circ v=c_{f,g}(\xi)\cdot g^*(u)\cdot v.
\]
Thus \(u\circ v\) is the composite
\[
  \zeta
  \lto{v}
  g^*(\eta)
  \lto{g^*(u)}
  g^*f^*(\xi)
  \lto{c_{f,g}(\xi)}
  (fg)^*(\xi).
\]
Also set
\[
  h(\overline\eta,\overline\xi)
  =
  \underset{f\in\Hom(T,S)}{\coprod}
  h_f(\overline\eta,\overline\xi).
\]
The preceding pairings then define pairings
\[
  h(\overline\eta,\overline\xi)
  \times
  h(\overline\zeta,\overline\eta)
  \longrightarrow
  h(\overline\zeta,\overline\xi),
\]
while the definition of \(h(\overline\eta,\overline\xi)\) gives an
evident map
\[
  p_{\overline\eta,\overline\xi}\colon
  h(\overline\eta,\overline\xi)
  \longrightarrow
  \Hom(T,S).
\]
One now verifies the following points:
\begin{enumerate}
\item[1)] Composition between elements of
\ifthenelse{\boolean{orig}}
{\(h(\eta,\xi)\)}
{\(h(\overline\eta,\overline\xi)\)}
is \emph{associative}.

\item[2)]
% page: 177
\marginpar{177}
%
For every \(\overline\xi=(S,\xi)\) in~\(\cal F_\circ\), consider the
\ifthenelse{\boolean{orig}}
{}
{identity}
element of
\[
\ifthenelse{\boolean{orig}}
{h_{\id_S}(\overline\xi,\overline\xi)
 =\Hom_{\cal F_S}(\id_S^*(\xi),\xi)
 =\Hom_{\cal F_S}(\xi,\xi),}
{h_{\id_S}(\overline\xi,\overline\xi)
 =\Hom_{\cal F(S)}(\id_S^*(\xi),\xi)
 =\Hom_{\cal F(S)}(\xi,\xi),}
\]
and its image in \(h(\overline\xi,\overline\xi)\). This element is a
left and right \emph{identity} for composition between elements of the
\(h(\overline\eta,\overline\xi)\).

This already shows that one obtains a category \(\cal F\) by setting
\[
  \Ob(\cal F)=\cal F_\circ,
  \qquad
  \Fl(\cal F)
  =
  \ifthenelse{\boolean{orig}}
  {\underset{\xi,\overline\eta\in\cal F_\circ}{\coprod}
   h(\overline\eta,\overline\xi).}
  {\underset{\overline\xi,\overline\eta\in\cal F_\circ}{\coprod}
   h(\overline\eta,\overline\xi).}
\]
(N.B. One cannot simply take \(\Fl(\cal F)\) to be the
\emph{union} of the sets \(h(\overline\eta,\overline\xi)\), since these
sets need not be disjoint.) Moreover:

\item[3)] The maps
\[
  p_\circ\colon\Ob(\cal F)\to\Ob(\cal E),
  \qquad
  p_1=(p_{\overline\eta,\overline\xi})
  \colon\Fl(\cal F)\to\Fl(\cal E)
\]
define a \emph{functor} \(p\colon\cal F\to\cal E\). Thus \(\cal F\)
becomes a category over~\(\cal E\), and the evident map
\[
  h_f(\overline\eta,\overline\xi)
  \longrightarrow
  \Hom(\overline\eta,\overline\xi)
\]
induces a \emph{bijection}
\[
  h_f(\overline\eta,\overline\xi)
  \isomto
  \Hom_f(\overline\eta,\overline\xi).
\]

\item[4)] The evident maps
\[
  \Ob(\cal F(S))\longrightarrow\cal F_\circ=\Ob(\cal F),
  \qquad
  \Fl(\cal F(S))\longrightarrow\Fl(\cal F),
\]
where the second is defined by the evident maps
\[
  \Hom_{\cal F(S)}(\xi,\xi')
  =
  h_{\id_S}(\overline\xi,\overline{\xi'})
  \longrightarrow
  \Hom(\overline\xi,\overline{\xi'}),
\]
define an \emph{isomorphism}
\[
  i_S\colon\cal F(S)\isomto\cal F_S.
\]

\item[5)] For every object
\(\overline\xi=(S,\xi)\) of~\(\cal F\), and every morphism
\(f\colon T\to S\) of~\(\cal E\), consider
% page: 178
\marginpar{178}
%
the element \(\overline\eta=(T,\eta)\) of~\(\cal F_T\), where
\(\eta=f^*(\xi)\), and the element \(\alpha_f(\xi)\) of
\(\Hom(\overline\eta,\overline\xi)\) that is the image of
\(\id_{f^*(\xi)}\) under the map
\[
  \Hom_{\cal F(T)}(f^*(\xi),f^*(\xi))
  =
  h_f(\overline\eta,\overline\xi)
  \longrightarrow
  \Hom_f(\overline\eta,\overline\xi).
\]
\emph{This element is Cartesian, and is the identity of
\(\overline\xi\) when \(f=\id_S\)}. In other words, the collection of
the \(\alpha_f(\xi)\) defines a \emph{normalized cleavage of
\(\cal F\) over~\(\cal E\)}. Moreover, by construction, the diagram
of functors
\[
\xymatrix{
\cal F(S) \ar[r]^{f^*} \ar[d]_{i_S}
&
\cal F(T) \ar[d]^{i_T}
\\
\cal F_S \ar[r]^{f^*_{\cal F}}
&
\cal F_T
}
\]
is commutative, where \(f^*_{\cal F}\) is the inverse-image functor
under \(f\) relative to the cleavage just defined on~\(\cal F\).
Finally:

\item[6)] Under the isomorphisms \(i_S\), the homomorphisms
\(c_{f,g}\) supplied with the pseudofunctor are carried to the
functorial homomorphisms \(c_{f,g}\) associated with the cleavage
of~\(\cal F\).
\end{enumerate}

We give only the verification of 1), which is, if anything, less
immediate than the others. It suffices to prove associativity of
composition between elements of sets of the form
\(h_f(\overline\eta,\overline\xi)\). Consider morphisms in~\(\cal E\)
\[
  S\xleftarrow{f}T\xleftarrow{g}U\xleftarrow{h}V,
\]
objects
\[
  \xi,\ \eta,\ \zeta,\ \tau
\]
in \(\cal F(S),\cal F(T),\cal F(U),\cal F(V)\), respectively, and
elements
\[
\begin{array}{ccc}
u\in h_f(\overline\eta,\overline\xi)
 &=&
 \Hom_{\cal F(T)}(\eta,f^*(\xi)),\\
v\in h_g(\overline\zeta,\overline\eta)
 &=&
 \Hom_{\cal F(U)}(\zeta,g^*(\eta)),\\
w\in h_h(\overline\tau,\overline\zeta)
 &=&
 \Hom_{\cal F(V)}(\tau,h^*(\zeta)).
\end{array}
\]
We
% page: 179
\marginpar{179}
%
must prove
\[
  (u\circ v)\circ w=u\circ(v\circ w),
\]
an equality in
\(\Hom_{\cal F(V)}(\tau,(fgh)^*(\xi))\). By the definitions, the two
sides are obtained by composing along the upper and lower boundaries,
respectively, of the following diagram:
\[
\xymatrix{
\tau
  \ar[r]^w
  \ar[drr]_{v\circ w}
&
h^*(\zeta)
  \ar[r]^{h^*(v)}
  \ar@/^1.4pc/[rrrrr]^{h^*(u\circ v)}
&
h^*g^*(\eta)
  \ar[rr]^{h^*g^*(u)}
  \ar[d]_{c_{g,h}(\eta)}
&&
h^*g^*f^*(\xi)
  \ar[rr]^{h^*(c_{f,g}(\xi))}
  \ar[d]^{c_{g,h}(f^*(\xi))}
&&
h^*(fg)^*(\xi)
  \ar[d]^{c_{fg,h}(\xi)}
\\
&&
(gh)^*(\eta)
  \ar[rr]^{(gh)^*(u)}
&&
(gh)^*f^*(\xi)
  \ar[rr]^{c_{f,gh}(\xi)}
&&
(fgh)^*(\xi).
}
\]
The middle square is commutative because \(c_{g,h}\) is a functorial
homomorphism, and the square on the right is commutative by
condition~\(B)\) for a pseudofunctor. This proves the assertion.

Of course, when the pseudofunctor in question already comes from a
normalized cloven category \(\cal F'\) over~\(\cal E\), one must still
specify the natural isomorphism between \(\cal F'\) and~\(\cal F\).
We leave the details to the reader.

We also leave to the reader the interpretation, in terms of
pseudofunctors, of the inverse image of a cloven category \(\cal F\)
over~\(\cal E\) under a base-change functor
\(\cal E'\to\cal E\).

\end{filecontents*}
% Original input: drafts/SGA1_VI_9_English.texfrag
\begin{filecontents*}[overwrite,noheader]{sga1-cumulative-input-087.tex}
\section[Example]{Example: a cloven category defined by a functor
\(\cal E^\circ\to\Cat\); split categories over~\(\cal E\)}
\label{VI.9}

Suppose that a functor
\[
\ifthenelse{\boolean{orig}}
{\phi\colon\cal E^\circ\to\Cat}
{\varphi\colon\cal E^\circ\to\Cat}
\]
is given. It defines a pseudofunctor by setting
\[
\ifthenelse{\boolean{orig}}
{\cal F(S)=\phi(S),}
{\cal F(S)=\varphi(S),}
\qquad
f^*=\varphi(f),
\qquad
c_{f,g}=\id_{(fg)^*}.
\]

% page: 180
\marginpar{180}%
%
The construction of the preceding section therefore gives a cloven
category \(\cal F\) over~\(\cal E\), called the category associated
with the functor~\(\varphi\). A cloven category over~\(\cal E\) is
isomorphic to a cloven category defined by a functor
\(\varphi\colon\cal E^\circ\to\Cat\) if and only if it satisfies
\[
  (fg)^*=g^*f^*,
  \qquad
  c_{f,g}=\id_{(fg)^*}.
\]
In terms of the set \(K\) of transport morphisms, this simply says
that \emph{the composite of two transport morphisms is a transport
morphism}. A cleavage of a category \(\cal F\) over~\(\cal E\)
satisfying this condition is called a \emph{splitting}
%
\index{splitting|hyperpage}%
%
of~\(\cal F\) over~\(\cal E\), and a category \(\cal F\) over
\(\cal E\) equipped with a splitting is called a \emph{split category
over~\(\cal E\)}.
%
\index{split category|hyperpage}%
%
This is a special case of the notion of a cloven category. The
category of split categories over~\(\cal E\) is therefore equivalent
to~\(\SheafHom(\cal E^\circ,\Cat)\). Note that a split category over
\(\cal E\) is, in particular, a cloven category over~\(\cal E\).

If \(\cal F\) is a fibered category over~\(\cal E\), a splitting of
\(\cal F\) need not exist. Suppose, for example, that
\(\Ob(\cal E)\) and \(\Ob(\cal F)\) each consist of one element and
that the endomorphisms of those elements form groups \(E\) and \(F\),
respectively. The projection functor \(p\) is then given by a group
homomorphism \(p\colon F\to E\), which is surjective because \(p\) is
fibered. One immediately verifies that the set of cleavages of
\(\cal F\) over~\(\cal E\) is in bijection with the set of maps
\[
  s\colon E\longrightarrow F
  \qquad\text{such that}\qquad
  ps=\id_E;
\]
these are the systems of representatives for the cosets modulo the
subgroup \(G=\Ker p\) of the surjective homomorphism
\(p\colon F\to E\). A cleavage is a splitting if and only if \(s\) is
a group homomorphism. Thus a splitting exists precisely when the
extension \(F\) of~\(E\) by~\(G\) is split; when \(G\) is commutative,
this is expressed by the vanishing of a certain cohomology class in
\(H^2(E,G)\), where \(G\) is regarded as a group on which \(E\)
acts.

Suppose, however, that \(\cal F\) is a fibered category over~\(\cal E\)
whose
% page: 181
\marginpar{181}
%
fiber categories \(\cal F_S\) are \emph{rigid}, \ie every object of
\(\cal F_S\) has trivial automorphism group. It is then easy to prove
that \(\cal F\) admits a splitting over~\(\cal E\). First, the
existence of a splitting is unchanged if \(\cal F\) is replaced by an
\(\cal E\)-equivalent category. We may therefore reduce to the case
in which the \(\cal F_S\) are both rigid and \emph{reduced}, meaning
that two isomorphic objects in~\(\cal F_S\) are equal. If \(G\) is a
rigid reduced category, every isomorphism between two functors
\(H\to G\), with \(H\) arbitrary, is an identity. Consequently, if
\(\cal F\) is fibered over~\(\cal E\) and its fiber categories are
rigid and reduced, then \(\cal F\) has a \emph{unique} cleavage, which
is necessarily a splitting. Thus \(\cal F\) is isomorphic to the
category defined by a functor
\[
  \varphi\colon\cal E^\circ\longrightarrow\Cat
\]
whose \(\varphi(S)\) are rigid \ifthenelse{\boolean{orig}}{discrete}{reduced} categories, and the functor
\(\varphi\) is determined up to isomorphism.
\end{filecontents*}
% Original input: drafts/SGA1_VI_OPENING_THROUGH_1_English.texfrag
\begin{filecontents*}[overwrite,noheader]{sga1-cumulative-input-088.tex}
% SGA 1, Exposé VI, French TeX lines 9929--10120.

\chapter{Fibered categories and descent}
\label{VI}
\marginpar{\raggedright\footnotesize French source p.~145}%

\setcounter{section}{-1}
\section{Introduction}

Contrary to what was announced in the introduction to the preceding
exposé, it has proved impossible to carry out descent in the category
of preschemes, even in special cases, without first developing with
sufficient care the language of descent in general categories.

The notion of “descent” provides the general framework for every
process of “gluing” objects, and consequently for gluing categories.
The most classical gluing situation begins with a topological space
\(X\) and an open cover \((X_i)\). Suppose that for every \(i\) we are
given a fibered space \(E_i\) over \(X_i\), and for every pair
\((i,j)\) an isomorphism
\[
  f_{ji}\colon E_i|_{X_{ij}}\xrightarrow{\sim}E_j|_{X_{ij}},
  \qquad X_{ij}=X_i\cap X_j,
\]
satisfying the familiar transitivity condition, written in abbreviated
form as \(f_{kj}f_{ji}=f_{ki}\). Then there exists a fibered space
\(E\) over \(X\), determined up to isomorphism by the existence of
isomorphisms
\[
  f_i\colon E|_{X_i}\xrightarrow{\sim}E_i
\]
satisfying \(f_{ji}=f_jf_i^{-1}\), with the usual abuse of notation.

Let \(X'=\coprod_iX_i\), which is a fibered space over \(X\), that is,
is equipped with a continuous map \(X'\to X\). The family \((E_i)\)
can be described more concisely as a fibered space \(E'\) over \(X'\).
The family of maps \(f_{ji}\) becomes an isomorphism between the two
pullbacks \(E''_1\) and \(E''_2\) of \(E'\) to
\[
  X''=X'\times_XX'
\]
by the two canonical projections. The gluing condition can then be
written as an identity between isomorphisms of the fibered spaces
\(E'''_1\) and \(E'''_3\) over the triple fiber product
\[
  X'''=X'\times_XX'\times_XX',
\]
where \(E'''_i\) is the pullback of \(E'\) to \(X'''\) by the
\(i\)-th canonical projection.

The construction of \(E\) from \(E'\) and the gluing isomorphism is a
typical descent process.
\marginpar{\raggedright\footnotesize French source p.~146}%
Conversely, starting with a fibered space \(E\) over \(X\), one says
that \(E\) is \emph{locally trivial with fiber \(F\)} if there is an
open cover \((X_i)\) of \(X\) such that every \(E|_{X_i}\) is
isomorphic to \(F\times X_i\). Equivalently, the pullback \(E'\) of
\(E\) to \(X'=\coprod_iX_i\) is isomorphic to \(X'\times F\).

Thus both the gluing of objects and the localization of a property are
connected with the study of certain kinds of base change \(X'\to X\).
In algebraic geometry, many other kinds of base change, notably
faithfully flat morphisms \(X'\to X\), must be regarded as providing
a localization process for preschemes, or other objects, over \(X\).
This kind of localization is used just as much as ordinary topological
localization, of which the latter is in fact a special case. The same
is true, to a lesser extent, in analytic geometry.

Most proofs, being reduced to verifications, are omitted or only
sketched. When useful, we spell out the less evident diagrams arising
in a proof.

\section{Universes, categories, and equivalences of categories}
\label{VI.1}

To avoid certain logical difficulties, we shall use the notion of a
\emph{universe}: a set large enough that the usual operations of set
theory do not lead outside it. A “universe axiom” guarantees that
every object belongs to some universe. For details, see the book in
preparation by C.~Chevalley and the lecturer.
\footnote{The eventual authors are C.~Chevalley and P.~Gabriel. The
book is due to appear in the year 3000. In the meantime, see also
SGA~4, Exposé~I.}
Thus \(\mathbf{Set}\) denotes not the category of all sets, a notion
\label{indnot:fb}
which has no meaning, but the category of sets belonging to a fixed
universe, which will not be indicated in the notation. Likewise,
\(\mathbf{Cat}\) denotes the category of categories belonging to that
\label{indnot:fc}
universe; by definition, the morphisms from an object \(X\) of
\(\mathbf{Cat}\) to another object \(Y\) are the functors from \(X\)
to \(Y\).

If \(\mathcal C\) is a category, we denote by
\(\operatorname{Ob}(\mathcal C)\)\label{indnot:fd} its \emph{set of objects}, and by
\(\operatorname{Fl}(\mathcal C)\)\label{indnot:fe} its \emph{set of arrows}, or
morphisms.
\marginpar{\raggedright\footnotesize French source p.~147}%
We shall therefore write \(X\in\operatorname{Ob}(\mathcal C)\), and
avoid the customary abuse \(X\in\mathcal C\). If \(\mathcal C\) and
\(\mathcal C'\) are categories, a \emph{functor} from \(\mathcal C\)
to \(\mathcal C'\) will always mean what is commonly called a
\emph{covariant} functor. Its specification includes its source and
target categories \(\mathcal C\) and \(\mathcal C'\). The functors
from \(\mathcal C\) to \(\mathcal C'\) form a set
\(\operatorname{Hom}(\mathcal C,\mathcal C')\), which is the set of
objects of the functor category
\(\operatorname{Fun}(\mathcal C,\mathcal C')\)\label{indnot:ff}. By definition, a
\emph{contravariant functor} from \(\mathcal C\) to \(\mathcal C'\)
is a functor from the opposite category \(\mathcal C^\circ\)\label{indnot:fg} to
\(\mathcal C'\).

We shall use the notions of projective and inductive limits of a
functor \(F\colon\mathcal I\to\mathcal C\), and in particular their
most common special cases: Cartesian products and fiber products, the
dual notions of coproducts and amalgamated sums, and the usual formal
properties of these operations.

For example, projective limits exist in the category \(\mathbf{Cat}\)
introduced above, provided the indexing categories \(\mathcal I\)
belong to the chosen universe. The set of objects, respectively the
set of arrows, of a projective-limit category \(\mathcal C\) of
categories \(\mathcal C_i\) is obtained by taking the projective limit
of the sets of objects, respectively arrows, of the categories
\(\mathcal C_i\). The best-known case is the product of a family of
categories. We shall constantly use below the fiber product of two
categories over a third.

For matters concerning categories and functors, pending the book
already mentioned, see the reference cited as
\hyperlink{cite.VI.1}{VI.1}. It is necessarily
very incomplete, even concerning the generalities sketched in this
section.

We take this opportunity to make explicit the notion of equivalence of
categories, which is not presented satisfactorily in that reference.
A functor
\[
  F\colon\mathcal C\longrightarrow\mathcal C'
\]
is called \emph{faithful}, respectively \emph{fully faithful}, if for
every pair of objects \(S,T\) of \(\mathcal C\), the map
\[
  \operatorname{Hom}(S,T)\longrightarrow
  \operatorname{Hom}(F(S),F(T)),\qquad u\longmapsto F(u),
\]
is injective, respectively bijective. The functor \(F\) is called an
\emph{equivalence of categories} if it is fully faithful and every
object \(S'\) of \(\mathcal C'\) is isomorphic to an object of the
form \(F(S)\).
\marginpar{\raggedright\footnotesize French source p.~148}%

Equivalently, there is a functor
\[
  G\colon\mathcal C'\longrightarrow\mathcal C
\]
quasi-inverse to \(F\), with \(GF\) isomorphic to
\(\operatorname{id}_{\mathcal C}\) and \(FG\) isomorphic to
\(\operatorname{id}_{\mathcal C'}\). The data of a functor
\(G\colon\mathcal C'\to\mathcal C\) and an isomorphism
\[
  \varphi\colon FG\xrightarrow{\sim}\operatorname{id}_{\mathcal C'}
\]
amount to choosing, for every \(S'\in\operatorname{Ob}(\mathcal C')\),
an object \(S=G(S')\) of \(\mathcal C\) and an isomorphism
\(u=\varphi(S')\colon F(S)\xrightarrow{\sim}S'\), compatibly with
morphisms. More precisely, there is a unique functor
\(\mathcal C'\to\mathcal C\) having the prescribed map
\(S'\mapsto G(S')\) on objects and for which
\(S'\mapsto\varphi(S')\) is a natural isomorphism
\(FG\to\operatorname{id}_{\mathcal C'}\).

Finally, if \(G\) is quasi-inverse to \(F\), and if isomorphisms
\[
  \varphi\colon FG\xrightarrow{\sim}\operatorname{id}_{\mathcal C'},
  \qquad
  \psi\colon GF\xrightarrow{\sim}\operatorname{id}_{\mathcal C}
\]
are chosen, then the two compatibility conditions on \(\varphi\) and
\(\psi\) stated in \hyperlink{cite.VI.1}{VI.1}, I.1.2 are equivalent. For every chosen
\(\varphi\), there is a unique \(\psi\) for which those conditions
hold.
\end{filecontents*}
% Original input: drafts/SGA1_V_1_5_THROUGH_END_English.texfrag
\begin{filecontents*}[overwrite,noheader]{sga1-cumulative-input-089.tex}
% SGA 1, Exposé V, French TeX lines 7554--7794.

\begin{corollary}
\label{V.1.5}
The prescheme \(X\) is affine, respectively separated, over \(Z\) if
and only if \(Y\) is. If \(X\) is of finite type over \(Z\), then it
is finite over \(Y\). If, in addition, \(Z\) is locally Noetherian,
then \(Y\) is of finite type over \(Z\).
\end{corollary}

Since \(X\) is affine, and hence separated, over \(Y\), it is affine,
respectively separated, over \(Z\) whenever \(Y\) is. Conversely,
suppose first that \(X\) is affine over \(Z\), and let us prove that
\(Y\) is. By Corollary~\ref{V.1.4}, we may suppose that \(Z\) is
affine. We are then reduced to proving that \(Y\) is affine when \(X\)
is affine, and this follows from the explicit description
\(Y=\operatorname{Spec}(A^G)\) in Proposition~\ref{V.1.1}.

Likewise, since \(p\colon X\to Y\) is integral, hence universally
closed, and surjective, it follows that \(Y\) is separated over \(Z\)
whenever \(X\) is. Indeed, in the diagram
\[
\begin{array}{ccc}
X\times_ZX & \xrightarrow{\ p\times_Zp\ } & Y\times_ZY \\
\big\uparrow{\scriptstyle\Delta_{X/Z}} &&
\big\uparrow{\scriptstyle\Delta_{Y/Z}} \\
X & \xrightarrow{\ p\ } & Y,
\end{array}
\]
the morphism \(X\times_ZX\to Y\times_ZY\) is closed. It therefore
maps the closed diagonal of \(X\) to a closed subset of
\(Y\times_ZY\), and this image is precisely the diagonal of \(Y\),
because \(p\) is surjective.

If \(X\) is of finite type over \(Z\), it is a fortiori of finite type
over \(Y\), and is therefore finite over \(Y\), since it is already
integral over \(Y\). Suppose further that \(Z\) is locally Noetherian;
we prove that \(Y\) is of finite type over \(Z\). By
Corollary~\ref{V.1.4}, we may suppose that \(Z\) is affine. The
topological space \(X\) is quasi-compact and \(p\colon X\to Y\) is
surjective, so \(Y\) is also quasi-compact and hence is a finite union
of affine open subsets.
\marginpar{\raggedright\footnotesize French source p.~108}%
Using Corollary~\ref{V.1.4} again, we reduce to the case in which \(Y\)
and \(X\) are affine. The ring \(A\) of \(X\) is then a finitely
generated algebra over the Noetherian ring \(C\) of \(Z\). It is known
that \(B=A^G\) is likewise a finitely generated \(C\)-algebra. Indeed,
\(A\) is integral, hence finite, over a \(C\)-subalgebra
\(B'\subseteq B\) of finite type; since \(B'\) is Noetherian, \(B\) is
also finite over \(B'\), and hence of finite type over \(C\).

\begin{corollary}
\label{V.1.6}
The prescheme \(X\) is affine, respectively a scheme, if and only if
\(Y\) is.
\end{corollary}

\begin{definition}
\label{V.1.7}
\marginpar{\raggedright\footnotesize French source p.~109}%
Let \(X\) be a prescheme on which a finite group \(G\) acts on the
right. The action of \(G\) is called \emph{admissible} if there exists
a morphism \(p\colon X\to Y\) having the properties of
Proposition~\ref{V.1.3}. This implies that \(X/G\) exists and is
isomorphic to \(Y\).
\end{definition}

\begin{proposition}
\label{V.1.8}
Let \(X\) be a prescheme on which the finite group \(G\) acts on the
right. The action is admissible if and only if \(X\) is a union of
\(G\)-invariant affine open subsets, or equivalently, if and only if
every \(G\)-orbit in \(X\) is contained in an affine open subset.
\end{proposition}

The first condition plainly implies the second, and the second in turn
implies the first. Indeed, let \(T\) be a \(G\)-orbit and let \(U\) be
an affine open subset containing it. The intersection of the
translates \(gU\), for \(g\in G\), is a \(G\)-stable open subset
\(U'\) containing \(T\) and contained in the affine open \(U\). In an
affine scheme, every finite subset has a fundamental system of affine
open neighborhoods, so \(T\) has an affine open neighborhood
\(V\subseteq U'\). Its translates \(gV\) are affine and contained in
the separated open subset \(U'\). Their intersection is therefore an
affine, \(G\)-invariant open subset \(U''\) containing \(T\).

Now the condition in Proposition~\ref{V.1.8} is necessary: if a
quotient \(p\colon X\to Y\) as in Proposition~\ref{V.1.3} exists,
take the inverse images of an affine open cover of \(Y\). It is also
sufficient. Indeed, Proposition~\ref{V.1.1} constructs the quotients
\(Y_i=X_i/G\) of the invariant affine opens \(X_i\). In each \(Y_i\),
the image of \(X_i\cap X_j\) is an open subset \(Y_{ij}\), identified
with \(X_{ij}/G\) by Corollary~\ref{V.1.4}. We thereby obtain
isomorphisms \(Y_{ij}\xrightarrow{\sim}Y_{ji}\), with which the
\(Y_i\) can be glued to construct \(Y\). Serre instead constructs
directly the quotient topological space \(Y=X/G\), equips it with the
sheaf \(p_*(\mathcal O_X)^G\), and verifies that \(Y\) becomes a
prescheme satisfying the conditions of Proposition~\ref{V.1.3}.

% The French source deliberately returns to the number 1.7 here.
\setcounter{proposition}{6}

\begin{corollary}
\label{cor:V.1.7}
If the action of \(G\) on \(X\) is admissible, then the action of every
subgroup \(H\subseteq G\) is admissible as well; in particular, \(X/H\)
exists.
\end{corollary}

This can also be checked directly in the situation of
Proposition~\ref{V.1.3}. We may always suppose that \(X\) is affine
over some \(Z\) and that the elements of \(G\) act by
\(Z\)-automorphisms, for example by taking \(Z=Y\). Indeed:

\begin{corollary}
\label{cor:V.1.8}
Suppose that \(X\) is affine over \(Z\) and that the operations of
\(G\) are \(Z\)-automorphisms. Then the action of \(G\) on \(X\) is
admissible. If \(X\) is defined by a quasi-coherent sheaf of algebras
\(\mathcal A\), then \(Y\) is defined by the sheaf of invariant
algebras \(\mathcal A^G\).
\end{corollary}

\begin{proposition}
\label{V.1.9}
Suppose that \(G\) acts admissibly on \(X\), that \(X/G=Y\) is a
\(Z\)-prescheme, and let \(Z'\to Z\) be a base-change morphism. Put
\[
  X'=X\times_ZZ',\qquad Y'=Y\times_ZZ'.
\]
The group \(G\) acts on \(X'\) by transport of structure, and the
morphism \(p'\colon X'\to Y'\) is invariant. If \(Z'\) is flat over
\(Z\), then \(p'\) again satisfies the hypotheses of
Proposition~\ref{V.1.3}; that is,
\[
  \mathcal O_{Y'}\xrightarrow{\sim}p'_*(\mathcal O_{X'})^G.
\]
Thus \(G\) acts admissibly on \(X'\), and
\[
  (X/G)\times_ZZ'\xrightarrow{\sim}(X\times_ZZ')/G.
\]
\end{proposition}

We may plainly take \(Z=Y\), and are then reduced to the case in which
\(Y\) and \(Y'\) are affine. We must show the following. Let \(B\) be
the invariant subring for the action of \(G\) on \(A\), and let \(B'\)
be a flat \(B\)-algebra. Then \(B'\) is the invariant subring of
\(A'=A\otimes_BB'\). This is immediate, since the exact sequence
\[
  0\longrightarrow B\xrightarrow{i}A\xrightarrow{j}A^{(G)},
\]
where the last term denotes a power of \(A\) and
\(j(x)=(s\cdot x-x)_{s\in G}\), remains exact after tensoring with
\(B'\).

The flatness hypothesis is essential. In particular, if \(Y'\) is a
closed subscheme of \(Y\), even a closed point, and \(X'\) is its
inverse image in \(X\), then in general \(Y'\) is \emph{not}
identified with \(X'/G\). We shall see that it is so nevertheless when
\(X\) is étale over \(Y\).

We conclude with a convenient, if elementary, formalism. Let \(Y\) be
a prescheme. Since coproducts exist in the category of preschemes, for
every set \(E\) we may form the sum of a family \((Y_i)_{i\in E}\) of
preschemes all equal to \(Y\). We denote this prescheme by \(Y\times E\).
\marginpar{\raggedright\footnotesize French source p.~110}%
It is characterized by the formula
\begin{equation}
\label{eq:V.1.*}
  \operatorname{Hom}(Y\times E,Z)
  =\operatorname{Hom}\bigl(E,\operatorname{Hom}(Y,Z)\bigr).
\tag{*}
\end{equation}
Here the second \(\operatorname{Hom}\) on the right means, of course,
the set of maps from \(E\) to the set \(\operatorname{Hom}(Y,Z)\).
There is a canonical morphism
\[
  Y\times E\longrightarrow Y,
\]
making \(Y\times E\) a prescheme over \(Y\). Since fiber products
commute with coproducts in the category of preschemes, if \(Y\) is a
\(Z\)-prescheme and \(Z'\to Z\) is a base change, then
\[
  (Y\times E)\times_ZZ'=(Y\times_ZZ')\times E.
\]
This formula is especially useful when \(Z=Y\). The definition also
immediately gives
\[
  (Y\times E)\times F
  =Y\times(E\times F)
  =(Y\times E)\times_Y(Y\times F),
\]
the last equality also following from the compatibility with fiber
products just noted.

For fixed \(Y\), the construction \(E\mapsto Y\times E\) is a functor
from sets to preschemes over \(Y\). By the preceding formula it
commutes with finite products. Thus, for example, every ordinary group
\(G\) gives rise to a group prescheme \(Y\times G\) over \(Y\), which
is finite over \(Y\) when \(G\) is finite. More generally the functor
is ``left exact,'' although we shall not use this here. It plainly
commutes with coproducts and is also ``right exact,'' as follows at
once from the defining formula~\eqref{eq:V.1.*}. In particular, if a
finite group \(G\) acts on the right on a set \(E\), then it acts on
the right on \(Y\times E\), and
\[
  (Y\times E)/G=Y\times(E/G).
\]
The quotient on the left in fact satisfies the conditions of
Proposition~\ref{V.1.3}, as is immediate.
\end{filecontents*}
% Original input: drafts/SGA1_V_2_English.texfrag
\begin{filecontents*}[overwrite,noheader]{sga1-cumulative-input-090.tex}
\section{Decomposition and inertia groups. The \'etale case}
\label{V.2}
Let $G$
\ifthenelse{\boolean{orig}}{\ignorespaces}{be a}
% correction: insert `a'.
finite group acting on the right on the prescheme $X$. If $x\in X$, the
% page: 111
\marginpar{111}
%
stabilizer $G_d(x)$ of~$x$ is called the \emph{decomposition group
of~$x$}.
%
\index{decomposition (group)|hyperpage}%
\index{decomposition group|hyperpage}%
%
This group acts canonically (on the left) on the residue field
\ifthenelse{\boolean{orig}}
{$\kappa(x)$}
{$\kres(x)$}, and the set of elements of~$G_d(x)$ which act trivially
is called the \emph{inertia group}
%
\index{inertia group|hyperpage}%
\index{inertia (group)|hyperpage}%
%
of~$x$, denoted~$G_i(x)$.

Suppose that $G$ acts admissibly on~$X$ and that $Y$ is a prescheme
over a prescheme~$Z$. Fix $z\in Z$, and an algebraically closed
extension $\Omega$
\ifthenelse{\boolean{orig}}
{of~$\kappa(z)$ having transcendence degree greater than that of the
$\kappa(x)/\kappa(z)$}
{of~$\kres(z)$ having transcendence degree greater than that of the
$\kres(x)/\kres(z)$}, where $x$ is a point of~$X$ above~$z$. We may
regard $\Spec(\Omega)$ as a $Z$-scheme, and the $\Omega$-valued points
of~$X$ correspond to the homomorphisms
\ifthenelse{\boolean{orig}}
{of~$\kappa(z)$-algebras $\kappa(x)\to\Omega$}
{of~$\kres(z)$-algebras $\kres(x)\to\Omega$}, where $x$ is a point
of~$X$ above~$z$. Since $\Omega$ was chosen sufficiently large, every
point $x$ of~$X$ above~$z$ is the underlying point of an
$\Omega$-valued point of~$X$. If $X(\Omega)$ and $Y(\Omega)$ denote,
respectively, the sets of $\Omega$-valued points of~$X$ and~$Y$, there
is a natural map
$$
X(\Omega)\to Y(\Omega),
$$
and, moreover, $G$ acts on~$X(\Omega)$ and the preceding map is
$G$-invariant. With this understood, conclusions (ii) and (iii)
of~\Ref{V.1.3} may also be interpreted as follows: \emph{the preceding
map is surjective and identifies $Y(\Omega)$ with the quotient
$X(\Omega)/G$. Moreover, if $x$ is the underlying point of
$a\in X(\Omega)$, then the stabilizer of~$a$ in~$G$ is precisely the
inertia group $G_i(x)$}. All this is in fact true without assuming that
$\Omega$ is \og sufficiently large\fg; the latter hypothesis serves
only to ensure that the inertia group of every point of~$X$ above~$z$
can be characterized as a \og geometric\fg{} stabilizer. We immediately
deduce, for example:
\begin{proposition}
\label{V.2.1}
Make a base extension $Z'\to Z$, and hence put
$X'=X\times_ZZ'$. Let $x'$ be a point of~$X'$ and $x$ its image in
$X$. Then $G_i(x)=G_i(x')$.
\end{proposition}

In the preceding considerations, it is enough to take for $\Omega$ a
sufficiently large extension
\ifthenelse{\boolean{orig}}
{of~$\kappa(z')$}
{of~$\kres(z')$}
(where $z,z'$ are the images of~$x,x'$ in $Z,Z'$, respectively).

\begin{proposition}
\label{V.2.2}
Under the conditions of~\Ref{V.1.3}, suppose that $Y$ is locally
Noetherian and that $X$ is finite over~$Y$. Let $H$ be a subgroup
of~$G$, put $X'=X/H$ (\cf \Ref{cor:V.1.7}), and let
% page: 112
\marginpar{112}
%
$x\in X$, with image $x'$ in $X'$ and image $y$ in $Y$.
\begin{enumerate}
\item[(i)] If $H\supset G_d(x)$, then the homomorphism
$\cal{O}_y\to\cal{O}_{x'}$ induces an isomorphism on completions.
\item[(ii)] If $H\supset G_i(x)$, then the homomorphism
$\cal{O}_y\to\cal{O}_{x'}$ is \'etale, \ie $X'$ is
\'etale over~$Y$ at $x'$.
\end{enumerate}
\end{proposition}
Let $Y_1=\Spec(\widehat{\cal{O}_y})$, and make the base change
$Y_1\to Y$. We obtain $X_1=X\times_Y Y_1$, finite over~$Y_1$, on
which $G$ acts, with quotient $Y_1$ by~\Ref{V.1.9}. Let~$y_1$ be the
unique point of~$Y_1$ above~$y$. Since
\ifthenelse{\boolean{orig}}
{$\kappa(y)=\kappa(y_1)$}
{$\kres(y)=\kres(y_1)$},
the fiber of~$X$ at~$y$ is isomorphic to that of~$X_1$ at $y_1$;
hence there is a unique point $x_1$ of~$X_1$ above~$x$. Moreover,
by~\Ref{V.1.9}, we have $X_1/H=X'_1=X'\times_Y Y_1$. Let $x'_1$ be
the image of~$x_1$ in $X'_1$. It lies above~$x'$, and one checks
easily (since $X'$ is of finite type over~$Y$) that the homomorphism
$\cal{O}_{x'}\to\cal{O}_{x'_1}$ induces an isomorphism on completions.
We are therefore reduced to the case where $Y$ is the spectrum of a
complete local ring~$B$, and hence $X$ is the spectrum of a finite
$B$-algebra~$A$, which is a product of finitely many local rings~$A_x$
\ifthenelse{\boolean{orig}}{corresponding}{corresponding}
% correction: corresponding
to the points $x_i$ of~$X$ above~$Y$. If $A_0$ corresponds to
$x=x_0$, then $A$ identifies with the ring
$\Hom_{G_d}(G,A_0)$ of functions $f\colon G\to A_0$ such that
$f(st)=sf(t)$ for $s\in G_d$, the operations of~$G$ on these functions
being defined by $(uf)(t)=f(tu)$. Thus, if $H$ is any subgroup of~$G$,
then $A^H$ is the ring of functions $f\colon G\to A_0$ such that
$$
f(stu)=sf(t)\qquad\text{for}\ s\in G_d,\ u\in H
$$
It is therefore a semilocal ring whose local components correspond to
the double cosets $G_daH$ in $G$; to the double coset defined by
$a\in G$ corresponds (by means of the map $f\mto f(a)$) the subring
$A_0^{H(a)}$ of~$A_0$, where $H(a)=G_d\cap aHa^{-1}$. Moreover, the
local component of $A^H$ corresponding to the image $x'$ of~$x$ is
also the one corresponding to the double coset $G_dH$ of the identity
element; its local component is therefore $A_0^{G_d\cap H}$. Thus, if
\ifthenelse{\boolean{orig}}{$G_d\in H$,}{$G_d\subset H$,}
% correction: $G_d\subset H$,
we obtain $A_0^{G_d}=A^G=B$, which proves (i). To prove (ii), we may,
by passing to a suitable finite flat extension of~$A$ and using
\Ref{V.2.1}, reduce to the case where the residue-field extension
\ifthenelse{\boolean{orig}}
{$\kappa(x)/\kappa(y)$}
{$\kres(x)/\kres(y)$}
is trivial. But then $G_i(x)=G_d(x)$, and we are reduced to the
preceding case.

\begin{corollaire}
\label{V.2.3}
Under
% page: 113
\marginpar{113}
%
the conditions of~\Ref{V.2.2}, suppose that $G_i(x)=(e)$. Then $X$
is \'etale over~$Y$ at $x$. Consequently, if $G_i(x)=(e)$
for every $x\in X$, then $X\to Y$ is an \'etale morphism.
\end{corollaire}

There is a partial converse:

\begin{corollaire}
\label{V.2.4} Suppose that $X$ is connected and that the action of
$G$ on~$X$ is faithful. The morphism $p\colon X\to Y=X/G$ is
\'etale if and only if the inertia groups of the points
of~$X$ are reduced to the identity element. If this is so, $G$
identifies with the group of all $Y$-automorphisms of the
$Y$-scheme~$X$.
\end{corollaire}

In view of~\Ref{V.2.3}, we may suppose that $X$ is \'etale
over~$Y$. But if an $s\in G$ belongs to some $G_i(x)$, it then follows
from \ifthenelse{\boolean{orig}}{I~}{}\Ref{I.5.4} that $s$ acts trivially on~\ifthenelse{\boolean{orig}}{$G$}{$X$}, and hence is the
identity element since $G$ is faithful; this proves the first
assertion. Let $u$ be a $Y$-automorphism of~$X$, and let $x\in X$.
By Proposition~\Ref{V.1.3}, there exists an $s\in G$ such that
\ifthenelse{\boolean{orig}}{$s(x)=u(x)$}{$x'=s(x)=u(x)$} and inducing the same residue homomorphism
\ifthenelse{\boolean{orig}}
{$\kappa(x)\to\kappa(x')$}
{$\kres(x')\to\kres(x)$}
as $u$. By \loccit, we therefore have $s=u$, which completes the proof.

\begin{remarque}
\label{V.2.5}
The hypothesis that $G$ acts faithfully is clearly not superfluous in
Corollary~\Ref{V.2.4}. The same is true of the hypothesis that $X$ is
connected, as is seen, for example, by taking $X=Y\times E$, where
$E$ is a finite set, and taking $G$ to be the group of permutations
of~$E$: $G$ acts with considerable inertia, yet
$(Y\times E)/G=Y\times(E/G)=Y$, and $X$ is \'etale over~$Y$.
Taking for $G$ a group strictly smaller than the symmetric group of~$E$
but acting transitively on~$E$, one sees that there will also be
$Y$-automorphisms of~$X$ which do not come from~$G$.
\end{remarque}

The standard example of a group $G$ acting without inertia is that of
$Y\times G$, on which $G$ acts through its right-translation action on
the factor~$G$: a $Y$-prescheme $X$ with right operator group $G$ is
said to be \emph{trivial} if it is isomorphic to $Y\times G$.

The following considerations are useful for relating preschemes with
finite operator groups to the notion of a principal bundle in a
category (a relation which, although we shall not need it in the rest
of the seminar, is important in other contexts). We fix a base
prescheme $Y$ and work in the category of $Y$-preschemes. If $G$ is a
% page: 114
\marginpar{114}
%
finite group, we shall write $G_Y=Y\times G$ for short; this is thus a
finite group scheme over~$Y$ (\cf \numero 1), and if $X$ is a
$Y$-prescheme, then
$$
X \times_Y G_Y = X \times G
$$
(same reference). Giving a $Y$-morphism $X\times_Y G_Y\to X$ is
therefore equivalent to giving a $Y$-morphism $X\times G\to X$, \ie
to giving, for each $g\in G$, a $Y$-morphism $T_g\colon X\to X$.
One checks at once that the data of the $T_g$ define on~$X$ the
structure of a prescheme with right operator group $G$ (\ie
$T_{gg'}=T_{g'}T_g$, \ifthenelse{\boolean{orig}}{$T_e=\id_x$}{$T_e=\id_X$}) if and only if the corresponding
$Y$-morphism $X\times_YG_Y\to X$ defines on~$X$ the structure of a
$Y$-prescheme with operator $Y$-group scheme (in the general sense of
objects with a $\cal{C}$-operator group in a category $\cal{C}$).
Suppose this is so. Recall that $X$ is said to be \emph{formally
principal homogeneous}
%
\index{formally principal homogeneous (prescheme)|hyperpage}%
%
under~$G_Y$\kern1pt\footnote{One now says instead that $X$ is a
pseudo-torsor under $G_Y$.} if the canonical morphism
$$
X \times_Y G_Y \to X \times_Y X
$$
whose components are respectively $\pr_1$ and the multiplication
morphism $\pi\colon X\times_YG_Y\to X$, is an isomorphism. In the
present case, identifying the first member with $X\times G$, the
morphism in question is the one which associates to each $g\in G$ the
morphism
$$
(\id_X,T_g)=(\id_X \times_Y T_g)\Delta_{X/Y} \colon X \to X \times_Y X
$$
Consequently, to say that $X$ is formally principal homogeneous under
$G_Y$ also means that $X\times_YX$ is isomorphic to the disjoint sum
of the transforms of the diagonal by the elements $(e,g)$ of
$G\times G$ (acting on~$X\times_YX$ in the evident way), where $e$
denotes the identity element of~$G$. If one does not wish to distinguish
left from right, and wants a formula which remains applicable to a
product of more than two factors identical to~$X$, the condition may
be formulated by saying that the canonical morphism
$$
X \times_G (G \times G) \to X \times_Y X
$$
obtained by assigning to the pair $(g,g')$ the morphism
$$
(T_g,T_{g'})=(T_g \times_Y T_{g'})\Delta_{X/Y} \colon X \to X \times_Y
X
$$
and making $G$ act on the left on~$G\times G$ by the diagonal
homomorphism,
$$
s(g,g')=(sg,sg'),
$$
is an \emph{isomorphism}.

The
% page: 115
\marginpar{115}
%
notion of a \emph{principal homogeneous space}
%
\index{principal homogeneous (prescheme, bundle)|hyperpage}%
%
is obtained from that of a formally principal homogeneous space by
adding a further axiom ensuring that the \og quotient\fg{} of~$X$ by
$G_Y$ exists and is precisely the right unit object of the category,
here~$Y$. This axiom may vary with the context, and is often most
conveniently stated explicitly (in the yoga of \og descent\fg) by
requiring that the object with operators become \og trivial\fg, \ie
isomorphic to the product $X\times_YG_Y$ (in the present case
$X\times G$), after a suitable base change of a specified type (in
such a way, in practice, as to permit the descent technique; \cf
Grothendieck, Technique de descente et th\'eor\`emes d'existence en
G\'eom\'etrie Alg\'ebrique, S\'em. Bourbaki \No 190, pages 26 to
28)\footnote{\Cf Exp\ptbl \Ref{VIII} for the theory of flat descent.}.
In this vein, we record here the characterization of principal
homogeneous bundles with group $G$ (in the sense of \loccit):

\begin{proposition}
\label{V.2.6}
Let $Y$ be a locally Noetherian prescheme and let $X$ be a
$Y$-prescheme with a finite group $G$ of operators acting on the
right. The following conditions are equivalent:
\begin{enumerate}
\item[(i)] $X$ is finite over~$Y$, $Y=X/G$, and the inertia groups of
the points of~$X$ are reduced to the identity.
\item[(ii)] There exists a faithfully flat and quasi-compact base
change $Y_1\to Y$ such that $X_1=X\times_YY_1$ is a trivial
$Y_1$-prescheme with operators, \ie isomorphic to $Y_1\times G$.
\item[(ii~bis)] As in (ii), but with $Y_1\to Y$ finite,
\'etale, and surjective.
\item[(iii)] $X$ is formally principal homogeneous under $G_Y$, and
is faithfully flat and quasi-compact over~$Y$.
\end{enumerate}
\end{proposition}
\textit{Proof} (i)$\To$(ii~bis). Take $Y_1=X$, noting that
$X\to Y$ is indeed finite, \'etale by~\Ref{V.2.3}, and
surjective. Let us show that $X_1$ is then trivial over~$Y_1$, which
will follow from the

\begin{corollaire}
\label{V.2.7} If \textup{(i)} holds and $X$ admits a section over~$Y$,
then $X$ is a trivial space with operators.
\end{corollaire}

Indeed, this section makes it possible to define a $G$-morphism
$\ifthenelse{\boolean{orig}}{X}{Y}\times G\to X$, which is surjective because $G$ is transitive on the
fibers of~$X$, injective because $G$ acts without inertia, and finally
a local isomorphism by \ifthenelse{\boolean{orig}}{I~}{}\Ref{I.5.3}, since $X$ is \'etale
over~$Y$. It is therefore an isomorphism.

(ii~bis)
% page: 116
\marginpar{116}
%
trivially implies (ii), which implies (i), since the ingredients of
(i) are \og invariant\fg{} under faithfully flat quasi-compact base
extension (\cf the Bourbaki Seminar cited above for \og finite\fg;
for inertia groups apply~\Ref{V.2.1}, and for $Y=X/G$ use a converse
to~\Ref{V.1.9} in the case of a \emph{faithfully flat} base change,
which we had forgotten to state explicitly).

In proving (i)$\To$(ii~bis), we proved (i)$\To$(iii) along the way.
Finally, (iii)$\To$(ii), because the first hypothesis in (iii) means
precisely that $X$ becomes trivial after making the base change
$Y_1=X$; hence (ii), since $X$ is faithfully flat and quasi-compact
over~$Y$.

\begin{definition}
\label{V.2.8} A $Y$-prescheme $X$ with right operator group $G$
satisfying the equivalent conditions of~\Ref{V.2.6} is called a
\emph{principal covering of~$Y$, with Galois group~$G$}.
%
\index{Galois group of a principal covering|hyperpage}%
\index{principal (covering)|hyperpage}%
\index{principal covering|hyperpage}%
%
\end{definition}
\end{filecontents*}
% Original input: drafts/SGA1_V_3_English.texfrag
\begin{filecontents*}[overwrite,noheader]{sga1-cumulative-input-091.tex}
\section{Automorphisms and morphisms of étale coverings}
\label{V.3}

\begin{proposition}
\label{V.3.1}
Let \(X\) be separated, étale, and of finite type over a locally
noetherian prescheme~\(Y\), and let \(G\) be a finite group acting
on~\(X\) by \(Y\)-automorphisms. Then the action of \(G\) is admissible,
and the quotient prescheme \(X/G\) is étale over~\(Y\).
\end{proposition}

Although \(X\) is not assumed finite over~\(Y\), it is quasi-projective
over~\(Y\), so the existence of \(X/G\) follows from~\Ref{V.1.8}.
Let us first prove:

\begin{corollary}
\label{V.3.2}
The morphism \(X\to X/G\) is étale.
\end{corollary}

We may clearly suppose first that \(G\) acts transitively on the set of
connected components of~\(X\), and then, by considering the stabilizer
of one connected component, that \(X\) itself is connected. Finally, we
may suppose that \(G\) acts faithfully. But then, just as in~\Ref{V.2.4},
one sees that \(G\) acts without inertia; hence \Ref{V.2.3} implies that
\(X\to X/G\) is étale. The conclusion follows from:

\begin{lemma}[a regret concerning Exposé I]
\label{V.3.3}
Let \(X\to X'\to Y\) be morphisms of finite type, let \(x\) be a point
of~\(X\), and let \(x'\) and \(y\) be its images. Suppose that \(Y\) is
locally noetherian. If two of the three morphisms in question are étale
at the indicated points, then so is the third.
\end{lemma}

It remains only to consider the case in which \(X\to X'\) and
\(X\to Y\) are étale at \(x\), and to prove that \(X'\to Y\) is étale
at \(x'\) (which is the case needed for~\Ref{V.3.1}).
% page: 117
\marginpar{117}
%
After making a suitable flat extension of the base~\(Y\), we reduce to
the case in which the residue-field extension
\ifthenelse{\boolean{orig}}
{\(\kappa(x)/\kappa(y)\)}
{\(\kres(x)/\kres(y)\)}
is trivial. Consider the homomorphisms
\(\cal{O}_y\to\cal{O}_{x'}\to\cal{O}_x\), together with the
homomorphisms obtained by passing to completions. The hypothesis says
that
\[
  \widehat{\cal{O}}_y\longrightarrow\widehat{\cal{O}}_x
  \quad\hbox{and}\quad
  \widehat{\cal{O}}_{x'}\longrightarrow\widehat{\cal{O}}_x
\]
are isomorphisms. It follows at once that
\(\widehat{\cal{O}}_y\to\widehat{\cal{O}}_{x'}\) is also an
isomorphism, which proves the lemma.

\begin{corollary}
\label{V.3.4}
If \(X\) is finite and étale over~\(Y\), then \(X/G\) is finite and
étale over~\(Y\).
\end{corollary}

\begin{proposition}
\label{V.3.5}
Let \(X\) and \(X'\) be two étale coverings of~\(Y\). Then every
\(Y\)-morphism \(f\colon X\to X'\) factors as the composite of a
surjective étale morphism \(X\to X''\) and the canonical immersion
\(X''\to X'\) of a subset \(X''\) of~\(X'\) that is both open and
closed.
\end{proposition}

It is known (\ifthenelse{\boolean{orig}}{I~}{}\Ref{I.4.8}) that \(f\) is étale, and hence an open
morphism. On the other hand, since \(X\) is finite over~\(Y\), the
morphism \(f\) is closed. Thus \(f(X)=X''\) is both open and closed
in~\(X'\), and the proof is complete. (N.B. it would have been enough
for \(X'\), instead of being an étale covering, to be unramified
over~\(Y\).)

\begin{corollary}
\label{V.3.6}
With the preceding notation, \(X\to X''\) is a strict epimorphism in
the category of preschemes, and \(X''\to X'\) is a monomorphism (indeed,
a strict monomorphism) in the category of preschemes.
\end{corollary}

By definition, the first assertion means that the sequence of morphisms
\[
\xymatrix@C=.7cm{X \times_{X''} X \ar@<2pt>[r]^-{\pr_1}
\ar@<-2pt>[r]_-{\pr_2} & X \ar[r] & X''}
\]
is exact. This follows from the fact that \(X\to X''\) is finite and
faithfully flat, as is easily seen (\cf Grothendieck, \loccit). The
dual assertion for \(X''\to X'\) is still more immediate.

Corollary~\Ref{V.3.6} will be useful in the theory of the fundamental
group in the next section. Those who dislike the notion of strict
epimorphism may replace Corollary~\Ref{V.3.6} by whatever variant they
prefer. Let us merely take this opportunity to point out that a
factorization
% page: 118
\marginpar{118}
%
\(f=f'f''\), with \(f''\) a strict epimorphism and \(f'\) a
monomorphism, is necessarily unique up to unique isomorphism (in any
category). At the same time, however, there may exist a factorization
\(f=f_1f_2\) having the dual properties: \(f_2\) is an epimorphism and
\(f_1\) a strict monomorphism (again unique up to unique isomorphism),
which is not isomorphic to the preceding one. It is enough, for
example, to take the category of topological vector spaces (Hausdorff,
if desired) and a morphism \(u\colon X\to X'\) such that \(u(X)\) is
not closed.

\begin{proposition}
\label{V.3.7}
Let \(Y\) be a connected locally noetherian prescheme, let \(y\) be a
point of~\(Y\), and let \(\mathit{\Omega}\) be an algebraically closed
extension
\ifthenelse{\boolean{orig}}
{of~\(\kappa(y)\)}
{of~\(\kres(y)\)}.
For every \(X\) over~\(Y\), denote by \(X(\Omega)\) the set of points
of~\(X\) with values in~\(\Omega\). Let \(X\) and \(X'\) be étale
coverings of~\(Y\), and let \(u\colon X\to X'\) be a \(Y\)-morphism
such that the corresponding map
\[
  X(\Omega)\longrightarrow X'(\Omega)
\]
is a bijection. Then \(u\) is an isomorphism.
\end{proposition}

We reduce immediately to the case in which \(X'\) is connected. Since
\(X\to X'\) is finite and étale, the geometric number of points in a
fiber of \(X\to X'\) is constant, and it is equal to \(1\) if and only
if the morphism in question is an isomorphism. But this number is also
the number of elements in a fiber of
\(X(\Omega)\to X'(\Omega)\), and the conclusion follows.
\ifthenelse{\boolean{orig}}{}
{\enlargethispage{.5cm}}%

\end{filecontents*}
% Original input: drafts/SGA1_V_4_English.texfrag
\begin{filecontents*}[overwrite,noheader]{sga1-cumulative-input-092.tex}
% Source: smf_doc-math_3_01.tex, lines 8280--8762 inclusive (SGA 1, Expose V).

\section{Axiomatic conditions for a Galois theory}
\label{V.4}
Let $\cal{C}$ be a category and $F$ a covariant functor from
$\cal{C}$ to the category of finite sets. Assume that the following
conditions are satisfied:

\begin{enumerateb}

\item[(G 1)] $\cal{C}$ has a terminal object\footnote{Recall that an
object $e$ of~$\cal{C}$ is called a \emph{terminal object} if, for every
$X$ in $\cal{C}$, $\Hom(X,e)$ has exactly one element. An \emph{initial
object} of~$\cal{C}$ is defined dually.}, and the fiber product of two
objects over a third object in $\cal{C}$ exists (this axiom may also be
stated by saying that finite inverse limits exist in $\cal{C}$).

\item[(G 2)] Finite coproducts exist in $\cal{C}$ (and hence so does an
initial object $\emptyset_{\cal{C}}$, playing the role of the empty set),
as does the quotient of an object of~$\cal{C}$ by a finite group of
automorphisms.

\item[(G 3)] Let $u \colon X \to Y$ be a morphism in $\cal{C}$. Then
$u$ factors as a composite $X\lto{u'} Y'\lto{u''} Y$, where $u'$ is a
\emph{strict} epimorphism and $u''$ is a monomorphism which is an
isomorphism onto a direct summand of~$Y$.

\item[(G 4)]
% page: 119
\marginpar{119}
%
The functor $F$ is left exact (\ie it sends a terminal object to a
terminal object and commutes with
\ifthenelse{\boolean{orig}}{the fiber product}{fiber products}).

\item[(G 5)] $F$ commutes with finite coproducts, sends strict
epimorphisms to epimorphisms, and commutes with passage to the quotient
by a finite group of automorphisms.

\item[(G 6)] Let $u \colon X \to Y$ be a morphism in $\cal{C}$ such
that $F(u)$ is an isomorphism. Then $u$ is an isomorphism.
\end{enumerateb}

Our aim is to construct a topological group~$\pi$, an inverse limit of
finite groups, and an equivalence between the category $\cal{C}$ and
the category $\cal{C}(\pi)$
\label{indnot:eb}\oldindexnot{$\cal{C}(\pi)$|hyperpage}%
%
\emph{of finite sets on which} $\pi$ \emph{acts continuously} (\ie in
such a way that the stabilizer of a point is an open subgroup, or,
equivalently, that there is a discrete quotient group which already
acts on the set in question), the equivalence constructed sending the
given functor $F$ to the evident inclusion functor from~$\cal{C}(\pi)$
into the category of finite sets. We note at once that the category
$\cal{C}(\pi)$ constructed from a topological group~$\pi$, together
with the preceding inclusion functor, does indeed satisfy conditions
(G 1) through (G 6).

We proceed in several steps.
\begin{enumerateb}
\item[a)] \emph{Let} $u \colon X \to Y$ \emph{be a morphism in}
$\cal{C}$. \emph{Then} $u$ \emph{is a monomorphism if and only if}
$F(u)$ \emph{is a monomorphism}. (This uses (G 1), (G 4), and (G 6).)

Indeed, to say that $u$ is a monomorphism means that the projection
$\pr_1 \colon X \times_Y X \to X$ is an isomorphism.

\item[b)] \emph{Every object} $X$ \emph{of} $\cal{C}$ \emph{is
Artinian}.

Indeed, if $X' \to X'' \to X$ are monomorphisms such that $F(X')$ and
$F(X'')$ have the same image in $F(X)$, then by a), $F(X') \to
F(X'')$ is an isomorphism, and hence $X' \to X''$ is an isomorphism
by (G 6).

\item[c)] \emph{The functor} $F$ \emph{is strictly
pro-representable}. (\cf Grothendieck, Technique de descente et
th\'eor\`emes d'existence en G\'eom\'etrie Alg\'ebrique, II,
S\'eminaire Bourbaki 195, February 1960.)

Indeed, by \loccit prop\ptbl \Ref{V.3.1}, this follows from b) and
(G 4). Thus one can find an inverse system, indexed by a directed
ordered set~$I$,
$$
P=(P_i)_{i \in I}
$$
% page: 120
\marginpar{120}%
%
in $\cal{C}$, regarded as a pro-object of~$\cal{C}$, and a functorial
isomorphism
\begin{equation*}
\label{eq:V.4.*} \tag{$*$} F(X)=\Hom_{\Pro(\cal{C})}(P,X)=\varinjlim_{i} \Hom_{\cal{C}}(P_i,X).
\end{equation*}
This isomorphism is realized by an element
$$
\varphi \in \varprojlim_{i} F(P_i)=F(P).
$$
One may moreover assume that the transition homomorphisms
$\varphi_{ji} \colon P_i \to P_j$ ($i \geq j$) are
\emph{epimorphisms}, and that \emph{every epimorphism} $P_i \to P'$
is equivalent to an epimorphism $P_i \to P_j$ for a suitable $j \leq
i$ (which determines the inverse system~$P$ in an essentially unique
way).

An object $X \in \cal{C}$ is called \emph{connected}
%
\index{connected (object of a category)|hyperpage}%
%
if it is not isomorphic to the coproduct of two objects of~$\cal{C}$
neither of which is isomorphic to the initial object
$\emptyset_{\cal{C}}$.

\item[d)] \emph{The} $P_i$ \emph{are connected and are not isomorphic
to} $\emptyset_{\cal{C}}$.

If $X$ is an initial object, then $F(X)=\emptyset$ by (G 5), applied
to the coproduct of an empty family, and the converse follows from
(G 6). Thus if $X'$ is an object of~$\cal{C}$ which is not initial,
\ie such that $F(X') \neq \emptyset$, there is no morphism from~$X'$
to~$X$. Consequently, if some $P_i$ is initial, then $i$ is a greatest
element of the directed ordered index set $I$, and formula
\eqref{eq:V.4.*} would say that $F(X)=\Hom(P_i,X)$ is a singleton for
every $X$, which is absurd since
$F(\emptyset_{\cal{C}})=\emptyset$. The~$P_i$ are therefore not
isomorphic to $\emptyset_{\cal{C}}$.

Suppose that $P_i=A \amalg B$. By (G 5) this gives
$F(P_i)=F(A) \amalg F(B)$. In particular, the element $a_i$ of
$F(P_i)$ corresponding under \eqref{eq:V.4.*} to the identity
homomorphism $P_i \to P_i$ belongs to $F(A) \amalg F(B)$, say to
$F(A)$. This means that there is a $j \geq i$ such that
$\varphi_{ij} \colon P_j \to P_i$ factors as
$P_j \to A \to P_i=A \amalg B$, where the second arrow is the
canonical morphism. Hence $F(P_j) \to F(P_i)$ factors as
$F(P_j) \to F(A) \to F(P_i)=F(A) \amalg F(B)$; since $F(P_j) \to
F(P_i)$ is surjective by (G 5), it follows that $F(B)=\emptyset$, and
hence that $B$ is isomorphic to $\emptyset_{\cal{C}}$.

\item[e)]
% page: 121
\marginpar{121}
%
\emph{Every morphism} $u \colon X \to Y$ \emph{in} $\cal{C}$,
\emph{with} $X$ \emph{not isomorphic to} $\emptyset_{\cal{C}}$
\emph{and} $Y$ \emph{connected, is a strict epimorphism. Every
endomorphism of a connected object is an automorphism}.

Consider the factorization of~$u$ in (G 3). Since $X \neq
\emptyset_{\cal{C}}$, it follows from (G 6) that $F(X) \neq
\emptyset$, hence $F(Y') \neq \emptyset$, hence $Y' \neq
\emptyset_{\cal{C}}$. Since $Y$ is connected, $Y'$ is therefore
identified with~$Y$, and thus $u$ is a strict epimorphism. Suppose
that $u$ is an endomorphism of the connected object $X$; let us show
that it is an automorphism. We may assume that $X$ is not isomorphic
to $\emptyset_{\cal{C}}$, so by what precedes $u$ is a strict
epimorphism. Hence $F(u)$ is an epimorphism by (G 5), and since $F(X)$
is a finite set, $F(u)$ is bijective. Thus $u$ is an automorphism by
(G 6).

In particular, \emph{every endomorphism of a} $P_i$ \emph{is an
automorphism}.

\item[f)] \emph{The following conditions on a} $P_i$ \emph{are
equivalent}: (i) The natural injective map
$\Hom(P_i,P_i) \to \Hom(P,P_i) \simeq F(P_i)$ is also surjective,
\ie for every $u \colon P \to P_i$ there is a $v \colon P_i \to
P_i$ such that $u=v\varphi_i$ (where $\varphi_i$ is the canonical
homomorphism $P \to P_i$). (ii) The group $\Aut(P_i)$ acts
transitively on~$F(P_i)$. (iii)~The group $\Aut(P_i)$ acts simply
transitively on~$F(P_i)$.

Upon identifying $\Hom(P,P_i)$ with $F(P_i)$, the map in (i) is
precisely $v \mto F(v)(\varphi_i)$. The equivalence of the three
conditions then follows from the fact that
$\Hom(P_i,P_i)=\Aut(P_i)$ and that the preceding map is already
injective.

A $P_i$ satisfying the equivalent conditions
\ifthenelse{\boolean{orig}}
{(i) (ii) (iii)}
{(i), (ii), (iii)}
of~f) is called \emph{Galois}.
%
\index{Galois (object)|hyperpage}%
%

\item[g)] \emph{For every} $X$ \emph{in} $\cal{C}$, \emph{there is
a Galois} $P_i$ \emph{such that every} $u \in \Hom(P,X)$
\emph{factors as} $P \lto{\varphi_i} P_i \longrightarrow X$.

Let $J=\Hom(P,X)=F(X)$. This is a finite set, so there is a $P_j$
such that every $u \colon P \to X$ factors as
\ifthenelse{\boolean{orig}}{$P\lto{j}P_j\longrightarrow X$}{$P\lto{\varphi_j}P_j\longrightarrow X$}, or, equivalently, such that the
natural morphism
$$
P \to X^J \qquad (J=\Hom(P,X))
$$
factors as
$$
P \lto{\varphi_j} P_j \to X^J.
$$
By
% page: 122
\marginpar{122}
%
(G~3), the morphism $P_j \to X^J$ factors as the composite of a
strict epimorphism and a monomorphism, where the former may be taken
to be of the form $\varphi_{ij} \colon P_j \to P_i$. We are thus
reduced to proving that $P_i$ is Galois. Let $k$ be an index
$\geq j$ such that every morphism $P \to P_i$ factors through
$P \lto{\varphi_k} P_k \longrightarrow P_i$. Note that the natural
morphism $P_k \to X^J$ still factors as the composite
$$
P_k \lto{\varphi_{ik}} P_i \lto{U} X^J,
$$
where the first \ifthenelse{\boolean{orig}}{arrow}{arrow} is a strict
epimorphism by~e), and the second is a monomorphism. We wish to show
that, for a given morphism $\psi \colon P_k \to P_i$, there is an
endomorphism $v$ of~$P_i$ such that
$\psi = v \varphi_{ik}$. For every $u \in \Hom(P_i,X)$, consider
$u \psi \in \Hom(P_k,X)$. It is of the form $u'\varphi_{ik}$ for a
uniquely determined $u' \in \Hom(P_i,X)$. The map $u \mto u'$ from
$J$ to $J$ thus defined by $\psi$ is injective, since $\psi$ is an
epimorphism by e); it is therefore bijective because~$J$ is finite.
The bijection $u \mto u'$ from~$J$ to $J$ thus defines an isomorphism
$\alpha \colon X^J \isomto X^J$ making the diagram
$$
\xymatrix{P_k \ar[r]^-{\varphi_{ik}} \ar@{=}[d] & P_i \ar[r]^-U & X^J
\ar[d]^-{\alpha}_-{\simeq} \\ P_k \ar[r]^{\psi} & P_i \ar[r]^-U & X^J}
$$
commutative. By the uniqueness properties of the factorization of a
morphism as a composite of a strict epimorphism and a monomorphism,
it follows (since $\psi$ too is a strict epimorphism by e)) that one
can find a morphism $v \colon P_i \to P_i$ which makes the diagram
commutative, as required.

It follows in particular that \emph{the Galois} $P_i$ \emph{form a
cofinal subsystem of the system} $(P_j)$. Since for a Galois object
$P_i$ one has
$$
\Hom(P,P_i)=\Hom(P_i,P_i)=\Aut(P_i),
$$
passage to the limit therefore gives
$$
\Hom(P,P) = \varprojlim_{i} \Hom(P,P_i) = \varprojlim_{i} \Hom(P_i,P_i) =
\varprojlim_{i} \Aut(P_i),
$$
where
\ifthenelse{\boolean{orig}}
{\ignorespaces}
{\enlargethispage{.5cm}}
% page: 123
\marginpar{123}
%
the inverse limit is taken over the Galois $P_i$. Moreover, under the
identification $\Hom(P,P_i)=F(P_i)$, and taking account of the fact
that $F$ sends epimorphisms to epimorphisms, one sees that the
transition homomorphisms in the preceding inverse system are
surjective. From all this we conclude:

\item[h)] \emph{One has}
$$
\Hom(P,P) = \Aut(P) = \varprojlim_{i} F(P_i) = \varprojlim_{i} \Aut(P_i),
$$
\emph{where the inverse limit is taken over the Galois} $P_i$.

In particular, $\Aut(P)$ appears as the inverse limit of an inverse
system of finite groups (with surjective transition homomorphisms); we
endow it with the inverse-limit topology of the discrete topologies.
\emph{We denote by}~$\pi$
\label{indnot:ec}\oldindexnot{$\pi$|hyperpage}%
%
\emph{and call the fundamental group}
%
\index{fundamental group of a Galois category|hyperpage}%
%
(of~$\cal{C}$ equipped with~$F$) \emph{the group opposite
to}~$\Aut(P)$. Thus this group acts \emph{on the right} on~$P$; it is
the inverse limit of finite groups $\pi_i$ acting on the right on the
Galois $P_i$, where $\pi_i$ is the group opposite to $\Aut(P_i)$.

In view of the functorial isomorphism
$$
F(X)=\Hom(P,X)
$$
and the definition of~$\pi$, we therefore see that $\pi$ acts
\emph{on the left} on~$F(X)$, and moreover acts continuously by g)
(for, with the notation of g), it is in fact $\pi_i$ that acts on
$F(X)$). It is immediate that, for every morphism
$u \colon X \to Y$ in $\cal{C}$, the morphism
$F(u) \colon F(X) \to F(Y)$ is compatible with the actions of~$\pi$.
\emph{Henceforth one may therefore regard} $F$ \emph{as a covariant
functor}
$$
F \colon \cal{C} \to \cal{C}(\pi),
$$
\emph{where} $\cal{C}(\pi)$ \emph{is the category of finite sets on
which} $\pi$ \emph{acts continuously on the left}.

We now define a functor in the opposite direction:
$$
G \colon \cal{C} \from \cal{C}(\pi)
$$
by the formula
$$
G(E)=P \times_{\pi} E,
$$
where $P \times_{\pi} E$ is defined as a solution of the universal
problem summarized by
$$
\Hom_{\cal{C}}(P \times_{\pi} E,X) \isomto \Hom_{\pi}(E,\Hom(P,X))
$$
(where,
% page: 124
\marginpar{124}
%
in the right-hand side, $\Hom(P,X)=F(X)$ is regarded as a set on which
$\pi$ acts on the left). We must prove the existence of the object
$P \times_{\pi} E$.

\item[i)] \emph{Let} $Q$ \emph{be an object of} $\cal{C}$ \emph{on
which a finite group} $G$ \emph{acts on the right, and let} $E$
\emph{be a finite set on which} $G$ \emph{acts on the left. Then}
$Q \times_G E$ \emph{exists, and the canonical map}
$$
F(Q) \times_G E \to F(Q \times_G E)
$$
\emph{is an isomorphism}.

Since finite coproducts exist in $\cal{C}$ by (G 2), and $F$ commutes
with them by (G 5), one is immediately reduced to the case in which
$G$ acts transitively on~$E$: indeed, if the~$E_j$ are the orbits of
$G$ in $E$, then
$$
Q \times_G E = \underset {j}{\amalg}\ Q \times_G E_j.
$$
Let $a \in E$, and let $H$ be its stabilizer. It follows at once from
the definition that $Q \times_G E$ is identified with $Q/H$. Existence
thus follows from (G~2), and the compatibility property for~$F$ follows
from (G~5).

\item[j)] \emph{Let} $E$ \emph{be an object of} $\cal{C}(\pi)$,
\emph{and let} $P_i$ \emph{be Galois and such that} $\pi_i$
\emph{already acts on} $E$. \emph{Then} $P_i \times_{\pi_i} E$
\emph{exists and there is a canonical isomorphism}
$$
E \isomto F(P_i \times_{\pi_i} E).
$$
\emph{If} $j \geq i$ \emph{is such that} $P_j$ \emph{is Galois,
then the canonical homomorphism} $P_j \times_{\pi_j} E \to\allowbreak P_i
\times_{\pi_i} E$ \emph{is an isomorphism}.

The first assertion is a special case of i), taking account of the
fact that $\pi_i$ acts simply transitively on~$F(P_i)$, which is
equipped with the marked point $\varphi_i$; this gives an isomorphism
$F(P_i) \times_{\pi_i} E \simeq E$. For the second assertion one may,
for example, use (G 6).

For every $j$, let $\cal{C}_j$ be the full subcategory of $\cal{C}$
formed by those $X$ for which $\Hom(P_j,X) \to \Hom(P,X) \simeq F(X)$
is bijective. We know from g) that $\cal{C}$ \emph{is the filtered
union of the} $\cal{C}_j$. Thus, for $X \in \cal{C}_j$, one has
\begin{align*}
\Hom_{\pi}(E,\Hom(P,X)) & \simeq \Hom_{\pi}(E,\Hom(P_j,X)) \simeq
\Hom_{\pi_j}(E,\Hom(P_j,X))\\ & \simeq \Hom(P_j \times_{\pi_j} E,X),
\end{align*}
and,
% page: 125
\marginpar{125}
%
in view of the last assertion in j), one obtains an isomorphism
functorial in the object $X$ of~$\cal{C}_j$:
$$
\Hom_{\pi}(E,\Hom(P,X)) \simeq \Hom(P_i \times_{\pi_i} E,X).
$$
Since this is true for every $j$, and since these functorial
isomorphisms, as $j$ varies, induce one another, we conclude:

\item[k)] \emph{Under the conditions of} j), \emph{the composite of
the canonical morphisms}
$$
E \to \Hom(P_i,P_i \times_{\pi_i} E) \to \Hom(P,P_i \times_{\pi_i} E)
$$
\emph{makes} $P_i \times_{\pi_i} E$ \emph{a solution of the universal
problem defining} $P \times_{\pi} E$, \ie \emph{the latter exists and
there is an isomorphism}
$$
P \times_{\pi} E \isomto P_i \times_{\pi_i} E.
$$

This completes the construction of the functor $G(E)$. On the other
hand, there is a functorial homomorphism
$$
\alpha \colon \id_{\cal{C}(\pi)} \to FG,
$$
that is, a homomorphism functorial in the object $E$ of~$\cal{C}(\pi)$:
$$
\alpha(E) \colon E \to FG(E) = F(P \times_{\pi} E),
$$
namely the composite of the canonical morphisms
$$
E \to F(P) \times_{\pi} E \to F(P \times_{\pi} E)
$$
(where the first comes from the marked point
$\varphi \in F(P)$). Combining j) and k), one finds:

\item[l)] \emph{The homomorphism} $\alpha$ \emph{is an isomorphism}.

Similarly, one defines a functorial homomorphism
$$
\beta \colon GF \to \id_{\cal{C}},
$$
that is, a homomorphism functorial in the object $X$ of~$\cal{C}$:
$$
\beta(X) \colon P \times_{\pi} F(X) \to X,
$$
as the morphism
% page: 126
\marginpar{126}
%
associated with the $\pi$-homomorphism
$$
F(X) \to \Hom(P,X)
$$
inverse to the canonical isomorphism $\Hom(P,X) \isomto F(X)$.

\item[m)] \emph{The composites}
\begin{gather*}
F(X) \lto{\alpha(F(X))} FGF(X) \lto{F(\beta(X))} F(X)\\
G(E) \lto{G(\alpha(E))} GFG(E) \lto{\beta(G(E))} G(E)
\end{gather*}
\emph{are the identity isomorphisms}.

A trotting donkey.

In view of l), it follows that:

\item[n)] \emph{The homomorphism} $\beta$ \emph{is an isomorphism}.

We have thus obtained the promised result:
\end{enumerateb}

\begin{theoreme}
\label{V.4.1}
Let $\cal{C}$ be a category satisfying conditions \emph{(G 1)},
\emph{(G 2)}, and \emph{(G~3)} at the beginning of this section, and
let $F$ be a covariant functor from~$\cal{C}$ to the category of finite
sets, satisfying conditions \emph{(G 4)}, \emph{(G 5)}, and
\emph{(G 6)}. Then the preceding canonical constructions define
equivalences of categories $F\colon \cal{C}\to \cal{C}(\pi)$ and
$G\colon \cal{C}(\pi)\to \cal{C}$ which are quasi-inverse to one
another. More precisely, there is a pro-object~$P$ of~$\cal{C}$ and a
functorial isomorphism $F(X) \isomfrom \Hom(P, X)$; $\pi$ is the group
opposite to the automorphism group of~$P$, endowed with a suitable
topology, so that~$\pi$ acts continuously on the sets
$\Hom(P, X)\simeq F(X)$. Finally, $G$ is given by
$G(E)\simeq P\times_{\pi} E$.
\end{theoreme}

\begin{remarques}
\label{V.4.2}
The statement of conditions (G 1) through (G 6) becomes simpler and
more pleasing if (G 2) and (G 5) are replaced respectively by:
\begin{enumerateb}
\item[(G' 2)] Finite direct limits exist in
\ifthenelse{\boolean{orig}}{$C$}{$\cal{C}$}.
% correction: `C' should be `\cal{C}'.
\item[(G' 5)] The functor $F$ is right exact (\ie it commutes with
finite direct limits).
\end{enumerateb}
These conditions appear stronger than (G~2) and (G~5), but it follows
at once from the structure theorem~\Ref{V.4.1} that they are implied
by (G~1) through (G~6). Note, however, that in the cases of interest
to us, checking (G~2) and (G~5) does in fact seem simpler than checking
(G'~2) and (G'~5). I do not know whether, in condition (G~3), the
requirement that $u''$ be an isomorphism onto a direct summand of~$Y$
could be omitted.
\end{remarques}

\end{filecontents*}
% Original input: drafts/SGA1_V_5_English.texfrag
\begin{filecontents*}[overwrite,noheader]{sga1-cumulative-input-093.tex}
\section{Galois categories}
\label{V.5}
% page: 127
\marginpar{127}
%

\begin{definition}
\label{V.5.1}
A \emph{Galois category}
%
\index{Galois category|hyperpage}%
%
is a category $\cal{C}$ equivalent to a category $\cal{C}(\pi)$, where
$\pi$ is a compact group that is a projective limit of finite groups
(i.e.\ is totally disconnected).

For the definition of~$\cal{C}(\pi)$, see the beginning of
\No~\Ref{V.4}. By Theorem~\Ref{V.4.1}, $\cal{C}$ is Galois if and only if
it satisfies conditions (G~1) through (G~3), and there exists a functor
$F$ from
\ifthenelse{\boolean{orig}}{$C$}{$\cal{C}$}
% correction: `$C$' should be `$\cal{C}$'.
to the category of finite sets satisfying conditions (G~4) through
(G~6), that is, one which is \emph{exact} and \emph{conservative}, in
general terminology. Such a functor is called a \emph{fundamental
functor}
%
\index{fundamental functor|hyperpage}%
%
of the Galois category $\cal{C}$\kern1pt\footnote{It seems preferable to
adopt the more expressive term ``fiber functor.''};
%
\index{fiber functor|hyperpage}%
%
it is pro-representable by a pro-object which we denote by~$P_F$. A
pro-object $P$ for which the associated functor $F$ is fundamental is
called a \emph{fundamental pro-object}.
%
\index{fundamental pro-object|hyperpage}%
%
In this way, the category of fundamental functors on~$\cal{C}$ is
anti-equivalent to the category of fundamental pro-objects. If $F$ and
$P$ correspond, the group $\Aut F$ is therefore isomorphic to the
opposite of the group $\Aut P$; hence the group denoted by~$\pi$ in the
preceding section is none other than \ifthenelse{\boolean{orig}}{$\Aut P$}{the opposite of~$\Aut P$, equivalently~$\Aut F$}. Recall that in the
preceding section, starting from a \emph{given} fundamental functor~$F$,
we constructed an equivalence of~$\cal{C}$ with $\cal{C}(\pi)$, where
$\pi=\Aut(F)$, which carries $F$ to the canonical functor from
\ifthenelse{\boolean{orig}}{$\cal{C}(\pi)_x$}{$\cal{C}(\pi)$} to the category of finite sets. In the standard case
$\cal{C}=\cal{C}(\pi)$, with $F$ the canonical functor, the fundamental
pro-object associated with~$F$ is precisely the projective system of
discrete quotients $\pi_i$ of~$\pi$.
\end{definition}

It may be useful to describe explicitly the category of pro-objects of
$\cal{C}(\pi)$. We obtain:

\begin{proposition}
\label{V.5.2}
The category $\Prodash\cal{C}(\pi)$
\label{indnot:ed}\oldindexnot{$\Prodash\cal{C}(\pi)$|hyperpage}%
%
is canonically equivalent to the category $\cal{C}'(\pi)$ of spaces,
equipped with a topological group of operators~$\pi$, which are compact
and totally disconnected.
\end{proposition}

The latter category contains $\cal{C}(\pi)$ as a full subcategory,
corresponding to compact discrete spaces with operators, and projective
limits exist in it. There is therefore a canonical functor
$$
g\colon \Prodash\cal{C}(\pi)\to\cal{C}'(\pi),
$$
which associates to a projective system $Q=(Q_i)$ the object
$X=\varprojlim_iQ_i$ of~$\cal{C}'(\pi)$. To define a functor in the
opposite direction, it is enough to define a contravariant functor from
$\cal{C}'(\pi)$ to the category of left-exact functors
$\cal{C}\to\Ens$. For every $X\in\cal{C}'(\pi)$, take
$$
h(X)(E)=\Hom(X,E)
$$
% page: 128
\marginpar{128}
%
(where $\Hom$ is taken in $\cal{C}'(\pi)$). It follows immediately from
the definitions that $h$ and~$g$ are adjoint, and that $hg$ is
canonically isomorphic to the identity functor of
$\Prodash\cal{C}(\pi)$. To prove that $g$ and~$h$ are quasi-inverse, it
remains to show that every object of~$\cal{C}'(\pi)$ is isomorphic to an
object of the form $g(Q)$, with $Q\in\Prodash\cal{C}(\pi)$. In other
words, \emph{every compact, totally disconnected space~$X$ equipped with
a topological group of operators~$\pi$ is isomorphic to a projective
limit of finite discrete spaces with operators}.

As a topological space without operators, $X$ is the projective limit of
its finite discrete quotients. We are therefore reduced to showing that
among these quotients there is a cofinal system invariant under~$\pi$.
For a given quotient $X'$, it is enough to show that the set of its
transforms under the operations of~$\pi$ is finite; one then takes the
supremum of these transforms, which is an invariant quotient dominating
$X'$. Equivalently, it is enough to find an open normal subgroup $\pi'$
of~$\pi$ whose elements act trivially on~$X'$. The quotient $X'$ amounts
to a finite partition of~$X$ into open subsets~$X_i$. By continuity and
compactness of~$\pi$, there exists a neighborhood~$V$ of the identity of
$\pi$ such that $s\in V$ implies $s\cdot X_i\subset X_i$ for every~$i$;
thus $s$ acts trivially on~$X'$. The open normal subgroups of~$\pi$ form
a fundamental system of neighborhoods of the identity, which completes
the proof.

Notice, more simply, that the category $\Ind\cal{C}(\pi)$ is canonically
equivalent to the category of sets on which $\pi$ acts continuously. We
shall not need this fact here.

\begin{proposition}
\label{V.5.3}
Let $\cal{C}$ be a Galois category, let $F$ be a fundamental functor
\ifthenelse{\boolean{orig}}{on~$C$,}{on~$\cal{C}$,}
% correction: `$C$' should be `$\cal{C}$'.
and let $P=(P_i)$ be the associated pro-object, normalized in the usual
way. Let
\ifthenelse{\boolean{orig}}{$X\in C$;}{$X\in\cal{C}$;}
% correction: `$C$' should be `$\cal{C}$'.
then $X$ is connected if and only if $\pi$ acts transitively on
$E=F(X)$.
\end{proposition}

We reduce to the standard case $\cal{C}=\cal{C}(\pi)$, with $F$ the
canonical functor, where the assertion is immediate.

\begin{corollary}
\label{V.5.4}
The following conditions on~$X$ are equivalent: \textup{(i)} $X$ is
connected and $\not\simeq\emptyset_{\cal{C}}$; \textup{(ii)} the group
$\pi$ acts transitively on $E=F(X)$, and $F(X)\ne\emptyset$;
\textup{(iii)} $X$ is isomorphic to some~$P_i$.
\end{corollary}

The equivalence
% page: 129
\marginpar{129}
%
of \textup{(i)} and \textup{(iii)} also follows readily from
\No~\Ref{V.4}, e).

\begin{proposition}
\label{V.5.5}
Let $Q=(Q_i)_{i\in I}$ be a pro-object
\ifthenelse{\boolean{orig}}{of~$C$,}{of~$\cal{C}$,}
% correction: `$C$' should be `$\cal{C}$'.
normalized in the usual way, and let $G$ be the corresponding functor
$G(X)=\Hom(Q,X)$ from~$\cal{C}$
\ifthenelse{\boolean{orig}}{(to~$\Ens$.}{(to~$\Ens$).}
% correction: insert `)'.
The following conditions are equivalent:
\begin{enumerate}
\item[(i)] $G$ commutes with finite direct sums;
\item[(ii)] $G$ commutes with the direct sum of two objects;
\item[(iii)] the $Q_i$ are connected and
\ifthenelse{\boolean{orig}}{$\not\simeq\emptyset_C$}
{$\not\simeq\emptyset_\cal{C}$};
% correction: `$C$' should be `$\cal{C}$'.
\item[(iv)] $Q$ is isomorphic to $\pi/H$, where $H$ is a closed subgroup
of~$\pi$;
\item[(v)] the functor $G$ is isomorphic to the functor
$E\mto E^H$, the set of $H$-invariants, defined by a closed subgroup
$H$ of~$\pi$.
\end{enumerate}
\end{proposition}

N.B.\ In conditions \textup{(iv)} and \textup{(v)}, a fundamental
functor is understood to have been chosen, thereby identifying
$\cal{C}$ with the category~$\cal{C}(\pi)$.

\subsubsection*{Proof}
We may suppose that $\cal{C}=\cal{C}(\pi)$. The implication
\textup{(i)}$\To$\textup{(ii)} is immediate, and
\textup{(ii)}$\To$\textup{(iii)} is proved as in property d) of
\No~\Ref{V.4}. Let us prove \textup{(iii)}$\To$\textup{(iv)}. The set
$\varprojlim_iQ_i$ is nonempty, being the projective limit of nonempty
finite sets. Let $a\in\varprojlim_iQ_i$. It defines a homomorphism of
spaces with operators
$$
\pi\to Q
$$
which is \emph{surjective}: for every~$i$, the composite
$\pi\to Q\to Q_i$ is surjective, since $\pi$ acts transitively on~$Q_i$
by Proposition~\Ref{V.5.3}. If $H$ is the stabilizer of~$a$, we obtain an
isomorphism $\pi/H\isomto Q$. The implications
\textup{(iv)}$\To$\textup{(v)} and
\textup{(v)}$\To$\textup{(i)} are again immediate.

\begin{proposition}
\label{V.5.6}
Let $\cal{C}$ be a Galois category, let $P$ be a fundamental pro-object
of~$\cal{C}$, and let $F$ be the associated fundamental functor. Let
$P'=(P'_i)_{i\in I}$ be a pro-object of~$\cal{C}$ in normal form, and let
$F'$ be the associated functor $F'(X)=\Hom(P',X)$ from~$\cal{C}$ to
$\Ens$. The following conditions are equivalent:
\begin{enumerate}
\item[(i)] $P'\simeq P$, or equivalently $F'\simeq F$;
\item[(ii)] $P'$ is fundamental, or equivalently $F'$ is fundamental;
\item[(iii)] $F'$ carries a direct sum of two objects to their direct sum,
and $X\not\simeq\emptyset_{\cal{C}}$ implies
$F'(X)\ne\emptyset$;
\item[(iv)]
% page: 130
\marginpar{130}
%
the connected objects of~$\cal{C}$ distinct from
$\emptyset_{\cal{C}}$ are exactly those isomorphic to some~$P'_i$.
\end{enumerate}
\end{proposition}

We have trivially \textup{(i)}$\To$\textup{(iii)} and
\textup{(i)}$\To$\textup{(ii)}; moreover,
\textup{(ii)}$\To$\textup{(iv)} by Corollary~\Ref{V.5.4}, applied to
$P'$ in place of~$P$. Furthermore, either \textup{(iii)} or
\textup{(iv)}, by Proposition~\Ref{V.5.5}, implies that $P'$ is of the
form $\pi/H$, where $H$ is a closed subgroup of~$\pi$. In case
\textup{(iii)}, for every open normal subgroup $\pi'$ of~$\pi$ there is a
$\pi$-homomorphism $P'=\pi/H\to\pi/\pi'$. Hence $H\subset\pi'$, so
$H=(0)$, and consequently \textup{(i)} holds. This proves the assertion.

\begin{corollary}
\label{V.5.7}
Let $\cal{C}$ be a Galois category. Its fundamental pro-objects are
isomorphic, and its fundamental functors are isomorphic.
\end{corollary}

In other words, \emph{the category of fundamental functors is a connected
groupoid~$\Gamma$},
\label{indnot:ef}\oldindexnot{$\Gamma$|hyperpage}%
%
which may be called the \emph{fundamental groupoid}
%
\index{fundamental groupoid|hyperpage}%
%
of the Galois category~$\cal{C}$. If $\cal{C}=\cal{C}(\pi)$, the
automorphism group of an object of the fundamental groupoid is isomorphic
to~$\pi$, this isomorphism being well determined up to an inner
automorphism. (Recall that a \emph{groupoid} is a category in which every
morphism is an isomorphism, and a \emph{connected} groupoid is a groupoid
all of whose objects are isomorphic.) The fundamental pro-objects of
$\cal{C}$ form a connected groupoid equivalent to the \emph{opposite} of
the fundamental groupoid. If $F,F'$ are two fundamental functors,
associated with fundamental pro-objects $P,P'$, then
$\Hom(F,F')=\Isom(F,F')$ is sometimes denoted by~$\pi_{F',F}$ and plays
the role of a ``set of \emph{homotopy classes of paths} from~$F$ to~$F'$''.
In particular, $\pi_{F,F}=\pi_F$ is precisely the \emph{fundamental group
of~$\cal{C}$ at~$F$} constructed in the preceding section. The pro-object
$P$ associated with~$F$ plays the role of a \emph{universal covering
at~$F$} of the final object $e_{\cal{C}}$ of~$\cal{C}$.

It may be convenient to have a description of~$\cal{C}$, up to
equivalence, in terms of its fundamental groupoid~$\Gamma$, without
choosing a particular object~$F$ of the latter. To every object $X$ of
$\cal{C}$ is associated a functor $E_X$ on the fundamental groupoid,
defined by
$$
E_X(F)=F(X),
$$
with values in~$\Ens$. Such a functor is known in topology as a ``local
system'' on the groupoid. The set $F(X)=E_X(F)$ may be called the
\emph{fiber} of~$X$ at~$F$, and $E_X$ the fiber functor associated
with~$X$. The functor $E_X$ has the following property:
% page: 131
\marginpar{131}
%
\ifthenelse{\boolean{orig}}
{\emph{\underline{for every} $F$, $E_X(F)$
% remark: Proposition 5.8 refers to this underlined condition, so it
% must remain underlined in the original branch.
\underline{is a finite set on which the topological group}
$\pi_F=\Aut(F)$ \underline{acts continuously}}}
{\emph{for every $F$, $E_X(F)$ is a finite set on which the topological
group $\pi_F=\Aut(F)$ acts continuously}}.
For a given covariant functor~$\xi$ from the fundamental groupoid to
$\Ens$, the preceding condition is equivalent to the same condition for
any one fixed~$F$. With this understood:

\begin{proposition}
\label{V.5.8}
The functor $X\mto E_X$ is an equivalence from the category $\cal{C}$ to
the category of covariant functors from the fundamental groupoid
$\Gamma$ of~$\cal{C}$ to~$\Ens$ which satisfy the condition
\ifthenelse{\boolean{orig}}
{underlined above.}
{highlighted above.}
\end{proposition}

Indeed, let $F_0$ be an object of the fundamental groupoid and put
$\pi_0=\Aut(F_0)$. The functor $\xi\mto\xi(F_0)$ is an equivalence from
the second category in Proposition~\Ref{V.5.8} to the category
$\cal{C}(\pi_0)$, as is seen immediately. On the other hand, its composite
with the functor $X\mto E_X$ is the natural equivalence
$\cal{C}\to\cal{C}(\pi_0)$. It follows that $X\mto E_X$ itself is an
equivalence.

\begin{corollary}
\label{V.5.9}
The category $\Prodash\cal{C}$
% remark: the Pro-C does not come out well in print; the C should be
% moved slightly to the left.
is canonically equivalent to the category of covariant functors~$\xi$
from the fundamental groupoid~$\Gamma$ to the category of topological
spaces satisfying the following condition: for every object~$F$ of
$\Gamma$, $\xi(F)$ is a compact, totally disconnected space equipped
with a topological group of operators~$\pi_F$.
\end{corollary}

Here again, it is enough to verify this condition on~$\xi$ for a single
$F$. The proof is the same as that of Proposition~\Ref{V.5.8}.

\begin{remarques}
\label{V.5.10}
Let $(F_s)_{s\in S}$ be a family
\ifthenelse{\boolean{orig}}{of object}{of objects}
% correction: objects.
of the fundamental groupoid~$\Gamma$. For $s,s'\in S$, put
$$
\Hom(s,s')=\Hom(F_s,F_{s'}).
$$
The set~$S$ thereby becomes a connected groupoid, and the map
$s\mto F_s$ becomes a fully faithful functor $f\colon S\to\Gamma$.
Consider the functor $X\mto E_X\circ f$ from~$\cal{C}$ to the functor
category $\Hom(S,\Ens)$. This gives a variant of
Proposition~\Ref{V.5.8}, and of Corollary~\Ref{V.5.9}, with $\Gamma$
replaced by~$S$. The resulting statement reduces to
Theorem~\Ref{V.4.1} when $S$ consists of one point, and is precisely
Proposition~\Ref{V.5.8} itself when $S$ is the set of objects of~$\Gamma$.
\end{remarques}

We shall use Corollary~\Ref{V.5.9} to define a canonical pro-object of
$\cal{C}$. To this end, consider the functor from~$\Gamma$ to the category
of topological spaces
% page: 132
\marginpar{132}
%
and indeed of topological groups:
$$
f\colon F\mto\Aut(F)=\pi_F.
$$
This functor satisfies the condition in Proposition~\Ref{V.5.8}. The
space with operators $f(F)$ under~$\pi_F$ is simply $\pi_F$ itself,
regarded as a space with operators under itself by inner automorphisms.
Thus $f$ corresponds to a pro-object of~$\cal{C}$ determined up to unique
isomorphism. It is in fact a pro-group in~$\cal{C}$, called the
\emph{fundamental pro-group of~$\cal{C}$},
%
\index{fundamental pro-group of a Galois category|hyperpage}%
%
and plays the role of a local system of fundamental groups. Thus it is a
pro-group~$\Pi$ of~$\cal{C}$ defined by the existence of an isomorphism,
functorial in~$F$,
$$
F(\Pi)\simeq\pi_F.
$$

If $X$ is any pro-object of~$\cal{C}$, there is a canonical morphism
$$
\Pi\times X\to X
$$
which makes $X$ an object with a left group of operators~$\Pi$ in
$\Prodash\cal{C}$. It is enough to observe that, for variable~$F$, there
is a canonical map
$$
\Pi(F)\times X(F)\to X(F),
$$
that is,
$$
\Aut(F)\times E_X(F)\to E_X(F),\quad\text{or}\quad
\pi_F\times F(X)\to F(X),
$$
which is functorial in~$F$. It is also functorial in~$X$, so for every
morphism $X\to Y$ of pro-objects the corresponding diagram
$$
\xymatrix{
\Pi\times X \ar[r]\ar[d]&X\ar[d]\\
\Pi\times Y \ar[r]&Y
}
$$
is commutative.

\begin{remarques}
\label{V.5.11}
One must not confuse a fundamental pro-object~$P$, which has no group
structure and is connected, with the fundamental pro-group~$\Pi$, which
is a pro-\emph{group} and is generally not connected. More precisely,
$\Pi$ is connected if
\ifthenelse{\boolean{orig}}{\ignorespaces}{and }
% correction: insert `and'.
only if $\pi_F$, acting on itself by inner automorphisms, acts
transitively; that is, if and only if $\pi$ consists only of the identity,
or equivalently if $\cal{C}$ is equivalent to the category of finite
sets. Another essential difference is that
% page: 133
\marginpar{133}
%
$\Pi$ is determined up to unique isomorphism, whereas $P$ is determined
only up to nonunique isomorphism.

Let $E$ be a finite set and consider the constant functor on the groupoid
$\Gamma$ with value~$E$. By Proposition~\Ref{V.5.8}, it defines an object
of~$\cal{C}$, denoted by~$E_{\cal{C}}$, which may also be interpreted as
the direct sum of $E$ copies of the final object $e_{\cal{C}}$ of
$\cal{C}$. We may regard $E_{\cal{C}}$ as a functor of~$E$, from the
category of finite sets to~$\cal{C}$. This functor is \emph{exact}, and
hence carries finite groups to $\cal{C}$-groups,
\ifthenelse{\boolean{orig}}{etc...}{etc.}
If an object $X$ of~$\cal{C}$ carries a right action of a finite group
$G$, we may therefore regard $X$ as an object of~$\cal{C}$ with a right
group of operators, the $\cal{C}$-group~$G_{\cal{C}}$. Extending the
general terminology for objects with $\cal{C}$-groups of operators, we
say that $X$ is \emph{formally principal homogeneous}
%
\index{formally principal homogeneous object|hyperpage}%
%
under~$G$ if it is formally principal homogeneous under
$G_{\cal{C}}$, that is, if the canonical morphism
$$
X\times G_{\cal{C}}\to X\times X
$$
induced by the right action of~$G_{\cal{C}}$ on~$X$ is an isomorphism.
We say that $X$ is \emph{principal homogeneous}
%
\index{principal homogeneous object|hyperpage}%
%
under~$G$ if it is so under $G_{\cal{C}}$, that is, if it is formally
principal homogeneous and, in addition, the quotient
$X/G=X/G_{\cal{C}}$ is~$e_{\cal{C}}$.

Fix a fundamental functor, and hence an equivalence of~$\cal{C}$ with a
category $\cal{C}(\pi)$. The object~$X$ corresponds to a set
$E=F(X)$ on which $\pi$ acts continuously on the left. Giving a right
action of~$G$ on~$X$ is then equivalent to giving a right action of~$G$
on the set~$E$ which commutes with the action of~$\pi$. It follows at once
that $X$ is principal homogeneous under~$G$ if and only if $E$ is a
principal homogeneous space
\ifthenelse{\boolean{orig}}
{under $G$}
{under $G$,}
% correction: insert the comma.
that is, if and only if $G$ acts simply transitively on~$E$. Moreover,
$X$ is formally principal homogeneous $\sss$ $E$ is principal homogeneous
\emph{or} empty. Comparing with Proposition~\Ref{V.5.3}, we see that if
$X$ is principal homogeneous under~$G$ \emph{and connected}, then the
given homomorphism from~$G$ to the group opposite to~$\Aut(X)$ is an
\emph{isomorphism}. Furthermore, an object~$X$
\ifthenelse{\boolean{orig}}{of~$C$}{of~$\cal{C}$}
% correction: `$C$' should be `$\cal{C}$'.
is connected and principal homogeneous under the group opposite to
$\Aut(X)$ if and only if, with the notation of \No~\Ref{V.4}, it is
isomorphic to a \emph{Galois}~$P_i$. In the standard case
$\cal{C}=\cal{C}(\pi)$, this means that $X$ is isomorphic to a quotient
of~$\pi$ by a normal subgroup.

Continue to fix a fundamental functor~$F$. The data of
% page: 134
\marginpar{134}
%
an object~$X$ principal homogeneous under a finite group~$G$ acting on
the right, together with a point $a\in F(X)$, are equivalent to the data
of a homomorphism from~$\pi$ to~$G$. Indeed, to such a homomorphism we
associate the set $E=G$, letting $\pi$ act on the left through the given
homomorphism $\pi\to G$ and the left translations of~$G$, and letting
$G$ act on the right by right translations; the marked point~$a$ of~$E$
is the identity element of~$G$. By the preceding discussion, this gives,
in an essentially unique way, every triple $(X,G,a)$ with the properties
considered above,
\ifthenelse{\boolean{orig}}
{since a}
{since a}
pointed principal homogeneous space under a group~$G$ is identified with
$G$ itself. In this way we obtain a direct geometric interpretation of
the functor $G\mto\Hom(\pi,G)$ from the category of finite groups to
$\Ens$. This functor is pro-representable by~$\pi$, and its consideration
would therefore provide another construction of the group~$\pi$
associated with~$F$.
\end{remarques}
\end{filecontents*}
% Original input: drafts/SGA1_V_6_English.texfrag
\begin{filecontents*}[overwrite,noheader]{sga1-cumulative-input-094.tex}
\section{Exact functors from one Galois category to another}
\label{V.6}

\begin{proposition}
\label{V.6.1}
Let \(\cal{C}\) and \(\cal{C'}\) be two Galois categories, let
\(H\colon\cal{C}\to\cal{C'}\) be a covariant functor, let \(F'\) be a
fundamental functor on~\(\cal{C'}\), and set \(F=F'\circ H\). The
following conditions are equivalent:
\begin{enumerate}
\item[(i)] \(H\) is \emph{exact}, \ie both left exact and right exact.
\item[(ii)] \(H\) is left exact, carries finite sums to finite sums, and
epimorphisms to epimorphisms (or, equivalently, carries objects not
isomorphic to \(\emptyset_{\cal C}\) to objects not isomorphic to
\(\emptyset_{\cal C'}\)).
\item[(iii)] \(F\) is a fundamental functor on~\(\cal C\).
\end{enumerate}
\end{proposition}

The implication (i)\(\To\)(ii) is a general categorical fact.
Moreover, the first form of (ii) implies the second. Indeed, if \(X\)
is an object of~\(\cal C\), then \(X\not\approx\emptyset_{\cal C}\) if
and only if the morphism \(X\to e_{\cal C}\) is an epimorphism; note
also that, since \(F\) is assumed left exact, it carries \(e_{\cal C}\)
to~\(e_{\cal C'}\). The second form of (ii) implies~(iii), since \(F\),
being left exact and therefore pro-representable, satisfies criterion
(iii) of~\Ref{V.5.6}. Finally, (iii) implies~(i), because \(F\) is exact
and conservative (\ie it satisfies axiom (G~6) of~\no\Ref{V.4}).

Let \(\Gamma\) now be the fundamental groupoid of~\(\cal C\), and
\(\Gamma'\) that of~\(\cal C'\). Thus, if
% page: 135
\marginpar{135}
%
\(H\) is exact, then \(F'\mto F'\circ H\) is a functor from the
groupoid~\(\Gamma'\) to the groupoid~\(\Gamma\), which we denote
by~\(\tH\):
\[
  \tH(F')(X)=F'(H(X)).
\]
Using the notation \(F(X)=E_X(F)\) introduced in~\no\ifthenelse{\boolean{orig}}{\Ref{V.6}}{\Ref{V.5}}, this
may also be written
\[
\ifthenelse{\boolean{orig}}
{E_{H(X)}(F')=E_X(\tH(F'))}
{E_{H(X)}(F')=E_X(\tH(F')).}
\]
This last formula, together with~\Ref{V.5.8} or~\Ref{V.4.1}, shows that
the exact functor \(H\) is determined up to unique isomorphism once
the corresponding functor~\(\tH\) is known. Fix \(F'\), and set
\(F=\tH(F')\). Then \(\tH\) defines a homomorphism from
\(\pi_{F'}=\Aut(F')\) to \(\pi_F=\Aut(F)\):
\[
  \tH\colon\pi_{F'}\longrightarrow\pi_F
  \qquad(F=\tH(F')=F'\circ H).
\]
The preceding formula also shows, in view of~\Ref{V.5.8}, that this
homomorphism has the following property: for every finite set \(E\) on
which \(\pi_F\) acts continuously, the group \(\pi_{F'}\) also acts
\emph{continuously} through the preceding homomorphism
\(\pi_{F'}\to\pi_F\). Applying this to the quotients of~\(\pi_F\) by
its open normal subgroups, one sees that the preceding condition also
means that the homomorphism in question is continuous.

Conversely, let \(F\) be an object of~\(\Gamma\), let \(F'\) be an
object of~\(\Gamma'\), and let
\[
  u\colon\pi_{F'}\longrightarrow\pi_F
\]
be a continuous homomorphism. It determines a manifestly exact functor
from~\(\cal C(\pi)\) to~\(\cal C(\pi')\), and hence, by~\Ref{V.4.1},
an exact functor \(H\colon\cal C\to\cal C'\) for which
\(\tH\colon\pi_{F'}\to\pi_F\) is precisely~\(u\). Instead of starting
with a group homomorphism, one may start with a \emph{functor}
\[
  U\colon\Gamma'\longrightarrow\Gamma
\]
such that, for \emph{every} \(F'\in\Gamma'\) (or for \emph{one}
\(F'\in\Gamma'\), which amounts to the same thing), the corresponding
homomorphism \(\pi_{F'}\to\pi_F\) is continuous. Every such functor is
isomorphic to one of the form~\(\tH\), where
\(H\colon\cal C\to\cal C'\) is an exact functor determined up to unique
isomorphism. Thus:

\begin{corollary}
\label{V.6.2}
% page: 136
\marginpar{136}
%
A functor \(H\colon\cal C\to\cal C'\) between Galois categories is
exact if and only if there exist equivalences
\(\cal C(\pi)\to\cal C\) and
\ifthenelse{\boolean{orig}}
{\(\cal C'\to\cal C(\pi)\)}
{\(\cal C'\to\cal C(\pi')\)}
which carry \(H\) to the functor
\(\cal C(\pi)\to\cal C(\pi')\) associated with a homomorphism of
topological groups \(\pi'\to\pi\).
\end{corollary}

\begin{corollary}
\label{V.6.3}
Let \(\cal C\) and \(\cal C'\) be two Galois categories, with
fundamental groupoids \(\Gamma\) and \(\Gamma'\). Then the category of
exact functors from~\(\cal C\) to~\(\cal C'\) is equivalent to the
category of functors \(U\colon\Gamma'\to\Gamma\) having the following
property: for every \(F'\) in~\(\Gamma'\) (or for one \(F'\) in
\(\Gamma'\), which amounts to the same thing), if \(F=U(F')\), then
the homomorphism
\[
  \pi_{F'}=\Aut(F')\longrightarrow\pi_F=\Aut(F)
\]
defined by \(U\) is continuous.
\end{corollary}

Consider the fundamental pro-group \(\Pi\) of~\(\cal C\). An exact
functor \(H\) carries it to a pro-group \(H(\Pi)\) in~\(\cal C'\). We
shall define a homomorphism
\[
  \Pi'\longrightarrow H(\Pi),
\]
where \(\Pi'\) is the fundamental pro-group of~\(\cal C'\), by
requiring that, for every object \(F'\) of~\(\Gamma'\), the
corresponding homomorphism
\[
  F'(\Pi')=\pi_{F'}\longrightarrow F'(H(\Pi))=\pi_F
  \qquad
\ifthenelse{\boolean{orig}}
{\text{(where \(F=F'\circ H=\tH(F')\))};}
{\text{(where \(F=F'\circ H=\tH(F')\))}}
\]
be the natural homomorphism
\[
  \Aut(F')\longrightarrow\Aut(F'\circ H).
\]
Since the latter is functorial in~\(F'\), it does indeed define, by
\Ref{V.5.8}, a homomorphism of pro-objects, and in fact of pro-groups,
in~\(\cal C'\). This homomorphism is said to be \emph{associated} with
the functor~\(H\).

Now let \(H'\) be a second exact functor, from the Galois category
\(\cal C'\) to a Galois category~\(\cal C''\). It is immediate that
\[
  {}^t(H'H)=\tH\,\tH'
\]
(N.B. this is an identity of functors, not merely a canonical
isomorphism). There is an analogous transitivity property for the
associated homomorphisms of fundamental pro-groups.

We
% page: 137
\marginpar{137}
%
now interpret the properties of the exact functor \(H\) in terms of
the corresponding homomorphism
\[
  u\colon\pi_{F'}\longrightarrow\pi_F
  \qquad\text{(where \(F=F'\circ H\)).}
\]
It is convenient to introduce the notion of a \emph{pointed object}
%
\index{pointed object|hyperpage}%
%
of the Galois category~\(\cal C\), equipped with its fundamental
functor \(F\). By definition, this is an object \(X\) of~\(\cal C\)
together with an element \(a\) of~\(F(X)\). It may therefore be viewed
as a finite set on which \(\pi_F\) acts continuously on the left,
together with a point~\(a\). Consequently, by~\Ref{V.5.3}, the
\emph{connected} pointed objects of~\(\cal C\) are identified with the
open subgroups of~\(\pi_F\). If \(U\) and \(V\) are two such subgroups,
corresponding to connected pointed objects \(X\) and \(Y\) of~\(\cal C\),
then a pointed homomorphism from \(X\) to \(Y\) exists if and only if
\(U\subset V\), and it is then unique.

Of course, the functor \(H\) carries pointed objects to pointed objects,
since \(F=F'\circ H\). Also note that a closed subgroup of a group such
as \(\pi_F\) is the intersection of the open subgroups containing it.
Thus,
\ifthenelse{\boolean{orig}}{\ignorespaces}{if}
\(M\) and \(N\) are two closed subgroups, then \(M\subset N\) if and
only if every open subgroup containing \(N\) also contains \(M\).
These remarks easily yield the following results:

\begin{proposition}
\label{V.6.4}
Let \(X\) be a connected pointed object of~\(\cal C\), associated with
an open subgroup \(U\) of~\(\pi_F\). Then
\ifthenelse{\boolean{orig}}{\(u\)}{\(U\)}
contains \(u(\pi_{F'})\) (respectively, the closed normal subgroup
generated by~\(u(\pi_{F'})\)) if and only if \(H(X)\) admits a pointed
section (respectively, is completely decomposed).
\end{proposition}

By a \emph{section} of an object \(X\) of a Galois category~\(\cal C\)
we mean, implicitly over the final object, a morphism from the final
object \(e_{\cal C}\) to~\(X\). Equivalently, it is an element \(a\) of
\(F(X)\) fixed by~\(\pi_F\). If \(X\) is pointed, this is called a
\emph{pointed section} when it is compatible with the pointed
structures on \(X\) and \(e_{\cal C}\), \ie when \(a\) is precisely
the distinguished element of~\(F(X)\). Such a section is therefore
unique, and exists if and only if the distinguished element of~\(F(X)\)
is fixed by~\(\pi_F\). Finally, an object of a Galois category is said
to be \emph{completely decomposed}
%
\index{completely decomposed object|hyperpage}%
%
if it is isomorphic to a sum of final objects, \ie if \(\pi_F\) acts
trivially on~\(F(X)\). For a pointed \(X\), this condition is clearly
stronger than the existence of a pointed section.
Proposition~\Ref{V.6.4} follows immediately from the preceding
definitions and remarks.

\begin{corollary}
\label{V.6.5}
% page: 138
\marginpar{138}
%
The homomorphism \(u\) is trivial if and only if \(H(X)\) is completely
decomposed for every object \(X\) of~\(\cal C\).
\end{corollary}

\begin{proposition}
\label{V.6.6}
Let \(X'\) be a connected pointed object of~\(\cal C'\), associated
with an open subgroup \(U'\) of~\(\pi_{F'}\). Then \(U'\) contains
\(\Ker u\) if and only if there exist a connected pointed object \(X\)
of~\(\cal C\) and a pointed homomorphism from the distinguished
connected component \(X'_0\) of~\(H(X)\) to~\(X'\). Equivalently,
\(X'\) is isomorphic, as a pointed object, to a quotient of the
distinguished connected component of the inverse image of a pointed
object of~\(\cal C\). If \(u\) is surjective, the preceding condition
is also equivalent to the following: \(X'\) is isomorphic to an object
\(H(X)\), where \(X\) is a pointed object of~\(\cal C\).
\end{proposition}

(The \emph{distinguished connected component}
%
\index{distinguished connected component of a pointed object|hyperpage}%
%
of a pointed object \(X\) of a Galois category~\(\cal C\) is the
unique connected pointed subobject of~\(X\). By~\Ref{V.5.3}, it
corresponds to the orbit under \(\pi_F\) of the distinguished point
of~\(F(X)\).) Since the fact that \(U'\) contains \(\Ker u\) does not
depend on the chosen pointing of~\(X'\) (another pointing replaces
\(U'\) by a conjugate subgroup), we obtain:

\begin{corollary}
\label{V.6.7}
The subgroup \(U'\) contains \(\Ker u\) if and only if there exist an
object \(X\) of~\(\cal C\), which may be assumed connected, and a
morphism from a connected component of~\(H(X)\) to~\(X'\). If \(u\)
is surjective, this also says that \(X'\) is isomorphic to an object
of the form~\(H(X)\).
\end{corollary}

\begin{corollary}
\label{V.6.8}
The homomorphism \(u\) is injective if and only if, for every object
\(X'\) of~\(\cal C'\), there exist an object \(X\) of~\(\cal C\) and
a homomorphism from a connected component of~\(H(X)\) to~\(X'\).
\end{corollary}

\begin{proposition}
\label{V.6.9}
The following conditions are equivalent:
\begin{enumerate}
\item[(i)] The homomorphism \(u\colon\pi_{F'}\to\pi_F\) is surjective.
\item[(ii)] For every connected object \(X\) of~\(\cal C\), the object
\(H(X)\) is connected.
\item[(iii)] The functor \(H\) is fully faithful.
\end{enumerate}
\end{proposition}

The last condition means that, for any two objects \(X,Y\) of~\(\cal C\),
the natural map
\[
  \Hom(X,Y)\longrightarrow\Hom(H(X),H(Y))
\]
is bijective.

\begin{corollary}
\label{V.6.10}
% page: 139
\marginpar{139}
%
The homomorphism \(u\) is an isomorphism if and only if \(H\) is an
equivalence of categories, or equivalently if the following two
conditions hold:
\begin{enumerate}
\item[a)] for every connected object \(X\) of~\(\cal C\), the object
\(H(X)\) is connected;
\item[b)] every object of~\(\cal C'\) is isomorphic to an object of the
form~\(H(X)\).
\end{enumerate}
\end{corollary}

\begin{proposition}
\label{V.6.11}
Let \(H\colon\cal C\to\cal C'\) and
\(H'\colon\cal C'\to\cal C''\) be exact functors between Galois
categories. Let \(F''\) be a fundamental functor on~\(\cal C''\), set
\(F'=F''H'\) and \(F=F'H\), and consider the associated homomorphisms
\[
  u'\colon\pi_{F''}\longrightarrow\pi_{F'}
  \qquad\text{and}\qquad
  u\colon\pi_{F'}\longrightarrow\pi_F.
\]
The inclusion \(\Ker u\supset\Im u'\), or equivalently the triviality
of \(uu'\), holds if and only if \(H'(H(X))\) is completely decomposed
for every object \(X\) of~\(\cal C\). The inclusion
\(\Ker u\subset\Im u'\) holds if and only if, for every connected
pointed object \(X'\) of~\(\cal C'\) such that \(H'(X')\) admits a
pointed section, there exist an object \(X\) of~\(\cal C\) and a
homomorphism from a connected component of~\(H(X)\) to~\(X'\).
\end{proposition}

The first assertion follows from the last assertion of~\Ref{V.6.4}.
The second follows by combining~\Ref{V.6.4} and~\Ref{V.6.6}.

\begin{remarques}
\label{V.6.12}
Under the conditions of~\Ref{V.6.8}, it is not true in general that
\(X'\) is isomorphic to an object of the form~\(H(X)\). One can show
that every connected object, and hence every object, of~\(\cal C'\)
is isomorphic to an object of the form \(H(X)\) if and only if \(u\)
is an isomorphism from~\(\pi_{F'}\) onto a \emph{direct-factor}
subgroup of~\(\pi_F\). In practice, however, one constructs directly a
\ifthenelse{\boolean{orig}}{right}{left} inverse \(\pi_F\to\pi_{F'}\) of~\(u\), using a suitable exact
functor from~\(\cal C'\) to~\(\cal C\).
\end{remarques}

\begin{proposition}
\label{V.6.13}
Let \(\cal C\) be a Galois category equipped with a fundamental
functor \(F\), let \(S\) be a connected object of~\(\cal C\), and let
\(\cal C'\) be the category of objects of~\(\cal C\) over~\(S\). Then
\(\cal C'\) is a Galois category, and the functor
\[
  X\mto H(X)=X\times S
\]
from~\(\cal C\) to~\(\cal C'\) is exact. Let \(a\in F(S)\), and let
\(F'\) be the functor from~\(\cal C'\) to the category of finite sets
defined by
\[
  F'(X')=\text{the inverse image of \(a\) under }
  F(X')\longrightarrow F(S).
\]
Then there is an isomorphism \(F\cong F'\circ H\), and the
corresponding homomorphism
\[
  u\colon\pi_{F'}\longrightarrow\pi_F
\]
% page: 140
\marginpar{140}%
%
is an isomorphism from~\(\pi_{F'}\) onto the open subgroup \(U\) of
\(\pi_F\) stabilizing the distinguished element \(a\) of~\(F(S)\).
\end{proposition}

The proof is left to the reader.

\end{filecontents*}
% Original input: drafts/SGA1_V_7_THROUGH_9_English.texfrag
\begin{filecontents*}[overwrite,noheader]{sga1-cumulative-input-095.tex}
\section{The case of preschemes}
\label{V.7}

Let \(S\) be a connected locally noetherian prescheme, and let
\[
  a\colon\Spec(\Omega)\longrightarrow S
\]
be a geometric point of~\(S\), with values in an algebraically closed
field~\(\Omega\). Set
\[
  \cal C=\text{the category of étale coverings of \(S\),}
\]
and, for an object \(X\) of~\(\cal C\), \ie an étale covering \(X\)
of~\(S\), set
\[
  F(X)=\text{the set of geometric points of \(X\) lying over \(a\).}
\]
Thus \(F\) becomes a functor from~\(\cal C\) to the category of finite
sets. Properties (G~1) through (G~6) hold: (G~1) is included among the
standard consequences in \ifthenelse{\boolean{orig}}{I~}{}\Ref{I.4.6}; (G~2) follows from~\Ref{V.3.4};
(G~3) from~\Ref{V.3.5}; (G~4) is immediate from the definition;
(G~5) follows from~\Ref{V.3.5} and the beginning of~\no\Ref{V.2}; and
(G~6) is proved in~\Ref{V.3.7}. We may therefore apply the results of
\no\Ref{V.4}, \Ref{V.5}, and~\Ref{V.6}.

In particular, this defines a pro-object \(P\) of~\(\cal C\)
representing \(F\), called the \emph{universal covering of \(S\) at
the point \(a\)},
%
\index{universal covering of a prescheme at a point|hyperpage}%
%
and a topological group
\(\pi=\Aut(F)=\Aut^0P\), called the \emph{fundamental group of \(S\)
at \(a\)},
%
\index{fundamental group of a prescheme at a point|hyperpage}%
%
and denoted by~\(\pi_1(S,a)\).
\label{indnot:eg}\oldindexnot{\(\pi_1(S,a)\)|hyperpage}%
%
The functor \(F\) then defines an equivalence between~\(\cal C\) and
the category of finite sets on which \(\pi=\pi_1(S,a)\) acts
continuously. This equivalence interprets the usual finite projective-
and inductive-limit operations on coverings---products, fiber
products, sums, passage to quotients,
\ifthenelse{\boolean{orig}}{etc...}{etc.}---in terms of the analogous
operations in~\(\cal C(\pi)\), \ie the evident operations on finite
sets with a \(\pi\)-action. Since the topological connected components
of an étale covering are themselves étale coverings, note also that
\emph{an object \(X\) of~\(\cal C\) is connected in~\(\cal C\) if and
only if it is topologically connected}. By~\Ref{V.5.3}, this says that
\(\pi_1\) acts transitively on~\(F(X)\).

% page: 141
\marginpar{141}
%
For an object \(X\) of~\(\cal C\) to be faithfully flat and
quasi-compact over~\(S\)---it is already flat and quasi-compact
over~\(S\)---it is necessary and sufficient that \(X\to S\) be
surjective, \ie an epimorphism in~\(\cal C\), or equivalently that
\(X\ne\emptyset\). Criterion~\Ref{V.2.6}~(iii) therefore shows that
\emph{\(X\) is a principal covering of
\ifthenelse{\boolean{orig}}{\(X\)}{\(S\)}
with group \(G\) if and only if it is a principal homogeneous space
under \(G\) in the category~\(\cal C\)}, in the sense defined
in~\no\Ref{V.5}.

If \(a'\) is another geometric point of~\(S\), corresponding to an
algebraically closed field \(\Omega'\), which may differ from
\(\Omega\) and may even have a different characteristic, then it
defines another fiber functor \(F'=F_{a'}\) from~\(\cal C\) to the
category of finite sets. This functor is again exact, and hence is
isomorphic to \(F=F_a\). Consequently, the fundamental groups
\(\pi_1(S;a)\)
\label{indnot:eg1}\oldindexnot{\(\pi_1(S,a)\)|hyperpage}%
%
for varying \(a\) are isomorphic to one another.

Let \(\pi_1(S;a,a')\)
\label{indnot:eh}\oldindexnot{\(\pi_1(S;a,a')\)|hyperpage}%
%
denote the set of isomorphisms---or, equivalently, the set of
homomorphisms---from \(F_a\) to \(F_{a'}\) between the associated fiber
functors. We thereby obtain a \emph{groupoid} whose objects are the
geometric points of~\(S\), the fundamental groups being the
automorphism groups of the objects of this groupoid. The set
\(\pi_1(S;a',a)\) may be called the \emph{set of homotopy classes of
paths from \(a\) to \(a'\)}. These classes compose in the evident way.

Finally, one may define a pro-group \(\Pi_1^S\) in~\(\cal C\), called
the \emph{fundamental pro-group of \(S\)}
%
\index{fundamental pro-group of a prescheme|hyperpage}%
%
or the \emph{local system of fundamental groups on \(S\)},
%
\index{local system of fundamental groups|hyperpage}%
%
characterized up to unique isomorphism by the existence of an
isomorphism, functorial in the geometric point \(a\) of~\(S\),
\[
  F_a(\Pi_1^S)=\pi_1(S;a)
\]
(\cf Remark~\Ref{V.5.10}). In particular, if \(s\) is an ordinary
point of~\(S\), the fiber of~\(\Pi_1^S\) at~\(s\) is a pro-group over
\(\kres(s)\), a projective limit of finite étale groups
over~\(\kres(s)\). One might call this pro-group the
\emph{fundamental group of \(S\) at the ordinary point \(s\)} and
denote it by~\(\pi_1(S,s)\). By definition, its points with values in
an algebraically closed extension \(\Omega\) of~\(\kres(s)\) are the
elements of~\(\pi_1(S;a)\), where \(a\) is the geometric point of~\(S\)
defined by that extension.

In particular, by taking \(S\) to be the spectrum of a field, every
field \(k\) has canonically and functorially associated with it a
pro-group over~\(k\), which one might denote by~\(\pi_1(k)\). It is a
projective limit of finite étale groups over~\(k\), and its points in
an algebraically closed extension \(\Omega\) of~\(k\) identify
% page: 142
\marginpar{142}
%
with the elements of the topological Galois group of~\(\bar k/k\),
where \(\bar k\) is the separable closure of~\(k\) in~\(\Omega\)
(\cf \Ref{V.8.1}). This group \(\pi_1(k)\) does not yet seem to have
attracted the attention of algebraists.

Now let
\[
  f\colon S'\longrightarrow S
\]
be a morphism from one connected locally noetherian prescheme to
another, let \(a'\) be a geometric point of~\(S'\), and let
\(a=f(a')\) be its direct image in~\(S\). The inverse-image functor
induces a functor from the category \(\cal C(S)\)
\label{indnot:ei}\oldindexnot{\(\cal C(S)\)|hyperpage}%
%
of étale coverings of~\(S\) to the category \(\cal C(S')\) of étale
coverings of~\(S'\):
\[
  f^\bbullet\colon\cal C(S)\longrightarrow\cal C(S').
\]
There is an isomorphism of functors
\[
  F_a\cong F_{a'}\circ f^\bbullet,
\]
so \(f^\bbullet\) is an \emph{exact} functor to which the results
of~\no\Ref{V.6} apply. In particular, there is a canonical homomorphism
\[
  u=\pi_1(f;a')\colon\pi_1(S',a')\longrightarrow\pi_1(S;a)
  \qquad(a=f(a')).
\]
\label{indnot:ej}\oldindexnot{\(\pi_1(f;a')\)|hyperpage}%
%
This homomorphism recovers the inverse-image functor as restriction of
the acting groups. Thus the properties of \(f^\bbullet\) are expressed
simply in terms of the properties of the associated group
homomorphism, as explained in~\no\Ref{V.6}. In particular, if \(S'\)
is an étale covering of~\(S\), then \(u\) is an isomorphism from
\(\pi_1(S',a')\) onto the open subgroup of~\(\pi_1(S,a)\) defining the
pointed connected étale covering \(S'\) of~\(S\), \ie the stabilizer
\(U\) of \(a'\in F_a(S')\) in~\(\pi_1(S;a)\).

To interpret the homomorphisms \(\pi_1(f;a')\) as the geometric point
\(a'\) varies, one must, in accordance with~\no\Ref{V.6}, consider a
homomorphism
\[
  \Pi_1(f)\colon\Pi_1^{S'}\longrightarrow f^\bbullet(\Pi_1^S)
\]
of pro-groups over~\(S'\), and then take the corresponding
homomorphism on geometric fibers.

\section{The case of a normal base prescheme}
\label{V.8}
% page: 143
\marginpar{143}
%

\begin{proposition}
\label{V.8.1}
Let \(S\) be the spectrum of a field \(k\), and let \(\Omega\) be an
algebraically closed extension of~\(k\), defining a geometric point
\(a\) of~\(S\) with values in~\(\Omega\). Let \(\bar k\) be the
separable closure of~\(k\) in~\(\Omega\). Then there is a canonical
isomorphism from~\(\pi_1(S,a)\) onto the topological Galois group of
\(\bar k/k\).
\end{proposition}

Let \(k'\) be the algebraic closure of~\(k\) in~\(\Omega\); it defines
a geometric point \(b\) of~\(S\) with values in~\(k'\). The natural
homomorphism of functors \(F_b\to F_a\) is plainly an isomorphism,
because every \(k\)-homomorphism from a finite separable extension
of~\(k\) into~\(\Omega\) necessarily has its image in~\(\bar k\), and
hence in~\(k'\). The group \(\pi'\) of \(k\)-automorphisms of~\(k'\)
acts naturally on~\(F_b\), giving a homomorphism
\[
  \pi'\longrightarrow\Aut(F_b)\isomto\Aut(F_a)=\pi_1(S;a).
\]
It is well known that the natural homomorphism from~\(\pi'\) to the
group \(\pi\) of automorphisms of~\(\bar k/k\) is an isomorphism. We
thus obtain a canonical homomorphism
\(\pi\to\pi_1(S;a)\), and it remains to show that it is an
isomorphism. It is injective because an element of its kernel is an
automorphism of~\(\bar k/k\) inducing the identity on every finite
separable subextension, and is therefore trivial. It is surjective
because, if \(X\) is a \emph{connected} étale covering of~\(S\), hence
defined by a finite separable \emph{extension} \(L/k\), then \(\pi\)
acts transitively on the set of \(k\)-homomorphisms from~\(L\)
into~\(k'\), as is well known.

\begin{proposition}
\label{V.8.2}
Let \(S\) be a connected, locally noetherian, normal prescheme, let
\(K=\kres(s)\) be its function field, equal to the residue field at
its generic point \(s\), and let \(\Omega\) be an algebraically closed
extension of~\(K\), defining a geometric point \(a'\) of
\(S'=\Spec(K)\) and a geometric point \(a\) of~\(S\). Then the
homomorphism
\[
  \pi_1(S';a')\longrightarrow\pi_1(S;a)
\]
is surjective. If the first group is identified with the Galois group
of the separable closure \(\bar K\) of~\(K\) in~\(\Omega\)
(\cf \Ref{V.8.1}), then the kernel of the preceding homomorphism
corresponds under Galois theory to the subextension of~\(\bar K/K\)
obtained by composing the finite extensions of~\(K\) in~\(\Omega\)
that are unramified over~\(S\).
\end{proposition}

The first assertion says that the inverse image on~\(S'\) of a
connected étale covering \(X\) of~\(S\) is connected, \ie that \(X\)
is integral; this is precisely \ifthenelse{\boolean{orig}}{I~}{}\Ref{I.10.1}. The kernel of the
preceding homomorphism consists of those automorphisms of~\(\bar K/K\)
that induce the identity on the sets \(F_a(X)\),
% page: 144
\marginpar{144}
%
where the étale covering \(X\) of~\(S\) may be assumed connected. This
means exactly that the automorphism induces the identity on the finite
subextensions of~\(\bar K/K\) that are unramified over~\(S\), proving
the last assertion.

\begin{remarquestar}
Through this interpretation of the fundamental group of the normal
prescheme \(S\) in terms of ordinary Galois theory, the definition in
this case has long been known.
\end{remarquestar}

\section{The case of disconnected preschemes: multi-Galois categories}
\label{V.9}

Let \(S\) be a locally noetherian prescheme, and let
\((S_i)_{i\in I}\) be its connected components. The category
\(\cal C(S)\) of étale coverings of~\(S\) is equivalent to the product
category of the \(\cal C(S_i)\), which may be interpreted in terms of
the fundamental groups of the \(S_i\) once a geometric point has been
chosen in each~\(S_i\). In applying descent theory to étale morphisms,
it is sometimes inconvenient to choose a geometric point in every
\(S_i\). It is then preferable to use the natural generalization
of~\Ref{V.5.8}, interpreting \(\cal C(S)\) as a category of functors
on the groupoid of geometric points of~\(S\), regarded as the sum of
the groupoids corresponding to the connected components of~\(S\).
The functors in question take values in the category of finite sets
and satisfy the continuity condition analogous to that used
in~\Ref{V.5.8}.

In practice, one has a family \((a_t)_{t\in E}\) of geometric points
of~\(S\) such that every connected component \(S_i\) of~\(S\) contains
at least one of them. As in~\Ref{V.5.10}, one may then replace the
groupoid of all geometric points of~\(S\) by the analogous groupoid
whose underlying set is~\(E\). These considerations should of course
be incorporated into general definitions concerning categories
equivalent to product categories of categories of the form
\(\cal C(\pi)\), which may be called \emph{multi-Galois categories}.
%
\index{multi-Galois category|hyperpage}%
%
We leave the details to the reader.

\end{filecontents*}
% Original input: drafts/SGA1_V_OPENING_THROUGH_1_4_English.texfrag
\begin{filecontents*}[overwrite,noheader]{sga1-cumulative-input-096.tex}
% SGA 1, Exposé V, French TeX lines 7383--7553.

\chapter{The fundamental group: generalities}
\label{V}
\marginpar{\raggedright\footnotesize French source p.~105}%

\setcounter{section}{-1}
\section{Introduction}

The present Seminar is the continuation of the 1960 Seminar. We shall
refer to the latter by abbreviations such as \ifthenelse{\boolean{orig}}{I~}{}\ref{I.9.7}, meaning:
Seminar on Algebraic Geometry, \hyperref[I]{Exposé~I},
\hyperref[I.9.7]{no.~9.7}. The numbers of the
1961 exposés will follow those of 1960. We shall refer to the
\emph{Éléments de Géométrie Algébrique} of Dieudonné and Grothendieck
by abbreviations such as EGA~I, 8.7.3.

The present exposé summarizes, with slight additions, the final exposés
of 1960, which had not been written up.

As in 1961, we shall as a rule restrict ourselves to locally
Noetherian preschemes, although this restriction is often inessential.
In Exposé~\ref{VI} we shall assume the theory of faithfully flat
descent, summarized in Bourbaki Seminar no.~190. If appropriate, we
shall give a more detailed account in a later exposé,
\footnote{See Exposés~\ref{VI} and~\ref{VIII}.}
once the reader has had occasion to appreciate the usefulness of this
technique in the theory of the fundamental group.

\section{Preschemes with a finite group of operators; quotient preschemes}
\label{V.1}

Let \(X\) be a prescheme and let a finite group \(G\) act on \(X\) by
automorphisms, on the right for definiteness. If \(X\) is affine with
ring \(A\), then \(G\) acts on \(A\) by automorphisms on the left.

For every prescheme \(Z\), the group \(G\) acts on the left on the set
\(\operatorname{Hom}(X,Z)\). We may therefore consider the set
\[
  \operatorname{Hom}(X,Z)^G
\]
of \(G\)-invariant morphisms. It depends functorially on \(Z\), and we
may ask whether this functor is representable, that is, isomorphic to
a functor \(Z\mapsto\operatorname{Hom}(Y,Z)\).
\marginpar{\raggedright\footnotesize French source p.~106}%
This means finding a prescheme \(Y\) and a \(G\)-invariant morphism
\[
  p\colon X\longrightarrow Y
\]
such that, for every \(Z\), the corresponding map \(g\mapsto gp\),
\[
  \operatorname{Hom}(Y,Z)\longrightarrow
  \operatorname{Hom}(X,Z)^G,
\]
is bijective. The pair \((Y,p)\) is then called a \emph{quotient
prescheme} of \(X\) by \(G\); it is determined up to unique
isomorphism.

\begin{proposition}
\label{V.1.1}
Let \(A\) be a ring on which the finite group \(G\) acts on the left,
let \(B=A^G\) be the subring of invariants, let
\(X=\operatorname{Spec}(A)\) and \(Y=\operatorname{Spec}(B)\), and let
\(p\colon X\to Y\) be the canonical morphism, which is plainly
\(G\)-invariant. Then:
\begin{enumerate}
\item[(i)] \(A\) is integral over \(B\); equivalently, \(p\) is an
integral morphism.
\item[(ii)] The morphism \(p\) is surjective, its fibers are the
\(G\)-orbits, and the topology of \(Y\) is the quotient topology of
that of \(X\).
\item[(iii)] Let \(x\in X\), let \(y=p(x)\), and let \(G_x\) be the
stabilizer of \(x\). Then \(\kappa(x)\) is a quasi-Galois algebraic
extension of \(\kappa(y)\), and the canonical map
\[
  G_x\longrightarrow
  \operatorname{Gal}\bigl(\kappa(x)/\kappa(y)\bigr)
\]
to the group of \(\kappa(y)\)-automorphisms of \(\kappa(x)\) is
surjective.
\item[(iv)] The pair \((Y,p)\) is a quotient prescheme of \(X\) by
\(G\).
\end{enumerate}
\end{proposition}

Statements \textup{(i)}, \textup{(ii)}, and \textup{(iii)} are
well-known in commutative algebra,
\footnote{See N.~Bourbaki, \emph{Commutative Algebra}, Chapter~V,
§§1--2, Theorem~2.}
and are recalled only for reference, except for the assertion about
the topology. The latter follows from this general fact, an easy
consequence of the Cohen--Seidenberg theorem: an integral morphism is
closed, that is, it maps closed sets to closed sets. We note at once:

\begin{corollary}
\label{V.1.2}
Under the preceding hypotheses, the natural homomorphism
\[
  \mathcal O_Y\longrightarrow p_*(\mathcal O_X)^G
\]
is an isomorphism.
\end{corollary}

This follows immediately from the formula
\[
  (S^{-1}A)^G=S^{-1}(A^G),
\]
valid for every multiplicatively closed subset \(S\) of \(B=A^G\).
The formula has a module version, and more generally holds after a
flat base change \(A\to A'\). Apply it when \(S\) is generated by one
element \(f\in B\).

Assertion \textup{(ii)} and Corollary~\ref{V.1.2} readily imply
\textup{(iv)}. More generally:

\begin{proposition}
\label{V.1.3}
\marginpar{\raggedright\footnotesize French source p.~107}%
Let \(X\) be a prescheme with a finite group \(G\) of automorphisms,
and let \(p\colon X\to Y\) be an invariant affine morphism such that
\[
  \mathcal O_Y\xrightarrow{\sim}p_*(\mathcal O_X)^G.
\]
Then conclusions \textup{(i)}, \textup{(ii)}, \textup{(iii)}, and
\textup{(iv)} of Proposition~\ref{V.1.1} still hold.
\end{proposition}

For \textup{(i)}, \textup{(ii)}, and \textup{(iii)}, we may suppose
that \(Y\), and hence \(X\), is affine. If \(B\) and \(A\) are their
rings, the hypothesis implies \(B=A^G\), so Proposition~\ref{V.1.1}
applies. For \textup{(iv)}, use \textup{(ii)} together with
\(\mathcal O_Y=p_*(\mathcal O_X)^G\).

\begin{corollary}
\label{V.1.4}
Under the hypotheses of Proposition~\ref{V.1.3}, for every open
subset \(U\subseteq Y\), the prescheme \(U\) is a quotient of
\(X|_U=p^{-1}(U)\) by \(G\).
\end{corollary}

Indeed, the morphism \(p^{-1}(U)\to U\) induced by \(p\) satisfies the
same hypotheses as \(p\).

Now suppose that \(X\) is a \(Z\)-prescheme and that the operations of
\(G\) are \(Z\)-automorphisms. By \textup{(iv)}, \(Y\) is then a
\(Z\)-prescheme. Thus:
\end{filecontents*}
% Original input: drafts/SGA1_XIII_1_10_AND_1_11_English.texfrag
\begin{filecontents*}[overwrite,noheader]{sga1-cumulative-input-097.tex}
\begin{definitions}
\label{XIII.1.10}

\subsubsection{}
\label{XIII.1.10.1}
Let $E$ be a category, and consider a diagram
$$
\xymatrix@C=.6cm{ \Phi \ar[r]^p & \Phi_1 \ar@<2pt>[r]^{p_1}
\ar@<-2pt>[r]_{p_2} & \Phi_2 } \text{,}
% remark: same remark as above, the comma is
% misplaced.
$$
where $\Phi$, $\Phi_1$, $\Phi_2$ are categories fibered
over~$E$, the arrows are morphisms of fibered categories, and
$$
a \colon p_1 p \isomto p_2 p
$$
is an isomorphism of functors.

We say that the diagram above is exact if the following condition
is satisfied:

a) For every pair of objects $x$, $y$ of~$\Phi$ and every morphism $f
\colon p(x) \to p(y)$ such that $p_1(f) = p_2(f)$ ($p_1p$ and
$p_2p$ being identified by means of $a$), there exists a
unique morphism $g \colon x \to y$ such that $p(g)=f$.

\subsubsection{}
\label{XIII.1.10.2}
Consider the diagram
\ifthenelse{\boolean{orig}}
{$$
\xymatrix{ \Phi \ar[r]^p & \Phi_1 \ar@<0.6ex>[r]^{p_1}
\ar@<-0.6ex>[r]_{p_2} & \Phi_2 \ar@<2ex>[r]^{p_{23}} \ar[r]^{p_{31}}
\ar@<-0.6ex>[r]_{p_{12}} & \Phi_3 } \text{,}
% remark: the comma is misplaced.
$$}
{$$
\xymatrix@C=.5cm{ \Phi \ar[r]^p & \Phi_1 \ar@<2pt>[rr]^{p_1,p_2}
\ar@<-2pt>[rr] && \Phi_2 \ar@<4pt>[rrr]^{p_{12},p_{23},p_{31}} \ar[rrr]
\ar@<-4pt>[rrr] &&& \Phi_3 } \text{,}
% remark: the comma is misplaced.
$$}
where
% page: 358
\marginpar{358}
%
$\Phi$, $\Phi_i$, $1 \leqslant i \leqslant 3$, are categories
fibered over~$E$, and the arrows are morphisms of
fibered categories. Suppose that we are given isomorphisms of
functors
$$
a \colon p_1 p \isomto p_2 p
$$
$$
a_1 \colon p_{31} p_2 \isomto p_{12} p_1 \text{,}\quad a_2 \colon p_{12}
p_2 \isomto p_{23} p_1 \text{,}\quad a_3 \colon p_{23} p_2 \isomto
p_{31}p_1
$$
such that the following diagram commutes:
$$
\xymatrix@C=1.2cm{ p_{23} p_1 p \ar[r]^{\id.a} & p_{23} p_2 p \ar[r]^{a_3.\id}
& p_{31} p_1 p \ar[d]^{\id.a} \\ p_{12} p_2 p \ar[u]^{a_2.\id} &
p_{12} p_1 p \ar[l]^{\id.a} & p_{31} p_2 p \ar[l]^{a_1.\id} } \text{.}
% remark: the period is misplaced.
$$
We identify $p_1 p$ and $p_2 p$, $p_{31} p_2$ and $p_{12} p_1$, and so on.

We say that the diagram above is exact if the following conditions
are satisfied:
\begin{enumerate}
\item[a)] The analogue of condition a) of~\Ref{XIII.1.10.1}.
\item[b)] For every object $x_1$ of~$\Phi_1$ and every isomorphism
$u \colon p_1(x_1) \isomto p_2(x_1)$ such that
\begin{equation*}
\label{eq:XIII.1.10.2.1}
\tag{\thesubsubsection.1} { p_{31}(u) p_{23}(u) = p_{12}(u)^{-1}
\text{,} }
\end{equation*}
there exists an object $x$ of~$\Phi$ and an isomorphism $i
\colon p(x) \isomto x_1$ making the diagram
\begin{equation*}
\label{eq:XIII.1.10.2.2}
\tag{\thesubsubsection.2}
\begin{array}{c}
{ \xymatrix{ p_1p(x) \ar@{=}[r]
\ar[d]_{p_1(i)} & p_2p(x) \ar[d]^{p_2(i)} \\ p_1(x_1) \ar[r]^{\sim}_u
& p_2(x_1)
% remark: I have taken the liberty of placing the u below the
% arrow...
} }
\end{array}
\end{equation*}
commute.
\end{enumerate}

\subsubsection{}
\label{XIII.1.10.3}
The notion of a morphism of exact diagrams of fibered categories over a
category $E$ is defined in the evident way.
\end{definitions}

\subsubsection{}
\label{XIII.1.10.4}
We shall use
% page: 359
\marginpar{359}%
%
in particular the notion of an exact diagram in the
case where $E$ is a site and $\Phi$, $\Phi_i$, $1 \leqslant i
\leqslant 3$, are stacks on~$E$.

Let $f \colon E \to E'$ be a morphism of sites and let
\ifthenelse{\boolean{orig}}
{\begin{equation*}
\label{eq:XIII.1.10.4.1}
\tag{\thesubsubsection.1} { \xymatrix{ \Phi \ar[r]^p & \Phi_1
\ar@<0.6ex>[r]^{p_1} \ar@<-0.6ex>[r]_{p_2} & \Phi_2
\ar@<2ex>[r]^{p_{23}} \ar[r]^{p_{31}} \ar@<-0.6ex>[r]_{p_{12}} &
\Phi_3 } }
\end{equation*}}
{\begin{equation*}
\label{eq:XIII.1.10.4.1}
\tag{\thesubsubsection.1} { \xymatrix@C=.5cm{ \Phi \ar[r]^p & \Phi_1
\ar@<2pt>[rr]^{p_1,p_2} \ar@<-2pt>[rr] && \Phi_2
\ar@<4pt>[rrr]^{p_{12},p_{23},p_{31}} \ar[rrr] \ar@<-4pt>[rrr] &&&
\Phi_3 } }
\end{equation*}}
be an exact diagram of stacks on~$E$. Taking direct images gives a
diagram
$$
\xymatrix@C=.5cm{ f_* \Phi \ar[r] & f_* \Phi_1 \ar@<2pt>[r]
\ar@<-2pt>[r] & f_* \Phi_2 \ar@<4pt>[r]\ar[r] \ar@<-4pt>[r] &
f_* \Phi_3 }
$$
which is evidently exact.

If $h \colon E'' \to E$ is a morphism of sites, one likewise has an
exact diagram
\ifthenelse{\boolean{orig}}
{$$
\xymatrix{ h^* \Phi \ar[r]^{p''} & h^* \Phi_1
\ar@<0.6ex>[r]^{p''_1} \ar@<-0.6ex>[r]_{p''_2} & h^* \Phi_2
\ar@<2ex>[r]^{p''_{23}} \ar[r]^{p''_{31}} \ar@<-0.6ex>[r]_{p''_{12}} &
h^* \Phi_3 }
$$}
{$$
\xymatrix@C=.5cm{ h^* \Phi \ar[r]^{p''} & h^* \Phi_1
\ar@<2pt>[rr]^{p''_1,p''_2} \ar@<-2pt>[rr] && h^* \Phi_2
\ar@<4pt>[rrr]^{p''_{12},p''_{23},p''_{31}} \ar[rrr] \ar@<-4pt>[rrr] &&&
h^* \Phi_3 }
$$}

Let us first verify condition~a) of~\Ref{XIII.1.10.2}. Let
$F''$ be an object of~$E''$, let $x''$ and $y''$ be two objects of~$(h^*
\Phi)_{F''}$, let $x''_1$ and $y''_1$ be their respective images in
$h^* \Phi_1$, and let $x''_2$, $y''_2$ be their images in $h^*
\Phi_2$. Let $u''_1 \colon x''_1 \to y''_1$ be a morphism such that
$p''_1(u''_1) = p''_2(u''_1)$, and let us prove that $u''_1$ comes
from a unique morphism $u'' \colon x'' \to y''$. Since the question is
local on~$F''$, we may suppose that we have an object $F_1$ of~$E$, a
morphism from~$F''$ to the inverse image $F''_1$ of~$F_1$ under $h$, and
objects $x$, $y$ of~$\Phi_{F_1}$ whose inverse images on~$F''$
are $x''$ and $y''$. Let $x_1$, $y_1$ (\resp $x_2$, $y_2$) be the
images of~$x$, $y$ in $\Phi_1$ (\resp~$\Phi_2$). We may suppose
that $u''_1$ comes from a morphism $u_1 \colon x_1 \to y_1$ such that
$p_1(u_1) = p_2(u_1)$. By the exactness of
\eqref{eq:XIII.1.10.4.1}, we obtain a unique morphism $u \colon x
\to y$ whose inverse image under $h$ is the desired morphism $u''$.

The
% page: 360
\marginpar{360}
%
condition~b) of~\Ref{XIII.1.10.2} is verified in the same
way. Let $x''_1$ be an object of~$(h^* \Phi_1)_{F''}$ and let $u''
\colon p''_1(x''_1) \to p''_2(x''_1)$ be a morphism satisfying the
relation
$$
p''_{31}(u'')p''_{23}(u'')=p''_{12}(u'')^{-1} \text{,}
$$
and let us prove that there exists an object $x''$ of~$(h^* \Phi)_{F''}$ and
an isomorphism $i'' \colon p''(x'') \simeq x''_1$ making
a diagram analogous to \eqref{eq:XIII.1.10.2.2} commute. Since the
question is local on~$F''$, we may suppose that we have an object
$F_1$, a morphism $F'' \to F''_1$ as above, and an object $x_1$
of~$(\Phi_1)_{F_1}$ whose inverse image in $(h^* \Phi_1)_{F''}$
is $x''_1$. Likewise, we may suppose that $u''$ comes from a
morphism $u \colon p_1(x_1) \to p_2(x_1)$ satisfying
\eqref{eq:XIII.1.10.2.1}. The existence of an object $x$ of~$\Phi_{F_1}$
whose inverse image under $h$ is an object $x''$ with the required
property then follows from the exactness of
\eqref{eq:XIII.1.10.4.1}.

\begin{exemples}
\label{XIII.1.11}
\begin{enumerate}
\item[1)] Let $f \colon X_1 \to X$ be a \emph{descent morphism}
for the category of \'etale sheaves on variable schemes
(for example, a universally submersive morphism (SGA~4
VIII~9.3)). Let $X_2 = X_1 \times_X X_1$, let $g \colon X_2 \to X$ be the
canonical projection, and let $F$ be a sheaf of sets on~$X$. It
then follows from \loccit that there is an exact sequence of
sheaves of sets
\begin{equation*}
\label{eq:XIII.1.11.1}
\tag{\thesubsection.1} { \xymatrix@C=.5cm{ F \ar[r] & f_* f^* F
\ar@<2pt>[r] \ar@<-2pt>[r] & g_* g^* F \text{.}}
% remark: the period is misplaced.
}
\end{equation*}
If $\Phi$ is the stack in discrete categories associated with
$F$, and $\Phi_3$ is the final stack on~$X$, \ie the stack all of
whose fibers consist of a single object
whose only morphism is the identity morphism, saying that the sequence
\eqref{eq:XIII.1.11.1} is exact is equivalent to saying that the
diagram of stacks
$$
\xymatrix@C=.5cm{ \Phi \ar[r] & f_* f^* \Phi \ar@<2pt>[r]
\ar@<-2pt>[r] & g_* g^* \Phi \ar[r] \ar@<4pt>[r]
\ar@<-4pt>[r] & \Phi_3 } \text{.}
% remark: the period is misplaced.
$$
is exact.

\item[2)]
Let
% page: 361
\marginpar{361}
%
$f \colon X_1 \to X$ be an \emph{effective descent morphism}
for the category of \'etale sheaves on variable schemes
(for example, a proper surjective morphism, an integral surjective
morphism, or a faithfully flat morphism locally of
finite presentation (SGA~4 VIII~9.4)). Let $X_2 = X_1 \times_X
X_1$, let $g \colon X_2 \to X$ be the canonical projection, let $X_3 = X_1
\times_X X_1 \times_X X_1$, and let $h \colon X_3 \to X$ be the
canonical morphism. Let $\Phi$ be a stack on~$X$, let $\Phi_1 = f_* f^* \Phi$,
$\Phi_2 = g_* g^* \Phi$, and $\Phi_3 = h_* h^* \Phi$. We then have an
exact diagram
$$
\xymatrix@C=.5cm{ \Phi \ar[r] & \Phi_1 \ar@<2pt>[r] \ar@<-2pt>[r] &
\Phi_2 \ar[r] \ar@<4pt>[r] \ar@<-4pt>[r] & \Phi_3 \text{,} }
$$
where the arrows are the canonical morphisms associated with the
projections.

Indeed, regard $\Phi$ as a stack on the category
$\Sch_X$ of schemes over~$X$, endowed with the \'etale topology.
Then, by \cite[VII 2.2.8]{XIII.2}, $\Phi$ is also a
stack for the finest topology on~$\Sch_X$ for which the
covering morphisms are the effective descent morphisms for the
category of \'etale sheaves. The exactness of the diagram
above follows immediately.
\end{enumerate}
\end{exemples}
\end{filecontents*}
% Original input: drafts/SGA1_XIII_1_12_AND_1_13_English.texfrag
\begin{filecontents*}[overwrite,noheader]{sga1-cumulative-input-098.tex}
\begin{proposition}
\label{XIII.1.12}
Let $S$ be a scheme, and let $f\colon X \to Y$ be an $S$-morphism.
\begin{enumerate}
\item[1)] Let
$$
\xymatrix@C=.5cm{ \Phi \ar[r] & \Phi_1 \ar@<2pt>[r] \ar@<-2pt>[r] &
\Phi_2 }
$$
be an exact diagram of stacks on~$X$. Suppose that $(\Phi_1, f)$ is
cohomologically proper relative to~$S$ in dimension
$\leqslant 0$ and that $(\Phi_2, f)$ is cohomologically proper relative to~$S$ in dimension $\leqslant -1$. Then $(\Phi,f)$
is cohomologically proper relative to~$S$ in dimension
$\leqslant 0$.

\item[2)] Let
$$
\xymatrix@C=.5cm{ \Phi \ar[r] & \Phi_1 \ar@<2pt>[r] \ar@<-2pt>[r] &
\Phi_2 \ar@<4pt>[r] \ar[r] \ar@<-4pt>[r] & \Phi_3 }
$$
be an exact diagram of stacks on~$X$. Suppose that $(\Phi_1,f)$ is
cohomologically proper relative to~$S$ in dimension
$\leqslant 1$, that $(\Phi_2, f)$ is cohomologically proper
% page: 362
\marginpar{362}
%
relative to~$S$ in dimension $\leqslant 0$, and that
$(\Phi_3, f)$ is cohomologically proper
\ifthenelse{\boolean{orig}}
{}
{relative}
% correction: `relative' is missing.
to~$S$ in dimension $\leqslant -1$. Then $(\Phi,f)$ is
cohomologically proper relative to~$S$ in dimension
$\leqslant 1$.
\end{enumerate}
\end{proposition}

For every $S$-scheme $S'$, consider the following commutative
diagram, all of whose squares are cartesian:
$$
\xymatrix{ X' \ar[r]^{f'} \ar[d]_h & Y' \ar[r] \ar[d]_g & S' \ar[d] \\
X \ar[r]^f & Y \ar[r] & S }
$$

Let us prove 2), the proof of 1) being analogous. Since the direct
and inverse image functors take exact diagrams of stacks to exact
diagrams \eqref{XIII.1.10.4}, there is the following morphism of exact
diagrams of stacks:
$$
\xymatrix{ g^* f_* \Phi \ar[r]^\pi \ar[d]_{\varphi} & g^*
f_* \Phi_1 \ar@<2pt>[r]^{\pi_1} \ar@<-2pt>[r]_{\pi_2}
\ar[d]_{\varphi_1} & g^* f_* \Phi_2 \ar@<4pt>[r] \ar[r]
\ar@<-4pt>[r] \ar[d]_{\varphi_2} & g^* f_* \Phi_3
\ar[d]_{\varphi_3} \\ f'_* h^* \Phi \ar[r] & f'_* h^*
\Phi_1 \ar@<2pt>[r] \ar@<-2pt>[r] & f'_* h^* \Phi_2
\ar@<4pt>[r] \ar[r] \ar@<-4pt>[r] & f'_* h^* \Phi_3 }
% remark: misplaced period.
\text{.}
$$
By hypothesis, $\varphi_1$ is an equivalence of categories,
$\varphi_2$ is fully faithful, and $\varphi_3$ is faithful. It
therefore follows from the preceding diagram that $\varphi$ is an
equivalence.

\begin{proposition}
\label{XIII.1.13}
Let $f \colon X \to Y$ be an $S$-morphism.
\begin{enumerate}
\item[1)] Let
$$
\xymatrix@C=.5cm{ F \ar[r] & G \ar@<2pt>[r] \ar@<-2pt>[r] & H }
$$
be an exact diagram of sheaves of sets on~$X$. Suppose that
$(G,f)$ is cohomologically proper relative to~$S$ in
dimension $\leqslant 0$ and that $(H,f)$ is cohomologically proper
relative to~$S$ in dimension $\leqslant -1$. Then $(F,f)$ is
cohomologically proper relative to~$S$ in dimension
$\leqslant 0$.
\item[2)]
Let
% page: 363
\marginpar{363}
%
$F \to G$ be a monomorphism of sheaves of groups on~$X$. If
$Y_1$ is a scheme \`etale over~$Y$, put $X_1 = Y_1 \times_Y
X$, and denote by $f_1$ (\resp $F_1$, \resp $G_1$) the inverse image of~$f$
(\resp $F$, \resp $G$) over~$Y_1$ (\cf \Ref{eq:XIII.1.0.2}). Suppose
that $(G,f)$ is cohomologically proper relative to~$S$ in
dimension $\leqslant 0$ (\resp in dimension $\leqslant 1$) and that,
for every scheme $Y_1$ \`etale over~$Y$ and every torsor $Q$
under $G_1$, $(Q/F_1, f_1)$ is cohomologically proper relative
to~$S$ in dimension $\leqslant -1$ (\resp in dimension $\leqslant
0$). Then $(F,f)$ is cohomologically proper relative to~$S$
in dimension $\leqslant 0$ (\resp in dimension $\leqslant 1$).
\item[3)] Let $F \to G$ be a monomorphism of sheaves of groups on~$X$. Suppose that $(F,f)$ is cohomologically proper relative
to~$S$ in dimension $\leqslant 1$ and that $(G,f)$ is
cohomologically proper relative to~$S$ in dimension
$\leqslant 0$. Then, for every torsor $Q$ under $G$, $(Q/F, f)$ is
cohomologically proper relative to~$S$ in dimension
$\leqslant 0$.
\end{enumerate}
\end{proposition}

\subsubsection*{Proof}
1) Let $\Phi$ (\resp $\Phi_1$, \resp $\Phi_2$) be the stack in
discrete categories associated with $F$ (\resp $G$,
\resp $H$), and let $\Phi_3$ be the final stack on~$X$. There is then
an exact diagram
$$
\xymatrix@C=.5cm{ \Phi \ar[r] & \Phi_1 \ar@<2pt>[r] \ar@<-2pt>[r] &
\Phi_2 \ar@<4pt>[r] \ar[r] \ar@<-4pt>[r] & \Phi_3 }
% remark: misplaced period.
\text{.}
$$
By hypothesis, $(\Phi_1, f)$ is cohomologically proper
relative to~$S$ in dimension $\leqslant 1$, and $(\Phi_2, f)$ is
cohomologically proper relative to~$S$ in dimension
$\leqslant 0$ \eqref{XIII.1.2}. Since $(\Phi_3, f)$ is clearly
cohomologically proper relative to~$S$ in dimension
$\leqslant -1$, it follows from~\Ref{XIII.1.12} that $(\Phi, f)$ is
cohomologically proper relative to~$S$ in dimension
$\leqslant 1$, \ie that $(F,f)$ is cohomologically proper
relative to~$S$ in dimension $\leqslant 0$.

\smallskip
2) Let us first show that, if $(G,f)$ is cohomologically
proper relative to~$S$ in dimension $\leqslant 0$ and if the
pairs $(Q/F_1, f_1)$ are cohomologically
% page: 364
\marginpar{364}
%
proper relative to~$S$ in dimension $\leqslant -1$, then
$(F,f)$ is cohomologically proper relative to~$S$ in
dimension $\leqslant 0$. By~\Ref{XIII.1.3.1}, it suffices to
prove that, for every scheme $Y_1$ \`etale over~$Y$ and
every torsor $P$ on~$X_1$ with group $F_1$, $(\leftexp{P\mkern-4mu}{F}_1,
f_1)$ is cohomologically proper relative to~$S$ in dimension
$\leqslant 0$ when $\leftexp{P\mkern-4mu}{F}_1$ is regarded as a
sheaf of sets, and that the canonical morphism
$$
d \colon g^*(\R^1 f_* F) \to \R^1 f'_* F'
$$
is injective. The first assertion follows at once from 1),
for if $Q$ denotes the torsor obtained from~$P$ by extension
of structure group along $F_1 \to G_1$, there is an isomorphism
$\leftexp{Q}{G}_1 / \leftexp{P\mkern-4mu}{F}_1 \isomto Q/F_1$.

Let us show that $d$ is injective. It suffices to prove that, if $Y_1$ is
a scheme \`etale over~$Y$, and if $P$ and $\tilde{P}$ are
two torsors under $F_1$ whose inverse images $P'$ and~$\tilde{P}'$
on~$X'_1$ are isomorphic, then, after replacing $Y_1$ by a
surjective \`etale extension, $P$ and~$\tilde{P}$ become
isomorphic. Choose an isomorphism $p' \colon P' \isomto
\tilde{P}'$. Let $Q$ (\resp~$\tilde{Q}$) be the torsor obtained from
$P$ (\resp $\tilde{P}$) by extension of structure group along $F_1 \to
G_1$. The inverse images $Q'$ (\resp $\tilde{Q}'$) of~$Q$
(\resp $\tilde{Q}$) on~$X'_1$ are obtained from~$P'$
(\resp $\tilde{P}'$) by
\ifthenelse{\boolean{orig}}
{extension}
{the extension}
% correction: `the' is missing.
of the structure group $F'_1 \to G'_1$; let $q' \colon Q' \isomto
\tilde{Q}'$ be the isomorphism obtained in the same way from~$p'$.
Since $(G,f)$ is cohomologically proper relative to~$S$ in dimension $\leqslant 0$, after replacing $Y_1$ by a surjective
\`etale extension, we may suppose that $q'$ is the inverse image of an
isomorphism $q \colon Q \isomto \tilde{Q}$. Associated with the torsor~$P$
(\resp~$\tilde{P}$) is a section $x$ of~$Q/F_1$
(\resp a section $\tilde{x}$ of~$\tilde{Q}/F_1$), and, for $P$
and $\tilde{P}$ to be isomorphic, it is necessary and sufficient that there be an
isomorphism $Q \isomto \tilde{Q}$ such that the induced isomorphism
$$
e \colon \H^0(X_1, Q/F_1) \to \H^0(X_1,\tilde{Q}/F_1)
$$
carries $x$ to $\tilde{x}$. Take the isomorphism $q$. The sections
$e(x)$ and $\tilde{x}$ of
$\H^0(X_1,\tilde{Q}/F_1)$ have the same image in
$\H^0(X'_1,\tilde{Q}'/F'_1)$. Since $(\tilde{Q}/F_1,f_1)$ is
cohomologically proper relative to~$S$ in dimension
$\leqslant -1$,
% page: 365
\marginpar{365}
%
after replacing $Y_1$ by a surjective \`etale extension, we indeed have
$e(x) = \tilde{x}$, which proves the injectivity of~$d$.

To complete the proof, it remains to show that, if
$(G,f)$ is cohomologically proper relative to~$S$ in
dimension $\leqslant 1$, and if, for every scheme $Y_1$ \`etale
over~$Y$ and every torsor $Q$ on~$X_1$ with group $G_1$,
$(Q/F_1, f_1)$ is cohomologically proper relative to~$S$ in
dimension $\leqslant 0$, then the morphism $d$ is surjective. Let
$P'$ be a torsor on~$X'_1$ with group $F'_1$, and let $Q'$ be the torsor under
$G'_1$ obtained from~$P'$ by extension of structure
group. The data of~$P'$ are equivalent to those of
$Q'$ together with a section $x'$ of~$\H^0(X'_1,Q'/F'_1)$. It follows
from the surjectivity of the morphism
$$
g^*(\R^1 f_* G) \to \R^1 f'_* G'
$$
that, after replacing $Y_1$ by a surjective \`etale extension,
there is a torsor $Q$ under~$G_1$ whose inverse image on~$X'_1$
is isomorphic to $Q'$. Using the fact that $(Q/F_1, f_1)$ is
cohomologically proper relative to~$S$ in dimension
$\leqslant 0$, we may likewise suppose that there is an
element $x$ of~$\H^0(X_1, Q/F_1)$ whose image in
$\H^0(X'_1, Q'/F'_1)$ is $x'$. Together, $Q$ and $x$
determine a torsor $P$ under $F_1$ whose inverse image on~$X'_1$ is isomorphic to $P'$, which proves the
surjectivity of~$d$.

\smallskip
3) Let us show that $(Q/F,f)$ is cohomologically proper
relative to~$S$ in dimension $\leqslant -1$, \ie that, for
every $S$-scheme $S'$, and every scheme $Y_1$ \`etale
over~$Y$, if $x$, $\tilde{x}$ are two elements of
$\H^0(X_1,Q_1/F_1)$ whose images $x'$, $\tilde{x}'$ in
$\H^0(X'_1, Q'_1/F'_1)$ are equal, then, after a surjective
extension of~$Y_1$, we have $x = \tilde{x}$. Associated with~$x$ (\resp~$\tilde{x}$)
is a torsor $P$ (\resp~$\tilde{P}$) under $F_1$ such that
$Q_1$ is obtained from~$P$ (\resp~$\tilde{P}$) by extension of structure group along $F_1
\to G_1$. It follows from the relation $x' = \tilde{x}'$
that there is an isomorphism $u' \colon P' \isomto
\tilde{P}'$ such that the isomorphism induced on~$Q'_1$ by~$u'$ is
the identity. Since $(F,f)$ is cohomologically proper
relative to~$S$ in dimension $\leqslant 0$, it follows
that,
% page: 366
\marginpar{366}
%
after a surjective \`etale extension of~$Y_1$, there is an
isomorphism $u \colon P \to \tilde{P}$ lifting $u'$; the fact that
$(G,f)$ is cohomologically proper relative to~$S$ in
dimension $\leqslant -1$ then implies that $x =
\tilde{x}$.

Let us show that $(Q/F,f)$ is cohomologically proper relative to~$S$ in dimension $\leqslant 0$. Let $Y''$ be a scheme \`etale
over~$Y'$, and let $x''$ be an element of~$\H^0(X'',Q''/F'')$. Associated with~$x''$
is a torsor $P''$ on~$X''$ with group $F''$. Since
$(F,f)$ is cohomologically proper relative to~$S$ in
dimension $\leqslant 1$, one can find surjective \`etale
morphisms $Y''_1 \to Y''$ and $Y_1 \to Y$ such that there is a
morphism $Y''_1 \to Y'_1$, and a torsor $P$ on~$X_1$ with group
$F_1$ whose inverse image on~$X''_1$ is isomorphic to the inverse
image of~$P''$. It then follows from the fact that $(G,f)$ is
cohomologically proper relative to~$S$ in dimension
$\leqslant 0$ that one can even choose $Y''_1$ and $Y_1$ so that
the torsor obtained from~$P$ by extension of structure group along $F_1
\to G_1$ is isomorphic to $Q_1$; the torsor~$P$ corresponds to an
element $x$ of~$\H^0(X_1,Q_1/F_1)$ whose image in
$\H^0(X''_1,Q''_1/F''_1)$ is equal to the inverse image of
$x''$, which completes the proof.
\end{filecontents*}
% Original input: drafts/SGA1_XIII_1_14_THROUGH_1_16_English.texfrag
\begin{filecontents*}[overwrite,noheader]{sga1-cumulative-input-099.tex}
\begin{proposition}
\label{XIII.1.14}
Let $f \colon X \to S$ be the structure morphism of an $S$-scheme,
and let $F$ be a sheaf
of sets or groups on~$X$ (\resp a sheaf of
ind-$\LL$-groups, where $\LL$ is a set of prime
numbers). Suppose that $F$ is locally constant, that $(F,f)$ is
cohomologically proper in dimension $\leqslant 0$ (\resp in
dimension $\leqslant 1$), and that $f$ is locally $0$-acyclic
(\resp locally $1$-aspherical for $\LL$) \textup{(SGA~4 XV~1.11)}.
Then, for every specialization $\overline{s}_1 \to
\overline{s}_2$ of geometric points of~$S$, the specialization
morphism \textup{(SGA~4 VIII~7.1)}
$$
a_0 \colon (f_* F)_{\overline{s}_2} \to (f_*
F)_{\overline{s}_1}
$$
is an isomorphism, and, if $F$ is a sheaf of groups, the morphism
$$
a_1 \colon (\R^1 f_* F)_{\overline{s}_2} \to (\R^1 f_*
F)_{\overline{s}_1}
$$
is injective (\resp the morphisms $a_0$ and $a_1$ are
isomorphisms).
\end{proposition}

The
% page: 367
\marginpar{367}
%
proof is obtained by reproducing verbatim that of SGA~4
XVI~2.3, with the expression \og proper\fg replaced by the
expression \og cohomologically proper\fg.

\ifthenelse{\boolean{orig}}
{}
{\enlargethispage{.5cm}}

\begin{corollaire}
\label{XIII.1.15}
Let $f \colon X \to S$ be a morphism, let $\Phi$ be a stack on~$X$, and
let $\LL$ be a set of prime numbers. Suppose that, for every scheme
$X_1$ \'etale over~$X$ and every pair of objects $x$, $y$ of
$\Phi_{X_1}$, the sheaf $\SheafHom_{X_1}(x,y)$ is locally
constant, the sheaf $\SheafAut_{X_1}(x)$ is a locally constant
ind-$\LL$-group, and the sheaf $S\Phi$ of maximal subgerbes
of~$\Phi$ \textup{\cite[III 2.1.7]{XIII.1}} is locally constant.
Suppose that $(\Phi,f)$ is cohomologically proper in dimension
$\leqslant 1$ and that $f$ is locally $1$-aspherical for $\LL$.
Then, for every specialization $\overline{s}_1 \to
\overline{s}_2$ of geometric points of~$S$, the specialization
morphism
$$
a \colon (f_* \Phi)_{\overline{s}_2} \to (f_*
\Phi)_{\overline{s}_1}
$$
is an equivalence of categories.
\end{corollaire}

Let $\overline{S}_1$ (\resp $\overline{S}_2$) be the strict
localization of~$S$ at $\overline{s}_1$ (\resp the strict localization
of~$S$ at $\overline{s}_2$), and let $\overline{X}_2$,
$\overline{\Phi}_2$ (\resp $\overline{X}_1$, $\overline{\Phi}_1$) be the
inverse images of $X$, $\Phi$ on~$\overline{S}_2$ (\resp of~$X$,
$\Phi$ on~$\overline{S}_1$), and consider the cartesian square
$$
\xymatrix{ \overline{X}_1 \ar[r]^h \ar[d]_{\overline{f}_1} &
\overline{X}_2 \ar[d]^{\overline{f}_2} \\ \overline{S}_1 \ar[r]^g &
\overline{S}_2 }
% remark: the period is misplaced.
\text{.}
$$
We must show that the functor
$$
\varphi \colon \overline{\Phi}_2(\overline{X}_2) \to \overline{\Phi}_1
(\overline{X}_1)
$$
is an equivalence. The functor $\varphi$ is fully faithful;
\ifthenelse{\boolean{orig}}
{in}
{}
% correction: an extra `in'.
in
% page: 368
\marginpar{368}
%
fact, let $x$, $y$ be two objects of~$(\overline{\Phi}_2)_{\overline{X}_2}$;
the canonical morphism
$$
\Hom_{\overline{X}_2}(x,y) \to
\Hom_{\overline{X}_1}(\varphi(x),\varphi(y))
$$
is identified with the canonical morphism
$$
\H^0(\overline{X}_2, \SheafHom_{\overline{X}_2}(x,y)) \to
\H^0(\overline{X}_1,h^*(\SheafHom_{\overline{X}_2}(x,y))) \text{.}
$$
This morphism is an isomorphism by~\Ref{XIII.1.14}.

Let us show that $\varphi$ is an equivalence. Let $x_1$ be an object
of~$\overline{\Phi}_1(\overline{X}_1)$, and let $G_1$ be the maximal
subgerbe of~$\overline{\Phi}_1$ generated by $x_1$. The morphism
$$
\H^0(\overline{X}_2, S\overline{\Phi}_2) \to \H^0(\overline{X}_1,
h^*(S \overline{\Phi}_2)) =
\H^0(\overline{X}_1,S\overline{\Phi}_1)
$$
is bijective, and hence there exists a maximal subgerbe $G_2$ of
$\overline{\Phi}_2$ such that $h^* G_2$ is isomorphic to
$G_1$. It then suffices to prove that the functor
$$
G_2 \to h_* h^* G_2
$$
is an equivalence. In this form, however, the question is local for
the \'etale topology on~$\overline{X}_2$. We may therefore suppose
that $G_2$ is a gerbe of torsors under the automorphism group of an
object of~$G_2$, in which case the assertion follows
from~\Ref{XIII.1.14}.

\begin{corollaire}
\label{XIII.1.16}
The hypotheses are those of~\Ref{XIII.1.14}. Suppose in addition that
$F$ is a sheaf of sets (\resp of ind-$\LL$-groups) and that
$f_* F$ (\resp $\R^1 f_* F$) is constructible. Then
$f_* F$ (\resp $\R^1 f_* F$) is locally constant.
\end{corollaire}

The corollary follows from~\Ref{XIII.1.14} by SGA~4
IX~2.11.

\begin{remarque}
\label{XIII.1.17}
Recall that the condition that $f$ be locally $0$-acyclic is satisfied
if $f$ is flat with separable fibers and $X$ and $S$ are locally
Noetherian (SGA~4 XV~4.1), and that the condition that $f$ be locally
$1$-aspherical
% page: 369
\marginpar{369}
%
for $\LL$ is satisfied if $f$ is smooth, with $\LL$ the set
of prime numbers distinct from the residual characteristics
of~$S$ (SGA~4 XV~2.1).
\end{remarque}
\end{filecontents*}
% Original input: drafts/SGA1_XIII_1_5_AND_1_6_English.texfrag
\begin{filecontents*}[overwrite,noheader]{sga1-cumulative-input-100.tex}
\begin{remarques}
\label{XIII.1.5}
a) Let \(F\) be a sheaf of groups on~\(X\) such that \((F,f)\) is
cohomologically proper relative to~\(S\) in dimension
\(\leqslant-1\) (\resp \(\leqslant0\)). If \(F\) is regarded as a
sheaf of sets, \((F,f)\) is cohomologically proper relative to~\(S\)
in dimension \(\leqslant-1\) (\resp \(\leqslant0\)), but the
converse is false.

For example, let \(Y\) be the spectrum of a strictly local discrete
valuation ring, with closed point \(t\) and generic point \(s\);
let \(f\colon X\to Y\) be a nonempty scheme over~\(Y\) whose closed
fiber is empty; let \(F\) be a nontrivial constant sheaf of groups
on~\(X\); and let \(P\) be a torsor under \(F\) such that
\(\H^0(X_s,\leftexp{P\mkern-4mu}{F}|X_s)=1\). Then
\((\leftexp{P\mkern-4mu}{F},f)\) is cohomologically proper relative
to~\(Y\) in dimension \(\leqslant-1\) when
\(\leftexp{P\mkern-4mu}{F}\) is regarded as a sheaf of sets. If
\(\leftexp{P\mkern-4mu}{F}\) is regarded as a sheaf of groups, there
is an isomorphism
\({\vphantom{F}}^{P^\circ\mkern-7mu}(\leftexp{P\mkern-4mu}{F})
\simeq F\); since the canonical morphism
$$
\H^0(X,F)\to\H^0(X_t,F|X_t)=1
$$
is not injective, this proves that
\((\leftexp{P\mkern-4mu}{F},f)\) is not cohomologically proper
relative to~\(Y\) in dimension \(\leqslant-1\).


b) Suppose that \(f\) is coherent (\ie quasi-compact and
quasi-separated). Let \(F\) be a stack on~\(X\). For every geometric
point \(\overline{y}\) of~\(Y'\), denote by \(\overline{Y}\)
(\resp \(\overline{Y'}\)) the strict localization of~\(Y\)
(\resp \(Y'\)) at \(\overline{y}\), and put
\(\overline{X}=X\times_Y\overline{Y}\),
\(\overline{X'}=X'\times_{Y'}\overline{Y'}\), etc. The pair \((F,f)\)
is cohomologically proper relative to~\(S\) in dimension
\(\leqslant-1\) (\resp \(\leqslant0\), \resp \(\leqslant1\)) if and
only if, for every \(S\)-scheme \(S'\) and every geometric point
\(\overline{y}\) of~\(Y'\), the canonical functor
$$
\overline{F}(\overline{X})\to\overline{F'}(\overline{X'})
$$
is faithful (\resp fully faithful, \resp an equivalence).

Indeed,
% page: 353
\marginpar{353}
%
if \(S'\) is an \(S\)-scheme, the functor
$$
g^* f_* F \to f'_* F'
$$
is faithful (\resp fully faithful, \resp an equivalence) if and only
if the same is true of the functor induced on the fibers at the
various geometric points \(\overline{y}'\) of~\(\overline{Y'}\)
\cite[III 2.1.5.9]{XIII.2}. The assertion therefore follows from the
calculation of the geometric fibers of the direct image of a stack
under a coherent morphism \cite[VII 2.1.5]{XIII.2}.


c) Let \(F\) be a stack on~\(X\). The property that \((F,f)\) is
cohomologically proper relative to~\(S\) in dimension
\(\leqslant-1\) (\resp \(\leqslant0\), \resp \(\leqslant1\)) is
local on~\(Y\) for the étale topology.

Let \(S'\) be an \(S\)-scheme, and let \(F'\) be the inverse image
of~\(F\) on~\(X'\) (\cf \eqref{eq:XIII.1.0.1}). If \((F,f)\) is
cohomologically proper relative to~\(S\) in dimension
\(\leqslant1\), then the same is true of~\((F',f')\). But if
\((F,f)\) is cohomologically proper relative to~\(S\) in dimension
\(\leqslant-1\) (\resp \(\leqslant0\)), the same need not be true
of~\((F',f')\).

For example, let \(S'\) be a discrete valuation ring, let
\(f'\colon E_{S'}\to S'\) be affine space over~\(S'\), let \(x\) be
a closed point of~\(E_{S'}\) over the generic point of~\(S'\), and
let \(F'\) be the sheaf of sets on~\(E_{S'}\) whose restriction to
\(E_{S'}-\{x\}\) is the one-element constant sheaf and whose fiber
at a geometric point over~\(x\) has two elements. Then
\((F',f')\) is not cohomologically proper relative to~\(S'\) in
dimension \(\leqslant-1\). Let \(S=S'[Z]\), let
\(f\colon E_S\to S\) be affine space over~\(S\), and let \(T\) be a
closed subset of \(X=E_S\) that does not meet the closed subset
\(Z=0\) and such that \(f(T)\) contains the generic point of~\(S\).
Let \(G\) be the inverse image of~\(F'\) on~\(X\), and let \(F\) be
the sheaf on~\(X\) obtained by extending \(G|X-T\) by the empty
sheaf. Then \((F,f)\) is cohomologically proper relative to~\(S\) in
dimension \(\leqslant-1\), but this is no longer true after the base
change \(S'\to S\) defined by \(Z=0\).


d)
% page: 354
\marginpar{354}
%
Let \(F\) be a stack on~\(X\) such that \((F,f)\) is
cohomologically proper relative to~\(Y\) in dimension
\(\leqslant-1\) (\resp \(\leqslant0\), \resp \(\leqslant1\)).
Then, for every geometric point \(\overline{y}\) of~\(Y\), the
canonical functor
$$
(f_* F)_{\overline{y}} \to F(X_{\overline{y}})
$$
is faithful (\resp fully faithful, \resp an equivalence of
categories).
\end{remarques}

\begin{proposition}
\label{XIII.1.6}
Let \(f\colon X\to Y\) and \(g\colon Y\to Z\) be two
\(S\)-morphisms, and let \(\Phi\) be a stack on~\(X\).
\begin{enumerate}
\item[1)] Suppose that \((\Phi,f)\) and \((f_*\Phi,g)\) are
cohomologically proper relative to~\(S\) in dimension
\(\leqslant-1\) (\resp \(\leqslant0\), \resp \(\leqslant1\)).
Then the same is true of~\((\Phi,gf)\).
\item[2)] Suppose that \((\Phi,gf)\) is cohomologically proper
relative to~\(S\) in dimension \(\leqslant-1\) (\resp that
\((\Phi,gf)\) is cohomologically proper relative to~\(S\) in
dimension \(\leqslant0\) and \((\Phi,f)\) is cohomologically proper
relative to~\(S\) in dimension \(\leqslant-1\), \resp that
\((\Phi,gf)\) is cohomologically proper relative to~\(S\) in
dimension \(\leqslant1\) and \((\Phi,f)\) is cohomologically proper
relative to~\(S\) in dimension \(\leqslant0\)). Then
\((f_*\Phi,g)\) is cohomologically proper relative to~\(S\) in
dimension \(\leqslant-1\) (\resp in dimension \(\leqslant0\),
\resp in dimension \(\leqslant1\)).
\end{enumerate}
\end{proposition}

For every \(S\)-scheme \(S'\), consider the following diagram, all of
whose squares are cartesian:
\begin{equation*}
\label{eq:XIII.1.6.1}
\tag{\thesubsection.1}
\begin{array}{c}
\xymatrix{ X' \ar[r]^{f'} \ar[d]_h & Y'
\ar[r]^{g'} \ar[d]_k & Z' \ar[r] \ar[d]_{m} & S' \ar[d] \\ X \ar[r]^f
& Y \ar[r]^g & Z \ar[r] & \,S. }
\end{array}
\end{equation*}
\ifthenelse{\boolean{orig}}
{the}
{The}
% correction: initial capital letter
canonical morphism
$$
m^* (g_* f_* \Phi) \to g'_* f'_* (h^* \Phi)
$$
is identified
% page: 355
\marginpar{355}
%
with the composite of the canonical morphisms
$$
m^*(g_* f_* \Phi) \lto{i} g'_*(k^*f_* \Phi) \lto{j} g'_* f'_*(h^* \Phi).
$$
\begin{enumerate}
\item[1)] The hypothesis implies that \(i\) and \(j\) are faithful
(\resp fully faithful, \resp equivalences); hence the same is true
of~\(ji\).
\item[2)] The hypothesis implies that \(ji\) is faithful
(\resp that \(ji\) is
\ifthenelse{\boolean{orig}}
{fuly}
{fully}
% correction: fully instead of fuly
faithful and \(j\) is faithful, \resp that \(ji\) is an equivalence
and \(j\) is fully faithful); it follows that \(i\) is faithful
(\resp fully faithful, \resp an equivalence).
\end{enumerate}

\end{filecontents*}
% Original input: drafts/SGA1_XIII_1_7_THROUGH_1_9_English.texfrag
\begin{filecontents*}[overwrite,noheader]{sga1-cumulative-input-101.tex}
\begin{corollaire}
\label{XIII.1.7}
Let \(f\colon X\to Y\) and \(g\colon Y\to Z\) be two
\(S\)-morphisms, and let \(F\) be a sheaf of groups on~\(X\).
Suppose that \((F,gf)\) is cohomologically proper relative to~\(S\)
in dimension \(\leqslant-1\) (\resp that \((F,gf)\) is
cohomologically proper relative to~\(S\) in dimension
\(\leqslant0\) and that \((F,f)\) is cohomologically proper relative
to~\(S\) in dimension \(\leqslant-1\)). Then \((f_*F,g)\) is
cohomologically proper relative to~\(S\) in dimension
\(\leqslant-1\) (\resp in dimension \(\leqslant0\)).
\end{corollaire}

Retain the notation of \eqref{eq:XIII.1.6.1}, and, for every scheme
\(Y_1\) étale over~\(Y\), denote by \(f_1,F_1\) the respective
inverse images of \(f,F\) under the morphism \(Y_1\to Y\). Let
\(\Phi\) be the stack of torsors under \(F\), and let \(\Psi\) be the
stack of torsors under \(f_*F\). There is a canonical functor
$$
\varphi \colon \Psi \to f_* \Phi \text{,}
$$
obtained by associating with every scheme \(Y_1\) étale over~\(Y\)
and every torsor \(P\) on~\(Y_1\) with group \(f_{1*}F_1\), the
torsor \(\widetilde P\) on~\(X_1\) obtained from \(f_1^*P\) by
extension of structure group along
\(f_1^*f_{1*}F_1\to F_1\). The functor \(\varphi\) is fully faithful.
Indeed, if \(P\) and \(Q\) are two torsors on~\(Y_1\) with group
\(f_{1*}F_1\), there is a canonical morphism
$$
\SheafIsom_{f_{1 *}F_1}(P,Q) \to
f_{1*}(\SheafIsom_{F_1}(\widetilde P,\widetilde Q))
$$
which
% page: 356
\marginpar{356}
%
is an isomorphism because this is so locally. It follows that the
canonical morphism
$$
\Isom_{f_{1 *}F_1}(P,Q) \to \Isom_{F_1}(\widetilde P,\widetilde Q)
$$
is an isomorphism, and hence that \(\varphi\) is fully faithful.

There is a commutative diagram
$$
\xymatrix{ m^*(g_* \Psi) \ar[r]^i \ar[d]_{m^*
g_*(\varphi)} & g'_* k^* \Psi \ar[d]^{g'_*
k^*(\varphi)} \\ m^*(g_* f_* \Phi) \ar[r]^j &
g'_* k^*(f_* \Phi)\text{,} }
% remark: as above, the comma is misplaced.
$$
where \(i\) and \(j\) are the base-change morphisms. It follows from
\Ref{XIII.1.6} 2) that \(j\) is faithful (\resp fully faithful).
Since \(g'_*k^*(\varphi)\) and \(m^*g_*(\varphi)\) are fully
faithful, the diagram above shows that \(i\) is faithful
(\resp fully faithful).

\begin{corollaire}
\label{XIII.1.8}
Let \(f\colon X\to Y\) be a coherent \(S\)-morphism, let
\(g\colon Y\to Z\) be a proper \(S\)-morphism, and let \(\Phi\) be
an ind-finite stack on~\(X\) \textup{\cite[VII 2.2.1]{XIII.2}}.
Suppose that \((\Phi,f)\) is cohomologically proper relative to~\(S\)
in dimension \(\leqslant-1\) (\resp \(\leqslant0\),
\resp \(\leqslant1\)); then the same is true of~\((\Phi,gf)\).
\end{corollaire}

Since \(f\) is coherent, \(f_*\Phi\) is an ind-finite stack
(SGA~4 IX 1.6~(ii)). The corollary therefore follows from
\Ref{XIII.1.6}~1) and~\Ref{XIII.1.4}.

\begin{corollaire}
\label{XIII.1.9}
Let \(f\colon X\to Y\) be an integral \(S\)-morphism, and let
\(g\colon Y\to Z\) be an \(S\)-morphism. If \(F\) is a sheaf of sets
on~\(X\), then \((f_*F,g)\) is cohomologically proper relative
to~\(S\) in dimension \(\leqslant-1\) (\resp \(\leqslant0\)) if and
only if the same is true of~\((F,gf)\). If \(F\) is a sheaf of
groups on~\(X\), then \((f_*F,g)\) is cohomologically proper relative
to~\(S\) in dimension \(\leqslant-1\) (\resp \(\leqslant0\),
\resp \(\leqslant1\)) if and only if the same is true of~\((F,gf)\).
\end{corollaire}

The assertion
% page: 357
\marginpar{357}
%
for a sheaf of sets follows from~\Ref{XIII.1.6} and from the fact
that \((F,f)\) is cohomologically proper relative to~\(S\) in
dimension \(\leqslant0\). Let \(F\) be a sheaf of groups on~\(X\),
and let \(\Phi\) be the stack of torsors under \(F\). By
SGA~4 VIII~5.8, every torsor under \(F\) is locally trivial over~\(Y\).
It follows that the stack \(f_*\Phi\) is equivalent to the stack of
torsors under \(f_*F\), the equivalence being obtained by associating
with every scheme \(Y_1\) étale over~\(Y\) and every torsor \(P\)
on~\(X_1=X\times_Y Y_1\) with group \(F|X_1\), the torsor \(f_*P\)
with group \(f_*F|Y_1\). Since \((F,f)\) is cohomologically proper
relative to~\(S\) in dimension \(\leqslant1\), the corollary
therefore follows from~\Ref{XIII.1.6}.

\enlargethispage{.5cm}
\skpt

\end{filecontents*}
% Original input: drafts/SGA1_XIII_1_OPENING_THROUGH_1_4_English.texfrag
\begin{filecontents*}[overwrite,noheader]{sga1-cumulative-input-102.tex}
\section{Cohomological properness}
\label{XIII.1}
%
\setcounter{subsection}{-1}
\subsection{}
\label{XIII.1.0}
%
Let \(S\) be a scheme and \(f\colon X\to Y\) a morphism of
\(S\)-schemes. If \(S'\) is an \(S\)-scheme, consider the following
diagram, all of whose squares are cartesian:
% page: 347
\marginpar{347}
%
\begin{equation*}
\label{eq:XIII.1.0.1}
\tag{\thesubsection.1}
\begin{array}{c}
\xymatrix{ X \ar[d]_f & \ar[l]_h X'
\ar[d]^{f'} \\ Y \ar[d] & \ar[l]_g Y' \ar[d] \\ S & \ar[l] \,S'. }
\end{array}
\end{equation*}

If \(Y_1\) is a scheme étale over~\(Y\), put
\(X_1=X\times_Y Y_1\), \(Y'_1=Y'\times_Y Y_1\), and consider the
cartesian square
\begin{equation*}
\label{eq:XIII.1.0.2}
\tag{\thesubsection.2}
\begin{array}{c}
\xymatrix{ X_1 \ar[d]_{f_1} & \ar[l]_{h_1}
X_1' \ar[d]^{f_1'} \\ Y_1 & \ar[l]^{g_1} \,Y_1'. }
\end{array}
\end{equation*}

\begin{definition}
\label{XIII.1.1}
Let \(F\) be a stack on~\(X\). We say that \((F,f)\) is
\emph{cohomologically proper relative to~\(S\) in dimension
\(\leqslant-1\)}
%
\index{cohomologically proper|hyperpage}%
\index{cohomological properness|hyperpage}%
%
(\resp in dimension \(\leqslant0\), \resp in dimension
\(\leqslant1\)) if, for every \(S\)-scheme \(S'\), the canonical
functor (defined in the evident way by the universal property of the
inverse image of stacks)
$$
g^* f_* F \to f'_* h^* F \qquad\text{ (\cf 1.0.1)}
$$
is faithful (\resp fully faithful, \resp an equivalence of
categories).
\end{definition}

If there is no possible ambiguity about~\(S\), in particular if
\(S=Y\), we say cohomologically proper instead of cohomologically
proper relative to~\(S\).

\subsection{}
\label{XIII.1.2}
Let \(F\) be a sheaf of sets on~\(X\), and let \(\Phi\) be the
associated stack in discrete categories, \ie the stack whose fiber
over every scheme \(X_1\) étale over~\(X\) is the discrete category
with object set \(F(X_1)\). We say that \((F,f)\) is cohomologically
proper relative to~\(S\) in dimension \(\leqslant-1\) (\resp in
dimension \(\leqslant0\)) if \((\Phi,f)\) is cohomologically proper
relative to~\(S\) in dimension \(\leqslant0\) (\resp in dimension
\(\leqslant1\)).

The
% page: 348
\marginpar{348}
%
canonical morphism
\begin{equation*}
\label{eq:XIII.1.2.1}
\tag{\thesubsection.1} { g^* f_* F \to f'_* h^* F }
\end{equation*}
induces, upon passing to the associated stacks in discrete
categories, the canonical morphism
$$
g^* f_* \Phi \to f'_* h^* \Phi \text{.}
$$
Consequently, saying that \((F,f)\) is cohomologically proper
relative to~\(S\) in dimension \(\leqslant-1\) (\resp in dimension
\(\leqslant0\)) is equivalent to saying that, for every \(S\)-scheme
\(S'\), the morphism \eqref{eq:XIII.1.2.1} is injective
(\resp bijective).

\subsection{}
\label{XIII.1.3}
Let \(F\) be a sheaf of groups on~\(X\), and let \(\Phi\) be the
stack of torsors on~\(X\) with group \(F\)
\cite[II 2.3.2]{XIII.1}. We say that \((F,f)\) is cohomologically
proper relative to~\(S\) in dimension \(\leqslant-1\)
(\resp \(\leqslant0\), \resp \(\leqslant1\)) if \((\Phi,f)\) is
cohomologically proper relative to~\(S\) in dimension
\(\leqslant-1\) (\resp \(\leqslant0\), \resp \(\leqslant1\)).
The cohomological properness condition can be made explicit as
follows.

\begin{subproposition}
\label{XIII.1.3.1}
Use the notation of \eqref{eq:XIII.1.0.1} and
\eqref{eq:XIII.1.0.2}. Let \(F\) be a sheaf of groups on~\(X\).
Denote by \(F'\) (\resp \(F_1\), \resp \(F'_1\), etc.) the inverse
image of~\(F\) on~\(X'\) (\resp on~\(X_1\), \resp on~\(X'_1\),
etc.). Then the following conditions are equivalent:
\begin{enumerate}
\item[(i)] \((F,f)\) is cohomologically proper relative to~\(S\)
in dimension \(\leqslant-1\) (\resp \(\leqslant0\),
\resp \(\leqslant1\)).
\item[(ii)] For every morphism \(S'\to S\), every scheme \(Y_1\)
étale over~\(Y\), and every torsor \(P\) on~\(X_1\) with group
\(F_1\), if \(\leftexp{P\mkern-4mu}{F}_1\) denotes the twist of
\(F_1\) by \(P\) \textup{\cite[II 4.1.2.3]{XIII.1}}, the canonical
morphism
$$
a_0 \colon g_1^*(f_{1 *}(\leftexp{P\mkern-4mu}{F}_1)) \to f'_{1
*}(\leftexp{P'\mkern-6mu}{F}'_1)
$$
is injective (\resp \(a_0\) is bijective and the canonical morphism
$$
a_1 \colon g^*(\R^1 f_* F) \to \R^1 f'_* F'
$$
is
% page: 349
\marginpar{349}
%
injective, \resp \(a_0\) and \(a_1\) are bijective).
\item[(ii~bis)] For every morphism \(S'\to S\), every scheme \(Y_1\)
étale over~\(Y\), every torsor \(P\) on~\(X_1\) with group \(F_1\),
and every torsor \(R\) under \(\leftexp{P\mkern-4mu}{F}_1\), the
canonical morphism
$$
\alpha_0 \colon g_1^*(f_{1 *}R) \to f'_{1 *} R'
$$
is injective (\resp \(\alpha_0\) is bijective, \resp the morphisms
\(\alpha_0\) and
$$
\alpha_1 \colon g_1^*(\R^1 f_{1 *}(\leftexp{P\mkern-4mu}{F}_1)) \to \R^1
f'_{1 *}(\leftexp{P'\mkern-6mu}{F}'_1)
$$
are bijective).
\end{enumerate}
\end{subproposition}

\subsubsection*{Proof}
% remark: all expressions of the form an overset label over wedge
% could be replaced by a superscripted wedge if vertical spacing causes
% difficulty.
(i)\(\To\)(ii~bis). By \cite[II 4.2.5]{XIII.1}, every torsor \(R\)
with group \(\leftexp{P\mkern-4mu}{F}_1\) is of the form
\(R=Q\overset{F_1}{\wedge}P^\circ\), where \(Q\) is a torsor with
group \(F_1\) and \(P^\circ\) is the opposite of~\(P\). We then have
\(R'\simeq Q'\overset{F'_1}{\wedge}{P'}^\circ\). Let \(\Phi\) be
the stack of torsors under \(F\), and let \(x,y\)
(\resp \(x',y'\)) be the objects of the fiber category
\((g^*f_*\Phi)_{Y'_1}\) (\resp \((f'_*\Phi')_{Y'_1}\)) associated
with \(P,Q\) (\resp \(P',Q'\)). We have
$$
Q\overset{F_1}{\wedge} P^\circ \simeq \SheafHom_{F_1}(P,Q) \text{,}
$$
and it follows that there are canonical isomorphisms
$$
\SheafHom_{Y'_1}(x,y) \simeq g_1^* f_{1 *}(Q
\overset{F_1}{\wedge} P^\circ) \text{,}\quad \SheafHom_{Y_1'}(x',y')
\simeq f'_{1 *}(Q' \overset{F'_1}{\wedge}{P'}^\circ)\text{.}
$$
Consequently, the morphism \(\alpha_0\) is identified with the
morphism
$$
\SheafHom_{Y'_1}(x,y) \to \SheafHom_{Y'_1}(x',y') \text{,}
$$
from which it follows that if \((F,f)\) is cohomologically proper
relative to~\(S\) in dimension \(\leqslant-1\), then \(\alpha_0\)
is injective, and that if \((F,f)\) is cohomologically proper in
dimension \(\leqslant0\), then \(\alpha_0\) is bijective.

Now suppose that \((F,f)\) is cohomologically proper relative
to~\(S\) in dimension \(\leqslant1\), \ie that the canonical
morphism
% page: 350
\marginpar{350}
%
$$
\varphi \colon g^* f_* \Phi \to f'_* \Phi'
$$
is an equivalence. Let \(G\) be the sheaf of maximal subgerbes of the
stack \(f_*\Phi\) \hbox{\cite[III~2.1.8]{XIII.1}}; then there is an
isomorphism \(G\simeq\R^1f_*F\). Since \(g^*G\) is the sheaf of
maximal subgerbes of~\(g^*f_*\Phi\) \cite[III 2.1.5.5]{XIII.2}, the
morphism \(\alpha_1\) is obtained from \(\varphi|Y'_1\) by taking
the sheaves of maximal subgerbes, and is therefore an isomorphism.

(ii~bis)\(\To\)(ii). It is enough to show that if the morphisms
\(\alpha_0\) are bijective, then the morphisms \(a_1\) are
injective. Let \(Y'_1\) be a scheme étale over~\(Y'\), and let \(s\)
and \(t\) be two elements of~\(g^*(\R^1f_*F)(Y'_1)\) having the same
image in \(\R^1f'_*F'(Y'_1)\); we shall show that \(s=t\). The
assertion is local for the étale topology of~\(Y'_1\), and, in view
of the definition of the inverse image \(g^*(\R^1f_*F)\), we may
suppose that \(Y'_1\) is the inverse image of a scheme \(Y_1\) étale
over~\(Y\), and that \(s\) and \(t\) come from torsors \(P\) and
\(Q\) on~\(X_1\). The hypothesis on \(s\) and \(t\) then means that
the inverse images \(P'\) and \(Q'\) of \(P\) and \(Q\) on~\(X'_1\)
are locally isomorphic for the étale topology of~\(Y'_1\). If we
put \(R=\SheafHom_{F_1}(P,Q)\), the fact that the morphism
$$
g_1^* f_{1 *} R \to f'_{1 *} R'
$$
is bijective proves that \(P\) and \(Q\) are locally isomorphic for
the étale topology of~\(Y_1\), and hence that \(s=t\).

(ii)\(\To\)(i). To prove that \(\varphi\) is faithful
(\resp fully faithful), it is enough to show that if \(Y_1\) is a
scheme étale over~\(Y\), if \(P,Q\) are two torsors on~\(X_1\) with
group~\(F_1\), and if \(x,y\) (\resp \(x',y'\)) are the objects of
\((g^*f_*\Phi)_{Y'_1}\) (\resp \((f'_*\Phi')_{Y'_1}\)) associated
with \(P,Q\) (\resp \(P',Q'\)), then the morphism
$$
a \colon \Hom(x,y) \to \Hom(x',y')
$$
is injective (\resp bijective). But \(a\) is identified with the
canonical morphism
$$
\H^0(Y'_1,g^*_1 f_{1 *} (Q \overset{F_1}{\wedge} P^\circ)) \to
\H^0(Y'_1, f'_{1 *}(Q' \overset{F'_1}{\wedge} {P'}^\circ))\text{.}
$$
If
% page: 351
\marginpar{351}
%
\(\Hom(x,y)\neq\emptyset\), then
\(Q\overset{F_1}{\wedge}P^\circ\) is a torsor under
\(\leftexp{P\mkern-4mu}{F}_1\) that is locally trivial over~\(Y_1\).
It follows that \(f_{1*}(Q\overset{F_1}{\wedge}P^\circ)\) is a
torsor under \(f_{1*}(\leftexp{P\mkern-4mu}{F}_1)\), and that
\(g^*_1f_{1*}(Q\overset{F_1}{\wedge}P^\circ)\) is a trivial torsor.
The morphism \(a\) is then identified with the canonical morphism
$$
\H^0(Y'_1, g^*_1 f_{1 *}(\leftexp{P\mkern-4mu}{F}_1)) \to \H^0 (Y'_1,
f'_{1 *}(\leftexp{P'\mkern-6mu}{F}'_1)) \text{.}
$$
The same holds if \(\Hom(x',y')\neq\emptyset\) and \(a_1\) is
injective, for then
\(Q'\overset{F'_1}{\wedge}{P'}^\circ\) is trivial, and the
injectivity of~\(a_1\) implies that \(P\) and \(Q\) are locally
isomorphic over~\(Y_1\). We conclude that if \(a_0\) is injective
(\resp if \(a_0\) is bijective and \(a_1\) injective), then
\(\varphi\) is faithful (\resp fully faithful).

It remains to show that if \(a_0\) and \(a_1\) are bijective, the
functor \(\varphi\) is essentially surjective. Let \(Y''\) be a
scheme étale over~\(Y'\), put \(X''=X'\times_{Y'}Y''\), and let
\(P''\) be a torsor on~\(X''\) with group \(F''=F'|X''\). We shall
show that there is an element \(x\) of~\((g^*f_*\Phi)_{Y''}\) whose
image in \((f'_*\Phi')_{Y''}\) is isomorphic to \(P''\). Let \(p''\)
be the class of~\(P''\). Since \(a_1\) is surjective, we can find a
surjective étale morphism \(Y''_1\to Y''\), an étale morphism
\(Y_1\to Y\) such that there is a morphism \(Y''_1\to Y'_1\), and a
torsor \(P_1\) on~\(X_1\) with group \(F_1\), whose inverse image
\(P''_1\) on~\(X''_1\) is isomorphic to the inverse image of~\(P''\).
Using the fact that \(\varphi\) is fully faithful, we see that the
object \(x_1\) of~\((g^*f_*\Phi)_{Y''_1}\) corresponding to
\(P''_1\) is equipped with descent data relative to
\(Y''_1\to Y''\), and hence comes from an element \(x\) of
\((g^*f_*\Phi)_{Y''}\). Since the image of~\(x\) in
\((f'_*\Phi')_{Y''}\) is \(P''\), this proves that \(\varphi\) is
essentially surjective and completes the proof.

\begin{exemple}
\label{XIII.1.4}
Let \(f\colon X\to Y\) be a proper morphism. It follows from
\cite[VII 2.2.2]{XIII.2} that, for every ind-finite stack \(F\)
on~\(X\), the pair \((F,f)\) is cohomologically proper (relative
to~\(Y\)) in dimension \(\leqslant1\). In particular, for every
sheaf of sets (\resp every sheaf of groups, \resp every ind-finite
sheaf of groups) \(F\) on~\(X\), \((F,f)\) is cohomologically
% page: 352
\marginpar{352}
%
proper in dimension \(\leqslant0\) (\resp in dimension
\(\leqslant0\), \resp in dimension \(\leqslant1\)).
\end{exemple}

\end{filecontents*}
% Original input: drafts/SGA1_XIII_2_0_English.texfrag
\begin{filecontents*}[overwrite,noheader]{sga1-cumulative-input-103.tex}
\section[A special case of cohomological properness]{A special case of cohomological properness: relative normal crossings divisors}
\label{XIII.2}

\setcounter{subsection}{-1}

\subsection{}
\label{XIII.2.0}
Let \(R\) be a discrete valuation ring with fraction field \(K\),
and let \(L\) be an étale \(K\)-algebra. Then \(L\) is the direct
product of finitely many fields \(L_i\), where \(L_i\) is an étale
extension of~\(K\). If \(L'_i\) denotes the Galois extension
generated by \(L_i\) in an algebraic closure of~\(L_i\), we say that
\(L\) is \emph{tamely ramified}
%
\index{tamely ramified (algebra)|hyperpage}%
%
over~\(R\) if the \(L'_i\) are tamely ramified extensions in the
sense of~\ifthenelse{\boolean{orig}}{X~}{}\Ref{X.3}, \ie if an inertia group \(I_i\) of \(L'_i|K\)
has order prime to the residue characteristic \(p\) of~\(R\).

It is known that \(I_i\) is in any case an extension of a cyclic
group of order prime to \(p\) by a \(p\)-group. (This follows from
\cite[ch\ptbl IV prop\ptbl 7 cor\ptbl 4]{XIII.5} when the residue
extension of~\(R\) is separable. The proof given in \loccit extends
to the general case as follows. Retain the hypotheses and notation
of \loccit, but without assuming the residue extension separable.
Let \(H_i\) be the subgroup of the inertia group \(G_0\) consisting
of the elements \(s\) of~\(G_0\) such that
\(s\pi/\pi\in U^i\) for every uniformizer \(\pi\) of~\(A_L\). One
then checks that \(G_0/H_1\) is a group of order prime to \(p\) and
that, for \(i\geqslant1\), the \(H_i/H_{i+1}\) are \(p\)-groups,
which gives the asserted result.)

\subsubsection{}
\label{XIII.2.0.1}
When \(R\) is strictly local, there is the following simple
characterization: the \(K\)-algebra \(L\) is tamely ramified over~\(R\)
if and only if the \([L_i:K]\) are prime to \(p\). Moreover, if \(L\)
is tamely ramified over~\(R\), the \(L_i\) are cyclic extensions
of~\(K\). Indeed, when \(R\) is strictly local, \(I_i\) is equal to
the Galois group
% page: 370
\marginpar{370}
%
of~\(L'_i\) over~\(K\). As just recalled, \(I_i\) is an extension of
a cyclic group of order prime to \(p\) by a \(p\)-group. If \(L'_i\)
is assumed tamely ramified over~\(R\), then \(I_i\) is a cyclic
group of order prime to \(p\). It follows that \([L_i:K]\) is prime
to \(p\) and that \(L_i=L'_i\). Conversely, if \([L_i:K]\) is prime
to \(p\), then \(I_i\) cannot contain a nontrivial normal
\(p\)-subgroup; \(I_i\) is therefore a cyclic group of order prime
to \(p\), which proves that \(L_i\) is tamely ramified over~\(R\).

\subsubsection{}
\label{XIII.2.0.2}
Let \(R\) be a discrete valuation ring with fraction field \(K\),
let \(L\) be an étale \(K\)-algebra, let \(\overline R\) be a strict
localization of~\(R\), let \(\overline K\) be its fraction field,
and put \(\overline L=L\otimes_K\overline K\). Then \(L\) is tamely
ramified over~\(R\) if and only if \(\overline L\) is tamely
ramified over~\(\overline R\). We may indeed reduce to the case
where \(L\) is a field. Write
\(\overline L=\prod_i\overline L_i\), where the \(\overline L_i\)
are field extensions of~\(\overline K\). If \(L'\) is the Galois
extension generated by \(L\), and if
\ifthenelse{\boolean{orig}}
{\(L'=L'\otimes_K\overline K\),}
{\(\overline{L'}=L'\otimes_K\overline K\),}
% correction: overline L' instead of L'
then \(\overline{L'}\) likewise decomposes as a product of fields,
\ifthenelse{\boolean{orig}}
{\(L'=\prod_j\overline{L'_j}\)}
{\(\overline{L'}=\prod_j\overline{L'_j}\)}
% correction: overline L' instead of L'
and every \(\overline L_i\) is a subextension of at least one
\(\overline{L'_j}\). Since \(L'\) is a Galois extension of~\(K\),
the \(\overline{L'_j}\) are Galois extensions of~\(\overline K\).
Suppose that \(L\) is tamely ramified over~\(R\). Since the Galois
group of~\(\overline{L'_j}|\overline K\) is isomorphic to the
inertia group of~\(L'|K\), the \(\overline{L'_j}\) are also tamely
ramified over~\(\overline R\), and hence so are the
\(\overline L_i\). Conversely, suppose that \(\overline L\) is
tamely ramified over~\(\overline R\). For each \(j\), let \(v_j\)
be the discrete valuation of~\(\overline{L'_j}\) extending the
valuation of~\(\overline K\), and also denote by \(v_j\) the induced
valuation on~\(L'\). As \(j\) varies, \(v_j\) ranges over the
valuations of~\(L'\) extending the valuation of~\(K\). Let
\(G=\Gal(L'|K)\), \(H=\Gal(L'|L)\), and let \(I_j\) be the inertia
group of~\(L'|K\) at \(v_j\), and \(J_j\) the inertia group of
\(L'|L\) at \(v_j\). The group \(I_j\) is an extension of a cyclic
group of order prime to \(p\) by a \(p\)-group \(P_j\). Since the
\(\overline L_i\) are tamely ramified over~\(\overline R\),
\(I_j/J_j\) has order prime to \(p\), and hence \(P_j\subset J_j\).
Consequently, the group \(H\) contains every \(P_j\), and therefore
also the group \(P\) generated by
% page: 371
\marginpar{371}
%
the \(P_j\) as \(j\) varies. But \(P\) is invariant in \(G\), since
an inner automorphism of~\(G\) permutes the \(I_j\), and hence also
the \(P_j\). It follows that \(P\) is a subgroup of~\(H\) normal
in~\(G\), and therefore, since \(L'\) is the Galois extension
generated by \(L\), that \(P=1\). This proves that \(L\) is tamely
ramified over~\(R\).

More generally, let \(R\to R'\) be a morphism of discrete valuation
rings such that the image of a uniformizer \(\pi\) of~\(R\) is a
uniformizer \(\pi'\) of~\(R'\), and such that the residue extension
\(\kres(R')|\kres(R)\) is separable. Let \(K\) be the fraction field
of~\(R\), \(K'\) the fraction field of~\(R'\), \(L\) an étale
\(K\)-algebra, and \(L'=L\otimes_KK'\). Then \(L\) is tamely
ramified over~\(R\) if and only if \(L'\) is tamely ramified
over~\(R'\). We may indeed suppose that \(R\) and \(R'\) are
strictly local. By~\Ref{XIII.2.0.1}, it is enough to prove that when
\(L\) is a field, so is \(L'\). Let \(\widetilde R\) be the
normalization of~\(R\) in \(L\), let \(\widetilde\pi\) be a
uniformizer of~\(\widetilde R\), and put
\(\widetilde R'=\widetilde R\otimes_RR'\). Since the extension
\(\kres(\widetilde R)|\kres(R)\) is radicial and the extension
\(\kres(R')|\kres(R)\) is étale,
\(\kres(\widetilde R)\otimes_{\kres(R)}\kres(R')\) is a field
[EGA IV 4.3.2 and 4.3.5]. This proves that \(\widetilde R'\) is a
local ring, and, since \(\pi\) maps to \(\pi'\) in~\(R'\), we have
\(\kres(\widetilde R')=\widetilde R'/(\widetilde\pi)\).
Consequently, \(\widetilde R'\) is a discrete valuation ring
\cite[ch\ptbl I \S 2 prop\ptbl 2]{XIII.5}, and hence \(L'\) is a
field.

\subsubsection{}
\label{XIII.2.0.3}
By reducing to the strictly local case, we see that a subalgebra of
a tamely ramified algebra is tamely ramified; that the tensor product
of two tamely ramified algebras is tamely ramified; that a tamely
ramified algebra remains so after extending the discrete valuation
ring; and that an algebra which becomes tamely ramified after a
tamely ramified extension is itself tamely ramified.

\end{filecontents*}
% Original input: drafts/SGA1_XIII_2_1_English.texfrag
\begin{filecontents*}[overwrite,noheader]{sga1-cumulative-input-104.tex}
\subsection{}
\label{XIII.2.1}
Let
% page: 372
\marginpar{372}
%
$X$ be an $S$-scheme, and let $D$ be a divisor $\geq 0$ on~$X$.
Recall (SGA~5 II~4.2) that $D$ is said to have
\emph{strict normal crossings relative to}~$S$
%
\index{normal crossings (divisor with)|hyperpage}%
\index{divisor with normal crossings|hyperpage}%
\index{strict normal crossings (divisor)|hyperpage}%
%
if there exists a finite family $(f_i)_{i\in I}$ of elements of
$\Gamma(X,\cal{O}_X)$ such that $D=\sum_{i\in
I}\divisor(f_i)$ and the following condition is satisfied:

% remark: the structuring of the text could be improved: 2.1.0 is not
% a subsection, but more like a quotation.
\setcounter{subsubsection}{-1}
\subsubsection{}
\label{XIII.2.1.0}
For every point $x$ of~$\Supp D$, $X$ is smooth over~$S$ at $x$, and,
if $I(x)$ denotes the set of $i\in I$ such that $f_i(x)=0$, the
subscheme $V((f_i)_{i\in I(x)})$ is smooth over~$S$ of codimension
\ifthenelse{\boolean{orig}}
{$\mathrm{card}.\,I(x)$}
{$\card I(x)$}
in~$X$.

The divisor $D$ is said to have \emph{normal crossings relative
to} $S$ if, locally on~$X$ for the \'etale topology, it has strict
normal crossings.

Let $D$ be a divisor with normal crossings relative to~$S$. Put
$Y=\Supp D$, $U=X-Y$, and denote by $i\colon U\to X$ the canonical
immersion. For every geometric point $\overline{s}$ of~$S$ and every
maximal point $y$ of the geometric fiber $Y_{\overline{s}}$, the ring
$R=\cal{O}_{X_{\overline{s}},y}$ is a discrete valuation ring.

In the remainder of this number, we shall use the following technical
definition:

\begin{subdefinition}
\label{XIII.2.1.1}
Let $F$ be a sheaf of sets on~$U$. We say that $F$ is
\emph{tamely ramified on}~$X$
%
\index{tamely ramified (sheaf)|hyperpage}%
%
(\emph{along}~$D$) \emph{relative to}~$S$ if, for every geometric
point $\overline{s}$ of~$S$, the following condition is satisfied:

For every maximal point $y$ of~$Y_{\overline{s}}$, the restriction of
$F$ to the fraction field $K$ of~$\cal{O}_{X_{\overline{s}},y}$ is
representable by the spectrum of an \'etale $K$-algebra $L$ which is
tamely ramified over~\ifthenelse{\boolean{orig}}{$\cal{O}_{X_{\overline{s},y}}$}{$\cal{O}_{X_{\overline{s}},y}$}.
\end{subdefinition}

Most often, when no confusion can result, we shall omit mention
of~$D$ from the terminology.

\begin{subdefinition}
\label{XIII.2.1.2}
If
% page: 373
\marginpar{373}
%
$F$ is a sheaf of groups on~$U$, tamely ramified on~$X$ relative
to~$S$, we denote by
$$
\H^1_{\tame}(U,F)
$$
\label{indnot:mc}\oldindexnot{$\H^1_{\tame}(U,F)$|hyperpage}%
%
the subset of~$\H^1(U,F)$ consisting of the classes of torsors under
$F$ which are tamely ramified on~$X$ relative to~$S$.
\end{subdefinition}

Let
$$
\xymatrix{ U \ar[r]^i \ar[d]_g & X \ar[dl]^f \\ T}
$$
be a commutative diagram of~$S$-schemes, with $i$ as
in~\Ref{XIII.2.1}; we denote by
$$
\R^1_{\tame} g_* F
$$
\label{indnot:md}\oldindexnot{$\R^1_{\tame} g_* F$|hyperpage}%
%
the sheaf on~$T$ associated with the presheaf
$T'\mto\H^1_{\tame}(U',F)$, where $T'$ ranges over the schemes
\'etale over~$T$ and where $U'=U\times_T T'$; $\R^1_{\tame} g_*F$
is a subsheaf of~$\R^1g_*F$.

Note that, if $g$ is coherent, if $\overline{t}$ is a geometric point
of~$T$, if $\overline{T}$ is the strict localization of $T$ at
$\overline{t}$, and if $\overline{U}=U\times_T\overline{T}$, then there
is an isomorphism
\begin{equation*}
\label{eq:XIII.2.1.2.1}
\tag{\thesubsubsection.1} (\R^1_{\tame} g_*F)_{\overline{t}}\simeq
\H^1_{\tame}(\overline{U},\overline{F})
\end{equation*}

\subsubsection{}
\label{XIII.2.1.3}
Let $C_t((U,X)/S)$, or simply $C_t$,
\label{indnot:me}\oldindexnot{$C_t((U,X)/S)$ or $C_t$|hyperpage}%
%
be the category of \'etale coverings of~$U$ that are tamely ramified
on~$X$ relative to~$S$. Suppose that $U$ is connected, and let $a$ be
a geometric point of~$U$. Let $\Gamma_t$ be the functor which assigns
to an \'etale covering~$U'$ of~$U$ tamely ramified on~$X$ relative
to~$S$ the set of geometric points of~$U'$ above~$a$. It follows
from~\Ref{XIII.2.0} that the pair $(C_t,\Gamma_t)$ satisfies axioms
($G_1$) through ($G_6$) of \ifthenelse{\boolean{orig}}{V~}{}\Ref{V.4}. Consequently, $\Gamma_t$ is
representable by a pro-object called the
\emph{tamely ramified universal covering
% page: 374
\marginpar{374}
%
of~$(U,X)$ relative to~$S$ pointed at~$a$}.
%
\index{tamely ramified (universal covering)|hyperpage}%
\index{universal tamely ramified covering@tamely ramified universal covering|hyperpage}%
%
The opposite of the group of $U$-automorphisms of the tamely ramified
universal covering is called the \emph{tamely ramified fundamental
group}
%
\index{tamely ramified fundamental group|hyperpage}%
\index{tamely ramified (fundamental group)|hyperpage}%
%
and is denoted by
$$
\pi_1^{\tame}((U,X)/S,a)\text{ or simply }\pi_1^{\tame}(U,a)\text{ or
even }\pi_1^{\tame}(U).
$$
\label{indnot:mf}\oldindexnot{$\pi_1^{\tame}((U,X)/S,a)$ or $\pi_1^{\tame}(U,a)$ or $\pi_1^{\tame}(U)$|hyperpage}%
%
It is clearly a quotient of the fundamental group $\pi_1(U,a)$
(\ifthenelse{\boolean{orig}}{V~}{}\Ref{V.6.9}).

\subsubsection{}
\label{XIII.2.1.4}
Let $F$ be a sheaf of groups on~$U$, let $P$ be a right torsor with
group $F$, and let $Q$ be a left torsor with group $F$. Suppose that
$P$ and $Q$ are tamely ramified on~$X$ relative to~$S$. Then the same
is true of~$P\stackrel{F}{\wedge}Q$. Indeed, we reduce to showing that,
if $R$ is a discrete valuation ring with fraction field $K$, if~$F$ is
a finite \'etale group scheme over~$K$, and if $P$ and $Q$ are two
torsors under~$F$ tamely ramified over~$R$, then the same is true
of~$P\stackrel{F}{\wedge} Q$. Now
$T=P\stackrel{F}{\wedge} Q$ is a quotient of~$P\times_K Q$. If $L$,
$M$, $N$ denote the $K$-algebras representing $T$, $P$, $Q$,
respectively, then $L$ is a subalgebra of $M\otimes_K N$, and it
follows from~\Ref{XIII.2.0.3} that $L$ is tamely ramified over~$R$.

It follows from the preceding that, if $F$ is a sheaf of groups on~$U$
and if there exists a torsor with group $F$ tamely ramified on~$X$
relative to~$S$, then $F$ is tamely ramified on~$X$ relative to~$S$.
Indeed, the opposite torsor $P^{\circ}$ of~$P$ is tamely ramified
on~$X$ relative to~$S$, since it is isomorphic to $P$ as a sheaf of
sets. If $\leftexp{P\mkern-4mu}{F}$ is the group obtained by twisting
$F$ by $P$, there is an isomorphism
$$
F\simeq P^{\circ}\stackrel{\leftexp{P\mkern-4mu}{F}}{\wedge}P
$$
and consequently $F$ is tamely ramified on~$X$ relative to~$S$.

As before, we see that, if $F\to F'$ is a morphism of sheaves of
groups on~$U$ tamely ramified on~$X$ relative to~$S$, and if $P$ is a
torsor under $F$ tamely ramified on~$X$ relative
% page: 375
\marginpar{375}
%
to~$S$, then the torsor $P'$ deduced from~$P$ by extension of the
structure group $F\to F'$ is tamely ramified on~$X$ relative to~$S$.

In particular, the canonical morphism
$$
\H^1(U,F)\to \H^1(U,F')
$$
induces, by restriction to $\H^1_{\tame}(U,F)$, a canonical morphism
$$
\H^1_{\tame}(U,F)\to\H^1_{\tame}(U,F')
$$

\subsubsection{}
\label{XIII.2.1.5}
Let $S'\to S$ be a morphism, and denote by $U'$ (\resp $X'$, etc.)
the inverse image of~$U$ (\resp $X$, etc.) on~$S'$. If $F$ is a
sheaf of sets on~$U$ tamely ramified on~$X$ relative to~$S$, it
follows from Definition~\Ref{XIII.2.1.1} and from~\Ref{XIII.2.0.3}
that $F'$ is tamely ramified on~$X'$ relative to~$S'$.

If now $F$ is a sheaf of groups on~$U$, the inverse image on~$S'$ of
a torsor under $F$ tamely ramified on~$X$ relative to~$S$ is a torsor
under $F'$ tamely ramified on~$X'$ relative to~$S'$. In particular,
there is a canonical functor
\begin{equation*}
\label{eq:XIII.2.1.5.1}
\tag{\thesubsubsection.1} C_t((U,X)/S)\to C_t((U',X')/S').
\end{equation*}
Suppose that $U$ and $U'$ are connected, and let $a$ be a geometric
point of~$U$ and $a'$ a geometric point of~$U'$ above~$a$; the
preceding gives a canonical morphism
\begin{equation*}
\label{eq:XIII.2.1.5.2}
\tag{\thesubsubsection.2} \pi_1^{\tame}(U',a')\to\pi_1^{\tame}(U,a).
\end{equation*}

If $S'\to S$ is a morphism and $h\colon T'\to T$ is the canonical
projection, the morphism
$$
h^*(\R^1g_*F)\to\R^1g_*'F'
$$
restricts to a canonical morphism
\begin{equation*}
\label{eq:XIII.2.1.5.3}
\tag{\thesubsubsection.3} h^*(\R^1_{\tame}g_*F)\to\R^1_{\tame}g_*'F'
\end{equation*}

\subsubsection{}
\label{XIII.2.1.6}
Let
% page: 376
\marginpar{376}
%
$F$ be a sheaf of groups on~$U$, tamely ramified on~$X$ relative
to~$S$. With the notation of~\Ref{XIII.2.1.2}, there are canonical
exact sequences:
\begin{equation*}
\label{eq:XIII.2.1.6.1}
\tag{\thesubsubsection.1}
\begin{array}{ccccccc}
1 & \longrightarrow & \H^1(X,i_*F)& \longrightarrow & \H^1_{\tame}(U,F) &
\longrightarrow & \H^0(X,\R^1_{\tame} i_*F)\\[10pt] 1 &
\longrightarrow & \R^1f_*(i_*F) & \longrightarrow & \R^1_{\tame}g_*F &
\longrightarrow & f_*(\R^1_{\tame}i_*F).\end{array}
\end{equation*}

The first is obtained from the exact sequence (SGA~4 III~3.2):
$$
1\to \H^1(X,i_*F)\to \H^1(U,F)\to
\H^0(X,\R^1i_*F).
$$
Indeed, it suffices to show that the image of~$\H^1(X,i_*F)$ in
$\H^1(U,F)$ is in fact contained in $\H^1_{\tame}(U,F)$ and that the image
of $\H^1_{\tame}(U,F)$ in $\H^0(X,\R^1i_*F)$ is in fact contained in
$\H^0(X,\R^1_{\tame}i_*F)$. Now the inverse image on~$U$ of a torsor under
$i_*F$ is a torsor under $i^*i_*F$ which is clearly tamely ramified
on~$X$ relative to~$S$; the same is therefore true after extension of
the structure group $i^*i_*F\to F$, which proves the existence of the
arrow $\H^1(X,i_*F)\to\H^1_{\tame}(U,F)$. The fact that the image of
$\H^1_{\tame}(U,F)$ in $\H^0(X,\R^1i_*F)$ is contained in
$\H^0(X,\R^1_{\tame}i_*F)$ follows at once from the definition of
$\R^1_{\tame}i_*F$. This proves the existence of the first exact
sequence, and the second follows from it by localization.
\end{filecontents*}
% Original input: drafts/SGA1_XIII_2_2_OPENING_THROUGH_THEOREM_English.texfrag
\begin{filecontents*}[overwrite,noheader]{sga1-cumulative-input-105.tex}
\subsection{}
\label{XIII.2.2}
Retain the notation of~\Ref{XIII.2.1}. We shall define a notion of a
tamely ramified object of a stack \(\Phi\) on~\(U\) when the latter
is given, locally \emph{on} \(X\) and on~\(S\) for the étale
topology, as the \emph{inverse image of a stack} \(\Psi\)
\emph{on}~\(S\).

First let \(G\) be a gerbe on~\(U\), and suppose that we are given a
surjective étale morphism \(S_1\to S\), a surjective étale morphism
\(X_2\to X\times_SS_1\), a trivial gerbe \(H\) on~\(S_1\), and an
isomorphism
$$
G|{U_2}\to H|{U_2},
$$
where
% page: 377
\marginpar{377}
%
\(X_1=X\times_SS_1\), \(U_1=U\times_XX_1\), \(U_2=U\times_XX_2\). Upon choosing a
trivialization of~\(H|X_2\), the isomorphism above identifies
\(G|U_2\) with the stack of torsors under a sheaf of groups \(F\).
An object \(x\) of~\(G_U\) is said to be tamely ramified on~\(X\)
relative
\ifthenelse{\boolean{orig}}
{to~\(S\),}
{to~\(S\)}
if the restriction of~\(x\) to \(U_2\) is a torsor tamely ramified
on~\(X_2\) relative to~\(S_1\). By~\Ref{XIII.2.1.4}, this notion does
not depend on the chosen trivialization of~\(H|X_2\).

Now let \(\Phi\) be a stack on~\(U\), and suppose that we are given a
surjective étale morphism \(S_1\to S\), a surjective étale morphism
\(X_2\to X\times_SS_1\), a stack \(\Psi\) on~\(S_1\), and an
isomorphism
$$
i\colon \Phi|U_2\to \Psi|U_2.
$$
Let \(x\) be an object of~\(\Phi_U\), let \(G_x\) be the maximal
subgerbe of~\(\Phi\) generated by \(x\)
\ifthenelse{\boolean{orig}}
{[III~2.1.7]}
{\cite[III~2.1.7]{XIII.1}}, and let \(S\Phi\) be the sheaf of
maximal subgerbes of~\(\Phi\). The isomorphism \(i\) induces an
isomorphism
$$
S\Phi|U_2\to S\Psi|U_2.
$$
It follows from~\Ref{XIII.5.7} that, after replacing \(S_1\) by a
surjective étale extension if necessary, there is a unique maximal
subgerbe \(H\) of~\(\Psi\), which may be assumed trivial, such that
\(i\) defines an isomorphism
$$
G_x|U_2\to H|U_2.
$$
The object \(x\) is said to be tamely ramified on~\(X\) relative
to~\(S\) if it is so as an object of~\(G_x\) equipped with the
isomorphism above.

\subsubsection{}
\label{XIII.2.2.1}
Let \(\Phi\) be a stack on~\(U\) given, locally on~\(X\) and \(S\),
as the inverse image of a stack on~\(S\), and let
$$
\xymatrix{ U \ar[r]^i
\ar[d]_g & X \ar[dl]^f \\ T }
$$
be a diagram as in~\Ref{XIII.2.1.2}. For every scheme \(T'\) étale
over~\(T\), if \(X'=X\times_TT'\) and \(U'=U\times_TT'\), consider the subset
\((g_*^{\tame}\Phi)_{T'}\)
\label{indnot:mg}\oldindexnot{\((g_*^{\tame}\Phi)_{T'}\)|hyperpage}%
%
of~\((g_*\Phi)_{T'}=\Phi_{U'}\), consisting
% page: 378
\marginpar{378}
%
of the objects of~\(\Phi_{U'}\) that are tamely ramified
%
\index{tamely ramified object of a stack|hyperpage}%
%
on~\(X'\) relative to~\(S\). The \emph{tamely ramified direct image}
%
\index{tamely ramified direct image of a stack|hyperpage}%
\index{tamely ramified (direct image of a stack)|hyperpage}%
%
of~\(\Phi\) under~\(g\), denoted
$$
g_*^{\tame}\Phi,
$$
is the full subcategory of~\(g_*\Phi\) whose objects over a scheme
\(T'\) étale over~\(T\) are the elements of
\((g_*^{\tame}\Phi)_{T'}\). It is clear that
\(g_*^{\tame}\Phi\) is a substack of~\(g_*\Phi\).

\subsubsection{}
\label{XIII.2.2.2}
Retain the notation of~\Ref{XIII.2.1.5}. By reducing to the case of a stack of torsors, we see that if
\(\Phi\) is a stack on~\(U\) which, locally for the étale topology
of~\(S\) and \(X\), is the inverse image of a stack on~\(S\), then
the canonical morphism
$$
h^*(g_*\Phi)\to g_*'\Phi'
$$
restricts to a canonical morphism
$$
h^*(g_*^{\tame}\Phi)\to g_*^{\prime\tame}\Phi'.
$$

\begin{remarques}
\label{XIII.2.3}

a) If \(F\) is a constructible locally constant sheaf of sets
on~\(U\), then, for \(F\) to be tamely ramified on~\(X\) relative
to~\(S\), it is enough that the condition of~\Ref{XIII.2.1.1} hold
at the geometric points of~\(S\) lying over the maximal points
of~\(S\). To see this, we may suppose that \(D\) is a strict normal
crossings divisor. The sheaf \(F\) is represented by an étale
covering \(V\) of~\(U\). If \(\overline s\) is a geometric point
of~\(S\) and \(y\) a maximal point of~\(Y_{\overline s}\), choose a geometric point
\(\overline y\) above~\(y\), denote by \(\overline S\) the strict localization of~\(S\) at \(\overline s\),
by \(\overline X\) the strict localization of~\(X\) at
\(\overline y\), and put
\(\overline U=U\times_X\overline X\),
\(\overline V=V\times_X\overline X\). If the condition
of~\Ref{XIII.2.1.1} holds at geometric points over the maximal
points of~\(S\), it follows from~\Ref{XIII.5.5} below that
\(\overline V\) is a covering of~\(\overline U\) tamely ramified
on~\(\overline X\) relative to~\(\overline S\). Consequently, \(V\)
is an étale covering of~\(U\) tamely ramified on~\(X\) relative
to~\(S\).

b)
% page: 379
\marginpar{379}
%
Let \(F\) be a sheaf of groups on~\(U\) tamely ramified on~\(X\)
relative to~\(S\). If \(\overline s\) is a geometric point of~\(S\)
and \(y\) a maximal point of~\(Y_{\overline s}\), choose a geometric point
\(\overline y\) above~\(y\), and denote by \(K\)
the fraction field of~\(\cal{O}_{X_{\overline s},y}\). Suppose that,
for every point \(\overline s\) and every point \(y\), the
\(K\)-algebra \(L\) whose spectrum represents \(F|K\) has rank
prime to the residue characteristic \(p\) of
\(\cal{O}_{X_{\overline s},y}\). By an abuse of language, we shall
sometimes say that \(F\) is prime to the residue characteristics
of~\(S\).
%
\index{prime to the residue characteristics (sheaf)|hyperpage}%
%
When this is so, every torsor \(P\) under \(F\) is tamely ramified
on~\(X\) relative to~\(S\).

Indeed, let \(\overline R\) be the strict localization of
\(\cal{O}_{X_{\overline s},y}\) at \(\overline y\), let
\(\overline K\) be its fraction field, and let \(\overline F\) be
the inverse image of~\(F\) on~\(\overline K\). We show that \(\overline F\)
may be assumed constant. Since \(F\) is tamely ramified on~\(X\)
relative to~\(S\), \(\overline F\) is represented by the spectrum of
a \(\overline K\)-algebra \(L=\prod L_i\), where the \(L_i\) are
extensions of~\(\overline K\) of degree prime to \(p\). We can
therefore find an extension \(K'\) of~\(\overline K\) of degree
prime to \(p\) such that \(\overline F|K'\) is a constant sheaf.
By~\Ref{XIII.2.0.3}, to prove that \(P|\overline K\) is tamely
ramified over~\(\overline R\), it is enough to see that \(P|K'\) is
tamely ramified over the integral closure of~\(\overline R\)
in~\(K'\); this reduces us to the case where \(\overline F\) is
constant. Henceforth suppose that \(\overline F\) is constant. The
\(\overline K\)-algebra \(H\) representing \(P|\overline K\) is
then a product of mutually isomorphic extensions \(H_i\)
of~\(\overline K\). Since the rank of~\(H\) is prime to \(p\), the
same is true of~\([H_1:\overline K]\), which proves that \(H\) is
tamely ramified over~\(\overline R\).

c) Let \(X\) be a regular scheme, let \(D\) be a normal crossings
divisor of~\(X\) (SGA~5 I~3.1.5), let \(U=X-\Supp D\), and let \(F\)
be a sheaf of sets on~\(U\). If \(y\) is a maximal point of
\(\Supp D\), denote by \(K\) the fraction field of
\(\cal{O}_{X,y}\). We say that \(F\) is \emph{tamely ramified
relative to}~\(D\)
%
\index{tamely ramified (sheaf)|hyperpage}%
%
if, for every maximal point \(y\) of~\(\Supp D\), \(F|K\) is
represented by a \(K\)-algebra tamely ramified over
\(\cal{O}_{X,y}\).
\end{remarques}

\begin{theoreme}
\label{XIII.2.4}
Let
% page: 380
\marginpar{380}
%
\(f\colon X\to S\) be an \(S\)-scheme, let \(D\) be a divisor
on~\(X\) with normal crossings relative to~\(S\)
\eqref{XIII.2.1}, put \(Y=\Supp D\) and \(U=X-Y\), and let
\(i\colon U\to X\) be the canonical immersion. Let \(F\) be a sheaf
of sets (\resp groups) on~\(U\), satisfying one of the following
conditions:
\begin{enumerate}
\item[a)] Locally for the étale topology \emph{on} \(X\) and on~\(S\),
\(F\) is the inverse image of a sheaf of sets (\resp a constructible
sheaf of groups) on~\(S\).
\item[b)] \(F\) is constructible locally constant on~\(U\) and is
tamely ramified on~\(X\) relative to~\(S\).
\end{enumerate}
Then the following conclusions hold:
\begin{enumerate}
\item[1)] \((F,i)\) is cohomologically proper relative to~\(S\) in
dimension \(\leq0\) (\resp for every morphism \(h\colon S'\to S\),
if \(i'\colon U'\to X'\) is the inverse image of~\(i\) on~\(S'\),
if \(F'=F|U'\), and if \(k=h_{(X)}\), the canonical morphism
$$
\Psi\colon k^*(\R^1_{\tame}i_*F)\to\R^1_{\tame}i_*'F'
$$
is an isomorphism).

If \(F\) is a sheaf of groups prime to the residue characteristics
of~\(S\) (\Ref{XIII.2.3}.b) (and tamely ramified on~\(X\) relative
to~\(S\)), then \((F,i)\) is cohomologically proper relative
to~\(S\) in dimension \(\leq1\).
\item[2)] If \(F\) is a constructible sheaf of sets (\resp groups),
then \(i_*F\) (\resp \(\R^1_{\tame}i_*F\)) is constructible.
\end{enumerate}
\end{theoreme}

\subsubsection*{Proof}
For every \(S\)-scheme \(S'\), consider the following diagram, all
of whose squares are cartesian:
$$
\xymatrix{ U \ar[dd]_-g \ar[dr]^-i & & & U'\ar[lll] \ar[dd]_{\hbox{\raise10mm\hbox{$g'$}}}
\ar[dr]^-{i'}& \\ & X \ar[dl]_{f\mkern-8mu} & & & X'\ar[lll]_{k\hspace*{1cm}} \ar[dl]_{f'\mkern-10mu} \\ S
& & & S'\ar[lll]_-h & & }
$$
Since
% page: 381
\marginpar{381}
%
the question is local on~\(X\) for the étale topology, we may
suppose that \(D\) is a strict normal crossings divisor relative
to~\(S\) \eqref{XIII.2.1}; moreover, after restricting \(X\) to a
neighborhood of~\(Y\) if necessary, we may suppose that \(X\) is
smooth over~\(S\).

\end{filecontents*}
% Original input: drafts/SGA1_XIII_2_3_English.texfrag
\begin{filecontents*}[overwrite,noheader]{sga1-cumulative-input-106.tex}
\subsection{}
\label{XIII.2.10}
If $U$ is a connected scheme, $a$ is a geometric point of~$U$, and
$\LL$ is a set of prime numbers, we denote by
\begin{equation*}
\label{eq:XIII.2.10.0}
\tag{2.10.0} \pi_1^{\LL}(U,a)
\end{equation*}
\label{indnot:mi}\oldindexnot{$\pi_1^{\LL}(U,a)$|hyperpage}%
%
the inverse limit of the finite quotients of~$\pi_1(U,a)$ whose orders
have all their prime factors in~$\LL$.

We shall define specialization homomorphisms for the fundamental group,
generalizing X.\Ref{X.2}.

Let $g\colon U\to S$ be a coherent morphism with geometrically
connected fibers (\resp a morphism of the form $g= fi$, where
$f\colon X \to S$ is a proper morphism of finite presentation and
$i\colon U\to X$ is an open immersion such that $U$ is the complement
in $X$ of a divisor with normal crossings relative to~$S$
(\cf \Ref{XIII.2.8})). Let ${\LL}$ be a set of prime numbers and
suppose, in the first case, that for every finite constant ${\LL}$-group
$C$, $(C_U,g)$ is cohomologically proper relative to~$S$ in dimension
$\leq 1$. Let $\overline{s}_1 \to \overline{s}_2$ be a specialization
morphism of geometric points of~$S$, let $\overline{S}$ be the
strict localization of~$S$ at $\overline{s}_2$, and put
$\overline{U}=U\times_S\overline{S}$. There is a commutative diagram
$$
\xymatrix{ U_{\bar{s}_1} \ar[r]^{h_1} \ar[d] & \bar{U}
\ar[d]^{\bar{g}} & U_{\bar{s}_2} \ar[l]_{h_2} \ar[d] \\ \bar{s}_1
\ar[r] & \bar{S} & \,\bar{s}_2\,. \ar[l] }
$$
Let $a_1$ be a geometric point of~$U_{\bar{s}_1},\ a_2$ a
geometric point of~$U_{\bar{s}_2}$; the morphisms $h_1$ and $h_2$
define canonical homomorphisms
% page: 391
\marginpar{391}%
%
\begin{align*}
\pi_1\colon \pi_1^{\LL}(U_{\bar{s}_1},a_1)&\to
\pi_1^{\LL}(\bar{U},a_1) & \pi_2\colon
\pi_1^{\LL}(U_{\bar{s}_2},a_2)&\to \pi_1^{\LL}(\bar{U},a_2)\\
(\text{\resp~} \pi_1\colon \pi_1^\tame (U_{\bar{s}_1},a_1)&\to
\pi_1^\tame (\bar{U},a_1) & \pi_2\colon
\pi_1^\tame(U_{\bar{s}_2},a_2)&\to \pi_1^\tame(\bar{U},a_2))
\end{align*}
(\ifthenelse{\boolean{orig}}{V~}{}\Ref{V.7} and~\Ref{eq:XIII.2.1.5.2}). The hypotheses of
cohomological properness (\resp \Ref{XIII.2.7}) show that $\pi_2$
is an isomorphism. If we choose a path class from~$a_1$ to $a_2$, we
obtain an isomorphism
$$
\pi_{12}\colon \pi_1^{\LL}(\bar{U},a_1)\isomto
\pi_1^{\LL}(\bar{U},a_2) \quad (\text{\resp~} \pi_{12}\colon
\pi_1^\tame (\bar{U},a_1)\isomto \pi_1^\tame(\bar{U},a_2))
$$
% correction: This last isomorphism was written only as a morphism.
and hence a homomorphism $\pi=\pi_2^{-1}\pi_{12}\pi_1$
$$
\pi\colon \pi_1^{\LL}(U_{\bar{s}_1},a_1)\to
\pi_1^{\LL}(U_{\bar{s}_2},a_2) \quad (\text{\resp~}\pi\colon
\pi_1^\tame (U_{\bar{s}_1},a_1)\to \pi_1^\tame
(U_{\bar{s}_2},a_2)).
$$
Changing the path class from~$a_1$ to $a_2$ amounts to modifying
$\pi$ by an inner automorphism of $\pi_1^{\LL}(U_{\bar{s}_2},a_2)$
(\resp of $\pi_1^\tame(U_{\bar{s}_2},a_2)$). One calls any one of the
homomorphisms defined above the \emph{specialization homomorphism for
the fundamental group}
%
\index{specialization homomorphism for the fundamental group|hyperpage}%
\index{specialization for the fundamental group (homomorphism)|hyperpage}%
%
associated with the morphism $\bar{s}_1\to \bar{s}_2$, and writes
simply
$$
\pi\colon \pi_1^{\LL}(U_{\bar{s}_1})\to
\pi_1^{\LL}(U_{\bar{s}_2}) \quad \text{(\resp~} \pi\colon
\pi_1^{\tame}(U_{\bar{s}_1})\to \pi_1^{\tame}(U_{\bar{s}_2}))
$$
for one of the homomorphisms defined above.

\begin{lemme}
\label{XIII.2.11}
Let $f\colon X\to S$ be a proper morphism of finite presentation, let
$D$ be a divisor on~$X$ with normal crossings relative to~$S$, put
$Y=\Supp D$, $U=X-Y$, and let $i\colon U\to X$ be the canonical
morphism. Let $\bar{s}_1\to \bar{s}_2$ be a specialization morphism
of geometric points of~$S$, let $y_1$ be a geometric point of
$Y_{\bar{s}_1}$, and let $y_2$ be a geometric point of
$Y_{\bar{s}_2}$ such that the projection $z_1$ of~$y_1$ on~$X$ is a
generization of the projection $z_2$ of~$y_2$, and choose a compatible
specialization morphism $y_1\to y_2$. Suppose moreover that the sets of
irreducible components of $D$ containing $z_1$ and $z_2$ coincide. Let $I_{y_1}$ be an
inertia subgroup of~$\pi_1^{\tame}(\bar{U}_{\bar{s}_1})$ at $y_1.$
Then the image of~$I_{y_{1}}$ under the specialization homomorphism
$$
\pi\colon \pi_1^{\tame}(U_{\bar{s}_1})\to
\pi_1^{\tame}(U_{\bar{s}_2})
$$
is
% page: 392
\marginpar{392}
%
an inertia subgroup of~$\pi_1^{\tame}(U_{\overline{s}_2})$ at~$y_2$.
\end{lemme}

Indeed, let $\overline{X}$ (\resp $\widetilde{X}$) be the strict
localization of~$X$ at $y_2$ (\resp at $y_1$), and put
$\overline{U}=U \times_X \overline{X}$ (\resp
$\widetilde{U}=U\times_X \widetilde{X}$). The chosen specialization induces a
morphism $\widetilde{U}\to \overline{U}$, and it follows
from~\Ref{XIII.2.10} that there is a commutative diagram
$$
\xymatrix{ \pi_1^{\tame} \left(\widetilde{U}_{\overline{s}_1} \right)
\ar[r]^{\pi '} \ar[d] & \pi_1^{\tame} \left(
\overline{U}_{\overline{s}_2} \right) \ar[d] \\ \pi_1^{\tame} \left(
U_{\overline{s}_1} \right) \ar[r]^{\pi } & \pi_1^{\tame} \left(
U_{\overline{s}_2} \right) }
$$
where $\pi'$ is the composite of the canonical homomorphism
$\pi_1^{\tame}(\widetilde{U}_ {\overline{s}_1})\to
\pi_1^{\tame}(\overline{U}_{\overline{s}_1})$ and the specialization
homomorphism. Since
$\pi_1^{\tame}(\widetilde{U}_{\overline{s}_1})$
(\resp $\pi_1^{\tame}(\overline{U}_{\overline{s}_2})$)
is an inertia group of
$\pi_1^{\tame}(U_{\overline{s}_1})$ at~$y_1$ (\resp of
$\pi_1^{\tame}(U_{\overline{s}_2})$ at~$y_2$), it suffices to prove
that $\pi'$ is surjective. This follows from the expression obtained
in~\Ref{XIII.5.6}.

\begin{corollaire}
\label{XIII.2.12}
Let $X$ be a connected proper smooth curve of genus $g$ over a
separably closed field $k$ of characteristic $p\geq 0$. Let $U$ be the
open subset obtained by removing from $X$ the $n$ distinct closed points
$a_1,\dots,a_n$. Then the tamely ramified fundamental group
$\pi_1^{\tame}(U)$ \eqref{XIII.2.1.3} can be generated by $2g+n$
elements $x_i, y_i,\sigma_j$, with $1\leq i\leq g$,
$1\leq j\leq n$, such that $\sigma_j$ is a generator of an inertia
group corresponding to $a_j$, and the following relation holds:
\begin{equation*}
\label{eq:XIII.2.12.*}
\tag{$*$} \prod_{1\leq i\leq
g}(x_iy_ix^{-1}_iy^{-1}_i)\cdot\prod_{1\leq j\leq n}\sigma_j=1.
\end{equation*}
For every finite group $G$ \emph{of order prime to} $p$,
generated by elements $\overline{x}_i,
\overline{y}_i,\overline{\sigma}_j$ satisfying the relation
\eqref{eq:XIII.2.12.*}, there exists an \'etale covering of~$U$ with
group $G$, corresponding to a homomorphism
$\pi_1^{\tame}(U)\to G$ which
% page: 393
\marginpar{393}
%
sends $x_i, y_i,\sigma_j$ to~$\overline{x}_i,
\overline{y}_i,\overline{\sigma}_j$, respectively. In other words,
if $p'$ denotes the set of prime numbers distinct from~$p$, then
$\pi_1^{p'}(U)$ is the pro-$p'$-group generated by
$x_i, y_i,\sigma_j$ subject to the single relation
\eqref{eq:XIII.2.12.*}.
\end{corollaire}

\subsubsection*{Proof}
We may suppose that $k$ is algebraically closed. Suppose first that
$k$ has characteristic zero. There then exists an algebraically closed
subfield $k'$ of~$k$, of finite transcendence degree over~$\QQ$, such
that $X$ is obtained by extension of scalars from a proper smooth curve
$X'$ defined over~$k'$, and we may suppose that the points
$a_1,\dots,a_n$ come from rational points $a'_1,\dots,a'_n$ of~$X'$.
Since $k'$ has finite transcendence degree over~$\QQ$, there is an
embedding of~$k'$ into the field of complex numbers $\CC$; put
$\widetilde{U}=U'\times_{k'}\CC.$ Let $k''$ be an algebraically closed
extension of~$k'$ such that there are $k'$-morphisms from~$k$ and
from~$\CC$ to $k''$. If $g'\colon U'\to k'$ is the structural morphism,
and if $\cal{F}'$ is a finite constant sheaf of groups on~$U'$, it
follows from~\Ref{XIII.2.9} that the specialization morphisms
$$
(\R^1g'_*\cal{F}')_{\CC}\to(\R^1g'_*\cal{F}')_{k''}
\from(\R^1g'_*\cal{F}')_k
$$
are isomorphisms. In terms of fundamental groups, this shows that there
is an isomorphism, defined up to inner automorphism,
$$
\pi_1(U)\to\pi_1(\widetilde{U}),
$$
and it is clear that this isomorphism sends an inertia group relative
to a point of~$X'-U'$ to an inertia group relative to the same point.
We may therefore suppose that $k=\CC$. In this case, the Riemann
existence theorem (\ifthenelse{\boolean{orig}}{XII~}{}\Ref{XII.5.2}) shows that the fundamental group
$\pi_1(U)$ is precisely the completion, for the topology of subgroups
of finite index, of the fundamental group of the analytic space
associated with~$U$. The latter can be computed by transcendental
methods \cite[ch\ptbl 7
\S 47]{XIII.3};
% page: 394
\marginpar{394}
%
it can be generated by $2g+n$ elements $x_i,
y_i,\sigma_j$ such that $\sigma_j$ is the image of a generator of the
local fundamental group $\pi_1(D_j-\{a_j\})$ of a small disk centered at~$a_j$,
\ie a generator of an inertia group corresponding to the point $a_j$,
and these elements satisfy the single relation
\eqref{eq:XIII.2.12.*}.

Now suppose that $k$ has characteristic $p>0$. There exists a complete
discrete valuation ring $A$ with residue field $k$ and fraction field
$K$ of characteristic zero, and a connected scheme $X_1$, proper and
smooth over~$S=\Spec A$, such that
$X_1\times_S \Spec k\simeq X$ (\ifthenelse{\boolean{orig}}{III~}{}\Ref{III.7.4}). The points $a_j$ then
lift to sections $s_j$ of~$X_1$ over~$S$. Let $Y_{1j}$ be the reduced
closed subscheme with underlying space $s_j(S)$, let $Y_1$ be the union
of the $Y_{1j}$, put $U_1=X_1-Y_1$, and let
$g_1\colon U_1\to S$ be the structural morphism. Let $\overline{K}$ be
an algebraically closed extension of~$K$, and put
$\overline{X}=X_1\times_S\overline{K}$ and $\overline{U}= U_1\times_S\overline{K}$. If $C$ is a finite constant
group, it follows from~\Ref{XIII.2.8} that the specialization
morphism
$$
(\R^1g_{1*}C_{U_1})_k\to(\R^1g_{1*}C_{U_1})_{\overline{K}}
$$
is injective, and even bijective if $C$ has order prime to~$p$. In terms
of fundamental groups, this means that the specialization homomorphism
\eqref{XIII.2.10}
$$
\pi \colon \pi_1(\overline{U})\to\pi^{\tame}_1(U)
$$
is surjective, and that the specialization homomorphism
$$
\pi^{p'}_1(\overline{U})\to\pi^{p'}_1(U)
$$
is bijective. Finally, if $x_i, y_i,\sigma_j$ are generators of
$\pi_1(\overline{U})$ such that $\sigma_j$ is a generator of an inertia
group corresponding to the point $b_j=Y_{1j}(\overline{K})$
of~$\overline{X}$, then, by~\Ref{XIII.2.11}, $\pi(\sigma_j)$ is a
generator of an inertia group corresponding to $a_j$. This completes
the proof.


\ifthenelse{\boolean{orig}}
{}
{\begin{remarqueMR}
\label{XIII.2.13}
Let $k$ be an algebraically closed field of characteristic $p>0$, let
$X$ be a connected proper smooth algebraic curve over~$k$ of genus $g$,
and let $U$ be an affine open subset of~$X$, the complement of $r\geq1$
rational points of~$X$. We have the fundamental group $\pi_1(U)$, its
quotient $\pi_1^{\tame}(U)$, which classifies the finite \'etale
coverings of~$U$ tamely ramified at the points of~$X-U$, and the
quotient $\pi_1^{p'}(U)$ of~$\pi_1^{\tame}(U)$, which classifies the
Galois \'etale coverings of~$U$ whose Galois group has order prime
to~$p$.

In (Coverings of algebraic curves, Amer.\ J.~Math.\ \textbf{79}
(1957), p\ptbl825--856), S\ptbl Abhyankar formulated a number of
conjectures on the structure of the finite groups that occur as Galois
groups of finite \'etale coverings of~$U$.

The conjectures concerning Galois groups of connected \'etale coverings
of~$U$ whose order is prime to~$p$ are proved
and made precise in Corollary~\Ref{XIII.2.12}. Before turning to the
finite quotients of~$\pi_1(U)$, let us first give some indications about
the \og size\fg of this group. Let $U = \Spec(A)$. We know that
$\Hom_{\mathrm{cont}}(\pi_1(U), \ZZ/p\ZZ)$ is described by
Artin--Schreier theory (Corollary\ptbl \ifthenelse{\boolean{orig}}{XI~}{}\Ref{XI.6.9}). This group is
isomorphic to $A/\wp(A)$, where $\wp$ is the map from~$A$ to $A$ that
sends $a$ to~$a^p-a$. Suppose first that $U$ is the affine line
$\AA^1$, with ring $k[T]$. Let~$E$ be the set of elements of~$k[T]$ of
the form $\sum_i a_iT^i$, with $a_i = 0$ whenever $p$ divides~$i$.
It is immediate that the composite map
$E\to k[T]\to k[T]/\wp k[T]$ is bijective. Hence the coefficients
$a_i$, $(i,p)=1$, behave like coordinates on a space parametrizing
the cyclic coverings of degree $p$ of the affine line. In particular,
it follows that $\pi_1(\AA^1)$ is not topologically finitely generated
and that, if one makes an extension $k\to k'$ of algebraically closed
fields, the natural homomorphism
$\pi_1(\AA^1\times_kk')\to\pi_1(\AA^1)$, which is surjective, is not
bijective, unlike in the proper case. In general, the curve $U$ can be
realized as a finite scheme over the affine line, and we conclude that
the same phenomena occur for $\pi_1(U)$ (and indeed, more generally,
for $\pi_1(V)$ for every connected affine $k$-scheme $V$ of finite type
and dimension $\geq1$).

With this said, if $G$ is a finite group, let $G^{(p')}$ denote the
largest quotient group of~$G$ whose order is prime to~$p$. For a finite
group $G$ to be a topological quotient of~$\pi_1(U)$, it is necessary
that $G^{(p')}$ be a topological quotient of~$\pi_1^{p'}(U)$, a
condition that can in principle be decided using
Corollary~\Ref{XIII.2.12}. In the article cited above, Abhyankar
conjectured that this necessary condition is also sufficient. This
conjecture was proved by M\ptbl Raynaud in the case of the affine line
and by D\ptbl Harbater in the general case (Invent. Math.\ \textbf{116}
(1994), p\ptbl425--462 and \textbf{117}, p\ptbl1--25). For example,
for the affine line $\pi_1^{p'}(\AA^1)=1$, and it follows that a finite
group $G$ is the Galois group of a connected \'etale covering of~$\AA^1$
if and only if $G^{(p')} = 1$, that is, if and only if $G$ is generated
by its Sylow $p$-subgroups. Thus every finite simple group whose order
is a multiple of~$p$ has the required property.
\end{remarqueMR}}
\end{filecontents*}
% Original input: drafts/SGA1_XIII_2_4_PROOF_English.texfrag
\begin{filecontents*}[overwrite,noheader]{sga1-cumulative-input-107.tex}
\Subsubsection*{Proof of~\Ref{XIII.2.4}~\textup{1)}}
\subsubsection{}
\label{XIII.2.4.1}
\emph{Case of a sheaf of sets satisfying} a). We may
suppose that $F=g^*G$, where $G$ is a sheaf on~$S$. It
then follows from (SGA~4 XVI~3.2) that the canonical morphism
\begin{equation*}
\label{eq:XIII.2.4.1.1}
\tag{\thesubsubsection.1} f^*G\to i_*F
\end{equation*}
is an isomorphism. For every $S$-scheme $h\colon S'\to S$, there is
similarly an isomorphism $f'^*G'\to i_*'F'$; consequently, the canonical
morphism
$$
\varphi\colon k^*(i_*F)\to i_*'F'
$$
is identified with the natural isomorphism
$$
k^*f^*G\simeq f'^*G'
$$

\subsubsection{}
\label{XIII.2.4.2}
\emph{Case of a sheaf of sets satisfying} b). We must
show that $\varphi$ is an isomorphism, and it is enough to
check this at every geometric point
$\overline{x}'$ of~$X'$. Let $\overline{S}$ (\resp $\overline{X}$,
\resp $\overline{S'}$, \resp $\overline{X'}$) be the strict localization
of~$S$ (\resp~$X$, \resp~$S'$, \resp $X'$) at the corresponding image of $\overline{x}'$, and put
$\overline{U}=U_{(\overline{X})}$,
$\overline{U'}=U_{(\overline{X'})}$, etc. The morphism
$\varphi_{\overline{x}'}$ is identified with the canonical morphism
$$
\overline{\varphi}\colon \H^0(\overline{U},\overline{F})\to
\H^0(\overline{U'}, \overline{F'})
$$
One can find a principal covering $V$ of~$\overline{U}$, of the
type appearing in~\Ref{XIII.5.4}, such that the inverse image of
$\overline{F}$ on~$V$ is a constant sheaf with value $C$. If
$\Pi$ is the Galois group of~$V$ over~$\overline{U}$, then $\Pi$
acts on~$\overline{F}|V$, and there is an isomorphism
\begin{equation*}
\label{eq:XIII.2.4.2.1}
\tag{\thesubsubsection.1}
\ifthenelse{\boolean{orig}}
{\H^0(U,F))\simeq \H^0(V,C_V)^{\Pi}}
{\H^0(\overline{U},\overline{F})\simeq \H^0(V,C_V)^{\Pi}}
% correction: \H^0(\overline{U},\overline{F}))\simeq \H^0(V,C_V)^{\Pi}
\end{equation*}
\label{indnot:mh}\oldindexnot{$\H^0(V,C_V)^{\Pi}$|hyperpage}%
%
where the right-hand side denotes the set of
elements of~$\H^0(V,C_V)$ invariant
% page: 382
\marginpar{382}
%
under~$\Pi$. Since $V'=V\times_{\overline{U}}\overline{U'}$ is a
principal covering of~$\overline{U'}$ with Galois group
$\Pi'\simeq \Pi$, one sees that the morphism $\overline{\varphi}$
is obtained, by taking the invariants under $\Pi$, from the
canonical morphism
$$
\H^0(V,C_V)\to \H^0(V',C_{V'}).
$$
Since $V$ and $V'$ are connected \eqref{XIII.5.4}, this morphism,
and hence also $\overline{\varphi}$, is an isomorphism.

Note that if, in addition, $F$ is a sheaf of groups and $P$ is a
torsor on~$\overline{U}$ with group $\overline{F}$,
tamely ramified relative to $D$, it follows
from the preceding proof and from~\Ref{XIII.2.2} that
$(\leftexp{P}{\overline{F}},\overline{i})$ is cohomologically
proper relative to~$S$ in dimension $\leq 0$.

\subsubsection{}
\label{XIII.2.4.3}
\emph{Case of a sheaf of groups.} To show that $\Psi$ is an
isomorphism, it suffices to prove that, for every geometric
point $\overline{y}'$ of~$Y'$, the morphism
$$
\Psi_{\overline{y}'}\colon
(k^*(\R^1_{\tame}i_*F))_{\overline{y}'}\to (\R^1_{\tame}
i_*'F')_{\overline{y}'}
$$
is an isomorphism. Now, by~\Ref{XIII.2.1.2},
$\Psi_{\overline{y}'}$ is identified with the canonical morphism
$$
\overline{\Psi}\colon
\H^1_{\tame}(\overline{U},\overline{F})\to
\H^1_{\tame}(\overline{U'},\overline{F'}).
$$
Let $\widetilde{U}$ be the tamely ramified universal covering
of~$\overline{U}$ \eqref{XIII.2.1.3}, and let $\widetilde{F}$
be the inverse image of~$\overline{F}$ on~$\widetilde{U}$. It follows
from~\Ref{XIII.5.7} in case a) and from~\Ref{XIII.5.5} in case b)
that $\H^1_{\tame}(\overline{U},\overline{F})$ is identified with the
subset of $\H^1(\overline{U},\overline{F})$ consisting of the classes
of~$\overline{F}$-torsors whose inverse image on~$\widetilde{U}$
is trivial. On the other hand, a
\ifthenelse{\boolean{orig}}
{standard}
{standard (\cf \ifthenelse{\boolean{orig}}{IX~}{}\Ref{IX.5}, p\ptbl\pageref{page-300})}
argument shows that
the set of elements of~$\H^1(\overline{U},\overline{F})$
whose inverse image on~$\widetilde{U}$ is trivial is identified
with
$\H^1(\pi_1^{\tame}(\overline{U}),\H^0(\widetilde{U},\widetilde{F})).$
Thus one obtains a canonical isomorphism
\begin{equation*}
\label{eq:XIII.2.4.3.1} \tag{\thesubsubsection.1} {}
\H^1_{\tame}(\overline{U},\overline{F})\isomto
\H^1(\pi_1^{\tame}(\overline{U}),\H^0(\widetilde{U},\widetilde{F})).
\end{equation*}

Consequently, the morphism $\overline{\Psi}$ is identified with the
canonical morphism
$$
\H^1(\pi_1^{\tame}(\overline{U}),\H^0(\widetilde{U},\widetilde{F}))
\to
\H^1(\pi_1^{\tame}(\overline{U'}),\H^0(\widetilde{U}',\widetilde{F}')).
$$
Let us show
% page: 383
\marginpar{383}
%
that this morphism is an isomorphism. The morphism
$\pi_1^{\tame}(\overline{U'})\to\pi_1^{\tame}(\overline{U})$
is an isomorphism by~\Ref{XIII.5.6}, and the same is true of
the morphism $\H^0(\widetilde{U},\widetilde{F})\to\allowbreak
\H^0(\widetilde{U}',\widetilde{F}')$. Indeed, this is clear
in case b), since $\widetilde{F}$ is constant and $\widetilde{U}$ and
$\widetilde{U}'$ are connected. In case a), let $\overline{G}$ be a
constructible sheaf of groups on~$\overline{S}$ such that
$\overline{F}=\overline{g}^*\overline{G};$ since the morphisms
$\widetilde{U}\to\overline{S}$ and
$\widetilde{U}'\to\overline{S'}$ are $0$-acyclic
\eqref{XIII.5.7}, there are isomorphisms
$$
\H^0(\widetilde{U},\widetilde{F})\simeq
\H^0(\overline{S},\overline{G}) \simeq
\H^0(\overline{S'},\overline{G'})\simeq
\H^0(\widetilde{U}',\widetilde{F}'),
$$
which implies that $\Psi$ is an isomorphism. The last
assertion of \Ref{XIII.2.4}~1) follows from the foregoing,
in view of \Ref{XIII.2.3}~b).

\subsubsection*{Proof of~\Ref{XIII.2.4}~\textup{2)}}

The case of a constructible sheaf of sets satisfying a)
follows immediately from \eqref{eq:XIII.2.4.1.1}. Let $F=g^*G$ be a
sheaf of groups satisfying a), where $G$ is a constructible sheaf;
we may suppose that $S$ is affine. Let $(S_j)_{j\in J}$
be a finite family of reduced locally closed subschemes of~$S$
whose union covers $S$, such that the inverse image of~$G$ on~$S_j$ is a locally constant sheaf. In view of
\Ref{XIII.2.4}~1), in order to prove that $\R^1_{\tame}i_*F$
is constructible, it suffices to prove this after the base change
$S_j\to S$, for every $j\in J$. We are therefore
reduced to case~b), where $F$ is locally constant.

From now on, suppose that $F$ is a sheaf of sets or of
groups satisfying b). Since the question is local for the
\'etale topology on~$X$, we may suppose that $X$ is of finite
presentation over~$S$, and, by passage to the limit, we may suppose that $X$
and $S$ are noetherian.

Let $D=\sum_{1\leq i\leq r}\divisor f_i$, where, for every point
$x$ of~$\Supp D$, if $I(x)$ is the set of $i$ such that
$f_i(x)=0$, the closed subscheme $V((f_i)_{i\in I(x)})$ is smooth over~$S$ of codimension $\card I(x)$ in~$X$. Let $\mathcal{P}$
be the set of subsets of~$[ 1,r ]$, and, for every $I\in\mathcal{P}$,
put
$$
X_I=\textstyle\Big(\bigcap_{i\in I}V(f_i)\Big)\cap\Big(\bigcap_{i\not\in I}X_{f_i}\Big).
$$
Let
% page: 384
\marginpar{384}
%
$z$ be a point of~$X_I$. After first passing to an
\'etale neighborhood of~$z$, one can find an open subset $W$ of~$X$
containing $z$ and a principal covering $V$ of~$U\cap W$,
tamely ramified on~$W$ relative to~$S$, of the type
considered in~\Ref{XIII.5.6.1}, such that the inverse
image of~$F$ on~$V$ is a constant sheaf with
value~$C$. Let $\pi$ be the Galois group of the covering~$V$. For
every geometric point $\overline{x}$ of~$X_I\cap W$, one then has,
by \eqref{eq:XIII.2.4.2.1},
$$
\H^0(\overline{U},\overline{F})\simeq \H^0(V,C_V)^{\pi}.
$$
It follows that $i_*F|X_I\cap W$ is locally constant (SGA~4
IX~2.13), and hence $i_*F$ is constructible.

Finally, let us show that if $F$ is a locally
constant sheaf of groups, then $\R^1_{\tame}i_*F|_{X_I}$ is constructible. If
$\overline{x}$ is a geometric point of~$X_I$, we have obtained
in \eqref{eq:XIII.2.4.3.1} the expression
$$
\H^1_{\tame}(\overline{U},\overline{F})\isomto
\H^1(\pi^{\tame}_1(\overline{U}),\H^0(\widetilde{U},\widetilde{F})).
$$
If $p$ is the residue characteristic of~$\overline{X}$, then
by~\Ref{XIII.5.6},
$\pi_1^{\tame}(\overline{U})=\prod_{l\neq p }\ZZ_l(1)^{\card I}$. Let $\LL$ be the set of prime numbers that
divide the order of the finite group $\H^0(\widetilde{U},\widetilde{F})$, and
let $K=\prod_{l\in\LL -\{p\}\cap\LL }\ZZ_l(1)^{\card I}$. It
follows from \cite[I \S5 ex.2]{XIII.4} that there is an isomorphism
$$
\H^1_{\tame}(\overline{U},\overline{F})\simeq
\H^1(K,\H^0(\widetilde{U},\widetilde{F})).
$$
Since $K$ is topologically finitely generated and
$\H^0(\widetilde{U},\widetilde{F})$ is finite, we first deduce
that the fibers of the sheaf $\R^1_{\tame}i_*F|_{X_I}$ are
finite. Moreover, the set $\LL$ does not depend on the point
$\overline{x}$. For every $q\in\LL$, let $X_{I,q}$ be the closed subset of
$X_I$ defined by the equation $q=0$, and let $X_I'$ be the open subset of~$X_I$
complementary to the union of the $X_{I,q}$. Then
$\R^1_{\tame}i_*F|_{X_{I,q}}$ and $\R^1_{\tame}i_*F|_{X_I'}$ are
locally constant; indeed, a specialization morphism between
geometric points of~$X_{I,q}$ (\resp of~$X_I'$) induces an
isomorphism on the groups $K$ \eqref{XIII.5.6.1}, hence also on the
sets $\H^1_{\tame}(\overline{U},\overline{F})$, and one may apply
SGA~4 IX~2.13.
\end{filecontents*}
% Original input: drafts/SGA1_XIII_2_5_THROUGH_2_9_English.texfrag
\begin{filecontents*}[overwrite,noheader]{sga1-cumulative-input-108.tex}
\begin{corollaire}
\label{XIII.2.5}
Let
% page: 385
\marginpar{385}
%
$f\colon X\to S$ be a morphism, let $D$ be a divisor on~$X$
with normal crossings relative to~$S$ \eqref{XIII.2.1}, let
$Y=\Supp D$, $U=X- Y$, and let $i\colon U\to X$ be the canonical
immersion. Let $\Phi$ be a stack on~$U$, and suppose that we are given
surjective \'etale morphisms $S_1\to S$ and $X_2\to
X\times_{S}S_1$, a stack $\Psi$ on~$S_1$, and an isomorphism
$\Phi|U_2\simeq\Psi|U_2$ (\cf \Ref{XIII.2.2}).

Then, for every morphism $h\colon S'\to S$, if $k=h_{(X)}$, the
canonical functor
$$
\varphi\colon k^*i_*\Phi\to i_*'\Phi'
$$
is fully faithful. If $\Psi$ is constructible, the canonical
functor
$$
\psi\colon k^*i_*^{\tame}\Phi\to i_*^{\prime\tame}\Phi'
$$
is an equivalence of categories.

Moreover, if the stack $\Psi$ is constructible (\resp if $\Psi$ is
$1$-constructible \eqref{XIII.0} and $S$ is locally noetherian),
% remark: include some remark here that the 0 is a reference to
% XIII.0.?
$i_*\Phi$ is constructible (\resp $i_*^t\Phi$ is
$1$-constructible).
\end{corollaire}

Let us show that $\varphi$ is fully faithful. It is enough to see
that, for every geometric point $\overline{x}'$ of~$X'$, the same is
true of the functor
$$
\overline{\varphi}\colon
\Phi(\overline{U})\to\Phi'(\overline{U'})
$$
(we again use the notation of~\Ref{XIII.2.4.2}). Let $a$, $b$ be two
objects of~$\Phi(\overline{U})$, and let $a',\ b'$ be their images under
$\overline{\varphi}$. Since the morphism
$\overline{U}\to\overline{S}$ is locally $0$-acyclic
\eqref{XIII.5.7}, we have an isomorphism
$$
\H^0(\overline{U},\overline{S\Phi})
\overset{\sim}{\to}\H^0(\overline{S},\overline{S\Psi}).
$$
Consequently, $a$ and $b$ are inverse images of objects
of~$\Psi$, and the same is therefore true of
$F=\SheafHom_{\overline{U}}(a,b)$. Since the sheaf
$F'=\SheafHom_{\overline{U'}}(a',b')$ is the inverse image on~$\overline{U'}$ of~$F$, it follows from~\Ref{XIII.2.4.1} that the
canonical morphism
$$
\H^0(\overline{U},F)\to \H^0(\overline{U'},F')
$$
is an isomorphism, which proves that $\overline{\varphi}$ is
fully faithful.

Let us show
% page: 386
\marginpar{386}
%
that, for every geometric point $\overline{x}'$ of~$X'$, the
functor
$$
\overline{\psi}\colon i^{\tame}_*\Phi(\overline{X'})\to
i_*^{\prime\tame}\Phi'(\overline{X'})
$$
is an equivalence. By what precedes,
$\overline{\psi}$ is fully faithful. Let us show that
$\overline{\psi}$ is essentially surjective. Let $a'$ be an
object of~$\Phi'(\overline{U'})$ tamely
ramified on~$X'$ relative to~$S'$ \eqref{XIII.2.2}, and
let us show that it is the image of a tamely
ramified object of~$\Phi(\overline{U})$. It follows from~\Ref{XIII.2.4}
1) that the canonical morphism
$$
\H^0(\overline{U},\overline{S}\overline{\Phi})\to
\H^0(\overline{U'},\overline{S}\overline{\Phi}')
$$
is an isomorphism. Let $G'$ be the maximal subgerbe of~$\Phi'$
generated by $a'$; there is then a maximal subgerbe $G$ of
$\overline{\Phi}$, the inverse image of a gerbe on~$\overline{S}$, such
that
$$
\overline{m}^*G\simeq G',
$$
where $\overline{m}$ is the morphism $\overline{U'}\to
\overline{U}$. The canonical functor
\begin{equation*}
\label{eq:XIII.2.5.*}
\tag{$*$} {\overline{k}^*\overline{i}_*^{\tame}G\to
\overline{i}_*^{\prime\tame}G'}
\end{equation*}
is an equivalence, because $G$ identifies with a gerbe of
torsors under a constructible sheaf of groups coming from
$\overline{S}$, and we can apply \Ref{XIII.2.4}~1). It
then follows from \eqref{eq:XIII.2.5.*} that there exists an
object $a$ of~$G(\overline{U})$, tamely
ramified on~$X$ relative to~$S$, whose inverse image on~$\overline{U'}$ is $a'$. This proves that $\psi$ is an
equivalence.

If $\Psi$ is constructible, then so is~$i_*\Phi$; indeed,
an object $x$ of~$i_*\Phi$ is, locally for the \'etale topology
on~$S$, the inverse image of an object $y$ of~$\Psi$; it
therefore follows from~\Ref{eq:XIII.2.4.1.1} that $\SheafAut (x)$ is
the inverse image of~$\SheafAut (y)$, and hence is constructible. Finally, if
$\Psi$ is $1$-constructible, then so is
$i^{\tame}_*\Phi$ by~\Ref{XIII.6.3} below.

\begin{corollaire}
\label{XIII.2.6}
The
% page: 387
\marginpar{387}
%
notation is that of~\Ref{XIII.2.4}. Suppose that $S$ has
characteristic zero at every point $s$ such that $Y_s\neq
\emptyset.$ Then, if $\cal{ F}$ is a locally constant constructible
sheaf of groups on~$U$ (\resp a constructible stack on~$U$ which,
locally on~$X$ and~$S$, is the inverse image of a constructible stack
on~$S$), $(\cal{ F},i)$ is cohomologically proper relative
to~$S$ in dimension $\leq 1.$
\end{corollaire}

Since every constructible sheaf of sets on~$U$ is
tamely ramified on~$X$ relative to~$S,$ the
corollary follows from~\Ref{XIII.2.4} (\resp \Ref{XIII.2.5}).

\begin{corollaire}
\label{XIII.2.7}
The notation is that of~\Ref{XIII.2.4}, but in addition we are given
an $S$-scheme $T$, a proper morphism $p\colon X\to T,$
and we suppose that $X$ and $T$ are of finite presentation over~$S$; let
$q=pi$. Let $\cal{ F}$ be a constructible sheaf of sets on~$U$
satisfying one of conditions \textup{a)} or \textup{b)} of~\Ref{XIII.2.4}
(\resp a sheaf of groups satisfying one of conditions
\textup{a)} or \textup{b)} of~\Ref{XIII.2.4}, \resp a stack on~$U$ which, locally on~$X$ and~$S$, is the inverse image of a constructible stack $G$ on~$S$). Then
the following conclusions hold:
\begin{enumerate}
\item[1)] $(\cal{ F},q)$ is cohomologically proper relative to~$S$ in dimension $\leq 0$ (\resp for every morphism $h\colon
S^{'}\to S$, if $m=h_{(T)}$, the canonical morphism:
$$
\Theta\colon m^*(\R^1_{\tame}q_*\cal{ F})\to
\R^1_{\tame}q_*^{'}\cal{ F}^{'}
$$
is an isomorphism, \resp for every morphism $h\colon S^{'}\to S$, if $m=h_{(T)}$,
the canonical morphism
$$
\xi \colon m^*(q_*^{\tame}\cal{ F})\to
{q'}_*^{{\tame}}\cal{F}^{'}
$$
is an equivalence).
\item[2)] The sheaf $q_*\cal{ F}$ (\resp the sheaf
$\R_{\tame}^1q_*\cal{ F}$, \resp the stack $q_*^{\tame}\cal{ F}$)
is constructible. In the last case, if $S$ is assumed locally
noetherian and $G$ is $1$-constructible, then so is
$q_*^{\tame}\cal{ F}.$
\end{enumerate}
\end{corollaire}

The
% page: 388
\marginpar{388}
%
first part follows immediately from~\Ref{XIII.2.4},
\Ref{XIII.2.5}, and the proof
of~\Ref{XIII.1.8}. Let us prove~2). If $\cal{ F}$ is a constructible
sheaf of sets on~$U$ satisfying \Ref{XIII.2.4}~a)
or \Ref{XIII.2.4}~b), it follows from \Ref{XIII.2.4}~2) that
$i_*\cal{ F}$ is constructible; hence so is
$q_*\cal{ F}=p_*(i_*\cal{ F})$ (SGA~4 XIV~1.1).

Let $\cal{F}$ be a constructible sheaf of groups on~$U$
satisfying \Ref{XIII.2.4}~a) or \Ref{XIII.2.4}~b), and let us prove
that $\R_{\tame}^1q_*\cal{ F}$ is constructible. By passage to
the limit (EGA~IV 8.10.5 and 17.7.8) and using 1), we may
suppose that $S$ is noetherian. Let $\Phi$ then be the stack on~$X$ whose
fiber over every scheme $X^{'}$ \'etale over~$X$ consists
of the torsors on~$U^{'}=U\times_{X}X^{'}$, with group $\cal{F}|U'$, which
are tamely ramified on~$X'$ relative to~$S$; thus
$$
S(\Phi)\simeq \R_{\tame}^1i_*\cal{ F}
$$
and this sheaf is constructible by \Ref{XIII.2.4}~2).
% correction: The reference was wrong.
It follows from~\Ref{XIII.6.2} below that $S(p_*\Phi)$
is constructible, \ie that $\R^1_{\tame}q_*\cal{ F}$ is
constructible.

Finally, if $\cal{ F}$ is a stack on~$U$ which, locally on~$X$ and
$S$, is the inverse image of a constructible stack on~$S,$
$i_*^{\tame}\cal{ F}$ is constructible, and hence so is
$q_*^{\tame}\cal{ F}=p_*(i_*^{\tame}\cal{ F})$. If
in addition $S$ is locally noetherian and $SG$ is constructible,
$S(i_*^{\tame}\cal{ F})$ is constructible by
\hyperref[XIII.6.1.2]{6.3}; hence
the same is true of~$S(p_*i_*^{\tame}\cal{ F})$, \ie of
$S(q_*^{\tame}\cal{ F})$~\Ref{XIII.6.2}.

\begin{corollaire}
\label{XIII.2.8}
Let
$$
\xymatrix{ U \ar[r]^{i} \ar[d]_{g} & X \ar[ld]^{f} \\ S }
$$
be a commutative diagram of schemes in which $U$ is the open
complement
% page: 389
\marginpar{389}
%
in $X$ of a divisor with normal crossings relative to~$S$, and $f$ is a \emph {proper} morphism of finite presentation. Let
$\LL$ be a set of prime numbers. Suppose that~$g$ is locally
$0$-acyclic (\resp locally $1$-aspherical for $\LL$). Then, if
$\cal{ F}$ is a sheaf of sets on~$U$ (\resp a sheaf of
$\LL$-groups) that is locally constant and constructible on~$U,$ and
tamely ramified on~$X$ relative to~$S$,
$g_*\cal{ F}$ (\resp $\R^1_{\tame}g_*\cal{ F}$) is locally
constant constructible, and $(\cal{ F},\ifthenelse{\boolean{orig}}{f}{g})$ is cohomologically proper
relative to~$S$ in dimension $\leq 0$ (\resp the formation of
$\R^1_{\tame}g_*\cal{ F}$ commutes with every base change
$S^{'}\to S$). In the non-respective case, if $\cal{ F}$ is a
sheaf of groups, then for every specialization $\overline{s}_1
\to \overline{s}_2$ of geometric points of~$S$, the
specialization morphism
$$
(\R^1_{\tame}g_*\cal{ F})_{\overline{s}_2}\to
(\R^1_{\tame}g_*\cal{ F})_{\overline{s}_1}
$$
is injective.
\end{corollaire}
The corollary follows immediately from~\Ref{XIII.2.7} and
from~\Ref{XIII.1.14} (\resp from the analogue of~\Ref{XIII.1.14} for
$\R^1_{\tame}g_*\cal{ F}$, which is proved as in \loccit).

\begin{corollaire} \label{XIII.2.9}
Let
$$
\xymatrix{ U \ar[r]^{i} \ar[d]_{g} & X \ar[ld]^{f} \\ S }
$$
be a commutative diagram of schemes in which $U$ is the open
complement in~$X$ of a divisor with normal crossings
relative to~$S$, and $f$ is a proper smooth morphism of
finite presentation. Let~$\LL $ be the set of prime numbers
distinct from the residue characteristics of~$S$. Let~$\cal{F}$ be a locally constant constructible sheaf of~$\LL $-groups
on~$U$, tamely ramified on~$X$ relative
to~$S$. Then $\R^1 g_*\cal{ F}$ is locally constant
constructible, and $(\cal{ F},\ifthenelse{\boolean{orig}}{f}{g})$ is cohomologically proper
relative
% page: 390
\marginpar{390}
%
to~$S$ in dimension $\leq 1$.
\end{corollaire}
The corollary follows from~\Ref{XIII.2.8} and from the fact that
$\R_{\tame}^1g_*\cal{ F}=\R^1g_*\cal{ F}$ (\Ref{XIII.2.3}~b)).
\end{filecontents*}
% Original input: drafts/SGA1_XIII_3_CASE2_AND_CASE3_A_English.texfrag
\begin{filecontents*}[overwrite,noheader]{sga1-cumulative-input-109.tex}
\Subsection*{2) Case of a constructible sheaf of sets}

Let $F$ be a constructible sheaf of sets on~$X$. By
SGA~4 IX 2.14~(ii), one can find a finite family of morphisms
$p_i \colon Z_i \to X$ and, on each $Z_i$, a finite constant sheaf of sets
$C_i$, such that there is a monomorphism
% page: 399
\marginpar{399}
%
$$
j \colon F \to \prod_{i} {p_i}_* C_i =G \quoi.
$$
By~1)\kern2pt a),
after restricting $S$ to a nonempty open subset, we may suppose
that the $(C_i,fp_i)$ are cohomologically proper relative
to~$S$ in dimension $\le 0$, and that the $f_*{p_i}_*C_i$ are
constructible. We already conclude that $(G,f)$ is
cohomologically proper relative to~$S$ in dimension $\le 0$
\eqref{XIII.1.9}, hence that $(F,f)$ is cohomologically proper
relative to~$S$ in dimension $\le -1$, and that $f_*F$ is
constructible. Let~$K$ be the amalgamated sum $K= G \coprod_{F} G$;
since $F$ and $G$ are constructible, so is~$K$. We
therefore conclude from the preceding that, after restricting
$S$ to a nonempty open subset, we may suppose that $(K,f)$ is
cohomologically proper relative to~$S$ in dimension $\le
-1$. It then follows from \Ref{XIII.1.13}~1) that $(F,f)$ is
cohomologically proper relative to~$S$ in dimension $\le 0$.

\Subsection*{3) Case of constant sheaves of groups}

If $F$ is a constant sheaf of groups on~$X$, we denote by $\Phi$ the
stack of torsors under $F$.

3) 1. Let us first show that, after restricting $S$ to a
nonempty open subset, for every constant sheaf of groups $F$ on~$X$
whose order is prime to the residue characteristics of~$S$,
$(F,f)$ is cohomologically proper relative to~$S$ in
dimension $\le 0$, and $f_*\Phi$ is constructible.

For this, we reduce to the case where $X$ is smooth over~$S$. After
making a finite extension of~$\kres(s)$, which is permissible because it
can be regarded as the composite of an \'etale extension and
a radicial extension, one can find a proper surjective morphism
$p_s \colon X_{1s}\to X_s$, where $X_{1s}$ is a
scheme smooth over~$s$ of the same dimension as $X_s$, and, after
restricting $S$ to a nonempty open subset, we may suppose that
$p_s$ comes from a proper surjective morphism $p \colon X_1 \to X$,
where $X_1$ is a scheme smooth over~$S$
% page: 400
\marginpar{400}
%
(EGA~IV 9.6.1 and 12.1.6). Let $X_2=X_1 \times_X X_1$, and let $q \colon
X_2 \to X$ be the canonical morphism. There is an exact diagram of stacks
on~$X$
$$
\xymatrix@C=.5cm{\Phi \ar[r] & p_*p^* \Phi \ar@<2pt>[r]\ar@<-2pt>[r] &
q_*q^*\Phi }
$$
(\Ref{XIII.1.11}~2)). Assuming that the theorem has been
proved in the smooth case, we first see that we may
suppose $f_*p_*p^*\Phi$ constructible; hence the same is true of
$f_*\Phi$ (\Ref{XIII.3.1.1} below). Moreover, by \Ref{XIII.1.6}~2),
we may suppose $(p_*p^*\Phi,f)$ cohomologically proper
relative to~$S$ in dimension $\le 0$; it follows that
$(\Phi,f)$ is cohomologically proper relative to~$S$ in
dimension $\le -1$. We may therefore suppose that $(q_*q^*\Phi,f)$ is
cohomologically proper relative to~$S$ in dimension $\le -1$,
and this implies that $(\Phi,f)$ is cohomologically proper
relative to~$S$ in dimension $\le 0$ (\Ref{XIII.1.12}~1)).

We then reduce, as in 1)\kern2pt2, to the case where $X$ is smooth
and affine over~$S$. Let
$$
\xymatrix{ X \ar[r]^{i} \ar[d]_{f} & P \ar[ld]^{q} \\ Y }
$$
be a compactification of~$X$, where $i$ is a dominant open immersion
and $q$ is a proper morphism. Since $\dim P_s=\dim X_s$,
we may apply the resolution of singularities hypothesis
to $P_{\overline s}$. After making an
\'etale extension and a radicial extension of~$S$, one can
find a proper morphism $r \colon Z \to P$, where $Z$ is smooth
over~$S$, $r^{-1}(X)\simeq X$, and $r^{-1}(X)$ is the
complement in $Z$ of a divisor with normal crossings
relative to~$S$. Every torsor under $F$ is then
tamely ramified on~$Z$ relative to~$S$
(\Ref{XIII.2.3}~b)). It therefore follows from~\Ref{XIII.2.7} that $(F,f)$
is cohomologically proper relative to~$S$ in dimension $\le
0$, which proves our assertion.

\ifthenelse{\boolean{orig}}
{3) \textbf{2. Reduction\marginpar{401} to the case where $X$ is smooth over~$S$}}%
{\pagebreak[4] 3) 2. Reduction\marginpar{401} to the case where $X$ is smooth over~$S$.}%
% page: 401
%
\par\nopagebreak
After making a finite extension of~$\kres(s)$, one can find a
proper surjective morphism $p_s\colon X_{1s}\to X_s$, where~$X_{1s}$
is a scheme smooth over~$s$, and we may suppose that $p_s$
comes from a proper surjective morphism $p\colon X_1 \to X$, where
$X_1$ is smooth over~$S$. Suppose that the theorem has been
proved for $fp$, and let us prove it for $f$. Let $F$ be a
finite constant sheaf of groups on~$X$, of order prime to the
residue characteristics of~$S$, and let $\Phi$ be the stack of
torsors under $F$. Let $X_2=X_1 \times_X X_1$, $X_3=X_1 \times_X
X_1 \times_X X_1$, and let $q\colon X_2 \to X$, $r\colon X_3 \to X$ be the
canonical morphisms. By \Ref{XIII.1.11}~2), there is an
exact diagram of stacks
$$
\xymatrix@C=.5cm{\Phi \ar[r] & p_*p^*\Phi \ar@<2pt>[r]\ar@<-2pt>[r] &
q_*q^*\Phi \ar@<4pt>[r] \ar[r] \ar@<-4pt>[r] & r_*r^*\Phi}\quoi.
$$
To prove that $(\Phi,f)$ is cohomologically proper relative
to~$S$ in dimension $\le 1$, it is enough to show that the same is
true of~$(p_*p^*\Phi,f)$, that $(q_*q^*\Phi, f )$ is
cohomologically proper relative to~$S$ in dimension $\le 0$,
and that $(r_*r^*\Phi,f)$ is cohomologically proper relative
to~$S$ in dimension $\le -1$ (\Ref{XIII.1.12}~2)). By~3)\kern.3em 1
above, we may suppose that, for every finite constant sheaf of groups
$F$, $(q^*\Phi,fq)$, $(q^*\Phi,q)$, $(r^*\Phi,fr)$, $(r^*\Phi,r)$
are cohomologically proper relative to~$S$ in dimension
$\le 0$. It then follows from \Ref{XIII.1.6}~2) that
$(q_*q^*\Phi,f)$ and $(r_*r^*\Phi,f)$ are cohomologically proper
relative to~$S$ in dimension $\le 0$. Since the theorem
is assumed proved in the smooth case, $(p^*\Phi,
fp)$ and $(p^*\Phi,p)$, and hence also $(p_*p^*\Phi,f)$ (\Ref{XIII.1.6}~2)),
are cohomologically proper relative to~$S$ in dimension
$\le 1$. This indeed shows that $(\Phi,f)$ is cohomologically proper
relative to~$S$ in dimension $\le 1$.

Moreover, $f_*p_*p^*\Phi$ is $1$-constructible by hypothesis;
by~3)\kern.3em 1 above, we may suppose that $f_*q_*q^*\Phi$ is
constructible; it therefore follows from~\Ref{XIII.3.1.1} below that
$f_*\Phi$ is $1$-constructible.

3) 3. Reduction to the case where $X$ is smooth and affine over~$S$.
% page: 402
\marginpar{402}
%

By 3)\kern2pt2, we may suppose that $X$ is smooth over~$S$. Let $(U_i)_{i\in
I}$ be a finite covering of~$X$ by affine open subsets, let $X_1$
be the disjoint union of the $U_i$, and let $p\colon X_1 \to X$ be the
canonical morphism. Since $p$ is an effective descent morphism for the
category of \'etale sheaves of finite type on
variable schemes, we see as in 3)\kern2pt2 that, if the
theorem is assumed proved for the $U_i$, \ie for $X_1$, it
is also true for~$X$.

3) 4. The case where $X$ is smooth and affine over~$S$.

As in 3)1, after restricting $S$ to a
nonempty open subset and making a surjective \'etale extension and a
surjective radicial extension of~$S$, one can find a commutative
diagram
$$
\xymatrix{ X \ar[r]^{i} \ar[d]_{f} & P \ar[ld]^{q} \\ Y }
$$
where $P$ is a scheme smooth over~$S$, $X$ is the complement
in $P$ of a divisor with normal crossings relative to~$S$, and $q$ is a proper morphism. If $F$ is a constant sheaf
of order prime to the residue characteristics of~$S$, every
torsor under $F$ is tamely ramified on~$P$
relative to~$S$ (\Ref{XIII.2.3}~b)). The fact that $(F,f)$ is
cohomologically proper relative to~$S$ in dimension $\le 1$
and that $f_*\Phi$ is constructible then follows
from~\Ref{XIII.2.7}.

\ifthenelse{\boolean{orig}}
{4) \textbf{Proof of } 2)\kern2pt b)}
{\Subsection*{4) Proof of 2)\kern2pt b)}}

4) 1. The case where $\Phi$ is a gerbe.

One can find a surjective \'etale morphism of finite type $p
\colon X_1 \to X$ such that $p^*\Phi$ is a trivial gerbe. By
descent, as in 3)\kern2pt2, we see that it is enough to prove the
theorem for $X_1$, $X_1\times_X X_1$, and
% page: 403
\marginpar{403}
%
$X_1\times_X X_1\times_X X_1$. We are therefore reduced to the case where
$\Phi$ is the gerbe of torsors under a constructible sheaf of groups
$F$ whose fibers have order prime to the
residue characteristics of~$S$.

By SGA~4 IX~2.14, one can find a finite family of
finite morphisms $p_i \colon Z_i \to X$ and, for each $i$, a
finite constant sheaf of groups $C_i$, of order prime to the
residue characteristics of~$S$, such that there is a
\ifthenelse{\boolean{orig}}
{morphism}
{monomorphism}
% correction: monomorphism
$$
j\colon F \to \prod_{i} p_{i*}C_i=G \quoi.
$$
Let $\Phi_i$ be the stack of torsors under $C_i$, and let $\Psi$ be the stack of
torsors under $G$. It follows from 2)\kern2pt a) that, after
restricting $S$ to a nonempty open subset, we may suppose that the
$(C_i,fp_i)$ are cohomologically proper relative to~$S$ in
dimension $\le 1$, and that the stacks $f_*p_{i*}\Phi_i$ are
$1$-constructible. It then follows from~\Ref{XIII.1.9} that the
$(p_{i*}C_i,f)$ are cohomologically proper relative to~$S$
in dimension $\le 1$; hence the same is true of~$(G,f)$. Moreover,
since $p_{i*}\Phi_i$ is equivalent to the stack of torsors under the
group $p_{i*}C_i$ (SGA~4 VIII~5.8), we see that $f_*\Psi$ is
$1$-constructible.

Since $\R^1f_*G$ is constructible, one can find a sheaf
representable by an \'etale $Y$-scheme of finite type $T$ and
an epimorphism
$$
a \colon T \to \R^1f_*G
$$
(SGA~4 IX~2.7); moreover, we may suppose that the image of the
identity section of~$T(T)$ is defined by a torsor $Q$ on~$X\times_Y
T=X_T$, with group $G|X_T$. Let $f_T \colon X_T \to T$ be the
canonical morphism, and put $F_T=F|X_T$, etc. By 1)\kern2pt b), after
restricting $S$ to a nonempty open subset, we may suppose that
$(Q/F_T,f_T)$ is cohomologically proper relative to~$S$ in
dimension $\le 0$ and that $f_{T*}(Q/F_T)$ is constructible. It
then follows from~\Ref{XIII.3.1.2} that $f_*\Phi$ is $1$-constructible.

Let us show that $(F,f)$ is cohomologically proper relative to
% page: 404
\marginpar{404}
%
$S$ in dimension $\le 1$. By \Ref{XIII.1.13}~2), it is enough to
prove that, for every scheme $Y_1$ \'etale over~$Y$ and every
torsor $Q_1$ under~$G_1$ on~$X_1=X\times_Y Y_1$, if $f_1 \colon X_1 \to Y_1$ is
the canonical morphism, then $(Q_1/F_1,f_1)$ is cohomologically
proper relative to~$S$ in dimension $\le 0$. But, by
definition of~$T$, $Q_1$ is, locally for the \'etale topology
on~$Y_1$, the inverse image of~$Q$, which proves our
reduction.

4) 2. The general case.\par\nopagebreak
Using Lemma~\Ref{XIII.6.1.1}, 4)\kern2pt1, and 1)\kern2pt b), we see that,
after restricting $S$ to a nonempty open subset, we may
suppose that $S(f_*\Phi)$ is constructible and $(S\Phi,f)$ is cohomologically
proper relative to~$S$ in dimension $\le 0$. We can then
find a sheaf representable by an \'etale
$Y$-scheme of finite type $T$ and an epimorphism
$$
a \colon T \to S(f_*\Phi)\quoi.
$$
We again use the notation of 4)\kern2pt1, and in addition put $Z=T\times_Y T$,
$X_Z=X\times_Y Z$, and denote by $f_Z \colon X_Z \to Z$ the
canonical morphism. We may suppose that $T$ is chosen so that the image $q$ of the
identity section of~$T(T)$ under $a$ is defined by an object~$p$
of~$(f_*\Phi)_T=\Phi_{X_T}$. Let $p_1$ and $p_2$ (\resp~$q_1$ and
$q_2$) be the inverse images of~$p$ (\resp~$q$) under the two projections
from~$Z$ to $T$. After restricting $S$ to
a nonempty open subset, we may suppose that $(\SheafAut_{X_T}(p),f_T)$ is
cohomologically proper relative to~$S$ in dimension $\le 1$,
that $(\SheafHom_{X_Z}(p_1,p_2),f_Z)$ is cohomologically proper
relative to~$S$ in dimension $\le 0$, and that
$f_{T*}(\SheafAut_{X_T}(p))$, $f_{Z*}(\SheafHom_{X_Z}(p_1,p_2))$ are
constructible (1)\kern2pt b) and 4)\kern2pt1)).
\end{filecontents*}
% Original input: drafts/SGA1_XIII_3_CASE3_B_THROUGH_3_1_3_English.texfrag
\begin{filecontents*}[overwrite,noheader]{sga1-cumulative-input-110.tex}
It follows first that, for every scheme \(Y_1\) étale over~\(Y\)
and every pair of objects \(y,y_1\) of~\((f_*\Phi)_{Y_1}\), the
sheaf \(\SheafAut_{Y_1}(y)\)
(\resp \(\SheafHom_{Y_1}(y,y_1)\)) is constructible: locally for the
étale topology of~\(Y_1\), such a sheaf is the inverse image of
\(f_{T*}(\SheafAut_{X_T}(p))\)
(\resp of~\(f_{Z*}(\SheafHom_{X_Z}(p_1,p_2))\)).

It remains to prove that \((\Phi,f)\) is cohomologically proper
relative
% page: 405
\marginpar{405}
%
to~\(S\) in dimension \(\leq1\). For this, it is enough to show
that, for every \(S\)-scheme \(S'\) and every geometric point
\(\overline y'\) of~\(Y'\), if \(\overline y\) denotes its image in~\(Y\),
\(\overline Y\) and \(\overline Y'\) denote the strict localizations
of~\(Y\) at~\(\overline y\) and of~\(Y'\) at~\(\overline y'\), respectively, and
\(\overline X=X\times_Y\overline Y\), etc., then the canonical
functor
$$
\overline{\varphi} \colon \Phi(\overline{X}) \to
\Phi'(\overline{X'})
$$
is an equivalence of categories.

Let us show that \(\overline\varphi\) is fully faithful. Let
\(x,y\in\Phi(\overline X)\), and let \(x',y'\) be their images in
\(\Phi'(\overline{X'})\). We must show that the canonical morphism
\begin{equation*}
\label{eq:XIII.4.*} \tag{$*$} {\Hom_{\overline{X}}(x,y)
\to \Hom_{\overline{X'}}(x',y') }
\end{equation*}
is bijective. By the definition of~\(T\), there are two morphisms
from~\(\overline Y\) to \(T\) such that \(x\) and \(y\) are the
inverse images of~\(p\) under these two morphisms. This amounts to
giving a morphism \(\overline Y\to Z\), and hence a morphism
\(h\colon\overline X\to X_Z\) such that
$$
h^*(p_1)=x \qquad h^*(p_2)=y.
$$
Consequently, there is a canonical isomorphism
$$
\SheafHom_{\overline{X}}(x,y)=
h^*(\SheafHom_{X_Z}(p_1,p_2))\quoi.
$$
But, since \((\SheafHom_{X_Z}(p_1,p_2),f_Z)\) is cohomologically
proper relative to~\(S\) in dimension \(\leq0\), the same is true of
\((\SheafHom_{\overline X}(x,y),\overline f)\), which proves that
the morphism \eqref{eq:XIII.4.*} is bijective.

Let us show that \(\overline\varphi\) is essentially surjective. Let
\(x'\in\Phi'(\overline{X'})\). Since \((S\Phi,f)\) is
cohomologically proper relative to~\(S\) in dimension \(\leq0\),
the canonical morphism
$$
\H^0(\overline{X},\overline{S\Phi}) \to
\H^0(\overline{X'},\overline{S\Phi}')
$$
is bijective. Let \(G'\) be the maximal subgerbe of~\(\Phi'\)
generated by~\(x'\). There is then a maximal subgerbe \(G\)
of~\(\overline\Phi\) whose inverse image on~\(\overline{X'}\)
is~\(G'\). By 4)\kern2pt1, \((G,\overline f)\) is cohomologically
proper relative to~\(S\) in dimension \(\leq1\). Consequently, the
canonical functor
$$
G(\overline{X}) \to G'(\overline{X'})
$$
% page: 406
\marginpar{406}%
%
is an equivalence of categories. This proves the existence of an
object \(x\) of~\(\Phi(\overline X)\) whose image in
\(\Phi'(\overline{X'})\) is isomorphic to~\(x'\), and completes the proof of
the theorem.

\begin{sublemme}
\label{XIII.3.1.1}
Let \(S\) be a locally Noetherian scheme, and let
$$
\xymatrix@C=.5cm{ \Phi \ar[r]^p & \Phi_1 \ar@<2pt>[rr]^{p_1,p_2}
\ar@<-2pt>[rr] && \Phi_2 }
$$
be an exact diagram of stacks on~\(S\) \eqref{XIII.1.10.1}. If
\(\Phi_1\) is constructible, then so is~\(\Phi\). If, for every
scheme \(S'\) étale over~\(S\) and every pair
\ifthenelse{\boolean{orig}}
{of object}
{of objects}
% correction: objects instead of object
\(x_1,y_1\) of~\((\Phi_1)_{S'}\), the sheaf
\(\SheafHom_{S'}(x_1,y_1)\) is constructible, then for every pair of
objects \(x,y\) of~\(\Phi_{S'}\), the same is true of
\(\SheafHom_{S'}(x,y)\). Suppose that \(\Phi_1\) is
\(1\)-constructible (\Ref{XIII.0}) and that \(\Phi_2\) is
constructible; then \(\Phi\) is \(1\)-constructible.
\end{sublemme}

For every scheme \(S'\) étale over~\(S\) and every object \(x\)
of~\(\Phi_{S'}\), there is a monomorphism
$$
\SheafAut_{S'}(x)\to\SheafAut_{S'}(p(x)).
$$
It follows from SGA~4 IX~2.9 that if \(\Phi_1\) is constructible,
then so is~\(\Phi\), and the second assertion of the lemma is proved
in the same way.

Now suppose that \(\Phi_1\) is \(1\)-constructible and that \(\Phi_2\)
is constructible. The morphism \(p\) induces a morphism on the
sheaves of maximal subgerbes,
$$
\varphi\colon S\Phi\to S\Phi_1.
$$
Let \(G\) be the image of~\(S\Phi\) under \(\varphi\). By
SGA~4 IX~2.9, \(G\) is a constructible sheaf. We can therefore find
a sheaf represented by an étale \(S\)-scheme \(T\) of finite type,
and an epimorphism
$$
a\colon T\to G
$$
(SGA~4 IX~2.7).
% page: 407
\marginpar{407}
%
Moreover, \(T\) can be chosen so that the image \(y\) of the identity
section of~\(T(T)\) under \(a\) is defined by an object \(x_1\)
of~\((\Phi_1)_T\) of the form \(x_1=p(x)\), where
\(x\in\Phi_T\).

It is enough to show that, for every point \(s\) of~\(S\), there is
a nonempty open subset \(U\) of~\(\overline{\{s\}}\) such that
\(S\Phi|U\) is constructible locally constant. Let \(s\in S\), let
\(\overline s\) be a geometric point over~\(s\), and let
\(\overline y_1,\dots,\overline y_n\) be the elements
of~\(G_{\overline s}\). By the definition of~\(T\), there are
morphisms \(h_i\colon\overline s\to T\) such that
\(h_i^*(y)=\overline y_i\). Let \(S'\) be the fiber product over~\(S\)
of \(n\) schemes isomorphic to \(T\), let
\(h\colon\overline s\to S'\) be the fiber product of the \(h_i\),
and let \(y_i\) (\resp \(x_i\)) be the inverse image of \(y\)
(\resp \(x\)) under the \(i\)-th projection from~\(S'\) to~\(T\).
Let \(F_i\) be the subsheaf of~\(S\Phi|S'\) that is the inverse
image of~\(y_i\), and let us show that the \(F_i\) are
constructible. The sheaf \(F_i\) is a quotient of the sheaf \(F'_i\)
such that, for every scheme \(S''\) étale over~\(S'\),
\begin{multline*}
F'_i(S'') = \{\text{classes modulo isomorphism of objects \(z\) of
\(\Phi_{S''}\) equipped}\\
\ifthenelse{\boolean{orig}}
{\text{ with an isomorphism \(i\colon p(z)\simeq p(x_i|S'')\)}.}
{\text{ with an isomorphism \(i\colon p(z)\simeq p(x_i|S'')\)}\}.}
% correction: missing closing brace
\end{multline*}
It is enough to show that the \(F'_i\) are constructible. Put
\(z_i=p_1p(x_i)\), \(z'_i=p_2p(x_i)\). There is a monomorphism
$$
\Psi_i\colon F'_i\to\SheafIsom_{S'}(z_i,z_i'),
$$
obtained by associating with every scheme \(S''\) étale over~\(S'\)
and every object \(z\) of~\(\Phi_{S''}\) equipped with an isomorphism
\(i\colon p(z)\simeq p(x_i|S'')\), the isomorphism from~\(z_i\)
to~\(z'_i\) defined by requiring the diagram
$$
\xymatrix@C=1.2cm{
p_1p(z) \ar[r]^-{p_1(i)}\ar[d]_-j& z_i\ar[d]\\ p_2p(z) \ar[r]^-{p_2(i)}& z'_i
}
$$
to be commutative (\(j\) is the canonical morphism associated with
the exact diagram).
% page: 408
\marginpar{408}
%

The morphism \(\Psi_i\) is injective: to say that two objects
\(z,z'\), equipped with isomorphisms
\(i\colon p(z)\isomto p(x_i|S'')\) and
\(i'\colon p(z')\isomto p(x_i|S'')\), define the same element of
\(\SheafIsom_{S''}(z_i,z'_i)\) amounts to saying that
\(p_1(i'^{-1}i)=p_2(i'^{-1}i)\), \ie that \(i'^{-1}i\) comes from
an isomorphism \(z\to z'\). Since \(F'_i\) is a subsheaf of
\(\SheafIsom_{S'}(z_i,z'_i)\), it is constructible.

Since there is a nonempty open subset \(U\) of~\(\overline{\{s\}}\)
such that \(y_1|U,\dots,y_n|U\) generate \(G\), it follows from
Lemma~\Ref{XIII.6.1.2} below that \(S\Phi|U\) is constructible.
Thus, after restricting \(U\) if necessary, \(S\Phi|U\) is
constructible locally constant.

\begin{sublemme}
\label{XIII.3.1.2}
Let \(S\) be a locally Noetherian scheme, let \(f\colon X\to S\) be
a morphism, let \(F\to G\) be a monomorphism of sheaves of groups
on~\(X\), let \(\Psi\) (\resp \(\Psi_1\)) be the stack of torsors
under \(F\) (\resp under \(G\)), and put
\(\Phi=f_*\Psi\), \(\Phi_1=f_*\Psi_1\). Suppose that we are given a
sheaf on~\(S\) represented by an étale \(S\)-scheme \(T\) of finite
type and a surjective morphism
$$
a\colon T\to S\Phi_1\simeq\R^1f_*G,
$$
such that there is a torsor \(Q\) on
\(X_T=X\times_ST\) with group \(G|X_T\), defining in
\(\R^1f_*G(T)\) the image under \(a\) of the identity section
of~\(T(T)\). Let \(f_T\colon X_T\to T\) be the canonical morphism,
and put \(F_T=F|X_T\). Suppose that \(\Phi_1\) is
\(1\)-constructible and that \(f_{T*}(Q/F_T)\) is constructible.
Then \(\Phi\) is \(1\)-constructible.
\end{sublemme}

Indeed, it is enough to repeat the proof of~\Ref{XIII.3.1.1}, with
the existence of the morphisms \(\Psi_i\) replaced by the
isomorphisms
$$
F'_i\isomto f_{T*}(Q/F_T)|S'.
$$

\begin{subremarque}
\label{XIII.3.1.3}
Suppose
% page: 409
\marginpar{409}
%
that \(k=\kres(s)\) has characteristic zero. Then schemes of finite
type over~\(\overline k\), of dimension \(\leq\dim X\), admit strong
resolution of singularities, and the proof of~\Ref{XIII.3.1} gives
the following results:
\begin{enumerate}
\item[a)] There is a nonempty open subset \(S_1\) of~\(S\) such
that, for every scheme \(S'\) over~\(S_1\) whose maximal points have
characteristic zero and every constructible locally constant sheaf
of sets \(F\) on~\(X'=X\times_SS'\), \((F,f')\) is cohomologically
proper relative to~\(S'\) in dimension \(\leq0\).

\item[b)] If all the residue characteristics of~\(S\) are zero,
there is a nonempty open subset \(S_1\) of~\(S\) such that, for
every scheme \(S'\) over~\(S_1\) and every constructible locally
constant sheaf of groups \(F\) on~\(X'\), \((F,f')\) is
cohomologically proper relative to~\(S'\) in dimension
\(\leq1\).
\end{enumerate}

Indeed, it is enough to repeat the proof of
\Ref{XIII.3.1}~2)\kern2pt a). Proposition~\Ref{XIII.2.7}, used in
3)\kern2pt4, applies to a locally constant sheaf \(F\), because,
with the notation of 3)\kern2pt4, every torsor under \(F\) is tamely
ramified on~\(P\) relative to~\(S\), since all the residue
characteristics of~\(S\) are zero.
\end{subremarque}

\end{filecontents*}
% Original input: drafts/SGA1_XIII_3_COROLLARIES_English.texfrag
\begin{filecontents*}[overwrite,noheader]{sga1-cumulative-input-111.tex}
\begin{corollaire}
\label{XIII.3.2}
Let \(k\) be a field of characteristic \(p\geq0\), let \(p'\) be
the set of prime numbers distinct from~\(p\), and let
\(f\colon X\to k\) be a coherent morphism.
\begin{enumerate}
\item[1)] For every sheaf of sets \(F\), \((F,f)\) is
cohomologically proper in dimension \(\leq0\).
\item[2)] Suppose that one of the following two conditions holds:
\begin{enumerate}
\item[a)] \(f\) is of finite type, and schemes of finite type of
dimension \(\leq\dim X\) over an algebraic closure of~\(k\) admit
strong resolution of singularities.
\item[b)] Schemes of finite type over an algebraic closure of~\(k\)
% page: 410
\marginpar{410}
%
admit strong resolution of singularities.
\end{enumerate}
Then, for every sheaf of ind-\(p'\)-groups \(F\), \((F,f)\) is
cohomologically proper in dimension \(\leq1\).
\end{enumerate}
\end{corollaire}

Let \(F\) be a sheaf of sets (\resp of ind-\(p'\)-groups).
By SGA~4 IX~2.7.2, \(F\) can be written as a filtered inductive limit
$$
F=\varinjlim F_i,
$$
where the \(F_i\) are constructible sheaves of sets
(\resp of ind-\(p'\)-groups). Since \(f\) is coherent, \(f_*\)
(\resp \(\R^1f_*\)) commutes with inductive limits
(SGA~4 VII~3.3). If the \((F_i,f)\) are known to be
cohomologically proper in dimension \(\leq0\) (\resp if every sheaf
of sets is cohomologically proper in dimension \(\leq0\) and the
\((F_i,f)\) are cohomologically proper in dimension \(\leq1\)),
then the same is true of~\((F,f)\). We may therefore suppose that
\(F\) is constructible.

If \(f\) is assumed to be of finite type (\resp to satisfy a)), the
proposition follows from \Ref{XIII.3.1}~1)\kern2pt b)
(\resp from \Ref{XIII.3.1}~2)\kern2pt b)). We now prove
\Ref{XIII.3.2} without assuming that \(f\) is of finite type. For
every scheme \(S'\) over~\(k\) and every geometric point
\(\overline s\) of~\(S'\), denote by \(\overline k\)
(\resp \(\overline{S'}\)) the strict localization of~\(k\)
at~\(\overline s\) (\resp of~\(S'\) at~\(\overline s\)), let
\(\overline X\) be the inverse image of~\(X\) on~\(\overline k\),
and consider the cartesian square
$$
\xymatrix@C=1.2cm{
\overline{X}\ar[d]_-{\overline{f}}&\ar[l]_g \overline{X'}\ar[d]^-{\overline{f}'}\\ \overline{k}&\ar[l]\overline{S'}
}
$$
It is enough to prove that, for every \(S'\) and every
\(\overline s\), the canonical morphism
$$
\H^0(\overline{X},\overline{F})\to
\H^0(\overline{X'},\overline{F'})\qquad
\text{(\resp \(\H^1(\overline{X},\overline{F})\to
\H^1(\overline{X'},\overline{F'})\))}
$$
is an isomorphism. It is enough to show the relations
% page: 411
\marginpar{411}
%
\begin{equation*}
\label{eq:XIII.3.*}
\tag{$*$} \overline{F}\simeq g_*g^*\overline{F}\qquad
\ifthenelse{\boolean{orig}}
{\text{(\resp $\R^1g_*(g^*\overline{F}))=0$).}}
{\text{(\resp $\R^1g_*(g^*\overline{F})=0$).}}
% correction: one closing parenthesis too many
\end{equation*}
In this form, the question is local on~\(X\) for the étale topology.
We may therefore suppose that \(X\) is affine; by passage to the
limit, we may suppose that \(X\) is of finite type over~\(k\). It is
then known that \((F,f)\) is cohomologically proper in dimension
\(\leq0\) (\resp \(\leq1\)), and that the same is true when \(X\) is
replaced by a scheme étale and of finite type over~\(X\). This proves
\eqref{eq:XIII.3.*}.

\begin{theoreme}
\label{XIII.3.3}
Let \(S\) be an irreducible scheme with generic point \(s\), and let
\(f\colon X\to S\) be a morphism of finite presentation. Suppose
that schemes of finite type of dimension \(\leqslant\dim X_s\) over
an algebraic closure \(\bar k\) of~\(\kres(s)\) admit resolution of
singularities \textup{(EGA~IV 7.9.1)}. If \(\LL\) denotes the set
of prime numbers distinct from the residue characteristics of~\(S\),
then there is a nonempty open subset \(S_1\) of~\(S\) such that the
morphism \(f|S_1\) is universally locally \(1\)-aspherical
for~\(\LL\).
\end{theoreme}

We may suppose that \(S\) is integral and \(X\) is reduced
(SGA~4 VIII~1.1). By passage to the limit, we may suppose that \(S\)
is Noetherian. Moreover, to prove the theorem, it is enough to prove
it after a finite extension \(S_1\to S\), where \(S_1\) is an
integral scheme and \(S_1\to S\) is a composite of étale extensions
and surjective radicial extensions.

First let us show that, after restricting \(S\) to a nonempty open
subset if necessary, \(f\) is universally locally \(0\)-acyclic.
After a radicial extension of~\(\kres(s)\), we may suppose that
\((X_s)_\red\to s\) is separable. Thus, after restricting \(S\) to a
nonempty open subset and making a surjective radicial extension
of~\(S\), we may suppose that \(f\) is flat with geometrically
separable fibers (EGA~IV 12.1.1), which implies that \(f\) is
universally locally \(0\)-acyclic (SGA~4 XV~4.1).

Let us show that, after restricting \(S\) to a nonempty open subset
if necessary, \(f\) is universally locally \(1\)-aspherical
for~\(\LL\). After making a
% page: 412
\marginpar{412}
%
finite extension of~\(\kres(s)\), which is allowed because it may be
regarded as a composite of an étale extension and a radicial
extension, we can find a proper surjective morphism
\(p_s\colon Y_s\to X_s\), where \(Y_s\) is a scheme smooth over~\(s\)
and has the same dimension as \(X_s\). We may suppose that \(p_s\)
comes from a proper surjective morphism \(p\colon Y\to X\), where
\(Y\) is smooth over~\(S\) (EGA~IV 9.6.1 and 12.1.6). It is enough
to show that, after restricting \(S\) to a nonempty open subset if
necessary, for every diagram with cartesian squares
$$
\xymatrix{ S''\ar[d]_i & X ''\ar[d]^j\ar[l]_{f''} \\ S'\ar[d] &
X'\ar[d]\ar[l]_{f'} \\ S & X, \ar[l]_f }
$$
where \(i\) is étale of finite presentation, and every sheaf of
ind-\(\LL\)-groups \(F\) on~\(S''\), if \(\Phi\) is the stack of
torsors under \(F\), then the canonical morphism
$$
{f'}^*i_*\Phi \to j_*{f''}^*\Phi
$$
is an equivalence. Put \(Z=Y\times_XY\) and
\(T=Y\times_XY\times_XY\). There is a natural commutative diagram
$$
\xymatrix{ S'' \ar[d]_i & X''\ar[l]_{f''}\ar[d]_j &
Y''\ar[l]_{p''}\ar[d] & Z''\ar[d]\ar@<-.4ex>[l]\ar@<.4ex>[l] &
T''\ar[d]\ar[l]\ar@<-.8ex>[l]\ar@<.8ex>[l] \\ S' \ar[d] & X'
\ar[l]_{f'}\ar[d] & Y'\ar[d]\ar[l]_{p'}&
Z'\ar[d]\ar@<-.4ex>[l]\ar@<.4ex>[l] &
T'\ar[d]\ar[l]\ar@<-.8ex>[l]\ar@<.8ex>[l] \\ S & X\ar[l]_f & Y\ar[l]_p
& Z\ifthenelse{\boolean{orig}}{'}{}\ar@<-.4ex>[l]\ar@<.4ex>[l] & T\ifthenelse{\boolean{orig}}{'}{}\ar[l]\ar@<-.8ex>[l]\ar@<.8ex>[l]
}
$$
% remark: the spacing of the double and triple arrows could be
% harmonized with that in the preceding sections.
Let
% page: 413
\marginpar{413}%
%
\(q\colon Z\to X\) and \(r\colon T\to X\) be the canonical
morphisms, with \(q',r',q'',r''\) having the evident meanings.
By \Ref{XIII.1.11}~2), there is the following essentially
commutative diagram, whose rows are exact:
$$
\ifthenelse{\boolean{orig}}
{\xymatrix{ f'^* i_*\Phi \ar[r]\ar[d]_a &
p'_*p'^*(f'^*i_x\Phi)\ar[d]_b\ar@<-.4ex>[r]\ar@<.4ex>[r] &
q'_*q'^*(f'^*i_*\Phi)\ar[d]_c\ar[r]\ar@<-.8ex>[r]\ar@<.8ex>[r]&
r'_*r'^*(f'^*i_*\Phi)\ar[d]_d \\ j_*f''^*\Phi \ar[r]&
j_*p''_*p''^*(f''^*\Phi)\ar@<-.4ex>[r]\ar@<.4ex>[r] &
j_*q''_*q''^*(f''^*\Phi)\ar[r]\ar@<-.8ex>[r]\ar@<.8ex>[r]&
j_*r''_*r''^*(f''^*\Phi) }}
{\xymatrix{ f'^* i_*\Phi \ar[r]\ar[d]_a &
p'_*p'^*(f'^*i_*\Phi)\ar[d]_b\ar@<-.4ex>[r]\ar@<.4ex>[r] &
q'_*q'^*(f'^*i_*\Phi)\ar[d]_c\ar[r]\ar@<-.8ex>[r]\ar@<.8ex>[r]&
r'_*r'^*(f'^*i_*\Phi)\ar[d]_d \\ j_*f''^*\Phi \ar[r]&
j_*p''_*p''^*(f''^*\Phi)\ar@<-.4ex>[r]\ar@<.4ex>[r] &
j_*q''_*q''^*(f''^*\Phi)\ar[r]\ar@<-.8ex>[r]\ar@<.8ex>[r]&
j_*r''_*r''^*(f''^*\Phi) }}
% correction: in row 1, column 2, replace i_x by i_*.
$$
Since \(Y\) is smooth over~\(S\), the morphism \(fp\) is universally
locally \(1\)-aspherical for \(\LL\) (SGA~4 XV~2.1),
% source reference: SGA 4 XV 2.1
and it follows from \cite[VII 2.1.7]{XIII.2} that \(b\) is an
equivalence of categories. On the other hand, after restricting
\(S\) to a nonempty open subset if necessary, we may suppose that
the morphisms \(Z\to S\) and \(T\to S\) are universally locally
\(0\)-acyclic. It follows that the functors \(c\) and \(d\) are
fully faithful, and the diagram above then shows that \(a\) is an
equivalence, completing the proof.

\begin{corollaire}
\label{XIII.3.4}
Let \(k\) be a field of characteristic \(p\geqslant0\), let \(p'\)
be the set of prime numbers distinct from~\(p\), and let
\(f\colon X\to k\) be a coherent morphism. Suppose that one of the
following two conditions holds:
\begin{enumerate}
\item[a)] \(f\) is of finite type, and schemes of finite type of
dimension \(\leqslant\dim X\) over an algebraic closure of~\(k\)
admit resolution of singularities.
\item[b)] Schemes of finite type over an algebraic closure of~\(k\)
admit resolution of singularities.
\end{enumerate}
Then \(f\) is universally locally \(1\)-aspherical for \(p'\).
\end{corollaire}

Case a) follows from~\Ref{XIII.3.3}. In case b), since the question
is local on~\(X\), we may suppose that \(X\) is affine; passage to
the limit (SGA~4 XV~1.3) reduces us to the case where \(X\) is of
finite type over~\(k\).

\begin{corollaire}
\label{XIII.3.5}
Let
% page: 414
\marginpar{414}
%
\(S\) be an irreducible scheme with generic point \(s\), and let
\(f\colon X\to S\) be a morphism of finite presentation. Suppose
that schemes of finite type of dimension
\(\leqslant\dim X_s\) over an algebraic closure \(\overline k\)
of~\(\kres(s)\) admit strong resolution of singularities
\textup{(SGA~5 I~3.1.5)}. If \(\LL\) denotes the set of prime numbers
distinct from the residue characteristics of~\(S\), then there is a
nonempty open subset \(S_1\) of~\(S\) such that, for every
specialization \(\bar s_1\to\bar s_2\) of geometric points
of~\(S_1\), the specialization morphism \eqref{XIII.2.10}
$$
\pi_1^{\LL}(X_{\bar{s}_1})\to\pi_1^{\LL}(X_{\bar{s}_2})
$$
is bijective.
\end{corollaire}

By~\Ref{XIII.3.1} and~\Ref{XIII.3.3}, after restricting \(S\) to a
nonempty open subset if necessary, we may suppose that \(f\) is
locally \(1\)-aspherical for \(\LL\), and that, for every finite
constant sheaf of \(\LL\)-groups \(F\) on~\(X\), \((F,f)\) is
cohomologically proper in dimension \(\leqslant1\). It then follows
from~\Ref{XIII.1.14} that, for every specialization
\(\bar s_1\to\bar s_2\) of geometric points of~\(S_1\), the
specialization morphism
$$
\left(\R^1 f_* F\right)_{\bar{s}_2} \to \left(\R^1 f_*
F\right)_{\bar{s}_1}
$$
is bijective. The corollary is precisely the translation of the
preceding statement into the language of fundamental groups.

\end{filecontents*}
% Original input: drafts/SGA1_XIII_3_OPENING_AND_CASE1_English.texfrag
\begin{filecontents*}[overwrite,noheader]{sga1-cumulative-input-112.tex}
\section{Cohomological properness and generic local acyclicity}
\label{XIII.3}
% page: 395
\marginpar{395}
%

\begin{theoreme}
\label{XIII.3.1}
Let $S$ be an irreducible scheme with generic
point~$s$, let $X$ and $Y$ be two $S$-schemes of
finite presentation, and let $f\colon X \to Y$ be an $S$-morphism. For every
$S$-scheme~$S'$, denote by $Y'$, $X'$, etc. the inverse images of~$Y$,
$X$, etc. under the morphism $S' \to S$. The following properties
hold:

\begin{enumerate}
\item[1)]
{\upshape a)} One can find a nonempty open subset~$S'$ of~$S$ such that, for
every finite constant sheaf of sets $F'$ on~$X'$, $f'_*F'$ is
constructible, and $(F',f')$ is cohomologically proper
relative to~$S'$ in dimension $\le 0$.

{\upshape b)} Let $F$ be a constructible sheaf of sets on~$X$. Then
one can find a nonempty open subset~$S'$ of~$S$ (depending on~$F$)
such that $f'_*F'$ is constructible and $(F',f')$ is
cohomologically proper relative to~$S'$ in dimension $\le 0$.

\item[2)] Suppose that the schemes of finite type of dimension $\le
\dim X_s$ over an algebraic closure $\overline{k}$ of~$\kres(s)$
are strongly desingularizable \textup{(SGA~5 I~3.1.5)}. Then the
following additional properties hold:

{\upshape a)} One can find a nonempty open subset~$S'$ of~$S$ such that, for
every finite constant sheaf of groups $F'$ on~$X'$ whose order is prime
to the residual characteristics of~$S$, if $\Phi'$ is the
stack of torsors under $F'$, then $f'_*\Phi'$ is $1$-constructible, and
$(F',f')$ is cohomologically proper relative to~$S'$ in
dimension $\le 1$.

{\upshape b)} Let $\LL$ be the set of prime numbers distinct from the
residual characteristics of~$S$, and let $\Phi$ be an
ind-$\LL$-stack that is $1$-constructible on~$X$ \eqref{XIII.0} such that,
for every scheme~$X_1$ \'etale over~$X$ and every pair
of objects $x, x_1$ of~$\Phi_{X_1}$, the sheaf $\SheafHom
_{X_1}(x,x_1)$ is constructible. Suppose moreover that $S$ is locally
noetherian. Then one can find a nonempty open subset~$S'$ of~$S$
such that $f'_*\Phi'$ is $1$-constructible, such that, for every scheme $Y_1$ \'etale over~$Y'$ and every pair
of objects $y, y_1$ of~$(f'_*\Phi')_{Y_1}$,
$\SheafHom_{Y_1}(y,y_1)$ is constructible,
% page: 396
\marginpar{396}
%
and such that $(\Phi',f')$ is cohomologically proper relative to~$S'$ in dimension $\le 1$.
\end{enumerate}
\end{theoreme}

\subsubsection*{Proof}
We may suppose that $S$ is affine; by SGA~4 VIII~1.1, we may
suppose that $S$ is integral; finally, by passage to the limit, we may
suppose that $S$ is the spectrum of an algebra of finite type over~$\ZZ$; in particular, $S$ is then noetherian. Since the question
is local on~$Y$, we may suppose that $Y$ is affine. Moreover, in order
to prove the theorem, it suffices to prove it after a
finite extension $S'\to S$, where $S'$ is an integral scheme
and where $S'\to S$ is a composite of \'etale morphisms and
finite surjective radicial morphisms.

\Subsection*{1) Case of constant sheaves of sets}

1) 1. Reduction to the case where $X$ is normal over~$S$. Let
$X_{1\overline{s}}$ be the normalization of~$(X_{\overline{s}})_{\red}$;
after restricting $S$ to a nonempty open subset and making
a radicial extension of~$S$, we may suppose that
$X_{1\overline{s}}$ comes from a scheme $X_1$ normal over~$S$, and
that the morphism $X_{1\overline{s}}\to X_{\overline{s}}$ comes from a
finite surjective morphism $p\colon X_1\to X$ (EGA~IV 8.8.2 and
9.6.1). Suppose that the theorem has been proved for
$fp$. After restricting $S$ to an open subset, we may suppose
that, for every constant sheaf of sets $F$ on~$X$, $(p^*F,fp)$
is cohomologically proper relative to~$S$ in dimension $\le
0$ and that $f_*p_*(p^*F)$ is
constructible. By~\Ref{XIII.1.9}, $(p_*p^*F,f)$ is then
cohomologically proper relative to~$S$ in dimension $\le
0$. The morphism
$$
F \to p_* p^* F=G
$$
is a monomorphism. Since $f_*F$ is
a subsheaf of~$f_*G$, it already follows that $f_*F$ is constructible (SGA~4
IX~2.9~(ii)) and that $(F,f)$ is cohomologically proper relative
to~$S$ in dimension $\le -1$.

Let $X_2=X_1 \times_{X} X_1$, and let $q \colon X_2 \to X$ be the
canonical morphism. By \Ref{XIII.1.11}~1),
% page: 397
\marginpar{397}
%
there is an exact sequence
$$
\xymatrix@C=.5cm{F \ar[r] & G \ar@<2pt>[r]\ar@<-2pt>[r] & q_*q^*F }.
$$

By what we have just proved, applied to
$fq$ in place of~$f$, after restricting $S$
to a nonempty open subset, we may suppose that, for every constant
sheaf of sets $F$ on~$X$, $(q^*F,fq)$ is cohomologically proper relative
to~$S$ in dimension $\le -1$, and hence that $(q_*q^*F,f)$ is
cohomologically proper relative to~$S$ in dimension $\le
-1$. It then follows from \Ref{XIII.1.13}~1) that $(F,f)$ is
cohomologically proper relative to~$S$ in dimension $\le 0$.

1) 2. Reduction to the case where $X$ is normal affine over~$S$.

Let $U_s$ be an affine open subset of~$X_s$ dense in $X_s$. After
restricting $S$ to a nonempty open subset, we may suppose that $U_s
\to X_s$ lifts to an open immersion $i\colon U\to X$,
schematically dominant relative to~$S$ (EGA~IV
8.9.1). Since the morphism $X \to S$ is normal, by
SGA~2 XIV~1.18 we have
$$
\prof \et_{X- U} (X) \ge 2 \quoi ;
$$
consequently, for every constant sheaf $F$ on~$X$, the canonical
morphism
$$
F \to i_*i^*F
$$
is an isomorphism. It follows that, if the
theorem is assumed proved for $fi$ and $i$, then, after
restricting~$S$ to a nonempty open subset, $(i^*F,i)$ and $(i^*F,fi)$
are cohomologically proper relative to~$S$ in dimension
$\le 0$. The same is therefore true of~$(F,f)$ (\Ref{XIII.1.6}~2)). Moreover, $f_* F=(fi)_*(i^*F)$ is constructible, which
completes the reduction.

1) 3. End of the proof.

We may suppose that $S$ is normal (EGA~IV 7.8.3). One can find a
compactification of~$X_s$:
% page: 398
\marginpar{398}
%
$$
\xymatrix{ X_s \ar[r]^{j_s} \ar[d]_{f_s} & P_s \ar[ld]^{g_s} \\ Y_s }
$$
where $j_s$ is a dominant open immersion and $g_s$ is a proper morphism;
after making a radicial extension of~$\kres(s)$ and
replacing $P_s$ by its normalization, which does not change $X_s$, we
may suppose that $P_s$ is geometrically normal. After
restricting $S$ to a nonempty open subset and making a surjective
radicial extension, we may suppose that the diagram above
comes from a diagram
$$
\xymatrix{ X \ar[r]^{j} \ar[d]_{f} & P \ar[ld]^{g} \\ Y }
$$
where $P$ is a scheme normal over~$S$, $j$ is a schematically
dominant open immersion relative to~$S$, and $g$ is a
proper morphism (EGA~IV 6.9.1, 9.9.4 and 9.6.1). For every finite
constant sheaf of sets $F$ on~$X$ with value $C$, $j_*F$ is the
constant sheaf with value $C$ (SGA~4 IX~2.14.1), and the same is true
after every base change $S' \to S$; it
follows that $(F,j)$ is cohomologically proper relative
to~$S$ in dimension $\le 0$. The same is therefore true of~$(F,f)$,
since $g$ is proper \eqref{XIII.1.8}. Since $g$ is proper,
$f_*F=g_*C_P$ is constructible, which completes the
proof of 1)\kern2pt a).
\end{filecontents*}
% Original input: drafts/SGA1_XIII_4_1_English.texfrag
\begin{filecontents*}[overwrite,noheader]{sga1-cumulative-input-113.tex}
\ifthenelse{\boolean{orig}}{\subsection{}}{\sgaFixedHeading{\subsection}{XIII.4.5}}
\label{XIII.4.5}
We retain the notation and hypotheses of~\Ref{XIII.4.3}. If
$\bar{s}$ is a geometric point of~$S$ and $a=g(\bar{s})$, the section
$g$ defines a morphism
$$
w\colon \pi_1(S,a)\to \pi_1'(X,a),
$$
such that $\pi_1'(X,a)$ is identified with the semidirect product of
$\pi_1(S,a)$ by $\pi_1^\LL(X_{\bar{s}},a)$. The profinite group
$\pi_1(S,a)$ therefore acts on~$\pi_1^\LL(X_{\bar{s}}, a)$. Since
$\pi_1^\LL(X_{\bar{s}},a)$ is a strict inverse limit of groups
invariant under the action of~$\pi_1(S,a)$, specifying
$\pi_1^\LL(X_{\bar{s}},a)$ together with this action is equivalent
to specifying a strict inverse system of finite \'etale group schemes
over~$S$, which we denote by
$$
\pi_1^\LL(X/S,g,\bar{s})\quad\text{or
simply}\quad\pi_1^\LL(X/S,g).
$$
\label{indnot:mk}\oldindexnot{$\pi_1^\LL(X/S,g,\bar{s})$ or $\pi_1^\LL(X/S,g)$|hyperpage}%
%

The following properties then hold:
\setcounter{subsubsection}{0}

\ifthenelse{\boolean{orig}}{\subsubsection{}}{\sgaFixedHeading{\subsubsection}{XIII.4.5.1}}
\label{XIII.4.5.1}
For every finite \'etale group scheme $G$ over~$S$ whose fibers are
$\LL$-groups, the set $E$ of classes of torsors $P$ under the pullback
$G_X$ of~$G$ to~$X$, equipped with an isomorphism
$g^*P\isomto G$, is canonically isomorphic to the set
$$
\Hom_S(\pi_1^\LL(X/S,g,\bar{s}),G) \quad\text{modulo inner
automorphisms of~$G$.}
$$

\ifthenelse{\boolean{orig}}{\subsubsection{}}{\sgaFixedHeading{\subsubsection}{XIII.4.5.2}}
\label{XIII.4.5.2}
For every finite \'etale group scheme $G$ over~$S$ whose fibers are
$\LL$-groups, the sheaf $\R^1f_*G_X$ is canonically isomorphic to the
sheaf associated with the presheaf
$$
S'\mto \Hom_{S'}(\pi_1^\LL(X/S,g,\bar{s}),G)\
\text{modulo inner automorphisms of~$G$.}
$$
($S'$ denotes a scheme \'etale over~$S$).

\ifthenelse{\boolean{orig}}{\subsubsection{}}{\sgaFixedHeading{\subsubsection}{XIII.4.5.3}}
\label{XIII.4.5.3}
Let
% page: 421
\marginpar{421}
%
$S'$ be a connected
\ifthenelse{\boolean{orig}}
{$S$-cheme}
{$S$-scheme}
% correction: `sch\'ema'.
, let $\bar{s}$ be a geometric point of~$S'$, and let $X'$, $g'$ be the
respective pullbacks of~$X,g$ to~$S'$. Then
$\pi_1^\LL(X'/S',g',\bar{s})$ is canonically isomorphic to the
pullback of~$\pi_1^\LL(X/S,g,\bar{s})$ to~$S'$. For every geometric
point $\xi$ of~$S$, the fiber
$\pi_1^\LL(X/S,g,\bar{s})_\xi$ is isomorphic to $\pi_1^\LL(X_\xi)$.

Indeed, specifying~$G$ is equivalent to specifying an abstract
$\LL$-group $\mathbf{G}$
% remark: Cas non prevu dans la page web. J'ai choisi le gras par
% defaut.
on which $\pi_1(S,a)$ acts, thereby inducing an action of
$\pi_1'(X,a)$ on~$\mathbf{G}$. The isomorphism defined
in~\Ref{XIII.4.5.1} is then obtained, by restriction to the subset
$E$, from the canonical morphism
\[
\H^1(\pi_1'(X,a),\mathbf{G})\to
\H^1(\pi_1^\LL(X_{\bar{s}},a),\mathbf{G}) =
\Hom(\pi_1^\LL(X_{\bar{s}},a),\mathbf{G})/\text{inner aut. }\mathbf{G},
\]
with the set $E$ mapping bijectively onto the subset of morphisms
from $\pi_1^\LL(X_{\bar{s}},a)$ to $\mathbf{G}$ that are compatible
with the action of
$\pi_1(S,a)$. Assertion~\Ref{XIII.4.5.3} then follows from the
definition of~$\pi_1^\LL(X/S,g,\bar{s})$, taking into account the
exact homotopy sequence \eqref{eq:XIII.4.3.1}, and~\Ref{XIII.4.5.2}
follows from \Ref{XIII.4.5.1} and~\Ref{XIII.4.5.3}.

\begin{proposition}[K\"unneth formula]
\label{XIII.4.6}
%
\index{Kunneth formula@K\"unneth formula|hyperpage}%
%
Let $k$ be a separably closed field of characteristic
$p\geq 0$, let $X$ and $Y$ be two connected $k$-schemes, let $a$ be a
geometric point of~$X$, let $b$ be a geometric point of~$Y$, and let
$c$ be a geometric point of~$X\times_k Y$ lying over~$a$ and
$b$. Suppose that one of the following two conditions holds:
\begin{enumerate}
\item[a)] $X$ is of finite type over~$k$, and schemes of finite type
over an algebraic closure $\bar{k}$ of~$k$ of dimension
$\leq\dim X$ are strongly desingularizable. \textup{(SGA~5 I~3.1.5)}.
\item[b)] $X$ is quasi-compact and quasi-separated, and every scheme
of finite type over~$\bar{k}$ is strongly desingularizable.
\end{enumerate}

Then, if $p'$ is the set of prime numbers distinct from~$p$, the
morphism
\begin{equation*}
\label{eq:XIII.4.6.0} \tag{\thesubsection.0} {}
\pi_1^{p'}(X\times_k Y,c)\to\pi_1^{p'}(X,a)\times\pi_1^{p'}(Y,b),
\end{equation*}
induced
% page: 422
\marginpar{422}
%
by the homomorphisms on fundamental groups associated with the
projections
$$
X\times_k Y\to X \quad \text{and} \quad X\times_k Y\to Y,
$$
is an isomorphism.
\end{proposition}

We may assume that $k$ is algebraically closed and that $X$ is reduced
(SGA~4 VIII~1.1). Let $Z=X\times_k Y$, and let $g\colon X\to k$ and $f\colon
Z\to Y$ be the canonical morphisms. The morphism $g$, and hence also $f$,
is universally locally $1$-aspherical for $p'$ by~\Ref{XIII.3.4}.
Since $X$ is connected, $f$ is $0$-acyclic (SGA~4 XV~1.16). Moreover,
it follows from~\Ref{XIII.3.2} that, for every finite $p'$-group $C$,
$(C_X,g)$ is cohomologically proper relative to $k$ in dimension
$\leq 1$. It follows that $(C_Z,f)$ is cohomologically proper
relative to $Y$ in dimension $\leq 1$ (\Ref{XIII.1.5}~c)) and that
$f_*C_Z$ and $\R^1f_*C_Z$ are constant sheaves. Consequently, $f$
satisfies all the hypotheses of~\Ref{XIII.4.2}. We therefore have the
exact sequence
$$
\pi_1^{p'}(X_b,c)\to\pi_1^{p'}(Z,c)\to\pi_1^{p'}(Y,b)\to 1.
$$
Moreover, the composite morphism
$$
\pi_1^{p'}(X_b,c)\to\pi_1^{p'}(Z,c)\to\pi_1^{p'}(X,a)
$$
is an isomorphism, and we therefore have the exact sequence
$$
1\to\pi_1^{p'}(X,a)\to\pi_1^{p'}(Z,c)\to\pi_1^{p'}(Y,b)\to 1.
$$
On the other hand, the morphism \eqref{eq:XIII.4.6.0} defines a
morphism from this exact sequence to the exact sequence
$$
1\to\pi_1^{p'}(X,a)\to
\pi_1^{p'}(X,a)\times\pi_1^{p'}(Y,b)\to\pi_1^{p'}(Y,b)\to 1,
$$
and it follows that the morphism \eqref{eq:XIII.4.6.0} is an
isomorphism.
\end{filecontents*}
% Original input: drafts/SGA1_XIII_4_2_English.texfrag
\begin{filecontents*}[overwrite,noheader]{sga1-cumulative-input-114.tex}
\subsection{}
\label{XIII.4.7} Let
$$
\xymatrix@C=1.2cm{
X \ar[d]_-f&\ar[l] X_U \ar[d]_-{f_U}& \ar[l]X_{\bar{s}}\ar[d] \\
S&\ar[l] U &\ar[l]\bar{s}
}
$$
be a
% page: 423
\marginpar{423}
%
diagram whose squares are cartesian, where $S$ is an
arcwise connected scheme (SGA~4 IX~2.12), $U$ is a connected open subset of
$S$, and $\bar{s}$ is a geometric point of~$U$. Let $a$ be a
geometric point of~$X_{\bar{s}}$, and let $\LL$ be a set of
prime numbers. Let $g$ be a section of~$f$, and suppose that the
following conditions are satisfied:
\begin{enumerate}
\item[a)] The morphism $f$ is $0$-acyclic and locally $0$-acyclic,
and, for every \'etale covering $S'$ of~$S$ and every
\'etale covering $E$ of~$X\times_S S'$ which is a quotient of a
Galois covering whose group is an $\LL$-group, $(E,f_{(S')})$ is
cohomologically proper relative to~$S'$ in dimension $\leq
0$.
\item[b)] The morphism $f_U$ is locally $1$-aspherical for
$\LL$, and, for every finite constant sheaf of $\LL$-groups $F$ on~$X_U$, $(F,f_U)$ is cohomologically proper in dimension $\leq 1$, and
the fibers of~$\R^1 f_{U*}F$ are finite.
\end{enumerate}

We then deduce from~\Ref{XIII.4.1} and~\Ref{XIII.4.3} the following
commutative diagram, whose rows are exact:
\begin{equation*}
\label{eq:XIII.4.7.0} \tag{\thesubsection.0} {}
\begin{array}{c}
\xymatrix{
1 \ar[r]& \pi_1^{\LL}(X_{\bar{s}},a) \ar[r]\ar@{=}[d]& \pi'_1(X_U,a) \ar[r]\ar[d]& \pi_1(U,a)\ar[r]\ar[d]& 1 \\
&\pi_1^{\LL}(X_{\bar{s}},a) \ar[r]&\pi'_1(X,a) \ar[r]& \pi_1(S,a) \ar[r]& 1
}
\end{array}
\end{equation*}

Thanks to the section $g$, we have morphisms
\ifthenelse{\boolean{orig}}
{$\pi_1(U,a)\to\pi'_1(X_U,a)$ $\pi_1(S,a)\to\pi'_1(X,a)$;}
{$$\pi_1(U,a)\to\pi'_1(X_U,a),\quad \pi_1(S,a)\to\pi'_1(X,a)\,\;$$}
% correction: insert a comma.
% and put in display mode
from these we
deduce a morphism from the amalgamated sum of~$\pi_1(S,a)$ and
$\pi'_1(X_U,a)$ over~$\pi_1(U,a)$ into $\pi'_1(X,a)$:
\begin{equation*}
\label{eq:XIII.4.7.1} \tag{\thesubsection.1} {}
\varphi\colon\pi=\pi_1(S,a)\coprod_{\pi_1(U,a)}\pi'_1(X_U,a)\to\pi'_1(X,a).
\end{equation*}
Suppose that the following condition is satisfied:
\begin{enumerate}
\item[c)] If $T=S-U$, then $\prof \et_T(S)\geq 2$ (SGA~2 XIV~1.1).
\end{enumerate}
Then the functor which assigns to an \'etale covering of~$S$ its
restriction to $U$ is fully faithful
(SGA~2 XVI~1.4). It follows that the morphism
$\pi_1(U,a)\to\pi_1(S,a)$ is surjective (\ifthenelse{\boolean{orig}}{V~}{}\Ref{V.6.9}), and we
deduce from the
% page: 424
\marginpar{424}
%
diagram \eqref{eq:XIII.4.7.0} that the same is true of the morphism
$\pi'_1(X_U,a)\to\pi'_1(X,a)$; a fortiori, $\varphi$ is an
epimorphism. Let
\begin{equation*}
\label{eq:XIII.4.7.2} \tag{\thesubsection.2} {}
K=\Ker(\pi_1(U,a)\to\pi_1(S,a)).
\end{equation*}
The group $\pi$ of \eqref{eq:XIII.4.7.1} is identified with the quotient of
$\pi'_1(X_U,a)$ by the closed normal subgroup $L$ generated
by the image of~$K$ in $\pi'_1(X_U,a)$. Regard
$\pi'_1(X_U,a)$ as the semidirect product of~$\pi_1(U,a)$ by
$\pi_1^\LL (X_{\bar{s}},a)$. The group $K$ then acts by
inner automorphisms on~$\pi_1^\LL (X_{\bar{s}},a)$, and the
quotient $\pi=\pi'_1(X_U,a)/L$ is identified with the semidirect product of
$\pi_1(U,a)/K=\pi_1(S,a)$ by the group $\pi_1^\LL(X_{\bar{s}},a)_K$
of coinvariants of~$\pi_1^{\LL}(X_{\bar{s}},a)$ under~$K$. Finally, we have
an epimorphism
\begin{equation*}
\label{eq:XIII.4.7.3} \tag{\thesubsection.3} {}
\varphi\colon\pi=\pi_1^\LL(X_{\bar{s}},a)_K\cdot\pi_1(S,a)\to\pi'_1(X,a).
\end{equation*}
\label{indnot:ml}\oldindexnot{$\pi_1^\LL(X_{\bar{s}},a)_K$|hyperpage}%
%

The following proposition gives conditions under which the
morphism $\varphi$ is an isomorphism.

\setcounter{subsubsection}{3}

\begin{subproposition}
\label{XIII.4.7.4} The notation is that of~\Ref{XIII.4.7}. In
addition to conditions \textup{a), b), c)}, suppose that the following
conditions are satisfied:
\begin{enumerate}
\item[d)] For every point $t$ of~$T=S-U$, the morphism $f$ is
locally $1$-aspherical for $\LL$ at $g(t)$.
\item[e)] For every point $t\in T$, every irreducible component
of the fiber $X_t$ contains $g(t)$ and, for every point $x$ of
$X_t-\{g(t)\}$ which is not maximal, one has
$$
\prof\hop_x(X)\geq 3 \quad\textup{(SGA~2 XIV~1.2)}
$$
and the ring $\cal{O}_{X,x}$ is noetherian.
\end{enumerate}

Then the morphism \eqref{eq:XIII.4.7.3} is an isomorphism.
\end{subproposition}

As stated above, the group $\pi$ is identified with the
quotient of~$\pi'_1(X_U,a)$ by the closed normal subgroup $L$
generated by the image of~$K$ \eqref{eq:XIII.4.7.2} in
$\pi'_1(X_U,a)$. This amounts to saying that $\pi$ classifies the
% page: 425
\marginpar{425}
%
principal coverings $Z$ of~$X_U$ such that $g^{-1}_U(Z)$ extends
to an \'etale covering of~$S$, and which induce on~$X_{\bar{s}}$ a covering obtained by extension of the structure group
from a principal covering whose group is an
$\LL$-group. To prove that $\varphi$ is an isomorphism, it is
enough to show that such a covering $Z$ extends to all of $X$.

Let us first show that $Z$ extends to an \'etale covering
of an open subset containing $X_U$ and $g(S)$. Let $W$ be a scheme
\'etale over~$X$ whose image contains $X_U$ and $g(S)$, and put
$W_U=W\times_S U$. Since the morphism $W\to S$ is $0$-acyclic
and $\mathrm{prof}.\,\et_T S\geq 2$, we have
% remark: writing prof. et but prof hop (without a period) is a little
% ugly; why not omit the period everywhere?
$$
\mathrm{prof}.\,\et_{(W-W_U)} W\geq 2
$$
(SGA~2 XIV~1.13); consequently, if $Z|W_U$ extends to an
\'etale covering of~$W$, this extension is unique up to unique
isomorphism. It follows that the problem of
extending $Z$ to a neighborhood of~$g(S)\cup X_U$ is local for the
\'etale topology near the points of~$g(T)$. If $t$ is a
point of~$T$, put $x=g(t)$, and denote by $\bar{X}$ (\resp $\bar{S}$)
the strict localization of~$X$ at $\bar{x}$ (\resp of~$S$ at
$\bar{t}$), put $\bar{U}=U\times_S \bar{S}$, $\bar{X}_U=\bar{X}\times_X
X_U$, and let $\bar{g}\colon\bar{S}\to\bar{X}$ be the morphism induced by
$g$. It is enough to show that, for every point $t$ of~$T$, the inverse
image $\bar{Z}$ of~$Z$ on~$\bar{X}_U$ extends to $\bar{X}$
or, equivalently, is trivial. But, by definition
of~$Z$, the inverse image of~$\bar{Z}$ on~$\bar{U}$ is
trivial. To prove that $\bar{Z}$ is trivial, it is enough to show
that it is of the form $\bar{f}^*_U E$, where $E$ is a
principal covering of~$\bar{U}$; indeed, we will then have
$E\simeq\bar{g}^*_U\bar{f}^*_U E \simeq\bar{g}^*_U\bar{Z}$, whence
the result, since $\bar{g}^*_U\bar{Z}$ is trivial. Since the
morphism $\bar{f}_U$ is $0$-acyclic and locally $0$-acyclic,
to prove that $\bar{Z}$ comes from~$\bar{U}$, it is enough to
show that, for every algebraic geometric point over a
point of~$\bar{U}$, which we may suppose to be the point
$\bar{s}$, $\bar{Z}|\bar{X}_{\bar{s}}$ is trivial (SGA~4 XV~1.15). But
$\bar{Z}|\bar{X}_{\bar{s}}$, being obtained by extension of the
structure group from a principal covering whose group is an
$\LL$-group, is trivial because the morphism $\bar{f}$ is
$1$-aspherical for $\LL$.

We
% page: 426
\marginpar{426}
%
have therefore shown that there exists an open neighborhood $V$ of
$g(S)\cup X_U$ such that $Z$ extends to an \'etale covering
$Z_V$ of~$V$. Let us show that $Z_V$ extends to all of $X$.
It is enough to see that, for every point $x$ of~$X-V$, one has
$$
\prof\hop_x X\geq 3.
$$
But this follows from hypothesis e) and from the fact that a point $x$
of~$X-V$ cannot be maximal in its fiber $X_t$: since every
irreducible component of~$X_t$ contains $g(t)$, every maximal point
of~$X_t$ belongs to $V$.

\begin{corollaire} \label{XIII.4.8}
The hypotheses are those of~\Ref{XIII.4.7.4}, but suppose in
addition that $\pi_1(S,a)=1$. Then there is an isomorphism
\begin{equation*}
\label{eq:XIII.4.8.*} \tag{$*$}
{}\pi_1^\LL(X,a)\isomto\pi_1^\LL(X_{\bar{s}},a)_K.
\end{equation*}
In particular, if $\pi_1^\LL(X_{\bar{s}},a)$ is topologically
finitely presented and if $K$ acts on~$\pi_1^\LL(X_{\bar{s}},a)$ through a finitely
generated group, then $\pi_1^\LL(X,a)$ is topologically finitely
presented.
\end{corollaire}

The isomorphism \eqref{eq:XIII.4.8.*} was proved
in~\Ref{XIII.4.7.4}. Suppose that $\pi_1^\LL(X_{\bar{s}},a)$ is the quotient
of the free pro-$\LL$-group on $n$ generators
$L(x_1,\dots,x_n)$ by the closed normal subgroup generated
by the elements $y_1,\dots,y_p$ of~$L(x_1,\dots,x_n)$, and
suppose that $K$ acts through a group
generated by elements $k_1,\dots,k_q$. If, for every
$i\in[1,n]$, $j\in[1,q]$, $z_{ij}$ denotes an element of
$L(x_1,\dots,x_n)$ lifting the element $(k_j\cdot
x_i)x_i^{-1}$, then $\pi_1^\LL(X_{\bar{s}},a)_K$ is the quotient of
$L(x_1,\dots,x_n)$ by the closed normal subgroup generated
by the elements $(y_i)_{i\in[1,p]}$,
$(z_{ij})_{i\in[1,n],j\in[1,q]}$.

% remark: mistake in original text: numbering goes back to 4.6.
\ifthenelse{\boolean{orig}}
{\setcounter{subsection}{5}}
{}

\begin{remarques}
\label{rem:XIII.4.6}
a) Conditions a) through e) of~\Ref{XIII.4.7} are
satisfied when $S$ is a connected normal scheme, $U$ is a
dense retrocompact open subset of~$S$, and $f$ is a proper morphism of
finite presentation, with geometrically connected and irreducible fibers
at every point $t$ of~$T$, $f$ being moreover separable,
% page: 427
\marginpar{427}
%
smooth at the points of~$X_U\cup g(T)$, $\LL$ being the set of
prime numbers distinct from the residue characteristics of
$S$, and $X$ being regular at every point of~$X_t$. Indeed,
condition a) follows from SGA~4 XV~4.1 and 1.4; conditions
b) and d) follow from~\Ref{XIII.1.4} and SGA~4 XV~2.1 and
XVI~5.2. Finally, e) follows from SGA~2 XIV~1.11.
% remark: reference to SGA 4?
\par\smallskip
b) Corollary~\Ref{XIII.4.8} applies to compute the
fundamental group $\pi_1^{p'}(X)$ of a proper smooth surface $X$
over a separably closed field $k$ of characteristic $p$
($p'$~denoting the set of prime numbers distinct from
$p$). The method was communicated to us by
J.P\ptbl \textsc{Murre}; it consists in reducing, by blowing up
$X$, to the case where there is a fibration $X\to\PP^1_k$ and an open subset $U$
of~$\PP^1_k$ satisfying the hypotheses of~\Ref{XIII.4.7} (see
SGA~7 for more details). The same method may be
used more generally (\loccit) to prove that, if
$X$ is a connected $k$-scheme of finite type and if the schemes
of finite type of dimension $\leq\dim X$ over an algebraic
closure of~$k$ are strongly desingularizable
(SGA~5 I~3.1.5), then $\pi_1^{p'}(X)$ is topologically finitely
presented.
\end{remarques}
\end{filecontents*}
% Original input: drafts/SGA1_XIII_4_OPENING_AND_0_English.texfrag
\begin{filecontents*}[overwrite,noheader]{sga1-cumulative-input-115.tex}
\section{Exact homotopy sequences}\label{XIII.4}\setcounter{subsection}{-1}
\subsection{}
\label{XIII.4.0}
Let $X$ and $S$ be two connected schemes, let $f\colon X\to S$ be a
morphism, let $a$ be a geometric point of~$X$, and let $\LL$ be a set
of prime numbers. Let $K$ be the kernel of the canonical homomorphism
$\pi_1(X,a)\to\pi_1(S,a)$, and let $N$ be the smallest normal
pro-subgroup of~$K$ such that $K/N$ is the pro-$\LL$-group
$K^{\LL}$. Then $N$ is normal in $\pi_1(X,a)$, and we denote by
$$
\pi_1'(X,a)
$$
\label{indnot:mj}\oldindexnot{$\pi_1'(X,a)$|hyperpage}%
%
the
% page: 415
\marginpar{415}
%
quotient of~$\pi_1(X,a)$ by $N$. If $a$ is a geometric point
of a geometric fiber~$X_{\bar{s}}$, the canonical morphisms
$$
\pi_1(X_{\bar{s}},a)\to \pi_1(X,a) \to
\pi_1(S,a)
$$
define canonical morphisms
$$
\pi_1^{\LL}(X_{\bar{s}},a)\lto{u} \pi_1'(X,a) \lto{v} \pi_1(S,a)
$$
We have $vu=0$.

\begin{proposition}
\label{XIII.4.1}
Let $S$ be a connected scheme, and let $f\colon X\to S$ be a
locally $0$-acyclic morphism \textup{(SGA~4 XV~1.11)}; suppose moreover that $f$
is $0$-acyclic (which, when $f$ is coherent, amounts to
saying that the geometric fibers of~$f$ are connected \textup{(SGA~4
XV~1.16)}). Let $\LL$ be a set of prime numbers. If $S'$ is a
scheme \'etale over~$S$, denote by $X'$ and $f'$ the inverse images
of~$X$ and $f$ over~$S'$. Suppose that, for every \'etale covering
$S'$ of~$S$ and every \'etale covering $E$ of
$X'$ that is a quotient of a Galois covering whose group is an
$\LL$-group, $(E,f')$ is cohomologically proper relative
to~$S'$ in dimension $\leqslant 0$ and $f'_* E$ is
constructible. Then, if $\bar{s}$ is a geometric point of
$S$ and $a$ is a geometric point of the fiber $X_{\bar{s}}$, the
sequence of group homomorphisms
\begin{equation*}
\label{eq:XIII.4.1.1}
\tag{\thesubsection.1}
\pi_1^{\LL}(X_{\bar{s}},a)\lto{u} \pi_1'(X,a)\lto{v}\pi_1(S,a)\to1
\end{equation*}
is exact.
\end{proposition}

This statement generalizes X.\Ref{X.1.4}, whose proof we shall
repeat.

Let us first show that $v$ is surjective. It suffices to show that, for
every connected \'etale covering $S'$ of~$S$, $X'$ is also
connected (\ifthenelse{\boolean{orig}}{V~}{}\Ref{V.6.9}). Let $C$ be a set having at least two
elements. Since $f$ is $0$-acyclic, the canonical morphism
$$
\H^0(S',C_{S'})\to \H^0(X',C_{X'})
$$
is bijective, and hence $X'$ is connected, which proves the surjectivity
of~$v$.

By
% page: 416
\marginpar{416}
%
definition of~$K^{\LL}$ \eqref{XIII.4.0}, there is an exact sequence
$$
1\to K^{\LL} \to \pi_1'(X,a)\to \pi_1(S,a)\to 1
$$
Let $\widetilde{S}$ be the universal covering of~$S$, and put
$\widetilde{X}=\widetilde{S}\times_SX$; the group $K^{\LL}$ classifies the
Galois coverings $P$ whose group is an $\LL$-group and for which
there exist an \'etale covering $S'$ of~$S$ and a Galois covering
$Q$ of~$X'=X\times_SS'$ such that there is an isomorphism
$P\simeq Q\times_{X'}\widetilde{X}$. For the sequence
\eqref{eq:XIII.4.1.1} to be exact, it is necessary and sufficient that the
canonical morphism
$$
\pi_1^{\LL}(X_{\bar{s}},a)\to K^{\LL}
$$
be surjective. By the interpretation of~$K^{\LL}$, this
amounts to saying that, for every \'etale covering $S'$ of~$S$
and every Galois covering $Q$ of~$X'$ whose group is an
$\LL$-group, such that $P=Q\times_{X'}\widetilde{X}$ is connected,
$Q|X_{\bar{s}}$ is connected. Let us show that this last
condition is satisfied. Indeed, let $S'$ be an \'etale covering
of~$S$, and let $Q$ be a Galois covering of~$X'$ whose group is an
$\LL$-group $F$, such that $Q|X_{\bar{s}}$ is disconnected; let us show
that, after possibly replacing $S'$ by an \'etale covering, $Q$
becomes disconnected. There is a subgroup $G$ of~$F$ distinct from
$F$ and a torsor $R$ under $G|X_{\bar{s}}$ such that $Q|X_{\bar{s}}$
is obtained by extension of structure group along $G\to F$ from $R$.
The \'etale covering $E=Q/G$ of~$X'$ is such that
$E|X_{\bar{s}}$ has a section. By~\Ref{XIII.1.16}, $f'_*E$
is locally constant constructible and, after possibly replacing $S'$
by an \'etale covering, we may even suppose that $f'_*E$
is constant. Since $(E,f')$ is cohomologically proper relative
to~$S'$ in dimension $\leqslant0$ and $\H^0(X_{\bar{s}},
E|X_{\bar{s}})$ is nonempty, one sees that $E$ has a section. But this
proves that $Q$ is disconnected, which completes the proof.

From~\Ref{XIII.4.1}, we deduce the following lemma, which will be used
in~\Ref{XIII.4.6}.

\begin{lemme}
\label{XIII.4.2}
Let $S$ be a connected scheme, let $f\colon X\to S$ be a
locally $0$-acyclic and $0$-acyclic morphism, and let $\LL$ be a set of
prime numbers. Suppose
% page: 417
\marginpar{417}
%
that, for every finite constant sheaf of~$\LL$-groups $F$ on~$X$,
$(F,f)$ is cohomologically proper relative to~$S$ in
dimension $\leqslant1$, and that, for every \'etale covering
$S'$ of~$S$ and every \'etale covering $E$ of~$X'$ that is a
quotient of a Galois covering whose group is an $\LL$-group,
$f'_*E$ is constructible. Then, if $\bar{s}$ is a geometric point
of~$S$ and $a$ is a geometric point of the
fiber $X_{\bar{s}}$, the sequence of group homomorphisms
$$
\pi_1^{\LL}(X_{\bar{s}},a)\to\pi_1'(X,a)\to\pi_1(S,a)\to1
$$
is exact.
\end{lemme}

The hypotheses of~\Ref{XIII.4.1} are satisfied. Indeed, it follows
from~\Ref{XIII.1.13} 3) that, for every scheme $S'$
\'etale over~$S$ and every \'etale covering $E$ of~$X'$ that is a
quotient of a Galois covering whose group is an $\LL$-group,
$(E,f')$ is cohomologically proper relative to~$S'$ in
dimension $\leqslant 0$.
\begin{proposition}
\label{XIII.4.3}
Let $S$ be a connected scheme, let $\LL$ be a set of prime
numbers, let $f\colon X\to S$ be a $0$-acyclic morphism, locally
$1$-aspherical for $\LL$ \textup{(SGA~4 XV~1.11)}, and let $g\colon S\to X$ be a
section of~$f$. Let $\bar{s}$ be a geometric point of~$S$, and let
$a$ be a geometric point of the fiber $X_{\bar{s}}$. Suppose
that, for every constant sheaf of
\ifthenelse{\boolean{orig}}
{$\LL$-group}
{$\LL$-groups}
% correction: agreement error: groups
~$F$, $(F,f)$ is cohomologically proper in dimension
$\leqslant1$, that the direct image under~$f$ of the stack of torsors under
$F$ is a $1$-constructible stack \eqref{XIII.0}, and that, for every \'etale covering $S'$ of~$S$ and every
\'etale covering $E$ of~$X'=X\times_SS'$ that is a quotient of a Galois
covering whose group is an $\LL$-group, $f'_*E$ is constructible. Then
the sequence of group homomorphisms
\begin{equation*}
\label{eq:XIII.4.3.1}
\tag{\thesubsection.1}
1\to\pi_1^{\LL}(X_{\bar{s}},a)\lto{u} \pi_1'(X,a)\lto{v}\pi_1(S,a)\to1
\end{equation*}
is exact.
\end{proposition}

In view of~\Ref{XIII.4.2}, it suffices to prove the injectivity
of the morphism $u$, \ie to prove that, for every principal
covering $\overline{Z}$ of~$X_{\bar{s}}$
% page: 418
\marginpar{418}
%
whose group is an $\LL$-group~$C$, there exist an \'etale covering
$Z$ of~$X$ and a morphism from a connected component of~$Z|X_{\bar{s}}$
to~$\overline{Z}$ (\ifthenelse{\boolean{orig}}{V~}{}\Ref{V.6.8}). Thus let $\overline{Z}$ be a
principal covering of~$X_{\bar{s}}$ whose group is an $\LL$-group~$C$, and let~$\bar{z}$ be its class in
$\H^1(X_{\bar{s}},C_{X_{\bar{s}}})$. By \Ref{XIII.1.5}~d), there is
a canonical isomorphism
$$
\left(\R^1f_*C_X\right)_{\bar{s}}\isomto\H^1(X_{\bar{s}},C_{X_{\bar{s}}}),
$$
and by~\Ref{XIII.1.16}, $\R^1f_*C_X$ is a
locally constant constructible sheaf. Thus one can find an
\'etale covering $S'$ of~$S$ such that $\R^1f_*C_X|S'$ is
constant. If $\bar{s}\to S'$ is a geometric point
above the geometric point $\bar{s}\to S$, there is an
element~$z$ of~$\H^0(S', \R^1f_*C_X)$ whose image in
$\H^1(X_{\bar{s}},C_{X_{\bar{s}}})$ is $\bar{z}$. By the
lemma~\Ref{XIII.4.3.1} below, one can find an \'etale covering
$S'_1$ of~$S'$ and a torsor $P$ on~$X'_1$ with group~$C$
whose image in $\H^0(S'_1,\R^1f_*C)$ is equal to the
restriction of~$z$. The torsor~$P$ is represented by an
\'etale covering $Z$ of~$X'_1=X\times_SS'_1$ such that
$Z\times_{X'_1}X_{\bar{s}}$ is isomorphic to $\overline{Z}$. If
$Z$ is regarded as an \'etale covering of~$X$, there is then a
morphism from~$Z\times_XX_{\bar{s}}$ to $\overline{Z}$, which
completes the proof.

\begin{sublemme}
\label{XIII.4.3.1}
Let $f\colon X\to S$ be a $0$-acyclic and locally
$0$-acyclic morphism, and let $g$ be a section of~$f$. Let $C$ be a finite constant
group such that $(C_X,f)$ is cohomologically proper in dimension
$\leqslant0$ and the direct image under $f$ of the stack of torsors under
$C_X$ is constructible. Then, for every section $z$ of
$\H^0(S,\R^1f_*C_X)$, one can find an \'etale covering
$S_1$ of~$S$ and, if $X_1=X\times_SS_1$, an element of
$\H^1(X_1,C_{X_1})$ whose image under the canonical morphism
$$
\H^1(X_1,C_{X_1})\to\H^0(S_1,\R^1f_{1*}C_{X_1})
$$
is equal to the restriction of~$z$ to
$\H^0(S_1,\R^1f_{1*}C_{X_1})$.
\end{sublemme}

For every scheme $S'$ \'etale over~$S$, put $X'=X\times_SS'$,
and denote by $g'$ the inverse image of~$g$ under the morphism
% page: 419
\marginpar{419}
%
$S'\to S$. The presheaf $G$ on~$S$ defined by
\newlength{\tmpRacinet} \setlength{\tmpRacinet}{\textwidth}
\addtolength{\tmpRacinet}{-7em}
$$
G(S') = \left\{\text{\parbox{\tmpRacinet}{isomorphism classes of
		torsors $P$ on~$X'$ with group $C_{X'}$, endowed with an
		isomorphism $g'^*P\isomto C_{S'}$ }} \right\}
$$
% remark:
% Layout perhaps debatable. Advantage: it adapts
% automatically to a change of format, if the font size
% is adapted to the page size
is then a sheaf. Indeed, this follows by descent from the fact
that an automorphism of a torsor $P$ on~$X'$ is completely
determined by its restriction to $g'(S')$. Moreover, there is a
surjective morphism
$$
G\to \R^1f_*C_X.
$$
Let $z$ be an element of~$\H^0(S,\R^1f_*C_X)$, and let $H$ be the
subsheaf of~$G$ that is the inverse image of~$z$. It suffices to show
that $H$ is a locally constant constructible sheaf. Since this
property is local on~$S$, we may suppose that $z$
comes from an element of~$\H^1(X,C_X)$ represented by
a torsor $P$ such that $g^*P$ is isomorphic to $C_S$. To give an
isomorphism $i\colon g^*P\isomto C_S$ amounts to giving a
global section of~$\SheafAut_{C_S}(g^*P)$, and two isomorphisms $i$
and $i'$ define the same element of~$G(S)$ if and
only if $\ifthenelse{\boolean{orig}}{ii^{-1}i'}{i^{-1}i'}$ is the image of an element of
$\Aut_{C_X}(P)$. Consider the canonical injection
$$
f_*\SheafAut_{C_X}(P)\to\SheafAut_{C_S}(g^*P)\simeq C,
$$
so $H$ is identified with the quotient of~$\SheafAut_{C_S}(g^*P)$ by
$f_*\SheafAut_{C_X}(P)$. By~\Ref{XIII.1.16},
$f_*\SheafAut_{C_X}(P)$ is locally constant; hence the same is true
of~$H$, which completes the proof.

\begin{exemples}
\label{XIII.4.4}
Note that, if $S$ is a connected scheme, the hypotheses
of~\Ref{XIII.4.1} are satisfied when the morphism $f$ is proper,
flat, and of finite presentation, with separable, connected geometric
fibers, and $\LL$ is arbitrary
(\cf \ifthenelse{\boolean{orig}}{X~}{}\Ref{X.1.3}). The hypotheses of~\Ref{XIII.4.3} are
satisfied if, in addition, $f$ is smooth and has a section, with
$\LL$ denoting the set of prime numbers distinct from the
residual characteristics of~$S$ (SGA~4 XV~2.1 and XVI~5.2).

The hypotheses of~\Ref{XIII.4.1} are also satisfied if $S$ is connected,
and if there is a scheme~$Z$ that is proper, flat, and of finite presentation over~$S$, with
% page: 420
\marginpar{420}
%
separable, connected geometric fibers, such that $X$ is
the complement in $Z$ of a normal crossings divisor
relative to~$S$, with $\LL$ the set of
prime numbers distinct from the residual characteristics of~$S$
\eqref{XIII.2.9}. The hypotheses of~\Ref{XIII.4.3} are
satisfied if, in addition, $f$ is smooth and has a section.
\end{exemples}
\end{filecontents*}
% Original input: drafts/SGA1_XIII_APPENDIX_II_AND_BIBLIOGRAPHY_English.texfrag
\begin{filecontents*}[overwrite,noheader]{sga1-cumulative-input-116.tex}
\section[Appendix II: finiteness for direct
images of stacks]{Appendix II: finiteness theorem for direct
images of stacks}
\label{XIII.6}

\begin{proposition}
\label{XIII.6.1}
Let $S$ be a locally noetherian scheme and let $f\colon X\to S$
be a morphism. If $S'$ is an $S$-scheme, we denote by $X'$
(\resp $f'$, etc.) the pullback of~$X$ (\resp $f$, etc.) by the
morphism $S'\to S$. Suppose that, for every
% page: 436
\marginpar{436}
%
scheme $S'$ \'etale over~$S$ and every constructible sheaf of sets
$F$ on~$X'$, $f'_*F$ is constructible, and that, for every
constructible sheaf of groups $F$ on~$X'$, $\R^{1}f'_*F$ is
constructible. Let $\Phi$ be a $1$-constructible stack on~$X$~\eqref{XIII.0}.
Then $f_*\Phi$ is $1$-constructible.
\end{proposition}

For every scheme $S'$ \'etale over~$S$ and every object $x$ of
$(f_*\Phi)_{S'}$, there is an isomorphism
$$
\SheafAut_{S'}(x)\simeq f'_*\SheafAut_{X'}(x) \quoi,
$$
where, on the right-hand side, $x$ is regarded as an object
of~$\Phi_{X'}$. The hypotheses therefore imply that $f_*\Phi$ is
constructible. Let $S\Phi$ be the sheaf of maximal subgerbes of
$\Phi$~\cite[III~2.1.7]{XIII.1}. Since $f_*(S\Phi)$ is constructible, we
may apply (SGA~4 IX~2.7) to it, and the fact that $f_*\Phi$ is
$1$-constructible then follows from the sublemma below.

\begin{sublemme}
\label{XIII.6.1.1}
Let $S$ be a locally noetherian scheme, let $f\colon X\to S$
be a morphism, and let $\Phi$ be a stack on~$X$. Suppose we are given
a sheaf on~$S$, representable by an \'etale $S$-scheme of finite
type $T$, a surjective morphism
$$
a\colon T\to f_*(S\Phi )
$$
and an object $p$ of the fiber $\Phi_{X_T}$ (where $X_T=X\times_ST$),
defining in $f_*(S\Phi )(T)=\allowbreak S\Phi (X_T)$ an element equal
to the image $q$ under $a$ of the identity section of
$T(T)$. Let $f_T\colon X_T\to T$ be the canonical morphism, and suppose
that the sheaf $\R^1f_{T*}(\SheafAut_{X_T}(p))$ is constructible;
then the same holds for~$S(f_*\Phi )$.
\end{sublemme}

The canonical morphism $f^*f_*\Phi\to\Phi$ gives a morphism
$$
S(f^*f_*\Phi )\simeq f^*(S(f_*\Phi ))\to S\Phi ,
$$
and hence a canonical morphism
$$
\varphi\colon S(f_*\Phi )\to f_*(S\Phi ).
$$
Let
% page: 437
\marginpar{437}
%
$F=S(f_*\Phi )$, and let $G$ be the image of~$F$ under $\varphi$; by
(SGA~4 IX,~2.9), $G$ is a constructible sheaf.

It is enough to show that, for every point $s$ of~$S$, there exists a
nonempty open subset~$U$ of~$\overline{\{s\}}$ such that $F|U$ is locally
constant constructible. Let $s\in S$, let $\overline{s}$ be a
geometric point lying over~$s$, and let
$\overline{q}_1,\dots,\overline{q}_n$ be the elements of~$G_{\overline{s}}$. By
the definition of~$T$, there exist $S$-morphisms
$h_i\colon\overline{s}\to T$ such that
$\overline{q}_i=h_i^*(q)$. Let $S'$ be the fiber product over~$S$ of
$n$ schemes isomorphic to $T$, let $\overline{s}\to S'$ be the fiber
product of the $h_i$, let $X'=X\times_SS'$, and let $q_i$ (\resp $p_i$)
be the pullback of~$q$ (\resp $p$) by the $i$-th projection of~$S'$
onto~$T$. If $\Psi_i$ is the maximal subgerbe of~$\Phi|X'$ generated
by $p_i$, the sheaf
$F_i=\R^1f_*'(\SheafAut_{X'}(p_i))$ is precisely the sheaf
$S(f_*'\Psi_i)$ of maximal subgerbes of~$f_*'\Psi_i$. In
particular, the canonical injection $\Psi_i\to\Phi |X'$ gives a
morphism
$$
\alpha_i\colon F_i\to F|S'.
$$

We shall show that $\alpha_i$ is a bijection from~$F_i$ onto the
inverse image of~$q_i$ in $F|S'$. For every scheme $S''$
\'etale over~$S'$, every section $y$ of~$F_i(S'')$ has image
$q_i|S''$ in $F(S'')$, because, locally for the \'etale topology
on~$S''$, $y$ is defined by an object~$x$ of~$\Phi_{X''}$ that is
isomorphic to $p_i|X''$. Conversely, if $y\in F(S'')$ has
image $q_i|S''$ in $F(S'')$, then, locally for the \'etale topology
on~$S''$, $y$ is defined by an object~$x$ of~$\Phi_{X''}$ that is
isomorphic to $p_i$; consequently, $x$ is an object of~$\Psi_{iX''}$
and therefore $y\in F_i(S'')$.

The proof is completed using~\Ref{XIII.6.1.2} below. Indeed, one can
find an open neighborhood $U'$ of~$s$ such that
$q_1|U',\dots,q_n|U'$ are sections of~$G(U')$ and generate this
sheaf. Since the $F_i|U'$ and $G|U'$ are constructible, so is~$F|U'$
by~\Ref{XIII.6.1.2}; after replacing $U$ by a smaller open subset if
necessary, $F|U$ is locally constant, which completes the proof.

\begin{sublemme}
\label{XIII.6.1.2}
Let
% page: 438
\marginpar{438}
%
$S$ be a locally noetherian scheme, and let $F\to G$ be a
surjective morphism of sheaves of sets on~$S$. Let $q_i$ be a
finite family of sections of~$G$ on~$S$ that generate $G$, and, for
each $i$, let $F_i$ be the subsheaf of~$F$ that is the inverse image
of $q_i$. Then, if $G$ and the $F_i$ are constructible, so is~$F$.
\end{sublemme}

To prove that $F$ is constructible, it is enough to show that, for
every point~$s$ of~$S$, there exists an open neighborhood $U$ of~$s$
such that $F|U$ is locally constant constructible. Let $s$ be a
point of~$S$. Since the sheaves $F_i$ and $G$ are constructible, one
can find an open neighborhood $U$ of~$s$ such that $F_i|U$ and $G|U$
are locally constant. We now show that $F|U$ is locally constant. By
(SGA~4 IX~2.13~(i)), it is enough to see that, if
$\overline{s}$ is a geometric point lying over~$s$,
$\tilde{s}$ is a geometric point of~$U$, and
$\overline{s}\to\tilde{s}$ is a specialization morphism, the
canonical morphism
$$
F_{\tilde{s}}\to F_{\overline{s}}
$$
is bijective.

Consider the commutative diagrams
$$
\xymatrix{ (F_i)_{\tilde{s}} \ar[r]^-{\tilde{a}_i} \ar[d]^{\wr} &
F_{\tilde{s}} \ar[r]^-{\tilde{a}} \ar[d]_b & G_{\tilde{s}}
\ar[d]^{\wr}\\ (F_i)_{\overline{s}} \ar[r]^-{\overline{a}_i} &
F_{\overline{s}} \ar[r]^-{\overline{a}} & G_{\overline{s}} }
$$
Let $\overline{q}_i$ (\resp $\tilde{q}_i$) be the pullback of~$q_i$
in $G_{\overline{s}}$ (\resp $G_{\tilde{s}}$); the morphisms
$\overline{a}$ and $\tilde{a}$ are surjective, and the morphism
$\overline{a}_i$ (\resp $\tilde{a}_i$) induces a bijection from
$(F_i)_{\overline{s}}$ (\resp $(F_i)_{\tilde{s}}$) onto~$\overline{a}^{-1}(\overline{q}_i)$
(\resp onto~$\tilde{a}^{-1}(\tilde{q}_i)$). It therefore follows
from the diagram above that $b$ is an isomorphism.

\begin{corollaire}
\label{XIII.6.2}
Let
% page: 439
\marginpar{439}
%
$S$ be a locally noetherian scheme and let $f\colon X\to S$
be a proper morphism. If $\Phi$ is a $1$-constructible stack on~$X$,
then $f_*\Phi$ is a \hbox{$1$-constructible} stack.
\end{corollaire}

The proof of~\Ref{XIII.6.1} also proves the following result, in view
of \Ref{XIII.2.4}~2).

\begin{corollaire}
\label{XIII.6.3}
Let $S$ be a locally noetherian scheme, let $f\colon X\to S$ be a
morphism, let $D$ be a divisor on~$X$ with normal crossings
relative to~$S$ \eqref{XIII.2.1}, let $Y=\Supp D$, let $U=X-Y$, and
let $i\colon U\to X$ be the canonical immersion. Let $\Phi$ be a
stack on~$U$ given, locally for the \'etale topology on~$X$ and $S$,
as the pullback of a stack $\Psi$ that is $1$-constructible on~$S$. Then
the stack $i_*^{\tame}\Phi$ is $1$-constructible.
\end{corollaire}

%\section{Bibliographie}

\begin{thebibliography}{0}

\bibitem[1]{XIII.1} J\ptbl \textsc{Giraud}, Cohomologie non ab\'elienne de
degr\'e~$2$, thesis, Paris (1966).

\bibitem[2]{XIII.2} J\ptbl \textsc{Giraud},
\ifthenelse{\boolean{orig}}
{Cohomologie non ab\'elienne de degr\'e~$2$, forthcoming.}
{Cohomologie non ab\'elienne, Springer, Berlin-New\kern2pt York (1971).}

\bibitem[3]{XIII.3}
\ifthenelse{\boolean{orig}}
{\textsc{Seifert}-\textsc{Threlfall}, Lehrbuch der Topologie, Chelsea, New
York (1934).}
{H\ptbl\textsc{Seifert}, W\ptbl\textsc{Threlfall}, Lehrbuch der Topologie, Chelsea, New
York (1934).}


\bibitem[4]{XIII.4}
\ifthenelse{\boolean{orig}}{J.P\ptbl}{J-P\ptbl}%
\textsc{Serre}, Cohomologie galoisienne, Lecture Notes
\no 5, Springer, Berlin (1964).

\bibitem[5]{XIII.5}
\ifthenelse{\boolean{orig}}{J.P\ptbl}{J-P\ptbl}%
\textsc{Serre}, Corps locaux, Hermann, Paris (1962).
\end{thebibliography}

\backmatter
\def\indexname{Terminological Index}
\printindex

\def\indexname{Index of Notation}


\begin{theindex}

\medskip\item $\Delta_{X/Y}$ or simply $\Delta $\dotfill\pageref{indnot:ab}
\medskip\item $\mathit {\Omega }^1_{X/Y}$\dotfill\pageref{indnot:ac}
\medskip\item $\cal {P}_{X/Y}^n$\dotfill\pageref{indnot:ad}, \pageref{indnot:ad1}
\medskip\item $\Delta_{X/Y}^n$\dotfill\pageref{indnot:ae}
\medskip\item $\goth {d}_{X/Y}$\dotfill\pageref{indnot:af}
\medskip\item $d_{X/S}^n$\dotfill\pageref{indnot:bb}
\medskip\item $\goth {g}_{X/S}$\dotfill\pageref{indnot:cb}
\medskip\item $\cal {C}(\pi )$\dotfill\pageref{indnot:eb}
\medskip\item $\pi $\dotfill\pageref{indnot:ec}
\medskip\item $\Prodash \cal {C}(\pi)$\dotfill\pageref{indnot:ed}
\medskip\item $\Gamma $\dotfill\pageref{indnot:ef}
\medskip\item $\pi_1(S,a)$\dotfill\pageref{indnot:eg}
\medskip\item $\pi_1(S;a,a')$\dotfill\pageref{indnot:eh}
\medskip\item $\cal {C}(S)$\dotfill\pageref{indnot:ei}
\medskip\item $\pi_1(f;a')$\dotfill\pageref{indnot:ej}
\medskip\item $\Ens $\dotfill\pageref{indnot:fb}
\medskip\item $\Cat $\dotfill\pageref{indnot:fc}
\medskip\item $\Ob (\cal {C})$\dotfill\pageref{indnot:fd}
\medskip\item $\Fl (\cal {C})$\dotfill\pageref{indnot:fe}
\medskip\item $\SheafHom (\cal {C},\cal {C}')$\dotfill\pageref{indnot:ff}
\medskip\item $\cal {C}^\circ $\dotfill\pageref{indnot:fg}
\medskip\item $\Cat_{/\cal {E}}$\dotfill\pageref{indnot:fh}
\medskip\item $\SheafHom_{\cal {E}/-}(\cal {F},\cal {G})$\dotfill\pageref{indnot:fi}
\medskip\item $v*u$\dotfill\pageref{indnot:fj}
\medskip\item $\cal {F}\times_{\cal {E}}\cal {G}$\dotfill\pageref{indnot:fk}
\medskip\item $\lambda^*\colon \Cat_{/\cal {E}}\to \Cat_{/\cal {E}'}$\dotfill\pageref{indnot:fl}
\medskip\item $\bf {\Gamma }(\cal {G}/\cal {E})$ and $\Gamma (\cal {G}/\cal {E})$\dotfill\pageref{indnot:fm}
\medskip\item $\cal {F}_S$\dotfill\pageref{indnot:fn}
\medskip\item $f^*_{\cal {F}}(\xi )$ or $f^*(\xi )$\dotfill\pageref{indnot:fo}
\medskip\item $\alpha_f(\xi )$\dotfill\pageref{indnot:fp}
\medskip\item $\SheafHom_\cart (\cal {F},\cal {G})$\dotfill\pageref{indnot:fq}
\medskip\item $\Cat^\cart_{/\cal {E}}$\dotfill\pageref{indnot:fr}
\medskip\item $\varprojLim \cal {F}/\cal {E}$\dotfill\pageref{indnot:fs}
\medskip\item $f_*^{\cal {F}}$ or~$f_*$\dotfill\pageref{indnot:ft}
\medskip\item $\Sch $\dotfill\pageref{indnot:hb}
\medskip\item $\bbmu_{n\,S}$\dotfill\pageref{indnot:kb}
\medskip\item $X^\an $\dotfill\pageref{indnot:lb}
\medskip\item $f^\an $\dotfill\pageref{indnot:lc}
\medskip\item $SF$ or $S(F)$\dotfill\pageref{indnot:mb}
\medskip\item $\H^1_{\tame}(U,F)$\dotfill\pageref{indnot:mc}
\medskip\item $\R^1_{\tame} g_* F$\dotfill\pageref{indnot:md}
\medskip\item $C_t((U,X)/S)$ or $C_t$\dotfill\pageref{indnot:me}
\medskip\item $\pi_1^{\tame}((U,X)/S,a)$ or $\pi_1^{\tame}(U,a)$ or $\pi_1^{\tame}(U)$\dotfill\pageref{indnot:mf}
\medskip\item $(g_*^{\tame}\Phi )_{T'}$\dotfill\pageref{indnot:mg}
\medskip\item $\H^0(V,C_V)^{\Pi }$\dotfill\pageref{indnot:mh}
\medskip\item $\pi_1^{\LL }(U,a)$\dotfill\pageref{indnot:mi}
\medskip\item $\pi_1'(X,a)$\dotfill\pageref{indnot:mj}
\medskip\item $\pi_1^\LL (X/S,g,\bar {s})$ or $\pi_1^\LL (X/S,g)$\dotfill\pageref{indnot:mk}
\medskip\item $\pi_1^\LL (X_{\bar {s}},a)_K$\dotfill\pageref{indnot:ml}
\medskip\item $\ZZ_\ell [1]$\dotfill\pageref{indnot:mm}
\par
\end{theindex}
\end{filecontents*}
% Original input: drafts/SGA1_XIII_APPENDIX_I_A_English.texfrag
\begin{filecontents*}[overwrite,noheader]{sga1-cumulative-input-117.tex}
\section{Appendix I: Variations on Abhyankar's lemma}
\label{XIII.5}
%
\index{Abhyankar's lemma|hyperpage}%
%
This appendix contains several variants of Abhyankar's lemma.

\begin{proposition}
\label{XIII.5.1} Let \(X=\Spec A\) be a regular local scheme, and let
\(D=\sum_{1\leq i\leq r}\divisor f_i\) be a normal crossings
divisor, where the \(f_i\) are elements of the maximal ideal of~\(A\)
belonging to a regular system of parameters. Let \(n_i\),
\(1\leq i\leq r\), be integers \(>0\), and put
$$
X'=X[T_1,\dots,T_r]/(T_1^{n_1}-f_1,\dots,T_r^{n_r}-f_r),
$$
Put \(U=X-\Supp D\) and \(U'=U\times_XX'\). Then \(X'\) is regular, and \(U'\) is the
complement in~\(X'\) of the normal crossings divisor
\(\sum_{1\leq i\leq r}\divisor T_i\). If the integers \(n_i\) are
prime
% page: 428
\marginpar{428}
%
to the residue characteristic \(p\) of~\(X\), then \(U'\) is a
connected étale covering of~\(U\), tamely ramified relative to~\(D\)
\textup{(\Ref{XIII.2.3}~c))}.
\end{proposition}

Indeed, \(X'\) is the spectrum of a local ring \(A'\) whose maximal
ideal is generated by \(T_1,\dots,T_r\).
\ifthenelse{\boolean{orig}}
{\ignorespaces}
{(We may suppose that \(r=\dim(A)\).)}
% correction: one may suppose that r=dim(A)
Since \(A'\) is finite and flat over~\(A\), and hence has dimension
\(r\), \(A'\) is regular (EGA~\(0_{\textup{IV}}\)~17.1.1), and the
\(T_i\) form a regular system of parameters of~\(A'\). Suppose that
the \(n_i\) are prime to \(p\). Since all the \(f_i\) are invertible
on~\(U\), the fact that \(U'\) is étale over~\(U\) follows from
\ifthenelse{\boolean{orig}}{I~}{}\Ref{I.7.4}. Moreover, \(U'\) is tamely ramified relative to~\(D\).
Indeed, let \(x_i\) be the generic point of~\(V(f_i)\), let
\(\bar R\) be the strict localization of
\(R=\cal{O}_{X,x_i}\), and let \(\bar K\) be the fraction field
of~\(\bar R\). Then the \(\bar K\)-algebra representing
\(U'|\bar K\) is obtained from the field
\(\bar K[T_i]/(T_i^{n_i}-f_i)\) by an unramified extension; it is
therefore tamely ramified over~\(\bar R\).

\begin{proposition}[Absolute Abhyankar lemma]
\label{XIII.5.2}
%
\index{absolute Abhyankar lemma|hyperpage}%
%
Let \(X\) be a regular local scheme, let
$$
D=\sum_{1\leq i\leq r}\divisor f_i
$$
be a normal crossings divisor as in~\Ref{XIII.5.1}, put
\(Y=\Supp D\), \(U=X-Y\), and let \(V\) be an étale covering of~\(U\)
tamely ramified relative to~\(D\). If \(x_i\) is the generic point
of the closed subset \(V(f_i)\), then \(\cal{O}_{X,x_i}\) is a
discrete valuation ring with fraction field \(K_i\), and
\(V|K_i=\Spec(\prod_{j\in J_i}L_j)\), where the \(L_j\) are finite
separable extensions of~\(K_i\). Let \(n_j\) be the order of the
inertia group of a Galois extension generated by \(L_j\), and let
\(n_i\) be the least common multiple of the \(n_j\) as \(j\) ranges
over~\(J_i\). If
$$
X'=X[T_1,\dots,T_r]/(T_1^{n_1}-f_1,\dots,T_r^{n_r}-f_r),
$$
\(U'=U_{(X')}\), \(V'=V_{(X')}\), etc., then the étale covering
\(V'\) of~\(U'\) extends, uniquely up to unique isomorphism, to an
étale covering of~\(X'\), and the \(n_i\) are prime to the residue
characteristic \(p\) of~\(X\).
\end{proposition}

Uniqueness follows from the fact that \(X'\) is normal
\eqref{XIII.5.1}. Indeed,
% page: 429
\marginpar{429}
%
an étale covering of~\(X'\) extending \(V'\) is isomorphic to the
normalization of~\(X'\) in the fiber of~\(V'\) at the generic point
of~\(X'\) \eqref{I.10.2}. If \(\bar x'\) is a geometric point
of~\(Y'\), denote by \(\bar X'\) the strict localization of~\(X'\)
at~\(\bar x'\), and put \(\bar V'=V'_{(\bar X')}\), etc. By descent,
in view of uniqueness, it is enough to show that for every geometric
point \(\bar x'\) of~\(Y'\), the étale covering \(\bar V'\)
of~\(\bar U'\) extends to~\(\bar X'\). Since an étale covering of an
open subset of the regular scheme~\(X'\) containing all points
\(x'\) such that \(\dim\cal{O}_{X',x'}\leq1\) extends to all
of~\(X'\) (SGA~2 XIV~1.11), we may even restrict ourselves to the
points \(\bar x'\) lying over a maximal point of~\(Y'\). At such a
point \(\bar x'\), the fact that \(\bar V'\) extends to an étale
covering of~\(\bar X'\) follows from~\Ref{X.3.6}.

Let us show that the \(n_i\) are prime to \(p\). Otherwise, for
example, we would have \(p|n_1\). After replacing \(X\) by
\ifthenelse{\boolean{orig}}
{\(X[T_1,\dots,T_r]/(T_1^{n_1/p}-f_1,T_2^{n_2}-f_2,\dots,T_r^{n_r}-f_r)\),}
{$$X[T_1,\dots,T_r]/(T_1^{n_1/p}-f_1,T_2^{n_2}-f_2,\dots,T_r^{n_r}-f_r),$$}
we reduce to the case
\(X'=X[T_1]/(T_1^p-f_1)\). It is enough to show that \(V\) extends
to an étale covering of~\(X\), since we would then have \(n_1=1\),
contrary to the hypothesis. For this, we may suppose that \(X\) is
strictly local. Let \(Z\) be the closed subscheme of~\(X\) defined
by \(p=0\), and let \(Z_1=Z\cap X_{f_1}\); \(Z_1\) is a nonempty
open subset of~\(Z\). By what precedes, the étale covering \(V'\)
of~\(U'\) extends to an étale covering \(W'\) of~\(X'\). Let
\(W''_1\) and \(W''_2\) be the inverse images of~\(W'\) under the
two projections \(X''=X'\times_XX'\rightrightarrows X'\). We show
that the descent isomorphism
\(u\colon W''_1|U''\to W''_2|U''\) extends to an \(X''\)-morphism
\(W''_1\to W''_2\), which will necessarily be descent data on~\(W'\)
relative to \(X'\to X\). Let \(Z''\) (\resp \(Z''_1\)) be the
inverse image of~\(Z\) (\resp \(Z_1\)) in~\(X''\). Since the
morphism \(Z''\to Z\) is radicial, there is an isomorphism
\(v\colon W''_1|Z''\to W''_2|Z''\) extending the isomorphism
\(u|Z''_1\). But since \(X\) is Henselian, there is a bijection
$$
\Hom_{X''}(W''_1,W''_2)\simeq\Hom_{Z''}(W''_1|Z'',W_2''|Z''),
$$
and hence a morphism \(w\colon W''_1\to W''_2\) lifting \(v\). The
open-and-closed
% page: 430
\marginpar{430}
%
subscheme of~\(X''\) over which \(u\) and \(w\) coincide contains
\(Z''_1\), and therefore equals~\(X''\). Thus \(V\) extends to~\(X\).

\setcounter{subsection}{3} \setcounter{subsubsection}{-1}

\subsubsection{}
\label{XIII.5.3.0} Retain the hypotheses and notation
of~\Ref{XIII.5.2}, and suppose in addition that
\ifthenelse{\boolean{orig}}
{\(S\)}
{\(X\)}
% correction: X instead of S
is strictly local. It follows from \loccit that every connected
étale covering of~\(U\) tamely ramified relative to~\(D\) is a
quotient of a tamely ramified covering of the form
$$
U'=U[T_1,\dots,T_r]/(T_1^{n_1}-f_1,\dots,T_r^{n_r}-f_r),
$$
where the \(n_i\) are integers prime to \(p\). Let \(\bbmu_n\) be
the group of \(n\)-th roots of unity of~\(U\). The group of
\(U\)-automorphisms of~\(U'\) is precisely
\(\bbmu_{n_1}\times\dots\times\bbmu_{n_r}\), an \(n_i\)-th root of
unity \(\xi_i\) acting on~\(U'\) by sending \(T_i\) to \(\xi_iT_i\).
We therefore have the following statement:

\setcounter{subsection}{2}

\begin{corollaire}
\label{XIII.5.3} Let \(X\) be a regular strictly local scheme of
residue characteristic \(p\geq0\), let
\(D=\sum_{1\leq i\leq r}\divisor f_i\) be a normal crossings divisor
on~\(X\), and put \(U=X-\Supp D\). Let
$$
\widetilde U=\varprojlim_{(n_i)}
U[T_1,\dots,T_r]/(T_1^{n_1}-f_1,\dots,T_r^{n_r}-f_r),
$$
where the projective limit is taken over the filtered ordered set
(under divisibility) of families of integers \(n_i>0\) prime to
\(p\). Then \(\widetilde U\) is a universal tamely ramified covering
of~\(U\). Consequently, the tamely ramified fundamental group of~\(U\)
is
$$
\pi_1^{\tame}(U)\simeq \prod_{\ell\neq p} \ZZ_\ell [1]^r
\qquad \text{(canonical isomorphism)}
$$
where
\(\ZZ_\ell[1]=\varprojlim_{n>0}\bbmu_{\ell^n}\). The group
\(\pi_1^{\tame}(U)\) is noncanonically isomorphic to
\(\prod_{\ell\neq p}\ZZ_\ell^r\).
\label{indnot:mm}\oldindexnot{\(\ZZ_\ell[1]\)|hyperpage}%
%
\end{corollaire}

\begin{proposition}
\label{XIII.5.4} Let \(f\colon X\to S\) be a morphism of schemes,
and let \(D=\sum_{1\leq i\leq r}\divisor f_i\) be a divisor on~\(X\)
with normal crossings relative to~\(S\) \eqref{XIII.2.1}, where,
% page: 431
\marginpar{431}
%
for every point \(x\) of~\(Y=\Supp D\), if
\(I(x)\subset[1,r]\) is the set of \(i\) such that \(f_i(x)=0\), the
subscheme \(V((f_i)_{i\in I(x)})\) is smooth over~\(S\), of
codimension
\ifthenelse{\boolean{orig}}
{\(\mathrm{card}.\,I(x)\)}
{\(\card I(x)\)}
in \(X\). Let \(U=X-Y\). Let \(x\) be a point of~\(Y\), put
\(X_1=\Spec\cal{O}_{X,x}\), \(U_1=U\times_XX_1\), let \(n_i\),
\(i\in I(x)\), be integers \(>0\), and put
$$
X'=X[T_i]_{i\in I(x)}/(T_i^{n_i}-f_i).
$$
Then, if \(x'\) is the point of~\(X'\) over~\(x\), \(X'\) is smooth
over~\(S\) at~\(x'\). If the integers \(n_i\) are prime to the
characteristic \(p\) of~\(\kres(x)\), then
\(U'_1=U_1\times_XX'\) is a connected étale covering of~\(U_1\),
tamely ramified on~\(X_1\) relative to~\(S\)
\eqref{XIII.2.1.1}.
\end{proposition}

If \(s=f(x)\), the geometric fiber \(X'_{\bar s}\) is regular at the
point \(x'\) \eqref{XIII.5.1}. Since \(X'\) is flat over~\(S\) in a
neighborhood of the inverse image of~\(Y\) in~\(X'\), this proves that \(X'\) is smooth over~\(S\)
at~\(x'\) (EGA~IV~12.1.6). If the integers \(n_i\) are prime to
\(p\), \(U'_1\) is an étale covering of~\(U_1\)
\eqref{I.7.4}; it is tamely ramified on~\(X_1\) relative to~\(S\)
because this is so on the geometric fibers at every point of~\(S\)
\eqref{XIII.5.1}. Finally, the connectedness of~\(U'_1\) follows
from SGA~4 XVI~3.2.

\end{filecontents*}
% Original input: drafts/SGA1_XIII_APPENDIX_I_B_English.texfrag
\begin{filecontents*}[overwrite,noheader]{sga1-cumulative-input-118.tex}
\begin{proposition}[Relative Abhyankar lemma]
\label{XIII.5.5}
\index{relative Abhyankar lemma|hyperpage}%
Let $X$ be an $S$-scheme and let $D$ be a normal crossings divisor
relative to~$S$, as in~\eqref{XIII.5.4}. Let $Y=\Supp D$,
$U=X- Y$, let $x$ be a point of~$Y$, let $X_{1}$ be the strict
localization of~$X$ at a geometric point above~$x$, let
$U_{1}=U\times_{X}X_{1}$, and let $V_{1}$ be an \'etale covering of
$U_{1}$. Suppose that, for every maximal point $s$ of~$S$,
$V_{1\overline{s}}$ is tamely ramified on~$X_{1\overline{s}}$ relative to
$\overline{s}$. Then there exist integers $n_{i}$ prime to the
characteristic $p$ of~$\kres(x)$, with $i\in I(x)$, such that, upon
putting
$$
X'_{1}=X_{1}[T_{i}]_{i\in I(x)}\big/ (T_{i}^{n_{i}}-f_{i}) \quoi,
$$
$U'_{1}=U_{1}\times_{X_{1}}X'_{1}$, etc., the \'etale covering
$V'_{1}$ of~$U'_{1}$ extends uniquely, up to unique isomorphism, to
an \'etale covering of~$X'_{1}$. In particular,
% page: 432
\marginpar{432}
%
$V_{1}$ is tamely ramified on~$X_{1}$ relative to~$S$.
\end{proposition}

We may assume that $S$ is local noetherian with closed point~$f(x)$.
For every maximal point $s$ of~$S$ and every $i\in I(x)$, let $x_{i}$
be the generic point of the closed subset $V(f_{i})$ of the fiber
$X_{1\overline{s}}$. The local ring $(\cal{O}_{X_{1},x_{i}})_{\red}$
is a discrete valuation ring with fraction field $K_{i}$, and we have
$V_{1}|K_{i}=\Spec(\prod_{j\in J(x_{i})}L_{j})$,
where $L_{j}$ is a finite separable extension of~$K_i$; let $n_{j}$
denote the order of the inertia group of a Galois extension generated
by $L_{j}$, and let $n_{i}$ be the l.c.m. of the $n_{j}$ as $s$ ranges
over the maximal points of~$S$ and $j\in J(x_{i})$.

With the $n_{i}$ so chosen, we shall show that $V'_{1}$ extends
uniquely to an \'etale covering of $X'_{1}$. Uniqueness follows from
the fact that, since $X'$ is smooth over~$S$ at the points of~$Y'$,
we have $\prof\et_{Y'_{1}}(X'_{1})\geqslant 2$
(SGA~4 XVI~3.2 or SGA~2 XIV~1.19). Let $x'_{1}$ be a point of
$Y'_{1}$, let $\overline{x}'_{1}$ be a geometric point above
$x'_{1}$, and denote by $\overline{X'_{1}}$ the strict localization
of $X'_{1}$ at $\overline{x}'_{1}$,
$\overline{U'_{1}}=U'_{1(\overline{X'_{1}})}$,~etc. By descent, in
view of uniqueness, it suffices to show that $\overline{V}'_{1}$
extends to $\overline{X'_{1}}$. Moreover, it is enough to take for
$x'_{1}$ only the maximal points of $Y'_{1}$; indeed, we shall then
have an extension of~$V'_{1}$ over an open subset $W'_{1}$ of
$X'_{1}$ containing the maximal points of~$Y'_{1}$. Now, if
$Z'_{1}=X'_{1}- W'_{1}$, then
$\codim(Z'_{1s},X'_{1s})\geqslant 2$ if $s$ is a maximal point of
$S$, and $\codim(Z'_{1s},X'_{1s})\geqslant 1$, $\prof\et_{s}(S)\geqslant
1$ if $s$ is a point of~$S$ that is not maximal; the fact that
$V'_{1}$ extends to all of $X'_{1}$ then follows from
(SGA~2 XIV~1.20). But at a geometric point
$\overline{x}'_{1}$ above a maximal point of~$Y'_{1}$,
$\overline{X'_{1}}_{\red}$ is the spectrum of a discrete valuation ring, and the
fact that $\overline{V'_{1}}$ extends to $\overline{X'_{1}}$ then
follows from~\ifthenelse{\boolean{orig}}{X~}{}\Ref{X.3.6}.

Let us show that the $n_{i}$ are prime to $p$. Indeed, if this were
not so, there would be an index $i_{0}\in I(x)$ such that $p$
divides $n_{i_{0}}$. After replacing $X$ by
$X_{1}[T_{i_{0}},T_{i}]_{i\in I(x),\,i\ne i_{0}}\big/
(T_{i_{0}}^{n_{i_{0}}/p}-f_{i_{0}},(T_{i}^{n_{i}}-f_{i})_{i\in I(x),\,i\ne i_{0}})$, we are reduced
to the case $X'_{1}=X_{1}[T]/(T^{p}-f_{i_{0}})$. By what precedes, the
% page: 433
\marginpar{433}
%
\'etale covering $V'_{1}$ of~$U'_{1}$ extends to an \'etale covering
$E'_{1}$ of~$X'_{1}$. Let $\eta$ be the closed point of $S$; since
the morphism $X'_{1\eta}\to X_{1\eta}$ is radicial, $V_{1\eta}$
extends to an \'etale covering $E_{1\eta}$ of $X_{1\eta}$. It then
follows, as in~\Ref{XIII.5.2}, that $E'_{1}$ is endowed with descent
data relative to the morphism $X'_{1}\to X_{1}$ extending the
natural descent data on~$E'_{1}|U'_{1}$. It follows that $V_{1}$
extends to $X_{1}$; but this implies $n_{i_{0}}=1$, contrary to the
hypothesis $n_{i_{0}}=p$.

\begin{corollaire}
\label{XIII.5.6}
Let $X$ be an $S$-scheme, and let $D=\sum_{1\leqslant i\leqslant
r}\divisor f_{i}$ be a normal crossings divisor relative to~$S$, as
in~\Ref{XIII.5.4}. Let $\overline{x}$ be a geometric point of~$X$,
let $\overline{X}$ be the strict localization of $X$ at
$\overline{x}$, let $\overline{Y}=Y_{(\overline{X})}$,
$\overline{U}=\overline{X}-\overline{Y}$, and let
$$
\widetilde{U}= \varprojlim_{(n_{i})}
\overline{U}[T_{i}]_{i\in I(x)}\big/ (T_{i}^{n_{i}}-f_{i}) \quoi,
$$
where the projective limit is taken over the filtered set of families
of integers $n_{i}>0$ prime to the characteristic $p$ of~$\kres(x)$.
Then $\widetilde{U}$ is a universal tamely ramified covering of
$\overline{U}$ relative to~$S$. Consequently, the tamely ramified
fundamental group of~$\overline{U}$ is
$$
\pi_{1}^{\tame}(\overline{U})\simeq \prod_{\ell\neq
p}\ZZ_{\ell}[1]^{I(x)} \qquad \textrm{(canonical isomorphism)\quoi.}
$$
The group $\pi_{1}^{\tame}(\overline{U})$ is noncanonically
isomorphic to $\prod_{\ell\neq p}\ZZ_{\ell}^{I(x)}$.
\end{corollaire}

\begin{subremarque}
\label{XIII.5.6.1}
Let $X$ be an $S$-scheme, let $D=\sum_{1\leqslant i\leqslant
r}\divisor f_{i}$ be a normal crossings divisor relative to~$S$, as
in~\Ref{XIII.5.4}, and let $U=X-\Supp D$.
For every subset $I\subset[1,r]$, let
$$
X_{I}= \Big(\tbigcap_{i\in I}V(f_{i})\Big) \cap \Big(\tbigcap_{i\in\complement
I}X_{f_{i}}\Big)\quoi .
$$
Let $p$ be a prime integer or zero, and let $Z$ be a subset of
$X_{I}$ all of whose points have characteristic~$p$. Let
$$
\widetilde{U}_{I}= \varprojlim_{(n_{i})}
\ifthenelse{\boolean{orig}}
{U[T_{i}]}
{U[T_{i}]_{i\in I}}
% correction: it should be
% U[T_{i}]_{i\in I}\big/ (T_{i}^{n_{i}}-f_{i}) \quoi,
\big/ (T_{i}^{n_{i}}-f_{i}) \quoi,
$$
where
% page: 434
\marginpar{434}
%
the projective limit is taken over the filtered set of families of
integers $n_{i}>0$ prime to $p$. Then, for every geometric point
$\overline{x}$ of~$Z$, the pullback of $\widetilde{U}_{I}$ to
$\overline{U}$ identifies with the universal tamely ramified covering
of~$\overline{U}$.
\end{subremarque}

\begin{corollaire}
\label{XIII.5.7}
Use the notation of~\Ref{XIII.5.6}. Let $\overline{S}$ be the strict
localization of~$S$ at~$\overline{x}$, and let
$$
\overline{g}\colon\overline{U}\to\overline{S} \quad \text{and} \quad
\tilde{g}\colon\widetilde{U}\to\overline{S}
$$
be the canonical morphisms. Then the morphisms $\overline{g}$ and
$\tilde g$ are $0$-acyclic \textup{(SGA~4 XV~1.3)}. Let $G$ be a
constructible sheaf of groups on~$\overline{S}$, let
$F=\overline{g}^*G$, and let $P$ be a torsor under $F$. Then $P$ is
tamely ramified on~$\overline{X}$ relative to $\overline{S}$ if and
only if its pullback $\widetilde{P}$ to~$\widetilde{U}$ is trivial.
\end{corollaire}

Indeed, for every scheme
$\overline{X'}=\overline{X}[T_{i}]_{i\in I(x)}\big/
(T_{i}^{n_{i}}-f_{i})$, where the $n_{i}$ are integers~$>0$ prime
to $p$, the morphism
$\overline{f}'\colon\overline{X'}\to\overline{S}$ is $0$-acyclic.
The geometric fibers of~$\overline{f}'$ at the various points of
$\overline{S}$ are therefore connected, and indeed irreducible. The
same is consequently true of the geometric fibers of the morphisms
$\overline{g}'\colon\overline{U'}\to\overline{S}$, which proves that
the $\overline{g}'$, and hence also $\tilde g$, are
$0$-acyclic (SGA~4 XV~1.16).

It is clear that a torsor $P$ on~$\overline{U}$ under $F$ whose
pullback to~$\widetilde{U}$ is trivial is tamely ramified on
$\overline{X}$ relative to~$\overline{S}$. Conversely, suppose that
$P$ is tamely ramified on~$\overline{X}$ relative to
$\overline{S}$; we shall show that its pullback to~$\widetilde{U}$
is trivial.

It follows from (SGA~4 IX~2.14~(ii)) that there exist a finite
morphism $n\colon S_{1}\to\overline{S}$, a constant sheaf of groups
$C$ on~$S_{1}$, and a monomorphism $G\to n_*C$. Consider the following
commutative diagram formed by Cartesian squares:
$$
\xymatrix{ \widetilde{U}_{1} \ar[r] \ar[d]_{r} &U_{1} \ar[r]^{g_{1}}
\ar[d]_{q} &S_{1} \ar[d]_{n} \cr \widetilde{U} \ar[r] &\overline{U}
\ar[r]^{\overline{g}} &\overline{S} \rlap{\quoi.}\cr }
$$
Let
% page: 435
\marginpar{435}
%
$C_{1}$ (\resp $\widetilde{C}_{1}$) be the pullback of~$C$ to~$U_{1}$
(\resp to~$\widetilde{U}_{1}$). We have a commutative diagram in
which $i$ and $j$ are isomorphisms (SGA~4 VIII~5.8):
\begin{equation*}
\label{eq:XIII.5.7.*}
\tag{$*$}
\begin{array}{c}
\xymatrix{ \H^{1}(\overline{U},q_*C_{1}) \ar[r]^{i}
\ar[d] &\H^{1}(U_{1},C_{1}) \ar[d] \cr
\H^{1}(\widetilde{U},r_*\widetilde{C}_{1}) \ar[r]^{j}
&\H^{1}(\widetilde{U}_{1},\widetilde{C}_{1}) \rlap{\quoi.}\cr }
\end{array}
\end{equation*}
Let $Q$ be the torsor under $q_*C_{1}$ obtained from~$P$ by
extension of structure group along $F\to q_*C_{1}$.
By~\Ref{XIII.2.1.4}, $Q$ is tamely ramified on~$\overline{X}$
relative to $\overline{S}$. The torsor $Q$ corresponds, via $i$, to
a torsor $Q_{1}$ under $C_{1}$, and it is clear that $Q_{1}$ is
tamely ramified on~$X_{1}=\overline{X}\times_{\overline{S}}S_{1}$
relative to~$S_{1}$. It therefore follows from~\Ref{XIII.5.6} that
the pullback $\widetilde{Q}_{1}$ of~$Q_{1}$ to~$\widetilde{U}_{1}$
is trivial, and diagram~\eqref{eq:XIII.5.7.*} then shows that the
pullback $\widetilde{Q}$ of~$Q$ to~$\widetilde{U}$ is trivial.

Consider the following commutative diagram, whose second row is exact
(SGA~4 XII~3.1):
\begin{equation*}
\label{eq:XIII.5.7.**}
\tag{$**$}
\begin{array}{c}
\xymatrix{ \H^{0}(\overline{S},n_*C/G) \ar[r]
\ar[d]_{k} &\H^{1}(\overline{S},G)\rlap{$~=~1$} \ar[d] \cr
\H^{0}(\widetilde{U},r_*\widetilde{C}_{1}/\widetilde{F}) \ar[r]
&\H^{1}(\widetilde{U},\widetilde{F}) \ar[r]
&\H^{1}(\widetilde{U},r_*\widetilde{C}_{1}) \rlap{\quoi.} \cr }
\end{array}
\end{equation*}
Since the morphism $\widetilde{U}\to\overline{S}$ is $0$-acyclic,
$k$ is an isomorphism. It then follows from~\eqref{eq:XIII.5.7.**}
that $\widetilde{P}$ is trivial.
\end{filecontents*}
% Original input: drafts/SGA1_XIII_OPENING_AND_0_English.texfrag
\begin{filecontents*}[overwrite,noheader]{sga1-cumulative-input-119.tex}
\chapterspace{-4}
\chapter[Cohomological properness for sheaves of sets...]{Cohomological properness\\ for sheaves of sets and\\ sheaves of noncommutative groups}
\label{XIII}
% page: 344
\marginpar{344}
%

\begin{center}
{by Mme M.\ \textsc{Raynaud}\footnote{Based on unpublished notes by
A\ptbl Grothendieck.}}
\end{center}
\vspace*{1cm}

The purpose of this exposé is to use étale cohomology to generalize
certain results of \Ref{IX} and \Ref{X}. It also shows how the results
of~SGA~5~II that still make sense for sheaves of noncommutative groups
can be extended to such sheaves. We assume familiarity with the
notions of étale cohomology presented in~SGA~4.

The main result \eqref{XIII.2.4} gives an important example of a
nonproper morphism \(f\colon U\to S\) that is
\emph{cohomologically proper in dimension~\(\leq1\)}, that is, such
that for certain sheaves of groups \(F\) on~\(U\) (for the étale
topology), the formation of~\(f_*F\) and \(\R^1f_*F\) commutes with
every base change \(S'\to S\). This property is indeed satisfied by
the open subset \(U\) of a scheme \(X\) proper over~\(S\) that is the
complement of a divisor \(D\) with normal crossings relative to~\(S\),
at least if \(F\) is required to be finite constant, of order prime to
the residue characteristics of~\(S\). If \(F\) is no longer assumed
to have order prime to the residue characteristics of~\(S\), there is
an analogous result with \(\R^1f_*F\) replaced by the subsheaf
\(\R^1_{\tame}f_*F\) obtained by considering only torsors under~\(F\)
that are \emph{tamely ramified on~\(X\) relative to~\(S\)}. In
particular, this makes it possible to show that the tamely ramified
fundamental group of a proper smooth algebraic curve over a separably
closed field, with finitely many closed points removed, is
topologically of finite type \eqref{XIII.2.12}.

The
% page: 345
\marginpar{345}
%
\No \Ref{XIII.4} is devoted to the homotopy exact sequence and the
Künneth formula.

Finally, an appendix gives useful variants of Abhyankar's lemma
proved in X.\Ref{X.3.6}.

\setcounter{section}{-1}

\section{Review of the theory of stacks}
\label{XIII.0}

We shall use below the theory of stacks presented in \cite{XIII.1}
and \cite{XIII.2}. We restrict ourselves to the case of the étale site
of a scheme. Given a scheme \(X\), denote by \(X_\et\) the étale site
of~\(X\). Recall that a
\ifthenelse{\boolean{orig}}
{stack}
{stack \(F\)}
%
\index{stack|hyperpage}%
%
on~\(X\) is a category fibered over~\(X_\et\) such that, for every
scheme \(X'\) étale over~\(X\) and every pair
\ifthenelse{\boolean{orig}}
{of object}
{of objects}
% correction: objects instead of object
\(x,y\) of the fiber \(F_{X'}\), the presheaf
\(\SheafHom_{X'}(x,y)\) is a sheaf, and such that, for every
surjective étale morphism \(X''\to X'\), every object of~\(F_{X''}\)
equipped with descent data relative to \(X''\to X'\) is the inverse
image of an object of~\(F_{X'}\).

We denote by \(F(X')\) the category of cartesian sections of
\(F/X'\). More generally, if \(\Sch_X\) is the category of schemes
over~\(X\), endowed with the étale topology, the stack \(F\) can be
extended to a stack \(\cal{F}\) on~\(\Sch_X\), and for every morphism
\(f\colon X'\to X\), we again denote by \(F(X')\) the category of
cartesian sections of this stack \(\cal{F}\) over~\(X'\).

A gerbe
%
\index{gerbe|hyperpage}%
%
is a stack such that, for every scheme \(X'\) étale over~\(X\) and
every pair of objects \(x,y\) of~\(F_{X'}\), every morphism from~\(x\)
to~\(y\) is an isomorphism, \(x\) and \(y\) are locally isomorphic,
and the collection of objects \(X'\) of~\(X_\et\) for which
\(F_{X'}\) is nonempty forms a covering of~\(X_\et\). For example,
the stack of torsors under a sheaf of groups is a gerbe that,
moreover, has a cartesian section. Conversely, a gerbe that has a
section, \ie for which there exists an object \(x\) of~\(F_X\), is
equivalent to the stack of torsors under the sheaf
% page: 346
\marginpar{346}
%
of groups \(\SheafAut_X(x)\).

There is an evident notion of a subgerbe and of a maximal subgerbe of
a stack~\(F\). Given a cartesian section \(x\) of \(F(X)\), there is a
unique maximal subgerbe \(G_x\) of~\(F\) through which \(x\) factors.
We call \(G_x\) the subgerbe generated by~\(x\); it is, by definition,
a trivial gerbe.
%
\index{trivial gerbe|hyperpage}%
%
The presheaf \(SF\)
\label{indnot:mb}\oldindexnot{\(SF\) or \(S(F)\)|hyperpage}%
%
defined by
$$
SF(X') = \{ \text{maximal subgerbes of }F|X' \}
$$
is a sheaf, called the sheaf of maximal subgerbes of~\(F\).
%
\index{sheaf of maximal subgerbes|hyperpage}%
%
Let \(O\) be the presheaf defined by
$$
O(X') = \{ \text{classes of objects of }F_{X'}\text{ modulo
isomorphism} \} \text{.}
$$
By associating with every object \(x\) of~\(F_{X'}\) the maximal
subgerbe of~\(F|X'\) generated by~\(x\), we obtain a morphism
$$
O \to SF \text{;}
$$
by \cite[III 2.1.4]{XIII.2}, this morphism makes~\(SF\) a sheaf
associated with~\(O\).

A stack \(F\) is said to be \emph{constructible}
%
\index{constructible (stack)|hyperpage}%
%
(\emph{\resp ind-\(\LL\)-finite},
%
\index{stack ind-finite (\resp ind-\(l\)-finite)|hyperpage}%
%
where \(\LL\) is a set of prime numbers) if, for every scheme \(X'\)
étale over~\(X\) and every object \(x\) of~\(F_{X'}\), the sheaf
\(\SheafAut_{X'}(x)\) has the corresponding property
\cite[VII 2.2.1]{XIII.2}. A stack is said to be
\(1\)-constructible
%
\index{stack \(1\)-constructible (\resp constructible)|hyperpage}%
%
if it is constructible and its sheaf of maximal subgerbes is
constructible.

\end{filecontents*}
% Original input: drafts/SGA1_XII_2_English.texfrag
\begin{filecontents*}[overwrite,noheader]{sga1-cumulative-input-120.tex}
\section[Comparison of the properties of a scheme]{Comparison of the properties of a scheme and of its
\hbox{associated} analytic space}
\label{XII.2}

\begin{proposition}
\label{XII.2.1}
Let $X$ be a $\CC$-scheme locally of finite type, let $X^{\an}$ be
the associated analytic space, and let $n$ be an integer. Consider the
property $P$ of being
\begin{itemize}
\item[(i)] nonempty
\item[(i')] discrete
\item[(ii)] Cohen--Macaulay
\item[(iii)] $(\mathrm{S}_n)$
\item[(iv)] regular
\item[(v)] $(\mathrm{R}_n)$
\item[(vi)] normal
\item[(vii)] reduced
\item[(viii)] of dimension $n$.
\end{itemize}
Then $X$ has property $P$ if and only if
$X^{\an}$ does.
\end{proposition}

Let $\varphi\colon X^{\an}\to X$ be the canonical morphism.
(i)~follows from the equality $|X^{\an}|=X(\CC )$
\ifthenelse{\boolean{orig}}{\eqref{XIII.1.1}}{\eqref{XII.1.1}} and from the fact that $X$ is Jacobson (EGA~IV~10.4.8). To say that~$X$
(\resp $X^{\an}$) is discrete is equivalent to saying that $\dim
X=0$ (\resp $\dim X^{\an}=0$, by \cite[\no19~\S 4~cor\ptbl 6]{XII.4}); (i')
therefore follows from~(viii).

Let $P$ be one of the properties (ii) through (vii). The scheme $X$
has property $P$ if and only if $P$
holds at every closed point of~$X$; indeed, since $X$
is excellent (EGA~IV~7.8.6~(iii)), the set of points
% page: 317
\marginpar{317}
%
at which $X$ has property $P$ is open (\loccit), and if this open set
contains all the closed points, it is all of~$X$.
Thus, to say that $X$ (\resp $X^{\an}$) has property $P$
is equivalent to saying that, for every point $x$ of~$X^{\an}$,
the local ring $\cal{O}_{X,\varphi (x)}$ (\resp $\cal{O}_{X^{\an},x}$)
has property $P$. Since whether an excellent local ring
has property $P$ can be checked after passing to its
completion, the proposition follows from the isomorphisms
$\widehat{\cal{O}}_{X,\varphi(x)}\isomto\widehat{\cal{O}}_{X^{\an},x}$
in cases (ii) through (vii). The same is true in case
(viii), in view of the equalities
$$
\dim X=\sup_x~\dim \cal{O}_{X,\varphi (x)}\qquad \dim
X^{\an}=\sup_x~\dim \cal{O}_{X^{\an},x}
$$
where $x\in X^{\an}$. This completes the proof.

\begin{proposition}
\label{XII.2.2}
Let $X$ be a $\CC$-scheme locally of finite type,
let $\varphi\colon X^{\an}\to X$ be the canonical morphism, and let $T$ be a
locally constructible subset of~$X$. Then
$$
\varphi^{-1}(\overline{T})=\overline{\varphi^{-1}(T)}.
$$
\end{proposition}

We may assume that $T$ is a dense open subset of~$X$. Let $H$ be the
reduced closed subscheme of~$X$ whose underlying space is
$X-T$; the associated analytic space $H^{\an}$ is a closed analytic subspace
of~$X^{\an}$ whose underlying space is
$X^{\an}-\varphi^{-1}(T)$. We must show that every point $x$ of
$H^{\an}$ belongs to $\overline{\varphi^{-1}(T)}$. At such a
point $x$, the germ $(X^{\an},x)$ contains the
subgerm $(H^{\an},x)$, and the latter is defined by a non-nilpotent ideal
of~$\cal{O}_{X^{\an},x}$. It then follows from the
Nullstellensatz \cite[\no19~\S 4~cor\ptbl 3]{XII.4} that every open neighborhood of~$x$
contains points of~$X^{\an}$ which do not belong to
$H^{\an}$, proving that
$x\in\overline{\varphi^{-1}(T)}$.

\begin{corollaire}
\label{XII.2.3}
Let
% page: 318
\marginpar{318}
%
$X$ be a $\CC$-scheme locally of finite type, let $\varphi
\colon X^\an \to X$ be the canonical morphism, and let $T$ be a locally
constructible subset of~$X$. The subset $T$ is open (\resp
closed, \resp dense) if and only if
$\varphi^{-1} (T)$ has the corresponding property.
\end{corollaire}

The corollary follows from~\Ref{XII.2.2} and from the fact that, since $X$
is a Jacobson scheme (EGA~IV~10.4.8), two
locally constructible subsets of~$X$ having the same trace on the
very dense subset $X (\CC )$ are equal.

\begin{proposition}
\label{XII.2.4} Let $X$ be a $\CC$-scheme locally of finite
type. Then $X$ is connected (\resp irreducible) if and
only if the same is true of~$X^\an$.
\end{proposition}

Suppose that $X^\an $ is connected (\resp irreducible). The image $X (\CC)
$ of~$X^\an$ in $X$ is then connected (\resp irreducible). It
follows that $X$ is connected (\resp irreducible), since the
closed subsets of~$X$ and of $X (\CC)$ correspond bijectively
(EGA~IV~10.1.2).

Conversely, suppose that $X$ is connected (\resp irreducible), and
let us show that the same is true of~$X^\an$. We may reduce to the case
where $X$ is irreducible. Indeed, suppose that $X$ is connected.
Given a point $x$ of~$X$, the set of points $y \in X$
for which there exists a finite sequence of irreducible closed
subschemes $X_1, \dots, X_n$ of~$X$, with $x \in X_1$, $y \in
X_n$, $X_i \cap X_{i+1} \neq \emptyset $ for $1 \leq i \leq n-1$,
is both open and closed, and hence is all of
$X$. For a sequence $X_1, \dots, X_n$ as above, we also have
$X_{i}^\an \cap X_{i+1}^\an \neq
\emptyset $ for $1 \leq i \leq n-1 $; thus, once it is known
that the $X_{i}^\an$ are connected, it follows that
$X^\an$ is connected as well.

We
% page: 319
\marginpar{319}
%
now suppose that $X$ is irreducible. We may further assume that
$X$ is affine. Indeed, if $(U_i)_{i \in I}$ is a covering of~$X$
by affine open subsets, any two of these open subsets have nonempty
intersection, and the same property therefore holds for the
covering $(U_{i}^\an)_{i \in I} $ of~$X^\an$; once it is known
that the $U_{i}^\an $ are irreducible, it follows that
$X^\an$ is irreducible as well.

We may further assume that $X$ is normal. Indeed, let $\tilde X$
be the normalization of~$X$; since the morphism $\tilde X \to X$ is
surjective, the same is true of~$\tilde X^\an \to X^\an$. This
shows that if $\tilde X^\an $ is irreducible, then
$X^\an$ is irreducible as well.

We now suppose that $X$ is affine and normal. Since the local rings
of~$X^\an$ are integral domains, saying that
$X^\an$ is irreducible is equivalent to saying that it is connected. Indeed, if
$\cal{F}$ is a closed analytic subset of~$X^\an$, the set
of points $x$ of~$X^\an$ at which $\codim_x (\cal{F}, X^\an ) =
0$ is a closed analytic subset of~$X^\an$ \cite[\no 20 A
Cor\ptbl 1]{XII.4} that is also open; if $X^\an$ is connected, this shows that
if $\cal{F} \neq X^\an$, then $\cal{F}$ is nowhere dense, and hence that $X^\an$
is irreducible. We are thus reduced to showing that
$X^\an$ is connected.

Let
$$
i \colon X \to P
$$
be a compactification of~$X$, where $P$ is a normal projective
$\CC$-scheme and $i$ is a dominant open immersion. It
then follows from \cite[\no 12 th\ptbl 1]{XII.10} that $P^\an$ is connected. Since
$X^\an$ is obtained by removing from $P^\an$ a nowhere dense closed
analytic subset, it follows from~\Ref{XII.2.5} below that $X^\an$
is connected as well.

\begin{lemme}
\label{XII.2.5}
Let
% page: 320
\marginpar{320}
%
$\cal{P} $ be a connected normal analytic space and $\cal{Y}$ a
nowhere dense closed analytic subset. Then $\cal{X} = \cal{P} - \cal{Y}$
is connected.
\end{lemme}

When $\cal{Y}$ has codimension $\geq 2$, the proposition
follows from \cite[\no 3 prop\ptbl 4]{XII.11}. In the general case, after
removing from $\cal{P}$ a closed analytic subset of codimension $\geq 2$,
we may assume that $\cal{P}$ and $\cal{Y}$
(regarded as a reduced analytic subspace of~$\cal{P}$) are regular. By the
implicit function theorem, every point $y$ of~$\cal{Y}$ has a
neighborhood $\cal{U}$ isomorphic to a ball in an affine
space $\cal{E}^{n}$, in such a way that $\cal{U} \cap \cal{Y}$ is defined by
the vanishing of a certain number of coordinate functions. This
shows that $\cal{U} - \cal{U} \cap \cal{Y}$ is connected, and therefore
so is~$\cal{X}$.

\begin{corollaire}
\label{XII.2.6} Let $X$ be a $\CC$-scheme locally of finite
type; the morphism
$$
\pi_0 (X^{\an}) \to \pi_0 (X)
$$
induced by the canonical morphism $X^\an \to X$ is bijective.
\end{corollaire}
\end{filecontents*}
% Original input: drafts/SGA1_XII_3_English.texfrag
\begin{filecontents*}[overwrite,noheader]{sga1-cumulative-input-121.tex}
\section{Comparison of properties of morphisms}
\label{XII.3}

\begin{proposition}
\label{XII.3.1} Let $f \colon X \to Y$ be a morphism of
$\CC$-schemes locally of finite type, and let $f^\an \colon X^\an \to
Y^\an$ be the morphism induced by~$f$ on the associated analytic
spaces. Let $P$ be the property of being
\begin{enumerate}
\item[(i)] flat
\item[(ii)] net (\ie unramified)
\item[(iii)] \'etale
\item[(iv)] smooth
\item[(v)] normal
\item[(vi)]
% page: 321
\marginpar{321}
%
reduced
\item[(vii)] injective
\item[(viii)] separated
\item[(ix)] an isomorphism
\item[(x)] a monomorphism
\item[(xi)] an open immersion.
\end{enumerate}
Then $f$ has property $P$ if and only if $f^\an$ does.
\end{proposition}

Denote by $\varphi \colon X^\an \to X $ and $\psi \colon Y^\an \to Y$ the
canonical morphisms. Let $x$ be a point of~$X^\an$, and put $y = f^\an
(x)$. The morphisms $\cal{O}_{Y^\an, y} \to {\mathcal O}_{X^\an, x}$
and $\cal{O}_{Y, \psi(y)} \to \cal{O}_{X, \varphi(x)} $ induced by
\ifthenelse{\boolean{orig}}{$f$ and $f^{\an}$}{$f^{\an}$ and $f$} yield the same morphism after passage to
completions \eqref{XII.1.1}. By \cite[ch\ptbl 3 \S 5 prop\ptbl 4]{XII.2}
(\resp EGA~IV 17.4.4), saying that $f^\an $
satisfies property (i) (\resp (ii)) is therefore equivalent to saying that $f$
satisfies (i) (\resp (ii)) at every closed point of~$X$. Since
the set of points of~$X$ at which (i) (\resp (ii))
holds is open (EGA~IV 11.1.1 and I 3.3), this
proves (i) and (ii), and hence also (iii).

\ifthenelse{\boolean{orig}}{}
{\enlargethispage{.5cm}}%
Let $P$ be property (iv) (\resp (v), \resp (vi)). In view
of~\Ref{XII.2.1} \ifthenelse{\boolean{orig}}{((v), (vi), (vii))}{((iv), (vi), (vii))}, the geometric fibers
of~$f^\an$ at the various points
$y$ of~$Y^\an$ are regular (\resp normal, \resp reduced) if and only if
the same holds for the geometric fibers of
$f$ at the various closed points $\psi (y)$ of~$Y$. The cases
(iv) (\resp (v), \resp (vi)) then follow from (i) and from the fact that
the set of points of~$Y$ where the geometric fibers of
$f$ are regular is open (EGA~IV 12.1.7).

(vii). If $f$ is injective, then so is~$f^\an$.
Conversely, suppose that $f^\an$ is injective, and let us show that
$f$ is as well. We may suppose
% page: 322
\marginpar{322}
%
that $f$ is of finite type. Since $f^\an$ is injective, the fibers
of~$f$ over the closed points of~$Y$ are radicial; since the set
of points of~$Y$ whose fiber is radicial is locally
constructible (EGA~IV 9.6.1), and since $Y$ is a
Jacobson scheme, all the fibers of~$f$ are radicial, and hence the
morphism is injective.

(viii). Let $\Delta \colon X \to X \times_Y X $ and $\Delta^\an
\colon X^\an \to X^\an \times_{Y^\an} X^\an $ be the diagonal
immersions, and let $\Theta \colon X^\an \times_{Y^\an} X^\an \to X \times_Y X
$ be the canonical morphism. By~\Ref{XII.2.3}, saying that
$\Delta (X)$ is closed in $X \times_Y X $ is equivalent to saying that
$\Delta^\an (X^\an)$ is closed in $X^\an \times_{Y^\an}
X^\an$.

Since an open immersion is nothing but an \'etale
injective morphism (EGA~IV 17.9.1 and \cite[\no 13 \S 1]{XII.4}), (xi) follows from (iii)
and (vii). Since an isomorphism is the same thing as a
surjective open immersion, (ix) follows from (xi) and
(\Ref{XII.3.2}~(i)) below. To say that $f$ is a monomorphism
is equivalent to saying that the diagonal morphism $\Delta \colon X \to
X \times_Y X$ is an isomorphism, so (x) follows from (ix).

\begin{proposition}
\label{XII.3.2}
Let $X$ and $Y$ be two $\CC$-schemes locally of finite type,
let $f\colon X \to Y$ be a morphism of finite type, and let $f^\an \colon X^\an \to
Y^\an $ be the morphism induced by~$f$ on the associated analytic
spaces. Let $P$ be the property of being
\begin{enumerate}
\item[(i)] surjective
\item[(ii)] dominant
\item[(iii)] a closed immersion
\item[(iv)] an immersion
\item[(v)] proper\footnote{We shall say that a morphism of
analytic spaces is proper if it is proper in the sense of \cite[ch\ptbl 1 \S 10 \no 1]{XII.1} and
if it is separated.}
\item[(vi)]
% page: 323
\marginpar{323}
%
finite.
\end{enumerate}
Then $f$ has property $P$ if and only if $f^\an$ does.
\end{proposition}

Let $\varphi \colon X^\an \to X$ and $\psi \colon Y^\an \to Y$ be the
canonical morphisms.

(i). If $f$ is surjective, then for every point $y$ of~$Y^\an$, $f^{-1}
(\psi (y)) $ is a nonempty closed subset of~$X$; it therefore contains
at least one closed point, which proves that $f^\an$ is
surjective. Conversely, if $f^\an$ is surjective, $f(X)$ is a
locally constructible subset of~$Y$ (EGA~IV 1.8.4) containing
all the closed points of~$Y$; hence $f(X)=Y$.

(ii) follows from~\Ref{XII.2.2}.

(iii). If $f$ is a closed immersion, then so is
$f^\an$ by \Ref{XII.1.1}~a). Conversely, if $f^\an$ is a
closed immersion, then so is~$f$ by
\Ref{XII.3.1}~(x) and \Ref{XII.3.2}~(v), since this is equivalent to saying
that $f$ is a proper monomorphism (EGA~IV 8.11.5).

(iv). It is clear that if $f$ is an immersion, then so is
$f^\an$. Conversely, suppose that $f^\an$ is an
immersion, and let $T$ be the image of~$X$ in $Y$, and $\overline{T}$
the scheme-theoretic closure of~$f$. We have a factorization of~$f$,
$$
X \lto{i} \overline{T} \lto{j} \ifthenelse{\boolean{orig}}{J}{Y} \quoi,
$$
where $j$ is a closed immersion, $i$ is the canonical morphism,
and from this we deduce the following factorization of~$f^\an$:
$$
X^\an \lto{i^\an} \overline{T}^{\mkern1mu \an} \lto{j^\an} Y^\an \
.
$$
Since
% page: 324
\marginpar{324}
%
$T = f(X)$ is a locally constructible subset of~$Y$ (EGA~IV
1.8.4), we have, by~\Ref{XII.2.2}, $\overline{T}^{\mkern1mu \an} =
\overline{f^\an (X^\an )}$. It follows that $i^\an (X^\an )$ is
open in~$\overline{T}^{\mkern1mu \an}$, and hence that $i(X) $ is open in
$\overline{T}$. Consider the canonical factorization of~$i$:
$$
X \lto{i_1} i (X) \lto{i_2} \overline{T}.
$$
The morphism $i_1^\an $ is a proper monomorphism, and hence so is
$i_1$ by \Ref{XII.3.2}~(v) and \Ref{XII.3.1}~(x);
this proves that $i_1$, and hence also $f$, is an immersion.

(v). Suppose that $f$ is proper, and let us show that
$f^\an$ is as well. Since properness of $f^\an $ is local
on~$Y^\an$, we may suppose that $Y$ is affine. By Chow's lemma
(EGA~II 5.6.1), one can find a projective $Y$-scheme $X'$ and a
projective surjective morphism
$$
g \colon X' \to X.
$$
The morphism $(fg)^\an = f^\an g^\an $ is projective and hence proper,
$g^\an $ is surjective, and it follows from \cite[ch\ptbl 1 \S 10]{XII.1} that
$f^\an $ is proper.

Conversely, suppose that $f^\an $ is proper, and let us show that
$f$ is as well. By \Ref{XII.3.1}~(viii), $f$ is
separated. It remains to prove that $f$ is universally
closed, and it even suffices to show that $f$ is closed;
indeed, for every $Y$-scheme $Y'$ locally of finite type, the
morphism
$$
f_{(Y')} = h \colon X \times_Y Y' \to Y'
$$
will also be closed, since $h^\an $ is proper. Let $T$ be a
closed subset of~$X$; $f(T)$ is a locally constructible set, and
we have
$$
f^\an (\varphi^{-1} (T)) = \psi^{-1} (f(T)).
$$
Since
% page: 325
\marginpar{325}
%
$f^\an $ is proper, $\psi^{-1} (f (T))$ is a closed subset of
$Y^\an $, and it therefore follows from~\Ref{XII.2.2} that
$$
\psi^{-1} (\overline{f(T)}) = \psi^{-1} (f(T))\quoi .
$$
This implies that $f(T) = \overline{f(T)}$, \ie that $f$
is closed and hence that $f$ is proper.

(vi). To say that a morphism is finite is equivalent to saying that it
is proper with finite fibers (EGA~III 4.4.2 and \cite[\no 19 \S 5]{XII.4}).
Since the set of points where the fibers of~$f$ are finite is
locally constructible (EGA~IV 9.7.9), the fibers of~$f$ are finite
if and only if the fibers of~$f^\an $ are finite; (vi)
therefore follows from (v).

\begin{remarque}
\label{XII.3.3}
\begin{enumerate}
\item[a)] Let $f \colon X \to Y$ be a morphism of~$\CC$-schemes
locally of finite type. The fact that $f^\an $ is a local isomorphism
does not imply that the same is true of~$f$. Indeed,
if $f$ is \'etale, then $f^\an $ is \'etale and hence a local isomorphism
\hbox{\cite[\no 13 \S 1]{XII.4}}, but $f$ need not be one.
\item[b)]
\ifthenelse{\boolean{orig}}
{the statement~\Ref{XII.3.2}}
{The statement~\Ref{XII.3.2}}
is not true if $f$ is not assumed to be of finite type. For example,
$f^\an $ can be a closed immersion without $f$ being one. Indeed, it
suffices to take $X$ to be the disjoint union of~$\ZZ $ copies of
$\Spec \CC$, $Y$ to be the affine line, and $f$ to be the morphism
obtained by sending the points of~$X$ to distinct points of~$Y$
forming a discrete subset.
\end{enumerate}
\end{remarque}
\end{filecontents*}
% Original input: drafts/SGA1_XII_4_English.texfrag
\begin{filecontents*}[overwrite,noheader]{sga1-cumulative-input-122.tex}
\section[Cohomological comparison theorems]{Cohomological comparison and existence theorems}
\label{XII.4}
% page: 326
\marginpar{326}
%

The purpose of this section is to reprove the results of
\cite[\no 2 th\ptbl 5 and th\ptbl 6]{XII.3}; the latter generalize the
theorems established in \cite[\no 12]{XII.10} for projective $X$ to the
case of a proper scheme, and extend them to the relative case. More
general results concerning proper relative
schemes over an analytic space are
proved in \cite[ch\ptbl VIII \no 3]{XII.7}.

Recall that the {\v{C}}ech cohomology used in \cite[\no 12]{XII.10} coincides with the usual cohomology in both the
algebraic and the analytic cases (EGA~III 1.4.1. and \cite[II 5.10]{XII.5}).
\subsection{}
\label{XII.4.1}
Let $f \colon X \to Y $ be a morphism of~$\CC$-schemes
locally of finite type, and consider the commutative diagram
$$
\xymatrix{ X^\an \ar[r]^\varphi \ar[d]_{f^\an} & X \ar[d]^f
\\ Y^\an \ar[r]^\psi & Y \quoi.}
$$
If $F$ is an $\cal{O}_X$-Module, then, for every integer $p \geq
0$, there are morphisms
$$
\R^p f_* F \lto{i} \R^p f_* (\varphi_* F^\an )
\lto{j} \R^p (f\cdot\varphi)_* F^\an \lto{k} \psi_*
(\R^p f_*^\an F^\an ) \quoi,
$$
where $i$ is induced by the canonical morphism $F \to \varphi_*
F^\an $, and $j, k$ are ``edge homomorphisms'' of Leray spectral
sequences. The composite $k.j.i$ gives rise to a canonical
morphism
\begin{equation*}
\label{eq:XII.4.1.1}
\tag{4.1.1} {\theta_p \colon (\R^p f_* F)^\an \to \R^p f_*^\an
(F^\an )}
\end{equation*}

\begin{theoreme}
\label{XII.4.2}
Let
% page: 327
\marginpar{327}
%
$f \colon X \to Y $ be a proper morphism of~$\CC$-schemes
locally of finite type, and let $F$ be a coherent $\cal{O}_X$-Module.
Then, for every integer $p \geq 0$, the morphism
\eqref{eq:XII.4.1.1}
$$
\theta_p \colon (\R^p f_* F)^\an \to \R^p f_*^\an (F^\an)
$$
is an isomorphism.
\end{theoreme}

1) \emph{The case where $f$ is projective}. The proof is
analogous to that of \cite[\no 13]{XII.10}. Let us recall it briefly. One
reduces to the case where $X$ is a projective space of the form $\PP^r_Y
$ over~$Y$. Let $\cal{Y} = Y^\an$, $\cal{P} = \PP^r_\cal{Y}$;
one first proves that
$$
f_*^\an \cal{O}_\cal{P} = \cal{O}_\cal{Y} \quoi, \qquad \R^p f_*^\an
(\cal{O}_\cal{P}) = 0 \qquad \text{for } p>0
$$
Indeed, to verify the preceding relations, one may
reduce to the case where $\cal{Y}$ is a ball $\cal{B}$ in an affine
space $\cal{E}^n$. Consider the ``standard covering'' $\{
\cal{U}_i \}$ of~$\cal{P}$ by $r+1$ open subsets isomorphic to $\cal{B}
\times \cal{E}^r$. Since these open subsets are Stein, for every
integer $p \geq 0$ there are isomorphisms
$$
\H^p (\{ \cal{U}_i \}, \cal{O}_\cal{P}) \isomto \H^p (\cal{P},
\cal{O}_\cal{P})
$$
One can then express the sections of the structure sheaf
$\cal{O}_\cal{P}$ over the open subsets $\cal{U}_i $ and over their
intersections in terms of Laurent series; an easy computation
shows that
$$
\H^0 (\cal{P}, \cal{O}_\cal{P} ) \isomto \H^0 (\cal{Y},
\cal{O}_\cal{Y} ) \quoi, \qquad \H^p (\cal{P}, \cal{O}_\cal{P} ) = 0
\qquad \text{for } p>0
$$
The proof is then completed by repeating the argument of \cite[\no 12
lemma~5]{XII.10}, with the cohomology groups replaced by the
cohomology sheaves.

2)
% page: 328
\marginpar{328}
%
\emph{The case where $f$ is proper}. We use EGA~III 3.1.2 to
reduce to the projective case. Let $\cal{K}$ be the category of
coherent $\cal{O}_X$-Modules for which $\theta_p$ is an
isomorphism for every $p \geq 0$. It suffices to prove that, for
every exact sequence $0 \to F' \to F \to F'' \to 0$
two of whose terms belong to $\cal{K}$, the same is true of the
third; that a direct summand of an object of~$\cal{K}$ belongs to
$\cal{K}$; and that, for every point $x$ of~$X$, one can find an
object $F$ of~$\cal{K}$ such that $F_x \neq 0$.

The first condition follows by applying the five lemma to the
following commutative diagram, whose rows are exact:
$$
\xymatrix@=.7cm{ \ar[r] & (\R^p f_* F')^\an \ar[d] \ar[r] &(\R^p f_*
F)^\an \ar[d] \ar[r] & (\R^p f_* F'')^\an \ar[d] \ar[r] &
(\R^{p+1} f_* F' )^\an \ar[r] \ar[d] &\, \\ \ar[r] & \R^p
f_*^\an F'^\an \ar[r] &\R^p f_*^\an F^\an \ar[r] & \R^p
f_*^\an F''^\an \ar[r] & \R^{p+1} f_*^\an F'^\an \ar[r] &
\quoi,}
$$
and the second condition is verified similarly.

To verify the third condition, one may restrict to the case
where $X$ is an irreducible scheme with generic point~$x$.
We could have assumed from the outset that $Y$ is noetherian. By
Chow's lemma (EGA~II 5.6.1), one can find a projective $Y$-scheme~$X'$
and a projective surjective morphism $g \colon X' \to X$. Moreover,
there is an integer~$n$ such that $\R^p g_* (\cal{O}_{X'} (n)) = 0 $ for every $p>0$ and
the canonical morphism $g^* g_* (\cal{O}_{X'} (n)) \to
\cal{O}_{X'} (n)$ is surjective (EGA~III 2.2.1). If we put
$F = g_* (\cal{O}_{X'} (n)) $, then the sheaf $F$ has the required
property. Indeed, $F_x \neq 0$; moreover, since the Leray spectral
sequence
$$
\R^p f_* (\R^q g_* (\cal{O}_{X'} (n))) \To \R^{p+q}
(\ifthenelse{\boolean{orig}}{f.g}{f\cdot g})_* (\cal{O}_{X'} (n))
$$
degenerates, there is an isomorphism
$$
\R^p f_* F \isomto \R^p (\ifthenelse{\boolean{orig}}{f.g}{f\cdot g})_* (\cal{O}_{X'} (n))
$$
% page: 329
\marginpar{329}%
%
As in the algebraic case, there is a canonical isomorphism
$$\R^p f_*^\an F^\an \isomto
\R^p (\ifthenelse{\boolean{orig}}{f.g}{f\cdot g})_*^\an (\cal{O}_{X'}(n)^\an ),$$
and the diagram
$$
\xymatrix@C=1cm{
(\R^p f_* F)^\an \ar[d]_{\theta_p} \ar[r]^-{\sim} & (\R^p (\ifthenelse{\boolean{orig}}{f.g}{f\cdot g})_* (\cal{O}_{X'} (n)))^\an \ar[d]_{\psi_p} \\
\R^p f_*^\an F^\an \ar[r]^-{\sim} & \R^p (\ifthenelse{\boolean{orig}}{f.g}{f\cdot g})_*^\an
(\cal{O}_{X'}(n)^\an)}
$$
is commutative. By~1), $\psi_p $ is an isomorphism; hence so is
$\theta_p$, which completes the proof.

\begin{corollaire}
\label{XII.4.3}
Let $X$ be a proper $\CC$-scheme, and let $F$ be a coherent
$\cal{O}_X$-Module. Then, for every integer $p\geq 0$, the canonical
morphism
$$
\H^p(X,F)\to \H^p(X^{\an},F^{\an})
$$
is an isomorphism.
\end{corollaire}

\begin{theoreme}
\label{XII.4.4}
Let $X$ be a proper $\CC$-scheme. The functor which assigns to every
coherent $\cal{O}_X$-Module $F$ its inverse image
$F^{\an}$ on~$X^{\an}$ is an equivalence of categories.
\end{theoreme}

1) \emph{The functor is fully faithful}. Indeed, let $F$
and $G$ be two coherent $\cal{O}_X$-Modules. The canonical morphism
$$
\Hom_{\cal{O}_X}(F,G)\to \Hom_{\cal{O}_{X^{\an}}}(F^{\an},G^{\an})
$$
is identified with the canonical morphism
$$
\H^0(X,\SheafHom_{\cal{O}_X}(F,G))\to
\H^0(X^{\an},\ifthenelse{\boolean{orig}}{\SheafHom_{\cal{O}_X}(F,G)}{\SheafHom_{\cal{O}_{X^{\an}}}(F^{\an},G^{\an})})
$$
% page: 330
\marginpar{330}%
%
(EGA~$0_{\textup{I}}$~6.7.6). Since $\SheafHom_{\cal{O}_X}(F,G)$ is
coherent, it follows from~\Ref{XII.4.3} that this morphism is
bijective.

2) \emph{The functor is essentially surjective}. When $X$ is
projective, the assertion follows from \cite[\no 12~th\ptbl 3]{XII.10}. The
general case reduces to the preceding one by using
Chow's lemma (EGA~II~5.6.1). Indeed, let $X'$ be a projective
$\CC$-scheme, let $f\colon X'\to X$ be a projective surjective morphism, and let $U$ be a
dense open subset of~$X$ such that $f$ induces an isomorphism
$f^{-1}(U)\simeq U$. We argue by noetherian induction
on~$X$; we may therefore suppose that, for every coherent sheaf
$\cal{G}$ on~$X^{\an}$ for which there exists a closed subset $Y$
of~$X$, distinct from~$X$, satisfying
$Y^{\an}\supset \Supp \cal{G}$, there is a coherent sheaf $G$
on~$X$ such that $G^{\an}\simeq \cal{G}$.

Let $\cal{F}$ be a coherent sheaf of $\cal{O}_{X^{\an}}$-modules, and let $\cal{K}$ and $\cal{L}$ be the coherent sheaves
defined by the condition that the sequence
$$
0\to\cal{K}\to\cal{F}\to f_*^{\an}f^{\an *}\cal{F}\to\cal{L}\to 0
$$
be exact. Since $X'$ is projective, there is a coherent
$\cal{O}_{X'}$-Module $F'$ such that
$F'^{\an}\simeq {f^{\an}}^*\cal{F}$; it then follows
from~\Ref{XII.4.2} that there is an isomorphism $(f_*F')^{\an}\simeq
f_*^{\an}{f^{\an}}^*\cal{F}$. Since $\cal{K}|U^{\an}$ and
$\cal{L}|U^{\an}$ vanish, there are coherent $\cal{O}_X$-Modules
$K$ and $L$ such that
$K^{\an}\simeq \cal{K}$ and $L^{\an}\simeq \cal{L}$. By~1), the
morphism $f_*^{\an}{f^{\an}}^*\cal{F}\to\cal{L}$ comes from a unique
morphism $f_*F'\to L$; let $I=\Ker(f_*F'\to L)$. The sheaf
$\cal{F}$ is then an extension of~$I^{\an}$ by~$K^{\an}$, and it remains
to see that this extension is obtained
% page: 331
\marginpar{331}
%
by inverse image from an extension of~$I$ by~$K$. It therefore suffices to
prove that the canonical morphism
\begin{equation*}
\label{eq:XII.4.*}
\tag{$*$} \ifthenelse{\boolean{orig}}{\Ext_{\cal{O}_X}^q(I,K)^{\an}}{\Ext_{\cal{O}_X}^q(I,K)}\isomto
\Ext_{\cal{O}_{X^{\an}}}^q(I^{\an},K^{\an})\quad \ifthenelse{\boolean{orig}}{q\neq 1}{q=1}
\end{equation*}
is bijective. But there are isomorphisms
$\mathbf{Ext}_{\cal{O}_X}^q(I,K)^{\an}\isomto
\mathbf{Ext}_{\cal{O}_{X^{\an}}}^q(I^{\an},K^{\an})$ for every integer
$q\geq 0$ (EGA~$0_{\textup{III}}$~12.3.5), and a morphism of spectral
sequences
$$
\xymatrix@R=.7cm@C=.5cm{
\H^p(X,\mathbf{Ext}_{\cal{O}_X}^q(I,K))\ar@{=>}[r]\ar[d] &
\Ext_{\cal{O}_X}^{p+q}(I,K)\ar[d]\\
\H^p(X^{\an},\mathbf{Ext}_{\cal{O}_{X^{\an}}}^q(I^{\an},K^{\an}))\ar@{=>}[r] & \Ext_{\cal{O}_{X^{\an}}}^{p+q}(I^{\an},K^{\an}).
}
$$
This morphism is an isomorphism because, by~\Ref{XII.4.3}, it is
an isomorphism on the terms $E^{pq}_2$, and this proves the
bijectivity of~$(*)$.

\begin{corollaire}
\label{XII.4.5}
The functor which assigns to every proper $\CC$-scheme $X$ the
analytic space $X^{\an}$ is fully faithful.
\end{corollaire}

We must show that, if $X$ and $Y$ are two proper
$\CC$-schemes, the canonical map
$$
\Hom_{\CC}(X,Y)\to\Hom(X^{\an},Y^{\an})
$$
is bijective. Now, to give a morphism from~$X$ to~$Y$ (\resp from
$X^{\an}$ to $Y^{\an}$) is equivalent to giving its graph,
\ie a closed subscheme $Z$ of~$X\times Y$ (\resp a closed
analytic subspace $\othercal{Z}$ of~$X^{\an}\times
Y^{\an}$), such that the restriction of the first projection
$X\times Y\to X$ to~$Z$ (\resp of~$X^{\an}\times Y^{\an}\to
X^{\an}$ to~$\othercal{Z}$) is
% page: 332
\marginpar{332}
%
an isomorphism. Since giving a closed subscheme
of~$X\times Y$ (\resp a closed analytic subspace
of~$X^{\an}\times Y^{\an}$) is equivalent to giving a coherent
sheaf of ideals in~$\cal{O}_{X\times Y}$ (\resp in~$\cal{O}_{X^{\an}\times Y^{\an}}$), the corollary follows
from~\Ref{XII.4.4}.

\begin{corollaire}
\label{XII.4.6}
Let $X$ be a proper $\CC$-scheme. The functor which assigns to every
scheme $X'$ finite (\resp finite \'etale) over~$X$ the analytic space
$X'^{\an}$
is an equivalence from the category of schemes finite (\resp finite
\'etale) over~$X$ to the category of analytic spaces finite (\resp
finite \'etale) over~$X^{\an}$.
\end{corollaire}

Indeed, to give a finite morphism $X'\to X$ (\resp $X'^{\an}\to
X^{\an}$) is equivalent to giving a coherent sheaf
of algebras over~$\cal{O}_X$ (\resp over~$\cal{O}_{X^{\an}}$)
\cite[\no 19~\S 5~th\ptbl 2]{XII.4}. The corollary therefore follows from~\Ref{XII.4.4}
in the unparenthesized case, and the parenthesized case follows
from it in view of~\Ref{XII.3.1}~(iii).
\end{filecontents*}
% Original input: drafts/SGA1_XII_5_A_English.texfrag
\begin{filecontents*}[overwrite,noheader]{sga1-cumulative-input-123.tex}
\section{Comparison theorems for \'etale coverings}
\label{XII.5}
\setcounter{subsection}{-1}
\subsection{}
\label{XII.5.0}
Let us make precise the notion of a finite covering of an analytic space.
%
\index{finite covering of an analytic space@finite covering of an analytic space|hyperpage}%
%
If $\othercal{X}$ is an analytic space, an analytic space
$\othercal{X}'$ finite over~$\othercal{X}$ is said to be a finite covering of
$\othercal{X}$ if every irreducible component of~$\othercal{X}'$ dominates
an irreducible component of~$\othercal{X}$.

% remark: should the fonts for formal schemes and complex analytic
% varieties both be \goth ? Some people think they should not.

\begin{theoreme}[\og Riemann existence theorem\fg]
\label{XII.5.1}
%
\index{Riemann (existence theorem)@Riemann (existence theorem)|hyperpage}%
\index{Riemann existence theorem@Riemann existence theorem|hyperpage}%
%
Let $X$ be a $\CC$-scheme locally of finite type, and let $X^{\an}$
be the analytic space associated with~$X$. The functor $\Psi$ which
assigns to every finite \'etale covering $X'$ of~$X$ the space
$X^{\prime\an}$ is an equivalence
% page: 333
\marginpar{333}
%
from the category of finite \'etale coverings of~$X$ to
the category of finite \'etale coverings of~$X^{\an}$.
\end{theoreme}

1) \emph{The functor $\Psi$ is fully faithful.} Let $X'$
and $X''$ be two finite \'etale coverings of~$X$, and let us prove that
the canonical map
\begin{equation*}
\label{eq:XII.5.*}
\tag{$*$} { \Hom_{X}(X',X'') \to
\Hom_{X^{\an}}(X^{\prime\an},X^{\prime\prime\an}) }
\end{equation*}
is bijective. We may suppose that $X'$ is connected. To give an
$X$-morphism from~$X'$ to $X''$ is equivalent to giving a
connected component $X_{i}$ of~$X'\times_{X}X''$ such that the morphism
$X_{i}\to X'$ induced by the first projection is an
isomorphism. Since the connected components of~$X'\times_{X}X''$
correspond bijectively to the connected components of
$X^{\prime\an}\times_{X^{\an}}X^{\prime\prime\an}$~\eqref{XII.2.6}, and since a
morphism $X_{i}\to X'$ is an isomorphism if and only if
$X_{i}^{\an}\to X^{\prime\an}$ is an isomorphism, this proves the
bijectivity of~\eqref{eq:XII.5.*}.

2) \emph{The functor $\Psi$ is essentially surjective.} Let
$\othercal{X}'$ be a finite \'etale covering of~$X^{\an}$, and let us prove
that there exists an \'etale covering $X'$ of~$X$ together with
an isomorphism $X^{\prime\an}\isomto\othercal{X}'$. In view of~1), the
question is local on~$X$, and we may therefore suppose that $X$ is affine.

a) Reduction to the case where $X$ is normal. We may suppose that $X$ is
reduced. Indeed, suppose that the theorem has been proved
for $X_{\red}$. The functor which assigns to a finite \'etale
covering $X'$ of~$X$ the finite \'etale covering
$X^{\prime\an}_{\red}$ of~$X_{\red}^{\an}$ is then an equivalence.
Since it is obtained by composing $\Psi$ with the functor $\Theta$ which
assigns to a finite \'etale covering of~$X^{\an}$ its inverse image
over~$X_{\red}^{\an}$, and since $\Theta$ is fully
faithful, this shows that $\Psi$ is an equivalence of
categories.

We
% page: 334
\marginpar{334}
%
may suppose that $X$ is normal. Indeed, let $\widetilde{X}$ be the
normalization of~$X$, and let $p\colon\widetilde{X}\to X$ be the
canonical morphism. Since $p$ is finite, $p$ is an effective descent
morphism for the category of \'etale
coverings~(\ifthenelse{\boolean{orig}}{IX~}{}\Ref{IX.4.7}). Supposing that the theorem has been
proved for $\widetilde{X}$, put
$\widetilde{\othercal{X}'}=\othercal{X}'\times_{X^{\an}}\widetilde{X}^{\an}$.
There then exist an \'etale covering $\widetilde{X}'$ of
$\widetilde{X}$ and an isomorphism $\widetilde{X}'{}^{\an}\simeq
\widetilde{\othercal{X}'}$. It follows from~1) that the natural
descent datum on~$\widetilde{\othercal{X}'}$ lifts to a descent datum
on~$\widetilde{X}'$ relative to $\widetilde{X}\to X$; this proves the
existence of an \'etale covering $X'$ of~$X$ together with an isomorphism
$i\colon X^{\prime\an}\times_{X^{\an}}\widetilde{X}^{\an} \simeq
\widetilde{\othercal{X}'}$ whose inverse images under the two projections
from
$\widetilde{X}^{\an}\times_{X^{\an}}\widetilde{X}^{\an}$ are the
same. By~\ifthenelse{\boolean{orig}}{IX~}{}\Ref{IX.3.2}, whose proof is
valid in the analytic case, the morphism $\widetilde{X}^{\an}\to
X^{\an}$ is a descent morphism for the category of
\'etale coverings, and consequently $i$ comes from an
isomorphism $X^{\prime\an}\simeq
\othercal{X}'$.

b) Reduction to the case where $X$ is regular. Let $U$ be
the open subset of regular points of~$X$, and let $i\colon U\to X$ and
$i^{\an}\colon U^{\an}\to X^{\an}$ be the canonical morphisms; since
$X$ is normal, we have $\codim(X-U,X)\geq 2$. Suppose that there exists an
\'etale covering $U'$ of~$U$ such that
$U^{\prime\an}\simeq \othercal{X}'|U^{\an}$, and let us show that $U'$ then
extends to an \'etale covering $X'$ of~$X$ such that
$X^{\prime\an}\simeq \othercal{X}'$. It suffices to see that $U'$ extends
to an \'etale covering $X'$ of~$X$; indeed, we will then have an
isomorphism $X^{\prime\an}|U^{\an}\simeq \othercal{X}'|U^{\an}$. But if
$\cal{F}$ and $\cal{G}$ are the coherent sheaves of algebras
over~$\cal{O}_{X^{\an}}$ defining $\othercal{X}'$ and
$X^{\prime\an}$ respectively, the fact that $X$ is normal and that
$\codim(X-U,X)\geq 2$ implies that the canonical morphisms
$$
\cal{F}\to i_*^{\an}(\cal{F}|U^{\an}) \qquad
\cal{G}\to i_*^{\an}(\cal{G}|U^{\an})
$$
are
% page: 335
\marginpar{335}
%
isomorphisms \cite[\no 3~prop\ptbl 4]{XII.11}. It follows that $\cal{F}$
and $\cal{G}$, and hence also $X^{\prime\an}$ and $\othercal{X}'$, are isomorphic.

Let $\varphi\colon X^{\an}\to X$ be the canonical morphism. Since the
problem of extending $U'$ to $X$ is local on~$X$, it is enough
to prove that, for every point $y$ of~$X^{\an}- U^{\an}$, the \'etale covering
$U'_{\varphi(y)}=U'\times_{X}\Spec\cal{O}_{X,\varphi(y)}$ of
$U_{\varphi(y)}=U\times_{X}\Spec\cal{O}_{X,\varphi(y)}$ extends
to $\Spec\cal{O}_{X,\varphi(y)}$. Let $H$ be the coherent
$\cal{O}_{U}$-algebra defining $U'$. The canonical
morphism
$$
\alpha\colon (i_*H)^{\an} \to i_*^{\an}(H^{\an})
=\cal{F}
$$
defines a morphism of sheaves of modules on~$\Spec\cal{O}_{X^{\an},y}$:
$$
\alpha_{y}\colon (i_*H)^{\an}_{y} \to \cal{F}_{y}\quoi,
$$
whose restriction to $U_{y}=U_{\varphi(y)}
\times_{\Spec\cal{O}_{X,\varphi(y)}}\Spec\cal{O}_{X^{\an},y}$ is an
isomorphism. But this proves that $H|U_{y}$ is trivial, and hence that
$U'_{\varphi(y)}$ extends to $\Spec\cal{O}_{X,\varphi(y)}$.

c) The case where $X$ is affine and regular. Let
$$
j\colon X\to P
$$
be a compactification of~$X$, where $P$ is a projective
$\CC$-scheme and $j$ is a dominant open immersion. By the
resolution of singularities theorem~\cite{XII.8}, one can
find a regular scheme $R$ and a projective morphism
$r\colon R\to P$ such that $r$ induces an isomorphism $r^{-1}(X)\simeq
X$ and $r^{-1}(X)$ is the complement in $R$ of a
divisor with normal crossings. Let
% page: 336
\marginpar{336}
%
$$
k\colon X\to R
$$
be the canonical immersion. We shall show that there exists a finite
normal covering~\eqref{XII.5.0} $\othercal{R}'$ of~$R^{\an}$ extending the
\'etale covering $X^{\prime\an}$. By
Proposition~\Ref{XII.5.3} below, such a covering is
unique; the problem of extending $X^{\prime\an}$ is therefore local on~$R^{\an}$ in a neighborhood of~$R^{\an}- X^{\an}$. Now every
point of~$R^{\an}- X^{\an}$ has an open neighborhood
$\othercal{V}$ isomorphic to a ball in an affine space
$\othercal{E}^{n}$ such that $\othercal{V}- \othercal{V}\cap X^{\an}$
is defined by the vanishing of the first $p$ coordinate functions
$z_{1},\dots,z_{p}$, where $0\leq p\leq n$. The fundamental group
of~$\othercal{U}=\othercal{V}\cap X^{\an}$ is isomorphic to
$\ZZ^{p}$, and every \'etale covering of~$\othercal{U}$ is a quotient
of a covering of the form
$$
\othercal{U}''= \othercal{U}[T_{1},\dots,T_{p}]/
(T_{1}^{n_{1}}-z_{1},\dots,T_{p}^{n_{p}}-z_{p})\quoi,
$$
where the $n_{i}$ are integers $>0$, by a subgroup $H$ of the
Galois group $\ZZ/n_{1}\ZZ\times\cdots\times\ZZ/n_{p}\ZZ$ of
$\othercal{U}''$. Now $\othercal{U}''$ extends to the regular covering
$$
\othercal{V}''= \othercal{V}[T_{1},\dots,T_{p}]/
(T_{1}^{n_{1}}-z_{1},\dots,T_{p}^{n_{p}}-z_{p})\quoi,
$$
of~$\othercal{V}$ on which $H$ acts, and the quotient of~$\othercal{V}''$
by $H$ is the desired extension.

The proof is then completed by~\Ref{XII.4.6}.
The covering $\othercal{R}'$ comes from a finite covering $R'$
of~$R$; the restriction of~$R'$ to $X$ is a covering $X'$ of
$X$ such that $X^{\prime\an}\simeq \othercal{X}'$, and
by~\Ref{XII.3.1}(iii), $X'$ is an \'etale covering
of~$X$.

\begin{corollaire}
\label{XII.5.2}
Let
% page: 337
\marginpar{337}%
%
$X$ be a connected $\CC$-scheme locally of finite type,
let $\varphi\colon X^{\an}{\to} X$ be the canonical morphism, and let $x$ be a point of
$X^{\an}$. Let $\pi_{1}(X^{\an},x)$ be the fundamental group of
the topological space $X^{\an}$ at $x$, and let $\pi_{1}(X,\varphi(x))$
be the fundamental group of the scheme $X$ at $\varphi(x)$
\textup{(\ifthenelse{\boolean{orig}}{V~}{}\Ref{V.7})}. Then $\pi_{1}(X,\varphi(x))$ is canonically
isomorphic to the completion of~$\pi_{1}(X^{\an},x)$ with respect to the
topology of subgroups of finite index.
\end{corollaire}

Indeed, let $\cal{C}$ be the category of finite \'etale
coverings of~$X^{\an}$, let $F$ be the functor from~$\cal{C}$ to
$\Ens$ which assigns to every finite \'etale covering $\othercal{X}'$ of
$X^{\an}$ the set of points of~$\othercal{X}'$ above
$x$, and let $\widehat{\pi}_{1}(X^{\an},x)$ be the profinite group
associated with $\cal{C}$ and $F$ as in~\ifthenelse{\boolean{orig}}{V~}{}\Ref{V.4}.
Since every finite \'etale covering of~$X^{\an}$ is a quotient of the
universal covering by a subgroup of finite index,
$\widehat{\pi}_{1}(X^{\an},x)$ is precisely the completion of
$\pi_{1}(X^{\an},x)$ for the topology of subgroups of finite index.
The corollary therefore follows from~\Ref{XII.5.1} and~\ifthenelse{\boolean{orig}}{V~}{}\Ref{V.6.10}.
\end{filecontents*}
% Original input: drafts/SGA1_XII_5_B_English.texfrag
\begin{filecontents*}[overwrite,noheader]{sga1-cumulative-input-124.tex}
\begin{proposition}
\label{XII.5.3}
Let $\othercal{X}$ be a normal analytic space, and let $\othercal{Y}$ be a
closed analytic subset such that
$\othercal{U}=\othercal{X}-\othercal{Y}$ is dense in $\othercal{X}$.
Then the functor which assigns to every finite normal
covering~\eqref{XII.5.0} $\othercal{X}'$ of~$\othercal{X}$ its restriction
to $\othercal{U}$ is fully faithful.
\end{proposition}

Let $\othercal{X}'$ and $\othercal{X}''$ be two finite normal coverings
of~$\othercal{X}$. We must show that the canonical map
$$
\Hom_{\othercal{X}}(\othercal{X}',\othercal{X}'') \to
\Hom_{\othercal{U}}(\othercal{X}'|\othercal{U},\othercal{X}''|\othercal{U})
$$
is bijective. Let $u$, $v$ be two $\othercal{X}$-morphisms from
$\othercal{X}'$ to $\othercal{X}''$ whose restrictions to $\othercal{U}$
are the same, and let us prove that $u=v$. The morphisms $u$ and
% page: 338
\marginpar{338}
%
$v$ coincide on the dense open subset
$\othercal{U}\times_{\othercal{X}}\othercal{X}'$, and hence on the
underlying topological spaces. By \hbox{\cite[\no 19~\S 4~cor\ptbl 5]{XII.4}},
this proves that $u=v$.

Now let $u$ be a $\othercal{U}$-morphism from~$\othercal{X}'|\othercal{U}$
to $\othercal{X}''|\othercal{U}$, and let us show that it extends to all
of~$\othercal{X}'$. We may suppose that $\othercal{X}'$ is regular.
Indeed, since $\othercal{X}'$ is normal, one can find an open subset
$\othercal{V}$ of~$\othercal{X}$ whose complement is an analytic subset
of codimension $\geq 2$, such that
$\othercal{X}'\times_{\othercal{X}}\othercal{V}=\othercal{V}'$ is regular.
Let $\othercal{V}''=\othercal{X}''\times_{\othercal{X}}\othercal{V}$, and suppose
the proposition proved for $\othercal{V}$. Consider the commutative
diagram
$$
\xymatrix{ \othercal{V}' \ar[dd]_{g'}\ar[rd]\ar[rrr]^{i'} && &\othercal{X}'
\ar[dd]^{\hbox{\raise10mm\hbox{$f'$}}} \cr &\othercal{V}'' \ar[ld]_{g''\mkern-10mu}\ar[rrr]^-{i''~~} & &
&\othercal{X}'' \ar[ld]_{f''\mkern-12mu} \cr \othercal{V} \ar[rrr]^{i} && &\othercal{X} \cr
}
$$
The morphism $u$ determines a morphism of
$\cal{O}_{\othercal{V}}$-algebras $g''_*\cal{O}_{\othercal{V}''}\to
g'_*\cal{O}_{\othercal{V}'}$, from which one deduces a morphism
$$
i_*g''_*\cal{O}_{\othercal{V}''} \to
i_*g'_*\cal{O}_{\othercal{V}'}\quoi.
$$
In view of the isomorphisms
$i'_*\cal{O}_{\othercal{V}'}\simeq\cal{O}_{\othercal{X}'}$,
$i''_*\cal{O}_{\othercal{V}''}\simeq\cal{O}_{\othercal{X}''}$
\cite[\no 3~prop\ptbl 4]{XII.11}, one deduces a morphism of
$\cal{O}_{\othercal{X}}$-algebras
$$
f''_*\cal{O}_{\othercal{X}''} \to
f'_*\cal{O}_{\othercal{X}'}\quoi,
$$
and hence the desired morphism $\othercal{X}'\to\othercal{X}''$.

We
% page: 339
\marginpar{339}
%
henceforth suppose that $\othercal{X}'$ is regular. Let
$\othercal{U}'=\othercal{U}\times_{\othercal{X}}\othercal{X}'$,
$\othercal{Y}'=\othercal{X}'-\othercal{U}'$. We regard $\othercal{Y}'$ as a
reduced analytic subspace of~$\othercal{X}'$; if $\othercal{Y}'_{1}$
is the singular locus of~$\othercal{Y}'$, then $\dim \othercal{Y}'_{1} <
\dim \othercal{Y}'$~\cite[\no 20~D\ptbl Th\ptbl 3]{XII.4}. Thus, by induction
on the dimension of~$\othercal{Y}'$, we may suppose that $\othercal{Y}'$
is smooth. Since it suffices to extend $u$ to an open neighborhood of
each point of~$\othercal{Y}'$, the implicit function theorem allows us
to suppose that $\othercal{X}'$ is a ball in an affine space
$\othercal{E}^{n}$ and that $\othercal{Y}'$ is the closed subset defined
by the vanishing of the first $p$ coordinate functions
$z_{1},\dots,z_{p}$, with $0\leq p\leq n$.

The morphism $u$ determines a section $s$ of
$p\colon\othercal{X}'\times_{\othercal{X}}\othercal{X}''\to\othercal{X}'$ over
$\othercal{U}'$; after restricting $\othercal{X}'$ if necessary, we may
suppose that $p_*(\cal{O}_{\othercal{X}'\times_{\othercal{X}}\othercal{X}''})$
is generated by elements $x_{1},\dots,x_{q}$ of
$\Gamma(\othercal{X}',
p_*\cal{O}_{\othercal{X}'\times_{\othercal{X}}\othercal{X}''})$; let
$u_{1},\dots,u_{q}\in\Gamma(\othercal{U}',\cal{O}_{\othercal{X}'})$ be the images
under $s$ of~${x_{1}}|\othercal{U}',\dots,{x_{q}}|\othercal{U}'$. To say
that $s$ extends to $\othercal{X}'$ is equivalent to saying that
$u_{1},\dots,u_{q}$ extend to sections of
$\Gamma(\othercal{X}',\cal{O}_{\othercal{X}'})$. But since \ifthenelse{\boolean{orig}}{$f$}{$p$} is finite,
each $u_{i}$ is a Laurent series in $z_{1},\dots,z_{p}$, with
coefficients that are convergent power series in
$z_{p{+}1},\dots,z_{n}$, and satisfies integral-dependence relations.
It follows that $u_{i}$ is bounded, and hence is a convergent power
series in $z_{1},\dots,z_{n}$; consequently it extends to
$\othercal{X}'$.

One may ask whether the functor introduced in~\Ref{XII.5.3} is an
equivalence of categories. The \textsc{Grauert}-\textsc{Remmert} theorem
\cite{XII.6} answers this question; below we give a proof using
resolution of singularities. One could also have used the
\textsc{Grauert}-\textsc{Remmert} theorem to prove~\Ref{XII.5.1};
this is what was done before~\cite{XII.8} became available.

\ifthenelse{\boolean{orig}}
{}
{\enlargethispage{.5cm}}
\begin{theoreme}[\textsc{Grauert}-\textsc{Remmert} theorem]
\label{XII.5.4}
Let
% page: 340
\marginpar{340}
%
%
\index{Grauert-Remmert (theorem)|hyperpage}%
\index{Grauert-Remmert theorem@Grauert-Remmert theorem|hyperpage}%
%
$\othercal{X}$ be a normal analytic space, and let $\othercal{Y}$ be a
closed analytic subset such that
$\othercal{U}=\othercal{X}-\othercal{Y}$ is dense in~$\othercal{X}$. Let
$\othercal{U}'$ be a finite normal covering of~$\othercal{U}$; suppose
that there exists a nowhere dense closed analytic subset $\othercal{S}$
of~$\othercal{X}$ such that the restriction of~$\othercal{U}'$ to
$\othercal{U}-\othercal{U}\cap\othercal{S}$ is \'etale. Then there exists a
finite normal covering $\othercal{X}'$ of~$\othercal{X}$ which extends
$\othercal{U}'$, and $\othercal{X}'$ is unique up to isomorphism.
\end{theoreme}

Uniqueness follows from~\Ref{XII.5.3}. The problem of extending
$\othercal{U}'$ is therefore local on~$\othercal{X}$. We may suppose that
$\othercal{U}$ is regular and that $\othercal{U}'$ is \'etale over~$\othercal{U}$.
Indeed, the set of regular points of~$\othercal{U}$ is a dense open
subset $\othercal{V}$ of~$\othercal{X}$ whose complement is an analytic
subset \hbox{\cite[\no20~D th\ptbl 2]{XII.4}}, and it suffices to replace
$\othercal{U}$ by the open subset $\othercal{V}-\othercal{V}\cap\othercal{S}$.

Let $y$ be a point of~$\othercal{X}-\othercal{U}$, and let us show that
$\othercal{U}'$ can be extended to a neighborhood of~$y$. After
restricting $\othercal{X}$ to an open neighborhood of~$y$, it follows
from the resolution of singularities theorem \cite{XII.8} that one can
find a regular analytic space $\othercal{X}_1$ and a projective morphism
$f\colon\othercal{X}_1\to\othercal{X}$ which, upon restriction to~$\othercal{U}$,
induces an isomorphism
$\othercal{U}_1=f^{-1}(\othercal{U})\simeq\othercal{U}$, such that
$\othercal{U}_1$ is the complement in~$\othercal{X}_1$ of a normal
crossings divisor. Let us show that $\othercal{U}'$ extends to a normal
finite covering of~$\othercal{X}_1$. Since the question is local
on~$\othercal{X}_1$, we may suppose that $\othercal{X}_1$ is a ball in
an affine space $\othercal{E}^n$ and that
$\othercal{X}_1-\othercal{U}_1$ is defined by the vanishing of the first
$p$ coordinate functions
\ifthenelse{\boolean{orig}}
{$z_1,\dots,z_q$,}
{$z_1,\dots,z_p$,}
% correction: `z_q' should be `z_p'.
with $0\leq p\leq n$. The \'etale covering $\othercal{U}'$
of~$\othercal{U}_1$ is a quotient of a covering of the form
$$
\othercal{U}_2=\othercal{U}_1[T_1,\dots,T_p]/\left(T_1^{n_1}-z_1,\dots,
T_p^{n_p}-z_p\right)
$$
by a subgroup $H$ of the Galois group of~$\othercal{U}_2$. The
covering $\othercal{U}_2$ extends
% page: 341
\marginpar{341}
%
to the covering
$$
\othercal{X}_2=\othercal{X}_1[T_1,\dots,T_p]/\left(T_1^{n_1}-z_1,\dots,
T_p^{n_p}-z_p\right)
$$
of~$\othercal{X}_1$, on which $H$ acts, and $\othercal{X}_2/H$ extends
$\othercal{U}'$ to~$\othercal{X}_1$.

Let $\othercal{X}_1'$ denote the finite normal covering of~$\othercal{X}_1$
which extends $\othercal{U}'$, and let $\othercal{F}_1$ denote the
coherent $\cal{O}_{\othercal{X}_1}$-algebra \ifthenelse{\boolean{orig}}{defined by~$\othercal{F}_1$}{defining~$\othercal{X}_1'$}.
By the \textsc{Grauert}-\textsc{Remmert} finiteness theorem
\hbox{\cite[\no15 th\ptbl 1.1]{XII.4}}, $f_*\othercal{F}_1$ is a coherent
$\cal{O}_{\othercal{X}}$-algebra. It therefore corresponds to a finite
covering $\othercal{X}'$ of~$\othercal{X}$, which is moreover normal
since $\othercal{X}_1'$ is normal, and $\othercal{X}'$ is the desired
extension of~$\othercal{U}'$.

\begin{remarque}
\label{XII.5.5}
In the statement~\Ref{XII.5.4}, the hypothesis on the locus of points
where the morphism $\othercal{U}'\to\othercal{U}$ is not \'etale cannot be
omitted. For example, let $\othercal{X}$ be the unit disk in the complex
plane, let $\othercal{U}$ be the complement of the origin in~$\othercal{X}$,
and let $\othercal{U}'=\othercal{U}[T]/(T^2-\sin 1/z)$, where~$z$ is the
coordinate function on~$\othercal{X}$. Then $\othercal{U}'$ is a normal
finite covering of~$\othercal{U}$ which does not extend to~$\othercal{X}$.
Indeed, suppose that $\othercal{U}'$ extends to a finite covering
$\othercal{X}'$ of~$\othercal{X}$; the locus of points of~$\othercal{X}$
where the morphism $\othercal{X}'\to\othercal{X}$ is not \'etale would then
be a closed analytic subset containing all points $z$ such that
$\sin 1/z=0$, which is absurd.

The hypothesis on the singular locus of the morphism
$\othercal{U}'\to\othercal{U}$ can, however, be omitted when
$\codim(\othercal{X}-\othercal{U},\othercal{X})\geq~2$. Indeed, we may suppose
that $\othercal{U}$ is regular. The locus of points of~$\othercal{U}$ where
$\othercal{U}'\to\othercal{U}$ is not \'etale is a divisor
of~$\othercal{U}$, and it follows from the \textsc{Remmert}-\textsc{Stein}
theorem \cite[th\ptbl 3]{XII.9} that it is
% page: 342
\marginpar{342}
%
the trace on~$\othercal{U}$ of a divisor of~$\othercal{X}$. Now, in this
case, if $\othercal{A}$ is a coherent $\cal{O}_{\othercal{U}}$-algebra
such that $\othercal{U}'=\Spec \an(\othercal{A})$, and if
$i:\othercal{U}\to\othercal{X}$ is the canonical morphism, it suffices to
take $\othercal{X}'=\Spec\an(i_*\othercal{A})$; indeed, we know that
$i_*\othercal{A}$ is
\ifthenelse{\boolean{orig}}
{coherents}
{coherent}
% correction: `coherents' should be `coherent'.
\cite[\no 1 th\ptbl 1]{XII.11}.
\end{remarque}

\ifthenelse{\boolean{orig}}
{}
{\begin{remarqueMR}
\label{XII.5.6}
There exist nontrivial finitely presented groups $G$ which have no
finite-index subgroups distinct from~$G$, for example G\ptbl Higman's
group (\cf J-P\ptbl Serre, Arbres et amalgames, Ast\'erisque
N$^\circ$46, prop\ptbl 6, ch\ptbl I, \S 1). Consequently, if such a
group is realized as the topological fundamental group of a scheme $V$
over~$\CC$, say a smooth projective scheme, the topological space
$V_{\mathrm{top}}$ underlying $V$ is not simply connected, but $V$ is
algebraically simply connected. At present, no such $V$ is known.
Let us nevertheless mention that D\ptbl Toledo constructed smooth
projective schemes $V$ over~$\CC$ whose topological fundamental group is
not Hausdorff for the topology of finite-index subgroups (the natural
morphism $\pi_1(V_{\mathrm{top}})\to\pi_1(V_{\alg})$ is not injective).
[D\ptbl Toledo, Projective varieties with non-residually finite
fundamental group, Publ.\ Math.\ IHES \textbf{77} (1993),
p\ptbl103--119].
\end{remarqueMR}}


%\section{Bibliography}
\begin{thebibliography}{10}
% page: 343
\marginpar{343}
%

\bibitem[1]{XII.1} N\ptbl \textsc{Bourbaki},
Topologie G\'en\'erale, Hermann, Paris (1960).

\bibitem[2]{XII.2} N\ptbl \textsc{Bourbaki},
Alg\`ebre Commutative, Hermann, Paris (1961).

\bibitem[3]{XII.3} H\ptbl \textsc{Cartan}. S\'eminaire E.N.S., Paris (1956-57).

\bibitem[4]{XII.4} H\ptbl \textsc{Cartan}. S\'eminaire E.N.S., Paris (1960-61).

\bibitem[5]{XII.5} R\ptbl \textsc{Godement}, Th\'eorie des Faisceaux, Hermann, Paris
(1958).

\bibitem[6]{XII.6} H\ptbl \textsc{Grauert} and R\ptbl \textsc{Remmert},
\ifthenelse{\boolean{orig}}
{Komplexe R\"aume, Math.\ Ann., 136 (1958).}
{Komplexe R\"aume, Math.\ Ann. \textbf{136} (1958), p\ptbl 245--318.}

\bibitem[7]{XII.7} M\ptbl \textsc{Hakim}, Sch\'emas relatifs, thesis, Paris (1967).

\bibitem[8]{XII.8} H\ptbl \textsc{Hironaka},
\ifthenelse{\boolean{orig}}
{Resolution of singularities of an algebraic variety over a field of characteristic zero, Annals of Math., 39 (1964).}
{Resolution of singularities of an algebraic variety over a field of characteristic zero, Ann.\ of Math. \textbf{39} (1964), p\ptbl 109--236.}


\bibitem[9]{XII.9} R\ptbl \textsc{Remmert} and K\ptbl \textsc{Stein},
Ueber die wesentlichen Singularit\"aten analytischer Mengen,
\ifthenelse{\boolean{orig}}
{Math.\ Ann., 126 (1953) p\ptbl 263--306.}
{Math.\ Ann. \textbf{126} (1953), p\ptbl 263--306.}

\bibitem[10]{XII.10}
\ifthenelse{\boolean{orig}}{J.P\ptbl}{J-P\ptbl}%
\textsc{Serre},
\ifthenelse{\boolean{orig}}
{G\'eom\'etrie Alg\'ebrique et G\'eom\'etrie Analytique, Annales de l'Institut Fourier, VI (1956), p\ptbl 1--42.}
{G\'eom\'etrie alg\'ebrique et g\'eom\'etrie analytique, Ann.\ Inst.\ Fourier (Grenoble) \textbf{6} (1956), p\ptbl 1--42.}

\bibitem[11]{XII.11}
\ifthenelse{\boolean{orig}}{J.P\ptbl}{J-P\ptbl}%
\textsc{Serre},
\ifthenelse{\boolean{orig}}
{Prolongement des Faisceaux analytiques
coh\'erents, Annales de l'Institut Fourier, XVI (1966), p\ptbl 363--374.}
{Prolongement de faisceaux analytiques
coh\'erents, Ann.\ Inst.\ Fourier (Grenoble) \textbf{16} (1966), p\ptbl 363--374.}


\end{thebibliography}
\end{filecontents*}
% Original input: drafts/SGA1_XII_OPENING_AND_1_English.texfrag
\begin{filecontents*}[overwrite,noheader]{sga1-cumulative-input-125.tex}
\chapterspace{-4}
\chapter{Algebraic geometry and analytic geometry}
\label{XII}
% page: 311
\marginpar{311}
%

\begin{center}
{Mme M.\ \textsc{Raynaud}\footnote{Based on unpublished notes by
A\ptbl Grothendieck.}}
\end{center}
\vspace*{1cm}
% remark: in this section, the ideal sheaves I are not underlined,
% hence are typeset in normal math mode, i.e., not as \cal{I}. Maybe
% some time this should be changed.

Proceeding as in \cite{XII.10}, one associates with every scheme $X$
locally of finite type over the field of complex numbers $\CC$ an
analytic space~$X^\an$
\label{indnot:lb}\oldindexnot{$X^\an$|hyperpage}%
%
whose underlying set is $X(\CC)$.

In \No \Ref{XII.2} and~\Ref{XII.3} of this expos\'e, we give a
\og dictionary\fg between the usual properties of~$X$ and
$X^\an$, and between the properties of a morphism $f\colon X\to
Y$ and those of the associated morphism $f^\an\colon X^\an \to Y^\an$.
\label{indnot:lc}\oldindexnot{$f^\an$|hyperpage}%
%

We then show that the comparison theorems for coherent sheaves
on~$X$ and $X^\an$, established in \cite[\No
12]{XII.10} for a projective variety, remain valid when
$X$ is a proper scheme.

Finally, in \No \Ref{XII.5} we prove the equivalence between the category of
finite \'etale coverings of~$X$ and the category of finite \'etale
coverings of~$X^\an$. As a bonus for the reader, we give a new proof
of the Grauert-Remmert theorem \cite{XII.6}, using resolution of
singularities~\cite{XII.8}.

\section{Analytic space associated with a scheme}
\label{XII.1}
% page: 312
\marginpar{312}
%

Let $X$ be a scheme locally of finite type over~$\CC$. Let $\Phi$
be the functor from the category of analytic spaces \cite[\No 9]{XII.4} to
the category of sets which assigns to an analytic space~$\othercal{X}$ the set of morphisms of ringed spaces in
$\CC$-algebras $\Hom_\CC (\othercal{X},X)$. We have the following
theorem:

\begin{theoremedefinition}
\label{XII.1.1} The functor $\Phi$ is representable by an
analytic space $X^\an$ and a morphism $\varphi\colon X^\an\to
X$. We say that $X^\an$ is the analytic space associated with~$X$.
%
\index{associated analytic space|hyperpage}%
%

If $|X^\an|$ is the underlying set of $X^\an$, $\varphi$
induces a bijection from~$|X^\an|$ onto the set $X(\CC)$ of points
of~$X$ with values in $\CC$. Moreover, for every point $x$ of~$X^\an$, the morphism
$$
\varphi_x\colon\cal{O}_{X,\varphi(x)}\to\cal{O}_{X^\an,x},
$$
which is necessarily local, yields upon passage to the
completions an isomorphism
$$
\widehat{\varphi}_x\colon\widehat{\cal{O}}_{X,\varphi(x)}\isomto
\widehat{\cal{O}}_{X^\an,x}\ifthenelse{\boolean{orig}}{,}{.}
$$
In particular, the morphism $\varphi$ is flat.
\end{theoremedefinition}

Note that the fact that $\varphi$ induces a bijection from~$X^\an$ onto~$X(\CC)$ follows from the universal property of~$X^\an$.
We also have the following assertions:

\begin{enumerate}
\item [a)] If the theorem is true for a scheme $Y$, it
is also true for every subscheme $X$ of~$Y$. Suppose
first that $X$ is an open subscheme of~$Y$; if $\psi\colon
Y^\an\to Y$ is the canonical morphism, $\psi^{-1}(X)$ is an open subset
% page: 313
\marginpar{313}
%
of~$Y^\an$, which we endow with the analytic-space structure induced
by that of~$Y^\an$. Since every morphism from an analytic space
$\othercal{X}$ to $X$ factors through $Y^\an$ by the
universal property of the latter, and therefore through
$X^\an$, which is the fiber product $Y^\an\times_Y X$, $X^\an$ is
the analytic space associated with $X$. Finally, the assertion
concerning the $\varphi_x$ is clear.

It remains only to consider the case where $X$ is a
closed subscheme of~$Y$. Let~$I$ be the coherent $\cal{O}_Y$-ideal
defining $X$; then $I\cdot\cal{O}_{Y^\an}$ is a coherent
sheaf of ideals of~$\cal{O}_{Y^\an}$ which
defines a closed analytic subspace $X^\an$ of~$Y^\an$; as
in the case of an open subscheme, one sees that $X^\an$ is
the analytic space associated with $X$. Let $\varphi\colon
X^\an\to X$ be the canonical morphism. For every point $x$ of~$X^\an$, the
morphism $\varphi_x$ is none other than the morphism
$$
\cal{O}_{Y,\psi(x)}/I_{\psi(x)}\to
\cal{O}_{Y^\an,x}/I_{\psi(x)}\cdot\cal{O}_{Y^\an,x}
$$
induced by $\psi_x$; its completion
$$
\widehat{\varphi}_x\colon
\widehat{\cal{O}}_{Y,\psi(x)}/I_{\psi(x)}\cdot\widehat{\cal{O}}_{Y,\psi(x)}
\to\widehat{\cal{O}}_{Y^\an,x}/I_{\psi(x)}\cdot\widehat{\cal{O}}_{Y^\an,x}
$$
is an isomorphism since $\widehat\psi_x$ is one, which
proves a).

\item [b)] If we have two $\CC$-schemes $X_1$, $X_2$ such that
$X_1^\an$ and $X_2^\an$ exist, then
$(X_1\times X_2)^\an$ exists as well. Indeed, let $\varphi_1\colon X_1^\an\to
X_1$, $\varphi_2\colon X_2^\an\to X_2$ be the canonical morphisms,
and let $p_1$, $p_2$ be the two projections of~$X_1^\an\times X_2^\an$. It
follows formally from EGA~I~1.8.1 that $X_1\times X_2$ is the
product of~$X_1$ and $X_2$ in the category of locally
ringed spaces; consequently, the morphisms
$\varphi_1\cdot p_1$ and $\varphi_2\cdot p_2$ define a
% page: 314
\marginpar{314}
%
morphism $\varphi\colon X_1^\an\times X_2^\an\to X_1\times X_2$, and
the pair $(X_1^\an\times X_2^\an,\varphi)$ represents the
functor $\othercal{X}\mto\Hom_\CC(\othercal{X},X_1\times X_2)$.

\item [c)] If $\othercal{E}^1$ denotes affine space of dimension 1,
\ie the topological space $\CC$ endowed with the sheaf of
holomorphic functions, the functor%
\ifthenelse{\boolean{orig}}
{}
{{\renewcommand{\thefootnote}{*}\addtocounter{footnote}{-1}\footnote{\lcrochetbf Added in 2003: $E^1_\CC$ denotes the (algebraic) affine line over $\CC$.\rcrochetbf}}}
$\othercal{X}\mto\Hom_\CC (\othercal{X},E^1_\CC)$
is representable by $\othercal{E}^1$, the canonical morphism
$\varphi\colon\othercal{E}^1\to E^1_\CC$ being the evident morphism.
Indeed, to give a morphism from an analytic space
$\othercal{X}$ to $E^1_\CC$ is equivalent to giving an
element of~$\Gamma(\othercal{X},\cal{O}_\othercal{X})$, which is
also equivalent to giving a morphism from~$\othercal{X}$ to~$\othercal{E}^1$. We
plainly have a bijection $|\othercal{E}^1|\simeq E^1(\CC)$, and, for
every point $x\in\othercal{E}^1$, the morphism $\widehat\varphi_x$ is
none other than the identity morphism of a ring of formal power series
in one variable over~$\CC$.
\end{enumerate}

\ifthenelse{\boolean{orig}}{}
{\enlargethispage{.5cm}}%
It follows from b) and c) that the theorem is true for
affine space $E^n_\CC$, $n\geq 0$. Using a), one sees that it is
also true for every affine scheme $X$ locally of finite
type over~$\CC$. If $X$ is no longer assumed affine and if $(X_i)$ is
a covering of~$X$ by affine open subsets, it follows from the
universal property and from a) that the $X_i^\an$ glue together
and thus define the analytic space $X^\an$ associated
with~$X$.

\subsection{}
\label{XII.1.2}
Let $f\colon X\to Y$ be a morphism of~$\CC$-schemes locally of
finite type. If $\varphi\colon X^{\an}\to X$ and $\psi\colon Y^{\an}\to
Y$ are the canonical morphisms, it follows from the
universal property of~$Y^{\an}$ that there exists a unique
morphism $f^{\an}\colon X^{\an}\to Y^{\an}$ such that the diagram
$$
\xymatrix{
X^{\an} \ar[r]\ar[d]_-{f^{\an}}& X\ar[d]^-f\\ Y^{\an} \ar[r]& Y
}
$$
is
% page: 315
\marginpar{315}
%
commutative. We have therefore defined a functor $\Phi$ from the
category of $\CC$-schemes locally of finite type to the
category of analytic spaces.

The functor $\Phi$ commutes with finite inverse limits. It suffices
indeed to see that $\Phi$ commutes with fiber products. Now, if
$X$, $Y$, $Z$ are schemes locally of finite type over~$\CC$,
it follows from the fact that $X\times_Z Y$ is the fiber product of
$X$ and $Y$ over~$Z$ in the category of locally
ringed spaces that $X^{\an}\times_{Z^{\an}}Y^{\an}$
satisfies the universal property characterizing
$(X\times_Z Y)^{\an}$.

\subsection{}
\label{XII.1.3}
Let $X$ be a $\CC$-scheme locally of finite type, let $X^{\an}$ be
the associated analytic space, and let $\varphi\colon X^{\an}\to X$ be the
canonical morphism. If $F$ is an ${\cal{O}}_X$-module, the inverse
image $\varphi^*F=F^{\an}$ is a sheaf of modules over~${\cal{O}}_{X^{\an}}$. This defines a functor from the
category of $\cal{O}_X$-modules to the category of
modules on~$X^{\an}$. This functor commutes with inductive limits
(EGA~0~4.3.2). Since the sheaf $\cal{O}_{X^{\an}}$ is coherent
\cite[\no18~\S 2~th\ptbl 2]{XII.4}, it sends coherent sheaves to
coherent sheaves (EGA~0~5.3.11). Moreover:

\begin{subproposition}
\label{XII.1.3.1}
The functor which assigns to an \ifthenelse{\boolean{orig}}{$\cal{O}_{X^{\an}}$}{$\cal{O}_X$}-module $F$ its
inverse image $F^{\an}$ on~$X^{\an}$ is exact, faithful, and
conservative.
\end{subproposition}

Exactness follows from the fact that the morphism $\varphi\colon
X^{\an}\to X$ is flat \ifthenelse{\boolean{orig}}{\eqref{XIII.1.1}}{\eqref{XII.1.1}}. Let us prove that the functor $F\mto
F^{\an}$ is faithful. In view of exactness, it suffices to
show that if $F^{\an}$ is zero, then so is~$F$. But,
for every point $x$ of~$X^{\an}$, we then have $F_{\varphi
(x)}\otimes_{\cal{O}_{X,\varphi (x)}}\cal{O}_{X^{\an},x}=0$. Since the
morphism $\cal{O}_{X,\varphi (x)}\to \cal{O}_{X^{\an},x}$ is
faithfully flat, we have $F_{\varphi (x)}=0$ for every closed point
$\varphi (x)$ of~$X$, and, since $X$ is Jacobson
(EGA~IV~10.4.8), this implies that $F$ is zero.

The
% page: 316
\marginpar{316}
%
fact that the functor $F\mto F^{\an}$ is conservative follows formally
from exactness and faithfulness.
\end{filecontents*}
% Original input: drafts/SGA1_XI_2_English.texfrag
\begin{filecontents*}[overwrite,noheader]{sga1-cumulative-input-126.tex}
\section{Abelian varieties}
\label{XI.2}
Let $k$ be an algebraically closed field, let $A$ be an abelian
variety over~$k$, \ie a group scheme over~$k$ that is proper over~$k$,
smooth over~$k$, and connected, and finally let $G$ be a commutative
group scheme of finite type over~$k$. Denote by $\Ext(A,G)$ the group
of classes of commutative extensions of~$A$ by~$G$, and by
$\H^1(A,G)$ the group of classes of principal bundles over~$A$ with
group~$G$ (compare \No \Ref{XI.4} below), and consider the canonical
homomorphism
$$
\Ext(A,G)\to \H^1(A,G)\quoi.
$$
An
% page: 287
\marginpar{287}
%
argument of Serre~\cite[chap\ptbl VII, th\ptbl 5]{XI.5} shows that this is an
injective homomorphism whose image is the set of \og primitive
elements\fg of~$\H^1(A,G)$, \ie the elements $\xi$ for which
$$
\pi^*(\xi) = \pr_1^*(\xi) + \pr_2^*(\xi)\quoi,
$$
where the $\pr_i$ are the two projections from~$A\times A$ onto~$A$,
and $\pi\colon A\times A\to A$ is the composition law of~$A$
(N.B. Serre states his theorem only for $G$ linear and connected, and
of course smooth over~$k$, but, on simplifying the first part of his
argument, one sees that these restrictions are unnecessary: it is
enough to note that every morphism from~$A$ into a group scheme $E$
of finite type over~$k$ which takes the identity to the identity is a
group homomorphism, and to apply this to the sections over~$A$ of an
extension $E$ of~$A$ by~$G$).

We shall apply this result when $G$ is a finite separable group scheme
over~$k$, \ie an ordinary finite group, assumed commutative. Using
then $\pi_1(A\times A)\isomto\allowbreak \pi_1(A)\times
\pi_1(A)$ (\ifthenelse{\boolean{orig}}{X~}{}\Ref{X.1.7}) and interpreting $\H^1(X,G)$ as
$\Hom(\pi_1(X),G)$ for every algebraic scheme~$X$, in particular for
$X=A$ or $X=A\times A$, one sees that every class in~$\H^1(A,G)$ is
primitive, and hence there is an isomorphism
$$
\Ext(A,G)\isomto \H^1(A,G)\quoi,
$$
in other words, \emph{every principal covering of~$A$ with
commutative structural group~$G$, pointed above the origin of~$A$, is
endowed in a unique way with the structure of an algebraic group
having the marked point as its origin and such that $A'\to A$ is a
homomorphism of algebraic groups.} In particular, if $A'$ is
connected, it too is an abelian variety, isogenous to~$A$.

On the other hand, since the functor $X\mto \pi_1(X)$ from pointed
algebraic schemes $X$ to groups commutes with products
(\ifthenelse{\boolean{orig}}{IX~}{}\Ref{IX.1.7}), it carries a group object in the first category to
a group object in the category of groups, \ie to a \emph{commutative}
group. Thus, \emph{if $A$ is an abelian variety, $\pi_1(A)$
% page: 288
\marginpar{288}
%
is a commutative group.} Therefore, to determine $\pi_1(A)$, it
suffices to know the functor $G\mto
\H^1(A,G)=\Hom(\pi_1(A),G)$ as $G$ ranges over the finite
\emph{commutative} groups. Finally, recall that for every integer
$n>0$, the multiplication-by-$n$ homomorphism on~$A$,
$$
A\lto{n} A
$$
is surjective and hence has finite kernel, \ie it is an isogeny, and
it follows that every isogeny $A'\to A$ is a quotient of an isogeny
of the preceding type. From this and standard arguments (\cf for
example~\cite{XI.6}), one obtains:

\begin{theoreme}[Serre-Lang]
\label{XI.2.1}
Let $A$ be an abelian variety over an algebraically closed field~$k$.
For every integer $n>0$, consider the ordinary finite group $K_n$
underlying the kernel ${}_nA$ of multiplication by $n$ on~$A$, and,
for every prime number~$\ell$, set
$$
T_\ell(A)=\varprojlim_r K_{\ell^r}
$$
and
$$
T_{.}(A)=\underset{\ell}{\prod}T_\ell(A)=\varprojlim_n K_n
$$
(where, if $m$ is a multiple of~$n$, $m=ns$, the map from $K_m$ to
$K_n$ is multiplication by~$s$). Then the group $\pi_1(A)$ is
canonically isomorphic to $T_{.}(A)$; consequently, for every prime
number $\ell$, the $\ell$-primary component of~$\pi_1(A)$ is
canonically isomorphic to $T_\ell(A)$.
\end{theoreme}

These isomorphisms are functorial as $A$ varies. The module
$T_\ell(A)$ is called the \emph{$\ell$-adic Tate module} of the
abelian variety~$A$. It is an additive functor in~$A$; in particular,
it gives rise to a representation of the ring $\Hom(A,A)$ of
endomorphisms of~$A$ on $T_\ell(A)$, called the \emph{$\ell$-adic
Weil representation}, which plays an important role in the theory of
abelian varieties (\cf for example \hbox{\cite[chap~VII]{XI.4}}).
Theorem~\Ref{XI.2.1} interprets it as the natural representation on
the \emph{$\ell$-adic homology group} of~$A$,
$\H_1(A,\ZZ_\ell)=\pi_1(A)_\ell$, which
% page: 289
\marginpar{289}
%
is clearly more satisfactory a priori, especially from the point of
view of the Lefschetz formula, \cite[Chap~V]{XI.4}. Let us recall here
Weil's results on the structure of~$T_\ell(A)$:
\begin{enumerate}
\item[a)] If $n$ is prime to $\car(k)$, then $K_n$ is a free module
of rank $2\dim A$ over~$\ZZ/n\ZZ$; hence, if $\ell$ is a prime number
$\neq \car(k)$, $T_\ell(A)$ is a free module of rank $2\dim A$ over
the ring $\ZZ_\ell$ of $\ell$-adic integers;
\item[b)] If $n$ is a power of~$\car(k)=p$, then $K_n$ is a free
module of rank $\nu\leq\dim A$ over~$\ZZ/n\ZZ$, with $\nu$
independent of~$n$; hence $T_p(A)$ is a free module of rank
$\nu\leq \dim A$ over the ring $\ZZ_p$ of $p$-adic integers.
\end{enumerate}

This shows that, in the theory of the fundamental group developed
here, the fundamental group of a varying abelian variety does not
vary regularly with the parameter: its $p$-primary component may
drop abruptly at values of the parameter $t$ corresponding to
residual characteristic~$p$. The best-known case of this phenomenon
is that of elliptic curves. Note, however, that, whatever $n$ may be
(whether or not it is prime to the characteristic), the actual kernel
${}_nA$ in~$A$ of multiplication by~$n$ is a finite group scheme over
$k$ of degree $n^{2g}$, where $g=\dim A$, and it is nonseparable over
$k$ if $n$ is a multiple of~$p=\car(k)$. Moreover, when $A$ varies in
a family of abelian varieties, \ie when there is an abelian scheme
$A$ over a base scheme $S$, one proves more generally that ${}_n A$
is a finite flat group scheme over~$S$, of degree $n^{2g}$ over~$S$;
that is, provided one takes the infinitesimal parts of the kernels
${}_n A$ into account, they behave regularly for every~$n$. This
suggests that the \og true\fg fundamental group of an abelian variety
$A$ is the pro-algebraic group (the formal inverse limit of finite
groups over~$k$) $\varprojlim_n{{}_nA}$, where by the \og true
fundamental group\fg of an algebraic scheme~$X$ one means the
pro-group that classifies the principal coverings of~$X$ whose
structural group is an arbitrary finite group $G$ over~$k$ (not
necessarily separable over~$k$). In this way, for example, the
representations of~$\Hom(A,A)$ on the $p$-primary component of the
true fundamental group of~$A$ recover the Weil characteristic
polynomial, defined by Weil using the primes $\ell\neq p$, more naturally
than does Serre's construction~\cite{XI.8}.
\end{filecontents*}
% Original input: drafts/SGA1_XI_3_English.texfrag
\begin{filecontents*}[overwrite,noheader]{sga1-cumulative-input-127.tex}
\section{Projective cones, an example of Zariski}
\label{XI.3}
% page: 290
\marginpar{290}
%

Continue to assume, for simplicity, that $k$ is an algebraically closed
field, and let $V$ be a connected projective $k$-scheme, a closed subscheme
of~$\PP^r_k$, which may, if desired, be assumed nonsingular. Let $Y=\hat C$
be the projective cone over~$V$, let $y_0$ be its vertex, let $X=\hat C_V$ be
the usual projective completion of the vector bundle
$C_V=\VV(\cal{O}_V(1))$ associated with $\cal{O}_V(1)$, and finally let
$$
f\colon X\to Y
$$
be the canonical morphism, contracting the zero section $X_0$ of~$C_V$
in~$X$ to a point (EGA~II~8.6.4). Since $X$ is a locally trivial bundle
over~$V$, with fibers $\PP^1$ and hence simply connected fibers, the morphism
$p\colon X\to V$ induces by
\ifthenelse{\boolean{orig}}
{(XI~\Ref{XI.4}.)}
{(\ifthenelse{\boolean{orig}}{X~}{}\Ref{X.1.4})}
% correction: the reference should be to `X 1.4'.
an isomorphism:
$$
\pi_1(X)\isomto\pi_1(V)\quoi.
$$
Since $p$ induces an isomorphism $X_0\to V$, we conclude that \emph{an
\'etale covering of~$X$ is completely split if and only if its restriction
to $X_0$ is completely split}. Now, for every \'etale covering $Y'$ of~$Y$,
the inverse image $X'=X\times_Y Y'$ is an \'etale covering of~$X$ that is
completely split over the fiber~$X_0$, and hence is trivial. Since the
homomorphism $\pi_1(X)\to\pi_1(Y)$ is surjective (\ifthenelse{\boolean{orig}}{IX~}{}\Ref{IX.3.4}), we
conclude that
$$
\pi_1(Y)=(e)
$$
in other words, \emph{every projective cone is simply connected.} (N.B. in
characteristic~$0$, the same result holds when $Y$ is taken to be the affine
cone).

Now suppose that $V$ is regular, \ie smooth over~$k$; then $X$ is regular,
and for a suitable projective embedding of~$V$, one obtains a \emph{normal}
projective cone $Y$. If $V$ is not simply connected, and hence $X$ is not
simply connected, let $X'$ be a nontrivial connected \'etale covering
of~$X$. Since the fibers of~$X$ over points $y\in Y$ distinct from $y_0$ are
reduced to a point, we see that the restriction of $X'$ to these fibers (in
particular, to the generic
fiber)
% page: 291
\marginpar{291}
%
is trivial; nevertheless, \ifthenelse{\boolean{orig}}{$X$}{$X'$} is not obtained as the inverse image of an
\'etale covering of~$Y$, since $Y$ is simply connected and $X'$ would be
completely split. This shows that (\ifthenelse{\boolean{orig}}{X~}{}\Ref{X.1.3} and~\Ref{X.1.4}) become
false if the hypothesis that $f$ is separable is replaced by the weaker
hypothesis that its fibers are separable algebraic schemes (or even smooth
ones) over the~$\kres(s)$. Likewise, note that the fundamental groups of the
geometric fibers $\overline X_y$ for $y\neq y_0$ are obviously reduced to
$(e)$, since these fibers are reduced to a point, whereas
$\pi_1(X_0)\neq e$; thus the semicontinuity theorem (\ifthenelse{\boolean{orig}}{X~}{}\Ref{X.2.4}) likewise
fails for~$f$.

Finally, let us give the example, pointed out by Zariski, in which these same
theorems fail when the hypothesis that $f$ is separable is replaced by the
hypothesis that $f$ is flat. Let $f\colon X\to Y$ be a morphism from a
nonsingular projective surface to the rational line $Y=\PP^1$, such that
\ifthenelse{\boolean{orig}}
{$K=k(x)$}
{$K=k(X)$}
% correction: `k(x)' should be `k(X)'.
is a \og regular\fg extension of~$k(f)$, \ie primary and separable (\ie the
geometric generic fiber is connected and separable), and such that the
divisor $(f)=X_0-X_\infty$ is an
\(n\)-fold
% correction : write `$n$-i\`eme' instead of `$n$.\`eme'
multiple of a divisor (where $n$ is an integer prime to the characteristic).
It is possible to construct such examples in every characteristic. Let $X'$
be the normalization of~$X$ in $K(f^{1/n})$, where $K=k(X)$ is the function
field of~$X$. The hypothesis on~$(f)$ implies that $X'$ is \'etale over~$X$.
Let $Y'$ be the normalization of~$Y$ in $k(t)(t^{1/n})$; it is ramified
over~$Y$ exactly at the points $t=0$ and $t=\infty$\ifthenelse{\boolean{orig}}{, and the restriction}{. Put $U=Y-\{0,\infty\}$. The restriction}
$X'|f^{-1}(U)$ is isomorphic to the inverse image of~$Y'|U$. In particular,
the restriction of~$X'$ to the \emph{geometric} generic fiber of~$X$ over~$Y$
is completely split. Nevertheless, $X'$ is not isomorphic to the inverse
image of an \'etale covering of~$Y$, since one sees at once that the latter
would necessarily be~$Y'$, which is absurd because $Y'$ is ramified
over~$Y$\kern2pt\footnote{One may observe, from the viewpoint of the
\og \'etale topology\fg (SGA~4~VII), that in this example
$\R^1(f_{\et})_*(\ZZ/n\ZZ)$ is \og not separated\fg over~\ifthenelse{\boolean{orig}}{$S$}{$Y$}.}.

Here is (following Serre) a simple way to realize the conditions of this
example, inspired by \cite[\No 20]{XI.5}: take $n$ to be a prime number
$\geq 5$, distinct from the characteristic, and let $G=\ZZ/n\ZZ$ act on
$\underline{k}^4$%
\ifthenelse{\boolean{orig}}
{}
{{\renewcommand{\thefootnote}{*}\addtocounter{footnote}{-1}\kern1pt\footnote{\lcrochetbf Added in 2003: $\underline{k}^4$ denotes affine space of dimension $4$ over $k$.\rcrochetbf}}}
by multiplying the coordinates by four distinct characters of~$G$ (which is
possible since $n\geq 5$). Then $G$ acts on projective space $\PP^3_k$, and
the only points fixed by $G$ are the four points corresponding to the
coordinate axes. The surface $X'$ with equation
\ifthenelse{\boolean{orig}}
{$x^4+y^4+z^4+t^4=0$}
{$x^n+y^n+z^n+t^n=0$}
is smooth over
% page: 292
\marginpar{292}
%
$k$ (Jacobian criterion), and contains none of the fixed points; thus, since
$G$ has prime order, it acts on~\ifthenelse{\boolean{orig}}{$X$}{$X'$} \og without fixed points\fg, \ie \ifthenelse{\boolean{orig}}{$X$}{$X'$} is a
principal covering of~$X=X'/G$ with group $G$. Let $g=x/y$ in
$k(X')=K'$; it is a Kummer generator of~$K'$ over~$K=k(X)$ if the chosen
characters were~$\chi^i$, $i=0,1,2,3$, with $\chi$ a primitive character.
Let $f$ be its \(n\)-th power, which is an element of~$K$. We see at once that
$K'$ is a regular extension of~$k(g)$; this follows from the fact that the
plane curve with homogeneous equation in $U,T,Z$,
$T^n+Z^n+(1+g^n)U^n=0$, is smooth over~$k(g)$ (Jacobian criterion), and from
the fact that every plane curve is known to be connected. On the other hand,
$k(f)=K\cap k(g)$, since the right-hand side is an extension of $k(f)$
contained in $k(g)$, an extension of prime degree, and is different
from~$k(g)$ (since $g\not\in K$). This implies that $K$ is a regular extension
of~$k(f)$. Finally, the divisor of~$f$ on~$X$ is an \(n\)-fold multiple of a
divisor, since its inverse image on~$X'$ is the divisor of~$g^n$, and hence an
\(n\)-fold multiple, and one can descend because $X'$ is \'etale over~$X$. We
would be done if the rational map $f\colon X\to \PP^1$ were a morphism, that
is, if the divisors of the zeros and poles of~$f$ did not meet. In fact, one
checks easily (again by looking on~$X'$) that the two divisors in question are
the multiples by $n$ of two smooth curves over~$k$, meeting transversely at a
point~$a$. If we now replace $X$ by the scheme obtained by blowing up~$a$,
the preceding conditions ($\divisor(f)$ divisible by $n$, and
$k(X_1)=k(X)$ a regular extension of~\ifthenelse{\boolean{orig}}{$k(X)$}{$k(f)$}) remain satisfied; moreover, $f$
is a \emph{morphism} $X_1\to \PP^1$, so the required conditions hold.
\end{filecontents*}
% Original input: drafts/SGA1_XI_4_A_English.texfrag
\begin{filecontents*}[overwrite,noheader]{sga1-cumulative-input-128.tex}
\section{The cohomology exact sequence}
\label{XI.4}

Let \(S\) be a prescheme. The category \(\Sch_{/S}\) of preschemes
over~\(S\) is then determined, as is the notion of a group in this
category; such a group will also be called a \emph{group prescheme
over~\(S\)}, or simply an \emph{\(S\)-group}. To simplify the
exposition and fix ideas, we shall most often restrict ourselves
below to groups which are \emph{affine} and \emph{flat} over~\(S\)
\footnote{In fact, for what follows, quasi-affine instead of affine
would suffice; see footnote~\eqref{note-296},
p\ptbl\pageref{note-296}.}. This will suffice for the applications
we have in mind, although of course many cases arise in which neither
hypothesis holds.

Let
% page: 293
\marginpar{293}
%
\(G\) be such an \(S\)-group, and let \(P\) be a prescheme over~\(S\)
on which \(G\) acts on the right. In particular, there is a morphism
\[
  \pi\colon P\times_SG\longrightarrow P
\]
in the category \(\Sch_{/S}\), satisfying the familiar axioms. We say
that \(P\) is \emph{formally principal homogeneous under}~\(G\) if
the morphism
\[
  P\times_SG\longrightarrow P\times_SP
\]
whose components are \(\pr_1\) and~\(\pi\) is an isomorphism.
Equivalently, for every object \(S'\) of~\(\Sch_{/S}\), the set
\[
  P(S')=\Hom_S(S',P),
\]
regarded as a set with operator group
\[
  G(S')=\Hom_S(S',G),
\]
is either empty or principal homogeneous, \ie either empty or
isomorphic to \(G(S')\) with \(G(S')\) acting by right translations.

We say that \(P\) is \emph{trivial} if it is isomorphic to~\(G\), on
which \(G\) acts by right translations, or equivalently if each of
the operator sets \(P(S')\) under \(G(S')\) is trivial. One verifies,
for example by the standard procedure of passing to the set-theoretic
case, that \(P\) \emph{is trivial if and only
\ifthenelse{\boolean{orig}}
{if it}
{if it}
is formally principal homogeneous and has a section over~\(S\)}. The
latter condition may be stated categorically by saying that \(P\)
has a section over the final object \(e=S\) of~\(\Sch_{/S}\), \ie
that there is a morphism from~\(e\) to~\(P\).

To define a principal homogeneous bundle \(P\) under~\(G\), a notion
stronger than that of a formally principal homogeneous bundle, one
must first specify in~\(\Sch_{/S}\) a class of morphisms to be used
for descent and to play the role of localization morphisms that
trivialize bundles. The most suitable choice varies with the context,
and no one choice contains all the others\footnote{\label{note-293}
See SGA~3, IV, especially~\S4.}. Here it will be convenient to adopt
the following definition:

\begin{definition}
\label{XI.4.1}
Let \(G\) be an \(S\)-group. A \emph{principal homogeneous bundle}
on the right under~\(G\) is an \(S\)-prescheme \(P\) with a right
action of the \(S\)-group~\(G\), such that there exist a covering of
\(S\) by open subsets \(U_i\), and for every~\(i\) a faithfully flat
and quasi-compact base-change morphism \(S_i'\to U_i\), for which
\ifthenelse{\boolean{orig}}
{\(P_i'=P\times_SS_i'\) is a trivial prescheme with operators under
\(G'=G\times_SS'\).}
{\(P'=P\times_SS'\) is a trivial prescheme with operators under
\(G'=G\times_SS'\), where \(S'\) is the \(S\)-prescheme given by
the disjoint sum of the \(S_i'\).}
\end{definition}

Observe that
% page: 294
\marginpar{294}
%
the base-change functor \(X\mto X'=X\times_SS'\), being left exact,
sends groups to groups and objects with operator groups to objects
with operator groups. Definition~\Ref{XI.4.1} is stable under base
change. We also have:

\begin{proposition}
\label{XI.4.2}
Let \(G\) be an \(S\)-group, flat and quasi-compact over~\(S\), and
let \(P\) be an \(S\)-prescheme on which \(G\) acts on the right.
The following conditions are equivalent:
\begin{enumerate}
\item[(i)] \(P\) is a principal homogeneous bundle
\ifthenelse{\boolean{orig}}
{under~\(G\).}
{under~\(G\);}
\item[(ii)] \(P\) is formally principal homogeneous under~\(G\),
%
\index{formally principal homogeneous (prescheme)|hyperpage}%
%
and the structural morphism \(P\to S\) is faithfully flat and
quasi-compact.
\end{enumerate}
\end{proposition}

If \(P\) is principal homogeneous under~\(G\), then, with the notation
of~\Ref{XI.4.1}, \(P'\) is faithfully flat and quasi-compact
over~\(S'\), since \(G'\) is so and \(P'\) is \(S'\)-isomorphic
to~\(G'\). Hence \(P\) has the same properties over~\(S\). For
surjectivity and quasi-compactness, see \ifthenelse{\boolean{orig}}{VIII~}{}\Ref{VIII.3.1}; flatness
was omitted from the sorites of Expos\'e~\Ref{VIII}. Conversely, if
(ii) holds, take the base change \(S'=P\), which is faithfully flat
and quasi-compact over~\(S\). Then \(P'\) is formally principal
homogeneous over~\(S'\), since \(P\) is so over~\(S\) and the
base-change functor is left exact. Moreover, \(P'\) has a section
over~\(S'\),
\ifthenelse{\boolean{orig}}
{namely}
{namely}
the diagonal section. It is therefore a trivial principal bundle,
which completes the proof.

\begin{corollary}
\label{XI.4.3}
If \(G\) is affine and flat over~\(S\), every principal homogeneous
bundle \(P\) under~\(G\) is affine and flat over~\(S\).
\end{corollary}

Indeed, it becomes so after a faithfully flat and quasi-compact
extension of the base, and \ifthenelse{\boolean{orig}}{VIII~}{}\Ref{VIII.5.6} applies.

The usefulness of Definition~\Ref{XI.4.1} for \(S\)-groups
\emph{flat} and \emph{affine} over~\(S\) rests on
\ifthenelse{\boolean{orig}}{VIII~}{}\Ref{VIII.2.1}: the morphisms \(S'\to S\) considered in
\Ref{XI.4.1} are effective descent morphisms for the category of
preschemes affine over other preschemes. Thanks to this fact, the
verification of the statements outlined below is essentially
categorical\footnote{\Cf \loccit in footnote~\eqref{note-293},
p\ptbl\pageref{note-293}.}.

Let \(E\) be an \(S\)-prescheme on which the \(S\)-group~\(G\) acts
on the left, and let \(P\) be a principal homogeneous bundle on the
right under~\(G\). We wish to define an associated bundle
\(E^{(P)}\), locally
% page: 295
\marginpar{295}
%
isomorphic to~\(E\). Let \(G\) act on the right on \(P\times_SE\)
by
\[
  (x,y)\mto(xg,g^{-1}y),
\]
the formula describing such actions in the set-theoretic context,
and extended to categories by the standard procedure. Subject to its
existence, set
\[
  E^{(P)}=(P\times_SE)/G.
\]
Then \(P\times_SE\) is a prescheme over \(T=E^{(P)}\) with right
operator group \(G_T=G\times_ST\). To work comfortably, we would
like \(P\times_SE\) to be, in addition, a principal homogeneous
bundle over~\(T\) with group~\(G_T\).

To verify the existence of~\(E^{(P)}\) and this additional property,
take the \(S'\) of Definition~\Ref{XI.4.1} and pass to the inverse
image over~\(S'\) of the original situation. Since \(P'\) is
\ifthenelse{\boolean{orig}}
{trivial}
{trivial,}
\ie isomorphic to \(G'_d\), one sees immediately that
\(E^{\prime(P')}\) exists and has the required exactness property.
Indeed, \(E'\times_{S'}P'\) is \(G'\)-isomorphic to
\(E'\times_{S'}G'\), so \(E^{\prime(P')}\) is isomorphic to~\(E'\).

Moreover, formation of the associated bundle in the case of a
trivial object with operators commutes with every base extension.
Applying this to the various base extensions
\[
  \xymatrix@C=.5cm{
  S''\ar@<2pt>[r]\ar@<-2pt>[r]&S'
  }
  \qquad\text{and}\qquad
  \xymatrix@C=.5cm{
  S'''\ar@<4pt>[r]\ar@<0pt>[r]\ar@<-4pt>[r]&S',
  }
\]
where \(S''\) and \(S'''\) are the double and triple fiber products
of~\(S'\) over~\(S\), shows that \(E^{\prime(P')}\) \emph{is endowed
with a descent datum relative to \(S'\to S\), and that \(E^{(P)}\)
exists with the required properties if and only if this descent datum
is effective}. In that case \(E^{(P)}\) is, of course, precisely the
descended object. Here one uses the fact that \(S'\to S\) is a
descent morphism in the category of \(S\)-preschemes; see
\ifthenelse{\boolean{orig}}{VIII~}{}\Ref{VIII.5.2}. It follows that \emph{the associated bundle
exists if \(E\) is affine over~\(S\)}.

We shall apply this construction when an \(S\)-group homomorphism
\(G\to H\) is given, taking for~\(E\) the \(S\)-prescheme~\(H\) with
the left action of~\(G\) induced by the given morphism. The group
\(H\) acts on itself on the right in a way that commutes with the
action of~\(G\) on~\(H\), and formation of the associated bundle,
whenever it exists over~\(S\), commutes with base extension. It
follows readily that \(H\) acts on the right on~\ifthenelse{\boolean{orig}}{\(P^{(H)}\)}{\(H^{(P)}\)}, which is
therefore a principal homogeneous bundle under~\(H\) in the sense
of~\Ref{XI.4.1}, and
% page: 296
\marginpar{296}
%
is in fact trivialized by the same morphism \(S'\to S\) as~\(P\).
In particular, \emph{to every principal homogeneous bundle \(P\)
under~\(G\) and every homomorphism of \(S\)-groups \(G\to H\), with
\(H\) affine over~\(S\), there is functorially associated a principal
homogeneous bundle with group~\(H\)}. This construction is functorial
in \((G\to H)\) and compatible with arbitrary base changes
\(T\to S\).

\end{filecontents*}
% Original input: drafts/SGA1_XI_4_B_English.texfrag
\begin{filecontents*}[overwrite,noheader]{sga1-cumulative-input-129.tex}
\begin{definition}
\label{XI.4.4}
Let \(G\) be an \(S\)-prescheme. We denote by \(\H^0(S,G)\) the set
of sections of~\(G\) over~\(S\), regarded as a group when \(G\) is
an \(S\)-group. In that case, \(\H^1(S,G)\) denotes the set of
isomorphism classes of principal homogeneous bundles over~\(S\) with
group~\(G\), endowed with the distinguished point corresponding to
trivial bundles\footnote{\label{note-296}This notation is compatible
\ifthenelse{\boolean{orig}}
{with}
{with}
the general cohomological notation (SGA~4, V) only when descent
effectivity criteria are available. These are hardly assured unless
\(G\) is affine, or merely quasi-affine; see
\ifthenelse{\boolean{orig}}{VIII~}{}\Ref{VIII.7.9}.}.
\end{definition}

Thus \(\H^0(S,G)\) is a functor in the \(S\)-prescheme~\(G\), with
values in the category of sets. This functor is left exact, and a
fortiori commutes with finite products; it therefore sends groups to
groups and commutative groups to commutative groups. Similarly, by
formation of associated bundles, \(\H^1(S,G)\) is a functor in the
\emph{affine} \(S\)-group~\(G\), with values in the category of
pointed sets, and one readily verifies that it commutes with finite
products. In particular, it sends groups in the category of affine
\(S\)-groups, \ie \emph{commutative affine \(S\)-groups}, to groups
in the category of pointed sets, and indeed to commutative groups,
since the group objects in the first category are commutative. Thus,
\emph{if \(G\) is a commutative affine \(S\)-group, then
\(\H^1(S,G)\) is a commutative group}, and a homomorphism
\(G\to H\) of commutative affine \(S\)-groups gives rise to a group
homomorphism
\[
  \H^1(S,G)\longrightarrow\H^1(S,H).
\]

For simplicity, in what follows we restrict ourselves to
\emph{affine commutative} \(S\)-groups. Let
\[
  0\longrightarrow G'\lto{u}G\lto{v}G''\longrightarrow0
\]
be a sequence of morphisms of such groups. \emph{We shall say that
this sequence is exact if \(vu=0\)}, which allows \(G\) to be regarded
as a prescheme over~\(G''\) with right operator group
\(G'_{G''}\), \emph{and if \(G\) is a principal homogeneous bundle
over~\(G''\) with group
% page: 297
\marginpar{297}
%
\(G'_{G''}=G'\times_SG''\)}. This implies in particular that
\(u\colon G'\to G\) is a kernel of~\(v\), and hence that the
sequence
\[
  0\longrightarrow\H^0(S,G')
  \longrightarrow\H^0(S,G)
  \longrightarrow\H^0(S,G'')
\]
is exact. It also allows one to define a map
\[
  \partial\colon\H^0(S,G'')\longrightarrow\H^1(S,G')
\]
as follows. To every section of~\(G''\) over~\(S\), \ie every
\(S\)-morphism \(f\colon S\to G''\), associate the principal
homogeneous bundle \(P_f\) over~\(S\), with group
\(G'\simeq f^*(G'_{G''})\), obtained by inverse image of the
principal homogeneous bundle \(G\) over~\(G''\). As an
\(S\)-prescheme, this is simply the inverse image under
\(v\colon G\to G''\) of the subprescheme which is the image of~\(S\)
under the immersion~\(f\); the action of~\(G'\) on~\(P_f\) is induced
by the right action of~\(G'\)
\ifthenelse{\boolean{orig}}
{on~\(G\)).}
{on~\(G\).}
% correction: remove the extra closing parenthesis.

We also leave to the reader the verification of the following
proposition, which presents no difficulty apart from its formulation:

\begin{proposition}
\label{XI.4.5}
The map
\(\partial\colon\H^0(S,G'')\to\H^1(S,G')\) is a group
homomorphism. The following sequence of homomorphisms is exact:
\begin{multline*}
0\longrightarrow\H^0(S,G')
\lto{u^0}\H^0(S,G)
\lto{v^0}\H^0(S,G'')
\lto{\partial}\H^1(S,G')\\
\lto{u^1}\H^1(S,G)
\lto{v^1}\H^1(S,G'').
\end{multline*}
All the homomorphisms other than~\(\partial\) arise from the
functoriality of \(\H^0\), respectively \(\H^1\).
\end{proposition}

\begin{remarques}
\label{XI.4.6}
The point of view presented here for the study of principal
homogeneous bundles is visibly inspired by Serre~\cite{XI.7}, which
the reader is strongly advised to consult. For a formalism that also
applies to structural \(S\)-groups quasi-projective over~\(S\), so
as in particular to include projective abelian schemes, it is
preferable to modify~\Ref{XI.4.1} by requiring \(S'\) to be a
disjoint sum of preschemes \(S_i'\) which are finite and locally free
over open subsets \(S_i\) covering~\(S\). The preceding developments,
including in particular~\Ref{XI.4.5}, then remain valid after
replacing the affine hypothesis everywhere by the quasi-projective
hypothesis, and interpreting accordingly the definition given above
of an exact sequence of \(S\)-groups. Indeed, one need only replace
the reference to \ifthenelse{\boolean{orig}}{VIII~}{}\Ref{VIII.2.1}
% page: 298
\marginpar{298}
%
by \ifthenelse{\boolean{orig}}{VIII~}{}\Ref{VIII.7.7}: the morphisms \(S'\to S\) used here are
effective descent morphisms for the fibered category of preschemes
quasi-projective over other preschemes. One should note, however, that
this second notion of principal homogeneous bundle is more restrictive
than the first one in~\Ref{XI.4.1}.
\end{remarques}

\subsection{}
\label{XI.4.7}
An even more restrictive notion of principal homogeneous bundle is
obtained by requiring \(S\) to be covered by open subsets \(S_i\)
such that, for every~\(i\), \(P|S_i\) is a trivial bundle with
operators under~\(G|S_i\). We then say that \(P\) is a
\emph{locally trivial principal homogeneous bundle}.
%
\index{locally trivial principal homogeneous bundle|hyperpage}%
%
For fixed~\(G\), the classes of these bundles form a subset of
\(\H^1(S,G)\), in bijection with
\(\H^1(S,\cal O_S(G))\), where \(\cal O_S(G)\) is the sheaf, in the
ordinary sense, of sections of~\(G\) over~\(S\); see~\cite{XI.2}.
For these \(\H^1\) sets still to give rise to the cohomology exact
sequence~\Ref{XI.4.5}, one must plainly suppose that
\[
  0\longrightarrow G'\longrightarrow G\longrightarrow G''
  \longrightarrow0
\]
is exact in the appropriate sense for this new context: \(G\) must
be a locally trivial bundle over~\(G''\) with group \(G'_{G''}\).
Equivalently, \(u\colon G'\to G\) is a kernel of
\(v\colon G\to G''\), and \(G\) locally admits a section over~\(G''\).

\subsection{}
\label{XI.4.8}
It is plainly very desirable to continue the exact
sequence~\Ref{XI.4.5} by introducing the higher cohomology groups
\(\H^i(S,G)\). This is possible in the framework of \emph{Weil
cohomology}. Consider the category \(\cal B\) of preschemes
quasi-compact over~\(S\), equipped with the class \(\cal M\) of
faithfully flat and quasi-compact morphisms, called localizing
morphisms. An abelian \emph{Weil sheaf} on~\(S\), or more precisely
on \((\cal B,\cal M)\), is a contravariant functor \(\cal F\) from
\(\cal B\) to the category of abelian groups, sending sums to
products and every sequence
\[
  {\def\labelstyle{\scriptstyle}
  \xymatrix@C=.5cm{
  T''=T'\times_TT'\ar@<+2pt>[rr]^-{\pr_1,\pr_2}
  \ar@<-2pt>[rr]&&T'\ar[r]^f&T
  }},
  \qquad f\in\cal M,
\]
to an \emph{exact} diagram of sets
\[
  \xymatrix@C=.5cm{
  \cal F(T)\ar[r]&
  \cal F(T')\ar@<2pt>[r]\ar@<-2pt>[r]&
  \cal F(T'').
  }
\]
Weil sheaves form an abelian category with exact direct limits and a
generator, and hence with enough injective objects~\cite{XI.1}. The
right derived functors of
\[
  \Gamma(\cal F)=\cal F(S)
\]
are denoted by \(\H^i(S,\cal F)\). Moreover, every commutative
\(S\)-group plainly defines a Weil sheaf
(\ifthenelse{\boolean{orig}}{VIII~}{}\Ref{VIII.5.2}), whose \(\H^0\) and \(\H^1\) are precisely
\(\H^0(S,G)\) and \(\H^1(S,G)\). This gives a reasonable definition
of the remaining groups \(\H^i(S,G)\). One
% page: 299
\marginpar{299}
%
also proves that an exact sequence of \(S\)-groups defines an exact
sequence of Weil sheaves, recovering and extending the exact
sequence~\Ref{XI.4.5}\footnote{For a systematic study of this point
of view, \cf SGA~4, I through~IX.}.

\subsection{}
\label{XI.4.9}
It would be appropriate to develop the noncommutative variants
of~\Ref{XI.4.5}, as in~\cite{XI.2}. For a systematic development, in
the proper framework, of the various cohomological notions outlined
in this section, we refer to a work then in preparation by
J\ptbl \textsc{Giraud}\footnote{\Cf J\ptbl \textsc{Giraud},
\ifthenelse{\boolean{orig}}
{\emph{Alg\`ebre homologique non ab\'elienne}, forthcoming in}
{\emph{Cohomologie non ab\'elienne},}
Springer-Verlag, 1971.}.

\end{filecontents*}
% Original input: drafts/SGA1_XI_5_English.texfrag
\begin{filecontents*}[overwrite,noheader]{sga1-cumulative-input-130.tex}
\section{Special cases of principal bundles}
\label{XI.5}

Suppose now that \(S\) is connected and endowed with a geometric
point~\(a\). Its fundamental group \(\pi_1(S,a)\) classifies the
\'etale coverings of~\(S\): the category of \'etale coverings of~\(S\)
is equivalent to the category of finite sets on which \(\pi_1\) acts
continuously. It follows that a finite \'etale group scheme \(G\)
over~\(S\) is determined essentially by an ordinary finite group
\(\cal G\), on which \(\pi_1\) acts continuously by group
automorphisms. An \'etale covering \(P\) of~\(S\) on which \(G\)
acts on the right is determined essentially by a finite set \(\cal P\)
on which \(\pi_1\) acts continuously on the left, and on which
\(\cal G\) acts on the right compatibly with the action of~\(\pi_1\):
\[
  s(p\cdot g)=(sp)\cdot(sg)
  \qquad
  \text{for }s\in\pi_1,\ p\in\cal P,\ g\in\cal G.
\]
One verifies that \(P\) is a principal homogeneous bundle in the
sense of~\Ref{XI.4.1} if and only if \(\cal P\) is a principal
homogeneous set under~\(\cal G\), for example by using the criterion
\Ref{XI.4.2}. In other words, \emph{the category of principal
homogeneous bundles over~\(S\) with group~\(G\) is equivalent to the
category of principal homogeneous bundles with group~\ifthenelse{\boolean{orig}}{\(G\)}{\(\cal G\)} in the
category of finite sets on which \(\pi_1\) acts continuously}. In
particular, one obtains a canonical bijection, functorial in~\(G\):
\begin{equation*}
\label{eq:XI.5.etoile}
\tag{\(*\)}
  \H^1(S,G)\isomto\H^1(\pi_1,\cal G).
\end{equation*}
Here
% page: 300
\marginpar{300}
%
the right-hand side denotes the set of isomorphism classes of
principal homogeneous bundles under~\(\cal G\) in the category of
finite sets on which \(\pi_1\) acts; there is no need to specify
continuity. As is well known, this set is the quotient of the set
\(Z^1(\pi_1,\cal G)\) of \(1\)-cocycles
\(\varphi\colon\pi_1\to\cal G\), satisfying
\[
  \varphi(1)=1,
  \qquad
  \varphi(st)=\varphi(s)\bigl(s\cdot\varphi(t)\bigr),
\]
by the natural action of~\(\cal G\).

\label{page-300}
An important case is that in which \(\pi_1\) acts trivially on
\(\cal G\), \ie when \(G\) is a completely split covering of~\(S\),
isomorphic to the disjoint sum of \(\cal G\) copies of~\(S\). One
then also writes \(\H^1(S,\cal G)\) in place of \(\H^1(S,G)\), and
the bijection~\eqref{eq:XI.5.etoile} identifies this set with
\[
  \H^1(\pi_1,\cal G)
  =\Hom(\pi_1,\cal G)/
  \{\text{inner automorphisms of }\cal G\}.
\]
In this case, a principal homogeneous bundle over~\(S\) with
group~\(G\) is nothing other than a \emph{principal covering} of~\(S\)
with group~\(\cal G\) (\ifthenelse{\boolean{orig}}{V~}{}\Ref{V.2.7}). The preceding bijection is the
one deduced from the correspondence between principal coverings
of~\(S\) with group~\(\cal G\), \emph{pointed} above~\(a\), and
continuous homomorphisms from \(\pi_1(S,a)\) to~\(\cal G\); see the
end of V, \No \Ref{V.5}.

The usefulness of connecting the theory of \'etale coverings with
that of principal bundles, already implicit in A\ptbl Weil,
\emph{G\'en\'eralisation des Fonctions Ab\'eliennes}, and made
explicit by S\ptbl Lang in his geometric class-field theory; see
Serre~\cite{XI.5}, stems from the following fact. Every \(S\)-group
finite and \'etale over~\(S\) can be embedded into an \(S\)-group
\(H\), affine and smooth over~\(S\), with connected fibers, and
commutative when \(G\) is commutative. The exact
sequence~\Ref{XI.4.5}, together where necessary with its
noncommutative variants, then allows the discrete classification of
principal coverings with group~\(G\) to be studied through the
continuous classification of principal bundles with group~\(H\),
and conversely. For the idea of the general construction of the
embedding \(G\to H\), which appears to be used rather little in
practice, see~\cite[VI~2.8]{XI.5}. In the next section we shall
develop only two important classical special cases. We shall need
the following auxiliary result:

\begin{proposition}
\label{XI.5.1}
Let
% page: 301
\marginpar{301}
%
\(S\) be a prescheme, and let \(G\) be an \(S\)-group isomorphic to
\(\Gl(n)_S\), for example \(\GG_{m\,S}\), or to \(\GG_{a\,S}\).
Then every principal homogeneous bundle under~\(G\) is locally
trivial.
\end{proposition}

Here \(\Gl(n)_S\), with \(n\geq0\), denotes the \(S\)-group
representing the contravariant functor
\[
  T\mto\Gl(n,\Gamma(T,\cal O_T))
\]
in the \(S\)-prescheme~\(T\). In particular, the multiplicative
group \(\GG_{m\,S}\) represents the contravariant functor
\[
  T\mto\Gamma(T,\cal O_T^*),
\]
and hence, as a prescheme over~\(S\), is isomorphic to
\[
  \Spec(\cal O_S[t,t^{-1}]),
\]
where \(t\) is an indeterminate. Similarly, \(\GG_{a\,S}\) represents
the contravariant functor
\[
  T\mto\Gamma(T,\cal O_T),
\]
and is therefore isomorphic, as an \(S\)-prescheme, to
\(\Spec(\cal O_S[t])\). By d\'evissage, \Ref{XI.5.1} recovers
Rosenlicht's local-triviality result for the case in which \(G\) has
a composition series whose successive factors are groups of the kind
considered here. For a finer study of local triviality questions for
principal homogeneous bundles, see~\cite{XI.7} and~\cite{XI.3}.

The first assertion is proved by observing that
\[
  G(T)=\Aut(\cal O_T^n),
\]
and that the morphisms \(S'\to S\) occurring in~\Ref{XI.4.1}, \ie
the faithfully flat and quasi-compact morphisms, are effective descent
morphisms for the fibered category of modules locally isomorphic to
\(\cal O_T^n\), \ie locally free of rank~\(n\)
(\ifthenelse{\boolean{orig}}{VIII~}{}\Ref{VIII.1.12}). The second assertion is proved similarly,
noting that in this case
\[
  G(T)=\Aut(\cal E_T),
\]
where \(\cal E_T\) is the \emph{trivial extension} of~\(\cal O_T\)
by~\(\cal O_T\), and where the automorphisms must of course preserve
the extension structure. The morphisms \(S'\to S\) in~\Ref{XI.4.1}
are effective descent morphisms for the fibered category of
extensions of~\(\cal O_T\) by~\(\cal O_T\), as follows readily from
\ifthenelse{\boolean{orig}}{VIII~}{}\Ref{VIII.1.1}, and such extensions are automatically locally
trivial.

\begin{remarque}
\label{XI.5.2}
The same type of proof applies to the symplectic group
\(\Symp(2n)_S\), since a nondegenerate alternating form on a module
locally isomorphic to~\(\cal O_S^{2n}\), \ie a form defining an
isomorphism of that module with its dual, is locally isomorphic to the
standard form. The corresponding result for the orthogonal group is
false, however, already when \(S\) is the spectrum of a field, since
there may be quadratic forms over a field which are not isomorphic
% page: 302
\marginpar{302}
%
to the standard form. In fact, it is proved essentially
in~\cite{XI.3} that the groups \(\Gl\), \(\Symp\), \(\GG_a\), and
those that can be d\'eviss\'ed into such groups are, up to minor
qualifications, the only ones for which a local-triviality result of
the kind considered here holds.
\end{remarque}

\begin{corollary}
\label{XI.5.3}
There are canonical bijections
\[
  \H^1(S,\GL(n)_S)\isomfrom
  \H^1(S,\GL(n,\cal O_S)),
\]
in particular
\[
  \H^1(S,\GG_{m\,S})\isomfrom
  \H^1(S,\cal O_S^*),
\]
and
\[
  \H^1(S,\GG_{a\,S})\isomfrom
  \H^1(S,\cal O_S),
\]
where the right-hand sides denote the cohomology of the topological
space \(S\) with coefficients in ordinary sheaves.
\end{corollary}

In particular, \(\H^1(S,\GL(n)_S)\) is identified with the set of
isomorphism classes of locally free modules of rank~\(n\) on~\(S\),
and \(\H^1(S,\GG_{a\,S})\) with the set of extension classes of the
module \(\cal O_S\) by itself.

\end{filecontents*}
% Original input: drafts/SGA1_XI_6_A_English.texfrag
\begin{filecontents*}[overwrite,noheader]{sga1-cumulative-input-131.tex}
\section[Application to principal coverings]
{Application to principal coverings: Kummer and Artin--Schreier theories}
\label{XI.6}

% In this section the source uses a and m as subscripts for the
% additive and multiplicative groups.

\begin{proposition}
\label{XI.6.1}
Let \(S\) be a prescheme, let \(n>0\) be an integer, let
\[
  u_n\colon\GG_{m\,S}\longrightarrow\GG_{m\,S}
\]
be the \(n\)-th power homomorphism, and let
\(\bbmu_{n\,S}\)
\label{indnot:kb}\oldindexnot{$\bbmu_{n\,S}$|hyperpage}%
%
be its kernel. Then \(\bbmu_{n\,S}\) is finite and locally free of
rank~\(n\) over~\(S\), and is \'etale over~\(S\) if and only if the
characteristic of every \(s\in S\) is prime to~\(n\). The sequence
of homomorphisms
\[
  0\longrightarrow\bbmu_{n\,S}
  \longrightarrow\GG_{m\,S}
  \lto{u_n}\GG_{m\,S}\longrightarrow0
\]
is exact in the sense of Section~\Ref{XI.4}. It will be called the
\emph{Kummer exact sequence}
%
\index{Kummer exact sequence|hyperpage}%
%
over~\(S\), relative to the integer~\(n\).
\end{proposition}

We have
% page: 303
\marginpar{303}
%
\[
  \GG_m=\Spec\cal O_S[t,t^{-1}],
\]
and \(u_n\) corresponds to the homomorphism on affine
\(\cal O_S\)-algebras given by
\[
  u_n(t)=t^n.
\]
The unit section of~\(\GG_{m\,S}\), on the other hand, corresponds to
the augmentation homomorphism of \(\cal O_S\)-algebras given by
\[
  \varepsilon(t)=1,
\]
whose kernel is the principal ideal \((t-1)\). Its image under
\(u_n\) is therefore the principal ideal \((1-t^n)\), and hence
\[
  \bbmu_{n\,S}=\Spec\bigl(\cal O_S[t]/(1-t^n)\bigr).
\]
This shows in particular that \(\bbmu_{n\,S}\) is finite over~\(S\)
and is defined by an algebra over~\(S\) which is free of rank~\(n\),
with basis \(t^i\), \(0\leq i\leq n-1\).

For this group to be \'etale at \(s\in S\), it is necessary and
sufficient that the reduced algebra \(k[t]/(1-t^n)\), where
\(k=\kres(s)\), obtained by formally adjoining the \(n\)-th roots of
unity to~\(k\), be separable over~\(k\). Equivalently, the roots of
\(1-t^n\) in an algebraic closure of~\(k\) must all be distinct,
which is the same as \(n\) being prime to the characteristic.

Finally, by the criterion~\Ref{XI.4.2}, exactness of the sequence in
\Ref{XI.6.1} reduces to proving that
\ifthenelse{\boolean{orig}}
{\(v\)}
{\(u_n\)}
% correction: v should be u_n.
is faithfully flat. We may suppose \(S\) affine with ring~\(A\), so
\(\GG_{m\,S}\) is affine with ring \(B=A[t,t^{-1}]\). It is enough to
check that \(u_n\) makes \(B\) a free module of rank~\(n\) over~\(B\),
or equivalently that \(u_n\) is injective and that
\[
  A[t,t^{-1}]
\]
is a free module of rank~\(n\) over \(A[t^n,t^{-n}]\). The elements
\(t^i\), \(0\leq i\leq n-1\), plainly form a basis, completing the
proof.

\begin{definition}
\label{XI.6.2}
The group \(\bbmu_{n\,S}\) is called the \emph{Kummer group of
rank~\(n\) over~\(S\)}, and a \emph{Kummer principal covering of
rank
% page: 304
\marginpar{304}
%
\ifthenelse{\boolean{orig}}
{\(nS\)}
{\(n\) over~\(S\)}}
is a principal homogeneous bundle over~\(S\) whose group is the
Kummer group of rank~\(n\).
\end{definition}

The set of these coverings is a group, denoted
\(\H^1(S,\bbmu_{n\,S})\), or simply \(\H^1(S,\bbmu_n)\). Formation
of the Kummer group of rank~\(n\) over~\(S\) is compatible with base
extension, so \(\bbmu_{n\,S}\) comes by base extension from the
\emph{absolute Kummer group \(\bbmu_n\)} over~\(\Spec(\ZZ)\).

Write \((\ZZ/n\ZZ)_S\) for the \(S\)-group defined by the ordinary
finite group \(\ZZ/n\ZZ\). If \(G\) is any \(S\)-group, the
homomorphisms of \(S\)-groups
\[
  u\colon(\ZZ/n\ZZ)_S\longrightarrow G
\]
are in bijection, compatibly with base change, with the sections of
\(G\) over~\(S\) whose \(n\)-th power is the unit section. The
correspondence sends \(u\) to the image under~\(u\) of the section of
\((\ZZ/n\ZZ)_S\) over~\(S\) defined by the generator
\(1\bmod n\ZZ\). Thus:

\begin{corollary}
\label{XI.6.3}
If \(\bbmu_{n\,S}\) is \'etale over~\(S\), there is a bijection
between isomorphisms of \(S\)-groups
\[
  (\ZZ/n\ZZ)_S\isomto\bbmu_{n\,S}
\]
and sections of \(\cal O_S^*\) having exact order~\(n\) on every
connected component of~\(S\). Such a section is called a
\emph{primitive \(n\)-th root of unity over~\(S\)}. Consequently,
\(\bbmu_{n\,S}\) is isomorphic as an \(S\)-group to
\((\ZZ/n\ZZ)_S\) if and only if it is \'etale
\ifthenelse{\boolean{orig}}
{over~\(S\)}
{over~\(S\),}
\ie the residue characteristics of~\(S\) are prime to~\(n\), and
there exists a primitive \(n\)-th root of unity over~\(S\).
\end{corollary}

This explains the role in classical Kummer theory of the hypotheses
that the base field, playing the role of~\(S\), have characteristic
prime to~\(n\) and contain the \(n\)-th roots of unity, together with
the choice of a primitive \(n\)-th root of unity. Once the language
of schemes is available, these hypotheses need no longer encumber
the theory: one should reason directly with \(\bbmu_n\), rather than
with \(\ZZ/n\ZZ\). Combining~\Ref{XI.6.1}, \Ref{XI.4.5}, and
\Ref{XI.5.3} gives the following general relation between the theory
of Kummer principal coverings and that of Picard groups:

\begin{proposition}
\label{XI.6.4}
Let
% page: 305
\marginpar{305}
%
\(S\) be a prescheme. There is a canonical exact sequence
\[
  \textstyle
  0\to\H^0(S,\bbmu_n)
  \to\H^0(S,\cal O_S^*)
  \to\H^0(S,\cal O_S^*)
  \to\H^1(S,\bbmu_n)
  \to\H^1(S,\cal O_S^*)
  \to\H^1(S,\cal O_S^*).
\]
On setting
\ifthenelse{\boolean{orig}}
{\(\H^1(S,\cal O_S)=\Pic(S)\),}
{\(\H^1(S,\cal O_S^*)=\Pic(S)\),}
and, for every abelian group \(A\), denoting the kernel and cokernel
of multiplication by~\(n\) on~\(A\) by \({}_nA\) and \(A_n\),
respectively, one obtains the exact sequence
\[
  0\longrightarrow
  \H^0(S,\cal O_S^*)_n
  \longrightarrow\H^1(S,\bbmu_n)
  \longrightarrow{\sideset{_n}{}\Pic}(S)
  \longrightarrow0.
\]
\end{proposition}

We spell out two important cases in which one or the other outer term
of this exact sequence vanishes:

\begin{corollary}
\label{XI.6.5}
Suppose that \({\sideset{_n}{}\Pic}(S)=0\), for example if \(S\) is
the spectrum of a local ring or of a factorial ring, and set
\[
  A=\H^0(S,\cal O_S).
\]
Then there is a canonical isomorphism
\[
  \H^1(S,\bbmu_n)\isomto A^*/{A^*}^n.
\]
\end{corollary}

This is essentially the classical statement of Kummer theory when
\(S\) is the spectrum of a field.

\begin{corollary}
\label{XI.6.6}
Suppose that every element of \(\H^0(S,\cal O_S)\) is an \(n\)-th
power, for example if \(\H^0(S,\cal O_S)\) is a product of
algebraically closed fields, or if \(S\) is reduced and proper over an
algebraically closed field~\(k\). Then there is a canonical
isomorphism
\[
  \H^1(S,\bbmu_n)\isomto{\sideset{_n}{}\Pic}(S).
\]
\end{corollary}

In particular, when \(S\) is proper and connected over an
algebraically closed field~\(k\), this relates the fundamental group
of~\(S\) to the finite-order points of the Picard scheme \(P\) of~\(S\)
over~\(k\). Thus, for \(n\) prime to the characteristic, one has an
isomorphism
\[
  \Hom(\pi_1(S),\ZZ/n\ZZ)\simeq{}_nP(k),
\]
which is often used in algebraic
% page: 306
\marginpar{306}
%
geometry. If the connected component \(P^0\) of~\(P\) is a complete
group scheme of dimension~\(g\), the results recalled in
Section~\Ref{XI.2}, together with finiteness of the
\ifthenelse{\boolean{orig}}
{N\'eron--S\'ev\'eri}
{N\'eron--Severi}
torsion group, show that for every prime \(\ell\) different from the
characteristic, the \(\ell\)-primary component of the abelianized
fundamental group \(\pi_1(S)\) is a finite-type module of rank~\(2g\)
over the ring \(\ZZ_\ell\) of \(\ell\)-adic integers. It is moreover
free for all but at most finitely many values of~\(\ell\).

As Serre observed, this can be used under suitable hypotheses to prove
that, when \(X\) is flat and projective over connected~\(S\), the
Picard schemes of the fibers of~\(X\) all have the same dimension,
by applying the semicontinuity theorem \ifthenelse{\boolean{orig}}{X~}{}\Ref{X.2.3}. Serre's argument
applies whenever the Picard scheme of~\(X\) over~\(S\) exists and
the connected Picard schemes of the fibers of~\(X\) over~\(S\) are
proper group schemes. This holds, for example, when the geometric
fibers of~\(X\) over~\(S\) are normal, with \(X\) still flat and
projective over~\(S\), and in particular when \(X\) is smooth and
projective over~\(S\).

Now let \(p\) be a prime, and suppose that \(S\) is a prescheme of
characteristic~\(p\), \ie \(p\cal O_S=0\). The \(p\)-th power map on
\(\cal O_S\) is additive, and the corresponding morphism, obtained
by replacing \(S\) by a variable \(T\) over~\(S\),
\[
  F\colon\GG_{a\,S}\longrightarrow\GG_{a\,S},
\]
is therefore a homomorphism of \(S\)-groups, called the
\emph{Frobenius homomorphism}. Such a morphism is defined for every
\(S\)-prescheme \(G\) obtained by base extension from a prescheme
\(G_0\) over the prime field \(\ZZ/p\ZZ\), and it is a group
homomorphism if \(G_0\) is a group prescheme. We set
\[
  \wp=\id-F\colon\GG_{a\,S}\longrightarrow\GG_{a\,S}.
\]

\end{filecontents*}
% Original input: drafts/SGA1_XI_6_B_English.texfrag
\begin{filecontents*}[overwrite,noheader]{sga1-cumulative-input-132.tex}
Consider, on the other hand, the \(S\)-group \((\ZZ/p\ZZ)_S\)
defined by the ordinary finite group~\(\ZZ/p\ZZ\). We have said
that, for every \(S\)-group~\(G\), the homomorphisms of \(S\)-groups
from~\((\ZZ/p\ZZ)_S\) to~\(G\) correspond bijectively to the
sections of~\(G\) over~\(S\) whose \(p\)-th power is the identity
section. When \(G=\GG_{a\,S}\), they therefore correspond
% page: 307
\marginpar{307}
%
to arbitrary sections of~\(G\) over~\(S\). Taking in particular the
section of~\(\GG_{a\,S}\) over~\(S\) corresponding to the unit
section of the sheaf of rings \(\cal{O}_S\), we obtain a
homomorphism of \(S\)-groups
$$
\ifthenelse{\boolean{orig}}
{i\colon(\ZZ/p\ZZ)\to\GG_{a\,S}}
{i\colon(\ZZ/p\ZZ)_S\to\GG_{a\,S}}
% correction: (\ZZ/p\ZZ)_S
$$

\begin{proposition}
\label{XI.6.7} The sequence of homomorphisms of \(S\)-groups
$$
\ifthenelse{\boolean{orig}}
{0\to(\ZZ/p\ZZ)_S\to\GG_{a\,S}\to G_{a\,S}\to 0}
{0\to(\ZZ/p\ZZ)_S\to\GG_{a\,S}\to \GG_{a\,S}\to 0}
% correction: last 'G' should be '\GG'.
$$
is exact (in the sense of \No \Ref{XI.4}). (It is called the
\emph{Artin--Schreier exact sequence} over~\(S\).)
%
\index{Artin--Schreier (exact sequence)|hyperpage}%
%
\end{proposition}

It is enough to prove this over the prime field \(k=\ZZ/p\ZZ\). It is
enough to observe that the homomorphism
\(\wp^*\colon k[t]\to k[t]\) defined by
\(\wp^*(t)=t-t^p\) makes \(k[t]\) a free module of rank \(p\) over
\(k[t]\): more precisely, \(k[t]\) is a free module over \(k[s]\),
where \(s=t-t^p\), with basis given by the \(t^i\)
(\(0\leq i\leq p-1\)).

Using \Ref{XI.4.5} and~\Ref{XI.5.3}, we conclude:

\begin{proposition}
\label{XI.6.8} There is a canonical exact sequence
\begin{multline*}
0\to\H^0(S,\ZZ/p\ZZ)\to\H^0(S,\cal{O}_S)\to\H^0(S,\cal{O}_S)\to
\H^1(S,\ZZ/p\ZZ)\\
\to\H^1(S,\cal{O}_S)\to\H^1(S,\cal{O}_S),
\end{multline*}
and hence an exact sequence
$$
0\to\H^0(S,\cal{O}_S)/\wp\H^0(S,\cal{O}_S)\to
\H^1(S,\ZZ/p\ZZ)\to\H^1(S,\cal{O}_S)^F\to 0,
$$
(where the superscript \(F\) in the last term denotes the subgroup
of invariants under the endomorphism \(F\), equal to the kernel of
\(\wp=\id-F\)).
\end{proposition}

Let us spell out two extreme cases:

\begin{corollaire}
\label{XI.6.9} Suppose that \(\H^1(S,\cal{O}_S)^F=0\), for example
that \(S\) is an affine scheme. Then, putting
\(A=\H^0(S,\cal{O}_S)\), there is a canonical isomorphism
$$
\H^1(S,\ZZ/p\ZZ)\isomto A/\wp A.
$$
\end{corollaire}

This is \emph{Artin--Schreier theory} in its classical form, at
least when \(S\) is the spectrum of a field.

\begin{corollaire}
\label{XI.6.10}
Suppose
% page: 308
\marginpar{308}%
%
that \(\wp\H^0(S,\cal{O}_S)\mkern1mu{=}\mkern1mu
\H^0(S,\cal{O}_S)\), for example that
\(\H^0(S,\cal{O}_S)\) is a product of algebraically closed fields,
or that \(S\) is proper over an algebraically closed field. Then
there is a canonical isomorphism
$$
\H^1(S,\ZZ/p\ZZ)\isomto\H^1(S,\cal{O}_S)^F.
$$
\end{corollaire}

\begin{remarques}
\label{XI.6.11}
The last statement is due to
\ifthenelse{\boolean{orig}}{J.P\ptbl}{J-P\ptbl}%
Serre
\cite{XI.9}. It is also possible to develop an analogous theory for
the structure group \(\ZZ/p^n\ZZ\), for arbitrary \(n\), by using
the Witt group scheme \(\WW_n\) in place of~\(\GG_a\), \cf \loccit
Notice that in characteristic \(p>0\), Kummer theory no longer gives
information about principal coverings of order \(p\), since
\(\bbmu_p\) is then an \emph{infinitesimal} group, \ie radicial over
the base, and hence has no direct connection with \(\ZZ/p\ZZ\).
Thus, at first sight, the theory of these coverings can no longer be
handled (when, for definiteness, \(S\) is a proper scheme over an
algebraically closed field) by the theory of the Picard scheme as
in~\Ref{XI.6.6}. Nevertheless, recalling that the Zariski tangent
space at the origin of
\ifthenelse{\boolean{orig}}{\(\mathbf{Pic}_{S/K}\)}{\(\mathbf{Pic}_{S/k}\)}\kern1pt\footnote{For the definition of
\ifthenelse{\boolean{orig}}{\(\mathbf{Pic}_{S/K}\)}{\(\mathbf{Pic}_{S/k}\)}, \cf A. Grothendieck, Sém. Bourbaki
\No 232 (February 1962).} is identified with
\(\H^1(S,\cal{O}_S)\), we see that \emph{knowledge of the group
scheme \({_p}\mathbf{Pic}_{S/k}\), the kernel of multiplication by
\(p\) in \(\mathbf{Pic}_{S/k}\), determines
\(\H^1(S,\ZZ/p\ZZ)\) as well as \(\H^1(S,\bbmu_p)\); note that it
also determines \(\H^1(S,\bbalpha_p)\)}. Here \(\bbalpha_p\)
denotes the infinitesimal group scheme over the prime field that is
the kernel of \(F\colon\GG_a\to\GG_a\) (it may also be described as
the spectrum of the restricted enveloping algebra of the trivial
one-dimensional \(p\)-Lie algebra): indeed, the exact
sequence~\Ref{XI.4.5} gives in this case
$$
\H^1(S,\bbalpha_p)\simeq\Ker
(F\colon\H^1(S,\cal{O}_S)\to\H^1(S,\cal{O}_S)),
$$
and, more generally, if \(\bbalpha_{p^n}\) denotes the kernel in
\(\GG_a\) of the \(n\)-th iterate of~\(F\), then
$$
\H^1(S,\bbalpha_{p^n})\simeq\Ker(F^n\colon\H^1(S,\cal{O}_S)\to
\H^1(S,\cal{O}_S)).
$$

In fact, knowledge of~\({_p}\mathbf{Pic}_{S/k}\) is equivalent to
knowledge of~\(\H^1(S,G)\) for every finite commutative algebraic
% page: 309
\marginpar{309}
%
group \(G\) annihilated by \(p\). More generally, knowledge
of~\({_{p^n}}\mathbf{Pic}_{S/k}\) is equivalent to knowledge
of~\(\H^1(S,G)\) for every finite commutative algebraic group \(G\)
annihilated by \(p^n\), by virtue of the following theorem, which in
the situation under consideration encompasses both Kummer theory and
Artin--Schreier theory:

Let \(G\) be a finite algebraic group over~\(k\), and let
\(D(G)=\SheafHom_{k\textup{-groups}}(G,\GG_m)\) be its
\emph{Cartier dual} (whose affine algebra is carried by the vector
space dual to the affine algebra of~\(G\), \ie by the hyperalgebra
of~\(G\) in the sense of Dieudonné--Cartier). Then there is a
canonical isomorphism
\begin{equation}\label{eqXI.6.11}\tag{$*$}
\H^1(S,G)\simeq
\Hom_{k\textup{-groups}}(D(G),\mathbf{Pic}_{S/k}).
\end{equation}
(N.B. \(S\) is a proper scheme over the algebraically closed
field~\(k\), with \(\H^0(S,\cal{O}_S)=k\).) This formula can also be
expressed by saying that the \emph{true fundamental group} of~\(S\)
alluded to in \No \Ref{XI.2}, after abelianization, is isomorphic to
the projective limit of the \(D(P_i)\), where \(P_i\) ranges over the
\emph{finite} algebraic subgroups of~\(\mathbf{Pic}_{S/k}\). We
denote it by \ifthenelse{\boolean{orig}}{\(T^\bbullet(\mathbf{Pic}_{X/k})\)}{\(T^\bbullet(\mathbf{Pic}_{S/k})\)}. When \(S\) is an
abelian variety, we saw in~\Ref{XI.2.1} that this group is also
isomorphic to the \emph{true} Tate module
\(T_{\bbullet}(S)=\varprojlim\, {{_n}S}\), and the isomorphism
\eqref{eqXI.6.11} then takes the more striking form
$$
\Ext^1(A,G)\simeq\Hom(D(G),B),
$$
where \(A\) is an abelian variety, \(B\) its dual, and \(G\) a finite
algebraic group over~\(k\). The results just indicated can moreover
be generalized to the case where \(k\) is replaced by an arbitrary
base prescheme, and to coefficient groups \(G\) other than finite
groups.
\end{remarques}

%\section{Bibliography}
\begin{thebibliography}{10}
% page: 310
\marginpar{310}
%

\bibitem[1]{XI.1} A\ptbl \textsc{Grothendieck},
Sur quelques points d'algèbre homologique,
\ifthenelse{\boolean{orig}}
{Tohoku Math.\ Journal, pp\ptbl 119--221, Vol\ptbl 9, 1957.}
{T\^ohoku Math.~J.\ \textbf{9} (1957) p\ptbl 119--221.}

\bibitem[2]{XI.2} A\ptbl \textsc{Grothendieck}, A general theory of fibre spaces with
structure sheaf, University of Kansas, 1955.

\bibitem[3]{XI.3} A\ptbl \textsc{Grothendieck}, Torsion homologique et sections
rationnelles, Séminaire \hbox{Chevalley}, June 16, 1958.

\bibitem[4]{XI.4} S\ptbl \textsc{Lang}, Abelian varieties, Interscience tracts in pure
and applied mathematics, \No 7, New York.

\bibitem[5]{XI.5}
\ifthenelse{\boolean{orig}}{J.P\ptbl}{J-P\ptbl}%
\textsc{Serre}, Groupes algébriques et corps de classes,
Actualités Scient. et Ind.\ \No 1264, Hermann, Paris, 1959.

\bibitem[6]{XI.6}
\ifthenelse{\boolean{orig}}{J.P\ptbl}{J-P\ptbl}%
\textsc{Serre}, Groupes proalgébriques,
\ifthenelse{\boolean{orig}}
{Publications Mathématiques de l'IHES \No 7, pp\ptbl 1--67, 1960.}
{Publications Mathématiques de l'IHES \textbf{7} (1960), p\ptbl 1--67.}

\bibitem[7]{XI.7}
\ifthenelse{\boolean{orig}}{J.P\ptbl}{J-P\ptbl}%
\textsc{Serre}, Espaces fibrés algébriques,
Séminaire Chevalley, April 21, 1958.

\bibitem[8]{XI.8}
\ifthenelse{\boolean{orig}}{J.P\ptbl}{J-P\ptbl}%
\textsc{Serre}, Quelques propriétés des
variétés abéliennes en
\ifthenelse{\boolean{orig}}
{car.~$p$, American Journal of Math., vol\ptbl 80, pp\ptbl 715--739, 1958.}
{caractéristique~$p$, Amer.\ J.\ Math. \textbf{80} (1958), p\ptbl 715--739.}

\bibitem[9]{XI.9}
\ifthenelse{\boolean{orig}}{J.P\ptbl}{J-P\ptbl}%
\textsc{Serre}, Sur la topologie des variétés
algébriques en
\ifthenelse{\boolean{orig}}
{car.~$p$}
{caractéristique~$p$},
Symposium Internacional de \ifthenelse{\boolean{orig}}{Topologia}{Topolog\'ia}
Algebraica, 1958.

\bibitem[10]{XI.10}
\ifthenelse{\boolean{orig}}{J.P\ptbl}{J-P\ptbl}%
\textsc{Serre}, On the fundamental group of a unirational
variety,
\ifthenelse{\boolean{orig}}
{Journal London Math.\ Society vol\ptbl 34, pp\ptbl 481--484, 1959.}
{J.~London Math.\ Soc.\ \textbf{34} (1959), p\ptbl 481--484.}
\end{thebibliography}


\end{filecontents*}
% Original input: drafts/SGA1_XI_OPENING_AND_1_English.texfrag
\begin{filecontents*}[overwrite,noheader]{sga1-cumulative-input-133.tex}
\chapterspace{-2}
\chapter{Examples and complements}
\label{XI}
% page: 285
\marginpar{285}
%

\section{Projective spaces, unirational varieties}
\label{XI.1}

\begin{proposition}
\label{XI.1.1}
Let $k$ be an algebraically closed field, and let $X=\PP_k^r$ be
projective space of dimension $r$ over~$k$. Then $X$ is \emph{simply
connected,} \ie $\pi_1(X)=0$.
\end{proposition}

For $r=0$, this is trivial. If $r=1$, we must show that if $X'$ is a
nonempty connected \'etale covering of~$X=\PP_k^1$, then
$X'\isomto X$. The genus formula gives, in this case, if~$g$ and $g'$ are
the genera of~$X$ and~$X'$:
$$
1-g' = d(1-g)\quoi,
$$
where $d$ is the degree of~$X'$ over~$X$. Since $g=0$, we therefore
have $1-g'=d$, which requires $d=1$ since $g'\ge 0$, and this proves
$X'\isomto X$. When $r\ge 2$, we proceed by induction
on~$r$, assuming that $\PP^{r'}$ is simply connected for
$r'<r$. Applying this to a hyperplane in~$\PP^r$ and using
(\ifthenelse{\boolean{orig}}{X~}{}\Ref{X.2.10}), it indeed follows that $\PP^r$ is simply
connected. Another proof: by (\ifthenelse{\boolean{orig}}{X~}{}\Ref{X.1.7}), we have
$\pi_1(\PP^1\times\cdots\times \PP^1) = \pi_1(\PP^1)\times\cdots
\times \pi_1(\PP^1)$, and hence $(\PP^1)^r$ is simply connected
since $\PP^1$ is, so $\PP^r$ is simply connected by the birational
invariance of the fundamental group (\ifthenelse{\boolean{orig}}{X~}{}\Ref{X.3.4}). This proof shows
more generally:

\begin{corollaire}
\label{XI.1.2}
Let $X$ be a proper normal scheme over an algebraically closed
field~$k$; if $X$ is a rational variety, \ie it is integral and its
function field is a purely transcendental extension of~$k$, then $X$
is simply connected.
\end{corollaire}

This result applies in particular to Grassmann varieties and, more
% page: 286
\marginpar{286}%
%
generally, to varieties~$G/H$, where $G$ is a connected linear group
over~$k$ and $H$ is an algebraic subgroup containing a Borel subgroup
\ifthenelse{\boolean{orig}}
{of~$B$.}
{of~$G$.}
% correction: `$B$' should be `$G$'.

Recall that a \emph{unirational variety over~$k$}
%
\index{unirationnelle (vari\'et\'e)|hyperpage}%
%
is a proper integral scheme over~$k$ whose function field $K$ is
contained in a purely transcendental extension~$K'$ of~$k$ that is
finite over~$K$ (\ie has the same transcendence degree over~$k$
as~$K$), \ie if there exists a dominant rational map
$f\colon\PP_k^r\to X$, with $r=\dim X$. If $X$ is normal, it follows
from the considerations preceding \ifthenelse{\boolean{orig}}{X~}{}\Ref{X.3.4} that, for every
connected \'etale covering $X'$ of~$X$, with function field~$L/K$,
the algebra $L\otimes_K K'$ over~$K'$ is unramified on the
model~$\PP^r$, and is therefore completely split by~\Ref{X.1.1}; this
shows that $L$ is $K$-isomorphic to a subextension of~$K'/K$. Taking
\ifthenelse{\boolean{orig}}{V~}{}\Ref{V.8.2} into account, this therefore proves:

\begin{corollaire}
\label{XI.1.3}
The fundamental group of a normal unirational variety over an
algebraically closed field is finite.
\end{corollaire}

(N.B. Note that in the definition of \og unirational\fg, it was not
necessary for $K'/K$ to be finite).

\ifthenelse{\boolean{orig}}
{\begin{remarques}
\label{XI.1.4}
Of course, the results of this section are well known. On the other
hand,
\ifthenelse{\boolean{orig}}{J.P\ptbl}{J-P\ptbl}%
Serre showed~\cite{XI.10} that when $X$ is a smooth projective
unirational variety over an algebraically closed field of
\emph{characteristic zero,} $X$ is simply connected. His proof is
transcendental in that it uses the Hodge symmetry theorem, and it is
not known whether this result extends to characteristic~$p>0$.
Moreover, it seems that no example is known of a smooth unirational
variety over~$k$ that is not already rational.
\end{remarques}}
{\begin{remarques}
\label{XI.1.4}
Of course, the results of this section are well known. On the other
hand,
\ifthenelse{\boolean{orig}}{J.P\ptbl}{J-P\ptbl}%
Serre showed~\cite{XI.10} that when $X$ is a smooth projective
unirational variety over an algebraically closed field of
\emph{characteristic zero,} $X$ is simply connected. His proof is
transcendental in that it uses the Hodge symmetry theorem.

\lcrochetbf Added in 2003 (MR) Let us restrict ourselves to the case of an
algebraically closed field $k$ of characteristic $p\geq 0$. In characteristic
$p>0$, the definition of a unirational $k$-variety given above is that of a
weakly unirational $k$-variety, as opposed to that of a strongly unirational
$k$-variety, where one assumes in addition that $K'$ is a separable extension
of~$K$.

In dimension $2$, there exist smooth projective weakly unirational surfaces
whose fundamental group (finite by \hyperref[XI.1.3]{1.3}) is nontrivial, and which are therefore
not rational (T\ptbl Shioda, On unirationality of supersingular surfaces,
Math.~Ann. \textbf{225} (1977), p\ptbl155--159.). On the other hand, a strongly
unirational surface is rational: this follows from Castelnuovo's rationality
criterion, extended to characteristic $p>0$ by O\ptbl Zariski (\cf J-P\ptbl
Serre, S\'em. Bourbaki \no146, 1957 and {\OE}uvres--Collected papers,
vol\ptbl 1, p\ptbl 467).

Over the field $\CC$ of complex numbers, examples are known of smooth
projective varieties of \hbox{dimension $\geq 3$} that are unirational but not
rational (\cf P\ptbl Deligne, Vari\'et\'es unirationnelles non
rationnelles [d'apr\`es M\ptbl Artin et D\ptbl Mumford], S\'em.\ Bourbaki
\no402, 1971-72, Lecture Notes vol\ptbl 317). This is the case, in particular,
of a smooth cubic hypersurface in projective space $\PP^4$.
Some of these examples extend to characteristic $>0$ to give strongly
unirational nonrational varieties (\cf J.P\ptbl Murre, Reduction of the proof
of the non-rationality of a non singular cubic threefold to a result of
Mumford, Compositio Math. \textbf{27} (1973), p\ptbl63--82).

Let $V$ be a normal integral projective $k$-variety. We say that $V$ is
rationally connected if there exist an integral $k$-scheme $T$ of finite type
and a $k$-morphism $F : T\times\PP^1\to V$ such that the morphism
\begin{equation}\tag*{(1)}
\begin{split}
T \times\PP^1\times\PP^1&\to V\times V\\
(t, u, u')&\mto(F(t,u), F(t,u'))
\end{split}
\end{equation}
is dominant. In particular, a rational curve passes through any two
sufficiently general rational points of~$V$. The terminology is justified by
the fact that if $V$ is rationally connected, two rational points can be
joined by a finite chain of rational curves. If $k$ has characteristic $p>0$,
we strengthen the preceding definition by requiring the map (1) to be
generically smooth. This gives the notion of a separably rationally connected
variety. Thus unirational varieties are rationally connected, and in
characteristic $p>0$, strongly unirational varieties are separably rationally
connected. J\ptbl Koll\'ar showed that separably rationally connected varieties
have trivial algebraic fundamental group (\cf O\ptbl Debarre, Vari\'et\'es
rationnellement connexes [d'apr\`es T\ptbl Graber, J\ptbl Harris, J\ptbl Starr
et A.J\ptbl de~Jong], S\'em. Bourbaki \no906, 2001-2002).\rcrochetbf
\end{remarques}}
\end{filecontents*}
% Original input: drafts/SGA1_X_1_B_English.texfrag
\begin{filecontents*}[overwrite,noheader]{sga1-cumulative-input-134.tex}
\subsection{}
\ifthenelse{\boolean{orig}}{\label{X.I.6}}{\label{X.1.6}}
Continue to suppose that \(f\) is proper and separable, but make no
further assumption on \(f_*(\cal O_X)\). The latter corresponds to a
well-determined \'etale covering \(Y'\) of~\(Y\), pointed above~\(b\)
by the image \(b'\) of~\(a\). Applying~\Ref{X.1.4} to the canonical
morphism \(X\to Y'\), and assuming \(f\) surjective, the exact
sequence~\Ref{X.1.4} is replaced by the following sequence, analogous
to the homotopy exact sequence of fiber spaces in algebraic topology:
\[
  \pi_1(\overline X_y,\overline a)
  \longrightarrow\pi_1(X,a)
  \longrightarrow\pi_1(Y,b)
  \longrightarrow\pi_0(\overline X_y,\overline a)
  \longrightarrow\pi_0(X,a)
  \longrightarrow\pi_0(Y,b)\longrightarrow e.
\]

Of course, in~\Ref{X.1.4} one cannot in general assert that the
homomorphism
\[
  \pi_1(\overline X_y,\overline a)\longrightarrow\pi_1(X,a)
\]
is injective. In algebraic topology its kernel is the image of
\(\pi_2(Y,b)\), and in algebraic geometry as well one ought to
introduce homotopy groups in every dimension, together with the full
homotopy exact sequence for a proper morphism satisfying suitable
hypotheses, for example being
\ifthenelse{\boolean{orig}}
{smooth.}
{smooth).}
% correction: close the parenthesis.
At present no result of this kind is available, apart from a
reasonable, if not definitive, definition of the higher homotopy
groups.

\begin{corollary}
\label{X.1.7}
Let \(k\) be an algebraically closed field, and let \(X\) and \(Y\)
be two connected preschemes over~\(k\), with \(X\) proper over~\(k\)
and \(Y\) locally
% page: 265
\marginpar{265}
%
noetherian. Let \(a\) be a geometric point of~\(X\) and \(b\) a
geometric point of~\(Y\), both with values in the same algebraically
closed extension \(K\) of~\(k\). Let \(c=(a,b)\) be the corresponding
geometric point of~\(X\times_kY\). Then the homomorphism
\[
  \pi_1(X\times_kY,c)
  \longrightarrow\pi_1(X,a)\times\pi_1(Y,b)
\]
induced by the two projections \(X\times_kY\to X\) and
\(X\times_kY\to Y\) is an \emph{isomorphism}.
\end{corollary}

First suppose that \(K=k\). Set \(Z=X\times_kY\), consider the
projection \(f\colon Z\to Y\) and the underlying point \(y\) of the
geometric point \(b\) of~\(Y\), and apply~\Ref{X.1.4}. For this
purpose, observe that after replacing \(X\) by \(X_\red\), which does
not change the fundamental groups in question, we may assume that
\(X\) is reduced, and hence separable over~\(k\). Thus \(Z\) is
separable over~\(Y\), and plainly has geometrically connected fibers,
since \(X\) is connected. The geometric fiber of~\(Z\) at~\(b\) is
canonically isomorphic to \(X\otimes_kK=X\).

Moreover, since the composite of the morphisms \(X\to Z\to X\) is
the identity, the homomorphism
\(\pi_1(X,a)\to\pi_1(Z,c)\) is injective, and~\Ref{X.1.4} gives an
exact sequence
\[
  e\longrightarrow\pi_1(X,a)
  \longrightarrow\pi_1(Z,c)
  \longrightarrow\pi_1(Y,b)\longrightarrow e.
\]
There is also the canonical exact sequence
\[
  e\longrightarrow\pi_1(X,a)
  \longrightarrow\pi_1(X,a)\times\pi_1(Y,b)
  \longrightarrow\pi_1(Y,b)\longrightarrow e,
\]
in which the displayed homomorphisms are the canonical inclusion and
projection. Finally, the identity homomorphisms on the outer terms
together with the canonical homomorphism
\[
  \pi_1(Z,c)\longrightarrow\pi_1(X,a)\times\pi_1(Y,b)
\]
on the middle terms define a homomorphism from the first exact
sequence to the second. The resulting diagram is plainly commutative.
Since the homomorphisms on the outer terms are isomorphisms, so is the
homomorphism on the middle terms. This proves~\Ref{X.1.7} in this
case.

When \(K\) is no longer assumed equal to~\(k\), the argument gives
only an isomorphism
\[
  \pi_1(Z,c)\longrightarrow
  \pi_1(X\otimes_kK,a)\times\pi_1(Y,b),
\]
% page: 266
\marginpar{266}
%
and~\Ref{X.1.7} is then equivalent to the following special case:

\begin{corollary}
\label{X.1.8}
Let \(X\) be a proper connected scheme over an algebraically closed
field~\(k\), let \(k'\) be an algebraically closed extension of~\(k\),
let \(a'\) be a geometric point of \(X\otimes_kk'\), and let \(a\) be
its image in~\(X\). Then the canonical homomorphism
\[
  \pi_1(X\otimes_kk',a')\longrightarrow\pi_1(X,a)
\]
is an \emph{isomorphism}.
\end{corollary}

Surjectivity of this homomorphism is equivalent to the assertion that,
if \(X'\) is a connected \'etale covering of~\(X\), then
\(X'\otimes_kk'\) is connected as well. This follows at once from
the fact that \(k\) is algebraically closed, and is also a special
case
\ifthenelse{\boolean{orig}}
{IX~\Ref{IX.3.4}.}
{of \ifthenelse{\boolean{orig}}{IX~}{}\Ref{IX.3.4}.}
% correction: insert of.
The properness hypothesis on~\(X\) has not yet been used. Injectivity
of the homomorphism in question means the following: \emph{every
\'etale covering of~\(X\otimes_kk'\) is isomorphic to the inverse
image of an \'etale covering of~\(X\)}.

By an essentially formal argument, one can find a \(k\)-subalgebra
\(A\) of
\ifthenelse{\boolean{orig}}
{\(K\)}
{\(k'\),}
of finite type over~\(k\), and an \'etale covering of
\(X\otimes_kA\) whose inverse image over \(X\otimes_kk'\) is
isomorphic to the given covering. Let \(Y=\Spec(A)\). This is an
integral \(k\)-scheme of finite type, and hence has rational points
over~\(k\). Apply~\Ref{X.1.7} to the fundamental group of
\(X\times Y\) at a point \((a,b)\) rational over~\(k\). Every
connected \'etale covering of~\(X\times Y\) is then isomorphic to a
quotient of a covering \(X'\times Y'\), where \(X'\) and \(Y'\) are
Galois \'etale coverings of~\(X\) and~\(Y\), with groups \(G\) and
\(G'\), by a subgroup \(H\) of \(G\times G'\). It follows that the
inverse image of this covering of~\(X\times Y\) over
\(X\times Y'\) is isomorphic to a covering of the form
\(X_1'\times Y'\), where \(X_1'\) is an \'etale covering of~\(X\).

Let \(L\) be the function field of~\(Y\), equal to the fraction field
of~\(A\) inside~\(k'\). The \'etale covering of \(X\otimes_kL\)
induced by the given covering of~\(X\times_kY\) has the following
property: there is a finite separable extension \(L'/L\) such that
its inverse image over \(X\otimes_kL'\) is isomorphic to
\(X_1'\otimes_kL'\). Since \(k'\) is algebraically closed, we may
suppose that \(L'\) is contained in~\(k'\). This proves that the
given \'etale covering of \(X\otimes_kk'\) is isomorphic to
\(X_1'\otimes_kk'\), as required.

The explicit description just given for \'etale coverings of a
product
% page: 267
\marginpar{267}
%
\(X\times_kY\) immediately implies:

\begin{corollary}
\label{X.1.9}
Let \(k\) be an algebraically closed field, let \(X\) and \(Y\) be
two locally noetherian preschemes over~\(k\), let \(Z=X\times_kY\),
and let \(Z'\) be an \'etale covering of~\(Z\). For every point
\(y\in Y\) rational over~\(k\), let
\(i_y\colon\Spec(k)\to Y\) be the canonically associated morphism,
let \(j_y=\id_X\times_ki_y\colon X\to Z\), and let \(X_y'\) be the
\'etale covering of~\(X\) obtained from~\(Z'\) by inverse image
under~\(j_y\). Suppose that \(Y\) is connected, and that \(X\) or
\(Y\) is proper over~\(k\). Then all the coverings \(X_y'\) of~\(X\)
are isomorphic.
\end{corollary}

Figuratively, one may say that \emph{a family of \'etale coverings
of~\(X\), parametrized by a connected prescheme~\(Y\), is constant
if \(X\), or the parameter prescheme~\(Y\), is proper over~\(k\)}.

\begin{remarques}
\label{X.1.10}
Corollaries~\Ref{X.1.7} through~\Ref{X.1.9} are due to
Lang--Serre~\cite{X.2} in the case of normal algebraic schemes.
Their work was the initial motivation for the theory of the
fundamental group developed in this Seminar. As those authors
observed, these results become false when the properness hypothesis
is dropped, at least in characteristic \(p>0\). For example, take
\(X=\Spec(k[t])\), the affine line. The coverings of~\(X\)
parametrized by the affine line \(Y=\Spec(k[s])\), defined by the
equations
\[
  x^p-x=st,
\]
are \'etale and pairwise nonisomorphic. This contradicts
\Ref{X.1.9}, and a fortiori~\Ref{X.1.7}. Likewise, if \(s\) is
regarded as an element transcendental over~\(k\) in an algebraically
closed extension \(K\) of~\(k\), one obtains an \'etale covering
\(X'\) of \(X\otimes_kK\) which does not come from an \'etale
covering of~\(X\).
\end{remarques}

\end{filecontents*}
% Original input: drafts/SGA1_X_2_A_English.texfrag
\begin{filecontents*}[overwrite,noheader]{sga1-cumulative-input-135.tex}
\section[Application of the existence theorem for sheaves]
{Application of the existence theorem for sheaves: a semicontinuity theorem for the fundamental groups of the fibers of a proper and separable morphism}
\label{X.2}
% page: 268
\marginpar{268}
%

\begin{theorem}
\label{X.2.1}
Let \(Y\) be the spectrum of a \emph{complete} noetherian local ring
\(A\), with residue field~\(k\), let \(X\) be a proper \(Y\)-scheme,
let \(X_0=X\otimes_Ak\), let \(a_0\) be a geometric point of~\(X_0\),
and let \(a\) be the corresponding geometric point of~\(X\). Then
the canonical homomorphism
\[
  \pi_1(X_0,a_0)\longrightarrow\pi_1(X,a)
\]
is an \emph{isomorphism}.
\end{theorem}

This is merely the translation into the language of the fundamental
group of the result recalled in \ifthenelse{\boolean{orig}}{IX~}{}\Ref{IX.1.10}. It is here that the
existence theorem
\ifthenelse{\boolean{orig}}
{for}
{for}
sheaves in algebraic-formal geometry enters the theory of the
fundamental group in an essential way.

Now introduce an algebraic closure \(\overline k\) of the residue
field~\(k\), and the geometric fiber
\[
  \overline X_0=X_0\otimes_k\overline k.
\]
\ifthenelse{\boolean{orig}}{}{Choose a geometric point \(\overline a_0\) of~\(\overline X_0\) above~\(a_0\). }We have the exact sequence (\ifthenelse{\boolean{orig}}{IX~}{}\Ref{IX.6.1})
\[
\ifthenelse{\boolean{orig}}
{e\to\pi_1(\overline X_0,\overline a_0)
 \to\pi_1(X_0,a_0)\to\pi_1(k,\overline k)\to e,}
{e\to\pi_1(\overline X_0,\overline a_0)
 \to\pi_1(X_0,a_0)\to\pi_1(k,\overline k)\to e.}
\]
On the other hand, there is the isomorphism~\Ref{X.2.1}, together
with the analogous and more elementary isomorphism
\[
  \pi_1(k,\overline k)\longrightarrow\pi_1(Y,b),
\]
where \(b\) is the image of~\(a\) in~\(Y\). Thus we obtain:

\begin{corollary}
\label{X.2.2}
\ifthenelse{\boolean{orig}}
{With the preceding notation, suppose that \(\overline X_0\) is
connected, and let \(\overline a_0\) be a geometric point of
\(\overline X_0=X_0\otimes_k\overline k\), with image \(a_0\)
in~\(X\),}
{With the preceding notation, suppose that
\(\overline X_0=X_0\otimes_k\overline k\) is connected, and let
\(\overline a_0\) be a geometric point of~\(\overline X_0\), with
image \(a_0\) in~\(X\),}
and let \(b_0\) be its image in~\(Y\). Then the following sequence of
canonical homomorphisms
\[
\ifthenelse{\boolean{orig}}
{e\to\pi_1(\overline X_0,\overline a_0)
 \to\pi_1(X,a_0)\to\pi_1(Y,b_0)\to e}
{e\to\pi_1(\overline X_0,\overline a_0)
 \to\pi_1(X,a_0)\to\pi_1(Y,b_0)\to e}
\]
is exact.
\end{corollary}

Compare this sequence with the exact sequence~\Ref{X.1.4}, but note
that a) no flatness or fiberwise separability hypothesis was required
for \(X\to Y\); and b) one has the important additional fact that
\emph{the homomorphism
\(\pi_1(\overline X_0,\overline a_0)\to\pi_1(X,a_0)\) is injective}.

This last fact will allow us to compare the fundamental groups of the
other
% page: 269
\marginpar{269}
%
geometric fibers of~\(X\) over~\(Y\) with that of~\(\overline X_0\).
Let \(y_1\) be an arbitrary point of~\(Y\), let \(X_1\) be its fiber,
and let \(\overline X_1\) be its geometric fiber with respect to a
chosen algebraically closed extension of~\(\kres(y_1)\). Let
\(\overline a_1\) be a geometric point of~\(\overline X_1\), let
\(a_1\) be its image in~\(X\), and let \(b_1\) be its image in~\(Y\).
Choose a path class from~\(a_1\) to~\(a_0\), and hence a path class
from~\(b_1\) to~\(b_0\). This gives a commutative diagram of
homomorphisms
\[
  \xymatrix@=.5cm{
  &\pi_1(\overline X_1,\overline a_1)\ar[r]
  &\pi_1(X,a_1)\ar[r]\ar[d]
  &\pi_1(Y,b_1)\ar[r]\ar[d]&e\\
  e\ar[r]&\pi_1(\overline X_0,\overline a_0)\ar[r]
  &\pi_1(X,a_0)\ar[r]&\pi_1(Y,b_0)\ar[r]&e,
  }
\]
in which the two displayed vertical arrows are isomorphisms. Since
the second row is exact, we obtain a canonical homomorphism, which we
shall call the \emph{specialization homomorphism for the fundamental
group}:
\[
  \pi_1(\overline X_1,\overline a_1)
  \longrightarrow
  \pi_1(\overline X_0,\overline a_0).
\]
It depends only on the chosen path class from~\(a_1\) to~\(a_0\), and
is therefore \emph{defined up to an inner automorphism} of
\(\pi_1(X,a_0)\). When the first row above is exact as well, it
follows at once that the specialization homomorphism is surjective.
Taking~\Ref{X.1.4} into account, we therefore obtain:

\begin{corollary}
\label{X.2.3}
Under the hypotheses of~\Ref{X.2.1}, suppose in addition that
\(f\colon X\to Y\) is \emph{separable} \eqref{X.1.1} and that
\(\overline X_0\) is connected. Then
\(f_*(\cal O_X)=\cal O_Y\) by~\Ref{X.1.2}. For every geometric fiber
\(\overline X_1\) of~\(X\) over~\(Y\), endowed with a geometric
point \(\overline a_1\), the specialization homomorphism
%
\index{specialization homomorphism for the fundamental group|hyperpage}%
\index{specialization for the fundamental group (homomorphism)|hyperpage}%
%
defined above is \emph{surjective}.
\end{corollary}

This is a \emph{semicontinuity} result for the fundamental group,
which does not yet seem to have an analogue in algebraic topology.
It may moreover be stated under more general hypotheses:

\begin{corollary}
\label{X.2.4}
Let
% page: 270
\marginpar{270}
%
\(f\colon X\to Y\) be a proper morphism with geometrically connected
fibers, with \(Y\) locally noetherian. Let \(y_0,y_1\) be two points
of~\(Y\) such that \(y_0\in\overline{\{y_1\}}\), and let
\(\overline X_0,\overline X_1\) be the geometric fibers of~\(X\)
corresponding to chosen algebraically closed extensions of
\(\kres(y_0)\) and \(\kres(y_1)\). Let \(\overline a_0\), respectively
\(\overline a_1\), be a geometric point of \(\overline X_0\),
respectively \(\overline X_1\). Then one can naturally define a
specialization homomorphism
\[
  \pi_1(\overline X_1,\overline a_1)
  \longrightarrow
  \pi_1(\overline X_0,\overline a_0),
\]
well-defined up to inner automorphism. This homomorphism is
\emph{surjective} if \(f\) is separable \eqref{X.1.1}.
\end{corollary}

First, it follows from~\Ref{X.1.8} that~\Ref{X.2.4} is essentially
independent of the chosen algebraically closed extensions of the
residue fields \(\kres(y_0)\) and \(\kres(y_1)\). We may therefore
replace \(Y\) by a prescheme \(Y'\) over~\(Y\) having a point \(y_0'\)
above~\(y_0\), respectively a point \(y_1'\) above~\(y_1\). Take for
\(Y'\) the spectrum of the completion of the local ring of~\(Y\)
at~\(y_0\), and apply~\Ref{X.2.3}.

\begin{remarques}
\label{X.2.5}
The final conclusion of~\Ref{X.2.4}, concerning surjectivity of the
specialization homomorphism, and a fortiori the results
\Ref{X.1.3} and~\Ref{X.1.4} from which it follows, becomes false if
\(f\colon X\to Y\) is no longer assumed separable, even for projective
schemes over an algebraically closed field of characteristic~0. We
shall see examples below both when \(f\) is flat but has a
nonseparable fiber, although \(X\) and \(Y\) are smooth over~\(k\),
and when the fibers of~\(f\) are separable but \(f\) is not flat, for
example when \(f\colon X\to Y\) is a birational morphism of normal
integral schemes; see \ifthenelse{\boolean{orig}}{XI~}{}\Ref{XI.3}. In these examples, the
fundamental group of the geometric generic fiber may be trivial while
that of a suitable geometric special fiber is not.

Moreover, even if \(f\colon X\to Y\) is proper and separable as
in~\Ref{X.2.4}, the specialization homomorphism often fails to be
% page: 271
\marginpar{271}
%
an isomorphism. For example, it is easy to give cases in which
\(\overline X_1\) is a nonsingular elliptic curve, so its fundamental
group is commutative and its \(\ell\)-primary component, for a prime
\(\ell\) different from the characteristic, is isomorphic to
\(\ZZ_\ell^2\); see~\Ref{XI}. Meanwhile \(\overline X_0\) may consist
of two nonsingular rational curves meeting in two points, or of two
rational curves tangent at one point, or of a rational curve having
a singular point which is a cusp. For the complete classification of
degenerate elliptic curves, see the recent work of
Kodaira~\cite{X.1} and N\'eron. In these three cases the fundamental
group of~\(\overline X_0\) is respectively \(\widehat\ZZ\), \(e\),
and \(e\), and is therefore \emph{strictly smaller} than that of
\(\overline X_1\).

We shall nevertheless obtain below, when \(f\) is \emph{smooth}, a
bound on the kernel of the specialization homomorphism which implies
in particular that, if \(\kres(y_0)\) has characteristic~0, the
specialization homomorphism is an isomorphism. Even for a smooth
morphism, however, if \(\kres(y_0)\) has positive characteristic, the
specialization homomorphism need not be an isomorphism, as one sees
for example when \(X\) is an abelian scheme over~\(Y\), even of
relative dimension~1,
\ifthenelse{\boolean{orig}}
{\cf XI~\Ref{XI.2}).}
{\cf \ifthenelse{\boolean{orig}}{XI~}{}\Ref{XI.2}.}
% correction: delete the final parenthesis.
A satisfactory theory of specialization of the fundamental group
must take account of the \emph{continuous component} of the
\emph{true} fundamental group, corresponding to the classification
of principal coverings whose structural groups are infinitesimal
groups. One might then expect that the true fundamental groups of the
geometric fibers of a smooth proper morphism \(f\colon X\to Y\) form
a well-behaved local system on~\(X\), an inverse limit of finite flat
group schemes over~\(X\)\footnote{This extremely attractive
conjecture is unfortunately disproved by an unpublished example of
M\ptbl Artin, already when the fibers of~\(f\) are algebraic curves
of a fixed genus \(g\geq2\).}. We shall return later to this point of
view. Our present aim, by contrast, is to push as far as possible the
phenomena common to the topological and scheme-theoretic theories of
the fundamental group.
\end{remarques}

Now let \(X_0\) be a proper smooth connected curve of genus~\(g\)
over an algebraically closed field~\(k\). If \(k\) has characteristic
zero, its fundamental group may be determined by transcendental
methods as follows. We know that \(X_0\) is obtained by base extension
from a curve defined
% page: 272
\marginpar{272}
%
over an algebraically closed extension of finite transcendence degree
of the prime field~\(\QQ\). By~\Ref{X.1.8}, we may suppose that \(k\)
itself has finite transcendence degree over~\(\QQ\). We may therefore
regard \(k\) as a subfield of the complex numbers~\(\CC\), and a
further application of~\Ref{X.1.8} allows us to suppose \(k=\CC\).
It is then not difficult to verify that the fundamental group of
\(X_0\) is isomorphic to the profinite completion of the fundamental
group of the associated topological space \(\widetilde X_0\), a
compact oriented surface of genus~\(g\), with respect to the topology
defined by subgroups of finite index\footnote{This deduction was
explained in one of the oral expos\'es which were not written up.}.
It is classical that the topological fundamental group is generated
by \(2g\) generators \(s_i,t_i\), \(1\leq i\leq g\), subject to the
single relation
\[
  (s_1t_1s_1^{-1}t_1^{-1})\cdots
  (s_gt_gs_g^{-1}t_g^{-1})=1.
\]
Thus the fundamental group of~\(X_0\) has \(2g\) \emph{topological}
generators \(s_i,t_i\), \(1\leq i\leq g\), satisfying the preceding
relation.

Now suppose that \(k\) has characteristic \(p>0\). Let \(A\) be the
ring of Witt vectors constructed from~\(k\), and let \(K\) be an
algebraically closed extension of its fraction field. We saw in
\ifthenelse{\boolean{orig}}{III~}{}\Ref{III.7.4} that there exists a scheme \(X\) over
\(Y=\Spec(A)\), proper and smooth over~\(Y\), whose reduction is
\(X_0\). Applying~\Ref{X.2.3}, we obtain a \emph{surjective}
homomorphism
\[
  \pi_1(X_1)\longrightarrow\pi_1(X_0),
\]
where \(X_1=X\otimes_AK\). It is immediate\footnote{\Cf EGA~IV,
12.2.} that \(X_1\) is smooth over~\(K\), connected by~\eqref{X.1.2},
and of dimension~1, and that its genus equals~\(g\), by invariance of
the Euler--Poincar\'e characteristic; see EGA~III, 7. Since \(K\) has
characteristic~0, the preceding result applies. We have therefore
proved, by \emph{transcendental methods}:

\begin{theorem}
\label{X.2.6}
Let \(X_0\) be a smooth proper connected algebraic curve over an
% page: 273
\marginpar{273}
%
algebraically closed field~\(k\), and let \(g\) be its genus. Then
\(\pi_1(X_0)\) has a system of \(2g\) topological generators subject
to the relation displayed above. If \(k\) has characteristic~0,
\(\pi_1(X_0)\) is moreover the free group of Galois type on those
generators subject to that relation.
\end{theorem}

\begin{remarques}
\label{X.2.7}
At present, to the author's knowledge, there is no purely algebraic
proof of the preceding result, except in genera \(0\) and~\(1\). To
begin with, it is not at all clear how to distinguish \(2g\) elements
of \(\pi_1(X_0)\) which one might then expect to form a system of
topological generators. In this respect, the situation of the
rational line with \(n\) points removed, and the study of its
coverings tamely ramified at those points, is more congenial: the
ramification groups at the \(n\) points provide \(n\) elements of the
fundamental group under study, and one proves that they generate it
topologically, as we shall see later\footnote{\Cf
Expos\'e~\Ref{XII}. Even these elements are really determined only up
to conjugacy, and one must make a judicious \emph{simultaneous}
choice of representatives in their conjugacy classes.}. Yet even in
this particularly concrete case no purely algebraic proof seems to be
known. Such a proof would plainly be extremely interesting. The only
fact about the fundamental group of a curve that appears to admit a
purely algebraic proof, apart from the weak finiteness theorem
\Ref{X.2.12} below, proved algebraically by
Lang--Serre~\cite{X.2}, is the determination of its abelianization by
means of the Jacobian, mentioned in the last line of
\ifthenelse{\boolean{orig}}{IX~}{}\Ref{IX.5.8}.
\end{remarques}

\end{filecontents*}
% Original input: drafts/SGA1_X_2_B_English.texfrag
\begin{filecontents*}[overwrite,noheader]{sga1-cumulative-input-136.tex}
\subsection{}
\label{X.2.8}
The final assertion of~\Ref{X.2.6} no longer holds in
characteristic $p>0$, as one sees already in the case of elliptic
curves. As we have already pointed out, we do not know whether the
fundamental group of~$X_0$ is topologically finitely presented; this
seems wholly unlikely.

\begin{theoreme}
\label{X.2.9}
Let $k$ be an algebraically closed field, and let $X$ be a proper
connected scheme over~$k$. Then the fundamental group of~$X$ is
topologically finitely generated.
\end{theoreme}

We
% page: 274
\marginpar{274}
%
proceed by induction on~$n=\dim X$, the assertion being trivial for
$n\le 0$. Thus suppose that $n>0$, and that the theorem has been
proved in dimensions~$n'<n$. By Chow's lemma (EGA~II~5.6.2), there
exists a projective scheme $X'$ over~$k$ and a surjective morphism
$X'\to X$. We may clearly suppose that $X'$ is reduced and, on
passing to its normalization, normal. By descent theory, it suffices
to prove that the fundamental groups of the connected components
of~$X'$ are topologically finitely generated (\ifthenelse{\boolean{orig}}{IX~}{}\Ref{IX.5.2}). We
are thus reduced to the case in which $X$ is \emph{projective} and
\emph{normal}. If $n=1$, it is enough to apply~\Ref{X.2.6}. If
$n\ge 2$, consider a projective immersion $X\to \PP_k^r$ and a
hyperplane section $Y=X\cdot H$
\ifthenelse{\boolean{orig}}
{(equipped}
{(equipped}
% correction: `muni' should be `munie'.
with the induced reduced structure), such that
\ifthenelse{\boolean{orig}}
{$Y\ne X$}
{$Y\ne X$,}
\ie $H\not\supset X$. We then have $\dim Y<n$, and, taking the
induction hypothesis into account, it suffices to prove that
$\pi_1(Y)\to \pi_1(X)$ is \emph{surjective.} More generally:

\begin{lemme}
\label{X.2.10}
Let $X$ be a proper prescheme over an algebraically closed field~$k$,
and let $g\colon X\to \PP_k^r$ be a morphism. Suppose that $X$ is
irreducible and normal and that $\dim g(X)\ge 2$. Let $H$ be a
hyperplane of~$\PP_k^r$ and let $Y=X\times_{\PP_k^r} H$. Then $Y$ is
connected, and the homomorphism $\pi_1(Y)\to\pi_1(X)$ is surjective.
\end{lemme}

These assertions in fact follow from the following one:

\begin{corollaire}
\label{X.2.11}
Under the preceding conditions, let $X'$ be a connected étale
covering of~$X$, and let $Y'=X'\times_X Y=X'\times_{\PP_k^r} H$ be the
induced covering of~$Y$. Then $Y'$ is connected.
\end{corollaire}

Since $X$ is normal, $X'$ is normal; hence, being connected, $X'$ is
irreducible. Moreover, its image in $\PP_k^r$ has dimension~$\ge 2$.
A well-known lemma due to Zariski (called the \og \emph{Bertini
theorem}\fg) therefore implies that if $H_1'$ is the generic
hyperplane in~$\PP_k^r$, defined over an extension $K$ of~$k$, then
$X'\times_{\PP^r} H_1\ifthenelse{\boolean{orig}}{}{'}$ is universally irreducible, hence universally
connected over~$K$. Zariski's connectedness theorem (EGA~III~4) then
implies that for \emph{every} hyperplane $H$ (defined over an
arbitrary extension of~$k$), $X'\times_{\PP^r} H$ is geometrically
connected. This completes the proof of~\Ref{X.2.11}, and hence
of~\Ref{X.2.9}.

\begin{corollaire}[Lang-Serre]
\label{X.2.12}
Under
% page: 275
\marginpar{275}
%
the conditions of~\Ref{X.2.9}, for every finite group~$G$, the set of
isomorphism classes of principal coverings of~$X$ with group~$G$ is
finite.
\end{corollaire}

\begin{remarque}
\label{X.2.13}
Under the conditions of~\Ref{X.2.10}, when $\dim g(X)\ge 3$ (at
least when $g$ is an immersion and $X$ is regular), we shall prove a
more precise result, known in algebraic geometry as the
\og \emph{Lefschetz theorem}\fg: $\pi_1(Y)\to \pi_1(X)$ \emph{is an
isomorphism.}\footnote{\Cf the SGA~2 seminar (1962), which followed
this one%
\ifthenelse{\boolean{orig}}
{.}
{, Theorem X~3.10.}}
In the classical cases there are analogous statements for homology
groups and higher homotopy groups, which sooner or later will have to
be encompassed within abstract algebraic geometry. Even for Hodge
cohomology $\H^p(X,\Omega^q)$, the question does not seem to have
been studied yet; moreover, it is hardly likely that for the latter
the Lefschetz theorems remain unchanged in characteristic~$p>0$.
\end{remarque}

\ifthenelse{\boolean{orig}}
{}
{\begin{remarque}
\label{X.2.14}
Let $R$ be a complete discrete valuation ring with algebraically
closed residue field $\kres$ of characteristic $p>0$ and fraction
field~$K$, and let $Y$ be a proper smooth connected curve of genus
$g$ over~$R$. Thus there is a surjective specialization morphism
$\operatorname{sp} : \pi_1(Y_{\overline{K}})
\to\pi_1(Y_{\kres})$.
We have already noted that if $K$ has characteristic 0, $\operatorname{sp}$ is not
an isomorphism as soon as $g\geq1$. Suppose that $K$ has
characteristic $p$, so that $R$ is isomorphic to the formal power
series ring~$\kres[[T]]$.
In equal characteristic $p>0$, the fundamental group is not
determined by the genus~$g$, as one sees already from elliptic curves,
which may be ordinary or supersingular. We mention the recent result
of A\ptbl Tamagawa (not yet published). If~$G$ is a profinite group,
we denote by $G^{\res}$ the quotient profinite group of~$G$ given by
the inverse limit of the finite solvable topological quotients of~$G$.

\begin{theoremstar}[A\ptbl Tamagawa]
Suppose that $g\geq2$, that the special fiber $Y_{\kres}$ is
definable over a finite field, and that the morphism
$\operatorname{sp}^{\res}:\pi_1(Y_{\overline{K}})^{\res}\to\pi_1(Y_{\kres})^{\res}$ induced by
$\operatorname{sp}$ upon passing to the quotient is bijective. Then the curve $Y$
is constant over~$R$.
\end{theoremstar}

Note that the Galois group of~$\overline{K}/K$ is solvable. The
preceding statement may also be expressed as follows: suppose that
every finite étale covering $X_K\to Y_K$ of the generic fiber $Y_K$,
Galois with solvable Galois group, has potentially good reduction
(that is, extends to a finite étale covering of~$Y$ after possibly
replacing~$R$ by its normalization in a finite extension of~$K$).
Then $Y$ is constant over~$R$.
\end{remarque}}
\end{filecontents*}
% Original input: drafts/SGA1_X_3_English.texfrag
\begin{filecontents*}[overwrite,noheader]{sga1-cumulative-input-137.tex}
\section[Application of the purity theorem]{Application of the purity theorem:
continuity theorem for the fundamental groups of the
fibers of a proper and smooth morphism}
\label{X.3}

Recall without proof\footnote{For a proof, \cf SGA~2 \ifthenelse{\boolean{orig}}{X~}{}\Ref{X.3.4}.} the

\begin{theoremedepurete}[Zariski-Nagata]
\label{X.3.1}
%
\index{purity theorem@purity theorem|hyperpage}%
%
Let $f\colon X\to Y$ be a
quasi-finite and dominant morphism of integral preschemes, with $X$
normal and $Y$ regular and locally Noetherian, and let $Z$ be
the set of points of~$X$ where $f$ is not \'etale,
\ie where~$f$ is ramified (these are equivalent,
\textup{\ifthenelse{\boolean{orig}}{I~}{}\Ref{I.9.5}~(ii)}). If $Z\ne X$, then $Z$ has codimension~$1$ in
$X$ at all its points, \ie for every irreducible component
$Z'$ of~$Z$ with generic point~$z$, the Krull dimension of
$\cal{O}_{X,z}$ is equal to~$1$.
\end{theoremedepurete}

Recall that a prescheme is called \emph{normal}
\resp \emph{regular} if its local rings are normal
\resp regular, and that the relation $Z\ne X$ also means that
the finite extension \ifthenelse{\boolean{orig}}{$R(Z)/R(X)$}{$R(X)/R(Y)$} (where $R$ denotes the field of
rational functions) is \emph{separable}. Passing to
the generic point $z$ of a component $Z'$ of~$Z$, and
localizing at the point $y$ of~$Y$ below~$z$, one obtains
% page: 276
\marginpar{276}
%
the equivalent statement:

\begin{corollaire}
\label{X.3.2}
Let $A$ be a regular Noetherian local ring, and let $A\to B$ be an
injective local homomorphism such that $B$ is normal, a localization
of an algebra of finite type over~$A$, and \emph{quasi-finite} over~$A$;
suppose in addition that $\dim A (=\dim B)\ge 2$, and that for every
prime ideal $\goth{p}$ of~$B$ distinct from the maximal ideal,
$B$ is \'etale over~$A$ at~$\goth{p}$, \ie $B_\goth{p}$ is
\'etale over~$A_\goth{q}$ (where $\goth{q}=A\cap\goth{p}$). Then
$B$ is \'etale over~$A$.
\end{corollaire}

Moreover, it is not difficult to reduce this last
statement to the case where $A$ is a \emph{complete} local ring,
and hence where $B$ is \emph{finite} over~$A$. Zariski~\cite{X.5} gives a
simple proof of this result, valid in the equal-characteristic
case; the general case is due
to Nagata~\cite{X.3}, who relies on a delicate result of
Chow; the latter has not been verified by any of the
participants in the Seminar, and ought to be the subject of a
later lecture. We mention here only the very simple proof
in the special case where $\dim A=2$, which is
sufficient for the most important application we shall make of it
in the present section. Since $B$ is normal, it is a
$B$-module of depth (older terminology: cohomological
codimension) $\ge 2$, and hence it is an $A$-module of depth~$\ge
2$; since $A$ is regular of dimension $2$, it follows
that $B$ is a \emph{free module} over~$A$\kern1pt\footnote{\Cf EGA $0_{\textup{IV}}$
17.3.4}. It then follows from \ifthenelse{\boolean{orig}}{I~}{}\Ref{I.4.10} that the set
of prime ideals~$\goth{q}$ of~$A$ at which $B$ is
ramified over~$A$ is the subset of~$\Spec(A)$ defined by a
principal ideal (generated by the discriminant of a basis of
$B$ over~$A$), and is therefore empty if it is contained in the closed
point of~$\Spec(A)$, which proves~\Ref{X.3.2} when $\dim A=2$.

We shall chiefly use~\Ref{X.3.1} in the equivalent form:

\begin{corollaire}
\label{X.3.3}
Let $X$ be a regular locally Noetherian
prescheme, and let $U$ be an open subset of~$X$ complementary
to a closed subset $Z$ of~$X$ of codimension~$\ge 2$. Then the
functor $X'\mto X'\times_X U$ from the category of
\'etale coverings of~$X$ to the category of
\'etale coverings of~$U$ is an equivalence
% page: 277
\marginpar{277}
%
of categories; in particular, if $a$ is a geometric point
of~$U$, the canonical homomorphism $\pi_1(U,a)\to
\pi_1(X,a)$ is an isomorphism.
\end{corollaire}

The last assertion is plainly a consequence of the
first, and for the latter we may plainly suppose that $X$
is connected, and hence irreducible. The normality of~$X$
already implies that the functor $X'\mto X'\times_X U$ from
the category of locally free coverings (not
necessarily \'etale) of~$X\ifthenelse{\boolean{orig}}{'}{}$ to the category of
coverings of~$U$ is fully faithful, since the functor
$\cal{E}\mto\cal{E}\vert U$ from the category of locally free Modules
on~$X$ to the category of locally
free Modules on~$U$ is so. It therefore remains to prove that for every
\emph{\'etale} covering $U'$ of~$U$, there exists an
\'etale covering $X'$ of~$X$ (necessarily unique
by what precedes) such that $U'$ is isomorphic to
$X'\times_X U$. We may plainly suppose that $U'$ is connected, and hence
irreducible since ($U$~being normal) $U'$ is normal. Let
$K$ be the field of rational functions on~$X$, or on~$U$ (it is
the same), and let $K'$ be that of~$U'$; then $U'$ identifies with the normalization
of~$U$ in~$K'$ \ifthenelse{\boolean{orig}}{I~}{}\Ref{I.10.3}. Let $X'$ be the normalization of
$X$ in~$K'$ (EGA~II~6.3); then \ifthenelse{\boolean{orig}}{$X'=X\times_X U \cong U'$}{$X'\times_X U \cong U'$}, while
$X'$ is normal and integral, and the structural morphism $f\colon
X'\to X$ is \emph{finite} and dominant (since $X$ is normal and $K'/K$ is
a finite separable extension). It is \'etale on
\ifthenelse{\boolean{orig}}{$U'=f^{-1}(U)=X'=f^{-1}(Z)$}{$U'=f^{-1}(U)$, and $X'-U'=f^{-1}(Z)$}, and since $Z$ has codimension $\ge 2$
in $X$, $f^{-1}(Z)$ has codimension $\ge 2$ in~$X'$. It then follows
from~\Ref{X.3.1} that $X'$ is \'etale over~$X$, which completes
the proof.

\medskip
Now let $f\colon X\to Y$ be a rational map from a
regular locally Noetherian prescheme $X$
to a prescheme~$Y$, and suppose that $f$ is defined
on an open subset $U$ complementary to a closed subset of
codimension~$\ge 2$. Then~\Ref{X.3.3} yields a functor,
defined up to isomorphism, from the category of
\'etale coverings of~$Y$ to the category of
\'etale coverings of~$X$, and hence, for every geometric point
$a$ of~$U$, with image $b$ in~$Y$, a canonical homomorphism
$$
\pi_1(X,a)\to \pi_1(Y,b)
$$
(deduced
% page: 278
\marginpar{278}
%
from the canonical homomorphism $\pi_1(U,a)\to \pi_1(Y,b)$ by means
of the isomorphism $\pi_1(U,a)\isomto \pi_1(X,a))$. When $f$ is
a dominant morphism, with $X$ and $Y$ integral with function fields $K$
and~$L$, so that $K$ is an extension of~$L$, and with $Y$
normal, these correspondences may be made more precise in terms of field extensions
by noting that for every finite extension $L'$ of~$L$ that is
unramified over~$Y$, the algebra $K'=L'\otimes_L K$ over~$K$ is unramified
over~$X$.

In particular, these observations show that the fundamental group
of connected regular locally Noetherian preschemes,
pointed by geometric points located in codimension~$\le 1$, is a
\emph{functor} when the morphisms in this category are taken to be
dominant rational maps defined on the complements of
closed subsets of codimension~$\ge 2$. Recalling, for example,
that a rational map from a normal scheme over a field $k$
to a proper scheme over~$k$ is defined on the
complement of a subset of codimension~$\ge 2$, one obtains:

\begin{corollaire}[Birational invariance of the fundamental group]
\label{X.3.4}
%
\index{birational invariance of the fundamental group|hyperpage}%
\index{birational invariance theorem for the fundamental group@birational invariance theorem for the fundamental group|hyperpage}%
%
Let $k$ be a field, let $X$ and $Y$ be two regular schemes proper over~$k$,
let $f\colon X\to Y$ be a birational map from~$X$
to~$Y$, and let $\Omega$ be an algebraically closed extension of the function field
$K$ of~$X$, making it possible to define the fundamental group
of~$X$ and the fundamental group of~$Y$. These groups are then
canonically isomorphic.
\end{corollaire}

This also means that if a finite extension $K'$ of~$K$
is unramified over one proper nonsingular \og model\fg $X$
of~$K$, then it is unramified over every other proper nonsingular model.

\begin{remarque}
\label{X.3.5}
For other applications of the purity theorem, see the
works
\ifthenelse{\boolean{orig}}
{of \textsc{Abhyankar}}
{of \textsc{Abhyankar}}
presented in~\cite{X.4}, inspired by the
results of \textsc{Zariski} \cite[chap\ptbl VIII]{X.6}, proved by
topological methods. The latter are far from having been assimilated
by \og abstract\fg algebraic geometry and deserve
further effort.
\end{remarque}

We
% page: 279
\marginpar{279}
%
shall need a few elementary facts from ramification theory.
Let $V$ be a discrete valuation ring with fraction
field~$K$ and residue field~$k$, let $L$ be a Galois extension
of~$K$ with group~$G$, and let $V'$ be the normalization of~$V$
in~$L$, which is a free module of rank $n=[L{\colon }K]$ over~$V$;
let $\goth{m'}$ be a maximal ideal of~$V'$, let $G_d$ be the subgroup of~$G$
consisting of the elements that leave $\goth{m'}$ invariant, so
that $G_d$ acts on the residue extension
$k'=V'/\goth{m'}$ of~$k$, and let $G_i$ be the subgroup of elements
of~$G_d$ that act trivially (recall that $G_d$ and $G_i$ are
called, respectively, the \emph{decomposition} and
\emph{inertia} subgroups of~$G$). One says that $L$ is \emph{tamely
ramified} over~$V$ if $n_i=[G_i{\colon }e]$ is prime to the
characteristic $p$ of~$k$ (a condition that is always
satisfied if $k$ has characteristic~0). It is well
known that $G_i$ then embeds canonically into the group ${k'}^*$,
and is therefore isomorphic to the group of
\ifthenelse{\boolean{orig}}
{$n_i$th}
{$n_i$-th} roots of unity
in~${k'}^*$; this implies in particular that $G_i$ \emph{is
cyclic}. The standard example of this situation is obtained by setting
$L=K[t]/(t^n - u)$, where $u$ is a uniformizer of~$V$ and $n$ is an
integer prime to~$p$: if $K$ contains the \(n\)-th roots of
unity, $L$ is a totally ramified Galois extension
of~$K$, with Galois group $G=G_i$ isomorphic to~$\ZZ/n\ZZ$.

\begin{lemme}[%
\ifthenelse{\boolean{orig}}
{\textsc{Abhyankar}'s lemma}
{\textsc{Abhyankar}'s lemma}]
\label{X.3.6}
%
\index{Abhyankar (lemma of)@Abhyankar's lemma|hyperpage}%
\index{Abhyankar's lemma|hyperpage}%
%
Let $V$ be a discrete valuation ring with
fraction field~$K$, let $L$ and $K'$ be two Galois extensions
of~$K$ that are \emph{tamely ramified} over~$V$, and let $n$ and $m$ be the
orders of the corresponding inertia groups
\ifthenelse{\boolean{orig}}
{.}
{,}
% correction: replace . by ,
Let $L'$ be a composite extension
of~$L$ and $K'$ over~$K$. If $m$ is a multiple of~$n$,
then $L'$ is unramified over the localizations of the
integral closure $V'$ of~$V$ in~$K'$.
\end{lemme}

Indeed, let $W'$ be the normalization of~$V'$ in~$L'$, let $\goth{m'}$
be a maximal ideal of~$V'$, and let $\goth{n'}$ be a maximal ideal of~$W'$ lying above
\ifthenelse{\boolean{orig}}
{$\goth{n'}$,}
{$\goth{m'}$,}
% correction: `n'' should be `m''.
let $\goth{n}$ be the maximal ideal it induces on the normalization $W$
of~$V$ in~$L$, let $G,\ H,\ M$ be the Galois groups of~$L,\ K',\ L'$
over~$K$, and let $G_i,\ H_i,\ \ifthenelse{\boolean{orig}}{L_i'}{M_i}$ be the inertia groups corresponding to the
chosen maximal ideals. Then $M$ embeds into the product
$G\times H$ and $M_i$ into the product $G_i\times H_i$, in such a way that the projections $M\to G$ and $M\to H$, $M_i\to G_i$
% page: 280
\marginpar{280}
%
and $M_i\to H_i$ are surjective (the intermediate-field
argument). It already follows, since $G_i$ and
$H_i$ are by hypothesis cyclic of orders \ifthenelse{\boolean{orig}}{$m$ and $n$}{$n$ and $m$} prime
to~$p$, that $M_i$ has order prime to~$p$, and hence is cyclic; and
since $m$ is a multiple of~$n$, the elements of~$G_i\times
H_i$ have
\ifthenelse{\boolean{orig}}
{trivial $m$th}
{trivial $m$-th} power, so $M_i$ has order
dividing~$m$, and therefore has order equal to $m$ since $M_i\to H_i$ is
surjective. This last homomorphism is therefore also
injective. But its kernel is the inertia group of~$\goth{n'}$
over~$\goth{m'}$, which proves that $L'$ is unramified
over~$K'$ at~$\goth{n'}$. This proves the lemma.

Let us now place ourselves under the conditions of~\Ref{X.2.4},
where we have a specialization homomorphism
$$
\pi_1(\overline X_1,\overline a_1)\to\pi_1(\overline X_0,\overline
a_0)
$$
that is \emph{surjective}, relative to a proper and
separable morphism $f\colon X\to Y$. We wish to determine the kernel of
this homomorphism. Proceeding as in the proof
of~\Ref{X.2.4}, we see that for this question we may always
suppose that $Y$ is the spectrum of a \emph{complete discrete
valuation ring~$V$ with algebraically closed residue
field} (since one can always find such a ring
and a morphism from its spectrum $Y'$ to $Y$ whose image is
$\{y_0,y_1\}$). Then $X_0=\overline X_0$, $\kres(y_0)=k=$ the residue
field of~$V$, and $\kres(y_1)=K=$ the fraction field of~$V$. Let
$K_s$ be the separable closure of~$K$, let $\overline K$ be its
algebraic closure, and for every subring $W$ of~$\overline
K$ containing~$V$, put $X_W=X\otimes_V W$. In particular, we have
$$
X_V=\ifthenelse{\boolean{orig}}{V}{X}, \qquad X_K=X_1, \qquad X_{\overline K}=\overline X_1.
$$
Moreover, the canonical morphism $\overline X_1=X_{\overline K}\to
X_{K_s}$ induces an isomorphism on fundamental groups
(\ifthenelse{\boolean{orig}}{IX~}{}\Ref{IX.4.11}); hence, in view of
the isomorphism~\Ref{X.2.1} $\pi_1(X_0)\to\pi_1(X)$, we are reduced
to studying the surjective homomorphism
$$
\pi_1(X_{K_s})\to\pi_1(X)
$$
associated
% page: 281
\marginpar{281}
%
with the canonical morphism $X_{K_s}\to X$. Determining the kernel
of the latter amounts to solving the following problem:
\emph{given a connected principal covering $Z_{K_s}$ of
$X_{K_s}$ with group~$G$} (thus associated with a homomorphism from
$\pi_1(X_{K_s})$ to~$G$), \emph{determine under what
conditions it is isomorphic to the inverse image of a
principal covering $Z$ of~$X$ with group~$G$.}

First note that $K_s$ is the increasing filtered union of its
finite subextensions $K'$ over~$K$, and consequently $Z_{K_s}$
is isomorphic to the inverse image of a principal covering
$Z_{K'}$ of~$X_{K'}$ for a suitable $K'$ (note,
however, that for fixed $K'$, $Z_{K'}$ is not uniquely
determined). To say that $Z_{K_s}$ is isomorphic to the inverse
image of a principal covering $Z$ of~$X$ means that there exists
a finite subextension $K''\supset K'$ of~$K_s$ such that
$Z_{K''}=Z_{K'}\otimes_{K'} K''$ is isomorphic to $Z\otimes_V
K''$. For a finite subextension $K'$
of~$K_s$, now denote by $V'$ the normalization of~$V$ in~$K'$; this is a
complete discrete valuation ring with residue
field~$k$. Thus the canonical morphism $X_{V'}\to X_V$ induces
an isomorphism on the fibers over the closed points of
$Y=\Spec(V)$ and $Y'=\Spec(V')$, and it then follows from~\Ref{X.2.1},
applied to $X_V$ and $X_{V'}$, that the induced homomorphism on
fundamental groups $\pi_1(X_{V'})\to \pi_1(X_V)$ is an
isomorphism, or equivalently that every principal covering of~$X_{V'}$
is the inverse image of a principal covering of~$X_V$
determined up to isomorphism. This therefore implies the

\begin{lemme}
\label{X.3.7}
Let $Z_{K'}$ be a connected principal covering of~$X_{K'}$ with
group~$G$, and let $Z_{K\ifthenelse{\boolean{orig}}{}{_s}}$ be its inverse image over~$X_{K_s}$\quoi. This
latter covering is isomorphic to the inverse image of a
principal covering $Z$ of~$X$ if and only if there exists a
finite extension $K''\supset K'$ of~$K$ contained in $K_s$ such that the principal
covering $Z_{K''}$ of~$X_{K''}$ is induced by a principal covering
of~$X_{V''}$.
\end{lemme}

Suppose in particular that the $X_{V''}$ are normal (it suffices
for this, for example, that $X_0$ be normal, and a fortiori that $X_0$
be simple, \cf \ifthenelse{\boolean{orig}}{I~}{}\Ref{I.9.1}).
% page: 282
\marginpar{282}
%
Since they are connected, they are therefore irreducible. Let $L$ be the
field of rational functions of $X$ and $X_{K}$, let $L'$ be that of
$X_{V'}$ and $X_{K'}$, and let $L''$ be that of $X_{V''}$ and~$X_{K''}$\quoi. Then
under the conditions of~\Ref{X.3.7}, $Z_{K'}$ defines a finite
separable extension $R'$ of~$L'$, and $Z_{K''}$ defines the extension
$R''=R'\otimes_{L'} L'' = R'\otimes_{K'} K''$. The condition
considered in~\Ref{X.3.7} thus also means that there exists a
finite separable extension $K''$ of~$K'$ such that
$R''=R'\otimes_{K'} K''$ is \emph{unramified} over the
normal scheme $X_{V''}$ with function field $L''=L'\otimes_{K'} K''$, and not
only over the open subset $X_{K''}$ of~$X_{V''}$.

From now on, suppose that $f\colon X\to Y$ is a
\emph{smooth} morphism, so the morphisms $X_{V'}\to \Spec(V')$ are smooth,
and hence the schemes $X_{V'}$ are \emph{regular}. Note that the
fiber of the closed point of~$\Spec(V')$ in $X_{V'}$ is
irreducible and of codimension~$1$. Let $\goth{o'}$ be its local
ring; it is therefore a discrete valuation ring with fraction field
$L'$ and residue field isomorphic to the field of rational
functions of~$X_0$, and thus of the same characteristic as
$k$. Define $\goth{o''}$ similarly in $L''$; it is
plainly the normalization of
\ifthenelse{\boolean{orig}}
{$\goth{o''}$}
{$\goth{o'}$}
% correction: $\goth{o'}$
in~$L''$. It
then follows from the purity theorem~\Ref{X.3.1}
or~\Ref{X.3.3} that $R''$ is unramified over~$X_{V''}$ if and
only if $R''$ is unramified over~$\goth{o''}$,
the normalization of~$\goth{o'}$ in~$L''$.

\ifthenelse{\boolean{orig}}{Now note that if $u'$ is a uniformizer of~$V'$, it is}{Let $n$ be an integer prime to the characteristic $p$ of~$k$. After enlarging $K'$ within~$K_s$, assume that $K'$ contains the $n$-th roots of unity.}
\ifthenelse{\boolean{orig}}{also a uniformizer of~$\goth{o'}$. Thus, if $n$ is an integer}{If $u'$ is a uniformizer of~$V'$, it is also a uniformizer of~$\goth{o'}$.}
\ifthenelse{\boolean{orig}}{prime to the characteristic $p$ of~$k$, and if we take}{If we take}
$K''=K'[t]/(t^n-u\ifthenelse{\boolean{orig}}{}{'})$, then
\ifthenelse{\boolean{orig}}
{$K'$}
{$K''$}
% correction: $K'$ should be $K''$
is a finite Galois extension
of~$K'$, and $L''$ is isomorphic to $L'[t]/(t^n-u')$, and hence is
tamely ramified over~$\goth{o'}$ with inertia group
of order~$n$. Now suppose that $G$ \emph{has order prime to~$p$;} this
implies that $R'$ is tamely ramified over~$\goth{o'}$, and choose for $n$
a multiple, prime to $p$, of the order of the inertia group of~$R'$
over~$\goth{o'}$ (for example, $n=[G\colon e]$). Applying the
\ifthenelse{\boolean{orig}}
{\textsc{Abhyankar}'s lemma}
{\textsc{Abhyankar}'s lemma}~\Ref{X.3.6}, we
see that the condition considered in~\Ref{X.3.7} is
satisfied.

This proves the following theorem:

\begin{theoreme}
\label{X.3.8}
Let $f\colon X\to Y$ be a proper and smooth morphism with
% page: 283
\marginpar{283}%
%
geometrically connected fibers, with $Y$ locally Noetherian,
and let $y_0$ and $y_1$ be two points of~$Y$ such that
\ifthenelse{\boolean{orig}}
{$y_0\in\overline{y_1}$}
{$y_0\in\overline{\{y_1\}}$};
let $\overline{X}_0$ and $\overline{X}_1$ be the corresponding geometric fibers,
and consider the specialization homomorphism
\eqref{X.2.4} $\pi_1(\overline X_1)\to\pi_1(\overline X_0)$. This
homomorphism is surjective, and every continuous homomorphism from
$\pi_1(\overline X_1)$ to a finite group $G$ of order prime to
the characteristic $p$ of~$\kres(y_0)$ comes from a homomorphism from
$\pi_1(\overline{X}_0)$ to~$G$.
\end{theoreme}

In other words:

\begin{corollaire}
\label{X.3.9}
If $\kres(y_0)$ has characteristic zero, then the specialization
homomorphism is an isomorphism. If
\ifthenelse{\boolean{orig}}
{$p>0$}
{$\kres(y_0)$ has characteristic $p>0$},
then the kernel of
the specialization homomorphism is contained in the intersection
of the kernels of the continuous homomorphisms from~$\pi_1(\overline X_1)$ to
finite groups of order prime to~$p$ (or equivalently, in the closed
normal subgroup generated by a Sylow $p$-subgroup of the
group of Galois type $\pi_1(\overline X_1)$); thus, if
\ifthenelse{\boolean{orig}}
{$\pi_1(X_1)^{(p)}$}
{$\pi_1(\overline X_1)^{(p)}$}
% correction: $\pi_1(X_1)^{(p)}$ should be $\pi_1(\overline X_1)^{(p)}$
denotes the quotient group of
$\pi_1(\overline X_1)$ by the preceding closed subgroup,
and if $\pi_1(\overline X_0)^{(p)}$ is defined similarly, then
the specialization homomorphism induces an \emph{isomorphism}
$$
\pi_1(\overline X_1)^{(p)}\isomto \pi_1(\overline X_0)^{(p)}
$$
\end{corollaire}

Note that the proof of~\Ref{X.3.8} is purely algebraic.
Proceeding as in~\Ref{X.2.6}, we conclude \emph{by
transcendental methods:}

\begin{corollaire}
\label{X.3.10}
Let $X_0$ be a proper, smooth, and connected curve of genus $g$ over an
algebraically closed field of characteristic~$p$. With
the notation introduced in~\Ref{X.3.9}, the group $\pi_1(X_0)^{(p)}$
is isomorphic to $\Gamma^{(p)}$, where $\Gamma$ is the group of
Galois type generated by elements $s_i,~t_i$
($1\le i\le g$) subject to the relation
$$
(s_1t_1s_1^{-1}t_1^{-1})\cdots (s_gt_gs_g^{-1}t_g^{-1})=1
$$
\end{corollaire}

\begin{remarques}
\label{X.3.11}
% page: 284
\marginpar{284}
%
When $\kres(y_0)$ has characteristic zero, the
result~\Ref{X.3.9} is well known by transcendental methods. Note
that the proof of~\Ref{X.3.10} invokes the
purity theorem in the mixed-characteristic case,
but only for rings of dimension 2, where the proof of that theorem is
easy and was recalled in the text.
\end{remarques}

%\section{Bibliography}
\begin{thebibliography}{0}

\bibitem[1]{X.1} K\ptbl \textsc{Kodaira}, On compact analytic surfaces,
\ifthenelse{\boolean{orig}}
{Princeton University Press, 1960.}
{Ann.\ of Math.\ \textbf{71} (1960), p\ptbl 111--152.}

\bibitem[2]{X.2} S\ptbl \textsc{Lang} and
\ifthenelse{\boolean{orig}}{J.P\ptbl}{J-P\ptbl}%
\textsc{Serre}, Sur les rev\^etements non
ramifi\'es des vari\'et\'es alg\'ebriques,
\ifthenelse{\boolean{orig}}
{Amer.\ Journ. of Math., pp\ptbl 319--330, 1957.}
{Amer.\ J.\ Math.\ \textbf{79} (1957), p\ptbl319--330.}

\bibitem[3]{X.3} M\ptbl \textsc{Nagata},
On the purity of branch loci in regular local rings,
\ifthenelse{\boolean{orig}}
{Illinois Journal of Math., p\ptbl 328--333, 1959.}
{Illinois J.\ Math. \textbf{3} (1959), p\ptbl 328--333.}

\bibitem[4]{X.4}
\ifthenelse{\boolean{orig}}{J.P\ptbl}{J-P\ptbl}%
\textsc{Serre}, Rev\^etements ramifi\'es du plan
projectif (d'apr\`es S\ptbl Abhyankar), S\'eminaire Bourbaki, Mai
1960.

\bibitem[5]{X.5} O\ptbl \textsc{Zariski},
On the purity of the branch locus of algebraic functions,
\ifthenelse{\boolean{orig}}
{Proc.\ Nat.\ Acad.\ Sc.\ USA, t\ptbl 44, p\ptbl 791--796, 1958.}
{Proc.\ Nat.\ Acad.\ Sc.\ USA  \textbf{44} (1958), p\ptbl 791--796.}


\bibitem[6]{X.6} O\ptbl \textsc{Zariski},
Algebraic surfaces, Ergebnisse... 1948, Chelsea (New York).
\end{thebibliography}
\end{filecontents*}
% Original input: drafts/SGA1_X_OPENING_AND_1_A_English.texfrag
\begin{filecontents*}[overwrite,noheader]{sga1-cumulative-input-138.tex}
\chapter{Specialization theory of the fundamental group}
\label{X}
% page: 261
\marginpar{261}
%
In the present expos\'e we restrict ourselves to studying the
fundamental group of geometric fibers in a \emph{proper} morphism,
\ie the fundamental group of a varying \emph{proper} algebraic
scheme. In a later expos\'e we shall extend the technique used here
to \'etale coverings \emph{tamely ramified} at infinity. This will
give, for example, a solution of the \emph{three-point problem} for
Galois coverings of order prime to the characteristic (\ie a
determination of the Galois coverings of the line \(\PP_k^1\)
ramified at no more than three specified points and tamely ramified
at those points), together with its evident variants.

\section[The homotopy exact sequence for a proper morphism]
{The homotopy exact sequence for a proper and separable morphism}
\label{X.1}

\begin{definition}
\label{X.1.1}
A prescheme \(X\) over a field~\(k\) is called \emph{separable}, or
\emph{separable over~\(k\)}, if \(X\otimes_kK\) is reduced for every
extension \(K\) of~\(k\). If \(f\colon X\to Y\) is a morphism of
preschemes, we say that \(f\) is \emph{separable}, or that \(X\)
\emph{is separable over~\(Y\)}, if \(X\) is flat over~\(Y\) and, for
every \(y\in Y\), the fiber \(X\otimes_Y\kres(y)\) is separable
over~\(\kres(y)\).
\end{definition}

For a prescheme \(X\) over a field~\(k\), separability also means that
\(X\) is \emph{reduced} and that the fields \(\kres(x)\), for \(x\)
the generic point of an irreducible component of~\(X\), are separable
extensions of~\(k\). Thus, if \(k\) is perfect, this is equivalent to
saying that \(X\) is reduced. Observe that if \(X\) is separable
over~\(Y\), then after every base change \(Y'\to Y\),
\(X'=X\times_YY'\) is separable over~\(Y'\). Under suitable
finiteness hypotheses one can also prove that the composite of two
separable morphisms is
% page: 262
\marginpar{262}
%
separable. We shall need this only in the following form:
\emph{if \(X\) is separable over~\(Y\), and \(X'\) is \'etale
over~\(X\), then \(X'\) is separable over~\(Y\)}. This is an
immediate consequence of the definitions and \ifthenelse{\boolean{orig}}{I~}{}\Ref{I.9.2}. The
hypothesis of a separable morphism will otherwise enter through the
following proposition:

\begin{proposition}
\label{X.1.2}
Let \(f\colon X\to Y\) be a proper and separable morphism, with \(Y\)
locally noetherian, and consider its \emph{Stein} factorization
\[
  X\lto{f'}Y'\dpl\to Y,
\]
where \(f_*'(\cal O_X)=\cal O_{Y'}\), and where \(Y'\), finite
over~\(Y\), is isomorphic to the spectrum of the algebra
\(f_*(\cal O_X)\). Then \(Y'\) is an \emph{\'etale covering}
of~\(Y\).
\end{proposition}

This proposition appears in EGA~III, 7\footnote{\Cf EGA~III,
7.8.10~(i).}. Let us indicate the principle of the proof. One readily
reduces to the case where \(Y\) is the spectrum of a complete local
ring~\(A\). After making a suitable finite flat extension of~\(A\),
corresponding to a suitable residue-field extension, one may suppose
that the connected components of the fiber over the closed point
\(y\) are geometrically connected. This also means that
\(\H^0(X_y,\cal O_{X_y})\) decomposes as a product of fields equal
to \(k=\kres(y)\). We may then suppose \(X\) connected, in which case
\[
  \H^0(X_y,\cal O_{X_y})=k.
\]
Thus the homomorphism
\[
  A\longrightarrow\H^0(X_y,\cal O_{X_y})
\]
is \emph{surjective}. A general proposition of K\"unneth type then
shows that \(f_*(\cal O_X)\) is defined by an \(A\)-module~\(B\)
which is free over~\(A\), and that
\[
  B/\goth mB\longrightarrow\H^0(X_y,\cal O_{X_y})=k
\]
is bijective. In the present case \(B\) is therefore an \'etale
algebra over~\(A\), which completes the proof.

\begin{theorem}
\label{X.1.3}
Let \(f\colon X\to Y\) be a proper and separable morphism, with \(Y\)
locally noetherian and connected, and suppose that
\(f_*(\cal O_X)=\cal O_Y\). This implies that the fibers of~\(X\)
over~\(Y\) are geometrically connected, and the converse follows
from~\Ref{X.1.2}. Let \(y\) be a point of~\(Y\), let
\(\overline{\kres(y)}\) be an algebraic closure of~\(\kres(y)\), and
put
\[
  \overline X_y
  =X_y\otimes_{\kres(y)}\overline{\kres(y)}.
\]
Finally, let \(X'\) be a \emph{connected} \'etale covering of~\(X\),
and put
\[
  \overline X_y'
  =X_y'\otimes_{\kres(y)}\overline{\kres(y)}.
\]
There exist an \'etale covering \(Y'\) of~\(Y\) and an
\(X\)-isomorphism
% page: 263
\marginpar{263}
%
\[
  X'\isomto X\times_YY'
\]
if and only if \(\overline X_y'\) has a section over
\(\overline X_y\).
\end{theorem}

Set \(Y'=\Spec(h_*(\cal O_{X'}))\), where \(h\colon X'\to Y\) is the
composite \(X'\to X\to Y\). It is enough to prove that the canonical
\(Y\)-morphism
\[
  X'\longrightarrow X\times_YY'
\]
is an \emph{isomorphism}, and that \(Y'\) is \'etale over~\(Y\).
By~\Ref{X.1.2}, we already know that \(Y'\) is \'etale over~\(Y\).
Consequently \(X\times_YY'\) is \'etale over~\(X\), and hence
\(X'\to X\times_YY'\) is likewise \'etale
(\ifthenelse{\boolean{orig}}{I~}{}\Ref{I.4.8}). Moreover, \(Y'\) is connected, as the image of the
connected prescheme~\(X'\), so \(X\times_YY'\) is connected because
\(X\) has connected fibers over~\(Y\)
(\ifthenelse{\boolean{orig}}{IX~}{}\Ref{IX.3.4} and \ifthenelse{\boolean{orig}}{V~}{}\Ref{V.6.9}~(iii)). Thus, to prove that
\(X'\to X\times_YY'\) is an isomorphism, it is enough to show that
its degree at \emph{one} point of~\(X\times_YY'\) equals~1. This
follows readily from the hypothesis that \(\overline X_y'\) has a
section over~\(\overline X_y\), either by using \ifthenelse{\boolean{orig}}{IX~}{}\Ref{IX.6.6}, or
more simply by observing that it is enough to prove the existence of
such a point in \(X\times_YY'\) after the base change
\ifthenelse{\boolean{orig}}
{\(\Spec(\overline k)\to Y\)}
{\(\Spec(\overline{\kres(y)})\to Y\),}
where it is evident. This completes the proof of~\Ref{X.1.3}.

Taking account of \ifthenelse{\boolean{orig}}{IX~}{}\Ref{IX.3.4} and the dictionary
\ifthenelse{\boolean{orig}}{V~}{}\Ref{V.6.9} and \ifthenelse{\boolean{orig}}{V~}{}\Ref{V.6.11}, we may put~\Ref{X.1.3} in the
following equivalent form:

\begin{corollary}
\label{X.1.4}
With the preceding notation for \(f\colon X\to Y\) and
\(\overline X_y\), let \(\overline a\) be a geometric point of
\(\overline X_y\), let \(a\) be its image in~\(X\), and let \(b\) be
its image in~\(Y\). Then the following sequence of group
homomorphisms is \emph{exact}:
\[
  \pi_1(\overline X_y,\overline a)
  \longrightarrow\pi_1(X,a)
  \longrightarrow\pi_1(Y,b)\longrightarrow e.
\]
\end{corollary}

\begin{remarques}
\label{X.1.5}
Observe that the proof of~\Ref{X.1.3} uses~\Ref{X.1.2} in an
essential way, and hence the \emph{first comparison theorem} in
algebraic-formal geometry. By contrast, the descent theory of
Expos\'e~\Ref{IX}
% page: 264
\marginpar{264}
%
entered only through \ifthenelse{\boolean{orig}}{IX~}{}\Ref{IX.3.4}, for which a direct proof is
easy in the case of a \emph{proper} morphism \(f\colon X\to Y\)
such that \(f_*(\cal O_X)=\cal O_Y\). Indeed, let \(Y'\) be \'etale
over~\(Y\), and suppose that \(X'=X\times_YY'\) is the disjoint sum
of two nonempty open subsets. We prove that the same is true of~\(Y'\).
We may write \(Y'=\Spec(\cal A)\), and hence
\(X'=\Spec(\cal B)\), with
\[
  \cal B=\cal A\otimes_{\cal O_Y}\cal O_X.
\]
The decomposition of~\(X'\) as a disjoint sum corresponds to a
decomposition of~\(\cal B\) as a product of two nonzero algebras
\(\cal B_1\) and \(\cal B_2\). Since
\(f_*(\cal O_X)=\cal O_Y\), it follows readily that
\(f_*(\cal B)=\cal A\). Thus \(\cal A\) is a product of the two
algebras \(f_*(\cal B_1)\) and \(f_*(\cal B_2)\), which are likewise
nonzero because their unit sections are nonzero. This proves the
assertion.
\end{remarques}

\end{filecontents*}
\documentclass[11pt,oneside]{amsbook}

\usepackage[T1]{fontenc}
\usepackage[utf8]{inputenc}
\usepackage{lmodern}
\usepackage{microtype}
\usepackage{amsmath,amssymb,mathtools}
\usepackage{ifthen}
\usepackage[all]{xy}
\usepackage[
  a4paper,
  inner=31mm,
  outer=31mm,
  top=28mm,
  bottom=30mm,
  marginparwidth=22mm,
  marginparsep=5mm
]{geometry}
\usepackage[hidelinks,pdfencoding=auto,hypertexnames=false]{hyperref}
\usepackage{bookmark}
\usepackage{xspace}

\newtheorem{proposition}{Proposition}[section]
\newtheorem{theorem}[proposition]{Theorem}
\newtheorem{corollary}[proposition]{Corollary}
\newtheorem{lemma}[proposition]{Lemma}
\newtheorem{lemme}[proposition]{Lemma}
\newtheorem{definition}[proposition]{Definition}
\newtheorem{sectiondefinition}{Definition}[section]
\newtheorem{corollaries}[proposition]{Corollaries}
\newtheorem{corollaire}[proposition]{Corollary}
\newtheorem{theoreme}[proposition]{Theorem}
\newtheorem{example}[proposition]{Example}
\newtheorem{exemples}[proposition]{Examples}
\newtheorem{exemple}[proposition]{Example}
\newenvironment{definitions}
  {\refstepcounter{proposition}\par\medskip\noindent
   \textbf{Definitions \theproposition.}\ }
  {\par\medskip}
\newtheorem*{serrecriterion}{Serre's criterion}
\newtheorem*{theoremstar}{Theorem}
\newtheorem{theoremedepurete}[proposition]{Purity theorem}
\newtheorem{theoremedefinition}[proposition]{Theorem and definition}
\newtheorem{subproposition}[proposition]{Proposition}
\newtheorem{subdefinition}[proposition]{Definition}
\newtheorem{sublemme}[proposition]{Lemma}
\newtheorem{subremarque}[proposition]{Remark}
\theoremstyle{remark}
\newtheorem{remarques}[proposition]{Remarks}
\newtheorem{remarque}[proposition]{Remark}
\newtheorem*{remarquestar}{Remark}
\newtheorem*{remarquesstar}{Remarks}
\newenvironment{remarqueMR}
  {\begin{remarquestar}\lcrochetbf Added in 2003 (MR):\ }
  {\rcrochetbf\end{remarquestar}}
\theoremstyle{plain}
\newenvironment{enumerateb}{\begin{enumerate}}{\end{enumerate}}

% Compatibility for the translated II.4.14 bridge, which retains a few
% harmless source-level notation macros.
\newboolean{orig}
\setboolean{orig}{false}
\renewcommand{\Ref}[1]{\ref{#1}}
\DeclareMathOperator{\Spec}{Spec}
\newcommand{\goth}[1]{\mathfrak{#1}}
\newcommand{\cal}[1]{\mathcal{#1}}
\newcommand{\ZZ}{\mathbb Z}
\newcommand{\QQ}{\mathbb Q}
\newcommand{\CC}{\mathbb C}
\newcommand{\PP}{\mathbb P}
\renewcommand{\AA}{\mathbf A}
\newcommand{\VV}{\mathbf V}
\DeclareMathOperator{\gr}{gr}
\newcommand{\cf}{cf.\ }
\newcommand{\ie}{i.e.\ }
\newcommand{\resp}{\ifthenelse{\boolean{orig}}{respectively}{respectively\xspace}}
\newcommand{\og}{\lq\lq{}}
\newcommand{\fg}{\rq\rq{}\ifthenelse{\boolean{orig}}{}{\xspace}}
\providecommand{\index}[1]{}
\providecommand{\oldindexnot}[1]{}
\DeclareMathAlphabet{\othercal}{U}{eus}{m}{n}
\newcommand{\kres}{\kappa}
\DeclareMathOperator{\Hom}{Hom}
\DeclareMathOperator{\hop}{hop}
\DeclareMathOperator{\Gal}{Gal}
\renewcommand{\H}{\mathrm H}
\DeclareMathOperator{\Aut}{Aut}
\DeclareMathOperator{\Isom}{Isom}
\DeclareMathOperator{\Ker}{Ker}
\DeclareMathOperator{\Ind}{Ind}
\DeclareMathOperator{\Pro}{Pro}
\DeclareMathOperator{\Prodash}{Pro\text{-}}
\DeclareMathOperator{\Proj}{Proj}
\DeclareMathOperator{\prof}{prof}
\DeclareMathOperator{\Norm}{N}
\DeclareMathOperator{\Ext}{Ext}
\DeclareMathOperator{\Gl}{\mathbf{GL}}
\DeclareMathOperator{\GL}{\mathbf{GL}}
\DeclareMathOperator{\Pic}{Pic}
\DeclareMathOperator{\Symp}{\mathbf{Sp}}
\DeclareMathOperator{\Supp}{Supp}
\DeclareMathOperator{\codim}{codim}
\DeclareMathOperator{\car}{char}
\DeclareMathOperator{\card}{card}
\DeclareMathOperator{\divisor}{div}
\DeclareMathOperator{\diag}{diag}
\DeclareMathOperator{\Symm}{Sym}
\newcommand{\res}{\mathrm{res}}
\newcommand{\red}{\mathrm{red}}
\newcommand{\R}{\mathrm{R}}
\newcommand{\et}{\mathrm{\acute{e}t}}
\newcommand{\an}{\mathrm{an}}
\newcommand{\alg}{\mathrm{alg}}
\newcommand{\GG}{\mathbf{G}}
\newcommand{\WW}{\mathbf{W}}
\newcommand{\bbalpha}{\boldsymbol{\alpha}}
\newcommand{\bbmu}{\boldsymbol{\mu}}
\newcommand{\LL}{\mathbb L}
\newcommand{\tame}{\mathrm{t}}
\newcommand{\leftexp}[2]{{\vphantom{#2}}^{#1}{#2}}
\newcommand{\skpt}{}
\let\Subsection\subsection
\let\Subsubsection\subsubsection
\newcommand{\Ens}{\mathbf{Ens}}
\newcommand{\Univ}{\mathfrak U}
\newcommand{\Cat}{\mathbf{Cat}}
\newcommand{\Sch}{\mathbf{Sch}}
\DeclareMathOperator{\CatFl}{\mathbf{Fl}}
\DeclareMathOperator{\Ob}{Ob}
\DeclareMathOperator{\Fl}{Fl}
\DeclareMathOperator{\SheafHom}{\mathbf{Hom}}
\DeclareMathOperator{\SheafAut}{\mathbf{Aut}}
\DeclareMathOperator{\SheafIsom}{\mathbf{Isom}}
\newcommand{\quoi}{}
\newcommand{\cart}{\mathrm{cart}}
\newcommand{\Fib}{\mathrm{Fib}}
\newcommand{\Ouv}{\operatorname{Ouv}}
\newcommand{\Fer}{\operatorname{Fer}}
\newcommand{\varprojLim}{\varprojlim}
\newcommand{\id}{\mathrm{id}}
\newcommand{\from}{\longleftarrow}
\newcommand{\mto}{\longmapsto}
\newcommand{\To}{\Longrightarrow}
\newcommand{\lto}[1]{\xrightarrow{#1}}
\let\dpl\displaystyle
\newcommand{\chapterspace}[1]{}
\newcommand{\isomto}{\xrightarrow{\sim}}
\newcommand{\isomfrom}{\xleftarrow{\sim}}
\newcommand{\Cf}{Cf.\ }
\newcommand{\loccit}{loc.\ cit.\ }
\newcommand{\ptbl}{.\ }
\newcommand{\No}{no.~}
\newcommand{\no}{no.~}
\newcommand{\numero}{no.~}
\newcommand{\sss}{if and only if}
\newcommand{\tH}{{}^t\!H}
\newcommand{\tbigcap}{\mathop{\textstyle\bigcap}\limits}
\newcommand{\bbullet}{{\scriptscriptstyle\bullet}}
\newcommand{\pr}{\mathrm{pr}}
\newcommand{\lcrochetbf}{{\fontfamily{ptm}\bfseries\upshape[}}
\newcommand{\rcrochetbf}{{\fontfamily{ptm}\bfseries\upshape]}}

\hypersetup{
  pdftitle={SGA 1: Étale Coverings and the Fundamental Group — English Translation},
  pdfauthor={English translation from the official French arXiv/SMF TeX},
  pdfsubject={Complete English translation of SGA 1},
  pdfkeywords={SGA 1, étale coverings, fundamental group, English translation}
}

\renewcommand{\thechapter}{\Roman{chapter}}
\renewcommand{\thesection}{\thechapter.\arabic{section}}
\numberwithin{equation}{section}
\setcounter{secnumdepth}{1}
\setcounter{tocdepth}{1}
\makeatletter
\newcommand{\sgaFixedHeading}[2]{#1*{#2}\def\@currentlabel{#2}}
\makeatother

\begin{document}

\frontmatter
\begin{titlepage}
  \centering
  \vspace*{35mm}
  {\Huge\bfseries \textit{Étale Coverings and the Fundamental Group}\par}
  \vspace{3mm}
  {\LARGE (SGA~1)\par}
  \vspace{20mm}
  {\Large English translation\par}
  \vspace{3mm}
  {\Large Complete volume\par}
  \vfill
  {\large Translated from the official French arXiv/SMF \TeX{} source\par}
  \vspace{5mm}
  {\large 22 July 2026\par}
\end{titlepage}

\tableofcontents

\mainmatter
\setcounter{chapter}{0}
\chapter{Étale morphisms}
\label{I}
\marginpar{\raggedright\footnotesize French source p.~1}

To simplify the exposition, we assume that all preschemes under
consideration are locally noetherian, at least after
\hyperref[I.2]{no.~I.2}.

\section{Notions of differential calculus}
\label{I.1}

Let \(X\) be a prescheme over \(Y\), and let \(\Delta_{X/Y}\), or simply
\(\Delta\),\label{indnot:ab} be the diagonal morphism
\(X\to X\times_Y X\).  It is an
immersion, hence a closed immersion of \(X\) into an open subset \(V\) of
\(X\times_Y X\).  Let \(\mathcal I_X\) be the ideal of the closed
subprescheme corresponding to the diagonal in \(V\) (N.B. if one wishes to
proceed intrinsically, without assuming that \(X\) is separated over
\(Y\)---a hypothesis that would be farcical---one should consider the
set-theoretic inverse image of \(\mathcal O_{X\times X}\) in \(X\), and
denote by \(\mathcal I_X\) the augmentation ideal in the latter\ldots).
The sheaf \(\mathcal I_X/\mathcal I_X^2\) may be regarded as a
quasi-coherent sheaf on \(X\); we denote it by \(\Omega^1_{X/Y}\).  It is
\label{indnot:ac}
of finite type if \(X\to Y\) is of finite type.  It behaves well under base
change \(Y'\to Y\).  We also introduce the sheaves
\(\mathcal O_{X\times_Y X}/\mathcal I_X^{n+1}=\mathcal P^n_{X/Y}\);\label{indnot:ad} these
are sheaves of \emph{rings} on \(X\), making \(X\) into a prescheme that
may be denoted by \(\Delta^n_{X/Y}\)\label{indnot:ae} and called the \emph{\(n\)-th
infinitesimal neighborhood of \(X/Y\)}.  The resulting sorites is entirely
trivial, although rather long\footnote{Cf.\ EGA IV~16.3.}; it would be
prudent to speak of it only when there is something useful to say about it,
in connection with smooth morphisms.

\section{Quasi-finite morphisms}
\label{I.2}

\begin{proposition}
\label{I.2.1}
Let \(A\to B\) be a local homomorphism (N.B. the rings are now
noetherian), and let \(\mathfrak m\) be the maximal ideal of~\(A\).  The
following conditions are equivalent:
\begin{enumerate}
\item[(i)] \(B/\mathfrak mB\) is finite-dimensional over
  \(k=A/\mathfrak m\).
\item[(ii)] \(\mathfrak mB\) is an ideal of definition, and
  \(B/\mathfrak r(B)=\kappa(B)\) is an extension of
  \(k=\kappa(A)\).
\item[(iii)] The completion \(\widehat B\) is finite over the completion
  \(\widehat A\) of~\(A\).
\end{enumerate}
\end{proposition}

One then says\marginpar{\raggedright\footnotesize French source p.~2} that
\(B\) \emph{is quasi-finite} over~\(A\).  A morphism
\(f\colon X\to Y\) is said to be quasi-finite at~\(x\) (or the
\(Y\)-prescheme \(f\) is said to be quasi-finite at~\(x\)) if
\(\mathcal O_x\) is quasi-finite over \(\mathcal O_{f(x)}\).  This is also
equivalent to saying that \(x\) is \emph{isolated in its fiber}
\(f^{-1}(f(x))\).\footnote{The original printing
(\hyperref[I]{Exposé~I}, p.~2) and both
French arXiv renderings print \(f^{-1}(x)\).  Since \(x\in X\) for
\(f\colon X\to Y\), the type-correct fiber through \(x\) is
\(f^{-1}(f(x))\), which is used here with the source defect explicitly
disclosed.} A morphism is said to be quasi-finite if it is so at
every point.\footnote{In EGA~II~6.2.3 one assumes in addition that \(f\) is
of finite type.}

\begin{corollary}
\label{I.2.2}
If \(A\) is complete, quasi-finite is equivalent to finite.
\end{corollary}

One could give the usual chain of implications \textup{(i), (ii), (iii),
(iv), (v)} for quasi-finite morphisms, but this does not seem indispensable
here.

% SGA 1, Expose I, section I.3.
% French authority lines 653--799; source-audited bounded unit.

\section{Unramified or net morphisms}
\label{I.3}

\begin{proposition}
\label{I.3.1}
Let \(f\colon X\to Y\) be a morphism of finite type, let \(x\in X\), and
put \(y=f(x)\).  The following conditions are equivalent:
\begin{enumerate}
\item[(i)] 
  \(\mathcal O_x/\mathfrak m_y\mathcal O_x\) is a finite separable
  extension of~\(\kappa(y)\).
\item[(ii)] \(\Omega^1_{X/Y}\) vanishes at~\(x\).
\item[(iii)] The diagonal morphism \(\Delta_{X/Y}\) is an open immersion
  in a neighborhood of~\(x\).
\end{enumerate}
\end{proposition}

For the implication \textup{(i)}\(\Rightarrow\)\textup{(ii)}, Nakayama's
lemma immediately reduces us to the case
\(Y=\operatorname{Spec}(k)\), \(X=\operatorname{Spec}(k')\), where the
assertion is well known and, moreover, immediate from the definition of
separability.  The implication \textup{(ii)}\(\Rightarrow\)\textup{(iii)}
follows from a pleasant and easy characterization of open immersions using
Krull.  For \textup{(iii)}\(\Rightarrow\)\textup{(i)}, one is
again reduced to the case \(Y=\operatorname{Spec}(k)\), with the diagonal
morphism an open immersion everywhere.  One must then prove that \(X\) is
finite, with a separable coordinate ring over~\(k\); for this one reduces to
the case where \(k\) is algebraically closed.  But then every closed point of
\(X\) is isolated (it is precisely the inverse image of the diagonal under
the morphism \(X\to X\times_k X\) defined by~\(x\)), and hence \(X\) is
finite.  One may therefore suppose that \(X\) consists of a single point,
with ring~\(A\).  Thus \(A\otimes_kA\to A\) is an isomorphism, whence
\(A=k\), as required.

\begin{definition}
\label{I.3.2}
\begin{enumerate}
\item[(a)] One then says that \(f\) is \emph{net}, or equivalently
  \emph{unramified}, at~\(x\), or that \(X\) is net, or unramified, at
  \(x\) over~\(Y\).
\item[(b)] Let \(A\to B\) be a local homomorphism.  It is said to be
  \emph{net}, or \emph{unramified}, and \(B\) is said to be a net, or
  unramified, local algebra over~\(A\), if
  \(B/\mathfrak r(A)B\) is a finite separable extension of
  \(A/\mathfrak r(A)\); that is, if
  \(\mathfrak r(A)B=\mathfrak r(B)\) and \(\kappa(B)\) is a separable
  extension of~\(\kappa(A)\).\footnote{Cf.\ the regrets in
  \hyperref[III]{Expos\'e~III}, \hyperref[III.1.2]{no.~1.2}.}
\end{enumerate}
\end{definition}

\paragraph{Remarks.}
The fact\marginpar{\raggedright\footnotesize French source p.~3} that
\(B\) is unramified over~\(A\) can already be checked on the completions of
\(A\) and~\(B\).  Unramified implies quasi-finite.

\begin{corollary}
\label{I.3.3}
The set of points at which \(f\) is unramified is open.
\end{corollary}

\begin{corollary}
\label{I.3.4}
Let \(X'\) and \(X\) be two preschemes of finite type over~\(Y\), and let
\(g\colon X'\to X\) be a \(Y\)-morphism.  If \(X\) is unramified over
\(Y\), the graph morphism

\[
  \Gamma_g\colon X'\longrightarrow X'\times_Y X
\]

is an open immersion.\footnote{The original printing
(\hyperref[I]{Expos\'e~I}, p.~3) and
the French TeX give \(X\times_Y X\) as the graph target and later write
\(\operatorname{id}_{X'}\) in the pullback map.  The typed fiber products
force \(X'\times_Y X\) and \(\operatorname{id}_X\), respectively; those
type-correct readings are used here with both source defects disclosed.}
\end{corollary}

Indeed, it is the inverse image of the diagonal morphism
\(X\to X\times_Y X\) under
\[
g\times_Y\operatorname{id}_{X}\colon
X'\times_Y X\longrightarrow X\times_Y X.
\]

One may also introduce the annihilator ideal
\(\mathfrak d_{X/Y}\)\label{indnot:af} of \(\Omega^1_{X/Y}\), called the
\emph{different ideal} of~\(X/Y\).  It defines a closed subprescheme of
\(X\) whose underlying set is the set of points at which \(X/Y\) is
ramified, that is, not unramified.

\begin{proposition}
\label{I.3.5}
\begin{enumerate}
\item[(i)] An immersion is unramified.
\item[(ii)] The composite of two unramified morphisms is unramified.
\item[(iii)] A base change of an unramified morphism is unramified.
\end{enumerate}
\end{proposition}

This may be seen equally well from criterion \textup{(ii)} or criterion
\textup{(iii)} (the latter seems more amusing to me).  One can of course
also be more precise by giving pointwise statements; this is more general
only in appearance (except in the setting of Definition~\ref{I.3.2}(b)), and
tedious.  As usual, one obtains the following corollaries.

\begin{corollaries}
\label{I.3.6}
\begin{enumerate}
\item[(iv)] The Cartesian product of two unramified morphisms is unramified.
\item[(v)] If \(gf\) is unramified, then \(f\) is unramified.
\item[(vi)] If \(f\) is unramified, then \(f_{\mathrm{red}}\) is
  unramified.
\end{enumerate}
\end{corollaries}

\begin{proposition}
\label{I.3.7}
Let \(A\to B\) be a local homomorphism.  Suppose that the residue-field
extension \(\kappa(B)/\kappa(A)\) is trivial, or that \(\kappa(A)\) is
algebraically closed.  Then \(B/A\) is unramified if and only if
\(\widehat B\), as an \(\widehat A\)-algebra, is a quotient of
\(\widehat A\).
\end{proposition}

\paragraph{Remarks.}
\begin{itemize}
\item When the residue-field extension is not assumed to be trivial, one may
  reduce to that case by making a suitable finite flat extension over~\(A\)
  that kills the extension in question.
\item Take the example in which \(A\) is the local ring of an ordinary
  double point of a curve and \(B\) is the local ring at a point of the
  normalization.  Then \(A\subset B\), and \(B\) is unramified
  over\marginpar{\raggedright\footnotesize French source p.~4} \(A\) with
  trivial residue-field extension; the map
  \(\widehat A\to\widehat B\) is surjective but \emph{not injective}.
  We shall therefore strengthen the notion of unramifiedness.
\end{itemize}

\section{Étale morphisms. Étale coverings}
\label{I.4}

We shall take for granted whatever is needed concerning flat morphisms; these
facts will be proved later, if necessary\footnote{Cf.\
\hyperref[IV]{Expos\'e~IV}.}.

\begin{definition}
\label{I.4.1}
\begin{enumerate}
\item[(a)] Let \(f\colon X\to Y\) be a morphism of finite type.  One says
  that \(f\) is \emph{étale} at~\(x\) if \(f\) is flat at~\(x\) and net
  (equivalently, unramified) at~\(x\).  One says that \(f\) is étale if it
  is so at every point.  One says that \(X\) is étale at~\(x\) over~\(Y\),
  or that it is a \(Y\)-prescheme étale at~\(x\), and so on.
\item[(b)] Let \(f\colon A\to B\) be a local homomorphism.  One says that
  \(f\) is étale, or that \(B\) is étale over~\(A\), if \(B\) is flat and
  unramified over~\(A\).\footnote{Cf.\ the regrets in
  \hyperref[III]{Expos\'e~III},
  \hyperref[III.1.2]{no.~1.2}.}
\end{enumerate}
\end{definition}

\begin{proposition}
\label{I.4.2}
For \(B/A\) to be étale, it is necessary and sufficient that
\(\widehat B/\widehat A\) be étale.
\end{proposition}

Indeed, this holds separately for ``unramified'' and for ``flat.''

\begin{corollary}
\label{I.4.3}
Let \(f\colon X\to Y\) be of finite type, and let \(x\in X\).  Whether
\(f\) is étale at~\(x\) depends only on the local homomorphism
\(\mathcal O_{f(x)}\to\mathcal O_x\), and even only on the corresponding
homomorphism between the completions.
\end{corollary}

\begin{corollary}
\label{I.4.4}
Suppose that the residue-field extension
\(\kappa(A)\to\kappa(B)\) is trivial, or that \(\kappa(A)\) is
algebraically closed.  Then \(B/A\) is étale if and only if
\(\widehat A\to\widehat B\) is an isomorphism.
\end{corollary}

This follows by combining flatness with Proposition~\ref{I.3.7}.

\begin{proposition}
\label{I.4.5}
Let \(f\colon X\to Y\) be a morphism of finite type.  Then the set of
points at which it is étale is open.
\end{proposition}

Indeed, this holds separately for ``unramified'' and for ``flat.''

This proposition shows that one may dispense with ``pointwise'' statements
when studying morphisms of finite type that are étale at some point.

\begin{proposition}
\label{I.4.6}
\begin{enumerate}
\item[(i)] An open immersion is étale.
\item[(ii)] The composite of two étale morphisms is étale.
\item[(iii)] A base change of an étale morphism is étale.
\end{enumerate}
\end{proposition}

\marginpar{\raggedright\footnotesize French source p.~5}%
Indeed, \textup{(i)} is trivial, and for \textup{(ii)} and \textup{(iii)}
it is enough to note that the assertions hold for ``unramified'' and for
``flat.''  To tell the truth, there are also corresponding statements for
local homomorphisms (without a finiteness condition); in any event, they will
have to appear in the \emph{multiplodoque}, beginning with the unramified
case.

\begin{corollary}
\label{I.4.7}
The Cartesian product of two étale morphisms is likewise étale.
\end{corollary}

\begin{corollary}
\label{I.4.8}
Let \(X\) and \(X'\) be of finite type over~\(Y\), and let
\(g\colon X\to X'\) be a \(Y\)-morphism.  If \(X'\) is unramified over
\(Y\) and \(X\) is étale over~\(Y\), then \(g\) is étale.
\end{corollary}

Indeed, \(g\) is the composite of the graph morphism
\[
  \Gamma_g\colon X\longrightarrow X\times_Y X',
\]
which is an open immersion by Corollary~\ref{I.3.4}, and the projection
\(X\times_Y X'\to X'\), which is étale because it is obtained from the
étale morphism \(X\to Y\) by the base change \(X'\to Y\).

\begin{definition}
\label{I.4.9}
An \emph{étale covering} (respectively, a \emph{net covering}, also called
an unramified covering) of~\(Y\) is a \(Y\)-scheme \(X\) that is finite
over~\(Y\) and étale (respectively, unramified) over~\(Y\).
\end{definition}

The first condition means that \(X\) is defined by a coherent sheaf of
algebras \(\mathcal B\) on~\(Y\).  The second then means that \(\mathcal B\)
is locally free over~\(Y\) (with no corresponding local-freeness requirement
in the net case), and, moreover, that for every \(y\in Y\) the fiber
\[
  \mathcal B(y)=\mathcal B_y\otimes_{\mathcal O_y}\kappa(y)
\]
is a separable algebra over~\(\kappa(y)\), that is, a finite product of
finite separable field extensions of~\(\kappa(y)\).

\begin{proposition}
\label{I.4.10}
Let \(X\) be a flat covering of~\(Y\) of degree~\(n\) (the definition of
this term deserved to appear in Definition~\ref{I.4.9}), defined by a
coherent locally free sheaf of algebras \(\mathcal B\).  In the usual way one
defines the trace homomorphism \(\mathcal B\to\mathcal A\), which is a
homomorphism of \(\mathcal A\)-modules, where
\(\mathcal A=\mathcal O_Y\).  For \(X\) to be étale, it is necessary and
sufficient that the corresponding trace pairing
\((x,y)\longmapsto\operatorname{Tr}_{\mathcal B/\mathcal A}(xy)\) induce an
isomorphism \(\mathcal B\xrightarrow{\sim}\mathcal B^\vee\),\footnote{The diplomatic French and the
original printing instead give \(\mathcal B\) as the target.  An
\(\mathcal A\)-bilinear form induces a map
\(\mathcal B\to\mathcal B^\vee\); the dual is therefore the type-correct
target, used here with the source-original defect disclosed.} or,
equivalently,
that the \emph{discriminant section}
\[
  d_{X/Y}=d_{\mathcal B/\mathcal A}
  \in
  \Gamma\!\left(
    Y,
    \bigwedge\nolimits^n\mathcal B^\vee
    \otimes_{\mathcal A}
    \bigwedge\nolimits^n\mathcal B^\vee
  \right)
\]
be invertible, or, finally, that the discriminant ideal defined by this
section be the unit ideal.
\end{proposition}

Indeed, one reduces to the case \(Y=\operatorname{Spec}(k)\); the assertion
is then the well-known separability criterion, and becomes immediate after
passing to an algebraic closure of~\(k\).

\marginpar{\raggedright\footnotesize French source p.~6}%
\paragraph{Remark.}
There will be a less trivial statement below, when one does not assume a
priori that \(X\) is flat over~\(Y\), but instead imposes a normality
hypothesis.

\input{sga1-cumulative-input-049.tex}

\input{sga1-cumulative-input-050.tex}

\input{sga1-cumulative-input-051.tex}

\input{sga1-cumulative-input-052.tex}

\input{sga1-cumulative-input-059.tex}

\input{sga1-cumulative-input-060.tex}

\input{sga1-cumulative-input-061.tex}

\input{sga1-cumulative-input-062.tex}

\input{sga1-cumulative-input-063.tex}

\input{sga1-cumulative-input-064.tex}

\input{sga1-cumulative-input-065.tex}

\input{sga1-cumulative-input-066.tex}

\input{sga1-cumulative-input-053.tex}

\input{sga1-cumulative-input-054.tex}

\input{sga1-cumulative-input-055.tex}

\input{sga1-cumulative-input-056.tex}

\input{sga1-cumulative-input-057.tex}

\input{sga1-cumulative-input-058.tex}

\input{sga1-cumulative-input-036.tex}

\input{sga1-cumulative-input-037.tex}

\input{sga1-cumulative-input-038.tex}

\input{sga1-cumulative-input-039.tex}

\input{sga1-cumulative-input-040.tex}

\input{sga1-cumulative-input-041.tex}

\input{sga1-cumulative-input-043.tex}

\input{sga1-cumulative-input-042.tex}

\input{sga1-cumulative-input-044.tex}

\input{sga1-cumulative-input-045.tex}

\input{sga1-cumulative-input-046.tex}

\input{sga1-cumulative-input-047.tex}

\input{sga1-cumulative-input-048.tex}

\input{sga1-cumulative-input-012.tex}

\input{sga1-cumulative-input-013.tex}

\input{sga1-cumulative-input-014.tex}

\input{sga1-cumulative-input-018.tex}

\input{sga1-cumulative-input-016.tex}

\input{sga1-cumulative-input-017.tex}

\input{sga1-cumulative-input-015.tex}

\input{sga1-cumulative-input-021.tex}

\input{sga1-cumulative-input-019.tex}

\input{sga1-cumulative-input-020.tex}

\input{sga1-cumulative-input-022.tex}

\input{sga1-cumulative-input-002.tex}

\input{sga1-cumulative-input-001.tex}

\input{sga1-cumulative-input-003.tex}

\input{sga1-cumulative-input-005.tex}

\input{sga1-cumulative-input-006.tex}

\input{sga1-cumulative-input-008.tex}

\input{sga1-cumulative-input-007.tex}

\input{sga1-cumulative-input-009.tex}

\input{sga1-cumulative-input-010.tex}

\input{sga1-cumulative-input-011.tex}

\input{sga1-cumulative-input-028.tex}

\input{sga1-cumulative-input-023.tex}

\input{sga1-cumulative-input-024.tex}

\input{sga1-cumulative-input-025.tex}

\input{sga1-cumulative-input-026.tex}

\input{sga1-cumulative-input-027.tex}

\input{sga1-cumulative-input-096.tex}

\input{sga1-cumulative-input-089.tex}

\input{sga1-cumulative-input-090.tex}

\input{sga1-cumulative-input-091.tex}

\input{sga1-cumulative-input-092.tex}

\input{sga1-cumulative-input-093.tex}

\input{sga1-cumulative-input-094.tex}

\input{sga1-cumulative-input-095.tex}

\input{sga1-cumulative-input-088.tex}

\input{sga1-cumulative-input-079.tex}

\input{sga1-cumulative-input-080.tex}

\input{sga1-cumulative-input-081.tex}

\input{sga1-cumulative-input-082.tex}

\input{sga1-cumulative-input-083.tex}

\input{sga1-cumulative-input-084.tex}

\input{sga1-cumulative-input-085.tex}

\input{sga1-cumulative-input-086.tex}

\input{sga1-cumulative-input-087.tex}

\input{sga1-cumulative-input-076.tex}

\input{sga1-cumulative-input-077.tex}

\input{sga1-cumulative-input-078.tex}

\input{sga1-cumulative-input-067.tex}

\input{sga1-cumulative-input-068.tex}

\input{sga1-cumulative-input-069.tex}

\input{sga1-cumulative-input-070.tex}

\input{sga1-cumulative-input-071.tex}

\input{sga1-cumulative-input-072.tex}

\input{sga1-cumulative-input-073.tex}

\input{sga1-cumulative-input-074.tex}

\input{sga1-cumulative-input-075.tex}

\input{sga1-cumulative-input-035.tex}

\input{sga1-cumulative-input-029.tex}

\input{sga1-cumulative-input-030.tex}

\input{sga1-cumulative-input-031.tex}

\input{sga1-cumulative-input-032.tex}

\input{sga1-cumulative-input-033.tex}

\input{sga1-cumulative-input-034.tex}

\input{sga1-cumulative-input-138.tex}

\input{sga1-cumulative-input-134.tex}

\input{sga1-cumulative-input-135.tex}

\input{sga1-cumulative-input-136.tex}

\input{sga1-cumulative-input-137.tex}

\input{sga1-cumulative-input-133.tex}

\input{sga1-cumulative-input-126.tex}

\input{sga1-cumulative-input-127.tex}

\input{sga1-cumulative-input-128.tex}

\input{sga1-cumulative-input-129.tex}

\input{sga1-cumulative-input-130.tex}

\input{sga1-cumulative-input-131.tex}

\input{sga1-cumulative-input-132.tex}

\input{sga1-cumulative-input-125.tex}

\input{sga1-cumulative-input-120.tex}

\input{sga1-cumulative-input-121.tex}

\input{sga1-cumulative-input-122.tex}

\input{sga1-cumulative-input-123.tex}

\input{sga1-cumulative-input-124.tex}

\input{sga1-cumulative-input-119.tex}

\input{sga1-cumulative-input-102.tex}

\input{sga1-cumulative-input-100.tex}

\input{sga1-cumulative-input-101.tex}

\input{sga1-cumulative-input-097.tex}

\input{sga1-cumulative-input-098.tex}

\input{sga1-cumulative-input-099.tex}

\input{sga1-cumulative-input-103.tex}

\input{sga1-cumulative-input-104.tex}

\input{sga1-cumulative-input-105.tex}

\input{sga1-cumulative-input-107.tex}

\input{sga1-cumulative-input-108.tex}

\input{sga1-cumulative-input-106.tex}

\input{sga1-cumulative-input-112.tex}

\input{sga1-cumulative-input-109.tex}

\input{sga1-cumulative-input-110.tex}

\input{sga1-cumulative-input-111.tex}

\input{sga1-cumulative-input-115.tex}

\input{sga1-cumulative-input-113.tex}

\input{sga1-cumulative-input-114.tex}

\input{sga1-cumulative-input-117.tex}

\input{sga1-cumulative-input-118.tex}

\input{sga1-cumulative-input-116.tex}

\end{document}

% SGA-CUMULATIVE-END 01_SGA1_EN_COMPLETE_PUBLISHED_SCOPE_READER.tex

% SGA-CUMULATIVE-BEGIN 02_SGA2_EN_COMPLETE_READER.tex
\documentclass[11pt,a4paper,oneside]{book}

\usepackage[T1]{fontenc}
\usepackage[utf8]{inputenc}
\usepackage{lmodern}
\usepackage{microtype}
\usepackage{amsmath,amssymb,amsthm,amscd,mathtools,mathrsfs}
\usepackage{xspace}
\usepackage{xcolor}
\usepackage{graphicx}
\usepackage{array,booktabs,longtable}
\usepackage{enumitem}
\usepackage{tocloft}
\usepackage{tikz-cd}
\usepackage[all]{xy}
\usepackage[a4paper,margin=27mm,marginparwidth=18mm,marginparsep=4mm]{geometry}
\usepackage[hidelinks,pdfencoding=auto,hypertexnames=false]{hyperref}

% Modern compatibility layer for the translated bodies.  The legacy
% sga2-smf.sty is intentionally not loaded: it is incompatible with the
% current LaTeX kernel.  The local shim retains only the macros actually
% required by the source-aligned English reader.
% BEGIN INCLUDED active_source/sga2_reader_macros.tex
% Modern compatibility layer for the source-aligned English SGA 2 bodies.
% This deliberately replaces, rather than loads, the legacy sga2-smf.sty.

\makeatletter
\@ifundefined{c@pageoriginal}{\newcounter{pageoriginal}}{}
\@ifundefined{c@toto}{\newcounter{toto}}{}
\@ifundefined{c@yequation}{\newcounter{yequation}}{}
\@ifundefined{c@sourcecorollary}{\newcounter{sourcecorollary}[section]}{}
\@ifundefined{c@sourcecondition}{\newcounter{sourcecondition}}{}
\@ifundefined{c@footnoteAG}{\newcounter{footnoteAG}}{}
\@ifundefined{c@footnoteNDE}{\newcounter{footnoteNDE}}{}
\@ifundefined{c@aconjecture}{\newcounter{aconjecture}[section]}{}
\makeatother
\renewcommand{\theaconjecture}{\Alph{aconjecture}}

% The corrected/annotated branch is always selected.
\providecommand{\sisi}[2]{#2}
\providecommand{\pageoriginale}{%
  \refstepcounter{pageoriginal}%
  \marginpar{\raggedright\footnotesize\bfseries\thepageoriginal}%
  \xspace}
\providecommand{\srcpage}[1]{%
  \marginpar{\raggedright\scriptsize\color{gray}Original French p.~#1}}
\providecommand{\srcpdfpage}[1]{%
  \marginpar{\raggedright\scriptsize\color{gray}French source PDF p.~#1}}
\providecommand{\chapterspace}[1]{}

% Original and editorial notes.
\providecommand{\nde}[1]{\footnote{\textit{Editor's note.} #1}}
\providecommand{\de}[1]{\footnote{\textit{Editor's note.} #1}}
\providecommand{\ndemark}{\footnotemark}
\providecommand{\ndetext}[1]{\footnotetext{\textit{Editor's note.} #1}}
\providecommand{\sfootnote}[1]{\footnote{#1}}
\providecommand{\sfootnotemark}{\footnotemark}
\providecommand{\sfootnotetext}[1]{\footnotetext{#1}}
\providecommand{\markedfootnote}[2]{%
  \begingroup\renewcommand{\thefootnote}{#1}\footnote{#2}\endgroup}

% Common prose abbreviations and quotations retained from the source macros.
\providecommand{\og}{``}
\providecommand{\fg}{''\xspace}
\xspaceaddexceptions{\,\markedfootnote}
\providecommand{\cf}{cf.\kern.3em}
\providecommand{\Cf}{Cf.\kern.3em}
\providecommand{\ie}{i.e.\kern.3em}
\providecommand{\resp}{resp.\kern.3em}
\providecommand{\loccit}{\textit{loc. cit.}\xspace}
\providecommand{\EGA}{\textup{EGA}\kern.2em}
\providecommand{\SGA}{\textup{SGA}\kern.3em}
\providecommand{\Exp}{Exp.\kern.3em}
\providecommand{\ptbl}{.\kern.2em}
\providecommand{\numero}{no.\kern.3em}
\providecommand{\no}{no.\kern.3em}
\providecommand{\No}{no.\kern.3em}
\providecommand{\afortiori}{\textit{a fortiori}\xspace}
\providecommand{\quoi}{}
\providecommand{\refname}{Bibliography}

% Calligraphic, bold, and ideal alphabets.
\providecommand{\Aa}{\mathcal A}
\providecommand{\Bb}{\mathcal B}
\providecommand{\Cc}{\mathcal C}
\providecommand{\Dd}{\mathcal D}
\providecommand{\Ee}{\mathcal E}
\providecommand{\Ff}{\mathcal F}
\providecommand{\Gg}{\mathcal G}
\providecommand{\Hh}{\mathcal H}
\providecommand{\Ii}{\mathcal I}
\providecommand{\Jj}{\mathcal J}
\providecommand{\Kk}{\mathcal K}
\providecommand{\Ll}{\mathcal L}
\providecommand{\Mm}{\mathcal M}
\providecommand{\Nn}{\mathcal N}
\providecommand{\Oo}{\mathcal O}
\providecommand{\Pp}{\mathcal P}
\providecommand{\Qq}{\mathcal Q}
\providecommand{\Rr}{\mathcal R}
\providecommand{\Ss}{\mathcal S}
\providecommand{\Tt}{\mathcal T}
\providecommand{\Uu}{\mathcal U}
\providecommand{\Vv}{\mathcal V}
\providecommand{\Ww}{\mathcal W}
\providecommand{\Xx}{\mathcal X}
\providecommand{\Yy}{\mathcal Y}
\providecommand{\Zz}{\mathcal Z}
\providecommand{\CJ}{\underline{C}_J}
\providecommand{\DDJ}{\underline{D}_J}
\providecommand{\bbmu}{\boldsymbol{\mu}}
\providecommand{\CC}{\mathbf C}
\providecommand{\GG}{\mathbf G}
\providecommand{\KK}{\mathbf K}
\providecommand{\LL}{\mathbf L}
\providecommand{\NN}{\mathbf N}
\providecommand{\PP}{\mathbf P}
\providecommand{\QQ}{\mathbf Q}
\providecommand{\ZZ}{\mathbf Z}
\providecommand{\Ab}{\mathbf{Ab}}
\providecommand{\ideal}[1]{\mathfrak{#1}}
\providecommand{\aaa}{\ideal a}
\providecommand{\pp}{\ideal p}
\providecommand{\qq}{\ideal q}
\providecommand{\rr}{\ideal r}
\providecommand{\mm}{\ideal m}
\providecommand{\m}{\ideal m}
\providecommand{\mideal}{\mathfrak m}
\providecommand{\formelF}{\mathfrak F}
\providecommand{\formelG}{\mathfrak G}

% Sheaves, standard objects, and local aliases used by the English units.
\let\oldS\S
\providecommand{\sheaf}[1]{\underline{#1}}
\providecommand{\ffam}{\mathbf f}
\providecommand{\gfam}{\mathbf g}
\providecommand{\hfam}{\mathbf h}
\providecommand{\RR}{\mathrm R}
\providecommand{\uC}{\underline{\mathrm C}}
\providecommand{\Ccat}{\mathcal C}
\providecommand{\Cbar}{\overline{\mathcal C}}
\providecommand{\Ctors}{\mathcal C_{\mideal}}
\providecommand{\CatC}{\mathcal C}
\providecommand{\catSch}{\mathbf{Sch}}
\providecommand{\uGamma}{\mathop{\underline{\Gamma}}\nolimits}
\providecommand{\uHom}{\mathop{\underline{\mathrm{Hom}}}\nolimits}
\providecommand{\uExt}{\mathop{\underline{\mathrm{Ext}}}\nolimits}
\providecommand{\uH}{\mathop{\underline{\mathrm H}}\nolimits}
\providecommand{\UH}{\mathop{\underline{\mathrm H}}\nolimits}
\providecommand{\SheafGamma}{\underline{\Gamma}}
\providecommand{\SheafH}{\mathcal H}
\providecommand{\SheafHom}{\mathcal H\!om}
\providecommand{\SheafExt}{\mathcal E\!xt}
\providecommand{\SheafI}{\mathcal I}
\providecommand{\Sheafphi}{\underline{\varphi}}
\providecommand{\SheafPic}{\mathbf{Pic}}
\providecommand{\cH}{\mathcal H}
\providecommand{\cExt}{\mathcal E\!xt}
\providecommand{\cHom}{\mathcal H\!om}
\providecommand{\sH}{\mathcal H}
\providecommand{\sExt}{\mathcal E\!xt}
\providecommand{\cF}{\mathcal F}
\providecommand{\calF}{\mathcal F}
\providecommand{\Lcal}{\mathcal L}
\DeclareMathOperator{\Proj}{Proj}
\DeclareMathOperator{\SheafProj}{\mathbf{Proj}}
\DeclareMathOperator{\gr}{gr}
\DeclareMathOperator{\E}{E}
\DeclareMathOperator{\degtr}{deg\,tr}
\DeclareMathOperator{\Br}{Br}
\DeclareMathOperator{\cd}{cd}
\DeclareMathOperator{\D}{D}
\DeclareMathOperator{\geom}{g\acute{e}om}
\DeclareMathOperator{\hop}{hop}
\DeclareMathOperator{\Ker}{Ker}
\DeclareMathOperator{\tr}{tr}
\DeclareMathOperator{\profet}{prof\;\acute{e}t}
\DeclareMathOperator{\profgeom}{prof\;g\acute{e}om}
\DeclareMathOperator{\rg}{rg}
\providecommand{\h}{\mathop{\underline{H}}\nolimits}
\providecommand{\cG}{\mathcal G}
\providecommand{\cO}{\mathcal O}
\providecommand{\uF}{\underline{\mathrm F}}
\providecommand{\uT}{\underline{\mathrm T}}
\providecommand{\OX}{\mathcal O_X}
\providecommand{\OXx}{\mathcal O_{X,x}}
\providecommand{\OXY}{\mathcal O_{X,Y}}
\providecommand{\OXZ}[1]{\mathcal O_{X,#1}}
\providecommand{\OoXZ}{\mathcal O_{X,Z}}
\providecommand{\OY}[1]{\mathcal O_{Y_{#1}}}
\providecommand{\omx}{\Omega_{X/k}}
\providecommand{\Et}{\mathbf{Et}}
\providecommand{\Etc}{\mathbf{Et}}
\providecommand{\LLcat}{\mathbf L}
\providecommand{\Ib}{\mathcal I'}
\providecommand{\Jb}{\mathcal J}
\providecommand{\restr}[2]{#1\vert_{#2}}
\providecommand{\rest}[1]{_{|#1}}

% Operators.
\providecommand{\Ann}{\operatorname{Ann}}
\providecommand{\Annop}{\operatorname{Ann}}
\providecommand{\Ass}{\operatorname{Ass}}
\providecommand{\Assop}{\operatorname{Ass}}
\providecommand{\Supp}{\operatorname{Supp}}
\providecommand{\Suppop}{\operatorname{Supp}}
\providecommand{\Spec}{\operatorname{Spec}}
\providecommand{\Specop}{\operatorname{Spec}}
\providecommand{\Frac}{\operatorname{Frac}}
\providecommand{\Homop}{\operatorname{Hom}}
\providecommand{\Hom}{\operatorname{Hom}}
\providecommand{\Ob}{\operatorname{Ob}}
\providecommand{\Ext}{\operatorname{Ext}}
\providecommand{\Extop}{\operatorname{Ext}}
\providecommand{\Pic}{\operatorname{Pic}}
\providecommand{\Cl}{\operatorname{Cl}}
\providecommand{\prof}{\operatorname{prof}}
\providecommand{\codh}{\operatorname{codh}}
\providecommand{\codim}{\operatorname{codim}}
\providecommand{\codimn}{\operatorname{codim}}
\providecommand{\codimop}{\operatorname{codim}}
\providecommand{\depthop}{\operatorname{depth}}
\providecommand{\lengthop}{\operatorname{length}}
\providecommand{\pdop}{\operatorname{pd}}
\providecommand{\Hloc}{\operatorname H}
\providecommand{\HH}{\mathrm H}
\renewcommand{\H}{\mathrm H}
\providecommand{\R}{\mathrm R}
\providecommand{\DD}{\mathrm D}
\providecommand{\DA}{\mathrm D A}
\providecommand{\DhatA}{\mathrm D\widehat A}
\providecommand{\CA}{\mathrm C A}
\providecommand{\ChatA}{\mathrm C\widehat A}
\providecommand{\ellA}{\operatorname{length}_A}
\providecommand{\mA}{\mathfrak m_A}
\providecommand{\mB}{\mathfrak m_B}
\providecommand{\Lef}{\operatorname{Lef}}
\providecommand{\Leff}{\operatorname{Leff}}
\providecommand{\Lefc}{\operatorname{Lef}}
\providecommand{\Leffc}{\operatorname{Leff}}

% Arrows, equation controls, and display conveniences.
\providecommand{\from}{\longleftarrow}
\providecommand{\mto}{\longmapsto}
\providecommand{\To}{\Longrightarrow}
\providecommand{\lto}[1]{\xrightarrow{#1}}
\providecommand{\isomto}{\overset{\sim}{\longrightarrow}}
\providecommand{\isomfrom}{\overset{\sim}{\longleftarrow}}
\providecommand{\SSI}{\unskip~$\Longleftrightarrow$~\ignorespaces}
\providecommand{\ALORS}{\unskip~$\Longrightarrow$~\ignorespaces}
\providecommand{\tbigoplus}{\mathop{\textstyle\bigoplus}\limits}
\providecommand{\tbigotimes}{\mathop{\textstyle\bigotimes}\limits}
\providecommand{\boule}{\raisebox{0.2ex}{$\scriptscriptstyle\bullet$}}
\providecommand{\steco}{%
  \setcounter{equation}{\value{yequation}}\refstepcounter{yequation}}
\providecommand{\skipqed}{\renewcommand{\qedsymbol}{}}
\providecommand{\sourceoddity}[2]{%
  \par\smallskip\noindent
  \fcolorbox{gray}{gray!6}{\parbox{0.94\linewidth}{\footnotesize
  \textbf{Source note [#1].} #2}}\par\smallskip}
\providecommand{\danger}{\marginpar{\Large$\triangle$}}
\let\pageoriginaled\pageoriginale
% SMF's \skpt suppresses punctuation in a statement heading; with the
% standard-class compatibility environments the heading punctuation is fixed.
\providecommand{\skpt}{}
\providecommand{\noo}{\hbox{${}^\circ$}}
\providecommand{\et}{\textup{\'et}}
\providecommand{\Isom}{\textup{Isom}}
\providecommand{\llbracket}{[\![}
\providecommand{\rrbracket}{]\!]}
\definecolor{softgray}{gray}{0.94}
\definecolor{notegray}{gray}{0.25}
\definecolor{noteblue}{RGB}{232,241,250}
\definecolor{notebluetext}{RGB}{20,65,110}

% Human-readable statement environments used by the bounded English units.
\newenvironment{sgastatement}[1]
  {\par\addvspace{0.8\baselineskip}\noindent\textbf{#1.}\ \itshape}
  {\par\addvspace{0.35\baselineskip}}
\newenvironment{sgaremark}[1]
  {\par\addvspace{0.8\baselineskip}\noindent\textbf{#1.}\ }
  {\par\addvspace{0.35\baselineskip}}
\newenvironment{sgatheorem}[1]
  {\par\addvspace{0.8\baselineskip}\noindent\textbf{#1.}\ \itshape}
  {\par\addvspace{0.35\baselineskip}}
\newenvironment{sgalemma}[1]
  {\par\addvspace{0.8\baselineskip}\noindent\textbf{#1.}\ \itshape}
  {\par\addvspace{0.35\baselineskip}}
\newenvironment{sgaexample}[1]
  {\par\addvspace{0.8\baselineskip}\noindent\textbf{#1.}\ \itshape}
  {\par\addvspace{0.35\baselineskip}}
\newenvironment{sgaproof}
  {\par\addvspace{0.4\baselineskip}\noindent\textit{Proof.}\ }
  {\hfill$\square$\par\addvspace{0.4\baselineskip}}
\newenvironment{proofbridge}
  {\par\noindent\textit{Proof (reduction).}\ }
  {\par}
\newenvironment{enumeratea}
  {\begin{enumerate}[label=\alph*)]}
  {\end{enumerate}}
\newenvironment{enumeratei}
  {\begin{enumerate}[label=(\roman*)]}
  {\end{enumerate}}

% Source-like theorem names retained in the translated later exposés.
\theoremstyle{plain}
\newtheorem{theoreme}[subsection]{Theorem}
\newtheorem{theorem}[subsection]{Theorem}
\newtheorem{corollaire}[subsection]{Corollary}
\newtheorem{lemme}[subsection]{Lemma}
\newtheorem{lemma}{Lemma}[section]
\newtheorem{proposition}[subsection]{Proposition}
\newtheorem{propositionblah}[subsection]{Proposition}
\newtheorem{sublemme}[subsubsection]{Lemma}
\newtheorem{subsublemme}[subsubsection]{Lemma}
\newtheorem{conjecture}[subsection]{Conjecture}
\newtheorem{probleme}[subsection]{Problem}
\theoremstyle{definition}
\newtheorem{definition}[subsection]{Definition}
\newtheorem{exemple}[subsection]{Example}
\newtheorem{supdefinition}[subsection]{Definition}
\theoremstyle{remark}
\newtheorem{remarque}[subsection]{Remark}
\newtheorem{remarques}[subsection]{Remarks}
\newtheorem{subremarque}[subsubsection]{Remark}
\newtheorem*{remarquestar}{Remark}
\newtheorem*{theoremIXoneone}{Theorem 1.1}
\newtheorem*{sgaproposition}{Proposition}
\newtheorem*{sgacorollary}{Corollary}

\NewDocumentEnvironment{enonce*}{O{plain} m o}
  {\par\addvspace{0.8\baselineskip}\noindent\textbf{#2%
   \IfValueT{#3}{ (#3)}.}\ %
   \def\sgaenoncestyle{#1}\def\sgaremarkstyle{remark}%
   \ifx\sgaenoncestyle\sgaremarkstyle\normalfont\else\itshape\fi}
  {\par\addvspace{0.35\baselineskip}}
\renewcommand{\thesourcecorollary}{\thesection.\arabic{sourcecorollary}}
\NewDocumentEnvironment{sourcecorollary}{o}
  {\refstepcounter{sourcecorollary}\par\medskip\noindent
   \textbf{Corollary~\thesourcecorollary%
   \IfValueT{#1}{\markedfootnote{(3)}{#1}}.}\ \itshape}
  {\par\medskip}

% LaTeX kernels now define uppercase \Ref.  It must therefore be renewed,
% after hyperref is loaded, rather than introduced with \newcommand.
\renewcommand{\Ref}[1]{\textup{\ref{#1}}}
\renewcommand{\thechapter}{\Roman{chapter}}
\renewcommand{\chaptername}{Expos\'e}

\setlength{\emergencystretch}{2em}
\raggedbottom

% SGA2 reader: bounded LMMathExtension10 extraction support, proposal r1.
% Installed authority: pdfx/glyphtounicode-cmr.tex, whole-glyph base scalars.
% Encoding/TFM binding: lm-mathex.enc, pdftex.map, omxlmex.fd.
% This file is not a wholesale import of the pdfx mappings.
% Size variation selectors are deliberately omitted; size is typographic.
% Scope is only lmex10.tfm. No source formulas or glyph outlines are changed.
% Hats/tildes below are spacing accent identities, not base-order repair.
% No arrow tips, extenders, delimiter assemblies, mapsto pieces, or XY tips.
% Root must input this in the pdfLaTeX preamble and rebuild/recheck extraction.
\ifdefined\pdfglyphtounicode
  \pdfgentounicode=1
  % Whole, sized delimiters observed in the pinned reader.
  \pdfglyphtounicode{tfm:lmex10/parenleftbig}{0028}
  \pdfglyphtounicode{tfm:lmex10/parenrightbig}{0029}
  \pdfglyphtounicode{tfm:lmex10/braceleftbig}{007B}
  \pdfglyphtounicode{tfm:lmex10/bracerightbig}{007D}
  \pdfglyphtounicode{tfm:lmex10/parenleftBig}{0028}
  \pdfglyphtounicode{tfm:lmex10/parenrightBig}{0029}
  \pdfglyphtounicode{tfm:lmex10/parenleftbigg}{0028}
  \pdfglyphtounicode{tfm:lmex10/parenrightbigg}{0029}
  \pdfglyphtounicode{tfm:lmex10/parenleftBigg}{0028}
  \pdfglyphtounicode{tfm:lmex10/parenrightBigg}{0029}
  \pdfglyphtounicode{tfm:lmex10/braceleftBigg}{007B}
  \pdfglyphtounicode{tfm:lmex10/bracketleftBig}{005B}
  \pdfglyphtounicode{tfm:lmex10/bracketrightBig}{005D}
  % Whole n-ary operators; bigoplus is 2A01, not binary oplus 2295.
  \pdfglyphtounicode{tfm:lmex10/circleplustext}{2A01}
  \pdfglyphtounicode{tfm:lmex10/circleplusdisplay}{2A01}
  \pdfglyphtounicode{tfm:lmex10/circlemultiplytext}{2A02}
  \pdfglyphtounicode{tfm:lmex10/summationtext}{2211}
  \pdfglyphtounicode{tfm:lmex10/uniontext}{22C3}
  \pdfglyphtounicode{tfm:lmex10/intersectiontext}{22C2}
  \pdfglyphtounicode{tfm:lmex10/summationdisplay}{2211}
  \pdfglyphtounicode{tfm:lmex10/productdisplay}{220F}
  \pdfglyphtounicode{tfm:lmex10/uniondisplay}{22C3}
  % Spacing accents per pdfx; prevents b/c/e/f/g fallback identities.
  % Does not establish the logical order or scope of the accented expression.
  \pdfglyphtounicode{tfm:lmex10/hatwide}{02C6}
  \pdfglyphtounicode{tfm:lmex10/hatwider}{02C6}
  \pdfglyphtounicode{tfm:lmex10/tildewide}{02DC}
  \pdfglyphtounicode{tfm:lmex10/tildewider}{02DC}
  \pdfglyphtounicode{tfm:lmex10/tildewidest}{02DC}
\else
  \errmessage{SGA2 lmex Unicode support requires pdfTeX primitives}
\fi

% END INCLUDED active_source/sga2_reader_macros.tex


% Roman exposé numbers need wider number boxes than the standard book TOC.
\setlength{\cftchapnumwidth}{3.2em}
\setlength{\cftsecnumwidth}{4.2em}
\setlength{\cftsubsecnumwidth}{5.4em}

\hypersetup{
  pdftitle={SGA 2: Local Cohomology of Coherent Sheaves and Local and Global Lefschetz Theorems},
  pdfauthor={Alexander Grothendieck et al.; English translation project},
  pdfsubject={Source-aligned English reader assembled from the corrected arXiv TeX authority math/0511279v1},
  pdfkeywords={SGA 2, local cohomology, Lefschetz theorems, algebraic geometry, English translation}
}

\begin{document}
\frontmatter

\title{Seminar on Algebraic Geometry at Bois Marie\\
   1962\\[0.4em]
   Local Cohomology of Coherent Sheaves and\\
   Local and Global Lefschetz Theorems\\
   (SGA~2)}

\author{Alexander Grothendieck\\
  {\small (written up by a group of participants)}\\[0.5em]
  With an additional expos\'e by Mich\`ele Raynaud}

\maketitle

\chapter*{Abstract}
\addcontentsline{toc}{chapter}{Abstract}
This volume is a newly updated edition of the book
\emph{Cohomologie locale des faisceaux coh\'erents et th\'eor\`emes de
Lefschetz locaux et globaux (SGA~2)}, Advanced Studies in Pure Mathematics~2,
North-Holland Publishing Company, Amsterdam, 1968, by A.~Grothendieck
\emph{et al.}

The monograph gives necessary and sufficient conditions for the finiteness
of the local cohomology sheaves of coherent sheaves. These results yield
algebraization theorems which, together with purity results also proved in
the text, lead to Lefschetz-type theorems for both the fundamental group and
the Picard group.

\medskip
\noindent\textbf{Keywords.} Algebraization, Lefschetz theorem, local duality,
local cohomology, depth, homotopical depth, fundamental group, Picard group.

\smallskip
\noindent\textbf{2020 Mathematics Subject Classification.}
14-02, 14B10, 14B15, 14B20, 14C22, 14J70, 32S20.

\chapter*{Editor's Preface}
\addcontentsline{toc}{chapter}{Editor's Preface}

The present text is a new updated edition of the book
``Cohomologie locale des faisceaux coh\'erents et th\'eor\`emes de Lefschetz
locaux et globaux (SGA~2)'', Advanced Studies in Pure Mathematics~2,
North-Holland Publishing Company, Amsterdam, 1968, by A.~Grothendieck
\emph{et al.} It is the second part of the SGA project initiated by
B.~Edixhoven, who prepared a new edition of SGA~1. This version is intended
to reproduce the original text as closely as possible, while incorporating
minor typographical corrections and numbered editorial footnotes explaining
the current status of questions raised in the first edition. Further details
have also been supplied for some proofs. To avoid confusion, original notes
retain star-based numbering, whereas new editorial notes use Arabic numbers.
The page numbers of the original edition are shown in the margins.

I thank the mathematicians who carried out most of the initial typesetting
in \LaTeXe: L.~Bayle, N.~Borne, O.~Brinon, J.~Buresi, M.~Chardin,
F.~Ducrot, P.~Graftiaux, F.~Han, P.~Karwasz, L.~Koelblen, D.~Madore,
S.~Morel, D.~Naie, B.~Osserman, J.~Riou, and V.~S\'echerre, as well as
C.~Sabbah for adapting the text to the SMF format. I also thank J.-B.~Bost,
P.~Colmez, O.~Gabber, W.~Fulton, S.~Kleiman, F.~Orgogozo, M.~Raynaud, and
J.-P.~Serre for their comments and advice.

\begin{flushright}
Yves Laszlo, editor
\end{flushright}

\medskip
\noindent\textit{English-reader note.} This reader is translated from the
corrected and annotated TeX distributed as arXiv:math/0511279v1. Mathematical
notation, labels, equations, editorial notes, and source ordering are retained.
External English candidates were comparison witnesses only, not source
authority. The editable English Markdown lineage by Jacob Reinhold
(\url{https://github.com/jcreinhold/sga}, snapshot
\texttt{e7a259f3f8608ad3edf9bf6eead3fd504dd2d23e}) was consulted as one such
comparison witness and is credited under its stated CC BY 4.0 license for
the translation contribution; that license does not cover the underlying
French work.

\cleardoublepage
\tableofcontents

\mainmatter
% BEGIN INCLUDED active_source/parts/00_introduction.tex
\chapter*{Introduction} \label{I.0}
\markboth{INTRODUCTION}{INTRODUCTION}

We\pageoriginale present here, in revised and expanded form, a photo-offset
reissue of the second \sisi{ALGEBRAIC GEOMETRY SEMINAR of the INSTITUT
DES HAUTES \'ETUDES SCIENTIFIQUES}{Algebraic Geometry Seminar of the
Institut des Hautes \'Etudes Scientifiques}, held in 1962 and originally
issued in mimeographed form.

The reader is referred to the Introduction to the first of these seminars
(cited below as \SGA 1) for the aims pursued by these seminars and their
relation to the \sisi{\'EL\'EMENTS DE G\'{E}OM\'ETRIE
ALG\'EBRIQUE}{\'El\'ements de G\'eom\'etrie Alg\'ebrique}.

The text of Expos\'es~\Ref{I} through~\Ref{XI} was written up as the
seminar proceeded, from my oral lectures and handwritten notes, by a group
of participants comprising I\ptbl~\sisi{GIROGUITTI}{Giorgiutti},
J\ptbl~\sisi{GIRAUD}{Giraud}, Mlle M\ptbl~\sisi{JAFFE}{Jaffe} (now
Mme M\ptbl~\sisi{HAKIM}{Hakim}), and A\ptbl~\sisi{LAUDAL}{Laudal}.
These notes were originally regarded as provisional and intended for very
limited circulation, pending their absorption into the EGA---an absorption
that has now become problematic, to say the least, as it has for the other
parts of the SGA. As the warning to the original edition said, the
\og confidential\fg character of the notes was meant to excuse certain
\og stylistic weaknesses\fg, doubtless more apparent in the present
seminar, \SGA 2, than in the others. I have tried, as far as possible, to
remedy them in the present reissue by a fairly thorough revision of the
original text. In particular, I have harmonized the systems used to number
statements in the different Expos\'es, introducing everywhere the same
decimal system, already used in most of the original Expos\'es in the
present \SGA 2 as well as in all the other parts of the SGA. This led me,
in particular, to revise completely the numbering\nde{The original
numbering has been preserved as far as possible, with the adverb
\emph{bis} added when ambiguous duplicates occurred here and there.} of
the statements in Expos\'es~\Ref{III} through~\Ref{VIII}\pageoriginale
(and consequently the references to those Expos\'es)\sfootnote{It goes
without saying that every reference to \SGA 2 appearing in those parts of
the SGA published in the \sisi{SERIES IN PURE MATHEMATICS}{Series in Pure
Mathematics} refers to the present volume, not to the original edition of
\SGA 2!}. I have also tried to remove from the original text the principal
typing and syntactical errors, which were numerous and troublesome.
Moreover, Mme M\ptbl~\sisi{HAKIM}{Hakim} kindly undertook to rewrite
Expos\'e~\Ref{IV} in a less telegraphic style than the original. As in the
other SGA reissues, I have also added a number of footnotes, either to give
additional references or to indicate the status of questions on which
progress has been made since the original text was written. Finally, the
seminar has been enlarged by a new Expos\'e, \sisi{namely}{namely}
Expos\'e~\Ref{XIV}, written by Mme\sisi{.}{}
\sisi{MICH\`ELE RAYNAUD}{Mich\`ele Raynaud} in 1967, which takes up and
completes suggestions contained in the \og Comments on
Expos\'e~\Ref{XIII}\fg (\Ref{XIII}~\Ref{XIII.6}), written in March 1963.
This Expos\'e treats Lefschetz-type theorems from the viewpoint of \'etale
cohomology, using the results on \'etale cohomology presented in \SGA 4
and \SGA 5 (to appear in the same
\sisi{SERIES IN PURE MATHEMATICS}{Series in Pure Mathematics})\nde{In
fact, these seminars were published by Springer (numbers 269, 279, 305,
and 589), but unfortunately are out of print.}; it is therefore less
\og elementary\fg in nature than the other Expos\'es of the present
volume, which use little more than the substance of Chapters~I through~III
of the EGA. Here is an outline of the contents of the present volume.
Expos\'e~\Ref{I} contains the sorites of \og cohomology with support in
$Y$\fg, $H_Y^*(X,F)$, where $Y$ is a closed subset of a space $X$,
cohomology that may be interpreted as the cohomology of $X$ \emph{modulo}
the open subset $X-Y$, and that is the abutment of a very useful
\og \emph{local-to-global spectral sequence}\fg
\Ref{I}~\Ref{I.2.6}, involving the \emph{sheaves} of cohomology
\og with support in $Y$\fg, $\sheaf{H}_Y^*(F)$\nde{The underlined
notation has been retained for sheafified versions of functors, since the
calligraphic analogue of $\Gamma$ is not clear.}. In many questions this
formalism can\pageoriginale play a \og localization\fg role analogous to
that played by tubular neighborhoods of $Y$ in differential geometry.
Expos\'e~\Ref{II} studies the preceding notions for quasi-coherent sheaves
on preschemes, and Expos\'e~\Ref{III} relates them to the classical notion
of \emph{depth} (\Ref{III}~\Ref{III.3.3}).

Expos\'es~\Ref{IV} and~\Ref{V} give forms of \emph{local duality}, which
may be compared with Serre's projective duality theorem
(\Ref{XII}~\Ref{XII.1.1}). We note that both kinds of duality theorem are
substantially generalized in \sisi{HARTSHORNE}{Hartshorne}'s seminar
(cited in the footnote at the end of \Exp \Ref{IV})\refstepcounter{toto}
\de{\label{noteconrad}Hartshorne's book contains sign errors and, more
importantly, does not actually prove the compatibility of the trace with
base change. Conrad completely redid this work and proved this crucial and
highly nontrivial compatibility (\textit{Grothendieck duality and base
change}, Lect. Notes in Math., vol.~1750, Springer-Verlag, Berlin, 2000).
Unfortunately, errors remain (\cf two preprints: Conrad~B.,
\og Clarifications and corrections to ``\textit{Grothendieck duality and
base change}''\fg\ and \og An addendum to Chapter 5 of
``\textit{Grothendieck duality and base change}''\fg). For a more concrete
aspect, with particular attention to the notion of residue, see Lipman's
work, especially Lipman~J., \textit{Dualizing sheaves, differentials and
residues on algebraic varieties}, Ast\'erisque, vol.~117, Soci\'et\'e
math\'ematique de France, 1984. A categorical proof of the duality theorem,
based on Brown's representability theorem, was obtained by Neeman
(Neeman~A., \og The Grothendieck duality theorem via Bousfield's techniques
and Brown representability\fg, \emph{J.~Amer. Math. Soc.} \textbf{9}
(1996), \numero 1, p\ptbl~205--236).}

Expos\'es~\Ref{VI} and~\Ref{VII} give elementary technical notions used in
Expos\'e~\Ref{VIII} to prove the \emph{finiteness theorem}
(\Ref{VIII}~\Ref{VIII.2.3}). This theorem gives necessary and sufficient
conditions, for a coherent sheaf $F$ on a noetherian scheme $X$, for the
local cohomology sheaves $\sheaf{H}_Y^i(F)$ to be coherent for $i\leq n$
(or, equivalently, for the sheaves $\R^if_*(F|X-Y)$ to be coherent for
$i\leq n-1$, where $f\colon X-Y\to X$ is the inclusion). It is one of the
central technical results of the seminar, and in Expos\'e~\Ref{IX} we show
how a theorem of this kind can be used to establish a \og comparison
theorem\fg and an \og existence theorem\fg in formal geometry, following
and generalizing the use made in (\EGA III~\oldS\S 4 and~5) of the
finiteness theorem for a proper morphism.

These last results are applied in~\Ref{X} and~\Ref{XI}, devoted
respectively to Lefschetz-type theorems for the fundamental group and
for\pageoriginale the Picard group. These theorems compare, under certain
conditions, the invariants ($\pi_1$ or $\Pic$) associated respectively
with a scheme $X$ and a subscheme $Y$, which plays the role of a
hyperplane section, and in particular give conditions under which they
are isomorphic. Roughly speaking, the hypotheses allow one to pass from
$Y$ to the formal completion of $X$ along $Y$, and then to apply the
results of~\Ref{IX} in order to pass from there to an open
\emph{neighborhood} $U$ of $Y$ in $X$. To pass from $U$ to $X$, one also
needs information of a \og purity\fg or \og parafactoriality\fg type for
the local rings of $X$ at the points of $Z=X-U$, which is a finite
discrete set in the cases under consideration. This explains the
interaction, in the proofs of Expos\'es~\Ref{X}, \Ref{XI}, and \Ref{XII},
between local and global results, particularly in certain inductions. The
main results obtained in~\Ref{X} and~\Ref{XI} are the theorems of
\emph{local nature} \Ref{X}~\Ref{X.3.4} (the \emph{purity theorem}) and
\Ref{XI}~\Ref{XI.3.14} (the \emph{parafactoriality theorem}). Notice that
these theorems are proved by cohomological techniques of an essentially
global nature. In~\Ref{XII}, using the preceding local results, one
obtains global variants for projective schemes over a field, or more
generally over a more or less arbitrary base scheme; typical statements
include \Ref{XII}~\Ref{XII.3.5} and \Ref{XII}~\Ref{XII.3.7}.

In~\Ref{XIII}, we review some of the many problems and conjectures
suggested by the results and methods of the seminar. Perhaps the most
interesting concern cohomological and homotopical Lefschetz-type theorems
for complex analytic spaces; \cf \Ref{XIII}, pages 26 and following
\de{Essentially all the conjectures stated in~\Ref{XIII} and \Ref{XIV}
have now been proved; see the footnotes in those sections for references
and comments.}. In the context of the \'etale cohomology of
schemes\pageoriginale, the corresponding conjectures are proved
in~\Ref{XIV} by a duality technique that should apply equally in the
complex analytic case (\cf the comments in \Ref{XIII}, p\ptbl~25, and
\Ref{XIV}~\Ref{XIV.6.4}). But the corresponding homotopical statements
for analytic spaces, and more particularly those involving the
fundamental group, seem to require entirely new techniques
(\cf \Ref{XIV}~\Ref{XIV.6.4}).

I am happy to thank all those who, in various capacities, helped bring the
present volume to publication, including the collaborators already named
in this Introduction. In particular, I wish to thank Mlle
\sisi{CHARDON}{Chardon} for the good grace with which she carried out the
thankless task of physically preparing the final manuscript for
photo-offset reproduction.

\bigskip\hfill Bures-sur-Yvette, April~1968
\par\smallskip \hfill A\ptbl~Grothendieck.

% END INCLUDED active_source/parts/00_introduction.tex

% BEGIN INCLUDED active_source/parts/01_expose_I.tex
% SGA 2, Expose I — English mathematical body.
% Authority coverage: official corrected arXiv TeX lines 87--503.
% Mechanically extracted without alteration from the two source-aligned units
% named below. The repeated Expose heading and the Section 1 bibliography copy
% were removed; the complete bibliography from Section 2 is retained.

\makeatletter
\providecommand{\sgareaderlabel}[2]{%
  \begingroup
  \def\@currentlabel{#2}%
  \phantomsection\label{#1}%
  \endgroup}
\makeatother

\chapter[Global and Local Cohomological Invariants]
{Global and Local Cohomological Invariants Relative to a Closed Subspace}
\label{I}

\section*{1. The functors $\Gamma_Z$ and $\uGamma_Z$}
\sgareaderlabel{I.1}{1}
\addcontentsline{toc}{section}{1. The functors Gamma Z and underlined Gamma Z}

\srcpage{6}
Let $X$ be a topological space, and let $\uC_X$ be the category of
sheaves of abelian groups on $X$.  Let $\Phi$ be a family of supports in
the sense of Cartan.  The functor $\Gamma_\Phi$ on $\uC_X$ is defined by
\begin{equation}
  \Gamma_\Phi(F)
  =\{f\in\Gamma(F)\mid \operatorname{supp}(f)\in\Phi\}.
  \tag{1}\label{eq:I.1}
\end{equation}
If $Z$ is a closed subset of $X$, by an abuse of notation we denote by
$\Gamma_Z$ the functor $\Gamma_\Phi$, where $\Phi$ is the family of
closed subsets of $X$ contained in $Z$.\footnote{At this single occurrence
the French printed and TeX sources underline $\Gamma_Z$.  The defining
equation~\eqref{eq:I.2} and the systematic distinction used thereafter
identify the intended functor here as the global, non-underlined $\Gamma_Z$.
The source anomaly and this explicit resolution are recorded in the tranche
comparison ledger.}
Thus
\begin{equation}
  \Gamma_Z(F)
  =\{f\in\Gamma(F)\mid \operatorname{supp}(f)\subset Z\}.
  \tag{2}\label{eq:I.2}
\end{equation}

We wish to extend this definition to the case in which $Z$ is a
\emph{locally closed} subset of $X$, hence closed in a suitable open subset
$V$ of $X$.  In that case we set
\begin{equation}
  \Gamma_Z(F)=\Gamma_Z(\restr{F}{V}).
  \tag{3}\label{eq:I.3}
\end{equation}
We must verify that $\Gamma_Z(F)$ ``does not depend'' on the chosen open
set.  It is enough to show that, if $V'$ is open and
$V\supset V'\supset Z$, then the map
$\rho^V_{V'}:F(V)\to F(V')$ maps
$\Gamma_Z(\restr{F}{V})$ isomorphically onto
$\Gamma_Z(\restr{F}{V'})$.  But
\begin{equation}
  \Gamma_Z(\restr{F}{V})=\ker\rho^V_{V-Z}.
  \tag{4}\label{eq:I.4}
\end{equation}
Consequently, if $f\in\Gamma_Z(\restr{F}{V})$ and
$\rho^V_{V'}(f)=\rho^V_{V-Z}(f)=0$, then $f=0$, since
$(V',V-Z)$ is a covering of $V$.  Likewise, if
$f'\in\Gamma_Z(\restr{F}{V'})$, then $f'\in F(V')$ and
\srcpage{7}$0\in F(V-Z)$ define an element $f\in F(V)$ such that
$\rho^V_{V'}(f)=f'$ and $f\in\Gamma_Z(\restr{F}{V})$.
Thus $\rho^V_{V'}$ induces an isomorphism
\[
  \Gamma_Z(\restr{F}{V})\longrightarrow
  \Gamma_Z(\restr{F}{V'}).
\]

Note that every open subset $W$ of $Z$ is induced by an open subset $U$
of $X$ in which $W$ is closed.  It follows that
$W\mapsto\Gamma_W(F)$ defines a presheaf on $Z$; one verifies that it is
a sheaf, denoted by $i^!(F)$, where $i:Z\to X$ is the canonical
immersion.  We obtain
\begin{equation}
  \Gamma_Z(F)=\Gamma(i^!(F)).
  \tag{5}\label{eq:I.5}
\end{equation}
The sheaf $i^!(F)$ is a subsheaf of $i^*(F)$.  Indeed, the canonical
homomorphism
\[
  \Gamma(\restr{F}{U})=\Gamma(U,F)\longrightarrow
  \Gamma(U\cap Z,i^*(F))
\]
is injective on
$\Gamma_{U\cap Z}(\restr{F}{U})\subset\Gamma(\restr{F}{U})$.
In summary, we have the following result.

\begin{sgastatement}{Proposition 1.1}\sgareaderlabel{I.1.1}{1.1}
There exists a unique subsheaf $i^!(F)$ of $i^*(F)$ such that, for every
open subset $U$ of $X$ for which $U\cap Z$ is closed in $U$, the map
\[
  \Gamma(\restr{F}{U})=\Gamma(U,F)\longrightarrow
  \Gamma(U\cap Z,i^*(F))
\]
induces an isomorphism
\[
  \Gamma_{U\cap Z}(\restr{F}{U})\longrightarrow
  \Gamma(U\cap Z,i^!(F)).
\]
\end{sgastatement}

If $Z$ is open, then simply
\begin{equation}
  i^!(F)=i^*(F)=\restr{F}{Z},
  \qquad \Gamma_Z(F)=\Gamma(Z,F).
  \tag{6}\label{eq:I.6}
\end{equation}

Suppose again that $Z$ is arbitrary.  As $U$ ranges through the open
subsets of $X$, the assignment
\[
  U\longmapsto \Gamma_{U\cap Z}(\restr{F}{U})
  =\Gamma(U\cap Z,i^!(F))
\]
is a sheaf on $X$, which we denote by $\uGamma_Z(F)$.  More precisely,
the preceding formula, which expresses the fact that $i^!$ commutes with
restriction to open subsets, gives an isomorphism
\begin{equation}
  \uGamma_Z(F)=i_*(i^!(F)).
  \tag{7}\label{eq:I.7}
\end{equation}
By definition,\srcpage{8} for every open subset $U$ of $X$,
\begin{equation}
  \Gamma(U,\uGamma_Z(F))
  =\Gamma_{U\cap Z}(\restr{F}{U}).
  \tag{8}\label{eq:I.8}
\end{equation}

Here there is a characteristic difference between the case in which
$Z$ is closed and the case in which $Z$ is open.  In the first case,
formula~\eqref{eq:I.8} shows that $\uGamma_Z(F)$ may be regarded as a
subsheaf of $F$; hence there is a \emph{canonical immersion}
\begin{equation}
  \uGamma_Z(F)\lhook\joinrel\longrightarrow F.
  \tag{$8'$}\label{eq:I.8'}
\end{equation}
When $Z$ is open, on the other hand, equation~\eqref{eq:I.6} shows that
the right-hand side of equation~\eqref{eq:I.8} is
$\Gamma(U\cap Z,F)$ and therefore receives a homomorphism from
$\Gamma(U,F)$.  Thus there is a \emph{canonical homomorphism} in the
opposite direction:\footnote{The original reference was \hyperref[eq:I.8]{(8)}.}
\begin{equation}
  F\longrightarrow\uGamma_Z(F).
  \tag{$8''$}\label{eq:I.8''}
\end{equation}
This is none other than the canonical homomorphism
\[
  F\longrightarrow i_*i^*(F),
\]
in view of the isomorphism
\begin{equation}
  \uGamma_Z(F)\simeq i_*i^*(F)
  \tag{6 bis}\label{eq:I.6bis}
\end{equation}
deduced from equations~\eqref{eq:I.6} and~\eqref{eq:I.7}.

For variable $F$, the objects $\Gamma_Z(F)$, $\uGamma_Z(F)$, and
$i^!(F)$ may of course be regarded as functors in $F$, with values,
respectively, in the category of abelian groups, the category of sheaves
of abelian groups on $X$, and the category of sheaves of abelian groups
on $Z$.  It is sometimes convenient to interpret the functor
\[
  i^!: \uC_X\longrightarrow\uC_Z
\]
as the right adjoint of the familiar functor
\[
  i_!: \uC_Z\longrightarrow\uC_X
\]
defined by the following proposition.

\srcpage{9}
\begin{sgastatement}{Proposition 1.2}\sgareaderlabel{I.1.2}{1.2}
Let $G$ be a sheaf of abelian groups on $Z$.  There exists a unique
subsheaf of $i_*(G)$, denoted $i_!(G)$, such that for every open subset
$U$ of $X$ the (identity) isomorphism
\[
  \Gamma(U\cap Z,G)=\Gamma(U,i_*(G))
\]
induces an isomorphism
\[
  \Gamma_{\Phi_{U\cap Z,U}}(U\cap Z,G)=\Gamma(U,i_!(G)),
\]
where $\Phi_{U\cap Z,U}$ denotes the family of subsets of $U\cap Z$
that are closed in $U$.
\end{sgastatement}

The verification reduces to observing that the left-hand side is a
sheaf as $U$ varies: namely, for a section of $i_*(G)$ on $U$, viewed
as a section of $G$ on $U\cap Z$, the property of having support closed
\emph{in $U$} is \emph{local on $U$}.  The sheaf $i_!(G)$ just defined
is also known as the \emph{sheaf obtained from $G$ by extension by zero}
outside $Z$; see Godement~\cite{Godement}.  In particular, if $Z$ is
closed, then
\begin{equation}
  i_!(G)=i_*(G).
  \tag{9}\label{eq:I.9}
\end{equation}
In general, however, the canonical injection
$i_!(G)\to i_*(G)$ is not an isomorphism, as is already well known when
$Z$ is open.  Evidently $i_!(G)$ depends functorially on $G$ (and is in
fact an exact functor of $G$).  We therefore have:

\begin{sgastatement}{Proposition 1.3}\sgareaderlabel{I.1.3}{1.3}
There is an isomorphism of bifunctors in $G$ and $F$, where $G$ is a
sheaf of abelian groups on $Z$ and $F$ is a sheaf of abelian groups on
$X$:
\begin{equation}
  \operatorname{Hom}(i_!(G),F)
  =\operatorname{Hom}(G,i^!(F)).
  \tag{10}\label{eq:I.10}
\end{equation}
\end{sgastatement}

To define such an isomorphism is equivalent to defining functorial
homomorphisms
\[
  i_!i^!(F)\longrightarrow F,
  \qquad
  G\longrightarrow i^!i_!(G),
\]
satisfying the familiar compatibility conditions (see, for example,
Shih's expos\'e in the Cartan seminar on cohomological operations).

Recalling\srcpage{10} that $i_!$ is exact, and hence takes
monomorphisms to monomorphisms, we conclude:

\begin{sgastatement}{Corollary 1.4}\sgareaderlabel{I.1.4}{1.4}
If $F$ is injective, then $i^!(F)$ is injective; consequently
$\uGamma_Z(F)=i_*i^!(F)$ is also injective.
\end{sgastatement}

Replacing $X$ by a variable open subset $U$ of $X$, \hyperref[I.1.3]{Proposition~1.3}
also gives:

\begin{sgastatement}{Corollary 1.5}\sgareaderlabel{I.1.5}{1.5}
There is an isomorphism, functorial in $F$ and $G$,
\begin{equation}
  \uHom(i_!(G),F)=i_*\bigl(\uHom(G,i^!(F))\bigr).
  \tag{11}\label{eq:I.11}
\end{equation}
\end{sgastatement}

Taking $G$ to be the constant sheaf on $Z$ defined by $\ZZ$, denoted
$\ZZ_Z$, \hyperref[I.1.3]{Proposition~1.3} and \hyperref[I.1.5]{Corollary~1.5} specialize as follows.

\begin{sgastatement}{Corollary 1.6}\sgareaderlabel{I.1.6}{1.6}
There are isomorphisms, functorial in $F$,
\begin{equation}
  \begin{aligned}
    \Gamma_Z(F)&=\operatorname{Hom}(\ZZ_{Z,X},F),\\
    \uGamma_Z(F)&=\uHom(\ZZ_{Z,X},F),
  \end{aligned}
  \tag{12}\label{eq:I.12}
\end{equation}
where $\ZZ_{Z,X}=i_!(\ZZ_Z)$ is the sheaf of abelian groups on $X$
obtained from the constant sheaf $\ZZ_Z$ on $Z$ by extension by zero
outside $Z$.
\end{sgastatement}

\begin{sgaremark}{Remark 1.7}\sgareaderlabel{I.1.7}{1.7}
Suppose that $X$ is a ringed space, equip $Z$ with the sheaf of rings
$\mathcal O_Z=i^{-1}(\mathcal O_X)$, and now let $\uC_X$ and $\uC_Z$
denote the categories of modules on $X$ and $Z$, respectively.  The
preceding considerations extend word for word, with $F$ a module on $X$
and $G$ a module on $Z$, and with \hyperref[I.1.3]{statements 1.3}--\hyperref[I.1.6]{1.6} interpreted
accordingly.
\end{sgaremark}

To conclude these generalities, let us examine what happens when the
locally closed subset $Z$ is changed.  Let $Z'\subset Z$ be another
locally closed subset, and let
\[
  j:Z'\longrightarrow Z,
  \qquad i':Z'\longrightarrow X,
  \qquad i'=ij
\]
be the canonical inclusions.\srcpage{11}  There are functorial
isomorphisms
\begin{equation}
  (ij)^!=j^!i^!,
  \qquad
  (ij)_!=i_!j_!.
  \tag{13}\label{eq:I.13}
\end{equation}
The first isomorphism in~\eqref{eq:I.13} defines a functorial
isomorphism
\begin{equation}
  \Gamma_{Z'}(F)
  =\Gamma(Z',(ij)^!(F))
  \simeq\Gamma(Z',j^!(i^!(F)))
  =\Gamma_{Z'}(i^!(F)).
  \tag{14}\label{eq:I.14}
\end{equation}

Suppose now that $Z'$ is \emph{closed} in $Z$, and let
\[
  Z''=Z-Z'
\]
be its complement in $Z$.  This is open in $Z$ and hence locally closed
in $X$.  Applying the canonical injection~\eqref{eq:I.8'} to
$i^!(F)$ on $Z$ with the subset $Z'$, and using~\eqref{eq:I.14}, gives
an injective functorial canonical homomorphism
\begin{equation}
  \Gamma_{Z'}(F)\longrightarrow\Gamma_Z(F).
  \tag{15}\label{eq:I.15}
\end{equation}
If $Z'$ is replaced by $Z''$ in~\eqref{eq:I.14} and
equation~\eqref{eq:I.8''} is used, one obtains a functorial
canonical homomorphism
\begin{equation}
  \Gamma_Z(F)\longrightarrow\Gamma_{Z''}(F).
  \tag{$15'$}\label{eq:I.15'}
\end{equation}

\begin{sgastatement}{Proposition 1.8}\sgareaderlabel{I.1.8}{1.8}
Under the preceding hypotheses, the sequence of functorial
homomorphisms
\begin{equation}
  0\longrightarrow\Gamma_{Z'}(F)
  \longrightarrow\Gamma_Z(F)
  \longrightarrow\Gamma_{Z''}(F)
  \tag{16}\label{eq:I.16}
\end{equation}
is exact.  If $F$ is flasque, the sequence remains exact after a zero is
added on the right.
\end{sgastatement}

\begin{sgaproof}
Replace $X$ by an open subset $V$ in which $Z$ is closed.  This reduces
the assertion to the case in which $Z$, and hence $Z'$, is closed.
Then $Z''$ is closed in the open subset $X-Z'$, and there is a canonical
injection
\[
  \Gamma_{Z''}(F)\longrightarrow\Gamma(X-Z',F).
\]
The exactness of~\eqref{eq:I.16} then says simply that the sections of
$F$ with support in $Z'$ are precisely those whose restriction to
$X-Z'$ is zero.

\srcpage{12}If $F$ is flasque, every element of $\Gamma_{Z''}(F)$, regarded as a
section of $F$ on $X-Z'$, extends to a section of $F$ on $X$.  The
latter section plainly has support in $Z$, proving that the last
homomorphism in~\eqref{eq:I.16} is then surjective.
\end{sgaproof}

\begin{sgastatement}{Corollary 1.9}\sgareaderlabel{I.1.9}{1.9}
There is a functorial exact sequence
\begin{equation}
  0\longrightarrow\uGamma_{Z'}(F)
  \longrightarrow\uGamma_Z(F)
  \longrightarrow\uGamma_{Z''}(F),
  \tag{16 bis}\label{eq:I.16bis}
\end{equation}
and, if $F$ is flasque, the sequence remains exact after a zero is added
on the right.
\end{sgastatement}

\hyperref[I.1.8]{Proposition~1.8} may be interpreted, by means of \hyperref[I.1.6]{Corollary~1.6}, as a
statement about the functors $\operatorname{Hom}$ and $\uHom$, as
follows.  First note that if $G$ is a sheaf of abelian groups on $Z$,
inducing the sheaves $j^*(G)$ and $k^*(G)$ on $Z'$ and $Z''$,
respectively, where $j:Z'\to Z$ and $k:Z''\to Z$ are the canonical
injections, then there is a canonical exact sequence of sheaves on $X$:
\begin{equation}
  0\longrightarrow k^*(G)_X
  \longrightarrow G_X
  \longrightarrow j^*(G)_X
  \longrightarrow 0.
  \tag{17}\label{eq:I.17}
\end{equation}
Here, to simplify notation, the subscript $X$ denotes the sheaf on $X$
obtained by extension by zero outside the space on which the sheaf in
question is defined.  Sequence~\eqref{eq:I.17} generalizes a familiar
exact sequence for $Z=X$ (see Godement~\cite{Godement}); it follows from
that sequence by writing it on $Z$ and applying the functor $i_!$.
Taking $G=\ZZ_Z$, we obtain in particular:

\begin{sgastatement}{Proposition 1.10}\sgareaderlabel{I.1.10}{1.10}
Under the preceding hypotheses, there is an exact sequence of sheaves of
abelian groups on $X$:
\begin{equation}
  0\longrightarrow\ZZ_{Z'',X}
  \longrightarrow\ZZ_{Z,X}
  \longrightarrow\ZZ_{Z',X}
  \longrightarrow 0.
  \tag{18}\label{eq:I.18}
\end{equation}
The two exact sequences in \hyperref[I.1.8]{Proposition~1.8} and \hyperref[I.1.9]{Corollary~1.9} are exactly
the sequences obtained from~\eqref{eq:I.18} by applying, respectively,
the functors $\operatorname{Hom}(-,F)$ and $\uHom(-,F)$.
\end{sgastatement}

This also gives an immediate proof that sequences~\eqref{eq:I.16} and
\eqref{eq:I.16bis} remain exact after a zero is added on the right,
provided that $F$ is \emph{injective}.
% CONTINUATION: source-aligned Section 2 unit, French lines 280--503.
\section*{2. The functors $\mathrm H_Z^*(X,F)$ and $\uH_Z^*(F)$}
\sgareaderlabel{I.2}{2}
\addcontentsline{toc}{section}{2. The functors H Z star and underlined H Z star}

\srcpage{13}
\begin{sgastatement}{Definition 2.1}\sgareaderlabel{I.2.1}{2.1}
The derived functors, in the variable $F$, of the functors $\Gamma_Z(F)$
and $\uGamma_Z(F)$ are denoted by $\mathrm H_Z^*(X,F)$ and $\uH_Z^*(F)$,
respectively.
\end{sgastatement}

They are cohomological functors with values, respectively, in the category
of abelian groups and in the category of sheaves of abelian groups on $X$.
When $Z$ is closed, $\mathrm H_Z^*(X,F)$ is, by definition, simply
$\mathrm H_\Phi^*(X,F)$, where $\Phi$ denotes the family of closed subsets
of $X$ contained in $Z$. When $Z$ is open, we shall see that
$\mathrm H_Z^*(X,F)$ is simply
$\mathrm H^*(Z,F)=\mathrm H^*(Z,\restr{F}{Z})$, by the following
proposition.

\begin{sgastatement}{\hyperref[I.2.2]{Proposition 2.2} (Excision theorem)}
\sgareaderlabel{I.2.2}{2.2}
Let $V$ be an open subset of $X$ containing $Z$. There is an isomorphism of
cohomological functors in $F$:
\begin{equation}
  \mathrm H_Z^*(X,F)
  \longrightarrow
  \mathrm H_Z^*(V,\restr{F}{V}).
  \tag{19}\label{eq:I.19}
\end{equation}
\end{sgastatement}

Indeed, there is a functorial isomorphism
$\Gamma_Z^X\simeq\Gamma_Z^Vj^!$, where $j\colon V\to X$ is the inclusion
and $j^!$ is therefore the restriction functor (cf. \hyperref[eq:I.14]{equation~(14)}). This
last functor is exact and, by \hyperref[I.1.4]{Corollary~1.4}, takes injectives to injectives;
the isomorphism~\eqref{eq:I.19} follows immediately.

When $Z$ is open, we may take $V=Z$, and obtain:

\begin{sgastatement}{Corollary 2.3}\sgareaderlabel{I.2.3bof}{2.3}
Suppose that $Z$ is open. There is an isomorphism of cohomological functors
\begin{equation}
  \mathrm H_Z^*(X,F)=\mathrm H^*(Z,F).
  \tag{20}\label{eq:I.20}
\end{equation}
\end{sgastatement}

The isomorphisms of \hyperref[I.1.6]{Corollary~1.6} and the definitions give the following
(cf.\ \cite{Tohoku}).

\begin{sgastatement}{Proposition 2.3 bis}\sgareaderlabel{I.2.3bis}{2.3 bis}
There are\srcpage{14} isomorphisms of cohomological functors
\begin{align}
  \mathrm H_Z^*(X,F)
    &\simeq \operatorname{Ext}^*(X;\ZZ_{Z,X},F),
    \tag{21}\label{eq:I.21}\\
  \uH_Z^*(F)
    &\simeq \uExt^*(\ZZ_{Z,X},F).
    \tag{21 bis}\label{eq:I.21bis}
\end{align}
\end{sgastatement}

\noindent\textit{Editorial note.} This proposition is numbered 2.3 in the
original edition; the corrected branch designates it 2.3 bis.

We may therefore apply the results of \cite{Tohoku} on the $\operatorname{Ext}$
of modules. First, we record the following interpretation of the sheaves
$\uH_Z^*(F)$ in terms of the global groups $\mathrm H_Z^*(X,F)$.

\begin{sgastatement}{Corollary 2.4}\sgareaderlabel{I.2.4}{2.4}
$\uH_Z^*(F)$ is canonically isomorphic to the sheaf associated with the
presheaf
\[
  U\longmapsto
  \mathrm H_{Z\cap U}^*(U,\restr{F}{U}).
\]
\end{sgastatement}

In particular, using \hyperref[I.2.3bof]{Corollary~2.3}, we obtain:

\begin{sgastatement}{Corollary 2.5}\sgareaderlabel{I.2.5}{2.5}
Suppose that $Z$ is open. There is an isomorphism of cohomological functors
\begin{equation}
  \uH_Z^*(F)=\mathrm R^*i_*i^*(F),
  \tag{22}\label{eq:I.22}
\end{equation}
where $i\colon Z\to X$ is the inclusion.
\end{sgastatement}

The $\operatorname{Ext}$ spectral sequence gives the following important
spectral sequence.

\begin{sgastatement}{\hyperref[I.2.6]{Theorem 2.6}}
\sgareaderlabel{I.2.6}{2.6}
There is a spectral sequence, functorial in $F$, abutting to
$\mathrm H_Z^*(X,F)$ and having initial term
\begin{equation}
  E_2^{p,q}(F)=\mathrm H^p(X,\uH_Z^q(F)).
  \tag{23}\label{eq:I.23}
\end{equation}
\end{sgastatement}

\begin{sgaremark}{Remarks 2.7}\sgareaderlabel{I.2.7}{2.7}
It follows immediately from \hyperref[I.2.4]{Corollary~2.4} that the sheaves $\uH_Z^q(F)$
vanish on $X-\overline Z$, and also vanish in the interior of $Z$ when
$q\ne0$; hence, for such $q$, $\uH_Z^q(F)$ is supported on the boundary of
$Z$.

Consequently, the right-hand side of~\eqref{eq:I.23} may be interpreted as a
cohomology group on $\overline Z$. We shall use \hyperref[I.2.6]{Theorem~2.6} when $Z$ is
closed\srcpage{15} in $X$; in that case the right-hand side of
\eqref{eq:I.23} may be interpreted as a cohomology group computed on $Z$:
\begin{equation}
  E_2^{p,q}(F)=\mathrm H^p(Z,\uH_Z^q(F)).
  \tag{23 bis}\label{eq:I.23bis}
\end{equation}
The corresponding equation was numbered \hyperref[eq:I.23]{(23)} in the original edition.
\end{sgaremark}

When $Z$ is open, the spectral sequence of \hyperref[I.2.6]{Theorem~2.6} is simply the Leray
spectral sequence for the continuous map $i\colon Z\to X$, in view of the
interpretation in \hyperref[I.2.5]{Corollary~2.5} used to compute its initial term.

We return to the exact sequence~\hyperref[eq:I.18]{(18)}. It gives rise to an exact
$\operatorname{Ext}$ sequence (cf.\ \cite{Tohoku}).

\noindent\textit{Editorial note.} The original edition cited
Proposition~1.10 at this point; the corrected branch cites equation~(18).

\begin{sgastatement}{Theorem 2.8}\sgareaderlabel{I.2.8}{2.8}
Let $Z$ be a locally closed subset of $X$, let $Z'$ be a closed subset of
$Z$, and put $Z''=Z-Z'$. There is a long exact sequence, functorial in $F$:
\begin{equation}
\begin{gathered}
  0\longrightarrow \mathrm H_{Z'}^0(X,F)
  \longrightarrow \mathrm H_Z^0(X,F)
  \longrightarrow \mathrm H_{Z''}^0(X,F)
  \xrightarrow{\partial} \mathrm H_{Z'}^1(X,F)
  \longrightarrow \mathrm H_Z^1(X,F)\longrightarrow\cdots \\
  \cdots\longrightarrow \mathrm H_{Z'}^i(X,F)
  \longrightarrow \mathrm H_Z^i(X,F)
  \longrightarrow \mathrm H_{Z''}^i(X,F)
  \xrightarrow{\partial} \mathrm H_{Z'}^{i+1}(X,F)
  \longrightarrow\cdots .
\end{gathered}
\tag{24}\label{eq:I.24}
\end{equation}
\end{sgastatement}

Recall how this exact sequence is obtained. Let $C(F)$ be an injective
resolution of $F$. Then the exact sequence~\hyperref[eq:I.18]{(18)}\footnote{See the preceding
editorial note.} gives the exact sequence
\begin{equation}
  0\longrightarrow\Gamma_{Z'}(C(F))
  \longrightarrow\Gamma_Z(C(F))
  \longrightarrow\Gamma_{Z''}(C(F))
  \longrightarrow0.
  \tag{25}\label{eq:I.25}
\end{equation}
This is precisely the sequence defined in \hyperref[I.1.8]{Proposition~1.8}. Its cohomology
long exact sequence is~\eqref{eq:I.24}.

\noindent\textit{Source note.} The corrected French TeX prints
$\Gamma(C(F))$ in the middle term of~(25). Equation~(24), the construction
from (18), and Proposition~1.8 require $\Gamma_Z(C(F))$; the support
subscript is restored here and the emendation is not silent.

The most important case for us is the one in which $Z$ is closed (and we may
always reduce to it by replacing $X$ with an open subset $V$ in which $Z$ is
closed). Then $Z'$ is closed, $Z''$ is closed in the open subset $X-Z'$, and
we may write
\begin{equation}
  \mathrm H_{Z''}^i(X,F)
  =\mathrm H_{Z''}^i(X-Z',\restr{F}{X-Z'}).
  \tag{26}\label{eq:I.26}
\end{equation}
This permits us to write the exact sequence~\eqref{eq:I.24} in terms of
cohomology with support\srcpage{16} in a given closed subset. The most common
case is $Z=X$. Writing $Z'=A$ for simplicity, we obtain:

\begin{sgastatement}{Corollary 2.9}\sgareaderlabel{I.2.9}{2.9}
Let $A$ be a closed subset of $X$. There is a long exact sequence,
functorial in $F$:
\begin{equation}
\begin{gathered}
  0\longrightarrow \mathrm H_A^0(X,F)
  \longrightarrow \mathrm H^0(X,F)
  \longrightarrow \mathrm H^0(X-A,F)
  \xrightarrow{\partial} \mathrm H_A^1(X,F)\longrightarrow\cdots \\
  \cdots\longrightarrow \mathrm H_A^i(X,F)
  \longrightarrow \mathrm H^i(X,F)
  \longrightarrow \mathrm H^i(X-A,F)
  \xrightarrow{\partial} \mathrm H_A^{i+1}(X,F)
  \longrightarrow\cdots .
\end{gathered}
\tag{27}\label{eq:I.27}
\end{equation}
\end{sgastatement}

This exact sequence shows that $\mathrm H_A^i(X,F)$ plays the role of a
relative cohomology group of $X$ modulo $X-A$, with coefficients in $F$. It
is in this capacity that it arises naturally in applications. By
``sheafifying''~\eqref{eq:I.24} and~\eqref{eq:I.27}, or by arguing directly,
and taking into account that the sheaf associated with
$U\mapsto\mathrm H^i(U,F)$ is zero for $i>0$, we obtain:

\begin{sgastatement}{Corollary 2.10}\sgareaderlabel{I.2.10}{2.10}
Under the hypotheses of \hyperref[I.2.8]{Theorem~2.8} there is a long exact sequence,
functorial in $F$:
\begin{equation}
  \cdots\longrightarrow \uH_{Z'}^i(F)
  \longrightarrow \uH_Z^i(F)
  \longrightarrow \uH_{Z''}^i(F)
  \xrightarrow{\partial}\uH_{Z'}^{i+1}(F)
  \longrightarrow\cdots .
  \tag{24 bis}\label{eq:I.24bis}
\end{equation}
\end{sgastatement}

\begin{sgastatement}{Corollary 2.11}\sgareaderlabel{I.2.11}{2.11}
Let $A$ be a closed subset of $X$. There is an exact sequence, functorial in
$F$:
\begin{equation}
  0\longrightarrow\uH_A^0(F)
  \longrightarrow F
  \longrightarrow f_*(\restr{F}{X-A})
  \xrightarrow{\partial}\uH_A^1(F)
  \longrightarrow0,
  \tag{28}\label{eq:I.28}
\end{equation}
and, for $i\ge2$, canonical isomorphisms
\begin{equation}
  \uH_A^i(F)
  =\uH_{X-A}^{i-1}(F)
  =\mathrm R^{i-1}f_*(\restr{F}{X-A}),
  \tag{29}\label{eq:I.29}
\end{equation}
where $f\colon X-A\to X$ is the inclusion.
\end{sgastatement}

Thus $\uH_A^0(F)$ and $\uH_A^1(F)$ are defined, respectively, as the kernel
and cokernel of the canonical homomorphism
\[
  F\longrightarrow f_*f^*(F)=f_*(\restr{F}{X-A}),
\]
and $\uH_A^i(F)$ for $i\ge2$ is expressed in terms of the derived functors
of $f_*$.

\noindent\textit{Source note.} The French display~(28) omits the argument
$(F)$ from $\uH_A^1$. The explanatory sentence immediately following the
display supplies it; it is restored here and recorded as an emendation.

\begin{sgastatement}{Corollary 2.12}\sgareaderlabel{I.2.12}{2.12}
Let $F$\srcpage{17} be a sheaf of abelian groups on $X$. If $F$ is flasque,
then for every locally closed subset $Z$ of $X$ and every integer $i\ne0$,
\[
  \mathrm H_Z^i(X,F)=0,
  \qquad
  \uH_Z^i(F)=0.
\]
Conversely, if $\mathrm H_Z^1(X,F)=0$ for every closed subset $Z$ of $X$,
then $F$ is flasque.
\end{sgastatement}

Suppose that $F$ is flasque. It induces a flasque sheaf on every open
subset. To prove that $\mathrm H_Z^i(X,F)=0$ for $i>0$, we may therefore
assume that $Z$ is closed; the assertion then follows from the exact
sequence~\eqref{eq:I.27}. For every locally closed $Z$, ``sheafifying''---that
is, applying \hyperref[I.2.4]{Corollary~2.4}---gives $\uH_Z^i(F)=0$ for $i>0$. Conversely,
suppose that $\mathrm H_Z^1(X,F)=0$ for every closed $Z$. Then
\eqref{eq:I.27} shows that
$\mathrm H^0(X,F)\to\mathrm H^0(X-Z,F)$ is surjective for every such $Z$,
which means precisely that $F$ is flasque.

\noindent\textit{Editorial note.} At both corresponding occurrences in the
French proof, the original edition cited Corollary~2.9; the corrected branch
specifies equation~(27).

Combining \hyperref[I.2.6]{Theorems~2.6} and~\hyperref[I.2.8]{2.8}, we deduce:

\begin{sgastatement}{Proposition 2.13}\sgareaderlabel{I.2.13}{2.13}
Let $F$ be a sheaf of abelian groups on $X$, let $Z$ be a closed subset of
$X$, put $U=X-Z$, and let $N$ be an integer. The following conditions are
equivalent:
\begin{enumerate}[label=\textup{(\roman*)},leftmargin=2.7em]
  \item $\uH_Z^i(F)=0$ for $i\le N$.
  \item For every open subset $V$ of $X$, the canonical homomorphism
  \[
    \mathrm H^i(V,F)\longrightarrow\mathrm H^i(V\cap U,F)
  \]
  is
  \begin{enumerate}[label=\textup{\alph*)},leftmargin=2.5em]
    \item bijective for $i<N$;
    \item injective for $i=N$.
  \end{enumerate}
\end{enumerate}
If $N>0$, it is enough in \textup{(ii)} to require \textup{a)}.
\end{sgastatement}

To prove \textup{(i)} $\Rightarrow$ \textup{(ii)}, the local nature of the
sheaves $\uH_Z^i(F)$ reduces us to proving the following.

\begin{sgastatement}{Corollary 2.14}\sgareaderlabel{I.2.14}{2.14}
If condition \textup{(i)} of \hyperref[I.2.13]{Proposition~2.13} holds, then
\[
  \mathrm H^i(X,F)\longrightarrow\mathrm H^i(U,F)
\]
is bijective for $i<N$ and injective for $i=N$.
\end{sgastatement}

Indeed, by the exact sequence~\eqref{eq:I.27}, this is equivalent to
$\mathrm H_Z^i(X,F)=0$ for $i\le N$, and this\srcpage{18} relation is an
immediate consequence of the spectral sequence~\eqref{eq:I.23bis}.

\noindent\textit{Editorial note.} At both corresponding occurrences in the
French text, the original edition cited Theorem~2.6; the corrected branch
specifies equation~(23 bis).

Conversely, condition \textup{(ii)} of \hyperref[I.2.13]{Proposition~2.13} says that, for every
open subset $V$ of $X$,
\[
  \mathrm H_{Z\cap V}^i(V,\restr{F}{V})=0
  \qquad\text{for }i\le N,
\]
which implies \textup{(i)} by \hyperref[I.2.4]{Corollary~2.4}. Moreover, if $N>0$, condition
\textup{b)} is superfluous. Indeed, if $N=1$, condition \textup{a)} together
with~\eqref{eq:I.28} gives $\uH_Z^i(F)=0$ for $i\le N$. If $N>1$,
condition \textup{a)} for $i=N-1>0$, together with~\eqref{eq:I.29}, gives
$\uH_Z^i(F)=0$ for $i\le N$.

Taking \hyperref[I.2.11]{Corollary~2.11} into account, this also proves condition \textup{(i)}
of \hyperref[I.2.13]{Proposition~2.13}.

\begin{sgaremark}{Remark}
Let $Y\to X$ be a closed immersion, and suppose that locally it has the form
$\{0\}\times Y\subset R^n\times Y$. Suppose that $F$ is a locally constant
sheaf on $X$. Then
\begin{equation}
  \uH_Y^i(F)\simeq
  \begin{cases}
    0,& i\ne n,\\
    F\otimes\uT_{Y,X},& i=n,
  \end{cases}
  \qquad
  \uT_{Y,X}\simeq\uH_Y^n(\ZZ_X),
  \tag{30}\label{eq:I.30}
\end{equation}
where $\uT_{Y,X}$ is the extension to $X$ of a sheaf on $Y$ locally
isomorphic to $\ZZ_Y$; it is called the \emph{normal orientation sheaf} of
$Y$ in $X$.
\end{sgaremark}

\noindent\textit{Source note.} The corrected French prose prints lowercase
$x$ where the ambient space must be $X$, and its sentence defining
$\uT_{Y,X}$ is grammatically incomplete. The reading above preserves the
display and states its evident mathematical syntax explicitly.

Using the spectral sequence~\eqref{eq:I.23bis}, we obtain in this case
\begin{equation}
  \mathrm H_Y^i(X,F)
  \simeq
  \mathrm H^{i-n}(Y,F\otimes\uT_{Y,X}),
  \tag{31}\label{eq:I.31}
\end{equation}
and recover the \emph{Gysin homomorphism}\sgareaderlabel{hgysin}{13}
\begin{equation}
  \mathrm H^j(Y,F\otimes\uT_{Y,X})
  \longrightarrow
  \mathrm H^{j+n}(X,F).
  \tag{32}\label{eq:I.32}
\end{equation}

\begin{thebibliography}{Godement}
\bibitem[Godement]{Godement}
R.~Godement, \emph{Topologie alg\'ebrique et th\'eorie des faisceaux},
Act. Scient. et Ind., vol.~1252, Hermann, Paris.

\bibitem[T\^ohoku]{Tohoku}
A.~Grothendieck, ``Sur quelques points d'alg\`ebre homologique,''
\emph{T\^ohoku Math. J.} \textbf{9} (1957), 119--221.
\end{thebibliography}

% END INCLUDED active_source/parts/01_expose_I.tex

% BEGIN INCLUDED active_source/parts/02_expose_II.tex
% SGA 2, Expose II — English mathematical body.
% Authority coverage: official corrected arXiv TeX lines 505--708.
% Mechanically extracted without alteration from the source-aligned unit.

\makeatletter
\providecommand{\sgareaderlabel}[2]{%
  \begingroup
  \def\@currentlabel{#2}%
  \phantomsection\label{#1}%
  \endgroup}
\makeatother

\chapter[Quasi-coherent Sheaves on Preschemes]
{Application to Quasi-coherent Sheaves on Preschemes}
\label{II}

\begin{sgastatement}{Proposition 1}\sgareaderlabel{II.1}{1}
Let\srcpage{19} $X$ be a prescheme and let $Z$ be a locally closed subset of
the form $Z=U-V$, where $U$ and $V$ are open subsets of $X$ such that
$V\subset U$ and the canonical immersions $U\to X$ and $V\to X$ are
quasi-compact. Then, for every quasi-coherent $\OX$-module $F$, the sheaves
$\uH_Z^i(F)$ are quasi-coherent.
\end{sgastatement}

By Expos\'e~I, \hyperref[eq:I.24]{equation~(24)}, there is an exact sequence of relative
cohomology sheaves
\[
  \longrightarrow \uH_U^{i-1}(F)
  \longrightarrow \uH_V^i(F)
  \longrightarrow \uH_Z^i(F)
  \longrightarrow \uH_U^i(F)
  \longrightarrow \uH_V^{i+1}(F)
  \longrightarrow .
\]
By EGA~III, 1.4.17, for the sheaves $\uH_Z^i(F)$ to be quasi-coherent it
therefore suffices that the sheaves $\uH_U^i(F)$ and $\uH_V^i(F)$ be
quasi-coherent. We may thus suppose that $Z$ is open and that the canonical
immersion $j\colon Z\to X$ is quasi-compact.

Since $Z$ is open, Expos\'e~I, \hyperref[eq:I.22]{equation~(22)}, gives a canonical isomorphism
\[
  \uH_Z^i(F)\simeq \mathrm R^i j_*(F|_Z).
\]
But $j$ is separated (EGA~I, 5.5.1) and quasi-compact; hence, by EGA~III,
1.4.10, the sheaves
$\mathrm R^i j_*(F|_Z)=\uH_Z^i(F)$ are quasi-coherent. This completes the
proof.

\begin{sgastatement}{Corollary 2}\sgareaderlabel{II.2}{2}
Let $Z$ be a closed subset of $X$ such that the canonical immersion
$X-Z\to X$ is quasi-compact. Then the $\OX$-modules $\uH_Z^i(F)$ are
quasi-coherent.
\end{sgastatement}

\begin{sgastatement}{Corollary 3}\sgareaderlabel{II.3}{3}
If $X$ is locally noetherian, then, for every locally closed subset $Z$ of
$X$ and every quasi-coherent $\OX$-module $F$, the sheaves $\uH_Z^i(F)$ are
quasi-coherent.
\end{sgastatement}

This follows immediately from \hyperref[II.2]{Corollary~2} and EGA~I, 6.6.4.

\begin{sgastatement}{Corollary 4}\sgareaderlabel{II.4}{4}
Suppose\srcpage{20} that $X$ is the spectrum of a ring $A$. Let $U$ be a
quasi-compact open subset of $X$, put $Y=X-U$, and let $F$ be a
quasi-coherent $\OX$-module. There is an isomorphism of cohomological
functors in $F$:
\begin{equation}
  \uH_Y^i(F)=\widetilde{\bigl(\mathrm H_Y^i(X,F)\bigr)}.
  \tag{4.1}\label{eq:II.4.1}
\end{equation}
There is moreover an exact sequence, functorial in $F$,
\begin{equation}
  0\longrightarrow \mathrm H_Y^0(X,F)
  \longrightarrow \mathrm H^0(X,F)
  \longrightarrow \mathrm H^0(U,F)
  \longrightarrow \mathrm H_Y^1(X,F)
  \longrightarrow 0,
  \tag{4.2}\label{eq:II.4.2}
\end{equation}
and there are isomorphisms, functorial in $F$,
\begin{equation}
  \mathrm H_Y^i(X,F)\simeq \mathrm H^{i-1}(U,F),
  \qquad i\geq 2.
  \tag{4.3}\label{eq:II.4.3}
\end{equation}
\end{sgastatement}

By \hyperref[II.2]{Corollary~2}, the sheaves $\uH_Y^i(F)$ are quasi-coherent. Since $X$ is
affine, it follows that
$\mathrm H^p(X,\uH_Y^i(F))=0$ for $p>0$. The spectral sequence of
Expos\'e~I, \hyperref[eq:I.23]{equation~(23)}, therefore degenerates, and
\[
  \mathrm H_Y^i(X,F)=\Gamma(\uH_Y^i(F)).
\]
Equation~\eqref{eq:II.4.1} now follows from EGA~I, 1.1.3.7; equations
\eqref{eq:II.4.2} and \eqref{eq:II.4.3} follow from the cohomology exact
sequence of Expos\'e~I, \hyperref[eq:I.27]{equation~(27)}, and from the vanishing
$\mathrm H^i(X,F)=0$ for $i>0$, since $X$ is affine.

\noindent\textit{Editorial note.} The corrected branch replaces the
original citation of Proposition~1 by Corollary~2 in the first sentence of
this proof and regularizes both Expos\'e~I equation references; those
corrected references are used above.

Under the hypotheses of \hyperref[II.4]{Corollary~4},\footnote{The corrected French source
adds: ``For consistency and clarity, only equations have been numbered in
parentheses.''} $U$ is a finite union of affine open subsets $X_f$.
Consequently, one can find an ideal $I$, generated by finitely many elements
$f_\alpha$ and defining $Y$; write $\ffam=(f_\alpha)$. With the notation of
EGA~III, 1,\footnote{Recall that $\mathrm H^\bullet(\ffam,M)$ is the Koszul
cohomology $\mathrm H^\bullet(\operatorname{Hom}(K_\bullet(\ffam),M))$ of
$\ffam$ with values in $M$ (EGA~III, 1.1.2), and that
$\mathrm H^\bullet((\ffam),M)$ is the direct limit (ibid., 1.1.6.5)
$\varinjlim_n\mathrm H^\bullet(\ffam^n,M)$, whose transition morphisms are
induced by the natural morphisms
$K_\bullet(\ffam^{n+1})\to K_\bullet(\ffam^n)$ (ibid., 1.1.6).} we have the
following result.

\begin{sgastatement}{Proposition 5}\sgareaderlabel{II.5}{5}
Suppose that $X$ is the spectrum of a ring $A$. Let
$\ffam=(f_\alpha)$ be a finite family of elements of $A$, let $Y$ be the
closed subset of $X$ that they define, let $M$ be an $A$-module, and let $F$
be the associated sheaf. There are isomorphisms of $\partial$-functors in
$M$:
\begin{equation}
  \mathrm H^i((\ffam),M)\simeq \mathrm H_Y^i(X,F).
  \tag{5.1}\label{eq:II.5.1}
\end{equation}
\end{sgastatement}

We shall also write
$\mathrm H_J^i(M)=\mathrm H_Y^i(X,F)$ when $Y$ is the closed subset of
$X=\Spec A$ defined by an ideal $J$ of $A$.

For $i=0$\srcpage{21} and $i=1$, use the exact sequences
\eqref{eq:II.4.2} and EGA~III, 1.4.3.2; for $i\geq2$, use
\eqref{eq:II.4.3} and EGA~III, 1.4.3.1. This gives isomorphisms functorial in
$M$. One verifies that, up to a sign depending only on $i$, they are
compatible with the boundary operator; hence the isomorphism of
$\partial$-functors in \eqref{eq:II.5.1} exists.

\noindent\textit{Editorial note.} For the case $i\geq2$, the original
edition cited EGA~III, 1.4.3.2; the corrected branch cites 1.4.3.1.

Now let $X$ be a prescheme, let $Y$ be a closed subset of $X$, and let
$f\colon Y\to X$ be the inclusion. Let $I$ be a quasi-coherent ideal defining
$Y$ in $X$, and let $F$ be a sheaf on $X$.

We have seen that there are isomorphisms of $\partial$-functors in $F$
\begin{align}
  \Ext_{\OX}^i(X;f_*f^{-1}(\OX),F)
    &\longrightarrow \mathrm H_Y^i(X,F),
    \tag{$*$}\label{eq:II.5.*}\\
  \uExt_{\OX}^i(f_*f^{-1}(\OX),F)
    &\longrightarrow \uH_Y^i(F).
    \tag{$**$}\label{eq:II.5.**}
\end{align}

\noindent\textit{Source note.} The corrected French source calls these
arrows isomorphisms but displays plain right arrows in both $(*)$ and $(**)$.
The arrows have been preserved exactly; no unmarked replacement by
$\xrightarrow{\sim}$ has been made.

Let $n,m$ be integers such that $m\geq n\geq0$. Denote by $i_{n,m}$ the
canonical map
\[
  \mathcal O_{Y_m}=\OX/I^{m+1}
  \longrightarrow \OX/I^{n+1}=\mathcal O_{Y_n},
\]
and by $j_n$ the map
$f_*f^{-1}(\OX)\to\mathcal O_{Y_n}$. The pairs
$(\mathcal O_{Y_n},i_{n,m})$ form a projective system, and the maps $j_n$ are
compatible with the $i_{n,m}$.

Applying the functor $\Ext_{\OX}^i(X;\mathord\cdot,F)$ yields a morphism
\[
  \varphi'\colon
  \varinjlim_n\Ext_{\OX}^i(X;\mathcal O_{Y_n},F)
  \longrightarrow
  \Ext_{\OX}^i(X;f_*f^{-1}(\OX),F).
\]
It is readily seen to be a morphism of cohomological functors in $F$.
Consequently, the morphism
\[
  \varphi\colon
  \varinjlim_n\Ext_{\OX}^i(X;\mathcal O_{Y_n},F)
  \longrightarrow \mathrm H_Y^i(X,F),
\]
obtained by composing $\varphi'$ with~\eqref{eq:II.5.*}, is itself a
morphism of cohomological functors in $F$.

We define\srcpage{22} similarly
\[
  \underline{\varphi}\colon
  \varinjlim_n\uExt_{\OX}^i(\mathcal O_{Y_n},F)
  \longrightarrow \uH_Y^i(F).
\]

Our goal is the following theorem.

\begin{sgastatement}{Theorem 6}\sgareaderlabel{II.6}{6}
\textup{a)} Let $X$ be a locally noetherian prescheme, let $Y$ be a closed
subset of $X$ defined by a coherent ideal $I$, and let $F$ be a
quasi-coherent $\OX$-module. Then $\boldsymbol{\varphi}$ is an isomorphism.

\textup{b)} If $X$ is noetherian, then $\varphi$ is an isomorphism.
\end{sgastatement}

\noindent\textit{Source note.} The corrected source defines the sheaf-level
morphism immediately above as $\underline{\varphi}$ but denotes it by the
bold symbol $\boldsymbol{\varphi}$ in Theorem~6(a) and Lemma~7. Both graphic
forms are retained rather than silently normalized.

\hyperref[II.6]{Theorem~6} will follow from \hyperref[II.6]{Theorem~6}(a) and the following lemma.

\begin{sgastatement}{Lemma 7}\sgareaderlabel{II.7}{7}
If the underlying topological space of $X$ is noetherian and
$\boldsymbol{\varphi}$ is an isomorphism, then $\varphi$ is also an
isomorphism.
\end{sgastatement}

We first prove \hyperref[II.7]{Lemma~7}. There is a spectral sequence
\begin{equation}
  \mathrm H^p(X,\uH_Y^q(F))
  \Longrightarrow \mathrm H_Y^*(X,F).
  \tag{7.1}\label{eq:II.7.1}
\end{equation}
There is, on the other hand, a direct system of spectral sequences
\begin{equation}
  \mathrm H^p\bigl(X,\uExt_{\OX}^q(\mathcal O_{Y_n},F)\bigr)
  \Longrightarrow \Ext_{\OX}^*(X;\mathcal O_{Y_n},F).
  \tag{7.2$_n$}\label{eq:II.7.2n}
\end{equation}

\noindent\textit{Editorial note.} In both abutments above, the original
edition printed a superscript $m$; the corrected branch replaces it by the
superscript $*$, which is preserved here.

It follows from the definitions of $\varphi$ and $\underline{\varphi}$ that
these morphisms are associated with a homomorphism $\Phi$ of spectral
sequences, from the direct limit of~\eqref{eq:II.7.2n} to
\eqref{eq:II.7.1}. If the underlying space of $X$ is noetherian, then by
Godement, 4.12.1,\footnote{Cf. the first bibliographic reference at the end
of Expos\'e~I.}
\[
  \varinjlim_n
    \mathrm H^p\bigl(X,\uExt_{\OX}^q(\mathcal O_{Y_n},F)\bigr)
  \xrightarrow{\sim}
  \mathrm H^p\!\left(
    X,\varinjlim_n\uExt_{\OX}^q(\mathcal O_{Y_n},F)\right).
\]
Thus $\Phi_2$ may be written as a morphism
\[
  \mathrm H^p\!\left(
    X,\varinjlim_n\uExt_{\OX}^q(\mathcal O_{Y_n},F)\right)
  \longrightarrow \mathrm H^p(X,\uH_Y^q(F)),
\]
which is precisely the morphism deduced from $\underline{\varphi}$.

If $\underline{\varphi}$\srcpage{23} is an isomorphism, then so is
$\Phi_2$, and hence so is $\varphi$ by EGA~$0_{\mathrm{III}}$, 11.1.5.
This proves \hyperref[II.7]{Lemma~7}.

We now prove \hyperref[II.6]{Theorem~6}(a). This is a local question on $X$. By \hyperref[II.4]{Corollary~4}
and EGA~I, 1.3.9 and 1.3.12, we may suppose that $X$ is the spectrum of a
ring $A$. Write $M$ for the $A$-module whose associated sheaf is $F$. It is
therefore enough to prove that, under the hypotheses of \hyperref[II.6]{Theorem~6}(a), the
canonical homomorphism
\begin{equation}
  \varinjlim_n\Ext_A^i(A/I^n,M)
  \longrightarrow \mathrm H_Y^i(X,M)
  \tag{7.3}\label{eq:II.7.3}
\end{equation}
is an isomorphism.

\noindent\textit{Source note.} The corrected French text first uses $M$ in
display~(7.3) without introducing it. The affine correspondence with the
previously given quasi-coherent sheaf $F$ is made explicit above.

Let $f_\alpha$ be finitely many elements of $A$ that generate $I$, and put
$\ffam=(f_\alpha)$. The sequence of ideals $(\ffam^n)$ is decreasing and
cofinal with the sequence of ideals $I^n$, so that~\eqref{eq:II.7.3} is
equivalent to a morphism of $\partial$-functors in $M$:
\begin{equation}
  \varinjlim_n\Ext_A^i(A/(\ffam^n),M)
  \longrightarrow \mathrm H_Y^i(X,M).
  \tag{7.4}\label{eq:II.7.4}
\end{equation}
On the other hand, there are canonical isomorphisms
\begin{equation}
\begin{split}
  \varinjlim_n\operatorname{Hom}_A(A/(\ffam^n),M)
  &\simeq \varinjlim_n
    \bigl(m\in M\mid(\ffam^n)m=0\bigr) \\
  &\simeq \mathrm H^0((\ffam),M).
\end{split}
  \tag{7.5}\label{eq:II.7.5}
\end{equation}
Since $\varinjlim_n\Ext_A^i(A/(\ffam^n),M)$ is a universal
$\partial$-functor in $M$, there is a unique morphism of $\partial$-functors
in $M$
\begin{equation}
  \varinjlim_n\Ext_A^i(A/(\ffam^n),M)
  \longrightarrow \mathrm H^i((\ffam),M)
  \tag{7.6}\label{eq:II.7.6}
\end{equation}
that agrees in degree zero with~\eqref{eq:II.7.5}.

The comparison obtained from~\eqref{eq:II.7.3} and
\eqref{eq:II.5.1} is a morphism of $\partial$-functors in $M$ that agrees
with~\eqref{eq:II.7.6} in degree zero; it therefore agrees with
\eqref{eq:II.7.6} in every degree. \hyperref[II.6]{Theorem~6}(a) is consequently an immediate
consequence of the following lemma.

\noindent\textit{Source note.} The corrected French sentence calls the
comparison just used ``the composite of (7.3) and (5.1).'' As the displayed
arrows in (7.3) and (5.1) have the same target rather than composable source
and target, the intended comparison requires identifying that target through
the isomorphism (5.1), equivalently using its inverse. The mathematical
comparison is stated explicitly above; the source's directional anomaly is
recorded here rather than silently concealed.

\begin{sgastatement}{Lemma 8}\sgareaderlabel{II.8}{8}
Let $A$ be a noetherian ring, let $I$ be an ideal generated by a finite
family $\ffam=(f_\alpha)$ of elements, and let $M$ be an $A$-module. Then the
homomorphisms~\eqref{eq:II.7.6} are isomorphisms.
\end{sgastatement}

\begin{sgastatement}{Lemma 9}\sgareaderlabel{II.9}{9}
Let $A$\srcpage{24} be a ring, let $\ffam=(f_\alpha)$ be a finite family of
elements of $A$, let $I$ be the ideal generated by $\ffam$, and let $i>0$ be
an integer. The following conditions are equivalent:
\begin{enumerate}[label=\textup{\alph*)},leftmargin=2.5em]
  \item The homomorphism~\eqref{eq:II.7.6} is an isomorphism for every $M$.
  \item $\mathrm H^i((\ffam),M)=0$ for every injective $M$.
  \item The projective system
  $(\mathrm H_i(\ffam^n,A))=\mathrm H_{i,n}$ is essentially zero; that is,
  for every $n$ there exists $n'>n$ such that
  $\mathrm H_{i,n'}\to\mathrm H_{i,n}$ is zero.
\end{enumerate}
\end{sgastatement}

\noindent\textit{Source note.} Lemma~9(c), and later Lemma~11, print their
systems with the index $n\in\ZZ$, whereas the proofs explicitly take $n>0$
and use ordinary positive powers. The printed indexing is retained; this
translation does not introduce negative powers.

Condition a) trivially implies b).

Condition b) implies a). Indeed, b) implies that
$M\mapsto\mathrm H^i((\ffam),M)$ is a universal cohomological functor, so
\eqref{eq:II.7.6} is a morphism of universal cohomological functors. It is an
isomorphism in degree zero and hence in every degree.

Condition c) implies b). Indeed, if $M$ is injective, then for every $n$
\[
  \mathrm H^i(\ffam^n,M)
  =\operatorname{Hom}(\mathrm H_i(\ffam^n,A),M)
  =\operatorname{Hom}(\mathrm H_{i,n},M).
\]
Thus c) implies that, for every $i$, the direct system
$(\mathrm H^i(\ffam^n,M))_{n\in\ZZ}$ is essentially zero, whence b).

Condition b) implies c). Indeed, fix $n>0$, and let $j$ be a monomorphism
from $\mathrm H_{i,n}$ into an injective module $M$. For $n'\geq n$, let
$j_{n'}\in\operatorname{Hom}(\mathrm H_{i,n'},M)$ be the composite of $j$
with the transition homomorphism
$t_{n',n}\colon\mathrm H_{i,n'}\to\mathrm H_{i,n}$. The $j_{n'}$ define an
element of $\mathrm H^i((\ffam),M)$, which is zero by hypothesis. There is
therefore an $n_0$ such that $j_{n'}=0$ if $n'>n_0$. Since $j$ is a
monomorphism, $j_{n'}=0$ implies $t_{n',n}=0$, which proves c).

\noindent\textit{Editorial note.} In the last implication, the original
edition shortened the transition-map symbol to $t_{n'}$; the corrected
branch restores $t_{n',n}$, as retained above.

\begin{sgastatement}{Corollary 10}\sgareaderlabel{II.10}{10}
Suppose that the underlying space of $X=\Spec(A)$ is noetherian. The
preceding conditions hold for every finite family of elements of $A$ and
every $i>0$ (or, equivalently, for $i=1$) if and only if, for every injective
$A$-module $M$, the sheaf $F$ associated with $M$ is flasque.
\end{sgastatement}

The condition is necessary. Indeed, let $\ffam=(f_\alpha)$ be a finite family
of elements of $A$, let $Y$ be the closed subset defined by $\ffam$, and put
$U=X-Y$. We then have the exact sequence
\[
  \mathrm H^0(X,F)\longrightarrow \mathrm H^0(U,F)
  \longrightarrow \mathrm H^1((\ffam),M)\longrightarrow0,
\]
and \hyperref[II.9]{Lemma~9}(b) shows that
$\mathrm H^0(X,F)\to\mathrm H^0(U,F)$ is surjective.

The condition is sufficient\srcpage{25} by~\eqref{eq:II.5.1} and the fact
that, for every closed subset $Y$ of $X$ and every flasque sheaf $F$ on $X$,
$\mathrm H_Y^i(X,F)=0$ for $i>0$.

\begin{sgastatement}{Lemma 11}\sgareaderlabel{II.11}{11}
Under the hypotheses of \hyperref[II.9]{Lemma~9}, for every noetherian $A$-module $N$ and
every $i>0$, the projective system
$(\mathrm H_{i,n}(N))_{n\in\ZZ}$, where
$\mathrm H_{i,n}(N)=\mathrm H_i(\ffam^n,N)$, is essentially zero.
\end{sgastatement}

\begin{sgaproof}
We argue by induction on the number $m$ of elements of $\ffam$.

If $m=1$, the family $\ffam$ consists of a single element, say $f$.
Then $\mathrm H_{i,n}(N)=0$ for $i>1$, and $\mathrm H_{1,n}(N)$ is
canonically isomorphic to the annihilator $N(n)$ of $f^n$ in $N$. For
$n'\geq n$, the transition homomorphism $N(n')\to N(n)$ is multiplication
by $f^{n'-n}$. The $N(n)$ form an increasing sequence of submodules of $N$.
Since $N$ is noetherian, there is an $n_0$ such that
$N(n)=N(n_0)$ for $n\geq n_0$. Hence all the $N(n)$ are annihilated by
$f^{n_0}$, and every transition homomorphism $N(n')\to N(n)$ is zero when
$n'\geq n+n_0$. The lemma is thus proved for $m=1$.

Now suppose that $m>1$ and that the lemma has been proved for all integers
$m'<m$. Put
$\gfam=(f_1,\ldots,f_{m-1})$ and $\hfam=f_m$.

For every $n>0$, EGA~III, 1.1.4.1 gives an exact sequence
\[
  0\longrightarrow
  \mathrm H_0(\hfam^n,\mathrm H_i(\gfam^n,N))
  \longrightarrow \mathrm H_i(\ffam^n,N)
  \longrightarrow
  \mathrm H_1(\hfam^n,\mathrm H_{i-1}(\gfam^n,N))
  \longrightarrow0,
\]
and, as $n$ varies, a projective system of exact sequences. The induction
hypothesis shows that, for $i>0$, the modules
$\mathrm H_i(\gfam^n,N)$ form an essentially zero projective system. The same
is therefore true of
$\mathrm H_0(\hfam^n,\mathrm H_i(\gfam^n,N))$, which is identified with a
quotient of $\mathrm H_i(\gfam^n,N)$. For the terms on the right, factor the
transition morphisms from $n'$ to $n$ as
\[
\begin{split}
  \mathrm H_1(\hfam^{n'},\mathrm H_{i-1}(\gfam^{n'},N))
  &\longrightarrow
  \mathrm H_1(\hfam^{n'},\mathrm H_{i-1}(\gfam^n,N)) \\
  &\longrightarrow
  \mathrm H_1(\hfam^n,\mathrm H_{i-1}(\gfam^n,N)).
\end{split}
\]
Since $\mathrm H_{i-1}(\gfam^n,N)$ is a noetherian module, the case $m=1$
shows that, for fixed $n$, there exists $n'>n$ such that the second arrow is
zero. One therefore sees\srcpage{26} that, in this projective system of exact
sequences, the two outer projective systems are essentially zero. The middle
projective system is therefore essentially zero as well.
\end{sgaproof}

We have thus proved \hyperref[II.11]{Lemma~11}, hence \hyperref[II.8]{Lemma~8}, and consequently \hyperref[II.6]{Theorem~6}.

\begin{sgaremark}{Remark}
\hyperref[II.6]{Theorem~6} can also be obtained by proving the condition of \hyperref[II.10]{Corollary~10} with
the aid of the structure theorems for injective modules over a noetherian
ring (Matlis, Gabriel).
\end{sgaremark}

% END INCLUDED active_source/parts/02_expose_II.tex

% BEGIN INCLUDED active_source/parts/03_expose_III.tex
% SOURCE COMPONENT: 01_section_1.tex
\makeatletter
\providecommand{\sgareaderlabel}[2]{%
  \begingroup
  \def\@currentlabel{#2}%
  \phantomsection\label{#1}%
  \endgroup}
\makeatother

\chapter{Cohomological Invariants and Depth}
\label{III}

\section*{1. Review}
\sgareaderlabel{III.1}{1}

We\srcpage{27} shall state a few definitions and results which the reader may
find, for example, in Chapter~I of the course given by J.-P. Serre at the
Collège de France in 1957--58.\footnote{\sgareaderlabel{noteserre}{*}\textit{Editorial note.} The
reissue of Serre's text (J.-P. Serre, \emph{Algèbre locale.
Multiplicités}, course at the Collège de France, 1957--58, written up by
Pierre Gabriel, second edition, Lecture Notes in Mathematics, vol.~11,
Springer-Verlag, 1965) no longer contains proofs of these statements. One may
consult N.~Bourbaki, \emph{Algèbre commutative}, Masson, as Serre himself
suggests.}

\begin{sgastatement}{Definition 1.1}\sgareaderlabel{def:III.1.1}{1.1}
Let $A$ be a ring (commutative with identity, as throughout what follows), and
let $M$ be an $A$-module (unital, as throughout what follows). We call:
\begin{itemize}[leftmargin=2.2em]
  \item the \emph{annihilator of $M$}, denoted $\Ann M$, the set of
  $a\in A$ such that $am=0$ for every $m\in M$;
  \item the \emph{support of $M$}, denoted $\Supp M$, the set of prime
  ideals $\mathfrak p$ of $A$ such that the localization $M_{\mathfrak p}$
  is nonzero;
  \item the “\emph{assassin of $M$},” or the “\emph{set of prime ideals
  associated with $M$},” denoted $\Ass M$, the set of prime ideals
  $\mathfrak p$ for which there exists a nonzero element of $M$ whose
  annihilator is $\mathfrak p$.
\end{itemize}
If $\mathfrak a$ is an ideal of $A$, we denote by $\mathfrak r(\mathfrak a)$
the radical of $\mathfrak a$ in $A$, that is, the set of elements of $A$ some
power of which lies in $\mathfrak a$.
\end{sgastatement}

The following results hold under the assumption that \emph{$A$ is noetherian
and $M$ is finitely generated}.

\begin{sgastatement}{Proposition 1.1}\sgareaderlabel{III.1.1}{1.1}
\begin{enumerate}[label=\textup{(\roman*)},leftmargin=2.5em]
  \item $\Ass M$ is a finite set.
  \item An element of $A$ annihilates a nonzero element of $M$ if and only
  if\srcpage{28} it belongs to one of the prime ideals associated with $M$.
  \item The radical of the annihilator of $M$,
  $\mathfrak r(\Ann M)$, is the intersection of the prime ideals associated
  with $M$ that are minimal with respect to inclusion in $\Ass M$.
\end{enumerate}
\end{sgastatement}

\begin{sgastatement}{Proposition 1.2}\sgareaderlabel{III.1.2}{1.2}
Let $\mathfrak p$ be a prime ideal of $A$. The following assertions are
equivalent:
\begin{enumerate}[label=\textup{(\roman*)},leftmargin=2.8em]
  \item $\mathfrak p\in\Supp M$.
  \item There exists $\mathfrak q\in\Ass M$ such that
  $\mathfrak q\subset\mathfrak p$.
  \item $\mathfrak p\supset\Ann M$.
  \item[\textup{(iii bis)}] $\mathfrak p\supset\mathfrak r(\Ann M)$.
\end{enumerate}
\end{sgastatement}

\begin{sgastatement}{Proposition 1.3}\sgareaderlabel{III.1.3}{1.3}
Let $N$ be a finitely generated $A$-module. Then
\[
  \Ass\operatorname{Hom}_A(N,M)=\Supp N\cap\Ass M.
\]
\end{sgastatement}



% SOURCE COMPONENT: 02_section_2.tex
\section*{2. Depth}
\sgareaderlabel{III.2}{2}

Throughout\srcpage{28} this section, $A$ denotes a commutative ring, $I$ an ideal of
$A$, and $M$ and $N$ two $A$-modules. We denote by $X$ the prime spectrum of
$A$ (its structure sheaf will not be used in this section), and by $Y$ the
closed set defined by $I$, namely
$Y=\Supp(A/I)=\{\mathfrak p\in X\mid \mathfrak p\supset I\}$.

\begin{sgastatement}{Lemma 2.1}\sgareaderlabel{III.2.1}{2.1}
Suppose that $A$ is noetherian, that $M$ and $N$ are finitely generated, and
that $\Supp N=Y$. Then the following assertions are equivalent:
\begin{enumerate}[label=\textup{(\roman*)},leftmargin=2.6em]
\item $\Homop_A(N,M)=0$.
\item $\Supp N\cap\Ass M=\varnothing$.
\item The ideal $I$ is not a zero divisor on $M$, meaning that, for every
$m\in M$, the equality $Im=0$ implies $m=0$.
\item There exists an $M$-regular element in $I$. (An element $a\in A$ is
called $M$-regular if multiplication by $a$ on $M$ is injective.)
\item For every $\mathfrak p\in Y$, the maximal ideal
$\mathfrak pA_{\mathfrak p}$ of the local ring $A_{\mathfrak p}$ is not
associated with $M_{\mathfrak p}$; in symbols,
$\mathfrak pA_{\mathfrak p}\notin\Ass M_{\mathfrak p}$.
\end{enumerate}
\end{sgastatement}

\begin{proof}
The equivalence (i)$\Leftrightarrow$(ii) follows because
$\Ass\Homop_A(N,M)=\varnothing$ is equivalent to (ii) by \hyperref[III.1.3]{Proposition 1.3},
and to (i) by an immediate consequence of \hyperref[III.1.2]{Proposition 1.2}.

The implication\srcpage{29} (iii)$\Rightarrow$(ii) follows by contradiction. If
$\mathfrak p\in\Supp N\cap\Ass M$, then $\mathfrak p\supset I$ and there
is an $m\in M$ whose annihilator is $\mathfrak p$; hence
$Im\subset\mathfrak p m=0$, contrary to (iii).

The implication (iv)$\Rightarrow$(iii) is immediate.

Conditions (ii) and (iv) are equivalent because $\Supp N=Y$, so (ii) says
that $I$ is contained in no prime ideal associated with $M$. Since those
associated ideals are prime and finite in number, this is equivalent to
saying that $I$ is not contained in their union. By \hyperref[III.1.1]{Proposition 1.1}(ii),
that union is precisely the set of elements of $A$ which are not
$M$-regular.

To prove (i)$\Rightarrow$(v), let $\mathfrak p\in Y$. From
$\Homop_A(N,M)=0$ and the formula
\[
  \bigl(\Homop_A(N,M)\bigr)_{\mathfrak p}
  =\Homop_{A_{\mathfrak p}}(N_{\mathfrak p},M_{\mathfrak p})
\]
one obtains $\Homop_{A_{\mathfrak p}}(N_{\mathfrak p},M_{\mathfrak p})=0$.
Thus, by \hyperref[III.1.3]{Proposition 1.3},
\[
  \Supp N_{\mathfrak p}\cap\Ass M_{\mathfrak p}=\varnothing.
\]
But $\mathfrak pA_{\mathfrak p}\in\Supp N_{\mathfrak p}$, so
$\mathfrak pA_{\mathfrak p}\notin\Ass M_{\mathfrak p}$.

Finally, suppose (v), and let $\mathfrak p\in\Ass M$. There is an
$m\in M$ with annihilator $\mathfrak p$. Its canonical image in
$M_{\mathfrak p}$ is nonzero, and its annihilator contains
$\mathfrak pA_{\mathfrak p}$ and hence equals it. Thus
$\mathfrak pA_{\mathfrak p}\in\Ass M_{\mathfrak p}$, so (v) gives
$\mathfrak p\notin Y$. This proves (i).
\end{proof}

We shall now work with these conditions after replacing the
functor $\Homop$ by its derived functors.

\begin{sgastatement}{Theorem 2.2}\sgareaderlabel{III.2.2}{2.2}\srcpage{30}
Let $A$ be a commutative ring, $I$ an ideal of $A$, $M$ an $A$-module, and
$n$ an integer.
\begin{enumerate}[label=\textup{(\alph*)},leftmargin=2.5em]
\item If there is a sequence $f_1,\ldots,f_{n+1}$ of elements of $I$ which
is $M$-regular (that is, $f_1$ is $M$-regular and $f_{i+1}$ is regular on
$M/(f_1,\ldots,f_i)M$ for $i\le n$), then, for every $A$-module $N$
annihilated by a power of $I$,
\[
  \Extop_A^i(N,M)=0\qquad\text{for }i\le n.
\]
\item If, in addition, $A$ is noetherian, $M$ is finitely generated, and
there is a finitely generated $A$-module $N$ with $\Supp N=V(I)$ and
$\Extop_A^i(N,M)=0$ for $i\le n$, then there exists an $M$-regular
sequence $f_1,\ldots,f_{n+1}$ of elements of $I$.
\end{enumerate}
\end{sgastatement}

We first prove (a), by induction. If $n<0$, the assertion is empty. Suppose
$n\ge0$ and that (a) has been proved for every $n'<n$. By hypothesis there
is an $M$-regular element $f_1\in I$. Denote by $f_1^i$ multiplication by
$f_1$ on $\Extop_A^i(N,M)$ and by $f_1^M$ multiplication by $f_1$ on $M$.
The sequence
\begin{equation}\label{eq:III.2.1}
  0\longrightarrow M\xrightarrow{f_1^M}M\longrightarrow M/f_1M
  \longrightarrow0
\end{equation}
is exact, and therefore so is
\[
  \Extop_A^{i-1}(N,M/f_1M)\xrightarrow{\delta}\Extop_A^i(N,M)
  \xrightarrow{f_1^i}\Extop_A^i(N,M).
\]
By hypothesis, $I^rN=0$ for some $r$, so $f_1^0$ is nilpotent. Since
$\Extop^i$ is a universal cohomological functor, the same is true of
$f_1^i$ for every $i$. Moreover, $M/f_1M$ has a regular sequence of length
$n$. The induction hypothesis therefore gives
\[
  \Extop_A^{i-1}(N,M/f_1M)=0\qquad\text{for }i\le n.
\]
\footnote{\textit{Source correction.} The French authority prints
$\mathrm{Ext}_A^{i-1}(N,M)=0$ for $i\le n-1$. The quotient module and the
bound are both defective: applying the induction hypothesis to the length
$n$ sequence on $M/f_1M$ gives the displayed formula, which is exactly what
is needed below.}
Consequently, for $i\le n$, $f_1^i$ is both injective and nilpotent, so
$\Extop_A^i(N,M)=0$.

We\srcpage{31} now prove (b), again by induction. If $n<0$, the assertion is
empty. If $n=0$, it follows from \hyperref[III.2.1]{Lemma 2.1}, (i)$\Rightarrow$(iv). If $n>0$,
the case $n=0$ gives an $M$-regular element $f_1\in I$. From
\eqref{eq:III.2.1} one obtains the exact sequence
\begin{equation}\label{eq:III.2.2}
  \Extop_A^{i-1}(N,M)\longrightarrow
  \Extop_A^{i-1}(N,M/f_1M)\longrightarrow\Extop_A^i(N,M).
\end{equation}
The hypotheses of (b) therefore hold for $M/f_1M$ and $n-1$. By induction
there is a regular sequence of $n$ elements of $I$ on $M/f_1M$, and hence
an $M$-regular sequence of $n+1$ elements of $I$ beginning with $f_1$.

This theorem suggests the following generalization of the classical
definition of the depth of a finitely generated module over a noetherian
ring.

\begin{sgastatement}{Definition 2.3}\sgareaderlabel{III.2.3}{2.3}
Let $A$ be a commutative ring with identity, $M$ an $A$-module, and $I$ an
ideal of $A$. The \emph{$I$-depth of $M$}, denoted $\prof_I M$, is the
supremum in $\mathbb N\cup\{+\infty\}$ of the natural numbers $n$ such
that, for every finitely generated $A$-module $N$ annihilated by a power of
$I$,
\[
  \Extop_A^i(N,M)=0\qquad\text{for every }i<n.
\]
\end{sgastatement}

It follows from the preceding theorem that, if $n$ is the supremum of the
lengths of $M$-regular sequences of elements of $I$, then
$n\le\prof_I M$. More precisely:

\begin{sgastatement}{Proposition 2.4}\sgareaderlabel{III.2.4}{2.4}
Let $A$ be a commutative ring, $I$ an ideal of $A$, $M$ an $A$-module,
and $n\in\mathbb N$. Consider the assertions:
\begin{enumerate}[label=\textup{(\arabic*)},leftmargin=2.7em]
\item $n\le\prof_I M$.
\item For every finitely generated $A$-module $N$ annihilated by a power
of $I$, one has
\[
  \Extop_A^i(N,M)=0\qquad\text{for }i<n.
\]
\item There\srcpage{32} exists a finitely generated $A$-module $N$ such that
$\Supp N=V(I)$, and such that $\Extop_A^i(N,M)=0$ whenever $i<n$.
\item There is an $M$-regular sequence of length $n$ consisting of
elements of $I$.
\end{enumerate}
The following logical implications hold:
\[
\begin{array}{ccccc}
\hyperref[III.2.4]{(1)}&\Longleftrightarrow&\hyperref[III.2.4]{(2)}&\Longleftarrow&\hyperref[III.2.4]{(4)}\\[-0.2em]
&&\Big\Downarrow&&\\[-0.2em]
&&\hyperref[III.2.4]{(3)}&&
\end{array}
\]
If $A$ is noetherian and $M$ is finitely generated, all four conditions are
equivalent.
\end{sgastatement}

\begin{proof}
\hyperref[III.2.4]{Conditions (1)} and \hyperref[III.2.4]{(2)} are equivalent by definition, and
\hyperref[III.2.4]{(2)}$\Rightarrow$\hyperref[III.2.4]{(3)} follows by taking $N=A/I$. The implication
\hyperref[III.2.4]{(4)}$\Rightarrow$\hyperref[III.2.4]{(2)} is \hyperref[III.2.2]{Theorem 2.2}(a). Finally, when $A$ is noetherian and
$M$ is finitely generated, \hyperref[III.2.4]{(3)}$\Rightarrow$\hyperref[III.2.4]{(4)} is \hyperref[III.2.2]{Theorem 2.2}(b).
\end{proof}

From now until the end of the section, we assume that $A$ is
noetherian and $M$ is finitely generated.

\begin{sgastatement}{Corollary 2.5}\sgareaderlabel{III.2.5}{2.5}
If $f\in I$ is $M$-regular, then
\[
  \prof_I M=\prof_I(M/fM)+1.
\]
\end{sgastatement}

Indeed, if $n\le\prof_I(M/fM)$, there is an $(M/fM)$-regular sequence
$f_1,\ldots,f_n$ of elements of $I$. Hence $f,f_1,\ldots,f_n$ is
$M$-regular, so $n+1\le\prof_I M$ and
$\prof_I M\ge\prof_I(M/fM)+1$. Conversely, for every finitely generated
$A$-module $N$ annihilated by a power of $I$, the exact sequence
  \eqref{eq:III.2.2} shows that if $i<\prof_I M$, then
  $\Extop_A^{i-1}(N,M/fM)=0$.\footnote{\textit{Source correction.}
  The French authority prints $i\le\prof_I M$. In the displayed exact
  sequence, vanishing of both adjacent terms is guaranteed for
  $i<\prof_I M$; this is precisely the range needed for the following depth
  inequality.} Thus
$\prof_I M-1\le\prof_I(M/fM)$.

\begin{sgastatement}{Corollary 2.6}\sgareaderlabel{III.2.6}{2.6}
Every finite $M$-regular sequence consisting of elements of $I$ can be
extended to a maximal $M$-regular sequence, whose length is necessarily the
$I$-depth of $M$.
\end{sgastatement}

\begin{sgastatement}{Remark 2.7}\sgareaderlabel{III.2.7}{2.7}\srcpage{33}
One can scarcely refrain from saying that the greater the depth of an
$A$-module, the more beautiful it is. A module whose support does not meet
$V(I)$ is among the most beautiful: indeed, $\prof_I M$ is finite if and
only if $\Supp M\cap V(I)\ne\varnothing$.
\end{sgastatement}

\begin{sgastatement}{Remark 2.8}\sgareaderlabel{III.2.8}{2.8}
If $A$ is a semi-local ring, let $\mathfrak r(A)$ be its radical and
$k=A/\mathfrak r(A)$ its residue ring. The relevant notion of depth is
obtained by taking $I=\mathfrak r(A)$; accordingly, we write simply
$\prof M$ for the $\mathfrak r(A)$-depth of an $A$-module $M$. This is the
notion of ``homological codimension'' (cf. Serre, \emph{op. cit.}), formerly
denoted $\codh_A M$, defined as the infimum of the integers $i$ such that
$\Extop_A^i(k,M)\ne0$. Indeed, $\Supp k=V(\mathfrak r(A))$.
\end{sgastatement}

\begin{sgastatement}{Proposition 2.9}\sgareaderlabel{III.2.9}{2.9}
If $A$ is noetherian and $M$ is finitely generated, then
\[
  \prof_I M=\inf_{\mathfrak p\in V(I)}\prof M_{\mathfrak p}.
\]
\end{sgastatement}

\begin{sgastatement}{Corollary 2.10}\sgareaderlabel{III.2.10}{2.10}
If $A$ is a noetherian semi-local ring and $M$ is a finitely generated
$A$-module, then
\[
  \prof M=\inf_{\mathfrak m}\prof M_{\mathfrak m},
\]
where $\mathfrak m$ ranges over the maximal ideals of $A$.
\end{sgastatement}

The corollary follows immediately from \hyperref[III.2.9]{Proposition 2.9}, since the prime
ideals containing the radical are precisely the maximal ideals.

We now prove \hyperref[III.2.9]{Proposition 2.9}. Let $f\in I$. If $f$ is $M$-regular and
$\mathfrak p\supset I$, then the image $g$ of $f$ in $A_{\mathfrak p}$
belongs to the maximal ideal $\mathfrak pA_{\mathfrak p}$ and is
$M_{\mathfrak p}$-regular, as follows from the exact sequence
\begin{equation}\label{eq:III.2.3}
  0\longrightarrow M_{\mathfrak p}\xrightarrow{g'}M_{\mathfrak p}
  \longrightarrow(M/fM)_{\mathfrak p}\longrightarrow0,
\end{equation}
where $g'$ denotes multiplication by $g$ on $M_{\mathfrak p}$. This also
shows that $(M/fM)_{\mathfrak p}\simeq M_{\mathfrak p}/gM_{\mathfrak p}$.
Applying \hyperref[III.2.5]{Corollary 2.5} to $M$ and to $M_{\mathfrak p}$ gives, by induction,
$\prof_I M\le\nu(M)$,\srcpage{34} where
\[
  \nu(M)=\inf_{\mathfrak p\in V(I)}\prof M_{\mathfrak p}.
\]
More precisely, induction gives $\nu(M)=\nu(M/fM)+1$ whenever $f$ is
$M$-regular. It remains to prove that, if $\nu(M)\ge1$, there is an
$M$-regular element in $I$. Applying \hyperref[III.2.1]{Lemma 2.1} to
$M_{\mathfrak p}$, $A_{\mathfrak p}$, and
$\mathfrak pA_{\mathfrak p}$ for every $\mathfrak p\in V(I)$ shows that
$\mathfrak pA_{\mathfrak p}\notin\Ass M_{\mathfrak p}$. Applying the
same lemma to $A$, $M$, and $I$ gives the conclusion.

\begin{sgastatement}{Proposition 2.11}\sgareaderlabel{III.2.11}{2.11}
Let $u:A\to B$ be a homomorphism of noetherian rings, let $I$ be an ideal
of $A$, and let $M$ be a finitely generated $A$-module. Put
$I_B=I\otimes_A B$ and $M_B=M\otimes_A B$. If $B$ is flat over $A$, then
\[
  \prof_{I_B}M_B\ge\prof_I M;
\]
if $B$ is faithfully flat over $A$, equality holds.
\end{sgastatement}

Indeed, put $N=A/I$. Flatness gives $N\otimes_A B=B/I_B$; set
$N_B=N\otimes_A B$. Again by flatness and the noetherian hypotheses,
\[
  \Extop_B^i(N_B,M_B)=\Extop_A^i(N,M)\otimes_A B.
\]
Thus $\Extop_A^i(N,M)=0$ implies $\Extop_B^i(N_B,M_B)=0$, and the converse
holds when $B$ is faithfully flat over $A$.



% SOURCE COMPONENT: 03_lemma_3_1.tex
\section*{3. Depth and Topological Properties}
\sgareaderlabel{III.3}{3}

\begin{sgastatement}{Lemma 3.1}\sgareaderlabel{III.3.1}{3.1}\srcpage{34}
Let $X$ be a topological space, $Y$ a closed subspace, and $F$ a sheaf of
abelian groups on $X$. Put $U=X-Y$. If $n$ is an integer, the following
conditions are equivalent:
\begin{enumerate}[label=\textup{(\roman*)},leftmargin=2.6em]
\item $\cH_Y^i(X,F)=0$ for $i<n$.\srcpage{35}
\item For every open subset $V$ of $X$, the group homomorphism
\[
  \HH^i(V,F)\longrightarrow\HH^i(V\cap U,F)
\]
is bijective for $i<n-1$ and injective for $i=n-1$.
\item For every open subset $V$ of $X$,
\[
  \HH_{Y\cap V}^i\bigl(V,F|_V\bigr)=0\qquad\text{for }i<n.
\]
\end{enumerate}
\end{sgastatement}

\begin{proof}
First, (ii) and (iii) are equivalent. Indeed, let $V$ be an open subset of
$X$, and put $X'=V$, $Y'=Y\cap V$, $F'=F|_V$, and $U'=X'-Y'$. Since $Y'$ is
closed in $X'$, there is an exact sequence
\[
  \HH_{Y'}^i(X',F')\longrightarrow \HH^i(X',F')
  \xrightarrow{\rho_i}\HH^i(U',F')
  \longrightarrow\HH_{Y'}^{i+1}(X',F').
\]
If the two outer terms vanish, $\rho_i$ is bijective; if the term on the left
vanishes, $\rho_i$ is injective. Thus (iii) implies (ii). Conversely, if
$i<n$, then $\HH_{Y'}^i(X',F')$ vanishes because $\rho_i$ is injective and
$\rho_{i-1}$ is surjective.

Next, (i) implies (iii). The local-to-global spectral sequence gives
\[
  \HH^p\bigl(X,\cH_Y^q(X,F)\bigr)\Longrightarrow
  \HH_Y^{p+q}(X,F).
\]
By hypothesis, $\cH_Y^q(X,F)=0$ for $q<n$; hence
$\HH_Y^{p+q}(X,F)=0$ whenever $p+q<n$.

Finally, (iii) implies (i), since (iii) says that the presheaf
\[
  V\longmapsto\HH_{Y\cap V}^i\bigl(V,F|_V\bigr)
\]
is zero. Its associated sheaf is therefore zero; because $Y$ is closed, that
associated sheaf is $\cH_Y^i(X,F)$.
\end{proof}


% SOURCE COMPONENT: 04_remark_3_2_proposition_3_3.tex
% Source-aligned production fragment: corrected French lines 1022--1060.
% Requires statement/proof environments and operators \prof, \cH, \cExt.

\begin{sgastatement}{Remark 3.2}\sgareaderlabel{III.3.2}{3.2}\srcpage{36}
The equivalence of (i) and (ii) was proved in Proposition \hyperref[I.2.13]{I.2.13}. As was
observed there, if $n\ge2$, the condition that $\rho_{n-1}$ be injective may
be omitted.\footnote{The original edition repeated a proof here, but that
proof was not entirely correct.}
\end{sgastatement}

\begin{sgastatement}{Proposition 3.3}\sgareaderlabel{III.3.3}{3.3}
Let $X$ be a locally noetherian prescheme, $Y$ a closed subprescheme of $X$,
and $F$ a coherent $\mathcal O_X$-module. The conditions of \hyperref[III.3.1]{Lemma 3.1} are
equivalent to each of the following conditions:
\begin{enumerate}[label=\textup{(\roman*)},leftmargin=2.8em]
\setcounter{enumi}{3}
\item For every $x\in Y$, one has $\prof F_x\ge n$.
\item For every coherent $\mathcal O_X$-module $G$ on $X$ whose support is
contained in $Y$,
\[
  \cExt_{\mathcal O_X}^i(G,F)=0\qquad\text{for }i<n.
\]
\item There exists a coherent $\mathcal O_X$-module $G$ whose support equals
$Y$ and such that
\[
  \cExt_{\mathcal O_X}^i(G,F)=0\qquad\text{for }i<n.
\]
\end{enumerate}
\end{sgastatement}

If $X$ is affine, \hyperref[III.2.4]{Proposition 2.4} supplies what is needed to establish the
equivalence of the three conditions in \hyperref[III.3.3]{Proposition 3.3}. These conditions are
local except for (v)$\Rightarrow$(vi); for that implication one may take
$G=\mathcal O_Y$ and invoke \hyperref[III.2.4]{Proposition 2.4} again. It therefore suffices to
prove (i)$\Rightarrow$(vi) and (v)$\Rightarrow$(i).

Let $J$ be the ideal of $Y$; it is a coherent sheaf of ideals. Put
$\mathcal O_{Y_m}=\mathcal O_X/J^{m+1}$. This is a coherent
$\mathcal O_X$-module whose support equals $Y$, and Theorem \hyperref[II.6]{II.6(b)} gives
\[
  \cH_Y^i(X,F)=\varinjlim_m
  \cExt_{\mathcal O_X}^i(\mathcal O_{Y_m},F).
\]
Thus (v) implies (i). Moreover,\srcpage{37} the transition morphisms in the
projective system of the $\mathcal O_{Y_m}$ are epimorphisms.

If the functor $\cExt^i$ is left exact in its first argument---at least when
that argument lies in the category of coherent $\mathcal O_X$-modules with
support contained in $Y$---then the transition morphisms of the inductive
system obtained by applying $\cExt^i$ to the $\mathcal O_{Y_m}$ are
injective. Condition (i) makes the limit zero, and therefore makes
$\cExt_{\mathcal O_X}^i(\mathcal O_{Y_m},F)$ zero for every $m$.

We argue by induction. The assertion is trivial for $n<0$. Suppose that
(i)$\Rightarrow$(vi) has been proved in degrees below $q$. Then
(i)$\Rightarrow$(v), so the long exact $\cExt$ sequence shows that
$\cExt^q$ is left exact in its first argument. Hence
$\cExt_{\mathcal O_X}^q(\mathcal O_{Y_m},F)=0$ for every $m$. Thus
(i)$\Rightarrow$(vi) holds through degree $q$.
\hfill$\square$\par

% SOURCE COMPONENT: 05_example_3_4.tex
\begin{sgastatement}{Example 3.4}\sgareaderlabel{III.3.4}{3.4}\srcpage{37}
Let $A$ be a noetherian local ring, $\mathfrak m$ its maximal ideal, $M$ a
finitely generated $A$-module, and $n$ an integer. Put
$X=\Spec(A)$, $Y=\{\mathfrak m\}$, and $U=X-Y$. Let $F$ be the sheaf associated
with $M$. The following conditions are equivalent:
\begin{enumerate}[label=\textup{(\arabic*)},leftmargin=2.8em]
\item $\prof M\ge n$.
\item The natural homomorphism
\[
  \HH^i(X,F)\longrightarrow\HH^i(U,F)
\]
is injective if $i=n-1$ and bijective if $i<n-1$.
\item $\Extop_A^i(k,M)=0$ if $i<n$, where $k=A/\mathfrak m$.
\item $\HH_Y^i(X,F)=0$ if $i<n$.
\end{enumerate}
\end{sgastatement}

Taking \hyperref[III.3.2]{Remark 3.2} into account, one obtains the following corollary.\srcpage{38}


% SOURCE COMPONENT: 06_corollary_3_5_theorem_3_6.tex
\begin{sgastatement}{Corollary 3.5}\sgareaderlabel{III.3.5}{3.5}\srcpage{38}
Let $X$ be a locally noetherian prescheme, $Y$ a closed subprescheme of $X$,
and $F$ a coherent $\mathcal O_X$-module. The following conditions are equivalent:
\begin{enumerate}[label=\textup{(\arabic*)},leftmargin=2.8em]
\item For every $x\in Y$, $\prof F_x\ge2$.
\item For every open subset $V$ of $X$, the natural homomorphism
\[
  \HH^0(V,F)\longrightarrow\HH^0\bigl(V\cap(X-Y),F\bigr)
\]
is bijective.
\end{enumerate}
\end{sgastatement}

\begin{sgastatement}{Theorem 3.6 (Hartshorne)}\sgareaderlabel{III.3.6}{3.6}
Let $X$ be a locally noetherian prescheme and $Y$ a closed subprescheme of
$X$. Suppose that $\prof\mathcal O_{X,x}\ge2$ for every $x\in Y$. Then the
natural map
\[
  \pi_0(X)\longrightarrow\pi_0(X-Y)
\]
is bijective.
\end{sgastatement}

\begin{proof}
Since $X$ is locally noetherian, it is locally connected. It therefore
suffices to prove that $X$ is connected if and only if $X-Y$ is connected.
A locally ringed space $(X,\mathcal O_X)$ is connected if and only if
$\HH^0(X,\mathcal O_X)$ is not a direct product of two nonzero rings. By
\hyperref[III.3.5]{Corollary 3.5}, applied with $F=\mathcal O_X$, the hypothesis implies that
\[
  \HH^0(X,\mathcal O_X)\longrightarrow
  \HH^0(X-Y,\mathcal O_X)
\]
is an isomorphism, and the conclusion follows.
\end{proof}


% SOURCE COMPONENT: 07_corollary_3_7_reduction.tex
\begin{sgastatement}{Corollary 3.7}\sgareaderlabel{III.3.7}{3.7}
Let $X$ be a locally noetherian prescheme. Let $d$ be an integer such that
$\dim\mathcal O_{X,x}\ge d$ implies $\prof\mathcal O_{X,x}\ge2$. If $X$ is
connected, then it is connected in codimension $d-1$: namely, if $X'$ and
$X''$ are two irreducible components of $X$, there is a sequence of
irreducible components of $X$,
\[
  X'=X_0,X_1,\ldots,X_n=X'',
\]
such that, for every $i$ with $0\le i<n$, the codimension in $X$ of
$X_i\cap X_{i+1}$ is at most $d-1$.
\end{sgastatement}

First note\srcpage{39} that if $X$ is Cohen--Macaulay, then $d=2$ has the
property just invoked. Recall also that the codimension of $Y$ in $X$ is
defined as the infimum of the dimensions of the local rings of $X$ at the
points of $Y$.

\begin{proofbridge}
Let $\mathcal F$ be the collection of closed subsets of $X$ whose codimension
in $X$ is at least $d$. This is an antifilter of closed subsets of $X$. A
closed subset $Y$ belongs to $\mathcal F$ if and only if, for every $y\in Y$,
there is an open neighborhood $V$ of $y$ such that $Y\cap V$ has codimension
at least $d$ in $V$.\footnote{Source correction: the French TeX says “an open
neighborhood $V$ of $X$”; the quantified point $y$ and the local condition
require “of $y$.\!”} Finally, if $X$ is connected and $Y\in\mathcal F$, then
$X-Y$ is connected by Hartshorne's theorem. The corollary therefore follows
from the next lemma, which is purely topological.
\end{proofbridge}


% SOURCE COMPONENT: 08_lemma_3_8.tex
\begin{sgastatement}{Lemma 3.8}\sgareaderlabel{III.3.8}{3.8}
Let $X$ be a connected, locally noetherian topological space, and let
$\mathcal F$ be an antifilter of closed subsets of $X$. Suppose that every
closed subset $Y\subset X$ which belongs locally to $\mathcal F$---that is,
for every $x\in X$ there are an open neighborhood $V$ of $x$ in $X$ and a
$Y'\in\mathcal F$ such that $V\cap Y=V\cap Y'$---belongs to $\mathcal F$.
The following conditions are equivalent:
\begin{enumerate}[label=\textup{(\roman*)},leftmargin=2.7em]
\item For every $Y\in\mathcal F$, the complement $X-Y$ is connected.
\item If $X'$ and $X''$ are distinct irreducible components, there is a
sequence $X_0,X_1,\ldots,X_n$ of irreducible components of $X$ such that
$X'=X_0$, $X''=X_n$, and
\[
  X_i\cap X_{i+1}\notin\mathcal F\qquad(0\le i<n).
\]
\end{enumerate}
\end{sgastatement}

\noindent\textit{Source correction.} The corrected French branch prints
$1\le i<n$ in (ii), omitting the first adjacent pair. Both the chain's
indexing and the proof below require $0\le i<n$; the target displays that
range rather than silently reproducing the defect.

\begin{proof}
Assume (ii), and let $Y\in\mathcal F$. We must prove that $U=X-Y$ is
connected. If $U'$ and $U''$ are two irreducible components of $U$, there are
irreducible components $X'$ and $X''$ of $X$ such that
$X'\cap U=U'$ and $X''\cap U=U''$. Choose a sequence
$X_0,\ldots,X_n$ having the property in (ii), and put
$U_i=X_i\cap U$ for $0\le i\le n$. The $U_i$ are irreducible components of
$U$, and $U_i\cap U_{i+1}$ is nonempty for $0\le i<n$. Otherwise
$X_i\cap X_{i+1}\subset Y$ would belong to $\mathcal F$,\srcpage{40} contrary to the
choice of the $X_i$. Hence $U$ is connected.

Conversely, assume (i). Let
\[
  Y=\bigcup_{X',X''}(X'\cap X''),
\]
where the union is over distinct irreducible components $X',X''$ of $X$ for
which $X'\cap X''\in\mathcal F$. This family of intersections is locally
finite because $X$ is locally noetherian; its members are closed, so $Y$ is
closed. Moreover, $Y$ belongs locally to $\mathcal F$, and hence
$Y\in\mathcal F$. Thus $U=X-Y$ is connected.

Let $X'$ and $X''$ be distinct irreducible components of $X$, and let $U'$
and $U''$ be their traces on $U$; by the construction of $Y$, these traces
are nonempty. They are irreducible components of $U$. Since $U$ is connected
and locally noetherian, there is a sequence $U_0,\ldots,U_n$ of irreducible
components of $U$ such that $U_0=U'$, $U_n=U''$, and, for $0\le i<n$,
\[
  U_i\cap U_{i+1}\ne\varnothing
  \quad\text{and}\quad U_i\cap U_{i+1}\ne U_i.
\]
Let $X_0,\ldots,X_n$ be the irreducible components of $X$ satisfying
$X_i\cap U=U_i$. If $X_i\cap X_{i+1}\in\mathcal F$, then by the construction
of $Y$ one would have either $U_i\cap U_{i+1}=\varnothing$ or
$U_i=U_{i+1}$, contradicting the choice of the $U_i$.
\end{proof}


% SOURCE COMPONENT: 09_corollary_3_9.tex
\begin{sgastatement}{Corollary 3.9}\sgareaderlabel{III.3.9}{3.9}
Let $A$ be a noetherian local ring. Suppose that, for every prime ideal
$\mathfrak p$ of $A$,
\[
  \bigl(\dim A_{\mathfrak p}\ge2\bigr)
  \Longrightarrow
  \bigl(\prof A_{\mathfrak p}\ge2\bigr).
\]
Suppose also that $A$ satisfies the chain condition.\footnote{Cf. EGA
$0_{\mathrm{IV}}$, 14.3.2.} Then every minimal prime ideal $\mathfrak p$ of
$A$ satisfies $\dim(A/\mathfrak p)=\dim A$. Equivalently, all irreducible
components of $\Spec A$ have the same dimension, namely $\dim A$.
\end{sgastatement}

Indeed, put $X=\Spec A$. If $X'$ and $X''$ are two irreducible components of
$X$, \hyperref[III.3.7]{Corollary 3.7} joins them by a chain with the stated properties. It is
therefore enough to prove that two successive components have the same
dimension, which follows from the second hypothesis.


% SOURCE COMPONENT: 10_example_3_10_remark_3_11.tex
\begin{sgastatement}{Example 3.10}\sgareaderlabel{III.3.10}{3.10}\srcpage{41}
Let $X$ be the union of two complementary linear subspaces of respective
dimensions $2$ and $3$ in a vector space of dimension $5$. More precisely,
let $X=\Spec A$, where
\[
  A=B/(\mathfrak p\cap\mathfrak q),\qquad
  B=k[X_1,\ldots,X_5],
\]
$\mathfrak p=(X_1,X_2,X_3)$, and $\mathfrak q=(X_4,X_5)$. Removing the
intersection point $x$ of the two linear subspaces disconnects $X$.
Consequently, the depth of $\mathcal O_{X,x}$ is $1$, since it cannot be at
least $2$ by \hyperref[III.3.6]{Theorem 3.6}. The failure of the preceding corollary's
equidimensionality conclusion gives another reason.

More generally, let $X$ be the union of two linear subspaces of dimensions
$p,q\ge2$ in a vector space of dimension $p+q$. Under no embedding of $X$
into a regular scheme is $X$ even set-theoretically a complete intersection
at the origin. Indeed, after possibly modifying its scheme structure without
changing the underlying topological space near the origin, $X$ would then be
Cohen--Macaulay and hence of depth at least $2$ at the origin, which it is not.
\end{sgastatement}

\begin{sgastatement}{Remark 3.11}\sgareaderlabel{III.3.11}{3.11}
Let $X$ be a locally noetherian prescheme, $Y$ a closed subprescheme of $X$,
and $F$ an $\mathcal O_X$-module. Depth is a purely topological notion,
expressed in terms of the vanishing of $\cH_Y^i(X,F)$ for $i<n$. One also
wishes to study these sheaves for a fixed $i$, or for $i>n$. The following
result addresses this question.
\end{sgastatement}


% SOURCE COMPONENT: 11_lemma_3_12_to_end.tex
\begin{sgastatement}{Lemma 3.12}\sgareaderlabel{III.3.12}{3.12}
Let $m$ be an integer. The equality $\cH_Y^i(X,F)=0$ for every $i>m$ and
every coherent $F$ holds if and only if it holds for $F=\mathcal O_X$.
\end{sgastatement}

By passage to a direct limit, it then holds for every quasi-coherent sheaf.
For example, if $Y$ can be locally described by $m$ equations---or, as one
says, if $Y$ is locally a set-theoretic complete intersection (as happens,
for example, when $X$ and $Y$ are nonsingular)---the computation of
$\cH_Y^i(X,F)$ by the Koszul complex shows that these sheaves vanish for
$i>m$. This fact was used implicitly in \hyperref[III.3.10]{Example 3.10}. The cohomological
condition is nevertheless insufficient, as the following example shows.

\begin{sgastatement}{Example 3.13}\sgareaderlabel{III.3.13}{3.13}\srcpage{42}
Let $X=\Spec A$, where $A$ is a normal noetherian local ring of dimension
$2$, and let $Y$ be a curve in $X$. The complement $X-Y$ can be shown to be
affine. Hence, for $i>1$, the supported-cohomology exact sequence gives
\[
  \HH_Y^i(X,\mathcal O_X)
  \simeq \HH^{i-1}(X-Y,\mathcal O_X)=0.
\]
Consequently the corresponding local-cohomology sheaf
$\cH_Y^i(\mathcal O_X)$ vanishes. Nevertheless, one can construct a curve
which is not described by a single equation.
\end{sgastatement}

\noindent\textit{Source correction.} The corrected French branch and its
rendered PDF write a sheaf-local-cohomology term with support $X-Y$ in the
middle of the displayed isomorphism. The stated reason is instead the
vanishing of $\HH^{i-1}(X-Y,\mathcal O_X)$ because $X-Y$ is affine; the
target displays the supported/global exact-sequence formula that this reason
requires and then states the sheaf consequence separately.

We shall seek\footnote{Cf. Expos\'e VIII.} conditions ensuring that
$\cH_Y^i(X,F)$ is coherent for a fixed $i$. In general it is not: for
example, consider $\HH_{\mathfrak m}^n(A)$ when $A$ is a noetherian local
ring of dimension $n>0$. If $A$ is a discrete valuation ring with fraction
field $K$, then
\[
  \HH_{\mathfrak m}^1(A)\simeq K/A,
\]
which is not a finitely generated $A$-module.

To illuminate the question, it is equivalent to the following problem. Let
$f:U\to X$ be an open immersion and $G$ a coherent sheaf on $U$. Find
criteria ensuring that the higher direct images $\RR^i f_*G$ are coherent
sheaves on $X$ for a fixed $i$. Such conditions are necessary for the use of
formal geometry encountered in Expos\'e IX and the subsequent expos\'es.

% END INCLUDED active_source/parts/03_expose_III.tex

% BEGIN INCLUDED active_source/parts/04_expose_IV.tex
% SOURCE COMPONENT: 01_unit_IV_1_open\part.tex

\makeatletter
\providecommand{\sgareaderlabel}[2]{%
  \begingroup
  \def\@currentlabel{#2}%
  \phantomsection\label{#1}%
  \endgroup}
\makeatother

\chapter{Dualizing Modules and Functors}
\label{IV}
\section*{1. Generalities on Functors of Modules}
\sgareaderlabel{IV.1}{1}

Fix\srcpage{43}\sgareaderlabel{pIV.1}{1} a commutative noetherian ring $A$. Let $\Ccat$ be the
category of finitely generated $A$-modules, $\Ccat'$ the category of all
$A$-modules, and $\Ab$ the category of abelian groups. This section studies
certain properties of additive functors
\[
  T:\Ccat^{\mathrm{op}}\longrightarrow\Ab.
\]

For $M\in\Ob\Ccat$, the group $T(M)$ has a canonical $A$-module structure.
If $f_M:M\to M$ denotes multiplication by $f\in A$, then $f$ acts on $T(M)$
as $T(f_M)$. In other words, $T$ factors as
\[
  \Ccat^{\mathrm{op}}\xrightarrow{T_0}\Ccat'
  \longrightarrow\Ab,
\]
where the second arrow is the forgetful functor. We henceforth regard
$T(M)$ as endowed with this $A$-module structure.

Compose the isomorphism $M\xrightarrow{\sim}\Homop_A(A,M)$ with the map
\[
  \Homop_A(A,M)\longrightarrow\Homop_A\bigl(T(M),T(A)\bigr)
\]
induced by $T$. This gives the following mutually adjoint forms of the same
pairing:
\begin{align*}
  M&\longrightarrow\Homop_A\bigl(T(M),T(A)\bigr),\\
  M\times T(M)&\longrightarrow T(A),\\
  T(M)&\longrightarrow\Homop_A\bigl(M,T(A)\bigr).
\end{align*}
They define a natural morphism of contravariant functors
\[
  \varphi_T:T\longrightarrow\Homop_A\bigl(-,T(A)\bigr).
\]

\begin{sgastatement}{Proposition 1.1}\sgareaderlabel{IV.1.1}{1.1}\srcpage{44}
The following two properties are equivalent:
\begin{enumerate}[label=\textup{\roman*)},leftmargin=2.7em]
\item $\varphi_T$ is an isomorphism of functors.
\item $T$ is left exact.
\end{enumerate}
\end{sgastatement}

The implication (i)$\Rightarrow$(ii) is immediate. For the converse, use the
following fact: if $u:F\to F'$ is a morphism between additive left-exact
functors from $\Ccat^{\mathrm{op}}$ to $\Ab$ and $u(A)$ is an isomorphism,
then $u$ is an isomorphism. Here one uses that $A$ is noetherian, so every
finitely generated $A$-module is finitely presented.

\begin{sgastatement}{Remark 1.2}\sgareaderlabel{IV.1.2}{1.2}
In particular, the representable functors
$T:{\Ccat'}^{\mathrm{op}}\to\Ab$ are precisely those which commute with
arbitrary inverse limits in ${\Ccat'}^{\mathrm{op}}$ (indexed by a
preordered set, not necessarily filtered).
\end{sgastatement}

Let $\operatorname{Fun}(\Ccat^{\mathrm{op}},\Ab)_{\mathrm{lex}}$ denote the
full subcategory of additive functors that are left exact. We have proved an
equivalence of categories
\[
  \Ccat'\xrightarrow{\sim}
  \operatorname{Fun}(\Ccat^{\mathrm{op}},\Ab)_{\mathrm{lex}},
\]
with mutually quasi-inverse functors
\[
  H\longmapsto\Homop_A(-,H),
  \qquad
  T\longmapsto T(A).
\]

% SOURCE COMPONENT: 02_unit_IV_1_3\part.tex

Let $J$ be an ideal of $A$, put $Y=V(J)\subset\Spec A$, and let
$\Ccat_Y$ be the full subcategory of $\Ccat$ consisting of finitely generated
$A$-modules $M$ with $\Supp M\subset Y$. Then
\[
  \Ccat_Y=\bigcup_n\Ccat^{(n)},
\]
where $\Ccat^{(n)}$ is the full subcategory of modules satisfying $J^nM=0$.

\begin{sgastatement}{Proposition 1.3}\sgareaderlabel{IV.1.3}{1.3}\srcpage{45}
With this notation, let
$T:\Ccat_Y^{\mathrm{op}}\to\Ab$ be an additive functor, and put
\[
  H=\varinjlim_n T(A/J^n).
\]
There is a natural morphism
\[
  \varphi_T:T\longrightarrow\Homop_A(-,H),
\]
and the following conditions are equivalent:
\begin{enumerate}[label=\textup{\roman*)},leftmargin=2.7em]
\item $\varphi_T$ is an isomorphism.
\item $T$ is left exact.
\end{enumerate}
\end{sgastatement}

\noindent\textit{Source correction.} The proposition line prints
$T:\Ccat_Y\to\Ab$, while its target, its restrictions, and the concluding
category equivalence are contravariant. The target restores the required
opposite category.

\begin{proof}
First define $\varphi_T$. For $M\in\Ob\Ccat_Y$, choose $n$ with $J^nM=0$.
Then $M$ is an $A/J^n$-module. If $T_n$ is the restriction of $T$ to
${\Ccat^{(n)}}^{\mathrm{op}}$ and $H_n=T(A/J^n)$, the construction of
\hyperref[IV.1.1]{Proposition 1.1} gives
\[
  T_n\longrightarrow\Homop_A(-,H_n).
\]
Composing with $H_n\to H$ gives, for every such $M$,
\[
  T(M)=T_n(M)\longrightarrow\Homop_A(M,H),
\]
and these maps define $\varphi_T$.

It is clear that (i) implies (ii). Conversely, suppose $T$ is left exact and
let $M\in\Ob\Ccat^{(n)}$. \hyperref[IV.1.1]{Proposition 1.1} applied over $A/J^n$ gives
\[
  T_n(M)\xrightarrow{\sim}\Homop_A(M,H_n).
\]
For every $n'\ge n$, \hyperref[IV.1.1]{Proposition 1.1} applied over $A/J^{n'}$ gives a
compatible isomorphism with $H_{n'}$. Hence
\[
  T(M)\simeq\varinjlim_{n'\ge n}\Homop_A(M,H_{n'}).
\]
Because this is a filtered direct limit and $M$ is finitely presented,
\[
  \varinjlim_{n'\ge n}\Homop_A(M,H_{n'})
  \xrightarrow{\sim}
  \Homop_A\!\left(M,\varinjlim_{n'\ge n}H_{n'}\right)
  =\Homop_A(M,H).
\]

Consequently, if $\Ccat'_Y$ denotes the category of all $A$-modules supported
in $Y$, without a finite-generation condition, there is a natural equivalence
\[
  \Ccat'_Y\xrightarrow{\sim}
  \operatorname{Fun}(\Ccat_Y^{\mathrm{op}},\Ab)_{\mathrm{lex}}.
\]
\end{proof}

% SOURCE COMPONENT: 03_unit_IV_1_4\part.tex
\subsection*{Application}

Retain the preceding notation, and let\srcpage{46}
\[
  T^*:\Ccat_Y^{\mathrm{op}}\longrightarrow\Ab
\]
be an exact $\partial$-functor. For every $i\in\mathbb Z$, put
\[
  H_n^i=T^i(A/J^n),\qquad H^i=\varinjlim_n H_n^i.
\]

\begin{sgastatement}{Theorem 1.4}\sgareaderlabel{IV.1.4}{1.4}
Let $n\in\mathbb Z$. Suppose there is an $i_0\in\mathbb Z$ such that
$T^i=0$ for every $i<i_0$. Then the following three conditions are equivalent:
\begin{enumerate}[label=\textup{\roman*)},leftmargin=2.7em]
\item $T^i=0$ for every $i<n$.
\item $H^i=0$ for every $i<n$.
\item There exists an $M_0\in\Ccat_Y$ such that $\Supp M_0=Y$ and
$T^i(M_0)=0$ for every $i<n$.
\end{enumerate}
\end{sgastatement}

\begin{proof}
Clearly (i) implies (ii) and (iii), taking $M_0=A/J$. We prove
(ii)$\Rightarrow$(i) by induction on $n$. It is true for $n=i_0$; suppose it
has been proved through $n$. If $H^i=0$ for every $i<n+1$, the induction
hypothesis gives $T^i=0$ for $i<n$. The vanishing of $T^{n-1}$ makes $T^n$
left exact, so \hyperref[IV.1.3]{Proposition 1.3} gives
\[
  T^n\xrightarrow{\sim}\Homop_A(-,H^n)=0.
\]

It remains to prove (iii)$\Rightarrow$(ii), again by induction. The assertion
is true for $n=i_0$. Suppose it has been proved through $n$, and let
$M_0\in\Ccat_Y$ satisfy $\Supp M_0=Y$ and $T^i(M_0)=0$ for every $i<n+1$.
The induction hypothesis gives $H^i=0$ for $i<n$; it remains to show
$H^n=0$. The preceding implication gives $T^{n-1}=0$, and therefore
\[
  T^n\xrightarrow{\sim}\Homop_A(-,H^n).
\]
Moreover,
\[
  \Ass\Homop_A(M_0,H^n)
  =\Supp M_0\cap\Ass H^n
  =\Ass H^n,
\]
because
\[
  \Ass H^n\subset\Supp H^n\subset Y=\Supp M_0.
\]
Thus
\[
  T^n(M_0)=0
  \quad\Longleftrightarrow\quad
  \Ass H^n=\varnothing
  \quad\Longleftrightarrow\quad
  H^n=0,
\]
which completes the proof.
\end{proof}

% SOURCE COMPONENT: 04_unit_IV_2_1\part.tex

\section*{2. Characterization of Exact Functors}
\sgareaderlabel{IV.2}{2}

The ring $A$ remains commutative and noetherian. Retain the notation of
\hyperref[IV.1.3]{Proposition 1.3}:\srcpage{47}
\[
  Y=V(J),\qquad
  T:\Ccat_Y^{\mathrm{op}}\longrightarrow\Ab,\qquad
  H=\varinjlim_n T(A/J^n).
\]
Assume that $T$ is left exact, so that
\[
  T(M)\xrightarrow{\sim}\Homop_A(M,H).
\]

\begin{sgastatement}{Proposition 2.1}\sgareaderlabel{IV.2.1}{2.1}
The following properties are equivalent:
\begin{enumerate}[label=\textup{\roman*)},leftmargin=2.7em]
\item $T$ is exact.
\item $H$ is injective in $\Ccat'$, the category of all $A$-modules.
\end{enumerate}
\end{sgastatement}

\begin{proof}
Only (i)$\Rightarrow$(ii) needs proof. Suppose the restriction of
$\Homop_A(-,H)$ to $\Ccat_Y^{\mathrm{op}}$ is exact. Since $A$ is
noetherian, it is enough to prove the following extension property: whenever
$N\subset M$ are finitely generated $A$-modules, every homomorphism
$f:N\to H$ extends to a homomorphism $\bar f:M\to H$.

By the definition of $H$ and finite generation of $N$, there is an integer
$n$ such that $J^n f(N)=0$. Give $M$ and $N$ their $J$-adic topologies. By
Krull's theorem, the $J$-adic topology on $N$ is equivalent to the topology
induced from the $J$-adic topology on $M$. Hence there is a
$V=J^kM$ such that
\[
  U=V\cap N\subset J^nN.
\]
Thus $U\subset\ker f$, and $f$ factors as
\[
\begin{tikzcd}[column sep=large]
N \arrow[r,two heads] \arrow[dr,"f"'] & N/U \arrow[d,"u"]\\
& H.
\end{tikzcd}
\]
Both $N/U$ and $M/V$ belong to $\Ccat_Y$. Exactness therefore extends $u$
to a map $\bar u:M/V\to H$:
\[
\begin{tikzcd}[column sep=large]
N/U \arrow[r,hook] \arrow[dr,"u"'] & M/V \arrow[d,"\bar u"]\\
& H.
\end{tikzcd}
\]
The composite $M\to M/V\xrightarrow{\bar u}H$ is the required extension
$\bar f$.\srcpage{48}
\end{proof}

% SOURCE COMPONENT: 05_unit_IV_2_2\part.tex

Retain the notation $Y=V(J)$ and let $\Ccat_Y$ denote the category of
finitely generated $A$-modules whose support is contained in $Y$.

\begin{sgastatement}{Corollary 2.2}\sgareaderlabel{IV.2.2}{2.2}\srcpage{48}
Let $K$ be an injective $A$-module. Then the submodule
\[
  H^0_J(K)=\{x\in K: J^n x=0\text{ for some }n\geq 0\}
\]
of elements annihilated by a suitable power of $J$ is injective.
\end{sgastatement}

\begin{proof}
By \hyperref[IV.2.1]{Proposition 2.1}, it is enough to verify that the restriction to
$\Ccat_Y^{\mathrm{op}}$ of
$\Homop_A(-,H^0_J(K))$ is exact. Let $M\in\Ccat_Y$. There is an integer
$k$ such that $J^kM=0$, and the inclusion
\[
  \Homop_A(M,H^0_J(K))\longrightarrow \Homop_A(M,K)
\]
is an isomorphism: every map $M\to K$ has image annihilated by $J^k$ and
therefore lands in $H^0_J(K)$. The assertion follows because
$\Homop_A(-,K)$ is exact.
\end{proof}

% SOURCE COMPONENT: 06_unit_IV_3_1\part.tex

\section*{3. The Case Where $T$ Is Left Exact and Every $T(M)$ Is Finitely Generated}
\sgareaderlabel{IV.3}{3}

As above, let\srcpage{48}
\[
  T:\Ccat_Y^{\mathrm{op}}\longrightarrow\Ab.
\]
From now on, assume that $T$ is left exact and factors through
$\Ccat_Y$:
\[
\begin{tikzcd}[column sep=large,row sep=small]
\Ccat_Y^{\mathrm{op}} \arrow[rr,"T"] \arrow[dr,swap,"\widetilde T"]
  && \Ab \\
& \Ccat_Y \arrow[ur,swap,"\text{forget}"] &
\end{tikzcd}
\]
Thus $T(T(M))=(T\circ T)(M)$ is defined. Under the representation
$T(M)\simeq\Homop_A(M,H)$, the canonical evaluation map
\[
  M\longrightarrow
  \Homop_A\bigl(\Homop_A(M,H),H\bigr)
\]
defines a natural morphism\srcpage{49}
\[
  \eta_M:M\longrightarrow (T\circ T)(M).
\]

\begin{sgastatement}{Proposition 3.1}\sgareaderlabel{IV.3.1}{3.1}
Assume, as before, that $A$ is noetherian, and assume in addition that
$A/J$ is artinian. The following conditions are equivalent:
\begin{enumerate}[label=\textup{\roman*)},leftmargin=2.7em]
\item $T$ is left exact and, for every $M\in\Ccat_Y$, the module $T(M)$
is finitely generated and $\eta_M:M\to(T\circ T)(M)$ is an isomorphism.

\item $T$ is exact and, for every residue field $k=A/\mathfrak m$ of a
maximal ideal $\mathfrak m\supset J$, one has $T(k)\simeq k$.

\item There is an isomorphism
$T\simeq\Homop_A(-,H)$ with $H$ injective and, for every $k$ as in
\emph{(ii)}, one has $\Homop_A(k,H)\simeq k$.

\item $T$ is exact and $\ellA T(M)=\ellA M$ for every $M\in\Ccat_Y$.
\end{enumerate}
\end{sgastatement}

\begin{proof}
The equivalence of \emph{(ii)} and \emph{(iii)} follows from
\hyperref[IV.2.1]{Proposition 2.1}.

\smallskip\noindent
\emph{(ii)$\Rightarrow$(iv).}
If $M\in\Ccat_Y$, then $M$ is an $A/J^n$-module for some $n$; since
$A/J^n$ is artinian, $M$ has finite length. We argue by induction on
$\ellA M$. When $\ellA M=1$, the module $M$ is one of the residue fields
occurring in \emph{(ii)}, so the assertion holds. If $\ellA M>1$, choose
a nonzero proper submodule $M'\subset M$ and put $M''=M/M'$. Exactness
of $T$ turns
\[
  0\longrightarrow M'\longrightarrow M\longrightarrow M''\longrightarrow0
\]
into
\[
  0\longrightarrow T(M'')\longrightarrow T(M)
   \longrightarrow T(M')\longrightarrow0.
\]
Both end terms have smaller length, and therefore
\[
\begin{aligned}
  \ellA T(M)
    &=\ellA T(M')+\ellA T(M'')\\
    &=\ellA M'+\ellA M''=\ellA M.
\end{aligned}
\]

\smallskip\noindent
\emph{(ii)$\Rightarrow$(i).}
By \emph{(iv)}, $T(M)$ has finite length and hence is finitely generated.
It remains to prove that every $\eta_M$ is an isomorphism. Again use
induction on $\ellA M$. For a residue field, $T(k)$ and $(T\circ T)(k)$
both have length one; evaluation at $1\in k$ shows that $\eta_k$ is
nonzero, hence it is an isomorphism. For an exact sequence as above,
naturality gives the
commutative diagram with exact rows\srcpage{50}
\[
\begin{tikzcd}[column sep=small]
0 \arrow[r] & M' \arrow[r] \arrow[d,"\eta_{M'}"']
  & M \arrow[r] \arrow[d,"\eta_M"']
  & M'' \arrow[r] \arrow[d,"\eta_{M''}"'] & 0\\
0 \arrow[r] & (T\!\circ T)(M') \arrow[r]
  & (T\!\circ T)(M) \arrow[r]
  & (T\!\circ T)(M'') \arrow[r] & 0.
\end{tikzcd}
\]
The outer vertical arrows are isomorphisms by induction, hence so is
$\eta_M$.

\smallskip\noindent
\emph{(i)$\Rightarrow$(ii).}
Let
\[
  0\longrightarrow M'\longrightarrow M\longrightarrow M''\longrightarrow0
\]
be exact in $\Ccat_Y$, and let $Q$ be the cokernel of
$T(M)\to T(M')$. Left exactness gives the exact sequence
\[
  0\longrightarrow T(M'')\longrightarrow T(M)
   \longrightarrow T(M')\longrightarrow Q\longrightarrow0.
\]
Applying $T$ at the left end and using naturality of evaluation yields
\[
\begin{tikzcd}[column sep=small]
0 \arrow[r] & T(Q) \arrow[r] & (T\!\circ T)(M') \arrow[r]
  & (T\!\circ T)(M)\\
&& M' \arrow[u,"\eta_{M'}","\sim"'] \arrow[r,hook]
  & M \arrow[u,"\eta_M"',"\sim"] .
\end{tikzcd}
\]
Thus $T(Q)=0$, and condition \emph{(i)} gives
$Q\simeq(T\circ T)(Q)=0$. Therefore $T(M)\to T(M')$ is surjective and
$T$ is exact.

Now let $k=A/\mathfrak m$ with $\mathfrak m\supset J$. The $A$-module
$T(k)$ is a finite-dimensional $k$-vector space. It is nonzero because
$k\simeq(T\circ T)(k)$. Choose a decomposition $T(k)\simeq k\oplus V$.
Additivity and biduality give
\[
\begin{aligned}
  (T\circ T)(k)
    &\simeq T(k)\oplus T(V)\\
    &\simeq k\oplus V\oplus T(V)\simeq k,
\end{aligned}
\]
so comparison of finite lengths forces $V=0$. Hence $T(k)\simeq k$.

\smallskip\noindent
\emph{(iv)$\Rightarrow$(iii).}
\hyperref[IV.2.1]{Proposition 2.1} supplies an injective module $H$ and an isomorphism
\[
  T\simeq\Homop_A(-,H).
\]
It remains only to check the residue fields.
For $k=A/\mathfrak m$ as above, $T(k)$ has length one, hence is isomorphic
to a residue field $k'=A/\mathfrak m'$. Moreover
\[
  \Supp k'=\Supp\Homop_A(k,H)\subseteq\Supp k=\{\mathfrak m\}.
\]
Thus $\mathfrak m'=\mathfrak m$ and $k'\simeq k$, as required.
\end{proof}

% SOURCE COMPONENT: 07_unit_IV_3_2_4_2\part.tex

\begin{sgastatement}{Remark 3.2}\sgareaderlabel{IV.3.2}{3.2}
One can\srcpage{51} show that condition \emph{(iv)} is equivalent to the condition
\begin{quote}
\emph{(iv)$'$} For every $M\in\operatorname{Ob}\Ccat_Y$, one has
\[
  \lengthop T(M)=\lengthop M.
\]
\end{quote}
\end{sgastatement}

\section*{4. Dualizing Module. Dualizing Functor}
\sgareaderlabel{IV.4}{4}

\begin{sgastatement}{Definition 4.1}\sgareaderlabel{IV.4.1}{4.1}
Let $A$ be a noetherian local ring and let $\mathfrak m$ be its maximal
ideal. A \emph{dualizing functor} for $A$ is a functor
\[
  T:\Ccat_{\mathfrak m}^{\mathrm{op}}\longrightarrow\Ab,
\]
where $\Ccat_{\mathfrak m}$ denotes $\Ccat_Y$ when
$Y=V(\mathfrak m)$, that satisfies the equivalent conditions of
\hyperref[IV.3.1]{Proposition 3.1}. An $A$-module $I\in\operatorname{Ob}\Ccat'_{\mathfrak m}$ is said to be \emph{dualizing for
$A$} if the functor
\[
  M\longmapsto\Homop_A(M,I)
\]
is dualizing.
\end{sgastatement}

\hyperref[IV.4.1]{Definition 4.1} can be generalized without assuming that $A$ is local.

\begin{sgastatement}{Definition 4.2}\sgareaderlabel{IV.4.2}{4.2}
Let $A$ be a noetherian ring, and let $\Cbar$ be the full subcategory of
$\Ccat$ consisting of the $A$-modules of finite length. A
\emph{dualizing functor} is an exact $A$-linear functor
\[
  T:\Cbar^{\mathrm{op}}\longrightarrow\Cbar
\]
such that the morphism of functors
\[
  \mathrm{id}\longrightarrow T\circ T
\]
is an isomorphism.
\end{sgastatement}

We shall prove an existence theorem and also prove that the module $I$
representing such a functor is locally artinian. We shall further show
that, for every maximal ideal $\mathfrak m$ of $A$, the
$\mathfrak m$-primary component of the socle of $I$ has length one.

% SOURCE COMPONENT: 08_unit_IV_4_3_5\part.tex

\begin{sgastatement}{Proposition 4.3}\sgareaderlabel{IV.4.3}{4.3}
Let\srcpage{51} $A$ and $B$ be noetherian local rings with maximal ideals
$\mA$ and $\mB$, and suppose that $B$ is a finite $A$-algebra. If $I$
is a dualizing module for $A$, then $\Homop_A(B,I)$ is a dualizing
module for $B$.
\end{sgastatement}

\begin{proof}
Let\srcpage{52}
\[
  R:\Ccat_{\mB}\longrightarrow\Ccat_{\mA}
\]
be the restriction-of-scalars functor; it is exact. Let $T$ be a
dualizing functor for $A$,
\[
  T:\Ccat_{\mA}^{\mathrm{op}}\longrightarrow\Ab.
\]
It is exact, and for every $M\in\operatorname{Ob}\Ccat_{\mA}$ the
natural morphism
\[
  M\longrightarrow (T\circ T)(M)
\]
is an isomorphism. Consequently, the contravariant composite $T\circ R$
is a dualizing functor for $B$. If $I$ represents $T$, the classical
formula
\[
  \Homop_A(M,I)\simeq
  \Homop_B\bigl(M,\Homop_A(B,I)\bigr),
\]
valid for every $B$-module $M$, shows that $\Homop_A(B,I)$ is a
dualizing module for $B$.
\end{proof}

\begin{sgastatement}{Corollary 4.4}\sgareaderlabel{IV.4.4}{4.4}
Let $A$ be a noetherian local ring, let $\mathfrak a$ be an ideal of
$A$, and put $B=A/\mathfrak a$. If $I$ is a dualizing module for $A$,
then the annihilator of $\mathfrak a$ in $I$ is a dualizing module for
$B$.
\end{sgastatement}

\begin{sgastatement}{Lemma 4.5}\sgareaderlabel{IV.4.5}{4.5}
Let $A$ be a noetherian local ring and let $I$ be a locally artinian
$A$-module. There is a canonical isomorphism
\[
  I\xrightarrow{\sim}\widehat I=I\otimes_A\widehat A.
\]
\end{sgastatement}

\begin{proof}
Let $I_n$ be the annihilator of $\mathfrak m^n$ in $I$, where
$\mathfrak m$ is the maximal ideal of $A$. To say that $I$ is locally
artinian means that $I$ is the direct limit of the $I_n$ and that each
$I_n$ has finite length. Since tensor product commutes with direct
limits, it is enough to treat each $I_n$ separately. The relation
$\mathfrak m^nI_n=0$ gives, for every $p\geq n$,
\[
  I_n\xrightarrow{\sim}I_n\otimes_A A/\mathfrak m^p
      \xrightarrow{\sim}I_n\otimes_A\widehat A.
\]
The second isomorphism uses
$\widehat A/\mathfrak m^p\widehat A\simeq A/\mathfrak m^p$. Passing
to the direct limit over $n$, and again using commutation with tensor
product, gives
\[ I\xrightarrow{\sim}I\otimes_A\widehat A. \]

It follows that\srcpage{53} restriction of scalars from $\widehat A$ to
$A$ and extension of scalars induce mutually quasi-inverse
\emph{equivalences} between the category of locally artinian
$\widehat A$-modules and the category of locally artinian $A$-modules.
\end{proof}

% SOURCE COMPONENT: 09_unit_IV_4_6_9\part.tex

\begin{sgastatement}{Proposition 4.6}\sgareaderlabel{IV.4.6}{4.6}
Let\srcpage{53} $A$ be a noetherian local ring, let $\widehat A$ be its completion,
let $I$ be a dualizing module for $A$ (respectively, for $\widehat A$), and
let $J$ be the completed module\footnote{Here ``the completed module'' means
$\widehat I=I\otimes_A\widehat A$ (cf. \hyperref[IV.4.5]{Lemma 4.5}): this is $I$ endowed with
its canonical $\widehat A$-module structure, not its $\mideal$-adic
completion. For example, if $p$ is prime and
$A=\widehat A=\mathbb Z_p$ is the ring of $p$-adic integers, then the injective envelope of
$k=\mathbb Z/p\mathbb Z$ is the discrete $\mathbb Z_p$-module
$\mathbb Q_p/\mathbb Z_p$, whose $p$-adic completion is zero.}
of $I$ (respectively, the $A$-module obtained by restriction of scalars).
Then $J$ is a dualizing module for $\widehat A$ (respectively, for $A$).
Moreover, the underlying abelian groups of $I$ and $J$ are isomorphic.
\end{sgastatement}

\begin{proof}
One merely observes that the equivalence between the category of locally
artinian $A$-modules and the category of locally artinian
$\widehat A$-modules induces an isomorphism between the bifunctors
$\Homop_A(-,-)$ and $\Homop_{\widehat A}(-,-)$, and that the
characterization of a dualizing functor or a dualizing module involves only
these bifunctors.
\end{proof}

\begin{sgastatement}{Theorem 4.7}\sgareaderlabel{IV.4.7}{4.7}
Let $A$ be a noetherian local ring.
\begin{enumerate}[label=\textup{\alph*)},leftmargin=2.5em]
\item A dualizing module $I$ always exists.
\item Any two dualizing modules are isomorphic (by a noncanonical
isomorphism).
\item An $A$-module $I$ is dualizing if and only if it is an injective
envelope of the residue field $k$ of $A$.
\end{enumerate}
\end{sgastatement}

\begin{sgastatement}{Remark 4.8}\sgareaderlabel{IV.4.8}{4.8}
\hyperref[IV.4.6]{Proposition 4.6} reduces the question to the case of a complete noetherian
local ring. By a structure theorem of Cohen,\footnote{See I.~S. Cohen,
``On the structure and ideal theory of complete local rings,''
\emph{Trans. Amer. Math. Soc.} \textbf{59} (1946), 54--106.}
such a ring is a quotient of a regular ring. \hyperref[IV.4.3]{Proposition 4.3} then permits us
to assume that $A$ is regular. As we shall see later, this observation gives
an explicit calculation of the dualizing module.\footnote{This was the
method followed by Grothendieck in 1957. The method by injective envelopes
used below appears to be due to K.~Morita, ``Duality for modules and its
applications to the theory of rings with minimum conditions,''
  \emph{Sc. Rep. Tokyo Kyoiku Daigaku} \textbf{6} (1958/59), 83--142.
Morita's work is independent of Grothendieck's, considerably predates this
seminar, and is not restricted to commutative base rings.}
We shall nevertheless prove \hyperref[IV.4.7]{Theorem 4.7} by other means.
\end{sgastatement}

\subsection*{Recollections}

Before\srcpage{54} proving the theorem, we recall several facts about injective
envelopes; see Gabriel, \emph{Des Cat\'egories Ab\'eliennes}, thesis,
Paris, 1961, Chapter II, \S5.

Let $\Ccat$ be an abelian category in which direct limits exist and are
exact\footnote{More precisely, small filtered direct limits are assumed
exact; one should also assume the existence of a generator. See
  [T\^ohoku]. For the category of modules, which suffices here, one may
instead consult Chapter 10 of Bourbaki's \emph{Alg\`ebre}.}
(for example, the category of modules). Every object $M$ embeds in an
injective object, and an \emph{injective envelope} of $M$ means an injective
object containing $M$ that is minimal among such objects. The following
properties hold:
\begin{enumerate}[label=\textup{\roman*)},leftmargin=2.7em]
\item Every object $M$ has an injective envelope $I$.
\item If $I$ and $J$ are two injective envelopes of $M$, there is an
isomorphism $I\simeq J$ (generally not unique) that induces the identity on
$M$.
\item The object $I$ is an \emph{essential extension}\sgareaderlabel{pIV.13}{13} of $M$: if
$P\subset I$ and $P\cap M=\{0\}$, then $P=\{0\}$. Conversely, if $I$ is
injective and is an essential extension of $M$, then $I$ is an injective
envelope of $M$.
\end{enumerate}

Granting these facts, it plainly suffices to prove part~\emph{c)} of
\hyperref[IV.4.7]{Theorem 4.7}.

\begin{proof}
Suppose first that $I$ is a dualizing module for $A$. Then $I$ is injective
and $\Homop_A(k,I)\simeq k$. Composing this isomorphism with the inclusion
\[
  \Homop_A(k,I)\hookrightarrow\Homop_A(A,I)\simeq I
\]
gives an inclusion
\[
  k\hookrightarrow I.
\]
We show that $I$ is an injective envelope of $k$. Let $J$ be an injective
module such that
\[
  k\subset J\subset I.
\]
Since $J$ is injective, there is an injective $A$-submodule $J'$ of $I$ such
that $I=J\oplus J'$. We claim that $\Homop_A(k,J')=0$. Indeed,
\[
  \Homop_A(k,I)\simeq\Homop_A(k,J)\oplus\Homop_A(k,J').
\]
The first summand is\srcpage{55} a nonzero vector subspace of
$\Homop_A(k,I)\simeq k$ (it contains the inclusion $k\subset J$).
Consequently $\Homop_A(k,J)\simeq k$ and $\Homop_A(k,J')=0$.

Induction on length now gives $\Homop_A(M,J')=0$ for every finite-length
$A$-module $M$. The module $I$ is the direct limit of the modules
$\Homop_A(A/\mideal^n,I)$ (cf. \hyperref[IV.1.3]{Proposition 1.3}), which have finite length by
hypothesis. Hence the projection $I\to J'$ vanishes, and therefore $J'=0$.

Conversely, let $I$ be an injective envelope of $k$. By \hyperref[IV.2.1]{Propositions 2.1}
and \hyperref[IV.3.1]{3.1\emph{(ii)}}, to prove that $I$ is dualizing it suffices to show that
$V=\Homop_A(k,I)$ is isomorphic to $k$. There are inclusions
\[
  k\subset V\subset I.
\]
As a $k$-vector space, $V$ decomposes as $k\oplus V'$, where $V'$ is a
vector subspace of $I$ with $V'\cap k=0$. Since $I$ is an essential
extension of $k$, we have $V'=0$, and thus $V=k$.
\end{proof}

\begin{sgastatement}{Corollary 4.9}\sgareaderlabel{IV.4.9}{4.9}
Let $A$ be a noetherian local ring. Every dualizing module for $A$ is
locally artinian.
\end{sgastatement}

\begin{proof}
Let $I$ be a dualizing module; it is an injective envelope of $k$. Using the
notation and result of \hyperref[IV.2.2]{Corollary 2.2}, we have
\[
  k\subset H^0_{\mideal}(I)\subset I,
\]
and $H^0_{\mideal}(I)$ is injective. It follows that
$I=H^0_{\mideal}(I)$, and hence that $I$ is locally artinian.\footnote{As
already observed, one may alternatively note directly that $I$ is the
direct limit of the modules $\Homop_A(A/\mideal^n,I)$.}
\end{proof}

% SOURCE COMPONENT: 10_unit_IV_5_1\part.tex

\section*{5. Consequences of the theory of dualizing modules}\sgareaderlabel{IV.5}{5}\srcpage{56}

The functor
\[
  T=\Homop_A(-,I)\colon \Ctors\longrightarrow\Ctors
\]
is an anti-equivalence. Indeed, $T\circ T$ is isomorphic to the identity
functor, and the argument is formal from that point on.

\emph{One deduces the usual properties of the notion of orthogonality:}

Let $M^*=\Homop_A(M,I)=T(M)$, and let $N\subset M$ be a submodule. The
\emph{orthogonal of} $N$ is the submodule $N'$ of $M^*$ consisting of those
elements of $M^*$ that vanish on $N$. In this way one obtains a bijection
between the set of submodules of $M$ and the set of submodules of $M^*$,
reversing the order relation.

In particular:
\begin{itemize}[leftmargin=2.3em]
\item $\operatorname{length}_{M}N=
  \operatorname{colength}_{M^*}N'$.
\item Cyclic modules, that is, those for which $M/\mideal M$ has dimension
  $0$ or $1$, correspond under the duality to modules whose socle has length
  $0$ or $1$.
\item If $A$ is artinian, the ideals of $A$ correspond to the submodules of
  $I$.
\end{itemize}
And so on.

Let $A$ be a noetherian local ring. Write $\DA$ for the category of
$A$-modules $M$ such that, for every $n\in\mathbb N$,
\[
  M^{(n)}=M/\mideal^{n+1}M
\]
has finite length and such that
\[
  M=\varprojlim_n M^{(n)},
\]
and let $\widehat A$ be the completion of $A$. Restriction of scalars and
completion are quasi-inverse equivalences between $\DA$ and $\DhatA$; up to
isomorphism, they are compatible with taking the underlying abelian groups of
the modules in question. Write $\CA$ for the category of locally artinian
$A$-modules with finite-dimensional socle.

\begin{sgastatement}{Proposition 5.1}\sgareaderlabel{IV.5.1}{5.1}\srcpage{57}
Let $A$ be a noetherian local ring, and let $I$ be a dualizing module for
$A$. The functors
\[
  \Homop_A(-,I)\colon (\CA)^{\mathrm{op}}\longrightarrow\DA
\]
and
\[
  \Homop_{\widehat A}(-,I)\colon \DA\longrightarrow(\CA)^{\mathrm{op}}
\]
are equivalences of categories, quasi-inverse to one another.

Moreover, transporting these functors by the equivalences of categories
between $\DA$ and $\DhatA$ on the one hand, and between $\CA$ and $\ChatA$ on
the other, yields the functor $\Homop_{\widehat A}(-,I)$.
\end{sgastatement}

\begin{proof}
Let $X\in\operatorname{Ob}(\CA)$. By definition,
\[
  X=\varinjlim_{k\in\mathbb N}X_k,
  \qquad
  X_k=\Homop_A(A/\mideal^{k+1},X),
\]
and therefore
\[
  \Homop_A(X,I)=\varprojlim_k\Homop_A(X_k,I).
\]
Thus
\[
  Y=\Homop_A(X,I)=\varprojlim_k\Homop_A(X_k,I)
\]
is a finitely generated $\widehat A$-module, as follows from
\emph{EGA}~$0_{\mathrm I}$, 7.2.9. In this connection, observe that $\DhatA$
is also the category of finitely generated $\widehat A$-modules, or,
equivalently, that $\DA$ is the category of $A$-modules that are complete and
finitely generated over $\widehat A$.\footnote{The printed French and the
corrected TeX write $Y=\varprojlim_k X_k$.  But the $X_k$ were introduced as
an increasing direct system, whereas the immediately preceding display gives
$\Homop_A(X,I)=\varprojlim_k\Homop_A(X_k,I)$.  The target therefore restores
this latter module as $Y$, which is also the object required by the remainder
of the proof.}

Now let $Y$ be such a module, and let $f\colon Y\to I$ be a
$\widehat A$-homomorphism. The image of $f$ is a finitely generated
submodule, hence is annihilated by $\mideal^k$ for some $k$: indeed, every
$x\in I$ is annihilated by some power of $\mideal$. Consequently $f$ factors
through $Y/\mideal^kY$, and it follows that
\begin{align*}
  \Homop_{\widehat A}(Y,I)
    &=\varinjlim_k\Homop_{\widehat A}(Y^{(k)},I),
      &\text{where }Y^{(k)}&=Y/\mideal^{k+1}Y,\\
    &=\varinjlim_k (Y^{(k)})^*.
\end{align*}
This object belongs to $\operatorname{Ob}(\CA)$. It follows immediately that
the two functors in the statement are quasi-inverse to one another.
\end{proof}

% SOURCE COMPONENT: 11_unit_IV_5_2_4\part.tex

It follows from the preceding discussion that replacing $A$ by $\widehat A$
changes neither the categories or functors under consideration nor the
underlying abelian groups of the modules involved. \hyperref[IV.5.1]{Proposition 5.1} may
therefore be restated as follows.\srcpage{58}

The restriction of $\Homop_{\widehat A}(-,I)$ to the category of finitely
generated $\widehat A$-modules takes its values in the category of locally
artinian $\widehat A$-modules having finite-dimensional socle. It admits a
quasi-inverse, namely the restriction of $\Homop_{\widehat A}(-,I)$ in the
opposite direction. On the intersection of these two categories the two
functors coincide (of course) and define an autoduality of the category of
finite-length $\widehat A$-modules.

\begin{sgastatement}{Example 5.2 (Macaulay)}\sgareaderlabel{IV.5.2}{5.2}
Let $A$ be a local ring with residue field $k$. Let $k_0$ be a subfield of
$A$ such that $k$ is finite over $k_0$, and put $[k:k_0]=d$. Every
finite-length $A$-module may be regarded as a finite-dimensional
$k_0$-vector space of dimension
\[
  d\,\operatorname{length}(M).
\]
The functor
\[
  T(M)=\Homop_{k_0}(M,k_0)
\]
is exact and preserves length, hence is dualizing for $A$. Its associated
dualizing module is
\[
  A'=\varinjlim_n\Homop_{k_0}(A/\mideal^n,k_0),
\]
the topological dual of $A$ endowed with the $\mideal$-adic topology.
\end{sgastatement}

\begin{sgastatement}{Example 5.3}\sgareaderlabel{IV.5.3}{5.3}
Let $A$ be a regular noetherian local ring of dimension $n$, let
$\mideal$ be its maximal ideal, and let $k$ be its residue field. There is a
regular system of parameters $(x_1,x_2,\ldots,x_n)$ that generates
$\mideal$ and is an $A$-regular sequence. The Koszul complex therefore gives
\[
\begin{aligned}
  \Extop_A^i(k,A)&=0 &&\text{if }i\ne n,\\
  \Extop_A^n(k,A)&\simeq k.
\end{aligned}
\]
Since the depth\srcpage{59} of $A$ is $n$, one has
$\Extop_A^i(M,A)=0$ for $i<n$ whenever
$M$ is annihilated by a power of $\mideal$. One also has
$\Extop_A^i(M,A)=0$ for $i>n$, because the global cohomological dimension
of $A$ equals $n$. Thus $\Extop_A^n(-,A)$ is exact, and
$\Extop_A^n(k,A)\simeq k$. It follows that:
\end{sgastatement}

\begin{sgastatement}{Proposition 5.4}\sgareaderlabel{IV.5.4}{5.4}
If $A$ is a regular noetherian local ring of dimension $n$, the functor
\[
  M\longmapsto\Extop_A^n(M,A)
\]
is dualizing. The associated dualizing module is
\[
  I=\varinjlim_r\Extop_A^n(A/\mideal^r,A).
\]
It is isomorphic to $H_{\mideal}^n(A)$ (Expos\'e II, \hyperref[II.6]{Theorem 6}).\footnote{For
a ring $A$, an ideal $J\subset A$, an $A$-module $M$, and
$i\in\mathbb Z$, set $H_J^i(M)=H_Y^i(X,\widetilde M)$, where
$X=\Spec(A)$ and $Y=V(J)$.}
\end{sgastatement}

% SOURCE COMPONENT: 12_unit_IV_5_5\part.tex

\begin{sgastatement}{Remark 5.5}\sgareaderlabel{IV.5.5}{5.5}\srcpage{59}
If $A$ satisfies the hypotheses of both preceding examples, the two resulting
dualizing modules are isomorphic. Suppose, for example, that $A$ is regular
of dimension $n$, complete, and of equal characteristic. There is then a
coefficient field, say $K$. After choosing a system of parameters
$(x_1,\ldots,x_n)$ of $A$, one can construct an isomorphism
\[
  A\simeq K[[T_1,\ldots,T_n]],
\]
and hence, as we shall see, an explicit isomorphism between the two dualizing
modules
\[
  v:H_{\mideal}^n(A)\longrightarrow A'.
\]

This isomorphism has an intrinsic interpretation in terms of the module
\[
  \Omega^n=\Omega^n(A/K)
\]
of completed relative differentials of top degree. Indeed, $\Omega^n$ has a
basis consisting of
\[
  dx_1\wedge dx_2\wedge\cdots\wedge dx_n.
\]
There is therefore an isomorphism
\[
  u:H_{\mideal}^n(\Omega^n)\longrightarrow H_{\mideal}^n(A).
\]
A remarkable fact is that the composite
\[
  vu=w:H_{\mideal}^n(\Omega^n)\longrightarrow A'
\]
is\srcpage{60} independent of the chosen system of parameters and commutes
with change of the base field.

To construct $v$, compute $H_{\mideal}^n(A)$ from the Koszul complex
associated with the $x_i$. One obtains
\[
  H_{\mideal}^n(A)=
  \varinjlim_r A/(x_1^r,\ldots,x_n^r).
\]
The transition morphisms are as follows. Put
\[
  I_r=A/(x_1^r,\ldots,x_n^r),
\]
and let $e_{a_1,\ldots,a_n}^r$ be the image of
$x_1^{a_1}x_2^{a_2}\cdots x_n^{a_n}$ in $I_r$. The elements
$e_{a_1,\ldots,a_n}^r$, for $0\le a_i<r$, form a basis of $I_r$.

For an integer $s\geq0$, the transition morphism
\[
  t_{r,r+s}:I_r\longrightarrow I_{r+s}
\]
is multiplication by $x_1^s x_2^s\cdots x_n^s$; consequently,
\[
  t_{r,r+s}\bigl(e_{a_1,\ldots,a_n}^r\bigr)
    =e_{a_1+s,\ldots,a_n+s}^{r+s}.\footnote{The printed French and the
  corrected TeX write $u_{r,r+s}$ in this display, immediately after defining
  the transition map as $t_{r,r+s}$. The target emends the isolated $u$ to
  $t$; the domain, codomain, and multiplication rule make the reference
  unambiguous.}
\]

Giving an $A$-homomorphism $w:M\to A'$ is equivalent to giving a
$K$-linear form $w':M\to K$ that is continuous on every finitely generated
submodule. For $M=H_{\mideal}^n(\Omega^n)$, defining $w$ is therefore
equivalent to defining a linear form
\[
  \rho:H_{\mideal}^n(\Omega^n)\longrightarrow K,
\]
called the \emph{residue form}\sgareaderlabel{pagecithar}{46}.\footnote{\sgareaderlabel{cithar}{*}For a fuller treatment of residues,
see R.~Hartshorne, \emph{Residues and Duality}, Lecture Notes in Mathematics
20, Springer, 1966.} To construct $\rho$, it is enough to define compatible
forms $\rho_r:I_r\to K$ by
\[
  \rho_r\bigl(e_{a_1,\ldots,a_n}^r\bigr)=
  \begin{cases}
    1,&\text{if }a_i=r-1\text{ for }1\le i\le n,\\
    0,&\text{otherwise.}
  \end{cases}
\]
\end{sgastatement}

% END INCLUDED active_source/parts/04_expose_IV.tex

% BEGIN INCLUDED active_source/parts/05_expose_V.tex
% BEGIN COMPONENT SGA2-V-C01 :: unit_V_1_1_1_3_opening
\makeatletter
\providecommand{\sgareaderlabel}[2]{%
  \begingroup
  \def\@currentlabel{#2}%
  \phantomsection\label{#1}%
  \endgroup}
\makeatother

\chapter{\texorpdfstring{Local Duality and Structure of the $H^i(M)$}{Local Duality and Structure of the H\textasciicircum i(M)}}
\label{V}

\section*{1. Complexes of homomorphisms}
\sgareaderlabel{V.1}{1}

\subsubsection*{1.1}
\sgareaderlabel{V.1.1}{1.1}

Let $F^{\bullet}$ and $G^{\bullet}$ be two graded modules. Write\srcpage{61}
\begin{equation}\tag{1}\label{eq:V.1}
  \Homop^{\bullet}(F^{\bullet},G^{\bullet})
\end{equation}
for the graded module of homomorphisms of graded modules from $F^{\bullet}$
to $G^{\bullet}$. Thus
\begin{equation}\tag{2}\label{eq:V.2}
  \Homop^{s}(F^{\bullet},G^{\bullet})
    =\prod_k \Homop(F^k,G^{k+s}).
\end{equation}

Suppose that $F^{\bullet}$ (respectively $G^{\bullet}$) is a complex, and let
$d_1$ (respectively $d_2$) be its differential. For
$h\in\Homop^s(F^{\bullet},G^{\bullet})$, set\footnote{\textit{Editor's note.}
The original sign convention was different, but it is incompatible with the
convention of Expos\'e VIII. The latter seems more reasonable because in that
case the degree-zero cohomology is the set of homotopy classes of morphisms
from $F^{\bullet}$ to $G^{\bullet}$. The calculations below have been modified
accordingly.}
\begin{equation}\tag{3}\label{eq:V.3}
  d(h)=h\circ d_1+(-1)^{s+1}d_2\circ h.
\end{equation}
One checks immediately that $d\circ d=0$; hence
$\Homop^{\bullet}(F^{\bullet},G^{\bullet})$, equipped with $d$, is a complex.
The cohomology group of this complex is denoted by
\begin{equation}\tag{4}\label{eq:V.4}
  \UH^{\bullet}(F^{\bullet},G^{\bullet}).
\end{equation}
If $G^{\bullet}$ is injective in every degree, then
\[
  F^{\bullet}\longmapsto
  \UH^{\bullet}(F^{\bullet},G^{\bullet})
\]
is an exact $\partial$-functor. Likewise, for arbitrary $F^{\bullet}$,
\[
  G^{\bullet}\longmapsto
  \UH^{\bullet}(F^{\bullet},G^{\bullet})
\]
is an exact $\delta$-functor on the category of complexes $G^{\bullet}$ that
are injective in every degree.

\begin{sgaremark}{Remark 1.2}\sgareaderlabel{V.1.2}{1.2}
The cycles of $\Homop^{\bullet}(F^{\bullet},G^{\bullet})$ are the
homomorphisms from $F^{\bullet}$ to $G^{\bullet}$ that commute or
anticommute with the differentials, according to the degree. The boundaries
of $\Homop^{\bullet}(F^{\bullet},G^{\bullet})$ are the homomorphisms from
$F^{\bullet}$ to $G^{\bullet}$ that are homotopic to zero.
\end{sgaremark}

\newpage
Let $A$ be a ring, let $M$ (respectively $N$) be an $A$-module, and let
$R(M)$ (respectively $R(N)$) be an injective resolution of $M$ (respectively
$N$). There is then a canonical isomorphism\srcpage{62}\footnote{\textit{Editor's
note.} The strange original numbering has been preserved.}
\begin{equation*}\tag{1.3}\label{eq:V.1.3}
  \UH^s(R(M),R(N))\simeq\Extop^s(M,N).
\end{equation*}
Indeed, let $i\colon M\to R(M)$ be the canonical augmentation, and let
$h\in\Homop^s(R(M),R(N))$. Denote by $t_s$ the map
\[
  h\longmapsto h_0\circ i
\]
from $\Homop^s(R(M),R(N))$ to $\Homop(M,R(N)^s)$. The family
$(t_s)_{s\geq 0}$ defines a homomorphism of (ordinary) complexes\footnote{Here
$M$ also denotes the complex $M[0]$ consisting of $M$ placed in degree zero.}
\[
  t\colon \Homop^{\bullet}(R(M),R(N))
    \longrightarrow \Homop^{\bullet}(M,R(N));
\]
that is,
\[
  (dh)_0\circ i=d_2\circ h_0\circ i.
\]

One checks easily that, on passing to cohomology, $t$ induces an isomorphism.
It follows in particular that
\[
  \UH^{\bullet}(R(M),R(N))
\]
does not ``depend'' on the chosen injective resolution $R(M)$ (respectively
$R(N)$) of $M$ (respectively $N$).
% END COMPONENT SGA2-V-C01

% BEGIN COMPONENT SGA2-V-C02 :: unit_V_1_4_1_5_section_completion
To every exact sequence of $A$-modules\srcpage{62}
\begin{equation}\tag{5}\label{eq:V.5}
  0\longrightarrow M'\longrightarrow M\longrightarrow M''\longrightarrow 0
\end{equation}
one associates an exact sequence of injective resolutions
\begin{equation}\tag{6}\label{eq:V.6}
  0\longrightarrow R(M')\longrightarrow R(M)
    \longrightarrow R(M'')\longrightarrow 0.
\end{equation}
One checks that the isomorphism \emph{\hyperref[eq:V.1.3]{(1.3)}} commutes with the homomorphisms
\begin{align}
  \UH^s(R(M'),R(N))
    &\longrightarrow \UH^{s+1}(R(M''),R(N)),\tag{8}\label{eq:V.8}\\
  \Extop^s(R(M'),R(N))
    &\longrightarrow \Extop^{s+1}(R(M''),R(N)),\tag{9}\label{eq:V.9}
\end{align}
deduced from \emph{\hyperref[eq:V.6]{(6)}} and \emph{\hyperref[eq:V.5]{(5)}}.\footnote{\textit{Source note.}
The corrected French TeX and the direct compiled source page both use
resolution arguments in \emph{(9)}: $R(M')$ and $R(N)$ in the source term,
and $R(M'')$ and $R(N)$ in the target term. This bounded translation
reproduces that literal reading; it does not silently replace those arguments
by $M'$, $M''$, and $N$.}

Let $P$\srcpage{63} be a third $A$-module, and let $R(P)$ be an injective
resolution of $P$. Then composition of graded morphisms gives a pairing
\begin{equation}\tag{10}\label{eq:V.10}
  \Homop^i(R(N),R(M))\times\Homop^j(R(M),R(P))
    \longrightarrow \Homop^{i+j}(R(N),R(P)),
\end{equation}
which defines a pairing
\begin{equation}\tag{11}\label{eq:V.11}
  \UH^i(R(N),R(M))\times\UH^j(R(M),R(P))
    \longrightarrow \UH^{i+j}(R(N),R(P)),
\end{equation}
and hence a homomorphism of functors in $M$:
\begin{equation*}\tag{1.4}\label{eq:V.1.4}
  \UH^i(R(N),R(M))\longrightarrow
  \Homop\!\left(\UH^j(R(M),R(P)),\UH^{i+j}(R(N),R(P))\right).
\end{equation*}

We shall see that \emph{\hyperref[eq:V.1.4]{(1.4)}} is a homomorphism of $\delta$-functors in
$M$. The exact sequences \emph{\hyperref[eq:V.5]{(5)}} and \emph{\hyperref[eq:V.6]{(6)}} give a commutative
diagram:
\begin{center}
\resizebox{\textwidth}{!}{%
\begin{tikzcd}[ampersand replacement=\&,column sep=large,row sep=large]
\Homop^i(R(N),R(M'))
  \arrow[r] \arrow[d]
\& \Homop\!\left(\Homop^j(R(M'),R(P)),\Homop^{i+j}(R(N),R(P))\right)
  \arrow[d,"{\Homop(q,\mathrm{id})}"] \\
\Homop^i(R(N),R(M))
  \arrow[r] \arrow[d,"p"]
\& \Homop\!\left(\Homop^j(R(M),R(P)),\Homop^{i+j}(R(N),R(P))\right)
  \arrow[d] \\
\Homop^i(R(N),R(M''))
  \arrow[r]
\& \Homop\!\left(\Homop^j(R(M''),R(P)),\Homop^{i+j}(R(N),R(P))\right).
\end{tikzcd}%
}
\end{center}
Let $h\in\Homop^i(R(N),R(M''))$ (respectively
$g\in\Homop^j(R(M'),R(P))$) be a cycle, and let
$h'\in\Homop^i(R(N),R(M))$ (respectively
$g'\in\Homop^j(R(M),R(P))$) be such that $p(h')=h$ (respectively
$q(g')=g$). To say that \emph{\hyperref[eq:V.1.4]{(1.4)}} is a homomorphism of $\delta$-functors
in $M$ is to say that
\begin{equation}\tag{12}\label{eq:V.12}
  g\circ dh'-dg'\circ h
\end{equation}
is a coboundary in $\Homop^{\bullet}(R(N),R(P))$.

Now
\begin{align*}
  dh'&=h'\circ d_1+(-1)^{i+1}d_2\circ h',\\
  dg'&=g'\circ d_2+(-1)^{j+1}d_3\circ g',
\end{align*}
with the obvious notation. Thus \emph{\hyperref[eq:V.12]{(12)}} becomes
\[
  g\circ h'\circ d_1+(-1)^{i+1}g\circ d_2\circ h'
  -g'\circ d_2\circ h-(-1)^{j+1}d_3\circ g'\circ h.
\]

On the other hand,\srcpage{64} since $h$ and $g$ are cycles,
\begin{align*}
  g\circ d_2&=(-1)^j d_3\circ g,\\
  d_2\circ h&=(-1)^i h\circ d_1.
\end{align*}
Consequently \emph{\hyperref[eq:V.12]{(12)}} is finally written as
\[
  d\!\left(g\circ h'+(-1)^{i+1}g'\circ h\right),
\]
which completes the proof.

Thus \emph{\hyperref[eq:V.1.3]{(1.3)}} and \emph{\hyperref[eq:V.1.4]{(1.4)}} give a homomorphism of $\delta$-functors
in $M$:
\begin{equation*}\tag{1.5}\label{eq:V.1.5}
  \Extop^i(N,M)\longrightarrow
  \Homop\!\left(\Extop^j(M,P),\Extop^{i+j}(N,P)\right).
\end{equation*}
% END COMPONENT SGA2-V-C02

% BEGIN COMPONENT SGA2-V-C03 :: unit_V_2_local_duality_theorem
\section*{2. The local duality theorem for a regular local ring}
\sgareaderlabel{V.2}{2}

Let $A$ be a regular local ring of dimension $r$, let $\mathfrak m$ be the
maximal ideal of $A$, and let $M$ be a finitely generated $A$-module. Set
\srcpage{64}\sgareaderlabel{pV2.4}{2.4}
\[
  H^i(M)=H^i_{\mathfrak m}(M)
  \qquad\text{(so }H^i(M)=\varinjlim_k
    \Extop^i(A/\mathfrak m^k,M)\text{)}.
\]
We saw in \emph{\hyperref[IV.5.4]{IV, 5.4}} that $I=H^r(A)$ is a dualizing module for $A$;
denote the associated dualizing functor by $D$. In \emph{\hyperref[eq:V.1.5]{(1.5)}}, put
$N=A/\mathfrak m^k$ and $P=A$. We then obtain a homomorphism of
$\delta$-functors in $M$,
\begin{equation}\tag{13}\label{eq:V.13}
  \varphi_k\colon \Extop^i(A/\mathfrak m^k,M)
  \longrightarrow
  \Homop\!\left(\Extop^{r-i}(M,A),
    \Extop^r(A/\mathfrak m^k,A)\right).
\end{equation}
Passing to the direct limit over $k$ gives a homomorphism of
$\delta$-functors
\begin{equation}\tag{14}\label{eq:V.14}
  \varphi\colon H^i(M)\longrightarrow
    D\!\left(\Extop^{r-i}(M,A)\right).
\end{equation}

\medskip
\noindent\textbf{Theorem 2.1 (Local duality theorem).}\sgareaderlabel{V.2.1}{2.1}
The preceding functorial homomorphism $\varphi$ is an isomorphism.

\medskip
\noindent\textit{Proof.}
If $i>r$,\srcpage{65} the right-hand side of \emph{\hyperref[eq:V.14]{(14)}} is trivially zero,
and the left-hand side is zero because
\[
  H^i(M)=\varinjlim_k \Extop^i(A/\mathfrak m^k,M),
\]
and each $\Extop^i(A/\mathfrak m^k,M)$ is zero (syzygy theorem).

If $i=r$, the preceding discussion shows that the two functors in $M$,
$H^r(M)$ and $D(\Homop(M,A))$, are right exact. Since $A$ is noetherian
and $M$ is finitely generated, it is enough to verify the isomorphism for
$M=A$, where it is immediate.

To show that $\varphi$ is a functorial isomorphism, it now suffices to
proceed by descending induction on $i$ and observe that every finitely
generated module admits a finite presentation, while for $i<r$ both sides
of \emph{\hyperref[eq:V.14]{(14)}} vanish when $M$ is finitely generated and free. This is
clear for the right-hand side. Since $H^i$ commutes with finite direct sums,
for the left-hand side it is enough to show that $H^i(A)=0$ for $i<r$.
This follows from \emph{\hyperref[III.3.4]{III, 3.4}}, since $\prof(A)=r$. \hfill$\square$
% END COMPONENT SGA2-V-C03

% BEGIN COMPONENT SGA2-V-C04 :: unit_V_3_1_opening_draft
\section*{3. Application to the structure of the $H^i(M)$}
\sgareaderlabel{V.3}{3}

\begin{sgatheorem}{Theorem 3.1}\sgareaderlabel{V.3.1}{3.1}
Let $A$ be a Noetherian local ring, let $D$ be a dualizing functor for $A$,
and let $M\ne 0$ be a finitely generated $A$-module of dimension $n$.
Then:\srcpage{65}
\begin{enumerate}
  \item[(i)] $H^i(M)=0$ if $i<0$ or if $i>n$.

  \item[(ii)] $\widehat{D(H^i(M))}$ is a finitely generated
  $\widehat A$-module of dimension at most $i$.

  \item[(iii)] $H^n(M)\ne 0$; moreover, if $A$ is complete, then
  $D(H^n(M))$ has dimension $n$ and
  \[
    \Assop\bigl(D(H^n(M))\bigr)
      =\{\,\mathfrak p\in\Assop(M)\mid \dim(A/\mathfrak p)=n\,\}.
  \]
\end{enumerate}
\end{sgatheorem}

\par\addvspace{0.8\baselineskip}\noindent\textit{Proof.}\ 
Let $I$ be the dualizing module associated with $D$. We know that
$\widehat I$ is a dualizing module for $\widehat A$. On the other hand,
\begin{align*}
  \widehat{H^i(M)}&=H^i(\widehat M),\\
  \widehat{D(H^i(M))}&=\Homop\bigl(H^i(\widehat M),\widehat I\bigr),\\
  \dim\widehat M&=\dim M.
\end{align*}
Consequently,\srcpage{66} we may suppose that $A$ is complete. By Cohen's
theorem, every complete local ring is a quotient of a regular local ring.
To reduce to that case, we need the following lemma.
% END COMPONENT SGA2-V-C04

% BEGIN COMPONENT SGA2-V-C05 :: unit_V_3_2_lemma_draft
\begin{sgalemma}{Lemma 3.2}\sgareaderlabel{V.3.2}{3.2}\srcpage{66}
Let $X$ (respectively $Y$) be a ringed space, let $X'$ (respectively $Y'$)
be a closed subspace of $X$ (respectively $Y$), and let
$f\colon X\to Y$ be a morphism of ringed spaces such that
$f^{-1}(Y')=X'$. Let $F$ be an $\mathcal O_X$-module. Denote by $A$
(respectively $B$) the ring $\Gamma(\mathcal O_X)$ (respectively
$\Gamma(\mathcal O_Y)$), and by $\overline f\colon B\to A$ the ring
homomorphism corresponding to $f$. There is a spectral sequence of
$B$-modules with initial term
\begin{equation}\tag{15}\label{eq:V.15}
  E_2^{p,q}=H_{Y'}^p\bigl(Y,R^q f_*(F)\bigr),
\end{equation}
converging to the $B$-module
$H_{X'}^{\bullet}(X,F)_{[\overline f]}$.
\end{sgalemma}

\begin{proof}
Let $\mathcal O_{Y,Y'}$ be the sheaf
$\mathcal O_Y\vert_{Y'}$, extended by zero outside $Y'$ (see Expos\'e I).
There is an isomorphism of $B$-modules
\begin{equation}\tag{16}\label{eq:V.16}
  \Homop\bigl(\mathcal O_{Y,Y'},f_*(F)\bigr)
  \simeq
  \Homop\bigl(f^*(\mathcal O_{Y,Y'}),F\bigr)_{[\overline f]}.
\end{equation}
Now
\begin{equation}\tag{17}\label{eq:V.17}
  f^*(\mathcal O_{Y,Y'})=\mathcal O_{X,X'}.
\end{equation}
Moreover, if $G$ is an injective $\mathcal O_X$-module, then $f_*(G)$ is an
injective $\mathcal O_Y$-module, at least when $f$ is flat; one can easily
reduce to this case by replacing $\mathcal O_X$, and so forth, by the constant
sheaves of rings $\mathbb Z$. Hence the spectral sequence of the composite
functor
\[
  F\longmapsto\Homop\bigl(\mathcal O_{Y,Y'},f_*(F)\bigr),
\]
whose initial term is
\[
  E_2^{p,q}=\Extop^p\bigl(Y;\mathcal O_{Y,Y'},R^q f_*(F)\bigr),
\]
converges, in view of \textup{\hyperref[eq:V.16]{(16)}} and \textup{\hyperref[eq:V.17]{(17)}}, to
\[
  \Extop^{\bullet}\bigl(X;\mathcal O_{X,X'},F\bigr)_{[\overline f]}.
\]
The lemma then\srcpage{67} follows from Expos\'e~I,
\hyperref[I.2.3bof]{Proposition~2.3~bis}.
\end{proof}
% END COMPONENT SGA2-V-C05

% BEGIN COMPONENT SGA2-V-C06 :: unit_V_3_change_of_rings_reduction
\srcpage{67}Now let $f\colon B\to A$ be a surjective homomorphism of local rings. Let
\[
  f\colon \Specop(A)\longrightarrow\Specop(B)
\]
be the corresponding morphism of affine schemes. Set
$X=\Specop(A)$ (respectively $X'=\{\mathfrak m_A\}$) and
$Y=\Specop(B)$ (respectively $Y'=\{\mathfrak m_B\}$), and let $M$ be an
$A$-module and $\widetilde M$ the corresponding $\mathcal O_X$-module. Since
$R^qf_*(\widetilde M)=0$ for $q>0$, the spectral sequence~\textup{\hyperref[eq:V.15]{(15)}}
degenerates, and \hyperref[V.3.2]{Lemma~3.2} gives an isomorphism of $B$-modules
\begin{equation}\tag{18}\label{eq:V.18}
  H_{\{\mathfrak m_B\}}^n\bigl(Y,f_*(\widetilde M)\bigr)
  \simeq
  H_{\{\mathfrak m_A\}}^n\bigl(X,\widetilde M\bigr)_{[f]},
\end{equation}
hence an isomorphism of $B$-modules
\begin{equation}\tag{19}\label{eq:V.19}
  H_{\mathfrak m_B}^n\bigl(M_{[f]}\bigr)
  \simeq
  H_{\mathfrak m_A}^n(M)_{[f]}.
\end{equation}
On the other hand, if $D_A$ (respectively $D_B$) is the dualizing functor for
$A$ (respectively $B$), then
\begin{equation}\tag{20}\label{eq:V.20}
  D_A(M)_{[f]}\simeq D_B\bigl(M_{[f]}\bigr).
\end{equation}
Finally, since there is an isomorphism of rings
\begin{equation}\tag{21}\label{eq:V.21}
  B/\Annop M_{[f]}\simeq A/\Annop M,
\end{equation}
it follows that the change of base rings under consideration changes
nothing. We may therefore suppose that $A$ is regular of dimension $r$.
% END COMPONENT SGA2-V-C06

% BEGIN COMPONENT SGA2-V-C07 :: unit_V_3_equivalence_conditions_a_c_draft
By \textup{\hyperref[V.2.1]{(2.1)}},\srcpage{67} we have
\begin{equation}\tag{22}\label{eq:V.22}
  D\bigl(H^i(M)\bigr)=\Extop^{r-i}(M,A).
\end{equation}
We shall\srcpage{68} prove the equivalence of the following properties:
\begin{enumerate}
\item[(a)] $\dim \Extop^j(M,A)\leq r-j$;
\item[(b)] for every $\mathfrak p\in X=\operatorname{Spec}(A)$ such that
  $\dim A_{\mathfrak p}<j$, one has
  $\Extop^j(M,A)_{\mathfrak p}=0$;
\item[(c)]
  $\codimop\bigl(\Supp(\Extop^j(M,A)),X\bigr)\geq j$.
\end{enumerate}

To prove \textup{(a)} $\Longrightarrow$ \textup{(b)}, let
$\mathfrak p\in X$ with $\dim A_{\mathfrak p}<j$. Then
$\dim A/\mathfrak p>r-j$, and hence, by \textup{(a)},
\[
  \Ann\bigl(\Extop^j(M,A)\bigr)\not\subset\mathfrak p,
\]
which implies that $\Extop^j(M,A)_{\mathfrak p}=0$. Let
$\mathfrak p\in\Supp(\Extop^j(M,A))$. Then
$\Extop^j(M,A)_{\mathfrak p}\neq0$, so by \textup{(b)} we have
$\dim A_{\mathfrak p}\geq j$. Therefore
\[
  \codimop\bigl(\Supp(\Extop^j(M,A)),X\bigr)
  =\inf\bigl\{\dim A_{\mathfrak p}\mid
  \mathfrak p\in\Supp(\Extop^j(M,A))\bigr\}\geq j,
\]
that is, \textup{(b)} $\Longrightarrow$ \textup{(c)}. Finally,
\textup{(c)} implies \textup{(a)} trivially.
% END COMPONENT SGA2-V-C07

% BEGIN COMPONENT SGA2-V-C08 :: unit_V_3_1_proof_items_i_ii_draft
\noindent\srcpage{68}Let us now prove the theorem.
\begin{enumerate}
\item[(i)]
Let $x=(x_1,\dots,x_r)$ be a system of parameters for $A$ such that
$x_i\in\Ann M$ for $i=1,\dots,r-n$. Let
$K^\bullet((x^k),M)$ be the Koszul complex. It is easy to see that the map
\[
  K^i((x^k),M)\longrightarrow K^i((x^{k'}),M)
\]
for $k<k'$ is zero if $i>n$. It follows that
\[
  H^i(M)=\varinjlim H^i((x^k),M)=0
\]
if $i>n$. On the other hand, it is trivial that $H^i(M)=0$ if $i<0$;
hence \textup{(i)} is proved.

\item[(ii)]
Since $A$ is regular, $\dim A_{\mathfrak p}<j$ implies that the global
homological dimension of $A_{\mathfrak p}$ is strictly less than $j$, and
hence
\[
  \Extop^j(M,A)_{\mathfrak p}
  =\Extop^j_{A_{\mathfrak p}}(M_{\mathfrak p},A_{\mathfrak p})=0.
\]
Thus we have proved \textup{(b)}, and consequently \textup{(a)}.
Assertion \textup{(ii)} then follows from \textup{\hyperref[eq:V.22]{(22)}} and
\textup{(a)}.
\end{enumerate}
% END COMPONENT SGA2-V-C08

% BEGIN COMPONENT SGA2-V-C09 :: unit_V_3_theorem31_proof_iii
\medskip
\noindent\textup{(iii)} There exists a
$\mathfrak p\in\Suppop(M)$ such that $\dim A_{\mathfrak p}=r-n$ and such that
$\Suppop(M_{\mathfrak p})=\{\mathfrak m A_{\mathfrak p}\}$. Since
$A_{\mathfrak p}$ is regular if $A$ is, we have
$\depthop A_{\mathfrak p}=r-n$, and therefore
\begin{equation}\tag{23}\label{eq:V.23}
  \Extop_A^{r-n}(M,A)_{\mathfrak p}
  =\Extop_{A_{\mathfrak p}}^{r-n}(M_{\mathfrak p},A_{\mathfrak p})\neq 0.
\end{equation}
Taking~\textup{\hyperref[eq:V.22]{(22)}} into account, this\srcpage{69} implies, on the one hand,
\[
  H^n(M)\neq 0,
\]
and, on the other hand,
\[
  \dim D\bigl(H^n(M)\bigr)\geq n,
\]
hence, by~\textup{(ii)},
\[
  \dim D\bigl(H^n(M)\bigr)=n.
\]
% END COMPONENT SGA2-V-C09

% BEGIN COMPONENT SGA2-V-C10 :: unit_V_3_theorem31_item_iii_completion
\medskip
Now let $Y=\Suppop(M)$. By~\textup{(i)}, we know that
$D\bigl(H^n(M')\bigr)=\Extop^{r-n}(M',A)$ is a left exact functor in $M'$ on
the category $(\mathcal C_Y)^\circ$. Hence there exists an $A$-module $H$ and
an isomorphism of functors in $M'$:
\[
  \Extop^{r-n}(M',A)=\Homop(M',H).
\]

Let $Y_i$, $i=1,\ldots,k$, be the irreducible components of $Y$ of maximum
dimension. We shall see that the assertion
$\Extop^{r-n}(M',A)\neq 0$ is equivalent to the assertion that there exists an
$i$ such that $\Suppop M'\supset Y_i$. Indeed, if
$\Suppop M'\supset Y_i$, then $\dim(M')=n$, and hence
$\Extop^{r-n}(M',A)\neq 0$.

If $\Suppop M'\not\supset Y_i$ for every $i=1,\ldots,k$, then
$\dim M'<n$, and
\[
  D\bigl(H^n(M')\bigr)=\Extop^{r-n}(M',A)=0.
\]

Since
$\Assop\bigl(\Extop^{r-n}(M,A)\bigr)=\Suppop M\cap\Assop(H)$, we see that
the last assertion of~\textup{(iii)} follows from the following lemma:
% END COMPONENT SGA2-V-C10

% BEGIN COMPONENT SGA2-V-C11 :: unit_V_3_3_lemma_and_theorem31_closure
\begin{sgalemma}{Lemma 3.3}\sgareaderlabel{V.3.3}{3.3}
Let $X=\Specop(A)$, let $Y$ be a closed subset of $X$, let
\[
  T\colon(\mathcal C_Y)^\circ\longrightarrow\Ab
\]
be a left exact functor, and let $Y_i$, $i=1,\ldots,k$, be a family of
irreducible components of $Y$ such that the assertion $T(M)=0$ is equivalent
to the assertion that, for every $i$, $\Suppop M\not\supset Y_i$. Then $T$ is
representable by a module $H$ such that
\[
  \Assop(H)=\bigcup_{i=1}^k\{y_i\},
\]
where $y_i$ is the generic point of $Y_i$, for $i=1,\ldots,k$.
\end{sgalemma}

\begin{proof}
For $y\in Y$, construct\srcpage{70} an $A$-module $M(y)$ such that
\[
  \Suppop\bigl(M(y)\bigr)=\overline{\{y\}}.
\]
Suppose that $y\neq y_i$ for every $i=1,\ldots,k$. Then
$Y_i\not\subset\Suppop\bigl(M(y)\bigr)$ for every $i=1,\ldots,k$, and hence
$T\bigl(M(y)\bigr)=0$. It follows that
\[
  \Assop\bigl(T(M(y))\bigr)
  =\Suppop\bigl(M(y)\bigr)\cap\Assop(H)=\varnothing,
\]
and therefore $y\notin\Assop(H)$. If $y=y_i$, then
$Y_i\subset\Suppop\bigl(M(y)\bigr)$, so $T\bigl(M(y)\bigr)\neq0$, and hence
\[
  \Assop\bigl(T(M(y))\bigr)
  =\Suppop\bigl(M(y)\bigr)\cap\Assop(H)\neq\varnothing.
\]
By the first part of the proof, this implies $y\in\Assop(H)$, which proves
the lemma.\qedhere
\end{proof}
% END COMPONENT SGA2-V-C11

% BEGIN COMPONENT SGA2-V-C12 :: unit_V_3_4_example_affine_complement_codimension
\begin{sgaexample}{Example 3.4}\sgareaderlabel{V.3.4}{3.4}
Let $A$ be a noetherian ring, let $X=\Specop(A)$, and let $Y$ be a closed
subset of $X$ such that $X-Y$ is affine. Then, for every irreducible component
$Y_{\alpha}$ of $Y$, one has
\[
  \codimop(Y_{\alpha},X)\leq 1.
\]
\end{sgaexample}

\begin{proof}
Consider $X$ as a prescheme over $X$. Let $y_{\alpha}\in Y_{\alpha}$ be a
generic point, and consider the morphism
\[
  \Specop(\mathcal O_{X,y_{\alpha}})\longrightarrow X.
\]
The affine scheme obtained by base-changing $X$ to
$\Specop(\mathcal O_{X,y_{\alpha}})$ is canonically isomorphic to
$\Specop(\mathcal O_{X,y_{\alpha}})$.

By \textup{(EGA I 3.2.7)}, if $y_0$ is the unique closed point of
\[
  Y_0=\Specop(\mathcal O_{X,y_{\alpha}}),
\]
then $Y_0-y_0$ is affine. By \textup{(EGA III 1.3.1)}, one obtains
\[
  \HH^i(Y_0-y_0,\mathcal O_{Y_0})=0\qquad\text{if }i>0.
\]
Hence,\srcpdfpage{62 / running 54} by \textup{(\hyperref[I.2.9]{I 2.9)}},
\[
  \HH^{i-1}(\mathcal O_{X,y_{\alpha}})
  =\HH^i_{\{y_0\}}(Y_0,\mathcal O_{Y_0})=0
  \qquad\text{if }i\geq 2.
\]
By \hyperref[V.3.1]{3.1\textup{(iii)}}, it follows that
\[
  \dim\mathcal O_{X,y_{\alpha}}\leq 1,
\]
and therefore
\[
  \codimop(Y_{\alpha},X)
  =\inf_{y\in Y_{\alpha}}\dim\mathcal O_{X,y}\leq 1.
\]
\end{proof}
% END COMPONENT SGA2-V-C12

% BEGIN COMPONENT SGA2-V-C13 :: unit_V_3_5_finite_length_localization
Let\srcpage{71} $A$ be a noetherian local ring, let $\mathfrak m$ be its
maximal ideal, and let $M$ be a finitely generated $A$-module. Suppose that
$A$ is a quotient of a regular local ring. Set $X=\Specop(A)$ and, for every
$x\in X$, set $\mathfrak m_x=\mathfrak m A_x$.

\begin{sgaproposition}\sgareaderlabel{V.3.5}{3.5}
The following two conditions are equivalent:
\begin{enumerate}
  \item[(a)] $\HH^i(M)$ has finite length;
  \item[(b)] for every $x\in X-\{\mathfrak m\}$,
  \[
    \HH_{\mathfrak m_x}^{\,i-\dim\overline{\{x\}}}(M_x)=0.
  \]
\end{enumerate}
\end{sgaproposition}

\begin{proof}
\renewcommand{\qedsymbol}{}
In view of \textup{\hyperref[V.3.2]{(3.2)}}, we may suppose that $A$ is regular. By
\textup{\hyperref[V.2.1]{(2.1)}}, we have
\[
  \HH^i(M)=\DD\bigl(\Extop^{r-i}(M,A)\bigr),
\]
where $r=\dim A$. By \textup{(\hyperref[IV.4.7]{IV 4.7)}}, condition \textup{(a)} is
equivalent\footnote{\sgareaderlabel{dualnul}{*}Indeed, one must show that, if $E$ is a finitely generated
$A$-module and $\DD(E)$ has finite length, then $E$ has finite length. Let $K$
(respectively $Q$) be the kernel (respectively cokernel) of the canonical
morphism $\epsilon\colon E\to\DD\DD(E)$. The canonical morphism
$\gamma\colon\DD(E)\to\DD\DD\DD(E)$ satisfies
$\DD(\epsilon)\circ\gamma=\operatorname{id}_{\DD(E)}$. Since $\DD(E)$ has
finite length, $\gamma$ is an isomorphism, and hence so is
$\DD(\epsilon)$. Since $\DD$ is exact, it follows that $\DD(K)$ and
$\DD(Q)$ are zero. It is enough to prove that an $A$-module whose dual is zero
is itself zero, for then $E=\DD\DD(E)$ has finite length, just as $\DD(E)$
does. Indeed, let $M_0$ be a finitely generated submodule of such a module
$M$. Since $\DD$ is exact, $\DD(M_0)$ is a quotient of $\DD(M)$ and is
therefore zero. Again by exactness,
$\DD(M_0/\mathfrak m_A M_0)=0$; by biduality, the finite-length module
$M_0/\mathfrak m_A M_0$ is therefore zero. Nakayama's lemma then gives
$M_0=0$, and finally $M=0$.} to
\begin{equation}
  \Extop^{r-i}(M,A)\text{ has finite length}. \tag{24}\label{eq:V.24}
\end{equation}
Condition \textup{\hyperref[eq:V.24]{(24)}} is equivalent to
\begin{equation}
  \text{for every }x\in X-\{\mathfrak m\},\qquad
  \Extop_A^{r-i}(M,A)_x=0. \tag{25}\label{eq:V.25}
\end{equation}
On the other hand, $A_x$ is regular of dimension
$r-\dim\overline{\{x\}}$. Hence, by \textup{\hyperref[V.2.1]{(2.1)}},
\begin{align}
  \HH_{\mathfrak m_x}^{\,i-\dim\overline{\{x\}}}(M_x)
  &={\DD}\!\left(
    \Extop_{A_x}^{(r-\dim\overline{\{x\}})
      -(i-\dim\overline{\{x\}})}(M_x,A_x)
  \right) \notag\\
  &={\DD}\!\left(\Extop_{A_x}^{r-i}(M_x,A_x)\right). \tag{26}\label{eq:V.26}
\end{align}
Since $M$ is finitely generated, one has
\[
  \Extop_A^{r-i}(M,A)_x
  =\Extop_{A_x}^{r-i}(M_x,A_x),
\]
which proves the proposition.
\end{proof}
% END COMPONENT SGA2-V-C13

% BEGIN COMPONENT SGA2-V-C14 :: unit_V_3_6_corollary_exposeV_close
\begin{sgacorollary}\sgareaderlabel{V.3.6}{3.6}
For $\HH^i(M)$ to have finite length for $i\leq n$, it is necessary and
sufficient that
\[
  \depthop(M_x)>n-\dim\overline{\{x\}}
\]
for every $x\in X-\{\mathfrak m\}$.
\end{sgacorollary}

\begin{proof}
\renewcommand{\qedsymbol}{}
This follows from \textup{\hyperref[V.3.5]{(3.5)}} and \textup{(\hyperref[III.3.1]{III 3.1)}}.
\end{proof}
% END COMPONENT SGA2-V-C14

% END INCLUDED active_source/parts/05_expose_V.tex

% BEGIN INCLUDED active_source/parts/06_expose_VI.tex
\makeatletter
\providecommand{\sgareaderlabel}[2]{%
  \begingroup
  \def\@currentlabel{#2}%
  \phantomsection\label{#1}%
  \endgroup}
\makeatother

\chapter[\texorpdfstring{The Functors $\Ext_Z^{\bullet}$ and $\SheafExt_Z^{\bullet}$}{The Functors Ext over Z and sheaf Ext over Z (all degrees)}]
{The Functors $\Ext_Z^{\bullet}(X;F,G)$ and
$\SheafExt_Z^{\bullet}(F,G)$}
\label{VI}

\section*{1. Generalities}
\sgareaderlabel{VI.1}{1}
\subsection*{1.1}
\sgareaderlabel{VI.1.1}{1.1}
Let $(X,\OX)$ be a ringed space and let $Z$ be a locally closed subset of
$X$. Let $F$ and $G$ be $\OX$-modules. We shall denote by
$\Ext_Z^i(X;F,G)$ (respectively $\SheafExt_Z^i(F,G)$) the $i$th derived
functor of
\[
 G\longmapsto \Gamma_Z\bigl(\SheafHom_{\OX}(F,G)\bigr)
 \quad\text{(respectively }G\longmapsto
 \SheafGamma_Z\bigl(\SheafHom_{\OX}(F,G)\bigr)\text{)}.
\]

\begin{sgastatement}{Lemma 1.2}\sgareaderlabel{VI.1.2}{1.2}
The sheaf $\SheafExt_Z^i(F,G)$ is canonically isomorphic to the sheaf
associated with the presheaf
\[
 U\longmapsto \Ext_{Z\cap U}^i\bigl(U;F|_U,G|_U\bigr).
\]
\end{sgastatement}

This follows from \textup{([T\^ohoku], 3.7.2)} and from the fact that
\[
 \Gamma\!\left(U;\SheafGamma_Z\bigl(\SheafHom_{\OX}(F,G)\bigr)\right)
 \simeq
 \Gamma_{Z\cap U}\!\left(\SheafHom_{\OX|_U}(F|_U,G|_U)\right)
\]
canonically.

\begin{sgastatement}{\hyperref[VI.1.3]{Theorem 1.3} (Excision theorem)}
\sgareaderlabel{VI.1.3}{1.3}
Let $V$ be an open subset of $X$ containing $Z$. Then there is an
isomorphism of cohomological functors
\begin{equation}
 \Ext_Z^{\bullet}(X;F,G)
 \simeq
 \Ext_Z^{\bullet}\bigl(V;F|_V,G|_V\bigr).
 \tag{1.3.1}\label{VI.1.3.1}
\end{equation}
\end{sgastatement}

Indeed, if $G^{\bullet}$ is an injective resolution of $G$, then
$G^{\bullet}|_V$ is an injective resolution of $G|_V$. The theorem follows
immediately.

\subsection*{1.4}
\sgareaderlabel{VI.1.4}{1.4}
Let $\OoXZ$ be the $\OX$-module defined by the following conditions
\textup{([Godement], 2.9.2)}:
\[
 \OoXZ|_{X-Z}=0
 \qquad\text{and}\qquad
 \OoXZ|_Z=\OX|_Z.
\]
We have seen that, for every $\OX$-module $H$, there is a functorial
isomorphism
\[
 \Gamma_Z(H)\simeq \Homop_{\OX}(\OoXZ,H).
\]
Consequently, there are isomorphisms functorial in $F$ and $G$:
\begin{align}
 \Gamma_Z\bigl(\SheafHom_{\OX}(F,G)\bigr)
  &\simeq \Homop_{\OX}\bigl(\OoXZ,\SheafHom_{\OX}(F,G)\bigr),
   \tag{1.4.1}\label{VI.1.4.1}\\[-0.15em]
 \Gamma_Z\bigl(\SheafHom_{\OX}(F,G)\bigr)
  &\simeq \Homop_{\OX}\bigl(\OoXZ\otimes_{\OX}F,G\bigr),
   \tag{1.4.2}\label{VI.1.4.2}\\[-0.15em]
 \Gamma_Z\bigl(\SheafHom_{\OX}(F,G)\bigr)
  &\simeq \Homop_{\OX}\bigl(F,\SheafHom_{\OX}(\OoXZ,G)\bigr)
   =\Homop_{\OX}\bigl(F,\SheafGamma_Z(G)\bigr).
   \tag{1.4.3}\label{VI.1.4.3}
\end{align}
In particular, it follows from~\eqref{VI.1.4.2} that there is an
isomorphism that is $\partial$-functorial in $F$ and $\delta$-functorial in $G$,
\[
 \theta\colon
 \Ext_{\OX}^{i}\bigl(\OoXZ\otimes_{\OX}F,G\bigr)
 \xrightarrow{\ \sim\ }
 \Ext_Z^{i}(X;F,G).
\]

\subsection*{1.5}
\sgareaderlabel{VI.1.5}{1.5}
By definition, the functor
$G\longmapsto\Gamma_Z\bigl(\SheafHom_{\OX}(F,G)\bigr)$ is the composite of
$G\longmapsto\SheafHom_{\OX}(F,G)$ and $\Gamma_Z$. Since $\Gamma_Z$ is left
exact \textup{(\hyperref[I.1.9]{I 1.9)}}, since $\SheafHom_{\OX}(F,G)$ is flasque when $G$ is
injective, and since $\Gamma_Z$ is exact on flasque sheaves \textup{(\hyperref[I.2.12]{I 2.12)}},
it follows from \textup{([T\^ohoku], 2.4.1)} that there is a spectral functor
abutting to $\Ext_Z^{\bullet}(X;F,G)$ whose initial term is
\[
 H_Z^p\bigl(X,\SheafExt_{\OX}^q(F,G)\bigr).
\]

On the other hand, it follows from~\eqref{VI.1.4.3} that
$\Gamma_Z\bigl(\SheafHom_{\OX}(F,G)\bigr)$ is the composite of
$G\longmapsto\SheafGamma_Z(G)$ and
$H\longmapsto\Homop_{\OX}(F,H)$. Since $\SheafGamma_Z$ takes injectives to
injectives \textup{(\hyperref[I.1.4]{I 1.4)}}, it follows from \textup{([T\^ohoku], 2.4.1)}
that there is a spectral functor abutting to $\Ext_Z^{\bullet}(X;F,G)$ whose
initial term is
\[
 \Ext_{\OX}^p\bigl(X;F,\SheafH_Z^q(G)\bigr).
\]

Finally, it follows from~\eqref{VI.1.4.2} and the $\Ext$ spectral sequence
that there is a spectral functor abutting to $\Ext_Z^{\bullet}(X;F,G)$ whose
initial term is $H^p\bigl(X;\SheafExt_Z^{\bullet}(F,G)\bigr)$. Hence:

\begin{sgastatement}{Theorem 1.6}\sgareaderlabel{VI.1.6}{1.6}
There are three spectral functors abutting to
$\Ext_Z^{\bullet}(X;F,G)$ whose initial terms are respectively
\begin{align}
 &H_Z^p\bigl(X,\SheafExt_{\OX}^q(F,G)\bigr),
   \tag{1.6.1}\label{VI.1.6.1}\\[-0.15em]
 &H^p\bigl(X,\SheafExt_Z^q(F,G)\bigr),
   \tag{1.6.2}\label{VI.1.6.2}\\[-0.15em]
 &\Ext_{\OX}^p\bigl(X;F,\SheafH_Z^q(G)\bigr).
   \tag{1.6.3}\label{VI.1.6.3}
\end{align}
\end{sgastatement}

\subsection*{1.7}
\sgareaderlabel{VI.1.7}{1.7}
Let $Z'$ be a closed subset of $Z$, and set $Z''=Z-Z'$. There is an exact
sequence
\begin{equation}
 0\longrightarrow \OXZ{Z''}\longrightarrow \OXZ{Z}
 \longrightarrow \OXZ{Z'}\longrightarrow 0.
 \tag{1.7.1}\label{VI.1.7.1}
\end{equation}
It generalizes the exact sequence of \textup{([Godement], 2.9.3)}. This exact
sequence splits locally; hence, for every $\OX$-module $F$, there is another
exact sequence
\begin{equation}
 0\longrightarrow F\otimes_{\OX}\OXZ{Z''}
 \longrightarrow F\otimes_{\OX}\OXZ{Z}
 \longrightarrow F\otimes_{\OX}\OXZ{Z'}\longrightarrow 0.
 \tag{1.7.2}\label{VI.1.7.2}
\end{equation}
Now let $G$ be an $\OX$-module. Applying the functor
$\Homop_{\OX}(\mathbin{\bullet},G)$ to \hyperref[VI.1.7.2]{(1.7.2)}, one obtains from \hyperref[VI.1.4.2]{(1.4.2)} and
the long exact sequence of $\Ext$ the following theorem.

\begin{sgastatement}{Theorem 1.8}\sgareaderlabel{VI.1.8}{1.8}
Let $Z$ be a locally closed subset of $X$, let $Z'$ be a closed subset of
$Z$, and set $Z''=Z-Z'$. There is an exact sequence, functorial in $F$ and
$G$,
\[
\begin{aligned}
 0&\longrightarrow \Homop_{Z'}(F,G)\longrightarrow \Homop_Z(F,G)
 \longrightarrow \Homop_{Z''}(F,G)\longrightarrow \Ext_{Z'}^1(F,G)
 \longrightarrow\cdots \\
 &\cdots\longrightarrow \Ext_Z^i(F,G)\longrightarrow \Ext_{Z''}^i(F,G)
 \longrightarrow \Ext_{Z'}^{i+1}(F,G)\longrightarrow\cdots .
\end{aligned}
\]
\end{sgastatement}

\begin{sgastatement}{Corollary 1.9}\sgareaderlabel{VI.1.9}{1.9}
Let $Y$ be a closed subset of $X$, and set $U=X-Y$. There is an exact
sequence, functorial in $F$ and $G$,
\[
\begin{aligned}
 0&\longrightarrow \Homop_Y(F,G)\longrightarrow \Homop_{\OX}(F,G)
 \longrightarrow \Homop_{\OX|_U}(F|_U,G|_U)
 \longrightarrow \Ext_Y^1(F,G)\longrightarrow\cdots \\
 &\cdots\longrightarrow \Ext_{\OX}^i(F,G)
 \longrightarrow \Ext_{\OX|_U}^i(F|_U,G|_U)
 \longrightarrow \Ext_Y^{i+1}(F,G)\longrightarrow\cdots .
\end{aligned}
\]
\end{sgastatement}

This corollary is an immediate consequence of \hyperref[VI.1.3]{Theorems (1.3)} and \hyperref[VI.1.8]{(1.8)}.

\section*{2. Applications to Quasi-Coherent Sheaves on Preschemes}
\sgareaderlabel{VI.2}{2}

\begin{sgastatement}{Proposition 2.1}\sgareaderlabel{VI.2.1}{2.1}
Let $X$ be a locally Noetherian prescheme. For every locally closed subset
$Z$ of $X$, every coherent $\OX$-module $F$, and every quasi-coherent
$\OX$-module $G$, the sheaves $\SheafExt_Z^i(F,G)$ are quasi-coherent.
\end{sgastatement}

As in \hyperref[VI.1.6.3]{(1.6.3)}, one shows that the sheaves $\SheafExt_Z^i(F,G)$ are obtained
as the abutment of a spectral sequence whose initial term is
\[
 \SheafExt_{\OX}^p\bigl(F,\SheafH_Z^q(G)\bigr).
\]
By \textup{(II, \hyperref[II.3]{Corollary 3})}, the $\SheafH_Z^q(G)$ are quasi-coherent, and
therefore so are the $\SheafExt_{\OX}^p(F,\SheafH_Z^q(G))$, since $F$ is
coherent. The proposition follows immediately.

\subsection*{2.2}
\sgareaderlabel{VI.2.2}{2.2}
Let $Y$ now be a closed subprescheme of $X$, and let $\SheafI$ be an ideal
defining $Y$. Let $m$ and $n$ be integers such that $m\geq n\geq 0$. We
denote by $i_{n,m}$ the canonical map
\[
 \OY{m}=\OX/\SheafI^{m+1}\longrightarrow
 \OX/\SheafI^{n+1}=\OY{n}
\]
and by $j_n$ the map $\OXY\to\OY{n}$. The
$(\OY{n},i_{n,m})$ form a projective system, and the $j_n$ are compatible
with the $i_{n,m}$.

Applying the functor $\Ext_{\OX}^{i}(F\otimes\mathord\bullet,G)$, we obtain a
morphism
\[
 \varphi'\colon
 \varinjlim_n \Ext_{\OX}^{i}(X;F\otimes\OY{n},G)
 \longrightarrow \Ext_{\OX}^{i}(X;F\otimes\OXY,G);
\]
this is a morphism of cohomological functors in $G$. The morphism
\[
 \varphi\colon
 \varinjlim_n \Ext_{\OX}^{i}(X;F\otimes\OY{n},G)
 \longrightarrow \Ext_Y^{i}(X;F,G),
\]
which is the composite $\theta\circ\varphi'$ (cf.~\hyperref[VI.1.4]{(1.4)}), is therefore also a
morphism of cohomological functors in $G$.

In the same way one defines
\[
 \Sheafphi\colon
 \varinjlim_n \SheafExt_{\OX}^{i}(F\otimes\OY{n},G)
 \longrightarrow \SheafExt_Y^{i}(F,G).
\]

\begin{sgastatement}{Theorem 2.3}\sgareaderlabel{VI.2.3}{2.3}
Let $X$ be a locally Noetherian prescheme, let $Y$ be a closed subset of $X$
defined by a coherent ideal $\SheafI$, let $F$ be a coherent $\OX$-module,
and let $G$ be a quasi-coherent $\OX$-module. Then:
\begin{enumerate}
 \renewcommand{\labelenumi}{\alph{enumi})}
 \item $\Sheafphi$ is an isomorphism.
 \item If $X$ is Noetherian, $\varphi$ is an isomorphism.
\end{enumerate}
\end{sgastatement}

\subsection*{Proof of \hyperref[VI.2.3]{Theorem 2.3}}
The proof of b) is almost word for word the same as that of
\textup{(\hyperref[II.6]{II 6} b))}, using the spectral sequence \hyperref[VI.1.6.2]{(1.6.2)}, so we do not reproduce
it.

For the proof of a), by \hyperref[VI.2.1]{(2.1)} we may assume that $X$ is affine with ring $A$,
that $F$ (respectively $G$) is defined by an $A$-module $M$ (respectively
$N$), and that $\SheafI$ is defined by an ideal $I$. It is enough to prove
that the homomorphism
\begin{equation}
 \varinjlim_n \Ext_A^i(M/I^nM,N)
 \longrightarrow \Ext_Y^i(X,F,G)
 \tag{2.3.1}\label{VI.2.3.1}
\end{equation}
induced by $\Sheafphi$ is an isomorphism.

Indeed, for $i=0$, the two sides of \hyperref[VI.2.3.1]{(2.3.1)} may be canonically identified with
the submodule of $\Homop_A(M,N)$ consisting of the elements of
$\Homop_A(M,N)$ annihilated by a power of $I$. One then sees that the
homomorphism \hyperref[VI.2.3.1]{(2.3.1)} is simply the identity map.

The functor
\[
 N\longmapsto \varinjlim_n \Ext_A^{\bullet}(M/I^nM,N)
\]
is a universal $\delta$-functor. We shall show that the same holds for the
functor $N\longmapsto \Ext_Y^{\bullet}(M,N)$. Indeed, if $N$ is an injective
module, then by \textup{(\hyperref[II.9]{II, 9} and \hyperref[II.11]{11})}, $\SheafH_Y^q(N)=0$ for $q\ne0$; and
by \textup{(\hyperref[IV.2.2]{IV.2.2)}}, $H_Y^0(N)$ is injective.

It follows that
\[
 \Ext_{\OX}^{p}\bigl(X;M,\SheafH_Y^q(N)\bigr)=0
 \qquad\text{if }p+q\ne0.
\]
Thus, by \hyperref[VI.1.6.3]{(1.6.3)}, $\Ext_Y^i(M,N)=0$ for $i\ne0$ when $N$ is injective. This
completes the proof.

\section*{Bibliography}
The same references as those listed at the end of Expos\'e I, cited
respectively as [T\^ohoku] and [Godement].

% END INCLUDED active_source/parts/06_expose_VI.tex

% BEGIN INCLUDED active_source/parts/07_expose_VII.tex
\makeatletter
\providecommand{\sgareaderlabel}[2]{%
  \begingroup
  \def\@currentlabel{#2}%
  \phantomsection\label{#1}%
  \endgroup}
\makeatother

\chapter[Vanishing and Coherence Criteria]
{Vanishing Criteria and Coherence Conditions for the Sheaves
$\SheafExt_Y^i(F,G)$}
\label{VII}

% BEGIN COMPONENT SGA2-VII-CUM-C01
\section*{1. Study for $i<n$}
\sgareaderlabel{VII.1}{1}

We first prove a lemma.

\begin{sgastatement}{Lemma 1.1}\sgareaderlabel{VII.1.1}{1.1}
Let $X$ be a locally Noetherian prescheme, let $Y$ be a closed subset of $X$,
and let $G$ be a quasi-coherent $\OX$-module. Suppose that, for every coherent
$\OX$-module $F$ whose support is contained in $Y$, one has
\[
  \SheafExt^{n-1}(F,G)=0.
\]
Then, for every coherent $\OX$-module $F$ and every closed subset $Z$ of $X$
such that $Z\cap\Supp F=Y\cap\Supp F$, one has an isomorphism
\[
  \SheafExt_Z^n(F,G)\cong
  \SheafHom\bigl(F,\SheafH_Y^n(G)\bigr).
\]
\end{sgastatement}
% END COMPONENT SGA2-VII-CUM-C01

% BEGIN COMPONENT SGA2-VII-CUM-C02
\subsection*{Proof of \hyperref[VII.1.1]{Lemma 1.1}}

We first observe that
\[
  \SheafExt_Z^i(F,G)
  =\SheafExt_{Z\cap\Supp F}^i(F,G).
\]
(This is immediate; cf. Expos\'e VI.) We first prove the assertion for
$Z=X$, so that $\Supp F\subset Y$. On the category of coherent $\OX$-modules
whose support is contained in $Y$, the functor
\[
  F\longmapsto\SheafExt^n(F,G)
\]
is left exact. By (\hyperref[IV.1.3]{IV, 1.3)}, it is represented by
\[
  I=\varinjlim_k\SheafExt^n
    \bigl(\OX/\mathcal I^{k+1},G\bigr),
\]
where $\mathcal I$ is a defining ideal of $Y$. By (\hyperref[II.6]{II, 6)}, one has
\[
  \SheafH_Y^n(G)\cong
  \varinjlim_k\SheafExt^n
    \bigl(\OX/\mathcal I^{k+1},G\bigr).
\]
This proves the assertion when $Z=X$. Again by (\hyperref[VI.2.3]{VI, 2.3)},
\[
  \SheafExt_Z^n(F,G)\cong
  \varinjlim_k\SheafExt^n
    \bigl(F/\mathcal J^{k+1}F,G\bigr),
\]
where $\mathcal J$ is a defining ideal of $Z$. The support of
$F/\mathcal J^{k+1}F$ is contained in $Y$ whenever
$Z\cap\Supp F\subset Y$. What we have just proved therefore gives
\[
  \SheafExt_Z^n(F,G)\cong
  \varinjlim_k\SheafHom
    \bigl(F/\mathcal J^{k+1}F,\SheafH_Y^n(G)\bigr).
\]
It remains to show that the natural homomorphism
\[
  \varinjlim_k\SheafHom
    \bigl(F/\mathcal J^{k+1}F,\SheafH_Y^n(G)\bigr)
  \longrightarrow
  \SheafHom\bigl(F,\SheafH_Y^n(G)\bigr)
\]
is an isomorphism when $Z\cap\Supp F=Y\cap\Supp F$. Since $X$ admits a cover by
Noetherian affine open subsets, we reduce to the case where $X$ is affine
Noetherian. Then $F(X)$ is a finitely generated $\OX(X)$-module and
$\Supp\SheafH_Y^n(G)\subset Y$. Hence every homomorphism
\[
  u:F(X)\longrightarrow\SheafH_Y^n(G)(X)
\]
is annihilated by a power of $\mathcal I$, and therefore by a power of
$\mathcal J$. This proves the lemma.
% END COMPONENT SGA2-VII-CUM-C02

% BEGIN COMPONENT SGA2-VII-CUM-C03
\begin{sgastatement}{Proposition 1.2}\sgareaderlabel{VII.1.2}{1.2}
Let $X$ be a locally Noetherian prescheme, let $Y$ be a closed subset of $X$,
let $G$ be a quasi-coherent $\OX$-module, and let $n$ be an integer. For any
closed subsets $Z$ and $S$ of $X$ such that $Z\cap S=Y$, the following
conditions are equivalent:
\begin{enumerate}
  \renewcommand{\labelenumi}{(\roman{enumi})}
  \item $\SheafH_Y^i(G)=0$ for $i<n$;
  \item there exists a coherent $\OX$-module $F$ with support $S$ such that
  \[
    \SheafExt_Z^i(F,G)=0\qquad\text{for }i<n;
  \]
  \item for every coherent $\OX$-module $F$ whose support is contained in $S$
  (that is, $\Supp F\cap Z=\Supp F\cap Y$), one has
  \[
    \SheafExt_Z^i(F,G)=0\qquad\text{for }i<n;
  \]
  \item for every coherent $\OX$-module $F$, one has
  \[
    \SheafExt_Y^i(F,G)=0\qquad\text{for }i<n.
  \]
\end{enumerate}
Moreover, if these conditions hold, then for every coherent $\OX$-module $F$
and every closed subset $Z'$ of $X$ such that
$Z'\cap\Supp F=Y\cap\Supp F$, one has isomorphisms
\[
  \SheafExt_{Z'}^n(F,G)\cong
  \SheafExt_Y^n(F,G)\cong
  \SheafHom\bigl(F,\SheafH_Y^n(G)\bigr).
\]
\end{sgastatement}

\noindent\textit{Source note.} The quantified subset in the preceding
sentence is $Z'$. The prior maintained transcription and this English reader
dropped the prime from the first term of the final display. The primary scan
and the proof both read $Z'$. Canon correction SGA-SRC-001624 restores
the missing prime as a diplomatic transcription repair; this reader now
reproduces that scan-faithful reading.
% END COMPONENT SGA2-VII-CUM-C03

% BEGIN COMPONENT SGA2-VII-CUM-C04
\subsection*{Proof of \hyperref[VII.1.2]{Proposition 1.2}: induction}

We argue by induction. The proposition is immediate for $n<0$. Suppose that it
has been proved for $n<q$. If one of the conditions holds for $n=q$ and for two
subsets $Z$ and $S$ as in the statement, then the induction hypothesis gives,
for every closed subset $Z'$ of $X$ and every coherent $\OX$-module $F$ such
that $Z'\cap\Supp F=Y\cap\Supp F$, isomorphisms
\begin{equation}\tag{1.1}\label{eq:VII.1.1}
  \SheafExt_{Z'}^{q-1}(F,G)
  \cong\SheafHom\bigl(F,\SheafH_Y^{q-1}(G)\bigr)
  \cong\SheafExt_Y^{q-1}(F,G).
\end{equation}
Consequently:
\begin{itemize}
  \item[(i)$\Rightarrow$(iv)] take $Z'=Y$ in \hyperref[eq:VII.1.1]{(1.1)};
  \item[(iv)$\Rightarrow$(iii)] take $Z'=Z$ in \hyperref[eq:VII.1.1]{(1.1)};
  \item[(iii)$\Rightarrow$(ii)] take $F=\mathcal O_S$;
  \item[(ii)$\Rightarrow$(i)] take $Z'=Z$ in \hyperref[eq:VII.1.1]{(1.1)}. This gives
  $\SheafHom\bigl(F,\SheafH_Y^{q-1}(G)\bigr)=0$. Moreover,
  \[
    \Supp\SheafH_Y^{q-1}(G)
    \subset Y=Z\cap S\subset S=\Supp F,
  \]
  so the following lemma applies.
\end{itemize}

\begin{sgastatement}{Lemma 1.3}\sgareaderlabel{VII.1.3}{1.3}
Let $X$ be a prescheme, let $P$ be a coherent $\OX$-module, and let $H$ be a
quasi-coherent $\OX$-module such that
\[
  \SheafHom(P,H)=0
  \qquad\text{and}\qquad
  \Supp P\supset\Supp H.
\]
Then $H=0$.
\end{sgastatement}
% END COMPONENT SGA2-VII-CUM-C04

% BEGIN COMPONENT SGA2-VII-CUM-C05
\subsection*{Proof of \hyperref[VII.1.3]{Lemma 1.3}}
It suffices to prove the lemma when $X$ is affine, since the affine open
subsets form a basis for the topology of $X$ and the hypotheses are preserved
under restriction to an open subset. In that case, we are reduced to a problem
about $A$-modules, where $X=\Spec(A)$. We apply the formula (valid under the
sole hypothesis that $\sisi{M}{P}$ is of finite type):
\[
  \Assop\Homop_A(P,H)=\Supp P\cap\Assop H.
\]
We know that
\[
  \Assop H\subset\Supp H\subset\Supp P
  \qquad\text{and}\qquad
  \Assop\Homop_A(P,H)=\varnothing;
\]
hence $\Assop H=\varnothing$, and therefore $H=0$.

\noindent\textit{Source note.} The printed French text reads $M$ in the
finite-type hypothesis. The diplomatic French and source apparatus preserve
$M$; corrected French and this English edition use the accepted reading $P$,
the module in the displayed formula.

\subsection*{Close of the proof of \hyperref[VII.1.2]{Proposition 1.2}}
To complete the proof of the proposition, it remains to observe that
condition~(iv) permits the application of \hyperref[VII.1.1]{Lemma~1.1}.
% END COMPONENT SGA2-VII-CUM-C05

% BEGIN COMPONENT SGA2-VII-CUM-C06
\begin{sgastatement}{Corollary 1.4}\sgareaderlabel{VII.1.4}{1.4}
Let $G$ be a coherent Cohen--Macaulay $\OX$-module, and let $n\in\mathbb Z$.
The conditions of \hyperref[VII.1.2]{Proposition~1.2} are equivalent to:
\[
  \textup{(v)}\qquad
  \codimop(Y\cap\Supp G,\Supp G)\geq n.
\]
\end{sgastatement}

\subsection*{Proof}
Recall first that an $\OX$-module is said to be Cohen--Macaulay if, for every
$x\in X$, the stalk $G_x$ is a Cohen--Macaulay $\Oo_{X,x}$-module; that is,
for every $x\in S=\Supp G$,
\begin{equation}
  \depthop G_x=\dim G_x=\dim\Oo_{S,x}.
  \tag{1.2}\label{eq:VII.1.2}
\end{equation}
By Proposition~\hyperref[III.3.3]{III.3.3}, condition~(i) of \hyperref[VII.1.2]{Proposition~1.2} is equivalent to
\begin{equation}
  \depthop_Y G=\inf_{x\in Y}\depthop G_x\geq n,
  \tag{1.3}\label{eq:VII.1.3}
\end{equation}
and hence also to
\[
  \depthop_Y G=\inf_{x\in Y\cap S}\depthop G_x\geq n,
\]
since the depth of the zero module is infinite.

By definition,
\[
  \codimop(Y\cap S,S)=\inf_{x\in S\cap Y}\dim\Oo_{S,x};
\]
the conclusion follows by applying formula~\eqref{eq:VII.1.2}.
% END COMPONENT SGA2-VII-CUM-C06

% BEGIN COMPONENT SGA2-VII-CUM-C07
We shall now prove a result that will allow us to deduce the coherence
conditions we have in view from certain vanishing criteria.

\begin{sgastatement}{Lemma 1.5}\sgareaderlabel{VII.1.5}{1.5}
Let $X$ be a locally Noetherian prescheme. Let $T^{\ast}$ be an exact
contravariant $\partial$-functor defined on the category of coherent
$\OX$-modules and taking values in the category of $\OX$-modules. Let $Y$ be a
closed subset of $X$, and let $i\in\mathbb Z$. Suppose that, for every coherent
$\OX$-module $F$ whose support is contained in $Y$, the modules $T^iF$ and
$T^{i-1}F$ are coherent. Let $F$ be a coherent $\OX$-module. Then $T^iF$ is
coherent if and only if $T^iF''$ is coherent, where
\[
  F''=F/\sisi{\Gamma}{\SheafGamma}_Y(F).
\]
\end{sgastatement}

Indeed, $F'=\sisi{\Gamma}{\SheafGamma}_Y(F)$ is coherent because $X$ is locally Noetherian. The
long exact cohomology sequence of $T^{\ast}$ then gives
\[
  T^{i-1}F'\longrightarrow T^iF''\longrightarrow T^iF
  \longrightarrow T^iF',
\]
where the two outer terms are coherent; the conclusion follows.
% END COMPONENT SGA2-VII-CUM-C07

% BEGIN COMPONENT SGA2-VII-CUM-C08
\begin{sgastatement}{Lemma 1.6}\sgareaderlabel{VII.1.6}{1.6}
If $F$ and $G$ are coherent and $\Supp F\subset Y$, then
$\SheafExt_Y^i(F,G)$ is coherent.
\end{sgastatement}

Indeed, $\SheafExt_Y^i(F,G)$ is isomorphic to $\SheafExt^i(F,G)$. More
generally, this holds on any ringed space $X$: if $Z$ is a closed subset
containing $Y\cap\Supp F$, then $\SheafExt_Z^i(F,G)$ is isomorphic to
$\SheafExt_Y^i(F,G)$ (cf. Expos\'e~VI).

\begin{sgastatement}{Proposition 1.7}\sgareaderlabel{VII.1.7}{1.7}
Suppose that $F$ and $G$ are coherent, and put
\[
  \Supp F=S,\qquad S'=\overline{S\cap(X-Y)}.
\]
Suppose that $\depthop G_x\geq n$ for every $x\in Y\cap S'$. Then
$\SheafExt_Y^i(F,G)$ is coherent for $i<n$.
\end{sgastatement}

Indeed, \hyperref[VII.1.6]{Lemma~1.6} allows us to apply \hyperref[VII.1.5]{Lemma~1.5} to
\[
  T^{\ast}(F)=\SheafExt_Y^{\ast}(F,G).
\]
Putting $F''=F/\sisi{\Gamma}{\SheafGamma}_Y(F)$, we see that $\Supp F''=S'$. By
Proposition~\hyperref[III.3.3]{III.3.3}, the hypothesis on the depth of $G$ ensures that
$\SheafH^i_{Y\cap S'}(G)=0$ for $i<n$; \hyperref[VII.1.2]{Proposition~1.2} then gives
$T^iF''=0$ for $i<n$, and the conclusion follows from \hyperref[VII.1.5]{Lemma~1.5}.
% END COMPONENT SGA2-VII-CUM-C08

% BEGIN COMPONENT SGA2-VII-CUM-C09
\section*{2. Study for $i>n$}
\sgareaderlabel{VII.2}{2}

Let $X$ be a locally Noetherian \emph{regular} prescheme; that is, suppose that
all its local rings are regular. Let $Y$ be a closed subset of $X$, and let
$F$ and $G$ be two coherent $\OX$-modules. Put
\[
  S=\Supp F,\qquad S'=\overline{S\cap(X-Y)},
\]
and put
\begin{align*}
  m&=\sup_{x\in Y\cap S}\dim\Oo_{X,x},\\
  n&=\sup_{x\in Y\cap S'}\dim\Oo_{X,x}.
\end{align*}
Then $n\leq m$.

\begin{sgastatement}{Proposition 2.1}\sgareaderlabel{VII.2.1}{2.1}
In the situation described above:
\begin{enumerate}
\item $\SheafExt_Y^i(F,G)=0$ if $i>m$;
\item $\SheafExt_Y^i(F,G)$ is coherent if $i>n$.
\end{enumerate}
\end{sgastatement}
% END COMPONENT SGA2-VII-CUM-C09

% BEGIN COMPONENT SGA2-VII-CUM-C10
\subsection*{Proof of \hyperref[VII.2.1]{Proposition 2.1}}

First note that $\SheafExt_Y^i(F,G)$ is coherent for every $i$ when
$\Supp F\subset Y$. Moreover, putting $F''=F/\sisi{\Gamma}{\SheafGamma}_Y(F)$ as above, we see
that $\Supp F''=S'$. Thus \hyperref[VII.2.1]{assertion~(2)} follows from \hyperref[VII.2.1]{assertion~(1)} and
\hyperref[VII.1.5]{Lemma~1.5}.\footnote{The French TeX and printed page cite Lemma~1.3. The
coherence reduction used here is Lemma~1.5; Lemma~1.3 is the earlier
$\mathcal Hom$/support lemma.}

To prove \hyperref[VII.2.1]{assertion~(1)}, first observe that
\[
  \SheafExt_Y^i(F,G)\cong
  \varinjlim_k \SheafExt^i(F/\Jj^kF,G),
\]
where $\Jj$ is the ideal defining $Y$. On the other hand, Theorem~4.2.2 of
A.~Grothendieck, ``Sur quelques points d'alg\`ebre homologique,''
\emph{T\^ohoku Mathematical Journal} \textbf{9} (1957), 119--221, shows that
the $\SheafExt$ sheaves commute with the formation of stalks, at least when
$X$ is a locally Noetherian prescheme and the first argument is coherent.
Since direct limits do likewise, for every $x\in X$ there are
isomorphisms\footnote{The French TeX and printed stalk formula omit the
superscript $i$ on the local $\Extop$ term. The preceding direct-limit formula
and the following vanishing display both require $\Extop^i$.}
\[
  (\SheafExt_Y^i(F,G))_x\cong
  \varinjlim_k
  \Extop^i_{\Oo_{X,x}}((F/\Jj^kF)_x,G_x).
\]

Since
\[
  \Supp\SheafExt_Y^i(F,G)\subset S\cap Y,
\]
it is enough to note that $x\in Y\cap S$ implies
$\dim\Oo_{X,x}\leq m$, and hence
\[
  \Extop^i_{\Oo_{X,x}}((F/\Jj^kF)_x,G_x)=0
  \qquad\text{if }i>m,
\]
because the global cohomological dimension of a regular local ring is equal
to its dimension.\footnote{Cf. EGA $0_{\mathrm{IV}}$, 17.3.1.}
% END COMPONENT SGA2-VII-CUM-C10

% BEGIN COMPONENT SGA2-VII-CUM-C11
Let $X$ be a locally Noetherian prescheme. For every subset $P$ of $X$, set
\[
  D(P)=\left\{\dim\Oo_{X,p}\mid p\in P\right\}.
\]

\subsection*{\hyperref[VII.2.2]{Lemma 2.2}}
\sgareaderlabel{VII.2.2}{2.2}

If $P$ is the underlying space of a connected subprescheme of $X$, then
$D(P)$ is an interval.\footnote{The French TeX and printed witness read
``of $A$.'' The ambient prescheme introduced immediately before is $X$, and
no object $A$ occurs in this argument. The reader therefore uses $X$
transparently.}

Indeed, let $a$ and $b$ belong to $D(P)$, corresponding to points $p$ and $q$
of $P$. We show that there is a sequence of points of $P$,
\[
  p=p_1,\ldots,p_n=q,
\]
such that, for $1\leq i<n$,
\[
  \lvert\dim\Oo_{X,p_i}-\dim\Oo_{X,p_{i+1}}\rvert=1.
\]
It will follow that $D(P)$ contains the interval $[a,b]$.\footnote{The French
TeX originally introduced the numerical endpoints as ``$a$ and $B$''
and printed the interval $[p,q]$, although $p$ and $q$ are points of $P$.
The reader normalizes the two numerical endpoints to $a,b$ and the interval to
$[a,b]$. The absolute-value bars are present in the current French TeX and
printed witness and are preserved.}

For this, observe that $p$ and $q$ can be joined by a chain of irreducible
components of $P$ in which two successive components meet. We are reduced to
the case where $p$ is the generic point of an irreducible component $Q$ of
$P$, and where $q\in Q$, hence $q\supset p$ as ideals of $\Oo_q$; there the
assertion follows immediately from the definition of dimension.
% END COMPONENT SGA2-VII-CUM-C11

\clearpage
% BEGIN COMPONENT SGA2-VII-CUM-C12
\subsection*{\hyperref[VII.2.3]{Proposition 2.3}}
\sgareaderlabel{VII.2.3}{2.3}

Let $X$ be a locally Noetherian regular prescheme, let $Y$ be a closed subset
of $X$, and let $F$ be a coherent $\OX$-module. Put
\[
  P=Y\cap\overline{\Supp F\cap(X-Y)}.
\]
Let $n\in\mathbb Z$, and suppose that $n\notin D(P)$. Then
$\SheafExt_Y^n(F,\OX)$ is coherent.\footnote{The French TeX and printed
statement omit the closure and read
$P=Y\cap\Supp F\cap(X-Y)$, which is empty literally. Section~2 defines
$S'=\overline{\Supp F\cap(X-Y)}$, and Propositions~2.1 and~1.7 both use the
locus $Y\cap S'$. The reader restores that uniquely supported closure.}

\subsection*{Proof}

The conclusion is local, and the hypotheses are preserved under restriction
to an open subset. Since $P$ is closed, it is locally Noetherian and hence
locally connected. We may therefore suppose that $X$ is affine and Noetherian
and that $P$ is connected.

If $P=\varnothing$, the depth hypothesis of \hyperref[VII.1.7]{Proposition~1.7} with threshold
$n+1$ is vacuous, and the conclusion follows. Assume henceforth that
$P\neq\varnothing$. By \hyperref[VII.2.2]{Lemma~2.2}, $D(P)$ is an interval. Since
$n\notin D(P)$, exactly one of the following alternatives holds:\footnote{The
French TeX and print write $D(P)=[a,b[$ and then treat only $n>b$ and $n<a$,
leaving the allowed case $n=b$ uncovered. The print does not decide whether
the intended one-character repair was a closed upper endpoint or a non-strict
upper inequality. The reader states the order alternatives forced by
$n\notin D(P)$ and \hyperref[VII.2.2]{Lemma~2.2} and therefore does not guess that glyph.}
\[
  n>d\quad\text{for every }d\in D(P),
  \qquad\text{or}\qquad
  n<d\quad\text{for every }d\in D(P).
\]

In the first case, put $m=\sup D(P)$. Then $n>m$, and \hyperref[VII.2.1]{Proposition~2.1},
applied with $G=\OX$, shows that $\SheafExt_Y^n(F,\OX)$ is coherent.

In the second case, for every $x\in P$ regularity gives
\[
  n<\dim\Oo_{X,x}=\depthop\Oo_{X,x},
\]
and hence $\depthop\Oo_{X,x}\geq n+1$. \hyperref[VII.1.7]{Proposition~1.7} with threshold
$n+1$ therefore gives coherence in every degree $i<n+1$, in particular in
degree $i=n$.
% END COMPONENT SGA2-VII-CUM-C12

% END INCLUDED active_source/parts/07_expose_VII.tex

% BEGIN INCLUDED active_source/parts/08_expose_VIII.tex
% BEGIN COMPONENT SGA2-VIII-CUM-C01
\makeatletter
\providecommand{\sgareaderlabel}[2]{%
  \begingroup
  \def\@currentlabel{#2}%
  \phantomsection\label{#1}%
  \endgroup}
\makeatother

\chapter{The Finiteness Theorem}
\label{VIII}

\section*{1. A Biduality Spectral Sequence\footnote{A reader familiar with
Verdier's language of derived categories will recognize the spectral sequence
associated with a \emph{biduality isomorphism}. Cf. SGA~6, I.}}
\sgareaderlabel{VIII.1}{1}

Let us state the result we wish to establish:

\begin{sgastatement}{Proposition 1.1}\sgareaderlabel{VIII.1.1}{1.1}
Let $A$ be a Noetherian ring and let $I$ be an ideal of $A$. Put
$X=\operatorname{Spec}(A)$ and $Y=V(I)$. Let $M$ be a finitely generated
$A$-module of finite projective dimension. Let
$\mathcal F=\widetilde M$ be the $\OX$-module associated with $M$.
\begin{enumerate}
\item[\textup{1)}] There is a spectral sequence
\[
  \Hloc_Y(X,\mathcal F)\Longleftarrow
  \Ext_Y^p(\Ext^{-q}(M,A),A).
\]
\item[\textup{2)}] There is a spectral sequence
\[
  \SheafH_Y(X,\mathcal F)\Longleftarrow
  \SheafExt_Y^p(\SheafExt^{-q}(\mathcal F,\OX),\OX).
\]
\end{enumerate}
\end{sgastatement}
% END COMPONENT SGA2-VIII-CUM-C01

% VERIFIED SOURCE SEAM: French line 2526 is blank and excluded.

% INTEGRATION CONTROL BEFORE SGA2-VIII-CUM-C02
\let\cF\uF

% BEGIN COMPONENT SGA2-VIII-CUM-C02
Of course, 2) follows from 1) by observing that, if $M$ and $N$ are two
finitely generated $A$-modules and if we put $\mathcal F=\widetilde M$ and
$\mathcal G=\widetilde N$, then there are isomorphisms
\begin{align*}
  \SheafH_Y(\mathcal F)
    &\simeq \widetilde{\Hloc_Y(X,\mathcal F)}, \\
  \SheafExt_Y(\mathcal F,\mathcal G)
    &\simeq \widetilde{\Ext_Y(\mathcal F,\mathcal G)}, \\
  \SheafExt_{\OX}(\mathcal F,\mathcal G)
    &\simeq \widetilde{\Ext_A(M,N)}.
\end{align*}

Let $\CatC$ be the category of $A$-modules, and let $\Ab$ be the category of
abelian groups. Let $\cF$ be the functor
\[
  \begin{aligned}
    \cF\colon \CatC &\longrightarrow \Ab, \\
    M &\longmapsto \Gamma_Y(\widetilde M).
  \end{aligned}
\]
We know from Expos\'e~II that there is an isomorphism of $\partial$-functors
\[
  \Hloc_Y^*(X,\widetilde M)\simeq \R^*\cF(M).
\]

Furthermore, let $\Ext_Y^*$ denote the right derived functors, in the second
argument, of
\[
  \cF\circ\operatorname{Hom}\colon \CatC^{\circ}\times\CatC\longrightarrow\Ab.
\]
We know from Expos\'e~VI that there is an isomorphism of $\partial$-functors
\[
  \Ext_Y^*(M,N)\simeq \Ext_Y^*(\widetilde M,\widetilde N).
\]

Finally, let us record the following result from Expos\'e~VI: if $C$ is an
injective $A$-module and $N$ is a finitely generated $A$-module, the sheaf
\[
  \SheafHom(\widetilde N,\widetilde C)
    \simeq \widetilde{\operatorname{Hom}(N,C)}
\]
is flasque; hence
\[
  \R^1\cF(\operatorname{Hom}(N,C))=0.
\]
% END COMPONENT SGA2-VIII-CUM-C02

% VERIFIED SOURCE SEAM: French line 2552 is blank and excluded.

% BEGIN COMPONENT SGA2-VIII-CUM-C03
It remains for us to prove the following result.

\medskip
\sgareaderlabel{VIII.1.2}{1.2}\noindent\textbf{Lemma 1.2.} --- Let $A$ be a Noetherian ring and let $\CatC$
be the category of $A$-modules. Let
$\cF\colon\CatC\longrightarrow\Ab$ be a left exact additive functor such that,
for every finitely generated $A$-module $N$ and every injective $A$-module
$C$, one has
\[
  \R^1\cF(\operatorname{Hom}(N,C))=0.
\]
Let $M$ be a finitely generated $A$-module of finite projective dimension.
There is a spectral sequence
\[
  \R^\ast\cF(M)
    \Longleftarrow
  \Ext_{\cF}^{p}(\Ext^{-q}(M,A),A),
\]
where $\Ext_{\cF}^{p}$ denotes the $p$-th right derived functor of
$\cF\circ\operatorname{Hom}$.
% END COMPONENT SGA2-VIII-CUM-C03

% VERIFIED SOURCE SEAM: French line 2560 is blank and excluded.

% BEGIN COMPONENT SGA2-VIII-CUM-C04
\emph{We shall consider only complexes whose differential has degree $+1$.}
By the hypothesis on $M$, there exists a projective resolution of $M$ of
finite length,
\[
  u\colon L^\bullet\longrightarrow M,
\]
where, moreover, the $L^p$ are finitely generated modules and $L^p=0$ if
$p\notin[-n,0]$. On the other hand, let
$v\colon M\longrightarrow I^\bullet$ be an injective resolution of $M$. I
claim that
\begin{equation}
  v\circ u\colon L^\bullet\longrightarrow I^\bullet
  \tag{1.1}\label{eq:VIII.1.1}
\end{equation}
is an injective resolution of $L^\bullet$. We must specify what this means.

\medskip
\sgareaderlabel{VIII.1.3}{1.3}\noindent\textbf{Definition 1.3.} --- Let $X^\bullet$ be a complex of
$A$-modules. An \emph{injective resolution} of $X^\bullet$ is a homomorphism
of complexes
\[
  x\colon X^\bullet\longrightarrow CX^\bullet
\]
such that $CX^p$ is injective for every $p\in\mathbb Z$, and $x$ induces an
isomorphism in homology.
% END COMPONENT SGA2-VIII-CUM-C04

% VERIFIED SOURCE SEAM: French line 2572 is blank and excluded.

% BEGIN COMPONENT SGA2-VIII-CUM-C05
\sgareaderlabel{VIII.1.4}{1.4}\noindent\textbf{Proposition 1.4.} --- Every left-bounded complex, i.e., every
complex for which there exists $q\in\ZZ$ such that $X^p=0$ for $p<q$, admits
an injective resolution. Moreover, suppose that
$u\colon X^\bullet\longrightarrow Y^\bullet$ is a homomorphism of complexes
(both left-bounded), and that
$x\colon X^\bullet\longrightarrow CX^\bullet$ and
$y\colon Y^\bullet\longrightarrow CY^\bullet$ are injective resolutions of
$X^\bullet$ and $Y^\bullet$. Then there exists a homomorphism of complexes
\[
  Cu\colon CX^\bullet\longrightarrow CY^\bullet,
\]
unique up to homotopy, such that the diagram
\[
\begin{CD}
X^\bullet @>{x}>> CX^\bullet \\
@V{u}VV @VV{Cu}V \\
Y^\bullet @>{y}>> CY^\bullet
\end{CD}
\]
commutes up to homotopy.
% END COMPONENT SGA2-VIII-CUM-C05

% VERIFIED SOURCE SEAM: French line 2582 is blank and excluded.

% BEGIN COMPONENT SGA2-VIII-CUM-C06
The proof is left to the reader.\footnote{See also H.~Cartan and
S.~Eilenberg, \emph{Homological Algebra}, Princeton Mathematical Series,
vol.~19, Princeton University Press, 1956.}

Let us recall a notation introduced in Expos\'e V.

\medskip
\noindent\textbf{Notation.} --- Let $X^\bullet$ and $Y^\bullet$ be two
complexes. We denote by $\operatorname{Hom}^\bullet(X^\bullet,Y^\bullet)$ the
\emph{simple complex} whose component of degree $n$ is
\[
  (\operatorname{Hom}^\bullet(X^\bullet,Y^\bullet))^n
  =\prod_{-p+q=n}\operatorname{Hom}(X^p,Y^q),
\]
also denoted by $\operatorname{Hom}^n(X^\bullet,Y^\bullet)$, and whose
differential is given by
\begin{gather*}
  \partial_n\colon
  \operatorname{Hom}^n(X^\bullet,Y^\bullet)
  \longrightarrow
  \operatorname{Hom}^{n+1}(X^\bullet,Y^\bullet),\\
  \partial_n=d'+(-1)^{n+1}d'',
\end{gather*}
where $d'$ and $d''$ are the differentials (of degree $+1$) induced by those
of $X^\bullet$ and $Y^\bullet$.
% END COMPONENT SGA2-VIII-CUM-C06

% VERIFIED SOURCE SEAM: French line 2596 is blank and excluded.

% INTEGRATION CONTROL BEFORE SGA2-VIII-CUM-C07
\renewcommand{\theequation}{1.\arabic{equation}}\setcounter{equation}{1}

% BEGIN COMPONENT SGA2-VIII-CUM-C07
Let $A^\bullet$ then be the complex defined by $A^p=0$ if $p\ne 0$ and
$A^0=A$. Let
\[
  a\colon A^\bullet\longrightarrow CA^\bullet
\]
be an injective resolution of $A^\bullet$. Consider the double complex
\begin{equation}\label{eq:VIII.1.2}
  Q^{p,q}=\operatorname{Hom}(\operatorname{Hom}(L^{-q},A),CA^p).
\end{equation}
The first spectral sequence of the bicomplex $\cF Q^{\bullet\bullet}$ will
give the conclusion of \hyperref[VIII.1.2]{Lemma~1.2}.

Set
\begin{equation}\label{eq:VIII.1.3}
  {L'}^\bullet=\operatorname{Hom}^\bullet(L^\bullet,A^\bullet),
\end{equation}
and
\begin{equation}\label{eq:VIII.1.4}
  P^\bullet=\operatorname{Hom}^\bullet({L'}^\bullet,CA^\bullet).
\end{equation}
% END COMPONENT SGA2-VIII-CUM-C07

% VERIFIED SOURCE SEAM: French line 2610 is blank and excluded.

% INTEGRATION CONTROL BEFORE SGA2-VIII-CUM-C08
\renewcommand{\theequation}{1.\arabic{equation}}\setcounter{equation}{4}

% BEGIN COMPONENT SGA2-VIII-CUM-C08
One sees easily that $P^\bullet$ is the simple complex associated with
$Q^{\bullet\bullet}$. Let us compute the abutment of the spectral sequence,
that is, the homology of $\uF P^\bullet$. To this end, using the fact that
$L^\bullet$ is projective of finite type in every degree, one proves that
\[
  L^\bullet \cong \operatorname{Hom}^\bullet({L'}^\bullet,A^\bullet).
\]
From the homomorphism $a\colon A^\bullet\to CA^\bullet$, one obtains a
homomorphism
\[
  b\colon \operatorname{Hom}^\bullet({L'}^\bullet,A^\bullet)
  \longrightarrow \operatorname{Hom}^\bullet({L'}^\bullet,CA^\bullet),
\]
or equivalently a homomorphism
\begin{equation}\label{eq:VIII.1.5}
  c\colon L^\bullet\longrightarrow P^\bullet.
\end{equation}

This being said, it is easy to see, using the fact that ${L'}^\bullet$ is
projective of finite type in every degree and bounded on the left, that
\hyperref[eq:VIII.1.5]{(1.5)} is an injective resolution of $L^\bullet$. Using \hyperref[VIII.1.4]{Proposition~1.4}, one
concludes that $P^\bullet$ is homotopy equivalent to $I^\bullet$, where
$I^\bullet$ is the injective resolution of $M$ introduced above in \hyperref[eq:VIII.1.1]{(1.1)}.
It follows that the \emph{abutment} of the first spectral sequence of the
double complex $\uF Q^{\bullet\bullet}$, which is
$H^*(\uF P^\bullet)$, \emph{is isomorphic} to $R^*\uF(M)$.
% END COMPONENT SGA2-VIII-CUM-C08

% VERIFIED SOURCE SEAM: French line 2617 is blank and excluded.

% BEGIN COMPONENT SGA2-VIII-CUM-C09
The initial term of the first spectral sequence of the bicomplex
$\cF Q^{\bullet\bullet}$ is
\[
  E_2^{p,q}
  ={}'\!H^p({}''\!H^q(\cF Q^{\bullet\bullet})).
\]

For every $p\in\ZZ$, $CA^p$ is injective. By the hypothesis on $\cF$, the
functor (restricted to the category of finitely generated modules)
\[
  N\longmapsto \cF\operatorname{Hom}(N,CA^p)
\]
is exact. Hence one obtains isomorphisms
\[
  {}''\!H^q(\cF\operatorname{Hom}({L'}^\bullet,CA^p))
  \cong
  \cF\operatorname{Hom}(H^{-q}({L'}^\bullet),CA^p).
\]
By the definition of $\operatorname{Ext}^{\ast}_{\cF}$ as the right-derived functor
of $\cF\circ\operatorname{Hom}$, one therefore obtains isomorphisms
\[
  E_2^{p,q}
  \cong
  \operatorname{Ext}^{p}_{\cF}(H^{-q}({L'}^\bullet),A).
\]

\noindent\emph{Source note.} The French TeX and printed display have
$H^q({L'}^\bullet)$ in the last formula. The preceding displayed
isomorphism, the statement of Lemma~1.2, and the identity below require
$H^{-q}({L'}^\bullet)$. This target therefore uses $H^{-q}$ as a transparent
editorial emendation. It is not a restoration from an independent original
scan; the French authority remains unchanged pending editorial action.

Now
\[
  {L'}^\bullet=\operatorname{Hom}^\bullet(L^\bullet,A^\bullet),
\]
where $L^\bullet$ is a projective resolution of $M$; hence there are
isomorphisms
\[
  \operatorname{Ext}^{q}(M,A)\cong H^q({L'}^\bullet).
\]
Substituting $-q$ in this last identity gives the conclusion.\hfill$\square$
% END COMPONENT SGA2-VIII-CUM-C09

% VERIFIED SOURCE SEAM: French line 2639 is blank and excluded.

% INTEGRATION CONTROL BEFORE SGA2-VIII-CUM-C10
\let\cF\calF\renewcommand{\theequation}{\thesection.\arabic{equation}}

% BEGIN COMPONENT SGA2-VIII-CUM-C10
\setcounter{section}{1}
\section{The Finiteness Theorem}\label{VIII.2}

\begin{theorem}\label{VIII.2.1}
\markedfootnote{(1)}{\textit{Editor's note.} For an analogous statement, but
in a somewhat more general situation, see Mme Raynaud (Raynaud M., ``Th\'eor\`emes de
Lefschetz en cohomologie des faisceaux coh\'erents et en cohomologie
\'etale. Application au groupe fondamental,'' \emph{Ann. Sci. \'Ec. Norm.
Sup.} (4) \textbf{7} (1974), pp.~29--52, Proposition~II.2.1).}%
\kern-.18em\textbf{.}\enspace Let $X$ be a locally Noetherian prescheme, let $Y$ be a closed subset of $X$,
and let $\cF$ be a coherent $\OX$-module. Suppose that $X$ is locally
embeddable in a regular prescheme\markedfootnote{*}{This condition may be
generalized to the hypothesis that a ``dualizing complex,'' in the sense
defined in R.~Hartshorne, \emph{Residues and Duality} (cited in the note to
Expos\'e~IV, p.~46), exists locally on $X$.}. Let $i\in\ZZ$. Suppose that:
\begin{enumerate}[label=\textup{\alph*)},leftmargin=2.2em,itemsep=0.5\baselineskip]
\item for every $x\in U=X-Y$, one has
\begin{equation*}
  H^{i-c(x)}(\cF_x)=0,
\end{equation*}
where we have set\markedfootnote{(2)}{\textit{Editor's note.} As in
Expos\'e~V, $H^{\ast}(\cF_x)$ denotes the local cohomology
$H^{\ast}_{\mathfrak m_x}(\cF_x)$.}:
\begin{equation}\label{eq:VIII.2.1}
  c(x)=\operatorname{codim}
  (\overline{\{x\}}\cap Y,\overline{\{x\}}).
\end{equation}
Then:
\item $\mathcal H^i_Y(\cF)$ is coherent.
\end{enumerate}
\end{theorem}
% END COMPONENT SGA2-VIII-CUM-C10

% VERIFIED SOURCE SEAM: French line 2660 is blank and excluded.

% BEGIN COMPONENT SGA2-VIII-CUM-C11
\setcounter{section}{2}
\setcounter{sourcecorollary}{1}
\begin{sourcecorollary}{\textit{Editor's note.} Strictly speaking, this is a
corollary of the proof that follows, not of the statement. The implication
\textup{c)}$\Rightarrow$\textup{a)} is tautological. The converse is not, but
follows from the proof. More precisely, as below, one covers $X$ by open
subsets embeddable in regular schemes; as explained below, this permits one to
reduce to the case where $X=\operatorname{Spec}(A)$ is affine and regular and
$\cF=\widetilde M$, where $M$ is an $A$-module of finite projective dimension.
In this case it is shown that conditions \textup{a)} and \textup{c)} are
equivalent to the dual conditions \textup{a$'$)} and \textup{c$'$)}. One then
shows that \textup{c$'$)} implies condition \textup{d)} (see below), which in
turn implies \textup{a$'$)}. See the discussion following~\hyperref[VIII.2.4]{2.4}.}\label{VIII.2.2}
Under the hypotheses of the preceding theorem, condition \textup{a)} is
equivalent to:
\begin{enumerate}[label=\textup{\alph*)},start=3,leftmargin=2.2em]
\item for every $x\in U$ such that $c(x)=1$, one has
  $H^{i-1}(\cF_x)=0$.
\end{enumerate}
\end{sourcecorollary}
% END COMPONENT SGA2-VIII-CUM-C11

% VERIFIED SOURCE SEAM: French line 2669 is blank and excluded.

% BEGIN COMPONENT SGA2-VIII-CUM-C12
\setcounter{section}{2}
\setcounter{sourcecorollary}{2}
\begin{sourcecorollary}\label{VIII.2.3}
Let $X$ be a locally Noetherian prescheme that is locally embeddable in a
regular prescheme, let $Y$ be a closed subset of $X$, let $\cF$ be a coherent
$\mathcal O_X$-module, and let $n$ be an integer. The following conditions are
equivalent:
\begin{enumerate}[label=\textup{(\roman*)},leftmargin=2.8em]
\item for every $x\in U$, one has
  $\prof \cF_x>n-c(x)$;
\item for every $x\in U$ such that $c(x)=1$, one has
  $\prof \cF_x\geq n$;
\item for every $i\in\mathbb Z$, $\mathcal H^i_Y(\cF)$ is coherent if
  $i\leq n$;
\item $R^p i_*(\cF|_U)$ is coherent for $p<n$, where $i\colon U\to X$ is the canonical immersion%
  \markedfootnote{(4)}{\textit{Editor's note.} This condition appeared only
  in the body of the proof, but not in the statement of the corollary; since
  it is used in \hyperref[VIII.3]{Section~3}, it has been added.}.
\end{enumerate}
\end{sourcecorollary}
% END COMPONENT SGA2-VIII-CUM-C12

% VERIFIED SOURCE SEAM: French line 2682 is blank and excluded.

% BEGIN COMPONENT SGA2-VIII-CUM-C13
\medskip
Assume these results established when $X$ is the spectrum of a regular
Noetherian ring $A$ and when $\cF$ is the sheaf associated with an $A$-module
of finite projective dimension.

First observe that, if $(X_j)_{j\in J}$ is an open covering of $X$ by open
subsets embeddable in a regular scheme, each of the conditions above is
equivalent to the conjunction of the analogous conditions obtained by
replacing $X$ by $X_j$, $Y$ by $Y_j=Y\cap X_j$, and $\cF$ by
$\cF|_{X_j}$. Indeed, only the conditions involving $c(x)$ can present a
difficulty. Let $x\in U$. If $x\in X_j$, set
\[
  c_j(x)=\codim(X_j\cap\overline{\{x\}}\cap Y,
                X_j\cap\overline{\{x\}}).
\]
Necessarily $c_j(x)\geq c(x)$. Let $y\in\overline{\{x\}}\cap Y$ be a point
that ``realizes the codimension,'' that is, such that
$c(x)=\dim\mathcal O_{\overline{\{x\}},y}$, and let $X_j$ be a member of
the covering such that $y\in X_j$. Then $x\in X_j$, and hence
$c_j(x)=c(x)$. It follows that condition~(a) for the $X_j$ implies
condition~(a) for $X$. At this point we have only a partial converse: condition
(a) for $X$ implies condition~(a) for those $X_j$ such that
$c(x)=c_j(x)$, which suffices for our purposes.%
\markedfootnote{(5)}{\textit{Editor's note.} In fact, condition~(a) for $X$
implies condition~(a) for every $X_j$, as asserted in the original text, but
establishing this requires reading the proof below in detail. At this point
the implication does not appear formal. For a property (a), (b), or (c),
append a subscript $J$ to denote its conjunction over the $X_j$. The proof
\textit{below} establishes $(c_J)\Rightarrow(a_J)$ (this is the chain of implications
$(c')\Rightarrow(d)\Rightarrow(a')$). But tautologically
$(a)\Rightarrow(c)$, and $(c)\Longleftrightarrow(c_J)$; hence
$(a)\Rightarrow(a_J)$.}
% END COMPONENT SGA2-VIII-CUM-C13

% VERIFIED SOURCE SEAM: French line 2714 is blank and excluded.

% INTEGRATION CONTROL BEFORE SGA2-VIII-CUM-C14
\setcounter{section}{2}\setcounter{lemma}{3}

% BEGIN COMPONENT SGA2-VIII-CUM-C14
\medskip
Choose a covering of $X$ by open subsets embeddable in a regular prescheme.
Applying the preceding observation, we see that we may assume that $X$ is
closed in a regular $X'$. The reduction to $X'$ is then immediate.

We may therefore assume that $X$ is regular, and even affine, by covering
$X$ with affine open subsets. That we may assume that $\cF=\widetilde M$, where
$M$ has finite projective dimension, will follow from the next lemma.

\begin{lemma}\label{VIII.2.4}
Let $X$ be a regular Noetherian prescheme. Let $\cF$ be a coherent
$\mathcal O_X$-module. The function that assigns to each $x\in X$ the
projective dimension of $\cF_x$ is bounded above.
\end{lemma}
% END COMPONENT SGA2-VIII-CUM-C14

% VERIFIED SOURCE SEAM: French line 2722 is blank and excluded.

% BEGIN COMPONENT SGA2-VIII-CUM-C15
\setcounter{equation}{1}
\medskip
Indeed, let $x\in X$, and let $U$ be an affine open neighborhood of $x$.
Let $L^\bullet$ be a projective resolution of the module $\cF(U)$, where the
$L^i$ are finitely generated. By hypothesis, the ring $\cO_{X,x}$ is regular,
so the projective dimension of $\cF_x$ is finite; denote this integer by $d$.
Set
\[
  K=\operatorname{im}(L^{-d}\longrightarrow L^{-d+1})=\ker(L^{-d+1}\longrightarrow L^{-d+2}).
\]
The module $K_x$ is free, since $d$ is the projective dimension of $\cF_x$
([M], Ch.~VI, Prop.~2.1). By (EGA $0_{\mathrm I}$, 5.4.1 Errata), it follows
that the $\cO_U$-module $\widetilde K$ is free on a neighborhood $U'$ of $x$,
with $U'\subset U$. Choose $f\in\cO_X(U)$ such that
$x\in D(f)\subset U'$. We then have a projective resolution of $M_f$, where
$M=\cF(U)$:
\[
  0\longrightarrow K_f\longrightarrow (L^{-d+1})_f\longrightarrow
  \cdots\longrightarrow M_f\longrightarrow 0.
\]

\noindent\emph{Source note.} The printed French kernel and the earlier
maintained-English emendation both misplace the $d$-th syzygy. Exactness
gives $K=\operatorname{im}(L^{-d}\to L^{-d+1})$
$=\ker(L^{-d+1}\to L^{-d+2})$, so the truncated resolution begins with
$(L^{-d+1})_f$. The printed $M_{f'}$ is also unbound after only $f$ is
chosen and is corrected to $M_f$. Canon correction SGA-SRC-001691
preserves the printed proof diplomatically and uses these type-correct
terms in the current branch.

This proves that the function under consideration is upper semicontinuous.
Since $X$ is quasi-compact, the conclusion follows.
% END COMPONENT SGA2-VIII-CUM-C15

% VERIFIED SOURCE SEAM: French line 2732 is blank and excluded.

% BEGIN COMPONENT SGA2-VIII-CUM-C16
\medskip
We henceforth assume that $X$ is affine, Noetherian, and regular, and that
$\cF=\widetilde M$, where $M$ is a finitely generated $A$-module, necessarily
of finite projective dimension. We shall proceed in several steps. First, we
find a condition~(d), equivalent to~(a), and prove that it is also equivalent
to~(c). Then, using the spectral sequence of the preceding subsection, we
prove $(d)\Rightarrow(b)$. It remains to prove
$(\mathrm{iii})\Rightarrow(\mathrm{ii})$; indeed,
$(\mathrm{i})\Leftrightarrow(\mathrm{ii})\Rightarrow(\mathrm{iii})$ follows
immediately from $(a)\Leftrightarrow(c)\Rightarrow(b)$.

Let $x\in U$. By hypothesis, $\cO_{X,x}$ is a regular local ring. If $D$
denotes the dualizing functor for $\cO_{X,x}$, it follows from (\hyperref[V.2.1]{V, 2.1)} that
\[
  D\,H^{i-c(x)}(\cF_x)\simeq
  \operatorname{Ext}^{d(x)-i}_{\cO_{X,x}}(\cF_x,\cO_{X,x}),
\]
where
\begin{equation}\label{eq:VIII.2.2}
  d(x)=\dim\cO_{X,x}+c(x)
      =\dim\cO_{X,x}
       +\codim(\overline{\{x\}}\cap Y,\overline{\{x\}}).
  \tag{2.2}
\end{equation}
Since $X$ is Noetherian and $\cF$ is coherent,
\begin{equation}\label{eq:VIII.2.3}
  D\,H^{i-c(x)}(\cF_x)\simeq
  (\sExt^{d(x)-i}_{\cO_X}(\cF,\cO_X))_x.
  \tag{2.3}
\end{equation}
Moreover, a module vanishes if and only if its dual vanishes (cf. the
editor's note~(4) on p.~54). For every $q\in\mathbb Z$, set
\begin{equation}\label{eq:VIII.2.4}
  \begin{cases}
    S_q=\Supp\sExt^q_{\cO_X}(\cF,\cO_X),\\
    S'_q=S_q\cap U,\quad (U=X-Y),\\
    Z_q=\overline{S'_q}\cap Y.
  \end{cases}
  \tag{2.4}
\end{equation}
% END COMPONENT SGA2-VIII-CUM-C16

% BEGIN COMPONENT SGA2-VIII-CUM-C17
\medskip
It follows from \hyperref[eq:VIII.2.3]{formula~(2.3)} that conditions~(a) and~(c) are respectively
equivalent to:

\medskip
\noindent\textup{(a$'$)} for every $q\in\mathbb Z$ and every $x\in S'_q$,
one has $q+i\ne d(x)$;

\smallskip
\noindent\textup{(c$'$)} for every $q\in\mathbb Z$ and every $x\in S'_q$,
if $c(x)=1$, then $q+i\ne d(x)$.

\medskip
\noindent Here is condition~(d), promised above:

\medskip
\noindent\textup{(d)} for every $q\in\mathbb Z$ and every $y\in Z_q$,
one has $q+i\ne\dim\cO_{X,y}$.

\medskip
\noindent These conditions are equivalent.

\smallskip
\noindent\textup{(a$'$)} $\Rightarrow$ \textup{(c$'$)} is recorded for
completeness.

\smallskip
\noindent\textup{(d)} $\Rightarrow$ \textup{(a$'$)}. Indeed, let
$q\in\mathbb Z$ and $x\in S'_q$. Choose a point
$y\in\overline{\{x\}}\cap Y$ that realizes the
codimension,\markedfootnote{(6)}{\textit{Editor's note.} Throughout what
follows, closures of points are endowed with the reduced structure.} that is,
such that
\begin{equation}\label{eq:VIII.2.5}
  \dim\cO_{\overline{\{x\}},y}
  =\codim(\overline{\{x\}}\cap Y,\overline{\{x\}})=c(x).
  \tag{2.5}
\end{equation}
Since $X$ is regular at $y$, it follows that
\begin{equation}\label{eq:VIII.2.6}
  \dim\cO_{X,y}=d(x)\qquad\text{(cf.~\hyperref[eq:VIII.2.2]{(2.2)}).}
  \tag{2.6}
\end{equation}
But $y\in\overline{\{x\}}$, and hence $y\in Z_q$; the conclusion follows.

\smallskip
\noindent\textup{(c$'$)} $\Rightarrow$ \textup{(d)}. Let $q\in\mathbb Z$
and $y\in Z_q$. Assume provisionally that there exists $x\in S'_q$ such
that
\[
  y\in\overline{\{x\}}\qquad\text{and}\qquad
  \dim\cO_{\overline{\{x\}},y}=1.
\]
(One also says that $y$ is an immediate specialization of $x$.) It follows
that $c(x)=1$, since $y$ realizes the codimension of
$\overline{\{x\}}\cap Y$ in $\overline{\{x\}}$, while $x\notin Y$.
Condition~\textup{(c$'$)} therefore gives
\[
  q+i\ne d(x).
\]
Whence the conclusion, upon observing that $d(x)=\dim\cO_{X,y}$ by~\hyperref[eq:VIII.2.6]{(2.6)}.
The assertion provisionally assumed above is the content of the following
lemma.

\medskip
\sgareaderlabel{VIII.2.5}{2.5}\noindent\textbf{Lemma 2.5.} Let $X$ be a locally Noetherian prescheme and
$Y$ a closed subset of $X$. Set $U=X-Y$, and suppose that $U$ is dense in
$X$. For every $y\in Y$, there exists $x\in U$ such that $y$ is an immediate
specialization of $x$, that is,
\[
  y\in\overline{\{x\}}\qquad\text{and}\qquad
  \dim\cO_{\overline{\{x\}},y}=1.
\]

We have applied the lemma with $X$ equal to the prescheme
$\overline{S'_q}$ and $Y$ equal to the subset
$Y\cap\overline{S'_q}$.

\medskip
\noindent\textit{Proof of \hyperref[VIII.2.5]{Lemma 2.5}.} There exists $x\in U$ such that
$y\in\overline{\{x\}}$. Choose such an $x$ for which
$r=\dim\cO_{\overline{\{x\}},y}$ is minimal. We must prove that $r=1$.
By the choice of $x$, every
$z\in\Spec(\cO_{\overline{\{x\}},y})$ with $z\ne x$ lies in $Y$.
Consequently, $\{x\}$ is open in
$\Spec(\cO_{\overline{\{x\}},y})$, and the conclusion follows.\hfill$\square$
% END COMPONENT SGA2-VIII-CUM-C17

% VERIFIED SOURCE SEAM: French line 2795 is blank and excluded.

% BEGIN COMPONENT SGA2-VIII-CUM-C18
\medskip
\noindent\textit{The second step is to deduce \textup{(b)} from
\textup{(d)}.}

\medskip
Set
\[
  D(Z_q)=\{\dim\cO_{X,y}\mid y\in Z_q\}.
\]
By~\textup{(d)}, for every $q\in\mathbb Z$ one has
$q+i\notin D(Z_q)$. Applying Proposition~\hyperref[VII.2.3]{VII.2.3}, we see that
\[
  \sExt_Y^{q+i}(\sExt^q(\cF,\cO_X),\cO_X)
  \quad\text{is coherent.}
\]
The initial term of the spectral sequence of the preceding subsection is
\[
  E_2^{p,q}
  =\sExt_Y^p(\sExt^{-q}(\cF,\cO_X),\cO_X).
\]
It follows that $E_2^{p,q}$ is coherent for every $p,q\in\mathbb Z$ such
that $p+q=i$. Now \sisi{there are only finitely many pairs $(p,q)$ such that
$p+q=i$}{only finitely many pairs $(p,q)$ with $p+q=i$ contribute nonzero $E_2$ terms},\markedfootnote{[S]}{\textit{Source note.} Taken literally for
arbitrary integers, this sentence is false. In context it must be read as
concerning the terms that can contribute, in view of the finite-projective-dimension
setup above. The French authority is unchanged.} and this spectral
sequence converges to $\sH_Y^\ast(\cF)$; whence the conclusion.
% END COMPONENT SGA2-VIII-CUM-C18

% VERIFIED SOURCE SEAM: French line 2806 is blank and excluded.

% BEGIN COMPONENT SGA2-VIII-CUM-C19
\medskip
\noindent\textit{It remains to prove that \textup{(iii)} implies
\textup{(ii)}.} Let
\[
  i\colon U\longrightarrow X
\]
denote the canonical immersion of $U$ into $X$. In view of the exact homology
sequence associated with the closed subset $Y$ (\hyperref[I.2.11]{I, 2.11)}, one sees that
\textup{(iii)} is \textit{equivalent to}:

\medskip
\noindent\textup{(iv)} $R^i i_*(\cF|_U)$ is coherent for $i<n$.
\refstepcounter{sourcecondition}\label{condition}

\medskip
\noindent{\scriptsize\color{gray}French printed p.~94}\quad Indeed, there is
an exact sequence
\[
  0\longrightarrow \sH_Y^0(\cF)\longrightarrow \cF
  \longrightarrow i_*(\cF|_U)\longrightarrow \sH_Y^1(\cF)
  \longrightarrow 0.
\]
Now $\sH_Y^0(\cF)$ is a quasi-coherent subsheaf of the coherent sheaf $\cF$,
and is therefore coherent. Hence $\sH_Y^1(\cF)$ is coherent if and only if
$i_*(\cF|_U)$ is coherent. Moreover, for $p>0$, the exact cohomology sequence
associated with the closed subset $Y$ reduces to isomorphisms
\[
  R^p i_*(\cF|_U)\overset{\approx}{\longrightarrow}\sH_Y^{p+1}(\cF).
\]
% END COMPONENT SGA2-VIII-CUM-C19

% BEGIN COMPONENT SGA2-VIII-CUM-C20
\medskip
\noindent\textit{We shall prove that \textup{(iv)} implies
\textup{(ii)}.} To this end, recall \textup{(ii)}:

\medskip
\noindent\textup{(ii)} For every $x\in U$ such that $c(x)=1$, one has
\[
  \operatorname{prof}\cF_x\geq n.
\]

\noindent We argue by induction on $n$.

\medskip
\noindent If $n=0$, both conditions are vacuous.
% END COMPONENT SGA2-VIII-CUM-C20

% VERIFIED SOURCE SEAM: French line 2827 is blank and excluded.

% BEGIN COMPONENT SGA2-VIII-CUM-C21
\medskip
\noindent If $n=1$, assume that $i_*(\cF|_U)$ is coherent. We argue by
contradiction and suppose that there exists $x\in U$ such that $c(x)=1$ and
\[
  \prof\cF_x=0,
\]
i.e. $x\in\Ass\cF_x$. Let $y\in\overline{\{x\}}\cap Y$ be such that
\[
  \dim\cO_{\overline{\{x\}},y}=1.
\]
Set
\[
  A=\cO_{X,y}\qquad\text{and}\qquad X'=\Spec(A).
\]

\setcounter{equation}{6}
\noindent Perform the base change $v\colon X'\longrightarrow X$, which is
flat:
\begin{equation}
\begin{CD}
U'=X'\times_X U @>{v'}>> U\\
@V{i'}VV @VV{i}V\\
X' @>{v}>> X .
\end{CD}
\tag{2.7}\label{eq:VIII.2.7}
\end{equation}
The morphism $i$ is separated, since it is an immersion, and is of finite
type, since it is an open immersion and $X$ is locally Noetherian. Since the
base change is flat, EGA~III, 1.4.15 gives an isomorphism
\begin{equation}
  v^*(i_*(\cF|_U))
  \approx
  i'_*(v'^*(\cF|_U)).
\tag{2.8}\label{eq:VIII.2.8}
\end{equation}

\noindent{\scriptsize\color{gray}French printed p.~95}\quad Denote by
$\underline{x}$ (respectively, $\underline{y}$) the ideal of $A$
corresponding to $x$ (respectively, $y$).

\medskip
\noindent Set $\cG=v'^*(\cF|_U)$. Then $\cG$ is coherent and
$\underline{x}\in\Ass\cG$, so there exists a monomorphism
\[
  \cO_{\overline{\{x\}}}\longrightarrow\cG;
\]
consequently, $i'_*(\cO_{\overline{\{x\}}}|_{U'})$ is coherent. By the choice
of $y$, $\dim(A/\underline{x})=1$, and therefore the support of
$\cO_{\overline{\{x\}}}$ is precisely
\[
  \overline{\{\underline{x}\}}
  =\{\underline{x}\}\cup\{\underline{y}\},
\]
since $\overline{\{\underline{x}\}}=\Spec(A/\underline{x})$ as a scheme. It
follows that
\[
  (\cO_{\overline{\{x\}}}|_{U'})(U')
  =\Frac(A/\underline{x}),
\]
the field of fractions of $A/\underline{x}$, and
\[
  i'_*(\cO_{\overline{\{x\}}}|_{U'})(X')
  =\Frac(A/\underline{x}).
\]
But $\Frac(A/\underline{x})$ is not a finitely generated $A$-module, because
$\underline{x}$ is different from the maximal ideal of $A$. This is a
contradiction.
% END COMPONENT SGA2-VIII-CUM-C21

% VERIFIED SOURCE SEAM: French line 2855 is blank and excluded.

% BEGIN COMPONENT SGA2-VIII-CUM-C22
\medskip
\noindent Suppose that $n>1$, and assume the result established for every
$n'<n$. By the induction hypothesis, for every $x\in U$ such that $c(x)=1$,
one has $x\notin\Ass\cF_x$. Let $x$ be such a point, and let
$y\in\overline{\{x\}}\cap Y$ be an immediate specialization of $x$, i.e.
\[
  \dim\cO_{\overline{\{x\}},y}=1.
\]
Perform the base change
\[
  v\colon\Spec(\cO_{X,y})\longrightarrow X,
\]
retaining the notation of diagram~\textup{\hyperref[eq:VIII.2.7]{(2.7)}}. Applying EGA~III, 1.4.15,
one obtains isomorphisms
\[
  v^*(R^p i_*(\cF|_U))
  \simeq
  R^p i'_*(v'^*(\cF|_U)),
  \qquad p\in\mathbb Z.
\]

\noindent\emph{Thus we are reduced to the case in which $X$ is the spectrum
of a local ring $A$ in which $\underline{x}$ is a prime ideal of dimension
$1$}, i.e. $\dim(A/\underline{x})=1$. Set
\[
  \cF'=\SheafGamma_Y(\cF)
  \qquad\text{and}\qquad
  \cF''=\cF/\cF'.
\]
We have $\cF_x\simeq\cF''_x$ and $y\notin\Ass\cF''$. Moreover,
$\cF'|_U=0$; hence the long exact sequence of the $R^p i_*$ gives
isomorphisms
\[
  R^p i_*(\cF|_U)
  \simeq
  R^p i_*(\cF''|_U),
  \qquad p\in\mathbb Z.
\]

\noindent{\scriptsize\color{gray}French printed p.~96}\quad Since $n>1$,
neither $x$ nor $y$ belongs to $\Ass\cF''$. Now $\underline{x}$ and
$\underline{y}$ are the only prime ideals of $A$ containing $\underline{x}$.
It follows from Expos\'e~\hyperref[III.2.1]{III, 2.1} that there exists an element
$g\in\underline{x}$ that is $M$-regular, where we have set
\[
  \cF''=\widetilde M,
  \qquad M=\cF''(X).
\]
Thus there is an exact sequence
\[
  0\longrightarrow M\xrightarrow{g'}M\longrightarrow N\longrightarrow0,
\]
in which $g'$ denotes multiplication by $g$ on $M$. From the resulting long
exact sequence one sees that
\[
  R^p i_*(\widetilde N|_U)
  \quad\text{is coherent for }p<n-1.
\]
Therefore, by the induction hypothesis,
\[
  \prof(\widetilde N)_x\geq n-1,
  \qquad\text{and hence}\qquad
  \prof\cF_x\geq n.
\]
\hfill$\square$
% END COMPONENT SGA2-VIII-CUM-C22

% VERIFIED SOURCE SEAM: French line 2873 is blank and excluded.

% BEGIN COMPONENT SGA2-VIII-CUM-C23
\section*{3. Applications}
\sgareaderlabel{VIII.3}{3}

\noindent From these results one deduces a coherence condition for the higher
direct images of a coherent sheaf under a \emph{morphism that is not proper}.

\medskip
\sgareaderlabel{VIII.3.1}{3.1}\noindent\textbf{Theorem 3.1.}---Let $f\colon X\to Y$ be a morphism of
preschemes. Suppose that $Y$ is \emph{locally Noetherian} and that $f$ is
\emph{proper}. Suppose that $X$ is \emph{locally embeddable in a regular
prescheme}. Let $n\in\mathbb Z$. Let $U$ be an open subset of $X$, and let
$\cF$ be a \emph{coherent} $\cO_U$-module. Suppose that, for every $x\in U$
such that
\[
  \codimn(\overline{\{x\}}\cap(X\setminus U),
  \overline{\{x\}})=1,
\]
one has $\prof\cF_x\geq n$. \emph{Then} the $\cO_Y$-modules
\[
  R^p(f\circ g)_*(\cF)
\]
are coherent for $p<n$, where $g$ is the canonical immersion of $U$ into
$X$.

\medskip
\noindent Indeed, there is a Leray spectral sequence whose abutment is
$R^*(f\circ g)_*(\cF)$ and whose initial term is given by
\[
  E_2^{p,q}=R^p f_*(R^q g_*(\cF)).
\]
Moreover, there exists a \emph{coherent} $\cO_X$-module $\cG$ such that
$\cG|_U\simeq\cF$ (EGA~I, 9.4.3). It then follows from the preceding section
that condition~\textup{(iv)} on p.~74 is satisfied
{\scriptsize\color{gray}[French printed p.~97]}, i.e. that
$R^q g_*(\cG|_U)$ is coherent for $q<n$. Applying the finiteness theorem of
EGA~III, 3.2.1 to $f$ and the sheaves $R^q g_*(\cF)$, one finds that
$E_2^{p,q}$ is coherent for $q<n$, whence the conclusion.
% END COMPONENT SGA2-VIII-CUM-C23

% VERIFIED SOURCE SEAM: French line 2887 is blank and excluded.

% BEGIN COMPONENT SGA2-VIII-CUM-C24
\medskip
\sgareaderlabel{VIII.3.2}{3.2}\noindent\textbf{Proposition 3.2.}---Let $X$ be a locally Noetherian
prescheme, locally embeddable in a regular prescheme. Let $U$ be an open
subset of $X$, and let $i\colon U\to X$ be the canonical immersion. Let
$n\in\mathbb Z$. Finally, let $\cF$ be a coherent Cohen--Macaulay
$\cO_U$-module (on $U$!). The following conditions are equivalent:
\begin{enumerate}
\item[(a)] $R^p i_*(\cF)$ is coherent for $p<n$;
\item[(b)] for every irreducible component $S'$ of the closure $\overline S$
of the support $S$ of $\cF$, one has
\[
  \codimn(S'\cap(X\setminus U),S')>n.
\]
\end{enumerate}
% END COMPONENT SGA2-VIII-CUM-C24

% VERIFIED SOURCE SEAM: French line 2900 is blank and excluded.

% BEGIN COMPONENT SGA2-VIII-CUM-C25
\medskip
\noindent\textit{Proof, opening.}---Let $\cG$ be a coherent $\cO_X$-module
such that $\cG|_U\simeq\cF$ (EGA~I, 9.4.3). Applying \hyperref[VIII.2.3]{Corollary~2.3} to
$\cG$, we find that condition \textup{(a)} is equivalent to \textup{(c)}:
\begin{enumerate}
\item[(c)] for every $x\in S$, one has
\[
  \prof\cF_x>n-c(x),
  \qquad
  c(x)=\codimn(\overline{\{x\}}\cap (X-U),\overline{\{x\}}).
\]
\end{enumerate}
% END COMPONENT SGA2-VIII-CUM-C25

% VERIFIED SOURCE SEAM: French line 2908 is blank and excluded.

% BEGIN COMPONENT SGA2-VIII-CUM-C26
\medskip
\noindent\textit{Proof, \textup{(a)}$\Rightarrow$\textup{(b)}.}---Indeed,
let $S'$ be an irreducible component of $\overline S$, and let $s$ be its
generic point. Since $\cF$ is Cohen--Macaulay, one has
$\prof \cF_s=\dim\cO_{S,s}=0$. Moreover, $\overline{\{s\}}=S'$, and the
conclusion follows.
% END COMPONENT SGA2-VIII-CUM-C26

% VERIFIED SOURCE SEAM: French line 2910 is blank and excluded.

% BEGIN COMPONENT SGA2-VIII-CUM-C27
\medskip
\noindent\textit{Proof, \textup{(b)}$\Rightarrow$\textup{(a)}, opening.}---Let
$x\in S$, and let $S'$ be an irreducible component of $\overline S$ such that
$x\in S'$. Let $s$ be the generic point of $S'$. Since $\cF$ is
Cohen--Macaulay, we know that
\[
  \prof \cF_x=\dim\cO_{\overline{\{s\}},x}.
\]
% END COMPONENT SGA2-VIII-CUM-C27

% BEGIN COMPONENT SGA2-VIII-CUM-C28
\medskip
\noindent\textit{Proof, \textup{(b)}$\Rightarrow$\textup{(a)}, case
split.}---If $c(x)=+\infty$, there is nothing to prove. Otherwise, there exists
$y\in (X-U)\cap\overline{\{x\}}$ such that
\[
  c(x)=\dim\cO_{\overline{\{x\}},y}.
\]
% END COMPONENT SGA2-VIII-CUM-C28

% BEGIN COMPONENT SGA2-VIII-CUM-C29
\medskip
\noindent\textit{Proof, \textup{(b)}$\Rightarrow$\textup{(a)},
conclusion.}---Now $\cO_{X,y}$ is a quotient of a regular local ring by
hypothesis; hence
\[
  \dim\cO_{\overline{\{s\}},y}
  =\dim\cO_{\overline{\{s\}},x}
   +\dim\cO_{\overline{\{x\}},y}>n.
\]
\hfill$\square$
% END COMPONENT SGA2-VIII-CUM-C29

% VERIFIED SOURCE SEAM: French line 2923 is blank and excluded.

% INTEGRATION CONTROL BEFORE SGA2-VIII-CUM-C30
\begingroup

% BEGIN COMPONENT SGA2-VIII-CUM-C30
\renewcommand{\refname}{Bibliography}
\begin{thebibliography}{M}
\bibitem[M]{M}
\textsc{H.~Cartan \& S.~Eilenberg} -- \emph{Homological Algebra},
Princeton Math. Series vol.~19, Princeton University Press, 1956.
\end{thebibliography}
\endgroup
% END COMPONENT SGA2-VIII-CUM-C30

% END INCLUDED active_source/parts/08_expose_VIII.tex

% BEGIN INCLUDED active_source/parts/09_expose_IX.tex
\makeatletter
\providecommand{\sgareaderlabel}[2]{%
  \begingroup
  \def\@currentlabel{#2}%
  \phantomsection\label{#1}%
  \endgroup}
\makeatother

\chapterspace{-1}
\chapter{Algebraic Geometry and Formal Geometry}
\label{IX}

\medskip
The purpose of this expos\'e is to generalize Theorems 3.4.2 and 4.1.5 of
\emph{EGA III} to the case of a morphism that is not proper.

\section*{1. The Comparison Theorem}
\sgareaderlabel{IX.1}{1}

Let $f\colon X\to X'$ be a morphism of preschemes that is
\emph{separated} and of \emph{finite type}. Suppose that $X'$ is locally
noetherian. Let $Y'$ be a closed subset of $X'$, and let
$Y=f^{-1}(Y')$. Let $\widehat X$ and $\widehat{X'}$ be the formal
completions of $X$ and $X'$ along $Y$ and $Y'$, respectively. Let
$\widehat f$ be the morphism induced by $f$ upon passing to the completions.

\begin{equation}
\tag{1.1}\label{eq:IX.1.1}
\begin{tikzcd}[column sep=3.1em,row sep=2.6em]
X \arrow[d,"f"'] & Y \arrow[l] \arrow[d] &&
\widehat{X} \arrow[d,"\widehat f"'] \arrow[r,"j"] & X \arrow[d,"f"] \\
X' & Y' \arrow[l] & {,} &
\widehat{X'} \arrow[r,"i"] & X'.
\end{tikzcd}
\end{equation}

We denote by $j$ (resp. $i$) the morphism from $\widehat X$ to $X$
(resp. from $\widehat{X'}$ to $X'$). It is known that $i$ and $j$ are
\emph{flat}.

Let $\mathcal I'$ be an ideal of definition of $Y'$, and let
$\mathcal J=f^*(\mathcal I')\cdot\mathcal O_X$; this is an ideal of
definition of $Y$. Hence
\begin{equation}
\tag{1.2}\label{eq:IX.1.2}
\widehat{X'}=
\left(Y',\varprojlim_{k\in\mathbb N}
\mathcal O_{X'}/{\mathcal I'}^{k+1}\right),
\qquad
\widehat X=
\left(Y,\varprojlim_{k\in\mathbb N}
\mathcal O_X/\mathcal J^{k+1}\right).
\end{equation}

For every $k\in\mathbb N$, set
\begin{equation}
\tag{1.3}\label{eq:IX.1.3}
Y'_k=\left(Y',\mathcal O_{X'}/(\mathcal I')^{k+1}\right),
\qquad
Y_k=\left(Y,\mathcal O_X/\mathcal J^{k+1}\right).
\end{equation}

Let $F$ be a coherent $\mathcal O_X$-module. For every $k\in\mathbb N$, set
\begin{equation}
\tag{1.4}\label{eq:IX.1.4}
F_k=F/\mathcal J^{k+1}F,
\qquad
\widehat F=j^*(F)\simeq\varprojlim F_k.
\end{equation}

If we set
\begin{equation}
\tag{1.5}\label{eq:IX.1.5}
R^i f_*(F)^{\wedge}
=\varprojlim_{k\in\mathbb N}
\left(R^i f_*(F)\otimes_{\mathcal O_{X'}}\mathcal O_{Y'_k}\right),
\qquad i\in\mathbb Z,
\end{equation}
there is a natural homomorphism
\begin{equation}
\tag{1.6}\label{eq:IX.1.6}
r_i:i^*(R^i f_*(F))\longrightarrow R^i f_*(F)^{\wedge},
\end{equation}
which is an \emph{isomorphism} when $R^i f_*(F)$ is \emph{coherent}.

As explained in \emph{EGA III}, 4.1.1, there is a commutative diagram
\begin{equation}
\tag{1.7}\label{eq:IX.1.7}
\begin{tikzcd}[column sep=5.2em,row sep=3.2em]
i^*(R^i f_*(F))
  \arrow[r,"\rho_i"]
  \arrow[d,"r_i"]
& R^i\widehat f_*(\widehat F)
  \arrow[d,"\psi_i"] \\
R^i f_*(F)^{\wedge}
  \arrow[r,"\varphi_i"]
& \displaystyle\varprojlim_{k\in\mathbb N} R^i f_*(F_k).
\end{tikzcd}
\end{equation}

In the cited passage one obtains a commutative diagram because
$R^i f_*(F)$ is known to be coherent, and the source and target of~\hyperref[eq:IX.1.6]{(1.6)}
are identified. In our case, $R^i f_*(F)$ will be coherent only for certain
values of $i$, for which we shall study~\hyperref[eq:IX.1.7]{(1.7)}.

Consider the graded ring
\begin{equation}
\tag{1.8}\label{eq:IX.1.8}
\mathcal S=\bigoplus_{k\in\mathbb N}{\mathcal I'}^k,
\end{equation}
and the graded $\mathcal S$-module
\begin{equation}
\tag{1.9}\label{eq:IX.1.9}
\mathcal H^i=\bigoplus_{k\in\mathbb N}R^i f_*(\mathcal J^kF),
\qquad i\in\mathbb Z,
\end{equation}
whose $\mathcal S$-module structure is defined as follows.

The sheaf $R^i f_*(\mathcal J^kF)$ is associated with the presheaf that
assigns to each affine open subset $U'$ of $X'$
\begin{equation}
\tag{1.10}\label{eq:IX.1.10}
H^i\!\left(f^{-1}(U'),
  \left.\mathcal J^kF\right|_{f^{-1}(U')}\right).
\end{equation}
Thus let $U'$ be an affine open subset of $X'$, set
\[
U=f^{-1}(U'),
\]
and let $x'\in{\mathcal I'}^m(U')$. Let $x$ be the image of $x'$ in
$\mathcal J^m(U)$. Multiplication by $x$ on $F|_U$ maps
$\mathcal J^kF|_U$ into $\mathcal J^{k+m}F|_U$ and hence, by
functoriality, gives a morphism
\begin{equation}
\tag{1.11}\label{eq:IX.1.11}
\mu^i_{x',k}(U'):
H^i(U,\mathcal J^kF|_U)
\longrightarrow
H^i(U,\mathcal J^{k+m}F|_U),
\end{equation}
defined for every $i\in\mathbb Z$ and every $k\in\mathbb N$. Upon passage
to the associated sheaf, these morphisms give $\mathcal H^i$ its graded
$\mathcal S$-module structure.

\begin{theoremIXoneone}\sgareaderlabel{IX.1.1}{1.1}
Let $n$ be an integer. Suppose that the graded $\mathcal S$-module
$\mathcal H^i$ is finitely generated for $i=n-1$ and $i=n$. Then:
\begin{enumerate}
\renewcommand{\labelenumi}{(\arabic{enumi})}
\setcounter{enumi}{-1}
\item
$r_n$ and $r_{n-1}$ are isomorphisms, and
$R^i\widehat f_*(\widehat F)$ is coherent for $i=n-1$;

\item
for $i=n-1$, the maps $\rho_i$, $\varphi_i$, and $\psi_i$ are topological
isomorphisms. In particular, the filtration on $R^{n-1}f_*(F)$ defined by
the kernels of the homomorphisms
\begin{equation}
\tag{1.12}\label{eq:IX.1.12}
R^{n-1}f_*(F)\longrightarrow R^{n-1}f_*(F_k)
\end{equation}
is $\mathcal I'$-good;

\item
for $i=n$, the maps $\rho_i$, $\varphi_i$, and $\psi_i$ are monomorphisms;
moreover, the filtration on $R^n f_*(F)$ is $\mathcal I'$-good, and $\psi_n$
is an isomorphism;

\item
for $i=n-2,n-1$, the projective system $R^i f_*(F_k)$ satisfies the uniform
Mittag--Leffler condition; that is, there is a \emph{fixed} integer $k\geq0$
such that, for every $p\geq0$ and every $p'\geq p+k$,
\[
\operatorname{Im}\!\left[R^i f_*(F_{p'})\longrightarrow R^i f_*(F_p)\right]
=
\operatorname{Im}\!\left[R^i f_*(F_{p+k})\longrightarrow R^i f_*(F_p)\right].
\]
\end{enumerate}
\end{theoremIXoneone}

Proceeding as in \emph{EGA III}, 4.1.8, it is easy to \emph{reduce to the case
where $X'$ is the spectrum of a noetherian ring} $A$. In this case, it is
known that
\begin{equation}
\tag{1.13}\label{eq:IX.1.13}
R^i f_*(F)=\widetilde{H^i(X,F)}
\qquad\text{(cf.~\hyperref[eq:IX.1.10]{(1.10)}).}
\end{equation}
Let $I$ be the ideal of $A$ such that $\widetilde I=\mathcal I'$, and set
\begin{equation}
\tag{1.14}\label{eq:IX.1.14}
S=\bigoplus_{k\in\mathbb N}I^k,
\end{equation}
\begin{equation}
\tag{1.15}\label{eq:IX.1.15}
H^i=\bigoplus_{k\in\mathbb N}H^i(X,\mathcal J^kF),
\qquad i\in\mathbb Z,
\end{equation}
where $H^i$ carries the graded $S$-module structure defined by~\hyperref[eq:IX.1.11]{(1.11)}, with
$U'=X'$.
\footnote{The corrected French TeX and same-edition PDF print $U=X'$ here.
In the construction preceding~(1.11), however, $U=f^{-1}(U')$; specializing
the affine open to the whole base therefore requires $U'=X'$. The French
authority is left unchanged.}

The proof is modeled on that of \emph{EGA III}, 4.1.5; let us give a summary.

We now work with $\varphi_i$ and $\psi_i$, to which the following module
homomorphisms correspond:
\begin{equation}
\tag{1.16}\label{eq:IX.1.16}
\begin{tikzcd}[column sep=5.2em,row sep=3.0em]
& H^i(\widehat X,\widehat F)
    \arrow[d,"\psi_i"] \\
H^i(X,F)^{\wedge}
    \arrow[r,"\varphi_i"]
& \displaystyle\varprojlim_{k\in\mathbb N}H^i(X,F_k).
\end{tikzcd}
\end{equation}

\noindent\textup{(a)} Suppose only that $H^i$ is a graded $S$-module of finite
type. It follows that \emph{the filtration on $H^i(X,F)$ defined by the
modules}
\begin{equation}
\tag{1.17}\label{eq:IX.1.17}
R_k^i=\ker\bigl(H^i(X,F)\longrightarrow H^i(X,F_k)\bigr)
\end{equation}
\emph{is $I$-good}. Indeed, the exact cohomology sequence
\begin{equation}
\tag{1.18}\label{eq:IX.1.18}
H^i(X,\mathcal J^{k+1}F)\longrightarrow H^i(X,F)
  \longrightarrow H^i(X,F_k)
\end{equation}
shows that the \emph{graded $S$-module}
$\bigoplus_{k\in\mathbb N}R_k^i$ \emph{is a quotient} of the graded
$S$-submodule
\[
\bigoplus_{k\in\mathbb N}H^i(X,\mathcal J^{k+1}F)
\]
of $H^i$, and hence \emph{is of finite type}, since $S$ is noetherian. This
proves the first assertion.

Set
\begin{equation}
\tag{1.19}\label{eq:IX.1.19}
M^i=H^i(X,F),\qquad H_k^i=H^i(X,F_k).
\end{equation}
There is a commutative diagram
\begin{equation}
\tag{1.20}\label{eq:IX.1.20}
\begin{tikzcd}[column sep=5.5em,row sep=3.2em]
H^i(X,F)^{\wedge}
  \arrow[r,"s_i"]
  \arrow[dr,"\varphi_i"']
& \displaystyle\varprojlim_k(M^i/R_k^i)
  \arrow[d,"t_i"] \\
& \displaystyle\varprojlim_k H_k^i,
\end{tikzcd}
\end{equation}
in which $s_i$ is an \emph{isomorphism}, since the filtration of $H^i(X,F)$
is $I$-good. Moreover, $t_i$ is a \emph{monomorphism}: the functor
$\varprojlim$ is left exact and, for every $k\geq 0$, the natural morphism
$M^i/R_k^i\longrightarrow H_k^i$ is a monomorphism by the definition of
$R_k^i$.

To study the surjectivity of $t_i$, introduce
\begin{equation}
\tag{1.21}\label{eq:IX.1.21}
Q_k^i=\operatorname{coker}\bigl(H^i(X,F)\longrightarrow H^i(X,F_k)\bigr),
\end{equation}
which gives a projective system of exact sequences
\begin{equation}
\tag{1.22}\label{eq:IX.1.22}
0\longrightarrow R_k^i\longrightarrow M^i\longrightarrow H_k^i
 \longrightarrow Q_k^i\longrightarrow 0.
\end{equation}
Using the exact cohomology sequence
\begin{equation}
\tag{1.23}\label{eq:IX.1.23}
H^i(X,F)\longrightarrow H^i(X,F_k)
 \longrightarrow H^{i+1}(X,\mathcal J^{k+1}F),
\end{equation}
we see that the \emph{graded $S$-module}
\begin{equation}
\tag{1.24}\label{eq:IX.1.24}
Q^i=\bigoplus_{k\in\mathbb N}Q_k^i
\end{equation}
\emph{is a graded $S$-submodule of $H^{i+1}$}. Moreover, for every
$k\geq0$,
\begin{equation}
\tag{1.25}\label{eq:IX.1.25}
I^{k+1}Q_k^i=0,
\end{equation}
because $Q_k^i$ is the image of $H_k^i$.

\sgareaderlabel{stepb}{b}\noindent\textup{(b)} Suppose only that $H^{i+1}$ is of finite type, and
consider $t_i$ alone, disregarding $s_i$. Since $S$ is noetherian, $Q^i$ is of
finite type; as $I^{k+1}Q_k^i=0$, there exist integers $r\geq0$ and $k_0\geq0$
such that
\begin{equation}
\tag{1.26}\label{eq:IX.1.26}
I^rQ_k^i=0\qquad\text{for }k\geq k_0.
\end{equation}
It follows that the \emph{projective system $(Q_k^i)_{k\in\mathbb N}$ is
essentially zero}, and hence that the \emph{projective system
$(H_k^i)_{k\in\mathbb N}$ satisfies the uniform Mittag--Leffler condition}.
From the exact sequence~\hyperref[eq:IX.1.22]{(1.22)} we obtain the exact sequence
\begin{equation}
\tag{1.27}\label{eq:IX.1.27}
0\longrightarrow M^i/R_k^i\longrightarrow H_k^i
 \longrightarrow Q_k^i\longrightarrow0,
\end{equation}
and therefore the exact sequence
\begin{equation}
\tag{1.28}\label{eq:IX.1.28}
0\longrightarrow\varprojlim_k M^i/R_k^i
 \xrightarrow{t_i}\varprojlim_k H_k^i
 \longrightarrow\varprojlim_k Q_k^i.
\end{equation}
But the projective system $(Q_k^i)_{k\in\mathbb N}$ is essentially zero, so
$t_i$ is an \emph{isomorphism}.

\noindent\textup{(c)} Let us prove that if $H^i$ is of finite type, then
$\psi_i$ is an isomorphism. It suffices to apply \emph{EGA}~$0_{\mathrm{III}}$,
13.3.1, taking the affine open subsets as a basis of open subsets of $X$.
This is legitimate: by~(b), the projective system
$(H_k^{i-1})_{k\in\mathbb N}$ satisfies the Mittag--Leffler condition.

The theorem follows formally from~(a), (b), and~(c). Observe that, in fact, at
each step the proof uses finite generation of $H^i$ for only one value of $i$.

Let us give some examples in which the hypothesis of \hyperref[IX.1.1]{Theorem~1.1} is satisfied.

\medskip
\sgareaderlabel{IX.1.2}{1.2}\noindent\textbf{Corollary 1.2.}\quad
\emph{Suppose that $\mathcal I'$ is generated by a section $t'$ of
$\mathcal O_{X'}$, and denote by $t$ the corresponding section of
$\mathcal O_X$. Let $F$ be a coherent $\mathcal O_X$-module, and let $n$ be
an integer. Assume that:}
\begin{enumerate}
\item[(i)] \emph{$t$ is $F$-regular (that is, multiplication by $t$ on $F$
is a monomorphism);}
\item[(ii)] \emph{$R^i f_*(F)$ is coherent for $i=n-1$ and $i=n$.}
\end{enumerate}
\emph{Then the hypothesis of \hyperref[IX.1.1]{Theorem~1.1} is satisfied.}

Indeed, multiplication by $t^k$ defines an isomorphism
$F\xrightarrow{\sim}\mathcal J^kF$, from which it follows that
\begin{equation}
\tag{1.29}\label{eq:IX.1.29}
H^i\simeq R^i f_*(F)\otimes_{\mathcal O_{X'}}\mathcal O_{X'}[T],
\end{equation}
where $T$ is an indeterminate. This proves the result.

\sgareaderlabel{IX.1.3}{1.3}\noindent\textbf{Corollary 1.3.}\quad
\emph{Suppose that $X'=\operatorname{Spec}(A)$, where $A$ is a noetherian ring
that is separated and complete for the $I$-adic topology. Suppose that the
graded $S$-module $H^i$ is of finite type for $i=n-1$ and $i=n$
(cf.~\hyperref[eq:IX.1.14]{(1.14)} and~\hyperref[eq:IX.1.15]{(1.15)}). Then the hypotheses of \hyperref[IX.1.1]{Theorem~1.1} are satisfied, and
there is a commutative diagram of isomorphisms}
\begin{equation}
\tag{1.30}\label{eq:IX.1.30}
\begin{tikzcd}[column sep=4.6em,row sep=3.0em]
H^i(X,F)
  \arrow[rr,"\rho_i'"]
  \arrow[dr,"\varphi_i'"']
&& H^i(\widehat X,\widehat F)
  \arrow[dl,"\psi_i"] \\
& \displaystyle\varprojlim_k H^i(X,F_k)
&& \text{for }i=n-1.
\end{tikzcd}
\end{equation}
\noindent{\footnotesize\emph{Source-numbering note
(SGA2-IX-EQ130-RENDERDEF-001).} The corrected French TeX
labels this diagram \texttt{eq:IX.1.30}, the following proof cites (1.30), and
the next numbered formula is (1.31). The same-edition rendered PDF nevertheless
prints (1.1), apparently because its subsection step resets the equation
counter. The French files are left unchanged; the intended tag (1.30) is used
here.}

\section*{Editor's note (1)}
In the same vein, see the article by W.-L. Chow, ``Formal functions on
homogeneous spaces,'' \emph{Invent. Math.} \textbf{86} (1986), no.~1,
pp.~115--130. The author proves the following result. Let $X$ be an algebraic
variety over a field, homogeneous under an algebraic group $G$, and let $Z$ be
a complete subvariety of $X$ of dimension $>0$. Suppose that $Z$ generates $X$
in the following sense: given $p\in Z$, let $\Gamma_p$ be the set of elements
of $G$ that send $p$ into $Z$. One says that $Z$ generates $X$ if the group
generated by the identity component of $\Gamma_p$ is all of $G$. In that case,
every formal rational function on $X$ along $Z$ is algebraic; compare the
results of Hironaka and Matsumura cited in editor's note~(3), p.~138.

Along the lines of the techniques introduced by those authors, we also mention
the elegant algebraization result due to D.~Gieseker, ``On two theorems of
Griffiths about embeddings with ample normal bundle,'' \emph{Amer. J. Math.}
\textbf{99} (1977), no.~6, pp.~1137--1150, Theorems~4.1 and~4.2. Let $X$ be a
connected projective variety of dimension $>0$, locally a complete intersection
over an algebraically closed field. Suppose there are two embeddings of $X$
into smooth projective varieties $Y$ and $W$. If the formal completions of $X$
in $Y$ and $W$ are equivalent, then there exists a scheme $U$ containing $X$
as a closed subscheme which embeds into both $Y$ and $W$ as an étale
neighborhood of $X$. In other words, formal equivalence implies étale
equivalence.

See also G.~Faltings, ``Formale Geometrie und homogene Räume,''
\emph{Invent. Math.} \textbf{64} (1981), pp.~123--165.

\medskip
\noindent{\footnotesize\emph{Source-defect note
(SGA2-IX-L3156-PUNCT-SRCDEF-001;
SGA2-IX-L3157-FALTINGS-TITLE-SRCDEF-001).} The French TeX and its
same-edition PDF omit the sentence stop before ``Autrement dit.'' They also
print the Faltings title in the mixed and misspelled form ``Formale Geometrie
and homogene Raüme.'' The target restores the sentence boundary and the
published German title; the French authority remains unchanged.}

\noindent\emph{Proof.}\quad
We simply note that $H^i(X,F)$ is of finite type and hence isomorphic to its
completion. One obtains~\hyperref[eq:IX.1.30]{(1.30)} by transcribing the diagram of modules~\hyperref[eq:IX.1.7]{(1.7)}
into the category of $A$-modules and replacing the left-hand vertical by
$H^i(X,F)$.

\medskip
\noindent{\footnotesize\emph{Source-numbering note
(SGA2-IX-EQ130-RENDERDEF-001).} The corrected French
TeX cites its label \texttt{eq:IX.1.30}; the next numbered formula is (1.31).
The same-edition rendered PDF prints the referenced diagram as (1.1) because
of an apparent counter reset. The French authority remains unchanged; the
intended reference (1.30) is retained here.}

\sgareaderlabel{IX.1.4}{1.4}\noindent\textbf{Proposition 1.4.}\quad
{\itshape Let $A$ be a noetherian ring. Let $t\in A$, and suppose that $A$ is
separated and complete for the $(tA)$-adic topology. Set
\begin{equation}
\tag{1.31}\label{eq:IX.1.31}
X'=\operatorname{Spec}(A),\qquad Y'=V(t),\qquad I=(tA).
\end{equation}
Let $T$ be a closed subset of $X'$, and set
\begin{equation}
\tag{1.32}\label{eq:IX.1.32}
X=X'\setminus T,\qquad
Y=Y'\cap X=Y'\setminus(Y'\cap T).
\end{equation}
Let $F$ be a coherent $\mathcal O_X$-module. Finally, set
\begin{equation}
\tag{1.33}\label{eq:IX.1.33}
T'=\bigl\{x\in X'\mid
\codimop\bigl(\overline{\{x\}}\cap T,\overline{\{x\}}\bigr)=1\bigr\}.
\end{equation}
Suppose that
\begin{enumerate}
\renewcommand{\labelenumi}{\alph{enumi})}
\item $t$ is $F$-regular;
\item $\prof_{T'}(F)\ge n+1$;
\item $A$ is a quotient of a regular noetherian ring.
\end{enumerate}
Then, in diagram~\hyperref[eq:IX.1.30]{(1.30)}, the morphisms $\rho'_i$, $\varphi'_i$, and $\psi_i$
are isomorphisms for $i<n$ and monomorphisms for $i=n$. Moreover, $\psi_n$
is an isomorphism.}

\medskip
\noindent{\footnotesize\emph{Source-notation and numbering note
(SGA2-IX-L3177-CLOSURE-BRACING-SRCDEF-001;
SGA2-IX-EQ130-RENDERDEF-001).} In (1.33)
the French TeX and its same-edition PDF put braces around $\overline{x}$,
whereas the defining codimension formula in Expos\'e VIII, Theorem~2.1, uses
$\overline{\{x\}}$. The target makes that intended closure-of-the-point
notation explicit; French remains unchanged. French line 3189 cites
\texttt{eq:IX.1.30}; the target retains (1.30), rather than the same-edition
PDF's accidental counter-reset display (1.1).}

\noindent\emph{Proof.}\quad
By~\hyperref[eq:IX.1.3]{(1.3)} and~\hyperref[eq:IX.1.2]{(1.2)}, it suffices to prove that $R^i f_*(F)$ is coherent for
$i\le n$; this follows from the finiteness theorem, Expos\'e~VIII,
\hyperref[IX.2.1]{Theorem~2.1}.

\medskip
\noindent In particular:

\sgareaderlabel{IX.1.5}{1.5}\noindent\textbf{Example 1.5.}\quad
\hyperref[IX.1.4]{Proposition~1.4} applies when $A$ is a local ring and $t$ belongs to the
radical $\mathfrak r(A)$ of $A$. We then take
$T=\{\mathfrak r(A)\}$. In this case, for $n=1$, we obtain the following
statement:

\medskip
{\itshape If $A$ is noetherian, separated and complete for the $t$-adic
topology, and is a quotient of a regular ring (for example, if $A$ is
complete), if moreover $t$ is $F$-regular and $\prof F_x\ge2$ for every
$x\in\operatorname{Spec}(A)$ such that $\dim(A/x)=1$, then the natural
homomorphism
\[
  \Gamma(X,F)\longrightarrow\Gamma(\widehat X,\widehat F)
\]
is an isomorphism.}

\medskip
Indeed, retaining the notation of \hyperref[IX.1.4]{Proposition~1.4}, we have
$T=\{\mathfrak r(A)\}$, and \hyperref[eq:IX.1.33]{formula~(1.33)} says that
\[
  T'=\{x\in\operatorname{Spec}(A)\mid\dim(A/x)=1\}.
\]

\medskip
\noindent{\footnotesize\emph{Source-cross-reference note
(SGA2-IX-EXAMPLE15-SRCDEF-001).} French lines 3196
and 3202 use the key \texttt{eq:IX.1.4}, which formally names equation~(1.4),
although the grammar and the retained hypotheses/notation unambiguously refer
to Proposition~1.4. The same-edition PDF displays only the number `1.4' and
does not resolve the key. The target makes the semantic reference explicit;
the French authority remains unchanged.}

\setcounter{section}{1}
\section{The Existence Theorem}\label{IX.2}

Let us first state EGA~III, 3.4.2 in a slightly more general form.

Let $f:\mathfrak X\to\mathfrak X'$ be an adic
morphism\markedfootnote{(*)}{This hypothesis is not essential; cf.
Expos\'e~XII, p.~118.} of formal preschemes, with $\mathfrak X'$ noetherian.
Let $\mathcal I'$ be an ideal of definition of $\mathfrak X'$; since $f$ is
adic,
$f^*\mathcal I'=\mathcal J$ is\markedfootnote{(2)}{\textit{Editor's note.}
By the very definition; cf. EGA~I, 10.12.1.} an ideal of definition of
$\mathfrak X$.

For every $n\in\mathbb N$, set
\begin{equation}\label{eq:IX.2.1}
  \mathfrak X_n=
  \bigl(\mathfrak X,\mathcal O_{\mathfrak X}/\mathcal J^{n+1}\bigr);
\end{equation}
this is an ordinary prescheme with the same underlying topological space as
$\mathfrak X$.

Let $\mathfrak F$ be a coherent $\mathcal O_{\mathfrak X}$-module. For every
$k\in\mathbb N$, the $\mathcal O_{\mathfrak X_k}$-modules
\begin{equation}
\tag{2.2}\label{eq:IX.2.2}
F_k=\mathfrak F/\mathcal J^{k+1}\mathfrak F
\end{equation}
are coherent. For every $i$, there is a homomorphism
\begin{equation}
\tag{2.3}\label{eq:IX.2.3}
\psi_i:R^i f_*(\mathfrak F)
  \longrightarrow\varprojlim_k R^i f_*(F_k),
\end{equation}
induced by functoriality from the natural homomorphism
\begin{equation}
\tag{2.4}\label{eq:IX.2.4}
\mathfrak F\longrightarrow F_k.
\end{equation}

Set
\begin{equation}
\tag{2.5}\label{eq:IX.2.5}
\mathbf S
  =\operatorname{gr}_{\mathcal I'}\mathcal O_{\mathfrak X'}
  =\bigoplus_{k\in\mathbb N}\mathcal I'^{k}/\mathcal I'^{k+1},
\end{equation}
\begin{equation}
\tag{2.6}\label{eq:IX.2.6}
\operatorname{gr}_{\mathcal J}(\mathfrak F)
  =\bigoplus_{k\in\mathbb N}
    \mathcal J^{k}\mathfrak F/\mathcal J^{k+1}\mathfrak F,
\end{equation}
\begin{equation}
\tag{2.7}\label{eq:IX.2.7}
\mathbf K^{i}
  =R^{i}f_*\bigl(\operatorname{gr}_{\mathcal J}(\mathfrak F)\bigr)
  =\bigoplus_{k\in\mathbb N}
    R^{i}f_*\bigl(\mathcal J^{k}\mathfrak F/
      \mathcal J^{k+1}\mathfrak F\bigr).
\end{equation}
It is clear that $\mathbf K^{i}$ is endowed with the structure of a graded
$\mathbf S$-module.

\medskip
\noindent\textbf{Theorem 2.1.}\sgareaderlabel{IX.2.1}{2.1}\quad
{\itshape Suppose that $\mathbf K^i$ is a graded $\mathbf S$-module of finite
type for $i=n-1$, $i=n$, and $i=n+1$. Then:
\begin{enumerate}
\renewcommand{\labelenumi}{(\roman{enumi})}
\setlength{\itemsep}{0.35em}
\item $R^n f_*(\mathfrak F)$ is coherent.
\item The homomorphism $\psi_n$ in~\hyperref[eq:IX.2.3]{(2.3)} is a topological isomorphism. The
natural filtration on the right-hand side of~\hyperref[eq:IX.2.3]{(2.3)} is $\mathcal I'$-good.
\item The projective system of the $R^n f_*(F_k)$ satisfies the uniform
Mittag-Leffler condition.
\end{enumerate}}

\medskip
\noindent\textit{Proof.}\quad
The proof follows very easily from EGA~$0_{\mathrm{III}}$, 13.7.7
(cf. EGA~III, 3.4.2), provided that the text on page~78 is corrected as
indicated\sgareaderlabel{pagerectif}{11} in (EGA~III~2, $\mathrm{Err}_{\mathrm{III}}\,24$).

\medskip
\sgareaderlabel{IX.2.2}{2.2}\noindent\textbf{Theorem 2.2}\textsuperscript{(3)}.\quad
{\itshape Let $A$ be a noetherian adic ring, and let $I$ be an ideal of
definition of $A$. Let $T$ be a closed subset of
$X'=\operatorname{Spec}(A)$. Suppose that $I$ is generated by an element
$t\in A$. Retain the notation of~\hyperref[eq:IX.1.31]{(1.31)}, \hyperref[eq:IX.1.32]{(1.32)}, and~\hyperref[eq:IX.1.33]{(1.33)}. Let
$\mathfrak F$ be a coherent $\mathcal O_{\widehat X}$-module. Set
\begin{equation}
\tag{2.8}\label{eq:IX.2.8}
F_0=\mathfrak F/\mathcal J\mathfrak F,
\end{equation}
\par\smallskip
\noindent{\normalfont\footnotesize\color{notegray}
\textbf{Source-render note
[SGA2-IX-T22-EQ28-COUNTER-SRCDEF-001].}
The admitted French counter order renders this display as~(2.1), although
its label and the surrounding sequence require~(2.8). The English target
uses~(2.8); the French authority remains unchanged.}
\par\smallskip
where $\mathcal J=t\mathcal O_{\widehat X}$ is an ideal of definition of
$\widehat X$. \emph{Suppose that} $A$ is a quotient of a regular
noetherian ring and that:
\begin{enumerate}
\renewcommand{\labelenumi}{(\arabic{enumi})}
\setlength{\itemsep}{0.35em}
\item $t$ is $\mathfrak F$-regular;
\item $\operatorname{prof}_{T'}F_0\geq 2$.
\end{enumerate}
Then there exists a coherent $\mathcal O_X$-module $F$ such that
$\widehat F\simeq\mathfrak F$.}

\medskip
\noindent\colorbox{softgray}{\parbox{0.95\linewidth}{\small
\color{notegray}\textbf{Editor's note (3).}
Numerous algebraization statements have been obtained since, without
mentioning those cited below; see the articles by Faltings or by Mme Raynaud
cited, respectively, in editor's note~(22) on p.~155 and editor's note~(7)
on p.~203.

\smallskip
{\footnotesize\textbf{Source-defect note
[SGA2-IX-L3265-ONT-GRAMMAR-SRCDEF-001].}
The admitted French reads \emph{on \'et\'e obtenus}; the plural subject
requires \emph{ont \'et\'e obtenus}. The English phrase ``have been
obtained'' gives the intended plural sense; the French authority remains
unchanged.}

\smallskip
One has in mind in particular Artin's results (see especially
Artin M., ``Algebraization of formal moduli. I,'' in \emph{Global Analysis
(Papers in Honor of K.~Kodaira)}, Univ. Tokyo Press, Tokyo, 1969,
pp.~21--71), but also recent results on the algebraicity of leaves of
foliations; see especially Bost J.-B., ``Algebraic leaves of algebraic
foliations over number fields,'' \emph{Publ. Math. Inst. Hautes
\'Etudes Sci.} \textbf{93} (2001), pp.~161--221, and Chambert-Loir A.,
``Th\'eor\`emes d'alg\'ebricit\'e en g\'eom\'etrie diophantienne (d'apr\`es
J.-B. Bost, Y.~Andr\'e, D.~\& G.~Chudnovsky),'' in \emph{S\'eminaire
Bourbaki, Vol. 2000/2001}, \emph{Ast\'erisque}, vol.~282, Soci\'et\'e
math\'ematique de France, Paris, 2002, Exp.~886, pp.~175--209, and the
references cited there. In particular, the latter two articles contain
discussions of the connection between algebraization questions and the
theory of Diophantine approximation.}}

\medskip
\noindent\textit{Proof.}\quad
It suffices to prove that $\widehat f_*(\mathfrak F)$ is a coherent
$\mathcal O_{\widehat{X'}}$-module, where
$\widehat f\colon\widehat X\to\widehat{X'}$ is the morphism of formal
preschemes obtained from the inclusion of $X$ into $X'$ by completion with
respect to $t$. Indeed, $A$ is separated and complete for the $t$-adic
topology, so there exists an $A$-module $F'$ whose completion is isomorphic
to $\widehat f_*(\mathfrak F)$. Since $X$ is open in $X'$, we may take
\[
F=\widetilde{F'}\big|_X.
\]

\noindent\colorbox{softgray}{\parbox{0.95\linewidth}{\footnotesize
\color{notegray}\textbf{Source-defect note
[SGA2-IX-L3280-PUSHFORWARD-BASE-SRCDEF-001].}
The French TeX and the same-edition reader print
$\mathcal O_{\widehat X}$ in the first sentence. Since the pushforward along
$\widehat f\colon\widehat X\to\widehat{X'}$ is a sheaf on the target, the
English target uses the type-correct $\mathcal O_{\widehat{X'}}$. The French
authority remains unchanged under manager decision
EG-SGA2-IX-T22-SOURCE-DEFECT-ADJUDICATION-20260719-0001.}}

\medskip
It remains to show that \hyperref[IX.2.1]{Theorem~2.1} applies to the morphism of formal
preschemes $\widehat f$ and to $\mathfrak F$. By hypothesis~\hyperref[IX.2.2]{(1)}, for every
$k\in\mathbb N$ there is an isomorphism
\[
\mathcal J^k\mathfrak F/\mathcal J^{k+1}\mathfrak F
  \to \mathfrak F/\mathcal J\mathfrak F,
\]
and hence the hypothesis of \hyperref[IX.2.1]{Theorem~2.1} is satisfied provided that
\[
R^i f_*(F_0)\text{ is coherent for }i\leq 1.
\]
This follows from~\hyperref[IX.2.2]{(2)} and the finiteness theorem of Expos\'e~VIII,
\hyperref[IX.2.1]{Theorem~2.1}. This proves the result.

\medskip
\noindent\colorbox{softgray}{\parbox{0.95\linewidth}{\footnotesize
\color{notegray}\textbf{Source-notation observation
[SGA2-IX-L3283-PLAIN-ARROW-OBS-001].}
The corrected TeX and its same-edition reader use a plain $\to$ in the
display while the preceding prose explicitly calls the map an isomorphism.
This is an observation, not a confirmed defect. The target retains both
features; no $\sim$ has been added over the arrow.}}

\medskip
It remains to specialize this statement under the assumption that $A$ is a
local ring.

\medskip
\sgareaderlabel{IX.2.3}{2.3}\noindent\textbf{Corollary 2.3.}\quad
\emph{Let $A$ be a noetherian local ring, and let
$t\in\mathfrak r(A)$ be an element of the radical of $A$. Suppose that $A$
is separated and complete for the $t$-adic topology and, moreover, is a
quotient of a regular ring (for example, suppose that $A$ is a complete
noetherian local ring). Set}
\begin{equation}
\tag{2.9}\label{eq:IX.2.9}
X'=\operatorname{Spec}A,\qquad T=\{\mathfrak r(A)\},
\end{equation}
\emph{and retain the notation of \hyperref[eq:IX.1.31]{(1.31)}, \hyperref[eq:IX.1.32]{(1.32)}, and \hyperref[eq:IX.1.33]{(1.33)}. Let
$\mathfrak F$ be a coherent $\mathcal O_{\widehat X}$-module. Suppose that}
\begin{enumerate}
\item $t$ is $\mathfrak F$-regular;
\item $\prof_{T'}F_0\geq 2$, where
$F_0=\mathfrak F/\mathcal J\mathfrak F$ and
$\mathcal J=t\mathcal O_{\widehat X}$.
\end{enumerate}
\emph{Then there exists a coherent $\mathcal O_X$-module $F$ such that
$\widehat F\simeq\mathfrak F$.}

\medskip
\noindent\colorbox{softgray}{\parbox{0.95\linewidth}{\footnotesize
\color{notegray}\textbf{Source-structure note
(SGA2-IX-L3304-ENUM-STRUCTURE-SRCDEF-001).}
The French TeX closes its \texttt{enumerate} environment only after the
conclusion, syntactically nesting ``Alors il existe'' inside item~(2). The
same-edition PDF preserves the source order but visually masks this nesting:
the conclusion begins at the corollary margin, not at the item-text margin.
The conclusion is consequent on both hypotheses, not part of hypothesis~(2);
the English therefore closes the list before ``Then.'' The French authority
remains byte-identical; this disclosed repair follows the manager's final
adjudication.}}

\medskip
\noindent Note that here $T'$ is the set of prime ideals $\mathfrak p$ of
$A$ such that
\[
  \dim(A/\mathfrak p)=1.
\]

% END INCLUDED active_source/parts/09_expose_IX.tex

% BEGIN INCLUDED active_source/parts/10_expose_X.tex
\makeatletter
\providecommand{\sgareaderlabel}[2]{%
  \begingroup
  \def\@currentlabel{#2}%
  \phantomsection\label{#1}%
  \endgroup}
\makeatother

\chapter{Application to the Fundamental Group}
\label{X}

\medskip
Throughout this expos\'e, $X$ will denote a locally noetherian prescheme,
$Y$ a closed subset of $X$, $U$ a variable open neighborhood of $Y$ in $X$,
and $\widehat X$ the formal completion of $X$ along $Y$
(\textup{EGA}~I, 10.8). For every prescheme $Z$, let $\Et(Z)$ denote the
category of \'etale coverings of $Z$, and let $\LLcat(Z)$\sgareaderlabel{pX.1}{1} denote the category
of coherent locally free modules on $Z$.

\section*{1. Comparison of $\Et(\widehat X)$ and $\Et(Y)$}
\sgareaderlabel{X.1}{1}

Let $\mathcal I$ be an ideal of definition of $Y$ in $X$. For every
$n\in\mathbb N$, set
\[
  Y_n=\left(Y,\left.(\mathcal O_X/\mathcal I^{\,n+1})\right|_Y\right).
\]
The $Y_n$ form a direct system of ordinary preschemes, or also of formal
preschemes when the structure sheaves are endowed with the discrete topology.
We know (EGA~I, 10.6.2) that $\widehat X$ is the direct limit, in the category
of formal preschemes, of the direct system $(Y_n)$. We also know (EGA~I,
10.13) that to give a formal $\widehat X$-prescheme $R$ of \emph{finite type}
is to give a direct system of ordinary $Y_n$-preschemes $R_n$ of
\emph{finite type}, such that
\[
  R_n\simeq R_{n+1}\times_{Y_{n+1}}Y_n.
\]
Moreover, for $R$ to be an étale covering of $\widehat X$, it is necessary
and sufficient that, for every $n$, $R_n$ be an étale covering of $Y_n$.

That said, it is easy to see that nilpotent elements make no difference to
étale coverings (SGA~1, 8.3); in other words, the base-change functor
\[
  \Et(Y_{n+1})\longrightarrow\Et(Y_n)
\]
is an equivalence of categories for every $n\in\mathbb N$. Hence:

\medskip
\sgareaderlabel{X.1.1}{1.1}\noindent\textbf{Proposition 1.1.}\quad
\emph{With the notation introduced above, the natural functor
\[
  \Et(\widehat X)\longrightarrow\Et(Y)
\]
is an equivalence of categories (cf.\ SGA~1, 8.4).}

\section*{2. Comparison of $\Et(Y)$ and $\Et(U)$, for varying $U$}
\sgareaderlabel{X.2}{2}

We shall introduce two conditions from which the comparison theorem announced
above will follow easily. Let $X$ be a locally noetherian prescheme, and let
$Y$ be a closed subset of $X$. We say that the pair $(X,Y)$ satisfies the
\emph{Lefschetz condition}, written $\Lefc(X,Y)$,\sgareaderlabel{pX.2}{2} if, for every open subset
$U$ of $X$ containing $Y$ and every locally free coherent sheaf $E$ on $U$,
the natural homomorphism
\[
  \Gamma(U,E)\longrightarrow\Gamma(\widehat X,\widehat E)
\]
is an isomorphism.

We say that the pair $(X,Y)$ satisfies the \emph{effective Lefschetz
condition}, written $\Leffc(X,Y)$, if $\Lefc(X,Y)$ holds and, in addition,
for every locally free coherent sheaf $\mathcal E$ on $\widehat X$, there
exist an open neighborhood $U$ of $Y$, a locally free coherent sheaf $E$ on
$U$, and an isomorphism
\[
  \widehat E\simeq\mathcal E.
\]

These conditions hold in two important examples:

\medskip
\sgareaderlabel{X.2.1}{2.1}\noindent\textbf{\textit{Example 2.1}}\footnote{\sgareaderlabel{noteX.2.1}{*}\textit{Editor's note.}
One can slightly improve (i): see Mme Raynaud (Raynaud~M., ``Théorèmes de
Lefschetz en cohomologie des faisceaux cohérents et en cohomologie étale.
Application au groupe fondamental,'' \emph{Ann. Sci. Éc. Norm. Sup. (4)}
\textbf{7} (1974), pp.~29--52, Corollaries I.1.4 and I.5). Condition (ii) can
be improved so as to dispense with the depth conditions along $Y$ (see
Theorem~3.3 of \emph{loc. cit.} for a precise statement). The proof of this
latter point is very technical; moreover, the article cited above gives only
indications of proof and refers to an earlier, detailed version published in
the \emph{Bulletin de la Société mathématique de France}.}.\ ---\ 
Let $A$ be a noetherian ring, and let $t\in\mathfrak r(A)$ be an
$A$-regular element belonging to the radical $\mathfrak r(A)$ of $A$.
Assume that $A$ is a quotient of a regular local ring and that $A$ is
complete for the $t$-adic topology (for example, that $A$ is complete for
the $\mathfrak r(A)$-adic topology). Set
\[
  X'=\operatorname{Spec}(A),\qquad Y'=V(t),
\]
and, further, set
\[
  x=\mathfrak r(A),\qquad X=X'\setminus\{x\},\qquad
  Y=Y'\setminus\{x\}.
\]
Thus $X$ is open in $X'$ and $Y=X\cap Y'$. Then:

\begin{enumerate}[label=(\roman*),leftmargin=2.5em]
\item If, for every prime ideal $\mathfrak p$ of $A$ such that
$\dim A/\mathfrak p=1$ (i.e. for every closed point of $X$), we have
$\depthop A_{\mathfrak p}\geq 2$, then $\Lefc(X,Y)$ holds;

\item if, moreover, for every prime ideal $\mathfrak p$ of $A$ such that
$t\in\mathfrak p$ and $\dim A/\mathfrak p=1$ (i.e. for every closed point
of $Y$), we have $\depthop A_{\mathfrak p}\geq 3$, then $\Leffc(X,Y)$
holds.
\end{enumerate}

Let us first prove that, for every open neighborhood $U$ of $Y$ in $X$, the
complement of $U$ in $X$ is a union of finitely many closed points (in $X$).
Observe that $U$ is open in $X$, hence in $X'$, and therefore
$Z'=X'\setminus U$ is closed. Let $I$ be an ideal defining $Z'$; it is
enough to prove that $A/I$ has dimension $1$. Since
\[
  Z'\cap Y'=\{x\},
\]
the ring $A/(I+(t))$ is Artinian, and the conclusion follows from the
``Hauptidealsatz.''

The \emph{first hypothesis is equivalent to}: ``for every prime ideal
$\mathfrak p$ of $A$ with $\mathfrak p\ne\mathfrak r(A)$, one has
\[
  \operatorname{depth} A_{\mathfrak p}
  \geq 3-\dim(A/\mathfrak p).
\]
Indeed, $A$ is a quotient of a regular ring, so we may apply \hyperref[VIII.2.3]{VIII, 2.3} to
the prescheme $X'$, the closed subset $\{x\}$, and the coherent sheaf
$\mathcal O_{X'}$, observing that
$c(\mathfrak p)=\dim(A/\mathfrak p)$ for
$\mathfrak p\in X=X'\setminus\{x\}$ (since $x$ is the closed point of
$X'$).

\medskip
\noindent\colorbox{noteblue}{\parbox{0.95\linewidth}{\footnotesize
\textbf{Source oddity SGA2-X-L3357-U-VS-X-SRCDEF-001.}
The French authority prints $\mathfrak p\in U=X'-\{x\}$, although $U$ was
just quantified as an arbitrary open neighborhood and line 3345 defines
$X=X'-\{x\}$. The English therefore uses $X$ here. The French file remains
unchanged; the same-edition PDF confirms only the printed $U$, not the
repair.}}

Let $U$ be an open neighborhood of $Y$ in $X$, and let $E$ be a locally
free $\mathcal O_U$-module. Set $Z=X\setminus U$, and let
$u\colon U\to X$ be the canonical immersion. We shall first prove that
$u_*(E)$ is a coherent $\mathcal O_X$-module, or, equivalently, that
$\mathcal H^i_Z(E')$ is coherent for $i=0,1$, where $E'$ is a coherent
extension of $E$ to $X$. For this, we apply Theorem \hyperref[VIII.2.1]{VIII, 2.1} to the
prescheme $X$, the closed subset $Z$, and the coherent sheaf $E'$. It is
enough to verify that, for every point $p\in U$ such that $c(p)=1$, one has
$\operatorname{depth}E'_p\geq1$, where we have set
\[
  c(p)=\operatorname{codim}
  \bigl(\overline{\{p\}}\cap Z,\overline{\{p\}}\bigr).
\]
Now suppose $p\in U$ and $c(p)=1$, and let $\mathfrak p$ denote the ideal
of $A$ corresponding to $p$. Then $\dim(A/\mathfrak p)=2$, since the
complement of $U$ is a union of finitely many closed points and $A$ is a
quotient of a regular ring. Moreover, $E$ is locally free, so for every
$p\in\operatorname{Supp}E$ one has
$\operatorname{depth}E_p=\operatorname{depth}\mathcal O_{U,p}$. Finally,
if $p\in U$ and $c(p)=1$, then
\[
  \operatorname{depth}E'_p
  =\operatorname{depth}E_p
  =\operatorname{depth}\mathcal O_{U,p}
  =\operatorname{depth}A_{\mathfrak p}
  \geq3-2=1.
\]

We must now prove that the natural homomorphism
\begin{equation}\label{eq:X.1}
  \Gamma(U,E)\longrightarrow\Gamma(\widehat X,\widehat E)
\end{equation}
is an isomorphism. Set $\overline{\!E}=u_*(E)$. The sheaf
$\overline{\!E}$ is coherent and has depth at least $2$ at every closed
point of $X$. It follows that $R^if_*(\overline{\!E})$ is coherent for
$i=0,1$, where $f\colon X\to X'$ denotes the canonical immersion of $X$
into $X'=\operatorname{Spec}(A)$ (Exposé VIII). Applying \hyperref[IX.1.5]{IX, 1.5}, we
conclude that
\begin{equation}\label{eq:X.2}
  \Gamma(U,\overline{\!E})
  \longrightarrow
  \Gamma(\widehat X,\widehat{\overline{\!E}})
\end{equation}
is an isomorphism, since $A$ is complete for the $t$-adic topology.

There is a commutative diagram
\[
\xymatrix@C=18mm@R=9mm{
  \Gamma(X,\overline{\!E})
    \ar[rr]^{\simeq}
    \ar[dr]_{\simeq}
  && \Gamma(U,E)\ar[dl] \\
  & \Gamma(\widehat X,\widehat E) &
}
\]
and the conclusion follows.

\medskip
Now let $\mathcal E$ be a locally free coherent sheaf on $\widehat X$. If
one has proved that $\mathcal E$ is algebraizable, that is, isomorphic to the
formal completion of a coherent $\mathcal O_X$-module $E$, it is easy to see
that $E$ is locally free in a neighborhood of $Y$, and hence to prove
$\Leffc(X,Y)$. Let $\widehat{X'}$ be the formal spectrum of $A$ for the
$t$-adic topology; it is identified with the formal completion of $X'$ along
$Y'$. Denote by $f$ the canonical immersion of $X$ into $X'$, by $f'$ the
canonical immersion of $Y$ into $Y'$, and by $\widehat f$ the morphism
obtained by passing to the completions. For $\mathcal E$ to be algebraizable,
it suffices that $\widehat f_*(\mathcal E)$ be a coherent
$\mathcal O_{\widehat{X'}}$-module, since $A$ is complete for the $t$-adic
topology.
\sourceoddity{SGA2-X-L3384-PUSHFORWARD-BASE-SRCDEF-001}{The French authority
prints $\mathcal O_{\widehat X}$ as the coefficient ring of the pushforward.
Because $\widehat f\colon\widehat X\to\widehat{X'}$, the target uses
$\mathcal O_{\widehat{X'}}$. Only this coefficient ring is corrected.}
Set $\mathcal I=t\mathcal O_{\widehat X}$; this is an ideal of definition of
$\widehat X$.

For every $n\geq0$, set
\[
  \mathcal E_n=\mathcal E/\mathcal I^{n+1}\mathcal E.
\]
At every closed point $y\in Y$, the depth of $\mathcal E_0$ is at least two;
indeed, $t$ is an $A$-regular element, and therefore
\[
  \depthop\mathcal O_{Y,y}
  =\depthop\mathcal O_{X,y}-1\geq2.
\]
\sourceoddity{SGA2-X-L3386-Y0-VS-Y-SRCDEF-001}{The French authority prints
$\mathcal O_{Y_0,y}$, but $Y_0$ is undefined in this argument. Since
$Y=V(t)\cap X$, the regular-element quotient gives $\mathcal O_{Y,y}$.}
It follows that $\widehat f_*(\mathcal E)$ is coherent (\hyperref[IX.2.3]{IX~2.3)}.
\hfill\textsc{q.e.d.}

\medskip
\sgareaderlabel{X.2.2}{2.2}\noindent\textbf{\textit{Example 2.2}}\quad
\textup{(This will make it possible to compare the fundamental group of a
projective variety with that of a hyperplane section).}

Let $K$ be a field, and let $X$ be a proper $K$-prescheme. Let $\mathcal L$
be an ample invertible $\mathcal O_X$-module. Let
$t\in\Gamma(X,\mathcal L)$ be an $\mathcal O_X$-regular element; this means
that, for every open subset $U$ and every isomorphism
\[
  u\colon \mathcal L|_U \xrightarrow{\ \sim\ } \mathcal O_U,
\]
the section $u(t)$ is a non-zero-divisor in $\mathcal O_U$ (a condition that
does not depend on $u$). Let $Y=V(t)$ be the subscheme of $X$ defined by the
equation $t=0$\footnote{\textit{Editor's note.} Condition (ii) is
superfluous; see the editor's note to \hyperref[X.2.1]{Example~2.1} (original running p.~90).}.
Then:

\begin{enumerate}[label=(\roman*),leftmargin=2.5em]
\item If, for every closed point $x$ of $X$, one has
$\depthop\mathcal O_{X,x}\geq2$, then $\Lefc(X,Y)$ holds;

\item if, moreover, for every closed point $y\in Y$, one has
$\depthop\mathcal O_{X,y}\geq3$, then $\Leffc(X,Y)$ holds.
\end{enumerate}

\medskip
This example will be treated in detail in Exposé~XII.

\medskip
Let $S$ be a prescheme. It is known (EGA~II, 6.1.2) that the functor which
assigns to every finite flat covering $r\colon R\to S$ the
$\mathcal O_S$-algebra $r_*(\mathcal O_R)$ induces an equivalence between
the category of finite flat coverings of $S$ and the category of coherent
locally free $\mathcal O_S$-algebras.
\sourceoddity{SGA2-X-L3402-OS-VS-OX-SRCDEF-001}{The French authority prints
$\mathcal O_X$ in both coefficient-ring positions. Since $S$ is arbitrary
and $r\colon R\to S$, the direct image and the equivalent algebra category
are over $\mathcal O_S$.}
Let $U$ be an open neighborhood of $Y$, and let $r\colon R\to U$ be a
finite flat covering of $U$. Let $\widehat R$ be the finite flat covering of
$\widehat X$ obtained from it by base change. Then
\[
  \widehat r_*(\mathcal O_{\widehat R})
  \simeq \widehat{r_*(\mathcal O_R)}.
\]

Assume now that $\Lefc(X,Y)$ holds. This implies that, for every $U$, the
inverse-image functor
\[
  \mathsf L(U)\longrightarrow\mathsf L(\widehat X)
\]
is fully faithful. Indeed, let $E$ and $F$ be two coherent locally free
$\mathcal O_U$-modules; $\cHom(E,F)$ is likewise coherent and locally free.
\sourceoddity{SGA2-X-L3408-DUPLICATE-HOM-SRCDEF-001}{The French authority
prints the preceding sentence twice consecutively. The target states the
identical assertion once.}
By hypothesis, the natural map
\[
  \Gamma\bigl(U,\cHom(E,F)\bigr)
  \longrightarrow
  \Gamma\bigl(\widehat X,\widehat{\cHom(E,F)}\bigr)
\]
is an isomorphism, which proves the assertion, since $\cHom$ commutes with
formal completion when all the modules involved are locally free. Formal
completion also commutes with tensor products. It follows that the functor
which assigns to every coherent locally free $\mathcal O_U$-algebra
$\mathcal A$ the $\mathcal O_{\widehat X}$-algebra $\widehat{\mathcal A}$
is fully faithful. More precisely, if $E$ is a coherent locally free
$\mathcal O_U$-module, commutative $\mathcal O_U$-algebra structures on $E$
correspond bijectively to commutative $\mathcal O_{\widehat X}$-algebra
structures on $\widehat E$.
\sourceoddity{SGA2-X-L3412-BIJECTION-MISSING-FIRST-SIDE-SRCDEF-001}{The
French authority introduces a bijective correspondence but prints only its
$\mathcal O_{\widehat X}$ side. The target restores the omitted
$\mathcal O_U$-algebra structures on $E$, as required by the preceding
fully faithful tensor-compatible completion functor.}

\medskip
\sgareaderlabel{X.2.3}{2.3}\noindent\textbf{\textit{Proposition 2.3.}}
Let $X$ be a locally noetherian prescheme, and let $Y$ be a closed subset of
$X$. Let $\widehat X$ be the formal completion of $X$ along $Y$. For every
open subset $U$ of $X$ containing $Y$, let $L_U$ (respectively $P_U$,
respectively $E_U$) denote the functor that assigns to every locally free
coherent $\mathcal O_U$-module (respectively every finite flat covering of
$U$, respectively every \'{e}tale covering of $U$) its pullback along
$\widehat X\to X$.

\begin{enumerate}[label=(\roman*),leftmargin=2.5em]
\item If $\Lefc(X,Y)$ holds, then, for every open neighborhood $U$ of $Y$,
the functors $L_U$, $P_U$, and $E_U$ are fully faithful.

\item If $\Leffc(X,Y)$ holds, then, for every locally free coherent
$\mathcal O_{\widehat X}$-module $\mathcal E$, there are an open
neighborhood $U$ of $Y$ and a locally free coherent $\mathcal O_U$-module
$E$ such that
\[
  L_U(E)\simeq\mathcal E.
\]
The analogous assertions hold for finite flat coverings and for \'{e}tale
coverings, with $P_U$ and $E_U$, respectively.
\end{enumerate}

\medskip
\noindent\textit{Proof.}
Part~(i) has already been proved.

Part~(ii) follows from (i) and the hypothesis, at least for $L_U$ and $P_U$.
Moreover, if $R$ is an \'{e}tale covering of $\widehat X$, there are an open
neighborhood $U$ of $Y$ in $X$ and a finite flat covering $R'$ of $U$ such
that
\[
  \widehat{R'}\simeq R.
\]
By restriction this gives a covering $R''$ of $Y$, which is \'{e}tale by
\hyperref[X.1.1]{Proposition~1.1}. Hence the restriction of $R'$ to some neighborhood $U'$ of
$Y$ is \'{e}tale.
\hfill$\square$

\medskip
\sgareaderlabel{X.2.4}{2.4}\noindent\textbf{\textit{Corollary 2.4.}}
Suppose that $\Lefc(X,Y)$ holds. A finite flat covering $R$ of an open
neighborhood $U$ of $Y$ is connected if and only if
$R\times_U\widehat X$ is connected. In particular, $Y$ is connected if and
only if $U$ is connected, and this is also equivalent to $X$ being
connected.

\medskip
\noindent\textit{Proof.}
Indeed, a locally ringed space $(X,\mathcal O_X)$ is connected if and only if
$\Gamma(X,\mathcal O_X)$ is not a direct product of two nonzero rings. But,
by $\Lefc(X,Y)$, one has
\[
  \Gamma\bigl(U,r_*\mathcal O_R\bigr)
  \simeq
  \Gamma\bigl(\widehat X,\widehat r_*\mathcal O_{\widehat R}\bigr).
\]

\medskip
\sgareaderlabel{X.2.5}{2.5}\noindent\textbf{\textit{Corollary 2.5.}}
Suppose that $\Lefc(X,Y)$ holds. Then, for every open neighborhood $U$ of
$Y$, the functor
\[
  \Etc(U)\longrightarrow\Etc(Y)
\]
is fully faithful. Suppose that $\Leffc(X,Y)$ holds. Then, for every
\'{e}tale covering $R$ of $Y$, there are an open neighborhood $U$ of $Y$
and a covering $R'$ of $U$ such that
$R'\times_U Y\simeq R$.

\medskip
\sgareaderlabel{X.2.6}{2.6}\noindent\textbf{\textit{Corollary 2.6\textsuperscript{(3)}.}}
Suppose that $\Lefc(X,Y)$ holds and that $Y$ is connected. Then every open
neighborhood $U$ of $Y$ is connected, and the natural homomorphism
$\pi_1(Y)\longrightarrow\pi_1(U)$ is surjective. Suppose, moreover, that
$\Leffc(X,Y)$ holds. Then the natural homomorphism
\[
  \pi_1(Y)\longrightarrow\varprojlim_{U}\,\pi_1(U)
\]
is an isomorphism. (N.B. A ``base point'' in $Y$ is understood to have been
chosen and is also taken as a base point in $X$ for the definition of the
fundamental groups.)

\medskip
\noindent All of this follows immediately from \hyperref[X.1.1]{Propositions~1.1} and~\hyperref[X.2.3]{2.3}.

\section*{Editor's note (3)}
Combined with \hyperref[X.3.3]{Proposition~3.3} and criteria \hyperref[XII.2.4]{XII.2.4} and \hyperref[XII.3.4]{XII.3.4}, this yields the
following relative Lefschetz theorem. Let $f\colon X\to S$ be a projective and
flat morphism of connected noetherian schemes, and let $D$ be an effective
relative Cartier divisor on $X$ that is relatively ample. If, for every
$s\in S$, the depth of $X_s$ at every closed point is $\geq 2$, then $D$ is
connected and, for every open subset $U$ of $X$ containing $D$, the map
\[
  i_U\colon \pi_1(D)\longrightarrow\pi_1(U)
\]
is surjective. If, in addition, the depth of $X_s$ at every closed point of
$D_s$ is $\geq 3$ and the local rings of $X$ at its closed points are
pure---for example, are complete intersections (cf. \hyperref[X.3.4]{X, 3.4)}---then $i_X$ is an
isomorphism.

Cf. J.-B. Bost, ``Lefschetz theorem for Arithmetic Surfaces,''
\emph{Ann. Sci. \'Ec. Norm. Sup. (4)} \textbf{32} (1999), pp.~241--312,
Theorems~1.1 and~2.1.

When $X$ is merely a smooth projective surface that is geometrically connected
over a field, $D$ is always connected and
\[
  \pi_1(D)\longrightarrow\pi_1(U)
\]
is surjective (where $U$ is an open subset containing $D$), provided only that
$D$ is nef and has self-intersection $>0$; cf. \emph{loc. cit.}, Theorem~2.3,
and also Theorem~2.4 for surfaces assumed only normal and complete.

In the case of an arithmetic surface (normal and quasi-projective) $X$ over a
ring of integers $\mathcal O_K$, Bost, improving results of Y.~Ihara
(Y.~Ihara, ``Horizontal divisors on arithmetic surfaces associated with
Bely\u{\i} uniformizations,'' in \emph{The Grothendieck theory of dessins
d'enfants (Luminy, 1993)}, London Math. Soc. Lect. Note Series, vol.~200,
Cambridge Univ. Press, Cambridge, 1994, pp.~245--254, or \emph{loc. cit.},
Corollary~7.2), proved that if a point
$P\in X(\mathcal O_K)$, which plays the role of the divisor $D$ in the
geometric situation, satisfies certain positivity conditions, then the map
\[
  \pi_1(X)\longrightarrow\pi_1(\operatorname{Spec}\mathcal O_K)
\]
induced by the projection is an isomorphism whose inverse is the map
\[
  \pi_1(\operatorname{Spec}\mathcal O_K)\longrightarrow\pi_1(X)
\]
induced by $P$; cf. \emph{loc. cit.}, Theorem~1.2.

\medskip
\noindent{\footnotesize\emph{Accepted source-normalization note
(SGA2-X-L3456-AU-DESSUS-NORMALIZATION-001@1).} The corrected French control
selects \textup{au-dessus} from
\texttt{\string\sisi\{au dessus\}\{au-dessus\}}; the target renders this as
``over.'' This is an explicit source normalization, not a source defect.}

\section*{3. Comparison of $\pi_1(X)$ and $\pi_1(U)$}
\sgareaderlabel{X.3}{3}

\medskip
\sgareaderlabel{X.3.1}{3.1}\noindent\textbf{Definition 3.1.}\quad
Let $X$ be a prescheme and $Z$ a closed subset of $X$. Set
$U=X\setminus Z$. We say that the pair $(X,Z)$ is \emph{pure} if, for every
open subset $V$ of $X$, the functor
\[
\begin{aligned}
  \Et(V)&\longrightarrow\Et(V\cap U),\\
  V'&\longmapsto V'\times_V(V\cap U)
\end{aligned}
\]
is an equivalence of categories\renewcommand{\thefootnote}{*}\footnote{For a
notion that is more satisfactory in certain respects, cf. the commentary
\hyperref[XIV.1.6]{XIV, 1.6(d)}.}.

\medskip
\noindent{\footnotesize\emph{Accepted source-normalization note
(SGA2-X-L3464-MAPSTO-NORMALIZATION-001@1).} The corrected French control
selects \texttt{\string\mto} from
\texttt{\string\sisi\{\string\rightsquigarrow\}\{\string\mto\}}; the target
renders the object map with $\longmapsto$. This is an explicit source
normalization, not a source defect.}

\medskip
\sgareaderlabel{X.3.2}{3.2}\noindent\textbf{Definition 3.2.}\quad
Let $A$ be a noetherian local ring. Set $X=\Spec A$. Let $\rr(A)$ be the
radical of $A$, and let $x=\rr(A)$ be the closed point of $X$. We say that
$A$ is \emph{pure} if the pair $(X,\{x\})$ is pure.

\medskip
\noindent We leave to the reader the task of not proving the following
proposition:

\medskip
\sgareaderlabel{X.3.3}{3.3}\noindent\textbf{Proposition 3.3.}\quad
Let $X$ be a locally noetherian prescheme and let $Z$ be a closed subset of
$X$. The pair $(X,Z)$ is pure if and only if, for every $z\in Z$, the local
ring $\Oo_{X,z}$ is pure\renewcommand{\thefootnote}{**}\footnote{Compare the
noncommutative case of \hyperref[XIV.1.8]{XIV, 1.8}, whose proof is essentially the same as that
of \hyperref[X.3.3]{Proposition~3.3}.}.

\medskip
\noindent{\footnotesize\color{notebluetext}\emph{Source-defect note
(SGA2-X-L3477-FERME-VS-FERMEE-SRCDEF-001).} French line~3477 reads
\emph{une partie ferm\'e de $X$}; the masculine adjective \emph{ferm\'e}
disagrees with the feminine noun \emph{partie}. The target renders the
intended sense as ``a closed subset of $X$.'' The admitted French control is
unchanged.}

\medskip
\noindent That said, the following theorem is the main result of this
section:

\medskip
\sgareaderlabel{X.3.4}{3.4}\noindent\textbf{Theorem 3.4 (Purity theorem)}%
\renewcommand{\thefootnote}{(4)}\footnote{\textsc{N.D.E.}: For the history of
the methods employed, see Grothendieck's letter of October~1, 1961 to Serre,
\emph{Correspondance Grothendieck-Serre}, edited by Pierre Colmez and
Jean-Pierre Serre, Documents Math\'ematiques, vol.~2, Soci\'et\'e
Math\'ematique de France, Paris, 2001.}\textbf{.}

\begin{enumerate}
  \renewcommand{\labelenumi}{(\roman{enumi})}
  \item A regular noetherian local ring of dimension $\geq 2$ is pure
  (Zariski--Nagata purity theorem).
  \item A noetherian local ring of dimension $\geq 3$ that is a complete
  intersection is pure.
\end{enumerate}

\medskip
\noindent Recall that a local ring $A$ is said to be a \emph{complete
intersection} if and only if there exist a \emph{regular} noetherian local
ring $B$ and a $B$-regular sequence $(t_1,\ldots,t_k)$ of elements of the
radical $\mathfrak r(B)$ of $B$ such that
\[
  A \simeq B/(t_1,\ldots,t_k).
\]

\medskip
\noindent In this connection, let us note that it would be less ambiguous to
say that $A$ is an \emph{absolute} complete intersection, in contrast to the
situation already encountered, in which $X$ is a locally noetherian prescheme
(not necessarily regular) and $Y$ is a closed subset of $X$ that is said to
be ``locally a set-theoretic complete intersection \emph{in $X$}.''

\medskip
\noindent Let us first prove a few lemmas.

\medskip
{\itshape
\sgareaderlabel{X.3.5}{3.5}\noindent\textbf{Lemma 3.5.} --- Let $X$ be a locally noetherian prescheme,
and let $U$ be an open subset of $X$. Set $Z=X\setminus U$. Let
$i\colon U\to X$ be the canonical immersion of $U$ into $X$. The following
conditions are equivalent:
\begin{enumerate}[label=\textup{(\roman*)},leftmargin=2.2em,itemsep=0.45em,topsep=0.45em]
\item For every open subset $V$ of $X$, if $V'=V\cap U$, the restriction
functor $F\longmapsto F|_{V'}$ from the category of locally free coherent
$\mathcal O_V$-modules to the category of locally free coherent
$\mathcal O_{V'}$-modules is fully faithful;
\item the natural homomorphism $\mathcal O_X\to i_*\mathcal O_U$ is an
isomorphism;
\item for every $z\in Z$, one has $\prof\mathcal O_{X,z}\geq 2$.
\end{enumerate}
}

We have already seen (\Ref{III}~\Ref{III.3.3}) the equivalence of (ii)
and (iii). Let us show that (ii) implies (i). Let $F$ and $G$ be two
locally free coherent $\Oo_V$-Modules. Then $\SheafHom(F,G)$ is also
locally free and coherent, so
$$
\SheafHom(F,G)\longrightarrow
i_\ast\bigl(\SheafHom(\sisi{F|{|V'}}{F_{|V'}},
\sisi{G|{|V'}}{G_{|V'}})\bigr)
$$
is an isomorphism, and therefore
$$
\Hom(F,G)\simeq
\Hom(\sisi{F|{|V'}}{F_{|V'}},\sisi{G|{|V'}}{G_{|V'}}).
$$
Conversely, take $F=G=\OX$ and apply (i) to every open subset $V$ of
$X$.

\medskip
\noindent Here is a useful ``descent lemma'':

\medskip
{\itshape
\sgareaderlabel{X.3.6}{3.6}\noindent\textbf{Lemma 3.6.} --- Let $X$ be a locally noetherian prescheme,
and let $Z$ be a closed subset of $X$. Set $U=X\setminus Z$. Suppose that
the homomorphism $\mathcal O_X\to i_*\mathcal O_U$ is an isomorphism. Let
$f\colon X_1\to X$ be a faithfully flat and quasi-compact morphism. Set
$Z_1=f^{-1}(Z)$. If the pair $(X_1,Z_1)$ is pure, then so is $(X,Z)$.
}

\medskip
\noindent\textit{Proof of \hyperref[X.3.6]{Lemma 3.6}.}\quad
Observe that the hypothesis
$\Oo_X\simeq i_*\Oo_U$ is preserved under flat base extension, since $i$ is
quasi-compact and, in this case, formation of the direct image commutes with
inverse image. This hypothesis implies that the functor
\[
  \Et(V)\longrightarrow\Et(U\cap V)
\]
defined by
\[
  V'\longmapsto V'\times_V(V\cap U)
\]
is fully faithful, as follows from the interpretation of an étale covering
in terms of a coherent locally free algebra. It remains to prove effectivity.
One may, for example, form the fiber square $X_2$ and the fiber cube $X_3$ of
$X_1$ over $X$, and observe that a faithfully flat and quasi-compact
morphism is a morphism of universal effective descent for the fibered
category of étale coverings over the category of preschemes. The conclusion
then follows formally\renewcommand{\thefootnote}{(*)}\footnote{Cf. J.~Giraud,
\emph{Méthode de la descente}, Mémoire no.~2 du \emph{Bulletin de la Société
Mathématique de France} (1964).}.

\medskip
{\itshape
\sgareaderlabel{X.3.7}{3.7}\noindent\textbf{Remark 3.7.} --- We proved along the way that if
$\mathcal O_X\to i_*\mathcal O_U$ is an isomorphism, then $X$ is connected if
and only if $U$ is connected, and then $\pi_1(U)\to\pi_1(X)$ is surjective.
}

\medskip
{\itshape
\sgareaderlabel{X.3.8}{3.8}\noindent\textbf{Corollary 3.8.} --- Let $A$ be a noetherian local ring.
Assume that $\depthop A\geq 2$. If $\widehat{A}$ is pure, then $A$ is pure.
}

\medskip
\noindent\emph{Proof.}\quad This follows from \hyperref[X.3.5]{Lemma~3.5} and \hyperref[X.3.6]{Lemma~3.6}.

\medskip
\noindent The following lemma is the essential point in the proof of the purity theorem:

\medskip
{\itshape
\sgareaderlabel{X.3.9}{3.9}\noindent\textbf{Lemma 3.9.} --- Let $A$ be a noetherian local ring, and
let $t\in\mathfrak r(A)$ be an $A$-regular element. Suppose that $A$ is
complete for the $t$-adic topology and is moreover a quotient of a regular
local ring (for example, if $A$ is complete). Set $B=A/tA$.
\begin{enumerate}[label=\textup{(\roman*)},leftmargin=2.2em,itemsep=0.6em,topsep=0.5em]
\item If, for every prime ideal $\mathfrak p$ of $A$ such that
$\dim A/\mathfrak p=1$, one has $\depthop A_{\mathfrak p}\geq 2$, then
the purity of $B$ implies the purity of $A$.

\item If, for every prime ideal $\mathfrak p$ of $A$ such that
$\dim A/\mathfrak p=1$, one has $\depthop A_{\mathfrak p}\geq 2$, if
$A_{\mathfrak p}$ is pure whenever $t\notin\mathfrak p$, and
if\renewcommand{\thefootnote}{(5)}\footnote{\textsc{N.D.E.}: This last
condition can be improved; cf. the editor's note~(1) on p.~90.}
$\depthop A_{\mathfrak p}\geq 3$ whenever $t\in\mathfrak p$, then the
purity of $A$ implies the purity of $B$.
\sourceoddity{SGA2-X-L3551-MISSING-EST-SRCDEF-001}{The French authority
prints \emph{si $A_{\mathfrak p}$ pur lorsque $t\notin\mathfrak p$},
omitting the finite copula \emph{est}. The English supplies ``is'' in the
unambiguous intended condition. The French authority remains unchanged.}
\end{enumerate}
}

\medskip
Let $X'=\operatorname{Spec}(A)$ and let $Y'=V(t)$, which we identify with
the spectrum of $B$. Let $x=\mathfrak r(A)$, and set
$X=X'\setminus\{x\}$ and $Y=Y'\setminus\{x\}=X\cap Y'$. Denote by
$\widehat{X'}$ the formal spectrum of $A$ for the $t$-adic topology; it is
identified with the formal completion of $X'$ along $Y'$.

Since $A$ is complete for the $t$-adic topology, observe that
$\Et(X')\longrightarrow\Et(\widehat{X'})$ is an equivalence of categories.
Likewise, $\Et(\widehat{X'})\longrightarrow\Et(Y')$ is an equivalence by
\hyperref[X.1.1]{Proposition~1.1}; hence $\Et(X')\longrightarrow\Et(Y')$ is an equivalence of
categories.

We prove \textup{(i)}. Consider the diagrams
\[
\xymatrix@C=12mm@R=8mm{
  X' & X\ar[l] &&
  \Et(X')\ar[r]^{a}\ar[d]_{c} & \Et(X)\ar[d]^{b}\\
  Y'\ar[u] & Y\ar[l]\ar[u] &&
  \Et(Y')\ar[r]^{d} & \Et(Y)
}
\]
We have just seen that $c$ is an equivalence; $d$ is also an equivalence by
the hypothesis that $B$ is pure; and finally $b$ is fully faithful, as was
seen in \hyperref[X.2.1]{Example~2.1}; cf. \hyperref[X.2.3]{Proposition~2.3}\textup{(i)}.

We prove \textup{(ii)}. This time suppose that $A$ is pure, so $a$ is an
equivalence; likewise $c$. Let us show that $b$ is an equivalence. By
\hyperref[X.2.1]{Example~2.1}, we know that $\Leffc(X,Y)$ holds, so $b$ is already fully
faithful; let us prove that it is essentially surjective. Apply
\hyperref[X.2.3]{Proposition~2.3}\textup{(ii)}, noting that if $U$ is an open neighborhood of
$Y$ in $X$, then $X\setminus U$ is a union of finitely many closed points.
The pair $(X,X\setminus U)$ is therefore pure by \hyperref[X.3.3]{Proposition~3.3}, since at
any such point $\mathfrak p$ one has
$\mathcal O_{X,\mathfrak p}=A_{\mathfrak p}$, which is pure by hypothesis.
This proves the assertion.

\medskip
\textit{Proof of the purity theorem.}
We first prove \textup{(i)} by induction on the dimension. Let $A$ be a
noetherian local ring of dimension~$2$. Set
$X'=\Spec(A)$, $x=\mathfrak r(A)$, and $X=X'\setminus\{x\}$. We have
$\depthop A=2$. We may therefore apply \hyperref[X.3.5]{Lemma~3.5} to the pair
$(X',\{x\})$, and hence
$\Et(X')\longrightarrow\Et(X)$ is fully faithful.

Now let $r\colon R\longrightarrow X$ be an \'{e}tale covering defined by
the coherent, locally free, \'{e}tale $\mathcal O_X$-algebra
$\mathcal A=r_*(\mathcal O_R)$. Let $i\colon X\longrightarrow X'$ denote
the canonical immersion. I claim that
$i_*(\mathcal A)=\mathcal B$ is a coherent $\mathcal O_{X'}$-algebra.
Indeed, it suffices to apply the ``finiteness theorem,'' \hyperref[VIII.2.3]{VIII, 2.3}. I also
claim that this algebra has depth at least~$2$ at $x$: it is the direct image
of an $\mathcal O_X$-module, with $X=X'\setminus\{x\}$. Since $A$ is a
regular ring of dimension~$2$, we have
\[
  \pdop \mathcal B+\depthop \mathcal B=\dim A=2,
\]
where $\pdop \mathcal B$ denotes the projective dimension of $\mathcal B$.
Thus $\pdop \mathcal B=0$, so $\mathcal B$ is projective and therefore
free. Consequently, $\mathcal B$ defines a finite flat covering of
$X'=\Spec(A)$. The set of points of $X'$ at which this covering is not
\'{e}tale is a closed subset of $X'$ whose defining ideal is principal,
namely the discriminant ideal of $\mathcal B/A$. By construction this closed
subset is contained in $x=\mathfrak r(A)$; it is therefore empty because
$\dim A=2$.

Let $A$ now be a regular noetherian local ring with $\dim A=n\geq 3$.
Suppose that \textup{(i)} has been proved for rings of dimension less than
$n$. To prove that $A$ is pure, we may assume by \hyperref[X.3.8]{Corollary~3.8} that $A$ is
complete. Choose $t\in\mathfrak r(A)$ whose image in
$\mathfrak r(A)/\mathfrak r(A)^2$ is nonzero. Then $B=A/tA$ is a regular
noetherian local ring of dimension $n-1$, and is therefore pure because
$n-1\geq 2$. The conclusion follows from \hyperref[X.3.9]{Lemma~3.9}\textup{(i)}, which is
applicable because $A$ is complete.

\medskip
\textit{Let us prove \textup{(ii)}.}
Let $A$ be a noetherian local ring of dimension at least~$3$. Suppose that
there exist a regular noetherian local ring $B$ and a $B$-sequence
$(t_1,\ldots,t_k)$ such that
\[
  A\simeq B/(t_1,\ldots,t_k).
\]
We prove that $A$ is pure by induction on $k$. If $k=0$, this follows from
\textup{(i)}. Suppose that $k\geq1$ and that the result has been proved for
$k'<k$. By \hyperref[X.3.8]{Corollary~3.8}, we may assume that $A$, and hence also $B$, is
complete. Set
\[
  C=B/(t_1,\ldots,t_{k-1}).
\]
Then $A\simeq C/t_kC$, and $t_k$ is $C$-regular. By the induction hypothesis,
$C$ is pure; it therefore suffices to verify that \hyperref[X.3.9]{Lemma~3.9}\textup{(ii)} is
applicable. In the notation of that lemma, its $A$ and $B$ are here $C$ and
$A$, respectively. We have $\dim C\geq4$; hence, for every prime ideal
$\mathfrak p$ of $C$ such that $\dim(C/\mathfrak p)=1$, one has
$\depthop C_{\mathfrak p}\geq3$. Moreover, $C_{\mathfrak p}$ is a complete
intersection with $k'\leq k-1$, and is therefore pure by the induction
hypothesis.

\medskip
{\itshape
\sgareaderlabel{X.3.10}{3.10}\noindent\textbf{Theorem 3.10.} --- Let $X$ be a locally noetherian
prescheme, and let $Y$ be a closed subset of $X$. Suppose that
$\Leff(X,Y)$ holds (cf. \hyperref[X.2.1]{Examples~2.1} and~\hyperref[X.2.2]{2.2}). Suppose further that, for
every open neighborhood $U$ of $Y$ and every $x\in X\setminus U$, the
local ring $\mathcal O_{X,x}$ is either regular of dimension $\geq 2$ or
a complete intersection of dimension $\geq 3$. Then
\[
  \pi_0(Y)\longrightarrow\pi_0(X)
\]
is a bijection, and, if $X$ is connected,
\[
  \pi_1(Y)\longrightarrow\pi_1(X)
\]
is an isomorphism.
}

\medskip

There is nothing more to prove. Observe that in the two cited
\hyperref[X.2.1]{Examples~2.1} and~\hyperref[X.2.2]{2.2}, the complement of $U$ is a union of finitely many
closed points; it follows that the hypothesis on the dimension of
$\mathcal O_{X,x}$ is not farcical.

% END INCLUDED active_source/parts/10_expose_X.tex

% BEGIN INCLUDED active_source/parts/11_expose_XI.tex
\chapter{Application to the Picard Group}\label{XI}
\medskip
This expos\'e follows the pattern of the preceding one, but this time the result of no.~1 is weaker.

\medskip
Throughout this expos\'e, $X$ will denote a locally noetherian prescheme;
$\Ii$, a quasi-coherent ideal of $\OX$ (so $Y=V(\Ii)$ is a closed subset of
$X$); $U$, a variable open neighborhood of $Y$ in $X$; and $\widehat X$,
the formal completion of $X$ along $Y$.

For every ringed space $(Z,\mathcal O_Z)$, let $\PP(Z)$\refstepcounter{toto}\label{pXI.1}
denote the category of invertible $\mathcal O_Z$-modules, that is, modules locally
free of rank~$1$, and let $\Pic(Z)$ denote the group of isomorphism classes of
invertible modules on $Z$.

\section{\texorpdfstring{Comparison of $\Pic(\widehat X)$ and $\Pic(Y)$}{Comparison of Pic(X-hat) and Pic(Y)}}\label{XI.1}

For every $n\in\NN$, set
\[
X_n=(Y,\OX/\Ii^{n+1})
\qquad\text{and}\qquad
P_n=\Ii^{n+1}/\Ii^{n+2}.
\]
The sequence of sheaves of abelian groups on $Y$
\begin{equation}\label{eq:XI.1.1}
0\longrightarrow P_n\xrightarrow{u}\mathcal O_{X_{n+1}}^*
\xrightarrow{v}\mathcal O_{X_n}^*\longrightarrow 1
\end{equation}
is \emph{exact}. More precisely, the group structure on $P_n$ is the additive
structure, $u(x)=1+x$ for every $x\in P_n$, and $v$ is the homomorphism induced
by the injection $\Ii^{n+2}\longrightarrow\Ii^{n+1}$. To see that $v$ is
surjective, observe that, for every $y\in Y$, $\mathcal O_{X_n,y}$ is a local
ring, the quotient of $\mathcal O_{X_{n+1},y}$ by a nilpotent ideal; the rest is
equally trivial. From~\eqref{eq:XI.1.1} we deduce an exact cohomology sequence:
\begin{equation*}\tag{$*$}\label{eq:XI.1.*}
\mathrm H^1(Y,P_n)\xrightarrow{u^1}\mathrm H^1(Y,\mathcal O_{X_{n+1}}^*)
\xrightarrow{v^1}\mathrm H^1(Y,\mathcal O_{X_n}^*)
\xrightarrow{d}\mathrm H^2(Y,P_n).
\end{equation*}

Moreover, for every $n\in\NN$, we know how to identify
$\Pic(X_n)$ with $\mathrm H^1(Y,\mathcal O_{X_n}^*)$. Furthermore, if $E$
is an invertible $\mathcal O_{X_{n+1}}$-module corresponding to a cohomology
class $c(E)$\pageoriginale, the cohomology class corresponding to the pullback
of $E$ to $X_n$ is equal to $v^1(c(E))$. Hence the following proposition:

\begin{proposition}\label{XI.1.1}
Retain the notation introduced above. Let $p\in\NN$. The map
\[
  \Pic(\widehat X)\longrightarrow\Pic(X_n)
\]
has the following properties:

\sourceoddity{SGA2-XI-L3609-PIC-YN-VS-XN-SRCDEF-001}{The French prints
$Y_n$. Expos\'e~XI locally defines and uses $X_n$, while Expos\'e~X used
$Y_n$ for the same infinitesimal thickening. We therefore read $X_n$ here.}

\begin{enumerate}
\renewcommand{\labelenumi}{(\roman{enumi})}
\item it is injective for $n\geq p$ if $\mathrm H^1(Y,P_n)=0$ for every
$n\geq p$;
\item it is an isomorphism for $n\geq p$ if $\mathrm H^i(Y,P_n)=0$ for
every $n\geq p$ and $i=1,2$.
\end{enumerate}
\end{proposition}

Of course, the exact sequence~\textup{($\ast$)} contains more
information than the proposition above. The reader will have noticed that
nothing has been said about the functor
$\PP(\widehat X)\longrightarrow\PP(Y)$. Given two invertible
$\mathcal O_{\widehat X}$-modules $E,F$,
$H=\SheafHom(E,F)$ is also invertible. If reduction modulo
$\Ii^{n+1}$ is denoted by a subscript $n$, we obtain an exact sequence:
\[
0\longrightarrow H_0\otimes P_n
\longrightarrow \SheafHom(E_{n+1},F_{n+1})
\longrightarrow \SheafHom(E_n,F_n)
\longrightarrow 0.
\]
Hence there is an exact cohomology sequence that we shall not write down and
whose interpretation is evident; this observation can be used to study the
functor~$\PP$.

\section{\texorpdfstring{Comparison of $\Pic(X)$ and $\Pic(\widehat X)$}{Comparison of Pic(X) and Pic(X-hat)}}\label{XI.2}

The reader will find in Expos\'e~X, no.~2, the proof of the following:






\begin{proposition} \label{XI.2.1} Suppose that $\Lef(X, Y)$ holds; then, for every open neighborhood $U$ of $Y$ in $X$, the functor
\begin{equation} \label{eq:XI.2.1} \PP(U)\to\PP(\hat X)
\end{equation}
is fully faithful, and the map
\begin{equation} \label{eq:XI.2.2} \Pic(U)\to\Pic(\hat X)
\end{equation}
is\pageoriginale therefore injective. If $\Leff(X, Y)$ holds, the map~(\ref{eq:XI.2.3}) is an isomorphism:
\begin{equation} \label{eq:XI.2.3}
\varinjlim_U\Pic(U)\to\Pic(\hat X){\quoi}.
\end{equation}
\end{proposition}

\begin{corollaire} \label{XI.2.2} Suppose that $\Lef(X, Y)$ holds and that, for every integer $n\geq p$, one has $\H^1(Y, P_n)=0$; then, for every open subset $U\supset Y$, the maps
\[
\Pic(X)\to\Pic(U)\to\Pic(X_n)
\]
are injective if $n\geq p$. If $\Leff(X, Y)$ holds and if, moreover, for every integer $n\geq p$, one has $\H^i(Y, P_n)=0$ for $i=1$ and $i=2$, then the map
\[
\varinjlim_U\Pic(U)\to\Pic(X_n)
\]
is an isomorphism for $n\geq p$.
\end{corollaire}

\section{Comparison of $\PP(X)$ and $\PP(U)$} \label{XI.3}

A definition:
\refstepcounter{subsection}\label{XI.3.1}
\begin{enonce*}[remark]{Definition 3.1\sfootnotemark}\sfootnotetext{For a more detailed study of the notion of parafactoriality, and for the proof of~\Ref{XI.3.3}, \cf \EGA IV 21.13, 21.14.}
Let $X$ be a prescheme and let $Z$ be a closed subset of~$X$. Put $U=X-Z$. We say that $X$ is \emph{parafactorial} at the points of $Z$ if, for every open subset $V$ of~$X$, the functor $\PP(V)\to \PP(V\cap U)$ is an equivalence of categories. We also say that the pair $(X, Z)$ is parafactorial.
\end{enonce*}

Recall that $\PP(Z)$ denotes the category of locally free modules of rank~$1$ on~$Z$.

\begin{definition} \label{XI.3.2}
A noetherian local ring is said to be \emph{parafactorial} if the pair $(\Spec(A), \{\rr(A)\})$ is parafactorial.
\end{definition}

One proves the following proposition, which shows that the notion is \og pointwise\fg:

\begin{proposition} \label{XI.3.3}
Suppose\pageoriginale that $X$ is locally noetherian. The pair $(X, Z)$ is parafactorial if and only if, for every $z\in Z$, the local ring $\Oo_{X, z}$ is parafactorial.
\end{proposition}

Note that \og parafactorial\fg includes \og fully faithful\fg. As in Lemma~\Ref{X.3.5} of Expos\'e~\Ref{X}, one proves the following:

\begin{lemme} \label{XI.3.4}
If $X$ is a locally noetherian prescheme and if $Z=X-U$ is a closed subset of $X$, the following conditions are equivalent:
\begin{enumeratei}
\item
for every open subset $V$ of $X$, the functor $\PP(V)\to \PP(V\cap U)$ is fully faithful;
\item
the homomorphism $\OX \to i_*(\Oo_U)$ is an isomorphism;
\item
for every $z\in Z$, one has $\prof(\Oo_{X, z})\geq 2$.
\end{enumeratei}
\end{lemme}

Thus \og parafactorial\fg means that the conditions of~\Ref{XI.3.4} are satisfied and that, for every open subset $V$ of $X$, the homomorphism $\Pic(V)\to\Pic(V\cap U)$ is surjective. In particular, if $X$ is the spectrum of a noetherian local ring, one obtains:

\begin{proposition} \label{XI.3.5}
Let $A$ be a noetherian local ring; it is parafactorial if and only if $\prof A\geq 2$ and $\Pic(X'-\{x\})=0$, where $X'=\Spec(A)$ and $x$ is the unique closed point of $X'$.
\end{proposition}

Notice that a local ring of dimension $\leq 1$ is never parafactorial, since its depth is $\leq 1$. Thus \emph{factorial does not imply parafactorial}; the implication does, however, hold for noetherian local rings of dimension $\geq 2$, as we shall see below.

\begin{lemme} \label{XI.3.6} Let $X$ be a locally noetherian prescheme and let $Z$ be a closed subset of $X$. Let $f\colon X_1\to X$ be a faithfully flat and quasi-compact morphism. Put $Z_1=f^{-1}(Z)$. If $(X_1, Z_1)$ is parafactorial, then $(X, Z)$ is parafactorial.
\end{lemme}

First note\pageoriginale that, if $i\colon(X-Z)\to X$ denotes the canonical immersion of $U=X-Z$ into $X$, the formation of the direct image under $i$ of a quasi-coherent $\Oo_U$-module commutes with base change by $f$, since $f$ is flat. It is therefore equivalent to assume the equivalent conditions of Lemma~\Ref{XI.3.4} for $(X, Z)$ or for $(X_1, Z_1)$, since $f$ is a morphism of descent for the category of quasi-coherent sheaves. It remains to prove that, for every open subset $V$ of $X$, $\Pic(V)\to\Pic(V\cap U)$ is surjective. We make the base change $V\to X$, which changes nothing (\sisi{sic}{\textit{sic}}), and we are reduced to $V=X$. We then note that, if ${\Lcal}$ is an invertible $\Oo_U$-module and if ${\Lcal}$ admits a locally free extension, this extension is isomorphic to $i_*({\Lcal})$, by what we have just seen. It remains to prove that $i_*({\Lcal})$ is invertible. Using once again the fact that direct image under $i$ commutes with flat base change, and that \og locally free of rank~$1$\fg is a property that descends under faithfully flat and quasi-compact morphisms, we are done.

\begin{corollaire} \label{XI.3.7} Let $A$ be a noetherian local ring; if $\hat A$ is parafactorial, then $A$ is parafactorial.
\end{corollaire}

One should not believe that if $A$ is parafactorial, then $\hat A$ is parafactorial\nde{For a precise study of the relationship between the factoriality of $A$ and that of its completion, see (Heitmann~R., {\og Characterization of completions of unique factorization domains\fg}, \emph{Trans. Amer. Math. Soc.} \textbf{337} (1993), \numero 1, pp\ptbl 379--387).}.

\danger

Before turning to the statement of the main theorem of this \no, let us make the connection with divisor theory and the notion of a factorial ring\sfootnote{For the generalities that follow, \cf also \EGA IV~21.}.

Let $X$ be a \emph{noetherian} and \emph{normal} prescheme. Let $Z^1(X)$ be the free abelian group generated by the $x\in X$ such that $\dim\OXx=1$. The local ring at such a point is a discrete valuation ring. We denote the corresponding normalized valuation by $v_x$. Let $K(X)$ be the ring of rational functions on $X$, and let
\[
p\colon K(X)^*\to Z^1(X)
\]
be the map that associates to every $f\in K(X)^*$ the codimension-one cycle:
\[
(f) = \sum_{x\in X, \;\dim\OXx=1} v_x(f)\cdot x{\quoi}.
\]
The image\pageoriginale of $p$ is denoted by $P(X)$, and its elements are called \emph{principal divisors}\sfootnote{In accordance with the terminology of \EGA IV~21, we now prefer to reserve the term \og divisors\fg for \og locally principal divisors\fg or \og Cartier divisors\fg.}. We set
\[
\Cl(X) = Z^1(X)/P(X){\quoi}.
\]
Let $\sisi{Z}{Z^{\prime1}(X)}$ be the subgroup of $Z^1(X)$ whose elements are the \emph{locally principal divisors}. We know that
\[
\Pic(X)\simeq Z^{\prime1}(X)/P(X){\quoi},
\]
and consequently $\Pic(X)$ is identified with a subgroup of $\Cl(X)$.

Note that if $U$ is a dense open subset of $X$, then $K(X)\to K(U)$ is an isomorphism, and that if $\codim(X-U, X)\geq 2$, \ie if every $x\in X$ such that $\dim\OXx\leq 1$ belongs to $U$, the homomorphism $Z^1(X)\to Z^1(U)$, and consequently $\Cl(X)\to\Cl(U)$, is also an isomorphism. Finally, if every $x\in U$ is factorial, \ie if $\OXx$ is factorial, then $Z^1(U) = Z^{\prime1}(U)$, and hence $\Pic(U)\simeq\Cl(U)$.

\begin{propositionblah} \label{XI.3.7.1}
Let $X$ be a noetherian and normal prescheme. Let $(U_i)_{i\in I}$ be a family of open subsets of $X$ such that:
\begin{enumeratea}
\item
the $U_i$ form a filter base\nde{\ie a decreasing filtered family.};
\item
if we set $Y_i=X-U_i$, then $\codim(Y_i, X)\geq 2$ for every $i$,
\item
if $x\in U_i$ for every $i\in I$, then $\OXx$ is factorial.
\end{enumeratea}
\noindent
Then there is an isomorphism:
\[
\varinjlim_{i\in I}\Pic(U_i)\lto{\approx}\Cl(X){\quoi}.
\]
\end{propositionblah} Note that b) implies that every $x\in X$ such that $\dim\OXx\leq 1$ belongs to $U_i$ for every $i$. Thus the $U_i$ are dense, and moreover the homomorphism $Z^1(U_i)\to Z^1(X)$ is an isomorphism, as is $K(U_i)\to K(X)$. Thus $\Pic(U_i)\subset \Cl(U_i) \simeq \Cl(X)$. To prove the desired assertion, it is therefore enough to show that every $D\in Z^1(X)$ belongs to $Z^{\prime1}(U_i)$ for a suitable $i$. It is enough to do this for irreducible positive \og divisors\fg. Thus let $x\in X$ be such that\pageoriginaled $\dim\OXx=1$. It is enough to prove that there exists an $i\in I$ such that $\overline{\{x\}}$ is locally principal at the points of $U_i$. Let $\Ii$ be the largest ideal of definition of the closed subset $\overline{\{x\}}$. The set of points in a neighborhood of which $\Ii$ is free is an open subset $U$. Now $U\supset \bigcap_{i\in I}U_i$ by~c). If we set $Y=X-U$, then $Y\subset \bigcup_{i\in I}Y_i$\sisi{}{ with $Y_i=X-U_i$}; now $Y$ is closed, and therefore has only \emph{finitely} many generic points, so it is contained in the union of finitely many $Y_i$, and hence in $Y_j$ for some $j\in I$, since the $U_i$ form a filter base. Thus $U\supset U_j$.\qed

\begin{corollaire} \label{XI.3.8}
Let $X$ be a noetherian and normal prescheme, and let $Y$ be a closed subset of codimension $\geq 2$. Suppose that, for every $p\in X-Y$, $\Oo_{X, p}$ is factorial. Then
\[
\Pic(X-Y)\to\Cl(X-Y)\to\Cl(X)
\]
are isomorphisms.
\end{corollaire}

\begin{corollaire} \label{XI.3.9}
Let $A$ be a noetherian normal local ring. Set $X'=\Spec(A)$ and $x = \rr(A)$. For $A$ to be \emph{factorial}, it is necessary and sufficient that $\Pic(X'-\{x\})=0$ and that $\pp\in X'-\{x\}$ imply that $A_\pp$ is factorial.
\end{corollaire}
Indeed, for $A$ to be factorial, it is necessary and sufficient that $\Cl(X')=0$\,\nde{see Bourbaki, \emph{Alg\`ebre commutative} VII.1.4, corollary to Theorem~2, and VII.3.2, Theorem~1}.

\begin{corollaire} \label{XI.3.10}
Let $A$ be a noetherian local ring of dimension $\geq 2$. Set $X'=\Spec(A)$, and let $x = \rr(A)$. Set $X=X'-\{x\}$. The following conditions are equivalent:
\begin{enumeratei}
\item
$A$ is factorial;
\item
\textup{a)}\enspace for every $y\in X$, $\Oo_{X, y}$ is factorial, and
\par\hspace*{6.3mm}\textup{b)}\enspace $A$ is parafactorial, \ie $\prof A\geq 2$ and $\Pic(X)=0$.
\end{enumeratei}
\end{corollaire}

Before\pageoriginale proving this corollary, let us state the
\refstepcounter{subsection}\label{XI.3.11}
\begin{enonce*}{Serre's normality criterion 3.11\sfootnotemark}\sfootnotetext{\Cf \EGA IV~5.8.6.} Let $A$ be a noetherian local ring. For $A$ to be normal, it is necessary and sufficient that
\begin{enumeratei}
\item
for every prime ideal ${\pp}$ of $A$ such that $\dim A_{{\pp}}\leq 1$, $A_{{\pp}}$ is normal,
\item
for every prime ideal ${\pp}$ of $A$ such that $\dim A_{{\pp}}\geq 2$, one has $\prof A_{{\pp}}\geq 2$.
\end{enumeratei}
\end{enonce*}

Let us prove~\Ref{XI.3.10}.

(i)\ALORS (ii). Since a localization of a factorial ring is again factorial, we have (ii)~a). Moreover, $A$ is normal, and hence $\prof A\geq 2$, since $\dim A\geq 2$ (\Ref{XI.3.11}~(ii)). Finally, $A$ is parafactorial; indeed, $\Pic(X)\simeq\Cl(X')=0$ (\cf \Ref{XI.3.9}).

(ii)\ALORS (i). Let us first prove that $A$ is \emph{normal} by applying \sisi{SERRE}{Serre}'s criterion. Since $\dim A\geq 2$, condition~(i) of the criterion is among the hypotheses. Moreover, for every ${\pp}\in X$, $A_{{\pp}}$ is factorial, hence normal, and therefore of depth $\geq 2$, at least if $\dim A_{{\pp}}\geq 2$. Finally, $\prof A\geq 2$ by (ii)~b). It remains to apply~\Ref{XI.3.9}.

\smallskip
Let us summarize the foregoing:

\begin{proposition} \label{XI.3.12}
Let $X$ be a locally noetherian prescheme, and let $\Ii$ be a quasi-coherent ideal of $X$. Let $Y=V(\Ii)$. Let $p\in\NN$. Suppose that:
\begin{enumerate}
\item[1)] one \sisi{has}{has} $\Leff(X, Y)$ (\textup{Expos\'e} \Ref{X});
\item[2)] $\H^i(X, \Ii^{n+1}/\Ii^{n+2})=0$ if $i=1$ or $2$ and if $n\geq p$;
\item[3)] for every open neighborhood $U$ of $Y$ in $X$ and every $x\in X-U$, the ring $\OXx$ is parafactorial.
\end{enumerate}
\noindent
Then, for every $n\geq p$, and every open neighborhood $U$ of $Y$, the homomorphisms
\[
\xymatrix{\Pic(X)\ar[rr]\ar[dr]&&\Pic(X_n)\\&\Pic(U)\ar[ur]&}
\]
are isomorphisms.
\end{proposition}

We\pageoriginale know some parafactorial rings:

\skpt
\begin{theoreme} \label{XI.3.13}
\begin{enumeratei}
\item
(Auslander--Buchsbaum)\nde{compare with the following purity result, due to Gabber. Let $X$ be the spectrum of a regular local ring $A$ of dimension $3$, let $a$ have nonzero differential, \ie $a\in\mm-\mm^2$, and let $U$ be the complement of $V(a)$. Then a vector bundle on $U$ is free (for a simple proof, see (Swan~R.G., {\og A simple proof of Gabber's theorem on projective modules over a localized local ring\fg} \emph{Proc. Amer. Math. Soc.} \textbf{103} (1988), \numero 4, p\ptbl 1025--1030). The rank-$1$ case is a special case of Theorem~\Ref{XI.3.13}. For purity results concerning vector bundles of arbitrary rank, whether in the analytic or algebraic setting, see (Gabber~O., {\og On purity theorems for vector bundles\fg}, \emph{Internat. Math. Res. Notices} (2002), \numero 15, p\ptbl 783--788).} A regular noetherian local ring is factorial (and hence parafactorial if its dimension is $\geq 2$).
\item
A noetherian local ring of dimension $\geq 4$ that is a \emph{complete intersection} is parafactorial.
\end{enumeratei}
\end{theoreme}

\begin{corollaire}[Conjecture of\label{XI.3.14} \sisi{SAMUEL}{Samuel}\ndemark]\ndetext{for a proof in the same spirit, but more elementary, see Call~F.\ \&\ Lyubeznik~G., {\og A simple proof of Grothendieck's theorem on the parafactoriality of local rings\fg}, in \emph{Commutative algebra: syzygies, multiplicities, and birational algebra (South Hadley, MA, 1992)}, Contemp. Math., vol.~159, American Mathematical Society, Providence, RI, 1994, p\ptbl 15--18.} A noetherian local ring $A$ that is a complete intersection and is factorial in codimension $\geq 3$ (\ie $\dim A_{{\pp}}\leq 3$ implies that $A_{{\pp}}$ is factorial) is factorial.
\end{corollaire}

\emph{Proof of the corollary}

We argue by induction on the dimension of $A$.

If $\dim A\leq 3$, then $A$ is factorial by hypothesis.

If $\dim A>3$, then by the induction hypothesis, noting that a localization of a complete intersection is again a complete intersection, all localizations of $A$ other than $A$ are factorial. By Theorem~\Ref{XI.3.13}~(ii), $A$ is parafactorial, and therefore factorial by~\Ref{XI.3.10}.

Let us prove~\Ref{XI.3.13}~(i) (following \sisi{KAPLANSKY}{Kaplansky})\sfootnote{This is the proof reproduced in \EGA IV~21.11.1.}.

Let $A$ be a regular noetherian local ring, and set $\dim A=n$. If $n=0$ or $1$, the result is known. Suppose $n\geq 2$; we argue by induction on $n$, and suppose $n\geq 2$ and the theorem proved for rings of dimension $<n$. Set $X'=\Spec(A)$ and $X=X'-\{x\}$, where $x=\rr(A)$. The localizations of $A$ other than $A$ are regular and of dimension $<n$, and hence factorial. Moreover, $\prof A = \dim A\geq 2$. It is therefore enough to prove that $\Pic(X)=0$ (cor\ptbl\Ref{XI.3.10}). Thus let ${\Lcal}$ be an invertible $\OX$-Module; we know that it can be extended to a coherent $\Oo_{X'}$-Module ${\Lcal}'$. There exists a resolution of ${\Lcal}'$ by free $\Oo_{X'}$-Modules:
\[
0\from {\Lcal}'\from {\Lcal}'_1\from\cdots\from{\Lcal}'_n\from 0{\quoi},
\]
because\pageoriginale the cohomological dimension of $A$ is finite. By restriction to $X$ we obtain a finite free resolution. It is therefore enough to prove the following lemma:

\begin{lemme} \label{XI.3.15} Let $(X, \OX)$ be a ringed space, and let ${\Lcal}$ be a locally free $\OX$-Module that admits a finite resolution by \emph{finite-type free} modules. Then $\det({\Lcal}) \simeq \OX$.
\end{lemme}

Recall that $\det({\Lcal})$ is defined as the highest exterior power of ${\Lcal}$. In the case under consideration, $\det({\Lcal}) \simeq{\Lcal}$, because ${\Lcal}$ is invertible, and hence the lemma gives the conclusion. Let us prove this lemma. Let
\[
0\from {\Lcal}_0\from {\Lcal}_1\from {\Lcal}_2\from\cdots\from{\Lcal}_n\from 0{\quoi},
\]
be the asserted exact sequence, where ${\Lcal}_0={\Lcal}$. Since everything is locally free, we have:
\[
\tbigotimes_{0\leq i\leq n}(\det({\Lcal}_i))^{(-1)^i} \simeq \OX
\]
and all the ${\Lcal}_i$, $i>0$, are free, and hence so are their determinants, and consequently so is the determinant of ${\Lcal}_0 = {\Lcal}$.\qed

\medskip
It remains for us to prove part~(ii) of the theorem. First, let us prove a lemma that will allow us to proceed by induction:

\begin{lemme} \label{XI.3.16}
Let $A$ be a noetherian local ring that is a quotient of a regular ring. Let $t\in\rr(A)$ be an $A$-regular element. Suppose that $A$ is complete for the $t$-adic topology. Set $X'=\Spec(A)$, $Y'=V(t)\simeq \Spec(B)$, $B=A/tA$, $X=X'-\{x\}$, $Y=Y'-\{x\}$, and $x=\rr(A)$. Suppose that:
\begin{enumeratea}
\item
for every $y\in X$ closed in $X$, one has $\prof\Oo_{X, y}\geq 2$,
\item
$\prof A/tA\geq 3$,
\end{enumeratea}
\noindent
then\pageoriginale the map $\Pic(X)\to\Pic(Y)$ is injective. In particular, if $B$ is parafactorial, then $A$ is also parafactorial.
\end{lemme}

We know that a) implies $\Lef(X, Y)$ by \Ref{X}~\Ref{X.2.1}. If we prove that $\H^1(Y, P_n)=0$ for every $n\geq 0$, then we will know \sisi{\ignorespaces}{by} (\Ref{XI.2.2}) that $\Pic(X)\to\Pic(Y)$ is injective. If, moreover, $B$ is parafactorial, then we will know that $\Pic(Y)=0$ (\Ref{XI.3.5}), and hence $\Pic(X)=0$; moreover, $\prof(A)=\prof(B)+1\geq 3+1\geq 2$, since $t$ is $A$-regular, and thus $A$ will be parafactorial by~\Ref{XI.3.5}.

Let $\Ii=\widetilde{(tA)}$ be the $\Oo_{X'}$-Module associated with the ideal $tA$. In \numero \Ref{XI.1}, we set $P_n = \sisi{(\Ii^{n+1}/\Ii^{n+2}){|Y}}{(\Ii^{n+1}/\Ii^{n+2})_{|Y}}$ for every $n\geq 0$. Since $t$ is $A$-regular, $P_n\simeq \Oo_Y$. \emph{It therefore remains for us to prove that $\H^1(Y, \Oo_Y)=0$.} Now $Y = Y'-\{x\}$ is an open subset of $Y'$, and hence we have an exact sequence~(\Ref{I}~(\Ref{eq:I.27})):
\[
\H^1(Y', \Oo_{Y'})\to\H^1(Y, \Oo_{Y})\to\H^2_x(Y', \Oo_{Y'}){\quoi},
\]
whose right-hand term is zero by hypothesis~b), and whose left-hand term is zero because $Y'$ is affine.\qed

\begin{lemme} \label{XI.3.17}
Keeping the hypotheses of~\Ref{XI.3.16}, suppose in addition that:
\begin{itemize}
\item[c)] for every $y$ closed in $Y$, one has $\prof\Oo_{X, y}\geq 3$,
\item[d)] $\prof A/tA\geq 4$ (stronger than b)),
\item[e)] for every $y$ closed in $X$ with $y\in Y$, the ring $\Oo_{X, y}$ is parafactorial.
\end{itemize}
Then the map $\Pic(X)\to\Pic(Y)$ is an isomorphism; in particular, for $A$ to be parafactorial, it is necessary and sufficient that $B$ be parafactorial.
\end{lemme}

\enlargethispage{5mm}%
We know~(\Ref{X}~\Ref{X.2.1}) that a)~and~c) imply $\Leff(X, Y)$. Moreover, by the argument just given, d) implies that $\H^i(Y, P_n)=0$ for every $n\geq 0$ and $i=1$ or $i=2$. Moreover, for every open neighborhood $U$ of $Y$ in $X$, the complement of $U$ in $X$ consists of finitely many closed points. By~e) and Theorem~\Ref{XI.3.12}, we deduce that $\Pic(X)\to \Pic(Y)$ is an isomorphism. On the other hand, $\prof A\geq \prof B\geq 2$; by the criterion~\Ref{XI.3.5}, we deduce\pageoriginale that $A$ is parafactorial if and only if $B$ is.

Let us now prove~\Ref{XI.3.13}~(ii). Let $R$ be a regular noetherian local ring. Let $(t_1, \ldots, t_k)$ be an $R$-sequence. Set $B=R/(t_1, \ldots, t_k)$ and suppose that $\dim B\geq 4$. We must prove that $B$ is parafactorial. We argue by induction on~$k$. If $k=0$, then $B$ is regular, and hence factorial by \Ref{XI.3.13}~(i), and therefore parafactorial by~\Ref{XI.3.10}. Suppose $k\geq 1$ and the theorem proved for $k'<k$. Set $A=R/(t_1, \ldots, t_{k-1})$, so that $B=A/t_k A$. We may suppose that $B$ is complete by~\Ref{XI.3.7}. By the induction hypothesis, $A$ is parafactorial. Let us prove that Lemma~\Ref{XI.3.17} applies. We have assumed that $B$ is complete, and hence so is $A$, and therefore $A$ is complete for the $t_k$-adic topology. If $x\in X$ and $x$ is closed in $X$, then $A_x$ is a complete intersection of dimension $\geq 4$, with $k'<k$. By the induction hypothesis, $A_x$ is parafactorial, and moreover has depth $\geq 4$. This gives a), c), and~e). Moreover, $\dim A\geq 5$, which gives~d).\qed

\begin{theoreme} \label{XI.3.18}
Let $X$ be a locally noetherian prescheme, and let $\Ii$ be a coherent sheaf of ideals of $X$. Set $Y=V(\Ii)$. Let $n$ be an integer. Suppose that:
\begin{enumeratei}
\item
one has $\sisi{\leff}{\Leff}(X, Y)$ (\cf examples \Ref{X}~\Ref{X.2.1}~and~\Ref{X}~\Ref{X.2.2}),
\item
for every $p\geq n$, one has $\H^i(Y, \Ii^{p+1}/\Ii^{p+2})=0$ for $i=1$ and $i=2$,
\item
for every open subset $U\supset Y$ and every $x\in X-U$, the ring $\OXx$ is regular of dimension $\geq 2$ or a complete intersection of dimension $\geq 4$.
\end{enumeratei}
\noindent
Then, for every open subset $U\supset Y$ and every integer $p\geq n$, the homomorphisms
\[
\Pic(X)\to\Pic(U)\to\Pic(Y_p)
\]
are isomorphisms, where $Y_p$ denotes the prescheme $(Y, \OX/\Ii^{p+1})$.
\end{theoreme}

It is enough to combine \Ref{XI.3.12}~and~\Ref{XI.3.13}.

% END INCLUDED active_source/parts/11_expose_XI.tex

% BEGIN INCLUDED active_source/parts/12_expose_XII.tex
\chapter[Applications to projective algebraic schemes]
{Applications\\ to projective algebraic schemes} \label{XII}

\renewcommand{\theequation}{\arabic{equation}}

\section[Projective duality theorem and finiteness theorem]
{Projective duality theorem and finiteness theorem\protect\sfootnotemark}\label{XII.1}

The\pageoriginale following theorem,\sfootnotetext{The present expos\'e, written in January~1963, is considerably more detailed than was the oral presentation in June~1962.} essentially contained in FAC\sfootnote{J.-P.~Serre, \sisi{Faisceaux alg\'ebriques Coh\'erents, Ann. of Math. t.61 (1955) 197-278}{{\og Faisceaux alg\'ebriques coh\'erents\fg}, \emph{Ann. of Math.} \textbf{61} (1955), p\ptbl197-278}.} (except that, at the time, Serre did not have at his disposal the language of Ext for sheaves of modules\nde{Readers interested in the history of mathematics will profitably consult the letter that Grothendieck wrote to Serre on 15 December 1955 and Serre's reply of 22 December of the same year; see \emph{Correspondance Grothendieck-Serre}, edited by Pierre Colmez and Jean-Pierre Serre, Documents Math\'ematiques, vol.~2, Soci\'et\'e Math\'ematique de France, Paris, 2001.}), is the global analogue of the local duality theorem (Expos\'e \Ref{IV}), which was modeled on it.

\begin{theoreme} \label{XII.1.1}
Let $k$ be a field, $X=\PP_k^r$ \sisi{the standard projective scheme}{projective space} of dimension $r$ over $k$, and $F$~a variable coherent module on $X$; then there is an isomorphism of $\partial$-functors in~$F$:
\steco
\begin{equation} \label{eq:XII.1} {}\H^i(X, F)' \isomfrom \Ext^{r-i}(X;F, \omx^r),
\end{equation}
where we set
\steco
\begin{equation} \label{eq:XII.2} {}\Omega_{X/k}^r= \Oo_{\PP_k^r}(-r-1).
\end{equation}
\end{theoreme}

\begin{remarquestar}
This module is of course also the module of relative differentials of degree $r$ of $X$ over $k$. In this form, the theorem remains valid if $X$ is a proper smooth scheme over $k$ (for the projective case, see A. Grothendieck, \og \sisi{Th\'eor\`emes de dualit\'e pour les faisceaux alg\'ebriques coh\'erents}{\emph{Th\'eor\`emes de dualit\'e pour les faisceaux alg\'ebriques coh\'erents}}\fg, S\'eminaire Bourbaki, May 1957)\sfootnote{For a more general duality theorem, \cf the Hartshorne seminar cited at the end of \Exp \Ref{IV}.}. When $F$ is locally free, one recovers Serre's duality theorem $\H^i(X, F)' \isomto \H^{r-i}(\SheafHom_{\OX}(F, \omx^r))$. Theorem \Ref{XII.1.1} (which also recovers the case of a projective $X$ over $k$, as in {\loccit}) will suffice for our purposes.
\end{remarquestar} The homomorphism\pageoriginale (\Ref{eq:XII.1}) is deduced from the Yoneda pairing
\steco
\begin{equation} \label{eq:XII.3} {}\H^i(X, F) \times \Ext^{r-i}(X;F, \omx^r) \to \H^r(X, \omx^r),
\end{equation}
and from a well-known isomorphism (\cf FAC, or \EGA III 2.1.12):
\steco
\begin{equation} \label{eq:XII.4} {}\H^r(X, \omx^r) = \H^r(\PP_k^r, \Oo_{\PP_k^r}(-r-1)) \isomto k.
\end{equation}
To show that (\Ref{eq:XII.1}) is an isomorphism, one proceeds as in the case of the local duality theorem, noting that $\H^r(X, F)$, as a functor of $F$, is right exact (since $\H^n(X, F)=0$ for $n>r$), and that every coherent module is isomorphic to a quotient of a direct sum of modules of the form $\Oo(-m)$, with $m$ large. By descending induction on $i$, this reduces us to checking the assertion for a sheaf of the form $\Oo(-m)$, where it is contained in the well-known explicit computations (FAC, or \EGA III 2.1.12). We may moreover assume $-m \leq -r-1$, in which case $\H^i(X, \Oo(-m))=0$ for $i\neq r$.

\begin{corollaire} \label{XII.1.2} For a given coherent $F$, and for $m$ sufficiently large, there is a canonical isomorphism
\steco
\begin{equation} \label{eq:XII.5} {}\H^i(X, F(-m))'\isomto\H^0(X, \SheafExt_{\OX}^{r-i}(F, \omx^r)(m))
\end{equation}
(where $'$ denotes the vector-space dual).
\end{corollaire}
Indeed, on \sisi{a projective scheme}{projective space} $X$, for every pair of coherent sheaves $F,G$ and for $m$ sufficiently large, there is a canonical isomorphism:
\steco
\begin{equation} \label{eq:XII.6} {}\Ext^n(X;F(-m), G)\simeq \Ext^n(X;F, G(m)) \isomto \H^0(X, \SheafExt_{\OX}^{n}(F, G)(m)),
\end{equation}
(the isomorphism between the first two terms being trivially valid for every $m$), as follows from the global $\Ext$ spectral sequence
\[
\H^p(X, \SheafExt^q_{\OX}(F, G(m))) \To \Ext^{\boule}(X;F, G(m)),
\]
which\pageoriginale degenerates for $m$ sufficiently large, by virtue of the fact that
\[
\SheafExt^q_{\OX}(F, G(m)) \simeq \SheafExt^q_{\OX}(F, G)(m),
\]
and that the $\SheafExt^q_{\OX}(F, G)$ are coherent sheaves. Thus (\Ref{eq:XII.5}) follows from (\Ref{eq:XII.6}) and (\Ref{eq:XII.1}).

\begin{corollaire} \label{XII.1.3}
\sisi{Equivalent conditions for given $i,F$}{For given $i,F$, the following conditions are equivalent}:
\begin{enumeratei}
\item
$\H^i(X, F(-m))=0$ for $m$ large.
\item[\textup{(i bis)}] $\H^i(X, F(\cdot))= \sisi{\coprod}{\bigoplus}_{m\in\ZZ}\H^i(X, F(-m))$ is a finite-type $S$-module, where $S=k[t_0, \ldots, t_r]$.
\item
$\SheafExt^{r-i}_{\OX}(F, \omx^r)=0$.
\item[\textup{(ii bis)}] $\SheafExt^{r-i}_{\OX}(F, \OX)=0$.
\item
$\H^i_x(F_x)=0$ for every closed point $x$ of $X$.
\item
$\H^{i+1}_x(\widetilde{F}_x)=0$ for every closed point $x$ of the punctured affine cone $\widetilde{X}=\Spec(S)-\Spec(k)$ over $X$, where $\widetilde{F}$ denotes the inverse image of $F$ under the canonical morphism $\widetilde{X}\to X$.
\end{enumeratei}
\end{corollaire}

\begin{proof}
(i) \SSI (i bis), since the submodule of $\H^i(X, F(\cdot))$ that is the sum of the homogeneous components of degree $\geq \nu$ is of finite type over $S$ (in fact, for $i\neq 0$, it is even of finite type over $k$), (\cf FAC or \EGA III 2.2.1 and 2.3.2).

(i) \SSI (ii) by Corollary \Ref{XII.1.2}.

(ii) \SSI (ii bis), because $\Omega_{X/k}^r$ is locally isomorphic to $\OX$.

(ii bis) \SSI (iii) by the local duality theorem for $\OXx$ (which is a regular\pageoriginale local ring of dimension $r$), according to which the \og dual\fg of $\SheafExt^{r-i}_{\OX}(F, \OX)_x$ is identified with $\H^i_x(F_x)$ (\hyperref[V.2.1]{V 2.1)}.

(ii bis) is equivalent to the analogous relation
\[
\SheafExt^{r-i}_{\Oo_{\widetilde{X}}}(\widetilde{F}, \Oo_{\widetilde{X}})=0
\]
(because $\widetilde{X}\to X$ is faithfully flat, so the inverse image of $\SheafExt^{r-i}_{\OX}(F, \OX)$ is isomorphic to $\SheafExt^{r-i}_{\Oo_{\widetilde{X}}}(\widetilde{F}, \Oo_{\widetilde{X}})$), and this last relation is equivalent to (iv) by the local duality theorem for the local ring $\Oo_{\widetilde{X},x}$, which is regular of dimension $r+1$.
\skipqed
\end{proof}

In particular, applying this to all $i\leq n$, we find:

\begin{corollaire} \label{XII.1.4}
Equivalent conditions for given $n,F$:
\begin{enumeratei}
\item
$\H^i(X, F(-m))=0$ for $i\leq n$ and $m$ large.
\item[\textup{(i bis)}] $\H^i(X, F(\cdot))$ is a finite-type $S$-module for $i\leq n$.
\item
$\prof(F_x)>n$ for every closed point $x$ of $X$.
\item
$\prof(\widetilde{F}_x)>n+1$ for every closed point $x$ of $\widetilde{X}$.
\end{enumeratei}
\end{corollaire}

The point of Corollaries \Ref{XII.1.3} and \Ref{XII.1.4} is to express a global condition (i) or~(i~bis) in terms of local conditions, \sisi{\ignorespaces}{namely} the vanishing of local invariants such as $\H_x^i(X, F_x)$ or $\H_x^i(X, \widetilde{F}_x)$, or an inequality involving depth. In this form, these results remain trivially valid for an arbitrary projective scheme $X$, and a very ample invertible sheaf $\OX(1)$ on $X$, as one sees by means of a suitable projective immersion $X\to\PP_k^r$. (Of course, conditions \Ref{XII.1.3}~(ii) and \Ref{XII.1.3}~(ii bis) are no longer equivalent to the others in this general case, unless, for example, one assumes that $X$ is regular.) One can moreover generalize to the case of a projective morphism $X\to S$ as follows:

\begin{proposition} \label{XII.1.5}
Let\pageoriginaled $f:X\to S$ be a projective morphism, with $S$ noetherian, let $\OX(1)$ be an invertible module on $X$ that is very ample relative to $S$, let $F$ be a coherent module on $X$, flat over $S$, let $s$ be a point of $S$, let $X_s$ be the fiber of $X$ at $s$ (regarded as a projective scheme over $k(s)$), let $F_s$ be the sheaf induced on $X_s$ by $F$, and finally let $i$ be an integer. Suppose that, for every closed point $x$ of $X_s$, one has $\H^i_x(F_{s, x})=0$ (for example, $\prof(F_{s, x})>i$). Then there exists an open neighborhood $U$ of $s$ such that the same condition holds for every $s'\in U$. Moreover, for such a $U$, one has
\[
\R^if_\ast(F(-m))= 0 \; \text{for}\; m \; \text{large, }
\]
and if $\Ss$ is a quasi-coherent graded algebra on $S$, generated by $\Ss^1$, and defining~$X$ together with $\OX(1)$ as $X=\Proj(\Ss)$, $\OX(1)=\Proj(\Ss(1))$, then the $S$-module
\[
\R^if_\ast(F(\cdot)) = \sisi{\coprod}{\tbigoplus}_{m\in\ZZ}\R^if_\ast(F(m))
\]
is of finite type over $U$.
\end{proposition}

Embed $X$ in an $X'=\PP_S^r$ in such a way that $\OX(1)$ is induced by $\Oo_{\sisi{P}{{X'}}}(1)$ (this is possible provided that $S$ is replaced by an affine neighborhood of $s$). Put\refstepcounter{toto}\label{P5}\sfootnote{The first part of \Ref{XII.1.5} can also be obtained immediately by applying the purely local statement \EGA IV 12.3.4 to the preceding $\underline{E}^j$, thereby bypassing most of the proof that follows.} for every integer $j$, and every $t\in S$:
\steco
\begin{equation} \label{eq:XII.7} {}\underline{E}^j(t)= \SheafExt^j_{\Oo_{X'_t}} (F_t, \Oo_{X'_t}(-r-1))
\end{equation}
Thus $\underline{E}^j(t)$ is a coherent module on $X_t$; I claim that, as $t$ varies, the family of these modules is \og constructible\fg in the following sense: for every $t\in S$ there exists a nonempty open subset $V$ of the closure $\overline{t}$, endowed with the induced reduced structure, and a coherent module $\underline{E}^j(V)$ on $X_V=X\times_SV$, \emph{flat} over $V$, such that, for every $t'\in V$, $\underline{E}^j(t')$ is isomorphic to the module induced by $\underline{E}^j(V)$ on $X_{t'}$. To verify this assertion, putting $Z=\overline{t}$ with its induced structure, we consider the coherent modules\pageoriginale
\[
\underline{E}^j(Z)= \SheafExt^j_{\Oo_{X'_Z}} (F_Z, \Oo_{X'_Z}(-r-1))
\]
(where the subscript $Z$ still means that one takes the induced object \sisi{over}{over} $Z$), and we take for $V$ a nonempty open subset of $Z$ such that the modules $\underline{E}^j(Z)$ are flat \sisi{over}{over} $V$: this is possible, since one verifies immediately that $\underline{E}^j(Z)=0$ for $j$ outside the interval $[0, r]$, and one can apply \SGA 1 IV 6.11. We then take $\underline{E}^j(V)=\underline{E}^j(Z)|_{X_V}$, and one readily verifies that it has the required property.

It follows from the preceding remark that there exists a finite partition of $S$ formed by sets $V_\alpha$ of the form $V=V(t)$ as above (noetherian induction), and, applying Serre's theorem \EGA III 2.2.1 to the $\underline{E}^j(V_\alpha)$, one sees that there exists an integer $m_0$ such that
\[
\R^if_{V_\alpha\ast}(\underline{E}^j(V_\alpha)(m))=0\textup{ for } i\neq 0, m\geq m_0, \textup{ for all }j,
\]
from which it follows, using the flatness of $\underline{E}^j(V_\alpha)$ over $V_\alpha$ and straightforward K\"unneth-type relations (\cf \EGA III par\ptbl7), that
\[
\H^i(X_t, \underline{E}^j(t)(m))=0\textup{ for } i\neq 0, m\geq m_0, \textup{ for all }j,
\]
for every $t\in V_\alpha$, and hence for every $t$, since the $ V_\alpha$ cover $S$. From this and the spectral sequence for global Ext it follows, by virtue of \Ref{XII.1.1}, as in the proof of \Ref{XII.1.2}, that there is an isomorphism
\steco
\begin{equation} \label{eq:XII.8} {}\H^i(X_t, F_t(-m))' \isomfrom \H^0(X_t, \underline{E}^{r-i}(t)(m))\textup{ for } m\geq m_0,
\end{equation}
for every integer $i$, and every $t\in S$.

Now use the hypothesis on $F_s$, which can be written
\steco
\begin{equation} \label{eq:XII.9} {}\underline{E}^{r-i}(s)=0,
\end{equation}
and which, by virtue of (\Ref{eq:XII.8}), is equivalent to
\steco
\begin{equation} \label{eq:XII.10} {}\H^i(X_s, F_s(-m))=0 \; \text{for}\; m\geq m_0.
\end{equation}
Since $F$, and hence $F(-m)$, is flat over $S$, it follows from
the K\"unneth-type relations already invoked that (for a given $m$) the
same relation (\Ref{eq:XII.10}) holds when\pageoriginale $s$ is replaced by a
$t$ near $s$, in particular by every generization $t$ of
$s$. By virtue of (\Ref{eq:XII.8}), we shall therefore have, for such a
generization, \steco
\begin{equation} \label{eq:XII.11} {}\underline{E}^{r-i}(t)=0,
\end{equation}
and the set of $t\in \sisi{s}{S}$ for which this relation holds is plainly constructible (since it induces an open subset on each $V_\alpha$); as it contains the generizations of $s$, it contains an open neighborhood $U$ of $s$. This proves the first assertion of \Ref{XII.1.5}. Moreover, for $t\in U$, it follows from (\Ref{eq:XII.11}) and (\Ref{eq:XII.8}) that
\steco
\begin{equation} \label{eq:XII.12} {}\H^i(X_t, \sisi{F_t(-m)}{F_t(-m))}=0 \; \text{for}\; m\geq m_0, \; t\in U,
\end{equation}
which, by virtue of the K\"unneth-type relations, implies (and in fact is much stronger than)
\steco
\begin{equation} \label{eq:XII.13} {}\R^if_\ast(F(-m))=0 \; \text{on}\; U, \; \text{for} \; m\geq m_0.
\end{equation}
This proves the second assertion of \Ref{XII.1.5}. Finally, the last assertion follows immediately by proceeding as at the beginning of the proof of \Ref{XII.1.3}.

\refstepcounter{subsection}\label{XII.1.6}
\begin{enonce*}[remark]{Remark 1.6\sfootnotemark} \sfootnotetext{This remark is made more precise by footnote \sisi{5.}{\pageref{P5}.}} The proof simplifies considerably (eliminating all consideration of constructibility) when one assumes from the outset that the hypothesis made for $s\in S$ holds for every $s'\in S$. In fact, when one assumes that $F_s$ has depth $>i$ at the closed points of $X$, there is a general statement, \emph{local in nature on} $X$, which says that the same condition holds for all the $X_t$, provided that $X$ is replaced by a suitable open neighborhood of the fiber $X_s$ (in other words, a certain subset of $X$, defined by conditions on the modules induced by $F$ on the fibers, is open, \cf \EGA IV). Since $f$ is proper here, this neighborhood can therefore be taken to be of the form $f^{-1}(U)$, which recovers the first assertion of \Ref{XII.1.5} without laborious d\'evissage. In this general case, one can still prove by the method of \loccit that the first assertion of \Ref{XII.1.5} (proved here globally, using the fact that $X$ is projective over $S$) again follows from a purely local statement on $X$ (which the reader may formulate explicitly if useful).
\end{enonce*}

\section[Lefschetz theory for a projective morphism: comparison]
{Lefschetz theory for a projective morphism: Grauert comparison theorem} \label{XII.2}

\setcounter{equation}{13}
It is\pageoriginale the following theorem:

\begin{theoreme} \label{XII.2.1}
Let $f:X\to S$ be a projective morphism, with $S$ Noetherian, let $\OX(1)$ be an invertible module on $X$, relatively ample over $S$, let $Y$ be the prescheme of zeros of a section $t$ of $\OX(1)$, let $J$ be the ideal defining $Y$, let $X_n$ be the closed subprescheme of $X$ defined by $J^{n+1}$, let $\widehat{X}$ be the formal completion of $X$ along $Y$, let $\widehat{f}:\widehat{X}\to S$ be the composite morphism $\widehat{X}\to X\to S$, and let $F$ be a coherent module on $X$, flat over $S$. Assume in addition that, for every $s\in S$, the module $F_s$ induced on the fiber $X_s$ has depth $>n$ at the points of that fiber, and that $t$ is $F$-regular. Under these conditions:
\begin{enumeratei}
\item
The canonical homomorphism
\[
\R^if_\ast(F) \to \R^i\widehat{f}_\ast(\widehat{F})
\]
is an isomorphism for $i<n$, and a monomorphism for $i=n$.
\item
The canonical homomorphism
\[
\R^i\widehat{f}_\ast(\widehat{F}) \to \varprojlim_m \R^if_\ast(F_m)
\]
is an isomorphism for $i\leq n$.
\end{enumeratei}
\end{theoreme}

\begin{proof}
We immediately reduce to the case where $S$ is affine, and then to proving the

\begin{corollaire} \label{XII.2.2}
Under the conditions of \Ref{XII.2.1}, assume in addition that $S$ is affine. Then:
\begin{enumeratei}
\item
The canonical homomorphism
\[
\H^i(X, F) \to \H^i(\widehat{X}, \widehat{F})
\]
is\pageoriginale an isomorphism for $i<n$, and a monomorphism for $i=n$.
\item
The canonical homomorphism
\[
\H^i(\widehat{X}, \widehat{F}) \to \varprojlim_m \H^i(X_m, F_m)
\]
is an isomorphism for $i\leq n$.
\end{enumeratei}
\end{corollaire}

After replacing $\OX(1)$ by a tensor power, and $t$ by a power of~$t$, we may assume that $\OX(1)$ is very ample relative to $S$. On the other hand, since $t$, and hence $t^m$, is $F$-regular, multiplication by $t^m$, regarded as a homomorphism from $F(-m)$ to $F$, is injective; therefore, for every $m\geq 0$, there is an exact sequence:
\steco
\begin{equation} \label{eq:XII.14}
0\to F(-m) \lto{t^m} F \to F_m \to0,
\end{equation}
and hence an exact cohomology sequence
\[
\H^i(X, F(-m)) \to \H^i(X, F) \to \H^i(X, F_m) \to \H^{i+1}(X, F(-m)).
\]
But by \Ref{XII.1.5}, we have $\H^i(X, F(-m))=0$ for $i\leq n$ and $m$ sufficiently large, which proves the

\begin{lemme} \label{XII.2.3}
For $m$ large, the canonical homomorphism
\[
\H^i(X, F) \to \H^i(X, F_m)
\]
is bijective if $i<n$, and injective if $i=n$.
\end{lemme}

\enlargethispage{-2\baselineskip}%
This shows that, for $i<n$, the inverse system $(\H^i(X_m, F_m))_{m\geq 0}$ is essentially constant, and \afortiori satisfies the Mittag-Leffler condition; hence (\sisi{given that}{in view of} $\widehat{F}= \varprojlim F_m$), we obtain (ii) from $\EGA 0_{\textup{III}}$ 13.3. On the other hand, (i) follows trivially from this, taking \Ref{XII.2.3} into account.

\refstepcounter{subsection}\label{XII.2.4}
\begin{enonce*}{Corollary 2.4\ndemark}\ndetext{See Corollary I.1.4 of the article by Mme Raynaud (Raynaud~M., {\og Th\'eor\`emes de Lefschetz en cohomologie des faisceaux coh\'erents et en cohomologie \'etale. Application au groupe fondamental\fg}, \emph{Ann. Sci. \'Ec. Norm. Sup. (4)} \textbf{7} (1974), p\ptbl 29--52).}
Let\pageoriginale $f:X\to S$ be a projective flat morphism, with $S$ locally Noetherian, let $\OX(1)$ be an invertible module on $X$, relatively ample over $S$, let $t$ be a section of this module that is $\OX$-regular, let $Y$ be the closed subprescheme of zeros of $t$, and let $\widehat{X}$ be the formal completion of $X$ along $Y$. Assume that, for every $s\in S$, $X_s$ has depth $\geq 1$ (\resp depth $\geq 2$) at its closed points. Then, for every open neighborhood $U$ of $Y$, the functor
\[
F \mto \widehat{F}
\]
from the category of locally free coherent modules on $U$ to the category of locally free coherent modules on $\widehat{X}$ is faithful (\resp fully faithful, \ie the Lefschetz condition ($\Lef$) of \Ref{X}~\Ref{X.2} holds).
\end{enonce*}

For two locally free Modules $F$ and $G$ on $U$, introduce the Module
\[
H=\SheafHom_{\Oo_U}(F, G)
\]
we are reduced to proving that the canonical homomorphism
\steco
\begin{equation} \label{eq:XII.15} {}\H^0(U, H) \to \H^0(\widehat{U}, \widehat{H})
\end{equation}
is injective (\resp bijective). Now the Modules $H_s$ have depth $\geq 1$ (\resp $\geq 2$) at the closed points of $X_s$, and hence we can apply \Ref{XII.2.1}, which gives the conclusion \Ref{XII.2.4} in the case where $U=X$. In the case of an arbitrary $U$, note that the question is local on $S$, so we may suppose that $S$ is affine. Then every coherent Module on $X$ is a quotient of a coherent locally free Module (since $\OX(1)$ is an invertible \sisi{absolutely}{relatively} ample Module on $X$). Since the dual Module $H'=\SheafHom(H, \Oo_U)$ extends to a coherent Module on $X$, which is therefore isomorphic to a cokernel of a homomorphism of locally free Modules on $X$, it follows by transposition that one can find a homomorphism
\[
u':{L'}^0 \to {L'}^1
\]
of\pageoriginale locally free Modules on $X$, inducing a homomorphism
\[
u:L^0 \to L^1
\]
of locally free Modules on $U$, such that there is an exact sequence
\[
0\to H \to L^0 \lto{u} L^1.
\]
Using the five lemma (which becomes the three lemma), and the left exactness of the functor $H^0$, we are reduced to proving that (\Ref{eq:XII.15}) is injective (\resp bijective) when $H$ is replaced there by $L^0$, $L^1$, which brings us back to the case where $H$ is induced by a locally free Module $H'$ on $X$. Moreover, in the non-parenthetical case this reduction is even unnecessary, since the kernel of (\Ref{eq:XII.15}) consists in any event of the sections of $H$ on~$U$ that vanish in some suitable open neighborhood $V$ of $Y$; but the restriction homomorphism $\H^0(U, H)\to\H^0(V, H)$ is injective, since $H$ has depth $\geq 1$ at the points of every closed subset $Z$ of $X$ not meeting $Y$ (\cf the lemma below). In the parenthetical case, we are reduced to proving that
\[
H^0(X, H')\to\H^0(U, H')
\]
is bijective, which follows from the fact that $H'$ has depth $\geq 2$ at all the points of a closed subset $Z=X- U$ of $X$ not meeting $Y$. It therefore remains only to prove the

\begin{lemme} \label{XII.2.5} Let $F$ be a coherent Module on $X$, flat over $S$, such that for every $s\in S$, $F_s$ has depth $\geq n$ at every closed point of $X_s$. Then, for every closed subset $Z$ of $X$ not meeting $Y$, $F$ has depth $\geq n$ at all the points of $Z$.
\end{lemme}

Indeed, for every $x\in X$, setting $s=f(x)$, we have
\steco
\begin{equation} \label{eq:XII.16} {}\prof(F_x)\geq \prof(F_{s, x}),
\end{equation}
as can be seen by lifting arbitrarily a maximal $F_{s, x}$-regular sequence of elements of\pageoriginale $\rr(\Oo_{X_s, x})$, which gives an $F_x$-regular sequence by \SGA 1 IV 5.7. Now if $x$ belongs to a $Z$ as in Lemma~\Ref{XII.2.5}, then $x$ is necessarily closed in $X_s$; in other words, $Z$ is \emph{finite} over $S$. Indeed, $Z$ (endowed with a structure induced from $X$) is projective over $S$ as a closed subprescheme of $X$, which is projective, and $Z$ is affine over $S$ as a closed subprescheme of $X- Y$, which is affine.
\skipqed
\end{proof}

\begin{remarque} \label{XII.2.6}
Suppose that for every $s\in S$, the section $t_s$ of $\Oo_{X_s}(1)$ induced by $t$ is $\Oo_{X_s}$-regular (which implies by \SGA 1 IV 5.7 that $t$ is $\OX$-regular). Then the hypotheses made are stable under extension of the base $S'\to S$ ($S'$ locally noetherian). Thus the conclusion remains valid after every base change.
\end{remarque}

\section[Lefschetz theory for a projective morphism: existence]{Lefschetz theory for a projective morphism: existence theorem}\label{XII.3}%

\setcounter{equation}{16}%
\refstepcounter{subsection}\label{XII.3.1}%
\begin{enonce*}{Theorem 3.1\ndemark}
\ndetext{For a version without the flatness hypothesis, see (Raynaud~M., {\og Th\'eor\`emes de Lefschetz en cohomologie des faisceaux coh\'erents et en cohomologie \'etale. Application au groupe fondamental\fg}, \emph{Ann. Sci. \'Ec. Norm. Sup. (4)} \textbf{7} (1974), p\ptbl 29--52, Theorem II.3.3).}
Let $f:X\to S$ be a projective morphism, with $S$ noetherian, let $\OX(1)$ be an invertible Module on $X$ that is ample relative to $S$, let $X_0$ be the closed subprescheme of zeros of a section $t$ of $\OX(1)$, let $\widehat{X}$ be the formal completion of $X$ along $X_0$, let $\formelF$ be a coherent Module on $\widehat{X}$, and let $F_0$ be the Module that it induces on $X_0$. Assume moreover:
\begin{enumeratea}
\item\label{XII.3.1a}
$\formelF$ is flat over $S$.
\item\label{XII.3.1b}
For every $s\in S$, the section $t_s$ induced by $t$ on $X_s$ is $\formelF_s$-regular (which implies that $F_0$ is also flat over $S$, \cf \textup{\SGA 1 IV 5.7}).
\item\label{XII.3.1c}
For every $s\in S$, $F_{0s}$ has depth $\geq 2$ at the closed points of $X_{0s}$.
\end{enumeratea}
Assume moreover that $S$ admits an ample invertible sheaf. Under these conditions, there exists a coherent Module $F$ on $X$, and an isomorphism from its formal completion $\widehat{F}$ to $\formelF$.
\end{enonce*}

This statement will follow from the following one:

\begin{corollaire} \label{XII.3.2}
Under conditions \Ref{XII.3.1a}, \Ref{XII.3.1b}, \Ref{XII.3.1c} above, the following hold:
\begin{enumeratei}
\item
The\pageoriginale Module $\widehat{f}_\ast(\formelF)$ on $S$ is coherent; hence, for every $n$, the Module $\widehat{f}_\ast(\formelF(n))$ on~$S$ is coherent.
\item
For $n$ large, the canonical homomorphism $\widehat{f}^\ast\widehat{f}_\ast(\formelF(n))\to \formelF(n)$ is surjective.
\end{enumeratei}
\end{corollaire}

Assume the corollary for the moment, and let us prove \Ref{XII.3.1}. By virtue of the last hypothesis made in \Ref{XII.3.1}, we may reduce to the case $X=\PP^r_S$, after replacing $\OX(1)$ and $t$ by a suitable power. I claim that we may further suppose that $t_s\neq 0$ for every~$s$. Otherwise, indeed, $\formelF_s=0$ by b), or equivalently, by Nakayama, $F_{0s}=0$, \ie $s$ does not belong to the image of $\Supp F_0$ under the morphism $f_0:X_0\to S$ induced by $f$. Now this image $S'$ is open by virtue of a), b), since $F_0$ is flat over $S$, and it is clear that it suffices to prove the conclusion of \Ref{XII.3.1} in the situation obtained by restricting over $S'$, since the coherent Module $F'$ on $X|_{S'}$ thereby obtained will be the restriction of a coherent Module $F$ on $X$, which will have the required property. We may therefore suppose that, in addition to hypotheses a), b), c), the following hypotheses also hold:
\begin{enumerate}
\item[\textup{a$'$)}] $\OX$ is flat over $S$.
\item[\textup{b$'$)}] For every $s\in S$, the section $t_s$ is $\Oo_{X_s}$-regular.
\item[\textup{c$'$)}] For every $s\in S$, $\Oo_{X_{0s}}$ has depth $\geq 2$ at the closed points of $X_{0s}$. (It~suffices to choose $X=\PP_S^r$ with $r\geq 3$, which is permissible.)
\end{enumerate}
Now \Ref{XII.3.2} implies that one can find an epimorphism
\steco
\begin{equation} \label{eq:XII.17}
\widehat{L} \to \formelF \to 0,
\end{equation}
where $L$ is a Module on $X$ of the form $\sisi{f_\ast(G)}{f^\ast(G)}(-n)$, with $G$ a coherent locally free Module on $S$: indeed, for $n$ large, it suffices to express the coherent Module $\widehat{f}_\ast(\formelF)$ on $S$ as a quotient of such a $G$. On the other hand, hypotheses a), b), c) on $f$, $t$ imply that $\widehat{L}$ satisfies the same conditions a), b), c) as $\formelF$. It follows easily that the same is true of the kernel of (\Ref{eq:XII.17}), to which the same argument may therefore be applied, so that $\formelF$ is represented as the cokernel of a homomorphism
\steco
\begin{equation} \label{eq:XII.18}
\widehat{L'} \to \widehat{L},
\end{equation}
where $L$, $L'$ are locally free Modules on $X$. Now, by virtue of a$'$), the second part of c$'$), and either \Ref{XII.2.1} or \Ref{XII.2.4}, the homomorphism (\Ref{eq:XII.18}) comes from a homomorphism $L'\to L$ of Modules on $X$. It now suffices to take for $F$ the cokernel of $L'\to L$, and we are done.

It remains to prove \Ref{XII.3.2}. This had been done in the
seminar by a somewhat laborious expedient, consisting in interpreting
everything in terms of cohomology on the punctured projective cone
of $X$ relative to $S$, in order to reduce to
Theorem \Ref{IX.2.1}. A more direct and more
satisfactory method (although substantially identical) now seems to me
to be the following. It consists in noting that in \Ref{IX},
\numero \Ref{IX.2}\refstepcounter{toto}\label{XII.14} (and with the
notation of that expos\'e), the hypothesis that the morphism
$f:\Xx\to\Xx'$ is adic is not used anywhere in the
proof of \Ref{IX.2.1}, via $\EGA 0_{\textup{III}}$ 13.7.7;
it suffices instead to suppose that $\Xx$ is also adic,
and to choose two ideals of definition $\Jj$ for $\Xx'$, $\Ii$
for $\Xx$, such that $\Jj\Oo_\Xx\subset\Ii$, and to define
$\Ss=\gr_\Jj(\Oo_{\Xx'})$, and consider $\gr_\Ii(\formelF)$. In
any event, \Ref{IX.2.1} can be applied directly to the
morphism $\widehat{f}:\widehat{X}\to S$ considered in the present
subsection, where one simply takes $\Jj=0$. Thus, to verify that
$\widehat{f}_\ast(\formelF)$ is coherent, it suffices, by virtue of
\loccit, to verify that $\R^if_{0\ast}(\gr_\Ii(\formelF))$ is
coherent on $S$ for $i=0, 1$; for this, note that by virtue of
a) and b), the Module in question is none other than
$\sisi{\coprod}{\bigoplus}_{m\geq 0}\R^if_{0\ast}(F_0(-m))$, which
is indeed coherent by virtue of hypothesis c) and
\Ref{XII.1.5}.

This proves \Ref{XII.3.2}~(i). For \Ref{XII.3.2}~(ii), we shall need the

\begin{lemme} \label{XII.3.3}
Under\pageoriginaled the conditions of \Ref{XII.3.1a}, \Ref{XII.3.1b}, \Ref{XII.3.1c} of \Ref{XII.3.1}, set
\[
G_m= \widehat{f}_\ast(\formelF(\cdot)_m) = \sisi{\coprod}{\tbigoplus}_n\widehat{f}_\ast(\formelF_m(n))
\]
Then the projective system $(G_m)$ satisfies the Mittag-Leffler condition.
\end{lemme}

We may suppose that $S$ is affine, with ring $A$. Let $\Ss$ be a finitely generated graded $A$-algebra in positive degrees, and let $t'\in\Ss_1$, such that $X$ embeds in $\Proj(\Ss)$, with $\OX(1)$ induced by \sisi{$\SheafProj(\Ss(1))$}{$\Oo(1)$} and the section $t$ the image of $t'$. Endow $\Ss$ with the $\Jj$-adic filtration, where $\Jj=t'\Ss$, and consider the projective system of the $\formelF(\cdot)_m$ in the category of abelian sheaves on $X_0$. We are again under the preliminary conditions of $\EGA 0_{\textup{III}}$ 13.7.7\sfootnote{Corrected as indicated in \Ref{IX} p\ptbl\sisi{11}{\pageref{pagerectif}.}}, and moreover $\H^i(X_0, \gr(\formelF(\cdot))\sisi{}{)}$ is a finite-type Module over $\gr_\Jj(\Ss)$ for $i=0, 1$. Indeed, since $t'$ is $\formelF$-regular, one sees at once that, as a Module over $(\Ss /t'\Ss)[T]$ (of which $\gr_\Jj(\Ss)$ is a quotient), the Module in question identifies with $\H^i(X_0, F_0(\cdot))\otimes_{\Ss /t'\Ss}(\Ss /t'\Ss)[T]$; but by \Ref{XII.1.5}, $\H^i(X_0, F_0(\cdot))$ is of finite type over $\Ss$, and hence over $\Ss /t'\Ss$, for $i=0, 1$, which proves our assertion. Consequently, we are under the conditions for applying $0_{\text{III}}$ 13.7.7 with $n=1$, which proves \Ref{XII.3.3}.

This point being established (and still supposing $S$ affine, as we may for proving \Ref{XII.3.2}~(ii)), let $m_0$ be such that $m\geq m_0$ implies $\Im(G_m\to G_0)= \Im(G_{m_0}\to G_0)$, so that both sides are also equal to $\Im(\varprojlim G_m\to G_0) = \Im(\widehat{f}_\ast(\formelF(\cdot))\to f_\ast\formelF_0(\cdot))$. Now note that for large $n$, $\formelF_{m_0}(n)$ is generated by its sections, and therefore $\formelF_0(n)$ is generated by sections that lift to $\formelF_{m_0}$, and hence (by the choice of~$m_0$) that lift to $\formelF$. Thus the sections of $\formelF$ generate $\formelF_0$, and hence also~$\formelF$ by Nakayama. This proves \Ref{XII.3.2}~(ii), and therefore \Ref{XII.3.1}.

\begin{corollaire} \label{XII.3.4}
Let $f:X\to S$ be a projective flat morphism, with $S$ locally noetherian, let $\OX(1)$ be an invertible Module on $X$, relatively ample over $S$, and let $t$ be a section of this Module such that for every $s\in S$, the induced section $t_s$ on\pageoriginale the fiber $X_s$ is $\Oo_{X_s}$-regular; let $X_0$ be the closed subprescheme of zeros of $t$, and let $\widehat{X}$ be the formal completion of $X$ along~$X_0$. Suppose that for every $s\in S$, $X_{0s}$ has depth $\geq 2$ at its closed points (\ie $X_s$ has depth $\geq 3$ at the closed points of $X_{0s}$), and $X_s$ has depth $\geq 2$ at its closed points. Under these conditions, the pair $(X, X_0)$ satisfies the effective Lefschetz condition ($\Leff$) \sisi{of \Ref{X}~\Ref{X.2}}{of paragraph~\Ref{X.2} of Expos\'e~\emph{\Ref{X}}}, \ie:
\begin{enumeratea}
\item
For every open neighborhood $U$ of $X_0$ in $X$, the functor
\[
F \sisi{\rightsquigarrow}{\mto} \widehat{F}
\]
from the category of coherent locally free Modules on $U$ to the category of coherent locally free Modules on $\widehat{X}$ is fully faithful.
\item
For every coherent locally free Module $\formelF$ on $\widehat{X}$, there exist an open neighborhood~$U$ of $X_0$ and a coherent locally free Module $F$ on $U$ such that $\formelF$ is isomorphic to~$\widehat{F}$.
\end{enumeratea}
\end{corollaire}

Indeed, a) has already been noted in \Ref{XII.2.4} under weaker conditions. For b), we apply \Ref{XII.3.1}, which gives the conclusion, at least if $S$ is noetherian and admits an absolutely ample invertible Module, in particular if $S$ is affine. Indeed, if $F$ is a coherent module on $X$ such that $\widehat{F}$ is isomorphic to $\formelF$, and hence locally free, it follows that $F$ is locally free on a neighborhood $U$ of $X_0$, and $\sisi{F|U}{F|_U}$ will satisfy the required condition. But now note that by \Ref{XII.2.5}, for such an $F$, its direct image under the immersion $U\to X$ is coherent, and is moreover independent of the chosen solution $(U, F)$ (taking account \sisi{\ignorespaces}{of the fact} that two solutions coincide in a neighborhood of $X_0$, by a)). More precisely, one can find a coherent Module $F$ on $X$ and an isomorphism $\widehat{F}\isomto \formelF$, such that $F$ has depth $\geq 2$ at all the points of $X$ that are closed in their fiber, and this determines $F$ up to unique isomorphism. By this uniqueness property, the solutions of the problem obtained over the affine open subsets of $S$ glue together, giving a coherent $F$ on all of $X$ and an isomorphism $\widehat{F}\simeq \formelF$. Restricting $F$ to the open subset $U$ of points at which it is free gives what was required.

Thanks\pageoriginale to \Ref{XII.2.4} and \Ref{XII.3.4}, in the setting of a projective algebraic scheme and a \og hyperplane section\fg thereof, one can apply the general facts established in Expos\'es \Ref{X} and \Ref{XI} concerning conditions (Lef) and (Leff). Thus:

\begin{corollaire} \label{XII.3.5} Let $X$ be a projective algebraic scheme equipped with an ample invertible module $\OX(1)$, let $t$ be a section of this module that is $\OX$-regular, and let $X_0$ be the closed subscheme of zeros of $t$. Assume that $X$ has depth $\geq 2$ at its closed points (\resp and depth $\geq 3$ at the closed points of $X_0$). Then $\pi_0(X_0)\to \pi_0(X)$ is bijective; in particular, $X$ is connected if and only if $X_0$ is connected, and, after choosing a geometric base point in $X_0$, $\pi_1(X_0)\to \pi_1(X)$ is surjective. More generally, for every open $U\supset X_0$, the homomorphism $\pi_1(X_0)\to \pi_1(U)$ is surjective (\resp the homomorphism $\pi_1(X_0)\to \varprojlim_U\pi_1(U)$ is bijective). In the parenthesized case, if one assumes in addition that the local ring at every closed point of $X$ not in $X_0$ is pure (\Ref{X.3.2}) (for example, regular, or merely a complete intersection), then $\pi_1(X_0)\to \pi_1(X)$ is an isomorphism.
\end{corollaire}

Apply \Ref{X.2.4} and \Ref{XII.3.4}. Note that, in the parenthesized case, the hypothesis that $X$ have depth $\geq 3$ at the closed points of $X_0$ implies that every irreducible component of $X$ having dimension $\neq 0$ has dimension $\geq 3$ (as one sees by noting that such a component necessarily meets $X_0$, and by considering a closed point of the intersection).

\begin{remarquestar}
When $X$ is normal and has dimension $\geq 2$ at all its points, it has depth $\geq 2$ at its closed points, and the non-parenthesized hypotheses of \Ref{XII.3.5} hold. In this case there is a more elementary proof of the surjectivity of $\pi_1(X_0)\to \pi_1(X)$ using Bertini's theorem (\cf \SGA 1 X.2.10). If one assumes in addition that $X_0$ is normal and that $X$ has dimension $\geq 3$ at all its points, then the parenthesized hypotheses of \Ref{XII.3.5} hold. In this case, \Ref{XII.3.5} was established by Grauert (indeed, thanks to the normality hypothesis, one can then dispense with the existence theorem \Ref{XII.3.1} by various devices)\pageoriginale. This proof of Grauert's was the starting point of the \og Lefschetz theory\fg that is the subject of the present seminar.
\end{remarquestar}

\begin{corollaire} \label{XII.3.6} Let $X$, $\OX(1)$, $t$, and $X_0$ be as in \Ref{XII.3.5}. Assume that $X$ has depth $\geq 2$ at its closed points, and that
\[
\H^i(X_0, \Oo_{X_0}(-n))=0
\]
for $n>0$ and for $i=1$ (\resp for $i=1$ and $i=2$), which, by \Ref{XII.1.4}, implies that $X_0$ has depth $\geq 2$ (\resp $\geq 3$) at its closed points, \ie that $X$ has depth $\geq 3$ (\resp $\geq 4$) at the closed points of $X_0$. Under these conditions, for every open neighborhood $U$ of $X_0$, $\Pic(U)\to\Pic(X_0)$ is injective; in particular, $\Pic(X)\to\Pic(X_0)$ is injective (\resp $\varprojlim_U \Pic(U)\to\Pic(X_0)$ is bijective). In the parenthesized case, if one assumes in addition that the local ring of $X$ at every closed point not in $X_0$ is parafactorial (\Ref{XI.3.1}) (for example, regular, or more generally a complete intersection), then $\Pic(X)\to\Pic(X_0)$ is bijective.
\end{corollaire}
Apply \Ref{XI}~\Ref{XI.3.12} and \Ref{XI.3.13}, noting that the parenthesized hypothesis implies that the irreducible components of $X$ having dimension $\neq 0$ have dimension $\geq 4$. In particular, applying this when $X$ is a global complete intersection of dimension $\geq 4$ in projective space, one obtains:

\begin{corollaire} \label{XII.3.7} Let $X$ be an algebraic scheme of dimension $\geq 3$ that is a complete intersection in the scheme $\PP_k^r$. Then $\Pic(X)$ is the free group generated by the class of the sheaf $\OX(1)$.
\end{corollaire}
Proceed by induction on the number of hypersurfaces whose intersection is $X$, applying \Ref{XII.3.6} and noting that, for a complete intersection $X$ of dimension $\geq 3$, one has $\H^i(X, \OX(n))=0$ for $i=1,2$ and every $n$.

\begin{remarque} \label{XII.3.8}
In\pageoriginale the case where $X$ is a nonsingular hypersurface, \Ref{XII.3.7} is due to Andreotti. The result \Ref{XII.3.7} can also be expressed (when $X$ is nonsingular) by saying that the homogeneous coordinate ring of $X$ is factorial, and in this form it is contained in \Ref{XI}~\Ref{XI.3.13} (ii). We also note that Serre had given a proof of \Ref{XII.3.7} in the nonsingular case by transcendental methods, using a specialization argument to reduce to characteristic 0, where the Lefschetz theorem is available in its classical form. Of course, the fact that the purely algebraic proof given here makes it possible to dispense with nonsingularity hypotheses in the statement of the Lefschetz theorem invites one to reconsider it in the classical case as well. \Cf the following expos\'e, which proposes conjectures in this direction.

In Corollaries \Ref{XII.3.5} and \Ref{XII.3.7}, we have worked over a base field, whereas the key theorems \Ref{XII.2.4} and \Ref{XII.3.4} are valid over an arbitrary base. To generalize Corollaries \Ref{XII.3.5} and \Ref{XII.3.6} to a general $S$, we must give workable criteria for a point of $X$ (flat over $S$) to have a \og pure\fg \resp parafactorial local ring. This will be the subject of the following n\textsuperscript{o}.
\end{remarque}

\section{Formal completion and normal flatness} \label{XII.4}

\setcounter{equation}{18}

\begin{theoreme} \label{XII.4.1} Let $X$ be a locally Noetherian prescheme, locally embeddable in a regular scheme, $Y$ a closed subset of $X$, $U=X- Y$, $X_0$ a closed subprescheme of $X$ defined by an ideal $\Jj$, $\widehat{X}$ the formal completion of $X$ along $X_0$, $U_0$ the trace of $X_0$ on $U$, $\widehat{U}$ the formal completion of $U$ along $U_0$, $i:U\to X$ and $\widehat{i}:\widehat{U}\to \widehat{X}$ the canonical immersions, and $n$ an integer. Assume that
\begin{enumeratea}
\item
$X$ is normally flat along $X_0$ at the points of $Y\cap X_0$, \ie at these points the \sisi{m}{M}odules $\Jj^n/\Jj^{n+1}$ over $X_0$ are flat, \ie \sisi{\ignorespaces}{locally} free.
\item
For every $x\in Y\cap X_0$, one has $\prof \Oo_{X_0, x}\geq n+2$.
\end{enumeratea}
Under\pageoriginale these conditions, the following assertions hold:
\begin{enumerate}
\sisi{\item[$1^\circ)$]}{\item} Let $F$ be a coherent Module on $U$, and suppose that:
\begin{enumerate}
\item[\textup{c)}] For every $x\in Y- Y\cap X_0$, one has $\prof \Oo_{X, x}\geq n+2$.

\item[\textup{d)}] $F$ is free at the points of $U_0$, and has depth $\geq n+1$ at all the points of $U$ where it is not free.
\end{enumerate}
Then the graded Module
\[
\sisi{\coprod}{\tbigoplus}_{m\geq 0}\R^pi_\ast(\Jj^mF)
\]
over $\sisi{\coprod}{\bigoplus}_{m\geq 0}\Jj^m$ is of finite type for $p\leq n$. \sisi{\item[$2^\circ$]}{\item} Let $\formelF$ be a coherent \sisi{}{locally free }Module on $\widehat{U}$; then the graded Module
\[
\sisi{\coprod}{\tbigoplus}_{m\geq 0}\R^pi_\ast(\Jj^m\formelF/\Jj^{m+1}\formelF)
\]
over $\sisi{\coprod}{\bigoplus}_{m\geq 0}\Jj^m/\Jj^{m+1}=\gr_{\Jj}(\OX)$ is of finite type for $p\leq n$.
\end{enumerate}
\end{theoreme}

\begin{proof}
1) Let $X'\!=\!\Spec(\sisi{\coprod}{\bigoplus}_{m\geq 0}\Jj^m)$; the base change \hbox{$f:X'\!\to\! X$} then defines $U'=X'- Y'$, $X_0'$, $U_0'=X_0'\cap U'$, and immersions $i':U'\to X'$, \hbox{$i_0':U_0'\to X_0'$}. Thus we have a Cartesian square
\[
\xymatrix{
X'\ar[d]_f&\ar[l]_-{i'} U'\ar[d]^g \\
X&\ar[l]_-{i} U
}
\]
and we have
\steco
\begin{equation} \label{eq:XII.19}
\sisi{\coprod}{\tbigoplus}_{m\geq 0}\R^pi_\ast(\Jj^mF) = \R^pi_\ast \Big(\sisi{\coprod}{\tbigoplus}_{m\geq 0} \Jj^mF \Big) =\R^pi_\ast (g_\ast(F')),
\end{equation}
where\pageoriginale $F'=g^\ast F$, so that there is indeed a canonical isomorphism
\[
g_\ast(F') \isomto \sisi{\coprod}{\tbigoplus}_{m\geq 0}\Jj^mF,
\]
since this is true at the points of $U_0$, because $F$ is free there by d), and also at the points of $\sisi{U_0}{U-U_0}$, because there one has $\Jj^m=\Oo_U$, (so that in both cases, $\Jj^m\otimes_{\Oo_U}F\to \Jj^mF$ is an isomorphism).

On the other hand, since $f$, and hence $g$, are affine, one has
\steco
\begin{equation} \label{eq:XII.20} {}\R^pi_\ast(g_\ast(F')) = \R^p(ig)_\ast(F') = \R^p(fi')_\ast(F') = f_\ast(R^pi'_\ast(F'))
\end{equation}
and therefore, comparing (\Ref{eq:XII.19}) and (\Ref{eq:XII.20}), we see that \hyperref[XII.4.1]{assertion 1}\sisi{${}^\circ$}{} is equivalent to the following: $R^pi'_\ast(F')$ is a Module of finite type, \ie coherent, on $X'$, for every $p\leq n$. Now, since $X$ is locally embeddable in a regular scheme, the same is true of $X'$, which is of finite type over $X$, and we may apply the coherence criterion \sisi{\ignorespaces}{\Ref{VIII}}~\Ref{VIII.2.3} to a coherent extension $F^{''}$ of $F'$: we wish to express that $\SheafH^p\sisi{_Y}{_{Y'}}(F^{''})$ is coherent for $p\leq n+1$, and this is also equivalent to saying that, for every $x'\in U'$ such that
\begin{equation*} \label{eq:XII.20bis} \tag{20 bis} {}\codim (\overline{x'}\cap Y', \overline{x'})=1,
\end{equation*}
one has
\steco
\begin{equation} \label{eq:XII.21} {} \prof F'_{x'}\geq n+1.
\end{equation}
Now this condition is satisfied at the points $x'$ where $F'$ is not free, because for such an $x'$ one has $x'\notin U_0'$ by d), and therefore $g$ is an isomorphism there; again by d), $F$ has depth $\geq n+1$ at $g(x')$, and hence $F'$ has depth $\geq n+1$ at $x'$. It is therefore enough to verify condition (\Ref{eq:XII.21}) at those $x'\in U'$ satisfying \sisi{(\Ref{eq:XII.21})}{(\Ref{eq:XII.20bis})} at which $F'$ is free. For this, it is enough to prove that one has\pageoriginale
\begin{equation*} \label{eq:XII.21bis} \tag{21 bis}
\prof \Oo_{X', x'}\geq n+1
\end{equation*}
\enlargethispage{\baselineskip}%
at these points; \afortiori it is enough to establish this relation at all the points~$x$ of $U'$ satisfying (\Ref{eq:XII.20}\sisi{}{ bis}). Now, \sisi{by the criterion cited}{again by criterion~\Ref{VIII.2.3}} in Expos\'e \Ref{VIII}, this is equivalent to the assertion that the Modules
\[
\SheafH^p_{Y'}(\Oo_{X'}) \hspace{5mm} \text{for} \; p\leq n+1
\]
are coherent. In fact, we shall prove that they are even zero, or, equivalently by Expos\'e~\Ref{III}, that one has
\steco
\begin{equation} \label{eq:XII.22} {} \prof \Oo_{X', x'}\geq n+2 \hspace{5mm} \text{for every }\; x'\in Y'.
\end{equation}
For this, we distinguish two cases. If $x'\notin X'_0$, then $f$ is an immersion at $x'$, and we must verify that $\sisi{F}{\Oo_X}$ has depth \sisi{at}{$\geq n+2$ at } the image $x=f(x')$, which is precisely condition c). If, on the other hand, $x'\in X'_0$, \ie $x=f(x')\in X_0$ and hence $x\in Y\cap X_0$, we apply conditions a) and b) by means of the

\begin{lemme} \label{XII.4.2}
Let $X$ be a locally Noetherian prescheme, $X_0$ a closed subprescheme of $X$ defined by an ideal $\Jj$, $X'=\Spec(\sisi{\coprod}{\bigoplus}_{m\geq 0} \Jj^m)$, $X_0'=\Spec(\sisi{\coprod}{\bigoplus}_{m\geq 0} \Jj^m/\Jj^{m+1}) = X'\times_XX_0$, and $x$ a point of $X_0$ at which $X$ is normally flat along $X_0$, \ie such that $\gr_{\Jj}(\OX)= \sisi{\coprod}{\bigoplus}_{m\geq 0}\Jj^m/\Jj^{m+1}$ is flat there as a Module over $X_0$. Then, for every sequence of elements $f_i$ ($1\leq i\leq m$) of $\Oo_{X, x}$ whose images in $\Oo_{X_0, x}$ form an $\Oo_{X_0, x}$-regular sequence, and every $x'\in X'$ above $x$, the images of the $f_i$ in $\Oo_{X', x'}$ (respectively in $\Oo_{X'_0, x'}$) likewise form an $\Oo_{X', x'}$-regular sequence (respectively an $\Oo_{X'_0, x'}$-regular sequence); in particular, one has
\steco
\begin{equation} \label{eq:XII.23}
\prof \Oo_{X', x'} \geq \prof \Oo_{X_0, x}, \quad \prof \Oo_{X'_0, x'} \geq \prof \Oo_{X_0, x}.
\end{equation}
\end{lemme}

To\pageoriginale prove it, we may suppose that $X$ is local with closed point $x$, and hence affine with ring $A=\Oo_{X, x}$, with $\Jj$ defined by the Ideal $J$; it is enough to prove that, for every sequence $f_i$ ($1\leq i\leq m$) of elements of $A$ whose images in $A/J$ form an $A/J$-regular sequence, the $f_i$ likewise form a $(\sisi{\coprod}{\bigoplus}_{m\geq 0} J^m)$-regular sequence and a $(\sisi{\coprod}{\bigoplus}_{m\geq 0} J^m/J^{m+1})$-regular sequence, \ie for every $m$, it forms a $J^m$-regular sequence and a $J^m/J^{m+1}$-regular sequence. The second assertion is trivial, since $J^m/J^{m+1}$ is a free module over $A/J$. The first follows by considering the $J$-adic filtration of $J^m$ and noting that the sequence of the~$f_i$ is regular on the graded module associated with $J^m$ for that filtration.

This proves \Ref{XII.4.2}, and consequently \Ref{XII.4.1}~1\sisi{\noo}{}).

Let us prove \Ref{XII.4.1} 2\sisi{\noo}{}). For this, use the Cartesian square
\[
\xymatrix{
X'_0\ar[d]_{f_0}&\ar[l]_-{i'_0} U'_0\ar[d]^{g_0} \\
X_0&\ar[l]_-{i_0} U_0
}
\]
and, proceeding as at the beginning of the proof of 1\sisi{\noo}{}), we find that
\steco
\begin{equation} \label{eq:XII.24}
\sisi{\coprod}{\tbigoplus}_{m\geq 0} \R^pi_\ast(\Jj^m\formelF/\Jj^{m+1}\formelF) \simeq {f_0}_\ast (\R^p{\sisi{i}{i'}_0}_\ast(\formelF_0')),
\end{equation}
where $\formelF_0=\formelF/\Jj\formelF$ and $\formelF_0'=g_0^\ast(\formelF_0)$, (using the fact that $\formelF$ is locally free). Thus the conclusion of 2\sisi{\noo}{}) is equivalent to saying that, for $p\leq n$, $\R^p{i'_0}_\ast(\formelF_0')$ is a coherent Module.

\enlargethispage{\baselineskip}%
Here again, taking into account that $\formelF_0'$ is locally free, criterion \Ref{VIII}~\Ref{VIII.2.3} allows us to reduce to proving the same assertion after replacing $\formelF_0'$ by $\Oo_{U'_0}$, \ie to proving that the Modules\pageoriginale
\[
\SheafH^p_{Y_0'}(\Oo_{X_0'}) \hspace{5mm} \text{for} \; p\leq n+1 \hspace{5mm} \text{(where} \; Y'_0=Y'\cap X_0'=X_0'- U'_0 \text{)}
\]
are coherent. Once again, we prove that they are in fact zero, i.e., that one has
\steco
\begin{equation} \label{eq:XII.25} {}\prof \Oo_{X'_0, x'} \geq n+2 \hspace{5mm} \text{for every} \; x'\in Y'_0.
\end{equation}
This does indeed follow from conditions a) and b), in view of \Ref{XII.4.2}. This completes the proof of \Ref{XII.4.1}.
\skipqed
\end{proof}

\begin{remarque} \label{XII.4.3}
It follows immediately, by descent, that the hypothesis that $X$ be locally embeddable in a regular scheme can be replaced by the following weaker hypothesis: there exists a faithfully flat and quasi-compact morphism $\overline{X}\to X$ such that $\overline{X}$ is locally embeddable in a regular scheme.
\end{remarque}

Theorem \Ref{XII.4.1} enables us to apply the results of Expos\'e \Ref{IX} (the comparison and existence theorems). We shall be interested in particular in the

\begin{corollaire} \label{XII.4.4}
Assume that conditions a), b), and c) of Theorem \Ref{XII.4.1} hold, with $n=1$ and $X=\Spec(A)$, where $A$ is separated and complete for the $J$-adic topology. Then
\begin{enumerate}
\sisi{\item[\textup{1\noo})]}{\item} The functor $F\mto \widehat{F}$ from the category of locally free coherent modules on~$U$ to the category of locally free coherent modules on $\widehat{U}$ is fully faithful. \sisi{\item[\textup{2\noo})]}{\item} For every locally free coherent module $\formelF$ on $\widehat{U}$, there exists a coherent module $F$ on $U$ and an isomorphism $\widehat{F}\isomto \formelF$.
\end{enumerate}
In particular, if, for every $x\in U$ whose closure in $U$ does not meet $U_0$, \ie such that $\overline{x}\cap X_0\subset Y$, one has $\prof \Oo_{U, x}\geq 2$, then the pair $(U, U_0)$ satisfies the effective Lefschetz condition ($\Leff$) of Expos\'e \Ref{X}.
\end{corollaire}
\noindent
(For the last assertion, one proceeds as in \Ref{X}~\Ref{X.2.1}.)

\pagebreak[2]
Special case of \Ref{XII.4.4}:

\begin{corollaire} \label{XII.4.5}
Let\pageoriginaled $A$ be a Noetherian ring, let $J$ be an ideal of $A$ contained in the radical, and let $A_0=A/J$. Assume
\begin{enumeratei}
\item
$\prof A_0\geq 3$.
\item
$\gr_J(A)$ is a free $A_0$-module.
\item
$A$ is complete for the $J$-adic topology.
\end{enumeratei}

Let $X=\Spec(A)$, $X_0=\Spec(A_0)=V(J)$, let $a$ be the closed point of $X$, let $U=X-\{a\}$, $U_0=X_0-\{a\}$, and let $\widehat{U}$ be the formal completion of $U$ along $U_0$. Then the functor $F\mto \widehat{F}$ from the category of locally free coherent modules on $U$ to the category of locally free coherent modules on $\widehat{U}$ is fully faithful. Moreover, for every locally free coherent module $\formelF$ on $\widehat{U}$, there exists a coherent module (not necessarily locally free!) $F$ on $U$, and an isomorphism $\widehat{F}\simeq\formelF$.
\end{corollaire}

Note that, thanks to \Ref{XII.4.3}, we did not have to assume that $A$ is a quotient of a regular ring, because the completion of $A$ for the $\rr(A)$-adic topology satisfies this condition in any case.

Proceeding as in Expos\'es \Ref{X} and \Ref{XI}, we deduce from \Ref{XII.4.5}

\begin{corollaire} \label{XII.4.6}
Under the conditions of \Ref{XII.4.5}, one has the following:
\begin{enumeratea}
\item
$U$ and $U_0$ are connected (\Ref{III}~\Ref{III.3.1}).

Choosing a geometric base point in $U_0$, the homomorphism
\[
\pi_1(U_0)\to \pi_1(U)
\]
is surjective.
\item
The homomorphism
\[
\Pic(U)\to \Pic(U_0)
\]
is injective.
\end{enumeratea}
\end{corollaire}

To\pageoriginale prove b), in view of \Ref{XII.4.5}, this amounts to checking that every isomorphism $L_0'\isomto L_0$ lifts to an isomorphism $\widehat{L'}\isomto \widehat{L}$. For this, one lifts step by step to isomorphisms $L_n'\isomto L_n$; the obstructions lie in $\H^1(U_0, \Jj^n/\Jj^{n+1})$, and these modules are zero because $J^n/J^{n+1}$ is free and $\prof A_0\geq 3$.

We are now in a position to prove the
\enlargethispage{-1\baselineskip}%
\begin{theoreme} \label{XII.4.7}
Let $A$ be a noetherian local ring, let $J$ be an ideal of $A$ contained in its radical, and let $A_0=A/J$. Suppose that
\begin{enumeratei}
\item
$\prof A_0\geq 3$.
\item
$\gr_J(A)$ is a free module over $A_0$.
\end{enumeratei}
\noindent
Then, if $A_0$ is \og pure\fg (\Ref{X}~\Ref{X.3.1}) (\resp parafactorial (\Ref{XI}~\Ref{XI.3.1})), the same holds for~$A$.
\end{theoreme}

\begin{proof}
By descent, we may suppose that we also have
\begin{enumeratei}\setcounter{enumi}{2}
\item
$A$ is complete for the $J$-adic topology.
\end{enumeratei}

Indeed, by (i) and (ii), we have $\prof(A)\geq 3$, and hence $\prof(\widehat{A})\geq 3$, where $\widehat{A}$ is the completion of $A$ for the $J$-adic topology, and we apply \Ref{X}~\Ref{X.3.6} and \Ref{XI}~\Ref{XI.3.6}. We are therefore under the conditions of \Ref{XII.4.5}. Since $\prof(A)\geq 3\geq 2$, to say that $A$ is parafactorial means simply that $\Pic(U)=0$, and by \Ref{XII.4.6}~b), it is enough for this that $\Pic(U_0)=0$, \ie that $A_0$ be parafactorial. To prove that $A$ is \og pure\fg if $A_0$ is, we must prove that if $V$ is an \'etale covering of $U$, defined by an algebra $\Bb$ on~$U$, then $\H^0(U, \Bb)$ is a finite \'etale algebra over $A$. Now, since $A_0$ is pure, the same holds for the $A_n$ (which differ from it only by nilpotent elements), and therefore, for every $n$, $\sisi{}{B_n=}\H^0(U, \Bb_n)$ is an \'etale algebra over $A/J^{n+1}$; these algebras of course fit together, so that $\varprojlim B_n$ is an \'etale algebra over $A$. By \Ref{XII.4.5}, this algebra is none other than $\H^0(U, \Bb)$, which proves our assertion.
\skipqed
\end{proof}

\begin{corollaire} \label{XII.4.8}
Let\pageoriginale $f:X\to Y$ be a flat morphism of locally noetherian preschemes, let $x\in X$, and let $y=f(x)$. Suppose that $\Oo_{X_y, x}$ is a \og pure\fg (\resp parafactorial) local ring of depth $\geq 3$. Then the same holds for $\Oo_{X, x}$.
\end{corollaire}

This is the result of the kind promised at the end of the preceding \sisi{N\textsuperscript{o}}{\numero}, for generalizing Corollaries \Ref{XII.3.5} and following. Thus, using \Ref{XII.3.4}, we obtain the

\begin{corollaire} \label{XII.4.9}
Let $f:X\to S$ be a projective flat morphism, with $S$ locally Noetherian, let $\OX(1)$ be an invertible module on $X$, ample relative to $S$, let $t$ be a section of $\OX(1)$ such that, for every $s\in S$, the section $t_s$ induced on $X_s$ is $\Oo_{X_s}$-regular, let $X_0$ be the closed subscheme of zeros of $t$, and let $X_m$ be the closed sub-\sisi{pre}{}scheme of zeros of $t^{m+1}$. Assume that, for every $s\in S$, $X_s$ has depth $\geq 3$ at all its closed points. Then:
\begin{enumeratea}
\item
If the local rings at the closed points of $X_s-X_{0,s}$ ($s\in S$) are \og pure\fg\,---for example, if they are complete intersections---then the functor $X'\mto X_0'=X'\times_X X_0$ from the category of finite \`etale coverings of $X$ to the category of finite \`etale coverings of $X_0$ is an equivalence of categories; in particular, after choosing a geometric base point in $X_0$, the homomorphism
\[
\pi_1(X_0)\to \pi_1(X)
\]
is an isomorphism\refstepcounter{toto}\nde{\label{FH}Let us point out
the striking connectedness result obtained subsequently by Fulton and
Hansen, in the case where $S=\Spec(k)$ ($k$ an algebraically closed
field). Let $g:X\to {\PP^m_k}\times{\PP^m_k}$ be such that $\dim
g(X)>m$; then the inverse image of the diagonal is connected. Among
other things, this makes it possible to generalize
Corollary~\Ref{XII.4.9} when $f$ is the structural morphism
of projective space $\PP^m_k$ over $S=\Spec(k)$: more precisely, an
irreducible subvariety $X$ of ${\PP^m_k}$ of dimension $>m/2$ has
trivial fundamental group! (\cf Fulton~W. \& Hansen J., {\og
A connectedness theorem for projective varieties, with
applications to intersections and singularities of mappings\fg},
\emph{Ann. of Math. (2)} \textbf{110} (1979), \numero 1, p\ptbl
159-166). For generalizations to Grassmannians or
abelian varieties, see {Debarre~O.}, {\og Th\'eor\`emes de
connexit\'e pour les produits d'espaces projectifs et les
grassmanniennes\fg}, \emph{Amer. J.~Math.} \textbf{118} (1996),
\numero 6, p\ptbl 1347-1367 and {\og Th\'eor\`emes de connexit\'e et
vari\'et\'es ab\'eliennes\fg}, \emph{Amer. J.~Math.} \textbf{117}
(1995), \numero 3, p\ptbl 787--805. The triviality result for the
fundamental group of $X$ above was obtained independently
by Faltings, who moreover proves that
$\text{Pic}(X)$ has no torsion prime to the characteristic
of $k$, using methods of algebraization of formal vector bundles,
more in the spirit of Grothendieck's techniques, \cf
(Faltings~G., \og{Algebraization of some formal vector bundles\fg}
\emph{Ann. of Math. (2)} \textbf{110} (1979), \numero 3, p\ptbl
501--514).}

\item
If the local rings at the closed points of $X_s-X_{0,s}$ ($s\in S$) are \og parafactorial\fg---for example, regular, or complete intersections of dimension $\geq 4$---then, for every integer $m$ such that $\R^if_{0\ast}(\Oo_{X_0}(-n))=0$ for $n>m$ and $i=1,2$, the map $\Pic(X)\to \Pic(X_m)$ is bijective.

Moreover, if $S$ is Noetherian and the $X_{0,s}$ have depth $\geq 3$ at their closed points, such $m$ exist (\cf \Ref{XII.1.5}).
\end{enumeratea}
\end{corollaire}

\begin{remarque} \label{XII.4.10}
Under\pageoriginale the conditions of the last assertion of \Ref{XII.4.9} b), we saw in \Ref{XII.1.5} that there exists an $m$ such that $n>m$ even implies $\H^i(X_0, \Oo_{X_0}(-n))=0$ for $i=1,2$, and indeed for $i\leq 2$\sisi{)}{}. This condition is stronger than $\R^if_{\ast}(\Oo_{X_0}(-n))=0$ for $i=1,2$, and it has the further advantage of being stable under base change. The same is true of the depth hypotheses made in \Ref{XII.4.9}, and also of a hypothesis of the type \og the $X_s$ are locally complete intersections\fg. It follows under these conditions that \Ref{XII.4.9} b) also implies that the morphism of functors
\[
\SheafPic_{X/S}\to \SheafPic_{X_n/S}
\]
in $\catSch_{/S}$ is an isomorphism, and hence so is the morphism of relative Picard schemes, when these exist:
\[
\Pic_{X/S}\to \Pic_{X_n/S}.
\]
Even when $S$ is the spectrum of an algebraically closed field, this statement is considerably more precise than the statement that merely says that $\Pic(X)\to \Pic(X_n)$ is bijective.

One may ask whether one can always take $n=0$ in the preceding conclusions (thus assuming that the $X_{0,s}$ have depth $\geq 3$ at their closed points). When $X_0$ is smooth over $S$ and the residual characteristics of $S$ are zero, this is indeed so, by the Kodaira \og vanishing theorem\fg (proved by transcendental methods, using a K\"ahler metric), which implies that, for every connected smooth projective scheme of dimension $n$ over a field $k$ of characteristic zero and every ample invertible module $L$ on $X$, one has $\H^i(X,L^{-1})=0$ for $i\neq n$. It is not known\nde{As Raynaud observed, the decomposition result for the de Rham complex of Deligne and Illusie readily implies the vanishing of the group $H^i(X,L^{-1})$ (with $L$ ample on $X$, projective and smooth over $k$ of characteristic $p>0$) for $i<\inf(p,\dim(X))$, provided that $X$ can be lifted to a flat scheme over $W_2(k)$ (\cf Deligne~P. \& Illusie L., {\og Rel\`evements modulo~$p^2$ et d\'ecomposition du complexe de de Rham\fg}, \emph{Invent. Math.} \textbf{89} (1987), \numero 2, p\ptbl 247--270); this gives a purely algebraic proof of Kodaira's result for projective varieties in characteristic zero. If $X$ cannot be lifted, it is well known that the \og vanishing theorem\fg is false; \cf the example in (Raynaud~M., {\og Contre-exemple au "vanishing theorem" en caract\'eristique $p>0$\fg}, in \emph{C.P.~Ramanujam---a tribute}, Tata Inst. Fund. Res. Studies in Math., vol.~8, Springer, Berlin-New York, 1978, p\ptbl 273--278); see also the very elegant examples in (Haboush~W. \& Lauritzen~N., {\og Varieties of unseparated flags\fg}, in \emph{Linear algebraic groups and their representations (Los Angeles, CA, 1992)}, Contemp. Math., vol.~153, American Mathematical Society, Providence, RI, 1993, p\ptbl 35--57), simplified in (Lauritzen~N. \& Rao~A.P., {\og Elementary counterexamples to Kodaira vanishing in prime characteristic\fg}, \emph{Proc. Indian Acad. Sci. Math. Sci.} \textbf{107} (1997), \numero 1, p\ptbl 21--25). On the other hand, I know of no example in which the map $\Pic(X_{n+1})\to\Pic(X_n)$ is not surjective for $n>1$ in positive characteristic, where $X_n$ denotes a thickened hyperplane section of a smooth projective $X$ as above.} at present whether this theorem can be replaced by a generalization in characteristic $p>0$, and whether the smoothness hypothesis can be replaced by a hypothesis of a more general nature (involving depth, or of \og complete intersection\fg type, and so on).
\end{remarque}

\section[Universal finiteness conditions]
{Universal finiteness conditions for a nonproper morphism} \label{XII.5} Recall\pageoriginale for reference the \setcounter{equation}{25}

\begin{proposition} \label{XII.5.1}
Let $f:X\to S$ be a proper morphism of preschemes, with $S$ locally noetherian, let $U$ be an open subset of $X$, let $g:U\to X$ be the canonical immersion, let $h=fg:U\to S$, and let $F$ be a Module on $U$. Suppose that the Modules $\R^ig_\ast(F)$ are coherent for $i\leq n$ (a hypothesis local in nature on $X$, which in practice is checked using the criterion \Ref{VIII}~\sisi{\Ref{VII.2.3}}{\Ref{VIII.2.3}}). Then $\R^ih_\ast(F)$ is coherent for $i\leq n$.
\end{proposition}

This follows at once from the Leray spectral sequence
\steco
\begin{equation} \label{eq:XII.26} {} \E_2^{p, q}=\R^pf_\ast(\R^qg_\ast(F)) \To \R^\ast h_\ast(F),
\end{equation}
and from the fact that the higher direct images under $f$ of a coherent Module on $X$ are coherent (\EGA III 3.2.1).

\begin{proposition} \label{XII.5.2}
Let $S$ be a locally noetherian prescheme, let $\Ss$ be a quasi-coherent graded Algebra of finite type over $S$, generated by $\Ss_1$, let $X$ be a subprescheme of $\Proj(\Ss)$, let $\OX(1)$ be the invertible Module on $X$ induced by $\Proj(\Ss(1))$ and very ample relative to $S$, let $U$ be an open subset of $X$, let $g:U\to X$ be the canonical immersion, let $h=fg:U\to S$, and let $F$ be a quasi-coherent Module on $U$, giving twisted Modules $F(m)=F\otimes \OX(m)$\sisi{)}{} ($m\in\ZZ$); let $n$ be an integer and $m_0$ an integer. The following conditions are equivalent:
\begin{enumeratei}
\item
$\R^ig_\ast(F)$ is coherent for $i\leq n$.
\item
$\sisi{\coprod}{\bigoplus}_{m>m_0}\R^ih_\ast(F(m))$ is a finite-type $\Ss$-Module for $i\leq n$.
\end{enumeratei}
\end{proposition}

\begin{proof}
Replacing $F$ by $F(m)$ in the spectral sequence above, we
obtain a spectral sequence of graded $\Ss$-Modules\pageoriginale
\[
\E_2^{p, q}= \sisi{\coprod}{\tbigoplus}_{m\geq m_0} \R^pf_\ast(\R^qg_\ast(F(m))\sisi{}{)} \To \sisi{\coprod}{\tbigoplus}_{m\geq m_0} \R^\ast h_\ast(F(m)).
\]
Since we have
\[
\R^qg_\ast(F(m)) \simeq \R^qg_\ast(F) (m),
\]
we see that if the $\R^ig_\ast(F)$ are coherent, then $\E_2^{p, q}$ is of finite type over $\Ss$ for $q\leq n$, by part~a) of Lemma~\Ref{XII.5.3} below, which implies that the abutment is of finite type over $\Ss$ in degree $i\leq n$. This proves (i) \ALORS (ii). Moreover, reasoning in the abelian category of graded $\Ss$-Modules modulo the thick subcategory $\mathcal{C}$ of those that are quasi-coherent of finite type, we obtain from the preceding spectral sequence
\[
\sisi{\coprod}{\tbigoplus}_{m\geq m_0} \R^{n+1} h_\ast(F(m)) \simeq \sisi{\coprod}{\tbigoplus}_{m\geq m_0} f_\ast (\R^{n+1} g_\ast(F)(m)) \mod \mathcal{C},
\]
which proves that if the left-hand side is a finite-type $\Ss$-Module, then $\R^{n+1} g_\ast(F)$ is coherent, by part~b) of Lemma~\Ref{XII.5.3}. This proves the implication (ii) \ALORS (i) by induction on $n$. It remains to prove:

\begin{lemme} \label{XII.5.3} Let $S$, $\Ss$, $X$, $f$ be as in \Ref{XII.5.2}, let $G$ be a quasi-coherent Module on $X$, and let $m_0$ be an integer. Then
\begin{enumeratea}
\item
If $G$ is coherent, then for every integer $i$, the graded Module
\[
\sisi{\coprod}{\tbigoplus}_{m\geq m_0} \R^i f_\ast(G(m))
\]
over $\Ss$ is of finite type.
\item
Conversely, suppose that the Module $\sisi{\coprod}{\bigoplus}_{m\geq m_0} \sisi{\R^i }{}f_\ast(G(m))$ over $\Ss$ is of finite type; then $G$ is coherent.
\end{enumeratea}
\end{lemme}

\begin{proof}[Proof of \Ref{XII.5.3}]
For a), the case $i=0$ is given in \EGA III 2.3.2, and the case $i>0$ in \EGA III 2.2.1 (i)(ii), which says that the $ \R^i f_\ast G(m)$ are coherent\pageoriginale and zero for large $m$ (if we suppose $S$ noetherian, as we may). For b), note that $G$ is isomorphic to $\SheafProj (\sisi{\coprod}{\bigoplus}_{m\geq m_0}f_\ast(G(m))\sisi{}{)}$ (\EGA II 3.4.4 and 3.4.2), which proves that $G$ is coherent if $\sisi{\coprod}{\bigoplus}_{m\geq m_0}f_\ast(G(m))$ is of finite type over $\Ss$, by \loccit 3.4.4.
\skipqed
\end{proof}
\skipqed
\end{proof}

\begin{corollaire}[{{\normalfont$S$ noetherian}}] \label{XII.5.4}
Suppose that $\R^ig_\ast(F)$ is coherent for $i\leq n$. Then, for $i\leq n+1$ and large $m$, there is a canonical isomorphism:
\[
\R^ih_\ast(F(m)) \simeq f_\ast(\R^ig_\ast(\sisi{G}{F})(m)).
\]
\end{corollaire}

Indeed, the spectral sequence (\Ref{eq:XII.26}) for $F(m)$ then degenerates in degree $\leq \sisi{n}{n+1}$, by \EGA III 2.2.1~(ii), which gives the result at once (and in fact recovers the implication (ii) \ALORS (i) of~\hyperref[XII.5.2]{5.2}).

\enlargethispage{1.2\baselineskip}%
\begin{corollaire} \label{XII.5.5}
Under the preliminary conditions of \Ref{XII.5.2}, with $S$ noetherian, the following conditions are equivalent:
\begin{enumeratei}
\item
$\sisi{\coprod}{\bigoplus}_{m\geq m_0}h_\ast(F(m)\sisi{}{)}$ is of finite type over $\sisi{S}{\Ss}$, and $\R^ih_\ast(F(m))=0$ for $0<i\leq n$ and large $m$.
\item
$g_\ast(F)$ is coherent, and $\R^ig_\ast(F)=0$ for $0<i\leq n$.
\item[\textup{(ii bis)}] $g_\ast(F)$ is coherent, and $\prof_Y g_\ast(F)>n+1$.
\end{enumeratei}
\end{corollaire}

The equivalence of (ii) and (ii bis) is contained in \Ref{III}~\Ref{III.3.3}. Moreover, by \Ref{XII.5.2}, conditions (i) and (ii) both imply that the $\R^ig_\ast(F)$ ($i\leq n$) are coherent. The equivalence of (i) and (ii) then follows from \Ref{XII.5.4}, taking into account the fact that for every coherent Module $G$ on $X$, one has $G=0$ if and only if $f_\ast(G(m))=0$ for large $m$, for example by \EGA III 2.2.1 (iii).

\begin{remarque} \label{XII.5.6}
The\pageoriginale criteria \Ref{XII.5.2} and \Ref{XII.5.5} can be interpreted by saying that the \og simultaneous finiteness condition\fg \Ref{XII.5.2}~(ii) is expressed by local regularity properties (in terms of depth by \Ref{VIII}~\Ref{VIII.2.1}) of $F$ at the points of $U$ near $Y=X- U$, whereas the \og asymptotic vanishing condition\fg \Ref{XII.5.5}~(i) is distinctly stronger in nature and is expressed by local regularity conditions on $g_\ast(F)$ at the points of~$Y$ itself. It would be interesting, in order to generalize Lefschetz-type theorems for projective morphisms to quasi-projective morphisms, to find local criteria on $X$ that are necessary and sufficient for the $\Ss$-Modules $\sisi{\coprod}{\bigoplus}_{m\geq 0} \R^ih_\ast(F(m))$ for $i\leq n$ to be of finite type. When $S$ is the spectrum of a field (and doubtless more generally if it is the spectrum of an artinian ring) and $Y=X- U$ is finite, one can show that it is necessary and sufficient that the following conditions hold:
\begin{enumerate}
\item[\textup{1\noo)}] $\prof F_x >n$ for every closed point $x$ of $U$ (compare \Ref{XII.1.4}).
\item[\textup{2\noo)}] $\R^ig_\ast(F)$ is coherent for $i\leq n$, or equivalently, there exists an open neighborhood $V$ of $Y$ such that for every closed point $x$ of $U\cap V$, one has $\prof F_x >n+1$.
\end{enumerate}
\end{remarque}

\begin{proposition} \label{XII.5.7}
Let $S$ be a locally Noetherian prescheme, let $g:U\to X$ be a morphism of preschemes of finite type over $S$\sfootnote{It is in fact enough that $g$ be quasi-compact and quasi-separated (\textup{\EGA IV1.2.1}), with no condition on $U,X$.}, with structural morphisms $h$ and $f$, let $F$ be a quasi-coherent module on $U$, and let $n$ be an integer. The following conditions are equivalent:
\begin{enumeratei}
\item
For every base change $S'\to S$, with $S'$ Noetherian, the module $\R^ng'_\ast(F')$ on $X'$ is coherent.
\item
For every base change as above and every coherent ideal $\Jj$ on~$S'$, denoting by $\Ii$ the ideal $\Jj\Oo_{X'}$ on $X'$, the graded module
\[
\sisi{\coprod}{\tbigoplus}_{m\geq 0} \R^ng'_\ast(\Ii^mF')
\]
over $\sisi{\coprod}{\bigoplus}_{m\geq 0} \Ii^m$ is of finite type.
\item
For\pageoriginale every base change $S'\to S$, and $\Jj$ as above, the graded module
\[
\sisi{\coprod}{\tbigoplus}_{m\geq 0} \R^ng'_\ast(\Ii^mF'/\Ii^{m+1}F')
\]
over $\gr_\sisi{I}{\Ii}(\Oo_{X'})=\sisi{\coprod}{\bigoplus}_{m\geq 0} \Ii^m/\Ii^{m+1}$ is of finite type.
\end{enumeratei}
\end{proposition}

Obviously, (ii) \ALORS (i) and (iii) \ALORS (i), as one sees by taking $\Ii=0$ in conditions (ii) and (iii). The reverse implications are obtained by applying (i) to the composite base change $S^{''}\to S'\to S$, where $S^{''}$ is equal to $\Spec \sisi{\coprod}{\bigoplus}_{m\geq 0} \Jj^m$ \resp $\Spec \sisi{\coprod}{\bigoplus}_{m\geq 0} \Jj^m/\Jj^{m+1}$.

The point of this proposition is that conditions of the form (ii) are those that occur in the \og algebraic-formal comparison theorems\fg, whereas conditions of the form (iii) occur in the \og existence theorems\fg that complement them; \cf Expos\'e \Ref{IX}. A first interesting case is that in which $f:X\to S$ is the identity, so that one is dealing with conditions on a morphism $h:U\to S$ locally of finite type and a quasi-coherent module $F$ on $U$, flat over $S$. To obtain sufficient conditions, we shall assume that $U$ embeds by $g:U\to X$ as an open subprescheme of an $X$ that is \emph{proper} over $S$. Applying \Ref{XII.5.1}, we thus see:

\begin{corollaire} \label{XII.5.8}
Let $f:X\to S$ be a proper morphism, with $S$ locally Noetherian, let $U$ be an open subset of $X$, let $g:U\to X$ be the canonical immersion, let $h=fg:U\to S$, and let $F$ be a quasi-coherent module on $\sisi{X}{U}$, flat over $S$. Assume that, for every base change $S'\to S$, with $S'$ locally Noetherian, $\R^ig'_\ast(F')$ is coherent on $X'$ for $i\leq n$. Then the following hold:
\begin{enumeratei}
\item
For every base change $S'\to S$, with $S'$ locally Noetherian, $\R^ih'_\ast(F')$ is coherent on $S'$ for $i\leq n$.
\item
For every $S'\to S$ as above, and every coherent ideal $\Jj$ on $S'$, the graded modules\pageoriginale
\[
\sisi{\coprod}{\tbigoplus}_{m\geq 0} \R^ih'_\ast(\Jj^mF')
\]
over $\sisi{\coprod}{\bigoplus}_{m\geq 0}\Jj^m$ are of finite type for $i\leq n$.
\item
For every $S'\to S$ and $\Jj$ as above, the graded modules
\[
\sisi{\coprod}{\tbigoplus}_{m\geq 0} \R^ih'_\ast(\Jj^mF'/\Jj^{m+1}F')
\]
over $\gr_\Jj(\Oo_{S'})=\sisi{\coprod}{\bigoplus}_{m\geq 0}\Jj^m/\Jj^{m+1}$ are of finite type for $i\leq n$.

Moreover, under the conditions of (ii), and by the comparison theorem \Ref{IX.1.1}, denoting by $\widehat{S'}$ the formal completion of $S'$ with respect to $\Jj$, and by $\widehat{U'}$ that of $U'$ with respect to $\Jj\Oo_{U'}$, the canonical homomorphisms
\[
\widehat{\R^ih'_\ast(F')} \to \R^i\widehat{h'}_\ast(\widehat{F'}) \to \varprojlim_k\R^ih'_\ast(F'_k)
\]
are isomorphisms for $i\leq n-1$.
\end{enumeratei}
\end{corollaire}

\begin{remarque} \label{XII.5.9}
Assume in addition that the conditions of \Ref{XII.5.8} hold with $F$ \emph{coherent}, and consider a base change $S'\to S$ as in \sisi{\Ref{XII.5.9}}{\Ref{XII.5.8}} (i). Assume further that $S'$ is locally embeddable in a regular scheme, or, more generally, that there exists a faithfully flat and quasi-compact morphism $S^{''}\to \sisi{S}{S'}$ such that $S^{''}$ is locally embeddable in a regular scheme; this condition holds in particular if $S'$ is local. Then the conclusion of \Ref{XII.5.8} (i) and (ii) remains valid when $F'$ is replaced there by a module $G'$ on $U'$ such that every point of $U'$ has an open neighborhood on which $G'$ is isomorphic to a module of the form ${F'}^n$. Indeed, one reduces to the case where $S'$ itself is locally embeddable in a regular scheme, and hence so are $S^{''}=\Spec\big(\sisi{\coprod}{\bigoplus}_{m\geq 0}\Jj^m\big)$ and $X^{''}=X'\times_{S'}S^{''}=X\times_SS^{''}$, which are of finite type over it. One then applies the finiteness criterion \Ref{VIII}~\Ref{VIII.2.3} to the direct images for $i\leq n$ of $G^{''}$ under the immersion $U^{''}\to X^{''}$, noting that its conditions hold by hypothesis for $F^{''}$\pageoriginaled, and therefore also for $G^{''}$, since they are expressed in terms of depth and $G^{''}$ is locally isomorphic to an ${F^{''}}^n$. The same argument shows that, if $\Gg'$ is a coherent module on $\widehat{U'}$ (the completion of $U'$ with respect to the ideal $\Jj\Oo_{U'}$), such that $\Gg'_0=\Gg'/\Jj\Gg'$ is locally of the form $\sisi{F_0^{'n}}{{F'_0}^n}$, then the conclusion of (iii) remains valid when $F'$ is replaced there by $\Gg'$. Using the results of Expos\'e \Ref{IX}, one thus obtains the following result:
\end{remarque}

\begin{corollaire} \label{XII.5.10}
Let $f:X\to S$ be a proper morphism, with $S$ locally Noetherian, and let $U$ be an open subset of $X$; assume that $U$ is flat over $S$, and that, for every base change $S'\to S$, with $S'$ locally Noetherian, $\sisi{R^i}{\R^i}g'_\ast(\Oo_{U'})$ is coherent on $X'$ for $i=0,1$. Assume now that $S'$ is of the form $\Spec(A)$, where $A$ is a Noetherian ring equipped with an ideal $J$ such that $A$ is separated and complete for the $J$-adic topology. Under these conditions:
\begin{enumeratei}
\item
The functor $F\mto\widehat{F}$ from the category of locally free modules on $U'$ to the category of locally free modules on $\widehat{U'}$ is fully faithful.
\item
For every locally free module $\formelF$ on $\widehat{U'}$, there exists a coherent module $F$ on~$U'$ (not necessarily locally free, alas), and an isomorphism $\widehat{F}\simeq\formelF$.
\end{enumeratei}
\end{corollaire}

It remains only to prove (ii), by \Ref{XII.5.9}. But it follows from this remark and \Ref{IX.2.1} that $\formelF$ is induced by a coherent module $\formelG$ on $\widehat{X'}$. By the existence theorem \EGA III 5.1.4, $\formelG$ is of the form $\widehat{F}$, where $F$ is coherent on $\sisi{X}{X'}$, which proves the conclusion.

\begin{remarques} \label{XII.5.11}
1\textsuperscript{o}) Using \Ref{XII.5.10}, \Ref{XII.4.7}, and a suitable hypothesis saying that certain local rings of the geometric fibers of $X'\to S'$ are \og pure\fg \resp parafactorial, one should be able to obtain statements saying that the functor $Z'\mto\widehat{Z'}$ from the category of finite \`etale coverings of $X'$ to the category of finite \`etale coverings\pageoriginale of $\widehat{X'}$ (or, equivalently, of $X'_0$) is an equivalence of categories, \resp that the functor $L\mto\widehat{L}$ from the category of invertible modules on $X'$ to the category of invertible modules on $\widehat{X'}$ is an equivalence. Using recent results of Murre, it is likely that one should be able to deduce existence theorems for Picard schemes of certain \emph{nonproper} algebraic schemes\nde{Of course, in the projective case one should refer to Grothendieck's existence theorems in FGA; \cf Grothendieck~A., {\og Technique de descente et th\'eor\`emes d'existence en g\'eom\'etrie alg\'ebrique. VI. Les sch\'emas de Picard: propri\'et\'es g\'en\'erales\fg}, in \emph{S\'eminaire Bourbaki}, vol.~7, Soci\'et\'e math\'ematique de France, Paris, 1995, \Exp 236, p\ptbl 221--243 and {\og Technique de descente et th\'eor\`emes d'existence en g\'eom\'etrie alg\'ebrique. V. Les sch\'emas de Picard: th\'eor\`emes d'existence\fg}, in \emph{S\'eminaire Bourbaki}, vol.~7, Soci\'et\'e math\'ematique de France, Paris, 1995, \Exp 232, 143--161. The nine finiteness conjectures found there are proved in Expos\'es \Ref{XII} and \Ref{XIII} by Mme Raynaud and Kleiman in \SGA 6. For an excellent elementary account of the subject, see Kleiman's expository article (Kleiman~S., {\og The Picard scheme\fg}, forthcoming in Contemp. Math.). For an application of these techniques to the global generalized Jacobians of a relative smooth curve, see (Contou-Carr\`ere~C., {\og La jacobienne g\'en\'eralis\'ee d'une courbe relative; construction et propri\'et\'e universelle de factorisation\fg}, \textit{C.~R. Acad. Sci. Paris S\'er.~A-B} \textbf{289} (1979), \numero3, A203--A206 and {\og Jacobiennes g\'en\'eralis\'ees globales relatives\fg}, in \emph{The Grothendieck Festschrift, Vol. II}, Progr. Math., vol.~87, Birkh\"auser, Boston, 1990, p\ptbl 69--109). See also, by the same author and in the purely local setting, the construction and study of the \og local generalized Jacobian\fg functor ({\og Jacobienne locale, groupe de bivecteurs de Witt universel, et symbole mod\'er\'e\fg}, \textit{C.~R. Acad. Sci. Paris S\'er.~I Math.} \textbf{318} (1994), \numero 8, p\ptbl 743--746). Moreover, although in the case of a projective smooth morphism the connected components of the Picard scheme are proper, this is no longer so in the singular case. The problem of compactifying Picard schemes then arises naturally; this problem has been studied in detail, notably in (Altman~A.B. \& Kleiman~S., {\og Compactifying the Picard scheme\fg}, \emph{Adv. in Math.} \textbf{35} (1980), \numero 1, p\ptbl 50--112, and {\og Compactifying the Picard scheme. II\fg}, \emph{Amer. J.~Math.} \textbf{101} (1979), \numero 1, p\ptbl 10--41). The case of curves had been studied earlier (D'Souza~C., {\og Compactification of generalised Jacobians\fg}, \emph{Proc. Indian Acad. Sci. Sect. A Math. Sci.} \textbf{88} (1979), \numero 5, p\ptbl 419--457). One even knows exactly when the compactified Jacobian of a curve is irreducible (Rego~C.J., {\og The compactified Jacobian\fg}, \emph{Ann. Sci. \'Ec. Norm. Sup. (4)} \textbf{13} (1980), \numero 2, p\ptbl 211--223), the latter being the closure of the ordinary Jacobian when the curve is geometrically integral and locally planar; for a construction in families of compactified Jacobians, see (Esteves~E., {\og Compactifying the relative Jacobian over families of reduced curves\fg}, \emph{Trans. Amer. Math. Soc.} \textbf{353} (2001), \numero 8, p\ptbl 3045--3095). Since then, existence results for the Picard scheme in the proper case have advanced beyond the original edition of \SGA 2; \cf (Murre~J.P., {\og On contravariant functors from the category of pre-schemes over a field into the category of abelian groups (with an application to the Picard functor)\fg}, \emph{Publ. Math. Inst. Hautes \'Etudes Sci.} \textbf{23} (1964), p\ptbl 5--43), and especially (Artin~M., {\og Algebraization of formal moduli. I\fg}, in \emph{Global Analysis (Papers in Honor of K\ptbl Kodaira)}, Univ. Tokyo Press, Tokyo, 1969, p\ptbl 21--71). See also (Raynaud~M., {\og Sp\'ecialisation du foncteur de Picard\fg}, \emph{Publ. Math. Inst. Hautes \'Etudes Sci.} \textbf{38} (1970), p\ptbl 27--76) for a scheme proper over a discrete valuation ring but not necessarily flat. For a discussion of more recent results, in particular Artin's, for the Picard functor of proper flat schemes, especially in the cohomologically flat case in dimension $0$, see Chapter VIII of (Bosch~S., L\"utkebohmert~W. \& Raynaud~M., \emph{N\'eron models}, Ergebnisse der Mathematik und ihrer Grenzgebiete (3), vol.~21, Springer-Verlag, Berlin, 1990) and the references cited there. Much more recently, very precise results have been obtained for relative curves $f:X\to S$ \sisi{over}{over} the spectrum $S$ of a discrete valuation ring with perfect residue field. More precisely, one assumes that $f$ is proper and flat, that $X$ is regular, and that $f_*\OX=\Oo_S$. On the other hand, one does not assume that $f$ is cohomologically flat in dimension~$0$, \ie one does not assume that $H^1(X,\Oo)$ is torsion-free. The Picard functor is then not representable, either by a scheme or by an algebraic space, as soon as the torsion in question is nonzero. Let~$J$ be the N\'eron model of the generic fiber of $f$; it is a quotient of the Picard functor $P$. Raynaud then showed that the kernel of the tangent map $H^1(X,\Oo)=\mathrm{Lie}(P)\to\mathrm{Lie}(J)$ coincides with the torsion subgroup of $H^1$, and that the cokernel has the same length (see Theorem 3.1 of (Liu~Q., Lorenzini~D. \& Raynaud~M., {\og N\'eron models, Lie algebras, and reduction of curves of genus one\fg}, \emph{Invent. Math.} \textbf{157} (2004), p\ptbl 455--518)). This result enables those authors to study the relation between the Birch--Swinnerton-Dyer and Artin--Tate conjectures (see Th.\ptbl 6.6 of {\loccit}). As for the local Picard scheme, see Boutot's thesis, cited in editor's note~\eqref{Boutot} on page~\pageref{Boutot}.}. In general, the removal of purity hypotheses in various existence theorems, especially representability theorems for functors such as the Hilbert or Picard functors, by means of the techniques developed in this seminar, deserves systematic study.

2\textsuperscript{o}) One may seek manageable necessary and sufficient conditions, in terms of depth, for the universal finiteness condition considered in \Ref{XII.5.10} to hold. When $S$ is the spectrum of a field, it follows easily from \EGA III 1.4.15 that it is necessary and sufficient for the $\R^ig_\ast(F)$ ($i\leq n$) to be coherent, which is indeed expressed in terms of depth by \Ref{VIII}~\Ref{VIII.2.3}. In the general case, however, note that it is not enough to require the preceding condition to hold for all the fibers $U_s\subset X_s$ ($s\in S$), even when $n=0$. For example, take $X=S$, let $S$ be the spectrum of a discrete valuation ring, let $U$ be the open subset consisting only of the generic point, and let $F=\Oo_U$.

3\textsuperscript{o}) Here, however, is a \emph{sufficient} condition ensuring that the hypotheses of \Ref{XII.5.10} hold: it is enough that $f$ be flat and that, for every $s\in S$ and every $x\in Y_s=X_s-U_s$, one have
\[
\prof \Oo_{X_s,x}\geq n+2.
\]
Indeed, taking Lemma \Ref{XII.2.5} into account (\cf relation (\Ref{eq:XII.16}) following \Ref{XII.2.5}), it follows that $g_*(\Oo_U)\simeq\OX$ and $\R^ig_*(\Oo_U)=0$ for $0<i\leq n$, and the same relations will clearly hold after every base change $S'\to S$.
\end{remarques}

% END INCLUDED active_source/parts/12_expose_XII.tex

% BEGIN INCLUDED active_source/parts/13_expose_XIII.tex
\chapter{Problems and conjectures} \label{XIII}

\section[Relations between global and local results]
{Relations between global and local results. Affine problems connected with duality} \label{XIII.1}

It\pageoriginale is well known that many statements concerning a projective scheme
$X$ can be formulated in terms of statements concerning
a certain graded ring, or better, a complete local ring,
\sisi{\ignorespaces}{namely} the homogeneous coordinate ring
of $X$ (\ie the affine ring of the cone $\tilde{X}$ over
$X$), or its completion (\ie the completion of the local ring at the vertex
of $\tilde{X}$). The value of this reformulation is that,
starting from known global results, it often makes it possible to
conjecture, or even prove, analogous results for complete
Noetherian local rings more general than those which
actually occur in the global statement, for example for
local rings that are not necessarily of equal
characteristic. Thus, Serre's duality theorem for
projective space \Ref{XII}~\Ref{XII.1.1} suggested the useful
local duality theorem \sisi{V.5.6}{\Ref{V}~\Ref{V.2.1}}. Serre's
fundamental theorem on the cohomology of coherent
\sisi{m}{M}odules on projective space (finiteness,
a\sisi{s}{}symptotic behavior for large $n$ of $\H^i(X,
F(n))$, {\cf \EGA III~2.2.1}) generalizes to a
structure theorem for the local invariants $\H^i_{\mm}(M)$; see
\sisi{V.5.7}{\Ref{V}~\Ref{V.3}}. Likewise, the Lefschetz theorems
for the fundamental group and the Picard group (\og
equivalence criteria\fg), well known in the
classical case and subsequently extended to an arbitrary base field,
suggested the \og local\fg Lefschetz theorems of Exposés
\Ref{X} and \Ref{XI}. Of course, the local theorems are in
turn valuable tools for obtaining global statements.
For example, local duality allows a global asymptotic property
\Ref{XII}~\Ref{XII.1.3} (i) to be formulated by the vanishing
of certain local invariants $\H_x^i(F_x)$. More
substantially, the local Lefschetz theorems, implying for
example the \og purity\fg or parafactoriality of certain
complete-intersection local rings (\Ref{X}~\Ref{X.3.4} and
\Ref{XI}~\Ref{XI.3.13}), make it possible in the global
Lefschetz theorems to dispense with certain nonsingularity
hypotheses, as in \Ref{X}~\Ref{X.3.5}, \Ref{X.3.6},
\Ref{X.3.7}.

Another\pageoriginale useful generalization of theorems concerning projective schemes over a field $k$ consists in replacing $k$ by a general base scheme. Thus, the sequel to {\EGA III} will give\nde{in fact, this generalization is not found there; see the note below and the editor's note~\eqref{noteconrad} on page~\pageref{noteconrad}} a generalization in this direction of Serre duality\sfootnote{\cf the Hartshorne Seminar, cited at the end of {\Exp \Ref{IV}}.}; the finiteness and asymptotic-behavior theorems for $\H^i(X, F(n))$ were stated in \EGA III~2.2.1 over a general base scheme; finally, the Lefschetz theorems can likewise be developed for a projective morphism, as we saw in \Ref{XII}~\Ref{XII.4.9}, thanks to the local theorem \Ref{XII}~\Ref{XII.4.7}. Of course, working over a general base scheme also leads to essentially new statements, such as the \og comparison theorem\fg \EGA III~4.1.5 and the existence theorem for sheaves \EGA III~5.1.4 (which, as we saw elsewhere in \Ref{IX}, arise from the same key theorems of a cohomological nature as the Lefschetz theorems for $\pi_1$ and $\Pic$).

It is therefore necessary to isolate theorems that simultaneously encompass the two generalizations just mentioned of statements concerning projective schemes over a field. The natural objects for such a common generalization are the \emph{Noetherian rings separated and complete for an $I$-adic topology}. From this point of view, their study has not yet been seriously undertaken and seems to me at present the most interesting subject in the local theory of coherent sheaves. Here is a typical problem in this direction:

\refstepcounter{subsection}\label{XIII.1.1}
\begin{enonce*}{Conjecture 1.1}
 [\og Second affine finiteness theorem\fg
 \markedfootnote{**}{This conjecture, and Conjecture \hyperref[XIII.1.2]{1.2} below, are false, as was shown by \sisi{R.~HARTSHORNE (unpublished)}{R\ptbl Hartshorne, {\og Affine duality and cofinite modules\fg}, \emph{Invent. Math.} \textbf{9} (1969/70), p.~145-164, Section 3.}}
 \sisi{}{\setcounter{footnoteNDE}{1}}\ndemark]
\ndetext{However, if $A$ is complete local (respectively regular of positive characteristic) and $J$ is the maximal ideal, the statement is true for $M$ of finite type (respectively $M=A$), \cf (Hartshorne~R., {\og Affine duality and cofinite modules\fg}, \emph{Invent. Math.} \textbf{9} (1969/70), p\ptbl 145-164, Corollary 1.4) (respectively (Huneke~C. \& Sharp~R., {\og Bass numbers of local cohomology modules\fg}, \emph{Trans. Amer. Math. Soc.} \textbf{339} (1993), \numero 2, p\ptbl 765--779), which moreover contains much stronger results). For completely different methods ($D$-modules) for treating characteristic zero, see (Lyubeznik~G., {\og Finiteness properties of local cohomology modules (an application of $D$-modules to commutative algebra)\fg}, \emph{Invent. Math.} \textbf{113} (1993), \numero 1, p\ptbl 41--55); see also, by the same author, {\og Finiteness properties of local cohomology modules for regular local rings of mixed characteristic: the unramified case\fg}, \emph{Comm. Algebra} \textbf{28} (2000), \numero 12, p\ptbl 5867--5882, Special issue in honor of Robin Hartshorne, and {\og Finiteness properties of local cohomology modules: a characteristic-free approach\fg}, \emph{J.~Pure Appl. Algebra} \textbf{151} (2000), \numero 1, p\ptbl 43--50. The notion of a cofinite module has since evolved under Hartshorne's leadership. One says that $M$ is $J$-cofinite if its support is contained in $V(J)$ and if all the $\Ext^i_A(A/J, M)$ are of finite type. On this subject, see for example (Delfino~D. \& Marley~Th., {\og Cofinite modules and local cohomology\fg}, \emph{J. Pure Appl. Algebra} \textbf{121} (1997), \numero 1, p\ptbl 45--52).}
Let $M$ be a module of finite type over a Noetherian ring $A$ (which may, if necessary, be assumed to be a quotient of a regular ring), and let $J$ be an ideal of $A$. Prove that the modules $\H^i_J(M)$ are \og $J$-cofinite\fg, \ie that the modules
$$
\Hom_A(A/J, \H^i_J(M))$$
are of finite type.
\end{enonce*}

Recall that $\H^i_J (M)$ denotes the module $\H^i_Y (X, \tilde{M})$ (where $X=\Spec (A)$, $Y=V(J)$) of \Exp \Ref{I}, interpreted in \Ref{II} in terms of an inductive limit of the cohomology\pageoriginale of Koszul complexes, or, for $i\ge 2$, the module $\H^{i-1} (X-Y, \tilde{M})$. \sisi{T}{T}o tell the truth, \Ref{XIII.1.1} should be a consequence of a more precise statement, implying that the $\H^i_J ({M})$ belong to a suitable abelian subcategory ${\DDJ}$ of the category ${\CJ}$ of $A$-modules with support $\subset Y=V(J)$, such that $H\in \Ob {\DDJ}$ implies that $H$ is $J$-cofinite. (\sisi{NB}{N.B.} The category of modules $H$ with support contained in $V(J)$ and which are $J$-cofinite is unfortunately not stable under passage to quotients!). The essential problem would then consist in defining ${\DDJ}$. More precisely, the solution of problem \Ref{XIII.1.1} should emerge (at least if $A$ is a quotient of a regular ring) from a duality theory generalizing both local duality and the duality theory for projective morphisms alluded to above, and which would be of the following kind:

\refstepcounter{subsection} \label{XIII.1.2}
\begin{enonce*}{Conjecture 1.2}[\og Affine duality\fg\sfootnotemark]
\sfootnotetext{This conjecture, false as stated, was nevertheless established in a fairly similar form by \sisi{R.~HARTSHORNE, \emph{Affine duality and cofinite modules} (forthcoming).}{R.~Hartshorne, {\og Affine duality and cofinite modules\fg}, \emph{Invent. Math.} \textbf{9} (1969/70), p\ptbl 145-164.}} Suppose that $A$ is regular, separated, and complete for the $J$-adic topology. Let $C^{\boule}(A)$ be an injective resolution of $A$.
\begin{enumeratei}
\item
Prove that the functor
$$D_J: L_{\boule}\mto \Hom_J(L_{\boule}, C^{\boule}(A))$$
from the category of complexes of $A$-modules that are free of finite type in every dimension and bounded above (\sisi{where}{where} the morphisms are homomorphisms of complexes up to homotopy) to the category of complexes of $A$-modules $K^{\boule}$ that are injective in every dimension and bounded above (\sisi{where}{where} the homomorphisms are defined in the same way), is fully faithful.

\item
Prove that, for every $K^{\boule}$ of the form $D_J(L_{\boule})$, the $\H^i(K^{\boule})(= \Ext_Y^i (X; L_{\boule}, \OX))$ are $J$-cofinite.

\item
More precisely, prove that the $K^{\boule}$ which are homotopic to a complex of the form $D_J(L_{\boule})$ can be characterized by finiteness properties of the $\H^i(K^{\boule})$ stronger than the one considered in (ii), for example by\pageoriginale the property $\H^i(K^{\boule})\in \Ob {\DDJ}$, where ${\DDJ}$ is a suitable abelian category, as contemplated above.
\end{enumeratei}
\end{enonce*}

Note that the problem has an affirmative solution when $A$ is local and $J$ is an ideal of definition (\cf \Exp \Ref{IV}), and also when $J$ is the zero ideal. In these two exceptional cases, one may take for ${\DDJ}$ simply the category of \sisi{m}{M}odules with support $V(J)$ which are $J$-cofinite (which in the second case means simply that one takes the category of \sisi{m}{M}odules of finite type over $A$). An affirmative solution of Conjecture \Ref{XIII.1.2} in general would give one for \Ref{XIII.1.1}, by taking for $L_{\boule}$ the dual of a finite free resolution of $M$. On the other hand, an affirmative solution of \Ref{XIII.1.1} would give an affirmative answer to the first part of the following conjecture, which we formulate in \og global\fg form:

\begin{conjecture} \label{XIII.1.3}
Let $X\subset \PP^r_k$ be a closed subscheme of the standard projective scheme which is locally a complete intersection and each of whose irreducible components has codimension $\ge s$. Let $U=\PP^r_k -X$.
\begin{enumeratei}
\item
Prove that, for every coherent Module $F$ on $U$, one has
$$
\dim \H^i(U, F)<+\infty{\rm \ for\ } i\ge s.{}\sfootnote{Part\sisi{}{ie} (i) of this conjecture was proved by R.~\sisi{HARTSHORNE}{Hartshorne} when $\sisi{Y}{X}$ is smooth \sisi{in \emph{Ample Vector Bundles}, Pub. Math. \numero 29, (1966)}{\cf \emph{Ample subvarieties of algebraic varieties}, Notes written in collaboration with C\ptbl Musili, Lect. Notes in Math, vol.~156, Springer-Verlag, Berlin-New York, 1970, Theorem III.5.2}. This author also found an example for (ii), \cf \sisi{R.~HARTSHORNE, \emph{Cohomological dimension of algebraic varieties}, forthcoming in the Ann. of Math.}{R.~Hartshorne, {\og Cohomological dimension of algebraic varieties\fg}, \emph{Ann. of Math. (2)} \textbf{88} (1968), p\ptbl 403--450, example on page 449.}}
$$

\item
Give an example, with $X$ connected and regular, for which one has
$$
\H^s(U, F)\neq 0.
$$
\end{enumeratei}
\end{conjecture}

To see that (i) is a special case of \Ref{XIII.1.1}, consider
$$
\H^i(U, F(\cdot))=\sisi{\coprod}{\tbigoplus}_n \H^i (U, F(n))= \H^i(E^{r+1} -\sisi{X}{\tilde X}, \tilde{F})$$
as a module over the affine ring $k[t_0, \ldots, t_r]$ of the cone $E^{r+1}$ over $\PP^r$. This module is none other than $\H_J^{i+1}(M)$, where $J$ is the ideal of the cone $\tilde X$ over $X$ in $E^{r+1}$. On the other hand, the hypothesis on $X$, which implies that $\sisi{X}{\tilde X}$ is also a complete intersection of co\sisi{-}{}dimension $\ge s$ at every\pageoriginale point of $E^{r+1}$ distinct from the origin, implies that $\H^{i+1}_J(M)$ is zero away from the origin for $i\ge s$. Thus, if it is $J$-cofinite as required by \Ref{XIII.1.1}, it is \afortiori $\mm$-cofinite, which readily implies that it is finite-dimensional in every degree\refstepcounter{toto}\nde{\label{noteHM}Hartshorne proved (Hartshorne~R., {\og Cohomological dimension of algebraic varieties\fg}, \emph{Ann. of Math. (2)} \textbf{88} (1968), p\ptbl403--450) that the cohomology $H^{n-1}(\PP^n_k-X,F)$ is zero for $F$ coherent and $X$ of positive dimension ($k$ algebraically closed). In fact, thanks essentially to Serre duality and the Lichtenbaum theorem --- the vanishing of the cohomology of coherent sheaves in the maximal dimension of irreducible noncomplete quasi-projective varieties ---, one is reduced to proving that the formal completion $\hat \PP^n_k$ and $X$ have the same field of rational functions. This is the difficult point (Theorem 7.2 of {\loccit}); in other words, $\PP^n_k$ is $G3$ in the terminology of (Hironaka~H. \& Matsumura~H., {\og Formal functions and formal embeddings\fg}, \emph{J.~Math. Soc. Japan} \textbf{20} (1968), p\ptbl 52--82). These authors independently proved the preceding results, and in fact much more. They proved that $X$ is universally $G3$ and calculated the field of rational functions of the formal completion of an abelian variety along a positive-dimensional subvariety. It is in this article that the now-classical conditions $G1,G2,G3$ appear for the first time.} Note moreover that Conjecture \Ref{XIII.1.3} already arises for a nonsingular irreducible curve $X$ in $\PP^3$: it is unknown whether in this case the $\H^2(\PP^3-X, \OX(n))$ are finite-dimensional, or whether they are necessarily zero\sfootnote{The question has just been answered affirmatively by R\ptbl \sisi{HARTSHORNE}{Hartshorne}\sisi{}{ (Hartshorne~R., {\og Ample vector bundles\fg}, \emph{Publ. Math. Inst. Hautes \'Etudes Sci.} \textbf{29} (1966), p\ptbl63--94, Theorem 8.1)} and H. \sisi{HIRONAKA}{Hironaka}.}. It is not even known whether there exists an irreducible curve in $\PP^3$ which is not set-theoretically the intersection of two hypersurfaces\label{courbesurface}\nde{see Conjecture~\Ref{XII.3.5} and the corresponding note.}.

\begin{probleme} \label{XIII.1.4}
Give an affine variant of the \og comparison theorem\fg \textup{\EGA III~4.1.5} as a theorem on the commutation of the functors $\H^i_J$ with certain inverse limits.
\end{probleme}

Finally, in the present line of thought, I had posed the following problem: \sisi{L}{l}et $A$ be a \emph{complete} regular Noetherian local ring, and $K$ its field of fractions; prove that $\Ext^i_A(K, A)=0$ for every $i$. An affirmative answer was given on the spot by \sisi{Mr AUSLANDER}{M.~Auslander}; the regularity hypothesis may be replaced by the hypothesis that $A$ is integral, and in fact $\Ext^i_A(K, M)=0$ for every $i$ whenever $M$ is of finite type over $A$. This follows immediately from the following statement, due to \sisi{AUSLANDER}{Auslander}: if $A$ is a complete Noetherian local ring, then for every module of finite type $M$ over $A$, the functors $\Ext_A^i(., M)$ transform inductive limits into inverse limits\nde{write $K=\varinjlim_{a\neq 0}A[1/a]$ and observe that $a\neq 0$ is $A$-regular.}.

\section{\texorpdfstring{Problems related to $\pi_0$: local Bertini theorems}{Problems related to pi\textunderscore 0: local Bertini theorems}} \label{XIII.2}

Let $A$ be a complete noetherian local ring, $f$ an element of its
maximal ideal, $X=\Spec (A)$, $Y=\Spec(A/fA)$. The use of the local
\og Lefschetz\fg technique makes it possible to give criteria
for $Y'=X'\cap Y$ (where $X'=X-\{\mm\}$) to be connected, in terms
of hypotheses on $X'$. Thus, it suffices that: a) $X'$\pageoriginale
is connected\sisi{}{,} b) $\prof \Oo_{X', x} \ge 2$ for every closed point $x$
of $X'$\sisi{}{,} c) $f$ is $A$-regular. Note, however, that
hypotheses b) and c) are not purely
topological in nature; for example, they are not invariant\sisi{}{}
under replacement of $X$ by $X_{\text{red}}$. In the analogous situation
for a projective scheme $X'$ over a field and a hyperplane
section $Y'$ of $X'$, the use of the \og Bertini theorem\fg
and Zariski's \og connectedness theorem\fg in fact gives
results of a distinctly more satisfactory form,
which had led me, in the oral seminar, to state
a conjecture that I have since solved affirmatively.
Let us therefore state here:

\begin{theoreme} \label{XIII.2.1}
Let $A$ be a complete noetherian local ring, $X$ its spectrum, $\aaa$ the closed point of $X$, $X'=X-\{ \aaa\}$. Suppose that $X$ satisfies the following conditions (where $k$ denotes an integer $\ge 1$):
\begin{itemize}
\item[a$_k$)] The irreducible components of $X'$ have dimension $\ge k+1$.

\item[b$_k$)] $X'$ is connected in dimension $\ge k$, \ie it cannot be disconnected by a closed subset of dimension $<k$ (\cf \Ref{III}~\Ref{III.3.8}).
\end{itemize}

Let $m$ be an integer, $0\le m\le k$, and let $f_1, \dotsc, f_m\in \rr(A)$; set $B=A/\sum_i f_iA$, $Y=\Spec(B)=V(f_1)\cap\dotsb\cap V(f_m)$, $Y'=X'\cap Y=Y-\{\aaa\}$. Then $Y$ satisfies conditions a$_{k-m}$), b$_{k-m}$). In particular, for every sequence of $m\le k$ elements $f_1, \dotsc, f_m$ of $\rr(A)$, $Y'=X'\cap V(f_1)\cap\dotsb\cap V(f_m)$ is connected.
\end{theoreme}

It is moreover easy to see that if the last conclusion holds (it plainly suffices to take $m=k$ there), and if the case where $X$ is irreducible of dimension~$0$ or $1$ is excluded, it follows that the irreducible components of $X'$ have dimension $\ge k+1$, and $X'$ is connected in dimension $\geq k$, so that, in a sense, \Ref{XIII.2.1} is a \og best possible\fg result.

\bigskip

Let us give the principle of the proof of~\Ref{XIII.2.1}. Only condition b$_{k-m}$) for $Y$ presents a problem. For a given $k$, one is easily reduced to the case where $X$ is integral, and even (by passing to the normalization, which is finite over $X$) to\pageoriginale the case where $X$ is \emph{normal}. If $k=1$, hence $\dim X' \ge 2$, then $X'$ has depth $\ge 2$ at its closed points and one may apply the result recalled at the beginning of \sisi{N${}^\circ$}{n$^\circ$}, which shows that \hbox{$Y'=X'\cap V(f)$} is connected. In the case $k\ge 1$, suppose the theorem proved for \hbox{$k'\!<\!k$}. By induction on $m$, one is reduced to the case $m=1$, \ie to verifying that, for $f_1\in\nobreak \rr(A)$, $X'\cap V(f_1)$ is connected in dimension $\ge k-1$. If it were not, \ie if it were disconnected by a $Z'$ of dimension $< k-1$, there would exist a sequence $f_2, \dotsc, f_k$ such that \hbox{$X'\cap V(f_1)\cap\dotsb\cap V(f_k)$} is disconnected, and in this sequence one may choose $f_2\in \rr(A)$ arbitrarily, subject only to the condition that it not vanish at any point of a certain finite subset $F$ of $X'$ (\sisi{\ignorespaces}{namely }the set of maximal points of $Z'$). Moreover, one easily verifies, using the fact that $X'$ is normal and therefore satisfies Serre's condition ({$S_2$})\sfootnote{\Cf \EGA IV 5.7.2.}, that there exists a finite subset $F'$ of $X'$ such that $f\in \rr(A)$, $V(f)\cap F'=\emptyset$ implies that $V(f)\cap X'$ also satisfies condition ({$S_2$}). One may then choose $f_2$ so that $f_2$ vanishes neither on $F$ nor on $F'$, and hence so that $X'\cap V(f_2)$ satisfies ({$S_2$}). But then\sisi{}{,} by the theorem of \sisi{HARTSHORNE}{Hartshorne} \Ref{III}~\Ref{III.3.6}\sisi{}{,} $X'\cap V(f_2)$ is connected in codimension~$1$ and thus (since every component of $X'\cap V(f_2)$ has dimension $\ge k$) it is connected in dimension $\ge k-1$. Applying the induction hypothesis to $V(f_2)=\Spec(A/f_2A)$, it follows that $X'\cap V(f_2)\cap V(f_1)\cap V(f_3)\cap \dotsb\cap V(f_k)$ is connected, whereas it had been constructed disconnected, a contradiction.

Let us mention some interesting corollaries:

\begin{corollaire} \label{XIII.2.2}
Let $f\colon X\to Y$ be a proper morphism of locally noetherian preschemes, with $Y$ integral, $y_0\in Y$, and $y_1$ the generic point of $Y$. Assume that
\begin{enumeratea}
\item
$Y$ is unibranch at $y_0$, and every irreducible component of $X$ dominates $Y$.
\item
The irreducible components of $X_1$ have dimension $\ge k+1$, and $X_1$ is connected in dimension $\ge k$.
\end{enumeratea}

Then\pageoriginale the irreducible components of $X_0$ have dimension $\ge k+1$, and $X_0$ is connected in dimension $\ge k$.
\end{corollaire}

Indeed, Zariski's connectedness theorem (\cf\ \EGA III 4.3.1) implies that $X_0$ is connected; to show that it is not disconnected by a closed subset of dimension $<k$, one is reduced to showing that the local rings at points $x\in X_0$ such that $\dim \overline{x}<k$ have a spectrum that is not disconnected by $x$. But this is true without assuming either that $f$ is proper or that $Y$ is unibranch at $y_0$. To see this, one is reduced to the case where $X$ is integral and dominates $Y$, and, if desired, where $Y$ is affine of finite type over $\ZZ$, so that one is under the hypotheses of the dimension formula for $\Oo_{X, x}$ over $\Oo_{Y, y_0}$. Using the finiteness of the normalization in this case, one may even suppose $X$ \emph{normal}; hence, by a theorem of \sisi{NAGATA}{Nagata}\sfootnote{\Cf \EGA IV 7.8.3 (i) (ii) (v).}, the completion of a local ring $\Oo_{X, x}$ of $X$ is still normal, and therefore (if $\Oo_{X, x}$ has dimension $N$) $\Spec (\hat{\Oo}_{X, x})$ is connected in dimension $\ge N-1$. Let $n=\dim\Oo_{\sisi{X}{Y}, y_0}$; then $\degtr~k(x)/k(y_0)<k$ implies $\dim \Oo_{X, x}>n+(k+1)-k=n+1$, in view \sisi{that }{of }$\dim X_1\ge k+1$, and, taking a system $f_1, \dotsc, f_n$ of parameters of $\Oo_{Y, y_0}$ and lifting them to elements of $\Oo_{X, x}$, one sees from \Ref{XIII.2.1} that $\Spec (\hat{\Oo}_{X, x}/\sum f_i\hat{\Oo}_{X, x})$ is connected in dimension $\ge 1$, \ie is not disconnected by its closed point, or, equivalently, $\Spec(\hat{\Oo}_{X_0, x})$ is not disconnected by its closed point; \afortiori the same holds for $\Spec (\Oo_{X_0, x})$.

As in the case of the ordinary connectedness theorem, one may vary \Ref{XIII.2.2} by taking geometric fibers (over algebraic closures of the residue fields), provided that one assumes $Y$ to be \emph{geometrically} unibranch at $y_0$, or (with no hypothesis other than that $Y$ is noetherian) that $f$ is universally open. Applying this to the case where $Y$ is the dual scheme of a projective scheme $\PP_k^r$ over a field, one recovers a strengthened form of the global result that had inspired \Ref{XIII.2.1}, namely:

\refstepcounter{subsection}\label{XIII.2.3}
\begin{enonce*}{Corollary 2.3\ndemark}\ndetext{for a very beautiful direct proof, see
(Fulton~W. \& Lazarsfeld.~R., {\og Connectivity and its applications in algebraic geometry\fg},
in \emph{Algebraic geometry (Chicago, Ill., 1980)}, Lect. Notes in Math., vol.~862,
Springer, Berlin-New York, 1981, p\ptbl 26--92, Theorem 2.1). \Cf also \cite{HL}, cited in the editor's note~\eqref{HL}~on page~\pageref{HL}.}
Let $X$ be a closed subscheme of $\PP^r_k$ ($k$ a field); suppose that the irreducible components of $X$ have dimension $\ge \sisi{k}{l}+1$, and that $X$ is geometrically connected\pageoriginaled in dimension $\ge \sisi{k}{l}$. Then for every sequence $H_1, \dotsc, H_m$ of $m$ hyperplanes of $\PP^r_k$ ($0\le m\le \sisi{k}{l}-1$), $X\cap H_1\cap \dotsb \cap H_m$ satisfies the same condition with $\sisi{k}{l}-m$; in particular, it is geometrically connected in dimension $\ge \sisi{k}{l}-1$.
\end{enonce*}

One may moreover modify this statement in the
obvious way for the case where one is given a proper morphism
$X \to \PP^r_k$, which is not necessarily an
immersion, and an analogous extension is possible for
\Ref{XIII.2.1} (by considering a scheme proper over~$X'$). These
statements moreover follow formally from the statements given
here, in view of the ordinary connectedness theorem, which
reduces us to the case of a finite morphism.

\begin{corollaire} \label{XIII.2.4}
Let $A$ be a complete normal noetherian local ring of dimension $\ge k+2$. Let $X=\Spec (A)$, $X'=X-\{ \aaa\}$, and let $f_1, \dotsc, f_k$ be elements of $\rr(A)$; then $Y'=X'\cap V(f_1) \cap \dotsb \cap V(f_k)$ is connected, and $\pi_1(Y')\to \pi_1(X')$ is surjective.
\end{corollaire}

One proceeds as in \SGA 1~X~2.11.

Throughout the preceding discussion, only questions of \emph{connectedness} were involved. In the global case, however, well-known theorems assert that for an irreducible projective variety $X\subset \PP^r_k$, with $k$ algebraically closed, its intersection with a sufficiently \og general\fg hyperplane $H$ is irreducible (and not merely connected): this is the \emph{\sisi{BERTINI}{Bertini} theorem}, proved by \sisi{ZARISKI}{Zariski}, which in turn implies, by Zariski's connectedness theorem, that for \emph{every} $H$, $H\cap X$ is connected (although not necessarily irreducible). One may moreover proceed in the opposite direction, by proving the latter result using a Lefschetz-type technique and deducing the Bertini theorem by reducing to the case where $X$ is normal and using the following result: for a \og sufficiently general\fg $H$, $X\cap H$ is also normal. This suggests the

\refstepcounter{subsection}\label{XIII.2.5}
\begin{enonce*}{Conjecture 2.5\ndemark}\ndetext{see the following editor's note.}
Let $A$ be a complete normal noetherian local ring. Show that there exists a nonzero $f\in\rr(A)$ such that $Y'=X'\cap V(f)=Y-\{ \aaa\}$ (where\pageoriginale $Y=\Spec (A/fA)$) is normal (and hence irreducible by \Ref{XIII.2.1} if $\dim A\ge 3$).
\end{enonce*}

Ideally, one should show that, in a suitable sense, there are even \og many\fg elements $f$ having the property in question; for example, that $f$ may be chosen in an arbitrary power of the maximal ideal. Using Serre's normality criterion and the remark made above concerning Serre's property $(S_2)$, one sees that \Ref{XIII.2.5} would have an affirmative answer if the same were true for the

\refstepcounter{subsection}\label{XIII.2.6}
\begin{enonce*}{Conjecture 2.6\ndemark}
\ndetext{a proof of this conjecture can now be found in the literature, and the preceding conjecture must therefore also be regarded as proved, as indicated above. One can also find two earlier published but unfortunately unsuccessful proof attempts, by Flenner and Trivedi. See Trivedi V., {\og Erratum: ``A local Bertini theorem in mixed characteristic"\fg}, \emph{Comm. Algebra} \textbf{25} (1997), \numero 5, p\ptbl1685-1686. However, the editor has not verified that the proof is now complete.}
Let $A$ be a complete noetherian local ring, $U$ an open subset of its spectrum $X$, and $F$ a finite subset of $X'=X-\{ \aaa\}$. Suppose that $U$ is regular. Prove that there exists an $f \in\rr(A)$ such that $V(f)\cap U$ is regular, and $V(f)\cap F=\emptyset$.
\end{enonce*}

For a result of \og local Bertini\fg type, see \sisi{CHOW}{Chow}~\cite{XIII.2}

\section{\texorpdfstring{Problems related to $\pi_1$}{Problems related to pi\textunderscore 1}} \label{XIII.3}
Here again there are many questions, suggested by the global results or by
the transcendental results.

\refstepcounter{subsection}\label{XIII.3.1}
\begin{enonce*}{Conjecture 3.1\ndemark}
\ndetext{The analogous statement is true for (connected) schemes $X$ of
finite type over a separably closed field $k$, under the hypothesis of
strong desingularization for all $\bar k$-schemes of finite type; in
particular, this holds if $k$ has characteristic zero or if $X$ has
dimension $\leq 2$. One reduces to the case of quasi-projective surfaces
by Lefschetz-type techniques developed by Mme Raynaud; \cf the notes
\textit{supra}; see \SGA 7.1, Theorem II.2.3.1.}
Let $A$ be a complete noetherian local ring with algebraically closed
residue field, let $X=\Spec(A)$, let $X'=X-\{\aaa\}$, where $\aaa$ is the
closed point. Suppose that the irreducible components of $X$ have
dimension $\geq 2$ and that $X'$ is connected.
\begin{enumeratei}
\item
Prove that $\pi_1(X')$ is topologically finitely generated.
\item
If $p$ is the characteristic exponent of the residue field $k$ of $A$,
prove that the largest topological quotient group of $\pi_1(X')$ that is
\og of order prime to $p$\fg is finitely presented.
\end{enumeratei}
\end{enonce*}

For part (i), using the descent theory of \SGA 1~IX~5.2 and
Theorem\pageoriginale~\hyperref[XIII.2.4]{2.4}, one is reduced to the case in which $A$ is normal of
dimension $2$. In this case, a systematic method for studying the
fundamental group of $X'$, introduced by \sisi{MUMFORD}{Mumford}~\cite{XIII.5}
in the transcendental setting, consists in \emph{desingularizing} $X$,
\ie considering a projective birational morphism $Z\to X$, with $Z$
integral and regular, which induces an isomorphism
$Z'=Z|X'\to X'$. It is plausible that such a $Z$ always exists; in
\sisi{all}{any} case this is what \sisi{ABHYANKAR}{Abhyankar}'s method~\cite{XIII.1}
proves in the case of \og equal characteristic\fg\sfootnote{The
possibility of \og resolving\fg $A$ has now been proved in full
generality by Abhyankar~\cite{XIII.8}.}. Let $C$ be the fiber over the
closed point of $X$ under $Z\to X$; it is an algebraic curve over the
residue field $k$, connected by the connectedness theorem. The solution
of \Ref{XIII.3.1} then seems to be related to the following problem.

\begin{probleme} \label{XIII.3.2}
With the preceding notation, relate $\pi_1(X')$ to the topological
invariants of $C$, in particular to its fundamental group, so as to
exhibit the topological finite generation of $\pi_1(X')$, using for
example \SGA \sisi{}{\textup{I}, Theorem~}\textup{X~2.6}.
\end{probleme}

Another method would be to regard $A$ as a finite algebra over a
complete \emph{regular} local ring $B$ of dimension $2$, ramified along
a curve $C$ contained in $\Spec(B)=Y$. This leads to the following
problem.

\begin{probleme} \label{XIII.3.3}
Let $A$ be a complete regular local ring of dimension $2$, with
algebra\sisi{.}{ically} closed residue field $k$; let $X$ be its spectrum,
and let $C$ be a closed subset of $X$ of dimension~$1$. Define local
invariants of the embedded curve $C$ that have a meaning independent of
the residue characteristic and whose knowledge makes it possible to
calculate the fundamental group of $X-C$ by generators and relations
when $k$ has characteristic zero. Prove that when $k$ has characteristic
$p>0$, the \og tame\fg fundamental group of $X-C$ is a quotient of the
preceding one, and that the two fundamental groups (in
\sisi{char.}{characteristic}~$0$ and in
\sisi{char.}{characteristic} $p>0$) have the same maximal quotient of
order prime to $p$.
\end{probleme}

Of course, \Ref{XIII.3.3} shows that in \Ref{XIII.3.1} one should also
replace $X'$ by a scheme of the form $X-Y$, where $Y$ is a closed subset
of $X$ that\pageoriginale has codimension $\geq 2$ in every component of $X$ that
contains it. When this restriction on $Y$ is dropped, there should still
be an analogous finiteness result, provided that \og tame\fg-type
restrictions are imposed on the ramification at the maximal points of
the irreducible components of $Y$ that have codimension~$1$.

\begin{probleme} \label{XIII.3.4}
Let $A$ be a complete noetherian local ring of dimension $2$, with
algebraically closed residue field. Again let $X=\Spec(A)$ and
$X'=X-\{\aaa\}$. Find special structural properties of $\pi_1(X')$ in
the case where $A$ is a complete intersection.
\end{probleme}

A satisfactory solution of this problem might make it possible to solve
the following old problem:

\refstepcounter{subsection}\label{XIII.3.5}
\begin{enonce*}{Conjecture 3.5\ndemark}
\ndetext{As of autumn 2004, this problem is still open.}
Find an irreducible curve in $\PP^3_k$ (where $k$ is an algebraically
closed field), preferably nonsingular, which is not set-theoretically the
intersection of two hypersurfaces.
\sisi{(KNESER~\cite{XIII.4} shows that it can always be obtained as the
intersection of three hypersurfaces).}{\ignorespaces}
\end{enonce*}
\sisi{\ignorespaces}{\noindent(Kneser~\cite{XIII.4} shows that it can
always be obtained as the intersection of three hypersurfaces).}

\section[\texorpdfstring{Problems related to higher $\pi_i$}{Problems related to higher pi\textunderscore i}]{Problems related to higher
$\pi_i$: local and global Lefschetz theorems for complex analytic
spaces\protect\ndemark} \label{XIII.4}

Let $X$ be a scheme\ndetext{The statements are made more precise in the
Comments (Section~\Ref{XIII.6}). The conjectures appearing there have
become theorems; \cf the footnotes to Section~\Ref{XIII.6}.} locally of
finite type over the field of complex numbers $\CC$. One can associate
with it a complex analytic space $X^h$, and hence homotopy and homology
invariants $\pi_i(X^h)$, $\H_i(X^h)$, $\H^i(X^h)$, etc.
\sisi{\dotsc}{} Moreover, $X$ is known to be connected if and only if
$X^h$ is connected, so there is a bijection
$$
\pi_0(X^h)\to \pi_0(X).
$$
Similarly, since every \sisi{\'etale covering}{\'etale cover} $X'$ of
$X$ defines an \sisi{\'etale covering}{\'etale cover} $X^{\prime h}$ of
$X^h$\pageoriginale, there is a canonical homomorphism
$$
\pi_1(X^h)\to \pi_1(X),
$$
which, by a theorem of Grauert--Remmert, is known to identify the second
group with the completion of the first for the topology defined by
subgroups of finite index. This simply expresses the fact that
$X'\sisi{\rightsquigarrow}{\mto}X^{\prime h}$ is an equivalence from the
category of \sisi{\'etale coverings}{\'etale covers} of $X$ to the
category of finite \sisi{\'etale coverings}{\'etale covers} of $X^h$.
It follows that the results of this seminar, obtained by purely
algebraic methods, concerning $\pi_0(X)$ and $\pi_1(X)$ imply results
for $\pi_0(X^h)$ and $\pi_1(X^h)$, which are transcendental in nature.
Moreover, if $X$ is proper, the well-known exact sequence
$0\to\ZZ\to\CC\to\CC^*\to0$ shows that the N\'eron--Severi group of $X$
(the quotient of its Picard group by the connected component of the
identity) is isomorphic to a subgroup of $\H^2(X^h,\ZZ)$; in the
nonsingular K\"ahler case, it is the subgroup denoted
$\H^{(1,1)}(X^h,\ZZ)$ (classes of type $(1,1)$):
$$
\Pic(X)/\Pic^0(X)\subset \H^2(X^h,\ZZ).
$$

Consequently, the information that we have obtained about Picard groups
implies information, admittedly very partial, about the groups
$\H^2(X^h,\ZZ)$. It is tempting to complete all these fragmentary
results with conjectures.

Very precise indications pointing in the same direction as those just
mentioned are supplied by a classical theorem of
\sisi{LEFSCHETZ}{Lefschetz}~\cite{XIII.7}. It asserts that if $X$ is a
nonsingular irreducible projective analytic space of dimension $n$, and
if $Y$ is a nonsingular hyperplane section, then the inclusion
$$
Y^{n-1}\to X^n
$$
induces\pageoriginale a homomorphism
$$
\pi_i(Y^{n-1})\to \pi_i(X^n)
$$
that is an \emph{isomorphism for} $i\leq n-2$ and an
\emph{epimorphism for} $i=n-1$. The analogous statement follows for the
homomorphisms
$$
\H_i(Y^{n-1})\to \H_i(X^n)
$$
on homology (integral homology, to be definite), while in cohomology
$$
\H^i(X^n)\to\H^i(Y^{n-1})
$$
is an isomorphism in degree $i\leq n-2$ and a monomorphism in degree
$i=n-1$. We have obtained variants of these results in the setting of
schemes, for $\pi_0$, $\pi_1$, and $\Pic$, valid moreover, to a large
extent, without nonsingularity hypotheses; \cf \Exp \Ref{XII}.
Furthermore, in removing the nonsingularity hypotheses, we made
essential use of \og local\fg variants of these global Lefschetz
theorems. All this suggests the following problems, which will doubtless
have to be attacked simultaneously\sfootnote{The formulations
\Ref{XIII.4.1} to \Ref{XIII.4.3} below are provisional. See Conjectures
\ref{XIII.6.A} to \ref{XIII.6.D} below, in the \og Comments on
\Exp \Ref{XIII}\fg, for more satisfactory formulations, and also
\Exp \Ref{XIV}.}.

\begin{probleme} \label{XIII.4.1}
Let $X$ be an analytic space and $Y$ a closed analytic subset of $X$
(or merely a closed subset?)\nde{The meaning of this question is not
clear; indeed, the statement of the problem itself does not seem to make
sense in this case, since the codimension of $Y$ in $X$ is not defined
when $Y$ is no longer assumed analytic.} such that for every $x\in Y$,
the local ring $\OXx$ is a complete intersection. Let $n$ be the complex
codimension of $Y$ in $X$. Is the canonical homomorphism
$$
\pi_i(X-Y)\to \pi_i(X)
$$
an isomorphism for $i\leq n-2$ and an epimorphism for $i=n-1$?
\end{probleme}

In this problem and the \sisi{preceding}{following} ones, a base point
is of course implicitly understood to have been chosen in order to
define the homotopy groups. To state the following problem, one must
define, for an analytic space $X$ (more generally, for a locally
path-connected space) and a point $x\in X$, local invariants\pageoriginale
$\pi_i^x(X)$\sfootnote{If $i\geq2$. For the case $i\leq1$, \cf the
\emph{Comments} in \numero 6 below, page~\pageref{p25y}.}.
\refstepcounter{toto}\label{pagexpo} To do this, choose a nonconstant map
$f$ from the interval $[0,1]$ to $X$ such that $f(0)=x$ and $f(t)\neq x$
for $t\neq0$; such a map exists if $x$ is not an isolated point. Then,
for every neighborhood $U$ of $x$, there is an $\varepsilon>0$ such that
$0<t<\varepsilon$ implies $f(t)\in U$, and the homotopy groups
$\pi_i(U-x,f(t))$ are essentially independent of $t$: as $t$ varies,
they are connected by a transitive system of isomorphisms. They may be
denoted by $\pi_i(U-x,f)$. Set\refstepcounter{toto}\label{pip15}\label{pXIII.15}
$$
\pi_i^x(X)=\varprojlim_U\pi_{i-1}(U-x,f),
$$
where the inverse limit is taken over the system of open neighborhoods
$U$ of $x$. Strictly speaking, this limit depends on $f$ and should be
denoted $\pi_i^x(X,f)$, but one verifies that as $f$ varies these groups
are isomorphic to one another\sfootnote{At least if $x$ does not locally
disconnect $X$; \cf the \emph{Comments} below, page~\pageref{p25y}.};
more precisely, they form a local system on the space of paths of the
kind under consideration issuing from $x$. These invariants are the
homotopical counterpart of the local cohomology invariants $\H_x^i(F)$
for a sheaf $F$ on $X$, introduced in \Ref{I}, and should play the role
of \emph{relative local homotopy groups} of $X$ modulo $X-x$. Their
vanishing for $i\leq n$ and every $x\in Y$, where $Y$ is a closed subset
of $X$ of topological dimension $\leq d$, should imply that the
homomorphisms
$$
\pi_i(X-Y)\to\pi_i(X)
$$
are bijective for $i<n-d$ and surjective for
$i=n-d$\sfootnote{For a corrected formulation, \cf the \emph{Comments}
below, page~\pageref{p25y}.}. From this point of view,
\Ref{XIII.4.1} would imply, when $Y$ is reduced to a point, a conjecture
of purely local nature, expressed by\refstepcounter{toto}
$$
\pi_i^x(X)=0\ \rm{for}\ i\leq n-1\label{p15l}
$$
when $X$ is a complete intersection of dimension $n$ at $x$.

\sisi{As}{As} an example of the local invariants $\pi_i^x(X)$, note that
if $x$ is a nonsingular point of \sisi{\ignorespaces}{$X$ of} complex
dimension $n$, then\pageoriginale
$$
\pi_i^x(X)=\pi_{i-1}(S^{2n-1}),
$$
where $S^{2n-1}$ denotes the sphere of dimension $2n-1$. In particular,
in this case $\pi_i^x(X)=0$ for $i\leq2n-1$, which corresponds to the
fact that if a closed subset $Y$ of codimension $\geq m$ is removed from
a topological \emph{manifold} $X$, then
$\pi_i(X-Y)\to\pi_i(X)$ is an isomorphism for $i\leq m-2$ and an
epimorphism for $i=m-1$.

This being said:

\begin{probleme} \label{XIII.4.2}
Let $X$ be an analytic space, let $x\in X$, let $t$ be a section of
$\OX$ vanishing at $x$, and let $Y$ be the zero set of $t$. Suppose that
the following conditions are satisfied:
\begin{enumeratea}
\item
$t$ is regular at $x$ (\ie it is not a zero divisor at $x$; this
hypothesis may moreover be superfluous).
\item
At the points $x'$ of $X-Y$ near $x$, $\Oo_{X,x'}$ is a complete
intersection. If \Ref{XIII.4.1} is true, this hypothesis should be
replaceable by the following more general one: for $x'$ as above,
$\pi_i^{x'}(X)=0$ for $i\leq n-1$.
\item
At the points $y$ of $Y-\{x\}$ near $x$, one has
$$
\prof\Oo_{X,y}\geq n
$$
(it suffices, for example, that $\prof\Oo_{X,x}\geq n$).
\end{enumeratea}

Under these conditions, is the canonical homomorphism
$$
\pi_i^x(Y)\to\pi_i^x(X)
$$
an isomorphism for $i\leq n-2$ and an epimorphism for $i=n-1$?
\end{probleme}

Here, finally, is a global variant of \Ref{XIII.4.2}, which should
follow from it by considering the affine cone at its vertex and
which would generalize the classical Lefschetz theorems:

\begin{probleme} \label{XIII.4.3}
Let\pageoriginale $X$ be a projective analytic space, equipped with an
ample invertible module $L$; let $t$ be a section of $L$, and let $Y$ be
the zero set of $t$. Suppose that:
\begin{enumeratea}
\item
$t$ is a regular section (\sisi{\ignorespaces}{a hypothesis} that may
be superfluous).
\item
For every $x\in X-Y$, $\Oo_{X,x}$ is a complete intersection. This
should be replaceable by $\pi_i^x(X)=0$ for $i\leq n-1$.
\item
For every $x\in Y$, $\prof\Oo_{X,x}\geq n$.
\end{enumeratea}

Under these conditions, is the homomorphism
$$
\pi_i(Y)\to\pi_i(X)
$$
an isomorphism for $i\leq n-2$ and an epimorphism for $i=n-1$?
\end{probleme}

We leave it to the reader to state analogous conjectures of a
\emph{cohomological} nature\sfootnote{\Cf \Exp \Ref{XIV} for the
corresponding results in the theory of schemes.}, with hypotheses and
conclusions concerning local cohomological invariants (with coefficients
in a given group). In any event, the key result seems to be
\Ref{XIII.4.2}, with hypothesis b) taken in the more general form just
indicated, whether one adopts the point of view of homology or that of
homotopy.

We have stated these conjectures in the transcendental setting in the
hope of interesting topologists in them and convincing them that
\og Lefschetz\fg-type questions are far from closed. Of course, now that
we are about to possess a good theory of the cohomology of schemes (with
finite coefficients), thanks to the recent work of M.~\sisi{ARTIN}{Artin},
the same questions arise in the setting of schemes, and it is difficult
to doubt that they will receive an affirmative answer in the near
future\sfootnote{See the preceding note.}.

\section{Problems related to local Picard groups} \label{XIII.5}

A\pageoriginale first fundamental problem, pointed out for the first time by \sisi{MUMFORD}{Mumford}~\cite{XIII.5} in a particular case, is the following. Let $A$ be a complete local ring with residue field $k$, let $X=\Spec(A)$, and let $U=\Spec(A)-\{a\}$, where $a$ is the maximal ideal of $A$, \ie the closed point of $\Spec(A)$. The problem is to construct a strict projective system $G$ of locally algebraic groups $G_i$ over $k$, and a natural isomorphism
\begin{equation*} \label{eq:XIII.5.+} \tag{$+$} {\Pic(U)\simeq G(k)} \end{equation*}
where, of course, we set $G(k)=\varprojlim G_i(k)$. Heuristically, the problem is to \og put an algebraic-group structure\fg (or at least a pro-algebraic one, in a suitable sense) on the group $\Pic(\sisi{X'}{U})$.

It is clear that, as stated, the problem is not sufficiently precise, since the datum of an isomorphism~(\Ref{eq:XIII.5.+}) is far from characterizing the pro-object $G$. If $A$ contains a subfield, again denoted by $k$, that is a coefficient field, the problem can be made more precise by requiring that, for a varying extension $k'$ of $k$, there be an isomorphism, functorial in $k'$:
\begin{equation*} \label{eq:XIII.5.+'} \tag{$+'$} {\Pic(U')\simeq G(k')} \end{equation*}
where $U'$ is the open subset analogous to $U$ in $\Spec(A')$, with $A'=A\hat{\otimes}_k k'$. One can proceed analogously even if $A$ has no coefficient field, provided that $k$ is perfect, since one can then construct an $A'$ functorially \og by residue-field extension $k'/k$\fg. Moreover, when $A$ admits a coefficient field, the algebraic structure obtained on $\Pic(U)$ will depend essentially on the choice of that coefficient field (as one sees already for the projective cone over an elliptic curve); it therefore seems necessary to start from a \og pro-algebraic ring over $k$\fg in the sense of \sisi{GREENBERG}{Greenberg}~\cite{XIII.3}, in order to define the pro-object $G$. It is moreover conceivable that, when no coefficient field is given, one obtains only a projective system of\pageoriginaled \emph{quasi}-algebraic groups in the sense of Serre, or rather quasi-locally algebraic groups (the groups $G_i$ obtained will not in general be of finite type over $k$, but only locally of finite type over $k$). It is even possible that in general one obtains only a still weaker structure on $\Pic(U)$, of the kind encountered by \sisi{NERON}{N\'eron}~\cite{XIII.6} in his theory of degeneration of abelian varieties defined over local fields.

One method of attacking the problem, likewise introduced by \sisi{MUMFORD}{Mumford}, consists in resolving the singularities of $X$, \ie in considering a projective birational morphism $Y\to X$ with $Y$ regular. When $U$ is regular (\ie $a$ is an isolated singular point), one can often find $Y$ such that $Y\sisi{|U}{_|U}=V\to U$ is an isomorphism. In that case one therefore has
$$
\Pic(U)\simeq\Pic(V)\simeq\Pic(Y)/\Im \ZZ^I,
$$
where $I$ is the set of irreducible components of the fiber $Y_{a}$\sisi{, }{} (each of which defines an element of $\Pic(Y)$, since it is a locally principal divisor, by the regularity of $Y$). On the other hand, using the techniques of Formal Geometry in \EGA III~4~and~5, in particular the existence theorem, one finds
$$
\Pic(Y)\simeq\varprojlim\Pic(Y_n),
$$
where $Y_n=Y\otimes_A A_n$ and $A_n=A/\mm^{n+1}$. When $A$ admits a coefficient field $k$, the theory of Picard schemes of the projective $k$-schemes $Y_n$ is available, and hence
$$
\Pic(Y_n)\simeq \mathbf{Pic}_{Y_n/k}(k).
$$

This therefore gives a construction of a projective system of locally algebraic groups $\mathbf{Pic}_{Y_n/k}/\Im \ZZ^I$, which is the desired system.\refstepcounter{toto}\nde{\label{Boutot}The question was greatly clarified by the results of {Boutot} (Boutot~J.-F., \emph{Sch\'ema de {P}icard local}, Lect. Notes in Math., vol.~632, Springer, Berlin, 1978). In particular, if $A$ is a complete (noetherian) local $k$-algebra of depth $\geq 2$ such that $H^2_\m(A)$ has finite dimension over $k$, the local Picard group is a group scheme locally of finite type over $k$, whose tangent space at the origin is $H^2_\m(A)$. If $A$ is moreover normal of dimension $\geq 3$, Serre's normality criterion~\Ref{XI}~\Ref{XI.3.11}, together with Corollary~\Ref{V}~\Ref{V.3.6}, ensures the required finiteness and hence the existence of the local Picard scheme. See also (Lipman~J., {\og The Picard group of a scheme over an Artin ring\fg}, \emph{Publ. Math. Inst. Hautes \'Etudes Sci.} \textbf{46} (1976), p\ptbl 15--86) for an approach closer to Grothendieck's outline above.} In the case considered here, one can moreover see (using the fact that $a$ is an isolated singular point) that the connected components of the universal-image subgroups\pageoriginale in this projective system form an \emph{essentially constant} projective system; thus, in this case, one obtains a locally algebraic group $G$ as a solution of the problem. If one further assumes that $A$ is normal of dimension $2$, then a remark of \sisi{MUMFORD}{Mumford} (stating that the intersection matrix of the components of $Y_{a}$ in $Y$ is negative definite\nde{Mumford~D., {\og The topology of normal singularities of an algebraic surface and a criterion for simplicity\fg}, \emph{Publ. Math. Inst. Hautes \'Etudes Sci.} \textbf{9} (1961), p\ptbl 5--22.}) implies that $G$ is even an algebraic group, \ie of finite type over $k$ (the number of its connected components being equal, moreover, to the determinant of the intersection matrix just considered).

If, on the contrary, $a$ is not an isolated singularity, examples (with $A$ of dimension $2$) show that one obtains a projective system of algebraic groups which does not reduce to a single algebraic group.

\refstepcounter{toto}\label{pXIII20}%
Once a good notion of \og \emph{local Picard scheme}\fg were available, one should strengthen the notion of parafactoriality by saying that $A$ is \og \emph{geometrically parafactorial}\fg when not only $A$, and even $\widehat{A}$, is parafactorial, but the local Picard scheme $G(\widehat{A})$ is the trivial group (which, when the residue field is not algebraically closed, is stronger than saying that $G$ has no rational point over $k$ other than the identity). The necessity of a strengthened notion of parafactoriality becomes apparent on recalling that there exist complete normal local rings of dimension $2$ which are factorial but admit finite \'etale algebras which are not\nde{Factorial rings with nonfactorial henselization arise naturally in the study of moduli spaces of vector bundles. See, for example, (Dr\'ezet~J.-M., {\og Groupe de Picard des vari\'et\'es de modules de faisceaux semi-stables sur $\mathbf{P}\sb 2$\fg}, in \emph{Singularities, representation of algebras, and vector bundles (Lambrecht, 1985)},
Lect. Notes in Math., vol.~1273, Springer, Berlin, 1987, p\ptbl 337--362). \textit{Strictly speaking}, Dr\'ezet shows that the completion is not factorial, but in fact the proof gives the result for the henselization: the point is that Luna's \'etale slice theorem (Luna~D., {\og Slices \'etales\fg}, in \emph{Sur les groupes alg\'ebriques}, M\'em. Soc. math. France, vol.~33, Soci\'et\'e math\'ematique de France, Paris, 1973, p\ptbl 81--105) describes the local ring of a quotient in the sense of invariant theory near a semistable point, locally for the \'etale topology.}. A \og \emph{geometrically factorial}\fg local ring would then be a normal ring~$A$ such that all \sisi{t}{th}e localizations of dimension $\geq 2$ are geometrically parafactorial, or better, such that the localizations of $\widehat{A}$ are parafactorial\sfootnote{For a more flexible notion of \og geometrically factorial\fg local ring, \cf \emph{Comments}, page~\pageref{commpara}.}. Of course, it would be interesting to find a \og good\fg definition of these notions, independent of the still undeveloped theory of local Picard schemes.\nde{See page~\pageref{commpara}: a local ring is geometrically factorial (\resp parafactorial) if its strict henselization is factorial (\resp parafactorial).}

It is in any event plausible that these notions will be needed if one wishes to obtain statements of the following type: Let $A$ be a \og good ring\fg (for example an algebra of finite type over $\ZZ$, or over a complete local ring, for example over a field). Let~$U$ be the set of $x\in X=\Spec(A)$ such that $\OXx$ is \og geometrically\pageoriginale factorial\fg; is $U$ open\sisi{?.}{?} Or again: Let $f\colon X\to Y$ be a flat morphism of finite type with $Y$ locally noetherian, and let $U$ be the set of $x\in X$ such that $\Oo_{X_{f(x)}, x}$ is \og geometrically factorial\fg; is $U$ open, at least under suitable additional hypotheses on $f$\sisi{?.}{?} I doubt that true statements of this type exist with the usual notion of factorial ring.

We have raised here, in a particular case, the question of studying the geometric properties of \og varying\fg local rings, for example the $\OXx$ as $x$ ranges over a prescheme $X$. When $X$ is, for example, a scheme of finite type over a field, it is known\nde{See \EGA IV.16.} that there exists on $X$ a projective system of finite algebras $P^n_{X/k}$ (obtained by completing $X\times_k X$ along the diagonal), whose fiber at every point $x\in X$ rational over $k$ is isomorphic to the projective system of the $\OXx/\mm_X^{n+1}$. It is then natural to relate the study of the completions of the local rings $\OXx$, for varying $x$, to that of the \og algebraic family of complete local rings\fg given by the $P^n$, noting that for every $x\in X$ (rational over $k$ or not), one obtains a complete local ring
$$
P^{\infty}(x)=\varprojlim P^n(x)
$$
(where $P^n(x)=$~reduced~fiber~$P^n\otimes_{\OXx}k(x)$). Of particular interest, for example, will be the complete ring thus associated with the generic point, and one expects its algebro-geometric properties (expressed, for example, by its local Picard, homotopy, or homology groups) to be essentially those of the completions $\widehat{\Oo}_{X,x}$ for $x$ in a suitable dense open subset $U$.

More generally, one may undertake the simultaneous study of the complete local rings obtained in this way from an adic projective system $(P_n)$ of finite algebras over a given scheme $X$. It is plausible that, under certain regularity conditions (such as flatness of the $P_n$),\pageoriginale the local homotopy groups arise from a projective system of finite group schemes over $X$, and that analogous results hold for the local Picard groups. As regards the \sisi{latter}{latter}, a first interesting case deserving investigation is that in which one starts with an algebraic surface $X$ having singular curves, and seeks to study the local Picard schemes at varying points of these curves in terms of a suitable pro-group scheme defined over the singular locus.

\section[Comments]{Comments\protect\sfootnotemark} \label{XIII.6}

The\pageoriginale viewpoint of\sfootnotetext{Written in March~1963.} the \og \'etale cohomology\fg of schemes, and recent progress in this theory, lead\sisi{}{s} us both to make some of the problems posed more precise and to broaden them. For the notion of \og topology\fg and of the \og \'etale topology of a scheme\fg, I refer to M\ptbl\sisi{ARTIN}{Artin}, \sisi{Grothendieck Topologies}{\emph{Grothendieck Topologies}}, Harvard University 1962 (mimeographed notes)\sfootnote{Or preferably to \SGA 4.}.

By providing a finer notion of localization than that furnished by the traditional \og Zariski topology\fg, this theory gives particular importance to \emph{strictly local rings}\refstepcounter{toto}\label{pXIII.6}, \ie henselian local rings with separably closed residue field. For every local ring $A$ with residue field $k$, and every separable closure $k'$ of $k$, one can find a local homomorphism from $A$ to a strictly local ring $A'$, the \emph{strict local closure} of $A$, with residue field $k'$, having an evident universal property. The ring $A'$ is henselian and flat over $A$, and $A'\otimes_A k\simeq k'$; it is noetherian if and only if $A$ is. (\Cf \loccit Chap\ptbl III, section 4)\sfootnote{Or \EGA IV~18.8.}. If $X$ is a prescheme, $x$ a point of $X$, and $x'$ a point above $x$ which is the spectrum of a separable closure $k'$ of $k=k(x)$, one is led to define the \emph{strictly local ring} of $X$ at $x'$, $\Oo'_{X,x'}$, as the strict local closure of the usual local ring $\OXx$, relative to the residue-field extension $\sisi{K'/K}{k'/k}$. It is the \emph{strictly} local rings at the \og geometric\fg points of $X$ which, from the viewpoint of the \'etale topology, are meant to reflect the local properties of the prescheme $X$. They also play, in many respects, the role formerly assigned to the \emph{completions} of the local rings of $X$ (say, at points with algebraically closed residue field), while remaining \og closer\fg to $X$ and allowing an easier passage to \og neighboring points\fg.

One should therefore reconsider for noetherian henselian local rings (\resp noetherian strictly local rings) many questions generally posed for complete local rings (possibly restricted\pageoriginale to those having an algebraically closed residue field). Thus the topological problems raised in n$^{\textup{os}}$~\Ref{XIII.2} and \Ref{XIII.3} arise more generally for strictly local rings. One can moreover state conjecturally, for \og good\fg strictly local rings, certain simple-connectedness and acyclicity properties for the geometric fibers of the canonical morphism $\Spec(\widehat{A})\to\Spec(A)$; these would show that, for many properties of a \og topological\fg nature, it amounts to the same thing to prove them for the ring $A$ or for its completion $\widehat{A}$. Certain results already obtained in this direction\sfootnote{\Cf M\ptbl\sisi{ARTIN}{Artin} in \SGA 4~XIX} give reason to hope that complete results in this direction will soon be available.

\refstepcounter{toto}\label{commpara}%
The notion of \'etale localization provides a seemingly reasonable definition of a \og \emph{geometrically parafactorial}\fg or \og \emph{geometrically factorial}\fg local ring\sisi{,}{} (the need for which was indicated in \numero\Ref{XIII.5}, p\ptbl\sisi{20}{\pageref{pXIII20}}): this will mean a local ring whose strict local closure is parafactorial, \resp factorial. Hypotheses of this nature arise naturally in the study of the \'etale cohomology of preschemes\nde{See, for example, (Strano~R., {\og The Brauer group of a scheme\fg}, \emph{Ann. Mat. Pura Appl. (4)} \textbf{121} (1979), p\ptbl 157--169), where the hypothesis that the local rings of a scheme $X$ are geometrically parafactorial sometimes makes it possible to prove that the Brauer groups of $X$ (computed in terms of Azumaya algebras) coincide with the cohomological Brauer group of $X$.}. Thus, if $X$ is a locally noetherian prescheme whose strictly local rings are factorial (\ie whose ordinary local rings are \og geometrically factorial\fg), one proves that the $\H^i(X_{\et},\GG_m)$ are torsion groups for $i\geq 2$ (which sometimes makes it possible to express these groups in terms of cohomology groups with coefficients in the groups $\bbmu_n$ of $n$\sisi{th}{-th} roots of unity), and if $X$ is integral with fraction field $K$, the natural homomorphism $\H^2(X_{\et},\GG_m)\to\H^2(K,\GG_m)=\Br(K)$ is injective\nde{The connection between the Brauer group and the Picard group is a close one. In this connection, let us cite the following results of Saito (Saito~S., {\og Arithmetic on two-dimensional local rings\fg}, \emph{Invent. Math.} \textbf{85} (1986), \numero 2, p\ptbl 379--414) in the case of surfaces, the first local and the other global. Let $A$ be an excellent normal henselian local ring of dimension $2$ with \textit{finite} residue field, and let $X$ be the complement of the closed point in $\Spec(A)$. Then there is a perfect duality of \emph{torsion} groups $\Pic(X)\times\Br(X)\to\QQ/\ZZ$ --- by the Brauer group of $X$ we mean the \emph{cohomological} Brauer group $\textrm{Br}(X)=H^2_\et(X,G_m)$. In the global case one has the following generalization of a result of Lichtenbaum (Lichtenbaum~S., {\og Duality theorems for curves over $p$-adic fields\fg}, \emph{Invent. Math.} \textbf{7} (1969), p\ptbl 120--136): let $k$ be the fraction field of a complete discrete valuation ring $\Oo$ with finite residue field, and let $X$ be a projective, smooth, and geometrically complete curve over $k$. The group $\Pic^0(X)$ is endowed with the topology induced by the adic topology of $k$, and $\Pic(X)$ is the topological group for which $\Pic^0(X)$ is an open subgroup. Then there is a perfect duality of topological groups $\Pic(X)\times\Br(X)\to\QQ/\ZZ$. Note that this statement, which concerns curves, is of course proved by considering a regular (proper and flat) model of $X$ over $\Oo$: it is a result about surfaces.}; examples show that these conclusions may fail, even for local $X$, if one assumes only that $X$ is factorial rather than geometrically factorial\sfootnote{\Cf A\ptbl \sisi{GROTHENDIECK}{Grothendieck}, le groupe de Brauer~II (S\'eminaire Bourbaki \numero 297, Nov.~1965), notably 1.8 and 1.11~b.}.

As regards the problems of local and global Lefschetz type raised in \Ref{XIII.3.4}, and their analogues in scheme theory, the homological version of these questions has been considerably clarified: everything follows formally from three general theorems. The first concerns the cohomological dimension of\pageoriginale certain affine schemes (\resp Stein spaces), such as affine schemes $X$ of finite type over an algebraically closed field: their cohomological dimension is $\leq\dim X$ (\og affine Lefschetz theorem\fg)\sfootnote{\Cf \SGA 4~XIV.}; the second is a duality theorem for the cohomology (with discrete coefficients) of a projective morphism\sfootnote{\Cf \SGA 4~XVIII.}; and the last is a \emph{local duality} theorem of an analogous nature\sfootnote{\Cf \SGA 5~I.}. In Algebraic Geometry, only the last of these has not been proved at the time of writing (it is, however, proved in \sisi{char.~$0$}{characteristic $0$}, using resolution of singularities by \sisi{HIRONAKA}{Hironaka}). Moreover, in the transcendental setting, global and local duality are already available, having recently been proved by \sisi{VERDIER}{Verdier}\nde{See Verdier~J.-L., {\og Dualit\'e dans la cohomologie des espaces localement compacts\fg}, in \emph{S\'eminaire Bourbaki}, vol.~9, Soci\'et\'e math\'ematique de France, Paris, 1995, \Exp 300, p\ptbl 337--349.}. Let us merely indicate that, in the statements of the homological versions of Problems \Ref{XIII.4.2} and \Ref{XIII.4.3} (which henceforth deserve the name of conjectures), the \og at infinity\fg conditions a) and c) are certainly superfluous; only the \emph{local cohomological structure} of $X-Y$ is important, and one may assume, for example, that it is locally a complete intersection of dimension $\geq n$. Moreover, in \Ref{XIII.4.3}, say, the fact that $Y$ is a hyperplane section should play no role and should be replaceable by the sole hypothesis that $X$ is compact and $X-Y$ is Stein (\ie in Algebraic Geometry, $X$ is proper over $k$ and $X-Y$ is affine; as we said, the homological version of this conjecture is proved for algebraic spaces over the field $\CC$)\sfootnote{\Cf \Exp \Ref{XIV}.}.

\refstepcounter{toto}\label{p25y}%
In the definition (p\ptbl\sisi{15}{\pageref{pip15}}) of the $\pi_i^x(X)$, one must assume $i\geq 2$. For $i=0,1$, there is no reasonable definition of the $\pi_i^x(X)$\sisi{,}{;} they should be replaced by $\H_0^x(X)$ and $\H_1^x(X)$, defined respectively as the co\sisi{-}{}kernel and kernel in the natural homomorphism
$$
\varprojlim \H_0(U-(x),\ZZ)\to\varprojlim\H_0(U,\ZZ).
$$

\sisi{At}{At} a pinch, and for convenience of formulation, one may set $\pi_i^x(X)=\H_i^x(X)$ for $i\leq 1$; otherwise, subsequent assertions concerning the $\pi_i^x$ must be supplemented by the corresponding assertions for $\H_0^x,\H_1^x$. If $x$ is an isolated point of $X$, one should set $\pi_i^x(X)=0$ for $i\neq 0$ and $\pi_0^x(X)=\H_0^x(X)=\ZZ$.

The assertion\pageoriginale that the $\pi_i^x(U,f)$ are mutually isomorphic is true only when $X$ is not disconnected by $x$ in a neighborhood of $x$, \ie if $\pi_i^x(X)=0$ for $i=0,1$. In the general case $\pi_i^x(X)$ can designate only a \emph{family} of groups, not necessarily mutually isomorphic; nevertheless the notation $\pi_i^x(X)=0$ retains an evident meaning.

\refstepcounter{toto}\label{p26y}\sisi{Page~15, line~-~11}{Page~\pageref{p15l}}, where I predict that the vanishing of the local homotopy invariants $\pi_i^x(X)$ for $x\in Y$, $i\leq n$, should imply that $\pi_i(X-Y)\to\pi_i(X)$ is bijective for $i<n-d$ and surjective for $i=n-d$, one must be cautious, since in the present context (as in Algebraic Geometry) one does not have at one's disposal \og general\fg points at which the local conditions would also have to apply. For this reason it will doubtless be necessary to use relative local homotopy invariants
$$
\pi_i^Y(X,f)=\pi_i^Y(X,x)=\varprojlim_U\pi_{i-1}(U-U\cap Y,f(t))\qquad\text{for }i\geq 2\quoi,
$$
(with an {\sisi{ad-hoc}{\textit{ad hoc}}} definition as above for $i=0,1$), where $Y$ is a closed subset of~$X$; or else to compensate for the absence of general points by expressing the hypotheses on $X$ in terms of properties, topological in nature (for the \'etale topology), of the spectra of the local rings of $X$, which makes it possible to recover general points. The same reservation applies to the generalization of Conjectures \Ref{XIII.4.2} and \Ref{XIII.4.3} to the case in which $X-Y$ is not assumed to be locally a complete intersection, a generalization suggested in the statements of conditions~b) of those conjectures.

To formulate the expurgated versions of Conjectures \Ref{XIII.4.2} and \Ref{XIII.4.3} suggested by the results mentioned above, we make the following

\begin{supdefinition} \label{XIII.6.1}
Let $X$ be a topological space, let $Y$ be a locally closed subset of $X$, and let $n$ be an integer. One says that $X$ has \emph{homotopical depth} $\geq n$ along $Y$, and writes $\mathrm{prof\:hpt}_Y(X)\geq n$, if for every $x\in Y$ one has $\pi_i^Y(X,x)=0$ for $i<n$.
\end{supdefinition}

It should\pageoriginale be equivalent to say that, for every open subset $X'$ of $X$ and every $x\in X'\cap U=U'$ (where $U=X-Y$), the canonical homomorphism
$$
\pi_i(U',x)\to\pi_i(X',x)
$$
is an isomorphism for $i<n-1$ and a monomorphism for $i=n-1$\nde{When the pair $(X,Y)$ is moreover polyhedral, this equivalence holds; \cf (Eyral~C., {\og Profondeur homotopique et conjecture de Grothendieck\fg}, \emph{Ann. Sci. \'Ec. Norm. Sup. (4)} \textbf{33} (2000), \numero 6, p\ptbl 823--836).}.

\begin{supdefinition} \label{XIII.6.2}
Let $X$ be a complex analytic space and let $n$ be an integer. One says that the \emph{rectified homotopical depth} of $X$ is $n$\sfootnote{In the first edition of these notes, we used the term \og true homotopical depth\fg. In the present version we follow \EGA IV~10.8.1.} if, for every locally closed analytic subset $Y$ of $X$, one has
\begin{equation*} \label{eq:XIII.6.x} \tag{$x$} {\mathrm{prof\:hpt}_Y(X)\geq n-\dim Y}
\end{equation*}
(where, of course, $\dim Y$ denotes the \emph{complex} dimension of $Y$).
\end{supdefinition}

It should be equivalent to say that, for every irreducible locally closed analytic subset $Y$ of $X$, there exists a closed analytic subset $Z$ of $Y$, of dimension $<\dim Y$, such that relation~(\Ref{eq:XIII.6.x}) holds with $Y-Z$ in place of $Y$. This would make it possible, for example, to restrict in Definition~\Ref{XIII.6.2} to the case where $Y$ is nonsingular.\refstepcounter{toto}\nde{\label{HL}All the conjectures that follow, suitably rectified, if I may say so, became theorems through the work of Hamm and L\^e D\~ung Tr\'ang (Hamm~H.A. \& L\^e D\~ung Tr\'ang, {\og Rectified homotopical depth and Grothendieck conjectures\fg}, in \emph{The Grothendieck Festschrift, Vol. II}, Progr. Math., vol.~87, Birkh\"auser, Boston, 1990, p\ptbl 311--351), cited as \cite{HL} below. As for the two conjecturally equivalent definitions of rectified depth, they are in fact equivalent to a third, expressed in terms of Whitney stratification (\cf \loccit, Theorem 1.4).}

\enlargethispage{\baselineskip}%
The following conjecture, purely topological in nature, is in the spirit of a \og \emph{local Hurewicz theorem}\fg.

\refstepcounter{aconjecture} \label{XIII.6.A}
\begin{enonce*}{Conjecture A}[\og local Hurewic\sisi{s}{z} \sisi{th.}{theorem}\fg\ndemark]\ndetext{As observed in \cite{HL}, Example 3.1.3, this conjecture is already false for $X=\{\boldsymbol{z}\in\nobreak\CC^n\mid\nobreak z_1^2+z_2^3+\cdots+z_n^3=0\}, n\geq 4$, and $Y$ reduced to the origin. Suitably modified, however, it is true (Theorem 3.1.4 of {\loccit}).}
Let $X$ be a topological space and $Y$ a locally closed subset, subject if necessary to \og smoothness\fg conditions such as local triangulability of the pair $(X,Y)$, and let $n$ be an integer $\geq 3$. In order that $\mathrm{prof\:hpt}_Y(X)\geq n$, it is necessary and sufficient that\pageoriginale
$$
\H^i_Y(\ZZ_X)=0\quad\text{for }i<n
$$
(one then says that $X$ has \emph{cohomological depth} $\geq n$ along $Y$), and that the local fundamental groups
$$
\sisi{\pi^2_Y}{\pi_2^Y}(X,x)=\varprojlim_{U\ni x}\pi_1(U-U\cap Y)
$$
vanish (one then says that $X$ is \og pure\fg along $Y$).
\end{enonce*}

Note that if $X$ is an analytic space, $Y$ an analytic subspace, and $X$ is pure along $Y$, then for every $x\in Y$ the local ring $\OXx$, as well as its localizations with respect to prime ideals containing the ideal defining the germ $Y$ at $x$ (\ie in the inverse image $Y_x$ of $Y$ under $\Spec(\OXx)=X_x\to X$), are pure in the sense of \Exp \Ref{X}; it seems plausible that the converse is also true. Analogous remarks apply to cohomological depth, it being understood that one works with the \'etale topology on the $\Spec(\OXx)$.

Conjecture \Ref{XIII.4.1} then generalizes to

\refstepcounter{aconjecture} \label{XIII.6.B}
\begin{enonce*}{Conjecture B}[\og Purity\fg\ndemark]\ndetext{This conjecture is proved, even when $E$ is singular, in \cite{HL}: it is Theorem 3.2.1.}
Let $E$ be an analytic space and $X$ an analytic subset of $E$. Suppose that $E$ is nonsingular of dimension $N$ at $x\sisi{}{\in X}$, and that $X$ can be described by $p$ analytic equations in a neighborhood of every point. Then the rectified homotopical depth of $X$ is $\geq N-p$.
\end{enonce*}

In particular, a local complete intersection of dimension $n$ at every point would have rectified homotopical depth $\geq n$, which is precisely Conjecture \Ref{XIII.4.1}.

Conjectures \Ref{XIII.4.2} and \Ref{XIII.4.3} generalize respectively to:

\refstepcounter{aconjecture} \label{XIII.6.C}
\begin{enonce*}{Conjecture C}[\og Local Lefschetz\fg\ndemark]\ndetext{This conjecture is proved in \cite{HL}, indeed even in the strong form of the remark that follows; \cf Theorem 3.3.1 of~\loccit.}
Let $X$ be an analytic space, let $Y$ be a closed analytic subset, and let $x$ be a point of $Y$. Suppose that $X-Y$ is Stein in a neighborhood of $x$ (for example, that $Y$ is defined by one equation at $x$), and that $X-Y$ has rectified homotopical depth $\geq n$ in a neighborhood of $x$ (for example, that at every point of $X-Y$ near $x$ it is a complete intersection of dimension $\geq n$; \cf Conjecture \Ref{XIII.6.B}). Then the canonical homomorphism
$$
\pi_i^x(Y)\to\pi_i^x(X)
$$
is an isomorphism for $i<n-1$ and an epimorphism for $i=n-1$.
\end{enonce*}

\refstepcounter{aconjecture} \label{XIII.6.D}
\begin{enonce*}{Conjecture D}[\og Global Lefschetz\fg\ndemark]
\ndetext{This conjecture is again proved in \cite{HL}, indeed even in the strong form of the remark that follows; \cf Theorem 3.4.1 of~\loccit.}
Let\pageoriginaled $X$ be a \emph{compact} analytic space and let $Y$ be an analytic subspace of $X$ such that $U=X-Y$ is Stein and has rectified homotopical depth $\geq n$ (for example, is a complete intersection of dimension $\geq n$ at every point). Then the canonical homomorphism
$$
\pi_i(Y)\to\pi_i(X)
$$
is an isomorphism for $i<n-1$ and an epimorphism for $i=n-1$.
\end{enonce*}

\begin{remarquestar}
When, in the statements \Ref{XIII.6.C} and \Ref{XIII.6.D}, the hypothesis that $X-Y$ is Stein is replaced by the hypothesis that $X-Y$ is a union of $c+1$ Stein open subsets (which plays the role of a topological \og concavity\fg hypothesis), the conclusions should simply be modified by replacing $n$ there by $n-c$.\nde{Finally, let us mention the following result of Fulton, to be compared with the Fulton--Hansen result cited in the editor's note~\eqref{FH} on page~\pageref{FH}: let $X$ and $H$ be closed subschemes of $\PP^m_\CC$, let $n$ be the dimension of $X$, and let $d$ be the codimension of $H$. Then the map
$$
\pi_i(X,X\cap H)\to\pi_i(\PP^m_\CC,H)
$$
is an isomorphism if $i\leq n-d$ and is surjective if $i=n-d-1$; see (Fulton~W. \& Lazarsfeld.~R., {\og Connectivity and its applications in algebraic geometry\fg}, in \emph{Algebraic geometry (Chicago, Ill., 1980)}, Lect. Notes in Math., vol.~862, Springer, Berlin--New York, 1981, p\ptbl 26--92).}
\end{remarquestar}

Finally, let us make explicit, in the \og global case\fg \Ref{XIII.6.D}, the conjecture concerning the fundamental group (obtained by taking $n=3$):

\refstepcounter{aconjecture}\label{XIII.6.D'}
\begin{enonce*}{Conjecture D$\boldsymbol'$}[Global Lefschetz for the fundamental group\ndemark]
\ndetext{This conjecture is proved in \cite{HL}; \cf Theorem 3.5.1 of~\loccit.}
Let $X$ be a compact analytic space over the field of complex numbers and let $Y$ be a closed analytic subset such that $U=X-Y$ is Stein. Suppose moreover that the following conditions hold:
\begin{enumeratei}
\item
For every $x\in U$, the local fundamental group $\pi_2^x(X,x)$ vanishes (\ie $X$ \emph{is \og pure at }$x$\fg), or merely the local ring $\OXx$ is pure.
\item
The local rings at the points of $U$ are \og connected in dimension $\geq 2$\fg.
\item
The local rings at the points of $U$ have dimension $\geq 3$.
\end{enumeratei}

Under these conditions, for every $x\in Y$, the homomorphism
$$
\pi_1(Y,x)\to\pi_1(X,x)
$$
is\pageoriginale an isomorphism (and $\pi_2(Y,x)\to\pi_2(X,x)$ is an epimorphism).
\end{enonce*}

Note that the local conditions (i), (ii), and (iii) on $U$ hold if $U$ is locally a complete intersection of dimension $\geq 3$. From the viewpoint of Algebraic Geometry (when $U$ arises from a scheme, again denoted by $U$), conditions (i) through (iii) correspond to hypotheses on the local invariants $\pi_i^x(U)$, \sisi{\ignorespaces}{namely} \hbox{$\pi_i^x(U)=0$} for $i<3-\degtr k(x)/k$, at points $x$ for which one has respectively $\degtr k(x)/k=0,1,2$. The global condition on $U$ ($U$ Stein) is satisfied if $X$ is projective and $Y$ is a hyperplane section.

\begin{thebibliography}{HL}\refstepcounter{toto}\pageoriginale\label{XIII.7}

\bibitem[1]{XIII.1} \sisi{%
S. Abhyankar, Local uniformisation on algebraic surfaces, Annals of Math. vol.~63, \No~3, 1956}
{%
\textsc{S. Abhyankar} -- {\og Local uniformisation on algebraic surfaces over ground fields of characteristic $p\neq0$\fg}, \emph{Annals of Math.} \textbf{63} (1956), p\ptbl 491--526.
}

\bibitem[2]{XIII.2} \sisi{%
W.L. Chow, On the theorem of Bertini for local domains, Proc. Nat. Acad. Sc. Vol.~44, \No~6, pp\ptbl 580--584, 1958}
{%
\textsc{W.L. Chow} -- {\og On the theorem of Bertini for local domains\fg}, \emph{Proc. Nat. Acad. Sci. U.S.A.} \textbf{44} (1958), p\ptbl 580--584.
}

\bibitem[3]{XIII.3} \sisi{%
M. Greenberg, Schemata over local rings, Annals of Math., vol~73, 1961}
{%
\textsc{M. Greenberg} -- {\og Schemata over local rings\fg}, \emph{Annals of Math.} \textbf{73} (1961), p\ptbl 624--648.
}

\bibitem[4]{XIII.4} \sisi{%
M. Kneser, Uber die Darstellung algebraischer Raumkurven als Durchschnitte von Fl\"achen, Archiv der Math. Vol.~XI, fasc.~3, 1960}
{%
\textsc{M. Kneser} -- {\og \"Uber die Darstellung algebraischer Raumkurven als Durchschnitte von Fl\"achen\fg}, \emph{Archiv der Math.} \textbf{XI} (1960), p\ptbl 157--158.
}

\bibitem[5]{XIII.5} \sisi{%
D. Mumford, The topology of normal singularities of an algebraic surface, and a criterion for simplicity, Publications Math\'ematiques \No~9, 1961}
{%
\textsc{D. Mumford} -- {\og The topology of normal singularities of an algebraic surface, and a criterion for simplicity\fg}, \emph{Publ. Math. Inst. Hautes \'Etudes Sci.} \textbf{9} (1961), p\ptbl 5--22.
}

\bibitem[6]{XIII.6} \sisi{%
A. N\'eron, Mod\`eles minimaux des vari\'et\'es ab\'eliennes sur les corps locaux et globaux, to appear in Publications Math\'ematiques (available as mimeographed notes in limited circulation...)}
{%
\textsc{A. N\'eron} -- {\og Mod\`eles minimaux des vari\'et\'es ab\'eliennes sur les corps locaux et globaux\fg}, \emph{Publ. Math. Inst. Hautes \'Etudes Sci.} \textbf{21} (1964), p\ptbl 5-128.
}

\bibitem[7]{XIII.7} \sisi{%
A.H. Wallace, Homology Theory of algebraic Varieties, Pergamon Press, 1958.}
{%
\textsc{A.H. Wallace} -- \emph{Homology Theory of algebraic Varieties}, Pergamon Press, 1958.
}

\bibitem[8]{XIII.8} \sisi{%
S. Abhyankar, Resolution of singularities of arithmetical surfaces, p\ptbl 111--152, Arithmetical Algebraic Geometry, Harper and Row, New-York 1963.}
{%
\textsc{S. Abhyankar} -- {\og Resolution of singularities of arithmetical surfaces\fg}, in \emph{Arithmetical Algebraic Geometry}, Harper and Row, New-York, 1965, p\ptbl 111--152.
}

\sisi{}{%
\bibitem[HL]{HL}
\textsc{H.A. Hamm \& L\^e D\~ung Tr\'ang} -- {\og Rectified homotopical depth and Grothendieck conjectures\fg}, in \emph{The Grothendieck Festschrift, Vol.~II}, Progr. Math., vol.~87, Birkhäuser, Boston, 1990, p\ptbl 311--351.
}
\end{thebibliography}

\chapterspace{-4}

% END INCLUDED active_source/parts/13_expose_XIII.tex

% BEGIN INCLUDED active_source/parts/14_expose_XIV.tex
\chapter[Depth and Lefschetz theorems]
{Depth and Lefschetz theorems in~\'{e}tale cohomology} \label{XIV}
\begin{center}
by Mrs M. Raynaud\sfootnote{Based on unpublished notes by A.~Grothendieck.}
\end{center}

\vspace*{1cm}
In\pageoriginale \numero \Ref{XIV.1}, we define a notion of \og \textit{\'{e}tale depth}\fg that is the analogue in \'{e}tale cohomology of the notion of depth studied in \Ref{III}, in the cohomology of coherent sheaves. After a technical part, we prove \og Lefschetz theorems\fg in \numero \Ref{XIV.4}, the central theorem being \Ref{XIV.4.2}. Let $X$ be a scheme, $Y$ a closed subset of $X$, $U$ the complementary open subset $X-Y$, and $F$ an abelian sheaf on $X$ for the \'{e}tale topology; in general, the purpose of Lefschetz theorems is to show that, if $F$ satisfies certain local conditions on $U$, expressible in terms of \'{e}tale depth at the points of $U$, then, under certain additional conditions of a global nature on $U$ (for example, $U$ \textit{affine}), the natural map of \'{e}tale cohomology groups
$$
\H^i(X, F)\to \H^i(Y, \sisi{F{|Y}}{F_{|Y}})$$
is an isomorphism for values $i<n$, where $n$ is a certain specified integer. By taking $F$ to be a constant sheaf, one thereby obtains conditions for $\pi_0(X)$ to equal $\pi_0(Y)$ and conditions for the abelianized fundamental groups of $X$ and $Y$ to be the same. In \numero \Ref{XIV.5}, the introduction of a notion of \og geometric depth\fg makes it possible to give useful special cases of the Lefschetz theorems (\Ref{XIV.5.7}). Finally, in \numero \Ref{XIV.6}, we mention some conjectures, concerning in particular \og noncommutative\fg variants of the theorems obtained.

\section{Cohomological and homotopical depth}\label{XIV.1}

\setcounter{subsection}{-1}

\subsection{\texorpdfstring{}{XIV.1.0}}\label{XIV.1.0}
Fix\pageoriginale the following notation. Let $X$ be a scheme\sfootnote{In accordance with the new terminology (\cf the reissue of \EGA I), here we shall call \og scheme\fg what was formerly called \og prescheme\fg, and \og separated scheme\fg what was called \og scheme\fg.}, $Y$ a closed subset of~$X$, $U$ the complementary open subset, and $i:Y=X-U\to X$ the canonical immersion. Let~$\underline{\Gamma}_Y$ be the functor that associates to an abelian sheaf on $X$ the \og sheaf of sections with support in $Y$\fg, that is, $\underline{\Gamma}_Y=i_*i^!$ (\cf \SGA 4 IV 3.8 and VIII 6.6), and let $\Gamma_Y$ be the functor $\Gamma\cdot \underline{\Gamma}_Y$ (where $\Gamma$ is the \og global sections\fg functor). Consider the derived category $D^+(X)$ and the derived functor $\R\underline{\Gamma}_Y$ (\resp $\R\Gamma_Y$) of $\underline{\Gamma}_Y$ (\resp of $\Gamma_Y$) (\cf \cite{XIV.3}). Given a complex of abelian sheaves $F$ on $X$, bounded below, it may be regarded as an element of $D^+(X)$; we then denote by ${\h}_Y^p(F)$ the $p$-th cohomology sheaf of $\R\underline{\Gamma}_Y(F)$ and by $\H_Y^p(X, F)$ the $p$-th cohomology group of $\R\Gamma_Y(F)$. The results of (\SGA 4 V 4.3 and 4.4) extend trivially to ${\h}_Y^p(F)$ and $\H_Y^p(X, F)$.

\enlargethispage{\baselineskip}%
\begin{proposition} \label{XIV.1.1}
Let $X$ be a scheme, $Y$ a closed subset of $X$, $U$ the complementary open subset, and $i:U\to X$ the canonical immersion. Let $F$ denote either a sheaf of sets on $X$, a sheaf of groups on $X$, or a complex of abelian sheaves on $X$, bounded below. Fix the following notation: if $X'\to X$ is a morphism, $U'$ and $F'$ denote the inverse images of $U$ and $F$ on $X'$; moreover, if $\bar{y}$ is a geometric point of $X$, $\bar{X}$ denotes the strict localization of $X$ at $\bar{y}$, and $\bar{U}$ and $\bar{F}$ the inverse images of $U$ and $F$ in $\bar{X}$.

$1^\circ)$ Let\pageoriginale $F$ be a sheaf of sets on $X$ and $n$ an integer $\leq 2$; then the following conditions are equivalent:

\textup{(i)} The canonical morphism
$$F\to i_*i^*F$$
is injective if $n\geq 1$, and bijective if $n\geq 2$.

\textup{(ii)} For every scheme $X'$ \'{e}tale over $X$, the canonical morphism
$$
\H^0(X', F')\to \H^0(U', F')
$$
is injective if $n\geq 1$, and bijective if $n\geq 2$.

Suppose in addition that the open subset $U$ is retrocompact in $X$; then the preceding conditions are equivalent to the following one:

\textup{(iii)} For every geometric point $\bar{y}$ of $Y$, the canonical morphism
$$
\H^0(\bar{X}, \bar{F})\to \H^0(\bar{U}, \bar{F})$$
is injective if $n\geq 1$, and bijective if $n\geq 2$.

$2^\circ)$ Let $F$ be a sheaf of groups on $X$ and $n$ an integer $\leq 3$; then the following conditions are equivalent:

\textup{(i)} The canonical morphism
$$F\to i_*i^*F$$
is injective if $n\geq 1$, bijective if $n\geq 2$, and if $n\geq 3$, in addition to the preceding conditions, the sheaf of pointed sets $\R^1i_*(i^*F)$ is zero.

\textup{(ii)} For every scheme $X'$ \'{e}tale over $X$, the canonical morphism
$$
\H^0(X', F')\to \H^0(U', F')
$$
is\pageoriginale injective if $n\geq 1$, bijective if $n\geq 2$, and moreover the canonical morphism
$$
\H^1(X', F')\to \H^1(U', F')
$$
is injective if $n\geq 2$, and bijective if $n\geq 3$.

\textup{(ii bis)} The same as \textup{(ii)}, except in the case $n=2$, where it is only assumed that $\H^0(X', F')\to \H^0(U', F')$ is bijective.

Suppose in addition that $U$ is retrocompact in $X$; then the preceding conditions are also equivalent to the following one:

\textup{(iii)} For every geometric point $\bar{y}$ of $Y$, the canonical morphism
$$
\H^0(\bar{X}, \bar{F})\to \H^0(\bar{U}, \bar{F})
$$
is injective if $n\geq 1$, and bijective if $n\geq 2$; finally, if $n\geq 3$, in addition to the preceding conditions, $\H^1(\bar{U}, \bar{F})$ is zero.

$3^\circ)$ Let $F$ be a complex of abelian sheaves, bounded below, and let $n$~be an integer; then the following conditions are equivalent:

\textup{(i)} We have ${\h}_Y^p(F)=0$ for $p<n$ (\cf \hyperref[XIV.1.0]{1.0}).

\textup{(ii)} For every scheme $X'$ \'{e}tale over $X$, the canonical morphism
$$
\H^p(X', F')\to \H^p(U', F')
$$
is bijective for $p<n-1$, and injective for $p=n-1$.

Suppose that $U$ is retrocompact in $X$; then the preceding conditions are also equivalent to the following one:

\textup{(iii)} For every geometric point $\bar{y}$ of $Y$, the canonical morphism
$$
\H^p(\bar{X}, \bar{F})\to \H^p(\bar{U}, \bar{F})$$
is bijective for $p<n-1$, and injective for $p=n-1$.

In\pageoriginale the case where $F$ is an abelian sheaf and $n\geq 2$, conditions \textup{(i)} and \textup{(ii)} are also equivalent to the following one:

\textup{(ii bis)} For every scheme $X'$ \'{e}tale over $X$, the canonical morphism
$$
\H^p(X', F')\to \H^p(U', F')$$
is bijective for $p<n-1$.
\end{proposition}

\skpt
\begin{proof}
$1^\circ)$ It is clear that \textup{(i)}\SSI \textup{(ii)}.
Let us show that, if $U$ is quasi-compact in $X$, \textup{(i)}\SSI
\textup{(iii)}. Indeed, \textup{(i)} is equivalent to saying that, for
every geometric point $\bar{y}$ of $X$, the morphism
$F_{\bar{y}}\to (i_*i^*F)_{\bar{y}}$ is injective if $n\geq 1$ and
bijective if $n\geq 2$ (\SGA 4 VIII 3.6). Since this morphism is in
any case bijective when $\bar{y}$ is a geometric point of $U$, it is
enough to consider the geometric points $\bar{y}$ of $Y$. Now it
follows from the fact that $i$ is quasi-compact and from
(\SGA 4 VIII 5.3) that the morphism
$$F_{\bar{y}}\to (i_*i^*F)_{\bar{y}}$$
is canonically identified with the morphism
$$
\H^0(\bar{X}, \bar{F})\to \H^0(\bar{U}, \bar{F}),
$$
which proves the equivalence of \textup{(i)} and \textup{(iii)}.

$2^\circ)$ \textup{(i)}\ALORS \textup{(ii)}. The assertions concerning $\H^0$ follow from $1^\circ)$. Let $i'$ be the canonical immersion of $U'$ into $X'$; the assertions concerning $\H^1$ follow from the exact sequence (\SGA 4 XII 3.2)
$$0\to \H^1(X', i'_*i'^*F')\to \H^1(U', F')\to \H^0(X', \R^1i'_*(i'^*F')).
$$

(ii bis)\ALORS \textup{(i)}\pageoriginale. By $1^\circ)$, it is enough to show that, for $n\geq 3$, one has $\R^1i_*(i^*F)=0$. Now $\R^1i_*(i^*F)$ is the sheaf associated with the presheaf $X'\mto \H^1(U', F')$, that is, by hypothesis, the sheaf associated with the presheaf $X'\mto \H^1(X', F')$, and the latter is zero.

\textup{(i)}\SSI \textup{(iii)}. Taking $1^\circ)$ into account, the only thing left to verify is that the relation $\R^1i_*(i^*F)=0$ is equivalent to the condition that $\H^1(\bar U, \bar F)=0$ for every geometric point $\bar y$ of $Y$. Since $\R^1i_*(i^*F)$ is zero outside $Y$, saying that $\R^1i_*(i^*F)=0$ is equivalent to saying that $(\R^1i_*(i^*F))_{\bar y}=0$ for every geometric point $\bar y$ of $Y$. It is then enough to note that, since $i$ is quasi-compact, one has $(\R^1i_*(i^*F))_{\bar y}=\H^1(\bar U, \bar F)$ (\SGA 4 VIII 5.3).

$3^\circ)$ \textup{(i)}\ALORS \textup{(ii)}. Let $X'$ be a scheme \`etale over $X$; one has the exact sequence (\SGA 4 V 4.5)
\begin{equation*} \label{eq:XIV.1.1.*} \tag{$*$} {\to \H^p_{Y'}(X', F')\to H^p(X', F')\to H^p(U', F')\to};
\end{equation*}
hence \textup{(ii)} is equivalent to $H^p_{Y'}(X', F')=0$ for $p<n$ and for every scheme $X'$ \`etale over $X$. Consider now the spectral sequence
$$
\E_2^{pq}=\H^p(X', {\h}_Y^q(F))\To \H_{Y'}^*(X', F');$$
by hypothesis, ${\h}_Y^q(F)=0$ for $q<n$, whence $\E_2^{pq}=0$ for $p+q<n$ and consequently $\H_{Y'}^p(X', F')=0$ for $p<n$.

\textup{(ii)}\ALORS \textup{(i)}. The sheaf ${\h}_Y^p(F)$ is associated with the presheaf $X'\mto \H_{Y'}^p(X', F')$; since we have already observed that \textup{(ii)} is equivalent to the relation $\H_{Y'}^p(X', F')=0$ for $p<n$ and for every scheme $X'$ \`etale over $X$, it follows that ${\h}_Y^p(F)=0$ for $p<n$.

\textup{(i)}\SSI \textup{(iii)}. The sheaves ${\h}_Y^p(F)$ are concentrated on $Y$; consequently, saying that ${\h}_Y^p(F)=0$ is equivalent to saying that, for every geometric point\pageoriginale $\bar y$ of $Y$, the fiber $({\h}_Y^p(F))_{\bar y}$ is zero. Since $i$ is quasi-compact, (\SGA 4 VIII 5.2) gives $({\h}_Y^p(F))_{\bar y}={\h}_{\bar Y}^p(\bar X, \bar F)$. The equivalence of \textup{(i)} and \textup{(iii)} follows, in view of the analogue on $\bar X$ of the exact sequence (\Ref{eq:XIV.1.1.*}).

(ii bis)\ALORS \textup{(ii)} when $F$ is an abelian sheaf. The only thing left to show is that ${\h}_Y^{n-1}(F)=0$. Now, for $n>2$, the sheaf ${\h}_Y^{n-1}(F)$ is associated with the cokernel of the morphism of presheaves $X'\mto \H^{n-2}(X', F')\to\H^{n-2}(U', F')$, and is therefore zero. The case $n=2$ follows from the fact that ${\h}_Y^1(F)$ is the cokernel of the morphism $F\to i_*i^*F$.
\skipqed
\end{proof}

\begin{definition} \label{XIV.1.2}
The notation is that of \Ref{XIV.1.1}. One says that $F$ has \`etale $Y$-depth $\geq n$ and writes
$$
\prof_Y(F)\geq n$$
if $F$ satisfies the equivalent conditions \textup{(i)} and \textup{(ii)} of \Ref{XIV.1.1}. If $F$ is a complex of abelian sheaves, the \`etale $Y$-depth of $F$ is the supremum of the~$n$ for which $\prof_Y(F)\geq n$; the same notation will be used if $F$ is a sheaf of sets \resp of not necessarily commutative groups (so that one then has $0\leq \prof_Y(F)\leq 2$\sisi{}{,} \resp $0\leq \prof_Y(F)\leq 3$, when the context does not make clear which of the three variants considered here is being used).

If $\LL$ is a set of prime numbers, one says that the \`etale $Y$-depth of $X$ for $\LL$ is $\geq n$ and writes
$$
\prof_Y^{\LL}(X)\geq n
$$
if, for every constant sheaf of the form $\ZZ/\ell\ZZ$ with $\ell\in \LL$, one has $\prof_Y(\ZZ/\ell\ZZ)\geq n$\pageoriginale. The \`etale $Y$-depth of $X$ for $\LL$ is defined in the evident way. If $\LL=\PP$, the set of all prime numbers, and there is no risk of confusion with the notation of \textup{(\EGA IV 5.7.1)} (which concerns the case $F=\OX$), $\LL$ is omitted from the notation; otherwise one writes $\prof \et(X)$.

Finally, one says that $X$ has homotopical $Y$-depth \sisi{\ignorespaces}{$\geq 3$} for $\LL \sisi{\geq 3}{}$ and writes
$$
\prof \hop_Y^{\LL}(X)\geq 3$$
if, for every constant finite sheaf $F$ of $\LL$-groups on $X$, one has $\prof_Y(F)\geq 3$. If $\LL=\PP$, $\LL$ is omitted from the notation.
\end{definition}

\begin{corollaire} \label{XIV.1.3}
Under the conditions of \Ref{XIV.1.1}, if $\prof_Y(F)\geq n$, then, for every closed subset~$Z$ of $Y$, one has
$$
\prof_Z(F)\geq n.
$$
\end{corollaire}

For example, let us give the argument when $F$ is a complex of abelian sheaves bounded below. We use \Ref{XIV.1.1} $3^\circ)$ \textup{(ii)}. Let $V=X-Z$ and, for every integer $p$, consider the sequence of morphisms
$$
\H^p(X, F)\lto{f} \H^p(V, F)\lto{g} \H^p(U, F).
$$

By hypothesis, $g$ and $g\circ f$ are bijective for $p<n-1$ and injective for $p=n-1$; hence the same is true of $f$. Since the argument remains valid when $X$ is replaced by a scheme $X'$ \`etale over $X$, this proves \Ref{XIV.1.3}.

\begin{corollaire} \label{XIV.1.4} The notation is that of \Ref{XIV.1.1} $2^\circ)$. If $X'$ is a scheme over\pageoriginale $X$, let $\Phi'$ denote the functor that assigns to an \`etale covering of $X'$ its restriction to~$U'$, and let $\Phi'_{F'}$ denote the functor that assigns to a torsor\sfootnote{That is, a \og homogeneous principal bundle\fg in older terminology.} under $F'$ its restriction to~$U'$. Then the following conditions are equivalent:

\textup{(i)} One has $\prof_Y(F)\geq 1$ (\resp $\prof_Y(F)\geq 2$, \resp $\prof_Y(F)\geq 3$).

\textup{(ii)} For every scheme $X'$ \`etale over $X$, the functor $\Phi'_{F'}$ is faithful (\resp fully faithful, \resp an equivalence of categories).

In particular, one has $\prof_Y({X})\geq 1$ (resp.\ $\prof_Y({X})\geq 2$, resp.\ $\prof \hop_Y(X)\geq 3$) if and only if the functor $\Phi'$ is faithful (resp.\ fully faithful, \resp an equivalence of categories).
\end{corollaire}

This follows from \Ref{XIV.1.1} $2^\circ)$ \textup{(ii)}, in view of the interpretation of $\H^1(X', F')$ as the set of classes (up to isomorphism) of torsors under $F'$ (\SGA 4 VII~2), and of \`etale coverings $Z$ of degree $n$ of a scheme as being associated with principal Galois coverings whose group is the symmetric group $\mathfrak{S}_n$: to $Z$ one associates the covering $\underline{\Isom}_X(Z_0, Z)$, where $Z_0$ is the trivial covering of $X$ of degree $n$.

\begin{corollaire} \label{XIV.1.5} Under the conditions of \Ref{XIV.1.1} $3^\circ)$, suppose that $\prof_Y(F)\geq n$; then
$$
\H_Y^n(X, F)\simeq \H^0(X, {\h}_Y^n(F)).
$$
\end{corollaire}

The corollary follows from the spectral sequence
$$
\E_2^{pq}=H^p(X, {\h}_Y^q(F))\To H_Y^*(X, F).
$$

Indeed, by hypothesis $\E_2^{pq}=0$ for $q<n$, and hence
$$
\H_Y^n(X, F)=\E_2^{0n}=\H^0(X, {\h}_Y^n(F)).
$$

\skpt
\begin{remarques} \label{XIV.1.6}
a) The notion\pageoriginale of $Y$-depth, in the form of the equivalent conditions \textup{(i)} and \textup{(ii)} of \Ref{XIV.1.1}, makes sense on any site. In the particular case in which $X$ is a locally Noetherian scheme equipped with the Zariski topology and $F$ is a sheaf of coherent $\OX$-modules, one recovers the usual notion of $Y$-depth as the infimum of the depths at the points of $Y$ (\Ref{III}).

b) For $n\leq 2$, the notion of \`etale $Y$-depth of $X$ is independent of $\LL$. For $n=1$, it simply means that $U$ is dense in $X$. Indeed, this condition is necessary for $\prof_Y(F)\geq 1$, and it is also sufficient because one may assume that $X$ is reduced, in which case the condition that $U$ be dense in $X$ is preserved by \`etale base change \textup{(\EGA IV 11.10.5) \textup{(ii)} b))}. If $U$ is quasi-compact in $X$, the relation $\prof_Y(X)\geq 1$ is also equivalent to saying that $Y$ contains no generic point of~$X$ (\EGA I 6.6.5). For $n=2$ and $U$ quasi-compact in $X$, the condition $\prof_Y(X)\geq 2$ is equivalent to the condition that, for every geometric point $\bar y$ of $Y$, $\bar U$ be nonempty and connected, that is, that \og $Y$ does not disconnect $X$ locally for the \`etale topology\fg.

c) If $X$ has $Y$-depth $\geq n$ for $\LL$ and $U$ is quasi-compact in $X$, then for every locally constant abelian sheaf of $\LL$-torsion $F$ on $X$, one has $\prof_Y(F)\geq n$. Indeed, since the property $\prof_Y(F)\geq n$ is local for the \`etale topology, one may assume that $F$ is constant; then $F$ is a filtered direct limit of sheaves that are finite direct sums of sheaves of the form $\ZZ/p^m\ZZ$, where $m$ is an integer $>0$ and $p\in \LL$. Using \Ref{XIV.1.1} \textup{(iii)} and \textup{(\SGA 4 VII 3.3)}, one sees that one may reduce to the case $F=\ZZ/p^m\ZZ$, and then\pageoriginale, by induction on $m$, to the case $F=\ZZ/p\ZZ$, for which the assertion follows from the definition.

d) By \Ref{XIV.1.4}, if $\prof_Y(X)\geq 3$, the pair $(X, Y)$ is \textit{pure} in the sense of \Ref{X}~\Ref{X.3.1}. In fact, the pure pairs encountered in practice (\cf \Ref{X}~\Ref{X.3.4}) satisfy the stronger condition of homotopical depth $\geq 3$, and this notion can therefore advantageously replace that of a pure pair.

e) Let $F$ be a complex of abelian sheaves and $T(F)$ the complex obtained by applying the translation functor to $F$ (\cite{XIV.3}); then one plainly has:
$$
\prof_Y(T(F))=\prof_Y(F)-1.
$$

f) Let us point out that the recent work of Artin--Mazur (\cite{XIV.1}) makes it possible to define the notion of homotopical depth $\geq n$ for every integer $n$ (not only for $n\geq 3$).

g) Under the conditions of \Ref{XIV.1.1} 3), one has $\prof_Y(F)=\infty$ if and only if $F\isomto Ri_*(i^*F)$ in $D^+(X)$. Indeed, the ${\h}_Y^p(F)$ are the cohomology sheaves of the cone (= mapping cylinder) of the canonical morphism $F\to Ri_*i^*(F)$.
\end{remarques}

\begin{definition} \label{XIV.1.7}
Let $X$ be a scheme, $x$ a point of $X$, $\bar x$ a geometric point above $x$, and $\bar X$ the strict localization of $X$ at $\bar x$. As before, $F$ denotes either a sheaf of sets on $X$, a sheaf of groups on $X$, or a bounded-below complex of abelian sheaves on $X$; $\bar F$ denotes its inverse image on~$\bar X$, and $\LL$ a set of prime numbers. We say that $F$ has \emph{étale depth $\geq n$ at the point $x$}\pageoriginale (respectively that the \emph{étale depth for $\LL$ of $X$ at $x$ is $\geq n$}, respectively that the \emph{homotopical depth for $\LL$ of $X$ at $x$ is $\geq 3$}) and write $\prof_x(F)\geq n$ (respectively $\prof_x^{\LL}(X)\geq n$, respectively $\prof \hop_x^{\LL}(X)\geq \sisi{3}{3}$) if $\prof_{\bar x}(\bar F)\geq n$ (respectively $\prof_{\bar x}^{\LL}(\bar X)\geq\nobreak n$, respectively $\prof \hop_{\bar x}^{\LL}(\bar X)\geq \sisi{3}{3}$). The integer $\prof_x^{\LL}(X)$ and, if $F$ is a complex of abelian sheaves, the integer $\prof_x(F)$ are defined in the evident way. If $\LL$ is the set of all prime numbers, we omit $\LL$ from the notation $\prof_x^{\LL}(X)$, unless there is a risk of confusion with the notation of \textup{(\EGA IV 5.7.1)}, in which case we write $\prof \et_x(X)$.
\end{definition}

We then have the following \textit{pointwise characterization} of depth:

\begin{theoreme} \label{XIV.1.8} Let $X$ be a scheme and $Y$ a closed subset of $X$ such that the open subset $U=X-Y$ is retrocompact in $X$. If $F$ is either a sheaf of sets on $X$, a sheaf of groups on $X$, or a bounded-below complex of abelian sheaves on $X$, then
$$
\prof_Y(F)=\inf_{y\in Y} \prof_y(F).
$$
\end{theoreme}

\subsubsection{} \label{XIV.1.8.1} Let us first show that, for every point $y$ of $Y$, one has the inequality $\prof_y(F)\geq \prof_Y(F)$. Indeed, let $\bar y$ be a geometric point above $y$, $\bar X$ the strict localization of $X$ at $\bar y$, and $\bar F$ and $\bar Y$ the inverse images of $F$ and $Y$ on $\bar X$. By \Ref{XIV.1.7} and \Ref{XIV.1.3}, one has
$$
\prof_y(F)=\prof_{\bar y}(\bar F)\geq \prof_{\bar Y}(\bar F)\geq \prof_Y(F),
$$
the last inequality using the hypothesis \og $U$ retrocompact\fg, through conditions \textup{(iii)} in \Ref{XIV.1.1} and transitivity in the formation of strict localizations.

\subsubsection{} \label{XIV.1.8.2}
Conversely, suppose that, for every point $y$ of $Y$, one has $\prof_y(\sisi{Y}{F})\geq n$ ($n$~an integer), and let us show that $\prof_Y(F)\geq n$.

First recall\pageoriginale the following well-known results \textup{(\SGA 4 VIII)}:

\begin{subsublemme} \label{XIV.1.8.2.1}
Let $X$ be a scheme and $F$ a sheaf of sets on $X$ (respectively let $G\to F$ be a monomorphism of sheaves of sets on $X$). Then two sections $s$ and $s'$ of $F$ coincide (respectively a section $s$ of~$F$ comes from a section of~$G$) if and only if this is so locally. In particular, if~$s$ and~$s'$ are two sections of $F$, there is a largest open subset $V$ of $X$ on which they coincide (respectively, if $s$ is a section of $F$ on $X$, there is a largest open subset $V$ of $X$ such that $s|V$ comes from a section of $G$ on $V$). This open subset is also the set of points~$x$ of $X$ such that, denoting by $\bar x$ a geometric point above~$x$, the sections~$s$ and~$s'$ have the same image in the fiber $F_{\bar x}$ (respectively the image of $s$ in $F_{\bar x}$ comes from an element of $G_{\bar x}$).
\end{subsublemme}

Let us return to the proof of \Ref{XIV.1.8}.

$1^\circ)$ Case where $F$ is a sheaf of sets. If $n=1$, it is enough to show that the canonical morphism
$$
\H^0(X, F)\to \H^0(U, F)$$
is injective, since the result also applies when $X$ is replaced by a scheme étale over $X$. Let $s$ and $s'$ be two sections of $F$ over $X$ which have the same image in $\H^0(U, F)$, and let $V$ be the largest open subset over which they are equal; clearly $V\supset U$. Suppose $V\neq X$, and let $y$ be a maximal point of $X-V$, $\bar y$ a geometric point above\pageoriginale~$y$, $\bar X$ the strict localization of $X$ at $\bar y$, and $\bar V$ and $\bar F$ the inverse images of $V$ and $F$ on $\bar X$. By the choice of $y$, one has $\bar X-\bar y=\bar V$, and by hypothesis the morphism
$$
\H^0(\bar X, \bar F)\to \H^0(\bar X-\bar y, \bar F)=\H^0(\bar V, \bar F)$$
is injective. It follows that $s$ and $s'$ coincide at the point $y$, which is absurd. If $n=2$, it is enough, in view of the preceding, to show that the morphism
$$
\H^0(X, F)\to \H^0(U, F)=\H^0(X, i_*i^*F)$$
is surjective (where $i$ is the canonical immersion of $U$ into $X$). Let $s$ be a section of $i_*i^*F$ over $X$, and let $V$ be the largest open subset over which it comes from a section of $F$. Suppose $V\neq X$, and let $y$ be a maximal point of $X-V$; with the preceding notation, it follows from the hypothesis that the canonical morphism
$$
\H^0(\bar X, \bar F)\to \H^0(\bar X-\bar y, \bar F)=\H^0(\bar V, \bar F)
$$
is bijective; consequently $s|\bar V$ extends to the point $\bar y$, which is absurd and completes the proof in case $1^\circ)$.

$2^\circ)$ Case where $F$ is a sheaf of groups. In view of $1^\circ)$, the only thing left to show is that, in the case $n=3$, the morphism
$$
\H^1(X, F)=\H^1(X, i_*i^*F)\to \H^1(U, F)$$
is bijective. We already know that it is injective, by $1^\circ)$ and \Ref{XIV.1.1} $2^\circ)$ (ii bis). For surjectivity, we use the exact sequence (\SGA 4 XII 3.2)
$$
0\to \H^1(X, i_*i^*F)\to \H^1(U, F)\lto{d}\H^0(X, \R^1i_*(i^*F)).
$$
Let $s\in \H^1(U, F)$, and let $V\supset U$ be the largest open subset over which $d(s)=0$; this is\pageoriginale also the largest open subset such that $s$ comes from an element of $\H^1(V, F)$. Suppose $V\neq X$, and let $y$ be a maximal point of $X-V$; if $\bar X$ is the strict localization of $X$ at a geometric point $\bar y$ above~$y$, then, with the evident notation, we have the exact sequence
$$
0\to \H^1(\bar X, \bar i_*(\bar i^*\bar F))\to \H^1(\bar U, \bar F)\lto{\bar d}\H^0(\bar X, \R^1\bar i_*(\bar i^*\bar F)).
$$
Since $i:U\to X$ is quasi-compact, $\R^1\bar i_*(\bar i^*\bar F)$ is the inverse image of $\R^1i_*(i^*F)$ under the morphism $\bar X\to X$, whence $\H^0(\bar X, \R^1\bar i_*(\bar i^*\bar F))=(\R^1i_*(i^*F))_{\bar y}$. By hypothesis, and since $y\in Y$, the morphism
$$
\H^1(\bar X, \bar F)\to \H^1(\bar V, \bar F)
$$
is bijective. The image $\bar s$ of $s$ in $\H^1(\bar U, \bar F)$, which extends to $\bar V$ by the definition of $V$, therefore also extends to $\bar X$; it follows that $\bar d(\bar s)=0$, and hence the image of $d(s)$ in the geometric fiber $(\R^1i_*(i^*F))_{\bar y}$ is zero; but this contradicts the definition of $V$, proving case~$2^\circ)$.

$3^\circ)$ Case where $F$ is a bounded-below complex of abelian sheaves. We argue by induction on $n$. The conclusion holds for sufficiently small $n$, since $F$ is bounded below. Suppose, therefore, that $\prof_Y(F)\geq n-1$, and let us show that $\prof_Y(F)\geq n$, knowing that for every point $y$ of $Y$ one has $\prof_y(F)\geq n$. It is enough to see that the canonical morphism
\begin{equation*} \label{eq:XIV.1.8.*} \tag{$*$} {\H^{n-2}(X, F)\to \H^{n-2}(U, F)}
\end{equation*}
is surjective and that
\begin{equation*} \label{eq:XIV.1.8.**} \tag{$**$} {\H^{n-1}(X, F)\to \H^{n-1}(U, F)}
\end{equation*}
is injective (the result also applying when $X$ is replaced by a scheme étale over\pageoriginale $X$).

a) Surjectivity of (\Ref{eq:XIV.1.8.*}). The proof is analogous to that of $2^\circ)$. In view of \Ref{XIV.1.5} and (\SGA 4 V 4.5), one has the exact sequence
$$
\H^{n-2}(X, F)\to \H^{n-2}(U, F)\lto{d}\H^{n-1}_Y(X, F)=\H^0(X, {\h}_Y^{n-1}(F)).
$$
Let $s\in \H^{n-2}(U, F)$, and let $V\supset U$ be the largest open subset over which $d(s)=0$, which is also the largest open subset such that $s$ extends to $\H^{n-2}(V, F)$. Suppose $V\neq X$, and let $y$ be a maximal point of $X-V$ and $\bar X$ the strict localization of $X$ at a geometric point $\bar y$ above~$y$. Since $i:U\to X$ is quasi-compact, the formation of ${\h}_Y^{n-1}(F)$ commutes with the base change $\bar X\to X$, and hence one has (with the evident notation) the exact sequence
$$
\H^{n-2}(\bar X, \bar F)\to \H^{n-2}(\bar U, \bar F)\lto{\bar d}\H^{n-1}_{\bar Y}(\bar X, \bar F)=({\h}_Y^{n-1}(F))_{\bar y},
$$
the last equality following from the retrocompactness hypothesis on $U$.

But by hypothesis one has the isomorphism
$$
\H^{n-2}(\bar X, \bar F)\isomto \H^{n-2}(\bar X-\bar y, \bar F)=\H^{n-2}(\bar V, \bar F);
$$
consequently the image $\bar s$ of $s$ in $\H^{n-2}(\bar U, \bar F)$, which extends (by definition of $V$) to $\H^{n-2}(\bar V, \bar F)$, also extends to $\H^{n-2}(\bar X, \bar F)$; but this shows that $\bar d(\bar s)=0$, that is, that $d(s)$ is zero at $\bar y$, which is absurd.

b) Injectivity of (\Ref{eq:XIV.1.8.**}). Using the surjectivity of (\Ref{eq:XIV.1.8.*}), one obtains the exact sequence
$$0\to \H^0(X, {\h}_Y^{n-1}(F))\to \H^{n-1}(X, F)\to \H^{n-1}(U, F)$$
and we must show that every element $s\in \H^0(X, {\h}_Y^{n-1}(F))$ is zero. Let $V$\pageoriginale be the largest open subset over which $s=0$. Suppose that $V\neq X$, and let $y$ be a maximal point of $X-V$ and $\bar X$ a strict localization of $X$ at a geometric point $\bar y$ above~$y$. By the induction hypothesis and \Ref{XIV.1.8.1}, one has $\prof_{\bar Y}(\bar F)\geq \prof_Y(F)\geq n-1$, and hence the map $e$ in the following diagram is injective:
$$
\H^0(\bar X, {\h}_Y^{n-1}(\bar F))= ({\h}_Y^{n-1}(F))_{\bar y}\lto{e} \H^{n-1}(\bar X, \bar F) \lto{f} \H^{n-1}(\bar V, \bar F).
$$
The same is true of $f$ by the hypothesis; the equality on the left follows from the retrocompactness hypothesis on $U$. Let $\bar s$ be the image of $s$ in $({\h}_Y^{n-1}(F))_{\bar y}$; since $s$ vanishes over $V$, one has $f\cdot e(\bar s)=0$, whence $\bar s=0$, which contradicts the choice of $y$ and completes the proof.

\begin{remarque} \label{XIV.1.9}
A result analogous to \Ref{XIV.1.8} is probably valid when the étale topos of a scheme $X$ is replaced by a \og locally finite-type topos\fg, that is, one definable by a locally finite-type site (\SGA 4 VI 1.1). To see this, one must use a result of P. Deligne (\SGA \sisi{4 XVII}{4 VI.9}) asserting that there are \og enough fiber functors\fg in such a topos.
\end{remarque}

We shall deduce from \Ref{XIV.1.8} important cases in which the étale depth can be determined.

\refstepcounter{subsection}\label{XIV.1.10}
\begin{enonce*}{Theorem 1.10}[Cohomological semipurity theorem\ndemark]\ndetext{\label{gabpur} Gabber subsequently proved --- in 1994 --- Grothendieck's absolute cohomological purity conjecture: if $Y$ is a closed subscheme of absolute Noetherian schemes of pure codimension $c$ and $n$ is an integer invertible on $X$, then $\SheafH_Y^q(\Lambda)$ is zero if $q\neq 2c$ and is $\Lambda_Y(-c)$ (Tate twist) otherwise, where $\Lambda=\ZZ/n\ZZ$. See (Fujiwara~K., {\og A Proof of the Absolute Purity Conjecture (after Gabber)\fg}, in \emph{Algebraic geometry 2000, Azumino (Hotaka)}, Adv. Stud. in Pure Math., vol.~36, 2002, p\ptbl 153-183). For applications to the existence of the dualizing complex, see (\SGA 5, Lect. Notes in Math., vol.~589, Springer-Verlag, 1977, p\ptbl 16--72), Exposé 1 and \loccit, \S8. This conjecture had been proved in the case $n=\ell^\nu$, with $\ell$ a sufficiently large prime invertible on $X$, by making crucial use of $K$-theory (Thomason~R.W., {\og Absolute cohomological purity\fg}, \emph{Bull. Soc. Math. France} \textbf{112} (1984), \numero 3, p\ptbl 397--406). $K$-theory reappears in Gabber's proof through the Atiyah-Hirzebruch-Thomason spectral sequence relating étale cohomology and $K$-theory, a method already used in Thomason's approach. Besides this result, the other fundamental argument is the generalization of the Lefschetz theorem cited in note~\eqref{gabaf}, page~\pageref{gabaf}.}
Let~$X$ denote either a scheme smooth over a field $k$, or an excellent regular scheme \textup{(\EGA IV 7.8.2)} of characteristic zero (N.B. if resolution of singularities in the sense of \textup{(\SGA 4 XIX)} is assumed, it is enough, more generally, to suppose that $X$ is an excellent regular scheme of equal characteristic).\pageoriginaled Let $Y$ be a closed subset of $X$ and $\LL$ the set of prime numbers distinct from the characteristic of $X$. Then
$$
\prof_Y^{\LL}(X)=2\codim(Y, X).
$$
\end{enonce*}

\begin{proof}
It follows from \Ref{XIV.1.8} that
$$
\prof_Y^{\LL}(X)=\inf_{y\in Y} \prof_y(X).
$$

On the other hand, since $\codim(Y, X)=\inf_{y\in Y} \dim \Oo_{X, y}$, we are reduced to showing that
$$
\prof_y^{\LL}(X)=2 \dim \Oo_{X, y},
$$
which follows from (\SGA 4 XVI 3.7 and XIX 3.2).
\skipqed
\end{proof}

\begin{theoreme}[Homotopical purity theorem\ndemark] \label{XIV.1.11}
If $X$ is a locally noetherian scheme that is regular (\resp whose local rings are complete intersections), and $Y$ is a closed subset of $X$ such that $\codim(Y, X)\geq 2$ (resp.\ \hbox{$\codim(Y, X)\geq3$}), then one has\ndetext{\label{jongoort} recently, de Jong and Oort obtained the following purity statement: let $\tilde S\to S$ be a resolution of the singularities of the spectrum $S$ of a normal noetherian local ring of dimension $2$, and let~$U$ be the complement of the closed point $s$ in $S$. Suppose in addition that $k(s)$ is algebraically closed. Then, for every prime number $p$, in particular if $S$ has characteristic $p$, the restriction morphism $H^1_\et(\tilde S,\QQ_p)\to H^1_\et(U,\QQ_p)$ is bijective (de Jong A.J. \& Oort~F, {\og Purity of the stratification by Newton polygons\fg}, \emph{J.~Amer. Math. Soc.} \textbf{13} (2000), \numero 1, p\ptbl 209--241, Theorem 3.2). If $k=\CC$ and $A$ is the completion of a surface singularity, this result is due to Mumford (see page~\pageref{XIII.7}, \cite{XIII.5}).}
$$
\prof \hop_Y(X)\geq 3.
$$
\end{theoreme}

Indeed, it follows from \Ref{XIV.1.8} that $\prof \hop_Y(X)=\inf_{y\in Y} \prof \hop_y(X)$. But the strictly local rings of $X$ at the various points of $Y$ are regular rings of dimension $\geq 2$ (\resp complete intersections of dimension $\geq 3$). It then follows from the purity theorem \Ref{X}~\Ref{X.3.4} that $\prof \hop_y(X)\geq 3$, which proves the theorem.

\begin{exemple} \label{XIV.1.12}
Let $X$ be a locally noetherian scheme, $Y$ a closed subset of~$X$, and $n=1$ or $2$. Then, if $\prof_Y(\OX)\geq n$ ($\prof_Y(\OX)$ denoting the $Y$-depth in the sense of coherent sheaves (\cf \hyperref[XIV.1.6]{1.6} a)), one also has $\prof_Y(X)\geq n$; this is obvious for $n=1$, and for $n=2$ it is\pageoriginale precisely Hartshorne's theorem (\Ref{III}~\Ref{III.1}). On the other hand, the analogous assertion is false for $n\geq 3$. For example, take an affine space of dimension $\geq 3$ over a field of characteristic $\neq 2$ and let the group $\ZZ/2\ZZ$ act by reflection through the origin. Let $X$ be the quotient and $Y=\{x\}$ the image of the origin in $X$. Then $\Oo_{X, x}$ is a Cohen--Macaulay ring, so $\prof_x(\OX)\geq 3$; but the affine space with the origin removed is an \'{e}tale covering of $X-\{x\}$ that does not extend to an \'{e}tale covering of $X$; hence, by \Ref{XIV.1.4}, $\prof_Y(X)=2$.
\end{exemple}

The following theorem is the analogue of (\EGA IV 6.3.1):

\begin{theoreme} \label{XIV.1.13}
Let $f:X\to S$ be a morphism of schemes, $Y$ a closed subset of $X$, and $Z$ a closed subset of $S$, such that $f(Y)\subset Z$. Suppose that the local rings of $X$ at the various points of $Y$ are noetherian and that the open subsets $X-Y$ and $S-Z$ are retrocompact in $X$ and $S$, respectively. Let $p$, $q$, $r$ be integers such that $p\geq -r$, $q\geq 0$, let $\LL$ be a set of prime numbers, and let $F$ be a complex of abelian sheaves of $\LL$-torsion on $S$, such that the cohomology sheaves $\H^i(F)$ vanish for $i<-r$. Suppose that
\begin{enumeratea}
\item
The morphism $f$ is locally ($p+q+r-2$)-acyclic for $\LL$ \textup{(\SGA 4 XV 1.11)}.
\item
One has
$$
\prof_Z(F)\geq p.
$$
\item
For every point $s$ of $Z$, one has
$$
\prof_{Y_s}^{\LL}(X_s)\geq q.
$$
\end{enumeratea}
\noindent
Then one has\pageoriginale
$$
\prof_Y(f^*F)\geq p+q.
$$
\end{theoreme}

We shall need the following lemma:

\begin{sublemme} \label{XIV.1.13.1}
Let $\LL$ be a set of prime numbers, let $n$ and $r$ be integers, and let $f:X\to S$ be a morphism locally $n$-acyclic for $\LL$. Let $F$ be a complex of abelian sheaves whose cohomology sheaves are of $\LL$-torsion, such that $\H^i(F)=0$ for $i<-r$, let $Z$ be a closed subset of $S$ such that $S-Z$ is retrocompact in $S$, and let $T=f^{-1}(Z)$. Then the canonical morphism
$$f^*({\h}_Z^i(F))\to {\h}_T^i(f^*F)$$
is bijective for $i<n-r+2$ and injective for $i=n-r+2$.
\end{sublemme}

Set $U=S-Z$ and $V=X-T$, so that one has the cartesian square
$$
\xymatrix{
V \ar[r]^-g\ar[d]_k& U\ar[d]^j\\
X \ar[r]^-f&S
}$$
Consider the following commutative diagram whose rows are exact
$$
\xymatrix{ \ar[r]&f^*({\h}_Z^i(F))\ar[r]\ar[d]& f^*({\h}^i(F))\ar[r]\ar[d]^{\wr} &f^*({\h}^i(\R j_*(j^*F)))\ar[r]\ar[d]& \\
\ar[r] &{\h}_T^i(f^*F)\ar[r]& {\h}^i(f^*F) \ar[r]&{\h}^i(\R k_*(k^*f^*F)) \ar[r]&\text{;}}$$
it follows that one is reduced to showing that the morphism
$$f^*({\h}^i(\R j_*(j^*F)))\to {\h}^i(\R k_*(k^*f^*F))$$
is\pageoriginale bijective for $i<n-r+1$ and injective for $i=n-r+1$. But such a morphism arises from the following morphism between hypercohomology spectral sequences
$$
\xymatrix{f^*E_2^{p, q}=f^*(\R^pj_*({\h}^q(j^*F)))\ar@{=>} [r]\ar[d]&f^*({\h}^*(\R j_*(j^*F)))\ar[d]\\
E'^{p, q}_2=\R^pk_*({\h}^q(k^*f^*F))\ar@{=>}[r]&{\h}^*(\R k_*(k^*f^*F)).}
$$
Since $j$ is quasi-compact, it follows from (\SGA 4 XV 1.10) that the morphism $f^*(E^{p, q}_2)\to E'^{p, q}_2$ is bijective for $p\leq n$ and injective for $p=n+1$; in particular, it is bijective for $p+q\leq n-r$ and injective for $p+q=n-r+1$. The conclusion follows at once.

Let us return to the proof of \Ref{XIV.1.13}. Let $T=f^{-1}(Z)$. By \Ref{XIV.1.13.1} and condition a), the canonical morphism $f^*(\h_Z^i(F))\to\h_T^i(f^*F)$ is an isomorphism for $i\leq p+q$. It therefore follows from b) that $\h^i_T(f^*F)=0$ for $i<p$ and, for $i<p+q$, the restriction of $\h^i_T(f^*F)$ to $T$ is the inverse image of a sheaf $G^i$ on $Z$. Let
$$f_T:T\to Z$$
be the restriction of $f$ to $T$. It then follows from c) and the corollary below that $\h^j_Y(f^*_T(G^i))=0$ for $j<q$. We conclude that
$$
\h^j_Y(\h^i_T(f^*F))=0\text{ for } i+j<p+q,
$$
because the inequality $i+j<p+q$ implies either $i<p$, in which case $\h^i_T(f^*F)=0$, or $j<q$, in which case $\h_Y^j(\h_T^i(f^*F))=0$. Given that, with the notation of \Ref{XIV.1.0}, $\underline{\Gamma}_Y=\underline{\Gamma}_Y.\underline{\Gamma}_T$, one has the spectral sequence
\begin{equation*} \label{eq:XIV.1.13.2} \tag{1.13.2}
E_2^{i, j}=\h_Y^j(\h_T^i(f^*F))\To\h^*_Y(f^*F);
\end{equation*}
since\pageoriginale $E_2^{i, j}=0$ for $i+j<p+q$, one sees that $\h_Y^k(f^*F)=0$ for $k<p+q$. The theorem will therefore be proved if we prove the following corollary (which is the special case of \hyperref[XIV.1.13]{1.13} obtained by taking $Z=S, \ r=p=0$ and $F$ concentrated in degree $0$).

\begin{corollaire}\label{XIV.1.14}
Let $f:X\to S$ be a morphism, let $Y$ be a closed subset such that the complementary open subset $X-Y$ is quasi-compact in $X$ and the local rings of~$X$ at the various points of~$Y$ are Noetherian. Let $\LL$ be a set of \hbox{prime numbers}, let $q$ be an integer, and let $F$ be an abelian sheaf of $\LL$-torsion on $S$. Suppose that~$f$ is locally $(q-2)$-acyclic for $\LL$ and that, for every point $s$ of $S$, one has \hbox{$\prof_{Y_s}^\LL(X_s)\!\geq\! q$}. Then $\prof_Y(f^*F)\geq q$.
\end{corollaire}

$1^\circ)$ Reduction to the case in which $X$ and $S$ are strictly local schemes, $f$ is a local morphism, and $Y$ consists of a closed point of $X$.

By \Ref{XIV.1.8}, to prove \Ref{XIV.1.14}, one must show that, for every point $y$ of $Y$, one has:$$
\prof_y(f^*F)\geq q.
$$

Let $s=f(y)$, let $\bar s$ be a geometric point above $s$, let $\bar y$ be a geometric point above~$y$ and $\bar s$, let $\bar X$ and $\bar S$ be the strict localizations of $X$ and $S$ at $\bar y$ and $\bar s$, respectively, let $\bar f:\bar X\to\bar S$ be the canonical morphism, and let $\bar F$ be the inverse image of $F$ on $\bar S$. Since $\prof_y(f^*F)=\prof_{\bar y}(\bar f^*\bar F)$, it is enough to show that the hypotheses of \Ref{XIV.1.14} are preserved when $f$ (\resp $Y$, \resp $F$) is replaced by $\bar f$ (\resp $\{\bar y\}$, \resp $\bar F$). The quasi-compactness condition follows from the Noetherian hypothesis on $\Oo_{X,y}$, which implies that $\bar X$ is Noetherian\pageoriginale. By (\SGA 4 XV~1.10 \textup{(i)}), $\bar f$ is still locally $(q-2)$-acyclic for $\LL$. Moreover, the fiber $(\bar X)_{\bar s}$ of $\bar X$ above $\bar s$ identifies with the strict localization of $X_s$ at $\bar y$, and therefore satisfies $\prof_{\bar y}((\bar X)_{\bar s})\geq q$. Since an analogous relation is trivially satisfied for the fibers of the $\bar S$-scheme $\bar X$ other than the closed fiber, this completes the reduction.

$2^\circ)$ The case in which $X$ and $S$ are strictly local, $f$ is a local morphism, and $Y$ consists of the closed point of $X$. Let
$$g:U=X-\{y\}\to S$$
be the structural morphism of $U$. We must show that the canonical morphism
$$u_i:H^i(X, f^*F)\to H^i(U, f^*F)$$
is bijective for $i\leq q-2$ and injective for $i=q-1$. Consider the commutative diagram
$$
\xymatrix{H^i(X, f^*F)\ar[rr]^{u_i}&& H^i(U, f^*F)\\&H^i(S, F)\ar[lu]^{v_i}\ar[ru]_{w_i}}$$
The morphism $v_i$ is plainly bijective for every $i$. Moreover, $g$ is locally $(q-2)$-acyclic for $\LL$; its fibers are also $(q-2)$-acyclic for $\LL$, as follows from the fact that $\prof_y(X_s)\geq q$ and that the fibers of $f$ are $(q-2)$-acyclic for~$\LL$. Since $g$ is quasi-compact because $X$ is Noetherian, (\SGA 4 XV 1.16) shows that $g$ is $(q-2)$-acyclic for $\LL$. Consequently $w_i$, and hence also $u_i$, is bijective for $i\leq q-2$ and injective for $i=q-1$, which completes the proof\pageoriginale of \Ref{XIV.1.14}.

\begin{corollaire} \label{XIV.1.15}
Let $f:X\to S$ be a morphism of schemes, let $\LL$ be a set of prime numbers, let $m$ and $r$ be integers, and let $F$ be a complex of abelian sheaves of $\LL$-torsion on $S$ such that $\h^i(F)=0$ for $i<-r$. Let $x$ be a point of $X$, let $s=f(x)$, and suppose that the local ring $\Oo_{X,x}$ is Noetherian. Then, if $f$ is locally $m$-acyclic for $\LL$, one has
\begin{equation*} \label{eq:XIV.115..*} \tag{$*$} {\prof_x(f^*F)\geq\inf(\prof_s(F)+\prof^\LL_x(X_s), n)\text{ where }n=m-r+2.}
\end{equation*}
In particular, if $n\geq\prof_s(F)+\prof_x^\LL(X_s)$, for example if $f$ is locally acyclic for $\LL$, then
\begin{equation*} \label{eq:XIV.15.**} \tag{$**$} {\prof_x(f^*F)\geq\prof_s(F)+\prof_x^\LL(X_s).}
\end{equation*}
If $\LL$ consists of a single element $\ell$ and $n\geq \prof_s(F)+\prof_x^\LL(X_s)$, the preceding inequality is an equality.
\end{corollaire}

One reduces to the case in which $s$ and $x$ are closed points by taking the strict localizations of $S$ and $X$ at geometric points $\bar s$ above $s$ and $\bar x$ above $x$ and $\bar s$. If $n\geq\prof_s(F)+\prof_x^\LL(X_s)$, then ($*$) follows from \Ref{XIV.1.13} by taking $p=\prof_s(F)$ and $q=\prof_x^\LL(X_s)$ there (the hypothesis that $S-\{s\}$ is quasi-compact in $S$ follows from the fact that $X-X_s$ is quasi-compact in $X$ and that $f$ is surjective, since it is $(-1)$-acyclic, except possibly when the conclusion of \hyperref[XIV.1.15]{1.15} is empty). If $n<\prof_s(F)+\prof_x^\LL(X_s)$, inequality \eqref{eq:XIV.115..*} again follows from \Ref{XIV.1.13}, for example by taking $p=\prof_s(F)$ and $q=n-p$. It remains to prove\pageoriginale the last assertion. Let $p=\prof_s(F)$ and $q=\prof_x(X_s)$; it follows from \eqref{eq:XIV.1.13.2} that
$$
\h_x^{p+q}(f^*(F))\simeq \h_x^q(f^*(\h^p_s(F))).
$$
Since $\prof_s(F)=p$, the sheaf $\h^p_s(F)$ is a nonzero sheaf of $\ell$-torsion, constant on~$s$. Consequently the sheaf $G=f^*(\h_s^p(F))$ is a nonzero sheaf of $\ell$-torsion, constant on $X_s$, and therefore contains a subsheaf isomorphic to $\ZZ/\ell\ZZ$; since $\h_x^q(\ZZ/\ell\ZZ)$ is nonzero, one indeed has $\h_x^q(G)\neq 0$.

\begin{corollaire} \label{XIV.1.16}
Let $f:X\to S$ be a regular morphism of excellent schemes (\textup{\EGA IV 7.8.2}) of characteristic zero, let $\ell$ be a prime number, and let $F$ be a complex of sheaves of $\ell$-torsion on $S$. Let $x\in X$ and $s=f(x)$; then
$$
\prof_x(f^*F)=\prof_s(F)+2\dim(\Oo_{X_s,x}).
$$
\end{corollaire}

Indeed, $f$ \sisi{is regular by the definition of \og excellent\fg and hence is}{is} locally acyclic (\SGA 4 XIX 4.1). It then follows from \Ref{XIV.1.15} that
$$
\prof_x(f^*F)=\prof_s(F)+\prof_x(X_s).
$$
But by \Ref{XIV.1.10},
$$
\prof_x(X_s)=2\dim\Oo_{X_s,x},
$$
which gives the result.

\begin{remarque} \label{XIV.1.17}
It follows from \Ref{XIV.1.15} that \Ref{XIV.1.13} remains valid when b) and c) are replaced\pageoriginale by the conditions:

b$'$) For every point $s\in f(Y)$, one has $\prof_s(F)\geq p$.

c$'$) For every point $x\in Y$, if $s=f(x)$, one has $\prof^\LL_x(X_s)\geq q$.
\end{remarque}

For a sheaf of sets or groups, one has the following theorem, analogous to \Ref{XIV.1.13}.

\begin{theoreme} \label{XIV.1.18}
Let $f:X\to S$ be a morphism of schemes and let $Y$ be a closed subset of $X$ such that $X-Y$ is quasi-compact in $X$ and, for every point $x$ of $Y$, the local ring $\Oo_{X,x}$ is Noetherian.

$1^\circ)$ Let $F$ be a sheaf of sets on $S$ and let $n$ be an integer equal to $1$ or $2$. Suppose that $f$ is locally $(n-2)$-acyclic and that, for every point $s$ of $f(Y)$,
$$
\prof_{Y_s}(X_s)+\prof_s(F)\geq n.
$$
Then:$$
\prof_Y(f^*F)\geq n.
$$

$2^\circ)$ Let $\LL$ be a set of prime numbers and let $F$ be a sheaf of ind-$\LL$-groups. Suppose that $f$ is locally $1$-aspherical for $\LL$ (\textup{\SGA 4 XV 1.11}) and that, for every point $s$ of $f(Y)$,
$$
\prof\hop^\LL_{Y_s}(X_s)+\prof_s(F)\geq 3.
$$
Then:$$
\prof_Y(f^*F)\geq 3.
$$
\end{theoreme}

\begin{proof}
We\pageoriginale reduce, as in \Ref{XIV.1.14} and \Ref{XIV.1.15}, to the case where $X$ and $S$ are strictly local schemes, $f$ is a local morphism, and $Y$ is the closed point $x$ of $X$. Let $s=f(x)$ be the closed point of $S$; we have the commutative diagram:
$$
\xymatrix{ X-X_s\ar[r]^i\ar[d]_h&X-\{x\}\ar[rd]_g\ar[rr]^j&&X\ar[ld]^f\\
S-\{s\}\ar[rr]_k&&S. }$$

$1^\circ)$ a) \textit{Case} $n=1$.

If $\prof_s(F)\geq 1$, then the morphism $F\to k_*k^*F$ is injective, and hence the morphism $f^*F\to f^*(k_*k^*F)$ is also injective. On the other hand, it follows from the fact that $f$ is locally $(-1)$-acyclic that the morphism $f^*(k_*k^*F)\to (j.i)_*(f^*\sisi{F{|X-X_s}}{F_{|X-X_s}}))$ is injective. Finally, the composite morphism $f^*F\to (j.i)_*(f^*\sisi{F{|X-X_s}}{F_{|X-X_s}}))$ is injective, which shows that $\prof_{X_s}(f^*F)\geq 1$, and consequently also $\prof_x(f^*F)\geq 1$.

If $\prof_x(X_s)\geq 1$, consider the commutative diagram
\begin{equation*} \label{eq:XIV.18.*} \tag{$*$}
\begin{array}{c}
\xymatrix{
H^0(X, f^*F)\ar[r]^-{v}\ar[d]_\wr&H^0(X-\{x\}, f^*F)\ar[d]\\
H^0(X_s, f^*F)\ar[r]^-{v'}&H^0(X-\{x\}, f^*F);
}
\end{array}
\end{equation*}
By hypothesis, $v'$ is injective, and therefore so is $v$.

b) \textit{Case} $n=2$. Consider the commutative diagram
\begin{equation*} \label{eq:XIV.1.18.**} \tag{$**$}
\begin{array}{c}
\xymatrix{ H^0(S, F)\ar[d]_{m}^{\wr}\ar[rr]^u&&H^0(S-\{s\}, F)\ar[d]^{\wr}_{ n}\\
H^0(X, f^*F)\ar[r]^-v&H^0(X-\{x\}, f^*F)\ar[r]^w&H^0(X-X_s, f^*F);
}
\end{array}
\end{equation*}
we must\pageoriginaled show that $v$ is bijective. The morphism $m$ is obviously bijective, and, since $f$ is $0$-acyclic, $n$ is also bijective.

If $\prof_s(F)\geq 2$, $u$ is bijective. As we saw in a), the sole hypothesis $\prof_s(F)\geq 1$ implies the relation $\prof_{X_s}(f^*F)\geq 1$; consequently $v$ and $w$ are injective, and it then follows from \eqref{eq:XIV.1.18.**} that $v$ is bijective.

If $\prof_x(X_s)\geq 2$, then $g$ is $0$-acyclic (because it is locally $0$-acyclic and its fibers are $0$-acyclic). It follows that $v\cdot m$ is bijective, and therefore $v$ is bijective.

If $\prof_s(F)\geq 1$ and $\prof_x(X_s)\geq 1$, then we already know that $v$ and $w$ are injective. Let $z$ be a maximal point of $X_s-\{x\}$ (such a point exists by the hypothesis $\prof_x(X_s)\geq 1$), let $\bar Z$ be the strict localization of $X$ at a geometric point above $z$, and let $\overline{f^*F}$ be the inverse image of $f^*F$ on $\bar Z$. Consider the commutative diagram
\begin{equation*}
\xymatrix@C=0mm{
&H^0(S, F)\ar[rr]\ar[ld]_{m}^{\sim}&&H^0(S-\{s\}, F)\ar[rd]_{\sim}^{n}\\
H^0(X, f^*F)\ar[rr]^-{v}\ar[rd]_{m'}^{\sim}&&H^0(X-\{x\}, f^*F)\ar[rr]^-{w}\ar[ld]^r&& H^0(X-X_s, f^*F)\ar[ld]^{n'}_{\sim}\\
&H^0(\bar Z, f^*F)\ar[rr]&&H^0(\bar Z-\{\bar z\}, \overline{f^*F})}
\end{equation*}
the morphism $m'\cdot m$ is obviously bijective, and it follows from the fact that $f$ is locally $0$-acyclic that $n'\cdot n$ is bijective; consequently $m'$ and $n'$ are also bijective. Since $w$ is injective, $r$ is also injective, and therefore $v$ is bijective.

$2^\circ)$ In view of b), we already know that $\prof_x(f^*F)\geq 2$.

If\pageoriginale $\prof_s(F)\geq 3$, then $\R^1k_*(k^*F)=1$\nde{the trivial torsor is subsequently denoted successively by $0$ or $1$; we have retained this double notation, which, on reflection, causes no ambiguity.}. Since $f$ is locally $1$-aspherical, one has $\R^1(j\cdot i)_*(f^*\sisi{F{|X-X_s}}{F_{|X-X_s}})\!=\!f^*(\R^1k_*(k^*F))\!=\!1$. Thus \hbox{$\prof_{X_s}(f^*F)\!\geq\! 3$} and consequently $\prof_x(f^*F)\geq 3$.

If $\prof_x(X_s)\geq 3$, then $g$ is $1$-aspherical (because $g$ is locally $1$-aspherical and its fibers are $1$-aspherical). Thus $H^1(X-\{x\}, f^*F)=H^1(S, F)=1$, and consequently $\prof_x(f^*F)\geq 3$.

If $\prof_s(F)\geq 2$ and $\prof_x(X_s)\geq 1$, we use the exact sequence (\SGA 4 XII 3.2):
$$1\to\R^1j_*(i_*(f^*\sisi{F{|X-X_s}}{F_{|X-X_s}}))\to\R^1(j.i)_*(f^*\sisi{F{|X-X_s}}{F_{|X-X_s}}) \to j_*(\R^1i_*(f^*\sisi{F{|X-X_s}}{F_{|X-X_s}})).
$$
Since $f$ and $g$ are locally $1$-aspherical, one has
\begin{align*}
\R^1(j\cdot i)_*(f^*\sisi{F{|X-X_s}}{F_{|X-X_s}})&\simeq f^*(\R^1k_*(k^*F)) \\
\R^1i_*(f^*\sisi{F{|X-X_s}}{F_{|X-X_s}})&\simeq g^*(\R^1k_*(k^*F));
\end{align*}
the preceding exact sequence can then be written in the form
\enlargethispage{\baselineskip}%
\begin{equation} \label{eq:XIV.18.***}\tag{${*}{*}{*}$}
1\to\R^1j_*(i_*(f^*\sisi{F{|X-X_s}}{F_{|X-X_s}}))\to f^*(\R^1k_*(k^*F))\lto{a} j_*(j^*(f^*(\R^1k_*(k^*F))).
\end{equation}
The hypothesis $\prof_s(F)\geq 2$ shows that the morphism $F\to k_*k^*F$ is bijective; applying $g^*$, and taking into account the fact that $g$ is locally $0$-acyclic, we find $f^*\sisi{F{|X-\{x\}}}{F_{|X-\{x\}}}=i_*(f^*\sisi{F{|X-X_s}}{F_{|X-X_s}})$. The hypothesis $\prof_x(X_s)\geq 1$ shows that the morphism~$a$ is injective (note that $f^*(\R^1k_*(k^*F))$ is a sheaf equal to $1$ outside $X_s$ and constant on $X_s$). It then follows from \eqref{eq:XIV.18.***} that $\R^1j_*(f^*\sisi{F{|X-\{x\}}}{F_{|X-\{x\}}})=1$, and hence $\prof_x(f^*F)\geq 3$\pageoriginale.

\enlargethispage{.8\baselineskip}%
If $\prof_s(F)\geq 1$ and $\prof_x(X_s)\geq 2$, consider the sheaf of homogeneous spaces $G$ defined by the exact sequence
$$
1\to F\to k_*k^*F\to G\to 1.
$$
Applying the exact functor $g^*$ to this exact sequence and using (\SGA 4 XII 3.1), we obtain the following commutative diagram with exact rows:
$$
\xymatrix@R=5mm{f^*(k_*k^*F)\ar[r]\ar[d]&f^*G\ar[r]\ar[d]^b&1\\
j_*(g^*(k_*k^*F))\sisi{\ar[r]}{\ar[r]^-{u}}&j_*(g^*G)\ar[r]&\R^1j_*(g^*F) \ar[r]&\R^1j_*(g^*(k_*k^*F)).}
$$
Since $\prof_x(X_s)\geq 2$, the morphism $b$ is bijective, so $u$ is surjective, and we thus have a map whose kernel is reduced to the neutral element:
$$
1\to\R^1j_*(g^*F)\to\R^1j_*(g^*(k_*k^*F))=R.
$$
Since $g^*(k_*k^*F)\simeq i_*(f^*\sisi{F{|X-X_s}}{F_{|X-X_s}})$ (because $g$ is locally $0$-acyclic), $R$ is identified with the first term of the exact sequence \eqref{eq:XIV.18.***}; but in the preceding case we saw that $R=1$ as soon as $\prof_x(X_s)\geq 1$, which proves that $\prof_x(f^*F)\geq 3$ and completes the proof of \Ref{XIV.1.18}.
\skipqed
\end{proof}

The following corollaries are generalizations of (\SGA 4 XVI 3.2 and 3.3).

\begin{corollaire} \label{XIV.1.19}
Let $f:X\to S$ be a flat morphism with separable fibers between locally Noetherian schemes, and let $Y$ be a closed subset of $X$. Suppose\pageoriginale that, for every point $s\in f(Y)$, the fiber $Y_s$ is rare\nde{\og rare\fg = \og having empty interior\fg, \cf Bourbaki TG~IX.52.} in $X_s$ and that one of the following two conditions holds:
\begin{enumeratea}
\item
the closure of $f(Y)$ is rare in $S$.
\item
$X_s$ is geometrically unibranch at the points of $Y_s$.
\par\noindent
Then
$$
\prof_Y(X)\geq 2.
$$
\end{enumeratea}
\end{corollaire}

Indeed, it follows from the hypothesis on $f$ that $f$ is locally $0$-acyclic (\SGA 4 XV 4.1). We then apply \Ref{XIV.1.13}. The hypothesis that $Y_s$ is rare in $X_s$ (respectively that $\overline{f(Y)}$ is rare in~$S$) is equivalent, by \Ref{XIV.1.6} b), to the relation $\prof_{Y_s}(X_s)\geq 1$ (respectively $\prof_{\overline{f(Y)}}(S)\geq 1$). The hypothesis that $X_s$ is geometrically unibranch at every point of $Y_s$ is equivalent to saying that the strict localization of $X_s$ at a geometric point of $Y_s$ is irreducible; knowing that $Y_s$ is rare in $X_s$, this clearly implies $\prof_{Y_s}(X_s)\geq 2$, by \Ref{XIV.1.8}. In either case \Ref{XIV.1.13} indeed gives $\prof_Y(X)\geq 2$.

\begin{corollaire} \label{XIV.1.20}
Let $f:X\to S$ be a regular morphism (\textup{\EGA IV 6.8.1}) between locally Noetherian schemes, and let $Y$ be a closed subset of $X$. Suppose that, for every point $s\in f(Y)$, one of the following conditions holds:
\begin{enumeratea}
\item
One has $\codim(Y_s, X_s)\geq 2$.
\item
One has $\codim(Y_s, X_s)\geq 1$ and $\prof_s(S)\geq 1$.
\item
One has $\prof\hop_s(S)\geq 3$.
\end{enumeratea}
\noindent
Then\pageoriginale
$$
\prof\hop_Y(X)\geq 3.
$$
\end{corollaire}

Indeed, this follows from \Ref{XIV.1.18}, since hypothesis a) implies $\prof\hop_{Y_s}(X_s)\geq 3$ (\cf \Ref{XIV.1.11}), and the condition $\codim(Y_s, X_s)\geq 1$ clearly implies $\prof_Y(X)\geq 2$.

\section{Technical lemmas} \label{XIV.2}

\subsection{\texorpdfstring{}{XIV.2.1}} \label{XIV.2.1}
Let $S$ be a locally noetherian scheme, $f:X\to S$ a morphism locally of finite type, and $t$ a point of $S$. If $x\in X$ is such that $s=f(x)\in\Spec\Oo_{S, t}$, set
$$
\delta_t(x)=\degtr k(x)/k(s)+\dim(\overline{\{s\}}),
$$
where $\overline{\{s\}}$ denotes the closure of $s$ in $\Spec\Oo_{S, t}$, and $k(x)$ and $k(s)$ the residue fields of $x$ and $s$, respectively. If $S$ is a local ring with closed point $t$, we also write $\delta(x)$ instead of $\delta_t(x)$ (\cf \SGA 4 XIV 2.2).
\begin{sublemme} \label{XIV.2.1.1} Let there be a cartesian square
$$
\xymatrix{X'\ar[r]^{h}\ar[d]_{f'}&X\ar[d]^f\\S'\ar[r]^g&S},
$$
where $S$ and $S'$ are noetherian local rings, with closed points $t$ and $t'$, respectively, $g$ is a faithfully flat morphism such that $g^{-1}(t)=t'$\pageoriginale, and $f$ is a morphism locally of finite type. Let $x'\in X',\ x=h(x'),\ s=f(x),\ s'=f'(x')$; then one has
$$
\delta(x')\leq\delta(x).
$$
Moreover, the preceding inequality is an equality if and only if one has:
$$
\degtr k(x)/k(s)=\degtr k(x')/k(s')\text{ and }\dim(\overline{\{s\}})=\dim(\overline{\{s'\}}).
$$
In particular, given $x\in X$, one can find $x'$ such that $\delta(x)=\delta(x')$.
\end{sublemme}

Indeed, one has (\EGA IV 6.11)
$$
\dim(\overline{\{s\}})=\dim g^{-1}(\overline{\{s\}}).
$$
It follows that, for every point $s'$ of $g^{-1}(s)$, one has the relation $\dim(\overline{\{s'\}})\leq \dim(\overline{\{s\}})$, and that, given $s$, one can find $s'\in g^{-1}(s)$ for which equality holds. Let $Z$ then denote the schematic closure of $x$ in the fiber $X_s$ of $X$ at~$s$, and let $Z'=Z\times_{\Spec k(s)}\Spec k(s')$. Then $Z'$ is equidimensional of dimension $\degtr k(x)/k(s)$; therefore, for every point $x'\in Z'_{x}$, one has
$$
\degtr k(x')/k(s')\leq\degtr k(x)/k(s)\text{, with equality}$$
when $x'$ is a maximal point of $Z'_x$. The asserted conclusion follows at once.

\subsection{\texorpdfstring{}{XIV.2.2}} \label{XIV.2.2} Let $f:X\to S$ be a morphism locally of finite type and $T$ a closed subset of $S$. Let $x\in X, s=f(x)$; we set
$$
\delta_T(x)=\degtr k(x)/k(s)+\codim(\overline{\{s\}}\cap T, \overline{\{s\}})=\inf_{t\in T\cap\overline{\{s\}}}\delta_t(x).
$$

\begin{sublemme} \label{XIV.2.2.1}
Let\pageoriginale there be a cartesian square
$$
\xymatrix{X'\ar[r]^{h}\ar[d]_{f'}&X\ar[d]^f\\S'\ar[r]^g&S},
$$
where the schemes $S$ and $S'$ are locally noetherian and catenary, the morphism $f$ is locally of finite type, and $g$ is faithfully flat. Let $T$ be a closed subset of $S$ and $T'$ a closed subset of $S'$, such that $g(T')\subset T$, let $x'$ be an element of $X'$, let $x=h(x')$, and let
$$h_{x'}:\Spec\Oo_{X', x'}\to\Spec\Oo_{X, x}$$
be the morphism induced by $h$. Then one has:
$$
\delta_T(x)-\delta_{T'}(x')\leq\dim h_{x'}^{-1}(x).
$$
\end{sublemme}

Let $s'=f'(x'), s=f(x)$. By definition, one has:
\begin{multline*}
\delta_T(x)-\delta_{T'}(x')=\degtr k(x)/k(s)-\degtr k(x')/k(s')\\
+\codim(\overline{\{s\}}\cap T, \overline{\{s\}})-\codim(\overline{\{s'\}}\cap T', \overline{\{s'\}}).
\end{multline*}
Since $g$ is faithfully flat, it follows from (\EGA IV 6.1.4) that one has
\begin{align*} \label{eq:XIV.2.2.*} \tag{$*$}
\codim(\overline{\{s\}}\cap T, \overline{\{s\}})&=\codim(g^{-1}(\overline{\{s\}})\cap g^{-1}(T), g^{-1}(\overline{\{s\}}))\\&\leq\codim(g^{-1}(\overline{\{s\}})\cap T', g^{-1}(\overline{\{s\}}));
\end{align*}
since $S'$ is catenary, one has, by ($\EGA 0_{\textup{IV}}$ 14.3.2 b)):
\begin{multline*}
\codim(\overline{\{s'\}}\cap T', g^{-1}(\overline{\{s\}}))= \codim(\overline{\{s'\}}\cap T', \overline{\{s'\}})+ \codim(\overline{\{s'\}}, g^{-1}(\overline{\{s\}}))\\
= \codim(\overline{\{s'\}}\cap T', g^{-1}(\overline{\{s\}})\cap T')+ \codim(g^{-1}(\overline{\{s\}})\cap T', g^{-1}(\overline{\{s\}})).
\end{multline*}
From\pageoriginale this relation and \eqref{eq:XIV.2.2.*}, one deduces
$$
\delta_T(x)-\delta_{T'}(x')\leq\degtr k(x)/k(s)-\degtr k(x')/k(s')+\codim(\overline{\{s'\}}, g^{-1}(\overline{\{s\}})).
$$
Let us calculate $\codim(\overline{\{s'\}}, g^{-1}(\overline{\{s\}}))=\dim\Oo_{S'_s, s'}$ (where $S'_s$ is the fiber of $S'$ at $s$). Let $Z$ be the closed image of $x$ in $X_s$ and let $Z'\subset X'_s$ be the scheme defined by the cartesian square
$$
\xymatrix{Z'\ar[r]\ar[d]&Z\ar[d]\\S'_s\ar[r]&\Spec k(s)}.
$$
The morphism $Z\to \Spec k(s)$ is flat, locally of finite type, and $\dim Z=\degtr k(x)/k(s)$. It then follows from (\EGA IV 6.1.2) that

$$
\dim(\Oo_{Z', x'})=\dim(\Oo_{S'_s, s'})+\degtr k(x)/k(s)-\degtr k(x')/k(s');$$
taking into account the fact that $Z'_{s'}\simeq Z\otimes_{k(s)}k(s')$, one then obtains:
$$
\delta_T(x)-\delta_{T'}(x')\leq\dim(\Oo_{Z', x'}).
$$
But $\Spec(\Oo_{Z', x'})$ identifies with the fiber at $x$ of the morphism
$$(\Spec(\Oo_{X', x'}))_s\to(\Spec(\Oo_{X, x}))_s,
$$
and hence also with the fiber at $x$ of $h_{x'}$, which proves the sublemma.

\subsection{\texorpdfstring{}{XIV.2.3}} \label{XIV.2.3}
The proofs of the theorems in \numero \Ref{XIV.4} are based on duality theory; they use the following lemmas. Let $m$ be an\pageoriginale integer that is a power of a prime number $\ell$; if $X$ is a scheme, all the sheaves considered on $X$ are sheaves of $\ZZ/m\ZZ$-modules; there is then a notion of a dualizing complex on $X$ (\SGA 5 I 1.7). Suppose that such a complex $K$ exists on $X$; then, for every geometric point~$\bar x$ \sisi{over}{over} a point $x$ of $X$, one deduces from $K$ (\cf \SGA 5 I 4.5) a dualizing complex~$K_{\bar x}$ on $\Spec k(\bar x)$, so that $K_{\bar x}\simeq\ZZ/m\ZZ[n]$ (the brackets denoting the translation functor) for a certain integer $n$ depending only on $x$. We set
$$
\delta^K(x)=n.
$$
If $K$ is normalized at the point $x$ (\SGA 5 I 4.5), then $n=0$.

\begin{sublemme} \label{XIV.2.3.1} Let $X$ be a locally noetherian scheme endowed with a dualizing complex $K$. If $x$ and $x'$ are two points of $X$, such that $x$ is a specialization of $x'$ and $\codim(\overline{\{x\}}, \overline{\{x'\}})=1$, then one has
$$
\delta^K(x)=\delta^{K}(x')-2.
$$
\end{sublemme}

One may first reduce to the case where $X$ is a strictly local scheme. Indeed, let $\bar X$ be the strict localization of $X$ at a geometric point $\bar x$ over $x$, let $i:\bar X\to X$ be the canonical morphism, and let $\bar x'$ be a geometric point of $\bar X$ over $x'$. Then $i^*K$ is a dualizing complex on $\bar X$, and one has (\SGA 5 I 4.5)
$$
(i^*K)_{\bar x}\simeq K_{\bar x}\text{ and } (i^*K)_{\bar x'}\simeq K_{\bar x'},
$$
which\pageoriginale completes the reduction to the strictly local case.

If $j:\overline{\{x'\}}\to X$ denotes the immersion of the reduced closed subscheme of $X$ whose underlying space is $\overline{\{x'\}}$, then $\R^!j(K)$ is a dualizing complex on $\overline{\{x'\}}$, and one sees at once, using (\SGA 5 I 4.5), that it suffices to prove the lemma for $\overline{\{x'\}}$. One is thus reduced to the case where $X$ is an integral strictly local scheme of dimension $1$.

Let $X'$ then be the normalization of $X$ and let $f:X'\to X$ be the canonical morphism; $f$ is integral, surjective, and radicial, and it follows that $f^*K$ is a dualizing complex on $X'$ and that it suffices to prove the lemma for $X'$ and for the points \sisi{over}{over}~$x$ and~$x'$. One is thus reduced to the case where $X$ is a local, integral, regular scheme of dimension $1$, but it is known (\cf \SGA 5 I 4.6.2 and 5.1) that then $\mu_m[2]$ and $\ZZ/m\ZZ$ are dualizing complexes, normalized at the points $x$ and $x'$, respectively; the lemma follows at once.

\begin{sublemme} \label{XIV.2.3.2} Let $S$ be a noetherian local scheme and $f:X\to S$ a morphism of finite type. If $K$ is a dualizing complex on $S$, normalized at the closed point $t$ of $S$, and if $\R^!f(K)=K'$ is a dualizing complex on $X$ (\cf \SGA 5 I 3.4.3), one has, for every point $x$ of $X$:$$
\delta^{\sisi{K}{{K'}}}(x)=2\delta(x).
$$
\end{sublemme}

Indeed, let $s=f(x)$ and let $x'$ be a closed point of the fiber $X_s$; then $\delta^{K'}(x')=\delta^K(s)$ and, by \Ref{XIV.2.3.1},
$$
\delta^K(s)=2\codim(\overline{\{t\}}, \overline{\{s\}})=2\dim(\overline{\{s\}}).
$$
Since\pageoriginaled $x'$ may be chosen to be a specialization of $x$, one has, by \Ref{XIV.2.3.1},
$$
\delta^{K'}(x)=\delta^{K'}(x')+2\codim(\overline{\{x'\}}, \overline{\{x\}})=\delta^{K'}(x')+2\degtr k(x)/k(s);$$
the lemma follows at once.

The following lemma will be used only for the converse of the Lefschetz theorem, in \numero\Ref{XIV.4}:

\begin{sublemme} \label{XIV.2.3.3}
Let there be a cartesian square
$$
\xymatrix{X'\ar[r]^{h}\ar[d]_{f'}&X\ar[d]^f\\S'\ar[r]^g&S},
$$
where $S$ is an excellent strictly local scheme of characteristic zero, $S'$ is the completion of $S$, and $f$ is a morphism of finite type. Let $\ell$ be a prime number, $x\in X$, let $Z$ be the schematic closure of $X'_x$ in $X'$, and let $i:X'_x\to Z$, $j:Z\to X'$ be the canonical morphisms. Then, if $k:X'\to R$ is a closed immersion of $X'$ into an excellent regular scheme $R$ of characteristic zero, the complex
$$K'=i^*(R^!(k.j)(\ZZ/\ell\ZZ))$$
is a constant dualizing complex on $X'_x$ (that is, it has only one nonzero cohomology sheaf, isomorphic to $\ZZ/\ell\ZZ$).
\end{sublemme}

\enlargethispage{\baselineskip}%
In view of (\SGA 5 I 3.4.3), the only thing to prove is that $K'$ is constant. But, since $Z$ is excellent, the set of points of $Z$ whose\pageoriginale local rings are regular is an open subset $U$ (\EGA IV 7.8.3 \textup{(iv)}), and $U$ plainly contains $X'_x$, which is regular. Let
$$u:U\to R$$
then be the canonical immersion of $U$ into $R$; it follows from the purity theorem (\SGA 4 XIX 3.2 and 3.4) and the isomorphism
$$(\mu_\ell)_S\simeq(\ZZ/\ell\ZZ)_S$$
($S$ strictly local) that one has
$$
\R^!u(\ZZ/\ell\ZZ)\simeq\ZZ/\ell\ZZ[2c],
$$
where $c$ is a locally constant function on $U$, necessarily constant in a neighborhood of $X'_x$, because the fibers of $g$ are geometrically integral by (\EGA IV 18.9.1), and hence $X'_x$ is integral. The lemma follows at once.

\section{Converse to the affine Lefschetz theorem} \label{XIV.3}

The present section will be used in \numero \Ref{XIV.4} to prove a converse to the \og Lefschetz theorem\fg; a reader interested only in the direct part of that theorem may therefore omit the present section.

\subsection{\texorpdfstring{}{XIV.3.1}} \label{XIV.3.1}

Recall the statement of the affine Lefschetz theorem\nde{\label{gabaf}Gabber proved the following generalization. Let $Y$ be a strictly local scheme of arithmetic type over a regular Noetherian scheme $S$ of dimension $\leq 1$. Let $f:X\to Y$ be an affine morphism of finite type, $\Lambda=\ZZ/n\ZZ$ with $n$ invertible on $X$, and $F$ a $\Lambda$-sheaf. Then $H^q(X,F)=0$ if $q>\delta(F)$. The following local Lefschetz theorem follows. Let $\Oo$ be strictly local of arithmetic type over $S$. For every non-zero-divisor $f\in \Oo$ and every $\Lambda$-sheaf $F$ on $\Spec(\Oo[f^{-1}])$, one has $H^q(\Spec(\Oo[f^{-1}]),F)=0$ for $q>\dim(\Oo)$. \Cf (Fujiwara~K., {\og A Proof of the Absolute Purity Conjecture (after Gabber)\fg}, in \emph{Algebraic geometry 2000, Azumino (Hotaka)}, Adv. Stud. in Pure Math., vol.~36, 2002, p\ptbl 153-183, $\S$5), and especially Illusie's article (Illusie~L., {\og Perversité et variation\fg}, \emph{Manuscripta Math.} \textbf{112} (2003), p\ptbl 271-295). This result is one of the crucial points used by Gabber in his proof of Grothendieck's purity theorem (\cf note~\eqref{gabpur}, page~\pageref{gabpur}).} (\SGA 4~XIX~6.1~bis):

Let $S$ be a strictly local excellent scheme of characteristic zero, $f:X\to S$ an affine morphism of finite type, and $F$ a torsion sheaf\pageoriginale on $X$. Then, setting
$$
\delta(F)=\sup\{\delta(x)|x\in X \text{ and } F_{\bar{x}} \neq 0\},
$$
one has
$$
H^q(X, F)=0 \text{ for } q>\delta(F).
$$

Before stating the converse, let us prove a few lemmas.

\begin{lemme} \label{XIV.3.2} Let $K$ be a field, $\ell$ a prime number distinct from the characteristic of $K$, and $F$ a nonzero constructible sheaf of $\ell$-torsion on $K$. Suppose that the $\ell$-cohomological dimension of $K$ (\textup{\SGA 4 X~1}) is equal to $n$ (this holds, for example, if $K$ is the field of fractions of an integral strictly local excellent ring of characteristic zero and dimension $n$ (\textup{\SGA 4~XIX~6.3}), or if $K$ is an extension of finite type and transcendence degree $n$ over a separably closed field (\textup{\SGA 4~X~2.1})). Then one can find a finite separable extension $L$ of $K$ such that:
$$H^n(L, \sisi{F{|L}}{F_{|L}}) \neq 0.
$$
\end{lemme}

One can find a finite extension $K'$ of $K$ such that the restrictions of $F$ and $\mu_\ell$ to $\Spec K'$ are constant sheaves. Then $\cd_{\ell}(K') = \cd_{\ell}(K) = n$ (\SGA 4~X~2.1), and it follows from (\cite{XIV.2}~II~\S 3~Prop\ptbl4~\textup{(iii)}) that one can find a finite extension $L$ of $K'$ such that\pageoriginale
$$
H^n(L, \mu_{\ell})\neq 0, \quad\text {\ie } H^n(L, \ZZ/\ell\ZZ) \neq 0.
$$
Now the functor $H^n(L,\cdot)$ is right exact on the category of sheaves of $\ell$-torsion, since $\cd_{\ell}(L)=n$; as $F$ has a quotient isomorphic to $\ZZ/\ell\ZZ$, one also has $H^n(L, \sisi{F{|L}}{F_{|L}}) \neq 0$.

\begin{corollaire} \label{XIV.3.3}
Let $k$ be a field, $K$ an extension of finite type and transcendence degree $n$ over $k$, and $F$ a nonzero constructible sheaf of $\ell$-torsion on $K$, with $\ell$ prime to the characteristic of $k$. Then one can find a finite separable extension~$L$ of~$K$ such that, if $u: \Spec L \to \Spec k$ denotes the canonical morphism, one has
$$R^n u_*(\sisi{F{|\Spec L}}{F_{|\Spec L}}) \neq 0.
$$
\end{corollaire}

When the field $k$ is separab\sisi{le}{ly} closed, the corollary is a special case of \Ref{XIV.3.2}. In the general case, one can find a finite separable extension $k_1$ of $k\sisi{'}{}$ such that the irreducible components of $K \otimes_k k_1$ are geometrically irreducible (\EGA IV~4.5.11); let $K_1$ be one of them. If $k'$ is a separable closure of $k_1$, then $K' = K_1 \otimes_{k_1} k'$ is a field, and by (\EGA IV~4.2) one has
$$
\deg \tr K'/k' = \deg \tr K/k = n.
$$
It then follows from \Ref{XIV.3.2} that one can find a finite separable extension $L'$ of $K'$ such that $H^n(L', \sisi{F{|L'}}{F_{|L'}}) \neq 0$. But $k' = \varinjlim_i k_i$, where $k_i$ runs through the finite extensions of $k_1$ contained in $k'$, and consequently\pageoriginale $K' = \varinjlim_i (k_i \otimes_{k_1} K_1)$. It follows that one can find an index $i$ and a finite separable extension $L$ of $k_i \otimes_{k_1} K_1 =K_i$ such that $L' \simeq L \otimes_{K_i} K'$. The extension $L$ of $K$ has the required property; indeed, it follows from the commutative diagram
$$
\xymatrix{
&{\Spec L}\ar[dl]_v\ar[dr]^u&\\
{\Spec k_i} \ar[rr]^w &{}&{\Spec k}\,,
}
$$
with $w$ finite and hence $R^q w_* = 0$ if $q>0$, that
$$R^n u_*(\sisi{F{|\Spec L}}{F_{|\Spec L}}) \simeq w_*(R^n v_*(\sisi{F{|\Spec L}}{F_{|\Spec L}})).
$$
Now $R^nv_*(\sisi{F{|\Spec L}}{F_{|\Spec L}}) \neq 0$, since $H^n(L', \sisi{F{|L'}}{F_{|L'}}) \neq 0$; therefore one also has $R^nu_*(\sisi{F{|\Spec L}}{F_{|\Spec L}}) \neq 0$.

Recall the following familiar lemma (\cf $\EGA 0_{\textup{III}}$~10.3.1.2 and \EGA IV~18.2.3):

\begin{lemme} \label{XIV.3.4}
Let $X$ be a scheme, $x$ a point of $X$, and $K$ a finite separable extension of $k(x)$. Then there exist a scheme $X_1$ étale over $X$ and affine, and a point $x_1 \in X_1$ above $x$, such that $k(x_1)$ is $k(x)$-isomorphic to $K$.
\end{lemme}

In \numero \Ref{XIV.4} we shall use the following technical form of the converse to \Ref{XIV.3.1}.

\begin{proposition} \label{XIV.3.5}
Let\pageoriginale there be a Cartesian square
$$
\xymatrix{
{X'}\ar[r]^h\ar[d]_{f'} &{X}\ar[d]^f\\
{S'}\ar[r]^g &{S}\,,
}
$$
where the schemes $S$ and $S'$ are strictly local excellent schemes of characteristic zero, the morphism $f$ is locally of finite type, $g$ is regular (\textup{\EGA IV~6.8.1}) and surjective, and the closed fiber of $g$ is reduced to the closed point of $S'$. Given an $S$-scheme $X_1$ (respectively an $S$-morphism $f_1$, etc.), we denote by $X_1'$ (respectively $f_1'$, etc.) the scheme $X_1 \times_S S'$ (respectively the morphism $(f_1)_{(S')}$, etc.). Let $F$ be a sheaf of $\ZZ/m\ZZ$-modules on $X'$ ($m$ a power of a prime number $\ell$), constructible and satisfying the following conditions:
\begin{enumeratei}
\item
For every point $x \in X$, one can find a finite separable extension~$K$ of $k(x)$ such that the restriction of $F$ to the fiber $(X')_{(\Spec K)}$ comes by inverse image from a constructible sheaf on $\Spec K$.
\item
For every morphism $f_1: X_1 \to S$, with $X_1$ affine and étale over $X$, for every point $s \in S$, and for every integer $q>0$, one can find a finite separable extension~$K$ of $k(s)$ such that the restriction of $R^q f'_{1*}(F|X'_1)$ to the fiber $S'_{(\Spec K)}$ comes by inverse image from a constructible sheaf on $\Spec K$.
\end{enumeratei}

Let $n$ be an integer, and suppose that for every affine scheme $X_1$ étale over $X$ one has\pageoriginale
$$
H^i(X_1', F)=0 \text{ for } i>n.
$$
Then, if $\bar{x}'$ is a geometric point above the point $x' \in X'$ such that $F_{\bar{x}'} \neq 0$, one has
$$
\delta(x') \leq n.
$$
\end{proposition}

Let $Z'$ be the set of points $x'$ of $X'$ such that $F_{\bar{x}'} \neq 0$. Then, if $Z=h(Z')$, one has, by \textup{(i)}, $Z' = h^{-1}(Z)$; let $x' \in X'$, $x=h(x')$, $s'=f'(x')$, $s=f(x)$. It follows from \Ref{XIV.2.1.1} and the fact that the function $\delta$ decreases under specialization that it suffices to prove the inequality $\delta(x') \leq n$ when $x$ is a maximal point of $Z$ and $x'$ is such that
$$r=\deg \tr k(x)/k(s) = \deg \tr k(x')/k(s') \text{ and }d=\dim \overline{\{s\}} = \dim\overline{\{s'\}}$$
Let $x'$ be such a point; it suffices to show that one can find an affine scheme $X_1$ \'{e}tale over~$X$ such that
$$
H^{d+r}(X_1', F) \neq 0.
$$

The set $Z'$ is constructible (\SGA 4~IX~2.4), and hence so is $Z$ (\EGA IV~1.9.12); one may then suppose, after restricting $X$ to a neighborhood of~$x$, that $Z$ is an irreducible closed subset with generic point $x$. Let $T=f(Z)$; $T$ is a constructible set contained in $\overline{\{s\}}$; one can therefore find an affine open subset $U$ of~$S$ such that $s \in U$ and $T \cap U = T_U$ is an irreducible closed subset of $U$ with generic point~$s$.

Let\pageoriginale $V$ then be an affine scheme \'{e}tale over $X$, whose image in $X$ contains $x$ and whose image in $S$ is contained in $U$; let $Z_V$ be the inverse image of $Z$ in $V$ and let $u: Z_V \to T_U$ be the canonical morphism. Let\sisi{}{} $W$ be an affine scheme \'{e}tale over $U$; denote by $T_W$ the inverse image of $T_U$ in $W$ and let $X_1 = W\times_U V$. Since $F$ is zero outside $Z'$, one has the spectral sequence
$$
E^{pq}_2 = H^p((T_W)', R^qu_*'(F\sisi{}{_{|(Z_V)'}})) \To H^*(X'_1, F).
$$
We shall show that $V$ and $W$ may be chosen so that
\begin{enumeratea}
\item
$E^{pq}_2=0$ for $p>d$ and for $q>r$.
\item
$E^{dr}_2 \neq 0$.
\end{enumeratea}
\noindent
The relation $H^{d+r}(X_1', F) \neq 0$ will then follow from the spectral sequence.
\begin{enumerate}
\item[$1^\circ)$]
Set $G_q = R^q u'_*(F\sisi{}{_{|(Z_V)'}})$; then one has:
$$(G_q)_{\bar{s}'}=H^q((Z_V)'_{\bar{s}'}, F\sisi{}{{_{|(Z_V)'}}}_{\bar{s}'}),
$$
because $s'$ is a maximal point of $(T_U)'$. Since the fiber $(Z_V)'_{\bar{s}'}$ is an affine scheme of finite type of dimension $r$ over a separably closed field, it follows from \Ref{XIV.3.1} that
$$(G_q)_{\bar{s}'}=0 \text{ for } q>r.
$$

For $q>r$, let $Y'_q$ be the set of points of $(T_U)'$ where the geometric fiber\pageoriginale of $G_q$ is $\neq 0$, and let $Y_q=g(Y'_q)$; then $Y'_q = g^{-1}(Y_q)$ by \textup{(ii)}, so $Y_q$ is a constructible subset of $T_U$ (\SGA 4~XIX~5.1 and \EGA 1.9.12) that does not contain $s$; after restricting $U$ to an open neighborhood of $s$, one may suppose that $G_q=0$ for $q>r$, and hence that $E_2^{pq}=0$ for $q>r$.

Moreover, since $(T_W)'$ is an affine scheme of finite type over $g^{-1}(\overline{\{s\}})$, one has, for every $q$ (\cf \Ref{XIV.3.1}):
$$H^p((T_W)', G_q)=0 \text{ for } p>\dim g^{-1}(\overline{\{s\}})=d,
$$
whence condition a).

\item[$2^\circ)$]
Let us show that $V$ may be chosen so that $(G_r)_{\bar{s}'} \neq 0$. By~\textup{(i)}, there exists a constructible sheaf $I$, defined over a finite separable extension $K$ of $k(x)$, whose inverse image on $(X')_{(\Spec K)}$ is isomorphic to $\sisi{F{|(X')_{(\Spec K)}}}{F_{|(X')_{(\Spec K)}}}$. By \Ref{XIV.3.3}, one finds a finite separable extension $L$ of $K$ such that, if $v:\Spec L \to \Spec k(s)$ denotes the canonical morphism, then $R^r v_*(I) \neq 0$. Since the morphism $S'_s \to \Spec k(s)$ is regular, one has, by (\SGA 4~XIX~4.2):
$$R^rv_*'(\sisi{F{|(\Spec L)'}}{F_{|(\Spec L)'}}) \simeq (R^rv_*(I))' \neq 0.
$$
By Lemma \Ref{XIV.3.4}, one can find an affine scheme $X_2$ \'{e}tale over $X$, and a point~$x_2$ of $X_2$ over $x$, such that $L$ is $k(x)$-isomorphic to $k(x_2)$, and one may suppose that $X_2$ lies over $U$. Since $x$ is a maximal point of $Z$, one has\pageoriginale
$$
\Spec L \simeq \varprojlim_V Z_V,
$$
where $V$ ranges over the affine open neighborhoods of $x_2$. By passage to the limit (\SGA 4~VII~5.8), after restriction to the geometric fiber at $\bar{s}'$, one deduces:
$$
(R^rv_*'(\sisi{F{|\Spec L}}{F_{|\Spec L}})')_{\bar{s}'} = \varinjlim_V (R^ru_*'(F\sisi{}{_{|(Z_V)'}})\sisi{}{)}_{\bar{s}'},
$$
which indeed shows that one can find $V$ such that $(G_r)_{\bar{s}'} \neq 0$.

\item[$3^\circ)$]
With the scheme $V$ having been chosen in $2^\circ)$, let us show that the scheme~$W$ may be chosen so that
$$
E^{dr}_2 = H^d((T_W)', G_r) \neq 0.
$$
By \textup{(ii)}, there exists a constructible sheaf $J$, defined over a finite separable extension $K$ of $k(s)$, whose inverse image on $(S')_{(\Spec K)}$ is isomorphic to $\sisi{{G_r}{|(S')_{(\Spec K)}}}{{G_r}_{|(S')_{(\Spec K)}}}$. By Lemma \Ref{XIV.3.2}, one can find a finite separable extension $L$ of $K$ such that $H^d(\Spec L, J) \neq 0$. Since the morphism $(S')_{(\Spec L)} \to \Spec L$ is acyclic (\SGA 4~XIX~4.1 and XV~1.10 and 1.16), one has
$$
H^d((\Spec L)', \sisi{{G_r}{|(S')_{(\Spec L)'}}}{{G_r}_{|(S')_{(\Spec L)'}}}) = H^d(\Spec L, J) \neq 0.
$$
By \Ref{XIV.3.4}, one can find an affine scheme $U_1$ \'{e}tale over $U$, and a point~$s_1$ over $s$, such that $k(s_1)$ is $k(s)$-isomorphic to $L$. Since $s$ is a maximal point of~$T_U$, one has
$$
\Spec L \simeq \varprojlim_W T_W,
$$
where\pageoriginaled $W$ ranges over the affine open neighborhoods of $s_1$. One deduces that $(\Spec L)' \simeq \varprojlim_W(T_W)'$, and by passage to the limit (\SGA 4~VII~5.8):
$$
H^d((\Spec L)', \sisi{{G_r}{|(\Spec L)'}}{{G_r}_{|(\Spec L)'}}) \simeq \varinjlim_W H^d((T_W)', \sisi{{G_r}{|(T_W)'}}{{G_r}_{|(T_W)'}});
$$
consequently, one can find $W$ such that
$$
H^d((T_W)', \sisi{{G_r}{|(T_W)'}}{{G_r}_{|(T_W)'}})\neq 0,
$$
which completes the proof of the theorem.
\end{enumerate}

\begin{corollaire} \label{XIV.3.6}
The hypotheses concerning $S, S', f, f', m$ are those of \Ref{XIV.3.5}. Let $F$ now denote a complex of sheaves of $\ZZ/m\ZZ$-modules on $X'$, bounded below and with constructible cohomology, and whose cohomology sheaves satisfy conditions \textup{(i)} and \textup{(ii)} of \Ref{XIV.3.5}. Let $n$ be an integer, and suppose that, for every affine scheme $X_1$ \'{e}tale over $X$, one has
$$
H^i(X'_1, F)=0\text{ for } i>n.
$$
Then, if $\bar{x}'$ is a geometric point lying over a point $x'$ of $X'$ such that, for an integer $j$, $(\underline{H}^j(F))_{\bar{x}'} \neq 0$, one has
$$
\delta(x')\leq n-j.
$$
\end{corollaire}

Let $T'$ be the set of points of $X'$ where the conclusion of \Ref{XIV.3.7} fails, and suppose $T' \neq \emptyset$; let $T=f(T')$, let $x$ be a maximal point of $T$, and let $x'$ be a point of $X'$ lying over~$x$. Let $j$ be the largest\pageoriginale integer such that $(\underline{H}^j(F))_{\bar{x}'} \neq 0$; thus $r=\delta(x)>n-j$. Let $Z'_q$ be the set of points where the geometric fiber of $\underline{H}^q(F)$ is $\neq 0$, and let $Z_q =h(Z'_q)$; as in the proof of \Ref{XIV.3.5}, one sees that $Z_q$ is constructible. Clearly $Z'_q=\emptyset$ for $q>n$ and for $q$ sufficiently small. The other values of $q$ fall into three subsets. Let
$$
Q_1=\{q\mid x\in Z_q\text{ and a generization of }x, \text{ distinct from }x, \not\in Z_q\}.
$$
One has $j \in Q_1$, and one can find an affine open neighborhood $U_1$ of $x$ such that, for every $q \in Q_1$, $U_1 \cap Z_q$ is an irreducible closed subset with generic point $x$. If $q \in Q_1$, one has
\begin{equation*} \label{eq:XIV.3.6.*} \tag{$*$} \delta(\sisi{\underline{H}^q(F){|U_1'}}{\underline{H}^q(F)_{|U_1'}})=\delta(x)\quad\text{(for the definition of $\delta(\underline{H}^q(F))$, \cf \Ref{XIV.3.1})}.
\end{equation*}
Let
$$
Q_2=\{q\mid\text{no generization of $x$ belongs to }Z_q\}.
$$
Then, if $j<q\leq n$, one has $q\in Q_2$, and one can find an affine open neighborhood $U_2$ of $x$ such that, for every $q \in Q_2$, $Z_q \cap U_2 = \emptyset$; thus one has
\begin{equation*} \label{eq:XIV.3.6.**} \tag{$**$} \sisi{\underline{H}^q(F){|U'_2}}{\underline{H}^q(F)_{|U'_2}}=0\text{ for }q\in Q_2.
\end{equation*}
Finally, let
$$
Q_3= \{q\mid Z_q \text{ contains strict generizations of }x\}.
$$
Then one can find an affine open neighborhood of $x$, $U_3$, such that, for every $q \in Q_3$, all the maximal points of $Z_q \cap U_3$ are generizations of $x$. If $q \in Q_3$, one has
\begin{equation*} \label{eq:XIV.3.6.***} \tag{${*}{*}{*}$} \delta(\sisi{\underline{H}^q(F){|U'_3}}{\underline{H}^q(F)_{|U'_3}})\leq n-q.
\end{equation*}

For\pageoriginale every affine scheme $X_1$ \'{e}tale over $U_1 \cap U_2 \cap \sisi{{}}{U}_3$, consider the \sisi{}{h}ypercohomology spectral sequence
$$
E^{pq}_2= H^p(X_1', \underline{H}^q(F)) \To H^*(X'_1, F).
$$
One has $E^{pq}_2=0$ for $q\sisi{>}{\in Q_2}$ by \eqref{eq:XIV.3.6.**}. One has $E^{pq}_2=0$ for $p+q \geq r+j$ except possibly for $p=r, q=j$. Indeed, this is clear if $q\in Q_2$; if $q\in Q_1$, then $p>r$ except when $p=r, q=j$, and this follows from \Ref{XIV.3.1} in view of (\Ref{eq:XIV.3.6.*}); finally, if $q\in Q_3$, since $r>n-j$, one has $p>n-q$, and the assertion follows from \Ref{XIV.3.1} in view of (\Ref{eq:XIV.3.6.***}). Since $H^{r+j}(X_1', F)=0$, it follows from the spectral sequence that
$$H^r(X_1', \underline{H}^j(F))=0;$$
but, by \Ref{XIV.3.5}, this implies $\delta(x)<r$, which is absurd.

\begin{corollaire} \label{XIV.3.7}
Let $S$ be an excellent strictly local scheme of characteristic zero, $f:X \to S$ a morphism locally of finite type, $m$ a power of a prime number, $F$ a complex of sheaves of $\ZZ/m\ZZ$-modules on $X$, bounded\sisi{}{} below, with constructible cohomology, and $n$ an integer. Then the following conditions are equivalent:
\begin{enumeratei}
\item
For every affine scheme $X_1$ \'{e}tale over $X$, one has
$$
H^i(X_1, F)=0\text{ for }i>n.
$$

\item
For\pageoriginale every geometric point $\bar{x}$ lying over the point $x$ of $X$, and for every integer~$j$ such that $(\underline{H}^j(F))_{\bar{x}} \neq 0$, one has
$$
\delta(x) \leq n-j.
$$
\end{enumeratei}
\end{corollaire}

\textup{(i)} \ALORS \textup{(ii)} is the special case of \Ref{XIV.3.6} obtained by taking $S=S'$.

\textup{(ii)} \ALORS \textup{(i)} follows immediately from \Ref{XIV.3.1}, using the hypercohomology spectral sequence
$$
H^p(X_1, \underline{H}^q(F)) \To H^*(X_1, F).
$$

\section{Main theorem and variants} \label{XIV.4}

\setcounter{subsection}{-1}

\subsection{\texorpdfstring{}{XIV.4.0}} \label{XIV.4.0}
Let $g\colon X\to S$ be a separated morphism of finite type, let $T$ be a closed subset of $S$, let \hbox{$Z=g^{-1}(T)$}, and let $F$ be a complex of abelian sheaves on $X$, bounded below. We call the group
$$
\H_{Z!}^{i}(X/S, F)=\H_{T}^{i}(S, \R_{!}g(F))
$$
the \emph{$i$-th cohomology group of $F$, with proper support and support in~$Z$}, where $\R_{!}g$ denotes the \og direct image with proper support\fg (\SGA 4~XVII). In the particular case in which $g$ is proper, one simply has
$$
\H_{Z!}^{i}(X/S, F)=\H_{Z}^{i}(X, F).
$$

\begin{proposition} \label{XIV.4.1} Let $f\colon U\to S$ be a morphism of finite type and let $F$ be a complex of abelian sheaves on $U$, bounded below. Suppose\pageoriginale that a factorization of $f$ is given:
$$
\xymatrix@C=15pt{ U\ar[rr]^-{i}\ar[rd]_-{f} & & X\ar[ld]^-{g}\\
& S & }$$
where $i$ is an open immersion and $g$ is a separated morphism of finite type, and let $G$ be a complex of abelian sheaves on $X$, bounded below, that extends $F$. Let $Y$ be a closed subscheme of $X$ whose underlying space is $X-U$, so that one has a commutative diagram:
$$
\xymatrix@C=15pt{ Y\ar[rr]^-{j}\ar[rd]_-{h} & & X\ar[ld]^-{g}\\
& S & }$$
Finally, let $n$ be an integer and let $T$ be a closed subset of $S$. Then the following conditions are equivalent:
\begin{enumeratei}
\item
One has $\prof_{T}(\R_{!}f(F))\geq n$.
\item
The canonical morphism
$$
\underline{H}_{T}^{i}(\R_{!}g(G))\to\underline{H}_{T}^{i}(\R_{!}h(j^{\ast}G))
$$
is bijective for $i<n-1$ and injective for $i=n-1$.
\item
For every scheme $S^{\prime}$ \`etale over $S$, if $X^{\prime}$ (\resp $f^{\prime}$, \resp etc.) denotes the scheme $X\times_{S}S^{\prime}$ (\resp the morphism $f_{(S^{\prime})}$, \resp etc.), the canonical morphism\pageoriginale
$$
\H_{g^{\prime -1}(T^{\prime})!}^{i}(X^{\prime}/S^{\prime}, G^{\prime})\to\H_{h^{\prime -1}(T^{\prime})!}^{i}(Y^{\prime}/S^{\prime}, j^{\prime\ast}G^{\prime})
$$
is bijective for $i<n-1$ and injective for $i=n-1$.
\end{enumeratei}
\end{proposition}

In the derived category $\D^{+}(X)$ (\cf \cite{XIV.3}), consider the distinguished triangle
$$
\xymatrix@C=15pt{ & j_{\ast}j^{\ast}G\ar[ld] & \\
i_{!}F\ar[rr] & & G.\ar[lu]}$$
Applying the functor $\R_{!}g$ to this triangle gives the triangle
\begin{equation*} \label{eq:XIV.4.1.*} \tag{$*$}
\begin{array}{c}
\xymatrix@C=0pt{
& \R_{!}h(j^{\ast}G)\ar[ld] & \\
\R_{!}f(F)\ar[rr] & & \R_{!}g(G).\ar[lu]
}
\end{array}
\end{equation*}

Let us show that \textup{(i)}\SSI\textup{(ii)}. Indeed, by Definition \Ref{XIV.1.2}, \textup{(i)} is equivalent to the relation
$$
\underline{H}_{T}^{i}(\R_{!}f(F))=0\quad\textrm{ for }i<n\textrm{;}$$
while \eqref{eq:XIV.4.1.*} gives the exact sequence of sheaves
$$
\to\underline{H}_{T}^{i}(\R_{!}f(F))\to\underline{H}_{T}^{i}(\R_{!}g(G))\to\underline{H}_{T}^{i}(\R_{!}h(j^{\ast}G))\to,
$$
which proves\pageoriginale the equivalence of \textup{(i)} and \textup{(ii)}.

\textup{(i)}\SSI\textup{(iii)}. Indeed, \textup{(i)} is equivalent to saying that, for every scheme $S^{\prime}$ \`etale over $S$, one has
\begin{equation*} \label{eq:XIV.4.1.**} \tag{$**$} {\H_{T^{\prime}}^{i}(S^{\prime}, \R_{!}f^{\prime}(F^{\prime}))=0\quad\textrm{ for }i<n.}
\end{equation*}
But \eqref{eq:XIV.4.1.*} gives the exact sequence of abelian groups
$$
\to\H_{T^{\prime}}^{i}(S^{\prime}, \R_{!}f^{\prime}(F^{\prime}))\to\H_{T^{\prime}}^{i}(S^{\prime}, \R_{!}g^{\prime}(G^{\prime}))\to \H_{T^{\prime}}^{i}(S^{\prime}, \R_{!}h^{\prime}(j^{\prime\ast}G^{\prime}))\to\;;$$
in view of \Ref{XIV.4.0}, this exact sequence takes the form
$$
\to\H_{T^{\prime}}^{i}(S^{\prime}, \R_{!}f^{\prime}(F^{\prime}))\to\H_{g^{\prime -1}(T^{\prime})!}^{i}(X^{\prime}/S^{\prime}, G^{\prime})\to\H_{h^{\prime -1}(T^{\prime})!}^{i}(Y^{\prime}/S^{\prime}, j^{\prime\ast}G^{\prime})\sisi{}{\to.}$$
The equivalence of \textup{(i)} and \textup{(iii)} follows, taking into account the form \eqref{eq:XIV.4.1.**} of \textup{(i)}.

\setcounter{subsubsection}{-1} \setcounter{subsection}{2}

\subsubsection{} \label{XIV.4.2.0}
When $f\colon U\to S$ is affine, we shall give \emph{local} conditions on $F$ under which conditions \textup{(i)} through \textup{(iii)} of \Ref{XIV.4.1} hold. In what follows, all schemes considered are excellent schemes of characteristic zero, and all sheaves are sheaves of $\ZZ/m\ZZ$-modules, where $m$ is a power of a prime number. If resolution of singularities in the sense of (\SGA 4~XIX) were available, the stated results and their proofs would remain valid for excellent schemes of equal characteristic, with $m$ prime to the characteristic.

\setcounter{subsection}{1}

\begin{theoreme} \label{XIV.4.2}
Let\pageoriginale $S$ be an excellent scheme of characteristic zero and let \hbox{$f\colon U\!\to\! S$} be a separated morphism of finite type. Let $F$ be a bounded-below complex of sheaves of $\ZZ/m\ZZ$-modules on $U$ with constructible cohomology, let $n$ be an integer, and let~$T$ be a closed subset of $S$. Then the following conditions are equivalent:
\begin{enumeratei}
\item
For every scheme $U_{1}$ \`etale over $U$ and affine over $S$, denoting by~$f_{1}$ the structural morphism of $U_{1}$ and by $F_{1}$ the restriction of $F\sisi{_{1}}{}$ to $U_{1}$, one has
$$
\prof_{T}(\R_{!}f_{1}(F_{1}))\geq n$$
(\cf Proposition\ptbl \Ref{XIV.4.1} for the meaning of this relation).
\item
For every point $u$ of $U$, one has
$$
\prof_{u}(F)\geq n-\delta_{T}(u),
$$
where, writing $s=f(u)$, one sets (\cf \Ref{XIV.2.2}): $\delta_{T}(u)=\degtr(k(u)/k(s))+\codim\big(\overline{\{ s\} }\cap T, \overline{\{ s\} } \big)$.
\end{enumeratei}
\end{theoreme}

\begin{proof}
$1^\circ)$ Let $t$ be a point of $T$, let $\overline{S}$ be the strict localization of $S$ at a geometric point above $t$, and let $S^{\prime}$ be the completion of $\overline{S}$, with closed point $t^{\prime}$; then $\overline{S}$ is excellent by (\EGA IV~7.9.5), so $S^{\prime}$ is an excellent complete strictly local scheme. Given a scheme $U$ over $S$ (respectively an $S$-morphism $f$, respectively etc.), we shall denote by $U^{\prime}$ (respectively $f^{\prime}$, respectively etc.) the scheme $U\times_{S}S^{\prime}$ (respectively the morphism $f_{(S^{\prime})}$, respectively etc.). We have the Cartesian square
$$
\xymatrix{ U^{\prime}\ar[r]^-{h}\ar[d]_-{f^{\prime}} & U\ar[d]^-{f}\\
S^{\prime}\ar[r]^-{g} & S, }
$$
in\pageoriginale which the morphism $g$ is regular (\EGA IV~7.8.2). Let us show that it is enough to prove that (for every point $t\in T$) the following two properties are equivalent:
\begin{itemize}
\item[$\textup{(i)}_{t}$] For every scheme $U_{1}$ étale over $U$ and affine over $S$, setting $f_{1}\colon U_{1}\to S$, one has
$$
\prof_{t^{\prime}}(\R_{!}f_{1}^{\prime}(F_{1}^{\prime}))\geq n.
$$

\item[$\textup{(ii)}_{t}$] For every point $u^{\prime}$ of $U^{\prime}$, one has
$$
\prof_{u^{\prime}}(F^{\prime})\geq n-\delta_{t^{\prime}}(u^{\prime}).
$$
\end{itemize}
It is enough to prove the following lemma:

\begin{sublemme} \label{XIV.4.2.1}
One has \textup{(i)}\SSI$\textup{(i)}_{t}$ for every $t\in T$ and \textup{(ii)}\SSI$\textup{(ii)}_{t}$ for every $t\in T$.
\end{sublemme}

\textup{(i)}\SSI$\textup{(i)}_{t}$ for every $t\in T$. Indeed, \textup{(i)} is equivalent to saying that, for every scheme $U_{1}$ étale over $U$ and affine over $S$, one has
$$
\prof_{T}(\R_{!}f_{1}(F_{1}))\geq n\;;$$
now, by \Ref{XIV.1.8},
$$
\prof_{T}(\R_{!}f_{1}(F_{1}))=\inf_{t\in T}\prof_{t}(\R_{!}f_{1}(F_{1})).
$$
Since $g^{\ast}(\R_{!}f_{1}(F_{1}))\simeq\R_{!}f_{1}^{\prime}(F_{1}^{\prime})$ (\SGA 4~XVII), one has, by \Ref{XIV.1.16},
$$
\prof_{t}(\R_{!}f_{1}(F_{1}))=\prof_{t^{\prime}}(\R_{!}f_{1}^{\prime}(F_{1}^{\prime})),
$$
and hence\pageoriginale \textup{(i)} is equivalent to saying that, for every $t\in T$, one has $\prof_{t^{\prime}}(\R_{!}f_{1}^{\prime}(F_{1}^{\prime}))\geq n$, which is precisely $\textup{(i)}_{t}$.

\textup{(ii)}$_{t}$ for every $t\in T$\ALORS\textup{(ii)}. Indeed, let $u\in U$; we must show the relation
$$
\prof_{u}(F)\geq n-\delta_{T}(u),
$$
where $\delta_{T}(u)=\inf_{t\in T\cap\overline{\{ s\} }}\delta_{t}(u)$ (\cf \Ref{XIV.2.2}); we are therefore reduced to showing that, for every $t\in T\cap\overline{\{ s\} }$, one has
$$
\prof_{u}(F)\geq n-\delta_{t}(u).
$$
Let $u^{\prime}$ be a point of $U^{\prime}$ such that $h(u^{\prime})=u$ and $\delta_{t^{\prime}}(u^{\prime})=\delta_{t}(u)$ (\cf \Ref{XIV.2.1.1}). Since $h$ is locally acyclic (\SGA 4~XIX~4.1), it follows from \Ref{XIV.1.16} and the fact that~$u^{\prime}$ is a generic point of $U_{u}^{\prime}$ that
$$
\prof_{u^{\prime}}(F^{\prime})=\prof_{u}(F).
$$
But by $\textup{(ii)}_{t}$ one has $\prof_{u}(F)=\prof_{u^{\prime}}(F^{\prime})\geq n-\delta_{t}(u)$, which proves \textup{(ii)}.

\textup{(ii)}\ALORS$\textup{(ii)}_{t}$ for every $t$. With the notation of \Ref{XIV.2.2.1}, for every point $u^{\prime}$ of $U^{\prime}$ one has, by \Ref{XIV.1.16},
$$
\prof_{u^{\prime}}(F^{\prime})\geq\prof_{u}(F)+2\dim h_{u^{\prime}}^{-1}(u)\geq\prof_{u}(F)+\dim h_{u^{\prime}}^{-1}(u).
$$
Taking into account \Ref{XIV.2.2.1} and \textup{(ii)}, we obtain
$$
\prof_{u^{\prime}}(F^{\prime})\geq n-\delta_{T}(u)+\dim h_{u^{\prime}}^{-1}(u)\geq n-\delta_{t^{\prime}}(u^{\prime}),
$$
which\pageoriginaled is precisely $\textup{(ii)}_{t}$.

$2^\circ)$ $\textup{(ii)}_{t}$\SSI$\textup{(i)}_{t}$. We reduce immediately to the case where $F$ is bounded, by truncating $F$ at a sufficiently high degree. We can realize $S^{\prime}$ as a closed subscheme of a complete regular local scheme, hence an excellent one; it then follows from (\SGA 5~I~3.4.3) that there exists a dualizing complex $K$ on $S^{\prime}$ and that $\R^{!}f^{\prime}(K)=K^{\prime}$ is a dualizing complex on $U^{\prime}$. We choose $K$ so that $\delta^{K}(t^{\prime})=0$ (for the definition of $\delta^{K}(t^{\prime})$, \cf \Ref{XIV.2.3}), and denote by $\D F^{\prime}$ the dual of $F^{\prime}$ with respect to $K^{\prime}$. We can reformulate the hypothesis $\textup{(ii)}_{t}$ as follows:

\begin{sublemme} \label{XIV.4.2.2}
Let $u^{\prime}$ be a point of $U^{\prime}$; then the following conditions are equivalent:
\begin{enumeratei}
\item
One has $\prof_{u^{\prime}}(F^{\prime})\geq n-\delta_{t^{\prime}}(u^{\prime})$.
\item
One has $\big(\underline{H}^{q}(\D F^{\prime}) \big)_{\overline{u}^{\prime}}=0$ for $q>-n-\delta_{t^{\prime}}(u^{\prime})$ ($\overline{u}^{\prime}$ a geometric point above $u^{\prime}$).
\end{enumeratei}
\end{sublemme}

Let $\overline{U}^{\prime}$ be the strict localization of $U^{\prime}$ at $\overline{u}^{\prime}$, and let $\overline{F}^{\prime}$ be the inverse image of $F^{\prime}$ under the morphism $\overline{U}^{\prime}\to U^{\prime}$. The relation $\prof_{u^{\prime}}(F^{\prime})\geq n-\delta_{t^{\prime}}(u^{\prime})$ is, by definition, equivalent to the following:
\begin{equation*} \label{eq:XIV.4.2.*} \tag{$*$} {\H_{\overline{u}^{\prime}}^{i}(\overline{F}^{\prime})=0\quad\textrm{ for }i<n-\delta_{t^{\prime}}(u^{\prime}).}
\end{equation*}
Let $\D\big(\H_{\overline{u}^{\prime}}^{i}(\overline{F}^{\prime})\big)$ be the dual of the abelian group $\H_{\overline{u}^{\prime}}^{i}(\overline{F}^{\prime})$ with respect to $\ZZ/m\ZZ$. By \Ref{XIV.2.3.2}, $K^{\prime}[-2\delta_{t^{\prime}}(u^{\prime})]=K^{\prime\prime}$ satisfies $\delta^{K^{\prime\prime}}(u^{\prime})=0$; since $F^{\prime}$ has constructible cohomology, one has $\overline{\D F^{\prime}}=\D(\overline{F}^{\prime})$, and the local\pageoriginale duality theorem (\SGA 5~I~4.5.3) then shows that
$$
\D\big(\H_{\overline{u}^{\prime}}^{i}(\overline{F}^{\prime})\big) \simeq\big(\H^{-i-2\delta_{t^{\prime}}(u^{\prime})}(\D F^{\prime})\big)_{\overline{u}^{\prime}}.
$$
Thus \eqref{eq:XIV.4.2.*} is equivalent to the relation
\begin{equation*} \label{eq:XIV.4.2.**} \tag{$**$}
\big(\underline{H}^{q}(\D F^{\prime})\big)_{\overline{u}^{\prime}}=0\quad\text{for }q>-n-\delta_{t^{\prime}}(u^{\prime}).
\end{equation*}

We are now in a position to prove the theorem. The relation \textup{(ii)}$_{t}$ is equivalent to relation \eqref{eq:XIV.4.2.**}. Let $G^{q}=\underline{H}^{q}(\D F^{\prime})$; the affine Lefschetz theorem (\Ref{XIV.3.1}) implies in particular that, for every scheme $U_{1}$ étale over $U$ and affine over $S$, one has
$$
\H^{p}(U_{1}^{\prime}, G^{q})=0\quad\textrm{ for }p>\delta(G^{q}),
$$
where $\delta(G^{q})$ is the upper bound of the $\delta_{t^{\prime}}(u^{\prime})$ for those $u^{\prime}$ such that $G_{\overline{u}^{\prime}}^{q}\neq 0$; by \eqref{eq:XIV.4.2.**}, one has $\delta(G^{q})\leq -n-q$, so $\textup{(ii)}_{t}$ implies the relation
$$
\H^{p}(U_{1}^{\prime}, \underline{H}^{q}(\D F^{\prime}))=0\quad\textrm{ for }p>-q-n.
$$
Taking into account the hypercohomology spectral sequence for the functor \og sections on $U_{1}^{\prime}$\fg with respect to the complex $\D F^{\prime}$:
$$
\E_{2}^{pq}=\H^{p}(U_{1}^{\prime}, \underline{H}^{q}(\D F^{\prime}))\To\H^{\ast}(U_{1}^{\prime}, \D F^{\prime}),
$$
we obtain the relation
\begin{equation*} \label{eq:XIV.4.2.***} \tag{${*}{*}{*}$}
\H^{i}(U_{1}^{\prime}, \D F^{\prime})=0\quad\textrm{ for }i>-n.
\end{equation*}

Conversely,\pageoriginale suppose that the preceding relation holds for every $U_{1}$ étale over $U$ and affine over $S$. Apply Proposition \Ref{XIV.3.6}, replacing $S$ there by $\overline{S}^{\prime}$; the hypotheses of \Ref{XIV.3.6} concerning $\sisi{S}{\overline{S}}$ are satisfied, since, for every affine scheme $\overline{U}_{1}$ étale over $\overline{U}$, one can find a scheme above $\overline{U}_{1}$ which comes by inverse image from a scheme étale over $U$ and affine over $S$; as for the hypotheses concerning $F$, they are satisfied by \Ref{XIV.2.3.3}. Thus, for every point $u^{\prime}$ of $U^{\prime}$ such that $\big(\underline{H}^{q}(\D F^{\prime})\big)_{\overline{u}^{\prime}}\neq 0$, one has:
$$
\delta_{t^{\prime}}(u^{\prime})\leq -n-q,
$$
which is precisely relation \eqref{eq:XIV.4.2.**}; we have therefore proved the equivalence
$$
\text{\textup{(ii)}$_{t}$\SSI\eqref{eq:XIV.4.2.***}}.
$$
We shall transform relation \eqref{eq:XIV.4.2.***}; first, one has
$$
\H^{i}(U_{1}^{\prime}, \D F^{\prime})=\big(\underline{H}^{i}(\R f_{1\ast}^{\prime}(\D F_{1}^{\prime}))\big)_{t^{\prime}}\;;$$
but by (\SGA 5~I~1.12), there is a canonical isomorphism
$$
\R f_{1\ast}^{\prime}(\D F_{1}^{\prime})\simeq\D\big(\R_{!}f_{1}^{\prime}(F_{1}^{\prime})\big),
$$
where $\D\big(\R_{!}f_{1}^{\prime}(F_{1}^{\prime})\big)$ denotes the dual of $\R_{!}f_{1}^{\prime}(F_{1}^{\prime})$ with respect to $K$. We thus see that $\textup{(ii)}_{t}$ is equivalent to
$$
\big(\underline{H}^{i}(\D(\R_{!}f_{1}^{\prime}(F_{1}^{\prime})))\big)_{t^{\prime}}=0\quad\textrm{ for }i>-n.
$$
Applying\pageoriginale the local duality theorem (\SGA 5~I~4.5.3) once again, but this time at the point $t^{\prime}$, we find that
$$
\big(\underline{H}^{i}(\D(\R_{!}f_{1}^{\prime}(F_{1}^{\prime})))\big)_{t^{\prime}}\simeq\D\big(\H_{t^{\prime}}^{-i}(\R_{!}f_{1}^{\prime}(F_{1}^{\prime}))\big),
$$
and finally $\textup{(ii)}_{t}$ is equivalent to the relation
$$
\H_{t^{\prime}}^{i}(\R_{!}f_{1}^{\prime}(F_{1}^{\prime}))=0\quad\textrm{ for }i<n,
$$
that is, $\prof_{t^{\prime}}(\R_{!}f_{1}^{\prime}(F_{1}^{\prime}))\geq n$, which completes the proof of the theorem.
\skipqed
\end{proof}

\begin{subremarque} \label{XIV.4.2.3}
The argument simplifies quite considerably when one assumes that $S$ admits (at least locally) a dualizing complex (for example, that it is locally immersible in a regular scheme). This avoids recourse to a completion (the passage to the case where $S$ is strictly local being immediate), to \Ref{XIV.2.3.3}, and to the rather unappealing technical statement \Ref{XIV.3.6}, which may then be replaced by the more congenial reference \Ref{XIV.3.7}.
\end{subremarque}

\begin{corollaire} \label{XIV.4.3}
Let $S$ be an excellent scheme of characteristic zero and $f\colon U\to S$ a separated morphism of finite type such that $U$ is the \emph{union of $c+1$ open subsets affine over~$S$}. Let $F$ be a complex of sheaves of $\ZZ/m\ZZ$-modules, bounded below and with constructible cohomology, let $n$ be an integer, and let $T$ be a closed subset of $S$. Suppose that, for every point $u\in U$, one has
$$
\prof_{u}(F)\geq n-\delta_{T}(u).
$$
Then one has\pageoriginale
$$
\prof_{T}(\R_{!}f(F))\geq n-c.
$$
\end{corollaire}

Indeed, let $U_{\sisi{i}{j}}$, $0\leq \sisi{i}{j}\leq c$, be a covering of $U$ by open subsets $U_{\sisi{i}{j}}$ affine over~$S$. Adopting again the notation of the proof of \Ref{XIV.4.2}, one has, for every $\sisi{i}{j}$,
$$
\H^{i}(U_{\sisi{i}{j}}^{\prime}, \underline{H}^{q}(\D F^{\prime}))=0\quad\textrm{ for }i>-n-q.
$$
Using the spectral sequence that relates the cohomology of $U$ to that of the covering formed by the $U_{\sisi{i}{j}}$ (\SGA 4~V~2.4), the preceding relation shows that
$$
\H^{i}(U^{\prime}, \underline{H}^{q}(\D F^{\prime}))=0\quad\textrm{ for }i>-n-q+c.
$$
The corollary then follows from the end of the proof of \Ref{XIV.4.2}.

\begin{corollaire} \label{XIV.4.4}
Let $S$ be an excellent scheme of characteristic zero, $g\colon X\to S$ a morphism, $U$ an open subset of $X$ that is the union of $c+1$ open subsets affine over $S$, $Y$ a closed subscheme with underlying space $X-U$, and $j\colon Y\to X$ the natural morphism. Let~$F$ be a complex of sheaves of $\ZZ/m\ZZ$-modules on $X$, bounded below and with constructible cohomology, let $T$ be a closed subset of $S$, and let $n$ be an integer. Suppose that, for every point $u$ of $U$, one has
$$
\prof_{u}(F)\geq n-\delta_{T}(u).
$$
Then\pageoriginale the canonical morphism
$$
\H_{g^{-1}(T)!}^{i}(X/S, F)\to\H_{(g^{-1}(T)\cap Y)!}^{i}(Y/S, j^{\ast}F)$$
is bijective for $i<n-c-1$, and injective for $i=n-c-1$.
\end{corollaire}

This follows immediately from \Ref{XIV.4.1} and \Ref{XIV.4.3}.

\begin{corollaire}[Local Lefschetz theorem] \label{XIV.4.5}
Let $S$ be an excellent henselian local scheme of characteristic zero, $t$ the closed point of $S$, $X$ a scheme proper over $S^{\prime}=S-\{ t\}$, and $U$ an open subset of $X$ that is the union of $c+1$ affine open subsets. Let $Y$ be a closed subscheme of $X$ with underlying space $X-U$, let $j\colon Y\to X$ be the canonical morphism, let $F$ be a complex of sheaves of $\ZZ/m\ZZ$-modules on $X$, bounded below and with constructible cohomology, and let $n$ be an integer. Suppose that, for every point $u$ of~$U$, one has
$$
\prof_{u}(F)\geq n-\delta_{t}^{\prime}(u), \quad\text{where } \delta_{t}^{\prime}(u)=\delta_{t}(u)-1.
$$
Then the canonical morphism
$$
\H^{i}(X, F)\to\H^{i}(Y, j^{\ast}F)$$
is bijective for $i<n-c-1$, and injective for $i=n-c-1$.
\end{corollaire}

Let $f\colon U\to S$ be the canonical morphism; it follows from
\Ref{XIV.4.2}, applied with $n$ replaced by $n+1$, that
$$
\prof_{t}(\R_{!}f(F\sisi{|U}{_{|U}}))\geq n+1-c.
$$
The preceding\pageoriginale relation shows that the canonical morphism
$$
\H^{i}(S, \R_{!}f(F\sisi{|U}{_{|U}}))\to\H^{i}(S^{\prime}, \R_{!}f(F\sisi{|U}{_{|U}}))$$
is bijective for $i<n-c$, and injective for $i=n-c$. Since $\R_{!}f(F\sisi{|U}{_{|U}})$ is zero outside $S^{\prime}$, one has $\H^{i}(S, \R_{!}f(F\sisi{|U}{_{|U}}))\simeq\big(\underline{H}^{i}(\R_{!}f(F\sisi{|U}{_{|U}}))\big)_{t}=0$, and consequently
\begin{equation*} \label{eq:XIV.4.5.*} \tag{$*$}
\H^{i}(S^{\prime}, \R_{!}f(F\sisi{|U}{_{|U}}))=0\quad\text{for }i<n-c.
\end{equation*}
Let $g\colon X\to S^{\prime}$, $h\colon Y\to S^{\prime}$, $f^{\prime}\colon U\to S^{\prime}$ be the canonical morphisms. It follows from the distinguished triangle
$$
\xymatrix@C=0pt{ & \R h_{\ast}(j^{\ast}F)\ar[ld] & \\
\R_{!}f^{\prime}(F\sisi{|U}{_{|U}})\ar[rr] & & \R g_{\ast}(F)\ar[lu] }$$
that condition \eqref{eq:XIV.4.5.*} is equivalent to the assertion that the morphism
$$
\H^{i}(S^{\prime}, \R g_{\ast}(F))\to\H^{i}(S^{\prime}, \R h_{\ast}(j^{\ast}F))$$
is bijective for $i<n-c-1$, and injective for $i=n-c-1$. Since this morphism identifies canonically with the morphism
$$
\H^{i}(X, F)\to\H^{i}(Y, j^{\ast}F),
$$
the conclusion follows at once.

\begin{corollaire}[Global Lefschetz theorem] \label{XIV.4.6}
Let\pageoriginale $S$ be the spectrum of a field, $X$ a scheme proper over $S$, and $U$ an open subset of $X$ that is the union of $c+1$ affine open subsets. Let~$Y$ be a closed subscheme of $X$ with underlying space $X-U$, let $j\colon Y\to X$ be the canonical morphism, let $F$ be a complex of sheaves of $\ZZ/m\ZZ$-modules on $X$, bounded below and with constructible cohomology, and let $n$ be an integer. Suppose that, for every point $u$ of $U$, one has
$$
\prof_{u}(F)\geq n-\dim\big(\overline{\{ u\} }\big).
$$
Then the canonical morphism
$$
\H^{i}(X, F)\to\H^{i}(Y, j^{\ast}F)$$
is bijective for $i<n-c-1$, and injective for $i=n-c-1$.

More generally, if $g\colon X\to S$ is a separated morphism of finite type, with the hypotheses on $S$, $U$, $Y$, and $F$ being the same as above, then the canonical morphism
$$
\H_{!}^{i}(X/S, F)\to\H_{!}^{i}(Y/S, j^{\ast}F)$$
(where $\H_{!}^{i}$ denotes cohomology with proper support, that is, $\H_{!}^{i}(X\sisi{}{/S}, F)=\H^{i}(S, \R_{!}g(F))$) is bijective for $i<n-c-1$, and injective for $i=n-c-1$.
\end{corollaire}

The corollary is a special case of \Ref{XIV.4.4}, with $T=S$.

Here is a partial converse to \Ref{XIV.4.3}:

\begin{proposition} \label{XIV.4.7}
Let\pageoriginale $S$ be a Noetherian scheme and let $f\colon U\to S$ be a morphism of finite type. Suppose that there exists a dualizing complex $K$ on $S$ and that $\R^{!}f(K)$ is a dualizing complex on $U$. Let $T$ be a closed subset of $S$ and let $c$ be an integer. Then the following conditions are equivalent:
\begin{enumeratei}
\item
For every bounded-below complex $F$ of sheaves of $\ZZ/m\ZZ$-modules on $U$ with constructible cohomology, and every integer $n$ such that, for every point $u$ of $U$,
$$
\prof_{u}(F)\geq n-\delta_{T}(u),
$$
one has
$$
\prof_{T}(\R_{!}f(F))\geq n-c.
$$

\item
For every constructible sheaf $G$ of $\ZZ/m\ZZ$-modules on $U$, and every point $t\in T$, one has
$$
\big(\R^{p}f_{\ast}(G)\big)_{\overline{t}}=0\quad\textrm{ for }p>\delta(G, f, t)+c$$
(recall from (\textup{\SGA \sisi{X~}{}4~XIX~6.0}) that $\delta(G, f, t)=\sup\{ \delta_{t}(u)|t\in\overline{\{ u\} }\textrm{ and }G_{\overline{u}}\neq 0\}$).
\end{enumeratei}
\end{proposition}

N.B. Condition \textup{(ii)} is satisfied by \Ref{XIV.3.1} if $f$ is separated and, locally on $S$ for the \`etale topology, $U$ is a union of $c+1$ open subsets affine over $S$; thus \Ref{XIV.4.7} contains \Ref{XIV.4.3}\sfootnote{At least in the \sisi{case where}{case where} $S$ locally admits a dualizing complex, \sisi{e.g.}{for example} when $S$ is locally embeddable in a regular scheme.}.

We may plainly assume that $S$ is local and that $T$ is the closed point $t$ of~$S$. The proof of \textup{(ii)}\ALORS\textup{(i)} is essentially identical\pageoriginale to part $2^\circ)$ of the proof of \Ref{XIV.4.2}. Let us briefly prove that \textup{(i)}\ALORS \textup{(ii)}. Set $n=\delta(G,f,t)$. The local duality theorem (\SGA 5~I~4.3.2), applied to $\D G$, shows that
$$
\D\big(\H_{\overline{u}}^{i}(\D G)\big) \simeq\big(\underline{H}^{-i-2\delta_{t}(u)}(G)\big)_{\overline{u}}.
$$
Since $G$ is concentrated in degree $0$, one therefore has $\H_{\overline{u}}^{i}(\D G)=0$ except possibly for $i=-2\delta_{t}(u)$; more precisely,
\sisi{\begin{align*}
\prof_{u}(\D G)=-2\delta_{t}(u) &\;\textrm{ if }\; G_{\overline{u}}\neq 0, \\
\prof_{u}(\D G)=\infty &\;\textrm{ if }\; G_{\overline{u}}=0.
\end{align*}}
{\[
\prof_{u}(\D G)=
\begin{cases}
-2\delta_{t}(u) &\text{if } G_{\overline{u}}\neq 0, \\
\infty &\text{if } G_{\overline{u}}=0.
\end{cases}
\]}
It follows that, for every $u\in U$,
$$
\prof_{u}(\D G)\geq -n-\delta_{t}(u).
$$
Hypothesis \textup{(i)} then gives $\prof_{t}(\R_{!}f(\D G))\geq -n-c$. We transform this relation using the isomorphism $\R_{!}f(\D G)\simeq\D(\R f_{\ast}(G))$ (\SGA 5 I~1.12) and applying the local duality theorem at $t$; this gives
$$
\big(\underline{H}^{i}(\R f_{\ast}(G)\big)_{\overline{t}}=0\quad\textrm{ for }i>n+c,
$$
which is precisely \textup{(ii)}.

\subsection{\texorpdfstring{}{XIV.4.8}} \label{XIV.4.8} Under the hypotheses of \Ref{XIV.4.4} with $g$ proper (\resp \Ref{XIV.4.5}, \resp \Ref{XIV.4.6} with $g$ proper), if $V$ is an open neighborhood of $Y$ in $X$, the\pageoriginaled morphism
$$
\H^{i}(V, F)\to\H^{i}(Y, j^{\ast}F)$$
is bijective for $i<n-c-1$ and injective for $i=n-c-1$. Indeed, if $\sisi{i}{\iota}\colon V\to X$ is the canonical morphism, it is enough to apply \Ref{XIV.4.4} (\resp \Ref{XIV.4.5}, \resp \Ref{XIV.4.6}) to the complex $\R \sisi{i}{\iota}_{\ast}(\sisi{F{|V}}{F_{|V}})$. One may ask whether the preceding morphism is bijective for $i=n-c-1$ and injective for $i=n-c$. It plainly suffices that the hypotheses hold with $n$ replaced by $n+1$; the following proposition shows that slightly less is enough.

\begin{proposition} \label{XIV.4.9}
Let $S$ be an excellent local scheme of characteristic zero with closed point $t$ (\resp assume in addition to the preceding conditions that $S$ is Henselian), let $f\colon X\to S$ be a proper morphism (\resp proper over $S-\{ \sisi{s}{t}\}$), and let $U$ be an open subset of $X$ that is the union of $c+1$ affine open subsets. Let $Y$ be a closed subscheme of $X$ whose underlying space is $X-U$, let $j\colon Y\to X$ be the canonical morphism, let $F$ be a bounded-below complex of sheaves of $\ZZ/m\ZZ$-modules on $X$ with constructible cohomology, and let $n$ be an integer. Suppose that, for every point $u$ of $U$,
$$
\prof_{u}(F)\geq\inf(n-1, n-\delta_{t}(u))\quad(\textrm{\resp }\;\prof_{u}(F)\geq\inf(n-1, n+1-\delta_{t}(u))).
$$
Then, for every open neighborhood $V$ of $Y$ in $X$, the canonical morphism\pageoriginale
$$
\H_{f^{-1}(t)}^{i}(V, F)\to\H_{f^{-1}(t)\cap Y}^{i}(Y, j^{\ast}F)\quad(\textrm{\resp }\;\H^{i}(V, F)\to\H^{i}(Y, j^{\ast}F))
$$
is bijective for $i<n-c-2$ and injective for $i=n-c-2$. Moreover, there exists an open neighborhood $V_{0}$ of $Y$ in $X$ such that, for every such $V$ with $V\subset V_{0}$, the canonical morphism
$$
\H_{f^{-1}(t)\cap V}^{i}(V, F)\to\H_{f^{-1}(t)\cap Y}^{i}(Y, j^{\ast}F)\quad(\textrm{\resp }\;\H^{i}(V, F)\to\H^{i}(Y, j^{\ast}F))
$$
is bijective for $i<n-c-1$ and injective for $i=n-c-1$.
\end{proposition}

\begin{proof}
For simplicity, set $\delta_{t}^{\prime}(u)=\delta_{t}(u)$ (\resp $\delta_{t}^{\prime}(u)=\delta_{t}(u)-1$). The first assertion of \Ref{XIV.4.9} follows from \Ref{XIV.4.8}, since the hypotheses of \Ref{XIV.4.4} (\resp \Ref{XIV.4.5}) hold there with $n$ replaced by $n-1$. They also hold for $n$ itself, except at points $u$ such that $\delta_{t}^{\prime}(u)=0$. Now, for $u\in U$, the condition $\delta_{t}^{\prime}(u)=0$ is equivalent to saying that $u$ is a \emph{closed point of $U_t$} (\resp a \emph{closed point of $X$}). Let $E$ be the set of points $u$ of $U$ such that $\delta_{t}^{\prime}(u)=0$; let us show that \emph{for all but finitely many points $u\in E$}, one has $\prof_{u}(F)\geq n$. Let $\overline{S}$ be the strict localization of $S$ at $t$, let $S^{\prime}$ be the completion of~$\overline{S}$, with closed point $t^{\prime}$, and consider the Cartesian square
$$
\xymatrix{ U^{\prime}\ar[r]^-{h}\ar[d]_-{f^{\prime}} & U\ar[d]^-{f}\\
S^{\prime}\ar[r]^-{g} & S.}$$
The depth hypotheses at the points of $U$ are preserved when\pageoriginale $U$ is replaced by $U^{\prime}$ and $F$ by its inverse image $F^{\prime}$ on $U^{\prime}$. Indeed, let $u^{\prime}\in U^{\prime}$ and $u=h(u^{\prime})$. If $u\not\in E$, one has $\prof_{u}(F)\geq n-\delta_{t}^{\prime}(u)$, and \Ref{XIV.4.2.1} shows that this implies $\prof_{u^{\prime}}(F^{\prime})\geq n-\delta_{t^{\prime}}^{\prime}(u^{\prime})$. If $u\in E$, then $u^{\prime}$ is a closed point of $U_{t}^{\prime}$ (\resp a point of $X^{\prime}=X\times_{S}S^{\prime}$ closed in $X^{\prime}$), and, since the fiber $U_{u}^{\prime}$ of $h$ at $u$ has dimension zero and $h$ is regular, \Ref{XIV.1.16} gives $\prof_{u^{\prime}}(F^{\prime})=\prof_{u}(F)\geq n-1$.

Let $K$ be a dualizing complex on $S^{\prime}$ normalized at the closed point $t^{\prime}$, and let $\D F^{\prime}$ be the dual of $F^{\prime}$ with respect to $\R^{!}f(K)$. By \Ref{XIV.4.2.2}, the \`etale-depth hypotheses at the points of $U^{\prime}$ translate into the relations
$$
\big(\underline{H}^{q}(\D F^{\prime})\big)_{\overline{u}^{\prime}}=0\quad\textrm{ for }q>-n-\delta_{t^{\prime}}^{\prime}(u^{\prime})\quad(\textrm{\resp }q>-n-2-\delta_{t^{\prime}}^{\prime}(u^{\prime})),
$$
if $u^{\prime}$ is not a point of $E^{\prime}=h^{-1}(E)$, and
$$
\big(\underline{H}^{q}(\D F^{\prime})\big)_{\overline{u}^{\prime}}=0\quad\textrm{ for }q>-(n-1)\quad(\textrm{\resp }q>-n-1), \textrm{ if }u^{\prime}\in E^{\prime}.
$$

Let $G=\underline{H}^{-(n-1)}(\D F^{\prime})$ (\resp $G=\underline{H}^{-n-1}(\D F^{\prime})$). Since $G$ is a constructible sheaf, the set of points at which its geometric fiber is nonzero is constructible (\SGA 4~IX~2.4~\textup{(iv)}). By hypothesis this set is contained in the set $E^{\prime}$ of closed points of $U_{t^{\prime}}^{\prime}$ (\resp of points of $U^{\prime}$ closed in $X^{\prime}$); it follows from \Ref{XIV.4.9.1} below that this set consists of finitely many points. Applying \Ref{XIV.4.2.2}, one sees that $\prof_{u^{\prime}}(F^{\prime})\geq n$ for all but finitely many points of $E^{\prime}$. By \Ref{XIV.1.16}, it follows that, for all but finitely many\pageoriginale points of $E^{\sisi{\prime}{}}$,
$$
\prof_{u}(F)\geq n.
$$

Let $V$ be an open neighborhood of $Y$ in $X$, contained in the complement in $X$ of the finite set of point\sisi{}{s} $u$ of $E$ for which $\prof_{u}(F)=n-1$. If $\sisi{i}{\iota}\colon V\to X$ is the canonical immersion, set
$$F_{1}=\R \sisi{i}{\iota}_{\ast}(\sisi{F{|V}}{F_{|V}})\;;$$
then $F_{1}$ is a bounded-below complex of sheaves on $X$ with constructible cohomology (\SGA 4~XIX~5.1). We shall see that, for every point $u$ of $U$, the complex $F_{1}$ satisfies
\begin{equation*} \label{eq:XIV.4.9.*} \tag{$*$}
\prof_{u}(F_{1})\geq n-\delta_{t}^{\prime}(u).
\end{equation*}
If $u\in U\cap V$, then $\prof_{u}(F_{1})=\prof_{u}(F)$, and relation \eqref{eq:XIV.4.9.*} holds by hypothesis at the points of $U$ that do not belong to $E$; at the latter points, it also holds by the choice of $V$. Finally, if $u\in U$ and $u\not\in V$, then by \Ref{XIV.1.6} g), $\prof_{u}(F_{1})=\infty$. Apply \Ref{XIV.4.4} (\resp \Ref{XIV.4.5}) with $F$ replaced by $F_{1}$; this gives the asserted result, since, for every $i$,
$$
\H_{f^{-1}(t)}^{i}(X, \R \sisi{i}{\iota}_{\ast}(\sisi{F{|V}}{F_{|V}}))\simeq\H_{f^{-1}(t)\cap V}^{i}(V, F)\quad(\textrm{\resp }\H^{i}(X, \R \sisi{i}{\iota}_{\ast}(\sisi{F{|V}}{F_{|V}}))\simeq\H^{i}(V, F).
$$
\skipqed
\end{proof}

\begin{sublemme} \label{XIV.4.9.1} A constructible set $E$ contained in the set of closed points of a Noetherian scheme $X$ consists of finitely many points\pageoriginale.
\end{sublemme}

Indeed, $E$ is a \sisi{finite union}{finite union} \sisi{of a set}{of sets} of the form $U\cap\complement V$, where $U$ and $V$ are open subsets of $X$; by hypothesis, all the points of $U\cap\complement V$ are generic points of this set, so there are only finitely many of them.

\section{Geometric depth} \label{XIV.5}

To apply \Ref{XIV.4.2} and its corollaries in practice, one needs a convenient criterion for verifying the étale-depth hypotheses at the points of $U$. We shall give such a criterion, using the local Lefschetz theorem~\sisi{(\Ref{XIV.4.5})}{\Ref{XIV.4.5}}.

\subsection{\texorpdfstring{}{XIV.5.1}} \label{XIV.5.1} \sisi{Let}{Let} $A$ be a Noetherian local ring; when we speak of the étale depth of $A$, we mean the depth at the closed point. We shall introduce a notion of \og geometric depth of $A$\fg and use \Ref{XIV.4.5} to compare it with the étale depth $\profet(A)$.

\begin{proposition} \label{XIV.5.2} Let $A$ be a Noetherian local ring; suppose that $A$ is isomorphic to a quotient of a regular local ring $B$ by an ideal $I$ (this is true, for example, when $A$ is complete, by Cohen's theorem (\textup{$\EGA 0_{\textup{IV}}$~19.8.8})). Let $q$ be the minimal number of generators of $I$; then the number $\dim(B)-q$ is independent of the choice of $B$.
\end{proposition}

The minimal number of generators of $I$ is also equal to the dimension of the $k$-vector space $I\otimes_{B}k$, where $k$ denotes the residue field of $A$. We reduce at once to the case where $A$ is complete, since $\widehat{A}\simeq\widehat{B}/\widehat{I}$ with $\dim \widehat{B}=\dim B$ and $\rg_{k}(I\otimes_{B}k)=\rg_{k}(\widehat{I}\otimes_{\widehat{B}}k)$; for the same reason we\pageoriginale may suppose that th\sisi{ese}{e} ring\sisi{s}{} $B$ \sisi{are}{is} complete\sisi{}{}. Let $B$ and $B^{\prime}$ be two complete regular local rings, $f\colon B\to A$ and $f^{\prime}\colon B^{\prime}\to A$ two surjective homomorphisms, and $I=\Ker(f)$, $I^{\prime}=\Ker(f^{\prime})$. We must show that
$$
\dim B-\rg_{k}(I\otimes_{B}k)=\dim B^{\prime}-\rg_{k}(I^{\prime}\otimes_{B^{\prime}}k).
$$

First consider the case in which there is a factorization
of the form
$$
\xymatrix@C=15pt{ B\ar[rr]^-{f}\ar[rd]_-{g} & & A\\
& B^{\prime}\ar[ru]_-{f^{\prime}} & }$$
with $g$ surjective. Let $J=\Ker(g)$; then $J\subset I$ and $I/J=I^{\prime}$. Since $B^{\prime}$ is regular, $\dim(B^{\prime})=\dim(B)-\rg_{k}(J\otimes_{B}k)$, and $J$ is generated by elements belonging to a regular system of parameters of $B$. It follows that there is an exact sequence
$$0\to J\otimes_{B}k\to I\otimes_{B}k\to I/J\otimes_{B^{\prime}}k\to 0,
$$
and consequently
$$
\dim B-\rg_{k}(I\otimes_{B}k)=\dim B-\rg_{k}(J\otimes_{B}k)-\rg_{k}(I/J\otimes_{B^{\prime}}k)=\dim B^{\prime}-\rg_{k}(I^{\prime}\otimes_{B^{\prime}}k).
$$

The general case reduces to the preceding one; to see this, it is enough to show that one can find a complete regular local ring $B^{\prime\prime}$ and surjective homomorphisms $g\colon B^{\prime\prime}\to B$ and $g^{\prime}\colon B^{\prime\prime}\to B^{\prime}$ making the diagram\pageoriginale commutative
\begin{equation*} \label{eq:XIV.5.2.*} \tag{$*$}
\begin{array}{c}
\xymatrix@C=10pt@R=15pt{ & B\ar[rd]^-{f} & \\
B^{\prime\prime}\ar[ru]^-{g}\ar[rd]_-{g^{\prime}} & & A\\
& B^{\prime}.\ar[ru]_-{f^{\prime}} &
}
\end{array}
\end{equation*}
Now, if $W$ is a Cohen ring with residue field $k$, there is a local morphism $W\to A$ which lifts to $B$ and $B^{\prime}$ (\EGA IV~19.8.6), so that one has the commutative diagram
$$
\xymatrix@R=15pt@C=10pt{ & B\ar[rd] & \\
W\ar[ru]\ar[rd] & & A\\
& B^{\prime}.\ar[ru] & }$$
One can find integers $n$ and $n^{\prime}$ and surjective morphisms \hbox{$h\colon W\llbracket T_{1},\dots, T_{n}\rrbracket\!\to\! B$} and $h^{\prime}\colon W\llbracket T_{1}^{\prime},\dots, T_{n^{\prime}}^{\prime}\rrbracket\to B^{\prime}$ which are morphisms of $W$-algebras ($\EGA 0_{\textup{IV}}$ 19.8.8); if we then set $B^{\prime\prime}=W\llbracket T_{1},\dots, T_{n}, T_{1}^{\prime},\dots, T_{n^{\prime}}^{\prime}\rrbracket$ and define $g$ and~$g^{\prime}$ as morphisms of $W$-algebras such that
$$g(T_{i})=h(T_{i}), \quad g(T_{i}^{\prime})=b_{i}, \quad g^{\prime}(T_{i})=b_{i}^{\prime}, \quad g^{\prime}(T_{i}^{\prime})=h^{\prime}(T_{i}^{\prime}),
$$
where $b_{i}$ (respectively $b_{i}^{\prime}$) is an element of $B$ (respectively of $B^{\prime}$) lifting \sisi{$h^{\prime}.f^{\prime}(T_{i}^{\prime})$ (respectively $h.f(T_{i})$)}{$(f^{\prime}\circ h^{\prime})(T_{i}^{\prime})$ (respectively $(f\circ\nobreak h)(T_{i})$)}, the diagram \eqref{eq:XIV.5.2.*} is indeed commutative.

Proposition \Ref{XIV.5.2} justifies the following definition:

\begin{definition} \label{XIV.5.3}
Let\pageoriginale $A$ be a Noetherian local ring and $\widehat{A}$ its completion, which is thus isomorphic to the quotient of a complete regular local ring $B$ by an ideal $I$; if $q$ is the minimal number of generators of $I$, the \emph{geometric depth} of $A$ is the number
$$
\profgeom(A)=\dim B-q
$$
\end{definition}

\begin{proposition} \label{XIV.5.4} Let $A$ be a Noetherian local ring. Then
$$
\profgeom(A)\leq\dim A,
$$
with equality if and only if $A$ is \sisi{a}{a} complete intersection.
\end{proposition}

We may suppose that $A$ is complete. Let $A=B/I$, where $B$ is a complete regular local ring and $I$ an ideal of $B$. If $(x_{1},\dots, x_{q})$ is a minimal system of generators of $I$, then $\dim A\geq\dim B-q$, and saying that $\dim A=\dim B-q$ is equivalent to saying that $(x_{1},\dots, x_{q})$ belongs to a system of parameters of $B$ ($\EGA 0_{\textup{IV}}$~16.3.7); the proposition follows immediately.

\begin{proposition} \label{XIV.5.5} Let $A$ and $A^{\prime}$ be Noetherian local rings, and $f\colon A\to A^{\prime}$ a local homomorphism. Suppose that $f$ is flat and that, denoting by $k$ the residue field of $A$, $A^{\prime}\otimes_{A}k$ is a field, separable over $\sisi{K}{k}$. Then
$$
\profgeom(A)=\profgeom(A^{\prime}).
$$
\end{proposition}

After\pageoriginale replacing $A$ and $A^{\prime}$ by their completions, we may suppose that $A$ and $A^{\prime}$ are complete (it follows from ($\EGA 0_{\textup{III}}$~10.2.1) that the flatness hypothesis is preserved, and this is evident for the other hypotheses). Let $A=B/I$, where $B$ is a regular local ring and $I$ an ideal of $B$. Since $A^{\prime}$ is formally smooth over $A$ ($\EGA 0_{\textup{IV}}$~19.8.2), it follows from ($\EGA 0_{\textup{IV}}$~19.7.2) that one can find a complete Noetherian local ring $B^{\prime}$ and a local homomorphism $B\to B^{\prime}$ such that $B^{\prime}$ is a flat $B$-module and $B^{\prime}\otimes_{B}A\simeq A^{\prime}$. Thus $A^{\prime}\simeq B^{\prime}/IB^{\prime}$; moreover the ring $B^{\prime}$ is regular: indeed, if $\mm$ is the maximal ideal of $B$, $\mm B^{\prime}$ is the maximal ideal of $B^{\prime}$, and, since $\mm$ is generated by a regular sequence by the definition of \og regular\fg, $\sisi{B^{\prime}}{\mm B^{\prime}}$ is generated by a $B^{\prime}$-regular sequence ($\EGA 0_{\textup{IV}}$~15.1.14). Since one clearly has $\dim B=\dim B^{\prime}$, and since $I$ and $IB^{\prime}$ have the same minimal number of generators, the assertion follows.

\refstepcounter{subsection}\label{XIV.5.6}
\begin{enonce*}{Theorem 5.6\ndemark}\ndetext{Illusie subsequently proved the inequality $\prof_x(\ZZ/\ell^\nu\ZZ)\geq \profgeom_x(X/S)-\delta(x)+1$ for $x$ a point of $X$, where $X$ is a scheme of finite type over a trait $S$ whose residual characteristic is prime to $\ell$, and $\nu\geq 1$. If $S$ has characteristic zero, this is a consequence of Theorem~\Ref{XIV.5.6}; see (Illusie~L., {\og Perversité et variation\fg}, \emph{Manuscripta Math.} \textbf{112} (2003), p\ptbl 271-295).}
Let $A$ be an excellent local ring of characteristic zero. Then
$$
\profet(A)\geq\profgeom(A).
$$
\end{enonce*}

One may suppose that $A$ is strictly local and complete, since geometric depth and \'{e}tale depth are preserved under passage to the strict henselization and to the completion by \Ref{XIV.5.5} and \Ref{XIV.1.16}. Let $A\simeq B/I$, where $B$ is a complete regular local ring, and let $(f_{1},\dots, f_{q})$ be a minimal system of generators of the ideal $I$. Thus one has\pageoriginale
$$
\pi=\profgeom(A)=\dim B-q.
$$
Consider the closed immersion
$$Y=\Spec A\to X=\Spec B,
$$
and let\sisi{}{} $U=X-Y=\bigcup_{1\leq i\leq q}X_{f_{i}}$. If $a$ denotes the closed point of $X$, one must show that, for every prime number $p$, one has
$$
\H_{a}^{i}(Y, \ZZ/p\ZZ)=0\quad\textrm{ for }i<\pi.
$$
Since $B$ is regular and excellent, one has $\profet(B)=2\dim X$ (\cf \Ref{XIV.1.10}) and consequently $\H_{a}^{i}(X, \ZZ/p\ZZ)=0$ for $i<2\dim X$. Thus, to prove the theorem, it suffices to prove that the morphism
\begin{equation*} \label{eq:XIV.5.6.*} \tag{$*$}
\H_{a}^{i}(X, \ZZ/p\ZZ)\to\H_{a}^{i}(Y, \ZZ/p\ZZ)
\end{equation*}
is bijective for $i<\pi$. For this, apply the local Lefschetz theorem \sisi{(\Ref{XIV.4.5})}{\Ref{XIV.4.5}} with $n=\pi+q$, $c=q-1$, and hence $n-c-1=\pi$. Note that $U=X-Y$ is the union of the $q$ affine open subsets $X_{f_{i}}$. Let us show that, for every point $u$ of $U$, one has:
$$
\profet_{u}(X)\geq\pi+q-\dim\big(\overline{\{ u\} }\big)=\dim\Oo_{X, u}$$
(where $\overline{\{ u\} }$ denotes the closure of $u$ \emph{in} $X-\overline{\{ a\} }$). Indeed, it follows from \Ref{XIV.1.10} that
$$
\profet_{u}(X)=2\dim\Oo_{X, u}\geq\dim\Oo_{X, u}.
$$
Using\pageoriginaled \Ref{XIV.4.5}, one sees that \eqref{eq:XIV.5.6.*} is bijective for $i<\pi$, which completes the proof of the theorem.

\begin{corollaire} \label{XIV.5.7} Let $S$ be the spectrum of a field of characteristic zero (\resp an excellent henselian local scheme of characteristic zero), and let $f\colon X\to S$ exhibit $X$ as a scheme proper over $S$ (\resp over $S-\{ s\}$). Let $U$ be a union of $c+1$ open subsets of $X$ that are \emph{affine}, let $Y$ be a closed subscheme with underlying space $X-U$, and let $n$ and $m$ be integers $>0$. Suppose that, for every point $u$ of $U$, one has
$$
\profgeom(\Oo_{X, u})\geq n-\dim\big(\overline{\{ u\} }\big)$$
($\overline{\{ u\} }$ being the closure of $u$ in $X$). Then the canonical morphism
$$
\H^{i}(X, \ZZ/m\ZZ)\to\H^{i}(Y, \ZZ/m\ZZ)$$
is bijective for $i<n-c-1$, and injective for $i=n-c-1$.
\end{corollaire}

Apply \Ref{XIV.4.5} and \Ref{XIV.4.6}. The hypotheses on \'{e}tale depth at the points of $U$ are satisfied because, by \Ref{XIV.5.6}, one has \sisi{$$
\profet_{u}(X)\geq\profgeom\Oo_{X, u}\geq n-\dim\big(\overline{\{ u\} }\big).
$$
}{$$
\profet_{u}(X)\geq\profgeom(\Oo_{X, u})\geq n-\dim\big(\overline{\{ u\} }\big).
$$
}
\enlargethispage{\baselineskip}%

\section{Open questions} \label{XIV.6}

\subsection{\texorpdfstring{}{XIV.6.1}} \label{XIV.6.1} One may ask whether the implication \textup{(ii)} \ALORS \textup{(i)} of \Ref{XIV.4.2} holds more generally for torsion sheaves $F$, not necessarily annihilated by a given integer $m$ and not necessarily constructible\refstepcounter{pageoriginal}\unskip\marginpar{\raggedright\footnotesize\bfseries\thepageoriginal}\nobreak. In the case where $S$ is not of characteristic zero, it seems possible that this implication remains valid, even for sheaves of $p$-torsion ($p$ being the residue characteristic). Finally, it is also unclear whether the hypothesis that $S$ is excellent can be removed.

\subsection{\texorpdfstring{}{XIV.6.2}} \label{XIV.6.2}
Let $X$ be a scheme proper over a field $k$, or else the complement of the closed point of a henselian local scheme, and let $j:Y \to X$ be a closed subscheme of $X$ whose complement $U$ is affine. Then, if $F$ is a sheaf of sets on $X$ or a sheaf of not necessarily commutative groups, the statements \Ref{XIV.4.5} or \Ref{XIV.4.6} and \Ref{XIV.4.9} still make sense for such an $F$, provided that one restricts to small values of $n$. If $u$ is a point of $U$, denote by $\bar{u}$ a geometric point lying over~$u$, by $X(\bar{u})$ the strict localization of $X$ at $\bar{u}$, and by $F_{\bar{u}}$ the fiber of $F$ at $\bar{u}$. Then, possibly making certain hypotheses on $X$ and $F$, for example by assuming $X$ excellent (possibly of characteristic zero, or of equal characteristic using resolution of singularities) and $F$ ind-finite (or, if necessary, even $\LL$-ind-finite with $\LL$ prime to the characteristic of $X$), one would like to prove the following statements:
\begin{enumeratea}
\item
Let $F$ be a sheaf of sets (\resp a sheaf of groups), and suppose that, for every point $u$ of $U$, one has
$$F_{\bar{u}} \to H^0(X(\bar{u})-\bar{u}, F)\text{ injective if } \dim(\overline{\{u\}}) \leq 1$$
(that is, for such a $u$, one has $\prof_u(F) \geq 1$). Then, as $V$ ranges over the set of open neighborhoods of $Y$, the canonical morphism\pageoriginale
$$
\varinjlim_V H^0(V, F) \to H^0(Y, j^*F)
$$
is bijective (\resp one has the preceding conclusion and, moreover, the morphism $\varinjlim_V H^1(V, F) \to H^1(Y, j^* F)$ is injective). If $F$ is constructible, one may replace the $\varinjlim$ by the cohomology of $V$ for $V$ \og sufficiently small\fg.
\item
Let $F$ be a sheaf of sets (\resp a sheaf of groups), and suppose that, for every point $u$ of $U$, one has $\prof_u(F) \geq 2 -\dim(\overline{\{u\}})$, which is also expressed by the relations
\begin{gather*}
F_{\bar{u}}\to H^0(X(\bar{u})-\bar{u}, F)\text{ is bijective if } \dim(\overline{\{u\}})=0\\
F_{\bar{u}}\to H^0(X(\bar{u})-\bar{u}, F)\text{ is injective if } \dim(\overline{\{u\}})=1.
\end{gather*}
Then the canonical morphism
$$H^0(X, F) \to H^0(Y, j^*F)$$
is bijective (\resp one has the preceding conclusion and, moreover, the morphism $H^1(X, F) \to H^1(Y, j^*F)$ is injective).
\item
Let $F$ be an ind-finite sheaf of groups. Suppose that, for every point $u$ of $U$, one has
\begin{gather*}
F_{\bar{u}}\to H^0(X(\bar{u})-\bar{u}, F)\text{ bijective if } \dim(\overline{\{u\}})=0\text{ or }1,\\
F_{\bar{u}}\to H^0(X(\bar{u})-\bar{u}, F)\text{ injective if } \dim(\overline{\{u\}})=2.
\end{gather*}
Then\pageoriginale, as $V$ ranges over the set of open neighborhoods of $Y$, the canonical morphisms
$$
\varinjlim_V H^0(V, F) \to H^0(Y, j^*F) \text{ and } \varinjlim_V H^1(V, F) \to H^1(Y, j^*F)
$$
are bijective. If $F$ is constructible, one may replace the $\varinjlim$ by the cohomology of~$V$ for $V$ sufficiently small.
\item
Let $F$ be a sheaf of groups. Suppose that, for every point $u$ of $U$, one has $\prof_u(F) \geq 3-\dim(\overline{\{u\}})$, which is also expressed by the conditions
\begin{gather*}
F_{\bar{u}}\to H^0(X(\bar{u})-\bar{u}, F)\text{ bijective, and } H^1(X(\bar{u})-\bar{u}, F)=0 \text{ if } \dim(\overline{\{u\}})=0,\\
F_{\bar{u}}\to H^0(X(\bar{u})-\bar{u}, F)\text{ bijective if } \dim(\overline{\{u\}})=1,\\
F_{\bar{u}}\to H^0(X(\bar{u})-\bar{u}, F)\text{ injective if } \dim(\overline{\{u\}})=2.
\end{gather*}
Then the canonical morphisms
$$H^0(X, F) \to H^0(Y, j^*F)\text{ and } H^1(X, F) \to H^1(Y, j^*F)$$
are bijective.
\end{enumeratea}

As evidence in favor of these statements\refstepcounter{toto}\nde{\label{MiR}all the statements of~\Ref{XIV.6.2}, apart from the constructible variants, were proved by Mrs Raynaud; see (Raynaud~M., {\og Th\'eor\`emes de Lefschetz en cohomologie des faisceaux coh\'erents et en cohomologie \'etale. Application au groupe fondamental\fg}, \emph{Ann. Sci. \'Ec. Norm. Sup. (4)} \textbf{7} (1974), p\ptbl 29--52, Corollary III.1.3).}, let us mention \Ref{XIII}~\Ref{XIII.2.1}, \Ref{X}~\Ref{X.3.4}, and \Ref{XII}~\Ref{XII.3.5}. Note that, by the argument of \Ref{XIV.4.8} and \Ref{XIV.4.9}, statement a) (\resp c)) would follow from b) (\resp d)).

\subsection{\texorpdfstring{}{XIV.6.3}} \label{XIV.6.3}
The following statement, analogous to \Ref{XIV.5.6}, would follow from d): if $A$ is\pageoriginale a noetherian local ring (possibly excellent) and $\prof \geom (A) \geq 3$, then $\prof \hop (A) \geq 3$. To see this, realize $Y'=\Spec A$ as a closed subscheme of a regular local scheme $X'=\Spec B$, whose complement is the union of $q$ affine open subsets, with the relation $\dim B -q=\prof \geom(A)$. For every point $x$ of $X'$, if $n=\dim B$, one has $\prof \hop_x(X') \geq \inf(3, n-\dim(\overline{\{x\}}))$ (\cf \Ref{XIV.1.11}), and one deduces from d) that this implies $\prof \hop_y(Y') \geq \inf(3, n-q-\dim(\overline{\{y\}}))$ for every point $y$ of $Y'$. The result is then obtained by taking $y$ to be the closed point of $Y'$.

\subsection{\texorpdfstring{}{XIV.6.4}} \label{XIV.6.4}
A variant of \Ref{XIV.4.2}, at least of the implication \textup{(ii)} \ALORS \textup{(i)}, should also hold in the complex analytic case, provided one works with \og analytically constructible\fg sheaves (\cf \Ref{XIII}); the proof would be analogous to that of \Ref{XIV.4.2}, using the duality theory of J.-L. Verdier. Note, however, that for the complex analytic analogue of the noncommutative variants mentioned in \Ref{XIV.6.2}, not even a method of approach is available for the statements concerning the fundamental group suggested by the results of Expos\'{e}s \Ref{X}, \Ref{XII}, and \Ref{XIII}, recalled at the end of \Ref{XIV.6.2}. Indeed, the methods of the Seminar seem inextricably tied to the case of \emph{finite} coverings (which may be studied in terms of coherent sheaves of \sisi{Algebras}{algebras}).

\begin{thebibliography}{9}
\bibitem[1]{XIV.1}\sisi{%
M\ptbl Artin et B\ptbl Mazur, Homotopy of Varieties in Etale Topology, in Proceedings of a Conference on Local Fields, Springer (1967).}
{%
\textsc{M\ptbl Artin \& B\ptbl Mazur} -- {\og Homotopy of Varieties in Etale Topology\fg}, in \emph{Proceedings of a Conference on Local Fields}, Springer, 1967.
}
\bibitem[2]{XIV.2}\sisi{%
J.-P\ptbl Serre, Cohomologie Galoisienne, Springer-Verlag, 1964.}
{%
\textsc{J.-P\ptbl Serre} -- \emph{Cohomologie Galoisienne}, Springer-Verlag, 1964.
}
\bibitem[3]{XIV.3}\sisi{%
J.-L\ptbl Verdier, Alg\`ebre homologique et Cat\'egories d\'eriv\'ees, North.\ Holland Pub.\ Cie.\, to appear}
{%
\textsc{J.-L\ptbl Verdier} -- \emph{Des cat\'egories d\'eriv\'ees des cat\'egories ab\'eliennes}, with a preface by Luc Illusie, edited by Georges Maltsiniotis, Ast\'erisque, vol.~239, Soci\'et\'e math\'ematique de France, Paris, 1996.
}
\end{thebibliography}

\chapter*{Index of notation}
\markboth{INDEX OF NOTATION}{INDEX OF NOTATION}

\begin{tabular}{p{9cm}p{3cm}}
$\Gamma_{Z}, \ \underline{\Gamma}_{Z}$\ \dotfill & \Ref{I}~\Ref{I.1}\\
$i_{!}, \ i^{!}$\ \dotfill & \Ref{I}~\Ref{I.1}\\
$\ZZ_{Z, X}$\ \dotfill & \Ref{I}~\Ref{I.1.6}\\
$H^{i}_{Z}(X, F), \underline{H}^{i}_{Z}(F)$\ \dotfill & \Ref{I}~\Ref{I.2}\\
$H^{i}_{J}(M), H^{i}((f), M), H^{i}(M)$\ \dotfill & \Ref{IV}~\Ref{IV.5.4}, \Ref{V}~\Ref{V.2} (p\ptbl\sisi{4}{\pageref{pV2.4}})\\
$\Ass M$, $\Supp M$, $\Ann \ M\sisi{}{,} \rr(\aaa)$\ \dotfill & \Ref{III}~\Ref{III.1.1}\\
$\prof_{I}(M), \prof_{Y}(F)$\ \dotfill & \Ref{III}~\Ref{III.2.3}, \Ref{III.2.8}\\
$\Ab$\ \dotfill & \Ref{IV}~\Ref{IV.1} (p\ptbl\sisi{1}{\pageref{pIV.1}})\\
$\Hom^{.}(F^{.}, G^{.})$\ \dotfill & \Ref{V}~\Ref{V.1.1}\\
$\Ext^{i}_{Z}(X;F, G), \underline{\Ext}^{i}_{Z}(F, G)$\ \dotfill & \Ref{VI}~\Ref{VI.1.1}\\
${\Et}(X), \LL(Z)$\ \dotfill & \Ref{X} (p\ptbl\sisi{1}{\pageref{pX.1}})\\
$\Lef(X, Y), \Leff(X, Y)$\ \dotfill & \Ref{X}~\Ref{X.2} (p\ptbl\sisi{2}{\pageref{pX.2}})\\
$\underline{P}(Z), \Pic(Z)$\ \dotfill & \Ref{XI} (p\ptbl\sisi{1}{\pageref{pXI.1}})\\
$\Pi_{i}^{x}(X)$\ \dotfill & \Ref{XIII} (p\ptbl\sisi{15}{\pageref{pXIII.15}})\\
$D^{+}(X)$\ \dotfill & \Ref{XIV}~\Ref{XIV.1.0}\\
$\prof_{Y}(F) \geq n$\ \dotfill & \Ref{XIV}~\Ref{XIV.1.2}\\
$\prof_{Y}^{\LL}(X) \geq n$\ \dotfill & \Ref{XIV}~\Ref{XIV.1.2}\\
$\profet_{Y}(X)$\ \dotfill & \Ref{XIV}~\Ref{XIV.1.2}\\
$\LL$\ \dotfill & \Ref{XIV}~\Ref{XIV.1.2}\\
$\prof_{x}(F) \geq n$\ \dotfill & \Ref{XIV}~\Ref{XIV.1.7}\\
$\prof_{x}^{\LL}(X)\geq n$\ \dotfill & \Ref{XIV}~\Ref{XIV.1.7}\\
$\prof\ \hop^{\LL}_{x}(X) \geq 3$\ \dotfill & \Ref{XIV}~\Ref{XIV.1.7}\\
$\profet_{x}(X)$\ \dotfill& \Ref{XIV}~\Ref{XIV.1.7}\\
$\delta_{t}(x), \ \delta (n)$\ \dotfill & \Ref{XIV}~\Ref{XIV.2.1}\\
$H^{i}_{Z!}(X/S, F)$\ \dotfill & \Ref{XIV}~\Ref{XIV.4.0}
\end{tabular}

\chapter*{Terminological index}
\markboth{TERMINOLOGICAL INDEX}{TERMINOLOGICAL INDEX}

\begin{footnotesize}
\noindent
\begin{tabular}{lll}
Auslander-Buchsbaum (th\'eor\`eme de -) & & \Ref{XI}~\Ref{XI.3.13}\\
comparaison (th\'eor\`eme de -) & comparison theorem & \Ref{IX}~\Ref{IX.1.1}, \Ref{XII}~\Ref{XII.2.1}\\
dualit\'e locale (th\'eor\`eme de -) & local duality theorem & \Ref{V}~\Ref{V.2.1}\\
dualit\'e projective (th\'eor\`eme de -) & projective duality theorem & \Ref{XII}~\Ref{XII.1.1} \\
enveloppe injective & injective envelope & \Ref{IV} (p\ptbl\sisi{15}{\pageref{pIV.13}})\\
excision (th\'eor\`eme d' -) & excision theorem & \Ref{I}~\Ref{I.2.2}, \Ref{VI}~\Ref{VI.1.3}\\
existence (th\'eor\`eme d' -) & existence theorem & \Ref{IX}~\Ref{IX.2.1}, \Ref{XII}~\Ref{XII.3.1}\\
extension essentielle & essential extension & \Ref{IV} (p\ptbl\sisi{15}{\pageref{pIV.13}})\\
finitude (th\'eor\`eme de -) & finiteness theorem & \Ref{VIII}~\Ref{VIII.2.1}, \Ref{XII}~\Ref{XII.1.5}\\
foncteur dualisant & dualizing functor & \Ref{IV}~\Ref{IV.4.1}, \Ref{IV.4.2}\\
forme r\'esidu & residue form & \Ref{IV}~\Ref{IV.5.5}\\
g\'eom\'etriquement factoriel, \resp & geometrically factorial (\resp &\\
g\'eom. parafactoriel (anneau local -) & geom. parafactorial) local ring & \Ref{XIII} (p\ptbl20, 24)\\
groupes d'homotopie locale & local homotopy groups & \Ref{XIII} (p\ptbl\sisi{15}{\pageref{pXIII.15}})\\
Gysin (homomorphisme de -) & Gysin homomorphism & \Ref{I} (p\ptbl\sisi{13}{\pageref{hgysin}})\\
Hartshorne (th\'eor\`eme de -) & & \Ref{III}~\Ref{III.3.6}\\
Lefschetz (condition de -) & Lefschetz condition & \Ref{X}~\Ref{X.2} (p\ptbl\sisi{2}{\pageref{pX.2}})\\
Lefschetz (condition de - effective) & effective Lefschetz condition & \Ref{X}~\Ref{X.2} (p\ptbl\sisi{2}{\pageref{pX.2}})\\
Lefschetz (th\'eor\`eme de - affine) & affine Lefschetz theorem & \Ref{XIV}~\Ref{XIV.3.1}\\
module dualisant & dualizing module & \Ref{IV}~\Ref{IV.4.1}\\
\sisi{module orthogonal d'un module}{orthogonal d'un sous-module} & \sisi{module orthogonal of a module}{orthogonal of a submodule} & \Ref{IV}~\Ref{IV.5}\\
parafactoriel (couple - de pr\'e\-sch\'emas) & parafactorial pair of preschemes & \Ref{XI}~\Ref{XI.3.1}\\
parafactoriel (anneau local -) & parafactorial local ring & \Ref{XI}~\Ref{XI.3.2}\\
profondeur & depth & \Ref{III}~\Ref{III.2.3}\\
profondeur \'etale & etale depth & \Ref{XIV}~\Ref{XIV.1.2}\\
profondeur g\'eom\'etrique d'un anneau & geometrical depth of a noetherian &\\
local noeth\'erien & local ring & \Ref{XIV}~\Ref{XIV.5.3}\\
profondeur homotopique & homotopical depth & \Ref{XIII}~\Ref{XIII.6} def\ptbl1 (p\ptbl\sisi{4}{\pageref{XIII.6.1}})\\
profondeur homotopique (pour $\LL$) & homotopical depth (with respect to $\LL$) & \Ref{XIV}~\Ref{XIV.1.2}\\
profondeur homotopique rectifi\'ee & rectified homotopical depth & \Ref{XIII} def\ptbl2 (p\ptbl\sisi{5}{\pageref{XIII.6.2}})\\
\end{tabular}
\newpage
\begin{tabular}{lll}
pur (anneau local -) & pure local ring & \Ref{X}~\Ref{X.3.2} \\
pure (couple - de pr\'esch\'emas) & pure pair of preschemes & \Ref{X}~\Ref{X.3.1}\\
puret\'e (th\'eor\`eme de - de Zariski-Nagata) & purity theorem of Zariski-Nagata & \Ref{X}~\Ref{X.3.4}\\
puret\'e cohomologique (th\'eor\`eme de -) & cohomological purity theorem & \Ref{XIV}~\Ref{XIV.1.11}\\
r\'egulier ($M$-r\'egulier) & $M$-regular & \Ref{III.2.1}\\
Samuel (conjecture de -) & & \Ref{XI}~\Ref{XI.3.14}\\
semi-puret\'e cohomologique (th\'eor\`eme de -) & cohomological semi-purity theorem & \Ref{XIV}~\Ref{XIV.1.10}\\
strictement local (anneau -) & strictly local ring & \Ref{XIII}~\Ref{XIII.6} (p\ptbl\sisi{23}{\pageref{pXIII.6}})\\
Zariski-Nagata (th\'eor\`eme de puret\'e de -) & purity theorem of Zariski-Nagata & \Ref{X}~\Ref{X.3.4}
\end{tabular}
\end{footnotesize}

\backmatter

% END INCLUDED active_source/parts/14_expose_XIV.tex


\end{document}

% SGA-CUMULATIVE-END 02_SGA2_EN_COMPLETE_READER.tex

% SGA-CUMULATIVE-BEGIN 03_SGA3_EN_COMPLETE_READER.tex
\documentclass[11pt,a4paper,oneside]{book}

\usepackage[T1]{fontenc}
\usepackage[utf8]{inputenc}
\usepackage[english]{babel}
\usepackage{lmodern}
\usepackage[a4paper,inner=30mm,outer=27mm,top=28mm,bottom=30mm]{geometry}
\usepackage{amsmath,amssymb,amsthm,mathrsfs,mathtools,stmaryrd}
\usepackage{tikz,tikz-cd}
\usetikzlibrary{
  arrows,
  arrows.meta,
  calc,
  decorations.pathmorphing,
  matrix,
  positioning,
  shapes,
  shapes.geometric
}
\usepackage{graphicx,float,adjustbox,enumitem,multicol}
\usepackage{etoolbox}
\usepackage{xcolor}
\usepackage[hidelinks,hypertexnames=false,pdfstartview={},linktoc=page,
  pdftitle={SGA 3 - Complete English Translation and Native-Diagram Reconstruction},
  pdfauthor={SGA 3 English reconstruction project},
  pdfsubject={Cumulative English reader for SGA 3: Introduction, Exposes I-XXVI, Tome-I subject index, Tome-III guide, and terminal index},
  pdfkeywords={SGA 3, group schemes, English translation, native TeX diagrams}]{hyperref}

\setlength{\emergencystretch}{3em}
\raggedbottom

\theoremstyle{plain}
\newtheorem{theorem}{Theorem}[section]
\newtheorem{proposition}[theorem]{Proposition}
\newtheorem{corollary}[theorem]{Corollary}
\newtheorem{lemma}[theorem]{Lemma}
\newtheorem{scholium}[theorem]{Scholium}
\newtheorem{application}[theorem]{Application}
\newtheorem{complements}[theorem]{Complements}

\theoremstyle{definition}
\newtheorem{definition}[theorem]{Definition}
\newtheorem{example}[theorem]{Example}

\theoremstyle{remark}
\newtheorem{remark}[theorem]{Remark}

\providecommand{\Ga}{\mathbb G_a}
\providecommand{\Gm}{\mathbb G_m}
\providecommand{\Norm}{\operatorname{Norm}}
\providecommand{\Centr}{\operatorname{Centr}}
\providecommand{\Ker}{\operatorname{Ker}}
\providecommand{\Dom}{\operatorname{Dom}}
\providecommand{\rank}{\operatorname{rank}}
\providecommand{\Spec}{\operatorname{Spec}}
\providecommand{\Lie}{\operatorname{Lie}}
\providecommand{\intaut}{\operatorname{int}}
\providecommand{\Fib}{\operatorname{Fib}}
\providecommand{\Aut}{\operatorname{Aut}}
\providecommand{\Autext}{\operatorname{Autext}}
\providecommand{\rgred}{\operatorname{rgred}}
\providecommand{\rgss}{\operatorname{rgss}}
\providecommand{\GL}{\operatorname{GL}}
\providecommand{\Qp}{\mathbb{Q}_{+}}
\providecommand{\cR}{\mathcal{R}}
\providecommand{\VR}{\mathcal{V}}
\providecommand{\Nil}{\operatorname{Nil}}

\providecommand{\SGAParagraph}[1]{%
  \par\medskip\noindent\textbf{#1.}---\quad
}
\providecommand{\NDE}[1]{%
  \footnote{\textit{Editors' note.} #1}%
}
\newenvironment{sgastatement}[1]
  {\par\medskip\noindent\textbf{#1.}---\ \itshape}
  {\par\medskip}
\providecommand{\SGAQED}{\hfill\(\square\)\par\smallskip}
\providecommand{\Idx}[2]{%
  \par\noindent\hangindent=1.2em\hangafter=1
  #1\if\relax\detokenize{#2}\relax\else,\ #2\fi\par
}
\providecommand{\SGASection}[2]{%
  \section*{#1.\ #2}%
  \addcontentsline{toc}{section}{#1.\ #2}%
}
\providecommand{\SGASubsection}[2]{%
  \subsection*{#1.\ #2}%
  \addcontentsline{toc}{subsection}{#1.\ #2}%
}
\providecommand{\SGAStatement}[2]{%
  \par\medskip\noindent\textbf{#1.}---\ {\itshape #2\par}\medskip
}
\providecommand{\SGAProof}[1]{%
  \par\smallskip\noindent\textit{#1.}\quad
}

\makeatletter
\providecommand{\SGARefTarget}[1]{\phantomsection\label{#1}}
\providecommand{\SGAInlineRefTarget}[1]{%
  \begingroup
    \def\@currentHref{#1}%
    \label{#1}%
  \endgroup}
\providecommand{\SGAInlineEquationTarget}[1]{%
  \begingroup
    \def\@currentHref{#1}%
    \ltx@label{#1}%
  \endgroup}
\let\SGA@PendingSectionTarget\@empty
\providecommand{\SGANextSectionTarget}[1]{%
  \gdef\SGA@PendingSectionTarget{#1}}
\providecommand{\SGA@UsePendingSectionTarget}{%
  \ifx\SGA@PendingSectionTarget\@empty\else
    \expandafter\SGAInlineRefTarget\expandafter{%
      \SGA@PendingSectionTarget}%
    \global\let\SGA@PendingSectionTarget\@empty
  \fi}
\renewcommand{\SGASection}[2]{%
  \section*{\SGA@UsePendingSectionTarget#1.\ #2}%
  \addcontentsline{toc}{section}{#1.\ #2}%
}
\renewcommand{\SGASubsection}[2]{%
  \subsection*{\SGA@UsePendingSectionTarget#1.\ #2}%
  \addcontentsline{toc}{subsection}{#1.\ #2}%
}
\pdfstringdefDisableCommands{%
  \def\SGAInlineRefTarget#1{}%
  \def\SGAInlineEquationTarget#1{}%
  \def\SGANextSectionTarget#1{}%
}
\providecommand{\SGARefLink}[2]{%
  \@ifundefined{r@#1}{#2}{\hyperref[#1]{#2}}%
}
\providecommand{\GuideChapter}[2]{%
  \@dottedtocline{0}{0em}{0em}{#1}{#2}}
\providecommand{\GuideChapterSplit}[3]{%
  \par\noindent\textbf{#1}\par\nobreak
  \@dottedtocline{0}{2em}{0em}{#2}{#3}}
\providecommand{\GuideSection}[2]{%
  \@dottedtocline{1}{1.45em}{2.25em}{#1}{#2}}
\newcommand{\SGAInputRoot}[1]{%
  \def\input@path{{inputs/#1/}}%
}
\newcommand{\SGATagsLeft}{%
  \tagsleft@true
  \let\veqno\@@leqno
}
\newcommand{\SGATagsRight}{%
  \tagsleft@false
  \let\veqno\@@eqno
}
\makeatother

\newcommand{\SGAStartExpose}[1]{%
  \clearpage
  \setcounter{section}{0}%
  \setcounter{subsection}{0}%
  \setcounter{theorem}{0}%
  \setcounter{footnote}{0}%
  \setcounter{figure}{0}%
  \SGAInputRoot{#1}%
}

\title{Seminar on Algebraic Geometry 3\\
  \large Group Schemes\\[0.75em]
  \normalsize Complete English Translation and Native-Diagram
  Reconstruction\\
  \small Introduction; Expos\'es I--XXVI; Tome-I subject index;\\
  Tome-III mathematical guide and terminal index}
\author{SGA 3 English reconstruction project}
\date{29 July 2026}

\begin{document}

\frontmatter
\maketitle

\chapter*{Reconstruction, provenance, and rights status}
\addcontentsline{toc}{chapter}{Reconstruction, provenance, and rights status}
This cumulative reader contains the complete English reconstruction of
the Introduction and Expos\'es~I--XXVI, together with the Tome-I subject
index, the Tome-III mathematical guide, and the terminal index.  All
delivered mathematical diagrams are native \TeX{} constructions; the
active source contains no raster-diagram call.

The current Polo--Gille re-edition PDFs control the French text,
mathematics, numbering, notes, and diagram appearance.  No public
editable Polo--Gille source was available when this reconstruction was
undertaken, so the English text, formulas, document structure, and
diagrams were reconstructed in \LaTeX{} from those PDF authorities.
Existing OCR was used only as read-only locator and drafting material;
it is not source authority.

Jacob Reinhold's CC~BY~4.0 English Markdown is credited as a comparison
and drafting lineage.  It is not the source authority.  That CC~BY~4.0
statement applies to the comparison translation and does not by itself
grant rights in the underlying French edition or in this cumulative
reconstruction.  This is an internal native-source checkpoint pending
fresh cumulative reference, privacy, rights, packaging, and independent
release review; no publication or public-license claim is made here.

\tableofcontents

\mainmatter

\SGAInputRoot{INTRO}
% BEGIN INCLUDED SGA3_English_current_buildable_source_20260823/source/inputs/INTRO/body.tex
% Exact R15 special component.
% BEGIN INCLUDED SGA3_English_current_buildable_source_20260823/source/inputs/INTRO/tex/components/00_introduction.tex
% SGA 3, Editorial Notice and Introduction (Loop 1 English reconstruction)
% Source component: AaIntro-18fev11.pdf, local pages 1--6.
% Combined Polo--Gille reader: physical PDF pages 5--10.
% OCR witness: ocr_fo/AaIntro/full/AaIntro-18fev11/AaIntro-18fev11.md.
% The PDF controls structure, mathematical notation, and footnotes.

\begingroup
\newcommand{\SGAIIIIntroFootnote}[2]{%
  \begingroup
  \renewcommand{\thefootnote}{#1}%
  \footnote{#2}%
  \endgroup
}

\chapter*{Editorial Notice}
\addcontentsline{toc}{chapter}{Editorial Notice}
\markboth{EDITORIAL NOTICE}{EDITORIAL NOTICE}

We present here a slightly revised edition of the original Seminar, whose purpose
and contents are described in the Introduction.  The revision consisted mainly of
correcting typographical errors, adding a few supplementary remarks or references
in footnotes, dividing the work in its present form into three volumes, each with a
detailed table of contents and an index of notation, and adding a terminological index at
the end of Volume~3.  In addition, Expos\'e~VI\textsubscript{B}, by J.-E. Bertin,
was partially rewritten by the author, especially \SGARefLink{sga3:intro:section:5}{Sections~5} and~10, so some
references to that Expos\'e differ from the references to the original version.  A
list of the Expos\'es in the Seminar appears at the beginning of this volume.

Since the first edition of the present Seminar appeared, all of Chapter~IV of the
\emph{\'El\'ements de G\'eom\'etrie Alg\'ebrique} has been published, making certain
passages of the Seminar unnecessary.  We have occasionally indicated in footnotes
the relevant references to EGA~IV that allow those passages to be bypassed.

For another treatment of algebraic groups that systematically uses the language of
schemes, we refer to M. Demazure and P. Gabriel, \emph{Groupes Alg\'ebriques}
(North-Holland \& Masson et Cie).  Unlike the present Seminar, that book assumes no
knowledge of algebraic geometry, but includes all the required preliminaries from
scheme theory; it can therefore serve as an introduction to the study of this
Seminar.  (It also treats subjects not covered in the Seminar, such as the
Dieudonn\'e-style structure theory of commutative affine algebraic groups, in
\emph{loc. cit.}, Chapter~V.)

\begin{flushright}
Bures-sur-Yvette, March 1970
\end{flushright}

\clearpage
\thispagestyle{empty}
\mbox{}
\clearpage

\chapter*{Introduction}
\addcontentsline{toc}{chapter}{Introduction}
\markboth{INTRODUCTION}{INTRODUCTION}

\section*{\SGAInlineRefTarget{sga3:intro:section:1}1. The purpose of this Seminar is twofold.}

On the one hand, we aim to provide convenient foundations for the theory of group
schemes in general.  Expos\'es~I--IV will supply the indispensable preliminary
exercises in schematic and categorical syntax.  To obtain a language that fits
geometric intuition without effort, and to avoid circumlocutions that become
intolerable in the long run, we always identify a prescheme~$X$ over another
prescheme~$S$ with the functor
\[
  (\mathbf{Sch})_{/S}^{\circ}\longrightarrow \mathbf{Set}
\]
that it represents\SGAIIIIntroFootnote{\ensuremath{\ast}}{Such a viewpoint seems
to have been considered for the first time eight or nine years ago, in connection
with the theory of formal groups, by P. Cartier, who unfortunately did not take the
trouble to state it precisely and systematize it as it deserved.}.  It is therefore
necessary to formulate many definitions so that they apply to arbitrary functors,
whether representable or not.  Moreover, nearly all the functors we shall use will
be ``sheaves'' (for the ``faithfully flat quasi-compact topology'').  Expos\'e~IV,
which treats groups only incidentally, sketches the language of ``localization'' and
sheaves; this language also proves very convenient in questions concerning the
representability of functors.  Above all, that Expos\'e will provide the most
convenient framework for the questions of passage to a quotient that arise later.

Expos\'e~V gives some general results on the existence of quotients.  They are taken
up again in Expos\'e~VI\textsubscript{A} for the quotient of an algebraic group over
a field (or, more generally, over an Artinian ring) by a
subgroup\SGAIIIIntroFootnote{\ensuremath{\ast\ast}}{For a more thorough study of
passage to the quotient, especially by reductive groups, see D. Mumford's important
work \emph{Geometric Invariant Theory}, Ergebnisse der Mathematik, vol.~34,
Springer, 1965.  We note that on one important point the terminology of that book
does not agree with ours: over a field of characteristic $p>0$, the groups Mumford
calls ``reductive''\textsuperscript{\SGARefLink{sga3:intro:section:1}{(1)}} are the smooth groups of multiplicative
type in the sense of the Seminar (cf. Expos\'e~IX).  One may reasonably suppose
that Mumford adopted this meaning of ``reductive,'' which loses its significance
over a base that is not a field, provisionally and as a kind of stopgap (and this is
approximately what Mumford explains, for other reasons, as early as the second
paragraph of his preface).  \textsuperscript{\SGARefLink{sga3:intro:section:1}{(1)}}\emph{Editor's note:} See the
remarks added at the end of this Introduction.}.  That Expos\'e and the following
Expos\'e~VI\textsubscript{B} also contain various elementary results specific to
algebraic groups over a field that are used routinely later.

Expos\'e~VII studies certain phenomena related to the characteristic of the base
field and, in particular, develops with the appropriate generality the
correspondence between radicial group schemes of height~1 and restricted
$p$-Lie algebras.

Finally, Expos\'e~XVIII contains the scheme-theoretic generalization of Weil's
theorem on the ``birational'' definition of algebraic groups.

On the other hand, we intend to generalize the Borel--Chevalley structure theory of
affine algebraic groups to groups over an arbitrary base prescheme.  In fact, while
the Seminar notes were being written, it became apparent that the affine hypothesis
was unnecessary for many results of the theory.  The most complete results are of
course obtained for ``semisimple group schemes,'' or more generally ``reductive
group schemes,'' with which we shall be exclusively concerned beginning with
Expos\'e~XIX.  Chevalley himself had already constructed the ``T\^ohoku groups''
over the ring of integers, and that construction will be taken up in the present
Seminar.  The principal uniqueness
theorem\SGAIIIIntroFootnote{\SGARefLink{sga3:intro:section:2}{(2)}}{\emph{Editor's note:} This refers to
Corollary~XXIII.5.6, which follows easily from the uniqueness theorem for split
reductive groups (cf. XXIII, \SGARefLink{sga3:XXIII:4.1}{Theorem~4.1} and \SGARefLink{sga3:XXIII:5.2}{Corollary~5.2}), since every semisimple
$S$-group is a ``twisted form'' (for the \(\acute{e}\)tale topology) of a T\^ohoku
group (cf. XXII, \SGARefLink{sga3:xxii:cor:2.3}{Corollary~2.3}).} gives a simple characterization of the ``twisted''
variants of these T\^ohoku groups over a base prescheme~$S$: they are the affine,
smooth group schemes over~$S$ whose geometric fibers are connected semisimple
groups in the usual sense\SGAIIIIntroFootnote{\ensuremath{\ast\ast\ast}}{This is
the essential result of M. Demazure's thesis, \emph{Sch\'emas en groupes
r\'eductifs}, Bull. Soc. Math. France \textbf{93} (1965), 369--413.}%
\textsuperscript{\SGARefLink{sga3:intro:section:2}{(2)}}.

\section*{\SGAInlineRefTarget{sga3:intro:section:2}2.}

As in the familiar theory over an algebraically closed field, the subtori of the
group schemes under consideration play a crucial role.  Thus the preliminary study
of tori and, more generally, of ``group schemes of multiplicative type''---both
intrinsically and as multiplicative-type subgroups of a given group---occupies a
fairly large part of this Seminar (Expos\'es~VIII--XII).  Their remarkable rigidity
(in some respects even greater than that of abelian schemes or semisimple group
schemes) makes them highly effective tools for studying certain more general groups.

\section*{\SGAInlineRefTarget{sga3:intro:section:3}3.}

Beginning with Expos\'e~XII (apart from Expos\'e~XVIII, already mentioned), we shall
routinely use the theory of affine algebraic groups over an algebraically closed
field, which the reader will find in the 1956 Chevalley Seminar, especially in
Expos\'es~IV--IX of that Seminar.  We shall also use, though to a lesser extent, the
later Expos\'es of the Chevalley Seminar devoted to the structure of semisimple
algebraic groups.  Indeed, we shall redevelop Chevalley's theory directly within the
framework of schemes; it will be seen that this approach makes the exposition
simpler and more precise, even over a base field.

\section*{\SGAInlineRefTarget{sga3:intro:section:4}4.}

The principal purpose of this Seminar is, of course, to develop techniques that
apply to the study of group schemes over an arbitrary base, that is, essentially,
to the study of families of algebraic groups.  The infinitesimal properties of such
families, and in particular the case of an Artinian base scheme, therefore play an
important role.  These properties arise even in the study of algebraic groups over
a field~$k$ when the field is not perfect, especially when applying descent in the
non-Galois case.  Among the new results obtained in this setting, let us mention the
existence of maximal tori and Cartan subgroups in an arbitrary smooth algebraic
group, the rationality of the variety of maximal tori, and various related results
(Expos\'e~XIV\SGAIIIIntroFootnote{\SGARefLink{sga3:intro:section:3}{(3)}}{\emph{Editor's note:} In particular,
\SGARefLink{sga3:xiv:1.1}{Theorems~1.1} and~\SGARefLink{sga3:xiv:6.1}{6.1}.}); another example is the correspondence between the ``forms''
of a semisimple group and homogeneous principal bundles under a suitable semisimple
algebraic group (generally not connected) (\SGARefLink{sga3:xxiv:1.16}{Expos\'e~XXIV, 1.16}--\SGARefLink{sga3:xxiv:1.20}{1.20}).  In general,
one may say that the methods required for work over a nonperfect base field are
essentially those used over arbitrary base preschemes and thus lie outside the
framework of classical algebraic geometry.

\section*{\SGAInlineRefTarget{sga3:intro:section:5}5.}

It did not seem useful to indicate, at the beginning of the written Expos\'es, the
date or dates of the corresponding oral presentations in the Seminar.  We shall
merely say that the order of the mimeographed Expos\'es (I--XXVI) does indeed agree
with the order of the oral presentations.  Moreover, the final written text was
sometimes prepared considerably later than the oral presentation and often differs
from it quite substantially: the written text is generally more detailed and more
complete (as in Expos\'es~IV and VII\textsubscript{B}), or even appreciably more
general (as in Expos\'es~XII and VII\textsubscript{B}), than the oral presentation.
Other written Expos\'es correspond to no oral presentation
(VI\textsubscript{B}, VII\textsubscript{A}, XV, XVI, XVII, and most of XXVI); they
were written and inserted into the mimeographed Seminar either to provide convenient
references for various other Expos\'es (notably in the case of
VI\textsubscript{B}) or because they form a natural continuation of the concepts
and techniques already developed.  Consequently, Expos\'es~VII\textsubscript{A},
VII\textsubscript{B}, XV, XVI, and XVII need not be read in order to study the rest
of the Seminar, particularly the part devoted to reductive group schemes.

\section*{\SGAInlineRefTarget{sga3:intro:section:6}6.}

From scheme theory we shall use primarily the general language of schemes set out in
EGA~I; the notions of flat, \(\acute{e}\)tale, and smooth morphisms set out in
SGA~1, I--V; and the theory of faithfully flat descent in SGA~1, VIII.  As far as
possible, we have avoided imposing unnecessary Noetherian hypotheses.  In return,
this has required us to replace the usual hypothesis ``of finite type'' with the
hypothesis ``of finite presentation.''  For the notion of a morphism of finite
presentation, the reader should consult EGA~IV, 1.4 and~1.6.  The results of
SGA~1, I--IV, which are most often stated in the Noetherian setting, will be
developed in general in EGA~IV\SGAIIIIntroFootnote{\ensuremath{\ast\ast\ast\ast}}{Since
this Introduction was written, all four parts (Sections~1--21) of EGA~IV have
appeared.}; standard methods for reducing certain kinds of statements involving
finite-presentation hypotheses to the Noetherian case will also be developed there
in detail (EGA~IV, Sections~8, 9, and~11).  A reader who does not wish to assume
these results from EGA~IV may simplify certain statements or their proofs by
assuming the base prescheme to be locally Noetherian.  This does, however, lead to
difficulties when a proof proceeds by descent from~$\widehat A$ to~$A$, where
$\widehat A$ is the completion of a Noetherian local ring~$A$, because that method
introduces the generally non-Noetherian ring
\[
  \widehat A\otimes_A\widehat A.
\]

\section*{\SGAInlineRefTarget{sga3:intro:section:7}7.}

References follow the usual decimal system: the reference 5.7.11 denotes the
proposition (or lemma, definition, and so forth) with that number in the same
Expos\'e; in a reference such as XVII~7.8, the Roman numeral denotes the Expos\'e.
We use the following abbreviations for our standard references:

\begin{description}
\item[\emph{Bible}] Chevalley Seminar, \emph{Groupes de Lie alg\'ebriques},
  1956/58.
\item[EGA $X$, $x.y.z$] J. Dieudonn\'e and A. Grothendieck,
  \emph{\'El\'ements de G\'eom\'etrie Alg\'ebrique}, Chapter~$X$, statement
  $x.y.z$ (or subsection $x.y$, and so forth).
\item[SGA $n$, $X$ $y.z$] \emph{S\'eminaire de G\'eom\'etrie Alg\'ebrique du
  Bois-Marie}, year~$n$, statement $y.z$ of Expos\'e~$X$.
\item[TDTE] A. Grothendieck, \emph{Techniques de descente et Th\'eor\`emes
  d'existence en G\'eom\'etrie Alg\'ebrique}, Expos\'es delivered in the Bourbaki
  Seminar between 1959 and 1962.
\end{description}

\begin{quote}\small
\noindent\textsuperscript{\SGARefLink{sga3:intro:section:1}{(1)}}\emph{Editor's note.}
Concerning quotients by a reductive group, the situation has evolved considerably
since A. Grothendieck wrote Note~$(\ast\ast)$.  For an affine algebraic group over
an arbitrary field~$k$, the appropriate notion, introduced by Mumford, is that of a
``geometrically reductive'' group.  (On the other hand, $G$ is said to be ``linearly
reductive'' if every rational representation of~$G$ is completely reducible; as
noted in Note~$(\ast\ast)$, however, that condition is too restrictive when
$\operatorname{char}(k)>0$.)  By results of M. Nagata and W. J. Haboush
([Na64], [NM64], [Ha75]), the geometrically reductive $k$-groups are precisely the
reductive $k$-groups in the sense of the present Seminar.  For all of this, see the
later editions of Mumford's book ([MF82]).  The notion of geometric reductivity over
an arbitrary base, together with its consequences for passage to a quotient, was
extended to this setting by C. S. Seshadri ([Se77]); additions due to M. Raynaud
appear in the article [CTS79] by J.-L. Colliot-Th\'el\`ene and J.-J. Sansuc.
Finally, for more recent developments on passage to the quotient by an algebraic
group or, more generally, by a groupoid (cf. Expos\'e~V of the present Seminar), we
refer to the articles [Ko97] and [KM97] by J. Koll\'ar and by S. Keel and S. Mori.
\end{quote}

\section*{Bibliography}

\begin{description}
\item[{[CTS79]}] J.-L. Colliot-Th\'el\`ene and J.-J. Sansuc,
  \emph{Fibr\'es quadratiques et composantes connexes r\'eelles}, Math. Ann.
  \textbf{244} (1979), 105--134.
\item[{[Ha75]}] W. J. Haboush, \emph{Reductive groups are geometrically
  reductive}, Ann. of Math. \textbf{102} (1975), no.~1, 67--83.
\item[{[KM97]}] S. Keel and S. Mori, \emph{Quotient by groupoids}, Ann. of Math.
  \textbf{145} (1997), no.~1, 193--213.
\item[{[Ko97]}] J. Koll\'ar, \emph{Quotient spaces modulo algebraic groups}, Ann.
  of Math. \textbf{145} (1997), no.~1, 33--79.
\item[{[MF82]}] D. Mumford and J. Fogarty, \emph{Geometric invariant theory},
  2nd ed., Springer-Verlag, 1982; 3rd ed., with F. Kirwan, 1994.
\item[{[Na64]}] M. Nagata, \emph{Invariants of a group in an affine ring},
  J. Math. Kyoto Univ. \textbf{3} (1964), no.~3, 369--377.
\item[{[NM64]}] M. Nagata and T. Miyata, \emph{Note on semi-reductive groups},
  J. Math. Kyoto Univ. \textbf{3} (1964), no.~3, 379--382.
\item[{[Se77]}] C. S. Seshadri, \emph{Geometric reductivity over an arbitrary
  base}, Adv. Math. \textbf{26} (1977), no.~3, 225--274.
\end{description}

\endgroup

% END INCLUDED SGA3_English_current_buildable_source_20260823/source/inputs/INTRO/tex/components/00_introduction.tex


% END INCLUDED SGA3_English_current_buildable_source_20260823/source/inputs/INTRO/body.tex


\SGAStartExpose{I}
% BEGIN INCLUDED SGA3_English_current_buildable_source_20260823/source/inputs/I/body.tex
% Mechanically extracted after \mainmatter from the frozen standalone master.

% BEGIN INCLUDED SGA3_English_current_buildable_source_20260823/source/inputs/I/tex/components/01_expose_I_opening_pp001_005.tex
% SGA 3, Expose I, Loop 1 English reconstruction.
% Scope: printed pages 1--5 through the end of Section 1.5.
% Authority: Polo--Gille Exp1-13oct24.pdf, component pages 1--5
% (combined-reader pages 11--15). OCR is locator material only.

\setcounter{footnote}{0}
\chapter*{Expos\'e I\\Algebraic Structures. Group Cohomology}
\addcontentsline{toc}{chapter}{Expos\'e I. Algebraic Structures. Group Cohomology}
\markboth{EXPOS\'E I. ALGEBRAIC STRUCTURES. GROUP COHOMOLOGY}%
  {EXPOS\'E I. ALGEBRAIC STRUCTURES. GROUP COHOMOLOGY}

\begin{center}
  \textit{by M. Demazure}
\end{center}

This expos\'e consists of two parts: the first collects a number of general
definitions and establishes notation that will often be used in the sequel;
the second treats group cohomology and culminates in \SGARefLink{sga3:i:theorem:5-3-3}{Theorem~5.3.3}, on the
vanishing of the cohomology of diagonalizable groups.
\footnote{Editor's note: version of October 13, 2024.}

We choose once and for all a Universe.\footnote{Editor's note: see SGA~4,
Expos\'e~I, \S0 and the Appendix; see also the discussion in [DG70], p.~xxvi.}
All definitions and constructions will be relative to this Universe. We shall
systematically allow ourselves the following abuse of language. To define a
functor \(f:\mathscr C\to\mathscr C'\), we shall be content to define the
object \(f(S)\) of \(\mathscr C'\) for every object \(S\) of \(\mathscr C\),
whenever there is no ambiguity about how to define \(f(h)\) for an arrow \(h\)
of \(\mathscr C\). In practice, we shall say: let
\(f:\mathscr C\to\mathscr C'\) be the functor defined by
\(f(S)=\cdots\).

\SGARefTarget{sga3:i:section:1}
\section*{1. Generalities}
\addcontentsline{toc}{section}{1. Generalities}

\subsection*{\SGAInlineRefTarget{sga3:i:subsection:1-1}1.1.}

Let \(\mathscr C\) be a category. We denote by \(\widehat{\mathscr C}\) the
category
\[
  \operatorname{Hom}\bigl(\mathscr C^{\circ},\mathbf{Set}\bigr)
\]
of contravariant functors from \(\mathscr C\) to the category \(\mathbf{Set}\)
of sets.\footnote{Editor's note: this is called the category of presheaves on
\(\mathscr C\); see IV, \SGARefLink{sga3:iv:definition:4:3:1}{4.3.1}.} There is a canonical functor
\[
  \mathbf h:\mathscr C\longrightarrow\widehat{\mathscr C}
\]
which associates with every \(X\in\operatorname{Ob}(\mathscr C)\) the functor
\(\mathbf h_X\) such that
\[
  \mathbf h_X(S)=\operatorname{Hom}(S,X).
\]
For every functor \(\mathbf F\in\operatorname{Ob}(\widehat{\mathscr C})\),
one has (see, for example, EGA~\(0_{\mathrm{III}}\), 8.1.4) a bijection
\[
  \operatorname{Hom}(\mathbf h_X,\mathbf F)
    \xrightarrow{\ \sim\ } \mathbf F(X).
\]
\footnote{Editor's note: this result is often called the Yoneda lemma; we use
that terminology in other editor's notes.}

In particular, for every pair \(X,Y\) of objects of \(\mathscr C\), the
canonical map
\[
  \operatorname{Hom}_{\mathscr C}(X,Y)
    \xrightarrow{\ \sim\ }
  \operatorname{Hom}_{\widehat{\mathscr C}}
    (\mathbf h_X,\mathbf h_Y)
\]
is bijective; that is, the functor \(\mathbf h\) is fully faithful. It
therefore defines an isomorphism of \(\mathscr C\) onto a full subcategory of
\(\widehat{\mathscr C}\), and an equivalence of \(\mathscr C\) with the full
subcategory of \(\widehat{\mathscr C}\) formed by the representable functors
(that is, the functors isomorphic to one of the form \(\mathbf h_X\)). In the
sequel we shall often identify \(X\) with \(\mathbf h_X\). The following
paragraphs are intended to show that this identification can be made safely.

\medskip
\noindent\SGAInlineRefTarget{sga3:i:remark:1-1-1}\textbf{Remark 1.1.1.}\footnote{Editor's note: this remark has been
added.}
We sometimes need the following variant. Let \(\mathscr D\) be a full
subcategory of \(\mathscr C\), and let
\(X,Y\in\operatorname{Ob}(\mathscr D)\). Denote by \(\mathbf h'_X\) and
\(\mathbf h'_Y\) the restrictions of \(\mathbf h_X\) and \(\mathbf h_Y\) to
\(\mathscr D\). Then
\[
  \operatorname{Hom}_{\mathscr C}(X,Y)
  =\operatorname{Hom}_{\mathscr D}(X,Y)
  =\operatorname{Hom}_{\widehat{\mathscr D}}
      (\mathbf h'_X,\mathbf h'_Y).
\]
Thus, to give a morphism \(X\to Y\) ``is the same thing'' as to give, for
every \(T\in\operatorname{Ob}(\mathscr D)\), a map
\[
  \phi(T):\operatorname{Hom}(T,X)\longrightarrow\operatorname{Hom}(T,Y)
\]
functorially in \(T\).

\subsection*{\SGAInlineRefTarget{sga3:i:subsection:1-2}1.2.}

We shall say that \(\mathbf F\) is a \emph{subobject} (or a subfunctor) of
\(\mathbf G\) if \(\mathbf F(S)\) is a subset of \(\mathbf G(S)\) for every
\(S\).\footnote{Editor's note: the order has been changed so that fiber
products are introduced before monomorphisms; see editor's note (9).}

In \(\widehat{\mathscr C}\), arbitrary inverse limits exist and are computed
pointwise:
\[
  \left(\varprojlim_i\mathbf F_i\right)(S)
    =\varprojlim_i\mathbf F_i(S).
\]
\footnote{Editor's note: similarly, arbitrary direct limits exist and are
computed pointwise:
\(\bigl(\varinjlim_i\mathbf F_i\bigr)(S)=
\varinjlim_i\mathbf F_i(S)\). In general, however, the functor \(\mathbf h\)
does not commute with direct limits.}
In particular, fiber products are defined by
\[
  (\mathbf F\mathbin{\times_{\mathbf G}}\mathbf F')(S)
    =\mathbf F(S)\mathbin{\times_{\mathbf G(S)}}\mathbf F'(S).
\]
\footnote{Editor's note: in particular, the equalizer of a pair of morphisms
\(u,v:\mathbf F\rightrightarrows\mathbf G\) is the subfunctor
\(\operatorname{Ker}(u,v)\) of \(\mathbf F\) defined by
\(\operatorname{Ker}(u,v)(S)=
\{x\in\mathbf F(S)\mid u(x)=v(x)\}\).}

As a final object of \(\widehat{\mathscr C}\), we choose the functor
\(\underline{\mathbf e}\) such that
\(\underline{\mathbf e}(S)=\{\varnothing\}\).\footnote{Editor's note:
\(\{\varnothing\}\), the power set of the empty set, is a one-element set.}
Every \(\mathbf F\in\operatorname{Ob}(\widehat{\mathscr C})\) has a unique
morphism to \(\underline{\mathbf e}\), and we set
\[
  \mathbf F\times\mathbf F'
    =\mathbf F\mathbin{\times_{\underline{\mathbf e}}}\mathbf F'.
\]

The functor \(\mathbf h\) commutes with inverse limits. In particular, for
\(X\times X'\) to exist, where
\(X,X'\in\operatorname{Ob}(\mathscr C)\), respectively for \(\mathscr C\) to
admit a final object \(e\), it is necessary and sufficient that
\(\mathbf h_X\times\mathbf h_{X'}\), respectively
\(\underline{\mathbf e}\), be representable; in that case
\[
  \mathbf h_X\times\mathbf h_{X'}\simeq\mathbf h_{X\times X'},
  \qquad
  \mathbf h_e\simeq\underline{\mathbf e}.
\]

A monomorphism in \(\widehat{\mathscr C}\) is precisely a morphism
\(\mathbf F\to\mathbf G\) such that, for every
\(S\in\operatorname{Ob}(\mathscr C)\), the corresponding map of sets
\(\mathbf F(S)\to\mathbf G(S)\) is injective.\footnote{Editor's note: if
\(\mathbf F(S)\to\mathbf G(S)\) is injective for every \(S\), it is clear
that \(\mathbf F\to\mathbf G\) is a monomorphism. The converse follows by
considering
\(\mathbf F\times_{\mathbf G}\mathbf F\rightrightarrows\mathbf F\to
\mathbf G\). Thus \(\mathbf F\to\mathbf G\) is a monomorphism if and only if
the diagonal morphism
\(\mathbf F\to\mathbf F\times_{\mathbf G}\mathbf F\) is an isomorphism
(see EGA~I, 5.3.8). Similarly, a pointwise surjective morphism is an
epimorphism, and the converse follows by considering the pushout
\(\mathbf G\amalg_{\mathbf F}\mathbf G\); see the proof of \SGARefLink{sga3:i:paragraph:4-4-4}{Lemma~4.4.4} in
Expos\'e~IV.}

\medskip
\noindent\textit{The functor \(\Gamma\).}
For every \(\mathbf F\in\operatorname{Ob}(\widehat{\mathscr C})\), set
\[
  \Gamma(\mathbf F)=\operatorname{Hom}
    (\underline{\mathbf e},\mathbf F).
\]
An element of \(\Gamma(\mathbf F)\) is therefore a family
\((\gamma_S)_{S\in\operatorname{Ob}(\mathscr C)}\), with
\(\gamma_S\in\mathbf F(S)\), such that for every arrow
\(f:S'\to S''\) of \(\mathscr C\),
\[
  \mathbf F(f)(\gamma_{S''})=\gamma_{S'}.
\]
For \(X\in\operatorname{Ob}(\mathscr C)\), set
\(\Gamma(X)=\Gamma(\mathbf h_X)\). If \(\mathscr C\) has a final object
\(e\), then
\[
  \Gamma(X)\simeq\operatorname{Hom}(e,X).
\]

\subsection*{\SGAInlineRefTarget{sga3:i:subsection:1-3}1.3.}

Let \(S\in\operatorname{Ob}(\mathscr C)\). We denote by
\(\mathscr C_{/S}\) the category of objects of \(\mathscr C\) over \(S\):
its objects are the arrows \(f:T\to S\) of \(\mathscr C\), and
\(\operatorname{Hom}(f,f')\) is the subset of
\(\operatorname{Hom}(T,T')\) formed by the arrows \(u\) such that
\(f=f'\circ u\). If \(\mathscr C\) has a final object \(e\), then
\(\mathscr C_{/e}\) is isomorphic to \(\mathscr C\). The category
\(\mathscr C_{/S}\) has a final object, namely the identity arrow
\(S\to S\).

If \(f:T\to S\) is an object of \(\mathscr C_{/S}\), one may form the
category \((\mathscr C_{/S})_{/f}\), which by abuse of language we denote by
\((\mathscr C_{/S})_{/T}\), and there is a canonical isomorphism
\[
  \mathscr C_{/T}\simeq(\mathscr C_{/S})_{/T}.
\]
This construction also applies to \(\widehat{\mathscr C}\); in particular,
one obtains the category \(\widehat{\mathscr C}_{/\mathbf h_S}\). On the
other hand, one may form the category \(\widehat{\mathscr C_{/S}}\).

If \(f:T\to S\) is an object of \(\mathscr C_{/S}\), then \(\Gamma(f)\)
is identified with the set \(\Gamma(T/S)\) of sections of \(T\) over \(S\),
that is, arrows \(S\to T\) that are right inverses to \(f\). Moreover,
\(\mathbf h_f:\mathbf h_T\to\mathbf h_S\) is then an object of
\(\widehat{\mathscr C}_{/\mathbf h_S}\), and
\[
  \Gamma(\mathbf h_f)
    \simeq\Gamma(\mathbf h_T/\mathbf h_S)
    \simeq\Gamma(T/S)
    \simeq\Gamma(f).
\]

\subsection*{\SGAInlineRefTarget{sga3:i:subsection:1-4}1.4.}

We now define an equivalence between the categories
\(\widehat{\mathscr C_{/S}}\) and
\(\widehat{\mathscr C}_{/\mathbf h_S}\). In other words, we shall prove that
to give a functor on the category of objects of \(\mathscr C\) over \(S\) is
``the same thing'' as to give a functor on \(\mathscr C\) equipped with a
morphism to \(\mathbf h_S\).

\medskip
\noindent\textbf{(i) Construction of
\(\alpha_S:\widehat{\mathscr C}_{/\mathbf h_S}
\to\widehat{\mathscr C_{/S}}\).}

First let \(H:\mathbf F\to\mathbf h_S\) be an object of
\(\widehat{\mathscr C}_{/\mathbf h_S}\). We must define a functor
\(\alpha_S(H)\) on \(\mathscr C_{/S}\). For an object
\(f:T\to S\) of \(\mathscr C_{/S}\), define \(\alpha_S(H)(f)\) to be the
inverse image of \(f\in\mathbf h_S(T)\) under the map
\[
  H(T):\mathbf F(T)\longrightarrow\mathbf h_S(T).
\]
\footnote{Editor's note: for example, if \(\mathbf F=\mathbf h_X\), then
\(H\) corresponds to a morphism \(h:X\to S\), and
\(H(T):\operatorname{Hom}_{\mathscr C}(T,X)\to
\operatorname{Hom}_{\mathscr C}(T,S)\) is the map \(g\mapsto h\circ g\).
Thus \(\alpha_S(\mathbf h_X)=\operatorname{Hom}_{\mathscr C_{/S}}(-,X)\).}

Next, let \(u:f\to f'\) be an arrow of \(\mathscr C_{/S}\). Then
\(\mathbf F(u):\mathbf F(T')\to\mathbf F(T)\) induces a map from
\(\alpha_S(H)(f')\) to \(\alpha_S(H)(f)\), denoted \(\alpha_S(H)(u)\). It
is immediate that the assignments
\[
  f\longmapsto\alpha_S(H)(f),
  \qquad
  u\longmapsto\alpha_S(H)(u)
\]
define a functor on \(\mathscr C_{/S}\), hence an object \(\alpha_S(H)\) of
\(\widehat{\mathscr C_{/S}}\).

Finally, let \(H:\mathbf F\to\mathbf h_S\) and
\(H':\mathbf F'\to\mathbf h_S\) be two objects of
\(\widehat{\mathscr C}_{/\mathbf h_S}\), and let \(U:H\to H'\) be a
morphism in that category:

% Authority: Exp1-13oct24.pdf p. 4; combined reader p. 14.
\begin{center}
  % BEGIN INCLUDED SGA3_English_current_buildable_source_20260823/source/inputs/I/native_diagrams/exp1/diag_001_alpha_morphism_triangle.tex
\begin{tikzcd}[column sep=5.2em,row sep=3.2em]
  \mathbf F
    \arrow[rr,"U"]
    \arrow[dr,"H"']
  &&
  \mathbf F'
    \arrow[dl,"H'"]
  \\
  &
  \mathbf h_S
  &
  {}.
\end{tikzcd}

% END INCLUDED SGA3_English_current_buildable_source_20260823/source/inputs/I/native_diagrams/exp1/diag_001_alpha_morphism_triangle.tex

\end{center}

For every \(f:T\to S\), the map
\(U(T):\mathbf F(T)\to\mathbf F'(T)\) induces a map
\[
  \alpha_S(U)(f):\alpha_S(H)(f)
    \longrightarrow\alpha_S(H')(f),
\]
which defines a morphism of functors
\[
  \alpha_S(U):\alpha_S(H)\longrightarrow\alpha_S(H').
\]
It is easy to verify that the assignments
\[
  H\longmapsto\alpha_S(H),
  \qquad
  U\longmapsto\alpha_S(U)
\]
define a functor
\[
  \alpha_S:\widehat{\mathscr C}_{/\mathbf h_S}
    \longrightarrow\widehat{\mathscr C_{/S}}.
\]

\medskip
\SGARefTarget{sga3:i:proposition:1:4:1}
\noindent\textbf{(ii) \SGARefLink{sga3:i:proposition:1:4:1}{Proposition 1.4.1}.}
\textit{The functor
\(\alpha_S:\widehat{\mathscr C}_{/\mathbf h_S}
\to\widehat{\mathscr C_{/S}}\) is an equivalence of categories.}

We indicate only the principle of the construction of a quasi-inverse functor
\[
  \beta_S:\widehat{\mathscr C_{/S}}
    \longrightarrow\widehat{\mathscr C}_{/\mathbf h_S}.
\]
Let \(\mathbf G\) be a functor on \(\mathscr C_{/S}\). For every object
\(T\) of \(\mathscr C\), set
\[
  \beta_S(\mathbf G)(T)
    =\coprod_{f\in\operatorname{Hom}(T,S)}\mathbf G(f),
  \qquad \operatorname{Hom}(T,S)=\mathbf h_S(T).
\]
This defines a functor \(\beta_S(\mathbf G)\) on \(\mathscr C\), equipped
with the evident projection to \(\mathbf h_S\).

\subsection*{\SGAInlineRefTarget{sga3:i:subsection:1-5}1.5.}

The equivalence \(\alpha_S\) commutes with the functors \(\Gamma\). In other
words, if \(H:\mathbf F\to\mathbf h_S\) is an object of
\(\widehat{\mathscr C}_{/\mathbf h_S}\) and \(\alpha_S(H)\) is the
corresponding object of \(\widehat{\mathscr C_{/S}}\), then
\[
  \Gamma(\alpha_S(H))
    \simeq\Gamma(H)
    \simeq\Gamma(\mathbf F/\mathbf h_S).
\]
The equivalence \(\alpha_S\) also commutes with the functors \(\mathbf h\).
That is, if \(f:T\to S\) is an object of \(\mathscr C_{/S}\), then
\(\mathbf h_f:\mathbf h_T\to\mathbf h_S\) is an object of
\(\widehat{\mathscr C}_{/\mathbf h_S}\), and its image under \(\alpha_S\)
is precisely \(\mathbf h_{\mathscr C_{/S}}(f)\), where
\[
  \mathbf h_{\mathscr C_{/S}}:
    \mathscr C_{/S}\longrightarrow\widehat{\mathscr C_{/S}}
\]
is the canonical functor.\footnote{Editor's note: see editor's note (10).}
Consequently:

\medskip
\noindent\SGAInlineRefTarget{sga3:i:proposition:1-5-1}\textbf{Proposition 1.5.1.}
\textit{Let \(H:\mathbf F\to\mathbf h_S\) be an object of
\(\widehat{\mathscr C}_{/\mathbf h_S}\). The functor
\[
  \alpha_S(H):(\mathscr C_{/S})^{\circ}\longrightarrow\mathbf{Set}
\]
is representable if and only if
\(\mathbf F:\mathscr C^{\circ}\to\mathbf{Set}\) is representable. If
\(\mathbf F\simeq\mathbf h_T\), then \(\alpha_S(H)\) is represented by the
object \(T\to S\) of \(\mathscr C_{/S}\).}

The equivalence \(\alpha_S\) is transitive in \(S\). If \(f:T\to S\) is an
object of \(\mathscr C_{/S}\), there is a commutative diagram of
equivalences:

% Authority: Exp1-13oct24.pdf p. 5; combined reader p. 15.
\begin{center}
  % BEGIN INCLUDED SGA3_English_current_buildable_source_20260823/source/inputs/I/native_diagrams/exp1/diag_002_alpha_transitivity.tex
\begin{tikzcd}[column sep=4.0em,row sep=3.0em]
  \bigl(\widehat{\mathscr C}_{/\mathbf h_S}\bigr)_{/\mathbf h_T}
    \arrow[r,"\alpha_{S/\mathbf h_T}"]
    \arrow[dr,"\simeq"']
  &
  \bigl(\widehat{\mathscr C_{/S}}\bigr)_{/\mathbf h_T}
    \arrow[r,"\alpha_f"]
  &
  \widehat{(\mathscr C_{/S})_{/T}}
    \arrow[dl,"\simeq"]
  \\
  &
  \widehat{\mathscr C}_{/\mathbf h_T}
    \arrow[r,"\alpha_T"]
  &
  \widehat{\mathscr C_{/T}}
  \qquad ,
\end{tikzcd}

% END INCLUDED SGA3_English_current_buildable_source_20260823/source/inputs/I/native_diagrams/exp1/diag_002_alpha_transitivity.tex

\end{center}

Here \(\alpha_{S/\mathbf h_T}\) denotes, provisionally, the restriction
(see \SGARefLink{sga3:i:subsection:1-6}{1.6}) of the functor \(\alpha_S\) to the objects over
\(\mathbf h_T\).

% End of assigned opening component. Section 1.6 begins in the next unit.

% END INCLUDED SGA3_English_current_buildable_source_20260823/source/inputs/I/tex/components/01_expose_I_opening_pp001_005.tex

% BEGIN INCLUDED SGA3_English_current_buildable_source_20260823/source/inputs/I/tex/components/02_expose_I_s16_next_pp005_010.tex
% SGA 3, Expose I, Loop 1 English reconstruction.
% Scope: Section 1.6 through the end of Remark 1.7.3.
% Authority: Polo--Gille Exp1-13oct24.pdf, component pages 5--10
% (combined-reader pages 15--20). OCR is drafting/locator material only.

\subsection*{\SGAInlineRefTarget{sga3:i:subsection:1-6}1.6. Base change in a functor}

For every \(S\in\operatorname{Ob}(\mathscr C)\), there is a canonical functor
\[
  i_S:\mathscr C_{/S}\longrightarrow\mathscr C
\]
defined by \(i_S(f)=T\) when \(f\) is the arrow \(T\to S\). If
\(f:T\to S\) is an object of \(\mathscr C_{/S}\), we denote by
\(i_{T/S}=i_f\) the functor
\[
  i_{T/S}:(\mathscr C_{/S})_{/T}\longrightarrow\mathscr C_{/S},
\]
and there is a commutative diagram:

% Authority: Exp1-13oct24.pdf p. 5; combined reader p. 15.
\begin{center}
  % BEGIN INCLUDED SGA3_English_current_buildable_source_20260823/source/inputs/I/native_diagrams/exp1/diag_003_base_change_triangle.tex
\begin{tikzcd}[column sep=4.7em,row sep=3.2em]
  (\mathscr C_{/S})_{/T}
    \arrow[r,"i_{T/S}"]
    \arrow[dr,"\simeq"']
  &
  \mathscr C_{/S}
    \arrow[r,"i_S"]
  &
  \mathscr C
  \\
  &
  \mathscr C_{/T}
    \arrow[ur,"i_T"']
  &
  {} ,
\end{tikzcd}

% END INCLUDED SGA3_English_current_buildable_source_20260823/source/inputs/I/native_diagrams/exp1/diag_003_base_change_triangle.tex

\end{center}

Thus, identifying \((\mathscr C_{/S})_{/T}\) with \(\mathscr C_{/T}\), as
we shall henceforth do,
\[
  i_S\circ i_{T/S}=i_T.
\]
Likewise, if \(\mathscr C\) is identified with \(\mathscr C_{/e}\) when
\(\mathscr C\) has a final object \(e\), then
\(i_{S/e}:\mathscr C_{/S}\to\mathscr C_{/e}\) is identified with \(i_S\).

For \(X\in\operatorname{Ob}(\mathscr C)\) (respectively,
\(Y\in\operatorname{Ob}(\mathscr C_{/S})\)), let \(p_S(X)\)
(respectively, \(p_{T/S}(Y)\)) be the object of \(\mathscr C_{/S}\)
(respectively, of \(\mathscr C_{/T}\)), when it exists, defined by
\(X\times S\) (respectively, \(Y\times_S T\)) equipped with its second
projection:

% Authority: Exp1-13oct24.pdf p. 5; combined reader p. 15.
\begin{center}
  % BEGIN INCLUDED SGA3_English_current_buildable_source_20260823/source/inputs/I/native_diagrams/exp1/diag_004_base_change_projections.tex
\begin{tikzcd}[column sep=6.5em,row sep=3.2em]
  X\times S
    \arrow[d,"p_S(X)"']
  &
  \textup{resp.}
  &
  Y\times_S T
    \arrow[d,"p_{T/S}(Y)"]
  \\
  S
  &&
  T
  \qquad .
\end{tikzcd}

% END INCLUDED SGA3_English_current_buildable_source_20260823/source/inputs/I/native_diagrams/exp1/diag_004_base_change_projections.tex

\end{center}

The partially defined functor \(p_S\) (respectively, \(p_{T/S}\)) is called
the \emph{base-change functor}. By the definition of a product
(respectively, a fiber product), it is right adjoint to \(i_S\)
(respectively, to \(i_{T/S}\)).
\footnote{Editor's note: that is, if \(g:U\to S\) (respectively,
\(h:V\to T\)) is an object of \(\mathscr C_{/S}\) (respectively, of
\(\mathscr C_{/T}\)), then \(i_S(g)=U\), while \(i_{T/S}(h)\) is the
object \(f\circ h:V\to S\) of \(\mathscr C_{/S}\), and
\[
  \operatorname{Hom}_{\mathscr C_{/S}}(U,X\times S)
    \simeq \operatorname{Hom}_{\mathscr C}(U,X),
\]
respectively,
\[
  \operatorname{Hom}_{\mathscr C_{/T}}(V,X\times_S T)
    \simeq \operatorname{Hom}_{\mathscr C_{/S}}(V,X).
\]}
We also write
\[
  p_S(X)=X_S,
  \qquad
  p_{T/S}(Y)=Y_T.
\]

The functor \(i_S\) defines a restriction functor
\[
  i_S^*:\widehat{\mathscr C}\longrightarrow
    \widehat{\mathscr C_{/S}}.
\]
We write
\[
  \mathbf F_S=i_S^*(\mathbf F)=\mathbf F\circ i_S.
\]
Plainly,
\[
  i_{T/S}^*\circ i_S^*=i_T^*,
\]
that is, for every functor
\(\mathbf F\in\operatorname{Ob}(\widehat{\mathscr C})\),
\[
  (\mathbf F_S)_T=\mathbf F_T.
\]
The notation is justified by the following result.

\medskip
\noindent\SGAInlineRefTarget{sga3:i:proposition:1-6-1}\textbf{Proposition 1.6.1.}
\textit{The functor
\[
  (\mathbf h_X)_S:(\mathscr C_{/S})^{\circ}\longrightarrow\mathbf{Set}
\]
is representable if and only if the product \(X\times S\) exists. In that
case
\[
  (\mathbf h_X)_S\simeq\mathbf h_{X_S}.
\]}

This shows that \(\mathbf F_S\) has two interpretations: it is the
restriction of \(\mathbf F\) to \(\mathscr C_{/S}\), and it is the functor
obtained by the base change \(e\leftarrow S\). This leads to the following
notation:

% Authority: Exp1-13oct24.pdf p. 6; combined reader p. 16.
\begin{center}
  % BEGIN INCLUDED SGA3_English_current_buildable_source_20260823/source/inputs/I/native_diagrams/exp1/diag_005_functor_base_change_notation.tex
\begin{tikzcd}[column sep=3.8em,row sep=3.0em]
  \mathbf F
    \arrow[r,no head]
    \arrow[d,no head]
  &
  \mathbf F_S
    \arrow[r,no head]
    \arrow[d,no head]
  &
  \mathbf F_T
    \arrow[d,no head]
  \\
  e
  &
  S
    \arrow[l]
  &
  T
    \arrow[l]
\end{tikzcd}

% END INCLUDED SGA3_English_current_buildable_source_20260823/source/inputs/I/native_diagrams/exp1/diag_005_functor_base_change_notation.tex

\end{center}

This diagram reflects both of the preceding interpretations. Notice that
\[
  \Gamma(\mathbf F_S)
    \simeq\operatorname{Hom}(\mathbf h_S,\mathbf F)
    \simeq\mathbf F(S),
\]
and, in particular,
\[
  \Gamma(X_S)\simeq\operatorname{Hom}(S,X).
\]

\subsection*{\SGAInlineRefTarget{sga3:i:subsection:1-7-0}1.7.0.}

Let \(\mathbf E\) be an object of \(\widehat{\mathscr C}\).
\footnote{Editor's note: the lemma below has been added (see SGA~4,
I.3.4); it is used in the proof of \SGARefLink{sga3:i:proposition:1-7-1}{1.7.1} and will be useful repeatedly in
the sequel.}
Consider the category \(\mathscr C_{/\mathbf E}\) of \emph{objects of
\(\mathscr C\)} over \(\mathbf E\). Its objects are pairs \((V,\rho)\)
formed by an object \(V\) of \(\mathscr C\) and a morphism
\(\rho:\mathbf h_V\to\mathbf E\) in \(\widehat{\mathscr C}\), that is,
\(\rho\in\mathbf E(V)\). A morphism from \((V,\rho)\) to
\((V',\rho')\) is a morphism \(f:V\to V'\) in \(\mathscr C\) such that
\[
  \rho'\circ f=\rho,
  \qquad\text{that is,}\qquad
  \mathbf E(f)(\rho')=\rho.
\]
Let \(\mathbf L\) be the functor
\[
  \mathbf L=\varinjlim_{(V,\rho)\in\mathscr C_{/\mathbf E}}\mathbf h_V.
\]
Thus, for every \(S\in\operatorname{Ob}(\mathscr C)\),
\[
  \mathbf L(S)
    =\varinjlim_{(V,\rho)\in\mathscr C_{/\mathbf E}}\mathbf h_V(S)
\]
is the set of equivalence classes of triples \((V,\rho,v)\), where
\(v:S\to V\) is a morphism in \(\mathscr C\), and where
\((V,\rho,v)\) is identified with \((V',\rho',f\circ v)\) for every
morphism \(f:V\to V'\) in \(\mathscr C\) such that
\(\rho'\circ f=\rho\).

The map \(\phi_{\mathbf E}(S)\) which associates with the class of
\((V,\rho,v)\) the element \(\rho\circ v\) of \(\mathbf E(S)\) is well
defined and gives a morphism of functors
\[
  \phi_{\mathbf E}:
  \varinjlim_{(V,\rho)\in\mathscr C_{/\mathbf E}}\mathbf h_V
    \longrightarrow\mathbf E.
\]

\medskip
\noindent\textbf{Lemma.}
\textit{The morphism \(\phi_{\mathbf E}\) is an isomorphism.}

Indeed, let \(S\in\operatorname{Ob}(\mathscr C)\). Every
\(x\in\mathbf E(S)\) is the image under \(\phi_{\mathbf E}(S)\) of the
triple \((S,x,\operatorname{id}_S)\); hence \(\phi_{\mathbf E}(S)\) is
surjective. Conversely, let
\(\ell_1=(V_1,\rho_1,v_1)\) and
\(\ell_2=(V_2,\rho_2,v_2)\) be two elements of \(\mathbf L(S)\) having
the same image in \(\mathbf E(S)\), and set
\[
  \gamma=\rho_1\circ v_1=\rho_2\circ v_2.
\]
Then \(\ell_1\) and \(\ell_2\) are both equal in \(\mathbf L(S)\) to the
class of \((S,\gamma,\operatorname{id}_S)\). Thus
\(\phi_{\mathbf E}(S)\) is injective.

\medskip
\noindent\textbf{Corollary.}
\textit{For every object \(\mathbf F\) of \(\widehat{\mathscr C}\),
\[
  \operatorname{Hom}(\mathbf E,\mathbf F)
    =\varprojlim_{(V,\rho)\in\mathscr C_{/\mathbf E}}\mathbf F(V).
\]}

\SGARefTarget{sga3:i:subsection:1:7}
\subsection*{1.7. The objects \(\underline{\operatorname{Hom}}\),
\(\underline{\operatorname{Isom}}\), etc.}

Let \(\mathbf F\) and \(\mathbf G\) be two objects of
\(\widehat{\mathscr C}\). We define another object of
\(\widehat{\mathscr C}\) by
\begin{align*}
  \underline{\operatorname{Hom}}(\mathbf F,\mathbf G)(S)
  &={}\operatorname{Hom}_{\widehat{\mathscr C_{/S}}}
       (\mathbf F_S,\mathbf G_S)\\
  &\simeq\operatorname{Hom}_{\widehat{\mathscr C}_{/\mathbf h_S}}
       (\mathbf F\times\mathbf h_S,\mathbf G\times\mathbf h_S)\\
  &\simeq\operatorname{Hom}_{\widehat{\mathscr C}}
       (\mathbf F\times\mathbf h_S,\mathbf G).
\end{align*}
The object \(\underline{\operatorname{Hom}}(\mathbf F,\mathbf G)\) has
the following properties:
\begin{enumerate}
  \item[(i)]
    \(\underline{\operatorname{Hom}}(\underline{\mathbf e},\mathbf G)
      \simeq\mathbf G\).
    \footnote{Editor's note: if \(\mathbf E\) is a third object of
    \(\widehat{\mathscr C}\), then
    \[
      \underline{\operatorname{Hom}}(\mathbf E,\mathbf F\times\mathbf G)
      \simeq
      \underline{\operatorname{Hom}}(\mathbf E,\mathbf F)
        \times
      \underline{\operatorname{Hom}}(\mathbf E,\mathbf G).
    \]}
  \item[\SGAInlineRefTarget{sga3:i:subsection:1:7:item:ii}(ii)] The formation of \(\underline{\operatorname{Hom}}\) commutes
    with base extension:
    \[
      \underline{\operatorname{Hom}}(\mathbf F_S,\mathbf G_S)
        \simeq
      \underline{\operatorname{Hom}}(\mathbf F,\mathbf G)_S.
    \]
  \item[(iii)]
    \((\mathbf F,\mathbf G)\mapsto
      \underline{\operatorname{Hom}}(\mathbf F,\mathbf G)\)
    is a bifunctor, contravariant in \(\mathbf F\) and covariant in
    \(\mathbf G\).
\end{enumerate}
These three properties follow immediately from the definitions.

We shall prove that, for every object \(\mathbf E\) of
\(\widehat{\mathscr C}\),
\[
  \operatorname{Hom}
    \bigl(\mathbf E,
      \underline{\operatorname{Hom}}(\mathbf F,\mathbf G)\bigr)
    \simeq
  \operatorname{Hom}(\mathbf E\times\mathbf F,\mathbf G).
\]
Let \(\phi:\mathbf E\times\mathbf F\to\mathbf G\). We must associate
with it a morphism from \(\mathbf E\) to
\(\underline{\operatorname{Hom}}(\mathbf F,\mathbf G)\). For an arrow
\(S'\to S\) of \(\mathscr C\), there are maps
\[
  \mathbf E(S)\times\mathbf F(S')
    \longrightarrow
  \mathbf E(S')\times\mathbf F(S')
    \xrightarrow{\ \phi(S')\ }
  \mathbf G(S').
\]
Every element \(e\) of \(\mathbf E(S)\) therefore defines, for each
\(S'\to S\), a map \(\mathbf F(S')\to\mathbf G(S')\) functorial in
\(S'\); this is an element \(\theta_\phi(e)\) of
\(\underline{\operatorname{Hom}}(\mathbf F,\mathbf G)(S)\). We have thus
obtained a map
\[
  \phi\longmapsto\theta_\phi,
  \qquad
  \operatorname{Hom}(\mathbf E\times\mathbf F,\mathbf G)
    \longrightarrow
  \operatorname{Hom}
    \bigl(\mathbf E,
      \underline{\operatorname{Hom}}(\mathbf F,\mathbf G)\bigr),
\]
which is functorial in \(\mathbf E\).

\medskip
\noindent\SGAInlineRefTarget{sga3:i:proposition:1-7-1}\textbf{Proposition 1.7.1.}
\footnote{Editor's note: part \SGARefLink{sga3:i:proposition:1-7-1:item:b}{(b)}, which will be useful in II.1 and
II.3.11, has been added. A second, more direct proof of part \SGARefLink{sga3:i:proposition:1-7-1:item:a}{(a)} has also
been sketched.}
\textit{Let
\(\mathbf E,\mathbf F,\mathbf G\in
\operatorname{Ob}(\widehat{\mathscr C})\).
\begin{enumerate}
  \item[\SGAInlineRefTarget{sga3:i:proposition:1-7-1:item:a}(a)] The map \(\phi\mapsto\theta_\phi\) is a bijection:
    \[
      \operatorname{Hom}(\mathbf E\times\mathbf F,\mathbf G)
        \xrightarrow{\ \sim\ }
      \operatorname{Hom}
        \bigl(\mathbf E,
          \underline{\operatorname{Hom}}(\mathbf F,\mathbf G)\bigr).
    \]
  \item[\SGAInlineRefTarget{sga3:i:proposition:1-7-1:item:b}(b)] Moreover, there is an isomorphism of functors
    \[
      \underline{\operatorname{Hom}}
        \bigl(\mathbf E,
          \underline{\operatorname{Hom}}(\mathbf F,\mathbf G)\bigr)
        \simeq
      \underline{\operatorname{Hom}}
        (\mathbf E\times\mathbf F,\mathbf G).
    \]
\end{enumerate}}

For \SGARefLink{sga3:i:proposition:1-7-1:item:a}{(a)}, regard both sides as functors of \(\mathbf E\). The asserted
result is true when \(\mathbf E=\mathbf h_X\), since it is then precisely
the definition of the functor
\(\underline{\operatorname{Hom}}(\mathbf F,\mathbf G)\). Moreover, both
sides, as functors of \(\mathbf E\), take direct limits to inverse limits.
Finally, by \SGARefLink{sga3:i:subsection:1-7-0}{Lemma~1.7.0}, every object \(\mathbf E\) of
\(\widehat{\mathscr C}\) is isomorphic to the direct limit of the
\(\mathbf h_X\), as \(X\) ranges over \(\mathscr C_{/\mathbf E}\). This
proves \SGARefLink{sga3:i:proposition:1-7-1:item:a}{(a)}.

Let us also sketch a direct proof of \SGARefLink{sga3:i:proposition:1-7-1:item:a}{(a)}. To every
\[
  \theta\in\operatorname{Hom}
    \bigl(\mathbf E,
      \underline{\operatorname{Hom}}(\mathbf F,\mathbf G)\bigr)
\]
associate \(\phi_\theta\in
\operatorname{Hom}(\mathbf E\times\mathbf F,\mathbf G)\) as follows. For
every \(S\in\operatorname{Ob}(\mathscr C)\), there is a map
\[
  \theta(S):\mathbf E(S)\longrightarrow
    \operatorname{Hom}(\mathbf F\times S,\mathbf G),
\]
functorial in \(S\). If
\((e,f)\in\mathbf E(S)\times\mathbf F(S)\), then \(f\) is a morphism
\(S\to\mathbf F\), so \(f\times\operatorname{id}_S\) is a morphism
\(S\to\mathbf F\times S\). On the other hand, \(\theta(S)(e)\) is a
morphism \(\mathbf F\times S\to\mathbf G\). Composition therefore gives
a morphism
\[
  \theta(S)(e)\circ(f\times\operatorname{id}_S):S\longrightarrow\mathbf G,
\]
that is, an element \(\phi_\theta(S)(e,f)\) of \(\mathbf G(S)\). It is
easy to check that \(S\mapsto\phi_\theta(S)\) is functorial in \(S\), and
therefore defines a morphism \(\phi_\theta\) from
\(\mathbf E\times\mathbf F\) to \(\mathbf G\). We leave to the reader the
verification that \(\theta\mapsto\phi_\theta\) and
\(\phi\mapsto\theta_\phi\) are mutually inverse bijections.

For \SGARefLink{sga3:i:proposition:1-7-1:item:b}{(b)}, if \(S\in\operatorname{Ob}(\mathscr C)\), then by \SGARefLink{sga3:i:subsection:1:7:item:ii}{1.7(ii)} and
part \SGARefLink{sga3:i:proposition:1-7-1:item:a}{(a)}, applied to \(\mathscr C_{/S}\),
\begin{align*}
  &\underline{\operatorname{Hom}}
    \bigl(\mathbf E,
      \underline{\operatorname{Hom}}(\mathbf F,\mathbf G)\bigr)(S)\\
  &\quad\simeq
    \operatorname{Hom}_S
      \bigl(\mathbf E_S,
        \underline{\operatorname{Hom}}_S(\mathbf F_S,\mathbf G_S)\bigr)\\
  &\quad\simeq
    \operatorname{Hom}_S
      (\mathbf E_S\times_S\mathbf F_S,\mathbf G_S)\\
  &\quad\simeq
    \operatorname{Hom}(\mathbf E\times\mathbf F\times S,\mathbf G)\\
  &\quad\simeq
    \underline{\operatorname{Hom}}
      (\mathbf E\times\mathbf F,\mathbf G)(S),
\end{align*}
and these isomorphisms are functorial in \(S\).

\medskip
\noindent\SGAInlineRefTarget{sga3:i:corollary:1-7-2}\textbf{Corollary 1.7.2.}
\textit{There are isomorphisms
\[
  \operatorname{Hom}
    \bigl(\mathbf E,
      \underline{\operatorname{Hom}}(\mathbf F,\mathbf G)\bigr)
    \simeq
  \operatorname{Hom}
    \bigl(\mathbf F,
      \underline{\operatorname{Hom}}(\mathbf E,\mathbf G)\bigr),
\]
and
\[
  \underline{\operatorname{Hom}}
    \bigl(\mathbf E,
      \underline{\operatorname{Hom}}(\mathbf F,\mathbf G)\bigr)
    \simeq
  \underline{\operatorname{Hom}}
    \bigl(\mathbf F,
      \underline{\operatorname{Hom}}(\mathbf E,\mathbf G)\bigr).
\]}

In particular, taking \(\mathbf E=\underline{\mathbf e}\) and using
\(\underline{\operatorname{Hom}}(\underline{\mathbf e},\mathbf G)
\simeq\mathbf G\), we obtain
\[
  \Gamma\bigl(\underline{\operatorname{Hom}}(\mathbf F,\mathbf G)\bigr)
    \simeq\operatorname{Hom}(\mathbf F,\mathbf G).
\]
Composition supplies functorial morphisms
\[
  \underline{\operatorname{Hom}}(\mathbf F,\mathbf G)
    \times
  \underline{\operatorname{Hom}}(\mathbf G,\mathbf H)
    \longrightarrow
  \underline{\operatorname{Hom}}(\mathbf F,\mathbf H).
\]

If \(\mathbf F\) and \(\mathbf G\) are two objects of
\(\widehat{\mathscr C}\), let \(\operatorname{Isom}(\mathbf F,\mathbf G)\)
denote the subset of \(\operatorname{Hom}(\mathbf F,\mathbf G)\) formed by
the isomorphisms from \(\mathbf F\) onto \(\mathbf G\). Define a subobject
\(\underline{\operatorname{Isom}}(\mathbf F,\mathbf G)\) of
\(\underline{\operatorname{Hom}}(\mathbf F,\mathbf G)\) by
\[
  \underline{\operatorname{Isom}}(\mathbf F,\mathbf G)(S)
    =\operatorname{Isom}(\mathbf F_S,\mathbf G_S).
\]
Then
\[
  \Gamma\bigl(\underline{\operatorname{Isom}}
    (\mathbf F,\mathbf G)\bigr)
    \simeq\operatorname{Isom}(\mathbf F,\mathbf G),
  \qquad
  \underline{\operatorname{Isom}}(\mathbf F,\mathbf G)
    \simeq
  \underline{\operatorname{Isom}}(\mathbf G,\mathbf F).
\]
In the special case \(\mathbf F=\mathbf G\), set
\begin{align*}
  \underline{\operatorname{End}}(\mathbf F)
    &=\underline{\operatorname{Hom}}(\mathbf F,\mathbf F),
  &
  \operatorname{End}(\mathbf F)
    &=\operatorname{Hom}(\mathbf F,\mathbf F)
      \simeq\Gamma\bigl(\underline{\operatorname{End}}(\mathbf F)\bigr),\\
  \underline{\operatorname{Aut}}(\mathbf F)
    &=\underline{\operatorname{Isom}}(\mathbf F,\mathbf F),
  &
  \operatorname{Aut}(\mathbf F)
    &=\operatorname{Isom}(\mathbf F,\mathbf F)
      \simeq\Gamma\bigl(\underline{\operatorname{Aut}}(\mathbf F)\bigr).
\end{align*}
Formation of the objects
\(\underline{\operatorname{Hom}}\),
\(\underline{\operatorname{Isom}}\),
\(\underline{\operatorname{Aut}}\), and
\(\underline{\operatorname{End}}\) commutes with base change.

\medskip
\noindent\SGAInlineRefTarget{sga3:i:remark:1-7-3}\textbf{Remark 1.7.3.}
\footnote{Editor's note: the numbering \SGARefLink{sga3:i:remark:1-7-3}{1.7.3} has been added for later
references.}
An object isomorphic to
\(\underline{\operatorname{Isom}}(\mathbf F,\mathbf G)\) may also be
constructed as follows. There is a morphism
\[
  \underline{\operatorname{Hom}}(\mathbf F,\mathbf G)
    \times
  \underline{\operatorname{Hom}}(\mathbf G,\mathbf F)
    \longrightarrow
  \underline{\operatorname{End}}(\mathbf F).
\]
Interchanging \(\mathbf F\) and \(\mathbf G\) gives a morphism
\[
  \underline{\operatorname{Hom}}(\mathbf F,\mathbf G)
    \times
  \underline{\operatorname{Hom}}(\mathbf G,\mathbf F)
    \longrightarrow
  \underline{\operatorname{End}}(\mathbf F)
    \times
  \underline{\operatorname{End}}(\mathbf G).
\]
On the other hand, the identity morphism of \(\mathbf F\) is an element of
\(\operatorname{End}(\mathbf F)\), and hence defines a morphism
\(\underline{\mathbf e}\to
\underline{\operatorname{End}}(\mathbf F)\). Doing the same for
\(\mathbf G\) and taking the product gives a morphism
\[
  \underline{\mathbf e}\longrightarrow
  \underline{\operatorname{End}}(\mathbf F)
    \times
  \underline{\operatorname{End}}(\mathbf G).
\]
It is then immediate that the fiber product of \(\underline{\mathbf e}\)
and
\[
  \underline{\operatorname{Hom}}(\mathbf F,\mathbf G)
    \times
  \underline{\operatorname{Hom}}(\mathbf G,\mathbf F)
\]
over
\(\underline{\operatorname{End}}(\mathbf F)
\times\underline{\operatorname{End}}(\mathbf G)\)
is isomorphic to
\(\underline{\operatorname{Isom}}(\mathbf F,\mathbf G)\).

All these definitions apply, in particular, when
\(\mathbf F=\mathbf h_X\) and \(\mathbf G=\mathbf h_Y\). If
\(\underline{\operatorname{Hom}}(\mathbf h_X,\mathbf h_Y)\) is
representable by an object of \(\mathscr C\), we denote that object by
\(\underline{\operatorname{Hom}}(X,Y)\). It has the following property:
if \(Z\times X\) exists, then
\[
  \operatorname{Hom}
    \bigl(Z,\underline{\operatorname{Hom}}(X,Y)\bigr)
    \simeq
  \operatorname{Hom}(Z\times X,Y).
\]
This property characterizes it when products exist in \(\mathscr C\).

Similarly, one defines, when they exist, the objects
\[
  \underline{\operatorname{Isom}}(X,Y),
  \qquad
  \underline{\operatorname{End}}(X),
  \qquad
  \underline{\operatorname{Aut}}(X).
\]
We merely note that the construction above shows that
\(\underline{\operatorname{Isom}}(X,Y)\) exists whenever fiber products
exist in \(\mathscr C\) and the objects
\[
  \underline{\operatorname{Hom}}(X,Y),
  \quad
  \underline{\operatorname{Hom}}(Y,X),
  \quad
  \underline{\operatorname{End}}(X),
  \quad\text{and}\quad
  \underline{\operatorname{End}}(Y)
\]
exist.

Everything above also applies to categories of the form \(\mathscr C_{/S}\).
The corresponding objects will be denoted as explicitly as possible by
appropriate symbols. For example, if \(T\) and \(T'\) are two objects of
\(\mathscr C\) over \(S\), we denote by
\(\underline{\operatorname{Hom}}_S(T,T')\) the object
\[
  \underline{\operatorname{Hom}}_{\mathscr C_{/S}}(T/S,T'/S).
\]

% End of assigned component. Section 1.8 begins in the next unit.

% END INCLUDED SGA3_English_current_buildable_source_20260823/source/inputs/I/tex/components/02_expose_I_s16_next_pp005_010.tex

% BEGIN INCLUDED SGA3_English_current_buildable_source_20260823/source/inputs/I/tex/components/03_expose_I_s18_s22_pp010_015.tex
% SGA 3, Expose I, Loop 1 English reconstruction.
% Scope: Section 1.8 through the end of Section 2.2, including the transition
% sentence immediately before Section 2.3.
% Authority: Polo--Gille Exp1-13oct24.pdf, component pages 10--15
% (combined-reader pages 20--25). OCR and the Reinhold English Markdown are
% drafting/comparison witnesses only.

\SGARefTarget{sga3:i:subsection:1:8}
\subsection*{1.8. Constant objects}

Let \(\mathscr C\) be a category in which direct sums and fiber products
exist and in which direct sums commute with base change (for example, the
category of schemes).\footnote{Editor's note: the former terminology
``preschemes/schemes'' has been replaced by the current terminology
``schemes/separated schemes.''} For every set \(E\) and every object \(S\) of
\(\mathscr C\), set
\begin{equation}
  \tag{1.8.1}\label{sga3:i:subsection:1-8:equation:1-8-1}
  E_S=
  \left\{
    \begin{aligned}
      &\text{the direct sum of a family }(S_i)_{i\in E}\text{ of objects}\
      &\text{of }\mathscr C\text{, all isomorphic to }S.
    \end{aligned}
  \right.
\end{equation}
This object is characterized by the formula
\begin{equation}
  \tag{1.8.2}\label{sga3:i:subsection:1-8:equation:1-8-2}
  \operatorname{Hom}_{\mathscr C}(E_S,T)
  =\operatorname{Hom}_{\mathbf{Set}}
      \bigl(E,\operatorname{Hom}_{\mathscr C}(S,T)\bigr),
\end{equation}
for every \(T\in\operatorname{Ob}(\mathscr C)\).

The object \(E_S\) has a canonical projection to \(S\), so that
\(E\mapsto E_S\) is in fact a functor from \(\mathbf{Set}\) to
\(\mathscr C_{/S}\).\footnote{Editor's note: in the next three paragraphs,
the order of the sentences has been changed and a few details have been added
concerning the role of condition \((*)\) below.} If \(S'\to S\) is an arrow
of \(\mathscr C\), then, because direct sums commute with base change,
\[
  E_{S'}=(E_S)_{S'}.
\]
In particular, if \(\mathscr C\) has a final object \(e\), then
\[
  E_S=(E_e)_S.
\]

The functor \(E\mapsto E_S/S\), from \(\mathbf{Set}\) to
\(\mathscr C_{/S}\), commutes with finite products. It is enough to prove
that
\begin{equation}
  \tag{\ensuremath{\times}}
  E_S\mathbin{\times_S}F_S=(E\times F)_S.
\end{equation}
Indeed, by the results of \SGARefLink{sga3:i:subsection:1:7}{1.7} applied to \(\mathscr C_{/S}\), for every
\(T\in\operatorname{Ob}(\mathscr C_{/S})\) there are natural isomorphisms
(all unspecified \(\operatorname{Hom}\)'s being taken in
\(\mathscr C_{/S}\))
\begin{align*}
  \operatorname{Hom}((E\times F)_S,T)
  &\simeq
    \operatorname{Hom}_{\mathbf{Set}}
      \bigl(E\times F,\operatorname{Hom}(S,T)\bigr)\\
  &\simeq
    \operatorname{Hom}_{\mathbf{Set}}
      \bigl(E,
        \operatorname{Hom}_{\mathbf{Set}}
          (F,\operatorname{Hom}(S,T))\bigr)\\
  &\simeq
    \operatorname{Hom}_{\mathbf{Set}}
      \bigl(E,\operatorname{Hom}(F_S,T)\bigr),
\end{align*}
and
\begin{align*}
  \operatorname{Hom}(E_S\mathbin{\times_S}F_S,T)
  &\simeq
    \operatorname{Hom}
      \bigl(E_S,\underline{\operatorname{Hom}}(F_S,T)\bigr)\\
  &\simeq
    \operatorname{Hom}_{\mathbf{Set}}
      \bigl(E,
        \operatorname{Hom}
          (S,\underline{\operatorname{Hom}}(F_S,T))\bigr).
\end{align*}
But
\[
  \operatorname{Hom}
    \bigl(S,\underline{\operatorname{Hom}}(F_S,T)\bigr)
  \simeq \operatorname{Hom}(F_S,T),
\]
which proves \((\times)\).

Suppose that, with \(\varnothing\) denoting an initial object of
\(\mathscr C\), the following diagram is cartesian:

% Authority: Exp1-13oct24.pdf p. 11; combined reader p. 21.
\begin{center}
  \makebox[0.62\textwidth][l]{%
    \textup{(*)}\hfill
    % BEGIN INCLUDED SGA3_English_current_buildable_source_20260823/source/inputs/I/native_diagrams/exp1/diag_006_constant_objects_cartesian_square.tex
\begin{tikzcd}[column sep=4.2em,row sep=3.0em]
  \varnothing
    \arrow[r]
    \arrow[d]
  &
  S
    \arrow[d]
  \\
  S
    \arrow[r]
  &
  S\amalg S
\end{tikzcd}

% END INCLUDED SGA3_English_current_buildable_source_20260823/source/inputs/I/native_diagrams/exp1/diag_006_constant_objects_cartesian_square.tex
%
    \hfill\phantom{\textup{(*)}}%
  }
\end{center}
This is the case for the category of schemes.\footnote{Editor's note:
Condition \((*)\) is not satisfied if \(\mathscr C\) is the category whose
arrows are \(A\to B\), \(\operatorname{id}_A\), and
\(\operatorname{id}_B\). In that case
\(B\coprod B=B\) and \(B\times_BB=B\not\simeq A\).}
Then the functor \(E\mapsto E_S\) commutes with finite inverse limits.

Indeed, in view of \((\times)\), it is enough to show that
\(E\mapsto E_S\) commutes with fiber products. Let \(u:E\to G\) and
\(v:F\to G\) be two maps of sets. Since direct sums in \(\mathscr C\)
commute with base change, one has
\begin{equation}
  \tag{1}\label{sga3:i:subsection:1-8:equation:1}
  F_S\mathbin{\times_{G_S}}E_S
  \simeq
  \coprod_{f\in F}
    \bigl(S_f\mathbin{\times_{G_S}}E_S\bigr)
  \simeq
  \coprod_{\substack{f\in F\\x\in E}}
    \bigl(S_f\mathbin{\times_{G_S}}S_x\bigr).
\end{equation}
If \(v(f)\ne u(x)\), there is a morphism in \(\mathscr C_{/S}\)
\[
  S_f\mathbin{\times_{G_S}}S_x
  \longrightarrow
  S_f\mathbin{\times_{S_{v(f)}\coprod S_{u(x)}}}S_x.
\]
By condition \((*)\), the term on the right is \(\varnothing\). Hence
\begin{equation}
  \tag{2}\label{sga3:i:subsection:1-8:equation:2}
  S_f\mathbin{\times_{G_S}}S_x\simeq\varnothing
  \qquad\text{if }v(f)\ne u(x).
\end{equation}
On the other hand, if \(v(f)=u(x)\), there is a morphism in
\(\mathscr C_{/S}\)
\[
  S\longrightarrow S_f\mathbin{\times_{G_S}}S_x.
\]
Since \(S\xrightarrow{\operatorname{id}}S\) is the final object of
\(\mathscr C_{/S}\), it follows that
\begin{equation}
  \tag{3}\label{sga3:i:subsection:1-8:equation:3}
  S_f\mathbin{\times_{G_S}}S_x\simeq S
  \qquad\text{if }v(f)=u(x).
\end{equation}
Combining \SGARefLink{sga3:i:proposition:1:4:1}{(1)}, \SGARefLink{sga3:i:section:2}{(2)}, and \SGARefLink{sga3:i:section:3}{(3)}, we obtain an isomorphism functorial in \(S\):
\[
  F_S\mathbin{\times_{G_S}}E_S
  \simeq \coprod_{F\times_G E}S
  =(F\mathbin{\times_G}E)_S.
\]

An object of the form \(E_S\) is called a \emph{constant object}. There is
a morphism, functorial in \(E\),
\[
  E\longrightarrow\Gamma(E_S/S),
\]
which associates with each \(i\in E\) the section of \(E_S\) over \(S\)
defined by the isomorphism from \(S\) onto \(S_i\). Suppose condition
\((*)\) is satisfied for every object \(S\) of \(\mathscr C\). Then
\(E\to\Gamma(E_S/S)\) is a monomorphism whenever
\(S\not\simeq\varnothing\).

If \(\mathscr C\) is the category of schemes, then
\(\Gamma(E_S/S)\) is identified with the locally constant maps from the
topological space \(S\) to the set \(E\), the preceding map associating with
each element of \(E\) the corresponding constant map. It follows that
\(E_S\) may equivalently be defined as representing the functor that assigns
to every \(S'\) over \(S\) the set of locally constant functions from the
topological space \(S'\) to \(E\).\footnote{Editor's note: the diagonal
morphism
\(E_S\to E_S\times_SE_S=(E\times E)_S\) is also a closed immersion; that
is, \(E_S\) is separated over \(S\).}

\section*{\SGAInlineRefTarget{sga3:i:section:2}2. Algebraic structures}
\addcontentsline{toc}{section}{2. Algebraic structures}

Given a species of algebraic structure in the category of sets, we shall
extend it to the category \(\mathscr C\). We first treat the example of
groups.

\subsection*{\SGAInlineRefTarget{sga3:i:subsection:2-1}2.1. Group structures}

We retain the notation of the preceding paragraph.

\medskip
\noindent\SGAInlineRefTarget{sga3:i:definition:2-1-1}\textbf{Definition 2.1.1.}
Let \(\mathbf G\in\operatorname{Ob}(\widehat{\mathscr C})\). A
\emph{\(\widehat{\mathscr C}\)-group structure} on \(\mathbf G\) is the
datum, for every \(S\in\operatorname{Ob}(\mathscr C)\), of a group structure
on the set \(\mathbf G(S)\), such that for every arrow
\(f:S'\to S''\) of \(\mathscr C\), the map
\[
  \mathbf G(f):\mathbf G(S'')\longrightarrow\mathbf G(S')
\]
is a group homomorphism. If \(\mathbf G\) and \(\mathbf H\) are two
\(\widehat{\mathscr C}\)-groups, a \emph{morphism of
\(\widehat{\mathscr C}\)-groups} from \(\mathbf G\) to \(\mathbf H\) is a
morphism \(u\in\operatorname{Hom}(\mathbf G,\mathbf H)\) such that
\(u(S):\mathbf G(S)\to\mathbf H(S)\) is a group homomorphism for every
\(S\in\operatorname{Ob}(\mathscr C)\).

We denote by
\(\operatorname{Hom}_{\widehat{\mathscr C}\text{-}\mathrm{Gr}}
  (\mathbf G,\mathbf H)\)
the set of these morphisms and by
\((\widehat{\mathscr C}\text{-}\mathrm{Gr})\) their category.

\medskip
\noindent\textbf{Examples.}
For \(\mathbf E\in\operatorname{Ob}(\widehat{\mathscr C})\), the object
\(\underline{\operatorname{Aut}}(\mathbf E)\) has an evident
\(\widehat{\mathscr C}\)-group structure. The final object
\(\underline{\mathbf e}\) has a unique
\(\widehat{\mathscr C}\)-group structure, making it a final object of
\((\widehat{\mathscr C}\text{-}\mathrm{Gr})\).

For every \(S\in\operatorname{Ob}(\mathscr C)\), let
\(e_{\mathbf G}(S)\) be the identity element of \(\mathbf G(S)\). The family
of elements \(e_{\mathbf G}(S)\) defines
\[
  e_{\mathbf G}\in\Gamma(\mathbf G)
  =\operatorname{Hom}(\underline{\mathbf e},\mathbf G).
\]
It is a morphism of \(\widehat{\mathscr C}\)-groups
\(\underline{\mathbf e}\to\mathbf G\), called the \emph{unit section} of
\(\mathbf G\).

To give a \(\widehat{\mathscr C}\)-group structure on \(\mathbf G\) is
equivalently to give a composition law on \(\mathbf G\), that is, a
\(\widehat{\mathscr C}\)-morphism
\[
  \pi_{\mathbf G}:\mathbf G\times\mathbf G\longrightarrow\mathbf G
\]
such that \(\pi_{\mathbf G}(S)\) equips \(\mathbf G(S)\) with a group
structure for every \(S\in\operatorname{Ob}(\mathscr C)\).

Similarly, \(f:\mathbf G\to\mathbf H\) is a morphism of
\(\widehat{\mathscr C}\)-groups if and only if the following diagram is
commutative:

% Authority: Exp1-13oct24.pdf p. 13; combined reader p. 23.
\begin{center}
  % BEGIN INCLUDED SGA3_English_current_buildable_source_20260823/source/inputs/I/native_diagrams/exp1/diag_007_group_morphism_square.tex
\begin{tikzcd}[column sep=4.5em,row sep=3.0em]
  \mathbf G\times\mathbf G
    \arrow[r,"\pi_{\mathbf G}"]
    \arrow[d,"{(f,f)}"']
  &
  \mathbf G
    \arrow[d,"f"]
  \\
  \mathbf H\times\mathbf H
    \arrow[r,"\pi_{\mathbf H}"]
  &
  \mathbf H
  \qquad .
\end{tikzcd}

% END INCLUDED SGA3_English_current_buildable_source_20260823/source/inputs/I/native_diagrams/exp1/diag_007_group_morphism_square.tex

\end{center}

A subobject \(\mathbf H\) of \(\mathbf G\) such that \(\mathbf H(S)\) is a
subgroup of \(\mathbf G(S)\) for every \(S\) has the evident
\(\widehat{\mathscr C}\)-group structure induced by that of \(\mathbf G\).
It is the unique structure for which the monomorphism
\(\mathbf H\to\mathbf G\) is a morphism of
\(\widehat{\mathscr C}\)-groups. With this structure, \(\mathbf H\) is
called a \emph{sub-\(\widehat{\mathscr C}\)-group} of \(\mathbf G\).

If \(\mathbf G\) and \(\mathbf H\) are two
\(\widehat{\mathscr C}\)-groups, then \(\mathbf G\times\mathbf H\) has the
evident product structure: for every \(S\), the set
\(\mathbf G(S)\times\mathbf H(S)\) is given the product of the prescribed
group structures. The resulting object is called the
\emph{product \(\widehat{\mathscr C}\)-group} of \(\mathbf G\) and
\(\mathbf H\); it is their product in the category of
\(\widehat{\mathscr C}\)-groups.

If \(\mathbf G\) is a \(\widehat{\mathscr C}\)-group, then
\(\mathbf G_S\) is a \(\widehat{\mathscr C_{/S}}\)-group for every
\(S\in\operatorname{Ob}(\mathscr C)\). For two
\(\widehat{\mathscr C}\)-groups \(\mathbf G\) and \(\mathbf H\), define
the object
\(\underline{\operatorname{Hom}}_{
  \widehat{\mathscr C}\text{-}\mathrm{Gr}}(\mathbf G,\mathbf H)\)
of \(\widehat{\mathscr C}\) by
\[
  \underline{\operatorname{Hom}}_{
    \widehat{\mathscr C}\text{-}\mathrm{Gr}}(\mathbf G,\mathbf H)(S)
  =
  \operatorname{Hom}_{
    \widehat{\mathscr C_{/S}}\text{-}\mathrm{Gr}}
    (\mathbf G_S,\mathbf H_S).
\]
Notice that
\(\underline{\operatorname{Hom}}_{
  \widehat{\mathscr C}\text{-}\mathrm{Gr}}\)
is not in general a \(\widehat{\mathscr C}\)-group, still less the internal
\(\underline{\operatorname{Hom}}\) object in the category
\((\widehat{\mathscr C}\text{-}\mathrm{Gr})\).

One similarly defines
\[
  \underline{\operatorname{Isom}}_{
    \widehat{\mathscr C}\text{-}\mathrm{Gr}}(\mathbf G,\mathbf H),
  \qquad
  \underline{\operatorname{End}}_{
    \widehat{\mathscr C}\text{-}\mathrm{Gr}}(\mathbf G),
  \qquad
  \underline{\operatorname{Aut}}_{
    \widehat{\mathscr C}\text{-}\mathrm{Gr}}(\mathbf G).
\]

\medskip
\noindent\SGAInlineRefTarget{sga3:i:definition:2-1-2}\textbf{Definition 2.1.2.}
Let \(G\in\operatorname{Ob}(\mathscr C)\). A
\emph{\(\mathscr C\)-group structure} on \(G\) is a
\(\widehat{\mathscr C}\)-group structure on
\(\mathbf h_G\in\operatorname{Ob}(\widehat{\mathscr C})\). A morphism from
the \(\mathscr C\)-group \(G\) to the \(\mathscr C\)-group \(H\) is an
element
\[
  u\in\operatorname{Hom}(G,H)
  \simeq\operatorname{Hom}(\mathbf h_G,\mathbf h_H)
\]
that defines a morphism of \(\widehat{\mathscr C}\)-groups from
\(\mathbf h_G\) to \(\mathbf h_H\).

We denote the category of \(\mathscr C\)-groups by
\((\mathscr C\text{-}\mathrm{Gr})\). There is a cartesian square in
\(\mathbf{Cat}\):

% Authority: Exp1-13oct24.pdf p. 13; combined reader p. 23.
\begin{center}
  % BEGIN INCLUDED SGA3_English_current_buildable_source_20260823/source/inputs/I/native_diagrams/exp1/diag_008_representable_group_cartesian_square.tex
\begin{tikzcd}[column sep=5.0em,row sep=3.2em]
  (\mathscr C\text{-}\mathrm{Gr.})
    \arrow[r]
    \arrow[d]
  &
  (\widehat{\mathscr C}\text{-}\mathrm{Gr.})
    \arrow[d]
  \\
  \mathscr C
    \arrow[r,"\mathbf h"]
  &
  \widehat{\mathscr C}
  \qquad .
\end{tikzcd}

% END INCLUDED SGA3_English_current_buildable_source_20260823/source/inputs/I/native_diagrams/exp1/diag_008_representable_group_cartesian_square.tex

\end{center}

All the preceding definitions and constructions therefore carry over at
once to \((\mathscr C\text{-}\mathrm{Gr})\) whenever the functors they use
(products, internal \(\underline{\operatorname{Hom}}\) objects, and so on)
are representable. They also apply to the categories \(\mathscr C_{/S}\).
In that case we write
\(\underline{\operatorname{Hom}}_{S\text{-}\mathrm{Gr}}\) for
\(\underline{\operatorname{Hom}}_{\mathscr C_{/S}\text{-}\mathrm{Gr}}\),
and similarly for the other objects.

\subsection*{\SGAInlineRefTarget{sga3:i:subsection:2-2}2.2.}

More generally, suppose \((T)\) is a species of structure on \(n\) base sets
defined by finite inverse limits---for example, by commutative diagrams
constructed from cartesian products, as in the structures of a monoid, a
group, a set with operators, a module over a ring, or a Lie algebra over a
ring. The preceding construction defines the notion of a ``structure of
species \((T)\)'' on \(n\) objects
\(\mathbf F_1,\ldots,\mathbf F_n\) of \(\widehat{\mathscr C}\). Such a
structure consists, for every \(S\in\operatorname{Ob}(\mathscr C)\), of a
structure of species \((T)\) on the sets
\(\mathbf F_1(S),\ldots,\mathbf F_n(S)\), in such a way that for every arrow
\(S'\to S''\) of \(\mathscr C\), the family of maps
\[
  \bigl(\mathbf F_i(S'')\bigr)_i
  \longrightarrow
  \bigl(\mathbf F_i(S')\bigr)_i
\]
is a polyhomomorphism for the species \((T)\).

Morphisms of structures of species \((T)\) are defined similarly, giving a
category
\[
  (\widehat{\mathscr C}\times\cdots\times\widehat{\mathscr C})^{(T)}.
\]
The fully faithful functor \(\mathbf h\times\cdots\times\mathbf h\) then
defines, by inverse image, the category
\[
  (\mathscr C\times\cdots\times\mathscr C)^{(T)}.
\]
Since it commutes with inverse limits, all the functorial properties,
notions, and notation introduced in \(\widehat{\mathscr C}\) carry over to
this category.

Now suppose fiber products exist in \(\mathscr C\), and let \((T)\) be a
species of algebraic structure specified by certain morphisms between
cartesian products, subject to axioms consisting of commutative diagrams
formed with those morphisms. A structure of species \((T)\) on a family of
objects of \(\mathscr C\) is then defined by corresponding morphisms between
cartesian products satisfying the corresponding commutativity conditions.
Consequently, if \(\mathscr C\) and \(\mathscr C'\) are categories with
products and \(f:\mathscr C\to\mathscr C'\) is a functor commuting with
products, then a structure of species \((T)\) on a family
\((F_i)\) of objects of \(\mathscr C\) induces one on the family
\((f(F_i))\) in \(\mathscr C'\). Thus every \(\mathscr C\)-group becomes a
\(\mathscr C'\)-group, every pair consisting of a \(\mathscr C\)-ring and a
module over it becomes the analogous pair in \(\mathscr C'\), and so on.

In particular, let \(\mathscr C\) satisfy the conditions of \SGARefLink{sga3:i:subsection:1:8}{1.8}.\footnote{
Editor's note: including condition \((*)\).} The functor
\(E\mapsto E_S\) defined there commutes with finite inverse limits. It
therefore carries groups to \(S\)-groups (that is,
\(\mathscr C_{/S}\)-groups), rings to \(S\)-rings, and so forth.

\medskip
\noindent\textbf{Remark.}
It is worth observing that the preceding construction, when applied to
\(\widehat{\mathscr C}\), gives back precisely the notions already defined
there. In other words, to give an object of \(\widehat{\mathscr C}\) a
structure of species \((T)\) when it is viewed as a functor on
\(\mathscr C\) is the same as to give the representable functor on
\(\widehat{\mathscr C}\) defined by that object a structure of species
\((T)\).\footnote{Editor's note: for group structures, for example, let
\(\mathbf G\in\operatorname{Ob}(\widehat{\mathscr C})\). If the functor
\[
  \widehat{\mathscr C}\longrightarrow\mathbf{Set},
  \qquad
  \mathbf F\longmapsto
  \operatorname{Hom}_{\widehat{\mathscr C}}(\mathbf F,\mathbf G)
\]
has a group structure, then so does its restriction to \(\mathscr C\),
namely \(X\mapsto\mathbf G(X)\). Conversely, if \(\mathbf G\) is a
\(\widehat{\mathscr C}\)-group, its multiplication morphism
\(\pi_{\mathbf G}:\mathbf G\times\mathbf G\to\mathbf G\) induces, for every
\(\mathbf F\in\operatorname{Ob}(\widehat{\mathscr C})\), a group structure
on
\(\operatorname{Hom}_{\widehat{\mathscr C}}(\mathbf F,\mathbf G)\),
functorially in \(\mathbf F\).}

We shall treat two further special cases of this construction: structures
with operator groups and modules.

% End of component. Section 2.3 begins in the next unit.

% END INCLUDED SGA3_English_current_buildable_source_20260823/source/inputs/I/tex/components/03_expose_I_s18_s22_pp010_015.tex

% BEGIN INCLUDED SGA3_English_current_buildable_source_20260823/source/inputs/I/tex/components/04_expose_I_s23_operator_groups.tex
% SGA 3, Expose I, Loop 1 English reconstruction.
% Scope: Section 2.3, from the heading through Proposition 2.3.8 and
% editor's notes (23)--(26); stops immediately before Section 3.
% Authority: Polo--Gille Exp1-13oct24.pdf, component local pages 15--17
% (combined-reader pages 25--27). OCR is drafting/locator material only.

\subsection*{\SGAInlineRefTarget{sga3:i:subsection:2-3}2.3. Structures with operator groups}

\medskip
\noindent\SGAInlineRefTarget{sga3:i:definition:2-3-1}\textbf{Definition 2.3.1.}
Let \(\mathbf E\in\operatorname{Ob}(\widehat{\mathscr C})\) and
\(\mathbf G\in
\operatorname{Ob}(\widehat{\mathscr C}\text{-}\mathbf{Gr})\).
A structure on \(\mathbf E\) of an object with
\(\widehat{\mathscr C}\)-operator group \(\mathbf G\) (or of a
\(\mathbf G\)-object) consists in giving \(\mathbf E(S)\), for every
\(S\in\operatorname{Ob}(\mathscr C)\), the structure of a set with
operator group \(\mathbf G(S)\), in such a way that, for every arrow
\(S'\to S''\) in \(\mathscr C\), the map of sets
\(\mathbf E(S'')\to\mathbf E(S')\) is compatible with the homomorphism
of operator groups \(\mathbf G(S'')\to\mathbf G(S')\).

As usual, this is equivalent to giving a morphism
\[
  \mu:\mathbf G\times\mathbf E\longrightarrow\mathbf E
\]
which, for every \(S\), equips \(\mathbf E(S)\) with the structure of a
set on which \(\mathbf G(S)\) acts. But
\[
  \operatorname{Hom}(\mathbf G\times\mathbf E,\mathbf E)
    \simeq
  \operatorname{Hom}
    (\mathbf G,\underline{\operatorname{End}}(\mathbf E)),
\]
so \(\mu\) defines a morphism
\(\mathbf G\to\underline{\operatorname{End}}(\mathbf E)\). It is
immediate that this morphism maps \(\mathbf G\) into
\(\underline{\operatorname{Aut}}(\mathbf E)\) and is a morphism of
\(\widehat{\mathscr C}\)-groups. Consequently, giving \(\mathbf E\) the
structure of an object with \(\widehat{\mathscr C}\)-operator group
\(\mathbf G\) is equivalent to giving a morphism of
\(\widehat{\mathscr C}\)-groups
\[
  \rho:\mathbf G\longrightarrow
    \underline{\operatorname{Aut}}(\mathbf E).
\]
In particular, every element \(g\in\mathbf G(S)\) defines an
automorphism \(\rho(g)\) of the functor \(\mathbf E_S\), that is, an
automorphism of \(\mathbf E\times\mathbf h_S\) commuting with the
projection \(\mathbf E\times\mathbf h_S\to\mathbf h_S\), and hence, in
particular, an automorphism of the set \(\mathbf E(S')\) for every
\(S'\to S\).

\medskip
\noindent\SGAInlineRefTarget{sga3:i:definition:2-3-2}\textbf{Definition 2.3.2.}
We denote by \(\mathbf E^{\mathbf G}\) the subobject of \(\mathbf E\)
defined as follows:
\[
  \mathbf E^{\mathbf G}(S)
    =\bigl\{x\in\mathbf E(S)\mid
      x_{S'}\text{ is invariant under }\mathbf G(S')
      \text{ for every }S'\to S\bigr\},
\]
where \(x_{S'}\) denotes the image of \(x\) under
\(\mathbf E(S)\to\mathbf E(S')\).

Thus \(\mathbf E^{\mathbf G}\), the ``subobject of invariants of
\(\mathbf G\),'' is the largest subobject of \(\mathbf E\) on which
\(\mathbf G\) acts trivially.

\medskip
\SGARefTarget{sga3:i:definition:2:3:3}
\noindent\textbf{Definition 2.3.3.}
Let \(\mathbf F\) be a subobject of \(\mathbf E\). We denote by
\(\underline{\operatorname{Norm}}_{\mathbf G}\mathbf F\) and
\(\underline{\operatorname{Centr}}_{\mathbf G}\mathbf F\) the
sub-\(\widehat{\mathscr C}\)-groups of \(\mathbf G\) defined by
\begin{align*}
  (\underline{\operatorname{Norm}}_{\mathbf G}\mathbf F)(S)
    &={}
      \{g\in\mathbf G(S)\mid \rho(g)\mathbf F_S=\mathbf F_S\}\\
    &={}
      \{g\in\mathbf G(S)\mid
        \rho(g)\mathbf F(S')=\mathbf F(S')
        \text{ for every }S'\to S\},\\
  (\underline{\operatorname{Centr}}_{\mathbf G}\mathbf F)(S)
    &={}
      \{g\in\mathbf G(S)\mid
        \rho(g)|_{\mathbf F_S}=\operatorname{id}\}\\
    &={}
      \{g\in\mathbf G(S)\mid
        \rho(g)|_{\mathbf F(S')}=\operatorname{id}
        \text{ for every }S'\to S\}.
\end{align*}
Here the vertical bar after \(\rho(g)\) denotes restriction.
\footnote{Editor's note: the original has been corrected by replacing
the inclusion \(\rho(g)\mathbf F_S\subset\mathbf F_S\) with an equality,
so as to ensure that
\(\underline{\operatorname{Norm}}_{\mathbf G}(\mathbf F)\) is indeed a
group. See VI\textsubscript{B}, \SGARefLink{prop:sga3-VIB-6.4}{6.4} for conditions under which the
``transporter'' coincides with the ``strict transporter.''}

\medskip
\noindent\SGAInlineRefTarget{sga3:i:scholium:2-3-3-1}\textbf{Scholium 2.3.3.1.}
\footnote{Editor's note: \SGARefLink{sga3:i:scholium:2-3-3-1}{Scholium 2.3.3.1} and \SGARefLink{sga3:i:remark:2-3-3-2}{Remark 2.3.3.2} have been
added.}
In particular, let \(x\in\Gamma(\mathbf E)\), that is (cf. \SGARefLink{sga3:i:subsection:1-2}{1.2}), a
collection of elements \(x_S\in\mathbf E(S)\),
\(S\in\operatorname{Ob}(\mathscr C)\), such that for every arrow
\(f:S'\to S\) one has
\(\mathbf E(f)(x_S)=x_{S'}\). If \(\mathscr C\) has a final object
\(S_0\), then \(\Gamma(\mathbf E)=\mathbf E(S_0)\). The element \(x\)
then defines a subfunctor of \(\mathbf E\), denoted by \(\mathbf x\),
and
\[
  \underline{\operatorname{Norm}}_{\mathbf G}\mathbf x
    =
  \underline{\operatorname{Centr}}_{\mathbf G}\mathbf x.
\]
We denote this functor by
\(\underline{\operatorname{Stab}}_{\mathbf G}(x)\) and call it the
\emph{stabilizer} of \(x\); thus, for every
\(S\in\operatorname{Ob}(\mathscr C)\),
\[
  \underline{\operatorname{Stab}}_{\mathbf G}(x)(S)
    =\{g\in\mathbf G(S)\mid\rho(g)x_S=x_S\}.
\]

Suppose that fiber products exist in \(\mathscr C\). If
\(\mathbf G=\mathbf h_G\) (respectively,
\(\mathbf E=\mathbf h_E\)), where \(G\) is a \(\mathscr C\)-group
(respectively, \(E\in\operatorname{Ob}(\mathscr C)\)), and if
\(\mathscr C\) has a final object \(S_0\), so that \(x\) is a morphism
\(S_0\to E\), then
\(\underline{\operatorname{Stab}}_{\mathbf G}(x)\) is represented by
the fiber product \(G\times_E S_0\), where the morphism \(G\to E\) is
the composite of
\[
  \operatorname{id}_G\times x:
  G=G\times S_0\longrightarrow G\times E
\]
and
\[
  \mu:G\times E\longrightarrow E.
\]

\medskip
\noindent\SGAInlineRefTarget{sga3:i:remark:2-3-3-2}\textbf{Remark 2.3.3.2.}\footnotemark[\value{footnote}]
The formation of
\(\mathbf E^{\mathbf G}\),
\(\underline{\operatorname{Norm}}_{\mathbf G}\mathbf F\), and
\(\underline{\operatorname{Centr}}_{\mathbf G}\mathbf F\) commutes
with base change; that is, for every
\(S\in\operatorname{Ob}(\mathscr C)\),
\[
  (\mathbf E^{\mathbf G})_S=(\mathbf E_S)^{\mathbf G_S},
  \qquad
  (\underline{\operatorname{Norm}}_{\mathbf G}\mathbf F)_S
    \simeq
  \underline{\operatorname{Norm}}_{\mathbf G_S}\mathbf F_S,
  \qquad
  (\underline{\operatorname{Centr}}_{\mathbf G}\mathbf F)_S
    \simeq
  \underline{\operatorname{Centr}}_{\mathbf G_S}\mathbf F_S.
\]

\medskip
\noindent\SGAInlineRefTarget{sga3:i:definition:2-3-4}\textbf{Definition 2.3.4.}
If \(G\) is a \(\mathscr C\)-group and \(\mathbf E\) is an object of
\(\widehat{\mathscr C}\) (respectively, \(E\) is an object of
\(\mathscr C\)), a \(G\)-object structure on \(\mathbf E\)
(respectively, on \(E\)) is an \(\mathbf h_G\)-object structure on
\(\mathbf E\) (respectively, on \(\mathbf h_E\)).

In view of this definition, all the notions and notation defined above
carry over to \(\mathscr C\) whenever they involve only representable
functors. For example, if
\(\underline{\operatorname{Norm}}_{\mathbf h_G}(\mathbf h_F)\) is
representable, then there is a unique subobject of \(G\) representing
it; it is then a sub-\(\mathscr C\)-group of \(G\), and is denoted by
\(\operatorname{Norm}_G(F)\), and so forth.

\medskip
\noindent\SGAInlineRefTarget{sga3:i:definition:2-3-5}\textbf{Definition 2.3.5.}
\begin{enumerate}
  \item[(a)]
  We say that the \(\widehat{\mathscr C}\)-group \(\mathbf G\) acts on
  the \(\widehat{\mathscr C}\)-group \(\mathbf H\) if \(\mathbf H\)
  is equipped with a \(\mathbf G\)-object structure such that, for
  every \(g\in\mathbf G(S)\), the automorphism of \(\mathbf H(S)\)
  defined by \(g\) is a group automorphism.

  Equivalently, for every \(g\in\mathbf G(S)\), the automorphism
  \(\rho(g)\) of \(\mathbf H_S\) is an automorphism of
  \(\widehat{\mathscr C}_{/S}\)-groups; or equivalently, the morphism of
  \(\widehat{\mathscr C}\)-groups
  \(\mathbf G\to\underline{\operatorname{Aut}}(\mathbf H)\) maps
  \(\mathbf G\) into
  \[
    \underline{\operatorname{Aut}}_{
      \widehat{\mathscr C}\text{-}\mathbf{Gr}}(\mathbf H).
  \]

  \item[(b)]
  In the situation above, there is a unique
  \(\widehat{\mathscr C}\)-group structure on
  \(\mathbf H\times\mathbf G\) such that, for every \(S\),
  \((\mathbf H\times\mathbf G)(S)\) is the semidirect product of the
  groups \(\mathbf H(S)\) and \(\mathbf G(S)\) for the given action of
  \(\mathbf G(S)\) on \(\mathbf H(S)\). We denote this
  \(\widehat{\mathscr C}\)-group by \(\mathbf H\mathbin{\cdot}\mathbf G\)
  and call it the \emph{semidirect product of \(\mathbf H\) by
  \(\mathbf G\)}. Thus, by definition,
  \[
    (\mathbf H\mathbin{\cdot}\mathbf G)(S)
      =\mathbf H(S)\mathbin{\cdot}\mathbf G(S).
  \]
\end{enumerate}

Let \(\mathbf G\) be a \(\widehat{\mathscr C}\)-group. For every arrow
\(S'\to S\) in \(\mathscr C\) and every \(g\in\mathbf G(S)\), let
\(\operatorname{Int}(g)\) be the automorphism of \(\mathbf G(S')\)
defined by
\[
  \operatorname{Int}(g)h=ghg^{-1}.
\]
This definition extends to a morphism of
\(\widehat{\mathscr C}\)-groups
\[
  \operatorname{Int}:\mathbf G\longrightarrow
  \underline{\operatorname{Aut}}_{
    \widehat{\mathscr C}\text{-}\mathbf{Gr}}(\mathbf G)
  \subset
  \underline{\operatorname{Aut}}(\mathbf G).
\]

\SGARefLink{sga3:i:definition:2:3:3}{Definition 2.3.3} therefore applies, and for every subobject
\(\mathbf E\) of \(\mathbf G\) one has the sub-
\(\widehat{\mathscr C}\)-groups of \(\mathbf G\)
\[
  \underline{\operatorname{Norm}}_{\mathbf G}(\mathbf E)
  \qquad\text{and}\qquad
  \underline{\operatorname{Centr}}_{\mathbf G}(\mathbf E).
\]

\medskip
\SGARefTarget{sga3:i:definition:2:3:6}
\noindent\textbf{Definition 2.3.6.}
The \emph{center} of \(\mathbf G\), denoted
\(\underline{\operatorname{Centr}}(\mathbf G)\), is the
sub-\(\widehat{\mathscr C}\)-group
\(\underline{\operatorname{Centr}}_{\mathbf G}(\mathbf G)\) of
\(\mathbf G\). We say that \(\mathbf G\) is commutative if
\(\underline{\operatorname{Centr}}(\mathbf G)=\mathbf G\), or,
equivalently, if \(\mathbf G(S)\) is commutative for every \(S\).

We say that the sub-\(\widehat{\mathscr C}\)-group \(\mathbf H\) of
\(\mathbf G\) is invariant in \(\mathbf G\) if
\(\underline{\operatorname{Norm}}_{\mathbf G}(\mathbf H)=\mathbf G\),
or, equivalently, if \(\mathbf H(S)\) is invariant in \(\mathbf G(S)\)
for every \(S\).
\footnote{Editor's note: in addition, \(\mathbf H\) is said to be
\emph{central} in \(\mathbf G\) if
\(\underline{\operatorname{Centr}}_{\mathbf G}(\mathbf H)=\mathbf G\),
or, equivalently, if \(\mathbf H(S)\) is central in \(\mathbf G(S)\)
for every \(S\).}

\medskip
\noindent\SGAInlineRefTarget{sga3:i:definition:2-3-6-1}\textbf{Definition 2.3.6.1.}
\footnote{Editor's note: \SGARefLink{sga3:i:definition:2-3-6-1}{Definitions 2.3.6.1} and \SGARefLink{sga3:i:definition:2-3-6-2}{2.3.6.2} have been
added.}
Let \(f:\mathbf G\to\mathbf G'\) be a morphism of
\(\widehat{\mathscr C}\)-groups. The \emph{kernel} of \(f\), denoted
\(\operatorname{Ker}f\), is the sub-\(\widehat{\mathscr C}\)-group of
\(\mathbf G\) defined by
\[
  (\operatorname{Ker}f)(S)
    =\{x\in\mathbf G(S)\mid f(S)(x)=1\}
    =\operatorname{Ker}f(S)
\]
for every \(S\in\operatorname{Ob}(\mathscr C)\); it is an invariant
sub-\(\widehat{\mathscr C}\)-group.

Notice that if \(\mathbf G=\mathbf h_G\) and
\(\mathbf G'=\mathbf h_{G'}\), if \(\mathscr C\) has a final object
\(S_0\), and if fiber products exist in \(\mathscr C\), then
\(\operatorname{Ker}(f)\) is represented by \(S_0\times_{G'}G\).

\medskip
\noindent\SGAInlineRefTarget{sga3:i:definition:2-3-6-2}\textbf{Definition 2.3.6.2.}\footnotemark[\value{footnote}]
Let \(\mathbf E\in\widehat{\mathscr C}\), and let \(\mathbf G\) be a
\(\widehat{\mathscr C}\)-group acting on \(\mathbf E\). We say that
the action of \(\mathbf G\) on \(\mathbf E\) is \emph{faithful} if the
kernel of the morphism
\(\mathbf G\to\underline{\operatorname{Aut}}(\mathbf E)\) is trivial;
that is, if, for every \(S\in\operatorname{Ob}(\mathscr C)\) and
\(g\in\mathbf G(S)\), the condition
\[
  g_{S'}\cdot x=x
  \qquad
  \text{for every }S'\to S\text{ and every }x\in\mathbf E(S')
\]
implies \(g=1\).

Many definitions and propositions from elementary group theory carry
over readily. We record only the following one, which will be useful to
us.

\medskip
\noindent\SGAInlineRefTarget{sga3:i:proposition:2-3-7}\textbf{Proposition 2.3.7.}
\textit{Let \(f:\mathbf W\to\mathbf G\) be a morphism of
\(\widehat{\mathscr C}\)-groups. Put
\(\mathbf H(S)=\operatorname{Ker}f(S)\). Let
\(u:\mathbf G\to\mathbf W\) be a morphism of
\(\widehat{\mathscr C}\)-groups that is a section of \(f\) (and is
therefore necessarily a monomorphism). Then \(\mathbf W\) is identified
with the semidirect product of \(\mathbf H\) by \(\mathbf G\) for the
action of \(\mathbf G\) on \(\mathbf H\) defined by
\[
  (g,h)\longmapsto\operatorname{Int}(u(g))h
\]
for \(g\in\mathbf G(S)\), \(h\in\mathbf H(S)\), and
\(S\in\operatorname{Ob}(\mathscr C)\).}

All these definitions and propositions carry over, as usual, to
\(\mathscr C\). In particular, the semidirect product of two
\(\mathscr C\)-groups \(H\) and \(G\), with \(G\) acting on \(H\), is
defined whenever the Cartesian product \(H\times G\) exists, and the
analog of \SGARefLink{sga3:i:proposition:2-3-7}{Proposition 2.3.7} takes the following form.

\medskip
\noindent\SGAInlineRefTarget{sga3:i:proposition:2-3-8}\textbf{Proposition 2.3.8.}
\textit{Let
\[
  H\xrightarrow{i}W\xrightarrow{f}G
\]
be a sequence of morphisms of \(\mathscr C\)-groups such that, for every
\(S\in\operatorname{Ob}(\mathscr C)\), the pair \((H(S),i(S))\) is a
kernel of \(f(S):W(S)\to G(S)\). Let \(u:G\to W\) be a morphism of
\(\mathscr C\)-groups that is a section of \(f\). Then \(W\) is
identified with the semidirect product of \(H\) by \(G\) for the action
of \(G\) on \(H\) such that, if \(S\in\operatorname{Ob}(\mathscr C)\),
\(g\in G(S)\), and \(h\in H(S)\), then
\[
  \operatorname{Int}(u(g))i(h)=i({}^{g}h).
\]}

% END INCLUDED SGA3_English_current_buildable_source_20260823/source/inputs/I/tex/components/04_expose_I_s23_operator_groups.tex

% BEGIN INCLUDED SGA3_English_current_buildable_source_20260823/source/inputs/I/tex/components/05_expose_I_s3_modules.tex
% SGA 3, Expose I, Loop 1 English reconstruction.
% Scope: Section 3 in full, from its heading through the paragraph immediately
% before Section 4; printed/component pages 18--20.
% Authority: Polo--Gille Exp1-13oct24.pdf, component pages 18--20
% (combined-reader pages 28--30). Force-OCR is drafting/locator material only.

\section*{\SGAInlineRefTarget{sga3:i:section:3}3. The category of \(\mathbf O\)-modules, the category of
\(\mathbf G\)-\(\mathbf O\)-modules}
\addcontentsline{toc}{section}{3. The category of O-modules, the category of G-O-modules}

\medskip
\noindent\SGAInlineRefTarget{sga3:i:definition:3-1}\textbf{Definition 3.1.}
Let \(\mathbf O\) and \(\mathbf F\) be two objects of
\(\widehat{\mathscr C}\). We say that \(\mathbf F\) is a
\(\widehat{\mathscr C}\)-module over the
\(\widehat{\mathscr C}\)-ring \(\mathbf O\), or, for short, an
\(\mathbf O\)-module, if, for every
\(S\in\operatorname{Ob}(\mathscr C)\), the set \(\mathbf O(S)\) has been
endowed with a ring structure and \(\mathbf F(S)\) with a module structure
over that ring, in such a way that for every arrow
\(S'\to S''\) in \(\mathscr C\), the map
\(\mathbf O(S'')\to\mathbf O(S')\) is a ring homomorphism and
\(\mathbf F(S'')\to\mathbf F(S')\) is a homomorphism of abelian groups
compatible with that ring homomorphism.

Once \(\mathbf O\) is fixed, one defines in the usual way a morphism between
two \(\mathbf O\)-modules \(\mathbf F\) and \(\mathbf F'\), hence the
abelian group
\(\operatorname{Hom}_{\mathbf O}(\mathbf F,\mathbf F')\), and the category
of \(\mathbf O\)-modules, denoted \((\mathbf O\text{-Mod.})\).

\medskip
\noindent\SGAInlineRefTarget{sga3:i:lemma:3-1-1}\textbf{Lemma 3.1.1.}\footnotemark
\textit{The category \((\mathbf O\text{-Mod.})\) has the structure of an
abelian category, defined objectwise. Moreover,
\((\mathbf O\text{-Mod.})\) satisfies axiom \(\mathrm{(AB\ 5)}\)
(see \([\mathrm{Gr}57,\,1.5]\)); that is, arbitrary direct sums exist, and
if \(\mathbf M\) is an \(\mathbf O\)-module, \(\mathbf N\) a submodule,
and \((\mathbf F_i)_{i\in I}\) an increasing filtered family of submodules
of \(\mathbf M\), then
\[
  \bigcup_{i\in I}(\mathbf F_i\cap\mathbf N)
    =\left(\bigcup_{i\in I}\mathbf F_i\right)\cap\mathbf N.
\]}

Indeed, let \(f:\mathbf F\to\mathbf F'\) be a morphism of
\(\mathbf O\)-modules. Define the \(\mathbf O\)-modules
\(\operatorname{Ker}f\) (respectively,
\(\operatorname{Im}f\) and \(\operatorname{Coker}f\)) by setting, for every
\(S\in\operatorname{Ob}(\mathscr C)\),
\[
  (\operatorname{Ker}f)(S)=\operatorname{Ker}f(S)
\]
(respectively, the analogous formulas). Then \(\operatorname{Ker}f\)
(respectively, \(\operatorname{Coker}f\)) is a kernel (respectively, a
cokernel) of \(f\), and there is an isomorphism of \(\mathbf O\)-modules
\[
  \mathbf F/\operatorname{Ker}f
    \xrightarrow{\ \sim\ }\operatorname{Im}f.
\]
This proves that \((\mathbf O\text{-Mod.})\) is an abelian category.

Arbitrary direct sums exist and are defined objectwise. Finally, if
\(\mathbf M\) is an \(\mathbf O\)-module, \(\mathbf N\) a submodule, and
\((\mathbf F_i)_{i\in I}\) an increasing filtered family of submodules of
\(\mathbf M\), then the inclusion
\[
  \bigcup_{i\in I}(\mathbf F_i\cap\mathbf N)
    \subset
  \left(\bigcup_{i\in I}\mathbf F_i\right)\cap\mathbf N
\]
is an equality. Indeed, if \(S\in\operatorname{Ob}(\mathscr C)\) and
\[
  x\in\mathbf N(S)\cap\bigcup_i\mathbf F_i(S),
\]
there is an \(i\in I\) such that
\(x\in\mathbf N(S)\cap\mathbf F_i(S)\).

\medskip
\noindent\SGAInlineRefTarget{sga3:i:proposition:3-1-2}\textbf{Proposition 3.1.2.}\footnotemark[\value{footnote}]
\textit{Suppose that the category \(\mathscr C\) is small, that is, that
\(\operatorname{Ob}(\mathscr C)\) is a set. Then
\((\mathbf O\text{-Mod.})\) has a generator, and consequently has enough
injective objects.}
\footnotetext{Editor's note: Statements \SGARefLink{sga3:i:lemma:3-1-1}{3.1.1} and \SGARefLink{sga3:i:proposition:3-1-2}{3.1.2} have been added to
show that the category \((\mathbf O\text{-Mod.})\) is abelian and satisfies
axiom \(\mathrm{(AB\ 5)}\), and moreover has enough injective objects when
the category \(\mathscr C\) is small. An error in \SGARefLink{sga3:i:proposition:3-1-2}{3.1.2}, pointed out in 2016
by Linyuan Liu, is corrected here.}

\textit{Proof.}
For every object \(S\) of \(\mathscr C\), define the
\(\mathbf O\)-module \(\mathbf O_S\) as follows. For every object \(S'\),
\(\mathbf O_S(S')\) is a direct sum of copies of \(\mathbf O(S')\) indexed
by \(\operatorname{Hom}(S',S)\), that is,
\[
  \mathbf O_S(S')
    =\bigoplus_{g:S'\to S}\mathbf O(S')g.
\]
For every morphism \(f:S''\to S'\), writing \(f^*\) for the ring
homomorphism \(\mathbf O(S')\to\mathbf O(S'')\), the morphism
\[
  \mathbf O_S(f):\mathbf O_S(S')\longrightarrow\mathbf O_S(S'')
\]
sends each element \(ag\) of \(\mathbf O_S(S')\), where
\(a\in\mathbf O(S')\) and \(g:S'\to S\), to the element
\(f^*(a)(g\circ f)\) of \(\mathbf O_S(S'')\). It is readily checked that
this defines a functor in \(\mathbf O\)-modules.

Since \(\operatorname{Ob}(\mathscr C)\) is a set by hypothesis, we may form
the direct sum
\[
  \mathbf G=\bigoplus_{S\in\operatorname{Ob}(\mathscr C)}\mathbf O_S.
\]

\medskip
\noindent\SGAInlineRefTarget{sga3:i:lemma:3-1-2-1}\textbf{Lemma 3.1.2.1.}
\textit{Let \(\mathbf F\) be an \(\mathbf O\)-module. For every
\(S\in\operatorname{Ob}(\mathscr C)\), every element
\(x\in\mathbf F(S)\) defines a unique morphism of \(\mathbf O\)-modules
\[
  \mathbf x:\mathbf O_S\longrightarrow\mathbf F
\]
such that \(\mathbf x(1\operatorname{id}_S)=x\).}

\textit{Proof.}
Let \(\phi:\mathbf O_S\to\mathbf F\) be a morphism of
\(\mathbf O\)-modules such that
\(\phi(1\operatorname{id}_S)=x\). Let \(g:S'\to S\) and
\(a\in\mathbf O(S')\). The \(\mathbf O(S')\)-linearity and the
commutativity of the diagram

% Authority: Exp1-13oct24.pdf component/printed p. 19; combined reader p. 29.
\begin{center}
  % BEGIN INCLUDED SGA3_English_current_buildable_source_20260823/source/inputs/I/native_diagrams/exp1/diag_009_OS_g_naturality_square.tex
\begin{tikzcd}[column sep=4.5em,row sep=3.0em]
  \mathbf O_S(S)
    \arrow[r,"\phi"]
    \arrow[d,"\mathbf O_S(g)"']
  &
  \mathbf F(S)
    \arrow[d,"\mathbf F(g)"]
  \\
  \mathbf O_S(S')
    \arrow[r,"\phi"]
  &
  \mathbf F(S')
\end{tikzcd}

% END INCLUDED SGA3_English_current_buildable_source_20260823/source/inputs/I/native_diagrams/exp1/diag_009_OS_g_naturality_square.tex

\end{center}

imply that
\[
  \phi(ag)
    =a\phi(1g)
    =a\mathbf F(g)\bigl(\phi(1\operatorname{id}_S)\bigr)
    =a\mathbf F(g)(x).
\]
Thus \(\phi\), if it exists, is determined by the preceding equality.
Conversely, if \(\phi\) is defined in this way, then for every
\(f:S''\to S'\) one has
\begin{align*}
  \phi\circ\mathbf O_S(f)(ag)
    &=\phi\bigl(f^*(a)(g\circ f)\bigr)\\
    &=f^*(a)\mathbf F(g\circ f)(x)\\
    &=\mathbf F(f)\bigl(a\mathbf F(g)(x)\bigr)\\
    &=\mathbf F(f)\circ\phi(ag),
\end{align*}
that is, the diagram

% Authority: Exp1-13oct24.pdf component/printed p. 19; combined reader p. 29.
\begin{center}
  % BEGIN INCLUDED SGA3_English_current_buildable_source_20260823/source/inputs/I/native_diagrams/exp1/diag_010_OS_f_naturality_square.tex
\begin{tikzcd}[column sep=4.5em,row sep=3.0em]
  \mathbf O_S(S')
    \arrow[r,"\phi"]
    \arrow[d,"\mathbf O_S(f)"']
  &
  \mathbf F(S')
    \arrow[d,"\mathbf F(f)"]
  \\
  \mathbf O_S(S'')
    \arrow[r,"\phi"]
  &
  \mathbf F(S'')
\end{tikzcd}

% END INCLUDED SGA3_English_current_buildable_source_20260823/source/inputs/I/native_diagrams/exp1/diag_010_OS_f_naturality_square.tex

\end{center}

is commutative. This proves the lemma.

Now let \(\mathbf F\) be an \(\mathbf O\)-module. For each
\(S\in\operatorname{Ob}(\mathscr C)\), let \(F_0(S)\) be a set of
generators of the \(\mathbf O(S)\)-module \(\mathbf F(S)\), and let
\(\mathbf L_S\) be the direct sum, indexed by \(F_0(S)\), of copies of
\(\mathbf O_S\). By the lemma, we obtain a morphism of
\(\mathbf O\)-modules
\[
  \mathbf L_S\longrightarrow\mathbf F
\]
such that the morphism \(\mathbf L_S(S)\to\mathbf F(S)\) is surjective,
and hence is an epimorphism in the category of
\(\mathbf O(S)\)-modules.

Set
\[
  I=\bigsqcup_{S\in\operatorname{Ob}(\mathscr C)}F_0(S).
\]
We then obtain a morphism of \(\mathbf O\)-modules
\[
  \mathbf G^{\oplus I}
    \longrightarrow
  \bigoplus_{S\in\operatorname{Ob}(\mathscr C)}\mathbf L_S
    \longrightarrow
  \mathbf F
\]
which is an epimorphism objectwise, and hence an epimorphism of
\((\mathbf O\text{-Mod.})\). This shows that \(\mathbf G\) is a generator
of \((\mathbf O\text{-Mod.})\) (see \([\mathrm{Gr}57,\,1.9.1]\)). Since
\((\mathbf O\text{-Mod.})\) satisfies \(\mathrm{(AB\ 5)}\), it follows
from \([\mathrm{Gr}57,\,1.10.1]\) that
\((\mathbf O\text{-Mod.})\) has enough injective objects.

\medskip
\noindent\SGAInlineRefTarget{sga3:i:remark:3-1-3}\textbf{Remark 3.1.3.}
If \(\mathbf O_0\) is the \(\widehat{\mathscr C}\)-ring defined by
\[
  \mathbf O_0(S)=\mathbb Z
\]
(which must not be confused with the functor associated with the constant
object \(\mathbb Z\)), then the category of \(\mathbf O_0\)-modules is
isomorphic to the category of commutative
\(\widehat{\mathscr C}\)-groups.

\medskip
\noindent\SGAInlineRefTarget{sga3:i:definition:3-1-4}\textbf{Definition 3.1.4.}
Notice that if \(\mathbf F\) is an \(\mathbf O\)-module, then for every
\(S\in\operatorname{Ob}(\mathscr C)\), \(\mathbf F_S\) is an
\(\mathbf O_S\)-module. We may therefore define an abelian
\(\widehat{\mathscr C}\)-group
\(\underline{\operatorname{Hom}}_{\mathbf O}(\mathbf F,\mathbf F')\) by
\[
  \underline{\operatorname{Hom}}_{\mathbf O}(\mathbf F,\mathbf F')(S)
    =\operatorname{Hom}_{\mathbf O_S}(\mathbf F_S,\mathbf F'_S).
\]
In the same way one defines the objects
\[
  \underline{\operatorname{Isom}}_{\mathbf O}(\mathbf F,\mathbf F'),
  \qquad
  \underline{\operatorname{End}}_{\mathbf O}(\mathbf F),
  \qquad\text{and}\qquad
  \underline{\operatorname{Aut}}_{\mathbf O}(\mathbf F),
\]
the last of which is a \(\widehat{\mathscr C}\)-group under the structure
induced by composition of automorphisms.

\medskip
\noindent\SGAInlineRefTarget{sga3:i:definition:3-2}\textbf{Definition 3.2.}
Let \(\mathbf O\) be a \(\widehat{\mathscr C}\)-ring,
\(\mathbf F\) an \(\mathbf O\)-module, and \(\mathbf G\) a
\(\widehat{\mathscr C}\)-group. A \(\mathbf G\)-\(\mathbf O\)-module
structure on \(\mathbf F\) is a \(\mathbf G\)-object structure such that,
for every \(S\in\operatorname{Ob}(\mathscr C)\) and every
\(g\in\mathbf G(S)\), the automorphism of \(\mathbf F(S)\) defined by
\(g\) is an automorphism of its \(\mathbf O(S)\)-module structure.

Equivalently, the morphism of \(\widehat{\mathscr C}\)-groups
\[
  \rho:\mathbf G\longrightarrow
    \underline{\operatorname{Aut}}(\mathbf F)
\]
defined in \SGARefLink{sga3:i:subsection:2-3}{2.3} maps \(\mathbf G\) into the
\(\widehat{\mathscr C}\)-subgroup
\(\underline{\operatorname{Aut}}_{\mathbf O}(\mathbf F)\) of
\(\underline{\operatorname{Aut}}(\mathbf F)\). Thus, giving a
\(\mathbf G\)-\(\mathbf O\)-module structure on the
\(\mathbf O\)-module \(\mathbf F\) is the same as giving a morphism of
\(\widehat{\mathscr C}\)-groups
\[
  \rho:\mathbf G\longrightarrow
    \underline{\operatorname{Aut}}_{\mathbf O}(\mathbf F).
\]
One defines naturally the abelian group
\(\operatorname{Hom}_{\mathbf G\text{-}\mathbf O}
  (\mathbf F,\mathbf F')\), and hence the additive category of
\(\mathbf G\)-\(\mathbf O\)-modules, denoted
\((\mathbf G\text{-}\mathbf O\text{-Mod.})\).

\medskip
\noindent\SGAInlineRefTarget{sga3:i:remark:3-2-1}\textbf{Remark 3.2.1.}
The reader will notice that \((\mathbf G\text{-}\mathbf O\text{-Mod.})\)
may also be defined as the category of \(\mathbf O[\mathbf G]\)-modules,
where \(\mathbf O[\mathbf G]\) is the algebra of the
\(\widehat{\mathscr C}\)-group \(\mathbf G\) over the
\(\widehat{\mathscr C}\)-ring \(\mathbf O\), whose definition is clear.
\footnote{Editor's note: the sentence which follows has been added; it will
be used in \SGARefLink{sga3:i:section:5}{Section~5}.}
Consequently, by \SGARefLink{sga3:i:lemma:3-1-1}{3.1.1} and \SGARefLink{sga3:i:proposition:3-1-2}{3.1.2},
\((\mathbf G\text{-}\mathbf O\text{-Mod.})\) is an abelian category
satisfying axiom \(\mathrm{(AB\ 5)}\); moreover, if \(\mathscr C\) is
small, \((\mathbf G\text{-}\mathbf O\text{-Mod.})\) has enough injective
objects.

All the constructions of this section specialize immediately to the case
where \(\mathbf G\) (or \(\mathbf O\), or both) is representable by an
object of \(\mathscr C\), which is thereby endowed with the corresponding
algebraic structure.

We have treated briefly the principal algebraic structures that occur later
in this seminar. For the others (the structure of an \(\mathbf O\)-Lie
algebra, for example), we believe that the reader has now seen enough
examples in this section to make the general mechanism sketched in \SGARefLink{sga3:i:subsection:2-2}{2.2} work
in each particular case.

We shall now apply what we have just done to the category of schemes, denoted
\((\mathbf{Sch})\), and more generally to the categories
\((\mathbf{Sch})_{/S}\), also denoted \((\mathbf{Sch}_{/S})\).

% End of assigned component. Section 4 begins on Exp1 component/printed p. 21,
% combined-reader p. 31, in the next unit.

% END INCLUDED SGA3_English_current_buildable_source_20260823/source/inputs/I/tex/components/05_expose_I_s3_modules.tex

% BEGIN INCLUDED SGA3_English_current_buildable_source_20260823/source/inputs/I/tex/components/06_expose_I_s41_s43_basic_scheme_groups.tex
% SGA 3, Expose I, Loop 1 English reconstruction.
% Scope: Section 4 opening through Definition 4.3.3.1, stopping immediately
% before Section 4.4 (Diagonalizable groups).
% Authority: Polo--Gille Exp1-13oct24.pdf, component-local pages 21--23
% (combined-reader pages 31--33). Force-OCR and jcreinhold are drafting and
% comparison witnesses only.

% Source page: Exp1 component-local p. 21; combined-reader p. 31.
% Re-edition printed p. 21; original-edition marginal pp. 23--24 also occur
% on this physical page.
\section*{\SGAInlineRefTarget{sga3:i:section:4}4. Algebraic structures in the category of schemes}
\addcontentsline{toc}{section}{4. Algebraic structures in the category of schemes}

Whenever there is no ambiguity, we shall permit ourselves the following
abuses of language. We shall denote by \(T\) the object
\(T\xrightarrow{f}S\) of \(\mathscr C_{/S}\), with the datum of \(f\)
(the structural morphism of \(T\)) understood, and we shall identify
\(\mathscr C\) with a subcategory of \(\widehat{\mathscr C}\). In view of
the compatibilities stated in the preceding sections, these identifications
may be made without danger.

We shall also simplify terminology according to the following pattern: a
\((\mathbf{Sch})\)-group will also be called a \emph{group scheme}; a
\((\mathbf{Sch})_{/S}\)-group will be called a \emph{group scheme over
\(S\)}, an \emph{\(S\)-group scheme}, or an \emph{\(S\)-group}; and it
will be called an \emph{\(A\)-group} when \(S\) is the spectrum of the
ring \(A\).

\subsection*{\SGAInlineRefTarget{sga3:i:subsection:4-1}4.1. Constant schemes}

The category of schemes satisfies the conditions required in \SGARefLink{sga3:i:subsection:1:8}{1.8}. Constant
objects may therefore be defined in it. For every set \(E\), there is a
scheme \(E_{\mathbb Z}\), and for every scheme \(S\), an \(S\)-scheme
\[
  E_S=(E_{\mathbb Z})_S.
\]
Recall (cf. \emph{loc. cit.}) that for every \(S\)-scheme \(T\),
\[
  \operatorname{Hom}_S(T,E_S)
\]
is the set of locally constant maps from the topological space underlying
\(T\) to \(E\).

The functor \(E\mapsto E_S\) commutes with finite inverse limits. In
particular, if \(G\) is a group, then \(G_S\) is an \(S\)-group scheme;
if \(O\) is a ring, then \(O_S\) is an \(S\)-ring scheme; and similarly
for other algebraic structures.

\subsection*{\SGAInlineRefTarget{sga3:i:subsection:4-2}4.2. Affine \(S\)-groups}

Let us recall some facts about affine \(S\)-schemes (EGA~II, Section~1).
An \(S\)-scheme \(T\) is said to be \emph{affine over \(S\)} if the
inverse image of every affine open subset of \(S\) is affine. The
\(\mathcal O_S\)-algebra \(f_*\mathcal O_T\), denoted by
\(\mathscr A(T)\), is then quasi-coherent, where \(f\) denotes the
structural morphism of \(T\). Conversely, to every quasi-coherent
\(\mathcal O_S\)-algebra \(\mathscr A\), one can associate an
\(S\)-scheme affine over \(S\), denoted \(\operatorname{Spec}(\mathscr A)\).
The functors
\[
  T\longmapsto\mathscr A(T),
  \qquad
  \mathscr A\longmapsto\operatorname{Spec}(\mathscr A)
\]
are quasi-inverse equivalences between the category of \(S\)-schemes
affine over \(S\) and the category opposite to that of quasi-coherent
\(\mathcal O_S\)-algebras.

It follows that giving an algebraic structure on an \(S\)-scheme \(T\)
affine over \(S\) is equivalent to giving the corresponding structure on
\(\mathscr A(T)\) in the category opposite to that of quasi-coherent
\(\mathcal O_S\)-algebras. In particular, if \(G\) is an \(S\)-group
affine over \(S\), then \(\mathscr A(G)\) is endowed with the structure of
an augmented \(\mathcal O_S\)-bialgebra: there are two morphisms of
\(\mathcal O_S\)-algebras
\[
  \Delta:\mathscr A(G)\longrightarrow
    \mathscr A(G)\otimes_{\mathcal O_S}\mathscr A(G),
  \qquad
  \varepsilon:\mathscr A(G)\longrightarrow\mathcal O_S,
\]
corresponding to the morphisms of \(S\)-schemes
\[
  \pi:G\times G\longrightarrow G,
  \qquad
  e_G:S\longrightarrow G.
\]
The maps \(\Delta\) and \(\varepsilon\) satisfy the following conditions,
which express that \(G\) is an \(S\)-monoid.
\footnote{Editor's note: of course, the inversion morphism
\(G\to G\) induces a morphism of \(\mathcal O_S\)-algebras
\(\tau:\mathscr A(G)\to\mathscr A(G)\) which, together with
\(\Delta\) and \(\varepsilon\), makes \(\mathscr A(G)\) into a Hopf
\(\mathcal O_S\)-algebra.}

% Source page: Exp1 component-local p. 22; combined-reader p. 32.
% Re-edition printed p. 22; original-edition marginal p. 25 occurs below.
\medskip
\noindent\textbf{(HA 1)}
\(\Delta\) is coassociative: the following diagram is commutative.

\begin{figure}[H]
  \centering
  % BEGIN INCLUDED SGA3_English_current_buildable_source_20260823/source/inputs/I/native_diagrams/exp1/diag_011_hopf_coassociativity.tex
\begin{tikzcd}[column sep=5.5em,row sep=3.8em]
  &
  \mathscr A(G)\otimes_{\mathcal O_S}\mathscr A(G)
    \arrow[dr,"\operatorname{Id}\otimes\Delta"]
  &
  \\
  \mathscr A(G)
    \arrow[ur,"\Delta"]
    \arrow[dr,"\Delta"']
  &&
  \mathscr A(G)\otimes_{\mathcal O_S}\mathscr A(G)
    \otimes_{\mathcal O_S}\mathscr A(G)
  \\
  &
  \mathscr A(G)\otimes_{\mathcal O_S}\mathscr A(G)
    \arrow[ur,"\Delta\otimes\operatorname{Id}"']
  &
  {} .
\end{tikzcd}

% END INCLUDED SGA3_English_current_buildable_source_20260823/source/inputs/I/native_diagrams/exp1/diag_011_hopf_coassociativity.tex

  \phantomsection\label{fig:sga3-exp1-p022-hopf-coassociativity}
\end{figure}

\noindent\textbf{(HA 2)}
\(\Delta\) is compatible with \(\varepsilon\): the following two composites
are the identity:
\[
  \mathscr A(G)
    \xrightarrow{\Delta}
  \mathscr A(G)\otimes_{\mathcal O_S}\mathscr A(G)
    \xrightarrow{\operatorname{Id}\otimes\varepsilon}
  \mathscr A(G)\otimes_{\mathcal O_S}\mathcal O_S
    \xrightarrow{\ \sim\ }
  \mathscr A(G),
\]
\[
  \mathscr A(G)
    \xrightarrow{\Delta}
  \mathscr A(G)\otimes_{\mathcal O_S}\mathscr A(G)
    \xrightarrow{\varepsilon\otimes\operatorname{Id}}
  \mathcal O_S\otimes_{\mathcal O_S}\mathscr A(G)
    \xrightarrow{\ \sim\ }
  \mathscr A(G).
\]

Let us take this opportunity to note once again that it follows from the
definition of an \(S\)-group structure that, in order to give such a
structure on an \(S\)-scheme \(G\) affine over \(S\), it is not necessary
to verify anything on \(\mathscr A(G)\); one need only endow every
\(G(S')\), for \(S'\) over \(S\), with a group structure functorial in
\(S'\). This observation applies \emph{mutatis mutandis} to morphisms.

\subsection*{\SGAInlineRefTarget{sga3:i:subsection:4-3}4.3. The groups \(\mathbb G_a\) and \(\mathbb G_m\). The ring
\(\mathbf O\)}

\subsubsection*{4.3.1.}

Let \(\mathbb G_a\) be the functor from \((\mathbf{Sch})^{\circ}\) to
\(\mathbf{Set}\) defined by
\[
  \mathbb G_a(S)=\Gamma(S,\mathcal O_S),
\]
endowed with the \(\widehat{(\mathbf{Sch})}\)-group structure defined by
the additive-group structure of the ring \(\Gamma(S,\mathcal O_S)\). It
is representable by an affine scheme, which we denote by \(\mathbb G_a\),
and which is therefore a group scheme:
\[
  \mathbb G_a=\operatorname{Spec}\mathbb Z[T].
\]
Indeed,
\[
  \mathbb G_a(S)
    =\operatorname{Hom}(S,\mathbb G_a)
    =\operatorname{Hom}_{\mathrm{Alg.}}
       \bigl(\mathbb Z[T],\Gamma(S,\mathcal O_S)\bigr)
    \simeq\Gamma(S,\mathcal O_S).
\]

For every scheme \(S\), one therefore has an \(S\)-group affine over
\(S\), denoted \(\mathbb G_{a,S}\), corresponding to the
\(\mathcal O_S\)-bialgebra \(\mathcal O_S[T]\), whose diagonal map and
augmentation are defined by
\[
  \Delta(T)=T\otimes1+1\otimes T,
  \qquad
  \varepsilon(T)=0.
\]

\subsubsection*{4.3.2.}

Let \(\mathbb G_m\) be the functor from \((\mathbf{Sch})^{\circ}\) to
\(\mathbf{Set}\) defined by
\[
  \mathbb G_m(S)=\Gamma(S,\mathcal O_S)^{\times},
\]
% Source page: Exp1 component-local p. 23; combined-reader p. 33.
% Re-edition printed p. 23; original-edition marginal p. 26 occurs below.
where \(\Gamma(S,\mathcal O_S)^{\times}\) denotes the multiplicative
group of invertible elements of the ring \(\Gamma(S,\mathcal O_S)\),
endowed with its natural \(\widehat{(\mathbf{Sch})}\)-group structure.
It is representable by an affine scheme denoted by \(\mathbb G_m\):
\[
  \mathbb G_m
    =\operatorname{Spec}\mathbb Z[T,T^{-1}]
    =\operatorname{Spec}\mathbb Z[\mathbb Z],
\]
where \(\mathbb Z[\mathbb Z]\) denotes the group algebra of \(\mathbb Z\)
over the ring \(\mathbb Z\). Indeed,
\[
  \mathbb G_m(S)
    =\operatorname{Hom}_{\mathrm{Alg.}}
       \bigl(\mathbb Z[T,T^{-1}],\Gamma(S,\mathcal O_S)\bigr)
    \simeq\Gamma(S,\mathcal O_S)^{\times}.
\]

For every scheme \(S\), one therefore has an \(S\)-group affine over
\(S\), denoted \(\mathbb G_{m,S}\), corresponding to the
\(\mathcal O_S\)-algebra \(\mathcal O_S[\mathbb Z]\), whose diagonal map
and augmentation are defined by
\[
  \Delta(x)=x\otimes x,
  \qquad
  \varepsilon(x)=1,
  \qquad x\in\mathbb Z.
\]

\subsubsection*{4.3.3.}

Let \(\mathbf O\) be the functor \(\mathbb G_a\) endowed with its
\(\widehat{(\mathbf{Sch})}\)-ring structure. It is represented by the
scheme \(\operatorname{Spec}\mathbb Z[T]\), which we denote by
\(\mathbf O\) when it is regarded as endowed with its ring-scheme
structure.

For every scheme \(S\),
\[
  \mathbf O_S
    =S\times_{\operatorname{Spec}\mathbb Z}
       \operatorname{Spec}\mathbb Z[T]
    =\operatorname{Spec}(\mathcal O_S[T])
\]
is therefore an \(S\)-ring scheme affine over \(S\). (Note: in EGA~II,
1.7.13, \(\mathbf O_S\) is denoted by \(S[T]\).)

\medskip
\noindent\SGAInlineRefTarget{sga3:i:paragraph:4-3-3-1}\textbf{4.3.3.1.}
\footnote{Editor's note: this paragraph has been added; it will be useful
later (cf. II, \SGARefLink{sga3:i:subsection:1-3}{1.3}).}
For every object \(X\) of \(\widehat{(\mathbf{Sch})}\),
\[
  \mathbf O(X)=\operatorname{Hom}(X,\mathbf O)
\]
has a ring structure, functorial in \(X\). In particular, if \(X'\) is a
scheme and morphisms \(x:X'\to X\) and \(f:X\to\mathbf O\) are given
(that is, \(f\in\mathbf O(X)\)), then
\[
  f(x)=f\circ x
    \in\mathbf O(X')
    =\Gamma(X',\mathcal O_{X'}).
\]

\medskip
\noindent\textbf{Definition.}
Let \(\pi:M\to X\) be a morphism in \(\widehat{(\mathbf{Sch})}\), and set
\[
  \mathbf O_X=\mathbf O\times X.
\]
We say that \(M\) is an \(\mathbf O_X\)-module if, for every \(X\)-scheme
\(X'\), the set
\[
  \operatorname{Hom}_X(X',M)
\]
has been endowed with an \(\mathbf O(X')\)-module structure, functorial in
\(X'\).

This is equivalent to giving an \(X\)-abelian-group law
\[
  \mu:M\times_X M\longrightarrow M
\]
and an \emph{external law}
\[
  \mathbf O\times M
    =\mathbf O_X\times_X M
    \longrightarrow M,
  \qquad
  (f,m)\longmapsto f\cdot m,
\]
which is an \(X\)-morphism, that is,
\(\pi(f\cdot m)=\pi(m)\), and which, for every \(x\in X(S)\), endows
\[
  M(x)=\{m\in M(S)\mid\pi(m)=x\}
\]
with an \(\mathbf O(S)\)-module structure.

In that case, for every
\(Y\in\operatorname{Ob}\bigl(\widehat{(\mathbf{Sch})}_{/X}\bigr)\),
not necessarily representable,
\[
  \operatorname{Hom}_X(Y,M)=\Gamma(M_Y/Y)
\]
is an \(\mathbf O(Y)\)-module, functorially in \(Y\).

% End of assigned component. Section 4.4 (Diagonalizable groups) begins next,
% on Exp1 component-local p. 23 / combined-reader p. 33.

% END INCLUDED SGA3_English_current_buildable_source_20260823/source/inputs/I/tex/components/06_expose_I_s41_s43_basic_scheme_groups.tex

% BEGIN INCLUDED SGA3_English_current_buildable_source_20260823/source/inputs/I/tex/components/07_expose_I_s44_s47_scheme_groups.tex
% SGA 3, Expose I, Loop 1 English reconstruction.
% Scope: Section 4.4 through the end of Section 4.7, stopping immediately
% before Section 5 (Group cohomology).
% Authority: Polo--Gille Exp1-13oct24.pdf, component-local pages 23--31
% (combined-reader pages 33--41; p. 31 is included only through Section 4.7).
% Force-OCR and jcreinhold are drafting and comparison witnesses only.

% Source page: Exp1 component-local p. 23; combined-reader p. 33.
% Re-edition printed p. 23; Section 4.4 begins near the foot of the page.
\subsection*{\SGAInlineRefTarget{sga3:i:subsection:4-4}4.4. Diagonalizable groups}

\medskip
\noindent\SGAInlineRefTarget{sga3:i:paragraph:4-4-1}\textbf{4.4.1.}
The construction of \(\mathbb G_m\) generalizes as follows. Let \(M\)
be a commutative group, and let \(M_{\mathbb Z}\) be the group scheme
associated with it by the procedure of \SGARefLink{sga3:i:subsection:4-1}{4.1}. Consider the functor defined
by
\[
  \mathbf D(M)(S)
    =\operatorname{Hom}_{\mathrm{groups}}(M,\mathbb G_m(S))
    \simeq
      \underline{\operatorname{Hom}}_{S\text{-}\mathbf{Gr}}
        (M_S,\mathbb G_{m,S}).
\]

% Source page transition: Exp1 component-local p. 24;
% combined-reader p. 34; re-edition printed p. 24.
This is a commutative \(\widehat{(\mathbf{Sch})}\)-group represented by a
group scheme denoted \(D(M)\); thus, by definition,
\[
  D(M)\simeq
  \underline{\operatorname{Hom}}_{(\mathbf{Sch})\text{-}\mathbf{Gr}}
    (M_{\mathbb Z},\mathbb G_m).
\]
Indeed, set
\[
  D(M)=\operatorname{Spec}\mathbb Z[M],
\]
where \(\mathbb Z[M]\) is the group algebra of \(M\) over \(\mathbb Z\).
Then
\[
  D(M)(S)
    =\operatorname{Hom}_{\mathrm{Alg.}}
       \bigl(\mathbb Z[M],\Gamma(S,\mathcal O_S)\bigr)
    \simeq
      \operatorname{Hom}_{\mathrm{groups}}
       \bigl(M,\Gamma(S,\mathcal O_S)^{\times}\bigr)
\]
by the very definition of the algebra \(\mathbb Z[M]\).

\medskip
\noindent\SGAInlineRefTarget{sga3:i:paragraph:4-4-2}\textbf{4.4.2.}
For every scheme \(S\), one therefore has an \(S\)-group scheme affine
over \(S\),
\[
  D_S(M)=D(M)_S
    =\bigl(
      \underline{\operatorname{Hom}}_{(\mathbf{Sch})\text{-}\mathbf{Gr}}
        (M_{\mathbb Z},\mathbb G_m)
      \bigr)_S
    =\underline{\operatorname{Hom}}_{S\text{-}\mathbf{Gr}}
       (M_S,\mathbb G_{m,S}).
\]
It is associated with the \(\mathcal O_S\)-bialgebra
\(\mathcal O_S[M]\), endowed with the diagonal map and augmentation
defined by
\[
  \Delta(x)=x\otimes x,
  \qquad
  \varepsilon(x)=1,
  \qquad x\in M.
\]

\medskip
\noindent\SGAInlineRefTarget{sga3:i:paragraph:4-4-3}\textbf{4.4.3.}
If \(f:M\to N\) is a homomorphism of commutative groups, one defines in
the evident way a morphism of \(S\)-groups
\[
  D_S(f):D_S(N)\longrightarrow D_S(M),
\]
and hence a functor
\[
  D_S:M\longmapsto D_S(M)
\]
from the category of abelian groups to the category of \(S\)-groups
affine over \(S\). It may also be defined as the composite of the
functor \(M\mapsto M_S\) and the functor
\[
  M_S\longmapsto
  \underline{\operatorname{Hom}}_{S\text{-}\mathbf{Gr}}
    (M_S,\mathbb G_{m,S}).
\]
This functor commutes with base extension.

An \(S\)-group isomorphic to a group of the form \(D_S(M)\) is said to
be \emph{diagonalizable}. Notice that the elements of \(M\) are
interpreted as certain characters of \(D_S(M)\), that is, as certain
elements of
\[
  \operatorname{Hom}_{S\text{-}\mathbf{Gr}}
    (D_S(M),\mathbb G_{m,S}).
\]
(Compare VIII, 1.)

\medskip
\noindent\SGAInlineRefTarget{sga3:i:paragraph:4-4-4}\textbf{4.4.4.}
Let us give some examples of diagonalizable groups. First,
\[
  D(\mathbb Z)=\mathbb G_m,
  \qquad
  D(\mathbb Z^r)=(\mathbb G_m)^r.
\]
Set
\[
  \boldsymbol\mu_n=D(\mathbb Z/n\mathbb Z),
\]
and call it the \emph{group of \(n\)-th roots of unity}. Indeed,
\[
  \boldsymbol\mu_n(S)
    =\operatorname{Hom}_{\mathrm{groups}}
       \bigl(\mathbb Z/n\mathbb Z,\Gamma(S,\mathcal O_S)^{\times}\bigr)
    =\{f\in\Gamma(S,\mathcal O_S)\mid f^n=1\}.
\]
The \(S\)-group \(\boldsymbol\mu_{n,S}\) corresponds to the
\(\mathcal O_S\)-algebra \(\mathcal O_S[T]/(T^n-1)\). Suppose, in
particular, that \(S\) is the spectrum of a field \(k\) of
characteristic \(p=n\). On setting \(T-1=s\), one obtains
\[
  k[T]/(T^p-1)=k[s]/(s^p),
\]
which shows that the underlying topological space of
\(\boldsymbol\mu_{p,k}\) consists of one point, whose local ring is the
Artinian \(k\)-algebra \(k[s]/(s^p)\). In the same vein, note that the
\(S\)-schemes \(\mathbb G_{a,S}\), \(\mathbb G_{m,S}\), and
\(\mathbf O_S\) are smooth over \(S\), and that \(D_S(M)\) is

% Source page transition: Exp1 component-local p. 25;
% combined-reader p. 35; re-edition printed p. 25.
flat over \(S\). Moreover, it is formally smooth (respectively, smooth)
over \(S\) if and only if no residue characteristic of \(S\) divides
the torsion of \(M\) (respectively, and if in addition \(M\) is of
finite type); see \SGARefLink{sga3:viii:proposition:2:1}{Expos\'e VIII, 2.1}.

\subsection*{\SGAInlineRefTarget{sga3:i:subsection:4-5}4.5. Other examples of groups}

The preceding procedure applies to the ``classical groups'' (linear
groups \(\operatorname{GL}_n\), symplectic groups
\(\operatorname{Sp}_n\), orthogonal groups \(\operatorname O_n\), and
so forth). For example, one defines \(\operatorname{GL}_n\) as
representing the functor \(\mathbf{GL}_n\) such that
\[
  \mathbf{GL}_n(S)
    =\operatorname{GL}\bigl(n,\Gamma(S,\mathcal O_S)\bigr)
    =\operatorname{Aut}_{\mathcal O_S}(\mathcal O_S^n).
\]
It may be constructed, for example, as the open subscheme of
\(\operatorname{Spec}\mathbb Z[T_{ij}]\), \(1\leq i,j\leq n\), defined
by the function
\[
  f=\det\bigl((T_{ij})_{i,j=1}^n\bigr),
\]
that is, as \(\operatorname{Spec}\mathbb Z[T_{ij},f^{-1}]\).

\subsection*{\SGAInlineRefTarget{sga3:i:subsection:4-6}4.6. Functor-modules in the category of schemes}

We shall associate with every \(\mathcal O_S\)-module on a scheme \(S\)
a \(\mathbf O_S\)-module, where \(\mathbf O_S\) denotes the functor-ring
introduced in 4.3.3. This can be done in two different ways, as follows.

\medskip
\noindent\SGAInlineRefTarget{sga3:i:definition:4-6-1}\textbf{Definition 4.6.1.}
Let \(S\) be a scheme. For every \(\mathcal O_S\)-module \(\mathscr F\),
denote by \(\mathbf V(\mathscr F)\) and \(\mathbf W(\mathscr F)\) the
contravariant functors on \((\mathbf{Sch})_{/S}\) defined by
\begin{align*}
  \mathbf V(\mathscr F)(S')
    &=\operatorname{Hom}_{\mathcal O_{S'}}
       (\mathscr F\otimes_{\mathcal O_S}\mathcal O_{S'},
        \mathcal O_{S'}),\\
  \mathbf W(\mathscr F)(S')
    &=\Gamma(S',\mathscr F\otimes_{\mathcal O_S}\mathcal O_{S'}),
\end{align*}
where \(\mathscr F\otimes_{\mathcal O_S}\mathcal O_{S'}\) denotes the
inverse image on \(S'\) of the \(\mathcal O_S\)-module \(\mathscr F\).

Then \(\mathbf V(\mathscr F)\) and \(\mathbf W(\mathscr F)\) are
endowed in the evident way with \(\mathbf O_S\)-module structures
(recall that
\(\mathbf O_S(S')=\Gamma(S',\mathcal O_{S'})
=\mathbf W(\mathcal O_S)(S')\)). Thus one has in fact defined two
functors \(\mathbf V\) and \(\mathbf W\) from the category of
\(\mathcal O_S\)-modules to the category of \(\mathbf O_S\)-modules,
with \(\mathbf V\) contravariant and \(\mathbf W\) covariant.

For the remainder of this paragraph, we restrict to the case in which
the \(\mathcal O_S\)-modules in question are quasi-coherent. Thus we
regard \(\mathbf V\) and \(\mathbf W\) as functors
\[
  \mathbf V:(\mathcal O_S\text{-}\mathbf{Mod.q.c.})^{\circ}
    \longrightarrow(\mathbf O_S\text{-}\mathbf{Mod.}),
  \qquad
  \mathbf W:(\mathcal O_S\text{-}\mathbf{Mod.q.c.})
    \longrightarrow(\mathbf O_S\text{-}\mathbf{Mod.}).
\]

\medskip
\noindent\SGAInlineRefTarget{sga3:i:remark:4-6-1-1}\textbf{Remark 4.6.1.1.}
\footnote{Editor's note: the numbering \SGARefLink{sga3:i:remark:4-6-1-1}{4.6.1.1} has been added for later
references.}
The reader will notice that, in what follows, every proposition involving
only the functor \(\mathbf W\) remains valid for arbitrary
\(\mathcal O_S\)-modules, not necessarily quasi-coherent.

\medskip
\noindent\SGAInlineRefTarget{sga3:i:proposition:4-6-2}\textbf{Proposition 4.6.2.}
\begin{enumerate}
  \item[\SGAInlineRefTarget{sga3:i:proposition:4-6-2:item:i}(i)]
  The functors \(\mathbf V\) and \(\mathbf W\) commute with base
  extension: if \(S'\) lies over \(S\) and \(\mathscr F\) is a
  quasi-coherent \(\mathcal O_S\)-module, then
  \[
    \mathbf V(\mathscr F\otimes\mathcal O_{S'})
      \simeq\mathbf V(\mathscr F)_{S'},
    \qquad
    \mathbf W(\mathscr F\otimes\mathcal O_{S'})
      \simeq\mathbf W(\mathscr F)_{S'}.
  \]

% Source page transition: Exp1 component-local p. 26;
% combined-reader p. 36; re-edition printed p. 26.
  \item[\SGAInlineRefTarget{sga3:i:proposition:4-6-2:item:ii}(ii)]
  The functors \(\mathbf V\) and \(\mathbf W\) are fully faithful: the
  canonical maps
  \[
    \operatorname{Hom}_{\mathcal O_S}(\mathscr F,\mathscr F')
      \longrightarrow
    \operatorname{Hom}_{\mathbf O_S}
      (\mathbf V(\mathscr F'),\mathbf V(\mathscr F)),
  \]
  \[
    \operatorname{Hom}_{\mathcal O_S}(\mathscr F,\mathscr F')
      \longrightarrow
    \operatorname{Hom}_{\mathbf O_S}
      (\mathbf W(\mathscr F),\mathbf W(\mathscr F'))
  \]
  are bijective.

  \item[\SGAInlineRefTarget{sga3:i:proposition:4-6-2:item:iii}(iii)]
  The functors \(\mathbf V\) and \(\mathbf W\) are additive:
  \[
    \mathbf V(\mathscr F\oplus\mathscr F')
      \simeq
    \mathbf V(\mathscr F)\times_S\mathbf V(\mathscr F'),
    \qquad
    \mathbf W(\mathscr F\oplus\mathscr F')
      \simeq
    \mathbf W(\mathscr F)\times_S\mathbf W(\mathscr F').
  \]
\end{enumerate}

Parts \SGARefLink{sga3:i:proposition:4-6-2:item:i}{(i)} and \SGARefLink{sga3:i:proposition:4-6-2:item:iii}{(iii)} are immediate from the definitions. For \SGARefLink{sga3:i:proposition:4-6-2:item:ii}{(ii)}, take
the \(S'\) to be open subschemes of \(S\). We leave the proof to the
reader (for \(\mathbf V\), use EGA~II, 1.7.14).

Recall (cf. \SGARefLink{sga3:i:definition:3-1-4}{3.1.4}) that if \(\mathbf F\) and \(\mathbf F'\) are
\(\mathbf O_S\)-modules, then
\(\underline{\operatorname{Hom}}_{\mathbf O_S}(\mathbf F,\mathbf F')\)
denotes the \(S\)-functor in abelian groups that assigns to every
\(S'\to S\) the group
\[
  \operatorname{Hom}_{\mathbf O_{S'}}(\mathbf F_{S'},\mathbf F'_{S'}).
\]

\medskip
\noindent\SGAInlineRefTarget{sga3:i:proposition:4-6-3}\textbf{Proposition 4.6.3.}
\textit{There are canonical morphisms in
\((\mathbf O_S\text{-}\mathbf{Mod.})\):}
\footnote{Editor's note: the isomorphism
\(\underline{\operatorname{Hom}}_{\mathbf O_S}
(\mathbf W(\mathscr F),\mathbf W(\mathscr F'))
\simeq
\underline{\operatorname{Hom}}_{\mathbf O_S}
(\mathbf V(\mathscr F'),\mathbf V(\mathscr F))\)
has been added.}

\begin{figure}[H]
  \centering
  % BEGIN INCLUDED SGA3_English_current_buildable_source_20260823/source/inputs/I/native_diagrams/exp1/diag_012_prop463_canonical_morphisms.tex
\begin{tikzcd}[column sep=3.2em,row sep=3.5em]
  \underline{\operatorname{Hom}}_{\mathbf O_S}
    \bigl(\mathbf W(\mathscr F),\mathbf W(\mathscr F')\bigr)
    \arrow[rr,equal]
  &&
  \underline{\operatorname{Hom}}_{\mathbf O_S}
    \bigl(\mathbf V(\mathscr F'),\mathbf V(\mathscr F)\bigr)
  \\
  &
  \mathbf W\!\left(
    \mathscr Hom_{\mathcal O_S}(\mathscr F,\mathscr F')
  \right)
    \arrow[ul]
    \arrow[ur]
  &
  {} .
\end{tikzcd}

% END INCLUDED SGA3_English_current_buildable_source_20260823/source/inputs/I/native_diagrams/exp1/diag_012_prop463_canonical_morphisms.tex

  \phantomsection\label{fig:sga3-exp1-p026-prop463-canonical-morphisms}
\end{figure}

This follows immediately from \SGARefLink{sga3:i:proposition:4-6-2:item:i}{4.6.2(i)} and \SGARefLink{sga3:i:proposition:4-6-2:item:ii}{(ii)}.

\medskip
\noindent\SGAInlineRefTarget{sga3:i:notation:4-6-3-1}\textbf{Notation 4.6.3.1.}
Let \(\mathscr F\) be a quasi-coherent \(\mathcal O_S\)-module. By
EGA~II, 1.7.8, the \(S\)-functor \(\mathbf V(\mathscr F)\) is
represented by an \(S\)-scheme affine over \(S\), denoted
\(\mathbb V(\mathscr F)\), and called the \emph{vector fibration}
\footnote{Editor's note: ``vector bundle'' has been replaced by ``vector
fibration.'' Current usage calls a vector fibration that is locally
trivial of rank \(r\) a ``vector bundle of rank \(r\),'' that is, one
whose sheaf of sections is locally isomorphic to
\(\mathcal O_S^{\oplus r}\).}
defined by \(\mathscr F\):
\[
  \mathbb V(\mathscr F)=\operatorname{Spec}(\mathcal S(\mathscr F)),
\]
where \(\mathcal S(\mathscr F)\) denotes the symmetric algebra of the
\(\mathcal O_S\)-module \(\mathscr F\).
\footnote{Editor's note: the articles [Ni04] and [Ni02], respectively,
show that if \(S\) is Noetherian and \(\mathscr F\) is a coherent
\(\mathcal O_S\)-module, then \(\mathbf W(\mathscr F)\), respectively
the \(S\)-group assigning
\(\operatorname{Aut}_{\mathcal O_T}
(\mathscr F\otimes\mathcal O_T)\) to every \(T\to S\), is representable
if and only if \(\mathscr F\) is locally free.}

\medskip
\noindent\SGAInlineRefTarget{sga3:i:proposition:4-6-4}\textbf{Proposition 4.6.4.}
\textit{Let \(\mathscr F\) and \(\mathscr F'\) be two quasi-coherent
\(\mathcal O_S\)-modules, and let \(\mathscr A\) be a quasi-coherent
\(\mathcal O_S\)-algebra. There is a functorial isomorphism
\[
  \operatorname{Hom}_S
   \bigl(\operatorname{Spec}(\mathscr A),
     \underline{\operatorname{Hom}}_{\mathbf O_S}
       (\mathbf W(\mathscr F'),\mathbf W(\mathscr F))\bigr)
  \xrightarrow{\ \sim\ }
  \operatorname{Hom}_{\mathcal O_S}
    (\mathscr F',\mathscr F\otimes_{\mathcal O_S}\mathscr A).
\]}

% Source page transition: Exp1 component-local p. 27;
% combined-reader p. 37; re-edition printed p. 27.
Indeed, put \(X=\operatorname{Spec}(\mathscr A)\). The left-hand side is
canonically isomorphic to
\[
  \underline{\operatorname{Hom}}_{\mathbf O_S}
    (\mathbf W(\mathscr F'),\mathbf W(\mathscr F))(X),
\]
that is, by definition, to
\begin{align*}
  &\operatorname{Hom}_{\mathbf O_X}
    (\mathbf W(\mathscr F')_X,\mathbf W(\mathscr F)_X)\\
  &\qquad\simeq
  \operatorname{Hom}_{\mathbf O_X}
    \bigl(\mathbf W(\mathscr F'\otimes\mathcal O_X),
          \mathbf W(\mathscr F\otimes\mathcal O_X)\bigr)
\end{align*}
(cf. \SGARefLink{sga3:i:proposition:4-6-2:item:i}{4.6.2(i)}). By \SGARefLink{sga3:i:proposition:4-6-2:item:ii}{4.6.2(ii)}, this may also be written as
\[
  \operatorname{Hom}_{\mathcal O_X}
    (\mathscr F'\otimes\mathcal O_X,\mathscr F\otimes\mathcal O_X)
  =
  \operatorname{Hom}_{\mathcal O_S}
    (\mathscr F',\pi_*\pi^*(\mathscr F)),
\]
where \(\pi:X\to S\) is the structural morphism. But EGA~II, 1.4.7
gives
\[
  \pi_*\pi^*(\mathscr F)\simeq\mathscr F\otimes\mathscr A,
\]
which completes the proof.

\medskip
\noindent\SGAInlineRefTarget{sga3:i:corollary:4-6-4-1}\textbf{Corollary 4.6.4.1.}
\textit{There is a canonical isomorphism
\[
  \mathbf W(\mathscr F\otimes\mathscr A)
    \simeq
  \underline{\operatorname{Hom}}_S
    (\operatorname{Spec}(\mathscr A),\mathbf W(\mathscr F)).
\]}

Indeed,
\footnote{Editor's note: the original has been expanded in what follows.}
let \(f:S'\to S\) be an \(S\)-scheme and put
\(X'=X\times_S S'\). There is a Cartesian square

\begin{figure}[H]
  \centering
  % BEGIN INCLUDED SGA3_English_current_buildable_source_20260823/source/inputs/I/native_diagrams/exp1/diag_013_cor4641_cartesian_square.tex
\begin{tikzcd}[column sep=4.0em,row sep=3.0em]
  X'
    \arrow[r,"f'"]
    \arrow[d,"\pi'"']
  &
  X
    \arrow[d,"\pi"]
  \\
  S'
    \arrow[r,"f"]
  &
  S
\end{tikzcd}

% END INCLUDED SGA3_English_current_buildable_source_20260823/source/inputs/I/native_diagrams/exp1/diag_013_cor4641_cartesian_square.tex

  \phantomsection\label{fig:sga3-exp1-p027-cor4641-cartesian-square}
\end{figure}

By EGA~II, 1.5.2, \(X'\) is affine over \(S'\) and
\(\pi'_*\mathcal O_{X'}=f^*(\mathscr A)\). Hence
\begin{align*}
  &\underline{\operatorname{Hom}}_S
     (\operatorname{Spec}(\mathscr A),\mathbf W(\mathscr F))(S')\\
  &\qquad=
  \operatorname{Hom}_{S'}
     \bigl(\operatorname{Spec}(f^*\mathscr A),
           \mathbf W(f^*\mathscr F)\bigr).
\end{align*}
By \SGARefLink{sga3:i:proposition:4-6-4}{4.6.4}, applied to \(f^*\mathscr F\),
\(\mathscr F'=\mathcal O_{S'}\), and \(f^*\mathscr A\), this is equal
to
\begin{align*}
  \Gamma(S',f^*\mathscr F\otimes f^*\mathscr A)
    &=\Gamma(S',f^*(\mathscr F\otimes\mathscr A))\\
    &=\mathbf W(\mathscr F\otimes\mathscr A)(S').
\end{align*}

\medskip
\noindent\SGAInlineRefTarget{sga3:i:proposition:4-6-5}\textbf{Proposition 4.6.5.}
\textit{If \(\mathscr F\) and \(\mathscr F'\) are two
\(\mathcal O_S\)-modules locally free of finite type, the morphisms of
\SGARefLink{sga3:i:proposition:4-6-3}{4.6.3} are isomorphisms.}

Indeed, for every \(S'\to S\),
\begin{align*}
  &\mathbf W
    \bigl(\mathscr{H}om_{\mathcal O_S}(\mathscr F,\mathscr F')\bigr)(S')\\
  &\quad=
    \Gamma\bigl(S',
      \mathscr{H}om_{\mathcal O_S}(\mathscr F,\mathscr F')
        \otimes\mathcal O_{S'}\bigr)\\
  &\quad=
    \operatorname{Hom}_{\mathcal O_{S'}}
      (\mathscr F\otimes\mathcal O_{S'},
       \mathscr F'\otimes\mathcal O_{S'}).
\end{align*}
The last member is isomorphic both to
\[
  \underline{\operatorname{Hom}}_{\mathbf O_S}
    (\mathbf W(\mathscr F),\mathbf W(\mathscr F'))(S')
\]
and to
\[
  \underline{\operatorname{Hom}}_{\mathbf O_S}
    (\mathbf V(\mathscr F'),\mathbf V(\mathscr F))(S')
\]
by \SGARefLink{sga3:i:proposition:4-6-2:item:i}{4.6.2(i)} and \SGARefLink{sga3:i:proposition:4-6-2:item:ii}{(ii)}.

\medskip
\noindent\SGAInlineRefTarget{sga3:i:corollary:4-6-5-1}\textbf{Corollary 4.6.5.1.}
\textit{Let \(\mathscr F\) be an \(\mathcal O_S\)-module locally free
of finite type. Put
\(\mathscr F^{\vee}
=\mathscr{H}om_{\mathcal O_S}(\mathscr F,\mathcal O_S)\). There are
canonical isomorphisms
\[
  \mathbf W(\mathscr F^{\vee})
    \simeq
  \underline{\operatorname{Hom}}_{\mathbf O_S}
    (\mathbf W(\mathscr F),\mathbf O_S)
    \simeq\mathbf V(\mathscr F),
\]
\[
  \mathbf V(\mathscr F^{\vee})
    \simeq
  \underline{\operatorname{Hom}}_{\mathbf O_S}
    (\mathbf V(\mathscr F),\mathbf O_S)
    \simeq\mathbf W(\mathscr F).
\]}

Finally, we have the following proposition.

\medskip
\noindent\SGAInlineRefTarget{sga3:i:proposition:4-6-6}\textbf{Proposition 4.6.6.}
\textit{Let \(f:\mathscr F\to\mathscr F'\) be a morphism of
\(\mathcal O_S\)-modules locally free of finite rank. The morphism
\[
  \mathbf W(f):\mathbf W(\mathscr F)\longrightarrow\mathbf W(\mathscr F')
\]
is a monomorphism if and only if \(f\) identifies \(\mathscr F\)
locally with a direct summand of \(\mathscr F'\).}

% Source page transition: Exp1 component-local p. 28;
% combined-reader p. 38; re-edition printed p. 28.
The direct assertion is essentially contained in EGA~0\textsubscript{I},
5.5.5.
\footnote{Editor's note: the following sentence has been expanded.}
Conversely, if \(\mathscr F\) is locally a direct summand of
\(\mathscr F'\), then for every \(\pi:S'\to S\), the module
\(\pi^*\mathscr F\) is a submodule of \(\pi^*\mathscr F'\), since this
holds locally as a direct summand. Thus
\[
  \mathbf W(\mathscr F)(S')=\Gamma(S',\pi^*\mathscr F)
\]
is a submodule of
\[
  \mathbf W(\mathscr F')(S')=\Gamma(S',\pi^*\mathscr F').
\]

\subsection*{\SGAInlineRefTarget{sga3:i:subsection:4-7}4.7. The category of \(G\)-\(\mathcal O_S\)-modules}

Let \(G\) be an \(S\)-group and \(\mathscr F\) an
\(\mathcal O_S\)-module; \(\mathbf W(\mathscr F)\) has a
\(\mathbf O_S\)-module structure.

\medskip
\noindent\SGAInlineRefTarget{sga3:i:definition:4-7-1}\textbf{Definition 4.7.1.}
A \emph{\(G\)-\(\mathcal O_S\)-module structure} on \(\mathscr F\) is
an \(\mathbf h_G\)-\(\mathbf O_S\)-module structure on
\(\mathbf W(\mathscr F)\) (cf. \SGARefLink{sga3:i:definition:3-2}{3.2}). A morphism of
\(G\)-\(\mathcal O_S\)-modules is, by definition, a morphism of the
associated \(\mathbf h_G\)-\(\mathbf O_S\)-modules. This gives the
category \((G\text{-}\mathcal O_S\text{-}\mathbf{Mod.})\); we denote by
\((G\text{-}\mathcal O_S\text{-}\mathbf{Mod.q.c.})\) its full
subcategory consisting of those \(G\)-\(\mathcal O_S\)-modules that are
quasi-coherent as \(\mathcal O_S\)-modules.

By \SGARefLink{sga3:i:definition:3-2}{3.2}, giving a \(G\)-\(\mathcal O_S\)-module structure on
\(\mathscr F\) is therefore equivalent to giving a morphism of
\(\widehat{(\mathbf{Sch})}_{/S}\)-groups
\[
  \rho:\mathbf h_G\longrightarrow
    \underline{\operatorname{Aut}}_{\mathbf O_S}
      (\mathbf W(\mathscr F)).
\]

\medskip
\noindent\SGAInlineRefTarget{sga3:i:remark:4-7-1-0}\textbf{Remark 4.7.1.0.}
\footnote{Editor's note: this remark has been added from [DG], II,
\SGARefLink{sga3:ii:section:2}{Section 2}, \SGARefLink{sga3:i:subsection:1-1}{1.1}.}
Since \SGARefLink{sga3:i:proposition:4-6-2}{4.6.2} gives an anti-isomorphism of \(S\)-functors in groups
\[
  (\dagger)\qquad
  \underline{\operatorname{Aut}}_{\mathbf O_S}
    (\mathbf W(\mathscr F))
  \simeq
  \underline{\operatorname{Aut}}_{\mathbf O_S}
    (\mathbf V(\mathscr F)),
\]
it is equivalent to give an \(\mathbf h_G\)-\(\mathbf O_S\)-module
structure on \(\mathbf W(\mathscr F)\) or on \(\mathbf V(\mathscr F)\).
Indeed, let
\[
  \rho:\mathbf h_G\longrightarrow
  \underline{\operatorname{Aut}}_{\mathbf O_S}
    (\mathbf W(\mathscr F))
\]
be as above. For every \(T\to S\) and \(g\in G(T)\), denote by
\(\rho^*(g)\) the image of \(\rho(g)\) under the anti-isomorphism
\((\dagger)\). Then
\[
  \rho^*(gh)=\rho^*(h)\circ\rho^*(g),
\]
so \(\rho^*\) defines a \emph{right}
\(\mathbf h_G\)-\(\mathbf O_S\)-module structure on
\(\mathbf V(\mathscr F)\). Setting
\[
  \rho^{\vee}(g)=\rho^*(g^{-1}),
\]
one obtains an \(\mathbf h_G\)-\(\mathbf O_S\)-module structure on
\(\mathbf V(\mathscr F)\), whose datum is equivalent to that of
\(\rho\).

\medskip
\noindent\SGAInlineRefTarget{sga3:i:remark:4-7-1-1}\textbf{Remark 4.7.1.1.}
One may say that the categories just defined have been constructed by
the following Cartesian squares.

\begin{figure}[H]
  \centering
  % BEGIN INCLUDED SGA3_English_current_buildable_source_20260823/source/inputs/I/native_diagrams/exp1/diag_014_rem4711_category_squares.tex
\begin{tikzcd}[column sep=3.3em,row sep=3.2em]
  (G\text{-}\mathcal O_S\text{-}\mathbf{Mod.q.c.})
    \arrow[r,hook]
    \arrow[d]
  &
  (G\text{-}\mathcal O_S\text{-}\mathbf{Mod.})
    \arrow[r]
    \arrow[d]
  &
  (\mathbf h_G\text{-}\mathbf O_S\text{-}\mathbf{Mod.})
    \arrow[d,"\mathrm{forget}"]
  \\
  (\mathcal O_S\text{-}\mathbf{Mod.q.c.})
    \arrow[r,hook]
  &
  (\mathcal O_S\text{-}\mathbf{Mod.})
    \arrow[r,"\mathbf W"]
  &
  (\mathbf O_S\text{-}\mathbf{Mod.})
  \qquad .
\end{tikzcd}

% END INCLUDED SGA3_English_current_buildable_source_20260823/source/inputs/I/native_diagrams/exp1/diag_014_rem4711_category_squares.tex

  \phantomsection\label{fig:sga3-exp1-p028-rem4711-category-squares}
\end{figure}

The categories \((\mathcal O_S\text{-}\mathbf{Mod.})\) and
\((\mathbf O_S\text{-}\mathbf{Mod.})\) are abelian, but one should
remember that in general the functor \(\mathbf W\) is neither left nor
right exact.
\footnote{Editor's note: the original has been corrected by deleting the
assertion that the category
\((G\text{-}\mathcal O_S\text{-}\mathbf{Mod.})\) is abelian; see
\SGARefLink{sga3:i:corollary:4-7-2-1}{4.7.2.1} below.}

% Source page transition: Exp1 component-local p. 29;
% combined-reader p. 39; re-edition printed p. 29.
\medskip
\noindent\SGAInlineRefTarget{sga3:i:remark:4-7-1-2}\textbf{Remark 4.7.1.2.}
\footnote{Editor's note: this remark has been added.}
Let \(\mathscr F\) be a \(G\)-\(\mathcal O_S\)-module. The subsheaf of
invariants \(\mathscr F^G\) is defined as follows: for every open subset
\(U\) of \(S\),
\[
  \mathscr F^G(U)=\mathbf W(\mathscr F)^G(U)
   =\bigl\{x\in\mathscr F(U)\mid
       g\cdot x_{S'}=x_{S'}
       \text{ for every }S'\xrightarrow{f}U,
       \ g\in G(S')\bigr\},
\]
where \(x_{S'}\) denotes the image of \(x\) in
\[
  \Gamma(S',f^*\mathscr F)=\Gamma(U,f_*f^*\mathscr F).
\]
One should note that the natural morphism
\[
  \mathbf W(\mathscr F^G)\longrightarrow\mathbf W(\mathscr F)^G
\]
is not an isomorphism in general. For example, let
\(S=\operatorname{Spec}(\mathbb Z)\), and let \(G\) be the constant
group \(\mathbb Z/2\mathbb Z=\{1,\tau\}\) acting on
\(\mathscr F=\mathcal O_S\) by \(\tau\cdot1=-1\). Then
\(\mathscr F^G=0\), whereas, if \(R\) is an \(\mathbb F_2\)-algebra,
\[
  \mathbf W(\mathscr F)^G(\operatorname{Spec}(R))=R.
\]

\medskip
\noindent\SGAInlineRefTarget{sga3:i:paragraph:4-7-2}\textbf{4.7.2.}
From now until the end of \SGARefLink{sga3:i:subsection:4-7}{4.7}, suppose that \(G\) is affine over
\(S\).
\footnote{Editor's note: see VI\textsubscript{B}, \SGARefLink{lem:sga3-VIB-11.1}{Sections 11.1}--\SGARefLink{par:sga3-VIB-11.6}{11.6}
for the extension of the results of \SGARefLink{sga3:i:proposition:4-7-2}{4.7.2} to the case where \(G\) is not
necessarily affine but \(G\) and \(\mathscr F\) are assumed flat over
\(S\).}
By \SGARefLink{sga3:i:proposition:4-6-4}{4.6.4}, giving a morphism of \(S\)-functors
\[
  \rho:\mathbf h_G\longrightarrow
    \underline{\operatorname{End}}_{\mathbf O_S}
      (\mathbf W(\mathscr F))
\]
is equivalent to giving a morphism of \(\mathcal O_S\)-modules
\[
  \mu:\mathscr F\longrightarrow
    \mathscr F\otimes_{\mathcal O_S}\mathscr A(G).
\]
Saying that \(\rho\) is a morphism of
\(\widehat{(\mathbf{Sch})}_{/S}\)-groups is then equivalent to saying
that \(\mu\) satisfies the following axioms.

\medskip
\noindent\textbf{(CM 1)} The following diagram is commutative.

\begin{figure}[H]
  \centering
  % BEGIN INCLUDED SGA3_English_current_buildable_source_20260823/source/inputs/I/native_diagrams/exp1/diag_015_cm1_coaction_square.tex
\begin{tikzcd}[column sep=5.3em,row sep=3.5em]
  \mathscr F
    \arrow[r,"\mu"]
    \arrow[d,"\mu"']
  &
  \mathscr F\otimes_{\mathcal O_S}\mathscr A(G)
    \arrow[d,"\operatorname{Id}\otimes\Delta"]
  \\
  \mathscr F\otimes_{\mathcal O_S}\mathscr A(G)
    \arrow[r,"\mu\otimes\operatorname{Id}"]
  &
  \mathscr F\otimes_{\mathcal O_S}\mathscr A(G)
    \otimes_{\mathcal O_S}\mathscr A(G)
  \qquad .
\end{tikzcd}

% END INCLUDED SGA3_English_current_buildable_source_20260823/source/inputs/I/native_diagrams/exp1/diag_015_cm1_coaction_square.tex

  \phantomsection\label{fig:sga3-exp1-p029-cm1-coaction-square}
\end{figure}

\noindent\textbf{(CM 2)} The following composite is the identity.

\begin{figure}[H]
  \centering
  % BEGIN INCLUDED SGA3_English_current_buildable_source_20260823/source/inputs/I/native_diagrams/exp1/diag_016_cm2_counit_composite.tex
\begin{tikzcd}[column sep=4.5em]
  \mathscr F
    \arrow[r,"\mu"]
  &
  \mathscr F\otimes\mathscr A(G)
    \arrow[r,"\operatorname{Id}\otimes\varepsilon"]
  &
  \mathscr F\otimes\mathcal O_S
    \arrow[r,"\sim"]
  &
  \mathscr F
  \qquad .
\end{tikzcd}

% END INCLUDED SGA3_English_current_buildable_source_20260823/source/inputs/I/native_diagrams/exp1/diag_016_cm2_counit_composite.tex

  \phantomsection\label{fig:sga3-exp1-p029-cm2-counit-composite}
\end{figure}

These axioms (CM 1) and (CM 2) are those of a right \emph{comodule}
structure over the bialgebra \(\mathscr A(G)\).
\footnote{Editor's note: left \(G\)-\(\mathcal O_S\)-modules correspond
naturally to right \(\mathscr A(G)\)-comodules.}

Put \(\mathscr A=\mathscr A(G)\). If \(\mathscr F\) and
\(\mathscr F'\) are \(\mathscr A\)-comodules, a morphism of comodules
\(f:\mathscr F\to\mathscr F'\) is a morphism of
\(\mathcal O_S\)-modules such that the following diagram is

% Source page transition: Exp1 component-local p. 30;
% combined-reader p. 40; re-edition printed p. 30.
commutative.

\begin{figure}[H]
  \centering
  % BEGIN INCLUDED SGA3_English_current_buildable_source_20260823/source/inputs/I/native_diagrams/exp1/diag_017_comodule_morphism_square.tex
\begin{tikzcd}[column sep=4.5em,row sep=3.0em]
  \mathscr F
    \arrow[r,"f"]
    \arrow[d,"\mu_{\mathscr F}"']
  &
  \mathscr F'
    \arrow[d,"\mu_{\mathscr F'}"]
  \\
  \mathscr F\otimes\mathscr A
    \arrow[r,"f\otimes\operatorname{Id}"]
  &
  \mathscr F'\otimes\mathscr A
  \qquad .
\end{tikzcd}

% END INCLUDED SGA3_English_current_buildable_source_20260823/source/inputs/I/native_diagrams/exp1/diag_017_comodule_morphism_square.tex

  \phantomsection\label{fig:sga3-exp1-p030-comodule-morphism-square}
\end{figure}

This gives the category \((\mathscr A\text{-}\mathbf{Comod.})\); we
denote by \((\mathscr A\text{-}\mathbf{Comod.q.c.})\) its full
subcategory consisting of those \(\mathscr A\)-comodules that are
quasi-coherent as \(\mathcal O_S\)-modules. We have therefore proved the
following statement.

\medskip
\noindent\SGAInlineRefTarget{sga3:i:proposition:4-7-2}\textbf{Proposition 4.7.2.}
\textit{Let \(G\) be an \(S\)-group affine over \(S\). There are
equivalences of categories
\[
  (G\text{-}\mathcal O_S\text{-}\mathbf{Mod.})
    \simeq(\mathscr A(G)\text{-}\mathbf{Comod.}),
\]
\[
  (G\text{-}\mathcal O_S\text{-}\mathbf{Mod.q.c.})
    \simeq(\mathscr A(G)\text{-}\mathbf{Comod.q.c.}).
\]}

If, in addition, \(S=\operatorname{Spec}(\Lambda)\) is affine and one
sets \(\Lambda[G]=\Gamma(S,\mathscr A(G))\), there is an equivalence of
categories
\[
  (\mathscr A(G)\text{-}\mathbf{Comod.q.c.})
    \simeq(\Lambda[G]\text{-}\mathbf{Comod.}).
\]
\footnote{Editor's note: the preceding sentence has been added.}

Suppose, moreover, that \(\mathscr A=\mathscr A(G)\) is a flat
\(\mathcal O_S\)-module.
\footnote{Editor's note: the following paragraph has been added from
[Se68, \SGARefLink{sga3:i:subsection:1-3}{Section 1.3}].}
Let \(\mathscr E\) be an \(\mathscr A\)-comodule and \(\mathscr F\) an
\(\mathcal O_S\)-submodule of \(\mathscr E\). Since \(\mathscr A\) is
flat over \(\mathcal O_S\), one may identify
\(\mathscr F\otimes\mathscr A\) (respectively,
\(\mathscr F\otimes\mathscr A\otimes\mathscr A\)) with an
\(\mathcal O_S\)-submodule of \(\mathscr E\otimes\mathscr A\)
(respectively, of
\(\mathscr E\otimes\mathscr A\otimes\mathscr A\)). Suppose that
\(\mu_{\mathscr E}\) maps \(\mathscr F\) into
\(\mathscr F\otimes\mathscr A\). Then its restriction
\[
  \mu_{\mathscr F}:\mathscr F\longrightarrow
    \mathscr F\otimes\mathscr A
\]
endows \(\mathscr F\) with an \(\mathscr A\)-comodule structure; one
says that \(\mathscr F\) is a \emph{subcomodule} of \(\mathscr E\).
Passing to the quotient, \(\mu_{\mathscr E}\) defines a morphism of
\(\mathcal O_S\)-modules
\[
  \mathscr E/\mathscr F\longrightarrow
  (\mathscr E/\mathscr F)\otimes\mathscr A,
\]
which endows \(\mathscr E/\mathscr F\) with an
\(\mathscr A\)-comodule structure. If
\(f:\mathscr E\to\mathscr E'\) is a morphism of
\(\mathscr A\)-comodules, then \(\operatorname{Ker}f\) (respectively,
\(\operatorname{Im}f\)) is an \(\mathscr A\)-subcomodule of
\(\mathscr E\) (respectively, of \(\mathscr E'\)), and \(f\) induces
an isomorphism of \(\mathscr A\)-comodules
\[
  \mathscr E/\operatorname{Ker}f
    \xrightarrow{\ \sim\ }
  \operatorname{Im}f.
\]
Moreover, if \(\mathscr E\) and \(\mathscr E'\) are quasi-coherent
\(\mathcal O_S\)-modules, then so are \(\operatorname{Ker}f\) and
\(\operatorname{Im}f\). Consequently,
\((\mathscr A\text{-}\mathbf{Comod.})\) and
\((\mathscr A\text{-}\mathbf{Comod.q.c.})\) are abelian categories.

\medskip
\noindent\SGAInlineRefTarget{sga3:i:corollary:4-7-2-1}\textbf{Corollary 4.7.2.1.}
\textit{Suppose that \(G\) is affine and flat over \(S\). Then the
category \((G\text{-}\mathcal O_S\text{-}\mathbf{Mod.q.c.})\)
(respectively, \((G\text{-}\mathcal O_S\text{-}\mathbf{Mod.})\)),
equivalent to the category of \(\mathscr A(G)\)-comodules that are
quasi-coherent over \(\mathcal O_S\) (respectively, of all
\(\mathscr A(G)\)-comodules), is abelian.}

\medskip
\noindent\SGAInlineRefTarget{sga3:i:paragraph:4-7-3}\textbf{4.7.3.}
Suppose now that \(G\) is a \emph{diagonalizable group}, that is, that
\(\mathscr A(G)\) is the algebra of a commutative group \(M\) over the
sheaf of rings \(\mathcal O_S\). If \(\mathscr F\) is an
\(\mathcal O_S\)-module, then
\[
  \mathscr F\otimes\mathscr A(G)
    =\coprod_{m\in M}\mathscr F\otimes m\mathcal O_S.
\]
Giving a morphism of \(\mathcal O_S\)-modules
\[
  \mu:\mathscr F\longrightarrow\mathscr F\otimes\mathscr A(G)
\]

% Source page transition: Exp1 component-local p. 31;
% combined-reader p. 41; re-edition printed p. 31. Only the continuation
% of Section 4.7 is included; Section 5 below it is deliberately excluded.
is therefore equivalent to giving \(\mathcal O_S\)-endomorphisms
\((\mu_m)_{m\in M}\) of \(\mathscr F\) such that, for every section
\(x\) of \(\mathscr F\) on an open subset of \(S\), the family
\((\mu_m(x))\) is a section of the direct sum
\(\coprod_{m\in M}\mathscr F\). This means that on every sufficiently
small open subset, only finitely many restrictions of the
\(\mu_m(x)\) are nonzero.

For \(\mu\) defined by
\[
  \mu(x)=\sum_{m\in M}\mu_m(x)\otimes m,
\]
condition (CM 1), respectively (CM 2), holds if and only if
\[
  \mu_m\circ\mu_{m'}=\delta_{mm'}\mu_m,
  \qquad
  \left(\text{respectively, }\sum_{m\in M}\mu_m
    =\operatorname{Id}_{\mathscr F}\right).
\]
Thus the \(\mu_m\) are pairwise orthogonal projectors whose sum is the
identity. We have proved the following result.

\medskip
\noindent\SGAInlineRefTarget{sga3:i:proposition:4-7-3}\textbf{Proposition 4.7.3.}
\textit{If \(G=D_S(M)\) is a diagonalizable \(S\)-group, the category
of quasi-coherent \(G\)-\(\mathcal O_S\)-modules (respectively, of all
\(G\)-\(\mathcal O_S\)-modules) is equivalent to the category of
quasi-coherent \(\mathcal O_S\)-modules (respectively, of all
\(\mathcal O_S\)-modules) graded by \(M\).}

\medskip
\noindent\textbf{Remark.}
If \(\mathscr F\) has an \(\mathcal O_S\)-algebra structure preserved
by the operations of \(G\), then the grading of \(\mathscr F\) is an
algebra grading. More precisely:

\medskip
\noindent\SGAInlineRefTarget{sga3:i:corollary:4-7-3-1}\textbf{Corollary 4.7.3.1.}
\textit{The functor
\(\mathscr A\mapsto\operatorname{Spec}(\mathscr A)\) induces an
equivalence between the category of quasi-coherent
\(\mathcal O_S\)-algebras graded by \(M\) and the category opposite to
that of \(S\)-schemes affine over \(S\) with \(S\)-operator group
\(G=D_S(M)\).}

\medskip
\noindent\SGAInlineRefTarget{sga3:i:proposition:4-7-4}\textbf{Proposition 4.7.4.}
\textit{Let \(G\) be a diagonalizable \(S\)-group. If
\[
  0\longrightarrow\mathscr F_1\longrightarrow\mathscr F_2
    \longrightarrow\mathscr F_3\longrightarrow0
\]
is an exact sequence of quasi-coherent \(G\)-\(\mathcal O_S\)-modules
that splits as a sequence of \(\mathcal O_S\)-modules, then it also
splits as a sequence of \(G\)-\(\mathcal O_S\)-modules.}

Indeed, if \(G=D_S(M)\), each \(\mathscr F_i\) is graded by the
\((\mathscr F_i)_m\), and for every \(m\in M\) the sequence
\[
  0\longrightarrow(\mathscr F_1)_m\longrightarrow(\mathscr F_2)_m
    \longrightarrow(\mathscr F_3)_m\longrightarrow0
\]
of \(\mathcal O_S\)-modules is split. The preceding proposition then
gives the result.

% End of assigned component. Section 5 (Group cohomology) begins next on
% Exp1 component-local p. 31 / combined-reader p. 41.

% END INCLUDED SGA3_English_current_buildable_source_20260823/source/inputs/I/tex/components/07_expose_I_s44_s47_scheme_groups.tex

% BEGIN INCLUDED SGA3_English_current_buildable_source_20260823/source/inputs/I/tex/components/08_expose_I_s5_group_cohomology.tex
% SGA 3, Expose I, Loop 1 English reconstruction.
% Scope: all of Section 5, "Group cohomology", stopping immediately before
% Section 6, "G-equivariant objects and modules".
% Authority: Polo--Gille Exp1-13oct24.pdf, component-local/printed pages 31--38
% (combined-reader pages 41--48). Force-OCR and the jcreinhold Markdown are
% drafting and comparison witnesses only; the PDF controls text, formulas,
% numbering, editorial notes, page boundaries, and the single diagram.

% Source page: Exp1 component-local/printed p. 31; combined-reader p. 41.
% Original-edition marginal p. 37 begins at Section 5.1 on this page.
\section*{\SGAInlineRefTarget{sga3:i:section:5}5. Group cohomology}
\addcontentsline{toc}{section}{5. Group cohomology}

\subsection*{\SGAInlineRefTarget{sga3:i:subsection:5-1}5.1. The standard complex}

Let \(\mathscr C\) be a category, \(\mathbf G\) a
\(\widehat{\mathscr C}\)-group, \(\mathbf O\) a
\(\widehat{\mathscr C}\)-ring, and \(\mathbf F\) a
\(\mathbf G\)-\(\mathbf O\)-module.\footnote{Editor's note: This complex
is often called the Hochschild complex; see, for example, Section~II.3 of
\([\mathrm{DG}70]\).} For \(n\geq 0\), set
\[
  C^n(\mathbf G,\mathbf F)
    =\operatorname{Hom}(\mathbf G^n,\mathbf F),
  \qquad
  \underline C^{\,n}(\mathbf G,\mathbf F)
    =\underline{\operatorname{Hom}}(\mathbf G^n,\mathbf F),
\]
where \(\mathbf G^0\) is the final object \(\underline{\mathbf e}\).
Then \(\underline C^{\,n}(\mathbf G,\mathbf F)\) (respectively,
\(C^n(\mathbf G,\mathbf F)\)) is endowed in the evident way with an
\(\mathbf O\)-module structure (respectively, a
\(\Gamma(\mathbf O)\)-module structure), and
\[
  C^n(\mathbf G,\mathbf F)
    \simeq \Gamma\bigl(\underline C^{\,n}(\mathbf G,\mathbf F)\bigr),
  \qquad
  \underline C^{\,n}(\mathbf G,\mathbf F)(S)
    =C^n(\mathbf G_S,\mathbf F_S).
\]

% --- Authority page transition: component/printed p. 31 -> p. 32;
% combined-reader p. 41 -> p. 42. ---
To give an element of \(C^n(\mathbf G,\mathbf F)\) is to give, for every
\(S\in\operatorname{Ob}(\mathscr C)\), an \(n\)-cochain of
\(\mathbf G(S)\) with values in \(\mathbf F(S)\), functorially in \(S\).
Recall that the boundary operator
\[
  \partial:C^n\bigl(\mathbf G(S),\mathbf F(S)\bigr)
    \longrightarrow
    C^{n+1}\bigl(\mathbf G(S),\mathbf F(S)\bigr)
\]
is given by
\begin{align*}
  \partial f(g_1,\ldots,g_{n+1})
    ={}&g_1f(g_2,\ldots,g_{n+1}) \\
     &+\sum_{i=1}^{n}(-1)^i
       f(g_1,\ldots,g_i g_{i+1},\ldots,g_{n+1}) \\
     &+(-1)^{n+1}f(g_1,\ldots,g_n).
\end{align*}
It is functorial in \(S\), and hence defines a homomorphism
\[
  \partial:C^n(\mathbf G,\mathbf F)
    \longrightarrow C^{n+1}(\mathbf G,\mathbf F)
\]
such that \(\partial\circ\partial=0\). We have thus defined a complex of
abelian groups (indeed, of \(\Gamma(\mathbf O)\)-modules), denoted
\(C^*(\mathbf G,\mathbf F)\). In the same way one defines the complex of
\(\mathbf O\)-modules \(\underline C^{\,*}(\mathbf G,\mathbf F)\), and
\[
  C^*(\mathbf G,\mathbf F)
    =\Gamma\bigl(\underline C^{\,*}(\mathbf G,\mathbf F)\bigr).
\]
We denote by \(H^n(\mathbf G,\mathbf F)\) (respectively,
\(\underline H^{\,n}(\mathbf G,\mathbf F)\)) the groups (respectively,
the \(\widehat{\mathscr C}\)-groups) of cohomology of
\(C^*(\mathbf G,\mathbf F)\) (respectively,
\(\underline C^{\,*}(\mathbf G,\mathbf F)\)). In particular,
\[
  \underline H^{\,0}(\mathbf G,\mathbf F)=\mathbf F^{\mathbf G},
  \qquad
  H^0(\mathbf G,\mathbf F)=\Gamma(\mathbf F^{\mathbf G}).
\]

\medskip
\noindent\SGAInlineRefTarget{sga3:i:remark:5-1-1}\textbf{Remark 5.1.1.}
The ``set-theoretic'' description of \(\partial\) given above is convenient
for checking that \(\partial\circ\partial=0\). One may also define
\(\partial\) in terms of the multiplication
\(m:\mathbf G\times\mathbf G\to\mathbf G\) and the action
\(\mu:\mathbf G\times\mathbf F\to\mathbf F\) as follows: for every
\(f\in C^n(\mathbf G,\mathbf F)\),
\[
  \partial f
    =\mu\circ(\operatorname{id}_{\mathbf G}\times f)
     +\sum_{i=1}^{n}(-1)^i
       f\circ
       \bigl(\operatorname{id}_{\mathbf G^{i-1}}\times m\times
             \operatorname{id}_{\mathbf G^{n-i}}\bigr)
     +(-1)^{n+1}f\circ\operatorname{pr}_{[1,n]},
\]
where \(\operatorname{pr}_{[1,n]}\) denotes the projection of
\(\mathbf G^{n+1}=\mathbf G^n\times\mathbf G\) onto \(\mathbf G^n\).
Similarly, for every \(S\in\operatorname{Ob}(\mathscr C)\) and
\(f\in\underline C^{\,n}(\mathbf G,\mathbf F)(S)
      =C^n(\mathbf G_S,\mathbf F_S)\),
\[
  \partial f
    =\mu_S\circ(\operatorname{id}_{\mathbf G_S}\times f)
     +\sum_{i=1}^{n}(-1)^i
       f\circ
       \bigl(\operatorname{id}_{\mathbf G_S^{i-1}}\times m_S\times
             \operatorname{id}_{\mathbf G_S^{n-i}}\bigr)
     +(-1)^{n+1}f\circ\operatorname{pr}_{[1,n]},
\]
where \(m_S\) and \(\mu_S\) are obtained from \(m\) and \(\mu\) by base
change.
\footnote{Editor's note: This remark has been added.}

\subsection*{\SGAInlineRefTarget{sga3:i:subsection:5-2}5.2.}

Recall (see \SGARefLink{sga3:i:section:3}{Section~3}) that
\((\mathbf G\text{-}\mathbf O\text{-Mod.})\) is an abelian category,
with its structure defined objectwise. Thus
\[
  0\longrightarrow\mathbf F'
   \longrightarrow\mathbf F
   \longrightarrow\mathbf F''
   \longrightarrow0
\]
is an exact sequence of \(\mathbf G\)-\(\mathbf O\)-modules if and only if
the sequence of abelian groups
\[
  0\longrightarrow\mathbf F'(S)
   \longrightarrow\mathbf F(S)
   \longrightarrow\mathbf F''(S)
   \longrightarrow0
\]
is exact for every \(S\in\operatorname{Ob}(\mathscr C)\).
\footnote{Editor's note: The preceding reminder has been added.}

% --- Authority page transition: component/printed p. 32 -> p. 33;
% combined-reader p. 42 -> p. 43. ---
Suppose that \(\mathscr C\) is small. Then, by \SGARefLink{sga3:i:remark:3-2-1}{3.2.1},
\((\mathbf G\text{-}\mathbf O\text{-Mod.})\) has enough injective objects,
so the derived functors of the left exact functors
\(\underline H^{\,0}\) and \(H^0\) are defined. We shall now show that
the functors \(\underline H^{\,n}\) (respectively, \(H^n\)) are indeed the
derived functors of \(\underline H^{\,0}\) (respectively, \(H^0\)).
\footnote{Editor's note: The preceding sentence has been added.}

\medskip
\noindent\SGAInlineRefTarget{sga3:i:definition:5-2-0}\textbf{Definition 5.2.0.}
For every \(\mathbf O\)-module \(\mathbf P\), let \(E(\mathbf P)\) denote
the object \(\underline{\operatorname{Hom}}(\mathbf G,\mathbf P)\) of
\(\widehat{\mathscr C}\), endowed with the following
\(\mathbf G\)-\(\mathbf O\)-module structure. For every
\(S\in\operatorname{Ob}(\mathscr C)\),
\[
  \underline{\operatorname{Hom}}(\mathbf G,\mathbf P)(S)
    =\operatorname{Hom}_S(\mathbf G_S,\mathbf P_S),
\]
and \(g\in\mathbf G(S)\) and \(a\in\mathbf O(S)\) act on
\(\phi\in\operatorname{Hom}_S(\mathbf G_S,\mathbf P_S)\) by
\[
  (g\cdot\phi)(h)=\phi(hg),
  \qquad
  (a\cdot\phi)(h)=a\phi(h),
\]
for every \(h\in\mathbf G(S')\), \(S'\to S\). Moreover, for every
\(\phi\in\operatorname{Hom}_S(\mathbf G_S,\mathbf P_S)\), set
\[
  \varepsilon(\phi)=\phi(1)\in\mathbf P(S),
\]
where \(1\) denotes the unit element of \(\mathbf G(S)\).
This defines a functor
\[
  E:(\mathbf O\text{-Mod.})
    \longrightarrow(\mathbf G\text{-}\mathbf O\text{-Mod.})
\]
and a natural transformation \(\varepsilon:E\to\operatorname{Id}\), where
\(\operatorname{Id}\) denotes the identity functor of
\((\mathbf O\text{-Mod.})\).
\footnote{Editor's note: The original has been modified so as to introduce
\SGARefLink{sga3:i:remark:5-2-0-1}{5.2.0.1} and \SGARefLink{sga3:i:lemma:5-2-0-2}{5.2.0.2}, which will be useful in the proof of \SGARefLink{sga3:i:theorem:5-3-1}{Theorem~5.3.1}.}

\medskip
\noindent\SGAInlineRefTarget{sga3:i:remark:5-2-0-1}\textbf{Remark 5.2.0.1.}
In what follows, let \(\mathbf G_1\) and \(\mathbf G_2\) denote two copies
of \(\mathbf G\). The morphism
\[
  \mathbf G_1\times E(\mathbf P)\longrightarrow E(\mathbf P),
  \qquad
  (g_1,\phi)\longmapsto\bigl(g_2\longmapsto\phi(g_2g_1)\bigr)
\]
corresponds, via the isomorphisms
\begin{align*}
  \operatorname{Hom}\bigl(\mathbf G_1\times E(\mathbf P),E(\mathbf P)\bigr)
    &\simeq
      \operatorname{Hom}\Bigl(E(\mathbf P),
        \underline{\operatorname{Hom}}\bigl(
          \mathbf G_1,
          \underline{\operatorname{Hom}}(\mathbf G_2,\mathbf P)
        \bigr)\Bigr)\\
    &\simeq
      \operatorname{Hom}\Bigl(E(\mathbf P),
        \underline{\operatorname{Hom}}(
          \mathbf G_2\times\mathbf G_1,\mathbf P)
      \Bigr),
\end{align*}
to the morphism
\[
  \phi\longmapsto\bigl((g_2,g_1)\longmapsto\phi(g_2g_1)\bigr),
\]
that is, to the morphism
\[
  \underline{\operatorname{Hom}}(\mathbf G,\mathbf P)
    \longrightarrow
  \underline{\operatorname{Hom}}(\mathbf G_2\times\mathbf G_1,\mathbf P)
\]
induced by the multiplication
\[
  \mu_{\mathbf G}:\mathbf G\times\mathbf G\longrightarrow\mathbf G,
  \qquad (g_2,g_1)\longmapsto g_2g_1.
\]

\medskip
\noindent\SGAInlineRefTarget{sga3:i:lemma:5-2-0-2}\textbf{Lemma 5.2.0.2.}
\textit{\SGAInlineRefTarget{sga3:i:lemma:5-2-0-2:item:i}(i) The functor \(E\) is right adjoint to the forgetful functor
\[
  (\mathbf G\text{-}\mathbf O\text{-Mod.})
    \longrightarrow(\mathbf O\text{-Mod.}).
\]
More precisely, \(\varepsilon:E\to\operatorname{Id}\) induces, for every
\(\mathbf M\in(\mathbf G\text{-}\mathbf O\text{-Mod.})\) and
\(\mathbf P\in\operatorname{Ob}(\mathbf O\text{-Mod.})\), a bijection}
\[
  \operatorname{Hom}_{\mathbf G\text{-}\mathbf O\text{-Mod.}}
    \bigl(\mathbf M,E(\mathbf P)\bigr)
  \xrightarrow{\ \sim\ }
  \operatorname{Hom}_{\mathbf O\text{-Mod.}}(\mathbf M,\mathbf P)
\]
\textit{functorial in \(\mathbf M\) and \(\mathbf P\).}

\textit{\SGAInlineRefTarget{sga3:i:lemma:5-2-0-2:item:ii}(ii) Consequently, if \(\mathbf I\) is an injective object of
\((\mathbf O\text{-Mod.})\), then \(E(\mathbf I)\) is an injective object
of \((\mathbf G\text{-}\mathbf O\text{-Mod.})\).}

% --- Authority page transition: component/printed p. 33 -> p. 34;
% combined-reader p. 43 -> p. 44. ---
\medskip
\noindent\textit{Proof.}
To every \(\mathbf O\)-morphism \(f:\mathbf M\to\mathbf P\), associate
the element \(\phi_f\) of
\(\operatorname{Hom}_{\mathbf O}(\mathbf M,E(\mathbf P))\) defined as
follows. For every \(S\in\operatorname{Ob}(\mathscr C)\) and
\(m\in\mathbf M(S)\), the element \(\phi_f(m)\) of
\(\operatorname{Hom}_S(\mathbf G_S,\mathbf P_S)\) is defined, for every
\(g\in\mathbf G(S')\), \(S'\to S\), by
\[
  \phi_f(m)(g)=f(gm)\in\mathbf P(S').
\]
Then, for every \(h\in\mathbf G(S)\), one has
\(\phi_f(hm)=h\cdot f(m)\); that is, \(\phi_f\) is an element of
\[
  \operatorname{Hom}_{\mathbf G\text{-}\mathbf O\text{-Mod.}}
    \bigl(\mathbf M,E(\mathbf P)\bigr).
\]
If
\(\phi\in\operatorname{Hom}_{\mathbf G\text{-}\mathbf O\text{-Mod.}}
  (\mathbf M,E(\mathbf P))\), and if for every \(m\in\mathbf M(S)\) we
write \(f(m)=\phi(m)(1)\), then
\[
  \phi_f(m)(g)
    =f(gm)
    =\phi(gm)(1)
    =(g\cdot\phi(m))(1)
    =\phi(m)(g);
\]
thus \(\phi_f=\phi\). Conversely, it is clear that
\(\phi_f(m)(1)=f(m)\). This proves \SGARefLink{sga3:i:lemma:5-2-0-2:item:i}{(i)}, and \SGARefLink{sga3:i:lemma:5-2-0-2:item:ii}{(ii)} follows at once.

\medskip
\noindent\SGAInlineRefTarget{sga3:i:definition:5-2-0-3}\textbf{Definition 5.2.0.3.}
Let \(\mathbf M\) be a \(\mathbf G\)-\(\mathbf O\)-module. The identity
map of \(\mathbf M\), considered as an \(\mathbf O\)-module, corresponds
under the adjunction to the morphism of
\(\mathbf G\)-\(\mathbf O\)-modules
\[
  \mu_{\mathbf M}:\mathbf M\longrightarrow E(\mathbf M)
\]
such that, for every \(S\in\operatorname{Ob}(\mathscr C)\) and
\(m\in\mathbf M(S)\), \(\mu_{\mathbf M}(m)\) is the morphism
\(\mathbf G_S\to\mathbf M_S\) defined as follows: for every
\(S'\to S\) and \(g\in\mathbf G(S')\),
\[
  \mu_{\mathbf M}(m)(g)=g\cdot m_{S'}\in\mathbf M(S').
\]
Notice that \(\mu_{\mathbf M}\) is a monomorphism. Indeed,
\(\varepsilon_{\mathbf M}:E(\mathbf M)\to\mathbf M\) is a morphism of
\(\mathbf O\)-modules such that
\(\varepsilon_{\mathbf M}\circ\mu_{\mathbf M}
  =\operatorname{id}_{\mathbf M}\). Consequently, as an
\(\mathbf O\)-module, \(\mathbf M\) is a direct summand of
\(E(\mathbf M)\).

\medskip
\noindent\SGAInlineRefTarget{sga3:i:proposition:5-2-1}\textbf{Proposition 5.2.1.}
\textit{Suppose that \(\mathscr C\) is small, that finite products exist
in \(\mathscr C\), and that \(\mathbf G\) is representable. Then the
functors \(\underline H^{\,n}(\mathbf G,-)\) (respectively,
\(H^n(\mathbf G,-)\)) are the derived functors of the left exact functor
\(\underline H^{\,0}(\mathbf G,-)\) (respectively,
\(H^0(\mathbf G,-)\)) on the category of
\(\mathbf G\)-\(\mathbf O\)-modules.}

By the well-known general results,
\footnote{Editor's note: See \([\mathrm{Gr}57,\,2.2.1]\) and
\([\mathrm{Gr}57,\,2.3]\). The original has also been expanded in what
follows.}
it suffices to verify that the \(\underline H^{\,n}(\mathbf G,-)\)
(respectively, \(H^n(\mathbf G,-)\)) form an effaceable cohomological
functor in degrees \(>0\). Let
\[
  0\longrightarrow\mathbf F'
   \longrightarrow\mathbf F
   \longrightarrow\mathbf F''
   \longrightarrow0
\]
be an exact sequence of \(\mathbf G\)-\(\mathbf O\)-modules, and let
\(S\in\operatorname{Ob}(\mathscr C)\). By hypothesis, \(\mathbf G\) is
represented by an object \(G\in\operatorname{Ob}(\mathscr C)\), and finite
products exist in \(\mathscr C\); in particular, \(\mathscr C\) has a
final object \(e\). Hence every \(\mathbf G^n\times\mathbf h_S\) is
represented by \(G^n\times S\) (with \(G^0=e\)), and the sequence
\[
  0\longrightarrow\mathbf F'(G^n\times S)
   \longrightarrow\mathbf F(G^n\times S)
   \longrightarrow\mathbf F''(G^n\times S)
   \longrightarrow0
\]
is exact. This shows that the sequence of \(\mathbf O\)-modules
\[
  0\longrightarrow\underline C^{\,n}(\mathbf h_G,\mathbf F')
   \longrightarrow\underline C^{\,n}(\mathbf h_G,\mathbf F)
   \longrightarrow\underline C^{\,n}(\mathbf h_G,\mathbf F'')
   \longrightarrow0
\]
is exact. It follows that \(\underline C^{\,*}(\mathbf G,-)\), regarded as
a functor on \((\mathbf G\text{-}\mathbf O\text{-Mod.})\) with values in
the category of complexes of \((\mathbf O\text{-Mod.})\), is exact. Thus
the functors \(\underline H^{\,n}(\mathbf G,-)\) are cohomological.

% --- Authority page transition: component/printed p. 34 -> p. 35;
% combined-reader p. 44 -> p. 45. ---
Since \(\Gamma\) is exact, the same is true of the functors
\(H^n(\mathbf G,-)\). It remains only to prove the following.

\medskip
\noindent\SGAInlineRefTarget{sga3:i:lemma:5-2-2}\textbf{Lemma 5.2.2.}
\textit{For every \(\mathbf P\in\operatorname{Ob}(\mathbf O\text{-Mod.})\),
one has}
\[
  \underline H^{\,n}\bigl(
    \mathbf G,\underline{\operatorname{Hom}}(\mathbf G,\mathbf P)
  \bigr)=0,
  \qquad
  H^n\bigl(
    \mathbf G,\underline{\operatorname{Hom}}(\mathbf G,\mathbf P)
  \bigr)=0,
  \qquad n>0.
\]

It is enough to show that
\(\underline C^{\,*}
  (\mathbf G,\underline{\operatorname{Hom}}(\mathbf G,\mathbf P))\) and
\(C^*(\mathbf G,\underline{\operatorname{Hom}}(\mathbf G,\mathbf P))\)
are homotopically trivial in degrees \(>0\). It is enough, in fact, to do
this for the second complex, since the corresponding result for the first
follows by base change.
\footnote{Editor's note: The original has been expanded in what follows.}
For every \(n\geq0\), define a morphism
\[
  \sigma:
  C^{n+1}\bigl(
    \mathbf G,\underline{\operatorname{Hom}}(\mathbf G,\mathbf P)
  \bigr)
  \longrightarrow
  C^n\bigl(
    \mathbf G,\underline{\operatorname{Hom}}(\mathbf G,\mathbf P)
  \bigr)
\]
as follows. Let
\(f\in C^{n+1}
  (\mathbf G,\underline{\operatorname{Hom}}(\mathbf G,\mathbf P))\).
For every \(S\in\operatorname{Ob}(\mathscr C)\) and
\(g_1,\ldots,g_n\in\mathbf G(S)\), the element
\(\sigma(f)(g_1,\ldots,g_n)\) of
\(\operatorname{Hom}_S(\mathbf G_S,\mathbf P_S)\) is defined, for every
\(S'\to S\) and \(x\in\mathbf G(S')\), by
\[
  \sigma(f)(g_1,\ldots,g_n)(x)
    =f(x,g_1,\ldots,g_n)(e)\in\mathbf P(S'),
\]
where \(e\) denotes the unit element of \(\mathbf G(S')\). Then \(\sigma\)
is a homotopy operator in degrees \(>0\). Indeed, for every
\(g_1,\ldots,g_{n+1}\in\mathbf G(S)\) and \(x\in\mathbf G(S')\), on the
one hand,
\begin{align*}
  \partial\sigma(f)(g_1,\ldots,g_{n+1})(x)
    ={}&f(xg_1,g_2,\ldots,g_{n+1})(e)\\
     &+\sum_{i=1}^{n}(-1)^i
       f(x,g_1,\ldots,g_i g_{i+1},\ldots,g_{n+1})(e)\\
     &+(-1)^{n+1}f(x,g_1,\ldots,g_n)(e),
\end{align*}
and, on the other hand,
\begin{align*}
  \sigma(\partial f)(g_1,\ldots,g_{n+1})(x)
    ={}&\bigl(xf(g_1,g_2,\ldots,g_{n+1})\bigr)(e)
       -f(xg_1,g_2,\ldots,g_{n+1})(e)\\
     &+\sum_{i=1}^{n}(-1)^{i+1}
       f(x,g_1,\ldots,g_i g_{i+1},\ldots,g_{n+1})\\
     &+(-1)^{n+2}f(x,g_1,\ldots,g_n)(e).
\end{align*}
Consequently,
\[
  \bigl(\partial\sigma(f)+\sigma(\partial f)\bigr)
    (g_1,\ldots,g_{n+1})(x)
  =f(g_1,\ldots,g_{n+1})(x);
\]
that is, \(\partial\sigma+\sigma\partial\) is the identity map of
\(C^{n+1}
  (\mathbf G,\underline{\operatorname{Hom}}(\mathbf G,\mathbf P))\) for
every \(n\geq0\).

\medskip
\noindent\SGAInlineRefTarget{sga3:i:remark:5-2-3}\textbf{Remark 5.2.3.}
The hypothesis that \(\mathscr C\) be small is used only to ensure the
existence of the derived functors \(R^n\underline H^{\,0}\) and \(R^nH^0\).
In every case, what precedes shows that the functors
\(\underline H^{\,n}(\mathbf G,-)\) (respectively,
\(H^n(\mathbf G,-)\)) form an effaceable cohomological functor in degrees
\(>0\), and hence are the (right) satellite functors of the left exact
functor \(\underline H^{\,0}(\mathbf G,-)\) (respectively,
\(H^0(\mathbf G,-)\)), in the sense of \([\mathrm{Gr}57,\,2.2]\). If
\((\mathbf G\text{-}\mathbf O\text{-Mod.})\) has enough injective objects
(as it does when \(\mathscr C\) is small), they coincide with the derived
functors (\emph{loc. cit.}, \SGARefLink{sga3:i:subsection:2-3}{2.3}).
\footnote{Editor's note: This remark has been added.}

% --- Authority page transition: component/printed p. 35 -> p. 36;
% combined-reader p. 45 -> p. 46. ---
\subsection*{\SGAInlineRefTarget{sga3:i:subsection:5-3}5.3. Cohomology of \(G\)-\(\mathscr O_S\)-modules}

Let \(S\) be a scheme, \(G\) an \(S\)-group, and \(\mathscr F\) a
quasi-coherent \(G\)-\(\mathscr O_S\)-module. The cohomology groups of
\(G\) with values in \(\mathscr F\) are defined by
\[
  H^n(G,\mathscr F)
    =H^n\bigl(\mathbf h_G,\mathbf W(\mathscr F)\bigr)
\]
(for the notation, see \SGARefLink{sga3:i:subsection:4-6}{4.6}).

Suppose that \(G\) is affine over \(S\). In view of \SGARefLink{sga3:i:proposition:4-6-4}{Proposition~4.6.4},
this cohomology is computed as follows: \(H^n(G,\mathscr F)\) is the
\(n\)-th cohomology group of the complex \(C^*(G,\mathscr F)\), whose
\(n\)-th term is
\[
  C^n(G,\mathscr F)
    =\Gamma\Bigl(
       S,
       \mathscr F\otimes
       \underbrace{\mathscr A(G)\otimes\cdots\otimes\mathscr A(G)}_
                   {n\text{ times}}
     \Bigr).
\]
If \(f\) (respectively, \(a_i\)) is a section of \(\mathscr F\)
(respectively, of \(\mathscr A(G)\)) over an open subset of \(S\), then
\begin{align*}
  \partial(f\otimes a_1\otimes\cdots\otimes a_n)
    ={}&\mu_{\mathscr F}(f)\otimes a_1\otimes a_2\otimes\cdots\otimes a_n\\
     &+\sum_{i=1}^{n}(-1)^i
       f\otimes a_1\otimes\cdots\otimes\Delta a_i\otimes\cdots\otimes a_n\\
     &+(-1)^{n+1}
       f\otimes a_1\otimes a_2\otimes\cdots\otimes a_n\otimes1,
\end{align*}
where
\[
  \Delta:\mathscr A(G)\longrightarrow\mathscr A(G)\otimes\mathscr A(G),
  \qquad
  \mu_{\mathscr F}:\mathscr F\longrightarrow
    \mathscr F\otimes\mathscr A(G)
\]
describe the coalgebra structure of \(\mathscr A(G)\) and the comodule
structure of \(\mathscr F\). Notice in passing that the cohomology of
\(G\) with values in \(\mathscr F\) therefore depends only on the
comodule structure of \(\mathscr F\), and in particular only on the
\(S\)-monoid structure of \(G\).

In particular,
\[
  H^0(G,\mathscr F)=\Gamma(S,\mathscr F^G),
\]
where \(\mathscr F^G\), the \emph{sheaf of invariants of \(\mathscr F\)},
is the sheaf whose sections over an open subset \(U\) of \(S\) are those
sections of \(\mathscr F\) over \(U\) whose inverse image in every
\(S'\) over \(U\) is invariant under \(G(S')\) (see \SGARefLink{sga3:i:remark:4-7-1-2}{4.7.1.2}).

\medskip
\noindent\SGAInlineRefTarget{sga3:i:theorem:5-3-1}\textbf{Theorem 5.3.1.}
\textit{Let \(S\) be an affine scheme and \(G\) an \(S\)-group that is
affine and flat over \(S\). The functors \(H^n(G,-)\) are the derived
functors of \(H^0(G,-)\) on the category of quasi-coherent
\(G\)-\(\mathscr O_S\)-modules.}

\medskip
\noindent\textit{Proof.}
Since \(G\) is affine and flat over \(S\), \SGARefLink{sga3:i:corollary:4-7-2-1}{4.7.2.1} shows that the category
\((G\text{-}\mathscr O_S\text{-Mod.q.c.})\) is equivalent to the category
\((\mathscr A(G)\text{-Comod.q.c.})\) of quasi-coherent
\(\mathscr A(G)\)-comodules over \(\mathscr O_S\), and is therefore
abelian. Moreover, since \(\mathscr A(G)\) is a flat
\(\mathscr O_S\)-module, each functor
\[
  \mathscr F\longmapsto
    \mathscr F\otimes_{\mathscr O_S}\mathscr A(G)^{\otimes n}
\]
is exact; since \(S\) is affine, it follows that \(C^*(G,-)\) is an exact
functor on \((G\text{-}\mathscr O_S\text{-Mod.q.c.})\).
\footnote{Editor's note: The original has been modified to show that the
category \((G\text{-}\mathscr O_S\text{-Mod.q.c.})\) is abelian and has
enough injective objects. Compare \([\mathrm{Ja}03]\), Part~I,
3.3--3.4, 3.9, \SGARefLink{sga3:i:subsection:4-2}{4.2}, and 4.14--4.16 (bearing in mind that ``\(k\)-group
scheme'' there means ``affine \(k\)-group scheme''; see \emph{loc. cit.},
\SGARefLink{sga3:i:subsection:2-1}{2.1}).}

Let \(\Delta\) (respectively, \(\eta\)) denote the comultiplication
(respectively, the augmentation) of \(\mathscr A(G)\). For every
quasi-coherent \(\mathscr O_S\)-module \(\mathscr P\), let
\[
  \operatorname{Ind}(\mathscr P)
    =\mathscr P\otimes_{\mathscr O_S}\mathscr A(G),
\]
endowed with the \(\mathscr A(G)\)-comodule structure defined by
\[
  \operatorname{id}_{\mathscr P}\otimes\Delta:
  \mathscr P\otimes_{\mathscr O_S}\mathscr A(G)
    \longrightarrow
  \mathscr P\otimes_{\mathscr O_S}\mathscr A(G)
    \otimes_{\mathscr O_S}\mathscr A(G).
\]

% --- Authority page transition: component/printed p. 36 -> p. 37;
% combined-reader p. 46 -> p. 47. ---
This defines a functor
\[
  \operatorname{Ind}:(\mathscr O_S\text{-Mod.q.c.})
    \longrightarrow(G\text{-}\mathscr O_S\text{-Mod.q.c.}).
\]
It follows from \SGARefLink{sga3:i:corollary:4-6-4-1}{4.6.4.1}, \SGARefLink{sga3:i:definition:5-2-0}{5.2.0}, and \SGARefLink{sga3:i:remark:5-2-0-1}{5.2.0.1} that there is an isomorphism of
\(\mathbf G\)-\(\mathbf O_S\)-modules
\[
  \tag{*}\label{sga3:i:theorem:5-3-1:equation:star}
  \mathbf W\bigl(\operatorname{Ind}(\mathscr P)\bigr)
    \simeq E\bigl(\mathbf W(\mathscr P)\bigr)
    =\underline{\operatorname{Hom}}
       \bigl(\mathbf G,\mathbf W(\mathscr P)\bigr).
\]
Under this identification, the morphism
\(\varepsilon:E(\mathbf W(\mathscr P))\to\mathbf W(\mathscr P)\)
corresponds to the morphism of \(\mathscr O_S\)-modules
\[
  \operatorname{id}_{\mathscr P}\otimes\eta:
    \operatorname{Ind}(\mathscr P)\longrightarrow\mathscr P.
\]

We have already used the fact that the functor
\[
  \mathbf W:(\mathscr O_S\text{-Mod.})
    \longrightarrow(\mathbf O_S\text{-Mod.})
\]
is fully faithful. By \SGARefLink{sga3:i:definition:4-7-1}{Definition~4.7.1}, the same is true of its restriction
to \((G\text{-}\mathscr O_S\text{-Mod.})\): if \(\mathscr M\) and
\(\mathscr M'\) are \(G\)-\(\mathscr O_S\)-modules, there is a functorial
isomorphism
\[
  \operatorname{Hom}_{G\text{-}\mathscr O_S\text{-Mod.}}
    (\mathscr M,\mathscr M')
  \simeq
  \operatorname{Hom}_{\mathbf G\text{-}\mathbf O_S\text{-Mod.}}
    \bigl(\mathbf W(\mathscr M),\mathbf W(\mathscr M')\bigr).
\]
Consequently, \SGARefLink{sga3:i:lemma:5-2-0-2}{Lemma~5.2.0.2} gives the following result.

\medskip
\noindent\SGAInlineRefTarget{sga3:i:corollary:5-3-1-1}\textbf{Corollary 5.3.1.1.}
\textit{\SGAInlineRefTarget{sga3:i:corollary:5-3-1-1:item:i}(i) The functor \(\operatorname{Ind}\) is right adjoint to the
forgetful functor}
\[
  (G\text{-}\mathscr O_S\text{-Mod.q.c.})
    \longrightarrow(\mathscr O_S\text{-Mod.q.c.}).
\]
\textit{More precisely, the map
\(\operatorname{id}_{\mathscr P}\otimes\eta:
  \operatorname{Ind}(\mathscr P)\to\mathscr P\)
induces, for every object \(\mathscr M\) of
\((G\text{-}\mathscr O_S\text{-Mod.q.c.})\), a bijection}
\[
  \operatorname{Hom}_{G\text{-}\mathscr O_S\text{-Mod.}}
    \bigl(\mathscr M,\operatorname{Ind}(\mathscr P)\bigr)
  \xrightarrow{\ \sim\ }
  \operatorname{Hom}_{\mathscr O_S}(\mathscr M,\mathscr P).
\]

\textit{\SGAInlineRefTarget{sga3:i:corollary:5-3-1-1:item:ii}(ii) Consequently, if \(\mathscr I\) is an injective object of
\((\mathscr O_S\text{-Mod.q.c.})\), then
\(\operatorname{Ind}(\mathscr I)\) is an injective object of
\((G\text{-}\mathscr O_S\text{-Mod.q.c.})\).}

Let \(\mathscr F\) be a \(G\)-\(\mathscr O_S\)-module, and let
\(\mu_{\mathscr F}:\mathscr F\to\operatorname{Ind}(\mathscr F)\) be the
map defining its \(\mathscr A(G)\)-comodule structure. It follows from
\SGARefLink{sga3:i:definition:5-2-0-3}{5.2.0.3} (or directly from axioms (CM~1) and (CM~2) of \SGARefLink{sga3:i:paragraph:4-7-2}{4.7.2}) that
\(\mu_{\mathscr F}\) is a morphism of
\(G\)-\(\mathscr O_S\)-modules, and that
\[
  (\operatorname{id}_{\mathscr F}\otimes\eta)\circ\mu_{\mathscr F}
    =\operatorname{id}_{\mathscr F}.
\]
Thus \(\mathscr F\) is a direct summand of
\(\operatorname{Ind}(\mathscr F)\) as an \(\mathscr O_S\)-module; in
particular, \(\mu_{\mathscr F}\) is a monomorphism. By (*) and \SGARefLink{sga3:i:lemma:5-2-2}{5.2.2},
\[
  H^n\bigl(\mathbf G,\mathbf W(\operatorname{Ind}(\mathscr F))\bigr)
    \simeq
  H^n\Bigl(
    \mathbf G,
    \underline{\operatorname{Hom}}_S
      \bigl(\mathbf G,\mathbf W(\mathscr F)\bigr)
  \Bigr)
    =0
  \qquad(n>0).
\]
Therefore \(H^n(G,-)\) is effaceable for \(n>0\).

Finally, since \(S\) is affine,
\((\mathscr O_S\text{-Mod.q.c.})\) has enough injective objects. Let
\(\mathscr F\hookrightarrow\mathscr I\) be a monomorphism of
\(\mathscr O_S\)-modules, where \(\mathscr I\) is injective in
\((\mathscr O_S\text{-Mod.q.c.})\). Since \(\mathscr A(G)\) is flat over
\(\mathscr O_S\), \(\operatorname{Ind}(\mathscr F)\) is a
\(G\)-\(\mathscr O_S\)-submodule of \(\operatorname{Ind}(\mathscr I)\).
Hence:

\medskip
\noindent\SGAInlineRefTarget{sga3:i:corollary:5-3-1-2}\textbf{Corollary 5.3.1.2.}
\textit{The abelian category \((G\text{-}\mathscr O_S\text{-Mod.q.c.})\)
has enough injective objects.}

Taking account of \([\mathrm{Gr}57,\,2.2.1\text{ and }2.3]\), already used
in the proof of \SGARefLink{sga3:i:proposition:5-2-1}{5.2.1}, completes the proof of \SGARefLink{sga3:i:theorem:5-3-1}{Theorem~5.3.1}.

\medskip
\noindent\SGAInlineRefTarget{sga3:i:remark:5-3-1-3}\textbf{Remark 5.3.1.3.}
\SGARefLink{sga3:i:corollary:5-3-1-1}{Corollary~5.3.1.1} may also be proved by the following calculation. To every
morphism of \(G\)-\(\mathscr O_S\)-modules
\[
  \phi:\mathscr M\longrightarrow
    \mathscr P\otimes_{\mathscr O_S}\mathscr A(G)
\]
associate the \(\mathscr O_S\)-morphism
\[
  (\operatorname{id}_{\mathscr P}\otimes\eta)\circ\phi:
    \mathscr M\longrightarrow\mathscr P.
\]
Conversely, to every \(\mathscr O_S\)-morphism
\(f:\mathscr M\to\mathscr P\), associate the morphism of
\(G\)-\(\mathscr O_S\)-modules
\[
  (f\otimes\operatorname{id}_{\mathscr A(G)})\circ\mu_{\mathscr M}:
    \mathscr M\longrightarrow
    \mathscr P\otimes_{\mathscr O_S}\mathscr A(G).
\]
One sees at once that
\[
  (\operatorname{id}_{\mathscr P}\otimes\eta)
    \circ(f\otimes\operatorname{id}_{\mathscr A(G)})
    \circ\mu_{\mathscr M}
  =(f\otimes\operatorname{id}_{\mathscr O_S})
    \circ(\operatorname{id}_{\mathscr P}\otimes\eta)
    \circ\mu_{\mathscr M}
  =f.
\]

% --- Authority page transition: component/printed p. 37 -> p. 38;
% combined-reader p. 47 -> p. 48. ---
On the other hand, for every \(\phi\), the following diagram is
commutative.

\begin{figure}[H]
  \centering
  % BEGIN INCLUDED SGA3_English_current_buildable_source_20260823/source/inputs/I/native_diagrams/exp1/diag_018_induction_adjunction_square.tex
\begin{tikzcd}[column sep=5.2em,row sep=3.3em]
  \mathscr M
    \arrow[r,"\phi"]
    \arrow[d,"\mu_{\mathscr M}"']
  &
  \mathscr P\otimes_{\mathcal O_S}\mathscr A(G)
    \arrow[d,"\operatorname{Id}_{\mathscr P}\otimes\Delta"]
  \\
  \mathscr M\otimes_{\mathcal O_S}\mathscr A(G)
    \arrow[r,"\phi\otimes\operatorname{Id}_{\mathscr A(G)}"]
  &
  \mathscr P\otimes_{\mathcal O_S}\mathscr A(G)
    \otimes_{\mathcal O_S}\mathscr A(G)
  \qquad .
\end{tikzcd}

% END INCLUDED SGA3_English_current_buildable_source_20260823/source/inputs/I/native_diagrams/exp1/diag_018_induction_adjunction_square.tex

  \phantomsection\label{fig:sga3-exp1-p038-induction-adjunction-square}
\end{figure}

It follows that
\begin{align*}
 &\Bigl(
    \bigl((\operatorname{id}_{\mathscr P}\otimes\eta)\circ\phi\bigr)
      \otimes\operatorname{id}_{\mathscr A(G)}
  \Bigr)\circ\mu_{\mathscr M}\\
 &\qquad=
  (\operatorname{id}_{\mathscr P}\otimes\eta
     \otimes\operatorname{id}_{\mathscr A(G)})
  \circ(\phi\otimes\operatorname{id}_{\mathscr A(G)})
  \circ\mu_{\mathscr M}\\
 &\qquad=
  (\operatorname{id}_{\mathscr P}\otimes\eta
     \otimes\operatorname{id}_{\mathscr A(G)})
  \circ(\operatorname{id}_{\mathscr P}\otimes\Delta)
  \circ\phi
  =\phi.
\end{align*}
This proves \SGARefLink{sga3:i:corollary:5-3-1-1:item:i}{Corollary~5.3.1.1(i)}, and \SGARefLink{sga3:i:corollary:5-3-1-1:item:ii}{(ii)} follows.

Let \(\mathscr F\) be a \(G\)-\(\mathscr O_S\)-module. We saw above that
axiom (CM~2) of \SGARefLink{sga3:i:proposition:4-7-2}{4.7.2} shows that, as an \(\mathscr O_S\)-module,
\(\mathscr F\) is a direct summand of \(E(\mathscr F)\). This gives:

\medskip
\noindent\SGAInlineRefTarget{sga3:i:proposition:5-3-2}\textbf{Proposition 5.3.2.}
\textit{Let \(S\) be an affine scheme and \(G\) an \(S\)-group that is
affine and flat. Suppose that every exact sequence
\[
  0\longrightarrow\mathscr F_1
   \longrightarrow\mathscr F_2
   \longrightarrow\mathscr F_3
   \longrightarrow0
\]
of quasi-coherent \(G\)-\(\mathscr O_S\)-modules which splits as a sequence
of \(\mathscr O_S\)-modules also splits as a sequence of
\(G\)-\(\mathscr O_S\)-modules. Then the functors \(H^n(G,-)\) vanish for
\(n>0\) (or, equivalently, the functor \(H^0(G,-)\) is exact).}

Indeed, by hypothesis, the sequence of \(G\)-\(\mathscr O_S\)-modules
\[
  0\longrightarrow\mathscr F
   \longrightarrow E(\mathscr F)
   \longrightarrow E(\mathscr F)/\mathscr F
   \longrightarrow0
\]
is split. Thus \(\mathscr F\) is a direct summand of \(E(\mathscr F)\) as
a \(G\)-\(\mathscr O_S\)-module, while the cohomology of the latter
vanishes.

Combining this with \SGARefLink{sga3:i:proposition:4-7-4}{Proposition~4.7.4} immediately gives:

\medskip
\noindent\SGAInlineRefTarget{sga3:i:theorem:5-3-3}\textbf{Theorem 5.3.3.}
\textit{Let \(S\) be an affine scheme and \(G\) a diagonalizable
\(S\)-group. For every quasi-coherent \(G\)-\(\mathscr O_S\)-module
\(\mathscr F\), one has}
\[
  H^n(G,\mathscr F)=0\qquad(n>0).
\]

\medskip
\noindent\textbf{Remark.}
\SGARefLink{sga3:i:proposition:5-3-2}{Proposition~5.3.2} remains valid when \(G\) is not necessarily flat over
\(S\); the proof then uses relative cohomology.
\footnote{Editor's note: The editors have not attempted to develop this
remark.}

% End of assigned component. Section 6 begins on Exp1
% component-local/printed p. 39 (combined-reader p. 49), in the next unit.

% END INCLUDED SGA3_English_current_buildable_source_20260823/source/inputs/I/tex/components/08_expose_I_s5_group_cohomology.tex

% BEGIN INCLUDED SGA3_English_current_buildable_source_20260823/source/inputs/I/tex/components/09_expose_I_s6_equivariant_objects.tex
% SGA 3, Expose I, Loop 1 English reconstruction.
% Scope: Section 6, "G-equivariant objects and modules," through the
% bibliography and end of Expose I.
% Authority: Polo--Gille Exp1-13oct24.pdf, component-local pages 39--48
% (combined-reader pages 49--58). Force-OCR and the English Markdown by
% Jacob Reinhold are drafting and comparison witnesses only.

% Source page: Exp1 component-local p. 39; combined-reader p. 49.
\section*{\SGAInlineRefTarget{sga3:i:section:6}6. \(G\)-equivariant objects and modules}
\addcontentsline{toc}{section}{6. G-equivariant objects and modules}

Let \(\mathscr C\) be a category having a final object and in which fiber
products exist.\footnote{Editor's note: This section has been added.}
Let \(G\) be a \(\widehat{\mathscr C}\)-group,
\(\pi:M\to X\) a morphism in \(\widehat{\mathscr C}\), and
\(\lambda=\lambda_X:G\times X\to X\) an action of \(G\) on \(X\).
In what follows, \(Y\times_f M\) denotes the fiber product of
\(\pi:M\to X\) and an \(X\)-functor \(f:Y\to X\).

For every \(U\in\operatorname{Ob}(\mathscr C)\) and \(x\in X(U)\), write
\(M_x=U\times_x M\); thus, for every \(\phi:U'\to U\),
\[
  M_x(U')
  =\{m\in M(U')\mid \pi(m)=x_{U'}=\phi^*(x)\}.
\]
Finally, if \(g\in G(U)\), we also write \(gx\) for the element
\(\lambda(g,x)\) of \(X(U)\).

\medskip
\noindent\SGAInlineRefTarget{sga3:i:definition:6-1}\textbf{Definition 6.1.}
\textit{\SGAInlineRefTarget{sga3:i:definition:6-1:item:a}(a) We say that \(M\) is a \(G\)-equivariant \(X\)-object if an
action \(\Lambda:G\times M\to M\) of \(G\) on \(M\), lifting
\(\lambda\), has been given; that is, if the following square commutes.}

\begin{figure}[H]
  \centering
  % BEGIN INCLUDED SGA3_English_current_buildable_source_20260823/source/inputs/I/native_diagrams/exp1/diag_019_equivariant_action_square.tex
\begin{tikzcd}[column sep=4.7em,row sep=3.2em]
  G\times M
    \arrow[r,"\Lambda"]
    \arrow[d,"\operatorname{id}_G\times\pi"']
  &
  M
    \arrow[d,"\pi"]
  \\
  G\times X
    \arrow[r,"\lambda"]
  &
  X
\end{tikzcd}

% END INCLUDED SGA3_English_current_buildable_source_20260823/source/inputs/I/native_diagrams/exp1/diag_019_equivariant_action_square.tex

  \phantomsection\label{fig:sga3-exp1-p039-equivariant-action-square}
\end{figure}

\textit{Equivalently, for every morphism
\((g,x):U\to G\times X\), one is given maps}
\[
  \Lambda_x^U(g):M_x(U)\longrightarrow M_{gx}(U),
  \qquad m\longmapsto g\cdot m,
\]
\textit{which satisfy \(1\cdot m=m\) and
\(g\cdot(h\cdot m)=(gh)\cdot m\), and which are functorial in the
\((G\times X)\)-object \(U\). Equivalently again, one is given morphisms
of \(U\)-objects}
\[
  \Lambda_x(g):M_x\longrightarrow M_{gx}
\]
\textit{satisfying
\(\Lambda_x(1)=\operatorname{id}\) and
\(\Lambda_{hx}(g)\circ\Lambda_x(h)=\Lambda_x(gh)\).}

\textit{(b) Now let \(\mathbf O\) be a
\(\widehat{\mathscr C}\)-ring and put
\(\mathbf O_X=\mathbf O\times X\). Under the hypotheses of~\SGARefLink{sga3:i:definition:6-1:item:a}{(a)}, we say
that \(M\) is a \(G\)-equivariant \(\mathbf O_X\)-module if it is an
\(\mathbf O_X\)-module (cf. \SGARefLink{sga3:i:paragraph:4-3-3-1}{Definition~4.3.3.1}, which is valid for every
ring functor on a category \(\mathscr C\)) and if the action \(\Lambda\)
is compatible with the \(\mathbf O_X\)-module structure on \(M\). Thus,
for every \((g,x)\in G(U)\times X(U)\),
\(U\in\operatorname{Ob}(\mathscr C)\), the map}
\[
  \Lambda_x(g):M_x\longrightarrow M_{gx},
  \qquad m\longmapsto g\cdot m,
\]
\textit{is a morphism of \(\mathbf O_U\)-modules.}

\medskip
\noindent\SGAInlineRefTarget{sga3:i:remark:6-2}\textbf{Remark 6.2.}
(i) In~\SGARefLink{sga3:i:definition:6-1:item:a}{(a)}, the conditions
\(\Lambda_x(1)=\operatorname{id}\) and
\(\Lambda_{hx}(g)\circ\Lambda_x(h)=\Lambda_x(gh)\) plainly imply that
each \(\Lambda_x(g)\) is an isomorphism, with inverse
\(\Lambda_{gx}(g^{-1})\). Conversely, if each \(\Lambda_x(g)\) is an
isomorphism, applying the second condition with \(h=1\) gives
\(\Lambda_x(1)=\operatorname{id}\).

(ii) Let \(M\) be an \(\mathbf O_X\)-module. Giving a morphism
\(\Lambda:G\times M\to M\) which makes the diagram of \SGARefLink{sga3:i:definition:6-1}{Definition~6.1}
commute and for which every morphism

% Source page: Exp1 component-local p. 40; combined-reader p. 50.
\[
  \Lambda_x(g):M_x\longrightarrow M_{gx},
  \qquad m\longmapsto g\cdot m,
\]
is an isomorphism of \(\mathbf O_U\)-modules is equivalent to giving an
isomorphism of \(\mathbf O_{G\times X}\)-modules
\[
  \theta:
  G\times M=(G\times X)\times_{\mathrm{pr}_X}M
  \xrightarrow{\ \sim\ }(G\times X)\times_\lambda M,
  \qquad
  (g,x,m)\longmapsto(g,x,g\cdot m).
\]
Since each \(\Lambda_x(g)\) is assumed to be an isomorphism, the equality
\(\Lambda_x(1)=\operatorname{id}\) follows from
\(\Lambda_{hx}(g)\circ\Lambda_x(h)=\Lambda_x(gh)\) with \(h=1\).
Consequently, \(\Lambda\) is an action of \(G\) on \(M\), that is,
\(g\cdot(h\cdot m)=(gh)\cdot m\), if and only if the following diagram
of \((G\times G\times X)\)-isomorphisms commutes. Here \(\mu\) denotes
multiplication in \(G\), and \(f^*(\theta)\) is obtained from \(\theta\)
by base change along \(f:G\times G\times X\to G\times X\).

\begin{figure}[H]
  \centering
  % BEGIN INCLUDED SGA3_English_current_buildable_source_20260823/source/inputs/I/native_diagrams/exp1/diag_020_equivariance_cocycle.tex
\begin{tikzcd}[column sep=1.7em,row sep=3.2em]
  (G\times G\times X)
    \mathop{\times}\limits_{\mathrm{pr}_X\circ\mathrm{pr}_{2,3}} M
    \arrow[rr,"\mathrm{pr}_{2,3}^{*}(\theta)","\sim"']
    \arrow[d,equal]
  &&
  (G\times G\times X)
    \mathop{\times}\limits_{\lambda\circ\mathrm{pr}_{2,3}} M
    \arrow[d,equal]
  \\
  (G\times G\times X)
    \mathop{\times}\limits_{\mathrm{pr}_X\circ(\mu\times\operatorname{id}_X)} M
    \arrow[d,"(\mu\times\operatorname{id}_X)^*(\theta)"',"\sim"]
  &&
  (G\times G\times X)
    \mathop{\times}\limits_{\mathrm{pr}_X\circ(\operatorname{id}_G\times\lambda)} M
    \arrow[d,"(\operatorname{id}_G\times\lambda)^*(\theta)","\sim"']
  \\
  (G\times G\times X)
    \mathop{\times}\limits_{\lambda\circ(\mu\times\operatorname{id}_X)} M
    \arrow[rr,equal]
  &&
  (G\times G\times X)
    \mathop{\times}\limits_{\lambda\circ(\operatorname{id}_G\times\lambda)} M
  \qquad .
\end{tikzcd}

% END INCLUDED SGA3_English_current_buildable_source_20260823/source/inputs/I/native_diagrams/exp1/diag_020_equivariance_cocycle.tex

  \phantomsection\label{fig:sga3-exp1-p040-equivariance-cocycle}
\end{figure}

Thus, giving a \(G\)-equivariant \(\mathbf O_X\)-module structure on
\(M\) is equivalent to giving an isomorphism \(\theta\) of
\(\mathbf O_{G\times X}\)-modules as above for which this diagram
commutes.

(iii) Everything above extends to the case in which \(G\) is merely a
\(\widehat{\mathscr C}\)-monoid. In that case, giving an action
\(\Lambda:G\times M\to M\) lifting \(\lambda\), such that every
\(\Lambda_x(g):M_x\to M_{gx}\) is a morphism of
\(\mathbf O_U\)-modules, is equivalent to giving a morphism \(\theta\)
of \(\mathbf O_{G\times X}\)-modules as in~(ii), for which the diagram
above commutes without the isomorphism signs under the arrows and which
satisfies
\[
  p_M\circ\theta\circ(\varepsilon_G\times\operatorname{id}_M)
  =\operatorname{id}_M,
\]
where \(\varepsilon_G\) is the unit section of \(G\) and \(p_M\) is the
projection onto \(M\).

\subsection*{\SGAInlineRefTarget{sga3:i:subsection:6-3}6.3. \(G\)-equivariant morphisms}

Let \(Y\) be a second object of \(\widehat{\mathscr C}\), endowed with
an action \(\lambda_Y:G\times Y\to Y\), and let \(N\) be a second
\(G\)-equivariant \(\mathbf O_X\)-module. A morphism
\(f:Y\to X\) in \(\widehat{\mathscr C}\) (respectively, a morphism of
\(\mathbf O_X\)-modules \(\phi:M\to N\)) is called
\emph{\(G\)-equivariant} if it commutes with the action of \(G\). In
set-theoretic notation, this says
\[
  f(g\cdot y)=g\cdot f(y)
  \qquad
  \bigl(\text{respectively, }\phi(g\cdot m)=g\cdot\phi(m)\bigr),
\]
or, equivalently,
\[
  f\circ\lambda_Y=\lambda_X\circ(\operatorname{id}_G\times f),
  \qquad
  \phi\circ\Lambda_M=\Lambda_N\circ(\operatorname{id}_G\times\phi).
\]
The following lemma is immediate.

% Source page: Exp1 component-local p. 41; combined-reader p. 51.
\medskip
\noindent\SGAInlineRefTarget{sga3:i:lemma:6-3-1}\textbf{Lemma 6.3.1.}
\textit{Let \(f:Y\to X\) be a \(G\)-equivariant morphism and let \(M\)
be a \(G\)-equivariant \(\mathbf O_X\)-module. Then the inverse image
\(f^*(M)=Y\times_fM\) is a \(G\)-equivariant
\(\mathbf O_Y\)-module.}

On the other hand, \(G\) acts on
\(\operatorname{Hom}_{\widehat{\mathscr C}}(Y,X)\). Indeed, let
\(T\in\operatorname{Ob}(\mathscr C)\), \(g\in G(T)\), and let \(\phi\)
be a \(T\)-morphism \(Y_T\to X_T\). Then \(g\) defines automorphisms
\(\lambda_Y(g)\) and \(\lambda_X(g)\) of \(Y_T\) and \(X_T\), and we
write \(g\cdot\phi\), or \(g\phi g^{-1}\), for the morphism
\[
  \lambda_X(g)\circ\phi\circ\lambda_Y(g^{-1}).
\]
This defines an action of \(G(T)\) on
\(\operatorname{Hom}_T(Y_T,X_T)\), functorially in \(T\).

\medskip
\noindent\SGAInlineRefTarget{sga3:i:definition:6-3-2}\textbf{Definition 6.3.2.}
\textit{If \(\phi:Y\to X\) is an arbitrary morphism in
\(\widehat{\mathscr C}\), we may therefore consider its stabilizer
\(\operatorname{Stab}_G(\phi)\) (cf. \SGARefLink{sga3:i:scholium:2-3-3-1}{2.3.3.1}). For every
\(T\in\operatorname{Ob}(\mathscr C)\),
\(\operatorname{Stab}_G(\phi)(T)\) is the subgroup of \(G(T)\) formed
by those \(g\) for which \(g\circ\phi_T=\phi_T\circ g\), that is, for
which the following diagram commutes.}

\begin{figure}[H]
  \centering
  % BEGIN INCLUDED SGA3_English_current_buildable_source_20260823/source/inputs/I/native_diagrams/exp1/diag_021_stabilizer_square.tex
\begin{tikzcd}[column sep=4.0em,row sep=3.0em]
  Y_T
    \arrow[r,"\phi_T"]
    \arrow[d,"g"']
  &
  X_T
    \arrow[d,"g"]
  \\
  Y_T
    \arrow[r,"\phi_T"]
  &
  X_T
\end{tikzcd}

% END INCLUDED SGA3_English_current_buildable_source_20260823/source/inputs/I/native_diagrams/exp1/diag_021_stabilizer_square.tex

  \phantomsection\label{fig:sga3-exp1-p041-stabilizer-square}
\end{figure}

\textit{Then \(\phi:Y\to X\) is equivariant under
\(\operatorname{Stab}_G(\phi)\).}

\subsection*{\SGAInlineRefTarget{sga3:i:subsection:6-4}6.4. Global sections}

Let \(M\) be a \(G\)-equivariant \(\mathbf O_X\)-module. Let \(S_0\)
denote the final object of \(\mathscr C\), and denote by
\(\prod_{X/S_0}M\) the \emph{functor of sections of \(M\) over \(X\)}
(cf. \SGARefLink{sga3:ii:definition:1-1}{Expos\'e~II, 1.1}): to every \(T\in\operatorname{Ob}(\mathscr C)\)
it assigns
\[
  \operatorname{Hom}_X(X_T,M)
  =\operatorname{Hom}_{X_T}(X_T,M_T)
  =\Gamma(M_T/X_T).
\]
Recall also that every morphism \(g:Z\to Y\) of
\(\widehat{\mathscr C}\)-objects over \(X\) induces a morphism of
abelian groups
\[
  M(g):M(Y)\longrightarrow M(Z),
\]
compatible with the ring homomorphism \(g^*:\mathbf O(Y)\to\mathbf
O(Z)\). In particular, when \(Z=Y\), so that \(g\) is an
\(X\)-endomorphism of \(Y\), one obtains a morphism of abelian groups
\[
  M(g):M(Y)\longrightarrow M(Y)
\]
which is not generally a morphism of \(\mathbf O(Y)\)-modules, but
which satisfies
\[
  M(g)(\alpha\cdot m)=g^*(\alpha)\cdot M(g)(m)
\]
for every \(m\in M(Y)\) and \(\alpha\in\mathbf O(Y)\). In what follows,
we simply write \(g^*\) instead of \(M(g)\).

Let \(T\in\operatorname{Ob}(\mathscr C)\), put \(X_T=X\times T\), and
let \(\mathrm{pr}_X:X_T\to X\) be the projection. For
\(\alpha\in\mathbf O(X_T)\) and \(g\in G(T)\), set
\[
  g(\alpha)=(g^{-1})^*(\alpha).
\]
Set-theoretically, this element of \(\mathbf O(X_T)\) is defined, for
every \(S\in\operatorname{Ob}(\mathscr C)\), \(x\in X(S)\), and
\(t\in T(S)\), by
\[
  g(\alpha)(x,t)=\alpha(g^{-1}x,t).
\]
This gives a left action of \(G(T)\) by ring automorphisms on
\(\mathbf O(X_T)\), functorially in \(T\), and \(g(\alpha)=\alpha\)
when \(\alpha\) is the image in \(\mathbf O(X_T)\) of an element of
\(\mathbf O(T)\).

Let \(\phi\) now be the identity morphism of \(X_T\) (cf. \SGARefLink{sga3:i:subsection:6-3}{6.3} and the
generalization below to a morphism \(\phi:Y\to X\)). Denote
\(\operatorname{Hom}_X(X_T,M)\) by \(M(\phi g^{-1})\) or \(M(\phi)\)
according as \(X_T\) is regarded as an \(X\)-object through
\(\mathrm{pr}_X\circ\lambda(g^{-1})_T\) or through
\(\mathrm{pr}_X\), respectively.

% Source page: Exp1 component-local p. 42; combined-reader p. 52.
Then \(\lambda(g^{-1})_T\) is an \(X\)-morphism between these two
\(X\)-objects. Hence there is a morphism of abelian groups
\[
  (g^{-1})^*:M(\phi)\longrightarrow M(\phi g^{-1}),
  \qquad
  m\longmapsto m\circ\lambda(g^{-1})_T,
\]
and
\[
  (g^{-1})^*(\alpha\cdot m)
  =(g\alpha)\cdot(g^{-1})^*(m)
\]
for \(m\in M(\phi)\) and \(\alpha\in\mathbf O(X_T)\). If \(m\) is a
section of \(M_T\) over \(X_T\), then \((g^{-1})^*(m)\) is the section
defined set-theoretically by
\((x,t)\mapsto m(g^{-1}x,t)\). In particular, \((g^{-1})^*\) is a
morphism of \(\mathbf O(T)\)-modules.

The functoriality of the \(\mathbf O(X_T)\)-module morphisms
\(\Lambda_x(g)\) gives the following commutative diagram.

\begin{figure}[H]
  \centering
  % BEGIN INCLUDED SGA3_English_current_buildable_source_20260823/source/inputs/I/native_diagrams/exp1/diag_022_section_naturality_square.tex
\begin{tikzcd}[column sep=4.5em,row sep=3.3em]
  M(\phi)
    \arrow[r,"(g^{-1})^*"]
    \arrow[d,"\Lambda_\phi(g)"']
  &
  M(\phi g^{-1})
    \arrow[d,"\Lambda_{\phi g^{-1}}(g)"]
  \\
  M(g\phi)
    \arrow[r,"(g^{-1})^*"]
  &
  M(g\phi g^{-1})
\end{tikzcd}

% END INCLUDED SGA3_English_current_buildable_source_20260823/source/inputs/I/native_diagrams/exp1/diag_022_section_naturality_square.tex

  \phantomsection\label{fig:sga3-exp1-p042-section-naturality-square}
\end{figure}

Moreover, \(g\phi g^{-1}=\phi\), because \(\phi\) is the identity.
Writing
\[
  A(g)=(g^{-1})^*\circ\Lambda_\phi(g),
\]
we obtain a morphism of abelian groups
\[
  A(g):M(\phi)\longrightarrow M(\phi)
\]
which is compatible with the action of \(G(T)\) on
\(\mathbf O(X_T)\), namely
\[
  A(g)(\alpha\cdot m)=(g\alpha)\cdot A(g)(m).
\]

Finally, if \(h\) is a second element of \(G(T)\), functoriality of
\(M\) and of the morphisms \(\Lambda_x(g)\) gives the following
commutative diagram.

\begin{figure}[H]
  \centering
  % BEGIN INCLUDED SGA3_English_current_buildable_source_20260823/source/inputs/I/native_diagrams/exp1/diag_023_section_action_composition.tex
\begin{tikzcd}[column sep=4.5em,row sep=3.3em]
  M(\phi)
    \arrow[d,"\Lambda_\phi(g)"']
    \arrow[dr,"\Lambda_\phi(hg)"]
  &
  \\
  M(g\phi)
    \arrow[r,"\Lambda_\phi(h)"]
    \arrow[d,"(g^{-1})^*"']
  &
  M(hg\phi)
    \arrow[d,"(g^{-1})^*"']
    \arrow[dr,"((hg)^{-1})^*"]
  &
  \\
  M(\phi)
    \arrow[r,"\Lambda_\phi(h)"]
  &
  M(h\phi)
    \arrow[r,"(h^{-1})^*"]
  &
  M(\phi)
\end{tikzcd}

% END INCLUDED SGA3_English_current_buildable_source_20260823/source/inputs/I/native_diagrams/exp1/diag_023_section_action_composition.tex

  \phantomsection\label{fig:sga3-exp1-p042-section-action-composition}
\end{figure}

It follows that \(A(h)\circ A(g)=A(hg)\), and hence:

\medskip
\noindent\SGAInlineRefTarget{sga3:i:proposition:6-4-1}\textbf{Proposition 6.4.1.}
\textit{For every \(T\), the \(\mathbf O(X_T)\)-module
\[
  \left(\prod_{X/S_0}M\right)(T)=\Gamma(M_T/X_T)
\]
is endowed, functorially in \(T\), with an action of \(G(T)\) compatible
with the action of \(G(T)\) on \(\mathbf O(X_T)\); that is,}
\[
  A(g)(\alpha\cdot m)=g(\alpha)\cdot A(g)(m).
\]

Since \(g(\alpha)=\alpha\) for \(\alpha\in\mathbf O(T)\), it follows in
particular that \(\prod_{X/S_0}M\) is a
\(G\)-\(\mathbf O_{S_0}\)-module.

More generally, let \(Y\) be a second \(G\)-equivariant
\(\widehat{\mathscr C}\)-object, let \(\phi:Y\to X\) be a morphism in
\(\widehat{\mathscr C}\), and put \(H=\operatorname{Stab}_G(\phi)\), as
in~\SGARefLink{sga3:i:subsection:6-3}{6.3}. Then
\[
  M_\phi=Y\times_\phi M
\]
is an \(H\)-equivariant \(\mathbf O_Y\)-module. We obtain:

% Source page: Exp1 component-local p. 43; combined-reader p. 53.
\medskip
\noindent\SGAInlineRefTarget{sga3:i:corollary:6-4-2}\textbf{Corollary 6.4.2.}
\textit{The functor \(\prod_{Y/S_0}M_\phi\) is a
\(\operatorname{Stab}_G(\phi)\)-\(\mathbf O_{S_0}\)-module.}

\subsection*{\SGAInlineRefTarget{sga3:i:subsection:6-5}6.5. \(G\)-equivariant \(\mathcal O_X\)-modules}

Apply the preceding discussion in the following situation. Let \(S\)
be a scheme, let \(G\) be an \(S\)-group scheme acting on an
\(S\)-scheme \(X\), and let \(\mathcal F\) be an
\(\mathcal O_X\)-module, not necessarily quasi-coherent.

\medskip
\noindent\SGAInlineRefTarget{sga3:i:definition:6-5-1}\textbf{Definition 6.5.1.}
\textit{We say that \(\mathcal F\) is a \(G\)-equivariant
\(\mathcal O_X\)-module if the \(\mathbf O_X\)-module
\(M=W(\mathcal F)\) is \(G\)-equivariant. Thus, for every morphism
\((g,x):U\to G\times_S X\), one is given isomorphisms of
\(\mathcal O_U\)-modules}
\[
  \Lambda_x(g):
  W\bigl(x^*(\mathcal F)\bigr)
  \xrightarrow{\ \sim\ }
  W\bigl((gx)^*(\mathcal F)\bigr),
\]
\textit{functorial in \(U\) and satisfying
\(\Lambda_{hx}(g)\circ\Lambda_x(h)=\Lambda_x(gh)\).}
Because \(W\) is fully faithful (cf. \SGARefLink{sga3:i:remark:4-6-1-1}{4.6.1.1} and \SGARefLink{sga3:i:proposition:4-6-2}{4.6.2}), these yield
isomorphisms of \(\mathcal O_U\)-modules
\[
  \Lambda_x(g):x^*(\mathcal F)
  \xrightarrow{\ \sim\ }(gx)^*(\mathcal F),
\]
where \(gx\) denotes \(\lambda\circ(g,x):U\to X\). Applying this to
the identity morphism of \(G\times_S X\) gives an isomorphism of
\(\mathcal O_{G\times_S X}\)-modules
\[
  (\star)\qquad
  \theta:\mathrm{pr}_X^*(\mathcal F)
  \xrightarrow{\ \sim\ }\lambda^*(\mathcal F)
\]
such that the following diagram \((\star\star)\) of morphisms of
sheaves on \(G\times_SG\times_SX\) commutes.

\begin{figure}[H]
  \centering
  % BEGIN INCLUDED SGA3_English_current_buildable_source_20260823/source/inputs/I/native_diagrams/exp1/diag_024_sheaf_cocycle.tex
\begin{tikzcd}[column sep=1.8em,row sep=3.5em]
  \mathrm{pr}_{2,3}^{*}\circ\mathrm{pr}_X^{*}(\mathcal F)
    \arrow[rr,"\mathrm{pr}_{2,3}^{*}(\theta)","\sim"']
    \arrow[d,equal]
  &&
  \mathrm{pr}_{2,3}^{*}\circ\lambda^{*}(\mathcal F)
    \arrow[rr,equal]
  &&
  (\operatorname{id}_G\times\lambda)^{*}\circ\mathrm{pr}_X^{*}(\mathcal F)
    \arrow[d,"(\operatorname{id}_G\times\lambda)^*(\theta)","\sim"']
  \\
  (\mu\times\operatorname{id}_X)^{*}\circ\mathrm{pr}_X^{*}(\mathcal F)
    \arrow[rr,"(\mu\times\operatorname{id}_X)^*(\theta)","\sim"']
  &&
  (\mu\times\operatorname{id}_X)^{*}\circ\lambda^{*}(\mathcal F)
    \arrow[rr,equal]
  &&
  (\operatorname{id}_G\times\lambda)^{*}\circ\lambda^{*}(\mathcal F)
  \qquad .
\end{tikzcd}

% END INCLUDED SGA3_English_current_buildable_source_20260823/source/inputs/I/native_diagrams/exp1/diag_024_sheaf_cocycle.tex

  \phantomsection\label{fig:sga3-exp1-p043-sheaf-cocycle}
\end{figure}

If \(\mathcal E\) is a second \(G\)-equivariant
\(\mathcal O_X\)-module, a morphism
\(\phi:\mathcal E\to\mathcal F\) of \(\mathcal O_X\)-modules is called
\emph{\(G\)-equivariant} if \(W(\phi):W(\mathcal E)\to W(\mathcal F)\)
is \(G\)-equivariant. With the notation above, this is equivalent to
\[
  \theta_{\mathcal F}\circ\mathrm{pr}_X^*(\phi)
  =\lambda^*(\phi)\circ\theta_{\mathcal E}.
\]

\medskip
\noindent\SGAInlineRefTarget{sga3:i:remark:6-5-2}\textbf{Remark 6.5.2.}
If \(f:Y\to X\) is \(G\)-equivariant, \SGARefLink{sga3:i:lemma:6-3-1}{Lemma~6.3.1} shows that
\(f^*(\mathcal F)\) is a \(G\)-equivariant
\(\mathcal O_Y\)-module.

\medskip
\noindent\SGAInlineRefTarget{sga3:i:remark:6-5-3}\textbf{Remark 6.5.3.}
Suppose in addition that \(\mathcal F\) is quasi-coherent. Giving a
\(G\)-equivariant \(\mathcal O_X\)-module structure on
\(M=W(\mathcal F)\) is then equivalent to giving such a structure on
\(V(\mathcal F)\), and this in turn is equivalent to giving an action
of \(G\) on the vector fibration \(V(\mathcal F)\), compatible with
the action on \(X\) and linear on the fibers.

Indeed, let
\[
  \varphi:\lambda^*(\mathcal F)
  \xrightarrow{\ \sim\ }\mathrm{pr}_X^*(\mathcal F)
\]
be the inverse of \(\theta\). For every morphism
\((g,x):U\to G\times_S X\), the isomorphism of
\(\mathcal O_U\)-modules
\[
  \varphi_x(g)
  =\Lambda_x(g)^{-1}
  =\Lambda_{gx}(g^{-1})
\]
from \((gx)^*(\mathcal F)\) to \(x^*(\mathcal F)\) induces an
isomorphism of \(\mathcal O_U\)-modules
\[
  V\bigl((gx)^*(\mathcal F)\bigr)
  \xrightarrow{\ \sim\ }V\bigl(x^*(\mathcal F)\bigr).
\]

% Source page: Exp1 component-local p. 44; combined-reader p. 54.
Denote the latter by \({}^{t}\!\varphi_x(g)\). Then
\[
  {}^{t}\!\varphi_{gx}(h)\circ{}^{t}\!\varphi_x(g)
  =
  {}^{t}\!\bigl(\Lambda_{gx}(h)\circ\varphi_x(g)\bigr)^{-1}
  =
  \Lambda_x(hg)^{-1}
  =
  {}^{t}\!\varphi_x(hg).
\]
For every \(X\)-scheme \(f:Y\to X\), one has
\[
  V(\mathcal F)\times_fY\simeq V\bigl(f^*(\mathcal F)\bigr)
  \qquad\text{(cf. 4.6.2).}
\]
Therefore the isomorphism
\[
\begin{aligned}
  {}^{t}\!\varphi:\quad
  G\times_SV(\mathcal F)
  &=(G\times_SX)\times_{\mathrm{pr}_X}V(\mathcal F)
   =V\bigl(\mathrm{pr}_X^*(\mathcal F)\bigr)\\
  &\xrightarrow{\ \sim\ }
   V\bigl(\lambda^*(\mathcal F)\bigr)
   =(G\times_SX)\times_\lambda V(\mathcal F)
\end{aligned}
\]
endows \(V(\mathcal F)\) with a \(G\)-equivariant
\(\mathcal O_X\)-module structure. Identifying every functor \(V(-)\)
with the vector fibration which represents it, and composing
\({}^{t}\!\varphi\) with the projection onto \(V(\mathcal F)\), gives
an action of \(G\) on \(V(\mathcal F)\), compatible with the action on
\(X\) and linear on the fibers.

This recovers the definition in [GIT, Chapter~1, \SGARefLink{sga3:i:section:3}{Section~3}], except that
Mumford there considers a locally free \(\mathcal O_X\)-module
\(\mathcal E\) of finite rank and defines an action of \(G\) on
\(V(\mathcal E)=W(\mathcal E^\vee)\). Indeed, the preceding diagram
\((\star\star)\), with \(\theta\) replaced by \(\varphi\) and the
directions of the arrows reversed, is exactly the diagram in
\emph{loc. cit.}, Definition~1.6, and the isomorphism
\({}^{t}\!\varphi\) above coincides with the isomorphism \(\Phi\) of
\emph{loc. cit.}, p.~31.

\medskip
\noindent\SGAInlineRefTarget{sga3:i:remark:6-5-4}\textbf{Remark 6.5.4.}
Consider in particular the case \(X=S\), with the trivial action of
\(G\). A \(G\)-equivariant \(\mathcal O_S\)-module \(\mathcal F\) is
then the same thing as a \(G\)-\(\mathcal O_S\)-module (cf. \SGARefLink{sga3:i:definition:4-7-1}{4.7.1}).
If \(f:G\to S\) denotes the morphism which here equals both
\(\mathrm{pr}_X\) and \(\lambda\), the isomorphism
\[
  (\star)\qquad
  \theta:f^*(\mathcal F)\xrightarrow{\ \sim\ }f^*(\mathcal F)
\]
is an element of
\[
\begin{aligned}
  \operatorname{Hom}_{\mathcal O_G}
    \bigl(f^*(\mathcal F),f^*(\mathcal F)\bigr)
  &=
  \operatorname{Hom}_{\mathcal O_G}
    \bigl(W(f^*(\mathcal F)),W(f^*(\mathcal F))\bigr)\\
  &=
  \operatorname{End}_{\mathcal O_S}\bigl(W(\mathcal F)\bigr)(G).
\end{aligned}
\]
It is precisely the morphism
\[
  \rho:G\longrightarrow
  \operatorname{End}_{\mathcal O_S}\bigl(W(\mathcal F)\bigr)
\]
which defines the action of \(G\) on \(W(\mathcal F)\). Moreover,
\(\theta\) corresponds by adjunction to the morphism of
\(\mathcal O_S\)-modules
\[
  f_*(\theta)\circ\tau:
  \mathcal F\longrightarrow f_*f^*(\mathcal F),
\]
where \(\tau:\mathcal F\to f_*f^*(\mathcal F)\) is the unit morphism
of the adjunction. This will be used in Expos\'e~VI\textsubscript{B},
\SGARefLink{cor:sga3-VIB-11.10}{11.10}~bis.

\subsection*{\SGAInlineRefTarget{sga3:i:subsection:6-6}6.6. The functors
\(\prod_{X/S}W(\mathcal F)\) and \(W(p_*(\mathcal F))\)}

Let \(S\) be a scheme, let \(G\) be an \(S\)-group scheme acting on
\(S\)-schemes \(X\) and \(Y\) through
\[
  \lambda:G\times_SX\to X,
  \qquad
  \mu:G\times_SY\to Y,
\]
and let \(p:X\to Y\) be a \(G\)-equivariant morphism. Suppose that
\(p:X\to Y\) and \(\nu_Y:Y\to S\) are quasi-compact and
quasi-separated and that \(\pi:G\to S\) is flat.

The projection \(\mathrm{pr}_Y:G\times_SY\to Y\) is then flat, as is
\(\mu\), since \(\mu\) is the composite of \(\mathrm{pr}_Y\) and the
automorphism \((g,y)\mapsto(g,gy)\). For a variable \(S\)-scheme
\(f:T\to S\), denote by \(p_T:X_T\to T\) and \(f_X:X_T\to X\) the
morphisms obtained from \(p\) and \(f\) by base change.

Let \(\mathcal F\) be a quasi-coherent \(G\)-equivariant
\(\mathcal O_X\)-module, and let
\[
  \theta:\mathrm{pr}_X^*(\mathcal F)
  \xrightarrow{\ \sim\ }\lambda^*(\mathcal F)
\]
be the isomorphism of \SGARefLink{sga3:i:definition:6-5-1}{6.5.1} \((\star)\). Since \(p\) is quasi-compact
and quasi-separated, \(p_*(\mathcal F)\) is quasi-coherent
(cf. EGA~I, 9.2.1).

% Source page: Exp1 component-local p. 45; combined-reader p. 55.
\medskip
\noindent\SGAInlineRefTarget{sga3:i:lemma:6-6-1}\textbf{Lemma 6.6.1.}
\textit{The quasi-coherent \(\mathcal O_Y\)-module
\(p_*(\mathcal F)\) is \(G\)-equivariant.}

Indeed, one has the following two cartesian squares.

\begin{figure}[H]
  \centering
  % BEGIN INCLUDED SGA3_English_current_buildable_source_20260823/source/inputs/I/native_diagrams/exp1/diag_025_base_change_squares.tex
\begin{tikzcd}[column sep=3.8em,row sep=3.2em]
  G\times_S X
    \arrow[r,"\mathrm{pr}_X"]
    \arrow[d,"q"']
  &
  X
    \arrow[d,"p"]
  &
  G\times_S X
    \arrow[l,"\lambda"']
    \arrow[d,"q"]
  \\
  G\times_S Y
    \arrow[r,"\mathrm{pr}_Y"]
  &
  Y
  &
  G\times_S Y
    \arrow[l,"\mu"']
  \qquad .
\end{tikzcd}

% END INCLUDED SGA3_English_current_buildable_source_20260823/source/inputs/I/native_diagrams/exp1/diag_025_base_change_squares.tex

  \phantomsection\label{fig:sga3-exp1-p045-base-change-squares}
\end{figure}

Since \(p\) is quasi-compact and quasi-separated and
\(\mathrm{pr}_Y\) and \(\mu\) are flat, EGA~III, 1.4.15, completed by
EGA~IV\textsubscript{1}, 1.7.21, gives
\[
  q_*\mathrm{pr}_X^*(\mathcal F)
    =\mathrm{pr}_Y^*p_*(\mathcal F),
  \qquad
  q_*\lambda^*(\mathcal F)
    =\mu^*p_*(\mathcal F).
\]
Consequently, \(\theta\) induces an isomorphism
\[
  \theta_Y:\mathrm{pr}_Y^*p_*(\mathcal F)
  \xrightarrow{\ \sim\ }\mu^*p_*(\mathcal F).
\]
The same argument shows that \(\theta_Y\) satisfies the cocycle
condition \SGARefLink{sga3:i:definition:6-5-1}{6.5.1} \((\star\star)\), because the morphisms
\(G\times_SG\times_SY\to G\times_SY\) involved are flat. This proves
the lemma.

In particular, take \(Y=S\) with the trivial action of \(G\). By
\SGARefLink{sga3:i:remark:6-5-4}{Remark~6.5.4}, \(p_*(\mathcal F)\) is then a
\(G\)-\(\mathcal O_S\)-module. If, in addition, \(G\) is affine over
\(S\), and if
\[
  \mathscr A(G)=\pi_*(\mathcal O_G),
\]
then \(p_*(\mathcal F)\) is an \(\mathscr A(G)\)-comodule, by~\SGARefLink{sga3:i:proposition:4-7-2}{4.7.2}.

On the other hand, by \SGARefLink{sga3:i:proposition:6-4-1}{Proposition~6.4.1}, the functor
\(\prod_{X/S}W(\mathcal F)\), which assigns to every \(f:T\to S\)
\[
  W\bigl(f_X^*(\mathcal F)\bigr)(X_T)
  =\Gamma\bigl(X_T,f_X^*(\mathcal F)\bigr)
  =\Gamma\bigl(T,p_{T*}f_X^*(\mathcal F)\bigr),
\]
is a \(G\)-\(\mathcal O_S\)-module. There is moreover a canonical
morphism
\[
  \tau:W\bigl(p_*(\mathcal F)\bigr)
  \longrightarrow\prod_{X/S}W(\mathcal F),
\]
which, for each \(f:T\to S\), is given by the canonical morphism
\[
  \Gamma\bigl(T,f^*p_*(\mathcal F)\bigr)
  \longrightarrow
  \Gamma\bigl(T,p_{T*}f_X^*(\mathcal F)\bigr).
\]
It is an isomorphism when restricted to the full subcategory of schemes
\(T\) flat over \(S\), and it is immediate that \(\tau\) is a morphism
of \(G\)-\(\mathcal O_S\)-modules. We have therefore proved the
following proposition; for part~(ii), compare [GIT], p.~32.

\medskip
\noindent\SGAInlineRefTarget{sga3:i:proposition:6-6-2}\textbf{Proposition 6.6.2.}
\textit{Let \(S\) be a scheme, let \(G\) be an \(S\)-group scheme
acting on an \(S\)-scheme \(X\), and let \(\mathcal F\) be a
quasi-coherent \(G\)-equivariant \(\mathcal O_X\)-module. Suppose that
\(\pi:G\to S\) is flat and \(p:X\to S\) is quasi-compact and
quasi-separated.}

\textit{(i) Then \(p_*(\mathcal F)\) is a quasi-coherent
\(G\)-\(\mathcal O_S\)-module. The canonical morphism}
\[
  W\bigl(p_*(\mathcal F)\bigr)
  \longrightarrow\prod_{X/S}W(\mathcal F)
\]
\textit{is a morphism of \(G\)-\(\mathcal O_S\)-modules, and the two
functors coincide on the category of flat \(S\)-schemes.}

\textit{(ii) If, in addition, \(G\) is affine over \(S\), and if
\(\mathscr A(G)=\pi_*(\mathcal O_G)\), then \(p_*(\mathcal F)\) has
the structure of an \(\mathscr A(G)\)-comodule.}\footnote{Editor's
note: The same holds if \(G\) is flat, quasi-compact, and
quasi-separated over \(S\), and if \(p_*(\mathcal F)\) is a flat
\(\mathcal O_S\)-module; cf. Expos\'e~VI\textsubscript{B}, \SGARefLink{concl:sga3-VIB-11.6.1-ii}{11.6.1(ii)}.}

% Source page: Exp1 component-local p. 46; combined-reader p. 56.
\subsection*{\SGAInlineRefTarget{sga3:i:subsection:6-7}6.7. Stabilizers}

Let \(S\) be a scheme, let \(G\) be an \(S\)-group scheme acting on an
\(S\)-scheme \(X\), and let \(\mathcal F\) be a quasi-coherent
\(G\)-equivariant \(\mathcal O_X\)-module. Let \(Y\) be a second
\(S\)-scheme with an action of \(G\), possibly the trivial action, let
\(\phi:Y\to X\) be an \(S\)-morphism, not necessarily
\(G\)-equivariant, and let \(\operatorname{Stab}_G(\phi)\) be the
stabilizer of \(\phi\) in \(G\) (cf. \SGARefLink{sga3:i:definition:6-3-2}{6.3.2}).

Let us point out at once (see Expos\'e~VI\textsubscript{B},
\SGARefLink{sec:sga3-VIB-6}{Section~6}) that \(\operatorname{Stab}_G(\phi)\) is representable by a
closed subgroup scheme \(H\) of \(G\) if \(X\) is separated over
\(S\) and \(Y\) is essentially free over \(S\) (cf.
\emph{loc. cit.}, \SGARefLink{def:sga3-VIB-6.2.1}{Definition~6.2.1}). Indeed, consider
\[
  r:G\times_SY\longrightarrow X\times_SX,
  \qquad
  r(g,y)=\bigl(\phi(y),g\phi(g^{-1}y)\bigr),
\]
put \(P=G\times_SY\), and let \(P'\) be the inverse image under \(r\)
of the diagonal \(\Delta_{X/S}\). Then (cf. \emph{loc. cit.},
\SGARefLink{ex:sga3-VIB-6.2.4-a}{6.2.4(a)})
\[
  \operatorname{Stab}_G(\phi)=\prod_{P/G}P'.
\]
It follows that \(\operatorname{Stab}_G(\phi)\) is representable by a
closed subgroup scheme \(H\) of \(G\) if \(X\) is separated over
\(S\) and \(Y\) is essentially free over \(S\). The latter condition
is automatic if \(S\) is the spectrum of a field or if \(Y=S\).
Under these hypotheses, \(\phi^*(\mathcal F)\) is a quasi-coherent
\(H\)-equivariant \(\mathcal O_Y\)-module (cf. \SGARefLink{sga3:i:remark:6-5-2}{6.5.2}).

If, in addition, \(\pi:G\to S\) and \(p:Y\to S\) are quasi-compact
and quasi-separated over \(S\), and \(\pi\) is flat, then
\(p_*\phi^*(\mathcal F)\) is an \(H\)-\(\mathcal O_S\)-module by
\SGARefLink{sga3:i:proposition:6-6-2}{Proposition~6.6.2}. In particular:

\medskip
\noindent\SGAInlineRefTarget{sga3:i:corollary:6-7-1}\textbf{Corollary 6.7.1.}
\textit{Let \(S\) be a scheme, let \(G\) be an \(S\)-group scheme
acting on an \(S\)-scheme \(X\), and let \(\mathcal F\) be a
quasi-coherent \(G\)-equivariant \(\mathcal O_X\)-module. Suppose that
\(\pi:G\to S\) is flat and \(X\to S\) is quasi-compact and separated.
Let \(\tau:S\to X\) be a section of \(X\) over \(S\).}

\textit{(i) The stabilizer \(H=\operatorname{Stab}_G(\tau)\) is a
closed subgroup scheme of \(G\), and \(\tau^*(\mathcal F)\) is a
quasi-coherent \(H\)-\(\mathcal O_S\)-module.}

\textit{(ii) If \(H\) is moreover affine over \(S\), for example if
\(G\) is, and if \(\mathscr A(H)=\pi_*(\mathcal O_H)\), then
\(\tau^*(\mathcal F)\) is an \(\mathscr A(H)\)-comodule.}

\subsection*{\SGAInlineRefTarget{sga3:i:subsection:6-8}6.8. \(G\)-equivariant sheaves on \(G\)}

To conclude, we give two results, \SGARefLink{sga3:i:proposition:6-8-1}{6.8.1} and \SGARefLink{sga3:i:proposition:6-8-6}{6.8.6}, which will be used
in Expos\'es~II and~III, in particular in III, \SGARefLink{sga3:iii:lemma:4-25}{4.25}.

\medskip
\noindent\SGAInlineRefTarget{sga3:i:proposition:6-8-1}\textbf{Proposition 6.8.1.}
\textit{Let \(S\) be a scheme, let \(\pi:G\to S\) be an \(S\)-group
scheme, and let \(\varepsilon:S\to G\) be the unit section. Consider
the action of \(G\) on itself by left translations. The functors}
\[
  \mathcal E\longmapsto\pi^*(\mathcal E),
  \qquad
  \mathcal F\longmapsto\varepsilon^*(\mathcal F)
\]
\textit{induce mutually quasi-inverse equivalences between the category
of quasi-coherent \(\mathcal O_S\)-modules and the category of
quasi-coherent \(G\)-equivariant \(\mathcal O_G\)-modules.}

\textit{Proof.}
Let \(\mu\) denote multiplication in \(G\), and let
\(\mathrm{pr}_2:G\times_SG\to G\) be the second projection. Since
\(\pi\circ\mu=\pi\circ\mathrm{pr}_2\), every quasi-coherent
\(\mathcal O_S\)-module \(\mathcal E\) has a canonical isomorphism
\[
  \mu^*\pi^*(\mathcal E)
  =\mathrm{pr}_2^*\pi^*(\mathcal E).
\]
It is immediate that this isomorphism satisfies the cocycle condition
\SGARefLink{sga3:i:definition:6-5-1}{6.5.1} \((\star\star)\), so \(\pi^*(\mathcal E)\) is a
\(G\)-equivariant \(\mathcal O_G\)-module. Since
\(\varepsilon^*\pi^*(\mathcal E)=\mathcal E\), the functor
\(\mathcal E\mapsto\pi^*(\mathcal E)\) is

% Source page: Exp1 component-local p. 47; combined-reader p. 57.
fully faithful. It remains to show that, for every \(G\)-equivariant
\(\mathcal O_G\)-module \(\mathcal F\),
\[
  \mathcal F\simeq\pi^*\varepsilon^*(\mathcal F).
\]
By hypothesis, there is an isomorphism
\[
  \theta:\mathrm{pr}_2^*(\mathcal F)
  \xrightarrow{\ \sim\ }\mu^*(\mathcal F).
\]
Pulling \(\theta\) back along the morphism
\[
  \tau:G\longrightarrow G\times_SG,
  \qquad
  \tau=(\operatorname{id}_G,\varepsilon\circ\pi),
\]
gives an isomorphism
\(\mathcal F\xrightarrow{\sim}\pi^*\varepsilon^*(\mathcal F)\).

\medskip
\noindent\SGAInlineRefTarget{sga3:i:remark:6-8-2}\textbf{Remarks 6.8.2.}
\SGARefLink{sga3:i:definition:6-1:item:a}{(a)} Consider the action of \(G\times_SG\) on \(G\) defined by
\[
  (g_1,g_2)\cdot g=g_1gg_2^{-1}.
\]
The stabilizer of the unit section \(\varepsilon:S\to G\) is the
diagonal subgroup \(H\) of \(G\times_SG\). Consequently, if
\(\mathcal F\) is a quasi-coherent
\((G\times_SG)\)-equivariant \(\mathcal O_G\)-module, then
\(\mathcal E=\varepsilon^*(\mathcal F)\) has the structure of an
\(H\)-\(\mathcal O_S\)-module, by~\SGARefLink{sga3:i:remark:6-5-2}{6.5.2}.

(b) One can show that
\(\mathcal F\mapsto\varepsilon^*(\mathcal F)\) is an equivalence
between the category of quasi-coherent
\((G\times_SG)\)-equivariant \(\mathcal O_G\)-modules and the category
of quasi-coherent \(H\)-\(\mathcal O_S\)-modules. This is a special
case of more general descent results (cf. Expos\'e~IV, Section~2, and
SGA~1, VIII); see, for example, [Th87], 1.2--1.3.

\medskip
\noindent\SGAInlineRefTarget{sga3:i:remark:6-8-3}\textbf{Remark 6.8.3.}
Retain the notation of \SGARefLink{sga3:i:proposition:6-8-1}{6.8.1}. For every quasi-coherent
\(\mathcal O_S\)-module \(\mathcal E\), let
\(\pi_*^G\pi^*(\mathcal E)\) be the \(\mathcal O_S\)-submodule of
\(\pi_*\pi^*(\mathcal E)\) whose sections over an open subset \(V\)
of \(S\) are the elements
\[
  \gamma\in\Gamma\bigl(\pi^{-1}(V),\pi^*(\mathcal E)\bigr)
\]
such that \(g\cdot\gamma_{S'}=\gamma_{S'}\) for every \(S'\to V\) and
\(g\in G(S')\). Then the natural morphism
\[
  \mathcal E\longrightarrow\pi_*^G\pi^*(\mathcal E)
\]
is an isomorphism. This is immediate if
\[
  \pi_*\pi^*(\mathcal E)
  =\mathcal E\otimes_{\mathcal O_S}\pi_*(\mathcal O_G),
\]
for example if \(G\to S\) is affine, or if \(G\to S\) is
quasi-compact and quasi-separated and \(\mathcal E\) is flat; the
general case is verified without difficulty.

\medskip
\noindent\SGAInlineRefTarget{sga3:i:remark:6-8-4}\textbf{Remark 6.8.4.}
Let \(S\) be a scheme, let \(H\) be an \(S\)-group scheme acting on an
\(S\)-scheme \(X\), and let \(\mathcal F\) be an
\(\mathcal O_X\)-module. Suppose \(H\) is flat over \(S\), and let
\(W_P(\mathcal F)\) denote the restriction of \(W(\mathcal F)\) to the
full subcategory of flat \(S\)-schemes. By~\SGARefLink{sga3:i:definition:6-5-1}{6.5.1}, endowing
\(\mathcal F\) with an \(H\)-equivariant module structure is equivalent
to giving an isomorphism
\[
  \theta:\mathrm{pr}_X^*(\mathcal F)\longrightarrow\lambda^*(\mathcal F)
\]
of sheaves on \(H\times_SX\) which satisfies the cocycle condition
\((\star\star)\). Hence, to endow \(\mathcal F\) with an
\(H\)-equivariant module structure, it suffices to endow
\(W_P(\mathcal F)\) with such a structure.

\medskip
\SGAInlineRefTarget{sga3:i:recollection:6-8-5}\noindent\textbf{Recollection \SGARefLink{sga3:i:recollection:6-8-5}{6.8.5}.}
Let \(X\) be an \(S\)-scheme and \(Y\) a closed \(S\)-subscheme of
\(X\). Let \(\mathcal N_{Y/X}\) denote the conormal sheaf of the
immersion \(i:Y\hookrightarrow X\) (cf. EGA~IV\textsubscript{4},
16.1.2). If \(S'\to S\) is flat and
\(i':Y'\hookrightarrow X'\) is obtained from \(i\) by base change,
then \emph{loc. cit.}, 16.2.2(iii), gives
\[
  \mathcal N_{Y/X}\otimes_{\mathcal O_Y}\mathcal O_{Y'}
  =\mathcal N_{Y'/X'}.
\]

\medskip
\noindent\SGAInlineRefTarget{sga3:i:proposition:6-8-6}\textbf{Proposition 6.8.6.}
\textit{Let \(S\) be a scheme, let \(X\) be an \(S\)-group scheme, and
let \(Y\) be an \(S\)-subgroup scheme of \(X\). Suppose that \(Y\) is
flat over \(S\). Then the conormal sheaf \(\mathcal N_{Y/X}\) is a
\((Y\times_SY)\)-equivariant \(\mathcal O_Y\)-module.}

Indeed, \(H=Y\times_SY\) is flat over \(S\). By \SGARefLink{sga3:i:remark:6-8-4}{Remark~6.8.4}, it is
therefore enough to endow \(W_P(\mathcal N_{Y/X})\) with an
\(H\)-equivariant module structure. Let \(S'\) be a flat \(S\)-scheme,
let \(Y'\hookrightarrow X'\) be the immersion obtained by base change,
and put
\[
  \mathcal N'
  =\mathcal N_{Y/X}\otimes_{\mathcal O_Y}\mathcal O_{Y'}.
\]
By~\SGARefLink{sga3:i:recollection:6-8-5}{6.8.5}, \(\mathcal N'=\mathcal N_{Y'/X'}\).
Every \(h\in H(S')\) induces an automorphism of \(Y'\), and hence, for
every \(y\in Y'(S')\), isomorphisms
\[
  \Gamma\bigl(S',y^*(\mathcal N')\bigr)
  \xrightarrow{\ \sim\ }
  \Gamma\bigl(S',y^*h^*(\mathcal N')\bigr)
\]

% Source page: Exp1 component-local p. 48; combined-reader p. 58.
which endow \(W_P(\mathcal N_{Y/X})\) with an \(H\)-equivariant module
structure (cf. \SGARefLink{sga3:i:definition:6-1}{6.1}).

\section*{Bibliography}
\addcontentsline{toc}{section}{Bibliography}
\footnote{Editor's note: These are additional references cited in this
Expos\'e.}

\begin{description}
\item[{[DG70]}] M. Demazure and P. Gabriel,
  \emph{Groupes Alg\'ebriques}, Masson \& North-Holland, 1970.
\item[{[Gr57]}] A. Grothendieck,
  "Sur quelques points d'alg\`ebre homologique,"
  \emph{T\^ohoku Math. J.} \textbf{9} (1957), 119--221.
\item[{[Ja03]}] J. C. Jantzen,
  \emph{Representations of Algebraic Groups}, Academic Press, 1987;
  second edition, American Mathematical Society, 2003.
\item[{[GIT]}] D. Mumford,
  \emph{Geometric Invariant Theory}, Springer-Verlag, 1965;
  second edition with J. Fogarty, 1982; third edition with J. Fogarty
  and F. Kirwan, 1994.
\item[{[Ni02]}] N. Nitsure,
  "Representability of \(\mathrm{GL}_{\mathcal E}\),"
  \emph{Proc. Indian Acad. Sci.} \textbf{112} (2002), no.~4, 539--542.
\item[{[Ni04]}] N. Nitsure,
  "Representability of Hom implies flatness,"
  \emph{Proc. Indian Acad. Sci.} \textbf{114} (2004), no.~1, 7--14.
\item[{[Se68]}] J.-P. Serre,
  "Groupes de Grothendieck des sch\'emas en groupes r\'eductifs
  d\'eploy\'es,"
  \emph{Publ. Math. I.H.\'E.S.} \textbf{34} (1968), 37--52.
\item[{[Th87]}] R. W. Thomason,
  "Equivariant resolution, linearization, and Hilbert's fourteenth
  problem over arbitrary base schemes,"
  \emph{Adv. Math.} \textbf{65} (1987), 16--34.
\end{description}

% End of Expose I and of the assigned component.

% END INCLUDED SGA3_English_current_buildable_source_20260823/source/inputs/I/tex/components/09_expose_I_s6_equivariant_objects.tex



% END INCLUDED SGA3_English_current_buildable_source_20260823/source/inputs/I/body.tex


\makeatletter
\providecommand{\SGAThreeNativeInput}[2][0.92\textwidth]{%
  \begin{adjustbox}{max width=#1,center,varwidth=\maxdimen}%
    \input{#2}%
  \end{adjustbox}%
}
\providecommand{\SGAThreeNativeDiagram}[5][0.92\textwidth]{%
  \begin{figure}[H]
    \centering
    \SGAThreeNativeInput[#1]{#2}%
    \caption{#3}
    \label{#4}
    \par\smallskip
    {\footnotesize\itshape Native \TeX{} reconstruction; source locator: #5.}
  \end{figure}%
}
\makeatother

\SGAStartExpose{II}
% BEGIN INCLUDED SGA3_English_current_buildable_source_20260823/source/inputs/II/body.tex
% Mechanically extracted after \mainmatter from the frozen standalone master.
% BEGIN INCLUDED SGA3_English_current_buildable_source_20260823/source/inputs/II/tex/components/10_expose_II_s1_s2_hom_dual_numbers.tex
% SGA 3, Expose II, Loop 1 English reconstruction.
% Scope: opening through Nota 2.2.2, stopping immediately before Section 3.
% Authority: Polo--Gille Exp2-14oct24.pdf, component-local pages 1--7
% (combined-reader pages 59--65). Force-OCR is a drafting locator; English
% wording was compared with Jacob Reinhold's CC BY 4.0 Markdown lineage.

\setcounter{footnote}{-1}
\chapter*{Expos\'e II\\Tangent Bundles -- Lie Algebras}
\addcontentsline{toc}{chapter}{Expos\'e II. Tangent Bundles -- Lie Algebras}
\markboth{EXPOS\'E II. TANGENT BUNDLES -- LIE ALGEBRAS}%
  {EXPOS\'E II. TANGENT BUNDLES -- LIE ALGEBRAS}

\begin{center}
  \textit{by M. Demazure}
\end{center}

We propose in this expos\'e to construct, in the theory of schemes, the
analogues of the tangent bundles and Lie algebras of classical theory. It
will, however, be useful not to restrict ourselves to schemes proper, but
also to consider certain functors on the category of schemes that are not
necessarily representable (for example, the functors
\(\underline{\operatorname{Hom}}\), \(\underline{\operatorname{Norm}}\),
and so forth). As announced in the preceding expos\'e (cf. \SGARefLink{sga3:i:subsection:1-1}{I~1.1}), we shall
identify a scheme with its associated functor.\footnote{Editor's note:
version of October 14, 2024.}

On the other hand, the constructions presented below go beyond the setting
of scheme theory. They are equally valid, for example, in the theory of
analytic spaces with nilpotent elements, apart from a few changes of detail.

Before beginning this construction, we must give a few general definitions
that supplement those of \SGARefLink{sga3:i:subsection:1:7}{I~1.7}.

% Source page: Exp2 component-local p. 1; combined-reader p. 59.
\SGARefTarget{sga3:ii:section:1}\section*{1. The functors
  \(\underline{\operatorname{Hom}}_{Z/S}(X,Y)\)}
\addcontentsline{toc}{section}{1. The functors Hom Z/S (X,Y)}

We resume the notation of \SGARefLink{sga3:i:subsection:1-1}{I~1.1}. We identify the category \(\mathscr C\)
with a full subcategory of
\[
  \widehat{\mathscr C}
  =\operatorname{Hom}(\mathscr C^{\circ},\mathbf{Set})
\]
(in particular, we suppress the underlinings that allowed us to distinguish
graphically an object of \(\widehat{\mathscr C}\) from an object of
\(\mathscr C\)).\footnote{Editor's note: in this edition these are rendered
by bold letters \(\mathbf F,\mathbf G,\mathbf V,\mathbf W\); cf. Expos\'e~I.
For functors such as \(\underline{\operatorname{Norm}}\) and
\(\underline{\operatorname{Centr}}\) (cf. \SGARefLink{sga3:ii:subsection:5-2}{5.2}), however, the underlinings of
the original have been retained, and one has been added to the functor
\(\underline{\operatorname{Lie}}\), in order to distinguish the functor
\(\underline{\operatorname{Lie}}(G/S)\) from the
\(\Gamma(S,\mathcal O_S)\)-Lie algebra
\(\operatorname{Lie}(G/S)=\underline{\operatorname{Lie}}(G/S)(S)\), used,
for example, in Expos\'e~VII\textsubscript{A}.}

Consider the following situation: four objects of
\(\widehat{\mathscr C}\), denoted by \(S,X,Y,Z\), the first of which is in
fact an object of \(\mathscr C\), with \(X\) and \(Y\) over \(Z\), and
\(Z\) over \(S\).

\SGAThreeNativeDiagram[0.24\textwidth]
  {native_diagrams/exp2/exp2_p002_XY_over_ZS.tex}
  {The objects \(X\) and \(Y\) over \(Z\), and \(Z\) over \(S\).}
  {fig:exp2-xy-over-zs}
  {Polo--Gille Exp2-14oct24.pdf, component-local p. 2; combined-reader p. 60}

% Source page: Exp2 component-local p. 2; combined-reader p. 60.
\SGAInlineRefTarget{sga3:ii:definition:1-1}\subsection*{\SGARefLink{sga3:ii:definition:1-1}{Definition 1.1}.}
We define an object of \(\widehat{\mathscr C}_{/S}\), denoted
\(\underline{\operatorname{Hom}}_{Z/S}(X,Y)\), by
\begin{equation}
  \underline{\operatorname{Hom}}_{Z/S}(X,Y)(S')
  =\operatorname{Hom}_{Z_{S'}}(X_{S'},Y_{S'})
  =\operatorname{Hom}_{Z}(X\times_S S',Y),
  \tag{1}\label{sga3:ii:section:1:equation:1}
\end{equation}
for every object \(S'\) of \(\mathscr C_{/S}\). One sees at once that
\(\underline{\operatorname{Hom}}_{Z/S}(X,Y)\) is simply the subobject of
\(\underline{\operatorname{Hom}}_S(X,Y)\) formed by the morphisms compatible
with \(p_X\) and \(p_Y\); that is, the kernel of the pair of morphisms
\[
  \underline{\operatorname{Hom}}_S(X,Y)
  \rightrightarrows
  \underline{\operatorname{Hom}}_S(X,Z),
\]
the first defined by composition with \(p_Y\), and the second being the
constant morphism whose ``image'' is \(p_X\).

As in \SGARefLink{sga3:i:subsection:1:7}{I~1.7}, for every object \(T\) of \(\widehat{\mathscr C}\) over
\(S\), there is a natural bijection
\begin{equation}
  \operatorname{Hom}_S
    \bigl(T,\underline{\operatorname{Hom}}_{Z/S}(X,Y)\bigr)
  \simeq \operatorname{Hom}_Z(X\times_S T,Y).
  \tag{2}\label{sga3:ii:section:1:equation:2}
\end{equation}
Moreover, by \SGARefLink{sga3:i:proposition:1-7-1}{I~1.7.1}, if \(E,F\) are objects of
\(\widehat{\mathscr C}\) over \(Z\), then
\[
  \operatorname{Hom}_Z\bigl(E,\underline{\operatorname{Hom}}_Z(F,Y)\bigr)
  \simeq \operatorname{Hom}_Z(E\times_ZF,Y)
  \simeq
  \operatorname{Hom}_Z\bigl(F,\underline{\operatorname{Hom}}_Z(E,Y)\bigr).
\]
Applying this with \(E=X\) and \(F=Z\times_ST\), one obtains natural
bijections, for every object \(T\) of
\(\widehat{\mathscr C}_{/S}\),
\begin{equation}
\begin{aligned}
  \operatorname{Hom}_S
    \bigl(T,\underline{\operatorname{Hom}}_{Z/S}(X,Y)\bigr)
  &\simeq \operatorname{Hom}_Z(X\times_ST,Y)\\
  &\simeq
  \begin{cases}
    \operatorname{Hom}_Z\bigl(Z\times_ST,
      \underline{\operatorname{Hom}}_Z(X,Y)\bigr),\\
    \operatorname{Hom}_Z\bigl(X,
      \underline{\operatorname{Hom}}_Z(Z\times_ST,Y)\bigr).
  \end{cases}
\end{aligned}
\tag{3}\label{sga3:ii:section:1:equation:3}
\end{equation}
These bijections are functorial in \(T\), and hence give isomorphisms of
\(S\)-functors.\footnote{Editor's note: points \SGARefLink{sga3:ii:section:2}{(2)}, \SGARefLink{sga3:ii:section:3}{(3)}, and \SGARefLink{sga3:ii:section:4}{(4)} have been
expanded; in particular, \SGARefLink{fig:exp2-hom-adjunction-triangle}{(4)} will be used in \SGARefLink{sga3:ii:proposition:3-11}{3.11}.}

\SGAThreeNativeDiagram[0.68\textwidth]
  {native_diagrams/exp2/exp2_p002_hom_adjunction_triangle.tex}
  {The functorial isomorphisms (4).}
  {fig:exp2-hom-adjunction-triangle}
  {Polo--Gille Exp2-14oct24.pdf, component-local p. 2; combined-reader p. 60}

Let us point out two special cases of the definition. If \(Z=S\), then
\[
  \underline{\operatorname{Hom}}_{S/S}(X,Y)
  =\underline{\operatorname{Hom}}_S(X,Y).
\]

% Source page: Exp2 component-local p. 3; combined-reader p. 61.
On the other hand, when \(X=Z\), we set
\begin{equation}
  \prod_{Z/S}Y
  =\underline{\operatorname{Hom}}_{Z/S}(Z,Y),
  \tag{5}\label{sga3:ii:section:1:equation:5}
\end{equation}
so that, by definition,
\[
  \left(\prod_{Z/S}Y\right)(S')
  =\operatorname{Hom}_Z(Z\times_SS',Y)
  \simeq \Gamma(Y_{S'}/Z_{S'}).
\]
The functor
\(\prod_{Z/S}:\widehat{\mathscr C}_{/Z}\to
\widehat{\mathscr C}_{/S}\) is right adjoint to the base-change functor
from \(S\) to \(Z\): for every \(S\)-functor \(U\),
\[
  \operatorname{Hom}_S\left(U,\prod_{Z/S}Y\right)
  =\operatorname{Hom}_Z(U\times_SZ,Y).
\]
If \(\mathscr C=\mathbf{Sch}\) and \(Z\) is an \(S\)-scheme, the functor
\(\prod_{Z/S}\) is called the ``Weil restriction of scalars.''
\footnote{Editor's note: the preceding two sentences have been added.}
One also has an isomorphism
\begin{equation}
  \underline{\operatorname{Hom}}_{Z/S}(X,Y)
  \simeq
  \underline{\operatorname{Hom}}_{X/S}(X,Y\times_ZX)
  =\prod_{X/S}(Y\times_ZX),
  \tag{6}\label{sga3:ii:section:1:equation:6}
\end{equation}
which gives in particular, when \(Z=S\), an isomorphism
\begin{equation}
  \underline{\operatorname{Hom}}_S(X,Y)
  \simeq \prod_{X/S}Y_X.
  \tag{7}\label{sga3:ii:section:1:equation:7}
\end{equation}

\SGAInlineRefTarget{sga3:ii:remark:1-2}\subsection*{\SGARefLink{sga3:ii:remark:1-2}{Remark 1.2}.}
The functor
\(Y\mapsto\underline{\operatorname{Hom}}_{Z/S}(X,Y)\) commutes with
products in the following sense: there is a functorial isomorphism
\[
  \underline{\operatorname{Hom}}_{Z/S}(X,Y\times_ZY')
  \simeq
  \underline{\operatorname{Hom}}_{Z/S}(X,Y)
  \times_S
  \underline{\operatorname{Hom}}_{Z/S}(X,Y').
  \tag{*}\label{sga3:ii:section:1:equation:star}
\]
It follows that if \(Y\) is a \(Z\)-group, respectively a \(Z\)-ring,
and so on, then \(\underline{\operatorname{Hom}}_{Z/S}(X,Y)\) is an
\(S\)-group, respectively an \(S\)-ring, and so on.

\SGAInlineRefTarget{sga3:ii:remark:1-3}\subsection*{\SGARefLink{sga3:ii:remark:1-3}{Remark 1.3}.}
Let \(\pi:M\to Y\) be a \(Y\)-functor in \(\mathcal O_Y\)-modules
(cf. \SGARefLink{sga3:i:paragraph:4-3-3-1}{I~4.3.3.1}), and set
\(H=\underline{\operatorname{Hom}}_{Z/S}(X,Y)\). Then
\(\underline{\operatorname{Hom}}_{Z/S}(X,M)\) has a natural
\(\mathcal O_H\)-module structure; more precisely, for every \(H'\) over
\(H\), the set
\[
  \operatorname{Hom}_H
    \bigl(H',\underline{\operatorname{Hom}}_{Z/S}(X,M)\bigr)
\]
has a natural \(\mathcal O(X\times_SH')\)-module structure.
\footnote{Editor's note: this remark has been added; it is useful for
Corollary~3.11.1.}

Indeed, let \(m:M\times_YM\to M\) and
\(\lambda:\mathcal O_Y\times_YM\to M\) denote the morphisms defining
the structures of a \(Y\)-abelian group and an \(\mathcal O_Y\)-module.
Let \(H'\) be an \(S\)-scheme over
\(H=\underline{\operatorname{Hom}}_{Z/S}(X,Y)\); equivalently, one has
been given a \(Z\)-morphism \(f:X\times_SH'\to Y\), which makes
\(X\times_SH'\) a \(Y\)-object. Then
\[
  \operatorname{Hom}_H
    \bigl(H',\underline{\operatorname{Hom}}_{Z/S}(X,M)\bigr)
\]
is the set of \(Z\)-morphisms \(\phi:X\times_SH'\to M\) such that
\(\pi\circ\phi=f\), that is, the set of \(Y\)-morphisms
\(X\times_SH'\to M\).

% Source page: Exp2 component-local p. 4; combined-reader p. 62.
If \(\phi,\psi\) are two such morphisms, define \(\phi+\psi\) to be the
composite \(Y\)-morphism
\[
  X\times_SH'\xrightarrow{\ \phi\times\psi\ }
  M\times_YM\xrightarrow{\ m\ }M.
\]
This makes \(\underline{\operatorname{Hom}}_{Z/S}(X,M)\) an abelian group
over \(H=\underline{\operatorname{Hom}}_{Z/S}(X,Y)\).
\footnote{Editor's note: if \(M\) is a \(Y\)-group, not necessarily
abelian, and one sets \(\phi\cdot\psi=m\circ(\phi\times\psi)\), this makes
\(\underline{\operatorname{Hom}}_{Z/S}(X,M)\) a group over
\(\underline{\operatorname{Hom}}_{Z/S}(X,Y)\).}

Similarly, if \(a\in\mathcal O(X\times_SH')\), that is, if
\(a:X\times_SH'\to\mathcal O_S\) is an \(S\)-morphism, define
\(a\phi\) to be the composite \(\lambda\circ(a\times\phi)\), where
\(a\times\phi\) is the \(Y\)-morphism from \(X\times_SH'\) to
\(\mathcal O_Y\times_YM\simeq\mathcal O_S\times_SM\) with components
\(a\) and \(\phi\). This gives
\(\operatorname{Hom}_H(H',\underline{\operatorname{Hom}}_{Z/S}(X,M))\)
an \(\mathcal O(X\times_SH')\)-module structure, functorial in the
\(H\)-object \(H'\).

\section*{\SGAInlineRefTarget{sga3:ii:section:2}2. The schemes \(I_S(\mathcal M)\)}
\addcontentsline{toc}{section}{2. The schemes I S (M)}

\SGAInlineRefTarget{sga3:ii:definition:2-1}\subsection*{\SGARefLink{sga3:ii:definition:2-1}{Definition 2.1}.}
Let \(S\) be a scheme and \(\mathcal M\) a quasi-coherent
\(\mathcal O_S\)-module. Denote by
\(\mathscr D_{\mathcal O_S}(\mathcal M)\) the quasi-coherent
\(\mathcal O_S\)-algebra \(\mathcal O_S\oplus\mathcal M\), where
\(\mathcal M\) is regarded as an ideal of square zero. Denote by
\[
  I_S(\mathcal M)=\operatorname{Spec}
    \mathscr D_{\mathcal O_S}(\mathcal M)
\]
the corresponding \(S\)-scheme.\footnote{Editor's note: the underlying
space of \(I_S(\mathcal M)\) is the same as that of \(S\).}
In particular, put
\[
  \mathscr D_{\mathcal O_S}
  =\mathscr D_{\mathcal O_S}(\mathcal O_S),
  \qquad I_S=I_S(\mathcal O_S),
\]
and call these, respectively, the algebra of dual numbers over \(S\) and
the scheme of dual numbers over \(S\).

The assignment \(\mathcal M\mapsto I_S(\mathcal M)\) is a contravariant
functor from quasi-coherent \(\mathcal O_S\)-modules to \(S\)-schemes.
In particular, the morphisms \(0\to\mathcal M\) and
\(\mathcal M\to0\) define, respectively, the structural morphism
\[
  \rho:I_S(\mathcal M)\longrightarrow I_S(0)=S
\]
and a section \(\varepsilon_{\mathcal M}\) of it, called the
\emph{zero section}.\footnote{Editor's note: compare \SGARefLink{sga3:ii:remark:3-1-1}{3.1.1}.}

\subsection*{\SGAInlineRefTarget{sga3:ii:subsection:2-1-1}2.1.1.}
Since \(\mathcal M\mapsto I_S(\mathcal M)\) is contravariant, every
\(a\in\operatorname{End}_{\mathcal O_S}(\mathcal M)\) defines an
\(S\)-endomorphism \(a^*\) of \(I_S(\mathcal M)\), and
\[
  1^*=\operatorname{id},\qquad (ab)^*=b^*\circ a^*,\qquad
  0^*=\varepsilon_{\mathcal M}\circ\rho,\qquad
  a^*\circ\varepsilon_{\mathcal M}=\varepsilon_{\mathcal M}.
\]
Consequently, \(I_S(\mathcal M)\) carries a right action of the
multiplicative monoid of
\(\operatorname{End}_{\mathcal O_S}(\mathcal M)\), commuting with the
\(S\)-morphisms \(I_S(\mathcal M)\to I_S(\mathcal M')\) arising from
morphisms \(\mathcal M'\to\mathcal M\); in particular, the operators
\(a^*\) preserve the zero section.\footnote{Editor's note: what follows
expands the original; the numbering \SGARefLink{sga3:ii:subsection:2-1-1}{2.1.1}--\SGARefLink{sga3:ii:subsection:2-1-3}{2.1.3} has also been added.}

For every \(a\in\operatorname{End}_{\mathcal O_S}(\mathcal M)\), every
\(f:S'\to S\), and every \(m\in I_S(\mathcal M)(S')\), write
\(m\cdot a=a^*(m)\). Then
\[
  m\cdot1=m,\qquad (m\cdot a)\cdot b=m\cdot(ab),\qquad
  m\cdot0=\varepsilon_{\mathcal M}(\rho(m)),
\]
and if \(m=\varepsilon_{\mathcal M}(f)\), then \(m\cdot a=m\).
\footnote{Editor's note: later we shall be chiefly interested in
\(a\in\mathcal O(S)\), acting by homotheties on \(\mathcal M\). For
example, if \(\mathcal M\) is free of rank \(r\), then for every
\(S'\to S\), \(\operatorname{Hom}_S(S',I_S(\mathcal M))\) identifies
with the set of \(r\)-tuples
\((\epsilon_1,\ldots,\epsilon_r)\in\mathcal O(S')^r\) satisfying
\(\epsilon_i\epsilon_j=0\) for all \(i,j\), and
\((\epsilon_1,\ldots,\epsilon_r)\cdot a
=(a\epsilon_1,\ldots,a\epsilon_r)\) for every \(a\in\mathcal O(S)\).}

% Source page: Exp2 component-local p. 5; combined-reader p. 63.
\SGAInlineRefTarget{sga3:ii:remark:2-1-2}\subsection*{\SGARefLink{sga3:ii:remark:2-1-2}{Remark 2.1.2}.}
Formation of \(I_S(\mathcal M)\) commutes with base extension: there are
canonical isomorphisms
\[
  I_S(\mathcal M)_{S'}
  \simeq I_{S'}(\mathcal M\otimes_{\mathcal O_S}\mathcal O_{S'}).
\]
For simplicity write \(I_{S'}(\mathcal M)=I_S(\mathcal M)_{S'}\). More
generally, if \(X\) is an \(S\)-functor, not necessarily representable,
write \(I_X(\mathcal M)=I_S(\mathcal M)\times_SX\).

\subsection*{\SGAInlineRefTarget{sga3:ii:subsection:2-1-3}2.1.3.}
By the preceding discussion, the multiplicative monoid of
\(\mathcal O(S')\) acts on the \(S'\)-scheme
\(I_{S'}(\mathcal M)\), functorially in \(\mathcal M\). In other words,
the \(S\)-scheme \(I_S(\mathcal M)\) is an object with an operator-monoid
structure \(\mathcal O_S\), functorial in \(\mathcal M\). Thus there is
a morphism of \(S\)-schemes
\[
  \lambda:I_S(\mathcal M)\times_S\mathcal O_S
  \longrightarrow I_S(\mathcal M),
\]
satisfying the evident conditions. For every \(S\)-functor \(X\), base
change gives a morphism of \(X\)-functors
\[
  \lambda_X:I_X(\mathcal M)\times_S\mathcal O_S
  \longrightarrow I_X(\mathcal M),
\]
making \(I_X(\mathcal M)\) an object with operator monoid
\(\mathcal O(X)\). Every
\(a\in\mathcal O(X)=\operatorname{Hom}_S(X,\mathcal O_S)\) defines the
\(X\)-endomorphism
\[
  a^*=\lambda_X\circ
  \bigl(\operatorname{id}_{I_X(\mathcal M)}\times(a\circ\operatorname{pr}_X)\bigr).
\]
Explicitly, if \(x\in X(S')\) and
\(m\in I_S(\mathcal M)(S')=I_{S'}(\mathcal M)(S')\), then
\(a(x)=a\circ x\in\mathcal O(S')\) and
\[
  (m,x)\cdot a=(m\cdot a(x),x).
\]
This operation is functorial in \(\mathcal M\) and preserves the zero
section \(\varepsilon_{\mathcal M}:X\to I_X(\mathcal M)\); thus
\(a^*\circ\varepsilon_{\mathcal M}=\varepsilon_{\mathcal M}\).
\footnote{Editor's note: this paragraph has been added; it will be useful
in \SGARefLink{sga3:ii:remark:3-4-2}{3.4.2} and \SGARefLink{sga3:ii:lemma:4-6-2}{4.6.2}.}

The operation is also ``functorial in \(X\)'' in the following sense. If
\(\pi:Y\to X\) is a morphism of \(S\)-functors and
\(u:\mathcal O(X)\to\mathcal O(Y)\) the corresponding ring morphism,
so that \(u(a)=a\circ\pi\), then the following diagram is commutative.

\SGAThreeNativeDiagram[0.30\textwidth]
  {native_diagrams/exp2/exp2_p005_operator_functoriality_square.tex}
  {Functoriality of the operator action under \(Y\to X\).}
  {fig:exp2-operator-functoriality}
  {Polo--Gille Exp2-14oct24.pdf, component-local p. 5; combined-reader p. 63}

\subsection*{\SGAInlineRefTarget{sga3:ii:subsection:2-2}2.2.}
Let \(\mathcal M\) and \(\mathcal N\) be two quasi-coherent
\(\mathcal O_S\)-modules. The commutative diagram below induces the
corresponding diagram of \(S\)-schemes on the following page.

\SGAThreeNativeDiagram[0.26\textwidth]
  {native_diagrams/exp2/exp2_p005_module_sum_diamond.tex}
  {The direct-sum diagram of \(\mathcal O_S\)-modules.}
  {fig:exp2-module-sum-diamond}
  {Polo--Gille Exp2-14oct24.pdf, component-local p. 5; combined-reader p. 63}

% Source page: Exp2 component-local p. 6; combined-reader p. 64.
\SGAThreeNativeDiagram[0.42\textwidth]
  {native_diagrams/exp2/exp2_p006_IS_sum_scheme_diagram.tex}
  {The induced commutative diagram of \(S\)-schemes.}
  {fig:exp2-is-sum-schemes}
  {Polo--Gille Exp2-14oct24.pdf, component-local p. 6; combined-reader p. 64}

\subsection*{Proposition 2.2.}
For every \(S\)-scheme \(X\), the diagram of functors over \(S\) obtained
by applying \(\underline{\operatorname{Hom}}_S(-,X)\) to the preceding
diagram is cartesian.

\SGAThreeNativeDiagram[0.58\textwidth]
  {native_diagrams/exp2/exp2_p006_Hom_cartesian_diagram.tex}
  {The cartesian diagram of \(\underline{\operatorname{Hom}}_S\)-functors.}
  {fig:exp2-hom-cartesian}
  {Polo--Gille Exp2-14oct24.pdf, component-local p. 6; combined-reader p. 64}

Indeed, for every \(S'\to S\), one must verify that the diagram of sets
obtained by evaluating the functors on \(S'\) is cartesian. Since formation
of \(I_S(\mathcal P)\) commutes with base extension, it is enough to take
\(S'=S\), and hence to verify that the following diagram of sets is
cartesian.

\SGAThreeNativeDiagram[0.58\textwidth]
  {native_diagrams/exp2/exp2_p006_points_cartesian_diagram.tex}
  {The cartesian diagram after evaluation on \(S\).}
  {fig:exp2-points-cartesian}
  {Polo--Gille Exp2-14oct24.pdf, component-local p. 6; combined-reader p. 64}

If \(x\in X(S)\), it follows from SGA~1, III~5.1 that
\(X(\varepsilon_{\mathcal M})^{-1}(x)\) is isomorphic, functorially in
\(\mathcal M\), to
\[
  \operatorname{Hom}_{\mathcal O_S}
    \bigl(x^*(\Omega^1_{X/S}),\mathcal M\bigr),
\]
where \(\Omega^1_{X/S}\) denotes the sheaf of relative differentials of
\(X\) over \(S\). This latter functor in \(\mathcal M\) plainly carries a
direct sum of \(\mathcal O_S\)-modules to the product of the corresponding
sets, which proves the proposition.\footnote{Editor's note: see also the
addition \SGARefLink{sga3:iii:theorem:0-1-8}{0.1.8} in Expos\'e~III.}

% Source page: Exp2 component-local p. 7; combined-reader p. 65.
\SGAInlineRefTarget{sga3:ii:corollary:2-2-1}\subsection*{\SGARefLink{sga3:ii:corollary:2-2-1}{Corollary 2.2.1}.}
Let \(X\) be an \(S\)-scheme and \(\mathcal M\) a finite free
\(\mathcal O_S\)-module. The functor
\(\underline{\operatorname{Hom}}_S(I_S(\mathcal M),X)\), as a functor
over \(X\), is isomorphic to a finite product over \(X\) of copies of
\(\underline{\operatorname{Hom}}_S(I_S,X)\).

\SGAInlineRefTarget{sga3:ii:nota:2-2-2}\subsection*{Nota 2.2.2.}
It follows from the proof of the proposition that
\(\underline{\operatorname{Hom}}_S(I_S,X)\), as an \(X\)-functor, is
isomorphic to \(\mathbf V(\Omega^1_{X/S})\) (\SGARefLink{sga3:i:definition:4-6-1}{I~4.6.1}), and is therefore
represented by the vector fibration\footnote{Editor's note: cf. editor's
note (33) of Expos\'e~I.}
\[
  V(\Omega^1_{X/S}).
\]

% END INCLUDED SGA3_English_current_buildable_source_20260823/source/inputs/II/tex/components/10_expose_II_s1_s2_hom_dual_numbers.tex

% BEGIN INCLUDED SGA3_English_current_buildable_source_20260823/source/inputs/II/tex/components/11_expose_II_s3_tangent_bundle_condition_E.tex
% SGA 3, Expose II, Loop 1 English reconstruction.
% Scope: Section 3, immediately after Nota 2.2.2 through 3.14, stopping
% immediately before Section 4.
% Authority: Polo--Gille Exp2-14oct24.pdf, component-local pages 7--21,
% with the Section 4 boundary verified on p. 22 (combined-reader pp. 65--80).
% Force-OCR is a drafting/locator witness only.
% English comparison witness: Jacob C. Reinhold, jcreinhold/sga commit
% e7a259f3f8608ad3edf9bf6eead3fd504dd2d23e (CC BY 4.0); consulted for
% comparison and drafting, not treated as source authority.

% Source page: Exp2 component-local p. 7; combined-reader p. 65.
\section*{\SGAInlineRefTarget{sga3:ii:section:3}3. The tangent bundle, condition (E)}
\addcontentsline{toc}{section}{3. The tangent bundle, condition (E)}

Throughout this section, unless otherwise stated, the letters
\(\mathcal M,\mathcal M',\mathcal N,\ldots\) will denote finite free
\(\mathcal O_S\)-modules (that is, modules isomorphic to finite direct sums
of copies of \(\mathcal O_S\)).

We shall systematically use the identifications established in Expos\'e~I.
Thus, by a ``functor over \(S\)'' we shall mean indifferently either a
functor equipped with a morphism to \(S=h_S\), or a functor on the category
of objects over \(S\). We shall use ``group functor over \(S\),'' and
similar expressions, in the same way.

\SGAInlineRefTarget{sga3:ii:definition:3-1}\subsection*{\SGARefLink{sga3:ii:definition:3-1}{Definition 3.1}.}
Let \(S\) be a scheme, let \(\mathcal M\) be a finite free
\(\mathcal O_S\)-module, and let \(X\) be a functor over \(S\). The
\emph{tangent bundle of \(X\) over \(S\), relative to the
\(\mathcal O_S\)-module \(\mathcal M\)}, is the \(S\)-functor
\[
  T_{X/S}(\mathcal M)
  =\underline{\operatorname{Hom}}_S(I_S(\mathcal M),X).
\]
In particular, the \emph{tangent bundle of \(X\) over \(S\)} is
\footnote{Editor's note: when \(S=\operatorname{Spec}(k)\) and \(X\) is a
\(k\)-scheme, one has
\(T_{X/S}(S')=\operatorname{Hom}_S(S'\otimes_k k[\epsilon],X)
=X(S'\otimes_k k[\epsilon])\); this recovers one of the usual definitions
of the tangent bundle.}
\[
  T_{X/S}=T_{X/S}(\mathcal O_S)
  =\underline{\operatorname{Hom}}_S(I_S,X).
\]

The assignment \(\mathcal M\mapsto T_{X/S}(\mathcal M)\) is a covariant
functor from finite free \(\mathcal O_S\)-modules to \(S\)-functors. In
particular, the morphisms \(\mathcal M\to0\) and \(0\to\mathcal M\)
define, respectively, an \(S\)-morphism
\[
  \pi_{\mathcal M}:T_{X/S}(\mathcal M)\longrightarrow
  T_{X/S}(0)\simeq X
\]
and a section \(\tau_0\) of it, called the \emph{zero section}.
\footnote{Editor's note: the original order ``\(0\to\mathcal M\) and
\(\mathcal M\to0\)'' has been corrected to ``\(\mathcal M\to0\) and
\(0\to\mathcal M\).''}

\SGAInlineRefTarget{sga3:ii:remark:3-1-1}\subsection*{\SGARefLink{sga3:ii:remark:3-1-1}{Remark 3.1.1}.}
\footnote{Editor's note: paragraphs 3.1.1 and \SGARefLink{sga3:ii:scholium:3-1-2}{3.1.2} have been added.}
The projection
\(\pi_{\mathcal M}:T_{X/S}(\mathcal M)\to X\) is induced by the zero
section \(\varepsilon_{\mathcal M}:S\to I_S(\mathcal M)\), whereas the
zero section \(\tau_0:X\to T_{X/S}(\mathcal M)\) is induced by the
structural morphism \(\rho:I_S(\mathcal M)\to S\). More explicitly, let
\(t\in T_{X/S}(\mathcal M)(S')\) correspond to the \(S\)-morphism
\(f:S'\times_S I_S(\mathcal M)\to X\), and let \(x\in X(S')\)
correspond to \(g:S'\to X\). Then
\[
  \pi(t)=f\circ(\operatorname{id}_{S'}\times\varepsilon_{\mathcal M}),
  \qquad
  \tau_0(x)=g\circ(\operatorname{id}_{S'}\times\rho).
\]

It follows from \SGARefLink{sga3:ii:definition:3-1}{3.1} that \(\mathcal M\mapsto T_{X/S}(\mathcal M)\) is a
covariant functor from finite free \(\mathcal O_S\)-modules to functors
over \(X\). In particular, \(\mathcal O(S)\) is an operator monoid of the
\(X\)-functor \(T_{X/S}(\mathcal M)\), compatible with functoriality in
\(\mathcal M\).

% Source page: Exp2 component-local p. 8; combined-reader p. 66.
\SGAInlineRefTarget{sga3:ii:scholium:3-1-2}\subsection*{Scholium 3.1.2.}
\footnotemark[15]
The preceding discussion means, in particular, the following. For every
\(S\)-morphism \(\varphi:X'\to X\), put
\[
  \Sigma(X',\mathcal M)
  =\operatorname{Hom}_X\bigl(X',T_{X/S}(\mathcal M)\bigr).
\]
The multiplicative monoid
\(\operatorname{End}_{\mathcal O_S}(\mathcal M)\) acts on
\(\Sigma(X',\mathcal M)\); write the action as
\((\lambda,x)\mapsto\lambda*x\). Then
\[
  \lambda*(\mu*x)=(\lambda\mu)*x,\qquad
  1*x=x,\qquad 0*x=\tau_0\circ\varphi.
\]
There is likewise an action of
\(\operatorname{End}_{\mathcal O_S}(\mathcal M\oplus\mathcal M)\) on
\(\Sigma(X',\mathcal M\oplus\mathcal M)\).
\footnote{Editor's note: in fact, only the actions of \(\mathcal O(S)\)
(respectively \(\mathcal O(S)\times\mathcal O(S)\)) on
\(\Sigma(X',\mathcal M)\) (respectively
\(\Sigma(X',\mathcal M\oplus\mathcal M)\)) will be used; see below and
the proof of \SGARefLink{sga3:ii:proposition:3-6}{3.6}.}

Let \(m:\mathcal M\oplus\mathcal M\to\mathcal M\) be addition, and let
\(\delta:\mathcal M\to\mathcal M\oplus\mathcal M\) be the diagonal.
Write
\[
  m_{X'}:\Sigma(X',\mathcal M\oplus\mathcal M)	o
    \Sigma(X',\mathcal M),
  \qquad
  \delta_{X'}:\Sigma(X',\mathcal M)	o
    \Sigma(X',\mathcal M\oplus\mathcal M)
\]
for the induced maps. For \(\lambda,\mu\in\mathcal O(S)\), let
\(\ell_\lambda\) denote multiplication by \(\lambda\) on \(\mathcal M\),
and let \(\ell_{\lambda,\mu}\) denote multiplication by
\((\lambda,\mu)\) on \(\mathcal M\oplus\mathcal M\). Since
\[
  m\circ\ell_{\lambda,\lambda}=\ell_\lambda\circ m,
  \qquad
  m\circ\ell_{\lambda,\mu}\circ\delta=\ell_{\lambda+\mu},
\]
we have, for every \(z\in\Sigma(X',\mathcal M\oplus\mathcal M)\) and
\(x\in\Sigma(X',\mathcal M)\),
\begin{equation}
  \lambda*m(z)=m((\lambda,\lambda)*z),
  \qquad
  m((\lambda,\mu)*\delta(x))=(\lambda+\mu)*x.
  \tag{$\dagger$}\label{sga3:ii:section:3:equation:dagger}
\end{equation}

\SGAInlineRefTarget{sga3:ii:definition:3-2}\subsection*{\SGARefLink{sga3:ii:definition:3-2}{Definition 3.2}.}
Let \(u\in X(S)=\operatorname{Hom}_S(S,X)=\Gamma(X/S)\). The
\emph{tangent space of \(X\) over \(S\) at \(u\), relative to
\(\mathcal M\)}, denoted \(L^u_{X/S}(\mathcal M)\), is the \(S\)-functor
obtained from the \(X\)-functor \(T_{X/S}(\mathcal M)\) by pullback along
\(u:S\to X\).

\SGAThreeNativeDiagram[0.36\textwidth]
  {native_diagrams/exp2/exp2_p008_tangent_space_cartesian_square.tex}
  {The cartesian square defining \(L^u_{X/S}(\mathcal M)\).}
  {fig:exp2-tangent-space-cartesian}
  {Polo--Gille Exp2-14oct24.pdf, component-local p. 8; combined-reader p. 66}

In particular, \(L^u_{X/S}(\mathcal O_S)\) is denoted \(L^u_{X/S}\) and
is called the \emph{tangent space of \(X\) over \(S\) at \(u\)}.

\SGAInlineRefTarget{sga3:ii:remark:3-2-1}\subsection*{\SGARefLink{sga3:ii:remark:3-2-1}{Remark 3.2.1}.}
\footnote{Editor's note: this remark has been added.}
It follows from \SGARefLink{sga3:ii:remark:3-1-1}{3.1.1} that, for every \(t:S'\to S\),
\(L^u_{X/S}(\mathcal M)(S')\) is the set of \(S\)-morphisms
\(f:I_{S'}(\mathcal M)\to X\) such that
\(f\circ\varepsilon_{\mathcal M}=u\circ t\), where
\(\varepsilon_{\mathcal M}:S'\to I_{S'}(\mathcal M)\) is the zero
section.

\SGAInlineRefTarget{sga3:ii:proposition:3-3}\subsection*{\SGARefLink{sga3:ii:proposition:3-3}{Proposition 3.3}.}
If \(X\) is represented by an \(S\)-scheme \(\widetilde X\), then
\(T_{X/S}(\mathcal M)\) and \(L^u_{X/S}(\mathcal M)\) are representable.
In particular, \(T_{X/S}\) and \(L^u_{X/S}\) are represented by the
vector fibrations
\[
  V(\Omega^1_{\widetilde X/S})
  \qquad\text{and}\qquad
  V(u^*\Omega^1_{\widetilde X/S})
\]
over \(\widetilde X\) and \(S\), respectively.

It plainly suffices to prove the assertion for \(T_{X/S}(\mathcal M)\),
since the corresponding assertion for \(L^u_{X/S}(\mathcal M)\) follows
by pullback. By \SGARefLink{sga3:ii:corollary:2-2-1}{Corollary~2.2.1}, it even suffices to treat \(T_{X/S}\),
in which case the result is exactly Nota~\SGARefLink{sga3:ii:nota:2-2-2}{2.2.2}.

% Source page: Exp2 component-local p. 9; combined-reader p. 67.
\SGAInlineRefTarget{sga3:ii:remark:3-3-1}\subsection*{\SGARefLink{sga3:ii:remark:3-3-1}{Remark 3.3.1}.}
The proposition gives a particularly simple description of the vector
fibration representing \(L^u_{X/S}\). The image of the section \(u\) of
\(X\) over \(S\) is locally closed,
\footnote{Editor's note: cf. EGA~I, 5.3.11.}
and hence is defined by a quasi-coherent ideal \(\mathfrak m\) in the
scheme induced on an open subscheme of \(X\). The quotient
\(\mathfrak m/\mathfrak m^2\) may be regarded as a quasi-coherent module
on \(S\); it is this module that defines the vector fibration in question.

For example, let \(X\) be an algebraic scheme over a field \(k\), and let
\(u\) be a \(k\)-rational point of \(X\). Let \(\mathfrak m\) be the
maximal ideal of the local ring \(\mathcal O_{X,u}\), and let
\(\mathfrak t\) be the \(k\)-vector space dual to
\(\mathfrak m/\mathfrak m^2\); this is the Zariski tangent space of
\(\mathcal O_{X,u}\) at \(u\). In the notation of \SGARefLink{sga3:i:corollary:4-6-5-1}{I~4.6.5.1},
\[
  L^u_{X/k}=V(\mathfrak m/\mathfrak m^2)=W(\mathfrak t).
\]

Returning to the general situation, observe first that
\(L^u_{X/S}(\mathcal M)\) is covariant from finite free
\(\mathcal O_S\)-modules to functors over \(S\). In particular,
\(\mathcal O(S)\) is a set of operators of the \(S\)-functor
\(L^u_{X/S}(\mathcal M)\), compatible with functoriality in
\(\mathcal M\).
\footnote{Editor's note: cf. \SGARefLink{sga3:ii:scholium:3-1-2}{3.1.2}.}

\SGAInlineRefTarget{sga3:ii:proposition:3-4}\subsection*{\SGARefLink{sga3:ii:proposition:3-4}{Proposition 3.4}.}
Formation of \(T_{X/S}(\mathcal M)\) and \(L^u_{X/S}(\mathcal M)\)
commutes with base extension. For every \(S\)-scheme \(S'\), there are
isomorphisms, functorial in \(\mathcal M\),
\[
  T_{X_{S'}/S'}(\mathcal M\otimes_{\mathcal O_S}\mathcal O_{S'})
    \xrightarrow{\sim} T_{X/S}(\mathcal M)_{S'},
\]
\[
  L^{u'}_{X_{S'}/S'}(\mathcal M\otimes_{\mathcal O_S}\mathcal O_{S'})
    \xrightarrow{\sim} L^u_{X/S}(\mathcal M)_{S'},
  \qquad u'=u_{S'}.
\]
This follows immediately from the compatibility of
\(\underline{\operatorname{Hom}}\) with base extension.

\SGAInlineRefTarget{sga3:ii:corollary:3-4-1}\subsection*{\SGARefLink{sga3:ii:corollary:3-4-1}{Corollary 3.4.1}.}
The \(X\)-functor \(T_{X/S}(\mathcal M)\), respectively the \(S\)-functor
\(L^u_{X/S}(\mathcal M)\), is naturally an object with operators
\(\mathcal O_X\), respectively \(\mathcal O_S\), and these structures
are functorial in \(\mathcal M\).

For \(L^u_{X/S}(\mathcal M)\), this is immediate. For each \(S'\) over
\(S\), the ring \(\mathcal O(S')\) acts on
\(\mathcal M\otimes\mathcal O_{S'}\), hence on
\[
  L^{u'}_{X_{S'}/S'}(\mathcal M\otimes\mathcal O_{S'})
  =L^u_{X/S}(\mathcal M)_{S'};
\]
this action is functorial in \(S'\).

The case of \(T_{X/S}(\mathcal M)\) requires a little more detail. For an
\(X\)-object \(f:X'\to X\), put
\(T_{X/S}(\mathcal M)_{X'}=T_{X/S}(\mathcal M)\times_X X'\). We must
equip
\[
  T_{X/S}(\mathcal M)_{X'}(X')
  =\operatorname{Hom}_X\bigl(X',T_{X/S}(\mathcal M)\bigr)
\]
with an action of \(\mathcal O(X')\), functorially in \(X'\). Let
\(X_{X'}=X\times_S X'\), and let \(f'\) be the section of \(X_{X'}\)
over \(X'\) defined by \(f\). The following diagram provides the required
identification.

% Source page: Exp2 component-local p. 10; combined-reader p. 68.
\SGAThreeNativeDiagram[0.39\textwidth]
  {native_diagrams/exp2/exp2_p010_base_change_operator_diagram.tex}
  {The base-change diagram used to define the \(\mathcal O_X\)-operator structure.}
  {fig:exp2-base-change-operator}
  {Polo--Gille Exp2-14oct24.pdf, component-local p. 10; combined-reader p. 68}

Together with \SGARefLink{sga3:ii:remark:3-2-1}{3.2.1}, the diagram identifies
\(T_{X/S}(\mathcal M)_{X'}(X')\) with
\[
  L^{f'}_{X_{X'}/X'}(\mathcal M)(X')
  =\left\{
    \begin{array}{c|c}
      X'\text{-morphisms }\psi:I_{X'}(\mathcal M)\to X_{X'}
      &\psi\circ\varepsilon_{\mathcal M}=f'
    \end{array}
  \right\}.
\]
Every \(a\in\mathcal O(X')\) acts through its action on
\(I_{X'}(\mathcal M)\): in the notation of \SGARefLink{sga3:ii:subsection:2-1-1}{2.1.1},
\[
  a\psi=\psi\circ a^*,
  \qquad (a\psi)(x)=\psi(x\cdot a)
\]
for every \(X''\to X'\) and every \(x\in I_{X'}(\mathcal M)(X'')\).
This construction is plainly functorial in \(X'\).
\footnote{Editor's note: the original has been expanded in what follows.}
The isomorphisms of \SGARefLink{sga3:ii:proposition:3-4}{Proposition~3.4} are therefore, by construction,
isomorphisms for the corresponding \(\mathcal O_{X_{S'}}\)- and
\(\mathcal O_{S'}\)-operator structures.

\SGAInlineRefTarget{sga3:ii:remark:3-4-2}\subsection*{\SGARefLink{sga3:ii:remark:3-4-2}{Remark 3.4.2}.}
\footnote{Editor's note: Remarks 3.4.2 and \SGARefLink{sga3:ii:remark:3-4-3}{3.4.3} have been added.}
The action of \(\mathcal O_X\) on \(T_{X/S}(\mathcal M)\) can be seen more
directly as follows. For every \(f:X'\to X\),
\[
  \operatorname{Hom}_X\bigl(X',T_{X/S}(\mathcal M)\bigr)
  =\left\{\varphi\in
    \operatorname{Hom}_S(I_{X'}(\mathcal M),X)
    \ \middle|\ \varphi\circ\varepsilon_{\mathcal M}=f\right\}.
\]
By \SGARefLink{sga3:ii:subsection:2-1-3}{2.1.3}, the \(S\)-functor \(I_{X'}(\mathcal M)\) carries an action of
the monoid \(\mathcal O(X')\) preserving the zero section. Thus, if
\(a^*\) denotes the \(X'\)-endomorphism defined by
\(a\in\mathcal O(X')\), then
\[
  a\varphi=\varphi\circ a^*,
  \qquad
  (a\varphi)(m,x')=\varphi(m\cdot a(x'),x')
\]
for every \(S'\to S\) and every
\((m,x')\in\operatorname{Hom}_S
  (S',I_S(\mathcal M)\times_S X')\). Notice that
\(a^*\circ\varepsilon_{\mathcal M}=\varepsilon_{\mathcal M}\), so that
\((a\varphi)\circ\varepsilon_{\mathcal M}=f\).

Likewise, for \(t:S'\to S\), the set
\(L^u_{X/S}(\mathcal M)(S')\) consists of the \(S\)-morphisms
\(\varphi:I_{S'}(\mathcal M)\to X\) satisfying
\(\varphi\circ\varepsilon_{\mathcal M}=u\circ t\), and
\(a\varphi=\varphi\circ a^*\) for \(a\in\mathcal O(S')\).

\SGAInlineRefTarget{sga3:ii:remark:3-4-3}\subsection*{Remark 3.4.3.}
\footnotemark[21]
The observations of \SGARefLink{sga3:ii:scholium:3-1-2}{3.1.2} apply equally to the action of \(\mathcal O_S\)
on \(L^u_{X/S}(\mathcal M)\) and to that of \(\mathcal O_X\) on
\(T_{X/S}(\mathcal M)\).

% Source page: Exp2 component-local p. 11; combined-reader p. 69.
\SGAInlineRefTarget{sga3:ii:definition:3-5}\subsection*{\SGARefLink{sga3:ii:definition:3-5}{Definition 3.5}.}
Let \(S\) be a scheme and \(X\) an \(S\)-functor. We say that \(X\)
\emph{satisfies condition (E) relative to \(S\)} if, for every \(S'\)
over \(S\) and every pair of finite free \(\mathcal O_{S'}\)-modules
\(\mathcal M,\mathcal N\), the diagram of sets obtained by applying \(X\)
to diagram \((*)\) of \SGARefLink{sga3:ii:definition:2-1}{2.1} is cartesian.
\footnote{Editor's note: for examples of \(S\)-functors and \(S\)-groups
that do not satisfy (E), see \SGARefLink{sga3:ii:subsection:6-2}{6.2} and \SGARefLink{sga3:ii:subsection:6-3}{6.3}.}

\SGAThreeNativeDiagram[0.46\textwidth]
  {native_diagrams/exp2/exp2_p011_condition_E_cartesian_diamond.tex}
  {The cartesian diagram in condition (E).}
  {fig:exp2-condition-e}
  {Polo--Gille Exp2-14oct24.pdf, component-local p. 11; combined-reader p. 69}

\subsection*{\SGAInlineRefTarget{sga3:ii:subsection:3-5-1}3.5.1.}
\footnote{Editor's note: what follows has been added.}
Equivalently, the functor
\(\mathcal M\mapsto T_{X/S}(\mathcal M)\) carries direct sums of finite
free \(\mathcal O_S\)-modules to products of \(X\)-functors. In that case,
the same is true of
\[
  \mathcal M\longmapsto L^u_{X/S}(\mathcal M)
  =S\times_XT_{X/S}(\mathcal M)
\]
for every \(u\in\Gamma(X/S)\).

\SGARefLink{sga3:ii:subsection:2-2}{Proposition~2.2} shows that every representable functor satisfies
condition (E). We shall sometimes abbreviate ``\(X\) satisfies condition
(E) relative to \(S\)'' to ``\(X/S\) satisfies condition (E).''

If \(X/S\) satisfies (E), then
\(\mathcal M\mapsto T_{X/S}(\mathcal M)\) commutes with products and hence
carries groups to groups. In particular, \(T_{X/S}(\mathcal M)\) is a
commutative \(X\)-group; for the same reason,
\(L^u_{X/S}(\mathcal M)\) is a commutative \(S\)-group.

\SGAInlineRefTarget{sga3:ii:proposition:3-6}\subsection*{\SGARefLink{sga3:ii:proposition:3-6}{Proposition 3.6}.}
If \(X/S\) satisfies (E), then the abelian-group structure on
\(T_{X/S}(\mathcal M)\), respectively on \(L^u_{X/S}(\mathcal M)\),
together with the action of \(\mathcal O_X\), respectively of
\(\mathcal O_S\), makes it an \(\mathcal O_X\)-module, respectively an
\(\mathcal O_S\)-module.

The operator action is functorial in \(\mathcal M\), and therefore
respects the abelian-group structure obtained by functoriality from that
of \(\mathcal M\).
\footnote{Editor's note: what follows has been expanded.}
Indeed, retaining the notation of \SGARefLink{sga3:ii:scholium:3-1-2}{3.1.2}, the abelian \(X\)-group law on
\(T_{X/S}(\mathcal M)\) is the composite
\[
  T_{X/S}(\mathcal M)\times_XT_{X/S}(\mathcal M)
  \simeq T_{X/S}(\mathcal M\oplus\mathcal M)
  \xrightarrow{\ m\ }T_{X/S}(\mathcal M),
\]
whereas
\[
  T_{X/S}(\mathcal M)\xrightarrow{\ \delta\ }
  T_{X/S}(\mathcal M\oplus\mathcal M)
  \simeq T_{X/S}(\mathcal M)\times_XT_{X/S}(\mathcal M)
\]
is the diagonal morphism.

% Source page: Exp2 component-local p. 12; combined-reader p. 70.
By \SGARefLink{sga3:ii:remark:3-4-3}{Remark~3.4.3} and the identities \((\dagger)\) of \SGARefLink{sga3:ii:scholium:3-1-2}{3.1.2},
\[
  \lambda(x+y)=\lambda x+\lambda y,
  \qquad
  (\lambda+\mu)x=\lambda x+\mu x
\]
for every \(f:X'\to X\), every
\(x,y\in\operatorname{Hom}_X(X',T_{X/S}(\mathcal M))\), and every
\(\lambda,\mu\in\mathcal O(X')\).

\SGAInlineRefTarget{sga3:ii:remark:3-6-1}\subsection*{\SGARefLink{sga3:ii:remark:3-6-1}{Remark 3.6.1}.}
If \(X\) is representable, then it satisfies (E), and \(T_{X/S}\) and
\(L^u_{X/S}\) are represented by vector fibrations. The preceding laws
are then exactly those induced by the vector-fibration structures
(cf. \SGARefLink{sga3:i:subsection:4-6}{I~4.6}).
\footnote{Editor's note: explicitly, for every \(S\)-morphism
\(f:X'\to X\), the action of \(\mathcal O(X')\) on
\(\operatorname{Hom}_X(X',T_{X/S}(\mathcal M))\) corresponds, under the
identification
\[
  \operatorname{Hom}_X(X',T_{X/S}(\mathcal M))
  \simeq
  \operatorname{Hom}_{\mathcal O_{X'}}
    (f^*\Omega^1_{X/S},\mathcal M\otimes_{\mathcal O_S}\mathcal O_{X'}),
\]
to the natural action on the right. This follows from the proof of SGA~1,
III~5.1; see also the addition 0.1.8 in Expos\'e~III.}

\SGAInlineRefTarget{sga3:ii:proposition:3-4-bis}\subsection*{\SGARefLink{sga3:ii:proposition:3-4-bis}{Proposition 3.4 bis}.}
If \(X/S\) satisfies (E), then \(X_{S'}/S'\) satisfies (E), and the
isomorphisms of 3.4 respect the \(\mathcal O_{X_{S'}}\)-module and
\(\mathcal O_{S'}\)-module structures, respectively.

No further comment is needed.

\SGAInlineRefTarget{sga3:ii:proposition:3-7}\subsection*{\SGARefLink{sga3:ii:proposition:3-7}{Proposition 3.7}.}
The functors \(T_{X/S}(\mathcal M)\) and \(L^u_{X/S}(\mathcal M)\) are
functorial in \(X\). Thus, if \(f:X\to X'\) is an \(S\)-morphism, the
following diagrams commute.

\SGAThreeNativeDiagram[0.74\textwidth]
  {native_diagrams/exp2/exp2_p012_functoriality_T_L_diagrams.tex}
  {Functoriality of tangent bundles and tangent spaces in \(X\).}
  {fig:exp2-tangent-functoriality}
  {Polo--Gille Exp2-14oct24.pdf, component-local p. 12; combined-reader p. 70}

Moreover, if \(f\) is a monomorphism, then so are \(T(f)\) and \(L(f)\).
\footnote{Editor's note: the preceding sentence has been added.}
The existence of \(T(f)\) and \(L(f)\), and the last assertion, follow
immediately from the definitions. Commutativity follows from functoriality
in \(\mathcal M\) and from \(T_{X/S}(0)=X\).

\SGAInlineRefTarget{sga3:ii:remark:3-7-1}\subsection*{\SGARefLink{sga3:ii:remark:3-7-1}{Remark 3.7.1}.}
\footnote{Editor's note: the hypothesis that \(X\) and \(X'\) be
representable has been added, and what follows has been expanded.}
Suppose that \(X\) and \(X'\) are representable, and let \(r\) be the rank
of \(\mathcal M\). By \SGARefLink{sga3:ii:nota:2-2-2}{2.2.2}, \(T_{X/S}(\mathcal M)\) is isomorphic to the
product over \(X\) of \(r\) copies of \(V(\Omega^1_{X/S})\), and similarly
for \(T_{X'/S}(\mathcal M)\). Hence the square in \SGARefLink{sga3:ii:proposition:3-7}{3.7} is cartesian when
\(f\) is an open immersion, or more generally when
\(f^*\Omega^1_{X'/S}=\Omega^1_{X/S}\), for example when \(f\) is
\'etale. Under these conditions there is an isomorphism of \(S\)-functors
\[
  L^u_{X/S}(\mathcal M)
  \xrightarrow{\sim}L^{f\circ u}_{X'/S}(\mathcal M).
\]

% Source page: Exp2 component-local p. 13; combined-reader p. 71.
In general, the square in \SGARefLink{sga3:ii:proposition:3-7}{3.7} defines the following morphism of
\(X\)-functors.

\SGAThreeNativeDiagram[0.42\textwidth]
  {native_diagrams/exp2/exp2_p013_general_X_functor_triangle.tex}
  {The induced morphism over \(X\) in the general case.}
  {fig:exp2-general-x-functor}
  {Polo--Gille Exp2-14oct24.pdf, component-local p. 13; combined-reader p. 71}

\SGAInlineRefTarget{sga3:ii:proposition:3-7-bis}\subsection*{\SGARefLink{sga3:ii:proposition:3-7-bis}{Proposition 3.7 bis}.}
Let \(f:X\to X'\) be an \(S\)-morphism. If \(X\) and \(X'\) satisfy
(E) relative to \(S\), then
\[
  T_{X/S}(\mathcal M)\xrightarrow{\ T(f)\ }
  T_{X'/S}(\mathcal M)_X,
  \qquad\text{respectively}\qquad
  L^u_{X/S}(\mathcal M)\xrightarrow{\ L(f)\ }
  L^{f\circ u}_{X'/S}(\mathcal M),
\]
is a morphism of \(\mathcal O_X\)-modules, respectively of
\(\mathcal O_S\)-modules.

This follows from \SGARefLink{sga3:ii:proposition:3-7}{Proposition~3.7} by functoriality in \(\mathcal M\).

\SGAInlineRefTarget{sga3:ii:proposition:3-8}\subsection*{\SGARefLink{sga3:ii:proposition:3-8}{Proposition 3.8}.}
For two functors \(X,Y\) over \(S\), there are isomorphisms, functorial in
\(\mathcal M\),
\begin{align}
  T_{X/S}(\mathcal M)\times_ST_{Y/S}(\mathcal M)
    &\xrightarrow{\sim}T_{(X\times_SY)/S}(\mathcal M), \tag{3.8.1}\label{sga3:ii:subsection:3-5-1:equation:3-8-1}\\
  L^u_{X/S}(\mathcal M)\times_SL^v_{Y/S}(\mathcal M)
    &\xrightarrow{\sim}L^{(u,v)}_{(X\times_SY)/S}(\mathcal M). \tag{3.8.2}\label{sga3:ii:subsection:3-5-1:equation:3-8-2}
\end{align}
The first isomorphism follows from \SGARefLink{sga3:ii:remark:1-2}{1.2} \((*)\);
\footnote{Editor's note: what follows has been expanded.}
the second follows from it
by base change along \((u,v):S\to X\times_SY\).

\SGAInlineRefTarget{sga3:ii:remark:3-8-0}\subsection*{\SGARefLink{sga3:ii:remark:3-8-0}{Remark 3.8.0}.}
The isomorphism (\SGARefLink{sga3:ii:subsection:3-5-1:equation:3-8-1}{3.8.1}) may also be read as an isomorphism of
\(X\times_SY\)-functors
\[
  \bigl(T_{X/S}(\mathcal M)\times_X(X\times_SY)\bigr)
  \mathop{\times}_{X\times_SY}
  \bigl(T_{Y/S}(\mathcal M)\times_Y(X\times_SY)\bigr)
  \xrightarrow{\sim}T_{(X\times_SY)/S}(\mathcal M).
\]

\subsection*{Corollary 3.8.1.}
If \(X/S\) carries an algebraic structure defined by finite cartesian
products, then \(T_{X/S}(\mathcal M)\) carries a structure of the same
kind, and the projection \(T_{X/S}(\mathcal M)\to X\) is a morphism for
that structure.

\SGAInlineRefTarget{sga3:ii:proposition:3-8-bis}\subsection*{\SGARefLink{sga3:ii:proposition:3-8-bis}{Proposition 3.8 bis}.}
If \(X/S\) and \(Y/S\) satisfy (E), then \((X\times_SY)/S\) satisfies
(E), and (\SGARefLink{sga3:ii:subsection:3-5-1:equation:3-8-1}{3.8.1}), respectively (\SGARefLink{sga3:ii:subsection:3-5-1:equation:3-8-2}{3.8.2}), is an isomorphism of
\(\mathcal O_{X\times_SY}\)-modules, respectively of
\(\mathcal O_S\)-modules.

\emph{Proof.}
\footnote{Editor's note: this proof has been added.}
By \SGARefLink{sga3:ii:subsection:3-5-1}{3.5.1} and (\SGARefLink{sga3:ii:subsection:3-5-1:equation:3-8-1}{3.8.1}), the product \((X\times_SY)/S\) satisfies (E). Let
us prove the assertion about (\SGARefLink{sga3:ii:subsection:3-5-1:equation:3-8-1}{3.8.1}). Given an \(S\)-morphism
\((x,y):Z\to X\times_SY\), \SGARefLink{sga3:ii:remark:3-4-2}{Remark~3.4.2} reduces the claim to the
\(\mathcal O(Z)\)-linearity of the map
\begin{multline*}
  \{\varphi\in\operatorname{Hom}_S(I_Z(\mathcal M),X)
       \mid\varphi\circ\varepsilon_{\mathcal M}=x\}
  \times
  \{\psi\in\operatorname{Hom}_S(I_Z(\mathcal M),Y)
       \mid\psi\circ\varepsilon_{\mathcal M}=y\}\\
  \longrightarrow
  \{\theta\in\operatorname{Hom}_S(I_Z(\mathcal M),X\times_SY)
       \mid\theta\circ\varepsilon_{\mathcal M}=(x,y)\},
\end{multline*}
which sends \((\varphi,\psi)\) to \(\varphi\times\psi\). This is clear,
since for \(a\in\mathcal O(Z)\),
\[
  (\varphi\circ a^*)\times(\psi\circ a^*)
  =(\varphi\times\psi)\circ a^*
  =a\cdot(\varphi\times\psi).
\]
The assertion about (\SGARefLink{sga3:ii:subsection:3-5-1:equation:3-8-2}{3.8.2}) follows similarly from \SGARefLink{sga3:ii:remark:3-2-1}{3.2.1}.

% Source page: Exp2 component-local p. 14; combined-reader p. 72.
\subsection*{\SGAInlineRefTarget{sga3:ii:subsection:3-9-0}3.9.0.}
\footnote{Editor's note: this paragraph has been added to explain the
introduction of the notion of an \(S\)-\(H\)-functor.}
If \(X\) is an \(S\)-group with unit section \(e:S\to X\), write
\[
  \underline{\operatorname{Lie}}(X/S,\mathcal M)
  =L^e_{X/S}(\mathcal M).
\]
Thus \(\underline{\operatorname{Lie}}(X/S,\mathcal M)\) is defined by the
following cartesian square.

\SGAThreeNativeDiagram[0.38\textwidth]
  {native_diagrams/exp2/exp2_p014_lie_kernel_cartesian_square.tex}
  {The cartesian square defining \(\underline{\operatorname{Lie}}(X/S,\mathcal M)\).}
  {fig:exp2-lie-kernel}
  {Polo--Gille Exp2-14oct24.pdf, component-local p. 14; combined-reader p. 72}

By \SGARefLink{sga3:ii:subsection:3-5-1:equation:3-8-1}{Corollary~3.8.1}, the projection
\(p:T_{X/S}(\mathcal M)\to X\) is a morphism of \(S\)-groups. Hence
\(\underline{\operatorname{Lie}}(X/S,\mathcal M)\) is an \(S\)-group,
identified by \(i\) with the kernel of \(p\).

If \(X/S\) also satisfies (E), \SGARefLink{sga3:ii:proposition:3-9}{Proposition~3.9} will show that this group
structure, induced by that of \(X\), agrees with the abelian-group
structure induced by functoriality in \(\mathcal M\) (cf. \SGARefLink{sga3:ii:subsection:3-5-1}{3.5.1}). In fact,
the result holds under the weaker assumption that \(X\) be an
\(S\)-functor in monoids, or more generally an \(S\)-functor in
\(H\)-sets, as defined next.

\subsection*{\SGAInlineRefTarget{sga3:ii:definition:3-9-0-1}\SGARefLink{sga3:ii:definition:3-9-0-1}{Definitions 3.9.0.1}.}
\begin{enumerate}
\item[(a)] An \emph{\(H\)-set} is a set \(X\) equipped with a composition
law having a two-sided identity, denoted \(e_X\) or simply \(e\).
\footnote{Editor's note: this terminology is inspired by the notion of an
\(H\)-space in topology.}
If \(f:X\to Y\) is a morphism of \(H\)-sets, its kernel is
\(\ker f=f^{-1}(e_Y)\), a sub-\(H\)-set of \(X\).

\item[(b)] An \emph{\(H\)-object} in a category \(\mathscr C\) is an
object \(X\) equipped with a morphism \(X\times X\to X\) and a section
of \(X\) over the final object that is a two-sided identity. Every
\(\mathscr C\)-monoid, and in particular every \(\mathscr C\)-group, is
a \(\mathscr C\)-\(H\)-object. An \(H\)-object in the category of functors
over \(S\) will be called an \emph{\(S\)-\(H\)-functor}.

% Source page: Exp2 component-local p. 15; combined-reader p. 73.
\item[(c)] If \(X\) is an \(S\)-\(H\)-functor (for example, an
\(S\)-group), and \(e:S\to X\) is its identity section, write
\[
  \underline{\operatorname{Lie}}(X/S,\mathcal M)=L^e_{X/S}(\mathcal M),
  \qquad
  \underline{\operatorname{Lie}}(X/S)
  =\underline{\operatorname{Lie}}(X/S,\mathcal O_S).
\]
\end{enumerate}

We now spell out the following special case of \SGARefLink{sga3:ii:subsection:3-5-1:equation:3-8-1}{3.8.1}.
\footnote{Editor's note: Corollary 3.9.0.2 has been added; it will be used
in \SGARefLink{sga3:ii:subsection:4-1}{4.1}.}

\SGAInlineRefTarget{sga3:ii:corollary:3-9-0-2}\subsection*{\SGARefLink{sga3:ii:corollary:3-9-0-2}{Corollary 3.9.0.2}.}
If \(X\) is an \(S\)-\(H\)-functor, respectively an \(S\)-group, then
\(T_{X/S}(\mathcal M)\) and
\(\underline{\operatorname{Lie}}(X/S,\mathcal M)\) are likewise
\(S\)-\(H\)-functors, respectively \(S\)-groups, and there are morphisms
of \(S\)-\(H\)-functors, respectively of \(S\)-groups, as follows.

\SGAThreeNativeDiagram[0.50\textwidth]
  {native_diagrams/exp2/exp2_p015_lie_tangent_split_sequence.tex}
  {The inclusion, projection, and zero section for the tangent functor.}
  {fig:exp2-lie-tangent-split}
  {Polo--Gille Exp2-14oct24.pdf, component-local p. 15; combined-reader p. 73}

Here \(i\) identifies \(\underline{\operatorname{Lie}}(X/S,\mathcal M)\)
with \(\ker p\), and \(s\) is a section of \(p\).

\SGAInlineRefTarget{sga3:ii:proposition:3-9}\subsection*{\SGARefLink{sga3:ii:proposition:3-9}{Proposition 3.9}.}
Let \(X\) be an \(S\)-\(H\)-functor satisfying (E) relative to \(S\).
The \(S\)-\(H\)-functor structure on
\(\underline{\operatorname{Lie}}(X/S,\mathcal M)\) induced by that of
\(X\) agrees with the \(S\)-group structure defined in \SGARefLink{sga3:ii:subsection:3-5-1}{3.5.1}.

The preceding discussion shows that
\(\underline{\operatorname{Lie}}(X/S,\mathcal M)\) is an \(H\)-object in
the category of \(\mathcal O_S\)-modules. The proposition follows from
the next lemma.
\footnote{Editor's note: this is inspired by the standard proof that the
fundamental group of an \(H\)-space is abelian.}

\SGAInlineRefTarget{sga3:ii:lemma:3-10}\subsection*{\SGARefLink{sga3:ii:lemma:3-10}{Lemma 3.10}.}
Let \(\mathscr C\) be a category. Let \(G\) be an \(H\)-object in the
category of \(\mathscr C\)-\(H\)-objects. Thus \(G\) is a
\(\mathscr C\)-\(H\)-object, whose composition law we denote
\(f:G\times G\to G\), together with a morphism of
\(\mathscr C\)-\(H\)-objects \(h:G\times G\to G\).
\footnote{Editor's note: \(G\times G\) is equipped with the law
\((G\times G)\times(G\times G)\to G\times G\),
\(((x,z),(y,t))\mapsto(f(x,y),f(z,t))\).}
Then \(f=h\), and \(f\) is commutative.

Evaluating the functors on a variable argument reduces the proof, as
usual, to the case in which \(\mathscr C\) is the category of sets. We
then have a set \(G\) and two maps \(f,h:G\times G\to G\) satisfying
\[
  h(f(x,y),f(z,t))=f(h(x,z),h(y,t)).
\]
Let \(e,u\in G\) be the respective identities for \(f\) and \(h\), so
that
\[
  f(e,x)=f(x,e)=x,
  \qquad h(u,x)=h(x,u)=x.
\]
First,
\[
  h(f(u,y),f(x,u))=f(x,y)=h(f(x,u),f(u,y)).
\]
Putting \(y=e\), respectively \(x=e\), gives
\begin{align*}
  x&=f(x,e)=h(f(u,e),f(x,u))=h(u,f(x,u))=f(x,u),\\
  y&=f(e,y)=h(f(e,u),f(u,y))=h(u,f(u,y))=f(u,y).
\end{align*}
Substituting these identities into the original relation yields
\[
  h(y,x)=f(x,y)=h(x,y).
\]
This proves the lemma and \SGARefLink{sga3:ii:proposition:3-9}{Proposition~3.9}. The following corollaries now
follow from \SGARefLink{sga3:ii:proposition:3-9}{3.9}.

\SGAInlineRefTarget{sga3:ii:corollary:3-9-1}\subsection*{\SGARefLink{sga3:ii:corollary:3-9-1}{Corollary 3.9.1}.}
If \(X\) is an \(S\)-\(H\)-functor satisfying (E) relative to \(S\),
every element of \(X(I_S(\mathcal M))\) that maps to the identity element
of \(X(S)\) is invertible.

% Source page: Exp2 component-local p. 16; combined-reader p. 74.
\SGAInlineRefTarget{sga3:ii:corollary:3-9-2}\subsection*{\SGARefLink{sga3:ii:corollary:3-9-2}{Corollary 3.9.2}.}
If \(X\) is an \(S\)-monoid satisfying (E) relative to \(S\), an element
of \(X(I_S(\mathcal M))\) is invertible if and only if its image in
\(X(S)\) is invertible.

\SGAInlineRefTarget{sga3:ii:corollary:3-9-3}\subsection*{\SGARefLink{sga3:ii:corollary:3-9-3}{Corollary 3.9.3}.}
If \(X\) is an \(S\)-group satisfying (E) relative to \(S\), the two
\(S\)-group laws on
\(\underline{\operatorname{Lie}}(X/S,\mathcal M)\) coincide.

\SGAInlineRefTarget{sga3:ii:corollary:3-9-4}\subsection*{\SGARefLink{sga3:ii:corollary:3-9-4}{Corollary 3.9.4}.}
\footnote{Editor's note: this corollary has been added; it amplifies
Remark 4.1.1.2 below.}
Let \(G\) be an \(S\)-group satisfying (E) relative to \(S\). For
\(n\in\mathbb Z\), let \(n_G:G\to G\) be the morphism of \(S\)-functors
defined by \(g\mapsto g^n\). Then the derived morphism
\[
  L(n_G):\underline{\operatorname{Lie}}(G/S)\longrightarrow
  \underline{\operatorname{Lie}}(G/S)
\]
is multiplication by \(n\); that is, it sends
\(x\in\underline{\operatorname{Lie}}(G/S)(S')\) to \(nx\).

The map \(n_G\) is not generally a group homomorphism, but it preserves
the identity section \(e:S\to G\), so its derivative indeed carries
\(\underline{\operatorname{Lie}}(G/S)=L^e_{G/S}\) to itself. If \(i\)
denotes the inclusion
\(\underline{\operatorname{Lie}}(G/S)\hookrightarrow T_{G/S}\), then
\(L(n_G)\) is characterized by
\[
  i(L(n_G)(x))=i(x)^n.
\]
By \SGARefLink{sga3:ii:proposition:3-9}{3.9}, the group laws on \(\underline{\operatorname{Lie}}(G/S)\) coming
from (E) and from the law on \(G\) agree. Hence
\(i(x)^n=i(nx)\), and therefore \(L(n_G)(x)=nx\).

Before drawing further consequences from Proposition~3.9, we prove one
more functoriality result.

\SGAInlineRefTarget{sga3:ii:proposition:3-11}\subsection*{\SGARefLink{sga3:ii:proposition:3-11}{Proposition 3.11}.}
In the situation of \SGARefLink{sga3:ii:section:1}{Section~1}, there is an isomorphism, functorial in
\(\mathcal M\),
\[
  T_{\underline{\operatorname{Hom}}_{Z/S}(X,Y)/S}(\mathcal M)
  \xrightarrow{\sim}
  \underline{\operatorname{Hom}}_{Z/S}
    \bigl(X,T_{Y/Z}(\mathcal M)\bigr).
\]
Indeed, by definition,
\[
  T_{\underline{\operatorname{Hom}}_{Z/S}(X,Y)/S}(\mathcal M)
  =\underline{\operatorname{Hom}}_S
    \bigl(I_S(\mathcal M),\underline{\operatorname{Hom}}_{Z/S}(X,Y)\bigr).
\]
Applying the isomorphism (4) of \SGARefLink{sga3:ii:definition:1-1}{1.1} with \(T=I_S(\mathcal M)\) gives
\begin{multline*}
  \underline{\operatorname{Hom}}_S
    \bigl(I_S(\mathcal M),\underline{\operatorname{Hom}}_{Z/S}(X,Y)\bigr)\\
  \simeq
  \underline{\operatorname{Hom}}_{Z/S}
    \bigl(X,\underline{\operatorname{Hom}}_Z
      (Z\times_SI_S(\mathcal M),Y)\bigr).
\end{multline*}
Since \(Z\times_SI_S(\mathcal M)\simeq I_Z(\mathcal M)\), this is
\[
  \underline{\operatorname{Hom}}_{Z/S}
    \bigl(X,\underline{\operatorname{Hom}}_Z(I_Z(\mathcal M),Y)\bigr)
  =\underline{\operatorname{Hom}}_{Z/S}
    \bigl(X,T_{Y/Z}(\mathcal M)\bigr).
\]

\SGAInlineRefTarget{sga3:ii:corollary:3-11-1}\subsection*{\SGARefLink{sga3:ii:corollary:3-11-1}{Corollary 3.11.1}.}
If \(Y/Z\) satisfies (E), then
\(\underline{\operatorname{Hom}}_{Z/S}(X,Y)/S\) satisfies (E), and the
isomorphism of \SGARefLink{sga3:ii:proposition:3-11}{3.11} respects the module structures over
\(\underline{\operatorname{Hom}}_{Z/S}(X,Y)\).
\footnote{Editor's note: writing
\(H=\underline{\operatorname{Hom}}_{Z/S}(X,Y)\), this implies in
particular that \(T_{H/S}(\mathcal M)\) is naturally a module over
\(\prod_{X\times_SH/H}\mathcal O_H\). This result appeared somewhat
surprising to the editors, so the proof has been given in detail.}

\emph{Proof.}
Let \(\mathcal M,\mathcal N\) be finite free \(\mathcal O_S\)-modules.
If \(Y/Z\) satisfies (E), then
\[
  T_{Y/Z}(\mathcal M\oplus\mathcal N)
  \simeq T_{Y/Z}(\mathcal M)\times_YT_{Y/Z}(\mathcal N).
\]
The right-hand side is a subfunctor of
\(T_{Y/Z}(\mathcal M)\times_ST_{Y/Z}(\mathcal N)\), and \SGARefLink{sga3:ii:remark:1-2}{1.2} \((*)\)
therefore gives
\begin{multline*}
  \underline{\operatorname{Hom}}_{Z/S}
    \bigl(X,T_{Y/Z}(\mathcal M\oplus\mathcal N)\bigr)\\
  \simeq
  \underline{\operatorname{Hom}}_{Z/S}\bigl(X,T_{Y/Z}(\mathcal M)\bigr)
  \mathop{\times}_{\underline{\operatorname{Hom}}_{Z/S}(X,Y)}
  \underline{\operatorname{Hom}}_{Z/S}\bigl(X,T_{Y/Z}(\mathcal N)\bigr).
\end{multline*}
Together with \SGARefLink{sga3:ii:proposition:3-11}{3.11}, this says that
\(\underline{\operatorname{Hom}}_{Z/S}(X,Y)\) satisfies (E), proving the
first assertion.

For the second, put
\(H=\underline{\operatorname{Hom}}_{Z/S}(X,Y)\), and let
\(\Delta:H'\to H\) be an \(S\)-morphism. Equivalently, let
\(\delta:H'\times_SX\to Y\) be the corresponding \(Z\)-morphism; it
makes \(H'\times_SX\) a \(Y\)-object. On the one hand, we have the
following commutative diagram.

% Source page: Exp2 component-local p. 17; combined-reader p. 75.
\SGAThreeNativeDiagram[0.82\textwidth]
  {native_diagrams/exp2/exp2_p017_hom_module_comparison_diagram.tex}
  {The first comparison diagram in the proof of \SGARefLink{sga3:ii:corollary:3-11-1}{Corollary 3.11.1}.}
  {fig:exp2-hom-module-comparison}
  {Polo--Gille Exp2-14oct24.pdf, component-local p. 17; combined-reader p. 75}

By \SGARefLink{sga3:ii:remark:1-3}{1.3}, the action of \(\alpha\in\mathcal O(H'\times_SX)\) on
\(\Psi\in\operatorname{Hom}_Y(H'\times_SX,T_{Y/Z}(\mathcal M))\) is
given, for every \(U\to S\) and every
\((h,x)\in\operatorname{Hom}_S(U,H'\times_SX)\), by
\[
  (\alpha\Psi)(h,x)=\alpha(h,x)\Psi(h,x),
\]
where \(U\) is viewed over \(Y\) through \(\delta\circ(h,x)\), and
\(\alpha(h,x)\in\mathcal O(U)\) acts through the
\(\mathcal O_Y\)-module structure on \(T_{Y/Z}(\mathcal M)\). By \SGARefLink{sga3:ii:remark:3-4-2}{3.4.2},
under the identification
\[
  \operatorname{Hom}_Y(H'\times_SX,T_{Y/Z}(\mathcal M))
  =\{\psi\in\operatorname{Hom}_Z(I_{H'\times_SX}(\mathcal M),Y)
       \mid\psi\circ\varepsilon_{\mathcal M}=\delta\},
\]
this action is
\begin{equation}
  (\alpha\psi)(m,h,x)
  =\psi(m\cdot\alpha(h,x),h,x)
  \tag{1}\label{sga3:ii:subsection:3-9-0:equation:1}
\end{equation}
for every
\((m,h,x)\in\operatorname{Hom}_S
  (U,I_S(\mathcal M)\times_SH'\times_SX)\).

% Source page: Exp2 component-local p. 18; combined-reader p. 76.
On the other hand, consider
\(T_{H/S}(\mathcal M)=\underline{\operatorname{Hom}}_S
  (I_S(\mathcal M),H)\). We have a second commutative diagram.

\SGAThreeNativeDiagram[0.84\textwidth]
  {native_diagrams/exp2/exp2_p018_tangent_hom_comparison_diagram.tex}
  {The second comparison diagram in the proof of \SGARefLink{sga3:ii:corollary:3-11-1}{Corollary 3.11.1}.}
  {fig:exp2-tangent-hom-comparison}
  {Polo--Gille Exp2-14oct24.pdf, component-local p. 18; combined-reader p. 76}

The bijection \((*)\) in the diagram is described as follows (cf. \SGARefLink{sga3:ii:section:1:equation:2}{1.1(2)}
and \SGARefLink{sga3:i:proposition:1-7-1}{I~1.7.1}). For every \(U\to S\) and
\((m,h,x)\in\operatorname{Hom}_S
  (U,I_S(\mathcal M)\times_SH'\times_SX)\), where \(U\) lies over \(Z\)
through \(U\xrightarrow{x}X\to Z\), one has
\(\Phi(m,h)\in\operatorname{Hom}_Z(X\times_SU,Y)\) and
\begin{equation}
  \varphi(m,h,x)=\Phi(m,h)\circ(x\times\operatorname{id}_U)
  \in\operatorname{Hom}_Z(U,Y).
  \tag{$\dagger$}\label{sga3:ii:subsection:3-9-0:equation:dagger}
\end{equation}
By \SGARefLink{sga3:ii:remark:3-4-2}{3.4.2}, with \(X\) replaced by \(H\) and \(X'\) by \(H'\), the
action of \(a\in\mathcal O(H')\) on
\(\Phi\in\operatorname{Hom}_S(I_{H'}(\mathcal M),H)\) is
\[
  (a\Phi)(m,h)=\Phi(m\cdot a(h),h).
\]
Consequently, if \(\varphi\), respectively \(a\varphi\), corresponds to
\(\Phi\), respectively \(a\Phi\), then \((\dagger)\) gives
\begin{equation}
  (a\varphi)(m,h,x)
  =\Phi(m\cdot a(h),h)\circ(x\times\operatorname{id}_U)
  =\varphi(m\cdot a(h),h,x).
  \tag{2}\label{sga3:ii:subsection:3-9-0:equation:2}
\end{equation}
Together with (1), this proves that the isomorphism
\[
  T_{H/S}(\mathcal M)
  \xrightarrow{\sim}
  \underline{\operatorname{Hom}}_{Z/S}
    (X,T_{Y/Z}(\mathcal M))
\]
of \SGARefLink{sga3:ii:corollary:3-11-1}{3.11.1} is an isomorphism of \(\mathcal O(H)\)-modules. Moreover, for
every \(H'\to H\), the \(\mathcal O(H')\)-module structure on
\(\operatorname{Hom}_H(H',T_{H/S}(\mathcal M))\) extends, functorially in
\(H'\), to an \(\mathcal O(H'\times_SX)\)-module structure.

Taking \(Z=S\) yields the next corollary.

\SGAInlineRefTarget{sga3:ii:corollary:3-11-2}\subsection*{\SGARefLink{sga3:ii:corollary:3-11-2}{Corollary 3.11.2}.}
There is an isomorphism, functorial in \(\mathcal M\),
\[
  T_{\underline{\operatorname{Hom}}_S(X,Y)/S}(\mathcal M)
  \xrightarrow{\sim}
  \underline{\operatorname{Hom}}_S(X,T_{Y/S}(\mathcal M)).
\]
If \(Y/S\) satisfies (E), then
\(\underline{\operatorname{Hom}}_S(X,Y)/S\) satisfies (E), and this
isomorphism respects the module structures over
\(\underline{\operatorname{Hom}}_S(X,Y)\).

Let \(u:X\to Y\) be an \(S\)-morphism, identified with the constant
morphism \(u:S\to\underline{\operatorname{Hom}}_S(X,Y)\) such that
\(u(f)=u\) for every \(f:S'\to S\). The fiber product of \(u\) with
\[
  \underline{\operatorname{Hom}}_S(X,T_{Y/S}(\mathcal M))
  \longrightarrow\underline{\operatorname{Hom}}_S(X,Y)
\]
identifies with
\(\underline{\operatorname{Hom}}_{Y/S}(X,T_{Y/S}(\mathcal M))\), where
\(X\) lies over \(Y\) through \(u\). Hence the definition of the tangent
space and \SGARefLink{sga3:ii:corollary:3-11-2}{Corollary~3.11.2} give the following result.
\footnote{Editor's note: the preceding sentences have been added.}

% Source page: Exp2 component-local p. 19; combined-reader p. 77.
\SGAInlineRefTarget{sga3:ii:corollary:3-11-3}\subsection*{\SGARefLink{sga3:ii:corollary:3-11-3}{Corollary 3.11.3}.}
For an \(S\)-morphism \(u:X\to Y\), there is an isomorphism, functorial
in \(\mathcal M\),
\[
  L^u_{\underline{\operatorname{Hom}}_S(X,Y)/S}(\mathcal M)
  \xrightarrow{\sim}
  \underline{\operatorname{Hom}}_{Y/S}(X,T_{Y/S}(\mathcal M)),
\]
where on the right \(X\) lies over \(Y\) through \(u\). If \(Y/S\)
satisfies (E), this is an isomorphism of \(\mathcal O_S\)-modules.
\footnote{Editor's note: the preceding sentence has been added.}

In particular, for \(Y=X\), the functor
\(\underline{\operatorname{End}}_S(X)\) is an \(S\)-functor in monoids,
hence an \(S\)-\(H\)-functor. Recall that
\(\underline{\operatorname{Lie}}(\underline{\operatorname{End}}_S(X)/S,
\mathcal M)\) denotes
\(L^e_{\underline{\operatorname{End}}_S(X)/S}(\mathcal M)\), with \(e\)
the identity section.

\SGAInlineRefTarget{sga3:ii:corollary:3-11-4}\subsection*{\SGARefLink{sga3:ii:corollary:3-11-4}{Corollary 3.11.4}.}
There is an isomorphism, functorial in \(\mathcal M\),
\[
  \underline{\operatorname{Lie}}
    (\underline{\operatorname{End}}_S(X)/S,\mathcal M)
  \xrightarrow{\sim}
  \prod_{X/S}T_{X/S}(\mathcal M).
\]
If \(X/S\) satisfies (E), this is an isomorphism of
\(\mathcal O_S\)-modules.
\footnotemark[38]

\SGAInlineRefTarget{sga3:ii:remark:3-11-5}\subsection*{\SGARefLink{sga3:ii:remark:3-11-5}{Remark 3.11.5}.}
\footnote{Editor's note: this remark has been added.}
Suppose that \(X/S\) satisfies (E). Then
\[
  \prod_{X/S}T_{X/S}(\mathcal M)
  =\underline{\operatorname{Hom}}_{X/S}(X,T_{X/S}(\mathcal M))
\]
is a module over \(\prod_{X/S}\mathcal O_X\). Explicitly, for every
\(S'\to S\),
\begin{multline*}
  \underline{\operatorname{Hom}}_{X/S}
    (X,T_{X/S}(\mathcal M))(S')\\
  =\{\psi\in\operatorname{Hom}_X(I_{S'}(\mathcal M)\times_SX,X)
    \mid\psi\circ(\varepsilon_{\mathcal M}\times\operatorname{id}_X)
      =\operatorname{pr}_X\}
\end{multline*}
is functorially an \(\mathcal O(X\times_SS')\)-module. This follows either
from \SGARefLink{sga3:ii:proposition:3-6}{3.6} and the properties of \(\prod_{X/S}\) (cf. \SGARefLink{sga3:ii:remark:1-2}{1.2}), or from the
proof of \SGARefLink{sga3:ii:corollary:3-11-1}{3.11.1}.

We now give a geometric interpretation of the tangent bundle. Let \(U\)
be an \(S\)-functor.
\footnote{Editor's note: the original has been simplified in what
follows.}
By \SGARefLink{sga3:i:corollary:1-7-2}{I~1.7.2}, there are isomorphisms, functorial in
\(\mathcal M\),
\begin{align*}
  T_{X/S}(\mathcal M)(U)
  &=\operatorname{Hom}_S
    \bigl(U,\underline{\operatorname{Hom}}_S(I_S(\mathcal M),X)\bigr)\\
  &\simeq\operatorname{Hom}_S
    \bigl(I_S(\mathcal M),\underline{\operatorname{Hom}}_S(U,X)\bigr)\\
  &\simeq\operatorname{Hom}_{I_S(\mathcal M)}
    \bigl(U_{I_S(\mathcal M)},X_{I_S(\mathcal M)}\bigr).
\end{align*}
In particular, the morphism \(\mathcal M\to0\) gives the following
commutative diagram.

\SGAThreeNativeDiagram[0.63\textwidth]
  {native_diagrams/exp2/exp2_p019_geometric_tangent_hom_square.tex}
  {The geometric interpretation of the projection \(\pi_{\mathcal M}\).}
  {fig:exp2-geometric-tangent-hom}
  {Polo--Gille Exp2-14oct24.pdf, component-local p. 19; combined-reader p. 77}

The second vertical arrow is obtained by base change along
\(\varepsilon_{\mathcal M}:S\to I_S(\mathcal M)\).
\footnote{Editor's note: via
\[
  \operatorname{Hom}_{I_S(\mathcal M)}
    (U_{I_S(\mathcal M)},X_{I_S(\mathcal M)})
  \simeq\operatorname{Hom}_S(U\times_SI_S(\mathcal M),X),
\]
this is equivalent to \SGARefLink{sga3:ii:remark:3-1-1}{Remark~3.1.1}. Likewise, \SGARefLink{sga3:ii:proposition:3-12}{3.12} is equivalent to the
first part of \SGARefLink{sga3:ii:remark:3-4-2}{Remark~3.4.2}.}

% Source page: Exp2 component-local p. 20; combined-reader p. 78.
\SGAInlineRefTarget{sga3:ii:proposition:3-12}\subsection*{\SGARefLink{sga3:ii:proposition:3-12}{Proposition 3.12}.}
Let \(h_0:U\to X\) be an \(S\)-morphism. Then
\(\operatorname{Hom}_X(U,T_{X/S}(\mathcal M))\) identifies with the set of
\(I_S(\mathcal M)\)-morphisms
\[
  U_{I_S(\mathcal M)}\longrightarrow X_{I_S(\mathcal M)}
\]
whose restriction to \(U\) is \(h_0\), where \(U\) is regarded as a
subobject of \(U\times_SI_S(\mathcal M)\) through
\(\operatorname{id}_U\times_S\varepsilon_{\mathcal M}\).

Taking \(U=X\) and \(h_0=\operatorname{id}_X\) gives the following.
\footnote{Editor's note: the original has been expanded in what follows.}

\SGAInlineRefTarget{sga3:ii:corollary:3-12-1}\subsection*{\SGARefLink{sga3:ii:corollary:3-12-1}{Corollary 3.12.1}.}
The set \(\Gamma(T_{X/S}(\mathcal M)/X)\) identifies with the set of
\(I_S(\mathcal M)\)-endomorphisms \(\phi\) of
\(X_{I_S(\mathcal M)}\) inducing the identity on \(X\); equivalently,
the following diagram commutes.
\footnote{Editor's note: when \(\mathcal M=\mathcal O_X\), these are
called \emph{infinitesimal endomorphisms} of \(X\); see \SGARefLink{sga3:ii:subsection:6-2}{6.2} at the end of
this expos\'e.}

\SGAThreeNativeDiagram[0.27\textwidth]
  {native_diagrams/exp2/exp2_p020_infinitesimal_endomorphism_square.tex}
  {An infinitesimal endomorphism inducing the identity on \(X\).}
  {fig:exp2-infinitesimal-endomorphism}
  {Polo--Gille Exp2-14oct24.pdf, component-local p. 20; combined-reader p. 78}

On the other hand, by \SGARefLink{sga3:ii:corollary:3-11-2}{3.11.2},
\[
  \Gamma(T_{X/S}(\mathcal M)/X)
  \simeq
  \underline{\operatorname{Lie}}
    (\underline{\operatorname{End}}_S(X)/S,\mathcal M)(S).
\]
If \(X/S\) satisfies (E), then
\(\underline{\operatorname{End}}_S(X)/S\) satisfies (E), so the latter
Lie functor is an \(\mathcal O_S\)-module by \SGARefLink{sga3:ii:proposition:3-6}{3.6} (indeed, by \SGARefLink{sga3:ii:remark:3-11-5}{3.11.5}, it
is a \(\prod_{X/S}\mathcal O_X\)-module). Applying \SGARefLink{sga3:ii:proposition:3-9}{3.9} gives the next
proposition.
\footnote{Editor's note: the original has been expanded in what follows.}

\SGAInlineRefTarget{sga3:ii:proposition:3-13}\subsection*{\SGARefLink{sga3:ii:proposition:3-13}{Proposition 3.13}.}
If \(X/S\) satisfies (E), the abelian group
\(\Gamma(T_{X/S}(\mathcal M)/X)\) identifies with the set of
\(I_S(\mathcal M)\)-endomorphisms of \(X_{I_S(\mathcal M)}\) that induce
the identity on \(X\). Consequently, every such endomorphism is an
automorphism.

\SGAInlineRefTarget{sga3:ii:corollary:3-13-1}\subsection*{\SGARefLink{sga3:ii:corollary:3-13-1}{Corollary 3.13.1}.}
Let \(u:X\to Y\) be an \(S\)-isomorphism, with \(Y/S\) satisfying (E).
Every \(I_S(\mathcal M)\)-morphism
\(X_{I_S(\mathcal M)}\to Y_{I_S(\mathcal M)}\) extending \(u\) is an
isomorphism.

\SGAInlineRefTarget{sga3:ii:corollary:3-13-2}\subsection*{\SGARefLink{sga3:ii:corollary:3-13-2}{Corollary 3.13.2}.}
If \(Y/S\) satisfies (E), the monomorphism
\[
  \underline{\operatorname{Isom}}_S(X,Y)
  \longrightarrow\underline{\operatorname{Hom}}_S(X,Y)
\]
induces, for every \(u\in\underline{\operatorname{Isom}}_S(X,Y)\), an
isomorphism
\[
  L^u_{\underline{\operatorname{Isom}}_S(X,Y)/S}(\mathcal M)
  \xrightarrow{\sim}
  L^u_{\underline{\operatorname{Hom}}_S(X,Y)/S}(\mathcal M).
\]

\emph{Proof.}
\footnote{Editor's note: the following proof has been added.}
For every \(S'\to S\), we must show that this map is a bijection on
\(S'\)-points. By base change (cf. \SGARefLink{sga3:ii:proposition:3-4}{3.4}), it suffices to take \(S'=S\).
The right-hand side then parametrizes the
\(I_S(\mathcal M)\)-morphisms
\(X_{I_S(\mathcal M)}\to Y_{I_S(\mathcal M)}\) extending \(u\), whereas
the left-hand side parametrizes those that are isomorphisms. The assertion
is therefore \SGARefLink{sga3:ii:corollary:3-13-1}{Corollary~3.13.1}.

% Source page: Exp2 component-local p. 21; combined-reader p. 79.
\SGAInlineRefTarget{sga3:ii:corollary:3-13-3}\subsection*{\SGARefLink{sga3:ii:corollary:3-13-3}{Corollary 3.13.3}.}
If \(X/S\) satisfies (E), the monomorphism
\(\underline{\operatorname{Aut}}_S(X)\to
\underline{\operatorname{End}}_S(X)\) induces, for every
\(u\in\underline{\operatorname{Aut}}_S(X)\), an isomorphism
\[
  L^u_{\underline{\operatorname{Aut}}_S(X)/S}(\mathcal M)
  \xrightarrow{\sim}
  L^u_{\underline{\operatorname{End}}_S(X)/S}(\mathcal M).
\]
In particular,
\[
  \underline{\operatorname{Lie}}
    (\underline{\operatorname{Aut}}_S(X)/S,\mathcal M)
  \xrightarrow{\sim}
  \underline{\operatorname{Lie}}
    (\underline{\operatorname{End}}_S(X)/S,\mathcal M)
  \xrightarrow{\sim}
  \prod_{X/S}T_{X/S}(\mathcal M),
\]
so that
\(\underline{\operatorname{Lie}}
  (\underline{\operatorname{Aut}}_S(X)/S,\mathcal M)\)
is a \(\prod_{X/S}\mathcal O_X\)-module.
\footnote{Editor's note: the preceding assertion has been added.}

\subsection*{\SGAInlineRefTarget{sga3:ii:subsection:3-14}3.14.}
Finally, suppose that \(X\) is representable.
\footnote{Editor's note: the reminder of \SGARefLink{sga3:ii:nota:2-2-2}{2.2.2} has been added, and the
sequel has been expanded.}
By \SGARefLink{sga3:ii:nota:2-2-2}{2.2.2}, the \(X\)-functor \(T_{X/S}\) is represented by
\(V(\Omega^1_{X/S})\), whence bijections
\[
  \Gamma(T_{X/S}/X)
  \simeq\operatorname{Hom}_X(\Omega^1_{X/S},\mathcal O_X)
  \simeq\operatorname{Der}_{\mathcal O_S}(\mathcal O_X).
\]
This also follows from the preceding results. By \SGARefLink{sga3:ii:proposition:3-13}{3.13},
\(\Gamma(T_{X/S}/X)\) identifies with the set of infinitesimal
endomorphisms of \(X\), that is, the \(I_S\)-endomorphisms of
\(X_{I_S}\) that induce the identity on \(X\). The schemes \(X\) and
\(X_{I_S}\) have the same underlying topological space, and their
structure sheaves are \(\mathcal O_X\) and
\[
  \mathscr D_{\mathcal O_X}=\mathcal O_X\oplus\mathcal M,
  \qquad \mathcal M=\mathcal O_X,
\]
where \(\mathcal M\) is a square-zero ideal. Let
\(\pi:\mathscr D_{\mathcal O_X}\to\mathcal O_X\) be the morphism of
\(\mathcal O_X\)-algebras vanishing on \(\mathcal M\). Giving an
infinitesimal endomorphism of \(X\) is equivalent to giving a morphism of
\(\mathcal O_S\)-algebras
\(\phi:\mathcal O_X\to\mathscr D_{\mathcal O_X}\) satisfying
\(\pi\circ\phi=\operatorname{id}_{\mathcal O_X}\), hence to giving an
\(\mathcal O_S\)-derivation of \(\mathcal O_X\).

If \(D,D'\in\operatorname{Der}_{\mathcal O_S}(\mathcal O_X)\), and
\(\phi_D\) denotes the infinitesimal endomorphism corresponding to \(D\),
then
\[
  \phi_{D+D'}=\phi_D\circ\phi_{D'}.
\]
Thus the identification
\[
  \{\text{infinitesimal endomorphisms of }X\}
  \simeq\operatorname{Der}_{\mathcal O_S}(\mathcal O_X)
\]
is an isomorphism of abelian groups. In view of \SGARefLink{sga3:ii:proposition:3-13}{3.13} and \SGARefLink{sga3:ii:remark:3-11-5}{3.11.5}, we have
therefore constructed an isomorphism of abelian groups, indeed of
\(\mathcal O(X)\)-modules,
\[
  \Gamma(T_{X/S}/X)
  \xrightarrow{\sim}\operatorname{Der}_{\mathcal O_S}(\mathcal O_X),
\]
recovering the classical interpretation of tangent vector fields as
derivations of the structure sheaf.
\footnote{Editor's note: this isomorphism is also functorial in \(S\). If
\(S'\to S\) and \(X'=X_{S'}\), then
\[
  \underline{\operatorname{Lie}}
    (\underline{\operatorname{Aut}}_S(X)/S)(S')
  \simeq\Gamma(T_{X'/S'}/X')
  \simeq\operatorname{Der}_{\mathcal O_{S'}}(\mathcal O_{X'}).
\]}
Finally,
\[
  \Gamma(T_{X/S}/X)=H^0(X,\mathfrak g_{X/S}),
\]
where \(\mathfrak g_{X/S}\) is the dual of \(\Omega^1_{X/S}\).

% Section 4 begins on Exp2 component-local p. 22; combined-reader p. 80.
% It is intentionally excluded from this component.

% END INCLUDED SGA3_English_current_buildable_source_20260823/source/inputs/II/tex/components/11_expose_II_s3_tangent_bundle_condition_E.tex

% BEGIN INCLUDED SGA3_English_current_buildable_source_20260823/source/inputs/II/tex/components/12_expose_II_s41_s44_lie_algebras.tex
% SGA 3, Expose II, Loop 1 English reconstruction.
% Scope: Section 4.1 through the end of Example 4.4.2, stopping immediately
% before Proposition 4.5.
% Authority: Polo--Gille Exp2-14oct24.pdf, component-local pages 22--34
% (combined-reader pages 80--92; printed pages 70--82). Force-OCR is a
% drafting locator; English wording was compared with Jacob Reinhold's
% CC BY 4.0 Markdown lineage.

\setcounter{footnote}{48}

% Source page: Exp2 component-local p. 22; combined-reader p. 80; printed p. 70.
\section*{\SGAInlineRefTarget{sga3:ii:section:4}4. Tangent Space to a Group --- Lie Algebras}
\addcontentsline{toc}{section}{4. Tangent Space to a Group --- Lie Algebras}

\subsection*{\SGAInlineRefTarget{sga3:ii:subsection:4-1}4.1.}
Let \(G\) be a group functor over \(S\). By \SGARefLink{sga3:ii:corollary:3-9-0-2}{3.9.0.2},
\(T_{G/S}(\mathcal M)\) and
\(\underline{\operatorname{Lie}}(G/S,\mathcal M)\) are endowed with
group structures over \(S\), and there are group morphisms

\SGAThreeNativeDiagram[0.52\textwidth]
  {native_diagrams/exp2/exp2_p022_split_tangent_sequence.tex}
  {The split tangent sequence of the group functor \(G\).}
  {fig:exp2-split-tangent-sequence}
  {Polo--Gille Exp2-14oct24.pdf, component-local p. 22; combined-reader p. 80; printed p. 70}

By definition, \(i\) identifies
\(\underline{\operatorname{Lie}}(G/S,\mathcal M)\) with the kernel of
\(p\), and \(s\) is a section of \(p\). It follows from \SGARefLink{sga3:i:proposition:2-3-7}{I~2.3.7} that this
sequence identifies \(T_{G/S}(\mathcal M)\) with the semidirect product
of \(G\) by \(\underline{\operatorname{Lie}}(G/S,\mathcal M)\).

\SGAInlineRefTarget{sga3:ii:definition:4-1-a}\subsection*{\SGARefLink{sga3:ii:definition:4-1-a}{Definition 4.1.A}.}
The corresponding action of \(G\) on
\(\underline{\operatorname{Lie}}(G/S,\mathcal M)\) is denoted
\[
  \operatorname{Ad}:G\longrightarrow
  \underline{\operatorname{Aut}}_{S\text{-gr.}}
    \bigl(\underline{\operatorname{Lie}}(G/S,\mathcal M)\bigr)
\]
and is called the \emph{adjoint representation} of \(G\), relative to
\(\mathcal M\). Thus, by definition, for \(x\in G(S')\) and
\(X\in\underline{\operatorname{Lie}}(G/S,\mathcal M)(S')\),
\[
  \operatorname{Ad}(x)X
  =i^{-1}\bigl(s(x)i(X)s(x)^{-1}\bigr).
\]
If \(G\) is commutative, then \(T_{G/S}(\mathcal M)\) is commutative as
well and \(\operatorname{Ad}(x)X=X\).
\footnote{Editor's note: the numbering \SGARefLink{sga3:ii:definition:4-1-a}{4.1.A} and \SGARefLink{sga3:ii:definition:4-1-b}{4.1.B} has been added in
order to bring these definitions into relief.}

\SGAInlineRefTarget{sga3:ii:definition:4-1-b}\subsection*{\SGARefLink{sga3:ii:definition:4-1-b}{Definition 4.1.B}.}
Let \(G,H\) be group functors over \(S\), and let \(f:G\to H\) be a group
morphism. Functoriality gives a morphism of exact sequences compatible
with the canonical sections.

\SGAThreeNativeDiagram[0.66\textwidth]
  {native_diagrams/exp2/exp2_p022_derived_exact_sequences.tex}
  {Functoriality of the split tangent sequence.}
  {fig:exp2-derived-exact-sequences}
  {Polo--Gille Exp2-14oct24.pdf, component-local p. 22; combined-reader p. 80; printed p. 70}

The morphism \(L(f)\), also denoted
\(\underline{\operatorname{Lie}}(f)\), is called the \emph{derived
morphism} of \(f\).
\footnote{Editor's note: if \(f\) is a monomorphism, then so are \(L(f)\)
and \(T(f)\); cf. \SGARefLink{sga3:ii:proposition:3-7}{3.7}.}

\SGAInlineRefTarget{sga3:ii:remark:4-1-c}\subsection*{\SGARefLink{sga3:ii:remark:4-1-c}{Remark 4.1.C}.}
If \(G/S\) and \(H/S\) satisfy condition \(E\), then \(L(f)\) respects the
\(\mathcal O_S\)-module structures obtained from ``functoriality in
\(\mathcal M\)'' (cf. \SGARefLink{sga3:ii:proposition:3-6}{3.6}).
\footnote{Editor's note: this remark has been added.}

\SGAInlineRefTarget{sga3:ii:proposition:4-1-1}\subsection*{\SGARefLink{sga3:ii:proposition:4-1-1}{Proposition 4.1.1}.}
Let \(g\in G(S)\). Then
\[
  \operatorname{Ad}(g):
  \underline{\operatorname{Lie}}(G/S,\mathcal M)
  \longrightarrow
  \underline{\operatorname{Lie}}(G/S,\mathcal M)
\]
is the derived morphism of
\(\operatorname{Int}(g):G\to G\).

Indeed,
\(\operatorname{Ad}(g)X=i^{-1}(\operatorname{Int}(g)i(X))\), which is
precisely \(L(\operatorname{Int}(g))(X)\) by the definition of the
derived morphism.

Suppose that \(G/S\) satisfies \(E\). By \SGARefLink{sga3:ii:proposition:3-9}{Proposition~3.9}, the group
structure just defined on
\(\underline{\operatorname{Lie}}(G/S,\mathcal M)\) is exactly the one
induced by its \(\mathcal O_S\)-module structure (which is defined by
(E)). \SGARefLink{sga3:ii:proposition:4-1-1}{Proposition~4.1.1} and the functoriality of the
\(\mathcal O_S\)-action (cf. \SGARefLink{sga3:ii:proposition:3-6}{3.6}) therefore yield the following
corollary.

% Source page: Exp2 component-local p. 23; combined-reader p. 81; printed p. 71.
\SGAInlineRefTarget{sga3:ii:corollary:4-1-1-1}\subsection*{\SGARefLink{sga3:ii:corollary:4-1-1-1}{Corollary 4.1.1.1}.}
Suppose that \(G/S\) satisfies \(E\). Then \(\operatorname{Ad}\) maps \(G\)
into the subgroup
\[
  \underline{\operatorname{Aut}}_{\mathcal O_S\text{-mod.}}
    \bigl(\underline{\operatorname{Lie}}(G/S,\mathcal M)\bigr)
\]
of
\(\underline{\operatorname{Aut}}_{S\text{-gr.}}
(\underline{\operatorname{Lie}}(G/S,\mathcal M))\). In other words, for
every \(g\in G(S')\), the automorphism \(\operatorname{Ad}(g)\) respects
the \(\mathcal O_{S'}\)-module structure on
\(\underline{\operatorname{Lie}}(G'/S',\mathcal M)\), where
\(G'=G\times_SS'\). Thus \(\operatorname{Ad}\) is a \emph{linear
representation} (cf. \SGARefLink{sga3:i:definition:3-2}{I~3.2}) of \(G\) on the \(\mathcal O_S\)-module
\(\underline{\operatorname{Lie}}(G/S,\mathcal M)\).

\SGAInlineRefTarget{sga3:ii:remark:4-1-1-2}\subsection*{\SGARefLink{sga3:ii:remark:4-1-1-2}{Remarks 4.1.1.2}.}
\textit{a)} The functor \(G/S\) satisfies \(E\) if and only if, for every
pair \((\mathcal M,\mathcal N)\) of finite free
\(\mathcal O_S\)-modules, the following diagram, obtained by applying
\(\underline{\operatorname{Lie}}(G/S,-)\) to diagram \((*)\) of \SGARefLink{sga3:ii:definition:2-1}{2.1}, is
cartesian.

\SGAThreeNativeDiagram[0.45\textwidth]
  {native_diagrams/exp2/exp2_p023_condition_e_cartesian.tex}
  {The cartesian criterion for condition \(E\).}
  {fig:exp2-condition-e-cartesian}
  {Polo--Gille Exp2-14oct24.pdf, component-local p. 23; combined-reader p. 81; printed p. 71}

\textit{b)} Suppose that \(G/S\) satisfies \(E\). The derived morphism of
the group law
\(\pi:G\times_SG\to G\) is the addition law on
\(\underline{\operatorname{Lie}}(G/S,\mathcal M)\). Notice that
\(\pi\) is not a group morphism, but \(\pi(e,e)=e\); hence the derived
morphism \(L(\pi)\) maps
\[
  T^{(e,e)}_{(G\times_SG)/S}(\mathcal M)
  =\underline{\operatorname{Lie}}(G/S,\mathcal M)
   \times_S\underline{\operatorname{Lie}}(G/S,\mathcal M)
\]
to \(\underline{\operatorname{Lie}}(G/S,\mathcal M)\); cf. \SGARefLink{sga3:ii:proposition:3-7}{3.7} and
\SGARefLink{sga3:ii:proposition:3-8}{3.8}. Likewise, for \(n\in\mathbb Z\), let \(n_G:G\to G\) denote the
\(S\)-functor morphism \(g\mapsto g^n\). Then \(L(n_G)\) is
multiplication by \(n\) on \(\underline{\operatorname{Lie}}(G/S)\); cf.
\SGARefLink{sga3:ii:corollary:3-9-4}{3.9.4}.

\subsection*{\SGAInlineRefTarget{sga3:ii:subsection:4-1-2-0}4.1.2.0.}
Consider the \(S\)-functor
\(
  \underline{\operatorname{Hom}}_{G/S}
    (G,T_{G/S}(\mathcal M)).
\)\footnote{Editor's note: the original has been expanded in what follows
to display the action of \(G\) on the functors
\(\underline{\operatorname{Hom}}_{G/S}(G,T_{G/S}(\mathcal M))\) and
\(\underline{\operatorname{Hom}}_S
(G,\underline{\operatorname{Lie}}(G/S,\mathcal M))\).}
For every \(S'\to S\), one has
\(T_{G/S}(\mathcal M)_{S'}\simeq T_{G_{S'}/S'}(\mathcal M)\)
(cf. \SGARefLink{sga3:ii:proposition:3-4}{3.4}), and therefore
\[
\begin{aligned}
 &\underline{\operatorname{Hom}}_{G/S}
   (G,T_{G/S}(\mathcal M))(S')\\
 &\qquad\simeq
 \operatorname{Hom}_{G_{S'}}
   \bigl(G_{S'},T_{G_{S'}/S'}(\mathcal M)\bigr)
 =\Gamma\bigl(T_{G_{S'}/S'}(\mathcal M)/G_{S'}\bigr).
\end{aligned}
\]
First note the isomorphism, functorial in \(S'\),
\begin{equation}
 \operatorname{Hom}_{S'}
  \bigl(G_{S'},\underline{\operatorname{Lie}}(G_{S'}/S',\mathcal M)\bigr)
 \xrightarrow{\sim}
 \Gamma\bigl(T_{G_{S'}/S'}(\mathcal M)/G_{S'}\bigr),
 \tag{*}\label{sga3:ii:subsection:4-1-2-0:equation:star}
\end{equation}
which sends \(f:G_{S'}\to
\underline{\operatorname{Lie}}(G_{S'}/S',\mathcal M)\) to the section
\(s_f:G_{S'}\to T_{G_{S'}/S'}(\mathcal M)\) defined, for every
\(S''\to S'\) and \(g\in G(S'')\), by
\[
  s_f(g)=i(f(g))s(g).
\]

% Source page: Exp2 component-local p. 24; combined-reader p. 82; printed p. 72.
Let \(h\) be an automorphism of the functor \(G_{S'}\) over \(S'\), not
necessarily preserving the group structure. To every section \(\tau\) of
\(T_{G_{S'}/S'}(\mathcal M)\) one associates, by transport of structure,
a section \(h(\tau)\): it is the unique section that makes the following
diagram commute.

\SGAThreeNativeDiagram[0.43\textwidth]
  {native_diagrams/exp2/exp2_p024_transport_square.tex}
  {Transport of a tangent section by an automorphism \(h\).}
  {fig:exp2-transport-square}
  {Polo--Gille Exp2-14oct24.pdf, component-local p. 24; combined-reader p. 82; printed p. 72}

In particular, take \(h=t_x\), right translation by an element
\(x\in G(S')\); thus \(h(g)=t_x(g)=g\cdot x\) for all
\(g\in G(S'')\), \(S''\to S'\). One then has
\[
  t_x(s_f)=s_{t_x(f)},
\]
where \(t_x(f):G_{S'}\to
\underline{\operatorname{Lie}}(G_{S'}/S',\mathcal M)\) is given by
\[
  t_x(f)(g)=f(g\cdot x^{-1}).
\]

Let \(G\) act by right translations on
\[
 \underline{\operatorname{Hom}}_{G/S}(G,T_{G/S}(\mathcal M))
 \quad\text{and}\quad
 \underline{\operatorname{Hom}}_S
   (G,\underline{\operatorname{Lie}}(G/S,\mathcal M))
\]
as follows. For \(S'\to S\) and \(x\in G(S')\), let
\[
\begin{aligned}
 \sigma&\in\Gamma(T_{G_{S'}/S'}(\mathcal M)/G_{S'}),\\
 f&\in\operatorname{Hom}_{S'}
 (G_{S'},\underline{\operatorname{Lie}}(G_{S'}/S',\mathcal M)).
\end{aligned}
\]
Set
\[
 (\sigma\cdot x)(g)=\sigma(g\cdot x^{-1})\cdot s(x),
 \qquad
 (f\cdot x)(g)=f(g\cdot x^{-1})
\]
for every \(g\in G(S'')\), \(S''\to S'\). The isomorphism \((*)\)
respects these \(G\)-actions.

Consequently, under this isomorphism, the elements of
\[
 \underline{\operatorname{Hom}}_{G/S}
   (G,T_{G/S}(\mathcal M))^G(S')
\]
correspond to the constant morphisms from \(G_{S'}\) to
\(\underline{\operatorname{Lie}}(G_{S'}/S',\mathcal M)\), that is, to
the morphisms factoring through \(G_{S'}\to S'\); equivalently, they
correspond to the elements of
\(\underline{\operatorname{Lie}}(G/S,\mathcal M)(S')\).

\paragraph{Terminology.}
The elements of
\(\underline{\operatorname{Hom}}_{G/S}
(G,T_{G/S}(\mathcal M))^G(S')\) are called the sections of
\(T_{G_{S'}/S'}(\mathcal M)\) invariant under right translation.

\SGAInlineRefTarget{sga3:ii:proposition:4-1-2}\subsection*{\SGARefLink{sga3:ii:proposition:4-1-2}{Proposition 4.1.2}.}
The map
\[
 \underline{\operatorname{Lie}}(G/S,\mathcal M)(S)
 \longrightarrow \Gamma(T_{G/S}(\mathcal M)/G),
 \qquad X\longmapsto (x\mapsto X\cdot x),
\]
is a bijection onto the set of sections invariant under right
translation.
\footnote{Statements 4.1.2, 4.1.3, and 4.1.4 can be obtained more simply
by observing that the automorphisms of \(G\) invariant under right
translations are the left translations. Editor's note: see, for example,
the proofs of Propositions 4.6 and 2.5 of \textup{[DG70]}, \S II.4.}

% Source page: Exp2 component-local p. 25; combined-reader p. 83; printed p. 73.
Likewise, let \(G\) act on the right on
\(\underline{\operatorname{End}}_{I_S(\mathcal M)/S}
(G_{I_S(\mathcal M)})\) by
\[
 (u\cdot x)(g)=u(g\cdot x^{-1})\cdot x,
\]
for \(S'\to S\), \(x\in G(S')\),
\(u\in\operatorname{End}_{I_{S'}(\mathcal M)}
(G_{I_{S'}(\mathcal M)})\), and \(g\in G(S'')\),
\(S''\to I_{S'}(\mathcal M)\). The morphism of 3.12.1
\[
 \underline{\operatorname{Hom}}_{G/S}(G,T_{G/S}(\mathcal M))
 \longrightarrow
 \underline{\operatorname{End}}_{I_S(\mathcal M)/S}
   (G_{I_S(\mathcal M)})
\]
respects the \(G\)-actions. For each \(S'\to S\), it therefore induces a
bijection between
\(\Gamma(T_{G_{S'}/S'}(\mathcal M)/G_{S'})\) and the set of
\(I_{S'}(\mathcal M)\)-endomorphisms \(u\) of
\(G_{I_{S'}(\mathcal M)}\) that induce the identity on \(G\) and
``commute with right translations,'' meaning that
\(u_{S''}\cdot x=u_{S''}\) for every \(S''\to S'\) and
\(x\in G(S'')\).

\SGAInlineRefTarget{sga3:ii:proposition:4-1-3}\subsection*{\SGARefLink{sga3:ii:proposition:4-1-3}{Proposition 4.1.3}.}
There is a bijection, functorial in \(G\), between
\(\underline{\operatorname{Lie}}(G/S,\mathcal M)(S)\) and the set of
\(I_S(\mathcal M)\)-endomorphisms of \(G_{I_S(\mathcal M)}\) that
induce the identity on \(G\) and commute with the right translations of
\(G\), in the sense just specified.

Taking \SGARefLink{sga3:ii:proposition:3-13}{3.13} into account gives the following theorem.

\SGAInlineRefTarget{sga3:ii:theorem:4-1-4}\subsection*{\SGARefLink{sga3:ii:theorem:4-1-4}{Theorem 4.1.4}.}
Let \(G\) be an \(S\)-functor in groups, and suppose that \(G/S\)
satisfies \(E\). Then
\(\underline{\operatorname{Lie}}(G/S,\mathcal M)(S)\) is identified,
functorially in \(G\), with the group of
\(I_S(\mathcal M)\)-automorphisms of \(G_{I_S(\mathcal M)}\) that
induce the identity on \(G\) and commute with the right translations of
\(G\).

For \(\mathcal M=\mathcal O_S\), this recovers one of the classical
definitions of the Lie algebra of a group.

\subsection*{\SGAInlineRefTarget{sga3:ii:subsection:4-2-0}4.2.0.}
Before proceeding, let us establish further consequences of
\SGARefLink{sga3:ii:proposition:3-11}{3.11}.\footnote{Editor's note: \SGARefLink{sga3:ii:subsection:4-2-0}{paragraph 4.2.0} has been added; its results
are used implicitly in \SGARefLink{sga3:ii:subsection:4-2}{4.2} and \SGARefLink{sga3:ii:theorem:4-7}{4.7} of the original.} Let
\(X,Y\) be over \(Z\), with \(Z\) over \(S\), as in \SGARefLink{sga3:ii:section:1}{Section~1}. As was
seen in \SGARefLink{sga3:ii:proposition:3-11}{3.11}, the isomorphisms \SGARefLink{fig:exp2-hom-adjunction-triangle}{1.1(4)}

\SGAThreeNativeDiagram[0.78\textwidth]
  {native_diagrams/exp2/exp2_p025_hom_adjunction_triangle.tex}
  {The adjunction isomorphisms of \SGARefLink{fig:exp2-hom-adjunction-triangle}{1.1(4)}.}
  {fig:exp2-hom-adjunction-p25}
  {Polo--Gille Exp2-14oct24.pdf, component-local p. 25; combined-reader p. 83; printed p. 73}

induce the following isomorphism \(\theta\).

\SGAThreeNativeDiagram[0.78\textwidth]
  {native_diagrams/exp2/exp2_p025_tangent_adjunction_triangle.tex}
  {The induced tangent-space isomorphism \(\theta\).}
  {fig:exp2-tangent-adjunction-p25}
  {Polo--Gille Exp2-14oct24.pdf, component-local p. 25; combined-reader p. 83; printed p. 73}

% Source page: Exp2 component-local p. 26; combined-reader p. 84; printed p. 74.
By \SGARefLink{sga3:ii:remark:1-3}{1.3}, if \(Y\) is a \(Z\)-group, then
\(\underline{\operatorname{Hom}}_Z(V,Y)\) is a \(Z\)-group for every
\(V\to Z\), in particular for \(V=I_Z(\mathcal M)\). Explicitly, if
\(Z''\to Z'\to Z\) and
\(\varphi,\psi\in\operatorname{Hom}_Z(V_{Z'},Y)(Z')\), their product is
defined by
\[
  (\varphi\cdot\psi)(v)=\varphi(v)\psi(v)
  \qquad (v\in V_{Z'}(Z'')).
\]

Suppose now that \(X\) and \(Y\) are \(Z\)-groups.

\SGAInlineRefTarget{sga3:ii:definition:4-2-0-1}\subsection*{\SGARefLink{sga3:ii:definition:4-2-0-1}{Definition 4.2.0.1}.}
We denote by
\(\underline{\operatorname{Hom}}_{(Z/S)\text{-gr.}}(X,Y)\) the
subfunctor of \(\underline{\operatorname{Hom}}_{Z/S}(X,Y)\) defined, for
every \(S'\to S\), by
\begin{equation}
 \underline{\operatorname{Hom}}_{(Z/S)\text{-gr.}}(X,Y)(S')
 =\operatorname{Hom}_{Z_{S'}\text{-gr.}}(X_{S'},Y_{S'}).
 \tag{3}\label{sga3:ii:subsection:4-2-0:equation:3}
\end{equation}
The same definition applies with \(Y\) replaced by the \(Z\)-group
\(T_{Y/Z}(\mathcal M)\).

Under the isomorphisms (2) above,
\(T_{\underline{\operatorname{Hom}}_{(Z/S)\text{-gr.}}(X,Y)/S}
(\mathcal M)(S')\) corresponds to the \(Z_{S'}\)-morphisms
\[
 \varphi:X_{S'}\times_{S'}I_{S'}(\mathcal M)\longrightarrow Y_{S'}
\]
that are multiplicative in \(X\), that is,
\[
 \varphi(x_1x_2,m)=\varphi(x_1,m)\varphi(x_2,m).
\]
These in turn correspond to the morphisms of \(Z_{S'}\)-groups
\(X_{S'}\to T_{Y/Z}(\mathcal M)_{S'}\). Hence:

\SGAInlineRefTarget{sga3:ii:proposition:4-2-0-2}\subsection*{\SGARefLink{sga3:ii:proposition:4-2-0-2}{Proposition 4.2.0.2}.}
Let \(X,Y\) be \(Z\)-groups, where \(Z\) is over \(S\). There are
isomorphisms of \(S\)-functors, functorial in \(\mathcal M\),
\[
 T_{\underline{\operatorname{Hom}}_{(Z/S)\text{-gr.}}(X,Y)/S}
   (\mathcal M)
 \xrightarrow{\sim}
 \underline{\operatorname{Hom}}_{(Z/S)\text{-gr.}}
   (X,T_{Y/Z}(\mathcal M)).
\]

Taking \(Z=S\) gives the following consequence. If \(Y\) is a
commutative \(S\)-group, then so is \(T_{Y/S}(\mathcal M)\), and hence
so are
\(H=\underline{\operatorname{Hom}}_{S\text{-gr.}}(X,Y)\),
\(\underline{\operatorname{Hom}}_{S\text{-gr.}}
(X,T_{Y/S}(\mathcal M))\), and \(T_{H/S}(\mathcal M)\).

\SGAInlineRefTarget{sga3:ii:corollary:4-2-0-3}\subsection*{\SGARefLink{sga3:ii:corollary:4-2-0-3}{Corollary 4.2.0.3}.}
Let \(X,Y\) be \(S\)-groups. There are isomorphisms of \(S\)-functors,
functorial in \(\mathcal M\),
\[
 T_{\underline{\operatorname{Hom}}_{S\text{-gr.}}(X,Y)/S}
   (\mathcal M)
 \xrightarrow{\sim}
 \underline{\operatorname{Hom}}_{S\text{-gr.}}
   (X,T_{Y/S}(\mathcal M)).
\]
If \(Y\) is commutative, these are moreover isomorphisms of abelian
\(S\)-groups.

\SGAInlineRefTarget{sga3:ii:definition:4-2-0-4}\subsection*{\SGARefLink{sga3:ii:definition:4-2-0-4}{Definition 4.2.0.4}.}
If \(Y\) is an \(\mathcal O_S\)-module, the functor
\(T_{Y/S}(\mathcal M)\), respectively
\(\underline{\operatorname{Lie}}(Y/S,\mathcal M)\), carries an
\(\mathcal O_S\)-module structure induced by that of \(Y\). Equipped
with this structure, it will be denoted
\(T'_{Y/S}(\mathcal M)\), respectively
\(\underline{\operatorname{Lie}}'(Y/S,\mathcal M)\).
\footnote{Editor's note: the prime is used to distinguish the
\(\mathcal O_S\)-module structure induced by that of \(Y\) from the one
that will be obtained from condition \(E\).}

Consequently, if \(X,Y\) are \(\mathcal O_S\)-modules, then
\(T'_{Y/S}(\mathcal M)=
\underline{\operatorname{Hom}}_S(I_S(\mathcal M),Y)\). If
\(H=\underline{\operatorname{Hom}}_{\mathcal O_S\text{-mod.}}(X,Y)\),
then
\(\underline{\operatorname{Hom}}_{\mathcal O_S\text{-mod.}}
(X,T'_{Y/S}(\mathcal M))\) and \(T'_{H/S}(\mathcal M)\) also carry
\(\mathcal O_S\)-module structures.

% Source page: Exp2 component-local p. 27; combined-reader p. 85; printed p. 75.
\SGAInlineRefTarget{sga3:ii:corollary:4-2-0-5}\subsection*{\SGARefLink{sga3:ii:corollary:4-2-0-5}{Corollary 4.2.0.5}.}
If \(X,Y\) are \(\mathcal O_S\)-modules, there are isomorphisms of
\(\mathcal O_S\)-modules, functorial in \(\mathcal M\),
\[
 T'_{\underline{\operatorname{Hom}}_{\mathcal O_S\text{-mod.}}(X,Y)/S}
   (\mathcal M)
 \xrightarrow{\sim}
 \underline{\operatorname{Hom}}_{\mathcal O_S\text{-mod.}}
   (X,T'_{Y/S}(\mathcal M)).
\]

\SGAInlineRefTarget{sga3:ii:definition:4-2-a}\subsection*{\SGARefLink{sga3:ii:definition:4-2-a}{Definition 4.2.A}.}
Let \(X,L\) be \(S\)-groups, with \(X\) acting on \(L\) by group
automorphisms (cf. \SGARefLink{sga3:i:definition:2-3-5}{I~2.3.5}). We define the subfunctor
\(\underline Z^1_S(X,L)\) of
\(\underline{\operatorname{Hom}}_S(X,L)\) by
\[
\begin{aligned}
 \underline Z^1_S(X,L)(S')
 =\{\,\varphi\in\operatorname{Hom}_{S'}(X_{S'},L_{S'})\mid{}
 &\varphi(x_1x_2)
   =\varphi(x_1)\bigl(x_1\cdot\varphi(x_2)\bigr),\\[-2pt]
 &x_1,x_2\in X(S''),\ S''\to S'\,\}.
\end{aligned}
\]
It is called the \emph{functor of crossed homomorphisms} from \(X\) to
\(L\).
\footnote{Editor's note: the original has been expanded below; in
particular, Definition 4.2.A and Remark 4.2.B have been added.}

\SGAInlineRefTarget{sga3:ii:remark:4-2-b}\subsection*{\SGARefLink{sga3:ii:remark:4-2-b}{Remark 4.2.B}.}
\textit{a)} If \(L'\) is another \(S\)-group on which \(X\) acts by group
automorphisms, then
\[
 \underline Z^1_S(X,L\times_SL')
 \simeq \underline Z^1_S(X,L)\times_S\underline Z^1_S(X,L').
\]

\textit{b)} If \(L\) is an \(X\)-\(\mathcal O_S\)-module, then
\(\underline Z^1_S(X,L)\) is the kernel of the differential
\[
 \partial:\underline{\operatorname{Hom}}_S(X,L)
 \longrightarrow\underline{\operatorname{Hom}}_S(X^2,L)
\]
defined in \SGARefLink{sga3:i:subsection:5-1}{I~5.1}. In particular, \(\underline Z^1_S(X,L)\) is then an
\(\mathcal O_S\)-module.

\subsection*{\SGAInlineRefTarget{sga3:ii:subsection:4-2}4.2.}
Let \(u:X\to Y\) be a morphism of \(S\)-groups. We describe
\(L^u_{\underline{\operatorname{Hom}}_{S\text{-gr.}}(X,Y)/S}
(\mathcal M)\). First, \SGARefLink{sga3:ii:corollary:3-11-3}{3.11.3} gives an isomorphism of \(S\)-functors,
functorial in \(\mathcal M\),
\begin{equation}
 L^u_{\underline{\operatorname{Hom}}_S(X,Y)/S}(\mathcal M)
 \xrightarrow{\sim}
 \underline{\operatorname{Hom}}_{Y/S}(X,T_{Y/S}(\mathcal M)).
 \tag{\(\dagger\)}\label{sga3:ii:subsection:4-2:equation:dagger}
\end{equation}
Since \(Y\) is an \(S\)-group,
\[
 T_{Y/S}(\mathcal M)
 =\underline{\operatorname{Lie}}(Y/S,\mathcal M)\cdot Y
 =\underline{\operatorname{Lie}}(Y/S,\mathcal M)_Y,
\]
and one obtains an isomorphism, functorial in \(\mathcal M\),
\[
 L^u_{\underline{\operatorname{Hom}}_S(X,Y)/S}(\mathcal M)
 \xrightarrow{\sim}
 \underline{\operatorname{Hom}}_S
   (X,\underline{\operatorname{Lie}}(Y/S,\mathcal M)).
\]

For \(S'\to S\), let \(u':X'\to Y'\) be the base change of \(u\). The
\(S\)-functor\footnote{Editor's note: this is \SGARefLink{sga3:ii:definition:4-2-0-1}{Definition 4.2.0.1} applied
with \(Z=Y\), to the \(Y\)-groups \(u:X\to Y\) and
\(T_{Y/S}(\mathcal M)\to Y\).}
\[
 \underline{\operatorname{Hom}}_{(Y/S)\text{-gr.}}
   \bigl(X,\underline{\operatorname{Lie}}(Y/S,\mathcal M)\cdot Y\bigr)
\]
is defined by
\[
\begin{aligned}
 &\underline{\operatorname{Hom}}_{(Y/S)\text{-gr.}}
 \bigl(X,\underline{\operatorname{Lie}}(Y/S,\mathcal M)\cdot Y\bigr)(S')\\
 &\quad=
 \operatorname{Hom}_{Y'\text{-gr.}}
 \bigl(X',(\underline{\operatorname{Lie}}(Y/S,\mathcal M)\cdot Y)_{S'}\bigr)\\
 &\quad=
 \operatorname{Hom}_{Y'\text{-gr.}}
 \bigl(X',\underline{\operatorname{Lie}}(Y'/S',\mathcal M)\cdot Y'\bigr).
\end{aligned}
\]
Thus \((\dagger)\) induces an isomorphism
\begin{equation}
 L^u_{\underline{\operatorname{Hom}}_{S\text{-gr.}}(X,Y)/S}
   (\mathcal M)
 \xrightarrow{\sim}
 \underline{\operatorname{Hom}}_{(Y/S)\text{-gr.}}
 \bigl(X,\underline{\operatorname{Lie}}(Y/S,\mathcal M)\cdot Y\bigr).
 \tag{\(\dagger'\)}\label{sga3:ii:subsection:4-2:equation:dagger-prime}
\end{equation}

On the other hand, the composite

\SGAThreeNativeDiagram[0.50\textwidth]
  {native_diagrams/exp2/exp2_p027_adjoint_action_chain.tex}
  {The action of \(X\) on
  \(\underline{\operatorname{Lie}}(Y/S,\mathcal M)\).}
  {fig:exp2-adjoint-action-chain}
  {Polo--Gille Exp2-14oct24.pdf, component-local p. 27; combined-reader p. 85; printed p. 75}

defines an action of \(X\) on
\(L=\underline{\operatorname{Lie}}(Y/S,\mathcal M)\) by group
automorphisms.

% Source page: Exp2 component-local p. 28; combined-reader p. 86; printed p. 76.
Let
\(\Phi\in\underline{\operatorname{Hom}}_{Y/S}
(X,\underline{\operatorname{Lie}}(Y/S,\mathcal M)\cdot Y)\). For
\(S''\to S'\to S\) and \(x\in X(S'')\), there is a unique expression
\[
 \Phi(S')(x)=\varphi(S')(x)\cdot u'(x),
 \qquad
 \varphi(S')(x)\in
 \underline{\operatorname{Lie}}(Y'/S',\mathcal M)(S'').
\]
This determines an element
\(\varphi\in\underline{\operatorname{Hom}}_S
(X,\underline{\operatorname{Lie}}(Y/S,\mathcal M))\). The map
\(\Phi(S')\) is a group morphism if and only if, for all
\(x_1,x_2\in X(S'')\),
\[
\begin{aligned}
 \varphi(S')(x_1x_2)
 &=\varphi(S')(x_1)
   \bigl(u(x_1)\varphi(S')(x_2)u(x_1)^{-1}\bigr)\\
 &=\varphi(S')(x_1)
   \bigl(x_1\cdot\varphi(S')(x_2)\bigr),
\end{aligned}
\]
that is, precisely when
\(\varphi\in\underline Z^1_S
(X,\underline{\operatorname{Lie}}(Y/S,\mathcal M))\).

\subsection*{Proposition 4.2.}
Let \(u:X\to Y\) be a morphism of \(S\)-groups. There is an isomorphism
of \(S\)-functors, functorial in \(\mathcal M\),
\[
 L^u_{\underline{\operatorname{Hom}}_{S\text{-gr.}}(X,Y)/S}
   (\mathcal M)
 \xrightarrow{\sim}
 \underline Z^1_S
   \bigl(X,\underline{\operatorname{Lie}}(Y/S,\mathcal M)\bigr).
\]

Suppose in addition that \(Y/S\) satisfies \(E\). It follows from
\SGARefLink{sga3:ii:corollary:4-2-0-3}{4.2.0.3}, just as in the proof of \SGARefLink{sga3:ii:corollary:3-11-1}{3.11.1}, that
\(\underline{\operatorname{Hom}}_{S\text{-gr.}}(X,Y)/S\) satisfies
\(E\). Thus (cf. \SGARefLink{sga3:ii:subsection:3-5-1}{3.5.1})
\[
\begin{aligned}
 L^u_{\underline{\operatorname{Hom}}_{S\text{-gr.}}(X,Y)/S}
 (\mathcal M\oplus\mathcal N)
 \simeq{}&
 L^u_{\underline{\operatorname{Hom}}_{S\text{-gr.}}(X,Y)/S}
 (\mathcal M)\\
 &\times_S
 L^u_{\underline{\operatorname{Hom}}_{S\text{-gr.}}(X,Y)/S}
 (\mathcal N).
\end{aligned}
\]
This also follows directly from \SGARefLink{sga3:ii:subsection:4-2}{Proposition~4.2} and \SGARefLink{sga3:ii:remark:4-2-b}{Remark~4.2.B}\(a\).
Consequently, both sides of the isomorphism in \SGARefLink{sga3:ii:subsection:4-2}{Proposition~4.2} carry
\(\mathcal O_S\)-module structures obtained from functoriality in
\(\mathcal M\), and that isomorphism respects them.

\SGAInlineRefTarget{sga3:ii:proposition:4-2-bis}\subsection*{\SGARefLink{sga3:ii:proposition:4-2-bis}{Proposition 4.2 bis}.}
Let \(u:X\to Y\) be a morphism of \(S\)-groups, and suppose that \(Y/S\)
satisfies \(E\). There is an isomorphism of \(\mathcal O_S\)-modules,
functorial in \(\mathcal M\),
\[
 L^u_{\underline{\operatorname{Hom}}_{S\text{-gr.}}(X,Y)/S}
   (\mathcal M)
 \xrightarrow{\sim}
 \underline Z^1_S
   \bigl(X,\underline{\operatorname{Lie}}(Y/S,\mathcal M)\bigr).
\]
\footnote{Editor's note: the original has been expanded from this point;
in particular, Proposition 4.2 bis has been added.}

When \(Y/S\) satisfies \(E\), \SGARefLink{sga3:ii:corollary:3-13-1}{3.13.1} also gives, just as in the proof of
\SGARefLink{sga3:ii:corollary:3-13-2}{3.13.2}, an isomorphism functorial in \(\mathcal M\), for every
\(u\in\underline{\operatorname{Isom}}_{S\text{-gr.}}(X,Y)\),
\begin{equation}
 L^u_{\underline{\operatorname{Isom}}_{S\text{-gr.}}(X,Y)/S}
   (\mathcal M)
 \xrightarrow{\sim}
 L^u_{\underline{\operatorname{Hom}}_{S\text{-gr.}}(X,Y)/S}
   (\mathcal M).
 \tag{*}\label{sga3:ii:subsection:4-2:equation:star}
\end{equation}
This yields the following two corollaries.

\SGAInlineRefTarget{sga3:ii:corollary:4-2-1}\subsection*{\SGARefLink{sga3:ii:corollary:4-2-1}{Corollary 4.2.1}.}
Let \(u:X\to Y\) be a morphism of \(S\)-groups. If \(Y/S\) satisfies
\(E\), there is an isomorphism of \(\mathcal O_S\)-modules, functorial in
\(\mathcal M\),
\[
 L^u_{\underline{\operatorname{Isom}}_{S\text{-gr.}}(X,Y)/S}
   (\mathcal M)
 \xrightarrow{\sim}
 \underline Z^1_S
   \bigl(X,\underline{\operatorname{Lie}}(Y/S,\mathcal M)\bigr).
\]

% Source page: Exp2 component-local p. 29; combined-reader p. 87; printed p. 77.
\SGAInlineRefTarget{sga3:ii:corollary:4-2-2}\subsection*{\SGARefLink{sga3:ii:corollary:4-2-2}{Corollary 4.2.2}.}
Let \(X\) be an \(S\)-group. If \(X/S\) satisfies \(E\), there is an
isomorphism of \(\mathcal O_S\)-modules, functorial in \(\mathcal M\),
\[
 \underline{\operatorname{Lie}}
   \bigl(\underline{\operatorname{Aut}}_{S\text{-gr.}}(X)/S,
         \mathcal M\bigr)
 \xrightarrow{\sim}
 \underline Z^1_S
   \bigl(X,\underline{\operatorname{Lie}}(X/S,\mathcal M)\bigr).
\]

If \(Y\) is commutative, its adjoint action on
\(L=\underline{\operatorname{Lie}}(Y/S,\mathcal M)\) is trivial, so
\(\underline Z^1_S(X,L)=
\underline{\operatorname{Hom}}_{S\text{-gr.}}(X,L)\).
\footnote{Editor's note: the following sentence has been added.}

\SGAInlineRefTarget{sga3:ii:corollary:4-2-3}\subsection*{\SGARefLink{sga3:ii:corollary:4-2-3}{Corollary 4.2.3}.}
Let \(Y\) be a commutative \(S\)-group. There is an isomorphism of
\(S\)-functors, functorial in \(\mathcal M\),
\[
 L^u_{\underline{\operatorname{Hom}}_{S\text{-gr.}}(X,Y)/S}
   (\mathcal M)
 \xrightarrow{\sim}
 \underline{\operatorname{Hom}}_{S\text{-gr.}}
   \bigl(X,\underline{\operatorname{Lie}}(Y/S,\mathcal M)\bigr).
\]

\subsection*{\SGAInlineRefTarget{sga3:ii:subsection:4-3}4.3.}
We now consider the case in which \(X\) and \(Y\) are
\(\mathcal O_S\)-modules. Recall (cf. \SGARefLink{sga3:ii:definition:4-2-0-4}{4.2.0.4}) that
\(T'_{Y/S}(\mathcal M)\), respectively
\(\underline{\operatorname{Lie}}'(Y/S,\mathcal M)\), denotes the
functor \(T_{Y/S}(\mathcal M)\), respectively
\(\underline{\operatorname{Lie}}(Y/S,\mathcal M)\), equipped with the
\(\mathcal O_S\)-module structure induced by that of \(Y\).

When \(Y/S\) satisfies \(E\), we shall always write
\(\underline{\operatorname{Lie}}(Y/S,\mathcal M)\) for this functor
equipped instead with the \(\mathcal O_S\)-module structure defined for
every functor satisfying \(E\). By \SGARefLink{sga3:ii:proposition:3-9}{3.9}, the underlying abelian-group
structures of
\(\underline{\operatorname{Lie}}(Y/S,\mathcal M)\) and
\(\underline{\operatorname{Lie}}'(Y/S,\mathcal M)\) coincide, but their
module structures need not coincide (see the counterexample in \SGARefLink{sga3:ii:subsection:6-3}{6.3}). For
\(S'\to S\), \(a\in\mathcal O(S')\), and
\(m\in\underline{\operatorname{Lie}}(Y/S,\mathcal M)(S')\), we write
\(a\mathbin{\cdot'}m\), respectively \(a\cdot m\), for the action coming
from \(Y\), respectively the action obtained from condition \(E\). We use
the same convention for the two actions on \(T_{Y/S}(\mathcal M)\).

There is an isomorphism of \(\mathcal O_S\)-modules
\[
 T'_{Y/S}(\mathcal M)
 \simeq \underline{\operatorname{Lie}}'(Y/S,\mathcal M)\oplus Y.
\]
Exactly as for \SGARefLink{sga3:ii:subsection:4-2}{Propositions~4.2} and \SGARefLink{sga3:ii:proposition:4-2-bis}{4.2 bis}, one therefore obtains:

\subsection*{Proposition 4.3.}
Let \(u:X\to Y\) be a morphism of \(\mathcal O_S\)-modules. There is an
isomorphism of \(S\)-functors, functorial in \(\mathcal M\),
\begin{equation}
 L^u_{\underline{\operatorname{Hom}}_{\mathcal O_S\text{-mod.}}(X,Y)/S}
   (\mathcal M)
 \xrightarrow{\sim}
 \underline{\operatorname{Hom}}_{\mathcal O_S\text{-mod.}}
   \bigl(X,\underline{\operatorname{Lie}}'(Y/S,\mathcal M)\bigr).
 \tag{*}\label{sga3:ii:subsection:4-3:equation:hom-lie-isomorphism-star}
\end{equation}
If \(Y/S\) satisfies \(E\), then
\(\underline{\operatorname{Hom}}_{\mathcal O_S\text{-mod.}}(X,Y)/S\)
satisfies \(E\), and \((*)\) is an isomorphism of
\(\mathcal O_S\)-modules when both sides carry the structures obtained
from functoriality in \(\mathcal M\).
\footnote{Editor's note: the following sentence has been added.}
\footnote{Editor's note: on the right-hand side, the scalar action is
\((a\varphi)(x)=a\varphi(x)\), where \(a\varphi(x)\) denotes the action
of \(a\) on \(\varphi(x)\in
\underline{\operatorname{Lie}}(Y/S,\mathcal M)\). This differs from the
action \((a\mathbin{\cdot'}\varphi)(x)
=a\mathbin{\cdot'}\varphi(x)=\varphi(ax)\), but the two agree when
\(\underline{\operatorname{Lie}}(Y/S,\mathcal M)
=\underline{\operatorname{Lie}}'(Y/S,\mathcal M)\).}

\SGAInlineRefTarget{sga3:ii:remark:4-3-bis}\subsection*{\SGARefLink{sga3:ii:remark:4-3-bis}{Remark 4.3 bis}.}
Let \(u:X\to Y\) be a morphism of
\(\mathcal O_S\)-modules.\footnote{Editor's note: this remark has been
added; it will be useful in \SGARefLink{sga3:ii:scholium:4-5-1}{4.5.1}.} Denote by
\(\tau_u\) the map that sends each \(\mathcal O_S\)-linear morphism
\(\varphi:X\to\underline{\operatorname{Lie}}'(Y/S,\mathcal M)\) to
\[
 u\oplus\varphi:X\longrightarrow
 T'_{Y/S}(\mathcal M)
 =Y\oplus\underline{\operatorname{Lie}}'(Y/S,\mathcal M).
 \]

% Source page: Exp2 component-local p. 30; combined-reader p. 88; printed p. 78.
The isomorphism of \SGARefLink{sga3:ii:subsection:4-3}{4.3} belongs to the following commutative diagram,
functorial in \(\mathcal M\).

\SGAThreeNativeDiagram[0.70\textwidth]
  {native_diagrams/exp2/exp2_p030_tau_u_square.tex}
  {Compatibility of the isomorphism of \SGARefLink{sga3:ii:subsection:4-3}{4.3} with \(\tau_u\).}
  {fig:exp2-tau-u-square}
  {Polo--Gille Exp2-14oct24.pdf, component-local p. 30; combined-reader p. 88; printed p. 78}

Moreover, when \(Y/S\) satisfies \(E\), \SGARefLink{sga3:ii:corollary:3-13-1}{3.13.1} gives, just as in the proof
of \SGARefLink{sga3:ii:corollary:3-13-2}{3.13.2}, for every
\(u\in\underline{\operatorname{Isom}}_{\mathcal O_S\text{-mod.}}(X,Y)\),
\begin{equation}
 L^u_{\underline{\operatorname{Isom}}_{\mathcal O_S\text{-mod.}}(X,Y)/S}
   (\mathcal M)
 =
 L^u_{\underline{\operatorname{Hom}}_{\mathcal O_S\text{-mod.}}(X,Y)/S}
   (\mathcal M).
 \tag{*}\label{sga3:ii:subsection:4-3:equation:isom-hom-lie-equality-star}
\end{equation}

\SGAInlineRefTarget{sga3:ii:corollary:4-3-1}\subsection*{\SGARefLink{sga3:ii:corollary:4-3-1}{Corollary 4.3.1}.}
Let \(X\) be an \(\mathcal O_S\)-module satisfying \(E\) over \(S\). There
is an isomorphism, functorial in \(\mathcal M\),
\[
 \underline{\operatorname{Lie}}
  \bigl(\underline{\operatorname{Aut}}_{\mathcal O_S\text{-mod.}}(X)/S,
        \mathcal M\bigr)
 \xrightarrow{\sim}
 \underline{\operatorname{Hom}}_{\mathcal O_S\text{-mod.}}
  \bigl(X,\underline{\operatorname{Lie}}'(X/S,\mathcal M)\bigr),
\]
and this isomorphism respects the \(\mathcal O_S\)-module structures
obtained from functoriality in \(\mathcal M\). In particular,
\(\underline{\operatorname{Aut}}_{\mathcal O_S\text{-mod.}}(X)/S\)
satisfies \(E\).
\footnote{Editor's note: compare editor's note (61). The final sentence
and its proof have also been added.}

\paragraph{Proof.}
The first assertion follows from \((*)\) and \SGARefLink{sga3:ii:subsection:4-3}{4.3}. For the second, since
\(X/S\) satisfies \(E\), there are isomorphisms of
\(\mathcal O_S\)-modules
\[
 \underline{\operatorname{Lie}}'(X/S,\mathcal M\oplus\mathcal N)
 \simeq
 \underline{\operatorname{Lie}}'(X/S,\mathcal M)
 \times_S
 \underline{\operatorname{Lie}}'(X/S,\mathcal N).
\]
Consequently,
\[
\begin{aligned}
 &\underline{\operatorname{Lie}}
  \bigl(\underline{\operatorname{Aut}}_{\mathcal O_S\text{-mod.}}(X)/S,
  \mathcal M\oplus\mathcal N\bigr)\\
 &\quad\simeq
 \underline{\operatorname{Lie}}
  \bigl(\underline{\operatorname{Aut}}_{\mathcal O_S\text{-mod.}}(X)/S,
  \mathcal M\bigr)
 \times_S
 \underline{\operatorname{Lie}}
  \bigl(\underline{\operatorname{Aut}}_{\mathcal O_S\text{-mod.}}(X)/S,
  \mathcal N\bigr).
\end{aligned}
\]
\SGARefLink{sga3:ii:remark:4-1-1-2}{Remark~4.1.1.2}\(a\) now proves the assertion.

\subsection*{\SGAInlineRefTarget{sga3:ii:subsection:4-3-2}4.3.2.}
Before continuing in this direction, let us examine more closely the
relations among \(Y\), \(\underline{\operatorname{Lie}}(Y/S)\), and
\(\underline{\operatorname{Lie}}'(Y/S)\). First,
\begin{equation}
 \underline{\operatorname{Lie}}(\mathcal O_S/S,\mathcal M)
 =\underline{\operatorname{Lie}}'(\mathcal O_S/S,\mathcal M)
 =\mathbf W(\mathcal M),
 \tag{1}\label{sga3:ii:subsection:4-3-2:equation:1}
\end{equation}
where \(\mathbf W(\mathcal M)\) is defined in \SGARefLink{sga3:i:subsection:4-6}{I~4.6}. There is therefore a
canonical isomorphism
\begin{equation}
 d:\mathcal O_S\xrightarrow{\sim}
 \underline{\operatorname{Lie}}(\mathcal O_S/S).
 \tag{2}\label{sga3:ii:subsection:4-3-2:equation:2}
\end{equation}

Let \(F\) be an \(\mathcal O_S\)-module. For every
\(S_2\to S_1\to S\), there is a dihomomorphism
\begin{equation}
 \left\{
 \begin{aligned}
   F(S_1)&\longrightarrow F(S_2),\\
   \mathcal O(S_1)&\longrightarrow\mathcal O(S_2),
 \end{aligned}
 \right.
 \tag{3}\label{sga3:ii:subsection:4-3-2:equation:3}
\end{equation}
and hence a morphism of \(\mathcal O(S_2)\)-modules
\[
 F(S_1)\otimes_{\mathcal O(S_1)}\mathcal O(S_2)
 \longrightarrow F(S_2).
\]
\footnote{Editor's note: \(S'\to S\) has been replaced by
\(S_2\to S_1\to S\), since both \(S_2\) and \(S_1\) must vary below.
The original has also been expanded in what follows.}

% Source page: Exp2 component-local p. 31; combined-reader p. 89; printed p. 79.
In particular, take \(S_1=S'\) and \(S_2=I_{S'}(\mathcal M)\). One
obtains morphisms of \(\mathcal O(S')\)-modules, functorial in
\(\mathcal M\),
\[
 F(S')\otimes_{\mathcal O(S')}
 T_{\mathcal O_S/S}(\mathcal M)(S')
 \longrightarrow T'_{F/S}(\mathcal M)(S').
\]
As \(S'\) varies, these give morphisms of \(\mathcal O_S\)-modules
\begin{equation}
 F\otimes_{\mathcal O_S}T_{\mathcal O_S/S}(\mathcal M)
 \longrightarrow T'_{F/S}(\mathcal M),
 \tag{4}\label{sga3:ii:subsection:4-3-2:equation:4}
\end{equation}
functorial in \(\mathcal M\). Compatibility with the tangent-bundle
projections onto their bases then yields morphisms
\begin{equation}
 F\otimes_{\mathcal O_S}
 \underline{\operatorname{Lie}}(\mathcal O_S/S,\mathcal M)
 \longrightarrow
 \underline{\operatorname{Lie}}'(F/S,\mathcal M),
 \tag{5}\label{sga3:ii:subsection:4-3-2:equation:5}
\end{equation}
and the following diagram commutes.

\SGAThreeNativeDiagram[0.78\textwidth]
  {native_diagrams/exp2/exp2_p031_tangent_exact_rows.tex}
  {The morphisms (4) and (5) in the split tangent sequences.}
  {fig:exp2-tangent-exact-rows}
  {Polo--Gille Exp2-14oct24.pdf, component-local p. 31; combined-reader p. 89; printed p. 79}

The morphisms (5) may be regarded as morphisms of abelian \(S\)-groups
\begin{equation}
 F\otimes_{\mathcal O_S}
 \underline{\operatorname{Lie}}(\mathcal O_S/S,\mathcal M)
 \longrightarrow
 \underline{\operatorname{Lie}}(F/S,\mathcal M).
 \tag{6}\label{sga3:ii:subsection:4-3-2:equation:6}
\end{equation}
Tensoring \(F\) with the isomorphism
\(d:\mathcal O_S\xrightarrow{\sim}
\underline{\operatorname{Lie}}(\mathcal O_S/S)\) gives, when
\(\mathcal M=\mathcal O_S\), a morphism of abelian \(S\)-groups
\begin{equation}
 F\xrightarrow{\sim}
 F\otimes_{\mathcal O_S}
 \underline{\operatorname{Lie}}(\mathcal O_S/S)
 \longrightarrow \underline{\operatorname{Lie}}(F/S),
 \tag{7}\label{sga3:ii:subsection:4-3-2:equation:7}
\end{equation}
also denoted \(d:F\to\underline{\operatorname{Lie}}(F/S)\).

\SGAInlineRefTarget{sga3:ii:remark:4-3-3}\subsection*{\SGARefLink{sga3:ii:remark:4-3-3}{Remark 4.3.3}.}
When \(F/S\) satisfies \(E\), the morphisms \SGARefLink{sga3:ii:section:6}{(6)} and (7) need not be
\(\mathcal O_S\)-linear if both sides are equipped with the module
structures obtained from \(\mathcal M\) by condition \(E\).

For example, let \(k\) be a field of characteristic \(p>0\), let
\(S=\operatorname{Spec}(k)\), and let \(F\) be the
\(\mathcal O_S\)-module that assigns to every \(S\)-scheme \(T\) the
set \(F(T)=\Gamma(T,\mathcal O_T)\), with scalars acting through the
\(p\)-th power:
\[
 r\cdot f=r^pf
 \qquad(r\in\mathcal O(T),\ f\in F(T)).
\]
As an \(S\)-functor in groups, \(F\simeq\mathbf G_{a,S}\). Hence \(F\)
satisfies \(E\) and
\(\underline{\operatorname{Lie}}(F/S)\simeq
\underline{\operatorname{Lie}}(\mathbf G_{a,S}/S)\simeq\mathcal O_S\).
For every \(T\to S\), the canonical morphism
\(d:F(T)\to\mathcal O(T)\) is the identity: it respects the abelian-group
structure but not the \(\mathcal O_S\)-module structure.
\footnote{Editor's note: the original asserted that, when \(F/S\)
satisfies \(E\), (6), and hence (7), are morphisms of
\(\mathcal O_S\)-modules. The counterexample in \SGARefLink{sga3:ii:subsection:6-3}{6.3} contradicts this.
That assertion has been removed, the preceding example inserted, and
Definition 4.4 modified accordingly. This does not affect what follows.}

% Source page: Exp2 component-local p. 32; combined-reader p. 90; printed p. 80.
\SGAInlineRefTarget{sga3:ii:remark:4-3-4}\subsection*{\SGARefLink{sga3:ii:remark:4-3-4}{Remark 4.3.4}.}
The morphisms (4) and (5) can be made explicit as follows. The morphism
\[
 \Theta:F\otimes_{\mathcal O_S}T_{\mathcal O_S/S}(\mathcal M)
 \longrightarrow T'_{F/S}(\mathcal M)
 =\underline{\operatorname{Hom}}_S(I_S(\mathcal M),F)
\]
is defined, for \(S'\to S\),
\(\alpha\in\mathcal O(I_{S'}(\mathcal M))\), and \(f:S'\to F\), by
\[
 \Theta(f\otimes\alpha)
 =\alpha\cdot(\tau_0\circ f)
 =\alpha\cdot(f\circ\rho),
\]
where \(\tau_0:F\to T'_{F/S}(\mathcal M)\) is the zero section and
\(\rho:I_{S'}(\mathcal M)\to S'\) is the structural morphism.

The morphism \(\Theta\) induces
\[
 \theta:F\otimes_{\mathcal O_S}
 \underline{\operatorname{Lie}}(\mathcal O_S/S,\mathcal M)
 \longrightarrow
 \underline{\operatorname{Lie}}'(F/S,\mathcal M).
\]
This follows from the functoriality in \(\mathcal M\) noted after (4),
and can also be seen directly. Indeed,
\[
 \underline{\operatorname{Lie}}'(F/S,\mathcal M)(S')
 =\{\,\varphi\in\operatorname{Hom}_S(I_{S'}(\mathcal M),F)
       \mid \varphi\circ\varepsilon_{\mathcal M}=e\,\},
\]
where \(e:S'\to F\) is the unit section. On the other hand,
\(\underline{\operatorname{Lie}}(\mathcal O_S/S,\mathcal M)(S')\)
equals \(\Gamma(S',\mathcal M)\), the kernel of the augmentation
\[
 \eta:\mathcal O(I_{S'}(\mathcal M))\longrightarrow\mathcal O(S').
\]
It remains to see that if \(f\in F(S')\) and
\(\alpha\in\Gamma(S',\mathcal M)\), then
\(\alpha\cdot(f\circ\rho)\circ\varepsilon_{\mathcal M}=e\).

Apply the dihomomorphism (3) to the \(S\)-morphism
\(\varepsilon_{\mathcal M}:S'\to I_{S'}(\mathcal M)\). The following
diagram commutes.

\SGAThreeNativeDiagram[0.43\textwidth]
  {native_diagrams/exp2/exp2_p032_epsilon_action_square.tex}
  {Compatibility of scalar action with
  \(F(\varepsilon_{\mathcal M})\).}
  {fig:exp2-epsilon-action-square}
  {Polo--Gille Exp2-14oct24.pdf, component-local p. 32; combined-reader p. 90; printed p. 80}

Thus, for every \(\varphi:I_{S'}(\mathcal M)\to F\),
\[
 (\alpha\cdot\varphi)\circ\varepsilon_{\mathcal M}
 =\eta(\alpha)\cdot(\varphi\circ\varepsilon_{\mathcal M}),
\]
and therefore
\((\alpha\cdot\varphi)\circ\varepsilon_{\mathcal M}=e\) when
\(\eta(\alpha)=0\).

In particular, taking \(\mathcal M=t\mathcal O_S\), one has
\(\underline{\operatorname{Lie}}(\mathcal O_S/S)=t\mathcal O_S\), and
the morphism \(F\to\underline{\operatorname{Lie}}'(F/S)\) is
\(f\mapsto t\cdot(f\circ\rho)\).

\footnote{Editor's note: this remark has been added; it will be useful in
\SGARefLink{sga3:ii:lemma:4-6-2}{4.6.2}.}

\SGAInlineRefTarget{sga3:ii:remark:4-3-5}\subsection*{\SGARefLink{sga3:ii:remark:4-3-5}{Remark 4.3.5}.}
Let \(F\) be an \(\mathcal O_S\)-module, set
\[
 E=\underline{\operatorname{End}}_{\mathcal O_S\text{-mod.}}(F),
\]
and denote by \(d_F\) and \(d_E\) the \(\mathcal O_S\)-linear morphisms
given by \SGARefLink{sga3:ii:subsection:4-3-2:equation:5}{4.3.2(5)}:
\[
\begin{aligned}
 d_F:
 F\otimes_{\mathcal O_S}
 \underline{\operatorname{Lie}}(\mathcal O_S/S,\mathcal M)
 &\longrightarrow
 \underline{\operatorname{Lie}}'(F/S,\mathcal M),\\
 d_E:
 E\otimes_{\mathcal O_S}
 \underline{\operatorname{Lie}}(\mathcal O_S/S,\mathcal M)
 &\longrightarrow
 \underline{\operatorname{Lie}}'(E/S,\mathcal M).
\end{aligned}
\]
\SGARefLink{sga3:ii:remark:4-3-4}{Remark~4.3.4} gives the following commutative diagram of
\(\mathcal O_S\)-linear morphisms.

\SGAThreeNativeDiagram[0.76\textwidth]
  {native_diagrams/exp2/exp2_p032_differentials_square.tex}
  {Compatibility of \(d_E\) and \(d_F\).}
  {fig:exp2-differentials-square}
  {Polo--Gille Exp2-14oct24.pdf, component-local p. 32; combined-reader p. 90; printed p. 80}

The right vertical arrow is the isomorphism \((*)\) of \SGARefLink{sga3:ii:subsection:4-3}{4.3}. Thus, by
that result, if \(F/S\) satisfies \(E\), then \(E/S\) satisfies \(E\), and
\((*)\) is also an isomorphism of \(\mathcal O_S\)-modules when the
right-hand terms are equipped with the module structures obtained from
\(E\).
\footnote{Editor's note: this remark has been added; it will be useful in
\SGARefLink{sga3:ii:proposition:4-5}{4.5} and \SGARefLink{sga3:ii:theorem:4-7}{4.7}.}

% Source page: Exp2 component-local p. 33; combined-reader p. 91; printed p. 81.
\SGAInlineRefTarget{sga3:ii:remark:4-4-0}\subsection*{\SGARefLink{sga3:ii:remark:4-4-0}{Remark 4.4.0}.}
In \SGARefLink{sga3:ii:subsection:4-3-2}{4.3.2}, the morphisms (4) are isomorphisms if and only if the
morphisms \SGARefLink{sga3:ii:section:5}{(5)} are. If these equivalent conditions hold, then \(F/S\)
satisfies \(E\). Indeed, it is enough to verify that
\[
 \underline{\operatorname{Lie}}'(F/S,\mathcal M\oplus\mathcal N)
 \simeq
 \underline{\operatorname{Lie}}'(F/S,\mathcal M)
 \times_S
 \underline{\operatorname{Lie}}'(F/S,\mathcal N).
\]
Under the hypothesis, the horizontal arrows in the following commutative
diagram are isomorphisms.

\SGAThreeNativeDiagram[0.84\textwidth]
  {native_diagrams/exp2/exp2_p033_good_module_square.tex}
  {The direct-sum comparison used to verify condition \(E\).}
  {fig:exp2-good-module-square}
  {Polo--Gille Exp2-14oct24.pdf, component-local p. 33; combined-reader p. 91; printed p. 81}

The second vertical arrow is therefore an isomorphism as well, which is
exactly condition \(E\) for \(F/S\).
\footnote{Editor's note: in view of the changes made in Definition 4.4
(cf. editor's note (65)), this remark has been inserted here; in the
original it occurred in the proof of \SGARefLink{sga3:ii:corollary:4-4-1}{4.4.1}.}

\SGAInlineRefTarget{sga3:ii:definition:4-4}\subsection*{\SGARefLink{sga3:ii:definition:4-4}{Definition 4.4}.}
An \(\mathcal O_S\)-module \(F\) is called \emph{good} if the morphisms
\[
 F\otimes_{\mathcal O_S}T_{\mathcal O_S/S}(\mathcal M)
 \longrightarrow T_{F/S}(\mathcal M),
\]
or, equivalently,
\[
 F\otimes_{\mathcal O_S}
 \underline{\operatorname{Lie}}(\mathcal O_S/S,\mathcal M)
 \longrightarrow
 \underline{\operatorname{Lie}}(F/S,\mathcal M),
\]
are isomorphisms of abelian \(S\)-groups (so that \(F/S\) satisfies
(E)), and if moreover they respect the \(\mathcal O_S\)-module
structures obtained from condition \(E\).

\SGAInlineRefTarget{sga3:ii:corollary:4-4-1}\subsection*{\SGARefLink{sga3:ii:corollary:4-4-1}{Corollary 4.4.1}.}
Let \(F\) be an \(\mathcal O_S\)-module, and consider the following
conditions:\footnote{Editor's note: the original proved that (i) implies
(ii) and (iii); the implication (ii) implies (i) has been added.}
\begin{enumerate}
  \item[(i)] \(F\) is a good \(\mathcal O_S\)-module;
  \item[(ii)] \(F/S\) satisfies \(E\), and
    \(d:F\to\underline{\operatorname{Lie}}(F/S)\) is an isomorphism of
    \(\mathcal O_S\)-modules;
  \item[(iii)]
    \(\underline{\operatorname{Lie}}(F/S,\mathcal M)
      =\underline{\operatorname{Lie}}'(F/S,\mathcal M)\).
\end{enumerate}
Then
\[
  \textup{(i)}\Longleftrightarrow\textup{(ii)}
  \Longrightarrow\textup{(iii)}.
\]

\paragraph{Proof.}
The implication
\(\textup{(i)}\Rightarrow\textup{(ii)}\) follows from the definition. To
prove \(\textup{(ii)}\Rightarrow\textup{(i)}\), it remains to show that
the morphisms of
abelian \(S\)-groups, functorial in \(\mathcal M\),
\[
 F\otimes_{\mathcal O_S}
 \underline{\operatorname{Lie}}(\mathcal O_S/S,\mathcal M)
 \xrightarrow{\sim}
 \underline{\operatorname{Lie}}(F/S,\mathcal M)
\]
are isomorphisms of \(\mathcal O_S\)-modules. Since \(F/S\) satisfies
\(E\), both sides carry finite direct sums of copies of \(\mathcal O_S\)
to finite products of abelian \(S\)-groups. We are therefore reduced to
the case \(\mathcal M=\mathcal O_S\), which is precisely the hypothesis.

Finally, \(\textup{(i)}\Rightarrow\textup{(iii)}\) follows from the
definition and from the
fact that the isomorphisms
\[
 F\otimes_{\mathcal O_S}
 \underline{\operatorname{Lie}}(\mathcal O_S/S,\mathcal M)
 \xrightarrow{\sim}
 \underline{\operatorname{Lie}}'(F/S,\mathcal M)
\]
of \SGARefLink{sga3:ii:subsection:4-3-2:equation:5}{4.3.2(5)} are \(\mathcal O_S\)-linear.

\SGAInlineRefTarget{sga3:ii:examples:4-4-2}\SGAInlineRefTarget{sga3:ii:example:4-4-2}\subsection*{\SGARefLink{sga3:ii:example:4-4-2}{Examples 4.4.2}.}
For every quasi-coherent \(\mathcal O_S\)-module \(\mathcal E\), the
\(\mathcal O_S\)-modules \(\mathbf V(\mathcal E)\) and
\(\mathbf W(\mathcal E)\) defined in \SGARefLink{sga3:i:subsection:4-6}{I~4.6} are good.

% Source page: Exp2 component-local p. 34; combined-reader p. 92; printed p. 82.
Indeed, for every \(f:S'\to S\), the morphisms
\[
\begin{aligned}
 \mathbf V(\mathcal E)(S')\otimes_{\mathcal O(S')}
   \mathcal O(I_{S'}(\mathcal M))
 &\longrightarrow T_{\mathbf V(\mathcal E)/S}(\mathcal M)(S'),\\
 \mathbf W(\mathcal E)(S')\otimes_{\mathcal O(S')}
   \mathcal O(I_{S'}(\mathcal M))
 &\longrightarrow T_{\mathbf W(\mathcal E)/S}(\mathcal M)(S')
\end{aligned}
\]
correspond respectively to
\[
\begin{aligned}
 &\operatorname{Hom}_{\mathcal O_{S'}\text{-mod.}}
   \bigl(f^*(\mathcal E),\mathcal O_{S'}\bigr)
 \otimes_{\mathcal O(S')}
 \Gamma\bigl(S',D_{\mathcal O_{S'}}(\mathcal M)\bigr)\\
 &\hspace{5em}\longrightarrow
 \operatorname{Hom}_{\mathcal O_{S'}\text{-mod.}}
   \bigl(f^*(\mathcal E),D_{\mathcal O_{S'}}(\mathcal M)\bigr)
 \qquad\textup{(V)}
\\[-2pt]
 &\Gamma(S',f^*(\mathcal E))
 \otimes_{\mathcal O(S')}
 \Gamma\bigl(S',D_{\mathcal O_{S'}}(\mathcal M)\bigr)\\
 &\hspace{5em}\longrightarrow
 \Gamma\bigl(S',f^*(\mathcal E)
   \otimes_{\mathcal O_{S'}}D_{\mathcal O_{S'}}(\mathcal M)\bigr)
 \qquad\textup{(W)}.
\end{aligned}
\]
These are isomorphisms because, as an \(\mathcal O_{S'}\)-module,
\(D_{\mathcal O_{S'}}(\mathcal M)\) is isomorphic to a finite direct
sum of copies of \(\mathcal O_{S'}\).
\footnote{Editor's note: the preceding proof has been added.}

% End of the assigned scope. Proposition 4.5 begins immediately here in
% Exp2-14oct24.pdf, component-local p. 34 (combined-reader p. 92).

% END INCLUDED SGA3_English_current_buildable_source_20260823/source/inputs/II/tex/components/12_expose_II_s41_s44_lie_algebras.tex

% BEGIN INCLUDED SGA3_English_current_buildable_source_20260823/source/inputs/II/tex/components/13_expose_II_s45_s411_lie_structure.tex
% SGA 3, Expose II, Loop 1 English reconstruction.
% Scope: Proposition 4.5 through Corollary 4.11.8, stopping immediately
% before Section 5.
% Authority: Polo--Gille Exp2-14oct24.pdf, component-local pages 34--43
% (combined-reader pages 92--101; printed pages 82--91).
% OCR was used only as a locator. English wording was drafted independently
% against the French authority and compared with Jacob Reinhold's CC BY 4.0
% English Markdown lineage; that lineage is comparison material, not source
% authority. The PDF controls text, formulas, numbering, notes, and diagrams.

\setcounter{footnote}{70}

% Source page: Exp2 component-local p. 34; combined-reader p. 92; printed p. 82.
\medskip
\noindent\SGAInlineRefTarget{sga3:ii:proposition:4-5}\textbf{Proposition 4.5.}
\textit{Let \(F\) be a good \(\mathcal O_S\)-module. Then:}

\textit{(i) The \(S\)-group
\(\underline{\operatorname{Aut}}_{\mathcal O_S\text{-mod.}}(F)\)
satisfies condition \((E)\), and there is a functorial isomorphism}
\[
 \underline{\operatorname{Lie}}\!
   \left(\underline{\operatorname{Aut}}_{\mathcal O_S\text{-mod.}}(F)/S,
         \mathcal M\right)
 \xrightarrow{\sim}
 \underline{\operatorname{Hom}}_{\mathcal O_S\text{-mod.}}
   \left(F,\underline{\operatorname{Lie}}(F/S,\mathcal M)\right),
\]
\textit{which respects the \(\mathcal O_S\)-module structures obtained from
condition \((E)\). In particular, there is an isomorphism of
\(\mathcal O_S\)-modules}
\[
 \underline{\operatorname{Lie}}\!
   \left(\underline{\operatorname{Aut}}_{\mathcal O_S\text{-mod.}}(F)/S\right)
 \xrightarrow{\sim}
 \underline{\operatorname{End}}_{\mathcal O_S\text{-mod.}}(F).
\]

\textit{\SGAInlineRefTarget{sga3:ii:proposition:4-5:item:ii}(ii)\footnote{Editor's note: part \SGARefLink{sga3:ii:proposition:4-5:item:ii}{(ii)}, an immediate consequence
of what precedes, has been added.} Moreover,
\(\underline{\operatorname{End}}_{\mathcal O_S\text{-mod.}}(F)\)
is a good \(\mathcal O_S\)-module.}

Indeed, by \SGARefLink{sga3:ii:corollary:4-4-1}{4.4.1}, \(F/S\) satisfies \((E)\), and
\begin{equation*}
 \underline{\operatorname{Lie}}(F/S,\mathcal M)
 =\underline{\operatorname{Lie}}'(F/S,\mathcal M)
 \simeq
 F\otimes_{\mathcal O_S}
 \underline{\operatorname{Lie}}(\mathbf O_S/S,\mathcal M).
 \tag{1}\label{sga3:ii:proposition:4-5:equation:1}
\end{equation*}
Assertion (i) now follows from \SGARefLink{sga3:ii:corollary:4-3-1}{4.3.1}. Put
\(E=\underline{\operatorname{End}}_{\mathcal O_S\text{-mod.}}(F)\).
By (1) and \SGARefLink{sga3:ii:remark:4-3-5}{Remark 4.3.5}, one has the following commutative diagram of
morphisms of abelian \(S\)-groups.

\SGAThreeNativeDiagram[0.90\textwidth]
  {native_diagrams/exp2/exp2_p034_endomorphism_differential_square.tex}
  {The comparison diagram for the differential \(d_E\).}
  {fig:exp2-p45-endomorphism-differential}
  {Polo--Gille Exp2-14oct24.pdf, component-local p. 34; combined-reader p. 92; printed p. 82}

Here \(d_F\) and the vertical isomorphism \((*)\) are isomorphisms of
\(\mathcal O_S\)-modules; consequently, so is \(d_E\). This proves \SGARefLink{sga3:ii:proposition:4-5:item:ii}{(ii)}.

\medskip
\noindent\SGAInlineRefTarget{sga3:ii:scholium:4-5-1}\textbf{Scholium 4.5.1.}
\footnote{Editor's note: this scholium, implicit in the original, has been
added for use in \SGARefLink{sga3:ii:subsection:4-7-1}{4.7.1}. Here and below, \(t,t'\), and so forth denote
square-zero variables; the letter \(\varepsilon\) is reserved for the unit
section of a group.}
Put (cf. \SGARefLink{sga3:ii:definition:2-1}{2.1})
\(\mathcal O_{I_S}=\mathcal O_S\oplus t\mathcal O_S\), with \(t^2=0\),
and let \(F\) be a good \(\mathcal O_S\)-module. For every \(S'\to S\),
the morphism
\[
 F(S')\oplus tF(S')
 =F(S')\otimes_{\mathcal O(S')}\mathcal O(I_{S'})
 \longrightarrow
 F(I_{S'})
 =F(S')\oplus\underline{\operatorname{Lie}}(F/S)(S'),
\]
which is the identity on \(F(S')\), induces an isomorphism of
\(\mathcal O(S')\)-modules
\[
 tF(S')\xrightarrow{\sim}\underline{\operatorname{Lie}}(F/S)(S').
\]
As \(S'\) varies, this gives an isomorphism that we may write
\[
 \underline{\operatorname{Lie}}(F/S)\simeq tF.
\]

% Source page: Exp2 component-local p. 35; combined-reader p. 93; printed p. 83.
Thus, for every \(S'\to S\), \SGARefLink{sga3:ii:proposition:4-5}{Proposition 4.5} gives a commutative
diagram.

\SGAThreeNativeDiagram[0.91\textwidth]
  {native_diagrams/exp2/exp2_p035_dual_number_automorphism_square.tex}
  {Endomorphisms and automorphisms over the dual numbers.}
  {fig:exp2-p451-dual-number-automorphisms}
  {Polo--Gille Exp2-14oct24.pdf, component-local p. 35; combined-reader p. 93; printed p. 83}

It follows from 4.3 bis that every
\[
 X\in
 \underline{\operatorname{End}}_{\mathcal O_{S'}\text{-mod.}}(F_{S'})
\]
corresponds to
\(\mathrm{id}+tX\in
\underline{\operatorname{Aut}}_{\mathcal O_{I_{S'}}\text{-mod.}}
(F_{I_{S'}})\).

\medskip
\noindent\SGAInlineRefTarget{sga3:ii:definition:4-6}\textbf{Definition 4.6.}
An \(S\)-functor in groups \(G\) is called \emph{good} if \(G/S\)
satisfies condition \((E)\) and
\(\underline{\operatorname{Lie}}(G/S)\) is a good
\(\mathcal O_S\)-module.

\footnote{Editor's note: the following sentence has been added.}
Notice that if \(F\) is a good \(\mathcal O_S\)-module, then it is a
good \(S\)-group: indeed, \(F/S\) satisfies \((E)\), and
\(\underline{\operatorname{Lie}}(F/S)\simeq F\) is a good
\(\mathcal O_S\)-module.

\medskip
\noindent\SGAInlineRefTarget{sga3:ii:example:4-6-1}\textbf{Example 4.6.1.}
\textit{If \(G\) is representable, it is good. Indeed, \(G/S\)
satisfies \((E)\), and \(\underline{\operatorname{Lie}}(G/S)\) has
the form \(V(E)\), hence is good by \SGARefLink{sga3:ii:examples:4-4-2}{4.4.2}.}

\medskip
\noindent\SGAInlineRefTarget{sga3:ii:lemma:4-6-2}\textbf{Lemma 4.6.2.}
\footnote{Editor's note: this lemma has been added.}
\textit{Let \(G\) be an \(S\)-functor in groups such that \(G/S\)
satisfies \((E)\), and put
\(F=\underline{\operatorname{Lie}}(G/S)\). Then \(F\) satisfies
\((E)\), and the morphism of abelian groups}
\[
 d:F\longrightarrow\underline{\operatorname{Lie}}(F/S)
\]
\textit{respects the \(\mathcal O_S\)-module structures. Consequently,
\(G\) is good if and only if
\(F\to\underline{\operatorname{Lie}}(F/S)\) is bijective.}

\textit{Proof.}
Suppose that \(G/S\) satisfies \((E)\). Let
\(\mathcal M,\mathcal N\) be finite free \(\mathcal O_S\)-modules,
write
\[
 F(\mathcal N)=\underline{\operatorname{Lie}}(G/S,\mathcal N),
\]
and let \(e\) denote the unit section of \(G\). For every \(S'\to S\),
\[
 F(\mathcal N)(S')
 =
 \left\{
  g\in\operatorname{Hom}_S(I_{S'}(\mathcal N),G)
  \ \middle|\ g\circ\varepsilon_{\mathcal N}=e
 \right\}.
\]
Moreover,
\(\underline{\operatorname{Lie}}(F(\mathcal N)/S,\mathcal M)(S')\)
identifies with the set of morphisms
\[
\left\{
\begin{array}{@{}l@{\ }|@{\ }l@{}}
 \Phi\in
 \operatorname{Hom}_S
   (I_{S'}(\mathcal N)\times_{S'}I_{S'}(\mathcal M),G)
 &
 \begin{aligned}
 \Phi\circ
  (\varepsilon_{\mathcal N}\times
   \mathrm{id}_{I_{S'}(\mathcal M)})
   &=e\in G(I_{S'}(\mathcal M)),\\
 \Phi\circ
  (\mathrm{id}_{I_{S'}(\mathcal N)}\times
   \varepsilon_{\mathcal M})
   &=e\in G(I_{S'}(\mathcal N)).
 \end{aligned}
\end{array}
\right\}.
\]
This proves
\begin{equation*}
 \underline{\operatorname{Lie}}(F(\mathcal N)/S,\mathcal M)
 \simeq
 \underline{\operatorname{Lie}}(F(\mathcal M)/S,\mathcal N).
 \tag{1}\label{sga3:ii:lemma:4-6-2:equation:1}
\end{equation*}
Since \(G/S\) satisfies \((E)\), it follows that
\[
\begin{aligned}
 \underline{\operatorname{Lie}}
   (F(\mathcal N)/S,\mathcal M_1\oplus\mathcal M_2)
 &\simeq
 \underline{\operatorname{Lie}}
   (F(\mathcal M_1\oplus\mathcal M_2)/S,\mathcal N)\\
 &\simeq
 \underline{\operatorname{Lie}}
   ((F(\mathcal M_1)\times_SF(\mathcal M_2))/S,\mathcal N).
\end{aligned}
\]
By \SGARefLink{sga3:ii:proposition:3-8}{3.8}, the last term is isomorphic to
\[
\begin{aligned}
 &\underline{\operatorname{Lie}}(F(\mathcal M_1)/S,\mathcal N)
 \times_S
 \underline{\operatorname{Lie}}(F(\mathcal M_2)/S,\mathcal N)\\
 &\qquad\simeq
 \underline{\operatorname{Lie}}(F(\mathcal N)/S,\mathcal M_1)
 \times_S
 \underline{\operatorname{Lie}}(F(\mathcal N)/S,\mathcal M_2).
\end{aligned}
\]

% Source page: Exp2 component-local p. 36; combined-reader p. 94; printed p. 84.
Hence
\begin{equation*}
 \underline{\operatorname{Lie}}
   (F(\mathcal N)/S,\mathcal M_1\oplus\mathcal M_2)
 \simeq
 \underline{\operatorname{Lie}}(F(\mathcal N)/S,\mathcal M_1)
 \times_S
 \underline{\operatorname{Lie}}(F(\mathcal N)/S,\mathcal M_2),
 \tag{2}\label{sga3:ii:lemma:4-6-2:equation:2}
\end{equation*}
so \(F(\mathcal N)/S\) satisfies \((E)\), by Remark 4.1.1.2(a).

It remains to show that the morphism of abelian groups
\[
 d:F(\mathcal N)\longrightarrow
 \underline{\operatorname{Lie}}(F(\mathcal N)/S)
\]
respects the \(\mathcal O_S\)-module structures. Consider the free
\(\mathcal O_S\)-module \(\mathcal M=t\mathcal O_S\). Then
\[
 \mathcal O(I_{S'}(\mathcal M))
 =\mathcal O(S')[t]/(t^2)
 =\mathcal O(S')\oplus t\mathcal O(S'),
 \qquad
 \underline{\operatorname{Lie}}(\mathbf O_S/S)\simeq t\mathcal O_S.
\]
Let \(\rho_t:I_S(\mathcal M)\to S\) be the structural morphism,
corresponding to the injection
\[
 u_t:\mathcal O_S\hookrightarrow
 D_{\mathcal O_S}(\mathcal M)
 =\mathcal O_S\oplus t\mathcal O_S.
\]
Recall (cf. 3.4.2) that for every \(S\)-functor \(X\),
\(F(\mathcal N)(X)\) is the set of \(S\)-morphisms
\(\phi:I_X(\mathcal N)\to G\) such that
\(\phi\circ\varepsilon_{\mathcal N}=e\), and the action of
\(a\in\mathcal O(X)\) is
\[
 a\cdot\phi=\phi\circ a^*,
\]
where \(a^*\) is the endomorphism of \(I_X(\mathcal N)\) associated
with \(a\); see \SGARefLink{sga3:ii:subsection:2-1-3}{2.1.3}.

Consequently, by \SGARefLink{sga3:ii:remark:4-3-4}{4.3.4}, the composite
\[
 d:F(\mathcal N)\xrightarrow{\sim}
 F(\mathcal N)\otimes_{\mathcal O_S}t\mathcal O_S
 \longrightarrow
 \underline{\operatorname{Lie}}(F(\mathcal N)/S)
\]
is given, for \(S'\to S\) and
\(f\in\operatorname{Hom}_S(S',F(\mathcal N))\), by
\[
 f\longmapsto f\otimes t
 \longmapsto t\cdot(f\circ\rho_t)
 =f\circ\rho_t\circ t^*.
\]
Here \(f\circ\rho_t\in F(\mathcal N)(I_{S'}(\mathcal M))\) is viewed
as an \(S\)-morphism
\[
 \phi:I_{I_{S'}(\mathcal M)}(\mathcal N)\longrightarrow G
\]
such that \(\phi\circ\varepsilon_{\mathcal N}=e\). We must prove that
\(d\) is \(\mathcal O_S\)-linear, or equivalently that
\[
 d(a\cdot f)=u_t(a)\cdot d(f)
 \qquad(a\in\mathcal O(S')).
\]
Viewing \(f\in F(\mathcal N)(S')\) as a morphism
\(I_{S'}(\mathcal N)\to G\), one likewise has
\(a\cdot f=f\circ a^*\). Functoriality in \(X\) of the action of
\(\mathcal O(X)\) on \(I_X(\mathcal N)\) gives the following
commutative diagram.

\SGAThreeNativeDiagram[0.57\textwidth]
  {native_diagrams/exp2/exp2_p036_scalar_action_square.tex}
  {Compatibility of scalar action with the structural morphism \(\rho_t\).}
  {fig:exp2-p462-scalar-action}
  {Polo--Gille Exp2-14oct24.pdf, component-local p. 36; combined-reader p. 94; printed p. 84}

Therefore
\[
\begin{aligned}
 d(a\cdot f)
 &=f\circ a^*\circ\rho_t\circ t^*\\
 &=f\circ\rho_t\circ u_t(a)^*\circ t^*\\
 &=f\circ\rho_t\circ t^*\circ u_t(a)^*
  =u_t(a)\cdot d(f).
\end{aligned}
\]
The penultimate equality follows from the commutativity of
\(\mathcal O\). This completes the proof of \SGARefLink{sga3:ii:lemma:4-6-2}{Lemma 4.6.2}.

\medskip
\noindent\SGAInlineRefTarget{sga3:ii:theorem:4-7}\textbf{Theorem 4.7.}
\textit{If \(F\) is a good \(\mathcal O_S\)-module, then the
\(S\)-group
\(\underline{\operatorname{Aut}}_{\mathcal O_S\text{-mod.}}(F)\)
is good.}

\footnote{Editor's note: the original has been simplified here by using
the additions made in \SGARefLink{sga3:ii:corollary:4-3-1}{4.3.1} and \SGARefLink{sga3:ii:proposition:4-5}{4.5}.}
Indeed, by \SGARefLink{sga3:ii:proposition:4-5}{4.5},
\(\underline{\operatorname{Aut}}_{\mathcal O_S\text{-mod.}}(F)/S\)
satisfies \((E)\), while
\[
 \underline{\operatorname{Lie}}\!
  \left(\underline{\operatorname{Aut}}_{\mathcal O_S\text{-mod.}}(F)/S\right)
 \simeq
 \underline{\operatorname{End}}_{\mathcal O_S\text{-mod.}}(F)
\]
is a good \(\mathcal O_S\)-module.

% Source page: Exp2 component-local p. 37; combined-reader p. 95; printed p. 85.
\subsection*{\SGAInlineRefTarget{sga3:ii:subsection:4-7-1}4.7.1.}
\footnote{Editor's note: the numbering \SGARefLink{sga3:ii:subsection:4-7-1}{4.7.1}, \SGARefLink{sga3:ii:definition:4-7-2}{4.7.2}, and \SGARefLink{sga3:ii:subsection:4-7-3}{4.7.3} has
been added.}
Let \(G\) be an \(S\)-group and \(F\) a good
\(\mathcal O_S\)-module. Suppose given a linear representation of
\(G\) in \(F\), that is (\SGARefLink{sga3:i:definition:4-7-1}{I~4.7.1}), a morphism of \(S\)-groups
\[
 \rho:G\longrightarrow
 \underline{\operatorname{Aut}}_{\mathcal O_S\text{-mod.}}(F).
\]
If \(G/S\) satisfies \((E)\), then \SGARefLink{sga3:ii:remark:4-1-c}{4.1.C} and \SGARefLink{sga3:ii:proposition:4-5}{4.5} give a morphism of
\(\mathcal O_S\)-modules, denoted \(\rho'\) or \(d\rho\),
\[
 \underline{\operatorname{Lie}}(G/S)
 \longrightarrow
 \underline{\operatorname{Lie}}\!
  \left(\underline{\operatorname{Aut}}_{\mathcal O_S\text{-mod.}}(F)/S\right)
 \simeq
 \underline{\operatorname{End}}_{\mathcal O_S\text{-mod.}}(F).
\]
\footnote{Editor's note: the following sentence has been added.}
Furthermore, put
\(\mathcal O_{I_S}=\mathcal O_S\oplus t\mathcal O_S\), with
\(t^2=0\). It follows from \SGARefLink{sga3:ii:scholium:4-5-1}{4.5.1} that, if \(S'\to S\) and
\[
 X\in\underline{\operatorname{Lie}}(G/S)(S')\subset G(I_{S'}),
\]
then in
\(\underline{\operatorname{Aut}}_{\mathcal O_{I_{S'}}\text{-mod.}}
(F_{I_{S'}})\) one has
\begin{equation*}
 \rho(X)=\mathrm{id}+t\rho'(X).
 \tag{*}\label{sga3:ii:subsection:4-7-1:equation:star}
\end{equation*}
In other words, for every \(S''\to I_{S'}\) and \(f\in F(S'')\),
\[
 \rho(X)(f)=f+t\rho'(X)(f)
 \quad\text{in }F(S'').
\]

\medskip
\noindent\SGAInlineRefTarget{sga3:ii:definition:4-7-2}\textbf{Definition 4.7.2.}
Let \(G\) be a good \(S\)-group. Then
\(\underline{\operatorname{Lie}}(G/S)\) is a good
\(\mathcal O_S\)-module, and there is a morphism of \(S\)-groups
\[
 \operatorname{Ad}:G\longrightarrow
 \underline{\operatorname{Aut}}_{\mathcal O_S\text{-mod.}}
   \bigl(\underline{\operatorname{Lie}}(G/S)\bigr).
\]
By \SGARefLink{sga3:ii:subsection:4-7-1}{4.7.1}, it induces a morphism of \(\mathcal O_S\)-modules
\[
 \operatorname{ad}:
 \underline{\operatorname{Lie}}(G/S)
 \longrightarrow
 \underline{\operatorname{End}}_{\mathcal O_S\text{-mod.}}
   \bigl(\underline{\operatorname{Lie}}(G/S)\bigr),
\]
or, equivalently, an \(\mathcal O_S\)-bilinear morphism
\[
 \underline{\operatorname{Lie}}(G/S)
 \times_S
 \underline{\operatorname{Lie}}(G/S)
 \longrightarrow
 \underline{\operatorname{Lie}}(G/S),
 \qquad
 (x,y)\longmapsto[x,y]=\operatorname{ad}(x)\cdot y.
\]
Here \(x,y\) are arbitrary elements of
\[
 \underline{\operatorname{Lie}}(G/S)(S')
 =
 \underline{\operatorname{Lie}}(G_{S'}/S')(S').
\]
If \(G\) is commutative, then \([x,y]=0\).

\subsection*{\SGAInlineRefTarget{sga3:ii:subsection:4-7-3}4.7.3.}
The bracket admits the following equivalent description. It is enough
to consider sections over \(S\), that is,
\[
 x,y\in\underline{\operatorname{Lie}}(G/S)(S).
\]
There is a canonical isomorphism
\(I_S\times_SI_S\simeq I_{I_S}\). To avoid confusion, let \(I\) and
\(I'\) denote two copies of \(I_S\), and put
\[
 \mathcal O_I=\mathcal O_S[t],
 \qquad
 \mathcal O_{I'}=\mathcal O_S[t'],
 \qquad
 t^2=0=(t')^2.
\]
There is then a commutative square

\SGAThreeNativeDiagram[0.32\textwidth]
  {native_diagrams/exp2/exp2_p037_dual_number_base_square.tex}
  {The two dual-number structures on \(I\times_S I'\).}
  {fig:exp2-p473-dual-number-base}
  {Polo--Gille Exp2-14oct24.pdf, component-local p. 37; combined-reader p. 95; printed p. 85}

in which the two arrows from \(I\times I'\) identify it with the
scheme of dual numbers over \(I\) or over \(I'\), respectively.
Writing \(L=\underline{\operatorname{Lie}}(G/S)\), one obtains the
following commutative diagram of groups.

% Source page: Exp2 component-local p. 38; combined-reader p. 96; printed p. 86.
\SGAThreeNativeDiagram[0.69\textwidth]
  {native_diagrams/exp2/exp2_p038_exact_sequence_diagram_1.tex}
  {The first exact-sequence diagram used to construct the bracket.}
  {fig:exp2-p473-exact-sequences-one}
  {Polo--Gille Exp2-14oct24.pdf, component-local p. 38; combined-reader p. 96; printed p. 86}

The ninth object in this diagram is
\(\underline{\operatorname{Lie}}(L/S)(S)\). If \(G\) is good, this
is \(L(S)\), and hence one obtains the next commutative diagram. Its
rows and columns are exact sequences of groups; the five groups
denoted \(L(-)\) are commutative.\footnote{Editor's note: the original
has been expanded in what follows.} Under the identifications
\[
 L(I)=L(S)\oplus tL(S),
 \qquad
 L(I')=L(S)\oplus t'L(S),
\]
the injections \(L(S)\hookrightarrow L(I)\) and
\(L(S)\hookrightarrow L(I')\) are \(u\mapsto tu\) and
\(u\mapsto t'u\), respectively.

\SGAThreeNativeDiagram[0.69\textwidth]
  {native_diagrams/exp2/exp2_p038_exact_sequence_diagram_2.tex}
  {The exact-sequence diagram defining \([x,y]\).}
  {fig:exp2-p473-exact-sequences-two}
  {Polo--Gille Exp2-14oct24.pdf, component-local p. 38; combined-reader p. 96; printed p. 86}

In such a diagram, take \(x,y\) as indicated and choose arbitrary
lifts \(\widetilde x\in L(I)\) of \(x\) and \(y'\in L(I')\) of \(y\).
The commutator
\[
 \widetilde x\,y'\,\widetilde x^{-1}(y')^{-1}
 \quad\text{in }G(I\times I')
\]
is independent of the chosen lifts and is the image of a uniquely
determined element \(z\) as shown. One has \(z=[x,y]\).
\footnote{Editor's note: the original said in a footnote that the author,
at Gabriel's request, acknowledged that the exercise was not immediate,
which was precisely why it had not been placed in the text. The details
have been supplied here, taking account of the addition in \SGARefLink{sga3:ii:subsection:4-7-1}{4.7.1}; see
also [DG70], \(\S\)II.4, \SGARefLink{sga3:ii:subsection:4-2}{4.2}.}

Indeed, use the same letters \(x\) and \(y\) for their images under
the canonical sections \(L(S)\to L(I)\) and \(L(S)\to L(I')\).
Then
\[
 \widetilde x=xu,\qquad y'=yv,
 \qquad u,v\in L(S)=L(I)\cap L(I').
\]
Since \(L(I)\) and \(L(I')\) are commutative,
\[
\begin{aligned}
 \widetilde x\,y'\,\widetilde x^{-1}(y')^{-1}
 &=xuyvu^{-1}x^{-1}v^{-1}y^{-1}\\
 &=xuyu^{-1}vx^{-1}v^{-1}y^{-1}\\
 &=xyx^{-1}y^{-1}.
\end{aligned}
\]

% Source page: Exp2 component-local p. 39; combined-reader p. 97; printed p. 87.
This element maps to the unit in both \(G(I)\) and \(G(I')\), and
therefore comes from the unique \(z\) indicated above. Finally,
viewing \(y\) as an element of \(L_{I'}(I')\), and \(x\) as an
element of \(L(S)\subset G(I')\), \SGARefLink{sga3:ii:subsection:4-7-1}{formula 4.7.1} \((*)\) gives
\[
 xyx^{-1}
 =\operatorname{Ad}(x)(y)
 =(\mathrm{id}+t'\operatorname{ad}(x))(y)
 =y+t'[x,y].
\]
Thus \(xyx^{-1}y^{-1}\in L(I')\) is the image of
\(z=[x,y]\in L(S)\).

Two properties are now apparent from this construction:

\textit{\SGAInlineRefTarget{sga3:ii:subsection:4-7-3:item:i}(i)} the bracket is functorial in \(G\). More precisely,
\[
 G\longmapsto\underline{\operatorname{Lie}}(G/S)
\]
is a functor from good \(S\)-groups to good \(\mathcal O_S\)-modules
equipped with an \(\mathcal O_S\)-bilinear composition law;

\textit{(ii)} \([x,y]+[y,x]=0\), since the diagram is symmetric about
its main diagonal.\footnote{Editor's note: that is, \(x\) and \(y\)
play symmetric roles. The image of \([y,x]\) in \(G(I\times I')\) is
the commutator
\(y'\widetilde x(y')^{-1}\widetilde x^{-1}\), the inverse of
\(\widetilde x y'\widetilde x^{-1}(y')^{-1}=[x,y]\). On the other
hand (cf. \SGARefLink{sga3:ii:definition:4-10}{4.10} below), it is not known at this point whether
\([x,x]=0\) necessarily holds: if \(x\) is represented by
\(\widetilde x\in L(I)\) and by \(x'\in L(I')\), it is not a priori
clear that \(\widetilde x\) and \(x'\) commute.}

\medskip
\noindent\SGAInlineRefTarget{sga3:ii:proposition:4-8}\textbf{Proposition 4.8.}
\textit{Let \(F\) be a good \(\mathcal O_S\)-module. Under the
identification}
\[
 \underline{\operatorname{Lie}}\!
  \left(\underline{\operatorname{Aut}}_{\mathcal O_S\text{-mod.}}(F)/S\right)
 =
 \underline{\operatorname{End}}_{\mathcal O_S\text{-mod.}}(F),
\]
\textit{one has}
\[
 \operatorname{Ad}(g)\cdot Y=g\circ Y\circ g^{-1},
 \qquad
 [X,Y]=X\circ Y-Y\circ X,
\]
\textit{for every \(S'\to S\),
\(g\in
\underline{\operatorname{Aut}}_{\mathcal O_{S'}\text{-mod.}}(F_{S'})\),
and}
\[
 X,Y\in
 \underline{\operatorname{Lie}}\!
  \left(\underline{\operatorname{Aut}}_{\mathcal O_S\text{-mod.}}(F)/S\right)
  (S')
 =
 \underline{\operatorname{End}}_{\mathcal O_{S'}\text{-mod.}}(F_{S'}).
\]

\textit{Proof.}\footnote{Editor's note: the original has been expanded
and simplified in what follows.}
After base change, we may assume \(S'=S\). Put \(I=I_S\) and
\(\mathcal O_I=\mathcal O_S[t]\), with \(t^2=0\). Recall from \SGARefLink{sga3:ii:scholium:4-5-1}{4.5.1}
that the inclusion
\[
 i:\underline{\operatorname{End}}_{\mathcal O_S\text{-mod.}}(F)
 \hookrightarrow
 \underline{\operatorname{Aut}}_{\mathcal O_I\text{-mod.}}(F_I)
\]
sends \(Y\) to \(\mathrm{id}+tY\). By the definition of
\(\operatorname{Ad}(g)\) in 4.1.A,
\[
 \mathrm{id}+t\operatorname{Ad}(g)(Y)
 =g\circ(\mathrm{id}+tY)\circ g^{-1}
 =\mathrm{id}+t(g\circ Y\circ g^{-1}),
\]
and therefore
\(\operatorname{Ad}(g)(Y)=g\circ Y\circ g^{-1}\).

Let \(I'\) be a second copy of \(I_S\), and put
\(\mathcal O_{I'}=\mathcal O_S[t']\), with \((t')^2=0\). Apply \SGARefLink{sga3:ii:subsection:4-7-3}{4.7.3}
to
\[
 G=\underline{\operatorname{Aut}}_{\mathcal O_S\text{-mod.}}(F),
 \qquad
 L=\underline{\operatorname{Lie}}(G/S)
  =\underline{\operatorname{End}}_{\mathcal O_S\text{-mod.}}(F).
\]
Here \(L=\operatorname{End}(F)\) is the differential identification of
\SGARefLink{sga3:ii:proposition:4-5}{Proposition 4.5}. The re-edition prints \(\operatorname{Aut}(F)\) at
this point, but \(\operatorname{End}(F)\) is the type-correct object
used by the ensuing endomorphism calculation. Identify \(X\) with its image under the
canonical section \(L(S)\hookrightarrow L(I)\). Its image in
\(G(I\times I')\) is \(\mathrm{id}+t'X\), whose inverse is
\(\mathrm{id}-t'X\). Similarly, \(Y\) lifts to
\(\mathrm{id}+tY\), with inverse \(\mathrm{id}-tY\). Their
commutator is
\[
\begin{aligned}
 &(\mathrm{id}+t'X)\circ(\mathrm{id}+tY)
 \circ(\mathrm{id}-t'X)\circ(\mathrm{id}-tY)\\
 &\hspace{5em}
 =\mathrm{id}+tt'(X\circ Y-Y\circ X).
\end{aligned}
\]
This is the image in \(G(I\times I')\) of
\(Z=X\circ Y-Y\circ X\in L(S)\): first \(Z\) maps to
\(tZ\in L(I)\), and then to
\(\mathrm{id}+t'tZ\in G(I\times I')\). By \SGARefLink{sga3:ii:subsection:4-7-3}{4.7.3},
\([X,Y]=X\circ Y-Y\circ X\).

% Source page: Exp2 component-local p. 40; combined-reader p. 98; printed p. 88.
\medskip
\noindent\SGAInlineRefTarget{sga3:ii:corollary:4-8-1}\textbf{Corollary 4.8.1.}
\textit{Let \(G\) be a good \(S\)-group and let
\(x,y,z\in\underline{\operatorname{Lie}}(G/S)(S')\). Then}
\[
 [x,[y,z]]+[y,[z,x]]+[z,[x,y]]=0.
\]
\footnote{Editor's note: the original has been expanded in what follows.}
Indeed, since \(G\) is good,
\(\underline{\operatorname{Lie}}(G/S)\) is a good
\(\mathcal O_S\)-module; hence \SGARefLink{sga3:ii:theorem:4-7}{Theorem 4.7} shows that
\[
 \underline{\operatorname{Aut}}_{\mathcal O_S\text{-mod.}}
   \bigl(\underline{\operatorname{Lie}}(G/S)\bigr)
\]
is a good \(S\)-group. The morphism
\[
 \operatorname{Ad}:G\longrightarrow
 \underline{\operatorname{Aut}}_{\mathcal O_S\text{-mod.}}
  \bigl(\underline{\operatorname{Lie}}(G/S)\bigr)
\]
and the functoriality in \SGARefLink{sga3:ii:subsection:4-7-3:item:i}{4.7.3(i)} yield
\[
 \operatorname{ad}[x,y]
 =[\operatorname{ad}x,\operatorname{ad}y].
\]
Together with \SGARefLink{sga3:ii:proposition:4-8}{Proposition 4.8}, this gives
\[
 \operatorname{ad}[x,y]
 =\operatorname{ad}x\circ\operatorname{ad}y
  -\operatorname{ad}y\circ\operatorname{ad}x.
\]
Applying both sides to \(z\) gives the Jacobi identity.

\medskip
\noindent\SGAInlineRefTarget{sga3:ii:corollary:4-8-2}\textbf{Corollary 4.8.2.}
\textit{Let \(G\) be a good \(S\)-group acting linearly on a good
\(\mathcal O_S\)-module \(F\); equivalently, let \(F\) be a
\(G\)-\(\mathcal O_S\)-module, with \(G\) and \(F\) good. Then the
linear map}
\[
 \rho':\underline{\operatorname{Lie}}(G/S)
 \longrightarrow
 \underline{\operatorname{End}}_{\mathcal O_S\text{-mod.}}(F)
\]
\textit{is a representation; that is,}
\[
 \rho'([x,y])
 =\rho'(x)\circ\rho'(y)-\rho'(y)\circ\rho'(x).
\]

\medskip
\noindent\SGAInlineRefTarget{sga3:ii:scholium:4-9}\textbf{Scholium 4.9.}
To every good \(S\)-group \(G\) (for example, every representable
one) we have functorially associated a good \(\mathcal O_S\)-module
\(\underline{\operatorname{Lie}}(G/S)\) equipped with a bilinear map
satisfying
\[
 [x,y]+[y,x]=0,
 \qquad
 [x,[y,z]]+[y,[z,x]]+[z,[x,y]]=0.
\]
We shall call \(\underline{\operatorname{Lie}}(G/S)\), with this
structure, the \emph{Lie algebra} of \(G\) over \(S\). The qualification
is that it is not yet known whether it is an
\(\mathcal O_S\)-Lie algebra in the strict sense.\footnote{Editor's
note: the unresolved point is whether \([x,x]=0\); see \SGARefLink{sga3:ii:definition:4-10}{4.10}.}
Every linear representation of \(G\) in a good
\(\mathcal O_S\)-module \(F\) induces a representation of its Lie
algebra. In particular, the adjoint representation of \(G\) induces
the adjoint representation of its Lie algebra.

\medskip
\noindent\SGAInlineRefTarget{sga3:ii:definition:4-10}\textbf{Definition 4.10.}
An \(S\)-functor in groups \(G\) is called \emph{very good} if it is
good and \(\underline{\operatorname{Lie}}(G/S)\) is an
\(\mathcal O_S\)-Lie algebra, that is, if \([x,x]=0\) identically.

\medskip
\noindent\SGAInlineRefTarget{sga3:ii:example:4-10-1}\textbf{Examples 4.10.1.}
The following \(S\)-groups are very good:
\(\underline{\operatorname{Aut}}_{\mathcal O_S\text{-mod.}}(F)\)
for every good \(\mathcal O_S\)-module \(F\) (by \SGARefLink{sga3:ii:theorem:4-7}{4.7} and \SGARefLink{sga3:ii:proposition:4-8}{4.8});
every representable group (as shown below); every good \(S\)-group
admitting a monomorphism into a very good \(S\)-group, in particular
every good subgroup functor of a representable group; and every good
\(S\)-group admitting a faithful linear representation in a good
\(\mathcal O_S\)-module, for example every good \(S\)-group for which
\(\operatorname{Ad}\) is a monomorphism.

% Source page: Exp2 component-local p. 41; combined-reader p. 99; printed p. 89.
\subsection*{\SGAInlineRefTarget{sga3:ii:subsection:4-11}4.11.}
Suppose now that \(G\) is a group scheme over \(S\). By 4.1.4,
\(\underline{\operatorname{Lie}}(G/S)(S)\) identifies with the group
of infinitesimal automorphisms of \(G/S\) invariant under right
translation; by \SGARefLink{sga3:ii:subsection:3-14}{3.14}, this is the group of derivations of
\(\mathcal O_G\) over \(\mathcal O_S\) invariant under right
translation. This identification respects the module structure and
is an anti-isomorphism of Lie algebras.\footnote{Editor's note: a sign
error in the original has been corrected here; cf. [DG70],
\(\S\)II.4, 4.4 and \SGARefLink{sga3:ii:definition:4-6}{4.6}.}

To see this, argue as in \SGARefLink{sga3:ii:subsection:4-7-3}{4.7.3}. Put
\[
 \mathcal O_I=\mathcal O_S[t],
 \qquad
 \mathcal O_{I'}=\mathcal O_S[t'],
 \qquad
 t^2=0=(t')^2,
\]
and let \(x\in L(I)\), \(y\in L(I')\), where
\(L=\underline{\operatorname{Lie}}(G/S)\). Left translation
\(\lambda_x\), respectively \(\lambda_y\), is an \(S\)-automorphism
of \(G_{I\times I'}\) inducing the identity on \(G_{I'}\),
respectively on \(G_I\). It corresponds to the
\(\mathcal O_S\)-automorphism
\[
 u=\mathrm{id}+t\,d_x,
 \qquad\text{respectively}\qquad
 v=\mathrm{id}+t'\,d_y
\]
of
\[
 \mathcal O_{G_{I\times I'}}
 =\mathcal O_G[t,t']/(t^2,(t')^2),
\]
where \(d_x,d_y\) are \(\mathcal O_S\)-derivations of
\(\mathcal O_G\) invariant under right translation. The
correspondence between \(S\)-automorphisms of \(G_{I\times I'}\)
and \(\mathcal O_S\)-automorphisms of
\(\mathcal O_{G_{I\times I'}}\) is contravariant. Therefore
\[
 \lambda_x\lambda_y\lambda_x^{-1}\lambda_y^{-1}
\]
corresponds to
\[
 v^{-1}u^{-1}vu
 =\mathrm{id}+tt'(d_yd_x-d_xd_y).
\]
It follows from \SGARefLink{sga3:ii:subsection:4-7-3}{4.7.3} that
\[
 x\longmapsto-d_x
\]
is an isomorphism of Lie algebras; for further details, see [DG70],
\(\S\)II.4, 4.4 and \SGARefLink{sga3:ii:definition:4-6}{4.6}. The same argument applies to
\[
 \underline{\operatorname{Lie}}(G/S)(S')
 =
 \underline{\operatorname{Lie}}(G_{S'}/S')(S')
\]
for every \(S'\to S\). Thus one recovers the classical
definition.\footnote{Editor's note: the original has been expanded in
what follows. In particular, the numbering \SGARefLink{sga3:ii:scholium:4-11-1}{4.11.1}--\SGARefLink{sga3:ii:corollary:4-11-8}{4.11.8} has been
added for later reference.}

\medskip
\noindent\SGAInlineRefTarget{sga3:ii:scholium:4-11-1}\textbf{Scholium 4.11.1.}
\textit{Via the isomorphism \(x\mapsto-d_x\),
\(\underline{\operatorname{Lie}}(G/S)\) identifies with the functor
that assigns to each \(S'\to S\) the \(\mathcal O(S')\)-Lie algebra
of derivations of \(G_{S'}\) relative to \(S'\) invariant under right
translation.}

Since every representable group is already known to be good by
\SGARefLink{sga3:ii:example:4-6-1}{4.6.1}, the preceding discussion gives:

\medskip
\noindent\SGAInlineRefTarget{sga3:ii:corollary:4-11-2}\textbf{Corollary 4.11.2.}
\textit{Every representable group is very good.}

Let \(\varepsilon:S\to G\) be the unit section of \(G\), and put
\[
 \omega^1_{G/S}=\varepsilon^*(\Omega^1_{G/S}).
\]
Recall from \SGARefLink{sga3:ii:proposition:3-3}{3.3} that
\(\underline{\operatorname{Lie}}(G/S)\) is represented by the vector
fibration \(V(\omega^1_{G/S})\).

\medskip
\noindent\SGAInlineRefTarget{sga3:ii:scholium:4-11-3}\textbf{Scholium 4.11.3.}
\textit{Thus to every \(S\)-group scheme \(G\) we have functorially
associated a vector fibration}
\[
 \operatorname{Lie}(G/S)=V(\omega^1_{G/S})
\]
\textit{over \(S\), representing the functor
\(\underline{\operatorname{Lie}}(G/S)\), and hence endowed with the
structure of an \(S\)-scheme in \(\mathcal O_S\)-Lie algebras.
Moreover, this construction commutes with base extension and finite
products (cf. \SGARefLink{sga3:ii:proposition:3-4}{3.4} and \SGARefLink{sga3:ii:proposition:3-8}{3.8}).}

\medskip
\noindent\SGAInlineRefTarget{sga3:ii:remark:4-11-4}\textbf{Remarks 4.11.4.}
\footnote{Editor's note: what follows has been expanded, taking account
of the additions made in Expose I, \(\S\)\SGARefLink{sga3:i:subsection:6-8}{6.8}.}
Let \(\pi:G\to S\) be the structural morphism.

\textit{(a)} The \(\mathcal O_G\)-module \(\Omega^1_{G/S}\) is
naturally \((G\times_SG)\)-equivariant (cf. I, \(\S\)6). Hence, by
\SGARefLink{sga3:i:proposition:6-8-1}{I~6.8.1},
\[
 \Omega^1_{G/S}\simeq\pi^*(\omega^1_{G/S}).
\]
For example, \(\Omega^1_{G/S}\) is locally free, respectively locally
free of finite type, whenever \(\omega^1_{G/S}\) is. This holds in
particular when \(S\) is the spectrum of a field, respectively when
\(S\) is the spectrum of a field and \(G\) is locally of finite type
over \(S\).

% Source page: Exp2 component-local p. 42; combined-reader p. 100; printed p. 90.
\textit{(b)} By \SGARefLink{sga3:i:remark:6-8-2}{I~6.8.2}, \(\omega^1_{G/S}\) has a canonical
\(G\)-\(\mathcal O_S\)-module structure, which induces the adjoint
action on
\[
 V(\omega^1_{G/S})=\operatorname{Lie}(G/S).
\]
\footnote{Editor's note: for this reason, the linear action of \(G\) on
\(\omega^1_{G/S}\) could be called the pre-adjoint action; by an abuse
of terminology it will still be called the adjoint action. A slightly
more general construction, used in Expose III (especially \SGARefLink{sga3:iii:paragraph:0-8}{III~0.8}), is
as follows. Suppose that to every \(Y\to S\) one assigns an
\(\mathcal O_Y\)-module \(\mathcal N(Y)\), and to every
\(S\)-morphism \(\phi:Z\to Y\) a functorial morphism of
\(\mathcal O_Z\)-modules
\(\mathcal N(\phi):\phi^*\mathcal N(Y)\to\mathcal N(Z)\).
For example, if \(\mathcal J\) is a quasi-coherent ideal of
\(\mathcal O_S\), one may take
\(\mathcal N(Y)=\mathcal J\mathcal O_Y\). The action of \(G\) on
\(\omega^1_{G/S}\) then induces a generalized adjoint action of \(G\)
on the \(S\)-functor in abelian groups that assigns to
\(f:Y\to S\) the group
\(\operatorname{Hom}_{\mathcal O_Y}
(f^*\omega^1_{G/S},\mathcal N(Y))\).}

\textit{(c)} By EGA~I, 5.3.11 and 5.4.6,
\(\varepsilon:S\to G\) is an immersion; it is a closed immersion if
\(G\) is separated over \(S\). Thus \(\omega^1_{G/S}\) identifies
with \(\mathcal I/\mathcal I^2\), where \(\mathcal I\) is the
quasi-coherent ideal defining \(\varepsilon(S)\) in an open
subscheme \(U\subset G\) in which \(\varepsilon(S)\) is closed. If
\(G\to S\) is separated, one may take \(U=G\). If
\(G=\operatorname{Spec}A(G)\) is affine over \(S\), then
\(\mathcal I\) is the augmentation ideal of \(A(G)\), namely the
kernel of
\[
 \varepsilon^\natural:A(G)\longrightarrow\mathcal O_S
\]
(cf. \SGARefLink{sga3:i:subsection:4-2}{I~4.2}).\footnote{Editor's note: the preceding description of
\(\omega^1_{G/S}\), which is useful later (for example in
VII\textsubscript{A}, \SGARefLink{sga3:VIIA:5.5}{5.5}), has been added.}

\medskip
\noindent\SGAInlineRefTarget{sga3:ii:remark:4-11-5}\textbf{Remark 4.11.5.}
The isomorphism
\(\Omega^1_{G/S}\simeq\pi^*(\omega^1_{G/S})\) identifies the
\(\mathcal O_S\)-module \(\omega^1_{G/S}\) with the sheaf
\[
 \pi_*(\Omega^1_{G/S})^G
\]
of differentials of \(G\) relative to \(S\) invariant under right
translation: its sections over an open \(U\subset S\) are the
sections of \(\Omega^1_{G/S}\) over \(\pi^{-1}(U)\) invariant under
right translation (cf. \SGARefLink{sga3:i:remark:6-8-3}{I~6.8.3}; compare also
VII\textsubscript{A}, \SGARefLink{sga3:VIIA:2.4}{2.4}).\footnote{Editor's note: this description of
\(\omega^1_{G/S}\) in terms of invariant differentials is not used
below.}

\medskip
\noindent\SGAInlineRefTarget{sga3:ii:notation:4-11-6}\textbf{Notation 4.11.6.}
Write \(\mathop{\mathscr L\!ie}(G/S)\) for the sheaf of sections of the vector
fibration \(\operatorname{Lie}(G/S)\to S\). It is the
\(\mathcal O_S\)-module
\[
 \mathop{\mathscr L\!ie}(G/S)
 =(\omega^1_{G/S})^\vee
 =\mathcal Hom_{\mathcal O_S}(\omega^1_{G/S},\mathcal O_S),
\]
dual to \(\omega^1_{G/S}\) (cf. EGA~II, 1.7.9), and it carries a
structure of \(\mathcal O_S\)-Lie algebra.

In general this construction does not commute with base extension.
Consequently, the Lie-algebra structure on this module does not
determine the structure of an \(S\)-scheme in
\(\mathcal O_S\)-Lie algebras on
\(\operatorname{Lie}(G/S)\).\footnote{Editor's note: the reason is
that \(\mathop{\mathscr L\!ie}(G/S)\simeq(\omega^1_{G/S})^\vee\) need not
determine \(\omega^1_{G/S}\). For example, let
\(S=\mathbb A^1_k\), where \(k\) is a field, and let \(G\) be the
\(S\)-group whose fiber at \(s=0\) is \(\mathbb G_{a,k}\), while all
other fibers are the unit group. Then
\(\mathop{\mathscr L\!ie}(G/S)=0\), but
\(\operatorname{Lie}(G_s/k)(k)=k\).}

Nevertheless:

\medskip
\noindent\SGAInlineRefTarget{sga3:ii:lemma:4-11-7}\textbf{Lemma 4.11.7.}
\textit{Suppose that \(\omega^1_{G/S}\) is locally free of finite type.
This is the case, in particular, if \(G\) is smooth over \(S\)
(cf. EGA~IV\textsubscript{4}, 17.2.4), or if \(S\) is the spectrum
of a field and \(G\) is locally of finite type over \(S\). Then}
\[
 \mathop{\mathscr L\!ie}(G/S)^\vee
 \simeq(\omega^1_{G/S})^{\vee\vee}
 \simeq\omega^1_{G/S},
\]
\textit{and therefore}
\[
 \operatorname{Lie}(G/S)
 =V(\omega^1_{G/S})
 =V\bigl(\mathop{\mathscr L\!ie}(G/S)^\vee\bigr)
 =W\bigl(\mathop{\mathscr L\!ie}(G/S)\bigr),
\]
\textit{where the final equality follows from \SGARefLink{sga3:i:corollary:4-6-5-1}{I~4.6.5.1}.}

% Source page: Exp2 component-local p. 43; combined-reader p. 101; printed p. 91.
Finally, let \(G\to H\) be a monomorphism of functors in groups.
Then
\[
 \operatorname{Lie}(G/S)\longrightarrow\operatorname{Lie}(H/S)
\]
is also a monomorphism (cf. \SGARefLink{sga3:ii:proposition:3-7}{3.7}). Every vector monomorphism of vector
fibrations is a closed immersion.\footnote{Editor's note: let
\(f:M\to N\) be a morphism of \(\mathcal O_S\)-modules and let
\(P=\operatorname{Coker}(f)\). If \(V(N)\to V(M)\) is a monomorphism,
the surjective morphism
\(\operatorname{Sym}(N)\to\operatorname{Sym}(P)\) factors through
\(\mathcal O_S\); hence \(P=0\).}
It follows that:

\medskip
\noindent\SGAInlineRefTarget{sga3:ii:corollary:4-11-8}\textbf{Corollary 4.11.8.}
\textit{Let \(G\hookrightarrow H\) be a monomorphism of
\(S\)-group schemes.}

\textit{(i) The morphism}
\[
 \operatorname{Lie}(G/S)\longrightarrow\operatorname{Lie}(H/S)
\]
\textit{is a closed immersion; consequently}
\[
 \omega^1_{H/S}\longrightarrow\omega^1_{G/S}
\]
\textit{is an epimorphism.}

\textit{(ii) If \(\omega^1_{G/S}\) is locally free of finite type,
then the corresponding morphism}
\[
 \mathop{\mathscr L\!ie}(G/S)\longrightarrow
 \mathop{\mathscr L\!ie}(H/S)
\]
\textit{identifies \(\mathop{\mathscr L\!ie}(G/S)\) with an
\(\mathcal O_S\)-submodule of \(\mathop{\mathscr L\!ie}(H/S)\) that is locally
a direct summand.}

% Section 5 begins immediately after this point on source local p. 43
% and is intentionally excluded from this component.

% END INCLUDED SGA3_English_current_buildable_source_20260823/source/inputs/II/tex/components/13_expose_II_s45_s411_lie_structure.tex

% BEGIN INCLUDED SGA3_English_current_buildable_source_20260823/source/inputs/II/tex/components/14_expose_II_s5_s6_examples_misc.tex
% SGA 3, Expose II, Loop 1 English reconstruction.
% Scope: Sections 5 and 6 through the bibliography/end, beginning after
% Corollary 4.11.8 and stopping at the end of the Expose.
% Authority: Polo--Gille Exp2-14oct24.pdf, component-local pages 43--52
% (combined-reader pages 101--110). Force-OCR is a drafting locator; English
% wording was compared with Jacob Reinhold's CC BY 4.0 Markdown lineage.
% The PDF controls text, formulas, numbering, notes, and all diagram images.

% Source page: Exp2 component-local p. 43; combined-reader p. 101.
\section*{\SGAInlineRefTarget{sga3:ii:section:5}5. Computation of some Lie algebras}
\addcontentsline{toc}{section}{5. Computation of some Lie algebras}

\subsection*{\SGAInlineRefTarget{sga3:ii:subsection:5-1}5.1. Examples of Lie algebras: diagonalizable groups}

Let \(G=D_S(M)\) be a diagonalizable group over \(S\) (\SGARefLink{sga3:i:subsection:4-4}{I~4.4}). Since
the formation of \(\underline{\operatorname{Lie}}(G/S)\) commutes with
base extension, it is enough to carry out the construction for
\(G=D(M)\). We then have
\[
 G(I_S)=\operatorname{Hom}_{\mathrm{gr}}
   \bigl(M,\Gamma(I_S,\mathscr O_{I_S})^\times\bigr)
 =\operatorname{Hom}_{\mathrm{gr}}
   \bigl(M,\Gamma(S,\mathscr D_{\mathscr O_S})^\times\bigr).
\]
There is a split exact sequence
\[
 1\longrightarrow\Gamma(S,\mathscr O_S)
  \longrightarrow\Gamma(S,\mathscr D_{\mathscr O_S})^\times
  \longrightarrow\Gamma(S,\mathscr O_S)^\times
  \longrightarrow1,
\]
which shows that \(\underline{\operatorname{Lie}}(G)(S)\) identifies
with \(\operatorname{Hom}_{\mathrm{gr}}(M,\mathscr O(S))\), endowed
with its evident \(\mathscr O(S)\)-module structure. After base change
we therefore obtain:

\medskip
\noindent\SGAInlineRefTarget{sga3:ii:proposition:5-1}\textbf{Proposition 5.1.}
\textit{There are isomorphisms}
\[
 \underline{\operatorname{Hom}}_{S\text{-}\mathrm{gr}}
      (M_S,\mathbf O_S)
 \xrightarrow{\sim}
 \underline{\operatorname{Lie}}(D_S(M)/S)
\]
\textit{and}
\[
 \mathscr Hom_{\mathrm{gr}}(\widetilde M_S,\mathscr O_S)
 \xrightarrow{\sim}
 \mathscr Lie(D_S(M)/S),
\]
\textit{where, in the second isomorphism, \(\widetilde M_S\) denotes
the constant sheaf of groups on \(S\) defined by \(M\), and
\(\mathscr Hom_{\mathrm{gr}}\) denotes the sheaf of homomorphisms of
sheaves of groups.}

\medskip
\noindent\SGAInlineRefTarget{sga3:ii:corollary:5-1-1}\textbf{Corollary 5.1.1.}
\textit{If \(M\) is free of finite type---or, as we shall say later,
if \(D_S(M)\) is a split torus---then, with \(W\) as in \SGARefLink{sga3:i:proposition:4-6-5}{I~4.6.5},}
\[
 W\bigl(\mathscr Lie(D_S(M)/S)\bigr)
   \xrightarrow{\sim}\underline{\operatorname{Lie}}(D_S(M)/S),
 \qquad
 M^\vee\otimes_{\mathbb Z}\mathscr O_S
   \xrightarrow{\sim}\mathscr Lie(D_S(M)/S).
\]
\textit{Here \(M^\vee\) denotes the dual of the abelian group \(M\).
In particular,}
\[
 \mathbf O_S\xrightarrow{\sim}
   \underline{\operatorname{Lie}}(\mathbb G_{m,S}/S),
 \qquad
 \mathscr O_S\xrightarrow{\sim}
   \mathscr Lie(\mathbb G_{m,S}/S).
\]

\subsection*{\SGAInlineRefTarget{sga3:ii:subsection:5-2}5.2. Normalizers and centralizers}

We first prove a few lemmas. Recall (\SGARefLink{sga3:i:lemma:3-1-1}{I~3.1.1}) that a sequence
\[
 0\longrightarrow F'\longrightarrow F\longrightarrow F''\longrightarrow0
\]
of \(\mathbf O_S\)-modules is called \emph{exact} if, for every
\(S'\to S\), the sequence
\[
 0\longrightarrow F'(S')\longrightarrow F(S')
  \longrightarrow F''(S')\longrightarrow0
\]
of \(\mathbf O(S')\)-modules is exact. Likewise, a sequence
\(1\to G'\to G\to G''\to1\) of \(S\)-groups is called exact if the
corresponding sequence of groups on \(S'\)-points is exact for every
\(S'\to S\).

\medskip
\noindent\SGAInlineRefTarget{sga3:ii:lemma:5-2-1}\textbf{Lemma 5.2.1.}
\textit{Let \(1\to G'\to G\to G''\to1\) be an exact sequence of
\(S\)-groups.}
\footnote{Editor's note: The statement and proof of this lemma have
been expanded; this will be useful in \SGARefLink{sga3:ii:theorem:5-3-1}{5.3.1}.}
\begin{enumerate}
\renewcommand{\labelenumi}{(\roman{enumi})}
\item \textit{The sequences}
\[
 1\longrightarrow T_{G'/S}(\mathscr M)
 \longrightarrow T_{G/S}(\mathscr M)
 \longrightarrow T_{G''/S}(\mathscr M)\longrightarrow1
\]
\textit{and}
\[
 1\longrightarrow\underline{\operatorname{Lie}}(G'/S,\mathscr M)
 \longrightarrow\underline{\operatorname{Lie}}(G/S,\mathscr M)
 \longrightarrow\underline{\operatorname{Lie}}(G''/S,\mathscr M)
 \longrightarrow1
\]
\textit{are exact.}
\item \textit{Let \(1\to H'\to H\to H''\to1\) be a second sequence
of groups. It is exact if and only if}
\[
 1\longrightarrow G'\times_S H'
 \longrightarrow G\times_S H
 \longrightarrow G''\times_S H''\longrightarrow1
\]
\textit{is exact.}
\item\SGAInlineRefTarget{sga3:ii:lemma:5-2-1:item:iii} \textit{If two of \(G',G,G''\) satisfy condition (E), then so
does the third.}
\item\SGAInlineRefTarget{sga3:ii:lemma:5-2-1:item:iv} \textit{If \(0\to F'\to F\to F''\to0\) is exact and two of the
\(\mathbf O_S\)-modules \(F',F,F''\) are good, then so is the third.}
\item \textit{If two of the \(S\)-groups \(G',G,G''\) are good, then
so is the third.}
\end{enumerate}

The first part of (i) is immediate, and the second follows from it.
Likewise, (ii) is immediate. We prove (iii). For brevity, write
\(L(\mathscr M)=\underline{\operatorname{Lie}}(G,\mathscr M)\),
\(L'(\mathscr M)=\underline{\operatorname{Lie}}(G',\mathscr M)\),
and similarly for \(L''\). We have the following commutative diagram.
\footnote{Editor's note: The proofs of (iii) and (v) have been added.}

\SGAThreeNativeDiagram[0.93\textwidth]
  {native_diagrams/exp2/exp2_p044_lemma521_conditionE_diagram.tex}
  {The exact-sequence diagram used to prove \SGARefLink{sga3:ii:lemma:5-2-1:item:iii}{Lemma~5.2.1(iii)}.}
  {fig:sga3-exp2-p044-lemma521-conditionE}
  {Polo--Gille Exp2-14oct24.pdf, component-local p. 44;
   combined-reader p. 102}

The upper row is exact by (i), and the lower row is exact by (i) and
(ii). Assertion (iii) follows.

We prove (iv). There is a commutative diagram

\SGAThreeNativeDiagram[0.96\textwidth]
  {native_diagrams/exp2/exp2_p044_lemma521_good_modules_diagram.tex}
  {The exact-sequence diagram used to prove \SGARefLink{sga3:ii:lemma:5-2-1:item:iv}{Lemma~5.2.1(iv)}.}
  {fig:sga3-exp2-p044-lemma521-good-modules}
  {Polo--Gille Exp2-14oct24.pdf, component-local p. 44;
   combined-reader p. 102}

with exact rows: the first because
\(T_{\mathbf O_S/S}(\mathscr M)\) is a free, hence flat,
\(\mathbf O_S\)-module, and the second by (i). Thus, if two of
\(F',F,F''\) are good, so is the third. Finally, (v) follows from
(iii) and (iv).

\medskip
\noindent\SGAInlineRefTarget{sga3:ii:lemma:5-2-2}\textbf{Lemma 5.2.2.}
\textit{Let \(G\) be an \(S\)-group, let \(E,H\) be two \(G\)-objects,
and let \(F\) be a \(G\)-\(\mathbf O_S\)-module.}
\begin{enumerate}
\renewcommand{\labelenumi}{(\roman{enumi})}
\item\SGAInlineRefTarget{sga3:ii:lemma:5-2-2:item:i} \textit{The canonical homomorphism
\(E^G\times_S H^G\to(E\times_S H)^G\) is an isomorphism.}
\item\SGAInlineRefTarget{sga3:ii:lemma:5-2-2:item:ii} \textit{If \(F\) is a good \(\mathbf O_S\)-module, then so is
\(F^G\).}
\end{enumerate}

The first assertion is immediate; we prove the second. For every
\(S'\to S\), there is a commutative diagram

\SGAThreeNativeDiagram[0.72\textwidth]
  {native_diagrams/exp2/exp2_p045_lemma522_invariants_square.tex}
  {The invariants-and-base-change square in the proof of \SGARefLink{sga3:ii:lemma:5-2-2}{Lemma~5.2.2}.}
  {fig:sga3-exp2-p045-lemma522-invariants}
  {Polo--Gille Exp2-14oct24.pdf, component-local p. 45;
   combined-reader p. 103}

and we must prove that \(\phi\) is bijective. It is plainly injective.
To prove surjectivity, let \((t_0,\ldots,t_n)\) be a basis of the free
\(\mathbf O(S')\)-module \(\mathbf O(I_{S'}(\mathscr M))\), and let
\[
 u=\sum_{i=0}^{n}f_i\otimes t_i
\]
be an element of
\(F(S')\otimes_{\mathbf O(S')}\mathbf O(I_{S'}(\mathscr M))\) such
that \(\phi_1(u)\) belongs to \(F^G(I_{S'}(\mathscr M))\). We show that
each \(f_i\) belongs to \(F^G(S')\).

Let \(S''\to S'\). Via the zero section \(\varepsilon_{\mathscr M}\),
we may regard \(S''\) as lying over \(I_{S'}(\mathscr M)\). Then, for
every \(g\in G(S'')\),
\[
 u_{S''}=g\cdot u_{S''}
   =\sum_{i=0}^{n}g\cdot(f_i)_{S''}\otimes(t_i)_{S''}.
\]
Since the \((t_i)_{S''}\) form a basis of
\[
 \mathbf O(I_{S''}(\mathscr M))
 \simeq \mathbf O(I_{S'}(\mathscr M))
        \otimes_{\mathbf O(S')}\mathbf O(S'')
\]
over \(\mathbf O(S'')\), it follows that
\(g\cdot(f_i)_{S''}=(f_i)_{S''}\), hence \(f_i\in F^G(S')\) for all
\(i\).

\medskip
\noindent\textbf{Notation.}
If \(E\) is an \(S\)-group and \(F\) a sub-\(S\)-group, write \(E/F\)
for the \(S\)-functor sending \(S'\to S\) to the set
\(E(S')/F(S')\) of classes \(\bar x=xF(S')\). If \(E\) is commutative,
then \(E/F\) is a commutative \(S\)-group.
\footnote{Editor's note: These notations have been expanded.}

Now let \(G\) be an \(S\)-group and \(K\) a sub-\(S\)-group, and set
\(E=\underline{\operatorname{Lie}}(G/S,\mathscr M)\) and
\(F=\underline{\operatorname{Lie}}(K/S,\mathscr M)\). The adjoint
action of \(K\) on \(E\) stabilizes \(F\), and therefore induces an
action on \(E/F\). For every \(S'\to S\),
\[
\begin{split}
 (E/F)^K(S')=\{\bar x\in E(S')/F(S')\mid{}&
  \text{for every }f:S''\to S'\text{ and }k\in K(S''),\\[-2pt]
 &f^*(\bar x^{-1})\operatorname{Ad}(k)(f^*(\bar x))\in F(S'')\},
\end{split}
\]
where \(f^*(\bar x)\) denotes the image of \(\bar x\) in \(E(S'')\).

\medskip
\noindent\SGAInlineRefTarget{sga3:ii:theorem:5-2-3}\textbf{Theorem 5.2.3.}
\textit{Let \(G\) be an \(S\)-group and \(K\) a sub-\(S\)-group. Put
(\SGARefLink{sga3:i:definition:2:3:3}{I~2.3.3})}
\[
 N=\underline{\operatorname{Norm}}_G(K),
 \qquad
 Z=\underline{\operatorname{Centr}}_G(K),
\]
\textit{and let \(K\) act on
\(\underline{\operatorname{Lie}}(G/S,\mathscr M)\) through the adjoint
representation of \(G\).}
\begin{enumerate}
\renewcommand{\labelenumi}{(\roman{enumi})}
\item\SGAInlineRefTarget{sga3:ii:theorem:5-2-3:item:i} \textit{If the group law of
\(\underline{\operatorname{Lie}}(G/S,\mathscr M)\) is
commutative,\footnote{This condition is automatic when \(G\) satisfies
(E), for example when \(G\) is representable.}\footnote{Editor's note:
For an example in which the \(S\)-group
\(\underline{\operatorname{Lie}}(G/S)\) is noncommutative, see \SGARefLink{sga3:ii:subsection:6-3}{6.3}.}
then}
\[
 \frac{\underline{\operatorname{Lie}}(N/S,\mathscr M)}
      {\underline{\operatorname{Lie}}(K/S,\mathscr M)}
 =
 \left(
 \frac{\underline{\operatorname{Lie}}(G/S,\mathscr M)}
      {\underline{\operatorname{Lie}}(K/S,\mathscr M)}
 \right)^K .
\]
\item\SGAInlineRefTarget{sga3:ii:theorem:5-2-3:item:ii} \textit{If the same group law is commutative, then}
\[
 \underline{\operatorname{Lie}}(Z/S,\mathscr M)
   =\underline{\operatorname{Lie}}(G/S,\mathscr M)^K .
\]
\item \textit{If \(G\) satisfies (E), then \(Z\) satisfies (E); if
\(G\) and \(K\) satisfy (E), then \(N\) satisfies (E).}
\item \textit{If \(G\) is good, then \(Z\) is good; if, moreover,
\(G\) is very good, then \(Z\) is very good.}
\item \textit{If \(G\) and \(K\) are good, then \(N\) is good; if,
moreover, \(G\) is very good, then \(N\) is very good.}
\end{enumerate}

To prove (i) and (ii),\footnote{Editor's note: The proof has been
expanded, and in particular \SGARefLink{sga3:ii:lemma:5-2-3-0}{Lemma~5.2.3.0}, implicit in the original,
has been added.} we use the following lemma, which follows from the
diagram of exact sequences considered in \SGARefLink{sga3:ii:definition:4-1-b}{4.1.B}, with \(G\) and \(H\)
interchanged.

\medskip
\noindent\SGAInlineRefTarget{sga3:ii:lemma:5-2-3-0}\textbf{Lemma 5.2.3.0.}
\textit{Let \(G\) be an \(S\)-group, \(H\) a sub-\(S\)-group, and
\(\mathscr M\) a free \(\mathscr O_S\)-module of finite type. Then
\(T_{H/S}(\mathscr M)\) and
\(\underline{\operatorname{Lie}}(G/S,\mathscr M)\) are sub-\(S\)-groups
of \(T_{G/S}(\mathscr M)\), and}
\[
 T_{H/S}(\mathscr M)\cap
 \underline{\operatorname{Lie}}(G/S,\mathscr M)
 =
 \underline{\operatorname{Lie}}(H/S,\mathscr M),
\]
\textit{where}
\[
 T_{H/S}(\mathscr M)\cap
 \underline{\operatorname{Lie}}(G/S,\mathscr M)
 \mathrel{\mathop:}=
 T_{H/S}(\mathscr M)\times_{T_{G/S}(\mathscr M)}
 \underline{\operatorname{Lie}}(G/S,\mathscr M).
\]

Since the functors in (i) and (ii) commute with base extension, it is
enough to prove the equalities on \(S\)-points. Put \(H=N\)
(respectively, \(H=Z\)) and let \(\alpha=\pm1\). By the preceding
lemma and the definitions of \(N\) and \(Z\) (\SGARefLink{sga3:i:definition:2:3:3}{I~2.3.3}),
\begin{equation}
\begin{split}
 \underline{\operatorname{Lie}}(H/S,\mathscr M)(S)
 =\bigl\{X\in{}&
 \underline{\operatorname{Lie}}(G/S,\mathscr M)(S)
 \subset G(I_S(\mathscr M))\ \bigm|\ \\
 &\text{for every }f:S'\to I_S(\mathscr M)
 \text{ and }u\in K(S'),\\
 &f^*(X^\alpha)\,u\,f^*(X^{-\alpha})\,u^{-1}\in K(S'),\\[-2pt]
 &\text{respectively }
 f^*(X)\,u\,f^*(X^{-1})\,u^{-1}=1
 \bigr\}.
\end{split}
\tag{*}\label{sga3:ii:lemma:5-2-3-0:equation:star}
\end{equation}
Here \(f^*(X)\) denotes the image of \(X\) in \(G(S')\).

For brevity, write
\[
 \mathfrak g=\underline{\operatorname{Lie}}(G/S,\mathscr M),\quad
 \mathfrak k=\underline{\operatorname{Lie}}(K/S,\mathscr M),\quad
 \mathfrak n=\underline{\operatorname{Lie}}(N/S,\mathscr M),\quad
 \mathfrak z=\underline{\operatorname{Lie}}(Z/S,\mathscr M).
\]
If \(X\in\underline{\operatorname{Lie}}(H/S,\mathscr M)(S)\), the
equalities (*) hold for every \(f:S'\to S\), because \(f\) factors
through \(I_S(\mathscr M)\) as
\[
 f=\rho\circ\varepsilon_{\mathscr M}\circ f,
\]
where \(\rho:I_S(\mathscr M)\to S\) is the structural morphism and
\(\varepsilon_{\mathscr M}:S\to I_S(\mathscr M)\) is the zero section.
Consequently,
\[
 \mathfrak n(S)/\mathfrak k(S)
   \subset(\mathfrak g/\mathfrak k)^K(S),
 \qquad
 \mathfrak z(S)\subset\mathfrak g^K(S).
\]

For the reverse inclusions, assume henceforth that
\(\mathfrak g\) is commutative. Then
\(\mathfrak g^K\) and \((\mathfrak g/\mathfrak k)^K\) are commutative
\(S\)-groups. Let \(X\in\mathfrak g(S)\), and write \(\bar X\) for its
image in \((\mathfrak g/\mathfrak k)(S)\). Suppose that
\(\bar X\in(\mathfrak g/\mathfrak k)^K(S)\) (respectively,
\(X\in\mathfrak g^K(S)\)); we show that
\(X\in\mathfrak n(S)\) (respectively, \(X\in\mathfrak z(S)\)).

Let \(f:S'\to I_S(\mathscr M)\). Consider the cartesian square

\SGAThreeNativeDiagram[0.48\textwidth]
  {native_diagrams/exp2/exp2_p047_theorem523_cartesian_square.tex}
  {The cartesian square used in the reverse-inclusion argument.}
  {fig:sga3-exp2-p047-theorem523-cartesian}
  {Polo--Gille Exp2-14oct24.pdf, component-local p. 47;
   combined-reader p. 105}

and let \(h\) be the section of \(\rho'\) defined by \(f\). It is
enough to show that, for every \(v\in K(I_{S'}(\mathscr M))\),
\begin{equation}
 \left.
 \begin{aligned}
 p^*(X^\alpha)\,v\,p^*(X^{-\alpha})\,v^{-1}
     &\in K(I_{S'}(\mathscr M)),\\
 \text{respectively}\qquad
 p^*(X)\,v\,p^*(X^{-1})\,v^{-1}&=1.
 \end{aligned}
 \right\}
\tag{**}\label{sga3:ii:lemma:5-2-3-0:equation:double-star}
\end{equation}
Indeed, take \(v=\rho'^*(u)\) and apply \(h^*\) to (**). This gives
(*), because \(\rho'\circ h=\operatorname{id}_{S'}\) and
\(p\circ h=f\).

We now prove (**), writing \(X\) for \(p^*(X)\). Every
\(v\in K(I_{S'}(\mathscr M))\) is uniquely of the form \(Yk\), with
\(Y\in\underline{\operatorname{Lie}}(K/S,\mathscr M)(S')\) and
\(k\in K(S')\). The expression
\(X^\alpha vX^{-\alpha}v^{-1}\) becomes
\[
 X^\alpha Y\bigl(kX^{-\alpha}k^{-1}\bigr)Y^{-1}
 =X^\alpha\operatorname{Ad}(k)(X^{-\alpha}),
\]
because \(\underline{\operatorname{Lie}}(G/S,\mathscr M)\) is
commutative. This element lies in
\(\underline{\operatorname{Lie}}(G/S,\mathscr M)(S')\). By
\SGARefLink{sga3:ii:lemma:5-2-3-0}{Lemma~5.2.3.0}, condition (**) is therefore equivalent, for every
\(k\in K(S')\), to
\[
 \begin{cases}
 \operatorname{Ad}(k)(X)=X, & H=Z,\\
 X^\alpha\operatorname{Ad}(k)(X^{-\alpha})
   \in\underline{\operatorname{Lie}}(K/S,\mathscr M)(S'), & H=N.
 \end{cases}
\]
For \(H=Z\), this follows from \(X\in\mathfrak g^K(S)\). For \(H=N\),
the condition may equally be written
\begin{equation}
 \operatorname{Ad}(k)(\bar X)=\bar X,
 \qquad
 \operatorname{Ad}(k)(\bar X^{-1})=\bar X^{-1}.
\tag{*'}\label{sga3:ii:lemma:5-2-3-0:equation:star-prime}
\end{equation}
The second equality follows from the first because the action of \(K\)
on \(\mathfrak g/\mathfrak k\) respects its group structure. Thus (**)
also follows from
\(\bar X\in(\mathfrak g/\mathfrak k)^K(S)\), proving (i) and (ii).

For (iii)--(v), put
\[
 \mathfrak g(\mathscr M)
   =\underline{\operatorname{Lie}}(G/S,\mathscr M)
\]
and define \(\mathfrak k(\mathscr M)\), \(\mathfrak z(\mathscr M)\),
and \(\mathfrak n(\mathscr M)\) similarly. If \(G/S\) satisfies (E),
then
\[
 \mathfrak g(\mathscr M\oplus\mathscr N)
 \simeq\mathfrak g(\mathscr M)\times_S\mathfrak g(\mathscr N),
\]
and hence, by \SGARefLink{sga3:ii:lemma:5-2-2:item:i}{Lemma~5.2.2(i)},
\[
 \mathfrak g(\mathscr M\oplus\mathscr N)^K
 \simeq
 \mathfrak g(\mathscr M)^K\times_S\mathfrak g(\mathscr N)^K.
\]
It follows that
\(\mathfrak z(\mathscr M\oplus\mathscr N)
\simeq\mathfrak z(\mathscr M)\times_S\mathfrak z(\mathscr N)\), so
\(Z\) satisfies (E). If \(K/S\) also satisfies (E), then
\[
 (\mathfrak g/\mathfrak k)(\mathscr M\oplus\mathscr N)
 \simeq
 (\mathfrak g/\mathfrak k)(\mathscr M)
 \times_S(\mathfrak g/\mathfrak k)(\mathscr N),
\]
\[
 (\mathfrak g/\mathfrak k)^K(\mathscr M\oplus\mathscr N)
 \simeq
 (\mathfrak g/\mathfrak k)(\mathscr M)^K
 \times_S(\mathfrak g/\mathfrak k)(\mathscr N)^K,
\]
and we have the following commutative diagram.

\SGAThreeNativeDiagram[0.96\textwidth]
  {native_diagrams/exp2/exp2_p048_theorem523_exact_diagram.tex}
  {The exact-sequence comparison proving condition (E) for \(N\).}
  {fig:sga3-exp2-p048-theorem523-exact}
  {Polo--Gille Exp2-14oct24.pdf, component-local p. 48;
   combined-reader p. 106}

It follows that
\[
 \mathfrak n(\mathscr M\oplus\mathscr N)
 \simeq\mathfrak n(\mathscr M)\times_S\mathfrak n(\mathscr N),
\]
so \(N\) satisfies (E).

Now write
\(\mathfrak g=\mathfrak g(\mathscr O_S)
=\underline{\operatorname{Lie}}(G/S)\), and similarly
\(\mathfrak k,\mathfrak z,\mathfrak n\). If \(G\) is good, then
\(\mathfrak g\) is a good \(\mathbf O_S\)-module, and \SGARefLink{sga3:ii:lemma:5-2-2:item:ii}{Lemma~5.2.2(ii)}
shows that \(\mathfrak z\simeq\mathfrak g^K\) is good; since \(Z\)
satisfies (E), \(Z\) is good. If \(K\) is also good, then
\(\mathfrak k\) is good, and Lemmas~5.2.1(iv) and 5.2.2(ii) show that
\(\mathfrak g/\mathfrak k\) and
\((\mathfrak g/\mathfrak k)^K\) are good. From the exact sequence
\[
 0\longrightarrow\mathfrak k\longrightarrow\mathfrak n
 \longrightarrow(\mathfrak g/\mathfrak k)^K\longrightarrow0
\]
and Lemma~5.2.1(iv), \(\mathfrak n\) is good. Finally, if \(G\) is
very good, so that \([x,x]=0\) identically for \(x\in\mathfrak g\),
then \(Z\) and \(N\) are plainly very good. This proves (iii)--(v).

\medskip
\noindent\SGAInlineRefTarget{sga3:ii:corollary:5-2-3-1}\textbf{Corollary 5.2.3.1.}
\textit{If the group law of
\(\underline{\operatorname{Lie}}(G/S)\) is commutative, then}
\[
 \underline{\operatorname{Lie}}
   \bigl(\underline{\operatorname{Centr}}(G)/S\bigr)
 =\underline{\operatorname{Lie}}(G/S)^G.
\]

\medskip
\noindent\SGAInlineRefTarget{sga3:ii:corollary:5-2-3-2}\textbf{Corollary 5.2.3.2.}
\textit{If the group law of
\(\underline{\operatorname{Lie}}(G/S)\) is commutative and \(K\) is a
normal subgroup of \(G\), then}
\[
 \left(
 \underline{\operatorname{Lie}}(G/S)/
 \underline{\operatorname{Lie}}(K/S)
 \right)^K
 =
 \underline{\operatorname{Lie}}(G/S)/
 \underline{\operatorname{Lie}}(K/S).
\]

\subsection*{\SGAInlineRefTarget{sga3:ii:subsection:5-3}5.3. Linear representations}

Let \(G\) be a good \(S\)-group acting linearly on a good
\(\mathbf O_S\)-module \(F\) through
\[
 \rho:G\longrightarrow
 \underline{\operatorname{Aut}}_{\mathbf O_S\text{-}\mathrm{mod}}(F).
\]
In \SGARefLink{sga3:ii:subsection:4-7-1}{4.7.1} and \SGARefLink{sga3:ii:corollary:4-8-2}{4.8.2} we defined the corresponding linear representation
\[
 \rho':
 \underline{\operatorname{Lie}}(G/S)
 \longrightarrow
 \underline{\operatorname{End}}_{\mathbf O_S\text{-}\mathrm{mod}}(F).
\]
The sub-\(S\)-groups
\(\underline{\operatorname{Norm}}_G(E)\) and
\(\underline{\operatorname{Centr}}_G(E)\) are defined for every subset
\(E\) of \(F\), and in particular for every
\(\mathbf O_S\)-submodule \(E\).

\medskip
\noindent\SGAInlineRefTarget{sga3:ii:definition:5-3-0}\textbf{Definition 5.3.0.}
For every \(S'\to S\), put
\[
\begin{split}
 \underline{\operatorname{Norm}}_{\underline{\operatorname{Lie}}(G/S)}
   (E)(S')
 &=\{X\in\underline{\operatorname{Lie}}(G/S)(S')
      \mid \rho'(X)E_{S'}\subset E_{S'}\},\\
 \underline{\operatorname{Centr}}_{\underline{\operatorname{Lie}}(G/S)}
   (E)(S')
 &=\{X\in\underline{\operatorname{Lie}}(G/S)(S')
      \mid \rho'(X)E_{S'}=0\}.
\end{split}
\]
This construction works for every linear representation of an
\(\mathbf O_S\)-Lie algebra in the sense of \SGARefLink{sga3:ii:scholium:4-9}{4.9}, and the two
subobjects obtained are \(\mathbf O_S\)-submodules stable under the
bracket.

\medskip
\noindent\SGAInlineRefTarget{sga3:ii:theorem:5-3-1}\textbf{Theorem 5.3.1.}
\textit{Let \(G\) be a good \(S\)-group acting linearly on a good
\(\mathbf O_S\)-module \(F\), and let \(E\) be an
\(\mathbf O_S\)-submodule of \(F\).}
\begin{enumerate}
\renewcommand{\labelenumi}{(\roman{enumi})}
\item \textit{One has}
\[
 \underline{\operatorname{Lie}}
   \bigl(\underline{\operatorname{Centr}}_G(E)/S\bigr)
 =
 \underline{\operatorname{Centr}}_
   {\underline{\operatorname{Lie}}(G/S)}(E),
\]
\textit{and \(\underline{\operatorname{Centr}}_G(E)\) is a good
\(S\)-group; it is very good if \(G\) is very good.}
\item \textit{Suppose that \(E\) is a good \(\mathbf O_S\)-module.
Then}\footnote{Editor's note: The original stated the following
equality without a hypothesis on \(E\); this has been corrected. See
\SGARefLink{sga3:ii:paragraph:6-5-1}{6.5.1} for counterexamples.}
\[
 \underline{\operatorname{Lie}}
   \bigl(\underline{\operatorname{Norm}}_G(E)/S\bigr)
 =
 \underline{\operatorname{Norm}}_
   {\underline{\operatorname{Lie}}(G/S)}(E),
\]
\textit{and \(\underline{\operatorname{Norm}}_G(E)\) is a good
\(S\)-group; it is very good if \(G\) is very good.}
\end{enumerate}
The proof is left to the reader.

\medskip
\noindent\SGAInlineRefTarget{sga3:ii:paragraph:5-3-2}\textbf{5.3.2.}
Let \(G\) be a good \(S\)-group. The preceding results apply in
particular to the adjoint representation. Let \(E\) be a good
\(\mathbf O_S\)-submodule of
\(\underline{\operatorname{Lie}}(G/S)\).\footnote{Editor's note: This
hypothesis has been added; cf. \SGARefLink{sga3:ii:paragraph:6-5-1}{6.5.1}.} It has a centralizer and a
normalizer in \(G\), and by \SGARefLink{sga3:ii:theorem:5-3-1}{5.3.1} their Lie algebras are the usual
centralizer and normalizer of \(E\), computed using the bracket:
\[
\begin{split}
 \underline{\operatorname{Centr}}_
   {\underline{\operatorname{Lie}}(G/S)}(E)(S')
 &=\{X\in\underline{\operatorname{Lie}}(G/S)(S')
       \mid[X,E_{S'}]=0\},\\
 \underline{\operatorname{Norm}}_
   {\underline{\operatorname{Lie}}(G/S)}(E)(S')
 &=\{X\in\underline{\operatorname{Lie}}(G/S)(S')
       \mid[X,E_{S'}]\subset E_{S'}\}.
\end{split}
\]

\medskip
\noindent\SGAInlineRefTarget{sga3:ii:paragraph:5-3-3}\textbf{5.3.3.}
Let \(K\) be a sub-\(S\)-group of \(G\). Then
\(\underline{\operatorname{Lie}}(K/S)\) is an
\(\mathbf O_S\)-submodule of
\(\underline{\operatorname{Lie}}(G/S)\). Suppose that it is a good
\(\mathbf O_S\)-module,\footnote{Editor's note: This hypothesis has
been added; cf. \SGARefLink{sga3:ii:paragraph:6-5-1}{6.5.1}.} as happens when \(K\) is a good \(S\)-group.
Plainly,
\[
 \underline{\operatorname{Centr}}_G(K)
 \subset
 \underline{\operatorname{Centr}}_G
   \bigl(\underline{\operatorname{Lie}}(K/S)\bigr),
\]
\[
 \underline{\operatorname{Norm}}_G(K)
 \subset
 \underline{\operatorname{Norm}}_G
   \bigl(\underline{\operatorname{Lie}}(K/S)\bigr).
\]
Hence, by \SGARefLink{sga3:ii:theorem:5-3-1}{5.3.1},
\[
 \underline{\operatorname{Lie}}
   \bigl(\underline{\operatorname{Centr}}_G(K)/S\bigr)
 \subset
 \underline{\operatorname{Centr}}_
   {\underline{\operatorname{Lie}}(G/S)}
   \bigl(\underline{\operatorname{Lie}}(K/S)\bigr),
\]
\[
 \underline{\operatorname{Lie}}
   \bigl(\underline{\operatorname{Norm}}_G(K)/S\bigr)
 \subset
 \underline{\operatorname{Norm}}_
   {\underline{\operatorname{Lie}}(G/S)}
   \bigl(\underline{\operatorname{Lie}}(K/S)\bigr).
\]
None of these four inclusions is an equality in general; many examples
will appear later. In particular, if \(K\) is a normal subgroup of
\(G\), then \(\underline{\operatorname{Lie}}(K/S)\) is an ideal of
\(\underline{\operatorname{Lie}}(G/S)\).

% Source page: Exp2 component-local p. 50; combined-reader p. 108.
\section*{\SGAInlineRefTarget{sga3:ii:section:6}6. Miscellaneous remarks}
\addcontentsline{toc}{section}{6. Miscellaneous remarks}

\subsection*{\SGAInlineRefTarget{sga3:ii:subsection:6-1}6.1.}
The bracket of two infinitesimal automorphisms can be defined for an
\(S\)-functor \(X\) which is not necessarily a group: apply the
results of this expos\'e to \(\underline{\operatorname{Aut}}_S(X)\).
To obtain a convenient formalism, one is led to assume that \(X\) is
good, meaning that the \(\mathbf O_X\)-module \(T_{X/S}\) is good. If
\(X\) is an \(S\)-group, this agrees with \SGARefLink{sga3:ii:definition:4-6}{Definition~4.6}.

\subsection*{\SGAInlineRefTarget{sga3:ii:subsection:6-2}6.2.}
There exist functors with infinitesimal endomorphisms that are not
automorphisms and hence, a fortiori, do not satisfy condition (E).
\footnote{Editor's note: The original has been expanded here. In
particular, O. Gabber pointed out the correct definition of the free
abelian monoid on a pointed set.}
For a pointed set \((E,x_0)\), let \(M_A(E)\) be the free abelian
monoid generated by \(E\), and let \(M_{A,P}(E,x_0)\) be the quotient
of \(M_A(E)\) by the equivalence relation generated by
\(m\sim x_0+m\). The assignment
\[
 (E,x_0)\longmapsto M_{A,P}(E,x_0)
\]
is left adjoint to the forgetful functor from abelian monoids to
pointed sets; \(M_{A,P}(E,x_0)\) is the free abelian monoid on the
pointed set \((E,x_0)\).

Now let \(X\) be the functor assigning to every scheme \(S\) the free
abelian monoid on \(\mathbf O(S)\), pointed by zero. Every morphism
\[
 f:S\longrightarrow I_{\mathbb Z}
   =\operatorname{Spec}(\mathbb Z[\varepsilon])
\]
corresponds to a square-zero element \(u_f\in\mathbf O(S)\), and hence
defines an endomorphism of \(X(S)\) by
\(x\mapsto x+u_f\), where the sum is taken in
\(M_{A,P}(\mathbf O(S),0)\). Thus we obtain an endomorphism
\(\phi\) of
\[
 X_{I_{\mathbb Z}}=X\times_{\mathbb Z}I_{\mathbb Z}
\]
defined, for \(f\in I_{\mathbb Z}(S)\) and \(x\in X(S)\), by
\[
 \phi((x,f))=(x+u_f,f).
\]
If \(f_0:S\to I_{\mathbb Z}\) is the composite of the structural
morphism \(S\to\operatorname{Spec}(\mathbb Z)\) and the zero section,
then \(u_{f_0}=0\), so
\(\phi((x,f_0))=(x,f_0)\). Thus \(\phi\) induces the identity on
\(X\): it is an infinitesimal endomorphism of \(X\) which is plainly
not an automorphism.

\subsection*{\SGAInlineRefTarget{sga3:ii:subsection:6-3}6.3.}
There exist modules which are not good. First, the preceding
counterexample can be modified slightly.\footnote{Editor's note: The
following observation has been added.} Let \(F(S)\) be the free
\(\mathbf O(S)\)-module with basis the elements of \(\mathbf O(S)\).
Then \(F\) does not satisfy (E) relative to
\(\operatorname{Spec}(\mathbb Z)\). Moreover, let \(G\) be the
\(\mathbb Z\)-group defined by
\[
 G(S)=\operatorname{GL}_n(R(S)),\qquad n\geq2,
\]
where \(R(S)\) is the \(\mathbf O(S)\)-algebra of the abelian group
\(\mathbf O(S)\). Then the \(\mathbb Z\)-group
\(\underline{\operatorname{Lie}}(G/\mathbb Z)\) is not commutative.

There are also the following counterexamples.\footnote{Editor's note:
Example (a) has been expanded and example (b) has been added.} Let
\(S=\operatorname{Spec}(k)\), where \(k\) is a field of characteristic
\(p>0\).

\medskip
\noindent\textit{(a)}
Let \(F\) be the \(\mathbf O_S\)-module assigning to every \(S\)-scheme
\(T\) the group
\[
 F(T)=\Gamma(T,\mathscr O_T),
\]
with its \(\mathbf O(T)\)-module structure defined by
\(r\cdot f=r^pf\). As an \(S\)-functor of groups, \(F\) is isomorphic
to \(\mathbb G_{a,S}\). Thus \(F/S\) satisfies (E), and
\[
 \underline{\operatorname{Lie}}(F/S)
 \simeq
 \underline{\operatorname{Lie}}(\mathbb G_{a,S}/S)
 \simeq\mathbf O_S.
\]
For each \(T\to S\), the canonical morphism
\(F\to\underline{\operatorname{Lie}}(F/S)\) is the identity map
\(F(T)\to\mathbf O(T)\). It respects the abelian-group structure but
not the module structure, so \(F\) is not good (4.4.1).

\medskip
\noindent\textit{(b)}
Let \(\alpha_{p,k}\) be the \(k\)-functor of groups defined by
\[
 \alpha_{p,k}(T)=\{x\in\mathbf O(T)\mid x^p=0\}.
\]
It is represented by
\(\operatorname{Spec}(k[X]/(X^p))\), hence is a very good \(S\)-group.
It also carries an \(\mathbf O_S\)-module structure, but is not a good
\(\mathbf O_S\)-module, because the canonical morphism
\[
 \alpha_{p,k}\longrightarrow
 \underline{\operatorname{Lie}}(\alpha_{p,k}/k)=\mathbb G_{a,k}
\]
is not bijective.

\subsection*{\SGAInlineRefTarget{sga3:ii:subsection:6-4}6.4.}
For a group-valued functor \(G\) over \(S\), the definitions give
\[
 (G/S\text{ satisfies (E)})
 \Longleftarrow(G\text{ is good})
 \Longleftarrow(G\text{ is very good}).
\]
It would be interesting to prove or disprove the reverse implications.
\footnote{Editor's note: A representable \(G\) is very good (\SGARefLink{sga3:ii:subsection:4-11}{4.11}).
For representability criteria see, for example, [MO67]. For
automorphisms of an affine algebraic group over an algebraically closed
field of characteristic zero, see [HM69].}

\subsection*{\SGAInlineRefTarget{sga3:ii:subsection:6-5}6.5.}
Let \(\operatorname{Nil}\) be the \(\mathbb Z\)-functor of groups
defined by
\[
 \operatorname{Nil}(S)
 =\{x\in\mathbf O(S)\mid
       \text{there exists }n\in\mathbb N\text{ with }x^n=0\}.
\]
\footnote{Editor's note: This paragraph has been added.}
The functor \(\operatorname{Nil}\) is very good but not representable.
Let \(\operatorname{Nil}^2\), \(\mathbf O_{\mathrm{red}}\), and \(F\)
be the \(\mathbb Z\)-functors of groups assigning to \(S\),
respectively, the ideal \(\operatorname{Nil}(S)^2\) and the groups
\[
 \mathbf O_{\mathrm{red}}(S)
   =\mathbf O(S)/\operatorname{Nil}(S),
 \qquad
 F(S)=\mathbf O(S)/\operatorname{Nil}(S)^2.
\]
It is easy to see that
\[
 \underline{\operatorname{Lie}}
   (\mathbf O_{\mathrm{red}}/\mathbb Z)=0,
\]
so the \(\mathbf O_{\mathbb Z}\)-module
\(\mathbf O_{\mathrm{red}}\) is not good, although it is a good
\(\mathbb Z\)-group. If \(\mathscr M,\mathscr N\) are free
\(\mathbb Z\)-modules of finite rank, then
\[
\begin{split}
 \operatorname{Nil}^2(I_S(\mathscr M\oplus\mathscr N))
 \simeq{}&
 \operatorname{Nil}^2(S)
 \oplus\operatorname{Nil}(S)\otimes_{\mathbb Z}\mathscr M\\
 &\oplus\operatorname{Nil}(S)\otimes_{\mathbb Z}\mathscr N,
\end{split}
\]
and hence
\[
\begin{split}
 F(I_S(\mathscr M\oplus\mathscr N))
 \simeq{}&F(S)
 \oplus\mathbf O_{\mathrm{red}}(S)\otimes_{\mathbb Z}\mathscr M\\
 &\oplus\mathbf O_{\mathrm{red}}(S)\otimes_{\mathbb Z}\mathscr N.
\end{split}
\]
Thus \(F\) satisfies condition (E), while
\[
 \underline{\operatorname{Lie}}(F/\mathbb Z)
   =\mathbf O_{\mathrm{red}}.
\]
Since \(\mathbf O_{\mathrm{red}}\) is not a good
\(\mathbf O_{\mathbb Z}\)-module, \(F\) is a \(\mathbb Z\)-group
which satisfies (E) but is not good.

\medskip
\noindent\SGAInlineRefTarget{sga3:ii:paragraph:6-5-1}\textbf{6.5.1.}
We now give the counterexamples announced in \SGARefLink{sga3:ii:theorem:5-3-1}{5.3.1}--\SGARefLink{sga3:ii:paragraph:5-3-3}{5.3.3}. Let \(S\)
be a scheme, let \(F=\mathbf O_S^{\oplus2}\) with the natural action
of the good \(S\)-group \(G=\operatorname{GL}_{2,S}\), and let \(E\)
be the \(\mathbf O_S\)-submodule of \(F\) formed by the pairs
\((x_1,x_2)\) for which \(x_2\) is nilpotent. Put
\[
 N=\underline{\operatorname{Norm}}_G(E).
\]
Then
\[
 \underline{\operatorname{Lie}}(N/S)
   =\underline{\operatorname{Lie}}(G/S),
\]
whereas, for every \(S'\to S\),
\[
 \underline{\operatorname{Norm}}_
   {\underline{\operatorname{Lie}}(G/S)}(E)(S')
 =
 \left\{
 \begin{pmatrix}a&b\\x&c\end{pmatrix}
 \,\middle|\,
 a,b,c,x\in\mathbf O(S'),\ x\text{ nilpotent}
 \right\}.
\]
Consequently,
\[
 \underline{\operatorname{Lie}}
   \bigl(\underline{\operatorname{Norm}}_G(E)/S\bigr)
 \neq
 \underline{\operatorname{Norm}}_
   {\underline{\operatorname{Lie}}(G/S)}(E).
\]

Taking the semidirect product \(G_1=F\rtimes G\) gives an analogous
counterexample in which \(E\) is an \(\mathbf O_S\)-submodule of
\(\underline{\operatorname{Lie}}(G_1/S)\). Moreover, with the notation
above,
\[
 E=\underline{\operatorname{Lie}}(K/S),
\]
where \(K\) is the subgroup \(\mathbf O_S\oplus\operatorname{Nil}^2\)
of \(F\); that is, \(K(S')\) consists of the pairs \((x_1,x_2)\) with
\(x_2\in\operatorname{Nil}(S')^2\).

\section*{Bibliography}
\addcontentsline{toc}{section}{Bibliography}
\footnote{Editor's note: Additional references cited in this expos\'e.}

\begin{description}
\item[{[DG70]}] M. Demazure and P. Gabriel,
  \emph{Groupes Alg\'ebriques}, Masson \& North-Holland, 1970.
\item[{[HM69]}] G. Hochschild and G. D. Mostow,
  \emph{Automorphisms of affine algebraic groups},
  \emph{J. Algebra} \textbf{13} (1969), 535--543.
\item[{[MO67]}] H. Matsumura and F. Oort,
  \emph{Representability of group functors and automorphisms of
  algebraic schemes}, \emph{Invent. Math.} \textbf{4} (1967), 1--25.
\end{description}

% End of Expose II.

% END INCLUDED SGA3_English_current_buildable_source_20260823/source/inputs/II/tex/components/14_expose_II_s5_s6_examples_misc.tex



% END INCLUDED SGA3_English_current_buildable_source_20260823/source/inputs/II/body.tex


\SGAStartExpose{III}
% BEGIN INCLUDED SGA3_English_current_buildable_source_20260823/source/inputs/III/body.tex
% Mechanically extracted after \mainmatter from the frozen standalone master.

% BEGIN INCLUDED SGA3_English_current_buildable_source_20260823/source/inputs/III/tex/components/15_expose_III_opening_s016.tex
% SGA 3, Expose III, Loop 1 English reconstruction.
% Scope: opening through Remark 0.1.6, stopping before Reminders 0.1.7.
% Authority: Polo--Gille Exp3-14oct24.pdf, component-local pages 1--5
% (combined-reader pages 111--115). Force-OCR is a drafting locator; English
% wording was compared with Jacob Reinhold's CC BY 4.0 Markdown lineage.
% The PDF controls text, formulas, numbering, notes, and diagram images.

\setcounter{footnote}{-1}
\chapter*{Expos\'e III\\Infinitesimal Extensions}
\addcontentsline{toc}{chapter}{Expos\'e III. Infinitesimal Extensions}
\markboth{EXPOS\'E III. INFINITESIMAL EXTENSIONS}%
  {EXPOS\'E III. INFINITESIMAL EXTENSIONS}

\begin{center}
  \textit{by M. Demazure}
\end{center}

% Source page: Exp3 component-local p. 1; combined-reader p. 111.
In this expos\'e we work in the following general situation. We have a
scheme \(S\) and a coherent nilpotent ideal \(\mathscr I\) on \(S\).
We denote by \(S_n\) the closed subscheme of \(S\) defined by the ideal
\(\mathscr I^{\,n+1}\), for \(n\geq 0\). In particular, \(S_0\) is
defined by \(\mathscr I\). Since \(\mathscr I\) is nilpotent, \(S_n=S\)
for all sufficiently large \(n\), and the \(S_i\) have the same
underlying topological space. A typical example is the following:
\(S\) is the spectrum of a local Artinian ring \(A\), and \(\mathscr I\)
is the ideal defined by the radical of \(A\), so that \(S_0\) is the
spectrum of the residue field of \(A\).\footnote{Editor's note: version
of October 14, 2024.}

In this situation one is given certain data over \(S_0\) and seeks data
over \(S\) that lift them, that is, data whose base change from \(S\) to
\(S_0\) recovers the original data. This is done step by step, through
the schemes \(S_n\). At each step we shall define the obstructions that
arise and classify the resulting solutions when they exist.

The passage from \(S_n\) to \(S_{n+1}\) may be generalized as follows.
We have a scheme \(S\) and two ideals \(\mathscr I\supset\mathscr J\)
such that \(\mathscr I\mathscr J=0\). In the preceding case \(S\),
\(\mathscr I\), and \(\mathscr J\) are respectively
\[
 S_{n+1},\qquad
 \mathscr I/\mathscr I^{\,n+2},\qquad
 \mathscr I^{\,n+1}/\mathscr I^{\,n+2}.
\]
We denote by \(S_0\), respectively \(S_{\mathscr J}\), the closed
subscheme of \(S\) defined by \(\mathscr I\), respectively
\(\mathscr J\), and pose a problem of extending data from
\(S_{\mathscr J}\) to \(S\).

% Source page: Exp3 component-local p. 2; combined-reader p. 112.
SGA~1, Expos\'e III, treated the problems of extending morphisms of
schemes and of extending schemes. Here we consider the problems of
extending group morphisms, group structures, and subgroups.

In \S0 we have collected those results of SGA~1, Expos\'e III, that
will be useful to us, putting them into the form best suited to our
purpose and sparing the reader constant reference to that
expos\'e.\footnote{Editor's note: the reader may prefer to begin with
\S1, which is easier and may serve as motivation and a guide to the
results of \S0.} \SGARefLink{sga3:iii:section:1}{Section~1} collects computations in group cohomology
that will be used later and have nothing to do with scheme theory.
\SGARefLink{sga3:iii:section:2}{Sections~2} and~3 treat, respectively, the extension of group morphisms
and the extension of group structures. In \S4 we briefly recall the
proof of a result stated in TDTE~IV concerning the extension of
subschemes, and apply it to the problem of extending subgroups. For the
remainder of the Seminar, only the result of \S2 concerning the
extension of group morphisms is indispensable.\footnote{Editor's note:
nevertheless, \SGARefLink{sga3:iii:corollary:3-10}{3.10} is used in \SGARefLink{sga3:xxiv:1.13}{XXIV~1.13}. Moreover, \SGARefLink{sga3:iii:corollary:4-37}{4.37} is used in the
proof of \SGARefLink{sga3:ix:theorem:3:2}{IX~3.2}~bis and \SGARefLink{sga3:iii:remark:3-6}{3.6}~bis, although those results also follow from
results of Expos\'e~X that do not use \SGARefLink{sga3:iii:corollary:4-37}{4.37}.}

The idea of reducing problems of infinitesimal extension to the usual
cohomological computations in group extensions was suggested by
J.~Giraud during the oral expos\'e, whose computations were appreciably
more complicated and less transparent. Unfortunately, this method seems
to work well only for the first two problems studied, and we have been
unable to avoid rather laborious computations for extensions of
subgroups.

To simplify the language, we shall call a \(Y\)-functor, respectively a
\(Y\)-scheme, and so on, a functor, respectively a scheme, and so on,
equipped with a morphism to the functor \(Y\). This extends the
definitions of Expos\'e~I, which concerned only representable \(Y\).

\section*{\SGAInlineRefTarget{sga3:iii:section:0}0. Reminders from SGA 1, Expos\'e III, and miscellaneous remarks}
\addcontentsline{toc}{section}{0. Reminders from SGA 1, Expos\'e III}

% Source page: Exp3 component-local p. 3; combined-reader p. 113.
We begin with a general definition.

\medskip
\SGARefTarget{sga3:iii:definition:0:1}\noindent\textbf{Definition 0.1.}
\textit{Let \(\mathscr C\) be a category, \(X\) an object of
\(\widehat{\mathscr C}\), and \(G\) a
\(\widehat{\mathscr C}\)-group acting on \(X\). We say that \(X\) is
\emph{formally principal homogeneous}\footnote{Editor's note: one also
calls such an object a \emph{pseudo-torsor}; see
EGA~IV\textsubscript{4}, 16.5.15. The more general notion of a formally
homogeneous object, not necessarily principal homogeneous, is defined
in \SGARefLink{sga3:iv:definition:6:7:1}{Expos\'e~IV, 6.7.1}.} under \(G\) if the following equivalent
conditions hold:}
\begin{enumerate}
\item[(i)] \textit{for every object \(T\) of \(\mathscr C\), the set
  \(X(T)\) is either empty or principal homogeneous under \(G(T)\);}
\item[(ii)] \textit{the morphism of functors
  \[
    G\times X\longrightarrow X\times X,
    \qquad (g,x)\longmapsto(gx,x),
  \]
  is an isomorphism.}
\end{enumerate}

We shall now put the results of SGA~1, III~\S5\footnote{Editor's note:
see also EGA~IV\textsubscript{4}, 16.5.14--18.} into the form most
useful here. Throughout this section we use the following notation. On a
scheme \(S\) we have two quasi-coherent ideals \(\mathscr I\) and
\(\mathscr J\) such that
\[
 \mathscr I\supset\mathscr J,
 \qquad
 \mathscr I\mathscr J=0.
\]
In particular, \(\mathscr J^2=0\). We denote by \(S_0\), respectively
\(S_{\mathscr J}\), the closed subscheme of \(S\) defined by
\(\mathscr I\), respectively \(\mathscr J\). For every \(S\)-functor
\(X\), we systematically denote by \(X_0\) and \(X_{\mathscr J}\) the
functors obtained by base change from \(S\) to \(S_0\) and
\(S_{\mathscr J}\). We use the same notation for a morphism.

% Source page: Exp3 component-local p. 4; combined-reader p. 114.
\medskip
\noindent\SGAInlineRefTarget{sga3:iii:definition:0-1-1}\textbf{Definition 0.1.1.}\footnote{Editor's note: the
numbering \SGARefLink{sga3:iii:definition:0-1-1}{0.1.1} through \SGARefLink{sga3:iii:corollary:0-1-13}{0.1.13} has been added to make the definitions
and results contained here more visible.}
\textit{Let \(X\) be an \(S\)-functor. Define a functor \(X^+\) over
\(S\) by}
\[
 \operatorname{Hom}_S(Y,X^+)
 =\operatorname{Hom}_{S_{\mathscr J}}
    (Y_{\mathscr J},X_{\mathscr J})
 =\operatorname{Hom}_S(Y_{\mathscr J},X)
\]
\textit{for a variable \(S\)-scheme \(Y\).} In the notation of
Expos\'e~II, \S1,
\[
 X^+\simeq\prod_{S_{\mathscr J}/S}X_{\mathscr J}.
\]
The identity morphism of \(X_{\mathscr J}\) defines, by construction, an
\(S\)-morphism
\[
 p_X:X\longrightarrow X^+.
\]
Explicitly,\footnote{Editor's note: the following sentence has been
added, and the next two remarks have been expanded.} for every
\(S\)-scheme \(Y\), the map
\[
 p_X(Y):
 \operatorname{Hom}_S(Y,X)\longrightarrow
 \operatorname{Hom}_S(Y,X^+)
 =\operatorname{Hom}_S(Y_{\mathscr J},X)
\]
is induced by the morphism \(Y_{\mathscr J}\to Y\).

\medskip
\noindent\SGAInlineRefTarget{sga3:iii:remark:0-1-2}\textbf{Remark 0.1.2.}
If \(X\) is an \(S\)-group functor, then \(X_{\mathscr J}\) is an
\(S_{\mathscr J}\)-group functor. The defining formula for \(X^+\)
therefore equips \(X^+\) with an \(S\)-group-functor structure, and the
description above shows that
\[
 p_X:X\longrightarrow X^+
\]
is a morphism of \(S\)-group functors.

\medskip
\noindent\SGAInlineRefTarget{sga3:iii:remark:0-1-3}\textbf{Remark 0.1.3.}
For every \(S\)-group functor \(Y\), one has
\[
 \operatorname{Hom}_{S\text{-}\mathrm{gr}}(Y,X^+)
 =
 \operatorname{Hom}_{S_{\mathscr J}\text{-}\mathrm{gr}}
   (Y_{\mathscr J},X_{\mathscr J}).
\]
Indeed, let \(f\in\operatorname{Hom}_S(Y,X^+)\) correspond to
\(f_{\mathscr J}\in
\operatorname{Hom}_{S_{\mathscr J}}
(Y_{\mathscr J},X_{\mathscr J})\). The condition that \(f\) be an
\(S\)-group morphism is that, for every \(T\to S\) and every
\(y,y'\in Y(T)\),
\[
 f(y\cdot y')=f(y)\cdot f(y'),
\]
which is equivalent to
\[
 f_{\mathscr J}(y_{\mathscr J})
 \cdot f_{\mathscr J}(y'_{\mathscr J})
 =f_{\mathscr J}\bigl((y\cdot y')_{\mathscr J}\bigr).
\]
Since
\((y\cdot y')_{\mathscr J}
=y_{\mathscr J}\cdot y'_{\mathscr J}\), this is precisely the
condition that \(f_{\mathscr J}\) be a group morphism. Applying this to
\(Y=X\), we again see that \(p_X\), which corresponds to
\(\operatorname{id}_{X_{\mathscr J}}\), is a morphism of
\(S\)-group functors.

Return now to the general case and suppose that \(X\) is an \(S\)-scheme.
An \(S\)-morphism from a variable \(S\)-scheme \(Y\) to \(X^+\) is, by
definition, an \(S_{\mathscr J}\)-morphism
\(g_{\mathscr J}:Y_{\mathscr J}\to X_{\mathscr J}\). We define an
\(X^+\)-functor in abelian groups \(L_X\) as follows.

\medskip
\noindent\SGAInlineRefTarget{sga3:iii:scholium:0-1-4}\textbf{Scholium 0.1.4.}\footnote{Editor's note: this scholium
has been added.}
If \(\pi:Y\to S\) is a morphism of schemes and \(M\) is an
\(\mathscr O_S\)-module, we denote the inverse image \(\pi^*M\) by
\[
 M\otimes_{\mathscr O_S}\mathscr O_Y.
\]
If \(\mathscr J\) is an ideal of \(\mathscr O_S\), we denote by
\(\mathscr J\mathscr O_Y\) the ideal of \(\mathscr O_Y\) that is the
image of
\[
 \pi^*(\mathscr J)
 =\mathscr J\otimes_{\mathscr O_S}\mathscr O_Y
 \longrightarrow\mathscr O_Y.
\]
For every morphism of \(S\)-schemes \(f:Z\to Y\), there is an
epimorphism of \(\mathscr O_Z\)-modules
\begin{equation}
 f^*(\mathscr J\mathscr O_Y)\longrightarrow
 \mathscr J\mathscr O_Z.
 \tag{0.1.4}\label{sga3:iii:scholium:0-1-4:equation:0-1-4}
\end{equation}
It follows from the commutative diagram below.

\SGAThreeNativeDiagram[0.70\textwidth]
  {native_diagrams/exp3/diagram_001_exp3_p004_scholie014_tensor_diagram.tex}
  {The tensor-product diagram yielding the epimorphism (\SGARefLink{sga3:iii:scholium:0-1-4:equation:0-1-4}{0.1.4}).}
  {fig:exp3-scholie-014-tensor}
  {Polo--Gille Exp3-14oct24.pdf, component-local p. 4; combined-reader p. 114}

\medskip
\noindent\SGAInlineRefTarget{sga3:iii:definition:0-1-5}\textbf{Definition 0.1.5.}\footnote{Editor's note: introducing
the functor \(L_X\) gives an extension, functorial in \(Y\) and in
\(X\), of SGA~1, III, Proposition~5.1; see \SGARefLink{sga3:iii:theorem:0-1-8}{0.1.8} and \SGARefLink{sga3:iii:corollary:0-1-9}{0.1.9} below, as
well as \SGARefLink{sga3:iii:corollary:0-1-10}{0.1.10} and \SGARefLink{sga3:iii:proposition:0-2}{0.2}. Moreover, when \(X\) is an \(S\)-group functor,
\(L_X\) gives rise, under certain hypotheses, to an \(S\)-group functor
\(L'_X\) and to an exact sequence
\(1\to L'_X\to X\to X^+\); see \SGARefLink{sga3:iii:remark:0-4}{0.4}, \SGARefLink{sga3:iii:corollary:0-9}{0.9}, and \SGARefLink{sga3:iii:proposition:0-11}{0.11}. These play an
essential role in this expos\'e; see \SGARefLink{sga3:iii:theorem:2-1}{Theorems~2.1}, \SGARefLink{sga3:iii:theorem:3-5}{3.5}, and \SGARefLink{sga3:iii:theorem:4-21}{4.21}.}
Let \(X\) be an \(S\)-scheme. For every \(X^+\)-scheme \(Y\), given by
a morphism
\[
 g_{\mathscr J}:Y_{\mathscr J}\longrightarrow X_{\mathscr J},
\]
set
\[
 \operatorname{Hom}_{X^+}(Y,L_X)
 =
 \operatorname{Hom}_{\mathscr O_{Y_0}}
 \bigl(g_0^*(\Omega^1_{X_0/S_0}),\mathscr J\mathscr O_Y\bigr),
\]
where \(\Omega^1_{X_0/S_0}\) is the module of relative differentials of
\(X_0\) over \(S_0\), and \(\mathscr J\mathscr O_Y\) is regarded as an
\(\mathscr O_{Y_0}\)-module because it is annihilated by
\(\mathscr I\).

Then \(L_X\) is an \(X^+\)-functor in abelian
groups.\footnote{Editor's note: what follows has been added.} Indeed,
for every \(X^+\)-morphism \(f:Z\to Y\), the functor \(f_0^*\) and the
morphism
\[
 f_0^*(\mathscr J\mathscr O_Y)\longrightarrow
 \mathscr J\mathscr O_Z
\]
of (\SGARefLink{sga3:iii:scholium:0-1-4}{0.1.4}) induce a natural homomorphism \(L_X(f)\):
\[
\begin{aligned}
 &\operatorname{Hom}_{\mathscr O_{Y_0}}
   \bigl(g_0^*(\Omega^1_{X_0/S_0}),\mathscr J\mathscr O_Y\bigr)\\
 &\quad\longrightarrow
 \operatorname{Hom}_{\mathscr O_{Z_0}}
   \bigl(f_0^*g_0^*(\Omega^1_{X_0/S_0}),
         f_0^*(\mathscr J\mathscr O_Y)\bigr)\\
 &\quad\longrightarrow
 \operatorname{Hom}_{\mathscr O_{Z_0}}
   \bigl(f_0^*g_0^*(\Omega^1_{X_0/S_0}),
         \mathscr J\mathscr O_Z\bigr).
\end{aligned}
\]

Finally, \(L_X(f)\) has the following simple local description. The
schemes \(Y\) and \(Y_{\mathscr J}\) have the same underlying
topological space, as do \(V\) and \(V_{\mathscr J}\) for every open
subset \(V\) of \(Y\). Let
\[
 U=\operatorname{Spec}(A),\qquad
 V=\operatorname{Spec}(B),\qquad
 W=\operatorname{Spec}(C)
\]
be affine opens such that \(U\) lies over
\(\operatorname{Spec}(\Lambda)\subset S\),
\(g_{\mathscr J}(V_{\mathscr J})\subset U\), and
\(W\subset f^{-1}(V)\). Write \(I,J\) for the ideals of \(\Lambda\)
corresponding to \(\mathscr I,\mathscr J\). Then \(f\), respectively
\(g_{\mathscr J}\), induces a morphism of \(\Lambda\)-algebras
\[
 \theta:B\longrightarrow C,
 \qquad\text{respectively}\qquad
 \varphi:A\longrightarrow B/JB,
\]
and clearly \(\theta(JB)\subset JC\). On the other hand,
\[
 m\in
 \operatorname{Hom}_{\mathscr O_{Y_0}}
 \bigl(g_0^*(\Omega^1_{X/S}),\mathscr J\mathscr O_Y\bigr)
\]
induces an element \(D\) of
\[
\begin{aligned}
 \operatorname{Hom}_{\mathscr O_{V_0}}
 \bigl(g_0^*(\Omega^1_{U/S}),\mathscr J\mathscr O_V\bigr)
 &=
 \operatorname{Hom}_{B/IB}
 \bigl(\Omega^1_{A/\Lambda}\otimes_A B/IB,JB\bigr)\\
 &=\operatorname{Der}_{\Lambda}(A,JB),
\end{aligned}
\]
and the image of \(L_X(f)(m)\) in
\[
\begin{aligned}
 \operatorname{Hom}_{\mathscr O_{W_0}}
 \bigl(f_0^*g_0^*(\Omega^1_{U/S}),\mathscr J\mathscr O_Z\bigr)
 &=
 \operatorname{Hom}_{C/IC}
 \bigl(\Omega^1_{A/\Lambda}\otimes_A C/IC,JC\bigr)\\
 &=\operatorname{Der}_{\Lambda}(A,JC)
\end{aligned}
\]
is simply \(\theta\circ D\).

% Source page: Exp3 component-local p. 5; combined-reader p. 115.
\medskip
\noindent\SGAInlineRefTarget{sga3:iii:remark:0-1-6}\textbf{Remark 0.1.6.}\footnote{Editor's note: this remark has
been expanded.}
Let \(f:X\to W\) be an \(S\)-morphism. It induces an \(S\)-morphism
\[
 f^+:X^+\longrightarrow W^+
\]
defined as follows. If \(g\in\operatorname{Hom}_S(Y,X^+)\)
corresponds to an \(S\)-morphism
\(g_{\mathscr J}:Y_{\mathscr J}\to X\), then \(f^+(g)\) is the
element \(f\circ g_{\mathscr J}\) of
\(\operatorname{Hom}_S(Y,W^+)\). The following diagram is commutative.

\SGAThreeNativeDiagram[0.44\textwidth]
  {native_diagrams/exp3/diagram_002_exp3_p005_remark016_naturality_square.tex}
  {Naturality of the morphisms \(p_X\) and \(p_W\).}
  {fig:exp3-remark-016-naturality}
  {Polo--Gille Exp3-14oct24.pdf, component-local p. 5; combined-reader p. 115}

% Continue with Reminders 0.1.7.

% END INCLUDED SGA3_English_current_buildable_source_20260823/source/inputs/III/tex/components/15_expose_III_opening_s016.tex

% BEGIN INCLUDED SGA3_English_current_buildable_source_20260823/source/inputs/III/tex/components/16_expose_III_s017_s0110_functorial_liftings.tex
% SGA 3, Expose III, Loop 1 English reconstruction.
% Scope: Reminders 0.1.7 through Corollary 0.1.10, beginning immediately
% after Remark 0.1.6 and stopping immediately before Lemma 0.1.11.
% Authority: Polo--Gille Exp3-14oct24.pdf, component-local pages 5--12
% (combined-reader pages 115--122). Force-OCR is a drafting locator; English
% wording was compared with Jacob Reinhold's CC BY 4.0 Markdown lineage.
% The PDF controls text, formulas, numbering, notes, and all diagram images.

% Source page: Exp3 component-local p. 5; combined-reader p. 115.
\medskip
\noindent\SGAInlineRefTarget{sga3:iii:paragraph:0-1-7}\textbf{Reminders 0.1.7.}\footnote{Editor's note: This
paragraph has been added (see also [DG70], \S I.4, nos.~1--2).}
In this paragraph, given an \(S\)-scheme \(X\), we ``recall'' certain
functorial properties of the module of differentials
\(\Omega^1_{X/S}\) and of the first infinitesimal neighborhood of the
diagonal, \(\Delta^{(1)}_{X/S}\). These properties follow readily from
EGA~IV\textsubscript{4}, \S\S16.1--16.4, but do not appear there
explicitly.

\smallskip
\noindent\textit{(a)} Let us begin by recalling the following facts
(cf. EGA~II, \S\S1.2--1.5). Let \(g:Y\to X\) be a morphism of schemes,
let \(\pi:X'\to X\) be an affine \(X\)-scheme, and let
\(\mathscr B=\pi_*(\mathscr O_{X'})\) be the corresponding
quasi-coherent \(\mathscr O_X\)-algebra. Then the \(Y\)-scheme
\(Y\times_X X'\) is affine and corresponds to the quasi-coherent
\(\mathscr O_Y\)-algebra \(g^*(\mathscr B)\), and there is a
commutative diagram of bijections:

\SGAThreeNativeDiagram[0.82\textwidth]
  {native_diagrams/exp3/diagram_003_exp3_p005_reminders017_affine_bijections.tex}
  {The affine-scheme/algebra bijections after base change by \(g\).}
  {fig:exp3-reminders-017-affine-bijections}
  {Polo--Gille Exp3-14oct24.pdf, component-local p. 5; combined-reader p. 115}

Moreover, these bijections are functorial in the pair \((X,X')\). More
precisely, let \(W'\) be an affine \(W\)-scheme corresponding to the
quasi-coherent \(\mathscr O_W\)-algebra \(\mathscr A\), and suppose that
one has a commutative diagram of morphisms of schemes as follows.

\SGAThreeNativeDiagram[0.38\textwidth]
  {native_diagrams/exp3/diagram_004_exp3_p005_reminders017_scheme_functoriality.tex}
  {The scheme morphisms used to express functoriality in \((X,X')\).}
  {fig:exp3-reminders-017-scheme-functoriality}
  {Polo--Gille Exp3-14oct24.pdf, component-local p. 5; combined-reader p. 115}

Denote by
\[
 \varphi:\mathscr A\longrightarrow f_*(\mathscr B),
 \qquad
 \varphi^\sharp:f^*(\mathscr A)\longrightarrow\mathscr B
\]
the algebra morphisms associated with \(f'\), and by
\[
 \theta:\mathscr B\longrightarrow g_*(\mathscr O_Y),
 \qquad
 \theta^\sharp:g^*(\mathscr B)\longrightarrow\mathscr O_Y
\]
those associated with a variable \(X\)-morphism \(g':Y\to X'\).
Then the following diagram is commutative, where \(Y\) is viewed as a
\(W\)-scheme through \(f\circ g\).

% Source page: Exp3 component-local p. 6; combined-reader p. 116.
\SGAThreeNativeDiagram[0.96\textwidth]
  {native_diagrams/exp3/diagram_005_exp3_p006_reminders017_hom_functoriality.tex}
  {Functoriality of the affine-scheme/algebra correspondences.}
  {fig:exp3-reminders-017-hom-functoriality}
  {Polo--Gille Exp3-14oct24.pdf, component-local p. 6; combined-reader p. 116}

\smallskip
\noindent\SGAInlineRefTarget{sga3:iii:paragraph:0-1-7:item:b}\textit{(b)} Let now \(X\) be an \(S\)-scheme. Let
\(\Omega^1_{X/S}\) be the module of differentials of \(X\) over \(S\),
and let \(\Delta^{(1)}_{X/S}\) be the first infinitesimal neighborhood
of the diagonal immersion
\(\delta_X:X\to X\times_S X\). It is a subscheme of
\(X\times_S X\), in which the diagonal \(\Delta_{X/S}\) is a closed
subscheme. Write \(\operatorname{pr}^i_X\), \(i=1,2\), for the two
projections \(X\times_S X\to X\), and write \(\pi_X\) for the
restriction of \(\operatorname{pr}^1_X\) to
\(\Delta^{(1)}_{X/S}\).

On the one hand, every morphism \(f:X\to W\) of \(S\)-schemes induces
an \(S\)-morphism
\[
 \Delta^{(1)}f:\Delta^{(1)}_{X/S}\longrightarrow
 \Delta^{(1)}_{W/S}
\]
for which the following diagram is commutative.

\SGAThreeNativeDiagram[0.86\textwidth]
  {native_diagrams/exp3/diagram_006_exp3_p006_reminders017_delta_functoriality.tex}
  {Functoriality of the first infinitesimal neighborhood of the diagonal.}
  {fig:exp3-reminders-017-delta-functoriality}
  {Polo--Gille Exp3-14oct24.pdf, component-local p. 6; combined-reader p. 116}

On the other hand, through the projection \(\pi_X\), the scheme
\(\Delta^{(1)}_{X/S}\) is an affine \(X\)-scheme: it is the spectrum
of the augmented quasi-coherent \(\mathscr O_X\)-algebra
\[
 \mathscr P^1_{X/S}=\mathscr O_X\oplus\Omega^1_{X/S},
\]
where \(\Omega^1_{X/S}\) is a square-zero ideal. The augmentation is
the \(\mathscr O_X\)-algebra morphism
\[
 \varepsilon_X:\mathscr P^1_{X/S}\longrightarrow\mathscr O_X
\]
that vanishes on \(\Omega^1_{X/S}\) and corresponds to the closed
immersion \(\delta_X:X\hookrightarrow\Delta^{(1)}_{X/S}\). Hence every
morphism of \(S\)-schemes \(f:X\to W\) induces a morphism of augmented
\(\mathscr O_X\)-algebras
\[
 f^*(\mathscr P^1_{W/S})
 =\mathscr O_X\oplus f^*(\Omega^1_{W/S})
 \longrightarrow
 \mathscr P^1_{X/S}=\mathscr O_X\oplus\Omega^1_{X/S},
\]
or, equivalently, a morphism of \(\mathscr O_X\)-modules
\[
 f_{X/W/S}:f^*(\Omega^1_{W/S})\longrightarrow\Omega^1_{X/S};
\]
see EGA~IV\textsubscript{4}, (16.4.3.6), and (16.4.18.2) for the
notation \(f_{X/W/S}\).

% Source page: Exp3 component-local p. 7; combined-reader p. 117.
Since \(\pi_X:\Delta^{(1)}_{X/S}\to X\) is affine, part~(a) shows
that \(\Delta^{(1)}f\) is completely determined by \(f_{X/W/S}\).
For every \(X\)-scheme \(g:Y\to X\), the set
\[
 \operatorname{Hom}_X(Y,\Delta^{(1)}_{X/S})
 \simeq
 \operatorname{Hom}_{\mathscr O_Y\text{-}\mathrm{alg}}
 \bigl(\mathscr O_Y\oplus g^*(\Omega^1_{X/S}),\mathscr O_Y\bigr)
\]
identifies with a subset of
\(\operatorname{Hom}_{\mathscr O_Y}
 (g^*(\Omega^1_{X/S}),\mathscr O_Y)\), namely
\[
 \widetilde{\operatorname{Hom}}_{\mathscr O_Y}
 \bigl(g^*(\Omega^1_{X/S}),\mathscr O_Y\bigr),
\]
the subset formed by the \(\mathscr O_Y\)-morphisms
\(\psi:g^*(\Omega^1_{X/S})\to\mathscr O_Y\) for which
\(\operatorname{Im}(\psi)\) is a square-zero ideal of
\(\mathscr O_Y\).\footnote{Editor's note: In the proof of SGA~1,
III~5.1, the sentence ``But the algebra homomorphisms
\(\mathscr A\to\mathscr O_Y\) correspond \ldots'' must be corrected
accordingly; the remainder of the cited passage is unchanged.}

Consequently, apply part~(a) to the following diagram.

\SGAThreeNativeDiagram[0.48\textwidth]
  {native_diagrams/exp3/diagram_007_exp3_p007_reminders017_delta_application.tex}
  {The diagram to which the affine functoriality statement is applied.}
  {fig:exp3-reminders-017-delta-application}
  {Polo--Gille Exp3-14oct24.pdf, component-local p. 7; combined-reader p. 117}

Taking into account that \(\Delta^{(1)}f\) is the restriction of
\(f\times f\) to \(\Delta^{(1)}_{X/S}\), one obtains the following
commutative diagram, functorial in the \(X\)-scheme \(g:Y\to X\).

\SGAThreeNativeDiagram[0.97\textwidth]
  {native_diagrams/exp3/diagram_008_exp3_p007_reminders017_star_diagram.tex}
  {The functorial comparison of first-neighborhood morphisms and differential maps, equation \(0.1.7\,(*)\).}
  {fig:exp3-reminders-017-star}
  {Polo--Gille Exp3-14oct24.pdf, component-local p. 7; combined-reader p. 117}

\medskip
\noindent\SGAInlineRefTarget{sga3:iii:remark:0-1-7-1}\textbf{Remark 0.1.7.1.}
Let us conclude this paragraph with the following remark, which will be
useful later (cf. \SGARefLink{sga3:iii:corollary:0-1-10}{0.1.10}). Denote by \(\widetilde L_X\) the
\(X\)-functor that associates to every \(X\)-scheme \(g:Y\to X\) the
set
\[
 \operatorname{Hom}_{\mathscr O_Y}
 \bigl(g^*(\Omega^1_{X/S}),\mathscr O_Y\bigr),
\]
and denote by \(\widetilde L_f:\widetilde L_X\to\widetilde L_W\) the
morphism of functors defined above: it sends
\(\psi\in\widetilde L_X(Y)\) to
\(\psi\circ g^*(f_{X/W/S})\). The preceding discussion gives a
commutative diagram of functors.

\SGAThreeNativeDiagram[0.74\textwidth]
  {native_diagrams/exp3/diagram_009_exp3_p007_remark0171_functor_diagram.tex}
  {Functoriality of \(\widetilde L_X\) and the first infinitesimal neighborhood.}
  {fig:exp3-remark-0171-functor-diagram}
  {Polo--Gille Exp3-14oct24.pdf, component-local p. 7; combined-reader p. 117}

% Source page: Exp3 component-local p. 8; combined-reader p. 118.
\medskip
\noindent\SGAInlineRefTarget{sga3:iii:theorem:0-1-8}\textbf{Theorem 0.1.8.} (SGA~1, III~5.1)\footnote{Editor's
note: The statement and proof of SGA~1, III~5.1 are reproduced here,
since they are needed below to prove \SGARefLink{sga3:iii:proposition:0-2}{Proposition~0.2}; see already the
following \SGARefLink{sga3:iii:corollary:0-1-9}{Corollary~0.1.9}. Moreover, \(\mathscr J\) (respectively
\(\mathscr I\)) is another calligraphic form of the letter \(J\)
(respectively \(I\)). Returning to the earlier notation, with
\(\mathscr J\subset\mathscr I\) ideals of \(\mathscr O_S\) such that
\(\mathscr I\mathscr J=0\), this result will be applied with
\(\mathscr J=\mathscr J\mathscr O_Y\) (respectively
\(\mathscr I=\mathscr I\mathscr O_Y\)).}
\textit{Let \(Y\) and \(X\) be two \(S\)-schemes, let \(\mathscr J\)
be a quasi-coherent square-zero ideal of \(\mathscr O_Y\), let
\(Y_{\mathscr J}\) be the closed subscheme of \(Y\) defined by
\(\mathscr J\), and let \(g_{\mathscr J}:Y_{\mathscr J}\to X\) be an
\(S\)-morphism.}
\begin{enumerate}
\item[\SGAInlineRefTarget{sga3:iii:theorem:0-1-8:item:a}(a)] \textit{The set \(P(g_{\mathscr J})\) of \(S\)-morphisms
  \(g:Y\to X\) extending \(g_{\mathscr J}\) is either empty or
  principal homogeneous under the abelian group}
  \[
   \operatorname{Hom}_{\mathscr O_{Y_{\mathscr J}}}
   \bigl(g_{\mathscr J}^*(\Omega^1_{X/S}),\mathscr J\bigr).
  \]
\item[\SGAInlineRefTarget{sga3:iii:theorem:0-1-8:item:b}(b)] \textit{Suppose that \(i:Y_0\hookrightarrow Y_{\mathscr J}\)
  is the closed immersion defined by a quasi-coherent ideal
  \(\mathscr I\supset\mathscr J\) such that
  \(\mathscr I\mathscr J=0\), and set
  \(g_0=g_{\mathscr J}\circ i\). Then the preceding abelian group is
  isomorphic to}
  \[
   \operatorname{Hom}_{\mathscr O_{Y_0}}
   \bigl(g_0^*(\Omega^1_{X/S}),\mathscr J\bigr).
  \]
\end{enumerate}

\noindent\textit{Proof.}
Part~\SGARefLink{sga3:iii:theorem:0-1-8:item:b}{(b)} follows at once from part~\SGARefLink{sga3:iii:theorem:0-1-8:item:a}{(a)}. Indeed, since \(\mathscr J\) is
annihilated by \(\mathscr I\), it may be regarded as an
\(\mathscr O_{Y_0}\)-module, and adjunction gives
\[
 \operatorname{Hom}_{\mathscr O_{Y_{\mathscr J}}}
  \bigl(g_{\mathscr J}^*(\Omega^1_{X/S}),\mathscr J\bigr)
 =
 \operatorname{Hom}_{\mathscr O_{Y_0}}
  \bigl(i^*g_{\mathscr J}^*(\Omega^1_{X/S}),\mathscr J\bigr).
\]

To prove part~\SGARefLink{sga3:iii:theorem:0-1-8:item:a}{(a)}, we may suppose that
\(P(g_{\mathscr J})\ne\emptyset\), that is, that an \(S\)-morphism
\(g:Y\to X\) extending \(g_{\mathscr J}\) exists. Write
\(j:Y_{\mathscr J}\hookrightarrow Y\) for the immersion. Then
\(P(g_{\mathscr J})\) is the set of \(S\)-morphisms \(g':Y\to X\)
such that \(g'\circ j=g_{\mathscr J}\). Giving such a \(g'\) is
equivalent to giving an \(S\)-morphism
\[
 h:Y\longrightarrow X\times_S X
\]
such that \(\operatorname{pr}_1\circ h=g\) and
\(h_{\mathscr J}=\delta\circ g_{\mathscr J}\), where
\(h_{\mathscr J}=h\circ j\) and
\(\delta:X\hookrightarrow X\times_S X\) is the diagonal immersion.

\SGAThreeNativeDiagram[0.56\textwidth]
  {native_diagrams/exp3/diagram_010_exp3_p008_theorem018_factorization.tex}
  {The morphism \(h\) determined by two extensions of \(g_{\mathscr J}\).}
  {fig:exp3-theorem-018-factorization}
  {Polo--Gille Exp3-14oct24.pdf, component-local p. 8; combined-reader p. 118}

Since \(h_{\mathscr J}\) factors through \(\delta\), and since \(Y\)
lies in the first infinitesimal neighborhood of the immersion
\(j:Y_{\mathscr J}\to Y\) because \(\mathscr J^2=0\), functoriality
(EGA~IV\textsubscript{4}, 16.2.2(i)) shows that the desired morphisms
\(h\) factor uniquely through \(\Delta^{(1)}_{X/S}\); see 0.1.7. Put
\[
 Y'=\Delta^{(1)}_{X/S}\times_XY,
 \qquad
 Y'_{\mathscr J}=Y'\times_YY_{\mathscr J}
 =\Delta^{(1)}_{X/S}\times_XY_{\mathscr J}.
\]
The desired morphisms \(h\) are then in bijection with the sections
\(u\) of \(Y'\to Y\) extending the section
\(u_{\mathscr J}=(\delta\circ g_{\mathscr J},\operatorname{id})\) of
\(Y'_{\mathscr J}\to Y_{\mathscr J}\). On the other hand, \(Y'\)
(respectively \(Y'_{\mathscr J}\)) is affine over \(Y\)
(respectively \(Y_{\mathscr J}\)) and corresponds to the
quasi-coherent algebra
\[
 \mathscr A=\mathscr O_Y\oplus g^*(\Omega^1_{X/S}),
 \qquad\text{respectively}\qquad
 \mathscr A_{\mathscr J}
 =\mathscr A\otimes_{\mathscr O_Y}\mathscr O_{Y_{\mathscr J}}
 =\mathscr O_{Y_{\mathscr J}}
   \oplus g_{\mathscr J}^*(\Omega^1_{X/S}).
\]
Let \(\varepsilon:\mathscr A\to\mathscr O_Y\) be the canonical
augmentation, that is, the \(\mathscr O_Y\)-algebra morphism that
vanishes on \(g^*(\Omega^1_{X/S})\), and define
\(\varepsilon_{\mathscr J}:\mathscr A_{\mathscr J}
 \to\mathscr O_{Y_{\mathscr J}}\) similarly. Then
\[
 \Gamma(Y'/Y)
 \simeq\operatorname{Hom}_{\mathscr O_Y\text{-}\mathrm{alg}}
       (\mathscr A,\mathscr O_Y),
 \qquad
 \Gamma(Y'_{\mathscr J}/Y_{\mathscr J})
 \simeq
 \operatorname{Hom}_{\mathscr O_{Y_{\mathscr J}}\text{-}\mathrm{alg}}
       (\mathscr A_{\mathscr J},\mathscr O_{Y_{\mathscr J}}).
\]

% Source page: Exp3 component-local p. 9; combined-reader p. 119.
Under these isomorphisms, the section
\(u=(\delta\circ g,\operatorname{id})\) (respectively
\(u_{\mathscr J}\)) corresponds to \(\varepsilon\) (respectively
\(\varepsilon_{\mathscr J}\)). Consequently,
\(P(g_{\mathscr J})\) is in bijection with the algebra morphisms
\(\mathscr A\to\mathscr O_Y\) that reduce to
\(\varepsilon_{\mathscr J}\); under this bijection, \(g\) corresponds
to \(\varepsilon\).

Put \(M=g^*(\Omega^1_{X/S})\). Then
\(\operatorname{Hom}_{\mathscr O_Y\text{-}\mathrm{alg}}
 (\mathscr A,\mathscr O_Y)\) identifies with the set of
\(\mathscr O_Y\)-morphisms \(\psi:M\to\mathscr O_Y\) whose image is a
square-zero ideal. We are interested in those morphisms that induce the
zero morphism
\(M\to\mathscr O_{Y_{\mathscr J}}=\mathscr O_Y/\mathscr J\), that is,
those that carry \(M\) into \(\mathscr J\). Conversely, since
\(\mathscr J^2=0\), every \(\mathscr O_Y\)-morphism
\(\phi:M\to\mathscr J\) comes from a unique algebra morphism
\(\mathscr A\to\mathscr O_Y\) reducing to
\(\varepsilon_{\mathscr J}\). Finally,
\[
 \operatorname{Hom}_{\mathscr O_Y}
   \bigl(g^*(\Omega^1_{X/S}),\mathscr J\bigr)
 =
 \operatorname{Hom}_{\mathscr O_{Y_{\mathscr J}}}
   \bigl(g_{\mathscr J}^*(\Omega^1_{X/S}),\mathscr J\bigr)
\]
because \(\mathscr J^2=0\), as in the proof of part~\SGARefLink{sga3:iii:theorem:0-1-8:item:b}{(b)}. Thus one
obtains a bijection
\[
 P(g_{\mathscr J})\simeq
 \operatorname{Hom}_{\mathscr O_{Y_{\mathscr J}}}
 \bigl(g_{\mathscr J}^*(\Omega^1_{X/S}),\mathscr J\bigr)
\]
under which \(g\) corresponds to the zero morphism.

For every
\[
 m\in
 \operatorname{Hom}_{\mathscr O_{Y_{\mathscr J}}}
 \bigl(g_{\mathscr J}^*(\Omega^1_{X/S}),\mathscr J\bigr),
\]
write \(m\cdot g\) for the element of \(P(g_{\mathscr J})\) associated
with \(g\) and \(m\) by the preceding bijection. We have already seen
that \(0\cdot g=g\); it remains to prove
\begin{equation}
 m'\cdot(m\cdot g)=(m+m')\cdot g.
 \tag{0.1.8\,(*)}\label{sga3:iii:theorem:0-1-8:equation:0-1-8-star}
\end{equation}
This may be checked locally.\footnote{Editor's note: The preceding
argument is reproduced almost word for word from SGA~1, III~5.1; what
follows has been expanded.} Indeed, the two preceding morphisms
\(Y\to X\) induce the same continuous map as \(g\) on the underlying
topological spaces. It is therefore enough to check that, for every
affine open \(U=\operatorname{Spec}(A)\) of \(X\) lying over an affine
open \(\operatorname{Spec}(\Lambda)\) of \(S\), and every affine open
\(V=\operatorname{Spec}(B)\) of \(g^{-1}(U)\), they induce the same
morphism of \(\Lambda\)-algebras \(A\to B\).

Set \(J=\Gamma(V,\mathscr J)\), and let
\(\varphi,\psi,\eta:A\to B\) be the morphisms induced respectively by
\(g\), \(m\cdot g\), and \(m'\cdot(m\cdot g)\). They agree modulo
\(J\). There are unique expressions
\[
 \psi=\varphi+D,
 \qquad
 \eta=\psi+D',
\]
where \(D\) belongs to
\[
 \operatorname{Der}_{\varphi}(A,J)
 =\{\delta\in\operatorname{Hom}_{\Lambda}(A,J)
   \mid \delta(ab)=\varphi(a)\delta(b)+\varphi(b)\delta(a)\},
\]
and \(D'\) belongs to \(\operatorname{Der}_{\psi}(A,J)\). But
\(\operatorname{Der}_{\varphi}(A,J)
 =\operatorname{Der}_{\psi}(A,J)\), since \(J^2=0\), and both identify
with
\[
 \operatorname{Hom}_{B/J}
 \bigl(\Omega^1_{A/\Lambda}\otimes_AB/J,J\bigr).
\]
Under this identification, \(D\) corresponds to \(m\) and \(D'\) to
\(m'\). Hence \(\eta=\varphi+D+D'\), while \(D+D'\) corresponds to
\(m+m'\), proving (\SGARefLink{sga3:iii:theorem:0-1-8}{0.1.8}\,(*)).

\medskip
\noindent\SGAInlineRefTarget{sga3:iii:corollary:0-1-9}\textbf{Corollary 0.1.9.}\footnote{Editor's note: This
corollary has been added to SGA~1, III~5.1; it proves part~(i) of
\SGARefLink{sga3:iii:proposition:0-2}{Proposition~0.2} below.}
\textit{Let \(X\) be an \(S\)-scheme, and retain the notation of
\SGARefLink{sga3:iii:definition:0-1-5}{0.1.5}. Then \(X\) carries a left action of the abelian
\(X^+\)-group \(L_X\), making \(X\) formally principal homogeneous
under \(L_X\) over \(X^+\). In other words, there is an isomorphism of
\(X^+\)-functors}
\[
 L_X\times_{X^+}X\xrightarrow{\sim}X\times_{X^+}X,
 \qquad
 (m,x)\longmapsto(x,m\cdot x).
\]

% Source page: Exp3 component-local p. 10; combined-reader p. 120.
\noindent\textit{Proof.}
Let \(i_0:X_0\hookrightarrow X\) be the immersion. Since
\(X_0=X\times_SS_0\), one first has
\[
 i_0^*(\Omega^1_{X/S})\simeq\Omega^1_{X_0/S_0}
\]
(EGA~IV, 16.4.5).

Let \(Y\) be an \(X^+\)-scheme, given by an \(S\)-morphism
\(g_{\mathscr J}:Y_{\mathscr J}\to X\), and let
\(g_0:Y_0\to X_0\) be obtained by base change. By \SGARefLink{sga3:iii:theorem:0-1-8}{0.1.8}, if
\(\operatorname{Hom}_{X^+}(Y,X)\) is nonempty, it is principal
homogeneous under
\[
 \operatorname{Hom}_{\mathscr O_{Y_0}}
 \bigl(g_0^*i_0^*(\Omega^1_{X/S}),\mathscr J\mathscr O_Y\bigr),
\]
which identifies with
\[
 \operatorname{Hom}_{\mathscr O_{Y_0}}
 \bigl(g_0^*(\Omega^1_{X_0/S_0}),\mathscr J\mathscr O_Y\bigr)
 =L_X(Y).
\]
Thus there is a bijection
\[
 L_X(Y)\times\operatorname{Hom}_{X^+}(Y,X)
 \xrightarrow{\sim}
 \operatorname{Hom}_{X^+}(Y,X\times_{X^+}X),
 \qquad
 (m,g)\longmapsto(g,m\cdot g).
\]
We show that this is functorial in \(Y\). Let \(f:Z\to Y\) be a
morphism of \(S\)-schemes. It is enough to prove that the following
diagram is commutative.

\SGAThreeNativeDiagram[0.88\textwidth]
  {native_diagrams/exp3/diagram_011_exp3_p010_corollary019_functoriality.tex}
  {Functoriality in \(Y\) of the principal-homogeneous bijection.}
  {fig:exp3-corollary-019-functoriality}
  {Polo--Gille Exp3-14oct24.pdf, component-local p. 10; combined-reader p. 120}

If \(\operatorname{Hom}_{X^+}(Y,X)=\emptyset\), there is nothing to
prove. It is therefore enough to show that, for every \(S\)-morphism
\(g:Y\to X\) extending \(g_{\mathscr J}\) and every \(m\in L_X(Y)\),
one has
\begin{equation}
 (m\cdot g)\circ f=L_X(f)(m)\cdot(g\circ f).
 \tag{0.1.9\,(*)}\label{sga3:iii:corollary:0-1-9:equation:0-1-9-star}
\end{equation}
These two \(S\)-morphisms \(Z\to X\) agree on \(Z_{\mathscr J}\) with
\(g_{\mathscr J}\circ f_{\mathscr J}\); in particular, they induce
the same continuous map as \(g\circ f\) on the underlying topological
spaces. It therefore suffices to check the equality on local rings.
Let \(z\in Z\), put \(y=f(z)\), \(x=g(y)\), and let \(A,B,C\) denote
respectively the local rings
\(\mathscr O_{X,x},\mathscr O_{Y,y},\mathscr O_{Z,z}\). Let \(s\) be
the image of \(x\) in \(S\), put \(\Lambda=\mathscr O_{S,s}\), let
\(J\) and \(I\) be the ideals of \(\Lambda\) corresponding to
\(\mathscr J\) and \(\mathscr I\), and let
\(\varphi,\psi:A\to B\) and \(\theta:B\to C\) be the
\(\Lambda\)-algebra morphisms induced by \(g\), \(m\cdot g\), and
\(f\), respectively. Then \(m\) induces an element \(D\) of
\[
 \operatorname{Hom}_{B/IB}
 \bigl(\Omega^1_{A/\Lambda}\otimes_AB/IB,JB\bigr)
 =\operatorname{Der}_{\Lambda}(A,JB),
\]
and \(\psi=\varphi+D\). Thus \((m\cdot g)\circ f\) corresponds to
\[
 \theta\circ\psi=\theta\circ\varphi+\theta\circ D.
\]
By \SGARefLink{sga3:iii:definition:0-1-5}{0.1.5}, \(\theta\circ D\) is the image of \(L_X(f)(m)\) in
\[
 \operatorname{Hom}_{C/IC}
 \bigl(\Omega^1_{A/\Lambda}\otimes_AC/IC,JC\bigr)
 =\operatorname{Der}_{\Lambda}(A,JC).
\]
Consequently, \(\theta\circ\varphi+\theta\circ D\) is the image of
\(L_X(f)(m)\cdot(g\circ f)\). This proves (\SGARefLink{sga3:iii:corollary:0-1-9}{0.1.9}\,(*)).

\medskip
\noindent\SGAInlineRefTarget{sga3:iii:corollary:0-1-10}\textbf{Corollary 0.1.10.}\footnote{Editor's note: The
functoriality of \(L_X\) in \(X\) has been made explicit, in particular
in part~\SGARefLink{sga3:iii:corollary:0-1-10:item:b}{(b)} below.}
\begin{enumerate}
\item[\SGAInlineRefTarget{sga3:iii:corollary:0-1-10:item:a}(a)] \textit{The functor \(L_X\) depends functorially on
  \(X\). For every \(S\)-morphism \(f:X\to W\), there is an
  \(S\)-morphism \(L_f:L_X\to L_W\) that is a morphism of abelian
  groups over \(f^+\); that is, the following diagram is commutative,}

  % Source page: Exp3 component-local p. 11; combined-reader p. 121.
  \SGAThreeNativeDiagram[0.46\textwidth]
    {native_diagrams/exp3/diagram_012_exp3_p011_corollary0110_lf_square.tex}
    {The morphism \(L_f\) over \(f^+\).}
    {fig:exp3-corollary-0110-lf-square}
    {Polo--Gille Exp3-14oct24.pdf, component-local p. 11; combined-reader p. 121}

  \textit{and, for every \(Y\to X^+\), the map}
  \[
   L_f(Y):\operatorname{Hom}_{X^+}(Y,L_X)
   \longrightarrow\operatorname{Hom}_{W^+}(Y,L_W),
  \]
  \textit{where \(Y\) lies over \(W^+\) through \(f^+\), is a
  homomorphism of abelian groups.}

\item[\SGAInlineRefTarget{sga3:iii:corollary:0-1-10:item:b}(b)] \textit{Moreover, the following diagram is commutative.}

  \SGAThreeNativeDiagram[0.66\textwidth]
    {native_diagrams/exp3/diagram_013_exp3_p011_corollary0110_action_square.tex}
    {Compatibility of \(L_f\) with the formally principal homogeneous actions.}
    {fig:exp3-corollary-0110-action-square}
    {Polo--Gille Exp3-14oct24.pdf, component-local p. 11; combined-reader p. 121}
\end{enumerate}

\noindent\textit{Proof.}
For part~\SGARefLink{sga3:iii:corollary:0-1-10:item:a}{(a)}, \(L_f\) is induced by the morphism of
\(\mathscr O_{X_0}\)-modules
\[
 f_{X_0/W_0/S_0}:f_0^*(\Omega^1_{W_0/S_0})
 \longrightarrow\Omega^1_{X_0/S_0}
\]
(cf. \SGARefLink{sga3:iii:paragraph:0-1-7:item:b}{0.1.7(b)}). For every \(X^+\)-scheme \(Y\), given by an
\(S\)-morphism \(g_{\mathscr J}:Y_{\mathscr J}\to X\), one has the
following commutative diagram, functorial in \(Y\).

\SGAThreeNativeDiagram[0.96\textwidth]
  {native_diagrams/exp3/diagram_014_exp3_p011_corollary0110_hom_square.tex}
  {The homomorphism \(L_f(Y)\) on the differential-module descriptions.}
  {fig:exp3-corollary-0110-hom-square}
  {Polo--Gille Exp3-14oct24.pdf, component-local p. 11; combined-reader p. 121}

Here \(L_f(Y)\) is the map
\[
 \psi\longmapsto
 \psi\circ g_0^*(f_{X_0/W_0/S_0}),
\]
which is indeed a homomorphism of abelian groups.\footnote{Editor's
note: Observe that \(L_f\) depends only on
\(f_0:X_0\to W_0\).}

We now prove part~\SGARefLink{sga3:iii:corollary:0-1-10:item:b}{(b)}. Let \(Y\) be an \(X^+\)-scheme. If
\(\operatorname{Hom}_{X^+}(Y,X)=\emptyset\), there is nothing to
prove. Thus let \(g\in\operatorname{Hom}_{X^+}(Y,X)\). We must show
that, for every \(m\in L_X(Y)\),
\begin{equation}
 f\circ(m\cdot g)=L_f(Y)(m)\cdot(f\circ g).
 \tag{0.1.10\,(*)}\label{sga3:iii:corollary:0-1-10:equation:0-1-10-star}
\end{equation}
With \(g\) fixed,
\(\operatorname{Hom}_X(Y,X\times_{X^+}X)\) is a subset of
\(\operatorname{Hom}_X(Y,\Delta^{(1)}_{X/S})\), and
\[
\begin{aligned}
 L_X(Y)
 &=\operatorname{Hom}_{\mathscr O_{Y_0}}
   \bigl(g_0^*(\Omega^1_{X_0/S_0}),\mathscr J\mathscr O_Y\bigr)\\
 &=\operatorname{Hom}_{\mathscr O_Y}
   \bigl(g^*(\Omega^1_{X/S}),\mathscr J\mathscr O_Y\bigr),
\end{aligned}
\]
which is a subset of \(\widetilde L_X(Y)\); see \SGARefLink{sga3:iii:paragraph:0-1-7}{0.1.7}. Moreover,
\(L_f(Y)\) is the restriction of \(\widetilde L_f(Y)\) to
\(L_X(Y)\). Finally, the bijection
\[
 L_X(Y)\xrightarrow{\sim}
 \operatorname{Hom}_X(Y,X\times_{X^+}X),
 \qquad
 m\longmapsto(g,m\cdot g),
\]
is the inverse of the restriction to
\(L_X(Y)\subset\widetilde L_X(Y)\) of the bijection
\[
 \operatorname{Hom}_X(Y,\Delta^{(1)}_{X/S})
 \xrightarrow{\sim}\{g\}\times\widetilde L_X(Y)
\]
considered in \SGARefLink{sga3:iii:remark:0-1-7-1}{0.1.7.1}. Consequently, equality (\SGARefLink{sga3:iii:corollary:0-1-10}{0.1.10}\,(*)) follows
% Source page: Exp3 component-local p. 12; combined-reader p. 122.
from (\SGARefLink{sga3:iii:paragraph:0-1-7}{0.1.7}\,(*)). Indeed, let \(g'\) be the \(X\)-morphism
\(Y\to\Delta^{(1)}_{X/S}\) defined by \((g,m\cdot g)\). The element
of \(\widetilde L_W(Y)\) corresponding to
\((f\circ g,f\circ g')\) is
\(\widetilde L_f(Y)(m)=L_f(Y)(m)\). Thus
\[
 L_f(m)\cdot(f\circ g)=f\circ(m\cdot g).
\]

% Continue with Lemma 0.1.11.

% END INCLUDED SGA3_English_current_buildable_source_20260823/source/inputs/III/tex/components/16_expose_III_s017_s0110_functorial_liftings.tex

% BEGIN INCLUDED SGA3_English_current_buildable_source_20260823/source/inputs/III/tex/components/17_expose_III_s0111_s02_products_action.tex
% SGA 3, Expose III, Loop 1 English reconstruction.
% Scope: Lemma 0.1.11 through Proposition 0.2, stopping before Remark 0.3.
% Authority: Polo--Gille Exp3-14oct24.pdf, component-local pages 12--14
% (combined-reader pages 122--124; printed pages 112--114, with original
% pagination 88--90 also shown by the re-edition). Force-OCR is a drafting
% locator; English wording was compared with Jacob Reinhold's CC BY 4.0
% Markdown lineage. The PDF controls formulas, notes, and diagram images.

\setcounter{footnote}{17}

% Source page: Exp3 component-local p. 12; combined-reader p. 122.
\medskip
\noindent\SGAInlineRefTarget{sga3:iii:lemma:0-1-11}\textbf{Lemma 0.1.11.}
\textit{Let \(X,X'\) be two \(S\)-schemes. There is a commutative
diagram}

\SGAThreeNativeDiagram[0.58\textwidth]
  {native_diagrams/exp3/diagram_015_exp3_p012_lemma0111_product_square.tex}
  {Compatibility of \(L_X\) and \(X^+\) with products.}
  {fig:exp3-lemma-0111-product}
  {Polo--Gille Exp3-14oct24.pdf, component-local p. 12; combined-reader p. 122}

\textit{Proof.}\footnote{Editor's note: the proof has been expanded.}
For every \(S\)-scheme \(Y\),
\[
\begin{aligned}
 \operatorname{Hom}_S(Y,X^+\times_S X'^+)
 &=
 \operatorname{Hom}_S(Y,X^+)
 \times\operatorname{Hom}_S(Y,X'^+)\\
 &\simeq
 \operatorname{Hom}_S(Y_{\mathscr J},X)
 \times\operatorname{Hom}_S(Y_{\mathscr J},X')\\
 &=
 \operatorname{Hom}_S\bigl(Y,(X\times_S X')^+\bigr).
\end{aligned}
\]
Thus
\[
 X^+\times_S X'^+\simeq(X\times_S X')^+.
\]

Now let \(Y\) be a scheme over \(X^+\times_S X'^+\), represented by a
morphism
\[
 h:Y_{\mathscr J}\longrightarrow X\times_S X'.
\]
Let \(p,q\) be the projections from \(X\times_S X'\) to \(X,X'\), and
put \(f=p\circ h\), \(g=q\circ h\). Since
\[
 \Omega^1_{(X_0\times_{S_0}X'_0)/S_0}
 \simeq
 p_0^*(\Omega^1_{X_0/S_0})
 \oplus q_0^*(\Omega^1_{X'_0/S_0})
\]
(EGA~IV\textsubscript{4}, 16.4.23), there is a natural isomorphism
\[
\begin{aligned}
 &\operatorname{Hom}_{\mathscr O_{Y_0}}
   \bigl(f_0^*(\Omega^1_{X_0/S_0}),\mathscr J\mathscr O_Y\bigr)
 \times
 \operatorname{Hom}_{\mathscr O_{Y_0}}
   \bigl(g_0^*(\Omega^1_{X'_0/S_0}),\mathscr J\mathscr O_Y\bigr)\\
 &\qquad\simeq
 \operatorname{Hom}_{\mathscr O_{Y_0}}
 \bigl(h_0^*(\Omega^1_{(X_0\times_{S_0}X'_0)/S_0}),
       \mathscr J\mathscr O_Y\bigr).
\end{aligned}
\]
In other words,
\[
 L_X(Y)\times L_{X'}(Y)\simeq L_{X\times_S X'}(Y).
\]

\medskip
\noindent\SGAInlineRefTarget{sga3:iii:remark:0-1-12}\textbf{Remark 0.1.12.}\footnote{Editor's note: this remark
and the corollary that follows have been added.}
Let \(\mathscr C\) be a category stable under fiber products, let \(S\)
be an object of \(\mathscr C\), and let \(T_1,T_2\) be two objects over
\(S\). For \(i=1,2\), let \(L_i\) and \(X_i\) be two objects over
\(T_i\), as in the diagram below.

\SGAThreeNativeDiagram[0.52\textwidth]
  {native_diagrams/exp3/diagram_016_exp3_p012_remark0112_fiber_objects.tex}
  {The objects \(L_i,X_i\) over \(T_i\), and \(T_i\) over \(S\).}
  {fig:exp3-remark-0112-objects}
  {Polo--Gille Exp3-14oct24.pdf, component-local p. 12; combined-reader p. 122}

There is then a natural isomorphism
\[
 (L_1\times_{T_1}X_1)\times_S(L_2\times_{T_2}X_2)
 \simeq
 (L_1\times_S L_2)
 \times_{T_1\times_S T_2}(X_1\times_S X_2).
\]
Consequently, the preceding lemma gives:

% Source page: Exp3 component-local p. 13; combined-reader p. 123.
\medskip
\noindent\SGAInlineRefTarget{sga3:iii:corollary:0-1-13}\textbf{Corollary 0.1.13.}
\textit{Let \(X_1,X_2\) be two \(S\)-schemes. There is a commutative
diagram of isomorphisms:}

\SGAThreeNativeDiagram[0.96\textwidth]
  {native_diagrams/exp3/diagram_017_exp3_p013_corollary0113_isomorphisms.tex}
  {The product isomorphisms of \SGARefLink{sga3:iii:corollary:0-1-13}{Corollary 0.1.13}.}
  {fig:exp3-corollary-0113}
  {Polo--Gille Exp3-14oct24.pdf, component-local p. 13; combined-reader p. 123}

We can now state:

\medskip
\noindent\SGAInlineRefTarget{sga3:iii:proposition:0-2}\textbf{Proposition 0.2.}
\textit{For every \(S\)-scheme \(X\), one can define a left action of
the \(X^+\)-abelian group \(L_X\) on the \(X^+\)-object \(X\), with
the following properties:}
\begin{enumerate}
\item[\SGAInlineRefTarget{sga3:iii:proposition:0-2:item:i}(i)] \textit{this action makes \(X\) formally principal
homogeneous under \(L_X\) over \(X^+\); that is, the morphism}
\[
 L_X\times_{X^+}X\longrightarrow X\times_{X^+}X
\]
\textit{is an isomorphism of \(X^+\)-functors;}

\item[\SGAInlineRefTarget{sga3:iii:proposition:0-2:item:ii}(ii)] \textit{the action is functorial in the \(S\)-scheme
\(X\): for every \(S\)-morphism \(f:X\to W\), the following diagram
is commutative;}

\SGAThreeNativeDiagram[0.55\textwidth]
  {native_diagrams/exp3/diagram_018_exp3_p013_proposition02_functoriality_square.tex}
  {Functoriality of the action in \SGARefLink{sga3:iii:proposition:0-2:item:ii}{Proposition 0.2(ii)}.}
  {fig:exp3-proposition-02-functoriality}
  {Polo--Gille Exp3-14oct24.pdf, component-local p. 13; combined-reader p. 123}

\item[\SGAInlineRefTarget{sga3:iii:proposition:0-2:item:iii}(iii)] \textit{the action commutes with fiber products: for all
\(S\)-schemes \(X_1,X_2\), the following diagram is commutative.}
\end{enumerate}

\SGAThreeNativeDiagram[0.99\textwidth]
  {native_diagrams/exp3/diagram_019_exp3_p013_proposition02_product_diagram.tex}
  {Compatibility of the action with fiber products.}
  {fig:exp3-proposition-02-products}
  {Polo--Gille Exp3-14oct24.pdf, component-local p. 13; combined-reader p. 123}

\textit{Proof.}\footnote{Editor's note: the original has been modified
and expanded to account for the preceding additions; in particular,
those additions incorporate, with details, the content of original
page 89.}
Parts \SGARefLink{sga3:iii:proposition:0-2:item:i}{(i)} and \SGARefLink{sga3:iii:proposition:0-2:item:ii}{(ii)} follow respectively from \SGARefLink{sga3:iii:corollary:0-1-9}{Corollaries 0.1.9} and
\SGARefLink{sga3:iii:corollary:0-1-10}{0.1.10}. To prove \SGARefLink{sga3:iii:proposition:0-2:item:iii}{(iii)}, write
\[
 P(X)=L_X\times_{X^+}X
\]
for every \(S\)-scheme \(X\). Applying \SGARefLink{sga3:iii:proposition:0-2:item:ii}{(ii)} to the projections
\[
 p_i:X_1\times_S X_2\longrightarrow X_i
\]
gives the following commutative squares.

\SGAThreeNativeDiagram[0.55\textwidth]
  {native_diagrams/exp3/diagram_020_exp3_p013_proposition02_projection_square.tex}
  {The square induced by each projection \(p_i\).}
  {fig:exp3-proposition-02-projection-square}
  {Polo--Gille Exp3-14oct24.pdf, component-local p. 13; combined-reader p. 123}

% Source page: Exp3 component-local p. 14; combined-reader p. 124.
For \(i=1,2\), these yield the commutative square

\SGAThreeNativeDiagram[0.55\textwidth]
  {native_diagrams/exp3/diagram_021_exp3_p014_proposition02_final_square.tex}
  {The product square in the proof of \SGARefLink{sga3:iii:proposition:0-2:item:iii}{Proposition 0.2(iii)}.}
  {fig:exp3-proposition-02-final-square}
  {Polo--Gille Exp3-14oct24.pdf, component-local p. 14; combined-reader p. 124}

Combining this square with \SGARefLink{sga3:iii:corollary:0-1-13}{Corollary 0.1.13} shows that its vertical
arrow is an isomorphism and gives precisely the commutative diagram
asserted in \SGARefLink{sga3:iii:proposition:0-2:item:iii}{(iii)}.

% Continue with Remark 0.3.

% END INCLUDED SGA3_English_current_buildable_source_20260823/source/inputs/III/tex/components/17_expose_III_s0111_s02_products_action.tex

% BEGIN INCLUDED SGA3_English_current_buildable_source_20260823/source/inputs/III/tex/components/18_expose_III_s03_s07_flat_liftings.tex
% SGA 3, Expose III, Loop 1 English reconstruction.
% Scope: Remark 0.3 through Corollary 0.7, stopping immediately before
% Corollary 0.7 bis.
% Authority: Polo--Gille Exp3-14oct24.pdf, component-local pages 14--17
% (combined-reader pages 124--127; printed pages 114--117). Force-OCR is a
% drafting locator; English wording was compared with Jacob Reinhold's
% CC BY 4.0 Markdown lineage. The PDF controls text, formulas, numbering,
% notes, and diagram images.

\setcounter{footnote}{20}

% Source page: Exp3 component-local p. 14; combined-reader p. 124; printed p. 114.
\medskip
\noindent\SGAInlineRefTarget{sga3:iii:remark:0-3}\textbf{Remark 0.3.}
Suppose that the \(X^+\)-scheme \(Y\) is flat over \(S\) (cf. SGA~1,
IV). We may then write
\[
 \operatorname{Hom}_{X^+}(Y,L_X)
 =
 \operatorname{Hom}_{\mathscr O_{Y_0}}
 \bigl(g_0^*(\Omega^1_{X_0/S_0}),
       \mathscr J\otimes_{\mathscr O_{S_0}}\mathscr O_{Y_0}\bigr).
\]

\medskip
\noindent\SGAInlineRefTarget{sga3:iii:remark:0-4}\textbf{Remark 0.4.}
Let \(\pi_0:X_0\to S_0\) be the structure morphism, and suppose that
there is an \(\mathscr O_{S_0}\)-module \(\omega^1_{X_0/S_0}\) such
that
\[
 \Omega^1_{X_0/S_0}\simeq
 \pi_0^*(\omega^1_{X_0/S_0}).
\]
(This occurs, in particular, when \(X_0\) is an \(S_0\)-group; cf.
\SGARefLink{sga3:ii:subsection:4-11}{II~4.11}.) Define a functor \(L'_X\) over \(S\) by
\begin{equation}
 \operatorname{Hom}_S(Y,L'_X)
 =
 \operatorname{Hom}_{\mathscr O_{Y_0}}
 \bigl(\omega^1_{X_0/S_0}\otimes_{\mathscr O_{S_0}}\mathscr O_{Y_0},
       \mathscr J\mathscr O_Y\bigr).
 \tag{0.4.1}\label{sga3:iii:remark:0-4:equation:0-4-1}
\end{equation}
Then
\[
 \operatorname{Hom}_{X^+}(Y,L_X)=\operatorname{Hom}_S(Y,L'_X)
\]
for every \(X^+\)-scheme \(Y\); in other words,
\[
 L_X=L'_X\mathbin{\times_S}X^+.
\]

Since
\(L_X\mathbin{\times_{X^+}}X=L'_X\mathbin{\times_S}X\), the action
of \(L_X\) on \(X\) induces an action of \(L'_X\) on \(X\), and
this action respects \(p_X:X\to X^+\).\footnote{Editor's note: what
follows has been expanded.} Indeed, if \(Y\) is an \(S\)-scheme,
\(h:Y\to X\) is an \(S\)-morphism, and \(m\in L'_X(Y)\), then
\(h\) and \(m\cdot h\) have the same restriction to \(Y_{\mathscr J}\);
that is,
\[
 p_X(m\cdot h)=p_X(h).
\]

\medskip
\noindent\SGAInlineRefTarget{sga3:iii:remark:0-5}\textbf{Remark 0.5.}
Keep the hypotheses and notation of \SGARefLink{sga3:iii:remark:0-4}{0.4}, and suppose in addition that
\(Y\) is an \(S\)-scheme flat over \(S\). Then
\[
 \operatorname{Hom}_{X^+}(Y,L_X)
 =\operatorname{Hom}_S(Y,L'_X)
 =\operatorname{Hom}_{S_0}(Y_0,L_{0X}),
\]
where the \(S_0\)-functor in abelian groups \(L_{0X}\) is defined by
the following identity (with \(T\) the variable \(S_0\)-scheme):
\begin{equation}
 \operatorname{Hom}_{S_0}(T,L_{0X})
 =
 \operatorname{Hom}_{\mathscr O_T}
 \bigl(\omega^1_{X_0/S_0}\otimes_{\mathscr O_{S_0}}\mathscr O_T,
       \mathscr J\otimes_{\mathscr O_{S_0}}\mathscr O_T\bigr).
 \tag{0.5.1}\label{sga3:iii:remark:0-5:equation:0-5-1}
\end{equation}
Thus, in the notation of II~1, we have shown that \(L'_X\) and
\(\prod_{S_0/S}L_{0X}\) have the same restriction to the full
subcategory of \((\mathbf{Sch})/S\) whose objects are the \(S\)-schemes
flat over \(S\).

\medskip
\noindent\SGAInlineRefTarget{sga3:iii:remark:0-6}\textbf{Remark 0.6.}
Keep the hypotheses and notation of \SGARefLink{sga3:iii:remark:0-5}{0.5},\footnote{Editor's note: the
following hypothesis, which was implicit in the original, has been
added.} and suppose in addition that \(\pi_0:X_0\to S_0\) has a
section \(\varepsilon_0\). Then
\[
 \omega^1_{X_0/S_0}\simeq
 \varepsilon_0^*(\Omega^1_{X_0/S_0}).
\]
First, independently of this last hypothesis, one has
\[
 \operatorname{Hom}_{S_0}(T,L_{0X})
 =
 \Gamma\!\left(
   T,
   \mathscr{H}om_{\mathscr O_T}
   \bigl(\omega^1_{X_0/S_0}\otimes_{\mathscr O_{S_0}}\mathscr O_T,
         \mathscr J\otimes_{\mathscr O_{S_0}}\mathscr O_T\bigr)
 \right).
\]

% Source page: Exp3 component-local p. 15; combined-reader p. 125; printed p. 115.
Suppose now that \(\omega^1_{X_0/S_0}\) admits a finite presentation
(cf. EGA~0\textsubscript{I}, 5.2.5), as is in particular the case when
\(X_0\) is locally of finite presentation over \(S_0\) (cf.
EGA~IV\textsubscript{4}, 16.4.22). If \(T\) is flat over \(S_0\), it
then follows from EGA~0\textsubscript{I}, 6.7.6 that
\[
 \mathscr{H}om_{\mathscr O_T}
 \bigl(\omega^1_{X_0/S_0}\otimes_{\mathscr O_{S_0}}\mathscr O_T,
       \mathscr J\otimes_{\mathscr O_{S_0}}\mathscr O_T\bigr)
 \simeq
 \mathscr{H}om_{\mathscr O_{S_0}}
 \bigl(\omega^1_{X_0/S_0},\mathscr J\bigr)
 \otimes_{\mathscr O_{S_0}}\mathscr O_T,
\]
whence
\[
 \operatorname{Hom}_{S_0}(T,L_{0X})
 =
 \Gamma\!\left(
   T,
   \mathscr{H}om_{\mathscr O_{S_0}}
   \bigl(\omega^1_{X_0/S_0},\mathscr J\bigr)
   \otimes_{\mathscr O_{S_0}}\mathscr O_T
 \right).
\]
Using the notation \(\mathbf W(\,\cdot\,)\) of \SGARefLink{sga3:i:definition:4-6-1}{I~4.6.1}, we have
therefore proved that, for every \(S_0\)-scheme \(T\) flat over
\(S_0\),
\[
 \operatorname{Hom}_{S_0}(T,L_{0X})
 =
 \operatorname{Hom}_{S_0}
 \left(T,
 \mathbf W\!\left(
   \mathscr{H}om_{\mathscr O_{S_0}}
   (\omega^1_{X_0/S_0},\mathscr J)
 \right)\right).
\]

In summary, if \(\omega^1_{X_0/S_0}\) admits a finite presentation
and one restricts to the category of \(S\)-schemes flat over \(S\),
then
\begin{equation}
 L'_X
 =
 \prod_{S_0/S}
 \mathbf W\!\left(
   \mathscr{H}om_{\mathscr O_{S_0}}
   (\omega^1_{X_0/S_0},\mathscr J)
 \right),
 \tag{0.6.1}\label{sga3:iii:remark:0-6:equation:0-6-1}
\end{equation}
and
\(\mathscr{H}om_{\mathscr O_{S_0}}
(\omega^1_{X_0/S_0},\mathscr J)\) is a quasi-coherent
\(\mathscr O_{S_0}\)-module, by EGA~I, 9.1.1.

Finally, suppose moreover that \(\omega^1_{X_0/S_0}\) is locally free
of finite rank, for example that \(X_0\) is smooth over \(S_0\) (in
which case it is automatically locally of finite presentation over
\(S_0\)). Then
\begin{equation}
 \mathscr{H}om_{\mathscr O_{S_0}}
   (\omega^1_{X_0/S_0},\mathscr J)
 \simeq
 \underline{\operatorname{Lie}}(X_0/S_0)
 \otimes_{\mathscr O_{S_0}}\mathscr J,
 \tag{0.6.2}\label{sga3:iii:remark:0-6:equation:0-6-2}
\end{equation}
where, by an abuse of language (since \(X_0\) need not be an
\(S_0\)-group),
\(\underline{\operatorname{Lie}}(X_0/S_0)\) denotes the dual of the
\(\mathscr O_{S_0}\)-module \(\omega^1_{X_0/S_0}\).\footnote{Editor's
note: this is justified by \SGARefLink{sga3:ii:proposition:3-3}{II~3.3} and \SGARefLink{sga3:iii:paragraph:4-11}{4.11}. Indeed, the functor
\(L^{\varepsilon_0}_{X_0/S_0}\), the ``tangent space to \(X_0\) over
\(S_0\) at \(\varepsilon_0\),'' is represented over \(S_0\) by the
vector bundle
\(\mathbf V(\varepsilon_0^*(\Omega^1_{X_0/S_0}))
=\mathbf V(\omega^1_{X_0/S_0})\). Its sheaf of sections is the dual
module
\(\mathscr{H}om_{\mathscr O_{S_0}}
(\omega^1_{X_0/S_0},\mathscr O_{S_0})\); this is
\(\underline{\operatorname{Lie}}(X_0/S_0)\) when \(X_0\) is an
\(S_0\)-group and \(\varepsilon_0\) is the identity section.}

\SGARefLink{sga3:iii:proposition:0-2}{Proposition~0.2} (and its proof) has two important
corollaries.\footnote{Editor's note: \SGARefLink{sga3:iii:corollary:0-7-bis}{Corollary 0.7 bis}, which was used
implicitly in the proof of \SGARefLink{sga3:iii:proposition:0-8}{0.8}, has been added to \SGARefLink{sga3:iii:corollary:0-7}{0.7}. We point out
that \SGARefLink{sga3:iii:proposition:0-8}{0.8}--\SGARefLink{sga3:iii:corollary:0-12}{0.12}, as well as \SGARefLink{sga3:iii:corollary:0-17}{0.17}, which play an important technical role
later in this expos\'e, are consequences of \SGARefLink{sga3:iii:corollary:0-7}{0.7} and \SGARefLink{sga3:iii:corollary:0-7-bis}{0.7 bis}.}

\medskip
\noindent\SGAInlineRefTarget{sga3:iii:corollary:0-7}\textbf{Corollary 0.7.}
\textit{Let \(X\) be an \(S\)-scheme.}
\begin{enumerate}
 \item[\SGAInlineRefTarget{sga3:iii:corollary:0-7:item:a}(a)] \textit{Every \(S\)-endomorphism of \(X\) inducing the
 identity on \(X_{\mathscr J}\) is an automorphism.}

 \item[\SGAInlineRefTarget{sga3:iii:corollary:0-7:item:b}(b)] \textit{There is an exact sequence of groups}
 \[
  0\longrightarrow
  \operatorname{Hom}_{\mathscr O_{X_0}}
    (\Omega^1_{X_0/S_0},\mathscr J\mathscr O_X)
  \xrightarrow{\ i\ }
  \operatorname{Aut}_S(X)
  \longrightarrow
  \operatorname{Aut}_{S_{\mathscr J}}(X_{\mathscr J}).
 \]

 \item[\SGAInlineRefTarget{sga3:iii:corollary:0-7:item:c}(c)] \textit{Let \(\operatorname{Aut}_S(X)\) act on the first
 group by transport of structure. Then, for every
 \(u\in\operatorname{Aut}_S(X)\) and
 \(m\in\operatorname{Hom}_{\mathscr O_{X_0}}
 (\Omega^1_{X_0/S_0},\mathscr J\mathscr O_X)\),}
 \[
  i(u\cdot m)=u\,i(m)\,u^{-1}.
 \]
\end{enumerate}

% Source page: Exp3 component-local p. 16; combined-reader p. 126; printed p. 116.
\noindent\textit{Proof.}
By \SGARefLink{sga3:iii:proposition:0-2:item:i}{0.2(i)}, \(\operatorname{Hom}_{X^+}(X,X)\) is a principal
homogeneous set under \(\operatorname{Hom}_{X^+}(X,L_X)\), since it is
certainly nonempty: it contains the marked point given by the identity
automorphism of \(X\).\footnote{Editor's note: what follows has been
expanded.} Consequently, the map \(m\mapsto m\cdot\operatorname{id}_X\)
induces a bijection
\[
 \operatorname{Hom}_{\mathscr O_{X_0}}
   (\Omega^1_{X_0/S_0},\mathscr J\mathscr O_X)
 =L_X(X)
 \xrightarrow{\sim}
 \operatorname{Hom}_{X^+}(X,X).
\]

Let \(m\in L_X(X)\), and let
\(f=m'\cdot\operatorname{id}_X\) be an element of
\(\operatorname{Hom}_{X^+}(X,X)\). Applying \SGARefLink{sga3:iii:proposition:0-2:item:ii}{0.2(ii)} to \(f\) gives
\[
 f\circ(m\cdot\operatorname{id}_X)
 =L_f(X)(m)\cdot f
 =L_f(X)(m)\cdot(m'\cdot\operatorname{id}_X).
\]
On the other hand, because \(f\) is an \(X^+\)-endomorphism of \(X\),
one has \(f_{\mathscr J}=\operatorname{id}_{X_{\mathscr J}}\), hence
\(f_0=\operatorname{id}_{X_0}\). Since \(L_f\) depends only on
\(f_0\) (cf. the editor's note (17) in \SGARefLink{sga3:iii:corollary:0-1-10}{0.1.10}), it follows that
\(L_f(X)(m)=m\). Thus the preceding equality becomes
\[
 (m'\cdot\operatorname{id}_X)\circ
 (m\cdot\operatorname{id}_X)
 =m\cdot(m'\cdot\operatorname{id}_X)
 =(m+m')\cdot\operatorname{id}_X.
\]
This shows that the bijection
\(m\mapsto m\cdot\operatorname{id}_X\) carries the group law of
\(\operatorname{Hom}_{X^+}(X,L_X)\) to composition of
\(X^+\)-endomorphisms of \(X\).

It follows first that every element of
\(\operatorname{Hom}_{X^+}(X,X)\) is invertible, proving part~\SGARefLink{sga3:iii:corollary:0-7:item:a}{(a)}, and
then that there is an exact sequence
\[
 0\longrightarrow \operatorname{Hom}_{X^+}(X,L_X)
 \xrightarrow{\ i\ }\operatorname{Aut}_S(X)
 \longrightarrow\operatorname{Aut}_{S_{\mathscr J}}(X_{\mathscr J}),
\]
which proves part~\SGARefLink{sga3:iii:corollary:0-7:item:b}{(b)}.

The morphism \(i\) just defined is functorial in \(X\) with respect to
isomorphisms, because it is defined structurally from the action of
\(L_X\) on \(X\) over \(X^+\), an action that is itself functorial in
\(X\) by \SGARefLink{sga3:iii:proposition:0-2:item:ii}{Proposition~0.2(ii)}.\footnote{Editor's note: what follows has
been expanded.} Hence every automorphism \(u\) of \(X\) over \(S\)
induces, by transport of structure, isomorphisms
\[
 h:\operatorname{Hom}_{X^+}(X,L_X)
   \xrightarrow{\sim}\operatorname{Hom}_{X^+}(X,L_X)
 \quad\text{and}\quad
 f:\operatorname{Aut}_S(X)
   \xrightarrow{\sim}\operatorname{Aut}_S(X)
\]
such that the following diagram commutes.

\SGAThreeNativeDiagram[0.70\textwidth]
  {native_diagrams/exp3/diagram_022_exp3_p016_cor07_functoriality_square.tex}
  {Functoriality of \(i\) under transport of structure.}
  {fig:exp3-cor07-functoriality}
  {Polo--Gille Exp3-14oct24.pdf, component-local p. 16; combined-reader p. 126; printed p. 116}

Thus \(f\circ i=i\circ h\). On the other hand, \(f\) is described by
the following commutative diagram.

\SGAThreeNativeDiagram[0.40\textwidth]
  {native_diagrams/exp3/diagram_023_exp3_p016_cor07_conjugation_square.tex}
  {The automorphism \(f\) acts by conjugation with \(u\).}
  {fig:exp3-cor07-conjugation}
  {Polo--Gille Exp3-14oct24.pdf, component-local p. 16; combined-reader p. 126; printed p. 116}

% Source page: Exp3 component-local p. 17; combined-reader p. 127; printed p. 117.
In other words,
\[
 f(a)=u\circ a\circ u^{-1}
 \qquad\text{for every }a\in\operatorname{Aut}_S(X).
\]
Writing \(i(h(m))=f(i(m))\) now gives the asserted formula.

% Continue with Corollary 0.7 bis.

% END INCLUDED SGA3_English_current_buildable_source_20260823/source/inputs/III/tex/components/18_expose_III_s03_s07_flat_liftings.tex

% BEGIN INCLUDED SGA3_English_current_buildable_source_20260823/source/inputs/III/tex/components/19_expose_III_s07bis_s010_group_actions.tex
% SGA 3, Expose III, Loop 1 English reconstruction.
% Scope: Corollary 0.7 bis through Corollary 0.10 and its transition,
% stopping immediately before Proposition 0.11.
% Authority: Polo--Gille Exp3-14oct24.pdf, component-local pages 17--19
% (combined-reader pages 127--129; printed pages 117--119). Force-OCR is
% a drafting locator; English wording was compared with Jacob Reinhold's
% CC BY 4.0 Markdown lineage at commit
% e7a259f3f8608ad3edf9bf6eead3fd504dd2d23e. The PDF controls prose,
% mathematics, numbering, notes, and diagram images.

\setcounter{footnote}{26}

% Source page: Exp3 component-local p. 17; combined-reader p. 127; printed p. 117.
\medskip
\noindent\SGAInlineRefTarget{sga3:iii:corollary:0-7-bis}\textbf{Corollary 0.7 bis.}
\textit{Let \(X\) be an \(S\)-scheme such that \(X_{\mathscr J}\) is an
\(S_{\mathscr J}\)-monoid. Then \(L_X\) has the structure of an
\(S\)-monoid, and there is a split exact sequence of \(S\)-monoids:}

\SGAThreeNativeDiagram[0.50\textwidth]
  {native_diagrams/exp3/diagram_024_exp3_p017_corollary07bis_split_sequence.tex}
  {The split exact sequence of \(S\)-monoids in \SGARefLink{sga3:iii:corollary:0-7-bis}{Corollary 0.7 bis}.}
  {fig:exp3-corollary-07bis-split-sequence}
  {Polo--Gille Exp3-14oct24.pdf, component-local p. 17; combined-reader p. 127; printed p. 117}

\textit{The monoid law induced on \(L'_X\) coincides with its abelian-group
structure. In particular, if \(X_{\mathscr J}\) is an
\(S_{\mathscr J}\)-group, then \(L_X\) is an \(S\)-group and is the
semidirect product of \(X^+\) and \(L'_X\).}

\noindent\textit{Proof.}
Since \(X_{\mathscr J}\) is an \(S_{\mathscr J}\)-monoid,
\[
 X^+=\prod_{S_{\mathscr J}/S}X_{\mathscr J}
\]
is an \(S\)-monoid; indeed,
\(X^+(Y)=X_{\mathscr J}(Y_{\mathscr J})\) for every \(Y\to S\).
For every \(S\)-scheme \(Y\), let \(\widetilde Y_{\mathscr J}\) denote
the affine \(Y_{\mathscr J}\)-scheme corresponding to the quasi-coherent
\(\mathscr O_{Y_{\mathscr J}}\)-algebra
\[
 \mathscr O_{Y_{\mathscr J}}\oplus\mathscr J\mathscr O_Y,
\]
that is, the graded algebra associated with the filtration
\(\mathscr O_Y\supset\mathscr J\mathscr O_Y\). Then \(L_X(Y)\) identifies
with \(X_{\mathscr J}(\widetilde Y_{\mathscr J})\), and \(L'_X(Y)\)
identifies with the kernel of the morphism
\[
 p:X_{\mathscr J}(\widetilde Y_{\mathscr J})
   \longrightarrow X_{\mathscr J}(Y_{\mathscr J})
\]
induced by the \emph{zero section}
\(Y_{\mathscr J}\to\widetilde Y_{\mathscr J}\). Equivalently, that
section is induced by the morphism of
\(\mathscr O_{Y_{\mathscr J}}\)-algebras
\[
 \mathscr O_{\widetilde Y_{\mathscr J}}
   \longrightarrow\mathscr O_{Y_{\mathscr J}}
\]
that vanishes on the ideal \(\mathscr J\mathscr O_Y\). Thus, for every
\(Y\to S\), there is a split exact sequence of monoids, functorial in
\(Y\):

\SGAThreeNativeDiagram[0.60\textwidth]
  {native_diagrams/exp3/diagram_025_exp3_p017_corollary07bis_pointwise_split_sequence.tex}
  {The pointwise split exact sequence for \(L_X(Y)\).}
  {fig:exp3-corollary-07bis-pointwise-sequence}
  {Polo--Gille Exp3-14oct24.pdf, component-local p. 17; combined-reader p. 127; printed p. 117}

It remains to show that the monoid law induced on \(L'_X\) coincides
with its abelian-group structure. Let \(\mu\) be the monoid law of
\(L_X\), and let \(e\) be its unit section. We must prove that, for
every \(m,m'\in L'_X(Y)\),
\[
 \mu(m\mathbin{\cdot}e,m'\mathbin{\cdot}e)
   =(m+m')\mathbin{\cdot}e.
\]

This may be seen in either of two ways. First, repeat the proof of
(\SGARefLink{sga3:iii:corollary:0-1-10}{0.1.10} \((\ast)\)), replacing the morphism \(f:X\to W\) there by
\[
 \mu:L_X\times_S L_X\longrightarrow L_X.
\]
Identify \(X^+(Y)=X_{\mathscr J}(Y_{\mathscr J})\) with its image under
\(s\) in \(L_X(Y)=X_{\mathscr J}(\widetilde Y_{\mathscr J})\). For all
\(g,g'\in X_{\mathscr J}(Y_{\mathscr J})\) and
\(m,m'\in L'_X(Y)\), one then obtains
\begin{equation}
 \mu(m\mathbin{\cdot}g,m'\mathbin{\cdot}g')
 =L_\mu^{(g,g')}(m,m')\mathbin{\cdot}\mu(g,g').
 \tag{$\star$}\label{sga3:iii:corollary:0-7-bis:equation:star}
\end{equation}
Here \(L_\mu^{(g,g')}\) denotes the morphism derived from \(\mu\) at the
point \((g,g')\); in other words,
\(\widetilde Y_{\mathscr J}\) lies over \(L_X\times_S L_X\) through
\((g,g')\). In particular,
\[
 \mu(m\mathbin{\cdot}e,m'\mathbin{\cdot}e)
 =L_\mu^{(e,e)}(m,m')\mathbin{\cdot}e.
\]
But
\[
 L_\mu^{(e,e)}(m,m')=L_{\ell_e}(m')+L_{r_e}(m),
\]
where \(\ell_e\), respectively \(r_e\), denotes left, respectively
right, translation by \(e\). Each is the identity map of
\(X_{\mathscr J}\), and therefore
\(L_\mu^{(e,e)}(m,m')=m+m'\).

Alternatively, one may argue as follows (cf. the proof of [DG70],
\S II.4, \SGARefLink{sga3:iii:theorem:3-5}{Theorem~3.5}). By \SGARefLink{sga3:iii:lemma:0-1-11}{Lemma~0.1.11}, formation of \(X^+\) and of
\(L_X\) commutes with products, and hence so does formation of \(L'_X\).
It follows that the morphism
\[
 \mu':L'_X\times L'_X\longrightarrow L'_X
\]
induced by \(\mu\) is a homomorphism for the abelian-group structure.
\SGARefLink{sga3:iii:corollary:3-10}{Lemma~3.10} of Expos\'e~II then shows that \(\mu'\) coincides with the
abelian-group law.

% Source page: Exp3 component-local p. 18; combined-reader p. 128; printed p. 118.
\medskip
\noindent\SGAInlineRefTarget{sga3:iii:paragraph:0-8}\textbf{0.8.}\footnote{Editor's note: the number 0.8 has been
inserted here to mark the return to the original text.}
Now let \(X\) be an \(S\)-scheme such that \(X_{\mathscr J}\) is an
\(S_{\mathscr J}\)-group. Suppose that there is an \(S\)-morphism
\[
 P:X\times_S X\longrightarrow X
\]
whose base change
\[
 P_{\mathscr J}:
 X_{\mathscr J}\times_{S_{\mathscr J}}X_{\mathscr J}
 \longrightarrow X_{\mathscr J}
\]
is the group law of \(X_{\mathscr J}\). An important special case is
that in which \(X\) is an \(S\)-group and \(P\) is its group law. We
obtain a morphism
\[
 L_P:L_X\times_S L_X\simeq L_{X\times_S X}\longrightarrow L_X.
\]
In fact, this morphism does not depend on \(P\), since it can be
computed from the group law \(P_{\mathscr J}\) of \(X_{\mathscr J}\),
as we now show.\footnote{Editor's note: the original has been expanded
in what follows.}

By \SGARefLink{sga3:iii:proposition:0-2:item:ii}{Proposition~0.2(ii)} and (iii), for every \(Y\to S\), every
\(x,x'\in X(Y)\), and every \(m,m'\in L'_X(Y)\),
\[
\begin{aligned}
 P(m\mathbin{\cdot}x,m'\mathbin{\cdot}x')
 &=P\bigl((m,m')\mathbin{\cdot}(x,x')\bigr)\\
 &=L_P^{(x,x')}(m,m')\mathbin{\cdot}\mu(g,g'),
\end{aligned}
\]
where \(g\), respectively \(g'\), is the image of \(x\), respectively
\(x'\), in \(X^+(Y)\). Moreover (cf. the proof of \SGARefLink{sga3:iii:corollary:0-10}{0.10}),
\(L_P^{(x,x')}=L_\mu^{(g,g')}\), and by \SGARefLink{sga3:iii:corollary:0-7-bis}{0.7 bis} \((\star)\) this is the
element of \(L'_X(Y)\) defined by the equality
\[
 L_\mu^{(g,g')}(m,m')\mathbin{\cdot}\mu(g,g')
   =\mu(m\mathbin{\cdot}g,m'\mathbin{\cdot}g')
\]
in \(L_X(Y)\).

Write \(\times\), rather than \(\mu\), for the group law of \(L_X\),
and write \(\operatorname{Ad}\) for the \emph{adjoint action} of \(X^+\)
on \(L'_X\). This action factors through \(X_0\) and is induced by the
adjoint action of \(X_0\) on \(\omega^1_{X_0/S_0}\). Then
\[
\begin{aligned}
 L_\mu^{(g,g')}(m,m')\times g\times g'
 &=m\times g\times m'\times g'\\
 &=\bigl(m\times\operatorname{Ad}(g)(m')\bigr)\times g\times g',
\end{aligned}
\]
and hence
\[
 L_P^{(x,x')}(m,m')=m\times\operatorname{Ad}(g)(m').
\]
We have therefore proved:

\medskip
\noindent\SGAInlineRefTarget{sga3:iii:proposition:0-8}\textbf{Proposition 0.8.}
\textit{Let \(P:X\times_S X\to X\) be an \(S\)-morphism such that
\(P_{\mathscr J}\) equips \(X_{\mathscr J}\) with an
\(S_{\mathscr J}\)-group structure. Let \(\times\) denote the group law
of \(L'_X\), let
\((m,x)\mapsto m\mathbin{\cdot}x\) denote the morphism
\(L'_X\times_S X\to X\) defining the action of \(L'_X\) on \(X\), and
let}
\[
 \operatorname{Ad}:X^+\longrightarrow
 \operatorname{Aut}_{S\text{-}\mathrm{gr}}(L'_X)
\]
\textit{be the adjoint action of \(X^+\) on \(L'_X\), induced by the
adjoint action of \(X_0\) on \(\omega^1_{X_0/S_0}\). Then, for every
\(S'\to S\), every \(x,x'\in X(S')\), and every
\(m,m'\in L'_X(S')\),}
\begin{equation}
 P(m\mathbin{\cdot}x,m'\mathbin{\cdot}x')
 =\bigl(m\times\operatorname{Ad}(p_X(x))(m')\bigr)
    \mathbin{\cdot}P(x,x').
 \tag{0.8.1}\label{sga3:iii:proposition:0-8:equation:0-8-1}
\end{equation}

\textit{If \(X\) is an \(S\)-group, we shall write \(\ast\) for its law,
\(e\) for its unit section, and \(i\) for the \(S\)-morphism defined by}
\[
 i(m)=m\mathbin{\cdot}e
\]
\textit{for every \(S'\to S\) and every \(m\in L'_X(S')\).}

% Source page: Exp3 component-local p. 19; combined-reader p. 129; printed p. 119.
\medskip
\noindent\SGAInlineRefTarget{sga3:iii:corollary:0-9}\textbf{Corollary 0.9.}
\textit{Let \(X\) be an \(S\)-group. Then \(X^+\) is naturally an
\(S\)-group, and \(p_X\) is a morphism of \(S\)-groups. Moreover, the
\(S\)-morphism}
\[
 i:L'_X\longrightarrow X,
 \qquad m\longmapsto m\mathbin{\cdot}e,
\]
\textit{is an isomorphism of \(S\)-groups from \(L'_X\) onto
\(\operatorname{Ker}(p_X)\), and, for every \(S'\to S\), every
\(x'\in X(S')\), and every \(m\in L'_X(S')\),}
\begin{equation}
 m\mathbin{\cdot}x'
 =(m\mathbin{\cdot}e)\ast x'
 =i(m)\ast x'.
 \tag{0.9.1}\label{sga3:iii:corollary:0-9:equation:0-9-1}
\end{equation}

The first two assertions were already proved in \SGARefLink{sga3:iii:remark:0-1-2}{0.1.2}. Since \(X\) is
formally principal homogeneous over \(X^+\) under
\(L_X=L'_X\times_S X^+\), the morphism \(i\) is indeed an isomorphism
of \(S\)-functors from \(L'_X\) onto the kernel of \(p_X\). The fact
that \(i\) is a group morphism and formula (\SGARefLink{sga3:iii:corollary:0-9:equation:0-9-1}{0.9.1}) follow from (\SGARefLink{sga3:iii:proposition:0-8:equation:0-8-1}{0.8.1}),
applied respectively with \(x=x'=e\), and with \(x=e,\ m'=1\).

\medskip
\noindent\SGAInlineRefTarget{sga3:iii:corollary:0-10}\textbf{Corollary 0.10.}
\textit{Let \(X\) be an \(S\)-group. With the preceding notation, for
every \(S'\to S\), every \(x\in X(S')\), and every
\(m'\in L'_X(S')\),}
\begin{equation}
 x\ast i(m')\ast x^{-1}
 =i\bigl(\operatorname{Ad}(p_X(x))(m')\bigr).
 \tag{0.10.1}\label{sga3:iii:corollary:0-10:equation:0-10-1}
\end{equation}
This follows from the equality
\(i(m')\ast x^{-1}=m'\mathbin{\cdot}x^{-1}\) and from (\SGARefLink{sga3:iii:proposition:0-8:equation:0-8-1}{0.8.1}), applied with
\(m=1\) and \(x'=x^{-1}\).

Thus, when \(X\) is an \(S\)-group, we have explicitly determined the
kernel of \(X\to X^+\) and the action of the inner automorphisms of
\(X\) on this kernel. We shall now see that the same can be done for
certain, not necessarily representable, \(S\)-group functors. One case
will be useful to us, namely the functors
\(\underline{\operatorname{Aut}}\) of \SGARefLink{sga3:i:subsection:1:7}{Expos\'e~I, 1.7}. We state it at
once in the next proposition.

% Continue with Proposition 0.11.

% END INCLUDED SGA3_English_current_buildable_source_20260823/source/inputs/III/tex/components/19_expose_III_s07bis_s010_group_actions.tex

% BEGIN INCLUDED SGA3_English_current_buildable_source_20260823/source/inputs/III/tex/components/20_expose_III_s011_end_s0_automorphisms.tex
% SGA 3, Expose III, Loop 1 English reconstruction.
% Scope: Proposition 0.11 through Corollary 0.17, completing section 0
% and stopping immediately before the section 1 heading.
% Authority: Polo--Gille Exp3-14oct24.pdf, component-local pages 19--22
% (combined-reader pages 129--132; printed pages 119--122). Force-OCR is
% a drafting locator; English wording was compared with Jacob Reinhold's
% CC BY 4.0 Markdown lineage at commit
% e7a259f3f8608ad3edf9bf6eead3fd504dd2d23e. The PDF controls prose,
% mathematics, numbering, notes, and the two temporary diagram images.

\setcounter{footnote}{28}

% Source page: Exp3 component-local p. 19; combined-reader p. 129; printed p. 119.
\medskip
\noindent\SGAInlineRefTarget{sga3:iii:proposition:0-11}\textbf{Proposition 0.11.}
\textit{Let \(E\) be an \(S\)-scheme, and put
\(X=\underline{\operatorname{Aut}}_S(E)\). The kernel of the morphism of
\(S\)-group functors}
\[
 p_X:X\longrightarrow X^+
\]
\textit{is canonically identified with the commutative \(S\)-group functor
\(L'_X\) defined by}
\[
 \operatorname{Hom}_S(Y,L'_X)
 =\operatorname{Hom}_{\mathscr O_{E_0\times_{S_0}Y_0}}
 \left(
   \Omega^1_{E_0/S_0}\otimes_{\mathscr O_{S_0}}\mathscr O_{Y_0},
   \mathscr J\mathscr O_{E\times_S Y}
 \right),
\]
\textit{where \(Y\) denotes a variable \(S\)-scheme.}

Indeed, for a variable \(S\)-scheme \(Y\), we have
\[
 \operatorname{Hom}_S(Y,X)=\operatorname{Aut}_Y(E\times_S Y)
\]
and
\begin{align*}
 \operatorname{Hom}_S(Y,X^+)
  &=\operatorname{Hom}_S(Y_{\mathscr J},X)\\
  &=\operatorname{Aut}_{Y_{\mathscr J}}(E\times_S Y_{\mathscr J})\\
  &=\operatorname{Aut}_{Y_{\mathscr J}}
    \bigl((E\times_S Y)\times_Y Y_{\mathscr J}\bigr).
\end{align*}
Applying \SGARefLink{sga3:iii:corollary:0-7}{0.7}~\SGARefLink{sga3:iii:corollary:0-7:item:b}{(b)} to the \(Y\)-scheme \(E\times_S Y\), we obtain an
isomorphism of groups
\[
 \operatorname{Hom}_S(Y,L'_X)
 \simeq
 \operatorname{Ker}\!\left(
   \operatorname{Hom}_S(Y,X)\longrightarrow
   \operatorname{Hom}_S(Y,X^+)
 \right).
\]
This isomorphism is readily checked to be functorial in the \(S\)-scheme
\(Y\). Hence it yields an isomorphism of \(S\)-groups
\[
 L'_X\simeq\operatorname{Ker}(X\longrightarrow X^+),
\]
which completes the proof of \SGARefLink{sga3:iii:proposition:0-11}{Proposition 0.11}.

% Source page: Exp3 component-local p. 20; combined-reader p. 130; printed p. 120.
\medskip
\noindent\SGAInlineRefTarget{sga3:iii:corollary:0-12}\textbf{Corollary 0.12.}\footnote{Editor's note: ``Remark'' has
been changed to ``Corollary''.}
\textit{Retain the notation of \SGARefLink{sga3:iii:proposition:0-11}{0.11}: \(E\) is an \(S\)-scheme and
\(X=\underline{\operatorname{Aut}}_S(E)\). There is a natural action \(f\)
of \(X\) on \(L'_X\), defined as follows. For every \(S\)-scheme \(Y\),}
\[
 \operatorname{Hom}_S(Y,X)=\operatorname{Aut}_Y(E\times_S Y)
\]
\textit{and}
\[
 \operatorname{Hom}_S(Y,L'_X)
 =\operatorname{Hom}_{\mathscr O_{E_0\times_{S_0}Y_0}}
 \left(
   \Omega^1_{E_0\times_{S_0}Y_0/Y_0},
   \mathscr J\mathscr O_{E\times_S Y}
 \right).
\]
\textit{Here}
\[
 \Omega^1_{E_0/S_0}\otimes_{\mathscr O_{S_0}}\mathscr O_{Y_0}
 \simeq \Omega^1_{E_0\times_{S_0}Y_0/Y_0}
\]
\textit{(see EGA~IV\textsubscript{4}, 16.4.5). The first group acts on
the second by transport of structure, and this action is functorial in
\(Y\). One then has}
\begin{equation}
 x\,i(m)\,x^{-1}=i\bigl(f(x)m\bigr),
 \tag{0.12.1}\label{sga3:iii:corollary:0-12:equation:0-12-1}
\end{equation}
\textit{for every \(Y\to S\), \(x\in\operatorname{Hom}_S(Y,X)\), and
\(m\in\operatorname{Hom}_S(Y,L'_X)\).}

Indeed, this follows from \SGARefLink{sga3:iii:corollary:0-7}{0.7}~\SGARefLink{sga3:iii:corollary:0-7:item:c}{(c)}, applied to the \(Y\)-scheme
\(E\times_S Y\).

\medskip
\noindent\SGAInlineRefTarget{sga3:iii:paragraph:0-13}\textbf{Reminder 0.13.}
The direct image of a quasi-coherent module under a morphism of finite
presentation is quasi-coherent. Under the same hypotheses, formation of
the direct image commutes with flat base change. In the situation

\SGAThreeNativeDiagram[0.50\textwidth]
  {native_diagrams/exp3/diagram_026_exp3_p020_reminder013_flat_base_change_square.tex}
  {The flat base-change square in Reminder \SGARefLink{sga3:iii:paragraph:0-13}{0.13}.}
  {fig:exp3-reminder-013-flat-base-change}
  {Polo--Gille Exp3-14oct24.pdf, component-local p. 20; combined-reader p. 130; printed p. 120}

if \(f\) (and hence \(f'\)) is of finite presentation and \(g\) (and
hence \(g'\)) is flat, then for every quasi-coherent
\(\mathscr O_T\)-module \(\mathscr F\),
\[
 f_*(\mathscr F)\otimes_{\mathscr O_S}\mathscr O_{S'}
 =f'_*\bigl(\mathscr F\otimes_{\mathscr O_S}\mathscr O_{S'}\bigr),
\]
or, more elegantly,
\[
 g^*\bigl(f_*\mathscr F\bigr)=f'_*\bigl(g'^*\mathscr F\bigr).
\]
These two facts hold more generally for a quasi-compact and
quasi-separated morphism \(f\); see EGA~I, 9.2.1, and
EGA~III\textsubscript{1}, 1.4.15 in the quasi-compact separated case
(taking EGA~III\textsubscript{2}, Err\textsubscript{III} 25 into
account), as well as EGA~IV\textsubscript{1}, 1.7.4 and 1.7.21.

\medskip
\noindent\SGAInlineRefTarget{sga3:iii:remark:0-14}\textbf{Remark 0.14.}\footnote{Editor's note: The first part of
this remark has been expanded, and a second part has been added.}
Return to the notation of \SGARefLink{sga3:iii:proposition:0-11}{0.11}: let \(E\) be an \(S\)-scheme, put
\(X=\underline{\operatorname{Aut}}_S(E)\), and let \(L'_X\) be the
commutative \(S\)-group functor defined by
\begin{align*}
 L'_X(Y)
 &=\operatorname{Hom}_{\mathscr O_{E_0\times_{S_0}Y_0}}
   \left(
     \Omega^1_{E_0/S_0}\otimes_{\mathscr O_{S_0}}\mathscr O_{Y_0},
     \mathscr J\mathscr O_{E\times_S Y}
   \right)\\
 &=\operatorname{Hom}_{\mathscr O_{E\times_S Y}}
   \left(
     \Omega^1_{E/S}\otimes_{\mathscr O_S}\mathscr O_Y,
     \mathscr J\mathscr O_{E\times_S Y}
   \right)\\
 &=\Gamma\!\left(
     E\times_S Y,
     \widetilde{\operatorname{Hom}}_{\mathscr O_{E\times_S Y}}
     \left(
       \Omega^1_{E/S}\otimes_{\mathscr O_S}\mathscr O_Y,
       \mathscr J\mathscr O_{E\times_S Y}
     \right)
   \right).
\end{align*}

% Source page: Exp3 component-local p. 21; combined-reader p. 131; printed p. 121.
Suppose that \(Y\) is flat over \(S\). There are then isomorphisms
\[
 \mathscr J\mathscr O_{E\times_S Y}
 \xleftarrow{\ \sim\ }
 (\mathscr J\mathscr O_E)\otimes_{\mathscr O_S}\mathscr O_Y
 \xrightarrow{\ \sim\ }
 (\mathscr J\mathscr O_E)\otimes_{\mathscr O_{S_0}}\mathscr O_{Y_0}.
\]
Suppose in addition that \(E\) is of finite presentation over \(S\).
Then \(\Omega^1_{E/S}\) is an \(\mathscr O_E\)-module of finite
presentation (EGA~IV\textsubscript{4}, 16.4.22), and therefore, by
EGA~0\textsubscript{I}, 6.7.6,
\begin{align*}
 &\widetilde{\operatorname{Hom}}_{\mathscr O_{E\times_S Y}}
 \left(
   \Omega^1_{E/S}\otimes_{\mathscr O_S}\mathscr O_Y,
   (\mathscr J\mathscr O_E)\otimes_{\mathscr O_S}\mathscr O_Y
 \right)\\
 &\qquad\simeq
 \widetilde{\operatorname{Hom}}_{\mathscr O_E}
 \left(\Omega^1_{E/S},\mathscr J\mathscr O_E\right)
 \otimes_{\mathscr O_S}\mathscr O_Y.
\end{align*}
Let \(\pi:E\to S\) and \(g:Y\to S\) be the structural morphisms.
Apply \SGARefLink{sga3:iii:paragraph:0-13}{0.13} to the diagram

\SGAThreeNativeDiagram[0.42\textwidth]
  {native_diagrams/exp3/diagram_027_exp3_p021_remark014_structural_base_change_square.tex}
  {The structural base-change square used in \SGARefLink{sga3:iii:remark:0-14}{Remark 0.14}.}
  {fig:exp3-remark-014-structural-base-change}
  {Polo--Gille Exp3-14oct24.pdf, component-local p. 21; combined-reader p. 131; printed p. 121}

and to the \(\mathscr O_E\)-module
\[
 \mathscr F=
 \widetilde{\operatorname{Hom}}_{\mathscr O_E}
 \left(\Omega^1_{E/S},\mathscr J\mathscr O_E\right).
\]
We obtain
\begin{align*}
 \Gamma(E\times_S Y,g'^*\mathscr F)
 &=\Gamma(Y,\pi'_*g'^*\mathscr F)\\
 &=\Gamma(Y,g^*\pi_*\mathscr F)\\
 &=W(\pi_*\mathscr F)(Y).
\end{align*}
Thus, if \(E\) is of finite presentation over \(S\), we have
\begin{equation}
 L'_X=
 W\!\left(
   \pi_*\widetilde{\operatorname{Hom}}_{\mathscr O_E}
   (\Omega^1_{E/S},\mathscr J\mathscr O_E)
 \right)
 \tag{0.14.1}\label{sga3:iii:remark:0-14:equation:0-14-1}
\end{equation}
on the category of \(S\)-schemes flat over \(S\). Notice also that the
module to which \(W\) is applied is quasi-coherent, by EGA~I, 9.1.1
and 9.2.1.

Let \(L_0\) be the \(S_0\)-functor\footnote{Editor's note: What follows
has been added.}
\[
 L_0=
 W\!\left(
   \pi_{0*}\widetilde{\operatorname{Hom}}_{\mathscr O_{E_0}}
   (\Omega^1_{E_0/S_0},\mathscr J\mathscr O_E)
 \right).
\]
Returning to the definition of \(L'_X(Y)\), and using the isomorphism
\[
 \mathscr J\mathscr O_{E\times_S Y}
 \simeq(\mathscr J\mathscr O_E)
       \otimes_{\mathscr O_{S_0}}\mathscr O_{Y_0},
\]
the same reasoning gives
\[
 L'_X(Y)=L_0(Y_0)=L_0(Y\times_S S_0)
 =\left(\prod_{S_0/S}L_0\right)(Y).
\]
Consequently, on the category of \(S\)-schemes flat over \(S\),
\[
 L'_X=
 \prod_{S_0/S}
 W\!\left(
   \pi_{0*}\widetilde{\operatorname{Hom}}_{\mathscr O_{E_0}}
   (\Omega^1_{E_0/S_0},\mathscr J\mathscr O_E)
 \right).
\]
It is not obvious that the action of \(X\) on \(L'_X\) defined in \SGARefLink{sga3:iii:corollary:0-12}{0.12}
comes from an action of
\(X_0=\underline{\operatorname{Aut}}_{S_0}(E_0)\) on \(L_0\).
This is nevertheless the case when \(E\) is also flat over \(S\).

% Source page: Exp3 component-local p. 22; combined-reader p. 132; printed p. 122.
Indeed, in that case there are canonical isomorphisms
\[
 \mathscr J\mathscr O_E
 \simeq\mathscr J\otimes_{\mathscr O_S}\mathscr O_E
 \simeq\mathscr J\otimes_{\mathscr O_{S_0}}\mathscr O_{E_0},
\]
and hence
\[
 L_0\simeq
 W\!\left(
   \pi_{0*}\widetilde{\operatorname{Hom}}_{\mathscr O_{E_0}}
   \left(
     \Omega^1_{E_0/S_0},
     \mathscr J\otimes_{\mathscr O_{S_0}}\mathscr O_{E_0}
   \right)
 \right),
\]
and
\begin{equation}
 L'_X=
 \prod_{S_0/S}
 W\!\left(
   \pi_{0*}\widetilde{\operatorname{Hom}}_{\mathscr O_{E_0}}
   \left(
     \Omega^1_{E_0/S_0},
     \mathscr J\otimes_{\mathscr O_{S_0}}\mathscr O_{E_0}
   \right)
 \right).
 \tag{0.14.2}\label{sga3:iii:remark:0-14:equation:0-14-2}
\end{equation}
For every \(S_0\)-scheme \(T\), we then have
\[
 L_0(T)\simeq
 \operatorname{Hom}_{\mathscr O_{E_0\times_{S_0}T}}
 \left(
   \Omega^1_{E_0\times_{S_0}T/T},
   \mathscr J\otimes_{\mathscr O_{S_0}}
     \mathscr O_{E_0\times_{S_0}T}
 \right).
\]
Moreover,
\[
 \operatorname{Hom}_{S_0}(T,X_0)
 =\operatorname{Aut}_T(E_0\times_{S_0}T)
\]
acts on \(L_0(T)\) by transport of structure, functorially in \(T\).
Finally, for every \(S\)-scheme \(Y\) flat over \(S\), the action by
transport of structure of
\[
 \operatorname{Hom}_S(Y,X)=\operatorname{Aut}_Y(E\times_S Y)
\]
on \(L'_X(Y)=L_0(Y_0)\) factors through
\(\operatorname{Aut}_{Y_0}(E_0\times_{S_0}Y_0)\).

Finally, let us extract the following two propositions from SGA~1,
Expos\'e~III.

\medskip
\noindent\SGAInlineRefTarget{sga3:iii:proposition:0-15}\textbf{Proposition 0.15.}
\textit{(SGA~1, Expos\'e~III, 6.8)}\footnote{Editor's note: This
``global'' lifting result uses the ``local'' lifting result of loc.
cit., 4.1, which is stated for \(S\) locally Noetherian; see
EGA~IV\textsubscript{4}, 18.1.1 for the general case.}
\textit{For every affine \(S_{\mathscr J}\)-scheme \(Y\) smooth over
\(S_{\mathscr J}\), there exists an \(S\)-scheme \(X\), smooth over
\(S\), such that \(X\times_S S_{\mathscr J}\simeq Y\); such an \(X\) is
unique up to a non-unique isomorphism.}

\medskip
\noindent\SGAInlineRefTarget{sga3:iii:proposition:0-16}\textbf{Proposition 0.16.}
\textit{(SGA~1, Expos\'e~III, 5.5)}\footnote{Editor's note: This is an
immediate consequence of the definition of ``(formally) smooth''
adopted in EGA~IV\textsubscript{4}, 17.3.1 (and 17.1.1).}
\textit{Let \(X\) be an \(S\)-scheme smooth over \(S\). For every affine
\(S\)-scheme \(Y\), the canonical map}
\[
 p_X(Y):\operatorname{Hom}_S(Y,X)
 \longrightarrow\operatorname{Hom}_S(Y,X^+)
 =\operatorname{Hom}_{S_{\mathscr J}}(Y_{\mathscr J},X_{\mathscr J})
\]
\textit{is surjective.}

\medskip
\noindent\SGAInlineRefTarget{sga3:iii:corollary:0-17}\textbf{Corollary 0.17.}
\textit{Let \(E\) be an affine \(S\)-scheme smooth over \(S\), and put
\(X=\underline{\operatorname{Aut}}_S(E)\). For every affine \(S\)-scheme
\(Y\), the canonical map}
\begin{align*}
 \operatorname{Aut}_Y(E\times_S Y)
  =\operatorname{Hom}_S(Y,X)
 &\longrightarrow \operatorname{Hom}_S(Y,X^+)\\
  &=\operatorname{Aut}_{Y_{\mathscr J}}
    (E_{\mathscr J}\times_{S_{\mathscr J}}Y_{\mathscr J})
\end{align*}
\textit{is surjective.}

Indeed, \(Y\times_S E\) is affine over the affine scheme \(Y\), hence
is affine. Applying \SGARefLink{sga3:iii:proposition:0-16}{0.16}, we see that every \(S_{\mathscr J}\)-morphism
\[
 Y_{\mathscr J}\times_{S_{\mathscr J}}E_{\mathscr J}
 \longrightarrow E_{\mathscr J}
\]
extends to an \(S\)-morphism
\[
 Y\times_S E\longrightarrow E.
\]
In other words,\footnote{Editor's note: The original has been modified
slightly so that the hypotheses of \SGARefLink{sga3:iii:corollary:0-7}{0.7} apply exactly.} every
\(Y_{\mathscr J}\)-endomorphism of
\(Y_{\mathscr J}\times_{S_{\mathscr J}}E_{\mathscr J}\) lifts to a
\(Y\)-endomorphism of \(Y\times_S E\). Then \SGARefLink{sga3:iii:corollary:0-7}{0.7}~\SGARefLink{sga3:iii:corollary:0-7:item:a}{(a)} shows that every
\(Y_{\mathscr J}\)-automorphism of
\(Y_{\mathscr J}\times_{S_{\mathscr J}}E_{\mathscr J}\) lifts to a
\(Y\)-automorphism of \(Y\times_S E\), which is the asserted property.

% Continue with section 1 (combined-reader p. 133).

% END INCLUDED SGA3_English_current_buildable_source_20260823/source/inputs/III/tex/components/20_expose_III_s011_end_s0_automorphisms.tex

% BEGIN INCLUDED SGA3_English_current_buildable_source_20260823/source/inputs/III/tex/components/21_expose_III_s11_s12_extensions_cohomology.tex
% SGA 3, Expose III, Loop 1 English reconstruction.
% Scope: Section 1 through Proposition 1.2.4(iii), stopping immediately
% before Section 1.3. The final item occupies the top of source p. 28.
% Authority: Polo--Gille Exp3-14oct24.pdf, component-local pages 23--28
% (p. 28 partial; combined-reader pages 133--138, p. 138 partial; printed
% pages 123--128). Force-OCR is a drafting locator; English wording was
% compared with Jacob Reinhold's CC BY 4.0 Markdown lineage at commit
% e7a259f3f8608ad3edf9bf6eead3fd504dd2d23e. The PDF controls prose,
% mathematics, numbering, notes, and diagram images.

\setcounter{footnote}{34}

% Source page: Exp3 component-local p. 23; combined-reader p. 133; printed p. 123.
\section*{\SGAInlineRefTarget{sga3:iii:section:1}1. Extensions and Cohomology}
\addcontentsline{toc}{section}{1. Extensions and Cohomology}

\subsection*{\SGAInlineRefTarget{sga3:iii:subsection:1-1}1.1.}
Let \(\mathscr C\) be a category stable under fiber
products.\footnote{Editor's note: the hypothesis that \(\mathscr C\)
is stable under fiber products has been added. It holds in the
applications in which \(\mathscr C\) is the category of schemes
\((\mathbf{Sch})\), or the category of functors
\((\mathbf{Sch})^{\circ}\to(\mathbf{Ens})\). If this hypothesis is
omitted, one must instead assume below that the fiber products
\(G\times_S\cdots\times_S G\) are representable.}
Let \(S\) be an object of \(\mathscr C\), let \(G\) be a representable
\(S\)-group, and let \(F\) be an \(S\)-functor in commutative groups
on which \(G\) acts. The cohomology groups \(H^n(G,F)\) were defined
in \SGARefLink{sga3:i:subsection:5-1}{I~5.1}. Recall that they are the cohomology groups of a complex
denoted \(C^\ast(G,F)\). Writing
\[
 (G/S)^n=G\times_S\cdots\times_S G
 \qquad(n\text{ factors}),
\]
one has
\[
 C^n(G,F)=\operatorname{Hom}_S((G/S)^n,F).
\]
Since \(G\), and hence each \((G/S)^n\), is representable, one also has
\[
 C^n(G,F)=F((G/S)^n).
\]
It follows from this and from the definition of the coboundary operator
that the complex \(C^\ast(G,F)\) depends only on the restriction of
\(F\) to the full subcategory of \(\mathscr C/S\) whose objects are
the Cartesian powers of \(G\) over \(S\). Consequently:

\medskip
\noindent\SGAInlineRefTarget{sga3:iii:lemma:1-1-1}\textbf{Lemma 1.1.1.}
\textit{Let \(\mathscr C\) be a category stable under fiber products,
let \(S\) be an object of \(\mathscr C\), and let \(G\) be a
representable \(S\)-group. Denote by \(\mathscr C(G)\) the full
subcategory of \(\mathscr C/S\) whose objects are the Cartesian powers
of \(G\) over \(S\). Let \(F\) and \(F'\) be two \(S\)-functors in
commutative groups on which \(G\) acts. If \(F\) and \(F'\) have the
same restriction to \(\mathscr C(G)\), there is a canonical
isomorphism}
\[
 H^\ast(G,F)\xrightarrow{\sim}H^\ast(G,F').
\]

\subsection*{\SGAInlineRefTarget{sga3:iii:subsection:1-1-2}1.1.2. Cohomology and restriction of scalars.\protect\footnotemark}
\footnotetext{Editor's note: the heading of this paragraph has been added.}
We record another comparison result. Let \(T\to S\) be a morphism in
\(\mathscr C\). If \(F\) is a \(T\)-functor in commutative groups,
then the functor obtained by restriction of scalars (cf. \SGARefLink{sga3:ii:section:1}{Expos\'e~II,
1}),
\[
 F_1=\prod_{T/S}F,
\]
is an \(S\)-functor in commutative groups, and there is a morphism of
\(S\)-functors in groups\footnotemark
\[
 u:\prod_{T/S}
 \underline{\operatorname{Aut}}_{T\text{-gr.}}(F)
 \longrightarrow
 \underline{\operatorname{Aut}}_{S\text{-gr.}}(F_1).
\]
\footnotetext{Editor's note: explicitly, let \(S'\to S\) and
\(\alpha\in
\underline{\operatorname{Aut}}_{(T\times_S S')\text{-gr.}}
(F_{T\times_S S'})\). For every \(U\to T\times_S S'\), let
\(\alpha_U\in
\operatorname{Aut}_{\mathrm{gr.}}(F_{T\times_S S'}(U))\) be the
element defined by \(\alpha\). Then \(u(S')(\alpha)\) is the element
\(\beta\) of
\[
 \underline{\operatorname{Aut}}_{S\text{-gr.}}(F_1)(S')
 =
 \underline{\operatorname{Aut}}_{S'\text{-gr.}}
 \left(\prod_{T\times_S S'/S'}F_{T\times_S S'}\right)
\]
such that \(\beta_{S''}=\alpha_{T\times_S S''}\) for every
\(S''\to S'\).}

Now let \(G\) be an \(S\)-functor in groups, and suppose an action of
\(G_T\) on \(F\) is given by
\[
 G_T\longrightarrow
 \underline{\operatorname{Aut}}_{T\text{-gr.}}(F).
\]
By the definition of \(\prod_{T/S}\), this gives a morphism of
\(S\)-functors in groups
\[
 G\longrightarrow
 \prod_{T/S}\underline{\operatorname{Aut}}_{T\text{-gr.}}(F).
\]
Composing with \(u\) therefore gives an action of \(G\) on
\(F_1=\prod_{T/S}F\).\footnote{Editor's note: explicitly, for every
\(S'\to S\), the action of \(G(S')\) on
\(F_1(S')=F(T\times_S S')\) is given by the action of
\(G_T(T\times_S S')=G(T\times_S S')\) on \(F(T\times_S S')\),
together with the group morphism
\(G(S')\to G(T\times_S S')\) corresponding to the projection
\(T\times_S S'\to S'\). Equivalently, the action of \(G_T\) on
\(F\) is a morphism \(G_T\times_T F\to F\). Applying
\(\prod_{T/S}\), and writing \(G_1=\prod_{T/S}G_T\), gives, by
\SGARefLink{sga3:ii:remark:1-2}{II~1.2}, a morphism \(G_1\times_S F_1\to F_1\), hence an action of
\(G_1\) on \(F_1\). The action of \(G\) is obtained through the
natural morphism \(G\to G_1\), which corresponds by adjunction to
\(\operatorname{id}_G\).}

% Source page: Exp3 component-local p. 24; combined-reader p. 134; printed p. 124.
\medskip
\noindent\SGAInlineRefTarget{sga3:iii:lemma:1-1-2}\textbf{Lemma 1.1.2.}
\textit{Under the preceding assumptions, there is a canonical
isomorphism}
\[
 H^\ast\!\left(G,\prod_{T/S}F\right)
 \simeq H^\ast(G_T,F).
\]
Indeed, by the definition of cohomology, the standard complexes are
canonically isomorphic.

\subsection*{\SGAInlineRefTarget{sga3:iii:subsection:1-2}1.2. Lifting group morphisms.\protect\footnotemark}
\footnotetext{Editor's note: this heading has been added.}
Following the general principles, we make the following definition.

\medskip
\noindent\SGAInlineRefTarget{sga3:iii:definition:1-2-1}\textbf{Definition 1.2.1.}
\textit{Let}
\[
 1\longrightarrow M\xrightarrow{u}E\xrightarrow{v}G
\]
\textit{be a sequence of morphisms of \(\widehat{\mathscr C}\)-groups.
It is called \emph{exact} if the following equivalent conditions hold:}
\begin{enumerate}
 \item[(i)] \textit{for every \(S\in\operatorname{Ob}\mathscr C\),
 the following sequence of ordinary groups is exact:}
 \[
  1\longrightarrow M(S)\xrightarrow{u(S)}E(S)
   \xrightarrow{v(S)}G(S);
 \]

 \item[\SGAInlineRefTarget{sga3:iii:definition:1-2-1:item:ii}(ii)] \textit{for every object \(H\) of
 \(\widehat{\mathscr C}\), the following sequence of ordinary groups
 is exact:}
 \[
  1\longrightarrow\operatorname{Hom}(H,M)
   \xrightarrow{u(H)}\operatorname{Hom}(H,E)
   \xrightarrow{v(H)}\operatorname{Hom}(H,G).
 \]
\end{enumerate}

Taking \(H=G\) in \SGARefLink{sga3:iii:definition:1-2-1:item:ii}{(ii)}, one sees that the set of sections of \(v\)
(which need not preserve the group structures) is either empty or
principal homogeneous under \(\operatorname{Hom}(G,M)\). Suppose that
it is nonempty, and let
\[
 s:G\longrightarrow E
\]
be a section of \(v\). For every \(S\in\operatorname{Ob}\mathscr C\)
and every \(x\in G(S)\), the element \(s(x)\in E(S)\) defines an
inner automorphism of \(E_S\) that normalizes \(M_S\) (more precisely,
the image of \(M_S\) under \(u_S\)), and hence an automorphism of
\(M_S\).

% Source page: Exp3 component-local p. 25; combined-reader p. 135; printed p. 125.
\medskip
\noindent\SGAInlineRefTarget{sga3:iii:scholium:1-2-1-1}\textbf{Scholium 1.2.1.1.}
\footnote{Editor's note: the numbering \SGARefLink{sga3:iii:scholium:1-2-1-1}{1.2.1.1} and \SGARefLink{sga3:iii:definition:1-2-1-2}{1.2.1.2} has been
added to highlight the notions introduced here.}
\textit{If \(M\) is commutative, one sees set-theoretically that this
automorphism is independent of the chosen section and depends only on
\(x\), and that it depends multiplicatively on \(x\). In summary, every
exact sequence}
\[
 (E)\qquad
 1\longrightarrow M\xrightarrow{u}E\xrightarrow{v}G
\]
\textit{in which \(M\) is commutative and \(v\) admits a section
determines a morphism of \(\widehat{\mathscr C}\)-groups}
\[
 G\longrightarrow
 \underline{\operatorname{Aut}}_{\widehat{\mathscr C}\text{-gr.}}(M),
\]
\textit{called the \emph{action of \(G\) on \(M\) defined by the
extension \((E)\)}.}

\medskip
\noindent\SGAInlineRefTarget{sga3:iii:definition:1-2-1-2}\textbf{Definition 1.2.1.2.}
We saw in \SGARefLink{sga3:i:proposition:2-3-7}{I~2.3.7} that \(v\) admits a section that is a morphism of
\(\widehat{\mathscr C}\)-groups if and only if the extension \((E)\)
is isomorphic, as an extension, to the semidirect product of \(M\) by
\(G\) for the preceding action. Such a section of \(v\) is called a
\emph{section of the extension \((E)\)}, or simply a \emph{section of
\((E)\)}.

If \(s\) is a section of \((E)\) and
\[
 m\in\Gamma(M)\simeq
 \operatorname{Ker}\bigl(\Gamma(E)\to\Gamma(G)\bigr)
\]
(for the definition of \(\Gamma\), see \SGARefLink{sga3:i:subsection:1-2}{I~1.2}), then the morphism
\(G\to E\) defined by
\[
 x\longmapsto u(m)s(x)u(m)^{-1}
\]
is also a section of \((E)\).\footnote{Editor's note: in the
right-hand side, \(x\) has been corrected to \(s(x)\).} It is called
the section deduced from \(s\) by the inner automorphism defined by
\(m\), or by \(u(m)\).

\medskip
\noindent\SGAInlineRefTarget{sga3:iii:lemma:1-2-2}\textbf{Lemma 1.2.2.}
\textit{Let}
\[
 (E)\qquad
 1\longrightarrow M\xrightarrow{u}E\xrightarrow{v}G
\]
\textit{be an exact sequence of \(\widehat{\mathscr C}\)-groups in
which \(M\) is commutative and \(v\) admits a section. Let \(G\) act
on \(M\) in the manner defined by \((E)\).}
\begin{enumerate}
 \item[\SGAInlineRefTarget{sga3:iii:lemma:1-2-2:item:i}(i)] \textit{The extension \((E)\) canonically defines a class
 \(c(E)\in H^2(G,M)\), whose vanishing is necessary and sufficient for
 the existence of a section of \((E)\).}

 \item[\SGAInlineRefTarget{sga3:iii:lemma:1-2-2:item:ii}(ii)] \textit{If \(c(E)=0\), the set of sections of \((E)\) is
 principal homogeneous under \(Z^1(G,M)\), and the set of sections of
 \((E)\), modulo the action of the inner automorphisms defined by the
 elements of \(\Gamma(M)\), is principal homogeneous under
 \(H^1(G,M)\).}

 \item[\SGAInlineRefTarget{sga3:iii:lemma:1-2-2:item:iii}(iii)]\footnote{Editor's note: this item has been added as the
 counterpart of \SGARefLink{sga3:iii:proposition:1-2-4:item:iii}{1.2.4(iii)}. The assertion is valid for every section
 of \(v\); see the proof.}
 \textit{Let \(s\) be a section of \((E)\). The set of conjugates of
 \(s\) by the inner automorphisms defined by \(\Gamma(M)\) is in
 bijection with \(\Gamma(M)/H^0(G,M)\).}
\end{enumerate}

The proof is exactly the same as for ordinary groups; starting from a
section of \(v\) ensures the functoriality of the set-theoretic
computations. We briefly indicate the principal steps.

% Source page: Exp3 component-local p. 26; combined-reader p. 136; printed p. 126.
\smallskip
\noindent\textbf{(a)}
To every section \(s\) of \(v\), associate the morphism
\[
 D_s:G\times G\longrightarrow M
\]
defined set-theoretically by
\[
 u(D_s(x,y))=s(xy)s(y)^{-1}s(x)^{-1}.
\]
The following computation shows that \(D_s\) is a
2-cocycle.\footnote{Editor's note: the original has been corrected to
make it compatible with \SGARefLink{sga3:i:subsection:5-1}{I~5.1}.} By the definition of the differential
in the standard complex (\SGARefLink{sga3:i:subsection:5-1}{I~5.1}),
\[
\begin{aligned}
 (\partial^2D_s)(x,y,z)
  ={}&\bigl(s(x)D_s(y,z)s(x)^{-1}\bigr)
     \cdot D_s(x,y)^{-1}\cdot D_s(xy,z)^{-1}\cdot D_s(x,yz).
\end{aligned}
\]
Substituting the definition of \(D_s\) into this formula gives,
without using commutativity,
\[
 D_s\in Z^2(G,M).
\]

\smallskip
\noindent\textbf{(b)}
If \(s\) and \(s'\) are two sections of \(v\), there is a morphism
\(f:G\to M\) such that \(s(x)=f(x)s'(x)\). Then
\[
\begin{aligned}
 D_{s'}(x,y)
 &=f^{-1}(xy)D_s(x,y)s(x)f(y)s(x)^{-1}f(x)\\
 &=D_s(x,y)\cdot f^{-1}(xy)
   \cdot\bigl(s(x)f(y)s(x)^{-1}\bigr)\cdot f(x),
\end{aligned}
\]
and therefore
\[
 D_{s'}=D_s\cdot\partial^1f.
\]
This shows that the class of \(D_s\) in \(H^2(G,M)\) is independent
of the chosen section \(s\) of \(v\); it is the class \(c(E)\) of the
extension \((E)\).\footnote{Editor's note: the preceding sentence has
been added.}

\smallskip
\noindent\textbf{\SGAInlineRefTarget{sga3:iii:lemma:1-2-2:item:c}(c)}
Let \(s\) and \(s'\) be two sections of \(v\), and let
\(m\in\Gamma(M)\). The equality
\[
 s(x)=m^{-1}s'(x)m
 \qquad(x\in G(S),\ S\in\operatorname{Ob}\mathscr C)
\]
is equivalent to
\[
 s(x)=m^{-1}s'(x)m\,s'(x)^{-1}s'(x),
 \qquad\text{that is,}\qquad
 s=\partial^0m\cdot s'.
\]
In particular, the stabilizer of \(s\) in \(\Gamma(M)\) is the
subgroup of elements \(m\) such that
\(\partial^0m=e_M\), namely \(H^0(G,M)\). This already proves \SGARefLink{sga3:iii:lemma:1-2-2:item:iii}{(iii)}.

\smallskip
\noindent\textbf{(d)}
The argument is now familiar.\footnote{Editor's note: what follows has
been expanded.} Let \(s_0\) be an arbitrary section of \(v\). There
exists a section \(s\), necessarily of the form \(s=f\cdot s_0\), that
is a group morphism, that is, satisfies \(D_s=0\), if and only if
\[
 (D_{s_0})^{-1}=\partial^1f,
\]
for some \(f:G\to M\), or equivalently if and only if \(c(E)=0\). This
proves \SGARefLink{sga3:iii:lemma:1-2-2:item:i}{(i)}.

In that case, the sections of \((E)\) are precisely the sections
\(s'=h\cdot s\), where \(h:G\to M\) satisfies
\(\partial^1h=0\), that is, \(h\in Z^1(G,M)\). Moreover, by \SGARefLink{sga3:iii:lemma:1-2-2:item:c}{(c)}, two
such sections \(h_1\cdot s\) and \(h_2\cdot s\) are conjugate under
\(\Gamma(M)\) if and only if \(h_1\) and \(h_2\) have the same image
in \(H^1(G,M)\). This proves \SGARefLink{sga3:iii:lemma:1-2-2:item:ii}{(ii)}.

Still let
\[
 (E)\qquad
 1\longrightarrow M\xrightarrow{u}E\xrightarrow{v}G
\]
be an exact sequence of \(\widehat{\mathscr C}\)-groups with \(M\)
commutative, and let
\[
 f:H\longrightarrow G
\]
be a morphism of \(\widehat{\mathscr C}\)-groups.

% Source page: Exp3 component-local p. 27; combined-reader p. 137; printed p. 127.
Consider \(E_f=H\times_G E\). It is a
\(\widehat{\mathscr C}\)-group, and the projection
\(v_f:E_f\to H\) is a morphism of
\(\widehat{\mathscr C}\)-groups; the same holds for
\(e_f:E_f\to E\). Sending \(M\) into \(E\) by \(u\) and into \(H\)
by the unit morphism also defines a morphism of
\(\widehat{\mathscr C}\)-groups \(u_f:M\to E_f\). We have therefore
constructed the following commutative diagram of
\(\widehat{\mathscr C}\)-groups.

\SGAThreeNativeDiagram[0.82\textwidth]
  {native_diagrams/exp3/diagram_028_exp3_p027_pullback_extension_diagram.tex}
  {The pullback extension \((E_f)\) and its morphism to \((E)\).}
  {fig:exp3-pullback-extension}
  {Polo--Gille Exp3-14oct24.pdf, component-local p. 27; combined-reader p. 137; printed p. 127}

It follows immediately that:

\medskip
\noindent\SGAInlineRefTarget{sga3:iii:lemma:1-2-3}\textbf{Lemma 1.2.3.}
\begin{enumerate}
 \item[(i)] \textit{The sequence \((E_f)\) is exact.}

 \item[(ii)] \textit{The map \(s\mapsto e_f\circ s=f'\) gives a
 bijective correspondence between the morphisms}
 \[
  s:H\longrightarrow E_f
 \]
 \textit{such that \(v_f\circ s=\operatorname{id}\), that is, the
 sections of \(v_f\), and the morphisms}
 \[
  f':H\longrightarrow E
 \]
 \textit{such that \(v\circ f'=f\), that is, the morphisms \(f'\)
 lifting \(f\).}

 \item[(iii)] \textit{Under this correspondence, the sections of
 \((E_f)\) correspond to the group morphisms \(f'\) lifting \(f\).}
\end{enumerate}

Applying \SGARefLink{sga3:iii:lemma:1-2-2}{Lemma~1.2.2} to the extension \((E_f)\), and taking \SGARefLink{sga3:iii:lemma:1-2-3}{1.2.3}
into account, gives the following proposition, which formally contains
\SGARefLink{sga3:iii:lemma:1-2-2}{Lemma~1.2.2}.

\medskip
\noindent\SGAInlineRefTarget{sga3:iii:proposition:1-2-4}\textbf{Proposition 1.2.4.}
\textit{Let}
\[
 (E)\qquad 1\longrightarrow M\longrightarrow E\xrightarrow{v}G
\]
\textit{be an exact sequence of \(\widehat{\mathscr C}\)-groups with
\(M\) commutative. Let \(f:H\to G\) be a morphism of
\(\widehat{\mathscr C}\)-groups, and suppose that it lifts to a
morphism, not necessarily a group morphism, \(f':H\to E\). Let \(H\)
act on \(M\) by the composite morphism, which is multiplicative and
independent of the choice of \(f'\),}
\[
 H\xrightarrow{f'}E\xrightarrow{\operatorname{int}}
 \underline{\operatorname{Aut}}_{\widehat{\mathscr C}\text{-gr.}}(M).
\]
\begin{enumerate}
 \item[(i)] \textit{The morphism \(f\) canonically defines a class
 \(c(f)\in H^2(H,M)\), whose vanishing is necessary and sufficient for
 the existence of a morphism of \(\widehat{\mathscr C}\)-groups}
 \[
  f':H\longrightarrow E
 \]
 \textit{lifting \(f\).}

 \item[(ii)] \textit{If \(c(f)=0\), the set of morphisms of
 \(\widehat{\mathscr C}\)-groups \(f'\) lifting \(f\), modulo the
 action of the inner automorphisms defined by elements of \(\Gamma(M)\)
 (that is, by elements \(m\in\Gamma(E)\) such that \(v(m)=e\)), is
 principal homogeneous under \(H^1(H,M)\).}

 % Source page: Exp3 component-local p. 28 (top only); combined-reader
 % p. 138 (top only); printed p. 128.
 \item[\SGAInlineRefTarget{sga3:iii:proposition:1-2-4:item:iii}(iii)] \textit{If \(f':H\to E\) is a group morphism lifting
 \(f\), the set of transforms of \(f'\) by the inner automorphisms
 defined by the elements of \(\Gamma(M)\) is isomorphic to}
 \[
  \Gamma(M)/\Gamma(M^H)=\Gamma(M)/H^0(H,M).
 \]
\end{enumerate}

% Continue with Section 1.3 on the same source page.

% END INCLUDED SGA3_English_current_buildable_source_20260823/source/inputs/III/tex/components/21_expose_III_s11_s12_extensions_cohomology.tex

% BEGIN INCLUDED SGA3_English_current_buildable_source_20260823/source/inputs/III/tex/components/22_expose_III_s13_group_law_extensions.tex
% SGA 3, Expose III, Loop 1 English reconstruction.
% Scope: Section 1.3, from its heading on Exp3 component-local p. 28
% through the end of the section on p. 31, stopping immediately before
% Section 2 (combined-reader pp. 138--141; printed pp. 128--131).
% Authority: Polo--Gille Exp3-14oct24.pdf. Jacob Reinhold's CC BY 4.0
% Markdown lineage at commit e7a259f3f8608ad3edf9bf6eead3fd504dd2d23e
% was used only as a comparison. The PDF controls the prose, mathematics,
% numbering, editorial notes, and the two temporary diagram images.

\setcounter{footnote}{45}

% Source page: Exp3 component-local p. 28; combined-reader p. 138; printed p. 128.
\subsection*{\SGAInlineRefTarget{sga3:iii:subsection:1-3}1.3. Extensions of group laws.}
Consider the following situation: one has a morphism in
\(\widehat{\mathscr C}\),
\[
 (\dagger)\qquad p:X\longrightarrow Y,
\]
and a commutative \(\widehat{\mathscr C}\)-group \(M\) acting on
\(X\), such that \(X\) is formally principal homogeneous over \(Y\)
under \(M_Y\).

If \(g:Y\to Z\) is an arbitrary morphism in
\(\widehat{\mathscr C}\), then \(g\circ p:X\to Z\) is invariant
under \(M\): for every \(S\in\operatorname{Ob}\mathscr C\), the map
\((g\circ p)(S)\) is invariant under the action of \(M(S)\) on
\(X(S)\). Conversely, we shall assume that the following condition
holds for \(n=1,2,3,4\):

\medskip
\noindent\((+)_{n}\): \textit{Every morphism from \(X^n\) to \(M\)
that is invariant under the action of \(M^n\) on \(X^n\) factors
uniquely through \(p^n:X^n\to Y^n\), where the superscript \(n\)
denotes a Cartesian power.}

\medskip
\noindent\SGAInlineRefTarget{sga3:iii:lemma:1-3-1}\textbf{Lemma 1.3.1.}
\begin{enumerate}
 \item[(i)] \textit{If \(h\) is a morphism from \(Y\) to \(M\),
 the automorphism \(u_h\) of \(X\), defined set-theoretically by
 \(x\mapsto h(p(x))\cdot x\), preserves the fibers of \(p\) and
 commutes with the operations of \(M\) on \(X\).\footnote{Editor's
 note: here and below, ``commutes with \(p\) and the operations of
 \(M\)'' has been replaced by ``preserves the fibers of \(p\) and
 commutes with the operations of \(M\).'' The equalities that follow
 have also been added.} Thus, for every
 \(S\in\operatorname{Ob}\mathscr C\), \(x\in X(S)\), and
 \(m\in M(S)\), one has}
 \[
  p\bigl(h(p(x))\cdot x\bigr)=p(x),
  \qquad
  m\cdot h(p(x))\cdot x
   =h(p(m\cdot x))\cdot m\cdot x.
 \]

 \item[(ii)] \textit{This construction gives a bijective
 correspondence between morphisms from \(Y\) to \(M\) and
 automorphisms of \(X\) that preserve the fibers of \(p\) and commute
 with the operations of \(M\).}
\end{enumerate}

The first assertion is clear, since \(p(m\cdot x)=p(x)\) and \(M\)
is commutative. Conversely, an automorphism \(u\) of \(X\) preserving
the fibers of \(p\) can be written set-theoretically as
\(x\mapsto g(x)x\), where \(g\) is a certain morphism from \(X\) to
\(M\). If \(u\) commutes with the operations of \(M\), then \(g\) is
invariant under \(M\),\footnote{Editor's note: indeed,
\(mg(x)x=mu(x)=u(mx)=g(mx)mx\), and hence \(g(mx)=g(x)\).} and the
conclusion follows from condition \((+)_{1}\).

Suppose now, in addition, that a group law on \(Y\) and an action of
\(Y\) on \(M\) are given, that is, a morphism of
\(\widehat{\mathscr C}\)-groups
\[
 (\ddagger)\qquad
 f:Y\longrightarrow
 \underline{\operatorname{Aut}}_{\widehat{\mathscr C}\text{-gr.}}(M).
\]

\medskip
\noindent\SGAInlineRefTarget{sga3:iii:definition:1-3-2}\textbf{Definition 1.3.2.}
\textit{A composition law on \(X\),}
\[
 P:X\times X\longrightarrow X,
\]
\textit{is called \emph{admissible} if it satisfies the following two
conditions:}

% Source page: Exp3 component-local p. 29; combined-reader p. 139; printed p. 129.
\begin{enumerate}
 \item[\SGAInlineRefTarget{sga3:iii:definition:1-3-2:item:i}(i)] \textit{\(P\) lifts the group law of \(Y\); that is, the
 following diagram is commutative.}

 \SGAThreeNativeDiagram[0.28\textwidth]
  {native_diagrams/exp3/diagram_029_exp3_p029_admissible_law_square.tex}
  {The compatibility of an admissible composition law with the group
   law on \(Y\).}
  {fig:exp3-p029-admissible-law-square}
  {Polo--Gille Exp3-14oct24.pdf, component-local p. 29; combined-reader
   p. 139; printed p. 129}

 \item[\SGAInlineRefTarget{sga3:iii:definition:1-3-2:item:ii}(ii)] \textit{For every \(S\in\operatorname{Ob}\mathscr C\)
 and all \(x,y\in X(S)\), \(m,n\in M(S)\), one has}
 \[
  (++)\qquad
  P(m\cdot x,n\cdot y)
   =m\cdot f(p(x))(n)\cdot P(x,y).
 \]
\end{enumerate}

\medskip
\noindent\SGAInlineRefTarget{sga3:iii:proposition:1-3-3}\textbf{Proposition 1.3.3.}
\textit{For a group law \(\ast\) on \(X\) to be admissible, it is
necessary and sufficient that the following four conditions hold:}
\begin{enumerate}
 \item[(i)] \textit{\(p:X\to Y\) is a group morphism.}

 \item[(ii)] \textit{The morphism \(i:M\to X\), defined by
 \(i(m)=m\cdot e_X\), is a group isomorphism from \(M\) onto
 \(\ker(p)\); in other words, set-theoretically,}
 \[
  (m\cdot e_X)\ast(n\cdot e_X)=(mn)\cdot e_X.
 \]

 \item[(iii)] \textit{One has
 \(m\cdot x=(m\cdot e_X)\ast x=i(m)\ast x\) for every
 \(m\in M(S)\) and \(x\in X(S)\).}

 \item[(iv)] \textit{The inner automorphisms of \(X\) act on
 \(\ker(p)\) according to the set-theoretic formula}
 \[
  x\ast i(m)\ast x^{-1}=i\bigl(f(p(x))m\bigr).
 \]
\end{enumerate}
The proof is immediate.

\medskip
\noindent\SGAInlineRefTarget{sga3:iii:lemma:1-3-4}\textbf{Lemma 1.3.4.}\footnote{Editor's note: the statement
and its proof have been expanded with a view to step~\SGARefLink{sga3:iii:paragraph:2-0:equation:e}{(e)} in the proof
of \SGARefLink{sga3:iii:proposition:1-3-6}{1.3.6}.}
\textit{Let \(h:Y\to M\) be a morphism and let \(u_h\) be the
automorphism \(x\mapsto h(p(x))\cdot x\) of \(X\) (cf.~\SGARefLink{sga3:iii:lemma:1-3-1}{1.3.1}). Let
\(P\) be an admissible composition law (respectively, an admissible
group law) on \(X\), and let \(P'\) be the composition law on \(X\)
obtained from \(P\) by means of \(u_h\), that is,}
\[
 P'(x,y)=u_h^{-1}\bigl(P(u_h(x),u_h(y))\bigr).
\]
\textit{Then:}
\begin{enumerate}
 \item[\SGAInlineRefTarget{sga3:iii:lemma:1-3-4:item:i}(i)] \textit{\(P'\) is an admissible composition law
 (respectively, an admissible group law).}

 \item[\SGAInlineRefTarget{sga3:iii:lemma:1-3-4:item:ii}(ii)] \textit{For every \(x,y\in X(S)\), with
 \(S\in\operatorname{Ob}\mathscr C\), set \(v=p(x)\) and
 \(w=p(y)\). Then}
 \[
  P'(x,y)=h(vw)^{-1}\cdot h(v)\cdot f(v)(h(w))\cdot P(x,y)
   =(\partial^1h)(p(x),p(y))\cdot P(x,y).
 \]
\end{enumerate}

\noindent\textit{Proof.}
One has \(u_h^{-1}=u_{h^\vee}\), where \(h^\vee:Y\to M\) is
defined by \(h^\vee(y)=h(y)^{-1}\). By \SGARefLink{sga3:iii:definition:1-3-2:item:i}{1.3.2(i)} and \SGARefLink{sga3:iii:lemma:1-3-4:item:ii}{(ii)},
\[
 P(h(v)\cdot x,h(w)\cdot y)
  =h(v)\cdot f(v)(h(w))\cdot P(x,y),
 \qquad p(P(x,y))=vw,
\]
and hence
\[
 P'(x,y)=h(vw)^{-1}\cdot h(v)\cdot f(v)(h(w))\cdot P(x,y)
  =(\partial^1h)(p(x),p(y))\cdot P(x,y).
\]
It follows immediately that \(P'\) satisfies conditions \SGARefLink{sga3:iii:lemma:1-3-4:item:i}{(i)} and \SGARefLink{sga3:iii:lemma:1-3-4:item:ii}{(ii)}
of \SGARefLink{sga3:iii:definition:1-3-2}{1.3.2}.

\medskip
\noindent\SGAInlineRefTarget{sga3:iii:definition:1-3-5}\textbf{Definition 1.3.5.}
\textit{Two admissible composition laws obtained from one another by
the procedure of \SGARefLink{sga3:iii:lemma:1-3-4}{1.3.4} are called \emph{equivalent}.}\footnote{Editor's
note: this is indeed an equivalence relation, since
\(u_h^{-1}=u_{h^\vee}\).}

% Source page: Exp3 component-local p. 30; combined-reader p. 140; printed p. 130.
\medskip
\noindent\SGAInlineRefTarget{sga3:iii:proposition:1-3-6}\textbf{Proposition 1.3.6.}
\textit{Suppose that an admissible composition law on \(X\) exists.
Then:}
\begin{enumerate}
 \item[\SGAInlineRefTarget{sga3:iii:proposition:1-3-6:item:i}(i)] \textit{There exists a canonically determined class
 \(c\in H^3(Y,M)\), whose vanishing is necessary and sufficient for
 the existence of an admissible associative composition law on \(X\).}

 \item[\SGAInlineRefTarget{sga3:iii:proposition:1-3-6:item:ii}(ii)] \textit{If \(c=0\), the set of admissible associative
 composition laws on \(X\) (respectively, the set of their equivalence
 classes) is principal homogeneous under \(Z^2(Y,M)\) (respectively,
 \(H^2(Y,M)\)).}
\end{enumerate}

The proof proceeds in several steps.

\smallskip
\noindent\textbf{(a)}
Let \(P\) be an admissible composition law on \(X\). Since \(P\)
lifts the associative composition law of \(Y\), there is a unique
morphism \(a:X^3\to M\) such that
\[
 (*)\qquad
 P(x,P(y,z))=a(x,y,z)\cdot P(P(x,y),z).
\]
Applying conditions \SGARefLink{sga3:iii:definition:1-3-2:item:i}{1.3.2(i)} and \SGARefLink{sga3:iii:proposition:1-3-6:item:ii}{(ii)}, one sees at once that \(a\) is
invariant under the action of \(M^3\) on \(X^3\).\footnote{Editor's
note: for \(x,y,z\in X(S)\), set \(u=p(x)\), \(v=p(y)\), and
\(w=p(z)\). Then
\[
\begin{aligned}
 P(mx,P(m'y,m''z))
  &=P(mx,m'f(v)(m'')P(y,z))\\
  &=mf(u)(m')f(uv)(m'')P(x,P(y,z)).
\end{aligned}
\]
On the other hand, since \(p(P(x,y))=uv\), one also has
\[
\begin{aligned}
 P(P(mx,m'y),m''z)
  &=P(mf(u)(m')P(x,y),m''z)\\
  &=mf(u)(m')f(uv)(m'')P(P(x,y),z).
\end{aligned}
\]
Comparison of these equalities with \((*)\) gives
\(a(mx,m'y,m''z)=a(x,y,z)\).}
Condition \((+)_{3}\) therefore gives:

\smallskip
\noindent\textit{(1) There exists a unique morphism
\(DP:Y^3\to M\) such that}
\[
 P(x,P(y,z))=DP(p(x),p(y),p(z))\cdot P(P(x,y),z),
\]
\textit{and \(P\) is associative if and only if \(DP=0\).}

\smallskip
\noindent\textbf{(b)}
Let us compute \(P(P(P(x,y),z),t)\) step by step using the preceding
formula. Put
\[
 p(x)=u,\qquad p(y)=v,\qquad p(z)=w,\qquad p(t)=h.
\]
One obtains\footnote{Editor's note: the diagram that follows has been
added.} the following pentagonal diagram, in which an arrow
\(a\xrightarrow{m}b\) means that \(b=m\cdot a\).

\SGAThreeNativeDiagram[0.82\textwidth]
  {native_diagrams/exp3/diagram_030_exp3_p030_associativity_pentagon.tex}
 {The pentagonal comparison used to compute the associativity
  obstruction.}
 {fig:exp3-p030-associativity-pentagon}
 {Polo--Gille Exp3-14oct24.pdf, component-local p. 30; combined-reader
  p. 140; printed p. 130}

Thus one finds
\[
\begin{aligned}
 &DP(u,v,w)\cdot DP(u,vw,h)\cdot f(u)(DP(v,w,h))\\
 &\qquad\cdot DP(u,v,wh)^{-1}\cdot DP(uv,w,h)^{-1}=e_M,
\end{aligned}
\]
that is, \(\partial^3DP(u,v,w,h)=e_M\). Moreover, the left-hand side
of the preceding formula can be expressed, by means of \(P\) and
\(a\), as the value at \((x,y,z,t)\) of a certain morphism
\(X^4\to M\).

% Source page: Exp3 component-local p. 31; combined-reader p. 141; printed p. 131.
It follows from the uniqueness assertion in \((+)_{4}\) that
\(\partial^3DP\) and \(e_M\), which factor the same morphism, are
equal. Hence:

\smallskip
\noindent\textit{(2) \(DP\) is a cocycle; that is,
\(DP\in Z^3(Y,M)\).}

\smallskip
\noindent\textbf{(c)}
If \(P\) and \(P'\) are two admissible composition laws on \(X\),
there is a unique morphism
\[
 b:X^2\longrightarrow M
\]
such that \(P'(x,y)=b(x,y)P(x,y)\). Applying \SGARefLink{sga3:iii:definition:1-3-2:item:ii}{1.3.2(ii)} to \(P\) and
\(P'\), one sees that \(b\) is invariant under \(M^2\), and hence,
by \((+)_{2}\):

\smallskip
\noindent\textit{\SGARefLink{sga3:iii:section:3}{(3)} For every pair \((P,P')\) of admissible
composition laws there exists a unique morphism
\(d(P,P'):Y^2\to M\) such that}
\[
 P'(x,y)=d(P,P')(p(x),p(y))\cdot P(x,y),
\]
\textit{and the set of admissible composition laws is thereby made
principal homogeneous under}
\[
 \operatorname{Hom}(Y^2,M)=C^2(Y,M).
\]

\smallskip
\noindent\textbf{(d)}
Under the preceding conditions one has
\[
 (4)\qquad DP'-DP=\partial^2d(P,P').
\]

\smallskip
\noindent\textbf{\SGARefLink{sga3:iii:paragraph:2-0:equation:e}{(e)}}
The laws \(P\) and \(P'\) are equivalent if and only if there is a
morphism
\[
 h\in C^1(Y,M)=\operatorname{Hom}(Y,M)
\]
such that \(d(P,P')=\partial^1h\). This follows from the definition of
equivalence and from \SGARefLink{sga3:iii:lemma:1-3-4:item:ii}{1.3.4(ii)}.

\smallskip
\noindent\textbf{(f)}
It remains only to conclude. We seek an associative \(P'\), that is,
one for which \(DP'=e_M\). Now \(DP\) is a cocycle whose class in
\(H^3(Y,M)\) does not depend on the chosen admissible composition law
\(P\), by \SGARefLink{sga3:iii:section:3}{(3)} and \SGARefLink{sga3:iii:section:4}{(4)}. This class is the required obstruction \(c\).
An appropriate \(P'\) can be chosen if and only if \(c=0\). Indeed,
after choosing an arbitrary \(P\), one must solve, by \SGARefLink{sga3:iii:section:1}{(1)},
\[
 0=DP'=DP+\partial^2d(P,P'),
\]
which is possible, by \SGARefLink{sga3:iii:section:3}{(3)} and \SGARefLink{sga3:iii:section:4}{(4)}, if and only if \(c=0\). Again by
\SGARefLink{sga3:iii:section:3}{(3)} and \SGARefLink{sga3:iii:section:4}{(4)}, the set of associative laws \(P'\) is principal
homogeneous under \(Z^2(Y,M)\). By \SGARefLink{sga3:iii:paragraph:2-0:equation:e}{(e)}, the set of associative laws
\(P'\) modulo equivalence is principal homogeneous under
\(H^2(Y,M)\).

% Continue with Section 2 on Exp3 component-local p. 31; combined-reader
% p. 141; printed p. 131.

% END INCLUDED SGA3_English_current_buildable_source_20260823/source/inputs/III/tex/components/22_expose_III_s13_group_law_extensions.tex

% BEGIN INCLUDED SGA3_English_current_buildable_source_20260823/source/inputs/III/tex/components/23_expose_III_s2_infinitesimal_group_morphisms.tex
% SGA 3, Expose III, Loop 1 English reconstruction.
% Scope: section 2, from the heading through Corollary 2.14, stopping
% immediately before section 3.
% Authority: Polo--Gille Exp3-14oct24.pdf, component-local pp. 31--37
% (combined-reader pp. 141--147; printed pp. 131--137), with only the
% section-2 portions of the first and last pages admitted. Force-OCR is a
% locator only. English wording was compared with Jacob C. Reinhold's CC BY
% 4.0 Markdown lineage at commit e7a259f3f8608ad3edf9bf6eead3fd504dd2d23e.
% The Polo--Gille PDF controls prose, mathematics, numbering, and notes.

\setcounter{footnote}{51}

% Source page: Exp3 local p. 31; combined p. 141; printed p. 131.
\section*{\SGAInlineRefTarget{sga3:iii:section:2}2. Infinitesimal extensions of a morphism of group schemes}
\addcontentsline{toc}{section}{2. Infinitesimal extensions of a morphism of group schemes}

Retain the notation of \S0. Let \(Y\) and \(X\) be two \(S\)-group
functors, and let \(M\) be the kernel of the group homomorphism
\(p_X:X\to X^+\). Thus there is an exact sequence of \(S\)-group
functors
\[
 1\longrightarrow M\longrightarrow X\xrightarrow{p_X}X^+.
\]
By the definition of \(X^+\), there are isomorphisms
\begin{align*}
 \operatorname{Hom}_S(Y,X^+)
  &\xrightarrow{\sim}\operatorname{Hom}_{S_{\mathscr J}}(Y_{\mathscr J},X_{\mathscr J}),\\
 \operatorname{Hom}_{S\text{-}\mathrm{gr}}(Y,X^+)
  &\xrightarrow{\sim}
    \operatorname{Hom}_{S_{\mathscr J}\text{-}\mathrm{gr}}
    (Y_{\mathscr J},X_{\mathscr J}).
\end{align*}
The map
\[
 \operatorname{Hom}_S(Y,p_X):
 \operatorname{Hom}_S(Y,X)\longrightarrow\operatorname{Hom}_S(Y,X^+)
\]
sends an \(S\)-morphism \(f:Y\to X\) to the \(S\)-morphism
\(f^+:Y\to X^+\) corresponding under these isomorphisms to the
\(S_{\mathscr J}\)-morphism
\(f_{\mathscr J}:Y_{\mathscr J}\to X_{\mathscr J}\) obtained from \(f\)
by base change. If \(M\) is commutative, \SGARefLink{sga3:iii:proposition:1-2-4}{Proposition 1.2.4} applies to
this situation.

\medskip
\noindent\SGAInlineRefTarget{sga3:iii:paragraph:2-0}\textbf{2.0.}\footnote{Editor's note: The numbering \SGARefLink{sga3:iii:paragraph:2-0}{2.0} has
been added for later references.}
In what follows we shall consider the case in which \(Y\) is flat over
\(S\), and \(X\) is an \(S\)-group functor of one of the following two
types:
\begin{enumerate}
\item[\(a\)] \(X\) is an \(S\)-group scheme;
\item[\(b\)]\SGAInlineRefTarget{sga3:iii:paragraph:2-0:item:b} \(X=\underline{\operatorname{Aut}}_S(E)\), where \(E\) is an
  \(S\)-scheme of finite presentation over \(S\).
\end{enumerate}
Write \((\mathrm{Flat})/S\) for the category of \(S\)-schemes flat over
\(S\). In case \(a\), respectively \(b\), the \(S\)-group functor
\(M=\operatorname{Ker}(X\to X^+)\), its restriction \(L\) to
\((\mathrm{Flat})/S\), and the actions on \(M\) of the inner
automorphisms of \(X\) were calculated in \SGARefLink{sga3:iii:corollary:0-9}{0.9}, \SGARefLink{sga3:iii:remark:0-5}{0.5}, and \SGARefLink{sga3:iii:corollary:0-10}{0.10},
respectively in \SGARefLink{sga3:iii:proposition:0-11}{0.11}, \SGARefLink{sga3:iii:remark:0-14}{0.14}, and \SGARefLink{sga3:iii:corollary:0-12}{0.12}.

More explicitly, in case \(a\), let \(L_0\) be the \(S_0\)-functor in
commutative groups defined, for every \(S_0\)-scheme \(T_0\), by
\[
 \operatorname{Hom}_{S_0}(T_0,L_0)
 =\operatorname{Hom}_{\mathscr O_{T_0}}
 \left(
  \omega^1_{X_0/S_0}\otimes_{\mathscr O_{S_0}}\mathscr O_{T_0},
  \mathscr J\otimes_{\mathscr O_{S_0}}\mathscr O_{T_0}
 \right).
\]
The group \(X_0\) acts on this functor through its adjoint
representation on \(\omega^1_{X_0/S_0}\), and
\[
 L=\prod_{S_0/S}L_0;
 \qquad L(T)=L_0(T\times_S S_0)
\]
for every \(S\)-scheme \(T\).

In case \(b\), let \(\pi:E\to S\) be the structural morphism. Then \(L\)
is the abelian-group functor on \((\mathrm{Flat})/S\) defined by
\[
 \operatorname{Hom}_S(T,L)
 =\Gamma\!\left(
  T,
  \pi_*\operatorname{Hom}_{\mathscr O_E}
   (\Omega^1_{E/S},\mathscr J\mathscr O_E)
  \otimes_{\mathscr O_S}\mathscr O_T
 \right),
\]
on which \(X\), regarded as a functor on \((\mathrm{Flat})/S\), acts as
described in \SGARefLink{sga3:iii:corollary:0-12}{0.12}. In either case there is an exact sequence of group
functors on \((\mathrm{Flat})/S\):
\[
 \tag{E}\label{sga3:iii:paragraph:2-0:equation:e}
 1\longrightarrow L\longrightarrow X\longrightarrow X^+.
\]

Since \(Y\) is flat over \(S\), \SGARefLink{sga3:iii:lemma:1-1-1}{Lemma 1.1.1} shows that the groups
\(H^i(Y,M)\) depend only on the restriction \(L\) of \(M\) to
\((\mathrm{Flat})/S\). In case \(a\), where
\(L=\prod_{S_0/S}L_0\), \SGARefLink{sga3:iii:lemma:1-1-2}{Lemma 1.1.2} gives isomorphisms
\[
 H^i(Y,L)\simeq H^i(Y_0,L_0).
\]

% Source pages: Exp3 local pp. 32--33; combined pp. 142--143;
% printed pp. 132--133.
Taking all this into account, \SGARefLink{sga3:iii:proposition:1-2-4}{Proposition 1.2.4} yields the following
result.\footnote{Editor's note: The definitions that appeared in the
original statement of \SGARefLink{sga3:iii:theorem:2-1}{Theorem 2.1} and occupied most of original page
113 have been moved above.}

\medskip
\noindent\SGAInlineRefTarget{sga3:iii:theorem:2-1}\textbf{Theorem 2.1.}
\textit{Let \(S\) be a scheme, and let \(\mathscr I\supset\mathscr J\)
be two quasi-coherent ideals satisfying
\(\mathscr I\mathscr J=0\), defining the closed subschemes \(S_0\) and
\(S_{\mathscr J}\). Let}
\begin{itemize}
\item[] \textit{\(X\) be an \(S\)-group functor of type \(a\) or \(b\),
  with \(L_0\) and \(L\) as above;}
\item[] \textit{\(Y\) be an \(S\)-group scheme flat over \(S\), and
  \(f_{\mathscr J}:Y_{\mathscr J}\to X_{\mathscr J}\) a morphism of
  \(S_{\mathscr J}\)-groups.}
\end{itemize}
\textit{Then:}
\begin{enumerate}
\item[\(i\)] \textit{The morphism \(f_{\mathscr J}\) lifts to a
  morphism of \(S\)-groups \(Y\to X\) if and only if the following two
  conditions hold:}
  \begin{enumerate}
  \item[\(i_1\)] \textit{\(f_{\mathscr J}\) lifts to a morphism of
    \(S\)-functors \(Y\to X\). By \SGARefLink{sga3:iii:proposition:1-2-4}{Proposition 1.2.4} this defines an
    action of \(Y\) on \(L\), independent of the chosen lift. In case
    \(a\), the resulting action of \(Y_0\) on \(L_0\) comes from
    \(f_0:Y_0\to X_0\) and the adjoint action of \(X_0\) on \(L_0\).}
  \item[\(i_2\)] \textit{The canonically defined obstruction
    \(c(f_{\mathscr J})\in H^2(Y,L)\) vanishes. In case \(a\), this group
    is isomorphic to \(H^2(Y_0,L_0)\).}
  \end{enumerate}
\item[\SGAInlineRefTarget{sga3:iii:theorem:2-1:item:ii}(ii)] \textit{If the conditions in \(i\) hold, the set
  \(\mathcal E\) of morphisms of \(S\)-group functors \(Y\to X\)
  extending \(f_{\mathscr J}\) is principal homogeneous under
  \(Z^1(Y,L)\), isomorphic in case \(a\) to \(Z^1(Y_0,L_0)\). Moreover,
  the quotient of \(\mathcal E\) by the action of inner automorphisms
  of \(X\) defined by sections over \(S\) that induce the unit section
  of \(X_{\mathscr J}\) over \(S_{\mathscr J}\) is principal
  homogeneous under \(H^1(Y,L)\), isomorphic in case \(a\) to
  \(H^1(Y_0,L_0)\).}
\item[\SGAInlineRefTarget{sga3:iii:theorem:2-1:item:iii}(iii)] \textit{If \(f:Y\to X\) is a morphism of \(S\)-group
  functors extending \(f_{\mathscr J}\), the set of morphisms \(Y\to X\)
  obtained from \(f\) by those inner automorphisms is isomorphic to}
  \[
   \Gamma(L)/H^0(Y,L),
  \]
  \textit{and in case \(a\) to
  \(\Gamma(L_0)/H^0(Y_0,L_0)\).}
\end{enumerate}

\medskip
\noindent\SGAInlineRefTarget{sga3:iii:remark:2-1-1}\textbf{Remark 2.1.1.}\footnote{Editor's note: This remark,
analogous to \SGARefLink{sga3:iii:remark:4-5-1}{4.5.1}, has been added to introduce the notation
\(d(f,f')\) and \(\bar d(f,f')\), used in \SGARefLink{sga3:iii:paragraph:4-38}{4.38}. Consequently, the
part of \SGARefLink{sga3:iii:theorem:2-1:item:ii}{2.1(ii)} concerning \(\mathcal E\) itself was also added.}
If \(f,f':Y\to X\) are morphisms of \(S\)-group functors extending
\(f_{\mathscr J}\), one obtains a cocycle
\[
 d(f,f')\in Z^1(Y,L)
 \quad\bigl(\simeq Z^1(Y_0,L_0)\text{ in case (a)}\bigr)
\]
such that
\begin{equation}
 f'=d(f,f')\cdot f.
 \tag{*}\label{sga3:iii:remark:2-1-1:equation:star}
\end{equation}
\footnote{Editor's note: We follow the sign conventions of the
original. One could instead write \(f'=d(f',f)\cdot f\), but then, in
order to retain the equality
\(d_1(d(i,i'))=d(Y,i'(Y))\) in \SGARefLink{sga3:iii:lemma:4-27}{4.27} for two immersions
\(i,i':Y\hookrightarrow X\), the sign of the class \(d(Y,Y')\)
introduced in \SGARefLink{sga3:iii:remark:4-5-1}{4.5.1} would have to be changed. This would in turn alter
the signs in \SGARefLink{sga3:iii:proposition:4-8}{4.8}, \SGARefLink{sga3:iii:corollary:4-14}{4.14}, and \SGARefLink{sga3:iii:lemma:4-17}{4.17}. We have preferred to retain all the
signs of the original, which are correct.}
Write \(\bar d(f,f')\) for the image of \(d(f,f')\) in
\(H^1(Y,L)\), or equivalently in \(H^1(Y_0,L_0)\) in case \(a\).

% Source page: Exp3 local p. 33; combined p. 143; printed p. 133.
\medskip
\noindent\SGAInlineRefTarget{sga3:iii:remark:2-2}\textbf{Remark 2.2.}
Keep the preceding notation; in particular, \(Y\) is flat over \(S\).
In case \(b\), formula (\SGARefLink{sga3:iii:remark:0-14:equation:0-14-1}{0.14.1}) identifies \(L\) with the restriction to
\((\mathrm{Flat})/S\) of
\[
 \mathbf W\!\left(
  \pi_*\operatorname{Hom}_{\mathscr O_E}
  (\Omega^1_{E/S},\mathscr J\mathscr O_E)
 \right).
\]
\SGAInlineRefTarget{sga3:iii:remark:2-2:item:a}In case \(a\), suppose in addition that \(X\) is locally of finite
presentation over \(S\). Formula (\SGARefLink{sga3:iii:remark:0-6:equation:0-6-1}{0.6.1}) then identifies \(L\) with the
restriction to \((\mathrm{Flat})/S\) of
\[
 \prod_{S_0/S}\mathbf W\!\left(
  \operatorname{Hom}_{\mathscr O_{S_0}}
   (\omega^1_{X_0/S_0},\mathscr J)
 \right).
\]
In both cases the module to which \(\mathbf W\) is applied is
quasi-coherent, by EGA I, 9.1.1. Suppose further that \(Y\) is affine
over \(S\).\footnote{Editor's note: This hypothesis and the reference
to I, \SGARefLink{sga3:i:subsection:5-3}{5.3} have been added.} By I, \SGARefLink{sga3:i:subsection:5-3}{5.3}, one obtains
\begin{align*}
 H^i(Y,L)=H^i(Y_0,L_0)
  &=H^i\!\left(Y_0,
    \operatorname{Hom}_{\mathscr O_{S_0}}
    (\omega^1_{X_0/S_0},\mathscr J)\right)
  &&\text{in case (a)},\\
 H^i(Y,L)
  &=H^i\!\left(Y,
    \pi_*\operatorname{Hom}_{\mathscr O_E}
    (\Omega^1_{E/S},\mathscr J\mathscr O_E)\right)
  &&\text{in case (b)}.
\end{align*}

% Source page: Exp3 local p. 34; combined p. 144; printed p. 134.
\medskip
\noindent\SGAInlineRefTarget{sga3:iii:remark:2-3}\textbf{Remark 2.3.}
\begin{enumerate}
\item[(1)] By \SGARefLink{sga3:iii:proposition:0-16}{0.16} and \SGARefLink{sga3:iii:corollary:0-17}{0.17}, condition \(i_1\) is automatic when
  \(Y\) is affine and
  \[
   \begin{cases}
    X\text{ is smooth over }S, &\text{in case (a)},\\
    E\text{ is smooth and affine over }S, &\text{in case (b)}.
   \end{cases}
  \]
\item[(2)] Under these hypotheses, with \(Y\) still flat over \(S\)
  as in \SGARefLink{sga3:iii:paragraph:2-0}{2.0}, \SGARefLink{sga3:iii:remark:2-2:item:a}{formulas 2.2(a)} and (\SGARefLink{sga3:iii:remark:0-6:equation:0-6-2}{0.6.2}) give in case \(a\)
  \[
   H^i(Y,L)=H^i(Y_0,L_0)
   =H^i\!\left(Y_0,
    \operatorname{Lie}(X_0/S_0)
    \otimes_{\mathscr O_{S_0}}\mathscr J\right),
  \]
  and in case \(b\), by (\SGARefLink{sga3:iii:remark:0-14:equation:0-14-2}{0.14.2}), \SGARefLink{sga3:iii:lemma:1-1-2}{1.1.2}, and I, \SGARefLink{sga3:i:subsection:5-3}{5.3},
  \[
   H^i(Y,L)=H^i\!\left(Y_0,
    \pi_{0*}\operatorname{Hom}_{\mathscr O_{E_0}}
    (\Omega^1_{E_0/S_0},
     \mathscr J\otimes_{\mathscr O_{S_0}}\mathscr O_{E_0})\right).
  \]
  \footnote{Editor's note: The preceding material has been added; see
  Editor's Note 31 in \SGARefLink{sga3:iii:remark:0-14}{0.14}.}
\end{enumerate}

We now state several corollaries for a diagonalizable \(S\)-group
\(Y\) (I, \SGARefLink{sga3:i:subsection:4-4}{4.4}). If \(S\) is affine, then, by I, \SGARefLink{sga3:i:theorem:5-3-3}{5.3.3},
\(H^n(Y,F)=0\) for every \(n>0\) and every quasi-coherent
\(\mathscr O_S\)-module \(F\).

\medskip
\noindent\SGAInlineRefTarget{sga3:iii:corollary:2-4}\textbf{Corollary 2.4.}
\textit{Let \(S_0\hookrightarrow S\) be a closed subscheme defined by
a nilpotent ideal, and let \(Y\) be a diagonalizable \(S\)-group. Let}
\begin{enumerate}
\item[\(a\)] \textit{\(X\) be an \(S\)-group locally of finite
  presentation over \(S\), or}
\item[\(b\)]\SGAInlineRefTarget{sga3:iii:corollary:2-4:item:b} \textit{\(X=\underline{\operatorname{Aut}}_S(E)\), where
  \(E\) is an \(S\)-scheme locally of finite presentation.}
\end{enumerate}
\textit{If \(f:Y\to X\) is a morphism of \(S\)-groups whose base
change \(f_0:Y_0\to X_0\) is the unit morphism, then \(f\) is the unit
morphism.}

Indeed, the question is local on \(S\), and in case \(b\) also local on
\(E\). We may therefore suppose that \(S\) is affine and, in case \(b\),
that \(E\) is of finite presentation over \(S\). Introduce the closed
subschemes \(S_n\subset S\) defined by the powers of the ideal defining
\(S_0\). This reduces the assertion to the case in which \(S_0\) is
defined by a square-zero ideal, where it follows from \SGARefLink{sga3:iii:theorem:2-1}{Theorem 2.1} and
\SGARefLink{sga3:iii:remark:2-2}{Remark 2.2}.

If \(f_0\) is not required to be the unit morphism, one has the
following statement.

\medskip
\noindent\SGAInlineRefTarget{sga3:iii:corollary:2-5}\textbf{Corollary 2.5.}
\textit{Let \(S\) and \(S_0\) be as in \SGARefLink{sga3:iii:corollary:2-4}{2.4}, with \(S\) affine. Let
\(Y\) be a diagonalizable \(S\)-group, \(X\) an \(S\)-group functor,
and \(f_0:Y_0\to X_0\) a morphism of \(S_0\)-group functors.}
\begin{enumerate}
\item[\(i\)]\footnote{Editor's note: The original stated: ``Suppose one
of the following two properties holds: \(a\) \(X\) is an \(S\)-group
locally of finite presentation; \(b\)
\(X=\underline{\operatorname{Aut}}_S(E)\), where \(E\) is of finite
presentation over \(S\). Then \(f_0\) extends to a morphism of
\(S\)-groups \(Y\to X\) if and only if it extends to a morphism of
\(S\)-functors \(Y\to X\).'' It added that this followed step by step
from part \SGARefLink{sga3:iii:corollary:2-5:item:ii}{(ii)} of the theorem. That argument appears insufficient. For
example, if \(\mathscr I^3=0\) and
\(\mathscr J=\mathscr I^2\), and if \(f:Y\to X\) is an \(S\)-functor
morphism lifting \(f_0\), then \(f_0\) lifts to an
\(S_{\mathscr J}\)-group morphism
\(g_{\mathscr J}:Y_{\mathscr J}\to X_{\mathscr J}\); but must
\(g_{\mathscr J}\) then lift to an \(S\)-functor morphism \(g:Y\to X\)?
In any event, assertion \(i\) is not used later, where \(X\) is always
assumed smooth over \(S\).}
\item[\SGAInlineRefTarget{sga3:iii:corollary:2-5:item:ii}(ii)] \textit{Suppose either that \(X\) is a smooth \(S\)-group,
or that \(X=\underline{\operatorname{Aut}}_S(E)\) with \(E\) smooth
and affine over \(S\). Then \(f_0\) extends to a morphism of \(S\)-groups
\(Y\to X\), and any two extensions are conjugate by an inner
automorphism of \(X\) defined by a section over \(S\) inducing the unit
section of \(X_0\) over \(S_0\).}
\end{enumerate}

% Source page: Exp3 local p. 35; combined p. 145; printed p. 135.
Introduce \(S_n\) as above.\footnote{Editor's note: What follows has
been slightly modified and expanded from the original.} For \SGARefLink{sga3:iii:corollary:2-5:item:ii}{(ii)}, a
scheme smooth over \(S\) is necessarily locally of finite presentation.
Thus, in case \(b\), a smooth affine \(E/S\) is of finite presentation,
so the hypothesis in \SGARefLink{sga3:iii:paragraph:2-0:item:b}{2.0(b)} is indeed satisfied.

Condition \(i_1\) of \SGARefLink{sga3:iii:theorem:2-1}{Theorem 2.1} is automatic by \SGARefLink{sga3:iii:proposition:0-16}{0.16} and \SGARefLink{sga3:iii:corollary:0-17}{0.17}.
Moreover, every section of \(X_{S_n}\) over \(S_n\) lifts to a section
of \(X_{S_{n+1}}\) over \(S_{n+1}\): this follows from smoothness in
case \(a\), and from \SGARefLink{sga3:iii:corollary:0-17}{0.17} in case \(b\). Consequently, if \(f\) and \(f'\)
are two lifts of \(f_0\), one may arrange step by step that
\(f_n=f'_n\) by lifting the inner automorphism supplied by
\SGARefLink{sga3:iii:theorem:2-1:item:ii}{Theorem 2.1(ii)}. This proves the assertion.

The same argument, together with \SGARefLink{sga3:iii:remark:2-3}{Remark 2.3}, gives:

\medskip
\noindent\SGAInlineRefTarget{sga3:iii:corollary:2-6}\textbf{Corollary 2.6.}
\textit{Let \(\mathscr I\) be a nilpotent ideal of \(S\) defining
\(S_0\), let \(Y\) be an affine \(S\)-group flat over \(S\), and let
\(X\) be a smooth \(S\)-group.}
\begin{enumerate}
\item[\(i\)] \textit{If, for every \(n\geq0\),}
 \[
  H^2\!\left(Y_0,
   \operatorname{Lie}(X_0/S_0)\otimes_{\mathscr O_{S_0}}
   \mathscr I^{n+1}/\mathscr I^{n+2}\right)=0,
 \]
 \textit{then every \(S_0\)-group morphism \(f_0:Y_0\to X_0\) lifts
 to an \(S\)-group morphism \(f:Y\to X\).}
\item[(ii)] \textit{If, for every \(n\geq0\),}
 \[
  H^1\!\left(Y_0,
   \operatorname{Lie}(X_0/S_0)\otimes_{\mathscr O_{S_0}}
   \mathscr I^{n+1}/\mathscr I^{n+2}\right)=0,
 \]
 \textit{then any two such lifts are conjugate by an inner
 automorphism of \(X\) defined by a section over \(S\) inducing the
 unit section of \(X_0\) over \(S_0\).}
\end{enumerate}

\medskip
\noindent\SGAInlineRefTarget{sga3:iii:lemma:2-7}\textbf{Lemma 2.7.}
\textit{Let \(S\) be affine, \(G\) an affine \(S\)-group,
\(\mathscr F\) a quasi-coherent \(\mathscr O_S\)-module, and
\(\mathscr L\) a locally free \(\mathscr O_S\)-module. Suppose \(G\)
acts on \(\mathscr F\) in the sense of Expos\'e I, hence acts on
\(\mathscr F\otimes_{\mathscr O_S}\mathscr L\). Let \(\Lambda\) be
the ring of \(S\), and let \(L\) be the projective \(\Lambda\)-module
defining \(\mathscr L\). Then there is a canonical isomorphism}
\[
 H^*(G,\mathscr F\otimes_{\mathscr O_S}\mathscr L)
 \simeq H^*(G,\mathscr F)\otimes_\Lambda L.
\]
\footnote{Editor's note: Here \(\mathscr L\) is given the trivial
action of \(G\).}

Let \(\mathscr A\) be the \(\mathscr O_S\)-algebra
\(\mathscr A(G)\), and consider the complex \(\mathscr C\) of quasi-coherent
\(\mathscr O_S\)-modules
\[
 0\longrightarrow\mathscr F
 \longrightarrow\mathscr F\otimes\mathscr A
 \longrightarrow\mathscr F\otimes\mathscr A\otimes\mathscr A
 \longrightarrow\cdots.
\]
\footnote{Editor's note: The original proof has been expanded and
 simplified. It also invoked the isomorphisms
 \(H^n(\mathscr C\otimes\mathscr L)\simeq
 H^n(\mathscr C)\otimes\mathscr L\) and
\(H_n(\Gamma(-))\simeq\Gamma(H_n(-))\), where
\(H^n(\mathscr C)\) denotes the cohomology sheaf of the complex
\(\mathscr C\).}

% Source page: Exp3 local p. 36; combined p. 146; printed p. 136.
By I, \SGARefLink{sga3:i:subsection:5-3}{5.3}, \(H^*(G,\mathscr F)\), respectively
\(H^*(G,\mathscr F\otimes\mathscr L)\), is the cohomology of
\(\Gamma(S,\mathscr C)\), respectively
\(\Gamma(S,\mathscr C\otimes\mathscr L)\). Since \(S\) is affine,
EGA I, 1.3.12 gives
\[
 \Gamma(S,\mathscr C\otimes\mathscr L)
 \simeq\Gamma(S,\mathscr C)\otimes_\Lambda L.
\]
Because \(L\) is projective, hence flat,
\[
 H^*(\Gamma(S,\mathscr C)\otimes_\Lambda L)
 \simeq H^*(\Gamma(S,\mathscr C))\otimes_\Lambda L,
\]
which proves the lemma.

Using the lemma, \SGARefLink{sga3:iii:corollary:2-6}{Corollary 2.6} becomes:

\medskip
\noindent\SGAInlineRefTarget{sga3:iii:corollary:2-8}\textbf{Corollary 2.8.}
\textit{Let \(S\) be affine and let \(\mathscr I\) be a nilpotent
ideal defining \(S_0\). Suppose every
\(\mathscr I^{n+1}/\mathscr I^{n+2}\) is locally free on \(S_0\). Let
\(Y\) be an affine \(S\)-group flat over \(S\), let \(X\) be a smooth
\(S\)-group, and let \(f_0:Y_0\to X_0\) be an \(S_0\)-group morphism.}
\begin{enumerate}
\item[\(i\)]\SGAInlineRefTarget{sga3:iii:corollary:2-8:item:i} \textit{If
\(H^2(Y_0,\operatorname{Lie}(X_0/S_0))=0\), then \(f_0\) lifts to an
\(S\)-group morphism \(Y\to X\).}
\item[\SGAInlineRefTarget{sga3:iii:corollary:2-8:item:ii}(ii)] \textit{If
\(H^1(Y_0,\operatorname{Lie}(X_0/S_0))=0\), then any two such lifts are
conjugate by an inner automorphism of \(X\) defined by a section over
\(S\) inducing the unit section of \(X_0\) over \(S_0\).}
\end{enumerate}

Taking \(Y=X\) gives:

\medskip
\noindent\SGAInlineRefTarget{sga3:iii:corollary:2-9}\textbf{Corollary 2.9.}
\textit{Let \(S\) and \(S_0\) be as above, and let \(X\) be a smooth
affine \(S\)-group.}
\begin{enumerate}
\item[\(i\)] \textit{If
\(H^1(X_0,\operatorname{Lie}(X_0/S_0))=0\), every endomorphism of \(X\)
over \(S\) that induces the identity on \(X_0\) is the inner
automorphism defined by a section over \(S\) inducing the unit section
of \(X_0\) over \(S_0\).}
\item[(ii)] \textit{If
\(H^2(X_0,\operatorname{Lie}(X_0/S_0))=0\), every
\(S_0\)-automorphism of \(X_0\) extends to an \(S\)-automorphism of
\(X\).}\footnote{Editor's note: Let \(f_0\) be an
\(S_0\)-automorphism of \(X_0\), with inverse \(g_0\). By \SGARefLink{sga3:iii:corollary:2-8:item:i}{2.8(i)}, lift
\(f_0\) and \(g_0\) to \(S\)-endomorphisms \(f\) and \(g\) of \(X\).
Then \(g\circ f\) and \(f\circ g\) induce the identity on \(X_0\), so
by \SGARefLink{sga3:iii:corollary:0-7}{0.7} they are \(S\)-automorphisms. Hence \(f\) and \(g\) are
\(S\)-automorphisms as well.}
\end{enumerate}

\medskip
\noindent\SGAInlineRefTarget{sga3:iii:remark:2-10}\textbf{Remark 2.10.}
The assertions involving \(H^1\) have converses by \SGARefLink{sga3:iii:theorem:2-1}{Theorem 2.1}. For
example, let \(S=I_{S_0}\) be the dual-number scheme over \(S_0\)
(II, \SGARefLink{sga3:iii:theorem:2-1}{2.1}), and let \(X\) be a flat \(S\)-group such that every
automorphism over \(S\) inducing the identity over \(S_0\) is an inner
automorphism defined by a section over \(S\) inducing the unit section
of \(X_0\). Then
\[
 H^1(X_0,\operatorname{Lie}(X_0/S_0))=0.
\]
\footnote{Editor's note: This is used in \SGARefLink{sga3:xxiv:1.13}{Expos\'e XXIV, 1.13}.}

\medskip
\noindent\SGAInlineRefTarget{sga3:iii:corollary:2-11}\textbf{Corollary 2.11.}
\textit{Let \(S,\mathscr I,\mathscr J\) be as in \SGARefLink{sga3:iii:theorem:2-1}{2.1}. Let \(Y\) be an
\(S\)-group scheme flat over \(S\), let \(X\) be an \(S\)-group scheme,
and let \(f:Y\to X\) be an \(S\)-group morphism. The set of morphisms
from \(Y\) to \(X\) obtained from \(f\) by conjugation with elements
\(x\in X(S)\) inducing the unit of \(X(S_{\mathscr J})\) is isomorphic
to the quotient}
\[
 \mathcal E=
 \operatorname{Hom}_{\mathscr O_{S_0}}(\omega^1_{X_0/S_0},\mathscr J)
 \big/
 \operatorname{Hom}_{\mathscr O_{S_0}}(\omega^1_{X_0/S_0},\mathscr J)
 ^{\operatorname{ad}(Y_0)}.
\]
\textit{Here the second group consists of those
\(\mathscr O_{S_0}\)-morphisms
\(\omega^1_{X_0/S_0}\to\mathscr J\) which, after every base change
\(S'\to S_0\), yield morphisms}
\[
 \omega^1_{X_0/S_0}\otimes_{\mathscr O_{S_0}}\mathscr O_{S'}
 \longrightarrow
 \mathscr J\otimes_{\mathscr O_{S_0}}\mathscr O_{S'}
\]
\textit{invariant under the action of \(Y_0(S')\) on the first
factor.}

By \SGARefLink{sga3:iii:theorem:2-1:item:iii}{Theorem 2.1(iii)}, this set is
\(\Gamma(L_0)/H^0(Y_0,L_0)\). Now
\[
 \Gamma(L_0)=
 \operatorname{Hom}_{\mathscr O_{S_0}}(\omega^1_{X_0/S_0},\mathscr J),
\]
and \(H^0(Y_0,L_0)\) is precisely
\(\Gamma(L_0)^{\operatorname{ad}(Y_0)}\) in the sense just stated.

% Source page: Exp3 local p. 37; combined p. 147; printed p. 137.
\medskip
\noindent\SGAInlineRefTarget{sga3:iii:corollary:2-12}\textbf{Corollary 2.12.}
\textit{Under the hypotheses of \SGARefLink{sga3:iii:corollary:2-11}{2.11}, suppose in addition that
\(\omega^1_{X_0/S_0}\) is locally free of finite rank. Then}
\[
 \mathcal E\simeq
 \Gamma\!\left(S_0,
  \operatorname{Lie}(X_0/S_0)\otimes_{\mathscr O_{S_0}}\mathscr J
 \right)
 \big/
 H^0\!\left(Y_0,
  \operatorname{Lie}(X_0/S_0)\otimes_{\mathscr O_{S_0}}\mathscr J
 \right).
\]
Indeed,\footnote{Editor's note: The following sentence has been added.}
local freeness of finite rank gives
\[
 \operatorname{Hom}_{\mathscr O_{S_0}}
 (\omega^1_{X_0/S_0},\mathscr J)
 \simeq
 \operatorname{Lie}(X_0/S_0)\otimes_{\mathscr O_{S_0}}\mathscr J.
\]

\medskip
\noindent\SGAInlineRefTarget{sga3:iii:corollary:2-13}\textbf{Corollary 2.13.}
\textit{Suppose in addition that \(Y_0\) is diagonalizable. Then}
\[
 \mathcal E\simeq
 \Gamma\!\left(S_0,
  \operatorname{Lie}(X_0/S_0)\otimes_{\mathscr O_{S_0}}\mathscr J
 \right)
 \big/
 \Gamma\!\left(S_0,
  \operatorname{Lie}(X_0/S_0)^{\operatorname{ad}(Y_0)}
  \otimes_{\mathscr O_{S_0}}\mathscr J
 \right),
\]
\textit{where
\(\operatorname{Lie}(X_0/S_0)^{\operatorname{ad}(Y_0)}\) is the factor
of the decomposition in I, \SGARefLink{sga3:i:proposition:4-7-3}{4.7.3} corresponding to the zero
character of \(Y_0\).}

Indeed, if \(Y_0\simeq D_{S_0}(M)\), the cited result gives a direct-sum
decomposition
\[
 \operatorname{Lie}(X_0/S_0)
 =\operatorname{Lie}(X_0/S_0)_0
 \oplus
 \bigoplus_{\substack{m\in M\\m\ne0}}
 \operatorname{Lie}(X_0/S_0)_m.
\]
Tensoring with \(\mathscr J\) gives the analogous decomposition for
\(\operatorname{Lie}(X_0/S_0)\otimes\mathscr J\), whence
\[
 H^0\!\left(Y_0,
  \operatorname{Lie}(X_0/S_0)\otimes\mathscr J\right)
 \simeq
 \Gamma\!\left(S_0,
  \operatorname{Lie}(X_0/S_0)_0
  \otimes_{\mathscr O_{S_0}}\mathscr J\right).
\]

\medskip
\noindent\SGAInlineRefTarget{sga3:iii:corollary:2-14}\textbf{Corollary 2.14.}
\textit{Suppose in addition that \(S_0\) is affine. Then}
\[
 \mathcal E\simeq
 \Gamma\!\left(
  S_0,
  \left[
   \operatorname{Lie}(X_0/S_0)
   \big/
   \operatorname{Lie}(X_0/S_0)^{\operatorname{ad}(Y_0)}
  \right]
  \otimes_{\mathscr O_{S_0}}\mathscr J
 \right).
\]

% Continue with section 3 on Exp3 local p. 37 / combined p. 147.



% END INCLUDED SGA3_English_current_buildable_source_20260823/source/inputs/III/tex/components/23_expose_III_s2_infinitesimal_group_morphisms.tex

% BEGIN INCLUDED SGA3_English_current_buildable_source_20260823/source/inputs/III/tex/components/24_expose_III_s3_infinitesimal_group_structures.tex
% SGA 3, Expose III, Loop 1 English reconstruction.
% Scope: complete section 3, from its heading through Corollary 3.10,
% stopping immediately before section 4.
% Authority: Polo--Gille Exp3-14oct24.pdf, component-local pp. 37--42
% (combined-reader pp. 147--152; original printed pp. 121--126), with
% only the section-3 portions of the first and last pages admitted.
% Force-OCR is a locator only. English wording was compared with Jacob C.
% Reinhold's CC BY 4.0 Markdown lineage at commit
% e7a259f3f8608ad3edf9bf6eead3fd504dd2d23e.
% The Polo--Gille PDF controls prose, mathematics, numbering, and notes.

\setcounter{footnote}{64}

% Source page: Exp3 local p. 37; combined p. 147; original printed p. 121.
\section*{\SGAInlineRefTarget{sga3:iii:section:3}3. Infinitesimal extensions of a group scheme}
\addcontentsline{toc}{section}{3. Infinitesimal extensions of a group scheme}

Still using the notation of \S0 (\(S,\mathscr I,\mathscr J\), and so on),
let \(X\) be an \(S\)-scheme, and suppose that \(X_{\mathscr J}\) is
endowed with a group structure. We propose to determine the \(S\)-group
structures on \(X\) that induce the given structure on
\(X_{\mathscr J}\).

From now on, assume that \(X\) is flat over \(S\). Let \(\mathscr C\) be
the category of \(S\)-schemes flat over \(S\); thus
\(X\in\operatorname{Ob}\mathscr C\). Denote by \(Y\), respectively \(M\),
the functor on \(\mathscr C\) defined by \(X^+\), respectively \(L'_X\).
The canonical morphism \(p_X:X\to X^+\) defines a morphism in
\(\widehat{\mathscr C}\)
\[
 p:X\longrightarrow Y,
\]
and the action of \(L'_X\) on \(X\) in
\(\widehat{(\mathrm{Sch})/S}\) defines an action of \(M\) on \(X\) in
\(\widehat{\mathscr C}\). One checks at once that \(X\) is thereby
formally principal homogeneous under \(M_Y\) over \(Y\) (cf. \SGARefLink{sga3:iii:proposition:0-2:item:i}{0.2(i)}
and \SGARefLink{sga3:iii:remark:0-4}{0.4}).

% Source page: Exp3 local p. 38; combined p. 148; original printed p. 122.
The action of \(X^+\) on \(L'_X\) defined in \SGARefLink{sga3:iii:proposition:0-8}{0.8} (denoted
\(\operatorname{Ad}\) there) defines an action, denoted \(\widetilde f\),
of \(Y\) on \(M\). On the other hand, by \SGARefLink{sga3:iii:remark:0-5}{0.5},
\[
 \operatorname{Hom}_{\widehat{\mathscr C}}(Z,M)
 \simeq \operatorname{Hom}_{S_0}(Z_0,L_0),
 \qquad Z\in\operatorname{Ob}\mathscr C,
\]
where \(L_0\) is the functor defined in \SGARefLink{sga3:iii:remark:0-5}{0.5}.

\medskip
\noindent\SGAInlineRefTarget{sga3:iii:lemma:3-1}\textbf{Lemma 3.1.}
\begin{enumerate}
\item[(i)] \textit{Condition \((+)_{n}\) of \SGARefLink{sga3:iii:subsection:1-3}{1.3} holds for every positive
integer \(n\).}
\item[\SGAInlineRefTarget{sga3:iii:lemma:3-1:item:ii}(ii)] \textit{Let the \(S_0\)-group \(X_0\) act on the
\(S_0\)-functor \(L_0\) through its adjoint representation. Then there is
a canonical isomorphism}
\[
 H^*(X_0,L_0)\simeq H^*(Y,M),
\]
\textit{where the first cohomology is computed in
\((\mathrm{Sch})/S_0\), and the second in \(\mathscr C\).}
\end{enumerate}

Both parts of the lemma follow from the relation
\begin{align*}
 \operatorname{Hom}_{\widehat{\mathscr C}}(Y,M)
 &\simeq
 \operatorname{Hom}_{\widehat{(\mathrm{Sch})/S_0}}
       (X^+\times_S S_0,L_0)\\
 &\simeq \operatorname{Hom}_{S_0}(X_0,L_0)\\
 &\simeq \operatorname{Hom}_{\widehat{\mathscr C}}(X,M),
\end{align*}
which follows at once from the definition of \(M\) as
\(\prod_{S_0/S}L_0\). More generally, the same relation holds with
\(X,Y\) replaced by \(X^n,Y^n\). Consequently, every morphism
\(X^n\to M\) factors uniquely through \(Y^n\), which gives \((+)_{n}\).
It also gives the relation
\(C^*(Y,M)=C^*(X_0,L_0)\), and hence \SGARefLink{sga3:iii:lemma:3-1:item:ii}{(ii)}.

We may therefore apply the constructions of \SGARefLink{sga3:iii:subsection:1-3}{1.3}. In particular:

\medskip
\noindent\SGAInlineRefTarget{sga3:iii:lemma:3-2}\textbf{Lemma 3.2.}
\textit{Let \(P:X\times_S X\to X\) be a morphism. Then \(P\) induces
the group law of \(X_{\mathscr J}\) if and only if \(P\) is an
admissible composition law on \(X\) (cf. \SGARefLink{sga3:iii:definition:1-3-2}{1.3.2}).}

Indeed, \(P\) induces the group law of \(X_{\mathscr J}\) if and only if
it lifts the group law of \(X^+\), or equivalently that of \(Y\). Thus it
remains only to show that every morphism \(P\) lifting the group law of
\(X_{\mathscr J}\) satisfies identity ((++)) of \SGARefLink{sga3:iii:definition:1-3-2:item:ii}{1.3.2(ii)}, which is
exactly what was proved in \SGARefLink{sga3:iii:proposition:0-8}{0.8}.

\medskip
\noindent\SGAInlineRefTarget{sga3:iii:proposition:3-3}\textbf{Proposition 3.3.}
\textit{Let \(S\) be a scheme and \(S_0\) a closed subscheme defined by
a nilpotent ideal. Let \(X\) be an \(S\)-scheme that is flat and either
quasi-compact or locally of finite presentation over \(S\). Let
\(P:X\times_S X\to X\) be a composition law on \(X\). Then \(P\) is a
group law if and only if the following two conditions hold:}
\begin{enumerate}
\item[\SGAInlineRefTarget{sga3:iii:proposition:3-3:item:i}(i)] \textit{\(P\) is associative;}
\item[\SGAInlineRefTarget{sga3:iii:proposition:3-3:item:ii}(ii)] \textit{\(P\) induces a group law on
\(X_0=X\times_S S_0\).}
\end{enumerate}

These conditions are clearly necessary. We prove that they are
sufficient. Suppose first that \(X\to S\) has a section. Since
\(X(S')\) is then nonempty for every \(S'\to S\), it suffices
\footnote{Editor's note: The original has been corrected by deleting
the inappropriate reference to an exercise of Bourbaki on semigroups
(cf. [BAlg], \S I.2, Exercises 9--13) and by indicating the role of the
left and right translations; see the following editor's note.}
to show that, for every \(x\in X(S')\), left and right translation by
\(x\) are isomorphisms of \(X_{S'}\).
\footnote{Editor's note: Let \(E\) be a nonempty set endowed with an
associative composition law such that every left translation
\(\ell_x\) is bijective, and fix \(x_0\in E\). There is a unique
\(e\in E\) such that \(x_0e=x_0\); then
\(x_0ex=x_0x\) implies \(ex=x\) for every \(x\in E\). Moreover, for
every \(x\) there is a unique \(x'\) such that \(xx'=e\). Suppose in
addition that there is a \(b\in E\) for which right translation
\(r_b\) is injective. Then \(xeb=xb\) gives \(xe=x\), so \(e\) is a
unit, and \(xx'x=x=xe\) gives \(x'x=e\); hence \(x'\) is both a left
and a right inverse of \(x\), and \(E\) is a group.

The hypothesis that \(r_b\) be injective is necessary: on any set
\(E\), define \(xy=y\) for all \(x,y\in E\). Every left translation is
then the identity (and the law is associative), whereas
\(r_y(E)=\{y\}\) for every \(y\); thus \(E\) is not a group when
\(|E|>1\).}

% Source page: Exp3 local p. 39; combined p. 149; original printed pp. 123--124.
We may plainly assume \(S'=S\). The translation in question, \(t\),
induces on \(X_0\) a translation \(t_0\), which is an automorphism since
\(X_0\) is a group. The conclusion follows by flatness (SGA 1 III,
4.2).
\footnote{Editor's note: Since \(X\) and \(X_0\) have the same
underlying topological space and \(t_0\) is an automorphism, \(t\) is a
homeomorphism, hence an affine morphism; cf. Expos\'e VI\textsubscript{B},
2.9.1, or EGA IV\textsubscript{4}, 18.12.7.1. It therefore suffices to
show the following. Let \(\mathscr J\) be a nilpotent ideal of a ring
\(\Lambda\), and let \(\varphi:A\to B\) be a morphism of
\(\Lambda\)-algebras, with \(B\) flat over \(\Lambda\), such that
\(\varphi\otimes_\Lambda(\Lambda/\mathscr J)\) is bijective. Then
\(\varphi\) is bijective. By the ``nilpotent Nakayama lemma,''
\(\varphi\) is surjective; moreover, since \(B\) is flat over
\(\Lambda\), one also has
\(\operatorname{Ker}(\varphi)\otimes_\Lambda
(\Lambda/\mathscr J)=0\), whence
\(\operatorname{Ker}(\varphi)=0\). Thus \(\varphi\) is bijective.}

Now drop the assumption that \(X\) has a section over \(S\), and
suppose that there exists \(S'\to S\) such that \(X_{S'}\) has a
section over \(S'\). By what has just been proved, \(X_{S'}\) is then an
\(S'\)-group; let \(e'\) be its unit section. The inverse image of
\(e'\) under
\(\operatorname{pr}_i:S'\times_S S'\to S'\), \(i=1,2\), is the unit
section of \(X_{S''}\) for the group law obtained by pulling
\(P_{S'}\) back along \(\operatorname{pr}_i\). But since \(P\) is
``defined over \(S\),'' these two group laws coincide, and so do their
unit sections. Hence
\[
 \operatorname{pr}_1^*(e')=\operatorname{pr}_2^*(e').
\]

If \(S'\to S\) is a descent morphism (cf. Expos\'e IV, no.~2), there is
a section of \(X\) whose base extension is \(e'\), and the proof is
complete. Since \(X_X\) has a section over \(X\), namely the diagonal
section, it is now enough to prove that \(X\to S\) is a descent
morphism. It is flat and surjective, and is either quasi-compact or
locally of finite presentation; it is therefore an (fpqc) covering,
hence a descent morphism (Expos\'e IV, no.~6).

\medskip
\noindent\textbf{Remark.}
In fact, the hypothesis that \(X\to S\) be quasi-compact or locally of
finite presentation is superfluous, by the following result, which the
reader may prove as an exercise on Expos\'e IV:

Under the stated hypotheses on \(S\) and \(S_0\), if \(X\to S\) is
flat and \(X_0\to S_0\) is an (fpqc) covering, then \(X\to S\) is a
descent morphism.

\medskip
\noindent\SGAInlineRefTarget{sga3:iii:lemma:3-4}\textbf{Lemma 3.4.}
\textit{Two admissible composition laws on \(X\) are equivalent
(cf. \SGARefLink{sga3:iii:definition:1-3-5}{1.3.5}) if and only if one is obtained from the other by an
automorphism of \(X\) over \(S\) that induces the identity on
\(X_{\mathscr J}\).}

% Source page: Exp3 local p. 40; combined p. 150; original printed p. 125.
Indeed, the morphisms constructed in \SGARefLink{sga3:iii:lemma:1-3-1}{1.3.1} are precisely those in the
preceding statement, by \SGARefLink{sga3:iii:corollary:0-7}{0.7}.
\footnote{Editor's note: Indeed, by the proof of \SGARefLink{sga3:iii:corollary:0-7}{0.7}, the
\(S\)-endomorphisms of \(X\) inducing the identity on
\(X_{\mathscr J}\) are the automorphisms \(m\cdot\operatorname{id}_X\),
where \(m\) ranges over
\(M(X)=\operatorname{Hom}_S(X,M)\). For every \(S'\to S\) and every
\(x\in X(S')\), one has
\((m\cdot\operatorname{id}_X)(x)=m(x)\cdot x\). By the proof of \SGARefLink{sga3:iii:lemma:3-1}{3.1},
each \(m:X\to M\) factors uniquely through a morphism
\(h:Y=X^+\to M\); hence \(m\cdot\operatorname{id}_X\) is the
automorphism \(u_h\) introduced in \SGARefLink{sga3:iii:lemma:1-3-1}{1.3.1}. The lemma now follows from
the definition of equivalence; cf. \SGARefLink{sga3:iii:lemma:1-3-4}{1.3.4} and \SGARefLink{sga3:iii:definition:1-3-5}{1.3.5}.}

Taking account of all the preceding results, \SGARefLink{sga3:iii:proposition:1-3-6}{Proposition 1.3.6} gives:

\medskip
\noindent\SGAInlineRefTarget{sga3:iii:theorem:3-5}\textbf{Theorem 3.5.}
\textit{Let \(S\) be a scheme, and let \(\mathscr I\) and
\(\mathscr J\) be two ideals on \(S\) such that
\(\mathscr I\supset\mathscr J\) and
\(\mathscr I\mathscr J=0\). Let \(S_0\) and \(S_{\mathscr J}\) be the
closed subschemes of \(S\) that they define. Let \(X\) be an
\(S\)-scheme flat over \(S\) (and either locally of finite presentation
or quasi-compact over \(S\)), and let \(X_0\) and
\(X_{\mathscr J}\) be the schemes obtained by base change. Suppose
\(X_{\mathscr J}\) is endowed with an \(S_{\mathscr J}\)-group
structure, and let \(L_0\) denote the \(S_0\)-functor in commutative
groups defined by}
\[
 \operatorname{Hom}_{S_0}(T,L_0)
 =\operatorname{Hom}_{\mathscr O_T}
 \left(
  \omega^1_{X_0/S_0}\otimes_{\mathscr O_{S_0}}\mathscr O_T,
  \mathscr J\otimes_{\mathscr O_{S_0}}\mathscr O_T
 \right),
\]
\textit{on which \(X_0\) acts through its adjoint representation.}
\begin{enumerate}
\item[\SGAInlineRefTarget{sga3:iii:theorem:3-5:item:i}(i)] \textit{An \(S\)-group structure exists on \(X\) that
induces the given structure on \(X_{\mathscr J}\) if and only if the
following conditions hold:}
  \begin{enumerate}
  \item[\((i_1)\)] \textit{there is a morphism of \(S\)-schemes
  \(X\times_S X\to X\) inducing the group law of
  \(X_{\mathscr J}\);}
  \item[\((i_2)\)] \textit{a certain obstruction class in
  \(H^3(X_0,L_0)\), canonically defined by \(X\) and the group law of
  \(X_{\mathscr J}\), is zero.}
  \end{enumerate}
\item[\SGAInlineRefTarget{sga3:iii:theorem:3-5:item:ii}(ii)] \textit{If the conditions in \SGARefLink{sga3:iii:theorem:3-5:item:i}{(i)} hold, the set
\(\mathcal E\) of group laws on \(X\) inducing the given law on
\(X_{\mathscr J}\) is principal homogeneous under
\(Z^2(X_0,L_0)\). Moreover, the quotient of \(\mathcal E\) by the
\(S\)-automorphisms of \(X\) inducing the identity on
\(X_{\mathscr J}\) is principal homogeneous under \(H^2(X_0,L_0)\).}
\end{enumerate}

\footnote{Editor's note: The following argument has been added.}
Indeed, by \SGARefLink{sga3:iii:lemma:3-2}{3.2} every morphism of \(S\)-schemes
\(f:X\times_S X\to X\) inducing the group law of
\(X_{\mathscr J}\) is an admissible composition law on \(X\). By
\SGARefLink{sga3:iii:proposition:1-3-6:item:i}{1.3.6(i)}, the existence of an associative admissible composition law
\(P:X\times_S X\to X\) is equivalent to the vanishing of a certain
class \(c(f)\in H^3(X_0,L_0)\); in that case, \SGARefLink{sga3:iii:proposition:3-3}{3.3} shows that \(P\) is a
group law. This proves \SGARefLink{sga3:iii:proposition:3-3:item:i}{(i)}, and \SGARefLink{sga3:iii:proposition:3-3:item:ii}{(ii)} then follows from \SGARefLink{sga3:iii:proposition:3-3}{3.3} and
\SGARefLink{sga3:iii:proposition:1-3-6:item:ii}{1.3.6(ii)}.

\medskip
\noindent\SGAInlineRefTarget{sga3:iii:remark:3-5-1}\textbf{Remark 3.5.1.}
\footnote{Editor's note: This remark, analogous to \SGARefLink{sga3:iii:remark:4-5-1}{4.5.1}, has been
added to introduce the notation \(\delta(\mu,\mu')\), or
\(\delta(X,X')\), used in \SGARefLink{sga3:iii:paragraph:4-38}{4.38}. Consequently, the part of \SGARefLink{sga3:iii:theorem:3-5:item:ii}{3.5(ii)}
concerning \(\mathcal E\) itself has also been added.}
If \(\mu\) and \(\mu'\) are group laws on \(X\) inducing the given law
on \(X_{\mathscr J}\), one obtains a cocycle
\(\delta(\mu,\mu')\in Z^2(X_0,L_0)\). The chosen sign convention is
\(\mu'=\delta(\mu,\mu')\cdot\mu\); that is, for every \(S'\to S\) and
\(x,y\in X(S')\),
\[
 \mu'(x,y)=\delta(\mu,\mu')(x_0,y_0)\cdot\mu(x,y).
\]
\footnote{Editor's note: We follow the sign conventions of the
original in order to have, in \SGARefLink{sga3:iii:paragraph:4-38:item:5}{4.38(5)}, the equality
\(\partial^1\bar d(X,X')=\bar\delta(X,X')\); see also Editor's
Note 54.}

% Source page: Exp3 local p. 41; combined p. 151; original printed p. 126.
Write \(\bar\delta(\mu,\mu')\) for the image of
\(\delta(\mu,\mu')\) in \(H^2(X_0,L_0)\). Finally, if \(X\) endowed
with \(\mu\), respectively \(\mu'\), is denoted simply by \(X\),
respectively \(X'\), write \(\delta(X,X')\) in place of
\(\delta(\mu,\mu')\), and similarly for \(\bar\delta(X,X')\).

\medskip
\noindent\SGAInlineRefTarget{sga3:iii:remark:3-6}\textbf{Remark 3.6.}
Let \(X_{\mathscr J}\) be a smooth affine \(S_{\mathscr J}\)-scheme.
By \SGARefLink{sga3:iii:proposition:0-15}{0.15}, up to isomorphism there is a unique \(S\)-scheme \(X\), smooth
over \(S\), whose reduction modulo \(\mathscr J\) is
\(X_{\mathscr J}\). If \(X_{\mathscr J}\) is endowed with an
\(S_{\mathscr J}\)-group structure, condition \((i_1)\) is automatic
by \SGARefLink{sga3:iii:proposition:0-16}{0.16}. Moreover, \SGARefLink{sga3:iii:remark:0-6}{0.6} simplifies the definition of \(L_0\) and gives:

\medskip
\noindent\SGAInlineRefTarget{sga3:iii:corollary:3-7}\textbf{Corollary 3.7.}
\textit{Let \(S,\mathscr I,\mathscr J\) be as in \SGARefLink{sga3:iii:lemma:3-1}{3.1}, and let
\(X_{\mathscr J}\) be a smooth affine \(S_{\mathscr J}\)-group.}
\begin{enumerate}
\item[(i)] \textit{The set of smooth \(S\)-groups whose reduction
modulo \(\mathscr J\) is \(X_{\mathscr J}\), up to isomorphism
inducing the identity on \(X_{\mathscr J}\), is either empty or
principal homogeneous under}
\[
 H^2\!\left(
  X_0,\operatorname{Lie}(X_0/S_0)
  \otimes_{\mathscr O_{S_0}}\mathscr J
 \right).
\]
\item[(ii)] \textit{A smooth \(S\)-group whose reduction modulo
\(\mathscr J\) is \(X_{\mathscr J}\) exists if and only if a certain
obstruction in}
\[
 H^3\!\left(
  X_0,\operatorname{Lie}(X_0/S_0)
  \otimes_{\mathscr O_{S_0}}\mathscr J
 \right)
\]
\textit{is zero.}
\end{enumerate}

As usual, one obtains the following corollaries.

\medskip
\noindent\SGAInlineRefTarget{sga3:iii:corollary:3-8}\textbf{Corollary 3.8.}
\textit{Let \(S\) be a scheme and \(S_0\) a closed subscheme defined by
a nilpotent ideal \(\mathscr I\). Let \(X_0\) be a smooth affine
\(S_0\)-group.}
\begin{enumerate}
\item[(i)] \textit{If, for every \(n\geq0\),}
\[
 H^2\!\left(
  X_0,\operatorname{Lie}(X_0/S_0)
  \otimes_{\mathscr O_{S_0}}
  \mathscr I^{n+1}/\mathscr I^{n+2}
 \right)=0,
\]
\textit{then any two smooth \(S\)-groups whose reduction is \(X_0\)
are isomorphic by an isomorphism inducing the identity on \(X_0\).}
\item[(ii)] \textit{If, for every \(n\geq0\),}
\[
 H^3\!\left(
  X_0,\operatorname{Lie}(X_0/S_0)
  \otimes_{\mathscr O_{S_0}}
  \mathscr I^{n+1}/\mathscr I^{n+2}
 \right)=0,
\]
\textit{then there exists a smooth \(S\)-group whose reduction is
\(X_0\).}
\end{enumerate}

\medskip
\noindent\SGAInlineRefTarget{sga3:iii:corollary:3-9}\textbf{Corollary 3.9.}
\textit{Let \(S\) be affine and \(S_0\) a closed subscheme defined by a
nilpotent ideal \(\mathscr I\). Suppose that every
\(\mathscr I^{n+1}/\mathscr I^{n+2}\) is locally free on \(S_0\). Let
\(X_0\) be a smooth affine \(S_0\)-group.}
\begin{enumerate}
\item[(i)] \textit{If
\(H^2(X_0,\operatorname{Lie}(X_0/S_0))=0\), then any two smooth
\(S\)-groups whose reduction is \(X_0\) are isomorphic.}
\item[(ii)] \textit{If
\(H^3(X_0,\operatorname{Lie}(X_0/S_0))=0\), then there exists a smooth
\(S\)-group whose reduction is \(X_0\).}
\end{enumerate}

% Source page: Exp3 local p. 42; combined p. 152; original printed p. 126.
\medskip
\noindent\SGAInlineRefTarget{sga3:iii:corollary:3-10}\textbf{Corollary 3.10.}
\textit{Let \(S_0\) be a scheme and let \(S=I_{S_0}\) be the
dual-number scheme over \(S_0\). Let \(X_0\) be a smooth \(S_0\)-group.
Then every smooth \(S\)-group \(Y\) for which \(Y_0\) is
\(S_0\)-isomorphic to \(X_0\) is \(S\)-isomorphic to
\(X=X_0\times_{S_0}S\) if and only if}
\[
 H^2(X_0,\operatorname{Lie}(X_0/S_0))=0.
\]
\footnote{Editor's note: This result is used in XXIV, \SGARefLink{sga3:xxiv:1.13}{1.13}.}

Indeed, by \SGARefLink{sga3:iii:theorem:3-5}{3.5} the set of classes of such groups \(Y\), up to an
\(S\)-group isomorphism inducing the identity on \(X_0\), is in
bijection with \(H^2(X_0,\operatorname{Lie}(X_0/S_0))\). Hence their
classes up to an arbitrary \(S\)-group isomorphism are in bijection
with
\[
 H^2(X_0,\operatorname{Lie}(X_0/S_0))/\Gamma_0,
\]
where
\[
 \Gamma_0=\operatorname{Aut}_{S_0\text{-}\mathrm{gr.}}(X_0),
\]
which acts in the evident way on \(H^2\). The conclusion follows at
once.
\footnote{Editor's note: Indeed,
\(\operatorname{Aut}_{S_0\text{-}\mathrm{gr.}}(X_0)\) acts by group
automorphisms on the abelian group
\(H^2(X_0,\operatorname{Lie}(X_0/S_0))\), so the orbit of \(0\) is the
singleton \(\{0\}\). Consequently, the quotient set is a singleton if
and only if
\(H^2(X_0,\operatorname{Lie}(X_0/S_0))=\{0\}\).}

% Continue with section 4. Its heading begins at the bottom of Exp3 local
% p. 42 (combined p. 152; original printed p. 127).

% END INCLUDED SGA3_English_current_buildable_source_20260823/source/inputs/III/tex/components/24_expose_III_s3_infinitesimal_group_structures.tex

% BEGIN INCLUDED SGA3_English_current_buildable_source_20260823/source/inputs/III/tex/components/25_expose_III_s41_s44_closed_subgroups.tex
% SGA 3, Expose III, Loop 1 English reconstruction.
% Scope: section 4 from its heading through Complement 4.4.1, stopping
% immediately before 4.5.0.
% Authority: Polo--Gille Exp3-14oct24.pdf, component-local pp. 42--47
% (combined-reader pp. 152--157; printed pp. 127--131), with only the
% section-4 portions of the first and last pages admitted. Force-OCR is a
% locator only. English wording was compared with Jacob C. Reinhold's CC BY
% 4.0 Markdown lineage at commit e7a259f3f8608ad3edf9bf6eead3fd504dd2d23e.
% The Polo--Gille PDF controls prose, mathematics, numbering, notes, and the
% seven source-derived Loop-1 diagram images included below.

\setcounter{footnote}{73}

% Source page: Exp3 local p. 42; combined p. 152; printed p. 127.
\section*{\SGAInlineRefTarget{sga3:iii:section:4}4. Infinitesimal extensions of closed subgroups}
\addcontentsline{toc}{section}{4. Infinitesimal extensions of closed subgroups}

We first state a result valid in an arbitrary abelian category.

\medskip
\noindent\SGAInlineRefTarget{sga3:iii:lemma:4-1}\textbf{Lemma 4.1.}
\textit{Let
\[
 0\longrightarrow A'\xrightarrow{i}A\xrightarrow{p}A''\longrightarrow0
\]
be an exact sequence, let \(\phi:A'\to Q\) be a morphism, and let
\(\pi:A''\to P\) be an epimorphism with kernel \(C\). Let \(E\) be the
set, up to isomorphism, of quadruples \((B,f,g,h)\) such that the sequence
\[
 0\longrightarrow Q\xrightarrow{f}B\xrightarrow{g}P\longrightarrow0
\]
is exact and the following diagram commutes:}

% Raster stand-in for Loop 1; replace by native TeX/TikZ in Loop 2.
% Authority: Exp3-14oct24.pdf local p. 42; combined reader p. 152.
\begin{center}
  \SGAThreeNativeInput[0.78\textwidth]
    {native_diagrams/exp3/diagram_031_exp3_p042_lemma41_commutative.tex}
\end{center}

\begin{enumerate}
\item[(i)] \textit{The set \(E\) is nonempty if and only if the image in
  \(\operatorname{Ext}^1(C,Q)\) of the element \(A\) of
  \(\operatorname{Ext}^1(A'',A')\) is zero.}
\item[\SGAInlineRefTarget{sga3:iii:lemma:4-1:item:ii}(ii)] \textit{Under these conditions, \(E\) is a principal homogeneous
  set under the abelian group \(\operatorname{Hom}(C,Q)\).}
\end{enumerate}

Introduce the amalgamated sum
\[
 B'=A\mathop{\sqcup}\limits_{A'}Q.
\]
There is then a commutative diagram with exact rows:
\footnote{Editor's note: One has
\(\operatorname{Coker}(j)=B'\mathop{\sqcup}\limits_Q0
=A\mathop{\sqcup}\limits_{A'}0=A''\), and
\(\operatorname{Ker}(j)\simeq\operatorname{Ker}(i)=0\), as one sees by
reasoning ``as if the category were a category of modules.'' For a proof
entirely in terms of arrows, see for example [Fr64], Theorem 2.5.4 (*).}

% Raster stand-in for Loop 1; replace by native TeX/TikZ in Loop 2.
% Authority: Exp3-14oct24.pdf local p. 43; combined reader p. 153.
\begin{center}
  \SGAThreeNativeInput[0.78\textwidth]
    {native_diagrams/exp3/diagram_032_exp3_p043_lemma41_pushout.tex}
\end{center}

% Source page: Exp3 local p. 43; combined p. 153; printed pp. 128--129.
It is clear that the category of solutions to the problem is canonically
isomorphic to the category of solutions to the corresponding problem for
the sequence
\[
 0\longrightarrow Q\longrightarrow B'\longrightarrow A''\longrightarrow0
\]
and the morphisms \(\operatorname{id}_Q\) and \(\pi:A''\to P\).
\footnote{Editor's note: In what follows, \(A\) has been replaced by
\(B'\), and the end of the argument has been expanded.}
In this case, \(E\) is in bijection with the set of subobjects \(N\) of
\(B'\) for which \(B'\to A''\) induces an isomorphism of \(N\) with the
kernel \(C\) of \(A''\to P\); equivalently, it is the set of morphisms
\(e:C\to B'\) lifting the canonical morphism \(C\to A''\). The abelian
group \(G=\operatorname{Hom}(C,Q)\) acts on \(E\) by
\(g\cdot e=g+e\), with addition in \(\operatorname{Hom}(C,B')\); if
\(E\ne\varnothing\), this makes \(E\) a principal homogeneous set under
\(G\).

We deduce the following.

\medskip
\noindent\SGAInlineRefTarget{sga3:iii:proposition:4-2}\textbf{Proposition 4.2.}
\footnote{Editor's note: The statement has been rewritten so as to fit
exactly the application made in \SGARefLink{sga3:iii:proposition:4-3}{4.3}; the proof has also been expanded,
following indications supplied by M. Demazure.}
\textit{Let \(S\) be a scheme, let \(S_{\mathscr J}\) be the closed
subscheme defined by a quasi-coherent square-zero ideal \(\mathscr J\),
let \(X\) be an \(S\)-scheme, and let \(\mathscr F\) be an
\(\mathscr O_X\)-module. Put
\[
 X_{\mathscr J}=X\times_S S_{\mathscr J},\qquad
 \mathscr F_{\mathscr J}
 =\mathscr F\otimes_{\mathscr O_S}\mathscr O_{S_{\mathscr J}},
\]
and let
\(\mathscr G_{\mathscr J}=\mathscr F_{\mathscr J}/\mathscr H_{\mathscr J}\)
be a quotient module of \(\mathscr F_{\mathscr J}\). Suppose given a
morphism of \(\mathscr O_{X_{\mathscr J}}\)-modules
\[
 f:\mathscr J\otimes_{\mathscr O_{S_{\mathscr J}}}
       \mathscr G_{\mathscr J}\longrightarrow\mathscr Q.
\]
Let \(\mathcal E\) be the sheaf of sets on \(X\) defined as follows: for
each open subset \(U\) of \(X\), \(\mathcal E(U)\) is the set of quotient
modules \(\mathscr G\) of \(\mathscr F|_U\) such that
\(\mathscr G/\mathscr J\mathscr G=\mathscr G_{\mathscr J}|_U\) and for
which there exists an isomorphism
\[
 h:\mathscr J\mathscr G\xrightarrow{\sim}\mathscr Q|_U
\]
making the following diagram commute:}

% Raster stand-in for Loop 1; replace by native TeX/TikZ in Loop 2.
% Authority: Exp3-14oct24.pdf local p. 43; combined reader p. 153.
\begin{center}
  \SGAThreeNativeInput[0.43\textwidth]
    {native_diagrams/exp3/diagram_033_exp3_p043_prop42_triangle.tex}
\end{center}

\textit{The isomorphism \(h\) is then unique, since
\(
 \mathscr J\otimes_{\mathscr O_{S_{\mathscr J}}}
 (\mathscr G_{\mathscr J}|_U)\to\mathscr J\mathscr G
\)
is an epimorphism. Then \(\mathcal E\) is a sheaf formally principal
homogeneous under the sheaf of commutative groups
\[
 \mathscr A=
 \operatorname{Hom}_{\mathscr O_X}(\mathscr H_{\mathscr J},\mathscr Q)
 =\operatorname{Hom}_{\mathscr O_{X_{\mathscr J}}}
   (\mathscr H_{\mathscr J},\mathscr Q).
\]}

\noindent\textit{Proof.}
If \(\mathcal E(U)=\varnothing\), there is nothing to prove; we may
therefore suppose that \(\mathcal E(U)\) contains an element
\(\widetilde{\mathscr G}\). In the following diagram, \(h\) is an
isomorphism and all the arrows are epimorphisms:

% Raster stand-in for Loop 1; replace by native TeX/TikZ in Loop 2.
% Authority: Exp3-14oct24.pdf local p. 44; combined reader p. 154.
\begin{center}
  \SGAThreeNativeInput[0.70\textwidth]
    {native_diagrams/exp3/diagram_034_exp3_p044_prop42_epimorphisms.tex}
\end{center}

% Source page: Exp3 local p. 44; combined p. 154; printed p. 129.
Thus the morphism
\[
 \mathscr J\otimes_{\mathscr O_{S_{\mathscr J}}}
 (\mathscr F_{\mathscr J}|_U)\longrightarrow\mathscr Q|_U
\]
induces a necessarily unique epimorphism
\(\phi:\mathscr J\mathscr F|_U\to\mathscr Q|_U\). If \(\mathscr G\)
is an \(\mathscr O_U\)-module such that
\(\mathscr G/\mathscr J\mathscr G=\mathscr G_{\mathscr J}|_U\) and there
is a commutative diagram with exact rows

% Raster stand-in for Loop 1; replace by native TeX/TikZ in Loop 2.
% Authority: Exp3-14oct24.pdf local p. 44; combined reader p. 154.
\begin{center}
  \SGAThreeNativeInput[0.69\textwidth]
    {native_diagrams/exp3/diagram_035_exp3_p044_prop42_exact_rows.tex}
\end{center}

\noindent(where \(p_{\mathscr J}\) is the projection
\(\mathscr G\to\mathscr G/\mathscr J\mathscr G
=\mathscr G_{\mathscr J}|_U\), so that
\(\mathscr Q|_U=\operatorname{Ker}(p_{\mathscr J})
=\mathscr J\mathscr G\)), then \(\mathscr G\) may be identified with a
quotient module of \(\mathscr F|_U\). Consequently, by \SGARefLink{sga3:iii:lemma:4-1:item:ii}{4.1(ii)}, the set
\(\mathcal E(U)\) is principal homogeneous under the abelian group
\[
 \operatorname{Hom}_{\mathscr O_X}
  (\mathscr H_{\mathscr J},\mathscr Q)(U)
 =\operatorname{Hom}_{\mathscr O_{X_{\mathscr J}}}
  (\mathscr H_{\mathscr J},\mathscr Q)(U).
\]

\medskip
\noindent\SGAInlineRefTarget{sga3:iii:proposition:4-3}\textbf{Proposition 4.3.} \textit{(TDTE \SGARefLink{sga3:iv:subsection:5:1}{IV 5.1}) Let \(S\) be a
scheme, let \(S_{\mathscr J}\) be the closed subscheme defined by a
quasi-coherent square-zero ideal \(\mathscr J\), let \(X\) be an
\(S\)-scheme, and let \(\mathscr F\) be a quasi-coherent
\(\mathscr O_X\)-module. Put
\[
 X_{\mathscr J}=X\times_S S_{\mathscr J},\qquad
 \mathscr F_{\mathscr J}
 =\mathscr F\otimes_{\mathscr O_S}\mathscr O_{S_{\mathscr J}}.
\]
Let
\(\mathscr G_{\mathscr J}=\mathscr F_{\mathscr J}/\mathscr H_{\mathscr J}\)
be a quasi-coherent quotient module of \(\mathscr F_{\mathscr J}\), flat
over \(S_{\mathscr J}\). For each open subset \(U\) of \(X\), let
\(\mathcal E(U)\) be the set of quasi-coherent quotient modules
\(\mathscr G\)\footnote{Editor's note: To lighten the statement, the
hypothesis that \(\mathscr G\) be quasi-coherent has been added here; the
remark that this hypothesis is automatically satisfied has been moved to
the proof, which has been expanded accordingly.} of \(\mathscr F|_U\),
flat over \(S\), such that
\(\mathscr G/\mathscr J\mathscr G\simeq
\mathscr G_{\mathscr J}|_U\). Then the \(\mathcal E(U)\) form a sheaf of
sets \(\mathcal E\) on \(X\), formally principal homogeneous under the
sheaf of commutative groups
\[
 \mathscr A=
 \operatorname{Hom}_{\mathscr O_{X_{\mathscr J}}}
  \left(
   \mathscr H_{\mathscr J},
   \mathscr J\otimes_{\mathscr O_{S_{\mathscr J}}}
   \mathscr G_{\mathscr J}
  \right).
\]}

\noindent\textit{Proof.}
Let \(\pi:X\to S\) be the structural morphism. Let \(U\) be an open
subset of \(X\), and let \(\mathscr G\) be an \(\mathscr O_U\)-module
flat over \(S\), such that
\(\mathscr G/\mathscr J\mathscr G\simeq
\mathscr G_{\mathscr J}|_U\). For every \(x\in U\), the stalk
\(\mathscr G_x\) is a flat module over the local ring
\(\mathscr O_{S,s}\), where \(s=\pi(x)\). Hence the morphism
\[
 \mathscr J_s\otimes_{\mathscr O_{S,s}}
  (\mathscr G/\mathscr J\mathscr G)_x
 =\mathscr J_s\otimes_{\mathscr O_{S,s}}\mathscr G_x
 \longrightarrow(\mathscr J\mathscr G)_x
\]
is bijective. We therefore have an exact sequence
\[
 0\longrightarrow
 \mathscr J\otimes_{\mathscr O_S}(\mathscr G_{\mathscr J}|_U)
 \longrightarrow\mathscr G\longrightarrow
 \mathscr G_{\mathscr J}|_U\longrightarrow0.
\]

% Source page: Exp3 local p. 45; combined p. 155; printed p. 130.
Since
\(\mathscr J\otimes_{\mathscr O_S}(\mathscr G_{\mathscr J}|_U)\) and
\(\mathscr G_{\mathscr J}|_U\) are quasi-coherent
\(\mathscr O_U\)-modules, \(\mathscr G\) is quasi-coherent as well
(EGA III, 1.4.17).

Conversely, since \(\mathscr G_{\mathscr J}\) is flat over
\(S_{\mathscr J}\), if \(\mathscr G\) is a quasi-coherent
\(\mathscr O_U\)-module such that
\(\mathscr G/\mathscr J\mathscr G\simeq\mathscr G_{\mathscr J}\) and
the morphism
\[
 \mathscr J\otimes_{\mathscr O_{S_{\mathscr J}}}
 \mathscr G_{\mathscr J}\longrightarrow\mathscr J\mathscr G
\]
is bijective, then \(\mathscr G\) is flat over \(S\), by the
``fundamental criterion of flatness'' (SGA 1, IV, 5.5).
\footnote{Editor's note: See also [BAC], \S III.5, Theorem 1.}

Consequently, the set \(\mathcal E(U)\) considered here is the set
considered in \SGARefLink{sga3:iii:proposition:4-2}{4.2} when \(f\) is the identity morphism of
\(\mathscr J\otimes_{\mathscr O_{S_{\mathscr J}}}
\mathscr G_{\mathscr J}\); the conclusion therefore follows from \SGARefLink{sga3:iii:proposition:4-2}{4.2}.
\hfill\(\square\)

Retain the preceding notation.
\footnote{Editor's note: The discussion that follows has been expanded and
\SGARefLink{sga3:iii:corollary:4-3-1}{Corollary 4.3.1} has been added. Recall also that ``pseudo-torsor'' is
synonymous with ``formally principal homogeneous.''}
Let \(Y_{\mathscr J}\) be a closed subscheme of \(X_{\mathscr J}\),
defined by a quasi-coherent ideal \(\mathscr I_{Y_{\mathscr J}}\), and
suppose that \(Y_{\mathscr J}\) is flat over \(S_{\mathscr J}\). Applying
\SGARefLink{sga3:iii:proposition:4-3}{4.3} to \(\mathscr F=\mathscr O_X\) and
\[
 \mathscr G_{\mathscr J}=\mathscr O_{Y_{\mathscr J}}
 =\mathscr O_{X_{\mathscr J}}/\mathscr I_{Y_{\mathscr J}},
\]
we obtain the following corollary.

\medskip
\noindent\SGAInlineRefTarget{sga3:iii:corollary:4-3-1}\textbf{Corollary 4.3.1.}
\textit{Let \(S,S_{\mathscr J},\mathscr J,X,X_{\mathscr J},
Y_{\mathscr J}\), and \(\mathscr I_{Y_{\mathscr J}}\) be as above, and
suppose that \(Y_{\mathscr J}\) is flat over \(S_{\mathscr J}\). Let
\(\mathscr A_{\mathscr J}\) be the sheaf of commutative groups
\[
 \operatorname{Hom}_{\mathscr O_{X_{\mathscr J}}}
 \left(
  \mathscr I_{Y_{\mathscr J}},
  \mathscr J\otimes_{\mathscr O_{S_{\mathscr J}}}
  \mathscr O_{Y_{\mathscr J}}
 \right)
\]
on \(X_{\mathscr J}\), and put
\(\mathscr A=i_*(\mathscr A_{\mathscr J})\), where
\(i:X_{\mathscr J}\hookrightarrow X\) is the immersion. For each open
subset \(U\) of \(X\), let \(\mathcal E(U)\) be the set of closed
subschemes \(Y\) of \(U\), flat over \(S\), such that
\[
 Y\times_S S_{\mathscr J}=Y_{\mathscr J}\cap U.
\]
Then \(\mathcal E\) is an \(\mathscr A\)-pseudo-torsor.}

\textit{If, moreover, such a \(Y\) exists locally---that is, if every
\(x\in X\) has an open neighborhood \(U\) with
\(\mathcal E(U)\ne\varnothing\)---then \(\mathcal E\) is an
\(\mathscr A\)-torsor.} The \(\mathscr A\)-torsors on \(X\) are
parametrized by
\[
 H^1(X,\mathscr A)=H^1(X_{\mathscr J},\mathscr A_{\mathscr J})
\]
(see, for example, EGA IV\textsubscript{4}, 16.5.15), and
\(\mathcal E\) has a global section---that is,
\(\mathcal E(X)\ne\varnothing\)---if and only if its corresponding
cohomology class is zero. Thus we obtain:

\medskip
\noindent\SGAInlineRefTarget{sga3:iii:corollary:4-4}\textbf{Corollary 4.4.}
\textit{Let \(S,S_{\mathscr J},\mathscr J,X,X_{\mathscr J},
Y_{\mathscr J}\), and \(\mathscr I_{Y_{\mathscr J}}\) be as above, and
suppose that \(Y_{\mathscr J}\) is flat over \(S_{\mathscr J}\). Let
\(E\) be the set of closed subschemes \(Y\) of \(X\), flat over \(S\),
such that \(Y\times_S S_{\mathscr J}=Y_{\mathscr J}\).}
\begin{enumerate}
\item[(i)] \textit{The set \(E\) is either empty or principal homogeneous
under the abelian group}
\[
 H^0(X,\mathscr A)=H^0(X_{\mathscr J},\mathscr A_{\mathscr J})
 =\operatorname{Hom}_{\mathscr O_{X_{\mathscr J}}}
 \left(
  \mathscr I_{Y_{\mathscr J}},
  \mathscr J\otimes_{\mathscr O_{S_{\mathscr J}}}
  \mathscr O_{Y_{\mathscr J}}
 \right).
\]

% Source page: Exp3 local p. 46; combined p. 156; printed p. 131.
\item[(ii)] \textit{The set \(E\) is nonempty if and only if both of the
following conditions hold:}
\begin{enumerate}
\item[(a)] \textit{The problem has a solution locally on \(X\).}
\item[(b)] \textit{A certain obstruction vanishes; it lies in}
\[
 H^1\!\left(
  X_{\mathscr J},
  \operatorname{Hom}_{\mathscr O_{X_{\mathscr J}}}
  \left(
   \mathscr I_{Y_{\mathscr J}},
   \mathscr J\otimes_{\mathscr O_{S_{\mathscr J}}}
   \mathscr O_{Y_{\mathscr J}}
  \right)
 \right).
\]
\end{enumerate}
\end{enumerate}

\medskip
\noindent\SGAInlineRefTarget{sga3:iii:complement:4-4-1}\textbf{Complement 4.4.1.}
\footnote{Editor's note: This complement has been added; it is useful in
proving Proposition 4.8(ii).}
Retain the notation of \SGARefLink{sga3:iii:corollary:4-4}{4.4}, and suppose that \(E\) contains an element
\(Y\). Let \(\mathscr I_Y\) be the ideal of \(\mathscr O_X\) defining
\(Y\), and let \(\mathscr I_{Y_{\mathscr J}}\) be its image in
\(\mathscr O_{X_{\mathscr J}}\). As we saw in the proof of \SGARefLink{sga3:iii:proposition:4-2}{4.2}, there is
a commutative diagram

% Raster stand-in for Loop 1; replace by native TeX/TikZ in Loop 2.
% Authority: Exp3-14oct24.pdf local p. 46; combined reader p. 156.
\begin{center}
  \SGAThreeNativeInput[0.44\textwidth]
    {native_diagrams/exp3/diagram_036_exp3_p046_complement441_square.tex}
\end{center}

\noindent and hence an epimorphism of \(\mathscr O_X\)-modules
\[
 \phi:\mathscr J\mathscr O_X\longrightarrow
 \mathscr J\otimes_{\mathscr O_{S_{\mathscr J}}}
 \mathscr O_{Y_{\mathscr J}}.
\]
Let \(\mathscr K\) be its kernel. For every element \(Y'\) of \(E\),
the morphism \(\mathscr O_X\to\mathscr O_{Y'}\) then factors through
\(\mathscr O_X/\mathscr K\), which is the amalgamated sum \(B'\) in the
proof of \SGARefLink{sga3:iii:lemma:4-1}{Lemma 4.1}. If \(\mathscr I_{Y'}\) denotes the ideal of \(Y'\)
in \(\mathscr O_X\), there is a commutative diagram

% Raster stand-in for Loop 1; replace by native TeX/TikZ in Loop 2.
% Authority: Exp3-14oct24.pdf local p. 46; combined reader p. 156.
\begin{center}
  \SGAThreeNativeInput[0.80\textwidth]
    {native_diagrams/exp3/diagram_037_exp3_p046_complement441_exact_diagram.tex}
\end{center}

Replacing \(X\) by the closed subscheme defined by \(\mathscr K\), we
are thus reduced to the case \(\mathscr K=0\). The datum of \(Y'\) is
then equivalent to that of the \(\mathscr O_X\)-submodule
\(\mathscr I_{Y'}\) of \(\mathscr O_X\) which maps bijectively onto
\(\mathscr I_{Y_{\mathscr J}}\) under the projection
\(p:\mathscr O_X\to\mathscr O_{X_{\mathscr J}}\). Let
\[
 f':\mathscr I_{Y_{\mathscr J}}\xrightarrow{\sim}\mathscr I_{Y'},
 \qquad
 f:\mathscr I_{Y_{\mathscr J}}\xrightarrow{\sim}\mathscr I_Y
\]
denote the respective inverse isomorphisms. Then \(f'-f\) is an element
of
\[
 \operatorname{Hom}_{\mathscr O_{X_{\mathscr J}}}
 \left(
  \mathscr I_{Y_{\mathscr J}},
  \mathscr J\otimes_{\mathscr O_{S_{\mathscr J}}}
  \mathscr O_{Y_{\mathscr J}}
 \right)
 =\operatorname{Hom}_{\mathscr O_{X_{\mathscr J}}}
  (\mathscr I_{Y_{\mathscr J}},\mathscr J\mathscr O_Y),
\]
which we denote by \(d(Y',Y)\). Notice that
\(d(Y,Y')=-d(Y',Y)\).

% Source page: Exp3 local p. 47; combined p. 157; printed p. 131.
For our fixed \(Y\) and variable \(Y'\), consider the morphism
\[
 \mathscr I_{Y'}\longrightarrow\mathscr O_X\longrightarrow
 \mathscr O_Y=\mathscr O_X/\mathscr I_Y.
\]
Since its composite with \(\mathscr O_Y\to\mathscr O_{Y_{\mathscr J}}\)
is zero, it takes values in
\[
 \mathscr J\mathscr O_Y
 =\mathscr J\otimes_{\mathscr O_{S_{\mathscr J}}}
  \mathscr O_{Y_{\mathscr J}}.
\]
More precisely, let \(V\) be an open subset of \(X\), let \(x'\) be a
section of \(\mathscr I_{Y'}\) on \(V\), and let \(x_{\mathscr J}\) be
its image in \(\Gamma(V,\mathscr I_{Y_{\mathscr J}})\). Then
\[
 x'=f'(x_{\mathscr J})
   =f(x_{\mathscr J})+(f'-f)(x_{\mathscr J})
   =f(x_{\mathscr J})+d(Y',Y)(x_{\mathscr J}).
\]
Consequently, the morphism
\(\mathscr I_{Y'}\to\mathscr J\mathscr O_Y\) is given by
\(d(Y',Y)\).

% Continue with 4.5.0 on Exp3 local p. 47 / combined p. 157 / printed p. 131.

% END INCLUDED SGA3_English_current_buildable_source_20260823/source/inputs/III/tex/components/25_expose_III_s41_s44_closed_subgroups.tex

% BEGIN INCLUDED SGA3_English_current_buildable_source_20260823/source/inputs/III/tex/components/26_expose_III_s450_s47_lci_koszul.tex
% SGA 3, Expose III, Loop 1 English reconstruction.
% Scope: 4.5.0 through the 4.7 suppression notice, stopping immediately
% before Remark 4.8.0.
% Authority: Polo--Gille Exp3-14oct24.pdf, component-local pp. 47--51
% (combined-reader pp. 157--161; original side-pages 131--132), with only the
% 4.5.0--4.7 portions of the first and last pages admitted. Force-OCR is a
% locator only. English wording was compared with Jacob C. Reinhold's CC BY
% 4.0 Markdown lineage at commit e7a259f3f8608ad3edf9bf6eead3fd504dd2d23e.
% The Polo--Gille PDF controls prose, mathematics, numbering, notes, and the
% source-derived Loop-1 bicomplex image included below.

\setcounter{footnote}{80}

% Source page: Exp3 local p. 47; combined p. 157; original side-pages 131--132.
\medskip
\noindent\SGAInlineRefTarget{sga3:iii:paragraph:4-5-0}\textbf{4.5.0.}\footnote{Editor's note: The numbering \SGARefLink{sga3:iii:paragraph:4-5-0}{4.5.0} has
been added to mark the return to the original.}
Retain the notation of \SGARefLink{sga3:iii:corollary:4-3-1}{4.3.1} and \SGARefLink{sga3:iii:corollary:4-4}{4.4}, and make the following
transformations. The quasi-coherent \(\mathscr O_{X_{\mathscr J}}\)-module
\[
 \mathscr I_{Y_{\mathscr J}}/\mathscr I_{Y_{\mathscr J}}^2
\]
is annihilated by \(\mathscr I_{Y_{\mathscr J}}\), and is therefore the
direct image of a quasi-coherent \(\mathscr O_{Y_{\mathscr J}}\)-module,
denoted by \(\mathscr N_{Y_{\mathscr J}/X_{\mathscr J}}\), and called the
\emph{conormal sheaf} of \(Y_{\mathscr J}\) in
\(X_{\mathscr J}\).\footnote{Editor's note: The following sentence has
been corrected.}
Since
\(\mathscr J\otimes_{\mathscr O_{S_{\mathscr J}}}
\mathscr O_{Y_{\mathscr J}}\) is annihilated by
\(\mathscr I_{Y_{\mathscr J}}\), the sheaf of commutative groups
\(\mathscr A_{\mathscr J}\) of \SGARefLink{sga3:iii:corollary:4-3-1}{4.3.1} identifies with
\[
 \begin{split}
 &\operatorname{Hom}_{\mathscr O_{Y_{\mathscr J}}}
 \left(
  \mathscr I_{Y_{\mathscr J}}/\mathscr I_{Y_{\mathscr J}}^2,
  \mathscr J\otimes_{\mathscr O_{S_{\mathscr J}}}
   \mathscr O_{Y_{\mathscr J}}
 \right) \\
 &\qquad={}
 \operatorname{Hom}_{\mathscr O_{Y_{\mathscr J}}}
 \left(
  \mathscr N_{Y_{\mathscr J}/X_{\mathscr J}},
  \mathscr J\otimes_{\mathscr O_{S_{\mathscr J}}}
   \mathscr O_{Y_{\mathscr J}}
 \right),
 \end{split}
\]
and hence, for every \(i\geq0\),
\[
 H^i(X_{\mathscr J},\mathscr A_{\mathscr J})
 =H^i\!\left(
  Y_{\mathscr J},
  \operatorname{Hom}_{\mathscr O_{Y_{\mathscr J}}}
  \left(
   \mathscr N_{Y_{\mathscr J}/X_{\mathscr J}},
   \mathscr J\otimes_{\mathscr O_{S_{\mathscr J}}}
    \mathscr O_{Y_{\mathscr J}}
  \right)
 \right).
\]

\footnote{Editor's note: The discussion that follows has been expanded.}
One may now drop the hypothesis that \(Y\) be closed, as follows. First
observe that every open subset \(U_{\mathscr J}\) of
\(X_{\mathscr J}\) is obtained by base change from the open subscheme
\(U\) of \(X\) having the same underlying topological space as
\(U_{\mathscr J}\). Let \(Y_{\mathscr J}\) now be a closed subscheme of
\(U_{\mathscr J}\), flat over \(S_{\mathscr J}\), and let
\(\mathscr I_{Y_{\mathscr J}}\) be the quasi-coherent ideal of
\(\mathscr O_{U_{\mathscr J}}\) defining \(Y_{\mathscr J}\). If
\(Y_{\mathscr J}\) lifts to a subscheme \(Y\) of \(X\), then \(Y\),
having the same underlying topological space as \(Y_{\mathscr J}\), is a
closed subscheme of \(U\). Consequently, the obstruction to lifting
\(Y_{\mathscr J}\) to a subscheme, flat over \(S\), of \(X\) or of
\(U\) is the same; it lies in
\[
 H^1\!\left(
  Y_{\mathscr J},
  \operatorname{Hom}_{\mathscr O_{Y_{\mathscr J}}}
  \left(
   \mathscr N_{Y_{\mathscr J}/X_{\mathscr J}},
   \mathscr J\otimes_{\mathscr O_{S_{\mathscr J}}}
    \mathscr O_{Y_{\mathscr J}}
  \right)
 \right).
\]

Finally, return to the notation of \S0. Let \(\mathscr I\) be a
quasi-coherent ideal of \(\mathscr O_S\) such that
\(\mathscr J\subset\mathscr I\) and \(\mathscr I\mathscr J=0\), and let
\(S_0\) be the closed subscheme of \(S_{\mathscr J}\) defined by
\(\mathscr I\). For every \(S\)-scheme \(Z\), write
\[
 Z_{\mathscr J}=Z\times_S S_{\mathscr J},
 \qquad Z_0=Z\times_S S_0.
\]
Since \(\mathscr J\) is annihilated by \(\mathscr I\), the notation of
\SGARefLink{sga3:iii:corollary:4-4}{4.4} gives
\[
 \mathscr J\otimes_{\mathscr O_{S_{\mathscr J}}}
  \mathscr O_{Y_{\mathscr J}}
 =\mathscr J\otimes_{\mathscr O_{S_0}}\mathscr O_{Y_0},
\]
and
\[
 \begin{split}
 &\operatorname{Hom}_{\mathscr O_{Y_{\mathscr J}}}
 \left(
  \mathscr N_{Y_{\mathscr J}/X_{\mathscr J}},
  \mathscr J\otimes_{\mathscr O_{S_{\mathscr J}}}
   \mathscr O_{Y_{\mathscr J}}
 \right) \\
 &\qquad={}
 \operatorname{Hom}_{\mathscr O_{Y_0}}
 \left(
  \mathscr N_{Y_{\mathscr J}/X_{\mathscr J}}
   \otimes_{\mathscr O_{Y_{\mathscr J}}}\mathscr O_{Y_0},
  \mathscr J\otimes_{\mathscr O_{S_0}}\mathscr O_{Y_0}
 \right),
 \end{split}
\]
and similarly for the other terms. Thus one obtains:

% Source page: Exp3 local p. 48; combined p. 158; original side-page 132.
\medskip
\noindent\SGAInlineRefTarget{sga3:iii:proposition:4-5}\textbf{Proposition 4.5.}
\textit{Let \(S\) be a scheme, and let \(S_0\) and
\(S_{\mathscr J}\) be the closed subschemes defined by quasi-coherent
ideals \(\mathscr I\) and \(\mathscr J\), where
\(\mathscr I\supset\mathscr J\) and
\(\mathscr I\mathscr J=0\). Let \(X\) be an \(S\)-scheme and let
\(Y_{\mathscr J}\) be a subscheme of \(X_{\mathscr J}\), flat over
\(S_{\mathscr J}\). Define the \(\mathscr O_{Y_0}\)-module
\[
 \mathscr A_0=
 \operatorname{Hom}_{\mathscr O_{Y_0}}
 \left(
  \mathscr N_{Y_{\mathscr J}/X_{\mathscr J}}
   \otimes_{\mathscr O_{Y_{\mathscr J}}}\mathscr O_{Y_0},
  \mathscr J\otimes_{\mathscr O_{S_0}}\mathscr O_{Y_0}
 \right).
\]
Then:}
\begin{enumerate}
\item[\SGAInlineRefTarget{sga3:iii:proposition:4-5:item:i}(i)] \textit{A subscheme \(Y\) of \(X\), reducing to
  \(Y_{\mathscr J}\) and flat over \(S\), exists if and only if the
  following two conditions hold:}
  \begin{enumerate}
  \item[\SGAInlineRefTarget{sga3:iii:proposition:4-5:item:a}(a)] \textit{Such a \(Y\) exists locally on \(X\).}
  \item[\SGAInlineRefTarget{sga3:iii:proposition:4-5:item:b}(b)] \textit{A certain obstruction in
    \(R^1\Gamma(Y_0,\mathscr A_0)\) vanishes.}\footnote{Editor's note:
    Here \(R^1\Gamma(Y_0,\mathscr A_0)\) denotes the ``coherent''
    cohomology group \(H^1(Y_0,\mathscr A_0)\) of the
    \(\mathscr O_{Y_0}\)-module \(\mathscr A_0\), in order to distinguish
    it from the ``Hochschild'' cohomology groups \(H^i(Y_0,M_0)\), where
    \(Y_0\) is an \(S_0\)-group and \(M_0\) an
    \(\mathscr O_{S_0}\)-module, which will be considered beginning in
    \SGARefLink{sga3:iii:paragraph:4-16}{4.16}.}
  \end{enumerate}
\item[\SGAInlineRefTarget{sga3:iii:proposition:4-5:item:ii}(ii)] \textit{Under these conditions, the set of subschemes \(Y\)
  satisfying the stated requirements is principal homogeneous under the
  abelian group \(\Gamma(Y_0,\mathscr A_0)\).}
\end{enumerate}

\medskip
\noindent\SGAInlineRefTarget{sga3:iii:remark:4-5-1}\textbf{Remark 4.5.1.}\footnote{Editor's note: This remark has
been placed here; it replaces \SGARefLink{sga3:iii:paragraph:4-7}{Remark 4.7} of the original.}
It follows from \SGARefLink{sga3:iii:proposition:4-5:item:ii}{4.5(ii)} that for every pair \((Y,Y')\) of
subschemes\footnote{Editor's note: ``closed subschemes'' has been corrected
to ``subschemes.''} of \(X\), flat over \(S\) and reducing to
\(Y_{\mathscr J}\), there is a ``deviation''
\[
 \begin{split}
 d(Y',Y)&\in\Gamma(Y_0,\mathscr A_0)\\
 &=\operatorname{Hom}_{\mathscr O_{Y_0}}
 \left(
  \mathscr N_{Y_{\mathscr J}/X_{\mathscr J}}
   \otimes_{\mathscr O_{Y_{\mathscr J}}}\mathscr O_{Y_0},
  \mathscr J\otimes_{\mathscr O_{S_0}}\mathscr O_{Y_0}
 \right).
 \end{split}
\]
With the sign convention adopted in \SGARefLink{sga3:iii:complement:4-4-1}{4.4.1}, \(d(Y',Y)\) corresponds to
the morphism of \(\mathscr O_X\)-modules
\[
 \mathscr I_{Y'}\lhook\joinrel\longrightarrow
 \mathscr O_X\longrightarrow\mathscr O_Y,
\]
which takes values in
\(\mathscr J\mathscr O_Y\simeq
\mathscr J\otimes_{\mathscr O_{S_0}}\mathscr O_{Y_0}\), and factors
through
\(\mathscr I_{Y'_{\mathscr J}}=\mathscr I_{Y_{\mathscr J}}\), then
through \(\mathscr N_{Y_{\mathscr J}/X_{\mathscr J}}\).

\medskip
\noindent\SGAInlineRefTarget{sga3:iii:remark:4-6}\textbf{Remark 4.6.}\footnote{Editor's note: We have retained,
for reference, \SGARefLink{sga3:iii:remark:4-6}{Remark 4.6} of the original, which did not include the
definition of ``locally complete intersection.'' It is followed by the
``correct'' definition, taken from SGA 6, VII 1.4 (and replacing that of
EGA IV\textsubscript{4}, 16.9.2), and by proofs of the three assertions
made in the remark.}
\textit{If \(X\) is flat over \(S\) and \(Y_{\mathscr J}\) is locally
complete intersection in \(X_{\mathscr J}\), then condition \SGARefLink{sga3:iii:proposition:4-5:item:a}{(a)} is
always satisfied, and every \(Y\) flat over \(S\) lifting
\(Y_{\mathscr J}\) is locally complete intersection in \(X\). If,
moreover, \(Y_0\) is affine, condition \SGARefLink{sga3:iii:proposition:4-5:item:b}{(b)} is also satisfied.}

\medskip
\noindent\SGAInlineRefTarget{sga3:iii:definition:4-6-1}\textbf{Definition 4.6.1.}
\textit{(cf. SGA 6, VII 1.1) Let \(B\) be a commutative ring, let
\(f:E\to B\) be a \(B\)-linear morphism, where \(E\) is a free
\(B\)-module of finite rank \(d\), and let \(I=f(E)\). If a basis of
\(E\) is chosen, \(f\) is given by a \(d\)-tuple
\((f_1,\ldots,f_d)\) of elements of \(B\), and \(I\) is the ideal
generated by the \(f_i\). The \emph{Koszul complex} \(K_\bullet(f)\)
is the graded \(B\)-module \(\bigwedge_B^\bullet E\), equipped with the
differential of degree \(-1\)}
\[
 x_1\wedge\cdots\wedge x_i
 \longmapsto
 \sum_{j=1}^i(-1)^{j-1}f(x_j)
 x_1\wedge\cdots\wedge\widehat{x_j}\wedge\cdots\wedge x_i.
\]
Thus there is an augmented chain complex, with \(B/I\) in degree
\(-1\),
\[
 \cdots\longrightarrow\bigwedge^2 E\longrightarrow E
 \xrightarrow{f}B\longrightarrow B/I\longrightarrow0,
\]
which is exact in degree \(0\) by definition, since \(I=f(E)\). The
morphism \(f\) is said to be \emph{regular} if \(K_\bullet(f)\) is
acyclic in degrees \(>0\), that is, if the augmented complex above is a
resolution of \(C=B/I\).

In that case, the proof of SGA 6, VII 1.2(b) shows that the
\(C\)-modules \(I^n/I^{n+1}\), for \(n\in\mathbb N\), are free, with
\(I/I^2\) of rank \(d\).

% Source page: Exp3 local p. 49; combined p. 159; original side-page 132.
\medskip
\noindent\SGAInlineRefTarget{sga3:iii:definition:4-6-2}\textbf{Definition 4.6.2.}
\textit{(cf. SGA 6, VII 1.4) Let \(X\) be a scheme, \(Y\) a subscheme,
and \(U\) an open subset of \(X\) such that \(Y\) is a closed subscheme
of \(U\), defined by the quasi-coherent ideal \(\mathscr I_Y\).}

The subscheme \(Y\) is said to be \emph{locally complete intersection}
in \(X\) if \(Y\hookrightarrow X\) is a regular immersion in the sense
of SGA 6, VII 1.4; that is, if for every \(y\in Y\) there exist an affine
open neighborhood \(V\) of \(y\) in \(U\), a finite free
\(\mathscr O_V\)-module \(\mathscr E\), and a regular morphism
\(f:\mathscr E\to\mathscr O_V\) with image
\(\mathscr I_Y|_V\), such that \(K_\bullet(f)\) is a resolution of
\(\mathscr O_{Y\cap V}\).

This implies that the immersion \(Y\hookrightarrow X\) is locally of
finite presentation and, by \SGARefLink{sga3:iii:definition:4-6-1}{4.6.1}, that the conormal sheaf
\[
 \mathscr N_{Y/X}=\mathscr I_Y/\mathscr I_Y^2
\]
is a finite locally free \(\mathscr O_Y\)-module.

\medskip
\noindent\SGAInlineRefTarget{sga3:iii:lemma:4-6-3}\textbf{Lemma 4.6.3.}\footnote{Editor's note: To prove the
results stated in \SGARefLink{sga3:iii:remark:4-6}{Remark 4.6}, \SGARefLink{sga3:iii:lemma:4-6-3}{Lemmas 4.6.3} and \SGARefLink{sga3:iii:lemma:4-6-4}{4.6.4}, \SGARefLink{sga3:iii:proposition:4-6-5}{Proposition 4.6.5},
and \SGARefLink{sga3:iii:remark:4-6-6}{Remark 4.6.6} have been added.}
\textit{Let \(A\) be a ring, let \(J\) be a square-zero ideal of
\(A\), put \(\overline A=A/J\), let \(B\) be a flat \(A\)-algebra, let
\(E\) be a free \(B\)-module of finite rank, and let \(f:E\to B\) be a
morphism of \(B\)-modules. Suppose that the morphism
\[
 g:\overline E=E\otimes_A\overline A
 \longrightarrow\overline B=B\otimes_A\overline A
\]
induced by \(f\) is regular and that
\(\overline B/g(\overline E)\) is flat over \(\overline A\). Then \(f\)
is regular and \(B/f(E)\) is flat over \(A\).}

\noindent\textit{Proof.}
Put
\[
 C=B/f(E),
 \qquad
 \overline C=C\otimes_A\overline A
 =\overline B/g(\overline E).
\]
First, the \(\bigwedge_B^i E\) are free \(B\)-modules, and hence flat
\(A\)-modules, since \(B\) is flat over \(A\). Since
\[
 \bigwedge_B^\bullet E\otimes_A\overline A
 \simeq\bigwedge_{\overline B}^\bullet\overline E,
\]
there is an exact sequence of complexes
\[
 0\longrightarrow
 J\otimes_A\bigwedge_B^\bullet E
 \longrightarrow\bigwedge_B^\bullet E
 \longrightarrow\bigwedge_{\overline B}^\bullet\overline E
 \longrightarrow0.
\]
Moreover, since \(J^2=0\), for every \(A\)-module \(M\) one has
\[
 J\otimes_A M
 \simeq J\otimes_{\overline A}(\overline A\otimes_A M).
\]
Writing the augmentation arrows as dashed arrows, and denoting the rank
of \(E\) by \(d\), one obtains the following bicomplex, whose rows are
exact.

% Raster stand-in for Loop 1; replace by a native TeX diagram in Loop 2.
% Authority: Exp3 local p. 50; combined reader p. 160; original side-page 132.
\begin{center}
 \SGAThreeNativeInput[0.78\textwidth]
   {native_diagrams/exp3/diagram_038_exp3_p150_koszul_bicomplex.tex}
\end{center}

% Source page: Exp3 local p. 50; combined p. 160; original side-page 132.
The right and left columns are also exact, since \(K_\bullet(g)\) is a
resolution of \(\overline C\), and \(\overline C\) is flat over
\(\overline A\). The long exact homology sequence associated with the
exact sequence of unaugmented complexes therefore shows that
\(K_\bullet(f)\) is acyclic in degrees \(>0\), and in degree \(0\) gives
an exact sequence
\[
 0\longrightarrow J\otimes_{\overline A}\overline C
 \longrightarrow C\longrightarrow\overline C\longrightarrow0.
\]
It follows from the ``fundamental criterion of flatness'' that \(C\) is
flat over \(A\) (cf. [BAC], \S III.5, Theorem 1).
\hfill\(\square\)

\medskip
\noindent\SGAInlineRefTarget{sga3:iii:lemma:4-6-4}\textbf{Lemma 4.6.4.}\footnotemark[88]
\textit{Let \(A\) be a commutative ring, let \(J\) be a nilpotent ideal,
and let \(N\subset M\) be \(A\)-modules such that \(M/N\) is flat over
\(A\). If \(x_1,\ldots,x_n\) are elements of \(N\) whose images generate
the image \(\overline N\) of \(N\) in \(M/JM\), then they generate
\(N\).}

Indeed, let \(N'\) be the submodule of \(N\) generated by the \(x_i\),
and put \(Q=N/N'\). The morphism
\(N'\otimes_A(A/J)\to\overline N\) is surjective. On the other hand,
since \(M/N\) is flat over \(A\), the morphism
\(N\otimes_A(A/J)\to\overline N\) is bijective. Thus
\(Q\otimes_A(A/J)=0\), whence \(Q=0\) by the ``nilpotent Nakayama
lemma'':
\[
 Q=JQ=J^2Q=\cdots=0.
\]

We can now prove the following proposition.

\medskip
\noindent\SGAInlineRefTarget{sga3:iii:proposition:4-6-5}\textbf{Proposition 4.6.5.}\footnotemark[88]
\textit{Let \(S,\mathscr I,\mathscr J\), and
\(X,Y_{\mathscr J}\) be as in \SGARefLink{sga3:iii:proposition:4-5}{4.5}. Suppose in addition that \(X\) is
flat over \(S\) and that \(Y_{\mathscr J}\) is locally complete
intersection in \(X_{\mathscr J}\).}
\begin{enumerate}
\item[\SGAInlineRefTarget{sga3:iii:proposition:4-6-5:item:a}(a)] \textit{Condition \SGARefLink{sga3:iii:proposition:4-6-5:item:a}{(a)} of \SGARefLink{sga3:iii:proposition:4-5:item:i}{4.5(i)} is then satisfied; moreover,
  every \(Y\) flat over \(S\) lifting \(Y_{\mathscr J}\) is locally
  complete intersection in \(X\).}
\item[\SGAInlineRefTarget{sga3:iii:proposition:4-6-5:item:b}(b)] \textit{If, in addition, \(Y_0\) is affine, condition \SGARefLink{sga3:iii:proposition:4-6-5:item:b}{(b)} of
  the cited result is also satisfied.}
\end{enumerate}

% Source page: Exp3 local p. 51; combined p. 161; original side-page 132.
\noindent\textit{Proof.}
The first assertion of \SGARefLink{sga3:iii:proposition:4-6-5:item:a}{(a)} follows from \SGARefLink{sga3:iii:lemma:4-6-3}{Lemma 4.6.3}, and the second then
follows from \SGARefLink{sga3:iii:lemma:4-6-4}{Lemma 4.6.4}. Moreover, the hypothesis implies (cf. \SGARefLink{sga3:iii:definition:4-6-2}{4.6.2})
that \(\mathscr N_{Y_{\mathscr J}/X_{\mathscr J}}\) is a finite locally
free \(\mathscr O_{Y_{\mathscr J}}\)-module. Hence the
\(\mathscr O_{Y_0}\)-module
\[
 \mathscr A_0=
 \operatorname{Hom}_{\mathscr O_{Y_0}}
 \left(
  \mathscr N_{Y_{\mathscr J}/X_{\mathscr J}}
   \otimes_{\mathscr O_{Y_{\mathscr J}}}\mathscr O_{Y_0},
  \mathscr J\otimes_{\mathscr O_{S_0}}\mathscr O_{Y_0}
 \right)
\]
is quasi-coherent (cf. EGA I, 1.3.12), and therefore
\(R^1\Gamma(Y_0,\mathscr A_0)=0\) when \(Y_0\) is affine.
\hfill\(\square\)

\medskip
\noindent\SGAInlineRefTarget{sga3:iii:remark:4-6-6}\textbf{Remark 4.6.6.}\footnotemark[88]
Let us conclude this paragraph with an example showing that, under the
hypotheses of \SGARefLink{sga3:iii:lemma:4-6-3}{Lemma 4.6.3}, a regular sequence
\((\overline g_1,\overline g_2)\) generating the ideal
\(\overline I=g(\overline E)\) need not lift to a regular sequence in
\(B\).

Let \(k\) be a field, put \(\overline A=k[X,Y]\), let
\(k\varepsilon\) denote the \(\overline A\)-module
\(\overline A/(X,Y)\), so that
\[
 P\varepsilon=P(0,0)\varepsilon
 \qquad(P\in\overline A),
\]
and let \(A=\overline A\oplus k\varepsilon\), where
\(J=k\varepsilon\) is a square-zero ideal. Thus
\(A/J=\overline A\).

The algebra \(B=A\otimes_k k[Z,T]\) is free, and hence flat, over
\(A\), and \(\overline B=k[X,Y,Z,T]\). Put
\[
 \overline g_1=XZ-YT,
 \qquad
 \overline g_2=XZ-1.
\]
Since \(\overline g_1\) is irreducible,
\(\overline B/(\overline g_1)\) is an integral domain; hence
\((\overline g_1,\overline g_2)\) is a regular sequence in
\(\overline B\), generating the ideal
\[
 \overline I=(XZ-1,YT-1).
\]
Therefore
\[
 \overline C=\overline B/\overline I
 =k[X,Y,X^{-1},Y^{-1}]
 =\overline A[X^{-1},Y^{-1}]
\]
is a flat \(\overline A\)-algebra, and also a flat \(A\)-algebra. But
every lift of \(\overline g_1\) to \(B\) has the form
\[
 XZ-YT+\lambda\varepsilon,
 \qquad \lambda\in k[Z,T],
\]
and therefore annihilates \(\varepsilon\).

\smallskip
\noindent{\footnotesize\textit{Source note.} The Polo--Gille PDF prints
\(XY-ZT+\lambda\varepsilon\) for this lift. Since the element being
lifted is \(\overline g_1=XZ-YT\), the displayed lift above restores
that same polynomial before adding \(\lambda\varepsilon\).}

\medskip
\noindent\SGAInlineRefTarget{sga3:iii:paragraph:4-7}\textbf{4.7.}
\SGARefLink{sga3:iii:paragraph:4-7}{Remark 4.7} has been suppressed here and placed in \SGARefLink{sga3:iii:remark:4-5-1}{4.5.1}.

% Continue with Remark 4.8.0 on Exp3 local p. 51 / combined p. 161.

% END INCLUDED SGA3_English_current_buildable_source_20260823/source/inputs/III/tex/components/26_expose_III_s450_s47_lci_koszul.tex

% BEGIN INCLUDED SGA3_English_current_buildable_source_20260823/source/inputs/III/tex/components/27_expose_III_s480_s410_obstruction_factorization.tex
% SGA 3, Expose III, Loop 1 English reconstruction.
% Scope: Remark 4.8.0 through the 4.10 suppression notice, stopping
% immediately before 4.11.
% Authority: Polo--Gille Exp3-14oct24.pdf, component-local pp. 51--55
% (combined-reader pp. 161--165; original side-pages 133--136), with only
% the 4.8.0--4.10 portions of the first and last pages admitted. Force-OCR
% is a locator only. English wording was compared with Jacob C. Reinhold's
% CC BY 4.0 Markdown lineage at commit
% e7a259f3f8608ad3edf9bf6eead3fd504dd2d23e.
% The Polo--Gille PDF controls prose, mathematics, numbering, notes, and the
% four source-derived Loop-1 diagram images included below.

\setcounter{footnote}{88}

% Source page: Exp3 local p. 51; combined p. 161; original side-page 133.
\medskip
\noindent\SGAInlineRefTarget{sga3:iii:remark:4-8-0}\textbf{Remark 4.8.0.}\footnote{Editor's note: This remark,
used in the proposition that follows, has been inserted here; it appeared
in \SGARefLink{sga3:iii:paragraph:4-10}{4.10} of the original.}
\textit{Let \(S\) be a scheme, \(S'\) a closed subscheme of \(S\), \(X\) an
\(S\)-scheme, \(Y\) an \(S\)-subscheme of \(X\), and put
\(X'=X\times_S S'\) and \(Y'=Y\times_S S'\). Then there is a
surjective morphism of \(\mathscr O_{Y'}\)-modules}
\[
 \mathscr N_{Y/X}\otimes_{\mathscr O_Y}\mathscr O_{Y'}
 \twoheadrightarrow \mathscr N_{Y'/X'}.
\]

Indeed, after replacing \(X\) by a suitable open subscheme, one may
suppose that \(Y\) is closed and is defined by an ideal
\(\mathscr I_Y\) of \(\mathscr O_X\). The image of \(\mathscr I_Y\) in
\(\mathscr O_{X'}\) is then the ideal \(\mathscr I_{Y'}\) defining
\(Y'\), and there is a surjective morphism of
\(\mathscr O_{Y'}\)-modules
\[
 \pi:
 (\mathscr I_Y/\mathscr I_Y^2)
 \otimes_{\mathscr O_Y}\mathscr O_{Y'}
 \twoheadrightarrow
 \mathscr I_{Y'}/\mathscr I_{Y'}^2.
\]
Suppose in addition that
\(\mathscr O_Y=\mathscr O_X/\mathscr I_Y\) is flat over
\(\mathscr O_S\). The natural morphism
\[
 \mathscr I_Y\otimes_{\mathscr O_X}\mathscr O_{X'}
 \longrightarrow \mathscr I_{Y'}
\]
is then bijective (cf. EGA IV\textsubscript{2}, 2.1.8).

% Source page: Exp3 local p. 52; combined p. 162; original side-page 133.
There is therefore the following commutative diagram with exact rows:

% Raster stand-in for Loop 1; replace by a native TeX diagram in Loop 2.
% Authority: Exp3 local p. 52; combined reader p. 162.
\begin{center}
 \SGAThreeNativeInput[0.96\textwidth]
   {native_diagrams/exp3/diagram_039_exp3_p052_remark480_exact_rows.tex}
\end{center}

It follows from the snake lemma that\footnote{Editor's note: In the
original, this was stated in \SGARefLink{sga3:iii:paragraph:4-10}{Remark 4.10} under the additional hypothesis
that \(Y'\) be locally complete intersection in \(X'\). Consequently,
that hypothesis also appeared in statements \SGARefLink{sga3:iii:lemma:4-12}{4.12}--\SGARefLink{sga3:iii:corollary:4-14}{4.14}; it seems in fact
superfluous, and has been deleted from those statements.}
\begin{equation}
 \mathscr N_{Y/X}\otimes_{\mathscr O_Y}\mathscr O_{Y'}
 \xrightarrow{\sim}\mathscr N_{Y'/X'}
 \qquad\text{if \(Y\) is flat over \(S\).}
 \tag{4.8.0}\label{sga3:iii:remark:4-8-0:equation:4-8-0}
\end{equation}

\medskip
\noindent\SGAInlineRefTarget{sga3:iii:proposition:4-8}\textbf{Proposition 4.8.}
\textit{Let \(S,S_0,S_{\mathscr J}\), and
\(\mathscr I,\mathscr J\) be as in \SGARefLink{sga3:iii:proposition:4-5}{4.5}.}\footnote{Editor's note: The
hypothesis that \(\mathscr I\) be nilpotent has been deleted, since it
appears to be superfluous (cf. the proof).}
\textit{Let \(X\) be an \(S\)-scheme, let \(Y\) be a subscheme of
\(X\), and let \(i:Y\hookrightarrow X\) be the immersion.}
\begin{enumerate}
\item[\SGAInlineRefTarget{sga3:iii:proposition:4-8:item:i}(i)] \textit{For every \(S\)-morphism \(f:T\to X\) such that
  \(f_{\mathscr J}:T_{\mathscr J}\to X_{\mathscr J}\) factors through
  \(Y_{\mathscr J}\), one can define an obstruction}
  \begin{equation}
   c(X,Y,f)\in
   \operatorname{Hom}_{\mathscr O_{T_0}}
   \left(
    f_0^*\bigl(\mathscr N_{Y/X}
      \otimes_{\mathscr O_Y}\mathscr O_{Y_0}\bigr),
    \mathscr J\mathscr O_T
   \right)
   \tag{*}\label{sga3:iii:proposition:4-8:equation:star}
  \end{equation}
  \textit{whose vanishing is equivalent to the existence of a
  factorization of \(f\) through \(Y\).}

\item[\SGAInlineRefTarget{sga3:iii:proposition:4-8:item:ii}(ii)] \textit{Let \(Y'\) be a second subscheme of \(X\). Suppose
  that \(Y'_{\mathscr J}=Y_{\mathscr J}\) and that \(Y\) and \(Y'\) are
  flat over \(S\). Then there are isomorphisms (cf. \SGARefLink{sga3:iii:remark:4-8-0}{4.8.0})}
  \[
   \mathscr J\mathscr O_Y
   \simeq
   \mathscr J\otimes_{\mathscr O_{S_0}}\mathscr O_{Y_0}
   \simeq
   \mathscr J\mathscr O_{Y'},
   \qquad
   \mathscr N_{Y/X}\otimes_{\mathscr O_Y}
     \mathscr O_{Y_{\mathscr J}}
   \xrightarrow{\sim}
   \mathscr N_{Y_{\mathscr J}/X_{\mathscr J}},
  \]
  \textit{and hence an isomorphism}
  \[
   \begin{split}
   u:\ &\operatorname{Hom}_{\mathscr O_{Y_0}}
   \left(
    \mathscr N_{Y_{\mathscr J}/X_{\mathscr J}}
      \otimes_{\mathscr O_{Y_{\mathscr J}}}\mathscr O_{Y_0},
    \mathscr J\mathscr O_Y
   \right)\\
   &\xrightarrow{\sim}
   \operatorname{Hom}_{\mathscr O_{Y_0}}
   \left(
    \mathscr N_{Y/X}\otimes_{\mathscr O_Y}\mathscr O_{Y_0},
    \mathscr J\mathscr O_{Y'}
   \right).
   \end{split}
  \]
  \textit{Writing \(i':Y'\hookrightarrow X\) for the canonical
  immersion and \(d(Y,Y')\) for the deviation of \SGARefLink{sga3:iii:remark:4-5-1}{4.5.1}, one has}\footnote{Editor's
  note: See also \SGARefLink{sga3:iii:lemma:4-27}{4.27} below.}
  \begin{equation}
   c(X,Y,i')=u\bigl(d(Y,Y')\bigr).
   \tag{**}\label{sga3:iii:proposition:4-8:equation:double-star}
  \end{equation}

\item[\SGAInlineRefTarget{sga3:iii:proposition:4-8:item:iii}(iii)] \textit{The canonical morphism}
  \[
   \mathscr N_{Y/X}\xrightarrow{D}i^*(\Omega^1_{X/S})
  \]
  \textit{(cf. SGA 1, II, formula 4.3)}\footnote{Editor's note: See also
  EGA IV\textsubscript{4}, 16.4.21. Recall that if \(U\) is an affine
  open subset of \(X\) such that \(Y\cap U\) is defined by the ideal
  \(I\) of \(A=\mathscr O_X(U)\), if \(d\) denotes the differential
  \(A\to\Gamma(U,\Omega^1_{X/S})\), and if \(x\in I\), then
  \(D(x+I^2)\) is the element \(d(x)\otimes1\) of
  \(\Gamma(U,\Omega^1_{X/S})\otimes_A(A/I)\).}
  \textit{induces a morphism}
  \[
   D_0:
   \mathscr N_{Y/X}\otimes_{\mathscr O_Y}\mathscr O_{Y_0}
   \longrightarrow
   \Omega^1_{X_0/S_0}\otimes_{\mathscr O_{X_0}}\mathscr O_{Y_0},
  \]
  \textit{and hence, for every \(S\)-morphism \(f:T\to X\) as in
  \SGARefLink{sga3:iii:proposition:4-8:item:i}{(i)}, a morphism}
  \[
   \begin{split}
   v_{f_0}:\ &
   \operatorname{Hom}_{\mathscr O_{T_0}}
   \left(f_0^*(\Omega^1_{X_0/S_0}),\mathscr J\mathscr O_T\right)
   \\[-0.15em]
   &\longrightarrow
   \operatorname{Hom}_{\mathscr O_{T_0}}
   \left(
    f_0^*\bigl(\mathscr N_{Y/X}
      \otimes_{\mathscr O_Y}\mathscr O_{Y_0}\bigr),
    \mathscr J\mathscr O_T
   \right),\\
   &\hspace{8em}a\longmapsto a\circ f_0^*(D_0),
   \end{split}
  \]
  \textit{where the first group above is
  \(\operatorname{Hom}_{X^+}(T,L_X)\), cf. \SGARefLink{sga3:iii:definition:0-1-5}{0.1.5}. For
  \(a\in\operatorname{Hom}_{X^+}(T,L_X)\), one has}
  \begin{equation}
   c(X,Y,a\cdot f)-c(X,Y,f)=v_{f_0}(a),
   \tag{***}\label{sga3:iii:proposition:4-8:equation:double-starstar}
  \end{equation}
  \textit{where \(a\cdot f\) denotes the composite morphism}
  \[
   T\xrightarrow{a\times f}L_X\times_{X^+}X\longrightarrow X.
  \]
\end{enumerate}

We shall prove part (i) of the proposition, leaving it to the reader to
(not) verify assertions (ii) and (iii); this verification proceeds by
reduction to the affine case and then by comparison of the explicit
definitions.\footnote{Editor's note: These verifications have been
carried out below.}

Let us therefore prove \SGARefLink{sga3:iii:proposition:4-8:item:i}{(i)}. The morphism \(f:T\to X\) defines a
morphism of sheaves of rings\footnote{Editor's note: On the one hand,
the hypothesis that \(\mathscr I\) be nilpotent---that is, that \(X_0\)
have the same underlying topological space as \(X\)---has been deleted;
on the other hand, the sentence that follows has been expanded.}
\[
 \phi:\mathscr O_X\longrightarrow f_*(\mathscr O_T).
\]
Let \(U\) be an open subscheme of \(X\) in which \(Y\) is closed. Since
\(T\), respectively \(Y_{\mathscr J}\), has the same underlying space
as \(T_{\mathscr J}\), respectively \(Y\), the continuous map underlying
\(f\) sends \(T\) into \(U\). Since \(U\) is open in \(X\), the morphism
\(\phi\) induces a morphism of sheaves of rings
\[
 \mathscr O_U=\mathscr O_X|_U\longrightarrow f_*(\mathscr O_T),
\]
that is, \(f\) factors through \(U\).

% Source page: Exp3 local p. 53; combined p. 163;
% original side-pages 134--135.
We may therefore restrict to the case in which \(Y\) is closed, hence
defined by a sheaf of ideals \(\mathscr I_Y\). The morphism \(f\)
factors through \(Y\) if and only if the composite
\[
 \mathscr I_Y\longrightarrow\mathscr O_X
 \longrightarrow f_*(\mathscr O_T)
\]
is zero. Since \(f_{\mathscr J}\) factors through
\(Y_{\mathscr J}\), the composite
\[
 \mathscr I_{Y_{\mathscr J}}\longrightarrow
 \mathscr O_{X_{\mathscr J}}\longrightarrow
 f_*(\mathscr O_{T_{\mathscr J}})
\]
is zero. Consider the following commutative diagram, whose first row is
exact:

% Raster stand-in for Loop 1; replace by a native TeX diagram in Loop 2.
% Authority: Exp3 local p. 53; combined reader p. 163.
\begin{center}
 \SGAThreeNativeInput[0.82\textwidth]
   {native_diagrams/exp3/diagram_040_exp3_p053_prop48_obstruction_diagram.tex}
\end{center}

It follows that \(\phi\) maps \(\mathscr I_Y\) into
\(f_*(\mathscr J\mathscr O_T)\).\footnote{Editor's note: The argument
that follows has been expanded.} Since \(\mathscr J^2=0\), the
\(\mathscr O_X\)-module \(f_*(\mathscr J\mathscr O_T)\), with its
\(\mathscr O_X\)-module structure induced by \(\phi\), is annihilated by
\(\mathscr I_Y\). Consequently, \(\phi\) induces a morphism of
\(\mathscr O_X\)-modules
\[
 h:i_*(\mathscr N_{Y/X})
   =\mathscr I_Y/\mathscr I_Y^2
 \longrightarrow f_*(\mathscr J\mathscr O_T).
\]

% Source page: Exp3 local p. 54; combined p. 164;
% original side-page 135.
On the other hand, one has the following cartesian squares:

% Raster stand-in for Loop 1; replace by a native TeX diagram in Loop 2.
% Authority: Exp3 local p. 54; combined reader p. 164.
\begin{center}
 \SGAThreeNativeInput[0.46\textwidth]
   {native_diagrams/exp3/diagram_041_exp3_p054_prop48_cartesian_squares.tex}
\end{center}

Here \(\tau_{T_0},\tau_{X_0},\tau_{Y_0}\) are the closed immersions
obtained by base change from \(S_0\hookrightarrow S\). Since
\(\mathscr J\mathscr O_T\) is a quasi-coherent
\(\mathscr O_T\)-module annihilated by \(\mathscr I\), there is an
isomorphism
\[
 \mathscr J\mathscr O_T
 \simeq
 (\tau_{T_0})_*\tau_{T_0}^*(\mathscr J\mathscr O_T),
\]
and hence
\[
 f_*(\mathscr J\mathscr O_T)
 \simeq
 (\tau_{X_0})_*(f_0)_*\tau_{T_0}^*(\mathscr J\mathscr O_T).
\]
Thus, by adjunction, \(h\) corresponds to a morphism of
\(\mathscr O_{T_0}\)-modules
\[
 h_0:
 f_0^*\tau_{X_0}^*i_*(\mathscr N_{Y/X})
 \longrightarrow
 \tau_{T_0}^*(\mathscr J\mathscr O_T).
\]
\smallskip
\noindent{\footnotesize\textit{Source notation note.} The
Polo--Gille prose and this \(h_0\)-display print \(i_{T_0}\) for the
closed immersion labeled \(\tau_{T_0}\) in the cartesian square and the
preceding base-change formulas. We use \(\tau_{T_0}\) consistently.}

Now
\[
 \tau_{X_0}^*i_*(\mathscr N_{Y/X})
 \simeq
 (i_0)_*\tau_{Y_0}^*(\mathscr N_{Y/X})
 =\mathscr N_{Y/X}\otimes_{\mathscr O_Y}\mathscr O_{Y_0}.
\]
Returning to the customary abuse of notation
\(\tau_{T_0}^*(\mathscr J\mathscr O_T)=\mathscr J\mathscr O_T\), the
morphism \(h_0\) identifies with a morphism of
\(\mathscr O_{T_0}\)-modules
\[
 h_0:
 f_0^*\bigl(\mathscr N_{Y/X}
   \otimes_{\mathscr O_Y}\mathscr O_{Y_0}\bigr)
 \longrightarrow\mathscr J\mathscr O_T,
\]
which is the required obstruction \(c(X,Y,f)\). This proves \SGARefLink{sga3:iii:proposition:4-8:item:i}{(i)}.

When \(f\) is the immersion \(i':Y'\hookrightarrow X\), one sees that
\(c(X,Y,i')\) comes from the morphism
\[
 \mathscr I_Y\lhook\joinrel\longrightarrow\mathscr O_X
 \longrightarrow\mathscr O_{Y'},
\]
and hence corresponds, by \SGARefLink{sga3:iii:complement:4-4-1}{4.4.1} and \SGARefLink{sga3:iii:remark:4-5-1}{4.5.1}, to the class \(d(Y,Y')\).
This proves (ii).

\smallskip
\noindent{\footnotesize\textit{Source notation note.} At this point the
Polo--Gille proof prints \(i':Y\hookrightarrow X\), while part \SGARefLink{sga3:iii:proposition:4-8:item:ii}{(ii)}
defines \(i'\) as the immersion of \(Y'\). The latter domain is used
here.}

Let us prove (iii). First,
\(D:\mathscr N_{Y/X}\to i^*(\Omega^1_{X/S})\) induces a morphism
\[
 \begin{split}
 D_0:\quad
 \tau_{Y_0}^*(\mathscr N_{Y/X})
 &\longrightarrow
 \tau_{Y_0}^*i^*(\Omega^1_{X/S})\\
 &=i_0^*\tau_{X_0}^*(\Omega^1_{X/S}).
 \end{split}
\]
Since \(X_0=X\times_S S_0\), one has
\[
 \tau_{X_0}^*(\Omega^1_{X/S})\simeq\Omega^1_{X_0/S_0}
 \qquad\text{(cf. EGA IV\textsubscript{4}, 16.4.5).}
\]
This gives the asserted morphism
\[
 D_0:
 \mathscr N_{Y/X}\otimes_{\mathscr O_Y}\mathscr O_{Y_0}
 \longrightarrow
 \Omega^1_{X_0/S_0}\otimes_{\mathscr O_{X_0}}\mathscr O_{Y_0}.
\]
We shall verify equality \((***)\) after the following remark.

\medskip
\noindent\SGAInlineRefTarget{sga3:iii:remark:4-9}\textbf{Remark 4.9.}
The obstruction \(c(X,Y,f)\) is computed locally on \(T\). Let
\(U=\operatorname{Spec}(C)\) be an affine open subset of \(T\) lying
over an affine open subset \(\operatorname{Spec}(A)\) of \(X\), itself
lying over an affine open subset \(\operatorname{Spec}(\Lambda)\) of
\(S\). Let \(J\subset I\subset\Lambda\), respectively
\(I_Y\subset A\), be the ideals corresponding to
\(\mathscr J\subset\mathscr I\), respectively \(\mathscr I_Y\), put
\(B=A/I_Y\), and let \(\phi:A\to C\) be the morphism of
\(\Lambda\)-algebras corresponding to \(f:T\to X\). Since
\(f(T_{\mathscr J})\subset Y_{\mathscr J}\), one has
\(\phi(I_Y)\subset JC\), and \(\phi\) therefore induces a morphism of
\(\Lambda\)-algebras
\[
 B\longrightarrow C/JC\longrightarrow C_0=C/IC.
\]
The obstruction \(c=c(X,Y,f)\) is then computed by the following
commutative diagram:

% Raster stand-in for Loop 1; replace by a native TeX diagram in Loop 2.
% Authority: Exp3 local p. 54; combined reader p. 164.
\begin{center}
 \SGAThreeNativeInput[0.60\textwidth]
   {native_diagrams/exp3/diagram_042_exp3_p054_remark49_local_obstruction.tex}
\end{center}

% Source page: Exp3 local p. 55; combined p. 165;
% original side-page 136.
That is, over \(U\), it is the unique element of
\[
 \operatorname{Hom}_{C_0}
 \left(I_Y/I_Y^2\otimes_B C_0,JC\right)
 =
 \operatorname{Hom}_{B_0}
 \left(I_Y/I_Y^2\otimes_B B_0,JC\right)
\]
such that, with the evident notation,
\[
 c(\overline x\otimes_B1)=\phi(x)
 \qquad\text{for every \(x\in I_Y\).}
\]

\footnote{Editor's note: What follows has been added.}
We can now complete the proof of \SGARefLink{sga3:iii:proposition:4-8:item:iii}{4.8(iii)}. Equality \((***)\) may be
checked locally on \(T\), so we are reduced to the affine situation
described above. Let \(d_{A/\Lambda}:A\to\Omega^1_{A/\Lambda}\) denote
the differential. Over \(U\), the element \(a\) corresponds to an
element \(a_U\) of
\[
 \begin{split}
 &\operatorname{Hom}_{C_0}
 \left(\Omega^1_{A_0/\Lambda_0}\otimes_{A_0}C_0,JC\right)
 \\[-0.15em]
 &\qquad\simeq
 \operatorname{Hom}_{B_0}
 \left(\Omega^1_{A/\Lambda}\otimes_A B_0,JC\right)
 \simeq
 \operatorname{Hom}_{A}(\Omega^1_{A/\Lambda},JC).
 \end{split}
\]
On the one hand, \(v_{f_0}(a)\) corresponds over \(U\) to
\(a_U\circ D_0\), where \(D_0\) is the morphism of
\(B\)-modules\footnote{Editor's note: Cf. Editor's Note 93.}
\[
 \begin{aligned}
 I_Y/I_Y^2&\longrightarrow\Omega^1_{A/\Lambda}\otimes_A B_0,\\
 x+I_Y^2&\longmapsto d_{A/\Lambda}(x)\otimes1.
 \end{aligned}
\]
On the other hand (cf. the proofs of \SGARefLink{sga3:iii:theorem:0-1-8}{0.1.8} and \SGARefLink{sga3:iii:corollary:0-1-9}{0.1.9}), the morphism of
\(\Lambda\)-algebras \(\phi':A\to C\) corresponding to \(a\cdot f\)
differs from \(\phi\) by the \(\Lambda\)-derivation \(A\to JC\)
associated with \(a_U\); that is,
\[
 \phi'=\phi+a_U\circ d_{A/\Lambda}
 =\phi+a_U\circ(d_{A/\Lambda}\otimes1).
\]
Consequently, writing \(c'=c(X,Y,a\cdot f)\), for every \(x\in I_Y\),
and writing \(\overline x\) for its image in \(I_Y/I_Y^2\), one has
\[
 \begin{split}
 (c'-c)(\overline x\otimes1)
 &=a_U\bigl(d_{A/\Lambda}(x)\otimes1\bigr)\\
 &=(a_U\circ D_0)(\overline x)
 =v_{f_0}(a)(\overline x).
 \end{split}
\]
This proves that \(c'-c=v_{f_0}(a)\).

\medskip
\noindent\SGAInlineRefTarget{sga3:iii:paragraph:4-10}\textbf{4.10.}
\SGARefLink{sga3:iii:paragraph:4-10}{Remark 4.10} of the original has been suppressed, since it was made
obsolete by the addition of \SGARefLink{sga3:iii:remark:4-8-0}{Remark 4.8.0}.

% Continue with 4.11 on Exp3 local p. 55 / combined p. 165 /
% original side-page 136.

% END INCLUDED SGA3_English_current_buildable_source_20260823/source/inputs/III/tex/components/27_expose_III_s480_s410_obstruction_factorization.tex

% BEGIN INCLUDED SGA3_English_current_buildable_source_20260823/source/inputs/III/tex/components/28_expose_III_s411_s4161_functorial_obstructions.tex
% SGA 3, Expose III, Loop 1 English reconstruction.
% Scope: 4.11 through Definition 4.16.1, stopping immediately before
% Lemma 4.17.
% Authority: Polo--Gille Exp3-14oct24.pdf, component-local pp. 55--62
% (combined-reader pp. 165--172; original printed pp. 136--142), with only
% the section-4.11 portion of local p. 55 admitted. Force-OCR is a locator
% only. English wording was compared with Jacob C. Reinhold's CC BY 4.0
% Markdown lineage at commit e7a259f3f8608ad3edf9bf6eead3fd504dd2d23e.
% The Polo--Gille PDF controls prose, mathematics, numbering, notes, and the
% thirteen source-derived Loop-1 diagram images included below.

\setcounter{footnote}{98}

% Source page: Exp3 local p. 55; combined p. 165; original printed p. 136.
\medskip
\noindent\SGAInlineRefTarget{sga3:iii:paragraph:4-11}\textbf{4.11.}
We now propose to study the following situation. Let \(S\),
\(S_{\mathscr J}\), and \(S_0\) be as in \SGARefLink{sga3:iii:proposition:4-8}{4.8}. We have three
\(S\)-schemes \(X,X',T\), a subscheme \(Y\) of \(X\) (respectively a
subscheme \(Y'\) of \(X'\)), and morphisms \(f:T\to X'\) and
\(g:X'\to X\).

% Raster stand-in for Loop 1; replace by native TeX/TikZ in Loop 2.
% Authority: Exp3-14oct24.pdf local p. 55; combined reader p. 165.
\begin{center}
  \SGAThreeNativeInput[0.62\textwidth]
    {native_diagrams/exp3/diagram_043_exp3_p055_411_setup.tex}
\end{center}

% Source page: Exp3 local p. 56; combined p. 166; original printed p. 137.
We suppose that, after reduction modulo \(\mathscr J\), this diagram can
be completed to a commutative diagram

% Raster stand-in for Loop 1; replace by native TeX/TikZ in Loop 2.
% Authority: Exp3-14oct24.pdf local p. 56; combined reader p. 166.
\begin{center}
  \SGAThreeNativeInput[0.72\textwidth]
    {native_diagrams/exp3/diagram_044_exp3_p056_411_reduction.tex}
\end{center}

By \SGARefLink{sga3:iii:proposition:4-8}{4.8} we therefore have the obstructions
\begin{align*}
 c(X,Y,g\circ i')&\in
 \operatorname{Hom}_{\mathscr O_{Y'_0}}
 \bigl(i_0'^*g_0^*(\mathscr N_{Y/X}
   \otimes_{\mathscr O_Y}\mathscr O_{Y_0}),
   \mathscr J\mathscr O_{Y'}\bigr),\\
 c(X',Y',f)&\in
 \operatorname{Hom}_{\mathscr O_{T_0}}
 \bigl(f_0^*(\mathscr N_{Y'/X'}
   \otimes_{\mathscr O_{Y'}}\mathscr O_{Y'_0}),
   \mathscr J\mathscr O_T\bigr),\\
 c(X,Y,g\circ f)&\in
 \operatorname{Hom}_{\mathscr O_{T_0}}
 \bigl(f_0^*g_0^*(\mathscr N_{Y/X}
   \otimes_{\mathscr O_Y}\mathscr O_{Y_0}),
   \mathscr J\mathscr O_T\bigr),
\end{align*}
whose relations we seek to compute.
\footnote{Editor's note: Starting with \SGARefLink{sga3:iii:lemma:4-17}{4.17}, this will be applied when
\(X\) is an \(S\)-group, \(g:X\times_S X\to X\) is multiplication,
\(Y\) is a subscheme of \(X\) such that \(Y_{\mathscr J}\) is a
subgroup of \(X_{\mathscr J}\), \(Y'=Y\times_S Y\), and the two
morphisms \(Y^3\to X^2\) send \((y_1,y_2,y_3)\) respectively to
\((y_1y_2,y_3)\) and \((y_1,y_2y_3)\). In that case, comparison of the
obstructions above will show that the obstruction to \(Y\) being a
subgroup of \(X\) lies in a certain Hochschild cohomology group
\(H^2(Y_0,N_0)\).}

\medskip
\noindent\SGAInlineRefTarget{sga3:iii:lemma:4-12}\textbf{Lemma 4.12.}
\textit{Suppose that \(Y'\) is flat over \(S\), so that
\[
 \mathscr J\otimes_{\mathscr O_{S_0}}\mathscr O_{Y'_0}
 =\mathscr J\mathscr O_{Y'}.
\]
Then:}
\begin{enumerate}
\item[(i)] \textit{There is a natural morphism}
\begin{align*}
 b_{f_0}:&\
 \operatorname{Hom}_{\mathscr O_{Y'_0}}
 \bigl(i_0'^*g_0^*(\mathscr N_{Y/X}
   \otimes_{\mathscr O_Y}\mathscr O_{Y_0}),
   \mathscr J\mathscr O_{Y'}\bigr)\\
 &\longrightarrow
 \operatorname{Hom}_{\mathscr O_{T_0}}
 \bigl(f_0^*g_0^*(\mathscr N_{Y/X}
   \otimes_{\mathscr O_Y}\mathscr O_{Y_0}),
   \mathscr J\mathscr O_T\bigr).
\end{align*}
\item[(ii)] \textit{There is also a natural morphism, functorial in
\(T\),}
\footnote{Editor's note: We have removed the hypothesis that \(Y'_0\)
be locally complete intersection in \(X'_0\), which is superfluous by
\SGARefLink{sga3:iii:remark:4-8-0}{4.8.0}. We have also added that \(a_{g_0}(f_0)\) is
\lq\lq functorial in \(T\)\rq\rq, a fact that plays a crucial role in
the proof of \SGARefLink{sga3:iii:lemma:4-17}{4.17}.}
\begin{align*}
 a_{g_0}(f_0):&\
 \operatorname{Hom}_{\mathscr O_{T_0}}
 \bigl(f_0^*(\mathscr N_{Y'/X'}
   \otimes_{\mathscr O_{Y'}}\mathscr O_{Y'_0}),
   \mathscr J\mathscr O_T\bigr)\\
 &\longrightarrow
 \operatorname{Hom}_{\mathscr O_{T_0}}
 \bigl(f_0^*g_0^*(\mathscr N_{Y/X}
   \otimes_{\mathscr O_Y}\mathscr O_{Y_0}),
   \mathscr J\mathscr O_T\bigr).
\end{align*}
\end{enumerate}

\noindent\textit{Proof.}
\footnote{Editor's note: The proof has been expanded to make the
\lq\lq functoriality in \(T\)\rq\rq\ of \(a_{g_0}\) explicit.}
First note that, with \(X,X',Y,Y'\) fixed, giving a \(T\) as above is
equivalent to giving a morphism
\((f,f_{\mathscr J}):T\to X'\times_{X'^{+}}Y'^{+}\). Put
\(Z=X'\times_{X'^{+}}Y'^{+}\), and let \(M\) and \(M'\) be the
\(Z\)-functors defined, for every \(f:T\to Z\), by
\begin{align*}
 M(T)&=
 \operatorname{Hom}_{\mathscr O_{T_0}}
 \bigl(f_0^*g_0^*(\mathscr N_{Y/X}
   \otimes_{\mathscr O_Y}\mathscr O_{Y_0}),
   \mathscr J\mathscr O_T\bigr),\\
 M'(T)&=
 \operatorname{Hom}_{\mathscr O_{T_0}}
 \bigl(f_0^*(\mathscr N_{Y'/X'}
   \otimes_{\mathscr O_{Y'}}\mathscr O_{Y'_0}),
   \mathscr J\mathscr O_T\bigr).
\end{align*}
\footnote{Editor's note: The situation simplifies from \SGARefLink{sga3:iii:paragraph:4-16}{4.16} onward:
we shall restrict to schemes flat over \(S\), \(Y\) will be a flat
\(S\)-group, and \(Y'=Y\times_S Y\); this will give \(S_0\)-functors
\(N_0\) and \(N'_0\).}

% Source page: Exp3 local p. 57; combined p. 167; original printed p. 138.
In all cases there is a commutative diagram

% Raster stand-in for Loop 1; replace by native TeX/TikZ in Loop 2.
% Authority: Exp3-14oct24.pdf local p. 57; combined reader p. 167.
\begin{center}
  \SGAThreeNativeInput[0.65\textwidth]
    {native_diagrams/exp3/diagram_045_exp3_p057_412_tensor_square.tex}
\end{center}

Since \(Y'\) is flat over \(S\), the left vertical arrow is an
isomorphism, and hence we obtain a morphism of
\(\mathscr O_{T_0}\)-modules
\[
 f_0^*(\mathscr J\mathscr O_{Y'})\longrightarrow
 \mathscr J\mathscr O_T.
\]
It induces a morphism of abelian groups
\begin{multline*}
 \operatorname{Hom}_{\mathscr O_{T_0}}
 \bigl(f_0^*g_0^*(\mathscr N_{Y/X}
  \otimes_{\mathscr O_Y}\mathscr O_{Y_0}),
  f_0^*(\mathscr J\mathscr O_{Y'})\bigr)\\
 \longrightarrow M(T),
\end{multline*}
and, on composing with the morphism
\begin{multline*}
 M(Y')\longrightarrow
 \operatorname{Hom}_{\mathscr O_{T_0}}
 \bigl(f_0^*g_0^*(\mathscr N_{Y/X}
  \otimes_{\mathscr O_Y}\mathscr O_{Y_0}),
  f_0^*(\mathscr J\mathscr O_{Y'})\bigr)
\end{multline*}
induced by \(f_0^*\), we obtain the morphism
\(b_{f_0}:M(Y')\to M(T)\).

Likewise, in all cases there is a diagram

% Raster stand-in for Loop 1; replace by native TeX/TikZ in Loop 2.
% Authority: Exp3-14oct24.pdf local p. 57; combined reader p. 167.
\begin{center}
  \SGAThreeNativeInput[0.67\textwidth]
    {native_diagrams/exp3/diagram_046_exp3_p057_412_conormal_square.tex}
\end{center}

Since \(Y'\) is flat over \(S\), the second vertical arrow is an
isomorphism by \SGARefLink{sga3:iii:remark:4-8-0:equation:4-8-0}{4.8.0}. We therefore obtain an
\(\mathscr O_{Y'_0}\)-morphism
\[
 i_0'^*g_0^*(\mathscr N_{Y/X}
   \otimes_{\mathscr O_Y}\mathscr O_{Y_0})
 \longrightarrow
 \mathscr N_{Y'/X'}
   \otimes_{\mathscr O_{Y'}}\mathscr O_{Y'_0}.
\]
This induces a morphism
\(a_{g_0}(\operatorname{id}_{Y'}):M'(Y')\to M(Y')\) and, for every
\(f:T\to Z\), a morphism \(a_{g_0}(f):M'(T)\to M(T)\), such that the
following diagram commutes:

% Raster stand-in for Loop 1; replace by native TeX/TikZ in Loop 2.
% Authority: Exp3-14oct24.pdf local p. 57; combined reader p. 167.
\begin{center}
  \SGAThreeNativeInput[0.50\textwidth]
    {native_diagrams/exp3/diagram_047_exp3_p057_412_functor_square.tex}
\end{center}

Here \(b'_{f_0}\) is defined in the same way as \(b_{f_0}\).
\hfill\textsc{q.e.d.}

\medskip
\noindent\SGAInlineRefTarget{sga3:iii:remark:4-12-1}\textbf{Remark 4.12.1.}
\footnote{Editor's note: This remark has been added for use in the
proof of \SGARefLink{sga3:iii:lemma:4-17}{4.17}.}
\textit{Let \(M_0\) and \(M'_0\) be the \(Y'_0\)-functors defined, for
every \(f:T_0\to Y'_0\), by}
\begin{align*}
 M_0(T)&=
 \operatorname{Hom}_{\mathscr O_{T_0}}
 \bigl(f_0^*g_0^*(\mathscr N_{Y/X}
  \otimes_{\mathscr O_Y}\mathscr O_{Y_0}),
  \mathscr J\otimes_{\mathscr O_{S_0}}\mathscr O_{T_0}\bigr),\\
 M'_0(T)&=
 \operatorname{Hom}_{\mathscr O_{T_0}}
 \bigl(f_0^*(\mathscr N_{Y'/X'}
  \otimes_{\mathscr O_{Y'}}\mathscr O_{Y'_0}),
 \mathscr J\otimes_{\mathscr O_{S_0}}\mathscr O_{T_0}\bigr).
\end{align*}

% Source page: Exp3 local p. 58; combined p. 168; original printed p. 139.
Note at once that \(Z_0=Y'_0\), and that on the category of
\(Z\)-schemes \(T\) flat over \(S\), the functors \(M\) and \(M'\)
coincide respectively with
\(\prod_{S_0/S}M_0\) and \(\prod_{S_0/S}M'_0\). In that case,
\(b_{f_0}\) is simply the morphism
\[
 f_0^*:M_0(Y'_0)\longrightarrow M(T_0)
\]
induced by \(f_0\). Moreover, for every morphism \(u:U\to T\), if
\(h=f\circ u\), the following diagram commutes:

% Raster stand-in for Loop 1; replace by native TeX/TikZ in Loop 2.
% Authority: Exp3-14oct24.pdf local p. 58; combined reader p. 168.
\begin{center}
  \SGAThreeNativeInput[0.50\textwidth]
    {native_diagrams/exp3/diagram_048_exp3_p058_4121_naturality_square.tex}
\end{center}

Thus \(a_{g_0}\) becomes a morphism of functors
\[
 \prod_{S_0/S}M'_0\longrightarrow\prod_{S_0/S}M_0.
\]

\medskip
\noindent\SGAInlineRefTarget{sga3:iii:proposition:4-13}\textbf{Proposition 4.13.}
\textit{Suppose that \(Y'\) is flat over \(S\). Then}
\[
 c(X,Y,g\circ f)=a_{g_0}\bigl(c(X',Y',f)\bigr)
   +b_{f_0}\bigl(c(X,Y,g\circ i')\bigr).
\]

Since the definitions of the various obstructions and of the morphisms
\(a_{g_0}\) and \(b_{f_0}\) are local, it plainly suffices to verify the
formula when all the schemes in question are affine. Write
\[
 \begin{gathered}
 S=\operatorname{Spec}(\Lambda),\qquad
 S_{\mathscr J}=\operatorname{Spec}(\Lambda/\mathscr J),\qquad
 S_0=\operatorname{Spec}(\Lambda/\mathscr I),\\
 T=\operatorname{Spec}(C),\qquad X=\operatorname{Spec}(A),\qquad
 Y=\operatorname{Spec}(A/I_Y)=\operatorname{Spec}(B),\\
 X'=\operatorname{Spec}(A'),\qquad
 Y'=\operatorname{Spec}(A'/I_{Y'})=\operatorname{Spec}(B').
 \end{gathered}
\]
There is therefore a diagram of rings and ideals
\footnote{Editor's note: We retain the notation of the original,
writing \(f:A'\to C\) and \(g:A\to A'\) for the ring homomorphisms
corresponding to \(f:T\to X'\) and \(g:X'\to X\). This explains the
formula \(c(X,Y,g\circ f)(\bar x\otimes1)=f(g(x))\), for \(x\in I_Y\).}

% Raster stand-in for Loop 1; replace by native TeX/TikZ in Loop 2.
% Authority: Exp3-14oct24.pdf local p. 58; combined reader p. 168.
\begin{center}
  \SGAThreeNativeInput[0.48\textwidth]
    {native_diagrams/exp3/diagram_049_exp3_p058_413_ring_ideal_diagram.tex}
\end{center}

Let us study the terms of the formula to be proved. In what follows,
for \(x\in I_Y\) (respectively \(u\in I_{Y'}\)), write \(\bar x\)
(respectively \(\bar u\)) for its image in \(I_Y/I_Y^2\)
(respectively \(I_{Y'}/I_{Y'}^2\)). If \(m\) belongs to a
\(\Lambda\)-module \(M\), write \(\bar m_0\) for its image in
\(M_0=M/\mathscr I M\).

We saw that \(c=c(X,Y,g\circ f)\) is the unique \(C_0\)-morphism
\[
 I_Y/I_Y^2\otimes_B C_0\longrightarrow\mathscr J C
\]
such that \(c(\bar x\otimes1)=f(g(x))\) for every \(x\in I_Y\).

% Source page: Exp3 local p. 59; combined p. 169; original printed p. 140.
Fix \(x\in I_Y\). Since \(g_{\mathscr J}(Y'_{\mathscr J})\subset
Y_{\mathscr J}\), one has \(g(x)\in I_{Y'}+\mathscr J A'\). Write
\[
 g(x)=x'+\sum_i\lambda_i a'_i,
 \qquad x'\in I_{Y'},\quad \lambda_i\in\mathscr J,\quad a'_i\in A'.
\]
Then
\begin{equation*}
 \tag{1}\label{sga3:iii:proposition:4-13:equation:1}
 c(X,Y,g\circ f)(\bar x\otimes1)
 =f(g(x))=f(x')+\sum_i\lambda_i f(a'_i).
\end{equation*}

Consider \(a_{g_0}(c(X',Y',f))\). By the definitions just given, it is
defined by the following diagram:

% Raster stand-in for Loop 1; replace by native TeX/TikZ in Loop 2.
% Authority: Exp3-14oct24.pdf local p. 59; combined reader p. 169.
\begin{center}
  \SGAThreeNativeInput[0.86\textwidth]
    {native_diagrams/exp3/diagram_050_exp3_p059_413_ag0_obstruction_diagram.tex}
\end{center}

Thus
\(a_{g_0}(c(X',Y',f))(\bar x\otimes1)=f(u)\), where \(u\in I_{Y'}\)
has image \(\bar u\) in \(I_{Y'_0}/I_{Y'_0}^2\) satisfying
\[
 \bar u_0\otimes1=\bar g_0(\bar x_0)\otimes1
 =g_0(\bar x_0)\otimes1.
\]
We may therefore take \(u=x'\), which gives
\begin{equation*}
 \tag{2}\label{sga3:iii:proposition:4-13:equation:2}
 a_{g_0}(c(X',Y',f))(\bar x\otimes1)=f(x').
\end{equation*}

Finally consider \(b_{f_0}(c(X,Y,g\circ i'))\). By hypothesis, the
morphism of \(\Lambda_0\)-algebras \(f_0:A'_0\to C_0\) factors through
\(B'_0\). Since \(Y'\) is flat, and hence
\[
 \mathscr J\otimes_{\Lambda_0}B'_0\xrightarrow{\sim}\mathscr J B',
\]
we obtain a morphism of \(B'_0\)-modules
\(\psi:\mathscr J B'\to\mathscr J C\) making the following diagram
commute:

% Raster stand-in for Loop 1; replace by native TeX/TikZ in Loop 2.
% Authority: Exp3-14oct24.pdf local p. 59; combined reader p. 169.
\begin{center}
  \SGAThreeNativeInput[0.74\textwidth]
    {native_diagrams/exp3/diagram_051_exp3_p059_413_tensor_J_square.tex}
\end{center}

Let
\[
 \varphi:\mathscr J B'\otimes_{B'_0}C_0\longrightarrow\mathscr J C
\]
be the morphism of \(C_0\)-modules induced by \(\psi\). Then, for every
\(a'\in A'\) and \(\lambda\in\mathscr J\),
\begin{equation*}
\tag{\(\dagger\)}\label{sga3:iii:proposition:4-13:equation:dagger}
 \varphi(\lambda\pi'(a')\otimes1)=\lambda f(a').
\end{equation*}

% Source page: Exp3 local p. 60; combined p. 170; original printed p. 141.
The element \(b_{f_0}(c(X,Y,g\circ i'))\) is then defined by the
commutative diagram

% Raster stand-in for Loop 1; replace by native TeX/TikZ in Loop 2.
% Authority: Exp3-14oct24.pdf local p. 60; combined reader p. 170.
\begin{center}
  \SGAThreeNativeInput[0.62\textwidth]
    {native_diagrams/exp3/diagram_052_exp3_p060_413_bf0_obstruction_diagram.tex}
\end{center}

Consequently,
\begin{equation*}
 \tag{3}\label{sga3:iii:proposition:4-13:equation:3}
 b_{f_0}(c(X,Y,g\circ i'))(\bar x\otimes1)
 =\varphi\left(\sum_i\lambda_i\pi'(a'_i)\otimes1\right)
 =\sum_i\lambda_i f(a'_i),
\end{equation*}
the last equality following from \((\dagger)\). Comparison of the three
explicit formulas \SGARefLink{sga3:iii:proposition:4-13:equation:1}{(1)}, \SGARefLink{sga3:iii:proposition:4-13:equation:2}{(2)}, and \SGARefLink{sga3:iii:proposition:4-13:equation:3}{(3)} gives the required formula.

\medskip
\noindent\SGAInlineRefTarget{sga3:iii:corollary:4-14}\textbf{Corollary 4.14.}
\textit{Let \(Y,Y'\) be two flat subschemes of \(X\), both reducing to
\(Y_{\mathscr J}\), and suppose that \(Y_0\) is locally complete
intersection in \(X_0\). If \(f:T\to X\) is an \(S\)-morphism such that
\(f_{\mathscr J}\) factors through \(Y_{\mathscr J}\to X_{\mathscr J}\),
then}
\[
 c(X,Y,f)-c(X,Y',f)=b_{f_0}\bigl(d(Y,Y')\bigr).
\]

Indeed, applying the preceding formula to the diagram

% Raster stand-in for Loop 1; replace by native TeX/TikZ in Loop 2.
% Authority: Exp3-14oct24.pdf local p. 60; combined reader p. 170.
\begin{center}
  \SGAThreeNativeInput[0.62\textwidth]
    {native_diagrams/exp3/diagram_053_exp3_p060_414_comparison_diagram.tex}
\end{center}

we find
\(c(X,Y,f)-c(X,Y',f)=b_{f_0}(c(X,Y,i'))\). Moreover, by \SGARefLink{sga3:iii:proposition:4-8:item:ii}{4.8(ii)},
\(c(X,Y,i')=d(Y,Y')\).

\medskip
\noindent\SGAInlineRefTarget{sga3:iii:proposition:4-15}\textbf{Proposition 4.15.}
\textit{Let \(X\) be an \(S\)-group smooth over \(S\), and let \(Y\)
be a sub-\(S\)-group flat and locally of finite presentation over \(S\).
Then \(Y\) is locally complete intersection in \(X\) (cf. \SGARefLink{sga3:iii:definition:4-6-2}{4.6.2}).}

% Source page: Exp3 local p. 61; combined p. 171; original printed p. 142.
\noindent\textit{Proof.}
\footnote{Editor's note: We have added to the statement the hypothesis
that \(Y\) be locally of finite presentation over \(S\), and supplied
the proof that follows, which is more direct than the one sketched in
the original. For completeness, we also spell out the latter. As in
the proof above, first reduce to the case \(S=\operatorname{Spec}(k)\),
where \(k\) is algebraically closed. By EGA IV\textsubscript{4},
16.9.10 and 19.3.2, it suffices to see that, for every \(y\in Y\), the
completion of the local ring \(\mathscr O_{Y,y}\) is a quotient of a
complete noetherian local ring by a regular sequence. By loc. cit.,
19.3.3, the set of \(y\in Y\) with this property is an open subset
\(U\subset Y\). Since \(Y\) is of finite type over \(k\), it suffices
to show that \(U\) contains every closed point. Because \(Y\) is a
\(k\)-group, a translation argument reduces this to the completion of
\(\mathscr O_{Y,e}\), that is, to the formal group \(\widehat Y\)
corresponding to \(Y\) (cf. Expos\'e VII\textsubscript{B}). Since \(X\)
is smooth, the affine algebra \(\mathscr A(\widehat X)\) is a power
series algebra \(k[[X_1,\ldots,X_n]]\). The Dieudonn\'e structure
theorem then shows that \(\mathscr A(\widehat Y)\) is isomorphic to a
quotient
\[
 k[[X_1,\ldots,X_{r+s}]]/
 (X_1^{p^{n_1}},\ldots,X_r^{p^{n_r}}),
\]
where \(p\) is the characteristic exponent of \(k\) and \(r+s\leq n\);
cf. Expos\'e VII\textsubscript{B}, Remark 5.5.2(b).}
We show that the immersion \(Y\to X\) is regular in the sense of EGA
IV\textsubscript{4}, 16.9.2. By EGA IV\textsubscript{4}, 19.5.1, this
implies that it is regular also in the sense of 4.6.2 (indeed, by loc.
cit., the two definitions are equivalent when \(S\) is locally
noetherian). Thus \lq\lq regular immersion\rq\rq\ will here mean regular
in the sense of EGA IV\textsubscript{4}, 16.9.2. Since \(X\) and \(Y\)
are flat and locally of finite presentation over \(S\), EGA
IV\textsubscript{4}, 19.2.4 reduces the assertion to proving that
\(Y_s\to X_s\) is a regular immersion for every \(s\in S\). By EGA
IV\textsubscript{4}, 19.1.5(ii), it is enough to check this on the
geometric fibers of \(S\), and hence when \(S\) is the spectrum of an
algebraically closed field \(k\).

By Expos\'e VI\textsubscript{A}, \SGARefLink{thm:sga3-VIA-3.2}{3.2}, the quotient \(X/Y\) exists and is
smooth, the morphism \(\pi:X\to X/Y\) is flat, and the following square
is cartesian:

% Raster stand-in for Loop 1; replace by native TeX/TikZ in Loop 2.
% Authority: Exp3-14oct24.pdf local p. 61; combined reader p. 171.
\begin{center}
  \SGAThreeNativeInput[0.35\textwidth]
    {native_diagrams/exp3/diagram_054_exp3_p061_415_cartesian_square.tex}
\end{center}

Here \(\bar e\) is the image in \(X/Y\) of the unit point of \(X\).
By flat base change (cf. EGA IV\textsubscript{4}, 19.1.5(iii)), it
therefore suffices to show that \(i\) is a regular immersion. This is
immediate: the noetherian local ring \(\mathscr O_{X/Y,\bar e}\) is
smooth (hence regular), so its maximal ideal is generated by a regular sequence.

\medskip
\noindent\SGAInlineRefTarget{sga3:iii:paragraph:4-16}\textbf{4.16.}
\footnote{Editor's note: We have reorganized \SGARefLink{sga3:iii:paragraph:4-16}{4.16} by collecting here,
on the one hand, the hypotheses stated at the end of \SGARefLink{sga3:iii:proposition:4-15}{4.15} and, on the
other, the definition of the obstruction \(DY\).}
Let \(X\) be an \(S\)-group smooth over \(S\), and denote its group law
by \(\mu:X\times_S X\to X\). Let \(Y_{\mathscr J}\) be a
sub-\(S_{\mathscr J}\)-group of \(X_{\mathscr J}\), flat and locally of
finite presentation over \(S_{\mathscr J}\). By \SGARefLink{sga3:iii:proposition:4-15}{4.15},
\(Y_{\mathscr J}\) is locally complete intersection in \(X_{\mathscr J}\). Hence,
by \SGARefLink{sga3:iii:proposition:4-6-5}{4.6.5}, every flat
\footnote{Editor's note: The original has been corrected by adding the
word \lq\lq flat\rq\rq.}
\(S\)-scheme \(Y\) lifting \(Y_{\mathscr J}\) is locally complete
intersection in \(X\). For such a \(Y\), \SGARefLink{sga3:iii:remark:4-8-0}{4.8.0} gives
\begin{equation*}
 \tag{4.16.1}\label{sga3:iii:paragraph:4-16:equation:4-16-1}
 \mathscr N_{Y/X}\otimes_{\mathscr O_Y}\mathscr O_{Y_0}
 =\mathscr N_{Y_0/X_0}
 =\mathscr N_{Y_{\mathscr J}/X_{\mathscr J}}
   \otimes_{\mathscr O_{Y_{\mathscr J}}}\mathscr O_{Y_0}.
\end{equation*}

% Source page: Exp3 local p. 62; combined p. 172; original printed p. 142.
On the other hand, let \(\varepsilon_0:S_0\to Y_0\) be the unit
section of \(Y_0\), and let \(\mathbf n_{Y_0/X_0}\) be the
quasi-coherent \(\mathscr O_{S_0}\)-module
\[
\mathbf n_{Y_0/X_0}=\varepsilon_0^*(\mathscr N_{Y_0/X_0}).
\]
Since \(Y_0\) and \(X_0\) are \(S_0\)-groups, it is immediate that
\(\mathscr N_{Y_0/X_0}\) is invariant under, say, left translations by
\(Y_0\). Hence
\footnote{Editor's note: See \SGARefLink{sga3:iii:lemma:4-25}{4.25} below.}
it is the inverse image of \(\mathbf n_{Y_0/X_0}\) under
\(Y_0\to S_0\); that is,
\begin{equation*}
 \tag{4.16.2}\label{sga3:iii:paragraph:4-16:equation:4-16-2}
 \mathscr N_{Y_0/X_0}
=\mathbf n_{Y_0/X_0}\otimes_{\mathscr O_{S_0}}\mathscr O_{Y_0}.
\end{equation*}

Taking (\SGARefLink{sga3:iii:paragraph:4-16:equation:4-16-1}{4.16.1}) and (\SGARefLink{sga3:iii:paragraph:4-16:equation:4-16-2}{4.16.2}) into account, we deduce first from \SGARefLink{sga3:iii:proposition:4-5}{4.5} that
the set of subschemes \(Y\) of \(X\), flat over \(S\) and lifting
\(Y_{\mathscr J}\), is either empty or principal homogeneous under
\begin{equation*}
 \tag{4.16.3}\label{sga3:iii:paragraph:4-16:equation:4-16-3}
 \operatorname{Hom}_{\mathscr O_{Y_0}}
\bigl(\mathbf n_{Y_0/X_0}\otimes_{\mathscr O_{S_0}}\mathscr O_{Y_0},
       \mathscr J\otimes_{\mathscr O_{S_0}}\mathscr O_{Y_0}\bigr).
\end{equation*}
On the other hand, \SGARefLink{sga3:iii:proposition:4-8:item:i}{4.8(i)} shows that, for every such \(Y\) and every
\(S\)-morphism \(f:T\to X\) such that
\(f_{\mathscr J}:T_{\mathscr J}\to X_{\mathscr J}\) factors through
\(Y_{\mathscr J}\), the obstruction \(c(X,Y,f)\) to \(f\) factoring
through \(Y\) belongs to
\[
 \operatorname{Hom}_{\mathscr O_{T_0}}
\bigl(\mathbf n_{Y_0/X_0}\otimes_{\mathscr O_{S_0}}\mathscr O_{T_0},
       \mathscr J\mathscr O_T\bigr).
\]
If \(T\) is also flat over \(S\), this last group is equal to
\[
 \operatorname{Hom}_{\mathscr O_{T_0}}
\bigl(\mathbf n_{Y_0/X_0}\otimes_{\mathscr O_{S_0}}\mathscr O_{T_0},
       \mathscr J\otimes_{\mathscr O_{S_0}}\mathscr O_{T_0}\bigr).
\]
This leads us to introduce the following group functor \(N_0\).

\medskip
\noindent\SGAInlineRefTarget{sga3:iii:definition:4-16-1}\textbf{Definition 4.16.1.}
\textit{Let \(N_0\) be the \(S_0\)-functor in commutative groups
defined, for every \(Z\in\operatorname{Ob}(\mathbf{Sch})_{/S_0}\), by}
\begin{equation*}
 \tag{*}\label{sga3:iii:definition:4-16-1:equation:star}
 \operatorname{Hom}_{S_0}(Z,N_0)
 =\operatorname{Hom}_{\mathscr O_Z}
\bigl(\mathbf n_{Y_0/X_0}\otimes_{\mathscr O_{S_0}}\mathscr O_Z,
       \mathscr J\otimes_{\mathscr O_{S_0}}\mathscr O_Z\bigr).
\end{equation*}
\textit{Then the set of subschemes \(Y\) of \(X\), flat over \(S\) and
lifting \(Y_{\mathscr J}\), is either empty or principal homogeneous
under}
\[
 \operatorname{Hom}_{S_0}(Y_0,N_0)=C^1(Y_0,N_0).
\]

For every such \(Y\), consider the diagram

% Raster stand-in for Loop 1; replace by native TeX/TikZ in Loop 2.
% Authority: Exp3-14oct24.pdf local p. 62; combined reader p. 172.
\begin{center}
  \SGAThreeNativeInput[0.52\textwidth]
    {native_diagrams/exp3/diagram_055_exp3_p062_4161_DY_diagram.tex}
\end{center}

and let
\[
 DY=c\bigl(X,Y,\mu\circ(i,i)\bigr)
\]
be the obstruction to \(\mu\circ(i,i)\) factoring through \(Y\), that
is, to \(Y\) being stable under the group law of \(X\). By what
precedes, \(DY\) is an element of
\[
 N_0(Y_0\times_{S_0}Y_0)=C^2(Y_0,N_0).
\]

% END INCLUDED SGA3_English_current_buildable_source_20260823/source/inputs/III/tex/components/28_expose_III_s411_s4161_functorial_obstructions.tex

% BEGIN INCLUDED SGA3_English_current_buildable_source_20260823/source/inputs/III/tex/components/29_expose_III_s417_s422_subgroup_obstruction.tex
% SGA 3, Expose III, Loop 1 English reconstruction.
% Scope: Lemma 4.17 through Remark 4.22, stopping immediately before 4.23.
% Authority: Polo--Gille Exp3-14oct24.pdf, component-local pp. 63--67
% (combined-reader pp. 173--177; printed pp. 143--148), with only the
% 4.17--4.22 portion of the last page admitted. OCR is a locator only.
% English wording was compared with Jacob C. Reinhold's CC BY 4.0 Markdown
% lineage at commit e7a259f3f8608ad3edf9bf6eead3fd504dd2d23e. The
% Polo--Gille PDF controls prose, mathematics, numbering, editorial notes,
% and the four source-derived Loop-1 diagram images included below.

\setcounter{footnote}{108}

% Source page: Exp3 local p. 63; combined p. 173; printed pp. 143--144.
\medskip
\noindent\SGAInlineRefTarget{sga3:iii:lemma:4-17}\textbf{Lemma 4.17.}\footnote{Editor's note: \SGARefLink{sga3:iii:lemma:4-17}{Sections 4.17} and
\SGARefLink{sga3:iii:paragraph:4-18}{4.18} have been modified to take account of the additions made in \SGARefLink{sga3:iii:paragraph:4-16}{4.16}.}
\textit{Let \(X\) be an \(S\)-group smooth over \(S\), and let
\(Y_{\mathscr J}\) be a sub-\(S_{\mathscr J}\)-group of
\(X_{\mathscr J}\), flat and locally of finite presentation over
\(S_{\mathscr J}\). For every subscheme \(Y\) of \(X\), flat over \(S\)
and lifting \(Y_{\mathscr J}\), consider the obstruction defined in
\SGARefLink{sga3:iii:definition:4-16-1}{4.16.1}:}
\[
 DY\in
 \operatorname{Hom}_{S_0}
 \bigl(Y_0\mathbin{\times_{S_0}}Y_0,N_0\bigr)
 =C^2(Y_0,N_0).
\]
\begin{enumerate}
\item[\SGAInlineRefTarget{sga3:iii:lemma:4-17:item:i}(i)] \textit{For \(Y\) to be a sub-\(S\)-group of \(X\), it is
  necessary and sufficient that \(DY=0\).}
\item[\SGAInlineRefTarget{sga3:iii:lemma:4-17:item:ii}(ii)] \textit{If \(Y_0\) is made to act on \(N_0\), by
  functoriality, through the inner automorphisms of \(Y_0\), then
  \(DY\in Z^2(Y_0,N_0)\).}
\item[\SGAInlineRefTarget{sga3:iii:lemma:4-17:item:iii}(iii)] \textit{If \(Y\) and \(Y'\) are two subschemes of \(X\),
  flat over \(S\), lifting \(Y_{\mathscr J}\) (so that the deviation
  \(d(Y,Y')\in C^1(Y_0,N_0)\) is defined; cf. \SGARefLink{sga3:iii:remark:4-5-1}{4.5.1}), then}
  \[
   DY'-DY=\partial^1d(Y,Y').
  \]
  \footnote{Editor's note: The original has
  \(DY'-DY=-\partial d(Y,Y')\), but the operator
  \(\partial\) there is the negative of the differential \(\partial^1\)
  defined in I, \SGARefLink{sga3:i:subsection:5-1}{5.1}.}
\end{enumerate}

We prove these assertions in succession.

\medskip
\noindent\SGAInlineRefTarget{sga3:iii:paragraph:4-18}\textbf{4.18. --- Proof of \SGARefLink{sga3:iii:lemma:4-17:item:i}{4.17(i)}.}
By definition, \(DY=0\) if and only if \(Y\) is stable under
the group law of \(X\). Thus \(DY=0\) if \(Y\) is a subgroup
of \(X\). Conversely, if \(DY=0\), then \(Y\) is endowed with
the induced law \(\mu_Y\), which is associative and reduces modulo
\(\mathscr J\) to the group law on \(X_{\mathscr J}\). Since \(Y\) is
flat and locally of finite presentation over \(S\), it follows from \SGARefLink{sga3:iii:proposition:3-3}{3.3}
that \(\mu_Y\) is a group law.

\medskip
\noindent\SGAInlineRefTarget{sga3:iii:paragraph:4-19}\textbf{4.19. --- Proof of \SGARefLink{sga3:iii:lemma:4-17:item:ii}{4.17(ii)}.}
We compare the two values of
\(u=c(X,Y,\mu^2\circ(i,i,i))\) computed in the following two diagrams
\((D_j)\), \(j=1,2\):

% Raster stand-in for Loop 1; replace by a native TeX diagram in Loop 2.
% Authority: Exp3 local p. 63; combined reader p. 173; printed p. 144.
\begin{center}
 \SGAThreeNativeInput[0.94\textwidth]
   {native_diagrams/exp3/diagram_056_exp3_p063_Dj_associativity.tex}
\end{center}

Here \(f_1=(1,\pi)\), \(f_2=(\pi,1)\), and \(\mu^2\) denotes the
morphism
\[
 \mu\circ f_1=\mu\circ f_2:
 X\mathbin{\times_S}X\mathbin{\times_S}X\longrightarrow X.
\]
\footnote{Editor's note: The notation has been slightly modified, and
the beginning of the argument has been expanded.}
Put
\[
 \mu_Y=\mu\circ(i,i),\qquad
 f_j^Y=f_j\circ(i,i,i),\qquad
 \mu^{2,Y}=\mu^2\circ(i,i,i).
\]
For \(j=1,2\), let \(a_j\) and \(b_j\) be the morphisms
\[
 a_j=a_{\mu_0}\bigl((f_j^Y)_0\bigr),
 \qquad
 b_j=b_{(f_j^Y)_0},
\]
associated by \SGARefLink{sga3:iii:lemma:4-12}{Lemma 4.12} with the pair of morphisms \((f_j^Y,\mu)\).
Thus
\begin{equation}
\tag{\(\dagger\)}\label{sga3:iii:paragraph:4-19:equation:dagger}
\left\{
\begin{aligned}
 a_j:\;&
 \operatorname{Hom}_{\mathscr O_{Y_0^3}}
 \left((f_j^Y)_0^*
  \mathscr N_{Y_0\times Y_0/X_0\times X_0},
  \mathscr J\mathscr O_{Y_0^3}\right)
 \longrightarrow\\[-0.15em]
 &\qquad
 \operatorname{Hom}_{\mathscr O_{Y_0^3}}
 \left((\mu^{2,Y})_0^*\mathscr N_{Y_0/X_0},
  \mathscr J\mathscr O_{Y_0^3}\right),\\[0.35em]
 b_j:\;&
 \operatorname{Hom}_{\mathscr O_{Y_0^2}}
 \left((\mu_Y)_0^*\mathscr N_{Y_0/X_0},
  \mathscr J\mathscr O_{Y_0^2}\right)
 \longrightarrow\\[-0.15em]
 &\qquad
 \operatorname{Hom}_{\mathscr O_{Y_0^3}}
 \left((\mu^{2,Y})_0^*\mathscr N_{Y_0/X_0},
  \mathscr J\mathscr O_{Y_0^3}\right).
\end{aligned}
\right.
\end{equation}

% Source page: Exp3 local p. 64; combined p. 174; printed pp. 144--145.
Since
\[
 \mathscr N_{Y_0\times Y_0/X_0\times X_0}
 \simeq
 \operatorname{pr}_1^*\mathscr N_{Y_0/X_0}
 \oplus
 \operatorname{pr}_2^*\mathscr N_{Y_0/X_0}
\]
(because \(X_0\) and \(Y_0\) are flat over \(S_0\)), and
\[
 \mathscr N_{Y_0/X_0}
 \simeq
 \mathbf n_{Y_0/X_0}\otimes_{\mathscr O_{S_0}}\mathscr O_{Y_0},
\]
we have
\[
 (f_j^Y)_0^*\mathscr N_{Y_0\times Y_0/X_0\times X_0}
 \simeq
 (\mathbf n_{Y_0/X_0}\oplus \mathbf n_{Y_0/X_0})
 \otimes\mathscr O_{Y_0^3},
\]
and likewise
\[
 \begin{aligned}
 (\mu^{2,Y})_0^*\mathscr N_{Y_0/X_0}
 &\simeq \mathbf n_{Y_0/X_0}\otimes\mathscr O_{Y_0^3},\\
 (\mu_Y)_0^*\mathscr N_{Y_0/X_0}
 &\simeq \mathbf n_{Y_0/X_0}\otimes\mathscr O_{Y_0^2}.
 \end{aligned}
\]
Moreover, since \(Y_0^2\) and \(Y_0^3\) are flat over \(S_0\),
\((\dagger)\) may be rewritten as
\begin{equation}
\tag{\(\ddagger\)}\label{sga3:iii:paragraph:4-19:equation:ddagger}
\left\{
\begin{aligned}
 a_j&:\operatorname{Hom}_{S_0}(Y_0^3,N_0\oplus N_0)
       \longrightarrow\operatorname{Hom}_{S_0}(Y_0^3,N_0),\\
 b_j&:\operatorname{Hom}_{S_0}(Y_0^2,N_0)
       \longrightarrow\operatorname{Hom}_{S_0}(Y_0^3,N_0).
\end{aligned}
\right.
\end{equation}

Applying \SGARefLink{sga3:iii:proposition:4-13}{4.13} twice to \(c(X,Y,\mu^{2,Y})=u\), we obtain
\[
 \begin{aligned}
 a_1\bigl(c(X^2,Y^2,f_1)\bigr)
 +b_1\bigl(c(X,Y,\mu_Y)\bigr)
 &=u\\
 &=a_2\bigl(c(X^2,Y^2,f_2)\bigr)
 +b_2\bigl(c(X,Y,\mu_Y)\bigr).
 \end{aligned}
\]
Now \(c(X,Y,\mu_Y)=DY\), and, since
\(f_1=(1,\mu)\) and \(f_2=(\mu,1)\), we have, with the evident notation,
\[
 c(X^2,Y^2,f_1)=(0,DY),
 \qquad
 c(X^2,Y^2,f_2)=(DY,0).
\]
Consequently,
\[
 u=a_1((0,DY))+b_1(DY)
  =a_2((DY,0))+b_2(DY).
\]
The first point to note is that \(b_j\) is simply
\(\operatorname{Hom}_{S_0}((f_j^Y)_0,N_0)\), that is, the morphism
deduced from \((f_j^Y)_0\) by functoriality. The preceding identity
therefore becomes
\[
 \begin{aligned}
 a_1((0,DY))
 &-\operatorname{Hom}((\mu,1),N_0)(DY)
 +\operatorname{Hom}((1,\mu),N_0)(DY)\\
 &\hspace{9em}-a_2((DY,0))=0.
 \end{aligned}
\]
The two middle terms are the \(DY(xy,z)\) and
\(DY(x,yz)\) parts of the formula for the \(2\)-coboundary.
It remains only to identify the other two terms.

We first compute the map \(a_j\). It is obtained, by pullback along
\((f_j^Y)_0\), from the morphism of \(\mathscr O_{Y_0^2}\)-modules
\[
 P:\mathbf n_{Y_0/X_0}\otimes\mathscr O_{Y_0^2}
 \longrightarrow
 (\mathbf n_{Y_0/X_0}\oplus \mathbf n_{Y_0/X_0})
 \otimes\mathscr O_{Y_0^2}
\]
induced by multiplication in \(Y_0\). Consider the vector bundle
\(V=\mathbf V(\mathbf n_{Y_0/X_0})\). By duality, \(P\) gives
\[
 \mathbf V(P):
 V\times_{S_0}V\times_{S_0}Y_0\times_{S_0}Y_0
 \longrightarrow V\times_{S_0}Y_0\times_{S_0}Y_0,
\]
which is expressed set-theoretically by
\[
 \mathbf V(P)(u,v,a,b)=(u+\operatorname{Ad}(a)v,ab,b).
\]
\footnote{Editor's note: The entries \(a,b\) have been replaced by
\(ab,b\) to show that \(\mathbf V(P)\) is obtained, by pullback to
\(Y_0^2\), from the multiplication morphism
\(V_{Y_0}\times_{S_0}V_{Y_0}\to V_{Y_0}\).}
This is proved exactly as is the corresponding fact for Lie algebras,
that is, for the module \(\omega^1_{Y_0/S_0}\). First observe that,
by functoriality in \(Y_0\), \(V\) is endowed with a group structure in
the category of vector bundles over \(S_0\). By the lemma already used
for Lie algebras (\SGARefLink{sga3:ii:lemma:3-10}{Expos\'e II, 3.10}), this structure coincides with the
group structure underlying its \(\mathscr O_S\)-module structure. Next,
\[
 \mathbf V\bigl(\mathbf n_{Y_0/X_0}
   \otimes_{\mathscr O_{S_0}}\mathscr O_{Y_0}\bigr)
 =\mathbf V(\mathscr N_{Y_0/X_0})
\]
is itself endowed with an \(S_0\)-group structure, namely the
semidirect product of the structure on \(V\) by that on \(Y_0\). It
remains only to identify the operations of \(Y_0\) on \(V\) to obtain
the desired formula.

% Source page: Exp3 local p. 65; combined p. 175; printed pp. 145--146.
We now compute the two remaining terms. Consider first
\(a_1((0,DY))\). It is computed by the following diagram, in
which \(\mathbf n\) denotes \(\mathbf n_{Y_0/X_0}\):

% Raster stand-in for Loop 1; replace by a native TeX diagram in Loop 2.
% Authority: Exp3 local p. 65; combined reader p. 175; printed p. 146.
\begin{center}
 \SGAThreeNativeInput[0.51\textwidth]
   {native_diagrams/exp3/diagram_057_exp3_p065_module_pullback.tex}
\end{center}

Regard the vector bundles defined by these modules as schemes over
\(S_0\), and take points with values in an arbitrary scheme. If
\((u,x,y,z)\) denotes a point of
\(\mathbf V(\mathscr J)\times Y_0^3\), we obtain the following diagram:

% Raster stand-in for Loop 1; replace by a native TeX diagram in Loop 2.
% Authority: Exp3 local p. 65; combined reader p. 175; printed p. 146.
\begin{center}
 \SGAThreeNativeInput[0.78\textwidth]
   {native_diagrams/exp3/diagram_058_exp3_p065_vector_bundle_points.tex}
\end{center}

Thus
\[
 a_1((0,DY))(x,y,z)
 =\operatorname{Ad}(x)DY(y,z),
\]
which is indeed the first term of the coboundary. Likewise,
\(a_2((DY,0))(x,y,z)=DY(x,y)\), and therefore
\[
 \begin{aligned}
 0={}&\operatorname{Ad}(x)DY(y,z)-DY(xy,z)
      +DY(x,yz)-DY(x,y)\\
   ={}&(\partial^2DY)(x,y,z).
 \end{aligned}
\]
\footnote{Editor's note: The signs have been changed to make them
compatible with I, \SGARefLink{sga3:i:subsection:5-1}{5.1}.}

% Source page: Exp3 local p. 66; combined p. 176; printed pp. 146--147.
\medskip
\noindent\SGAInlineRefTarget{sga3:iii:paragraph:4-20}\textbf{4.20. --- Proof of \SGARefLink{sga3:iii:lemma:4-17:item:iii}{4.17(iii)}.}
\footnote{Editor's note: The argument that follows has been expanded
from the original.}
We compare the two values of
\(v=c(X,Y,\mu\circ(i',i'))\) computed in the following two diagrams:

% Raster stand-in for Loop 1; replace by native TeX diagrams in Loop 2.
% Authority: Exp3 local p. 66; combined reader p. 176; printed p. 147.
\begin{center}
 \SGAThreeNativeInput[0.78\textwidth]
   {native_diagrams/exp3/diagram_059_exp3_p066_comparison_diagrams.tex}
\end{center}

Let \(f=\mu\circ(i',i')\). Diagram \((*)\) gives
\begin{equation}
 v=DY'+f_0^*\bigl(c(X,Y,i')\bigr).
 \tag{1}\label{sga3:iii:paragraph:4-20:equation:1}
\end{equation}
Now \(Y'_0=Y_0\), and \(f_0\) is the multiplication
\(Y_0^2\to Y_0\); hence
\begin{equation}
 f_0^*\bigl(c(X,Y,i')\bigr)(x_0,y_0)
 =c(X,Y,i')(x_0y_0).
 \tag{2}\label{sga3:iii:paragraph:4-20:equation:2}
\end{equation}
Put \(c=c(X,Y,i')\). Under the identification
\(N'_0\simeq N_0\oplus N_0\), the element
\[
 c(X\times_S X,Y\times_S Y,(i',i'))
\]
identifies with the pair \((c,c)\). Thus, putting \(h=(i',i')\),
diagram \((\dagger)\) gives
\begin{equation}
 v=h_0^*(DY)+a_{\mu_0}(c,c).
 \tag{3}\label{sga3:iii:paragraph:4-20:equation:3}
\end{equation}
Since \(h_0\) is the identity map on \(Y_0^2\), we have
\(h_0^*(DY)=DY\). Finally, by the preceding
calculation of \(a_{\mu_0}\), for every \(S'\to S\) and every
\(x_0,y_0\in Y_0(S'_0)\),
\begin{equation}
 a_{\mu_0}(c,c)(x_0,y_0)
 =c(x_0)+\operatorname{Ad}(x_0)\bigl(c(y_0)\bigr).
 \tag{4}\label{sga3:iii:paragraph:4-20:equation:4}
\end{equation}
It follows that
\begin{equation}
\begin{aligned}
 (DY'-DY)(x_0,y_0)
  ={}&\operatorname{Ad}(x_0)\bigl(c(X,Y,i')(y_0)\bigr)
      -c(X,Y,i')(x_0y_0)+c(X,Y,i')(x_0)\\
  ={}&\bigl(\partial^1c(X,Y,i')\bigr)(x_0,y_0).
\end{aligned}
\tag{5}\label{sga3:iii:paragraph:4-20:equation:5}
\end{equation}
Since \(c(X,Y,i')=d(Y,Y')\) (cf. \SGARefLink{sga3:iii:proposition:4-8:item:iii}{4.8(iii)}), this proves that
\[
 DY'-DY=\partial^1d(Y,Y').
\]

\medskip
\noindent\SGAInlineRefTarget{sga3:iii:theorem:4-21}\textbf{Theorem 4.21.}
\textit{Let \(S\) be a scheme and let \(\mathscr I\) and
\(\mathscr J\) be two ideals\footnote{Editor's note: The hypothesis
that \(\mathscr I\) be nilpotent, which appears superfluous, has been
removed.} on \(S\), such that \(\mathscr I\supset\mathscr J\) and
\(\mathscr I\mathscr J=0\). Let \(X\) be an \(S\)-group smooth over
\(S\), and let \(Y_{\mathscr J}\) be a sub-\(S_{\mathscr J}\)-group of
\(X_{\mathscr J}\), flat and locally of finite presentation over
\(S_{\mathscr J}\). Consider the \(S_0\)-functor of commutative groups
\(N_0\) defined by}
\[
 \begin{aligned}
 \operatorname{Hom}_{S_0}(T,N_0)
 ={}&\operatorname{Hom}_{\mathscr O_T}
 \left(
  \mathbf n_{Y_0/X_0}\otimes_{\mathscr O_{S_0}}\mathscr O_T,
  \mathscr J\otimes_{\mathscr O_{S_0}}\mathscr O_T
 \right),\\[-0.2em]
 &\hspace{17em}T\in\operatorname{Ob}(\mathbf{Sch}/S_0),
 \end{aligned}
\]
\textit{on which \(Y_0\) acts through the inner automorphisms of
\(X_0\).}
\begin{enumerate}
\item[(i)] \textit{For there to exist a sub-\(S\)-group of \(X\), flat
  over \(S\), which reduces to \(Y_{\mathscr J}\), it is necessary and
  sufficient that the following two conditions hold:}
  \begin{enumerate}
  \item[\((i_1)\)] \textit{There exists a subscheme \(Y\) of \(X\),
    flat over \(S\), lifting \(Y_{\mathscr J}\). This condition is
    automatically satisfied when \(Y_0\) is affine (cf. \SGARefLink{sga3:iii:proposition:4-6-5}{4.6.5}).}
  \item[\((i_2)\)] \textit{A certain canonical obstruction, an element
    of \(H^2(Y_0,N_0)\), vanishes.}
  \end{enumerate}
\item[\SGAInlineRefTarget{sga3:iii:theorem:4-21:item:ii}(ii)] \textit{If the conditions in (i) hold, the set of
  sub-\(S\)-groups \(Y\) of \(X\), flat over \(S\) and reducing to
  \(Y_{\mathscr J}\), is a principal homogeneous set under
  \(Z^1(Y_0,N_0)\).}\footnote{Editor's note: The question whether this
  set, modulo conjugation by the elements \(x\in X(S)\) inducing the
  identity element of \(X(S_{\mathscr J})\), is a principal homogeneous
  set under \(H^1(Y_0,N_0)\), occupies \SGARefLink{sga3:iii:paragraph:4-23}{4.23}--\SGARefLink{sga3:iii:proposition:4-36}{4.36}.}
\end{enumerate}

% Source page: Exp3 local p. 67; combined p. 177; printed pp. 147--148.
Indeed, condition \((i_1)\) is necessary. Suppose that it holds, and let
\(Y\), flat over \(S\), lift \(Y_{\mathscr J}\). We must find a
\(Y'\), also flat over \(S\) and lifting \(Y_{\mathscr J}\), such that
\(DY'=0\),\footnote{Editor's note: The expression
\(\partial DY'\) in the original has been corrected to
\(DY'\).} cf. \SGARefLink{sga3:iii:lemma:4-17:item:i}{4.17(i)}. By \SGARefLink{sga3:iii:lemma:4-17:item:iii}{4.17(iii)}, this amounts to finding
an element \(d(Y',Y)\in C^1(Y_0,N_0)\) such that
\[
 DY=\partial^1d(Y',Y).
\]
\footnote{Editor's note: See the editor's note to \SGARefLink{sga3:iii:lemma:4-17:item:iii}{4.17(iii)}.}
Let \(c\in H^2(Y_0,N_0)\) be the class represented by \(DY\),
which is a cocycle by \SGARefLink{sga3:iii:lemma:4-17:item:ii}{4.17(ii)}. By \SGARefLink{sga3:iii:lemma:4-17:item:iii}{4.17(iii)}, this class does not depend
on the choice of \(Y\), and its vanishing is necessary and sufficient
for the existence of \(d(Y',Y)\) satisfying the preceding equation.
This proves (i).

If \(Y\) has now been chosen so that \(DY=0\), the equation to
be solved becomes \(\partial^1d(Y',Y)=0\), which proves \SGARefLink{sga3:iii:theorem:4-21:item:ii}{(ii)}.

\medskip
\noindent\SGAInlineRefTarget{sga3:iii:remark:4-22}\textbf{Remark 4.22.}
Retain the notation of \SGARefLink{sga3:iii:theorem:4-21}{4.21}. By \SGARefLink{sga3:iii:proposition:4-15}{4.15}, \(Y_0\) is locally a complete
intersection in \(X_0\); hence \(\mathscr N_{Y_0/X_0}\) is a finite
locally free \(\mathscr O_{Y_0}\)-module, and therefore
\[
 \mathbf n_{Y_0/X_0}
 =\varepsilon_0^*(\mathscr N_{Y_0/X_0})
\]
is a finite locally free \(\mathscr O_{S_0}\)-module. Thus, writing
\[
 \mathbf n_{Y_0/X_0}^{\vee}
 =\operatorname{Hom}_{\mathscr O_{Y_0}}
  (\mathbf n_{Y_0/X_0},\mathscr O_{Y_0}),
\]
we have
\[
 \begin{aligned}
 &\operatorname{Hom}_{\mathscr O_T}
 \left(
  \mathbf n_{Y_0/X_0}\otimes_{\mathscr O_{S_0}}\mathscr O_T,
  \mathscr J\otimes_{\mathscr O_{S_0}}\mathscr O_T
 \right)\\
 &\qquad\simeq
 \mathbf n_{Y_0/X_0}^{\vee}
 \otimes_{\mathscr O_{S_0}}\mathscr J
 \otimes_{\mathscr O_{S_0}}\mathscr O_T
 \end{aligned}
\]
for every \(T\to S_0\).\footnote{Editor's note: The preceding argument
has been expanded. It shows that the following isomorphism is valid
without a flatness hypothesis. On the other hand, since \SGARefLink{sga3:iii:paragraph:4-16}{4.16} the
discussion has been restricted to \(S\)-schemes \(f:T\to S\) flat over
\(S\), in order to ensure that the group
\[
 \operatorname{Hom}_{\mathscr O_{T_0}}
 \left(
  \mathbf n_{Y_0/X_0}\otimes_{\mathscr O_{S_0}}
   \mathscr O_{T_0},
  \mathscr J\mathscr O_T
 \right),
\]
in which the obstruction \(c(X,Y,f)\) lies, coincides with
\(N_0(T_0)\); see the end of \SGARefLink{sga3:iii:paragraph:4-16}{4.16}.}
Consequently, the \(S_0\)-functor \(N_0\) is isomorphic to
\[
 \mathbf W\left(
  \mathbf n_{Y_0/X_0}^{\vee}
  \otimes_{\mathscr O_{S_0}}\mathscr J
 \right)
 \simeq
 \mathbf W\left(
  \operatorname{Hom}_{\mathscr O_{S_0}}
  (\mathbf n_{Y_0/X_0},\mathscr J)
 \right).
\]
It follows that there are isomorphisms\footnote{Editor's note: With
the notation of I, \SGARefLink{sga3:i:subsection:5-3}{5.3}, and assuming that \(Y_0\) is affine over
\(S_0\).}
\[
 \begin{aligned}
 H^2(Y_0,N_0)
 &\simeq H^2\!\left(
   Y_0,\operatorname{Hom}_{\mathscr O_{S_0}}
   (\mathbf n_{Y_0/X_0},\mathscr J)\right)\\
 &\simeq H^2\!\left(
   Y_0,\mathbf n_{Y_0/X_0}^{\vee}
   \otimes_{\mathscr O_{S_0}}\mathscr J\right),\\[0.4em]
 Z^1(Y_0,N_0)
 &\simeq Z^1\!\left(
   Y_0,\operatorname{Hom}_{\mathscr O_{S_0}}
   (\mathbf n_{Y_0/X_0},\mathscr J)\right)\\
 &\simeq Z^1\!\left(
   Y_0,\mathbf n_{Y_0/X_0}^{\vee}
   \otimes_{\mathscr O_{S_0}}\mathscr J\right).
 \end{aligned}
\]

% Continue with 4.23 on Exp3 local p. 67 / combined p. 177 /
% printed p. 148.

% END INCLUDED SGA3_English_current_buildable_source_20260823/source/inputs/III/tex/components/29_expose_III_s417_s422_subgroup_obstruction.tex

% BEGIN INCLUDED SGA3_English_current_buildable_source_20260823/source/inputs/III/tex/components/30_expose_III_s423_s430_conjugacy_coboundaries.tex
% SGA 3, Expose III, Loop 1 English reconstruction.
% Scope: 4.23 through the discussion following Corollary 4.30, stopping
% immediately before Corollary 4.31.
% Authority: Polo--Gille Exp3-14oct24.pdf, component-local pp. 67--70
% (combined-reader pp. 177--180; re-edition pp. 167--170; original
% side-pages 148--152). OCR is a locator only. English wording was compared
% with Jacob C. Reinhold's CC BY 4.0 Markdown lineage at commit
% e7a259f3f8608ad3edf9bf6eead3fd504dd2d23e. The Polo--Gille PDF controls
% prose, mathematics, numbering, editorial notes, and the two source-derived
% Loop-1 diagram images included below.

\setcounter{footnote}{120}

% Source page: Exp3 local p. 67; combined p. 177; re-edition p. 167;
% original side-page 148.
\medskip
\noindent\SGAInlineRefTarget{sga3:iii:paragraph:4-23}\textbf{4.23. ---}
Still under the hypotheses of \SGARefLink{sga3:iii:theorem:4-21}{4.21}, we shall now study how the set of
subgroups \(Y\) lifting \(Y_{\mathscr J}\) behaves under conjugation by
sections of \(X\). If \(x\) is a section of \(X\) over \(S\) inducing the
unit section of \(X_{\mathscr J}\), the inner automorphism
\(\operatorname{Int}(x)\) defined by \(x\) sends flat subgroups of \(X\)
lifting \(Y_{\mathscr J}\) to flat subgroups of \(X\) lifting
\(Y_{\mathscr J}\). Under the conditions of \SGARefLink{sga3:iii:theorem:4-21:item:ii}{4.21(ii)}, the set of these
subgroups is a principal homogeneous space under \(Z^1(Y_0,N_0)\). We
shall see that there is then a subgroup
\(\Gamma\) of \(B^1(Y_0,N_0)\)\footnote{Editor's note: \(Z^1(Y_0,N_0)\)
has been replaced by \(B^1(Y_0,N_0)\), since the proof shows that
\(\Gamma\) is a subgroup of \(B^1(Y_0,N_0)\); cf. \SGARefLink{sga3:iii:lemma:4-27}{4.27}--\SGARefLink{sga3:iii:corollary:4-29}{4.29}.} such that
two subgroups of \(X\), flat over \(S\) and lifting \(Y_{\mathscr J}\),
are conjugate by an \(x\in X(S)\) inducing the unit of
\(X(S_{\mathscr J})\) if and only if their \emph{difference} in
\(Z^1(Y_0,N_0)\) belongs to \(\Gamma\). In the best cases we shall show
that \(\Gamma=B^1(Y_0,N_0)\). Thus the set of flat subgroups of \(X\)
lifting \(Y_{\mathscr J}\), modulo conjugation by the elements
\(x\in X(S)\) inducing the unit of \(X(S_{\mathscr J})\), is either empty
or a principal homogeneous space under \(H^1(Y_0,N_0)\) (cf. \SGARefLink{sga3:iii:corollary:4-29}{4.29} and
\SGARefLink{sga3:iii:proposition:4-36}{4.36}).

% Source page: Exp3 local p. 68; combined p. 178; re-edition p. 168;
% original side-pages 148--150.
\medskip
\noindent\SGAInlineRefTarget{sga3:iii:paragraph:4-24}\textbf{4.24. ---}
We retain the notation of \SGARefLink{sga3:iii:theorem:4-21}{4.21}. Let \(Y\) be a flat subgroup of \(X\)
whose reduction is \(Y_{\mathscr J}\). Recall that in \SGARefLink{sga3:iii:remark:0-5}{0.5} we introduced
the functor \(L_0^X\) (respectively \(L_0^Y\)) defined, for a variable
\(S_0\)-scheme \(T\), by the identity
\[
 \operatorname{Hom}_{S_0}(T,L_0^X)
 =\operatorname{Hom}_{\mathscr O_T}
 \left(
  \omega^1_{X_0/S_0}\otimes_{\mathscr O_{S_0}}\mathscr O_T,
  \mathscr J\otimes_{\mathscr O_{S_0}}\mathscr O_T
 \right),
\]
and similarly with \(Y\) in place of \(X\), as well as the functor
\[
 L'_X=\prod_{S_0/S}L_0^X.
\]
We then have:

\medskip
\noindent\SGAInlineRefTarget{sga3:iii:lemma:4-25}\textbf{Lemma 4.25.}
\textit{There is a canonical exact sequence of
\(Y_0\)-\(\mathscr O_{S_0}\)-modules}
\[
 \tag{+}\label{sga3:iii:lemma:4-25:equation:plus}
 \mathbf n_{Y_0/X_0}
 \xrightarrow{d}
 \omega^1_{X_0/S_0}
 \longrightarrow
 \omega^1_{Y_0/S_0}
 \longrightarrow 0
\]
\textit{having the following properties:}
\begin{enumerate}
\item[\SGAInlineRefTarget{sga3:iii:lemma:4-25:item:i}(i)] \textit{After pullback to \(Y_0\), \(d\) gives the morphism
  \(\overline D_0\) of \SGARefLink{sga3:iii:proposition:4-8:item:iii}{4.8(iii)}.}
\item[\SGAInlineRefTarget{sga3:iii:lemma:4-25:item:ii}(ii)] \textit{If \(X_0\) and \(Y_0\) are smooth over \(S_0\), then
  \(d\) is injective. Since the two modules of differentials are then
  locally free of finite type, the same is true of
  \(\mathbf n_{Y_0/X_0}\), and the sequence is locally split.}
\end{enumerate}

\textit{Proof.}\footnote{Editor's note: The proof has been expanded to
take account of the additions made in Expos\'e I, Section 6.8.}
Let \(\pi_0:Y_0\to S_0\) denote the structure morphism. By SGA 1 II,
formula (4.3) (see also EGA IV\textsubscript{4}, 16.4.21), there is a
canonical exact sequence of \(\mathscr O_{Y_0}\)-modules
\[
 \tag{\(\dagger\)}\label{sga3:iii:lemma:4-25:equation:dagger}
 \mathscr N_{Y_0/X_0}
 \xrightarrow{\overline D_0}
 \Omega^1_{X_0/S_0}\otimes_{\mathscr O_{X_0}}\mathscr O_{Y_0}
 \longrightarrow
 \Omega^1_{Y_0/S_0}
 \longrightarrow 0.
\]
Since this sequence consists of
\((Y_0\mathbin{\times_S}Y_0)\)-equivariant modules and morphisms, its
pullback \((+)\) by \(\varepsilon_0^*\) is an exact sequence of
\(Y_0\)-\(\mathscr O_{S_0}\)-modules, and \((\dagger)\) is the pullback
of \((+)\) by \(\pi_0^*\) (cf. \SGARefLink{sga3:i:subsection:6-8}{Expos\'e I, Section 6.8}). This proves \SGARefLink{sga3:iii:lemma:4-25:item:i}{(i)}.

Suppose in addition that \(X_0\) and \(Y_0\) are smooth over \(S_0\).
Then, by SGA 1 II, 4.10 (see also EGA IV\textsubscript{4}, 17.2.3(i)
and 17.2.5), \(\overline D_0\) is injective and the sequence
\((\dagger)\) consists of \(\mathscr O_{Y_0}\)-modules locally free of
finite type (and is therefore locally split). By the equivalence of
categories I, \SGARefLink{sga3:i:proposition:6-8-1}{6.8.1}, \(d\) is also injective, and consequently the
sequence \((+)\) has the asserted properties.

% Source page: Exp3 local p. 69; combined p. 179; re-edition p. 169;
% original side-pages 150--151.
\medskip
\noindent\SGAInlineRefTarget{sga3:iii:paragraph:4-26}\textbf{4.26. ---}\footnote{Editor's note: The original has
been expanded to make clear that this is an exact sequence of
\(Y_0\)-\(\mathscr O_{S_0}\)-modules.}
For every \(S_0\)-scheme \(f:T\to S_0\), the sequence \((+)\) gives an
exact sequence of \(Y_0(T)\)-\(\mathscr O(T)\)-modules
\[
\begin{aligned}
0&\longrightarrow
 \operatorname{Hom}_{\mathscr O_T}
 \left(f^*\omega^1_{Y_0/S_0},f^*\mathscr J\right)
 \longrightarrow
 \operatorname{Hom}_{\mathscr O_T}
 \left(f^*\omega^1_{X_0/S_0},f^*\mathscr J\right)\\
&\hspace{10em}\longrightarrow
 \operatorname{Hom}_{\mathscr O_T}
 \left(f^*\mathbf n_{Y_0/X_0},f^*\mathscr J\right).
\end{aligned}
\]
Thus there is an exact sequence of
\(Y_0\)-\(\mathscr O_{S_0}\)-modules
\[
 \tag{4.26.1}\label{sga3:iii:paragraph:4-26:equation:4-26-1}
 0\longrightarrow L_0^Y\longrightarrow L_0^X
 \xrightarrow{d}N_0.
\]
It gives an exact sequence of complexes of abelian groups
\[
 0\longrightarrow C^*(Y_0,L_0^Y)
 \longrightarrow C^*(Y_0,L_0^X)
 \xrightarrow{d^*}C^*(Y_0,N_0),
\]
and, in particular, the following commutative diagram with exact rows:

% Raster stand-in for Loop 1; replace by a native TeX diagram in Loop 2.
% Authority: Exp3 local p. 69; combined p. 179; re-edition p. 169;
% original side-page 150.
\begin{center}
 \SGAThreeNativeInput[0.88\textwidth]
   {native_diagrams/exp3/diagram_060_exp3_p069_exact_complex_rows.tex}
\end{center}

Observe that \(C^0(Y_0,L_0^Y)\) (respectively \(C^0(Y_0,L_0^X)\)) is
just
\[
 \operatorname{Hom}_{S_0}(S_0,L_0^Y)
 =\operatorname{Hom}_{S}(S,L'_Y)
\]
(respectively the analogous group with \(X\) in place of \(Y\)); by
\SGARefLink{sga3:iii:corollary:0-9}{0.9}, it is the group of sections of \(Y\) (respectively \(X\)) over
\(S\) inducing the unit section of \(X_{\mathscr J}\). Also,
\(d^1\) is precisely the morphism \(v_{i_{Y_0}}\) of \SGARefLink{sga3:iii:proposition:4-8:item:iii}{4.8(iii)}, where
\(i_{Y_0}:Y_0\to X_0\) is the canonical immersion.\footnote{Editor's
note: This follows from the definition of
\(d:\mathbf n_{Y_0/X_0}\to\Omega^1_{X_0/S_0}\) (cf. \SGARefLink{sga3:iii:lemma:4-25}{4.25}) and from the
definition of \(v_{i_{Y_0}}\) (cf. \SGARefLink{sga3:iii:proposition:4-8}{4.8}).}

\medskip
\noindent\SGAInlineRefTarget{sga3:iii:lemma:4-27}\textbf{Lemma 4.27.}
\textit{Under the conditions of \SGARefLink{sga3:iii:theorem:4-21}{4.21} for \(S\), \(\mathscr I\),
\(\mathscr J\), and \(X\), let \(Y\) be a subgroup of \(X\), flat over
\(S\) and lifting \(Y_{\mathscr J}\). Let \(i:Y\hookrightarrow X\)
denote the canonical immersion.}\footnote{Editor's note: Part \SGARefLink{sga3:iii:lemma:4-27:item:i}{(i)} below
has been added; it will be useful in \SGARefLink{sga3:iii:proposition:4-35-1}{4.35.1} and then in \SGARefLink{sga3:iii:paragraph:4-38:item:4}{4.38(4)} and \SGARefLink{sga3:iii:paragraph:4-20:equation:5}{(5)}.}
\begin{enumerate}
\item[\SGAInlineRefTarget{sga3:iii:lemma:4-27:item:i}(i)] \textit{Let \(i':Y\to X\) be a morphism of \(S\)-schemes
  lifting \(i_0\) (so that \(i'\) is also an immersion), let
  \(Y'=i'(Y)\), and let \(d(i,i')\) be the element of
  \(C^1(Y_0,L_0^X)\) such that \(i'=d(i,i')\cdot i\) (cf. \SGARefLink{sga3:iii:proposition:1-2-4}{1.2.4}).
  Then the deviation \(d(Y,Y')\in C^1(Y_0,N_0)\) (cf. \SGARefLink{sga3:iii:remark:4-5-1}{4.5.1}) is given
  by}
  \[
   d(Y,Y')=d^1\bigl(d(i,i')\bigr).
  \]
\item[\SGAInlineRefTarget{sga3:iii:lemma:4-27:item:ii}(ii)] \textit{Let \(x\in C^0(Y_0,L_0^X)\) be a section of \(X\)
  over \(S\) inducing the unit section of \(X_{\mathscr J}\) over
  \(S_{\mathscr J}\). Then the deviation
  \(d(Y,\operatorname{Int}(x)Y)\in C^1(Y_0,N_0)\) (cf. \SGARefLink{sga3:iii:remark:4-5-1}{4.5.1}) is given
  by}
  \[
   -d\bigl(Y,\operatorname{Int}(x)Y\bigr)
   =d^1\partial x=\partial d^0x.
  \]
\end{enumerate}

% Source page: Exp3 local p. 70; combined p. 180; re-edition p. 170;
% original side-page 152.
Indeed, \(Y'\) is the image of \(Y\) under the composite morphism
\footnote{Editor's note: Recall that
\(L'_X=\prod_{S_0/S}L_0^X\).}
\[
 Y\xrightarrow{(d(i,i'),i)}L'_X\mathbin{\times_S}X
 \longrightarrow X,
\]
which is denoted \(d(i,i')\cdot i\) in \SGARefLink{sga3:iii:proposition:4-8:item:iii}{4.8(iii)}. By loc. cit. and the
equality \(v_{i_0}=d^1\), we therefore have
\[
 \begin{aligned}
 &c\bigl(X,Y',d(i,i')\cdot i\bigr)-c(X,Y',i)\\
 &\qquad=v_{i_0}\bigl(d(i,i')\bigr)
 =d^1\bigl(d(i,i')\bigr).
 \end{aligned}
\]
But \(d(i,i')\cdot i=i'\) factors through \(Y'\) by definition, so the
first term is zero. Moreover, by \SGARefLink{sga3:iii:proposition:4-8:item:ii}{4.8(ii)},
\[
 c(X,Y',i)=d(Y',Y)=-d(Y,Y').
\]
Hence \(d(Y,Y')=d^1(d(i,i'))\), which proves \SGARefLink{sga3:iii:lemma:4-27:item:i}{(i)}.

Now let \(x\) be as in \SGARefLink{sga3:iii:lemma:4-27:item:ii}{(ii)}. From the formula
\[
 xyx^{-1}=xyx^{-1}y^{-1}y
 =(x-\operatorname{Ad}(y)x)y=(-\partial x)(y)\cdot y,
\]
we see that \(Y'\) is the image of \(Y\) under the immersion
\(i'=(-\partial x)\cdot i_Y\). Part \SGARefLink{sga3:iii:lemma:4-27:item:i}{(i)} therefore gives
\[
 -d\bigl(Y,\operatorname{Int}(x)Y\bigr)
 =d^1\partial x=\partial d^0x.
\]

\medskip
\noindent\SGAInlineRefTarget{sga3:iii:corollary:4-28}\textbf{Corollary 4.28.}
\textit{Two subgroups \(Y\) and \(Y'\) of \(X\), flat over \(S\) and
lifting \(Y_{\mathscr J}\), are conjugate by a section of \(X\) over
\(S\) inducing the unit section of \(X_{\mathscr J}\) if and only if}
\[
 d(Y,Y')\in
 \partial d^0 C^0(Y_0,L_0^X)
 \subset \partial C^0(Y_0,N_0)
 =B^1(Y_0,N_0).
\]
\noindent\emph{Source note.} The French re-edition prints \(G\) here;
the ambient group throughout this section is \(X\).

\medskip
\noindent\SGAInlineRefTarget{sga3:iii:corollary:4-29}\textbf{Corollary 4.29.}
\textit{If \(d^0\) is surjective, then \(Y\) and \(Y'\) as above are
conjugate by a section of \(X\) over \(S\) inducing the unit section of
\(X_{\mathscr J}\) if and only if \(d(Y,Y')\in B^1(Y_0,N_0)\).}

\medskip
\noindent\SGAInlineRefTarget{sga3:iii:corollary:4-30}\textbf{Corollary 4.30.}
\textit{Let \(Y\) be as in \SGARefLink{sga3:iii:lemma:4-27}{4.27}. The set of conjugates of \(Y\) by
sections of \(X\) over \(S\) inducing the unit section of
\(X_{\mathscr J}\) is isomorphic to}
\[
 d^1\partial C^0(Y_0,L_0^X)
 =C^0(Y_0,L_0^X)\big/\ker(d^1\partial).
\]

Observe now that
\(C^0(Y_0,L_0^X)/\ker(d^1\partial)\) is computed using only the left-hand
square of the commutative diagram in \SGARefLink{sga3:iii:paragraph:4-26}{4.26}. In particular, it can also
be computed in any diagram of the same type having the same left-hand
square. Consider the functor \(L_0^X/L_0^Y\) over \(S_0\), defined by
\[
 \operatorname{Hom}_{S_0}(T,L_0^X/L_0^Y)
 =\operatorname{Hom}_{S_0}(T,L_0^X)
  \big/\operatorname{Hom}_{S_0}(T,L_0^Y).
\]
There is a commutative diagram

% Raster stand-in for Loop 1; replace by a native TeX diagram in Loop 2.
% Authority: Exp3 local p. 70; combined p. 180; re-edition p. 170;
% original side-page 152.
\begin{center}
 \SGAThreeNativeInput[0.91\textwidth]
   {native_diagrams/exp3/diagram_061_exp3_p070_quotient_exact_complex_rows.tex}
\end{center}

whence, by the preceding observation:

% Continue with Corollary 4.31 on Exp3 local p. 71 / combined p. 181 /
% re-edition p. 171 / original side-page 153.

% END INCLUDED SGA3_English_current_buildable_source_20260823/source/inputs/III/tex/components/30_expose_III_s423_s430_conjugacy_coboundaries.tex

% BEGIN INCLUDED SGA3_English_current_buildable_source_20260823/source/inputs/III/tex/components/31_expose_III_s431_s4351_conjugacy_smooth.tex
% SGA 3, Expose III, Loop 1 English reconstruction.
% Scope: Corollary 4.31 through Proposition 4.35.1, stopping immediately
% before Proposition 4.36.
% Authority: Polo--Gille Exp3-14oct24.pdf, component-local pp. 71--73
% (combined-reader pp. 181--183; re-edition pp. 171--173; original
% side-pages 153--155). OCR is a locator only. English wording was compared
% with Jacob C. Reinhold's CC BY 4.0 Markdown lineage at commit
% e7a259f3f8608ad3edf9bf6eead3fd504dd2d23e. The Polo--Gille PDF controls
% prose, mathematics, numbering, editorial notes, and the source-derived
% Loop-1 diagram image included below.

\setcounter{footnote}{126}

% Source page: Exp3 local p. 71; combined p. 181; re-edition p. 171;
% original side-page 153.
\medskip
\noindent\SGAInlineRefTarget{sga3:iii:corollary:4-31}\textbf{Corollary 4.31.}
\textit{Let \(Y\) be as in \SGARefLink{sga3:iii:lemma:4-27}{4.27}. The set of conjugates of \(Y\) by
sections of \(X\) over \(S\) inducing the unit section of
\(X_{\mathscr J}\) is isomorphic to}
\[
 E=\partial C^0(Y_0,L_0^X/L_0^Y)
 =C^0(Y_0,L_0^X/L_0^Y)\big/H^0(Y_0,L_0^X/L_0^Y).
\]

\medskip
\noindent\SGAInlineRefTarget{sga3:iii:corollary:4-32}\textbf{Corollary 4.32.}
\textit{Suppose in addition that \(S_0\) is affine and that
\(\omega^1_{Y_0/S_0}\) is finite locally free.\footnote{Editor's note:
\(\operatorname{Lie}(Y_0/S_0)\) has been corrected to
\(\omega^1_{Y_0/S_0}\).} If}
\[
 \mathscr F_0=
 \bigl(\operatorname{Lie}(X_0/S_0)/
       \operatorname{Lie}(Y_0/S_0)\bigr)
 \otimes_{\mathscr O_{S_0}}\mathscr J,
\]
\textit{then
\(E=\Gamma(S_0,\mathscr F_0)/H^0(Y_0,\mathscr F_0)\).}

Indeed,\footnote{Editor's note: The original has been expanded in what
follows.} since \(\omega^1_{Y_0/S_0}\) is finite locally free, as is
\(\omega^1_{X_0/S_0}\) (because \(X\) is assumed smooth over \(S\)),
\SGARefLink{sga3:iii:remark:0-6}{0.6} gives
\[
 L_0^Y=\mathbf W\bigl(
   \operatorname{Lie}(Y_0/S_0)\otimes_{\mathscr O_{S_0}}\mathscr J
 \bigr),
 \qquad
 L_0^X=\mathbf W\bigl(
   \operatorname{Lie}(X_0/S_0)\otimes_{\mathscr O_{S_0}}\mathscr J
 \bigr).
\]
On the other hand, by \SGARefLink{sga3:iii:lemma:4-25}{4.25} there is an exact sequence of
\(Y_0\)-\(\mathscr O_{S_0}\)-modules
\[
 0\longrightarrow \mathscr K\longrightarrow \omega^1_{X_0/S_0}
 \xrightarrow{\varphi}\omega^1_{Y_0/S_0}\longrightarrow0,
 \qquad \mathscr K=\ker(\varphi).
\]
Since \(\omega^1_{Y_0/S_0}\) and \(\omega^1_{X_0/S_0}\) are finite
locally free, there is a locally split exact sequence
\[
 0\longrightarrow
 \operatorname{Lie}(Y_0/S_0)\otimes_{\mathscr O_{S_0}}\mathscr J
 \longrightarrow
 \operatorname{Lie}(X_0/S_0)\otimes_{\mathscr O_{S_0}}\mathscr J
 \longrightarrow \mathscr F_0\longrightarrow0.
\]
It follows that there is an exact sequence of
\(Y_0\)-\(\mathscr O_{S_0}\)-modules
\[
 0\longrightarrow L_0^Y\longrightarrow L_0^X
 \longrightarrow\mathbf W(\mathscr F_0).
\]
By the argument used to prove \SGARefLink{sga3:iii:corollary:4-31}{4.31}, we may compute \(E\) as the image
of the composite map
\[
 C^0(Y_0,L_0^X)\xrightarrow{\pi}
 C^0(Y_0,\mathbf W(\mathscr F_0))\xrightarrow{\partial}
 C^1(Y_0,\mathbf W(\mathscr F_0)).
\]
Now \(\pi\) is the map
\[
 \Gamma\bigl(S_0,
   \operatorname{Lie}(X_0/S_0)\otimes_{\mathscr O_{S_0}}\mathscr J
 \bigr)
 \longrightarrow\Gamma(S_0,\mathscr F_0).
\]
Since \(S_0\) is affine, \(\pi\) is surjective, and this proves the
assertion.

\medskip
\noindent\SGAInlineRefTarget{sga3:iii:corollary:4-33}\textbf{Corollary 4.33.}
\textit{Let \(S,\mathscr I,\mathscr J\), and \(X\) be as in \SGARefLink{sga3:iii:theorem:4-21}{4.21},
and let \(Y\) be a diagonalizable subgroup of \(X\). Suppose that
\(\omega^1_{Y_0/S_0}\) is finite locally free and that \(S_0\) is
affine.\footnote{Editor's note: The hypothesis on
\(\omega^1_{Y_0/S_0}\) has been added, and the hypothesis ``\(S\)
affine'' has been replaced by ``\(S_0\) affine.''} The set of
subgroups of \(X\) conjugate to \(Y\) by a section of \(X\) over
\(S\) inducing the unit section of \(X_{\mathscr J}\) is isomorphic
to}
\[
 E=\Gamma\!\left(
 S_0,
 \frac{\operatorname{Lie}(X_0/S_0)}
      {\operatorname{Lie}(X_0/S_0)^{\operatorname{ad}(Y_0)}}
 \otimes_{\Gamma(S_0,\mathscr O_{S_0})}\Gamma(S_0,\mathscr J)
 \right),
\]
\textit{that is, to
\(L_0^X(Y_0)/H^0(Y_0,L_0^X)\).}\footnote{Editor's note: What follows
has been added; cf. \SGARefLink{sga3:iii:corollary:4-34}{4.34}.}

Indeed, by I, \SGARefLink{sga3:i:proposition:4-7-3}{4.7.3} (cf. \SGARefLink{sga3:iii:corollary:2-13}{2.13}), write
\[
 \operatorname{Lie}(X_0/S_0)
 =\operatorname{Lie}(X_0/S_0)^{\operatorname{ad}(Y_0)}
  \oplus\mathscr R.
\]
Since \(Y_0\) is commutative,
\[
 \operatorname{Lie}(Y_0/S_0)
 \subset\operatorname{Lie}(X_0/S_0)^{\operatorname{ad}(Y_0)},
\]
and hence
\[
 \begin{aligned}
 \mathscr F_0={}&
 \left(
  \frac{\operatorname{Lie}(X_0/S_0)^{\operatorname{ad}(Y_0)}}
       {\operatorname{Lie}(Y_0/S_0)}
 \right)\otimes_{\mathscr O_{S_0}}\mathscr J
 \ \oplus\ \mathscr R\otimes_{\mathscr O_{S_0}}\mathscr J,\\
 \mathscr F_0^{\operatorname{ad}(Y_0)}={}&
 \left(
  \frac{\operatorname{Lie}(X_0/S_0)^{\operatorname{ad}(Y_0)}}
       {\operatorname{Lie}(Y_0/S_0)}
 \right)\otimes_{\mathscr O_{S_0}}\mathscr J.
 \end{aligned}
\]
By \SGARefLink{sga3:iii:corollary:4-32}{4.32}, therefore,
\(E\simeq
\Gamma(S_0,\mathscr R\otimes_{\mathscr O_{S_0}}\mathscr J)\).
Returning to the definition of \(\mathscr R\) completes the proof.

% Source page: Exp3 local p. 72; combined p. 182; re-edition p. 172;
% original side-page 154.
\medskip
\noindent\SGAInlineRefTarget{sga3:iii:corollary:4-34}\textbf{Corollary 4.34.}
\textit{Let \(S,\mathscr I,\mathscr J\), and \(X\) be as in \SGARefLink{sga3:iii:theorem:4-21}{4.21},
and let \(Y\) be a diagonalizable subgroup of \(X\). Suppose that
\(\omega^1_{Y_0/S_0}\) is finite locally free and that \(S_0\) is
affine.\footnote{Editor's note: Cf. Editor's Note 129.} If
\(x\in X(S)\) induces the unit section of \(X_{\mathscr J}\) and
normalizes \(Y\), then it centralizes \(Y\).}

This follows immediately by comparing the preceding corollary with
\SGARefLink{sga3:iii:corollary:2-14}{2.14}. Indeed, \SGARefLink{sga3:iii:corollary:4-33}{4.33} shows that the elements of \(C^0(Y_0,L_0^X)\)
which globally preserve \(Y\) are the elements of
\(H^0(Y_0,L_0^X)\), and \SGARefLink{sga3:iii:corollary:2-14}{2.14} shows that these are precisely the
elements which act trivially on the canonical immersion
\(Y\longrightarrow X\).

\medskip
\noindent\SGAInlineRefTarget{sga3:iii:paragraph:4-35}\textbf{4.35. ---}
Return to the general situation of \SGARefLink{sga3:iii:theorem:4-21}{4.21} and suppose that
\(Y_{\mathscr J}\) is smooth over \(S_{\mathscr J}\). By \SGARefLink{sga3:iii:lemma:4-25:item:ii}{4.25(ii)},
there is then an exact sequence of
\(Y_0\)-\(\mathscr O_{S_0}\)-modules
\[
 \tag{*}\label{sga3:iii:paragraph:4-35:equation:star}
 0\longrightarrow\operatorname{Lie}(Y_0/S_0)
 \longrightarrow\operatorname{Lie}(X_0/S_0)
 \longrightarrow\mathbf n_{Y_0/X_0}^{\vee}\longrightarrow0,
\]
and these are finite locally free \(\mathscr O_{S_0}\)-modules.

On the other hand, by SGA~1, II, \SGARefLink{sga3:ii:definition:4-10}{4.10}, every subscheme \(Y\) of
\(X\) lifting \(Y_{\mathscr J}\) and flat over \(S\) is smooth over
\(S\).\footnote{Editor's note: What follows, together with
\SGARefLink{sga3:iii:proposition:4-35-1}{Proposition 4.35.1}, has been added; it is implicit in the original.
Cf. 4.38(5).} Suppose in addition that \(S_0\) and
\(Y_{\mathscr J}\) are affine. Since
\(\mathbf n_{Y_0/X_0}^{\vee}\) is a locally free
\(\mathscr O_{S_0}\)-module, the sequence \((*)\) remains exact after
tensoring with \(\mathscr J\) over \(\mathscr O_{S_0}\), and then
pulling back to \(Y_0^n\). Since the \(Y_0^n\) are affine, this gives
an exact sequence of complexes of abelian groups
\[
 0\longrightarrow C^*(Y_0,L_0^Y)
 \longrightarrow C^*(Y_0,L_0^X)
 \xrightarrow{d^*}C^*(Y_0,N_0)\longrightarrow0,
\]
and, in particular, a commutative diagram with exact rows:

% Raster stand-in for Loop 1; replace by a native TeX diagram in Loop 2.
% Authority: Exp3 local p. 72; combined reader p. 182; re-edition p. 172;
% original side-page 154.
\begin{center}
 \SGAThreeNativeInput[0.94\textwidth]
   {native_diagrams/exp3/diagram_062_exp3_p072_435_cochain_diagram.tex}
\end{center}

Let now \(Y\) and \(Y'\) be two subgroups of \(X\) lifting
\(Y_{\mathscr J}\) and flat, hence smooth, over \(S\). Since
\(Y_{\mathscr J}\) is affine, \SGARefLink{sga3:iii:proposition:0-15}{0.15} shows that \(Y\) and \(Y'\) are
isomorphic as schemes lifting \(Y_{\mathscr J}\); that is, there is
an isomorphism of \(S\)-schemes
\[
 f:Y\xrightarrow{\sim}Y'
\]

% Source page: Exp3 local p. 73; combined p. 183; re-edition p. 173;
% original side-page 155.
inducing the identity on \(Y_{\mathscr J}\). On the one hand, by
\SGARefLink{sga3:iii:proposition:1-2-4}{1.2.4}, \(f\) defines an element \(a\in C^1(Y_0,L_0^X)\) such that
\[
 f(y)=a(y_0)y
\]
for every \(y\in Y(S')\), \(S'\to S\), and by \SGARefLink{sga3:iii:lemma:4-27:item:i}{4.27(i)},
\[
 d^1(a)=d(Y,Y').
\]
Moreover, because \(Y\) and \(Y'\) are subgroups of \(X\), the last
element belongs to \(Z^1(Y_0,N_0)\) (cf. \SGARefLink{sga3:iii:theorem:4-21}{4.21}). Thus
\(\partial a\) is an element of \(Z^2(Y_0,L_0^Y)\), whose image
\(\overline{\partial a}\) in \(H^2(Y_0,L_0^Y)\) depends only on the
class \(\overline{d(Y,Y')}\in H^1(Y_0,N_0)\). This is precisely the
definition of the connecting map
\[
 \partial^1:H^1(Y_0,N_0)\longrightarrow H^2(Y_0,L_0^Y),
\]
so that
\[
 \partial^1\bigl(\overline{d(Y,Y')}\bigr)
 =\overline{\partial a}.
\]

On the other hand, transport the group structure of \(Y'\) by \(f\),
and let \(Y_1\) be the resulting group, whose underlying scheme is
therefore \(Y\). Thus its group law \(\mu_1\) is defined, for every
\(S'\to S\) and \(x,y\in Y(S')\), by
\[
 \mu_1(x,y)=f^{-1}\bigl(f(x)f(y)\bigr).
\]
By \SGARefLink{sga3:iii:remark:3-5-1}{3.5.1}, \(Y_1\) defines a cocycle
\[
 \delta(Y,Y_1)\in
 Z^2\bigl(Y_0,\operatorname{Lie}(Y_0/S_0)\bigr)
\]
such that, for every \(S'\to S\) and \(x,y\in Y(S')\),
\[
 \delta(Y,Y_1)(x_0,y_0)xy
 =\mu_1(x,y)=f^{-1}\bigl(f(x)f(y)\bigr).
\]
Put \(b=\delta(Y,Y_1)\). For every \(S'\to S\) and
\(x,y\in Y(S')\), we have
\((b(x_0,y_0)xy)_0=x_0y_0\). Hence
\(f(b(x_0,y_0)xy)\) is equal, on the one hand, to
\[
 a(x_0y_0)b(x_0,y_0)xy,
\]
and, on the other hand, to
\[
 f(x)f(y)=a(x_0)x\,a(y_0)y
 =a(x_0)\operatorname{Ad}(x_0)\bigl(a(y_0)\bigr)xy.
\]
Comparing these two expressions, we find that \(b(x_0,y_0)\) equals
\[
 \begin{aligned}
 &a(x_0y_0)^{-1}a(x_0)
   \operatorname{Ad}(x_0)\bigl(a(y_0)\bigr)\\
 &\qquad={}
 \operatorname{Ad}(x_0)\bigl(a(y_0)\bigr)-a(x_0y_0)+a(x_0)
 =(\partial a)(x_0,y_0).
 \end{aligned}
\]
Thus \(\delta(Y,Y_1)=\partial a\), and we have proved:

\medskip
\noindent\SGAInlineRefTarget{sga3:iii:proposition:4-35-1}\textbf{Proposition 4.35.1.}\footnotemark[132]
\textit{Under the hypotheses of \SGARefLink{sga3:iii:theorem:4-21}{4.21}, suppose in addition that
\(S_0\) is affine and that \(Y_{\mathscr J}\) is smooth over
\(S_{\mathscr J}\) and affine. Let \(Y,Y'\) be two subgroups of
\(X\) lifting \(Y_{\mathscr J}\) and flat (hence smooth) over
\(S\). Let \(f:Y\xrightarrow{\sim}Y'\) be an isomorphism of
\(S\)-schemes inducing the identity on \(Y_{\mathscr J}\), and let
\(Y_1\) denote the group obtained by transporting the group structure
of \(Y'\) by \(f\). Then}
\[
 \partial^1\bigl(\overline{d(Y,Y')}\bigr)
 =\overline{\delta(Y,Y_1)}.
\]

% Continue with Proposition 4.36 on Exp3 local p. 73 / combined p. 183 /
% re-edition p. 173 / original side-page 155.

% END INCLUDED SGA3_English_current_buildable_source_20260823/source/inputs/III/tex/components/31_expose_III_s431_s4351_conjugacy_smooth.tex

% BEGIN INCLUDED SGA3_English_current_buildable_source_20260823/source/inputs/III/tex/components/32_expose_III_s436_s438_bibliography.tex
% SGA 3, Expose III, Loop 1 English reconstruction.
% Scope: Proposition 4.36 through section 4.38 and the Expose III
% bibliography/end matter.
% Authority: Polo--Gille Exp3-14oct24.pdf, component-local pp. 73--78
% (combined-reader pp. 183--188; re-edition pp. 173--177; original-side
% pp. 155--159). Local p. 78 / combined p. 188 is the trailing blank verso.
% OCR is a locator only. English wording was compared with Jacob C.
% Reinhold's CC BY 4.0 Markdown lineage at commit
% e7a259f3f8608ad3edf9bf6eead3fd504dd2d23e. The Polo--Gille PDF controls
% prose, mathematics, numbering, and editorial notes. There are no genuine
% graphical diagrams in this component; its exact sequences are equations
% reconstructed natively in TeX.

\setcounter{footnote}{132}

% Source page: Exp3 local p. 73; combined p. 183; re-edition p. 173;
% original-side p. 155.
\medskip
\noindent\SGAInlineRefTarget{sga3:iii:proposition:4-36}\textbf{Proposition 4.36.}
\textit{Under the hypotheses of \SGARefLink{sga3:iii:theorem:4-21}{4.21}, suppose in addition that
\(Y_{\mathscr J}\) is smooth over \(S_{\mathscr J}\) and that \(S_0\)
is affine. The set of sub-\(S\)-groups \(Y\) of \(X\), flat (or smooth)
over \(S\), reducing to \(Y_{\mathscr J}\), modulo conjugation by
sections of \(X\) over \(S\) inducing the unit section of
\(X_{\mathscr J}\), is either empty or a principal homogeneous set under
the group}
\[
 H^1\!\left(
 Y_0,
 \bigl[\operatorname{Lie}(X_0/S_0)/\operatorname{Lie}(Y_0/S_0)\bigr]
 \otimes_{\mathscr O_{S_0}}\mathscr J
 \right).
\]

It is enough to verify that \SGARefLink{sga3:iii:corollary:4-29}{Corollary 4.29} applies, that is, that
\[
 d^0:
 \operatorname{Hom}_{\mathscr O_{S_0}}
   (\omega^1_{X_0/S_0},\mathscr J)
 \longrightarrow
 \operatorname{Hom}_{\mathscr O_{S_0}}
   (\mathbf n_{Y_0/X_0},\mathscr J)
\]
is surjective. This follows from the fact that the sequence \((+)\) of
\SGARefLink{sga3:iii:lemma:4-25:item:ii}{4.25(ii)} is split, since \(S_0\) is affine.
\footnote{Editor's note: This also follows from the proof of \SGARefLink{sga3:iii:corollary:4-32}{4.32}.}

Let us finally state a corollary common to \SGARefLink{sga3:iii:theorem:4-21}{4.21} and \SGARefLink{sga3:iii:proposition:4-36}{4.36}. In fact, this
will be the only form in which we shall subsequently use the general
results of this section.
\footnote{Editor's note: For example, \SGARefLink{sga3:iii:corollary:4-37}{4.37} is used in Expos\'e IX to
prove statements \SGARefLink{sga3:ix:theorem:3:2:bis}{3.2 bis} and \SGARefLink{sga3:ix:theorem:3:6:bis}{3.6 bis}.}

% Source page: Exp3 local p. 74; combined p. 184; re-edition p. 174;
% original-side p. 156.
\medskip
\noindent\SGAInlineRefTarget{sga3:iii:corollary:4-37}\textbf{Corollary 4.37.}
\textit{Let \(S\) be a scheme and let \(S_0\) be the closed subscheme
defined by a nilpotent ideal \(\mathscr I\). Let \(X\) be an \(S\)-group
smooth over \(S\), and let \(Y_0\) be a sub-\(S_0\)-group of \(X_0\),
flat over \(S_0\).}
\begin{enumerate}
\item[\SGAInlineRefTarget{sga3:iii:corollary:4-37:item:i}(i)] \textit{If \(S_0\) is affine, \(Y_0\) is smooth over \(S_0\),
and if}
\[
 H^1\!\left(
 Y_0,
 \bigl[\operatorname{Lie}(X_0/S_0)/\operatorname{Lie}(Y_0/S_0)\bigr]
 \otimes_{\mathscr O_{S_0}}
 \mathscr I^{n+1}/\mathscr I^{n+2}
 \right)=0
\]
\textit{for every \(n\geq 0\), then any two sub-\(S\)-groups of \(X\),
flat (or smooth) over \(S\) and reducing to \(Y_0\), are conjugate by a
section of \(X\) over \(S\) inducing the unit section of \(X_0\).}

\item[\SGAInlineRefTarget{sga3:iii:corollary:4-37:item:ii}(ii)] \textit{If \(Y_0\) is affine and of finite presentation,
and if}\footnote{Editor's note: We have replaced
\(\operatorname{Hom}_{\mathscr O_{S_0}}
(\mathbf n_{Y_0/X_0},\mathscr I^{n+1}/\mathscr I^{n+2})\) by
\(\mathbf n_{Y_0/X_0}^{\vee}\otimes_{\mathscr O_{S_0}}
\mathscr I^{n+1}/\mathscr I^{n+2}\), in accordance with \SGARefLink{sga3:iii:remark:4-22}{Remark 4.22}.}
\[
 H^2\!\left(
 Y_0,
 \mathbf n_{Y_0/X_0}^{\vee}
 \otimes_{\mathscr O_{S_0}}
 \mathscr I^{n+1}/\mathscr I^{n+2}
 \right)=0
\]
\textit{for every \(n\geq 0\), then there exists a sub-\(S\)-group of
\(X\), flat over \(S\), reducing to \(Y_0\).}
\end{enumerate}

\medskip
\noindent\SGAInlineRefTarget{sga3:iii:paragraph:4-38}\textbf{4.38.}
It remains to relate the three constructions made in this Expos\'e. To
avoid inessential complications, we shall place ourselves in the
following situation: \(S_0\) is the spectrum of a field \(k\), \(S\) is
the spectrum of the dual numbers \(D(k)\), \(G\) is an \(S\)-group
smooth over \(S\), and \(K\) is a sub-\(S\)-group, smooth over \(S\) and
affine.
\footnote{Editor's note: We have slightly modified the original in what
follows. In particular, we have replaced \(X\) by \(G\) and \(Y\) by
\(K\), and denoted their Lie algebras by \(\mathfrak g_0\) and
\(\mathfrak k_0\). We have also written \(H^i(K_0,{\cdot})\) explicitly
in place of the abbreviation \(H^i({\cdot})\) used in the original.}

Put
\[
 \mathfrak g_0=\operatorname{Lie}G_0
 =\Gamma(S_0,\operatorname{Lie}G_0)
 =\operatorname{Lie}(G_0/S_0)(S_0),
 \qquad
 \mathfrak k_0=\operatorname{Lie}K_0.
\]
We have an exact sequence of \(k\)-vector spaces
\footnote{Editor's note: Equipped with the adjoint action of \(K_0\).}
\[
 0\longrightarrow\mathfrak k_0
 \xrightarrow{\ i\ }\mathfrak g_0
 \xrightarrow{\ d\ }\mathbf n^{\vee}_{K_0/G_0}
 \longrightarrow 0,
\]
which gives rise to an exact cohomology sequence
\[
\begin{aligned}
0&\longrightarrow H^0(K_0,\mathfrak k_0)
 \xrightarrow{\ i^0\ }H^0(K_0,\mathfrak g_0)
 \xrightarrow{\ d^0\ }H^0(K_0,\mathfrak g_0/\mathfrak k_0)\\
&\xrightarrow{\ \partial^0\ }H^1(K_0,\mathfrak k_0)
 \xrightarrow{\ i^1\ }H^1(K_0,\mathfrak g_0)
 \xrightarrow{\ d^1\ }H^1(K_0,\mathbf n^{\vee}_{K_0/G_0})
 \xrightarrow{\ \partial^1\ }H^2(K_0,\mathfrak k_0).
\end{aligned}
\]

These various groups all have a geometric meaning.

\begin{enumerate}
\item[\textbf{a)}]
\(H^0(K_0,\mathfrak k_0)
=\operatorname{Lie}\underline{\operatorname{Centr}}(K_0)\)
\footnote{Editor's note: Since the formation of centralizers and
normalizers commutes with base change (cf. I, \SGARefLink{sga3:i:scholium:2-3-3-1}{2.3.3.1}), we have written
\(\underline{\operatorname{Centr}}(K_0)\) in place of
\(\underline{\operatorname{Centr}}(K)_0\) in the original, and likewise
\(\underline{\operatorname{Centr}}_{G_0}(K_0)\) and
\(\underline{\operatorname{Norm}}_{G_0}(K_0)\) in place of
\(\underline{\operatorname{Centr}}_G(K)_0\) and
\(\underline{\operatorname{Norm}}_G(K)_0\).}
by II, \SGARefLink{sga3:ii:theorem:5-2-3}{5.2.3}.

% Source page: Exp3 local p. 75; combined p. 185; re-edition p. 175;
% original-side pp. 157--158.
\item[\textbf{b)}]
\(H^0(K_0,\mathfrak g_0)
=\operatorname{Lie}\underline{\operatorname{Centr}}_{G_0}(K_0)\)
\footnotemark[138] (idem).

\item[\textbf{c)}]
\(H^0(K_0,\mathfrak g_0/\mathfrak k_0)
=\operatorname{Lie}\underline{\operatorname{Norm}}_{G_0}(K_0)
/\mathfrak k_0\)\footnotemark[138] (idem).

\item[\textbf{d)}]
\[
 H^1(K_0,\mathfrak k_0)
 =\operatorname{Lie}\!\left(
   \underline{\operatorname{Aut}}_{S_0\text{-gr.}}(K_0)
  \right)/\operatorname{Im}(\mathfrak k_0),
\]
where \(\operatorname{Im}(\mathfrak k_0)\) denotes the image of
\(\mathfrak k_0\) under the morphism \(\operatorname{Lie}(\operatorname{Int}_0)\)
deduced from
\[
 \operatorname{Int}_0:K_0\longrightarrow
 \underline{\operatorname{Aut}}_{S_0\text{-gr.}}(K_0).
\]
Indeed, it follows from \SGARefLink{sga3:iii:theorem:2-1:item:ii}{2.1(ii)}, applied with \(Y=X=K\) and
\(f_0=\operatorname{id}_{K_0}\), that \(Z^1(K_0,\mathfrak k_0)\) is the
group of infinitesimal automorphisms of the \(S_0\)-group \(K_0\), and
that \(H^1(K_0,\mathfrak k_0)\) is obtained by quotienting by inner
infinitesimal automorphisms, that is, by the image of \(\mathfrak k_0\).
Moreover, by II, \SGARefLink{sga3:ii:corollary:4-2-2}{4.2.2}, we also have
\[
 Z^1(K_0,\mathfrak k_0)
 =\operatorname{Lie}\!\left(
   \underline{\operatorname{Aut}}_{S_0\text{-gr.}}(K_0)
  \right).
\]
\footnote{Editor's note: This is the Lie algebra
\(\operatorname{Der}_k(\mathfrak k_0)\) of derivations of
\(\mathfrak k_0\); hence \(H^1(K_0,\mathfrak k_0)\) is the quotient of
\(\operatorname{Der}_k(\mathfrak k_0)\) by the inner derivations (that
is, by the image of
\(\operatorname{ad}:\mathfrak k_0\to\operatorname{Der}_k(\mathfrak k_0)\)).}

\item[\textbf{e)}]
By \SGARefLink{sga3:iii:theorem:2-1:item:ii}{2.1(ii)}, \(H^1(K_0,\mathfrak g_0)\) is the group of deviations
between homomorphisms \(K\to G\) extending the canonical immersion
\(i_0:K_0\to G_0\), modulo the deviations obtained through the action
of the inner automorphisms of \(G\) defined by elements of \(G(S)\)
which give the unit of \(G(S_0)\) (that is, elements of
\(\mathfrak g_0\)).

\item[\textbf{f)}]
By 4.36, \(H^1(K_0,\mathbf n^{\vee}_{K_0/G_0})\) is the group of
deviations between subgroups \(K'\) of \(G\) extending \(K_0\) and
flat over \(S\) (hence smooth over \(S\), cf. SGA 1, II, 4.10), modulo
the deviations obtained through the action of the inner automorphisms
of \(G\) constructed as above.

\item[\textbf{g)}]
By 3.5(ii), \(H^2(K_0,\mathfrak k_0)\) is the group of deviations
between group structures on \(K\) extending that of \(K_0\), modulo
the \(S\)-automorphisms of \(K\) inducing the identity on \(K_0\).
\end{enumerate}

We now propose to make explicit the six morphisms of the preceding
exact sequence in terms of the geometric interpretation just given.

\begin{enumerate}
\item[\textbf{1)}]
The morphisms \(i^0\) and \(d^0\) are simply those obtained by passing
to the Lie algebra (and then, for \(d^0\), to the quotient) in the
canonical monomorphisms
\[
 \underline{\operatorname{Centr}}(K_0)
 \longrightarrow
 \underline{\operatorname{Centr}}_{G_0}(K_0)
 \longrightarrow
 \underline{\operatorname{Norm}}_{G_0}(K_0).
\]
This follows immediately from the definition of the identifications
\textbf{a)}, \textbf{b)}, and \textbf{c)}.

\item[\textbf{2)}]
The morphism \(\partial^0\) is constructed as follows. Let
\[
 \bar x\in
 \operatorname{Lie}\underline{\operatorname{Norm}}_{G_0}(K_0)
 /\mathfrak k_0.
\]
Lift it to an element
\[
 x\in\operatorname{Lie}\underline{\operatorname{Norm}}_{G_0}(K_0)
 \subset\underline{\operatorname{Norm}}_G(K)(S).
\]
Then \(\operatorname{Int}(x)\) defines an automorphism of \(K\)
inducing the identity on \(K_0\), hence an element of
\(\operatorname{Lie}\underline{\operatorname{Aut}}_{S_0\text{-gr.}}(K_0)\).
Denote by \(\overline{\operatorname{Int}(x)}\) the image of this element
in
\[
 \operatorname{Lie}\underline{\operatorname{Aut}}_{S_0\text{-gr.}}(K_0)
 /\operatorname{Im}(\mathfrak k_0).
\]
Then
\begin{equation}
 \partial^0(\bar x)
 =-\overline{\operatorname{Int}(x)}
 =\overline{\operatorname{Int}(x^{-1})}.
 \tag{*}\label{sga3:iii:paragraph:4-38:equation:star}
\end{equation}

Indeed, let us compute the element of
\(\operatorname{Lie}\underline{\operatorname{Aut}}_{S_0\text{-gr.}}(K_0)\)
defined by \(\operatorname{Int}(x)\). By definition it corresponds to
an element \(a\in Z^1(K_0,\mathfrak k_0)\) such that
\[
 xyx^{-1}=a(y_0)y,
 \qquad y\in K(S'),\quad S'\longrightarrow S.
\]
But this can also be written
\[
 a(y_0)=xyx^{-1}y^{-1}=x-\operatorname{Ad}(y)x
 =-\partial(x)(y_0),
\]
whence \(a=-\partial(x)\).

% Source page: Exp3 local p. 76; combined p. 186; re-edition p. 176;
% original-side p. 159.
On the other hand,\footnote{Editor's note: We have added the following
sentence.} the image under \(\partial\) of
\[
 x\in\operatorname{Lie}\underline{\operatorname{Norm}}_{G_0}(K_0)
 \subset\mathfrak g_0
\]
is an element of \(Z^1(K_0,\mathfrak k_0)\), whose image
\(\overline{\partial(x)}\) in \(H^1(K_0,\mathfrak k_0)\) depends only
on \(\bar x\). By definition of the connecting map \(\partial^0\),
\[
 \partial^0(\bar x)=\overline{\partial(x)}.
\]
Combined with \(a=-\partial(x)\), this proves \((*)\).

\item[\textbf{3)}]
\footnote{Editor's note: We have expanded the original in what follows,
and in \((**)\) we have corrected \(u\circ i\) to \(i\circ u\).}
Let \(i:K\to G\) denote the canonical immersion. Let
\(\bar u\in H^1(K_0,\mathfrak k_0)\) be the image of an element
\[
 u\in
 \operatorname{Lie}\underline{\operatorname{Aut}}_{S_0\text{-gr.}}(K_0)
 \subset\underline{\operatorname{Aut}}_{S\text{-gr.}}(K).
\]
Then
\begin{equation}
 i^1(\bar u)=\bar d(i,i\circ u),
 \tag{**}\label{sga3:iii:paragraph:4-38:equation:double-star}
\end{equation}
where \(\bar d(i,i\circ u)\) is the class defined in \SGARefLink{sga3:iii:remark:2-1-1}{2.1.1}.

Indeed, \(\bar u\) is the image of an element
\(v\in Z^1(K_0,\mathfrak k_0)\) such that
\(u(y)=v(y_0)y\), and \(i^1(\bar u)\) is the image in
\(H^1(K_0,\mathfrak g_0)\) of the cocycle
\(i\circ v\in Z^1(K_0,\mathfrak g_0)\). Since \(i\) is a group
morphism, the equality \(u(y)=v(y_0)y\) implies
\[
 i(u(y))=i(v(y_0))i(y).
\]
It follows that \(i\circ v=d(i,i\circ u)\), proving \((**)\).

\item[\textbf{4)}]\SGAInlineRefTarget{sga3:iii:paragraph:4-38:item:4}
Let \(i':K\to G\) be a group morphism lifting \(i_0\), let
\(\bar d(i,i')\) be the class defined in \SGARefLink{sga3:iii:remark:2-1-1}{2.1.1}, and let
\[
 d\bigl(K,i'(K)\bigr)
 \in C^1(K_0,\mathbf n^{\vee}_{K_0/G_0})
\]
be the deviation defined in \SGARefLink{sga3:iii:remark:4-5-1}{4.5.1}. By \SGARefLink{sga3:iii:theorem:4-21}{4.21},
\(d(K,i'(K))\) belongs to
\(Z^1(K_0,\mathbf n^{\vee}_{K_0/G_0})\). Denote its image in
\(H^1(K_0,\mathbf n^{\vee}_{K_0/G_0})\) by
\(\bar d(K,i'(K))\). Then, by \SGARefLink{sga3:iii:lemma:4-27:item:i}{4.27(i)},
\begin{equation}
 d^1\bigl(\bar d(i,i')\bigr)=\bar d\bigl(K,i'(K)\bigr).
 \tag{\(\dagger\)}\label{sga3:iii:paragraph:4-38:equation:dagger}
\end{equation}
\par\smallskip
\noindent\emph{Source note.} In these three adjacent occurrences the
French re-edition prints \(\mathbf n_{K_0/X_0}\), without the dual. The
English uses \(\mathbf n^{\vee}_{K_0/G_0}\), as required by the exact
sequence and cohomology sequence established at the start of \SGARefLink{sga3:iii:paragraph:4-38}{4.38}.

\item[\textbf{5)}]\SGAInlineRefTarget{sga3:iii:paragraph:4-38:item:5}
Finally, let \(K'\) be a subgroup of \(G\) lifting \(K_0\) and flat,
hence smooth, over \(S\). We have assumed that \(K_0\) is affine. We
therefore know that \(K\) and \(K'\) are isomorphic as schemes extending
\(K_0\) (cf. \SGARefLink{sga3:iii:proposition:0-15}{0.15}), so there exists an isomorphism of \(S\)-schemes
\[
 f:K\xrightarrow{\ \sim\ }K'
\]
inducing the identity on \(K_0\). Transport the group structure of
\(K'\) by \(f\), and let \(K_1\) be the group thereby obtained (whose
underlying scheme is thus \(K\)); that is, the group law \(\mu_1\) of
\(K_1\) is defined, for every \(S'\to S\) and every
\(x,y\in K(S')\), by
\[
 \mu_1(x,y)=f^{-1}\bigl(f(x)f(y)\bigr).
\]
By \SGARefLink{sga3:iii:remark:3-5-1}{3.5.1},\footnote{Editor's note: We have modified the original in what
follows, taking account of the additions made in \SGARefLink{sga3:iii:remark:3-5-1}{3.5.1} and \SGARefLink{sga3:iii:proposition:4-35-1}{4.35.1}.}
\(K_1\) defines a cocycle
\(\delta(K,K_1)\in Z^2(K_0,\mathfrak k_0)\) such that, for every
\(S'\to S\) and every \(x,y\in K(S')\),
\[
 \delta(K,K_1)(x_0,y_0)xy
 =\mu_1(x,y)=f^{-1}\bigl(f(x)f(y)\bigr).
\]
Then, by \SGARefLink{sga3:iii:proposition:4-35-1}{4.35.1},
\begin{equation}
 \partial^1\bar d(K,K')=\bar\delta(K,K_1).
 \tag{\(\ddagger\)}\label{sga3:iii:paragraph:4-38:equation:ddagger}
\end{equation}
\end{enumerate}

% Source page: Exp3 local p. 77; combined p. 187; re-edition p. 177.
% The bibliography carries no original-side page number in the re-edition.
\section*{Bibliography}
\addcontentsline{toc}{section}{Bibliography}
\footnote{Editor's note: Additional references cited in this Expos\'e.}

\begin{description}
\item[{[BAlg]}] N. Bourbaki, \emph{Alg\`ebre}, Chap. I--III,
  Hermann, 1970.
\item[{[BAC]}] N. Bourbaki, \emph{Alg\`ebre commutative}, Chap. I--IV,
  Masson, 1985.
\item[{[DG70]}] M. Demazure and P. Gabriel,
  \emph{Groupes alg\'ebriques}, Masson \& North-Holland, 1970.
\item[{[Fr64]}] P. Freyd, \emph{Abelian categories}, Harper and Row,
  1964.
\end{description}

% Exp3 local p. 78 / combined p. 188 is blank.
% End of Expose III.

% END INCLUDED SGA3_English_current_buildable_source_20260823/source/inputs/III/tex/components/32_expose_III_s436_s438_bibliography.tex



% END INCLUDED SGA3_English_current_buildable_source_20260823/source/inputs/III/body.tex


\SGAStartExpose{IV}
% BEGIN INCLUDED SGA3_English_current_buildable_source_20260823/source/inputs/IV/body.tex
% Mechanically extracted after \mainmatter from the frozen standalone master.
% BEGIN INCLUDED SGA3_English_current_buildable_source_20260823/source/inputs/IV/tex/components/33_expose_IV_complete.tex
% SGA 3 Loop-1 integration unit: complete Expose IV.
% Source components are preserved byte-for-byte from
% sga3_expose_iv_english_translation_20260723/tex/components.
% BEGIN INCLUDED SGA3_English_current_buildable_source_20260823/source/inputs/IV/tex/components/01_expose_IV_opening_s11_s16.tex
% SGA 3, Expose IV, English reconstruction.
% Scope: opening through Lemma 1.6 and its proof, stopping before Lemma 1.7.
% Authority: Polo--Gille Exp4-13oct24.pdf, component-local pages 1--3
% (combined-reader pages 189--191; printed pages 160--163).
% Jacob Reinhold's CC BY 4.0 Markdown lineage was used only as a phrasing
% comparison. The Polo--Gille PDF controls prose, formulas, numbering, and notes.

\setcounter{footnote}{0}
\SGARefTarget{sga3:iv:chapter:iv}\chapter*{Expos\'e IV\\Topologies and Sheaves}
\addcontentsline{toc}{chapter}{Expos\'e IV. Topologies and Sheaves}
\markboth{EXPOS\'E IV. TOPOLOGIES AND SHEAVES}%
  {EXPOS\'E IV. TOPOLOGIES AND SHEAVES}

\begingroup
\renewcommand{\thefootnote}{\fnsymbol{footnote}}
\begin{center}
  \textit{by M. Demazure}\footnote{This text develops the substance of two
  lectures given by A.~Grothendieck, supplementing them at several important
  points that had been passed over in silence or only briefly touched upon.}
\end{center}
\endgroup

\setcounter{footnote}{-1}

% Source page: Exp4 component-local p. 1; combined-reader p. 189; printed p. 160.
This expos\'e is intended to acquaint the reader with the essentials of the
language of topologies and sheaves (without cohomology), a language
particularly convenient for questions concerning the formation of quotients,
among others.\footnote{Editor's note: version of October 13, 2024.}

The first three sections develop the language of forming quotients. The
fourth, which is the central part, presents the theory of sheaves, oriented
mainly toward its application to questions about quotients; the fifth applies
this theory to the formation of quotients in groups and to principal
homogeneous bundles. The final section concerns the category of schemes more
specifically and defines several useful topologies on that category.

The reader will profitably consult \SGARefLink{sga3:iv:bib:as}{[AS]}, \SGARefLink{sga3:iv:bib:ma}{[MA]}, \SGARefLink{sga3:iv:bib:d}{[D]}, and \SGARefLink{sga3:iv:bib:sga:4}{SGA~4}; in particular,
\SGARefLink{sga3:iv:bib:d}{[D]} for the applications of topologies to descent theory, and \SGARefLink{sga3:iv:bib:sga:4}{SGA~4} for
questions concerning universes (which are treated especially poorly in this
expos\'e).

\SGARefTarget{sga3:iv:section:1}\section*{1. Universal Effective Epimorphisms}
\addcontentsline{toc}{section}{1. Universal Effective Epimorphisms}

Throughout the remainder of this expos\'e, a category \(\mathscr C\) is fixed.

\medskip
\noindent\SGARefTarget{sga3:iv:definition:1:1}\textbf{Definition 1.1.}
\textit{A morphism \(u:T\to S\) is called an \emph{epimorphism} if, for every
object \(X\), the corresponding map}
\[
  X(S)=\operatorname{Hom}(S,X)
  \longrightarrow
  X(T)=\operatorname{Hom}(T,X)
\]
\textit{is injective.}\footnote{Editor's note: that is, if \(u\) is
right-cancellable.}
\textit{It is called a \emph{universal epimorphism} if, for every morphism
\(S'\to S\), the fiber product \(T'=T\times_S S'\) exists and
\(u':T'\to S'\) is an epimorphism.}

% Source page: Exp4 component-local p. 2; combined-reader p. 190; printed p. 162.
\medskip
\noindent\SGARefTarget{sga3:iv:definition:1:2}\textbf{Definition 1.2.}
\textit{A diagram of maps of sets}
\[
  A\xrightarrow{u}B
  \overset{v_1}{\underset{v_2}{\rightrightarrows}}C
\]
\textit{is said to be \emph{exact} if \(u\) is injective and its image consists
of those elements \(b\in B\) for which \(v_1(b)=v_2(b)\). A diagram of the
same form in \(\mathscr C\) is said to be \emph{exact} if, for every object
\(X\) of \(\mathscr C\), the corresponding diagram of sets}
\[
  A(X)\longrightarrow B(X)\rightrightarrows C(X)
\]
\textit{is exact; one then also says that \(u\) makes \(A\) a \emph{kernel}
of the pair of arrows \((v_1,v_2)\). Dually, a diagram in \(\mathscr C\)}
\[
  C\overset{v_1}{\underset{v_2}{\rightrightarrows}}B
  \xrightarrow{u}A
\]
\textit{is said to be exact if it is exact as a diagram in the opposite
category \(\mathscr C^{\circ}\); that is, if for every object \(X\) of
\(\mathscr C\), the corresponding diagram of sets}
\[
  X(A)\longrightarrow X(B)\rightrightarrows X(C)
\]
\textit{is exact.}\footnote{Editor's note: this implies, in particular, that
\(u\) is an epimorphism.}
\textit{One also says that \(u\) makes \(A\) a \emph{cokernel} of the pair of
arrows \((v_1,v_2)\).}

\medskip
\noindent\SGARefTarget{sga3:iv:definition:1:3}\textbf{Definition 1.3.}
\textit{A morphism \(u:T\to S\) is called an \emph{effective epimorphism} if
the fiber square \(T\times_S T\) exists and the diagram}
\[
  T\times_S T
  \overset{\operatorname{pr}_1}{\underset{\operatorname{pr}_2}{\rightrightarrows}}
  T\xrightarrow{u}S
\]
\textit{is exact; that is, if \(u\) makes \(S\) a cokernel of
\((\operatorname{pr}_1,\operatorname{pr}_2)\). It is called a
\emph{universal effective epimorphism} if, for every morphism \(S'\to S\),
the fiber product \(T'=T\times_S S'\) exists and the morphism
\(u':T'\to S'\) is an effective epimorphism.}

There are evidently implications
\[
\begin{array}{ccc}
 \text{universal effective epimorphism}
   &\Longrightarrow& \text{effective epimorphism}\\
 \Big\Downarrow && \Big\Downarrow\\
 \text{universal epimorphism}
   &\Longrightarrow& \text{epimorphism},
\end{array}
\]
but in general no other implication holds.\footnote{Editor's note: for
example, if \(\mathscr C=\mathbf{Sch}\) is the category of schemes, it is
easy to see that every universal epimorphism is surjective. Let
\[
  T=\coprod_{p\ \mathrm{prime}}\operatorname{Spec}(\mathbf F_p),
  \qquad S=\operatorname{Spec}(\mathbf Z).
\]
Then the morphism \(u:T\to S\) is an epimorphism that is not universal. On
the other hand, \(T\times_S T\) identifies with \(T\), so the two
projections \(T\times_S T\to T\) coincide. Since
\(\operatorname{id}_T\) does not descend to a morphism \(S\to T\), this
shows that \(u\) is not an effective epimorphism.}

% Source page: Exp4 component-local p. 3; combined-reader p. 191; printed p. 163.
\medskip
\noindent\SGARefTarget{sga3:iv:definition:1:4:0}\textbf{Definition 1.4.0.}\footnote{Editor's note: the numbering
\SGARefLink{sga3:iv:definition:1:4:0}{1.4.0} has been added for later reference.}
We ``recall'' that a morphism \(u:T\to S\) is said to be
\emph{squarable} if, for every morphism \(S'\to S\), the fiber product
\(T\times_S S'\) exists.

\medskip
\noindent\SGARefTarget{sga3:iv:lemma:1:4}\textbf{Lemma 1.4.}
\textit{Consider morphisms \(U\xrightarrow{v}T\xrightarrow{u}S\). Then:}
\begin{enumerate}
  \item[(a)]\SGARefTarget{sga3:iv:lemma:1:4:item:a:01} \textit{if \(u\) and \(v\) are epimorphisms, then \(uv\) is an
    epimorphism; and if \(uv\) is an epimorphism, then \(u\) is an
    epimorphism;}
  \item[(b)]\SGARefTarget{sga3:iv:lemma:1:4:item:b:01} \textit{if \(u\) and \(v\) are universal epimorphisms, then
    \(uv\) is a universal epimorphism; and if \(uv\) is a universal
    epimorphism and \(u\) is squarable, then \(u\) is a universal
    epimorphism.}
\end{enumerate}

Lemma~\SGARefLink{sga3:iv:lemma:1:4}{1.4} follows immediately from the definitions. We deduce:

\medskip
\noindent\SGARefTarget{sga3:iv:corollary:1:5}\textbf{Corollary 1.5.}
\textit{Let \(u:X\to Y\) and \(u':X'\to Y'\) be universal epimorphisms,
and suppose that \(Y\times Y'\) exists. Then \(X\times X'\) exists and}
\[
  u\times u':X\times X'\longrightarrow Y\times Y'
\]
\textit{is a universal epimorphism.}

We also record the following.

\medskip
\noindent\SGARefTarget{sga3:iv:definition:1:6:0}\textbf{Definition 1.6.0.}\footnote{Editor's note: the numbering
\SGARefLink{sga3:iv:definition:1:6:0}{1.6.0} has been added for later reference.}
An object \(S\) of \(\mathscr C\) is called \emph{squarable} if its product
with every object of \(\mathscr C\) exists. (If \(\mathscr C\) has a final
object \(e\), this is equivalent to saying that the morphism \(S\to e\) is
squarable; cf. \SGARefLink{sga3:iv:definition:1:4:0}{1.4.0}.)

\medskip
\noindent\SGARefTarget{sga3:iv:lemma:1:6}\textbf{Lemma 1.6.}
\textit{Let \(u:X\to Y\) be a morphism in \(\mathscr C_{/S}\). For it to be
an epimorphism (respectively, a universal epimorphism, an effective
epimorphism, or a universal effective epimorphism), it is sufficient that the
corresponding morphism in \(\mathscr C\) have the same property; it is also
necessary if \(S\) is assumed to be a squarable object of \(\mathscr C\).}

The proof is immediate and is left to the reader. The hypothesis that \(S\)
is squarable is used to interpret the \(\mathscr C\)-morphisms from an object
\(Y\) of \(\mathscr C_{/S}\) to an object \(Z\) of \(\mathscr C\) as the
\(\mathscr C_{/S}\)-morphisms from \(Y\) to \(Z\times S\).

% END INCLUDED SGA3_English_current_buildable_source_20260823/source/inputs/IV/tex/components/01_expose_IV_opening_s11_s16.tex

% BEGIN INCLUDED SGA3_English_current_buildable_source_20260823/source/inputs/IV/tex/components/02_expose_IV_s17_through_s2.tex
% SGA 3, Expose IV, English reconstruction.
% Scope: Lemma 1.7 through the end of section 2 (Remark 2.7), stopping
% immediately before the section 3 heading.
% Authority: Polo--Gille Exp4-13oct24.pdf, component-local pages 3--9
% (combined-reader pages 191--197).
% Jacob Reinhold's CC BY 4.0 Markdown lineage was used only as a phrasing
% comparison. The Polo--Gille PDF controls prose, formulas, numbering, notes,
% and diagrams.

% Source page: Exp4 component-local p. 3; combined-reader p. 191.
\medskip
\noindent\SGARefTarget{sga3:iv:lemma:1:7}\textbf{Lemma 1.7.}
\textit{With the notation of \SGARefLink{sga3:iv:lemma:1:4}{1.4}, if \(u\) and \(v\) are effective
epimorphisms and \(v\) is a universal epimorphism, then \(uv\) is an
effective epimorphism.}

To see this, consider the diagram
\[
\begin{array}{ccccc}
  S & \xleftarrow{u} & T & \leftleftarrows & T\times_S T \\
    &                 & \scriptstyle v\,\Big\uparrow
                      &                  & \scriptstyle v\times_Sv\,\Big\uparrow \\[-0.25em]
    &                 & U & \leftleftarrows & U\times_S U \\
    &                 & \substack{\uparrow\\[-0.6ex]\uparrow}
                      & \nearrow & \\
    &                 & U\times_T U & &
\end{array}
\tag{1.7.1}\label{sga3:iv:lemma:1:7:eq:1:7:1:01}
\]
By hypothesis, the first row and the first column are exact; moreover,
by \SGARefLink{sga3:iv:corollary:1:5}{1.5} and \SGARefLink{sga3:iv:lemma:1:6}{1.6}, \(v\times_Sv\) is an epimorphism, since \(v\) is a
universal epimorphism. The conclusion follows by an evident diagram chase.
Indeed, if an element of \(X(U)\) has the same two images in
\(X(U\times_SU)\), then it has, \textit{a fortiori}, the same two images in
\(X(U\times_TU)\), and hence comes from an element of \(X(T)\), since the
first column is exact.

% Source page: Exp4 component-local p. 4; combined-reader p. 192.
Since the first row is exact, it remains only to verify that the element in
question has the same two images in \(X(T\times_ST)\). Because
\(v\times_Sv\) is an epimorphism, it is enough to verify that their images in
\(X(U\times_SU)\) are the same, which is indeed the case.

\medskip
\noindent\SGARefTarget{sga3:iv:proposition:1:8}\textbf{Proposition 1.8.}
\textit{Consider morphisms \(U\xrightarrow{v}T\xrightarrow{u}S\). If \(u\)
and \(v\) are universal effective epimorphisms, then \(uv\) is a universal
effective epimorphism. If \(uv\) is a universal effective
epimorphism and \(u\) is squarable, then \(u\) is a universal effective
epimorphism.}

The first implication follows at once from \SGARefLink{sga3:iv:lemma:1:7}{1.7}. For the second, consider \SGARefLink{sga3:iv:proposition:1:8:eq:1:8:1:01}{the
following diagram} (of ``bisimplicial'' type):
\[
\begin{array}{ccccc}
  S & \xleftarrow{u} & T & \leftleftarrows & T\times_ST \\
  \scriptstyle uv\,\Big\uparrow
    & & \Big\uparrow & & \Big\uparrow \\[-0.25em]
  U & \longleftarrow & U\times_ST
    & \leftleftarrows & U\times_ST\times_ST \\
  \substack{\uparrow\\[-0.6ex]\uparrow}
    & & \substack{\uparrow\\[-0.6ex]\uparrow}
    & & \substack{\uparrow\\[-0.6ex]\uparrow} \\[-0.25em]
  U\times_SU & \longleftarrow & U\times_SU\times_ST
    & \leftleftarrows & U\times_SU\times_ST\times_ST .
\end{array}
\tag{1.8.1}\label{sga3:iv:proposition:1:8:eq:1:8:1:01}
\]
The first, second, and third columns are exact by the hypothesis that \(uv\)
is a universal effective epimorphism. The second row is exact because
\(U\times_ST\to U\) is an effective epimorphism (it has a section over
\(U\)), and the same is true of the third row, for the same reason. An
evident diagram chase then shows that the first row is exact, that is, that
\(u\) is an effective epimorphism. Since the hypotheses are invariant under
an arbitrary base change \(S'\to S\), it follows that \(u\) is in fact a
universal effective epimorphism.

\medskip
\noindent\SGARefTarget{sga3:iv:corollary:1:9}\textbf{Corollary 1.9.}
\textit{Let \(u:X\to Y\) and \(u':X'\to Y'\) be universal effective
epimorphisms, and suppose that \(Y\times Y'\) exists. Then \(X\times X'\)
exists and}
\[
  u\times u':X\times X'\longrightarrow Y\times Y'
\]
\textit{is a universal effective epimorphism.}

The proof is the same as for \SGARefLink{sga3:iv:corollary:1:5}{1.5}, using the diagram
\[
\begin{array}{ccccc}
   & & Y' & \xleftarrow{u'} & X' \\
   & & \Big\uparrow & & \Big\uparrow \\[-0.25em]
  X & \longleftarrow & X\times Y' & \longleftarrow & X\times X' \\
  \scriptstyle u\,\Big\downarrow
    & & \Big\downarrow & \swarrow_{\scriptstyle u\times u'} & \\[-0.25em]
  Y & \longleftarrow & Y\times Y' & & .
\end{array}
\tag{1.9.1}\label{sga3:iv:corollary:1:9:eq:1:9:1:01}
\]

% Source page: Exp4 component-local p. 5; combined-reader p. 193.
\medskip
\noindent\SGARefTarget{sga3:iv:corollary:1:10}\textbf{Corollary 1.10.}
\textit{Let \(u:T\to S\) be a squarable morphism, and let \(S'\to S\) be a
base-change morphism that is a universal effective epimorphism. Then \(u\)
is a universal effective epimorphism if and only if}
\[
  u':T'=T\times_SS'\longrightarrow S'
\]
\textit{is one:}
\[
\begin{array}{ccc}
  T & \xleftarrow{v'} & T' \\
  \scriptstyle u\,\Big\downarrow & & \Big\downarrow\,\scriptstyle u' \\[-0.25em]
  S & \xleftarrow{v} & S'.
\end{array}
\tag{1.10.1}\label{sga3:iv:corollary:1:10:eq:1:10:1:01}
\]
Only sufficiency requires proof. If \(u'\) is a universal effective
epimorphism, then so is \(vu'\), by \SGARefLink{sga3:iv:proposition:1:8}{1.8}. Since \(vu'=uv'\), another
application of \SGARefLink{sga3:iv:proposition:1:8}{1.8} shows that \(u\) is a universal effective epimorphism.

\medskip
\noindent\SGARefTarget{sga3:iv:remark:1:11}\textbf{Remark 1.11.}
\textit{The same reasoning shows that in \SGARefLink{sga3:iv:corollary:1:10}{1.10} one may replace ``universal
effective epimorphism'' by ``universal epimorphism,'' by ``universal and
effective epimorphism,'' or simply by ``epimorphism'' (and in the last case
the hypothesis that \(u\) be squarable is plainly unnecessary).}

In the proof of \SGARefLink{sga3:iv:proposition:1:8}{1.8} we used \SGARefLink{sga3:iv:proposition:1:12}{the following result}, which deserves to be stated
explicitly.

\medskip
\noindent\SGARefTarget{sga3:iv:proposition:1:12}\textbf{Proposition 1.12.}
\textit{Let \(u:T\to S\) be a morphism admitting a section. Then \(u\) is an
epimorphism; if \(T\times_ST\) exists, it is an effective epimorphism; and if,
in addition, \(u\) is squarable, it is a universal effective epimorphism.}

The first assertion is contained in \SGARefLink{sga3:iv:lemma:1:4:item:a:01}{1.4(a)}, and the third follows at once
from the second, so it is enough to establish the second. In fact, there is a
stronger conclusion: for every functor
\(F:\mathscr C^{\circ}\to\mathbf{Set}\), not necessarily representable, the
diagram of sets
\[
  F(S)\longrightarrow F(T)\rightrightarrows F(T\times_ST)
\]
is exact. This may be regarded as a special case of the formalism of
\v{C}ech cohomology (in dimension \(0\)!), which we merely recall here.
Assume only that \(T\times_ST\) exists, and set
\[
  \check H^{0}(T/S,F)
  =\operatorname{Ker}\bigl(F(T)\rightrightarrows F(T\times_ST)\bigr).
\]

One may plainly regard \(\check H^{0}(T/S,F)\) as a contravariant functor of
the variable \(T\) in \(\mathscr C_{/S}\): every \(S\)-morphism \(T'\to T\)
defines a map
\[
\tag{+}\label{sga3:iv:proposition:1:12:eq:unnamed:01}
  \check H^{0}(T/S,F)\longrightarrow\check H^{0}(T'/S,F).
\]
Fix \(T,T'\) in \(\mathscr C_{/S}\). A familiar calculation shows that, if
there exists an \(S\)-morphism \(T'\to T\), the corresponding map \SGARefLink{sga3:iv:proposition:1:12:eq:unnamed:01}{(+)} is
independent of the choice of that morphism.\footnote{Editor's note: this is
the following argument, communicated by M.~Demazure. Let
\(f,g:T'\to T\), and let \(\phi:T'\to T\times_ST\) be the morphism with
components \(f\) and \(g\), so that \(p_1\circ\phi=f\) and
\(p_2\circ\phi=g\). Then
\(F(\phi):F(T\times_ST)\to F(T')\) satisfies
\(F(\phi)\circ F(p_1)=F(f)\) and
\(F(\phi)\circ F(p_2)=F(g)\). For every
\(x\in\check H^0(T/S,F)\), however, one has
\(F(p_1)(x)=F(p_2)(x)\). Applying \(F(\phi)\) to both sides gives
\(F(f)(x)=F(g)(x)\), which shows that \(f\) and \(g\) induce the same map.}

% Source page: Exp4 component-local p. 6; combined-reader p. 194.
Thus \(\check H^{0}(T/S,F)\) may be regarded as a functor on the category
associated with the set \(\operatorname{Ob}(\mathscr C_{/S})\), preordered
by the relation of domination: \(T'\) dominates \(T\) if there exists an
\(S\)-morphism \(T'\to T\). In particular, if \(T\) and \(T'\) are
isomorphic in this latter category, that is, if each dominates the other,
then \SGARefLink{sga3:iv:proposition:1:12:eq:unnamed:01}{(+)} is an isomorphism of sets. This applies in particular when \(T'\)
is the final object of \(\mathscr C_{/S}\), essentially \(S\) itself. The
object \(T\) always dominates \(T'=S\), while the converse holds precisely
when \(T/S\) has a section. This proves \SGARefLink{sga3:iv:proposition:1:12}{1.12} in the stronger form announced.

\medskip
\noindent\SGARefTarget{sga3:iv:remark:1:13}\textbf{Remark 1.13.}
\textit{For various applications, the notions introduced and the results
stated in this expos\'e must be developed more generally for a family of
morphisms \(u_i:T_i\to S\) with the same target, rather than for a single
morphism \(u:T\to S\). Such a family is called \emph{epimorphic} if, for
every object \(X\) of \(\mathscr C\), the corresponding map}
\[
  X(S)\longrightarrow\prod_i X(T_i)
\]
\textit{is injective. One similarly introduces the notion of an effective
epimorphic family and the universal variants of these notions. Whenever
needed below, we shall take for granted that the results of this expos\'e
extend to this more general situation.}

\medskip
\noindent\SGARefTarget{sga3:iv:remark:1:14}\textbf{Remark 1.14.}
\textit{Every morphism that is both a monomorphism and an effective
epimorphism is an isomorphism.}

Indeed, in the notation of \SGARefLink{sga3:iv:definition:1:3}{1.3}, the fact that \(T\to S\) is a monomorphism
implies that the two morphisms
\[
  \operatorname{pr}_1,\operatorname{pr}_2:T\times_ST
  \rightrightarrows T
\]
are equal (and are isomorphisms). But a diagram
\[
  C\overset{v}{\underset{v}{\rightrightarrows}}B\xrightarrow{u}A
\]
is exact only if \(u\) is an isomorphism, as follows immediately from \SGARefLink{sga3:iv:definition:1:3}{the
definition}.\footnote{Editor's note: a monomorphism that is an epimorphism
need not be an isomorphism. For example, in
\(\mathscr C=\mathbf{Sch}\), the morphism
\[
  \operatorname{Spec}(\mathbf F_p)\amalg
  \operatorname{Spec}(\mathbf Z[1/p])
  \longrightarrow\operatorname{Spec}(\mathbf Z)
\]
is a monomorphism and a surjective epimorphism, but is not an isomorphism.}

\SGARefTarget{sga3:iv:section:2}\section*{2. Descent Morphisms}
\addcontentsline{toc}{section}{2. Descent Morphisms}

Recall the following definitions.

\medskip
\noindent\SGARefTarget{sga3:iv:definition:2:1}\textbf{Definition 2.1.}
\textit{Let \(f:S'\to S\) be a morphism such that
\(S''=S'\times_SS'\) exists, and let \(u':X'\to S'\) be an object over
\(S'\). A \emph{gluing datum} on \(X'/S'\), relative to \(f\), is an
\(S''\)-isomorphism}
\[
  c:X''_1\xrightarrow{\sim}X''_2,
\]
\textit{where \(X''_i\), for \(i=1,2\), denotes the inverse image, assumed
to exist, of \(X'/S'\) under the projection
\(\operatorname{pr}_i:S''\to S'\). The gluing datum \(c\) is called a
\emph{descent datum} if it satisfies the ``cocycle condition''}
\[
  \operatorname{pr}_{3,1}^{*}(c)
  =\operatorname{pr}_{3,2}^{*}(c)\,
   \operatorname{pr}_{2,1}^{*}(c),
\]
\textit{where \(\operatorname{pr}_{i,j}\), for
\(1\leq j<i\leq3\), are the canonical projections from}
\[
  S'''=S'\times_SS'\times_SS'
\]
\textit{to \(S''\) (and we now assume that \(S'''\) also exists), where
\(\operatorname{pr}_{i,j}^{*}(c)\) is the inverse image of \(c\), viewed
as an \(S'''\)-morphism from \(X'''_j\) to \(X'''_i\), and where, for each
integer \(k\) from \(1\) to \(3\), \(X'''_k\) denotes the inverse image,
assumed to exist, of \(X'/S'\) under the \(k\)-th projection
\(q_k:S'''\to S'\).}

% Source page: Exp4 component-local p. 7; combined-reader p. 195.
In the second part of \SGARefLink{sga3:iv:definition:2:1}{the definition} we have therefore used customary
identifications and abuses of notation,\footnote{Editor's note: for example,
we have identified, on the one hand,
\(\operatorname{pr}_{2,1}^{*}(X''_1)
 =\operatorname{pr}_{3,1}^{*}(X''_1)\), and, on the other hand,
\(\operatorname{pr}_{2,1}^{*}(X''_2)
 =\operatorname{pr}_{3,2}^{*}(X''_2)\).} which experience shows to be
harmless, but which must plainly be avoided in a rigorous exposition of
descent theory---a theory that must itself justify, to some extent, these
customary abuses of language. Such a formal exposition (\SGARefLink{sga3:iv:bib:d}{[D]}) was written by
Giraud, with a view to justifying and making precise \SGARefLink{sga3:iv:bib:sga:1}{SGA~1}, Expos\'e~VII,
which was never written. For a detailed account of the faithfully flat
descent results used constantly in this Seminar, one may consult \SGARefLink{sga3:iv:bib:sga:1}{SGA~1},
Expos\'e~VIII.

Continue to let \(f:S'\to S\) be a morphism for which
\(S''=S'\times_SS'\) exists, and let \(X\) be an object over \(S\) such
that \(X'=X\times_SS'\) and \(X''=X\times_SS''\) exist. The inverse images
of \(X'\) under \(\operatorname{pr}_i\), for \(i=1,2\), then exist and are
canonically isomorphic; consequently, \(X'/S'\) carries a canonical gluing
datum relative to \(f\). When \(S'''\) and \(X'''=X\times_SS'''\) exist,
this is in fact a descent datum. If \(Y\) is another object over \(S\)
satisfying the same conditions as \(X/S\), then for every \(S\)-morphism
\(X\to Y\), the corresponding \(S'\)-morphism \(X'\to Y'\) is compatible
with the canonical gluing data on \(X'\) and \(Y'\). In particular, if
\(S'\to S\) is squarable, then
\[
  X\longmapsto X'=X\times_SS'
\]
is a functor from \(\mathscr C_{/S}\) to the ``category of objects over
\(S'\) equipped with a descent datum relative to \(f\)''---a category whose
definition is left to the reader and which is a full subcategory of the
``category of objects over \(S'\) equipped with a gluing datum relative to
\(f\).''

With this understood, we make \SGARefLink{sga3:iv:definition:2:2}{the following definition}.

\medskip
\noindent\SGARefTarget{sga3:iv:definition:2:2}\textbf{Definition 2.2.}
\textit{A morphism \(f:S'\to S\) is called a \emph{descent morphism}
(respectively, an \emph{effective descent morphism}) if \(f\) is squarable
(that is, \(X\times_SS'\) exists for every \(X\to S\)) and the preceding
functor}
\[
  X\longmapsto X'=X\times_SS'
\]
\textit{from the category \(\mathscr C_{/S}\) of objects over \(S\) to the
category of objects over \(S'\) equipped with a descent datum relative to
\(f\) is fully faithful (respectively, an equivalence of categories).}

The first of these notions can be expressed using gluing data alone, and
therefore without invoking the triple fiber product \(S'''\): \(f\) is a
descent morphism precisely when it is squarable and \(X\mapsto X'\) is a
fully faithful functor from \(\mathscr C_{/S}\) to the category of objects
over \(S'\) equipped with a gluing datum relative to \(f\). Written out, this
means that for any two objects \(X,Y\) over \(S\), \SGARefLink{sga3:iv:definition:2:2:eq:x:01}{the following diagram} of
sets is exact:
\[
\tag{x}\label{sga3:iv:definition:2:2:eq:x:01}
  \operatorname{Hom}_S(X,Y)
  \longrightarrow \operatorname{Hom}_{S'}(X',Y')
  \overset{p_1}{\underset{p_2}{\rightrightarrows}}
  \operatorname{Hom}_{S''}(X'',Y'').
\]

% Source page: Exp4 component-local p. 8; combined-reader p. 196.
Here \(p_i(u')\) denotes the inverse image of
\(u'\in\operatorname{Hom}_{S'}(X',Y')\) under the projection
\(\operatorname{pr}_i:S''\to S'\), for \(i=1,2\). Indeed, by definition,
the kernel of the pair \((p_1,p_2)\) is precisely the set of
\(S'\)-morphisms \(X'\to Y'\) compatible with the canonical gluing data.

By the definition of the inverse images \(Y'\) and \(Y''\), there are
canonical bijections
\[
  \operatorname{Hom}_{S'}(X',Y')\simeq\operatorname{Hom}_S(X',Y),
  \qquad
  \operatorname{Hom}_{S''}(X'',Y'')\simeq\operatorname{Hom}_S(X'',Y).
\]
Thus exactness of \SGARefLink{sga3:iv:definition:2:2:eq:x:01}{(x)} is equivalent to exactness of
\[
\tag{xx}\label{sga3:iv:definition:2:2:eq:xx:01}
  \operatorname{Hom}_S(X,Y)
  \longrightarrow\operatorname{Hom}_S(X',Y)
  \rightrightarrows\operatorname{Hom}_S(X'',Y),
\]
which is obtained by applying \(\operatorname{Hom}_S(-,Y)\) to the diagram
in \(\mathscr C_{/S}\)
\[
\tag{xxx}\label{sga3:iv:definition:2:2:eq:xxx:01}
  X''\rightrightarrows X'\longrightarrow X,
\]
itself obtained from
\[
  S''\rightrightarrows S'\longrightarrow S
\]
by the base change \(X\to S\). In view of \SGARefLink{sga3:iv:lemma:1:6}{1.6}, this proves the first part of
\SGARefLink{sga3:iv:proposition:2:3}{the following proposition}.

\medskip
\noindent\SGARefTarget{sga3:iv:proposition:2:3}\textbf{Proposition 2.3.}
\textit{Let \(f:S'\to S\) be a morphism. If \(f\) is a universal effective
epimorphism, then it is a descent morphism. Conversely, every descent
morphism is a universal effective epimorphism if \(S\) is a squarable object
of \(\mathscr C\), that is, if its product with every object of
\(\mathscr C\) exists.}

It remains to prove the converse. If \(f\) is a descent morphism, we must
show that it is a universal effective epimorphism: for every morphism
\(X\to S\), diagram \SGARefLink{sga3:iv:definition:2:2:eq:xxx:01}{(xxx)} must be exact, or equivalently, its image under
\(\operatorname{Hom}(-,Z)\) must be an exact diagram of sets for every
object \(Z\) of \(\mathscr C\). By hypothesis, \(Z\times S\) exists. Let
\(Y\) be the object of \(\mathscr C_{/S}\) that it defines. The image of
\SGARefLink{sga3:iv:definition:2:2:eq:xxx:01}{(xxx)} under \(\operatorname{Hom}(-,Z)\) is then isomorphic to its image
under \(\operatorname{Hom}_S(-,Y)\), and the latter is exact by the
hypothesis on \(f\).

The results concerning descent morphisms may therefore be applied to
universal effective epimorphisms. For example:

\medskip
\noindent\SGARefTarget{sga3:iv:proposition:2:4}\textbf{Proposition 2.4.}
\textit{Let \(f:S'\to S\) be a descent morphism (for example, a universal
effective epimorphism). Then:}
\begin{enumerate}
  \item[(a)]\SGARefTarget{sga3:iv:proposition:2:4:item:a:01} \textit{For every \(S\)-morphism \(u:X\to Y\), the morphism
    \(u\) is an isomorphism (respectively, a monomorphism) if and only if
    \(u':X'\to Y'\) is one.}
  \item[(b)]\SGARefTarget{sga3:iv:proposition:2:4:item:b:01} \textit{Let \(X,Y\) be two subobjects of \(S\), and let
    \(X',Y'\) be their inverse images in \(S'\). Then \(X\) is contained in
    \(Y\) (respectively, is equal to \(Y\)) if and only if the same is true
    of \(X'\) and \(Y'\).}
\end{enumerate}

% Source page: Exp4 component-local p. 9; combined-reader p. 197.
For \SGARefLink{sga3:iv:proposition:2:4:item:a:01}{(a)}, \SGARefLink{sga3:iv:definition:2:2}{the definition} shows that if \(u'\) is an isomorphism in the
category of objects with gluing datum, then \(u\) is an isomorphism. Every
isomorphism of objects over \(S'\) that is compatible with gluing data is,
however, an isomorphism of objects with gluing datum: its inverse is also
compatible with the gluing data. For \SGARefLink{sga3:iv:proposition:2:4:item:b:01}{(b)}, it remains to show that if \(X'\)
is contained in \(Y'\), that is, if there exists an \(S'\)-morphism
\(X'\to Y'\), then the same is true of \(X,Y\) over \(S\). Since
\(Y'\to S'\), and hence \(Y''\to S''\), is a monomorphism, the morphism
\(X'\to Y'\) is automatically compatible with the gluing data and therefore
comes from an \(S\)-morphism \(X\to Y\).

The proof applies more generally to two objects \(X,Y\) over \(S\), with
\(Y\to S\) a monomorphism, when one asks whether \(X\to S\) factors through
\(Y\): it is enough that \(X'\to S'\) factor through \(Y'\).

\medskip
\noindent\SGARefTarget{sga3:iv:corollary:2:5}\textbf{Corollary 2.5.}
\textit{Let \(f:S'\to S\) be a universal effective epimorphism and
\(g:S\to T\) a morphism such that \(S\times_TS\) exists. Suppose that
\(S''=S'\times_SS'\) is also a fiber product of \(S'\) with itself over
\(T\), that is, that the canonical morphism}
\[
  S'\times_SS'\xrightarrow{\sim}S'\times_TS'
\]
\textit{is an isomorphism. Then \(g:S\to T\) is a monomorphism (and the
converse is of course true).}

Indeed, consider the cartesian diagram
\[
\begin{array}{ccc}
  S'\times_SS' & \xrightarrow{\sim} & S'\times_TS' \\
  \Big\downarrow & & \Big\downarrow\,\scriptstyle f \\[-0.25em]
  S & \longrightarrow & S\times_TS ,
\end{array}
\tag{2.5.1}\label{sga3:iv:corollary:2:5:eq:2:5:1:01}
\]
where the second horizontal arrow is the diagonal morphism. By \SGARefLink{sga3:iv:corollary:1:9}{1.9}, the
second vertical arrow \(f\) is a universal effective epimorphism; by
hypothesis, the first horizontal arrow is an isomorphism. It follows from
\SGARefLink{sga3:iv:proposition:2:4:item:a:01}{2.4(a)}, or equally from \SGARefLink{sga3:iv:proposition:2:4:item:b:01}{2.4(b)},\footnote{Editor's note: applied to
\(f:S'\times_TS'\to S\times_TS=Y\).} that
\(S\to S\times_TS\) is also an isomorphism, which says precisely that
\(g:S\to T\) is a monomorphism.

\medskip
\noindent\SGARefTarget{sga3:iv:remark:2:6}\textbf{Remark 2.6.}
\textit{The notions introduced in \SGARefLink{sga3:iv:definition:2:1}{2.1} in terms of morphisms between certain
inverse limits have an evident expression in terms of the contravariant
functors defined by the objects \(S,S',X'\) under consideration. Subject to
the existence of the fiber products considered in \SGARefLink{sga3:iv:definition:2:1}{2.1}, there is a one-to-one
correspondence between gluing data (respectively, descent data) on \(X'/S'\)
relative to \(S'\to S\), and gluing data (respectively, descent data) for the
corresponding objects of}
\[
  \widehat{\mathscr C}
  =\operatorname{Hom}(\mathscr C^{\circ},\mathbf{Set}).
\]
\textit{This makes it possible, if desired, to extend these notions to the
case in which no existence hypothesis is imposed on fiber products in
\(\mathscr C\).}

\medskip
\noindent\SGARefTarget{sga3:iv:remark:2:7}\textbf{Remark 2.7.}
\textit{The notions introduced in \SGARefLink{sga3:iv:section:2}{this section} generalize to families of
morphisms. They can also be presented more intrinsically by means of the
notion of a sieve (\SGARefLink{sga3:iv:subsection:4:1}{4.1}). For these matters the reader is referred to \SGARefLink{sga3:iv:bib:d}{[D]}.}

% END INCLUDED SGA3_English_current_buildable_source_20260823/source/inputs/IV/tex/components/02_expose_IV_s17_through_s2.tex

% BEGIN INCLUDED SGA3_English_current_buildable_source_20260823/source/inputs/IV/tex/components/03_expose_IV_s3.tex
% SGA 3, Expose IV, English reconstruction.
% Scope: all of Section 3, stopping immediately before the Section 4 heading.
% Authority: Polo--Gille Exp4-13oct24.pdf, component-local pages 10--20
% (combined-reader pages 198--208; printed pages 186--196).
% Jacob Reinhold's CC BY 4.0 Markdown lineage was used only as a phrasing
% comparison. The Polo--Gille PDF controls prose, formulas, numbering, and notes.

\setcounter{footnote}{9}

% Source page: Exp4 component-local p. 10; combined-reader p. 198; printed p. 186.
\SGARefTarget{sga3:iv:section:3}\section*{3. Universal Effective Equivalence Relations}
\addcontentsline{toc}{section}{3. Universal Effective Equivalence Relations}

\SGARefTarget{sga3:iv:subsection:3:1}\subsection*{3.1. Equivalence Relations: Definitions}
\addcontentsline{toc}{subsection}{\SGARefLink{sga3:iv:subsection:3:1}{3.1}. Equivalence Relations: Definitions}

\noindent\SGARefTarget{sga3:iv:definition:3:1:1}\textbf{Definition 3.1.1.}
\textit{A \(\mathscr C\)-equivalence relation in an object
\(X\in\operatorname{Ob}\mathscr C\) is a representable subfunctor \(R\) of
the functor \(X\times X\) such that, for every
\(S\in\operatorname{Ob}\mathscr C\), \(R(S)\) is the graph of an equivalence
relation on \(X(S)\).}

\SGARefLink{sga3:iv:definition:3:1:1}{This definition} applies in particular to the category
\(\widehat{\mathscr C}\). If \(X\) is regarded as an object of
\(\widehat{\mathscr C}\), then a \(\widehat{\mathscr C}\)-equivalence relation
in \(X\) is simply a subfunctor \(R\) of \(X\times X\) such that \(R(S)\) is
the graph of an equivalence relation on \(X(S)\), for every object \(S\) of
\(\mathscr C\).\footnote{Editor's note: the condition is plainly necessary.
Conversely, if \(R(S)\) is the graph of an equivalence relation for every
\(S\in\operatorname{Ob}\mathscr C\), this equivalence relation extends to
\(R(F)\) for every \(F\in\operatorname{Ob}\widehat{\mathscr C}\): two
morphisms \(\varphi,\psi:F\to R\) are declared equivalent if, for every
\(S\in\operatorname{Ob}\mathscr C\) and every \(x\in F(S)\), the elements
\(\varphi(x)\) and \(\psi(x)\) are equivalent in \(X(S)\).}

If \(R\) is a \(\mathscr C\)-equivalence relation in \(X\), let \(p_i\)
denote the morphism from \(R\) to \(X\) induced by the projection
\(\operatorname{pr}_i:X\times X\to X\). Thus there is a diagram
\[
  p_1,p_2:R\rightrightarrows X.
\]

\medskip
\noindent\SGARefTarget{sga3:iv:definition:3:1:2}\textbf{Definition 3.1.2.}
\textit{A morphism \(u:X\to Z\) is said to be \emph{compatible with \(R\)}
if \(up_1=up_2\). An object of \(\mathscr C\) that is a cokernel of the pair
\((p_1,p_2)\) is also called a \emph{quotient object} of \(X\) by \(R\), and
is denoted \(X/R\).}

Thus one has an exact diagram
\[
  R\overset{p_1}{\underset{p_2}{\rightrightarrows}}
  X\xrightarrow{p}X/R,
\]
and \(X/R\) represents the covariant functor
\[
  \operatorname{Hom}_{\mathscr C}(X/R,Z)
  =\{\text{morphisms from \(X\) to \(Z\) compatible with \(R\)}\}.
\]
Since quotient objects are taken in \(\mathscr C\), the quotient \(X/R\) is
unique whenever it exists.

These definitions extend at once to a
\(\widehat{\mathscr C}\)-equivalence relation in \(X\). One should note,
however, that the identification of the objects of \(\mathscr C\) with their
images in \(\widehat{\mathscr C}\) does not commute with the formation of
quotients: the quotient \(X/R\) of \(X\) by \(R\) in \(\mathscr C\) is not
\emph{a priori} a quotient of \(X\) by \(R\) in
\(\widehat{\mathscr C}\). One must therefore avoid identifying
\(\mathscr C\) uncritically with its image in \(\widehat{\mathscr C}\) in
questions involving passage to a quotient.\footnote{Editor's note: here is an
illustration that also previews the sequel of this expos\'e. Let \(G\) be a
\(\mathscr C\)-group, let \(H\) be a \(\mathscr C\)-subgroup, and let \(R\)
be the equivalence relation in \(G\) defined by
\(G\times H\to G\times G\), \((g,h)\mapsto(g,gh)\) (cf. \SGARefLink{sga3:iv:subsection:3:2}{3.2}). The functor
\(Q\) defined by \(Q(S)=G(S)/H(S)\) is a quotient in
\(\widehat{\mathscr C}\) (by \SGARefLink{sga3:iv:proposition:4:4:9}{4.4.9} applied to the coarsest topology; cf.
\SGARefLink{sga3:iv:scholium:4:4:2}{4.4.2}), but it is not in general the quotient one wants. For example, when
\(\mathscr C=\mathbf{Sch}\), there is an exact sequence of affine group
schemes
\[
  1\longrightarrow\mu_2\longrightarrow\mathbf G_m
  \xrightarrow{p}\mathbf G_m\longrightarrow1
\]
that identifies \(\mathbf G_m\) with the quotient \(\mathbf G_m/\mu_2\).
Moreover, because \(p\) is finite and locally free, \(\mathbf G_m\) is the
sheaf quotient of \(\mathbf G_m\) by \(\mu_2\) in the larger category of
sheaves for the fppf topology; cf. \SGARefLink{sga3:iv:theorem:4:6:6:item:ii:01}{4.6.6(ii)} and \SGARefLink{sga3:iv:corollary:6:3:2}{6.3.2}. By contrast, the
quotient \(Q\) in \(\widehat{\mathscr C}\) is not isomorphic to
\(\mathbf G_m\), since, for example, \(Q(\mathbf Z)=\{1\}\), whereas
\(\mathbf G_m(\mathbf Z)=\{\pm1\}\). Thus \(Q\) is not an fppf sheaf and,
\emph{a fortiori}, is not representable.}

% Source page: Exp4 component-local p. 11; combined-reader p. 199; printed p. 187.
In what follows, we shall simply say \emph{equivalence relation} to mean a
\(\widehat{\mathscr C}\)-equivalence relation; whenever necessary, we shall
specify that a \(\mathscr C\)-equivalence relation is intended.\footnote{Editor's
note: even for a \(\widehat{\mathscr C}\)-equivalence relation in \(X\), the
question of interest is the existence of a quotient in \(\mathscr C\).}

\medskip
\noindent\SGARefTarget{sga3:iv:definition:3:1:3}\textbf{Definition 3.1.3.}
\textit{If \(X\) is an object of \(\mathscr C\) over \(S\), an
\emph{equivalence relation in \(X\) over \(S\)} is an equivalence relation
\(R\) in \(X\) for which the structure morphism \(X\to S\) is compatible
with \(R\).}

The canonical monomorphism \(R\to X\times X\) then factors through the
monomorphism
\[
  X\times_S X\longrightarrow X\times X
\]
and defines an equivalence relation in the object \(X\to S\) of
\(\mathscr C_{/S}\). Whenever the quotient \(X/R\) exists, it is equipped
with a canonical morphism to \(S\), and the corresponding object of
\(\mathscr C_{/S}\) is a quotient of \(X\in\operatorname{Ob}\mathscr C_{/S}\)
by the preceding equivalence relation. Conversely, if \(S\) is a squarable
object of \(\mathscr C\) and \(Y\to S\) is a quotient of \(X\) by this
equivalence relation in \(\mathscr C_{/S}\), then \(Y\) is a quotient of
\(X\) by \(R\) in \(\mathscr C\). In any event, we shall never have to
consider quotients in \(\mathscr C_{/S}\) that are not already quotients in
\(\mathscr C\).

\medskip
\noindent\SGARefTarget{sga3:iv:definition:3:1:4}\textbf{Definition 3.1.4.}
\textit{Let \(X\) (respectively \(X'\)) be an object of \(\mathscr C\)
equipped with an equivalence relation \(R\) (respectively \(R'\)). A morphism
\[
  u:X\longrightarrow X'
\]
is said to be \emph{compatible with \(R\) and \(R'\)} if the following
equivalent conditions hold:}
\begin{enumerate}
  \item[(i)]\SGARefTarget{sga3:iv:definition:3:1:4:item:i:01} \textit{for every \(S\in\operatorname{Ob}\mathscr C\), two
  points of \(X(S)\) congruent modulo \(R(S)\) are carried by \(u\) to two
  points of \(X'(S)\) congruent modulo \(R'(S)\);}
  \item[(ii)]\SGARefTarget{sga3:iv:definition:3:1:4:item:ii:01} \textit{there exists a necessarily unique morphism \(R\to R'\)
  making the diagram}
  \[
    \begin{array}{ccc}
      R & \longrightarrow & R'\\
      \big\downarrow && \big\downarrow\\
      X\times X & \xrightarrow{u\times u} & X'\times X'
    \end{array}
  \]
  \textit{commute.}
\end{enumerate}

% Source page: Exp4 component-local p. 12; combined-reader p. 200; printed p. 188.
By the universal property of \(X/R\), whenever \(X/R\) and \(X'/R'\)
exist there is a unique morphism \(v\) making the diagram
\[
  \begin{array}{ccc}
    X & \xrightarrow{p} & X/R\\
    {\scriptstyle u}\big\downarrow && \big\downarrow{\scriptstyle v}\\
    X' & \xrightarrow{p'} & X'/R'
  \end{array}
\]
commute.

\medskip
\noindent\SGARefTarget{sga3:iv:definition:3:1:5}\textbf{Definition 3.1.5.}
\textit{A subobject (that is, a representable subfunctor) \(Y\) of \(X\) is
said to be \emph{stable under the equivalence relation \(R\)} if the following
equivalent conditions hold:}
\begin{enumerate}
  \item[(i)]\SGARefTarget{sga3:iv:definition:3:1:5:item:i:01} \textit{for every \(S\in\operatorname{Ob}\mathscr C\), the
  subset \(Y(S)\) of \(X(S)\) is stable under \(R(S)\);}
  \item[(ii)]\SGARefTarget{sga3:iv:definition:3:1:5:item:ii:01} \textit{the two subobjects of \(R\) that are the inverse
  images of \(Y\) under \(p_1\) and \(p_2\) are identical.}
\end{enumerate}

An important special case of a stable subobject of \(X\) is the following:
the quotient \(X/R\) exists and \(Y\) is the inverse image in \(X\) of a
subobject of \(X/R\).

\medskip
\noindent\SGARefTarget{sga3:iv:definition:3:1:6}\textbf{Definition 3.1.6.}
\textit{Let \(R\) be an equivalence relation in \(X\), and let \(X'\to X\)
be a morphism. The equivalence relation \(R'\) in \(X'\) obtained from the
cartesian diagram}
\[
  \begin{array}{ccc}
    R' & \longrightarrow & R\\
    \big\downarrow && \big\downarrow\\
    X'\times X' & \longrightarrow & X\times X
  \end{array}
\]
\textit{is called the equivalence relation in \(X'\) that is the
\emph{inverse image} of the equivalence relation \(R\) in \(X\). In
particular, if \(X'\) is a subobject of \(X\), it is called the equivalence
relation \emph{induced} on \(X'\) by \(R\), and is denoted \(R_{X'}\).}

The morphism \(X'\to X\) is compatible with \(R'\) and \(R\); hence, when
the quotients exist, it induces a morphism \(X'/R'\to X/R\) (\SGARefLink{sga3:iv:definition:3:1:4}{3.1.4}). If
\(X'\) is a subobject of \(X\), we shall see later that in certain cases one
can prove that \(X'/R'\to X/R\) is a monomorphism and therefore identifies
\(X'/R'\) with a subobject of \(X/R\). In that situation, the inverse image
of this subobject in \(X\) is a subobject containing \(X'\) and stable under
\(R\): the \emph{saturation} of \(X'\) with respect to the equivalence
relation \(R\).

\medskip
\noindent\SGARefTarget{sga3:iv:proposition:3:1:7}\textbf{Proposition 3.1.7.}
\textit{If the subobject \(Y\) of \(X\) is stable under \(R\), then for
\(i=1,2\) the following two squares are cartesian:}
\[
  \begin{array}{ccc}
    R_Y & \longrightarrow & R\\
    {\scriptstyle p_i}\big\downarrow && \big\downarrow{\scriptstyle p_i}\\
    Y & \longrightarrow & X.
  \end{array}
\]

The proof is immediate.

% Source page: Exp4 component-local p. 13; combined-reader p. 201; printed p. 189.
\SGARefTarget{sga3:iv:subsection:3:2}\subsection*{3.2. Equivalence Relation Defined by a Group Acting Freely}
\addcontentsline{toc}{subsection}{\SGARefLink{sga3:iv:subsection:3:2}{3.2}. Equivalence Relation Defined by a Group Acting Freely}

\noindent\SGARefTarget{sga3:iv:definition:3:2:1}\textbf{Definition 3.2.1.}
\textit{Let \(X\) be an object of \(\mathscr C\), and let \(H\) be a
\(\mathscr C\)-group acting on \(X\), that is, equipped with a morphism of
\(\widehat{\mathscr C}\)-groups}
\[
  H\longrightarrow\operatorname{Aut}(X).
\]
\textit{The group \(H\) is said to act \emph{freely} on \(X\) if the
following equivalent conditions hold:}
\begin{enumerate}
  \item[(i)]\SGARefTarget{sga3:iv:definition:3:2:1:item:i:01} \textit{for every \(S\in\operatorname{Ob}\mathscr C\), the
  group \(H(S)\) acts freely on \(X(S)\);}
  \item[(ii)]\SGARefTarget{sga3:iv:definition:3:2:1:item:ii:01} \textit{the morphism of functors}
  \[
    H\times X\longrightarrow X\times X
  \]
  \textit{defined on sets by \((h,x)\mapsto(hx,x)\) is a monomorphism.}
\end{enumerate}

In \SGARefLink{sga3:iv:definition:3:2:1}{the preceding definition} the group \(H\) was assumed to act on \(X\) on
the left. There is of course an analogous notion when the group acts on the
right, that is, when one is given a group morphism from the opposite group
\(H^{\circ}\) to \(\operatorname{Aut}(X)\).

If \(H\) acts freely on \(X\), the image of \(H\times X\) under the
morphism in \SGARefLink{sga3:iv:definition:3:2:1:item:ii:01}{(ii)} is an equivalence relation in \(X\), called the
\emph{equivalence relation defined by the action of \(H\) on \(X\)}. When
the quotient of \(X\) by this equivalence relation exists, it is denoted
\(H\backslash X\) (or \(X/H\) when \(H\) acts on the right). It represents
the following covariant functor: for every object \(Z\) of \(\mathscr C\),
\[
  \operatorname{Hom}(H\backslash X,Z)
  =\{\text{morphisms from \(X\) to \(Z\) invariant under \(H\)}\},
\]
where a morphism \(f:X\to Z\) is said to be invariant under \(H\) if, for
every \(S\in\operatorname{Ob}\mathscr C\), the corresponding morphism
\(X(S)\to Z(S)\) is invariant under the group \(H(S)\).

\medskip
\noindent\SGARefTarget{sga3:iv:lemma:3:2:2}\textbf{Lemma 3.2.2.}
\textit{Under the assumptions of \SGARefLink{sga3:iv:definition:3:2:1}{3.2.1}, let \(Y\) be a subobject of \(X\).
The following conditions are equivalent:}
\begin{enumerate}
  \item[(i)]\SGARefTarget{sga3:iv:lemma:3:2:2:item:i:01} \textit{\(Y\) is stable under the equivalence relation defined
  by \(H\) (\SGARefLink{sga3:iv:definition:3:1:5}{3.1.5});}
  \item[(ii)]\SGARefTarget{sga3:iv:lemma:3:2:2:item:ii:01} \textit{for every \(S\in\operatorname{Ob}\mathscr C\), the
  subset \(Y(S)\) of \(X(S)\) is stable under \(H(S)\);}
  \item[(iii)]\SGARefTarget{sga3:iv:lemma:3:2:2:item:iii:01} \textit{there exists a necessarily unique morphism \(f\)
  making the diagram}
  \[
    \begin{array}{ccc}
      H\times Y & \xrightarrow{f} & Y\\
      \big\downarrow && \big\downarrow\\
      H\times X & \longrightarrow & X
    \end{array}
  \]
  \textit{commute.}
\end{enumerate}
\textit{Under these conditions, \(f\) defines a morphism of
\(\widehat{\mathscr C}\)-groups}
\[
  H\longrightarrow\operatorname{Aut}(Y),
\]
\textit{and the equivalence relation in \(Y\) defined by this action of
\(H\) is precisely the equivalence relation induced on \(Y\) by the
equivalence relation in \(X\) defined by the action of \(H\).}

% Source page: Exp4 component-local p. 14; combined-reader p. 202; printed p. 190.
The proof is immediate. There is of course an analogous statement for a
right action. The action of \(H\) on \(Y\) defined above will be called the
action induced on \(Y\) by the given action of \(H\) on \(X\).

Consider now two \(\mathscr C\)-groups \(H\) and \(G\), together with a
group morphism
\[
  u:H\longrightarrow G.
\]
Then \(H\) acts on \(G\) by translations (set-theoretically,
\(hg=u(h)g\)), and it acts freely if and only if \(u\) is a monomorphism.
When it exists, the quotient of \(G\) by this action of \(H\) is denoted
\(H\backslash G\). Similarly, one defines a right action of \(H\) on \(G\)
and a quotient \(G/H\). These quotients are functorial with respect to the
groups involved. More precisely, one has \SGARefLink{sga3:iv:lemma:3:2:3}{the following lemma}, stated for
right quotients.

\medskip
\noindent\SGARefTarget{sga3:iv:lemma:3:2:3}\textbf{Lemma 3.2.3.}
\textit{Let \(u:H\to G\) and \(u':H'\to G'\) be monomorphisms of
\(\mathscr C\)-groups, and suppose that a morphism of \(\mathscr C\)-groups
\[
  f:G\longrightarrow G'
\]
is given. The following conditions are equivalent:}
\begin{enumerate}
  \item[(i)]\SGARefTarget{sga3:iv:lemma:3:2:3:item:i:01} \textit{\(f\) is compatible with the equivalence relations in
  \(G\) and \(G'\) defined by \(H\) and \(H'\);}
  \item[(ii)]\SGARefTarget{sga3:iv:lemma:3:2:3:item:ii:01} \textit{for every \(S\in\operatorname{Ob}\mathscr C\), one
  has \(fu(H(S))\subseteq u'(H'(S))\);}
  \item[(iii)]\SGARefTarget{sga3:iv:lemma:3:2:3:item:iii:01} \textit{there exists a necessarily unique multiplicative
  morphism \(g:H\to H'\) such that the diagram}
  \[
    \begin{array}{ccc}
      H & \xrightarrow{g} & H'\\
      {\scriptstyle u}\big\downarrow && \big\downarrow{\scriptstyle u'}\\
      G & \xrightarrow{f} & G'
    \end{array}
  \]
  \textit{commutes.}
\end{enumerate}
\textit{Under these conditions, if the quotients \(G/H\) and \(G'/H'\)
exist, there is a unique morphism \(\overline f\) making the diagram}
\[
  \begin{array}{ccc}
    G & \xrightarrow{f} & G'\\
    {\scriptstyle p}\big\downarrow && \big\downarrow{\scriptstyle p'}\\
    G/H & \xrightarrow{\overline f} & G'/H'
  \end{array}
\]
\textit{commute.}

The first assertion is proved by reduction to the set-theoretic case. The
second follows immediately from \SGARefLink{sga3:iv:lemma:3:2:3:item:i:01}{(i)}.

The notions introduced above for general equivalence relations could all be
translated into the present setting. We record only \SGARefLink{sga3:iv:lemma:3:2:4}{the following lemma},
whose proof is immediate by reduction to the set-theoretic case.

% Source page: Exp4 component-local p. 15; combined-reader p. 203; printed p. 191.
\medskip
\noindent\SGARefTarget{sga3:iv:lemma:3:2:4}\textbf{Lemma 3.2.4.}
\textit{Let \(u:H\to G\) be a monomorphism of \(\mathscr C\)-groups, and
let \(G'\) be a \(\mathscr C\)-subgroup of \(G\). The subobject \(G'\) of
\(G\) is stable under the equivalence relation defined by \(H\) if and only
if \(u\) factors through the canonical monomorphism \(G'\to G\). Under \SGARefLink{sga3:iv:lemma:3:2:4}{this
condition}, the action of \(H\) on \(G'\) induced by the action on \(G\)
defined by \(u\) is precisely the action arising from the monomorphism
\(H\to G'\) through which \(u\) factors.}

\SGARefTarget{sga3:iv:subsection:3:3}\subsection*{3.3. Universally Effective Equivalence Relations}
\addcontentsline{toc}{subsection}{\SGARefLink{sga3:iv:subsection:3:3}{3.3}. Universally Effective Equivalence Relations}

\noindent\SGARefTarget{sga3:iv:definition:3:3:1}\textbf{Definition 3.3.1.}
\textit{Let \(f:X\to Y\) be a morphism. The \emph{equivalence relation
defined by \(f\) in \(X\)}, denoted \(\mathscr R(f)\), is the
\(\widehat{\mathscr C}\)-equivalence relation in \(X\) that is the image of
the canonical monomorphism}
\[
  X\times_Y X\longrightarrow X\times X.
\]

\medskip
\noindent\SGARefTarget{sga3:iv:definition:3:3:2}\textbf{Definition 3.3.2.}
\textit{An equivalence relation \(R\) in \(X\) is said to be
\emph{effective} if:}
\begin{enumerate}
  \item[(i)]\SGARefTarget{sga3:iv:definition:3:3:2:item:i:01} \textit{\(R\) is representable, that is, it is a
  \(\mathscr C\)-equivalence relation;}
  \item[(ii)]\SGARefTarget{sga3:iv:definition:3:3:2:item:ii:01} \textit{the quotient \(Y=X/R\) exists in
  \(\mathscr C\);}\footnote{Editor's note: the words ``in
  \(\mathscr C\)'' have been added.}
  \item[(iii)]\SGARefTarget{sga3:iv:definition:3:3:2:item:iii:01} \textit{the diagram}
  \[
    R\overset{p_1}{\underset{p_2}{\rightrightarrows}}
    X\xrightarrow{p}Y
  \]
  \textit{makes \(R\) the fiber square of \(X\) over \(Y\), that is,
  \(R\) is the equivalence relation defined by \(p\).}
\end{enumerate}

\medskip
\noindent\SGARefTarget{sga3:iv:scholium:3:3:2:1}\textbf{Scholium 3.3.2.1.}\footnote{Editor's note: the numbering
\SGARefLink{sga3:iv:scholium:3:3:2:1}{3.3.2.1} (respectively \SGARefLink{sga3:iv:scholium:3:3:3:1}{3.3.3.1}) has been added for later reference; Remark
\SGARefLink{sga3:iv:remark:3:3:3:2}{3.3.3.2} has also been added.}
\textit{If \(R\) is an effective equivalence relation in \(X\), then \(p\)
is an effective epimorphism (\SGARefLink{sga3:iv:definition:1:3}{1.3}). If \(f:X\to Y\) is an effective
epimorphism, then \(\mathscr R(f)\) is an effective equivalence relation in
\(X\), one quotient of which is \(Y\). Thus there is a one-to-one
``Galois'' correspondence between effective equivalence relations in \(X\)
and effective quotients of \(X\), that is, equivalence classes of effective
epimorphisms with source \(X\).}

\medskip
\noindent\SGARefTarget{sga3:iv:definition:3:3:3}\textbf{Definition 3.3.3.}
\textit{An equivalence relation \(R\) in \(X\) is said to be
\emph{universally effective} if the quotient \(Y=X/R\) exists and, for
every \(Y'\to Y\), the fiber products \(X'=X\times_Y Y'\) and
\(R'=R\times_Y Y'\) exist and \(R'\) is the fiber square of \(X'\) over
\(Y'\). Equivalently, \(R\) is effective and \(p:X\to X/R\) is a universal
effective epimorphism.}

\medskip
\noindent\SGARefTarget{sga3:iv:scholium:3:3:3:1}\textbf{Scholium 3.3.3.1.}\footnotemark[14]
\textit{As above, there is a one-to-one correspondence between universally
effective equivalence relations in \(X\) and universal effective quotients
of \(X\).}

% Source page: Exp4 component-local p. 16; combined-reader p. 204; printed p. 192.
\medskip
\noindent\SGARefTarget{sga3:iv:remark:3:3:3:2}\textbf{Remark 3.3.3.2.}\footnotemark[14]
Suppose that \(\mathscr C\) is the category of \(S\)-schemes, and write
\(\mathbf A^1_S\) for affine one-space over \(S\). Let
\(R\subset X\times_S X\) be a universally effective equivalence relation,
and let \(p:X\to Y\) be its quotient. Then, for every open subset \(U\) of
\(Y\),
\[
  \mathcal O(U)=\operatorname{Hom}_S(U,\mathbf A^1_S)
\]
is the set of elements \(\varphi\) of
\[
  \mathcal O(p^{-1}(U))
  =\operatorname{Hom}_S(p^{-1}(U),\mathbf A^1_S)
\]
such that \(\varphi\circ\operatorname{pr}_1
=\varphi\circ\operatorname{pr}_2\). In particular, if \(R\) is defined by
the action of a group \(H\) acting freely on the right on \(X\) (cf. \SGARefLink{sga3:iv:definition:3:2:1}{3.2.1}),
then \(\mathcal O(U)\) is the set of all
\(\varphi\in\mathcal O(p^{-1}(U))\) such that
\[
  \varphi(xh)=\varphi(x)
\]
for every \(S'\to S\), \(x\in X(S')\), and \(h\in H(S')\).

\medskip
\noindent\SGARefTarget{sga3:iv:proposition:3:3:4}\textbf{Proposition 3.3.4.}
\textit{Let \(R\) be a universally effective equivalence relation in
\(X\). Let \(f:X\to Z\) be a morphism compatible with \(R\), hence
factoring through a morphism \(g:X/R\to Z\). The following conditions are
equivalent:}
\begin{enumerate}
  \item[(i)]\SGARefTarget{sga3:iv:proposition:3:3:4:item:i:01} \textit{\(g\) is a monomorphism;}
  \item[(ii)]\SGARefTarget{sga3:iv:proposition:3:3:4:item:ii:01} \textit{\(R\) is the equivalence relation defined by \(f\).}
\end{enumerate}

Indeed, \SGARefLink{sga3:iv:proposition:3:3:4:item:i:01}{(i)} plainly implies \SGARefLink{sga3:iv:proposition:3:3:4:item:ii:01}{(ii)}, while the converse is precisely \SGARefLink{sga3:iv:corollary:2:5}{2.5}.

\medskip
\noindent\SGARefTarget{sga3:iv:definition:3:3:5}\textbf{Definition 3.3.5.}
\textit{Let \(H\) be a \(\mathscr C\)-group acting freely on \(X\). One
says that \(H\) acts \emph{effectively}, or that the action of \(H\) on
\(X\) is \emph{effective} (respectively \emph{universally effective}), if
the equivalence relation in \(X\) defined by the action of \(H\) is
effective (respectively universally effective).}

\SGARefTarget{sga3:iv:subsection:3:4}\subsection*{3.4. \((M)\)-Effectivity}
\addcontentsline{toc}{subsection}{\SGARefLink{sga3:iv:subsection:3:4}{3.4}. (M)-Effectivity}

In practice, universal effective epimorphisms are most often difficult to
characterize. Nevertheless, one frequently has at hand a number of
morphisms of this kind--in scheme theory, for example, faithfully flat
quasi-compact morphisms. This leads to the developments below.

\medskip
\noindent\SGARefTarget{sga3:iv:para:3:4:1}\textbf{3.4.1.}
We first state several conditions on a family \((M)\) of morphisms of
\(\mathscr C\):
\begin{enumerate}
  \item[(a)]\SGARefTarget{sga3:iv:para:3:4:1:item:a:01} \((M)\) is stable under base change: every
  \(u:T\to S\) belonging to \((M)\) is squarable (cf. \SGARefLink{sga3:iv:definition:1:4:0}{1.4.0}), and for every
  \(S'\to S\), the base change
  \[
    u':T\times_S S'\longrightarrow S'
  \]
  belongs to \((M)\);
  \item[(b)]\SGARefTarget{sga3:iv:para:3:4:1:item:b:01} the composite of two elements of \((M)\) belongs to \((M)\);
  \item[(c)]\SGARefTarget{sga3:iv:para:3:4:1:item:c:01} every isomorphism belongs to \((M)\);
  \item[(d)]\SGARefTarget{sga3:iv:para:3:4:1:item:d:01} every element of \((M)\) is an effective epimorphism.
\end{enumerate}

Conditions \SGARefLink{sga3:iv:para:3:4:1:item:a:01}{(a)} and \SGARefLink{sga3:iv:para:3:4:1:item:b:01}{(b)} imply:
\begin{enumerate}
  \item[(a$'$)]\SGARefTarget{sga3:iv:para:3:4:1:item:a:prime:01} the cartesian product of two elements of \((M)\) belongs
  to \((M)\). More precisely, let \(u:X\to Y\) and \(u':X'\to Y'\) be two
  \(S\)-morphisms belonging to \((M)\). If \(Y\times_S Y'\) exists, then
  \(X\times_S X'\) exists and \(u\times_Su'\) belongs to \((M)\).
\end{enumerate}

% Source page: Exp4 component-local p. 17; combined-reader p. 205; printed p. 193.
This follows from the diagram
\[
\begin{array}{ccccc}
  &&Y'&\xleftarrow{u'}&X'\\
  &&\big\uparrow&&\big\uparrow\\[-2pt]
  X&\longleftarrow&X\times_S Y'&\longleftarrow&X\times_S X'\\
  {\scriptstyle u}\big\downarrow&&\big\downarrow&&
    \underset{\scriptstyle u\times_Su'}{\swarrow}\\[-2pt]
  Y&\longleftarrow&Y\times_S Y'&&
\end{array}
\]
Likewise, \SGARefLink{sga3:iv:para:3:4:1:item:a:01}{(a)} and \SGARefLink{sga3:iv:para:3:4:1:item:d:01}{(d)} imply:
\begin{enumerate}
  \item[(d$'$)]\SGARefTarget{sga3:iv:para:3:4:1:item:d:prime:01} every element of \((M)\) is a universal effective
  epimorphism.
\end{enumerate}

\medskip
\noindent\SGARefTarget{sga3:iv:para:3:4:2}\textbf{3.4.2.}
The family \((M_0)\) of universal effective epimorphisms satisfies
conditions \SGARefLink{sga3:iv:para:3:4:1:item:a:01}{(a)}--\SGARefLink{sga3:iv:para:3:4:1:item:d:01}{(d)} of \SGARefLink{sga3:iv:para:3:4:1}{3.4.1}. Indeed, \SGARefLink{sga3:iv:para:3:4:1:item:a:01}{(a)}, \SGARefLink{sga3:iv:para:3:4:1:item:c:01}{(c)}, and \SGARefLink{sga3:iv:para:3:4:1:item:d:01}{(d)} hold by definition,
and \SGARefLink{sga3:iv:para:3:4:1:item:b:01}{(b)} follows from \SGARefLink{sga3:iv:proposition:1:8}{1.8}. Henceforth we assume given a family \((M)\) of
morphisms of \(\mathscr C\) satisfying these conditions; our results apply
in particular to the family \((M_0)\).

\medskip
\noindent\SGARefTarget{sga3:iv:definition:3:4:3}\textbf{Definition 3.4.3.}
An equivalence relation \(R\) in \(X\) is said to be \emph{of type \((M)\)}
if it is representable and \(p_1\in(M)\); conditions \SGARefLink{sga3:iv:para:3:4:1:item:b:01}{(b)} and \SGARefLink{sga3:iv:para:3:4:1:item:c:01}{(c)} then imply
that \(p_2\in(M)\).

The relation \(R\) is said to be \emph{\((M)\)-effective} if it is effective
and the canonical morphism \(X\to X/R\) belongs to \((M)\).

A quotient \(Y\) of \(X\) is said to be \emph{\((M)\)-effective} if the
canonical morphism \(X\to Y\) belongs to \((M)\).

The following consequences are immediate from \SGARefLink{sga3:iv:definition:3:4:3}{the definition}.\footnote{Editor's
note: the numbering \SGARefLink{sga3:iv:proposition:3:4:3:1}{3.4.3.1} has been added for later reference, and the proof
of part \SGARefLink{sga3:iv:proposition:3:4:3:1:item:i:01}{(i)} has been expanded.}

\medskip
\noindent\SGARefTarget{sga3:iv:proposition:3:4:3:1}\textbf{Proposition 3.4.3.1.}
\begin{enumerate}
  \item[(i)]\SGARefTarget{sga3:iv:proposition:3:4:3:1:item:i:01} \textit{An \((M)\)-effective equivalence relation is of type
  \((M)\) and is universally effective.}
  \item[(ii)]\SGARefTarget{sga3:iv:proposition:3:4:3:1:item:ii:01} \textit{An \((M)\)-effective quotient is universally effective
  (cf. \SGARefLink{sga3:iv:definition:3:3:3}{3.3.3}).}
  \item[(iii)]\SGARefTarget{sga3:iv:proposition:3:4:3:1:item:iii:01} \textit{The maps \(R\mapsto X/R\) and
  \(p\mapsto\mathscr R(p)\) establish a one-to-one correspondence between
  \((M)\)-effective equivalence relations in \(X\) and
  \((M)\)-effective quotients of \(X\).}
  \item[(iv)]\SGARefTarget{sga3:iv:proposition:3:4:3:1:item:iv:01} \textit{To be \((M_0)\)-effective is equivalent to being
  universally effective.}
\end{enumerate}

We prove \SGARefLink{sga3:iv:proposition:3:4:3:1:item:i:01}{(i)}. Since \(R\) is \((M)\)-effective, there is a cartesian square
\[
  \begin{array}{ccc}
    R & \xrightarrow{p_2} & X\\
    {\scriptstyle p_1}\big\downarrow && \big\downarrow{\scriptstyle p}\\
    X & \xrightarrow{p} & X/R,
  \end{array}
\]
and \(p\in(M)\). By \SGARefLink{sga3:iv:para:3:4:1:item:a:01}{3.4.1(a)}, \(p_1\) and \(p_2\) belong to \((M)\), so
\(R\) is of type \((M)\).

% Source page: Exp4 component-local p. 18; combined-reader p. 206; printed p. 194.
Set \(Y=X/R\), and let \(Y'\to Y\) be any morphism. By \SGARefLink{sga3:iv:para:3:4:1:item:a:01}{3.4.1(a)}, the fiber
products \(X'=X\times_Y Y'\) and \(R'=R\times_Y Y'\) exist, and the
morphisms \(X'\to Y'\) and \(p_i':R'\to X'\) belong to \((M)\). Finally,
since \(R=X\times_Y X\), associativity of the fiber product gives
\[
  R'=X\times_Y X\times_Y Y'=X'\times_{Y'}X'.
\]
This shows that \(R'\) is \((M)\)-effective; in particular, \(R\) is
universally effective. This proves \SGARefLink{sga3:iv:proposition:3:4:3:1:item:i:01}{(i)}, as well as \SGARefLink{sga3:iv:proposition:3:4:3:1:item:iv:01}{(iv)}. Parts \SGARefLink{sga3:iv:proposition:3:4:3:1:item:ii:01}{(ii)} and \SGARefLink{sga3:iv:proposition:3:4:3:1:item:iii:01}{(iii)}
follow, taking Definition \SGARefLink{sga3:iv:definition:3:3:2}{3.3.2} into account.

\medskip
\noindent\SGARefTarget{sga3:iv:para:3:4:4}\textbf{3.4.4.}
Let \(H\) be an \(S\)-group whose structure morphism belongs to \((M)\). If
\(H\) acts freely on the \(S\)-object \(X\), then it defines an equivalence
relation of type \((M)\) in \(X\). Indeed, by \SGARefLink{sga3:iv:para:3:4:1:item:a:01}{(a)}, the fiber product
\(H\times_S X\) exists and
\(\operatorname{pr}_2:H\times_S X\to X\) belongs to \((M)\). The action of
\(H\) is said to be \emph{\((M)\)-effective} if the equivalence relation in
\(X\) defined by this action is \((M)\)-effective.

\medskip
\noindent\SGARefTarget{sga3:iv:proposition:3:4:5}\textbf{Proposition 3.4.5 (\((M)\)-effectivity and base change).}
\textit{Let \(R\) be an \((M)\)-effective equivalence relation in \(X\)
over \(S\), and set \(Y=X/R\). Let \(S'\to S\) be a base change for which
\(Y'=Y\times_S S'\) exists. Then \(X'=X\times_S S'\) exists,
\(R'=R\times_S S'\) is an \((M)\)-effective equivalence relation in \(X'\)
over \(S'\), and}
\[
  X'/R'\simeq (X/R)'.
\]

Indeed, the canonical morphisms \(X\to Y\) and \(R\to Y\) belong to
\((M)\); hence, by \SGARefLink{sga3:iv:para:3:4:1:item:a:prime:01}{(a$'$)}, \(X'\) and \(R'\) are representable. By
associativity of the fiber product, \(R'\) is the equivalence relation in
\(X'\) defined by the canonical morphism \(X'\to Y'\), which belongs to
\((M)\). This proves \SGARefLink{sga3:iv:proposition:3:4:5}{the assertion}.

\medskip
\noindent\SGARefTarget{sga3:iv:proposition:3:4:6}\textbf{Proposition 3.4.6 (\((M)\)-effectivity and cartesian
products).}
\textit{Let \(R\) (respectively \(R'\)) be an \((M)\)-effective equivalence
relation in \(X\) (respectively \(X'\)) over \(S\). If
\((X/R)\times_S(X'/R')\) exists, then \(X\times_SX'\) exists,
\(R\times_SR'\) is an \((M)\)-effective equivalence relation in
\(X\times_SX'\) over \(S\), and}
\[
  (X\times_SX')/(R\times_SR')
  \simeq (X/R)\times_S(X'/R').
\]

Set \(Y=X/R\) and \(Y'=X'/R'\). By \SGARefLink{sga3:iv:para:3:4:1:item:a:prime:01}{(a$'$)}, the fiber product
\(X\times_SX'\) exists and the canonical morphism
\[
  q:X\times_SX'\longrightarrow Y\times_SY'
\]
belongs to \((M)\). The formula
\[
  (X\times_SX')\times_{(Y\times_SY')}(X\times_SX')
  \simeq (X\times_YX)\times_S(X'\times_{Y'}X')
\]
(all unlabelled fiber products being taken over \(S\)) shows that
\(R\times_SR'\) is the equivalence relation in \(X\times_SX'\) defined by
\(q\), which completes the proof.

% Source page: Exp4 component-local p. 19; combined-reader p. 207; printed p. 195.
\medskip
\noindent\SGARefTarget{sga3:iv:para:3:4:7}\textbf{3.4.7.}\footnote{Editor's note: \SGARefLink{sga3:iv:para:3:4:7}{this paragraph} has been
added.}
Suppose that \(\mathscr C\) has a final object \(e\), and let
\(f:G\to G'\) be a morphism of \(\mathscr C\)-groups belonging to \((M)\).
By \SGARefLink{sga3:iv:para:3:4:1:item:a:01}{3.4.1(a)}, the kernel \(\operatorname{Ker}(f)\) is represented by
\(e\times_{G'}G\), and the morphism \(\operatorname{Ker}(f)\to e\) belongs
to \((M)\).

On the other hand, the equivalence relation defined by \(f\) is the same as
that defined by the action of \(\operatorname{Ker}(f)\), say on the right,
on \(G\): it is the image of the morphism
\[
  G\times\operatorname{Ker}(f)\longrightarrow G\times G
\]
defined on sets by \((g,h)\mapsto(g,gh)\). Consequently, Scholium \SGARefLink{sga3:iv:scholium:3:3:2:1}{3.3.2.1}
gives \SGARefLink{sga3:iv:corollary:3:4:7:1}{the following corollary}.

\medskip
\noindent\SGARefTarget{sga3:iv:corollary:3:4:7:1}\textbf{Corollary 3.4.7.1.}
\textit{Suppose that \(\mathscr C\) has a final object \(e\), and let
\(f:G\to G'\) be a morphism of \(\mathscr C\)-groups belonging to \((M)\).
Then the action of \(\operatorname{Ker}(f)\) on \(G\) is
\((M)\)-effective, and \(G'\) is the quotient
\(G/\operatorname{Ker}(f)\).}

\SGARefTarget{sga3:iv:subsection:3:5}\subsection*{3.5. Construction of Quotients by Descent}
\addcontentsline{toc}{subsection}{\SGARefLink{sga3:iv:subsection:3:5}{3.5}. Construction of Quotients by Descent}

It frequently happens that one cannot construct a quotient directly, but
can do so after a suitable base change. This subsection gives a useful
criterion for that situation.

In \SGARefLink{sga3:iv:definition:2:1}{2.1} we defined a descent datum on an object \(X'\) over \(S'\), relative
to a morphism \(S'\to S\).

\medskip
\noindent\SGARefTarget{sga3:iv:definition:3:5:1}\textbf{Definition 3.5.1.}
\textit{A descent datum on \(X'\), relative to \(S'\to S\), is said to be
\emph{effective} if \(X'\), equipped with this descent datum, is isomorphic
to the inverse image over \(S'\) of an object \(X\) over \(S\), equipped
with its canonical descent datum.}

If \(S'\to S\) is a descent morphism (\SGARefLink{sga3:iv:definition:2:2}{2.2}), then the object \(X\) in \SGARefLink{sga3:iv:definition:3:5:1}{the
definition} is unique up to unique isomorphism. To say that \(S'\to S\) is
an effective descent morphism (\SGARefLink{sga3:iv:definition:2:2}{2.2}) means that it is a descent morphism and
that every descent datum relative to it is effective.

Now consider an equivalence relation \(R\) in an object \(X\) over \(S\).
Let \(X'\), \(X''\), and \(X'''\) be the inverse images of \(X\) over
\[
  S',\qquad S''=S'\times_SS',\qquad
  S'''=S'\times_SS'\times_SS',
\]
respectively, and let \(R'\), \(R''\), and \(R'''\) be the equivalence
relations obtained from \(R\) by inverse image. Suppose that \(R'\) in
\(X'\) is \((M)\)-effective, and consider the quotient
\(Y'=X'/R'\), an object over \(S'\). By \SGARefLink{sga3:iv:proposition:3:4:5}{3.4.5}, its two inverse images over
\(S''\) are isomorphic to \(X''/R''\). Thus the \(S'\)-object \(Y'\) has a
canonical gluing datum. Applying in the same way the uniqueness of
\(X'''/R'''\), one sees that this is in fact a descent datum. (In this
argument it has been tacitly assumed that all the displayed fiber products
exist. This holds, in particular, if \(S'\to S\) is squarable, for example
if it is a descent morphism.)

\medskip
\noindent\SGARefTarget{sga3:iv:proposition:3:5:2}\textbf{Proposition 3.5.2.}
\textit{Let \(R\) be an equivalence relation in the object \(X\) over
\(S\), and let \(S'\to S\) be a universal effective epimorphism. Assume
that every \(S\)-morphism whose inverse image over \(S'\) belongs to
\((M)\) itself belongs to \((M)\). The following conditions are equivalent:}
\begin{enumerate}
  \item[(i)]\SGARefTarget{sga3:iv:proposition:3:5:2:item:i:01} \textit{\(R\) is \((M)\)-effective in \(X\);}
  % Source page: Exp4 component-local p. 20; combined-reader p. 208; printed p. 196.
  \item[(ii)]\SGARefTarget{sga3:iv:proposition:3:5:2:item:ii:01} \textit{\(R'\) is \((M)\)-effective in \(X'\), and the
  canonical descent datum on \(X'/R'\) is effective.}
\end{enumerate}
\textit{When these conditions hold, the object descended from \(X'/R'\) is
canonically isomorphic to \(X/R\).}

The implication \SGARefLink{sga3:iv:proposition:3:5:2:item:i:01}{(i)}\(\Rightarrow\)\SGARefLink{sga3:iv:proposition:3:5:2:item:ii:01}{(ii)} is simply \SGARefLink{sga3:iv:para:3:4:4}{3.4.4} translated into the
language of descent. Once the converse is proved, the last assertion of \SGARefLink{sga3:iv:proposition:3:5:2}{the
proposition} follows from the fact that a universal effective epimorphism is
a descent morphism (\SGARefLink{sga3:iv:proposition:2:3}{2.3}), and hence the descended object is unique up to
unique isomorphism.

We prove \SGARefLink{sga3:iv:proposition:3:5:2:item:ii:01}{(ii)}\(\Rightarrow\)\SGARefLink{sga3:iv:proposition:3:5:2:item:i:01}{(i)}. Let \(Y'=X'/R'\), and let \(Y\) be the
descended object. The canonical morphism
\(p':X'\to X'/R'=Y'\) is compatible by construction with the descent data:
its two inverse images over \(S''\) coincide with the canonical morphism
\(X''\to X''/R''\). It is therefore the inverse image over \(S'\) of an
\(S\)-morphism \(p:X\to Y\). Since \(p'\) belongs to \((M)\), the assumption
on \(S'\to S\) implies that \(p\) also belongs to \((M)\). Moreover,
because \(p'\) is compatible with the equivalence relation \(R'\), the
morphism \(p\) is compatible with \(R\), again because a universal
effective epimorphism is a descent morphism. Hence there is a morphism
\[
  R\longrightarrow X\times_YX.
\]
To prove that \(R\) is \((M)\)-effective and that \(Y\) is isomorphic to
\(X/R\), it suffices to prove that this morphism is an isomorphism. After
base change from \(S\) to \(S'\), it becomes an isomorphism because \(R'\)
is effective; it is therefore an isomorphism for the same reason as in \SGARefLink{sga3:iv:proposition:2:4}{2.4}.

Notice that the hypothesis in \SGARefLink{sga3:iv:proposition:3:5:2}{the proposition} is satisfied if \((M)\) is
the family \((M_0)\) of universal effective epimorphisms and if
\(\mathscr C\) has fiber products (\SGARefLink{sga3:iv:corollary:1:10}{1.10}). We obtain the following.

\medskip
\noindent\SGARefTarget{sga3:iv:corollary:3:5:3}\textbf{Corollary 3.5.3.}
\textit{Suppose that \(\mathscr C\) has fiber products (fiber products over
\(S\) would suffice). Let \(R\) be an equivalence relation in \(X\) over
\(S\), and let \(S'\to S\) be a universal effective epimorphism. The
following conditions are equivalent:}
\begin{enumerate}
  \item[(i)]\SGARefTarget{sga3:iv:corollary:3:5:3:item:i:01} \textit{\(R\) is universally effective in \(X\);}
  \item[(ii)]\SGARefTarget{sga3:iv:corollary:3:5:3:item:ii:01} \textit{\(R'\) is universally effective in \(X'\), and the
  canonical descent datum on \(X'/R'\) is effective.}
\end{enumerate}
\textit{When these conditions hold, the object descended from \(X'/R'\) is
canonically isomorphic to \(X/R\).}

% End of assigned scope. The Section 4 heading begins immediately here in
% the authority, on Exp4 component-local p. 20 (combined-reader p. 208).

% END INCLUDED SGA3_English_current_buildable_source_20260823/source/inputs/IV/tex/components/03_expose_IV_s3.tex

% BEGIN INCLUDED SGA3_English_current_buildable_source_20260823/source/inputs/IV/tex/components/04_expose_IV_s41_s43.tex
% SGA 3, Expose IV, English reconstruction.
% Scope: section 4 heading and sections 4.1--4.3 in full, stopping
% immediately before the section 4.4 heading.
% Authority: Polo--Gille Exp4-13oct24.pdf, component-local pages 20 (partial)
% through 35 (partial), combined-reader pages 208--223, printed pages 196--211.
% Jacob Reinhold's CC BY 4.0 Markdown lineage was used only as a phrasing
% comparison. The Polo--Gille PDF controls prose, formulas, numbering, notes,
% and diagrams.

% Source page: Exp4 component-local p. 20 (partial); printed p. 196.
\SGARefTarget{sga3:iv:section:4}\section*{4. Topologies and Sheaves}
\addcontentsline{toc}{section}{4. Topologies and Sheaves}

The notion of a sieve, and the presentation of the notion of a topology
adopted here (\SGARefLink{sga3:iv:definition:4:2:1}{4.2.1}), which is more intrinsic and in many respects more
convenient than the presentation by covering families in \SGARefLink{sga3:iv:bib:ma}{[MA]}, are due to
J.~Giraud \SGARefLink{sga3:iv:bib:as}{[AS]}.

\SGARefTarget{sga3:iv:subsection:4:1}\subsection*{4.1. Sieves}
\addcontentsline{toc}{subsection}{\SGARefLink{sga3:iv:subsection:4:1}{4.1}. Sieves}

\medskip
\noindent\SGARefTarget{sga3:iv:definition:4:1:1}\textbf{Definition 4.1.1.}
\textit{A \emph{sieve} of the category \(\mathscr C\) is a subfunctor
\(C\) of the final functor}
\[
  e:\mathscr C^{\circ}\longrightarrow\mathbf{Set}.
\]

% Source page: Exp4 component-local p. 21; printed p. 197.
To every sieve \(C\) of \(\mathscr C\) one associates the set \(E(C)\) of
objects \(X\) of \(\mathscr C\) such that \(C(X)\ne\varnothing\), that is,
such that the structural morphism \(X\to e\) factors through \(C\). Thus
\[
\tag{+}\label{sga3:iv:definition:4:1:1:eq:unnamed:01}
\left\{
\begin{aligned}
 X\in E(C)&\quad\Longleftrightarrow\quad C(X)=e(X)=\{\varnothing\},\\
 X\notin E(C)&\quad\Longleftrightarrow\quad C(X)=\varnothing.
\end{aligned}
\right.
\]
The set \(E=E(C)\) has \SGARefLink{sga3:iv:definition:4:1:1:eq:unnamed:02}{the following property}:
\[
\tag{++}\label{sga3:iv:definition:4:1:1:eq:unnamed:02}
 X\in E\ \text{and}\ \operatorname{Hom}(Y,X)\ne\varnothing
 \quad\Longrightarrow\quad Y\in E.
\]
(If the set \(\operatorname{Ob}\mathscr C\) is endowed with its natural
preorder, with \(Y\) dominating \(X\) when there is an arrow from \(Y\) to
\(X\), then the sets \(E\) satisfying \SGARefLink{sga3:iv:definition:4:1:1:eq:unnamed:02}{(++)} are precisely the subsets of
\(\operatorname{Ob}\mathscr C\) that contain every object dominating one of
their elements.)\footnote{Editor's note: Here ``dominating'' is meant in the
sense of the preorder just mentioned: \(Y\) dominates \(X\) if there is an
arrow \(Y\to X\). On the other hand, if \(X,Y\) are two subobjects of an
object \(Z\), one says (cf. \SGARefLink{sga3:iv:proposition:2:4}{2.4}) that \(Y\) contains \(X\) if
\(X\subset Y\). To avoid any ambiguity between these usages, the sequel uses
``dominating'' in the first case and ``containing'' in the second.}

Conversely, if \(E\) is a subset of \(\operatorname{Ob}\mathscr C\) having
property \SGARefLink{sga3:iv:definition:4:1:1:eq:unnamed:02}{(++)}, then \(E\) can be written uniquely in the form \(E(C)\), with
\(C\) defined by the formulas \SGARefLink{sga3:iv:definition:4:1:1:eq:unnamed:01}{(+)}. Thus sieves of \(\mathscr C\) correspond
bijectively to subsets of \(\operatorname{Ob}\mathscr C\) satisfying \SGARefLink{sga3:iv:definition:4:1:1:eq:unnamed:02}{(++)}.
By abuse of language, we shall sometimes call the set \(E(C)\) itself a
sieve of \(\mathscr C\).

Let \(C\) and \(C'\) be two sieves of \(\mathscr C\).\footnote{Editor's
note: The following sentence has been expanded.} Since both are subfunctors
of the final functor \(e\), it is equivalent to say that \(C\) dominates
\(C'\) (for the domination relation in
\(\operatorname{Ob}\widehat{\mathscr C}\)), that \(C\) is a subfunctor of
\(C'\), or that \(E(C)\subset E(C')\). In this case we say that \(C\) is
\emph{finer} than \(C'\). This is an order relation on the set of sieves of
\(\mathscr C\). Moreover,
\[
 E(C)\cap E(C')=E(C\times C'),
\]
so the set of sieves of \(\mathscr C\) is filtered for the order relation
``is finer than.''

Every subset \(E\) of \(\operatorname{Ob}\widehat{\mathscr C}\), and in
particular every subset of \(\operatorname{Ob}\mathscr C\), defines a sieve
\(C(E)\): the set of \(X\in\operatorname{Ob}\mathscr C\) such that
\(F(X)\ne\varnothing\) for at least one \(F\in E\) satisfies \SGARefLink{sga3:iv:definition:4:1:1:eq:unnamed:02}{(++)}, and hence
defines the desired sieve.

This sieve may also be defined as the image of the family of morphisms
\(\{F\to e\mid F\in E\}\), in the sense of \SGARefLink{sga3:iv:definition:4:1:2}{the following definition}.

\medskip
\noindent\SGARefTarget{sga3:iv:definition:4:1:2}\textbf{Definition 4.1.2.}
\textit{Let \(\{F_i\to F\}\) be a family of morphisms of
\(\widehat{\mathscr C}\) with common target \(F\). The \emph{image} of this
family is the subfunctor of \(F\) defined by}
\[
 S\longmapsto\bigcup_i\operatorname{Im}F_i(S)\subset F(S).
\]

\medskip
\noindent\SGARefTarget{sga3:iv:proposition:4:1:3}\textbf{Proposition 4.1.3.}
\textit{Formation of the image commutes with base change. In the preceding
notation let \(I\) be the image of the family \(\{F_i\to F\}\). For every
morphism \(G\to F\) in \(\widehat{\mathscr C}\), the image of the family}
\[
 \{F_i\times_FG\longrightarrow G\}
\]
\textit{is the subfunctor \(I\times_FG\) of \(G\).}

% Source page: Exp4 component-local p. 22; printed p. 198.
\medskip
\noindent\SGARefTarget{sga3:iv:definition:4:1:4:0}\textbf{Definition 4.1.4.0.}\footnote{Editor's note: The numbering
\SGARefLink{sga3:iv:definition:4:1:4:0}{4.1.4.0} has been added to highlight \SGARefLink{sga3:iv:definition:4:1:4:0}{this definition}.}
\textit{Let \(C\) be a sieve of \(\mathscr C\). If \(E\) is a subset of
\(\operatorname{Ob}\mathscr C\) such that \(C(E)=C\), then \(E\) is called
a \emph{base} of \(C\). Every sieve \(C\) has a base, for example
\(E(C)\).}

We now describe \(\operatorname{Hom}(C,F)\), where \(C\) is a sieve of
\(\mathscr C\) and \(F\in\operatorname{Ob}\widehat{\mathscr C}\), by means
of a base \(\{S_i\}\) of \(C\). For every pair \((i,j)\) there is a
diagram in \(\widehat{\mathscr C}\):
\[
\begin{array}{ccc}
 S_i\times S_j&\longrightarrow&S_i\\
 \big\downarrow&&\big\downarrow\\
 S_j&\longrightarrow&e.
\end{array}
\]
It gives a diagram of sets
\[
 \Gamma(F)=\operatorname{Hom}(e,F)
 \xrightarrow{\sigma}
 \prod_i\operatorname{Hom}(S_i,F)
 \mathrel{\substack{\xrightarrow{\tau_1}\\[-0.5ex]
                     \xrightarrow[\tau_2]{}}}
 \prod_{i,j}\operatorname{Hom}(S_i\times S_j,F),
\]
with \(\tau_1\sigma=\tau_2\sigma\). Hence there is a morphism
\[
 \operatorname{Hom}(e,F)\longrightarrow
 \operatorname{Ker}\!\left(
 \prod_i\operatorname{Hom}(S_i,F)
 \rightrightarrows
 \prod_{i,j}\operatorname{Hom}(S_i\times S_j,F)
 \right).
\]
One checks immediately:

\medskip
\noindent\SGARefTarget{sga3:iv:proposition:4:1:4}\textbf{Proposition 4.1.4.}
\textit{There is an isomorphism, functorial in \(F\),}
\[
 \operatorname{Hom}(C,F)\xrightarrow{\sim}
 \operatorname{Ker}\!\left(
 \prod_i\operatorname{Hom}(S_i,F)
 \rightrightarrows
 \prod_{i,j}\operatorname{Hom}(S_i\times S_j,F)
 \right),
\]
\textit{for which the diagram}
\[
\begin{array}{ccc}
 \operatorname{Hom}(e,F)&\longrightarrow&
 \operatorname{Ker}\!\left(
 \prod_i\operatorname{Hom}(S_i,F)\rightrightarrows
 \prod_{i,j}\operatorname{Hom}(S_i\times S_j,F)\right)\\
 \Big\Vert&&\overset{\sim}{\Big\uparrow}\\
 \operatorname{Hom}(e,F)&\longrightarrow&\operatorname{Hom}(C,F)
\end{array}
\]
\textit{is commutative; the bottom arrow is induced by the canonical
morphism \(C\to e\).}

\medskip
\noindent\SGARefTarget{sga3:iv:corollary:4:1:5}\textbf{Corollary 4.1.5.}
\textit{Suppose that the fiber products \(S_i\times S_j\) are representable,
for example that the \(S_i\) are squarable. Then, for every
\(F\in\operatorname{Ob}\widehat{\mathscr C}\), there is an isomorphism}
\[
 \operatorname{Hom}(C,F)\xrightarrow{\sim}
 \operatorname{Ker}\!\left(
 \prod_iF(S_i)\rightrightarrows
 \prod_{i,j}F(S_i\times S_j)\right).
\]

% Source page: Exp4 component-local p. 23; printed p. 199.
\medskip
\noindent\SGARefTarget{sga3:iv:remark:4:1:6}\textbf{Remark 4.1.6.}
Let \(R\) be a sieve of \(\mathscr C\), let \(\overline R\) denote the full
subcategory of \(\mathscr C\) whose object set is \(E(R)\), and let
\[
 i_R:\overline R\longrightarrow\mathscr C
\]
be the inclusion functor. There is an isomorphism, functorial in
\(F\in\operatorname{Ob}\widehat{\mathscr C}\),
\[
 \operatorname{Hom}(R,F)\xrightarrow{\sim}\Gamma(F\circ i_R),
\]
for which the diagram
\[
\begin{array}{ccc}
 \operatorname{Hom}(e,F)&\longrightarrow&\operatorname{Hom}(R,F)\\
 \overset{\sim}{\Big\downarrow}&&
 \overset{\sim}{\Big\downarrow}\\
 \Gamma(F)&\longrightarrow&\Gamma(F\circ i_R)
\end{array}
\]
is commutative; the bottom arrow is induced by \(i_R\).

\medskip
\noindent\SGARefTarget{sga3:iv:definition:4:1:7}\textbf{Definition 4.1.7.}
\textit{Let \(\mathscr C\) be a category. A \emph{sieve of the object}
\(S\) of \(\mathscr C\) is a sieve of the category
\(\mathscr C_{/S}\).}

A sieve of \(S\) is therefore a subobject of \(S\) in
\(\widehat{\mathscr C}\). It corresponds canonically to a subset of
\(\operatorname{Ob}\mathscr C_{/S}\) that contains the source of every arrow
whose target it contains. By abuse of language, such a subset will also be
called a sieve of \(S\).

\SGARefTarget{sga3:iv:subsection:4:2}\subsection*{4.2. Topologies: Definitions}
\addcontentsline{toc}{subsection}{\SGARefLink{sga3:iv:subsection:4:2}{4.2}. Topologies: Definitions}

\medskip
\noindent\SGARefTarget{sga3:iv:definition:4:2:1}\textbf{Definition 4.2.1.}
\textit{A \emph{topology} on a category \(\mathscr C\) consists, for every
object \(S\) of \(\mathscr C\), of a set \(J(S)\) of sieves of \(S\), called
\emph{covering sieves} or \emph{refinements} of \(S\), satisfying the
following axioms:}
\begin{enumerate}
 \item[(T 1)]\SGARefTarget{sga3:iv:definition:4:2:1:item:t:1:01} \textit{For every refinement \(R\) of \(S\) and every
 morphism \(T\to S\), the sieve \(R\times_ST\) of \(T\) is covering
 (``stability under base change'').}
 \item[(T 2)]\SGARefTarget{sga3:iv:definition:4:2:1:item:t:2:01} \textit{If \(R,C\) are two sieves of \(S\), if \(R\) is
 covering, and if for every \(T\in\operatorname{Ob}\mathscr C\) and every
 morphism \(T\to R\), the sieve \(C\times_ST\) of \(T\) is covering, then
 \(C\) is a refinement of \(S\).}\footnote{Editor's note: In other words,
 if \(C\) is covering ``locally with respect to the covering sieve \(R\),''
 then \(C\) is covering.}
 \item[(T 3)]\SGARefTarget{sga3:iv:definition:4:2:1:item:t:3:01} \textit{If \(C\supset R\) are two sieves of \(S\) and
 \(R\) is covering, then \(C\) is covering.}
 \item[(T 4)]\SGARefTarget{sga3:iv:definition:4:2:1:item:t:4:01} \textit{For every \(S\), \(S\) is a refinement of \(S\).}
\end{enumerate}

% Source page: Exp4 component-local p. 24; printed p. 200.
These axioms may be reformulated as follows. Suppose that a topology
\(S\mapsto J(S)\) on \(\mathscr C\) is given. For every
\(F\in\operatorname{Ob}\widehat{\mathscr C}\), let \(J(F)\) be the set of
subfunctors \(R\) of \(F\) such that, for every morphism \(T\to F\) in
\(\widehat{\mathscr C}\) with \(T\) representable, the sieve
\(R\times_FT\) of \(T\) is covering. By \SGARefLink{sga3:iv:definition:4:2:1:item:t:1:01}{(T~1)}, this notation agrees with
the preceding one. An \(R\in J(F)\) is again called a \emph{refinement} of
\(F\). The preceding axioms immediately imply:
\begin{enumerate}
 \item[\((T'0)\)]\SGARefTarget{sga3:iv:definition:4:2:1:item:tprime0:01} If \(F\supset G\) are objects of
 \(\widehat{\mathscr C}\), and if
 \(G\times_FS\in J(S)\) for every \(S\in\operatorname{Ob}\mathscr C\) and
 every morphism \(S\to F\), then \(G\in J(F)\).
 \item[\((T'1)\)]\SGARefTarget{sga3:iv:definition:4:2:1:item:tprime1:01} If \(G\in J(F)\) and \(H\to F\) is a morphism in
 \(\widehat{\mathscr C}\), then \(G\times_FH\in J(H)\).
 \item[\((T'2)\)]\SGARefTarget{sga3:iv:definition:4:2:1:item:tprime2:01} If \(F\supset G\supset H\), if \(G\in J(F)\),
 and if \(H\in J(G)\), then \(H\in J(F)\).
 \item[\((T'3)\)]\SGARefTarget{sga3:iv:definition:4:2:1:item:tprime3:01} If \(F\supset G\supset H\) and \(H\in J(F)\),
 then \(G\in J(F)\).
 \item[\((T'4)\)]\SGARefTarget{sga3:iv:definition:4:2:1:item:tprime4:01} For every
 \(F\in\operatorname{Ob}\widehat{\mathscr C}\), one has \(F\in J(F)\).
\end{enumerate}
Conversely, if for every \(F\in\operatorname{Ob}\widehat{\mathscr C}\)
one is given a set \(J(F)\) of subobjects of \(F\) satisfying
\SGARefLink{sga3:iv:definition:4:2:1:item:tprime0:01}{\(T'0\)}--\SGARefLink{sga3:iv:definition:4:2:1:item:tprime4:01}{\(T'4\)}, then \(S\mapsto J(S)\) defines a topology on
\(\mathscr C\), and these two constructions are inverse to each other.

Properties \SGARefLink{sga3:iv:definition:4:2:1:item:tprime1:01}{\(T'1\)}, \SGARefLink{sga3:iv:definition:4:2:1:item:tprime2:01}{\(T'2\)}, and \SGARefLink{sga3:iv:definition:4:2:1:item:tprime3:01}{\(T'3\)} imply the further
property\footnote{Editor's note: \SGARefLink{sga3:iv:definition:4:2:1:item:tprime1:01}{\(T'1\)} and \SGARefLink{sga3:iv:definition:4:2:1:item:tprime2:01}{\(T'2\)} suffice:
\(G\cap H=G\times_FH\) belongs to \(J(H)\) by \SGARefLink{sga3:iv:definition:4:2:1:item:tprime1:01}{\(T'1\)}, and hence to
\(J(F)\) by \SGARefLink{sga3:iv:definition:4:2:1:item:tprime2:01}{\(T'2\)}.}:
\begin{enumerate}
 \item[\((T'5)\)]\SGARefTarget{sga3:iv:definition:4:2:1:item:tprime5:01} If \(G\) and \(H\) are two subobjects of \(F\) and
 \(G,H\in J(F)\), then \(G\cap H\in J(F)\).
\end{enumerate}
Thus \(J(F)\), ordered by \(\supset\), is filtered. This observation will
be useful below.

\medskip
\noindent\SGARefTarget{sga3:iv:para:4:2:2}\textbf{4.2.2.}
The topology defined by \(J\) is said to be \emph{finer} than the topology
defined by \(J'\) if
\(J(S)\supset J'(S)\) for every
\(S\in\operatorname{Ob}\mathscr C\). Equivalently,
\(J(F)\supset J'(F)\) for every
\(F\in\operatorname{Ob}\widehat{\mathscr C}\).

Every set of topologies on \(\mathscr C\) has a greatest lower bound. More
precisely, let \(I\) be an index set and, for each \(i\in I\), let
\(S\mapsto J_i(S)\) be a topology on \(\mathscr C\). Then
\[
 J(S)=\bigcap_{i\in I}J_i(S)
\]
defines a topology, and it is the greatest lower bound of the given set.

In particular, suppose that for every
\(S\in\operatorname{Ob}\mathscr C\) a set \(E(S)\) of sieves of \(S\) is
given. The \emph{topology generated by these sets} is the least fine
topology for which every element of \(E(S)\) is a refinement of \(S\).

\medskip
\noindent\SGARefTarget{sga3:iv:definition:4:2:3}\textbf{Definition 4.2.3.}
\textit{Let \(\{F_i\to F\}\) be a family of morphisms in
\(\widehat{\mathscr C}\), and let \(G\subset F\) be its image (\SGARefLink{sga3:iv:definition:4:1:2}{4.1.2}). The
family is \emph{covering} if \(G\in J(F)\). A morphism is covering if the
one-member family consisting of that morphism is covering.}

\SGARefLink{sga3:iv:definition:4:2:3}{This definition} applies in particular to an inclusion: a sieve \(C\) of
\(S\) is covering if and only if the canonical morphism \(C\to S\) is
covering.

% Source page: Exp4 component-local p. 25; printed p. 201.
Axioms \SGARefLink{sga3:iv:definition:4:2:1:item:tprime0:01}{\(T'0\)}--\SGARefLink{sga3:iv:definition:4:2:1:item:tprime5:01}{\(T'5\)} give the following properties of covering
families:
\begin{enumerate}
 \item[(C 0)]\SGARefTarget{sga3:iv:definition:4:2:3:item:c:0:01} If \(\{F_i\to F\}\) is a family in
 \(\widehat{\mathscr C}\) such that every representable base change
 \(\{F_i\times_FS\to S\}\) is covering, then the original family is
 covering.
 \item[(C 1)]\SGARefTarget{sga3:iv:definition:4:2:3:item:c:1:01} If \(\{F_i\to F\}\) is covering and \(G\to F\) is any
 morphism, then \(\{F_i\times_FG\to G\}\) is covering (``stability under
 base change'').
 \item[(C 2)]\SGARefTarget{sga3:iv:definition:4:2:3:item:c:2:01} If \(\{F_i\to F\}\) is covering and, for every \(i\),
 \(\{F_{ij}\to F_i\}\) is covering, then the composite family
 \(\{F_{ij}\to F\}\) is covering (``stability under composition'').
 \item[(C 3)]\SGARefTarget{sga3:iv:definition:4:2:3:item:c:3:01} If \(\{G_j\to F\}\) is covering and \(\{F_i\to F\}\) is
 a family such that for every \(j\) there is an \(i\) for which
 \(G_j\to F\) factors through \(F_i\to F\), then \(\{F_i\to F\}\) is
 covering (``saturation'').
 \item[(C 4)]\SGARefTarget{sga3:iv:definition:4:2:3:item:c:4:01} Every one-member family consisting of an isomorphism is
 covering.
\end{enumerate}
Notice that \SGARefLink{sga3:iv:definition:4:2:3:item:c:2:01}{(C~2)} and \SGARefLink{sga3:iv:definition:4:2:3:item:c:3:01}{(C~3)} also imply:
\begin{enumerate}
 \item[(C 5)]\SGARefTarget{sga3:iv:definition:4:2:3:item:c:5:01} If \(\{F_i\to F\}\) is a family for which there exists a
 covering family \(\{G_j\to F\}\) such that
 \(\{F_i\times_FG_j\to G_j\}\) is covering for every \(j\), then
 \(\{F_i\to F\}\) is covering (``a locally covering family is
 covering'').
\end{enumerate}

\medskip
\noindent\SGARefTarget{sga3:iv:para:4:2:4}\textbf{4.2.4.}
Conversely, let \(\mathscr C\) be a category having fiber products, and
suppose that for every \(S\in\operatorname{Ob}\mathscr C\) one is given a
set of families of morphisms in \(\mathscr C\) with target \(S\), called
covering families, satisfying \SGARefLink{sga3:iv:definition:4:2:3:item:c:1:01}{(C~1)}--\SGARefLink{sga3:iv:definition:4:2:3:item:c:4:01}{(C~4)}, and hence also \SGARefLink{sga3:iv:definition:4:2:3:item:c:5:01}{(C~5)}. Let
\(J(S)\) be the set of sieves of \(S\) having a covering base (or,
equivalently by \SGARefLink{sga3:iv:definition:4:2:3:item:c:3:01}{(C~3)}, all of whose bases are covering). Then
\(S\mapsto J(S)\) defines a topology on \(\mathscr C\), and the two
preceding constructions are inverse to one another.

In applications it is inconvenient to consider all covering families, since
one often has a fairly simple description of a ``sufficient'' number of
them. This leads to the following definitions.

\medskip
\noindent\SGARefTarget{sga3:iv:definition:4:2:5:0}\textbf{Definition 4.2.5.0.}\footnote{Editor's note: \SGARefLink{sga3:iv:definition:4:2:5:0}{This
definition}, which appeared after \SGARefLink{sga3:iv:definition:4:2:5}{4.2.5} in the original, has been placed here
because it will be used in \SGARefLink{sga3:iv:proposition:6:2:1}{6.2.1} in a slightly more general setting than
that of \SGARefLink{sga3:iv:definition:4:2:5}{4.2.5}.}
\textit{Let \(\mathscr C\) be a category, and suppose given, for every
\(S\in\operatorname{Ob}\mathscr C\), a set \(P(S)\) of families of
morphisms in \(\mathscr C\) with target \(S\). The \emph{topology generated
by \(P\)} is the least fine topology for which the given families are
covering.}

\medskip
\noindent\SGARefTarget{sga3:iv:definition:4:2:5}\textbf{Definition 4.2.5.}
\textit{A \emph{pretopology} on a category \(\mathscr C\) consists, for
every \(S\in\operatorname{Ob}\mathscr C\), of a set \(R(S)\) of
families \(\{S_i\to S\}\), called \emph{covering for the pretopology},
satisfying:}
\begin{enumerate}
 \item[(P 1)]\SGARefTarget{sga3:iv:definition:4:2:5:item:p:1:01} \textit{For every \(\{S_i\to S\}\in R(S)\) and
 every morphism \(T\to S\), the fiber products \(S_i\times_ST\) exist and
 \(\{S_i\times_ST\to T\}\in R(T)\).}
 \item[(P 2)]\SGARefTarget{sga3:iv:definition:4:2:5:item:p:2:01} \textit{If \(\{S_i\to S\}\in R(S)\) and, for each
 \(i\), \(\{T_{ij}\to S_i\}\in R(S_i)\), then the composite
 family \(\{T_{ij}\to S\}\) belongs to \(R(S)\).}
 \item[(P 3)]\SGARefTarget{sga3:iv:definition:4:2:5:item:p:3:01} \textit{Every one-member family consisting of an isomorphism
 is covering.}
\end{enumerate}

% Source page: Exp4 component-local p. 26; printed p. 202.
\medskip
\noindent\SGARefTarget{sga3:iv:proposition:4:2:6}\textbf{Proposition 4.2.6.}
\textit{For each \(S\), let \(J(S)\) be the set of sieves of \(S\) that
are covering for the topology generated by the pretopology \(R\).
Let \(J_R(S)\) be the subset of \(J(S)\) consisting of the
sieves defined by the families in \(R(S)\). Then
\(J_R(S)\) is cofinal in \(J(S)\): every refinement of \(S\)
contains a sieve defined by a family in \(R(S)\).}

For every \(S\), let \(J'(S)\) be the set of sieves of \(S\) containing a
sieve in \(J_R(S)\). Clearly \(J'(S)\subset J(S)\). To prove
\(J(S)=J'(S)\), it suffices to prove that the \(J'(S)\) define a topology,
that is, satisfy \SGARefLink{sga3:iv:definition:4:2:1:item:t:1:01}{(T~1)}--\SGARefLink{sga3:iv:definition:4:2:1:item:t:4:01}{(T~4)}. Axioms \SGARefLink{sga3:iv:definition:4:2:1:item:t:1:01}{(T~1)}, \SGARefLink{sga3:iv:definition:4:2:1:item:t:3:01}{(T~3)}, and \SGARefLink{sga3:iv:definition:4:2:1:item:t:4:01}{(T~4)} are immediate;
only \SGARefLink{sga3:iv:definition:4:2:1:item:t:2:01}{(T~2)} remains.\footnote{Editor's note: Typographical errors in the
original have been corrected in what follows.}

Let \(U\in J'(S)\), let \(C\) be a sieve of \(S\), and suppose that
\(C\times_ST\in J'(T)\) for every \(T\to U\). We must prove
\(C\in J'(S)\). By definition, \(U\) contains a refinement \(U'\) defined
by a family \(\{S_i\to S\}\in R(S)\). Since \SGARefLink{sga3:iv:definition:4:2:1:item:t:3:01}{(T~3)} has already been
verified, it is enough to prove \(U'\cap C\in J'(S)\); hence we may suppose
\(U=U'\). By hypothesis,
\(C\times_SS_i\in J'(S_i)\) for every \(i\). Thus for each \(i\) there is
a covering family \(\{T_{ij}\to S_i\}\in R(S_i)\) such that
\(T_{ij}\to S_i\) factors through \(C\times_SS_i\to S_i\). Consequently
\(T_{ij}\to S\) factors through \(C\to S\). Hence \(C\) contains the sieve
defined by the composite family \(\{T_{ij}\to S\}\), and \SGARefLink{sga3:iv:definition:4:2:5:item:p:2:01}{(P~2)} completes
the proof.

Axioms \SGARefLink{sga3:iv:definition:4:2:5:item:p:1:01}{(P~1)}--\SGARefLink{sga3:iv:definition:4:2:5:item:p:3:01}{(P~3)} are those of \SGARefLink{sga3:iv:bib:ma}{[MA]}. Because pretopologies are so useful
in practice, we shall interpret every important result in terms of a
pretopology defining the topology under consideration.

\medskip
\noindent\SGARefTarget{sga3:iv:remark:4:2:7}\textbf{Remark 4.2.7.}
One may introduce a slightly more general notion: for every \(S\), one
specifies a set of covering families satisfying \SGARefLink{sga3:iv:definition:4:2:5:item:p:1:01}{(P~1)}, \SGARefLink{sga3:iv:definition:4:2:5:item:p:3:01}{(P~3)}, and the
conclusion of Proposition~\SGARefLink{sga3:iv:proposition:4:2:6}{4.2.6}. This occurs in particular when the given
families satisfy \SGARefLink{sga3:iv:definition:4:2:5:item:p:1:01}{(P~1)}, \SGARefLink{sga3:iv:definition:4:2:5:item:p:3:01}{(P~3)}, and \SGARefLink{sga3:iv:definition:4:2:3:item:c:5:01}{(C~5)}. See \SGARefLink{sga3:iv:bib:d}{[D]}.

\medskip
\noindent\SGARefTarget{sga3:iv:definition:4:2:8}\textbf{Definition 4.2.8.}
\textit{Let \(\mathscr C\) be equipped with a topology, let \(S\) be an
object of \(\mathscr C\), and let \(P(S')\) be a condition on an argument
\(S'\in\operatorname{Ob}\mathscr C_{/S}\). Suppose that
\(\operatorname{Hom}(S'',S')\ne\varnothing\) implies
\(P(S')\Rightarrow P(S'')\). The condition \(P\) is said to hold
\emph{locally on \(S\)} for the topology under consideration if the
following equivalent conditions hold:}
\begin{enumerate}
 \item[(i)]\SGARefTarget{sga3:iv:definition:4:2:8:item:i:01} \textit{The set of \(S'\to S\) for which \(P(S')\) holds is a
 refinement of \(S\).}
 \item[(ii)]\SGARefTarget{sga3:iv:definition:4:2:8:item:ii:01} \textit{There is a refinement of \(S\) such that \(P(S')\)
 holds for every \(S'\) in that refinement.}
 \item[(iii)]\SGARefTarget{sga3:iv:definition:4:2:8:item:iii:01} \textit{If the topology is defined by a pretopology, there
 is a covering family for that pretopology such that \(P(S')\) holds for
 every \(S'\) in the family.}
\end{enumerate}

% Source page: Exp4 component-local p. 27; printed p. 203.
\medskip
\noindent\SGARefTarget{sga3:iv:example:4:2:9}\textbf{Example 4.2.9.}
\textit{Let \(f:X\to Y\) be an \(S\)-morphism. We say that \(f\) is
\emph{locally an isomorphism} if there is a covering family
\(\{S_i\to S\}\) such that \(f\times_SS_i\) is an isomorphism for every
\(i\). Equivalently, there is a refinement \(R\) of \(S\) such that
\(X(T)\to Y(T)\) is an isomorphism for every \(T\to R\).}

Many other examples of this ``local'' language will appear below.

\SGARefTarget{sga3:iv:subsection:4:3}\subsection*{4.3. Presheaves, Sheaves, and the Sheaf Associated to a Presheaf}
\addcontentsline{toc}{subsection}{\SGARefLink{sga3:iv:subsection:4:3}{4.3}. Presheaves, Sheaves, and the Sheaf Associated to a Presheaf}

\medskip
\noindent\SGARefTarget{sga3:iv:definition:4:3:1}\textbf{Definition 4.3.1.}
\textit{Let \(\mathscr C\) be a category. A \emph{presheaf of sets} on
\(\mathscr C\) is a contravariant functor from \(\mathscr C\) to the
category of sets. The category}
\[
 \widehat{\mathscr C}
 =\operatorname{Hom}(\mathscr C^{\circ},\mathbf{Set})
\]
\textit{is called the \emph{category of presheaves} on \(\mathscr C\). If
\(\mathscr C\) is equipped with a topology, a presheaf \(P\) is
\emph{separated} (respectively, is a \emph{sheaf}) if, for every
\(S\in\operatorname{Ob}\mathscr C\) and every \(R\in J(S)\), the canonical
map}
\[
\tag{+}\label{sga3:iv:definition:4:3:1:eq:unnamed:01}
 P(S)=\operatorname{Hom}(S,P)\longrightarrow\operatorname{Hom}(R,P)
\]
\textit{is injective (respectively, bijective). The full subcategory of
\(\widehat{\mathscr C}\) whose objects are the sheaves is called the
\emph{category of sheaves} and is denoted by
\(\widetilde{\mathscr C}\).}\footnote{Editor's note: If \(Q\) is a
subpresheaf of a separated presheaf \(P\), then \(Q\) is separated. Indeed,
for every sieve \(R\) of \(S\), the composite
\(Q(S)\hookrightarrow P(S)\hookrightarrow P(R)\) is injective and factors
through \(Q(S)\to Q(R)\).}

\medskip
\noindent\SGARefTarget{sga3:iv:proposition:4:3:2}\textbf{Proposition 4.3.2.}
\textit{Let \(P\) be a separated presheaf (respectively, a sheaf). For
every \(H\in\operatorname{Ob}\widehat{\mathscr C}\) and every
\(R\in J(H)\), the canonical map}
\[
\tag{+}\label{sga3:iv:proposition:4:3:2:eq:unnamed:01}
 \operatorname{Hom}(H,P)\longrightarrow\operatorname{Hom}(R,P)
\]
\textit{is injective (respectively, bijective).}

Let \(P\) be separated, let \(H\) be a presheaf, let \(R\in J(H)\), and
let \(u,v:H\to P\) satisfy \(uj=vj\). For every \(f:S\to H\), with
\(S\in\operatorname{Ob}\mathscr C\), the sieve \(R\times_HS\) is a
refinement of \(S\), and \(ufj_S=vfj_S\). This is expressed by the diagram
\[
\begin{array}{ccccc}
 R\times_HS&\xrightarrow{j_S}&S&&\\
 \overset{f_R}{\big\downarrow}&&
 \overset{f}{\big\downarrow}&&\\
 R&\xrightarrow{j}&H&
 \mathrel{\substack{\xrightarrow{u}\\[-0.5ex]\xrightarrow[v]{}}}&P.
\end{array}
\]
Since \(P\) is separated, it follows that \(uf=vf\). This holds for every
representable \(S\), so \(u=v\).

% Source page: Exp4 component-local p. 28; printed p. 204.
Now suppose that \(P\) is a sheaf. Let \(g:R\to P\); we show that it
factors through \(H\). For every \(f:S\to H\), with
\(S\in\operatorname{Ob}\mathscr C\), the morphism
\(g\circ f_R:R\times_HS\to P\) factors uniquely through \(S\), thereby
defining a morphism \(h:S\to P\). By uniqueness, \(h\) is plainly functorial
in \(f\):
\[
\begin{tikzcd}[column sep=4.5em,row sep=3.1em]
 R\times_HS
   \arrow[r,"f_R"]
   \arrow[d,"j_S"']
 &
 R
   \arrow[r,"g"]
   \arrow[d,"j"']
 &
 P
 \\
 S
   \arrow[r,"f"']
   \arrow[rru,bend right=38,"h"']
 &
 H
   \arrow[ru,dashed]
 &
 {}
\end{tikzcd}
\]
Thus, for every \(S\), we have defined a map \(H(S)\to P(S)\) functorial
in \(S\), and hence a morphism \(H\to P\) with the required properties.

\medskip
\noindent\SGARefTarget{sga3:iv:corollary:4:3:2:1}\textbf{Corollary 4.3.2.1.}\footnote{Editor's note: \SGARefLink{sga3:iv:corollary:4:3:2:1}{This
corollary} has been added.}
\textit{Let \(R\) and \(F\) be two sheaves. If \(R\) is a refinement of
\(F\), then \(R=F\).}

Indeed, suppose that \(R\) is a refinement of \(F\), and write
\(j:R\hookrightarrow F\) for the inclusion. By \SGARefLink{sga3:iv:proposition:4:3:2}{4.3.2},
\(\operatorname{Hom}(F,R)=\operatorname{End}(R)\); hence there is a
\(\pi:F\to R\) such that \(\pi j=\operatorname{id}_R\). Likewise,
\(\operatorname{End}(F)=\operatorname{Hom}(R,F)\), and
\(j\pi j=j\) implies \(j\pi=\operatorname{id}_F\). Thus \(j\) is an
isomorphism.

\medskip
\noindent\SGARefTarget{sga3:iv:proposition:4:3:3}\textbf{Proposition 4.3.3 (\SGARefLink{sga3:iv:bib:as}{[AS]}, \SGARefLink{sga3:iv:definition:1:3}{1.3}).}
\textit{Let \(\mathscr C\) be a category and \(P\) a presheaf on
\(\mathscr C\). For every \(S\in\operatorname{Ob}\mathscr C\), let
\(J(S)\) be the set of sieves \(R\) of \(S\) such that, for every
\(T\to S\), the map}
\[
\tag{+}\label{sga3:iv:proposition:4:3:3:eq:unnamed:01}
 \operatorname{Hom}(T,P)\longrightarrow
 \operatorname{Hom}(R\times_ST,P)
\]
\textit{is injective (respectively, bijective). Then the \(J(S)\) define a
topology on \(\mathscr C\); that is, they satisfy \SGARefLink{sga3:iv:definition:4:2:1:item:t:1:01}{(T~1)}--\SGARefLink{sga3:iv:definition:4:2:1:item:t:4:01}{(T~4)}.}

\medskip
\noindent\SGARefTarget{sga3:iv:corollary:4:3:4}\textbf{Corollary 4.3.4.}
\textit{For every \(S\in\operatorname{Ob}\mathscr C\), let \(K(S)\) be a
family of sieves satisfying \SGARefLink{sga3:iv:definition:4:2:1:item:t:1:01}{(T~1)}, and let \(P\) be a presheaf on
\(\mathscr C\). Then \(P\) is separated (respectively, a sheaf) for the
topology generated by the \(K(S)\) if and only if, for every
\(S\in\operatorname{Ob}\mathscr C\) and every \(R\in K(S)\), the canonical
map}
\[
\tag{+}\label{sga3:iv:corollary:4:3:4:eq:unnamed:01}
 \operatorname{Hom}(S,P)\longrightarrow\operatorname{Hom}(R,P)
\]
\textit{is injective (respectively, bijective).}

\medskip
\noindent\SGARefTarget{sga3:iv:corollary:4:3:5}\textbf{Corollary 4.3.5.}
\textit{For each \(S\in\operatorname{Ob}\mathscr C\), let
\(R(S)\) be a set of families of morphisms in \(\mathscr C\) with
target \(S\), satisfying \SGARefLink{sga3:iv:definition:4:2:5:item:p:1:01}{(P~1)}, for example a pretopology. Let \(P\) be a
presheaf on \(\mathscr C\). Then \(P\) is separated (respectively, a sheaf)
for the topology generated by \(R\) if and only if, for every
\(S\in\operatorname{Ob}\mathscr C\) and every family
\(\{S_i\to S\}\in R(S)\), the map}
\[
 P(S)\longrightarrow\prod_iP(S_i)
\]
\textit{is injective (respectively, the diagram}
\[
 P(S)\longrightarrow\prod_iP(S_i)
 \rightrightarrows\prod_{i,j}P(S_i\times_SS_j)
\]
\textit{is exact).}

% Source page: Exp4 component-local p. 29; printed p. 205.
\medskip
\noindent\SGARefTarget{sga3:iv:definition:4:3:6}\textbf{Definition 4.3.6.}
\textit{The \emph{canonical topology} on a category \(\mathscr C\) is the
finest topology for which every representable functor is a sheaf.}

\medskip
\noindent\SGARefTarget{sga3:iv:corollary:4:3:7}\textbf{Corollary 4.3.7.}
\textit{A sieve \(R\) of \(S\) is a refinement for the canonical topology
if and only if, for every morphism \(T\to S\) in \(\mathscr C\) and every
\(X\in\operatorname{Ob}\mathscr C\), the canonical map}
\[
 \operatorname{Hom}(T,X)\longrightarrow
 \operatorname{Hom}(R\times_ST,X)
\]
\textit{is bijective.}

\medskip
\noindent\SGARefTarget{sga3:iv:definition:4:3:8}\textbf{Definition 4.3.8.}
\textit{A sieve covering for the canonical topology is called a
\emph{universal effective epimorphic sieve}.}

\medskip
\noindent\SGARefTarget{sga3:iv:corollary:4:3:9}\textbf{Corollary 4.3.9.}
\textit{A universal effective epimorphic family defines a universal
effective epimorphic sieve. Conversely, every squarable family defining a
universal effective epimorphic sieve is universal effective epimorphic.}

Return now to a category \(\mathscr C\) equipped with an arbitrary topology
and to the construction of the sheaf associated to a presheaf \(P\). Let
\(S\) be an object of \(\mathscr C\). If \(R\supset R'\) are two
refinements of \(S\), there is a diagram
\[
\begin{array}{ccc}
 \operatorname{Hom}(S,P)&\longrightarrow&\operatorname{Hom}(R,P)\\
 &\searrow&\big\downarrow\\
 &&\operatorname{Hom}(R',P).
\end{array}
\]
As already observed, the ordered set \(J(S)\) is filtered. Since
\(S\in J(S)\), there is an evident morphism
\[
 \operatorname{Hom}(S,P)\longrightarrow
 \varinjlim_{R\in J(S)}\operatorname{Hom}(R,P).
\]

% Source page: Exp4 component-local p. 30; printed p. 206.
\medskip
\noindent\SGARefTarget{sga3:iv:definition:4:3:10:0}\textbf{Definition 4.3.10.0.}\footnote{Editor's note: The
numbering \SGARefLink{sga3:iv:definition:4:3:10:0}{4.3.10.0} has been added for later reference. Moreover, it follows
from \SGARefLink{sga3:iv:definition:4:3:10:0}{the definition} that if \(Q\to P\) is a monomorphism, then the same is
true of \(LP\to LQ\); thus \(L\) ``preserves monomorphisms'' (see also
\SGARefLink{sga3:iv:remark:4:3:16}{4.3.16} for the more general statement that \(L\) ``commutes with finite
inverse limits'').}
Set
\[
 \check H^0(S,P)=
 \varinjlim_{R\in J(S)}\operatorname{Hom}(R,P).
\]
One checks that \(\check H^0(S,P)\) depends functorially on \(S\), and hence
defines a functor \(LP\) by
\[
\tag{++}\label{sga3:iv:definition:4:3:10:0:eq:unnamed:01}
 \operatorname{Hom}(S,LP)=\check H^0(S,P)
 =\varinjlim_{R\in J(S)}\operatorname{Hom}(R,P).
\]
By construction there are morphisms
\[
 \ell_P:P\longrightarrow LP,
 \qquad
 z_R:\operatorname{Hom}(R,P)\longrightarrow\operatorname{Hom}(S,LP).
\]

\medskip
\noindent\SGARefTarget{sga3:iv:lemma:4:3:10}\textbf{Lemma 4.3.10.}
\begin{enumerate}
 \item[(i)]\SGARefTarget{sga3:iv:lemma:4:3:10:item:i:01} \textit{For every refinement \(R\) of \(S\) and every
 \(u:R\to P\), the diagram}
 \[
 \begin{array}{ccc}
 P&\xrightarrow{\ell_P}&LP\\
 \overset{u}{\Big\uparrow}&&
 \overset{z_R(u)}{\Big\uparrow}\\
 R&\xrightarrow{i_R}&S
 \end{array}
 \]
 \textit{is commutative.}
 \item[(ii)]\SGARefTarget{sga3:iv:lemma:4:3:10:item:ii:01} \textit{For every morphism \(v:S\to LP\), there exist a
 refinement \(R\) of \(S\) and a morphism \(u:R\to P\) such that
 \(v=z_R(u)\).}
 \item[(iii)]\SGARefTarget{sga3:iv:lemma:4:3:10:item:iii:01} \textit{Let \(Q\) be a functor and let \(u,v:Q\to P\)
 satisfy \(\ell_Pu=\ell_Pv\). Then the kernel of the pair \((u,v)\) is a
 refinement of \(Q\).}
 \item[(iv)]\SGARefTarget{sga3:iv:lemma:4:3:10:item:iv:01} \textit{Let \(u:R\to P\) and \(u':R'\to P\). The equality
 \(z_R(u)=z_{R'}(u')\) holds if and only if there is a refinement
 \(R''\subset R\cap R'\) of \(S\) on which \(u\) and \(u'\) coincide.}
\end{enumerate}

\noindent\textit{Proof.}
For \SGARefLink{sga3:iv:lemma:4:3:10:item:i:01}{(i)}, we must prove \(z_R(u)i_R=\ell_Pu\). It is enough to compare the
composites of these two morphisms with every morphism \(g:T\to R\), where
\(T\) is representable. Put \(f=i_Rg\) and
\(R'=R\times_ST\).\footnote{Editor's note: Let \(p:R'\to R\) be the
projection. Since \(i_R\) is a monomorphism, \(p=gi_{R'}\). Let \(g'\) be
the section of \(i_{R'}\) defined by \(g\). Then \(pg'=g\), hence
\(pg'i_{R'}=gi_{R'}=p\). This implies \(g'i_{R'}=\operatorname{id}_{R'}\),
so \(i_{R'}:R'\to T\) is an isomorphism with inverse \(g'\).} We have
\[
\begin{array}{ccccc}
 P&\xrightarrow{\ell_P}&LP&&\\
 \overset{u}{\Big\uparrow}&&
 \overset{z_R(u)}{\Big\uparrow}&&\\
 R&\xrightarrow{i_R}&S&&\\
 \overset{p}{\Big\uparrow}&\overset{g}{\nwarrow}&
 \overset{f}{\Big\uparrow}&&\\
 R'&\xrightarrow[\sim]{i_{R'}}&T.&&
\end{array}
\]
By \SGARefLink{sga3:iv:definition:4:3:10:0}{the definition} of \(\ell_P\),
\(\ell_Pug=z_R(u)f\); this is the special case of \SGARefLink{sga3:iv:lemma:4:3:10:item:i:01}{the assertion} in which
\(i_R\) is an isomorphism. But \(z_R(u)f=z_R(u)i_Rg\), proving \SGARefLink{sga3:iv:lemma:4:3:10:item:i:01}{(i)}.

Assertions \SGARefLink{sga3:iv:lemma:4:3:10:item:ii:01}{(ii)} and \SGARefLink{sga3:iv:lemma:4:3:10:item:iv:01}{(iv)} merely express \SGARefLink{sga3:iv:definition:4:3:10:0}{the definition} of
\(\operatorname{Hom}(S,LP)\) as a direct limit.

For \SGARefLink{sga3:iv:lemma:4:3:10:item:iii:01}{(iii)}, let \(K\) be the kernel of \((u,v)\). For every morphism
\(f:S\to Q\), with \(S\) representable, \(K\times_QS\) is a subfunctor of
the kernel of \(uf,vf:S\to P\). By \SGARefLink{sga3:iv:definition:4:2:1:item:tprime0:01}{\(T'0\)}, it is therefore enough to
prove \SGARefLink{sga3:iv:lemma:4:3:10:item:iii:01}{the assertion} when \(Q=S\) is representable. In that case, \SGARefLink{sga3:iv:lemma:4:3:10:item:ii:01}{(ii)} and
\SGARefLink{sga3:iv:lemma:4:3:10:item:iv:01}{(iv)} show that \(K\) contains a refinement of \(S\), and hence is itself a
refinement of \(S\).

One finally checks that \(P\mapsto LP\) defines a functor
\[
 L:\widehat{\mathscr C}\longrightarrow\widehat{\mathscr C},
\]
and that \(P\mapsto\ell_P\) defines a morphism of functors
\[
 \ell:\operatorname{Id}_{\widehat{\mathscr C}}\longrightarrow L.
\]

% Source page: Exp4 component-local p. 31; printed p. 207.
We now state the essential result.

\medskip
\noindent\SGARefTarget{sga3:iv:proposition:4:3:11}\textbf{Proposition 4.3.11.}
\begin{enumerate}
 \item[(i)]\SGARefTarget{sga3:iv:proposition:4:3:11:item:i:01} \textit{For every presheaf \(P\), the presheaf \(LP\) is
 separated and \(\ell_P:P\to LP\) is covering (\SGARefLink{sga3:iv:definition:4:2:3}{4.2.3}).}
 \item[(ii)]\SGARefTarget{sga3:iv:proposition:4:3:11:item:ii:01} \textit{If \(P\) is a sheaf, then \(P\to LP\) is an
 isomorphism.}
 \item[(iii)]\SGARefTarget{sga3:iv:proposition:4:3:11:item:iii:01} \textit{For every presheaf \(P\) and every separated
 presheaf (respectively, sheaf) \(F\), the map}
 \[
  \operatorname{Hom}(\ell_P,F):
  \operatorname{Hom}(LP,F)\longrightarrow\operatorname{Hom}(P,F)
 \]
 \textit{is injective (respectively, bijective).}
 \item[(iv)]\SGARefTarget{sga3:iv:proposition:4:3:11:item:iv:01} \textit{If \(P\) is separated, then
 \(\ell_P:P\to LP\) is a covering monomorphism (so \(P\) is a refinement
 of \(LP\)), and \(LP\) is a sheaf.}
\end{enumerate}

\noindent\textit{Proof.}
\SGARefLink{sga3:iv:proposition:4:3:11:item:i:01}{(i)} First, \(P\to LP\) is covering. Indeed, for every \(v:S\to LP\),
Lemma~\SGARefLink{sga3:iv:lemma:4:3:10:item:i:01}{4.3.10(i)},\SGARefLink{sga3:iv:lemma:4:3:10:item:ii:01}{(ii)} supplies a refinement \(R\) of \(S\) and a commutative
diagram (the map from \(R\) to the fiber product is induced by
\(v':R\to P\) and \(i_R:R\to S\))
\[
\begin{tikzcd}[column sep=4.2em,row sep=2.8em]
 {}
 &
 P
   \arrow[r,"\ell_P"]
 &
 LP
 \\
 {}
 &
 P\times_{LP}S
   \arrow[u]
   \arrow[r]
 &
 S
   \arrow[u,"v"']
 \\
 R
   \arrow[ru]
   \arrow[ruu,bend left=10,"v'"]
   \arrow[rru,bend right=8,"i_R"']
 &
 {}
 &
 {}
\end{tikzcd}
\]
so \(P\times_{LP}S\to S\) is covering. By \SGARefLink{sga3:iv:definition:4:2:3:item:c:0:01}{(C~0)}, \(P\to LP\) is
covering.

% Source page: Exp4 component-local p. 32; printed p. 208.
Now suppose that two morphisms \(v_1,v_2:S\to LP\) induce the same
morphism on a refinement \(R\) of \(S\). There are refinements \(R_i\),
\(i=1,2\), and morphisms \(u_i:R_i\to P\) with
\(z_{R_i}(u_i)=v_i\). Replacing \(R\) by a smaller refinement, we may
suppose \(R_1=R_2=R\). The commutative diagram of Lemma~\SGARefLink{sga3:iv:lemma:4:3:10:item:i:01}{4.3.10(i)} then
gives \(\ell_Pu_1=\ell_Pu_2\). By \SGARefLink{sga3:iv:lemma:4:3:10:item:iii:01}{4.3.10(iii)}, \(u_1\) and \(u_2\)
coincide on a refinement of \(R\), hence on a refinement of \(S\); by
\SGARefLink{sga3:iv:lemma:4:3:10:item:iv:01}{4.3.10(iv)}, \(z_R(u_1)=z_R(u_2)\). Thus \(LP\) is separated.

\SGARefLink{sga3:iv:proposition:4:3:11:item:ii:01}{(ii)} This is clear: if \(P\) is a sheaf, then
\(\operatorname{Hom}(S,P)\to\operatorname{Hom}(R,P)\) is already an
isomorphism for every refinement \(R\) of \(S\).

\SGARefLink{sga3:iv:proposition:4:3:11:item:iii:01}{(iii)} Let \(u,v:LP\to F\) satisfy \(u\ell_P=v\ell_P\). To show \(u=v\),
it is enough to prove \(uf=vf\) for every \(f:S\to LP\) with \(S\)
representable. There are a refinement \(R\) of \(S\) and a morphism
\(g:R\to P\) with \(f=z_R(g)\). The morphisms \(uf\) and \(vf\) agree on
\(R\), where both equal \(u\ell_Pg=v\ell_Pg\). If \(F\) is separated,
they therefore agree on \(S\), proving injectivity. If \(F\) is a sheaf,
the commutative diagram
\[
\begin{tikzcd}[column sep=large,row sep=large]
 P \arrow[r,"\ell_P"] \arrow[d] &
 LP \arrow[d] \arrow[dl]\\
 F \arrow[r,"\sim"] & LF
\end{tikzcd}
\]
shows that \(\operatorname{Hom}(\ell_P,F)\) is also surjective.

\SGARefLink{sga3:iv:proposition:4:3:11:item:iv:01}{(iv)} Suppose that \(P\) is separated. To show that \(P\to LP\) is a
monomorphism, it is enough to consider \(u,v:S\to P\), with \(S\)
representable, satisfying \(\ell_Pu=\ell_Pv\). By \SGARefLink{sga3:iv:lemma:4:3:10:item:iii:01}{4.3.10(iii)}, \(u\) and
\(v\) agree on a refinement of \(S\), and therefore agree because \(P\) is
separated. Thus \(P\to LP\) is a monomorphism; since it is covering by
\SGARefLink{sga3:iv:proposition:4:3:11:item:i:01}{(i)}, \(P\) is a refinement of \(LP\).

It remains to show that \(LP\) is a sheaf. We already know from \SGARefLink{sga3:iv:proposition:4:3:11:item:i:01}{(i)} that it
is separated. Let \(S\in\operatorname{Ob}\mathscr C\), let \(R\) be a
refinement of \(S\), and let \(h:R\to LP\). We must find \(u:S\to LP\)
with \(ui_R=h\). Since \(P\) is a refinement of \(LP\),
\(R'=P\times_{LP}R\) is a refinement of \(R\), and hence of \(S\). Put
\(u=z_{R'}(h')\); the situation is
\[
\begin{array}{ccccc}
 P&\xrightarrow{\ell_P}&LP&&\\
 \overset{h'}{\Big\uparrow}&&
 \overset{h}{\Big\uparrow}&\overset{u}{\nwarrow}&\\
 R'&\xrightarrow{j}&R&\xrightarrow{i_R}&S.
\end{array}
\]
Here the diagonal arrow is \(u:S\to LP\).
We have \(ui_{R'}=\ell_Ph'=hj\), hence \(ui_Rj=hj\). Since \(R'\) is a
refinement of \(R\) and \(LP\) is separated, Proposition~\SGARefLink{sga3:iv:proposition:4:3:2}{4.3.2} gives
\(ui_R=h\).

% Source page: Exp4 component-local p. 33; printed p. 209.
\medskip
\noindent\SGARefTarget{sga3:iv:corollary:4:3:12}\textbf{Corollary 4.3.12.}\footnote{Editor's note: The statement
and proof of \SGARefLink{sga3:iv:corollary:4:3:12}{this corollary} have been expanded.}
\textit{Let \(F\) be a sheaf and let \(R\) be a subobject of \(F\) in
\(\widehat{\mathscr C}\). Then \(R\) is a separated presheaf,
\(\ell_R:R\to LR\) is a covering monomorphism, and there is a commutative
diagram}
\[
\begin{array}{ccc}
 R&\xrightarrow{i}&F\\
 &\overset{\ell_R}{\searrow}
 &\overset{j}{\Big\uparrow}\\[-0.5ex]
 &&LR.
\end{array}
\]
\textit{Consequently, \(R\) is a refinement of \(F\) if and only if
\(j\) is an isomorphism.}

We already observed that \(R\) is separated and that \(j\) is a
monomorphism; see editor notes (24) and (26). By Proposition~\SGARefLink{sga3:iv:proposition:4:3:11:item:iv:01}{4.3.11(iv)},
\(\ell_R\) is a covering monomorphism. Thus, if \(j\) is an isomorphism,
\(R\) is a refinement of \(F\). Conversely, if \(i\) is covering, then so
is \(j\), and Proposition~\SGARefLink{sga3:iv:corollary:4:3:2:1}{4.3.2.1} shows that \(j\) is an isomorphism.

\medskip
\noindent\SGARefTarget{sga3:iv:remark:4:3:13}\textbf{Remark 4.3.13.}
If \(J'(S)\) is a cofinal subset of \(J(S)\), then
\SGARefTarget{sga3:iv:supp:remark:4:3:13:display:01}\[
 \operatorname{Hom}(S,LP)=
 \varinjlim_{R\in J'(S)}\operatorname{Hom}(R,P).
\]
In particular, let \(S\mapsto R(S)\) be a pretopology generating
the given topology. The functor \(L\) can be described using the covering
families belonging to \(R(S)\). Writing out \SGARefLink{sga3:iv:supp:remark:4:3:13:display:01}{the preceding formula}
recovers the construction of \SGARefLink{sga3:iv:bib:ma}{[MA]}.

Let \(i:\widetilde{\mathscr C}\to\widehat{\mathscr C}\) be the inclusion
functor. Proposition~\SGARefLink{sga3:iv:proposition:4:3:11}{4.3.11} gives \SGARefLink{sga3:iv:theorem:4:3:14}{the following theorem}.

\medskip
\noindent\SGARefTarget{sga3:iv:theorem:4:3:14}\textbf{Theorem 4.3.14.}
\textit{There is a unique functor
\(a:\widehat{\mathscr C}\to\widetilde{\mathscr C}\) for which the diagram}
\[
\begin{array}{ccc}
 \widehat{\mathscr C}&\xrightarrow{L}&\widehat{\mathscr C}\\
 \overset{a}{\Big\downarrow}&&
 \overset{L}{\Big\downarrow}\\
 \widetilde{\mathscr C}&\xrightarrow{i}&\widehat{\mathscr C}
\end{array}
\]
\textit{is commutative; that is, \(L(L(P))\) is a sheaf for every
presheaf \(P\). The functors \(i\) and \(a\) are adjoint: for every
presheaf \(P\) and every sheaf \(F\), there is an isomorphism, functorial in
\(P\) and \(F\),}
\[
 \operatorname{Hom}_{\widehat{\mathscr C}}(P,i(F))
 \simeq
 \operatorname{Hom}_{\widetilde{\mathscr C}}(a(P),F),
\]
\textit{or simply}
\[
 \operatorname{Hom}(P,F)\simeq\operatorname{Hom}(a(P),F).
\]

% Source page: Exp4 component-local p. 34; printed p. 210.
\medskip
\noindent\SGARefTarget{sga3:iv:definition:4:3:15}\textbf{Definition 4.3.15.}
\textit{The sheaf \(a(P)\) is called the sheaf \emph{associated} to the
presheaf \(P\).}

\medskip
\noindent\SGARefTarget{sga3:iv:remark:4:3:16}\textbf{Remark 4.3.16.}
Since \(L\) is constructed using inverse limits and filtered direct limits,
it commutes with finite inverse limits.\footnote{Editor's note: In
particular, if \(K\) is the kernel of a pair of morphisms of presheaves
\(u,v:Q\rightrightarrows P\), then \(LK\) is the kernel of
\(Lu,Lv:LQ\rightrightarrows LP\). This will be used in \SGARefLink{sga3:iv:lemma:4:4:5}{4.4.5}.}

Moreover, under the identification \(L(P\times P)=LP\times LP\), the
morphism \(\ell_{P\times P}\) is identified with
\(\ell_P\times\ell_P\). Thus, for example, if \(P\) is a presheaf of
groups, then \(LP\) is canonically a presheaf of groups and the canonical
morphism \(P\to LP\) is a morphism of groups. The same holds for the
functor \(a\). Consequently, if \(P\) is a presheaf of groups and \(F\) a
sheaf of groups, there is an isomorphism
\[
 \operatorname{Hom}_{\widehat{\mathscr C}\text{-gr.}}(P,i(F))
 \simeq
 \operatorname{Hom}_{\widetilde{\mathscr C}\text{-gr.}}(a(P),F).
\]
See \SGARefLink{sga3:iv:bib:d}{[D]} for further details.

\medskip
\noindent\SGARefTarget{sga3:iv:para:4:3:17}\textbf{4.3.17.}
Let \(\mathscr V\) be any category. A \emph{presheaf on \(\mathscr C\)
with values in \(\mathscr V\)} is a contravariant functor from
\(\mathscr C\) to \(\mathscr V\). Before defining sheaves with values in
\(\mathscr V\), we recall the definition of the inverse limit of a
functor. Let \(\mathscr R\) and \(\mathscr V\) be categories, and let
\[
 F:\mathscr R^{\circ}\longrightarrow\mathscr V
\]
be a contravariant functor. Write \(\varprojlim F\) for the object of
\(\widehat{\mathscr V}\) defined by
\[
 \begin{aligned}
 (\varprojlim F)(X)
 &=\operatorname{Hom}_{\widehat{\mathscr V}}
     (X,\varprojlim F)\\
 &=\varprojlim_{U\in\operatorname{Ob}\mathscr R}
     \operatorname{Hom}_{\mathscr V}(F(U),X)\\
 &=\operatorname{Hom}(c_X,F),
 \end{aligned}
\]
where \(X\) varies over the objects of \(\mathscr V\), where \(c_X\)
denotes the contravariant functor from \(\mathscr R\) to \(\mathscr V\)
that sends every object to \(X\) and every arrow to
\(\operatorname{id}_X\), and where the last \(\operatorname{Hom}\) is
taken in \(\operatorname{Hom}(\mathscr R^{\circ},\mathscr V)\). If
\(\mathscr R\) has a final object \(e_{\mathscr R}\), then
\(\varprojlim F=F(e_{\mathscr R})\). If \(\mathscr V=\mathbf{Set}\), the
functor \(\varprojlim\) is identified with \(\Gamma\).

% Source page: Exp4 component-local p. 35 (partial); printed p. 211.
Let \(S\) be an object of \(\mathscr C\) and \(R\) a sieve of \(S\). Let
\(\overline R\) be the full subcategory of \(\mathscr C_{/S}\) whose
object set is \(E(R)\), and let
\(i_R:\overline R\to\mathscr C_{/S}\) be the canonical functor. A
presheaf \(P\) on \(\mathscr C\) with values in \(\mathscr V\) restricts
to a functor
\[
 P_S:(\mathscr C_{/S})^{\circ}\longrightarrow\mathscr V.
\]
The functor \(i_R\) induces a morphism in \(\widehat{\mathscr V}\):
\[
 P(S)=P_S(S)=\varprojlim P_S
 \longrightarrow\varprojlim(P_S\circ i_R).
\]
We write
\[
 \varprojlim_RP=\varprojlim(P_S\circ i_R)
\]
for the latter object of \(\widehat{\mathscr V}\). In view of \SGARefLink{sga3:iv:remark:4:1:6}{4.1.6},
Definition~\SGARefLink{sga3:iv:definition:4:3:1}{4.3.1} generalizes as follows.

\medskip
\noindent\SGARefTarget{sga3:iv:definition:4:3:18}\textbf{Definition 4.3.18.}
\textit{A presheaf \(P\) on \(\mathscr C\) with values in \(\mathscr V\)
is \emph{separated} (respectively, a \emph{sheaf}) if, for every
\(S\in\operatorname{Ob}\mathscr C\) and every \(R\in J(S)\), the canonical
morphism in \(\widehat{\mathscr V}\)}
\[
 P(S)\longrightarrow\varprojlim_RP
\]
\textit{is a monomorphism (respectively, an isomorphism).}

When \(\mathscr V\) is the category \(\mathbf{Gr}\) of groups (or any
other category of sets endowed with algebraic structures defined by finite
inverse limits), the following notions are equivalent (cf. \SGARefLink{sga3:iv:bib:d}{[D]}): a
presheaf on \(\mathscr C\) with values in \(\mathbf{Gr}\) whose underlying
presheaf of sets is a sheaf, and a group object in the category of sheaves
of sets. With these identifications understood, we shall always regard
sheaves with values in a category of sets endowed with algebraic structures
defined by finite inverse limits as sheaves of sets equipped, in
\(\widetilde{\mathscr C}\), with the corresponding algebraic structure.

% Scope ends here, immediately before section 4.4 on source page 35.

% END INCLUDED SGA3_English_current_buildable_source_20260823/source/inputs/IV/tex/components/04_expose_IV_s41_s43.tex

% BEGIN INCLUDED SGA3_English_current_buildable_source_20260823/source/inputs/IV/tex/components/05_expose_IV_s44_opening_through_444.tex
% SGA 3, Expose IV, English reconstruction.
% Scope: Section 4.4 heading through the proof of Lemma 4.4.4, stopping
% immediately before editor's note 32 and Lemma 4.4.5.
% Authority: Polo--Gille Exp4-13oct24.pdf, component-local pages 35--37
% (pp. 35 and 37 partial; combined-reader pages 223--225; printed pp. 205--213
% under the re-edition's dual foliation).
% Jacob Reinhold's CC BY 4.0 Markdown lineage was used only as a phrasing
% comparison. The Polo--Gille PDF controls prose, formulas, numbering, and notes.

\setcounter{footnote}{29}

\SGARefTarget{sga3:iv:subsection:4:4}\subsection*{4.4. Exactness Properties of the Category of Sheaves}
\addcontentsline{toc}{subsection}{\SGARefLink{sga3:iv:subsection:4:4}{4.4}. Exactness Properties of the Category of Sheaves}

\medskip
\noindent\SGARefTarget{sga3:iv:theorem:4:4:1}\textbf{Theorem 4.4.1.}
\begin{enumerate}
  \item[(i)]\SGARefTarget{sga3:iv:theorem:4:4:1:item:i:01} \textit{Arbitrary inverse limits exist in
    \(\widetilde{\mathscr C}\), and \textquotedblleft are computed in
    \(\widehat{\mathscr C}\)\textquotedblright; that is, the inclusion functor
    \[
      i:\widetilde{\mathscr C}\longrightarrow\widehat{\mathscr C}
    \]
    commutes with inverse limits. More explicitly, if \((X_\alpha)\) is an
    inverse system of sheaves, the presheaf}
    \[
      \varprojlim_\alpha i(X_\alpha):
      S\longmapsto\varprojlim_\alpha X_\alpha(S)
    \]
    \textit{is a sheaf, and}
    \[
      i\!\left(\varprojlim_\alpha X_\alpha\right)
      =\varprojlim_\alpha i(X_\alpha).
    \]

  \item[(ii)]\SGARefTarget{sga3:iv:theorem:4:4:1:item:ii:01} \textit{Arbitrary direct limits exist in
    \(\widetilde{\mathscr C}\). If \((X_\alpha)\) is a direct system of
    sheaves, then}
    \[
      \varinjlim_\alpha X_\alpha
      =a\!\left(\varinjlim_\alpha i(X_\alpha)\right),
    \]
    \textit{where \(\varinjlim_\alpha i(X_\alpha)\) is the presheaf direct
    limit of the \(i(X_\alpha)\):}
    \[
      \varinjlim_\alpha i(X_\alpha):
      S\longmapsto\varinjlim_\alpha X_\alpha(S).
    \]

  \item[(iii)]\SGARefTarget{sga3:iv:theorem:4:4:1:item:iii:01} \textit{The functor}
    \[
      a:\widehat{\mathscr C}\longrightarrow\widetilde{\mathscr C}
    \]
    \textit{commutes with arbitrary direct limits and finite inverse limits.}
\end{enumerate}

Assertions \SGARefLink{sga3:iv:theorem:4:4:1:item:i:01}{(i)} and \SGARefLink{sga3:iv:theorem:4:4:1:item:ii:01}{(ii)} follow formally from the adjunction formula
(\SGARefLink{sga3:iv:theorem:4:3:14}{4.3.14}), and\footnote{Editor's note: the original has been modified here.}
the first assertion of \SGARefLink{sga3:iv:theorem:4:4:1:item:iii:01}{(iii)} follows from \SGARefLink{sga3:iv:theorem:4:4:1:item:ii:01}{(ii)}. The second assertion of \SGARefLink{sga3:iv:theorem:4:4:1:item:iii:01}{(iii)}
was already noted in \SGARefLink{sga3:iv:remark:4:3:16}{4.3.16}.

\medskip
\noindent\SGARefTarget{sga3:iv:scholium:4:4:2}\textbf{Scholium 4.4.2.}
\textit{\SGARefLink{sga3:iv:theorem:4:4:1}{This theorem} makes it possible to use the following method for proving
in \(\widetilde{\mathscr C}\) a statement involving both arbitrary direct
limits and finite inverse limits (for example,
\textquotedblleft every epimorphism is universally
effective\textquotedblright; cf. below).
One begins by proving the corresponding statement in the category of sets,
then extends it pointwise to the category of presheaves. One then uses \SGARefLink{sga3:iv:theorem:4:4:1}{the
preceding theorem} to pass from the category of presheaves to the category of
sheaves. Many examples of this method will occur below (\SGARefLink{sga3:iv:proposition:4:4:3}{4.4.3}, \SGARefLink{sga3:iv:lemma:4:4:6}{4.4.6}, \SGARefLink{sga3:iv:proposition:4:4:9}{4.4.9},
and so forth).}

Finally, observe that statements concerning the category of presheaves are
formally corollaries of the corresponding statements concerning the category
of sheaves. It is enough to take the coarsest
(\textquotedblleft chaotic\textquotedblright)
topology, namely the topology defined by
\[
  J(S)=\{S\}
  \qquad\text{for every }S\in\operatorname{Ob}\mathscr C;
\]
every functor is a sheaf for this topology.

\medskip
\noindent\SGARefTarget{sga3:iv:proposition:4:4:3}\textbf{Proposition 4.4.3.}
\textit{Let \(\mathscr F=\{F_i\to F\}\) be a family of morphisms of sheaves.
The following conditions are equivalent:}
\begin{enumerate}
  \item[(i)]\SGARefTarget{sga3:iv:proposition:4:4:3:item:i:01} \textit{\(\mathscr F\) is an epimorphic family;}
  \item[(ii)]\SGARefTarget{sga3:iv:proposition:4:4:3:item:ii:01} \textit{\(\mathscr F\) is a universal effective epimorphic
    family (\SGARefLink{sga3:iv:remark:1:13}{1.13});}
  \item[(iii)]\SGARefTarget{sga3:iv:proposition:4:4:3:item:iii:01} \textit{\(\mathscr F\) is covering (\SGARefLink{sga3:iv:definition:4:2:3}{4.2.3});}
  \item[(iv)]\SGARefTarget{sga3:iv:proposition:4:4:3:item:iv:01} \textit{the sheaf image of \(\mathscr F\), that is, the sheaf
    associated with the presheaf image of \(\mathscr F\) (\SGARefLink{sga3:iv:definition:4:1:2}{4.1.2}), is \(F\).}
\end{enumerate}

The equivalence of \SGARefLink{sga3:iv:proposition:4:4:3:item:iii:01}{(iii)} and \SGARefLink{sga3:iv:proposition:4:4:3:item:iv:01}{(iv)} follows from \SGARefLink{sga3:iv:corollary:4:3:12}{4.3.12}. The remaining
equivalences will follow from the next lemmas.

\medskip
\noindent\SGARefTarget{sga3:iv:lemma:4:4:4}\textbf{Lemma 4.4.4.}
\textit{Let \(f:X\to Y\) be a monomorphism of sheaves that is also an
epimorphism. Then \(f\) is an isomorphism.}

\SGARefLink{sga3:iv:lemma:4:4:4}{The lemma} is clear first of all in the category of sets. Let us next prove it
in the category of presheaves. Consider the presheaf
\[
  V:S\longmapsto Y(S)\amalg_{X(S)}Y(S);
\]
this is the amalgamated sum of \(Y\) and \(Y\) over \(X\) in the category of
presheaves.\footnote{Editor's note: the remainder of the proof has been
slightly modified.}
Let \(i_1,i_2:Y\rightrightarrows V\) be the two coordinate morphisms. If
\(X\to Y\) is an epimorphism in the category of presheaves, then \(i_1=i_2\).
In that case, for every \(S\), the map \(X(S)\to Y(S)\) is surjective; since
it is also injective, it is bijective, and hence \(X\to Y\) is an
isomorphism.

Finally, work in the category \(\widetilde{\mathscr C}\) of sheaves. By
\SGARefLink{sga3:iv:theorem:4:4:1:item:ii:01}{4.4.1(ii)}, the amalgamated sum in \(\widetilde{\mathscr C}\) of the sheaves
\(Y\) and \(Y\) over \(X\) is the sheaf \(Z\) associated with the presheaf
\(V\). Consider the diagram
\[
  X\xrightarrow{f}Y
  \overset{i_1}{\underset{i_2}{\rightrightarrows}}
  V\xrightarrow{\tau}Z=a(V).
\]
We have \(i_1f=i_2f\), hence \(\tau i_1f=\tau i_2f\), and therefore
\(\tau i_1=\tau i_2\), since \(f\) is an epimorphism in
\(\widetilde{\mathscr C}\). By part \SGARefLink{sga3:iv:lemma:4:4:5:item:iii:01}{(iii)} of \SGARefLink{sga3:iv:lemma:4:4:5}{the following lemma}, the
presheaf \(V\) is separated, so \(\tau\) is a monomorphism (\SGARefLink{sga3:iv:proposition:4:3:11:item:iv:01}{4.3.11(iv)}).
Thus \(i_1=i_2\), and we saw above that this implies that \(f:X\to Y\) is an
isomorphism. This proves Lemma~\SGARefLink{sga3:iv:lemma:4:4:4}{4.4.4}, once it has been verified that \(V\) is
separated.

% END INCLUDED SGA3_English_current_buildable_source_20260823/source/inputs/IV/tex/components/05_expose_IV_s44_opening_through_444.tex

% BEGIN INCLUDED SGA3_English_current_buildable_source_20260823/source/inputs/IV/tex/components/06_expose_IV_lemma_445.tex
% SGA 3, Expose IV, English reconstruction.
% Scope: editor's note 32 and Lemma 4.4.5 with its complete proof, stopping
% immediately before Lemma 4.4.6.
% Authority: Polo--Gille Exp4-13oct24.pdf, component-local pages 37--39
% (pp. 37 and 39 partial; combined-reader pages 225--227).
% The external English Markdown was used only as a comparison. The PDF controls
% prose, formulas, numbering, and editorial notes.

\setcounter{footnote}{31}

Let \(R_{\varnothing}\) denote the presheaf that assigns the empty set to
every \(S\in\operatorname{Ob}\mathscr C\); it is an initial object of
\(\widehat{\mathscr C}\).\footnote{Editor's note: part \SGARefLink{sga3:iv:lemma:4:4:5:item:i:01}{(i)} of Lemma \SGARefLink{sga3:iv:lemma:4:4:5}{4.4.5} has
been corrected, and the proofs of all three parts have been given in detail.}

\medskip
\noindent\SGARefTarget{sga3:iv:lemma:4:4:5}\textbf{Lemma 4.4.5.}
\begin{enumerate}
  \item[(i)]\SGARefTarget{sga3:iv:lemma:4:4:5:item:i:01} \textit{Suppose that
    \(R_{\varnothing}\notin J(S)\) for every
    \(S\in\operatorname{Ob}\mathscr C\). If \((X_i)\) is a family of
    separated presheaves, then the coproduct presheaf
    \(\coprod_i X_i\) is separated.}\footnote{Editor's note: in general, the
    coproduct of two sheaves \(F,G\) is not a sheaf. Indeed, let
    \(S_1,S_2\in\operatorname{Ob}\mathscr C\). Suppose that the coproduct
    \(S=S_1\coprod S_2\) exists in \(\mathscr C\), and that the fiber product
    \(S_1\times_S S_2\) is an initial object \(\varnothing\) of
    \(\mathscr C\) (cf. I, \SGARefLink{sga3:i:subsection:1:8}{1.8}). Let \(R\) be the sieve of \(S\) with basis
    \(\{S_1,S_2\}\). Then \((F\coprod G)(R)\) is the disjoint union of
    \(F(S)\coprod G(S)\) and the sets
    \(F(S_i)\times G(S_j)\), for \(i\ne j\); thus \(F\coprod G\) is not a
    sheaf in general. On the other hand, let \(\mathscr C\) be the category
    with a single object \(S\) and \(\operatorname{id}_S\) as its only
    morphism, equipped with the topology defined by
    \(J(S)=\{R_{\varnothing},S\}\). Then the only separated presheaves are
    \(R_{\varnothing}\) and \(X=\mathbf h_S\) (which is a sheaf), and
    \(X\coprod X\) is not separated.}

  \item[(ii)]\SGARefTarget{sga3:iv:lemma:4:4:5:item:ii:01} \textit{Consider an equivalence relation in the category of
    presheaves}
    \[
      X\overset{u}{\underset{v}{\rightrightarrows}}Y,
    \]
    \textit{and let \(w:Y\to Z\) be the quotient. If \(X\) is a sheaf and
    \(Y\) is separated, then \(Z\) is separated.}

  \item[(iii)]\SGARefTarget{sga3:iv:lemma:4:4:5:item:iii:01} \textit{Consider an amalgamated sum in the category of
    presheaves, in which \(u\) and \(u'\) are monomorphisms:}
    \[
    \begin{array}{ccc}
      X & \xrightarrow{\ u\ } & Y\\
      {\scriptstyle u'}\big\downarrow && \big\downarrow\\
      Y' & \longrightarrow & V.
    \end{array}
    \]
    \textit{If \(Y\) and \(Y'\) are separated and \(X\) is a sheaf, then
    \(V\) is separated.}
\end{enumerate}

\noindent\textit{Proof.}
For \SGARefLink{sga3:iv:lemma:4:4:5:item:i:01}{(i)}, set \(X=\coprod_i X_i\). Let
\(S\in\operatorname{Ob}\mathscr C\), let
\(\tau:R\hookrightarrow S\) be a refinement of \(S\), and let
\(x_1,x_2\in X(S)\) satisfy \(x_1\circ\tau=x_2\circ\tau\).
There are indices \(i,j\) such that \(x_1\in X_i(S)\) and
\(x_2\in X_j(S)\). Since \(R\ne R_{\varnothing}\), there is a morphism
\(\phi:T\to R\) with \(T\in\operatorname{Ob}\mathscr C\). Then
\[
  x_1\circ\tau\circ\phi=x_2\circ\tau\circ\phi.
\]
Since \(X(T)\) is the disjoint union of the \(X_k(T)\), it follows that
\(i=j\). Because \(X_i\) is separated and is a subobject of \(X\), the map
\[
  X_i(S)\longrightarrow X_i(R)\longrightarrow X(R)
\]
is injective. Hence \(x_1=x_2\), which proves that \(X\) is separated.

We next prove \SGARefLink{sga3:iv:lemma:4:4:5:item:iii:01}{(iii)}. Let \(i:Y\to V\) and \(j:Y'\to V\) be the canonical
morphisms, and let \(K\) be the kernel of the pair
\[
  Y\times Y'
  \overset{p}{\underset{q}{\rightrightarrows}}V,
  \qquad
  p=i\circ\operatorname{pr}_1,\quad
  q=j\circ\operatorname{pr}_2.
\]
For \(S\in\operatorname{Ob}\mathscr C\), the monomorphisms \(u,u'\) allow us
to identify \(X(S)\) with its image in \(Y(S)\), respectively in \(Y'(S)\).
Write \(Z(S)\), respectively \(Z'(S)\), for the complement. Then \(V(S)\)
identifies with the disjoint union of \(Z(S)\), \(Z'(S)\), and \(X(S)\).
It follows at once that
\[
  i(S):Y(S)\longrightarrow V(S),
  \qquad
  j(S):Y'(S)\longrightarrow V(S)
\]
are injective and that
\[
  (u\times u')(S):X(S)\longrightarrow K(S)
\]
is bijective. Consequently \(i\) and \(j\) are monomorphisms, and
\(u\times u':X\to K\) is an isomorphism.

Set \(U=Y\coprod Y'\). For a refinement
\(\tau:R\hookrightarrow S\), there is a commutative diagram
\[
\begin{array}{ccc}
  U(S)&\xrightarrow{\ U(\tau)\ }&U(R)\\
  \big\downarrow&&\big\downarrow\\
  V(S)&\xrightarrow{\ V(\tau)\ }&V(R).
\end{array}
\]
Let \(v_1,v_2\in V(S)\) have the same image in \(V(R)\). By the definition
of \(V\), they lift to elements \(y_1,y_2\in U(S)\); let \(z_1,z_2\) denote
their images in \(U(R)\). Then \(z_1,z_2\) have the same image in \(V(R)\).

Since \(i:Y\to V\) and \(j:Y'\to V\) are monomorphisms, the maps
\(Y(R)\to V(R)\) and \(Y'(R)\to V(R)\) are injective. Thus, because \(Y\)
and \(Y'\) are separated, if \(y_1,y_2\) both belong to \(Y(S)\), or both
belong to \(Y'(S)\), then \(y_1=y_2\). Otherwise we may suppose that
\(y_1\in Y(S)\) and \(y_2\in Y'(S)\), so that \(z_1\in Y(R)\) and
\(z_2\in Y'(R)\). Since \(i\circ z_1=j\circ z_2\), the morphism
\[
  z_1\boxtimes z_2:R\longrightarrow Y\times Y'
\]
factors through \(K=X\). Moreover, since \(X\) is a sheaf, the restriction
map identifies \(X(S)\) with \(X(R)\). Hence there is an \(x\in X(S)\) such
that
\[
  u(x)\circ\tau=z_1=y_1\circ\tau,
  \qquad
  u'(x)\circ\tau=z_2=y_2\circ\tau.
\]
Since \(Y\) and \(Y'\) are separated, \(u(x)=y_1\) and \(u'(x)=y_2\), and
therefore \(v_1=v_2\). This proves that \(V\) is separated.

Finally, we prove \SGARefLink{sga3:iv:lemma:4:4:5:item:ii:01}{(ii)}. First observe that the morphism of presheaves
\[
  X\xrightarrow{\,u\boxtimes v\,}K,
\]
where \(K\) is the kernel of the pair
\[
  w\circ\operatorname{pr}_i:Y\times Y\rightrightarrows Z,
\]
is an isomorphism. Indeed, for every \(T\in\operatorname{Ob}\mathscr C\),
\(X(T)\) is an equivalence relation on \(Y(T)\), and hence the diagram
\[
  X(T)\xrightarrow{\,u\boxtimes v\,}Y(T)\times Y(T)
  \overset{w\circ\operatorname{pr}_1}
    {\underset{w\circ\operatorname{pr}_2}{\rightrightarrows}}
  Z(T)
\]
is exact in the category of sets.

Now let \(S\in\operatorname{Ob}\mathscr C\), and let
\(g_1,g_2:S\to Z\) be two morphisms that agree on a refinement
\(\tau:R\hookrightarrow S\). By the construction of \(Z\),
\[
  Z(S)=Y(S)/X(S),
\]
so there are morphisms \(f_1,f_2:S\to Y\) such that \(wf_i=g_i\). We have
\(wf_1\tau=wf_2\tau\); by the preceding observation, there is a morphism
\(\phi:R\to X\) such that
\[
  u\phi=f_1\tau,\qquad v\phi=f_2\tau.
\]
Since \(X\) is a sheaf, there is a morphism \(\psi:S\to X\) with
\(\phi=\psi\tau\). Thus, in \(Y(R)\),
\[
  u\psi\tau=f_1\tau,\qquad v\psi\tau=f_2\tau.
\]
Because \(Y(S)\to Y(R)\) is injective, \(Y\) being separated, it follows
that \(u\psi=f_1\) and \(v\psi=f_2\), hence \(g_1=g_2\). Thus \(Z\) is
separated.


% END INCLUDED SGA3_English_current_buildable_source_20260823/source/inputs/IV/tex/components/06_expose_IV_lemma_445.tex

% BEGIN INCLUDED SGA3_English_current_buildable_source_20260823/source/inputs/IV/tex/components/07_expose_IV_lemma_446.tex
% SGA 3, Expose IV, English reconstruction.
% Scope: Lemma 4.4.6 and its complete proof, stopping before Lemma 4.4.7.
% Authority: Polo--Gille Exp4-13oct24.pdf, component-local pages 39--40
% (both partial; combined-reader pages 227--228).
% The external English Markdown was used only as a comparison. The PDF controls
% prose, formulas, numbering, and editorial notes.

\setcounter{footnote}{33}

\medskip
\noindent\SGARefTarget{sga3:iv:lemma:4:4:6}\textbf{Lemma 4.4.6.}\footnote{Editor's note: the statement of \SGARefLink{sga3:iv:lemma:4:4:6}{the
lemma} and its proof have been given in detail.}
\begin{enumerate}
  \item[(i)]\SGARefTarget{sga3:iv:lemma:4:4:6:item:i:01} \textit{Let \(\{F_i\xrightarrow{f_i}F\}\) be a family of
    morphisms of presheaves, and let \(G\subset F\) be its presheaf image.
    Then \SGARefLink{sga3:iv:supp:lemma:4:4:6:diagram:01}{the following diagram} in \(\widehat{\mathscr C}\) is exact:}
    \SGARefTarget{sga3:iv:supp:lemma:4:4:6:diagram:01}\[
      \left(\coprod_i F_i\right)\times_G
      \left(\coprod_j F_j\right)
      \overset{\operatorname{pr}_1}
        {\underset{\operatorname{pr}_2}{\rightrightarrows}}
      \coprod_i F_i\longrightarrow G.
    \]
    \textit{Equivalently, for every presheaf \(H\), \SGARefLink{sga3:iv:lemma:4:4:6:eq:unnamed:01}{the following diagram} of
    sets is exact:}
    \[
    \tag{*}\label{sga3:iv:lemma:4:4:6:eq:unnamed:01}
      \operatorname{Hom}(G,H)
      \longrightarrow
      \prod_i\operatorname{Hom}(F_i,H)
      \overset{p}{\underset{q}{\rightrightarrows}}
      \prod_{i,j}\operatorname{Hom}(F_i\times_G F_j,H).
    \]

  \item[(ii)]\SGARefTarget{sga3:iv:lemma:4:4:6:item:ii:01} \textit{Every covering family of morphisms of sheaves is a
    universal effective epimorphic family.}
\end{enumerate}

\noindent\textit{Proof.}
For \SGARefLink{sga3:iv:lemma:4:4:6:item:i:01}{(i)}, let \(H\) be a presheaf. The map
\[
  f^*:\operatorname{Hom}(G,H)
  \longrightarrow\prod_i\operatorname{Hom}(F_i,H),
  \qquad
  \phi\longmapsto(\phi\circ f_i)_i,
\]
is injective. Indeed, for each \(S\in\operatorname{Ob}\mathscr C\), the map
\(\phi(S)\) is determined by the maps \((\phi\circ f_i)(S)\), since the family
\[
  f_i(S):F_i(S)\longrightarrow G(S)
\]
is surjective. It is clear that the image of \(f^*\) is contained in
\(\ker(p,q)\).

Conversely, let \(\phi_i:F_i\to H\) be a family of morphisms such that, for
every \(i,j\), the diagram
\[
\begin{array}{ccc}
  F_i\times_G F_j&\longrightarrow&F_j\\
  \big\downarrow&&\big\downarrow{\scriptstyle\phi_j}\\
  F_i&\xrightarrow{\ \phi_i\ }&H
\end{array}
\]
commutes. For every \(S\), the map
\[
  \coprod_i F_i(S)\longrightarrow H(S)
\]
then factors uniquely through a map \(\phi(S):G(S)\to H(S)\). These maps
define a morphism \(\phi:G\to H\) satisfying
\(f^*(\phi)=(\phi_i)\). This proves the exactness of \(\SGARefLink{sga3:iv:lemma:4:4:6:eq:unnamed:01}{(*)}\), and hence \SGARefLink{sga3:iv:lemma:4:4:6:item:i:01}{(i)}.

We prove \SGARefLink{sga3:iv:lemma:4:4:6:item:ii:01}{(ii)}. Since covering families are stable under base change, it is
enough to show that every covering family is effective epimorphic. Let
\(\{F_i\to F\}\) be a covering family of morphisms of sheaves, and let \(G\)
be its presheaf image. Because the family is covering, the monomorphism
\(G\hookrightarrow F\) is covering as well; therefore, by \SGARefLink{sga3:iv:corollary:4:3:12}{4.3.12},
\[
  a(G)=F.
\]
Moreover, since \(G\hookrightarrow F\) is a monomorphism, the fiber products
\(F_i\times_G F_j\) and \(F_i\times_F F_j\) coincide. Part \SGARefLink{sga3:iv:lemma:4:4:6:item:i:01}{(i)} therefore
shows that, for every presheaf \(H\), the diagram
\[
\tag{**}\label{sga3:iv:lemma:4:4:6:eq:unnamed:02}
  \operatorname{Hom}(G,H)
  \longrightarrow
  \prod_i\operatorname{Hom}(F_i,H)
  \overset{p}{\underset{q}{\rightrightarrows}}
  \prod_{i,j}\operatorname{Hom}(F_i\times_F F_j,H)
\]
is exact. If, in addition, \(H\) is a sheaf, then
\[
  \operatorname{Hom}(G,H)
  =\operatorname{Hom}(a(G),H)
  =\operatorname{Hom}(F,H).
\]
Thus \(\SGARefLink{sga3:iv:lemma:4:4:6:eq:unnamed:02}{(**)}\) shows that \(\{F_i\to F\}\) is an effective epimorphic family
in the category \(\widetilde{\mathscr C}\) of sheaves.


% END INCLUDED SGA3_English_current_buildable_source_20260823/source/inputs/IV/tex/components/07_expose_IV_lemma_446.tex

% BEGIN INCLUDED SGA3_English_current_buildable_source_20260823/source/inputs/IV/tex/components/08_expose_IV_s447_s448bis.tex
% SGA 3, Expose IV, English reconstruction.
% Scope: Lemma 4.4.7, completion of the proof of Proposition 4.4.3, and
% Remarks 4.4.8 and 4.4.8 bis, stopping before editor's note 38 and Lemma
% 4.4.8.1.
% Authority: Polo--Gille Exp4-13oct24.pdf, component-local pages 40--41
% (both partial; combined-reader pages 228--229).
% The external English Markdown was used only as a comparison.

\setcounter{footnote}{34}

\medskip
\noindent\SGARefTarget{sga3:iv:lemma:4:4:7}\textbf{Lemma 4.4.7.}
\textit{Every family of morphisms of sheaves \(\{F_i\to F\}\) factors as a
covering family \(\{F_i\to G\}\) followed by a monomorphism \(G\to F\).}

Indeed, it is enough to take \(G\) to be the sheaf image of the given family.

\medskip
\noindent\textit{Proof of Proposition \SGARefLink{sga3:iv:proposition:4:4:3}{4.4.3}.}
We saw in \SGARefLink{sga3:iv:lemma:4:4:6}{4.4.6} that \SGARefLink{sga3:iv:proposition:4:4:3:item:iii:01}{(iii)} implies \SGARefLink{sga3:iv:proposition:4:4:3:item:ii:01}{(ii)}, while \SGARefLink{sga3:iv:proposition:4:4:3:item:ii:01}{(ii)} evidently implies \SGARefLink{sga3:iv:proposition:4:4:3:item:i:01}{(i)}.
Finally, let \(\{F_i\to F\}\) be an epimorphic family. By Lemma~\SGARefLink{sga3:iv:lemma:4:4:7}{4.4.7} it
factors as a covering family followed by a monomorphism. The latter
monomorphism, being dominated\footnote{Editor's note: recall (cf. editor's
note 17) that a morphism \(g:G\to H\) is said to be \emph{dominated} by a
family of morphisms \(F_i\to H\) if that family factors through \(g\).}
by an epimorphic family, is itself an epimorphism; hence it is an isomorphism
by \SGARefLink{sga3:iv:lemma:4:4:4}{4.4.4}.

\medskip
\noindent\SGARefTarget{sga3:iv:remark:4:4:8}\textbf{Remark 4.4.8.}\footnote{Editor's note: the numbering of the
original, which had two paragraphs numbered \SGARefLink{sga3:iv:lemma:4:4:7}{4.4.7}, has been corrected to agree
with later references.}
Because the presheaf image of the family
\(\mathscr F=\{F_i\to F\}\) is separated, the construction of the associated
sheaf shows that the conditions of Proposition~\SGARefLink{sga3:iv:proposition:4:4:3}{4.4.3} are also equivalent to
the following:
\begin{enumerate}
  \item[(v)]\SGARefTarget{sga3:iv:remark:4:4:8:item:v:01} for every \(S\in\operatorname{Ob}\mathscr C\), each
    \(f\in F(S)\) lies locally in the image of \(\mathscr F\); that is,

  \item[(vi)]\SGARefTarget{sga3:iv:remark:4:4:8:item:vi:01} for every \(S\in\operatorname{Ob}\mathscr C\) and every
    \(f\in F(S)\), the set of morphisms \(S'\to S\) such that the image of
    \(f\) in \(F(S')\) belongs to the image of one of the \(F_i(S')\) is a
    refinement of \(S\);

  \item[(vii)]\SGARefTarget{sga3:iv:remark:4:4:8:item:vii:01} if the topology is defined by a pretopology
    \(S\mapsto\mathscr R(S)\), then, for every
    \(S\in\operatorname{Ob}\mathscr C\) and every \(f\in F(S)\), there is a
    family \(\{S_j\to S\}\in\mathscr R(S)\) such that, for each \(j\), the
    image \(f_j\) of \(f\) in \(F(S_j)\) belongs to the image of one of the
    \(F_i(S_j)\).
\end{enumerate}

\medskip
\noindent\SGARefTarget{sga3:iv:remark:4:4:8:bis}\textbf{Remark 4.4.8 bis.}
If the sheaf \(F\) is representable, the preceding conditions are also
equivalent to:
\begin{enumerate}
  \item[(viii)]\SGARefTarget{sga3:iv:remark:4:4:8:bis:item:viii:01} the set of morphisms \(T\to F\), with
    \(T\in\operatorname{Ob}\mathscr C\), for which there exist an \(i\) and
    a commutative diagram
    \[
    \begin{array}{ccc}
       &F_i&\\[-0.4ex]
       &\big\downarrow&\\[-0.4ex]
      T&\longrightarrow&F
    \end{array}
    \]
    is a refinement of \(F\).
\end{enumerate}

Indeed, if \SGARefLink{sga3:iv:remark:4:4:8:bis:item:viii:01}{(viii)} holds, the presheaf image of the \(F_i\) is dominated
by\footnote{Editor's note: \textquotedblleft contains\textquotedblright{} has
been replaced by \textquotedblleft is dominated by\textquotedblright; cf.
editor's note 17.}
a refinement of \(F\), which implies that the family is covering.
Conversely, apply \SGARefLink{sga3:iv:remark:4:4:8:item:vi:01}{(vi)} to \(\operatorname{id}_F\in F(F)\).

\SGARefTarget{sga3:iv:remark:4:4:8:bis:item:b:01}In pictorial language, \SGARefLink{sga3:iv:remark:4:4:8:bis:item:viii:01}{this condition} says: \textit{locally on \(F\), there
is an \(i\) for which \(F_i\to F\) has a section}. In particular, a morphism
\[
  G\longrightarrow F,
\]
where \(G\) is a sheaf and \(F\) is a representable sheaf, is covering if
and only if it has a section locally on \(F\).

% END INCLUDED SGA3_English_current_buildable_source_20260823/source/inputs/IV/tex/components/08_expose_IV_s447_s448bis.tex

% BEGIN INCLUDED SGA3_English_current_buildable_source_20260823/source/inputs/IV/tex/components/09_expose_IV_lemmas_4481_4482.tex
% SGA 3, Expose IV, English reconstruction.
% Scope: editor's note 38 and Lemmas 4.4.8.1--4.4.8.2, stopping before
% Proposition 4.4.9.
% Authority: Polo--Gille Exp4-13oct24.pdf, component-local page 41 partial
% (combined-reader page 229; printed page 217).
% The external English Markdown was used only as a comparison.

\setcounter{footnote}{37}

The following lemmas will be useful in VI\textsubscript{B}, \SGARefLink{par:sga3-VIB-8.1}{8.1} and
\SGARefLink{prop:sga3-VIB-8.2}{8.2}.\footnote{Editor's note: Lemmas \SGARefLink{sga3:iv:lemma:4:4:8:1}{4.4.8.1} and \SGARefLink{sga3:iv:lemma:4:4:8:2}{4.4.8.2} have been added.}

\medskip
\noindent\SGARefTarget{sga3:iv:lemma:4:4:8:1}\textbf{Lemma 4.4.8.1.}
\textit{Let \(Q\subset P\) be presheaves of groups on \(\mathscr C\), with
\(Q\) normal in \(P\). Suppose that \(P\) is separated and that \(Q\) is a
sheaf. Then the presheaf of groups \(P/Q\) is separated. Consequently,}
\[
  a(P/Q)=L(P/Q)
  =\varinjlim_{R\in J(S)}(P/Q)(R).
\]

It is enough to prove the first assertion, since the second follows from it
by \SGARefLink{sga3:iv:proposition:4:3:11:item:iv:01}{4.3.11(iv)} and \SGARefLink{sga3:iv:proposition:4:3:11:item:ii:01}{(ii)}. Let \(S\in\operatorname{Ob}\mathscr C\), let \(R\)
be a sieve of \(S\), and let \(y\in P(S)/Q(S)\) have identity image in
\((P/Q)(R)\). We must show that \(y=1\).

By hypothesis, the morphism \(P(S)\to P(R)\) is injective, while
\(Q(S)\to Q(R)\) is an isomorphism, and the upper row in the following
commutative diagram is exact:
\[
\begin{array}{ccccccccc}
1&\longrightarrow&Q(S)&\longrightarrow&P(S)&\longrightarrow&P(S)/Q(S)
  &\longrightarrow&1\\
&&\big\downarrow{\scriptstyle\sim}&&\big\downarrow&&\big\downarrow\\
1&\longrightarrow&Q(R)&\longrightarrow&P(R)&\longrightarrow&(P/Q)(R).&&
\end{array}
\]
\SGARefLink{sga3:iv:lemma:4:4:8:1}{The result} follows if the lower row is exact. Let \(f:R\to P\) have identity
image in \((P/Q)(R)\), and let \(\phi:T\to R\) with
\(T\in\operatorname{Ob}\mathscr C\). Then \(f\circ\phi\) is the identity in
\[
  (P/Q)(T)=P(T)/Q(T),
\]
that is, \(f\circ\phi\in Q(T)\). Hence
\(\phi\mapsto f\circ\phi\) is a functorial map \(R(T)\to Q(T)\), and so
defines a morphism of functors \(\pi:R\to Q\) such that
\(\pi(\operatorname{id}_R)=f\). Thus \(f\in Q(R)\). This proves the
exactness of the lower row and establishes \SGARefLink{sga3:iv:lemma:4:4:8:1}{the lemma}.

\medskip
\noindent\SGARefTarget{sga3:iv:lemma:4:4:8:2}\textbf{Lemma 4.4.8.2.}
\textit{Let \(H\subset G\) be presheaves of groups on \(\mathscr C\).}
\begin{enumerate}
  \item[(i)]\SGARefTarget{sga3:iv:lemma:4:4:8:2:item:i:01} \textit{If \(H\) is normal in \(G\), then \(L(H)\) is normal in
    \(L(G)\).}
  \item[(ii)]\SGARefTarget{sga3:iv:lemma:4:4:8:2:item:ii:01} \textit{If \(H\) is central in \(G\), then \(L(H)\) is central
    in \(L(G)\).}
\end{enumerate}

Let \(S\in\operatorname{Ob}\mathscr C\). We must show (cf. I~\SGARefLink{sga3:i:definition:2:3:6}{2.3.6}) that
\(L(H)(S)\) is normal, respectively central, in \(L(G)(S)\). Let
\[
  h\in L(H)(S),\qquad g\in L(G)(S).
\]
There are a sieve \(R\) of \(S\) and elements \(h'\in H(R)\),
\(g'\in G(R)\) such that
\[
  h=z_R(h'),\qquad g=z_R(g')
\]
in the notation of \SGARefLink{sga3:iv:lemma:4:3:10}{4.3.10}. Since \(z_R\) is a group morphism,
\[
  ghg^{-1}h^{-1}=z_R(g'h'g'^{-1}h'^{-1}).
\]
In case \SGARefLink{sga3:iv:lemma:4:4:8:2:item:i:01}{(i)},
\[
  g'h'g'^{-1}h'^{-1}\in H(R),
\]
and therefore \(ghg^{-1}h^{-1}\in L(H)(S)\). In case \SGARefLink{sga3:iv:lemma:4:4:8:2:item:ii:01}{(ii)},
\(g'h'g'^{-1}h'^{-1}=1\), and hence \(ghg^{-1}h^{-1}=1\).


% END INCLUDED SGA3_English_current_buildable_source_20260823/source/inputs/IV/tex/components/09_expose_IV_lemmas_4481_4482.tex

% BEGIN INCLUDED SGA3_English_current_buildable_source_20260823/source/inputs/IV/tex/components/10_expose_IV_props_449_4410.tex
% SGA 3, Expose IV, English reconstruction.
% Scope: Propositions 4.4.9--4.4.10 and their proofs, stopping before
% Proposition 4.4.11.
% Authority: Polo--Gille Exp4-13oct24.pdf, component-local page 42
% (combined-reader page 230; printed pages 218 and 211).
% The external English Markdown was used only as a comparison.

\setcounter{footnote}{38}

\medskip
\noindent\SGARefTarget{sga3:iv:proposition:4:4:9}\textbf{Proposition 4.4.9.}
\textit{Every equivalence relation in \(\widetilde{\mathscr C}\) is
universally effective (\SGARefLink{sga3:iv:definition:3:3:3}{3.3.3}). More precisely, let \(R\) be a
\(\widetilde{\mathscr C}\)-equivalence relation on the sheaf \(X\). Then the
sheaf associated with the separated presheaf}
\[
  i(X)/i(R):S\longmapsto X(S)/R(S)
\]
\textit{is a universal effective quotient of \(X\) by \(R\).}

Let \(X/R\) denote the sheaf quotient of \(X\) by \(R\), which exists by
\SGARefLink{sga3:iv:theorem:4:4:1:item:ii:01}{4.4.1(ii)}:
\[
  X/R=a\bigl(i(X)/i(R)\bigr).
\]
We must show that \(X\to X/R\) is a universal effective epimorphism and that
the morphism
\[
  f:R\longrightarrow X\times_{X/R}X
\]
is an isomorphism. The first assertion has already been proved in \SGARefLink{sga3:iv:proposition:4:4:3}{4.4.3}.
As for \(f\), it is obtained by applying the functor \(a\) to the morphism
\[
  i(R)\longrightarrow i(X)\times_{i(X/R)}i(X),
\]
or, since \(i(X)/i(R)\) is separated by \SGARefLink{sga3:iv:lemma:4:4:5:item:ii:01}{4.4.5(ii)}, so that
\[
  i(X)/i(R)\longrightarrow i(X/R)
\]
is a monomorphism, to the canonical morphism
\[
  i(R)\longrightarrow
  i(X)\times_{\,i(X)/i(R)}i(X).
\]
Thus we are reduced to proving the same statement in the category of
presheaves. But \(i(X)/i(R)\) is the presheaf
\[
  S\longmapsto X(S)/R(S),
\]
so \SGARefLink{sga3:iv:proposition:4:4:9}{the assertion} reduces to its set-theoretic analogue, where it is
immediate.

\medskip
\noindent\SGARefTarget{sga3:iv:proposition:4:4:10}\textbf{Proposition 4.4.10.}
\textit{Under the hypotheses of \SGARefLink{sga3:iv:proposition:4:4:9}{4.4.9}, let \(Y\) be a subsheaf of \(X\), and
write \(R_Y\) for the equivalence relation induced on \(Y\) by \(R\). Then the
canonical morphism (\SGARefLink{sga3:iv:definition:3:1:6}{3.1.6})}
\[
  Y/R_Y\longrightarrow X/R
\]
\textit{is a monomorphism. It identifies \(Y/R_Y\) with a subsheaf of
\(X/R\), namely the sheaf image of the composite}
\[
  Y\longrightarrow X\longrightarrow X/R.
\]

The morphism of presheaves
\[
  i(Y)/i(R_Y)
  =i(Y)/i(R)_{i(Y)}
  \longrightarrow i(X)/i(R)
\]
is a monomorphism. Since the functor \(a\) is left exact (\SGARefLink{sga3:iv:remark:4:3:16}{4.3.16}), it carries
monomorphisms to monomorphisms; hence
\[
  Y/R_Y\longrightarrow X/R
\]
is a monomorphism. The final assertion follows from the commutative diagram
\[
\begin{array}{ccc}
  Y&\longrightarrow&X\\
  \big\downarrow&&\big\downarrow\\
  Y/R_Y&\longrightarrow&X/R
\end{array}
\]
and the fact that \(Y\to Y/R_Y\) is covering.

In view of \SGARefLink{sga3:iv:proposition:4:4:10}{this proposition}, we shall always identify \(Y/R_Y\) with a
subsheaf of \(X/R\).


% END INCLUDED SGA3_English_current_buildable_source_20260823/source/inputs/IV/tex/components/10_expose_IV_props_449_4410.tex

% BEGIN INCLUDED SGA3_English_current_buildable_source_20260823/source/inputs/IV/tex/components/11_expose_IV_prop_4411_end_s44.tex
% SGA 3, Expose IV, English reconstruction.
% Scope: Proposition 4.4.11 and paragraph 4.4.12, completing Section 4.4 and
% stopping immediately before the Section 4.5 heading.
% Authority: Polo--Gille Exp4-13oct24.pdf, component-local page 43 partial
% (combined-reader page 231; printed pages 219 and 212).
% The external English Markdown was used only as a comparison.

\setcounter{footnote}{38}

\medskip
\noindent\SGARefTarget{sga3:iv:proposition:4:4:11}\textbf{Proposition 4.4.11.}
\textit{Let \(R\) be a \(\widetilde{\mathscr C}\)-equivalence relation on the
sheaf \(X\). For every subsheaf \(Y\) of \(X\) stable under \(R\), let
\(Y'=Y/R_Y\), regarded as a subsheaf of \(X'=X/R\). Then}
\[
  Y=Y'\times_{X'}X,
\]
\textit{and the maps}
\[
  Y\longmapsto Y/R_Y,
  \qquad
  Y'\longmapsto Y'\times_{X'}X
\]
\textit{define a bijective correspondence between the subsheaves \(Y\) of
\(X\) stable under \(R\) and the subsheaves \(Y'\) of \(X'\).}

If \(Y'\) is a subsheaf of \(X'\), then \(Y'\times_{X'}X\) is a subsheaf of
\(X\) stable under \(R\).\footnote{Editor's note: moreover,
\((Y'\times_{X'}X)/R=Y'\).}
If \(Y'\) is obtained by passage to the quotient from a subsheaf \(Y\) of
\(X\), then \(Y\) is a subobject of \(Y'\times_{X'}X\). It is therefore
enough to prove the following: if \(Y\) and \(Y_1\) are two subsheaves of
\(X\), both stable under \(R\), with \(Y_1\) containing \(Y\), and if their
quotients \(Y/R_Y\) and \(Y_1/R_{Y_1}\) are identical, then \(Y=Y_1\).

We may plainly reduce to the case \(Y_1=X\). Write
\[
  P=i(X)/i(R),\qquad Q=i(Y)/i(R_Y)
\]
for the corresponding presheaves. The diagram
\[
\begin{array}{ccc}
  X&\longrightarrow&P\\
  \big\uparrow&&\big\uparrow\\
  Y&\longrightarrow&Q
\end{array}
\]
is Cartesian. On the other hand, there is a commutative diagram
\[
\begin{array}{ccc}
  Q&\longrightarrow&a(Q)\\
  \big\downarrow&&\big\downarrow{\scriptstyle\sim}\\
  P&\longrightarrow&a(P).
\end{array}
\]
Since \(Q\hookrightarrow a(Q)\) is covering by \SGARefLink{sga3:iv:proposition:4:3:11}{4.3.11}, the monomorphism
\(Q\hookrightarrow P\) is covering, so \(Q\) is a refinement of \(P\).
After base change, \(Y\) is a refinement of \(X\). Because \(X\) and \(Y\)
are sheaves, \SGARefLink{sga3:iv:corollary:4:3:12}{4.3.12} then gives \(Y=X\).

\medskip
\noindent\SGARefTarget{sga3:iv:para:4:4:12}\textbf{4.4.12.}
In particular, if \(Y\) is a subsheaf of \(X\) and \(Y'=Y/R_Y\), the
preceding correspondence defines a subsheaf \(\overline Y\) of \(X\), stable
under \(R\), containing \(Y\), and minimal with these properties.
It is called the \emph{saturation} of \(Y\) for the equivalence relation
\(R\).

% END INCLUDED SGA3_English_current_buildable_source_20260823/source/inputs/IV/tex/components/11_expose_IV_prop_4411_end_s44.tex

% BEGIN INCLUDED SGA3_English_current_buildable_source_20260823/source/inputs/IV/tex/components/12_expose_IV_s45_opening_prop_451.tex
% SGA 3, Expose IV, English reconstruction.
% Scope: Section 4.5 opening through Proposition 4.5.1 and the definition of
% the induced topology on C/S, stopping before Proposition 4.5.2.
% Authority: Polo--Gille Exp4-13oct24.pdf, component-local pages 43--45
% (pp. 43 and 45 partial; combined-reader pages 231--233).
% The external English Markdown was used only as a comparison.

\setcounter{footnote}{39}

\SGARefTarget{sga3:iv:subsection:4:5}\subsection*{4.5. The Case of a Topology Coarser than the Canonical Topology}
\addcontentsline{toc}{subsection}{\SGARefLink{sga3:iv:subsection:4:5}{4.5}. A Topology Coarser than the Canonical Topology}

By \SGARefLink{sga3:iv:definition:4:3:6}{4.3.6} and \SGARefLink{sga3:iv:definition:4:3:8}{4.3.8}, the following conditions on a topology \(\mathcal T\) on
\(\mathscr C\) are equivalent:
\begin{enumerate}
  \item[(i)]\SGARefTarget{sga3:iv:subsection:4:5:item:i:01} \(\mathcal T\) is coarser than the canonical topology of
    \(\mathscr C\);
  \item[(ii)]\SGARefTarget{sga3:iv:subsection:4:5:item:ii:01} every representable presheaf is a sheaf for \(\mathcal T\);
  \item[(iii)]\SGARefTarget{sga3:iv:subsection:4:5:item:iii:01} every covering sieve for \(\mathcal T\) is universal effective
    epimorphic.
\end{enumerate}
If \(\mathcal T\) is defined by a pretopology
\(S\mapsto\mathscr R(S)\), these conditions are also equivalent to:
\begin{enumerate}
  \item[(iv)]\SGARefTarget{sga3:iv:subsection:4:5:item:iv:01} every family belonging to \(\mathscr R(S)\) is universal
    effective epimorphic.
\end{enumerate}

When these conditions hold, the canonical functor
\(\mathscr C\to\widehat{\mathscr C}\) factors through a functor
\[
  j_{\mathscr C}=j:\mathscr C\longrightarrow\widetilde{\mathscr C};
\]
we also write \(j(S)=\widetilde S\).\footnote{Editor's note: we shall also
write \(\mathbf h_S=\widehat S\); cf. the first commutative diagram in
\SGARefLink{sga3:iv:corollary:4:5:4}{4.5.4}.}

\medskip
\noindent\SGARefTarget{sga3:iv:proposition:4:5:1}\textbf{Proposition 4.5.1.}
\textit{The functor \(j:\mathscr C\to\widetilde{\mathscr C}\) is fully
faithful and commutes with arbitrary inverse limits. In particular, it is
left exact and therefore preserves algebraic structures defined by finite
inverse limits.}

This follows immediately from the commutative factorization
\[
\begin{array}{ccc}
 &\mathscr C&\\[-0.5ex]
 {\scriptstyle j}\swarrow&&\searrow\\[-0.5ex]
 \widetilde{\mathscr C}&\xrightarrow{\ i\ }&\widehat{\mathscr C}
\end{array}
\]
and from \SGARefLink{sga3:iv:theorem:4:4:1:item:i:01}{4.4.1(i)}.

Before giving further properties of \(j\), we must define the topology
induced on a category \(\mathscr C/S\). We no longer suppose here that the
given topology is necessarily coarser than the canonical topology. A sieve
of an object \(T\) in \(\mathscr C\), together with a morphism \(T\to S\),
naturally defines a sieve of \(T\) in \(\mathscr C/S\), because \SGARefLink{sga3:iv:definition:4:1:2}{the
definition} of a sieve of \(T\) depends only on the category
\[
  \mathscr C/T=(\mathscr C/S)/T.
\]
For example, if the sieve is defined by a family \(\{T_i\to T\}\), its image
in \(\mathscr C/S\) is defined by the same family, now regarded as a family
of morphisms in \(\mathscr C/S\).

Declaring these images of covering sieves to be covering defines a topology
on \(\mathscr C/S\), called the \emph{topology induced} by the given
topology. In the terminology of \SGARefLink{sga3:iv:bib:as}{[AS]}, \SGARefLink{sga3:iv:proposition:2:3}{2.3}, it is the coarsest topology on
\(\mathscr C/S\) for which the canonical functor
\[
  i_S:\mathscr C/S\longrightarrow\mathscr C
\]
is a comorphism.\footnote{Editor's note: that is, for every object
\(T\to S\) of \(\mathscr C/S\), every covering sieve \(R'\) of
\(i_S(T\to S)=T\), when regarded as a sieve of the object
\((T\to S)\in\operatorname{Ob}(\mathscr C/S)\), is covering.}
By definition, the identifications
\[
  (\mathscr C/S)/T=\mathscr C/T
\]
respect these topologies.

% END INCLUDED SGA3_English_current_buildable_source_20260823/source/inputs/IV/tex/components/12_expose_IV_s45_opening_prop_451.tex

% BEGIN INCLUDED SGA3_English_current_buildable_source_20260823/source/inputs/IV/tex/components/13_expose_IV_prop_452_cor_453.tex
% SGA 3, Expose IV, Proposition 4.5.2 through Corollary 4.5.3.
% French authority: Exp4-13oct24.pdf, local pages 44--45
% (combined-reader pages 232--233; printed pages 220--221).
% The external English Markdown was used only as a comparison.

\medskip
\noindent\SGARefTarget{sga3:iv:proposition:4:5:2}\textbf{Proposition 4.5.2.}
\textit{Let \(S\) be a representable sheaf on \(\mathscr C\), and let
\(F\to S\) be a morphism in \(\widehat{\mathscr C}\). The presheaf
\[
  S'\longmapsto \operatorname{Hom}_{S}(S',F)
\]
on \(\mathscr C/S\) is separated (respectively, is a sheaf) if and only if
\(F\) is a separated presheaf (respectively, a sheaf) on \(\mathscr C\).}

For every functor \(P\), one has (I, \SGARefLink{sga3:i:proposition:1:4:1}{1.4.1})
\[
  \operatorname{Hom}(P,F)
  =\coprod_{f\in\operatorname{Hom}(P,S)}
    \operatorname{Hom}_{f}(P,F).
\]
\footnote{Editor's note: the argument implicit in the original has been
expanded below.}
Let \(S'\in\operatorname{Ob}(\mathscr C)\), and let
\(\eta:C'\to S'\) be a covering sieve. Since \(S\) is a sheaf, the map
\(f\mapsto f\circ\eta\) is a bijection
\[
  \operatorname{Hom}(S',S)\xrightarrow{\sim}
  \operatorname{Hom}(C',S).
\]
Consequently, the map
\[
  \operatorname{Hom}_{\widehat{\mathscr C}}(S',F)
  \longrightarrow
  \operatorname{Hom}_{\widehat{\mathscr C}}(C',F)
\]
decomposes as the ``disjoint union'' of the maps
\[
  \operatorname{Hom}_{f}(S',F)
  \longrightarrow
  \operatorname{Hom}_{f\circ\eta}(C',F).
\]
\SGARefLink{sga3:iv:proposition:4:5:2}{The proposition} follows.

\medskip
\noindent\SGARefTarget{sga3:iv:corollary:4:5:3}\textbf{Corollary 4.5.3.}
\textit{If a topology on \(\mathscr C\) is coarser than the canonical
topology, then the topology it induces on \(\mathscr C/S\) is coarser than
the canonical topology of \(\mathscr C/S\).}

% END INCLUDED SGA3_English_current_buildable_source_20260823/source/inputs/IV/tex/components/13_expose_IV_prop_452_cor_453.tex

% BEGIN INCLUDED SGA3_English_current_buildable_source_20260823/source/inputs/IV/tex/components/14_expose_IV_cor_454_scholium_455.tex
% SGA 3, Expose IV, Corollary 4.5.4 through Scholium 4.5.5.
% French authority: Exp4-13oct24.pdf, local pages 45--46
% (combined-reader pages 233--234; printed pages 221--222).
% The external English Markdown was used only as a comparison.

\medskip
\noindent\SGARefTarget{sga3:iv:corollary:4:5:4}\textbf{Corollary 4.5.4.}
\textit{Suppose that the given topology on \(\mathscr C\) is coarser than
the canonical topology. For every \(S\in\operatorname{Ob}(\mathscr C)\),
there is an equivalence of categories}
\[
  \widetilde{\mathscr C}_{/\widetilde S}
  \xrightarrow{\ \sim\ }
  \widetilde{\mathscr C/S}.
\]
\textit{The following diagrams commute up to isomorphism; every unlabelled
arrow in them is an equivalence:}
\[
\begin{array}{c@{\;}c@{\;}c@{\;}c@{\;}c@{\;}c@{\;}c}
\mathscr C_{/S}
 &\xrightarrow{(j_{\mathscr C})_{/S}}&
 \widetilde{\mathscr C}_{/\widetilde S}
 &\xrightarrow{(i_{\mathscr C})_{/\widetilde S}}&
 \widehat{\mathscr C}_{/\widehat S}
 &\xrightarrow{(a_{\mathscr C})_{/\widehat S}}&
 \widetilde{\mathscr C}_{/\widetilde S}
 \\
 \big\Vert && \big\downarrow\wr && \big\downarrow\wr && \big\downarrow\wr
 \\
\mathscr C_{/S}
 &\xrightarrow{j_{\mathscr C/S}}&
 \widetilde{\mathscr C/S}
 &\xrightarrow{i_{\mathscr C/S}}&
 \widehat{\mathscr C/S}
 &\xrightarrow{a_{\mathscr C/S}}&
 \widetilde{\mathscr C/S}.
\end{array}
\]
If \(T\to S\) is an object of \(\mathscr C/S\), one also has the
commutative diagram
\[
\begin{array}{ccccc}
 \bigl(\widetilde{\mathscr C}_{/\widetilde S}\bigr)_{/\widetilde T}
 &\longrightarrow&
 \bigl(\widetilde{\mathscr C/S}\bigr)_{/\widetilde T}
 &&\\[-0.15em]
 \Big\downarrow\rlap{\(\scriptstyle\wr\)}
 &&& \searrow &
 \widetilde{(\mathscr C/S)_{/T}}\, .
 \\
 \widetilde{\mathscr C}_{/\widetilde T}
 &\longrightarrow&
 \widetilde{\mathscr C/T}
 &\mathrel{\overset{\sim}{\nearrow}}&
\end{array}
\]

The commutativity of the first two squares follows from the definition of
the equivalence
\(
  \widetilde{\mathscr C}_{/\widetilde S}
  \to\widetilde{\mathscr C/S}
\).
To prove the commutativity of the last square, it is enough to observe that
the following square commutes:
\[
\begin{array}{ccc}
 \widehat{\mathscr C}_{/\widehat S}
 &\xrightarrow{L'}&
 \widehat{\mathscr C}_{/\widehat S}
 \\
 \big\downarrow && \big\downarrow
 \\
 \widehat{\mathscr C/S}
 &\xrightarrow{L''}&
 \widehat{\mathscr C/S}.
\end{array}
\]
Here \(L'\) is the restriction of the functor \(L_{\mathscr C}\) to
\(\widehat{\mathscr C}_{/\widehat S}\), and \(L''\) is the functor
\(L_{\mathscr C/S}\). This is seen immediately by returning to \SGARefLink{sga3:iv:definition:4:3:10:0}{the
definition} of the functors \(L\), given after \SGARefLink{sga3:iv:corollary:4:3:9}{4.3.9}. The second diagram is
simply the restriction to the categories of sheaves of the corresponding
diagram for the categories of presheaves (Expos\'e I, \SGARefLink{sga3:i:section:1}{no.~1}), which is
commutative.

\medskip
\noindent\SGARefTarget{sga3:iv:scholium:4:5:5}\textbf{Scholium 4.5.5.}
The preceding assertions show that, when the given topology is coarser than
the canonical topology, one may identify \(\mathscr C\) with a full
subcategory of \(\widetilde{\mathscr C}\), itself a full subcategory of
\(\widehat{\mathscr C}\). Under these identifications one may make the
customary abuses of language concerning \(\mathscr C\to
\widehat{\mathscr C}\), justified by the preceding commutative diagrams.
In particular, the first diagram of \SGARefLink{sga3:iv:corollary:4:5:4}{4.5.4} shows that the functor \(a\) may
be used without any special precaution.

In the next subsection we shall see that the identification of
\(\mathscr C\) with a full subcategory of \(\widetilde{\mathscr C}\),
unlike the corresponding identification with
\(\widehat{\mathscr C}\), commutes with the formation of certain direct
limits; we shall then explain how this fact is to be used.

\emph{From now on, unless expressly stated otherwise, we assume that the
given topology is coarser than the canonical topology, and we
systematically make the identifications described above.}

% END INCLUDED SGA3_English_current_buildable_source_20260823/source/inputs/IV/tex/components/14_expose_IV_cor_454_scholium_455.tex

% BEGIN INCLUDED SGA3_English_current_buildable_source_20260823/source/inputs/IV/tex/components/15_expose_IV_props_456_459.tex
% SGA 3, Expose IV, Proposition 4.5.6 through Remark 4.5.9.
% French authority: Exp4-13oct24.pdf, local pages 46--47
% (combined-reader pages 234--235; printed pages 222--223).
% The external English Markdown was used only as a comparison.

\medskip
\noindent\SGARefTarget{sga3:iv:proposition:4:5:6}\textbf{Proposition 4.5.6.}
\textit{Let \(F\) and \(G\) be two sheaves over \(S\), and let
\(f:F\to G\) be an \(S\)-morphism. The following conditions on \(f\) are
equivalent:}
\begin{enumerate}
  \item[(i)]\SGARefTarget{sga3:iv:proposition:4:5:6:item:i:01} \textit{\(f\) is a monomorphism (respectively, an
    epimorphism, respectively, an isomorphism) in
    \(\widetilde{\mathscr C}\);}
  \item[(ii)]\SGARefTarget{sga3:iv:proposition:4:5:6:item:ii:01} \textit{\(f\) is a monomorphism (respectively, an
    epimorphism, respectively, an isomorphism) in
    \(\widetilde{\mathscr C}_{/S}=\widetilde{\mathscr C/S}\).}
\end{enumerate}

For monomorphisms and isomorphisms this is immediate, since it is a
question about presheaves. For epimorphisms it follows from the description
of epimorphisms as covering morphisms (\SGARefLink{sga3:iv:proposition:4:4:3}{4.4.3}) and from the fact that, by \SGARefLink{sga3:iv:subsection:4:5}{the
definition} of the induced topology, these are the same in
\(\widetilde{\mathscr C}\) and in \(\widetilde{\mathscr C/S}\).

\medskip
\noindent\SGARefTarget{sga3:iv:proposition:4:5:7}\textbf{Proposition 4.5.7.}
\textit{Let \(f:F\to G\) be a morphism of sheaves. The following conditions
are equivalent:}
\begin{enumerate}
  \item[(i)]\SGARefTarget{sga3:iv:proposition:4:5:7:item:i:01} \textit{\(f\) is a monomorphism (respectively, an
    epimorphism, respectively, an isomorphism);}
  \item[(ii)]\SGARefTarget{sga3:iv:proposition:4:5:7:item:ii:01} \textit{for every \(S\in\operatorname{Ob}(\mathscr C)\),
    the morphism \(f_S:F_S\to G_S\) is locally a monomorphism
    (respectively, an epimorphism, respectively, an isomorphism), that is:}
  \item[(iii)]\SGARefTarget{sga3:iv:proposition:4:5:7:item:iii:01} \textit{for every \(S\in\operatorname{Ob}(\mathscr C)\),
    the set of morphisms \(T\to S\) for which
    \(F_T\to G_T\) is a monomorphism (respectively, an epimorphism,
    respectively, an isomorphism) is a refinement of \(S\).}
\end{enumerate}
\textit{If the given topology is defined by a pretopology \(\mathscr R\),
these conditions are also equivalent to:}
\begin{enumerate}
  \item[(iv)]\SGARefTarget{sga3:iv:proposition:4:5:7:item:iv:01} \textit{for every \(S\in\operatorname{Ob}(\mathscr C)\),
    there is a covering family
    \(\{S_i\to S\}\in\mathscr R(S)\) such that, for every \(i\),
    \(F_{S_i}\to G_{S_i}\) is a monomorphism (respectively, an
    epimorphism, respectively, an isomorphism).}
\end{enumerate}
\textit{If \(\mathscr C\) has a final object \(e\), it is enough in
conditions \SGARefLink{sga3:iv:proposition:4:5:7:item:ii:01}{(ii)}, \SGARefLink{sga3:iv:proposition:4:5:7:item:iii:01}{(iii)}, and \SGARefLink{sga3:iv:proposition:4:5:7:item:iv:01}{(iv)} to take \(S=e\).}

Clearly \SGARefLink{sga3:iv:proposition:4:5:7:item:ii:01}{(ii)}, \SGARefLink{sga3:iv:proposition:4:5:7:item:iii:01}{(iii)}, and \SGARefLink{sga3:iv:proposition:4:5:7:item:iv:01}{(iv)} are equivalent. To prove the equivalence of
\SGARefLink{sga3:iv:proposition:4:5:7:item:i:01}{(i)} and \SGARefLink{sga3:iv:proposition:4:5:7:item:ii:01}{(ii)}, as well as \SGARefLink{sga3:iv:proposition:4:5:7}{the assertion} concerning a final object, it is
enough to show that \SGARefLink{sga3:iv:proposition:4:5:7:item:ii:01}{(ii)} implies \SGARefLink{sga3:iv:proposition:4:5:7:item:i:01}{(i)} and that the properties in question
are stable under base change. We prove the latter point first. For
monomorphisms and isomorphisms it is immediate, since it is a question
about presheaves. For epimorphisms it follows from the fact that every
epimorphism of sheaves is universal (\SGARefLink{sga3:iv:proposition:4:4:3}{4.4.3}).

It remains to show that \SGARefLink{sga3:iv:proposition:4:5:7:item:ii:01}{(ii)} implies \SGARefLink{sga3:iv:proposition:4:5:7:item:i:01}{(i)}. Suppose first that
\(f_S:F_S\to G_S\) is locally a monomorphism (respectively, an
isomorphism). There is then a covering sieve \(C\) of \(S\) such that
\(f_T\) is a monomorphism (respectively, an isomorphism) for every
\(T\to S\) belonging to \(C\). Since an inverse limit of monomorphisms
(respectively, isomorphisms) is again one,
\(\operatorname{Hom}(C,f)\) is a monomorphism (respectively, an
isomorphism); see \SGARefLink{sga3:iv:proposition:4:1:4}{4.1.4}. The commutative diagram
\[
\begin{array}{ccc}
 \operatorname{Hom}(C,F)
 &\xrightarrow{\operatorname{Hom}(C,f)}&
 \operatorname{Hom}(C,G)
 \\
 \big\uparrow\wr && \big\uparrow\wr
 \\
 F(S)&\xrightarrow{f(S)}&G(S)
\end{array}
\]
therefore shows that \(f(S)\) is injective (respectively, bijective).

Finally, suppose that \(f:F\to G\) is locally an epimorphism, and let
\(H\subset G\) be its image. For every \(S\to G\), the sheaf
\(H\times_GS\) is the image of
\[
  f\times_GS:F\times_GS\longrightarrow S.
\]
To prove that \(f\) is an epimorphism, it is enough to prove that \(H\) is
a refinement of \(G\), or equivalently that \(H\times_GS\) is a
refinement of \(S\) for every \(S\). But this holds after each base change
\(T\to S\) belonging to a refinement of \(S\), since
\(f\times_GS\) is locally covering. Hence \(H\times_GS\) is itself a
refinement of \(S\), by Axiom \SGARefLink{sga3:iv:definition:4:2:1:item:t:2:01}{(T~2)}.

\medskip
\noindent\SGARefTarget{sga3:iv:corollary:4:5:8}\textbf{Corollary 4.5.8.}
\textit{Let \(F\) and \(G\) be two sheaves over \(S\), and let
\(f:F\to G\) be an \(S\)-morphism. Then \(f\) is a monomorphism,
respectively an epimorphism, respectively an isomorphism, if and only if it
is so locally on \(S\).}

\medskip
\noindent\SGARefTarget{sga3:iv:remark:4:5:9}\textbf{Remark 4.5.9.}
The proof of \SGARefLink{sga3:iv:proposition:4:5:7}{the proposition} shows that its assertion about monomorphisms
(respectively, isomorphisms) remains valid if one assumes only that \(F\)
is a separated presheaf (respectively, a sheaf) and \(G\) is an arbitrary
presheaf (respectively, a separated presheaf).


% END INCLUDED SGA3_English_current_buildable_source_20260823/source/inputs/IV/tex/components/15_expose_IV_props_456_459.tex

% BEGIN INCLUDED SGA3_English_current_buildable_source_20260823/source/inputs/IV/tex/components/16_expose_IV_def_4510_through_4513.tex
% SGA 3, Expose IV, Definition 4.5.10 through Corollary 4.5.13.
% French authority: Exp4-13oct24.pdf, local pages 47--48
% (combined-reader pages 235--236; printed pages 223--224).
% The external English Markdown was used only as a comparison.

We return temporarily to the case of an arbitrary topology and make \SGARefLink{sga3:iv:definition:4:5:10}{the
following definition}.

\medskip
\noindent\SGARefTarget{sga3:iv:definition:4:5:10}\textbf{Definition 4.5.10.}
\textit{Let \(G\to F\) be a morphism in
\(\widehat{\mathscr C}\). We say that \(G\) is a \emph{relative sheaf
over \(F\)} if, for every \(F\)-functor \(H\) and every refinement \(R\)
of \(H\), the canonical map}
\[
\tag{+}\label{sga3:iv:definition:4:5:10:eq:unnamed:01}
  \operatorname{Hom}_{F}(H,G)
  \xrightarrow{\ \sim\ }
  \operatorname{Hom}_{F}(R,G)
\]
\textit{is bijective.}

Proposition \SGARefLink{sga3:iv:proposition:4:5:2}{4.5.2} immediately gives the following generalization.

\medskip
\noindent\SGARefTarget{sga3:iv:proposition:4:5:11}\textbf{Proposition 4.5.11.}
\textit{If \(F\) is a sheaf, then \(G\) is a relative sheaf over \(F\) if
and only if \(G\) is a sheaf.}

\medskip
\noindent\SGARefTarget{sga3:iv:lemma:4:5:12}\textbf{Lemma 4.5.12.}
\textit{In a sequence}
\[
  X\longrightarrow T\longrightarrow S,
\]
\textit{where \(X,T,S\) are three objects of
\(\widehat{\mathscr C}\), suppose that \(X\) is a relative sheaf over
\(T\). Then}
\[
  U=\prod_{T/S}X
\]
\textit{is a relative sheaf over \(S\).}

Indeed, for every \(S\)-functor \(Y\), one has
\[
  \operatorname{Hom}_{S}(Y,U)
  =\operatorname{Hom}_{T}(T\times_SY,X).
\]
If \(C\) is a sieve of \(Y\), then \(T\times_SC\) is a sieve of
\(T\times_SY\), and the conclusion is immediate.

\medskip
\noindent\SGARefTarget{sga3:iv:corollary:4:5:13}\textbf{Corollary 4.5.13.}
\textit{The presheaves
\(\operatorname{Hom}_{T/S}(X,Y)\),
\(\operatorname{Isom}_{S}(X,Y)\), and the analogous presheaves are
sheaves whenever all the arguments occurring in them are sheaves.}

Indeed, all these presheaves are constructed using fiber products and
presheaves \(\prod_{T/S}X\) (I, \SGARefLink{sga3:i:subsection:1:7}{1.7} and II, \SGARefLink{sga3:ii:section:1}{1}). It is therefore enough to
verify \SGARefLink{sga3:iv:corollary:4:5:13}{the assertion} for a presheaf \(\prod_{T/S}X\); in that case it
follows from \SGARefLink{sga3:iv:proposition:4:5:11}{4.5.11} and \SGARefLink{sga3:iv:lemma:4:5:12}{4.5.12}.


% END INCLUDED SGA3_English_current_buildable_source_20260823/source/inputs/IV/tex/components/16_expose_IV_def_4510_through_4513.tex

% BEGIN INCLUDED SGA3_English_current_buildable_source_20260823/source/inputs/IV/tex/components/18_expose_IV_s46.tex
% SGA 3, Expose IV, all of section 4.6.
% French authority: Exp4-13oct24.pdf, component-local pages 48--53
% (current-edition pages 224--229; original-print pages 220--227).
% Stops immediately before section 4.7.
% The external English Markdown was used only as a comparison.

\setcounter{footnote}{42}

\SGARefTarget{sga3:iv:subsection:4:6}\subsection*{4.6. Description of the Quotient of a Sheaf by an Equivalence Relation}
\addcontentsline{toc}{subsection}{\SGARefLink{sga3:iv:subsection:4:6}{4.6}. Quotient of a Sheaf by an Equivalence Relation}

Recall that the given topology \(\mathcal T\) is assumed to be coarser than
the canonical topology.

\medskip
\noindent\SGARefTarget{sga3:iv:proposition:4:6:1}\textbf{Proposition 4.6.1.}
\textit{Let
\(
 R \mathrel{\substack{\xrightarrow{\ p_1\ }\\[-0.6ex]
                       \xrightarrow[\ p_2\ ]{}}} X
\)
be a \(\widetilde{\mathscr C}\)-equivalence relation in the sheaf \(X\).
Define \(F\in\operatorname{Ob}(\widehat{\mathscr C})\) as follows: for
each \(S\in\operatorname{Ob}(\mathscr C)\),\footnote{Editor's note:
\(R\times S\) denotes the equivalence relation in \(X\times S\) defined
by \(R\times S_{\mathrm{diag}}\subset
X\times X\times S\times S\), and \(R_Z\) denotes the equivalence
relation that it induces on \(Z\) (cf. \SGARefLink{sga3:iv:definition:3:1:6}{3.1.6}).}}
\[
 F(S)=\left\{\begin{array}{l}
 \text{sub-\(S\)-sheaves \(Z\) of \(X_S\), stable under
 \(R\times S\), whose quotient}\cr
 \text{by \(R_Z\) is \(S\), that is, for which
 \(R_Z\rightrightarrows Z\to S\) is exact}
\end{array}\right\}.
\]
\textit{For every sheaf \(Y\), \(\operatorname{Hom}(Y,F)\) is naturally
identified with the set}
\[
 \left\{\begin{array}{l}
 \text{sub-\(Y\)-sheaves of \(X\times Y\), stable under \(R\times Y\),}\cr
 \text{and having quotient \(Y\)}
 \end{array}\right\}.
\]
\textit{In particular, the subsheaf \(R\subset X\times X\) corresponds
to an element \(p\in\operatorname{Hom}(X,F)\), and the exact diagram}
\[
 R\mathrel{\substack{\xrightarrow{\ p_1\ }\\[-0.6ex]
                     \xrightarrow[\ p_2\ ]{}}}X\xrightarrow{p}F
\]
\textit{identifies \(F\) with the sheaf quotient \(X/R\).}

Indeed, put \(Q=X/R\). For every sheaf \(Y\) and every morphism
\(f\in\operatorname{Hom}(Y,Q)\), corresponding to a section
\(s:Y\hookrightarrow Q\times Y\), consider the diagram
\[
\tag{*}\label{sga3:iv:proposition:4:6:1:eq:unnamed:01}
\begin{array}{ccccc}
 R\times Y&\rightrightarrows&X\times Y&\longrightarrow&Q\times Y\cr
 &&\big\uparrow&&\big\uparrow\!{\scriptstyle s}\cr
 &&Z&\longrightarrow&Y,
\end{array}
\]
in which the right-hand square is cartesian. By \SGARefLink{sga3:iv:proposition:4:4:11}{4.4.11}, \(Z\) is
immediately seen to be a sub-\(Y\)-sheaf of \(X\times Y\), stable under
\(R\times Y\), with quotient \(Y\). Conversely, every \(Z\) of this
kind arises from a unique section of \(Q\times Y\) over \(Y\). Taking
\(Y\) first to be representable gives an isomorphism \(Q\simeq F\);
taking arbitrary \(Y\) then gives the asserted description of
\(\operatorname{Hom}(Y,F)\). Finally, the canonical morphism \(X\to Q\)
corresponds to the sub-\(X\)-sheaf \(R\subset X\times X\), which completes
the proof.

\medskip
\noindent\SGARefTarget{sga3:iv:corollary:4:6:2}\textbf{Corollary 4.6.2.}
\textit{Let \(G\) be any subfunctor of \(F\) such that
\(\operatorname{Hom}(X,G)\subset\operatorname{Hom}(X,F)\) contains
\(R\). Then the canonical morphism \(p:X\to F\) factors through \(G\).}

Since \(p\) is covering (\SGARefLink{sga3:iv:proposition:4:4:9}{4.4.9} and \SGARefLink{sga3:iv:proposition:4:4:3}{4.4.3}), it follows that \(G\) is a
refinement of \(F\). In particular, every subsheaf \(G\) of \(F\)
satisfying \SGARefLink{sga3:iv:corollary:4:6:2}{the preceding condition} equals \(F\) (\SGARefLink{sga3:iv:corollary:4:3:12}{4.3.12}).

\medskip
\noindent\SGARefTarget{sga3:iv:para:4:6:3}\textbf{4.6.3.}
We now consider the case in which \(X\) and \(R\) are representable.
We begin with some terminology. Besides the conditions \SGARefLink{sga3:iv:para:3:4:1:item:a:01}{(a)}--\SGARefLink{sga3:iv:para:3:4:1:item:d:01}{(d)}
introduced in \SGARefLink{sga3:iv:para:3:4:1}{3.4.1}, we shall use the following conditions on a family
\((M)\) of morphisms of \(\mathscr C\); for completeness, we repeat
conditions \SGARefLink{sga3:iv:para:4:6:3:item:a:01}{(a)}--\SGARefLink{sga3:iv:para:4:6:3:item:c:01}{(c)}:
\begin{enumerate}
 \item[(a)]\SGARefTarget{sga3:iv:para:4:6:3:item:a:01} \((M)\) is stable under base change.
 \item[(b)]\SGARefTarget{sga3:iv:para:4:6:3:item:b:01} The composite of two elements of \((M)\) belongs to \((M)\).
 \item[(c)]\SGARefTarget{sga3:iv:para:4:6:3:item:c:01} Every isomorphism belongs to \((M)\).
 \item[\((d_{\mathcal T})\)]\SGARefTarget{sga3:iv:para:4:6:3:item:d:mathcal:t:01} Every element of \((M)\) is covering.
 \footnote{Editor's note: \(\mathcal T\) denotes the given topology on
 \(\mathscr C\), which is coarser than the canonical topology.}
 \item[\((e_{\mathcal T})\)]\SGARefTarget{sga3:iv:para:4:6:3:item:e:mathcal:t:01} Let \(f:X\to Y\) be a morphism of
 \(\mathscr C\). If there is a refinement \(R\) of \(Y\) such that,
 for every \(Y'\to R\), the morphism
 \(X\times_Y Y'\to Y'\) belongs to \((M)\), then \(f\) belongs to
 \((M)\).
\end{enumerate}
Conditions \SGARefLink{sga3:iv:para:4:6:3:item:a:01}{(a)} and \SGARefLink{sga3:iv:para:4:6:3:item:b:01}{(b)} imply
\begin{enumerate}
 \item[\((a')\)]\SGARefTarget{sga3:iv:para:4:6:3:item:aprime:01} the fiber product of two elements of \((M)\) belongs
 to \((M)\).
\end{enumerate}
On the other hand, by \SGARefLink{sga3:iv:corollary:4:3:9}{4.3.9}, conditions \SGARefLink{sga3:iv:para:4:6:3:item:a:01}{(a)} and
\SGARefLink{sga3:iv:para:4:6:3:item:d:mathcal:t:01}{\((d_{\mathcal T})\)} imply
\begin{enumerate}
 \item[\((d')\)]\SGARefTarget{sga3:iv:para:4:6:3:item:dprime:01} every element of \((M)\) is a universal effective
 epimorphism.
\end{enumerate}

\medskip
\noindent\SGARefTarget{sga3:iv:para:4:6:4}\textbf{4.6.4.}
When \(\mathscr C\) has fiber products, the family of covering morphisms,
denoted \((M_{\mathcal T})\), satisfies the preceding conditions. Indeed
(cf. \SGARefLink{sga3:iv:definition:4:2:3}{4.2.3}), \SGARefLink{sga3:iv:para:4:6:3:item:a:01}{(a)} follows from \SGARefLink{sga3:iv:definition:4:2:3:item:c:1:01}{(C~1)}, \SGARefLink{sga3:iv:para:4:6:3:item:b:01}{(b)} from \SGARefLink{sga3:iv:definition:4:2:3:item:c:2:01}{(C~2)}, \SGARefLink{sga3:iv:para:4:6:3:item:c:01}{(c)} from \SGARefLink{sga3:iv:definition:4:2:3:item:c:4:01}{(C~4)},
\SGARefLink{sga3:iv:para:4:6:3:item:d:mathcal:t:01}{\((d_{\mathcal T})\)} from \SGARefLink{sga3:iv:definition:4:2:3}{the definition}, and
\SGARefLink{sga3:iv:para:4:6:3:item:e:mathcal:t:01}{\((e_{\mathcal T})\)} from \SGARefLink{sga3:iv:definition:4:2:3:item:c:5:01}{(C~5)}. Thus the results established below for
a family satisfying these conditions apply in particular to
\((M_{\mathcal T})\). In particular, \(\mathcal T\) may be taken to be
the canonical topology and \((M)\) to be the family of universal
effective epimorphisms.

\medskip
\noindent\SGARefTarget{sga3:iv:lemma:4:6:5}\textbf{Lemma 4.6.5.}
\textit{Let \((M)\) be a family of morphisms satisfying conditions
\SGARefLink{sga3:iv:para:4:6:3:item:a:01}{(a)}--\SGARefLink{sga3:iv:para:4:6:3:item:e:mathcal:t:01}{\((e_{\mathcal T})\)} above. Let \(R\) be a
\(\mathscr C\)-equivalence relation of type \((M)\) in
\(X\in\operatorname{Ob}(\mathscr C)\). Let \(\widetilde X\) be the
sheaf defined by \(X\), let \(\widetilde R\) be the
\(\widetilde{\mathscr C}\)-equivalence relation in \(\widetilde X\)
defined by \(R\), and let \(\widetilde X/\widetilde R\) be the sheaf
quotient. Then \(R\) is \((M)\)-effective if and only if
\(\widetilde X/\widetilde R\) is representable. When these conditions
hold, \(\widetilde X/\widetilde R\) is represented by the quotient
\(X/R\).}

Suppose first that \(R\) is \((M)\)-effective, and put \(Y=X/R\). The
canonical morphism \(p:X\to Y\) belongs to \((M)\), hence is covering by
\SGARefLink{sga3:iv:para:4:6:3:item:d:mathcal:t:01}{\((d_{\mathcal T})\)}. The corresponding morphism
\[
 \widetilde p:\widetilde X\longrightarrow
 \widetilde X/\widetilde R
\]
is therefore a universal effective epimorphism in
\(\widetilde{\mathscr C}\) (\SGARefLink{sga3:iv:proposition:4:4:3}{4.4.3}). Consequently it identifies
\(\widetilde X/\widetilde R\) with the quotient of \(\widetilde X\) by
the equivalence relation \(R'\) defined by \(\widetilde p\). Since the
canonical functor \(\mathscr C\to\widetilde{\mathscr C}\) commutes with
fiber products, \(R'\) is precisely \(\widetilde R\), because \(R\) is
the equivalence relation defined by \(p\).

Conversely, suppose that \(\widetilde X/\widetilde R\) is represented by
an object \(Y\) of \(\mathscr C\). Let \(p:X\to Y\) be the morphism
deduced from the canonical morphism
\(\widetilde X\to\widetilde X/\widetilde R\); by \SGARefLink{sga3:iv:proposition:4:4:3}{4.4.3} it is covering.
As before, \(R\) is plainly the equivalence relation defined by \(p\).
It remains only to show that \(p\in(M)\). The cartesian diagram
\[
\begin{array}{ccccc}
 R&\xrightarrow{\ \sim\ }&X\times_YX&\longrightarrow&X\cr
 {\scriptstyle p_2}\searrow&&\big\downarrow\!{\scriptstyle\mathrm{pr}_2}
 &&\big\downarrow\!{\scriptstyle p}\cr
 &&X&\xrightarrow{p}&Y
\end{array}
\]
shows that, after the covering base change \(p\), the morphism \(p\)
becomes \(p_2\), which belongs to \((M)\). \SGARefLink{sga3:iv:lemma:4:6:5}{The assertion} follows from
\SGARefLink{sga3:iv:para:4:6:3:item:e:mathcal:t:01}{\((e_{\mathcal T})\)}.

\medskip
\noindent\SGARefTarget{sga3:iv:corollary:4:6:5:1}\textbf{Corollary 4.6.5.1.}
\footnote{Editor's note: \SGARefLink{sga3:iv:corollary:4:6:5:1}{This corollary} has been added.}
\textit{Let \((M)\) satisfy conditions \SGARefLink{sga3:iv:para:4:6:3:item:a:01}{(a)}--\SGARefLink{sga3:iv:para:4:6:3:item:e:mathcal:t:01}{\((e_{\mathcal T})\)} above,
and let \(f:G\to G'\) be a morphism of \(\mathscr C\)-groups such that
\(f\in(M)\). Suppose that \(\operatorname{Ker}(f)\) is representable
(as is the case when \(\mathscr C\) has a final object \(e\)). Then the
equivalence relation in \(G\) defined by
\(H=\operatorname{Ker}(f)\) is \((M)\)-effective, and \(G'\) represents
the sheaf quotient \(\widetilde G/\widetilde H\) for the topology
\(\mathcal T\).}

This follows from \SGARefLink{sga3:iv:lemma:4:6:5}{4.6.5} and \SGARefLink{sga3:iv:corollary:3:4:7:1}{3.4.7.1}.

We can now state the main result of this subsection.

\medskip
\noindent\SGARefTarget{sga3:iv:theorem:4:6:6}\textbf{Theorem 4.6.6.}
\textit{Let \((M)\) be a family of morphisms satisfying Axioms
\SGARefLink{sga3:iv:para:4:6:3:item:a:01}{(a)}--\SGARefLink{sga3:iv:para:4:6:3:item:e:mathcal:t:01}{\((e_{\mathcal T})\)} of \SGARefLink{sga3:iv:para:4:6:3}{4.6.3}. Let \(R\) be a
\(\mathscr C\)-equivalence relation of type \((M)\) (cf. \SGARefLink{sga3:iv:definition:3:4:3}{3.4.3}) in an
object \(X\) of \(\mathscr C\). Consider the functor
\(F\in\operatorname{Ob}(\widehat{\mathscr C})\) defined by}
\[
 F(S)=\left\{\begin{array}{l}
 \text{sub-\(S\)-sheaves \(Z\) of \(X_S\), stable under \(R\times S\),}\cr
 \text{whose quotient by \(R_Z\) is \(S\)}
 \end{array}\right\}.
\]
\textit{Let \(F_0\) be the subfunctor of \(F\) defined as follows:
\(F_0(S)\) consists of the representable \(Z\in F(S)\), that is,}
\[
 F_0(S)=\left\{\begin{array}{l}
 \text{sub-\(\mathscr C/S\)-objects \(Z\) of \(X_S\), stable under
 \(R\times S\), such that}\cr
 \text{\(R_Z\) is \((M)\)-effective with quotient \(S\), i.e. such that}\cr
 \text{\(Z\to S\) belongs to \((M)\) and
 \(R_Z\simeq Z\times_S Z\)}
 \end{array}\right\}.
\]
\textit{Then:}
\begin{enumerate}
 \item[(i)]\SGARefTarget{sga3:iv:theorem:4:6:6:item:i:01} \textit{The morphism
 \(p\in\operatorname{Hom}(X,F)=F(X)\) defined by the subobject
 \(R\subset X\times X\) identifies \(F\) with the sheaf quotient of
 \(X\) by \(R\).}
 \item[(ii)]\SGARefTarget{sga3:iv:theorem:4:6:6:item:ii:01} \textit{The following conditions are equivalent:}
 \begin{enumerate}
  \item[(a)]\SGARefTarget{sga3:iv:theorem:4:6:6:item:a:01} \textit{\(F\) is representable;}
  \item[(b)]\SGARefTarget{sga3:iv:theorem:4:6:6:item:b:01} \textit{\(F_0\) is representable;}
  \item[(c)]\SGARefTarget{sga3:iv:theorem:4:6:6:item:c:01} \textit{\(R\) is \((M)\)-effective.}
 \end{enumerate}
 \textit{Under these conditions, \(F=F_0=X/R\).}
 \item[(iii)]\SGARefTarget{sga3:iv:theorem:4:6:6:item:iii:01} \textit{Let \((N)\) be a family of morphisms stable under
 base change, such that for every covering family \(\{S_i\to S\}\) and
 every family \(\{T_i\to S_i\}\) of morphisms in \((N)\), every
 descent datum on the \(T_i\) relative to \(\{S_i\to S\}\) is
 effective. Suppose that \(X\) is squarable (cf. \SGARefLink{sga3:iv:definition:1:6:0}{1.6.0}) and that
 \(R\to X\times X\) belongs to \((N)\). Then \(F_0=F\).}
\end{enumerate}

\noindent\textit{Proof.}
Assertion \SGARefLink{sga3:iv:theorem:4:6:6:item:i:01}{(i)} was proved in \SGARefLink{sga3:iv:proposition:4:6:1}{4.6.1}.

For \SGARefLink{sga3:iv:theorem:4:6:6:item:ii:01}{(ii)}, we have already proved the equivalence of \SGARefLink{sga3:iv:theorem:4:6:6:item:a:01}{(a)} and \SGARefLink{sga3:iv:theorem:4:6:6:item:c:01}{(c)}, as well
as the equality \(F=X/R\). It remains to prove that either \SGARefLink{sga3:iv:theorem:4:6:6:item:b:01}{(b)} or \SGARefLink{sga3:iv:theorem:4:6:6:item:c:01}{(c)}
implies \(F_0=F\). First note, as stated in \SGARefLink{sga3:iv:theorem:4:6:6}{the theorem}, that \(F_0\) is
indeed a subfunctor of \(F\): if \(S\in\operatorname{Ob}(\mathscr C)\)
and \(Z\in F_0(S)\), then \(Z\to S\) is squarable, and hence
\(Z\times_SS'\in F_0(S')\) for every \(S'\to S\). Since
\(R\in F(X)\) belongs to \(F_0(X)\), Corollary \SGARefLink{sga3:iv:corollary:4:6:2}{4.6.2} shows that \SGARefLink{sga3:iv:theorem:4:6:6:item:b:01}{(b)}
implies \(F_0=F\).

Now suppose that \SGARefLink{sga3:iv:theorem:4:6:6:item:c:01}{(c)} holds, and let \(Q\) be an object of \(\mathscr C\)
representing \(X/R\). Then \(X\to Q\) belongs to \((M)\). For every
\(S\in\operatorname{Ob}(\mathscr C)\) and every \(Z\in F(S)\), diagram
\(\SGARefLink{sga3:iv:proposition:4:6:1:eq:unnamed:01}{(*)}\) of \SGARefLink{sga3:iv:proposition:4:6:1}{4.6.1} gives
\[
 Z=S\times_{Q\times S}(X\times S).
\]
Thus \(Z\) is representable and \(Z\to S\) belongs to \((M)\), so
\(Z\in F_0(S)\).

For \SGARefLink{sga3:iv:theorem:4:6:6:item:iii:01}{(iii)}, let \(f\in\operatorname{Hom}(S,F)\) correspond to
\(Z\in F(S)\). We must show that \(f\) factors through \(F_0\), or
equivalently that \(Z\) is representable. This is immediate when \(f\)
factors through \(X\), by \SGARefLink{sga3:iv:lemma:4:6:7}{the following lemma}.

\medskip
\noindent\SGARefTarget{sga3:iv:lemma:4:6:7}\textbf{Lemma 4.6.7.}
\textit{Let \(x_0\in X(S)\). The image of \(x_0\) in \(F(S)\)
corresponds to the subsheaf \(Z\subset X_S\) defined by the two
cartesian squares}
\[
\begin{array}{ccccc}
 X_S&\xrightarrow{\mathrm{id}_{X_S}\times x_0}&
 X_S\times_SX_S&\longrightarrow&X\times X\cr
 \big\uparrow&&\big\uparrow&&\big\uparrow\cr
 Z&\longrightarrow&R_S&\longrightarrow&R.
\end{array}
\]

This follows at once from the description of the morphism \(X\to F\).

Returning to the proof of \SGARefLink{sga3:iv:theorem:4:6:6}{the theorem}, suppose first that \(f\) factors
through \(X\). Then \(Z\) is representable; moreover, since
\(R\to X\times X\) belongs to \((N)\), so does \(Z\to X_S\).

In general, \(f\) need not factor through \(X\). Since \(X\to F\) is
covering (\SGARefLink{sga3:iv:proposition:4:4:3}{4.4.3}), however, \SGARefLink{sga3:iv:remark:4:4:8:item:vii:01}{4.4.8(vii)} gives a covering family
\(\{S_i\to S\}\) and, for each \(i\), a morphism \(S_i\to X\) making
the diagram
\SGARefTarget{sga3:iv:supp:lemma:4:6:7:diagram:01}\[
\begin{array}{ccc}
 X&\longrightarrow&F\cr
 \big\uparrow&&\big\uparrow\!{\scriptstyle f}\cr
 S_i&\longrightarrow&S
\end{array}
\]
commute. By the preceding argument, the morphism
\(f_i:S_i\to F\) defined by \SGARefLink{sga3:iv:supp:lemma:4:6:7:diagram:01}{this diagram} belongs to
\(\operatorname{Hom}(S_i,F_0)\) and corresponds to the subsheaf
\(Z\times_SS_i\subset X_{S_i}\). The morphism
\(Z\times_SS_i\to X_{S_i}\) belongs to \((N)\), and the family
\(\{X_{S_i}\to X_S\}\) is covering. It therefore remains only to prove
\SGARefLink{sga3:iv:proposition:4:6:8}{the following proposition}.

\medskip
\noindent\SGARefTarget{sga3:iv:proposition:4:6:8}\textbf{Proposition 4.6.8.}
\textit{Let \(\{S_i\to S\}\) be a covering family and let \(Z\) be a
sheaf over \(S\). Suppose that, for every \(i\), the
\(S_i\)-functor \(Z\times_SS_i\) is represented by an object \(T_i\).
Then the family \((T_i)\) carries a canonical descent datum relative to
\(\{S_i\to S\}\). The sheaf \(Z\) is representable if and only if this
datum is effective; in that case the descended object represents \(Z\).}

First observe that, by \SGARefLink{sga3:iv:proposition:4:4:3}{4.4.3}, \(\{S_i\to S\}\) is a universal effective
epimorphic family in \(\widetilde{\mathscr C}\), and hence a descent
family in \(\widetilde{\mathscr C}\) (\SGARefLink{sga3:iv:proposition:2:3}{2.3}). If \(Z\) is represented by
an object \(T\), then \(T\times_SS_i\), regarded as a sheaf, is
isomorphic to \(Z\times_SS_i\). Thus the descent datum on the \(T_i\) is
effective, and its uniquely determined descended object is isomorphic to
\(T\), hence represents \(Z\).

Conversely, suppose that the canonical descent datum on the \(T_i\) is
effective, and let \(T\) be the descended object. Because
\(\{S_i\to S\}\) is a descent family in
\(\widetilde{\mathscr C}\), there is an \(S\)-morphism \(T\to Z\) whose
base change to each \(S_i\) is the canonical morphism
\(T_i\to Z\times_SS_i\). This morphism is locally an isomorphism; since
both \(T\) and \(Z\) are sheaves, \SGARefLink{sga3:iv:corollary:4:5:8}{4.5.8} implies that it is an
isomorphism.

\medskip
\noindent\SGARefTarget{sga3:iv:corollary:4:6:9}\textbf{Corollary 4.6.9.}
\textit{Let \(R\) be an \((M)\)-effective equivalence relation in
\(X\). For every sheaf \(F\), the map}
\[
 \operatorname{Hom}(X/R,F)\longrightarrow\operatorname{Hom}(X,F)
\]
\textit{identifies the first set with the subset of the second consisting
of morphisms compatible with \(R\).}

\medskip
\noindent\SGARefTarget{sga3:iv:corollary:4:6:10}\textbf{Corollary 4.6.10.}
\textit{Let \(\mathcal T'\) be a topology coarser than \(\mathcal T\),
for which the morphisms in \((M)\) are covering. Under the hypotheses of
\SGARefLink{sga3:iv:theorem:4:6:6:item:iii:01}{4.6.6(iii)}, \(X/R\) is also the sheaf quotient of \(X\) by \(R\) for
every topology lying between \(\mathcal T'\) and the canonical topology.}

\medskip
\noindent\SGARefTarget{sga3:iv:remark:4:6:11}\textbf{Remark 4.6.11.}
If, in the statement of \SGARefLink{sga3:iv:theorem:4:6:6:item:iii:01}{4.6.6(iii)}, one assumes in addition that the
descended morphism \(T\to S\) belongs to \((N)\) whenever the hypotheses
there are satisfied, then the inclusion morphisms \(Z\hookrightarrow X_S\)
also belong to \((N)\). This follows immediately from the construction of
\(Z\) by descent.

\medskip
\noindent\SGARefTarget{sga3:iv:remark:4:6:12}\textbf{Remark 4.6.12.}
The implications
\[
 \text{\SGARefLink{sga3:iv:theorem:4:6:6:item:c:01}{(c)}}\Longrightarrow\text{\SGARefLink{sga3:iv:theorem:4:6:6:item:b:01}{(b)}}\Longrightarrow\text{\SGARefLink{sga3:iv:theorem:4:6:6:item:a:01}{(a)}}
 \qquad\text{and}\qquad
 \text{\SGARefLink{sga3:iv:theorem:4:6:6:item:c:01}{(c)}}\Longrightarrow[F_0=F=X/R]
\]
were proved without using the sufficiency direction of Lemma \SGARefLink{sga3:iv:lemma:4:6:5}{4.6.5}, the
only place where condition \SGARefLink{sga3:iv:para:4:6:3:item:e:mathcal:t:01}{\((e_{\mathcal T})\)} enters. They therefore
remain valid when \((M)\) satisfies only conditions
\SGARefLink{sga3:iv:para:4:6:3:item:a:01}{(a)}--\SGARefLink{sga3:iv:para:4:6:3:item:d:mathcal:t:01}{\((d_{\mathcal T})\)}. The family of squarable covering morphisms is
an example of such a family \((M)\) (compare \SGARefLink{sga3:iv:para:4:6:4}{4.6.4}). For the canonical
topology these are precisely the universal effective epimorphisms. Hence:

\medskip
\noindent\SGARefTarget{sga3:iv:corollary:4:6:13}\textbf{Corollary 4.6.13.}
\textit{Let \(R\) be a universal effective equivalence relation in
\(X\). Then the object \(X/R\) of \(\mathscr C\) is the sheaf quotient
of \(X\) by \(R\) for the canonical topology. It represents the
following functor: \((X/R)(S)\) is the set of
sub-\(\mathscr C/S\)-objects \(Z\) of \(X_S\), stable under
\(R\times S\), for which the induced equivalence relation is universal
effective and has quotient \(S\).}

Similarly, for an arbitrary topology:

\medskip
\noindent\SGARefTarget{sga3:iv:corollary:4:6:14}\textbf{Corollary 4.6.14.}
\textit{Let \((M)\) be the family of squarable covering morphisms. If
\(R\) is an \((M)\)-effective equivalence relation in \(X\), then the
object \(X/R\) of \(\mathscr C\) is the sheaf quotient of \(X\) by
\(R\), and it represents the functor \(F_0\) of \SGARefLink{sga3:iv:theorem:4:6:6}{4.6.6}.}

\medskip
\noindent\SGARefTarget{sga3:iv:scholium:4:6:15}\textbf{Scholium 4.6.15.}
We can now make the following clarifications to \SGARefLink{sga3:iv:scholium:4:5:5}{4.5.5}. In questions
involving inverse limits alone---fiber products, algebraic structures,
and so forth---the results of Expos\'e I and \SGARefLink{sga3:iv:scholium:4:5:5}{4.5.5} allow us to identify
\(\mathscr C\), without distinction, with a full subcategory of either
\(\widetilde{\mathscr C}\) or \(\widehat{\mathscr C}\). The same is not
true for questions that mix inverse and direct limits.

Whenever both inverse and direct limits occur, and especially when
passing to quotients (for example, when putting a group structure on the
quotient of a group by an invariant subgroup), we shall regard the given
category as embedded in the category of sheaves. Thus, if \(R\) is a
\(\mathscr C\)-equivalence relation in an object \(X\) of
\(\mathscr C\), the notation \(X/R\) will mean the sheaf quotient of
\(X\) by \(R\), formerly denoted \(j(X)/j(R)\), and, when this sheaf is
representable, also the object that represents it. The preceding results
show that, in the most important cases, a quotient in \(\mathscr C\) is
also a quotient in the category of sheaves. In any event, we shall not
use the notation \(X/R\) for a quotient in \(\mathscr C\) that does not
coincide with the quotient in \(\widetilde{\mathscr C}\)---for instance,
one that is not universal. This modifies the definitions of \SGARefLink{sga3:iv:section:3}{no.~3}.

To study a problem of this kind, one therefore works first in the
category of sheaves, where all the usual results hold (cf. \SGARefLink{sga3:iv:subsection:4:4}{no.~4.4}), and
then specializes the resulting statements to the original category,
using the results of the present subsection and, when available,
criteria for the effectivity of descent. Examples of this method will
appear in the following subsections.

% END INCLUDED SGA3_English_current_buildable_source_20260823/source/inputs/IV/tex/components/18_expose_IV_s46.tex

% BEGIN INCLUDED SGA3_English_current_buildable_source_20260823/source/inputs/IV/tex/components/19_expose_IV_s47.tex
% SGA 3, Expose IV, all of section 4.7.
% French authority: Exp4-13oct24.pdf, local pages 53--55
% (combined-reader pages 241--243; printed pages 229--231).
% The external English Markdown was used only as a comparison.

\setcounter{footnote}{45}

\SGARefTarget{sga3:iv:subsection:4:7}\subsection*{4.7. Use of Effectivity Criteria: Isomorphism Theorem}
\addcontentsline{toc}{subsection}{\SGARefLink{sga3:iv:subsection:4:7}{4.7}. Use of Effectivity Criteria: Isomorphism Theorem}

In this subsection we give an example of the use of effectivity criteria.
The initial data are a topology \(\mathcal T\) on \(\mathscr C\), still
assumed coarser than the canonical topology; a family \((M)\) of morphisms
of \(\mathscr C\) satisfying Axioms \SGARefLink{sga3:iv:para:4:6:3:item:a:01}{(a)}--\SGARefLink{sga3:iv:para:4:6:3:item:e:mathcal:t:01}{\((e_{\mathcal T})\)} of \SGARefLink{sga3:iv:para:4:6:3}{4.6.3};
and a family \((N)\) of morphisms of \(\mathscr C\) which may satisfy the
following axioms:
\begin{enumerate}
  \item[(a)]\SGARefTarget{sga3:iv:subsection:4:7:item:a:01} \((N)\) is stable under base change;
  \item[\((f_{\mathcal T})\)]\SGARefTarget{sga3:iv:subsection:4:7:item:f:mathcal:t:01} morphisms in \((N)\) descend for the given
    topology. More precisely, for every
    \(S\in\operatorname{Ob}(\mathscr C)\), every covering family
    \(\{S_i\to S\}\), and every family \(\{T_i\to S_i\}\) of morphisms
    in \((N)\), every descent datum on the \(T_i\) relative to
    \(\{S_i\to S\}\) is effective; if \(T\) denotes the descended
    object, then \(T\to S\) belongs to \((N)\).
\end{enumerate}

Since every element of \((M)\) is covering (condition \SGARefLink{sga3:iv:para:4:6:3}{4.6.3}
\SGARefLink{sga3:iv:para:4:6:3:item:d:mathcal:t:01}{\((d_{\mathcal T})\)}), Axiom \SGARefLink{sga3:iv:subsection:4:7:item:f:mathcal:t:01}{\((f_{\mathcal T})\)} implies \SGARefLink{sga3:iv:subsection:4:7:item:f:m:01}{the following
axiom}:\footnote{Editor's note: Axiom \SGARefLink{sga3:iv:subsection:4:7:item:f:m:01}{\((f_M)\)}, which originally appeared
before Proposition \SGARefLink{sga3:iv:proposition:4:7:2}{4.7.2}, has been moved here.}
\begin{enumerate}
  \item[\((f_M)\)]\SGARefTarget{sga3:iv:subsection:4:7:item:f:m:01} if \(Y\to X\) belongs to \((N)\) and
    \(X\to X_1\) belongs to \((M)\), every descent datum on \(Y\)
    relative to \(X\to X_1\) is effective; if \(Y_1\) denotes the
    descended object, then \(Y_1\to X_1\) belongs to \((N)\).
\end{enumerate}

Here is an example, to be treated later: \(\mathscr C\) is the category of
schemes, \(\mathcal T\) is the faithfully flat quasi-compact topology,
\((M)\) is the family of faithfully flat quasi-compact morphisms, and
\((N)\) is the family of closed immersions, or the family of quasi-compact
immersions.\footnote{Editor's note: cf. \SGARefLink{sga3:iv:subsection:6:4}{\S6.4}; see also VI\(_A\), \SGARefLink{sga3:iv:proposition:5:3:1}{5.3.1}.}

Taking \SGARefLink{sga3:iv:remark:4:6:11}{4.6.11} into account, the principal result of \SGARefLink{sga3:iv:theorem:4:6:6}{4.6.6} may be recalled
as follows.

\medskip
\noindent\SGARefTarget{sga3:iv:proposition:4:7:1}\textbf{Proposition 4.7.1.}
\textit{Let \(X\) be a squarable object of \(\mathscr C\), and let \(R\)
be an equivalence relation of type \((M)\) in \(X\). Suppose that
\(R\to X\times X\) belongs to \((N)\), where \((N)\) satisfies \SGARefLink{sga3:iv:subsection:4:7:item:a:01}{(a)} and
\SGARefLink{sga3:iv:subsection:4:7:item:f:mathcal:t:01}{\((f_{\mathcal T})\)}. Then the sheaf quotient \(X/R\) is given by}
\[
\begin{aligned}
 (X/R)(S)=\{Z\subseteq X_S\mid{}&
   Z\text{ is an }S\text{-subobject stable under }R\times S,\\
 & Z\to X_S\text{ belongs to }(N),\quad
   Z\to S\text{ is covering or belongs to }(M),\\
 & R_Z\simeq Z\times_S Z\,\}.
\end{aligned}
\]

We also have \SGARefLink{sga3:iv:proposition:4:7:2}{the following result}.

\medskip
\noindent\SGARefTarget{sga3:iv:proposition:4:7:2}\textbf{Proposition 4.7.2.}
\textit{Let \(X\in\operatorname{Ob}(\mathscr C)\), let \(R\) be an
\((M)\)-effective equivalence relation in \(X\), and let \((N)\) be a
family of morphisms satisfying \SGARefLink{sga3:iv:subsection:4:7:item:a:01}{(a)} and \SGARefLink{sga3:iv:subsection:4:7:item:f:m:01}{\((f_M)\)}. For every subobject
\(Y\) of \(X\) which is stable under \(R\) and for which \(Y\to X\)
belongs to \((N)\), the equivalence relation \(R_Y\) induced in \(Y\) by
\(R\) is \((M)\)-effective. Moreover, the quotient}
\[
  Y/R_Y=Y'
\]
\textit{is a subobject of \(X'=X/R\), and \(Y'\to X'\) belongs to
\((N)\).}

\textit{The map \(Y\mapsto Y'\) is a bijection from the set of subobjects
\(Y\) of \(X\) which are stable under \(R\) and for which \(Y\to X\)
belongs to \((N)\), onto the set of subobjects \(Y'\) of \(X'\) for which
\(Y'\to X'\) belongs to \((N)\). Its inverse is}
\[
  Y'\longmapsto Y'\times_{X'}X.
\]

Indeed, since \(R\) is \((M)\)-effective, the morphism \(X\to X'\)
belongs to \((M)\). Let \(Y'\) be a subobject of \(X'\) such that the
canonical morphism \(Y'\to X'\) belongs to \((N)\). Then
\(Y=Y'\times_{X'}X\) is a subobject of \(X\) stable under \(R\), and the
morphisms \(Y\to X\) and \(Y\to Y'\) belong respectively to \((N)\) and
\((M)\), since both families are stable under base change. Let \(R_Y\)
denote the equivalence relation induced in \(Y\) by \(R\). By \SGARefLink{sga3:iv:proposition:4:4:11}{4.4.11}, the
sheaf quotient \(Y/R_Y\) is represented by \(Y'\); hence \SGARefLink{sga3:iv:lemma:4:6:5}{4.6.5} shows that
\(R_Y\) is \((M)\)-effective.

Conversely, let \(Y\) be a subobject of \(X\), stable under \(R\), such
that \(Y\to X\) belongs to \((N)\). Since \(Y\) is stable under \(R\),
its two inverse images in
\[
  R=X\times_{X'}X
\]
are identical, and \(Y\) is equipped with a descent datum relative to
\(X\to X'\). The required conclusion follows from Axiom \SGARefLink{sga3:iv:subsection:4:7:item:f:m:01}{\((f_M)\)} for
the family \((N)\).

\medskip
\noindent\SGARefTarget{sga3:iv:corollary:4:7:3}\textbf{Corollary 4.7.3.}
\textit{Let \(X\in\operatorname{Ob}(\mathscr C)\), and let \(R\) be an
\((M)\)-effective equivalence relation in \(X\). Suppose in addition that
\(R\to X\times X\) belongs to \((N)\), where \((N)\) satisfies \SGARefLink{sga3:iv:subsection:4:7:item:a:01}{(a)} and
\SGARefLink{sga3:iv:subsection:4:7:item:f:mathcal:t:01}{\((f_{\mathcal T})\)}. Then, for every \(Y\) as in \SGARefLink{sga3:iv:proposition:4:7:2}{4.7.2},
\(R_Y\to Y\times Y\) also belongs to \((N)\), and consequently \SGARefLink{sga3:iv:proposition:4:7:1}{4.7.1}
gives}
\[
\begin{aligned}
 (Y/R_Y)(S)=\{Z\subseteq Y_S\mid{}&
   Z\text{ is an }S\text{-subobject stable under }R_Y\times S,\\
 & Z\to Y_S\text{ belongs to }(N)
   \text{ (hence so does }Z\to X_S),\\
 & Z\to S\text{ is covering},\quad R_Z\simeq Z\times_S Z\,\}.
\end{aligned}
\]

% END INCLUDED SGA3_English_current_buildable_source_20260823/source/inputs/IV/tex/components/19_expose_IV_s47.tex

% BEGIN INCLUDED SGA3_English_current_buildable_source_20260823/source/inputs/IV/tex/components/20_expose_IV_s5.tex
% SGA 3, Expose IV, English reconstruction.
% Scope: all of Section 5, stopping immediately before the Section 6 heading.
% Authority: Polo--Gille Exp4-13oct24.pdf, component-local pages 55--61
% (combined-reader pages 243--249).
% Jacob Reinhold's CC BY 4.0 Markdown lineage was used only as a phrasing
% comparison. The Polo--Gille PDF controls prose, formulas, numbering, and notes.

\setcounter{footnote}{47}

% Source page: Exp4 component-local p. 55; combined-reader p. 243.
\SGARefTarget{sga3:iv:section:5}\section*{5. Passage to the Quotient and Algebraic Structures}
\addcontentsline{toc}{section}{5. Passage to the Quotient and Algebraic Structures}

\SGARefTarget{sga3:iv:subsection:5:1}\subsection*{5.1. Principal Homogeneous Bundles}
\addcontentsline{toc}{subsection}{\SGARefLink{sga3:iv:subsection:5:1}{5.1}. Principal Homogeneous Bundles}

\noindent\SGARefTarget{sga3:iv:definition:5:1:0}\textbf{Definition 5.1.0.}\footnote{Editor's note: the numbering
\SGARefLink{sga3:iv:definition:5:1:0}{5.1.0} has been introduced for later reference.}
We recall (III \SGARefLink{sga3:iii:definition:0:1}{0.1}) that an object \(X\) equipped with a right action of
\(H\) is said to be \emph{formally principal homogeneous}\footnote{Editor's
note: one also says ``pseudo-torsor''; cf. EGA IV\textsubscript{4}, 16.5.15.
The more general notion of a formally homogeneous object (not necessarily
principal homogeneous) is defined in Addendum \SGARefLink{sga3:iv:definition:6:7:1}{6.7.1} at the end of this
expos\'e.} under \(H\) if the canonical morphism of functors
\[
  X\times H\longrightarrow X\times X,
  \qquad (x,h)\longmapsto(x,xh),
\]
is an isomorphism. Equivalently (cf. \emph{loc. cit.}), for every
\(S\in\operatorname{Ob}\mathscr C\), the set \(X(S)\) is formally principal
homogeneous under \(H(S)\): it is either empty or principal homogeneous
under \(H(S)\). In particular, when \(H\) acts on itself by right
translations, it is formally principal homogeneous under itself.

\medskip
\noindent\SGARefTarget{sga3:iv:definition:5:1:1}\textbf{Definition 5.1.1.}
An object \(X\) with an action of \(H\) is said to be \emph{trivial} if it
is isomorphic, as an \(H\)-object, to \(H\) with its translation action.

\medskip
\noindent\SGARefTarget{sga3:iv:proposition:5:1:2}\textbf{Proposition 5.1.2.}
\textit{If \(X\) is formally principal homogeneous under \(H\), there is an
isomorphism}
\[
  \Gamma(X)\xrightarrow{\sim}
  \operatorname{Isom}_{H\text{-obj.}}(H,X)
\]
\textit{of principal homogeneous sets under \(\Gamma(H)\).}

To a section \(x\) of \(X\), associate the morphism \(H\to X\) defined on
sets by \(h\mapsto xh\). \SGARefLink{sga3:iv:proposition:5:1:2}{The assertion} follows immediately by reduction to
the set-theoretic case.

\medskip
\noindent\SGARefTarget{sga3:iv:corollary:5:1:3}\textbf{Corollary 5.1.3.}
\textit{There is an isomorphism of \(H\)-objects}
\[
  X\xrightarrow{\sim}\operatorname{Isom}_{H\text{-obj.}}(H,X).
\]

\medskip
\noindent\SGARefTarget{sga3:iv:corollary:5:1:4}\textbf{Corollary 5.1.4.}
\textit{An object with a group action is trivial if and only if it is
formally principal homogeneous and has a section.}

% Source page: Exp4 component-local p. 56; combined-reader p. 244.
\medskip
\noindent\SGARefTarget{sga3:iv:definition:5:1:5}\textbf{Definition 5.1.5.}
Let \(\mathscr C\) be a category equipped with a topology. An \(S\)-object
\(X\) with an action of the \(S\)-group \(H\) is called a
\emph{principal homogeneous bundle under \(H\)} if it is locally trivial,
that is, if the following equivalent conditions hold:\footnote{Editor's
note: in this case one also says that \(X\) is an \(H\)-torsor.}
\begin{enumerate}
  \item[(i)]\SGARefTarget{sga3:iv:definition:5:1:5:item:i:01} the collection of morphisms \(T\to S\) for which the functor
  \(X\times_ST\) is trivial under \(H\times_ST\) is a refinement of \(S\);
  \item[(ii)]\SGARefTarget{sga3:iv:definition:5:1:5:item:ii:01} for a pretopology defining the given topology, there is a
  covering family \(\{S_i\to S\}\) such that, for every \(i\), the
  \(S_i\)-functor \(X\times_SS_i\) is trivial under the \(S_i\)-group
  functor \(H\times_SS_i\), equivalently, it has a section over \(S_i\).
\end{enumerate}


\medskip
\noindent\SGARefTarget{sga3:iv:proposition:5:1:6}\textbf{Proposition 5.1.6.}
\textit{Let \(\mathscr C\) be a category equipped with a topology
\(\mathcal T\), and let \((M)\) be a family of morphisms of \(\mathscr C\)
satisfying axioms \SGARefLink{sga3:iv:para:4:6:3:item:a:01}{(a)}--\SGARefLink{sga3:iv:para:4:6:3:item:e:mathcal:t:01}{\((e_{\mathcal T})\)} of \SGARefLink{sga3:iv:para:4:6:3}{4.6.3}. Let \(H\) be an
\(S\)-group whose structure morphism \(H\to S\) belongs to \((M)\), and let
\(P\) be an \(S\)-object with an action of \(H\). The following conditions
are equivalent:}
\begin{enumerate}
  \item[(i)]\SGARefTarget{sga3:iv:proposition:5:1:6:item:i:01} \textit{\(P\) is a principal homogeneous bundle under \(H\)
  (Definition \SGARefLink{sga3:iv:definition:5:1:5}{5.1.5});}
  \item[(ii)]\SGARefTarget{sga3:iv:proposition:5:1:6:item:ii:01} \textit{\(P\) is formally principal homogeneous under \(H\),
  and its structure morphism \(P\to S\) belongs to \((M)\);}
  \item[(iii)]\SGARefTarget{sga3:iv:proposition:5:1:6:item:iii:01} \textit{there exists a morphism \(S'\to S\) belonging to
  \((M)\) such that, after base change from \(S\) to \(S'\), \(P\) becomes
  trivial; that is, \(P\times_SS'\) is trivial under \(H\times_SS'\);}
  \item[(iv)]\SGARefTarget{sga3:iv:proposition:5:1:6:item:iv:01} \textit{\(H\) acts freely and \((M)\)-effectively on \(P\),
  and the quotient \(P/H\) is isomorphic to \(S\).}
\end{enumerate}

First, \SGARefLink{sga3:iv:proposition:5:1:6:item:ii:01}{(ii)} and \SGARefLink{sga3:iv:proposition:5:1:6:item:iv:01}{(iv)} are equivalent. In either case \(P\to S\) belongs to
\((M)\), hence is squarable, which ensures that the fiber products
\(H\times_SP\) and \(P\times_SP\) are representable. Condition \SGARefLink{sga3:iv:proposition:5:1:6:item:ii:01}{(ii)} plainly
implies \SGARefLink{sga3:iv:proposition:5:1:6:item:iii:01}{(iii)}: one may take \(S'=P\), since formal principal homogeneity
implies that \(P\times_SP\) is trivial under \(H\times_SP\) by \SGARefLink{sga3:iv:corollary:5:1:4}{5.1.4}, its
diagonal section providing a section. Condition \SGARefLink{sga3:iv:proposition:5:1:6:item:iii:01}{(iii)} plainly implies \SGARefLink{sga3:iv:proposition:5:1:6:item:i:01}{(i)},
because \(\{S'\to S\}\) is a covering family by axiom
\SGARefLink{sga3:iv:para:4:6:3:item:d:mathcal:t:01}{\((d_{\mathcal T})\)}.

It remains to prove \SGARefLink{sga3:iv:proposition:5:1:6:item:i:01}{(i)}\(\Rightarrow\)\SGARefLink{sga3:iv:proposition:5:1:6:item:ii:01}{(ii)}. The morphism of sheaves
\[
  P\times_SH\longrightarrow P\times_SP
\]
is locally an isomorphism and hence is an isomorphism (\SGARefLink{sga3:iv:corollary:4:5:8}{4.5.8}); thus \(P\) is
formally principal homogeneous. The structure morphism \(P\to S\) is
locally isomorphic to the structure morphism \(H\to S\), which belongs to
\((M)\). It therefore belongs to \((M)\) by axiom
\SGARefLink{sga3:iv:para:4:6:3:item:e:mathcal:t:01}{\((e_{\mathcal T})\)}.

The equivalence between \SGARefLink{sga3:iv:proposition:5:1:6:item:i:01}{(i)} and \SGARefLink{sga3:iv:proposition:5:1:6:item:iv:01}{(iv)} admits the following generalization.

\medskip
\noindent\SGARefTarget{sga3:iv:proposition:5:1:7}\textbf{Proposition 5.1.7.}
\textit{Under the same assumptions on \(\mathscr C\) and \((M)\), let
\(H\) be an \(S\)-group and let \(X\) be an \(S\)-object on which \(H\)
acts on the right. Suppose that \(H\to S\) belongs to \((M)\). The following
conditions are equivalent:}
\begin{enumerate}
  \item[(i)]\SGARefTarget{sga3:iv:proposition:5:1:7:item:i:01} \textit{\(H\) acts freely and \((M)\)-effectively on \(X\);}
  \item[(ii)]\SGARefTarget{sga3:iv:proposition:5:1:7:item:ii:01} \textit{there exists an \(S\)-morphism \(p:X\to Y\),
  compatible with the equivalence relation on \(X\) defined by the action
  of \(H\), such that the induced action of \(H\times_SY\) on \(X\) over
  \(Y\) makes \(X\) a principal homogeneous bundle under \(H_Y\) over
  \(Y\).}
\end{enumerate}
\textit{Under these conditions, \(p\) identifies \(Y\) with the quotient
\(X/H\).}

% Source page: Exp4 component-local p. 57; combined-reader p. 245.
If \(p:X\to Y\) is compatible with the action of \(H\), then the induced
action of \(H\times_SY\) on \(X\) over \(Y\) defines the same equivalence
relation on \(X\) as the original action of \(H\), by the isomorphism
\[
  H_Y\times_YX\xrightarrow{\sim}H\times_SX.
\]
\SGARefLink{sga3:iv:proposition:5:1:7}{The proposition} follows from this observation and the preceding equivalence
\SGARefLink{sga3:iv:proposition:5:1:6:item:iv:01}{(iv)}\(\Leftrightarrow\)\SGARefLink{sga3:iv:proposition:5:1:6:item:i:01}{(i)}.

\medskip
\noindent\SGARefTarget{sga3:iv:corollary:5:1:7:1}\textbf{Corollary 5.1.7.1.}\footnote{Editor's note: this special
case has been added as a corollary; it will be used several times in the
following expos\'es.}
\textit{Let \(\mathscr C\) be a category having a final object and fiber
products, and let it be equipped with a topology \(\mathcal T\) coarser than
the canonical topology. Let \(f:G\to H\) be a morphism of
\(\mathscr C\)-groups, put \(K=\operatorname{Ker}(f)\), and suppose that
\(f\) is covering for \(\mathcal T\). Then \(H\) represents the sheaf
quotient \(G/K\), and \(f\) is a \(K_H\)-torsor. Equivalently, one may say
that \(G\) is a \(K\)-torsor over \(H\).}

Indeed, because \(f\) is covering, it is a universal effective epimorphism
(\SGARefLink{sga3:iv:proposition:4:4:3}{4.4.3}). Hence, by \SGARefLink{sga3:iv:scholium:3:3:3:1}{3.3.3.1}, \(H\) is the quotient of \(G\) by the
equivalence relation
\[
  \mathscr R(f)=G\times_HG.
\]
Moreover, the morphism
\[
  G\times K\longrightarrow G\times_HG,
  \qquad (g,k)\longmapsto(g,gk),
\]
is an isomorphism of \(K_G=G\times_HK_H\)-objects; its inverse is
\((g,g')\mapsto(g,g^{-1}g')\). Thus \(\mathscr R(f)\) is the equivalence
relation defined by \(K\). Since \(f:G\to H\) is covering, it is a
\(K_H\)-torsor by \SGARefLink{sga3:iv:proposition:5:1:6:item:ii:01}{5.1.6(ii)}, or directly by Definition \SGARefLink{sga3:iv:definition:5:1:5:item:ii:01}{5.1.5(ii)}.

We can now sharpen Theorem \SGARefLink{sga3:iv:theorem:4:6:6}{4.6.6} in the case of a quotient by a group action.

\medskip
\noindent\SGARefTarget{sga3:iv:proposition:5:1:8}\textbf{Proposition 5.1.8.}
\textit{Under the assumptions of \SGARefLink{sga3:iv:proposition:5:1:7}{5.1.7}, let \(F_0\) denote the functor over
\(S\) defined as follows: for each \(S'\to S\), \(F_0(S')\) is the set of
representable \(S'\)-subfunctors \(Z\) of \(X\times_SS'\) that are stable
under \(H\times_SS'\) and are principal homogeneous bundles under this
\(S'\)-group for the induced action (\SGARefLink{sga3:iv:lemma:3:2:2}{3.2.2}).}
\begin{enumerate}
  \item[(i)]\SGARefTarget{sga3:iv:proposition:5:1:8:item:i:01} \textit{The following conditions are equivalent:}
  \begin{enumerate}
    \item[(a)]\SGARefTarget{sga3:iv:proposition:5:1:8:item:a:01} \textit{the action of \(H\) on \(X\) is
    \((M)\)-effective and free;}\footnote{Editor's note: the words ``and
    free'' have been added.}
    \item[(b)]\SGARefTarget{sga3:iv:proposition:5:1:8:item:b:01} \textit{\(F_0\) is representable.}
  \end{enumerate}
  \textit{Under these conditions, \(F_0=X/H\).}
  \item[(ii)]\SGARefTarget{sga3:iv:proposition:5:1:8:item:ii:01} \textit{Let \((N)\) be a family of morphisms stable under
  base change, such that for every covering family
  \(\{S_i'\to S'\}\) and every family of morphisms
  \(\{T_i\to S_i'\}\) in \((N)\), every descent datum on the \(T_i\)
  relative to \(\{S_i'\to S'\}\) is effective. Suppose that
  \(X\times_SH\to X\times_SX\) belongs to \((N)\), and that \(X\) is
  squarable. Then the element \(p\in\operatorname{Hom}(X,F_0)\)
  corresponding to the subobject \(X\times_SH\) of \(X\times_SX\)
  identifies \(F_0\) with the sheaf quotient \(X/H\).}
\end{enumerate}

% Source page: Exp4 component-local p. 58; combined-reader p. 246.
\SGARefTarget{sga3:iv:subsection:5:2}\subsection*{5.2. Group Structures and Passage to the Quotient}
\addcontentsline{toc}{subsection}{\SGARefLink{sga3:iv:subsection:5:2}{5.2}. Group Structures and Passage to the Quotient}

In this subsection we are concerned with the algebraic structures that can
be placed on the quotient (G/H) of a group by a subgroup. We first work in
the category of sheaves on \(\mathscr C\) for an arbitrary topology. Taking
the canonical topology and using \SGARefLink{sga3:iv:lemma:4:5:12}{4.5.12} will then give results concerning
universal effective passage to the quotient in \(\mathscr C\).

\medskip
\noindent\SGARefTarget{sga3:iv:proposition:5:2:1}\textbf{Proposition 5.2.1.}
\textit{Let \(u:H\to G\) be a monomorphism of sheaves of groups. There is a
unique \(G\)-object structure on the quotient sheaf \(G/H\) for which the
canonical morphism
\[
  p:G\longrightarrow G/H
\]
is a morphism of \(G\)-objects. This structure is functorial in the pair
\((G,H)\): given a commutative diagram
\[
  \begin{array}{ccc}
    H & \longrightarrow & G \\
    \downarrow && \downarrow\scriptstyle f \\
    H' & \longrightarrow & G',
  \end{array}
\]
the morphism \(f:G/H\to G'/H'\) of \SGARefLink{sga3:iv:lemma:3:2:3}{3.2.3} is compatible with the morphism
\(f\) between the acting groups.}

Indeed, \(G/H\) is the sheaf associated with the presheaf
\[
  i(G)/i(H):S\longmapsto G(S)/H(S).
\]
Since the functor \(a\) is left exact, it carries objects with a group
action to objects with a group action. The presheaf \(i(G)/i(H)\) is an
\(i(G)\)-object, and therefore
\[
  G/H=a\bigl(i(G)/i(H)\bigr)
\]
is an \(a(i(G))=G\)-object. This construction plainly has all the stated
properties.

\medskip
\noindent\SGARefTarget{sga3:iv:corollary:5:2:2}\textbf{Corollary 5.2.2.}
\textit{Let \(u:H\to G\) be a monomorphism of \(\mathscr C\)-groups, and
suppose that the action of \(H\) on \(G\) is universally effective. There
is a unique \(G\)-object structure on the quotient object
\(G/H\in\operatorname{Ob}\mathscr C\) for which \(p:G\to G/H\) is a
morphism of objects with an action. In the sense of \SGARefLink{sga3:iv:proposition:5:2:1}{the preceding
proposition}, this structure is functorial in the pair \((H,G)\), with
\(H\) acting universally effectively on \(G\).}

\medskip
\noindent\SGARefTarget{sga3:iv:proposition:5:2:3}\textbf{Proposition 5.2.3.}
\textit{Let \(u:H\to G\) be a monomorphism of sheaves of groups that
identifies \(H\) with an invariant subgroup sheaf of \(G\). There is a
unique sheaf-of-groups structure on the quotient sheaf \(G/H\) for which
the canonical morphism \(p:G\to G/H\) is a group morphism. This structure
is functorial in the pair \((H,G)\), with \(H\) invariant.}

The proof is similar to that of \SGARefLink{sga3:iv:proposition:5:2:1}{5.2.1}.

\medskip
\noindent\SGARefTarget{sga3:iv:corollary:5:2:4}\textbf{Corollary 5.2.4.}
\textit{Let \(u:H\to G\) be a monomorphism of \(\mathscr C\)-groups that
identifies \(H\) with an invariant subgroup of \(G\). Suppose that the
action of \(H\) on \(G\) is universally effective. There is a unique group
structure on the quotient object \(G/H\in\operatorname{Ob}\mathscr C\) for
which the canonical morphism \(G\to G/H\) is a group morphism. This
structure is functorial in the pair \((H,G)\), with \(H\) invariant and
acting universally effectively.}

% Source page: Exp4 component-local p. 59; combined-reader p. 247.
The group structure of \(G/H\) can be characterized more explicitly.

\medskip
\noindent\SGARefTarget{sga3:iv:proposition:5:2:5}\textbf{Proposition 5.2.5.}
\textit{Under the assumptions of \SGARefLink{sga3:iv:corollary:5:2:4}{5.2.4}, let \(K\) be a
\(\mathscr C\)-group and let \(f:G\to K\) be a morphism. The following
conditions are equivalent:}
\begin{enumerate}
  \item[(i)]\SGARefTarget{sga3:iv:proposition:5:2:5:item:i:01} \textit{\(f\) is a group morphism compatible with the
  equivalence relation defined by \(H\);}
  \item[(ii)]\SGARefTarget{sga3:iv:proposition:5:2:5:item:ii:01} \textit{\(f\) is a group morphism whose restriction
  \(H\to K\) is the trivial morphism;}
  \item[(iii)]\SGARefTarget{sga3:iv:proposition:5:2:5:item:iii:01} \textit{\(f\) factors through a group morphism
  \(G/H\to K\).}
\end{enumerate}
\textit{In particular, there is an isomorphism, functorial in the group
\(K\),}
\[
  \operatorname{Hom}_{\mathscr C\text{-gr.}}(G/H,K)
  \xrightarrow{\sim}
  \bigl\{f\in\operatorname{Hom}_{\mathscr C\text{-gr.}}(G,K)
  \mid f\circ u=e\bigr\}.
\]

The equivalence of \SGARefLink{sga3:iv:proposition:5:2:5:item:i:01}{(i)} and \SGARefLink{sga3:iv:proposition:5:2:5:item:ii:01}{(ii)} is proved set-theoretically, and
\SGARefLink{sga3:iv:proposition:5:2:5:item:iii:01}{(iii)}\(\Rightarrow\)\SGARefLink{sga3:iv:proposition:5:2:5:item:ii:01}{(ii)} is immediate. The equivalence of \SGARefLink{sga3:iv:proposition:5:2:5:item:iii:01}{(iii)} and \SGARefLink{sga3:iv:proposition:5:2:5:item:ii:01}{(ii)}
follows from the formula
\[
  \operatorname{Hom}(G/H,K)
  \simeq \operatorname{Hom}\bigl(i(G)/i(H),K\bigr)
\]
and the definition of the group structure on \(G/H\).

\medskip
\noindent\SGARefTarget{sga3:iv:remark:5:2:6}\textbf{Remark 5.2.6.}
In the preceding situation, if the kernel of \(f\) is exactly \(H\), then
the morphism \(G/H\to K\) through which \(f\) factors is a monomorphism.
This follows at once from \SGARefLink{sga3:iv:proposition:3:3:4}{3.3.4}.

For sheaves of groups, \SGARefLink{sga3:iv:proposition:4:4:11}{4.4.11} can be sharpened as follows.

\medskip
\noindent\SGARefTarget{sga3:iv:proposition:5:2:7}\textbf{Proposition 5.2.7.}
\textit{Let \(G\) be a sheaf of groups and \(H\) an invariant subgroup
sheaf. For every subgroup sheaf \(K\) of \(G\) containing \(H\), let
\(K'=K/H\), regarded as a subgroup of \(G'=G/H\). Then}
\[
  K=K'\times_{G'}G,
\]
\textit{and the maps}
\[
  K\longmapsto K/H,
  \qquad
  K'\longmapsto K'\times_{G'}G
\]
\textit{give a bijection between the subgroup sheaves of \(G\) containing
\(H\) and the subgroup sheaves of \(G'\). Under this correspondence, the
invariant subgroup sheaves of \(G\) and of \(G'\) correspond.}

The first assertion follows readily from \SGARefLink{sga3:iv:proposition:4:4:11}{4.4.11} and \SGARefLink{sga3:iv:lemma:3:2:4}{3.2.4}. It remains to
show that \(K\) is invariant in \(G\) if and only if \(K'\) is invariant in
\(G'\). If \(K\) is invariant in \(G\), then the presheaf
\(i(K)/i(H)\) is invariant in \(i(G)/i(H)\); by the usual argument, the
same is true of the associated sheaves. Conversely, if \(K'\) is invariant
in \(G'\), then the fiber product \(K'\times_{G'}G\) is invariant in
\(G\), as is immediate.

Now let \(L\) be an arbitrary subgroup sheaf of \(G\). Denote by
\(\overline L\) the saturation of \(L\) for the equivalence relation
defined by \(H\); we also write \(\overline L=L\cdot H\).

\medskip
\noindent\SGARefTarget{sga3:iv:proposition:5:2:8}\textbf{Proposition 5.2.8.}
\textit{Under the preceding assumptions, \(L\cdot H\) is a subgroup sheaf
of \(G\) containing \(H\), and the image of \(L\) in \(G/H\) is
identified with}
\[
  (L\cdot H)/H\simeq L/(H\cap L).
\]

% Source page: Exp4 component-local p. 60; combined-reader p. 248.
Indeed, let \(L'\) denote the image sheaf of \(L\) in \(G/H\). Under the
correspondence of \SGARefLink{sga3:iv:proposition:5:2:8}{the preceding proposition}, it is the subgroup sheaf of
\(G/H\) corresponding to \(L\cdot H\). Since \(L\to L'\) is covering,
and hence a universal effective epimorphism of sheaves, it follows from
\SGARefLink{sga3:iv:proposition:4:4:9}{4.4.9} that \(L'\) is identified with the quotient of \(L\) by the kernel
of \(L\to L'\), which is plainly \(H\cap L\).

Finally, consider a sheaf of groups \(G\), a subgroup sheaf \(K\), and an
invariant subgroup sheaf \(H\) of \(K\). We first define a right action of
the sheaf of groups \(H\backslash K\;(=K/H)\) on \(G/H\). The group
\(K\) acts on \(G\) by right translations. Since \(H\) is invariant in
\(K\), this action is compatible with the equivalence relation defined by
the action of \(H\), and therefore induces an action of \(K\) on \(G/H\),
that is, a morphism
\[
  K^{\circ}\longrightarrow\operatorname{Aut}(G/H)
\]
from the group opposite to \(K\). Since the target is a sheaf (\SGARefLink{sga3:iv:corollary:4:5:13}{4.5.13}) and
this morphism is trivial on \(H\), it factors through \(K/H\) and defines
the desired action. The right and left actions of \(G\) on itself commute;
hence the actions of \(G\) and \(K/H\) on \(G/H\) commute.

\medskip
\noindent\SGARefTarget{sga3:iv:proposition:5:2:9}\textbf{Proposition 5.2.9.}
\textit{Under the preceding assumptions, \(K/H\) acts freely on the right
on \(G/H\), and there is a canonical isomorphism of \(G\)-objects}
\[
  (G/H)/(K/H)\simeq G/K.
\]
\textit{If \(K\) is invariant in \(G\), in which case \(K/H\) is invariant
in \(G/H\) by \SGARefLink{sga3:iv:proposition:5:2:7}{5.2.7}, this isomorphism respects the group structures on
both sides.}

There is an isomorphism of presheaves
\[
  \bigl(i(G)/i(H)\bigr)\big/\bigl(i(K)/i(H)\bigr)
  \xleftarrow{\sim} i(G)/i(K),
\]
which respects the \(i(G)\)-object structures. Applying the functor \(a\)
to this relation gives the asserted result.

\medskip
\noindent\SGARefTarget{sga3:iv:corollary:5:2:10}\textbf{Corollary 5.2.10.}
\textit{Let \(G\) be a \(\mathscr C\)-group, \(K\) a
\(\mathscr C\)-subgroup of \(G\), and \(H\) an invariant
\(\mathscr C\)-subgroup of \(K\). Let \((M)\) be a family of morphisms of
\(\mathscr C\) satisfying axioms \SGARefLink{sga3:iv:para:4:6:3:item:a:01}{(a)}--\SGARefLink{sga3:iv:para:4:6:3:item:e:mathcal:t:01}{\((e_{\mathcal T})\)}. Suppose that
the right action of \(H\) on \(G\) (respectively, on \(K\)) is
\((M)\)-effective. Then \(K/H\) acts naturally and freely on the right on
\(G/H\), and this action commutes with that of \(G\). The following
conditions are equivalent:}
\begin{enumerate}
  \item[(i)]\SGARefTarget{sga3:iv:corollary:5:2:10:item:i:01} \textit{the action of \(K\) on \(G\) is
  \((M)\)-effective;}
  \item[(ii)]\SGARefTarget{sga3:iv:corollary:5:2:10:item:ii:01} \textit{the action of \(K/H\) on \(G/H\) is
  \((M)\)-effective.}
\end{enumerate}
\textit{Under these conditions, there is an isomorphism of
objects\footnote{Editor's note: of \(\mathscr C\).} with a \(G\)-action}
\[
  (G/H)/(K/H)\simeq G/K.
\]

% Source page: Exp4 component-local p. 61; combined-reader p. 249.
\SGARefTarget{sga3:iv:subsection:5:3}\subsection*{5.3. Use of Effectivity Criteria: Noether's Theorem}
\addcontentsline{toc}{subsection}{\SGARefLink{sga3:iv:subsection:5:3}{5.3}. Use of Effectivity Criteria: Noether's Theorem}
\footnote{Editor's note: in fact, the isomorphisms established in
\SGARefLink{sga3:iv:proposition:5:2:8}{5.2.8}--\SGARefLink{sga3:iv:corollary:5:2:10}{5.2.10} deserve to be called ``Noether isomorphism theorems.''}
Let \(\mathscr C\), \(\mathcal T\), and \((M)\) be as usual. Let \((N)\)
be a family of morphisms satisfying axioms \SGARefLink{sga3:iv:subsection:4:7:item:a:01}{(a)} and \SGARefLink{sga3:iv:subsection:4:7:item:f:m:01}{\((f_M)\)} of \SGARefLink{sga3:iv:subsection:4:7}{4.7}.
Combining \SGARefLink{sga3:iv:proposition:5:2:7}{5.2.7} and \SGARefLink{sga3:iv:proposition:4:7:2}{4.7.2} gives \SGARefLink{sga3:iv:proposition:5:3:1}{the following result}.

\medskip
\noindent\SGARefTarget{sga3:iv:proposition:5:3:1}\textbf{Proposition 5.3.1.}
\textit{Let \(G\) be a \(\mathscr C\)-group, and let \(H\) be an invariant
\(\mathscr C\)-subgroup of \(G\) that acts \((M)\)-effectively on \(G\).
For every \(\mathscr C\)-subgroup \(K\) of \(G\) containing \(H\) such
that \(K\to G\) belongs to \((N)\), the action of \(H\) on \(K\) is
\((M)\)-effective, and the quotient \(K/H=K'\) is a
\(\mathscr C\)-subgroup of \(G/H=G'\) such that \(K'\to G'\) belongs to
\((N)\). The map \(K\mapsto K'\) is a bijection between the set of
\(\mathscr C\)-subgroups \(K\) of \(G\) that contain \(H\) and for which
\(K\to G\) belongs to \((N)\), and the set of \(\mathscr C\)-subgroups
\(K'\) of \(G'\) for which \(K'\to G'\) belongs to \((N)\). Its inverse
is}
\[
  K'\longmapsto K\times_{G'}G.
\]
\textit{Under this correspondence, the invariant subgroups of \(G\) and
of \(G'\) correspond.}

\medskip
\noindent\SGARefTarget{sga3:iv:corollary:5:3:2}\textbf{Corollary 5.3.2.}
\textit{If \(H\to G\) belongs to \((N)\), then \(\mathscr C\) has a final
object \(e\), and the unit section \(e\to G/H\) belongs to
\((N)\).}\footnote{Editor's note: this follows by applying \SGARefLink{sga3:iv:proposition:5:3:1}{the proposition}
with \(K=H\).}

% END INCLUDED SGA3_English_current_buildable_source_20260823/source/inputs/IV/tex/components/20_expose_IV_s5.tex

% BEGIN INCLUDED SGA3_English_current_buildable_source_20260823/source/inputs/IV/tex/components/21_expose_IV_s6_and_bibliography.tex
% SGA 3, Expose IV, English reconstruction.
% Scope: section 6 in full, followed by the bibliography; stops at the end
% of Expose IV and contains no material from Expose V.
% Authority: Polo--Gille Exp4-13oct24.pdf, component-local pages 61--73
% (combined-reader pages 249--261; printed pages 239--250).
% Jacob Reinhold's CC BY 4.0 Markdown lineage was used only as a phrasing
% and structure comparison. The Polo--Gille PDF controls all prose,
% numbering, editorial notes, formulas, diagrams, and references.

% Source page: Exp4 component-local p. 61; combined-reader p. 249.
\SGARefTarget{sga3:iv:section:6}\section*{6. Topologies in the Category of Schemes}
\addcontentsline{toc}{section}{6. Topologies in the Category of Schemes}

\SGARefTarget{sga3:iv:subsection:6:1}\subsection*{6.1. The Zariski Topology}
\addcontentsline{toc}{subsection}{\SGARefLink{sga3:iv:subsection:6:1}{6.1}. The Zariski Topology}

This is the topology generated by the following pretopology: a family of
morphisms \(\{S_i\to S\}\) is covering if each morphism is an open
immersion and the union of the images of the \(S_i\) is all of \(S\). It is
denoted by \((\mathrm{Zar})\).

\medskip
\noindent\SGARefTarget{sga3:iv:definition:6:1:1}\textbf{Definition 6.1.1.}
\textit{A sheaf for the Zariski topology is also called a \emph{functor of
local nature}: it is a contravariant functor from \(\mathbf{Sch}\) to
\(\mathbf{Set}\) such that, for every scheme \(S\) and every covering of
\(S\) by open subschemes \(S_i\), the diagram}
\[
  F(S)\longrightarrow\prod_iF(S_i)
  \rightrightarrows\prod_{i,j}F(S_i\cap S_j)
\]
\textit{is exact.}

In particular, a functor of local nature carries direct sums to products.
Since every representable functor is a sheaf, the Zariski topology is
coarser than the canonical topology.

As a matter of terminology, whenever we say "local" or "locally" without
further qualification, the reference is to the Zariski topology and hence
to the usual meaning of these words.

% Source page: Exp4 component-local p. 62; combined-reader p. 250.
\SGARefTarget{sga3:iv:subsection:6:2}\subsection*{6.2. A Procedure for Constructing Topologies}
\addcontentsline{toc}{subsection}{\SGARefLink{sga3:iv:subsection:6:2}{6.2}. A Procedure for Constructing Topologies}

\medskip
\noindent\SGARefTarget{sga3:iv:proposition:6:2:1}\textbf{Proposition 6.2.1.}
\textit{Let \(\mathscr C\) be a category, let \(\mathscr C'\) be a full
subcategory, and let \(P\) be a set of families of morphisms in
\(\mathscr C\) with common target, stable under base change and composition
(that is, satisfying axioms \SGARefLink{sga3:iv:definition:4:2:5:item:p:1:01}{(P~1)} and \SGARefLink{sga3:iv:definition:4:2:5:item:p:2:01}{(P~2)} of \SGARefLink{sga3:iv:definition:4:2:5}{4.2.5}). Let \(P'\) be a set
of families of morphisms in \(\mathscr C'\) that contains every singleton
family consisting of an identity isomorphism. Equip \(\mathscr C\) with
the topology generated by \(P\) and \(P'\) (cf. \SGARefLink{sga3:iv:definition:4:2:5:0}{4.2.5.0}), and suppose that
the following three conditions hold:}

\begin{enumerate}
  \item[(a)]\SGARefTarget{sga3:iv:proposition:6:2:1:item:a:01} \textit{If \(\{S_i\to S\}\in P'\), so that
    \(S_i,S\in\operatorname{Ob}\mathscr C'\), and \(T\to S\) is a
    morphism in \(\mathscr C'\), then the fiber products
    \(S_i\times_ST\), formed in \(\mathscr C\), exist and the family
    \(\{S_i\times_ST\to T\}\) belongs to \(P'\); in particular,
    \(S_i\times_ST\in\operatorname{Ob}\mathscr C'\).}

    \textit{\SGARefLink{sga3:iv:proposition:6:2:1:item:a:01}{This condition} implies that \(P'\) is stable under base change
    in \(\mathscr C'\), but is not equivalent to that assertion: it also
    requires the inclusion functor \(\mathscr C'\hookrightarrow\mathscr C\)
    to commute with certain fiber products.}

  \item[(b)]\SGARefTarget{sga3:iv:proposition:6:2:1:item:b:01} \textit{For every \(S\in\operatorname{Ob}\mathscr C\), there
    exists a family \(\{S_i\to S\}\in P\) with
    \(S_i\in\operatorname{Ob}\mathscr C'\) for every \(i\).}

  \item[(c)]\SGARefTarget{sga3:iv:proposition:6:2:1:item:c:01} \textit{In a situation}
    \[
    \begin{array}{ccccc}
      & & S_{ijk} & &\\
      & & \scriptstyle(P')\,\Big\downarrow & &\\[-0.25em]
      S_i & \xleftarrow{\ (P)\ } & S_{ij} & &\\
      \scriptstyle(P')\,\Big\downarrow & & & &\\[-0.25em]
      S & & & &
    \end{array}
    \tag{6.2.1.c}\label{sga3:iv:proposition:6:2:1:eq:6:2:1:c:01}
    \]
    \textit{where \(S,S_i,S_{ij},S_{ijk}\) are objects of
    \(\mathscr C'\), where \(\{S_i\to S\}\in P'\),
    \(\{S_{ij}\to S_i\}\in P\) for each \(i\), and
    \(\{S_{ijk}\to S_{ij}\}\in P'\) for each pair \(ij\), there exists a
    family \(\{T_n\to S\}\in P'\), and for each \(n\) a multi-index
    \((ijk)\), fitting into a commutative diagram}
    \[
    \begin{array}{ccccc}
      S_{ijk} & \longleftarrow & T_n\\
      & \searrow & \Big\downarrow\\[-0.25em]
      & & S .
    \end{array}
    \tag{6.2.1.d}\label{sga3:iv:proposition:6:2:1:eq:6:2:1:d:01}
    \]
\end{enumerate}

\textit{Under these hypotheses, a sieve \(R\) of an object
\(S\in\operatorname{Ob}\mathscr C\) is covering if and only if there is a
composite family}
\[
\begin{array}{ccccc}
  R & \dashleftarrow & S_{ij}\\
  \Big\downarrow & & \scriptstyle(P')\,\Big\downarrow\\[-0.25em]
  S & \xleftarrow{\ (P)\ } & S_i
\end{array}
\tag{6.2.1.e}\label{sga3:iv:proposition:6:2:1:eq:6:2:1:e:01}
\]
\textit{where \(S_i,S_{ij}\in\operatorname{Ob}\mathscr C'\),
\(\{S_i\to S\}\in P\), and \(\{S_{ij}\to S_i\}\in P'\) for every \(i\),
such that the resulting morphisms \(S_{ij}\to S\) factor through \(R\).
Equivalently, the sieve generated by this composite family is contained in
\(R\).}

\smallskip
\noindent\textit{Proof.}
Families belonging to \(P\) and \(P'\) are covering, and therefore so is a
composite of such families, by \SGARefLink{sga3:iv:definition:4:2:3:item:c:2:01}{(C~2)}. A sieve of the displayed form is
thus covering because it contains a covering sieve.

% Source page: Exp4 component-local p. 63; combined-reader p. 251.
Conversely, it is enough to show that the sieves of this form constitute a
topology, that is, to verify axioms \SGARefLink{sga3:iv:definition:4:2:1:item:t:1:01}{(T~1)}--\SGARefLink{sga3:iv:definition:4:2:1:item:t:4:01}{(T~4)} of \SGARefLink{sga3:iv:definition:4:2:1}{4.2.1}.

\smallskip
\noindent\textit{Axiom \SGARefLink{sga3:iv:definition:4:2:1:item:t:4:01}{(T~4)}.}
Let \(S\in\operatorname{Ob}\mathscr C\). By \SGARefLink{sga3:iv:proposition:6:2:1:item:b:01}{(b)}, there exists a family
\(\{S_i\to S\}\in P\) with
\(S_i\in\operatorname{Ob}\mathscr C'\). The singleton families
\(\{S_i\xrightarrow{\operatorname{id}_{S_i}}S_i\}\) belong to \(P'\) by
hypothesis. Hence the maximal sieve \(S\) of \(S\) has the required form:
\[
\begin{array}{ccc}
  S_i & \longrightarrow & S\\
  \scriptstyle(P')\,\Big\downarrow\rlap{\(\scriptstyle\operatorname{id}\)}
      & & \Big\downarrow\rlap{\(\scriptstyle\operatorname{id}_S\)}\\[-0.25em]
  S_i & \xrightarrow{\ (P)\ } & S .
\end{array}
\tag{6.2.1.T4}\label{sga3:iv:proposition:6:2:1:eq:6:2:1:t4:01}
\]

\smallskip
\noindent\textit{Axiom \SGARefLink{sga3:iv:definition:4:2:1:item:t:3:01}{(T~3)}.}
This is evident.

\smallskip
\noindent\textit{Axiom \SGARefLink{sga3:iv:definition:4:2:1:item:t:2:01}{(T~2)}.}
Let \(R\) be a sieve of \(S\) of the required form, and let \(C\) be a
sieve\footnote{Editor's note: the original has been corrected here.} of
\(S\) such that, for every \(T\in\operatorname{Ob}\mathscr C\) and every
morphism \(T\to S\) factoring through \(R\), the sieve \(C\times_ST\) of
\(T\) has the required form. Since \(S_{ij}\to S\) factors through \(R\),
the sieve
\[
  C_{ij}=C\times_SS_{ij}
\]
of \(S_{ij}\) has the required form; it is represented schematically by
\[
\begin{array}{ccccc}
    & & S_{ij} & \longleftarrow & C_{ij}\\
    & \swarrow & \scriptstyle(P')\,\Big\downarrow & &\\[-0.25em]
    & & S_i & & \Bigg\downarrow\\[-0.25em]
    & & \scriptstyle(P)\,\Big\downarrow & &\\[-0.25em]
  R & \longrightarrow & S & \longleftarrow & C .
\end{array}
\tag{6.2.1.T2a}\label{sga3:iv:proposition:6:2:1:eq:6:2:1:t2a:01}
\]
Consequently, for every pair \(ij\), there is a diagram of the form
\[
\begin{array}{ccccc}
  S_{ijkl} & & & &\\
  \scriptstyle(P')\,\Big\downarrow & \searrow & & &\\[-0.25em]
  S_{ijk} & & & &\\
  \scriptstyle(P)\,\Big\downarrow & & & &\\[-0.25em]
  S_{ij} & \longleftarrow & C_{ij}. & &
\end{array}
\tag{6.2.1.T2b}\label{sga3:iv:proposition:6:2:1:eq:6:2:1:t2b:01}
\]
We have therefore produced a composite family
\[
  S_{ijkl}\xrightarrow{\ (P')\ }S_{ijk}
  \xrightarrow{\ (P)\ }S_{ij}
  \xrightarrow{\ (P')\ }S_i
  \xrightarrow{\ (P)\ }S
\tag{6.2.1.T2c}\label{sga3:iv:proposition:6:2:1:eq:6:2:1:t2c:01}
\]
belonging to \(P\circ P'\circ P\circ P'\), factoring through \(C\), and
having every object other than \(S\) in \(\mathscr C'\).

% Source page: Exp4 component-local p. 64; combined-reader p. 252.
Apply condition \SGARefLink{sga3:iv:proposition:6:2:1:item:c:01}{(c)} to each family \(\{S_{ijkl}\to S_i\}\). For every \(i\)
we obtain a family \(\{T_{in}\to S_i\}\in P'\) such that \(T_{in}\to S\)
factors through one of the \(S_{ijkl}\), and hence through \(C\):
\[
\begin{array}{ccccc}
  T_{in} & \longrightarrow & S_{ijkl}\\
  \scriptstyle(P')\,\Big\downarrow & & \Big\downarrow\\[-0.25em]
  S_i & & C\\
  \scriptstyle(P)\,\Big\downarrow & \swarrow &\\[-0.25em]
  S & &
\end{array}
\tag{6.2.1.T2d}\label{sga3:iv:proposition:6:2:1:eq:6:2:1:t2d:01}
\]
Thus the sieve \(C\) of \(S\) has the required form, completing the
verification of \SGARefLink{sga3:iv:definition:4:2:1:item:t:2:01}{(T~2)}.

\smallskip
\noindent\textit{Axiom \SGARefLink{sga3:iv:definition:4:2:1:item:t:1:01}{(T~1)}.}
Let \(R\) be a sieve of \(S\) of the prescribed form and let \(T\to S\) be
a morphism in \(\mathscr C\). We show that the sieve \(R\times_ST\) of
\(T\) again has the prescribed form. Consider the commutative configuration
\[
\begin{array}{ccccccc}
  & & S_{ij} & \multicolumn{3}{c}{\xleftarrow{\hspace{6em}}} & U_{ikj}\\
  & \swarrow & \scriptstyle(P')\,\Big\downarrow & & & &
    \scriptstyle(P')\,\Big\downarrow\\[-0.25em]
  & & S_i & \longleftarrow & T_i & \xleftarrow{\ (P)\ } & U_{ik}\\
  & & \scriptstyle(P)\,\Big\downarrow & &
    \scriptstyle(P)\,\Big\downarrow & \swarrow_{\scriptstyle(P)} &\\[-0.25em]
  R & \longrightarrow & S & \longleftarrow & T & & .
\end{array}
\tag{6.2.1.T1a}\label{sga3:iv:proposition:6:2:1:eq:6:2:1:t1a:01}
\]
Put \(T_i=S_i\times_ST\). The family \(\{T_i\to T\}\) belongs to \(P\),
by \SGARefLink{sga3:iv:definition:4:2:5:item:p:1:01}{(P~1)}. By \SGARefLink{sga3:iv:proposition:6:2:1:item:b:01}{(b)}, choose families \(\{U_{ik}\to T_i\}\in P\) with
\(U_{ik}\in\operatorname{Ob}\mathscr C'\). By \SGARefLink{sga3:iv:definition:4:2:5:item:p:2:01}{(P~2)},
\(\{U_{ik}\to T\}\in P\). Condition \SGARefLink{sga3:iv:proposition:6:2:1:item:a:01}{(a)} shows that
\[
  U_{ik}\times_{S_i}S_{ij}=U_{ikj}
\]
is an object of \(\mathscr C'\) and that, for each \(ik\),
\(\{U_{ikj}\to U_{ik}\}\in P'\). The commutative diagram
\[
\begin{array}{ccccc}
  R & \longleftarrow & U_{ikj}\\
  \Big\downarrow & & \scriptstyle(P')\,\Big\downarrow\\[-0.25em]
  S & \longleftarrow & T & \xleftarrow{\ (P)\ } U_{ik}
\end{array}
\tag{6.2.1.T1b}\label{sga3:iv:proposition:6:2:1:eq:6:2:1:t1b:01}
\]
shows that the morphisms \(U_{ikj}\to T\) factor through the sieve
\(T\times_SR\) of \(T\). This sieve therefore has the required form, and
the proof is complete.

\medskip
\noindent\SGARefTarget{sga3:iv:corollary:6:2:2}\textbf{Corollary 6.2.2.}
\textit{Let \(S\in\operatorname{Ob}\mathscr C'\), and let \(R\) be a
sieve of \(S\). Then \(R\) is covering if and only if there exists a family
\(\{T_i\to S\}\in P'\) that factors through \(R\).}

Such a sieve is plainly covering. Conversely, apply condition \SGARefLink{sga3:iv:proposition:6:2:1:item:c:01}{(c)} with the
family \(\{S_i\to S\}\) reduced to the identity isomorphism of \(S\); \SGARefLink{sga3:iv:proposition:6:2:1}{the
proposition} then shows that every covering sieve has the stated form.

% Source page: Exp4 component-local p. 65; combined-reader p. 253.
\medskip
\noindent\SGARefTarget{sga3:iv:corollary:6:2:3}\textbf{Corollary 6.2.3.}
\textit{A presheaf \(F\) on \(\mathscr C\) is separated (respectively, is
a sheaf) if and only if the map}
\[
  F(S)\longrightarrow\prod_iF(S_i)
\]
\textit{is injective (respectively, the diagram}
\[
  F(S)\longrightarrow\prod_iF(S_i)
  \rightrightarrows\prod_{i,j}F(S_i\times_SS_j)
\]
\textit{is exact) in each of the following two cases:}
\begin{enumerate}
  \item[(i)]\SGARefTarget{sga3:iv:corollary:6:2:3:item:i:01} \textit{\(\{S_i\to S\}\in P\);}
  \item[(ii)]\SGARefTarget{sga3:iv:corollary:6:2:3:item:ii:01} \textit{\(S,S_i\in\operatorname{Ob}\mathscr C'\) and
    \(\{S_i\to S\}\in P'\).}
\end{enumerate}

The conditions are necessary because the families in question are
covering. Let \(C\) be the sieve of \(S\) that is the image of a composite
family
\[
  \{S_{ij}\xrightarrow{\ (P')\ }S_i\xrightarrow{\ (P)\ }S\}.
\]
A straightforward diagram chase shows that the conditions of \SGARefLink{sga3:iv:corollary:6:2:3}{the
corollary} imply that
\[
  \operatorname{Hom}(S,F)
  \xrightarrow{\ g\ }\operatorname{Hom}(C,F)
\]
is injective, respectively bijective. Every covering sieve \(R\) of \(S\)
contains a sieve \(C\) of this form, and there is a commutative diagram
\[
\begin{array}{ccccc}
  \operatorname{Hom}(S,F) & \xrightarrow{\ f\ } &
    \operatorname{Hom}(R,F)\\
  & \searrow_{\scriptstyle g} & \Big\downarrow\rlap{\(\scriptstyle h\)}\\[-0.25em]
  & & \operatorname{Hom}(C,F).
\end{array}
\tag{6.2.3.1}\label{sga3:iv:corollary:6:2:3:eq:6:2:3:1:01}
\]
Since \(g\) is injective, so is \(f\), and hence \(F\) is separated.
Moreover, \(R\) refines \(C\), so \(h\) is also injective. If \(g\) is
bijective, then \(f\) is bijective as well, and \(F\) is a sheaf.

\medskip
\noindent\SGARefTarget{sga3:iv:remark:6:2:4}\textbf{Remark 6.2.4.}
\textit{\SGARefLink{sga3:iv:corollary:6:2:3}{The preceding corollary} does not follow from \SGARefLink{sga3:iv:corollary:4:3:5}{4.3.5}, since \(P'\) is
not stable under base extension.}

\medskip
\noindent\SGARefTarget{sga3:iv:remark:6:2:5}\textbf{Remark 6.2.5.}
\textit{Condition \SGARefLink{sga3:iv:proposition:6:2:1:item:c:01}{(c)} holds in particular when:}
\begin{enumerate}
  \item[(i)]\SGARefTarget{sga3:iv:remark:6:2:5:item:i:01} \textit{\(P'\) is stable under composition;}
  \item[(ii)]\SGARefTarget{sga3:iv:remark:6:2:5:item:ii:01} \textit{whenever \(\{S_j\to S\}\) is a family of morphisms in
    \(\mathscr C'\) belonging to \(P\), it contains a subfamily belonging
    to \(P'\).}
\end{enumerate}

\SGARefTarget{sga3:iv:subsection:6:3}\subsection*{6.3. Application to the Category of Schemes}
\addcontentsline{toc}{subsection}{\SGARefLink{sga3:iv:subsection:6:3}{6.3}. Application to the Category of Schemes}

Take \(\mathscr C\) to be the category of schemes, \(\mathscr C'\) the full
subcategory of affine schemes, and \(P\) the set of surjective families of
open immersions. We shall consider the following sets \(P'_i\):
\begin{align*}
 P'_1 &: \text{finite surjective families of flat morphisms};\\
 P'_2 &: \text{finite surjective families of flat, finitely presented,
          quasi-finite morphisms};\\
 P'_3 &: \text{finite surjective families of \'{e}tale morphisms};\\
 P'_4 &: \text{finite surjective families of finite \'{e}tale morphisms}.
\end{align*}
The description of \(P'_2\) is accompanied by the following
clarification.\footnote{Editor's note: let \(P''_2\) be the set of finite
surjective families of morphisms in \(\mathscr C'\) that are flat and of
finite presentation. By Proposition~\SGARefLink{sga3:iv:proposition:6:3:1}{6.3.1} below, the topology \(T_2\)
generated by \(P\) and \(P'_2\) coincides with the topology generated by
\(P\) and \(P''_2\). This follows from the results of
EGA~IV\textsubscript{4}, \S17.16 on quasi-sections; see the proof of \SGARefLink{sga3:iv:proposition:6:3:1}{6.3.1}.
We recall the following special case of EGA~IV\textsubscript{4}, 17.16.2.
Let \(S\) be affine and let \(f:X\to S\) be a surjective flat morphism
locally of finite presentation. Then there exist a faithfully flat,
finitely presented, quasi-finite morphism \(S'\to S\), with \(S'\) affine,
and an \(S\)-morphism \(S'\to X\).}

% Source page: Exp4 component-local p. 66; combined-reader p. 254.
For each \(P'_i\), except \(P'_4\), the conditions of Proposition~\SGARefLink{sga3:iv:proposition:6:2:1}{6.2.1}
hold. Condition \SGARefLink{sga3:iv:proposition:6:2:1:item:c:01}{(c)} follows from \SGARefLink{sga3:iv:remark:6:2:5}{6.2.5}: since an affine scheme is
quasi-compact, every family of morphisms in \(\mathscr C'\) belonging to
\(P\) contains a finite subfamily that still belongs to \(P\), and hence
belongs to \(P'_i\) for \(i=1,2,3\). The topology \(T_i\) generated by
\(P\) and \(P'_i\) is denoted and named as follows:\footnotemark[\value{footnote}]
\[
\begin{array}{ccl}
 T_1&=(\mathrm{fpqc})&
   \text{the faithfully flat quasi-compact topology},\\
 T_2&=(\mathrm{fppf})&
   \text{the faithfully flat topology (locally of finite presentation)},\\
 T_3&=(\mathrm{\acute et})&
   \text{the \'{e}tale topology},\\
 T_4&=(\mathrm{\acute etf})&
   \text{the finite \'{e}tale topology}.
\end{array}
\]
Since \(P'_1\supset P'_2\supset P'_3\supset P'_4\), one has
\[
  (\mathrm{fpqc})\geq(\mathrm{fppf})\geq(\mathrm{\acute et})
  \geq(\mathrm{\acute etf})\geq(\mathrm{Zar}).
\]

\medskip
\noindent\SGARefTarget{sga3:iv:proposition:6:3:1}\textbf{Proposition 6.3.1.}
\begin{enumerate}
  \item[(i)]\SGARefTarget{sga3:iv:proposition:6:3:1:item:i:01} \textit{A sieve \(R\) of \(S\) is covering for \(T_i\),
    \(1\leq i\leq3\), if and only if there exist a covering
    \((S_p)\) of \(S\) by affine open subschemes and, for every \(p\), a
    family \(\{S_{pq}\to S_p\}\in P'_i\), with all \(S_{pq}\) affine,
    such that every morphism \(S_{pq}\to S\) factors through \(R\).}
    \footnote{Editor's note: every family
    \(\{S_{pq}\to S_p\}\in P'_i\) is finite. Hence
    \(S'_p=\coprod_qS_{pq}\) is affine, and the family may be replaced by
    the morphism \(S'_p\to S_p\), which still belongs to \(P'_i\).}

  \item[(ii)]\SGARefTarget{sga3:iv:proposition:6:3:1:item:ii:01} \textit{A presheaf \(F\) on \(\mathbf{Sch}\) is a sheaf for
    \((\mathrm{fpqc})\) (respectively, \((\mathrm{fppf})\),
    \((\mathrm{\acute et})\), or \((\mathrm{\acute etf})\)) if and only
    if:}
    \begin{enumerate}
      \item[(a)]\SGARefTarget{sga3:iv:proposition:6:3:1:item:a:01} \textit{\(F\) is a sheaf for \((\mathrm{Zar})\), that is,
        a functor of local nature;}
      \item[(b)]\SGARefTarget{sga3:iv:proposition:6:3:1:item:b:01} \textit{for every faithfully flat (respectively,
        faithfully flat, finitely presented and quasi-finite; surjective
        \'{e}tale; or surjective finite \'{e}tale) morphism \(T\to S\),
        with \(T\) and \(S\) affine, the diagram}
        \[
          F(S)\longrightarrow F(T)\rightrightarrows F(T\times_ST)
        \]
        \textit{is exact.}
    \end{enumerate}

  \item[(iii)]\SGARefTarget{sga3:iv:proposition:6:3:1:item:iii:01} \textit{The topologies \(T_i\), for \(i=1,2,3,4\), are
    coarser than the canonical topology.}

  \item[(iv)]\SGARefTarget{sga3:iv:proposition:6:3:1:item:iv:01} \textit{Every surjective family of flat open morphisms
    (respectively, flat morphisms locally of finite presentation;
    \'{e}tale morphisms; or finite \'{e}tale morphisms) is covering for
    \((\mathrm{fpqc})\) (respectively, \((\mathrm{fppf})\),
    \((\mathrm{\acute et})\), or \((\mathrm{\acute etf})\)).}

  \item[(v)]\SGARefTarget{sga3:iv:proposition:6:3:1:item:v:01} \textit{Every finite surjective family of flat quasi-compact
    morphisms is covering for \((\mathrm{fpqc})\). In particular, every
    faithfully flat quasi-compact morphism is covering for
    \((\mathrm{fpqc})\).}
\end{enumerate}

\smallskip
\noindent\textit{Proof.}
Part \SGARefLink{sga3:iv:proposition:6:3:1:item:i:01}{(i)} follows from \SGARefLink{sga3:iv:proposition:6:2:1}{6.2.1} and part \SGARefLink{sga3:iv:proposition:6:3:1:item:ii:01}{(ii)} from \SGARefLink{sga3:iv:corollary:6:2:3}{6.2.3}, taking into account
that a sheaf for the Zariski topology carries direct sums to products.
Every representable functor is a sheaf for \((\mathrm{Zar})\) and satisfies
condition \SGARefLink{sga3:iv:proposition:6:3:1:item:ii:01}{(ii)}\SGARefLink{sga3:iv:proposition:6:3:1:item:b:01}{(b)}, by \SGARefLink{sga3:iv:bib:sga:1}{SGA~1}, VIII, 5.3. Thus \(T_1\) is coarser than the
canonical topology, which proves \SGARefLink{sga3:iv:proposition:6:3:1:item:iii:01}{(iii)}.

% Source page: Exp4 component-local p. 67; combined-reader p. 255.
We prove \SGARefLink{sga3:iv:proposition:6:3:1:item:iv:01}{(iv)}. Let \(\{S_i\to S\}\) be a family as in the statement. By
covering \(S\) with affine open subschemes, we reduce immediately to the
case in which \(S\) is affine.\footnote{Editor's note: the argument that
follows has been simplified by using the fact that \(S\) is now assumed
affine.}

First suppose that the morphisms \(S_i\to S\) are flat and open
(respectively, \'{e}tale). Cover each \(S_i\) by affine open subschemes
\(S_{ij}\). Since the morphisms are open, the images \(T_{ij}\) of the
\(S_{ij}\) form an open covering of \(S\). The affine scheme \(S\) is
quasi-compact, so finitely many \(T_{ij}\), indexed by a finite set \(E\),
cover it. Then
\[
  S'=\coprod_{(i,j)\in E}S_{ij}
\]
is affine, and \(S'\to S\) belongs to \(P'_1\), respectively \(P'_3\), and
is therefore covering. Since it factors through the given family, that
family is covering.

In the \((\mathrm{\acute etf})\) case, every \(S_i\) is finite over \(S\)
and hence affine. In the preceding argument one may take the singleton
covering \(\{S_i\}\) of \(S_i\), which gives \(S'\to S\in P'_4\).

Now suppose that the morphisms \(f_i:S_i\to S\) are flat and locally of
finite presentation. For every \(s\in S\), the proof of
EGA~IV\textsubscript{4}, 17.16.2 supplies an affine subscheme \(X(s)\) of
some \(S_i\), with \(s\in f_i(X(s))\), such that the restriction
\[
  g_i:X(s)\longrightarrow S
\]
is flat, finitely presented, and quasi-finite. Its image \(g_i(X(s))\) is
an open neighborhood \(U(s)\) of \(s\), by EGA~IV\textsubscript{2}, 2.4.6.
Since \(S\) is affine, finitely many such opens
\(U(s_j)\), \(j=1,\ldots,n\), cover \(S\). Consequently,
\[
  X'=\coprod_{j=1}^{n}X(s_j)
\]
is affine and \(X'\to S\) is surjective, flat, finitely presented, and
quasi-finite, so it belongs to \(P'_2\).\footnote{Editor's note: this shows
that if \(P''_2\) denotes the set of finite surjective families of flat,
finitely presented morphisms in \(\mathscr C'\), then the topology
generated by \(P\) and \(P''_2\) equals \(T_2\). Moreover, in the notation
used at the beginning of the proof of \SGARefLink{sga3:iv:proposition:6:3:1:item:iv:01}{(iv)}, if each \(S_i\) is covered by
affine opens of finite presentation over \(S\), then the resulting
morphism \(S'\to S\) belongs to \(P''_2\).}
This completes the proof of \SGARefLink{sga3:iv:proposition:6:3:1:item:iv:01}{(iv)}.

We prove \SGARefLink{sga3:iv:proposition:6:3:1:item:v:01}{(v)}. Let \(\{S_i\to S\}\) be a finite surjective family of flat
quasi-compact morphisms, and cover \(S\) by affine open subschemes \(T_j\).
The schemes
\[
  S_{ij}=T_j\times_SS_i
\]
are quasi-compact, and therefore admit finite affine open coverings
\(T_{ijk}\). Each \(T_{ijk}\to T_j\) is flat, and the family
\(\{T_{ijk}\to T_j\}\) is finite and surjective, hence covering for \(T_1\).
By composition, \(\{T_{ijk}\to S\}\) is covering as well. It factors
through the given family, which is therefore covering:
\[
\begin{array}{ccccc}
  S_i & \longleftarrow & S_{ij} & \longleftarrow & T_{ijk}\\
  \Big\downarrow & & \Big\downarrow & \swarrow &\\[-0.25em]
  S & \longleftarrow & T_j & &
\end{array}
\tag{6.3.1.v}\label{sga3:iv:proposition:6:3:1:eq:6:3:1:v:01}
\]

% Source page: Exp4 component-local p. 68; combined-reader p. 256.
\medskip
\noindent\SGARefTarget{sga3:iv:corollary:6:3:2}\textbf{Corollary 6.3.2.}
\textit{Let \((M_i)\) be the following families of morphisms:}
\begin{align*}
 (M_1)&:\ \text{faithfully flat quasi-compact morphisms};\\
 (M_2)&:\ \text{faithfully flat morphisms locally of finite presentation};\\
 (M_3)&:\ \text{surjective \'{e}tale morphisms};\\
 (M_4)&:\ \text{surjective finite \'{e}tale morphisms}.
\end{align*}
\textit{The family \((M_i)\) satisfies axioms \SGARefLink{sga3:iv:para:4:6:3:item:a:01}{(a)}, \SGARefLink{sga3:iv:para:4:6:3:item:b:01}{(b)}, \SGARefLink{sga3:iv:para:4:6:3:item:c:01}{(c)},
\SGARefLink{sga3:iv:para:4:6:3:item:d:mathcal:t:01}{\((d_{T_i})\)}, and \SGARefLink{sga3:iv:para:4:6:3:item:e:mathcal:t:01}{\((e_{T_i})\)} of \SGARefLink{sga3:iv:para:4:6:3}{4.6.3}.}
\footnote{Editor's note: the original has been corrected by adding
surjectivity to \((M_3)\) and \((M_4)\); it is automatically satisfied in
the other two cases.}

Conditions \SGARefLink{sga3:iv:para:4:6:3:item:a:01}{(a)}, \SGARefLink{sga3:iv:para:4:6:3:item:b:01}{(b)}, and \SGARefLink{sga3:iv:para:4:6:3:item:c:01}{(c)} are classical (see EGA and \SGARefLink{sga3:iv:bib:sga:1}{SGA~1},
passim).\footnote{Editor's note: see EGA~I, 6.6.4 for "quasi-compact";
EGA~II, 6.1.5 for "finite"; EGA~IV\textsubscript{1}, 1.6.2 for "locally
of finite presentation"; EGA~IV\textsubscript{2}, 2.2.13 for "faithfully
flat"; and EGA~IV\textsubscript{4}, 17.3.3 for "\'{e}tale".}
By \SGARefLink{sga3:iv:proposition:6:3:1:item:iv:01}{6.3.1(iv)},\SGARefLink{sga3:iv:proposition:6:3:1:item:v:01}{(v)}, \((M_i)\) satisfies \SGARefLink{sga3:iv:para:4:6:3:item:d:mathcal:t:01}{\((d_{T_i})\)}. It remains to verify
\SGARefLink{sga3:iv:para:4:6:3:item:e:mathcal:t:01}{\((e_{T_i})\)}. It is enough to verify \SGARefLink{sga3:iv:para:4:6:3:item:e:mathcal:t:01}{\((e_{T_1})\)}, which implies all the
others; this follows from \SGARefLink{sga3:iv:bib:sga:1}{SGA~1}, VIII, \S\S4--5.

\medskip
\noindent\SGARefTarget{sga3:iv:corollary:6:3:3}\textbf{Corollary 6.3.3.}
\textit{Let \(X\) be a scheme and let \(R\) be an equivalence relation in
\(X\) of type \((M_i)\). Then \(R\) is \((M_i)\)-effective if and only if
the sheaf quotient of \(X\) by \(R\) for \(T_i\) is representable; in that
case it is represented by the quotient \(X/R\).}

This is \SGARefLink{sga3:iv:lemma:4:6:5}{4.6.5}.

\SGARefTarget{sga3:iv:subsection:6:4}\subsection*{6.4. Effectivity Conditions}
\addcontentsline{toc}{subsection}{\SGARefLink{sga3:iv:subsection:6:4}{6.4}. Effectivity Conditions}

We now seek families \((N)\) of morphisms satisfying axiom \SGARefLink{sga3:iv:subsection:4:7:item:f:mathcal:t:01}{\((f_T)\)} of
\SGARefLink{sga3:iv:subsection:4:7}{4.7}. Observe first that \SGARefLink{sga3:iv:subsection:4:7:item:f:mathcal:t:01}{\((f_{T_1})\)} implies \SGARefLink{sga3:iv:subsection:4:7:item:f:mathcal:t:01}{\((f_{T_i})\)}, so it is
enough to consider the \((\mathrm{fpqc})\) topology.

\medskip
\noindent\SGARefTarget{sga3:iv:lemma:6:4:1}\textbf{Lemma 6.4.1.}
\textit{The following families of morphisms satisfy axiom \SGARefLink{sga3:iv:subsection:4:7:item:f:mathcal:t:01}{\((f_{T_1})\)}
of \SGARefLink{sga3:iv:subsection:4:7}{4.7}, that is, they "descend by \((\mathrm{fpqc})\)":}
\begin{align*}
 (N)&:\ \text{open immersions};\\
 (N')&:\ \text{closed immersions};\\
 (N'')&:\ \text{quasi-compact immersions}.
\end{align*}

By \SGARefLink{sga3:iv:proposition:6:3:1:item:ii:01}{6.3.1(ii)}, it is enough to verify that these families descend for the
Zariski topology and along a faithfully flat quasi-compact morphism. The
first assertion is clear. For \((N)\), the second is \SGARefLink{sga3:iv:bib:sga:1}{SGA~1}, VIII, 4.4; for
\((N')\), it is loc. cit., 1.9; for \((N'')\), argue as in loc. cit., 5.5,
using the two preceding results.

% Source page: Exp4 component-local p. 69; combined-reader p. 257.
\medskip
\noindent\SGARefTarget{sga3:iv:corollary:6:4:2}\textbf{Corollary 6.4.2.}
\textit{\SGARefLink{sga3:iv:lemma:6:4:1}{The same result} holds for quasi-compact open immersions.}

These results make the general results of \SGARefLink{sga3:iv:proposition:4:7:1}{4.7.1}, \SGARefLink{sga3:iv:proposition:4:7:2}{4.7.2}, \SGARefLink{sga3:iv:proposition:5:1:8}{5.1.8}, \SGARefLink{sga3:iv:proposition:5:3:1}{5.3.1}, and
so on applicable to the present situation. We state the first as an
example.

\medskip
\noindent\SGARefTarget{sga3:iv:corollary:6:4:3}\textbf{Corollary 6.4.3.}
\textit{\((=\SGARefLink{sga3:iv:proposition:4:7:1}{4.7.1}+\SGARefLink{sga3:iv:corollary:4:6:10}{4.6.10})\). Let \(X\) be a scheme and \(R\) an
equivalence relation in \(X\). Suppose that \(R\to X\) is faithfully flat
and quasi-compact, and that \(R\to X\times X\) is a closed immersion
(respectively, open, quasi-compact, or quasi-compact open). Then the sheaf
quotient \(X/R\) is the same for the \((\mathrm{fpqc})\) topology as for
the canonical topology. For every scheme \(S\),}
\[
(X/R)(S)=
\left\{
\begin{array}{l}
\text{closed (respectively, open, retrocompact, or open retrocompact)}
\\
\text{subschemes \(Z\) of \(X_S\), stable under \(R\times S\), such that}
\\
\text{\(Z\to S\) is faithfully flat and quasi-compact and}
\\
\text{\(R_Z\rightrightarrows Z\to S\) is exact}
\end{array}
\right\}.
\tag{6.4.3.1}\label{sga3:iv:corollary:6:4:3:eq:6:4:3:1:01}
\]
\footnote{Editor's note: recall that a subscheme \(Z\) of a scheme \(T\)
is called \emph{retrocompact} if the immersion \(Z\hookrightarrow T\) is
quasi-compact; see EGA~0\textsubscript{III}, 9.1.1.}

\SGARefTarget{sga3:iv:subsection:6:5}\subsection*{6.5. Principal Homogeneous Bundles}
\addcontentsline{toc}{subsection}{6.5. Principal Homogeneous Bundles}

We merely record the terminology:
\[
\begin{array}{c|l}
 \text{topology}&\text{principal homogeneous bundles}\\ \hline
 (\mathrm{fpqc})&\text{principal homogeneous bundles (without qualifier)}
 \\
 (\mathrm{\acute et})&\text{quasi-isotrivial principal homogeneous bundles}
 \\
 (\mathrm{\acute etf})&\text{locally isotrivial principal homogeneous bundles}
 \\
 (\mathrm{Zar})&\text{locally trivial principal homogeneous bundles}.
\end{array}
\tag{6.5.1}\label{sga3:iv:subsection:6:5:eq:6:5:1:01}
\]

\SGARefTarget{sga3:iv:subsection:6:6}\subsection*{6.6. Other Topologies}
\addcontentsline{toc}{subsection}{\SGARefLink{sga3:iv:subsection:6:6}{6.6}. Other Topologies}

Other topologies on the category of schemes are sometimes useful. One is
the \emph{global finite \'{e}tale topology}
\((\mathrm{\acute etfg})\), generated by the pretopology whose covering
families are the surjective families of finite \'{e}tale morphisms. It is
not finer than the Zariski topology. The corresponding principal
homogeneous bundles are called \emph{isotrivial}.

% Source page: Exp4 component-local p. 70; combined-reader p. 258.
The position of this topology, and of the chaotic topology, in the lattice
is as follows:
\[
\begin{array}{ccccc}
  & & (\text{canonical}) & &\\
  & & \Big\downarrow & &\\[-0.25em]
  & & (\mathrm{fpqc}) & &\\
  & & \Big\downarrow & &\\[-0.25em]
  & & (\mathrm{fppf}) & &\\
  & & \Big\downarrow & &\\[-0.25em]
  & & (\mathrm{\acute et}) & &\\
  & & \Big\downarrow & &\\[-0.25em]
  & & (\mathrm{\acute etf}) & &\\
  & \swarrow & & \searrow &\\[-0.25em]
  (\mathrm{Zar}) & & & & (\mathrm{\acute etfg})\\
  & \searrow & & \swarrow &\\[-0.25em]
  & & (\text{chaotic}) & &
\end{array}
\tag{6.6.1}\label{sga3:iv:subsection:6:6:eq:6:6:1:01}
\]
\footnote{Editor's note: recall (cf. \SGARefLink{sga3:iv:scholium:4:4:2}{4.4.2}) that the chaotic topology is
the coarsest topology, defined by \(J(S)=\{S\}\) for every
\(S\in\operatorname{Ob}\mathscr C\).}

\SGARefTarget{sga3:iv:subsection:6:7}\subsection*{6.7. Homogeneous Spaces}
\addcontentsline{toc}{subsection}{\SGARefLink{sga3:iv:subsection:6:7}{6.7}. Homogeneous Spaces}

Let\footnote{Editor's note: the numbering that follows has been added.}
\(G\) be an \(S\)-group scheme, let \(X\) be an \(S\)-scheme equipped
with a left action of \(G\), and let
\[
  \Phi:G\times_SX\longrightarrow X\times_SX
\]
be the morphism of \(S\)-schemes defined on points by
\[
  (g,x)\longmapsto(gx,x).
\]
Recall (cf. \SGARefLink{sga3:iv:definition:5:1:0}{5.1.0} and III.\SGARefLink{sga3:iii:definition:0:1}{0.1}) that \(X\) is called a
\emph{formally principal homogeneous space} under \(G\) if the following
equivalent conditions hold:
\begin{enumerate}
  \item[(i)]\SGARefTarget{sga3:iv:subsection:6:7:item:i:01} for every \(T\to S\), the set \(X(T)\) is empty or is a
    principal homogeneous space under \(G(T)\);
  \item[(ii)]\SGARefTarget{sga3:iv:subsection:6:7:item:ii:01} \(\Phi\) is an isomorphism of \(S\)-functors;
  \item[(iii)]\SGARefTarget{sga3:iv:subsection:6:7:item:iii:01} \(\Phi\) is an isomorphism of \(S\)-schemes.
\end{enumerate}
The equivalence \SGARefLink{sga3:iv:subsection:6:7:item:i:01}{(i)}\(\Longleftrightarrow\)\SGARefLink{sga3:iv:subsection:6:7:item:ii:01}{(ii)} is clear, and
\SGARefLink{sga3:iv:subsection:6:7:item:ii:01}{(ii)}\(\Longleftrightarrow\)\SGARefLink{sga3:iv:subsection:6:7:item:iii:01}{(iii)} holds because
\(\mathscr C=(\mathbf{Sch}/S)\) is a full subcategory of
\(\widehat{\mathscr C}\).

The notion of a \emph{formally homogeneous space}, not necessarily
principal homogeneous, is obtained by requiring \(\Phi\) to be an
epimorphism in the category of sheaves for an appropriate topology
\(\mathcal T\). Indeed, requiring \(\Phi\) to be an epimorphism of
\(S\)-functors is equivalent to requiring that, for every \(T\to S\), the
set \(X(T)\) be empty or homogeneous (not necessarily principal
homogeneous) under \(G(T)\). This is too restrictive, as the following
simple example shows.

Let
\[
  S=\operatorname{Spec}R,\qquad G=\mathbf G_{m,R},\qquad
  X=\mathbf G_{m,R},
\]
and let \(G\) act on \(X\) by \(t\cdot x=t^2x\). Then \(\Phi\) is
surjective, finite, and \'{e}tale, and hence is an epimorphism in the
category of sheaves for the \((\mathrm{\acute etf})\) topology (and,
\textit{a fortiori}, an epimorphism of \(S\)-schemes). On the other hand,
the points \(1\) and \(-1\) of \(X(R)\) are not conjugate under any element
of \(G(R)\). Thus
\[
  G(R)\times X(R)\longrightarrow X(R)\times X(R)
\]
is not surjective.\footnote{Editor's note: this difficulty plainly comes
from the following fact. Let \(\mathscr C'\) be a full subcategory of
\(\widehat{\mathscr C}\) containing \(\mathscr C\), for example the
category \(\widetilde{\mathscr C}_{\mathcal T}\) of sheaves on
\(\mathscr C\) for a topology \(\mathcal T\) coarser than the canonical
topology. If \(f:X\to Y\) is a morphism in \(\mathscr C\), then the
implications
\[
 \text{\(f\) epimorphic in \(\widehat{\mathscr C}\)}
 \Longrightarrow
 \text{\(f\) epimorphic in \(\mathscr C'\)}
 \Longrightarrow
 \text{\(f\) epimorphic in \(\mathscr C\)}
\]
are generally strict.}
We are therefore led to \SGARefLink{sga3:iv:definition:6:7:1}{the following definition}.

% Source page: Exp4 component-local p. 71; combined-reader p. 259.
\medskip
\noindent\SGARefTarget{sga3:iv:definition:6:7:1}\textbf{Definition 6.7.1.}
\textit{Let \(G\) be an \(S\)-group, let \(X\) be an \(S\)-scheme
equipped with an action of \(G\), and let \(\mathcal T\) be a topology on
\((\mathbf{Sch}/S)\) coarser than the canonical topology. The scheme \(X\)
is called a \emph{formally homogeneous space under \(G\)}, relative to
\(\mathcal T\), if the following equivalent conditions hold:}
\begin{enumerate}
  \item[(i)]\SGARefTarget{sga3:iv:definition:6:7:1:item:i:01} \textit{the morphism
    \(\Phi:G\times_SX\to X\times_SX\) is an epimorphism in the category of
    sheaves for \(\mathcal T\);}
  \item[(ii)]\SGARefTarget{sga3:iv:definition:6:7:1:item:ii:01} \textit{for every \(T\to S\) and every \(x,y\in X(T)\),
    there exist a \(\mathcal T\)-covering morphism \(T'\to T\) and an
    element \(g\in G(T')\) such that \(y_{T'}=g\cdot x_{T'}\).}
\end{enumerate}

\medskip
\noindent\SGARefTarget{sga3:iv:remark:6:7:2}\textbf{Remark 6.7.2.}
\textit{Condition \SGARefLink{sga3:iv:definition:6:7:1:item:i:01}{(i)} implies in particular that \(\Phi\) is a universal
effective epimorphism in \((\mathbf{Sch}/S)\), by \SGARefLink{sga3:iv:proposition:4:4:3}{4.4.3}. As is readily
seen, it follows that \(\Phi\) is surjective; cf. \SGARefLink{sga3:iv:definition:1:3}{1.3}, Editor's note (3).}

\medskip
\noindent\SGARefTarget{sga3:iv:propdef:6:7:3}\textbf{Proposition and Definition 6.7.3.}
\footnote{Editor's note: cf. \SGARefLink{sga3:iv:bib:ray70}{[Ray70]}, Definition VI.1.1.}
\textit{Let \(G\) be an \(S\)-group, let \(X\) be an \(S\)-scheme
equipped with an action of \(G\), and let \(\mathcal T\) be a topology on
\((\mathbf{Sch}/S)\) coarser than the canonical topology. The following
conditions are equivalent:}
\begin{enumerate}
  \item[(i)]\SGARefTarget{sga3:iv:propdef:6:7:3:item:i:01} \textit{\(X\) satisfies both of the following conditions:}
    \begin{enumerate}
      \item[(1)]\SGARefTarget{sga3:iv:propdef:6:7:3:item:1:01} \textit{the morphism
        \(\Phi:G\times_SX\to X\times_SX\) is covering, that is, \(X\) is a
        formally homogeneous \(G\)-space;}
      \item[(2)]\SGARefTarget{sga3:iv:propdef:6:7:3:item:2:01} \textit{the structural morphism \(X\to S\) is also
        covering, that is, locally for \(\mathcal T\) it has a section
        % The PDF has an unresolved typesetting token here,
        % ``RefIV.4.4.8bis.bis''; 4.4.8 bis, (b) is the intended reference.
        (cf. \SGARefLink{sga3:iv:remark:4:4:8:bis}{4.4.8 bis}, \SGARefLink{sga3:iv:remark:4:4:8:bis:item:b:01}{(b)}).}
    \end{enumerate}

  \item[(ii)]\SGARefTarget{sga3:iv:propdef:6:7:3:item:ii:01} \textit{Locally on \(S\) for the topology \(\mathcal T\),
    the scheme \(X\), with its \(G\)-action, is isomorphic to the
    quotient sheaf of \(G\) by a subgroup scheme \(H\), for
    \(\mathcal T\).
    In other words, there is a covering family \(\{S_i\to S\}\) such that
    every \(X\times_SS_i\) represents the sheaf quotient of
    \(G\times_SS_i\) by some subgroup scheme \(H_i\).}
\end{enumerate}
\textit{Under these conditions, \(X\) is called a
\emph{\(G\)-homogeneous space}, relative to \(\mathcal T\).}

\smallskip
\noindent\textit{Proof.}
Assume \SGARefLink{sga3:iv:propdef:6:7:3:item:ii:01}{(ii)}. Put
\[
  G_i=G\times_SS_i,\qquad X_i=X\times_SS_i.
\]
The scheme \(X_i\) has a section over \(S_i\), namely the composite of the
unit section \(\varepsilon_i:S_i\to G_i\) with the projection
\(\pi_i:G_i\to X_i=G_i/H_i\). It follows that \(X\to S\) is covering.

% Source page: Exp4 component-local p. 72; combined-reader p. 260.
Moreover, \(\pi_i\) is covering, so \(\pi_i\times\pi_i\) is covering as
well, by \SGARefLink{sga3:iv:definition:4:2:3}{4.2.3} \SGARefLink{sga3:iv:definition:4:2:3:item:c:1:01}{(C~1)},\SGARefLink{sga3:iv:definition:4:2:3:item:c:2:01}{(C~2)}. There is a commutative
diagram\footnote{Editor's note: the Polo--Gille re-edition prints
\(G_i\times_{S_i}X_i\) at the upper right. We write
\(X_i\times_{S_i}X_i\), since \(\Phi_i\), as the base change of
\(\Phi:G\times_SX\to X\times_SX\), and \(\pi_i\times\pi_i\) both have
this codomain.}
\[
\begin{array}{ccccc}
  G_i\times_{S_i}X_i & \xrightarrow{\ \Phi_i\ } &
    X_i\times_{S_i}X_i\\
  \scriptstyle\operatorname{id}\times\pi_i\,\Big\uparrow
    & & \scriptstyle\pi_i\times\pi_i\,\Big\uparrow\\[-0.25em]
  G_i\times_{S_i}G_i & \xrightarrow[\sim]{\ \Psi_i\ } &
    G_i\times_{S_i}G_i ,
\end{array}
\tag{6.7.3.1}\label{sga3:iv:propdef:6:7:3:eq:6:7:3:1:01}
\]
where \(\Phi_i\) is obtained from \(\Phi\) by the base change \(S_i\to S\),
and \(\Psi_i\) is the isomorphism defined on points by
\[
  (g,g')\longmapsto(gg',g).
\]
The composite \((\pi_i\times\pi_i)\circ\Psi_i\) is covering; hence
\(\Phi_i\) is covering by \SGARefLink{sga3:iv:definition:4:2:3}{4.2.3} \SGARefLink{sga3:iv:definition:4:2:3:item:c:3:01}{(C~3)}. Thus \(\Phi\) is locally covering
and is therefore covering, by \SGARefLink{sga3:iv:definition:4:2:3}{4.2.3} \SGARefLink{sga3:iv:definition:4:2:3:item:c:5:01}{(C~5)}. This proves
\SGARefLink{sga3:iv:propdef:6:7:3:item:ii:01}{(ii)}\(\Longrightarrow\)\SGARefLink{sga3:iv:propdef:6:7:3:item:i:01}{(i)}.

Conversely, assume \SGARefLink{sga3:iv:propdef:6:7:3:item:i:01}{(i)}, and first suppose in addition that the structural
morphism \(X\to S\) has a section \(\sigma\). By EGA~I, 5.3.13,
\(\sigma\) is an immersion. Define \(H=G\times_XS\) by \SGARefLink{sga3:iv:propdef:6:7:3:eq:6:7:3:2:01}{the following
diagram}, whose two squares are cartesian:
\[
\begin{array}{ccccc}
  H & \longrightarrow & G
    & \xrightarrow{\ \operatorname{id}_G\boxtimes\sigma\ } & G\times_SX\\
  \Big\downarrow & & \scriptstyle\pi\,\Big\downarrow
    & & \scriptstyle\Phi\,\Big\downarrow\\[-0.25em]
  S & \xrightarrow{\ \sigma\ } & X
    & \xrightarrow{\ \operatorname{id}_X\boxtimes\sigma\ } & X\times_SX .
\end{array}
\tag{6.7.3.2}\label{sga3:iv:propdef:6:7:3:eq:6:7:3:2:01}
\]
Here \(\pi\), \(\operatorname{id}_G\boxtimes\sigma\), and
\(\operatorname{id}_X\boxtimes\sigma\) are the morphisms defined on points,
for \(T\to S\), \(g\in G(T)\), and \(x\in X(T)\), by
\[
 \pi(g)=g\cdot\sigma_T,\qquad
 (\operatorname{id}_G\boxtimes\sigma)(g)=(g,\sigma_T),\qquad
 (\operatorname{id}_X\boxtimes\sigma)(x)=(x,\sigma_T).
\]
Then \(\pi\) is covering, and \(H\) is a subgroup scheme of \(G\)
representing the stabilizer \(\operatorname{Stab}_G(\sigma)\) of \(\sigma\)
(cf. I, \SGARefLink{sga3:i:definition:2:3:3}{2.3.3}). Thus, for every \(T\to S\),
\[
  H(T)=\{g\in G(T)\mid g\cdot\sigma_T=\sigma_T\}.
\]

Let \(G/H\) denote the presheaf
\[
  T\longmapsto G(T)/H(T),
\]
and let \(a(G/H)\) be its associated sheaf for \(\mathcal T\). The
preceding construction gives a commutative diagram of presheaves equipped
with a \(G\)-action:
\[
\begin{array}{ccccc}
  G & \xrightarrow{\ \pi\ } & X\\
  \Big\downarrow & \nearrow_{\scriptstyle\bar\pi} &\\[-0.25em]
  G/H & &
\end{array}
\tag{6.7.3.3}\label{sga3:iv:propdef:6:7:3:eq:6:7:3:3:01}
\]
where \(\bar\pi\) is a monomorphism. Since \(\pi\) is covering,
\(\bar\pi\) is covering as well. By \SGARefLink{sga3:iv:corollary:4:3:12}{4.3.12}, \(\bar\pi\) therefore induces
an isomorphism
\[
  a(G/H)\xrightarrow{\sim}X.
\]
We have proved: if \(X\) is a \(G\)-homogeneous space such that \(X\to S\)
has a section \(\sigma\), then \(X\) represents the sheaf quotient \(G/H\),
where \(H=G\times_XS\) is the stabilizer of \(\sigma\).

In general, by hypothesis there is a covering family \(\{S_i\to S\}\) such
that every morphism
\[
  X_i=X\times_SS_i\longrightarrow S_i
\]
has a section \(\sigma_i\). Put \(G_i=G\times_SS_i\). The morphism
\[
  \Phi_i:G_i\times_{S_i}X_i\longrightarrow X_i\times_{S_i}X_i
\]
obtained from \(\Phi\) by base change along \(S_i\to S\) is still covering.

% Source page: Exp4 component-local p. 73; combined-reader p. 261.
The preceding argument then gives
\[
  X_i\simeq G_i/H_i,
\]
where \(H_i\) is the stabilizer of \(\sigma_i\) in \(G_i\). This completes
the proof of \SGARefLink{sga3:iv:propdef:6:7:3:item:i:01}{(i)}\(\Longrightarrow\)\SGARefLink{sga3:iv:propdef:6:7:3:item:ii:01}{(ii)}.

\SGARefTarget{sga3:iv:section:bibliography}\section*{Bibliography}
\addcontentsline{toc}{section}{Bibliography}

\begin{description}
  \item[{[AS]}]\SGARefTarget{sga3:iv:bib:as} J.~Giraud, \emph{Analysis Situs}, S\'eminaire Bourbaki,
    Expos\'e 256, May 1963.

  \item[{[D]}]\SGARefTarget{sga3:iv:bib:d} J.~Giraud, \emph{M\'ethode de la descente}, M\'emoires de la
    Soci\'et\'e Math\'ematique de France, vol.~2 (1964), pp.~iii--viii and
    1--150.

  \item[{[MA]}]\SGARefTarget{sga3:iv:bib:ma} M.~Artin, \emph{Grothendieck Topologies}, mimeographed
    notes, Harvard, 1962.

  \item[{[SGA 1]}]\SGARefTarget{sga3:iv:bib:sga:1} \emph{S\'eminaire de G\'eom\'etrie Alg\'ebrique du
    Bois-Marie 1960--61: Rev\^etements \'etales et groupe fondamental},
    Lecture Notes in Mathematics 224 (1971); recomposed and annotated
    edition, Documents Math\'ematiques 3, Soci\'et\'e Math\'ematique de
    France, 2003.

  \item[{[SGA 4]}]\SGARefTarget{sga3:iv:bib:sga:4} \emph{S\'eminaire de G\'eom\'etrie Alg\'ebrique du
    Bois-Marie 1963--1964: Th\'eorie des topos et cohomologie \'etale des
    sch\'emas}, vols.~I--III, Lecture Notes in Mathematics 269, 270 (1972)
    and 305 (1973).

  \item[{[TDTE I]}]\SGARefTarget{sga3:iv:bib:tdte:i} A.~Grothendieck, \emph{Techniques de descente et
    th\'eor\`emes d'existence en g\'eom\'etrie alg\'ebrique I.
    G\'en\'eralit\'es. Descente par morphismes fid\`element plats},
    S\'eminaire Bourbaki, Expos\'e 190, December 1959.

  \item[{[Ray70]}]\SGARefTarget{sga3:iv:bib:ray70} M.~Raynaud, \emph{Faisceaux amples sur les sch\'emas en
    groupes et les espaces homog\`enes}, Lecture Notes in Mathematics 119,
    Springer-Verlag, 1970.\footnote{Editor's note: this reference has been
    added.}
\end{description}

% END INCLUDED SGA3_English_current_buildable_source_20260823/source/inputs/IV/tex/components/21_expose_IV_s6_and_bibliography.tex


% END INCLUDED SGA3_English_current_buildable_source_20260823/source/inputs/IV/tex/components/33_expose_IV_complete.tex



% END INCLUDED SGA3_English_current_buildable_source_20260823/source/inputs/IV/body.tex


\SGAStartExpose{V}
\makeatletter
\renewcommand{\SGAThreeNativeInput}[1]{%
  \filename@parse{#1}%
  \edef\SGAThreeNativeFile{%
    native_diagrams/exp5/\filename@base.\filename@ext.tex%
  }%
  \input{\SGAThreeNativeFile}%
}
\makeatother
\providecommand{\SGAThreeDiagram}[5][0.92\textwidth]{}
\renewcommand{\SGAThreeDiagram}[5][0.92\textwidth]{%
  \begin{figure}[H]
    \centering
    \resizebox{#1}{!}{\SGAThreeNativeInput{#2}}%
    \caption{#3}
    \label{#4}
    \par\smallskip
    {\footnotesize\itshape Native \LaTeX{} reconstruction from the
      corrected Polo--Gille diagram; source locator: #5.}
  \end{figure}%
}
% BEGIN INCLUDED SGA3_English_current_buildable_source_20260823/source/inputs/V/body.tex
% Mechanically extracted after \mainmatter from the frozen standalone master.
\setcounter{footnote}{0}

\IfFileExists{tex/components/00_expose_V_pp251_252_en.tex}{%
  % BEGIN INCLUDED SGA3_English_current_buildable_source_20260823/source/inputs/V/tex/components/00_expose_V_pp251_252_en.tex
% SGA 3, Expose V, Loop 1 English reconstruction.
% Scope: printed pages 251--252, plus the closing sentence of Lemma 1.1
% at the top of printed page 253; Lemma 1.2 is excluded.
% Authority: Polo--Gille Exp5-13oct24.pdf, component-local pages 1--3
% (combined-reader pages 263--265).  Page 3 controls only the first sentence.
% OCR and the Reinhold English Markdown are comparison witnesses only.

\setcounter{footnote}{0}
\chapter*{Expos\'e V\\Construction of Quotient Schemes}
\phantomsection\label{exp:sga3-V}
\addcontentsline{toc}{chapter}{Expos\'e V. Construction of Quotient Schemes}
\markboth{EXPOS\'E V. CONSTRUCTION OF QUOTIENT SCHEMES}%
  {EXPOS\'E V. CONSTRUCTION OF QUOTIENT SCHEMES}

\begin{center}
  \textit{by P. Gabriel}
\end{center}

The purpose of this expos\'e is to prove the theorems stated in TDTE~III.%
\footnote{\label{note:sga3-V-1}Editor's note: namely, \SGARefLink{rem:sga3-V-5.1}{Theorems~5.1}, 5.3,
\SGARefLink{lem:sga3-V-6.1}{6.1}, 6.2, and~\SGARefLink{lem:sga3-V-7.2}{7.2} of TDTE~III. The first two (respectively, the next two)
correspond to \hyperref[thm:sga3-V-4.1]{Theorem~4.1} (respectively, \hyperref[thm:sga3-V-7.1]{Theorems~7.1} and~\hyperref[thm:sga3-V-8.1]{8.1}) of this
expos\'e. \SGARefLink{lem:sga3-V-7.2}{Theorem~7.2} of TDTE~III is proved in Expos\'e~VI\textsubscript{A},
3.2 and~3.3.}
If \(X\) and \(T\) are two objects of a category~\(\mathscr C\), we write
\(X(T)\) instead of \(\operatorname{Hom}_{\mathscr C}(T,X)\). Similarly, if
\(\varphi:Y\to X\) is an arrow (respectively, if \(T\) is an object) of
\(\mathscr C\), then \(\varphi(T)\) denotes the map
\(g\mapsto\varphi\circ g\) from \(Y(T)\) to \(X(T)\):

\SGAThreeDiagram[0.23\textwidth]
  {figures/exp5/p251_phi_of_T_triangle.png}
  {The map \(\varphi(T):Y(T)\to X(T)\).}
  {fig:sga3-exp5-p251-phi-of-T}
  {Polo--Gille Exp5-13oct24.pdf, component-local p. 1; combined-reader p. 263; printed p. 251}

Likewise, \(T(\varphi)\) denotes the map \(g\mapsto g\circ\varphi\) from
\(T(X)\) to \(T(Y)\):

\SGAThreeDiagram[0.23\textwidth]
  {figures/exp5/p251_T_of_phi_triangle.png}
  {The map \(T(\varphi):T(X)\to T(Y)\).}
  {fig:sga3-exp5-p251-T-of-phi}
  {Polo--Gille Exp5-13oct24.pdf, component-local p. 1; combined-reader p. 263; printed p. 251}

Finally, if \(P\) is a scheme, we denote by \(\underline P\) the underlying
set of~\(P\).

As an exception, in the present expos\'e we do not follow the convention on
quotient notation stated in \SGARefLink{sga3:iv:scholium:4:6:15}{IV~4.6.15} (\emph{loc. cit.}, top of p.~227 of
the original), because we wish to give here a construction of quotients that
also applies to ``pre-equivalence relations''%
\footnote{\label{note:sga3-V-2}Editor's note: that is, groupoids with base
\(X\); see the terminology at the end of \hyperref[sec:sga3-V-1]{Section~1}. When \(\mathscr C\) is
the category of schemes, the quotient \(p:X\to Y\) of a groupoid \(X_\ast\)
with base \(X\) exists under certain hypotheses (see \hyperref[thm:sga3-V-4.1]{4.1}, \hyperref[lem:sga3-V-6.1]{6.1}, and~\hyperref[thm:sga3-V-7.1]{7.1}). If,
in addition, \(X_\ast\) is an equivalence relation, then under the same
hypotheses \(p\) is faithfully flat and quasi-compact, hence a universal
epimorphism; see \emph{loc. cit.}.}
that are not equivalence relations.

\section*{1. \(\mathscr C\)-groupoids}
\phantomsection\label{sec:sga3-V-1}
\addcontentsline{toc}{section}{1. \texorpdfstring{\(\mathscr C\)}{C}-groupoids}

\phantomsection\label{par:sga3-V-1-a}%
\noindent\textbf{a)} Let \(\mathscr C\) be a category in which products and
fiber products exist. Recall first that a diagram

\SGAThreeDiagram[0.39\textwidth]
  {figures/exp5/p252_exact_pair_cokernel.png}
  {A parallel pair followed by its proposed cokernel.}
  {fig:sga3-exp5-p252-exact-pair-cokernel}
  {Polo--Gille Exp5-13oct24.pdf, component-local p. 2; combined-reader p. 264; printed p. 252}

\noindent in \(\mathscr C\) is said to be \emph{exact} if
\(p d_0=p d_1\) and if, for every \(T\in\mathscr C\), the map \(T(p)\) is a
bijection from \(T(Y)\) onto the subset of \(T(X_0)\) consisting of the
arrows \(f:X_0\to T\) such that \(f d_0=f d_1\). One also says that
\((Y,p)\) is the \emph{cokernel} of \((d_0,d_1)\), and writes
\phantomsection\label{eq:sga3-V-1-cokernel}%
\begin{equation*}
  (Y,p)=\operatorname{Coker}(d_0,d_1).
\end{equation*}

\phantomsection\label{par:sga3-V-1-b}%
\noindent\textbf{b)} For example, let \(\mathscr C\) be the category
\((\mathbf{Esp.An})\) of ringed spaces. In this case a cokernel \((Y,p)\)
always exists, and it may be described as follows. The underlying topological
space of \(Y\) is obtained from \(X_0\) by identifying the points \(d_0(x)\)
and \(d_1(x)\), and by giving \(Y\) the quotient topology. The canonical map
\(\pi:X_0\to Y\), together with \(d_0,d_1\), then induces a pair of arrows
of sheaves of rings on \(Y\):

\SGAThreeDiagram[0.66\textwidth]
  {figures/exp5/p252_ringed_space_sheaf_pair.png}
  {The pair of induced arrows between direct-image sheaves.}
  {fig:sga3-exp5-p252-ringed-space-sheaf-pair}
  {Polo--Gille Exp5-13oct24.pdf, component-local p. 2; combined-reader p. 264; printed p. 252}

Here \(\mathcal O_i\) is the structure sheaf of \(X_i\). As the sheaf of
rings on \(Y\), one chooses the subsheaf of \(\pi_*(\mathcal O_0)\) whose
sections \(s\) satisfy \(\delta_0(s)=\delta_1(s)\). The arrow \(p\) is
defined in the evident way.

Consider%
\footnote{\label{note:sga3-V-3}Editor's note: \hyperref[lem:sga3-V-1.1]{Lemmas~1.1} and~\hyperref[sga3:exp5:lemma:1.2]{1.2} have been
added; they are used several times in Sections~\hyperref[sec:sga3-V-5]{5}--\hyperref[sec:sga3-V-9]{9}.}
the following diagram in \((\mathbf{Esp.An})\):

\SGAThreeDiagram[0.18\textwidth]
  {figures/exp5/p252_ringed_space_pair_inline.png}
  {The pair \((d_0,d_1):X_1\rightrightarrows X_0\).}
  {fig:sga3-exp5-p252-ringed-space-pair}
  {Polo--Gille Exp5-13oct24.pdf, component-local p. 2; combined-reader p. 264; printed p. 252}

\noindent Let \((Y,p)\) be its cokernel. An open subset \(U\) of \(X_0\)
is said to be \emph{saturated} if
\[
  d_0^{-1}(U)=d_1^{-1}(U),
\]
which is equivalent to saying \(U=p^{-1}(p(U))\). In this case, since \(Y\)
has the quotient topology, \(p(U)\) is an open subset of~\(Y\).

\medskip
\phantomsection\label{lem:sga3-V-1.1}%
\noindent\textbf{Lemma 1.1.}
\textit{Let \(U\) be a saturated open subset of \(X\), and set
\(V=p(U)\). Denote by
\(U_1=d_0^{-1}(U)=d_1^{-1}(U)\) the corresponding open subset of \(X_1\),
and by \(\widetilde d_0,\widetilde d_1\), and \(\widetilde p\) the
restrictions of \(d_0,d_1\) to \(U_1\), and of \(p\) to \(U\), respectively.
Then \((V,\widetilde p)\) is a cokernel in \((\mathbf{Esp.An})\) of:}%
\footnote{\label{note:sga3-V-4}Editor's note: this is not the case in the
category of schemes. For example, let \(S=\operatorname{Spec}(\mathbf C)\),
let
\(X_0=\mathbf A^2_S=\operatorname{Spec}(\mathbf C[x_1,x_2])\), let
\(d_1:\mathbf G_{m,S}\times_S\mathbf A^2_S\to\mathbf A^2_S\) be the
action of \(\mathbf G_{m,S}\) on \(\mathbf A^2_S\) by homotheties, let
\(d_0\) be the projection onto the second factor, and let
\(U=\mathbf A^2_S\setminus\{m\}\), where \(m=(0,0)\). Then the projective
line \(\mathbf P^1_S\) is the cokernel of
\((\widetilde d_0,\widetilde d_1)\) in both \((\mathbf{Esp.An})\) and
\((\mathbf{Sch})\), whereas the cokernel \(Y\) of \((d_0,d_1)\) in
\((\mathbf{Esp.An})\) is the union of \(\mathbf P^1_S\) and the point
\(y_0=\{p(m)\}\). The only open subset containing \(y_0\) is \(Y\), and
\(\Gamma(Y,\mathcal O_Y)=\mathbf C\). If
\(f:\mathbf A^2_S\to T\) is a morphism of \(S\)-schemes such that
\(f d_0=f d_1\), and if \(V=\operatorname{Spec}(A)\) is an affine open
subset of \(T\) containing \(t_0=f(y_0)\), then
\(f^{-1}(V)=\mathbf A^2\), and the ring homomorphism
\(A\to\mathbf C[x_1,x_2]\) factors through \(\mathbf C\). This shows that
\(S=\operatorname{Spec}(\mathbf C)\) is the cokernel of \((d_0,d_1)\) in
the category \((\mathbf{Sch}/S)\).}

\SGAThreeDiagram[0.40\textwidth]
  {figures/exp5/p252_lemma_1_1_restricted_cokernel.png}
  {The restricted pair and its cokernel in \hyperref[lem:sga3-V-1.1]{Lemma~1.1}.}
  {fig:sga3-exp5-p252-lemma-1-1-restricted-cokernel}
  {Polo--Gille Exp5-13oct24.pdf, component-local p. 2; combined-reader p. 264; printed p. 252}

The verification is easy.

% END INCLUDED SGA3_English_current_buildable_source_20260823/source/inputs/V/tex/components/00_expose_V_pp251_252_en.tex
%
}{%
  \chapter*{Expos\'e V opening (pending)}%
}

\IfFileExists{tex/components/01_expose_V_pp253_255_en.tex}{%
  % BEGIN INCLUDED SGA3_English_current_buildable_source_20260823/source/inputs/V/tex/components/01_expose_V_pp253_255_en.tex
% SGA 3, Expose V, Loop 1 English reconstruction.
% Scope: printed pp. 253--255, from Lemma 1.2 through the terminology at
% the end of section 1; section 2 is excluded.
% Authority: Polo--Gille Exp5-13oct24.pdf, component-local pp. 3--5
% (combined-reader pp. 265--267). The PDF controls text, formulas,
% numbering, notes, and the temporary Loop-1 diagram images.
% Jacob Reinhold's CC BY 4.0 Markdown is comparison/drafting material only.

% Source page: Exp5 component-local p. 3; combined-reader p. 265;
% printed p. 253.
\medskip
\phantomsection\label{sga3:exp5:lemma:1.2}
\noindent\textbf{Lemma 1.2.}---
\textit{Let \(d_0,d_1:X_1\rightrightarrows X_0\) be a diagram in
\((\mathbf{Sch})\), and let \((Y,p)\) be its cokernel in
\((\mathbf{Esp.An})\).
\begin{enumerate}
\item[(i)]\phantomsection\label{lem:sga3-V-1.2-i}%
If \(Y\) is a scheme and \(p\) is a morphism of schemes,
  then \((Y,p)\) is a cokernel of \((d_0,d_1)\) in
  \((\mathbf{Sch})\).
\item[(ii)]\phantomsection\label{lem:sga3-V-1.2-ii}%
Suppose that every point of \(X_0\) has a saturated open
  neighborhood \(U\) with the following property. Let
  \(\widetilde d_0\) and \(\widetilde d_1\) be the restrictions of
  \(d_0\) and \(d_1\) to
  \(d_0^{-1}(U)=d_1^{-1}(U)\), and let \((Q,q)\) be the cokernel of
  \((\widetilde d_0,\widetilde d_1)\) in \((\mathbf{Esp.An})\). If
  \(Q\) is a scheme and \(q\) is a morphism of schemes, then
  \((Y,p)\) is a cokernel of \((d_0,d_1)\) in
  \((\mathbf{Sch})\).
\end{enumerate}}

\hyperref[lem:sga3-V-1.2-i]{Part~(i)} is proved in \hyperref[proof:sga3-V-4.1-c]{\S4(c)}; since the proof is short, let us repeat
it here. Let \(f:X_0\to Z\) be a morphism of schemes such that
\(fd_0=fd_1\). By hypothesis there is a unique morphism of ringed
spaces \(r:Y\to Z\) such that \(f=rp\). It remains to show that, for
every \(y\in Y\), the homomorphism
\(\mathscr O_{r(y)}\to\mathscr O_y\) induced by \(r\) is local. This
follows because \(p\) is surjective, so \(y=p(x)\) for some \(x\), and
because the homomorphism
\(\mathscr O_{f(x)}\to\mathscr O_x\) induced by \(f\) is local.

\hyperref[lem:sga3-V-1.2-ii]{Part~(ii)} follows from \hyperref[lem:sga3-V-1.2-i]{part~(i)} and \hyperref[lem:sga3-V-1.1]{the preceding lemma}.

\medskip
\phantomsection\label{par:sga3-V-1-c}%
\noindent\textbf{c)} In this expos\'e we study the existence of
\(\operatorname{Coker}(d_0,d_1)\) when the double arrow
\((d_0,d_1)\) is placed in a richer context. More precisely, let
\[
  X_2=X_1\times_{d_1,d_0}X_1
\]
be the fiber product of the diagram
% Native Loop-2 reconstruction; printed p. 253.
\begin{equation}\SGAInlineEquationTarget{sga3:v:lemma:1-2:equation:star}
  \resizebox{0.20\textwidth}{!}{%
    \SGAThreeNativeInput{figures/exp5/p253_star_fiber_product.png}}
  \tag{*}\label{fig:sga3-exp5-star-fiber-product}
\end{equation}
and let \(d'_0,d'_2:X_2\to X_1\) be the two canonical projections.
Thus, by definition, the square
% Native Loop-2 reconstruction; printed p. 253.
\begin{equation}\SGAInlineEquationTarget{sga3:v:lemma:1-2:equation:0}
  \resizebox{0.21\textwidth}{!}{%
    \SGAThreeNativeInput{figures/exp5/p253_square_0.png}}
  \tag{0}\label{fig:sga3-exp5-square-0}
\end{equation}
is Cartesian.

In addition, suppose that a third arrow \(d'_1:X_2\to X_1\) is
given. We say that
\[
  (d_0,d_1:X_1\rightrightarrows X_0,d'_1)
\]
is a \(\mathscr C\)-\emph{groupoid} if, for every object \(T\) of
\(\mathscr C\), \(X_1(T)\) is the set of arrows of a groupoid
\(X_\ast(T)\) whose set of objects is \(X_0(T)\), whose source map is
\(d_1(T)\), whose target map is \(d_0(T)\), and whose composition map
is \(d'_1(T)\). As usual, we identify
\[
 (X_1\times_{d_1,d_0}X_1)(T)
 \simeq X_1(T)\times_{d_1(T),d_0(T)}X_1(T),
\]
and recall that a groupoid is a category in which every arrow is
invertible.\footnote{\label{note:sga3-V-5}Editor's note: in this case \(X_2(T)\) is the set
of pairs \((f_2,f_1)\) of composable arrows, that is, pairs such that
\(d_0(f_1)=d_1(f_2)\); the maps \(d'_0,d'_1,d'_2\) send
\((f_2,f_1)\), respectively, to \(f_2,f_2\circ f_1,f_1\).}

% Source page: Exp5 component-local p. 4; combined-reader p. 266;
% printed p. 254.
If \(\varphi\) is an arrow of the groupoid \(X_\ast(T)\), the map
\(f\mapsto\varphi\circ f\) is a bijection from the set of arrows \(f\)
whose target equals the source of \(\varphi\) onto the set of arrows
having the same target as \(\varphi\). This is equivalently expressed
by saying that the square
% Native Loop-2 reconstruction; printed p. 254.
\begin{equation}\SGAInlineEquationTarget{sga3:v:lemma:1-2:equation:1}
  \resizebox{0.19\textwidth}{!}{%
    \SGAThreeNativeInput{figures/exp5/p254_square_1.png}}
  \tag{1}\label{fig:sga3-exp5-square-1}
\end{equation}
is Cartesian.

Similarly, \(g\mapsto g\circ\varphi\) is a bijection from the set of
arrows \(g\) of \(X_\ast(T)\) whose source equals the target of
\(\varphi\) onto the set of arrows having the same source as
\(\varphi\). Equivalently, the square
% Native Loop-2 reconstruction; printed p. 254.
\begin{equation}\SGAInlineEquationTarget{sga3:v:lemma:1-2:equation:2}
  \resizebox{0.19\textwidth}{!}{%
    \SGAThreeNativeInput{figures/exp5/p254_square_2.png}}
  \tag{2}\label{fig:sga3-exp5-square-2}
\end{equation}
is Cartesian.

On the other hand, let \(s:X_0\to X_1\) be the unique arrow of
\(\mathscr C\) such that, for every \(T\), the map
\(s(T):X_0(T)\to X_1(T)\) sends each object of \(X_\ast(T)\) to its
identity arrow.\footnote{\label{note:sga3-V-6}Editor's note: \(T\mapsto s(T)\) defines an
element of \(\operatorname{Hom}(h_{X_0},h_{X_1})\), and the latter is
equal to \(\operatorname{Hom}(X_0,X_1)\) by the Yoneda lemma.}
The arrow \(s\) satisfies
\begin{align}
  d_1s&=\operatorname{id}_{X_0},
    \tag{3}\label{eq:sga3-exp5-3}\\
  d_0s&=\operatorname{id}_{X_0}.
    \tag{3 bis}\label{eq:sga3-exp5-3bis}
\end{align}
Finally, associativity of the composition maps \(d'_1(T)\) is
equivalent to commutativity of the diagram
% Native Loop-2 reconstruction; printed p. 254.
\begin{equation}\SGAInlineEquationTarget{sga3:v:lemma:1-2:equation:4}
  \resizebox{0.72\textwidth}{!}{%
    \SGAThreeNativeInput{figures/exp5/p254_associativity_4.png}}
  \tag{4}\label{fig:sga3-exp5-associativity-4}
\end{equation}

% Source page: Exp5 component-local p. 5; combined-reader p. 267;
% printed p. 255.
Conversely, conditions
\hyperref[sga3:v:lemma:1-2:equation:1]{(1)},
\hyperref[sga3:v:lemma:1-2:equation:2]{(2)}, and
\hyperref[sga3:v:lemma:1-2:equation:4]{(4)}, together with the
existence of an arrow \(s\) satisfying
\hyperref[eq:sga3-exp5-3]{(3)}, imply that
\((X_1\rightrightarrows X_0,d'_1)\) is a
\(\mathscr C\)-groupoid. Condition
\hyperref[eq:sga3-exp5-3]{(3)} is harmless: it merely ensures that
\(d_1(T):X_1(T)\to X_0(T)\) is surjective for every
\(T\in\mathscr C\). In the sequel we shall mainly use the Cartesian
squares
\hyperref[sga3:v:lemma:1-2:equation:0]{(0)},
\hyperref[sga3:v:lemma:1-2:equation:1]{(1)}, and
\hyperref[sga3:v:lemma:1-2:equation:2]{(2)}, summarized in the diagram
% Native Loop-2 reconstruction; printed p. 255.
\begin{equation}\SGAInlineEquationTarget{sga3:v:lemma:1-2:equation:0-1-2}
  \resizebox{0.34\textwidth}{!}{%
    \SGAThreeNativeInput{figures/exp5/p255_diagram_012.png}}
  \tag{0,1,2}\label{fig:sga3-exp5-diagram-012}
\end{equation}
The two left-hand squares in this diagram (that is, squares
\hyperref[sga3:v:lemma:1-2:equation:0]{(0)} and
\hyperref[sga3:v:lemma:1-2:equation:2]{(2)}) are Cartesian; the first row
is exact; and \(X_2\) is identified with the fiber product
\(X_1\times_{d_0,d_0}X_1\).

We use associativity only indirectly, for example to ensure the
existence of an arrow \(s\) satisfying
\hyperref[eq:sga3-exp5-3]{(3)} and
\hyperref[eq:sga3-exp5-3bis]{(3 bis)}, or to ensure the existence of
an arrow
\begin{equation}
  \sigma:X_1\longrightarrow X_1,
  \qquad d_0\sigma=d_1,\quad d_1\sigma=d_0.
  \tag{\ensuremath{\dagger}}\label{eq:sga3-exp5-dagger}
\end{equation}
Here \(\sigma\) is chosen so that
\(\sigma(T):X_1(T)\to X_1(T)\) sends every arrow of \(X_\ast(T)\) to
its inverse.\footnote{\label{note:sga3-V-7}Editor's note: the Yoneda lemma shows that
\(\sigma\) is an involutive automorphism of \(X_1\); this will be used,
for example, in \hyperref[par:sga3-V-3-e]{3(e)} and in \hyperref[thm:sga3-V-4.1]{Theorem~4.1}.}

By an abuse of language, we shall sometimes call a diagram
% Native Loop-2 reconstruction; printed p. 255.
\begin{equation*}
  \resizebox{0.40\textwidth}{!}{%
    \SGAThreeNativeInput{figures/exp5/p255_groupoid_triple_arrow.png}}
  \phantomsection\label{fig:sga3-exp5-groupoid-triple-arrow}
\end{equation*}
a \(\mathscr C\)-groupoid when
\hyperref[sga3:v:lemma:1-2:equation:0]{(0)},
\hyperref[sga3:v:lemma:1-2:equation:1]{(1)}, and
\hyperref[sga3:v:lemma:1-2:equation:2]{(2)} are Cartesian,
\hyperref[sga3:v:lemma:1-2:equation:4]{(4)} is commutative, and
there exists an \(s\) satisfying
\hyperref[eq:sga3-exp5-3]{(3)}. Thus \(X_2\) may be ``a'' fiber
product of \hyperref[sga3:v:lemma:1-2:equation:star]{(*)} without
being ``the'' fiber product of
\hyperref[sga3:v:lemma:1-2:equation:star]{(*)}.\footnote{\label{note:sga3-V-8}Editor's
note: see \hyperref[par:sga3-V-2-a]{Example~2(a)} below.}

\medskip
\phantomsection\label{sga3:exp5:terminology:c-groupoid}
\noindent\textit{Terminology.} In place of the
\(\mathscr C\)-groupoid \(X_\ast\), we shall also speak of the
\emph{groupoid \(X_\ast\) with base \(X_0\)}, or of the
\emph{pre-equivalence relation \(X_\ast\) in \(X_0\)}.

\section*{2. Examples of \(\mathscr C\)-groupoids}
\phantomsection\label{sec:sga3-V-2}
\addcontentsline{toc}{section}{2. Examples of
  \texorpdfstring{\(\mathscr C\)}{C}-groupoids}

\phantomsection\label{par:sga3-V-2-a}%
\noindent\textbf{a)} Let \(X\) be an object of \(\mathscr C\), and let
\(G\) be a \(\mathscr C\)-group acting on the left on \(X\). We denote
by \(d_0:G\times X\to X\) the arrow defining the action of \(G\) on
\(X\), by \(d_1:G\times X\to X\) the projection onto the second
factor, and by \(\mu:G\times G\to G\) the arrow defining

% END INCLUDED SGA3_English_current_buildable_source_20260823/source/inputs/V/tex/components/01_expose_V_pp253_255_en.tex
%
}{%
  \chapter*{Expos\'e V, printed pages 253--255 (pending)}%
}

\IfFileExists{tex/components/02_expose_V_pp256_260_en.tex}{%
  % BEGIN INCLUDED SGA3_English_current_buildable_source_20260823/source/inputs/V/tex/components/02_expose_V_pp256_260_en.tex
% SGA 3, Expose V, Loop 1 English reconstruction.
% Scope: component-local PDF pp. 6--10, whose page headers are printed
% pp. 256--260. This begins with the continuation of Section 2(a) and ends
% at the end of Section 3(f), including editor's notes 9--15. The legacy
% marginal page markers advance through 262 inside this five-page range.
% Component-local p. 11 (Section 4 / Theorem 4.1) is excluded.
% Authority: Polo--Gille Exp5-13oct24.pdf, combined-reader pp. 268--272.
% OCR and Jacob Reinhold's English Markdown are comparison witnesses only.

% Source page: Exp5 component-local p. 6; combined-reader p. 268;
% page header p. 256.
the \(\mathscr C\)-group structure on \(G\), and finally by
\(\operatorname{pr}_{2,3}\) the projection of
\(G\times G\times X=G\times(G\times X)\) onto the second factor. Then

\SGAThreeDiagram[0.47\textwidth]
  {figures/exp5/p256_action_groupoid_continuation.png}
  {The groupoid attached to the action of the \(\mathscr C\)-group \(G\) on \(X\).}
  {fig:sga3-exp5-p256-action-groupoid}
  {Polo--Gille Exp5-13oct24.pdf, component-local p. 6; combined-reader p. 268; page header p. 256}

\noindent is a \(\mathscr C\)-groupoid.

\medskip
\phantomsection\label{par:sga3-V-2-b}%
\noindent\textbf{b)} Let \(d_0,d_1:X_1\to X_0\) be an
\emph{equivalence pair}. In other words, if
\(d_0\boxtimes d_1:X_1\to X_0\times X_0\) is the arrow with components
\(d_0\) and \(d_1\), we suppose that, for every object \(T\) of
\(\mathscr C\), the map \((d_0\boxtimes d_1)(T)\) is a bijection from
\(X_1(T)\) onto the graph of an equivalence relation on \(X_0(T)\).
Consequently, \(X_1(T)\) is identified with the set of pairs \((x,y)\) of
elements of \(X_0(T)\) such that \(x\sim y\). Likewise,
\[
  X_2(T)=\bigl(X_1\times_{d_1,d_0}X_1\bigr)(T)
\]
is identified with the set of triples \((x,y,z)\) of elements of
\(X_0(T)\) such that \(x\sim y\) and \(y\sim z\). There is therefore a
unique arrow \(d'_1:X_2\to X_1\) making squares
\hyperref[sga3:v:lemma:1-2:equation:1]{(1)} and
\hyperref[sga3:v:lemma:1-2:equation:2]{(2)} commute: \(d'_1(T)\) must send
\((x,y,z)\in X_2(T)\) to \((x,z)\in X_1(T)\). With this choice of
\(d'_1\),
\[
  \bigl(d_0,d_1:X_1\rightrightarrows X_0,d'_1\bigr)
\]
is a \(\mathscr C\)-groupoid.

Conversely, consider a \(\mathscr C\)-groupoid \(X_\ast\) for which
\(d_0\boxtimes d_1:X_1\to X_0\times X_0\) is a monomorphism. Then
\((d_0,d_1)\) is an equivalence pair, and \(X_\ast\) can be reconstructed
from \((d_0,d_1)\) as explained \hyperref[par:sga3-V-2-b]{a few lines above}.%
\footnote{\label{note:sga3-V-9}Editor's note: in particular, if \(G\) is
a \(\mathscr C\)-group acting on the left on an object \(X\) of
\(\mathscr C\), and if \(X_\ast\) is the \(\mathscr C\)-groupoid defined
in \hyperref[par:sga3-V-2-a]{2(a)}, then \((d_0,d_1)\) is an equivalence pair if and only if \(G\)
acts freely on \(X\); cf. Expos\'e III, 3.2.1.}

\medskip
\phantomsection\label{par:sga3-V-2-c}%
\noindent\textbf{c)} If \(p:X\to Y\) is any arrow of \(\mathscr C\), and
if \(\operatorname{pr}_1\) and \(\operatorname{pr}_2\) are the two
projections of \(X\times_{p,p}X\) onto \(X\), then
\[
  (\operatorname{pr}_1,\operatorname{pr}_2):
  X\times_{p,p}X\rightrightarrows X
\]
is an equivalence pair. The arrow \(p\) is called an
\emph{effective epimorphism} if the diagram

\SGAThreeDiagram[0.36\textwidth]
  {figures/exp5/p256_effective_epimorphism.png}
  {The defining diagram for an effective epimorphism.}
  {fig:sga3-exp5-p256-effective-epimorphism}
  {Polo--Gille Exp5-13oct24.pdf, component-local p. 6; combined-reader p. 268; page header p. 256}

\noindent is exact, that is, if
\[
  (Y,p)=\operatorname{Coker}(\operatorname{pr}_1,\operatorname{pr}_2).
\]

For example, let \(S\) be a Noetherian scheme and let \(\mathscr C\) be
the category of schemes finite over \(S\). We show that an epimorphism in
\(\mathscr C\) need not be effective. Take
\[
  S=\operatorname{Spec}k[T^3,T^5],\qquad
  Y=S,\qquad X=\operatorname{Spec}k[T],
\]
where \(k\) is a commutative field. If \(i\) is the inclusion of
\(B=k[T^3,T^5]\) into \(A=k[T]\), take \(p=\operatorname{Spec}i\). Then
\(X\times_{p,p}X\) is identified with \(\operatorname{Spec}(A\otimes_B A)\),
and \(\operatorname{Coker}(\operatorname{pr}_1,\operatorname{pr}_2)\)
with \(\operatorname{Spec}B'\), where \(B'\) is the subring of \(A\)
consisting of the elements \(a\) such that
\(a\otimes_B1=1\otimes_Ba\). But
\[
  T^7\otimes_B1
  =(T^2T^5)\otimes_B1
  =T^2\otimes_BT^5
  =T^2\otimes_B(T^3T^2)
  =T^5\otimes_BT^2
  =1\otimes_BT^7.
\]
Thus \(T^7\) belongs to \(B'\), does not belong to \(B\), and
\(\operatorname{Spec}B'\) is different from \(\operatorname{Spec}B\),
which gives the counterexample.%
\footnote{\label{note:sga3-V-10}Editor's note: the same argument applies
to \(B=k[T^3,T^4]\) and \(T^5\otimes_B1\); more generally, it applies
to \(B=k[T^n,T^{n+r}]\) and the element \(T^{n+2r}\otimes_B1\), provided
that \(n\) does not divide \(2r\).}

% Source page: Exp5 component-local p. 7; combined-reader p. 269;
% page header p. 257.
\section*{3. Some sorites concerning \(\mathscr C\)-groupoids}
\phantomsection\label{sec:sga3-V-3}
\addcontentsline{toc}{section}{3. Some sorites concerning
  \texorpdfstring{\(\mathscr C\)}{C}-groupoids}

Here are, in no particular order, some remarks used below.

\medskip
\phantomsection\label{par:sga3-V-3-a}%
\noindent\textbf{a)} Let

\SGAThreeDiagram[0.31\textwidth]
  {figures/exp5/p257_x_groupoid.png}
  {The \(\mathscr C\)-groupoid \(X_\ast\).}
  {fig:sga3-exp5-p257-x-groupoid}
  {Polo--Gille Exp5-13oct24.pdf, component-local p. 7; combined-reader p. 269; page header p. 257}

\noindent be a \(\mathscr C\)-groupoid, and let \(f_0:Y_0\to X_0\) be
an arrow of \(\mathscr C\). We shall define a \(\mathscr C\)-groupoid
with base \(Y_0\),

\SGAThreeDiagram[0.31\textwidth]
  {figures/exp5/p257_y_induced_groupoid.png}
  {The induced \(\mathscr C\)-groupoid \(Y_\ast\).}
  {fig:sga3-exp5-p257-y-induced-groupoid}
  {Polo--Gille Exp5-13oct24.pdf, component-local p. 7; combined-reader p. 269; page header p. 257}

\noindent which will be said to be \emph{induced by \(X_\ast\) and
\(f_0\)}. We also say that \(Y_\ast\) is the \emph{inverse image of
\(X_\ast\) under the base change \(f_0\)}.

We choose \(Y_1\) to be the fiber product of the diagram

\SGAThreeDiagram[0.29\textwidth]
  {figures/exp5/p257_base_change_fiber_square.png}
  {The fiber-product definition of \(Y_1\).}
  {fig:sga3-exp5-p257-base-change-fiber-square}
  {Polo--Gille Exp5-13oct24.pdf, component-local p. 7; combined-reader p. 269; page header p. 257}

\noindent and take \(e_0\) and \(e_1\) to be the composites of the
canonical arrow \(Y_1\to Y_0\times Y_0\) with the first and second
projections of \(Y_0\times Y_0\), respectively. The morphism
\(Y_1\to Y_0\times Y_0\) is then \(e_0\boxtimes e_1\), and
\[
  f_0\circ e_i=d_i\circ f_1\qquad(i=0,1),
\]
where \(f_1\) denotes the projection of \(Y_1\) onto \(X_1\).

Set \(Y_2=Y_1\times_{e_0,e_1}Y_1\), cf.
\hyperref[par:sga3-V-1-c]{1(c)}. One may say that the pair \((e_0,e_1)\)
is defined so that, for every \(T\in\mathscr C\) and every pair \((y,x)\)
of elements of \(Y_0(T)\), there is a certain bijective correspondence
\[
  \psi\longmapsto{}_{y}\psi_{x}
\]
between the arrows \(\psi\) of \(X_\ast(T)\) with source \(f_0(x)\) and
target \(f_0(y)\), and the arrows \({}_{y}\psi_x\) of \(Y_\ast(T)\)
with source \(x\) and target \(y\). Thus \(e'_1:Y_2\to Y_1\) is
determined by defining, for every \(T\in\mathscr C\), the composition
of arrows in \(Y_\ast(T)\) by the formula
\[
  {}_{z}\psi_y\circ{}_{y}\varphi_x
  ={}_{z}(\psi\circ\varphi)_x.
\]
It is clear that this definition makes each \(Y_\ast(T)\) a groupoid.

\medskip
\phantomsection\label{par:sga3-V-3-b}%
\noindent\textbf{b)} Given the \(\mathscr C\)-groupoid \(X_\ast\) and
the base change \(f_0:Y_0\to X_0\), one can reconstruct the pair
\((e_0,e_1):Y_1\rightrightarrows Y_0\) in another way.%
\footnote{\label{note:sga3-V-11}Editor's note: this second viewpoint is
used in \hyperref[par:sga3-V-3-f]{3(f)} and in the proof of \hyperref[lem:sga3-V-6.1]{6.1}.}
Construct \(Y_0\times_{X_0}X_1\), \(\operatorname{pr}_1\), and
\(\operatorname{pr}_2\) so that the square

\SGAThreeDiagram[0.27\textwidth]
  {figures/exp5/p257_reconstruction_square.png}
  {The Cartesian square defining \(Y_0\times_{X_0}X_1\).}
  {fig:sga3-exp5-p257-reconstruction-square}
  {Polo--Gille Exp5-13oct24.pdf, component-local p. 7; combined-reader p. 269; page header p. 257}

% Source page: Exp5 component-local p. 8; combined-reader p. 270;
% page header p. 258.
\noindent is Cartesian. A straightforward reduction to the set-theoretic
case then shows that the following square is Cartesian:

\SGAThreeDiagram[0.27\textwidth]
  {figures/exp5/p258_cartesian_reconstruction_square.png}
  {The second Cartesian square in the reconstruction of \((e_0,e_1)\).}
  {fig:sga3-exp5-p258-cartesian-reconstruction-square}
  {Polo--Gille Exp5-13oct24.pdf, component-local p. 8; combined-reader p. 270; page header p. 258}

Here \(f_1\) denotes the canonical projection of
\[
  Y_1=(Y_0\times Y_0)\times_{(X_0\times X_0)}X_1
\]
onto \(X_1\).

\medskip
\phantomsection\label{par:sga3-V-3-c}%
\noindent\textbf{c)} We give two examples of inverse images of a
\(\mathscr C\)-groupoid. Take \(Y_0=X_1\) and \(f_0=d_0\). For every
object \(T\) of \(\mathscr C\), the set \(Y_1(T)\) is then identified
with the set of diagrams of the form

\SGAThreeDiagram[0.18\textwidth]
  {figures/exp5/p258_inverse_image_phi_diagram.png}
  {An arrow of the first inverse-image groupoid.}
  {fig:sga3-exp5-p258-inverse-image-phi}
  {Polo--Gille Exp5-13oct24.pdf, component-local p. 8; combined-reader p. 270; page header p. 258}

\noindent in \(X_\ast(T)\). The source of such a diagram is the arrow
\(f\), and its target is the arrow \(g\). These diagrams compose in the
evident way.

Now set \(Y'_0=X_1\) and \(f'_0=d_1\) (we add the primes% 
\footnote{\label{note:sga3-V-12}Editor's note: ``accents'' in the
original.}
to avoid any confusion with the preceding example). In this case, for
every \(T\in\mathscr C\), the set \(Y'_1(T)\) is identified with the set
of diagrams of the form

\SGAThreeDiagram[0.18\textwidth]
  {figures/exp5/p258_inverse_image_psi_diagram.png}
  {An arrow of the second inverse-image groupoid.}
  {fig:sga3-exp5-p258-inverse-image-psi}
  {Polo--Gille Exp5-13oct24.pdf, component-local p. 8; combined-reader p. 270; page header p. 258}

\noindent in the groupoid \(X_\ast(T)\). The source of such a diagram is
\(f\), its target is \(g\), and these diagrams again compose in the
evident way.

It is now clear that the identity map of \(Y_0(T)\), together with the
map

\SGAThreeDiagram[0.44\textwidth]
  {figures/exp5/p258_groupoid_isomorphism_map.png}
  {The map from \(Y_1(T)\) to \(Y'_1(T)\).}
  {fig:sga3-exp5-p258-groupoid-isomorphism-map}
  {Polo--Gille Exp5-13oct24.pdf, component-local p. 8; combined-reader p. 270; page header p. 258}

\noindent defines an isomorphism from the groupoid \(Y_\ast(T)\) onto
\(Y'_\ast(T)\). Moreover, this isomorphism depends functorially on \(T\),
so the \(\mathscr C\)-groupoids \(Y_\ast\) and \(Y'_\ast\) are
isomorphic.%
\footnote{\label{note:sga3-V-13}Editor's note: this plays a crucial role
in the proof of \hyperref[lem:sga3-V-6.1]{Lemma~6.1}.}

% Source page: Exp5 component-local p. 9; combined-reader p. 271;
% page header p. 259.
\medskip
\phantomsection\label{par:sga3-V-3-d}%
\noindent\textbf{d)}

\medskip
\phantomsection\label{prop:sga3-V-3.1}%
\noindent\textbf{Proposition 3.1.}---
\textit{Keep the notation of \hyperref[par:sga3-V-3-a]{3(a)}, and suppose that \(f_0\) is an
effective and universal epimorphism. Then
\(\operatorname{Coker}(d_0,d_1)\) exists if and only if
\(\operatorname{Coker}(e_0,e_1)\) exists.%
\footnote{\label{note:sga3-V-14}Editor's note: the original has been
modified to display explicitly the isomorphism below.}
Moreover, in this case \(f_0\) induces an isomorphism}
\phantomsection\label{eq:sga3-V-3.1-cokernel-isomorphism}%
\[
  \operatorname{Coker}(d_0,d_1)
  \xrightarrow{\ \sim\ }
  \operatorname{Coker}(e_0,e_1).
\]

Recall first that an epimorphism \(f_0:Y_0\to X_0\) is called
\emph{universal} if, in every Cartesian square

\SGAThreeDiagram[0.18\textwidth]
  {figures/exp5/p259_universal_epimorphism_square.png}
  {A base change of the universal epimorphism \(f_0\).}
  {fig:sga3-exp5-p259-universal-epimorphism-square}
  {Polo--Gille Exp5-13oct24.pdf, component-local p. 9; combined-reader p. 271; page header p. 259}

\noindent the arrow \(f'\) is an epimorphism. Now let \(C(d_0,d_1)\)
denote the covariant functor from \(\mathscr C\) to sets which assigns to
every \(T\in\mathscr C\) the kernel of the pair
\[
  T(d_0),T(d_1):T(X_0)\rightrightarrows T(X_1),
\]
and define \(C(e_0,e_1)\) in the same way. For every
\(T\in\mathscr C\), there is therefore a commutative diagram

\SGAThreeDiagram[0.51\textwidth]
  {figures/exp5/p259_functor_kernel_diagram.png}
  {The comparison of the functors \(C(d_0,d_1)\) and \(C(e_0,e_1)\).}
  {fig:sga3-exp5-p259-functor-kernel-diagram}
  {Polo--Gille Exp5-13oct24.pdf, component-local p. 9; combined-reader p. 271; page header p. 259}

Here \(T(f)\) is the injection induced by the injection \(T(f_0)\). If
we show that \(T(f)\) is surjective for every \(T\), we obtain a
functorial isomorphism
\[
  f:C(d_0,d_1)\xrightarrow{\ \sim\ }C(e_0,e_1).
\]
Thus one of these functors is representable if and only if the other is,
which proves \hyperref[prop:sga3-V-3.1]{the proposition}.

To prove that \(T(f)\) is surjective, consider the diagram

\SGAThreeDiagram[0.52\textwidth]
  {figures/exp5/p259_delta_proof_diagram.png}
  {The diagram used to prove the surjectivity of \(T(f)\).}
  {fig:sga3-exp5-p259-delta-proof-diagram}
  {Polo--Gille Exp5-13oct24.pdf, component-local p. 9; combined-reader p. 271; page header p. 259}

Here \(\Delta\) is the section of \(Y_1\to Y_0\times Y_0\) defined by
the morphism
\[
  s\circ f_0\circ\operatorname{pr}_1:Y_0\times Y_0\longrightarrow X_1,
\]
where \(s:X_0\to X_1\) satisfies
\hyperref[eq:sga3-exp5-3]{(3)} and
\hyperref[eq:sga3-exp5-3bis]{(3 bis)} of \hyperref[sec:sga3-V-1]{Section~1}.

% Source page: Exp5 component-local p. 10; combined-reader p. 272;
% page header p. 260. Legacy marginal markers on this page reach 262.
If \(g:Y_0\to T\) satisfies \(g\circ e_0=g\circ e_1\), then
\[
  g\circ e_0\circ\Delta=g\circ e_1\circ\Delta,
\]
so \(g\circ\operatorname{pr}_1=g\circ\operatorname{pr}_2\). Since
\(f_0\) is an effective epimorphism, \(g\) factors as \(h\circ f_0\)
for an arrow \(h:X_0\to T\); in other words,
\(g=T(f_0)(h)\). It remains to prove that
\(h\in C(d_0,d_1)(T)\), or equivalently that \(hd_0=hd_1\). But
\[
  hd_0f_1=hf_0e_0=ge_0=ge_1=hf_0e_1=hd_1f_1,
\]
and the desired equality follows because \(f_1\) is an epimorphism
(since \(f_0\) is a universal epimorphism).

\medskip
\phantomsection\label{par:sga3-V-3-e}%
\noindent\textbf{e)} Now let \(S\) be a scheme and take
\(\mathscr C=(\mathbf{Sch}/S)\). The data of a \(\mathscr C\)-groupoid

\SGAThreeDiagram[0.37\textwidth]
  {figures/exp5/p260_scheme_groupoid.png}
  {A \((\mathbf{Sch}/S)\)-groupoid.}
  {fig:sga3-exp5-p260-scheme-groupoid}
  {Polo--Gille Exp5-13oct24.pdf, component-local p. 10; combined-reader p. 272; page header p. 260}

\noindent defines an equivalence relation on the set \(\underline{X_0}\)
underlying the scheme \(X_0\): for \(x,y\in\underline{X_0}\), write
\(x\sim y\) if there exists \(z\in\underline{X_1}\) such that
\(x=d_1z\) and \(y=d_0z\). Reflexivity and symmetry are clear.%
\footnote{\label{note:sga3-V-15}Editor's note: reflexivity follows from
the existence of \(s:X_0\to X_1\), which is a section of both \(d_0\)
and \(d_1\); symmetry follows from the existence of the involution
\(\sigma\) of \(X_1\) which ``exchanges \(d_0\) and \(d_1\)'', that is,
which satisfies \(d_0\sigma=d_1\) and \(d_1\sigma=d_0\); cf. \hyperref[sec:sga3-V-1]{Section~1},
\hyperref[eq:sga3-exp5-3]{(3)},
\hyperref[eq:sga3-exp5-3bis]{(3 bis)}, and
\hyperref[eq:sga3-exp5-dagger]{(\(\dagger\))}.}
Let us prove transitivity. If \(x\sim y\) and \(y\sim z\), there exist
\(u,v\in\underline{X_1}\) such that
\[
  x=d_1u,\qquad y=d_0u,\qquad y=d_1v,\qquad z=d_0v.
\]
It follows that \((v,u)\) belongs to the set-theoretic fiber product
\(\underline{X_1}\times_{\underline{d_1},\underline{d_0}}
\underline{X_1}\). Since the canonical map

\SGAThreeDiagram[0.32\textwidth]
  {figures/exp5/p260_underlying_fiber_product_map.png}
  {The canonical map from the underlying set of the fiber product to the
  fiber product of the underlying sets.}
  {fig:sga3-exp5-p260-underlying-fiber-product-map}
  {Polo--Gille Exp5-13oct24.pdf, component-local p. 10; combined-reader p. 272; page header p. 260}

\noindent is surjective, \((v,u)\) is the image of some
\(w\in\underline{X_2}\). Hence
\[
  x=d_1d'_1w,\qquad z=d_0d'_1w,
\]
and therefore \(x\sim z\).

\medskip
\phantomsection\label{par:sga3-V-3-f}%
\noindent\textbf{f)} Keep the notation of \hyperref[par:sga3-V-3-a]{3(a)} and \hyperref[par:sga3-V-3-b]{3(b)}, still taking
\(\mathscr C=(\mathbf{Sch}/S)\). If \(x,y\) are points of \(Y_0\), then
\[
  x\sim y\quad\Longleftrightarrow\quad f_0(x)\sim f_0(y).
\]
Thus the inverse image of the equivalence relation defined by a groupoid
is the equivalence relation defined by the inverse image of that
groupoid.

Indeed, suppose that \(x\sim y\). There is a \(z\in\underline{Y_1}\)
such that \(x=e_1(z)\) and \(y=e_0(z)\). Since
\(f_0\circ e_i=d_i\circ f_1\) for \(i=0,1\), it follows that
\[
  f_0(x)=d_1f_1(z),\qquad f_0(y)=d_0f_1(z),
\]
and hence \(f_0(x)\sim f_0(y)\).

Conversely, suppose that \(f_0(x)\sim f_0(y)\), and let
\(z\in\underline{X_1}\) satisfy
\[
  f_0(y)=d_1(z),\qquad f_0(x)=d_0(z).
\]
Using the construction and notation of \hyperref[par:sga3-V-3-b]{3(b)}, there is a point
\(t\in\underline{Y_0\times_{X_0}X_1}\) such that
\(\operatorname{pr}_1(t)=x\) and \(\operatorname{pr}_2(t)=z\). Since
\(f_0(y)=d_1\operatorname{pr}_2(t)\), there is likewise an
\(s\in\underline{Y_1}\) such that
\[
  y=e_1(s),\qquad (e_0\boxtimes f_1)(s)=t.
\]
It follows that
\[
  e_0(s)=\operatorname{pr}_1(e_0\boxtimes f_1)(s)
  =\operatorname{pr}_1(t)=x,
\]
so \(x\sim y\).

% END INCLUDED SGA3_English_current_buildable_source_20260823/source/inputs/V/tex/components/02_expose_V_pp256_260_en.tex
%
  \IfFileExists{tex/components/03_expose_V_section4_theorem41_localp011_en.tex}{%
    % BEGIN INCLUDED SGA3_English_current_buildable_source_20260823/source/inputs/V/tex/components/03_expose_V_section4_theorem41_localp011_en.tex
% SGA 3, Expose V, Loop 1 English reconstruction.
% Scope: section 4 heading, Theorem 4.1, editor's notes 16--20, and the
% paragraph announcing the proof strategy.
% Authority: Polo--Gille Exp5-13oct24.pdf, component-local p. 11
% (combined-reader p. 273). This re-edited page bears printed header 261
% and internal original-page markers 263--264.
% Jacob Reinhold's CC BY 4.0 Markdown is comparison/drafting material only.

\section*{4. Passage to the quotient by a finite flat groupoid
  (proof of a special case)}
\phantomsection\label{sec:sga3-V-4}
\addcontentsline{toc}{section}{4. Passage to the quotient by a finite flat groupoid}

\medskip
\phantomsection\label{thm:sga3-V-4.1}
\noindent\textbf{Theorem 4.1.}---
\textit{Consider a \((\mathbf{Sch}/S)\)-groupoid}

% Native Loop-2 reconstruction; local p. 11 / combined p. 273.
\begin{equation*}
  \resizebox{0.40\textwidth}{!}{%
    \SGAThreeNativeInput{figures/exp5/p261_theorem41_groupoid.png}}
  \phantomsection\label{fig:sga3-exp5-theorem41-groupoid}
\end{equation*}

\textit{and suppose that the following conditions hold:}%
\footnote{\label{note:sga3-V-16}Editor's note: since
\(d_0=d_1\sigma\), where \(\sigma\) is an involutive automorphism of
\(X_1\), these two conditions are symmetric in \(d_1\) and \(d_0\).
Moreover,
\(d_0d_1^{-1}(x)=d_1d_0^{-1}(x)\).}
\begin{enumerate}
\item[\textit{a)}]\phantomsection\label{thm:sga3-V-4.1-a}%
  \textit{\(d_1\) is finite locally free;}
\item[\textit{b)}]\phantomsection\label{thm:sga3-V-4.1-b}%
  \textit{for every \(x\in X_0\), the set
  \(d_0d_1^{-1}(x)\) is contained in an affine open subset of
  \(X_0\).}%
  \footnote{\label{note:sga3-V-17}Editor's note: \hyperref[thm:sga3-V-4.1-b]{condition~b)} cannot
  be omitted. H.~Hironaka gave an example of an action of the finite
  group \(G=\mathbb Z/2\mathbb Z\) on a proper complex variety \(X_0\)
  for which the quotient \(X_0/G\) is an algebraic space that is not a
  scheme; see \(\hyperref[bib:sga3-V-Hi62]{[{\rm Hi62}]}\), and also
  \(\hyperref[bib:sga3-V-Mum65]{[{\rm Mum65}]}\), Chapter~4, \S3.}
\end{enumerate}
\textit{Then:}
\begin{enumerate}
\item[(i)]\phantomsection\label{thm:sga3-V-4.1-i}%
  \textit{there exists a cokernel \((Y,p)\) of \((d_0,d_1)\) in
  \((\mathbf{Sch}/S)\). Moreover, such a pair \((Y,p)\) is a
  cokernel of \((d_0,d_1)\) in the category of all ringed spaces;}
\item[\SGAInlineRefTarget{thm:sga3-V-4.1:item:ii}(ii)]\phantomsection\label{thm:sga3-V-4.1-ii}%
  \textit{\(p\) is integral and open, and \(Y\) is affine when
  \(X_0\) is affine;}%
  \footnote{\label{note:sga3-V-18}Editor's note: the assertion that
  \(p\) is open has been added by repeating the analogous argument
  given in \hyperref[lem:sga3-V-6.1]{6.1}.}
\item[(iii)]\phantomsection\label{thm:sga3-V-4.1-iii}%
  \textit{the morphism
  \[
    X_1\longrightarrow X_0\times_YX_0
  \]
  with components \(d_0\) and \(d_1\) is surjective;}
\item[(iv)]\phantomsection\label{thm:sga3-V-4.1-iv}%
  \textit{if \((d_0,d_1)\) is an equivalence pair, then
  \(X_1\to X_0\times_YX_0\) is an isomorphism}%
  \footnote{\label{note:sga3-V-19}Editor's note: in this case
  \(X_1\to X_0\times_SX_0\) is therefore an immersion
  (EGA~I, 5.3.10); see also Expos\'e~VI\textsubscript{B}, 9.2.1.
  For existence of the quotient---in the category of schemes or in
  that of algebraic spaces---under the weaker hypothesis that
  \(d_0\) and \(d_1\) are finite but not necessarily flat, see
  \(\hyperref[bib:sga3-V-An73]{[{\rm An73}]}\), \S1.1, \(\hyperref[bib:sga3-V-Fe03]{[{\rm Fe03}]}\), \(\hyperref[bib:sga3-V-Ko08]{[{\rm Ko08}]}\), and
  the references there.}
  \textit{and \(p:X_0\to Y\) is finite locally free.}%
  \footnote{\label{note:sga3-V-20}Editor's note: the consequences that
  follow have been made explicit; see \(\hyperref[bib:sga3-V-Ray67a]{[{\rm Ray67a}]}\),
  Theorem~1(iii), and \hyperref[proof:sga3-V-5-b-general]{the proof given below} at the end of \hyperref[sec:sga3-V-5]{Section~5}.}
  \textit{Moreover, \((Y,p)\) is a cokernel of \((d_0,d_1)\) in the
  category of sheaves for the fppf topology and, for every base change
  \(Y'\to Y\), the scheme \(Y'\) is the cokernel of the groupoid
  \(X_\ast\times_YY'\) obtained from \(X_\ast\) by the base change
  \(X_0\times_YY'\to X_0\).}
\end{enumerate}
\textit{In particular, for every base change \(S'\to S\), the scheme
\(Y'=Y\times_SS'\) is the cokernel of the \(S'\)-groupoid
\(X'_\ast=X_\ast\times_SS'\). Thus, in this case, formation of the
quotient commutes with base change.}

It follows immediately from \hyperref[thm:sga3-V-4.1-i]{(i)} that the underlying topological space
of \(Y\) is the quotient of the underlying topological space of
\(X_0\) by the equivalence relation defined by the
\((\mathbf{Sch}/S)\)-groupoid \(X_\ast\).

We shall first prove \hyperref[thm:sga3-V-4.1]{the theorem} when \(X_0\) is affine and \(d_1\) is
locally free of constant rank \(n\). We shall then see how to reduce
the general situation to this special case.
% SGA3-V-V2-EOL-PRESERVATION

% END INCLUDED SGA3_English_current_buildable_source_20260823/source/inputs/V/tex/components/03_expose_V_section4_theorem41_localp011_en.tex
%
    \IfFileExists{tex/components/04_expose_V_theorem41_proof_a_localp012_en.tex}{%
      % BEGIN INCLUDED SGA3_English_current_buildable_source_20260823/source/inputs/V/tex/components/04_expose_V_theorem41_proof_a_localp012_en.tex
% SGA 3, Expose V, Loop 1 English reconstruction.
% Scope: affine/ring setup for Theorem 4.1 and part a) of its proof,
% stopping immediately before part b).
% Authority: Polo--Gille Exp5-13oct24.pdf, component-local p. 12
% (combined-reader p. 274). This re-edited page bears printed header 262
% and internal original-page markers 264--265.
% Jacob Reinhold's CC BY 4.0 Markdown is comparison/drafting material only.

In the case now under consideration, \(X_0,X_1,X_2\) are affine. We
may therefore suppose that
\[
  X_i=\operatorname{Spec}A_i,\qquad
  d_j=\operatorname{Spec}\delta_j,\qquad
  d'_k=\operatorname{Spec}\delta'_k,
\]
where the \(A_i\) are commutative rings and the
\(\delta_j,\delta'_k\) are ring homomorphisms. The diagram
\hyperref[sga3:v:lemma:1-2:equation:0-1-2]{(0,1,2)} may then be replaced by
the following dual diagram:

% Native Loop-2 reconstruction; local p. 12 / combined p. 274.
\begin{equation*}
  \resizebox{0.57\textwidth}{!}{%
    \SGAThreeNativeInput{figures/exp5/p262_dual_cocartesian_012star.png}}
  \phantomsection\label{fig:sga3-exp5-dual-012star}
\end{equation*}
The two left-hand squares in this diagram are cocartesian.

Let \(B\) be the subring of \(A_0\) consisting of the elements
\(a\in A_0\) for which \(\delta_0(a)=\delta_1(a)\).

\medskip
\phantomsection\label{proof:sga3-V-4.1-a}
\noindent\textbf{a)} We first show that \(A_0\) is integral over
\(B\). For \(a\in A_0\), let
\begin{equation}
  P_{\delta_1}(T,\delta_0(a))
  =T^n-\sigma_1T^{n-1}+\cdots+(-1)^n\sigma_n
  \label{eq:sga3-V-4.1-characteristic-polynomial}
\end{equation}
be the characteristic polynomial of \(\delta_0(a)\), where \(A_1\)
is regarded as an \(A_0\)-algebra through \(\delta_1\); see
Bourbaki, \emph{Alg\`ebre}, VIII, \S12, and
\emph{Alg\`ebre commutative}, II, \S5, Exercise~9.
Since the two left-hand squares of the diagram
\hyperref[fig:sga3-exp5-dual-012star]{\((0,1,2)^\ast\)} are
cocartesian, one has
\begin{align}
 \delta_0\!\left(P_{\delta_1}(T,\delta_0(a))\right)
   &=P_{\delta'_2}(T,\delta'_0\delta_0(a)), \\
 \delta_1\!\left(P_{\delta_1}(T,\delta_0(a))\right)
   &=P_{\delta'_2}(T,\delta'_1\delta_0(a)).
\end{align}
Because \(\delta'_0\delta_0=\delta'_1\delta_0\), it follows that
\[
 \delta_0\!\left(P_{\delta_1}(T,\delta_0(a))\right)
 =
 \delta_1\!\left(P_{\delta_1}(T,\delta_0(a))\right);
\]
equivalently,
\[
  \delta_0(\sigma_i)=\delta_1(\sigma_i)
  \qquad\text{for every }i.
\]

On the other hand, the Cayley--Hamilton theorem gives
\[
 \delta_0(a)^n-\delta_1(\sigma_1)\delta_0(a)^{n-1}
 +\cdots+(-1)^n\delta_1(\sigma_n)=0.
\]
Since \(\delta_1(\sigma_i)=\delta_0(\sigma_i)\), one also has
\[
 \delta_0(a)^n-\delta_0(\sigma_1)\delta_0(a)^{n-1}
 +\cdots+(-1)^n\delta_0(\sigma_n)=0.
\]
Hence
\[
  a^n-\sigma_1a^{n-1}+\cdots+(-1)^n\sigma_n=0,
\]
because there is a homomorphism
\(\tau:A_1\to A_0\) such that
\(\tau\delta_0=\operatorname{id}_{A_0}\), so \(\delta_0\) is
injective. It follows that \(A_0\) is integral over \(B\).

% END INCLUDED SGA3_English_current_buildable_source_20260823/source/inputs/V/tex/components/04_expose_V_theorem41_proof_a_localp012_en.tex
%
      \IfFileExists{tex/components/05_expose_V_theorem41_proof_b_lemma411_en.tex}{%
        % BEGIN INCLUDED SGA3_English_current_buildable_source_20260823/source/inputs/V/tex/components/05_expose_V_theorem41_proof_b_lemma411_en.tex
% SGA 3, Expose V, Loop 1 English reconstruction.
% Scope: Theorem 4.1 proof part b), followed by Lemma 4.1.1 and its proof;
% stops immediately before part c), Proof of (i).
% Authority: Polo--Gille Exp5-13oct24.pdf, component-local pp. 12--13
% (combined-reader pp. 274--275). The re-edited page headers are 262--263;
% internal original-page markers in this range advance through 265--266.
% Jacob Reinhold's CC BY 4.0 Markdown is comparison/drafting material only.

\medskip
\phantomsection\label{proof:sga3-V-4.1-b}%
\noindent\textbf{b)} Consider two prime ideals \(x\) and \(y\) of
\(A_0\). We shall show that the equality
\[
  x\cap B=y\cap B
\]
implies the existence of a prime ideal \(z\) of \(A_1\) such that
\[
  x=d_0(z),\qquad y=d_1(z).
\]

Indeed, if this assertion were false, \(x\) would be distinct from
\(\delta_0^{-1}(t)\) for every prime ideal \(t\) of \(A_1\) such that
\(\delta_1^{-1}(t)=y\). For every such \(t\), one would have
\[
  \delta_0^{-1}(t)\cap B
  =\delta_1^{-1}(t)\cap B
  =y\cap B=x\cap B.
\]
It would follow from Cohen--Seidenberg (see Bourbaki,
\emph{Alg\`ebre commutative}, V, \S2, Corollary~1 to Proposition~1)
that \(x\) is contained in none of the ideals
\(\delta_0^{-1}(t)\).%
\footnote{\label{note:sga3-V-21}Editor's note: the original has been expanded in what follows;
in particular, Lemma~\hyperref[lem:sga3-V-4.1.1]{4.1.1}, taken from \hyperref[bib:sga3-V-DG70]{[DG70]}, III,
\S2.4, Lemma~4.3, has been added.}
There are at most \(n\) prime ideals \(t\) of \(A_1\) satisfying
\(\delta_1^{-1}(t)=y\) (see \emph{loc. cit.}, Proposition~3). Hence,
by the Prime Avoidance Lemma (\emph{loc. cit.}, II, \S1,
Proposition~3), there would be an element \(a\in x\) belonging to none
of the ideals \(\delta_0^{-1}(t)\).

Consequently, \(\delta_0(a)\) would belong to none of these prime
ideals \(t\), and \hyperref[lem:sga3-V-4.1.1]{the lemma below} would imply that the norm
\(N_{\delta_1}(\delta_0(a))\) does not belong to \(y\). Here the norm
is computed by regarding \(A_1\) as an \(A_0\)-algebra through
\(\delta_1\); with the notation of \hyperref[proof:sga3-V-4.1-a]{part~a)},
\[
  N_{\delta_1}(\delta_0(a))=\sigma_n.
\]
On the other hand,
\[
  (-1)^{n-1}\sigma_n
  =a^n+\sum_{i=1}^{n-1}(-1)^i\sigma_i a^{n-i},
\]
so this norm belongs to \(B\cap x=B\cap y\), a contradiction.

\medskip
\phantomsection\label{lem:sga3-V-4.1.1}%
\noindent\textbf{Lemma 4.1.1.}---
\textit{Let \(A\to A'\) be a homomorphism of commutative rings which
makes \(A'\) a projective \(A\)-module of rank \(n\). Let
\(p\in\operatorname{Spec}(A)\), let \(q_1,\ldots,q_r\) be the points of
\(\operatorname{Spec}(A')\) lying above \(p\), and let \(a\in A'\).
Then \(a\in q_1\cup\cdots\cup q_r\) if and only if its norm \(N(a)\)
belongs to \(p\).}

Indeed, after replacing \(A\) and \(A'\) by the localizations \(A_p\)
and \(A'_p\), we reduce to the case in which \((A,p)\) is local and
\(A'\) is semilocal, with
\[
  \operatorname{Spec}(A')=\{q_1,\ldots,q_r\}.
\]
In this case \(A'\) is a free \(A\)-module of rank \(n\) (see
Bourbaki, \emph{Alg\`ebre commutative}, II, \S3.2, Corollary~2 to
Proposition~5), and \(N(a)\) is the determinant of the endomorphism
\[
  \ell_a:A'\longrightarrow A',\qquad a'\longmapsto aa'.
\]
Thus one has the equivalences
\[
  N(a)\notin p
  \quad\Longleftrightarrow\quad N(a)\text{ is invertible}
  \quad\Longleftrightarrow\quad \ell_a\text{ is invertible}
  \quad\Longleftrightarrow\quad
  a\notin q_1\cup\cdots\cup q_r.
\]

% END INCLUDED SGA3_English_current_buildable_source_20260823/source/inputs/V/tex/components/05_expose_V_theorem41_proof_b_lemma411_en.tex
%
        \IfFileExists{tex/components/06_expose_V_theorem41_proof_cde_lemma412_en.tex}{%
          % BEGIN INCLUDED SGA3_English_current_buildable_source_20260823/source/inputs/V/tex/components/06_expose_V_theorem41_proof_cde_lemma412_en.tex
% SGA 3, Expose V, Loop 1 English reconstruction.
% Scope: Theorem 4.1 proof parts c), d), and opening of e), including
% Lemma 4.1.2 and its proof; stops before the proof of Theorem 4.1(iii).
% Authority: Polo--Gille Exp5-13oct24.pdf, component-local pp. 13--14
% (combined-reader pp. 275--276). Re-edited page headers 263--264;
% internal original-page markers advance through 266.
% Jacob Reinhold's CC BY 4.0 Markdown is comparison/drafting material only.

\medskip
\phantomsection\label{proof:sga3-V-4.1-c}%
\noindent\textbf{c) Proof of \hyperref[thm:sga3-V-4.1-i]{(i)}.}
Set \(Y=\operatorname{Spec}B\) and \(p=\operatorname{Spec}i\), where
\(i:B\hookrightarrow A_0\) is the inclusion. By \hyperref[proof:sga3-V-4.1-a]{part~a)}, the morphism
\(p:X_0\to Y\) is surjective. We first show that \((Y,p)\) is a
cokernel of \((d_0,d_1)\) in the category of all ringed spaces.
Indeed, \hyperref[proof:sga3-V-4.1-b]{part~b)} shows that the set underlying \(\operatorname{Spec}B\)
is obtained from the set underlying \(X_0\) by identifying the points
\(x\) and \(y\) for which there exists \(z\in X_1\) such that
\[
  d_1z=y,\qquad d_0z=x.
\]
Moreover, since \(i\) is integral, \(p=\operatorname{Spec}i\) is
closed, and therefore \(Y\) carries the quotient topology induced by
\(X_0\).

It follows that \(p\) is open. Indeed, let \(U'\) be any open subset
of \(X_0\). Since \(d_1\) is surjective and finite locally free, it is
faithfully flat and of finite presentation, hence open. Thus
\[
  U=d_1\!\left(d_0^{-1}(U')\right),
\]
the saturation of \(U'\) for the equivalence relation defined by
\(X_\ast\), is open. Since \(Y\) has the quotient topology,
\(p(U')=p(U)\) is open.

Finally, the choice of \(B\), together with the fact that \(p,d_0\),
and \(d_1\) are affine, shows that the canonical sequence of sheaves
of rings
\[
  \mathcal O_Y\longrightarrow p_\ast\mathcal O_{X_0}
  \underset{p_\ast(\delta_0)}{\overset{p_\ast(\delta_1)}
    {\rightrightarrows}}
  p_\ast d_{0\ast}\mathcal O_{X_1}
  =p_\ast d_{1\ast}\mathcal O_{X_1}
\]
is exact.

It remains to show that \((Y,p)\) is also the cokernel of
\((d_0,d_1)\) in the category of schemes---more generally, in the
category of locally ringed spaces. Let \(q:X_0\to Z\) be a morphism of
schemes such that \(qd_0=qd_1\). By what precedes, there is a unique
morphism of ringed spaces \(r:Y\to Z\) such that \(q=rp\). We must
show that, for every \(y\in Y\), the homomorphism
\[
  \mathcal O_{Z,r(y)}\longrightarrow\mathcal O_{Y,y}
\]
induced by \(r\) is local. Since \(p\) is surjective, \(y=p(x)\) for
some \(x\in X_0\), and the assertion follows from the fact that the
homomorphism
\[
  \mathcal O_{Z,q(x)}\longrightarrow\mathcal O_{X_0,x}
\]
induced by \(q\) is local.

\medskip
\phantomsection\label{proof:sga3-V-4.1-d}%
\noindent\textbf{d) Proof of \hyperref[thm:sga3-V-4.1-ii]{(ii)}.}
This follows from \hyperref[proof:sga3-V-4.1-a]{parts~a)} and~\hyperref[proof:sga3-V-4.1-c]{c)}.

\medskip
\phantomsection\label{proof:sga3-V-4.1-e}%
\noindent\textbf{e) Proof of \hyperref[thm:sga3-V-4.1-iii]{(iii)}.}
Recall that \(\underline P\) denotes the set underlying a scheme
\(P\), and that \(\underline d:\underline P\to\underline Q\) denotes
the map induced by a morphism \(d:P\to Q\).

\medskip
\phantomsection\label{lem:sga3-V-4.1.2}%
\noindent\textbf{Lemma 4.1.2.}---%
\footnote{\label{note:sga3-V-22}Editor's note: \hyperref[lem:sga3-V-4.1.2]{this lemma}, which is used repeatedly in this
Expos\'e and in later Expos\'es (\(\mathrm{VI_A}\), \(\mathrm{VI_B}\)),
has been inserted here. It appeared as Lemma \(\mathrm{VI_B}\), 4.5.1,
in the original 1965 edition of SGA~3.}
\textit{Let \((A,\mathfrak m)\) be a local ring, let \(k\) be its
residue field, and let \(K/k\) be a field extension. There exists a
local, flat \(A\)-algebra \(B\) such that \(B/\mathfrak mB\) is
\(k\)-isomorphic to \(K\). If \(K/k\) is finite, \(B\) may moreover
be chosen finite and free over \(A\).}

This is proved in EGA \(0_{\mathrm{III}}\), 10.3.1, where it is also
shown that \(B\) may be chosen Noetherian when \(A\) is Noetherian.
For the reader's convenience, we recall the proof.

Put \(A'=A[T]\), where \(T\) is an indeterminate. If \(K=k(T)\), let
\(\mathfrak p=\mathfrak mA'\) and \(B=A'_{\mathfrak p}\). Then
\[
  B/\mathfrak mB\simeq k[T]_{(0)}=k(T).
\]
The ring \(B\) is flat over \(A'\), which is a free \(A\)-module, so
\(B\) is flat over \(A\).

Suppose next that \(K=k(t)=k[t]\), where \(t\) is algebraic over
\(k\). Set \(B=A'/(F)\), where \(F\in A'\) is a monic polynomial whose
image in \(k[T]\) is the minimal polynomial \(f\) of \(t\) over \(k\).
Then \(B\) is a free \(A\)-module of finite rank
\[
  \deg(F)=\deg(f).
\]
In particular, \(B\) is integral over \(A\), and hence every maximal
ideal of \(B\) contains \(\mathfrak m\). Since
\[
  B/\mathfrak mB\simeq k[T]/(f)\simeq K,
\]
\(B\) is local, with maximal ideal \(\mathfrak mB\). This already
proves that \(B\) can be chosen finite and free over \(A\) when
\([K:k]<\infty\).

In the general case, choose a family \((t_i)_{i\in I}\) of generators
of \(K/k\), and well-order \(I\); that is, give \(I\) a total order
\(\leq\) in which every nonempty subset has a least element. For
\(i\in I\), let \(k_i\) (respectively \(k_{<i}\)) be the subfield of
\(K\) generated by the \(t_j\) with \(j\leq i\) (respectively
\(j<i\)). By adding one element if necessary, we may suppose that
\(I\) has a greatest element \(\xi\), so that \(K=k_\xi\).

Let \(J\subseteq I\) be the set of indices \(i\) for which there
exists a direct system \((A_j)_{j\leq i}\) of local, flat
\(A\)-algebras such that
\[
  A_j/\mathfrak mA_j\simeq k_j
\]
and \(A_j\) is flat over \(A_\ell\) whenever \(\ell<j\). Suppose that
\(I\setminus J\) is nonempty, let \(i\) be its least element, and put
\[
  A'=\varinjlim_{j<i}A_j.
\]
Since tensor products commute with direct limits, \(A'\) is flat over
\(A\) and over every \(A_j\), \(j<i\), and
\[
  A'/\mathfrak mA'
  \simeq A'\otimes_A(A/\mathfrak m)
  \simeq k_{<i}.
\]
Moreover, \(A'\) is local, with maximal ideal \(\mathfrak mA'\).
Indeed, if \(x=f_j(x_j)\) is not invertible, then \(x_j\) is not
invertible and hence belongs to the maximal ideal
\(\mathfrak mA_j\) of \(A_j\); therefore \(x\in\mathfrak mA'\).

The monogenic case treated above now gives a local, flat
\(A'\)-algebra \(A_i\) such that
\[
  A_i/\mathfrak mA_i\simeq k_{<i}(t_i)=k_i.
\]
The algebra \(A_i\) is flat over every \(A_j\), \(j<i\), so \(i\in J\),
contrary to the choice of \(i\). This contradiction proves that
\(J=I\), and \(A_\xi\) has the required properties. This proves
Lemma~\hyperref[lem:sga3-V-4.1.2]{4.1.2}.

% END INCLUDED SGA3_English_current_buildable_source_20260823/source/inputs/V/tex/components/06_expose_V_theorem41_proof_cde_lemma412_en.tex
%
          \IfFileExists{tex/components/07_expose_V_theorem41_proof_iii_lemma413_en.tex}{%
            % BEGIN INCLUDED SGA3_English_current_buildable_source_20260823/source/inputs/V/tex/components/07_expose_V_theorem41_proof_iii_lemma413_en.tex
% SGA 3, Expose V, Loop 1 English reconstruction.
% Scope: proof of Theorem 4.1(iii), then the opening of part f) through
% Lemma 4.1.3 and its proof; stops before the proof of Theorem 4.1(iv).
% Authority: Polo--Gille Exp5-13oct24.pdf, component-local pp. 15--16
% (combined-reader pp. 277--278). Re-edited page headers 265--266;
% internal original-page marker 267 occurs on local p. 15.
% Jacob Reinhold's CC BY 4.0 Markdown is comparison/drafting material only.

Let us now prove \hyperref[thm:sga3-V-4.1-iii]{Theorem~4.1(iii)}. Recall once more that
\(\underline P\) denotes the set underlying a scheme \(P\), and that
\(\underline d:\underline P\to\underline Q\) is the map induced by a
morphism \(d:P\to Q\). \hyperref[proof:sga3-V-4.1-b]{Part~b)} can then be restated by saying that the
map
\[
  \underline{d_0}\boxtimes\underline{d_1}:
  \underline{X_1}\longrightarrow
  \underline{X_0}\times_{\underline Y}\underline{X_0}
\]
with components \(\underline{d_0}\) and \(\underline{d_1}\) is
surjective. This map factors as
\[
  \underline{X_1}
  \xrightarrow{\ \underline{d_0}\boxtimes\underline{d_1}\ }
  \underline{X_0}\times_{\underline Y}\underline{X_0}
  \xrightarrow{\ q\ }
  \underline{X_0\times_YX_0},
\]
where \(q\) is the canonical map. The image of
\(\underline{d_0}\boxtimes\underline{d_1}\) therefore contains every
point \(v\in\underline{X_0\times_YX_0}\) such that
\[
  \{v\}=q^{-1}(q(v)).
\]
This condition is satisfied in particular when \(v\) is
\emph{rational over \(Y\)}, that is, when the residue field
\(\kappa(v)\) is identified with the residue field \(\kappa(w)\) of
the image \(w\) of \(v\) in \(Y\).%
\footnote{\label{note:sga3-V-23}Editor's note: the pages of Lecture Notes in Mathematics
151 are permuted here; their actual order is
265--266--268--269--267--270--271.}

If \(v\in\underline{X_0\times_YX_0}\) is not rational over \(Y\),
again let \(w\)
be its image in \(Y\). By
Lemma~\hyperref[lem:sga3-V-4.1.2]{4.1.2}, there exist a local ring
\(C\) with radical \(\mathfrak m\) and a local, flat homomorphism
\[
  f:\mathcal O_{Y,w}\longrightarrow C
\]
such that \(C/\mathfrak m\) is isomorphic to \(\kappa(v)\) as a
\(\kappa(w)\)-algebra. Put \(Y'=\operatorname{Spec}C\), and let
\(\pi:Y'\to Y\) be the morphism induced by \(f\). The canonical
projection
\[
  (X_0\times_YX_0)\times_YY'\longrightarrow X_0\times_YX_0
\]
sends some point \(v'\), rational over \(Y'\), to \(v\). Since
\[
  (X_0\times_YX_0)\times_YY'
  \simeq
  (X_0\times_YY')\times_{Y'}(X_0\times_YY'),
\]
and since the hypotheses of \hyperref[thm:sga3-V-4.1]{Theorem~4.1} and the preceding results,
in particular \hyperref[proof:sga3-V-4.1-b]{part~b)}, remain valid after the base change
\(\pi:Y'\to Y\), the point \(v'\) is the image of an element
\[
  u'\in\underline{X_1\times_YY'}
\]
under the base change of \(d_0\boxtimes d_1\). If \(u\) is the image
of \(u'\) in \(X_1\), then
\[
  v=(d_0\boxtimes d_1)(u),
\]
as required.

\medskip
\phantomsection\label{proof:sga3-V-4.1-f}%
\noindent\textbf{f) Proof of \hyperref[thm:sga3-V-4.1-iv]{(iv)}.}%
\footnote{\label{note:sga3-V-24}Editor's note: \hyperref[lem:sga3-V-4.1.3]{the following lemma} has been added from
\hyperref[bib:sga3-V-DG70]{[DG70]}, I, \S5, Proposition~1.5; see also the proof of EGA
\(\mathrm{IV}_3\), 8.11.5. It was used implicitly in the original and
explicitly in \hyperref[bib:sga3-V-DG70]{[DG70]}, III, \S2, 4.6. It is also useful in
\hyperref[thm:sga3-V-7.1]{Theorem~7.1 below}.}

\medskip
\phantomsection\label{lem:sga3-V-4.1.3}%
\noindent\textbf{Lemma 4.1.3.}---
\textit{If a monomorphism of schemes \(f:T\to Z\) is finite, it is a
closed immersion.}

Covering \(Z\) by affine open subsets \(Z_i\), and replacing \(f\) by
the induced morphisms \(f^{-1}(Z_i)\to Z_i\), reduces us---since a
finite morphism is affine---to the case
\[
  Z=\operatorname{Spec}B,\qquad T=\operatorname{Spec}A.
\]
Because \(f\) is a monomorphism, the diagonal morphism
\[
  T\longrightarrow T\times_ZT
\]
is an isomorphism (EGA~I, 5.3.8); equivalently, multiplication gives
an isomorphism
\[
  A\otimes_BA\xrightarrow{\ \sim\ }A.
\]
Consequently, for every maximal ideal \(\mathfrak m\) of \(B\), there
is an isomorphism
\[
  (A/\mathfrak mA)\otimes_k(A/\mathfrak mA)
  \xrightarrow{\ \sim\ }A/\mathfrak mA,
  \qquad k=B/\mathfrak m.
\]
Since \(A\) is finite over \(B\), the \(k\)-vector space
\(A/\mathfrak mA\) has finite dimension \(d\), and this isomorphism
implies \(d=0\) or \(d=1\). Hence the map
\[
  k=B/\mathfrak m\longrightarrow A/\mathfrak mA
\]
is surjective. Nakayama's lemma, applied after localization at
\(\mathfrak m\), shows that \(B_{\mathfrak m}\to A_{\mathfrak m}\)
is surjective. It follows that \(B\to A\) is surjective: its cokernel
\(C\) satisfies \(C_{\mathfrak m}=0\) for every maximal ideal
\(\mathfrak m\), and therefore vanishes. This proves \hyperref[lem:sga3-V-4.1.3]{the lemma}.

% END INCLUDED SGA3_English_current_buildable_source_20260823/source/inputs/V/tex/components/07_expose_V_theorem41_proof_iii_lemma413_en.tex
%
            \IfFileExists{tex/components/08_expose_V_theorem41_proof_iv_lemma42_en.tex}{%
              % BEGIN INCLUDED SGA3_English_current_buildable_source_20260823/source/inputs/V/tex/components/08_expose_V_theorem41_proof_iv_lemma42_en.tex
% SGA 3, Expose V, Loop 1 English reconstruction.
% Scope: continuation of Theorem 4.1(iv), including Lemma 4.2 and its
% application; stops immediately before the commutative diagram on local p. 17.
% Authority: Polo--Gille Exp5-13oct24.pdf, component-local p. 16
% (combined-reader p. 278), re-edited page header 266; internal marker 268.
% Jacob Reinhold's CC BY 4.0 Markdown is comparison/drafting material only.

\phantomsection\label{proof:sga3-V-4.1-iv-affine}%
Let us now prove \hyperref[thm:sga3-V-4.1-iv]{(iv)}. By hypothesis,
\[
  X_0=\operatorname{Spec}A_0,\qquad
  X_1=\operatorname{Spec}A_1,
\]
and, for \(i=0,1\), the homomorphism
\(\delta_i:A_0\to A_1\) makes \(A_1\) a finite \(A_0\)-module.
Therefore, a fortiori, the homomorphism
\[
  A_0\otimes_BA_0\longrightarrow A_1
\]
is finite.

We assume in addition that
\[
  d=d_0\boxtimes d_1:X_1\longrightarrow X_0\times_YX_0
\]
is a monomorphism. By Lemma~\hyperref[lem:sga3-V-4.1.3]{4.1.3}, the
corresponding homomorphism \(A_0\otimes_BA_0\to A_1\) is surjective.
We shall prove that it is an isomorphism; in the course of doing so,
we shall also prove that \(p:X_0\to Y\) is finite locally free.

It is enough to prove that, for every prime ideal \(\mathfrak p\) of
\(B\), the homomorphism
\[
  (A_0)_{\mathfrak p}\otimes_{B_{\mathfrak p}}(A_0)_{\mathfrak p}
  \longrightarrow(A_1)_{\mathfrak p},
\]
with components \(\delta_{0\mathfrak p}\) and
\(\delta_{1\mathfrak p}\), is bijective. Thus we may suppose that
\(B\) is local. \hyperref[proof:sga3-V-4.1-b]{Part~b)} then shows that \((A_0)_{\mathfrak p}\) is
semilocal. Indeed, if \(\mathfrak m\) is one of its maximal ideals,
the others have the form \(\delta_0^{-1}(\mathfrak n)\), where
\(\mathfrak n\) ranges over the prime ideals of \(A_1\) satisfying
\(\delta_1^{-1}(\mathfrak n)=\mathfrak m\). There are at most
\[
  n=[A_1:A_0]
\]
such prime ideals \(\mathfrak n\).

After a faithfully flat base change, we may also suppose that the
residue field of \(B\) is infinite.%
\footnote{\label{note:sga3-V-25}Editor's note: see Lemma~\hyperref[lem:sga3-V-4.1.2]{4.1.2}.}
We can then use \hyperref[lem:sga3-V-4.2]{the following lemma}.

\medskip
\phantomsection\label{lem:sga3-V-4.2}%
\noindent\textbf{Lemma 4.2.}---
\textit{Let \(B\) be a local ring with infinite residue field, let
\(A\) be a semilocal ring, and let \(i:B\to A\) be a homomorphism
which sends the maximal ideal \(\mathfrak n\) of \(B\) into the
radical \(\mathfrak r\) of \(A\). Let \(M\) be a free \(A\)-module of
rank \(n\), and let \(N\) be a \(B\)-submodule of \(M\) which
generates \(M\) as an \(A\)-module. Then \(N\) contains an
\(A\)-basis of \(M\).}

Recall that a sequence \(m_1,\ldots,m_n\) of elements of \(M\) is an
\(A\)-basis if and only if its canonical images in \(M/\mathfrak rM\)
form a basis over \(A/\mathfrak r\). We may therefore replace
\[
  M\text{ by }M/\mathfrak rM,\quad
  N\text{ by }N/(N\cap\mathfrak rM),\quad
  A\text{ by }A/\mathfrak r,\quad
  B\text{ by }B/\mathfrak n.
\]
In this case \hyperref[lem:sga3-V-4.2]{the lemma} is elementary. Namely, if
\[
  A=K_1\times\cdots\times K_r,
  \qquad
  M=K_1^n\times\cdots\times K_r^n,
\]
choose, for each \(j\), an element \(x_j\in N\) whose \(j\)-th
component is nonzero. Because the field \(B\) is infinite, a suitable
\(B\)-linear combination \(x\) of the \(x_j\) has all its components
nonzero. Replace \(M\) by \(M/Ax\), and proceed by induction on \(n\).

Apply \hyperref[lem:sga3-V-4.2]{the lemma} with \(A=A_0\), with \(i:B\hookrightarrow A_0\), with
\(M=A_1\) regarded as an \(A_0\)-module through \(\delta_1\), and with
\(N=\delta_0(A_0)\). Indeed, since
\[
  d_0\boxtimes d_1:X_1\longrightarrow X_0\times_YX_0
\]
is a closed immersion, the homomorphism
\[
  A_0\otimes_BA_0\longrightarrow A_1
\]
with components \(\delta_0,\delta_1\) is surjective. This says exactly
that \(\delta_0(A_0)\) generates the \(A_0\)-module \(A_1\).

Choose \(a_1,\ldots,a_n\in A_0\) so that
\(\delta_0(a_1),\ldots,\delta_0(a_n)\) form an \(A_0\)-basis of
\(A_1\). If \(a_1,\ldots,a_n\) form a \(B\)-basis of \(A_0\), then
\(A_0\otimes_BA_0\to A_1\) sends the basis
\[
  (1\otimes a_i)_{1\leq i\leq n}
\]
to the basis
\[
  (\delta_0(a_i))_{1\leq i\leq n},
\]
and is therefore bijective. Thus, if
\[
  \varepsilon:\mathbb Z^n\longrightarrow A_0
\]
is the homomorphism of abelian groups which sends the standard basis
of \(\mathbb Z^n\) to \(a_1,\ldots,a_n\), it remains only to prove
that the map
\[
  B\otimes_{\mathbb Z}\mathbb Z^n\longrightarrow A_0
\]
with components \(i\) and \(\varepsilon\) is bijective.

% END INCLUDED SGA3_English_current_buildable_source_20260823/source/inputs/V/tex/components/08_expose_V_theorem41_proof_iv_lemma42_en.tex
%
              \IfFileExists{tex/components/09_expose_V_theorem41_close_section5_opening_en.tex}{%
                % BEGIN INCLUDED SGA3_English_current_buildable_source_20260823/source/inputs/V/tex/components/09_expose_V_theorem41_close_section5_opening_en.tex
% SGA 3, Expose V, Loop 1 English reconstruction.
% Scope: final diagram and conclusion of the affine constant-rank proof of
% Theorem 4.1, followed by Section 5(a) and the opening sentence of 5(b).
% Authority: Polo--Gille Exp5-13oct24.pdf, component-local p. 17
% (combined-reader p. 279), re-edited page header 267; internal markers
% 269--270. Jacob Reinhold's Markdown is comparison/drafting material only.

The dual diagram \hyperref[fig:sga3-exp5-dual-012star]{\((0,1,2)^\ast\)} considered at the beginning of this
proof induces the following commutative diagram:

\SGAThreeDiagram[0.78\textwidth]
  {figures/exp5/p267_proof_iv_ring_modules_diagram.png}
  {The commutative ring-and-module diagram concluding the affine proof
  of \hyperref[thm:sga3-V-4.1-iv]{Theorem~4.1(iv)}.}
  {fig:sga3-exp5-p267-proof-iv-ring-modules}
  {Polo--Gille Exp5-13oct24.pdf, component-local p. 17;
  combined-reader p. 279; re-edited page header p. 267}

Here \(u_0,u_1,u_2\) have, respectively, the component pairs
\[
  (i,\varepsilon),\qquad
  (\delta_1,\delta_0\varepsilon),\qquad
  (\delta'_2,\delta'_0\delta_0\varepsilon).
\]
We know that \(u_1\) is an isomorphism. Since the two left-hand
squares of \((0,1,2)^\ast\) are cocartesian, \(u_2\) is bijective.
The two horizontal rows of the diagram are exact, and therefore
\(u_0\) is bijective.%
\footnote{\label{note:sga3-V-26}Editor's note: the argument which follows has been added.}
Thus \(A_0\) is a locally free \(B\)-module of rank \(n\). By the
preceding reductions, it follows that
\[
  \delta_0\otimes\delta_1:
  A_0\otimes_BA_0\xrightarrow{\ \sim\ }A_1
\]
is an isomorphism. This completes the proof of \hyperref[thm:sga3-V-4.1]{Theorem~4.1} in the
special case under consideration: \(X_0\) is affine and \(d_1\) is
locally free of constant rank \(n\).

\section*{5. Passage to the quotient by a finite flat groupoid
  (the general case)}
\phantomsection\label{sec:sga3-V-5}
\addcontentsline{toc}{section}{5. Passage to the quotient by a finite
  flat groupoid (the general case)}

\phantomsection\label{par:sga3-V-5-a}%
\noindent\textbf{a)}
Let \(U^{(n)}\) be the largest open subset of \(X_0\) over which
\(d_1\) is finite locally free of rank \(n\). The scheme \(X_0\) is
the disjoint union of the \(U^{(n)}\). The two Cartesian squares

\begin{figure}[H]
  \centering
  \resizebox{0.29\textwidth}{!}{%
    \SGAThreeNativeInput{figures/exp5/p267_section5_rank_square_d0.png}}
  \qquad\textit{and}\qquad
  \resizebox{0.29\textwidth}{!}{%
    \SGAThreeNativeInput{figures/exp5/p267_section5_rank_square_d1.png}}
  \caption{The two Cartesian squares controlling the rank strata
    \(U^{(n)}\).}
  \label{fig:sga3-exp5-p267-section5-rank-squares}
  \par\smallskip
  {\footnotesize\itshape Native Loop-2 reconstructions; source locator:
  Polo--Gille Exp5-13oct24.pdf, component-local p. 17;
  combined-reader p. 279; re-edited page header p. 267.}
\end{figure}

\noindent show that the inverse images of \(U^{(n)}\) under \(d_0\)
and \(d_1\) both coincide with the largest open subset of \(X_1\)
over which \(d'_2\) is locally free of rank \(n\).%
\footnote{\label{note:sga3-V-27}Editor's note: since \(d_0\) (respectively \(d_1\)) is
surjective, flat, and finite, hence faithfully flat and affine,
\(d'_2\) has rank \(n\) over a neighborhood of a point \(x\in X_1\)
if and only if \(d_1\) has rank \(n\) over a neighborhood of
\(d_0(x)\) (respectively \(d_1(x)\)).}
Thus
\[
  d_0^{-1}(U^{(n)})=d_1^{-1}(U^{(n)}),
\]
and the groupoid \(X_\ast\) is the disjoint union of the groupoids
\(X_\ast^{(n)}\) induced by \(X_\ast\) on the open-and-closed
subschemes \(U^{(n)}\). It is therefore enough to prove
\hyperref[thm:sga3-V-4.1]{Theorem~4.1} for each \(X_\ast^{(n)}\). We are reduced to the case in
which \(d_1\) is finite locally free of rank \(n\).

\medskip
\phantomsection\label{par:sga3-V-5-b}%
\noindent\textbf{b)}
\textit{We can now prove \hyperref[thm:sga3-V-4.1]{the theorem} in the general case.}

% END INCLUDED SGA3_English_current_buildable_source_20260823/source/inputs/V/tex/components/09_expose_V_theorem41_close_section5_opening_en.tex
%
                \IfFileExists{tex/components/10_expose_V_section5_general_case_en.tex}{%
                  % BEGIN INCLUDED SGA3_English_current_buildable_source_20260823/source/inputs/V/tex/components/10_expose_V_section5_general_case_en.tex
% SGA 3, Expose V, Loop 1 English reconstruction.
% Scope: Section 5(b), general-case construction through the first paragraph
% on the effective equivalence relation R; ends with local p. 18.
% Authority: Polo--Gille Exp5-13oct24.pdf, component-local p. 18
% (combined-reader p. 280), re-edited page header 268; internal marker 271.
% Jacob Reinhold's CC BY 4.0 Markdown is comparison/drafting material only.

\phantomsection\label{proof:sga3-V-5-b-general}%
By \hyperref[par:sga3-V-5-a]{part~a)}, we may suppose that \(d_1\) is locally free of rank \(n\).
Let \((Y,p)\) be a cokernel of \((d_0,d_1)\) in the category of all
ringed spaces. The argument at the end of \hyperref[proof:sga3-V-4.1-c]{4(c)} shows that, in order to
prove \hyperref[thm:sga3-V-4.1-i]{Theorem~4.1(i)}, it is enough to show that \(Y\) is a scheme and
that \(p:X_0\to Y\) is a morphism of schemes. By
Lemma~\hyperref[sga3:exp5:lemma:1.2]{1.2}, this question is local on \(Y\).
Let \(y\in Y\) and choose \(x\in X_0\) with \(p(x)=y\). If \(x\) has a
saturated affine open neighborhood \(U\), then \hyperref[sec:sga3-V-4]{Section~4} shows that
\(p(U)\) is an affine open subset of \(Y\), and
\(p|_U:U\to p(U)\) is a morphism of schemes. Thus it remains to prove
that every \(x\in X_0\) has a saturated affine open neighborhood. We
proceed as follows; the argument is taken from SGA~1, VIII,
Corollary~7.6.

\SGAThreeDiagram[0.90\textwidth]
  {figures/exp5/p268_section5_affine_saturated_chain.png}
  {The nested open subsets used to construct a saturated affine
  neighborhood of \(x\).}
  {fig:sga3-exp5-p268-affine-saturated-chain}
  {Polo--Gille Exp5-13oct24.pdf, component-local p. 18;
  combined-reader p. 280; re-edited page header p. 268}

By \hyperref[thm:sga3-V-4.1-b]{condition 4.1(b)}, there is an affine open subset \(V\subseteq X_0\)
which contains \(d_1(d_0^{-1}(x))\).%
\footnote{\label{note:sga3-V-28}Editor's note:
\(d_1(d_0^{-1}(x))=d_0(d_1^{-1}(x))\); see \hyperref[note:sga3-V-16]{editor's note~16} in
\hyperref[thm:sga3-V-4.1]{Theorem~4.1}.}
Put \(F=X_0\setminus V\). Since \(d_1\) is integral,
\(d_1(d_0^{-1}(F))\) is closed, and
\[
  V'=X_0\setminus d_1(d_0^{-1}(F))
\]
is the largest saturated open subset contained in \(V\). The set
\(V'\) is a neighborhood of the finite set
\(d_1(d_0^{-1}(x))\). Hence there is a section \(f\) of the structure
sheaf of \(V\), vanishing on \(V\setminus V'\), such that
\(d_1(d_0^{-1}(x))\) is contained in the principal open subset
\(V_f\subseteq V\). We shall prove that the largest saturated open
subset \((V_f)'\) contained in \(V_f\) is affine; it will then be the
required neighborhood.

Put \(Z(f)=V'\setminus V_f\). Then \(d_0^{-1}(Z(f))\) is the set of
points of
\[
  d_0^{-1}(V')=d_1^{-1}(V')
\]
at which \(d_0^\ast(f)\) vanishes. The morphism induced by \(d_1\)
from this common inverse image to \(V'\) is locally free of rank
\(n\).%
\footnote{\label{note:sga3-V-29}Editor's note: the reference to
Lemma~\hyperref[lem:sga3-V-4.1.1]{4.1.1} has been added; see \hyperref[bib:sga3-V-DG70]{[DG70]},
III, \S5.2, p.~313.}
Lemma~\hyperref[lem:sga3-V-4.1.1]{4.1.1} therefore shows that
\(d_1(d_0^{-1}(Z(f)))\) is precisely the zero locus of the norm \(N\)
of \(d_0^\ast(f)\) relative to \(d_1\). It follows that
\[
  (V_f)'=V'\setminus d_1(d_0^{-1}(Z(f)))
\]
is the set of points of \(V_f\) at which \(N\) does not vanish.
Consequently, \((V_f)'\) is affine.

This proves \hyperref[thm:sga3-V-4.1-i]{Theorem~4.1(i)}; assertions \hyperref[thm:sga3-V-4.1-ii]{(ii)}, \hyperref[thm:sga3-V-4.1-iii]{(iii)}, and the first part
of \hyperref[thm:sga3-V-4.1-iv]{(iv)} are then clear. It remains to prove the consequences stated
at the end of \hyperref[thm:sga3-V-4.1-iv]{(iv)}; see \hyperref[bib:sga3-V-Ray67a]{[Ray67a]}, Theorem~1(iii).

By hypothesis, the groupoid \(X_\ast\) comes from an equivalence
relation
\[
  i:R\longrightarrow X_0\times X_0,
\]
where \(i\) is an immersion (see \hyperref[note:sga3-V-19]{editor's note~19}). We have shown that
\(R\) is effective (see \SGARefLink{sga3:iv:definition:3:3:2}{Expos\'e IV, 3.3.2}), and that
\[
  p:X_0\longrightarrow Y=X_0/R
\]
is surjective and finite locally free; in particular, it is faithfully
flat and of finite presentation.

% END INCLUDED SGA3_English_current_buildable_source_20260823/source/inputs/V/tex/components/10_expose_V_section5_general_case_en.tex
%
                  \IfFileExists{tex/components/11_expose_V_section5_fppf_remark51_en.tex}{%
                    % BEGIN INCLUDED SGA3_English_current_buildable_source_20260823/source/inputs/V/tex/components/11_expose_V_section5_fppf_remark51_en.tex
% SGA 3, Expose V, Loop 1 English reconstruction.
% Scope: conclusion of Theorem 4.1's fppf/base-change assertions and all of
% Remark 5.1 on preservation of quasi-projectivity.
% Authority: Polo--Gille Exp5-13oct24.pdf, component-local p. 19
% (combined-reader p. 281), re-edited page header 269.
% Jacob Reinhold's CC BY 4.0 Markdown is comparison/drafting material only.

Consequently, if \((M)\) denotes the family of faithfully flat
morphisms locally of finite presentation, then \(R\) is
\((M)\)-effective. By \SGARefLink{sga3:iv:corollary:6:3:3}{Expos\'e IV, 6.3.3}, \((Y,p)\) therefore
represents the quotient sheaf of \(X_0\) by \(R\) for the fppf
topology, and the base-change assertions follow from \SGARefLink{sga3:iv:proposition:3:4:3:1}{Expos\'e IV,
3.4.3.1}.

\medskip
\phantomsection\label{rem:sga3-V-5.1}%
\noindent\textbf{Remark 5.1.}---%
\footnote{\label{note:sga3-V-30}Editor's note: this paragraph has been added.}
Keep the hypotheses and notation of \hyperref[thm:sga3-V-4.1]{Theorem~4.1}, and suppose in
addition that \(S\) is locally Noetherian and that
\(\pi_0:X_0\to S\) is quasi-projective. We shall prove that
\(\pi:Y\to S\) is quasi-projective.

The preceding hypotheses imply that \(Y\to S\) is of finite type; see
the proof of \hyperref[lem:sga3-V-6.1-ii]{6.1(ii)}. Let \(\mathscr A\) be an invertible
\(\mathcal O_{X_0}\)-module which is ample for \(\pi_0\). By EGA~II,
6.1.12, \(p_\ast\mathscr A\) is an invertible
\(p_\ast\mathcal O_{X_0}\)-module. Hence \(Y\) has an affine open
cover \((V_i)_{i\in I}\) such that \(\mathscr A\) is trivial over each
of the saturated affine open subsets
\[
  U_i=p^{-1}(V_i).
\]

For each \(i\), put
\[
  A_{i,0}=\mathcal O_{X_0}(U_i),
\]
and let \(A_{i,1}\) be the ring of the affine open subset
\(d_0^{-1}(U_i)=d_1^{-1}(U_i)\) of \(X_1\). Let
\(\delta_{i,0}\) (respectively \(\delta_{i,1}\)) be the homomorphism
\(A_{i,0}\to A_{i,1}\) induced by \(d_0\) (respectively \(d_1\)), and
put
\[
  B_i=\mathcal O_Y(V_i)
  =\{b\in A_{i,0}\mid\delta_{i,0}(b)=\delta_{i,1}(b)\}.
\]

Following EGA~II, \S6.5, consider the invertible
\(\mathcal O_{X_0}\)-module
\[
  N_{d_1}\!\left(d_0^\ast\mathscr A\right),
\]
the norm, relative to the finite locally free morphism
\(d_1:X_1\to X_0\), of the invertible
\(\mathcal O_{X_1}\)-module \(d_0^\ast\mathscr A\). Suppose that
\(\mathscr A\), relative to the cover \((U_i)_{i\in I}\), is given by
transition functions
\[
  c_{ij}\in\mathcal O_{X_0}(U_i\cap U_j)^\times.
\]
Then \(N_{d_1}(d_0^\ast\mathscr A)\) is given by the transition
functions
\[
  N_{\delta_1}(\delta_0(c_{ij}))
  \in\mathcal O_{X_0}(U_i\cap U_j)^\times.
\]
By \hyperref[proof:sga3-V-4.1-a]{4(a)}, these elements belong to
\(\mathcal O_Y(V_i\cap V_j)^\times\), and hence define an invertible
\(\mathcal O_Y\)-module \(\mathscr L\) such that
\[
  p^\ast\mathscr L=N_{d_1}(d_0^\ast\mathscr A).
\]
Moreover, for every \(n\in\mathbb N^\ast\),
\[
  p^\ast(\mathscr L^n)
  =N_{d_1}\!\left(d_0^\ast(\mathscr A^n)\right);
\]
see \emph{loc. cit.}, (6.5.2.1).

We now show that \(\mathscr L\) is ample for \(\pi:Y\to S\).
Replacing \(S\) by an affine open subset, we may suppose that \(S\)
is affine. Let \(y\in Y\), choose \(x\in X_0\) with \(p(x)=y\), let
\(V\) be an affine open neighborhood of \(y\) in \(Y\), and put
\(U=p^{-1}(V)\). Since \(\mathscr A\) is \(\pi_0\)-ample, there exist
\(n\in\mathbb N^\ast\) and a section
\[
  s\in\Gamma(X_0,\mathscr A^n)
\]
such that
\[
  x\in(X_0)_s\subseteq U.
\]
With the notation above, \(s\) is represented by sections
\[
  a_i\in A_{i,0}=\mathcal O_{X_0}(U_i),
  \qquad a_i=c_{ij}a_j\ \text{ on }\ U_i\cap U_j,
\]
and \((X_0)_s\) is the union of the open subsets
\[
  U'_i=\{\mathfrak p\in\operatorname{Spec}(A_{i,0})
    \mid a_i\notin\mathfrak p\}.
\]

For every \(i\), put
\[
  N(a_i)=N_{\delta_1}(\delta_0(a_i))\in B_i.
\]
\hyperref[thm:sga3-V-4.1-i]{Theorem~4.1(i)} and
Lemma~\hyperref[lem:sga3-V-4.1.1]{4.1.1} give
\[
  p(U'_i)
  =p\,d_1d_0^{-1}(U'_i)
  =p\,d_1\!\left(
    \{\mathfrak q\in\operatorname{Spec}(A_{i,1})
      \mid\delta_{i,0}(a_i)\notin\mathfrak q\}
    \right),
\]
while
\[
  d_1\!\left(
    \{\mathfrak q\in\operatorname{Spec}(A_{i,1})
      \mid\delta_{i,0}(a_i)\notin\mathfrak q\}
    \right)
  =
  \{\mathfrak p\in\operatorname{Spec}(A_{i,0})
    \mid N_{\delta_1}(\delta_{i,0}(a_i))\notin\mathfrak p\}.
\]
It follows that
\[
  p(U'_i)
  =
  \{\mathfrak p\in\operatorname{Spec}(B_i)
    \mid N(a_i)\notin\mathfrak p\}.
\]
Consequently,
\[
  p((X_0)_s)=Y_{N(s)},
\]
where \(N(s)\) is the section of \(\mathscr L^n\) on \(Y\) defined by
the sections \(N(a_i)\in\mathcal O_Y(V_i)\). Thus
\begin{equation}
  y\in p((X_0)_s)=Y_{N(s)}\subseteq p(U)=V.
  \tag{*}\label{eq:sga3-V-5.1-star}
\end{equation}
This proves that \(\mathscr L\) is ample for \(\pi:Y\to S\), and
hence completes the proof that \(\pi\) is quasi-projective.

% END INCLUDED SGA3_English_current_buildable_source_20260823/source/inputs/V/tex/components/11_expose_V_section5_fppf_remark51_en.tex
%
                    \IfFileExists{tex/components/12_expose_V_section6_quasisection_setup_en.tex}{%
                      % BEGIN INCLUDED SGA3_English_current_buildable_source_20260823/source/inputs/V/tex/components/12_expose_V_section6_quasisection_setup_en.tex
% SGA 3, Expose V, Loop 1 English reconstruction.
% Scope: Section 6 opening and quasi-section reduction through the groupoid V_*.
% Authority: Polo--Gille Exp5-13oct24.pdf, component-local p. 20
% (combined-reader p. 282), re-edited page header 270.
% Jacob Reinhold's CC BY 4.0 Markdown is comparison/drafting material only.

\section*{6. Passage to the quotient when a quasi-section exists}
\phantomsection\label{sec:sga3-V-6}
\addcontentsline{toc}{section}{6. Passage to the quotient when a quasi-section exists}

We shall now prove a technical lemma that will be used in the proof of
the two theorems we have in view. Let \(S\) be a scheme and let
\SGAThreeDiagram[0.54\textwidth]
  {figures/exp5/p270_section6_groupoid_X.png}
  {The groupoid \(X_\ast\).}
  {fig:sga3-V-6-X-groupoid}
  {Expos\'e V, component-local p.~20 / combined-reader p.~282 / re-edited header p.~270}
be a \((\mathbf{Sch}/S)\)-groupoid. We call a \emph{quasi-section} of
the groupoid \(X_\ast\) any subscheme \(U\) of \(X_0\) for which the
following two conditions hold:
\begin{enumerate}
\item[(1)]\phantomsection\label{def:sga3-V-6-quasisection-1}%
The restriction \(v\) of \(d_1\) to \(d_0^{-1}(U)\) is a
  finite, locally free, surjective morphism from \(d_0^{-1}(U)\) onto
  \(X_0\).
\item[(2)]\phantomsection\label{def:sga3-V-6-quasisection-2}%
Every subset \(E\) of \(U\) consisting of points that are
  pairwise equivalent for the equivalence relation defined by
  \(X_\ast\) (\hyperref[par:sga3-V-3-e]{\S3.e)}) is contained in an affine open
  subset of \(U\).\footnote{\label{note:sga3-V-31}Editor's note: If \(x,y\in E\), there is a
  \(z\in X_1\) such that \(d_1(z)=x\) and \(d_0(z)=y\); thus \(z\)
  belongs to \(v^{-1}(x)\), which is finite by \hyperref[def:sga3-V-6-quasisection-1]{(1)}. Hence \(E\) is
  contained in the finite set
  \(d_0(v^{-1}(x))=d_0d_1^{-1}(x)\cap U\). The source PDF prints
  \(d_1(x)=x\) in the first equality; the type-correct reading
  \(d_1(z)=x\) is used here.}
\end{enumerate}

If \(U\) is a quasi-section of \(X_\ast\), the
\((\mathbf{Sch}/S)\)-groupoid
\SGAThreeDiagram[0.48\textwidth]
  {figures/exp5/p270_section6_groupoid_U.png}
  {The groupoid \(U_\ast\) induced by \(X_\ast\).}
  {fig:sga3-V-6-U-groupoid}
  {Expos\'e V, component-local p.~20 / combined-reader p.~282 / re-edited header p.~270}
induced by \(X_\ast\) and the inclusion \(U\hookrightarrow X_0\)
(see \hyperref[par:sga3-V-3-a]{\S3.a)}) satisfies the hypotheses of
Theorem~\hyperref[thm:sga3-V-4.1]{4.1}. Indeed, put
\[
  V=d_0^{-1}(U),
\]
and let \(u\) and \(v\) be the morphisms with source \(V\) induced by
\(d_0\) and \(d_1\), respectively:
\SGAThreeDiagram[0.43\textwidth]
  {figures/exp5/p270_section6_span_X0_V_U.png}
  {The morphisms \(v:V\to X_0\) and \(u:V\to U\).}
  {fig:sga3-V-6-span-X0-V-U}
  {Expos\'e V, component-local p.~20 / combined-reader p.~282 / re-edited header p.~270}

By \hyperref[par:sga3-V-3-b]{\S3.b)}, there is a Cartesian square
\SGAThreeDiagram[0.27\textwidth]
  {figures/exp5/p270_section6_cartesian_square.png}
  {The Cartesian square defining \(U_1\).}
  {fig:sga3-V-6-U1-cartesian}
  {Expos\'e V, component-local p.~20 / combined-reader p.~282 / re-edited header p.~270}
so \(u_1\) is surjective and finite locally free by \hyperref[def:sga3-V-6-quasisection-1]{condition~(1)}.
Together with \hyperref[def:sga3-V-6-quasisection-2]{condition~(2)}, \hyperref[def:sga3-V-6-quasisection-1]{condition~(1)} therefore ensures that the
groupoid \(U_\ast\) satisfies the hypotheses of
Theorem~\hyperref[thm:sga3-V-4.1]{4.1}. In particular,
\(\operatorname{Coker}(u_0,u_1)\) exists in \((\mathbf{Sch}/S)\).
Moreover, \(d_0\) has a section, so \(u\) is a universal effective
epimorphism (see Expos\'e IV, 1.12).\footnote{The French re-edition
mistakenly cites Expos\'e III, which has no \S1.12; see Canon record
\texttt{SGA-SRC-000856}.}
Proposition~\hyperref[prop:sga3-V-3.1]{3.1}
then shows that \(\operatorname{Coker}(u_0,u_1)\) coincides with the
cokernel \(\operatorname{Coker}(v_0,v_1)\) of the groupoid \(V_\ast\):
\SGAThreeDiagram[0.48\textwidth]
  {figures/exp5/p270_section6_groupoid_V.png}
  {The groupoid \(V_\ast\).}
  {fig:sga3-V-6-V-groupoid}
  {Expos\'e V, component-local p.~20 / combined-reader p.~282 / re-edited header p.~270}

% END INCLUDED SGA3_English_current_buildable_source_20260823/source/inputs/V/tex/components/12_expose_V_section6_quasisection_setup_en.tex
%
                      \IfFileExists{tex/components/13_expose_V_lemma61_statement_proof_a_opening_en.tex}{%
                        % BEGIN INCLUDED SGA3_English_current_buildable_source_20260823/source/inputs/V/tex/components/13_expose_V_lemma61_statement_proof_a_opening_en.tex
% SGA 3, Expose V, Loop 1 English reconstruction.
% Scope: Section 6 reduction, Lemma 6.1, and the opening of proof (a).
% Authority: Polo--Gille Exp5-13oct24.pdf, component-local p. 21
% (combined-reader p. 283), re-edited page header 271.
% Jacob Reinhold's CC BY 4.0 Markdown is comparison/drafting material only.

This is the inverse image of \(U_\ast\) under the base change
\(u:V\to U\), and hence also the inverse image of \(X_\ast\) under
the base change
\SGAThreeDiagram[0.43\textwidth]
  {figures/exp5/p271_section6_base_change_d0.png}
  {The base change defining \(V_\ast\).}
  {fig:sga3-V-6-base-change-d0}
  {Expos\'e V, component-local p.~21 / combined-reader p.~283 / re-edited header p.~271}

By \hyperref[par:sga3-V-3-c]{\S3.c)}, \(V_\ast\) is isomorphic to the groupoid
\(V'_\ast\), the inverse image of \(X_\ast\) under the base change
\SGAThreeDiagram[0.49\textwidth]
  {figures/exp5/p271_section6_base_change_d1.png}
  {The base change defining \(V'_\ast\).}
  {fig:sga3-V-6-base-change-d1}
  {Expos\'e V, component-local p.~21 / combined-reader p.~283 / re-edited header p.~271}
Thus \(V'_\ast\) admits a cokernel in \((\mathbf{Sch}/S)\). Since
\(v:V\to X_0\) is flat, surjective, and finite, it is faithfully flat
and quasi-compact, and therefore a universal effective epimorphism by
Expos\'e IV, 6.3.2.\footnote{The French re-edition mistakenly cites
Expos\'e III, which ends at \S4; see Canon record
\texttt{SGA-SRC-000857}.}
Proposition~\hyperref[prop:sga3-V-3.1]{3.1} consequently
shows that \(X_\ast\) admits a cokernel
\(\operatorname{Coker}(d_0,d_1)\) in \((\mathbf{Sch}/S)\). We have
therefore proved the first assertion of \hyperref[lem:sga3-V-6.1-i]{part~(i)} of \hyperref[lem:sga3-V-6.1]{the following
lemma}.\footnote{\label{note:sga3-V-32}Editor's note: The argument that follows has been
slightly modified. In particular, the additional hypothesis that
\(d_0\) be flat has been moved to \hyperref[lem:sga3-V-6.1-iv]{part~(iv)}; the former part~(ii) has
been split into \hyperref[lem:sga3-V-6.1-ip]{parts~(i')} and~\hyperref[lem:sga3-V-6.1-ii]{(ii)}; and the second assertion of
\hyperref[lem:sga3-V-6.1-ip]{part~(i')} has been added.}

\medskip
\phantomsection\label{lem:sga3-V-6.1}%
\noindent\textbf{Lemma 6.1.}---\emph{Suppose that the
\((\mathbf{Sch}/S)\)-groupoid \(X_\ast\) possesses a quasi-section.
Then:}
\begin{enumerate}
\item[\SGAInlineRefTarget{lem:sga3-V-6.1:item:i}(i)]\phantomsection\label{lem:sga3-V-6.1-i}%
\emph{There exists a cokernel \((Y,p)\) of
  \((d_0,d_1)\) in \((\mathbf{Sch}/S)\). Moreover, every such
  \((Y,p)\) is a cokernel of \((d_0,d_1)\) in the category of all
  ringed spaces.}
\item[(i')]\phantomsection\label{lem:sga3-V-6.1-ip}%
\emph{The morphism \(p\) is surjective; it is open
  (respectively, universally closed) if \(d_0\) is.}
\item[(ii)]\phantomsection\label{lem:sga3-V-6.1-ii}%
\emph{Suppose that \(S\) is locally Noetherian and that
  \(X_0\) is locally of finite type (respectively, of finite type)
  over \(S\). Then \(p\) and \(Y\to S\) are locally of finite
  presentation (respectively, of finite presentation).}
\item[(iii)]\phantomsection\label{lem:sga3-V-6.1-iii}%
\emph{The morphism
  \(X_1\to X_0\times_Y X_0\) with components \(d_0\) and \(d_1\) is
  surjective.}
\item[(iv)]\phantomsection\label{lem:sga3-V-6.1-iv}%
\emph{If \((d_0,d_1)\) is an equivalence pair, then
  \(X_1\to X_0\times_Y X_0\) is an isomorphism. Moreover, if
  \(d_0:X_1\to X_0\) is flat, then \(p\) is faithfully flat.}
\end{enumerate}

Before proving the second assertion of \hyperref[lem:sga3-V-6.1-i]{part~(i)}, we shall prove
\hyperref[lem:sga3-V-6.1-ip]{parts~(i')}, \hyperref[lem:sga3-V-6.1-ii]{(ii)}, and \hyperref[lem:sga3-V-6.1-iii]{(iii)}.

\medskip
\phantomsection\label{par:sga3-V-6-a}%
\noindent\textbf{a) Proof of \hyperref[lem:sga3-V-6.1-ip]{(i')} and \hyperref[lem:sga3-V-6.1-ii]{(ii)}.}---%
We have just seen that \((Y,p)\) identifies both with
\(\operatorname{Coker}(v_0,v_1)\) and with
\(\operatorname{Coker}(u_0,u_1)\). Let \(q\) and \(r\) be the
canonical epimorphisms from \(U\) and \(V\), respectively, to \(Y\):
\SGAThreeDiagram[0.39\textwidth]
  {figures/exp5/p271_lemma61_cokernel_triangle.png}
  {The canonical epimorphisms to \(Y\).}
  {fig:sga3-V-6.1-cokernel-triangle}
  {Expos\'e V, component-local p.~21 / combined-reader p.~283 / re-edited header p.~271}
By hypothesis, \(v\) is surjective and finite locally free, hence
open.

% END INCLUDED SGA3_English_current_buildable_source_20260823/source/inputs/V/tex/components/13_expose_V_lemma61_statement_proof_a_opening_en.tex
%
                        \IfFileExists{tex/components/14_expose_V_lemma61_proofs_a_b_c_opening_en.tex}{%
                          % BEGIN INCLUDED SGA3_English_current_buildable_source_20260823/source/inputs/V/tex/components/14_expose_V_lemma61_proofs_a_b_c_opening_en.tex
% SGA 3, Expose V, Loop 1 English reconstruction.
% Scope: Lemma 6.1 proof (a), proof (b), and opening of proof (c).
% Authority: Polo--Gille Exp5-13oct24.pdf, component-local p. 22
% (combined-reader p. 284), re-edited page header 272.
% Jacob Reinhold's CC BY 4.0 Markdown is comparison/drafting material only.

On the other hand, if \(d_0:X_1\to X_0\) is open (respectively,
universally closed), then \(u\), being obtained from \(d_0\) by base
change, has the same property. By
Theorem~\hyperref[thm:sga3-V-4.1]{4.1}, \(q\) is surjective, integral,
and open. It follows that \(r\) is surjective and is open
(respectively, universally closed) if \(d_0\) is. Since \(v\) is
surjective, the same is true of \(p\). This proves \hyperref[lem:sga3-V-6.1-ip]{part~(i')}.

Now suppose that \(S\) is locally Noetherian and that \(X_0\) is
locally of finite type over \(S\), so that \(X_0\) is locally
Noetherian.

We show that \(Y\) is locally of finite presentation over \(S\). Let
\[
  S'=\operatorname{Spec}R
\]
be an affine open subset of \(S\), let
\[
  Y'=\operatorname{Spec}B
\]
be an affine open subset of \(Y\) mapping into \(S'\), and let
\[
  U'=\operatorname{Spec}A
\]
be the inverse image of \(Y'\) in \(U\). Since \(R\) is Noetherian,
it is enough to show that \(B\) is a finitely generated
\(R\)-algebra. By \hyperref[sec:sga3-V-4]{\S\S4} and \hyperref[sec:sga3-V-5]{5}, \(B\) is contained in \(A\), which is
a finitely generated \(R\)-algebra; the assertion therefore follows
because \(R\) is Noetherian and \(A\) is integral over \(B\).

Finally, since \(X_0\to S\) is locally of finite type, so is \(p\)
(EGA~I, 6.6.6). Hence \(p\) is locally of finite presentation because
\(Y\) is locally Noetherian.

It remains to prove the final assertion of \hyperref[lem:sga3-V-6.1-ii]{part~(ii)}. Suppose in
addition that \(X_0\) is of finite type over \(S\). Since \(p\) is
surjective, \(Y\) is then quasi-compact over \(S\), and hence of
finite type over \(S\). Because \(S\) is locally Noetherian, the
morphisms \(X_0\to S\) and \(Y\to S\) are of finite presentation, and
therefore so is \(p:X_0\to Y\) (EGA~IV\(_1\), 1.6.2(v)).

\medskip
\phantomsection\label{par:sga3-V-6-b}%
\noindent\textbf{b) Proof of \hyperref[lem:sga3-V-6.1-iii]{(iii)}.}---%
The groupoid \(V_\ast\) with base \(V\) is isomorphic both to the
inverse image of \(U_\ast\) under the base change \(u\) and to the
inverse image of \(X_\ast\) under the base change \(v\). We therefore
have a double Cartesian square:
\SGAThreeDiagram[0.65\textwidth]
  {figures/exp5/p272_lemma61_double_cartesian_square.png}
  {The double Cartesian square used in the proof of \hyperref[lem:sga3-V-6.1-iii]{Lemma~6.1(iii)}.}
  {fig:sga3-V-6.1-double-cartesian-square}
  {Expos\'e V, component-local p.~22 / combined-reader p.~284 / re-edited header p.~272}
Since \(u_0\boxtimes u_1\) is surjective, so is
\(v_0\boxtimes v_1\). Since \(v\times v\) is surjective, the
composite morphism
\[
  V_1\longrightarrow X_0\times_Y X_0
\]
is surjective, and hence so is \(d_0\boxtimes d_1\).

\medskip
\phantomsection\label{par:sga3-V-6-c}%
\noindent\textbf{c) Proof of \hyperref[lem:sga3-V-6.1-i]{(i)}.}---%
It remains to prove that \((Y,p)\) is a cokernel of \((d_0,d_1)\) in
the category of all ringed spaces. We first show that \(Y\) is
obtained from \(X_0\) by identifying those points \(x\) and \(y\) for
which there exists \(z\in X_1\) with
\[
  d_0(z)=x,\qquad d_1(z)=y.
\]
Indeed, \(p\) is surjective and \(p d_0=p d_1\). Conversely, if
\(p(x)=p(y)\), there is a point \(z'\in X_0\times_Y X_0\) whose first
projection is \(x\) and whose second projection is \(y\). If
\(z\in X_1\) satisfies
\[
  (d_0\boxtimes d_1)(z)=z',
\]
then \(d_0(z)=x\) and \(d_1(z)=y\), as required.

On the other hand, if \(W\) is a saturated open subset of \(X_0\),
then \(W\cap U\) is a saturated open subset of \(U\). By
Theorem~\hyperref[thm:sga3-V-4.1]{4.1}, \(q(W\cap U)\) is open in
\(Y\). Since \(q(W\cap U)=p(W)\), the topology on \(Y\) is the
quotient topology induced by \(X_0\).

It remains to show that the canonical sequence of sheaves of rings
\begin{equation}
  \mathcal O_Y
  \longrightarrow p_\ast\mathcal O_{X_0}
  \rightrightarrows
  p_\ast d_{0\ast}\mathcal O_{X_1}
  =
  p_\ast d_{1\ast}\mathcal O_{X_1}
  \label{eq:sga3-V-6.1-sheaf-exactness}
\end{equation}
is exact.

% END INCLUDED SGA3_English_current_buildable_source_20260823/source/inputs/V/tex/components/14_expose_V_lemma61_proofs_a_b_c_opening_en.tex
%
                          \IfFileExists{tex/components/15_expose_V_lemma61_proofs_c_d_en.tex}{%
                            % BEGIN INCLUDED SGA3_English_current_buildable_source_20260823/source/inputs/V/tex/components/15_expose_V_lemma61_proofs_c_d_en.tex
% SGA 3, Expose V, Loop 1 English reconstruction.
% Scope: completion of Lemma 6.1, proofs (c) and (d).
% Authority: Polo--Gille Exp5-13oct24.pdf, component-local p. 23
% (combined-reader p. 285), re-edited page header 273.
% Jacob Reinhold's CC BY 4.0 Markdown is comparison/drafting material only.

Let \(Y'\) be an open subset of \(Y\), and put
\[
  U'=q^{-1}(Y'),\qquad X'_0=p^{-1}(Y'),
\]
and so forth.\footnote{\label{note:sga3-V-33}Editor's note: The argument that follows has
been expanded, in particular the verification that \(U'\) is a
quasi-section of the groupoid induced on \(X'_0\).}
Then \(U'\) is an open subset of \(U\) saturated for the equivalence
relation defined by the groupoid \(U_\ast\). It follows from
\hyperref[lem:sga3-V-1.1]{Lemmas~1.1} and
\hyperref[sga3:exp5:lemma:1.2]{1.2} that \(Y'\) is the cokernel, both
in \((\mathbf{Sch}/S)\) and in the category of ringed spaces, of the
groupoid induced by \(U_\ast\) on \(U'\). Likewise, \(X'_0\) is an
open subset of \(X_0\) saturated for the equivalence relation defined
by \(X_\ast\), and we have the following commutative diagram, in which
both squares are Cartesian:
\SGAThreeDiagram[0.58\textwidth]
  {figures/exp5/p273_lemma61_restricted_groupoids_cartesian_corrected.png}
  {The restricted groupoids over \(X'_0\) and \(X_0\).}
  {fig:sga3-V-6.1-restricted-groupoids}
  {Expos\'e V, component-local p.~23 / combined-reader p.~285 / re-edited header p.~273; lower-right \(U\) corrected from source-PDF \(U'\)}
Here \(\widetilde d_1\) is surjective and finite locally free. On the
other hand, let \(x\in U'\). Since \(U\) is a quasi-section, the set
\[
  E=d_0d_1^{-1}(x)\cap U
\]
is finite and is contained in an affine open subset \(W\) of \(U\).
Then \(E'=E\cap U'\) is a finite set contained in the quasi-affine
open subset \(W\cap U'\). Consequently there is an affine open subset
\(W'\) of \(W\cap U'\) that contains \(E'\). This proves that \(U'\)
is a quasi-section of the groupoid \(X'_\ast\) induced by \(X_\ast\)
on \(X'_0\). The first assertion of \hyperref[lem:sga3-V-6.1-i]{part~(i)}, applied to \(X'_\ast\)
and \(U'\), now shows that \(Y'\) is the cokernel of \(X'_\ast\) in
\((\mathbf{Sch}/S)\).

In particular, for every \(S\)-scheme \(T\), the sequence
\begin{equation}
  T(Y')
  \xrightarrow{\,T(p|_{X'_0})\,}
  T(X'_0)
  \underset{T(d_0|_{X'_1})}
    {\overset{T(d_1|_{X'_1})}{\rightrightarrows}}
  T(X'_1)
  \label{eq:sga3-V-6.1-T-exact}
\end{equation}
is exact. If \(T\) is the affine line \(\mathbf G_{a,S}\)
(\SGARefLink{sga3:i:subsection:4-3}{Expos\'e I, 4.3}), this sequence identifies with
\begin{equation}
\begin{aligned}
  \Gamma(Y',\mathcal O_Y)
  &\longrightarrow
  \Gamma(p^{-1}(Y'),\mathcal O_{X_0})
  \underset{\delta_0}{\overset{\delta_1}{\rightrightarrows}}
  \Gamma(d_0^{-1}p^{-1}(Y'),\mathcal O_{X_1})\\
  &\hspace{9em}=
  \Gamma(d_1^{-1}p^{-1}(Y'),\mathcal O_{X_1}),
\end{aligned}
\label{eq:sga3-V-6.1-sections-exact}
\end{equation}
which is therefore exact for every open subset \(Y'\). This completes
the proof of Lemma~\hyperref[lem:sga3-V-6.1-i]{6.1(i)}.

\medskip
\phantomsection\label{par:sga3-V-6-d}%
\noindent\textbf{d) Proof of \hyperref[lem:sga3-V-6.1-iv]{(iv)}.}---%
If \((d_0,d_1)\) is an equivalence pair, then so is
\((u_0,u_1)\). It follows from
Theorem~\hyperref[thm:sga3-V-4.1]{4.1} that
\[
  u_0\boxtimes u_1:U_1\longrightarrow U\times_Y U
\]
is an isomorphism, and hence so is \(v_0\boxtimes v_1\); see the
Cartesian squares in part~\hyperref[par:sga3-V-6-b]{(b)}. Since
\(v\times v\) is faithfully flat and quasi-compact,
\(d_0\boxtimes d_1\) is an isomorphism (SGA~1, VIII, \SGARefLink{sga3:viii:corollary:5:4}{5.4}).

Moreover, if \(d_0\) is flat, then \(u\) is flat. The morphism \(q\)
is flat by Theorem~\hyperref[thm:sga3-V-4.1]{4.1}, so \(r\) is flat
as well. Since \(v\) is faithfully flat, \(p\) is flat; being also
surjective, it is faithfully flat.

% END INCLUDED SGA3_English_current_buildable_source_20260823/source/inputs/V/tex/components/15_expose_V_lemma61_proofs_c_d_en.tex
%
                            \IfFileExists{tex/components/16_expose_V_section7_theorem71_en.tex}{%
                              % BEGIN INCLUDED SGA3_English_current_buildable_source_20260823/source/inputs/V/tex/components/16_expose_V_section7_theorem71_en.tex
% SGA 3, Expose V, Loop 1 English reconstruction.
% Scope: Section 7 and Theorem 7.1 through the reduction to local
% quasi-sections, including the opening seam on component-local p. 25.
% Authority: Polo--Gille Exp5-13oct24.pdf, component-local pp. 24--25
% (combined-reader pp. 286--287), re-edited headers 274--275.
% Jacob Reinhold's CC BY 4.0 Markdown is comparison/drafting material only.

\section*{7. Quotient by a proper and flat groupoid}
\phantomsection\label{sec:sga3-V-7}
\addcontentsline{toc}{section}{7. Quotient by a proper and flat groupoid}

\medskip
\phantomsection\label{thm:sga3-V-7.1}%
\noindent\textbf{Theorem 7.1.}---%
\footnote{\label{note:sga3-V-34}Editor's note: We mention here the
article of S.~Keel and S.~Mori \hyperref[bib:sga3-V-KM97]{[KM97]}, which establishes the following
theorem. Let \(X\) be an algebraic space of finite type over a locally
Noetherian base \(S\), and let
\(j:R\to X\times_S X\) be a flat groupoid whose stabilizer
\(j^{-1}(\Delta_X)\) is finite over \(X\). There then exists an
algebraic space that is a geometric quotient of \(X\) by \(R\) and a
uniform categorical quotient; moreover, if \(j\) is separated, this
quotient is separated. In particular, if a flat \(S\)-group scheme
\(G\) acts properly on \(X\), with finite stabilizer (that is, if
\[
  G\times_S X\longrightarrow X\times_S X,\qquad
  (g,x)\longmapsto(x,g\cdot x)
\]
is proper and the stabilizer of the diagonal is finite over \(X\)),
then there exists a geometric quotient \(X\to X/G\). For a reductive
\(S\)-group scheme \(G\), this is a result of J.~Koll\'ar \hyperref[bib:sga3-V-Ko97]{[Ko97]}.}%
\emph{Let \(S\) be a locally Noetherian scheme and let}
\SGAThreeDiagram[0.56\textwidth]
  {figures/exp5/p274_theorem71_groupoid.png}
  {The groupoid occurring in \hyperref[thm:sga3-V-7.1]{Theorem~7.1}.}
  {fig:sga3-V-7.1-groupoid}
  {Expos\'e V, component-local p.~24 / combined-reader p.~286 / re-edited header p.~274}
\emph{be a \((\mathbf{Sch}/S)\)-groupoid such that \(d_1\) is proper
and flat, \(X_0\) is quasi-projective over \(S\),}%
\footnote{\label{note:sga3-V-35}Editor's note: This hypothesis on
\(X_0\) is necessary; see \hyperref[note:sga3-V-17]{editor's note~(17)} to
Theorem~\hyperref[thm:sga3-V-4.1]{4.1}.}
\emph{and the morphism
\[
  d:X_1\longrightarrow X_0\times_S X_0
\]
with components \(d_0\) and \(d_1\) is quasi-finite. Then:}
\begin{enumerate}
\item[(i)]\phantomsection\label{thm:sga3-V-7.1-i}%
  \emph{There exists a cokernel \((Y,p)\) of \((d_0,d_1)\) in
  \((\mathbf{Sch}/S)\). Moreover, every such \((Y,p)\) is a cokernel
  of \((d_0,d_1)\) in the category of all ringed spaces.}
\item[(ii)]\phantomsection\label{thm:sga3-V-7.1-ii}%
  \emph{The morphism \(p\) is surjective, open, proper, and of finite
  presentation, and \(Y\to S\) is separated and of finite
  presentation.}%
  \footnote{\label{note:sga3-V-36}Editor's note: TDTE~III,
  \SGARefLink{lem:sga3-V-6.1}{Theorem~6.1}, states that \(Y\to S\) is quasi-projective if \(S\) is
  Noetherian. The editors did not see how to deduce this from the
  local existence of quasi-sections.}
\item[(iii)]\phantomsection\label{thm:sga3-V-7.1-iii}%
  \emph{The morphism
  \[
    X_1\longrightarrow X_0\times_Y X_0
  \]
  with components \(d_0\) and \(d_1\) is surjective.}
\item[(iv)]\phantomsection\label{thm:sga3-V-7.1-iv}%
  \emph{If \((d_0,d_1)\) is an equivalence pair, then
  \(X_1\to X_0\times_Y X_0\) is an isomorphism and \(p\) is
  faithfully flat.}%
  \footnote{\label{note:sga3-V-37}Editor's note: The consequences
  that follow have been made explicit; see \hyperref[bib:sga3-V-Ray67a]{[Ray67a]},
  Theorem~1(iii), and the end of the proof of \hyperref[thm:sga3-V-7.1]{the theorem}. We also
  mention that another proof of \hyperref[thm:sga3-V-7.1]{Theorem~7.1} in the case of an
  equivalence relation, based on the existence of Hilbert schemes,
  is given in \hyperref[bib:sga3-V-AK80]{[AK80]}, Theorem~2.9; see also \hyperref[bib:sga3-V-BLR90]{[BLR90]}, \S8.2,
  Theorem~12. In that case it is furthermore shown that \(Y\to S\)
  is quasi-projective.}
  \emph{Moreover, \((Y,p)\) is a cokernel of \((d_0,d_1)\) in the
  category of sheaves for the fppf topology, and, for every base
  change \(Y'\to Y\), the scheme \(Y'\) is the cokernel of the
  groupoid \(X_\ast\times_Y Y'\) obtained from \(X_\ast\) by the
  base change \(X_0\times_Y Y'\to X_0\).}
\end{enumerate}

\emph{In particular, for every base change \(S'\to S\), the scheme
\[
  Y'=Y\times_S S'
\]
is the cokernel of the \(S'\)-groupoid
\[
  X'_\ast=X_\ast\times_S S'.
\]
Thus, in this case, \textquotedblleft{}formation of the quotient
commutes with base change\textquotedblright.}

Let \((Y,p)\) be the cokernel of \((d_0,d_1)\) in the category of all
ringed spaces. Lemma~\hyperref[sga3:exp5:lemma:1.2]{1.2} shows that,
to prove part~\hyperref[thm:sga3-V-7.1-i]{(i)}, it is enough to show
that every point \(z\in X_0\) has a saturated open neighborhood
\(U_z\) such that, if \(\widetilde d_0\) and
\(\widetilde d_1\) denote the restrictions of \(d_0\) and \(d_1\) to
\[
  d_0^{-1}(U_z)=d_1^{-1}(U_z),
\]
and if \((Q,q)\) denotes the cokernel of
\((\widetilde d_0,\widetilde d_1)\) in the category of ringed spaces,
then \(Q\) is a scheme and \(q\) is a morphism of schemes.

By Lemma~\hyperref[lem:sga3-V-6.1-i]{6.1(i)}, it is therefore enough
to show that every point \(z\in X_0\) has a saturated open
neighborhood \(U_z\) such that the groupoid induced by \(X_\ast\) on
\(U_z\) has a quasi-section. We may even assume that \(z\) is closed
in its fiber over \(S\) (we shall say that \(z\) is \emph{closed
relative to \(S\)}).%
\footnote{\label{note:sga3-V-38}Editor's note: Indeed, if such an open
neighborhood \(U_z\) has been constructed for every point \(z\) that
is closed relative to \(S\), then the union of these \(U_z\) covers
\(X_0\): every fiber over \(S\) of the closed complement is a
Noetherian scheme with no closed points, and is therefore empty.}
The existence of \(U_z\) follows from the lemmas below.
% SGA3-V-V2-EOL-PRESERVATION

% END INCLUDED SGA3_English_current_buildable_source_20260823/source/inputs/V/tex/components/16_expose_V_section7_theorem71_en.tex
%
                              \IfFileExists{tex/components/17_expose_V_lemma72_en.tex}{%
                                % BEGIN INCLUDED SGA3_English_current_buildable_source_20260823/source/inputs/V/tex/components/17_expose_V_lemma72_en.tex
% SGA 3, Expose V, Loop 1 English reconstruction.
% Scope: Lemma 7.2, including its closing seam on component-local p. 26.
% Authority: Polo--Gille Exp5-13oct24.pdf, component-local pp. 25--26
% (combined-reader pp. 287--288), re-edited headers 275--276.
% Jacob Reinhold's CC BY 4.0 Markdown is comparison/drafting material only.

\medskip
\phantomsection\label{lem:sga3-V-7.2}%
\noindent\textbf{Lemma 7.2.}---%
\emph{Let \(T\) be an affine Noetherian scheme, let \(X,Y,Z\) be
\(T\)-schemes of finite type, with \(X\) quasi-projective over \(T\),
and let}
\SGAThreeDiagram[0.30\textwidth]
  {figures/exp5/p275_lemma72_square.png}
  {The diagram in \hyperref[lem:sga3-V-7.2]{Lemma~7.2}.}
  {fig:sga3-V-7.2-square}
  {Expos\'e V, component-local p.~25 / combined-reader p.~287 / re-edited header p.~275}
\emph{be a diagram in \((\mathbf{Sch}/T)\). Let \(z\) be a point of
\(v(Y)\) that is closed relative to \(T\), and suppose that \(v\) is
flat at the points of \(v^{-1}(z)\). Then there exists a closed
subscheme \(F\) of \(X\) such that
\[
  u\bigl(u^{-1}(F)\cap v^{-1}(z)\bigr)
\]
is finite and nonempty, and such that the restriction of \(v\) to
\(u^{-1}(F)\) is flat at the points of \(v^{-1}(z)\).}

Put \(T=\operatorname{Spec}A\). We may suppose that
\[
  X=\operatorname{Proj}\mathscr S,
\]
where \(\mathscr S\) is the symmetric algebra of a finitely generated
\(A\)-module \(E\).

If \(u(v^{-1}(z))\) is finite, we may take \(F=X\). Otherwise, let
\[
  y_1,\ldots,y_n
\]
be the points of the fiber \(v^{-1}(z)\) associated with its
structure sheaf \(\mathcal O_{v^{-1}(z)}\). Thus, if
\(\mathcal O_i\) denotes the local ring of \(v^{-1}(z)\) at \(y_i\),
the maximal ideal of \(\mathcal O_i\) consists of zero divisors. Let
\(t\) be the image of \(z\) in \(T\). The set \(u(v^{-1}(z))\) is an
infinite constructible subset of the fiber of \(X\) over \(t\).
Consequently, that fiber contains a closed point
\(x\in u(v^{-1}(z))\) distinct from
\(u(y_1),\ldots,u(y_n)\). The open subset
\(X-\{x\}\) is a neighborhood of \(u(y_1),\ldots,u(y_n)\), and hence
contains an open neighborhood of the form \(D_+(f)\), where
\(f\in\mathscr S\) is homogeneous of degree \(d\) (we use the
notation of EGA~II, \S2.3).

The closed subscheme
\[
  X_1=V_+(f)
\]
defined by \(f\) therefore contains \(x\) and avoids the points
\(u(y_1),\ldots,u(y_n)\). Its inverse image
\[
  Y_1=u^{-1}\bigl(V_+(f)\bigr)
\]
is consequently distinct from \(Y\) and meets \(v^{-1}(z)\). We
shall in fact prove that the restriction \(v_1\) of \(v\) to \(Y_1\)
is flat at the points of \(v^{-1}(z)\). If
\(u(v_1^{-1}(z))\) is finite, we may then take \(F=X_1\). Otherwise,
we repeat the argument with \(Y_1,v_1\), and the morphism \(u_1\)
induced by \(u\) on \(Y_1\), in place of \(Y,v,u\). In this way we
obtain a descending sequence
\[
  X\supset X_1\supset\cdots
\]
of closed subschemes of \(X\). Since such a sequence terminates,
\[
  u\bigl(u^{-1}(X_n)\cap v^{-1}(z)\bigr)
\]
is finite and nonempty for some \(n\), and we then take \(F=X_n\).

It remains to show that \(v_1\) is flat at the points of
\(v^{-1}(z)\). Let \(y\in Y_1\) lie over \(z\), let
\(\mathcal O_y\) be the local ring of \(Y\) at \(y\), let
\(\overline{\mathcal O}_y\) be the local ring of \(v^{-1}(z)\) at
\(y\), and let \(\mathcal O_{v(y)}\) be the local ring of \(Z\) at
\(v(y)\). Choose \(g\in\mathscr S_1\) such that \(D_+(g)\) is a
neighborhood of \(u(y)\) in \(X\), and let \(\varphi\) and
\(\overline\varphi\) be the respective images of \(f/g^d\) in
\(\mathcal O_y\) and \(\overline{\mathcal O}_y\). By the construction
of \(f\), the element \(\overline\varphi\) is not a zero divisor in
\(\overline{\mathcal O}_y\). Since \(\mathcal O_y\) is flat over
\(\mathcal O_z=\mathcal O_{v(y)}\), the element \(\varphi\) is not a zero divisor in
\(\mathcal O_y\), and
\[
  \mathcal O_y/\mathcal O_y\varphi
\]
is flat over \(\mathcal O_z\) (SGA~1, IV, 5.7). This quotient is
precisely the local ring of \(Y_1\) at \(y\).
% SGA3-V-V2-EOL-PRESERVATION

% END INCLUDED SGA3_English_current_buildable_source_20260823/source/inputs/V/tex/components/17_expose_V_lemma72_en.tex
%
                                \IfFileExists{tex/components/18_expose_V_lemma73_opening_en.tex}{%
                                  % BEGIN INCLUDED SGA3_English_current_buildable_source_20260823/source/inputs/V/tex/components/18_expose_V_lemma73_opening_en.tex
% SGA 3, Expose V, Loop 1 English reconstruction.
% Scope: Lemma 7.3 through the finiteness argument, including its closing
% seam on component-local p. 27.
% Authority: Polo--Gille Exp5-13oct24.pdf, component-local pp. 26--27
% (combined-reader pp. 288--289), re-edited headers 276--277.
% Jacob Reinhold's CC BY 4.0 Markdown is comparison/drafting material only.

\medskip
\phantomsection\label{lem:sga3-V-7.3}%
\noindent\textbf{Lemma 7.3.}---%
\emph{Retain the notation and hypotheses of
Theorem~\hyperref[thm:sga3-V-7.1]{7.1}. Every point \(z\in X_0\)
that is closed relative to \(S\) has a saturated open neighborhood
\(U_z\) such that the groupoid induced by \(X_\ast\) on \(U_z\) has
a quasi-section.}

The assertion is local on \(S\), so we may suppose that \(S\) is
affine and Noetherian and apply \hyperref[lem:sga3-V-7.2]{the preceding lemma} to the diagram
\SGAThreeDiagram[0.30\textwidth]
  {figures/exp5/p276_lemma73_base_square.png}
  {The diagram to which \hyperref[lem:sga3-V-7.2]{Lemma~7.2} is applied.}
  {fig:sga3-V-7.3-base-square}
  {Expos\'e V, component-local p.~26 / combined-reader p.~288 / re-edited header p.~276}
in \((\mathbf{Sch}/S)\). Thus let \(F\) be a closed subscheme of
\(X_0\) such that
\[
  d_0\bigl(d_0^{-1}(F)\cap d_1^{-1}(z)\bigr)
\]
is finite and nonempty, and such that the restriction of \(d_1\) to
\(d_0^{-1}(F)\) is flat at the points of \(d_1^{-1}(z)\).

Let \(F_1\) and \(F_2\) be the inverse images of \(F\) under \(d_0\)
and \(d_0d'_0=d_0d'_1\), respectively, and let
\(\widetilde d_0,\widetilde d_1,\ldots\) denote the morphisms induced
by \(d_0,d_1,\ldots\). We then have a commutative diagram
\SGAThreeDiagram[0.67\textwidth]
  {figures/exp5/p276_lemma73_F2_diagram.png}
  {The induced diagram on \(F_2,F_1,F\).}
  {fig:sga3-V-7.3-F-diagram}
  {Expos\'e V, component-local p.~26 / combined-reader p.~288 / re-edited header p.~276}
in which the two left-hand squares are Cartesian, the first row is
exact (cf. the diagram
\hyperref[sga3:v:lemma:1-2:equation:0-1-2]{\((0,1,2)\)} in \hyperref[sec:sga3-V-1]{\S1}), and
\(q,\widetilde q\) are the structure morphisms.

We first show that only finitely many points of \(F_1\) lie over
\(z\).%
\footnote{\label{note:sga3-V-39}Editor's note: Details have been added,
and the role of the properness hypothesis on \(d_0\) and \(d_1\) in
Theorem~\hyperref[thm:sga3-V-7.1]{7.1} has been made explicit.
Compare the statement and proof of \hyperref[thm:sga3-V-8.1]{Theorem~8.1}, where this properness
hypothesis is omitted.}
Let \(s\) be the image of \(z\) in \(S\). Since \(F\) is of finite
type over \(S\), the fiber \(\widetilde q^{-1}(s)\) is a Noetherian
scheme. Since \(\widetilde d_0\) is proper,
\[
  \widetilde d_0\bigl(\widetilde d_1^{-1}(z)\bigr)
\]
is a closed subscheme of \(\widetilde q^{-1}(s)\) consisting of
finitely many points. Consequently (cf. EGA~I, 6.2.2), those points
are closed in \(\widetilde q^{-1}(s)\), and also in the fiber
\(q^{-1}(s)\) of \(X_0\), because \(F\) is closed in \(X_0\).

Let \(y\) be one of these points. Since \(q^{-1}(s)\) is of finite
type over \(\kappa(s)\), it has affine open neighborhoods
\[
  \operatorname{Spec}B\quad\text{and}\quad\operatorname{Spec}C
\]
of \(y\) and \(z\), respectively, with \(B\) and \(C\) finitely
generated \(\kappa(s)\)-algebras. The points \(y\) and \(z\)
correspond to maximal ideals
\[
  \mathfrak p\subset B,\qquad \mathfrak q\subset C.
\]
The fields \(B/\mathfrak p\) and \(C/\mathfrak q\) are finite over
\(\kappa(s)\), so
\[
  (B/\mathfrak p)\otimes_{\kappa(s)}(C/\mathfrak q)
\]
is a finite-dimensional \(\kappa(s)\)-algebra. Its maximal ideals
correspond precisely to the points of \(X_0\times_S X_0\) whose
second projection is \(z\) and whose first projection is \(y\).
Thus there are only finitely many points \(u\in X_0\times_S X_0\)
whose second projection is \(z\) and whose first projection belongs
to
\[
  \widetilde d_0\bigl(\widetilde d_1^{-1}(z)\bigr).
\]
Finally, since \(X_1\to X_0\times_S X_0\) has finite fibers, each
such point \(u\) is the image of only finitely many points of \(X_1\).
This proves the assertion.
% SGA3-V-V2-EOL-PRESERVATION

% END INCLUDED SGA3_English_current_buildable_source_20260823/source/inputs/V/tex/components/18_expose_V_lemma73_opening_en.tex
%
                                  \IfFileExists{tex/components/19_expose_V_lemma73_close_lemma74_opening_en.tex}{%
                                    % BEGIN INCLUDED SGA3_English_current_buildable_source_20260823/source/inputs/V/tex/components/19_expose_V_lemma73_close_lemma74_opening_en.tex
% SGA 3, Expose V, Loop 1 English reconstruction.
% Scope: completion of Lemma 7.3 and the page-27 opening of Lemma 7.4.
% Authority: Polo--Gille Exp5-13oct24.pdf, component-local p. 27
% (combined-reader p. 289), re-edited header 277.
% Jacob Reinhold's CC BY 4.0 Markdown is comparison/drafting material only.

The morphism \(\widetilde d_1\) is therefore quasi-finite and flat at
the points of \(F_1\) lying over \(z\). Since \(\widetilde d_1\) is
of finite type, SGA~1, IV, 6.10 and EGA~III, 4.4.10 imply that the set
\(\Phi\) of points of \(F_1\) at which \(\widetilde d_1\) is not both
flat and quasi-finite is closed in \(F_1\), and hence in \(X_1\)
because \(F_1\) is closed in \(X_1\).%
\footnote{\label{note:sga3-V-40}Editor's note: If \(f:X\to Y\) is
locally of finite type, the set of points \(x\in X\) that are isolated
in their fiber \(f^{-1}(f(x))\) is open in \(X\). In EGA~III, 4.4.10,
this is deduced from Zariski's Main Theorem when \(Y\) is locally
Noetherian. For arbitrary \(Y\), it follows from Chevalley's
semicontinuity theorem (EGA~IV\(_3\), 13.1.3 and 13.1.4).
Consequently, \(f\) is quasi-finite at \(x\) if and only if it is of
finite type at \(x\) and \(x\) is isolated in
\(f^{-1}(f(x))\). This will be used below; see \hyperref[note:sga3-V-42]{editor's note~(42)}.}
Since \(d_1\) is proper, \(\widetilde d_1(\Phi)\) is closed and, by
what precedes, does not contain \(z\). Put
\[
  W=\widetilde d_1(F_1)-\widetilde d_1(\Phi).
\]
The restriction of \(\widetilde d_1\) to
\(\widetilde d_1^{-1}(W)\) is of finite presentation (by the
Noetherian hypotheses), flat, proper, and quasi-finite.%
\footnote{\label{note:sga3-V-41}Editor's note: The original has been
expanded in what follows.}
It is therefore finite, locally free, and open, by EGA~III, 4.4.2 and
EGA~IV\(_2\), 2.1.12 and 2.4.6. Consequently,
\(\widetilde d_1(F_1)\) is a neighborhood of \(z\), and \(W\) is the
largest open subset of \(X_0\) contained in \(\widetilde d_1(F_1)\)
over which \(\widetilde d_1\) is both flat and quasi-finite.

We shall see in Lemma~\hyperref[lem:sga3-V-7.4]{7.4} that the inverse
images of \(\Phi\) under \(\widetilde d'_1\) and
\(\widetilde d'_0\) both identify with the set of points of \(F_2\)
at which \(\widetilde d'_2\) is not both flat and quasi-finite.
It follows that
\[
\begin{aligned}
  d_0^{-1}(W)
  &=
  \widetilde d'_2(F_2)
  -
  \widetilde d'_2\bigl((\widetilde d'_0)^{-1}(\Phi)\bigr),\\
  d_1^{-1}(W)
  &=
  \widetilde d'_2(F_2)
  -
  \widetilde d'_2\bigl((\widetilde d'_1)^{-1}(\Phi)\bigr)
\end{aligned}
\]
coincide. Thus \(W\) is saturated.

Set
\[
  W_1=\widetilde d_1^{-1}(W).
\]
The equality \(d_0^{-1}(W)=d_1^{-1}(W)\) implies
\[
  (\widetilde d'_2)^{-1}\bigl(d_0^{-1}(W)\bigr)
  =
  (\widetilde d'_2)^{-1}\bigl(d_1^{-1}(W)\bigr),
\]
that is,
\[
  (\widetilde d'_0)^{-1}(W_1)
  =
  (\widetilde d'_1)^{-1}(W_1).
\]
The morphism \(\widetilde d_0\) is faithfully flat and quasi-compact,
because \(d_0\), like \(d_1\), is surjective, proper, and flat.
Moreover, the square
\SGAThreeDiagram[0.30\textwidth]
  {figures/exp5/p277_lemma73_descent_square.png}
  {The Cartesian square used to descend \(W_1\).}
  {fig:sga3-V-7.3-descent-square}
  {Expos\'e V, component-local p.~27 / combined-reader p.~289 / re-edited header p.~277}
is Cartesian. It follows that
\[
  W_1=\widetilde d_0^{-1}(U)
\]
for an open subset \(U\) of \(F\) (SGA~1, VIII, \SGARefLink{sga3:viii:corollary:4:4}{4.4}). This open
subset \(U\) is a quasi-section of the groupoid with base \(W\)
induced by \(X_\ast\). We may therefore take \(U_z=W\).

It remains to state \hyperref[lem:sga3-V-7.4]{the following lemma}, whose proof is classical.

\medskip
\phantomsection\label{lem:sga3-V-7.4}%
\noindent\textbf{Lemma 7.4.}---%
\emph{Consider a Cartesian square of schemes}
\SGAThreeDiagram[0.30\textwidth]
  {figures/exp5/p277_lemma74_cartesian_square.png}
  {The Cartesian square in \hyperref[lem:sga3-V-7.4]{Lemma~7.4}.}
  {fig:sga3-V-7.4-cartesian-square}
  {Expos\'e V, component-local p.~27 / combined-reader p.~289 / re-edited header p.~277}
% The statement continues on component-local p. 28.
% SGA3-V-V2-EOL-PRESERVATION

% END INCLUDED SGA3_English_current_buildable_source_20260823/source/inputs/V/tex/components/19_expose_V_lemma73_close_lemma74_opening_en.tex
%
                                    \IfFileExists{tex/components/20_expose_V_localp28_31_sections7_9_en.tex}{%
                                      % BEGIN INCLUDED SGA3_English_current_buildable_source_20260823/source/inputs/V/tex/components/20_expose_V_localp28_31_sections7_9_en.tex
% SGA 3, Expose V, Loop 1 English reconstruction -- DISJOINT FUTURE STAGING.
% Scope: component-local pp. 28--31 (combined-reader pp. 290--293),
% re-edited headers 278--281. This begins in the continuation of Lemma 7.4,
% closes Section 7, contains Section 8, and opens Section 9 through 9(b).
% Authority: Polo--Gille Exp5-13oct24.pdf, SHA-256
% 9198200633F929FE1822520371A9200DA4F8CF2513EFAB3D6A0C7E9330DB84CF.
% Jacob Reinhold's CC BY 4.0 Markdown is comparison/drafting material only.
% Integration order: only after the separately owned local pp. 24--27 tranche.

and let \(x\) be a point of \(F_2\).

\begin{enumerate}
\item[\SGAInlineRefTarget{lem:sga3-V-7.4:item:i}(i)] If \(u\) is flat, \(d'\) is flat at \(x\) if and only if
  \(d\) is flat at \(v(x)\).
\item[(ii)] If \(d\) is locally of finite type, \(d'\) is quasi-finite
  at \(x\) if and only if \(d\) is quasi-finite at \(v(x)\).%
  \footnote{\label{note:sga3-V-42}Editor's note: The conditions are
  sufficient by base change (EGA~II, 6.2.4(iii), and
  EGA~IV\(_2\), 2.1.4). Conversely, put
  \(y=d'(x)\) and \(z=u(y)=d(v(x))\), and suppose that \(d'\) is flat
  at \(x\) and that \(u\) (hence also \(v\)) is flat. Then
  \(\mathcal O_{v(x)}\to\mathcal O_x\) is faithfully flat, as is
  \(\mathcal O_z\to\mathcal O_y\to\mathcal O_x\). Consequently,
  \(\mathcal O_z\to\mathcal O_{v(x)}\) is faithfully flat
  (EGA~IV\(_2\), 2.2.11(iv)). Finally, suppose that \(d\) is locally
  of finite type and that \(d'\) is quasi-finite at \(x\). Then
  \(v(x)\) is isolated in its fiber \(d^{-1}(z)\), since \(x\) is
  isolated in its fiber
  \(d'^{-1}(y)=d^{-1}(z)\otimes_{\kappa(z)}\kappa(y)\). By Chevalley's
  semicontinuity theorem there is therefore an open neighborhood of
  \(v(x)\) every point of which is isolated in its fiber
  (EGA~IV\(_3\), 13.1.3 and 13.1.4), and hence \(d\) is quasi-finite
  at \(v(x)\).}
\end{enumerate}

We have thus proved that \(X_0\) has a covering by saturated open
subsets \(W\) such that the groupoid \(W_\ast\) induced by \(X_\ast\)
on \(W\) possesses a quasi-section.%
\footnote{\label{note:sga3-V-43}Editor's note: The sequel has been
modified to take advantage of the additions made to
Lemma~\hyperref[lem:sga3-V-6.1]{6.1}.}

By Lemma~\hyperref[lem:sga3-V-6.1]{6.1} and the reductions stated
after Theorem~\hyperref[thm:sga3-V-7.1]{7.1}, this proves assertions
\hyperref[thm:sga3-V-7.1-i]{(i)} and
\hyperref[thm:sga3-V-7.1-iii]{(iii)} of that theorem, as well as the
facts that \(p\) is surjective and open and that \(p\) and \(Y\to S\)
are locally of finite presentation. Moreover, because \(X_0\to S\)
is quasi-projective, and hence separated and of finite type, \(p\) is
separated; the proof of Lemma~\hyperref[lem:sga3-V-6.1-ii]{6.1(ii)}
shows that \(p\) and \(Y\to S\) are of finite presentation.

To prove that \(p\) is proper, it remains to prove that it is
universally closed. Since this assertion is local on \(Y\), we may
work over a saturated open subset \(W\) such that the groupoid
\(W_\ast\) induced by \(X_\ast\) on \(W\) possesses a quasi-section
\(U\) (because \(X_0\) is covered by such open subsets). With the
notation of \hyperref[par:sga3-V-6-a]{6(a)}, there is a commutative
diagram
\SGAThreeDiagram[0.48\textwidth]
  {figures/exp5/exp5_localp28_WVU_Z_commutative.png}
  {The properness diagram over the open subset \(Z\subseteq Y\).}
  {fig:sga3-V-7.1-properness}
  {Expos\'e V, component-local p.~28 / combined-reader p.~290 / re-edited header p.~278}
where \(Z\) is an open subset of \(Y\), all the arrows are
surjective, and \(q\) is integral. By hypothesis
\(d_0:X_1\to X_0\) is proper, so \(u\), which is obtained from \(d_0\)
by base change, is proper as well. Consequently \(r\) is universally
closed, and hence so is \(p\), since \(v\) is surjective.

Finally, because \(p\) is surjective and universally closed and
\(X_0\) is quasi-projective, hence separated, the diagonal
\(\Delta_{Y/S}(Y)\) is closed in \(Y\times_S Y\): it is the image under
\(p\times p\) of the diagonal \(\Delta_{X_0/S}(X_0)\). Thus \(Y\) is
separated over \(S\). This completes the proof of
Theorem~\hyperref[thm:sga3-V-7.1-ii]{7.1(ii)}.

The assertions in Theorem~\hyperref[thm:sga3-V-7.1-iv]{7.1(iv)} are
local on \(Y\). Since the saturated open subsets \(U_z\) cover
\(X_0\), it is enough to verify them after replacing \(X\) and \(Y\)
by \(U_z\) and \(V_z=p(U_z)\), respectively. As observed at the
beginning of the proof of Theorem~\hyperref[thm:sga3-V-7.1]{7.1},
Lemmas~\hyperref[lem:sga3-V-1.1]{1.1},
\hyperref[sga3:exp5:lemma:1.2]{1.2}, and
\hyperref[lem:sga3-V-6.1-i]{6.1(i)} show that
\((V_z,p|_{U_z})\) is the cokernel, both in \((\mathbf{Sch})\) and in
the category of ringed spaces, of the groupoid induced by \(X_\ast\)
on \(U_z\). The hypothesis that \(d=(d_0,d_1)\) is a monomorphism is
preserved under the base change \(U_z\to X_0\). The first two
assertions of Theorem~\hyperref[thm:sga3-V-7.1-iv]{7.1(iv)} therefore
follow from Lemma~\hyperref[lem:sga3-V-6.1-iv]{6.1(iv)}.

Let us finally prove the consequences stated at the end of \hyperref[thm:sga3-V-7.1-iv]{part~(iv)}
(see \(\hyperref[bib:sga3-V-Ray67a]{[\mathrm{Ray67a}]}\), Theorem~1(iii)). By hypothesis, the
groupoid \(X_\ast\) comes from an equivalence relation
\(R\to X_0\times X_0\). We have proved that \(R\) is effective
(\SGARefLink{sga3:iv:definition:3:3:2}{Expos\'e~IV, 3.3.2}) and that
\(p:X_0\to Y=X_0/R\) is faithfully flat and of finite presentation.
Consequently, if \((M)\) denotes the family of faithfully flat
morphisms locally of finite presentation, then \(R\) is
\((M)\)-effective. By
\SGARefLink{sga3:iv:corollary:6:3:3}{Expos\'e~IV, 6.3.3}, the pair \((Y,p)\) therefore
represents the quotient sheaf of \(X_0\) by \(R\) for the fppf
topology, and the assertions concerning base change follow from
\SGARefLink{sga3:iv:proposition:3:4:3:1}{Expos\'e~IV, 3.4.3.1}.

\section*{8. Passage to the quotient by a flat groupoid that is not
  necessarily proper}
\phantomsection\label{sec:sga3-V-8}
\addcontentsline{toc}{section}{8. Passage to the quotient by a flat groupoid that is not necessarily proper}

\medskip
\phantomsection\label{thm:sga3-V-8.1}
\noindent\textbf{Theorem 8.1.}---%
\footnote{\label{note:sga3-V-44}Editor's note: There is a largest open
subset \(W\) of \(X_0\) satisfying the conclusions of \hyperref[thm:sga3-V-8.1]{the theorem}.
Indeed, let \(W\) be an open subset as in \hyperref[thm:sga3-V-8.1]{the theorem} and let
\(W^\sharp\) be a dense saturated open subset of \(W\). Since \(p\)
is open, \(V^\sharp=p(W^\sharp)\) is open in \(V\), and
\(W^\sharp=p^{-1}(V^\sharp)\), because \(W^\sharp\) is saturated. By
Lemma~\hyperref[lem:sga3-V-1.1]{1.1}, \(V^\sharp\) is a cokernel for
the groupoid induced on \(W^\sharp\). Thus two open subsets \(W\) and
\(W'\) satisfying the conclusions of \hyperref[thm:sga3-V-8.1]{the theorem} may be glued along
their intersection \(W^\sharp\), and conditions \hyperref[thm:sga3-V-8.1-i]{(i)}--\hyperref[thm:sga3-V-8.1-iv]{(iv)}, together
with the property that \(p\) and \(V\to S\) are locally of finite
presentation, are preserved. Conclusion (ii') follows, as in the
proof of Lemma~\hyperref[lem:sga3-V-6.1-ii]{6.1(ii)}, from the
hypothesis that \(X_0\) is of finite type over the Noetherian scheme
\(S\).

Lemmas~\hyperref[lem:sga3-V-1.1]{1.1} and
\hyperref[sga3:exp5:lemma:1.2]{1.2} also show that the union of all
saturated open subsets \(U\subseteq X_0\) such that the open subset
\(p(U)\subseteq Y\) is a scheme and
\(p|_U:U\to p(U)\) is a morphism of schemes is the largest saturated
open subset \(\Omega\subseteq X_0\) satisfying \hyperref[thm:sga3-V-8.1-i]{condition~(i)} of
Theorem~\hyperref[thm:sga3-V-8.1]{8.1}. \hyperref[thm:sga3-V-8.1]{The theorem} shows that
\(\Omega\) contains a dense open subset \(W\), but it is not immediate
that \(\Omega\) satisfies \hyperref[thm:sga3-V-8.1-ii]{properties (ii)}--\hyperref[thm:sga3-V-8.1-iv]{(iv)}. For more precise
results, and for the representability of the fppf quotient sheaf
\(\widetilde{X/R}\) (where \(R\) denotes the groupoid \(X_\ast\) with
base \(X=X_0\)), see \(\hyperref[bib:sga3-V-Ray67a]{[\mathrm{Ray67a}]}\), \(\hyperref[bib:sga3-V-Ray67b]{[\mathrm{Ray67b}]}\), and
Appendix~I of \(\hyperref[bib:sga3-V-An73]{[\mathrm{An73}]}\). These works use weaker hypotheses:
\(S\) may be arbitrary, \(X\) is locally of finite type over \(S\),
and \(R\) is an \(S\)-groupoid with base \(X\) such that \(d_0\) (and
hence \(d_1\)) is flat and of finite presentation. In particular, if
\(\widetilde{X/R}\) is represented by an \(S\)-scheme \(Y\), then \(Y\)
is also the cokernel in the category of ringed spaces. The converse
is false in general (see \(\hyperref[bib:sga3-V-Mum65]{[\mathrm{Mum65}]}\), Chapter~0, \S3,
Example~0.4, cited in \(\hyperref[bib:sga3-V-Ray67a]{[\mathrm{Ray67a}]}\), Remark~1), but it is true
when \(d=(d_0,d_1)\) is an immersion. Under this hypothesis,
\(p:\Omega\to Z:=\Omega/R_\Omega\) is faithfully flat and of finite
presentation. If, in addition, \(S\) is locally Noetherian, then a
codimension-one point \(x\) of \(X\) belongs to \(\Omega\) if and only
if the graph of the groupoid induced on
\(\operatorname{Spec}(\mathcal O_{X,x})\) is closed. See
\(\hyperref[bib:sga3-V-Ray67a]{[\mathrm{Ray67a}]}\), Proposition~1; \(\hyperref[bib:sga3-V-Ray67b]{[\mathrm{Ray67b}]}\),
Proposition~1 and Theorems~2, 1, and 4; and \(\hyperref[bib:sga3-V-An73]{[\mathrm{An73}]}\),
Theorems~5 and 6 on pages~66--67 and Proposition~3.3.1 on page~49.
For an action of an algebraic group over an algebraically closed field
\(k\), see also \(\hyperref[bib:sga3-V-DR81]{[\mathrm{DR81}]}\).}%
\textit{Let \(S\) be a Noetherian scheme and}
\SGAThreeDiagram[0.57\textwidth]
  {figures/exp5/exp5_localp29_theorem81_groupoid.png}
  {The \((\mathbf{Sch}/S)\)-groupoid of \hyperref[thm:sga3-V-8.1]{Theorem~8.1}.}
  {fig:sga3-V-8.1-groupoid}
  {Expos\'e V, component-local p.~29 / combined-reader p.~291 / re-edited header p.~279}
\textit{a \((\mathbf{Sch}/S)\)-groupoid such that \(d_1\) is flat and
of finite type, \(X_0\) is of finite type over \(S\), and the morphism
\(X_1\to X_0\times_S X_0\) with components \(d_0\) and \(d_1\) is
quasi-finite.}

\textit{There exists a dense saturated open subset \(W\subseteq X_0\)
having the following properties:}
\begin{enumerate}
\item[(i)]\phantomsection\label{thm:sga3-V-8.1-i}%
  \textit{if
  \(W_2\overset{w'_i}{\rightrightarrows}
    W_1\overset{w_j}{\rightrightarrows}W\)
  is the groupoid induced by \(X_\ast\) on \(W\), then
  \((w_0,w_1)\) has a cokernel \((V,p)\) in
  \((\mathbf{Sch}/S)\); moreover, \((V,p)\) is a cokernel of
  \((w_0,w_1)\) in the category of all ringed spaces;}
\item[(ii)]\phantomsection\label{thm:sga3-V-8.1-ii}%
  \textit{\(p\) is surjective and open;}
\item[(ii')]\phantomsection\label{thm:sga3-V-8.1-iip}%
  \textit{\(p\) and \(V\to S\) are of finite presentation;}
\item[(iii)]\phantomsection\label{thm:sga3-V-8.1-iii}%
  \textit{the morphism \(W_1\to W\times_VW\) with components \(w_0\)
  and \(w_1\) is surjective;}
\item[(iv)]\phantomsection\label{thm:sga3-V-8.1-iv}%
  \textit{if \((d_0,d_1)\) is an equivalence pair, then
  \(W_1\to W\times_VW\) is an isomorphism and \(p\) is faithfully
  flat.}
\end{enumerate}

We shall prove that \(W\) may be chosen so that the
\((\mathbf{Sch}/S)\)-groupoid \(W_\ast\) induced by \(X_\ast\)
possesses a quasi-section (see \hyperref[sec:sga3-V-7]{\S7}).
Theorem~\hyperref[thm:sga3-V-8.1]{8.1} will then follow from
Lemma~\hyperref[lem:sga3-V-6.1]{6.1}.

Suppose provisionally that for every point \(z\in X_0\) closed
relative to \(S\) (see \hyperref[sec:sga3-V-7]{\S7}) there is a
saturated open subset \(W_z\) that possesses a quasi-section and meets
every irreducible component of \(X_0\) passing through \(z\). Then the
complement \(X_0-W_z\) is saturated: its saturation
\(d_1(d_0^{-1}(X_0-W_z))\) is open and does not meet \(W_z\). If this
complement is nonempty, choose in it a point \(z'\) closed relative to
\(S\), and associate to \(z'\) an open subset \(W_{z'}\) as above. We
may further suppose that \(W_{z'}\subseteq X_0-W_z\). Then \(W_z\) and
\(W_{z'}\) are disjoint, and the groupoid induced by \(X_\ast\) on
\(W_z\cup W_{z'}\) possesses a quasi-section. The process must stop,
because \(X_0\) has only finitely many irreducible components. It
therefore remains to construct \(W_z\).

For this purpose we may suppose that \(S\) is affine. Let \(y\in X_1\)
satisfy \(d_1(y)=z\), let \(X\) be an affine open subset of \(X_0\)
containing \(d_0(y)\), let \(Y\) be the inverse image of \(X\) under
\(d_0:X_1\to X_0\), and let \(u:Y\to X\) and \(v:Y\to X_0\) be the
morphisms induced by \(d_0\) and \(d_1\). Since \(X\) is affine, hence
quasi-projective, Lemma~\hyperref[lem:sga3-V-7.2]{7.2} provides a
subscheme \(F\subseteq X_0\) such that
\(d_0^{-1}(F)\cap d_1^{-1}(z)\) is nonempty,
\(d_0(d_0^{-1}(F)\cap d_1^{-1}(z))\) is finite, and the restriction of
\(d_1\) to \(d_0^{-1}(F)\) is flat at the points of \(v^{-1}(z)\).
We can therefore resume the notation of
Lemma~\hyperref[lem:sga3-V-7.3]{7.3}, writing \(F_1\) and \(F_2\) for
the inverse images of \(F\) in \(X_1\) and \(X_2\), and so on:
\SGAThreeDiagram[0.56\textwidth]
  {figures/exp5/exp5_localp30_F_groupoid.png}
  {The induced diagram on \(F_2,F_1,F\) used in the proof of \hyperref[thm:sga3-V-8.1]{Theorem~8.1}.}
  {fig:sga3-V-8.1-F-groupoid}
  {Expos\'e V, component-local p.~30 / combined-reader p.~292 / re-edited header p.~280}

As in \hyperref[lem:sga3-V-7.3]{7.3}, one proves that
\(\widetilde d_1\) is quasi-finite at the points of
\(\widetilde d_1^{-1}(z)\). It is therefore natural to consider the
open subset \(F'_1\subseteq F_1\) consisting of the points at which
\(\widetilde d_1\) is both flat and quasi-finite. By
Lemma~\hyperref[lem:sga3-V-7.4]{7.4}, the inverse images of \(F'_1\)
under \(\widetilde d'_1\) and \(\widetilde d'_0\) both consist of the
points of \(F_2\) at which \(\widetilde d'_2\) is flat and
quasi-finite. Thus these inverse images coincide, and \(F'_1\) is of
the form \(\widetilde d_0^{-1}(F')\) for an open subset \(F'\subseteq F\)
(SGA~1, VIII, \SGARefLink{sga3:viii:corollary:4:4}{4.4}). After replacing \(F\) by \(F'\), we may therefore
suppose that \(\widetilde d_1\) is quasi-finite and flat. Let \(W_z\)
be the largest open subset of \(\widetilde d_1(F_1)\) over which
\(\widetilde d_1\) is finite and flat.

The open subset \(W_z\) need not contain \(z\), but it contains the
generic points of the irreducible components of \(X_0\) passing
through \(z\).%
\footnote{\label{note:sga3-V-45}Editor's note: Indeed, let \(\eta\) be
such a generic point. The hypotheses imply that
\(\mathcal O_{X_0,\eta}\) is an Artinian local ring and that
\(\mathcal O_{F_1,\eta}\) is a finitely generated
\(\mathcal O_{X_0,\eta}\)-module. By SGA~1, VIII, \SGARefLink{sga3:viii:examples:6:5}{6.5}, there is
therefore an open neighborhood of \(\eta\) over which
\(\widetilde d_1\) is finite.}
Since \(d_0\) (and likewise \(d_1\)) is faithfully flat and of finite
presentation, hence open, SGA~1, VIII, \SGARefLink{sga3:viii:corollary:5:7}{5.7} shows that
\(d_0^{-1}(W_z)\) and \(d_1^{-1}(W_z)\) both coincide with the largest
open subset of \(\widetilde d'_2(F_2)\) over which
\(\widetilde d'_2\) is finite and flat. As in
\hyperref[lem:sga3-V-7.3]{7.3}, the inverse images of
\(\widetilde d_1^{-1}(W_z)\) under \(\widetilde d'_0\) and
\(\widetilde d'_1\) therefore coincide. Hence
\(\widetilde d_1^{-1}(W_z)=\widetilde d_0^{-1}(U)\) for an open subset
\(U\subseteq F\), and \(U\) is a quasi-section for the groupoid induced
by \(X_\ast\) on \(W_z\).

\section*{9. Elimination of the Noetherian hypotheses in \hyperref[thm:sga3-V-7.1]{Theorem 7.1}}
\phantomsection\label{sec:sga3-V-9}
\addcontentsline{toc}{section}{9. Elimination of the Noetherian hypotheses in \texorpdfstring{\protect\hyperref[thm:sga3-V-7.1]{Theorem 7.1}}{\SGARefLink{thm:sga3-V-7.1}{Theorem 7.1}}}

\medskip
\phantomsection\label{par:sga3-V-9-a}%
\noindent\textbf{a)}\quad
Resume the notation and hypotheses of
Lemma~\hyperref[lem:sga3-V-6.1]{6.1}, and let
\(\pi:S'\to S\) be an arbitrary base change. If \(f:X\to Y\) is a
morphism of \(S\)-schemes, write \(f':X'\to Y'\) for the morphism of
\(S'\)-schemes obtained from \(f\) by extension of the base along
\(\pi\). With this convention, \(p':X'_0\to Y'\) is surjective, as is
the morphism
\[
  X'_1\longrightarrow X'_0\times_{Y'}X'_0
\]
with components \(d'_0\) and \(d'_1\). Thus the set underlying \(Y'\)
is the quotient of the set underlying \(X'_0\) by the equivalence
relation on \(X'_0\) defined by the \(S'\)-groupoid \(X'_\ast\).
Moreover, \(q':U'\to Y'\) is integral and surjective. The topology of
\(Y'\) is consequently the quotient topology of that of \(U'\), and
hence also of that of \(X'_0\); see the proof in
\hyperref[par:sga3-V-6-c]{\S6(c)}.

On the other hand, \(U'\) is plainly a quasi-section of the
\(S'\)-groupoid \(X'_\ast\), so Lemma~\hyperref[lem:sga3-V-6.1]{6.1}
applies. In particular, \(X'_\ast\) has a cokernel \((Y_1,p_1)\), and
the topological space underlying \(Y_1\) is obtained from that
underlying \(X'_0\) by identifying the points equivalent under the
relation defined by \(X'_\ast\). It follows that the canonical
morphism \(Y_1\to Y'\) is a homeomorphism. In fact, it is a universal
homeomorphism. Indeed, if \(S''\) is an \(S'\)-scheme, let \(Y_2\) be
the cokernel of
\((d_0\times_S S'',d_1\times_S S'')\). The preceding argument,
applied to the base changes \(S''\to S'\) and \(S''\to S\), shows that
\[
  Y_2\longrightarrow Y_1\times_{S'}S''
  \qquad\text{and}\qquad
  Y_2\longrightarrow Y\times_SS''\simeq Y'\times_{S'}S''
\]
are homeomorphisms. Hence
\[
  Y_1\times_{S'}S''\longrightarrow Y'\times_{S'}S''
\]
is a homeomorphism as well.

\medskip
\phantomsection\label{par:sga3-V-9-b}%
\noindent\textbf{b)}\quad
The same remarks plainly apply when the groupoid \(X_\ast\)
\emph{locally} possesses quasi-sections (see the proof of
Theorem~\hyperref[thm:sga3-V-7.1]{7.1}).%
\footnote{\label{note:sga3-V-46}Editor's note: What follows has been
expanded in order to bring out \hyperref[thm:sga3-V-9.0]{Theorem~9.0 below}.}
For example, we have \hyperref[thm:sga3-V-9.0]{the following theorem}:
% SGA3-V-V2-EOL-PRESERVATION

% END INCLUDED SGA3_English_current_buildable_source_20260823/source/inputs/V/tex/components/20_expose_V_localp28_31_sections7_9_en.tex
%
                                      \IfFileExists{tex/components/21_expose_V_localp32_35_sections9_10_en.tex}{%
                                        % BEGIN INCLUDED SGA3_English_current_buildable_source_20260823/source/inputs/V/tex/components/21_expose_V_localp32_35_sections9_10_en.tex
% SGA 3, Expose V, Loop 1 English reconstruction.
% Staging scope only: component-local PDF pp. 32--35 inclusive.
% Authority: Polo--Gille Exp5-13oct24.pdf, SHA-256
% 9198200633F929FE1822520371A9200DA4F8CF2513EFAB3D6A0C7E9330DB84CF.
% Locator mapping:
%   local p. 32 = combined-reader p. 294 = re-edited header p. 282.
%   local p. 33 = combined-reader p. 295 = re-edited header p. 283;
%     original-edition marginal marker 286 occurs after the algebra diagram.
%   local p. 34 = combined-reader p. 296 = re-edited header p. 284.
%   local p. 35 = combined-reader p. 297 = re-edited header p. 285.
% Jacob Reinhold's CC BY 4.0 Markdown is comparison/drafting material only.

\phantomsection\label{thm:sga3-V-9.0}%
\noindent\textbf{Theorem 9.0.}---%
\textit{Let \(S\) be an arbitrary scheme and}
\begin{equation*}
  \resizebox{0.34\textwidth}{!}{%
    \SGAThreeNativeInput{figures/exp5/exp5_localp032_theorem90_groupoid.png}}
  \phantomsection\label{fig:sga3-V-9.0-groupoid}
\end{equation*}
\textit{a \((\mathbf{Sch}/S)\)-groupoid such that \(X_0\) is of
finite presentation and quasi-projective over \(S\), \(d_1\) is of
finite presentation, proper, and flat, and the morphism}
\[
  d_0\boxtimes d_1:X_1\longrightarrow X_0\times_S X_0
\]
\textit{is quasi-finite. Then:}
\begin{enumerate}
\item[(1)]\phantomsection\label{thm:sga3-V-9.0-1}%
\textit{Every point \(x\) of \(X_0\) has a saturated open
neighborhood \(W\) such that the groupoid induced by \(X_\ast\) on
\(W\) has a quasi-section.}

\item[(2)]\phantomsection\label{thm:sga3-V-9.0-2}%
\textit{Let \((Y,p)\) be the cokernel of \((d_0,d_1)\) in the
category of all ringed spaces. Then \(Y\) is a scheme, \(p\) is a
morphism of schemes, and \((Y,p)\) is a cokernel of \((d_0,d_1)\) in
\((\mathbf{Sch}/S)\).}

\item[(3)]\phantomsection\label{thm:sga3-V-9.0-3}%
\textit{The morphism \(p\) is surjective, open, and universally
closed.}

\item[(4)]\phantomsection\label{thm:sga3-V-9.0-4}%
\textit{The morphism}
\[
  d:X_1\longrightarrow X_0\times_Y X_0
\]
\textit{with components \(d_0\) and \(d_1\) is surjective.}

\item[(5)]\phantomsection\label{thm:sga3-V-9.0-5}%
\textit{If \((d_0,d_1)\) is an equivalence pair, then:}
  \begin{enumerate}
  \item[(a)]\phantomsection\label{thm:sga3-V-9.0-5-a}%
  \textit{the morphism}
  \[
    d:X_1\longrightarrow X_0\times_Y X_0
  \]
  \textit{is an isomorphism and \(p\) is faithfully flat;}

  \item[(b)]\phantomsection\label{thm:sga3-V-9.0-5-b}%
  \textit{the morphisms \(p\) and \(Y\to S\) are of finite
  presentation, and \((Y,p)\) is a cokernel of \((d_0,d_1)\) in the
  category of sheaves for the fppf topology.}
  \end{enumerate}
\end{enumerate}

\noindent\textit{Proof.} For \hyperref[thm:sga3-V-9.0-1]{(1)}, the question is local on \(S\), so
we may suppose that \(S=\operatorname{Spec}B\) is affine. There then
exist a ring \(A\) of finite type over \(\mathbb Z\), a morphism
\[
  S\longrightarrow T=\operatorname{Spec}A,
\]
and a \((\mathbf{Sch}/T)\)-groupoid \(Z_\ast\) such that
\[
  X_\ast\simeq Z_\ast\times_T S.
\]
Indeed, apply EGA~IV\textsubscript{3}, 8.8.3 to
\(S_0=\operatorname{Spec}\mathbb Z\) and
\(S_i=\operatorname{Spec}A_i\), where the \(A_i\) range over the
finitely generated \(\mathbb Z\)-subalgebras of \(B\). Moreover, by
EGA~IV\textsubscript{3}, 8.10.5, we may suppose that \(Z_\ast\)
satisfies the hypotheses of
Theorem~\hyperref[thm:sga3-V-7.1]{7.1}. Consequently, \(Z_\ast\) has
quasi-sections locally.

By part~\hyperref[par:sga3-V-9-a]{(a)}, the same is therefore true of
\(X_\ast\), and assertions \hyperref[thm:sga3-V-9.0-2]{(2)}, \hyperref[thm:sga3-V-9.0-3]{(3)}, \hyperref[thm:sga3-V-9.0-4]{(4)}, and \hyperref[thm:sga3-V-9.0-5-a]{(5)(a)} follow from
Lemma~\hyperref[lem:sga3-V-6.1]{6.1}, just as in the proof of
Theorem~\hyperref[thm:sga3-V-7.1]{7.1}.

\medskip
\phantomsection\label{proof:sga3-V-9.0-c}%
\noindent\textbf{c) Proof of \hyperref[thm:sga3-V-9.0-5-b]{(5)(b)}.}---%
Let us show that \(Y\to S\) is of finite
presentation.\footnote{\label{note:sga3-V-47}Editor's note: The
original states that this follows from
Proposition~\hyperref[sga3:v:proposition:9:1]{9.1} below. We were unable to
reconstruct that argument. The proof that follows was communicated
to us by O.~Gabber.}
By hypothesis, \((d_0^{X_\ast},d_1^{X_\ast})\) is an equivalence
pair; that is,
\[
  d^{X_\ast}:X_1\longrightarrow X_0\times_S X_0
\]
is a monomorphism. By EGA~IV\textsubscript{3}, 8.10.5, after enlarging
\(A\) if necessary, we may suppose that
\[
  d^{Z_\ast}:Z_1\longrightarrow Z_0\times_T Z_0
\]
is a monomorphism. Since \(T=\operatorname{Spec}A\), with \(A\)
Noetherian, Theorem~\hyperref[thm:sga3-V-7.1]{7.1} then shows that the
groupoid \(Z_\ast\) has a cokernel \((Q,q)\) in
\((\mathbf{Sch}/T)\), that \(q\) and \(Q\to T\) are of finite
presentation, and, moreover, that
\[
  q:Z_0\longrightarrow Q
\]
is faithfully flat and \(d^{Z_\ast}\) induces an isomorphism
\[
  Z_1\xrightarrow{\ \sim\ }Z_0\times_Q Z_0.
\]
Put \(Q_S=Q\times_T S\). Since \(X_i\simeq Z_i\times_T S\), we obtain
an isomorphism
\[
  d^{Z_\ast}\times_T S:
  X_1\xrightarrow{\ \sim\ }X_0\times_{Q_S}X_0.
\]
Denote its inverse by \(\phi\), and let \(\pi\) be the canonical
morphism
\[
  X_0\times_Y X_0\longrightarrow X_0\times_{Q_S}X_0.
\]

Then \(\phi\circ\pi\) is the inverse of
\[
  d_0\boxtimes d_1:
  X_1\xrightarrow{\ \sim\ }X_0\times_Y X_0.
\]
It follows that the equivalence relation defined by \(X_\ast\), that
is, the monomorphism shown in
Figure~\hyperref[fig:sga3-V-9.0-equivalence-relation]{\ref*{fig:sga3-V-9.0-equivalence-relation}},
\SGAThreeDiagram[0.78\textwidth]
  {figures/exp5/exp5_localp033_equivalence_relation.png}
  {The equivalence-relation monomorphism defined by \(X_\ast\).}
  {fig:sga3-V-9.0-equivalence-relation}
  {Expos\'e V, component-local p.~33 / combined-reader p.~295 / re-edited header p.~283}
identifies with the equivalence relation \(\mathscr R(q_S)\) defined
by \(q_S:X_0\to Q_S\). Since \(q_S\) is faithfully flat and of finite
presentation, hence a universal effective epimorphism,
\(\mathscr R(q_S)\) has quotient \(Q_S\); see \SGARefLink{sga3:iv:definition:3:3:2}{Expos\'e IV, 3.3.2}.
Consequently,
\[
  Y\simeq Q\times_T S,
\]
so \(Y\to S\) and \(p\) are of finite presentation. Moreover, by
\SGARefLink{sga3:iv:corollary:6:3:3}{Expos\'e IV, 6.3.3}, \((Y,p)\) is also a cokernel of \((d_0,d_1)\) in
the category of sheaves for the fppf topology.

\medskip
\SGARefTarget{sga3:v:proposition:9:1}
\phantomsection\label{prop:sga3-V-9.1}%
\noindent\textbf{Proposition 9.1.}---%
\textit{Consider morphisms of schemes}
\[
  X_0\xrightarrow{\,p\,}Y\xrightarrow{\,q\,}S
\]
\textit{such that \(qp\) is of finite type (respectively, of finite
presentation) and \(p\) is faithfully flat and of finite
presentation. Then \(q\) is of finite type (respectively, of finite
presentation).}%
\begingroup
\renewcommand{\thefootnote}{*}%
\footnote{\label{note:sga3-V-9.1-generalization}See
EGA~IV\textsubscript{4}, 17.7.5 for a more general result.}
\endgroup
\addtocounter{footnote}{-1}%

\noindent Since \(p\) is surjective and \(qp\) is quasi-compact,
\(q\) is quasi-compact. We may therefore suppose that \(S\), \(Y\),
and \(X_0\) are affine, with rings \(A\), \(B\), and \(C\),
respectively. We have
\[
  B=\varinjlim_i B_i,
\]
where the \(B_i\) range over the finitely generated
\(A\)-subalgebras of \(B\). Since \(C\) is of finite presentation
over \(B\), there exist an index \(i_0\), a \(B_{i_0}\)-algebra
\(C_{i_0}\) of finite presentation, and an isomorphism
\[
  C\simeq C_{i_0}\otimes_{B_{i_0}}B.
\]
For \(i\geq i_0\), put
\[
  C_i=C_{i_0}\otimes_{B_{i_0}}B_i;
\]
then
\[
  C\simeq C_i\otimes_{B_i}B.
\]

\SGAThreeDiagram[0.22\textwidth]
  {figures/exp5/exp5_localp033_algebra_limit_diagram.png}
  {The compatible \(A\)-algebras used in the limit argument.}
  {fig:sga3-V-9.1-algebra-limit}
  {Expos\'e V, component-local p.~33 / combined-reader p.~295 / re-edited header p.~283}

Since \(C\) is faithfully flat over \(B\), EGA~IV\textsubscript{3},
11.2.6 and 8.10.5(vi) give an index \(i_1\geq i_0\) such that
\(C_{i_1}\) is faithfully flat over \(B_{i_1}\). Consequently,
\(C_i\) is faithfully flat over \(B_i\) for every \(i\geq i_1\).
For \(i\geq i_1\), the canonical map \(C_i\to C\) is then injective,
because it is obtained from \(B_i\to B\) by faithfully flat base
extension.

% The re-edited source places original-edition marginal marker 286 here.

If \(C\) is of finite type over \(A\), it follows that there is an
index \(j\geq i_1\) such that \(C_j=C\), whence \(B_j=B\), because
\(C_j\) is faithfully flat over \(B_j\). Thus \(B\) is of finite type
over \(A\).

Now suppose that \(C\) is of finite presentation over \(A\). By what
precedes, \(B\) is of finite type over \(A\), and hence has the form
\[
  B=\overline B/I,
\]
where \(\overline B\) is a polynomial algebra over \(A\) in finitely
many indeterminates and \(I\) is an ideal of \(\overline B\). The
ideal \(I\) is the union of its finitely generated subideals
\(I_\alpha\); consequently,
\[
  B=\varinjlim_\alpha B_\alpha,
  \qquad B_\alpha=\overline B/I_\alpha.
\]
Proceeding as above, there exist an index \(\alpha_0\), a
\(B_{\alpha_0}\)-algebra \(C_{\alpha_0}\) of finite presentation,
and an isomorphism
\[
  C\simeq C_{\alpha_0}\otimes_{B_{\alpha_0}}B.
\]
For \(\alpha\geq\alpha_0\), put
\[
  C_\alpha=C_{\alpha_0}\otimes_{B_{\alpha_0}}B_\alpha,
\]
so that
\[
  C\simeq C_\alpha\otimes_{B_\alpha}B
  \qquad(\alpha\geq\alpha_0).
\]
Again by EGA~IV\textsubscript{3}, 11.2.6 and 8.10.5(vi),
\(C_\alpha\) is faithfully flat over \(B_\alpha\) for all sufficiently
large \(\alpha\). In that case, the kernels of the maps
\(C_\alpha\to C\) and \(C_\alpha\to C_\beta\), respectively, identify
with
\[
  C_\alpha\otimes_{B_\alpha}(I/I_\alpha)
  \quad\text{and}\quad
  C_\alpha\otimes_{B_\alpha}(I_\beta/I_\alpha)
  \qquad(\beta\geq\alpha).
\]

Since \(C_\alpha\) and \(C\) are of finite presentation over \(A\)
and \(C_\alpha\to C\) is surjective,
\[
  C_\alpha\otimes_{B_\alpha}(I/I_\alpha)
\]
is a finitely generated ideal,\footnote{\label{note:sga3-V-48}Editor's
note: see EGA~IV\textsubscript{1}, 1.4.4.} and it is the union of the
ideals
\[
  C_\alpha\otimes_{B_\alpha}(I_\beta/I_\alpha).
\]
Thus, for all sufficiently large \(\beta\),
\[
  C_\alpha\otimes_{B_\alpha}(I_\beta/I_\alpha)
  =
  C_\alpha\otimes_{B_\alpha}(I/I_\alpha).
\]
Since \(C_\alpha\) is faithfully flat over \(B_\alpha\), this also
gives \(I_\beta=I\). Therefore \(B\) is of finite presentation over
\(A\).

\bigskip
\phantomsection\label{sec:sga3-V-10}%
\section*{10. Supplement: quotients by a group scheme}
\addcontentsline{toc}{section}{10. Supplement: quotients by a group scheme}

Sections \hyperref[sec:sga3-V-10.2]{10.2}--\hyperref[sec:sga3-V-10.4]{10.4} below, written following suggestions of
M.~Raynaud, apply the preceding theorems to the case of an action of
a group scheme. For the reader's convenience, \hyperref[sga3:v:paragraph:10-1]{Section~10.1} first
reproduces statements 2.1--2.3 of Expos\'e XVI.

\medskip
\phantomsection\label{subsec:sga3-V-10.1}%
\noindent\SGAInlineRefTarget{sga3:v:paragraph:10-1}\textbf{10.1. Representability theorems for quotients.}---%
We first recall \hyperref[thm:sga3-V-10.1.1]{the following result}.

\medskip
\phantomsection\label{thm:sga3-V-10.1.1}%
\noindent\textbf{Theorem 10.1.1.}---%
\textit{Let \(S\) be a scheme, let \(X\) and \(Y\) be two
\(S\)-schemes, and let \(f:X\to Y\) be an \(S\)-morphism. Suppose
that one of the following two conditions holds:}
\begin{enumerate}
\item[\(\alpha\))]\phantomsection\label{thm:sga3-V-10.1.1-alpha}%
\textit{the morphism \(f\) is locally of finite presentation;}
\item[\(\beta\))]\phantomsection\label{thm:sga3-V-10.1.1-beta}%
\textit{the scheme \(S\) is locally Noetherian and \(X\) is locally
of finite type over \(S\).}
\end{enumerate}
\textit{Then the following conditions are equivalent:}
\begin{enumerate}
\item[(i)]\phantomsection\label{thm:sga3-V-10.1.1-i}%
\textit{there exist an \(S\)-scheme \(X'\) and a factorization}
\[
  f:X\xrightarrow{\,f'\,}X'\xrightarrow{\,f''\,}Y,
\]
\textit{where \(f'\) is a faithfully flat \(S\)-morphism locally of
finite presentation and \(f''\) is a monomorphism;}

\item[(ii)]\phantomsection\label{thm:sga3-V-10.1.1-ii}%
\textit{the first projection}
\[
  p_1:X\times_Y X\longrightarrow X
\]
\textit{is flat.}
\end{enumerate}
\textit{Moreover, if these conditions hold, then \((X',f')\) is a
quotient of \(X\) by the equivalence relation defined by \(f\) for
the fppf topology. Thus the factorization \(f=f''\circ f'\) in \hyperref[thm:sga3-V-10.1.1-i]{(i)} is
unique up to isomorphism.}

The case in which \(Y\) is locally Noetherian and \(X\) is of finite
type over \(Y\) is treated in \hyperref[bib:sga3-V-Mur65]{[Mur65]}, Corollary~2 to Theorem~2. We
shall see that the general case reduces to that one.

We first make a few remarks.

\medskip
\phantomsection\label{par:sga3-V-10.1.1-a}%
\noindent\textbf{a)} The implication \hyperref[thm:sga3-V-10.1.1-i]{(i)}\(\Rightarrow\)\hyperref[thm:sga3-V-10.1.1-ii]{(ii)} is
immediate. Indeed, the first projection
\[
  p'_1:X\times_{X'}X\longrightarrow X
\]
factors through \(X\times_Y X\):
\[
  p'_1:
  X\times_{X'}X\xrightarrow{\,u\,}X\times_Y X
  \xrightarrow{\,p_1\,}X.
\]
The morphism \(u\) is an isomorphism because \(f''\) is a
monomorphism, and \(p'_1\) is flat because \(f'\) is flat; hence
\(p_1\) is flat.

\medskip
\phantomsection\label{par:sga3-V-10.1.1-b}%
\noindent\textbf{b)} The assertions of
Theorem~\hyperref[thm:sga3-V-10.1.1]{10.1.1} are local on \(Y\), and
hence local on \(S\). They are also local on \(X\), as follows readily
from the fact that a flat morphism locally of finite presentation is
open (EGA~IV\textsubscript{3}, 11.3.1).

\medskip
\phantomsection\label{par:sga3-V-10.1.1-c}%
\noindent\textbf{c)} Under hypothesis
\(\hyperref[thm:sga3-V-10.1.1-alpha]{\alpha}\) of
Theorem~\hyperref[thm:sga3-V-10.1.1]{10.1.1}, the preceding remarks
reduce us to the case where \(X\) and \(Y\) are affine and \(f\) is
of finite presentation. Replacing \(S\) by \(Y\), if necessary, we
may suppose that \(X\) and \(Y\) are of finite presentation over
\(S\). EGA~IV\textsubscript{3}, 11.2.6 then reduces us to the case
where \(S\) is Noetherian.

\medskip
\phantomsection\label{par:sga3-V-10.1.1-d}%
\noindent\textbf{d)} Under hypothesis
\(\hyperref[thm:sga3-V-10.1.1-beta]{\beta}\) of
Theorem~\hyperref[thm:sga3-V-10.1.1]{10.1.1}, we may suppose that
\(S\), \(X\), and \(Y\) are affine, that \(S\) is Noetherian, and
that \(X\) is of finite type over \(S\). Write \(Y\) as a filtered
inverse limit of affine schemes \(Y_i\) of finite type over \(S\).
The schemes
\[
  X\times_{Y_i}X
\]
form a filtered decreasing family of closed subschemes of
\(X\times_S X\), whose inverse limit is \(X\times_Y X\). Since
\(X\times_S X\) is Noetherian, for all sufficiently large \(i\) we
have
\[
  X\times_{Y_i}X=X\times_Y X.
\]
It follows that
\[
  f_i:X\xrightarrow{\,f\,}Y\longrightarrow Y_i
\]
satisfies \hyperref[thm:sga3-V-10.1.1-ii]{condition (ii)} of
Theorem~\hyperref[thm:sga3-V-10.1.1]{10.1.1} whenever \(f\) does.
The equivalence relation on \(X\) defined by \(f\) is the same as the
one defined by \(f_i\). Thus it suffices to prove
\hyperref[thm:sga3-V-10.1.1-ii]{(ii)}\(\Rightarrow\)\hyperref[thm:sga3-V-10.1.1-i]{(i)} for \(f_i\), reducing us to the case where
\(Y\) is of finite type over \(S\).

\medskip
\noindent\textit{Application to group schemes.}
Let \(S\) be a scheme, let \(G\) be an \(S\)-group scheme locally of
finite presentation over \(S\), and suppose that \(G\) acts on the
left on an \(S\)-scheme \(X\). If \(X\to S\) has a section \(\xi\),
recall that the stabilizer \(\operatorname{Stab}_G(\xi)\) is
representable by a subgroup scheme of \(G\); see \SGARefLink{sga3:i:definition:2:3:3}{Expos\'e I, 2.3.3}.

\medskip
\phantomsection\label{thm:sga3-V-10.1.2}%
\noindent\textbf{Theorem 10.1.2.}---%
\textit{Let \(S\) be a scheme and let \(G\) be an \(S\)-group scheme
locally of finite presentation over \(S\), acting on an \(S\)-scheme
\(X\). Suppose that \(X\to S\) has a section \(\xi\) whose stabilizer
\(H\) in \(G\) is flat over \(S\). If one of the following
hypotheses holds:}
\begin{enumerate}
\item[\textit{a)}]\phantomsection\label{thm:sga3-V-10.1.2-a}%
\textit{\(X\) is locally of finite type over \(S\);}
\item[\textit{b)}]\phantomsection\label{thm:sga3-V-10.1.2-b}%
\textit{\(S\) is locally Noetherian,}
\end{enumerate}
\textit{then the fppf quotient sheaf \(G/H\) is representable by an
\(S\)-scheme, locally of finite presentation over \(S\), and the
\(S\)-morphism}
\[
  f:G\longrightarrow X,\qquad g\longmapsto g\cdot\xi
\]

% Exact continuation cursor: component-local p. 36 / combined-reader p. 298
% / re-edited header p. 286, beginning "factors as".
% SGA3-V-V2-EOL-PRESERVATION

% END INCLUDED SGA3_English_current_buildable_source_20260823/source/inputs/V/tex/components/21_expose_V_localp32_35_sections9_10_en.tex
%
                                        \IfFileExists{tex/components/22_expose_V_section10_group_actions_through_theorem1041_en.tex}{%
                                          % BEGIN INCLUDED SGA3_English_current_buildable_source_20260823/source/inputs/V/tex/components/22_expose_V_section10_group_actions_through_theorem1041_en.tex
% SGA 3, Expose V, Loop 1 English reconstruction.
% Scope: component-local pp. 36--38, continuation of Theorem 10.1.2
% through the statement of Theorem 10.4.1.
% Authority: Polo--Gille Exp5-13oct24.pdf, combined-reader pp. 298--300,
% re-edited headers 286--288 (original running markers 289--291).
% Jacob Reinhold's CC BY 4.0 Markdown is comparison/drafting material only.

\textit{factors as}
\SGAThreeDiagram[0.42\textwidth]
  {figures/exp5/p286_theorem1012_factorization.png}
  {The factorization of the orbit morphism in
   Theorem~\hyperref[thm:sga3-V-10.1.2]{10.1.2}.}
  {fig:sga3-V-10.1.2-factorization}
  {Expos\'e V, component-local p.~36 / combined-reader p.~298 /
   re-edited header p.~286}
\textit{where \(p\) is the canonical projection, which is faithfully
flat and locally of finite presentation, and \(i\) is a monomorphism.}

\medskip
\noindent\textit{Proof.}---
The morphism \(f\) makes \(G\) into an \(X\)-scheme. By definition of
the stabilizer of \(\xi\), the morphism
\phantomsection\label{eq:sga3-V-10.1.2-stabilizer-isomorphism}%
\[
  G\times_S H\longrightarrow G\times_XG,
  \qquad (g,h)\longmapsto(g,gh)
\]
is an isomorphism. Since \(H\) is flat over \(S\), the scheme
\(G\times_SH\) is flat over \(G\); hence the first projection
\[
  p_1:G\times_XG\longrightarrow G
\]
is flat.%
\footnote{\label{source-note:sga3-V-10.1.2-fiber-product-base}The source PDF prints \(G\times_SG\) in this projection.
The \hyperref[eq:sga3-V-10.1.2-stabilizer-isomorphism]{preceding displayed isomorphism} and the application of
Theorem~\hyperref[thm:sga3-V-10.1.1]{10.1.1} require the fiber product \(G\times_XG\), which is the
reading adopted here.}
Moreover, if \(X\) is locally of finite type over \(S\), then \(f\) is
locally of finite presentation (EGA~IV\(_1\), 1.4.3(v)); otherwise,
\(S\) is assumed locally Noetherian. It therefore suffices to apply
Theorem~\hyperref[thm:sga3-V-10.1.1]{10.1.1} to \(f\). It remains only
to see that \(G/H\) is locally of finite presentation over \(S\), and
this follows immediately from
Proposition~\hyperref[sga3:v:proposition:9:1]{9.1}.

\medskip
\phantomsection\label{cor:sga3-V-10.1.3}%
\noindent\textbf{Corollary 10.1.3.}---
\textit{Let \(S\) be a scheme and \(u:G\to H\) a morphism of
\(S\)-group schemes. Suppose that \(G\) is locally of finite
presentation over \(S\), and that either \(H\) is locally of finite
type over \(S\) or \(S\) is locally Noetherian.}

\textit{If \(K=\ker(u)\) is flat over \(S\), then the quotient group
\(G/K\) is representable by an \(S\)-group scheme locally of finite
presentation over \(S\), and \(u\) factors as}
\SGAThreeDiagram[0.43\textwidth]
  {figures/exp5/p286_corollary1013_factorization.png}
  {The factorization of \(u\) through \(G/K\).}
  {fig:sga3-V-10.1.3-factorization}
  {Expos\'e V, component-local p.~36 / combined-reader p.~298 /
   re-edited header p.~286}
\textit{where \(p\) is the canonical projection and \(i\) is a
monomorphism.}

\medskip
\noindent\textit{Proof.}---
Apply Theorem~\hyperref[thm:sga3-V-10.1.2]{10.1.2} with \(X=H\) and
with the unit section of \(H\) as \(\xi\).

\section*{10.2. Stabilizer of the diagonal}
\phantomsection\label{sec:sga3-V-10.2}
\addcontentsline{toc}{section}{\SGARefLink{sec:sga3-V-10.2}{10.2}. Stabilizer of the diagonal}

Let \(S\) be a Noetherian scheme, \(X\) an \(S\)-scheme of finite
type, and \(G\) a flat \(S\)-group scheme of finite type acting on the
left on \(X\). Thus we have an \(S\)-action
\(d_0:G\times_SX\to X\). Let \(d_1:G\times_SX\to X\) be projection
onto the second factor. As in \hyperref[par:sga3-V-2-a]{\S2(a)}, we have the groupoid
\SGAThreeDiagram[0.67\textwidth]
  {figures/exp5/p286_section102_action_groupoid.png}
  {The action groupoid attached to \((X,G)\).}
  {fig:sga3-V-10.2-action-groupoid}
  {Expos\'e V, component-local p.~36 / combined-reader p.~298 /
   re-edited header p.~286}
whose cokernel, when it exists, is denoted \(G\backslash X\).

\medskip
\phantomsection\label{def:sga3-V-10.2.1}%
\noindent\textbf{Definition 10.2.1.}---
\textit{We denote by \(F\subset G\times_SX\) the \emph{stabilizer of
the diagonal section}; that is, \(F\) is the \(X\)-scheme defined by
the Cartesian square}
\SGAThreeDiagram[0.47\textwidth]
  {figures/exp5/p287_definition1021_cartesian_square.png}
  {The stabilizer \(F\) of the diagonal section.}
  {fig:sga3-V-10.2.1-stabilizer}
  {Expos\'e V, component-local p.~37 / combined-reader p.~299 /
   re-edited header p.~287}

Then \(F\) is an \(X\)-subgroup scheme of \(G\times_SX\). Since
\(G\times_SX\) is of finite type over the Noetherian scheme \(S\), it
is Noetherian; consequently \(F\) is of finite type both over \(S\)
and over \(X\) (EGA~I, 6.3.5 and 6.3.6). If \(X\to S\) is separated,
then \(F\) is a closed \(X\)-subgroup scheme of \(G\times_SX\).

Recall that \(G\) is said to \emph{act freely} on \(X\) when
\[
  (d_0,d_1):G\times_SX\longrightarrow X\times_SX
\]
is a monomorphism (Expos\'e~III, 3.2.1). Equivalently, \(F\) is the
trivial group scheme over \(X\).

\section*{10.3. The case in which \(F\) is quasi-finite over \(X\)}
\phantomsection\label{sec:sga3-V-10.3}
\addcontentsline{toc}{section}{\SGARefLink{sec:sga3-V-10.3}{10.3}. The case in which F is quasi-finite over X}

Since \(F\) is of finite type over \(X\), it is quasi-finite over
\(X\) if and only if the stabilizers of the geometric points of \(X\)
are finite.

\medskip
\phantomsection\label{thm:sga3-V-10.3.1}%
\noindent\textbf{Theorem 10.3.1.}---%
\footnote{\label{note:sga3-V-49}Editor's note: here too there is a
largest open subset \(U\subset X\) satisfying the conclusions of the
\hyperref[thm:sga3-V-10.3.1]{theorem}; see editor's \hyperref[note:sga3-V-44]{note~44}.}
\textit{Under the hypotheses of \hyperref[sec:sga3-V-10.2]{10.2},
suppose that \(F\) is quasi-finite over \(X\). Then there is a dense,
\(G\)-saturated open subset \(U\subset X\) having the following
properties:}
\begin{enumerate}
\item[(i)]\phantomsection\label{thm:sga3-V-10.3.1-i}%
  \textit{in \((\mathbf{Sch}/S)\), the cokernel
  \(V=G\backslash U\) exists; moreover, \(V\) is a quotient in the
  category of ringed spaces;}
\item[(ii)]\phantomsection\label{thm:sga3-V-10.3.1-ii}%
  \textit{\(p:U\to V\) is surjective, open, and of finite
  presentation;}
\item[(iii)]\phantomsection\label{thm:sga3-V-10.3.1-iii}%
  \textit{\(V\) is of finite presentation over \(S\);}
\item[(iv)]\phantomsection\label{thm:sga3-V-10.3.1-iv}%
  \textit{the morphism
  \[
    G\times_SU\longrightarrow U\times_VU,
    \qquad (g,x)\longmapsto(gx,x),
  \]
  is surjective;}
\item[(v)]\phantomsection\label{thm:sga3-V-10.3.1-v}%
  \textit{if, in addition, \(G\) acts freely on \(X\), then
  \(U\to V\) is a left \(G\)-torsor, locally trivial for the fppf
  topology. In particular, \(U\to V\) is faithfully flat.}%
  \footnote{\label{note:sga3-V-50}Editor's note: if one assumes in
  addition that \(G\) is a reductive \(S\)-group scheme and that its
  free action on \(X\) is linearizable, then \(G\backslash X\) is
  known to be representable and \(X\to G\backslash X\) is a left
  \(G\)-torsor. This follows from results of Raynaud and Seshadri and
  appears in \(\hyperref[bib:sga3-V-CTS79]{[\mathrm{CTS79}]}\), Proposition~6.11.}
\end{enumerate}

\noindent\textit{Proof.}---
By hypothesis, the morphism
\[
  G\times_SX\longrightarrow X\times_SX,
  \qquad (g,x)\longmapsto(gx,x),
\]
is quasi-finite. Theorem~\hyperref[thm:sga3-V-8.1]{8.1} therefore
applies to the groupoid defined by \((X,G)\). Hence there is a dense
saturated open subset \(U\subset X\) for which the quotient
\(G\backslash U\) exists, and it has properties \hyperref[thm:sga3-V-10.3.1-i]{(i)}, \hyperref[thm:sga3-V-10.3.1-ii]{(ii)}, and \hyperref[thm:sga3-V-10.3.1-iii]{(iii)};
the \hyperref[thm:sga3-V-8.1]{same theorem} also gives \hyperref[thm:sga3-V-10.3.1-iv]{(iv)}.

To prove \hyperref[thm:sga3-V-10.3.1-v]{(v)}, recall that \(G\) acts freely on \(X\) if and only if
\((d_0,d_1)\) is an equivalence pair (Expos\'e~III, 3.2.1).%
\footnote{\label{source-note:sga3-V-10.3.1-proof-part-v}At this point the source PDF says \emph{to prove (iv)}.
The cited Theorem~\hyperref[thm:sga3-V-8.1-iv]{8.1(iv)} and the conclusion about a torsor establish
\hyperref[thm:sga3-V-10.3.1-v]{part~(v)} in the present numbering; \hyperref[thm:sga3-V-10.3.1-iv]{part~(iv)} follows from
Theorem~\hyperref[thm:sga3-V-8.1-iii]{8.1(iii)}.}
In that case Theorem~\hyperref[thm:sga3-V-8.1-iv]{8.1(iv)} shows that
\[
  G\times_SU\longrightarrow U\times_VU
\]
is an isomorphism and that \(p\) is faithfully flat and of finite
presentation. Thus \(U\) is a \(G\)-torsor over \(V\), locally trivial
for the fppf topology.

\section*{10.4. The case in which \(F\) is flat over \(X\)}
\phantomsection\label{sec:sga3-V-10.4}
\addcontentsline{toc}{section}{\SGARefLink{sec:sga3-V-10.4}{10.4}. The case in which F is flat over X}

Write
\[
  d=(d_0,d_1):G\times_SX\longrightarrow X\times_SX,
  \qquad d(g,x)=(gx,x).
\]
Recall that the sheaf-theoretic graph \(\widetilde\Gamma\) of the
equivalence relation associated with \((X,G)\) is the fppf
\(S\)-subsheaf of \(X\times_SX\) given by the image of
\((d_0,d_1)\). It is the fppf sheaf associated with the graph functor
\[
  T\longmapsto
  \Gamma(T)=\bigl\{(x_0,x_1)\in X(T)\times X(T)
  \mathrel{|} x_0\in G(T)x_1\bigr\}.
\]
Put \(G_X=G\times_SX\). For every \(S\)-scheme \(T\) there is a
surjective map
\[
  G_X(T)\longrightarrow\Gamma(T),
  \qquad (g,x)\longmapsto(gx,x),
\]
which induces a bijection%
\phantomsection\label{eq:sga3-V-10.4-phiT-bijection}%
\[
  \phi(T):G_X(T)/F(T)\xrightarrow{\ \sim\ }\Gamma(T).
\]
Indeed, if \((g,x),(g',x')\in G_X(T)\) have the same image, then
\(x'=x\) and \(g^{-1}g'x=x\), so that
\((g^{-1}g',x)\in F(T)\), and \((g,x)\) and \((g',x)\) have the same
image in \(G_X(T)/F(T)\).%
\footnote{\label{source-note:sga3-V-10.4-ft-typing}The source PDF prints
\((g^{-1}g'x,x)\in F(T)\). Since
\(F(T)\subseteq G(T)\times X(T)\), the first coordinate must be the
group element \(g^{-1}g'\); the English adopts this type-correct
reading. In the same sentence the source writes \(F_X(T)\) once in
the denominator; the \hyperref[eq:sga3-V-10.4-phiT-bijection]{displayed bijection immediately above} and the
\hyperref[eq:sga3-V-10.4-quotient-presheaf]{presheaf immediately below} both use the defined stabilizer \(F(T)\),
which is retained consistently here.}

By definition (see \SGARefLink{sga3:iv:theorem:4:4:1:item:ii:01}{Expos\'e~IV, 4.4.1(ii)}, or the proof of 5.2.1),
the quotient sheaf \(G_X/F\) is the fppf sheaf associated with the
presheaf%
\phantomsection\label{eq:sga3-V-10.4-quotient-presheaf}%
\[
  T\longmapsto G_X(T)/F(T)\simeq\Gamma(T).
\]
We therefore obtain an isomorphism of sheaves
\[
  \phi:G_X/F\xrightarrow{\ \sim\ }\widetilde\Gamma.
\]

\medskip
\phantomsection\label{thm:sga3-V-10.4.1}%
\noindent\textbf{Theorem 10.4.1.}---%
\footnote{\label{note:sga3-V-51}Editor's note: this is part~(2) of
Theorem~3 of \(\hyperref[bib:sga3-V-Ray67b]{[\mathrm{Ray67b}]}\). That note also sketches another
proof of Theorem~\hyperref[thm:sga3-V-10.1.1]{10.1.1}.}
\textit{Under the hypotheses of \hyperref[sec:sga3-V-10.2]{10.2},
the following assertions hold:}
\begin{enumerate}
\item[\textit{a)}]\phantomsection\label{thm:sga3-V-10.4.1-a}%
  \textit{the sheaf \(\widetilde\Gamma\) is representable if and only
  if \(F\) is flat over \(X\);}
\item[\textit{b)}]\phantomsection\label{thm:sga3-V-10.4.1-b}%
  \textit{suppose that \(F\) is flat over \(X\). Then the morphisms
  induced by \(d_1\) and \(d_0\)}
  \SGAThreeDiagram[0.37\textwidth]
    {figures/exp5/p288_theorem1041_double_arrow.png}
    {The two maps induced by \(d_1\) and \(d_0\).}
    {fig:sga3-V-10.4.1-double-arrow}
    {Expos\'e V, component-local p.~38 / combined-reader p.~300 /
     re-edited header p.~288}
  \textit{are faithfully flat and of finite presentation.}
\end{enumerate}

% END INCLUDED SGA3_English_current_buildable_source_20260823/source/inputs/V/tex/components/22_expose_V_section10_group_actions_through_theorem1041_en.tex
%
                                          \IfFileExists{tex/components/23_expose_V_theorem1041_proof_through_bibliography_en.tex}{%
                                            % BEGIN INCLUDED SGA3_English_current_buildable_source_20260823/source/inputs/V/tex/components/23_expose_V_theorem1041_proof_through_bibliography_en.tex
% SGA 3, Expose V, Loop 1 English reconstruction.
% Scope: component-local pp. 39--41, proof of Theorem 10.4.1 through
% Corollary 10.4.3 and the complete bibliography.
% Authority: Polo--Gille Exp5-13oct24.pdf, combined-reader pp. 301--303,
% re-edited headers 289--291 (original running markers 292--294).
% Jacob Reinhold's CC BY 4.0 Markdown is comparison/drafting material only.

\medskip
\phantomsection\label{proof:sga3-V-10.4.1-a}%
\noindent\textit{Proof of \hyperref[thm:sga3-V-10.4.1-a]{a)}.}---
Suppose that the fppf sheaf \(G_X/F\) is represented by an \(X\)-scheme
\(Y\). Then, by \SGARefLink{sga3:iv:corollary:6:3:3}{Expos\'e~IV, 6.3.3}, the morphism
\(p:G_X\to Y\) is faithfully flat and locally of finite presentation,
and the second square in the \hyperref[fig:sga3-V-10.4.1-first-cartesian]{following diagram} is Cartesian:
\SGAThreeDiagram[0.66\textwidth]
  {figures/exp5/p289_theorem1041_first_cartesian.png}
  {The Cartesian square used in the proof of
   Theorem~\hyperref[thm:sga3-V-10.4.1-a]{10.4.1(a)}.}
  {fig:sga3-V-10.4.1-first-cartesian}
  {Expos\'e V, component-local p.~39 / combined-reader p.~301 /
   re-edited header p.~289}
The \hyperref[fig:sga3-V-10.4.1-first-cartesian]{first square} is obtained by base change along the unit section
\(e_X:X\to G_X\). Since \(p\) is faithfully flat and locally of
finite presentation, the same is true of \(F\to X\).

Conversely, suppose \(F\) is flat over \(X\), and put
\(X^2=X\times_SX\). The morphism \(d:G_X\to X^2\) gives the fiber
square
\SGAThreeDiagram[0.48\textwidth]
  {figures/exp5/p289_theorem1041_fiber_square.png}
  {The fiber square associated with \(d:G_X\to X^2\).}
  {fig:sga3-V-10.4.1-fiber-square}
  {Expos\'e V, component-local p.~39 / combined-reader p.~301 /
   re-edited header p.~289}
\phantomsection\label{eq:sga3-V-10.4.1-pr1-torsor}%
The first projection
\[
  \operatorname{pr}_1:G_X\times_{X^2}G_X\longrightarrow G_X
\]
is then an \(F\times_XG_X\)-torsor over \(G_X\), and is therefore flat
and of finite type, because \(F\) is.%
\footnote{\label{source-note:sga3-V-10.4.1-pr1-torsor}The source PDF instead prints the morphism
\(G_X\times_{X^2}G_X\to X^2\) and calls it an
\(F\times_XX^2\)-torsor over \(X^2\). That morphism is not generally
surjective. The \hyperref[eq:sga3-V-10.4.1-pr1-torsor]{first projection} to \(G_X\) is the stabilizer torsor
and is the flatness input required by Theorem~\hyperref[thm:sga3-V-10.1.1]{10.1.1}; the English uses
this type-correct reading.}
By
Theorem~\hyperref[thm:sga3-V-10.1.1]{10.1.1}, \(d\) factors uniquely
as
\SGAThreeDiagram[0.47\textwidth]
  {figures/exp5/p289_theorem1041_factorization.png}
  {The unique factorization of \(d\).}
  {fig:sga3-V-10.4.1-factorization}
  {Expos\'e V, component-local p.~39 / combined-reader p.~301 /
   re-edited header p.~289}
where \(\psi\) is faithfully flat and of finite type, and \(\tau\) is
a monomorphism of schemes.

It follows that the morphism of sheaves \(\psi:G_X\to Y\) is
\(F\)-invariant, and hence induces a morphism
\[
  \bar\psi:G_X/F\longrightarrow Y.
\]
Moreover, because \(\psi\) is faithfully flat and of finite type, the
monomorphism of sheaves \(\tau\) factors through the sheaf image of
\(d\), namely \(\widetilde\Gamma\). Thus the isomorphism
\(G_X/F\simeq\widetilde\Gamma\) factors through the monomorphism
\(Y\to\widetilde\Gamma\). We conclude that \(Y\) represents \(G_X/F\).

\medskip
\phantomsection\label{proof:sga3-V-10.4.1-b}%
\noindent\textit{Proof of \hyperref[thm:sga3-V-10.4.1-b]{b)}.}---
Suppose \(F\) is flat over \(X\). By \hyperref[thm:sga3-V-10.4.1-a]{part~(a)} and \hyperref[proof:sga3-V-10.4.1-a]{its proof}, \(G_X/F\)
is representable, and
\[
  p:G_X\longrightarrow G_X/F
\]
is faithfully flat and of finite presentation. On the other hand,
the morphisms \(d_i:G_X\to X\), for \(i=0,1\), are faithfully flat
and of finite presentation by hypothesis. Since
\(d_i=\bar d_i\circ p\), it follows from EGA~IV\(_2\), 2.2.13(iii),
and EGA~IV\(_3\), 11.3.16, that each \(\bar d_i\) is faithfully flat
and of finite presentation.

\medskip
\phantomsection\label{thm:sga3-V-10.4.2}%
\noindent\textbf{Theorem 10.4.2.}---%
\footnote{\label{note:sga3-V-52}Editor's note: here too there is a
largest open subset \(U\subset X\) satisfying the conclusions of the
\hyperref[thm:sga3-V-10.4.2]{theorem}. Moreover, a codimension-one point \(x\in X\) belongs to
\(U\) if and only if
\[
  (G_X/F)\times_X\operatorname{Spec}(\mathcal O_{X,x})
  \longrightarrow
  \operatorname{Spec}(\mathcal O_{X,x})
  \times_S\operatorname{Spec}(\mathcal O_{X,x})
\]
is a closed immersion; see editor's \hyperref[note:sga3-V-44]{note~44}.}
\textit{Under the hypotheses of \hyperref[sec:sga3-V-10.2]{10.2},
suppose that \(F\) is flat over \(X\). Then there is a dense saturated
open subset \(U\subset X\) such that the fppf quotient
\(V=G\backslash U\) is an \(S\)-scheme of finite type and
\(U\to V\) is faithfully flat and of finite presentation.}

\medskip
\noindent\textit{Proof.}---
Theorem~\hyperref[thm:sga3-V-10.4.1]{10.4.1} shows that
\(G_X/F\simeq\widetilde\Gamma\) is representable. The fppf sheaf
\(G\backslash X\) is therefore the quotient sheaf of
\SGAThreeDiagram[0.38\textwidth]
  {figures/exp5/p290_theorem1042_double_arrow.png}
  {The equivalence pair whose quotient is \(G\backslash X\).}
  {fig:sga3-V-10.4.2-double-arrow}
  {Expos\'e V, component-local p.~40 / combined-reader p.~302 /
   re-edited header p.~290}
By \hyperref[thm:sga3-V-10.4.1-b]{what precedes},
\[
  \bar d_i:G_X/F\longrightarrow X
\]
is faithfully flat and of finite presentation for \(i=0,1\), and the
morphism
\SGAThreeDiagram[0.53\textwidth]
  {figures/exp5/p290_theorem1042_monomorphism_chain.png}
  {The monomorphism from the quotient \(G_X/F\) to \(X\times_SX\).}
  {fig:sga3-V-10.4.2-monomorphism}
  {Expos\'e V, component-local p.~40 / combined-reader p.~302 /
   re-edited header p.~290}
is a monomorphism. In other words, \((\bar d_0,\bar d_1)\) is an
equivalence pair. Consequently,
Theorem~\hyperref[thm:sga3-V-8.1]{8.1} applies. There is therefore a
dense saturated open subset \(U\subset X\) such that the fppf quotient
\(V=G\backslash U\) is an \(S\)-scheme of finite type and \(U\to V\)
is faithfully flat and of finite presentation.

Taking the generic-flatness theorem into account
(EGA~IV\(_2\), 6.9.3), we obtain the \hyperref[cor:sga3-V-10.4.3]{following result}.

\medskip
\phantomsection\label{cor:sga3-V-10.4.3}%
\noindent\textbf{Corollary 10.4.3.}---
\textit{Under the hypotheses of \hyperref[sec:sga3-V-10.2]{10.2},
suppose that \(X\) is reduced. Then there is a dense saturated open
subset \(U\subset X\) such that the fppf quotient
\(G\backslash U\) is an \(S\)-scheme of finite type and
\(U\to G\backslash U\) is faithfully flat and of finite
presentation.}

\section*{Bibliography\protect\footnotemark}
\phantomsection\label{sec:sga3-V-bibliography}
\addcontentsline{toc}{section}{Bibliography}
\footnotetext{\label{note:sga3-V-53}Editor's note: additional references cited in this expos\'e.}

\begin{description}
\item[{[AK80]}]\phantomsection\label{bib:sga3-V-AK80} A.~B. Altman and S.~L. Kleiman,
  \emph{Compactifying the Picard Scheme}, Adv. Math. \textbf{35}
  (1980), 50--112.
\item[{[An73]}]\phantomsection\label{bib:sga3-V-An73} S.~Anantharaman,
  \emph{Sch\'emas en groupes, espaces homog\`enes et espaces
  alg\'ebriques sur une base de dimension 1}, M\'em. Soc. Math.
  France \textbf{33} (1973), 5--79.
\item[{[BLR90]}]\phantomsection\label{bib:sga3-V-BLR90} S.~Bosch, W.~L\"utkebohmert, and M.~Raynaud,
  \emph{N\'eron models}, Springer-Verlag, 1990.
\item[{[CTS79]}]\phantomsection\label{bib:sga3-V-CTS79} J.-L.~Colliot-Th\'el\`ene and J.-J.~Sansuc,
  \emph{Fibr\'es quadratiques et composantes connexes r\'eelles},
  Math. Ann. \textbf{244} (1979), 105--134.
\item[{[DG70]}]\phantomsection\label{bib:sga3-V-DG70} M.~Demazure and P.~Gabriel,
  \emph{Groupes alg\'ebriques}, Masson \& North-Holland, 1970.
\item[{[DR81]}]\phantomsection\label{bib:sga3-V-DR81} J.~Dixmier and M.~Raynaud,
  \emph{Sur le quotient d'une vari\'et\'e alg\'ebrique par un groupe
  alg\'ebrique}, pp.~327--344 in \emph{Mathematical Analysis and
  Applications} (L.~Schwartz 65th birthday, ed. L.~Nachbin), Adv.
  Math. Suppl. Stud., Vol.~7A, 1981.
\item[{[Fe03]}]\phantomsection\label{bib:sga3-V-Fe03} D.~Ferrand,
  \emph{Conducteur, descente et pincement}, Bull. Soc. Math. France
  \textbf{131} (2003), no.~4, 553--585.
\item[{[Hi62]}]\phantomsection\label{bib:sga3-V-Hi62} H.~Hironaka,
  \emph{An example of a non-K\"ahlerian complex analytic deformation
  of K\"ahlerian complex structures}, Ann. of Math. \textbf{75}
  (1962), no.~1, 190--208.
\item[{[KM97]}]\phantomsection\label{bib:sga3-V-KM97} S.~Keel and S.~Mori,
  \emph{Quotient by groupoids}, Ann. of Math. \textbf{145} (1997),
  no.~1, 193--213.
\item[{[Ko97]}]\phantomsection\label{bib:sga3-V-Ko97} J.~Koll\'ar,
  \emph{Quotient spaces modulo algebraic groups}, Ann. of Math.
  \textbf{145} (1997), no.~1, 33--79.
\item[{[Ko08]}]\phantomsection\label{bib:sga3-V-Ko08} J.~Koll\'ar,
  \emph{Quotients by finite equivalence relations}, arXiv:0812.3608.
\item[{[Mum65]}]\phantomsection\label{bib:sga3-V-Mum65} D.~Mumford,
  \emph{Geometric invariant theory}, Springer-Verlag, 1965; 2nd ed.
  with J.~Fogarty, 1982; 3rd ed. with J.~Fogarty and F.~Kirwan, 1994.
\item[{[Mur65]}]\phantomsection\label{bib:sga3-V-Mur65} J.~P. Murre,
  \emph{Representation of unramified functors. Applications}
  (according to unpublished results of A.~Grothendieck), S\'em.
  Bourbaki, Vol.~9, Exp.~294 (1965), Soc. Math. France, 1995.
\item[{[Ray67a]}]\phantomsection\label{bib:sga3-V-Ray67a} M.~Raynaud,
  \emph{Passage au quotient par une relation d'\'equivalence plate},
  pp.~78--85 in \emph{Proc. Conf. Local Fields (Driebergen)}
  (ed. T.~A. Springer), Springer-Verlag, 1967.
\item[{[Ray67b]}]\phantomsection\label{bib:sga3-V-Ray67b} M.~Raynaud,
  \emph{Sur le passage au quotient par un groupo\"ide plat},
  C. R. Acad. Sci. Paris (S\'er. A) \textbf{265} (1967), 384--387.
\end{description}

% END INCLUDED SGA3_English_current_buildable_source_20260823/source/inputs/V/tex/components/23_expose_V_theorem1041_proof_through_bibliography_en.tex
%
                                          }{}%
                                        }{}%
                                      }{}%
                                    }{}%
                                  }{}%
                                }{}%
                              }{}%
                            }{}%
                          }{}%
                        }{}%
                      }{}%
                    }{}%
                  }{}%
                }{}%
              }{}%
            }{}%
          }{}%
        }{}%
      }{}%
    }{}%
  }{}%
}{}%


% END INCLUDED SGA3_English_current_buildable_source_20260823/source/inputs/V/body.tex


\SGAStartExpose{VI}
% BEGIN INCLUDED SGA3_English_current_buildable_source_20260823/source/inputs/VI/body.tex
% Mechanically extracted after \mainmatter from the frozen standalone master.

% BEGIN INCLUDED SGA3_English_current_buildable_source_20260823/source/inputs/VI/tex_reference_v2/components/00_expose_VIA_title_through_corollary_0_5_2_loop2_r1_en.tex
% SGA 3, Expose VI_A, English reconstruction -- Loop 1.
% Scope: component-local pages 1--4, from the title through the end of
% Corollary 0.5.2, stopping before the Section 1 heading.
% Combined-reader pages: 304--307. Printed pages: 293--296.
% Authority: Polo--Gille Exp6A-13oct24.pdf,
% SHA-256 6C8795572D1FC3B21BD02D24B1BFB20A8FB447475702C2674C96939B1EB69582.
% Jacob Reinhold's CC BY 4.0 English Markdown was consulted only as a
% comparison lineage. The Polo--Gille PDF controls prose, formulas,
% numbering, footnotes, and the source-bound Loop-1 diagram.

\setcounter{footnote}{-1}
\chapter*{Expos\'e \texorpdfstring{VI\textsubscript{A}}{VI_A}\\
Generalities on Algebraic Groups}
\phantomsection
\label{sga3:VIA:title}
\addcontentsline{toc}{chapter}{Expos\'e VI-A. Generalities on Algebraic Groups}
\markboth{EXPOS\'E VI\textsubscript{A}. GENERALITIES ON ALGEBRAIC GROUPS}%
  {EXPOS\'E VI\textsubscript{A}. GENERALITIES ON ALGEBRAIC GROUPS}

\begin{center}
  \textit{by P. Gabriel}
\end{center}

% Authority: Exp6A local p. 1; combined-reader p. 304; printed p. 293.
\footnote{Editor's note: version of October 13, 2024.}Throughout this
chapter, \(A\) will denote an Artinian local ring with residue field \(k\).
A group scheme \(G\) over \(\operatorname{Spec} A\) will be called simply
an \(A\)-\emph{group}. This \(A\)-group is said to be \emph{locally of
finite type} if its underlying scheme is locally of finite type over \(A\);
it is said to be \emph{algebraic} if its underlying scheme is of finite
type over \(A\).

\section*{0. Preliminary Remarks}
\phantomsection
\label{sga3:VIA:0}
\addcontentsline{toc}{section}{0. Preliminary Remarks}

\medskip
\noindent\textbf{0.1.}
\phantomsection
\label{sga3:VIA:0.1}
Let us first consider a group scheme \(G\) over an arbitrary scheme \(S\).
We call the structural morphism
\[
  \mu:G\times_S G\longrightarrow G
\]
\emph{multiplication}, and we call the morphism \(c:G\to G\), defined by
\(c(T)(x)=x^{-1}\) for every scheme \(T\) over \(S\) and every
\(x\in G(T)\), \emph{inversion}. If \(U\) and \(V\) are subsets of the
underlying set of \(G\), we denote by \(U\cdot V\) the image under
multiplication of the subset of \(G\times_S G\) consisting of those points
whose first projection belongs to \(U\) and whose second projection belongs
to \(V\). Likewise, the notations \(U^{-1}\) and \(c(U)\) are equivalent.

Let \(\operatorname{pr}_1\) be the projection of \(G\times_S G\) onto the
first factor, and let
\[
  \sigma:G\times_S G\longrightarrow G\times_S G
\]
be the morphism with components \(\operatorname{pr}_1\) and \(\mu\). For
every \(S\)-scheme \(T\), the map \(\sigma(T)\) is
\[
  (x,y)\longmapsto (x,xy);
\]
it follows that \(\sigma\) is an automorphism. The composite of this
automorphism with the projection \(\operatorname{pr}_2\) of
\(G\times_S G\) onto the second factor is the multiplication morphism.
When \(G\) is flat over \(S\), \(\operatorname{pr}_2\), and hence \(\mu\),
are flat morphisms; when \(G\) is smooth over \(S\),
\(\operatorname{pr}_2\), and hence \(\mu\), are smooth morphisms, and so
on.

\medskip
\noindent\textbf{0.2.}
\phantomsection
\label{sga3:VIA:0.2}
From now on we assume that \(S\) is the spectrum of an Artinian local ring
\(A\) with residue field \(k\). We denote by
\((\mathbf{Sch}_{/k})_{\mathrm{red}}\) the category of reduced schemes over
\(k\). For every scheme \(X\) over \(A\), the reduced scheme
\(X_{\mathrm{red}}\) is an object of
\((\mathbf{Sch}_{/k})_{\mathrm{red}}\), and the functor
\[
  X\longmapsto X_{\mathrm{red}}
\]
is right adjoint to the inclusion of
\((\mathbf{Sch}_{/k})_{\mathrm{red}}\) into \(\mathbf{Sch}_{/A}\).
It follows that, for every \(A\)-group \(G\), \(G_{\mathrm{red}}\) is a
group in the category
\((\mathbf{Sch}_{/k})_{\mathrm{red}}\).\footnote{Editor's note: the
sentences that follow have been added.} In other words, for every reduced
\(k\)-scheme \(T\), \(G_{\mathrm{red}}(T)\) is endowed, functorially in
\(T\), with a group structure. One must bear in mind that
  \(G_{\mathrm{red}}\) is not necessarily a \(k\)-group, because its
  ``multiplication'' is only a morphism
\stepcounter{footnote}%
\[
  (G_{\mathrm{red}}\times_k G_{\mathrm{red}})_{\mathrm{red}}
  \longrightarrow
  G_{\mathrm{red}}\text{\hyperlink{sga3:VIA:fn:red-examples}{%
    \textsuperscript{\normalfont\thefootnote}}}.
\]
\footnotetext[\value{footnote}]{\hypertarget{sga3:VIA:fn:red-examples}{}Editor's note: for examples of group schemes \(G\) over a
nonperfect field \(k\) for which \(G_{\mathrm{red}}\) is not a group scheme
over \(k\), see \SGARefLink{ex:sga3-VIA-1.3.2}{1.3.2} below.}

\footnote{Editor's note: the remainder of \SGARefLink{sga3:VIA:0.2}{0.2}, as well as \SGARefLink{sga3:VIA:0.3}{0.3}, has been
slightly modified.}
However, if \(k\) is a perfect field, the inclusion of
\((\mathbf{Sch}_{/k})_{\mathrm{red}}\) into \(\mathbf{Sch}_{/k}\)
commutes with products. Groups in
\((\mathbf{Sch}_{/k})_{\mathrm{red}}\) may therefore be identified with
group schemes over \(k\) whose underlying schemes are reduced. In this
case, if \(G\) is a \(k\)-group, \(G_{\mathrm{red}}\) is a group subscheme
of \(G\); in general, this subgroup is not normal in \(G\).

For example, if \(k\) is a field of characteristic \(3\), the constant
group \((\mathbb Z/2\mathbb Z)_k\) acts nontrivially on the diagonalizable
group \(D_k(\mathbb Z/3\mathbb Z)\) (cf. \SGARefLink{sga3:i:subsection:4-1}{Expos\'e I, 4.1} and \SGARefLink{cor:sga3-VIB-4.4}{4.4}). If
\(G\) denotes the semidirect product of \(D_k(\mathbb Z/3\mathbb Z)\) by
\((\mathbb Z/2\mathbb Z)_k\) defined by this action, then
\(G_{\mathrm{red}}\) is identified with \((\mathbb Z/2\mathbb Z)_k\) and
is not normal in \(G\).\footnote{Editor's note: for an analogous example
with \(G\) connected, one may take, when
\(\operatorname{char}(k)=p>0\), the semidirect product of
\(\alpha_{p,k}\) and \(\mathbb G_{m,k}\).}

Let \(k\) be an arbitrary field, let \(k^{p^{-\infty}}\) be its perfect
closure, and let \(H\) be a group in
\((\mathbf{Sch}_{/k})_{\mathrm{red}}\). Then
\[
  (H\otimes_k k^{p^{-\infty}})_{\mathrm{red}}
\]
is a group scheme over the perfect field \(k^{p^{-\infty}}\). Since
\(H\otimes_k k^{p^{-\infty}}\) and
\((H\otimes_k k^{p^{-\infty}})_{\mathrm{red}}\) have the same underlying
topological space, groups in \((\mathbf{Sch}_{/k})_{\mathrm{red}}\) share
with \(k\)-groups certain topological properties that are invariant under
extension of the base field. For example, it will follow from \hyperref[sga3:VIA:0.3]{0.3} and the
preceding remarks that every group in
\((\mathbf{Sch}_{/k})_{\mathrm{red}}\) is separated.

In what follows, we shall encounter groups in
\((\mathbf{Sch}_{/k})_{\mathrm{red}}\) whenever we deal with a nonempty
locally closed subset \(U\) of an \(A\)-group \(G\) such that
\[
  U\cdot U=U
  \qquad\text{and}\qquad
  U^{-1}=U.
\]
Indeed, the reduced subscheme of \(G\) defined by \(U\) is then a group in
\((\mathbf{Sch}_{/k})_{\mathrm{red}}\).

% Authority: Exp6A local p. 2; combined-reader p. 305; printed p. 294.
\medskip
\noindent\textbf{0.3.}
\phantomsection
\label{sga3:VIA:0.3}
\textit{An \(A\)-group \(G\) is always separated, because the unit section
\(e:\operatorname{Spec}A\to G\) is a closed immersion.}
Indeed, let \(x\) be the unique point of \(\operatorname{Spec}A\), and let
\(\eta:G\to\operatorname{Spec}A\) be the structural morphism. Since
\(\eta\circ e=\operatorname{id}_{\operatorname{Spec}A}\), for every
affine open \(U=\operatorname{Spec}B\) of \(G\) containing \(e(x)\), the
homomorphism \(B\to A\) has a section and is therefore surjective. It
follows that \(e\) is a closed immersion.\footnote{Editor's note: this
argument is valid for every local ring \(A\) of dimension zero, and shows
that, if \(S\) is a discrete scheme, every \(S\)-group is separated; cf.
VI\textsubscript{B}, \SGARefLink{prop:sga3-VIB-5.2}{5.2}. On the other hand (cf.
VI\textsubscript{B}, \SGARefLink{ctr:sga3-VIB-5.6.4}{5.6.4}), if \(S\) contains a nonisolated closed point
\(s\), the \(S\)-scheme \(G\) obtained by gluing two copies of \(S\) along
the open subset \(S-\{s\}\) is not separated over \(S\), although it
carries an \(S\)-group structure.}

The diagonal in \(G\times_A G\) is identified with the functor from
\((\mathbf{Sch}_{/A})^\circ\) to \(\mathbf{Sets}\) which associates to
every \(A\)-scheme \(S\) the inverse image of the identity element of
\(G(S)\) under the map
\[
  \varphi(S):G(S)\times G(S)\longrightarrow G(S),
  \qquad (x,y)\longmapsto x\cdot y^{-1}.
\]
We therefore have the following cartesian square. Since the diagonal
morphism is obtained from a closed immersion by base change, it is itself
a closed immersion.

\begin{center}
  \phantomsection
  \label{sga3:VIA:0.3-cartesian-square}
  \begin{tikzcd}[row sep=large,column sep=large]
    G\times_S G
      \arrow[r,"\varphi"]
    &
    G
    \\
    G
      \arrow[u,"\text{diagonal morphism}"]
      \arrow[r]
    &
    \operatorname{Spec}A
      \arrow[u,"\text{unit section}"']
  \end{tikzcd}
\end{center}

% The exact source-PDF crop above prints G x_S G in its upper-left corner.
% The surrounding source prose identifies the fiber product as G x_A G.

% Authority: Exp6A local p. 3; combined-reader p. 306; printed p. 295.
\medskip
\noindent\textbf{0.4.}\footnote{Editor's note: the original has been
expanded in what follows; in particular, \hyperref[sga3:VIA:0.4.1]{Remark 0.4.1} has been added.}
\phantomsection
\label{sga3:VIA:0.4}
Let \(G\) be an \(A\)-scheme. We shall say that a point \(g\) of \(G\) is
\emph{strictly rational over \(A\)} if there exists an \(A\)-morphism
\[
  s:\operatorname{Spec}A\longrightarrow G
\]
which sends the unique point of \(\operatorname{Spec}A\) to \(g\); that
is, if the homomorphism \(A\to\mathcal O_{G,g}\) admits a retraction. One
then has \(\kappa(g)=k\), so \(g\) is a closed point of \(G\). Indeed, if
\(B\) is the ring of an affine open neighborhood of \(g\), and if
\(\mathfrak p\) is the prime ideal of \(B\) corresponding to \(g\), then
\(B/\mathfrak p\subset k\) is a finite \(A\)-algebra, hence an integral
Artinian ring, and therefore a field.

Suppose henceforth that \(G\) is an \(A\)-group. Such a morphism
\(s:\operatorname{Spec}A\to G\) defines an automorphism \(r_s\) of the
scheme \(G\) over \(A\), called \emph{right translation by \(s\)}. For
every morphism \(\pi:S\to\operatorname{Spec}A\), the automorphism
\(r_s(\pi)\) of \(G(S)\) is defined by
\[
  x\longmapsto x\cdot G(\pi)(s),
  \qquad x\in G(S).
\]
Similarly, \(\ell_s\) denotes \emph{left translation by \(s\)}, namely
the automorphism of \(G\) defined by
\[
  \ell_s(\pi)(x)=G(\pi)(s)\cdot x,
  \qquad x\in G(S).
\]

The schemes \(G\otimes_A k\) and \(G\) have the same underlying
topological space \(\underline G\), \(G\otimes_A k\) is a \(k\)-group,
and \(s\otimes_A k\) depends only on \(g\), not on \(s\). Consequently,
the automorphisms of \(\underline G\) induced by \(r_s\) and \(\ell_s\)
(or by \(r_{s\otimes k}\) and \(\ell_{s\otimes k}\)) depend only on \(g\),
not on \(s\). Thus, if \(P\) is a subset of \(\underline G\), we may write
\[
  r_g(P)\ \text{or}\ P\cdot g
  \qquad
  \bigl(\text{respectively, }\ell_g(P)\ \text{or}\ g\cdot P\bigr)
\]
instead of \(r_s(P)\) (respectively, \(\ell_s(P)\)), in agreement with
\hyperref[sga3:VIA:0.1]{0.1}.

\medskip
\noindent\textbf{Remark 0.4.1.}
\phantomsection
\label{sga3:VIA:0.4.1}
\textit{If \(g\) is a strictly rational point of \(G\), and if
\(A\to A'\) is a morphism of Artinian local rings, then
\(G'=G\otimes_A A'\) has a unique point \(g'\) above \(g\), and \(g'\) is
strictly rational over \(A'\). Moreover, if \(P'\) denotes the inverse
image of \(P\) in \(G'\), then \(P'\cdot g'\) is the inverse image of
\(P\cdot g\) (cf. EGA I, 3.4.8).}

\medskip
\noindent\SGAInlineRefTarget{sga3:vi:proposition:0-5}\textbf{Proposition 0.5.}\footnote{Editor's note: the original
stated this result under the hypothesis that \(G\) be locally of finite
type over \(A\). Since the result is useful in the general case, and
since the proof is essentially the same, it has been stated and proved
here in full generality. It will be used several times below.}
\phantomsection
\label{sga3:VIA:0.5}
\textit{Let \(G\) be an \(A\)-group, and let \(U,V\) be two dense open
subsets of \(G\). Then \(U\cdot V\), that is, the image of
\(U\times_A V\) under multiplication, is the entire underlying space of
\(G\).}

\begin{proof}
Since \(G\) and \(G\otimes_A k\) have the same underlying space, after
replacing \(A\) by \(k\) and \(G\) by \(G\otimes_A k\), if necessary, we
may assume that \(A=k\). Let \(g\in G\), and set \(K=\kappa(g)\). Then
left translation \(\ell_g\) is an automorphism of \(G_K\). Since the
projection \(G_K\to G\) is open, \(U_K\) and \(V_K\) are dense open
subsets of \(G_K\), as is the image of \(V_K\) under \(\ell_g\).
Consequently there exists \(v\in V_K\) such that
\(u=\ell_g(v)\) belongs to \(U_K\).

Let \(L\) be an extension of \(K\) containing \(\kappa(v)\), and hence
\(\kappa(u)\), and let \(g_L\) and \(v_L\) be the \(L\)-points of \(G_L\)
induced by \(g\) and \(v\). Then
\[
  g_L\cdot v_L=u'
\]
is a point of \(G_L\) above \(u\), and therefore
\[
  g_L=u'\cdot v_L^{-1}
\]
lies above \(U\cdot V\). Thus \(g\in U\cdot V\), proving the proposition.
\end{proof}

% Authority: Exp6A local p. 4; combined-reader p. 307; printed p. 296.
\medskip
\noindent\textbf{Corollary 0.5.1.}\footnote{Editor's note: this corollary
has been inserted here; cf. \SGARefLink{prop:sga3-VIA-2.4}{2.4} and \SGARefLink{cor:sga3-VIA-2.6.2}{2.6.2} below.}
\phantomsection
\label{sga3:VIA:0.5.1}
\textit{If \(G\) is an irreducible \(A\)-group, then \(G\) is
quasi-compact.}

\begin{proof}
Let \(U\) be a nonempty affine open subset of \(G\). Then \(U\) is dense
in \(G\), so by \hyperref[sga3:VIA:0.5]{0.5} the morphism
\[
  \mu:U\times_A U\longrightarrow G
\]
is surjective. Hence \(G\) is quasi-compact, since \(U\times_A U\) is.
\end{proof}

\medskip
\noindent\textbf{Corollary 0.5.2.}\footnote{Editor's note: this
corollary, used implicitly in VIII, \SGARefLink{sga3:viii:corollary:6:7}{6.7}, has been added; see also
VI\textsubscript{B}, \SGARefLink{cor:sga3-VIB-6.2.5}{6.2.5}.}
\phantomsection
\label{sga3:VIA:0.5.2}
\textit{Let \(G\) be an \(A\)-group scheme, and let \(H\) be a group
subscheme of \(G\) over \(A\). Then \(H\) is closed.}

\begin{proof}
Let \(k'\) be the perfect closure of the residue field \(k\) of \(A\).
Since the underlying topological spaces of \(G\) and \(H\) are unchanged
under the base change \(A\to k\to k'\), we may assume that \(A=k\) is a
perfect field. We may then assume that \(G\) and \(H\) are reduced, hence
geometrically reduced.

Let \(\overline H\) be the closure of \(H\). Then
\(\mu^{-1}(\overline H)\) is a closed subset of \(G\times G\) containing
\(H\times H\). The morphisms
\[
  H\longrightarrow\operatorname{Spec}k
  \qquad\text{and}\qquad
  \overline H\longrightarrow\operatorname{Spec}k
\]
are universally open. Since \(H\) is dense in \(\overline H\),
\(H\times H\) is dense in \(\overline H\times H\), and
\(\overline H\times H\) is dense in
\(\overline H\times\overline H\). Thus \(H\times H\) is dense in
\(\overline H\times\overline H\). Consequently
\[
  \mu(\overline H\times\overline H)\subset\overline H.
\]
Since \(\overline H\times\overline H\) is reduced, \(\mu\) induces a
morphism
\[
  \mu':\overline H\times\overline H\longrightarrow\overline H.
\]

Now let \(g\in\overline H\), and put \(K=\kappa(g)\). The projection
\(\overline H_K\to\overline H\) is open, so \(H_K\) and
\(\ell_g(H_K)\) are two dense open subsets of \(\overline H_K\). Hence
there exist \(u,v\in H_K\) such that \(\ell_g(v)=u\). As in the proof of
\SGARefLink{sga3:vi:proposition:0-5}{0.5}, it follows that \(g\) belongs to \(H\cdot H=H\). Therefore
\(\overline H=H\).
\end{proof}

% Exact continuation cursor: Exp6A local p. 4, immediately before
% "1. Local properties of an A-group locally of finite type".

% END INCLUDED SGA3_English_current_buildable_source_20260823/source/inputs/VI/tex_reference_v2/components/00_expose_VIA_title_through_corollary_0_5_2_loop2_r1_en.tex

% BEGIN INCLUDED SGA3_English_current_buildable_source_20260823/source/inputs/VI/tex_reference_v2/components/01_expose_VIA_section1_through_prop131_en.tex
% SGA 3, Expose VI_A, Loop 1 English reconstruction.
% Scope: Section 1 heading through the transition after Proposition 1.3.1;
% Exemples 1.3.2 is excluded.
% Authority: Polo--Gille Exp6A-13oct24.pdf, component-local pages 4--6
% (combined-reader pages 307--309; printed pages 296--298;
% current-rescribe indices 292--295).
% OCR and the Reinhold English Markdown are comparison witnesses only.

\section*{1. Local properties of an \(A\)-group locally of finite type}
\phantomsection\label{sec:sga3-VIA-1}
\addcontentsline{toc}{section}{1. Local properties of an
  \texorpdfstring{\(A\)}{A}-group locally of finite type}

We shall first see that, if \(G\) is locally of finite type and flat over
\(A\), one can ``make every closed point of \(G\) strictly rational'' by
means of a finite flat extension of the base.

\medskip
\phantomsection\label{par:sga3-VIA-1.1}%
\noindent\textbf{1.1.}
Unless expressly stated otherwise, from now on we assume that \(G\) is an
\(A\)-group locally of finite type. When \(A\) is a field \(k\), we shall
obtain in Expos\'e~VII\textsubscript{B} very precise results on the local
rings of \(G\).%
\footnote{\label{note:sga3-VIA-10}Editor's note: for example, they are
always complete intersections; see Expos\'e~VII\textsubscript{B}, \SGARefLink{sga3:VIIA:5.5.1}{5.5.1}.
Moreover, if \(\operatorname{char}(k)=0\), then \(G\) is smooth
(Expos\'e~VI\textsubscript{B}, \SGARefLink{cor:sga3-VIB-1.6.1}{1.6.1}; see also
Expos\'e~VII\textsubscript{B}, \SGARefLink{sga3:VIIA:3.3.1}{3.3.1}).}
Here we content ourselves with a few elementary results.

\medskip
\phantomsection\label{prop:sga3-VIA-1.1.1}%
\noindent\textbf{Proposition 1.1.1.}
\textit{Let \(x\) be a point of an \(A\)-group \(G\) locally of finite type
and flat over \(A\). Then the local ring \(\mathcal O_{G,x}\) is
Cohen--Macaulay, and there is a system of parameters
\(a_1,\ldots,a_n\) of \(\mathcal O_{G,x}\) such that
\(\mathcal O_{G,x}/(a_1,\ldots,a_n)\) is a finite flat \(A\)-module
(hence finite free).}

\medskip
We first suppose that \(A\) is equal to its residue field \(k\). It then
suffices to prove that \(\mathcal O_{G,x}\) is Cohen--Macaulay, and we may
restrict to the case in which \(x\) is a closed point
(cf.\ EGA~\(0_{\mathrm{IV}}\), 16.5.13). By
\hyperref[lem:sga3-VIA-1.1.2]{Lemma~1.1.2} below, \(G\) contains a closed
point \(y\) such that \(\mathcal O_{G,y}\) is Cohen--Macaulay. By
SGA~1, I, \S~9, this is equivalent to saying that, for every finite
extension \(K\) of the base field \(k\) and every point \(\overline y\) of
\(\overline G=G\otimes_k K\) lying above \(y\), the local ring
\(\mathcal O_{\overline G,\overline y}\) is Cohen--Macaulay. If \(K\) has
been chosen sufficiently large, that is, if \(K\) contains a normal
extension of \(k\) containing the residue fields \(\kappa(y)\) and
\(\kappa(x)\), then \(\overline y\) is (strictly) rational over \(K\), as
is every point \(\overline x\) of \(\overline G\) lying above \(x\).%
\footnote{\label{note:sga3-VIA-11}Editor's note: indeed, the hypothesis on
\(K\) implies that, for every extension \(L\) of \(K\), every
\(k\)-morphism \(\kappa(x)\to L\) (respectively,
\(\kappa(y)\to L\)) factors through \(K\). Consequently, all the points of
\(G\otimes_k K\) lying above \(x\) or \(y\) have \(K\) as residue field;
that is, they are (strictly) rational over \(K\).}
Since the automorphism
\[
  r_{\overline x}\circ(r_{\overline y})^{-1}
\]
sends \(\overline y\) to \(\overline x\), it follows that
\(\mathcal O_{\overline G,\overline x}\), and hence
\(\mathcal O_{G,x}\) (SGA~1, I, \S~9), are Cohen--Macaulay.

Now let \(A\) again be arbitrary. The preceding argument applies to
\(k\otimes_A G\), so that \(k\otimes_A\mathcal O_{G,x}\) is
Cohen--Macaulay. If \(a_1,\ldots,a_n\) is a sequence of elements of
\(\mathcal O_{G,x}\) whose image in
\(k\otimes_A\mathcal O_{G,x}\) is a system of parameters, it follows from
SGA~1, IV, 5.7, or EGA~\(0_{\mathrm{IV}}\), 15.1.16, that
\(a_1,\ldots,a_n\) is an \(\mathcal O_{G,x}\)-regular sequence and that
\(\mathcal O_{G,x}/(a_1,\ldots,a_n)\) is finite and flat (hence finite and
free) over \(A\).

\medskip
\phantomsection\label{lem:sga3-VIA-1.1.2}%
\noindent\textbf{Lemma 1.1.2.}
\textit{Every nonempty scheme \(X\), locally of finite type over an
Artinian ring \(A\), contains a closed point \(x\) whose local ring is
Cohen--Macaulay.}

\medskip
We may plainly suppose that \(X\) is affine, with algebra \(B\), and argue
by induction on \(\dim X\) (the assertion is clear if \(X\) is discrete,
since all its local rings are then Artinian). Since \(B\) is of finite type
over \(A\), if \(\dim B>0\), then \(B\) contains an element \(a\) that is
not invertible and is not a zero divisor.%
\footnote{\label{note:sga3-VIA-12}Editor's note: indeed, \(B\) is a
Noetherian Jacobson ring (cf.\ \textup{[BAC]}, V, \S~3.4). If every
noninvertible element is a zero divisor, then every prime ideal is an
associated prime ideal of \(B\); in particular, \(B\) has only finitely
many maximal ideals \(\mathfrak m_1,\ldots,\mathfrak m_n\). Since \(B\)
is a Jacobson ring, the intersection of the \(\mathfrak m_i\) is the
nilradical of \(B\). It follows that each \(\mathfrak m_i\) is a minimal
prime ideal of \(B\), and hence that \(\dim B=0\).}
The closed subscheme
\[
  X'=\operatorname{Spec} B/(a)
\]
of \(X\) then has dimension strictly less than \(\dim X\), and by the
induction hypothesis it contains a closed point \(x\) such that
\(\mathcal O_{X',x}\) is Cohen--Macaulay. Since
\(\mathcal O_{X',x}=\mathcal O_{X,x}/(a)\), and \(a\) is neither invertible
nor a zero divisor in \(\mathcal O_{X,x}\), the ring
\(\mathcal O_{X,x}\) is Cohen--Macaulay (see also EGA~IV\textsubscript{2},
6.11.3).

\medskip
\phantomsection\label{prop:sga3-VIA-1.2}%
\noindent\textbf{Proposition 1.2.}
\textit{Let \(A\) be an Artinian local ring, let \(G\) be an \(A\)-group
locally of finite type and flat over \(A\), and let \(x\) be a closed point
of \(G\). Then there exists a local \(A\)-algebra \(A'\), finite and free
over \(A\), such that every point of \(G\otimes_A A'\) lying above \(x\)
is strictly rational over \(A'\).}

\medskip
\noindent\footnote{\label{note:sga3-VIA-13}Editor's note: the original has
been expanded in what follows.}
Indeed, let \(k_1\) be a finite normal extension of \(k\) containing the
residue field \(\kappa(x)\) of \(x\). By \SGARefLink{lem:sga3-V-4.1.2}{Lemma~V.4.1.2}, there exists a
local \(A\)-algebra \(A_1\), finite and free over \(A\), with residue field
\(k_1\). In this case (cf.\ Editor's note~\ref{note:sga3-VIA-11}), all the
points \(g_1,\ldots,g_n\) of \(G\otimes_A A_1\) lying above \(x\in G\)
have \(k_1\) as residue field; that is, \(g_1,\ldots,g_n\) are rational
over \(A_1\) in the sense of Expos\'e~V, \S~4.e).

Let \(B_1,\ldots,B_n\) be the local rings of \(g_1,\ldots,g_n\).
By \hyperref[prop:sga3-VIA-1.1.1]{1.1.1}, the rings
\(B_1,\ldots,B_n\) have quotients \(B'_1,\ldots,B'_n\) that are Artinian
and finite and free over \(A_1\). Set
\[
  A'=B'_1\otimes_{A_1}\cdots\otimes_{A_1}B'_n.
\]
Then \(A'\) is local, finite, and free over \(A_1\), and, for each
\(i=1,\ldots,n\), there are surjective homomorphisms
\phantomsection\label{eq:sga3-VIA-1.2-surjections}%
\[
  B_i\otimes_{A_1}A'
    \twoheadrightarrow B'_i\otimes_{A_1}A'
    \twoheadrightarrow A',
\]
the second induced by the multiplication map
\[
  B'_i\otimes_{A_1}B'_i\twoheadrightarrow B'_i.
\]
Consequently, \(A'\) has the required property.

\medskip
\phantomsection\label{par:sga3-VIA-1.3}%
\noindent\textbf{1.3.}
Let \(e\) be the unit element (or origin) of \(G\), that is, the image of
the unique point of \(\operatorname{Spec}A\) under the unit section
\(\operatorname{Spec}A\to G\). By definition, \(e\) is strictly rational
over \(A\).

\par\medskip
\noindent\phantomsection\label{prop:sga3-VIA-1.3.1}%
\textbf{Proposition 1.3.1.}%
\footnote{\label{note:sga3-VIA-14}Editor's note: \hyperref[concl:sga3-VIA-1.3.1-1]{part~(1)}, which is useful
in \hyperref[ex:sga3-VIA-1.3.2]{Examples~1.3.2} below, has been added.}
\textit{Let \(G\) be a group locally of finite type and flat over an
Artinian ring \(A\), and let \(K\) (respectively, \(\overline k\)) be the
perfect closure (respectively, an algebraic closure) of the residue field
\(k\) of \(A\).}

\begin{enumerate}
\item[(1)]\phantomsection\label{concl:sga3-VIA-1.3.1-1} \textit{For every closed point \(x\) of
\(\overline G=G\otimes_A\overline k\),\footnote{Canon source note
(SGA-SRC-000859): before the explicit reduction \(A=k\), the re-edition
prints \(k\) as the tensor base here and again in the proof. Both objects
are defined over the Artinian ring \(A\); the later post-reduction tensor
bases remain \(k\).} the local rings
\(\mathcal O_{\overline G,e}\) and \(\mathcal O_{\overline G,x}\) are
isomorphic. In particular, the tangent spaces \(T_e\overline G\) and
\(T_x\overline G\) have the same dimension.}

\item[\SGAInlineRefTarget{prop:sga3-VIA-1.3.1:item:2}(2)] \textit{The following assertions are equivalent:}
  \begin{enumerate}
  \item[(i)] \textit{\(G\otimes_A K\) is reduced.}
  \item[(i bis)] \textit{\(\mathcal O_{G,e}\otimes_A K\) is reduced.}
  \item[\SGAInlineRefTarget{prop:sga3-VIA-1.3.1:item:ii}(ii)] \textit{\(G\) is smooth over \(A\).}
  \item[(ii bis)] \textit{\(G\) is smooth over \(A\) at the origin.}
  \end{enumerate}
\end{enumerate}

\noindent\textit{Proof.}
\SGARefLink{concl:sga3-VIA-1.3.1-1}{(1)} Let \(x\) be a closed point of \(\overline G\). There is exactly one
\(\overline k\)-morphism
\[
  s:\operatorname{Spec}\overline k\longrightarrow\overline G
\]
whose image is \(x\). Right translation \(r_s\) then induces an
isomorphism from
\[
  \mathcal O_{\overline G,e}
    =\mathcal O_{G,e}\otimes_A\overline k
\]
onto \(\mathcal O_{\overline G,x}\), which proves~\SGARefLink{concl:sga3-VIA-1.3.1-1}{(1)}.

We now prove~\SGARefLink{prop:sga3-VIA-1.3.1:item:2}{(2)}. By SGA~1, II.2.1, we reduce at once to the case in which
\(A\) is a field (\(A=k\)). The implications
\[
  \text{(i)}\Rightarrow\text{(i bis)},\qquad
  \text{(ii)}\Rightarrow\text{(ii bis)},\qquad
  \text{(ii)}\Rightarrow\text{(i)},\qquad
  \text{(ii bis)}\Rightarrow\text{(i bis)}
\]
are evident, so it remains only to prove that (i bis) implies~\SGARefLink{prop:sga3-VIA-1.3.1:item:ii}{(ii)}.

Now (i bis) implies that
\(\mathcal O_{G,e}\otimes_k\overline k\) is reduced. Hence, by~\SGARefLink{concl:sga3-VIA-1.3.1-1}{(1)},
\(\mathcal O_{\overline G,x}\) is reduced for every closed point \(x\) of
\(\overline G\), and therefore \(\overline G\) is reduced. Since
\(\overline G\) is locally of finite type over \(\overline k\), there is
at least one closed point \(y\) such that
\(\mathcal O_{\overline G,y}\) is regular. By~\SGARefLink{concl:sga3-VIA-1.3.1-1}{(1)}, the local rings at all
closed points of \(\overline G\) are isomorphic to
\(\mathcal O_{\overline G,e}\); thus all these local rings are regular.
Consequently, \(\overline G\) is smooth over \(\overline k\), and hence
\(G\) is smooth over \(k\).

\footnote{\label{note:sga3-VIA-15}Editor's note: the following sentence
and \hyperref[ex:sga3-VIA-1.3.2]{Examples~1.3.2} have been added.}
We may now give the examples below, pointed out by M.~Raynaud, of group
schemes \(G\) over a nonperfect field \(k\) for which \(G_{\mathrm{red}}\)
is not a \(k\)-group scheme.

% END INCLUDED SGA3_English_current_buildable_source_20260823/source/inputs/VI/tex_reference_v2/components/01_expose_VIA_section1_through_prop131_en.tex

% BEGIN INCLUDED SGA3_English_current_buildable_source_20260823/source/inputs/VI/tex_reference_v2/components/02_expose_VIA_examples132_through_cor241_en.tex
% SGA 3, Expose VI_A, Loop 1 English reconstruction.
% Scope: Examples 1.3.2 through the end of Corollary 2.4.1;
% paragraph 2.5 is excluded.
% Authority: Polo--Gille Exp6A-13oct24.pdf, component-local pages 7--11
% (combined-reader pages 310--314; printed pages 299--303).
% The supplied opening locator (local p. 6 / combined p. 309 / printed
% p. 298) is the immediately preceding page; the Examples 1.3.2 heading
% is visibly printed on the next page.
% OCR and the Reinhold English Markdown are comparison witnesses only.

\medskip
\phantomsection\label{ex:sga3-VIA-1.3.2}%
\noindent\textbf{Examples 1.3.2.}
Let \(k\) be a nonperfect field of characteristic \(p>0\), let
\(t\in k-k^p\), let \(\overline k\) be an algebraic closure of \(k\), and
let \(\alpha\in\overline k\) satisfy \(\alpha^p=t\).

\begin{enumerate}
\item[(1)] Consider the additive group
\(\mathbb G_{a,k}=\operatorname{Spec}k[X]\), and let \(G\) be the group
subscheme, finite over \(k\), defined by the additive polynomial
\(X^{p^2}-tX^p\). Then
\[
  G_{\mathrm{red}}
    =\operatorname{Spec}k[X]/X\bigl(X^{p(p-1)}-t\bigr)
\]
is \'etale at the origin. If it were a group scheme over \(k\), it would
be smooth (by \hyperref[prop:sga3-VIA-1.3.1]{1.3.1}); but \(G_{\mathrm{red}}\) is not geometrically
reduced, so it is not a group scheme over \(k\).

\item[(2)] Consider
\[
  \mathbb G^4_{a,k}=\operatorname{Spec}k[X,Y,U,V],
\]
and let \(G\) be the group subscheme defined by the ideal \(I\) generated
by the additive polynomials
\[
  P=X^p-tY^p
  \qquad\text{and}\qquad
  Q=U^p-tV^p.
\]
Then \(G\) is two-dimensional and irreducible, since
\[
  (G_{\overline k})_{\mathrm{red}}
    \simeq\operatorname{Spec}\overline k[Y,V]
\]
is irreducible.

Let \(A=k[X,Y,U,V]\), and let \(\mathfrak m\) be its augmentation ideal.
Denote by \(x,y,u,v\) the images of \(dX,dY,dU,dV\) in
\[
  \Omega^1_{A/k}\otimes_A(A/\mathfrak m),
\]
viewed as linear forms on the tangent space
\[
  k^4=T_0\mathbb G^4_{a,k}.
\]
We shall show that the subspace \(E=T_0G_{\mathrm{red}}\) equals \(k^4\).
Otherwise, there would be a linear form
\[
  f=ax+by+a'u+b'v,
  \qquad a,b,a',b'\in k,
\]
not all of whose coefficients are zero, that vanishes on \(E\). Recall
that the formation of \(\Omega^1_{A/k}\), and hence that of tangent spaces,
commutes with base change (cf.\ EGA~IV\textsubscript{4}, 16.4.5), and
identify \(f\) with its image in \((\overline k^4)^*\). Since
\[
  (G_{\overline k})_{\mathrm{red}}
    \subset (G_{\mathrm{red}})_{\overline k},
\]
the form \(f\) vanishes on the subspace
\(T_0(G_{\overline k})_{\mathrm{red}}\) of \(\overline k^4\), which is
defined by the equations
\[
  g_1=x-\alpha y
  \qquad\text{and}\qquad
  g_2=u-\alpha v.
\]
Thus
\[
  f=\lambda g_1+\mu g_2,
  \qquad \lambda,\mu\in\overline k.
\]
But \(\lambda g_1+\mu g_2\) belongs to \(k^4\) only if
\(\lambda=\mu=0\). This contradiction shows that \(E=k^4\), and hence
\[
  T_0(G_{\mathrm{red}})_{\overline k}=\overline k^4.
\]

On the other hand, \(R=XV-YU\) belongs to \(\sqrt I\), since
\[
  R^p=(X^p-tY^p)V^p-Y^p(U^p-tV^p).
\]
Consequently, the tangent space \(F\) at the point
\((\alpha,1,\alpha,1)\) of \((G_{\mathrm{red}})_{\overline k}\) is
contained in the hyperplane \(H\) of \(\overline k^4\) with equation
\[
  \alpha\,dV+dX-dU-\alpha\,dY=0,
\]
and therefore has dimension at most \(3\).%
\footnote{\label{note:sga3-VIA-16}Editor's note: in fact, it is easy to
see that \(\sqrt I\) is generated by \(P,Q,R\) at every point other than
\(0\), so \(F=H\).}
It follows from \hyperref[concl:sga3-VIA-1.3.1-1]{part~(1)} of \SGARefLink{prop:sga3-VIA-1.3.1}{1.3.1} that \(G_{\mathrm{red}}\) is not a group
scheme over \(k\).
\end{enumerate}

\section*{2. Connected components of an \(A\)-group locally of finite type}
\phantomsection\label{sec:sga3-VIA-2}
\addcontentsline{toc}{section}{2. Connected components of an
  \texorpdfstring{\(A\)}{A}-group locally of finite type}

\medskip
\phantomsection\label{par:sga3-VIA-2.1}%
\noindent\textbf{2.1.}
Let us first consider an arbitrary \(A\)-group \(G\), and let \(G'\) be
the connected component of the origin \(e\) in \(G\). This connected
component is plainly closed, so we may identify it with the reduced closed
subscheme of \(G\) having \(G'\) as its underlying space.%
\footnote{\label{note:sga3-VIA-17}Editor's note: we shall see below (\SGARefLink{par:sga3-VIA-2.2}{2.2})
that \(G'\) is a group in the category
\((\mathbf{Sch}_{/k})_{\mathrm{red}}\).}

\medskip
\phantomsection\label{prop:sga3-VIA-2.1.1}%
\noindent\textbf{Proposition 2.1.1.}
\textit{For every extension \(K\) of the residue field \(k\) of \(A\),
the underlying space of \(G'\otimes_AK\) is the connected component of the
origin in the \(K\)-group \(G\otimes_AK\); in other words, \(G'\) is
geometrically connected.}

\medskip
Indeed, let \((G\otimes_AK)'\) be the connected component of the origin in
\(G\otimes_AK\). Since the image of \((G\otimes_AK)'\) in \(G\) is
connected and contains the unit element of \(G\), that image is contained
in \(G'\). Thus \((G\otimes_AK)'\) is contained in the inverse image
\(G'\otimes_AK\) of \(G'\) in \(G\otimes_AK\). The proposition therefore
follows from the connectedness of \(G'\otimes_AK\), which is proved in the
following lemma.

\medskip
\phantomsection\label{lem:sga3-VIA-2.1.2}%
\noindent\textbf{Lemma 2.1.2.}
\textit{Let \(X\) and \(Y\) be two connected schemes over a field \(k\).
If \(X\) contains a rational point, then \(X\times_kY\) is connected.}

\medskip
We give a direct proof below of this result from EGA~IV\textsubscript{2}
(4.5.8 and 4.5.14).

Suppose first that \(Y\) is nonempty, connected, and affine, with algebra
\(B\). Then \(X\times_kY\) is the spectrum of the quasi-coherent
\(\mathcal O_X\)-algebra
\[
  \mathcal B=\mathcal O_X\otimes_kB.
\]
We wish to show that every nonempty subset \(U\) of \(X\times_kY\) that is
both open and closed equals \(X\times_kY\). Now \(U\) is affine over \(X\),
and its affine \(\mathcal O_X\)-algebra is a direct factor of
\(\mathcal B\). It follows from \hyperref[lem:sga3-VIA-2.1.3]{Lemma~2.1.3} below that the image of \(U\)
in \(X\) is open and closed, and hence is all of \(X\). In particular,
this image contains a rational point \(x\) of \(X\), so \(U\) meets the
inverse image of \(x\) in \(X\times_kY\). Since this inverse image is
isomorphic to \(Y\), and hence connected, \(U\) contains the whole inverse
image. The same conclusion would hold for the complement of \(U\) in
\(X\times_kY\) if \(U\) were different from \(X\times_kY\), which is
absurd.

If \(Y\) is now an arbitrary \(k\)-scheme, the preceding argument shows
that the fibers of the canonical projection
\[
  X\times_kY\longrightarrow Y
\]
are connected. If \(x\) is a rational point of \(X\), all these fibers
meet the subscheme \(\{x\}\times_kY\), which is itself connected. This
proves the assertion.

\medskip
\phantomsection\label{lem:sga3-VIA-2.1.3}%
\noindent\textbf{Lemma 2.1.3.}
\textit{Let \(X\) be a scheme, and let \(\mathcal A\) be a
quasi-coherent \(\mathcal O_X\)-algebra that is locally%
\footnote{\label{note:sga3-VIA-18}Editor's note: the word ``locally'' has
been added. The reference to Kaplansky is \emph{Projective modules},
Ann.\ of Math.\ \textbf{68} (1958), 372--377; see also
\textup{[BAC]}, \S~II.3, Ex.\ 3.}
a direct factor of a free \(\mathcal O_X\)-module. Then the image of
\(\operatorname{Spec}\mathcal A\) in \(X\) is open and closed.}

\medskip
Let \(V\) be this image. It is clear that \(V\) is contained in the
support of \(\mathcal A\). Conversely, if \(x\) belongs to the support of
\(\mathcal A\), then \(\mathcal A_x\) is nonzero and is a free
\(\mathcal O_{X,x}\)-module, since, by Kaplansky, every projective module
over a local ring is free. Consequently, the fiber of
\(\operatorname{Spec}\mathcal A\) at \(x\), which is affine with algebra
\[
  \mathcal A_x\otimes_{\mathcal O_{X,x}}\kappa(x),
\]
is nonempty. Thus the image of \(\operatorname{Spec}\mathcal A\)
coincides with the support of \(\mathcal A\).

If \(s\) is the unit section of \(\mathcal A\), the equality \(s_x=0\)
implies that \(s\), and hence \(\mathcal A\), vanish in a neighborhood of
\(x\). Therefore the support of \(\mathcal A\) is closed. On the other
hand, we may assume that \(\mathcal A\) is a direct factor of the direct
sum \(\mathcal L\) of a family
\((\mathcal O_X^\alpha)\) of copies of \(\mathcal O_X\).%
\footnote{\label{note:sga3-VIA-19}Editor's note: the original has been
expanded in what follows.}
For each \(\beta\), let \(\varphi^\beta\) denote the restriction to
\(\mathcal A\) of the canonical projection
\[
  \mathcal L=\bigoplus_\alpha\mathcal O_X^\alpha
    \longrightarrow\mathcal O_X^\beta.
\]
For \(x\in X\), let \(\mathcal A_x^*\) denote the dual module
\[
  \operatorname{Hom}_{\mathcal O_{X,x}}
    (\mathcal A_x,\mathcal O_{X,x}).
\]
Since \(\mathcal A_x\) is a direct factor of \(\mathcal L_x\), the
restriction map
\[
  \operatorname{Hom}_{\mathcal O_{X,x}}
    (\mathcal L_x,\mathcal O_{X,x})
  \longrightarrow
  \operatorname{Hom}_{\mathcal O_{X,x}}
    (\mathcal A_x,\mathcal O_{X,x})
  =\mathcal A_x^*
\]
is surjective. If \(\mathcal A_x\ne0\), then, since \(\mathcal A_x\) is
free and nonzero, there exist \(a\in\mathcal A_x\) and a linear form
\(\varphi\in\mathcal A_x^*\) such that \(\varphi(a)=1\). There is
therefore a family \((\xi^\alpha)\) of elements of
\(\mathcal O_{X,x}\) such that, for every
\(u\in\mathcal A_{X,x}\),
\[
  \varphi(u)=\sum_\alpha\xi^\alpha\varphi^\alpha(u);
\]
this sum is finite because \(\varphi^\alpha(u)=0\) except for finitely many
\(\alpha\). Applying this to \(u=a\), we find
\(\alpha_1,\ldots,\alpha_n\) such that
\[
  1=\varphi(a)
   =\xi^{\alpha_1}\varphi^{\alpha_1}(a)
    +\cdots+
    \xi^{\alpha_n}\varphi^{\alpha_n}(a).
\]
There is then an open neighborhood \(U\) of \(x\) on which \(a\) and the
\(\xi^{\alpha_i}\) come from sections
\[
  \widetilde a\in\mathcal A(U)
  \qquad\text{and}\qquad
  \widetilde\xi^{\alpha_i}\in\mathcal O_X(U).
\]
The equality
\[
  \sum_i
    \widetilde\xi^{\alpha_i}
    \varphi^{\alpha_i}(\widetilde a)=1
  \quad\text{on }U
\]
shows that \(\widetilde a_y\ne0\) for every \(y\in U\).
(``The support of a projective module is open.'')

\medskip
\phantomsection\label{par:sga3-VIA-2.2}%
\noindent\textbf{2.2.}
The notation remains that of \hyperref[par:sga3-VIA-2.1]{2.1}. It is clear that \(G'\) is a reduced
\(k\)-scheme. \hyperref[lem:sga3-VIA-2.1.2]{Lemma~2.1.2} shows that \(G'\times_kG'\) is connected, so
\((G'\times_kG')_{\mathrm{red}}\) is the reduced subscheme of
\(G\times_AG\) whose underlying space is the connected component of the
origin. In particular, the multiplication morphism
\[
  \mu:G\times_AG\longrightarrow G
\]
induces a morphism
\[
  \mu':(G'\times_kG')_{\mathrm{red}}\longrightarrow G'
\]
that makes \(G'\) a group in
\((\mathbf{Sch}_{/k})_{\mathrm{red}}\).

\medskip
\phantomsection\label{par:sga3-VIA-2.2bis}%
\noindent\textbf{2.2.bis.}%
\footnote{\label{note:sga3-VIA-20}Editor's note: this paragraph has been
added in order to make the connection with
Expos\'e~VI\textsubscript{B}, \S~3.}
Recall (cf.\ Expos\'e~V) that, if \(P\) is a scheme, \(\underline P\)
denotes the underlying topological space of \(P\). We define an
\(A\)-subfunctor \(G^0\) of \(G\) by setting, for every \(A\)-scheme
\(S\),
\[
  G^0(S)
    =\{u\in G(S)\mid u(\underline S)\subset\underline{G'}\}.
\]
Let \(c:G\to G\) be the inversion morphism. Since
\(c(\underline{G'})=\underline{G'}\), we have
\(c\circ u\in G^0(S)\) for every \(u\in G^0(S)\). On the other hand, if
\(u,v\in G^0(S)\), then \(u\boxtimes v\) sends \(\underline S\) into the
subspace of \(\underline{G\times_AG}\) consisting of the points whose two
projections belong to \(\underline{G'}\). This subspace is identified with
the underlying space of \(G'\times_AG'\), which is connected by
\hyperref[lem:sga3-VIA-2.1.2]{Lemma~2.1.2}. Consequently, \(\mu\circ(u\boxtimes v)\) sends
\(\underline S\) into \(\underline{G'}\). This proves that \(G^0\) is a
group subfunctor of \(G\) over \(A\).

If the connected component of \(e\) is open in \(G\), then the subfunctor
\(G^0\) is represented by the subscheme induced by \(G\) on this open
subset, which is therefore a group subscheme of \(G\); this subscheme will
also be denoted by \(G^0\). In this case, with the notation of \SGARefLink{prop:sga3-VIB-2.1}{2.1},
\[
  G'=(G^0)_{\mathrm{red}},
\]
and the topological spaces \(\underline{G'}\) and \(\underline{G^0}\)
coincide.%
\footnote{\label{note:sga3-VIA-21}Editor's note: in this expos\'e, the
notation \(G^0\) is reserved for the case in which the connected component
of \(e\) is open; in Expos\'e~VI\textsubscript{B}, \S~3, this connected
component will be denoted by \(\underline{G^0}\) in every case. This is a
slight abuse of notation, but it is compatible with the preceding
discussion when the connected component is the underlying topological
space of an open group subscheme \(G^0\) of \(G\).}

\medskip
\phantomsection\label{par:sga3-VIA-2.3}%
\noindent\textbf{2.3.}
In accordance with our convention in \SGARefLink{par:sga3-VIA-1.1}{1.1}, we again assume from now on that
\(G\) is locally of finite type over \(A\). Then \(G\) is locally
Noetherian, hence locally connected;%
\footnote{\label{note:sga3-VIA-22}Editor's note: indeed, in a Noetherian
space there are only finitely many connected components, and each is
therefore open; see also EGA~I, 6.1.9.}
thus every connected component of \(G\) is open.

We shall write \(G^0\) for the subscheme induced by \(G\) on the connected
component \(G'\) of the unit element. By \hyperref[par:sga3-VIA-2.2bis]{2.2.bis}, \(G^0\) is a group
subscheme of \(G\), which we shall call the \emph{neutral component} of
\(G\). For every \(A\)-scheme \(S\), therefore,
\[
  G^0(S)
    =\{u\in G(S)\mid
      u(\underline S)\subset\underline{G^0}=\underline{G'}\}.
\]

Let \(G^\alpha\) be an arbitrary connected component of \(G\), and let
\[
  \nu^\alpha:G^\alpha\times_AG^0\longrightarrow G
\]
be the morphism defined by
\[
  \nu^\alpha(S)(g,\gamma)=g\gamma g^{-1}
\]
for every \(S\in(\mathbf{Sch}_{/A})\), \(g\in G^\alpha(S)\), and
\(\gamma\in G^0(S)\). If \(e\) is the origin of \(G\), the restriction of
\(\nu^\alpha\) to \(G^\alpha\times_A\{e\}\) is the trivial morphism.
Since \(G^\alpha\times_AG^0\) is connected by \hyperref[lem:sga3-VIA-2.1.2]{2.1.2}, it follows that
\(\nu^\alpha\) factors through \(G^0\). Thus, for every \(A\)-scheme
\(S\), \(G^0(S)\) is a normal subgroup of \(G(S)\). We have proved the
following proposition.%
\footnote{\label{note:sga3-VIA-23}Editor's note: the number \SGARefLink{prop:sga3-VIA-2.3.1}{2.3.1} has
been added in order to emphasize this statement. Moreover, \(G^0\) is
even a characteristic subgroup of \(G\); see \SGARefLink{thm:sga3-VIA-2.6.5:item:ii}{2.6.5(ii)}.}

\medskip
\phantomsection\label{prop:sga3-VIA-2.3.1}%
\noindent\textbf{Proposition 2.3.1.}
\textit{Let \(G\) be an \(A\)-group locally of finite type. Then the
neutral component \(G^0\) is an open normal group subscheme of \(G\).}

\medskip
\phantomsection\label{prop:sga3-VIA-2.4}%
\noindent\textbf{Proposition 2.4.}
\textit{Let \(G\) be an \(A\)-group locally of finite type.}
\begin{enumerate}
\item[\SGAInlineRefTarget{prop:sga3-VIA-2.4:item:i}(i)] \textit{\(G^0\) is irreducible, and \(G^0\otimes_Ak\) is
geometrically irreducible over \(k\).}
\item[\SGAInlineRefTarget{prop:sga3-VIA-2.4:item:ii}(ii)] \textit{\(G^0\) is quasi-compact, and hence is of finite type
over \(A\).}
\end{enumerate}

\noindent\footnote{\label{note:sga3-VIA-24}Editor's note: in the
statement, ``\(G^0\) is geometrically irreducible'' has been replaced by
``\(G^0\otimes_Ak\) is geometrically irreducible over \(k\),'' and the
proof has been expanded.}
\textit{Proof.}
\SGARefLink{prop:sga3-VIA-2.4:item:i}{(i)} Since \(G^0\) and \(G^0\otimes_Ak\) have the same underlying
topological space, it suffices to prove the second assertion. Let
\(\overline k\) be an algebraic closure of \(k\). By \SGARefLink{prop:sga3-VIB-2.2}{2.2},
\((G\otimes_A\overline k)_{\mathrm{red}}\) is a reduced
\(\overline k\)-group locally of finite type, and hence is smooth over
\(\overline k\) by \SGARefLink{prop:sga3-VIA-1.3.1}{1.3.1}. A fortiori, the local rings of
\((G\otimes_A\overline k)_{\mathrm{red}}\) are integral domains.
Therefore,%
\footnote{\label{note:sga3-VIA-25}Editor's note: the original has been
expanded here, and the reference to EGA~I, 6.1.10, has been added.}
since \(G\otimes_A\overline k\) is locally Noetherian, the connected
components of \(G\otimes_A\overline k\) are irreducible
(cf.\ EGA~I, 6.1.10). In particular, the connected component
\(G^0\otimes_A\overline k\) (cf.\ 2.1.1) is irreducible.

\SGARefLink{prop:sga3-VIA-2.4:item:ii}{(ii)} We now show that \(G^0\) is of finite type over \(A\). Since \(G^0\)
is locally of finite type over \(A\), it suffices to prove that \(G^0\) is
quasi-compact. Since \(G^0\) is irreducible, this follows from \SGARefLink{sga3:VIA:0.5.1}{0.5.1}.

\medskip
\phantomsection\label{cor:sga3-VIA-2.4.1}%
\noindent\textbf{Corollary 2.4.1.}
\textit{Every connected component of \(G\) is irreducible,%
\footnote{\label{note:sga3-VIA-26}Editor's note: one should bear in mind
that a nonneutral connected component is not geometrically connected in
general. For example, if \(k=\mathbb R\), the group
\(\mu_{3,\mathbb R}\), represented by
\(\mathbb R[X]/(X^3-1)\), has two connected components:
\(\{e\}=\operatorname{Spec}\mathbb R\) and
\(C=\operatorname{Spec}\mathbb R[X]/(X^2+X+1)\); and
\(C\otimes_{\mathbb R}\mathbb C\) has two components.}
of finite type over \(A\),%
\footnote{\label{note:sga3-VIA-27}Editor's note: the assertion that
follows has been added; cf.\ Expos\'e~VI\textsubscript{B}, \SGARefLink{lem:sga3-VIB-1.5}{1.5}.}
and of the same dimension as \(G^0\).}

\medskip
Indeed, we may assume that \(A\) equals its residue field \(k\). Let \(C\)
be a connected component of \(G\), let \(x\) be a closed point of \(C\),
let \(\kappa(x)\) be the residue field of \(x\), and let \(k'\) be a
finite normal extension of \(k\) containing \(\kappa(x)\). The canonical
projection
\[
  \pi:C\otimes_kk'\longrightarrow C
\]
is open and closed.%
\footnote{\label{note:sga3-VIA-28}Editor's note: the original has been
expanded in what follows.}
Consequently, if \(C'\) is a connected component of \(C\otimes_kk'\), the
projection \(C'\to C\) is surjective. Thus \(C'\) contains a point
\(y\in\pi^{-1}(x)\), and such a point is rational over \(k'\)
(cf.\ the proof of \hyperref[prop:sga3-VIA-1.2]{1.2}); it follows that \(C'\) is the image of
\(G^0\otimes_kk'\) under right translation \(r_y\). Now
\(G^0\otimes_kk'\) is of finite type over \(k'\), by \hyperref[prop:sga3-VIA-2.4]{2.4}, and
\(\pi^{-1}(x)\) is finite, of cardinality at most \([k':k]\).
Consequently, \(C\otimes_kk'\) is of finite type over \(k'\), and hence
\(C\) is of finite type over \(k\).

Moreover, since \(G^0\otimes_kk'\) is irreducible by \SGARefLink{prop:sga3-VIA-2.4}{2.4}, the same is true
of \(C'\),%
\footnote{\label{note:sga3-VIA-29}Editor's note: the original has been
simplified here.}
and hence also of \(C\), because the projection \(C'\to C\) is
surjective.

Finally, we saw above that \(C\otimes_kk'\) is the disjoint union of
finitely many translates of \(G^0\otimes_kk'\). Since dimension is
invariant under extension of the base field (cf.\
EGA~IV\textsubscript{2}, 4.1.4), it follows that \(C\) has the same
dimension as \(G^0\). (Moreover, by EGA~IV\textsubscript{2}, 5.2.1, one
has \(\dim_gG=\dim G^0\) for every point \(g\in G\).)

% END INCLUDED SGA3_English_current_buildable_source_20260823/source/inputs/VI/tex_reference_v2/components/02_expose_VIA_examples132_through_cor241_en.tex

% BEGIN INCLUDED SGA3_English_current_buildable_source_20260823/source/inputs/VI/tex_reference_v2/components/03_expose_VIA_section25_lemma251_prop252_en.tex
% SGA 3, Expose VI_A, Loop 1 English reconstruction.
% Scope: paragraph 2.5 through the end of Proposition 2.5.2;
% Lemma 2.5.3 is excluded.
% Authority: Polo--Gille Exp6A-13oct24.pdf, component-local pages 11--13
% (combined-reader pages 314--316; printed pages 303--305).
% OCR and the Reinhold English Markdown are comparison witnesses only.

\medskip
\phantomsection\label{sec:sga3-VIA-2.5}%
\noindent\SGAInlineRefTarget{sga3:vi:paragraph:2-5}\textbf{2.5.}
This paragraph has been added. The results that follow appear in
Expos\'e~VI\textsubscript{B}, but could (or should) have appeared in
VI\textsubscript{A}; it is useful to have them available already, in
order to make \hyperref[thm:sga3-VIA-3.2]{Theorem~3.2} below more precise.

\medskip
\phantomsection\label{lem:sga3-VIA-2.5.1}%
\noindent\textbf{Lemma 2.5.1.}
\textit{Let \((A,\mathfrak m)\) be an Artinian local ring, and let
\(k=A/\mathfrak m\) be its residue field.}
\begin{enumerate}
\item[\SGAInlineRefTarget{lem:sga3-VIA-2.5.1:item:i}(i)] \textit{If \(X\) is an \(A\)-scheme such that
\(X\otimes_Ak\) is locally of finite type (respectively, of finite type)
over \(k\), then \(X\) has the same property over \(A\).}
\item[\SGAInlineRefTarget{lem:sga3-VIA-2.5.1:item:ii}(ii)] \textit{Let \(u:X\to Y\) be a morphism of \(A\)-schemes. If
\(u\otimes_Ak\) is an immersion (respectively, a closed immersion), then
so is \(u\).}
\end{enumerate}

\noindent\textit{Proof.}
\SGARefLink{lem:sga3-VIA-2.5.1:item:i}{(i)} Suppose that \(X\otimes_Ak\) is locally of finite type over \(k\).
Let \(U=\operatorname{Spec}B\) be an affine open subset of \(X\). By
hypothesis, there are elements \(x_1,\ldots,x_n\) of \(B\) whose images
generate \(B/\mathfrak mB\) as a \(k\)-algebra, and the ``nilpotent
Nakayama lemma'' then shows that the \(x_i\) generate \(B\) as an
\(A\)-algebra. This proves that \(X\) is locally of finite type over
\(A\). If, in addition, \(X\otimes_Ak\) is quasi-compact, then so is
\(X\), which has the same underlying topological space; hence \(X\) is
of finite type over \(A\). This proves \SGARefLink{lem:sga3-VIA-2.5.1:item:i}{(i)}.

We prove \SGARefLink{lem:sga3-VIA-2.5.1:item:ii}{(ii)}. Suppose that \(u\otimes_Ak\) is an immersion
(respectively, a closed immersion). Then \(u\) is a homeomorphism of
\(X\) onto a locally closed (respectively, closed) subset of \(Y\), and,
for every \(x\in X\), the ring homomorphism
\[
  \varphi_x:\mathcal O_{Y,u(x)}\longrightarrow\mathcal O_{X,x}
\]
is such that \(\varphi_x\otimes_Ak\) is surjective. The nilpotent
Nakayama lemma therefore shows that \(\varphi_x\) is surjective, so
\(u\) is an immersion (respectively, a closed immersion).

\medskip
\phantomsection\label{prop:sga3-VIA-2.5.2}%
\noindent\textbf{Proposition 2.5.2.}
\textit{Let \(A\) be an Artinian local ring with residue field \(k\), and
let \(u:G\to H\) be a quasi-compact morphism between \(A\)-groups locally
of finite type.}
\begin{enumerate}
\item[\SGAInlineRefTarget{prop:sga3-VIA-2.5.2:item:a}(a)] \textit{The set \(u(G)\) is closed in \(H\), and its connected
components are irreducible and all have the same dimension.}
\item[\SGAInlineRefTarget{prop:sga3-VIA-2.5.2:item:b}(b)] \textit{One has}
\[
  \dim G=\dim u(G)+\dim\operatorname{Ker}(u).
\]
\item[(c)] \textit{If \(u\) is a monomorphism, then it is a closed
immersion.}
\end{enumerate}

\noindent\textit{Proof.}
By \hyperref[lem:sga3-VIA-2.5.1]{the preceding lemma}, it suffices to
prove the proposition when \(A=k\). Moreover, the properties under
consideration are stable under fpqc descent, and dimension is invariant
under extension of the base field, so we may assume that \(k\) is
algebraically closed.

We prove \SGARefLink{prop:sga3-VIA-2.5.2:item:a}{(a)}. Let \(C\) be the reduced subscheme of \(H\) whose
underlying topological space is \(u(G)\). Since \(u(G)\) is stable under
the inversion morphism of \(H\), so is \(C\). Furthermore,
\(u:G\to C\) is quasi-compact and dominant. Thus, by
EGA~IV\textsubscript{2}, 2.3.7, the same is true of
\[
  u\times_k\operatorname{id}_G
  \qquad\text{and}\qquad
  \operatorname{id}_H\times_ku,
\]
and hence of their composite
\[
  u\times_ku:G\times_kG\longrightarrow C\times_kC.
\]
Consequently, multiplication on \(H\) sends \(C\times_kC\) into \(C\),
and therefore \(C\) is a group subscheme of \(H\).

Replacing \(H\) by \(C\), we are thus reduced to the case in which \(u\)
is dominant. Then \(G(k)\) is dense in \(H\), and hence meets every
connected component of \(H\); it therefore acts transitively on the set
of those connected components. It is consequently enough to show that
\(u(G)\) contains \(H^0\). Replacing \(G\) by \(u^{-1}(H^0)\), we may
therefore suppose that \(H=H^0\). In that case,
\hyperref[prop:sga3-VIA-2.4]{2.4} shows that \(H\) is irreducible and of
finite type over \(k\), and hence Noetherian. Moreover, \(u\) is locally
of finite type (cf.\ EGA~I, 6.6.6) and quasi-compact, and is therefore of
finite type. By Chevalley's constructibility theorem
(cf.\ EGA~IV\textsubscript{1}, 1.8.5), \(u(G)\) is consequently a
constructible (and dense) subset of
\[
  H=\overline{u(G)},
\]
and thus contains a dense open subset \(U\) of \(H\)
(cf.\ EGA~0\textsubscript{III}, 9.2.2). By
\hyperref[sga3:VIA:0.5]{0.5}, we then have
\[
  H=U\cdot U\subset u(G),
\]
whence \(u(G)=H\). Together with
\hyperref[cor:sga3-VIA-2.4.1]{2.4.1}, this proves \SGARefLink{prop:sga3-VIA-2.5.2:item:a}{(a)}.

We prove \SGARefLink{prop:sga3-VIA-2.5.2:item:b}{(b)}. Recall first that the functor
\(\operatorname{Ker}(u)\) (cf.\ I, \SGARefLink{sga3:i:definition:2-3-6-1}{2.3.6.1}) is represented by
\(u^{-1}(e)\), where \(e\) denotes the identity element of \(H\). Since
\(u\) is of finite type, \(\operatorname{Ker}(u)\) is of finite type
over \(k\). On the other hand, replacing \(H\) by the reduced closed
subscheme \(u(G)\), we may assume that \(u\) is surjective. Let \(u^0\)
denote the restriction of \(u\) to \(G^0\). Since \(G\) and
\(\operatorname{Ker}(u)\) are equidimensional, and since
\[
  \operatorname{Ker}(u)^0\subset\operatorname{Ker}(u^0),
\]
we are reduced to the case in which \(G\), and hence also \(H\), is
irreducible.

By EGA~IV\textsubscript{3}, 9.2.6.2 and 10.6.1(ii), the set of
\(y\in H\) such that
\[
  \dim u^{-1}(y)=\dim G-\dim H
\]
contains a nonempty open subset \(V\). Since \(u\) is surjective,
\(U=u^{-1}(V)\) is then a nonempty open subset of \(G\), and therefore
contains a closed point \(x\) of \(G\), because \(G\) is a Jacobson
scheme (cf.\ EGA~IV\textsubscript{3}, 10.4.8). Right translation \(r_x\)
is an isomorphism from \(\operatorname{Ker}(u)\) onto
\(u^{-1}(u(x))\), whence
\[
  \dim\operatorname{Ker}(u)
    =\dim u^{-1}(u(x))
    =\dim G-\dim H.
\]

We prove (c), following \textup{[DG70]}, I, \S~3.4. (Another proof is
given in Expos\'e~VI\textsubscript{B}, 1.4.2.) Suppose that \(u\) is a
monomorphism. If \(C\) is a connected component of \(G\), there is a
closed point \(x\in G\) such that \(C=r_x(G^0)\). If \(u_C\)
(respectively, \(u^0\)) denotes the restriction of \(u\) to \(C\)
(respectively, to \(G^0\)), then
\[
  u_C=r_{u(x)}\circ u^0\circ r_x^{-1}.
\]
It therefore suffices to show that \(u^0\) is a closed immersion. We may
thus suppose that \(G=G^0\), so that \(G\) is irreducible and of finite
type over \(k\).

Let \(\xi\) be the generic point of \(G\). Then
\(\mathcal O_{G,\xi}\) is an Artinian local ring; denote its maximal
ideal by \(\mathfrak m\). Let \(h=u(\xi)\), let \(\mathfrak n\) be the
maximal ideal of \(\mathcal O_{H,h}\), and set
\[
  A=\mathcal O_{G,\xi}/\mathfrak n\mathcal O_{G,\xi}.
\]
Since \(u\) is a monomorphism, so is the morphism
\[
  u_h:\operatorname{Spec}(A)\longrightarrow
      \operatorname{Spec}(\kappa(h))
\]
obtained by base change. The multiplication morphism
\[
  A\otimes_{\kappa(h)}A\longrightarrow A
\]
is therefore an isomorphism (cf.\ EGA~I, 5.3.8), whence
\(A=\kappa(h)\). Nakayama's lemma then shows that
\[
  \mathcal O_{H,h}\longrightarrow\mathcal O_{G,\xi}
\]
is surjective, since
\(\mathfrak n\mathcal O_{G,\xi}\) is contained in \(\mathfrak m\) and
is therefore nilpotent.

Let \(V\) be an affine open subset of \(H\) containing \(h\), let
\(U\ne\varnothing\) be an affine open subset of \(G\) contained in
\(u^{-1}(V)\), and let
\[
  \varphi:\mathcal O_H(V)\longrightarrow\mathcal O_G(U)
\]
be the homomorphism of \(k\)-algebras induced by \(u\). Let
\(\mathfrak p\) be the prime ideal of \(\mathcal O_G(U)\) corresponding
to \(\xi\),\begingroup
\renewcommand{\thefootnote}{\fnsymbol{footnote}}%
\footnote[1]{Translator's note: the French source prints \(\eta\) here;
the generic point introduced immediately above is \(\xi\).}%
\endgroup
and set
\(\mathfrak q=\varphi^{-1}(\mathfrak p)\). Since
\(G\) is of finite type over \(k\), the \(k\)-algebra
\(\mathcal O_G(U)\) is generated by finitely many elements
\(a_1,\ldots,a_n\). By the preceding argument, there are elements
\(b_1,\ldots,b_n,s\in\mathcal O_H(V)\), with
\(s\notin\mathfrak q\), such that in
\(\mathcal O_G(U)_{\mathfrak p}\) one has
\[
  a_i/1=\varphi(b_i)/\varphi(s).
\]
There are therefore elements
\[
  t_1,\ldots,t_n\in\mathcal O_G(U)-\mathfrak p
\]
such that
\[
  t_i\bigl(a_i\varphi(s)-\varphi(b_i)\bigr)=0.
\]
Set
\[
  t=t_1\cdots t_n\varphi(s)\in\mathcal O_G(U)-\mathfrak p.
\]
The equalities \(a_i/1=\varphi(b_i)/\varphi(s)\) then already hold in
\(\mathcal O_G(U)_t\). Since
\[
  t\in\mathcal O_G(U)=k[a_1,\ldots,a_n],
\]
there is an element \(b\in\mathcal O_H(V)\) such that, for some
\(r\in\mathbb N\),
\[
  t/1=\varphi(b)/\varphi(s)^r.
\]
Consequently, \(\varphi\) induces a surjection
\[
  \mathcal O_H(V)_{sb}\longrightarrow\mathcal O_G(U)_t,
\]
and hence \(u\) is a local immersion at \(\xi\).

The open subset \(W\) of \(G\) consisting of the points at which \(u\)
is a local immersion is therefore nonempty. Since \(G\) is a Jacobson
scheme, \(W\) contains a closed point \(y\). To prove that \(W=G\), it
is enough to show that every closed point of \(G\) belongs to \(W\). But
every closed point \(x\) is the image of \(y\) under the translation
\(r_x\circ r_y^{-1}\), and hence belongs to \(W\). Thus \(W=G\), which
proves that \(u\) is a local immersion.

Since \(G\) is irreducible, it follows that \(u\) is an immersion.
Indeed, for every \(x\in G\), let \(U_x\) and \(V_x\) be open subsets of
\(G\) and \(H\), respectively, such that \(x\in U_x\) and \(u\)
induces a closed immersion of \(U_x\) into \(V_x\). Since \(U_x\) is
dense in \(G\), the subset \(u(U_x)\) is dense in \(u(G)\cap V_x\).
But \(u(U_x)\) is closed in \(V_x\), so
\footnote{Canon source note (SGA-SRC-000858): the re-edition prints
\(u(U_x)=V_x\). The preceding density statement and the next inverse-image
identity force the intersection shown here.}
\[
  u(U_x)=u(G)\cap V_x.
\]
Moreover, since \(u\) is injective, \(U_x=u^{-1}(V_x)\). Thus \(u\)
induces a closed immersion of \(G\) into the open subscheme of \(H\)
covered by the \(V_x\). Hence \(u:G\to H\) is an immersion. We have
already seen that \(u(G)\) is closed in \(H\), so \(u\) is a closed
immersion.

% END INCLUDED SGA3_English_current_buildable_source_20260823/source/inputs/VI/tex_reference_v2/components/03_expose_VIA_section25_lemma251_prop252_en.tex

% BEGIN INCLUDED SGA3_English_current_buildable_source_20260823/source/inputs/VI/tex_reference_v2/components/04_expose_VIA_lemma253_prop254_loop2_r1_en.tex
% SGA 3, Expose VI_A, Loop 1 English reconstruction.
% Scope: Lemma 2.5.3 through the end of Proposition 2.5.4;
% paragraph 2.6 is excluded.
% Authority: Polo--Gille Exp6A-13oct24.pdf, component-local pages 13--15
% (combined-reader pages 316--318; printed pages 305--307).
% OCR and the Reinhold English Markdown are comparison witnesses only.
% The three diagrams are exact 180-dpi crops of the controlling PDF:
% - local p. 13: x=625,y=1495,w=230,h=165 on a 1488x2105 render;
% - local p. 14: x=590,y=550,w=320,h=160 on a 1488x2105 render;
% - local p. 15: x=620,y=580,w=300,h=135 on a 1488x2105 render.

\setcounter{footnote}{29}
\medskip
\phantomsection\label{lem:sga3-VIA-2.5.3}%
\noindent\textbf{Lemma 2.5.3.}%
\footnote{\label{note:sga3-VIA-30}Editor's note: no finiteness
hypotheses are made in this statement; this will be useful later
(cf.\ \SGARefLink{cor:sga3-VIA-6.2}{6.2} and Expos\'e~VI\textsubscript{B}, \S~12).}
\textit{Let \(A\) be an Artinian local ring, let \(k\) be its residue
field, let \(G\) be a flat \(A\)-group, and let \(X\) be an \(A\)-scheme
equipped with a left action
\[
  \mu:G\times_AX\longrightarrow X
\]
of \(G\) and with a section \(s_0:\operatorname{Spec}A\to X\).
(This is the case, for example, if \(G'\) is a second \(A\)-group and
one is given a morphism of \(A\)-groups \(G\to G'\).)}

\textit{Let \(\varphi\) be the morphism
\[
  \mu\circ(\operatorname{id}_G\times s_0):
  G=G\times_AA\longrightarrow X.
\]
If \(\varphi\) is flat at a point \(g\) of \(G\), then \(\varphi\) is
flat.}

\noindent\textit{Proof.}
Since \(G\) is flat over \(A\), the fiberwise criterion for flatness
(EGA~IV\textsubscript{3}, 11.3.10.2) shows that it is enough to prove
that \(\varphi\otimes_Ak\) is flat. We may therefore suppose that
\(A=k\). In that case, giving \(s_0\) is equivalent to giving a
\(k\)-point \(x_0\in X(k)\), and \(\varphi\) is the morphism
\(h\mapsto hx_0\).

Let \(h\in G\); we show that \(\varphi\) is flat at \(h\). Let \(K\) be
an extension of \(k\) containing copies of both \(\kappa(g)\) and
\(\kappa(h)\). We have a cartesian square
\begin{center}
  \phantomsection\label{diag:sga3-VIA-2.5.3-basechange}%
  \begin{tikzcd}[row sep=large,column sep=large]
    G_K \arrow[r,"\varphi_K"] \arrow[d] &
    X_K \arrow[d] \\
    G \arrow[r,"\varphi"'] &
    X
  \end{tikzcd}
\end{center}
in which the two vertical arrows are faithfully flat. Thus, by
\SGARefLink{lem:sga3-V-7.4:item:i}{Expos\'e~V, 7.4(i)}, \(\varphi_K\) is flat at every point
\(g'\in G_K\) lying over \(g\); and, in order to show that \(\varphi\)
is flat at \(h\), it is enough to show that \(\varphi_K\) is flat at
one point \(h'\) lying over \(h\). We are therefore reduced to the case
in which \(g\) and \(h\) are rational.

Set \(u=hg^{-1}\), and let \(\ell_u\) (respectively, \(\mu_u\)) be left
translation on \(G\) (respectively, on \(X\)) defined by \(u\). Since
\(\varphi\circ\ell_u=\mu_u\circ\varphi\), we obtain a commutative
square
\begin{center}
  \phantomsection\label{diag:sga3-VIA-2.5.3-translations}%
  \begin{tikzcd}[row sep=large,column sep=huge]
    \mathcal O_{G,h} \arrow[r,"\sim"] &
    \mathcal O_{G,g} \\
    \mathcal O_{X,hx_0} \arrow[u] \arrow[r,"\sim"'] &
    \mathcal O_{X,gx_0} \arrow[u]
  \end{tikzcd}
\end{center}
whose horizontal arrows are isomorphisms. Since the morphism
\[
  \mathcal O_{X,gx_0}\longrightarrow\mathcal O_{G,g}
\]
is flat, the morphism
\[
  \mathcal O_{X,hx_0}\longrightarrow\mathcal O_{G,h}
\]
is flat as well.

\medskip
\phantomsection\label{prop:sga3-VIA-2.5.4}%
\noindent\textbf{Proposition 2.5.4.}
\textit{Let \(A\) be an Artinian local ring, let \(k\) be its residue
field, let \(G\) be an \(A\)-group locally of finite type, and let
\(X\ne\varnothing\) be an \(A\)-scheme locally of finite type equipped
with a left action of \(G\). Suppose that the morphism
\[
  \varphi:G\times_AX\longrightarrow X\times_AX,
  \qquad (g,x)\longmapsto(gx,x),
\]
defined on underlying sets by the displayed rule, is surjective. Then:}
\begin{enumerate}
\item[(i)] \textit{The connected components of \(X\) are of finite
type, irreducible, and all of the same dimension.}
\item[(ii)]\phantomsection\label{concl:sga3-VIA-2.5.4-ii} \textit{More precisely, let \(\overline{k}\) be an
algebraic closure of \(k\), and let \(x\) be a closed point of
\(X\otimes_A\overline{k}\). Its stabilizer \(F\) is a closed group
subscheme of \(G\otimes_A\overline{k}\), and the dimension of the
irreducible components of \(X\) is
\[
  \dim G-\dim F.
\]}
\end{enumerate}

In view of \hyperref[lem:sga3-VIA-2.5.1]{Lemma~2.5.1}, we may suppose
that \(A=k\). Suppose first that \(k\) is algebraically closed. Then
\(G_{\mathrm{red}}\) is a \(k\)-group locally of finite type; hence,
replacing \(G\) by \(G_{\mathrm{red}}\) and \(X\) by
\(X_{\mathrm{red}}\), we may suppose that \(G\) and \(X\) are reduced.

Since \(G\times_kX\) is locally of finite type over \(k\), the morphism
\(\varphi\) is locally of finite type (cf.\ EGA~I, 6.6.6(v)), and is
therefore locally of finite presentation because \(X\times_kX\) is
locally Noetherian. Let \(x\) be a rational point of \(X\). The
morphism
\[
  \varphi_x:G\longrightarrow X
\]
obtained from \(\varphi\) by base change is surjective and locally of
finite presentation. If \(\eta\) is a maximal point of \(X\), then
\(\mathcal O_{X,\eta}\) is a field, since \(X\) is reduced; consequently,
\(\varphi_x\) is flat at every point of \(G\) lying over \(\eta\).
By \hyperref[lem:sga3-VIA-2.5.3]{Lemma~2.5.3}, \(\varphi_x\) is
therefore flat. Thus \(\varphi_x:G\to X\) is faithfully flat and
locally of finite presentation, and hence open (cf.\
EGA~IV\textsubscript{2}, 2.4.6).

Since \(G^0\) is open in \(G\), irreducible, and quasi-compact
(by \hyperref[prop:sga3-VIA-2.4]{2.4}), every orbit
\[
  G^0x=\varphi_x(G^0),
\]
as \(x\) ranges over the rational points of \(X\), is an open,
irreducible, and quasi-compact subset of \(X\), and is therefore of
finite type over \(k\), since \(X\) is locally of finite type over
\(k\).

Every nonempty open subset of \(X\) contains a rational point, so these
open subsets cover \(X\). Furthermore, any two of them are either
disjoint or equal. Indeed, if
\[
  \varphi_x(G^0)\cap\varphi_y(G^0)
\]
is nonempty, it contains a rational point \(z\). There are then
rational points \(g,h\in G^0\) such that
\[
  gx=z=hy,
\]
whence \(x=g^{-1}z\) and \(y=h^{-1}z\), and consequently
\[
  \varphi_x(G^0)=\varphi_z(G^0)=\varphi_y(G^0).
\]
It follows that the orbits \(\varphi_x(G^0)\) are also closed, and are
therefore at once the connected components and the irreducible
components of \(X\).

Finally, let \(x,y\) be two rational points of \(X\). Since
\(\varphi_x\) is surjective, there is a rational point \(g\in G\) such
that \(y=gx\). Since \(G^0\) is a normal subgroup of \(G\), the orbit
\(G^0y\) is the image of \(G^0x\)\begingroup
\renewcommand{\thefootnote}{\fnsymbol{footnote}}%
\footnote[1]{Translator's note: the French source prints \(G^0\) here;
as the object being translated is an orbit in \(X\), the intended
source expression is \(G^0x\).}%
\endgroup
under translation \(\ell_g\) on
\(X\); hence \(G^0y\) and \(G^0x\) have the same dimension.

Moreover, by \SGARefLink{sga3:i:scholium:2-3-3-1}{Expos\'e~I, 2.3.3.1}, the stabilizer of \(x\) is represented
by the closed subscheme \(F\) of \(G\) defined by the following
cartesian square:
\begin{center}
  \phantomsection\label{diag:sga3-VIA-2.5.4-stabilizer}%
  \begin{tikzcd}[row sep=large,column sep=large]
    F \arrow[r] \arrow[d] &
    G \arrow[d,"\varphi_x"] \\
    \operatorname{Spec}k \arrow[r] &
    X
  \end{tikzcd}
\end{center}
Thus \(F\) is a \(k\)-group locally of finite type, while
\(F\cap G^0\) is a \(k\)-group of finite type containing \(F^0\). By
\hyperref[cor:sga3-VIA-2.4.1]{2.4.1}, the schemes \(F\) and
\(F\cap G^0\) are equidimensional, of the same dimension as \(F^0\).
Let
\[
  C=\varphi_x(G^0)
\]
be the irreducible component of \(X\) containing \(x\). Arguing as in
the proof of \hyperref[prop:sga3-VIA-2.5.2]{part~(b) of 2.5.2}, we
obtain
\[
  \dim C=\dim G^0-\dim F^0=\dim G-\dim F.
\]

In the general case, that is, when \(k\) is an arbitrary field, let
\(\overline{k}\) be an algebraic closure of \(k\). Let \(C\) be a
connected component of \(X\), and let \(C'\) be a connected component
of \(C\otimes_k\overline{k}\). Then \(C'\) is a connected component of
\[
  X'=X\otimes_k\overline{k}.
\]
The morphism \(\pi:X'\to X\) is open (cf.\
EGA~IV\textsubscript{2}, 2.4.10), and, since it is integral, it is also
closed; therefore \(\pi(C')=C\). Since \(C'\) is irreducible and
quasi-compact, \(C\) is irreducible and quasi-compact, and hence is of
finite type over \(k\), because \(X\) is locally of finite type over
\(k\).

Finally, dimension is invariant under extension of the base field
(cf.\ EGA~IV\textsubscript{2}, 4.1.4), so
\(\dim C=\dim C'\). Since all the irreducible components of \(X'\) have
the same dimension, the same is true of those of \(X\).

% END INCLUDED SGA3_English_current_buildable_source_20260823/source/inputs/VI/tex_reference_v2/components/04_expose_VIA_lemma253_prop254_loop2_r1_en.tex

% BEGIN INCLUDED SGA3_English_current_buildable_source_20260823/source/inputs/VI/tex_reference_v2/components/05_expose_VIA_section26_through_prop264_en.tex
% SGA 3, Expose VI_A, Loop 1 English reconstruction.
% Scope: paragraph 2.6 through the end of Proposition 2.6.4;
% Theorem 2.6.5 is excluded.
% Authority: Polo--Gille Exp6A-13oct24.pdf, component-local pages 15--18
% (combined-reader pages 318--321; printed pages 307--310).
% OCR and the Reinhold English Markdown are comparison witnesses only.

\medskip
\phantomsection\label{sec:sga3-VIA-2.6}%
\noindent\SGAInlineRefTarget{sga3:vi:paragraph:2-6}\textbf{2.6. Complements.---}
This paragraph has been added; it is taken from [Per75], II,
\S\S~1--2, with complements due to O.~Gabber.%
\footnote{\label{note:sga3-VIA-31}Editor's note: these results were
communicated to us by O.~Gabber, in particular \SGARefLink{prop:sga3-VIA-2.6.6}{2.6.6}, which plays an
important role in \hyperref[sec:sga3-VIB-5]{Section~5} of Expos\'e~VI\textsubscript{B}.}
It shows that the preceding results are valid for every group scheme
\(G\) over a field \(k\). (This will be used in \hyperref[sec:sga3-VIB-5]{Sections~5}, 6, and 7
of Expos\'e~VI\textsubscript{B}.)

Fix a field \(k\). We begin with the following lemma
(loc.\ cit., II, 2.1.1), which does not appear explicitly in
EGA~IV\textsubscript{2}, \S~4.4 (although it can probably be read
between the lines at the beginning of loc.\ cit., \S~4.4.1).

\medskip
\phantomsection\label{lem:sga3-VIA-2.6.0}%
\noindent\textbf{Lemma 2.6.0.}
\textit{Let \(X\) be an irreducible \(k\)-scheme, let \(K\) be an
extension of \(k\), and let \(X'\) be an irreducible component of
\(X_K\). The projection \(X'\to X\) is surjective.}

\noindent\textit{Proof.}
Let \(B\) be a transcendence basis of \(K\) over \(k\), and set
\(L=k(B)\subset K\). By EGA~IV\textsubscript{2}, 4.3.2 and 4.4.1,
\(X_L\) is irreducible and \(X'\) dominates \(X_L\). Hence
\(X'\to X_L\) is integral and dominant, and therefore surjective.
Since \(X_L\to X\) is surjective (loc.\ cit., 4.4.1), so is
\(X'\to X\).

For the rest of \hyperref[sga3:vi:paragraph:2-6]{2.6}, fix a \(k\)-group scheme \(G\) and an action
\(\mu:G\times X\to X\) of \(G\) on a \(k\)-scheme \(X\), satisfying
the following condition:
\[
  \tag{\(\star\)}\label{sga3:vi:lemma:2-6-0:equation:star}
  \Phi:G\times X\longrightarrow X\times X,\qquad
  (g,x)\longmapsto(gx,x)
  \quad\text{is surjective.}
\]
(This holds in particular when \(G\) acts on itself by left
translations.) For short, we shall say: ``Let \(X\) be a
\(G\)-scheme satisfying \((\star)\).'' Finally, write \(k'\) for the
perfect closure of \(k\).

\medskip
\phantomsection\label{prop:sga3-VIA-2.6.1}%
\noindent\textbf{Proposition 2.6.1.}
\textit{Let \(X\) be a \(G\)-scheme satisfying \((\star)\).}
\begin{enumerate}
\item[\SGAInlineRefTarget{prop:sga3-VIA-2.6.1:item:i}(i)] \textit{\(X\) is geometrically pointwise irreducible over
\(k\): for every extension \(K\) of \(k\), each point \(x\in X_K\)
belongs to a unique irreducible component of \(X_K\).}
\item[\SGAInlineRefTarget{prop:sga3-VIA-2.6.1:item:ii}(ii)] \textit{Every local ring of \((X_{k'})_{\mathrm{red}}\)
is normal.}
\item[\SGAInlineRefTarget{prop:sga3-VIA-2.6.1:item:iii}(iii)] \textit{Let \(\eta\) be a maximal point of
\((X_{k'})_{\mathrm{red}}\), put
\(C=\overline{\{\eta\}}\), and let \(L\) be the algebraic closure of
\(k\) in \(\kappa(\eta)\). Then \(C\) is an \(L\)-scheme and is
geometrically irreducible over \(L\); that is,
\(C\otimes_LK\) is irreducible for every extension \(K\) of \(L\).}
\item[(iv)] \textit{In particular, if \(x\) is a rational point of
\(X\), then the irreducible component of \(X\) containing \(x\) is
geometrically irreducible over \(k\).}
\end{enumerate}

\noindent\textit{Proof.}
\SGARefLink{prop:sga3-VIA-2.6.1:item:i}{(i)} Since hypothesis \((\star)\) is preserved by every base change
\(k\to K\), it is enough to show that every \(x\in X\) belongs to a
unique irreducible component of \(X\). Since
\(\operatorname{Spec}(k')\to\operatorname{Spec}(k)\) is a universal
homeomorphism, we may further suppose that \(k\) is perfect. We may
then suppose that \(G\) and \(X\) are reduced.

Let \(\eta\) be a maximal point of \(X\), and let \(z\) be any point
of \(Z=\overline{\{\eta\}}\). Since \(X\) is reduced,
\(\mathcal O_{X,\eta}=\kappa(\eta)\); and, since \(k\) is perfect,
\(\kappa(\eta)\otimes_kK\) is normal for every extension \(K\) of
\(k\) (cf.\ EGA~IV\textsubscript{2}, 6.14.2). Thus every point of
\(X_K\) lying above \(\eta\) is normal.

Since \(\Phi\) is surjective, there is a point \(\gamma\) of
\(G\times X\) such that the two projections of \(\Phi(\gamma)\) are
\(z\) and \(\eta\). Set \(K=\kappa(\gamma)\). There are then rational
points \(g\in G_K\) and \(\eta'\in X_K\), with \(\eta'\) lying above
\(\eta\), such that \(z'=g\eta'\) lies above \(z\). By the preceding
paragraph, \(\eta'\) is a normal point of \(X_K\), and therefore so is
\(z'\). Since \(\pi:X_K\to X\) is flat, it follows that
\(z=\pi(z')\) is a normal point of \(X\)
(cf.\ EGA~IV\textsubscript{2}, 2.1.13). This proves (ii) and (i).

The first assertion of \SGARefLink{prop:sga3-VIA-2.6.1:item:iii}{(iii)} now follows from \SGARefLink{prop:sga3-VIA-2.6.1:item:ii}{(ii)}. Since \(L\) is
algebraically closed in \(\kappa(\eta)\), \(C\) is geometrically
irreducible over \(L\), by EGA~IV\textsubscript{2}, 4.5.9.

Finally, if \(X\) has a rational point \(x\), it follows from \SGARefLink{prop:sga3-VIA-2.6.1:item:iii}{(iii)}
that \(L=k\), and hence the irreducible component \(C\) containing
\(x\) is geometrically irreducible over \(k\).\begingroup
\renewcommand{\thefootnote}{\fnsymbol{footnote}}%
\footnote[1]{Translator's note: the French source prints \(X\) in
this conclusion and then switches from the point \(x\) to undefined
symbols \(e\) and \(e_K\). The statement and argument require the
component \(C\) and the points \(x\) and \(x_K\), respectively.}%
\endgroup
This can also be seen directly as follows (cf.\ [Per75], II, \SGARefLink{par:sga3-VIA-2.1}{2.1}).
Let \(C\) be the irreducible component of \(X\) containing \(x\), and
let \(K\) be an extension of \(k\). The scheme \(X_K\) has a unique
point \(x_K\) above \(x\), and by \SGARefLink{prop:sga3-VIA-2.6.1:item:i}{(i)}, \(x_K\) belongs to a unique
irreducible component \(C'\) of \(X_K\). On the other hand, by
\hyperref[lem:sga3-VIA-2.6.0]{2.6.0}, every irreducible component of
\(C_K\) contains \(x_K\), and therefore equals \(C'\).

\medskip
\phantomsection\label{not:sga3-VIA-C0}%
\noindent\textbf{Notation.---}
Write provisionally \(C^0\) for the reduced closed subscheme of \(G\)
whose underlying space is the unique irreducible component of \(G\)
containing the identity element \(e\).

\medskip
\phantomsection\label{cor:sga3-VIA-2.6.2}%
\noindent\textbf{Corollary 2.6.2.}
\textit{\(C^0\) is geometrically irreducible over \(k\) and is
set-theoretically stable under the group law; that is,
\((C^0_{k'})_{\mathrm{red}}\) is a subgroup of
\((G_{k'})_{\mathrm{red}}\). Consequently, \(C^0\) is
quasi-compact.}

\noindent\textit{Proof.}
By \hyperref[prop:sga3-VIA-2.6.1]{2.6.1}, \(C^0\) is geometrically
irreducible over \(k\), so \(C^0\times C^0\) is irreducible and
\[
  \nu(C^0\times C^0)\subset C^0,
\]
where \(\nu\) denotes the morphism \((g,h)\mapsto gh^{-1}\). Since
\(\operatorname{Spec}k'\to\operatorname{Spec}k\) is a universal
homeomorphism, the same conclusion holds for \(C^0_{k'}\), and then
for
\[
  H=(C^0_{k'})_{\mathrm{red}}.
\]
Since \(H\times_{k'}H\) is reduced, \(\nu\) induces a morphism
\[
  H\times_{k'}H\longrightarrow H;
\]
thus \(H\) is a subgroup of \((G_{k'})_{\mathrm{red}}\).
Consequently, by \hyperref[sga3:VIA:0.5.1]{0.5.1}, \(H\), and
hence \(C^0\), is quasi-compact.

\medskip
\phantomsection\label{rem:sga3-VIA-2.6.3}%
\noindent\textbf{Recall 2.6.3.---}
Let \(Y\) be a scheme. Recall (EGA~0\textsubscript{III}, 9.1.1) that
a subset \(E\) of \(Y\) is called \emph{retrocompact} if the
inclusion \(E\hookrightarrow Y\) is quasi-compact. Moreover, by
EGA~IV\textsubscript{1}, 1.9.5(v) and 1.10.1, if \(U\) is a
retrocompact open subset of \(Y\), its closure \(\overline U\) is the
union of the closures \(\overline{\{y\}}\) of the points \(y\in U\).
Of course, it is enough to let \(y\) run through the maximal points of
\(U\), that is, the maximal points of \(Y\) contained in \(U\).

Consequently, if \(Y\) is quasi-separated and \(\eta\) is a maximal
point of \(Y\), then the intersection of the closed neighborhoods of
\(\eta\) equals \(\overline{\{\eta\}}\). Indeed, if
\(y\in Y-\overline{\{\eta\}}\), then \(y\) belongs to an affine open
\(V\) not containing \(\eta\). Since \(Y\) is quasi-separated, \(V\)
is retrocompact, and therefore \(\overline V\) is the union of the
\(\overline{\{\xi\}}\), for \(\xi\) running through the maximal points
of \(Y\) belonging to \(V\). Hence \(\eta\notin\overline V\), so
\(Y-\overline V\) is a closed neighborhood of \(\eta\) not containing
\(y\).

\medskip
\phantomsection\label{prop:sga3-VIA-2.6.4}%
\noindent\textbf{Proposition 2.6.4.}
\textit{Let \(X\) be a \(G\)-scheme satisfying \((\star)\), and let
\(U\) be an open subset of \(X\).}
\begin{enumerate}
\item[\SGAInlineRefTarget{prop:sga3-VIA-2.6.4:item:i}(i)] \textit{\(C^0U\) is an open subset of \(X\), equal to the
union of the irreducible components of \(X\) whose generic point
belongs to \(U\).}
\item[(i\(^{\prime}\))] \textit{Consequently, if \(X\) is
irreducible, it is quasi-compact.}
\item[\SGAInlineRefTarget{prop:sga3-VIA-2.6.4:item:ii}(ii)] \textit{If, moreover, \(U\) is retrocompact in \(X\), then
\(C^0U=\overline U\), and this subset is both open and closed in
\(X\).}
\end{enumerate}

\noindent\textit{Proof.}
\SGARefLink{prop:sga3-VIA-2.6.4:item:i}{(i)} First, \(C^0U\) is open, since it is the union, for \(g\in C^0\),
of the projections of the open subsets
\[
  g\cdot U_{\kappa(g)}\subset X_{\kappa(g)},
\]
and every projection \(X_{\kappa(g)}\to X\) is open.

To prove the second assertion of \SGARefLink{prop:sga3-VIA-2.6.4:item:i}{(i)}, we may replace \(k\) by \(k'\)
(since \(\operatorname{Spec}k'\to\operatorname{Spec}k\) is a
universal homeomorphism), and therefore suppose that \(k\) is perfect.
We may then suppose that \(G\) and \(X\) are reduced, and hence
geometrically reduced.

Let \(\eta\) be a maximal point of \(X\) contained in \(U\), and put
\(Z=\overline{\{\eta\}}\). Consider
\[
  \mu':C^0\times Z\longrightarrow X.
\]
Since \(C^0\) is geometrically irreducible, \(C^0\times Z\) is
irreducible; write \(\gamma\) for its generic point. The morphism
\(\mu'\) sends the point \(\varepsilon(\eta)\), where \(\varepsilon\)
denotes the identity section of \(C^0\), to \(\eta\). Hence
\(\gamma\) is sent to a generization of \(\eta\), and therefore to
\(\eta\). Thus \(\mu'\) sends the underlying space of \(C^0\times Z\)
into \(Z\), and, since \(C^0\times Z\) is reduced, \(\mu'\) factors
through \(Z\).

Now let \(z\in Z\), and put \(K=\kappa(z)\). The morphism
\[
  \mu_z:G_K\longrightarrow X_K,\qquad h\longmapsto h\cdot z,
\]
is surjective. Let \(\alpha\) be a maximal point of \(X_K\). The local
ring \(\mathcal O_{X_K,\alpha}\) is a field, since \(X_K\) is reduced,
so \(\mu_z\) is flat at every point of \(G_K\) above \(\alpha\).
It follows from \hyperref[lem:sga3-VIA-2.5.3]{Lemma~2.5.3} that
\(\mu_z\) is flat.

On the other hand, \(\mu_z\) sends the generic point \(\omega\) of
\(C^0_K\) to a point \(t\in Z_K\). Let \(\beta\) be a maximal point
of \(Z_K\) such that \(t\in\overline{\{\beta\}}\). Since \(\mu_z\) is
flat, there is a generization \(\xi\) of \(\omega\) such that
\(\mu_z(\xi)=\beta\). But \(\omega\) is a maximal point of \(G_K\), so
necessarily \(\xi=\omega\). Thus \(\mu_z(\omega)=\beta\), and
\(\beta\) lies above \(\eta\), because \(Z_K\to Z\) is flat.

Put \(L=\kappa(\omega)\), and let \(\omega_L\) and \(z_L\) be the
\(L\)-points obtained from \(\omega\) and \(z\). Then
\[
  \omega_L\cdot z_L=\beta'
\]
is a point of \(Z_L\) lying above \(\beta\), and hence
\[
  z_L=\omega_L^{-1}\cdot\beta'\in C^0_L\cdot U_L.
\]
Therefore \(z\in C^0\cdot U\), which proves \SGARefLink{prop:sga3-VIA-2.6.4:item:i}{(i)}.

Since \(C^0\) is quasi-compact by
\hyperref[cor:sga3-VIA-2.6.2]{2.6.2}, assertion
(i\(^{\prime}\)) follows: if \(X\) is irreducible and \(U\) is a
nonempty affine open subset, then \(X=C^0U\), so \(X\) is the image of
the morphism \(C^0\times U\to X\), and is therefore quasi-compact.
Finally, if \(U\) is retrocompact in \(X\), then by
\hyperref[rem:sga3-VIA-2.6.3]{2.6.3}, \(\overline U\) is the union of
the \(\overline{\{\eta\}}\), where \(\eta\) runs through the maximal
points of \(X\) contained in \(U\). Hence \(\overline U=C^0U\).
This proves \SGARefLink{prop:sga3-VIA-2.6.4:item:ii}{(ii)}.

% END INCLUDED SGA3_English_current_buildable_source_20260823/source/inputs/VI/tex_reference_v2/components/05_expose_VIA_section26_through_prop264_en.tex

% BEGIN INCLUDED SGA3_English_current_buildable_source_20260823/source/inputs/VI/tex_reference_v2/components/06_expose_VIA_theorem265_through_remark267_loop2_r1_en.tex
% SGA 3, Expose VI_A, Loop 1 English reconstruction.
% Scope: Theorem 2.6.5 through Remark 2.6.7;
% Section 3 is excluded.
% Authority: Polo--Gille Exp6A-13oct24.pdf, component-local pages 18--20
% (combined-reader pages 321--323; printed pages 310--312).
% OCR and the Reinhold English Markdown are comparison witnesses only.
% The diagram is an exact 180-dpi crop of the controlling PDF:
% local p. 20, x=650,y=460,w=260,h=170 on a 1488x2105 render.

One then obtains the following result ([Per75], II, Th.~\SGARefLink{prop:sga3-VIA-2.4}{2.4}; see also
[Per76], Prop.~4.1.1).

\medskip
\phantomsection\label{thm:sga3-VIA-2.6.5}%
\noindent\textbf{Theorem 2.6.5.}
\textit{Let \(k\) be a field, and let \(G\) be a \(k\)-group scheme.}
\begin{enumerate}
\item[\SGAInlineRefTarget{thm:sga3-VIA-2.6.5:item:i}(i)] \textit{There is a unique group subscheme \(G^0\) of
\(G\), called the \emph{identity component} of \(G\), such that:}
  \begin{enumerate}
  \item[\SGAInlineRefTarget{thm:sga3-VIA-2.6.5:item:a}(a)] \textit{the underlying space of \(G^0\) is the
  irreducible component of the identity element;}
  \item[\SGAInlineRefTarget{thm:sga3-VIA-2.6.5:item:b}(b)] \textit{\(G^0\to G\) is a flat closed immersion; that is,
  \[
    \mathcal O_{G,g}=\mathcal O_{G^0,g}
    \qquad\text{for every }g\in G^0.
  \]}
  \end{enumerate}
\item[\SGAInlineRefTarget{thm:sga3-VIA-2.6.5:item:ii}(ii)] \textit{Moreover, \(G^0\) is quasi-compact,
geometrically irreducible, and is a characteristic subgroup of \(G\).}
\item[\SGAInlineRefTarget{thm:sga3-VIA-2.6.5:item:iii}(iii)] \textit{If \(G\) is connected, then \(G=G^0\).}
\end{enumerate}

\noindent\textit{Proof.}
\SGARefLink{thm:sga3-VIA-2.6.5:item:i}{(i)} Recall first that \(G\) is separated
(\hyperref[sga3:VIA:0.3]{0.3}), and hence, a fortiori,
quasi-separated. Let \(U\) be an affine open subset of \(G\) containing
the generic point \(\omega\) of \(C^0\). By
\hyperref[rem:sga3-VIA-2.6.3]{2.6.3} and
\hyperref[prop:sga3-VIA-2.6.4]{2.6.4},
\[
  \overline U=C^0U
\]
is both open and closed, and, set-theoretically,
\[
  C^0=\overline{\{\omega\}}
\]
is the intersection of these open-and-closed subsets.

As \(U\) runs through the affine open subsets containing \(\omega\),
the \(G\)-schemes \(\overline U\) form a projective system whose
transition morphisms are affine, since they are closed immersions.
We can therefore form its projective limit \(G^0\)
(cf.\ EGA~IV\textsubscript{3}, 8.2.2). Thus, for every affine open
\(V\) of \(G\), \(G^0\cap V\) is the spectrum of the algebra
\[
  \varinjlim_U\mathcal O_G(V\cap\overline U)
  =
  \mathcal O_G(V)\Big/\sum_U I_{\overline U}(V),
\]
where \(I_{\overline U}(V)\) denotes the kernel of
\[
  \mathcal O_G(V)\longrightarrow
  \mathcal O_G(V\cap\overline U).
\]
It follows that the underlying space of \(G^0\) is \(C^0\), and that
\(G^0\to G\) is a closed immersion.

Moreover, for every \(g\in G^0\), \(\mathcal O_{G^0,g}\) is the
inductive limit, as \(V\) runs through the affine open subsets of
\(G\) containing \(g\), of the \(k\)-algebras
\[
  \mathcal O_{G^0}(V\cap G^0)
  =
  \varinjlim_U\mathcal O_G(V\cap\overline U).
\]
This double inductive limit is identified with
\[
  \varinjlim_U\
  \varinjlim_{\substack{V\\g\in V\subset\overline U}}
  \mathcal O_G(V)
  =
  \mathcal O_{G,g}.
\]
Thus
\[
  \mathcal O_{G^0,g}=\mathcal O_{G,g}
\]
(see also EGA~IV\textsubscript{2}, 5.13.3(ii)), and consequently
\(i:G^0\to G\) is a flat closed immersion. Conversely, this condition
implies that \(i^*(\mathcal O_G)=\mathcal O_{G^0}\), so \(G^0\) is
uniquely determined by conditions \SGARefLink{thm:sga3-VIA-2.6.5:item:a}{(a)} and \SGARefLink{thm:sga3-VIA-2.6.5:item:b}{(b)}. This proves \SGARefLink{thm:sga3-VIA-2.6.5:item:i}{(i)}.

The first two assertions of \SGARefLink{thm:sga3-VIA-2.6.5:item:ii}{(ii)} follow from
\hyperref[cor:sga3-VIA-2.6.2]{2.6.2}. Finally, let \(S\) be a
\(k\)-scheme and let \(\phi\) be an automorphism of the \(S\)-group
\(G_S\). For every \(s\in S\), \(\phi_s\) sends
\(G^0_{\kappa(s)}\) into itself, and therefore
\(\phi(G^0_S)\subset G^0_S\). Moreover, the closed immersion
\[
  i_S:G^0_S\hookrightarrow G_S
\]
obtained from \(i\) by base change is flat, so
\[
  \mathcal O_{G_S,\phi(z)}
  =
  \mathcal O_{G^0_S,\phi(z)}
  \qquad\text{for every }z\in G^0_S.
\]
Hence \(\phi\circ i_S\) factors through \(G^0_S\). This proves that
\(G^0\) is a characteristic subgroup of \(G\), and proves \SGARefLink{thm:sga3-VIA-2.6.5:item:ii}{(ii)}.
Finally, \SGARefLink{thm:sga3-VIA-2.6.5:item:iii}{(iii)} is a special case of part \SGARefLink{thm:sga3-VIA-2.6.5:item:i}{(i)} of the following
proposition.%
\footnote{\label{note:sga3-VIA-32}Editor's note: this proposition, as
well as the preceding results, was communicated to us by O.~Gabber.
It will be used to correct the proof of \hyperref[thm:sga3-VIB-5.3]{Theorem~5.3} in
Expos\'e~VI\textsubscript{B}.}

\medskip
\phantomsection\label{prop:sga3-VIA-2.6.6}%
\noindent\textbf{Proposition 2.6.6.}
\textit{Let \(k\) be a field, and let \(G\) be a \(k\)-group acting
on a \(k\)-scheme \(X\) in such a way that
\[
  G\times X\longrightarrow X\times X,\qquad
  (g,x)\longmapsto(gx,x)
\]
is surjective. Suppose that \(X\) is quasi-separated. Then:}
\begin{enumerate}
\item[\SGAInlineRefTarget{prop:sga3-VIA-2.6.6:item:i}(i)] \textit{Every connected component \(C\) of \(X\) is
irreducible.}
\item[\SGAInlineRefTarget{prop:sga3-VIA-2.6.6:item:ii}(ii)] \textit{Let \(\eta\) be the generic point of \(C\), and
let \(L\) be the algebraic closure of \(k\) in \(\kappa(\eta)\).
Then \(C\) is an \(L\)-scheme, is geometrically irreducible over
\(L\), and the morphism
\[
  G^0_L\times_LC\longrightarrow C\times_LC
\]
is surjective.}
\item[\SGAInlineRefTarget{prop:sga3-VIA-2.6.6:item:iii}(iii)] \textit{In particular, if \(C\) contains a rational
point \(x\), then \(C\) is geometrically irreducible over \(k\), and
\[
  \phi:G^0\longrightarrow C,\qquad g\longmapsto gx,
\]
is surjective.}
\end{enumerate}

\noindent\textit{Proof.}
\SGARefLink{prop:sga3-VIA-2.6.6:item:i}{(i)} Let \(C\) be a connected component of \(X\), let \(Y\) be an
irreducible component of \(X\) contained in \(C\), let \(\eta\) be
the generic point of \(Y\), and let \(U\) be an affine open subset of
\(X\) containing \(\eta\). Since \(X\) is quasi-separated, \(U\) is
retrocompact in \(X\). Hence, by
\hyperref[prop:sga3-VIA-2.6.4]{2.6.4},
\[
  \overline U=G^0U
\]
is an open-and-closed subset of \(X\) which meets \(C\), and therefore
contains \(C\). By \hyperref[rem:sga3-VIA-2.6.3]{2.6.3}, the
intersection of the \(\overline U\), as \(U\) runs through a
fundamental system of affine open neighborhoods of \(\eta\), is
\(\overline{\{\eta\}}\). It follows that
\[
  C=\overline{\{\eta\}},
\]
which proves \SGARefLink{prop:sga3-VIA-2.6.6:item:i}{(i)}.

The first assertion of \SGARefLink{prop:sga3-VIA-2.6.6:item:ii}{(ii)}, and also that of \SGARefLink{prop:sga3-VIA-2.6.6:item:iii}{(iii)}, follows from
\hyperref[prop:sga3-VIA-2.6.1]{2.6.1}. We first prove the second
assertion of \SGARefLink{prop:sga3-VIA-2.6.1:item:iii}{(iii)}. Let \(\omega\), respectively \(\eta\), denote the
generic point of \(G^0\), respectively \(C\). Let \(z\in C\), and set
\(K=\kappa(z)\). Since \(G^0\), respectively \(C\), is geometrically
irreducible over \(k\), the scheme \(G^0_K\), respectively \(C_K\),
is irreducible; denote its generic point by \(\xi\), respectively
\(\beta\). We saw in the proof of
\hyperref[prop:sga3-VIA-2.6.4]{2.6.4} that
\[
  \mu_z:G_K\longrightarrow C_K,\qquad h\longmapsto h\cdot z,
\]
sends \(\xi\) to \(\beta\). Likewise, \(\phi(\omega)=\eta\).

Put \(L=\kappa(\xi)\), and let \(\xi_L\) and \(z_L\) be the \(L\)-points
obtained from \(\xi\) and \(z\). Then
\[
  \xi_L\cdot z_L=\beta'
\]
is a point of \(C_L\) lying above \(\beta\in C_K\), and hence also
above \(\eta\in C\). Consider the cartesian square
\begin{center}
  \phantomsection\label{diag:sga3-VIA-2.6.6-cartesian-square}%
  \begin{tikzcd}[row sep=large,column sep=large]
    G_L^0 \arrow[r,"\varphi_L"] \arrow[d,"\pi_{G^0}"'] &
    C_L \arrow[d,"\pi_C"] \\
    G^0 \arrow[r,"\varphi"'] &
    C
  \end{tikzcd}
\end{center}
Since
\[
  \phi_L\bigl(\pi_{G^0}^{-1}(\omega)\bigr)
  =
  \pi_C^{-1}(\eta)
\]
(cf.\ EGA~I, 3.4.8), there is a \(g\in G^0_L\) such that
\(\phi_L(g)=\beta'\). Therefore
\[
  z_L
  =
  \xi_L^{-1}\cdot\phi_L(g)
  =
  \phi_L(\xi_L^{-1}g),
\]
whence
\[
  \phi\bigl(\pi_{G^0}(\xi_L^{-1}g)\bigr)
  =
  \pi_C(z_L)
  =
  z.
\]
This proves that \(\phi\) is surjective.

We now prove the second assertion of \SGARefLink{prop:sga3-VIA-2.6.6:item:ii}{(ii)}. It is enough to show that,
for every \(z\in C\), the morphism
\[
  \mu_z:
  G^0_L\otimes_L\kappa(z)
  \longrightarrow
  C\otimes_L\kappa(z)
\]
is surjective. But this follows from \SGARefLink{prop:sga3-VIA-2.6.6:item:iii}{(iii)}, since \(z\) is a rational
point of \(C\otimes_L\kappa(z)\).

\medskip
\phantomsection\label{rem:sga3-VIA-2.6.7}%
\noindent\textbf{Remark 2.6.7.---}
\textit{Under the hypotheses of \hyperref[prop:sga3-VIA-2.6.6]{2.6.6}, if \(C(k)=\varnothing\), the
morphism
\[
  G^0\times_kC\longrightarrow C\times_kC
\]
need not be surjective. For example, let \(k=\mathbb R\) and
\(G=\{\pm1\}_{\mathbb R}\). The \(G\)-torsor
\[
  X=\operatorname{Spec}\mathbb R[X]/(X^2+1)
\]
is connected, but
\[
  G^0\times_{\mathbb R}X\longrightarrow X\times_{\mathbb R}X
\]
is not surjective. (Here \(L=\mathbb C\), and
\[
  G^0\times_{\mathbb R}X\longrightarrow X\times_{\mathbb C}X
\]
is an isomorphism.)}

% END INCLUDED SGA3_English_current_buildable_source_20260823/source/inputs/VI/tex_reference_v2/components/06_expose_VIA_theorem265_through_remark267_loop2_r1_en.tex

% BEGIN INCLUDED SGA3_English_current_buildable_source_20260823/source/inputs/VI/tex_reference_v2/components/07_expose_VIA_section31_theorem32_statement_loop2_r1_en.tex
% SGA 3, Expose VI_A, Loop 1 English reconstruction.
% Scope: Section 3 heading and paragraph 3.1 through the full statement
% of Theorem 3.2 and the two following setup paragraphs;
% paragraph 3.2.1 is excluded.
% Authority: Polo--Gille Exp6A-13oct24.pdf, component-local pages 20--21
% (combined-reader pages 323--324; printed pages 312--313).
% OCR and the Reinhold English Markdown are comparison witnesses only.
% The three diagrams are exact 180-dpi crops of the controlling PDF:
% local p. 20, x=530,y=1180,w=440,h=60;
% local p. 20, x=470,y=1320,w=550,h=120; and
% local p. 21, x=520,y=700,w=470,h=80, on 1488x2105 renders.

\section*{3. Construction of quotients \(F\backslash G\)
  (for \(G,F\) of finite type)}
\phantomsection\label{sec:sga3-VIA-3}
\addcontentsline{toc}{section}{3. Construction of quotients
  \texorpdfstring{\(F\backslash G\)}{F backslash G}
  (for G, F of finite type)}

\medskip
\phantomsection\label{par:sga3-VIA-3.1}%
\noindent\textbf{3.1.---}
Let \(A\) be an Artinian local ring, and let
\(u:F\to G\) be a homomorphism of \(A\)-groups. If
\[
  \mu:F\times_A F\longrightarrow F
  \qquad\text{and}\qquad
  \nu:G\times_A G\longrightarrow G
\]
denote the multiplication morphisms, and if \(\lambda\) denotes the
composite morphism
\begin{center}
  \phantomsection\label{diag:sga3-VIA-3.1-lambda-composite}%
  \begin{tikzcd}[column sep=large]
    F\times_A G
      \arrow[r,"u\times G"]
    &
    G\times_A G
      \arrow[r,"\nu"]
    &
    G
  \end{tikzcd},
\end{center}
recall that the \emph{left quotient} \(F\backslash G\) of \(G\) by
\(F\) is the cokernel of the \((\mathrm{Sch}/A)\)-groupoid \(G_*\)
described below:
\begin{center}
  \phantomsection\label{diag:sga3-VIA-3.1-groupoid}%
  \begin{tikzcd}[column sep=huge]
    F\times_A F\times_A G
      \arrow[r,shift left=2.1ex,"F\times\lambda"]
      \arrow[r,"\mu\times G" description]
      \arrow[r,shift right=2.1ex,"\operatorname{pr}_{2,3}"']
    &
    F\times_A G
      \arrow[r,shift left=1.2ex,"\lambda"]
      \arrow[r,shift right=1.2ex,"\operatorname{pr}_2"']
    &
    G
  \end{tikzcd}
\end{center}
Here \(\operatorname{pr}_2\) and \(\operatorname{pr}_{2,3}\) are the
projections of \(F\times_A G\) and
\(F\times_A(F\times_A G)\), respectively, onto the second factors.
We shall say that \(G_*\) is the \emph{groupoid with base \(G\)
defined by \(u\)} (cf.\ Expos\'e~V, \S~2.a; as in Expos\'e~V, we do
not follow in this expos\'e the convention of IV, \SGARefLink{sga3:iv:scholium:4:6:15}{4.6.15}).

Since the unique \(A\)-morphism \(F\to\operatorname{Spec}A\) is
universally open (EGA~IV\textsubscript{2}, 2.4.9),
\(\operatorname{pr}_2\) is an open morphism. The same therefore holds
for \(\lambda\), which is the composite of \(\operatorname{pr}_2\)
with the automorphism \(\sigma\) of \(F\times_A G\) defined by
\[
  \sigma(S)(x,y)=\bigl(x,u(S)(x)\cdot y\bigr),
\]
where \(S\) is a variable \(A\)-scheme and
\(x\in F(S)\), \(y\in G(S)\). One sees in the same way that
\(\operatorname{pr}_2\) and \(\lambda\) are flat when \(F\) is flat
over \(A\).

Let us also note, to conclude these preliminaries, that every
\(A\)-morphism \(s:\operatorname{Spec}A\to G\) defines an automorphism
of the groupoid \(G_*\) which induces on \(G\), \(F\times_A G\), and
\(F\times_A F\times_A G\), respectively, the automorphisms
\[
  r_s,\qquad
  \mathrm{id}_F\times_A r_s,\qquad
  \mathrm{id}_F\times_A\mathrm{id}_F\times_A r_s.
\]
We shall again write \(r_s\) for this automorphism of \(G_*\), and
shall say that \(r_s\) is the \emph{right translation defined by
\(s\)} (cf.\ \SGARefLink{sga3:VIA:0.4}{0.4}).

\medskip
\phantomsection\label{thm:sga3-VIA-3.2}%
\noindent\textbf{Theorem 3.2.}
\textit{Let \(F\) and \(G\) be flat groups locally of finite type over
an Artinian local ring \(A\). Let \(u:F\to G\) be a quasi-compact
homomorphism of \(A\)-groups whose kernel is finite over \(A\). Then:}%
\footnote{\label{note:sga3-VIA-33}Editor's note: we have added
(ii\(^{\prime}\)) and expanded point~\hyperref[concl:sga3-VIA-3.2-iv]{(iv)}, taking account of the
additions made in \SGARefLink{prop:sga3-VIA-2.5.2}{2.5.2}, \SGARefLink{prop:sga3-VIA-2.5.4}{2.5.4}, and \SGARefLink{lem:sga3-V-6.1}{Expos\'e~V, 6.1}.}
\begin{enumerate}
\item[\SGAInlineRefTarget{thm:sga3-VIA-3.2:item:i}(i)] \textit{The left quotient \(F\backslash G\) of \(G\) by
\(F\) exists in \((\mathrm{Sch}/A)\), and the sequence}
\begin{center}
  \phantomsection\label{diag:sga3-VIA-3.2-exact-sequence}%
  \begin{tikzcd}[column sep=large]
    F\times_A G
      \arrow[r,shift left=1.2ex,"\lambda"]
      \arrow[r,shift right=1.2ex,"\operatorname{pr}_2"']
    &
    G
      \arrow[r,"p"]
    &
    F\backslash G
  \end{tikzcd}
\end{center}
\textit{is exact in the category of all ringed spaces.}

\item[\SGAInlineRefTarget{thm:sga3-VIA-3.2:item:ii}(ii)] \textit{The canonical morphism
\(p:G\to F\backslash G\) is surjective and open, and
\(F\backslash G\) is a direct sum of schemes of finite type over
\(A\).}

\item[(ii\(^{\prime}\))]\phantomsection\label{concl:sga3-VIA-3.2-ii-prime} \textit{More precisely,
\(X=F\backslash G\) is equipped with a right action of \(G\) such
that \(p(e)\cdot g=p(g)\) for every \(g\in G\). Consequently, the
connected components of \(X\) are of finite type over \(A\),
irreducible, and all of dimension \(\dim G-\dim F\).}

\item[\SGAInlineRefTarget{thm:sga3-VIA-3.2:item:iii}(iii)] \textit{The canonical morphism
\[
  F\times_A G\longrightarrow G\times_{(F\backslash G)}G
\]
is surjective.}

\item[(iv)]\phantomsection\label{concl:sga3-VIA-3.2-iv} \textit{If \(u\) is a monomorphism,}%
\footnote{\label{note:sga3-VIA-34}Editor's note: in this case, the
hypothesis that \(G\) be flat can be omitted; cf.\ \SGARefLink{par:sga3-VIA-3.3}{subsection~3.3}.}
\textit{ then:}
\begin{enumerate}
\item[\SGAInlineRefTarget{thm:sga3-VIA-3.2:item:a}(a)] \textit{The morphism
\[
  F\times_A G
    \xrightarrow{(\lambda,\operatorname{pr}_2)}
  G\times_{(F\backslash G)}G
\]
is an isomorphism, and \(G\to F\backslash G\) is faithfully flat and
locally of finite presentation.}

\item[(a\(^{\prime}\))] \textit{\(F\backslash G\) represents the
\((\mathrm{fppf})\) quotient sheaf
\(\widetilde{F\backslash G}\), and \(G\to F\backslash G\) is an
\(F\)-torsor locally trivial for the \((\mathrm{fppf})\) topology.}

\item[(b)] \textit{\(F\backslash G\) is flat over \(A\), and is
smooth over \(A\) if \(G\) is so.}

\item[(c)] \textit{\(u:F\to G\) is a closed immersion, and
\(F\backslash G\) is separated.}

\item[(d)] \textit{If, moreover, \(F\) is a normal subgroup of
\(G\), there is a unique \(A\)-group structure on \(F\backslash G\)
for which \(p:G\to F\backslash G\) is a morphism of \(A\)-groups.}
\end{enumerate}
\end{enumerate}

In the proof of this theorem, \(A'\) will denote an \(A\)-algebra
which is local, finite, and free over \(A\). If \(R\) is a relation
involving \(A'\), we shall say that ``\(R(A')\) is true when \(A'\)
is sufficiently large'' if there is an \(A\)-algebra \(A_1\), local,
finite, and free over \(A\), such that the relation \(R(A')\) holds
for every \(A\)-algebra \(A'\) which is local, finite, and free over
\(A_1\).

We shall first prove the theorem when \(F\) and \(G\) are of finite
type over \(A\).

% END INCLUDED SGA3_English_current_buildable_source_20260823/source/inputs/VI/tex_reference_v2/components/07_expose_VIA_section31_theorem32_statement_loop2_r1_en.tex

% BEGIN INCLUDED SGA3_English_current_buildable_source_20260823/source/inputs/VI/tex_reference_v2/components/08_expose_VIA_section321_quasisection_descent_loop2_r1_en.tex
% SGA 3, Expose VI_A, Loop 1 English reconstruction.
% Scope: paragraph 3.2.1 in full; paragraph 3.2.2 is excluded.
% Authority: Polo--Gille Exp6A-13oct24.pdf, component-local pages 22--23
% (combined-reader pages 325--326; printed pages 314--315).
% OCR and the Reinhold English Markdown are comparison witnesses only.
% The diagram is an exact 180-dpi crop of the controlling PDF:
% local p. 23, x=490,y=440,w=580,h=170 on a 1488x2105 render.

\medskip
\phantomsection\label{par:sga3-VIA-3.2.1}%
\noindent\textbf{3.2.1.---}
Suppose for a moment that every point of \(G\) has an open saturated
neighborhood \(W\) such that the groupoid induced by \(G_*\) on \(W\)
has a quasi-section (cf.\ Expos\'e~V, \S~6). Then, by \SGARefLink{lem:sga3-V-6.1}{Expos\'e~V,
6.1}, assertions \SGARefLink{lem:sga3-V-6.1:item:i}{(i)}, (ii), (iii), and (iv)(a) hold, and
\(F\backslash G\) is of finite type over \(A\).\begingroup
\renewcommand{\thefootnote}{\fnsymbol{footnote}}%
\footnote[1]{Translator's note: the French source prints \(k\) here
and again in editor's \hyperref[note:sga3-VIA-39]{note 39}. In this part of the proof the base is
the Artinian local ring \(A\), and ``finite type over \(A\)'' is the
required statement.}%
\endgroup
Moreover, under the hypothesis of (iv), since
\[
  G\longrightarrow F\backslash G
\]
is faithfully flat and locally of finite presentation, assertion (b)
follows from EGA~IV, 2.2.14 and 17.7.7. Assertion (iv)(a\(^{\prime}\))
follows from (iv)(a), by Expos\'e~IV, 3.4.3.1, 5.2.2, and 5.1.6.
Finally, assertion (iv)(c) will be proved in \hyperref[par:sga3-VIA-3.2.5]{3.2.5}, while
(ii\(^{\prime}\)) and
(iv)(d) will be proved in \hyperref[sec:sga3-VIA-5]{Section~5}.

We now prove the following assertion:
\[
  \tag{\(\dagger\)}\label{eq:sga3-VIA-3.2.1-dagger}
  \text{every finite subset of \(F\backslash G\) is contained in an
  affine open subset.}
\]%
\footnote{\label{note:sga3-VIA-35}Editor's note: let us point out that
if \(A=k\) is a field, every quasi-compact open subset of
\(F\backslash G\) is quasi-projective (a result due to Chow for smooth
algebraic groups); cf.\ [Ray70], VI, \SGARefLink{cor:sga3-VIB-2.6}{2.6}. By contrast, over the
Artinian local ring \(A=\mathbb C[\varepsilon]/(\varepsilon^2)\),
there exist abelian \(A\)-schemes \(G\) which are not projective
(loc.\ cit., XII, 4.2).}

If \(U\) is a quasi-section of the groupoid induced by \(G_*\) on a
saturated open subset \(W\) of \(G\), then \(U\otimes_AA'\) is a
quasi-section of the \((\mathrm{Sch}/A')\)-groupoid induced by
\(G_*\otimes_AA'\) on \(W\otimes_AA'\). Moreover, if \(U_*\) is the
\((\mathrm{Sch}/A)\)-groupoid induced by \(G_*\) on \(U\), then
\(U_*\otimes_AA'\) is identified with the
\((\mathrm{Sch}/A')\)-groupoid induced by \(G_*\otimes_AA'\) on
\(U\otimes_AA'\).%
\footnote{\label{note:sga3-VIA-36}Editor's note: the preceding
assertion is valid for every base change \(A\to A'\).}

It follows from the proofs in Expos\'e~V that the construction of the
quotient
\[
  X=F\backslash G
\]
commutes with a base extension \(A\to A'\) of the kind considered
here.%
\footnote{\label{note:sga3-VIA-37}Editor's note: this is explained in
detail in \SGARefLink{par:sga3-VIA-4.6}{4.6} below. One must show that formation of direct images by
the morphisms \(p\), \(\lambda\), and \(\operatorname{pr}_2\)
commutes with flat base changes \(A\to A'\). Since \(F\) and \(G\)
are of finite type over the Artinian ring \(A\), the morphisms \(f\)
in question are all quasi-compact and quasi-separated, and
\[
  f_*(\mathcal O_X)\otimes_AA'
  =
  f'_*(\mathcal O_{X'})
\]
(with the evident notation) follows from
EGA~IV\textsubscript{1}, 1.7.21.}

Let \(x_1,\ldots,x_n\) be points of \(X=F\backslash G\), which we may
suppose to be closed,%
\footnote{\label{note:sga3-VIA-38}Editor's note: indeed, let
\(y_1,\ldots,y_n\) be arbitrary points of \(X\). Since \(X\) is of
finite type over \(A\), the closure of each \(y_i\) contains a closed
point \(x_i\), and every open subset containing \(x_i\) contains
\(y_i\).}
and let \(g_1,\ldots,g_n\) be closed points of \(G\) mapping to
\(x_1,\ldots,x_n\). Let \(V\) be an everywhere dense affine open
subset of \(X\),%
\footnote{\label{note:sga3-VIA-39}Editor's note: such an open subset
exists because \(X\) is of finite type over \(A\): \(X\) has only
finitely many irreducible components \(C_1,\ldots,C_p\), and it is
enough, for each \(i\), to choose a nonempty affine open subset
contained in
\[
  C_i-\bigcup_{j\ne i}C_j.
\]
(Here, moreover, one knows from (ii\(^{\prime}\)) that the \(C_i\)
are disjoint.)}
and let \(U\) be its inverse image in \(G\). By \SGARefLink{prop:sga3-VIA-1.2}{1.2}, there is an
\(A\)-algebra \(A'\), local, finite, and free over \(A\), such that
the points \(g'_1,\ldots,g'_p\) of
\[
  G'=G\otimes_AA'
\]
lying above \(g_1,\ldots,g_n\) are strictly rational over \(A'\).%
\footnote{\label{note:sga3-VIA-40}Editor's note: the original has
been expanded in what follows.}
Since the morphisms \(G'\to G\) and \(G\to X\) are open,
\[
  U'=U\otimes_AA'
\]
is dense in \(G'\). Hence the open subset
\[
  \bigcap_{i=1}^p (U')^{-1}\cdot g'_i
\]
is nonempty and therefore contains a closed point \(x\). By \hyperref[prop:sga3-VIA-1.2]{1.2}
(and \hyperref[sga3:VIA:0.4.1]{0.4.1}), after enlarging \(A'\) if necessary, we may suppose that
\(x\) is strictly rational over \(A'\). Since
\(x\in(U')^{-1}\cdot g'_i\), we have
\[
  g'_i\in U'\cdot x
  \qquad(1\leq i\leq p).
\]

Write \(V'\) for the inverse image of \(V\) in
\[
  X'=X\otimes_AA'.
\]
It is an affine open subset of \(X'\), and is also the image of \(U'\)
under the projection \(G'\to X'\). Since right translation \(r_x\)
is an automorphism of the groupoid \(G_*\otimes_AA'\), it induces an
automorphism, again denoted \(r_x\), of the quotient \(X'\).
Consequently,
\[
  V'\cdot x=r_x(V'),
\]
the image of \(U'\cdot x\) in \(X'\), is an affine open subset of
\(X'\) containing the images \(x'_1,\ldots,x'_p\) of
\(g'_1,\ldots,g'_p\).

Consider the equivalence relation on \(X'=X\otimes_AA'\) defined by
the projection \(X\otimes_AA'\to X\):
\begin{center}
  \phantomsection\label{diag:sga3-VIA-3.2.1-descent-equivalence}%
  \begin{tikzcd}[column sep=large]
    X\otimes_A A'\otimes_A A'
      \arrow[r,shift left=1.2ex,"d_1"]
      \arrow[r,shift right=1.2ex,"d_0"']
    &
    X\otimes_A A'
      \arrow[r]
    &
    X
  \end{tikzcd},
\end{center}
Here \(d_0\) and \(d_1\) are induced by the two canonical injections of
\(A'\) into \(A'\otimes_AA'\). Since \(A'\) is a finite free
\(A\)-algebra, say of rank \(n\), \(d_0\) and \(d_1\) are finite and
locally free of rank \(n\). We may therefore apply the argument of
\SGARefLink{par:sga3-V-5-b}{Expos\'e~V, 5.b}, taken from the proof of SGA~1, VIII.7.6. It follows
that \(x'_1,\ldots,x'_p\) are contained in a saturated affine open
subset \(W'\) contained in the affine open subset \(V'\cdot x\). The
image of \(W'\) in \(X\) then contains \(x_1,\ldots,x_n\), and is an
affine open subset of \(X\), by Expos\'e~V, 4.1(ii).

% END INCLUDED SGA3_English_current_buildable_source_20260823/source/inputs/VI/tex_reference_v2/components/08_expose_VIA_section321_quasisection_descent_loop2_r1_en.tex

% BEGIN INCLUDED SGA3_English_current_buildable_source_20260823/source/inputs/VI/tex_reference_v2/components/09_expose_VIA_section322_induction_en.tex
% SGA 3, Expose VI_A, Loop 1 English reconstruction.
% Scope: paragraph 3.2.2 in full; paragraph 3.2.3 is excluded.
% Authority: Polo--Gille Exp6A-13oct24.pdf, component-local page 23
% (combined-reader page 326; printed page 315).
% OCR and the Reinhold English Markdown are comparison witnesses only.

\medskip
\phantomsection\label{par:sga3-VIA-3.2.2}%
\noindent\textbf{3.2.2.---}
For every \(A\)-algebra \(A'\), local, finite, and free over \(A\),
let \(U(A')\) denote the set of points of \(G\otimes_AA'\) having an
open saturated neighborhood \(W\) such that the groupoid induced by
\(G_*\otimes_AA'\) on \(W\) has a quasi-section. It is clear that
\(U(A')\) is saturated under the operations of
\(G(\operatorname{Spec}A')\) on \(G\otimes_AA'\). We shall show that,
when \(A'\) is sufficiently large,
\[
  U(A')=G\otimes_AA'.
\]

By \SGARefLink{par:sga3-VIB-8.1}{Theorem~V.8.1}, \(U(A)\) is nonempty, and therefore contains a closed
point \(y\). We argue by induction on
\[
  \dim\bigl(G-U(A)\bigr).
\]
Let \(g_1,\ldots,g_n\) be closed points belonging to the various
irreducible components of \(G-U(A)\). By \SGARefLink{prop:sga3-VIA-1.2}{1.2}, there is an
\(A\)-algebra \(A'\), local, finite, and free over \(A\), such that
the points
\[
  g'_1,\ldots,g'_p
  \qquad
  \bigl(\text{respectively, }x=x_1,\ldots,x_r\bigr)
\]
of \(G'=G\otimes_AA'\), lying above \(g_1,\ldots,g_n\)
(respectively, above \(y\)), are strictly rational over \(A'\).
Then \(U(A')\) contains
\[
  \bigl(U(A)\otimes_AA'\bigr)\cdot x^{-1}g'_i
  \qquad\text{for every }i.
\]
Consequently, \(U(A')\) contains \(g'_1,\ldots,g'_p\), and
\[
  \dim\bigl(G'-U(A')\bigr)
  <
  \dim\bigl(G-U(A)\bigr).
\]

The induction hypothesis therefore gives an algebra \(A''\), local,
finite, and free over \(A'\), such that
\[
  U(A'')
  =
  G'\otimes_{A'}A''
  =
  G\otimes_AA''.
\]

% END INCLUDED SGA3_English_current_buildable_source_20260823/source/inputs/VI/tex_reference_v2/components/09_expose_VIA_section322_induction_en.tex

% BEGIN INCLUDED SGA3_English_current_buildable_source_20260823/source/inputs/VI/tex_reference_v2/components/10_expose_VIA_section323_descent_loop2_r1_en.tex
% SGA 3, Expose VI_A, Loop 1 English reconstruction.
% Scope: paragraph 3.2.3 in full; paragraph 3.2.4 is excluded.
% Authority: Polo--Gille Exp6A-13oct24.pdf, component-local pages 23--24
% (combined-reader pages 326--327; printed pages 315--316).
% OCR and the Reinhold English Markdown are comparison witnesses only.
% The commutative diagram is an exact 180-dpi crop of the controlling PDF:
% local p. 24, x=265,y=420,w=930,h=325 on a 1488x2105 render.

\medskip
\phantomsection\label{par:sga3-VIA-3.2.3}%
\noindent\textbf{3.2.3.---}
We are now in a position to prove the existence of
\(F\backslash G\) when \(F\) and \(G\) are of finite type over
\(A\). Let \(A'\) be sufficiently large over \(A\) that
\(U(A')\) coincides with \(G\otimes_AA'\) (cf.\ \hyperref[par:sga3-VIA-3.2.2]{3.2.2}). Set
\[
  A''=A'\otimes_AA',
\]
and, for every \(A\)-scheme \(X\), write \(X'\) and \(X''\) for
the fiber products \(X\otimes_AA'\) and \(X\otimes_AA''\).
By \hyperref[par:sga3-VIA-3.2.1]{3.2.1} and \hyperref[par:sga3-VIA-3.2.2]{3.2.2}, the quotients \(F'\backslash G'\) and
\(F''\backslash G''\) exist, and we have the following commutative
diagram, in which the first two rows and columns are exact:
\begin{center}
  \phantomsection\label{diag:sga3-VIA-3.2.3-exact-descent}%
  \begin{tabular}{c@{\qquad}c}
    \textup{(*)}
    &
    \begin{tikzcd}[row sep=large,column sep=large]
      F''\times_{A''}G''
        \arrow[r,shift left=1.2ex,"\operatorname{pr}_2''"]
        \arrow[r,shift right=1.2ex,"\lambda''"']
        \arrow[d,shift left=1.2ex,"w_2"]
        \arrow[d,shift right=1.2ex,"w_1"']
      &
      G''
        \arrow[r,"p''"]
        \arrow[d,shift left=1.2ex,"v_2"]
        \arrow[d,shift right=1.2ex,"v_1"']
      &
      F''\backslash G''
        \arrow[d,shift left=1.2ex,"u_2"]
        \arrow[d,shift right=1.2ex,"u_1"']
      \\
      F'\times_{A'}G'
        \arrow[r,shift left=1.2ex,"\operatorname{pr}_2'"]
        \arrow[r,shift right=1.2ex,"\lambda'"']
        \arrow[d,"h"']
      &
      G'
        \arrow[r,"p'"]
        \arrow[d,"g"]
      &
      F'\backslash G'
      \\
      F\times_A G
        \arrow[r,shift left=1.2ex,"\operatorname{pr}_2"]
        \arrow[r,shift right=1.2ex,"\lambda"']
      &
      G\mathpunct{.}
      &
      {}
    \end{tikzcd}
  \end{tabular}
\end{center}

In this diagram, \(\operatorname{pr}'_2\) and \(\lambda'\)
(respectively, \(\operatorname{pr}''_2\) and \(\lambda''\)) are
obtained from \(\operatorname{pr}_2\) and \(\lambda\) by the evident
base changes; the morphisms \(g\) and \(h\) are induced by the
canonical injection \(A\to A'\). We write \(p'\) and \(p''\) for the
canonical morphisms, while \(v_1,v_2\) and \(w_1,w_2\) are induced
by the two canonical injections of \(A'\) into \(A''\).

Finally, formation of the quotient \(F'\backslash G'\) commutes with
the two base changes
\[
  f_1,f_2:\operatorname{Spec}A''
    \rightrightarrows \operatorname{Spec}A'.
\]
Thus, if
\[
  \pi':F'\backslash G'\longrightarrow\operatorname{Spec}A'
\]
denotes the structure morphism, there are canonical isomorphisms,
for \(i=1,2\),
\[
  \tau_i:F''\backslash G''
    \xrightarrow{\ \sim\ }
    (F'\backslash G')\times_{\pi',\,f_i}\operatorname{Spec}A'',
\]
and \(u_i\) is the composite of \(\tau_i\) with the projection
\[
  (F'\backslash G')\times_{\pi',\,f_i}\operatorname{Spec}A''
    \longrightarrow F'\backslash G'.
\]

Now, for a diagram of type \((*)\) in which the first two rows and
columns are exact, one checks readily that
\(\operatorname{Coker}(\operatorname{pr}_2,\lambda)\) exists if and
only if \(\operatorname{Coker}(u_1,u_2)\) exists, and that these two
cokernels are then identified. The existence of \(F\backslash G\)
therefore follows from that of \(\operatorname{Coker}(u_1,u_2)\).

The compatibility of the formation of \(F\backslash G\) with the
base extensions considered here (cf.\ editor's
note~\ref{note:sga3-VIA-37} in \SGARefLink{par:sga3-VIA-3.2.1}{3.2.1} and \SGARefLink{par:sga3-VIA-4.6}{4.6} below) shows that the
composite
\begin{center}
  \phantomsection\label{diag:sga3-VIA-3.2.3-descent-datum}%
  \(
  (F'\backslash G')\times_{\pi',\,f_1}\operatorname{Spec}A''
    \xrightarrow{\ \tau_1^{-1}\ }
  F''\backslash G''
    \xrightarrow{\ \tau_2\ }
  (F'\backslash G')\times_{\pi',\,f_2}\operatorname{Spec}A''
  \)
\end{center}
is a descent datum on \(F'\backslash G'\) relative to
\[
  f:\operatorname{Spec}A'\longrightarrow\operatorname{Spec}A.
\]
By \hyperref[par:sga3-VIA-3.2.1]{3.2.1}~\eqref{eq:sga3-VIA-3.2.1-dagger} and
SGA~1, VIII.7.6, this descent datum is effective; that is,
\(\operatorname{Coker}(u_1,u_2)\) exists. (One could instead apply
\SGARefLink{thm:sga3-V-4.1}{Expos\'e~V, Theorem~4.1} directly.)

% END INCLUDED SGA3_English_current_buildable_source_20260823/source/inputs/VI/tex_reference_v2/components/10_expose_VIA_section323_descent_loop2_r1_en.tex

% BEGIN INCLUDED SGA3_English_current_buildable_source_20260823/source/inputs/VI/tex_reference_v2/components/11_expose_VIA_sections324_325_loop2_r1_en.tex
% SGA 3, Expose VI_A, Loop 1 English reconstruction.
% Scope: paragraphs 3.2.4 and 3.2.5 in full; subsection 3.3 is excluded.
% Authority: Polo--Gille Exp6A-13oct24.pdf, component-local pages 24--25
% (combined-reader pages 327--328; printed pages 316--317).
% OCR and the Reinhold English Markdown are comparison witnesses only.
% The Cartesian square is an exact 180-dpi crop of the controlling PDF:
% local p. 25, x=540,y=575,w=450,h=145 on a 1488x2105 render.

\medskip
\phantomsection\label{par:sga3-VIA-3.2.4}%
\noindent\textbf{3.2.4.---}
To complete the proof of assertions (i), (ii), (iii), and (iv)(a) of
\hyperref[thm:sga3-VIA-3.2]{3.2} in the case where \(F\) and \(G\) are of finite type over \(A\),
it remains to study the quotient \(F\backslash G\). By
\SGARefLink{lem:sga3-V-6.1}{Expos\'e~V, 6.1}, assertions \SGARefLink{thm:sga3-VIA-3.2:item:ii}{(ii)}, \SGARefLink{thm:sga3-VIA-3.2:item:iii}{(iii)}, and \SGARefLink{concl:sga3-VIA-3.2-iv}{(iv)}\SGARefLink{thm:sga3-VIA-3.2:item:a}{(a)} ``become true''
after the base change
\[
  f:\operatorname{Spec}A'\longrightarrow\operatorname{Spec}A.
\]
By EGA~IV\textsubscript{2}, 2.6.1, 2.6.2, and 2.7.1, these assertions
were therefore true before the base change. Finally, to prove the
second assertion of (i), namely that \(F\backslash G\) is the cokernel
of \((\operatorname{pr}_2,\lambda)\) in the category of all ringed
spaces, it suffices to refer to Expos\'e~V, \S6(c).

\medskip
\phantomsection\label{par:sga3-VIA-3.2.5}%
\noindent\textbf{3.2.5.---}%
\footnote{\label{note:sga3-VIA-41}Editor's note: we have added this
paragraph.}
We now prove assertion (iv)(c) of \SGARefLink{thm:sga3-VIA-3.2}{3.2}, following the proof of
Expos\'e~VI\textsubscript{B}, \SGARefLink{lem:sga3-VIB-9.2.1}{9.2.1}. Write
\[
  X=F\backslash G
\]
and let \(d\) be the morphism
\[
  F\times_A G\longrightarrow G\times_A G
\]
whose components are \(\lambda\) and \(\operatorname{pr}_2\).

Since \(u\) is a closed immersion by \hyperref[prop:sga3-VIA-2.5.2]{2.5.2}, and since
\[
  d=\sigma\circ(u\times\operatorname{id}_G),
  \qquad
  \sigma(x,y)=(xy,y),
\]
where \(\sigma\) is an automorphism of \(G\times_A G\), the morphism
\(d\) is a closed immersion. On the other hand, by (iv)\SGARefLink{prop:sga3-VIA-2.5.2:item:a}{(a)}, the
following square is Cartesian:
\begin{center}
  \phantomsection\label{diag:sga3-VIA-3.2.5-cartesian-square}%
  \begin{tikzcd}[row sep=large,column sep=large]
    F\times_A G \arrow[r,"d"] \arrow[d] &
    G\times_A G \arrow[d,"p\times p"] \\
    X \arrow[r,"\Delta_X"'] &
    X\times_A X
  \end{tikzcd}
\end{center}
Moreover, \(p\), and hence also \(p\times p\), is faithfully flat and
locally of finite presentation. Therefore, by fppf descent, since
\(d\) is a closed immersion, so is \(\Delta_X\); that is, \(X\) is
separated.

% END INCLUDED SGA3_English_current_buildable_source_20260823/source/inputs/VI/tex_reference_v2/components/11_expose_VIA_sections324_325_loop2_r1_en.tex

% BEGIN INCLUDED SGA3_English_current_buildable_source_20260823/source/inputs/VI/tex_reference_v2/components/12_expose_VIA_section33_theorems_corollary_loop2_r1_en.tex
% SGA 3, Expose VI_A, Loop 1 English reconstruction.
% Scope: section 3.3 in full, through Corollary 3.3.3;
% section 4 is excluded.
% Authority: Polo--Gille Exp6A-13oct24.pdf, component-local pages 25--26
% (combined-reader pages 328--329; printed pages 317--318).
% OCR and the Reinhold English Markdown are comparison witnesses only.
% The exact sequence is an exact 180-dpi crop of the controlling PDF:
% local p. 25, x=500,y=1546,w=540,h=78 on a 1488x2105 render.

\medskip
\phantomsection\label{par:sga3-VIA-3.3}%
\noindent\textbf{3.3.---}%
\footnote{\label{note:sga3-VIA-42}Editor's note: we have added this
subsection.}
In \hyperref[thm:sga3-VIA-3.2]{Theorem~3.2}, the hypothesis that \(G\) be flat can be omitted when
\(u:F\to G\) is a monomorphism. This generalization is mentioned in
Expos\'e~VI\textsubscript{B}, Remark~9.3(b), and also in
\textup{[Ray67a]}, Example~a)\SGARefLink{thm:sga3-VIA-3.2:item:i}{(i)}, p.~82. The proof, given in
\textup{[An73]}, \SGARefLink{sec:sga3-VIA-4}{Theorem~4}, follows from \hyperref[thm:sga3-VIA-3.2]{Theorem~3.2} and from the
following theorem of Grothendieck (mentioned in \textup{[Ray67a]},
Theorem~1(ii), and proved in \textup{[DG70]}, \S~III.2, \SGARefLink{prop:sga3-VIB-7.1}{7.1}).
If \(X\) is a scheme and \(R\) an equivalence relation on \(X\), we
write \(\widetilde{X/R}\) for the \((\mathrm{fppf})\) quotient sheaf
of \(X\) by \(R\) (cf.\ \SGARefLink{sga3:iv:proposition:4:4:9}{Expos\'e~IV, 4.4.9}).

\medskip
\phantomsection\label{thm:sga3-VIA-3.3.1}%
\noindent\textbf{Theorem 3.3.1 (Grothendieck).}
\textit{Let \(A\) be a ring, let \(X\) be an \(A\)-scheme, and let
\[
  d_0,d_1:R\rightrightarrows X
\]
be an \(A\)-equivalence relation on \(X\), with each \(d_i\)
faithfully flat and of finite presentation. Let \(X_0\) be a saturated
closed subscheme of \(X\), defined by a nilpotent ideal, and let
\(R_0\) be the equivalence relation induced by \(R\) on \(X_0\).
If the \((\mathrm{fppf})\) quotient sheaf
\(\widetilde{X_0/R_0}\) is representable by an \(A\)-scheme, then so
is \(\widetilde{X/R}\).}

For the proof, see \textup{[DG70]}, \S~III.2, \SGARefLink{prop:sga3-VIB-7.1}{7.1}. Let us now return
to the case where \(A\) is an Artinian local ring. Let
\(u:F\to G\) be a quasi-compact morphism between \(A\)-groups locally
of finite type, and suppose in addition that \(u\) is a monomorphism.
Then \(u\) is a closed immersion by \SGARefLink{prop:sga3-VIA-2.5.2}{2.5.2}.

We may now state the following variant of \hyperref[thm:sga3-VIA-3.2]{Theorem~3.2}.

\medskip
\phantomsection\label{thm:sga3-VIA-3.3.2}%
\noindent\textbf{Theorem 3.3.2.}
\textit{Let \(A\) be an Artinian local ring, let \(G\) be an
\(A\)-group locally of finite type, and let \(F\) be a closed
subgroup of \(G\), flat over \(A\).}%
\footnote{\label{note:sga3-VIA-43}Editor's note: for an example in
which \(F\) is not flat and \(\widetilde{F\backslash G}\) is not
representable, see \textup{[DG70]}, \S~III.3, no.~\SGARefLink{prop:sga3-VIB-3.3}{3.3}.}
\textit{Then:}
\begin{enumerate}
\item[\SGAInlineRefTarget{thm:sga3-VIA-3.3.2:item:i}(i)] \textit{The \((\mathrm{fppf})\) quotient sheaf
\(\widetilde{F\backslash G}\) is representable by a separated
\(A\)-scheme \(F\backslash G\), locally of finite type; moreover, the
sequence}
\begin{center}
  \phantomsection\label{diag:sga3-VIA-3.3.2-exact-sequence}%
  \begin{tikzcd}[column sep=large]
    F\times_A G
      \arrow[r,shift left=1.2ex]
      \arrow[r,shift right=1.2ex]
    &
    G
      \arrow[r,"p"]
    &
    F\backslash G
  \end{tikzcd}
\end{center}
\textit{is exact in the category of all ringed spaces.}

\item[\SGAInlineRefTarget{thm:sga3-VIA-3.3.2:item:ii}(ii)] \textit{The morphism
\[
  F\times_A G
    \xrightarrow{(\lambda,\operatorname{pr}_2)}
  G\times_{(F\backslash G)}G
\]
is an isomorphism, and
\(p:G\to F\backslash G\) is faithfully flat and locally of finite
presentation; hence \(p\) is an \(F\)-torsor locally trivial for the
\((\mathrm{fppf})\) topology.}

\item[\SGAInlineRefTarget{thm:sga3-VIA-3.3.2:item:iii}(iii)] \textit{If \(G\) is flat (respectively, of finite type;
respectively, smooth) over \(A\), then so is \(F\backslash G\).}

\item[(iv)]\phantomsection\label{concl:sga3-VIA-3.3.2-iv} \textit{\(X=F\backslash G\) is equipped with a right
action of \(G\) such that \(p(e)\cdot g=p(g)\) for every \(g\in G\).
Consequently, the connected components of \(X\) are of finite type
over \(A\), irreducible, and all of dimension
\(\dim G-\dim F\).}

\item[(v)]\phantomsection\label{concl:sga3-VIA-3.3.2-v} \textit{If, in addition, \(F\) is a normal subgroup of
\(G\), there is a unique \(A\)-group structure on \(F\backslash G\)
for which \(p:G\to F\backslash G\) is a morphism of \(A\)-groups.}
\end{enumerate}

Assertions \SGARefLink{thm:sga3-VIA-3.3.2:item:i}{(i)} and \SGARefLink{thm:sga3-VIA-3.3.2:item:ii}{(ii)} follow from \SGARefLink{thm:sga3-VIA-3.2}{3.2} and \SGARefLink{thm:sga3-VIA-3.3.1}{3.3.1}. Since
\(G\to F\backslash G\) is faithfully flat and locally of finite
presentation, assertion (iii) follows from EGA~IV, 2.2.14, 2.7.1,
and 17.7.7. Assertions (iv) and (v) will be proved in \hyperref[sec:sga3-VIA-5]{Section~5}.
Let us record at once the following corollary.

\medskip
\phantomsection\label{cor:sga3-VIA-3.3.3}%
\noindent\textbf{Corollary 3.3.3.}
\textit{Let \(A\) be an Artinian local ring, let \(G\) be an
\(A\)-group locally of finite type, and let \(H\) be a closed
subgroup of \(G\), flat over \(A\). Write \(p\) for the morphism
\(G\to G/H\), and write \(\lambda\) (respectively,
\(\operatorname{pr}_1\)) for the morphism \(G\times H\to G\) defined
by \(\lambda(g,h)=gh\) (respectively, the projection
\(G\times H\to G\)). Then, for every open subset \(U\) of \(G/H\),
\[
  \mathcal O(U)
  =
  \bigl\{
    \varphi\in\mathcal O\bigl(p^{-1}(U)\bigr)
    \,\bigm|\,
    \varphi\circ\lambda=\varphi\circ\operatorname{pr}_1
  \bigr\}.
\]
In other words, \(\mathcal O(U)\) is the set of
\(\varphi\in\mathcal O(p^{-1}(U))\) such that
\(\varphi(gh)=\varphi(g)\) for every \(A\)-scheme \(S\) and every
\(g\in G(S)\), \(h\in H(S)\).}

Indeed, \(p:G\to G/H\) is faithfully flat and locally of finite
presentation, hence a covering for the \((\mathrm{fppf})\) topology;
the result follows from \SGARefLink{sga3:iv:remark:3:3:3:2}{Expos\'e~IV, 3.3.3.2}.

% END INCLUDED SGA3_English_current_buildable_source_20260823/source/inputs/VI/tex_reference_v2/components/12_expose_VIA_section33_theorems_corollary_loop2_r1_en.tex

% BEGIN INCLUDED SGA3_English_current_buildable_source_20260823/source/inputs/VI/tex_reference_v2/components/13_expose_VIA_section4_opening_par41_en.tex
% SGA 3, Expose VI_A, Loop 1 English reconstruction.
% Scope: section 4 heading and introductory paragraph, followed by
% paragraph 4.1 in full; paragraph 4.2 is excluded.
% Authority: Polo--Gille Exp6A-13oct24.pdf, component-local pages 26--27
% (combined-reader pages 329--330; printed pages 318--319).
% OCR and the Reinhold English Markdown are comparison witnesses only.

\section*{4. Construction of quotients \(F\backslash G\)
  (general case)}
\phantomsection\label{sec:sga3-VIA-4}
\addcontentsline{toc}{section}{4. Construction of quotients
  \texorpdfstring{\(F\backslash G\)}{F backslash G}
  (general case)}

We now assume the hypotheses of \hyperref[thm:sga3-VIA-3.2]{Theorem~3.2}, without requiring
\(F\) and \(G\) to be of finite type over \(A\).

\medskip
\phantomsection\label{par:sga3-VIA-4.1}%
\noindent\textbf{4.1.---}
Let \(G^\alpha\) be a connected component of \(G\). We first show
that the saturation
\[
  \mathscr S(G^\alpha)
\]%
\footnote{\label{note:sga3-VIA-44}Editor's note: we have written
\(\mathscr S(G^\alpha)\), rather than \(\overline{G^\alpha}\), for
the saturation of \(G^\alpha\).}
of \(G^\alpha\) under the equivalence relation defined by the
groupoid \(G_*\) is an open-and-closed subset of \(G\); in other
words, it is a union of connected components of \(G\).

This saturation is the image of \(F\times_A G^\alpha\) under
\(\lambda\), and is therefore open in \(G\) (cf.\ \S~\hyperref[par:sga3-VIA-3.1]{3.1}). Let
\(k\) be the residue field of \(A\), and let \(\bar k\) be an
algebraic closure of \(k\). It remains to show that the image of
\[
  (F\times_A G^\alpha)\otimes_A k
\]
under \(\lambda\otimes_A k\) is closed in \(G\otimes_A k\), or,
equivalently by SGA~1, VIII.4.4, that the image of
\[
  (F\times_A G^\alpha)\otimes_A\bar k
\]
under \(\lambda\otimes_A\bar k\) is closed in
\(G\otimes_A\bar k\). Since \(G^\alpha\otimes_A\bar k\) is a union
of finitely many connected components of \(G\otimes_A\bar k\), we
are reduced to the case where \(A\) is an algebraically closed field,
which we now assume. In this case, \(\mathscr S(G^\alpha)\) is the
union of the images of \(G^\alpha\) under the left translations
\(\ell_{u(x)}\), where \(x\) runs through the closed points of \(F\).
The assertion follows because these images are connected components
of \(G\).

% END INCLUDED SGA3_English_current_buildable_source_20260823/source/inputs/VI/tex_reference_v2/components/13_expose_VIA_section4_opening_par41_en.tex

% BEGIN INCLUDED SGA3_English_current_buildable_source_20260823/source/inputs/VI/tex_reference_v2/components/14_expose_VIA_par42_saturation_identity_component_en.tex
% SGA 3, Expose VI_A, Loop 1 English reconstruction.
% Scope: paragraph 4.2 in full; paragraph 4.3 is excluded.
% Authority: Polo--Gille Exp6A-13oct24.pdf, component-local page 27
% (combined-reader page 330; printed page 319).
% OCR and the Reinhold English Markdown are comparison witnesses only.

\medskip
\phantomsection\label{par:sga3-VIA-4.2}%
\noindent\textbf{4.2.---}
Take \(G^\alpha\) in particular to be the connected component
\(G^0\) of the identity in \(G\). Then \(\mathscr S(G^0)\) plainly
contains the image of \(F\) under \(u\), which is precisely the
equivalence class of the identity. On the other hand, if \(F^\beta\)
is a connected component of \(F\), then
\(F^\beta\times_A G^0\) is connected by \hyperref[lem:sga3-VIA-2.1.2]{2.1.2}, so its image under
\(\lambda\) is contained in the connected component of
\(u(F^\beta)\) in \(G\). In other words,
\(\mathscr S(G^0)\) is the union of the connected components that
meet the image of \(F\).

Let us also observe that the open subscheme of \(G\) whose underlying
space is \(\mathscr S(G^0)\) is a subgroup of \(G\); we continue to
denote it by \(\mathscr S(G^0)\). Indeed, inversion on \(G\)
preserves the image of \(F\) and permutes the connected components of
\(G\) that meet this image. It therefore suffices to prove that
\[
  \nu:G\times_A G\longrightarrow G
\]
maps
\(\mathscr S(G^0)\times_A\mathscr S(G^0)\) into
\(\mathscr S(G^0)\). For this we may suppose that \(A\) is an
algebraically closed field: with the notation of \hyperref[par:sga3-VIA-4.1]{4.1},
\(\mathscr S(G^0)\otimes_A\bar k\) is the saturation of
\((G\otimes_A\bar k)^0\) for the equivalence relation defined by the
homomorphism \(u\otimes_A\bar k\).

If \(G^\gamma\) and \(G^\delta\) are connected components of
\(\mathscr S(G^0)\), then \(G^\gamma\times_A G^\delta\) is connected
and its image under \(\nu\) meets the image of \(F\). Consequently,
\[
  \nu(G^\gamma\times_A G^\delta)
\]
is contained in a connected component of \(G\) that meets
\(u(F)\).\begingroup
\renewcommand{\thefootnote}{\fnsymbol{footnote}}%
\footnote[1]{Translator's note: the French source prints
\(u(G^\gamma\times_A G^\delta)\) in this final sentence. The map
defined and used throughout the paragraph is the multiplication
\(\nu:G\times_A G\to G\); the target therefore uses \(\nu\).}%
\endgroup

% END INCLUDED SGA3_English_current_buildable_source_20260823/source/inputs/VI/tex_reference_v2/components/14_expose_VIA_par42_saturation_identity_component_en.tex

% BEGIN INCLUDED SGA3_English_current_buildable_source_20260823/source/inputs/VI/tex_reference_v2/components/15_expose_VIA_par43_identity_saturation_quotient_loop2_r1_en.tex
% SGA 3, Expose VI_A, Loop 1 English reconstruction.
% Scope: paragraph 4.3 in full; paragraph 4.4 is excluded.
% Authority: Polo--Gille Exp6A-13oct24.pdf, component-local pages 27--28
% (combined-reader pages 330--331; printed pages 319--320).
% OCR and the Reinhold English Markdown are comparison witnesses only.
% Loop 2 replaces the four exact 180-dpi source crops with native
% tikz-cd reconstructions controlled against the same PDF coordinates:
% local p. 27, x=490,y=1235,w=520,h=145;
% local p. 28, x=470,y=525,w=550,h=80;
% local p. 28, x=470,y=695,w=550,h=135;
% local p. 28, x=420,y=920,w=650,h=78;
% all source renders are 1488x2105 pixels.

\medskip
\phantomsection\label{par:sga3-VIA-4.3}%
\noindent\textbf{4.3.---}
It follows from the preceding discussion that the groupoid \(G_*\),
with base \(G\), defined by \(u\), is the direct sum of the groupoids
\(\mathscr S(G^\alpha_*)\) induced by \(G_*\) on the various
open-and-closed parts of \(G\) of the form
\(\mathscr S(G^\alpha)\). The cokernel of \(G_*\) is therefore the
direct sum of the cokernels of the groupoids
\(\mathscr S(G^\alpha_*)\), which we are thus led to study
separately.

First consider the groupoid \(\mathscr S(G^0_*)\) induced by \(G_*\)
on \(\mathscr S(G^0)\). It is clear that
\(\mathscr S(G^0_*)\) is the groupoid with base
\(\mathscr S(G^0)\) defined by the homomorphism
\[
  F\longrightarrow\mathscr S(G^0)
\]
induced by \(u\) (\S~\hyperref[par:sga3-VIA-3.1]{3.1}). The cokernel whose existence we wish to
prove is thus identified with \(F\backslash\mathscr S(G^0)\).
Consider, on the other hand, the groupoid
\begin{center}
  \phantomsection\label{diag:sga3-VIA-4.3-G0-groupoid}%
  \begin{tikzcd}[column sep=large]
    G^0_2
      \arrow[r,shift left=1.4ex,"\ell'_2"]
      \arrow[r,"\ell'_1"]
      \arrow[r,shift right=1.4ex,"\ell'_0"'] &
    G^0_1
      \arrow[r,shift left=0.8ex,"\ell_1"]
      \arrow[r,shift right=0.8ex,"\ell_0"'] &
    G^0_0=G^0
  \end{tikzcd}
\end{center}
induced by \(\mathscr S(G^0_*)\) on \(G^0\). Referring to the
explicit construction in Expos\'e~V, \S~3.b, the object denoted there
by \(Y_0\times_{X_0}X_1\) is precisely \(F\times_A G^0\). Hence
\(G^0_1\) is the inverse image of \(G^0\) under the morphism
\[
  F\times_A G^0\longrightarrow\mathscr S(G^0)
\]
induced by \(\lambda\).

I claim that this inverse image is \(F_0\times_A G^0\), where
\(F_0\) denotes the inverse image of \(G^0\) under \(u\). Indeed, if
\(F^\beta\) is a connected component of \(F_0\), then
\(F^\beta\times_A G^0\) is connected by \SGARefLink{lem:sga3-VIA-2.1.2}{2.1.2} and
\(\lambda(F^\beta\times_A G^0)\) is contained in \(G^0\).%
\begingroup
\renewcommand{\thefootnote}{\fnsymbol{footnote}}%
\footnote[1]{Translator's note: the French source prints
\(\lambda(F\times_A G^0)\) at this point. The statement concerns the
connected component \(F^\beta\) just introduced, and the converse
immediately below confirms the intended restriction to
\(F^\beta\times_A G^0\).}%
\endgroup
Conversely, if \(F^\beta\) is a connected component of \(F\) not
contained in \(F_0\), then the image of
\(F^\beta\times_A G^0\) is again connected and contains
\(u(F^\beta)\); if \(u(F^\beta)\) is not contained in \(G^0\), then
\(\lambda(F^\beta\times_A G^0)\) does not meet \(G^0\).

It follows that the groupoid \(G^0_*\) induced by \(G_*\) on \(G^0\)
is the groupoid with base \(G^0\) defined by the homomorphism
\[
  F_0\longrightarrow G^0
\]
induced by \(u\). Since \(G^0\), and hence \(F_0\), are of finite
type over \(A\), \hyperref[sec:sga3-VIA-3]{Section~3} shows that \(G^0_*\) has a cokernel,
namely \(F_0\backslash G^0\).%
\begingroup
\renewcommand{\thefootnote}{\fnsymbol{footnote}}%
\footnote[2]{Translator's note: the French source says
``by \hyperref[sec:sga3-VIA-4]{paragraph 4}.'' The finite-type quotient invoked here is the
result of \hyperref[sec:sga3-VIA-3]{Section~3}; \hyperref[sec:sga3-VIA-4]{Section~4} is the general-case construction
currently in progress.}%
\endgroup

I now claim that \(F_0\backslash G^0\) is identified with
\(F\backslash\mathscr S(G^0)\). The proof is analogous to that of
the first part of assertion~(i) of \SGARefLink{lem:sga3-V-6.1}{Expos\'e~V, Lemma~6.1}. Consider
the diagram
\begin{center}
  \phantomsection\label{diag:sga3-VIA-4.3-span}%
  \begin{tikzcd}[column sep=large]
    \mathscr S(G^0) &
    F\times_A G^0
      \arrow[l,"v"']
      \arrow[r,"\operatorname{pr}_2"] &
    G^0
  \end{tikzcd},
\end{center}
where \(v\) is the morphism induced by \(\lambda\). Since
\(\operatorname{pr}_2\) has a section, it is a universal effective
epimorphism. Consequently, \(F_0\backslash G^0\) coincides with
\(\operatorname{Coker}(v_0,v_1)\), where
\begin{center}
  \phantomsection\label{diag:sga3-VIA-4.3-V-groupoid}%
  \begin{tikzcd}[column sep=large]
    V_2
      \arrow[r,shift left=1.4ex,"v'_2"]
      \arrow[r,"v'_1"]
      \arrow[r,shift right=1.4ex,"v'_0"'] &
    V_1
      \arrow[r,shift left=0.8ex,"v_1"]
      \arrow[r,shift right=0.8ex,"v_0"'] &
    V=F\times_A G^0
  \end{tikzcd}
\end{center}
is the inverse image under \(\operatorname{pr}_2\) of the groupoid
\(G^0_*\) (cf.\ Expos\'e~V, \S~3.a), and hence also the inverse image
of \(\mathscr S(G^0_*)\) under the composite morphism
\begin{center}
  \phantomsection\label{diag:sga3-VIA-4.3-composite}%
  \begin{tikzcd}[column sep=large]
    F\times_A G^0
      \arrow[r,"\mathrm{inclusion}"] &
    F\times_A\mathscr S(G^0)
      \arrow[r,"\operatorname{pr}_2"] &
    \mathscr S(G^0)
  \end{tikzcd}.
\end{center}

Similarly, since \(v\) is faithfully flat and quasi-compact,
\(F\backslash\mathscr S(G^0)\) coincides with the cokernel of the
inverse image of \(\mathscr S(G^0_*)\) under the base change \(v\).
This inverse image is isomorphic to \(V_*\) by Expos\'e~V, \S~3.c.
It follows that the canonical inclusion of \(G^0_*\) in
\(\mathscr S(G^0_*)\) induces an isomorphism
\[
  F_0\backslash G^0
    \xrightarrow{\ \sim\ }
  F\backslash\mathscr S(G^0).
\]

Finally, formation of \(F\backslash\mathscr S(G^0)\) commutes with
finite locally free base changes, because the same holds for
\(F_0\backslash G^0\) (cf.\ editor's
note~\ref{note:sga3-VIA-37} and \SGARefLink{par:sga3-VIA-4.6}{4.6} below).

% END INCLUDED SGA3_English_current_buildable_source_20260823/source/inputs/VI/tex_reference_v2/components/15_expose_VIA_par43_identity_saturation_quotient_loop2_r1_en.tex

% BEGIN INCLUDED SGA3_English_current_buildable_source_20260823/source/inputs/VI/tex_reference_v2/components/16_expose_VIA_par44_arbitrary_component_descent_en.tex
% SGA 3, Expose VI_A, Loop 1 English reconstruction.
% Scope: paragraph 4.4 in full; paragraph 4.5 is excluded.
% Authority: Polo--Gille Exp6A-13oct24.pdf, component-local page 28
% (combined-reader page 331; printed page 320).
% OCR and the Reinhold English Markdown are comparison witnesses only.

\medskip
\phantomsection\label{par:sga3-VIA-4.4}%
\noindent\textbf{4.4.---}
It remains to construct the cokernel of the groupoid
\(\mathscr S(G^\alpha_*)\) when \(G^\alpha\) is an arbitrary
connected component of \(G\). Let \(A'\) be a sufficiently large
local \(A\)-algebra, finite and free over \(A\) (cf.\ \SGARefLink{thm:sga3-VIA-3.2}{3.2}). Then
\(G^\alpha\otimes_AA'\) is the union of finitely many connected
components
\[
  C^1,\ldots,C^n
\]
of \(G\otimes_AA'\), each of which has a strictly rational point.
For every \(i\), there is therefore a right translation \(r_i\) of
\(G\otimes_AA'\) which maps \(G^0\otimes_AA'\) onto \(C^i\).
This translation induces an isomorphism
\[
  \mathscr S(G^0_*)\otimes_AA'
    \xrightarrow{\ \sim\ }
  \mathscr S(C^i_*).
\]
Consequently, the groupoid induced by \(G_*\otimes_AA'\) on the
saturation of \(C^i\) has a cokernel.

Since \(\mathscr S(G^\alpha_*)\otimes_AA'\) is the direct sum of
some of the groupoids \(\mathscr S(C^i_*)\), it too has a cokernel.
This cokernel is a direct sum of some number of copies of
\[
  (F_0\otimes_AA')\backslash(G^0\otimes_AA').
\]
Thus every finite subset of it is contained in an affine open subset.
Moreover, formation of this cokernel commutes with finite locally
free extensions of the base (cf.\ editor's
note~\ref{note:sga3-VIA-37} and \SGARefLink{par:sga3-VIA-4.6}{4.6} below). As in \SGARefLink{par:sga3-VIA-3.2.3}{3.2.3}, it follows
that this cokernel has the form \(Y\otimes_AA'\), where \(Y\) is a
cokernel of \(\mathscr S(G^\alpha_*)\).

% END INCLUDED SGA3_English_current_buildable_source_20260823/source/inputs/VI/tex_reference_v2/components/16_expose_VIA_par44_arbitrary_component_descent_en.tex

% BEGIN INCLUDED SGA3_English_current_buildable_source_20260823/source/inputs/VI/tex_reference_v2/components/17_expose_VIA_par45_46_base_change_en.tex
% SGA 3, Expose VI_A, Loop 1 English reconstruction.
% Scope: paragraphs 4.5 and 4.6 in full, through Proposition 4.6.1;
% section 5 is excluded.
% Authority: Polo--Gille Exp6A-13oct24.pdf, component-local page 29
% (combined-reader page 332; printed page 321).
% OCR and the Reinhold English Markdown are comparison witnesses only.

\medskip
\phantomsection\label{par:sga3-VIA-4.5}%
\noindent\textbf{4.5.---}
We have now constructed \(F\backslash G\) and shown that it is a
direct sum of schemes of finite type over \(A\). The remaining
assertions of \hyperref[thm:sga3-VIA-3.2]{Theorem~3.2} reduce directly to assertions about the
groupoids \(\mathscr S(G^\alpha_*)\). As in Expos\'e~V, \S~6, the
second assertion in \SGARefLink{thm:sga3-VIA-3.2:item:i}{(i)} follows from the first together with \SGARefLink{thm:sga3-VIA-3.2:item:ii}{(ii)} and
\SGARefLink{thm:sga3-VIA-3.2:item:iii}{(iii)}; it is therefore enough to prove \SGARefLink{thm:sga3-VIA-3.2:item:ii}{(ii)}, (iii), and (iv)(a).

Since \(A'\) is a local \(A\)-algebra, finite and free over \(A\), the
morphism \(A\to A'\) is faithfully flat and of finite presentation.
By SGA~1, VIII (3.1, \SGARefLink{par:sga3-VIA-4.6}{4.6}, 5.4), it is therefore enough to verify the
corresponding assertions for the groupoid
\[
  \mathscr S(G^\alpha_*)\otimes_AA'.
\]
This groupoid is isomorphic to a direct sum of finitely many copies of
\(\mathscr S(G^0_*)\otimes_AA'\) (cf.\ \hyperref[par:sga3-VIA-4.4]{4.4}), so we are reduced to the
groupoid \(\mathscr S(G^0_*)\). For this last groupoid we continue to
follow the proof given in Expos\'e~V, \S~6, as we began to do in \SGARefLink{par:sga3-VIA-4.3}{4.3}.

\medskip
\phantomsection\label{par:sga3-VIA-4.6}%
\noindent\textbf{4.6.---}
Let us conclude this section with some remarks concerning \SGARefLink{lem:sga3-VIA-6.1}{Lemma~6.1}
and \S~9(a) of Expos\'e~V. Under the hypotheses and with the notation
of Expos\'e~V, \S~9(a), we seek a condition under which formation of
the cokernel of the \((\mathrm{Sch}/S)\)-groupoid \(X_*\) commutes
with a base extension
\[
  \pi:S'\longrightarrow S.
\]
The cokernels of \(X_*\) and \(X'_*\) are identified with the
cokernels of the groupoids \(U_*\) and \(U'_*\) induced by \(X_*\)
and \(X'_*\) on the quasi-sections \(U\) and \(U'\), respectively.
We are therefore reduced to the case of a
\((\mathrm{Sch}/S)\)-groupoid satisfying the hypotheses of
\SGARefLink{thm:sga3-V-4.1}{Expos\'e~V, Theorem~4.1}.

\footnote{\label{note:sga3-VIA-45}Editor's note: we have added the
following sentence.}%
Let \(Y\) be the cokernel of \(U_*\), put
\[
  Y'=Y\times_SS',
\]
and let \(Y_1\) be the cokernel of \(U'_*\); we saw in
Expos\'e~V, \S~9(a), that the canonical morphism
\[
  Y_1\longrightarrow Y'
\]
is a homeomorphism (indeed, a universal homeomorphism); we may thus
identify \(Y_1\) and \(Y'\) as topological spaces. If
\(p:U\to Y\) is the canonical morphism and
\(p':U'\to Y'\) is obtained from it by base change, we want the
following sequence of \(\mathcal O_{Y'}\)-modules to be exact:
\[
  \tag{*}\label{eq:sga3-VIA-4.6-star}
  \mathcal O_{Y'}
    \longrightarrow p'_*(\mathcal O_{U'})
    \rightrightarrows
  p'_*u'_{1*}(\mathcal O_{U'_1})
    =
  p'_*u'_{0*}(\mathcal O_{U'_1}).
\]
\footnote{\label{note:sga3-VIA-46}Editor's note: in what follows we
have modified the original; the additional hypotheses imposed there
on \(X_*\) are unnecessary.}%
Since we are under the hypotheses of \SGARefLink{thm:sga3-V-4.1}{Expos\'e~V, Theorem~4.1},
\(u_0\) and \(u_1\) are finite and locally free; moreover,
\(p\) is integral by \SGARefLink{thm:sga3-V-4.1:item:ii}{Expos\'e~V, 4.1(ii)}.
Thus \(p\) and \(p\circ u_i\) are affine, hence separated and
quasi-compact.

Consequently, if \(S'\) is flat over \(S\), EGA~III\textsubscript{1},
1.4.15 (taking account of correction Err\textsubscript{III}~25 in
EGA~III\textsubscript{2}) identifies the sequence
\eqref{eq:sga3-VIA-4.6-star} with the inverse image of the sequence
\[
  \tag{**}\label{eq:sga3-VIA-4.6-double-star}
  \mathcal O_Y
    \longrightarrow p_*(\mathcal O_U)
    \rightrightarrows
  p_*u_{1*}(\mathcal O_{U_1})
    =
  p_*u_{0*}(\mathcal O_{U_1}),
\]
which is exact.%
\footnote{\label{note:sga3-VIA-47}Editor's note: we have added both
the following sentence and the numbering \SGARefLink{prop:sga3-VIA-4.6.1}{4.6.1} below, in order to
highlight the stated result.}
An analogous argument applies when the groupoid
\(X_*\) has quasi-sections ``locally'' (cf.\ the proof of
\SGARefLink{thm:sga3-V-7.1}{Expos\'e~V, Theorem~7.1}).
We therefore obtain:

\medskip
\phantomsection\label{prop:sga3-VIA-4.6.1}%
\noindent\textbf{Proposition 4.6.1.}
\textit{Formation of the cokernel of \(X_*\) commutes with flat base
extensions when \(X_*\) has quasi-sections locally.}

% END INCLUDED SGA3_English_current_buildable_source_20260823/source/inputs/VI/tex_reference_v2/components/17_expose_VIA_par45_46_base_change_en.tex

% BEGIN INCLUDED SGA3_English_current_buildable_source_20260823/source/inputs/VI/tex_reference_v2/components/18_expose_VIA_par47_flat_base_change_en.tex
% SGA 3, Expose VI_A, Loop 1 English reconstruction.
% Scope: paragraph 4.7 in full, through Proposition 4.7.1;
% section 5 is excluded.
% Authority: Polo--Gille Exp6A-13oct24.pdf, component-local page 30
% (combined-reader page 333; printed page 322).
% OCR and the Reinhold English Markdown are comparison witnesses only.

\medskip
\phantomsection\label{par:sga3-VIA-4.7}%
\noindent\textbf{4.7.---}
Consider the groupoid \(G_*\) of \hyperref[thm:sga3-VIA-3.2]{Theorem~3.2}, supposing temporarily
that \(F\) and \(G\) are of finite type over \(A\). By \SGARefLink{par:sga3-VIA-3.2.2}{3.2.2} there
is a local \(A\)-algebra \(A'\), finite and free over \(A\), such
that \(G_*\otimes_AA'\) has quasi-sections ``locally.'' For every
base extension \(T\to\operatorname{Spec}A\), the sequence
\[
  (F''\backslash G'')\times_{\operatorname{Spec}A}T
    \rightrightarrows
  (F'\backslash G')\times_{\operatorname{Spec}A}T
    \longrightarrow
  (F\backslash G)\times_{\operatorname{Spec}A}T
\]
obtained from diagram \((*)\) of \hyperref[diag:sga3-VIA-3.2.3-exact-descent]{3.2.3} is exact.

If, in addition, \(T\) is flat over \(\operatorname{Spec}A\), then
\hyperref[par:sga3-VIA-4.6]{4.6} identifies
\[
  (F''\backslash G'')\times_{\operatorname{Spec}A}T
  \quad\text{and}\quad
  (F'\backslash G')\times_{\operatorname{Spec}A}T
\]
respectively with the cokernels of the groupoids
\[
  (G_*\otimes_AA'')\times_{\operatorname{Spec}A}T
  \quad\text{and}\quad
  (G_*\otimes_AA')\times_{\operatorname{Spec}A}T.
\]
The diagram obtained from \hyperref[diag:sga3-VIA-3.2.3-exact-descent]{3.2.3}\((*)\) by the base change
\(T\to\operatorname{Spec}A\) then shows that
\[
  (F\backslash G)\times_{\operatorname{Spec}A}T
\]
is the cokernel of
\(G_*\times_{\operatorname{Spec}A}T\). The same argument applies in
the general case, when \(F\) and \(G\) are merely locally of finite
type over \(A\). Thus:

\medskip
\phantomsection\label{prop:sga3-VIA-4.7.1}%
\noindent\textbf{Proposition 4.7.1.}
\textit{Under the hypotheses of \hyperref[thm:sga3-VIA-3.2]{Theorem~3.2}, for every flat
\(A\)-scheme \(T\), the scheme
\[
  (F\backslash G)\times_{\operatorname{Spec}A}T
\]
is identified with the left quotient of
\(G\times_{\operatorname{Spec}A}T\) by
\(F\times_{\operatorname{Spec}A}T\).}%
\footnote{\label{note:sga3-VIA-48}Editor's note: we have added the
numbering \SGARefLink{prop:sga3-VIA-4.7.1}{4.7.1} in order to highlight this result.}

% END INCLUDED SGA3_English_current_buildable_source_20260823/source/inputs/VI/tex_reference_v2/components/18_expose_VIA_par47_flat_base_change_en.tex

% BEGIN INCLUDED SGA3_English_current_buildable_source_20260823/source/inputs/VI/tex_reference_v2/components/19_expose_VIA_section5_opening_par51_loop2_r1_en.tex
% SGA 3, Expose VI_A, Loop 1 English reconstruction.
% Scope: section 5 heading and paragraph 5.1 in full;
% paragraph 5.2 is excluded.
% Authority: Polo--Gille Exp6A-13oct24.pdf, component-local pages 30--31
% (combined-reader pages 333--334; printed pages 322--323).
% OCR and the Reinhold English Markdown are comparison witnesses only.
% Loop 2 replaces the exact 180-dpi source crop with a native tikz-cd
% reconstruction controlled against local p. 30,
% x=520,y=1142,w=530,h=305 on a 1488x2105 render.

\section*{5. Connections with Expos\'e IV and consequences}
\phantomsection\label{sec:sga3-VIA-5}
\addcontentsline{toc}{section}{5. Connections with Expos\'e IV and
  consequences}

\medskip
\phantomsection\label{par:sga3-VIA-5.1}%
\noindent\textbf{5.1.---}%
\footnote{\label{note:sga3-VIA-49}Editor's note: we have changed the
title of this section, which was called ``Complements'' in the
original.}
We retain the notation of \S~3 and the hypotheses of \hyperref[thm:sga3-VIA-3.2]{Theorem~3.2}.
We then have the following commutative diagram:
\begin{center}
  \phantomsection\label{diag:sga3-VIA-5.1-right-action}%
  \begin{tikzcd}[row sep=huge,column sep=huge]
    F\times_A G\times_A G
      \arrow[r,"F\times\nu"]
      \arrow[d,shift right=0.9ex,"\operatorname{pr}_2\times G"']
      \arrow[d,shift left=0.9ex,"\lambda\times G"] &
    F\times_A G
      \arrow[d,shift right=0.9ex,"\operatorname{pr}_2"']
      \arrow[d,shift left=0.9ex,"\lambda"] \\
    G\times_A G
      \arrow[r,"\nu"]
      \arrow[d,"p\times G"'] &
    G
      \arrow[d,"p"] \\
    (F\backslash G)\times_A G
      \arrow[r,dashed,"\rho"] &
    F\backslash G
  \end{tikzcd}
\end{center}
It satisfies
\[
  \operatorname{pr}_2\circ(F\times\nu)
    =
  \nu\circ(\operatorname{pr}_2\times G),
  \qquad
  \lambda\circ(F\times\nu)
    =
  \nu\circ(\lambda\times G).
\]
Moreover, since \(G\) is flat over \(A\), the left-hand vertical
sequence is exact by \hyperref[par:sga3-VIA-4.7]{4.7}. Hence \(\nu\) induces a morphism of
\(A\)-schemes
\[
  \rho:(F\backslash G)\times_A G\longrightarrow F\backslash G.
\]

The morphism \(\rho\) makes \(G\) act on the right on
\(F\backslash G\), as is immediately checked; moreover, the
canonical morphism \(G\to F\backslash G\) commutes with the right
actions of \(G\) on \(G\) and on \(F\backslash G\).%
\footnote{\label{note:sga3-VIA-50}Editor's note: we have added what
follows.}
This proves the first assertion of \hyperref[concl:sga3-VIA-3.2-ii-prime]{Theorem~3.2}(ii\(^{\prime}\)).
By \SGARefLink{prop:sga3-VIA-2.5.4}{2.5.4}, the connected components of
\[
  X=F\backslash G
\]
are therefore of finite type, irreducible, and all of the same
dimension.

To compute this dimension, we may suppose that \(A=k\) and that
\(k\) is algebraically closed. By \SGARefLink{sga3:i:scholium:2-3-3-1}{Expos\'e~I, 2.3.3.1}, the
stabilizer of the \(k\)-point \(p(e)\) is represented by the fiber
\[
  H=p^{-1}\bigl(p(e)\bigr).
\]
Since \(F\backslash G\) is the quotient of \(G\) by \(F\) in the
category of ringed spaces, the underlying space of this fiber is
\(u(F)\). As \(\ker(u)\) is finite, we obtain
\[
  \dim H=\dim u(F)=\dim F.
\]
It follows from \hyperref[concl:sga3-VIA-2.5.4-ii]{2.5.4(ii)} that
\[
  \dim X=\dim G-\dim F.
\]
This proves \hyperref[concl:sga3-VIA-3.2-ii-prime]{Theorem~3.2}(ii\(^{\prime}\)), and hence also
\hyperref[concl:sga3-VIA-3.3.2-iv]{Theorem~3.3.2(iv)}.

% END INCLUDED SGA3_English_current_buildable_source_20260823/source/inputs/VI/tex_reference_v2/components/19_expose_VIA_section5_opening_par51_loop2_r1_en.tex

% BEGIN INCLUDED SGA3_English_current_buildable_source_20260823/source/inputs/VI/tex_reference_v2/components/20_expose_VIA_par52_group_structure_en.tex
% SGA 3, Expose VI_A, Loop 1 English reconstruction.
% Scope: paragraph 5.2 in full; paragraph 5.3 is excluded.
% Authority: Polo--Gille Exp6A-13oct24.pdf, component-local page 31
% (combined-reader page 334; printed page 323).
% OCR and the Reinhold English Markdown are comparison witnesses only.

\medskip
\phantomsection\label{par:sga3-VIA-5.2}%
\noindent\textbf{5.2.---}
When the homomorphism of \(A\)-groups \(u:F\to G\) is a
monomorphism, the result of \hyperref[par:sga3-VIA-5.1]{5.1} may also be recovered from the
results of Expos\'e~IV. Indeed, the canonical morphism
\[
  p:G\longrightarrow F\backslash G
\]
is faithfully flat and open by \hyperref[thm:sga3-VIA-3.2]{3.2}; it is therefore a covering for
the \((\mathrm{fpqc})\) topology by \SGARefLink{sga3:iv:proposition:6:3:1}{Expos\'e~IV, 6.3.1}, and we may
apply \SGARefLink{sga3:iv:corollary:5:2:2}{Expos\'e~IV, Corollaries~5.2.2} and 5.2.4.

In particular, suppose, in addition to the hypotheses of \hyperref[thm:sga3-VIA-3.2]{3.2}, that
\(u\) is the inclusion in \(G\) of a normal subgroup \(F\). Then
there is a unique \(A\)-group structure on \(F\backslash G\) for
which the canonical morphism
\[
  p:G\longrightarrow F\backslash G
\]
is a homomorphism of \(A\)-groups.%
\footnote{\label{note:sga3-VIA-51}Editor's note: we have added the
following sentence.}
This proves \hyperref[concl:sga3-VIA-3.3.2-v]{Theorem~3.3.2(v)}.

% END INCLUDED SGA3_English_current_buildable_source_20260823/source/inputs/VI/tex_reference_v2/components/20_expose_VIA_par52_group_structure_en.tex

% BEGIN INCLUDED SGA3_English_current_buildable_source_20260823/source/inputs/VI/tex_reference_v2/components/21_expose_VIA_par53_correspondences_en.tex
% SGA 3, Expose VI_A, Loop 1 English reconstruction.
% Scope: paragraph 5.3 and subparagraphs 5.3.1--5.3.3 in full;
% paragraph 5.4 is excluded.
% Authority: Polo--Gille Exp6A-13oct24.pdf, component-local pages 31--32
% (combined-reader pages 334--335; printed pages 323--324).
% OCR and the Reinhold English Markdown are comparison witnesses only.

\medskip
\phantomsection\label{par:sga3-VIA-5.3}%
\noindent\textbf{5.3.---}
We now review several statements from Expos\'e~IV.

\medskip
\phantomsection\label{par:sga3-VIA-5.3.1}%
\noindent\textbf{5.3.1.---}
The statements of \SGARefLink{sga3:iv:proposition:5:2:7}{Expos\'e~IV, 5.2.7} and \SGARefLink{par:sga3-VIA-5.3.1}{5.3.1} translate as follows.
Let \(F\) and \(G\) be groups locally of finite type and flat over
\(A\), with \(F\) a closed normal subgroup of \(G\). The maps
\[
  H\longmapsto F\backslash H,
  \qquad
  H'\longmapsto H'\times_{(F\backslash G)}G
\]
define a bijection between the flat \(A\)-subgroups of \(G\)
containing \(F\) and the flat \(A\)-subgroups of \(F\backslash G\).
Under this bijection, the closed (respectively, normal) subgroups of
\(G\) containing \(F\) correspond to the closed (respectively,
normal) subgroups of \(F\backslash G\).%
\footnote{\label{note:sga3-VIA-52}Editor's note: in addition to the
cited statements of Expos\'e~IV, we use the fact that, since
\(G\to F\backslash G\) is faithfully flat, an \(A\)-subgroup \(H\)
of \(G\) is flat over \(A\) if and only if \(F\backslash H\) is.}

\medskip
\phantomsection\label{par:sga3-VIA-5.3.2}%
\noindent\textbf{5.3.2.---}
\SGARefLink{sga3:iv:proposition:5:2:9}{Expos\'e~IV, Proposition~5.2.9}, gives the following result. Let
\(F,H,G\) be groups locally of finite type and flat over \(A\), and
suppose
\[
  F\subset H\subset G,
\]
with \(F\) closed in \(G\) and normal in \(H\). Then
\(F\backslash H\) acts freely on the left on \(F\backslash G\), the
quotient scheme
\[
  (F\backslash H)\backslash(F\backslash G)
\]
exists, and there is a canonical isomorphism of \(G\)-schemes
\[
  (F\backslash H)\backslash(F\backslash G)
    \xrightarrow{\ \sim\ }
  H\backslash G.
\]

\medskip
\phantomsection\label{par:sga3-VIA-5.3.3}%
\noindent\textbf{5.3.3.---}
Finally, \SGARefLink{sga3:iv:proposition:5:2:8}{Expos\'e~IV, 5.2.8}, gives the following assertion. Let
\(F,H,G\) be groups locally of finite type and flat over \(A\).
Suppose that \(F\) is contained in \(G\), closed and normal in \(G\),
that \(H\) is contained in \(G\), and that \(F\cap H\) is flat over
\(A\). Let
\[
  F\mathbin{\overset{\tau}{\times}}_A H
\]
be the \(A\)-group whose underlying scheme is \(F\times_AH\), with
multiplication defined by
\[
  \bigl((x,h),(y,h')\bigr)
    \longmapsto
  \bigl(xhyh^{-1},hh'\bigr).
\]
Likewise, let
\[
  u:H\cap F\longrightarrow
    F\mathbin{\overset{\tau}{\times}}_A H
\]
be the monomorphism \(x\mapsto(x^{-1},x)\), and let \(F\cdot H\) be
the quotient
\[
  (F\cap H)\backslash
    \bigl(F\mathbin{\overset{\tau}{\times}}_A H\bigr).
\]
Then there is a canonical isomorphism
\[
  F\backslash(F\cdot H)
    \xrightarrow{\ \sim\ }
  (F\cap H)\backslash H.
\]

% END INCLUDED SGA3_English_current_buildable_source_20260823/source/inputs/VI/tex_reference_v2/components/21_expose_VIA_par53_correspondences_en.tex

% BEGIN INCLUDED SGA3_English_current_buildable_source_20260823/source/inputs/VI/tex_reference_v2/components/22_expose_VIA_par54_abelian_categories_loop2_r1_en.tex
% SGA 3, Expose VI_A, Loop 1 English reconstruction.
% Scope: paragraph 5.4, Proposition 5.4.1, Theorem 5.4.2, and
% Corollary 5.4.3 in full; paragraph 5.5 is excluded.
% Authority: Polo--Gille Exp6A-13oct24.pdf, component-local pages 32--33
% (combined-reader pages 335--336; printed pages 324--325).
% OCR and the Reinhold English Markdown are comparison witnesses only.
% Loop 2 replaces the exact 180-dpi source crop with a native tikz-cd
% factorization controlled against the PDF, local page 32.

\medskip
\phantomsection\label{par:sga3-VIA-5.4}%
\noindent\textbf{5.4.---}%
\footnote{\label{note:sga3-VIA-53a}Editor's note: we have expanded
the original in what follows; in particular, we have added
\hyperref[prop:sga3-VIA-5.4.1]{Proposition~5.4.1}.}
Let \(u:G\to H\) be a quasi-compact morphism between \(A\)-groups
locally of finite type, such that the kernel \(N\) of \(u\) is flat
over \(A\). In this case, by \SGARefLink{thm:sga3-VIA-3.3.2}{3.3.2} and \SGARefLink{prop:sga3-VIB-5.2}{5.2}, the quotient \(A\)-group
\[
  C=N\backslash G
\]
exists, and the morphism \(p:G\to C\) is faithfully flat and locally
of finite presentation. On the other hand, by \SGARefLink{sga3:iv:remark:5:2:6}{Expos\'e~IV, 5.2.6},
\(u\) induces a monomorphism
\[
  v:C\longrightarrow H,
\]
which is quasi-compact (because \(u\) is quasi-compact and
\(G\to C\) is surjective; cf.\ EGA~IV\(_1\), 1.1.3), and is
therefore a closed immersion by \hyperref[prop:sga3-VIA-2.5.2]{2.5.2}. We have thus obtained the
following proposition.

\medskip
\phantomsection\label{prop:sga3-VIA-5.4.1}%
\noindent\textbf{Proposition 5.4.1.---}
\emph{Let \(u:G\to H\) be a quasi-compact morphism between
\(A\)-groups locally of finite type, such that
\(N=\ker u\) is flat over \(A\). Then \(u\) admits the factorization}
\begin{center}
  \phantomsection\label{diag:sga3-VIA-5.4.1-factorization}%
  \begin{tikzcd}[row sep=large,column sep=large]
    G \arrow[r,"u"] \arrow[d,"p"'] & H \\
    N\backslash G \arrow[ur,"i"']
  \end{tikzcd}
\end{center}
\emph{where \(p\) is faithfully flat and locally of finite
presentation, and \(i\) is a closed immersion.}
\par\medskip

Suppose in addition that \(G\) is flat over \(A\). Then, by \SGARefLink{thm:sga3-VIA-3.3.2}{3.3.2},
\(C=N\backslash G\) is flat over \(A\), and hence the quotient
\[
  X=C\backslash H
\]
exists in \((\mathbf{Sch}_{/A})\) and represents the quotient
\((\mathrm{fppf})\)-sheaf \(\widetilde C\backslash H\), while
\[
  q:H\longrightarrow X
\]
is a \(C\)-torsor. Consequently, if
\[
  e:\operatorname{Spec}A\longrightarrow H
\]
denotes the unit section of \(H\), then \(v\) induces an isomorphism
of \((\mathrm{fppf})\)-sheaves between \(\widetilde C\) and the fiber
product of \(q\) with
\[
  q\circ e:\operatorname{Spec}A\longrightarrow X,
\]
which is represented by a closed subscheme of \(H\). Consequently,
\(v\) is an isomorphism of \(C\) onto a closed group subscheme
\(K\) of \(H\), equal to the stabilizer of the \(A\)-point
\(q\circ e\) of \(X\).%
\begingroup
\renewcommand{\thefootnote}{\fnsymbol{footnote}}%
\footnote[1]{\label{note:sga3-VIA-5.4-K-ambient}Translator's note:
the French source declares \(e:\operatorname{Spec}A\to G\), although
the displayed composition \(q\circ e\) requires the unit section of
\(H\), and then prints ``a closed group subscheme \(K\) of \(G\).''
Here \(v:C\to H\), the displayed fiber product is represented inside
\(H\), and the next paragraph identifies \(K\) with
\(\ker(H\to\overline H)\); accordingly, \(e\) is taken to be the unit
section of \(H\), and the intended ambient group of \(K\) is \(H\).}%
\endgroup{}
(This also gives another proof that every quasi-compact monomorphism
\(v:C\to H\) between \(A\)-groups locally of finite type is a closed
immersion; cf.\ \SGARefLink{prop:sga3-VIA-2.5.2}{2.5.2} and Expos\'e~VI\(_B\), \SGARefLink{cor:sga3-VIB-1.4.2}{1.4.2}.)

Suppose further that \(C\) is a normal subgroup of \(H\). In this
case, the \(A\)-group
\[
  \overline H=C\backslash H
\]
is the cokernel, in the category of \(A\)-groups, of the morphism
\(u:G\to H\), and \(K\) is the kernel of the morphism
\[
  H\longrightarrow\overline H.
\]
When \(G\) and \(H\) are abelian \(A\)-groups, \(K\) is the image of
\(u\) in the category of abelian \(A\)-groups, whereas
\[
  C=(\ker u)\backslash G
\]
is the coimage of \(u\). Taking into account the isomorphism
\[
  C\xrightarrow{\ \sim\ }K
\]
just established, we obtain the following.%
\footnote[53]{\label{note:sga3-VIA-53b}Editor's note: we have added
the number \SGARefLink{thm:sga3-VIA-5.4.2}{5.4.2} to this theorem.}

\medskip
\phantomsection\label{thm:sga3-VIA-5.4.2}%
\noindent\textbf{Theorem 5.4.2.---}
\emph{Let \(k\) be a field. The category of commutative algebraic
\(k\)-groups is abelian.}

Indeed, when \(k\) is a field, \(\ker(u)\) is flat over \(k\),
whatever \(u\) may be.

\noindent
\footnote{\label{note:sga3-VIA-54}Editor's note: we have added what
follows.}
Note that the full subcategory of commutative affine algebraic
\(k\)-groups is thick. Indeed, consider an exact sequence of
commutative algebraic \(k\)-groups
\[
  1\longrightarrow N\longrightarrow G\longrightarrow
  G/N\longrightarrow 1.
\]
If \(G\) is affine, then \(N\) is clearly affine, and \(G/N\) is
also affine by a theorem of Chevalley; cf.\ Expos\'e~VI\(_B\),
\SGARefLink{thm:sga3-VIB-11.17}{11.17}. Conversely, if \(N\) and \(G/N\) are affine, then \(G\) is
also affine by Expos\'e~VI\(_B\), \SGARefLink{concl:sga3-VIB-9.2-viii}{9.2(viii)}. We therefore obtain the
following corollary.

\medskip
\phantomsection\label{cor:sga3-VIA-5.4.3}%
\noindent\textbf{Corollary 5.4.3.---}
\emph{Let \(k\) be a field. The category of commutative affine
algebraic \(k\)-groups is abelian.}

Let us also point out that the category of all commutative affine
\(k\)-groups (not necessarily of finite type) is abelian. This
follows from Expos\'e~VI\(_B\), \SGARefLink{thm:sga3-VIB-11.17}{11.17} and \SGARefLink{thm:sga3-VIB-11.18.2}{11.18.2}
(cf.\ \textup{[DG70]}, \S~III.3, \SGARefLink{prop:sga3-VIB-7.4}{7.4}); see also
Expos\'e~VII\(_B\), 2.4.2 for a proof using formal groups.

% END INCLUDED SGA3_English_current_buildable_source_20260823/source/inputs/VI/tex_reference_v2/components/22_expose_VIA_par54_abelian_categories_loop2_r1_en.tex

% BEGIN INCLUDED SGA3_English_current_buildable_source_20260823/source/inputs/VI/tex_reference_v2/components/23_expose_VIA_par55_56_loop2_r1_en.tex
% SGA 3, Expose VI_A, Loop 1 English reconstruction.
% Scope: paragraphs 5.5 and 5.6, through Proposition 5.6.2 and its
% proof in full; Section 6 is excluded.
% Authority: Polo--Gille Exp6A-13oct24.pdf, component-local pages 33--34
% (combined-reader pages 336--337; printed pages 325--326).
% OCR and the Reinhold English Markdown are comparison witnesses only.
% Loop 2 replaces the two exact 180-dpi source crops with native tikz-cd
% reconstructions controlled against the PDF, local page 34:
%   quotient square: x=605,y=710,w=300,h=155 on a 1488x2105 render;
%   local-flatness square: x=575,y=1095,w=350,h=150 on that render.

\medskip
\phantomsection\label{par:sga3-VIA-5.5}%
\noindent\textbf{5.5.---}
Let \(G\) be a group locally of finite type and flat over an
Artinian local ring \(A\). We know from \SGARefLink{cor:sga3-VIB-2.3}{2.3} that the connected
component \(G^0\) of the identity is a normal open group subscheme
of \(G\), hence is also flat over \(A\). It then follows from \hyperref[thm:sga3-VIA-3.2]{3.2}
and \hyperref[par:sga3-VIA-5.2]{5.2} that
\[
  G^0\backslash G
\]
is an \(A\)-group scheme, flat over \(A\). Moreover, since every
connected component \(G^\alpha\) of \(G\) is saturated for the
equivalence relation defined by \(G^0\), the quotient
\(G^0\backslash G\) is the disjoint union of the
\(G^0\backslash G^\alpha\) (cf.\ \hyperref[par:sga3-VIA-4.3]{4.3}). In particular, the connected
component of the identity in \(G^0\backslash G\) is precisely
\[
  G^0\backslash G^0\simeq\operatorname{Spec}A,
\]
and hence the morphism
\[
  G^0\backslash G\longrightarrow\operatorname{Spec}A
\]
is a local isomorphism at the identity. Consequently,
\(G^0\backslash G\) is \'etale over \(\operatorname{Spec}A\), by
Expos\'e~VI\(_B\), \SGARefLink{prop:sga3-VIB-1.3}{1.3}.%
\footnote{\label{note:sga3-VIA-55}Editor's note: we have expanded
the preceding discussion from the original and added
\hyperref[prop:sga3-VIA-5.5.1]{Proposition~5.5.1} in order to make this result explicit.}
We therefore obtain the following proposition (for assertion~\SGARefLink{prop:sga3-VIA-5.5.1:item:ii}{(ii)},
see \textup{[DG70]}, \S~II.5, 1.7--1.10).

\medskip
\phantomsection\label{prop:sga3-VIA-5.5.1}%
\noindent\textbf{Proposition 5.5.1.---}
\emph{Let \(A\) be an Artinian local ring and \(G\) an \(A\)-group
locally of finite type and flat.
\begin{enumerate}
\item[\textup{(i)}] \(G^0\backslash G\) is an \'etale \(A\)-group.
\item[\textup{\SGAInlineRefTarget{prop:sga3-VIA-5.5.1:item:ii}(ii)}] Consequently, if \(A=k\) is an algebraically
closed field, \(G^0\backslash G\) is a constant \(k\)-group acting
simply transitively on the set of connected components of \(G\).
Thus, if \(G\) is algebraic, \(G^0\backslash G\) is finite.
\end{enumerate}}

\medskip
\phantomsection\label{par:sga3-VIA-5.6}%
\noindent\textbf{5.6.---}
Now let \(k\) be a perfect field and \(G\) a \(k\)-group locally of
finite type. We saw in \SGARefLink{sga3:VIA:0.2}{0.2} that \(G_{\mathrm{red}}\) is then a group
subscheme of \(G\). Moreover, the equivalence class of the identity
of \(G\) under the left action of \(G_{\mathrm{red}}\) on \(G\) is
the whole underlying space of \(G\). \hyperref[thm:sga3-VIA-3.2]{Theorem~3.2} therefore gives
the following result.

\medskip
\phantomsection\label{prop:sga3-VIA-5.6.1}%
\noindent\textbf{Proposition 5.6.1.---}
\emph{Let \(k\) be a perfect field and \(G\) a \(k\)-group locally
of finite type. Then the \(k\)-scheme
\[
  G_{\mathrm{red}}\backslash G
\]
is the spectrum of a finite local \(k\)-algebra with residue field
\(k\).}

\noindent
\footnote{\label{note:sga3-VIA-56}Editor's note: we have added the
number 5.6.1, as well as the proof that follows.}
Indeed, by \SGARefLink{thm:sga3-VIA-3.2}{3.2}, the quotient \(G_{\mathrm{red}}\backslash G\) has a
single point, with residue field \(k\), and is a \(k\)-scheme of
finite type. It is therefore the spectrum of a finite-dimensional
local \(k\)-algebra (cf.\ EGA~I, 6.4.4).

\medskip
\phantomsection\label{prop:sga3-VIA-5.6.2}%
\noindent\textbf{Proposition 5.6.2.---}
\emph{Let \(u:F\to G\) be a morphism between groups locally of
finite type over a perfect field \(k\). The following assertions
are equivalent:
\begin{enumerate}
\item[\textup{\SGAInlineRefTarget{prop:sga3-VIA-5.6.2:item:i}(i)}] \(u\) is flat.
\item[\textup{\SGAInlineRefTarget{prop:sga3-VIA-5.6.2:item:ii}(ii)}] The morphism
\[
  u^0:F^0\longrightarrow G^0
\]
is dominant, and the morphism
\[
  v:F_{\mathrm{red}}\backslash F
    \longrightarrow G_{\mathrm{red}}\backslash G
\]
induced by \(u\) is flat.
\end{enumerate}}

\noindent
\footnote{\label{note:sga3-VIA-57}Editor's note: we have expanded
the proof of \textup{\SGARefLink{prop:sga3-VIA-5.6.2:item:ii}{(ii)}} \(\Rightarrow\) \textup{\SGARefLink{prop:sga3-VIA-5.6.2:item:i}{(i)}} and
simplified the diagram below.}
Indeed, consider the following commutative diagram:
\begin{center}
  \phantomsection\label{diag:sga3-VIA-5.6.2-quotients}%
  \begin{tikzcd}[row sep=large,column sep=large]
    F
      \arrow[r,"p"]
      \arrow[d,"u"'] &
    F_{\mathrm{red}}\backslash F
      \arrow[d,"v"] \\
    G
      \arrow[r,"q"'] &
    G_{\mathrm{red}}\backslash G
  \end{tikzcd}
\end{center}
where \(p\) and \(q\) denote the canonical projections. By
\SGARefLink{concl:sga3-VIA-3.2-iv}{3.2(iv)}, \(p\) and \(q\) are faithfully flat. Consequently, if
\(u\) is flat, then
\[
  q\circ u=v\circ p
\]
is flat, and hence \(v\) is flat as well.

Conversely, suppose that \(v\) is flat and \(u^0\) is dominant.
Since \(u^0\) is quasi-compact (the scheme \(F^0\) is of finite type
over \(k\) by \SGARefLink{prop:sga3-VIA-2.4}{2.4}, and hence Noetherian), it maps the generic point
\(\xi\) of \(F^0\) to the generic point \(\eta\) of \(G^0\). Let
\(R\) be the finite local \(k\)-algebra such that
\[
  G_{\mathrm{red}}\backslash G=\operatorname{Spec}R,
\]
and let \(\mathfrak m\) be its maximal ideal. We have local
homomorphisms of local rings
\[
  R\longrightarrow\mathcal O_{G,\eta}
   \longrightarrow\mathcal O_{F,\xi}.
\]
Observe that the following square is Cartesian:
\begin{center}
  \phantomsection\label{diag:sga3-VIA-5.6.2-local-flatness}%
  \begin{tikzcd}[row sep=large,column sep=large]
    G_{\mathrm{red}}
      \arrow[r]
      \arrow[d] &
    G
      \arrow[d,"q"] \\
    \operatorname{Spec}(R/\mathfrak m)
      \arrow[r] &
    \operatorname{Spec}(R)
  \end{tikzcd}
\end{center}
It follows that
\[
  \mathcal O_{G,\eta}/\mathfrak m\mathcal O_{G,\eta}
    \simeq \mathcal O_{G_{\mathrm{red}},\eta}
    =\kappa(\eta),
\]
so that
\[
  \mathcal O_{F,\xi}/\mathfrak m\mathcal O_{F,\xi}
\]
is flat over
\(\mathcal O_{G,\eta}/\mathfrak m\mathcal O_{G,\eta}\).

On the other hand, since \(q\) and \(v\circ p\) are flat, \(G\)
and \(F\) are flat over \(R\). The local criterion for flatness
(cf.\ EGA~IV\(_3\), 11.3.10.2) therefore shows that
\(\mathcal O_{F,\xi}\) is flat over \(\mathcal O_{G,\eta}\); in
other words, \(u\) is flat at \(\xi\). It follows from \hyperref[lem:sga3-VIA-2.5.3]{2.5.3} that
\(u\) is flat.

% END INCLUDED SGA3_English_current_buildable_source_20260823/source/inputs/VI/tex_reference_v2/components/23_expose_VIA_par55_56_loop2_r1_en.tex

% BEGIN INCLUDED SGA3_English_current_buildable_source_20260823/source/inputs/VI/tex_reference_v2/components/24_expose_VIA_section6_through_prop64_loop2_r1_en.tex
% SGA 3, Expose VI_A, Loop 1 English reconstruction.
% Scope: Section 6 opening through the proof of Proposition 6.4;
% Theorem 6.5 is excluded.
% Authority: Polo--Gille Exp6A-13oct24.pdf, component-local pages 34--36
% (combined-reader pages 337--339; printed pages 326--328).
% OCR and the Reinhold English Markdown are comparison witnesses only.
% Loop 2 replaces the exact 180-dpi Proposition 6.4 source crop with a
% native tikz-cd reconstruction controlled against PDF local page 35.

\section*{6. Complements on \(k\)-groups not necessarily of finite type}
\phantomsection\label{sec:sga3-VIA-6}
\addcontentsline{toc}{section}{6. Complements on \(k\)-groups not
  necessarily of finite type}

\footnote{\label{note:sga3-VIA-58}Editor's note: we have added the
results that follow, taken from \textup{[Per75]}. Note that
\hyperref[lem:sga3-VIA-6.1]{Lemma~6.1} can be expressed, in Weil's language, by saying that
``every point of \(G\) is a product of two generic points.''}
Let us also record the following results, which will be useful in
the addendum, Expos\'e~VI\(_B\), \S~12. We fix a base field \(k\).

\medskip
\phantomsection\label{lem:sga3-VIA-6.1}%
\noindent\textbf{Lemma 6.1.---}
\emph{Let \(G\) be a \(k\)-group. For every \(x\in G\), there
exists a point \(u\in G\times G\) such that \(\mu(u)=x\) and both
projections \(p_1(u)\) and \(p_2(u)\) are maximal points of \(G\).}

\emph{Proof.}
Set \(K=\kappa(x)\). Since the projection \(G_K\to G\) sends
maximal points to maximal points, we are reduced to the case in
which \(x\) is rational. Left translation \(\lambda_x\)
(respectively, right translation \(\rho_x\)) then gives a morphism
\[
  G\longrightarrow G\times G,\qquad
  g\longmapsto\bigl(\lambda_x(g^{-1}),g\bigr)
\]
(respectively,
\(g\longmapsto\bigl(g,\rho_x(g^{-1})\bigr)\)), which induces an
isomorphism from \(G\) onto \(\mu^{-1}(x)\), inverse to \(p_2\)
(respectively, \(p_1\)). Thus, if \(u\) is a maximal point of
\(\mu^{-1}(x)\), then \(p_1(u)\) and \(p_2(u)\) are maximal points
of \(G\), and \(u\) has the required properties.

\medskip
\phantomsection\label{cor:sga3-VIA-6.2}%
\noindent\textbf{Corollary 6.2.---}
\emph{Let \(f:G\to H\) be a quasi-compact dominant morphism of
\(k\)-groups.
\begin{enumerate}
\item[\textup{(i)}]\phantomsection\label{concl:sga3-VIA-6.2-i} \(f\) is surjective.
\item[\textup{(ii)}] If \(H\) is reduced, \(f\) is faithfully flat.
\end{enumerate}}

\emph{Proof.}
Write \(\mu_H\) (respectively, \(\mu_G\)) for the multiplication of
\(H\) (respectively, \(G\)). Let \(h\in H\). By \hyperref[lem:sga3-VIA-6.1]{6.1}, there exists
\(u\in H\times H\) such that \(\mu_H(u)=h\) and
\(\alpha=p_1(u)\), \(\beta=p_2(u)\) are maximal points of \(H\).
Since \(f\) is quasi-compact and dominant,
\(f^{-1}(\alpha)\) and \(f^{-1}(\beta)\) are nonempty
(cf.\ EGA~IV\(_1\), 1.1.5). Hence there exists \(v\in G\times G\)
such that
\[
  (f\times f)(v)=u
\]
(cf.\ EGA~I, 3.5.2). Then \(g=\mu_G(v)\) satisfies \(f(g)=h\).
This proves that \(f\) is surjective.

Suppose in addition that \(H\) is reduced. Then
\(\mathcal O_{H,\alpha}\) is a field, and we have seen above that
\(f^{-1}(\alpha)\ne\varnothing\). Thus \(f\) is flat at every
point \(\xi\in f^{-1}(\alpha)\), and consequently \(f\) is flat by
\hyperref[lem:sga3-VIA-2.5.3]{Lemma~2.5.3}.

\medskip
\phantomsection\label{rem:sga3-VIA-6.3}%
\noindent\textbf{Reminder 6.3.---}
Recall (cf.\ EGA~IV\(_3\), 11.10.1) that a morphism of schemes
\(f:X\to Y\) is said to be \emph{schematically dominant} if it
satisfies the following condition: for every open subset \(U\) of
\(Y\), if \(Z\) is a closed subscheme of \(U\) such that the
morphism \(f^{-1}(U)\to U\) factors through \(Z\), then \(Z=U\).
When \(f\) is quasi-compact and quasi-separated, this is equivalent
to saying that the closed image of \(X\) under \(f\) is \(Y\)
(cf.\ loc.\ cit., 11.10.3(iv), and EGA~I, 9.5.8).

\medskip
\phantomsection\label{prop:sga3-VIA-6.4}%
\noindent\textbf{Proposition 6.4.---}
\emph{Let \(f:H\to G\) be a quasi-compact morphism of \(k\)-groups.
Then the closed image of \(f\) is a group subscheme \(H'\) of
\(G\), and \(f\) factors as}
\begin{center}
  \phantomsection\label{diag:sga3-VIA-6.4-factorization}%
  \begin{tikzcd}[row sep=large,column sep=large]
    H \arrow[r,"f"] \arrow[d,"f'"'] & G \\
    H' \arrow[ur,"i"']
  \end{tikzcd}
\end{center}
\emph{where \(f'\) is schematically dominant, quasi-compact, and
surjective.}

\emph{Proof.}
Since \(H\) is separated (\hyperref[sga3:VIA:0.3]{0.3}), \(f\) is quasi-compact and
separated. Thus \(f_*(\mathcal O_H)\) is a quasi-coherent
\(\mathcal O_G\)-module, and the closed image \(H'\) of \(f\)
exists and is the closed subscheme of \(G\) defined by the
quasi-coherent ideal
\[
  \mathcal I
    =\ker\bigl(\mathcal O_G\longrightarrow f_*(\mathcal O_H)\bigr)
\]
(cf.\ EGA~I, \S~9.5).

Write \(c_G,\mu_G\) (respectively, \(c_H,\mu_H\)) for the inversion
and multiplication morphisms of \(G\) (respectively, \(H\)).
Then \(c_G\circ f=f\circ c_H\) factors through the closed subscheme
\(c_G(H')\), whence
\[
  H'\subset c_G(H')
  \qquad\text{and therefore}\qquad
  H'=c_G(H'),
\]
since \(c_G^2=\operatorname{id}_G\). Similarly,
\[
  f\circ\mu_H=\mu_G\circ(f\times f)
\]
factors through \(H'\), so \(f\times f\) factors through the closed
subscheme \(\mu_G^{-1}(H')\) of \(G\times G\).%
\begingroup
\renewcommand{\thefootnote}{\fnsymbol{footnote}}%
\footnote[1]{\label{note:sga3-VIA-6.4-multiplication-symbol}%
Translator's note: in these two occurrences the French source prints
\(\pi_H\) and \(\pi_G\), although the immediately preceding sentence
defines the multiplication morphisms as \(\mu_H\) and \(\mu_G\).
The multiplication symbols required by the identity are used here.}%
\endgroup{}
On the other hand, since formation of the closed image commutes
with flat base change (EGA~III, 1.4.15, and EGA~IV\(_1\), 1.7.21),
the closed image of \(f\times\operatorname{id}_H\)
(respectively, \(\operatorname{id}_{H'}\times f\)) is
\(H'\times H\) (respectively, \(H'\times H'\)). Hence, by
``transitivity of closed images'' (EGA~I, 9.5.5), the closed image
of \(f\times f\) is \(H'\times H'\). It is therefore contained in
\(\mu_G^{-1}(H')\); that is, the restriction of \(\mu_G\) to
\(H'\times H'\) factors through \(H'\). This proves that \(H'\) is
a closed group subscheme of \(G\). Write \(i:H'\hookrightarrow G\)
for the inclusion.

Now \(f=i\circ f'\), where \(f':H\to H'\) is schematically
dominant and quasi-compact (because \(f\) is quasi-compact and
\(i\) is separated). By \hyperref[cor:sga3-VIA-6.2]{6.2}, \(f'\) is therefore surjective.
This proves \hyperref[prop:sga3-VIA-6.4]{6.4}.

% END INCLUDED SGA3_English_current_buildable_source_20260823/source/inputs/VI/tex_reference_v2/components/24_expose_VIA_section6_through_prop64_loop2_r1_en.tex

% BEGIN INCLUDED SGA3_English_current_buildable_source_20260823/source/inputs/VI/tex_reference_v2/components/25_expose_VIA_theorem65_through_bibliography_loop2_r1_en.tex
% SGA 3, Expose VI_A, Loop 1 English reconstruction.
% Scope: Theorem 6.5 through Corollary 6.10 and the bibliography;
% this closes Expose VI_A. Exp6A local page 38 is substantive.
% Authority: Polo--Gille Exp6A-13oct24.pdf, component-local pages 36--38
% (combined-reader pages 339--341; printed pages 328--330).
% OCR and the Reinhold English Markdown are comparison witnesses only.
% Loop 2 replaces the two exact 180-dpi source crops with native tikz-cd
% reconstructions controlled against local page 37.

We can now state the following theorem (\textup{[Per75]}, V, \SGARefLink{prop:sga3-V-3.1}{3.1}
and \SGARefLink{thm:sga3-VIA-3.2}{3.2}; see also \textup{[Per76]}, 0.0 and \SGARefLink{sga3:VIA:0.1}{0.1}).

\medskip
\phantomsection\label{thm:sga3-VIA-6.5}%
\noindent\textbf{Theorem 6.5 (D. Perrin).---}
\emph{Let \(G\) be a quasi-compact \(k\)-group.
\begin{enumerate}
\item[\textup{(i)}]\phantomsection\label{concl:sga3-VIA-6.5-i} \(G\) is the projective limit of a filtered
  system \((G_i)\) of \(k\)-groups of finite type (whose transition
  morphisms \(u_{ij}:G_j\to G_i\) are affine for \(i\) sufficiently
  large), and the morphisms \(G\to G_i\) are faithfully flat (and
  affine for \(i\) sufficiently large).
\item[\textup{(ii)}]\phantomsection\label{concl:sga3-VIA-6.5-ii} Let \(H\) be a closed \(k\)-subgroup of \(G\).
  Then the \(\mathrm{fpqc}\) quotient sheaf
  \(\widetilde{G/H}\) is a \(k\)-scheme in either of the following
  two cases:
  \begin{enumerate}
  \item[\textup{(1)}] the immersion \(H\to G\) is of finite
    presentation; in this case \(G/H\) is of finite type over \(k\);
  \item[\textup{(2)}] \(H\) is normal in \(G\).
  \end{enumerate}
\end{enumerate}}

For the proof of this theorem, which rests on several intermediate
theorems, we refer to \textup{[Per75]}. For the reader's
convenience, however, we prove the two corollaries below; cf.\
\textup{[Per75]}, V, \SGARefLink{par:sga3-VIA-3.3}{3.3}--3.4, or \textup{[Per76]}, 4.2.3--4.2.5.

\medskip
\phantomsection\label{cor:sga3-VIA-6.6}%
\noindent\textbf{Corollary 6.6.---}
\emph{Let \(f:G\to H\) be a quasi-compact morphism of \(k\)-groups.
\begin{enumerate}
\item[\textup{(i)}] If \(f\) is schematically dominant, then it is
  faithfully flat.
\item[\textup{(ii)}] This is so, in particular, if \(H\) is affine
  and the morphism
  \[
    f^\natural:\mathcal O(H)\longrightarrow\mathcal O(G)
  \]
  is injective.
\end{enumerate}}

\emph{Proof.}
Suppose that \(f\) is schematically dominant. By \hyperref[concl:sga3-VIA-6.2-i]{6.2(i)}, \(f\) is
surjective, so by \hyperref[lem:sga3-VIA-2.5.3]{2.5.3} it is enough to prove that \(f\) is flat at
the identity element of \(G\). Since
\[
  \mathcal O_{H,e}=\mathcal O_{H^0,e}
\]
(cf.\ \hyperref[thm:sga3-VIA-2.6.5]{2.6.5}), we may replace \(H\) by \(H^0\), and hence suppose
that \(H\) is irreducible. Then \(H\) is quasi-compact
(loc.\ cit.); since \(f\) is quasi-compact, \(G\) is also
quasi-compact. By \hyperref[concl:sga3-VIA-6.5-i]{6.5(i)},
\[
  H=\varprojlim_i H_i,
\]
where each \(H_i\) is an algebraic \(k\)-group. Let \(N_i\) be the
kernel of \(G\to H_i\); it is a closed normal \(k\)-subgroup of
\(G\). Moreover, since the identity section
\(\{e\}\to H_i\) is of finite presentation, the immersion
\(N_i\to G\) is of finite presentation. Hence, by \hyperref[concl:sga3-VIA-6.5-ii]{6.5(ii)}, the
\(\mathrm{fpqc}\) quotient \(G/N_i\) is an algebraic \(k\)-group.
We then have a commutative diagram
\begin{center}
  \phantomsection\label{diag:sga3-VIA-6.6-square}%
  \begin{tikzcd}[row sep=large,column sep=large]
    G
      \arrow[r,"f"]
      \arrow[d,"q_i"'] &
    H
      \arrow[d,"p_i"] \\
    G/N_i
      \arrow[r,"f_i"'] &
    H_i
  \end{tikzcd}
\end{center}
in which \(f\) is schematically dominant and \(p_i,q_i\) are
faithfully flat, hence schematically dominant. Thus
\[
  f_i\circ q_i=p_i\circ f
\]
is schematically dominant, and therefore so is \(f_i\). On the
other hand, \(f_i\) is a monomorphism, hence a closed immersion
because \(G/N_i\) and \(H_i\) are algebraic (\hyperref[prop:sga3-VIA-2.5.2]{2.5.2}). It follows that
\(f_i\) is an isomorphism, and consequently \(G\to H_i\) is
faithfully flat. By \textup{[BAC]}, I, \S~\SGARefLink{prop:sga3-VIB-2.7}{2.7}, Proposition~9, the
morphism \(G\to H\) is then flat; it is also surjective by \hyperref[concl:sga3-VIA-6.2-i]{6.2(i)},
and is therefore faithfully flat. This proves \SGARefLink{prop:sga3-VIB-2.7:item:i}{(i)}. Assertion \SGARefLink{prop:sga3-VIB-2.7:item:ii}{(ii)}
follows, because if \(H\) is affine and \(f^\natural\) is
injective, then the closed image of \(f\) is \(H\), so \(f\) is
schematically dominant (cf.\ \hyperref[prop:sga3-VIA-6.4]{6.4}).

\medskip
\phantomsection\label{cor:sga3-VIA-6.7}%
\noindent\textbf{Corollary 6.7.---}
\emph{Let \(u:G\to H\) be a morphism of \(k\)-groups and set
\(N=\ker(u)\). Suppose that \(u\) is quasi-compact.
\begin{enumerate}
\item[\textup{\SGAInlineRefTarget{cor:sga3-VIA-6.7:item:i}(i)}] The \(\mathrm{fpqc}\) quotient sheaf
  \(\widetilde{G/N}\) is represented by a \(k\)-group scheme
  \(G/N\), and \(u\) factors as
  \begin{center}
    \phantomsection\label{diag:sga3-VIA-6.7-factorization}%
    \begin{tikzcd}[ampersand replacement=\&,row sep=large,column sep=large]
      G \arrow[r,"u"] \arrow[d,"p"'] \& H \\
      G/N \arrow[ur,"i"']
    \end{tikzcd}
  \end{center}
  where \(p\) is faithfully flat and \(i\) is a closed immersion.
\item[\textup{\SGAInlineRefTarget{cor:sga3-VIA-6.7:item:ii}(ii)}] In particular, if \(u\) is a monomorphism,
  then \(u\) is a closed immersion; and if \(u\) is schematically
  dominant, then \(u\) is faithfully flat.
\end{enumerate}}

\emph{Proof.}
\textup{\SGARefLink{cor:sga3-VIA-6.7:item:i}{(i)}} By \hyperref[prop:sga3-VIA-6.4]{6.4} and \hyperref[cor:sga3-VIA-6.6]{6.6}, the closed image \(G'\) of \(u\) is a
closed group subscheme of \(H\), and
\[
  p:G\longrightarrow G'
\]
is faithfully flat and quasi-compact.%
\begingroup
\renewcommand{\thefootnote}{\fnsymbol{footnote}}%
\footnote[1]{\label{note:sga3-VIA-6.7-image-ambient}Translator's
note: the French source prints ``a closed group subscheme of
\(G\)'' at this point. Since \(u:G\to H\), its closed image lies in
\(H\), as the displayed factorization also shows.}%
\endgroup{}
Clearly \(\ker(p)=\ker(u)=N\). Thus, by \SGARefLink{sga3:iv:scholium:3:3:2:1}{Expos\'e~IV, 3.3.2.1} and
5.1.7, \(G'\) represents the \(\mathrm{fpqc}\) quotient sheaf
\(\widetilde{G/N}\).

\textup{\SGARefLink{cor:sga3-VIA-6.7:item:ii}{(ii)}} The second assertion is contained in \hyperref[cor:sga3-VIA-6.6]{6.6}; let us
prove the first. If \(u\) is a monomorphism, then so is \(p\).
Thus \(p\) is both a monomorphism and an effective epimorphism, and
is therefore an isomorphism (cf.\ \SGARefLink{sga3:iv:remark:1:14}{Expos\'e~IV, 1.14}). This proves
\hyperref[cor:sga3-VIA-6.7]{6.7}.

Let us finally record the following corollaries
(cf.\ \textup{[Per76]}, 4.2.6--4.2.8).

\medskip
\phantomsection\label{cor:sga3-VIA-6.8}%
\noindent\textbf{Corollary 6.8.---}
\emph{The category of commutative quasi-compact \(k\)-group schemes
is abelian.}

Taking \hyperref[cor:sga3-VIA-6.7]{6.7} into account, the proof is analogous to that of \hyperref[thm:sga3-VIA-5.4.2]{5.4.2}.

\medskip
\phantomsection\label{cor:sga3-VIA-6.9}%
\noindent\textbf{Corollary 6.9.---}
\emph{If \(\operatorname{char}(k)=0\), every \(k\)-group scheme
\(G\) is geometrically reduced.}

Indeed, if \(g\in G\) and \(K\) is an algebraic closure of
\(\kappa(g)\), then
\[
  \mathcal O_{G,e}\otimes_k K
    \simeq
  \mathcal O_{G_K,g_K}.
\]
It is therefore enough to prove that
\(\mathcal O_{G,e}=\mathcal O_{G^0,e}\) is geometrically reduced.
We are thus reduced to the case in which \(G\) is connected, hence
quasi-compact by \hyperref[thm:sga3-VIA-2.6.5]{2.6.5}. The result then follows from \hyperref[thm:sga3-VIA-6.5]{6.5} and
Cartier's theorem for algebraic groups
(cf.\ Expos\'e~VI\(_B\), \SGARefLink{cor:sga3-VIB-1.6.1}{1.6.1}, or \textup{[DG70]}, \S~II.6,
\SGARefLink{par:sga3-VIA-1.1}{Theorem~1.1}).

\medskip
\phantomsection\label{cor:sga3-VIA-6.10}%
\noindent\textbf{Corollary 6.10.---}
\emph{Let \(G\) be a quasi-compact \(k\)-group, and suppose that
\(k\) is algebraically closed.
\begin{enumerate}
\item[\textup{(i)}] If \(f:G\to H\) is a faithfully flat morphism
  of \(k\)-groups, then the induced map \(G(k)\to H(k)\) is
  surjective.
\item[\textup{(ii)}] The set of rational points is dense in \(G\).
\end{enumerate}}

For the proof, we refer to \textup{[Per75]}, V, \SGARefLink{prop:sga3-VIB-3.7}{3.7} and 3.9.

\section*{Bibliography}
\phantomsection\label{bib:sga3-VIA}
\addcontentsline{toc}{section}{Bibliography}
\footnote{\label{note:sga3-VIA-59}Editor's note: additional
references cited in this Expos\'e.}

\begin{description}
\item[{[An73]}] S. Anantharaman,
  \emph{Sch\'emas en groupes, espaces homog\`enes et espaces
  alg\'ebriques sur une base de dimension 1},
  M\'em. Soc. Math. France \textbf{33} (1973), 5--79.
\item[{[BAC]}] N. Bourbaki,
  \emph{Alg\`ebre commutative}, Chapters I--IV and V--VII,
  Masson, 1985.
\item[{[DG70]}] M. Demazure and P. Gabriel,
  \emph{Groupes alg\'ebriques}, Masson \& North-Holland, 1970.
\item[{[Per75]}] D. Perrin,
  \emph{Sch\'emas en groupes quasi-compacts sur un corps},
  Publ. Math. Orsay no.~165--75.46 (part 1),
  \texttt{http://portail.mathdoc.fr/PMO/}.
\item[{[Per76]}] D. Perrin,
  \emph{Approximation des sch\'emas en groupes, quasi-compacts
  sur un corps}, Bull. Soc. Math. France \textbf{104} (1976),
  323--335.
\item[{[Ray67a]}] M. Raynaud,
  \emph{Passage au quotient par une relation d'\'equivalence
  plate}, pp.~78--85 in \emph{Proc. Conf. Local Fields
  (Driebergen)}, ed.\ T.~A. Springer, Springer-Verlag, 1967.
\item[{[Ray67b]}] M. Raynaud,
  \emph{Sur le passage au quotient par un groupo\"\i de plat},
  C. R. Acad. Sci. Paris (S\'er. A) \textbf{265} (1967),
  384--387.
\item[{[Ray70]}] M. Raynaud,
  \emph{Faisceaux amples sur les sch\'emas en groupes et les
  espaces homog\`enes}, Lecture Notes in Mathematics
  \textbf{119}, Springer-Verlag, 1970.
\end{description}

% End of Expose VI_A. Exp6B begins in the next component.

% END INCLUDED SGA3_English_current_buildable_source_20260823/source/inputs/VI/tex_reference_v2/components/25_expose_VIA_theorem65_through_bibliography_loop2_r1_en.tex

% BEGIN INCLUDED SGA3_English_current_buildable_source_20260823/source/inputs/VI/tex_reference_v2/components/26_expose_VIB_title_through_lemma121_en.tex
% SGA 3, Expose VI_B, Loop 1 English reconstruction.
% Scope: title and opening through Lemma 1.2.1 in full;
% Counterexample 1.2.2 is excluded.
% Authority: Polo--Gille Exp6B-13oct24.pdf, component-local pages 1--2
% (combined-reader pages 342--343; original-margin printed pages
% 319--321).
% OCR and the Reinhold English Markdown are comparison witnesses only.

\setcounter{footnote}{-1}
\chapter*{Expos\'e \texorpdfstring{VI\textsubscript{B}}{VI_B}\\
Generalities on Group Schemes}
\phantomsection
\label{sga3:VIB:title}
\addcontentsline{toc}{chapter}{Expos\'e VI-B. Generalities on Group Schemes}
\markboth{EXPOS\'E VI\textsubscript{B}. GENERALITIES ON GROUP SCHEMES}%
  {EXPOS\'E VI\textsubscript{B}. GENERALITIES ON GROUP SCHEMES}

\begin{center}
  \textit{by J.-E. Bertin}
\end{center}

\footnote{\label{note:sga3-VIB-0}Editor's note: version of
October 13, 2024.}
This Expos\'e, which does not correspond to any oral lecture in the
seminar, is intended to bring together a number of commonly used
technical results concerning group schemes.%
\begingroup
\renewcommand{\thefootnote}{\fnsymbol{footnote}}%
\footnote[1]{This Expos\'e was quite substantially revised after its
mimeographed edition; in particular, \hyperref[sec:sga3-VIB-10]{Sections 10} and 11 were completely
rewritten.\textsuperscript{(1)}}%
\endgroup
\footnotetext[1]{Editor's note: and also \hyperref[sec:sga3-VIB-5]{Section 5}; cf.\ Editor's
\hyperref[note:sga3-VIB-34]{notes 34} and 46.}
\setcounter{footnote}{1}

\section*{1. Morphisms of groups locally of finite type over a field}
\phantomsection
\label{sec:sga3-VIB-1}
\addcontentsline{toc}{section}{1. Morphisms of groups locally of finite
  type over a field}

\medskip
\phantomsection\label{par:sga3-VIB-1.1}%
\noindent\textbf{1.1.---}
Let \(A\) be an Artinian local ring, let \(G\) and \(H\) be two
\(A\)-groups,%
\footnote{\label{note:sga3-VIB-2}Editor's note: recall the convention
adopted at the beginning of Expos\'e~VI\(_A\): an \(A\)-group is a
group scheme over \(A\) (and every such scheme is separated;
cf.\ Expos\'e~VI\(_A\), \SGARefLink{sga3:VIA:0.3}{0.3}).}
and let \(u:G\to H\) be a morphism of \(A\)-groups. Then \(u\)
induces a group morphism
\[
  u(A):G(A)\longrightarrow H(A).
\]%
\footnote{\label{note:sga3-VIB-3}Editor's note: we have replaced
``homomorphism'' by ``morphism.''}
Since \(H(A)\) acts on \(H\) by right translation, restriction along
\(u(A)\) defines an action of \(G(A)\) on \(H\). This action is
compatible with \(u\) and with the action of \(G(A)\) on \(G\) by
right translation. Since \(G(A)\) acts transitively on the strictly
rational points of \(G\) (see Expos\'e~VI\(_A\), \SGARefLink{sga3:VIA:0.4}{0.4} for the
definition), these points ``all behave in the same way with respect
to \(u\).'' The following properties spring from this observation.

\medskip
\phantomsection\label{prop:sga3-VIB-1.2}%
\noindent\textbf{Proposition 1.2.---}
\emph{Let \(u:G\to H\) be a quasi-compact morphism between
\(A\)-groups locally of finite type over \(A\). Then the set
\(u(G)\) is closed in \(H\), and}
\[
  \dim G=\dim u(G)+\dim\ker u.
\]%
\footnote{\label{note:sga3-VIB-4}Editor's note: this statement also
appears in Expos\'e~VI\(_A\), \SGARefLink{prop:sga3-VIA-2.5.2}{2.5.2}.}

Since \(u\) commutes with the inversion morphisms of \(G\) and \(H\),
the image \(u(G)\) is invariant under inversion on \(H\); the same is
therefore true of its closure \(\overline{u(G)}\) in \(H\). On the
other hand, let \(L\) be the set of points of \(H\times_A H\) whose
two projections belong to \(u(G)\). It is clear that \(L\) is the
image of
\[
  u\times_Au:G\times_A G\longrightarrow H\times_A H.
\]
The multiplication morphism of \(H\) consequently sends \(L\) into
\(u(G)\); in other words,
\[
  u(G)\cdot u(G)=u(G).
\]
\hyperref[lem:sga3-VIB-1.2.1]{Lemma~1.2.1} below shows that the closure \(\overline L\) of \(L\)
in \(H\times_AH\) is the set of points whose two projections belong
to \(\overline{u(G)}\). Hence
\[
  \overline{u(G)}\cdot\overline{u(G)}
    =\overline{u(G)}.
\]
Thus the reduced subscheme of \(H\) whose underlying space is
\(\overline{u(G)}\) is naturally equipped with a group structure in
the category \((\mathbf{Sch}_{/k})_{\mathrm{red}}\), where \(k\) is
the residue field of \(A\) (cf.\ Expos\'e~VI\(_A\), \SGARefLink{sga3:VIA:0.2}{0.2}).

Let us prove the first assertion of \hyperref[prop:sga3-VIB-1.2]{1.2}. After replacing \(A\) by
an algebraic closure of its residue field \(k\), we may suppose that
\(A\) is an algebraically closed field \(k\)
(cf.\ EGA~IV\(_2\), 2.3.12). Replacing \(u\) by
\[
  u_{\mathrm{red}}:G_{\mathrm{red}}\longrightarrow H_{\mathrm{red}},
\]
we may suppose that \(G\) and \(H\) are reduced. In this case, as we
have just seen, \(u(G)\) is the underlying space of a reduced group
subscheme of \(H\);%
\begingroup
\renewcommand{\thefootnote}{\fnsymbol{footnote}}%
\footnote[1]{\label{note:sga3-VIB-1.2-image-ambient}Translator's
note: the French source prints ``a reduced group subscheme of \(G\)''
at this point. Since \(u(G)\subset H\), and the preceding paragraph
constructs the reduced subscheme inside \(H\), the ambient group is
\(H\).}%
\endgroup{}
we may therefore suppose that \(u\) is dominant. Then \(G(k)\) acts
transitively on the set of connected components of \(H\), and it is
enough to prove that \(u(G)\cap H^0\) is closed. We are thus reduced
to the case in which \(H\) is connected, hence irreducible and of
finite type (Expos\'e~VI\(_A\), \SGARefLink{cor:sga3-VIB-2.4}{2.4}). The morphism \(u\) is then of
finite type, since it is quasi-compact and locally of finite type.
As \(H\) is Noetherian, \(u(G)\) is constructible
(EGA~IV\(_1\), 1.8.5), and therefore contains an open subset \(V\)
of \(H\) (EGA~0\(_{\mathrm{III}}\), 9.2.2). By
Expos\'e~VI\(_A\), 0.5, we then have
\[
  H=V\cdot V\subset u(G)\cdot u(G)=u(G).
\]

Let us prove the second assertion. Recall first that the functor
\(\ker u\) (cf.\ \SGARefLink{sga3:i:definition:2-3-6-1}{Expos\'e~I, 2.3.6.1}) is represented by
\(u^{-1}(e)\), where \(e\) denotes the identity element of \(H\).
When \(u\) is locally of finite type, \(\ker u\) is therefore
locally of finite type over \(A\). As above, this time using
EGA~IV\(_2\), 4.1.4, we reduce to the case in which \(A\) is an
algebraically closed field \(k\).

We may further suppose that \(G\) and \(H\) are irreducible and of
finite type and that \(u\) is dominant. Indeed, since \(k\) is
algebraically closed, the connected components of \(G\), on whose
set \(G(k)\) acts transitively, clearly all have the same dimension;
and if \(u^0\) denotes the restriction of \(u\) to \(G^0\), then
\[
  (\ker u)^0\subset\ker u^0,
  \qquad
  \dim\ker u^0=\dim\ker u.
\]
Likewise, \(u\) is then of finite type over \(k\). If \(\eta\) is
the generic point of \(H\), then
\[
  \dim u^{-1}(\eta)=\dim G-\dim H
\]
(EGA~IV\(_3\), 10.6.1(ii)). By EGA~IV\(_3\), 9.2.3 and 9.2.6, the
set of \(y\in H\) such that
\[
  \dim u^{-1}(y)=\dim u^{-1}(\eta)
\]
contains a nonempty open subset \(V\). Since \(u\) is dominant,
\(U=u^{-1}(V)\) is a nonempty open subset of \(G\); it contains a
closed point \(x\), because \(G\) is a Jacobson scheme
(EGA~IV\(_3\), 10.4.7). Right translation \(r_x\) is an isomorphism
from \(\ker u\) onto \(u^{-1}(u(x))\), and hence
\[
  \dim\ker u
    =\dim u^{-1}(u(x))
    =\dim u^{-1}(\eta)
    =\dim G-\dim H.
\]

\medskip
\phantomsection\label{lem:sga3-VIB-1.2.1}%
\noindent\textbf{Lemma 1.2.1.---}
\emph{Let \(f:X'\to X\) and \(g:Y'\to Y\) be two quasi-compact
dominant morphisms of schemes over an Artinian local ring \(A\).
Then}
\[
  f\times_Ag:X'\times_A Y'\longrightarrow X\times_A Y
\]
\emph{is dominant (and quasi-compact).}

Indeed,
\[
  f\times_Ag
    =
  (\operatorname{id}_X\times_Ag)
    \circ
  (f\times_A\operatorname{id}_{Y'}).
\]
It is therefore enough to prove that
\(f\times_A\operatorname{id}_{Y'}\) and
\(\operatorname{id}_X\times_Ag\) are dominant. For this purpose we
may replace \(A\) by its residue field \(k\). Then \(X\) and \(Y'\)
are flat over \(A=k\), and
\(f\times_A\operatorname{id}_{Y'}\)
(respectively, \(\operatorname{id}_X\times_Ag\)) is obtained from
\(f\) (respectively, \(g\)) by the flat base change \(Y'\to A\)
(respectively, \(X\to A\)). It is therefore dominant (and
quasi-compact) by EGA~IV\(_2\), 2.3.7.

% END INCLUDED SGA3_English_current_buildable_source_20260823/source/inputs/VI/tex_reference_v2/components/26_expose_VIB_title_through_lemma121_en.tex

% BEGIN INCLUDED SGA3_English_current_buildable_source_20260823/source/inputs/VI/tex_reference_v2/components/27_expose_VIB_counter122_prop13_en.tex
% SGA 3, Expose VI_B, Loop 1 English reconstruction.
% Scope: Counterexample 1.2.2 and Proposition 1.3, including
% Lemma 1.3.0 and the complete proof of Proposition 1.3;
% Corollary 1.3.1 is excluded.
% Authority: Polo--Gille Exp6B-13oct24.pdf, component-local pages 3--4
% (combined-reader pages 344--345; original-margin printed page 322,
% ending at the seam before printed page 323).
% OCR and the Reinhold English Markdown are comparison witnesses only.

\medskip
\phantomsection\label{ctr:sga3-VIB-1.2.2}%
\noindent\textbf{Counter\SGARefLink{ctr:sga3-VIB-1.2.2}{example 1.2.2}.---}
\emph{Let \(k\) be a field of characteristic zero, let \(G\) be
the constant \(k\)-group \(\mathbb Z\), and let \(H\) be the
additive \(k\)-group \(\mathbb G_{a,k}\). Let \(u:G\to H\) be a
morphism of \(k\)-groups. If \(u\ne0\), then \(u(G)\) is not closed
in \(H\).}

\medskip
\phantomsection\label{prop:sga3-VIB-1.3}%
\noindent\textbf{Proposition 1.3.---}%
\footnote{\label{note:sga3-VIB-5}Editor's note: in the original,
\SGARefLink{prop:sga3-VIB-1.3}{1.3}, \SGARefLink{cor:sga3-VIB-1.3.1}{1.3.1}, and \SGARefLink{lem:sga3-VIB-1.3.1.1}{1.3.1.1} were stated with \(A\) a field. For
completeness, we treat the case of an Artinian local ring, and for
this purpose have added \hyperref[lem:sga3-VIB-1.3.0]{Lemma~1.3.0}.}
\emph{Let \(A\) be an Artinian local ring with residue field \(k\),
let \(G\) be an \(A\)-group locally of finite type and flat, and
let \(u:G\to H\) be a morphism of \(A\)-groups. The following
assertions are equivalent:
\begin{enumerate}
\item[\textup{(i)}]\phantomsection\label{concl:sga3-VIB-1.3-i} \(u\) is flat (respectively, quasi-finite,
  unramified, smooth, or \'etale) at some point of \(G\);%
  \footnote{\label{note:sga3-VIB-6}Editor's note: recall that a
  morphism of schemes \(f:X\to Y\) is said to be of finite type
  (respectively, of finite presentation or quasi-finite) at \(x\)
  if there are affine open subsets \(V\), containing \(f(x)\), and
  \(U\), containing \(x\), such that \(f(U)\subset V\) and
  \(\mathcal O_X(U)\) is an \(\mathcal O_Y(V)\)-algebra of finite
  type (respectively, of finite presentation, or with the
  additional condition that \(f^{-1}(f(x'))\) is finite for every
  \(x'\in U\)); see also Editor's note 40 of Expos\'e~V. The
  morphism is locally quasi-finite (respectively, locally of finite
  type or locally of finite presentation) if it has the property
  at every point \(x\). It is smooth (respectively, unramified or
  \'etale) at \(x\) if there is an open neighborhood \(U\) of \(x\)
  such that \(f|_U:U\to Y\) is smooth (respectively, unramified or
  \'etale). These definitions show that the loci where \(f\) is of
  finite presentation, of finite type, quasi-finite, smooth,
  unramified, or \'etale are open in \(X\). Moreover, because the
  flatness locus of a morphism locally of finite presentation is
  open (EGA~IV\(_3\), 11.3.1), the locus where \(f\) is both of
  finite presentation and flat is open as well. These facts will
  be used repeatedly below.}
\item[\textup{\SGAInlineRefTarget{prop:sga3-VIB-1.3:item:ii}(ii)}] \(u\) is flat (respectively, quasi-finite,
  unramified, smooth, or \'etale).
\end{enumerate}}

\emph{Proof.}
It is enough to prove that \textup{\SGARefLink{concl:sga3-VIB-1.3-i}{(i)}} implies \textup{\SGARefLink{prop:sga3-VIB-1.3:item:ii}{(ii)}}. We
first establish the following lemma.

\medskip
\phantomsection\label{lem:sga3-VIB-1.3.0}%
\noindent\textbf{Lemma 1.3.0.---}
\emph{Let \(A\to B\to C\) be morphisms of commutative rings, and
let \(\mathfrak n\) be a nilpotent ideal of \(A\). Suppose that
\(C/\mathfrak nC\) is a
\((B/\mathfrak nB)\)-algebra of finite type.
\begin{enumerate}
\item[\textup{\SGAInlineRefTarget{lem:sga3-VIB-1.3.0:item:i}(i)}] Then \(C\) is a \(B\)-algebra of finite type.
\item[\textup{\SGAInlineRefTarget{lem:sga3-VIB-1.3.0:item:ii}(ii)}] Moreover, if \(C\) is flat over \(A\) and
  \(C/\mathfrak nC\) is a
  \((B/\mathfrak nB)\)-algebra of finite presentation, then \(C\)
  is a \(B\)-algebra of finite presentation.
\end{enumerate}}

Indeed, let \(x_1,\ldots,x_n\) be elements of \(C\) whose images
generate \(C/\mathfrak nC\) as a
\((B/\mathfrak nB)\)-algebra. By the nilpotent Nakayama lemma, the
\(x_i\) generate \(C\) as a \(B\)-algebra. This proves \SGARefLink{lem:sga3-VIB-1.3.0:item:i}{(i)}. Let
\[
  \varphi:B[X_1,\ldots,X_n]\longrightarrow C
\]
be the resulting surjective morphism, and put \(I=\ker\varphi\).

Suppose now that \(C\) is flat over \(A\) and that
\(C/\mathfrak nC\) is of finite presentation over
\(B/\mathfrak nB\). On the one hand, \(I/\mathfrak nI\) identifies
with the kernel of
\[
  \overline\varphi
    =\varphi\otimes_A(A/\mathfrak n).
\]
On the other hand, by EGA~IV\(_1\), 1.4.4,
\(\ker\overline\varphi\) is an ideal of finite type. Choose
\(P_1,\ldots,P_s\in I\) whose images generate
\(I/\mathfrak nI\) as an ideal. The nilpotent Nakayama lemma shows
that the \(P_i\) generate \(I\). This proves \SGARefLink{lem:sga3-VIB-1.3.0:item:ii}{(ii)}.

We return to the proof of \hyperref[prop:sga3-VIB-1.3]{1.3}. Let \(x\) be an arbitrary point of
\(G\). Since \(G\) is flat over \(A\), the fiberwise flatness
criterion, in the form EGA~IV\(_3\), 11.3.10.2, shows that \(u\)
is flat at \(x\) if \(u\otimes_Ak\) is. Likewise, the preceding
lemma shows that \(u\) is of finite type (respectively, of finite
presentation) at \(x\) if \(u\otimes_Ak\) is. The remaining
properties can then be checked on the fibers
(cf.\ EGA~IV\(_4\), 17.4.1, 17.5.1, and 17.6.1 for the
unramified, smooth, and \'etale cases). We are therefore reduced to
the case \(A=k\).

Now let \(x\) be a point of \(G\) at which one of the conditions
in \hyperref[concl:sga3-VIB-1.3-i]{1.3(i)} holds. The properties under consideration are preserved
by \(\mathrm{fpqc}\) descent (cf.\ EGA~IV, 2.5.1, 2.7.1, and
17.7.1). Thus, after replacing \(k\) by an algebraic closure of
\(\kappa(x)\), we reduce to the case in which \(k\) is
algebraically closed and \(x\in G(k)\).

Since \(G\) is a Jacobson scheme (cf.\ EGA~IV\(_3\), 10.4.7), and
the set \(W\) of points of \(G\) at which \(u\) is flat
(respectively, quasi-finite, unramified, smooth, or \'etale) is
stable under generization%
\footnote{\label{note:sga3-VIB-7}Editor's note: the original stated
that the flat locus \(W\) is open, which does not seem evident
a priori.}
(respectively, open), it is enough to show that every
\(y\in G(k)\) belongs to \(W\). For such a point \(y\), the
translation
\[
  r_y\circ r_x^{-1}
\]
sends \(x\) to \(y\), so \(u\) has the required property at \(y\).
Thus \(y\in W\), and the proof is complete.

% END INCLUDED SGA3_English_current_buildable_source_20260823/source/inputs/VI/tex_reference_v2/components/27_expose_VIB_counter122_prop13_en.tex

% BEGIN INCLUDED SGA3_English_current_buildable_source_20260823/source/inputs/VI/tex_reference_v2/components/28_expose_VIB_cor131_cor132_en.tex
% SGA 3, Expose VI_B, Loop 1 English reconstruction.
% Scope: Corollary 1.3.1 through the complete proof of
% Corollary 1.3.2, including Lemma 1.3.1.1; Proposition 1.4 is
% excluded.
% Authority: Polo--Gille Exp6B-13oct24.pdf, component-local pages
% 4--5 (combined-reader pages 345--346; original-margin printed
% pages 323--324). This corrects the supplied page locator while
% preserving the exact theorem boundary.
% OCR and the Reinhold English Markdown are comparison witnesses
% only.

\medskip
\phantomsection\label{cor:sga3-VIB-1.3.1}%
\noindent\textbf{Corollary 1.3.1.---}
\emph{Let \(A\) be an Artinian local ring with residue field \(k\),
and let \(G\) be a flat \(A\)-group. The following assertions are
equivalent:}%
\footnote{\label{note:sga3-VIB-8}Editor's note: it is not enough to
suppose that \(G\) is flat over \(A\) at one point. For example,
suppose that \(A\ne k\), let \(H\) be the constant \(A\)-group
\((\mathbb Z/2\mathbb Z)_A\), and let \(G\) be the closed
\(A\)-subgroup of \(H\) whose nonidentity component is reduced.
Then the structural morphism \(G\to\operatorname{Spec}A\) is a
local isomorphism at the identity point \(e\), but is not flat at
the point \(g\ne e\).}
\begin{enumerate}
\item[\textup{\SGAInlineRefTarget{cor:sga3-VIB-1.3.1:item:i}(i)}] \emph{\(G\) is locally quasi-finite
  (respectively, unramified, smooth, or \'etale) over \(A\) at
  some point;}
\item[\textup{\SGAInlineRefTarget{cor:sga3-VIB-1.3.1:item:ii}(ii)}] \emph{\(G\) is locally quasi-finite
  (respectively, unramified, smooth, or \'etale) over \(A\).}
\end{enumerate}

\footnote{\label{note:sga3-VIB-9}Editor's note: we have added the
sentence that follows and simplified the proof of
\hyperref[lem:sga3-VIB-1.3.1.1]{Lemma~1.3.1.1} by invoking EGA~IV, 2.7.1 instead of
loc.\ cit., 17.7.5.}
\emph{Proof.}
If \(G\) satisfies one of the conditions in \textup{\SGARefLink{cor:sga3-VIB-1.3.1:item:i}{(i)}} at a
point \(x\), there is an open neighborhood \(U\) of \(x\) that is
of finite type over \(A\). Consequently, it is enough to apply
\hyperref[prop:sga3-VIB-1.3]{1.3} in the case where \(H\) is the unit \(A\)-group, taking the
following lemma into account.

\medskip
\phantomsection\label{lem:sga3-VIB-1.3.1.1}%
\noindent\textbf{Lemma 1.3.1.1.---}
\emph{Let \(A\) be an Artinian local ring and \(G\) an
\(A\)-group. If \(G\) contains a nonempty open subset of finite
type over \(A\), then \(G\) is locally of finite type over \(A\).}

\emph{Proof.}
By \hyperref[lem:sga3-VIB-1.3.0]{Lemma~1.3.0}, we may suppose that \(A\) equals its residue field
\(k\). Moreover, by \(\mathrm{fpqc}\) descent, we may suppose that
\(k\) is algebraically closed (cf.\ EGA~IV\(_2\), 2.7.1). Let
\(V\) be the open subset of \(G\) consisting of the points at
which \(G\) is of finite type over \(k\); by hypothesis,
\(V\ne\varnothing\). Since \(G\) is a Jacobson scheme, \(V\)
contains a closed point \(x\), and to prove that \(V=G\) it is
enough to show that every closed point \(y\) of \(G\) belongs to
\(V\). For such a point \(y\), the translation
\[
  r_y\circ r_x^{-1}
\]
sends \(x\) to \(y\), and hence \(y\in V\).

\medskip
\phantomsection\label{cor:sga3-VIB-1.3.2}%
\noindent\textbf{Corollary 1.3.2.---}
\emph{Let \(A\) be an Artinian local ring, and let \(u:G\to H\)
be a morphism between \(A\)-groups locally of finite type. The
following assertions are equivalent:
\begin{enumerate}
\item[\textup{\SGAInlineRefTarget{cor:sga3-VIB-1.3.2:item:i}(i)}] \(u\) is universally open;
\item[\textup{\SGAInlineRefTarget{cor:sga3-VIB-1.3.2:item:ii}(ii)}] \(u\) is open;
\item[\textup{\SGAInlineRefTarget{cor:sga3-VIB-1.3.2:item:iii}(iii)}] \(u\) is open at some point of \(G\);
\item[\textup{\SGAInlineRefTarget{cor:sga3-VIB-1.3.2:item:iv}(iv)}] the morphism
  \(u^0:G^0\to H^0\) induced by \(u\) is dominant;
\item[\textup{(iv\('\))}] \(u^0\) is surjective;
\item[\textup{\SGAInlineRefTarget{cor:sga3-VIB-1.3.2:item:v}(v)}] there is a connected component \(C\) of \(G\)
  such that, if \(D\) denotes the connected component of \(H\)
  containing \(u(C)\), the induced morphism
  \(u':C\to D\) is dominant.
\end{enumerate}}

\emph{Proof.}
The implications
\[
  \textup{(i)}\Longrightarrow\textup{(ii)}
  \Longrightarrow\textup{(iii)}
  \qquad\text{and}\qquad
  \textup{(iv\('\))}\Longrightarrow\textup{(iv)}
  \Longrightarrow\textup{(v)}
\]
are clear. Since \(G^0\) is of finite type over \(A\)
(Expos\'e~VI\(_A\), \SGARefLink{cor:sga3-VIB-2.4}{2.4}), and hence Noetherian, \(u^0\) is
quasi-compact. Thus \(u^0(G^0)\) is closed in \(H^0\) by \hyperref[prop:sga3-VIB-1.2]{1.2}, and
therefore \textup{\SGARefLink{cor:sga3-VIB-1.3.2:item:iv}{(iv)}} implies \textup{(iv\('\))}. On the other
hand, \(G^0\) (respectively, \(C\)) is open in \(G\)
(Expos\'e~VI\(_A\), \SGARefLink{cor:sga3-VIB-2.3}{2.3}), while \(H^0\) (respectively, \(D\)) is
irreducible (Expos\'e~VI\(_A\), \SGARefLink{cor:sga3-VIA-2.4.1}{2.4.1}). Hence \textup{\SGARefLink{cor:sga3-VIB-1.3.2:item:ii}{(ii)}}
implies \textup{\SGARefLink{cor:sga3-VIB-1.3.2:item:iv}{(iv)}} (respectively, \textup{\SGARefLink{cor:sga3-VIB-1.3.2:item:iii}{(iii)}} implies
\textup{\SGARefLink{cor:sga3-VIB-1.3.2:item:v}{(v)}}). It remains to prove that \textup{\SGARefLink{cor:sga3-VIB-1.3.2:item:v}{(v)}} implies
\textup{\SGARefLink{cor:sga3-VIB-1.3.2:item:i}{(i)}}.

We endow the open subset \(C\) (respectively, \(D\)) of \(G\)
(respectively, \(H\)) with its induced scheme structure, and let
\(u'\) denote the induced morphism \(C\to D\). Let \(k\) be the
residue field of \(A\).%
\footnote{\label{note:sga3-VIB-10}Editor's note: we have expanded
the original in what follows.}
Since the base change
\[
  \operatorname{Spec}k\longrightarrow\operatorname{Spec}A
\]
is a universal homeomorphism, we may suppose that \(A=k\). By
hypothesis \(u'\) is dominant; and since \(C\) is of finite type
over \(k\) (Expos\'e~VI\(_A\), 2.4.1), and hence Noetherian,
\(u'\) is quasi-compact. By EGA~IV\(_2\), 2.3.7,
\[
  u'\otimes_k\bar{k}
\]
is again quasi-compact and dominant, where \(\bar{k}\) denotes an
algebraic closure of \(k\). Since
\(C\otimes_k\bar{k}\) is a union of connected components of
\(G\otimes_k\bar{k}\), the morphism
\[
  u\otimes_k\bar{k}:
  G\otimes_k\bar{k}\longrightarrow H\otimes_k\bar{k}
\]
satisfies assertion \textup{(v)}. In view of
EGA~IV\(_2\), 2.6.4, we are thus reduced to the case where
\(A=k\) is an algebraically closed field.

In this case we may further replace \(u\) by \(u_{\mathrm{red}}\),
and are thereby reduced to the case where \(H\) is reduced. Let
\(\xi\) (respectively, \(\eta\)) be the generic point of \(C\)
(respectively, \(D\)). Since \(u'\) is quasi-compact and dominant,
\[
  u'(\xi)=\eta
\]
(cf.\ EGA~IV\(_1\), 1.1.5). Since \(H\) is reduced, the local
ring \(\mathcal O_{H,\eta}\) is a field, and hence \(u'\) is flat
at \(\xi\).%
\footnote{\label{note:sga3-VIB-11}Editor's note: the original
invoked the generic flatness theorem (EGA~IV\(_2\), 6.9.1),
which is not necessary here.}
It follows from \SGARefLink{prop:sga3-VIB-1.3}{1.3} that \(u\) is flat. Moreover, since \(u\) is
locally of finite type and \(H\) is locally Noetherian, \(u\) is
locally of finite presentation. Therefore \(u\) is universally
open by EGA~IV\(_2\), 2.4.6.


% END INCLUDED SGA3_English_current_buildable_source_20260823/source/inputs/VI/tex_reference_v2/components/28_expose_VIB_cor131_cor132_en.tex

% BEGIN INCLUDED SGA3_English_current_buildable_source_20260823/source/inputs/VI/tex_reference_v2/components/29_expose_VIB_prop14_cor142_en.tex
% SGA 3, Expose VI_B, Loop 1 English reconstruction.
% Scope: Proposition 1.4 through Corollary 1.4.2 in full;
% Counterexample 1.4.3 is excluded.
% Authority: Polo--Gille Exp6B-13oct24.pdf, component-local pages 5--7
% (combined-reader pages 346--348; original-margin printed pages
% 325--327).
% OCR and the Reinhold English Markdown are comparison witnesses only.

\medskip
\phantomsection\label{prop:sga3-VIB-1.4}%
\noindent\textbf{Proposition 1.4.---}
\emph{Let \(A\) be an Artinian local ring, and let \(u:G\to H\)
be a quasi-compact morphism between \(A\)-groups locally of finite
type. The following assertions are equivalent:
\begin{enumerate}
\item[\textup{\SGAInlineRefTarget{prop:sga3-VIB-1.4:item:i}(i)}] \(u\) is proper;
\item[\textup{\SGAInlineRefTarget{prop:sga3-VIB-1.4:item:ii}(ii)}] there exists \(h\in H\) such that the fiber
  \(u^{-1}(h)\) is nonempty and proper over \(\kappa(h)\);
\item[\textup{\SGAInlineRefTarget{prop:sga3-VIB-1.4:item:iii}(iii)}] \(\ker u\) is proper over \(A\).
\end{enumerate}}

\emph{Proof.}
It is clear that \textup{\SGARefLink{prop:sga3-VIB-1.4:item:i}{(i)}} implies \textup{\SGARefLink{prop:sga3-VIB-1.4:item:iii}{(iii)}}, and that
\textup{\SGARefLink{prop:sga3-VIB-1.4:item:iii}{(iii)}} implies \textup{\SGARefLink{prop:sga3-VIB-1.4:item:ii}{(ii)}}. The hypotheses imply that
\(u\) is of finite type; and since \(G\) is separated
(Expos\'e~VI\(_A\), 0.3), \(u\) is separated as well
(EGA~I, 5.5.1). It remains to prove that \textup{(ii)} implies
that \(u\) is universally closed. We may therefore suppose that
\(A\) equals its residue field \(k\).%
\footnote{\label{note:sga3-VIB-12}Editor's note: we have added the
following sentence.}
Let \(k'\) be an algebraic closure of \(\kappa(h)\), let
\[
  u':G'\longrightarrow H'
\]
be obtained from \(u\) by base change, and let \(h'\in H'\) be a
point above \(h\). Then
\[
  u'^{-1}(h')=u^{-1}(h)\times_{\kappa(h)}k'
\]
is nonempty and proper, and it is enough to prove that \(u'\) is
proper (EGA~IV\(_2\), 2.6.4). Thus we may suppose that \(k\) is
algebraically closed and \(h\in H(k)\).

By \hyperref[prop:sga3-VIB-1.2]{1.2}, \(u(G)\) is then the underlying set of a closed reduced
group subscheme of \(H\). Since every closed immersion is proper
(EGA~II, 5.4.2), we may suppose that \(u\) is surjective and that
\(H\) is reduced. The group \(G(k)\) then acts transitively on the
set of closed points of \(H\). Consequently, for every closed point
\(y\in H\), the fiber \(u^{-1}(y)\) is proper over \(\kappa(y)\).
By EGA~IV\(_3\), 9.6.1, the set of \(y\in H\) for which
\(u^{-1}(y)\) is not proper over \(\kappa(y)\) is locally
constructible. It contains no closed point, and is therefore empty
(cf.\ EGA~IV\(_3\), 10.3.1 and 10.4.7).

\footnote{\label{note:sga3-VIB-13}Editor's note: we have expanded
the original in what follows.}
Now consider the generic point \(\eta\) of \(H^0\). By what
precedes, the fiber
\[
  u^{-1}(\eta)=G\times_H\operatorname{Spec}\kappa(\eta)
\]
is proper over \(\kappa(\eta)\). Since \(H\) is reduced,
\(\kappa(\eta)=\mathcal O_{H,\eta}\). The latter ring is the direct
limit of the rings \(\mathcal O_H(V)\), as \(V\) ranges over the
nonempty open subsets of \(H^0\). It follows from
EGA~IV\(_3\), 8.1.2(a) and 8.10.5(xii), that there exists a
nonempty open subset \(V\subset H^0\) such that the restriction of
\(u\) over \(V\) is proper. The translates \(g\cdot V\), for
\(g\in G(k)\), form an open covering of \(H\); and for every
\(g\in G(k)\), the restriction of \(u\) over \(g\cdot V\) is
proper. Hence \(u\) is proper (cf.\ EGA~II, 5.4.1).

\medskip
\phantomsection\label{cor:sga3-VIB-1.4.1}%
\noindent\textbf{Corollary 1.4.1.---}
\emph{Let \(A\) be an Artinian local ring, and let \(u:G\to H\)
be a morphism between \(A\)-groups locally of finite type. The
following assertions are equivalent:
\begin{enumerate}
\item[\textup{(i)}] \(u\) is locally quasi-finite;
\item[\textup{(ii)}] \(u\) is quasi-finite at some point;
\item[\textup{(iii)}] \(\ker u\) is discrete;
\item[\textup{\SGAInlineRefTarget{cor:sga3-VIB-1.4.1:item:iv}(iv)}] the restriction of \(u\) to each connected
  component of \(G\) is finite.
\end{enumerate}
If \(u\) is quasi-compact, these assertions are also equivalent to
\begin{enumerate}
\item[\textup{(v)}] \(u\) is finite.
\end{enumerate}}

It is clear that \textup{(iv)} implies \textup{(iii)}, that
\textup{(iii)} implies \textup{(ii)} (EGA~I, 6.4.4), and that,
when \(u\) is quasi-compact, \textup{(iv)} and \textup{(v)} are
equivalent. We have already seen in \hyperref[prop:sga3-VIB-1.3]{1.3} that \textup{\SGARefLink{concl:sga3-VIB-1.3-i}{(i)}} and
\textup{\SGARefLink{prop:sga3-VIB-1.3:item:ii}{(ii)}} are equivalent.

It remains to prove that \textup{\SGARefLink{concl:sga3-VIB-1.3-i}{(i)}} implies \textup{\SGARefLink{cor:sga3-VIB-1.4.1:item:iv}{(iv)}}. Let
\(C\) be a connected component of \(G\). Since \(C\) is of finite
type over \(A\) (Expos\'e~VI\(_A\), \SGARefLink{cor:sga3-VIA-2.4.1}{2.4.1}) and \(G,H\) are
separated (Expos\'e~VI\(_A\), \SGARefLink{sga3:VIA:0.3}{0.3}),%
\begingroup
\renewcommand{\thefootnote}{\fnsymbol{footnote}}%
\footnote[1]{\label{note:sga3-VIB-1.4.1-separated-xref}Translator's
note: the French source cites Expos\'e~VI\(_A\), 0.2 here; the
separation statement is \SGARefLink{sga3:VIA:0.3}{0.3}, as also recalled in Editor's \hyperref[note:sga3-VIB-2]{note 2}
above.}%
\endgroup{}
EGA~I, 5.5.1 and 6.3.4 show
that the restriction \(u'\) of \(u\) to \(C\) is separated and of
finite type.%
\footnote{\label{note:sga3-VIB-14}Editor's note: we have expanded
the original in what follows.}
Its fibers are discrete, so \(u'\) is quasi-finite
(cf.\ EGA~II, 6.2.2). Since every quasi-finite proper morphism is
finite (cf.\ EGA~III\(_1\), 4.4.2), it is enough to prove that
\(u'\) is universally closed.

For this, we may suppose that \(A\) equals its residue field \(k\).
By \(\mathrm{fpqc}\) descent (cf.\ EGA~IV\(_2\), 2.6.4), it is
enough to prove that
\[
  u'\otimes_k\overline k
\]
is universally closed, where \(\overline k\) is an algebraic
closure of \(k\). Since \(C\) is of finite type over \(k\), the
scheme \(C\otimes_k\overline k\) is the disjoint union of finitely
many connected components \(C'_1,\ldots,C'_n\) of
\(G\otimes_k\overline k\), and it is enough to prove the assertion
for each \(C'_i\). We are thus reduced to the case in which \(k\)
is algebraically closed.

Let \(g\) be a closed point of \(C\), and let
\[
  u^0:G^0\longrightarrow H
\]
be the restriction of \(u\) to \(G^0\). Then
\[
  u'=r_{u(g)}\circ u^0\circ r_{g^{-1}},
\]
where \(r_g\) denotes right translation by \(g\). Thus, to prove
that \(u'\) is proper, it is enough to prove that \(u^0\) is
proper. By hypothesis \(u\) is locally quasi-finite, so the
nonempty fiber \(\ker u\) is discrete. We have seen that \(u^0\)
is of finite type, hence its fiber \(\ker u^0\) is finite
(cf.\ EGA~II, 6.2.2), and is therefore nonempty and proper.
\hyperref[prop:sga3-VIB-1.4]{Proposition~1.4} now shows that \(u^0\) is proper.

\medskip
\phantomsection\label{cor:sga3-VIB-1.4.2}%
\noindent\textbf{Corollary 1.4.2.---}
\emph{Let \(A\) be an Artinian local ring, and let \(u:G\to H\)
be a quasi-compact morphism between \(A\)-groups locally of finite
type. The following assertions are equivalent:%
\footnote{\label{note:sga3-VIB-15}Editor's note: the hypothesis
that \(G\) and \(H\) be locally of finite type may be removed,
because \textup{[Per76]}, 4.2.4, states that every quasi-compact
monomorphism \(u:G\to H\) between group schemes over a field \(k\)
is a closed immersion. See also \textup{[DG70]}, III.3.7.2(b), for
the case in which \(G\) and \(H\) are affine; in that case every
morphism \(G\to H\) is affine (EGA~II, 1.6.2(v)), hence
quasi-compact.}
\begin{enumerate}
\item[\textup{\SGAInlineRefTarget{cor:sga3-VIB-1.4.2:item:i}(i)}] \(u\) is a closed immersion;
\item[\textup{\SGAInlineRefTarget{cor:sga3-VIB-1.4.2:item:ii}(ii)}] \(u\) is a monomorphism;
\item[\textup{\SGAInlineRefTarget{cor:sga3-VIB-1.4.2:item:iii}(iii)}] \(\ker u\) is trivial, that is, isomorphic
  to the unit \(A\)-group.
\end{enumerate}
In particular, every group subscheme%
\footnote{\label{note:sga3-VIB-16}Editor's note: in this particular
case, see also Expos\'e~VI\(_A\), \SGARefLink{sga3:VIA:0.5.2}{0.5.2}, which holds without any
finiteness hypothesis.}
of \(H\) is closed.}

\begingroup
\renewcommand{\thefootnote}{\fnsymbol{footnote}}%
\footnote[1]{\label{note:sga3-VIB-1.4.2-unit-base}Translator's note:
the French source calls the trivial kernel the ``unit \(k\)-group''
both in the statement and in the proof. Here the base is the
Artinian local ring \(A\), and \(\ker u\) is an \(A\)-group; the
unit object required by the monomorphism criterion is therefore
\(\operatorname{Spec}A\), the unit \(A\)-group.}%
\endgroup{}
It is clear that \textup{\SGARefLink{cor:sga3-VIB-1.4.2:item:i}{(i)}} implies \textup{\SGARefLink{cor:sga3-VIB-1.4.2:item:ii}{(ii)}}. Considering
the functors represented by \(G\) and \(H\), one sees immediately
that \textup{\SGARefLink{cor:sga3-VIB-1.4.2:item:ii}{(ii)}} and \textup{\SGARefLink{cor:sga3-VIB-1.4.2:item:iii}{(iii)}} are equivalent. Finally, if
\(\ker u\) is the unit \(A\)-group, it is a proper nonempty fiber.
\hyperref[prop:sga3-VIB-1.4]{Proposition~1.4} therefore shows that \(u\) is a proper
monomorphism; it is of finite presentation because \(H\) is
locally Noetherian (EGA~IV\(_1\), 1.6.1), and hence is a closed
immersion (EGA~IV\(_3\), 8.11.5).

The final assertion follows because \(H\) is locally Noetherian,
so every immersion \(G\to H\) is quasi-compact (EGA~I, 6.6.4).

% END INCLUDED SGA3_English_current_buildable_source_20260823/source/inputs/VI/tex_reference_v2/components/29_expose_VIB_prop14_cor142_en.tex

% BEGIN INCLUDED SGA3_English_current_buildable_source_20260823/source/inputs/VI/tex_reference_v2/components/30_expose_VIB_counter143_through_cor161_en.tex
% SGA 3, Expose VI_B, Loop 1 English reconstruction.
% Scope: Counterexample 1.4.3 through Corollary 1.6.1 in full;
% Section 2 is excluded.
% Authority: Polo--Gille Exp6B-13oct24.pdf, component-local pages 7--8
% (combined-reader pages 348--349; original-margin printed pages
% 327--328).
% OCR and the Reinhold English Markdown are comparison witnesses only.

\medskip
\phantomsection\label{ctr:sga3-VIB-1.4.3}%
\noindent\textbf{Counter\SGARefLink{ctr:sga3-VIB-1.4.3}{example 1.4.3}.---}
\emph{Let \(k\) be a field of characteristic zero, let \(G\) be
the constant \(k\)-group \(\mathbb Z\), and let \(H\) be the
additive \(k\)-group \(\mathbb G_{a,k}\). Let \(u:G\to H\) be a
morphism of \(k\)-groups. If \(u\ne0\), then \(\ker u=0\), but
\(u\) is not a closed immersion (cf.\ \hyperref[ctr:sga3-VIB-1.2.2]{1.2.2}).}

We shall use below the following two results, which should have
appeared in Expos\'e~VI\(_A\).

\medskip
\phantomsection\label{lem:sga3-VIB-1.5}%
\noindent\textbf{Lemma 1.5.---}
\emph{Let \(k\) be a field and let \(G\) be a \(k\)-group locally
of finite type. Then:
\begin{enumerate}
\item[\textup{\SGAInlineRefTarget{lem:sga3-VIB-1.5:item:i}(i)}] all the irreducible components of \(G\) have
  the same dimension;
\item[\textup{\SGAInlineRefTarget{lem:sga3-VIB-1.5:item:ii}(ii)}] for every \(g\in G\),
  \[
    \dim_gG=\dim G.
  \]
\end{enumerate}}

\footnote{\label{note:sga3-VIB-17}Editor's note: the implication
\textup{(ii)} \(\Rightarrow\) \textup{(i)} is a general fact
(cf.\ EGA~0\(_{\mathrm{IV}}\), 14.1.6), while
\textup{(i)} \(\Rightarrow\) \textup{(ii)} follows from the fact
that if \(X\) is an irreducible scheme of finite type over a field,
then \(\dim X=\dim U\) for every nonempty open subset \(U\) of
\(X\). The essential point here is therefore assertion \textup{\SGARefLink{lem:sga3-VIB-1.5:item:i}{(i)}},
which was already established in an addition to
Expos\'e~VI\(_A\), \SGARefLink{cor:sga3-VIA-2.4.1}{2.4.1}. We have modified the statement and proof
accordingly.}
Assertion \textup{\SGARefLink{lem:sga3-VIB-1.5:item:i}{(i)}} has already been proved in
Expos\'e~VI\(_A\), \SGARefLink{cor:sga3-VIA-2.4.1}{2.4.1}, and \textup{\SGARefLink{lem:sga3-VIB-1.5:item:ii}{(ii)}} follows. Indeed, let
\(g\in G\), and let \(C\) be the connected component of \(G\)
containing \(g\). By definition
(EGA~0\(_{\mathrm{IV}}\), 14.1.2), \(\dim_gG\) is the infimum of
the integers \(\dim U\), where \(U\) ranges over the open
neighborhoods of \(g\). Hence
\[
  \dim_gG=\dim U_0
\]
for some \(U_0\), which we may suppose is contained in \(C\)
(because \(\dim V\leq\dim U\) when \(V\subset U\)). Since \(C\) is
irreducible and of finite type over \(k\)
(Expos\'e~VI\(_A\), \SGARefLink{cor:sga3-VIA-2.4.1}{2.4.1}),
\[
  \dim U_0
    =\dim C
    =\operatorname{trdeg}_k\kappa(\xi),
\]
where \(\xi\) is the generic point of \(C\)
(EGA~IV\(_2\), 5.2.1). Therefore
\[
  \dim_gG=\dim C=\dim G.
\]

\medskip
\phantomsection\label{prop:sga3-VIB-1.6}%
\noindent\textbf{Proposition 1.6.---}%
\footnote{\label{note:sga3-VIB-18}Editor's note: in the statement,
we have replaced ``along the identity section'' by ``at every point
of the identity section.'' At the end of the proof, we have also
made explicit the cited results from EGA~IV\(_4\).}
\emph{Let \(S\) be a scheme of characteristic zero, and let \(G\)
be an \(S\)-group scheme locally of finite presentation over \(S\)
at every point of the identity section \(\epsilon(S)\). The scheme
\(G\) is smooth over \(S\) at every point of the identity section
if and only if the \(\mathcal O_S\)-module}
\[
  \omega_{G/S}
    =\epsilon^*\bigl(\Omega^1_{G/S}\bigr),
\]
\emph{called the conormal module of the identity section of \(G\),
is locally free.}

Recall that a scheme \(S\) is said to be of characteristic zero if
the field \(\kappa(x)\) has characteristic zero for every closed
point \(x\in S\).

Recall also (\SGARefLink{sga3:ii:subsection:4-11}{Expos\'e~II, 4.11}) that, if
\(\pi:G\to S\) is the structural morphism, then
\[
  \Omega^1_{G/S}=\pi^*(\omega_{G/S}).
\]
Thus it is equivalent to say that the \(\mathcal O_S\)-module
\(\omega_{G/S}\) is locally free or that the
\(\mathcal O_G\)-module \(\Omega^1_{G/S}\) is locally free.

If there is an open neighborhood \(U\) of \(\epsilon(S)\) that is
smooth over \(S\), then by EGA~IV\(_4\), 17.2.3,
\(\Omega^1_{U/S}\) is locally free of finite type, and so is
\[
  \omega_{G/S}=\epsilon^*(\Omega^1_{U/S}).
\]

Conversely, if \(\omega_{G/S}\) is locally free, then so is
\[
  \Omega^1_{G/S}=\pi^*(\omega_{G/S}).
\]
Since \(S\) has characteristic zero, the Jacobian criterion
(EGA~IV\(_4\), 16.12.2) shows that \(G\) is differentially smooth
over \(S\). It follows from EGA~IV\(_4\), 17.12.5, that \(G\) is
smooth over \(S\) at every point of the identity section.

\medskip
\phantomsection\label{cor:sga3-VIB-1.6.1}%
\noindent\textbf{Corollary 1.6.1 (Cartier).---}
\emph{Let \(k\) be a field of characteristic zero. Every
\(k\)-group locally of finite type over \(k\) is smooth over
\(k\).}

Indeed, the \(k\)-module \(\omega_{G/k}\) is then plainly free.
By \SGARefLink{prop:sga3-VIB-1.6}{1.6}, \(G\) is smooth over \(k\) at the identity point \(e\), and
hence smooth over \(k\) by \hyperref[cor:sga3-VIB-1.3.1]{1.3.1}.

% END INCLUDED SGA3_English_current_buildable_source_20260823/source/inputs/VI/tex_reference_v2/components/30_expose_VIB_counter143_through_cor161_en.tex

% BEGIN INCLUDED SGA3_English_current_buildable_source_20260823/source/inputs/VI/tex_reference_v2/components/31_expose_VIB_section2_opening_loop2_r1_en.tex
% SGA 3, Expose VI_B, Loop 1 English reconstruction.
% Scope: Section 2 heading and paragraph 2.0 through Corollary 2.4
% in full; Proposition 2.5 is excluded.
% Authority: Polo--Gille Exp6B-13oct24.pdf, component-local pages
% 8--11 (combined-reader pages 349--352; original-margin printed
% pages 329--333).
% OCR and the Reinhold English Markdown are comparison witnesses only.
% The Lemma 2.1.1 cartesian square is an exact 180-dpi crop of the
% controlling PDF: local p. 9, x=650,y=1480,w=210,h=180 on a
% 1488x2105 render.

\section*{2. ``Open properties'' of groups and group morphisms
  locally of finite presentation}
\phantomsection\label{sec:sga3-VIB-2}
\addcontentsline{toc}{section}{2. ``Open properties'' of groups and
  group morphisms locally of finite presentation}

\medskip
\phantomsection\label{par:sga3-VIB-2.0}%
\noindent\textbf{2.0.---}
In all that follows, \(S\) will denote an arbitrary scheme; an
\(S\)-group scheme will be called an \(S\)-group. Given an
\(S\)-group \(G\), we shall denote by \(\epsilon\) the identity
section, by
\[
  c:G\longrightarrow G
\]
the inversion morphism, and by
\[
  \mu:G\times_S G\longrightarrow G
\]
the multiplication morphism. For every \(S\)-scheme \(X\), we
shall denote by \(\pi\), or by \(\pi_X\), the structural morphism
\[
  X\longrightarrow S.
\]

Given a property \(\mathscr P(u)\) of a morphism of \(S\)-schemes
\[
  u:X\longrightarrow Y,
\]
we shall say that \(\mathscr P(u)\) is \emph{stable under base
change} if, whenever \(u\) satisfies \(\mathscr P(u)\), the same is
true of the morphism \(u_{(Y')}\), for every \(S\)-morphism
\(Y'\to Y\). We say that \(\mathscr P(u)\) is \emph{local in
nature for the topology \(\mathcal T\)}
(cf.\ Expos\'e~IV, \S\S~4 and 6) if \(\mathscr P(u)\) satisfies the
following two conditions:
\begin{enumerate}
\item[\textup{a)}] \(\mathscr P(u)\) is stable under base change;
\item[\textup{b)}] whenever there is a family of \(S\)-morphisms
  \(Y_i\to Y\), covering for the topology \(\mathcal T\), such
  that each morphism \(u_{(Y_i)}\) satisfies \(\mathscr P(u)\),
  then \(u\) satisfies \(\mathscr P(u)\).
\end{enumerate}

\medskip
\phantomsection\label{prop:sga3-VIB-2.1}%
\noindent\textbf{Proposition 2.1.---}
\emph{Let \(\mathscr P(u)\) be a property of a morphism of
\(S\)-schemes, local in nature for the \(\mathrm{fpqc}\) topology.
Let \(u:G\to H\) be a morphism of \(S\)-groups. Suppose that \(G\)
is flat and universally open over \(S\).}

\emph{Let \(W\) be the largest open subset of \(H\) over which \(u\)
satisfies the property \(\mathscr P(u)\), and let
\(V=u^{-1}(W)\). Then
\[
  U=\pi_G(V)
\]
is open in \(S\), and \(V\) is an open group subscheme of
\(G|_U\).}

The existence of a largest open subset \(W\) of \(H\) over which
\(u\) satisfies \(\mathscr P(u)\) follows from the fact that
\(\mathscr P(u)\) is local in nature for the Zariski topology.
Since \(\pi_G\) is universally open, \(\pi_G(V)\) is an open
subset \(U\) of \(S\). It is enough to show that \(V\) is a group
subscheme of \(G|_U\). We may therefore suppose that \(U=S\).

Set
\[
  G'=G\times_SV,\qquad
  H'=H\times_SV,\qquad
  V'=V\times_SV,\qquad
  W'=W\times_SV,
\]
and \(u'=u_{(V)}\). Let \(W'_1\) be the largest open subset of
\(H'\) over which \(u'\) satisfies \(\mathscr P(u)\). Since \(V\)
is flat and universally open over \(S\), the same is true of
\(H'\) over \(H\), and \hyperref[lem:sga3-VIB-2.1.1]{Lemma~2.1.1} below shows that
\[
  W'_1=W'.
\]
Now consider the automorphism \(a\) (respectively, \(b\)) of the
\(V\)-scheme \(G'\) (respectively, \(H'\)), namely right
translation by the inverse of the diagonal section \(\delta\)
(respectively, by the inverse of \(u(\delta)\)), defined by
\[
  a(g,v)=(g\cdot v^{-1},v),
  \qquad
  \bigl(\text{respectively, }b(h,v)
    =(h\cdot u(v^{-1}),v)\bigr),
\]
for every morphism \(T\to S\), every \(g\in G(T)\),
\(v\in V(T)\), and \(h\in H(T)\).%
\begingroup
\renewcommand{\thefootnote}{\fnsymbol{footnote}}%
\footnote[1]{\label{note:sga3-VIB-2.1-b-formula}Translator's note:
the French source omits \(b\) before \((h,v)\) in the second
displayed definition. The preceding sentence defines the
automorphism \(b\), and the identity immediately below requires
the displayed map to be \(b(h,v)\); the symbol is supplied here.}%
\endgroup{}
It is clear that
\[
  u'\circ a=b\circ u',
\]
which shows that \(W'\) is stable under \(b\), hence that \(V'\)
is stable under \(a\), so that \(V\) is a group subscheme of \(G\).

\medskip
\phantomsection\label{lem:sga3-VIB-2.1.1}%
\noindent\textbf{Lemma 2.1.1.---}
\emph{Let \(\mathscr P(u)\) be a property of an \(S\)-morphism
\(u\), local in nature for the \(\mathrm{fpqc}\) topology.
Consider a cartesian square of morphisms of \(S\)-schemes:}
\begin{center}
  \phantomsection\label{diag:sga3-VIB-2.1.1-cartesian-square}%
  \begin{tikzcd}[row sep=large,column sep=large]
    X' \arrow[r,"f'"] \arrow[d] &
    Y' \arrow[d,"g"'] \\
    X \arrow[r,"f"] &
    Y
  \end{tikzcd}
\end{center}
\emph{where \(g\) is flat and open. Let \(W\) (respectively,
\(W'_1\)) be the largest open subset of \(Y\) (respectively,
\(Y'\)) over which \(f\) (respectively,
\(f'=f_{(Y')}\)) satisfies \(\mathscr P(u)\). Then}
\[
  \emph{\(W'_1=W\times_Y Y'\).}
\]

Set \(W'=W\times_Y Y'\). Since \(\mathscr P(u)\) is stable under
base change, we have
\[
  W'\subset W'_1.
\]
Since \(g\) is open,
\[
  W_1=g(W'_1)
\]
is an open subset of \(Y\). Set
\[
  V_1=f^{-1}(W_1),
  \qquad
  V'_1=V_1\times_{W_1}W'_1;
\]
it is clear that \(V'_1=f'^{-1}(W'_1)\). Since \(g\) is flat and
open, the morphism
\[
  W'_1\longrightarrow W_1=g(W'_1)
\]
induced by \(g\) is faithfully flat and open, hence is a covering
for the \(\mathrm{fpqc}\) topology
(cf.\ \SGARefLink{sga3:iv:proposition:6:3:1:item:iv:01}{Expos\'e~IV, 6.3.1(iv)}). Since the morphism
\[
  V'_1\longrightarrow W'_1
\]
induced by \(f'\) satisfies \(\mathscr P(u)\), the same is true of
the morphism \(V_1\to W_1\) induced by \(f\). Thus
\[
  W_1\subset W,
  \qquad
  W'_1\subset g^{-1}(W_1)
    \subset g^{-1}(W)=W',
\]
and therefore \(W'=W'_1\).

\medskip
\phantomsection\label{rem:sga3-VIB-2.1.2}%
\noindent\textbf{Remark 2.1.2.---}
\emph{A great many properties of a morphism are local in nature
for the \(\mathrm{fpqc}\) topology; among them are the properties
of being flat, (universally) open, (locally) of finite type, of
finite presentation, quasi-finite
(cf.\ EGA~IV\(_2\), 2.5.1, 2.6.1, and 2.7.1), smooth, \'etale,
and unramified (EGA~IV\(_4\), 17.7.3).
}

\emph{The proof of \hyperref[prop:sga3-VIB-2.1]{2.1} in fact uses only base changes by flat morphisms.
The proposition therefore applies to a property satisfying
condition \textup{b)} of \hyperref[par:sga3-VIB-2.0]{2.0} relative to the \(\mathrm{fpqc}\)
topology and stable under base change by flat morphisms (for
example, the property of being quasi-compact and dominant).
}

\emph{Of course, one can state an analogous proposition concerning
properties local in nature for a topology \(\mathcal T\) finer
than the Zariski topology, the condition to be verified by \(G\)
then being that \(\pi_G\) be universally open and covering for
the topology \(\mathcal T\).
}

\emph{In particular, if \(G\) is flat and locally of finite presentation
over \(S\), there is an analogous statement for properties stable
under base change by flat morphisms locally of finite presentation
and satisfying condition \textup{b)} of \SGARefLink{par:sga3-VIB-2.0}{2.0} relative to the
\(\mathrm{fpqc}\) topology (for example, the properties of being
regular, reduced, Cohen--Macaulay, and so on; cf.\
EGA~IV\(_2\), 6.8).}

\medskip
\phantomsection\label{prop:sga3-VIB-2.2}%
\noindent\textbf{Proposition 2.2.---}
\emph{Let \(G\) and \(H\) be two \(S\)-groups, and let
\(u:G\to H\) be a morphism of \(S\)-groups. Then:}%
\footnote{\label{note:sga3-VIB-19}Editor's note: in case
\textup{\SGARefLink{prop:sga3-VIB-2.2:item:i}{(i)}}, the open subset \(V\) consists of all the points of
\(G\) at which \(u\) is smooth, respectively \'etale,
respectively of finite presentation and flat; cf.\ editor's
\hyperref[note:sga3-VIB-6]{note~6}. By contrast, in case \textup{\SGARefLink{prop:sga3-VIB-2.2:item:ii}{(ii)}}, \(V\) is the largest
open subset contained in the set \(E\) of points of \(G\) at
which \(u\) is universally open, but \(E\) is not necessarily
open, as the following example shows
(EGA~IV\(_3\), 14.1.3(i)). Let \(k\) be a field, let
\(H=S=\operatorname{Spec}A\), where \(A=k[x]\), and let \(G\) be
the \(S\)-group \(\operatorname{Spec}A[y]/(xy)\). Then \(E\) is
the identity section of \(G\), which is not open.}
\begin{enumerate}
\item[\textup{\SGAInlineRefTarget{prop:sga3-VIB-2.2:item:i}(i)}] \emph{Suppose that \(G\) or \(H\) is flat over
  \(S\), and that \(G\) or \(H\) is locally of finite presentation
  over \(S\). Let \(V\) be the largest open subset of \(G\) such
  that the restriction of \(u\) to \(V\) is flat and locally of
  finite presentation (respectively, smooth, or \'etale). Then
  \(U=\pi_V(V)\) is open in \(S\), and \(V\) is an open group
  subscheme of \(G|_U\).}
\item[\textup{\SGAInlineRefTarget{prop:sga3-VIB-2.2:item:ii}(ii)}] \emph{Suppose that \(G\) or \(H\) is
  universally open over \(S\), and let \(V\) be the largest open
  subset of \(G\) such that the restriction of \(u\) to \(V\) is
  universally open. Then \(U=\pi_V(V)\) is open in \(S\), and
  \(V\) is an open group subscheme of \(G|_U\).}
\end{enumerate}

We first prove \textup{\SGARefLink{prop:sga3-VIB-2.2:item:i}{(i)}}. We show that the restriction
\(\pi_V\) of \(\pi_G\) to \(V\) is flat and locally of finite
presentation:
\begin{enumerate}
\item[\textup{a)}] if \(\pi_G\) is flat (respectively, locally of
  finite presentation), then so is \(\pi_V\);
\item[\textup{b)}] if \(\pi_H\) is flat (respectively, locally of
  finite presentation), then so is \(\pi_V\), as the composite of
  the restriction of \(u\) to \(V\) with \(\pi_H\).
\end{enumerate}
Thus, in the four cases contemplated in the statement, \(\pi_V\)
is flat and locally of finite presentation, hence universally open
(EGA~IV\(_2\), 2.4.6). Set \(U=\pi_V(V)\); then \(U\) is open in
\(S\). It is therefore enough to show that \(V\) is an open group
subscheme of \(G|_U\), and we may suppose that \(U=S\).

Set
\[
  G'=G\times_SV,\qquad
  H'=H\times_SV,\qquad
  V'=V\times_SV,
  \qquad\text{and}\qquad
  u'=u_{(V)}.
\]
Since \(V\) is flat and locally of finite presentation over \(S\),
the same is true of \(H'\) over \(H\). By EGA~IV\(_4\), 17.7.4,
\(V'\) is then the largest open subset of \(G'\) such that the
restriction of \(u'\) to \(V'\) is flat and locally of finite
presentation (respectively, smooth, or \'etale). With the
automorphisms \(a\) and \(b\) defined as in the proof of \SGARefLink{prop:sga3-VIB-2.1}{2.1}, it
is clear that \(V'\) is stable under \(a\), hence that \(V\) is a
group subscheme of \(G\).

We now prove \textup{\SGARefLink{prop:sga3-VIB-2.2:item:ii}{(ii)}}. The restriction \(\pi_V\) of \(\pi_G\)
to \(V\) is universally open, either because \(\pi_G\) is
universally open, or because it is the composite of the restriction
of \(u\) to \(V\) with \(\pi_H\), when \(\pi_H\) is universally
open. Set \(U=\pi_V(V)\); then \(U\) is open in \(S\). It is enough
to show that \(V\) is an open group subscheme of \(G|_U\), and we
may suppose that \(U=S\).

Set, as before,
\[
  G'=G\times_SV,\qquad
  H'=H\times_SV,\qquad
  V'=V\times_SV,
  \qquad\text{and}\qquad
  u'=u_{(V)}.
\]
Then \(\pi_V:V\to S\) is surjective and universally open; the same
is therefore true of \(G'\to G\). Consequently, \(V'\) is the
largest open subset of \(G'\) such that the restriction of \(u'\)
to \(V'\) is universally open, by EGA~IV\(_3\), 14.3.4(i),(ii).
With the automorphisms \(a\) and \(b\) defined as before, it is
clear that \(V'\) is stable under \(a\), hence that \(V\) is a
group subscheme of \(G\).

\medskip
\phantomsection\label{cor:sga3-VIB-2.3}%
\noindent\textbf{Corollary 2.3.---}
\emph{Let \(G\) be an \(S\)-group, and let \(V\) be the largest
open subset of \(G\) that is flat and locally of finite
presentation (respectively, smooth, \'etale, or universally open)
over \(S\). Then \(U=\pi_V(V)\) is open in \(S\), and \(V\) is an
open group subscheme of \(G|_U\).}

It is enough to apply \hyperref[prop:sga3-VIB-2.2]{2.2} when \(H\) is the identity \(S\)-group
and \(u\) is the unique morphism of \(S\)-groups \(G\to H\), since
\(\pi_H\) is then an isomorphism and
\[
  \pi_G=\pi_H\circ u.
\]

\medskip
\phantomsection\label{cor:sga3-VIB-2.4}%
\noindent\textbf{Corollary 2.4.---}
\emph{Let \(G\) be an \(S\)-group. If there is a neighborhood \(X\)
of the identity section having one (or more) of the following
properties:}
\begin{center}
  \emph{\(X\) is flat and locally of finite presentation
  (respectively, smooth, \'etale, or universally open) over \(S\),}
\end{center}
\emph{then there is an open group subscheme of \(G\) having the
same properties.}

It is enough to apply \hyperref[cor:sga3-VIB-2.3]{2.3}, observing that here, with the notation
of \hyperref[prop:sga3-VIB-2.2]{2.2}, we have
\[
  \epsilon(S)\subset V,
\]
and hence \(U=S\).

% END INCLUDED SGA3_English_current_buildable_source_20260823/source/inputs/VI/tex_reference_v2/components/31_expose_VIB_section2_opening_loop2_r1_en.tex

% BEGIN INCLUDED SGA3_English_current_buildable_source_20260823/source/inputs/VI/tex_reference_v2/components/32_expose_VIB_prop25_through_lemma271_en.tex
% SGA 3, Expose VI_B, Loop 1 English reconstruction.
% Scope: Proposition 2.5, Corollary 2.6, Proposition 2.7, and
% Lemma 2.7.1 in full; Corollary 2.8 is excluded.
% Authority: Polo--Gille Exp6B-13oct24.pdf, component-local pages
% 11--14 (combined-reader pages 352--355; original-margin printed
% pages 333--336, stopping before printed page 337).
% OCR and the Reinhold English Markdown are comparison witnesses only.

\medskip
\phantomsection\label{prop:sga3-VIB-2.5}%
\noindent\textbf{Proposition 2.5.---}%
\footnote{\label{note:sga3-VIB-20}Editor's note: we have expanded
the statement and proof, and in \textup{\SGARefLink{concl:sga3-VIB-2.5-i}{(i)}} we have weakened the
hypothesis on \(H\), replacing ``of finite presentation'' by
``of finite type.''}
\emph{Let \(u:G\to H\) be a morphism of \(S\)-groups.}

\begin{enumerate}
\item[\textup{(i)}]\phantomsection\label{concl:sga3-VIB-2.5-i} \emph{Suppose that \(G\) (respectively, \(H\))
  is of finite presentation and flat (respectively, of finite type)
  over \(S\) at the points of its identity section. Then the sets}
  \[
    U_{\mathrm{plat}}
      \supset U_{\mathrm{lisse}}
      \supset U_{\mathrm{\acute{e}t.}},
  \]
  \emph{formed by the points \(s\in S\) such that \(u_s\) is flat
  (respectively, smooth, or \'etale), are open in \(S\).}

  \emph{If, moreover, \(G\) (respectively, \(H\)) is locally of
  finite presentation and flat (respectively, locally of finite
  type) over \(S\), then the set \(V_{\mathrm{plat}}\)
  (respectively, \(V_{\mathrm{lisse}}\) or
  \(V_{\mathrm{\acute{e}t.}}\)) of points of \(G\) at which \(u\)
  is flat (respectively, smooth, or \'etale) is the inverse image
  under \(\pi_G\) of \(U_{\mathrm{plat}}\) (respectively,
  \(U_{\mathrm{lisse}}\) or \(U_{\mathrm{\acute{e}t.}}\)).}

\item[\textup{(ii)}]\phantomsection\label{concl:sga3-VIB-2.5-ii} \emph{Suppose that, for every \(s\in S\),
  \(G_s\) is locally of finite type over \(\kappa(s)\) (a condition
  satisfied if \(G\) is of finite type over \(S\) at the points of
  the identity section \textup{(\SGARefLink{lem:sga3-VIB-1.3.1.1}{1.3.1.1})}), and that \(u\) is
  locally of finite type (respectively, locally of finite
  presentation) at the points of the identity section of \(G\).
  Then the sets}
  \[
    U_{\mathrm{l.q.f.}}\supset U_{\mathrm{n.r.}},
  \]
  \emph{formed by the \(s\in S\) such that \(u_s\) is locally
  quasi-finite (respectively, unramified), are open in \(S\).}

  \emph{If, moreover, \(u\) is locally of finite type
  (respectively, locally of finite presentation), then the set
  \(V_{\mathrm{l.q.f.}}\) (respectively,
  \(V_{\mathrm{n.r.}}\)) of points of \(G\) at which \(u\) is
  quasi-finite (respectively, unramified) is the inverse image under
  \(\pi_G\) of \(U_{\mathrm{l.q.f.}}\) (respectively,
  \(U_{\mathrm{n.r.}}\)).}
\end{enumerate}

We prove \textup{\SGARefLink{concl:sga3-VIB-2.5-i}{(i)}}. First note (\SGARefLink{lem:sga3-VIB-1.3.1.1}{1.3.1.1}) that, for every
\(s\in S\), \(G_s\) is locally of finite type over \(\kappa(s)\).
Let \(Y\) be the open subset of \(H\) consisting of the points at
which \(\pi_H\) is of finite type, and let \(X_1\) be the open
subset of \(u^{-1}(Y)\) consisting of the points at which
\(\pi_G\) is of finite presentation. By EGA~IV\(_3\), 11.3.1, the
set \(X\) of points of \(X_1\) at which \(\pi_G\) is of finite
presentation and flat is open in \(X_1\), hence in \(G\).

Let
\[
  u_X:X\longrightarrow Y
\]
be the morphism induced by \(u\). Since \(Y\) (respectively, \(X\))
contains the identity section of \(H\) (respectively, \(G\)), the
restriction \(\pi_X\) of \(\pi_G\) to \(X\) is surjective, and the
morphism
\[
  S\longrightarrow X
\]
induced by \(\epsilon_G\) is an \(S\)-section of \(X\).

Consider the sets
\[
  W_{\mathrm{plat}}
    \supset W_{\mathrm{lisse}}
    \supset W_{\mathrm{\acute{e}t.}},
\]
consisting of the points of \(X\) at which \(u_X\) is flat
(respectively, smooth, or \'etale). It is clear that
\(W_{\mathrm{lisse}}\) and \(W_{\mathrm{\acute{e}t.}}\) are open
in \(X\); cf.\ Editor's \hyperref[note:sga3-VIB-6]{note~6}. On the other hand, since \(\pi_Y\)
is locally of finite type and
\[
  \pi_X=\pi_Y\circ u_X
\]
is locally of finite presentation, EGA~IV\(_1\), 1.4.3(v), shows
that \(u_X\) is locally of finite presentation. Consequently, the
flatness locus \(W_X\) of \(u_X\) is open in \(X\)
(EGA~IV\(_3\), 11.3.1).

Let \(x\in X\) and set \(s=\pi_X(x)\). Then, by
EGA~IV\(_2\), 11.3.10,%
\begingroup
\renewcommand{\thefootnote}{\fnsymbol{footnote}}%
\footnote[1]{\label{note:sga3-VIB-2.5-EGA-volume}Translator's note:
the source citation is retained. Section 11 occurs in EGA~IV\(_3\),
not EGA~IV\(_2\); the cited result is EGA~IV\(_3\), 11.3.10.}%
\endgroup{}
combined with EGA~IV\(_4\), 17.5.1 (respectively, 17.6.1),
\(x\) belongs to \(W_{\mathrm{plat}}\) (respectively,
\(W_{\mathrm{lisse}}\) or \(W_{\mathrm{\acute{e}t.}}\)) if and
only if \(u_s\) is flat (respectively, smooth, or \'etale) at the
point \(x\), or, equivalently by \hyperref[prop:sga3-VIB-1.3]{1.3}, if and only if \(u_s\) is
flat (respectively, smooth, or \'etale).

Consequently, for
\[
  \text{``\(\flat=\) flat, smooth, or \'etale,''}
\]
we have
\[
  U_\flat=\epsilon_G^{-1}(W_\flat),
  \qquad
  W_\flat=\pi_X^{-1}(U_\flat).
\]
Since \(W_\flat\) is open in \(X\), hence in \(G\), the first
equality shows that \(U_\flat\) is open in \(S\).

The second assertion of \textup{\SGARefLink{concl:sga3-VIB-2.5-i}{(i)}} follows from what we have just
seen, for then
\[
  Y=H,\qquad X=G,\qquad
  V_\flat=W_\flat=\pi_G^{-1}(U_\flat).
\]

We prove assertion \textup{\SGARefLink{concl:sga3-VIB-2.5-ii}{(ii)}}. Let
\[
  V_{\mathrm{l.t.f.}}\supset V_{\mathrm{l.q.f.}}
\]
be the open subsets of \(G\) consisting of the points at which
\(u\) is of finite type (respectively, quasi-finite). Set
\[
  X=V_{\mathrm{l.t.f.}},
\]
and let \(u_X\) denote the restriction of \(u\) to \(X\). By
hypothesis, \(X\) contains \(\epsilon(S)\). Let \(x\in X\) and set
\(s=\pi_X(x)\).

If \(u\) is quasi-finite at \(x\), then, by base change, \(u_s\) is
quasi-finite at \(x\); and since \(G_s\) is assumed to be locally
of finite type, \(u_s\) is locally quasi-finite by \SGARefLink{cor:sga3-VIB-1.3.1}{1.3.1}.

Conversely, if \(u_s\) is locally quasi-finite, then
\[
  u_X^{-1}\bigl(u_X(x)\bigr)
    \subset u_s^{-1}\bigl(u_s(x)\bigr)
\]
is a finite set. Since
\[
  u_X:X\longrightarrow H
\]
is locally of finite type, Chevalley's semicontinuity theorem
(EGA~IV\(_3\), 13.1.3) shows that there is an open neighborhood
\(W\) of \(x\) in \(X\) such that the fiber
\[
  u_X^{-1}\bigl(u_X(x')\bigr)
\]
is discrete for every \(x'\in W\). Hence \(u_X\) is quasi-finite
at \(x\), by EGA~II, 6.2.2 (and EGA~III\(_2\),
Err\(_{\mathrm{III}}\)~20).

Consequently,
\[
  U_{\mathrm{l.q.f.}}
    =\epsilon_G^{-1}\bigl(V_{\mathrm{l.q.f.}}\bigr),
  \qquad
  V_{\mathrm{l.q.f.}}
    =\pi_X^{-1}\bigl(U_{\mathrm{l.q.f.}}\bigr),
\]
and the first equality shows that \(U_{\mathrm{l.q.f.}}\) is open
in \(S\).

If, moreover, \(u\) is locally of finite type over \(S\), then
\(X=G\), and hence \(V_{\mathrm{l.q.f.}}\) is the inverse image
under \(\pi_G\) of \(U_{\mathrm{l.q.f.}}\).%
\begingroup
\renewcommand{\thefootnote}{\fnsymbol{footnote}}%
\footnote[1]{\label{note:sga3-VIB-2.5-u-loci}Translator's note:
in this sentence and in the analogous sentence for
\(V_{\mathrm{n.r.}}\) below, the French source prints \(G\) in
place of \(u\). The statement assumes that \(u\) is locally of
finite type (respectively, locally of finite presentation), and
\(X\) and \(Z\) are precisely the corresponding loci for \(u\);
the required \(u\) is used here.}%
\endgroup{}
This proves the first half of \textup{\SGARefLink{concl:sga3-VIB-2.5-ii}{(ii)}}.

Now consider the open subsets
\[
  V_{\mathrm{l.p.f.}}\supset V_{\mathrm{n.r.}},
\]
consisting of the points at which \(u\) is of finite presentation
(respectively, unramified), and suppose that
\[
  Z=V_{\mathrm{l.p.f.}}
\]
contains the identity section of \(G\). Let \(z\in Z\); set
\[
  s=\pi_Z(z),\qquad h=u_Z(z).
\]

If \(u\) is unramified at \(z\), then, by base change, \(u_s\) is
unramified at \(z\); and since \(G_s\) is assumed to be locally of
finite type, \(u_s\) is unramified by \SGARefLink{cor:sga3-VIB-1.3.1}{1.3.1}.

Conversely, if \(u_s\) is unramified at the point \(z\), then the
fiber
\[
  u_s^{-1}(h)=u^{-1}(h)
\]
is unramified over \(\kappa(h)\) at the point \(z\). Since \(Z\)
is an open subset of \(G\), the fiber
\[
  u_Z^{-1}(h)=u_s^{-1}(h)\cap Z
\]
is unramified over \(\kappa(h)\) at the point \(z\). Since \(u_Z\)
is locally of finite presentation, EGA~IV\(_4\), 17.4.1, implies
that \(u_Z\) is unramified at the point \(z\).

Thus
\[
  U_{\mathrm{n.r.}}
    =\epsilon_G^{-1}\bigl(V_{\mathrm{n.r.}}\bigr),
  \qquad
  V_{\mathrm{n.r.}}
    =\pi_Z^{-1}\bigl(U_{\mathrm{n.r.}}\bigr),
\]
and the first equality shows that \(U_{\mathrm{n.r.}}\) is open in
\(S\).

If, moreover, \(u\) is locally of finite presentation over \(S\),
then \(Z=G\), and hence \(V_{\mathrm{n.r.}}\) is the inverse image
under \(\pi_G\) of \(U_{\mathrm{n.r.}}\). This completes the proof
of assertion \textup{\SGARefLink{concl:sga3-VIB-2.5-ii}{(ii)}}.

\medskip
\phantomsection\label{cor:sga3-VIB-2.6}%
\noindent\textbf{Corollary 2.6.---}
\emph{Let \(u:G\to H\) be a morphism of \(S\)-groups that is
radicial (as is the case if \(u\) is a monomorphism;
EGA~I, 3.5.4). Suppose that \(G\) (respectively, \(H\)) is locally
of finite presentation and flat (respectively, locally of finite
type) over \(S\). Then the set \(U\) of \(s\in S\) such that
\(u_s\) is an open immersion is open in \(S\), and the restriction
of \(u\) over \(U\) is an open immersion.}

By \hyperref[concl:sga3-VIB-2.5-i]{2.5}\textup{\SGARefLink{concl:sga3-VIB-2.5-i}{(i)}}, the set \(U'\) of points \(s\in S\) such that
\(u_s\) is \'etale is open in \(S\). Since \(u\) is radicial, so
is \(u_s\), for every \(s\in S\); hence EGA~IV\(_4\), 17.9.1,
gives
\[
  U=U',
\]
which shows that \(U\) is open. Finally, by \hyperref[concl:sga3-VIB-2.5-i]{2.5}\textup{\SGARefLink{concl:sga3-VIB-2.5-i}{(i)}}, the
restriction of \(u\) over \(U\) is \'etale; since \(u\) is
radicial, this restriction is an open immersion
(EGA~IV\(_4\), 17.9.1).

\medskip
\phantomsection\label{prop:sga3-VIB-2.7}%
\noindent\textbf{Proposition 2.7.---}
\emph{Let \(G\) be an \(S\)-group. The following conditions are
equivalent:}
\begin{enumerate}
\item[\textup{\SGAInlineRefTarget{prop:sga3-VIB-2.7:item:i}(i)}] \emph{\(G\) is unramified over \(S\) at the
  points of the identity section;}
\item[\textup{\SGAInlineRefTarget{prop:sga3-VIB-2.7:item:ii}(ii)}] \emph{the identity section is an open
  immersion;}
\item[\textup{\SGAInlineRefTarget{prop:sga3-VIB-2.7:item:iii}(iii)}] \emph{\(G\) is of finite presentation over
  \(S\) at the points of the identity section, and, for every
  \(s\in S\), \(G_s\) is unramified over \(\kappa(s)\).}
\end{enumerate}
\emph{If, moreover, \(G\) is locally of finite presentation over
\(S\), then each of the three preceding conditions is equivalent
to the following:}
\begin{enumerate}
\item[\textup{(iv)}] \emph{\(G\) is unramified over \(S\).}
\end{enumerate}

The equivalence of assertions \textup{\SGARefLink{prop:sga3-VIB-2.7:item:i}{(i)}} and \textup{\SGARefLink{prop:sga3-VIB-2.7:item:ii}{(ii)}}
follows from the more general \hyperref[lem:sga3-VIB-2.7.1]{Lemma~2.7.1} below. Note (\SGARefLink{lem:sga3-VIB-1.3.1.1}{1.3.1.1})
that either condition \textup{\SGARefLink{prop:sga3-VIB-2.7:item:i}{(i)}} or condition \textup{\SGARefLink{prop:sga3-VIB-2.7:item:iii}{(iii)}}
implies that, for every \(s\in S\), \(G_s\) is locally of finite
type over \(\kappa(s)\). Thus, by EGA~IV\(_4\), 17.4.1, assertion
\textup{(i)} is equivalent to the assertion that, for every
\(s\in S\), \(G_s\) is unramified over \(\kappa(s)\) at the point
\(e_s\), the identity of \(G_s\), or, equivalently by \hyperref[cor:sga3-VIB-1.3.1]{1.3.1}, to
the assertion that \(G_s\) is unramified over \(\kappa(s)\).
Consequently, assertions \textup{\SGARefLink{cor:sga3-VIB-1.3.1:item:i}{(i)}} and \textup{\SGARefLink{prop:sga3-VIB-2.7:item:iii}{(iii)}} are
equivalent. Finally, if \(G\) is locally of finite presentation
over \(S\), assertions \textup{(iii)} and \textup{(iv)} are
equivalent (EGA~IV\(_4\), 17.4.1).

\medskip
\phantomsection\label{lem:sga3-VIB-2.7.1}%
\noindent\textbf{Lemma 2.7.1.---}
\emph{Let \(G\) be an \(S\)-scheme equipped with a section
\(\epsilon\). In order that \(G\) be unramified over \(S\) at the
points of this section, it is necessary and sufficient that
\(\epsilon\) be an open immersion.}

The condition is necessary by EGA~IV\(_4\), 17.4.1,
\textup{a)}\(\Rightarrow\)\textup{b\(^{\prime\prime}\)}).
Conversely, if \(\epsilon\) is an open immersion, then the
restriction to \(\epsilon(S)\) of the structural morphism
\[
  G\longrightarrow S
\]
is an isomorphism; hence \(G\) is unramified over \(S\) at the
points of \(\epsilon(S)\).

% END INCLUDED SGA3_English_current_buildable_source_20260823/source/inputs/VI/tex_reference_v2/components/32_expose_VIB_prop25_through_lemma271_en.tex

% BEGIN INCLUDED SGA3_English_current_buildable_source_20260823/source/inputs/VI/tex_reference_v2/components/33_expose_VIB_cor28_through_cor212_en.tex
% SGA 3, Expose VI_B, Loop 1 English reconstruction.
% Scope: Corollary 2.8 through Corollary 2.12 in full;
% Section 3 and paragraph 3.0 are excluded.
% Authority: Polo--Gille Exp6B-13oct24.pdf, component-local pages
% 14--16 (combined-reader pages 355--357; original-margin printed
% pages 337--340, stopping before Section 3 on local page 16).
% OCR and the Reinhold English Markdown are comparison witnesses only.

\medskip
\phantomsection\label{cor:sga3-VIB-2.8}%
\noindent\textbf{Corollary 2.8.---}%
\footnote{\label{note:sga3-VIB-21}Editor's note: we have changed
the presentation by separating the assertions concerning the
``quasi-finite'' case from those concerning the ``unramified''
case.}
\emph{Let \(u:G\to H\) be a morphism of \(S\)-groups. Suppose
that, for every \(s\in S\), \(G_s\) is locally of finite type over
\(\kappa(s)\).}
\begin{enumerate}
\item[\textup{a)}] \emph{If \(u\) is locally of finite type, the
  following conditions are equivalent:}
  \begin{enumerate}
  \item[\textup{(i)}] \emph{\(u\) is locally quasi-finite;}
  \item[\textup{\SGAInlineRefTarget{cor:sga3-VIB-2.8:item:ii}(ii)}] \emph{for every \(s\in S\), the morphism
    \(u_s:G_s\to H_s\) is locally quasi-finite;}
  \item[\textup{\SGAInlineRefTarget{cor:sga3-VIB-2.8:item:iii}(iii)}] \emph{\(\operatorname{Ker}u\) is locally
    quasi-finite over \(S\);}
  \item[\textup{(iv)}] \emph{the fibers of
    \(\operatorname{Ker}u\) are discrete.}
  \end{enumerate}
\item[\textup{b)}] \emph{If \(u\) is locally of finite
  presentation, the following conditions are equivalent:}
  \begin{enumerate}
  \item[\textup{(v)}] \emph{\(u\) is unramified;}
  \item[\textup{(vi)}] \emph{for every \(s\in S\), the morphism
    \(u_s:G_s\to H_s\) is unramified;}
  \item[\textup{\SGAInlineRefTarget{cor:sga3-VIB-2.8:item:vii}(vii)}] \emph{\(\operatorname{Ker}u\) is
    unramified;}
  \item[\textup{(viii)}] \emph{the identity section
    \(S\to\operatorname{Ker}u\) is an open immersion.}
  \end{enumerate}
\end{enumerate}

The equivalences
\[
  \textup{(i)}\Longleftrightarrow\textup{(ii)}
  \qquad\text{and}\qquad
  \textup{(v)}\Longleftrightarrow\textup{(vi)}
\]
follow from \hyperref[concl:sga3-VIB-2.5-ii]{2.5}\textup{\SGARefLink{cor:sga3-VIB-2.8:item:ii}{(ii)}}, and Reminder~\hyperref[rem:sga3-VIB-2.8.1]{2.8.1} below shows that
\[
  \textup{(i)}\Longrightarrow\textup{(iii)}
  \qquad\text{and}\qquad
  \textup{(v)}\Longrightarrow\textup{(vii)}.
\]

For every \(s\in S\), denote by \(e_s\) the identity element of
\(H_s\). Then \textup{\SGARefLink{cor:sga3-VIB-2.8:item:iii}{(iii)}} (respectively, \textup{\SGARefLink{cor:sga3-VIB-2.8:item:vii}{(vii)}})
implies that, for every \(s\in S\),
\[
  (\operatorname{Ker}u)_s
    =\operatorname{Ker}u_s
    =u_s^{-1}(e_s)
\]
is locally quasi-finite (respectively, unramified) over
\(\kappa(s)=\kappa(e_s)\). It follows that \(u_s\) is quasi-finite
(respectively, unramified) at the identity point of \(G_s\), and
hence, by \SGARefLink{prop:sga3-VIB-1.3}{1.3}, that \(u_s\) is locally quasi-finite (respectively,
unramified). Thus
\[
  \textup{(iii)}\Longrightarrow\textup{(ii)}
  \qquad\text{and}\qquad
  \textup{(vii)}\Longrightarrow\textup{(vi)}.
\]
Finally,
\[
  \textup{(ii)}\Longleftrightarrow\textup{(iv)}
  \quad\text{by 1.4.1,}\qquad
  \textup{(vii)}\Longleftrightarrow\textup{(viii)}
  \quad\text{by 2.7.}
\]

\medskip
\phantomsection\label{rem:sga3-VIB-2.8.1}%
\noindent\textbf{Reminder \SGARefLink{rem:sga3-VIB-2.8.1}{2.8.1}.---}
Recall (cf.\ \SGARefLink{sga3:i:definition:2-3-6-1}{Expos\'e~I, 2.3.6.1}) that, given a morphism
\(u:G\to H\) of \(S\)-groups, the \emph{kernel of \(u\)}, denoted
\(\operatorname{Ker}u\), is the group subfunctor of \(G\) defined,
for every morphism \(f:T\to S\), by
\[
  (\operatorname{Ker}u)(T)
    =\{a\in G(T)\mid u\circ a=\epsilon_H\circ f\}.
\]
By loc.\ cit.\ (or EGA~I, 4.4.1), this functor is represented by
the \(S\)-group
\[
  G\times_H S=u^{-1}\bigl(\epsilon_H(S)\bigr),
\]
which is denoted simply by \(\operatorname{Ker}u\). In particular,
the structural morphism
\[
  \operatorname{Ker}u\longrightarrow S
\]
is obtained from \(u\) by base change.

\medskip
\phantomsection\label{lem:sga3-VIB-2.9}%
\noindent\textbf{Lemma 2.9.---}
\emph{Let \(\pi:X\to S\) be a morphism admitting an \(S\)-section
\(\epsilon\).}
\begin{enumerate}
\item[\textup{(i)}] \emph{If \(\pi\) is injective, then it is
  integral.}%
  \begingroup
  \renewcommand{\thefootnote}{\ensuremath{(\ast)}}%
  \footnote[1]{\label{note:sga3-VIB-2.9-star}This is also a
  special case of EGA~IV\(_4\), 18.12.11, since \(\pi\) is
  plainly a universal homeomorphism.}%
  \endgroup
\item[\textup{(ii)}] \emph{If \(\pi\) is locally of finite type
  and, for every \(s\in S\), \(\pi_s\) is an isomorphism, then
  \(\pi\) is an isomorphism.}%
  \begingroup
  \renewcommand{\thefootnote}{\ensuremath{(\ast\ast)}}%
  \footnote[2]{\label{note:sga3-VIB-2.9-starstar}This is also an
  immediate consequence of EGA~IV\(_4\), 18.12.6.}%
  \endgroup
\end{enumerate}

First observe that, by \hyperref[lem:sga3-VIB-2.9.1]{Lemma~2.9.1} below, \(\pi\), being a
homeomorphism, is an affine morphism.

If \(\pi\) is injective, then \(\epsilon\) is surjective. Since
\(\epsilon\) is a surjective immersion, \(\epsilon(S)\) is
isomorphic to the closed subscheme of \(X\) defined by a nil ideal
\(\mathcal I\) of \(\mathcal O_X\). Since \(\epsilon\) is an
\(S\)-section of \(\pi\), there is a direct-sum decomposition
\[
  \mathcal O_X=\mathcal O_S\oplus\mathcal I
\]
of \(\mathcal O_S\)-modules. Since \(\mathcal I\) is a nil ideal
of \(\mathcal O_X\), it is plainly integral over \(\mathcal O_S\).
Thus \(\mathcal O_X\) is integral over \(\mathcal O_S\), and
\(\pi\) is integral.

Now suppose that \(\pi\) is locally of finite type. Since
\(\pi\circ\epsilon=\operatorname{id}_S\), the section \(\epsilon\)
is locally of finite presentation
(cf.\ EGA~IV\(_1\), 1.4.3(v)); hence \(\mathcal I\) is an ideal of
finite type in \(\mathcal O_X\) (EGA~IV\(_1\), 1.4.1). For every
\(s\in S\), one has
\[
  \mathcal O_{X_s}
    =\kappa(s)\oplus
      \mathcal I\otimes_{\mathcal O_S}\kappa(s).
\]
By hypothesis, \(\epsilon_s\) is an isomorphism, so
\[
  \mathcal I\otimes_{\mathcal O_S}\kappa(s)=0
\]
for every \(s\in S\), and therefore, a fortiori,
\[
  \mathcal I\otimes_{\mathcal O_X}\kappa(x)=0
\]
for every \(x\in X\). Nakayama's lemma now gives
\(\mathcal I=0\), and consequently \(\pi\) is an isomorphism.

\medskip
\phantomsection\label{lem:sga3-VIB-2.9.1}%
\noindent\textbf{Lemma 2.9.1.---}
\emph{Let \(f:X\to Y\) be a morphism of schemes that is a
homeomorphism. Then \(f\) is an affine morphism.}%
\begingroup
\renewcommand{\thefootnote}{\ensuremath{(\ast\ast\ast)}}%
\footnote[3]{\label{note:sga3-VIB-2.9.1-starstarstar}Cf.\
EGA~IV\(_4\), 18.12.7.1 for a slightly more general result, proved
by the same argument.}%
\endgroup

It is enough to show that every point \(y\in Y\) has an open
neighborhood \(W\) such that the restriction of \(f\) over \(W\)
is affine. Let \(y\in Y\), and let \(V\) be an affine open
neighborhood of \(y\) in \(Y\). Set
\[
  V'=f^{-1}(V).
\]
Then \(V'\) is an open neighborhood of \(x=f^{-1}(y)\) in \(X\).
There is an affine open neighborhood \(W'\) of \(x\) in \(X\)
contained in \(V'\). Set
\[
  W=f(W').
\]
Then \(W\) is an open neighborhood of \(y\) in \(Y\), contained
in the affine scheme \(V\), and therefore \(W\) is separated.
Since \(W'\) is affine, the restriction of \(f\) over \(W\) is an
affine morphism (EGA~II, 1.2.3).

\medskip
\phantomsection\label{cor:sga3-VIB-2.10}%
\noindent\textbf{Corollary 2.10.---}
\emph{Let \(G\) be an \(S\)-group locally of finite type. Suppose
that, for every \(s\in S\), \(G_s\) is the identity
\(\kappa(s)\)-group. Then \(G\) is the identity \(S\)-group.}

More generally:

\medskip
\phantomsection\label{cor:sga3-VIB-2.11}%
\noindent\textbf{Corollary 2.11.---}
\emph{Let \(f:X\to Y\) be an \(S\)-morphism locally of finite
type. In order that \(f\) be a monomorphism, it is necessary and
sufficient that \(f_s\) be a monomorphism for every \(s\in S\).}

The condition is plainly necessary; we prove that it is
sufficient. If \(f_s\) is a monomorphism for every \(s\in S\),
then, a fortiori, \(f_y\) is a monomorphism for every \(y\in Y\).
We may therefore assume that \(Y=S\).

By EGA~I, 5.3.8, to prove that \(f\) is a monomorphism it is enough
to show that
\[
  \Delta_f:X\longrightarrow X\times_SX
\]
is an isomorphism, or, equivalently, that the first projection
\[
  p:X\times_SX\longrightarrow X
\]
is an isomorphism. If \(f_s\) is a monomorphism, EGA~I, 5.3.8
likewise shows that the first projection
\[
  X_s\times_{\kappa(s)}X_s\longrightarrow X_s,
\]
which is identified with \(p_s\), is an isomorphism. But \(p\)
has the \(S\)-section \(\Delta_f\); hence \hyperref[lem:sga3-VIB-2.9]{Lemma~2.9} shows that, if
\(p_s\) is an isomorphism for every \(s\in S\), then \(p\) itself
is an isomorphism.

\medskip
\phantomsection\label{cor:sga3-VIB-2.12}%
\noindent\textbf{Corollary 2.12.---}
\emph{Let \(G\) be an \(S\)-scheme admitting an \(S\)-section
\(\epsilon\), and suppose that the structural morphism
\(\pi:G\to S\) is closed. Let \(s\in S\) be such that \(\pi\) is
of finite presentation at \(\epsilon(s)\), and such that}
\[
  \epsilon_s:\operatorname{Spec}\bigl(\kappa(s)\bigr)
    \longrightarrow G_s
\]
\emph{is an isomorphism (or, equivalently, such that}
\[
  \pi_s:G_s\longrightarrow
    \operatorname{Spec}\kappa(s)
\]
\emph{is an isomorphism). Then there exists an open neighborhood
\(U\) of \(s\) in \(S\) such that}
\[
  \epsilon|_U:U\longrightarrow G\times_SU
\]
\emph{is an isomorphism.}

Let \(V\) be the set of points of \(G\) at which \(\pi\) is
unramified. This set is known to be open
(cf.\ Editor's \hyperref[note:sga3-VIB-6]{note~6}), and it contains \(\epsilon(s)\). Hence
\[
  W=\epsilon^{-1}(V)
\]
is an open subset of \(S\) containing \(s\), and \(\pi\) is
unramified at \(\epsilon(t)\) for every \(t\in W\). Since \(\pi\)
is closed, the same is true of \(\pi|_W\), and we may therefore
suppose that \(W=S\).

By \hyperref[lem:sga3-VIB-2.7.1]{2.7.1},
\[
  X=G-\epsilon(S)
\]
is then a closed subset of \(G\) that does not meet \(G_s\).
Since \(\pi\) is closed,
\[
  F=\pi(X)
\]
is a closed subset of \(S\) that does not contain \(s\). Set
\[
  U=S-F.
\]
Then \(U\) is an open subset of \(S\), and \(\epsilon|_U\) is an
isomorphism from \(U\) onto \(G|_U\).

% END INCLUDED SGA3_English_current_buildable_source_20260823/source/inputs/VI/tex_reference_v2/components/33_expose_VIB_cor28_through_cor212_en.tex

% BEGIN INCLUDED SGA3_English_current_buildable_source_20260823/source/inputs/VI/tex_reference_v2/components/34_expose_VIB_section3_through_cor36_loop2_r1_en.tex
% SGA 3, Expose VI_B, Loop 1 English reconstruction.
% Scope: Section 3 heading and paragraph 3.0 through Corollary 3.6
% in full; Proposition 3.7 is excluded.
% Authority: Polo--Gille Exp6B-13oct24.pdf. The semantic boundary
% begins at the bottom of component-local page 16 and ends on local
% page 18 (combined-reader pages 357--359; original-margin printed
% pages 341--343). Proposition 3.7 begins on local page 19 /
% combined-reader page 360, before the printed-page 344 marker.
% OCR and the Reinhold English Markdown are comparison witnesses only.
% The Proposition 3.3 diagram is an exact 180-dpi crop of the
% controlling PDF: local p. 18, x=530,y=415,w=430,h=220 on a
% 1488x2105 render.

\section*{3. Identity component of a group locally of finite
  presentation}
\phantomsection\label{sec:sga3-VIB-3}
\addcontentsline{toc}{section}{3. Identity component of a group
  locally of finite presentation}

\medskip
\phantomsection\label{par:sga3-VIB-3.0}%
\noindent\textbf{3.0.---}
Given a subset \(A\) (respectively, \(B\)) of an \(S\)-scheme
\(X\) (respectively, \(Y\)), by abuse of notation
\[
  A\times_S B
\]
will denote the subset of \(X\times_S Y\) consisting of the points
whose first projection belongs to \(A\) and whose second projection
belongs to \(B\).

Given a subset \(A\) of an \(S\)-group \(G\), we shall say that
\(A\) is \emph{stable under the group law of \(G\)} if
\[
  c(A)\subset A
  \qquad\text{and}\qquad
  \mu(A\times_S A)\subset A.
\]

\medskip
\phantomsection\label{def:sga3-VIB-3.1}%
\noindent\textbf{Definition 3.1.---}
Let \(G\) be an \(S\)-functor in groups satisfying the following
condition:
\[
  \tag{+}\label{sga3:vi:definition:3-1:equation:plus}
  \text{for every \(s\in S\), the functor
  \(G_s=G\otimes_S\kappa(s)\) is representable.}
\]
We then denote by \(\underline{G_s^0}\) the connected component of
the identity element of the \(\kappa(s)\)-group \(G_s\).%
\footnote{\label{note:sga3-VIB-22}Editor's note: this is a slightly
abusive notation, but it is compatible with the notation of
Expos\'e~VI\(_A\), \SGARefLink{cor:sga3-VIB-2.3}{2.3} when this connected component is the
underlying topological space of an open group subscheme \(G^0\) of
\(G\); cf.\ Expos\'e~VI\(_A\), \SGARefLink{par:sga3-VIA-2.2bis}{2.2.b}is.}
We define a sub-\(S\)-functor in groups of \(G\), called the
\emph{identity component of \(G\)} and denoted by \(G^0\), by
setting, for every morphism \(T\to S\),
\[
  G^0(T)
    =
  \bigl\{u\in G(T)\bigm|
    \text{for every \(s\in S\), }
    u_s(T_s)\subset\underline{G_s^0}\bigr\}.
\]
Thus we have defined the functor
\[
  G\longmapsto G^0
\]
from
\[
  \bigl(\widehat{\mathbf{Sch}}_{/S}\bigr)\text{-gr.}
  \quad\text{to}\quad
  \bigl(\widehat{\mathbf{Sch}}_{/S}\bigr)\text{-gr.}
\]

\medskip
\phantomsection\label{rem:sga3-VIB-3.2}%
\noindent\textbf{Remark 3.2.---}
\textup{(i)} Let \(G\) be an \(S\)-functor in groups satisfying
condition \((+)\). Then
\[
  \operatorname{Lie}(G^0/S)=\operatorname{Lie}(G/S)
\]
by Expos\'e~II.%
\footnote{\label{note:sga3-VIB-23}Editor's note: we have added what
follows.}
Indeed, for every \(S\)-scheme \(T\), denote by \(I_T\) the
\(T\)-scheme of dual numbers over \(T\), and by
\[
  \tau:T\longrightarrow I_T
\]
the zero section (cf.\ \SGARefLink{sga3:ii:definition:2-1}{Expos\'e~II, 2.1}). By definition,
\[
  \operatorname{Lie}(G/S)(T)
    =
  \{u\in G(I_T)\mid u\circ\tau=e\},
\]
where \(e:T\to G\) denotes the composite of \(T\to S\) with the
identity section \(S\to G\). Likewise,
\[
  \begin{aligned}
  \operatorname{Lie}(G^0/S)(T)
    &=
  \{u\in G^0(I_T)\mid u\circ\tau=e\} \\
    &=
  \bigl\{u\in G(I_T)\bigm|
    u\circ\tau=e
    \ \text{and}\
    u_s\bigl((I_T)_s\bigr)\subset\underline{G_s^0}
    \text{ for every \(s\in S\)}\bigr\}.
  \end{aligned}
\]
For every \(s\in S\), the schemes \(T_s\) and
\[
  (I_T)_s=I_{T_s}
\]
have the same underlying set. Hence, if
\(u\in\operatorname{Lie}(G/S)(T)\), then
\[
  u_s(I_{T_s})=e_s,
\]
where \(e_s\) denotes the identity point of \(G_s\), and therefore
\(u\in\operatorname{Lie}(G^0/S)(T)\). Thus the inclusion
\[
  \operatorname{Lie}(G^0/S)(T)
    \subset
  \operatorname{Lie}(G/S)(T)
\]
is an equality for every \(T\), and consequently
\[
  \operatorname{Lie}(G^0/S)=\operatorname{Lie}(G/S).
\]

\textup{(ii)} Let \(G\) and \(G'\) be two \(S\)-functors in groups
satisfying \((+)\). Then:
\begin{enumerate}
\item[\textup{a)}] if \(G\subset G'\), then \(G^0\subset G'^0\);
\item[\textup{b)}] if \(G\subset G'\) and \(G'^0\subset G\), then
  \(G^0=G'^0\);
\item[\textup{c)}] if, for every \(s\in S\), \(G_s\) is locally of
  finite type over \(\kappa(s)\), then \(G^0\) satisfies condition
  \((+)\), by Expos\'e~VI\(_A\), \SGARefLink{cor:sga3-VIB-2.3}{2.3}, and
  \[
    (G^0)^0=G^0.
  \]
\end{enumerate}

\medskip
\phantomsection\label{prop:sga3-VIB-3.3}%
\noindent\textbf{Proposition 3.3.---}
\emph{Let \(G\) be an \(S\)-functor in groups satisfying condition
\((+)\), and let \(S'\) be an \(S\)-scheme. Then}
\[
  (G\times_S S')^0=G^0\times_S S'.
\]
\emph{In other words, the functor \(G\mapsto G^0\) commutes with
base change; that is, the following diagram is commutative:}
\begin{center}
  \phantomsection\label{diag:sga3-VIB-3.3-base-change}%
  \begin{tikzcd}[row sep=large,column sep=huge]
    \bigl(\widehat{\mathrm{Sch}}_{/S}\bigr)\text{-}\mathrm{gr.}
      \arrow[r,"{G\mapsto G^0}"]
      \arrow[d] &
    \bigl(\widehat{\mathrm{Sch}}_{/S}\bigr)\text{-}\mathrm{gr.}
      \arrow[d] \\
    \bigl(\widehat{\mathrm{Sch}}_{/S'}\bigr)\text{-}\mathrm{gr.}
      \arrow[r] &
    \bigl(\widehat{\mathrm{Sch}}_{/S'}\bigr)\text{-}\mathrm{gr.}
  \end{tikzcd}
\end{center}

It is enough to check that, for every \(s'\in S'\), whose image in
\(S\) we denote by \(s\),
\[
  \begin{aligned}
  \bigl((G\times_S S')\otimes_{S'}\kappa(s')\bigr)^0
    &=
  \bigl(G_s\otimes_{\kappa(s)}\kappa(s')\bigr)^0 \\
    &=
  G_s^0\otimes_{\kappa(s)}\kappa(s').
  \end{aligned}
\]
This follows from Expos\'e~VI\(_A\), \SGARefLink{rem:sga3-VIB-2.1.2}{2.1.2}. Note, for later use in
\SGARefLink{prop:sga3-VIB-4.2}{4.2}, that we have not used the group structure of \(G_s\), but
only the fact that \(G_s^0\) has a rational point, namely
\(\epsilon(s)\), and is therefore geometrically connected
(see also EGA~IV\(_2\), 4.5.14).

\medskip
\phantomsection\label{case:sga3-VIB-3.4}%
\noindent\textbf{Special Case 3.4.---}
Let \(G\) be an \(S\)-group scheme, and denote by
\(\underline{G^0}\) the subset of \(G\) equal to the union of the
\(\underline{G_s^0}\) as \(s\) runs through \(S\). Then
\(\underline{G^0}\) is a subset of \(G\) stable under the group law
of \(G\) (cf.\ \hyperref[par:sga3-VIB-3.0]{3.0}), and, for every morphism \(S'\to S\),
\[
  G^0(S')
    =
  \bigl\{u\in G(S')\bigm|
    u(S')\subset\underline{G^0}\bigr\}.
\]
When \(\underline{G^0}\) is an open subset of \(G\), it is clear
that \(G^0\) is represented by the subscheme of \(G\) induced on
the open subset \(\underline{G^0}\).

\medskip
\phantomsection\label{prop:sga3-VIB-3.5}%
\noindent\textbf{Proposition 3.5.---}
\emph{Let \(S\) be a quasi-compact and quasi-separated scheme, and
let \(G\) be an \(S\)-group with fibers locally of finite type.
Then there is a quasi-compact open subset \(U\) of \(G\) containing
\(\underline{G^0}\).}

Since the identity section \(\epsilon\) is an immersion,
\(\epsilon(S)\) is a quasi-compact subspace of \(G\). Thus there is
a quasi-compact open subset \(V\) of \(G\) containing
\(\epsilon(S)\). Since \(S\) is quasi-separated and \(V\) is
quasi-compact, \(V\) is quasi-compact over \(S\)
(EGA~IV\(_1\), 1.2.4). Hence \(V\times_SV\) is quasi-compact over
\(S\), and therefore quasi-compact. It follows that
\[
  V\cdot V=\mu(V\times_SV)
\]
is quasi-compact. Set
\[
  V_s=V\cap G_s,
  \qquad
  V_s^0=V\cap\underline{G_s^0}.
\]
Then \(V_s^0\) is open in \(\underline{G_s^0}\), and is dense in
\(\underline{G_s^0}\) because \(G_s^0\) is irreducible
(Expos\'e~VI\(_A\), \SGARefLink{cor:sga3-VIB-2.4}{2.4}). Therefore
\[
  V_s^0\cdot V_s^0=\underline{G_s^0}
  \qquad
  \text{(Expos\'e~VI\(_A\), 0.5),}
\]
which shows that
\[
  V_s\cdot V_s\supset\underline{G_s^0},
  \qquad\text{and hence}\qquad
  V\cdot V\supset\underline{G^0}.
\]
Finally, since \(V\cdot V\) is quasi-compact, there is a
quasi-compact open subset \(U\) of \(G\) containing \(V\cdot V\),
and therefore containing \(\underline{G^0}\).

\medskip
\phantomsection\label{cor:sga3-VIB-3.6}%
\noindent\textbf{Corollary 3.6.---}
\emph{Let \(G\) be an \(S\)-group whose fibers are locally of
finite type and connected. Then \(G\) is quasi-compact over \(S\).}

The assertion is local on \(S\) (EGA~I, 6.6.1), so we are reduced
to the case where \(S\) is affine. By 3.5, there is then a
quasi-compact open subset \(U\) of \(G\) containing
\[
  \underline{G^0}=G.
\]
Thus \(G\) is quasi-compact, and hence quasi-compact over the affine
scheme \(S\) (EGA~I, 6.6.4(v)).

% END INCLUDED SGA3_English_current_buildable_source_20260823/source/inputs/VI/tex_reference_v2/components/34_expose_VIB_section3_through_cor36_loop2_r1_en.tex

% BEGIN INCLUDED SGA3_English_current_buildable_source_20260823/source/inputs/VI/tex_reference_v2/components/35_expose_VIB_prop37_cor38_en.tex
% SGA 3, Expose VI_B, Loop 1 English reconstruction.
% Scope: Proposition 3.7 and Corollary 3.8 in full.
% Authority: Polo--Gille Exp6B-13oct24.pdf, component-local pages
% 19--20 (combined-reader pages 360--361; original-margin printed
% pages 343--345).
% OCR and the Reinhold English Markdown are comparison witnesses only.

\medskip
\phantomsection\label{prop:sga3-VIB-3.7}%
\noindent\textbf{Proposition 3.7.---}%
\footnote{\label{note:sga3-VIB-24}Editor's note: we recall the following
definitions and results (cf.\ EGA~0\(_{\mathrm{III}}\), \S9.1, and
EGA~IV\(_1\), \S\S1.8 and 1.9). Let \(X\) be a topological space.
\textup{(i)} An open subset \(U\) of \(X\) is called
\emph{retrocompact} if the inclusion \(U\hookrightarrow X\) is
quasi-compact, i.e.\ if \(U\cap V\) is quasi-compact for every
quasi-compact open subset \(V\subset X\).
\textup{\SGARefLink{prop:sga3-VIB-3.7:item:ii}{(ii)}} A subset \(C\) of \(X\) is called \emph{constructible}
if it is a finite union of subsets of the form \(U-V\), where \(U\)
and \(V\) are retrocompact open subsets of \(X\).
\textup{\SGARefLink{prop:sga3-VIB-3.7:item:iii}{(iii)}} A subset \(L\) of \(X\) is called \emph{locally
constructible} if, for every \(x\in X\), there is an open
neighborhood \(U\) of \(x\) in \(X\) such that \(L\cap U\) is
constructible in \(U\). (If \(X\) is quasi-compact and
quasi-separated, then \(L\) is constructible.)
\textup{(iv)} A subset \(I\) of \(X\) is called
\emph{ind-constructible} if, for every \(x\in X\), there is an open
neighborhood \(U\) of \(x\) in \(X\) such that \(I\cap U\) is a
union of locally constructible subsets of \(U\).
Now let \(f:X\to Y\) be a morphism of schemes. By Chevalley's
constructibility theorem (cf.\ EGA~IV\(_1\), 1.8.4 and
1.9.5(viii)), if \(f\) is of finite presentation (respectively,
locally of finite presentation), then the image under \(f\) of every
locally constructible subset is locally constructible
(respectively, ind-constructible).}
\emph{Let \(G\) be an \(S\)-group locally of finite presentation.
Then:}
\begin{enumerate}
\item[\textup{\SGAInlineRefTarget{prop:sga3-VIB-3.7:item:i}(i)}] \emph{\(\underline{G^0}\) is ind-constructible
  in \(G\).}
\item[\textup{\SGAInlineRefTarget{prop:sga3-VIB-3.7:item:ii}(ii)}] \emph{If, moreover, \(G\) is quasi-separated
  over \(S\), and \(S\) is quasi-compact and quasi-separated, then
  \(\underline{G^0}\) is constructible.}
\item[\textup{\SGAInlineRefTarget{prop:sga3-VIB-3.7:item:iii}(iii)}] \emph{Consequently, if \(G\) is
  quasi-separated over \(S\), then \(\underline{G^0}\) is locally
  constructible.}
\end{enumerate}

We first prove the first assertion. Since
\[
  \pi:G\longrightarrow S
\]
is locally of finite presentation, given \(s\in S\), there are an
open subset \(U\) of \(G\) containing \(\epsilon(s)\) and an open
subset \(V\) of \(S\) containing \(s\) such that
\[
  \pi(U)\subset V
\]
and the induced morphism
\[
  \pi':U\longrightarrow V
\]
is of finite presentation. Then
\[
  T=\epsilon^{-1}(U)
\]
is an open subset of \(S\). Set
\[
  W=(\pi')^{-1}(T),\qquad
  \pi''=\pi'|_W.
\]
Thus \(\pi'':W\to T\) is of finite presentation and admits as a
section the morphism
\[
  \epsilon'':T\longrightarrow W
\]
induced by \(\epsilon\).

For every \(t\in T\), since \(\underline{G_t^0}\) is irreducible
(VI\(_A\), \hyperref[prop:sga3-VIA-2.4]{2.4}), \(W\cap\underline{G_t^0}\) is dense in
\(\underline{G_t^0}\), hence irreducible and therefore connected.
It is thus the connected
component of \((\pi'')^{-1}(t)\) containing \(\epsilon''(t)\).
It follows from EGA~IV\(_3\), 9.7.12, that the union \(W^0\) of the
\(W\cap\underline{G_t^0}\), for \(t\in T\), is locally
constructible in \(W\).
On the other hand, VI\(_A\), \SGARefLink{sga3:vi:proposition:0-5}{0.5}, gives
\[
  \underline{G^0}\cap\pi^{-1}(T)=W^0\cdot W^0;
\]
that is, \(\underline{G^0}\cap\pi^{-1}(T)\) is the image of
\[
  W^0\times_T W^0
\]
under the morphism
\[
  \mu'':W\times_T W\longrightarrow\pi^{-1}(T)
\]
induced by \(\mu\).%
\footnote{\label{note:sga3-VIB-25}Editor's note: we have expanded the
original in what follows.}
Since \(W\times_T W\) (respectively, \(\pi^{-1}(T)\)) is of finite
presentation (respectively, locally of finite presentation) over
\(T\), the morphism \(\mu''\) is locally of finite presentation and
quasi-separated, by EGA~IV\(_1\), 1.4.3 and 1.2.2. If, moreover,
\(\pi^{-1}(T)\) is quasi-separated over \(T\), then \(\mu''\) is
quasi-compact (loc.\ cit., 1.2.4), hence of finite presentation.
Since \(W^0\times_T W^0\) is locally constructible in \(W\times_T W\)
(because \(W^0\) is locally constructible in \(W\)), Chevalley's
constructibility theorem (loc.\ cit., 1.8.4 and 1.9.5(viii)) shows
that \(\underline{G^0}\cap\pi^{-1}(T)\) is ind-constructible in
\(\pi^{-1}(T)\), and is locally constructible in \(\pi^{-1}(T)\)
if \(G\), and hence \(\pi^{-1}(T)\), is quasi-separated over \(T\).
This proves assertions \textup{\SGARefLink{prop:sga3-VIB-3.7:item:i}{(i)}} and \textup{\SGARefLink{prop:sga3-VIB-3.7:item:iii}{(iii)}}.

Now suppose that \(G\) is quasi-separated over \(S\), and that \(S\)
is quasi-compact and quasi-separated. By \hyperref[prop:sga3-VIB-3.5]{3.5}, there is a
quasi-compact open subset \(U\) of \(G\) containing
\(\underline{G^0}\). Since
\(G\) is quasi-separated over \(S\), \(G\) is quasi-separated;
hence the open subset \(U\) is retrocompact (EGA~IV\(_1\), 1.2.7),
and it suffices to show that \(\underline{G^0}\) is constructible
in \(U\)
(EGA~0\(_{\mathrm{III}}\), 9.1.8).

Moreover, \(U\) is quasi-compact, hence quasi-compact over \(S\)
(EGA~IV\(_1\), 1.2.4), and is quasi-separated over \(S\).
Therefore the restriction of \(\pi\) to \(U\) is of finite
presentation. By EGA~IV\(_3\), 9.7.12, \(\underline{G^0}\) is locally
constructible in \(U\), hence constructible in \(U\), since \(U\)
is quasi-compact and quasi-separated (EGA~IV\(_1\), 1.8.1).
This proves \textup{(ii)}, and it follows that, for every
quasi-compact and quasi-separated open subset \(T\) of \(S\)
(for example, every affine open subset of \(S\)),
\[
  \underline{G^0}\cap\pi^{-1}(T)
\]
is constructible.

\medskip
\phantomsection\label{cor:sga3-VIB-3.8}%
\noindent\textbf{Corollary 3.8.---}
\emph{Let \(S_0\) be a quasi-compact and quasi-separated scheme,
\(I\) an increasing filtered preordered set,
\((\mathcal A_i)_{i\in I}\) a direct system of commutative
quasi-coherent \(\mathcal O_{S_0}\)-algebras,}
\[
  \mathcal A=\varinjlim_i\mathcal A_i,\qquad
  S_i=\operatorname{Spec}\mathcal A_i\quad(i\in I),\qquad
  S=\operatorname{Spec}\mathcal A
\]
\emph{\textup{(cf.\ EGA~II, 1.3.1)}. Let \(G\) be an
\(S_0\)-group scheme locally of finite presentation. Then the
canonical map}
\[
  \varinjlim_i G^0(S_i)\longrightarrow G^0(S)
\]
\emph{is bijective.}

Since \(G\) is locally of finite presentation over \(S_0\),%
\begingroup
\renewcommand{\thefootnote}{\fnsymbol{footnote}}%
\footnote[1]{\label{note:sga3-VIB-3.8-base}Translator's note: the
French source prints ``over \(S\)'' here. The scheme \(G\) was
defined over \(S_0\), while \(S\) and the \(S_i\) are the test
schemes of the direct system; the required base is therefore
\(S_0\).}%
\endgroup{}
the canonical map
\[
  \varinjlim_i G(S_i)\longrightarrow G(S)
\]
is bijective, by EGA~IV\(_2\), 8.14.2(c). It follows immediately
that the canonical map
\[
  \varinjlim_i G^0(S_i)\longrightarrow G^0(S)
\]
is injective. We prove that it is surjective. Let
\[
  g\in G^0(S)\subset G(S).
\]
There is an \(i\in I\) such that \(g\) factors, by means of an
element \(g_i\in G(S_i)\), through \(S\to S_i\). By hypothesis,
\[
  g^{-1}(\underline{G^0})=S.
\]
But \(\underline{G^0}\) is ind-constructible in \(G\), by \hyperref[prop:sga3-VIB-3.7]{3.7}, so
\[
  g_i^{-1}(\underline{G^0})
\]
is ind-constructible in \(S_i\). It follows from EGA~IV\(_2\),
8.3.4, that there is an index \(j\geq i\) such that
\[
  g_j^{-1}(\underline{G^0})=S_j,
\]
where \(g_j\) is the map induced by \(g_i\) under the base change
\(S_j\to S_i\). Thus \(g_j\in G^0(S_j)\), and hence \(g\) comes
from an element of
\[
  \varinjlim_i G^0(S_i).
\]

% END INCLUDED SGA3_English_current_buildable_source_20260823/source/inputs/VI/tex_reference_v2/components/35_expose_VIB_prop37_cor38_en.tex

% BEGIN INCLUDED SGA3_English_current_buildable_source_20260823/source/inputs/VI/tex_reference_v2/components/36_expose_VIB_prop39_through_lemma3101_en.tex
% SGA 3, Expose VI_B, Loop 1 English reconstruction.
% Scope: Proposition 3.9, Theorem 3.10, and Lemma 3.10.1 in full.
% Authority: Polo--Gille Exp6B-13oct24.pdf, component-local pages
% 20--22 (combined-reader pages 361--363; original-margin printed
% pages 345--347, stopping before Proposition 3.11).
% OCR and the Reinhold English Markdown are comparison witnesses only.

\medskip
\phantomsection\label{prop:sga3-VIB-3.9}%
\noindent\textbf{Proposition 3.9.---}
\emph{Let \(G\) be an \(S\)-group locally of finite presentation.
Suppose that \(G^0\) is representable. Then the canonical morphism}
\[
  i:G^0\longrightarrow G
\]
\emph{is an open immersion; moreover, \(G^0\) is quasi-compact over
\(S\).}

Since \(G^0\) is a subfunctor of the functor \(G\), the morphism
\(i\) is a monomorphism, hence is radicial. From the definition of
the functor \(G^0\), one checks immediately from the definition
(EGA~IV\(_4\), 17.1.1) that \(i\) is formally \'etale (noting that
\(G^0\) is the image of \(i\) in \(G\)). Finally, by the
characterization (EGA~IV\(_3\), 8.14.2(c)) of \(S\)-schemes locally
of finite presentation in terms of the functors they represent, and
by \SGARefLink{cor:sga3-VIB-3.8}{3.8}, \(G^0\) is locally of finite presentation over \(S\), since
\(G\) is. Hence \(i\) is locally of finite presentation
(cf.\ EGA~IV\(_1\), 1.4.3). It is therefore a radicial \'etale
morphism, and hence an open immersion (EGA~IV\(_4\), 17.9.1).

The last assertion follows from \hyperref[cor:sga3-VIB-3.6]{3.6}.

\medskip
\phantomsection\label{thm:sga3-VIB-3.10}%
\noindent\textbf{Theorem 3.10.---}
\emph{Let \(G\) be an \(S\)-group. The following conditions are
equivalent:}
\begin{enumerate}
\item[\textup{(i)}] \emph{\(G\) is smooth over \(S\) at the points
  of the identity section.}
\item[\textup{(ii)}] \emph{\(G\) is flat and locally of finite
  presentation over \(S\) at the points of the identity section,
  and, for every \(s\in S\), \(G_s\) is smooth over \(\kappa(s)\).}
\item[\textup{\SGAInlineRefTarget{thm:sga3-VIB-3.10:item:iii}(iii)}] \emph{There is an open group subscheme \(G'\)
  of \(G\) that is smooth over \(S\).}
\item[\textup{\SGAInlineRefTarget{thm:sga3-VIB-3.10:item:iv}(iv)}] \emph{\(G^0\) is representable by an open
  subscheme of \(G\) that is smooth over \(S\).}
\end{enumerate}

It is clear that
\[
  \textup{(iv)}\Longrightarrow\textup{(iii)}
  \Longrightarrow\textup{(i)},
\]
and, by \hyperref[cor:sga3-VIB-1.3.1]{1.3.1} and \hyperref[cor:sga3-VIB-2.4]{2.4}, \textup{\SGARefLink{cor:sga3-VIB-1.3.1:item:i}{(i)}} implies both \textup{\SGARefLink{cor:sga3-VIB-1.3.1:item:ii}{(ii)}} and
\textup{(iii)}. Moreover, \textup{(ii)} implies \textup{(i)} by
EGA~IV\(_4\), 17.5.1.

It remains to prove that \textup{\SGARefLink{thm:sga3-VIB-3.10:item:iii}{(iii)}} implies \textup{\SGARefLink{thm:sga3-VIB-3.10:item:iv}{(iv)}}.
\hyperref[lem:sga3-VIB-3.10.1]{Lemma~3.10.1} below shows that \(G'\) contains \(\underline{G^0}\) and that
\[
  G'^0=G^0.
\]
It therefore suffices to show that \(\underline{G^0}\) is open in \(G\), for we
have already seen in \hyperref[case:sga3-VIB-3.4]{3.4} that \(G^0\) will then be represented by
the smooth group subscheme induced by \(G'\) on the open subset
\[
  \underline{G'^0}=\underline{G^0}.
\]
We may thus suppose that \(G'=G\).

To show that \(\underline{G^0}\) is open, it suffices to show that every
\(s\in S\) has a neighborhood \(T\) in \(S\) such that
\[
  \underline{G^0}\cap\pi^{-1}(T)
\]
is open in \(\pi^{-1}(T)\). Let \(s\in S\). Since \(G=G'\), the
morphism \(\pi\) is locally of finite presentation. As in the proof
of \SGARefLink{prop:sga3-VIB-3.7}{3.7}, we may therefore construct an open subset \(T\) of \(S\)
containing \(s\), and an open subset \(W\) of \(G\) containing
\(\epsilon(s)\), such that the induced morphism
\[
  \pi'':W\longrightarrow T
\]
is of finite presentation and admits as a section the morphism
\[
  \epsilon'':T\longrightarrow W
\]
induced by \(\epsilon\). For every \(t\in T\),
\[
  W\cap\underline{G_t^0}
\]
is then the connected component of \((\pi'')^{-1}(t)\) containing
\(\epsilon''(t)\). Since \(\pi\) is smooth, so is \(\pi''\), which
is therefore smooth and of finite presentation. By EGA~IV\(_3\),
15.6.5, the union \(W^0\) of the \(W\cap\underline{G_t^0}\), for \(t\in T\),
is open in \(W\).

On the other hand, VI\(_A\), \SGARefLink{sga3:vi:proposition:0-5}{0.5}, gives
\[
  W^0\cdot W^0=\underline{G^0}\cap\pi^{-1}(T),
\]
and we must show that this is open in \(\pi^{-1}(T)\). We may now
suppose that \(T=S\); it remains to show that \(W^0\cdot W^0\) is
open in \(G\). Since \(\pi\) is universally open, so is \(\mu\)
(VI\(_A\), \SGARefLink{sga3:VIA:0.1}{0.1}). Hence, since \(W^0\) is open in \(G\), so is
\[
  W^0\cdot W^0=\mu(W^0\times_S W^0).
\]

This result will be strengthened in \SGARefLink{cor:sga3-VIB-4.4}{4.4}.

\medskip
\phantomsection\label{lem:sga3-VIB-3.10.1}%
\noindent\textbf{Lemma 3.10.1.---}
\emph{Let \(G\) be an \(S\)-group whose fibers are locally of
finite type. Then every open group subscheme \(H\) of \(G\) contains
\(\underline{G^0}\) and satisfies}
\[
  G^0=H^0.
\]

Let \(s\in S\), and set
\[
  G'_s=H_s\cap\underline{G_s^0}.
\]
Then \(G'_s\) is an open subset of \(\underline{G_s^0}\), and is dense in
\(\underline{G_s^0}\) since \(\underline{G_s^0}\) is irreducible
(VI\(_A\), \hyperref[cor:sga3-VIB-2.4]{2.4}). Hence
\[
  G'_s\cdot G'_s=\underline{G_s^0}
\]
by VI\(_A\), \SGARefLink{sga3:vi:proposition:0-5}{0.5}, which shows that
\[
  \underline{G_s^0}=G'_s\cdot G'_s\subset H_s\cdot H_s=H_s.
\]
Thus \(\underline{G_s^0}=\underline{H_s^0}\) for every \(s\), whence
\[
  \underline{G^0}\subset H,\qquad G^0=H^0.
\]

% END INCLUDED SGA3_English_current_buildable_source_20260823/source/inputs/VI/tex_reference_v2/components/36_expose_VIB_prop39_through_lemma3101_en.tex

% BEGIN INCLUDED SGA3_English_current_buildable_source_20260823/source/inputs/VI/tex_reference_v2/components/37_expose_VIB_prop311_section4_prop41_en.tex
% SGA 3, Expose VI_B, Loop 1 English reconstruction.
% Scope: Proposition 3.11, the Section 4 heading, and Proposition 4.1.
% Authority: Polo--Gille Exp6B-13oct24.pdf, component-local page
% 22 (combined-reader page 363; original-margin printed
% pages 347--348).
% OCR and the Reinhold English Markdown are comparison witnesses only.

\medskip
\phantomsection\label{prop:sga3-VIB-3.11}%
\noindent\textbf{Proposition 3.11.---}%
\footnote{\label{note:sga3-VIB-26}Editor's note: compare
VI\(_A\), \SGARefLink{par:sga3-VIB-5.6}{5.6}.}
\emph{Let \(u:G\to H\) be a morphism between \(S\)-groups locally
of finite presentation. If \(u\) is flat, the induced map}
\[
  u^0:G^0\longrightarrow H^0
\]
\emph{is surjective. The converse holds if \(G\) is flat over \(S\)
and \(H\) has reduced fibers.}

Suppose that \(u\) is flat. Then, for every \(s\in S\), \(u_s\) is
flat and locally of finite presentation, hence open
(EGA~IV\(_2\), 2.4.6). Therefore
\[
  u_s^0:\underline{G_s^0}\longrightarrow\underline{H_s^0}
\]
is surjective, by \SGARefLink{cor:sga3-VIB-1.3.2}{1.3.2}.%
\footnote{\label{note:sga3-VIB-27}Editor's note: we have shortened
the original by referring here to \SGARefLink{cor:sga3-VIB-1.3.2}{1.3.2}. We have also simplified
the argument below by omitting a reference to the generic flatness
theorem (EGA~IV\(_2\), 6.9.1), which is not needed here.}
Thus the map
\[
  u^0:G^0\longrightarrow H^0
\]
is surjective.

Conversely, suppose that \(u^0:G^0\to H^0\) is surjective, that
\(G\) is flat over \(S\), and that \(H\) has reduced fibers. Then,
for every \(s\in S\), the morphism
\[
  u_s^0:\underline{G_s^0}\longrightarrow\underline{H_s^0}
\]
is surjective, and hence is flat at every point lying above the
generic point \(\eta\) of \(\underline{H_s^0}\), since
\[
  \mathcal O_{\underline{H_s^0},\eta}
\]
is a field. It follows from \hyperref[prop:sga3-VIB-1.3]{1.3} that \(u_s\) is flat. Therefore
\(u\) is flat, by the fiberwise flatness criterion
(EGA~IV\(_3\), 11.3.11).

\section*{4. Dimension of the fibers of groups locally of finite presentation}
\phantomsection\label{sec:sga3-VIB-4}%
\addcontentsline{toc}{section}{4. Dimension of the fibers of groups locally of finite presentation}

\medskip
\phantomsection\label{prop:sga3-VIB-4.1}%
\noindent\textbf{Proposition 4.1.---}
\emph{Let \(G\) be an \(S\)-scheme locally of finite type, equipped
with an \(S\)-section \(\epsilon\), and suppose that, for every
\(s\in S\),}
\[
  \dim G_s=\dim_{\epsilon(s)}G_s
\]
\emph{\textup{(as is the case when \(G\) is an \(S\)-group;
\hyperref[lem:sga3-VIB-1.5]{1.5})}.}
\begin{enumerate}
\item[\textup{\SGAInlineRefTarget{prop:sga3-VIB-4.1:item:i}(i)}] \emph{The function}
  \[
    s\longmapsto\dim G_s
  \]
  \emph{is upper semicontinuous on \(S\).}
\item[\textup{\SGAInlineRefTarget{prop:sga3-VIB-4.1:item:ii}(ii)}] \emph{If, moreover, \(G\) is locally of finite
  presentation over \(S\), this function is locally constructible.}
\end{enumerate}

Let
\[
  \pi:G\longrightarrow S
\]
be the structural morphism. Chevalley's semicontinuity theorem
(EGA~IV\(_3\), 13.1.3) asserts that the function
\[
  x\longmapsto\dim_x\pi^{-1}\bigl(\pi(x)\bigr)
\]
is upper semicontinuous on \(G\). For every \(s\in S\), however,
\[
  \dim G_s
    =\dim\pi^{-1}(s)
    =\dim_{\epsilon(s)}
       \pi^{-1}\bigl(\pi(\epsilon(s))\bigr).
\]
Since the map \(s\mapsto\epsilon(s)\) is continuous on \(S\), the
composite function
\[
  s\longmapsto\dim G_s
\]
is upper semicontinuous on \(S\).

Now suppose that \(G\) is locally of finite presentation over
\(S\). Arguing as above, to show that \(s\mapsto\dim G_s\) is
locally constructible, it is enough to show that the function
\[
  x\longmapsto\dim_x\pi^{-1}\bigl(\pi(x)\bigr)
\]
is locally constructible on \(G\), which follows from
EGA~IV\(_3\), 9.9.1.

% END INCLUDED SGA3_English_current_buildable_source_20260823/source/inputs/VI/tex_reference_v2/components/37_expose_VIB_prop311_section4_prop41_en.tex

% BEGIN INCLUDED SGA3_English_current_buildable_source_20260823/source/inputs/VI/tex_reference_v2/components/38_expose_VIB_prop42_en.tex
% SGA 3, Expose VI_B, Loop 1 English reconstruction.
% Scope: Proposition 4.2, including source editor's notes 28--30.
% Authority: Polo--Gille Exp6B-13oct24.pdf, component-local pages
% 22--24 (combined-reader pages 363--365; original-margin printed
% pages 348--350).
% OCR and the Reinhold English Markdown are comparison witnesses only.

\medskip
\phantomsection\label{prop:sga3-VIB-4.2}%
\noindent\textbf{Proposition 4.2.---}
\emph{Let \(\pi:G\to S\) be an \(S\)-scheme locally of finite
presentation, equipped with an \(S\)-section \(\epsilon\), and
satisfying the following two conditions \textup{(which hold when
\(G\) is an \(S\)-group, by \hyperref[lem:sga3-VIB-1.5]{1.5} and VI\(_A\), \SGARefLink{cor:sga3-VIA-2.4.1}{2.4.1})}:}
\begin{enumerate}
\item[\textup{\SGAInlineRefTarget{prop:sga3-VIB-4.2:item:a}(a)}] \emph{For every \(s\in S\) and every
  \(x\in G_s\), one has}
  \[
    \dim G_s=\dim_x G_s
  \]
  \emph{\textup{(or equivalently \textup{(\hyperref[lem:sga3-VIB-1.5]{1.5})}, for every
  \(s\in S\), all the irreducible components of \(G_s\) have the
  same dimension)}.}
\item[\textup{(b)}] \emph{For every \(s\in S\), if \(\underline{G_s^0}\)
  denotes the connected component of \(G_s\) containing
  \(\epsilon(s)\), then \(\underline{G_s^0}\) is geometrically irreducible.}
\end{enumerate}
\emph{Let \(s\in S\). The following conditions are equivalent:}
\begin{enumerate}
\item[\textup{\SGAInlineRefTarget{prop:sga3-VIB-4.2:item:i}(i)}] \emph{\(G\) is universally open over \(S\) at
  the points of \(\underline{G_s^0}\).}
\item[\textup{(i bis)}] \emph{\(G\) is universally open over \(S\)
  at every point of a neighborhood of \(\epsilon(s)\) in
  \(\underline{G_s^0}\).}%
  \footnote{\label{note:sga3-VIB-28}Editor's note: the original
  reads ``at the points of \(\epsilon(s)\).'' It is not enough to
  suppose that \(G\) is universally open over \(S\) at
  \(\epsilon(s)\), as the example given in editor's \hyperref[note:sga3-VIB-19]{note 19} shows.
  We have corrected the proof accordingly.}
\item[\textup{\SGAInlineRefTarget{prop:sga3-VIB-4.2:item:ii}(ii)}] \emph{The function}
  \[
    t\longmapsto\dim G_t
  \]
  \emph{is constant in a neighborhood of \(s\) in \(S\).}
\item[\textup{(iii)}]\phantomsection\label{concl:sga3-VIB-4.2-iii} \emph{\(\underline{G^0}\) is ``universally open over
  \(S\) at the points of \(\underline{G_s^0}\),'' that is, given
  \(S'\to S\), \(s'\in {S'}_s\), \(g\in\underline{G_{s'}^0}\), and an open
  neighborhood \(V\) of \(g\) in \(G'=G_{S'}\), then}
  \[
    \pi\bigl(V\cap\underline{{G'}^0}\bigr)
  \]
  \emph{is an open neighborhood of \(s'\) in \(S'\).}%
  \footnote{\label{note:sga3-VIB-29}Editor's note: we have added
  condition \textup{\SGARefLink{concl:sga3-VIB-4.2-iii}{(iii)}}, pointed out by O.\ Gabber; it will be
  useful later \textup{(\SGARefLink{thm:sga3-VIB-5.7}{5.7})}.}
\end{enumerate}

Clearly (i) implies (i bis). By EGA~IV\(_3\), 14.3.3.1(ii), the set
of points of \(\underline{G_s^0}\) at which \(\pi_G\) is universally open is
closed in \(\underline{G_s^0}\). Since \(\underline{G_s^0}\) is irreducible,
it follows that
(i bis) implies \SGARefLink{prop:sga3-VIB-4.2:item:i}{(i)}.

We show that \SGARefLink{prop:sga3-VIB-4.2:item:i}{(i)} implies \SGARefLink{prop:sga3-VIB-4.2:item:ii}{(ii)}. Let \(T\) be the set of
\(t\in S\) such that
\[
  \dim G_t=\dim G_s.
\]
By 4.1, \(T\) is locally constructible; hence, by
EGA~IV\(_1\), 1.10.1, to show that \(T\) is a neighborhood of
\(s\), it is enough to show that every generalization \(s'\) of
\(s\) belongs to \(T\).

Let \(x\) be the generic point of \(\underline{G_s^0}\), and let \(U\) be an
open subset of \(G\) containing \(x\). Since \(\pi_G\) is
universally open at \(x\), EGA~IV\(_3\), 14.3.13 gives, for every
generalization \(s'\) of \(s\),
\[
  \dim\bigl(U\cap G_{s'}\bigr)
    \geq \dim_x\bigl(U\cap G_s\bigr).
\]
In view of hypothesis \SGARefLink{prop:sga3-VIB-4.2:item:a}{(a)}, this implies
\[
  \dim G_{s'}\geq\dim G_s.
\]
On the other hand, by \hyperref[prop:sga3-VIB-4.1]{4.1} the function
\(s\mapsto\dim G_s\) is upper semicontinuous, so that
\[
  \dim G_{s'}\leq\dim G_s.
\]
Thus \(s'\in T\), proving that \SGARefLink{prop:sga3-VIB-4.1:item:i}{(i)} implies \SGARefLink{prop:sga3-VIB-4.1:item:ii}{(ii)}.

It is clear that \SGARefLink{concl:sga3-VIB-4.2-iii}{(iii)} implies \SGARefLink{prop:sga3-VIB-4.1:item:i}{(i)}; we prove that \SGARefLink{prop:sga3-VIB-4.2:item:ii}{(ii)} implies
\SGARefLink{concl:sga3-VIB-4.2-iii}{(iii)}. The dimension of the fibers is unchanged by extension of
the base field, and the formation of \(\underline{G^0}\) commutes
with base change (see the proof of \hyperref[prop:sga3-VIB-3.3]{3.3}). We may therefore assume
that \(S'=S\) and \(s'=s\). Moreover, every open subset \(V\) of
\(G\) meeting \(\underline{G_s^0}\) contains the generic point \(\eta\) of
\(\underline{G_s^0}\), so we may assume that \(g=\eta\).

We may further suppose that \(S\) is affine. Let \(U\) be an
affine open neighborhood of \(\epsilon(s)\); it is then of finite
presentation over \(S\). Replacing \(S\) by \(\epsilon^{-1}(U)\),
and then \(U\) by \(U\cap\pi^{-1}(S)\), we reduce to the case in
which
\[
  \pi:U\longrightarrow S
\]
is of finite presentation and admits
\(\epsilon:S\to U\) as a section. By EGA~IV\(_3\), 9.7.12,
\(\underline{G^0}\cap U\) is constructible in \(U\). Replacing
\(V\) by an affine open subset contained in \(V\cap U\), we find
that \(\underline{G^0}\cap V\) is constructible in \(V\). Since
\(\pi:V\to S\) is of finite presentation, Chevalley's
constructibility theorem (EGA~IV\(_1\), 1.8.4) shows that
\[
  \pi\bigl(\underline{G^0}\cap V\bigr)
\]
is locally constructible in \(S\).

Consequently, by \emph{loc.\ cit.}, 1.10.1, to show that
\(\pi(\underline{G^0}\cap V)\) is an open neighborhood of \(s\),
it is enough to show that, for every generalization \(t\) of
\(s\), there is a generalization \(\xi\) of \(\eta\) belonging to
\(\underline{G^0}\), and hence to
\(\underline{G^0}\cap V\). The generic point \(\xi\) of
\(\underline{G_t^0}\)
is such a generalization of \(\eta\). Indeed, \(\epsilon(s)\)
belongs to the closure \(X\) of \(\{\xi\}\) in \(G\); hence,
by Chevalley's semicontinuity theorem (EGA~IV\(_3\), 13.1.3),
\[
  \dim_{\epsilon(s)}X_s\geq\dim_\xi X_t.
\]
On the other hand, hypothesis \SGARefLink{prop:sga3-VIB-4.2:item:ii}{(ii)} implies
\[
  \dim_\xi\underline{G_t^0}=\dim_{\epsilon(s)}G_s.
\]
It follows that one of the irreducible components of \(X_s\)
containing \(\epsilon(s)\) is equal to \(\underline{G_s^0}\), whence
\[
  \eta\in\overline{\{\xi\}}.
\]
This proves that \SGARefLink{prop:sga3-VIB-4.2:item:ii}{(ii)} implies \SGARefLink{concl:sga3-VIB-4.2-iii}{(iii)}, and hence proves the
proposition.

\footnote{\label{note:sga3-VIB-30}Editor's note: the preceding
argument was communicated to us by O.\ Gabber. We have also
retained the proof of the implication \textup{\SGARefLink{prop:sga3-VIB-4.2:item:ii}{(ii)}}
\(\Rightarrow\)\textup{\SGARefLink{prop:sga3-VIB-4.2:item:i}{(i)}} given in the original.}%
One can also prove that \SGARefLink{prop:sga3-VIB-4.2:item:ii}{(ii)} implies \SGARefLink{prop:sga3-VIB-4.2:item:i}{(i)} as follows. Since \(\pi\)
is locally of finite presentation, there exist an open subset
\(U\) of \(G\) containing \(\epsilon(s)\) and an open subset
\(V\) of \(S\) containing \(s\), such that \(\pi(U)\subset V\)
and the morphism
\[
  \pi':U\longrightarrow V
\]
induced by \(\pi\) is of finite presentation. Set
\[
  T=\epsilon^{-1}(U),
  \qquad
  W={\pi'}^{-1}(T)=U\cap\pi^{-1}(T).
\]
The morphism
\[
  \pi'':W\longrightarrow T
\]
induced by \(\pi'\) is of finite presentation and admits as a
section the morphism
\[
  \epsilon'':T\longrightarrow W
\]
induced by \(\epsilon\). Moreover, for every \(t\in T\), since
\(\underline{G_t^0}\) is irreducible,
\(W\cap\underline{G_t^0}\) is dense in \(\underline{G_t^0}\),
hence irreducible and therefore connected. It is consequently the
connected component of \({\pi''}^{-1}(t)\) containing
\(\epsilon''(t)\).

Since \(W\cap\underline{G_t^0}\) is a dense open subset of
\(\underline{G_t^0}\), \hyperref[lem:sga3-VIB-1.5]{1.5}
and EGA~IV\(_2\), 5.2.1 give
\[
  \dim(W\cap\underline{G_t^0})=\dim\underline{G_t^0}=\dim G_t.
\]
Thus the function
\[
  t\longmapsto\dim(W\cap\underline{G_t^0})
\]
is constant in a neighborhood of \(s\) in \(T\). Finally, we show
that \(W\cap\underline{G_t^0}\) is geometrically irreducible for every
\(t\in T\). Let \(K\) be an extension of \(\kappa(t)\). Then
\[
  (W\cap\underline{G_t^0})\otimes_{\kappa(t)}K
    =
  \bigl(W\otimes_{\kappa(t)}K\bigr)
    \cap
  \bigl(\underline{G_t^0}\otimes_{\kappa(t)}K\bigr)
\]
is a nonempty open subset of
\(\underline{G_t^0}\otimes_{\kappa(t)}K\), and is therefore irreducible
because \(\underline{G_t^0}\otimes_{\kappa(t)}K\) is irreducible.

We are now in the situation of EGA~IV\(_3\), 15.6.6(ii), which
asserts that \(\pi''\), and hence \(\pi\), is universally open at
the points of \(W\cap\underline{G_s^0}\). But, by
EGA~IV\(_3\), 14.3.3.1(ii),
the subset of \(G_s\) consisting of the points at which \(\pi\)
is universally open is closed in \(G_s\). Since it contains
\(W\cap\underline{G_s^0}\), it therefore contains its closure
\(\underline{G_s^0}\).

% END INCLUDED SGA3_English_current_buildable_source_20260823/source/inputs/VI/tex_reference_v2/components/38_expose_VIB_prop42_en.tex

% BEGIN INCLUDED SGA3_English_current_buildable_source_20260823/source/inputs/VI/tex_reference_v2/components/39_expose_VIB_cor43_cor44_en.tex
% SGA 3, Expose VI_B, Loop 1 English reconstruction.
% Scope: Corollaries 4.3 and 4.4 in full, including editor's notes
% 31 and 32.
% Authority: Polo--Gille Exp6B-13oct24.pdf, component-local pages
% 24--25 (combined-reader pages 365--366; original-margin printed
% pages 350--351), ending immediately before the Section 5 heading.
% OCR and the Reinhold English Markdown are comparison witnesses only.

\medskip
\phantomsection\label{cor:sga3-VIB-4.3}%
\noindent\textbf{Corollary 4.3.---}
\emph{Let \(G\) be a flat \(S\)-group locally of finite
presentation. Then the function}
\[
  s\longmapsto\dim G_s
\]
\emph{is locally constant on \(S\).}

This follows immediately from \SGARefLink{prop:sga3-VIB-4.2}{4.2}, since every flat morphism
locally of finite presentation is universally open
(EGA~IV\(_2\), 2.4.6).

\medskip
\phantomsection\label{cor:sga3-VIB-4.4}%
\noindent\textbf{Corollary 4.4.---}
\emph{Let \(G\) be an \(S\)-group locally of finite presentation
over \(S\) at the points of the identity section. Consider the
following conditions:}
\begin{enumerate}
\item[\textup{\SGAInlineRefTarget{cor:sga3-VIB-4.4:item:i}(i)}] \emph{\(G\) is smooth over \(S\) at the points
  of the identity section \textup{(cf.\ \hyperref[thm:sga3-VIB-3.10]{3.10})}.}
\item[\textup{\SGAInlineRefTarget{cor:sga3-VIB-4.4:item:ii}(ii)}]\phantomsection\label{cond:sga3-VIB-4.4-ii} \emph{For every \(s\in S\), \(G_s\) is smooth
  over \(\kappa(s)\), and the function}
  \[
    s\longmapsto\dim G_s
  \]
  \emph{is locally constant on \(S\).}
\item[\textup{(iii)}] \emph{For every \(s\in S\), \(G_s\) is smooth
  over \(\kappa(s)\), and there exists a neighborhood \(V\) of the
  identity section such that}
  \[
    \pi_V:V\longrightarrow S
  \]
  \emph{is universally open.}
\item[\textup{\SGAInlineRefTarget{cor:sga3-VIB-4.4:item:iv}(iv)}] \emph{For every \(s\in S\), \(G_s^0\) is
  smooth over \(\kappa(s)\), and \(G^0\) is representable by an
  open group subscheme of \(G\), universally open over \(S\).}
\end{enumerate}
\emph{Then the following implications hold:}
\[
  \textup{(i)}\Longrightarrow\textup{(ii)}
  \Longleftrightarrow\textup{(iii)}
  \Longleftrightarrow\textup{(iv)}.
\]
\emph{If, moreover, \(S\) is reduced, then conditions
\textup{\SGARefLink{cor:sga3-VIB-4.4:item:i}{(i)}}--\textup{\SGARefLink{cor:sga3-VIB-4.4:item:iv}{(iv)}} are equivalent and imply that \(G^0\)
is smooth over \(S\).}%
\footnote{\label{note:sga3-VIB-31}Editor's note: the hypothesis that
\(S\) be reduced cannot be omitted here. Indeed, let \(k\) be a
field, let \(S=\operatorname{Spec}A\), where
\(A=k[\delta]\) with \(\delta^2=0\), and let \(G\) be the closed
subgroup
\[
  \operatorname{Spec}A[X]/(\delta X)
\]
of \(\mathbb G_{a,S}\) (which assigns to every \(A\)-algebra \(R\)
the subgroup of those \(r\in R\) such that \(\delta r=0\)).
Then \(G\) is an \(S\)-group satisfying \textup{\SGARefLink{cor:sga3-VIB-4.4:item:ii}{(ii)}}--\textup{\SGARefLink{cor:sga3-VIB-4.4:item:iv}{(iv)}},
but it is not flat, and hence not smooth, over \(S\).}

Let us show that \textup{\SGARefLink{cor:sga3-VIB-4.4:item:i}{(i)}} implies \textup{\SGARefLink{cor:sga3-VIB-4.4:item:ii}{(ii)}}. For every
\(x\in G\), by \hyperref[lem:sga3-VIB-1.5]{1.5},
\[
  \dim_x\pi^{-1}\bigl(\pi(x)\bigr)
    =\dim\pi^{-1}\bigl(\pi(x)\bigr).
\]
Consequently, by EGA~IV\(_4\), 17.10.2, the function
\[
  x\longmapsto
  \dim_x\pi^{-1}\bigl(\pi(x)\bigr)
    =\dim\pi^{-1}\bigl(\pi(x)\bigr)
\]
is continuous in a neighborhood of the identity section.
Therefore the function \(s\mapsto\dim G_s\) is continuous on \(S\),
and hence locally constant on \(S\). On the other hand, for every
\(s\in S\), \(G_s\) is smooth over \(\kappa(s)\), by \SGARefLink{cor:sga3-VIB-1.3.1}{1.3.1}.

Let us show that \textup{\SGARefLink{cor:sga3-VIB-1.3.1:item:ii}{(ii)}} implies \textup{\SGARefLink{cor:sga3-VIB-4.4:item:iv}{(iv)}}. It is enough
to show that \(\underline{G^0}\) is open in \(G\), for then, by
\hyperref[case:sga3-VIB-3.4]{3.4}, \(G^0\) will be representable by the group subscheme induced
on \(\underline{G^0}\), and the properties of \(G^0\) stated above
will follow from \hyperref[cor:sga3-VIB-2.4]{2.4} and \hyperref[prop:sga3-VIB-4.2]{4.2}. Given \(s\in S\), construct, as in the
proof of \hyperref[thm:sga3-VIB-3.10]{3.10},
\[
  W,\quad T,\quad \pi'',\quad \epsilon'',\quad\text{and}\quad W^0.
\]
Then, by EGA~IV\(_3\), 15.6.7, \(W^0\) is open in \(W\). On the
other hand, under hypothesis \hyperref[cond:sga3-VIB-4.4-ii]{4.4}\textup{(ii)}, it follows from \hyperref[prop:sga3-VIB-4.2]{4.2}
that \(\pi\) is universally open at every point of \(W^0\).
Therefore, by VI\(_A\), \SGARefLink{sga3:VIA:0.1}{0.1}, \(\mu\) is universally open at every
point of
\[
  W^0\times_S W^0,
\]
which shows that \(W^0\cdot W^0\) is open; the proof is completed
as in the proof of \hyperref[thm:sga3-VIB-3.10]{Theorem~3.10}.

It is clear that
\[
  \textup{(iv)}\Longrightarrow\textup{(iii)},
\]
and \(\textup{(iii)}\Longrightarrow\textup{(ii)}\) follows by
applying \SGARefLink{prop:sga3-VIB-4.2}{4.2} to \(V\).

Finally, suppose that \textup{\SGARefLink{prop:sga3-VIB-4.2:item:ii}{(ii)}}--\textup{\SGARefLink{thm:sga3-VIB-3.10:item:iv}{(iv)}} hold, and let us
show that \(G^0\) is smooth over \(S\) when \(S\) is reduced.%
\footnote{\label{note:sga3-VIB-32}Editor's note: we have expanded
the original in the argument that follows.}
For this, we may suppose that \(G=G^0\). Then \(G\) is of finite
presentation over \(S\), by \SGARefLink{cor:sga3-VIB-5.5}{5.5} below. Thus \(\pi_G\) is of finite
presentation, has geometrically integral fibers, and has locally
constant relative dimension on \(S\). It follows from
EGA~IV\(_3\), 15.6.7 that the morphism
\[
  G\times_S S_{\mathrm{red}}\longrightarrow S_{\mathrm{red}}
\]
induced by \(\pi_G\) is flat. Hence \(\pi_G\) is flat when \(S\) is
reduced. In that case, \(G=G^0\) is smooth over \(S\), by
\hyperref[thm:sga3-VIB-3.10]{Theorem~3.10}.

% END INCLUDED SGA3_English_current_buildable_source_20260823/source/inputs/VI/tex_reference_v2/components/39_expose_VIB_cor43_cor44_en.tex

% BEGIN INCLUDED SGA3_English_current_buildable_source_20260823/source/inputs/VI/tex_reference_v2/components/40_expose_VIB_section5_through_rem541_loop2_r1_en.tex
% SGA 3, Expose VI_B, Loop 1 English reconstruction.
% Scope: Section 5 heading through Remark 5.4.1, including source
% editor's notes 33--36 and the source-footnote Theorem 5.3A.
% Authority: Polo--Gille Exp6B-13oct24.pdf, component-local pages
% 25--26 (combined-reader pages 366--367; original-margin printed
% pages 351--352 and re-edition headers 355--356).
% Corollary 5.5 begins on local page 27 / combined-reader page 368.
% OCR and Reinhold English Markdown are comparison witnesses only.
% The Proposition 5.1 diagram is an exact 180-dpi source crop:
% local p. 25, x=565,y=1080,w=365,h=180 on a 1488x2105 render.

\section*{5. Separation of groups and homogeneous spaces}
\phantomsection\label{sec:sga3-VIB-5}%
\addcontentsline{toc}{section}{5. Separation of groups and homogeneous spaces}

\medskip
\phantomsection\label{prop:sga3-VIB-5.1}%
\noindent\textbf{Proposition 5.1.---}
\emph{For an \(S\)-group \(G\) to be separated, it is necessary
and sufficient that the identity section of \(G\) be a closed
immersion.}

The condition is necessary (EGA~I, 5.4.6); it is sufficient by
virtue of the following cartesian diagram (cf. VI\(_A\), 0.3):
\begin{center}
  \phantomsection\label{diag:sga3-VIB-5.1-cartesian}%
  \begin{tikzcd}[row sep=large,column sep=huge]
    G\times_S G
      \arrow[r,"{\mu\circ(\operatorname{id}_G\times c)}"] &
    G \\
    G
      \arrow[u,"\Delta_{G/S}"]
      \arrow[r,"\pi"] &
    S
      \arrow[u,"\varepsilon"']
  \end{tikzcd}\,.
\end{center}

\medskip
\phantomsection\label{prop:sga3-VIB-5.2}%
\noindent\textbf{Proposition 5.2.---}
\emph{If \(S\) is discrete, every \(S\)-group is separated.}

Indeed,
\[
  S=\coprod_{s\in S}\operatorname{Spec}\mathcal O_{S,s},
\]
and, by EGA~I, 5.5.5, it is enough to show that, for every
\(s\in S\),
\[
  G\times_S\operatorname{Spec}\mathcal O_{S,s}
\]
is separated. This follows from VI\(_A\), \SGARefLink{sga3:VIA:0.3}{0.3}, since
\(\mathcal O_{S,s}\) is a local ring of dimension \(0\).%
\footnote{\label{note:sga3-VIB-33}Editor's note: we have replaced
``Artinian'' by ``of dimension \(0\)'' and mentioned this
generalization in VI\(_A\), \SGARefLink{sga3:VIA:0.3}{0.3}.}

\medskip
\phantomsection\label{thm:sga3-VIB-5.3}%
\noindent\textbf{Theorem 5.3.---}
\emph{Let \(S\) be a scheme, let \(G\) be an \(S\)-group scheme
locally of finite presentation over \(S\) and universally open over
\(S\) in a neighborhood of the identity section, and let \(X\) be
an \(S\)-scheme on which \(G\) acts in such a way that the morphism}
\[
  \Phi:G\times_S X\longrightarrow X\times_S X,
  \qquad (g,x)\longmapsto(gx,x),
\]
\emph{is surjective. Suppose moreover that, for every \(s\in S\):}
\begin{enumerate}
\item[\textup{(i)}] \emph{there exists an open subscheme \(U\) of
  \(X\), separated over \(S\), such that \(U_s\) is dense in
  \(X_s\);}
\item[\textup{\SGAInlineRefTarget{thm:sga3-VIB-5.3:item:ii}(ii)}] \emph{the fiber \(X_s\) is locally of finite
  type over \(\kappa(s)\).}%
\footnote{\label{note:sga3-VIB-34}Editor's note: we have corrected
the original by assuming that \(G\) is universally open over \(S\)
in a neighborhood of the identity section and by adding hypothesis
\textup{\SGARefLink{thm:sga3-VIB-5.3:item:ii}{(ii)}}; see below for examples, due to O.\ Gabber, showing
that the statements \hyperref[thm:sga3-VIB-5.3]{5.3} and \hyperref[cor:sga3-VIB-5.4]{5.4} of the original are false without
additional hypotheses. We also point out that \hyperref[thm:sga3-VIB-5.3]{Theorem 5.3} is a
reworked version of \hyperref[thm:sga3-VIB-5.3A]{Theorem 5.3A} below, which appears in the 1965
edition of SGAD and is due to M.\ Raynaud; compare the notes
\((*)\) in \SGARefLink{sga3:x:lemma:8-5}{Expos\'e~X, 8.5} and 8.8.

\smallskip
\phantomsection\label{thm:sga3-VIB-5.3A}%
\noindent\textbf{Theorem 5.3A (Raynaud).---}
\emph{Let \(G\) be an \(S\)-group locally of finite presentation,
universally open over \(S\), and with connected fibers. Then \(G\)
is separated over \(S\). More generally, every \(S\)-scheme \(X\)
locally of finite presentation over \(S\), equipped with an action
of \(G\) such that the morphism
\(\Phi:G\times_SX\to X\times_SX\),
\((g,x)\mapsto(gx,x)\), is surjective, is separated over \(S\).}}
\end{enumerate}
\emph{Then \(X\) is separated over \(S\).}

\medskip
\phantomsection\label{cor:sga3-VIB-5.4}%
\noindent\textbf{Corollary 5.4.---}
\emph{Let \(S,G,X\) be as in the preliminary hypotheses of \SGARefLink{thm:sga3-VIB-5.3}{5.3}.
Suppose moreover that \(X\) has fibers locally of finite type and
connected. Then:}
\begin{enumerate}
\item[\textup{\SGAInlineRefTarget{cor:sga3-VIB-5.4:item:i}(i)}] \emph{\(X\) is separated over \(S\).}
\item[\textup{(ii)}] \emph{If there exists an open subset \(V\) of
  \(X\), quasi-compact over \(S\), which meets every nonempty fiber
  \(X_s\),}%
\footnote{\label{note:sga3-VIB-35}Editor's note: without this
hypothesis there is the following counterexample, pointed out by
O.\ Gabber. Let \(G=S\) be a local scheme with closed point \(s\),
such that \(S^\ast=S-\{s\}\) is not quasi-compact, and let \(X\)
be the disjoint union of \(\{s\}\) and \(S^\ast\).}
\emph{then \(X\) is quasi-compact over \(S\).}
\end{enumerate}

\emph{Proof.}
\textup{\SGARefLink{cor:sga3-VIB-5.4:item:i}{(i)}} Let \(s\in S\) be such that \(X_s\ne\varnothing\).
The morphism
\[
  \Phi_s:
  G_s\times_{\kappa(s)}X_s
  \longrightarrow
  X_s\times_{\kappa(s)}X_s
\]
obtained from \(\Phi\) by base change is surjective. Since \(X_s\)
is connected, it is irreducible by VI\(_A\), \SGARefLink{prop:sga3-VIA-2.5.4}{2.5.4}. Thus, if \(U\)
is an affine open subset of \(X\) such that \(U_s\) is nonempty,
then \(U_s\) is dense in \(X_s\), and the theorem applies.

To prove \textup{\SGARefLink{concl:sga3-VIA-2.5.4-ii}{(ii)}}, we may suppose that \(S\) is affine. Then
\(V\) is quasi-compact and, by \hyperref[prop:sga3-VIB-3.5]{3.5}, there exists a quasi-compact
open subset \(U\) of \(G\) containing \(\underline{G^0}\). Let
\(s\in S\) be such that \(X_s\ne\varnothing\). Then \(X_s\) is
irreducible (VI\(_A\), \SGARefLink{prop:sga3-VIA-2.6.6}{2.6.6}), and, since \(U_s\) contains
\(\underline{G_s^0}\), the morphism
\[
  U_s\times_{\kappa(s)}V_s\longrightarrow X_s,
  \qquad (g,x)\longmapsto gx,
\]
is surjective (VI\(_A\), \SGARefLink{prop:sga3-VIA-2.6.4}{2.6.4}). Consequently the morphism
\[
  U\times_SV\longrightarrow X
\]
is surjective, and \(X\) is therefore quasi-compact, since \(U\)
and \(V\) are. As \(S\) was assumed affine, hence separated, it
follows that \(X\) is quasi-compact over \(S\)
(cf. EGA~I, 6.6.4\textup{(v)}).

\medskip
\phantomsection\label{rem:sga3-VIB-5.4.1}%
\noindent\textbf{Remark 5.4.1.---}%
\footnote{\label{note:sga3-VIB-36}Editor's note: we have added this
remark; compare editor's \hyperref[note:sga3-VIB-34]{note 34}.}
We shall see in the course of the proof that the conclusion of
\hyperref[thm:sga3-VIB-5.3]{Theorem 5.3} remains true if one assumes only:

\begin{quote}
\textup{(i')} there exists an open subscheme \(U\) of \(X\),
separated over \(S\), such that \(U_s\) is dense in every
irreducible component of \(X_s\).
\end{quote}

This follows from \textup{\SGARefLink{cor:sga3-VIB-5.4:item:i}{(i)}} if \(X_s\) locally has a finite
number of irreducible components, as is the case under hypothesis
\textup{(ii)}. We shall also see later that \SGARefLink{cor:sga3-VIB-5.4}{5.4} remains true under
the hypothesis that every fiber \(X_s\) is quasi-separated and
connected.

% END INCLUDED SGA3_English_current_buildable_source_20260823/source/inputs/VI/tex_reference_v2/components/40_expose_VIB_section5_through_rem541_loop2_r1_en.tex

% BEGIN INCLUDED SGA3_English_current_buildable_source_20260823/source/inputs/VI/tex_reference_v2/components/41_expose_VIB_cor55_through_lemma560_loop2_r1_en.tex
% SGA 3, Expose VI_B, Loop 1 English reconstruction.
% Scope: Corollary 5.5, paragraph 5.6 opening, and Lemma 5.6.0
% in full, including source editor's notes 37 and 38.
% Authority: Polo--Gille Exp6B-13oct24.pdf, component-local pages
% 27--28 (combined-reader pages 368--369; original-margin printed
% pages 352--353 and re-edition headers 357--358).
% Lemma 5.6.1.0 begins on local page 28 / combined page 369.
% OCR and Reinhold English Markdown are comparison witnesses only.
% Both diagrams are exact 180-dpi source crops for Loop 1.

\medskip
\phantomsection\label{cor:sga3-VIB-5.5}%
\noindent\textbf{Corollary 5.5.---}
\emph{Let \(S\) be a scheme, and let \(G\) be an \(S\)-group
scheme, locally of finite presentation, with connected fibers, and
universally open over \(S\). Then \(G\) is separated and of finite
presentation over \(S\).}%
\footnote{\label{note:sga3-VIB-37}Editor's note: we point out here
the following result ([Ray70a], VI, \SGARefLink{prop:sga3-VIB-2.5}{2.5}). If \(S\) is normal,
\(G\) is smooth with connected fibers, and \(X\) is an
\textup{(fppf)} homogeneous \(G\)-space---that is, the morphisms
\(X\to S\) and
\(G\times_SX\to X\times_SX\) are coverings for the
\textup{(fppf)} topology---locally of finite type over \(S\), then
\(X\) is locally quasi-projective over \(S\). In particular, \(G\)
is quasi-projective over \(S\). See also editor's \hyperref[note:sga3-VIB-35]{note 35} in
VI\(_A\).}

Indeed, by \SGARefLink{cor:sga3-VIB-3.6}{3.6} and \SGARefLink{cor:sga3-VIB-5.4}{5.4}, \(G\) is quasi-compact and separated over
\(S\). Since \(G\) is locally of finite presentation over \(S\),
it is therefore of finite presentation over \(S\).

\medskip
\phantomsection\label{par:sga3-VIB-5.6}%
\noindent\textbf{5.6.---}
\emph{Proof of \hyperref[thm:sga3-VIB-5.3]{Theorem 5.3}.}
Before proving \SGARefLink{thm:sga3-VIB-5.3}{5.3}, let us establish a few lemmas.

\medskip
\phantomsection\label{lem:sga3-VIB-5.6.0}%
\noindent\textbf{Lemma 5.6.0.---}%
\footnote{\label{note:sga3-VIB-38}Editor's note: we have added this
lemma, communicated by O.\ Gabber. It improves
EGA~IV\(_3\), 14.4.1.2 and corrects the proof of
\emph{loc.\ cit.}, 14.4.1.3 without changing its hypotheses
\textup{(compare the erratum \textup{(ErrIV, 38)} in
EGA~IV\(_4\))}.}
\begin{enumerate}
\item[\textup{\SGAInlineRefTarget{lem:sga3-VIB-5.6.0:item:i}(i)}] \emph{Let \(A\subset B\) be integral domains,
  with \(B\) integral over \(A\), and let
  \(\mathfrak q\in\operatorname{Spec}(B)\) be such that \(A\) is
  unibranch at
  \(\mathfrak p=\mathfrak q\cap A\). Then the morphism}
  \[
    \pi:\operatorname{Spec}(B)\longrightarrow\operatorname{Spec}(A)
  \]
  \emph{is open at \(\mathfrak q\).}
\item[\textup{\SGAInlineRefTarget{lem:sga3-VIB-5.6.0:item:ii}(ii)}] \emph{Let \(X,Y\) be irreducible preschemes,
  let \(f:X\to Y\) be a dominant morphism, and let \(x\in X\) be a
  point at which \(f\) is quasi-finite, such that \(y=f(x)\) is a
  unibranch point of \(Y\). Then \(f\) is open at \(x\). In
  particular, if}
  \[
    Y=\operatorname{Spec}(\mathcal O_{Y,y})
  \]
  \emph{is a local prescheme with closed point \(y\), then
  \(f(U)=Y\) for every neighborhood \(U\) of \(x\).}
\end{enumerate}

\textup{\SGARefLink{lem:sga3-VIB-5.6.0:item:i}{(i)}}
Let \(K\) (respectively, \(L\)) be the fraction field of \(A\)
(respectively, \(B\)), let \(A'\) be the normalization of \(A\),
and let \(B'\) be the subring of \(L\) generated by \(A'\) and
\(B\). Then \(B'\) is integral over \(A'\). Set
\[
  Y=\operatorname{Spec}(A),\quad
  X=\operatorname{Spec}(B),\quad
  Y'=\operatorname{Spec}(A'),\quad
  X'=\operatorname{Spec}(B').
\]
We then have the commutative diagram
\begin{center}
  \phantomsection\label{diag:sga3-VIB-5.6.0-integral-square}%
  \begin{tikzcd}[row sep=large,column sep=large]
    X' \arrow[r,"\pi'"] \arrow[d] &
    Y' \arrow[d,"\phi"'] \\
    X \arrow[r,"\pi"] &
    Y
  \end{tikzcd}
\end{center}
in which all morphisms are integral and surjective.

Since \(A\) is unibranch at \(\mathfrak p\), \(Y'\) has a unique
point \(\mathfrak p'\) above \(\mathfrak p\). Consequently, if
\(U\) is an open neighborhood of \(\mathfrak p'\) in \(Y'\), then
\[
  \phi(Y'-U)
\]
is a closed subset not containing \(\mathfrak p\), so that its
open complement is contained in \(\phi(U)\). This shows that
\(\phi:Y'\to Y\) is open at \(\mathfrak p'\). It is therefore
enough to show that \(\pi'\) is open, and we are reduced to the
case where \(A=A'\) is normal.

Let \(N\) be a quasi-Galois extension of \(K\) containing \(L\),
let
\[
  \widetilde X=\operatorname{Spec}(\widetilde B),
\]
where \(\widetilde B\) is the integral closure of \(B\) in \(N\),
and let \(G=\operatorname{Aut}(N/K)\). Since
\(\widetilde X\to X\) is surjective, it is enough to show that
\[
  \widetilde\pi:\widetilde X\longrightarrow Y
\]
is open. Let \(U\) be an open subset of \(\widetilde X\), and set
\[
  U'=\bigcup_{g\in G}gU.
\]
Because \(G\) acts transitively on the fibers of
\(\widetilde X\to Y\)
(cf. [BAC], Chap.~V, \S2.3, Prop.~6), one has
\[
  \widetilde\pi(U)=\widetilde\pi(U'),
\]
and the latter is the open complement of the closed subset
\[
  \widetilde\pi(\widetilde X-U').
\]
This proves \textup{\SGARefLink{lem:sga3-VIB-5.6.0:item:i}{(i)}}.

\textup{\SGARefLink{lem:sga3-VIB-5.6.0:item:ii}{(ii)}}
We may suppose that \(Y\) and \(X\) are reduced. By Zariski's Main
Theorem (cf. EGA~IV\(_3\), 8.12.9), there exist affine open
neighborhoods \(U\) of \(x\) and \(V=\operatorname{Spec}(A)\) of
\(y\), such that \(f(U)\subset V\), and a factorization
\begin{center}
  \phantomsection\label{diag:sga3-VIB-5.6.0-zariski-factorization}%
  \begin{tikzcd}[row sep=large,column sep=large]
    U
      \arrow[r,hook,"j"]
      \arrow[dr,"f"'] &
    V' \arrow[d,"u"] \\
    {} &
    V
  \end{tikzcd}\,,
\end{center}
where \(j\) is an open immersion and \(u\) is finite. Replacing
\(V'\) by the closure of \(j(U)\), we may suppose that \(V'\) is
irreducible, hence that \(V'=\operatorname{Spec}(B)\), where \(B\)
is an integral \(A\)-algebra, finite over \(A\). Moreover, since
\(f\) is dominant, so is \(u\), and the morphism \(A\to B\) is
therefore injective.

Since \(A\) is unibranch at \(y\), it follows from \textup{\SGARefLink{lem:sga3-VIB-5.6.0:item:i}{(i)}}
that \(u\) is open at \(j(x)\), and hence that
\[
  f=u\circ j
\]
is open at \(x\). This proves the first assertion of
\textup{\SGARefLink{lem:sga3-VIB-5.6.0:item:ii}{(ii)}}. The second follows because, if \(Y\) is a local
prescheme with closed point \(y\), every open subset containing
\(y\) equals \(Y\). In the case
\(Y=\operatorname{Spec}(\mathcal O_{Y,y})\), one may also use,
instead of EGA~IV\(_3\), 8.12.9, the local form of Zariski's Main
Theorem found, for example, in [Pes66] or [Ray70b], Chap.~IV,
Th.~1.

% END INCLUDED SGA3_English_current_buildable_source_20260823/source/inputs/VI/tex_reference_v2/components/41_expose_VIB_cor55_through_lemma560_loop2_r1_en.tex

% BEGIN INCLUDED SGA3_English_current_buildable_source_20260823/source/inputs/VI/tex_reference_v2/components/42_expose_VIB_lemma5610_loop2_r1_en.tex
% SGA 3, Expose VI_B, Loop 1 English reconstruction.
% Scope: Lemma 5.6.1.0 in full, including source editor's note 39,
% and the transition sentence with source editor's note 40.
% Authority: Polo--Gille Exp6B-13oct24.pdf, component-local pages
% 28--29 (combined-reader pages 369--370; re-edition headers
% 358--359). This editor-added unit ends immediately before the
% original-margin printed-page-353 marker at Lemma 5.6.1.
% Lemma 5.6.1 itself begins on local page 29.
% OCR and Reinhold English Markdown are comparison witnesses only.
% The diagram is an exact 180-dpi source crop for Loop 1.

\medskip
\phantomsection\label{lem:sga3-VIB-5.6.1.0}%
\noindent\textbf{Lemma 5.6.1.0.---}%
\footnote{\label{note:sga3-VIB-39}Editor's note: we have made
explicit this lemma, which is used in the proof of \hyperref[lem:sga3-VIB-5.6.1]{Lemma 5.6.1}.}
\emph{Let \(f:X\to S\) be a morphism locally of finite
presentation, let \(s\in S\), and let \(x\in X_s\). Set
\(n=\dim_x(X_s)\), and let
\(q:\mathbb A^n_S\to S\) be the structure morphism.}

\emph{Suppose that a quasi-finite morphism
\(u:X\to\mathbb A^n_S\) is given such that \(f=q\circ u\), and
suppose that \(f\) is universally open at the generic point \(z\)
of an irreducible component \(\mathfrak z\) of \(X_s\), containing
\(x\) and of dimension \(n\). Then \(u\) is universally open at
the point \(x\in X\).}

First observe that
\[
  (\dagger)\qquad
  \text{\emph{the hypotheses are preserved under every base change
  \(\pi:S'\to S\) covering \(s\)}}
\]
(that is, such that \(\pi^{-1}(s)\ne\varnothing\)).
Indeed, let \(S'\to S\) be such a morphism, let
\begin{center}
  \phantomsection\label{diag:sga3-VIB-5.6.1.0-base-change}%
  \begin{tikzcd}[row sep=large,column sep=large]
    X'
      \arrow[r,"u'"]
      \arrow[dr,"f'"'] &
    \mathbb A^n_{S'} \arrow[d,"q'"] \\
    {} &
    S'
  \end{tikzcd}
\end{center}
be the diagram obtained by base change, and let \(s'\) be a point
of \(S'\) above \(s\) and \(x'\) a point of \(X'_{s'}\) above
\(x\). Since
\[
  X'_{s'}=X_s\otimes_{\kappa(s)}\kappa(s'),
\]
the lifting of generizations and the invariance of dimension under
field extension (EGA~IV\(_2\), 2.3.4\textup{(i)} and 4.1.4) show
that \(x'\) is contained in an irreducible component
\(\mathfrak z'\) of \(X'_{s'}\) whose generic point \(z'\) lies
above \(z\), and
\[
  n\leq\dim\mathfrak z'
    \leq\dim_{x'}X'_{s'}\leq n.
\]
Thus
\[
  \dim\mathfrak z'=n=\dim_{x'}X'_{s'}.
\]
Since \(f\) is universally open at \(z\), \(f'\) is universally
open at \(z'\).%
\footnote[1]{\label{note:sga3-VIB-5.6.1.0-generic-points}Translator's
note: the French source prints \(x\) and
\(x'\), respectively, in this sentence. The stated hypothesis and
the immediately preceding choice of generic points require \(z\)
and \(z'\).}

Set \(Y=\mathbb A^n_S\). By EGA~IV\(_3\), 14.3.3.1\textup{(i)},
to prove that \(u\) is universally open at \(x\), it is enough to
prove that, for every integer \(r\geq0\) and every point \(x'\) of
\[
  X'=X[T_1,\ldots,T_r]
\]
above \(x\), the morphism
\[
  u':X'\longrightarrow
  Y'=\mathbb A^n_S[T_1,\ldots,T_r]
\]
is open at \(x'\). Now set
\[
  S'=S[T_1,\ldots,T_r],
\]
and let \(q'\) be the projection
\[
  Y'\simeq\mathbb A^n_{S'}\longrightarrow S'.
\]
This is the situation obtained by the base change \(S'\to S\).
Thus, by the preceding observation, we are reduced to showing that
\(u\) is open at \(x\).

Set \(y=u(x)\). Since \(u\) is locally of finite presentation
(because \(f\) and \(q\) are; cf. EGA~IV\(_1\), 1.4.3),
EGA~IV\(_1\), 1.10.3 shows that it is enough to prove
\[
  u\bigl(\operatorname{Spec}\mathcal O_{X,x}\bigr)
  =\operatorname{Spec}\mathcal O_{Y,y}.
\]
For this purpose we may assume that \(S\) is affine and integral.
Let \(\pi:S'\to S\) be its normalization, and denote by
\(u':X'\to Y'\) the morphism obtained from \(u\) by base change.
Since
\[
  \pi_Y:Y'\longrightarrow Y
\]
is integral and surjective, one has
\[
  \operatorname{Spec}\mathcal O_{Y,y}
  =
  \bigcup_{y'}
  \pi_Y\bigl(\operatorname{Spec}\mathcal O_{Y',y'}\bigr),
\]
where the union ranges over all points \(y'\) of \(Y'\) above
\(y\). It is therefore enough to show that, for every such
\(y'\) and every \(x'\in X'_{y'}\) above \(x\),
\[
  u'\bigl(\operatorname{Spec}\mathcal O_{X',x'}\bigr)
  =\operatorname{Spec}\mathcal O_{Y',y'}.
\]
Since the hypotheses are preserved under the base change
\(\pi:S'\to S\), we are thereby reduced to the case of \(S'\);
that is, we may assume that \(S\) is integral and normal.

The hypotheses on \(f\) now imply, by EGA~IV\(_3\), 14.3.13, that
there is an irreducible component \(Z\) of \(X\), containing
\(\mathfrak z\) (and hence \(x\)), dominating \(S\), and such that
\[
  \dim_z(Z_s)=n=\dim(Z_\eta),
\]
where \(\eta\) is the generic point of \(S\). Let \(\xi\) be the
generic point of \(Z\). Since \(u\), and hence also
\[
  u_\eta:Z_\eta\longrightarrow\mathbb A^n_\eta,
\]
is quasi-finite, the closure in \(\mathbb A^n_\eta\) of the point
\[
  u_\eta(\xi)=u(\xi)
\]
has dimension \(n\). Thus \(u(\xi)\) is the generic point of
\(\mathbb A^n_\eta\), which is also the generic point of
\(Y=\mathbb A^n_S\). Consequently, if \(g\) denotes the
restriction of \(u\) to \(Z\), then
\[
  g:Z\longrightarrow Y
\]
is quasi-finite and dominant. Since \(Y\) is normal, \hyperref[lem:sga3-VIB-5.6.0]{Lemma 5.6.0}
shows that \(g\) is open. Hence \(u\) is open at every point of
\(Z\), and in particular at \(x\). Consequently,
\[
  u\bigl(\operatorname{Spec}\mathcal O_{X,x}\bigr)
  =\operatorname{Spec}\mathcal O_{Y,y},
\]
which completes the proof of \hyperref[lem:sga3-VIB-5.6.1.0]{Lemma 5.6.1.0}.

The following lemma advantageously replaces
EGA~IV\(_3\), 14.5.10,%
\footnote{\label{note:sga3-VIB-40}Editor's note: we have corrected
the original, which indicated 19.5.10.}
in that it is independent of noetherian hypotheses.

% END INCLUDED SGA3_English_current_buildable_source_20260823/source/inputs/VI/tex_reference_v2/components/42_expose_VIB_lemma5610_loop2_r1_en.tex

% BEGIN INCLUDED SGA3_English_current_buildable_source_20260823/source/inputs/VI/tex_reference_v2/components/43_expose_VIB_lemma561_loop2_r1_en.tex
% SGA 3, Expose VI_B, Loop 1 English reconstruction.
% Scope: Lemma 5.6.1 in full, including source editor's note 41.
% Authority: Polo--Gille Exp6B-13oct24.pdf, component-local pages
% 29--31 (combined-reader pages 370--372; original-margin printed
% pages 353--354 and re-edition headers 359--361).
% Lemma 5.6.2.0 begins on local page 31 / combined page 372.
% OCR and Reinhold English Markdown are comparison witnesses only.
% All four diagrams are exact 180-dpi source crops for Loop 1.

\medskip
\phantomsection\label{lem:sga3-VIB-5.6.1}%
\noindent\textbf{Lemma 5.6.1.---}
\emph{Let \(S\) be a scheme, let \(f:X\to S\) be an \(S\)-scheme
locally of finite presentation, let \(s\) be a point of \(S\), and
let \(x\) be a closed point of \(X_s\). Suppose that \(f\) is
universally open at the generic point \(z\) of an irreducible
component \(Z\) of \(X_s\) containing \(x\), and that}
\[
  \dim_x(X_s)=\dim Z.
\]
\emph{Then there exists a commutative diagram}
\begin{center}
  \phantomsection\label{diag:sga3-VIB-5.6.1-statement}%
  \begin{tikzcd}[row sep=large,column sep=large]
    X \arrow[r,"f"] &
    S \\
    S''
      \arrow[u,"h"]
      \arrow[r,"\pi"]
      \arrow[ur,"\phi"] &
    S' \arrow[u,"w"']
  \end{tikzcd}\enspace,
\end{center}
\emph{where \(S'\) is affine, \(w\) is \'etale, \(\pi\) is finite,
surjective, and of finite presentation, and
\(\phi^{-1}(s)\) consists of a single point \(s''\), such that
\(h(s'')=x\) and \(\kappa(s'')=\kappa(x)\).}%
\footnote{\label{note:sga3-VIB-41}Editor's note: consequently,
\(h\) induces a section \(\sigma\) of
\(X\times_S S''\to S''\) such that \(\sigma(s'')\) lies above
\(x\); compare EGA~IV\(_3\), 14.5.10.}

\emph{Proof.}
First, we may suppose that \(S=\operatorname{Spec}(A)\) and
\(X=\operatorname{Spec}(B)\), where \(B\) is an \(A\)-algebra of
finite presentation. Let \(\mathfrak p\) and \(\mathfrak q\) be
the prime ideals of \(A\) and \(B\) corresponding respectively to
\(s\) and \(x\), so that
\[
  \mathfrak p=A\cap\mathfrak q.
\]
Set \(n=\dim_x(X_s)\), and choose \(t_1,\ldots,t_n\in B\) whose
images in
\[
  \mathcal O_{X_s,x}=B_{\mathfrak q}/\mathfrak p B_{\mathfrak q}
\]
form a system of parameters. Then
\[
  \mathcal O_{X_s,x}/(t_1,\ldots,t_n)
\]
is finite-dimensional over \(\kappa(x)\), and hence also over
\(\kappa(s)\), because \(x\) is a closed point of the algebraic
\(\kappa(s)\)-scheme \(X_s\). Consequently, the \(S\)-morphism
\[
  u:X\longrightarrow\mathbb A^n_S
    =\operatorname{Spec}\bigl(A[T_1,\ldots,T_n]\bigr),
  \qquad T_i\longmapsto t_i,
\]
is of finite presentation, \(x\) is isolated in the fiber
\(u^{-1}(u(x))\), and
\[
  u(x)=\tau_0(s),
\]
where \(\tau_0:S\to\mathbb A^n_S\) denotes the zero section of
\(\mathbb A^n_S\to S\), corresponding to the \(A\)-algebra
homomorphism
\[
  A[T_1,\ldots,T_n]\longrightarrow A,\qquad T_i\longmapsto0.
\]
The set of points of \(X\) that are isolated in their fibers over
\(\mathbb A^n_S\) is open (EGA~IV\(_3\), 13.1.4). Thus, after
shrinking \(X\), we may suppose that \(u\) is quasi-finite and
that \(u^{-1}(u(x))=\{x\}\). By \hyperref[lem:sga3-VIB-5.6.1.0]{Lemma 5.6.1.0}, \(u\) is
universally open at \(x\).

Let
\[
  B_0=B/(t_1,\ldots,t_n),\qquad
  X_0=X\times_{\mathbb A^n_S}\tau_0(S)
      =\operatorname{Spec}(B_0),
\]
and let \(u_0:X_0\to S\) be obtained from \(u\) by the base change
\(\tau_0:S\hookrightarrow\mathbb A^n_S\). Then \(u_0\) is
quasi-finite and of finite presentation, is universally open at
\(x\), and \(x\) is the unique point of \(X_0\) above \(s\).

Let \(A'\) be the henselization of the local ring
\[
  A_{\mathfrak p}=\mathcal O_{S,s},
\]
and set
\[
  S'=\operatorname{Spec}(A'),\qquad
  B'_0=B_0\otimes_A A',\qquad
  X'_0=X_0\times_S S'=\operatorname{Spec}(B'_0).
\]
The closed point \(s'\) of \(S'\) is the unique point of \(S'\)
above \(s\), and \(\kappa(s')=\kappa(s)\). The scheme \(X'_0\)
has a unique point \(x'\) above \(s'\); moreover
\(\kappa(x')=\kappa(x)\), and \(x'\) is also the unique point of
\(X'_0\) above both \(s\) and \(x\). Since \(A'\) is henselian,
EGA~IV\(_4\), 18.5.11 gives a disjoint decomposition into open and
closed parts
\[
  \tag{*}\label{sga3:vi:lemma:5-6-1:equation:star}
  X'_0=V\amalg W,
  \qquad
  V=\operatorname{Spec}(\mathcal O_{X'_0,x'}).
\]
The local ring \(\mathcal O_{X'_0,x'}\) is finite and of finite
presentation over \(A'\). Hence the restriction \(\pi\) of
\(u'_0\) to \(V\) is finite and of finite presentation. Moreover,
because \(u'_0:X'_0\to S'\) is open at \(x'\),
\[
  u'_0(V)=\operatorname{Spec}(\mathcal O_{S',s'})=S',
\]
so \(\pi\) is surjective.

This proves the assertion when \(S=S'\). In general, \(A'\) is
the filtered direct limit of \(A\)-subalgebras \(A_i\) that are
\'etale over \(A\) and for which
\[
  S_i=\operatorname{Spec}(A_i)
\]
has a unique point \(s_i\) above \(s\), with
\(\kappa(s_i)=\kappa(s)\). We have
\[
  B'_0=\varinjlim_i B_i,
  \qquad
  B_i=B_0\otimes_A A_i.
\]
Set
\[
  X_i=\operatorname{Spec}(B_i)=X_0\times_S S_i,
  \qquad
  C=\mathcal O_{X'_0,x'}.
\]
By \((*)\), one has
\[
  C\simeq B'_0/fB'_0
\]
for some idempotent \(f\in B'_0\). There are an index \(i\) and
an element \(f_i\in B_i\) whose image in \(B'_0\) is \(f\). Put
\[
  C_i=B_i/f_iB_i,\qquad V_i=\operatorname{Spec}(C_i).
\]
Then
\[
  C=C_i\otimes_{A_i}A',
  \qquad
  V=V_i\times_{S_i}S',
\]
and \(C_i\) is an \(A_i\)-algebra of finite presentation, since
the same is true of \(B_i\). Consequently, the morphisms
\begin{center}
  \phantomsection\label{diag:sga3-VIB-5.6.1-limit-source}%
  \begin{tikzcd}[row sep=large,column sep=large]
    X'_0 \arrow[dr,"u'_0"] &
    {} \\
    V
      \arrow[u,"h'"]
      \arrow[r,"\pi"] &
    S'
  \end{tikzcd}
\end{center}
arise, by the base change \(S'\to S_i\), from morphisms of finite
presentation
\begin{center}
  \phantomsection\label{diag:sga3-VIB-5.6.1-limit-descent}%
  \begin{tikzcd}[row sep=large,column sep=large]
    X_i \arrow[dr,"u_i"] &
    {} \\
    V_i
      \arrow[u,"h_i"]
      \arrow[r,"\pi_i"] &
    S_i
  \end{tikzcd}
\end{center}

For every \(j\geq i\), let
\[
  V_j=V_i\times_{S_i}S_j
\]
and let \(\pi_j:V_j\to S_j\) be the morphism of finite
presentation obtained from \(\pi_i\) by base change. Since
\(\pi:V\to S'\) is finite and surjective, EGA~IV\(_2\), 8.10.5
shows that, for some \(j\), the morphism
\(\pi_j:V_j\to S_j\) is finite and surjective. Then
\(w:S_j\to S\) is affine and \'etale, \(S_j\) has a unique point
\(s_j\) above \(s\), and \(V_j\) has a unique point \(x_j\) above
\(s_j\), since \(x'\) is the unique point of
\[
  V=V_j\times_{S_j}S'
\]
above \(s'\):
\begin{center}
  \phantomsection\label{diag:sga3-VIB-5.6.1-final-composite}%
  {\small
  \begin{tikzcd}[row sep=large,column sep=large]
    X_j
      \arrow[r]
      \arrow[dr,"u_j"] &
    X_0
      \arrow[r]
      \arrow[dr,"u_0"] &
    X
      \arrow[d,"f"']
      \arrow[dr,"u"] &
    {} \\
    V_j
      \arrow[u,"h_j"]
      \arrow[r,"\pi_j"] &
    S_j
      \arrow[r,"w"] &
    S
      \arrow[r,"\tau_0"] &
    \mathbb A^n_S
  \end{tikzcd}}
\end{center}
Thus \(x_j\) is the unique point of \(V_j\) above \(s\), its image
under \(V_j\to X_j\to X\) is \(x\), and
\[
  \kappa(x_j)=\kappa(x).
\]

% END INCLUDED SGA3_English_current_buildable_source_20260823/source/inputs/VI/tex_reference_v2/components/43_expose_VIB_lemma561_loop2_r1_en.tex

% BEGIN INCLUDED SGA3_English_current_buildable_source_20260823/source/inputs/VI/tex_reference_v2/components/44_expose_VIB_lemma5620_en.tex
% SGA 3, Expose VI_B, Loop 1 English reconstruction.
% Scope: Lemma 5.6.2.0 in full, including source editor's note 42.
% Authority: Polo--Gille Exp6B-13oct24.pdf, component-local pages
% 31--32 (combined-reader pages 372--373; original-margin printed
% page 354 and re-edition headers 361--362).
% Corollary 5.6.2 begins on local page 32 / combined page 373.
% OCR and Reinhold English Markdown are comparison witnesses only.
% This unit contains no diagrams.

\medskip
\phantomsection\label{lem:sga3-VIB-5.6.2.0}%
\noindent\textbf{Lemma 5.6.2.0.---}%
\footnote{\label{note:sga3-VIB-42}Editor's note: we have added this
lemma, communicated by O.\ Gabber. It makes it possible to simplify
the proof of \SGARefLink{cor:sga3-VIB-5.6.2}{5.6.2} and to prove \hyperref[thm:sga3-VIB-5.3]{Theorem 5.3}, as well as \hyperref[cor:sga3-VIB-5.4]{5.4}, in a
more general form; see \hyperref[thm:sga3-VIB-5.7]{5.7} and \hyperref[sec:sga3-VIB-5.8]{5.8} below.}
\emph{Let \(k\) be a field and let \(G\) be a nonempty
\(k\)-scheme of finite type.}
\begin{enumerate}
\item[\textup{\SGAInlineRefTarget{lem:sga3-VIB-5.6.2.0:item:i}(i)}] \emph{Let \(K\) be an extension of \(k\), and
  let \(W\) be a dense open subset of \(G_K\). Denote by \(\pi\)
  the projection \(G_K\to G\). Then}
  \[
    U=\bigl\{g\in G\mathrel{\big|}W\text{ contains every maximal
      point of }\pi^{-1}(g)\bigr\}
  \]
  \emph{is a dense open subset of \(G\). \textup{(N.B. If \(g\)
  is a closed point of \(G\) belonging to \(U\), then
  \(\pi^{-1}(g)\subset W\).}}

\item[\textup{\SGAInlineRefTarget{lem:sga3-VIB-5.6.2.0:item:ii}(ii)}] \emph{Suppose in addition that \(G\) is
  geometrically irreducible. Let}
  \[
    \mu:G\times X\longrightarrow Y
  \]
  \emph{be a morphism of \(k\)-schemes, let \(x\) be a point of
  \(X\), and let \(\Omega\) be an open subset of \(Y\) such that}
  \[
    \mu_x^{-1}\bigl(\Omega_{\kappa(x)}\bigr)\ne\varnothing,
  \]
  \emph{where \(\mu_x\) denotes the morphism}
  \[
    G_{\kappa(x)}\longrightarrow Y_{\kappa(x)},
    \qquad g\longmapsto\mu(g,x).
  \]
  \emph{Then}
  \[
    U=\bigl\{g\in G\mathrel{\big|}\mu\text{ sends every maximal
      point of }g\times x\text{ into }\Omega\bigr\}
  \]
  \emph{is a dense open subset of \(G\). For every closed point
  \(g\) of \(G\) belonging to \(U\), one has}
  \[
    \mu(g\times x)\subset\Omega,
  \]
  \emph{and hence}
  \[
    \mu(g',x)\in\Omega_{\kappa(x)}
    \quad\text{\textup{(respectively, }}\,
    \mu(g,x')\in\Omega_{\kappa(g)}\text{\textup{)}},
  \]
  \emph{for every point \(g'\in G_{\kappa(x)}\) above \(g\)
  \textup{(respectively, \(x'\in X_{\kappa(g)}\) above \(x\))}.}

\item[\textup{\SGAInlineRefTarget{lem:sga3-VIB-5.6.2.0:item:iii}(iii)}] \emph{Suppose that \(G=H^0\), where \(H\)
  is a \(k\)-group scheme locally of finite type, acting on a
  nonempty \(k\)-scheme \(X\) in such a way that the morphism}
  \[
    H\times_S X\longrightarrow X\times_S X,
    \qquad (h,x)\longmapsto(hx,x)
  \]
  \emph{is surjective. Let \(U\) be an open subset of \(X\).
  Then:}
  \begin{enumerate}
  \item[\textup{\SGAInlineRefTarget{lem:sga3-VIB-5.6.2.0:item:a}(a)}] \emph{\(G\cdot U\) is an open subset of
    \(X\), equal to the union of the irreducible components of
    \(X\) whose generic point belongs to \(U\).}
  \item[\textup{\SGAInlineRefTarget{lem:sga3-VIB-5.6.2.0:item:b}(b)}] \emph{Suppose that \(U\) is dense in every
    irreducible component of \(X\). Then, for every finite subset
    \(A\) of \(X\), there exists a closed point \(g\in G\) such
    that}
    \[
      g\cdot A'\subset U_{\kappa(g)},
    \]
    \emph{where \(A'\) is the inverse image of \(A\) in
    \(X_{\kappa(g)}\). In particular, the morphism}
    \[
      G\times U\longrightarrow X
    \]
    \emph{is surjective.}
  \end{enumerate}
\end{enumerate}

\emph{Proof.}
\textup{\SGARefLink{lem:sga3-VIB-5.6.2.0:item:i}{(i)}}
Let \(\bar{k}\) be a separable closure of \(k\), and let \(L\) be
an extension of \(k\) containing a copy of \(\bar{k}\) and of
\(K\). Denote by \(\pi_L\) (respectively, \(\phi\)) the projection
\[
  G_L\longrightarrow G
  \qquad\text{(respectively, \(G_L\longrightarrow G_K\)).}
\]
For every \(g\in G\), the map \(\phi\) sends the maximal points of
\(\pi_L^{-1}(g)\) surjectively onto those of \(\pi^{-1}(g)\).
Consequently,
\[
  U=\bigl\{g\in G\mathrel{\big|}W_L\text{ contains every maximal
    point of }\pi_L^{-1}(g)\bigr\}.
\]
Thus, replacing \(K\) by \(L\), we are reduced to the case in
which \(K\) contains \(\bar{k}\).

Since the projection
\[
  p:G_K\longrightarrow G_{\bar{k}}
\]
is surjective and open, \(V=p(W)\) is a dense open subset of
\(G_{\bar{k}}\). Moreover,
\[
  p^{-1}(g)=\operatorname{Spec}
    \bigl(\kappa(g)\otimes_{\bar{k}}K\bigr)
\]
is irreducible for every \(g\in G_{\bar{k}}\) (cf.\
EGA~IV\(_2\), 4.3.3). Hence, for every \(g\in V\), the generic
point of \(p^{-1}(g)\) belongs to \(W\).

On the other hand, set
\[
  \mathscr G=\operatorname{Gal}(\bar{k}/k).
\]
The projection \(q:G_{\bar{k}}\to G\) is surjective, and
\(\mathscr G\) acts transitively on the fibers of \(q\).
Therefore
\[
  U=q(V'),
\]
where \(V'\) is the intersection of the
\(\mathscr G\)-conjugates of \(V\).

Let \(Z\) be the closed subset \(G_{\bar{k}}-V\), endowed with
its reduced closed subscheme structure. Since \(G_{\bar{k}}\),
and hence \(Z\), is of finite presentation over \(\bar{k}\), the
scheme \(Z\) is obtained by base change from a reduced closed
subscheme
\[
  Z_1\subset G\otimes_k k_1
\]
for some finite Galois extension \(k_1\) of \(k\). Thus the
\(\mathscr G\)-conjugates of \(Z\) are finite in number, and their
union is again a nowhere-dense closed subset \(F\) of
\(G_{\bar{k}}\). It follows that
\[
  V'=G_{\bar{k}}-F
\]
is a dense open subset of \(G_{\bar{k}}\). Since
\(q:G_{\bar{k}}\to G\) is surjective and open, \(U=q(V')\) is
therefore a dense open subset of \(G\).

Moreover, for every closed point \(g\) of \(G\), the fiber
\(\pi^{-1}(g)\) consists of finitely many closed points of \(G_K\).
Thus, if \(g\in U\), then
\[
  \pi^{-1}(g)\subset W.
\]
This proves \textup{\SGARefLink{lem:sga3-VIB-5.6.2.0:item:i}{(i)}}, and \textup{\SGARefLink{lem:sga3-VIB-5.6.2.0:item:ii}{(ii)}} follows from it.

Assertion \textup{\SGARefLink{lem:sga3-VIB-5.6.2.0:item:iii}{(iii)}\SGARefLink{lem:sga3-VIB-5.6.2.0:item:a}{(a)}} was proved in VI\(_A\), \SGARefLink{prop:sga3-VIA-2.6.4}{2.6.4}.
Finally, if \(U\) is dense in every irreducible component of
\(X\), then
\[
  X=G\cdot U.
\]
Hence, for every \(x\in X\),
\[
  \mu_x^{-1}\bigl(U_{\kappa(x)}\bigr)
\]
is a nonempty open subset of \(G_{\kappa(x)}\), and
\textup{\SGARefLink{lem:sga3-VIB-5.6.2.0:item:iii}{(iii)}\SGARefLink{lem:sga3-VIB-5.6.2.0:item:b}{(b)}} follows from \textup{\SGARefLink{lem:sga3-VIB-5.6.2.0:item:ii}{(ii)}}.

% END INCLUDED SGA3_English_current_buildable_source_20260823/source/inputs/VI/tex_reference_v2/components/44_expose_VIB_lemma5620_en.tex

% BEGIN INCLUDED SGA3_English_current_buildable_source_20260823/source/inputs/VI/tex_reference_v2/components/45_expose_VIB_cor562_en.tex
% SGA 3, Expose VI_B, Loop 1 English reconstruction.
% Scope: Corollary 5.6.2 in full, including source editor's note 43.
% Authority: Polo--Gille Exp6B-13oct24.pdf, component-local pages
% 32--33 (combined-reader pages 373--374; original-margin printed
% page 355 and re-edition headers 362--363).
% Lemma 5.6.3 begins on local page 33 / combined page 374.
% OCR and Reinhold English Markdown are comparison witnesses only.
% This unit contains no diagrams.

\medskip
\phantomsection\label{cor:sga3-VIB-5.6.2}%
\noindent\textbf{Corollary 5.6.2.---}
\emph{Let \(S\), \(G\), and \(X\) be as in the preliminary
hypotheses of \hyperref[thm:sga3-VIB-5.3]{5.3}, and let \(U\) be an open subset of \(X\), let
\(s\in S\), and let \(A\) be a finite subset of \(X_s\). Suppose
that \(U_s\) is dense in \(X_s\) and that \(X_s\) is locally of
finite type over \(\kappa(s)\).}%
\footnote{\label{note:sga3-VIB-43}Editor's note: we have added the
hypothesis on \(X_s\) and simplified (and corrected) the proof,
taking account of the addition of \SGARefLink{lem:sga3-VIB-5.6.2.0}{5.6.2.0}. Moreover, the proof
shows that the conclusion remains valid if one assumes only that
\(U_s\) is dense in every irreducible component of \(X_s\).}

\emph{Then there are morphisms \(f:S''\to S\) and \(h:S''\to G\),
where \(f\) factors as a finite surjective morphism
\(S''\to S'\) followed by an \'etale morphism \(S'\to S\), such
that the inverse image \(A''\) of \(A\) in \(X\times_S S''\)
\textup{(i.e., in
\(X_s\times_{\operatorname{Spec}\kappa(s)}S''_s\))} is
contained in}
\[
  \ell_h^{-1}(U_{S''}),
\]
\emph{where \(\ell_h\) denotes the translation of
\(X_{S''}=X\times_S S''\) defined by the element
\(h\in G(S'')\).}

\emph{Proof.}
Since \(X_s\) is locally of finite type over \(\kappa(s)\), the
connected components of \(X_s\) are open and irreducible (cf.\
VI\(_A\), \SGARefLink{prop:sga3-VIA-2.5.4}{2.5.4}). Hence \(U_s\) is dense in every irreducible
component of \(X_s\).

By \hyperref[lem:sga3-VIB-5.6.2.0]{Lemma 5.6.2.0}, there therefore exists a closed point
\[
  g\in G_s^0
\]
such that
\[
  g\cdot a'\in U_s\otimes_{\kappa(s)}\kappa(g)
\]
for every \(a'\in X_{\kappa(g)}\) lying above a point of \(A\).

By \hyperref[lem:sga3-VIB-5.6.1]{Lemma 5.6.1}, there are morphisms \(f:S''\to S\) and
\(h:S''\to G\), where \(f\) factors as a finite surjective
morphism \(S''\to S'\) followed by an \'etale morphism
\(S'\to S\), such that \(f^{-1}(s)\) consists of a single point
\(s_0\), and such that
\[
  h(s_0)=g,\qquad \kappa(s_0)=\kappa(g).
\]
Since \(S''_s\) is finite and has the single point \(s_0\), the
canonical closed immersion
\[
  X_s\otimes_{\kappa(s)}\kappa(s_0)
  \longrightarrow
  X_s\times_{\operatorname{Spec}\kappa(s)}S''_s
\]
is nilpotent and hence induces a homeomorphism of underlying spaces.
\footnote{Canon source note (SGA-SRC-000863): the re-edition identifies
these schemes. A finite one-point fibre need not be reduced; the
canonical map from its residue field is instead a nilpotent closed
immersion, which is exactly sufficient for this point-set argument.}
Denote by \(A''\) the inverse image of \(A\) in the underlying space
of the right-hand side. Via this homeomorphism and
\(\kappa(s_0)=\kappa(g)\), it is identified with the inverse image of
\(A\) in
\[
  X_s\otimes_{\kappa(s)}\kappa(s_0)
  =X_s\otimes_{\kappa(s)}\kappa(g).
\]
Then the translation \(\ell_h\) of \(X_{S''}\) sends \(A''\)
into
\[
  U_s\otimes_{\kappa(s)}\kappa(s_0).
\]

% END INCLUDED SGA3_English_current_buildable_source_20260823/source/inputs/VI/tex_reference_v2/components/45_expose_VIB_cor562_en.tex

% BEGIN INCLUDED SGA3_English_current_buildable_source_20260823/source/inputs/VI/tex_reference_v2/components/46_expose_VIB_lemma563_en.tex
% SGA 3, Expose VI_B, Loop 1 English reconstruction.
% Scope: Lemma 5.6.3 in full, including source editor's note 44,
% and the following deduction of Theorem 5.3.
% Authority: Polo--Gille Exp6B-13oct24.pdf, component-local pages
% 33--34 (combined-reader pages 374--375; original-margin printed
% pages 355--356 and re-edition headers 363--364).
% Counterexample 5.6.4 begins on local page 34 / combined page 375.
% OCR and Reinhold English Markdown are comparison witnesses only.
% This unit contains no diagrams.

\medskip
\phantomsection\label{lem:sga3-VIB-5.6.3}%
\noindent\textbf{Lemma 5.6.3.---}
\emph{Let \(X\) be an \(S\)-scheme. The following conditions are
equivalent:}
\begin{enumerate}
\item[\textup{\SGAInlineRefTarget{lem:sga3-VIB-5.6.3:item:i}(i)}]\phantomsection\label{cond:sga3-VIB-5.6.3-i}
  \emph{\(X\) is separated over \(S\).}

\item[\textup{\SGAInlineRefTarget{lem:sga3-VIB-5.6.3:item:ii}(ii)}]\phantomsection\label{cond:sga3-VIB-5.6.3-ii}
  \emph{For every \(S\)-scheme \(T\), every section
  \(\sigma:T\to X_T\) is a closed immersion.}

\item[\textup{\SGAInlineRefTarget{lem:sga3-VIB-5.6.3:item:iii}(iii)}]\phantomsection\label{cond:sga3-VIB-5.6.3-iii}
  \emph{For every reduced \(S\)-scheme \(T\), two \(S\)-morphisms
  \(f_1,f_2:T\to X\) that agree on a dense open subset \(U\) of
  \(T\) are equal.}

\item[\textup{\SGAInlineRefTarget{lem:sga3-VIB-5.6.3:item:iv}(iv)}]\phantomsection\label{cond:sga3-VIB-5.6.3-iv}
  \emph{For every \(s\in S\) and every pair of points
  \(x_1,x_2\) of \(X_s\), there exist a morphism}
  \[
    \phi:S''\longrightarrow S'\longrightarrow S
  \]
  \emph{and an open subscheme \(V\) of \(X_{S''}\), separated
  over \(S''\), such that:}
  \begin{enumerate}
  \item[\textup{(a)}]\phantomsection\label{cond:sga3-VIB-5.6.3-iv-a}
    \emph{\(S'\to S\) is open, \(S''\to S'\) is closed and
    surjective, and \(\phi^{-1}(s)\ne\varnothing\);}
  \item[\textup{(b)}]\phantomsection\label{cond:sga3-VIB-5.6.3-iv-b}
    \emph{the inverse image of \(\{x_1,x_2\}\) in \(X_{S''}\)
    is contained in \(V\).}
  \end{enumerate}

\item[\textup{(iv\(^{\prime}\))}]
  \phantomsection\label{cond:sga3-VIB-5.6.3-iv-prime}
  \emph{For every \(s\in S\), every pair of points
  \(x_1,x_2\in X_s\) is contained in an open subset \(V\) of
  \(X\) that is separated over \(S\).}%
  \footnote{\label{note:sga3-VIB-44}Editor's note: we have
  simplified the formulation of condition \textup{\SGARefLink{lem:sga3-VIB-5.6.3:item:iv}{(iv)}} and
  added condition \textup{(iv\(^{\prime}\))}. We have also added
  the proof of the implication \textup{\SGARefLink{lem:sga3-VIB-5.6.3:item:i}{(i)}}\(\Rightarrow\)\textup{\SGARefLink{lem:sga3-VIB-5.6.3:item:iii}{(iii)}},
  which is used in the proof of
  \textup{\SGARefLink{lem:sga3-VIB-5.6.3:item:iv}{(iv)}}\(\Rightarrow\)\textup{\SGARefLink{lem:sga3-VIB-5.6.3:item:iii}{(iii)}}. Observe moreover
  that, if
  \[
    T=\operatorname{Spec}B,\qquad
    B=k[\varepsilon,x]/(\varepsilon^2,\varepsilon x)
  \]
  \textup{(\(k\) a field)} and
  \[
    X=T_{\mathrm{red}}=\operatorname{Spec}k[x],
  \]
  then the morphisms
  \[
    \phi_\lambda:T\longrightarrow X,\qquad
    x\longmapsto x+\lambda\varepsilon
    \quad(\lambda\in k),
  \]
  agree on the dense open subset
  \(\operatorname{Spec}B_x\), but are not equal.}
\end{enumerate}

\emph{Proof.}
The implication \textup{\SGARefLink{lem:sga3-VIB-5.6.3:item:ii}{(ii)}}\(\Rightarrow\)\textup{\SGARefLink{lem:sga3-VIB-5.6.3:item:i}{(i)}} is
clear (take \(T=X\) and let \(\sigma\) be the diagonal section),
as are
\[
  \textup{(i)}\Rightarrow\textup{(iv\(^{\prime}\))}
  \Rightarrow\textup{(iv)}.
\]
On the other hand,
\(\textup{(i)}\Rightarrow\textup{(ii)}\) by
EGA~I, 5.4.6.

\smallskip
\noindent\textup{\SGARefLink{lem:sga3-VIB-5.6.3:item:iii}{(iii)}}\(\Rightarrow\)\textup{\SGARefLink{lem:sga3-VIB-5.6.3:item:ii}{(ii)}}.
Let \(\sigma:T\to X_T\) be a section of \(p:X_T\to T\).
By EGA~I, 5.3.13, \(\sigma\) is an immersion, that is, an
isomorphism of \(T\) onto a locally closed subscheme \(E\) of
\(X_T\). To show that \(E\) is closed, we may suppose that \(T\)
and \(E\) are reduced. Let \(\overline E\) be the reduced closed
subscheme of \(X_T\) whose underlying space is the closure of
\(E\), so that \(E\) is a dense open subscheme of
\(\overline E\). The immersion
\[
  i:\overline E\hookrightarrow X_T
\]
and \(\sigma\circ p\circ i\) agree on \(E\), hence on
\(\overline E\) by hypothesis \textup{\SGARefLink{lem:sga3-VIB-5.6.3:item:iii}{(iii)}}. Thus every point
of \(\overline E\) belongs to \(\sigma(T)=E\), and consequently
\(\overline E=E\).

\smallskip
\noindent\textup{\SGARefLink{lem:sga3-VIB-5.6.3:item:i}{(i)}}\(\Rightarrow\)\textup{\SGARefLink{lem:sga3-VIB-5.6.3:item:iii}{(iii)}}.
Suppose that \(X\) is separated over \(S\), and let \(T\) be a
reduced \(S\)-scheme and \(f_1,f_2:T\to X\) two \(S\)-morphisms
that agree on a dense open subset \(U\) of \(T\). Let
\[
  g:T\longrightarrow X\times_S X
\]
be the morphism with components \(f_1\) and \(f_2\). Since
\[
  D=\Delta_{X/S}(X)
\]
is closed, its inverse image under \(g\) is a closed subset of
\(T\) containing \(U\), and is therefore equal to \(T\). Since
\(T\) is reduced, \(g\) factors through \(D\) (cf.\
EGA~I, 5.2.2); consequently
\[
  f_1=p_1\circ g=p_2\circ g=f_2.
\]

\smallskip
\noindent\textup{\SGARefLink{lem:sga3-VIB-5.6.3:item:iv}{(iv)}}\(\Rightarrow\)\textup{\SGARefLink{lem:sga3-VIB-5.6.3:item:iii}{(iii)}}.
Let \(T\) be a reduced \(S\)-scheme, and let
\(f_1,f_2:T\to X\) be two \(S\)-morphisms that agree on a dense
open subset \(U\). Since \(T\) is reduced, in order to prove that
\(f_1=f_2\), it is enough to prove that \(f_1=f_2\)
set-theoretically. Indeed, suppose this has been established,
let \(t\in T\), let \(V\) be an affine open subset of \(X\)
containing \(f_1(t)=f_2(t)\), and let \(W\) be the open
neighborhood of \(t\) that is the inverse image of \(V\) under
the continuous map underlying both \(f_1\) and \(f_2\). Then the
morphisms
\[
  f_i|_W:W\longrightarrow V
\]
agree on the dense open subset \(U\cap W\) of \(W\). Since
\(V\) is separated and \(W\) is reduced, this implies that
\(f_1|_W=f_2|_W\), and hence that \(f_1=f_2\).

Now let \(t\in T\), let \(s\) be its image in \(S\), and let us
show that the points
\[
  x_1=f_1(t),\qquad x_2=f_2(t)
\]
of \(X_s\) are equal. Let
\[
  \phi:S''\longrightarrow S'\longrightarrow S
\]
and an open subset \(V\) of \(X\times_S S''\) be as in
\textup{\SGARefLink{lem:sga3-VIB-5.6.3:item:iv}{(iv)}}. Set
\[
  T'=T\times_S S',\qquad T''=T\times_S S'',
\]
and denote by
\[
  g:T''\longrightarrow T,\qquad
  f_i'':T''\longrightarrow X_{S''}\quad(i=1,2)
\]
the morphisms obtained by base change.

Since \(U\) is dense in \(T\) and \(T'\to T\) is open, the
inverse image \(U'\) of \(U\) in \(T'\) is dense in \(T'\).
Let \(U''\) be the inverse image of \(U'\) in \(T''\), and let
\(F\) be the reduced subscheme of \(T''\) whose underlying space
is \(\overline{U''}\). Since \(T''\to T'\) is surjective and
closed, the image of \(F\) contains \(U'\) and is closed, and is
therefore equal to \(T'\). Consequently,
\[
  F\cap g^{-1}(t)
\]
contains a point \(u\).

For \(i=1,2\), let \(h_i\) be the restriction of \(f_i''\) to
\(F\). Then \(h_i(u)\) is a point of \(X_{S''}\) above
\(f_i(t)=x_i\), and hence belongs to \(V\), since \(V\) contains
the inverse image of \(\{x_1,x_2\}\) in \(X_{S''}\).

Thus
\[
  W=h_1^{-1}(V)\cap h_2^{-1}(V)
\]
is an open subset of \(F\) containing \(u\), and the
\(S''\)-morphisms
\[
  h_i|_W:W\longrightarrow V
\]
agree on the dense open subset \(U''\cap W\) of \(W\). Since
\(V\) is separated over \(S''\), we have
\[
  h_1|_W=h_2|_W.
\]
It follows that \(h_1(u)=h_2(u)\), and hence that
\(f_1(t)=f_2(t)\). This proves
\(\textup{(iv)}\Rightarrow\textup{(iii)}\).

\medskip
\phantomsection\label{ded:sga3-VIB-theorem-5.3-from-5.6.2-5.6.3}%
\hyperref[thm:sga3-VIB-5.3]{Theorem 5.3} now follows from \hyperref[cor:sga3-VIB-5.6.2]{5.6.2} and the implication
\(\textup{(iv)}\Rightarrow\textup{(i)}\) of \hyperref[lem:sga3-VIB-5.6.3]{5.6.3}.

% END INCLUDED SGA3_English_current_buildable_source_20260823/source/inputs/VI/tex_reference_v2/components/46_expose_VIB_lemma563_en.tex

% BEGIN INCLUDED SGA3_English_current_buildable_source_20260823/source/inputs/VI/tex_reference_v2/components/47_expose_VIB_counter564_thm57_en.tex
% SGA 3, Expose VI_B, Loop 1 English reconstruction.
% Scope: Counterexample 5.6.4 and Theorem 5.7 in full, including
% source editor's notes 45 and 46.
% Authority: Polo--Gille Exp6B-13oct24.pdf, component-local pages
% 34--35 (combined-reader pages 375--376; original-margin printed
% pages 356--357 and re-edition headers 364--365).
% Corollary 5.7.1 begins on local page 35 / combined page 376 /
% original-margin printed page 357 and is excluded.
% OCR and Reinhold English Markdown are comparison witnesses only.
% This unit contains no diagrams.

\medskip
\phantomsection\label{ctr:sga3-VIB-5.6.4}%
\noindent\textbf{Counter\SGARefLink{ctr:sga3-VIB-5.6.4}{example 5.6.4}.---}
\emph{Not every \(S\)-group \(G\) is separated. Let \(S\) be a
scheme having a non-isolated closed point \(s\), and let \(G\) be
the scheme obtained by gluing two copies of \(S\) along the open
subset \(S-\{s\}\). It is easily seen that \(G\) is not separated
over \(S\), and that it carries a natural \(S\)-group structure
whose fibers are all trivial, except for the fiber \(G_s\), which
is isomorphic to the two-element group
\(\mathbb Z/2\mathbb Z\).}%
\footnote{\label{note:sga3-VIB-45}Editor's note: we reproduced this
example in VI\(_A\), \SGARefLink{sga3:VIA:0.3}{0.3}, editor's \SGARefLink{sec:sga3-VIB-5}{note 5}.}

\medskip
\phantomsection\label{thm:sga3-VIB-5.7}%
\noindent\textbf{Theorem 5.7.---}%
\footnote{\label{note:sga3-VIB-46}Editor's note: on the one hand,
we omitted Corollary 5.6.5, which repeated \SGARefLink{cor:sga3-VIB-5.5}{5.5}. On the other hand,
the original stated \hyperref[cor:sga3-VIB-5.7.1]{Corollary 5.7.1} below as \hyperref[cor:sga3-VIB-5.7.1]{Remark 5.7}, referring
for its proof to \SGARefLink{par:sga3-VIA-4.7}{4.7}, a number that does not exist in
\emph{Lecture Notes in Mathematics} 151 (but did appear in the
1965 edition of SGAD, whose numbers \SGARefLink{par:sga3-VIA-4.5}{4.5} and \SGARefLink{par:sga3-VIA-4.6}{4.6} became \SGARefLink{lem:sga3-VIB-5.6.1}{5.6.1} and
\SGARefLink{cor:sga3-VIB-5.6.2}{5.6.2}); we have added \hyperref[thm:sga3-VIB-5.7]{Theorem 5.7}, communicated by O.\ Gabber,
which makes the aforementioned statement of SGAD more precise.}
\emph{Let \(S\) be a scheme, let \(G\) be an \(S\)-group locally
of finite presentation over \(S\), such that the function
\[
  s\longmapsto\dim G_s
\]
is locally constant on \(S\), let \(X\) be an \(S\)-scheme on
which \(G\) acts, and let \(U\) be an open subset of \(X\),
separated over \(S\). Then
\(\underline{G^0}\mathbin{\cdot}U\) is an open subset of \(X\),
separated over \(S\).}

\emph{Proof.}
Let
\[
  \mu:G\times_SX\longrightarrow X
\]
denote the action of \(G\) on \(X\). It is the composite of the
automorphism
\[
  (g,x)\longmapsto(g,gx)
\]
of \(G\times_SX\) and the projection onto \(X\). Since
\(G\times_SU\) is an open subset of \(G\times_SX\), it follows
from \hyperref[concl:sga3-VIB-4.2-iii]{4.2}\textup{(iii)} that
\[
  V=\underline{G^0}\mathbin{\cdot}U
\]
is open in \(X\).

Then, arguing as in the proof of \hyperref[cor:sga3-VIB-5.6.2]{5.6.2}, we deduce from the
implication \textup{(iv)}\(\Rightarrow\)\textup{(i)} of \hyperref[lem:sga3-VIB-5.6.3]{5.6.3}
that \(V\) is separated over \(S\).

% END INCLUDED SGA3_English_current_buildable_source_20260823/source/inputs/VI/tex_reference_v2/components/47_expose_VIB_counter564_thm57_en.tex

% BEGIN INCLUDED SGA3_English_current_buildable_source_20260823/source/inputs/VI/tex_reference_v2/components/48_expose_VIB_cor571_through_cor582_en.tex
% SGA 3, Expose VI_B, Loop 1 English reconstruction.
% Scope: Corollary 5.7.1, Remark 5.7.2, the Section 5.8 heading,
% Proposition 5.8.1, and Corollary 5.8.2 in full, including source
% editor's note 47.
% Authority: Polo--Gille Exp6B-13oct24.pdf, component-local page 35
% (combined-reader page 376; original-margin printed page 357 and
% re-edition header 365); note 47 is checked through its continuation
% at the top of local page 36 / combined-reader page 377.
% Example 5.8.3 begins on local page 35 / combined-reader page 376.
% OCR and Reinhold English Markdown are comparison witnesses only.
% This unit contains no diagrams.

\medskip
\phantomsection\label{cor:sga3-VIB-5.7.1}%
\noindent\textbf{Corollary 5.7.1.---}
\emph{Let \(S\) and \(G\) be as in
\hyperref[thm:sga3-VIB-5.7]{5.7}, and let \(\sigma,\tau\) be two
\(S\)-sections of \(G^0\), that is,
\(\sigma,\tau\in G^0(S)\). Then the coincidence subscheme
\(S(\sigma,\tau)\subset S\) of \(\sigma\) and \(\tau\)
\textup{(that is, the inverse image of the diagonal of
\(G\times_S G\) under the morphism \((\sigma,\tau)\))} is closed.}

Indeed, for every \(s\in S\), let \(U\) be an affine open subset
of \(G\) containing \(\epsilon(s)\). Then
\[
  V=\epsilon^{-1}(U)
\]
is an open neighborhood of \(s\) in \(S\), and, since \(G^0U\)
is separated, \(S(\sigma,\tau)\cap V\) is closed in \(V\).

\medskip
\phantomsection\label{rem:sga3-VIB-5.7.2}%
\noindent\textbf{Remark 5.7.2.---}
Gabber points out that one can show the following: if \(S\) is
henselian local, with closed point \(s\), then the intersection
of the open subsets \(G^0U\), where \(U\) ranges over the affine
open neighborhoods of \(\epsilon(s)\), is an open group
subscheme of \(G\), separated over \(S\).

\section*{5.8. Complements}
\phantomsection\label{sec:sga3-VIB-5.8}%
\addcontentsline{toc}{section}{\SGARefLink{sec:sga3-VIB-5.8}{5.8}. Complements}

\noindent\footnote{\label{note:sga3-VIB-47}Editor's note: we have
added this paragraph of complements and counterexamples, all
communicated by O.\ Gabber.}%
We begin by recalling
\hyperref[prop:sga3-VIA-2.6.6]{Proposition VI\(_A\), 2.6.6}:

\medskip
\phantomsection\label{prop:sga3-VIB-5.8.1}%
\noindent\textbf{Proposition 5.8.1.---}
\emph{Let \(k\) be a field, let \(G\) be a \(k\)-group locally of
finite type acting on a \(k\)-scheme \(X\) in such a way that the
morphism}
\[
  G\times X\longrightarrow X\times X,
  \qquad (g,x)\longmapsto(gx,x),
\]
\emph{is surjective. Suppose that \(X\) is quasi-separated. Then
the connected components of \(X\) are irreducible.}

\medskip
\phantomsection\label{cor:sga3-VIB-5.8.2}%
\noindent\textbf{Corollary 5.8.2.---}
\emph{Let \(S,G,X\) be as in the preliminary hypotheses of
\hyperref[thm:sga3-VIB-5.3]{5.3}. Suppose moreover that every
fiber \(X_s\) is quasi-separated and connected. Then \(X\) is
separated and quasi-compact over \(S\).}

Indeed, by the preceding proposition, every fiber \(X_s\) is
irreducible, and the rest of the proof is identical to that of
\hyperref[cor:sga3-VIB-5.4]{5.4}.

% END INCLUDED SGA3_English_current_buildable_source_20260823/source/inputs/VI/tex_reference_v2/components/48_expose_VIB_cor571_through_cor582_en.tex

% BEGIN INCLUDED SGA3_English_current_buildable_source_20260823/source/inputs/VI/tex_reference_v2/components/49_expose_VIB_example583_en.tex
% SGA 3, Expose VI_B, Loop 1 English reconstruction.
% Scope: Example 5.8.3 in full.
% Authority: Polo--Gille Exp6B-13oct24.pdf, component-local pages
% 35--36 (combined-reader pages 376--377; original-margin printed
% page 357 and re-edition headers 365--366).
% Example 5.8.4 begins on local page 36 / combined page 377.
% OCR and Reinhold English Markdown are comparison witnesses only.
% This unit contains no diagrams or source editor's notes.

\medskip
\phantomsection\label{ex:sga3-VIB-5.8.3}%
\noindent\textbf{Example 5.8.3.---}
Fix an algebraically closed field \(k\). Recall first that every
``locally Boolean'' topological space \(X\), that is, every space
having a basis of compact open subsets, becomes a \(k\)-scheme
when equipped with the sheaf of locally constant \(k\)-valued
functions (cf.\ [DG70], I, \S1, \SGARefLink{cor:sga3-VIB-2.12}{2.12}). Such an \(X\) will be
called a \emph{locally Boolean \(k\)-scheme}.

We also note that every topological space \(E\) having a basis of
separated open subsets---and likewise every scheme \(X\)---has a
dense separated open subset. Indeed, every increasing union of
separated open subsets is separated (for a scheme \(X\), this
follows from \hyperref[lem:sga3-VIB-5.6.3]{5.6.3}), so there is a maximal such open subset \(U\).
It must be dense: otherwise there would be a nonempty open subset
\(V\) with \(U\cap V=\varnothing\); the open set \(V\) would
contain a nonempty separated open subset \(W\), and
\[
  U\cup W
\]
would still be separated, contradicting the maximality of \(U\).

Now let \(C\) be the Cantor space, viewed as the underlying space
of the group
\[
  (\mathbb Z/2\mathbb Z)^{\mathbb N}
\]
with its product topology. For every \(p\in C\), let \(C(p)\) be
another copy of \(C\), and let \(X\) be obtained by gluing each
\(C(p)\) to \(C\) along \(C-\{p\}\). Then \(X\) is a
nonseparated locally Boolean \(k\)-scheme, and \(C\) is a dense
open subset of \(X\).

Let \(G\) be the group of automorphisms of the \(k\)-scheme \(X\),
that is, of homeomorphisms of \(X\). Then \(G\) acts transitively
on \(X\). Indeed, \(X\) is the union of \(C\) and, for each point
\(p\in C\), a second point \(p'\). This second point is
characterized as the unique \(x\in X-\{p\}\) for which
\[
  (C-\{p\})\cup\{x\}
\]
is separated. Consequently, every automorphism \(\phi\) of \(C\)
extends to an automorphism \(\phi_X\) of \(X\) satisfying
\[
  \phi_X(p')=\phi_X(p)'
\]
for every \(p\). Moreover, the map
\[
  \theta_p:X\longrightarrow X
\]
that exchanges \(p\) and \(p'\) and fixes every other point is an
automorphism of \(X\). Finally, the automorphism group of \(C\)
acts transitively on \(C\): using the group structure of \(C\),
it is enough to consider translations.

Thus the discrete \(k\)-group \(G\), which is locally of finite
type, acts on the \(k\)-scheme \(X\) in such a way that
\[
  G\times X\longrightarrow X\times X,
  \qquad (g,x)\longmapsto(gx,x),
\]
is surjective, but \(X\) is not separated, even though \(C\) is a
dense separated open subset.

% END INCLUDED SGA3_English_current_buildable_source_20260823/source/inputs/VI/tex_reference_v2/components/49_expose_VIB_example583_en.tex

% BEGIN INCLUDED SGA3_English_current_buildable_source_20260823/source/inputs/VI/tex_reference_v2/components/50_expose_VIB_example584_en.tex
% SGA 3, Expose VI_B, Loop 1 English reconstruction.
% Scope: Example 5.8.4 in full.
% Authority: Polo--Gille Exp6B-13oct24.pdf, component-local pages
% 36--37 (combined-reader pages 377--378; re-edition headers
% 366--367). Section 6 begins later on local page 37.
% OCR and Reinhold English Markdown are comparison witnesses only.
% This unit contains no diagrams or source editor's notes.

\medskip
\phantomsection\label{ex:sga3-VIB-5.8.4}%
\noindent\textbf{Example 5.8.4.---}
Retain the notation of the preceding example. Using the
description of \(C\) as the set of sequences
\[
  (u_n)_{n\in\mathbb N},
  \qquad u_n\in\mathbb Z/2\mathbb Z,
\]
one sees that \(C\) with one point \(p\) removed is homeomorphic
to a countable disjoint union of copies of \(C\). Indeed, using
the group structure of \(C\), for example, translation reduces us
to the case in which \(p=0\), the zero sequence. Then
\[
  C-\{0\}=\coprod_{i\in\mathbb N}C_i,
\]
where
\[
  C_i=
  \bigl\{(u_n)_{n\in\mathbb N}\mathrel{\big|}
    u_i=1\ \text{and}\ u_j=0\ \text{for}\ j<i\bigr\},
\]
and every \(C_i\) is homeomorphic to \(C\). It follows by
induction on \(\lvert F\rvert\) that, for every nonempty finite
subset \(F\) of \(C\), the space \(C-F\) is homeomorphic to a
countable disjoint union of copies of \(C\), and hence to \(C\)
with one point removed.

For each \(F\subset C\) of cardinality \(2\), let \(C(F)\) be
another copy of \(C\), denote by \(q_F\) the point \(0\) of
\(C(F)\), and choose a homeomorphism
\[
  \phi_F:C(F)-\{q_F\}\xrightarrow{\sim}C-F.
\]
Let \(X'\) be obtained by gluing every \(C(F)\) to \(C\) by means
of \(\phi_F\). Then \(X'\) is a nonseparated locally Boolean
\(k\)-scheme. Moreover, the construction shows that every locally
constant function
\[
  f:X'\longrightarrow k
\]
is constant. Indeed, if \(x,y\in C\) and \(F=\{x,y\}\), then every
neighborhood of \(q_F\) meets every neighborhood of \(x\) and of
\(y\). Thus
\[
  f(x)=f(q_F)=f(y).
\]
If \(F'=\{z,t\}\), with \(z,t\in C\), then likewise
\[
  f(q_{F'})=f(z)=f(q_{\{z,x\}})=f(x).
\]
Consequently, \(X'\) is connected.

Every point \(x\in X'\) also has a pointed neighborhood
homeomorphic to \((C,0)\). More precisely, for every \(F\), fix a
homeomorphism of pointed topological spaces
\[
  h_F:(C,0)\xrightarrow{\sim}(C(F),q_F),
\]
and let \(T\) be the group of translations of \(C\). If
\(x=q_F\), use \(h_F\); if \(x\in C\), the translation
\(t_x\in T\) is a homeomorphism
\[
  t_x:(C,0)\xrightarrow{\sim}(C,x),
\]
and it is the unique element of \(T\) with this property.%
\footnote[1]{\label{note:sga3-VIB-5.8.4-Xprime-hF}Translator's
note: in these two passages the French source prints \(X\) in
place of the \(X'\) just constructed, and \(q_F\) in place of the
homeomorphism \(h_F\) just defined. The English uses \(X'\) and
\(h_F\), respectively.}

Let \(L\) be the free group generated by the \(h_F\), and set
\[
  G=T*L,
\]
the free product (coproduct) of \(T\) and \(L\). For
\[
  h\in\{h_F\}_{\lvert F\rvert=2}\cup T,
\]
let \(\sigma(h)\) denote the corresponding generator of \(G\),
and let \(S(h)\) and \(B(h)\) be respectively the source and the
target of \(h\). It is convenient also to set
\[
  \sigma(h_F^{-1})=\sigma(h_F)^{-1},\qquad
  S(h_F^{-1})=B(h_F),\qquad
  B(h_F^{-1})=S(h_F),
\]
and to let \(E\) be the set consisting of the elements of \(T\)
and all \(h_F,h_F^{-1}\).

On the product space
\[
  P=G\times X',
\]
where \(G\) carries the discrete topology, consider the
equivalence relation generated by
\[
  \bigl(g\sigma(h),x\bigr)\sim\bigl(g,h(x)\bigr)
  \qquad\text{when }x\in S(h),
\]
for every \(h\in E\), and let \(Z\) be the quotient space. Thus
\(Z\) is obtained by gluing open subsets of the disjoint union
\[
  \coprod_{g\in G}\{g\}\times X'.
\]
Consequently, the saturation \(\widetilde\Omega\) of every open
subset \(\Omega\) of \(P\) is open (cf.\ [BTop], I, \S5.1,
Example~2). Explicitly, since every open subset of \(P\) is the
union of its intersections with the slices
\(\{g\}\times X'\), it is enough to consider
\(\{1\}\times W\), with \(W\) open in \(X'\). Its saturation is
the union of \(\{1\}\times W\) and the sets
\[
  \{\sigma(h_1)\}\times
  h_1^{-1}\bigl(W\cap B(h_1)\bigr),
\]
\[
  \{\sigma(h_1)\sigma(h_2)\}\times
  h_2^{-1}\Bigl(
    h_1^{-1}\bigl(W\cap B(h_1)\bigr)\cap B(h_2)
  \Bigr),
\]
and so on, for all finite sequences
\(h_1,\ldots,h_n\) of elements of \(E\). It is therefore open,
and the projection
\[
  \pi:P\longrightarrow Z
\]
is open.

Moreover, the word
\[
  \sigma(h_1)\cdots\sigma(h_n)
\]
is reduced in \(G\), unless one of the \(h_i\) is the identity of
\(T\), two consecutive \(h_i\)'s belong to \(T\), or one of the
pairs \((h_i,h_{i+1})\) is
\[
  (h_F,h_F^{-1})
  \quad\text{or}\quad
  (h_F^{-1},h_F).
\]
Suppose, therefore, that \(x\in X'\) and that
\[
  \beta=
  \bigl(
    \sigma(h_1)\cdots\sigma(h_n),
    h_n^{-1}\cdots h_1^{-1}(x)
  \bigr)
\]
belongs to \(\{1\}\times X'\). We may then suppose that every
\(h_i\) is a translation \(t_i\). In that case,
\[
  \sigma(t_1\cdots t_n)=1
\]
implies \(t_1\cdots t_n=1\), and hence
\(\beta=(1,x)\). The equivalence relation is compatible with the
action of \(G\) on \(P=G\times X'\) by left translation on the
first factor. It follows that, for every \(g\in G\), the
restriction of \(\pi\) to \(\{g\}\times X'\) is injective.

Now let \(z\in Z\) and choose a representative \((g,x)\in P\).
By the preceding discussion,
\[
  U=\pi\bigl(\{g\}\times X'\bigr)
\]
is an open neighborhood of \(z\), and the continuous map
\[
  \{g\}\times X'\longrightarrow U
\]
induced by \(\pi\) is open and bijective, hence a homeomorphism.
Thus \(Z\) is locally isomorphic to \(X'\), and hence also to
\(C\); it is therefore again a locally Boolean \(k\)-scheme.

Finally, \(G\) acts transitively on \(Z\). Indeed, every \(z\in Z\)
is \(G\)-conjugate to an element of the form
\(\pi(1,x)\), so it is enough to show that every
\((1,x)\in P\) is equivalent to some \((\sigma(h),0)\). If
\(x\in C\), take \(h=t_x\); if \(x=q_F\), take \(h=h_F\). Thus
\(Z\) is a locally Boolean \(k\)-scheme endowed with a transitive
action of the discrete group \(G\), but \(Z\) is not separated.

% END INCLUDED SGA3_English_current_buildable_source_20260823/source/inputs/VI/tex_reference_v2/components/50_expose_VIB_example584_en.tex

% BEGIN INCLUDED SGA3_English_current_buildable_source_20260823/source/inputs/VI/tex_reference_v2/components/51_expose_VIB_definition61_en.tex
% SGA 3, Expose VI_B, Loop 1 English reconstruction.
% Scope: Section 6 heading, Definition 6.1 (i)--(iii), and
% Remark 6.1.1 in full, including the section star note and source
% editor's notes 48--50.
% Authority: Polo--Gille Exp6B-13oct24.pdf, component-local pages
% 37--39 (combined-reader pages 378--380; displayed re-edition
% headers 367--369; original-margin printed page 358 begins on
% local page 38).
% Proposition 6.2 begins on local page 39 / combined page 380 and
% is excluded.
% OCR and Reinhold English Markdown are comparison witnesses only.
% This unit contains no diagrams.

\begingroup
\renewcommand{\thefootnote}{*}%
\section*{6. Subfunctors and group subschemes%
  \texorpdfstring{\protect\footnotemark}{}}
\footnotetext{\label{note:sga3-VIB-section6-star}On the same
theme, see also the results of Expos\'e~XI, \S6, whose natural
place would be in the present Expos\'e~VI\(_B\). There one finds,
in particular, representability criteria for certain group
subfunctors of a given group scheme.
\phantomsection\label{note:sga3-VIB-48}\textup{(48) Editor's note: we have
inserted these results below
(\hyperref[thm:sga3-VIB-6.5.2]{6.5.2}--%
\hyperref[cor:sga3-VIB-6.5.5]{6.5.5}).}}
\endgroup
\phantomsection\label{sec:sga3-VIB-6}%
\addcontentsline{toc}{section}{6. Subfunctors and group subschemes}

\medskip
\phantomsection\label{def:sga3-VIB-6.1}%
\noindent\textbf{Definition 6.1.---}
\textup{(i)} Let \(X\) be an \(S\)-functor (that is, a functor
from \((\mathbf{Sch}/S)^\circ\) to \((\mathbf{Ens})\)), let \(G\)
be an \(S\)-functor in groups, and let \(u\) and \(v\) be two
\(S\)-morphisms from \(X\) to \(G\). The \emph{transporter of
\(u\) into \(v\)}, denoted
\(\underline{\operatorname{Transp}}(u,v)\), is the
sub-\(S\)-functor of \(G\) defined by
\phantomsection\label{eq:sga3-VIB-6.1-transp-uv}%
\[
\begin{aligned}
\underline{\operatorname{Transp}}(u,v)(S')
  &=
  \bigl\{g\in G(S')\bigm|
    (\operatorname{int}g)\circ u_{S'}=v_{S'}\bigr\}                                      \\
  &=
  \left\{g\in G(S')\mathrel{\left|\,\vphantom{\begin{array}{c}a\\a\end{array}}\right.}
    \begin{array}{l}
      g_{S''}u_{S''}(x)g_{S''}^{-1}=v_{S''}(x),\\[-1pt]
      \text{for every \(x\in X(S'')\) and every \(S''\to S'\)}
    \end{array}
  \right\}.
\end{aligned}
\]
In particular,
\(\underline{\operatorname{Transp}}(u,u)\) is a group
sub-\(S\)-functor of \(G\); it is called the
\emph{centralizer of \(u\)} and is denoted
\(\underline{\operatorname{Centr}}(u)\).

\smallskip
\noindent\textup{(ii)} Let \(G\) be an \(S\)-functor in groups,
and let \(X\) and \(Y\) be two sub-\(S\)-functors of \(G\). The
\emph{transporter of \(X\) into \(Y\)} (respectively, the
\emph{strict transporter of \(X\) into \(Y\)}), denoted
\(\underline{\operatorname{Transp}}_G(X,Y)\) (respectively,
\(\underline{\operatorname{Transp}}^{\mathrm{str}}_G(X,Y)\)),
is the sub-\(S\)-functor of \(G\) defined by
\phantomsection\label{eq:sga3-VIB-6.1-transp-XY}%
\[
\begin{aligned}
\underline{\operatorname{Transp}}_G(X,Y)(S')
  &=
  \bigl\{g\in G(S')\bigm|
    (\operatorname{int}g)(X_{S'})\subset Y_{S'}\bigr\}                 \\
  &=
  \bigl\{g\in G(S')\bigm|
    g_{S''}X(S'')g_{S''}^{-1}\subset Y(S'')
    \text{ for every \(S''\to S'\)}\bigr\},
\end{aligned}
\]
and, respectively,
\phantomsection\label{eq:sga3-VIB-6.1-transpstr-XY}%
\[
\begin{aligned}
\underline{\operatorname{Transp}}^{\mathrm{str}}_G(X,Y)(S')
  &=
  \bigl\{g\in G(S')\bigm|
    (\operatorname{int}g)(X_{S'})=Y_{S'}\bigr\}                       \\
  &=
  \bigl\{g\in G(S')\bigm|
    g_{S''}X(S'')g_{S''}^{-1}=Y(S'')
    \text{ for every \(S''\to S'\)}\bigr\}.
\end{aligned}
\]
Note that
\phantomsection\label{eq:sga3-VIB-6.1-inversion}%
\[
\underline{\operatorname{Transp}}^{\mathrm{str}}_G(X,Y)
 =
\underline{\operatorname{Transp}}_G(X,Y)
 \cap
c\bigl(\underline{\operatorname{Transp}}_G(Y,X)\bigr),
\]
where \(c\) denotes the inversion morphism of \(G\).%
\footnote{\label{note:sga3-VIB-49}Editor's note: if
\(u:X\to G\) and \(v:Y\to G\) are two arbitrary morphisms of
\(S\)-functors, let \(u(X)\) be the image functor of \(u\),
defined by
\[
  u(X)(S')=u(S')\bigl(X(S')\bigr)\subset G(S');
\]
it is a subfunctor of \(G\), as is the image functor \(v(Y)\).
One may then consider the transporter of the image of \(u\) into
the image of \(v\),
\(\underline{\operatorname{Transp}}_G(u(X),v(Y))\). Thus, in
part~\textup{(ii)} of the definition, it is not necessary to
restrict to subfunctors, that is, to the case in which \(u\) and
\(v\) are monomorphisms. In the original, this restriction
sometimes imposed repetitions in the hypotheses, such as:
``Let \(u,v:X\to G\) be morphisms of \(S\)-functors and let
\(w:H\to G\) and \(w':K\to G\) be monomorphisms; then
\(\underline{\operatorname{Transp}}(u,v)\),
\(\underline{\operatorname{Centr}}_G(w)
 =\underline{\operatorname{Centr}}_G H\), and
\(\underline{\operatorname{Transp}}_G(H,K)\) satisfy \(\ldots\).''
Such repetitions can be avoided by considering
\(\underline{\operatorname{Transp}}_G(u(X),v(Y))\). We have made
such modifications in \hyperref[prop:sga3-VIB-6.2]{6.2} and,
later, in \hyperref[prop:sga3-VIB-10.11]{10.11}.}

\smallskip
\noindent\textup{\SGARefLink{prop:sga3-VIB-6.2:item:iii}{(iii)}} Let \(G\) be an \(S\)-functor in groups,
let \(H\) be a sub-\(S\)-functor of \(G\), and let
\(i:H\to G\) be the canonical \(S\)-morphism. The
\emph{centralizer} and \emph{normalizer of \(H\) in \(G\)} are
the following group sub-\(S\)-functors of \(G\):
\phantomsection\label{eq:sga3-VIB-6.1-centralizer-normalizer}%
\[
\underline{\operatorname{Centr}}_G H
  =\underline{\operatorname{Centr}}(i)
  =\underline{\operatorname{Transp}}(i,i),
\qquad
\underline{\operatorname{Norm}}_G H
  =\underline{\operatorname{Transp}}^{\mathrm{str}}_G(H,H).
\]
Finally, the \emph{center of \(G\)} is the \(S\)-functor in
groups
\[
  \underline{\operatorname{Centr}}(\operatorname{id}_G)
  =\underline{\operatorname{Centr}}_G G;
\]
it is denoted \(\underline{\operatorname{Centr}}G\).

\medskip
\phantomsection\label{rem:sga3-VIB-6.1.1}%
\noindent\textbf{Remark 6.1.1.---}%
\footnote{\label{note:sga3-VIB-50}Editor's note: we have added
this remark, which is used implicitly in
\hyperref[prop:sga3-VIB-6.2]{Proposition~6.2}.}
It follows from the definitions that the functors
\(\underline{\operatorname{Transp}}(u,v)\) and
\(\underline{\operatorname{Transp}}_G(X,Y)\) (and hence also
\(\underline{\operatorname{Transp}}^{\mathrm{str}}_G(X,Y)\))
commute with base change: for every \(S'\to S\), if
\(G',X',Y',u',v'\) are obtained from \(G,X,Y,u,v\) by base
change, then
\phantomsection\label{eq:sga3-VIB-6.1.1-base-change}%
\[
\underline{\operatorname{Transp}}(u,v)_{S'}
  =\underline{\operatorname{Transp}}(u',v'),
\qquad
\underline{\operatorname{Transp}}(X,Y)_{S'}
  =\underline{\operatorname{Transp}}(X',Y').
\]

% END INCLUDED SGA3_English_current_buildable_source_20260823/source/inputs/VI/tex_reference_v2/components/51_expose_VIB_definition61_en.tex

% BEGIN INCLUDED SGA3_English_current_buildable_source_20260823/source/inputs/VI/tex_reference_v2/components/52_expose_VIB_prop62_through_prop622_en.tex
% SGA 3, Expose VI_B, Loop 1 English reconstruction.
% Scope: Proposition 6.2, Definition 6.2.1, and Proposition 6.2.2
% in full, including source editor's notes 51--54.
% Authority: Polo--Gille Exp6B-13oct24.pdf, component-local pages
% 39--40 (combined-reader pages 380--381; displayed re-edition
% headers 369--370). Theorem 6.2.3 begins on local page 40 and is
% excluded.
% OCR and Reinhold English Markdown are comparison witnesses only.
% This unit contains no diagrams.

\medskip
\phantomsection\label{prop:sga3-VIB-6.2}%
\noindent\textbf{Proposition 6.2.---}%
\footnote{\label{note:sga3-VIB-51}Editor's note: we have modified
the statement as indicated in editor's \hyperref[note:sga3-VIB-49]{note 49}.}
\emph{Let \(G\) be an \(S\)-group. For a subfunctor \(T\) of the
functor \(G\), consider the following property:}
\[
  (+_{\!f})\qquad
  \text{for every \(s\in S\), \(T_s\) is representable by a
  closed subscheme of \(G_s\).}
\]
\emph{Let \(u,w:X\to G\) and \(v:Y\to G\) be morphisms of
\(S\)-schemes. Then:}
\begin{enumerate}
\item[\textup{(i)}]
  \(\underline{\operatorname{Transp}}(u,w)\) and
  \(\underline{\operatorname{Centr}}(u)
    =\underline{\operatorname{Transp}}(u,u)\)
  satisfy condition \((+_{\!f})\).
\item[\textup{(ii)}]
  \(\underline{\operatorname{Transp}}_G(u(X),v(Y))\) and
  \(\underline{\operatorname{Norm}}_G v(Y)\) satisfy condition
  \((+_{\!f})\) if, for every \(s\in S\), \(v_s\) is a closed
  immersion.
\item[\textup{\SGAInlineRefTarget{prop:sga3-VIB-6.2:item:iii}(iii)}]
  \(\underline{\operatorname{Transp}}^{\mathrm{str}}_G(X,Y)\)
  satisfies condition \((+_{\!f})\) if, for every \(s\in S\),
  \(u_s\) and \(v_s\) are closed immersions.
\end{enumerate}

This follows from
\hyperref[rem:sga3-VIB-6.1.1]{Remark~6.1.1} and
\hyperref[cor:sga3-VIB-6.2.5]{Corollary~6.2.5} below.%
\footnote{\label{note:sga3-VIB-52}Editor's note: the original
referred to the results of Expos\'e VIII, \S6. For the reader's
convenience, we have reproduced these results (with the exception
of VIII, \SGARefLink{sga3:viii:proposition:6:3}{6.3} and \SGARefLink{sga3:viii:remark:6:8}{6.8}) in \SGARefLink{def:sga3-VIB-6.2.1}{6.2.1}--\SGARefLink{cor:sga3-VIB-6.2.5}{6.2.5} below. Moreover, this had
been suggested by A.\ Grothendieck in a note at the beginning of
VIII.6: ``The present section is independent of the theory of
diagonalizable groups; its natural place would be in
VI\(_B\).''}

\medskip
\phantomsection\label{def:sga3-VIB-6.2.1}%
\noindent\textbf{Definition 6.2.1.---}
Let \(f:X\to S\) be a morphism of schemes. One says that \(f\) is
\emph{essentially free}, or that \(X\) is \emph{essentially free
over \(S\)}, if one can find a covering of \(S\) by affine open
subsets \(S_i\), for every \(i\) an affine \(S_i\)-scheme \(S'_i\)
faithfully flat over \(S_i\), and a covering
\((X'_{ij})_j\) of
\[
  X'_i=X\times_S S'_i
\]
by affine open subsets \(X'_{ij}\), such that, for every
\((i,j)\), the ring of \(X'_{ij}\) is a projective%
\footnote{\label{note:sga3-VIB-53}Editor's note: on the one hand,
we have replaced the word ``free'' by ``projective'', as
indicated in VIII, \SGARefLink{sga3:viii:remark:6:8}{Remark 6.8}. On the other hand, the notion of an
essentially free \(S\)-scheme is closely related to the geometric
notion of an \(S\)-scheme that is flat and pure, introduced and
studied by M.\ Raynaud and L.\ Gruson ([RG71]); see the addition
\hyperref[par:sga3-VIB-6.9]{6.9} below.}
module over the ring of \(S'_i\).

\medskip
\phantomsection\label{prop:sga3-VIB-6.2.2}%
\noindent\textbf{Proposition 6.2.2.---}
\begin{enumerate}
\item[\textup{a)}]
  \emph{If \(X\) is essentially free over \(S\), it is flat over
  \(S\); the converse holds if \(S\) is Artinian.}
\item[\textup{b)}]
  \emph{If \(S\) is the spectrum of a field, every \(S\)-scheme
  is essentially free over \(S\).}
\item[\textup{c)}]
  \emph{If \(X\) is essentially free over \(S\), then
  \(X'=X\times_S S'\) is essentially free over \(S'\), for every
  \(S'\to S\). The converse holds if \(S'\to S\) is faithfully
  flat and quasi-compact.}
\end{enumerate}

The proof is immediate, using for the converse in \textup{a)} the
fact that a flat module over an Artinian local ring is free.%
\footnote{\label{note:sga3-VIB-54}Editor's note: indeed, let
\((A,\mathfrak m)\) be an Artinian local ring, \(k\) its residue
field, \(M\) an arbitrary \(A\)-module, and let
\((x_i)_{i\in I}\) be elements of \(M\) whose images form a basis
of \(M/\mathfrak mM\) over \(k\). Let \(F\) be the free
\(A\)-module with basis \((e_i)_{i\in I}\), and let
\(\varphi:F\to M\) be the \(A\)-morphism defined by
\(\varphi(e_i)=x_i\). Then \(Q=\operatorname{Coker}\varphi\)
satisfies \(Q=\mathfrak mQ\), whence, since \(\mathfrak m\) is
nilpotent, \(Q=0\). Suppose moreover that \(M\) is flat over
\(A\); then \(K=\operatorname{Ker}\varphi\) satisfies
\(K\otimes_Ak=0\), that is, \(K=\mathfrak mK\), whence \(K=0\).}

The introduction of
\hyperref[def:sga3-VIB-6.2.1]{Definition~6.2.1} is justified by
the following theorem.

% END INCLUDED SGA3_English_current_buildable_source_20260823/source/inputs/VI/tex_reference_v2/components/52_expose_VIB_prop62_through_prop622_en.tex

% BEGIN INCLUDED SGA3_English_current_buildable_source_20260823/source/inputs/VI/tex_reference_v2/components/53_expose_VIB_theorem623_en.tex
% SGA 3, Expose VI_B, Loop 1 English reconstruction.
% Scope: Theorem 6.2.3 in full, including its proof and source
% editor's note 55.
% Authority: Polo--Gille Exp6B-13oct24.pdf, component-local page 40
% (combined-reader page 381; re-edition header 370).
% Examples 6.2.4 begin on local page 41 / combined page 382 and
% are excluded.
% OCR and Reinhold English Markdown are comparison witnesses only.
% This unit contains no diagrams.

\medskip
\phantomsection\label{thm:sga3-VIB-6.2.3}%
\noindent\textbf{Theorem 6.2.3.---}
\emph{Let \(S\) be a scheme, let \(Z\) be an \(S\)-scheme that
is essentially free, and let \(Y\) be a closed subscheme of
\(Z\). Consider the functor}
\[
  F=\prod_{Z/S}Y/Z:
  (\mathbf{Sch}_{/S})^\circ\longrightarrow\mathbf{Sets}
\]
\emph{defined by the following condition: \(F(S')=\varnothing\)
when \(Y_{S'}\ne Z_{S'}\), and \(F(S')\) is a singleton
otherwise.}% 
\footnote{\label{note:sga3-VIB-55}Editor's note: cf.\
Expos\'e~II, \S1, where this functor is denoted
\(\prod_{Z/S}Y\). For every \(S'\to S\),
\[
  F(S')=\Gamma(Y_{S'}/Z_{S'}),
\]
which here is equal to
\(\{\operatorname{id}_{Z_{S'}}\}\) if
\(Y_{S'}=Z_{S'}\), and is empty otherwise. We have also added
assertion \textup{\SGARefLink{concl:sga3-VIB-6.2.3-i}{(ii)}}, which is used in the proof of
\hyperref[cor:sga3-VIB-6.5.3]{6.5.3}.}
\begin{enumerate}
\item[\textup{\SGAInlineRefTarget{thm:sga3-VIB-6.2.3:item:i}(i)}]\phantomsection\label{concl:sga3-VIB-6.2.3-i}
  \emph{This functor is representable by a closed subscheme
  \(T\) of \(S\).}
\item[\textup{(ii)}]\phantomsection\label{concl:sga3-VIB-6.2.3-ii}
  \emph{If, moreover, \(Y\to Z\) is of finite presentation,
  then so is \(T\to S\).}
\end{enumerate}

\emph{Proof.}
First observe that the functor in question is a sheaf for the
fpqc topology. Since \(F(S')\) is either \(\varnothing\) or
\(\{\mathrm{pt}\}\) for every \(S'\), this amounts to checking
the following assertions. If \((S_i)\) is an open covering of
\(S\) and every
\[
  Y_{S_i}\longrightarrow Z_{S_i}
\]
is an isomorphism, then \(Y\to Z\) is an isomorphism; and if
\(S'\to S\) is faithfully flat and quasi-compact and
\[
  Y_{S'}\longrightarrow Z_{S'}
\]
is an isomorphism, then again \(Y\to Z\) is an isomorphism. The
first assertion is clear, and the second follows from SGA~1,
VIII, 5.4, or EGA~IV\(_2\), 2.7.1.

Moreover, by SGA~1, VIII, 1.9, faithfully flat and
quasi-compact morphisms are effective descent morphisms for the
fibered category of closed-immersion arrows. With the
notation of \hyperref[def:sga3-VIB-6.2.1]{6.2.1}, this allows us
to restrict to the case \(S=S'_i\).

Let \((Z_j)\) be a covering of \(Z\) by affine open subsets such
that
\[
  \mathcal O(Z_j)
\]
is a projective module over \(A=\mathcal O(S)\). Put
\[
  Y_j=Y\cap Z_j,
\]
and let
\[
  F_j:(\mathbf{Sch}_{/S})^\circ\longrightarrow\mathbf{Sets}
\]
be the functor defined in terms of \((Z_j,Y_j)\) as \(F\) was
defined in terms of \((Z,Y)\). It is a subfunctor of the final
functor, and clearly
\[
  F=\bigcap_j F_j.
\]
It is therefore enough to prove that every \(F_j\) is
representable by a closed subscheme \(T_j\) of \(S\), for then
\(F\) is represented by the closed subscheme
\[
  T=\bigcap_j T_j
\]
of \(S\).

We may consequently suppose that \(Z\) is itself affine, say
\[
  Z=\operatorname{Spec}B,
\]
where \(B\) is a projective \(A\)-module. Let \(L\) be a free
\(A\)-module with basis \((e_\lambda)_{\lambda\in\Lambda}\) of
which \(B\) is a direct summand as an \(A\)-module, and let
\[
  \phi_\lambda:L\longrightarrow A
\]
be the coordinate forms with respect to this basis. Let \(E\)
be a set of generators for the ideal \(J\) of \(B\) defining
the closed subscheme \(Y\) of \(Z\), and let \(I\) be the ideal
of \(A\) generated by the coordinates
\[
  \phi_\lambda(x),
  \qquad x\in E,\quad \lambda\in\Lambda.
\]

For every \(A\)-algebra \(C\), the morphism
\[
  B\otimes_A C\longrightarrow(B/J)\otimes_A C
\]
is an isomorphism if and only if the image of \(x\otimes1\) in
\(L\otimes_A C\) is zero for every \(x\in E\). This, in turn,
is equivalent to saying that the kernel of \(A\to C\) contains
the ideal \(I\). Hence
\[
  T=V(I)=\operatorname{Spec}(A/I)
\]
has the required property, proving \textup{\SGARefLink{thm:sga3-VIB-6.2.3:item:i}{(i)}}.

Finally, if \(Y\to Z\) is of finite presentation, the set \(E\)
may be chosen finite. Then \(I\) is a finitely generated ideal
of \(A\), that is, the closed immersion \(T\to S\) is of finite
presentation. This proves \textup{\SGARefLink{concl:sga3-VIB-6.2.3-i}{(ii)}}.

% END INCLUDED SGA3_English_current_buildable_source_20260823/source/inputs/VI/tex_reference_v2/components/53_expose_VIB_theorem623_en.tex

% BEGIN INCLUDED SGA3_English_current_buildable_source_20260823/source/inputs/VI/tex_reference_v2/components/54_expose_VIB_examples624_through_cor625_en.tex
% SGA 3, Expose VI_B, Loop 1 English reconstruction.
% Scope: Examples 6.2.4(a)--(e) and Corollary 6.2.5 in full,
% including source editor's notes 56--58.
% Authority: Polo--Gille Exp6B-13oct24.pdf, component-local pages
% 41--43 (combined-reader pages 382--384; displayed re-edition
% headers 371--373). Corollary 6.2.6 begins on local page 43 and
% is excluded.
% OCR and Reinhold English Markdown are comparison witnesses only.
% The displayed constructions are reconstructed natively in TeX;
% this unit contains no raster diagrams.

\medskip
\phantomsection\label{ex:sga3-VIB-6.2.4}%
\noindent\textbf{Examples 6.2.4.---}
Let us give important examples of functors which reduce to
functors
\[
  \prod_{Z/S}Y/Z
\]
of the type considered in
\hyperref[thm:sga3-VIB-6.2.3]{6.2.3}, and for which it will be
useful to have representability criteria later on. We denote by
\(S\) a scheme, and by \(X,Y,Z,\ldots\) schemes over \(S\).

\medskip
\phantomsection\label{ex:sga3-VIB-6.2.4-a}%
\noindent\textbf{a)}
Suppose given an \(S\)-morphism
\begin{equation}
  q:X\longrightarrow\underline{\operatorname{Hom}}_S(Y,Z),
  \tag{x}\label{eq:sga3-VIB-6.2.4-x}
\end{equation}
(``\(X\) acts on \(Y\), with values in \(Z\)''), that is, a
morphism
\begin{equation}
  r:X\times_S Y\longrightarrow Z.
  \tag{xx}\label{eq:sga3-VIB-6.2.4-xx}
\end{equation}
Consider a subscheme \(Z'\) of \(Z\), and hence a monomorphism
\[
  \underline{\operatorname{Hom}}_S(Y,Z')
  \longrightarrow
  \underline{\operatorname{Hom}}_S(Y,Z),
\]
which makes the first functor a subfunctor of the second. Let
\(X'\) be the inverse image of this subfunctor under \(q\). Thus
\(X'\) is the subfunctor of \(X\) for which \(X'(T)\) is the set
of \(x\in X(T)\) such that
\[
  q(x):Y_T\longrightarrow Z_T
\]
factors through \(Z'_T\).

This functor \(X'\) may be described as follows. Put
\[
  P=X\times_S Y,
\]
and let \(P'\) be the inverse image of \(Z'\) under
\(r:P\to Z\). Then there is an evident isomorphism
\begin{equation}
  X'\simeq\prod_{P/X}P'/P.
  \tag{xxx}\label{eq:sga3-VIB-6.2.4-xxx}
\end{equation}
It follows that, if \(Y\) is essentially free over \(S\) and
\(Z'\) is closed in \(Z\), the subfunctor \(X'\) of \(X\) is
representable by a closed subscheme of \(X\).%
\footnote{\label{note:sga3-VIB-56}Editor's note: we have added the
sentence that follows.}
If, moreover, \(Z'\to Z\) is of finite presentation, then so is
\(X'\to X\).

\medskip
\phantomsection\label{ex:sga3-VIB-6.2.4-b}%
\noindent\textbf{b)}
Suppose given two ways of letting \(X\) act on \(Y\), with values
in \(Z\), that is, two morphisms
\[
  q_1,q_2:X\rightrightarrows
  \underline{\operatorname{Hom}}_S(Y,Z),
\]
and put
\[
  X'=\operatorname{Ker}(q_1,q_2).
\]
Thus \(X'\) is the subfunctor of \(X\) for which \(X'(T)\) is the
set of \(x\in X(T)\) such that the two morphisms
\[
  q_1(x),q_2(x):Y_T\rightrightarrows Z_T
\]
are equal. The data of \(q_1,q_2\) are equivalent to the data of
a morphism
\[
  q:X\longrightarrow
  \underline{\operatorname{Hom}}_S(Y,Z\times_S Z),
\]
or, equivalently, of a morphism
\[
  r:X\times_S Y\longrightarrow Z\times_S Z.
\]
Put \(U=Z\times_S Z\), and let \(U'\) be the diagonal subscheme
of \(Z\times_S Z\). Then \(X'\) is precisely the inverse image
under \(q\) of the subfunctor
\[
  \underline{\operatorname{Hom}}_S(Y,U')
  \subset
  \underline{\operatorname{Hom}}_S(Y,U).
\]
It can therefore be put in the form
\hyperref[eq:sga3-VIB-6.2.4-xxx]{\textup{(xxx)}}, with
\(P=X\times_S Y\) and \(P'\) the inverse image of the diagonal
under \(r\), that is, the kernel of
\[
  r_1,r_2:X\times_S Y\rightrightarrows Z.
\]
We are thus under the hypotheses of
\hyperref[ex:sga3-VIB-6.2.4-a]{\textup{(a)}}.

Consequently, if \(Y\) is essentially free over \(S\) and \(Z\)
is separated over \(S\), the subfunctor \(X'\) of \(X\) is
representable by a closed subscheme of \(X\).%
\footnotemark[56]
If, moreover, \(Z\to S\) is locally of finite type, then
\(X'\to X\) is of finite presentation.

\medskip
\phantomsection\label{ex:sga3-VIB-6.2.4-c}%
\noindent\textbf{c)}
Suppose given a morphism
\[
  q:X\longrightarrow\underline{\operatorname{Hom}}_S(Y,Y),
\]
that is, ``\(X\) acts on \(Y\).'' Let \(X'\) be the ``kernel'' of
this morphism: \(X'\) is the subfunctor of \(X\) for which
\(X'(T)\) is the set of \(x\in X(T)\) such that
\[
  q(x):Y_T\longrightarrow Y_T
\]
is the identity. This functor falls under
\hyperref[ex:sga3-VIB-6.2.4-b]{\textup{(b)}}, as one sees by
introducing a second morphism
\[
  q':X\longrightarrow\underline{\operatorname{Hom}}_S(Y,Y)
\]
by letting \(X\) act trivially on \(Y\). Hence, if \(Y\) is
essentially free and separated over \(S\), the subfunctor \(X'\),
the kernel of \(q\), is representable by a closed subscheme of
\(X\).%
\footnotemark[56]
If, moreover, \(Y\to S\) is locally of finite type, then
\(X'\to X\) is of finite presentation.

\medskip
\phantomsection\label{ex:sga3-VIB-6.2.4-d}%
\noindent\textbf{d)}
Under the hypotheses of \hyperref[ex:sga3-VIB-6.2.4-c]{(c)},
consider the subfunctor \(Y'\) of \(Y\) ``of the invariants under
\(X\).'' Thus \(Y'(T)\) is the set of \(y\in Y(T)\) such that
the corresponding morphism
\[
  \bar q(y):X_T\longrightarrow Y_T
\]
is the constant \(T\)-morphism with value \(y\). Introducing
\(\bar q'\) as in \hyperref[ex:sga3-VIB-6.2.4-c]{(c)}, and the
morphisms corresponding to \(\bar q\) and \(\bar q'\),
\[
  \bar q,\bar q':Y\rightrightarrows
  \underline{\operatorname{Hom}}_S(X,Y),
\]
one sees that
\[
  Y'=\operatorname{Ker}(\bar q,\bar q').
\]
It therefore falls again under
\hyperref[ex:sga3-VIB-6.2.4-b]{\textup{(b)}}, with the roles of
\(X\) and \(Y\) reversed and with \(Z=Y\).

Consequently, if \(X\) is essentially free over \(S\) and \(Y\)
is separated over \(S\), the subfunctor \(Y'\) of \(Y\) of the
invariants under \(X\) is representable by a closed subscheme of
\(Y\).%
\footnotemark[56]
If, moreover, \(Y\to S\) is locally of finite type, then
\(Y'\to Y\) is of finite presentation.

\medskip
\phantomsection\label{ex:sga3-VIB-6.2.4-e}%
\noindent\textbf{e)}
Constructions of the kind described in the preceding examples
occur especially often in group theory. Thus, when \(G\) is an
\(S\)-group scheme acting on the \(S\)-scheme \(X\),
\[
  q:G\longrightarrow\underline{\operatorname{Aut}}_S(X),
\]
the kernel of \(q\), the subgroup of \(G\) acting trivially, is a
closed subscheme of \(G\), provided \(X\) is essentially free and
separated over \(S\)
(\hyperref[ex:sga3-VIB-6.2.4-c]{example~(c)}). The subobject
\(X^G\) of invariants is a closed subscheme of \(X\), provided
\(G\) is essentially free over \(S\) and \(X\) is separated over
\(S\) (\hyperref[ex:sga3-VIB-6.2.4-d]{example~(d)}).

Let \(Y,Z\) be subschemes of \(X\), and consider the subfunctor
\(\underline{\operatorname{Transp}}_G(Y,Z)\) of \(G\), the
transporter of \(Y\) into \(Z\). Its \(T\)-valued points, for
\(T\) over \(S\), are the elements \(g\in G(T)\) such that the
corresponding automorphism of \(X_T\) satisfies
\[
  g(Y_T)\subset Z_T,
\]
that is, induces a morphism \(Y_T\to X_T\) which factors through
\(Y_T\to Z_T\). Hence, if \(Y\) is essentially free over \(S\)
and \(Z\) is closed in \(X\), then
\(\underline{\operatorname{Transp}}_G(Y,Z)\) is a closed
subscheme of \(G\)
(\hyperref[ex:sga3-VIB-6.2.4-a]{example~(a)}).

One may also consider the strict transporter of \(Y\) into
\(Z\),%
\footnote{\label{note:sga3-VIB-57}Editor's note: denoted
\(\underline{\operatorname{Transp}}^{\mathrm{str}}_G(Y,Z)\).}
whose \(T\)-valued points, for \(T\) over \(S\), are the elements
\(g\in G(T)\) such that \(g(Y_T)=Z_T\). This functor is precisely
\[
  \underline{\operatorname{Transp}}_G(Y,Z)
  \cap
  \sigma\bigl(\underline{\operatorname{Transp}}_G(Z,Y)\bigr),
\]
where \(\sigma\) is the inversion morphism of \(G\). Consequently,
if \(Y\) and \(Z\) are essentially free over \(S\) and closed in
\(X\), the strict transporter of \(Y\) into \(Z\) is a closed
subscheme of \(G\).

An important case is \(X=G\), with \(G\) acting on itself by
inner automorphisms. If \(H\) is a subscheme of \(G\), the strict
transporter of \(H\) into \(H\) is also called the
\emph{normalizer of \(H\) in \(G\)}, and is denoted
\(\underline{\operatorname{Norm}}_G H\). Hence, if \(H\) is a
closed group subscheme of \(G\), essentially free over \(S\),
then \(\underline{\operatorname{Norm}}_G H\) is representable by
a closed group subscheme of \(G\).

Finally, let \(Z\) be a subscheme of \(G\). Its centralizer
\(\underline{\operatorname{Centr}}_G(Z)\) in \(G\) is the group
subfunctor of \(G\) defined by the procedure of
\hyperref[ex:sga3-VIB-6.2.4-d]{(d)}, when one regards \(Z\) as
acting on \(G\) by the operations induced by those of \(G\).
Thus, if \(Z\) is essentially free over \(S\) and \(G\) is
separated over \(S\), then
\(\underline{\operatorname{Centr}}_G(Z)\) is a closed group
subscheme of \(G\). In particular, if \(G\) is essentially free
and separated over \(S\), then the center \(C\) of \(G\), which
is none other than \(\underline{\operatorname{Centr}}_G(G)\), is
a closed group subscheme of \(G\).

When \(S\) is the spectrum of a field,
\hyperref[prop:sga3-VIB-6.2.2]{6.2.2\textup{(b)}} shows that the
``essentially free'' conditions in examples
\hyperref[ex:sga3-VIB-6.2.4-a]{(a)}--%
\hyperref[ex:sga3-VIB-6.2.4-e]{(e)} above are automatically
satisfied, leaving only the separation conditions. Recalling
that a group scheme over a field is necessarily separated
(\hyperref[sga3:VIA:0.3]{VI\(_A\), 0.3}), we obtain, for example:

\medskip
\phantomsection\label{cor:sga3-VIB-6.2.5}%
\noindent\textbf{Corollary 6.2.5.---}
\emph{Let \(G\) be a group scheme over a field \(k\), and let
\(Y,Y'\) be two subschemes of \(G\). Then:}
\begin{enumerate}
\item[\textup{(i)}]\phantomsection\label{concl:sga3-VIB-6.2.5-i}
  \emph{The centralizer of \(Y\) in \(G\) is a closed group
  subscheme of \(G\); in particular, this holds for the center
  \(\underline{\operatorname{Centr}}_G(G)\) of \(G\).}
\item[\textup{(i\(^{\prime}\))}]
  \phantomsection\label{concl:sga3-VIB-6.2.5-iprime}
  \emph{More generally, if \(u,v:X\to G\) are morphisms of
  schemes, then \(\underline{\operatorname{Transp}}_G(u,v)\) is
  representable by a closed subscheme of \(G\).}
\item[\textup{(ii)}]\phantomsection\label{concl:sga3-VIB-6.2.5-ii}
  \emph{If \(Y\) is closed, the transporter
  \(\underline{\operatorname{Transp}}_G(Y',Y)\) is a closed
  subscheme of \(G\). If \(Y'\) is also closed, the same
  conclusion holds for
  \(\underline{\operatorname{Transp}}^{\mathrm{str}}_G(Y',Y)\).}
\item[\textup{(iii)}]\phantomsection\label{concl:sga3-VIB-6.2.5-iii}
  \emph{For every group subscheme}%
  \footnote{\label{note:sga3-VIB-58}Editor's note: indeed, over
  a field \(k\), every group subscheme of \(G\) is closed; see
  \hyperref[sga3:VIA:0.5.2]{VI\(_A\), 0.5.2}. Moreover, \SGARefLink{cor:sga3-VIB-6.2.5}{6.2.5}
  completes the insertion of the results taken from VIII, \S6.}
  \emph{\(H\) of \(G\),
  \(\underline{\operatorname{Norm}}_G(H)\) is a closed group
  subscheme of \(G\).}
\end{enumerate}

% END INCLUDED SGA3_English_current_buildable_source_20260823/source/inputs/VI/tex_reference_v2/components/54_expose_VIB_examples624_through_cor625_en.tex

% BEGIN INCLUDED SGA3_English_current_buildable_source_20260823/source/inputs/VI/tex_reference_v2/components/55_expose_VIB_cor626_en.tex
% SGA 3, Expose VI_B, Loop 1 English reconstruction.
% Scope: Corollary 6.2.6 and its proof in full, including source
% editor's note 59.
% Authority: Polo--Gille Exp6B-13oct24.pdf, component-local pages
% 43--44 (combined-reader pages 384--385; re-edition headers
% 373--374; original-margin printed page 359 begins on local 44).
% Remark 6.3 begins on local page 44 / combined page 385 and is
% excluded.
% OCR and Reinhold English Markdown are comparison witnesses only.
% This unit contains no diagrams.

\medskip
\phantomsection\label{cor:sga3-VIB-6.2.6}%
\noindent\textbf{Corollary 6.2.6.---}%
\footnote{\label{note:sga3-VIB-59}Editor's note: we have inserted
this corollary here (cf.\ \SGARefLink{sga3:xii:6.1}{Expos\'e~XII, 6.1}), as it will be useful
later. We return in \SGARefLink{rem:sga3-VIB-6.3}{6.3} to the original text of
Expos\'e~VI\(_B\).}
\emph{Let \(k\) be a field and let \(G\) be a connected
algebraic \(k\)-group. Then
\(\underline{\operatorname{Centr}}_G(G)\) is representable by a
closed group subscheme \(Z\) of \(G\), and \(G/Z\) is an affine
algebraic \(k\)-group.}

\emph{Proof.}
The first assertion is, of course, contained in
\hyperref[cor:sga3-VIB-6.2.5]{6.2.5\textup{(i)}}, but we shall
see that it also follows from the proof of the second assertion.
Indeed, \(G\) acts by the adjoint representation on the
finite-dimensional \(k\)-vector spaces
\[
  V_n=\mathcal O_{G,e}/\mathfrak m_e^{\,n+1},
\]
where \(\mathfrak m_e\) is the maximal ideal of
\(\mathcal O_{G,e}\). Let \(K_n\) be the kernel of the morphism
\[
  \rho_n:G\longrightarrow\operatorname{GL}(V_n).
\]
By \hyperref[prop:sga3-VIA-5.4.1]{VI\(_A\), 5.4.1},
\(\rho_n\) induces a closed immersion
\[
  G/K_n\lhook\joinrel\longrightarrow\operatorname{GL}(V_n),
\]
so every \(G/K_n\) is affine. Since \(G\) is Noetherian, the
intersection \(K\) of the \(K_n\) is equal to one of the \(K_n\);
hence \(G/K\) is affine.

On the other hand, if \(Z\) denotes the center of \(G\), it is
clear that
\[
  Z\subset K.
\]
We shall prove the reverse inclusion. Let
\(\widehat{\mathcal O}_{G,e}\) be the completion of
\(\mathcal O_{G,e}\) for the \(\mathfrak m_e\)-adic topology,
and put
\[
  \widehat S=\operatorname{Spec}\widehat{\mathcal O}_{G,e},
  \qquad
  S=\operatorname{Spec}\mathcal O_{G,e}.
\]
Since \(\widehat S\to S\) is faithfully flat, and since the two
morphisms
\[
  K\times_k S\longrightarrow S,
  \qquad
  (g,x)\longmapsto gxg^{-1}
  \quad\text{and}\quad
  (g,x)\longmapsto x
\]
coincide after the base change \(\widehat S\to S\), they
coincide. Thus \(K\) acts trivially on \(\mathcal O_{G,e}\).

By \hyperref[ex:sga3-VIB-6.2.4-e]{6.2.4}\textup{(e)}, the subobject \(G^K\) of points of \(G\)
fixed by \(K\), which is precisely
\(\underline{\operatorname{Centr}}_G(K)\), is a closed
subscheme of \(G\). It is therefore defined by a quasi-coherent
ideal \(\mathcal I\) of \(\mathcal O_G\). Since \(G^K\) contains
\[
  S=\operatorname{Spec}\mathcal O_{G,e}
\]
and \(\mathcal I\) is of finite type (because \(G\) is
Noetherian), there is an open neighborhood \(U\) of \(e\) such
that
\[
  \mathcal I|_U=0.
\]
The subgroup
\[
  G^K=\underline{\operatorname{Centr}}_G(K)
\]
therefore contains \(U\), and hence also \(U\cdot U\), which is
equal to \(G\) because \(G\) is irreducible
(\hyperref[sga3:VIA:0.5]{VI\(_A\), 0.5}). Consequently,
\[
  \underline{\operatorname{Centr}}_G(K)=G.
\]
Thus \(K\subset Z\), and therefore \(Z=K\), as required.

% END INCLUDED SGA3_English_current_buildable_source_20260823/source/inputs/VI/tex_reference_v2/components/55_expose_VIB_cor626_en.tex

% BEGIN INCLUDED SGA3_English_current_buildable_source_20260823/source/inputs/VI/tex_reference_v2/components/56_expose_VIB_rem63_through_def65_en.tex
% SGA 3, Expose VI_B, Loop 1 English reconstruction.
% Scope: Remark 6.3, Proposition 6.4, Lemmas 6.4.2--6.4.3,
% and Definition 6.5 in full, including source editor's note 60.
% Authority: Polo--Gille Exp6B-13oct24.pdf, component-local pages
% 44--45 (combined-reader pages 385--386; re-edition headers
% 374--375; original-margin printed pages 359--360).
% Proposition 6.5.1 begins on local page 45 / combined page 386
% and is excluded.
% OCR and Reinhold English Markdown are comparison witnesses only.
% This unit contains no diagrams.

\medskip
\phantomsection\label{rem:sga3-VIB-6.3}%
\noindent\textbf{Remark 6.3.---}
\emph{Let \(k\) be an algebraically closed field, let \(G\) be a
\(k\)-group, and let \(H\) be a group subscheme of \(G\). Suppose
that \(G\) and \(H\) are of finite type over \(k\) and are
reduced. Then
\(\underline{\operatorname{Norm}}_G H\) (respectively,
\(\underline{\operatorname{Centr}}_G H\)) is representable by a
group subscheme of \(G\), whose associated reduced subscheme is
precisely the normalizer (respectively, centralizer) of \(H\) in
\(G\) in the sense of the \emph{Bible}.}

\medskip
\phantomsection\label{prop:sga3-VIB-6.4}%
\noindent\textbf{Proposition 6.4.---}
\emph{Let \(G\) be an \(S\)-group, and let \(u:X\to G\) be a
monomorphism of \(S\)-schemes. Set}
\[
  T=\underline{\operatorname{Transp}}_G(X,X).
\]
\emph{The following conditions are equivalent:}
\begin{enumerate}
\item[\textup{(i)}]\phantomsection\label{cond:sga3-VIB-6.4-i}
  \emph{\(T\) is a group sub-\(S\)-functor of \(G\).}
\item[\textup{(ii)}]\phantomsection\label{cond:sga3-VIB-6.4-ii}
  \emph{\[
    T=\underline{\operatorname{Transp}}^{\mathrm{str}}_G(X,X)
     =\underline{\operatorname{Norm}}_G X.
  \]}
\end{enumerate}
\emph{These conditions hold in each of the following two cases:}
\begin{enumerate}
\item[\textup{a)}]\phantomsection\label{case:sga3-VIB-6.4-a}
  \emph{\(X\) is of finite presentation over \(S\).}
\item[\textup{b)}]\phantomsection\label{case:sga3-VIB-6.4-b}
  \emph{\(T\) is representable by a scheme of finite
  presentation over \(S\).}
\end{enumerate}

The equivalence of \textup{(i)} and \textup{(ii)} follows from
the fact that, for every morphism \(S'\to S\) and all
\(t,t'\in T(S')\), one has
\[
  tt'\in T(S'),
\]
and from the identity
\[
  \underline{\operatorname{Transp}}^{\mathrm{str}}_G(X,X)
  =T\cap c(T)
\]
(cf.\ \hyperref[def:sga3-VIB-6.1]{6.1\textup{(ii)}}).

Suppose first that \textup{a)} holds. If \(t\in T(S)\), then
\(\operatorname{int}(t)\) is a monomorphism from \(X\) to itself,
and therefore an \(S\)-automorphism of \(X\) by EGA~IV\(_4\),
17.9.6. Thus
\[
  t\in\underline{\operatorname{Transp}}^{\mathrm{str}}_G(X,X),
\]
which proves the assertion in case \textup{a)}.

In case \textup{b)}, it is clear that
\[
  \mu(T\times_S T)\subset T,
\]
and the assertion follows from the next lemma.

\medskip
\phantomsection\label{lem:sga3-VIB-6.4.2}%
\noindent\textbf{Lemma 6.4.2.---}%
\footnote{\label{note:sga3-VIB-60}Editor's note: we have retained
the numbering of the original; there is no number 6.4.1.}
\emph{Let \(G\) be an \(S\)-scheme of finite presentation,
equipped with an associative law (in the sense of
EGA~0\(_{\mathrm{III}}\), 8.2.5). Suppose that, for every
\(S\)-scheme \(S'\) and every \(g\in G(S')\), right and left
translation by \(g\) in the set \(G(S')\) are injective, and
suppose that \(G(S)\ne\varnothing\). Then \(G\) is an
\(S\)-group.}

\emph{Proof.}
It is enough to prove that \(G(S')\) is a group for every
\(S\)-scheme \(S'\). The hypothesis immediately implies that
right and left translation by every \(g\in G(S')\) on
\(G_{S'}\) are \(S'\)-monomorphisms
\[
  G_{S'}\longrightarrow G_{S'}.
\]
They are therefore \(S'\)-automorphisms, since \(G\) is of
finite presentation over \(S\) (EGA~IV\(_4\), 17.9.6).
Consequently, right and left translation by \(g\) in the set
\(G(S')\) are bijective, and the assertion reduces to the
following lemma.

\medskip
\phantomsection\label{lem:sga3-VIB-6.4.3}%
\noindent\textbf{Lemma 6.4.3.---}
\emph{Let \(G\) be a nonempty set with an associative law such
that right and left translations are bijective. Then \(G\) is a
group.}

The proof is left to the reader.

\medskip
\phantomsection\label{def:sga3-VIB-6.5}%
\noindent\textbf{Definition 6.5.---}
\emph{Let \(G\) be an \(S\)-group, let \(H\) be an
\(S\)-functor, and let \(u:H\to G\) be a monomorphism.}
\begin{enumerate}
\item[\textup{(i)}]
  The \emph{connected centralizer of \(H\) in \(G\)}, denoted
  \[
    \underline{\operatorname{Centr}}^0_G H,
  \]
  is the identity component of the functor
  \(\underline{\operatorname{Centr}}_G H\)
  (cf.\ \hyperref[def:sga3-VIB-3.1]{3.1} and
  \hyperref[def:sga3-VIB-6.1]{6.1\textup{(iii)}}).
\item[\textup{(i\(^{\prime}\))}]
  For every morphism \(u:X\to G\), one similarly defines the
  functor
  \[
    \underline{\operatorname{Centr}}^0(u)
  \]
  (cf.\ \SGARefLink{prop:sga3-VIB-6.2}{6.2}\textup{(iv)}).
\item[\textup{(ii)}]
  When \(u_s\) is a closed immersion for every \(s\in S\), the
  \emph{connected normalizer of \(H\) in \(G\)}, denoted
  \[
    \underline{\operatorname{Norm}}^0_G H,
  \]
  is the identity component of the functor
  \(\underline{\operatorname{Norm}}_G H\).
\end{enumerate}

\noindent\emph{N.B.}
By \hyperref[cor:sga3-VIB-1.4.2]{1.4.2}, the hypothesis in
\textup{\SGARefLink{cor:sga3-VIB-1.4.2:item:ii}{(ii)}} is satisfied when \(H\) is an \(S\)-group scheme,
\(G\) and \(H\) are locally of finite type over \(S\), and
\(u\) is a quasi-compact morphism of \(S\)-groups.

% END INCLUDED SGA3_English_current_buildable_source_20260823/source/inputs/VI/tex_reference_v2/components/56_expose_VIB_rem63_through_def65_en.tex

% BEGIN INCLUDED SGA3_English_current_buildable_source_20260823/source/inputs/VI/tex_reference_v2/components/57_expose_VIB_prop651_en.tex
% SGA 3, Expose VI_B, Loop 1 English reconstruction.
% Scope: Proposition 6.5.1 in full and source editor's note 61.
% Authority: Polo--Gille Exp6B-13oct24.pdf, component-local page
% 45 (combined-reader page 386; re-edition header 375;
% original-margin printed page 361).
% Theorem 6.5.2 begins on local page 46 / combined page 387
% and is excluded.
% OCR and Reinhold English Markdown are comparison witnesses only.
% This unit contains no diagrams.

\medskip
\phantomsection\label{prop:sga3-VIB-6.5.1}%
\noindent\textbf{Proposition 6.5.1.---}
\emph{Let \(G\) be an \(S\)-group locally of finite presentation
and quasi-separated over \(S\), let \(H\) be a smooth \(S\)-group
with connected fibers, and let \(u:H\to G\) be a monomorphism.
Let \(N\) be the normalizer of \(H\) in \(G\)
(cf.\ \hyperref[def:sga3-VIB-6.1]{6.1}). By
\hyperref[cor:sga3-VIB-6.5.5]{6.5.5} below, \(N\) is representable
by a closed group subscheme of \(G\), of finite presentation over
\(G\). Under these circumstances, the following conditions are
equivalent:}
\begin{enumerate}
\item[\textup{\SGAInlineRefTarget{prop:sga3-VIB-6.5.1:item:i}(i)}]\phantomsection\label{cond:sga3-VIB-6.5.1-i}
  \emph{The canonical morphism \(H\to N\) is an open immersion.}
\item[\textup{\SGAInlineRefTarget{prop:sga3-VIB-6.5.1:item:ii}(ii)}]\phantomsection\label{cond:sga3-VIB-6.5.1-ii}
  \emph{\(N^0=H\)
  (cf.\ \hyperref[thm:sga3-VIB-3.10]{3.10}).}
\item[\textup{\SGAInlineRefTarget{prop:sga3-VIB-6.5.1:item:iii}(iii)}]\phantomsection\label{cond:sga3-VIB-6.5.1-iii}
  \emph{For every \(s\in S\), one has
  \(H_s=(N_s)^0\).}
\end{enumerate}

Condition \textup{\SGARefLink{prop:sga3-VIB-6.5.1:item:i}{(i)}} implies \textup{\SGARefLink{prop:sga3-VIB-6.5.1:item:ii}{(ii)}} by
\hyperref[lem:sga3-VIB-3.10.1]{Lemma 3.10.1}, since \(H^0=H\).
It is also clear that \textup{\SGARefLink{prop:sga3-VIB-6.5.1:item:ii}{(ii)}} implies \textup{\SGARefLink{prop:sga3-VIB-6.5.1:item:iii}{(iii)}}, for
\[
  H_s=(N^0)_s=(N_s)^0.
\]

Finally, let us prove that \textup{\SGARefLink{prop:sga3-VIB-6.5.1:item:iii}{(iii)}} implies \textup{\SGARefLink{prop:sga3-VIB-6.5.1:item:i}{(i)}}.
Since \(H_s=(N^0)_s\) for every \(s\in S\), the morphism
\[
  H_s\longrightarrow N_s
\]
is an open immersion. Moreover, \(H\) and \(N\) are locally of
finite presentation over \(S\), and \(H\) is flat over \(S\).
Consequently, \(H\to N\) is an open immersion
(EGA~IV\(_4\), 17.9.5).

For the reader's convenience, the results \SGARefLink{sga3:xi:theorem:6:8}{6.8}--6.11 of
Exposé XI have been included below.\footnote{%
\label{note:sga3-VIB-61}Editor's note: \hyperref[thm:sga3-VIB-6.5.2]{Sections 6.5.2}--\SGARefLink{cor:sga3-VIB-6.5.5}{6.5.5},
taken from \SGARefLink{sga3:xi:theorem:6:8}{Exposé XI, 6.8}--6.11, have been inserted here. This
was in fact suggested by A.\ Grothendieck in a note at the
beginning of XI~6: ``The present section does not use the results
of \SGARefLink{sga3:xi:section:3}{Sections 3}, 4, and 5 [of XI]; its natural place would be in
VI\textsubscript{B}.''}

% END INCLUDED SGA3_English_current_buildable_source_20260823/source/inputs/VI/tex_reference_v2/components/57_expose_VIB_prop651_en.tex

% BEGIN INCLUDED SGA3_English_current_buildable_source_20260823/source/inputs/VI/tex_reference_v2/components/58_expose_VIB_thm652_cor653_en.tex
% SGA 3, Expose VI_B, Loop 1 English reconstruction.
% Scope: Theorem 6.5.2 and Corollary 6.5.3 in full, including
% their complete proofs and source editor's notes 62--64.
% Authority: Polo--Gille Exp6B-13oct24.pdf, component-local pages
% 46--47 (combined-reader pages 387--388; displayed re-edition
% headers 376--377). Corollary 6.5.4 begins on local page 47 and
% is excluded.
% OCR and Reinhold English Markdown are comparison witnesses only.
% This unit contains no diagrams.

\medskip
\phantomsection\label{thm:sga3-VIB-6.5.2}%
\noindent\textbf{Theorem 6.5.2.---}
\emph{Let \(X\) be a smooth scheme over \(S\), with geometrically
irreducible fibers.}%
\footnote{\label{note:sga3-VIB-62}Editor's note: on the one hand,
we have corrected the original by replacing ``connected'' with
``irreducible'' (see the proof). On the other hand, according to
[RG71], I, 3.3.4\textup{(iii)} and 4.1.1, it is enough to assume
that \(X\) is flat over \(S\), with geometrically irreducible
fibers and without embedded components.}
\begin{enumerate}
\item[\textup{(i)}]\phantomsection\label{concl:sga3-VIB-6.5.2-i}
  \emph{For every closed subscheme \(Y\) of \(X\), the functor
  \[
    \prod_{X/S}Y/X
  \]
  is representable by a closed subscheme \(T\) of \(S\).}
\item[\textup{(ii)}]\phantomsection\label{concl:sga3-VIB-6.5.2-ii}
  \emph{Moreover, if \(Y\to X\) is of finite presentation, then
  so is \(T\to S\).}
\end{enumerate}

\emph{Proof.}
Let \(f:X\to S\) denote the structure morphism. Since \(f\) is
faithfully flat and locally of finite
presentation, it is a covering for the fpqc topology. Moreover,
\[
  T=\prod_{X/S}Y/X
\]
is evidently a subsheaf of \(S\) for the fpqc topology. It follows
that the representability of \(T\) by a closed subscheme of \(S\)
is a question local on \(S\) for the fpqc topology,%
\footnote{\label{note:sga3-VIB-63}Editor's note: see, for example,
the addition 1.7 in Expos\'e VIII.}
and the same is true of the question whether \(T\) is of finite
presentation over \(S\) (EGA~IV\(_2\), 2.7.1).

After the base change \(S'\to S\) with \(S'=X\), we are reduced
to the case in which \(X\) admits a section
\[
  \epsilon:S\longrightarrow X.
\]
We may also suppose that \(S\) is affine, and hence
quasi-compact. We then have:

\medskip
\phantomsection\label{cor:sga3-VIB-6.5.3}%
\noindent\textbf{Corollary 6.5.3.---}
\emph{Under the hypotheses of
\hyperref[thm:sga3-VIB-6.5.2]{6.5.2}, suppose that \(S\) is
quasi-compact, that \(X\to S\) admits a section \(\epsilon\), and
that \(Y\to X\) is of finite presentation. Then there exists an
integer \(n\geq0\) such that}
\phantomsection\label{eq:sga3-VIB-6.5.3-infinitesimal}%
\[
  \prod_{X/S}Y/X
  =
  \prod_{X_n/S}Y_n/X_n,
\]
\emph{where \(X_n\) is the \(n\)-th infinitesimal neighborhood of
the immersion \(\epsilon:S\to X\), and
\(Y_n=Y\cap X_n\).}

When \(Y\) is of finite presentation over \(X\), this corollary
does indeed imply
\hyperref[thm:sga3-VIB-6.5.2]{6.5.2}. Since \(X\) is smooth over
\(S\), \(X_n\) is finite and locally free over \(S\), and is
therefore, a fortiori, essentially free over \(S\)
(see \hyperref[def:sga3-VIB-6.2.1]{6.2.1}). Hence
\[
  \prod_{X_n/S}Y_n/X_n
\]
is representable by a closed subscheme \(T\) of \(S\), of finite
presentation over \(S\), by
\hyperref[thm:sga3-VIB-6.2.3]{6.2.3}.

Let us first prove \hyperref[cor:sga3-VIB-6.5.3]{6.5.3}, and
hence \hyperref[thm:sga3-VIB-6.5.2]{6.5.2}, when \(S\) is
Noetherian. Put
\phantomsection\label{eq:sga3-VIB-6.5.3-Tn}%
\[
  T_n=\prod_{X_n/S}Y_n/X_n.
\]
The \(T_n\) form a decreasing sequence of closed subschemes of
\(S\), and since \(S\) is Noetherian this sequence is stationary.
Let
\phantomsection\label{eq:sga3-VIB-6.5.3-R}%
\[
  R=\bigcap_{n\geq0}\prod_{X_n/S}Y_n/X_n
\]
be their common value for all sufficiently large \(n\). Clearly
\(T\subset R\), and it is enough to prove \(R\subset T\).

After the base change \(R\to S\), we are reduced to the case
\[
  R=S,
\]
that is,
\[
  Y_n=X_n\quad\text{for every \(n\)},
  \qquad\text{or equivalently}\qquad
  Y\supset X_n\quad\text{for every \(n\)}.
\]
We must then prove \(T=S\), that is, \(Y=X\).

Now \(Y\supset X_n\) for every \(n\), together with the fact that
\(X\) is locally Noetherian, implies that in a neighborhood of
each point of \(\epsilon(S)\), \(Y\) is an induced open subscheme
of \(X\).%
\footnote{\label{note:sga3-VIB-64}Editor's note: see, for example,
the proof of \hyperref[cor:sga3-VIB-6.2.6]{6.2.6}.}
Consequently, there exists an induced open subscheme \(U\) of
\(X\), containing \(\epsilon(S)\), such that \(U\subset Y\). By
EGA~IV\(_3\), 11.10.10, since the fibers of \(X/S\) are integral,
\(U\) is schematically dense in \(X\). Since \(Y\) is a closed
subscheme containing \(U\), it follows that \(Y=X\). This proves
\hyperref[cor:sga3-VIB-6.5.3]{6.5.3}, and hence
\hyperref[thm:sga3-VIB-6.5.2]{6.5.2}, when \(S\) is Noetherian.

The general case is reduced to the preceding one. For every
\(s\in S\), there exist an affine open neighborhood \(U\) of
\(s\) and an affine open neighborhood \(V\) of \(\epsilon(s)\)
such that \(f(V)\subset U\). Then \(f(V)\) is an open
neighborhood of \(s\) contained in \(U\). Let \(S'\) be an affine
open neighborhood of \(s\) contained in
\[
  \epsilon^{-1}(V)\cap f(V),
\]
and put
\[
  X'=V\cap f^{-1}(S').
\]
Then \(X'\) and \(S'\) are affine open subsets of \(X\) and \(S\),
respectively, and \(X'/S'\) admits a section. By the local nature
of \SGARefLink{thm:sga3-VIB-6.5.2}{6.5.2} and \SGARefLink{cor:sga3-VIB-6.5.3}{6.5.3}, we may suppose that \(S=S'\).

Thus \(X'\) is an affine open subset of \(X\) containing
\(\epsilon(S)\). Since every fiber \(X_s\) is irreducible,
\(X'_s\) is schematically dense in \(X_s\). Consequently, by
EGA~IV\(_3\), 11.10.10, \(X'\) is schematically dense in \(X\);
likewise, after every base change \(S_1\to S\),
\[
  X'\times_S S_1
\]
is schematically dense in \(X\times_S S_1\).

It follows that
\phantomsection\label{eq:sga3-VIB-6.5.3-open-restriction}%
\[
  \prod_{X/S}Y/X
  =
  \prod_{X'/S}Y'/X',
  \qquad Y'=Y\cap X'.
\]
This reduces us to the case \(X=X'\), so we may suppose that both
\(S\) and \(X\) are affine. Moreover, if
\[
  X=\operatorname{Spec}B
\]
and \(J\) is the ideal of \(B\) defining \(Y\), then \(J\) is the
filtered colimit of its finitely generated subideals. Thus \(Y\)
is the intersection of closed subschemes \(Y_i\) of \(X\) that
are of finite presentation over \(X\), and consequently
\phantomsection\label{eq:sga3-VIB-6.5.3-finite-presentation-reduction}%
\[
  \prod_{X/S}Y/X
  =
  \bigcap_i\prod_{X/S}Y_i/X.
\]
To prove \hyperref[thm:sga3-VIB-6.5.2]{6.5.2}, we are therefore
reduced to the case in which \(Y\) is of finite presentation over
\(X\), with \(S\) and \(X\) affine.

In that case \(X\) and \(Y\) over \(S\) are obtained by base
change \(S\to S_0\) from an analogous situation \(X_0,Y_0\) over
a Noetherian scheme \(S_0\). This reduces us to the Noetherian
case already treated, and completes the proof of
\hyperref[thm:sga3-VIB-6.5.2]{6.5.2} and
\hyperref[cor:sga3-VIB-6.5.3]{6.5.3}.

% END INCLUDED SGA3_English_current_buildable_source_20260823/source/inputs/VI/tex_reference_v2/components/58_expose_VIB_thm652_cor653_en.tex

% BEGIN INCLUDED SGA3_English_current_buildable_source_20260823/source/inputs/VI/tex_reference_v2/components/59_expose_VIB_cor654_loop2_r1_en.tex
% SGA 3, Expose VI_B, Loop 1 English reconstruction.
% Scope: Corollary 6.5.4 and its proof in full, including source
% editor's note 65 and the source-bound commutative triangle.
% Authority: Polo--Gille Exp6B-13oct24.pdf, component-local pages
% 47--48 (combined-reader pages 388--389; re-edition headers
% 377--378).
% Corollary 6.5.5 begins on local page 48 and is excluded.
% OCR and Reinhold English Markdown are comparison witnesses only.
% Loop 2 replaces the source-bound Loop-1 PNG with a faithful
% native LaTeX reconstruction; the replacement is closed.

\medskip
\phantomsection\label{cor:sga3-VIB-6.5.4}%
\noindent\textbf{Corollary 6.5.4.---}
\emph{Let \(X\) be a smooth \(S\)-group scheme of finite
presentation with connected fibers, let \(Y\) be a group scheme
of finite presentation over \(S\), and let \(i:Y\to X\) be a
monomorphism of \(S\)-group schemes.}
\begin{enumerate}
\item[\textup{(i)}]\phantomsection\label{concl:sga3-VIB-6.5.4-i}
  \emph{The functor
  \[
    \prod_{X/S}Y/X
  \]
  is representable by a closed subscheme of \(S\) of finite
  presentation.}
\item[\textup{(ii)}]\phantomsection\label{concl:sga3-VIB-6.5.4-ii}
  \emph{If, moreover, \(S\) is quasi-compact, then for all
  sufficiently large \(n\),
  \[
    \prod_{X/S}Y/X
      =\prod_{X_n/S}Y_n/X_n,
  \]
  where \(X_n\) denotes the \(n\)-th infinitesimal neighborhood
  of the unit section \(\epsilon:S\to X\), and
  \(Y_n=X_n\cap Y\).}
\end{enumerate}

The proof is essentially that of
\hyperref[cor:sga3-VIB-6.5.3]{6.5.3}.%
\footnote{\label{note:sga3-VIB-65}Editor's note: the original has
been expanded in what follows, and references to EGA~IV have
been added.}
On the one hand, \(i\) is locally of finite presentation
(EGA~IV\(_1\), 1.4.3). On the other hand, the unit sections of
\(Y\) and \(X\) induce bijective immersions
\[
  S\longrightarrow Y_n
  \qquad\text{and}\qquad
  S\longrightarrow X_n,
\]
and hence isomorphisms of \(S_{\mathrm{red}}\) with
\((Y_n)_{\mathrm{red}}\) and \((X_n)_{\mathrm{red}}\).
Consequently, \(i_n\) is quasi-compact, hence of finite type, and
one has the commutative diagram
\begin{figure}[H]
  \centering
  \phantomsection\label{fig:sga3-VIB-6.5.4-reduced-triangle}%
  \begin{tikzcd}[row sep=large,column sep=large]
    (Y_n)_{\mathrm{red}}
      \arrow[r,"\tau"]
      \arrow[dr,"\sigma"'] &
    Y_n \arrow[d,"i_n"'] \\
    {} &
    X_n
  \end{tikzcd}
\end{figure}
where \(\sigma\) and \(\tau\) are closed immersions and \(\tau\)
is surjective. Since \(i_n\) is separated (being a
monomorphism), it follows that \(i_n\) is proper
(EGA~II, 5.4.3). Thus \(i_n\) is a proper monomorphism of finite
presentation, and therefore a closed immersion
(EGA~IV\(_3\), 8.11.5). It follows from
\hyperref[thm:sga3-VIB-6.2.3]{6.2.3} that
\[
  \prod_{X_n/S}Y_n/X_n
\]
is representable by a closed subscheme \(T_n\) of \(S\), of
finite presentation over \(S\). It remains to prove the last
assertion of \hyperref[cor:sga3-VIB-6.5.4]{6.5.4} when \(S\) is assumed affine as well.

We again reduce immediately to the case where \(S\) is
Noetherian. With the notation of the proof of \hyperref[cor:sga3-VIB-6.5.3]{6.5.3}, it is enough
to prove \(R=T\), or equivalently that
\[
  Y\supset X_n\quad\text{for every }n
  \qquad\Longrightarrow\qquad
  Y=X.
\]
The hypothesis implies that \(i:Y\to X\) is étale at the points
of the unit section of \(Y\) over \(S\). Thus \(Y\) is smooth
over \(S\) at those points, and the open subset \(U\) of points
of \(Y\) at which \(Y\) is smooth over \(S\) is an induced open
subgroup of \(Y\)
(cf.\ \hyperref[cor:sga3-VIB-2.3]{2.3}).
\begin{samepage}
The morphism
\[
  \tau:U\longrightarrow X
\]
is then an étale monomorphism by
\hyperref[prop:sga3-VIB-2.5]{2.5}, and hence an open immersion.
\end{samepage}
Since the fibers of \(X\) are connected and every open subgroup
of an algebraic group is also closed, \(\tau\) is a surjective
open immersion, that is, an isomorphism. Therefore \(U=X\), and
a fortiori \(Y=X\). This completes the proof of \hyperref[cor:sga3-VIB-6.5.4]{6.5.4}.

Proceeding as in
\hyperref[ex:sga3-VIB-6.2.4-e]{6.2.4\textup{(e)}}, one deduces
\hyperref[cor:sga3-VIB-6.5.5]{Corollary 6.5.5} from \hyperref[cor:sga3-VIB-6.5.4]{6.5.4}.

% END INCLUDED SGA3_English_current_buildable_source_20260823/source/inputs/VI/tex_reference_v2/components/59_expose_VIB_cor654_loop2_r1_en.tex

% BEGIN INCLUDED SGA3_English_current_buildable_source_20260823/source/inputs/VI/tex_reference_v2/components/60_expose_VIB_cor655_en.tex
% SGA 3, Expose VI_B, Loop 1 English reconstruction.
% Scope: Corollary 6.5.5 and its proof in full, including source
% editor's note 66.
% Authority: Polo--Gille Exp6B-13oct24.pdf, component-local pages
% 48--49 (combined-reader pages 389--390; re-edition headers
% 378--379).
% Definition 6.6 begins on local page 49 and is excluded.
% OCR and Reinhold English Markdown are comparison witnesses only.
% This unit contains no diagrams.

\medskip
\phantomsection\label{cor:sga3-VIB-6.5.5}%
\noindent\textbf{Corollary 6.5.5.---}
\emph{Let \(G\) be an \(S\)-group scheme locally of finite type
and quasi-separated over \(S\), let \(H\) be a smooth group
scheme of finite presentation over \(S\) with connected fibers,
and let \(i:H\to G\) be a monomorphism of \(S\)-groups. Then:}
\begin{enumerate}
\item[\textup{a)}]\phantomsection\label{concl:sga3-VIB-6.5.5-a}
  \emph{\(\underline{\operatorname{Centr}}_G(H)\) and
  \(\underline{\operatorname{Norm}}_G(H)\) are representable by
  closed subschemes of \(G\), of finite presentation over \(G\).}
\item[\textup{a\(^{\prime}\))}]
  \phantomsection\label{concl:sga3-VIB-6.5.5-aprime}
  \emph{Likewise, for every monomorphism \(j:K\to G\) of
  \(S\)-group schemes that is of finite presentation,
  \(\underline{\operatorname{Transp}}_G(H,K)\) is representable
  by a closed subscheme of \(G\), of finite presentation over
  \(G\).}
\item[\textup{b)}]\phantomsection\label{concl:sga3-VIB-6.5.5-b}
  \emph{If \(G\) is quasi-compact, there exists an integer
  \(n\geq 0\) such that, if \(H_n\) denotes the \(n\)-th
  infinitesimal neighborhood of the unit section of \(H\), then}
  \[
    \underline{\operatorname{Centr}}_G(H)
      =\underline{\operatorname{Centr}}_G(H_n),
    \qquad
    \underline{\operatorname{Norm}}_G(H)
      =\underline{\operatorname{Norm}}_G(H_n),
  \]
  \[
    \underline{\operatorname{Transp}}_G(H,K)
      =\underline{\operatorname{Transp}}_G(H_n,K)
      =\underline{\operatorname{Transp}}_G(H_n,K_n).
  \]
\end{enumerate}

\emph{Proof.}%
\footnote{\label{note:sga3-VIB-66}Editor's note: the argument that
follows has been expanded. This also completes the insertion of
XI, \SGARefLink{sga3:xi:theorem:6:8}{6.8}--6.11; at \SGARefLink{def:sga3-VIB-6.6}{6.6} below we return to the original text of
VI\textsubscript{B}.}
First observe that the hypothesis on \(G\) implies that the
monomorphism \(H\to G\) is of finite presentation
(EGA~IV\(_1\), 1.2.4 and 1.4.3), as is the diagonal immersion
\[
  \Delta_{G/S}:G\longrightarrow G\times_S G
\]
(loc.\ cit., 1.4.3.1). The case of
\(\underline{\operatorname{Norm}}_G(H)\) therefore reduces to
that of the transporter, by taking \(H=K\).

In view of
\hyperref[ex:sga3-VIB-6.2.4-e]{6.2.4\textup{(e)}}, we apply
\hyperref[cor:sga3-VIB-6.5.4]{6.5.4} to the group scheme
\[
  X=H_G=H\times_S G
\]
over the base scheme \(G\), and to the following group subscheme
\(Y\).

In the case of
\(\underline{\operatorname{Transp}}_G(H,K)\), take \(Y\) to be
the inverse image of \(K_G\) under the morphism of \(G\)-groups
\[
  H\times_S G\longrightarrow G\times_S G,
  \qquad
  (h,g)\longmapsto(ghg^{-1},g).
\]
In the case of
\(\underline{\operatorname{Centr}}_G(H)\), take \(Y\) to be the
inverse image of the diagonal subgroup
\[
  \Delta_{G/S}(G)_G
\]
of \((G\times_S G)_G\) under the morphism of \(G\)-groups
\[
  H\times_S G\longrightarrow G\times_S G\times_S G,
  \qquad
  (h,g)\longmapsto(h,ghg^{-1},g).
\]

% END INCLUDED SGA3_English_current_buildable_source_20260823/source/inputs/VI/tex_reference_v2/components/60_expose_VIB_cor655_en.tex

% BEGIN INCLUDED SGA3_English_current_buildable_source_20260823/source/inputs/VI/tex_reference_v2/components/61_expose_VIB_def66_through_par69_en.tex
% SGA 3, Expose VI_B, Loop 1 English reconstruction.
% Scope: Definition 6.6, Remark 6.7, Example 6.8, and paragraph
% 6.9 in full, including source editor's note 67.
% Authority: Polo--Gille Exp6B-13oct24.pdf, component-local page
% 49 (combined-reader page 390; re-edition header 379;
% original-margin printed page 362).
% Section 7 begins on local page 50 and is excluded.
% OCR and Reinhold English Markdown are comparison witnesses only.
% This unit contains no diagrams.

\medskip
\phantomsection\label{def:sga3-VIB-6.6}%
\noindent\textbf{Definition 6.6.---}
Let \(G\) be an \(S\)-functor in groups and let \(H\) be a group
sub-\(S\)-functor of \(G\). One says that \(H\) is \emph{invariant}
(respectively, \emph{central}, respectively,
\emph{characteristic}) in \(G\) if
\[
  \underline{\operatorname{Norm}}_G H=G
\]
(respectively, if
\(\underline{\operatorname{Centr}}_G H=G\), respectively, if for
every \(S\)-scheme \(T\) and every automorphism
\[
  a\in\operatorname{Aut}_{T\text{-}\mathrm{gr}}(G_T)
\]
one has \(a(H_T)\subset H_T\)). In other words, for every
\(S\)-scheme \(T\), the subgroup \(H(T)\) of \(G(T)\) is normal
in \(G(T)\) (respectively, central in \(G(T)\), respectively,
invariant under every automorphism of \(G_T\)).

\noindent\emph{N.B.}
If \(H\) is central (respectively, characteristic), then it is
invariant.

\medskip
\phantomsection\label{rem:sga3-VIB-6.7}%
\noindent\textbf{Remark 6.7.---}
Let \(G\) and \(H\) be two \(S\)-groups and let \(u:H\to G\) be
a monomorphism. The group \(H\) is invariant (respectively,
central) in \(G\) if and only if the morphism
\[
  \mu\circ
  \bigl(c\circ\operatorname{pr}_2,\,
        \mu\circ(u\times\operatorname{id}_G)\bigr)
  :H\times_S G\longrightarrow G
\]
(defined by
\[
  (h,g)\longmapsto g^{-1}hg
\]
for \(g\in G(S')\) and \(h\in H(S')\)) factors through \(u\)
(respectively, is equal to \(u\circ\operatorname{pr}_1\)).
The group \(H\) is characteristic in \(G\) if and only if, for
every \(S\)-scheme \(T\) and every \(T\)-group automorphism \(a\)
of \(G_T\), the morphism \(a\circ u(T)\) factors through \(u(T)\).

\medskip
\phantomsection\label{ex:sga3-VIB-6.8}%
\noindent\textbf{Example 6.8.---}
Let \(G\) be an \(S\)-functor in groups. Then
\(\underline{\operatorname{Centr}}G\) is characteristic and
central. If \(G_s\) is representable for every \(s\in S\), then
\(G^0\) is characteristic. This follows from the definitions and
\hyperref[prop:sga3-VIB-3.3]{3.3}.

\medskip
\phantomsection\label{par:sga3-VIB-6.9}%
\noindent\textbf{6.9.---}%
\footnote{\label{note:sga3-VIB-67}Editor's note: this subsection
has been added.}
In [RG71], I, \SGARefLink{cor:sga3-VIA-3.3.3}{3.3.3}, the authors introduce the geometric notion
of a \emph{pure \(S\)-scheme}, which is local on \(S\) for the
étale topology; we refer to \emph{loc.\ cit.} for the precise
definition. We merely record the following points:
\begin{enumerate}
\item[\textup{a)}]\phantomsection\label{point:sga3-VIB-6.9-a}
  (\emph{loc.\ cit.}, Theorem 3.3.5) If \(B\) is a flat,
  finitely presented \(A\)-algebra, then \(B\) is a projective
  \(A\)-module if and only if
  \(\operatorname{Spec}(B)\) is pure over
  \(\operatorname{Spec}(A)\).
\item[\textup{b)}]\phantomsection\label{point:sga3-VIB-6.9-b}
  Consequently, if \(X\to S\) is locally of finite
  presentation, flat, and pure, then \(X\) is essentially free
  over \(S\).
\item[\textup{c)}]\phantomsection\label{point:sga3-VIB-6.9-c}
  (\emph{loc.\ cit.}, 3.3.4\textup{(iii)}) If \(X\to S\) is
  locally of finite type, flat, has geometrically irreducible
  fibers, and has no embedded components, then \(X\) is pure over
  \(S\).
\end{enumerate}
Every group scheme locally of finite type over a field is
Cohen--Macaulay
(cf.\ \hyperref[prop:sga3-VIA-1.1.1]{VI\(_A\), 1.1.1}), and
therefore has no embedded components. We obtain in particular:
\begin{enumerate}
\item[\textup{d)}]\phantomsection\label{point:sga3-VIB-6.9-d}
  Every \(S\)-group scheme \(G\) that is locally of finite
  presentation, flat, and has connected fibers is pure over
  \(S\), and hence essentially free over \(S\).
\end{enumerate}
All the statements of
\hyperref[ex:sga3-VIB-6.2.4-e]{6.2.4\textup{(e)}} may therefore
be used with the results of \textup{(b)} and \textup{(d)} above
taken into account.

% END INCLUDED SGA3_English_current_buildable_source_20260823/source/inputs/VI/tex_reference_v2/components/61_expose_VIB_def66_through_par69_en.tex

% BEGIN INCLUDED SGA3_English_current_buildable_source_20260823/source/inputs/VI/tex_reference_v2/components/62_expose_VIB_section7_prop71_en.tex
% SGA 3, Expose VI_B, Loop 1 English reconstruction.
% Scope: Section 7 heading and Proposition 7.1 in full, including
% its complete proof and source editor's notes 68--69.
% Authority: Polo--Gille Exp6B-13oct24.pdf, component-local pages
% 50--51 (combined-reader pages 391--392; re-edition headers
% 380--381; original-margin printed page 363 at the seam).
% Corollary 7.1.1 begins on local page 51 and is excluded.
% OCR and Reinhold English Markdown are comparison witnesses only.
% This unit contains no diagrams.

\section*{7. Generated subgroups; commutator group}
\phantomsection\label{sec:sga3-VIB-7}%
\addcontentsline{toc}{section}{7. Generated subgroups; commutator group}

Throughout this section, \(k\) denotes a fixed field.

\medskip
\phantomsection\label{prop:sga3-VIB-7.1}%
\noindent\textbf{Proposition 7.1.---}
\emph{Let \(G\) be a \(k\)-group, let
\((X_i)_{i\in I}\) be a family of geometrically reduced
\(k\)-schemes,}%
\footnote{\label{note:sga3-VIB-68}Editor's note: here and below,
the little-used term ``separable'' has been replaced by the
customary term ``geometrically reduced''; see EGA~IV\(_2\),
4.6.2.}
\emph{and, for every \(i\in I\), let \(f_i:X_i\to G\) be a
\(k\)-morphism.}
\begin{enumerate}
\item[\textup{\SGAInlineRefTarget{prop:sga3-VIB-7.1:item:i}(i)}]\phantomsection\label{concl:sga3-VIB-7.1-i}
  \emph{There exists a smallest closed \(k\)-group subscheme of
  \(G\) through which every \(f_i\) factors. It is denoted
  \[
    \Gamma_G\bigl((f_i)_{i\in I}\bigr).
  \]
  It is a geometrically reduced \(k\)-scheme, and hence smooth
  when \(G\) is locally of finite type over \(k\)
  (\hyperref[cor:sga3-VIB-1.3.1]{1.3.1}).}
\item[\textup{\SGAInlineRefTarget{prop:sga3-VIB-7.1:item:ii}(ii)}]\phantomsection\label{concl:sga3-VIB-7.1-ii}
  \emph{Put
  \[
    X=\coprod_{i\in I}X_i,
  \]
  and let \(f:X\to G\) be the morphism whose restriction to
  \(X_i\) is \(f_i\), for every \(i\in I\). Put
  \[
    X^1=X\amalg X,
  \]
  and let \(f^1:X^1\to G\) be the morphism whose restrictions to
  the two copies of \(X\) are \(f\) and \(c\circ f\),
  respectively. For every \(n>1\), put
  \[
    X^n=X^1\times_kX^{n-1},
    \qquad
    f^n=\mu\circ(f^1\times_k f^{n-1}):X^n\longrightarrow G.
  \]
  Then \(\Gamma_G((f_i)_{i\in I})\) is the reduced subscheme of
  \(G\) whose underlying space is the closure of
  \[
    \bigcup_{n\geq1}f^n(X^n).
  \]}
\item[\textup{\SGAInlineRefTarget{prop:sga3-VIB-7.1:item:iii}(iii)}]\phantomsection\label{concl:sga3-VIB-7.1-iii}
  \emph{For every \(k\)-scheme \(S\),
  \(\Gamma_G((f_i)_{i\in I})_S\) is the smallest closed group
  subscheme of \(G_S\) through which every
  \[
    f_{i,S}:(X_i)_S\longrightarrow G_S
  \]
  factors.}
\item[\textup{(iii\(^{\prime}\))}]
  \phantomsection\label{concl:sga3-VIB-7.1-iiiprime}
  \emph{Moreover, \(\Gamma_G((f_i)_{i\in I})_S\) is the smallest
  group subscheme of \(G_S\), not necessarily closed, through
  which every \(f_{i,S}\) factors.}%
  \footnote{\label{note:sga3-VIB-69}Editor's note: the original
  stated \textup{(iii\(^{\prime}\))} under the additional
  hypothesis that \(G\) be locally of finite type over \(k\);
  this hypothesis may be omitted in view of
  \hyperref[sga3:VIA:0.5.2]{VI\(_A\), 0.5.2}.}
\end{enumerate}

First note that, to prove \textup{\SGARefLink{prop:sga3-VIB-7.1:item:i}{(i)}}, \textup{\SGARefLink{prop:sga3-VIB-7.1:item:iii}{(iii)}}, and
\textup{(iii\(^{\prime}\))}, after defining \(X\) and \(f\) as
in \textup{\SGARefLink{concl:sga3-VIB-7.1-ii}{(ii)}}, we are reduced to the case in which \(I\)
consists of a single element.

Let \(H\) be the reduced subscheme of \(G\) whose underlying set
is
\[
  \overline{\bigcup_{n\geq1}f^n(X^n)}.
\]
The family of morphisms
\[
  f^n:X^n\longrightarrow H
\]
is schematically dominant (cf.\ EGA~IV\(_3\), 11.10.4), and therefore
every closed subscheme of \(G\) through which all the \(f^n\)
factor also contains \(H\). Moreover, by \emph{loc.\ cit.},
11.10.7, \(H\) is geometrically reduced. To prove
\textup{\SGARefLink{prop:sga3-VIB-7.1:item:i}{(i)}} and \textup{\SGARefLink{prop:sga3-VIB-7.1:item:ii}{(ii)}}, it is therefore enough to show
that \(H\) is a group subscheme of \(G\), that is, that the
restriction of \(c\) to \(H\) and the restriction of \(\mu\) to
\(H\times_kH\) factor through the inclusion \(H\to G\).

Since \(H\) is geometrically reduced, \(H\times_kH\) is reduced,
and it is enough to verify, set-theoretically, that
\[
  c(H)\subset H,
  \qquad
  \mu(H\times_kH)\subset H.
\]
By EGA~IV\(_3\), 11.10.6, the union of the
\[
  f^n_{(H)}(X^n\times_kH)
\]
is schematically dense in \(H\times_kH\). Similarly, for every
\(n\geq1\), the union of the
\[
  f^m_{(X^n)}(X^n\times_kX^m),
  \qquad m\geq1,
\]
is schematically dense in \(X^n\times_kH\). It is therefore
enough to prove
\[
  \mu\!\left(
    f^n_{(H)}
      \bigl(f^m_{(X^n)}(X^n\times_kX^m)\bigr)
  \right)\subset H
\]
and \(c(f^n(X^n))\subset H\). But
\[
\begin{aligned}
  \mu\!\left(
    f^n_{(H)}
      \bigl(f^m_{(X^n)}(X^n\times_kX^m)\bigr)
  \right)
  &=
  \mu\!\left((f^n\times_k f^m)(X^n\times_kX^m)\right)\\
  &=f^{n+m}(X^{n+m})\subset H.
\end{aligned}
\]
Also, since \(c(f^1(X^1))\subset f^1(X^1)\), for every \(n\)
one has
\[
  c(f^n(X^n))\subset f^n(X^n)\subset H.
\]
This proves \textup{\SGARefLink{prop:sga3-VIB-7.1:item:i}{(i)}} and \textup{\SGARefLink{prop:sga3-VIB-7.1:item:ii}{(ii)}}.

We now prove \textup{\SGARefLink{concl:sga3-VIB-7.1-iii}{(iii)}}. Let \(G'\) be a closed group
subscheme of \(G_S\) through which \(f_S\) factors. We must show
that \(G'\) contains \(H_S\), or equivalently that
\[
  H_S=H_S\times_{G_S}G'.
\]
Put
\[
  H'_S=H_S\times_{G_S}G'=G'\cap H_S.
\]
Every \(f^n_S\) factors through both \(G'\) and \(H_S\), and
hence through \(H'_S\). Since the family
\[
  f^n:X^n\longrightarrow H
\]
is schematically dominant, EGA~IV\(_3\), 11.10.6 implies that
the base-changed family
\[
  f^n_S:X^n_S\longrightarrow H_S
\]
is schematically dominant as well. Because every \(f^n_S\)
factors through \(H'_S\), EGA~IV\(_3\), 11.10.1\textup{(c)}
gives
\[
  H'_S=H_S.
\]
This proves \textup{\SGARefLink{prop:sga3-VIB-7.1:item:iii}{(iii)}}.

Finally, we show that \(H_S\) is the smallest group subscheme,
not necessarily closed, of \(G_S\) through which \(f_S\)
factors. Let \(G'\) be such a group subscheme and again put
\[
  H'_S=H_S\times_{G_S}G'.
\]
We must show \(H'_S=H_S\). By \textup{\SGARefLink{concl:sga3-VIB-7.1-iii}{(iii)}}, it is enough to
show that \(H'_S\) is closed in \(H_S\). It is enough in turn to
show that \(H_S\) and \(H'_S\) have the same underlying set, and
for this it suffices to show, for every \(s\in S\), that
\[
  H'_s
  :=H'_S\times_S\kappa(s)
  =H_s\times_{G_s}G'_s
  =H_s.
\]
By \hyperref[sga3:VIA:0.5.2]{VI\(_A\), 0.5.2}, the
\(\kappa(s)\)-group subscheme \(G'_s\) is closed in \(G_s\).
Hence \(H'_s\) is closed in \(H_s\), and the preceding argument,
applied to \(H'_s\), \(H_s\), and the \(f^n_s\), shows that
\[
  H'_s=H_s.
\]
This proves \textup{(iii\(^{\prime}\))} and completes the proof.

% END INCLUDED SGA3_English_current_buildable_source_20260823/source/inputs/VI/tex_reference_v2/components/62_expose_VIB_section7_prop71_en.tex

% BEGIN INCLUDED SGA3_English_current_buildable_source_20260823/source/inputs/VI/tex_reference_v2/components/63_expose_VIB_cor711_through_cor721_en.tex
% SGA 3, Expose VI_B, Loop 1 English reconstruction.
% Scope: Corollary 7.1.1; Definitions and Remarks 7.2(i)--(v);
% and Corollary 7.2.1 with their proofs, including source
% editor's notes 70--72.
% Authority: Polo--Gille Exp6B-13oct24.pdf, component-local pages
% 51--52 (combined-reader pages 392--393; re-edition headers
% 381--382; original-margin printed pages 363--365).
% Definition 7.2.2 begins on local page 52 and is excluded.
% OCR and Reinhold English Markdown are comparison witnesses only.
% This unit contains no diagrams.

\medskip
\phantomsection\label{cor:sga3-VIB-7.1.1}%
\noindent\textbf{Corollary 7.1.1.---}%
\footnote{\label{note:sga3-VIB-70}Editor's note: this corollary
has been added to point out this special case of \SGARefLink{prop:sga3-VIB-7.1}{7.1}.}
\emph{Let \(G\) be a \(k\)-group locally of finite type, let
\(A\) and \(B\) be two \(k\)-subgroups of \(G\) that are smooth
and of finite type, and let \(i_A\) (respectively, \(i_B\)) be the
inclusion of \(A\) (respectively, \(B\)) into \(G\). Suppose that
\(B\) normalizes \(A\). Then}
\phantomsection\label{form:sga3-VIB-7.1.1-generated-product}%
\[
  A\cdot B=\Gamma_G(i_A,i_B).
\]

\emph{Proof.}
Let
\[
  H=A\rtimes B
\]
be the semidirect product of \(A\) and \(B\) (cf.\ I, \SGARefLink{sga3:i:definition:2-3-5}{2.3.5});
it is a smooth \(k\)-group of finite type. The group morphism
\[
  u:H\longrightarrow G,
  \qquad
  (a,b)\longmapsto ab,
\]
is therefore quasi-compact. Hence, by
\hyperref[prop:sga3-VIB-1.2]{1.2},
\[
  u(H)=A\cdot B
\]
is a closed reduced subscheme of \(G\), and is a group in the
category \((\mathrm{Sch}/k)_{\mathrm{red}}\). By
EGA~IV\(_3\), 11.10.7, \(A\cdot B=u(H)\) is geometrically
reduced, since \(H\) is; it is therefore a closed subgroup of
\(G\). The evident inclusion
\[
  A\cdot B\subset\Gamma_G(i_A,i_B)
\]
now proves the corollary.

\par
\ifdim\dimexpr\pagegoal-\pagetotal\relax<12\baselineskip
  \newpage
\fi
\medskip
\phantomsection\label{defrem:sga3-VIB-7.2}%
\noindent\textbf{Definitions and \SGARefLink{defrem:sga3-VIB-7.2}{Remarks 7.2}.---}%
\footnote{\label{note:sga3-VIB-71}Editor's note: item
\textup{(viii)} of the original has been placed in
\textup{(i\(^{\prime}\))}, and items \textup{(vi)} and
\textup{(vii)} have been singled out as \hyperref[cor:sga3-VIB-7.2.1]{Corollary 7.2.1} and
\hyperref[def:sga3-VIB-7.2.2]{Definition 7.2.2} below.}
\begin{enumerate}
\item[\textup{(i)}]\phantomsection\label{def:sga3-VIB-7.2-i}
  Let \(G\) be a \(k\)-group, let \((X_i)_{i\in I}\) be a
  family of geometrically reduced \(k\)-schemes, and, for every
  \(i\in I\), let \(f_i:X_i\to G\) be a \(k\)-morphism. The
  \emph{closed group subscheme of \(G\) generated by the family
  \((f_i)_{i\in I}\)} is the smallest closed group subscheme of
  \(G\) through which every \(f_i\) factors; in this subsection
  it is denoted
  \[
    \Gamma_G\bigl((f_i)_{i\in I}\bigr).
  \]
  If \(X\) is a subscheme of \(G\), geometrically reduced over
  \(k\), and \(f:X\hookrightarrow G\) is the immersion, we write
  \(\Gamma_G(X)\) instead of \(\Gamma_G(f)\).

\item[\textup{(i\(^{\prime}\))}]
  \phantomsection\label{def:sga3-VIB-7.2-iprime}
  With the notation of
  \hyperref[concl:sga3-VIB-7.1-ii]{7.1\textup{(ii)}}, we shall
  sometimes put
  \phantomsection\label{form:sga3-VIB-7.2-iprime-Gamma}%
  \[
    \Gamma'_G(f)=\bigcup_{n\geq1}f^n(X^n).
  \]
  Notice that \(\Gamma'_G(f)\) is a subset of \(G\) stable
  under the group law, in the sense of \hyperref[par:sga3-VIB-3.0]{3.0}.

\item[\textup{(ii)}]\phantomsection\label{rem:sga3-VIB-7.2-ii}
  If \(X_1\) and \(X_2\) are geometrically reduced
  \(k\)-schemes and \(f_1:X_1\to G\) and \(f_2:X_2\to G\) are
  \(k\)-morphisms such that
  \[
    f_1(X_1)=f_2(X_2)
  \]
  as sets, then it is clear that
  \[
    \Gamma_G(f_1)=\Gamma_G(f_2).
  \]

\item[\textup{(iii)}]\phantomsection\label{def:sga3-VIB-7.2-iii}
  Let \(E\) be a subset of \(G\) such that the reduced
  subscheme \(\overline E\) of \(G\) is geometrically reduced.
  The \emph{closed group subscheme of \(G\) generated by \(E\)}
  is the group subscheme
  \[
    \Gamma_G(E):=\Gamma_G(i),
  \]
  where \(i:\overline E\hookrightarrow G\) is the inclusion of
  that reduced subscheme.

\item[\textup{(iv)}]\phantomsection\label{rem:sga3-VIB-7.2-iv}
  Since every group subscheme of \(G\) is closed, by
  \hyperref[sga3:VIA:0.5.2]{VI\(_A\), 0.5.2},%
  \footnote{\label{note:sga3-VIB-72}Editor's note: the
  hypothesis that \(G\) be locally of finite type over \(k\)
  has been removed here.}
  we shall speak of the ``generated group subscheme'' rather
  than the ``generated closed group subscheme.''

\item[\textup{(v)}]\phantomsection\label{rem:sga3-VIB-7.2-v}
  Let \(X\) be a geometrically reduced \(k\)-scheme and let
  \(f:X\to G\) be a \(k\)-morphism. Suppose that \(f(X)\)
  contains the unit element \(e\) of \(G\). Put
  \[
    X'^{\,1}=X\times_kX,
    \qquad
    f'^{\,1}=\mu\circ\bigl(f\times_k(c\circ f)\bigr),
  \]
  and, for \(n>1\), put
  \phantomsection\label{form:sga3-VIB-7.2-v-recursion}%
  \[
    X'^{\,n}=X'^{\,1}\times_kX'^{\,n-1},
    \qquad
    f'^{\,n}
      =\mu\circ\bigl(f'^{\,1}\times_k f'^{\,n-1}\bigr).
  \]
  Then \(\Gamma_G(f)\) is the reduced subscheme of \(G\) whose
  underlying space is the closure of
  \[
    \bigcup_{n\geq1}f'^{\,n}(X'^{\,n}).
  \]

  Indeed, recall the notation of
  \hyperref[concl:sga3-VIB-7.1-ii]{7.1\textup{(ii)}}:
  \[
    X^1=X\amalg X,
    \qquad
    X^n=X^1\times_kX^{n-1}\quad(n\geq2).
  \]
  The hypothesis \(e\in f(X)\) gives, for every \(n\geq1\),
  \phantomsection\label{form:sga3-VIB-7.2-v-inclusions}%
  \[
    f^n(X^n)
      \subset f'^{\,n}(X'^{\,n})
      \subset f^{2n}(X^{2n}).
  \]
  Consequently, there exists an integer \(n\) such that
  \(f^n(X^n)=\Gamma_G(f)\) if and only if there exists an
  integer \(m\) such that
  \(f'^{\,m}(X'^{\,m})=\Gamma_G(f)\).
\end{enumerate}

From \hyperref[rem:sga3-VIB-7.2-v]{Remark 7.2\textup{(v)}} we
deduce the following result.

\medskip
\phantomsection\label{cor:sga3-VIB-7.2.1}%
\noindent\textbf{Corollary 7.2.1.---}
\emph{Let \(X\) be a geometrically reduced and geometrically
connected \(k\)-scheme, and let \(f:X\to G\) be a \(k\)-morphism
such that \(f(X)\) contains the unit element of \(G\). Then the
\(k\)-group \(\Gamma_G(f)\) is connected, and hence
irreducible.}

\emph{Proof.}
Each \(X'^{\,n}\) is connected, so the union of the
\[
  f'^{\,n}(X'^{\,n}),
\]
all of which contain the unit element, is connected. Its closure
\(\Gamma_G(f)\) is therefore connected as well. It follows that
\(\Gamma_G(f)\) is irreducible, by
\hyperref[prop:sga3-VIA-2.4]{VI\(_A\), 2.4} when \(G\), and
hence also \(\Gamma_G(f)\), is locally of finite type over
\(k\), and by
\hyperref[thm:sga3-VIA-2.6.5]{VI\(_A\), 2.6.5\textup{(iii)}}
in the general case.

% END INCLUDED SGA3_English_current_buildable_source_20260823/source/inputs/VI/tex_reference_v2/components/63_expose_VIB_cor711_through_cor721_en.tex

% BEGIN INCLUDED SGA3_English_current_buildable_source_20260823/source/inputs/VI/tex_reference_v2/components/64_expose_VIB_def722_through_cor73_en.tex
% SGA 3, Expose VI_B, Loop 1 English reconstruction.
% Scope: Definition 7.2.2, Recall 7.3.0, and Corollary 7.3
% with its complete proof, including source editor's notes 73--74.
% Authority: Polo--Gille Exp6B-13oct24.pdf, component-local pages
% 52--54 (combined-reader pages 393--395; re-edition headers
% 382--384; original-margin printed pages 365--367).
% Proposition 7.4 begins on local page 54 and is excluded.
% OCR and Reinhold English Markdown are comparison witnesses only.
% This unit contains no diagrams.

\medskip
\phantomsection\label{def:sga3-VIB-7.2.2}%
\noindent\textbf{Definition 7.2.2.---}
Let \(G\) be a \(k\)-group.
\begin{enumerate}
\item[\textup{a)}]\phantomsection\label{def:sga3-VIB-7.2.2-a}
  Let \(A\) and \(B\) be two geometrically reduced
  \(k\)-group subschemes of \(G\). The \emph{commutator group
  subscheme of \(A\) and \(B\) in \(G\)}, denoted
  \[
    (A,B)_G
    \qquad\text{or simply}\qquad
    (A,B),
  \]
  is the closed group subscheme of \(G\) generated by the
  morphism
  \phantomsection\label{form:sga3-VIB-7.2.2-commutator}%
  \[
    \nu:A\times_kB\longrightarrow G
  \]
  defined, for every \(k\)-scheme \(S\), by
  \[
    (a,b)\longmapsto aba^{-1}b^{-1},
    \qquad
    a\in A(S),\quad b\in B(S).
  \]

\item[\textup{b)}]\phantomsection\label{def:sga3-VIB-7.2.2-b}
  Suppose that \(G\) is geometrically reduced over \(k\). The
  \emph{derived group of \(G\)}, denoted \(\mathfrak{D}(G)\),%
  \footnote{\label{note:sga3-VIB-73}Editor's note: the notation
  \(\operatorname{Der}(G)\) of the original has been changed to
  avoid any risk of confusion with a space of derivations.}
  is the group
  \[
    \mathfrak{D}(G)=(G,G).
  \]
\end{enumerate}

\noindent\emph{N.B.}
The group \(G\) is commutative if and only if \(\mathfrak{D}(G)\) is the
unit \(k\)-group.

\medskip
\phantomsection\label{recall:sga3-VIB-7.3.0}%
\noindent\textbf{Reminders \SGARefLink{recall:sga3-VIB-7.3.0}{7.3.0}.---}%
\footnote{\label{note:sga3-VIB-74}Editor's note: these reminders
have been added, and the statement of \SGARefLink{cor:sga3-VIB-7.3}{7.3} has consequently been
modified and expanded.}
Let \(\varphi:Y\to Z\) be a morphism of \(S\)-schemes. Recall
that the image presheaf \(\varphi(Y)\) is the \(S\)-functor
which assigns to every \(S'\) over \(S\) the subset
\[
  \varphi\bigl(Y(S')\bigr)\subset Z(S').
\]
Suppose that \(T\) is a subscheme of \(Z\). If the inclusion of
presheaves
\[
  \varphi(Y)\hookrightarrow Z
\]
factors through \(T\), that is, if
\[
  \varphi\circ h\in T(S')
\]
for every \(S\)-morphism \(h:S'\to Y\), then, set-theoretically,
\(\varphi(Y)\subset T\): take \(h=\operatorname{id}_Y\).
Conversely, if \(Y\) is reduced and
\(\varphi(Y)\subset T\) set-theoretically, then \(\varphi\)
factors through \(T\).

Recall also that, by \hyperref[prop:sga3-VIB-6.2]{6.2}, if
\(H\) is a closed \(k\)-group subscheme of \(G\), then
\[
  \underline{\operatorname{Centr}}_G H
  \qquad\text{and}\qquad
  \underline{\operatorname{Norm}}_G H
\]
are representable by closed \(k\)-group subschemes of \(G\).

\medskip
\phantomsection\label{cor:sga3-VIB-7.3}%
\noindent\textbf{Corollary 7.3.---}
\emph{Let \(G\) be a \(k\)-group, let \(X\) be a geometrically
reduced \(k\)-scheme, and let \(f:X\to G\) be a
\(k\)-morphism.}
\begin{enumerate}
\item[\textup{\SGAInlineRefTarget{cor:sga3-VIB-7.3:item:i}(i)}]\phantomsection\label{concl:sga3-VIB-7.3-i}
  \emph{Let \(S\) be a \(k\)-scheme and let \(u\) be an
  endomorphism of the \(S\)-group \(G_S\).}
  \begin{enumerate}
  \item[\textup{(a)}]\phantomsection\label{concl:sga3-VIB-7.3-i-a}
    \emph{If}
    \[
      u\bigl(f_S(X_S)\bigr)\subset\Gamma_G(f)_S
    \]
    \emph{set-theoretically, then the morphism}
    \[
      u:\Gamma_G(f)_S\longrightarrow G_S
    \]
    \emph{factors through \(\Gamma_G(f)_S\).}

  \item[\textup{\SGAInlineRefTarget{cor:sga3-VIB-7.3:item:b}(b)}]\phantomsection\label{concl:sga3-VIB-7.3-i-b}
    \emph{If \(u\) is an automorphism of the \(S\)-group \(G_S\)
    and}
    \[
      u\bigl(f_S(X_S)\bigr)=f_S(X_S)
    \]
    \emph{set-theoretically, then \(u\) induces an automorphism
    of \(\Gamma_G(f)_S\). In particular, if \(g\in G(S)\)
    satisfies}
    \[
      \operatorname{int}(g)\bigl(f_S(X_S)\bigr)=f_S(X_S)
    \]
    \emph{set-theoretically, then}
    \[
      g\in
      \underline{\operatorname{Norm}}_{G_S}
        \bigl(\Gamma_G(f)_S\bigr)(S).
    \]

  \item[\textup{\SGAInlineRefTarget{cor:sga3-VIB-7.3:item:c}(c)}]\phantomsection\label{concl:sga3-VIB-7.3-i-c}
    \emph{If \(u\circ f_S=f_S\), then the restriction of \(u\)
    to the subgroup \(\Gamma_G(f)_S\) of \(G_S\) is the identity.
    In particular, if \(g\in G(S)\) satisfies}
    \[
      \operatorname{int}(g)\circ f_S=f_S,
    \]
    \emph{then}
    \[
      g\in
      \underline{\operatorname{Centr}}_{G_S}
        \bigl(\Gamma_G(f)_S\bigr)(S).
    \]
  \end{enumerate}

\item[\textup{\SGAInlineRefTarget{cor:sga3-VIB-7.3:item:ii}(ii)}]\phantomsection\label{concl:sga3-VIB-7.3-ii}
  \emph{Let \(H\) be a group subscheme of \(G\). Then \(H\)
  centralizes (respectively, normalizes) \(\Gamma_G(f)\) if and
  only if, for every \(k\)-scheme \(S\) and every \(h\in H(S)\),}
  \[
    \operatorname{int}(h)\circ f_S=f_S
  \]
  \emph{(respectively,}
  \[
    \operatorname{int}(h)\bigl(f_S(X_S)\bigr)
      \subset\Gamma_G(f)_S
  \]
  \emph{set-theoretically). Equivalently, for every
  \(S'\to S\) and every \(x\in X(S')\), the elements \(h_{S'}\)
  and \(f(x)\) of \(G(S')\) commute (respectively,}
  \[
    h_{S'}f(x)h_{S'}^{-1}\in\Gamma_G(f)(S').
  \]
  \emph{).}

\item[\textup{\SGAInlineRefTarget{cor:sga3-VIB-7.3:item:iii}(iii)}]\phantomsection\label{concl:sga3-VIB-7.3-iii}
  \emph{In particular, let \(Y\) be a second geometrically
  reduced \(k\)-scheme and let \(\varphi:Y\to G\) be a
  \(k\)-morphism. Suppose that, for every \(k\)-scheme \(S'\),
  every \(x\in X(S')\), and every \(y\in Y(S')\), the elements
  \(\varphi(y)\) and \(f(x)\) of \(G(S')\) commute
  (respectively,}
  \[
    \operatorname{int}\bigl(\varphi(y)\bigr)\bigl(f(x)\bigr)
      \in\Gamma_G(f)(S').
  \]
  \emph{).}
  \emph{Then \(\Gamma_G(\varphi)\) is a \(k\)-subgroup of}
  \[
    \underline{\operatorname{Centr}}_G\bigl(\Gamma_G(f)\bigr)
    \qquad\text{\emph{(respectively, of)}}\qquad
    \underline{\operatorname{Norm}}_G\bigl(\Gamma_G(f)\bigr).
  \]

\item[\textup{\SGAInlineRefTarget{cor:sga3-VIB-7.3:item:iv}(iv)}]\phantomsection\label{concl:sga3-VIB-7.3-iv}
  \emph{If \(g\) is a \(k\)-point of \(G\), then the \(k\)-group
  \(\Gamma_G(\{g\})\) is commutative.}

\item[\textup{(v)}]\phantomsection\label{concl:sga3-VIB-7.3-v}
  \emph{Let \(A\) and \(B\) be two group subschemes of \(G\),
  geometrically reduced over \(k\). If \(A\) and \(B\) are
  invariant (respectively, characteristic), then so is
  \((A,B)\).}
\end{enumerate}

\emph{Proof.}
We first prove \textup{\SGARefLink{cor:sga3-VIB-7.3:item:i}{(i)}\SGARefLink{concl:sga3-VIB-7.3-i}{(a)}}. Put
\[
  \Gamma_S=\Gamma_G(f)_S.
\]
Then \(u^{-1}(\Gamma_S)\) is a closed \(S\)-group subscheme of
\(G_S\). By
\hyperref[concl:sga3-VIB-7.1-iii]{7.1\textup{(iii)}}, it is
therefore enough to show that \(f_S\) factors through
\(u^{-1}(\Gamma_S)\). Since
\[
  u\bigl(f_S(X_S)\bigr)\subset\Gamma_S
\]
and \(X_S\) is reduced, \(u\circ f_S\) factors through
\(\Gamma_S\); hence \(f_S\) factors through
\(u^{-1}(\Gamma_S)\). This proves \textup{\SGARefLink{concl:sga3-VIB-7.3-i-a}{(a)}}.

The first assertion of \textup{\SGARefLink{cor:sga3-VIB-7.3:item:b}{(b)}} follows by applying
\textup{\SGARefLink{concl:sga3-VIB-7.3-i-a}{(a)}} to \(u\) and \(u^{-1}\). In fact, it is enough to
assume
\[
  u\bigl(f_S(X_S)\bigr)\subset\Gamma_S
  \qquad\text{and}\qquad
  u^{-1}\bigl(f_S(X_S)\bigr)\subset\Gamma_S.
\]
The second assertion is the special case \(u=\operatorname{int}(g)\).

We now prove \textup{\SGARefLink{concl:sga3-VIB-7.3-i-c}{(c)}}. Consider the morphism of \(S\)-groups
\[
  G_S\longrightarrow G_S\times_SG_S,
  \qquad
  g\longmapsto\bigl(u(g),g\bigr),
\]
and let \(H\) be the inverse image of the diagonal. It is a group
subscheme of \(G_S\). Since \(G\) is separated over \(k\), by
\hyperref[sga3:VIA:0.3]{VI\(_A\), 0.3}, \(G_S\) is separated
over \(S\), and hence \(H\) is closed in \(G_S\). By hypothesis,
\(f_S\) factors through \(H\); therefore \(H\) contains
\(\Gamma_G(f)_S\), by
\hyperref[concl:sga3-VIB-7.1-iii]{7.1\textup{(iii)}}. This proves
the first assertion of \textup{\SGARefLink{cor:sga3-VIB-7.3:item:c}{(c)}}, and the second is the
special case \(u=\operatorname{int}(g)\).

We next prove \textup{\SGARefLink{cor:sga3-VIB-7.3:item:ii}{(ii)}}. Put
\[
  \Gamma=\Gamma_G(f)
\]
and let \(i:\Gamma\hookrightarrow G\) denote the inclusion. The
group \(H\) is contained in
\[
  N=\underline{\operatorname{Norm}}_G(\Gamma)
\]
if and only if, for every \(k\)-scheme \(S\) and every
\(h\in H(S)\),
\[
  \operatorname{int}(h)(\Gamma_S)\subset\Gamma_S.
\]
Applying this condition to \(h\) and \(h^{-1}\) shows that the
inclusion is then an equality. By \textup{\SGARefLink{cor:sga3-VIB-7.3:item:i}{(i)}\SGARefLink{concl:sga3-VIB-7.3-i}{(a)}}, this condition
holds if and only if
\[
  \operatorname{int}(h)\bigl(f_S(X_S)\bigr)\subset\Gamma_S.
\]
Similarly, \(H\) is contained in
\[
  C=\underline{\operatorname{Centr}}_G(\Gamma)
\]
if and only if, for every \(k\)-scheme \(S\) and every
\(h\in H(S)\),
\[
  \operatorname{int}(h)\circ i_S=i_S.
\]
By \textup{\SGARefLink{cor:sga3-VIB-7.3:item:i}{(i)}\SGARefLink{cor:sga3-VIB-7.3:item:c}{(c)}}, this holds if and only if
\[
  \operatorname{int}(h)\circ f_S=f_S.
\]
This proves \textup{\SGARefLink{cor:sga3-VIB-7.3:item:ii}{(ii)}}.

We now prove \textup{\SGARefLink{cor:sga3-VIB-7.3:item:iii}{(iii)}}. In view of \textup{\SGARefLink{cor:sga3-VIB-7.3:item:ii}{(ii)}}, the
hypothesis implies that \(\varphi\) factors through \(C\)
(respectively, \(N\)). The latter is a closed subgroup of \(G\),
by \hyperref[prop:sga3-VIB-6.2]{6.2}, and therefore contains
\(\Gamma_G(\varphi)\), by
\hyperref[concl:sga3-VIB-7.1-i]{7.1\textup{(i)}}.%
\begingroup
\renewcommand{\thefootnote}{\fnsymbol{footnote}}%
\footnote[1]{\label{note:sga3-VIB-7.3-Gamma-phi}Translator's note:
the French source prints \(\Gamma_G(f)\) here. Since
\(\varphi\), rather than \(f\), is the morphism shown to factor
through \(C\) (respectively, \(N\)), \SGARefLink{prop:sga3-VIB-7.1}{7.1}\textup{\SGARefLink{concl:sga3-VIB-7.1-i}{(i)}} gives
\(\Gamma_G(\varphi)\), as required by the statement.}
\endgroup

Assertion \textup{\SGARefLink{cor:sga3-VIB-7.3:item:iv}{(iv)}} follows from \textup{\SGARefLink{cor:sga3-VIB-7.3:item:iii}{(iii)}} by taking
both \(f\) and \(\varphi\) to be the closed immersion
\[
  \operatorname{Spec}k\hookrightarrow G
\]
defined by \(g\). Indeed, for every \(k\)-scheme \(S\), the sets
\(X(S)\) and \(Y(S)\) then each have a single element, which is
sent by \(f\), respectively by \(\varphi\), to \(g\).

Finally, we prove \textup{\SGARefLink{concl:sga3-VIB-7.3-iv}{(v)}}. Let
\[
  \nu:G\times_kG\longrightarrow G,
  \qquad
  (g,h)\longmapsto ghg^{-1}h^{-1},
\]
and let \(\nu'\) be its restriction to \(A\times_kB\). By
\hyperref[def:sga3-VIB-7.2.2]{7.2.2},
\[
  (A,B)=\Gamma_G(\nu').
\]
Let \(S\) be a \(k\)-scheme, and let \(u\) be an inner
automorphism, respectively an arbitrary \(S\)-group
automorphism, of \(G_S\). Then
\[
  u\circ\nu_S=\nu_S\circ(u\times_Su).
\]
The hypothesis implies that \(u\) induces automorphisms \(u_1\)
of \(A_S\) and \(u_2\) of \(B_S\). Consequently,
\phantomsection\label{form:sga3-VIB-7.3-v-commutator-image}%
\[
  u\bigl(\nu'_S(A_S\times_SB_S)\bigr)
    =\nu'_S\bigl(u_1(A_S)\times_Su_2(B_S)\bigr)
    =\nu'_S(A_S\times_SB_S).
\]
By \textup{\SGARefLink{concl:sga3-VIB-7.3-i}{(i)}\SGARefLink{concl:sga3-VIB-7.3-i-b}{(b)}}, \(u\) therefore induces an automorphism of
\((A,B)_S\). This proves \textup{\SGARefLink{concl:sga3-VIB-7.3-v}{(v)}}.

% END INCLUDED SGA3_English_current_buildable_source_20260823/source/inputs/VI/tex_reference_v2/components/64_expose_VIB_def722_through_cor73_en.tex

% BEGIN INCLUDED SGA3_English_current_buildable_source_20260823/source/inputs/VI/tex_reference_v2/components/65_expose_VIB_prop74_en.tex
% SGA 3, Expose VI_B, Loop 1 English reconstruction.
% Scope: Proposition 7.4 and its complete proof.
% Authority: Polo--Gille Exp6B-13oct24.pdf, component-local pages
% 54--55 (combined-reader pages 395--396; re-edition headers
% 384--385; original-margin printed pages 367--369).
% The preceding tail of 7.3 on local page 54 is excluded.
% Lemma 7.5 begins on local page 56 and is excluded.
% OCR and Reinhold English Markdown are comparison witnesses only.
% This unit contains no diagrams.

\medskip
\phantomsection\label{prop:sga3-VIB-7.4}%
\noindent\textbf{Proposition 7.4.---}
\emph{Let \(G\) be a \(k\)-group locally of finite type, let
\(X\) be a finite-type \(k\)-scheme that is geometrically
reduced and geometrically connected, and let \(f:X\to G\) be a
\(k\)-morphism such that \(f(X)\) contains the unit element
\(e\) of \(G\). Then, with the notation of
\hyperref[concl:sga3-VIB-7.1-ii]{7.1\textup{(ii)}}, there
exists an integer \(N\) such that, set-theoretically,}
\phantomsection\label{form:sga3-VIB-7.4-stabilization}%
\[
  f^N(X^N)=\Gamma_G(f).
\]

\emph{Proof.}
By \hyperref[concl:sga3-VIB-7.1-iii]{7.1\textup{(iii)}} and
EGA~IV\(_2\), 2.6.1, we may assume that \(k\) is algebraically
closed. By
\hyperref[cor:sga3-VIB-7.2.1]{Corollary 7.2.1}, we may assume
that \(G=G^0\). Finally, with the notation of
\hyperref[rem:sga3-VIB-7.2-v]{7.2\textup{(v)}}, it is enough to
show that there exists an integer \(N\) such that
\[
  f'^{\,N}(X'^{\,N})=\Gamma_G(f).
\]

\emph{First case.}
Suppose that \(X\) is irreducible. The closed irreducible
subsets
\[
  \overline{f'^{\,n}(X'^{\,n})}
\]
then form an increasing sequence in the space \(G\), which is
Noetherian because \(G=G^0\) is of finite type over \(k\)
(\hyperref[prop:sga3-VIA-2.4]{VI\(_A\), 2.4}). Hence this
sequence is stationary, and there exists an integer \(m\) such
that
\[
  \overline{f'^{\,m}(X'^{\,m})}=\Gamma_G(f).
\]

Moreover, because \(X\) and \(G\) are of finite type over
\(k\), the morphisms \(f'^{\,n}\) are of finite type over
\(k\). Consequently, \(f'^{\,m}(X'^{\,m})\) is constructible
in \(G\) (EGA~IV\(_1\), 1.8.5), and therefore contains an open
subset \(U\) of its closure \(\Gamma_G(f)\)
(EGA~0\(_{\mathrm{IV}}\), 9.2.3). By
\hyperref[sga3:VIA:0.5]{VI\(_A\), 0.5}, one then has
\phantomsection\label{form:sga3-VIB-7.4-first-case-chain}%
\[
  \Gamma_G(f)
  \subset U\cdot U
  \subset f'^{\,2m}(X'^{\,2m})
  \subset \Gamma_G(f),
\]
and hence
\[
  f'^{\,2m}(X'^{\,2m})=\Gamma_G(f).
\]

\emph{Second case.}
Suppose that \(X\) has exactly two irreducible components
\(A_1\) and \(A_2\). Since \(X\) is connected and \(k\) is
algebraically closed, there exists
\[
  a\in(A_1\cap A_2)(k).
\]
The four irreducible subsets
\[
  A_i\times_kA_j
  \qquad (i,j=1,2)
\]
therefore cover \(X\times_kX\), and the image of each under the
morphism
\[
  f'^{\,1}=\mu\circ\bigl(f\times_k(c\circ f)\bigr)
\]
contains \(e\). Let \(f'^{\,1}_{ij}\) denote the restriction of
\(f'^{\,1}\) to the reduced subscheme \(A_i\times_kA_j\), and
put
\[
\begin{aligned}
  Y={}&
  (A_1\times_kA_1)
  \times_k(A_1\times_kA_2)
  \times_k(A_2\times_kA_2)
  \times_k(A_2\times_kA_1),\\
  g={}&
  \mu\circ\Bigl(
    \mu\circ
      (f'^{\,1}_{11}\times_k f'^{\,1}_{12})
    \times_k
    \mu\circ
      (f'^{\,1}_{22}\times_k f'^{\,1}_{21})
  \Bigr).
\end{aligned}
\]
The scheme \(Y\) is irreducible, reduced, and of finite type.
By the first case, there is therefore an integer \(m\) such
that
\[
  g'^{\,m}(Y'^{\,m})=\Gamma_G(g).
\]
For every \(n\geq1\), however,\begingroup
\renewcommand{\thefootnote}{\fnsymbol{footnote}}%
\footnote[1]{\label{note:sga3-VIB-7.4-Yprime}Translator's note:
the French source prints \(g'^{\,n}(Y^{\,n})\) in the middle
term. Since the source of \(g'^{\,n}\) is \(Y'^{\,n}\), the
prime is restored here; compare the preceding equality
\(g'^{\,m}(Y'^{\,m})=\Gamma_G(g)\).}%
\endgroup
\phantomsection\label{form:sga3-VIB-7.4-two-component-chain}%
\[
  f'^{\,n}(X'^{\,n})
  \subset g'^{\,n}(Y'^{\,n})
  \subset f'^{\,4n}(X'^{\,4n}).
\]
It follows that
\[
  \Gamma_G(f)=\Gamma_G(g)
  \qquad\text{and}\qquad
  f'^{\,4m}(X'^{\,4m})=\Gamma_G(f).
\]

\emph{General case.}
We argue by induction on the number of irreducible components
of \(X\); this number is finite because \(X\), being of finite
type over \(k\), is Noetherian. Suppose the proposition has
been proved when \(X\) has at most \(r-1\) irreducible
components, and suppose that \(X\) has \(r\) such components,
say \(A_1,\ldots,A_r\). The unit \(e\) belongs to the image of
one of the \(A_i\); suppose, for example, that
\[
  e\in f(A_1).
\]
Put
\[
  X_r=A_1\cup\cdots\cup A_{r-1},
  \qquad
  Y=\Gamma_G(f_r),
\]
where \(f_r\) denotes the restriction of \(f\) to the reduced
subscheme \(X_r\). We choose the numbering of the \(A_i\) so
that \(X_r\) is connected, which is always possible. Then
\(Y\) is a closed, reduced, irreducible subgroup of \(G\), by
\hyperref[cor:sga3-VIB-7.2.1]{Corollary 7.2.1}.

Put
\[
  Z=\overline{f(A_r)}
\]
and let \(T=Y\cup Z\), endowed with its closed reduced
subscheme structure; let \(i:T\hookrightarrow G\) be the
inclusion. It follows immediately from
\hyperref[rem:sga3-VIB-7.2-ii]{7.2\textup{(ii)}} that
\[
  \Gamma_G(i)=\Gamma_G(f).
\]
Moreover, \(T\) is connected: since \(X\) is connected,
\(X_r\) and \(A_r\) have a common point \(a\), and \(Y\) and
\(Z\) have the common point \(f(a)\). Also, \(e\in T\), and
\(T\) has at most two irreducible components, because \(Y\)
and \(Z\) are irreducible.

By the induction hypothesis, there exists an integer \(m\)
such that
\[
  f_r'^{\,m}(X_r'^{\,m})
  =\Gamma_G(f_r)=Y.
\]
Since \(X=X_r\cup A_r\), and since
\[
  Z=\overline{f(A_r)}
  \subset\overline{f'^{\,m}(X'^{\,m})}
\]
because \(e\in f(A_1)\), we obtain
\phantomsection\label{form:sga3-VIB-7.4-general-case-inclusion}%
\[
  f(X)
  \subset Y\cup Z
  \subset\overline{f'^{\,m}(X'^{\,m})}.
\]

On the other hand, because \(T\) has at most two irreducible
components, the second case gives an integer \(p\) such that
\[
  i'^{\,p}(T'^{\,p})
  =\Gamma_G(i)=\Gamma_G(f).
\]
Now
\[
  T\subset\overline{f'^{\,m}(X'^{\,m})},
\]
and hence
\[
  \overline{f'^{\,mp}(X'^{\,mp})}=\Gamma_G(f).
\]
As in the first case, this last equality implies
\[
  f'^{\,2mp}(X'^{\,2mp})=\Gamma_G(f).
\]
This completes the proof.

% END INCLUDED SGA3_English_current_buildable_source_20260823/source/inputs/VI/tex_reference_v2/components/65_expose_VIB_prop74_en.tex

% BEGIN INCLUDED SGA3_English_current_buildable_source_20260823/source/inputs/VI/tex_reference_v2/components/66_expose_VIB_lemma75_en.tex
% SGA 3, Expose VI_B, Loop 1 English reconstruction.
% Scope: Lemma 7.5 and its complete proof, including source
% editor's note 75.
% Authority: Polo--Gille Exp6B-13oct24.pdf, component-local page
% 56 (combined-reader page 397; re-edition header page 386;
% original-margin printed page 370).
% Remarque 7.6.0 begins later on local page 56 and is excluded.
% OCR and Reinhold English Markdown are comparison witnesses only.
% This unit contains no diagrams.

\medskip
\phantomsection\label{lem:sga3-VIB-7.5}%
\noindent\textbf{Lemma 7.5.---}
\emph{Let \(S\) be a scheme, let \(G\) be an \(S\)-group
scheme, let \(X\) be an \(S\)-prescheme, and let
\(f:X\to G\) be an \(S\)-morphism. Suppose that \(X\) and \(G\)
are locally of finite presentation over \(S\), and that, for
every \(s\in S\) and every maximal point \(g\) of \(G_s\),
there exists a point \(x\) of \(X\) such that \(f(x)=g\) and
\(f\) is flat at \(x\). Then the morphism}
\phantomsection\label{form:sga3-VIB-7.5-covering-morphism}%
\[
  \mu\circ(f\times_S f):
  X\times_S X\longrightarrow G
\]
\emph{is a covering for the \((\mathrm{fppf})\) topology.}

\footnote{\label{note:sga3-VIB-75}Editor's note: the argument
in the original has been expanded in what follows.}
By EGA~IV\(_3\), 11.3.1, the set \(V\) of points of \(X\) at
which \(f\) is flat is open, and \(f|_V\) is an open morphism.
Thus
\[
  U=f(V)
\]
is an open subset of \(G\). Moreover, by hypothesis,
\[
  U\cap G_s
\]
is dense in \(G_s\) for every \(s\in S\). Let \(\phi\) denote
the restriction of \(\mu\circ(f\times_S f)\) to
\(V\times_S V\).

It is enough to show that \(\phi\) is a covering for the
\((\mathrm{fppf})\) topology, and for this it is enough to show
that \(\phi\) is faithfully flat and of finite presentation
(\hyperref[sga3:iv:proposition:6:3:1:item:iv:01]
{Expos\'e~IV, 6.3.1\textup{(iv)}}). Now \(\phi\) is the
composite
\phantomsection\label{form:sga3-VIB-7.5-factorization}%
\[
  V\times_S V
  \xrightarrow{\ f|_V\times f|_V\ }
  U\times_S U
  \xrightarrow{\ \mu\ }
  G,
\]
where the first morphism is faithfully flat and locally of
finite presentation, since \(f|_V\) is. It therefore remains
only to prove the same assertions for the restriction of
\(\mu\) to \(U\times_S U\).

Since \(G\to S\) is locally of finite presentation and flat, so
is
\[
  \mu:G\times_S G\longrightarrow G.
\]
Indeed, this morphism is isomorphic to the morphism obtained
from \(G\to S\) by the base change \(G\to S\)
(\hyperref[sga3:VIA:0.1]{VI\(_A\), 0.1}). The induced morphism
\[
  U\times_S U\longrightarrow G
\]
therefore has the same properties. To prove that this last
morphism is surjective, it is enough to check fiber by fiber.
There the result follows from
\hyperref[sga3:VIA:0.5]{VI\(_A\), 0.5}, since
\(U\cap G_s\) is a dense open subset of \(G_s\) for every
\(s\in S\).

% END INCLUDED SGA3_English_current_buildable_source_20260823/source/inputs/VI/tex_reference_v2/components/66_expose_VIB_lemma75_en.tex

% BEGIN INCLUDED SGA3_English_current_buildable_source_20260823/source/inputs/VI/tex_reference_v2/components/67_expose_VIB_rem760_prop76_en.tex
% SGA 3, Expose VI_B, Loop 1 English reconstruction.
% Scope: Remark 7.6.0 and Proposition 7.6 in full, including the
% special-case paragraph, complete proof, and source editor's notes
% 76--77.
% Authority: Polo--Gille Exp6B-13oct24.pdf, component-local pages
% 56--57 (combined-reader pages 397--398; re-edition headers
% 386--387; original-margin printed page 370).
% The preceding Lemma 7.5 on local page 56 and the following
% Remark 7.6.1 on local page 57 are excluded.
% OCR and Reinhold English Markdown are comparison witnesses only.
% This unit contains no diagrams.

\medskip
\phantomsection\label{rem:sga3-VIB-7.6.0}%
\noindent\textbf{Remark 7.6.0.---}%
\footnote{\label{note:sga3-VIB-76}Editor's note: this definition
has been added; in the original it was contained in the statement
of \hyperref[prop:sga3-VIB-7.6]{Proposition 7.6}.}
\emph{Let \(S\) be a scheme, let \(H\) be an \(S\)-group, and
let \(f:X\to H\) be a morphism of \(S\)-schemes. The group
subpresheaf of \(H\) generated by the image of \(f\), denoted}
\[
  \bigl\langle\operatorname{Im}f\bigr\rangle,
\]
\emph{is the group subfunctor of \(H\) which assigns to every
\(S\)-scheme \(S'\) the subgroup of \(H(S')\) generated by
\(f(X(S'))\). Since every element of this subgroup can be written
as a finite product}
\phantomsection\label{form:sga3-VIB-7.6.0-finite-product}%
\[
  f(x_1)^{\varepsilon_1}\cdots f(x_n)^{\varepsilon_n},
  \qquad
  n\in\mathbb{N}^{*},\quad
  x_i\in X(S'),\quad
  \varepsilon_i=\pm1,
\]
\emph{it follows that, if one puts}
\[
  X^1=X\amalg X
\]
\emph{and defines the morphisms \(f^n:X^n\to H\) as in
\hyperref[concl:sga3-VIB-7.1-ii]{7.1}, then
\(\langle\operatorname{Im}f\rangle\) is precisely the image
presheaf of the morphism}
\phantomsection\label{form:sga3-VIB-7.6.0-finfty}%
\[
  f^\infty:X^\infty\longrightarrow H,
\]
\emph{where}
\[
  X^\infty=\coprod_{n\geq1}X^n
\]
\emph{and \(f^\infty:X^\infty\to H\) is the \(S\)-morphism
whose restriction to each \(X^n\) is \(f^n\).}

\medskip
\phantomsection\label{prop:sga3-VIB-7.6}%
\noindent\textbf{Proposition 7.6.---}%
\footnote{\label{note:sga3-VIB-77}Editor's note: the original
stated this result under the hypotheses of the special case, but
the more general form was used implicitly in the proof of \SGARefLink{prop:sga3-VIB-10.12}{10.12};
the statement has been rewritten accordingly.}
\emph{Let \(S\) be a scheme, let \(X\) be an \(S\)-scheme flat
and locally of finite presentation over \(S\), let \(H\) be an
\(S\)-group locally of finite presentation over \(S\) and with
reduced fibers, let \(f:X\to H\) be a morphism of \(S\)-schemes,
and let \(f^\infty:X^\infty\to H\) be the \(S\)-morphism
introduced above. The following conditions are equivalent:}
\begin{enumerate}
\item[\textup{\SGAInlineRefTarget{prop:sga3-VIB-7.6:item:i}(i)}]\phantomsection\label{cond:sga3-VIB-7.6-i}
  \emph{\(H\) represents the \(S\)-sheaf for the
  \((\mathrm{fppf})\) topology associated with the presheaf
  \(\langle\operatorname{Im}f\rangle\).}

\item[\textup{\SGAInlineRefTarget{prop:sga3-VIB-7.6:item:ii}(ii)}]\phantomsection\label{cond:sga3-VIB-7.6-ii}
  \emph{\(f^\infty\) is a covering for the
  \((\mathrm{fppf})\) topology.}

\item[\textup{\SGAInlineRefTarget{prop:sga3-VIB-7.6:item:iii}(iii)}]\phantomsection\label{cond:sga3-VIB-7.6-iii}
  \emph{\(f^\infty\) is surjective, that is,}
  \phantomsection\label{form:sga3-VIB-7.6-union}%
  \[
    H=\bigcup_{n\geq1}f^n(X^n).
  \]
\end{enumerate}
\emph{If, in addition, \(H\) is quasi-compact, these conditions
are also equivalent to the following one:}
\begin{enumerate}
\item[\textup{(iv)}]\phantomsection\label{cond:sga3-VIB-7.6-iv}
  \emph{There exists an integer \(n\) such that
  \(f^n:X^n\to H\) is a covering for the
  \((\mathrm{fppf})\) topology, and hence is surjective.}
\end{enumerate}

\emph{This applies in particular when \(X\) is a
geometrically reduced \(k\)-scheme locally of finite type,
\(G\) is a \(k\)-group locally of finite type,
\(\varphi:X\to G\) is a morphism of \(k\)-schemes,
\(H=\Gamma_G(\varphi)\), and \(f:X\to H\) is the morphism
induced by \(\varphi\).}

\emph{Proof.}
The sheaf considered in \textup{\SGARefLink{prop:sga3-VIB-7.6:item:i}{(i)}} is the image sheaf of
\(X^\infty\) under \(f^\infty\). Hence, by
\hyperref[sga3:iv:proposition:4:4:3]{IV, 4.4.3}, saying that
\(H\) represents this sheaf is equivalent to saying that
\(f^\infty\) is a covering; this in turn implies that
\(f^\infty\) is surjective. Conversely, suppose that
\(f^\infty\) is surjective. We shall show that it is then a
covering for the \((\mathrm{fppf})\) topology.

Let \(s\in S\), let \(\eta\) be a maximal point of the fiber
\(H_s\), and let \(x\in X^\infty_s\) be such that
\[
  f^\infty(x)=\eta;
\]
such an \(x\) exists because \(f^\infty\) is surjective. Since
the fiber \(H_s\) is reduced, \(\mathcal{O}_{H_s,\eta}\) is a
field, and therefore \(f^\infty_s\) is flat at \(x\). Since
\(X^\infty\) and \(H\) are locally of finite presentation over
\(S\), and \(X^\infty\) is flat over \(S\), the fiberwise
criterion for flatness (EGA~IV\(_3\), 11.3.10) shows that
\(f^\infty\) is flat at \(x\). Thus, by
\hyperref[lem:sga3-VIB-7.5]{Lemma 7.5}, the morphism
\phantomsection\label{form:sga3-VIB-7.6-finfty-product}%
\[
  \mu\circ
  \bigl(f^\infty\times_S f^\infty\bigr):
  X^\infty\times_S X^\infty\longrightarrow H
\]
is a covering for the \((\mathrm{fppf})\) topology.

Now \(X^n\times_S X^m\) is canonically isomorphic to
\(X^{n+m}\), and, under this isomorphism,
\[
  \mu\circ\bigl(f^n\times_S f^m\bigr)
\]
corresponds to \(f^{n+m}\). It is therefore clear that
\[
  \mu\circ
  \bigl(f^\infty\times_S f^\infty\bigr)
\]
factors through \(f^\infty\), so \(f^\infty\) is a covering
for the \((\mathrm{fppf})\) topology. This proves
\(\textup{(iii)}\Rightarrow\textup{(ii)}\), and hence the
equivalence of conditions \textup{\SGARefLink{prop:sga3-VIB-7.6:item:i}{(i)}}, \textup{\SGARefLink{prop:sga3-VIB-7.6:item:ii}{(ii)}}, and
\textup{\SGARefLink{prop:sga3-VIB-7.6:item:iii}{(iii)}}.

Furthermore, since the morphisms \(X\to S\) and \(H\to S\) are
locally of finite presentation, so is \(f:X\to H\)
(EGA~IV\(_1\), 1.4.3\textup{(v)}). The multiplication
\[
  \mu:H\times_S H\longrightarrow H
\]
is also locally of finite presentation
(\hyperref[sga3:VIA:0.1]{VI\(_A\), 0.1}), and it follows that
each \(f^n:X^n\to H\) is locally of finite presentation.
Consequently, by EGA~IV\(_1\), 1.9.5\textup{(viii)}, the sets
\[
  f^n(X^n)
\]
are ind-constructible subsets of \(H\).

Suppose, in addition, that \(H\) is quasi-compact; then \(S\) is
also quasi-compact, since \(H\to S\) is surjective. By
EGA~IV\(_1\), 1.9.9, there exists an integer \(p\) such that
\[
  f^p(X^p)=H.
\]
As above, the fact that the fibers of \(H\) are reduced,
together with \hyperref[lem:sga3-VIB-7.5]{Lemma 7.5}, shows
that the morphism
\phantomsection\label{form:sga3-VIB-7.6-f2p-covering}%
\[
  \mu\circ\bigl(f^p\times_S f^p\bigr):
  X^p\times_S X^p\longrightarrow H
\]
is a covering for the \((\mathrm{fppf})\) topology. Since this
morphism is
\[
  f^{2p}:X^{2p}\longrightarrow H,
\]
this completes the proof of \hyperref[prop:sga3-VIB-7.6]{Proposition 7.6}.

% END INCLUDED SGA3_English_current_buildable_source_20260823/source/inputs/VI/tex_reference_v2/components/67_expose_VIB_rem760_prop76_en.tex

% BEGIN INCLUDED SGA3_English_current_buildable_source_20260823/source/inputs/VI/tex_reference_v2/components/68_expose_VIB_rem761_lemma77_en.tex
% SGA 3, Expose VI_B, Loop 1 English reconstruction.
% Scope: Remark 7.6.1 and Lemma 7.7 with its complete proof.
% Authority: Polo--Gille Exp6B-13oct24.pdf, component-local pages
% 57--58 (combined-reader pages 398--399; re-edition headers
% 387--388; original-margin printed pages 370--371).
% Proposition 7.8 begins later on local page 58 and is excluded.
% OCR and Reinhold English Markdown are comparison witnesses only.
% This unit contains no diagrams.

\medskip
\phantomsection\label{rem:sga3-VIB-7.6.1}%
\noindent\textbf{Remark 7.6.1.---}
The equivalent conditions of
\hyperref[prop:sga3-VIB-7.6]{7.6} plainly imply that the sheaf
\(F\) under consideration is representable. The converse is
false in general:% 
\footnote{\label{note:sga3-VIB-78}Editor's note: the argument
in the original has been expanded in what follows.}
for example, if \(k\) has characteristic \(0\), let
\(G=\mathbb G_{a,k}\), and let
\[
  f:\operatorname{Spec}k\longrightarrow G
\]
be the morphism defined by the point \(1\) of \(G(k)\). Then
\(F\) is represented by the constant \(k\)-group
\(\mathbb Z_k\), whereas \(\Gamma_G(f)=G\); hence the
monomorphism
\[
  \mathbb Z_k\hookrightarrow G
\]
is not surjective.

For simplicity, suppose now that the hypotheses of the special
case in \hyperref[prop:sga3-VIB-7.6]{7.6} hold, and suppose
that \(F\) is representable. By EGA~IV\(_3\), 8.14.2, the
sheaf \(F\) is locally of finite presentation over \(k\).
The question is therefore whether the dominant monomorphism
\[
  F\longrightarrow\Gamma_G(f)
\]
is an isomorphism, or, equivalently, a closed immersion. By
\hyperref[cor:sga3-VIB-1.4.2]{1.4.2}, this is so if \(F\) is
quasi-compact; and by
\hyperref[sga3:VIA:0.5.1]{VI\(_A\), 0.5.1}, it is so if \(F\)
is connected. In particular, by
\hyperref[cor:sga3-VIB-7.2.1]{7.2.1}, it is so when \(X\) is
connected and \(f(X)\) contains the unit element of \(G\).

\medskip
\phantomsection\label{lem:sga3-VIB-7.7}%
\noindent\textbf{Lemma 7.7.---}
\emph{Let \(k\) be an algebraically closed field, let \(G\) be
a \(k\)-group locally of finite type, let \(X\) be a
geometrically reduced \(k\)-scheme locally of finite type, let
\(f:X\to G\) be a \(k\)-morphism, and let \(H\) be a group
subscheme of \(G\) such that \(H\subset f(X)\). Put}
\phantomsection\label{form:sga3-VIB-7.7-Gamma-prime}%
\[
  \Gamma'=\bigcup_{n\geq1}f^n(X^n),
  \qquad
  \Gamma'_0=\Gamma'\cap G(k),
  \qquad
  H_0=H(k),
\]
\emph{and suppose that \(H_0\) has finite index in
\(\Gamma'_0\). Then there exists an integer \(m\) such that}
\[
  f^m(X^m)=\Gamma_G(f)
\]
\emph{(cf.\ \hyperref[prop:sga3-VIB-7.6]{7.6}), and
\(\Gamma_G(f)\) is the union of finitely many translates of
\(H\).}

For every \(n\geq1\), the subset \(f^n(X^n)\) of \(G\) is
ind-constructible (EGA~IV\(_1\), 1.9.5\textup{(viii)}).
The same is therefore true of \(\Gamma'\). Since \(G\) is a
Jacobson scheme, \(\Gamma'_0\) is dense in \(\Gamma'\). By
hypothesis, there are finitely many points
\[
  a_1,\ldots,a_r\in\Gamma'_0
\]
such that
\[
  \Gamma'_0=a_1H_0\cup\cdots\cup a_rH_0.
\]
Consequently,
\phantomsection\label{form:sga3-VIB-7.7-closure-chain}%
\[
\begin{aligned}
  \Gamma_G(f)
    &=\overline{\Gamma'}
     =\overline{\Gamma'_0}
     =\overline{a_1H_0\cup\cdots\cup a_rH_0}\\
    &=\overline{a_1H_0}\cup\cdots\cup\overline{a_rH_0}\\
    &=a_1\overline{H_0}\cup\cdots\cup
      a_r\overline{H_0},
\end{aligned}
\]
where the last equality follows because translation by each
\(a_i\) is a homeomorphism of \(G\) onto itself. Thus
\[
  \Gamma_G(f)=a_1H\cup\cdots\cup a_rH.
\]

It is also clear that there exists an integer \(p\) such that
every \(a_i\), for \(1\leq i\leq r\), belongs to
\(f^p(X^p)\). Finally, since \(H\subset f(X)\), one has
\[
  a_iH\subset f^{p+1}(X^{p+1})
\]
for every \(i\). Hence
\[
  f^{p+1}(X^{p+1})=\Gamma_G(f).
\]

% END INCLUDED SGA3_English_current_buildable_source_20260823/source/inputs/VI/tex_reference_v2/components/68_expose_VIB_rem761_lemma77_en.tex

% BEGIN INCLUDED SGA3_English_current_buildable_source_20260823/source/inputs/VI/tex_reference_v2/components/69_expose_VIB_prop78_en.tex
% SGA 3, Expose VI_B, Loop 1 English reconstruction.
% Scope: Proposition 7.8 and its complete proof, including source
% editor's notes 79--80.
% Authority: Polo--Gille Exp6B-13oct24.pdf, component-local pages
% 58--59 (combined-reader pages 399--400; re-edition headers
% 388--389; original-margin printed pages 372--373).
% Proposition 7.9 begins on local page 60 and is excluded.
% OCR and Reinhold English Markdown are comparison witnesses only.
% This unit contains no diagrams.

\medskip
\phantomsection\label{prop:sga3-VIB-7.8}%
\noindent\textbf{Proposition 7.8.---}
\emph{Let \(G\) be a \(k\)-group locally of finite type, and
let \(A\) and \(B\) be two geometrically reduced \(k\)-group
subschemes of \(G\) (hence smooth at the generic points of
their irreducible components, and therefore smooth by
\(1.3\)). Suppose that one of the following conditions holds:}
\begin{enumerate}
\item[\textup{a)}]\phantomsection\label{cond:sga3-VIB-7.8-a}
  \emph{\(A\) and \(B\) are invariant and of finite type over
  \(k\);}
\item[\textup{b)}]\phantomsection\label{cond:sga3-VIB-7.8-b}
  \emph{\(A\) is connected and \(B\) is of finite type over
  \(k\).}
\end{enumerate}
\emph{Then \((A,B)\) is of finite type over \(k\), and
represents the sheaf associated, for the \((\mathrm{fppf})\)
topology (or the \((\mathrm{fpqc})\) topology), with the group
presheaf of commutators of \(A\) and \(B\) in \(G\).
Moreover,}%
\footnote{\label{note:sga3-VIB-79}Editor's note: the assertions
that follow have been made explicit, and the reduction to the
algebraically closed case has been expanded in the proof.}
\emph{the \(k\)-groups \((A,B^0)\) and \((A^0,B)\) are
connected, and}
\phantomsection\label{form:sga3-VIB-7.8-identity-component}%
\[
  (A,B)^0=(A,B^0)\cdot(A^0,B).
\]

\emph{Proof.}
By \hyperref[prop:sga3-VIB-7.6]{7.6}, to show that \((A,B)\)
is the desired associated sheaf, it is enough to show that
there exists an integer \(n\) such that
\phantomsection\label{form:sga3-VIB-7.8-finite-generation}%
\[
  \nu^n\bigl((A\times_kB)^n\bigr)=(A,B),
\]
with the notation of
\hyperref[def:sga3-VIB-7.2.2]{7.2.2}. To prove this and the
other two assertions, we may suppose that \(k\) is
algebraically closed.

Indeed, let \(\overline{k}\) be an algebraic closure of \(k\).
By \hyperref[concl:sga3-VIB-7.1-iii]{7.1\textup{(iii)}} and
\hyperref[prop:sga3-VIA-2.4]{VI\(_A\), 2.4}, one has, with
the evident notation,
\phantomsection\label{form:sga3-VIB-7.8-base-change}%
\[
\begin{gathered}
  (A,B)_{\overline{k}}
    =(A_{\overline{k}},B_{\overline{k}}),\\
  \bigl((A,B)^0\bigr)_{\overline{k}}
    =(A_{\overline{k}},B_{\overline{k}})^0,\\
  (A,B^0)_{\overline{k}}
    =(A_{\overline{k}},B_{\overline{k}}^0),
  \qquad\text{and so on.}
\end{gathered}
\]
Consequently, suppose that
\((A_{\overline{k}},B_{\overline{k}})\) is of finite type over
\(\overline{k}\); respectively, suppose that
\((A_{\overline{k}},B_{\overline{k}}^0)\) and
\((A_{\overline{k}}^0,B_{\overline{k}})\) are connected and
that the morphism\begingroup
\renewcommand{\thefootnote}{\fnsymbol{footnote}}%
\footnote[1]{\label{note:sga3-VIB-7.8-mixed-base-change}
Translator's note: the French source prints
\((A_{\overline{k}},B_{\overline{k}}^0)\) in both factors of
this product. The second factor must be
\((A_{\overline{k}}^0,B_{\overline{k}})\), as in the asserted
identity and throughout the proof; the mixed product is
restored here.}%
\endgroup
\[
  (A_{\overline{k}},B_{\overline{k}}^0)
  \times_{\overline{k}}
  (A_{\overline{k}}^0,B_{\overline{k}})
  \longrightarrow
  (A_{\overline{k}},B_{\overline{k}})^0
\]
is surjective. The analogous assertions then hold over \(k\),
by EGA~IV\(_2\), 2.7.1 and 2.6.1.

Let
\[
  B^1,\ldots,B^p
\]
be the connected components of \(B\) other than the identity
component \(B^0\). There are only finitely many, since \(B\)
is of finite type over \(k\), hence Noetherian. In case
\textup{a)}, likewise let
\[
  A^1,\ldots,A^q
\]
be the connected components of \(A\) other than \(A^0\). In
case \textup{b)}, only \(A^0=A\) is considered. Let
\(\nu_{ij}\) be the restriction of \(\nu\) to
\[
  A^i\times_kB^j.
\]
Each \(A^i\) and \(B^j\) is irreducible, by
\hyperref[cor:sga3-VIA-2.4.1]{VI\(_A\), 2.4.1}; hence so are
\(A^0\times_kB^j\) and \(A^i\times_kB^0\). Since the unit
element of \(G\) belongs to \(A^0\) and to \(B^0\), it belongs
to
\[
  \nu_{0j}(A^0\times_kB^j)
  \qquad\text{and}\qquad
  \nu_{i0}(A^i\times_kB^0).
\]
Thus every \(\Gamma_G(\nu_{0j})\) and every
\(\Gamma_G(\nu_{i0})\) is connected, by
\hyperref[cor:sga3-VIB-7.2.1]{7.2.1}.

Let \(u_{0j}\), respectively \(u_{i0}\), denote the inclusion
of \(\Gamma_G(\nu_{0j})\), respectively
\(\Gamma_G(\nu_{i0})\), in \(G\). Then
\phantomsection\label{form:sga3-VIB-7.8-generated-connected-groups}%
\[
  (A^0,B)=\Gamma_G\bigl((u_{0j})_{j=0}^{p}\bigr),
  \qquad
  (A,B^0)=\Gamma_G\bigl((u_{i0})_{i=0}^{q}\bigr),
\]
and these groups are connected. Moreover, it follows from
\hyperref[prop:sga3-VIB-7.4]{7.4} and the preceding
constructions that there is an integer \(r\) such that both
\((A^0,B)\) and \((A,B^0)\) are contained in
\[
  \nu^r\bigl((A\times_kB)^r\bigr).
\]
In case \textup{b)}, one has \((A,B)=(A^0,B)\), and the proof
is complete.

We now consider case \textup{a)}.%
\footnote{\label{note:sga3-VIB-80}Editor's note: the argument
that follows has been expanded to take account of the added
\hyperref[cor:sga3-VIB-7.1.1]{Corollary 7.1.1}.}
We already know that \((A^0,B)\) and \((A,B^0)\) are smooth
connected \(k\)-subgroups of \(G\), hence are of finite type
over \(k\)
(\hyperref[prop:sga3-VIA-2.4]{VI\(_A\), 2.4}). Since \(A^0\)
is a characteristic subgroup of \(A\)
(\hyperref[thm:sga3-VIA-2.6.5]{VI\(_A\), 2.6.5}), it is an
invariant subgroup of \(G\); hence
\(\hyperref[concl:sga3-VIB-7.3-v]{7.3\textup{(v)}}\) shows
that \((A^0,B)\) is an invariant subgroup of \(G\), and the
same holds for \((A,B^0)\).

By \hyperref[cor:sga3-VIB-7.1.1]{7.1.1}, the subgroup \(H\) of
\(G\) generated by \((A,B^0)\) and \((A^0,B)\) is therefore
\[
  H=(A,B^0)\cdot(A^0,B).
\]
In particular,
\phantomsection\label{form:sga3-VIB-7.8-H-contained}%
\[
  H\subset\nu^{2r}\bigl((A\times_kB)^{2r}\bigr).
\]

If \(T\) is a subset of \(G\) stable under the group law
(cf.\ \(3.0\)), write \(T_0\) for the group of \(k\)-points of
\(G\) belonging to \(T\). Put
\[
  \Gamma'
    =\bigcup_{q\geq1}\nu^q\bigl((A\times_kB)^q\bigr).
\]
By \hyperref[prop:sga3-VIB-7.9]{Proposition 7.9} below,
\phantomsection\label{form:sga3-VIB-7.8-rational-commutators}%
\[
  (A^0,B)_0=(A_0^0,B_0),
  \qquad
  (A,B^0)_0=(A_0,B_0^0),
  \qquad
  \Gamma'_0=(A_0,B_0).
\]
Therefore
\[
  H_0=(A_0^0,B_0)\cdot(A_0,B_0^0)
\]
has finite index in \(\Gamma'_0\) (Bible, Expos\'e~III,
Appendix), since \(A_0\) and \(B_0\) are invariant and
\(A_0^0\), respectively \(B_0^0\), has finite index in
\(A_0\), respectively \(B_0\).

The hypotheses of \hyperref[lem:sga3-VIB-7.7]{Lemma 7.7} now
hold. Since
\[
  H\subset\nu^{2r}\bigl((A\times_kB)^{2r}\bigr),
\]
there exists an integer \(m\) such that
\[
  \nu^{2rm}\bigl((A\times_kB)^{2rm}\bigr)=(A,B),
\]
and there are finitely many \(k\)-points
\(a_1,\ldots,a_n\) of \(G\) such that
\[
  (A,B)=a_1H\cup\cdots\cup a_nH.
\]
Since \(H\) is of finite type over \(k\), every \(a_iH\) is
quasi-compact. Their union \((A,B)\) is therefore
quasi-compact, hence of finite type over \(k\). Finally, \(H\)
is irreducible, so each \(a_iH\) is irreducible; and because
the unit \(e\) belongs to \(H\), it follows that
\[
  H=(A,B)^0.
\]

% END INCLUDED SGA3_English_current_buildable_source_20260823/source/inputs/VI/tex_reference_v2/components/69_expose_VIB_prop78_en.tex

% BEGIN INCLUDED SGA3_English_current_buildable_source_20260823/source/inputs/VI/tex_reference_v2/components/70_expose_VIB_prop79_cor710_en.tex
% SGA 3, Expose VI_B, Loop 1 English reconstruction.
% Scope: Proposition 7.9 and Corollary 7.10 in full.
% Authority: Polo--Gille Exp6B-13oct24.pdf, component-local page 60
% (combined-reader page 401; re-edition header 390; original-margin
% printed page 374).
% Section 8 begins later on local page 60 and is excluded.
% OCR and Reinhold English Markdown are comparison witnesses only.
% This unit contains no diagrams.

\medskip
\phantomsection\label{prop:sga3-VIB-7.9}%
\noindent\textbf{Proposition 7.9.---}%
\emph{Let \(k\) be an algebraically closed field, and let \(G\)
be a \(k\)-group locally of finite type.}
\begin{enumerate}
\item[\textup{\SGAInlineRefTarget{prop:sga3-VIB-7.9:item:i}(i)}]\phantomsection\label{concl:sga3-VIB-7.9-i}
  \emph{Let \(A\) and \(B\) be two ind-constructible subsets
  of \(G\). Denote by \(A_0\) the set of rational points of
  \(G\) that belong to \(A\). Then}
  \[
    (A\cdot B)_0=A_0\cdot B_0,
  \]
  \emph{where the product on the right is taken in the group
  \(G(k)\).}

\item[\textup{\SGAInlineRefTarget{prop:sga3-VIB-7.9:item:ii}(ii)}]\phantomsection\label{concl:sga3-VIB-7.9-ii}
  \emph{Let \(X\) be a geometrically reduced \(k\)-scheme
  locally of finite type, and let \(f:X\to G\) be a
  \(k\)-morphism. Put}
  \[
    \Gamma'_G(f)=
      \bigcup_{n\geq1}f^n(X^n).
  \]
  \emph{Then \(\Gamma'_G(f)_0\) is the subgroup of \(G(k)\)
  generated by \(f(X)_0\).}

\item[\textup{(iii)}]\phantomsection\label{concl:sga3-VIB-7.9-iii}
  \emph{In particular, let \(A\) and \(B\) be two smooth
  group subschemes of \(G\), and put}
  \[
    \Gamma'=
      \bigcup_{n\geq1}
        \nu^n\bigl((A\times_k B)^n\bigr)
  \]
  \emph{with the notation of
  \hyperref[def:sga3-VIB-7.2.2]{7.2.2}. Then
  \(\Gamma'_0\) is the commutator subgroup of \(A(k)\) and
  \(B(k)\) in \(G(k)\).}
\end{enumerate}

\emph{Proof.}
We first prove \textup{\SGARefLink{prop:sga3-VIB-7.9:item:i}{(i)}}. It is clear that
\[
  A_0\cdot B_0\subset(A\cdot B)_0.
\]
Conversely, let \(z\in(A\cdot B)_0\). Then
\(\mu^{-1}(z)\) is a closed subset of \(G\times_kG\), while
\(A\times_kB\) is an ind-constructible subset of
\(G\times_kG\) (cf.\ \hyperref[par:sga3-VIB-3.0]{3.0}).
Consequently,
\[
  \mu^{-1}(z)\cap(A\times_kB)
\]
is a nonempty ind-constructible subset of \(G\times_kG\).
By EGA~IV\(_3\), 10.4.8, it therefore contains a rational
point of \(G\times_kG\). Its projections \(x\) and \(y\) are
rational points of \(G\) such that \(x\in A\), \(y\in B\),
and \(x\cdot y=z\). Hence
\begingroup
\renewcommand{\thefootnote}{\fnsymbol{footnote}}%
\footnote[1]{\label{note:sga3-VIB-7.9-product}Translator's note:
the French proof prints \((A,B)_0=A_0\cdot B_0\) here.
The statement and the multiplication argument require
\((A\cdot B)_0=A_0\cdot B_0\), which is restored in the
English text.}%
\endgroup
\[
  (A\cdot B)_0=A_0\cdot B_0.
\]

To prove \textup{\SGARefLink{concl:sga3-VIB-7.9-ii}{(ii)}}, observe that, since \(f^n\) is locally
of finite type, \(f^n(X^n)\) is an ind-constructible subset
of \(G\) (EGA~IV\(_1\), 1.9.5\textup{(viii)})\begingroup
\renewcommand{\thefootnote}{\fnsymbol{footnote}}%
\footnote[2]{\label{note:sga3-VIB-7.9-EGA-volume}
Translator's note: the French source prints EGA~IV\(_4\) here.
The cited result 1.9.5\textup{(viii)} belongs to
EGA~IV\(_1\), which is restored in the English text.}%
\endgroup.
Assertion \textup{\SGARefLink{prop:sga3-VIB-7.9:item:i}{(i)}} then shows by induction that, if
\[
  A=f(X)_0,
\]
one has
\[
  f^n(X^n)_0=(A\cup A^{-1})^n.
\]
Consequently,
\phantomsection\label{form:sga3-VIB-7.9-generated-rational-points}%
\[
  \Gamma'_G(f)_0
    =\bigcup_{n\geq1}f^n(X^n)_0
    =\bigcup_{n\geq1}(A\cup A^{-1})^n,
\]
which is the subgroup of \(G(k)\) generated by
\(A=f(X)_0\).

Finally, \textup{(iii)} follows from \textup{\SGARefLink{prop:sga3-VIB-7.9:item:ii}{(ii)}} and the
definitions.

\medskip
\phantomsection\label{cor:sga3-VIB-7.10}%
\noindent\textbf{Corollary 7.10.---}%
\emph{Under the hypotheses of
\hyperref[prop:sga3-VIB-7.8]{7.8}, if \(k\) is algebraically
closed, then \((A,B)(k)\) is the commutator subgroup of
\(A(k)\) and \(B(k)\) in \(G(k)\).}

Indeed, it is enough to apply
\hyperref[concl:sga3-VIB-7.9-iii]{7.9\textup{(iii)}},
\hyperref[prop:sga3-VIB-7.8]{7.8}, and
\hyperref[prop:sga3-VIB-7.6]{7.6}.

% END INCLUDED SGA3_English_current_buildable_source_20260823/source/inputs/VI/tex_reference_v2/components/70_expose_VIB_prop79_cor710_en.tex

% BEGIN INCLUDED SGA3_English_current_buildable_source_20260823/source/inputs/VI/tex_reference_v2/components/71_expose_VIB_section8_lemma811_loop2_r1_en.tex
% SGA 3, Expose VI_B, Loop 1 English reconstruction.
% Scope: Section 8 heading, paragraph 8.1, and Lemma 8.1.1 in full,
% including source editor's note 81.
% Authority: Polo--Gille Exp6B-13oct24.pdf, component-local pages
% 60--61 (combined-reader pages 401--402; re-edition headers
% 390--391; original-margin printed pages 374--375).
% Lemma 8.1.2 begins on local page 61 and is excluded.
% OCR and Reinhold English Markdown are comparison witnesses only.
% Loop 2 replaces the source-bound commutative-square crop with a
% faithful native LaTeX reconstruction; the replacement is closed.

\section*{8. Solvable or nilpotent group schemes}
\phantomsection\label{sec:sga3-VIB-8}%
\addcontentsline{toc}{section}{8. Solvable or nilpotent group schemes}

\medskip
\phantomsection\label{par:sga3-VIB-8.1}%
\noindent\textbf{8.1.---}
Let \(\mathcal C\) be a category equipped with a topology \(T\)
(cf.\ Expos\'e~IV, \S4). For every presheaf \(P\) on
\(\mathcal C\), write \(P^\flat\) for the associated sheaf.

Let \(G\) be a presheaf of groups on \(\mathcal C\), let \(A\)
and \(B\) be two group subpresheaves of \(G\), and let
\(\operatorname{Comm}(A,B)\) be the group presheaf of
commutators of \(A\) and \(B\) in \(G\). Thus, for every
\(S\in\operatorname{Ob}\mathcal C\),
\(\operatorname{Comm}(A,B)(S)\) is the subgroup of \(G(S)\)
generated by the commutators
\[
  aba^{-1}b^{-1},
  \qquad a\in A(S),\quad b\in B(S).
\]
We write
\phantomsection\label{form:sga3-VIB-8.1-CommT}%
\[
  \operatorname{Comm}_T(A,B)
    =\operatorname{Comm}(A,B)^\flat.
\]

In the proof of \hyperref[prop:sga3-VIB-8.2]{8.2}, we shall
need the following results.%
\footnote{\label{note:sga3-VIB-81}Editor's note: these results
were stated without proof in the original. We have highlighted
them as \hyperref[lem:sga3-VIB-8.1.1]{Lemmas 8.1.1} and \SGARefLink{lem:sga3-VIB-8.1.2}{8.1.2} and supplied the details of the
proofs.}

\medskip
\phantomsection\label{lem:sga3-VIB-8.1.1}%
\noindent\textbf{Lemma 8.1.1.---}
\emph{Let \(A\subset G\) be sheaves of groups, with \(A\)
invariant in \(G\).}
\begin{enumerate}
\item[\textup{\SGAInlineRefTarget{lem:sga3-VIB-8.1.1:item:i}(i)}]\phantomsection\label{concl:sga3-VIB-8.1.1-i}
  \emph{\(\operatorname{Comm}_T(G,A)\) is the smallest
  invariant group subsheaf \(C\) of \(G\) such that the sheaf
  \((A/C)^\flat\), associated with the quotient presheaf
  \(A/C\), is central in \((G/C)^\flat\).}

\item[\textup{(ii)}]\phantomsection\label{concl:sga3-VIB-8.1.1-ii}
  \emph{In particular, \(\operatorname{Comm}_T(G,G)\) is the
  smallest invariant group subsheaf \(C\) of \(G\) such that
  the quotient sheaf \((G/C)^\flat\) is commutative.}
\end{enumerate}

Clearly, \textup{\SGARefLink{concl:sga3-VIB-8.1.1-ii}{(ii)}} is the special case \(A=G\) of
\textup{\SGARefLink{concl:sga3-VIB-8.1.1-i}{(i)}}, so it is enough to prove \textup{\SGARefLink{concl:sga3-VIB-8.1.1-i}{(i)}}. Let \(C\)
be a group subsheaf of \(A\), invariant in \(G\), such that the
quotient sheaf \((A/C)^\flat\) is central in
\((G/C)^\flat\). By Lemma IV, \SGARefLink{sga3:iv:lemma:4:4:8:1}{4.4.8.1}, the presheaves \(A/C\)
and \(G/C\) are separated. Hence, by IV, \SGARefLink{sga3:iv:proposition:4:3:11}{4.3.11}, all the
morphisms in the following diagram are monomorphisms:
\begin{center}
  \phantomsection\label{diag:sga3-VIB-8.1.1-monomorphisms}%
  \begin{tikzcd}[row sep=large,column sep=large]
    A/C \arrow[r,hook] \arrow[d,hook] &
    G/C \arrow[d,hook] \\
    (A/C)^\flat \arrow[r,hook] &
    (G/C)^\flat
  \end{tikzcd}
\end{center}
Since \((A/C)^\flat\) is central in \((G/C)^\flat\), the
presheaf \(A/C\) is central in \(G/C\). It follows that
\[
  \operatorname{Comm}(G,A)\subset C,
\]
and hence, by IV, \SGARefLink{sga3:iv:corollary:4:3:12}{4.3.12}, \(C\) contains
\(\operatorname{Comm}_T(G,A)\).

Conversely, \(\operatorname{Comm}(G,A)\) is an invariant,
separated group subpresheaf of \(A\)
(cf.\ IV, \SGARefLink{sga3:iv:definition:4:3:1}{4.3.1}, editor's note 24). Therefore, by IV,
\SGARefLink{sga3:iv:lemma:4:4:8:2}{4.4.8.2}\textup{\SGARefLink{lem:sga3-VIB-8.1.1:item:i}{(i)}} and IV, \SGARefLink{sga3:iv:proposition:4:3:11}{4.3.11},
\[
  C=\operatorname{Comm}_T(G,A)
\]
is an invariant group subsheaf of \(A\) containing
\(\operatorname{Comm}(G,A)\). Consequently, the presheaf
\(A/C\) is central in \(G/C\), and hence, by IV,
\SGARefLink{sga3:iv:lemma:4:4:8:2}{4.4.8.2}\textup{\SGARefLink{sga3:iv:lemma:4:4:8:2:item:ii:01}{(ii)}}, \((A/C)^\flat\) is central in
\((G/C)^\flat\). This proves \hyperref[lem:sga3-VIB-8.1.1]{Lemma 8.1.1}.

% END INCLUDED SGA3_English_current_buildable_source_20260823/source/inputs/VI/tex_reference_v2/components/71_expose_VIB_section8_lemma811_loop2_r1_en.tex

% BEGIN INCLUDED SGA3_English_current_buildable_source_20260823/source/inputs/VI/tex_reference_v2/components/72_expose_VIB_lemma812_loop2_r1_en.tex
% SGA 3, Expose VI_B, Loop 1 English reconstruction.
% Scope: Lemma 8.1.2 with its complete proof.
% Authority: Polo--Gille Exp6B-13oct24.pdf, component-local pages
% 61--62 (combined-reader pages 402--403; re-edition headers
% 391--392; original-margin printed pages 375--376).
% Proposition 8.2 begins on local page 62 and is excluded.
% OCR and Reinhold English Markdown are comparison witnesses only.
% Loop 2 replaces the source-bound composite-morphism crop with a
% faithful native LaTeX reconstruction; the replacement is closed.

\medskip
\phantomsection\label{lem:sga3-VIB-8.1.2}%
\noindent\textbf{Lemma 8.1.2.---}
\emph{Let \(G\) be a sheaf of groups, and let \(A,B\) be two
group subpresheaves of \(G\).}
\begin{enumerate}
\item[\textup{\SGAInlineRefTarget{lem:sga3-VIB-8.1.2:item:i}(i)}]\phantomsection\label{concl:sga3-VIB-8.1.2-i}
  \emph{The morphism}
  \[
    \tau:\operatorname{Comm}(A,B)
      \longrightarrow
      \operatorname{Comm}(A^\flat,B^\flat)
  \]
  \emph{is a covering monomorphism.}

\item[\textup{(ii)}]\phantomsection\label{concl:sga3-VIB-8.1.2-ii}
  \emph{Consequently, there is an isomorphism}
  \phantomsection\label{form:sga3-VIB-8.1.2-CommT-isomorphism}%
  \[
    \operatorname{Comm}_T(A,B)
      \xrightarrow{\ \sim\ }
    \operatorname{Comm}_T(A^\flat,B^\flat).
  \]
\end{enumerate}

\emph{Proof.}
\textup{\SGARefLink{concl:sga3-VIB-8.1.2-i}{(i)}} Since \(A\) (respectively, \(B\)) is a
subpresheaf of \(G\), the associated sheaf \(A^\flat\)
(respectively, \(B^\flat\)) is a subpresheaf of \(G\)
containing \(A\) (respectively, \(B\)). It follows that
\(\tau\) is a monomorphism.

We now show that \(\tau\) is a covering. Let
\(S\in\operatorname{Ob}\mathcal C\) and let
\[
  g\in\operatorname{Comm}(A^\flat,B^\flat)(S).
\]
There exist an integer \(n\geq1\) and, for \(1\leq i\leq n\),
elements
\[
  a'_i\in A^\flat(S),\qquad
  b'_i\in B^\flat(S),\qquad
  \varepsilon_i\in\{\pm1\},
\]
such that
\[
  g=(a'_1,b'_1)^{\varepsilon_1}\cdots
    (a'_n,b'_n)^{\varepsilon_n},
\]
where \((a,b)\) denotes the commutator \(aba^{-1}b^{-1}\).
There is a refinement \(R\) of \(S\) such that
\[
  a'_i\in A(R)
  \qquad\text{and}\qquad
  b'_i\in B(R)
\]
for every \(i\). Then \(g\) is the composite morphism
\begin{center}
  \phantomsection\label{diag:sga3-VIB-8.1.2-composite}%
  \begin{tikzcd}[column sep=large]
    R
      \arrow[r,"{(a'_1,\ldots,b'_n)}"] &
    (A\times B)^n
      \arrow[r,"{\Phi^{\varepsilon_1,\ldots,\varepsilon_n}}"] &
    \operatorname{Comm}(A,B)
  \end{tikzcd}
\end{center}
where
\(\Phi=\Phi^{\varepsilon_1,\ldots,\varepsilon_n}\) is the
morphism defined set-theoretically by
\phantomsection\label{form:sga3-VIB-8.1.2-Phi}%
\[
  \Phi(a_1,b_1,\ldots,a_n,b_n)
    =(a_1,b_1)^{\varepsilon_1}\cdots
     (a_n,b_n)^{\varepsilon_n}
\]
for every \(T\in\operatorname{Ob}\mathcal C\), with
\(a_i\in A(T)\) and \(b_i\in B(T)\). This shows that
\[
  g\in\operatorname{Comm}(A,B)(R).
\]
It follows, as in the proof of IV, \SGARefLink{sga3:iv:proposition:4:3:11}{4.3.11}\textup{\SGARefLink{lem:sga3-VIB-8.1.2:item:i}{(i)}}, that
\[
  \operatorname{Comm}(A,B)
    \longrightarrow
  \operatorname{Comm}(A^\flat,B^\flat)
\]
is a covering.

The morphism
\[
  \operatorname{Comm}(A^\flat,B^\flat)
    \longrightarrow
  \operatorname{Comm}(A^\flat,B^\flat)^\flat
\]
is likewise a covering monomorphism
(IV, \SGARefLink{sga3:iv:proposition:4:3:11}{4.3.11}\textup{\SGARefLink{sga3:iv:proposition:4:3:11:item:iv:01}{(iv)}}). Hence the same is true of
\[
  \operatorname{Comm}(A,B)
    \longrightarrow
  \operatorname{Comm}(A^\flat,B^\flat)^\flat.
\]
By IV, \SGARefLink{sga3:iv:corollary:4:3:12}{4.3.12}, we therefore obtain an isomorphism
\[
  \operatorname{Comm}(A,B)^\flat
    \xrightarrow{\ \sim\ }
  \operatorname{Comm}(A^\flat,B^\flat)^\flat.
\]
This proves \hyperref[lem:sga3-VIB-8.1.2]{Lemma 8.1.2}.

% END INCLUDED SGA3_English_current_buildable_source_20260823/source/inputs/VI/tex_reference_v2/components/72_expose_VIB_lemma812_loop2_r1_en.tex

% BEGIN INCLUDED SGA3_English_current_buildable_source_20260823/source/inputs/VI/tex_reference_v2/components/73_expose_VIB_prop82_def821_prop83_en.tex
% SGA 3, Expose VI_B, Loop 1 English reconstruction.
% Scope: Proposition 8.2, Definition 8.2.1, and Proposition 8.3
% in full.
% Authority: Polo--Gille Exp6B-13oct24.pdf, component-local pages
% 62--63 (combined-reader pages 403--404; re-edition headers
% 392--393; original-margin printed page 376).
% Proposition 8.4 begins on local page 63 and is excluded.
% OCR and Reinhold English Markdown are comparison witnesses only.
% This unit contains no diagrams.

\medskip
\phantomsection\label{prop:sga3-VIB-8.2}%
\noindent\textbf{Proposition 8.2.---}
\emph{Let \(\mathcal C\) be a category, let \(T\) be a
topology on \(\mathcal C\), let \(G\) be a sheaf of groups on
\(\mathcal C\), and let \(n\) be an integer with \(n\geq0\).
The following conditions are equivalent:}
\begin{enumerate}
\item[\textup{\SGAInlineRefTarget{prop:sga3-VIB-8.2:item:i}(i)}]\phantomsection\label{concl:sga3-VIB-8.2-i}
  \emph{Put \(K_0=G\), and, for \(p\geq1\), put}
  \[
    K_p=\operatorname{Comm}(K_{p-1},K_{p-1})
    \qquad
    \bigl(\textup{respectively, }\
      K_p=\operatorname{Comm}(G,K_{p-1})\bigr).
  \]
  \emph{Then \(K_n\) is the unit presheaf of groups.}

\item[\textup{\SGAInlineRefTarget{prop:sga3-VIB-8.2:item:ii}(ii)}]\phantomsection\label{concl:sga3-VIB-8.2-ii}
  \emph{Put \(K'_0=G\), and, for \(p\geq1\), put}
  \[
    K'_p=\operatorname{Comm}_T(K'_{p-1},K'_{p-1})
    \qquad
    \bigl(\textup{respectively, }\
      K'_p=\operatorname{Comm}_T(G,K'_{p-1})\bigr).
  \]
  \emph{Then \(K'_n\) is the unit sheaf of groups.}

\item[\textup{\SGAInlineRefTarget{prop:sga3-VIB-8.2:item:iii}(iii)}]\phantomsection\label{concl:sga3-VIB-8.2-iii}
  \emph{There is a sequence}
  \[
    G=H_0\supset H_1\supset\cdots\supset H_n
  \]
  \emph{of invariant group subsheaves of \(G\) such that, for
  every \(i\), the quotient sheaf \(H_i/H_{i+1}\) is
  commutative \textup{(respectively, central in
  \(G/H_{i+1}\))}, and \(H_n\) is the unit sheaf.}
\end{enumerate}

It is clear that \(K_n\subset K'_n\); consequently,
\textup{\SGARefLink{prop:sga3-VIB-8.2:item:ii}{(ii)}} implies \textup{\SGARefLink{prop:sga3-VIB-8.2:item:i}{(i)}}. We show that
\textup{\SGARefLink{prop:sga3-VIB-8.2:item:i}{(i)}} implies \textup{\SGARefLink{prop:sga3-VIB-8.2:item:ii}{(ii)}}. One has
\[
  K'_1=\operatorname{Comm}_T(G,G)=K_1^\flat,
\]
and induction using
\hyperref[lem:sga3-VIB-8.1.2]{Lemma 8.1.2} gives
\[
  K'_p=K_p^\flat
\]
for every \(p\). Consequently, if \(K_n\) is the unit
presheaf, then \(K'_n=K_n^\flat\) is the unit sheaf.

Finally, conditions \textup{\SGARefLink{prop:sga3-VIB-8.2:item:ii}{(ii)}} and \textup{\SGARefLink{prop:sga3-VIB-8.2:item:iii}{(iii)}} are
equivalent by
\hyperref[lem:sga3-VIB-8.1.1]{Lemma 8.1.1}.

\medskip
\phantomsection\label{def:sga3-VIB-8.2.1}%
\noindent\textbf{Definition 8.2.1.---}
\emph{When these conditions hold, the sheaf \(G\) is said to
be \textup{solvable of class \(n\)}
\textup{(respectively, nilpotent of class \(n\))}. If there is
an integer \(n\) for which these conditions hold, \(G\) is
said to be \textup{solvable}
\textup{(respectively, nilpotent)}.}

Observe that, by condition \textup{\SGARefLink{concl:sga3-VIB-8.2-i}{(i)}}, this does not depend
on the topology \(T\).

\medskip
\phantomsection\label{prop:sga3-VIB-8.3}%
\noindent\textbf{Proposition 8.3.---}
\emph{Let \(k\) be a field, let \(S\) be a nonempty
\(k\)-scheme, let \(\Omega\) be an algebraically closed
extension of \(k\), and let \(G\) be a smooth \(k\)-group of
finite type. The following conditions are equivalent:}
\begin{enumerate}
\item[\textup{\SGAInlineRefTarget{prop:sga3-VIB-8.3:item:i}(i)}]\phantomsection\label{concl:sga3-VIB-8.3-i}
  \emph{\(G\) is solvable of class \(n\)
  \textup{(respectively, nilpotent of class \(n\))}.}

\item[\textup{(ii)}]\phantomsection\label{concl:sga3-VIB-8.3-ii}
  \emph{\(G\times_kS\) is solvable of class \(n\)
  \textup{(respectively, nilpotent of class \(n\))}.}

\item[\textup{(iii)}]\phantomsection\label{concl:sga3-VIB-8.3-iii}
  \emph{The group \(G(\Omega)\) is solvable of class \(n\)
  \textup{(respectively, nilpotent of class \(n\))}.}

\item[\textup{(iv)}]\phantomsection\label{concl:sga3-VIB-8.3-iv}
  \emph{Put \(K_0=G\), and, for \(p\geq1\), consider the
  \(k\)-groups}
  \[
    K_p=(K_{p-1},K_{p-1})
    \qquad
    \bigl(\textup{respectively, }\ K_p=(G,K_{p-1})\bigr)
  \]
  \emph{\textup{(cf.\ \hyperref[def:sga3-VIB-7.2.2]{7.2.2})}.
  Then \(K_n\) is the unit \(k\)-group.}
\end{enumerate}

The equivalence of conditions \textup{\SGARefLink{prop:sga3-VIB-8.3:item:i}{(i)}} and \textup{\SGARefLink{concl:sga3-VIB-8.3-i}{(ii)}}
follows from
\hyperref[prop:sga3-VIB-8.2]{Proposition 8.2}, since the
formation of the group presheaf of commutators commutes with
base change (IV, \SGARefLink{sga3:iv:proposition:4:1:3}{4.1.3}).

The equivalence of \textup{\SGARefLink{prop:sga3-VIB-8.3:item:i}{(i)}} and \textup{\SGARefLink{concl:sga3-VIB-8.3-iii}{(iv)}} follows
from the fact that, by
\hyperref[prop:sga3-VIB-7.8]{7.8}, the \(k\)-group
\[
  K_p=(K_{p-1},K_{p-1})
  \qquad
  \bigl(\textup{respectively, }\ K_p=(G,K_{p-1})\bigr)
\]
represents the sheaf
\[
  \operatorname{Comm}_T(K_{p-1},K_{p-1})
  \qquad
  \bigl(\textup{respectively, }\
    \operatorname{Comm}_T(G,K_{p-1})\bigr),
\]
where \(T\) is the \textup{(fppf)} or \textup{(fpqc)}
topology.

To prove the equivalence of conditions \textup{\SGARefLink{concl:sga3-VIB-8.3-ii}{(iii)}} and
\textup{\SGARefLink{concl:sga3-VIB-8.3-iv}{(iv)}}, we may assume that \(k=\Omega\); the result then
follows from \hyperref[cor:sga3-VIB-7.10]{7.10}.

% END INCLUDED SGA3_English_current_buildable_source_20260823/source/inputs/VI/tex_reference_v2/components/73_expose_VIB_prop82_def821_prop83_en.tex

% BEGIN INCLUDED SGA3_English_current_buildable_source_20260823/source/inputs/VI/tex_reference_v2/components/74_expose_VIB_prop84_en.tex
% SGA 3, Expose VI_B, Loop 1 English reconstruction.
% Scope: Proposition 8.4 and its complete proof, including the
% source's starred note and editor's notes 82--84.
% Authority: Polo--Gille Exp6B-13oct24.pdf, component-local pages
% 63--65 (combined-reader pages 404--406; re-edition headers
% 393--395; original-margin printed pages 377--379).
% Section 9 begins on local page 65 and is excluded.
% OCR and Reinhold English Markdown are comparison witnesses only.
% This unit contains no diagrams.

\medskip
\phantomsection\label{prop:sga3-VIB-8.4}%
\noindent\textbf{Proposition 8.4.---}
\emph{Let \(G\) be an \(S\)-group of finite presentation such
that, for every \(s\in S\), the fiber \(G_s\) is smooth over
\(\kappa(s)\). Let \(T\) be the set of points \(s\in S\) for
which \(G_s\) is solvable \textup{(respectively, nilpotent)}.}
\begin{enumerate}
\item[\textup{(i)}]\phantomsection\label{concl:sga3-VIB-8.4-i}
  \emph{Then \(T\) is locally constructible in \(S\).}

\item[\textup{(ii)}]\phantomsection\label{concl:sga3-VIB-8.4-ii}
  \emph{If, in addition, \(G\) is assumed flat and separated
  over \(S\) \textup{(that is, when \(G\) is smooth,
  quasi-compact, and separated over \(S\))}, then \(T\) is
  closed in \(S\).}
\end{enumerate}

\emph{Proof.}
We may clearly suppose that \(S\) is affine, with ring \(A\).
By \hyperref[par:sga3-VIB-10.1]{10.1} and
\hyperref[par:sga3-VIB-10.10]{10.10\textup{b)}},%
\begingroup
\renewcommand{\thefootnote}{\fnsymbol{footnote}}%
\footnote[1]{\label{note:sga3-VIB-8.4-star}
In the course of this proof we shall use results established
in \hyperref[sec:sga3-VIB-10]{Section~10}; like \hyperref[sec:sga3-VIB-9]{Section~9}, they do not depend on the
present \hyperref[sec:sga3-VIB-8]{Section~8}.}%
\endgroup\ 
there exist a Noetherian subring \(A'\) of \(A\) and an
\(A'\)-group \(G'\) of finite type such that
\[
  G'\otimes_{A'}A\simeq G.
\]
By EGA~IV\(_3\), 11.2.6 and 8.10.5,%
\footnote{\label{note:sga3-VIB-82}Editor's note: the reference
10.8.5 has been corrected to 8.10.5.}
if \(G\) is flat and separated over \(S\), we may suppose that
\(G'\) is flat and separated over
\(S'=\operatorname{Spec}A'\).%
\footnote{\label{note:sga3-VIB-83}Editor's note: what follows
has been expanded.}
Since \(G\) is of finite presentation over \(S\), the set of
points \(s\in S\) for which \(G_s\) is geometrically reduced
---or, equivalently, smooth over \(\kappa(s)\)---is
locally constructible (EGA~IV\(_3\), 9.7.7). Hence, by
EGA~IV\(_3\), 9.3.3, we may suppose that \(G'_{s'}\) is smooth
over \(\kappa(s')\) for every \(s'\in S'\). On the other hand,
if \(s'\) denotes the image of \(s\) in \(S'\), then
\[
  G'_{s'}\otimes_{\kappa(s')}\kappa(s)\simeq G_s.
\]
It follows from \hyperref[prop:sga3-VIB-8.3]{8.3} that \(G_s\)
is solvable (respectively, nilpotent) if and only if the same
is true of \(G'_{s'}\). We are therefore reduced to the case
where \(S\) is a Noetherian affine scheme.

We first show that \(T\) is constructible. Applying the
criterion of EGA~0\(_{\mathrm{III}}\), 9.2.3, and arguing as
above, we are reduced to showing that, when \(S\) is
Noetherian and integral, either \(T\) or \(S-T\) contains a
nonempty open subset of \(S\).

Suppose, then, that \(S\) is integral and Noetherian, with
generic point \(\eta\). For every \(s\in S\), put
\(K_s^0=G_s\), and put, for \(p\geq1\),
\phantomsection\label{form:sga3-VIB-8.4-Ksp}%
\[
  K_s^p=(K_s^{p-1},K_s^{p-1})
  \qquad
  \bigl(\textup{respectively, }\
    K_s^p=(G_s,K_s^{p-1})\bigr).
\tag*{\(\dagger\)}
\]
\begingroup
\renewcommand{\thefootnote}{\fnsymbol{footnote}}%
\footnotetext[2]{\label{note:sga3-VIB-8.4-G-vs-Gs}
Translator's note: in the nilpotent branch of this display,
the French source prints \(G\). Since \(K_s^{p-1}\) is a
\(\kappa(s)\)-group and the subsequent formulas use \(G_s\),
the fiber subscript is restored here.}%
\endgroup
We first show that the sequence of closed subschemes
\(K_\eta^p\) is stationary. By
\hyperref[concl:sga3-VIB-7.3-v]{7.3\textup{(v)}}, each
\(K_\eta^p\) is invariant, and hence the sequence
\((K_\eta^p)_p\) is decreasing. It is therefore stationary,
since \(G_\eta\) is Noetherian. Thus there is an integer \(n\)
such that
\phantomsection\label{form:sga3-VIB-8.4-stationary}%
\[
  K_\eta^p=K_\eta^n
  \qquad\text{for every }p\geq n.
\]

On the other hand, by
\hyperref[cor:sga3-VIB-10.12.1]{10.12.1} and
\hyperref[cor:sga3-VIB-10.13]{10.13}, there exist a nonempty
open subset \(S'\) of \(S\) and an \(S'\)-group \(D\) of
finite presentation such that, for every \(s\in S'\),
\phantomsection\label{form:sga3-VIB-8.4-D-fibers}%
\[
  D_s=K_s^n,
  \qquad
  (D_s,D_s)=D_s
  \quad
  \bigl(\textup{respectively, }(G_s,D_s)=D_s\bigr).
\]
We may suppose that \(S'=S\). Then, for every \(s\in S\) and
every \(p\geq n\), one has \(D_s=K_s^p\). Consequently, \(G_s\)
is solvable (respectively, nilpotent) if and only if \(D_s\)
is isomorphic to the unit \(\kappa(s)\)-group.

By EGA~IV\(_3\), 9.6.1\textup{(xi)}, the set of points
\(s\in S\) for which the structural morphism
\[
  D_s\longrightarrow\operatorname{Spec}\kappa(s)
\]
is an isomorphism is constructible.%
\footnote{\label{note:sga3-VIB-84}Editor's note: here one uses
the fact that \(S\) is assumed affine, and is therefore
quasi-compact and quasi-separated
\textup{(cf.\ EGA~0\(_{\mathrm{III}}\), 9.1.12)}.}
It is therefore either nowhere dense or contains a nonempty
open subset of \(S\). We have thus proved that \(T\) is
locally constructible.

We now show that, if \(G\) is also flat and separated over
\(S\), then \(T\) is closed. Since \(T\) is locally
constructible, it is closed if and only if it is stable under
specialization (EGA~IV\(_1\), 1.10.1).

Let \(s\in S\), and let \(s'\) be a specialization of \(s\)
in \(S\). Since we have reduced to the case where \(S\) is
Noetherian, EGA~II, 7.1.9 provides a discrete valuation ring
\(A\) and a morphism
\[
  \operatorname{Spec}A\longrightarrow S
\]
such that \(s\), respectively \(s'\), is the image of the
generic point \(\alpha\), respectively the closed point \(a\),
of \(\operatorname{Spec}A\). It is enough to show that, if
\[
  G'=G\times_S\operatorname{Spec}A
\]
and \(G'_\alpha\) is solvable (respectively, nilpotent), then
the same is true of \(G'_a\). Since \(G\) is flat and
separated over \(S\), the \(A\)-group \(G'\) is flat and
separated.%
\begingroup
\renewcommand{\thefootnote}{\fnsymbol{footnote}}%
\footnote[3]{\label{note:sga3-VIB-8.4-Gprime-alpha}
Translator's note: the French source prints \(G'_\alpha\) as
being flat and separated over \(A\). The intended object is
the base-changed \(A\)-group \(G'\); its generic fiber
\(G'_\alpha\) is a group over \(\kappa(\alpha)\).}%
\endgroup\ 
We are thus reduced to the case where \(S\) is the spectrum of
a discrete valuation ring \(A\).

Since \(G_\alpha\) is assumed solvable (respectively,
nilpotent), there is an integer \(q\) such that
\(K_\alpha^q\), with the notation introduced above, is
isomorphic to the unit \(\kappa(\alpha)\)-group. For every
\(n\), let \(\overline{K_\alpha^n}\) denote the schematic
closure of \(K_\alpha^n\) in \(G\), in the sense of
EGA~IV\(_2\), 2.8.5. We show by induction on \(p\) that
\phantomsection\label{form:sga3-VIB-8.4-closure-induction}%
\[
  \bigl(\overline{K_\alpha^p}\bigr)_a
    \supset K_a^p.
\]
First, since \(G\) is flat over \(A\), one has
\(\overline{G_\alpha}=G\) by EGA~IV\(_2\), 2.8.5, and hence
\[
  \bigl(\overline{K_\alpha^0}\bigr)_a=K_a^0.
\]

Let \(p\geq0\), and suppose that
\[
  K_a^p\subset\bigl(\overline{K_\alpha^p}\bigr)_a.
\]
Following \hyperref[def:sga3-VIB-7.2.2]{7.2.2}, denote by
\(\nu_a,\nu_\alpha,\overline{\nu},\overline{\nu}_a\) the
following morphisms. In the solvable case they are
\phantomsection\label{form:sga3-VIB-8.4-nu-solvable}%
\[
\begin{aligned}
  \nu_a:\;&K_a^p\times_{\kappa(a)}K_a^p
    \longrightarrow G_a,\\
  \nu_\alpha:\;&K_\alpha^p\times_{\kappa(\alpha)}K_\alpha^p
    \longrightarrow G_\alpha,\\
  \overline{\nu}:\;&
    \overline{K_\alpha^p}\times_A\overline{K_\alpha^p}
    \longrightarrow G,\\
  \overline{\nu}_a:\;&
    \bigl(\overline{K_\alpha^p}\bigr)_a
    \times_{\kappa(a)}
    \bigl(\overline{K_\alpha^p}\bigr)_a
    \longrightarrow G_a;
\end{aligned}
\]
respectively, in the nilpotent case they are
\phantomsection\label{form:sga3-VIB-8.4-nu-nilpotent}%
\[
\begin{aligned}
  \nu_a:\;&G_a\times_{\kappa(a)}K_a^p
    \longrightarrow G_a,\\
  \nu_\alpha:\;&G_\alpha\times_{\kappa(\alpha)}K_\alpha^p
    \longrightarrow G_\alpha,\\
  \overline{\nu}:\;&G\times_A\overline{K_\alpha^p}
    \longrightarrow G,\\
  \overline{\nu}_a:\;&
    G_a\times_{\kappa(a)}
    \bigl(\overline{K_\alpha^p}\bigr)_a
    \longrightarrow G_a.
\end{aligned}
\]
Since \(\nu_\alpha\) factors through \(K_\alpha^{p+1}\),
\(\overline{\nu}\) factors through
\(\overline{K_\alpha^{p+1}}\). The latter is clearly a group
subscheme of \(G\), and therefore contains
\(\Gamma_G(\overline{\nu})\). By
\hyperref[concl:sga3-VIB-7.1-iii]{7.1\textup{(iii)}},
\[
  \Gamma_G(\overline{\nu})_a
    =\Gamma_{G_a}(\overline{\nu}_a).
\]
Moreover, the induction hypothesis gives
\phantomsection\label{form:sga3-VIB-8.4-induction-inclusions}%
\[
\begin{aligned}
  K_a^p\times_{\kappa(a)}K_a^p
  &\subset
  \bigl(\overline{K_\alpha^p}\bigr)_a
  \times_{\kappa(a)}
  \bigl(\overline{K_\alpha^p}\bigr)_a
  &&\text{in the solvable case},\\
  G_a\times_{\kappa(a)}K_a^p
  &\subset
  G_a\times_{\kappa(a)}
  \bigl(\overline{K_\alpha^p}\bigr)_a
  &&\text{in the nilpotent case}.
\end{aligned}
\]
Consequently,
\phantomsection\label{form:sga3-VIB-8.4-Gamma-chain}%
\[
  K_a^{p+1}
  =\Gamma_{G_a}(\nu_a)
  \subset\Gamma_{G_a}(\overline{\nu}_a)
  =\Gamma_G(\overline{\nu})_a
  \subset\bigl(\overline{K_\alpha^{p+1}}\bigr)_a.
\]

Finally, \(K_\alpha^q\) is isomorphic to the unit
\(\kappa(\alpha)\)-group. The unit \(A\)-group is flat over
\(A\) and is isomorphic to a closed subscheme of \(G\), since
\(G\) is separated over \(A\)
\textup{(cf.\ \hyperref[prop:sga3-VIB-5.1]{5.1})}. It follows
from EGA~IV\(_2\), 2.8.5 that
\(\overline{K_\alpha^q}\) is isomorphic to the unit
\(A\)-group. Since we have proved that
\[
  K_a^q\subset\bigl(\overline{K_\alpha^q}\bigr)_a,
\]
the group \(K_a^q\) is isomorphic to the unit
\(\kappa(a)\)-group. Thus \(G_a\) is solvable
\textup{(respectively, nilpotent)}.

% END INCLUDED SGA3_English_current_buildable_source_20260823/source/inputs/VI/tex_reference_v2/components/74_expose_VIB_prop84_en.tex

% BEGIN INCLUDED SGA3_English_current_buildable_source_20260823/source/inputs/VI/tex_reference_v2/components/75_expose_VIB_section9_def91_prop92_loop2_r1_en.tex
% SGA 3, Expose VI_B, Loop 1 English reconstruction.
% Scope: Section 9 heading and opening, Definition 9.1, and
% Proposition 9.2 with clauses (i)--(xiii) and complete proof.
% Authority: Polo--Gille Exp6B-13oct24.pdf, component-local pages
% 65--67 (combined-reader pages 406--408; re-edition headers
% 395--397; original-margin printed pages 380--383).
% Lemma 9.2.1 begins on local page 67 and is excluded.
% OCR and Reinhold English Markdown are comparison witnesses only.
% Loop 2 replaces both source-bound cartesian-diagram crop groups
% with faithful native LaTeX reconstructions; both are closed.

\section*{9. Quotient sheaves}
\phantomsection\label{sec:sga3-VIB-9}%
\addcontentsline{toc}{section}{9. Quotient sheaves}

The present section is essentially limited to recalling, in the
particular case of homogeneous spaces of groups, well-known general
facts concerning passage to the quotient by flat equivalence relations
(cf.\ Expos\'e~IV).

\medskip
\phantomsection\label{def:sga3-VIB-9.1}%
\noindent\textbf{Definition 9.1.---}
\emph{Given a monomorphism \(u:G'\to G\) of \(S\)-groups, we
write \(G/G'\) \textup{(respectively, \(G'\backslash G\))} and call it
the \textup{right \textup{(respectively, left)} quotient sheaf of
\(G\) by \(G'\)}, meaning the quotient sheaf, for the
\textup{(fpqc)} topology, of \(G\) by the equivalence relation defined
by the monomorphism}
\phantomsection\label{form:sga3-VIB-9.1-equivalence-relations}%
\[
\begin{aligned}
  G\times_S G'
    &\xrightarrow{\ \delta\circ(\operatorname{id}_G\times u)\ }
      G\times_S G,\\[-1mm]
  &\hspace{7em}\textup{\emph{\textup{(respectively,}}}\\[-1mm]
  G'\times_S G
    &\xrightarrow{\ \gamma\circ(u\times\operatorname{id}_G)\ }
      G\times_S G\textup{\emph{\textup{),}}}
\end{aligned}
\]
\emph{where \(\delta\) \textup{(respectively, \(\gamma\))} denotes
the automorphism of \(G\times_SG\) defined, for
\(g,h\in G(T)\), by}
\[
  (g,h)\longmapsto(g,gh)
  \qquad
  \bigl(\textup{respectively, }(h,g)\longmapsto(hg,g)\bigr).
\]

\medskip
\phantomsection\label{prop:sga3-VIB-9.2}%
\noindent\textbf{Proposition 9.2.---}
\emph{Let \(u:G'\to G\) be a monomorphism of \(S\)-groups.
Suppose that \(G/G'\) is representable by an \(S\)-scheme \(G''\).
Then:}
\begin{enumerate}
\item[\textup{(i)}]\phantomsection\label{concl:sga3-VIB-9.2-i}
  \emph{The canonical morphism \(p:G\to G''\) is a covering for
  the \textup{(fpqc)} topology.}

\item[\textup{(ii)}]\phantomsection\label{concl:sga3-VIB-9.2-ii}
  \emph{If \(\varepsilon''=p\circ\varepsilon\)
  \textup{(this morphism is called the unit section of \(G''\))},
  the following diagrams are cartesian:}
  \begin{center}
    \phantomsection\label{diag:sga3-VIB-9.2-cartesian}%
    \begin{tabular}{c@{\qquad}c}
      \begin{tikzcd}[row sep=large,column sep=large]
        G\times_SG'
          \arrow[r,"{\mu\circ(\operatorname{id}_G\times u)}"]
          \arrow[d,"\operatorname{pr}_1"'] &
        G \arrow[d,"p"] \\
        G \arrow[r,"p"'] &
        G''
      \end{tikzcd},
      &
      \begin{tikzcd}[row sep=large,column sep=large]
        G' \arrow[r,"u"] \arrow[d,"\pi'"'] &
        G \arrow[d,"p"] \\
        S \arrow[r,"\varepsilon''"'] &
        G''
      \end{tikzcd}.
    \end{tabular}
  \end{center}
  \emph{In particular, \(u\) is an immersion.}

\item[\textup{(iii)}]\phantomsection\label{concl:sga3-VIB-9.2-iii}
  \emph{There exists on \(G''\) a unique structure of an
  \(S\)-scheme with left operator group \(G\) such that \(p\) is a
  morphism of \(S\)-schemes with operator group \(G\).}

\item[\textup{(iv)}]\phantomsection\label{concl:sga3-VIB-9.2-iv}
  \emph{If, in addition, \(G'\) is assumed invariant in \(G\),
  there exists on \(G''\) a unique \(S\)-group structure such that
  \(p\) is a morphism of \(S\)-groups.}

\item[\textup{(v)}]\phantomsection\label{concl:sga3-VIB-9.2-v}
  \emph{Let \(S_0\) be an \(S\)-scheme, and put
  \(G_0=G\times_SS_0\) and \(G'_0=G'\times_SS_0\). Then
  \(G_0/G'_0\) is representable by
  \(G''_0=G''\times_SS_0\).}%
  \footnote{\label{note:sga3-VIB-85}Editor's note: that is, the
  quotient is universal; cf.\ Expos\'e~IV, \S3.}

\item[\textup{(vi)}]\phantomsection\label{concl:sga3-VIB-9.2-vi}
  \emph{Let \(\mathscr P\) be a property of an \(S\)-morphism.
  Suppose that \(\mathscr P\) is stable under base change. If
  \(p:G\to G''\) has property \(\mathscr P\), then the structural
  morphism \(\pi':G'\to S\) has property \(\mathscr P\) as well.}

\item[\textup{(vii)}]\phantomsection\label{concl:sga3-VIB-9.2-vii}
  \emph{Let \(\mathscr P\) be a property of an \(S\)-morphism.
  Suppose that \(\mathscr P\) is local for the \textup{(fpqc)}
  topology \textup{(cf.\
  \hyperref[par:sga3-VIB-2.0]{2.0} and
  \hyperref[rem:sga3-VIB-2.1.2]{2.1.2})}. Then \(p:G\to G''\)
  has property \(\mathscr P\) if and only if
  \(\pi':G'\to S\) has property \(\mathscr P\).}%
  \begingroup
  \renewcommand{\thefootnote}{\fnsymbol{footnote}}%
  \footnote[1]{\label{note:sga3-VIB-9.2-vii-base}
  Translator's note: the French source prints \(S'\) as the target of
  \(\pi'\) in this clause. No \(S'\) is defined here, and clause
  \textup{\SGARefLink{concl:sga3-VIB-9.2-vi}{(vi)}} identifies \(\pi'\) as the structural morphism
  \(G'\to S\); that target is used here.}%
  \endgroup

\item[\textup{(viii)}]\phantomsection\label{concl:sga3-VIB-9.2-viii}
  \emph{Let \(\mathscr P\) be a property of an \(S\)-morphism.
  Suppose that \(\mathscr P\) is local for the \textup{(fpqc)}
  topology and stable under composition. If the structural
  morphisms \(G''\to S\) and \(G'\to S\) have property
  \(\mathscr P\), then the structural morphism \(G\to S\) has
  property \(\mathscr P\) as well.}

\item[\textup{(ix)}]\phantomsection\label{concl:sga3-VIB-9.2-ix}
  \emph{If \(G\) is reduced, then \(G''\) is reduced.}

\item[\textup{(x)}]\phantomsection\label{concl:sga3-VIB-9.2-x}
  \emph{For \(G''\) to be separated over \(S\), it is necessary
  and sufficient that \(u\) \textup{(or, equivalently,
  \(\varepsilon''\))} be a closed immersion.}

\item[\textup{(xi)}]\phantomsection\label{concl:sga3-VIB-9.2-xi}
  \emph{For \(G'\) to be flat over \(S\), it is necessary and
  sufficient that \(p\) be flat \textup{(or, equivalently,
  faithfully flat)}.}

  \emph{In this case, for \(G''\) to be flat over \(S\), it is
  necessary and sufficient that \(G\) be flat over \(S\).}

\item[\textup{(xii)}]\phantomsection\label{concl:sga3-VIB-9.2-xii}
  \emph{For \(G'\) to be flat and locally of finite presentation
  over \(S\), it is necessary and sufficient that \(p:G\to G''\)
  be faithfully flat and locally of finite presentation.}

  \emph{In this case, for \(G''\) to be locally of finite
  presentation \textup{(respectively, locally of finite type, of
  finite type, smooth, \'etale, unramified, locally quasi-finite,
  quasi-finite)} over \(S\), it is enough that \(G\) have the same
  property over \(S\) \textup{(and the condition is also necessary
  in the first two cases; cf.\ \textup{(viii)})}.}

\item[\textup{(xiii)}]\phantomsection\label{concl:sga3-VIB-9.2-xiii}
  \emph{Suppose that \(G'\) is flat and of finite presentation
  over \(S\).}
  \begin{enumerate}
  \item[\textup{a)}]\phantomsection\label{concl:sga3-VIB-9.2-xiii-a}
    \emph{Then \(p\) is of finite presentation and faithfully
    flat.}
  \item[\textup{b)}]\phantomsection\label{concl:sga3-VIB-9.2-xiii-b}
    \emph{Moreover, for \(G\) to be of finite presentation over
    \(S\), it is necessary and sufficient that \(G''\) be so.}
  \end{enumerate}
\end{enumerate}

\emph{Proof.}
Recall that the equivalence relation under consideration is
universally effective (IV, \SGARefLink{sga3:iv:proposition:4:4:9}{4.4.9}). Assertions
\textup{(i)}, \textup{(iii)}, \textup{(iv)}, \textup{(v)}, and the
first assertion of \textup{(ii)} then follow, respectively, from
IV, \SGARefLink{sga3:iv:proposition:4:4:3}{4.4.3}, 5.2.2, 5.2.4, 3.4.5, and \SGARefLink{thm:sga3-VIA-3.3.2}{3.3.2}\textup{(iii)}. The second
assertion of \textup{(ii)} follows from the first, as shown by the
following cartesian diagram, since
\[
  (G\times_SG')\times_GS
\]
is isomorphic to \(G'\):
\begin{center}
  \phantomsection\label{diag:sga3-VIB-9.2-proof-cartesian}%
  {\small
  \begin{tikzcd}[row sep=large,column sep=large]
    G'
      \arrow[r,"{((\varepsilon\circ\pi'),\operatorname{id}_{G'})}"]
      \arrow[d,"\pi'"'] &
    G\times_SG'
      \arrow[r,"{\mu\circ(\operatorname{id}_G\times u)}"]
      \arrow[d,"\operatorname{pr}_1"] &
    G \arrow[d,"p"] \\
    S \arrow[r,"\varepsilon"'] &
    G \arrow[r,"p"'] &
    G''
  \end{tikzcd}}
\end{center}
Indeed, \(\varepsilon''\) is an \(S\)-section of \(G''\), hence an
immersion (EGA~I, 5.3.13); by the preceding cartesian diagram the same
is true of \(u\), which completes the proof of \textup{(ii)}.
Moreover, \textup{(vi)} is an immediate consequence of the second
cartesian diagram in \textup{(ii)}.

We prove \textup{\SGARefLink{concl:sga3-VIB-9.2-vii}{(vii)}}. By \textup{\SGARefLink{concl:sga3-VIB-9.2-i}{(i)}}, \(p\) is a covering for
the \textup{(fpqc)} topology. Hence, by \textup{(ii)}, in order to
show that \(p\) has property \(\mathscr P\), it is enough to show
that the first projection
\[
  \operatorname{pr}_1:G\times_SG'\longrightarrow G
\]
has property \(\mathscr P\). This follows from the stability of
\(\mathscr P\) under base change, since
\(\operatorname{pr}_1\) is obtained from \(\pi'\) by base change.

It is clear that \textup{(viii)} follows from \textup{(vii)}, since
\[
  \pi=\pi''\circ p,
\]
where \(\pi'':G''\to S\) denotes the structural morphism.

We prove \textup{\SGARefLink{concl:sga3-VIB-9.2-ix}{(ix)}}. By \textup{\SGARefLink{concl:sga3-VIB-9.2-i}{(i)}}, \(p\) is an epimorphism.
Since \(G\) is reduced, \(p\) factors through the immersion
\[
  G''_{\mathrm{red}}\longrightarrow G'',
\]
which is therefore also an epimorphism, and hence an isomorphism
(IV, \SGARefLink{sga3:iv:lemma:4:4:4}{4.4.4}).

We prove \textup{(x)}. If \(G''\) is separated over \(S\), then
\(\varepsilon''\) is a closed immersion by EGA~I, 5.4.6. On the
other hand, we saw in \textup{(ii)} that \(\varepsilon''\) is a
closed immersion if and only if \(u\) is. Finally, if \(u\) is a
closed immersion, then so is
\[
  \delta\circ(\operatorname{id}_G\times u):
    G\times_SG'\longrightarrow G\times_SG.
\]
It follows from
\hyperref[lem:sga3-VIB-9.2.1]{Lemma~9.2.1 below} that \(G''\) is
separated over \(S\).

Assertion \textup{(xi)} follows from \textup{(vii)} and
EGA~IV\(_2\), 2.2.13.

Assertion \textup{(xii)} follows from \textup{(vii)},
EGA~IV\(_4\), 17.7.5 and 17.7.7, and the fact that \(p\) is
universally open: for every \(s\in S\), if the underlying space of
\(G_s\) is discrete, then the same is true of the underlying space
of \(G''_s\).

Finally, assertion \textup{(xiii)} follows from \textup{(vii)},
\textup{(viii)}, and EGA~IV\(_4\), 17.7.5.

% END INCLUDED SGA3_English_current_buildable_source_20260823/source/inputs/VI/tex_reference_v2/components/75_expose_VIB_section9_def91_prop92_loop2_r1_en.tex

% BEGIN INCLUDED SGA3_English_current_buildable_source_20260823/source/inputs/VI/tex_reference_v2/components/76_expose_VIB_lemma921_rem93_loop2_r1_en.tex
% SGA 3, Expose VI_B, Loop 1 English reconstruction.
% Scope: Lemma 9.2.1 with complete proof, Remark 9.2.2, and
% Remark 9.3(a)--(b), including cases 1--3, editor's notes 86--89,
% and the source's starred warning.
% Authority: Polo--Gille Exp6B-13oct24.pdf, component-local pages
% 67--68 (combined-reader pages 408--409; re-edition headers
% 397--398; original-margin printed pages 383--385).
% Section 10 begins on local page 68 and is excluded.
% OCR and Reinhold English Markdown are comparison witnesses only.
% Loop 2 replaces the source-bound cartesian-square crop with a
% faithful native LaTeX reconstruction; the replacement is closed.

\medskip
\phantomsection\label{lem:sga3-VIB-9.2.1}%
\noindent\textbf{Lemma 9.2.1.---}
\emph{Let \(X\) be an \(S\)-scheme, and let \(R\) be an
equivalence relation on \(X\) defined by the monomorphism
\(v:R\to X\times_S X\). Suppose that \(R\) is effective.
Then:}
\begin{enumerate}
\item[\textup{(i)}]\phantomsection\label{concl:sga3-VIB-9.2.1-i}
  \emph{The morphism \(v\) is an immersion; it is a closed
  immersion if \(Y=X/R\) is separated.}

\item[\textup{(ii)}]\phantomsection\label{concl:sga3-VIB-9.2.1-ii}
  \emph{Suppose, in addition, that \(Y\) represents the
  \textup{(fpqc)} quotient sheaf of \(X\) by \(R\),}%
  \footnote{\label{note:sga3-VIB-86}Editor's note: as
  O.\ Gabber pointed out, this condition is used in the proof
  and must be included among the hypotheses.}
  \emph{and that \(v\) is a closed immersion. Then
  \(Y=X/R\) is separated over \(S\).}
\end{enumerate}

Recall (IV, \SGARefLink{sga3:iv:definition:3:3:2}{Definition~3.3.2}) that the hypothesis ``\(R\) is
effective'' means that there is a morphism of \(S\)-schemes
\[
  p:X\longrightarrow Y
\]
such that the natural morphism
\[
  R\longrightarrow X\times_Y X
\]
is an isomorphism. It follows (EGA~I, 5.3.5) that the following
diagram is cartesian:
\begin{center}
  \phantomsection\label{diag:sga3-VIB-9.2.1-effective-square}%
  \begin{tikzcd}[row sep=large,column sep=large]
    R \arrow[r,"v"] \arrow[d] &
    X\times_SX \arrow[d,"p\times p"] \\
    Y \arrow[r,"\Delta_{Y/S}"'] &
    Y\times_SY
  \end{tikzcd}
\end{center}
Since \(\Delta_{Y/S}\) is an immersion (EGA~I, 5.3.9), the same
is true of \(v\). Likewise, if \(Y\) is separated over \(S\),
then \(\Delta_{Y/S}\) is a closed immersion, and hence so is
\(v\).

Conversely, suppose that \(v\) is a closed immersion and that
\(Y\) represents the \textup{(fpqc)} quotient sheaf of \(X\)
by \(R\). Then \(p\) is a covering for the \textup{(fpqc)}
topology (IV, \SGARefLink{sga3:iv:proposition:4:4:3}{4.4.3}), and therefore \(p\times p\) is also a
covering: by base change, \(p\times\operatorname{id}_X\) and
\(\operatorname{id}_Y\times p\) are coverings, and so is their
composite \(p\times p\). Thus, by \textup{(fpqc)} descent
(cf.\ EGA~IV\(_2\), 2.7.1), \(\Delta_{Y/S}\) is a closed
immersion; in other words, \(Y\) is separated over \(S\).

\medskip
\phantomsection\label{rem:sga3-VIB-9.2.2}%
\noindent\textbf{Remark 9.2.2.---}
\emph{Under the general hypotheses of
\hyperref[prop:sga3-VIB-9.2]{9.2}, if \(G'\) is flat and
locally of finite presentation over \(S\), then \(p\) is a
covering for the \textup{(fppf)} topology,}%
\footnote{\label{note:sga3-VIB-87}Editor's note: \textup{(fpqc)}
has been corrected to \textup{(fppf)}.}
\emph{by \hyperref[concl:sga3-VIB-9.2-vii]{9.2\textup{(vii)}}.
Consequently, assertions
\hyperref[concl:sga3-VIB-9.2-vii]{9.2\textup{(vii)}} and
\hyperref[concl:sga3-VIB-9.2-viii]{9.2\textup{(viii)}} extend
to properties \(\mathscr P\) that are local for the
\textup{(fppf)} topology.}

\medskip
\phantomsection\label{rem:sga3-VIB-9.3}%
\noindent\textbf{Remark 9.3.---}
\phantomsection\label{rem:sga3-VIB-9.3-a}%
\textup{a)} The question whether a quotient \(G/G'\) is
representable is often delicate; in this seminar we prove the
representability of certain particular quotients.

In general, in order to conclude that \(G/G'\) is representable,
it is not enough to suppose that \(G\) and \(G'\) are of finite
presentation over \(S\) and that \(G'\) is flat over \(S\).
Indeed, suppose in addition that \(G\) is smooth with connected
fibers. In that case, if \(G/G'\) is a scheme, it is separated
by \hyperref[cor:sga3-VIB-5.4]{Corollary~5.4}; therefore
\(G'\hookrightarrow G\) is a closed immersion by
\hyperref[concl:sga3-VIB-9.2-x]{9.2\textup{(x)}}.
Consequently, if \(G'\) is not closed in \(G\), then \(G/G'\)
is not representable.

To obtain such a counterexample, take \(S\) to be the spectrum
of a discrete valuation ring and put \(G=\mathbb G_{m,S}\).
Let \(n>1\) be an integer invertible on \(S\). Then
\phantomsection\label{form:sga3-VIB-9.3-mu-n}%
\[
  \mu_n=\ker\bigl(G\xrightarrow{\ n\ }G\bigr)
\]
is a closed subgroup of \(G\), \'etale over \(S\)
(cf.\ Expos\'e~VII\(_A\)% 
\footnote{\label{note:sga3-VIB-88}Editor's note: see
Expos\'e~VII\(_A\), \SGARefLink{prop:sga3-VIB-8.4}{8.4}, or \SGARefLink{sga3:viii:proposition:2:1}{Expos\'e~VIII, 2.1}.}).
Let \(G'\) be the open subgroup of \(\mu_n\) obtained by
removing from \(\mu_n\) the closed subset of its closed fiber
that is complementary to the origin. Then \(G'\) is not closed
in \(G\), and hence \(G/G'\) is not representable. One can also
construct such examples in which \(G'\) is smooth with connected
fibers.

\phantomsection\label{rem:sga3-VIB-9.3-b}%
\textup{b)} It remains possible, however, that \(G/G'\) is
representable when \(G\) and \(G'\) are of finite presentation
over \(S\), and \(G'\) is flat over \(S\) and closed in \(G\)%
\begingroup
\renewcommand{\thefootnote}{\fnsymbol{footnote}}%
\footnote[1]{\label{note:sga3-VIB-9.3-star}This is too
optimistic, as M.\ Raynaud shows in his thesis
\textup{(loc.\ cit., X, 14)}.}%
\endgroup\hspace{0.15em}%
\footnote{\label{note:sga3-VIB-89}Editor's note: the starred
remark refers to counterexample X, 14 in \textup{[Ray70a]}.
The base there is a regular local scheme of dimension \(2\).}.
Under these hypotheses, \(G/G'\) is known to be representable
in the following particular cases:
\begin{enumerate}
\item[\textup{\(1^\circ\)---}]\phantomsection\label{case:sga3-VIB-9.3-b-1}
  \(S\) is the spectrum of an Artinian ring
  (cf.\ \hyperref[thm:sga3-VIA-3.2]{VI\(_A\), 3.2} and
  \hyperref[thm:sga3-VIA-3.3.2]{3.3.2}).

\item[\textup{\(2^\circ\)---}]\phantomsection\label{case:sga3-VIB-9.3-b-2}
  \(G'\) is proper over \(S\), and \(G\) is quasi-projective
  over \(S\) (cf.\ \hyperref[thm:sga3-V-7.1]{V, 7.1}).

\item[\textup{\(3^\circ\)---}]\phantomsection\label{case:sga3-VIB-9.3-b-3}
  \(S\) is locally Noetherian of dimension \(1\)
  (cf.\ \textup{[An73]}, Theorem~4.C).
\end{enumerate}

% END INCLUDED SGA3_English_current_buildable_source_20260823/source/inputs/VI/tex_reference_v2/components/76_expose_VIB_lemma921_rem93_loop2_r1_en.tex

% BEGIN INCLUDED SGA3_English_current_buildable_source_20260823/source/inputs/VI/tex_reference_v2/components/77_expose_VIB_section10_through_lemma102_en.tex
% SGA 3, Expose VI_B, Loop 1 English reconstruction.
% Scope: Section 10 heading, paragraphs 10.0 and 10.1, and
% Lemma 10.2 in full, including editor's notes 90--92.
% Authority: Polo--Gille Exp6B-13oct24.pdf, component-local pages
% 68--70 (combined-reader pages 409--411; re-edition headers
% 398--400; original-margin printed pages 385--387).
% Lemma 10.3 begins on local page 70 and is excluded.
% OCR and Reinhold English Markdown are comparison witnesses only.
% This unit contains no diagrams.

\section*{10. Passage to the projective limit in group schemes
  and schemes with an operator group}
\phantomsection\label{sec:sga3-VIB-10}%
\addcontentsline{toc}{section}{10. Passage to the projective limit
  in group schemes and schemes with an operator group}

\medskip
\phantomsection\label{par:sga3-VIB-10.0}%
\noindent\textbf{10.0.---}
Recall the essential result of EGA~IV\(_3\), \S 8.8.
Suppose that the following data are given: a quasi-compact and
quasi-separated scheme \(S_0\), an upward filtered preordered set
\(I\), a direct system
\((\mathcal A_i)_{i\in I}\) of quasi-coherent commutative
\(\mathcal O_{S_0}\)-algebras,
\[
  \mathcal A=\varinjlim_i\mathcal A_i,\qquad
  S_i=\operatorname{Spec}\mathcal A_i\quad(i\in I),
  \qquad
  S=\operatorname{Spec}\mathcal A\footnotemark.
\]
\footnotetext{\label{note:sga3-VIB-90}Editor's note: since \(S\)
is affine over \(S_0\), it is quasi-compact and
quasi-separated; cf.\ Editor's \hyperref[note:sga3-VIB-92]{note~92} below.}
Then the category of \(S\)-schemes of finite presentation is
determined, up to equivalence, by the categories of
\(S_i\)-schemes of finite presentation, the functors between
these categories
\[
  \rho_{ji}:X_i\longmapsto X_i\times_{S_i}S_j
  \qquad(i\leq j),
\]
and the transitivity isomorphisms
\[
  \rho_{kj}\circ\rho_{ji}\xrightarrow{\ \sim\ }\rho_{ki}.
\]

More precisely, given \(j\in I\) and an \(S_j\)-scheme
\(X_j\) of finite presentation, put, for every \(i\in I\)
with \(i\geq j\),
\[
  X_i=X_j\times_{S_j}S_i,
  \qquad
  X=X_j\times_{S_j}S.
\]
Then EGA~IV\(_3\), 8.8.2 gives:
\begin{enumerate}
\item[\textup{(i)}]\phantomsection\label{concl:sga3-VIB-10.0-i}
  Given \(j\in I\) and two \(S_j\)-schemes \(X_j\) and \(Y_j\)
  of finite presentation, the canonical map
  \phantomsection\label{form:sga3-VIB-10.0-Hom-limit}%
  \[
    \varinjlim_{i\geq j}
      \operatorname{Hom}_{S_i}(X_i,Y_i)
      \longrightarrow
      \operatorname{Hom}_S(X,Y)
  \]
  is bijective.

\item[\textup{(ii)}]\phantomsection\label{concl:sga3-VIB-10.0-ii}
  For every \(S\)-scheme \(X\) of finite presentation, there
  are an index \(j\in I\), an \(S_j\)-scheme \(X_j\) of finite
  presentation, and an \(S\)-isomorphism
  \[
    X\xrightarrow{\ \sim\ }X_j\times_{S_j}S.
  \]
\end{enumerate}

It follows from EGA~IV\(_3\), 8.8.3 that whenever a diagram
\(\mathcal D\) involves finitely many objects and arrows in
the category of \(S\)-schemes of finite presentation, one can
find an index \(i\in I\) and a diagram \(\mathcal D_i\) in the
category of \(S_i\)-schemes of finite presentation such that
\(\mathcal D\), up to isomorphism, is obtained from
\(\mathcal D_i\) by the base change \(S\to S_i\). One can even
find \(i\) and \(\mathcal D_i\) such that every cartesian square
in \(\mathcal D\) is obtained from a cartesian square in
\(\mathcal D_i\).

\medskip
\phantomsection\label{par:sga3-VIB-10.1}%
\noindent\textbf{10.1.---}
Moreover, many familiar properties of a morphism that are
stable under base change have the following property:
\begin{quote}
\emph{Let \(u:X\to Y\) be an \(S\)-morphism between
\(S\)-schemes of finite presentation. By
\hyperref[par:sga3-VIB-10.0]{10.0}, it is obtained by base
change from an \(S_j\)-morphism \(u_j:X_j\to Y_j\) between
\(S_j\)-schemes of finite presentation. Then \(u\) has the
property \(\mathcal P\) if and only if there is an \(i\geq j\)
such that}
\[
  u_i=u_j\times_{S_j}S_i
\]
\emph{has the property \(\mathcal P\).}
\end{quote}

This applies when \(\mathcal P\) is any of the following
properties of a morphism: separated, surjective, radicial,
affine, quasi-affine, finite, quasi-finite, proper,
projective, quasi-projective, an isomorphism, a monomorphism,
an immersion, an open immersion, a closed immersion
\textup{(EGA~IV\(_3\), 8.10.5)}, flat
\textup{(EGA~IV\(_3\), 11.2.6)}, smooth, unramified, or
\'etale \textup{(EGA~IV\(_4\), 17.7.8)}.%
\footnote{\label{note:sga3-VIB-91}Editor's note: the word
``flat'' has been added, and the reference 17.7.6 has been
corrected to 17.7.8.}

This also applies when \(\mathcal P\) is the property of being
a covering for the \textup{(fppf)} topology. Indeed, given
two \(S\)-schemes \(X\) and \(Y\) of finite presentation and
an \(S\)-morphism \(u:X\to Y\), it follows from
\hyperref[sga3:iv:proposition:6:3:1:item:i:01]%
  {\SGARefLink{sga3:iv:proposition:6:3:1:item:i:01}{Expos\'e~IV, 6.3.1\textup{(i)}}}%
\footnote{\label{note:sga3-VIB-92}Editor's note: one also uses
the fact that \(Y\), being of finite presentation over \(S\),
is quasi-compact.}
that \(u\) is a covering for the \textup{(fppf)} topology if
and only if there are an \(S\)-scheme \(Z\) and a faithfully
flat \(S\)-morphism \(v:Z\to Y\) of finite presentation that
factors through \(u\).

The purpose of this Section~\hyperref[sec:sga3-VIB-10]{10} is
to give variants of results of this kind for the category of
\(S\)-groups of finite presentation, for the category of
\(S\)-schemes with an operator group, and for certain
properties of group monomorphisms \textup{(being invariant,
being central with representable quotient sheaf, and so on)}.

The first two preliminary results of this kind follow. In
\hyperref[lem:sga3-VIB-10.2]{10.2} through
\hyperref[lem:sga3-VIB-10.9]{10.9}, we retain the notation
introduced in \hyperref[par:sga3-VIB-10.0]{10.0}.

\medskip
\phantomsection\label{lem:sga3-VIB-10.2}%
\noindent\textbf{Lemma 10.2.---}
\emph{Let \(G_j\) and \(H_j\) be two \(S_j\)-groups of finite
presentation. For every \(i\geq j\), put
\[
  G_i=G_j\times_{S_j}S_i,
  \qquad
  G=G_j\times_{S_j}S,
\]
and define \(H_i\) and \(H\) similarly. Then the canonical
map}
\phantomsection\label{form:sga3-VIB-10.2-Hom-limit}%
\[
  \varinjlim_{i\geq j}
    \operatorname{Hom}_{S_i\text{-}\mathrm{gr}}(G_i,H_i)
    \longrightarrow
    \operatorname{Hom}_{S\text{-}\mathrm{gr}}(G,H)
\]
\emph{is bijective.}

% END INCLUDED SGA3_English_current_buildable_source_20260823/source/inputs/VI/tex_reference_v2/components/77_expose_VIB_section10_through_lemma102_en.tex

% BEGIN INCLUDED SGA3_English_current_buildable_source_20260823/source/inputs/VI/tex_reference_v2/components/78_expose_VIB_lemma103_through_lemma108_loop2_r1_en.tex
% SGA 3, Expose VI_B, Loop 1 English reconstruction.
% Scope: Lemma 10.3 through Lemma 10.8, all intervening proofs,
% Corollary 10.5, Proposition 10.6, and editor's notes 93--94.
% Authority: Polo--Gille Exp6B-13oct24.pdf, component-local pages
% 70--71 (combined-reader pages 411--412; re-edition headers
% 400--401; original-margin printed pages 387--390).
% Lemma 10.9 begins on local page 71 and is excluded.
% Editor's note 95 belongs to 10.10 after the excluded Lemma 10.9
% and is therefore outside this unit.
% OCR and Reinhold English Markdown are comparison witnesses only.
% Loop 2 replaces the source-bound crop for Diagram (D) with a
% faithful native LaTeX reconstruction; the replacement is closed.

\medskip
\phantomsection\label{lem:sga3-VIB-10.3}%
\noindent\textbf{Lemma 10.3.---}
\emph{Let \(G\) be an \(S\)-group of finite presentation.
Then there are an index \(j\in I\), an \(S_j\)-group \(G_j\)
of finite presentation, and an isomorphism of \(S\)-groups}
\[
  G\xrightarrow{\ \sim\ }G_j\times_{S_j}S.
\]

Assertions \hyperref[lem:sga3-VIB-10.2]{10.2} and
\hyperref[lem:sga3-VIB-10.3]{10.3} are easy consequences of
\hyperref[par:sga3-VIB-10.0]{10.0} and
\hyperref[par:sga3-VIB-10.1]{10.1}, taking into account the
interpretation%
\footnote{\label{note:sga3-VIB-93}Editor's note: in terms of
commutative diagrams of \(S\)-morphisms.}
of the structure of an \(S\)-group given in EGA~0\(_{\mathrm{III}}\),
8.2.5 and 8.2.6.

\medskip
\phantomsection\label{lem:sga3-VIB-10.4}%
\noindent\textbf{Lemma 10.4.---}
\emph{Let \(u:G\to H\) be a morphism of \(S\)-groups between
\(S\)-groups of finite presentation. By
\hyperref[lem:sga3-VIB-10.3]{10.3} and
\hyperref[lem:sga3-VIB-10.2]{10.2}, \(u\) is obtained by base
change from a morphism \(u_j:G_j\to H_j\) between
\(S_j\)-groups of finite presentation. Then \(u\) is a central
monomorphism \textup{(respectively, an invariant monomorphism)}
if and only if there is an \(i\geq j\) such that}
\phantomsection\label{form:sga3-VIB-10.4-base-change}%
\[
  u_i=u_j\times_{S_j}S_i
\]
\emph{has the same property.}

This follows immediately from
\hyperref[par:sga3-VIB-10.0]{10.0} and
\hyperref[par:sga3-VIB-10.1]{10.1}, in view of the
characterization in \hyperref[rem:sga3-VIB-6.7]{6.7} of
central or invariant group monomorphisms.

\medskip
\phantomsection\label{cor:sga3-VIB-10.5}%
\noindent\textbf{Corollary 10.5.---}
\emph{Let \(G_j\) be an \(S_j\)-group of finite presentation.
Then \(G_j\times_{S_j}S\) is commutative if and only if there
is an \(i\geq j\) such that \(G_j\times_{S_j}S_i\) is
commutative.}

Indeed, to say that an \(S\)-group is commutative is the same
as saying that it is central when regarded as a group
subscheme of itself.

\medskip
\phantomsection\label{prop:sga3-VIB-10.6}%
\noindent\textbf{Proposition 10.6.---}
\emph{Let \(G_j\) be an \(S_j\)-group of finite presentation,
and let \(G'_j\) be a group subscheme of \(G_j\), flat and of
finite presentation over \(S_j\). Then the quotient}
\phantomsection\label{form:sga3-VIB-10.6-quotient-over-S}%
\[
  (G_j\times_{S_j}S)/(G'_j\times_{S_j}S)
\]
\emph{is representable for the \textup{(fpqc)} topology if
and only if there is an \(i\geq j\) such that}
\phantomsection\label{form:sga3-VIB-10.6-quotient-over-Si}%
\[
  (G_j\times_{S_j}S_i)/(G'_j\times_{S_j}S_i)
\]
\emph{is representable.}

This is a consequence of the following more general lemma.

\medskip
\phantomsection\label{lem:sga3-VIB-10.7}%
\noindent\textbf{Lemma 10.7.---}
\emph{Let \(X_j\) be an \(S_j\)-scheme of finite presentation,
and let \(R_j\) be an equivalence relation on \(X_j\), flat and
of finite presentation.}%
\footnote{\label{note:sga3-VIB-94}Editor's note: that is, such
that the composite
\[
  R_j\lhook\joinrel\longrightarrow X_j\times_{S_j}X_j
  \xrightarrow{\ \operatorname{pr}_1\ }X_j
\]
is flat and of finite presentation.}
\emph{Then the quotient sheaf}
\[
  (X_j\times_{S_j}S)/(R_j\times_{S_j}S)
\]
\emph{for the topology \(\mathcal T=\textup{(fppf)}\) or
\(\textup{(fpqc)}\) is representable if and only if there is
an \(i\geq j\) such that the quotient sheaf}
\[
  (X_j\times_{S_j}S_i)/(R_j\times_{S_j}S_i)
\]
\emph{for the topology \(\mathcal T\) is representable.}

In view of the statements EGA~IV\(_3\),\footnote{Canon source note
(SGA-SRC-000861): both the current re-edition and the official original
print IV\(_2\), but all four cited sections belong to EGA~IV\(_3\).}
8.8.2, 8.8.3,
8.10.5, and 11.2.6 recalled in
\hyperref[par:sga3-VIB-10.0]{10.0}, this lemma follows from
the next result.

\medskip
\phantomsection\label{lem:sga3-VIB-10.8}%
\noindent\textbf{Lemma 10.8.---}
\emph{Let \(\mathcal T\) be the \textup{(fppf)} or
\textup{(fpqc)} topology. Let \(X\) be an \(S\)-scheme of
finite presentation \textup{(respectively, locally of finite
presentation)}, and let \(R\) be an equivalence relation on
\(X\) defined by a monomorphism}
\[
  v:R\longrightarrow X\times_S X
\]
\emph{such that \(\operatorname{pr}_1\circ v\) is flat and of
finite presentation \textup{(respectively, flat and locally
of finite presentation)}. Then the following conditions are
equivalent:}
\begin{enumerate}
\item[\textup{(i)}]\phantomsection\label{concl:sga3-VIB-10.8-i}
  \emph{The quotient sheaf \(X/R\) for the topology
  \(\mathcal T\) is representable.}

\item[\textup{(ii)}]\phantomsection\label{concl:sga3-VIB-10.8-ii}
  \emph{There are an \(S\)-scheme \(Y\) of finite presentation
  \textup{(respectively, locally of finite presentation)} and
  a faithfully flat morphism \(p:X\to Y\) such that the
  diagram}
\end{enumerate}
\begin{center}
  \phantomsection\label{diag:sga3-VIB-10.8-D}%
  \begin{tabular}{c@{\qquad}c}
    \textup{(D)} &
    \begin{tikzcd}[row sep=large,column sep=large]
      R
        \arrow[r,"{\operatorname{pr}_1\circ v}"]
        \arrow[d,"{\operatorname{pr}_2\circ v}"'] &
      X \arrow[d,"p"] \\
      X \arrow[r,"p"'] &
      Y
    \end{tikzcd}
  \end{tabular}
\end{center}
\emph{is cartesian.}

First note that, by IV, \SGARefLink{sga3:iv:definition:3:3:2}{3.3.2} and 4.4.3, the sheaf \(X/R\) for
the topology \(\mathcal T\) is representable by \(Y\) if and
only if diagram~\hyperref[diag:sga3-VIB-10.8-D]{\textup{(D)}}
is cartesian and \(p\) is a covering for the topology
\(\mathcal T\).

Let us prove that
\hyperref[concl:sga3-VIB-10.8-i]{\textup{(i)}} implies
\hyperref[concl:sga3-VIB-10.8-ii]{\textup{(ii)}}.
Hypothesis~\hyperref[concl:sga3-VIB-10.8-i]{\textup{(i)}}
implies that diagram~\hyperref[diag:sga3-VIB-10.8-D]%
{\SGARefLink{diag:sga3-VIB-10.8-D}{\textup{(D)}}} is cartesian. Thus
\(\operatorname{pr}_1\circ v\) is obtained from \(p\) by base
change along \(p\), and \(p\) is a covering for the
\textup{(fpqc)} topology. Consequently, by \textup{(fpqc)}
descent (EGA~IV\(_2\), 2.7.1), since
\(\operatorname{pr}_1\circ v\) is faithfully flat and
\textup{(locally)} of finite presentation, the same is true
of \(p\). Then EGA~IV\(_4\), 17.7.5 shows that, since \(X\) is
\textup{(locally)} of finite presentation over \(S\), so is
\(Y\).

Let us prove that
\hyperref[concl:sga3-VIB-10.8-ii]{\textup{(ii)}} implies
\hyperref[concl:sga3-VIB-10.8-i]{\textup{(i)}}. It is enough
to show that \(p\) is a covering for the \textup{(fppf)}
topology. But \(p\) is faithfully flat by hypothesis, and it
is locally of finite presentation because \(X\) and \(Y\) are
locally of finite presentation over \(S\)
\textup{(EGA~IV\(_1\), 1.4.3\textup{(v)})}.

% END INCLUDED SGA3_English_current_buildable_source_20260823/source/inputs/VI/tex_reference_v2/components/78_expose_VIB_lemma103_through_lemma108_loop2_r1_en.tex

% BEGIN INCLUDED SGA3_English_current_buildable_source_20260823/source/inputs/VI/tex_reference_v2/components/79_expose_VIB_lemma109_through_prop1011_en.tex
% SGA 3, Expose VI_B, Loop 1 English reconstruction.
% Scope: Lemma 10.9, paragraph 10.10, and Proposition 10.11
% with its complete proof, including editor's notes 95--97.
% Authority: Polo--Gille Exp6B-13oct24.pdf, component-local
% pages 71--73 (combined-reader pages 412--414; re-edition
% headers 401--403; original-margin printed pages 390--392).
% Corollary 10.11.1 begins on local page 73 and is excluded.
% OCR and Reinhold English Markdown are comparison witnesses only.
% This bounded unit contains no genuine diagrams.

\medskip
\phantomsection\label{lem:sga3-VIB-10.9}%
\noindent\textbf{Lemma 10.9.---}
\emph{Let \(G_j\) be an \(S_j\)-group of finite presentation,
and put
\[
  G=G_j\times_{S_j}S.
\]
Then \(G^0\) is representable if and only if there is an
\(i\geq j\) such that
\[
  (G_i)^0=(G_j\times_{S_j}S_i)^0
\]
is representable.}

The condition is sufficient, since the functor
\(G\longmapsto G^0\) commutes with base change, by
\hyperref[prop:sga3-VIB-3.3]{3.3}.

Conversely, suppose that \(G^0\) is representable. Then, by
\hyperref[prop:sga3-VIB-3.9]{3.9}, \(G^0\) is open in \(G\)
and quasi-compact over \(S\), and hence is of finite
presentation over \(S\), since \(G\) is. By
\hyperref[lem:sga3-VIB-10.3]{10.3} and
\hyperref[par:sga3-VIB-10.1]{10.1}, there are an \(i\geq j\)
and an open group subscheme \(H_i\) of \(G_i\) such that
\[
  H_i\times_{S_i}S=G^0.
\]
The structural morphism \(G^0\to S\) has geometrically
connected fibers
\textup{(\hyperref[prop:sga3-VIA-2.1.1]{VI\(_A\), 2.1.1})}.
Consequently, by EGA~IV\(_3\), 9.3.3 and 9.7.7, after
increasing \(i\) we may suppose that the structural morphism
\(H_i\to S_i\) has geometrically connected fibers.%
\begingroup
\renewcommand{\thefootnote}{\fnsymbol{footnote}}%
\footnote[1]{\label{note:sga3-VIB-10.9-Hi-base}
Translator's note: the French source prints \(H_i\to S\).
Since \(H_i\) is an \(S_i\)-group and the argument spreads out
the connectedness of \(G^0\to S\), the base here is \(S_i\).}%
\endgroup\ 
It follows from
\hyperref[lem:sga3-VIB-3.10.1]{3.10.1} that the underlying
space of \(H_i\) is precisely \((G_i)^0\); hence \(H_i\)
represents \((G_i)^0\).

\medskip
\phantomsection\label{par:sga3-VIB-10.10}%
\noindent\textbf{10.10.---}
Recall two very useful special cases of the situation stated
in \hyperref[par:sga3-VIB-10.0]{10.0}
\textup{(cf.\ EGA~IV\(_3\), 8.1.2\textup{(a)} and
\textup{(c)})}:

\smallskip
\phantomsection\label{par:sga3-VIB-10.10-a}%
\noindent\textup{a)}
Given a point \(x\) of a scheme \(X\), put
\(S_0=\operatorname{Spec}\mathbb Z\), and consider the
decreasing filtered inverse system \((S_i)_{i\in I}\) of
affine open neighborhoods of \(x\). Then
\[
  S=\operatorname{Spec}\mathcal O_{X,x}.
\]
In particular, if \(x\) is the generic point of an integral
scheme \(X\), then
\[
  S=\operatorname{Spec}\kappa(x).
\]

\smallskip
\phantomsection\label{par:sga3-VIB-10.10-b}%
\noindent\textup{b)}
Put \(S_0=\operatorname{Spec}\mathbb Z\), and consider the
family \((\mathcal A_i)_{i\in I}\), preordered by inclusion,
of the finitely generated \(\mathbb Z\)-subalgebras of the
ring of an affine scheme \(S\). Since the
\(\mathcal A_i\) are Noetherian rings, this makes it possible
in many cases to pass from the Noetherian case to the general
case.

We shall now give two results concerning the special case
considered in \hyperref[par:sga3-VIB-10.10-a]{(a)}.%
\footnote{\label{note:sga3-VIB-95}Editor's note: observe that
the proof also uses case \textup{(b)}.}

\medskip
\phantomsection\label{prop:sga3-VIB-10.11}%
\noindent\textbf{Proposition 10.11.---}%
\footnote{\label{note:sga3-VIB-96}Editor's note: the statement
has been simplified, and the case of group subschemes has
been treated separately in
\hyperref[cor:sga3-VIB-10.11.1]{Corollary~10.11.1}.}
\begingroup\itshape
Let \(S\) be an integral scheme with generic point
\(\eta\). Let \(G\) be an \(S\)-group and let \(Y,Z\) be
\(S\)-schemes, all of finite presentation. Let
\[
  u,v,i:Y\longrightarrow G,
  \qquad
  j:Z\longrightarrow G
\]
be morphisms of \(S\)-schemes. Suppose that
\[
  i_\eta:Y_\eta\longrightarrow G_\eta,
  \qquad
  j_\eta:Z_\eta\longrightarrow G_\eta
\]
are closed immersions.

Then there is a nonempty open subset \(S'\) of \(S\) such
that the morphisms
\[
  i':Y'\longrightarrow G',
  \qquad
  j':Z'\longrightarrow G',
\]
obtained by base change, are closed immersions, and such that
the functors
\phantomsection\label{form:sga3-VIB-10.11-scheme-functors}%
\[
\begin{gathered}
  \underline{\operatorname{Transp}}(u',v'),\qquad
  \underline{\operatorname{Transp}}_{G'}
    \bigl(i'(Y'),j'(Z')\bigr),\\
  \underline{\operatorname{Transp}}^{\mathrm{str}}_{G'}
    \bigl(i'(Y'),j'(Z')\bigr)
\end{gathered}
\]
respectively, the functors
\phantomsection\label{form:sga3-VIB-10.11-group-functors}%
\[
  \underline{\operatorname{Centr}}(u'),
  \qquad
  \underline{\operatorname{Norm}}_{G'}i'(Y'),
\]
are representable by closed \(S'\)-subschemes
\textup{(respectively, closed \(S'\)-subgroups)} of \(G'\),
of finite presentation over \(S'\).
\par\endgroup

\emph{Proof.}
We shall apply the results of
\hyperref[par:sga3-VIB-10.1]{10.1}, first in the situation of
\hyperref[par:sga3-VIB-10.10-a]{10.10\textup{(a)}}, and then
in that of
\hyperref[par:sga3-VIB-10.10-b]{10.10\textup{(b)}}.
The schemes \(G_\eta,Y_\eta,Z_\eta\) are flat over the field
\(\kappa(\eta)\); the group \(G_\eta\) is separated over
\(\kappa(\eta)\)
\textup{(\hyperref[sga3:VIA:0.3]{VI\(_A\), 0.3})}; and
\(i_\eta,j_\eta\) are closed immersions. Therefore, by
\hyperref[par:sga3-VIB-10.1]{10.1}, there are an affine open
subset \(S'=\operatorname{Spec}(A')\) of \(S\), a Noetherian
subring \(A\) of \(A'\), \(A\)-schemes
\(G_A,Y_A,Z_A\), flat and of finite presentation over \(A\),
and morphisms
\[
  u_A,v_A,i_A:Y_A\longrightarrow G_A,
  \qquad
  j_A:Z_A\longrightarrow G_A,
\]
such that \(G_A\) is an \(A\)-group separated over \(A\),
\(i_A\) and \(j_A\) are closed immersions, and
\[
  G\times_SS'=G_A\otimes_AA',
\]
with analogous equalities for the remaining schemes and
morphisms. Since the functors under consideration over
\(S'\) are obtained by base change from the analogous
functors over \(\operatorname{Spec}(A)\), it is enough to
prove the result for the latter.

By EGA~IV\(_2\), 6.9.2, after replacing \(A\) by a
localization \(A_f\) \textup{(and hence \(S'\) by the affine
open subset \(S'_f\))}, we may suppose that
\(G_A,Y_A,Z_A\) are essentially free over \(A\), in the sense
of \hyperref[def:sga3-VIB-6.2.1]{6.2.1}.%
\footnote{\label{note:sga3-VIB-97}Editor's note: indeed, since
\(G_A\) is flat and of finite presentation over \(A\), it is
covered by affine open subsets \(G_1,\ldots,G_n\) such that
each \(\mathcal O(G_i)\) is a flat \(A\)-algebra of finite
presentation. By EGA~IV\(_2\), 6.9.2, there is an
\(f_i\in A\) such that
\(\mathcal O(G_i)_{f_i}\) is a free \(A_{f_i}\)-module. One
may therefore replace \(\operatorname{Spec}(A)\) by the
affine open subset \(D(f)\), where
\(f=f_1\cdots f_n\), and proceed in the same way for \(Y_A\)
and \(Z_A\).}
It then follows from
\hyperref[ex:sga3-VIB-6.2.4-b]{6.2.4\textup{(b)}} and
\hyperref[ex:sga3-VIB-6.2.4-e]{6.2.4\textup{(e)}} that,
under the hypotheses of the proposition, the functors under
consideration are representable by closed \(A\)-subschemes
of \(G_A\). These subschemes are of finite presentation over
\(A\), since \(A\) is Noetherian and \(G_A\) is of finite
presentation over \(A\); and in the cases of
\(\underline{\operatorname{Centr}}(u_A)\) and
\(\underline{\operatorname{Norm}}_{G_A}i_A(Y_A)\), they are
closed \(A\)-subgroups of \(G_A\).

% END INCLUDED SGA3_English_current_buildable_source_20260823/source/inputs/VI/tex_reference_v2/components/79_expose_VIB_lemma109_through_prop1011_en.tex

% BEGIN INCLUDED SGA3_English_current_buildable_source_20260823/source/inputs/VI/tex_reference_v2/components/80_expose_VIB_cor10111_through_cor1013_r2_en.tex
% SGA 3, Expose VI_B, Loop 1 English reconstruction.
% Scope: Corollary 10.11.1, Proposition 10.12 with its complete
% proof, Corollary 10.12.1, and Corollary 10.13 with its complete
% proof, including editor's notes 98--101.
% Authority: Polo--Gille Exp6B-13oct24.pdf, component-local pages
% 73--74 (combined-reader pages 414--415; re-edition headers
% 403--404; original-margin printed pages 392--394).
% Paragraph 10.14 begins on local page 74 and is excluded.
% OCR and Reinhold English Markdown are comparison witnesses only.
% This bounded unit contains no genuine diagrams.

\medskip
\phantomsection\label{cor:sga3-VIB-10.11.1}%
\noindent\textbf{Corollary 10.11.1.---}
\begingroup\itshape
Let \(S\) be an integral scheme with generic point \(\eta\).
Let \(G,H,K\) be \(S\)-groups of finite presentation, and let
\[
  i:H\longrightarrow G,
  \qquad
  j:K\longrightarrow G
\]
be two quasi-compact monomorphisms of \(S\)-groups. Then
there is a nonempty open subset \(S'\) of \(S\) such that the
morphisms
\[
  i':H'\longrightarrow G',
  \qquad
  j':K'\longrightarrow G',
\]
obtained by base change, are closed immersions, and such that
the functors
\phantomsection\label{form:sga3-VIB-10.11.1-transporter-functors}%
\[
  \underline{\operatorname{Transp}}_{G'}(H',K'),
  \qquad
  \underline{\operatorname{Transp}}^{\mathrm{str}}_{G'}(H',K'),
\]
respectively, the functors
\phantomsection\label{form:sga3-VIB-10.11.1-centralizer-normalizer}%
\[
  \underline{\operatorname{Centr}}_{G'}H',
  \qquad
  \underline{\operatorname{Norm}}_{G'}H',
\]
are representable by closed \(S'\)-subschemes
\textup{(respectively, closed \(S'\)-subgroups)} of \(G'\),
of finite presentation over \(S'\).
\par\endgroup

This follows from
\hyperref[prop:sga3-VIB-10.11]{the preceding proposition},
since the hypotheses imply, by
\hyperref[cor:sga3-VIB-1.4.2]{1.4.2}, that \(i_\eta\) and
\(j_\eta\) are closed immersions.

\medskip
\phantomsection\label{prop:sga3-VIB-10.12}%
\noindent\textbf{Proposition 10.12.---}
\begingroup\itshape
Let \(S\) be an integral scheme, let \(G\) be an \(S\)-group
of finite presentation, and let \(A\) and \(B\) be two group
subschemes of \(G\), of finite presentation over \(S\) and
with smooth generic fibers. Suppose moreover that one of the
following conditions holds:
\begin{enumerate}
\item[\textup{a)}]\phantomsection\label{cond:sga3-VIB-10.12-a}
  \(A\) has connected generic fiber;
\item[\textup{b)}]\phantomsection\label{cond:sga3-VIB-10.12-b}
  \(A\) and \(B\) are invariant in \(G\).
\end{enumerate}
Then there are a nonempty open subset \(S'\) of \(S\) and a
closed group subscheme \(D'\) of
\[
  G'=G|_{S'},
\]
of finite presentation over \(S'\) and with smooth fibers,
that represents the \textup{(fppf)} sheaf associated with
the presheaf of groups of commutators of
\[
  A'=A|_{S'}
  \qquad\text{and}\qquad
  B'=B|_{S'}
\]
in \(G'\). Moreover, \(D'\) has connected fibers in case
\textup{(a)}, and is invariant in \(G'\) in case \textup{(b)}.

In particular, for every \(s\in S'\),
\[
  D'_s=(A_s,B_s),
\]
with the notation of \hyperref[def:sga3-VIB-7.2.2]{7.2.2}.
\par\endgroup

\emph{Proof.}
Let \(\eta\) be the generic point of \(S\), and put
\[
  H_\eta=(A_\eta,B_\eta).
\]
Since \(A_\eta\) and \(B_\eta\) are smooth,
\hyperref[prop:sga3-VIB-7.8]{7.8} in case \textup{(a)} and
\hyperref[concl:sga3-VIB-7.3-v]{7.3\textup{(v)}} in case
\textup{(b)} show that \(H_\eta\) is connected,
respectively invariant in \(G_\eta\).

We are in the situation of
\hyperref[par:sga3-VIB-10.0]{10.0} corresponding to
\hyperref[par:sga3-VIB-10.10-a]{10.10\textup{(a)}}.
Consequently, by \hyperref[lem:sga3-VIB-10.3]{10.3} and
\hyperref[par:sga3-VIB-10.1]{10.1}, there are a nonempty open
subset \(S'\) of \(S\) and an \(S'\)-group subscheme \(D'\),
of finite presentation and closed in \(G'\), such that
\[
  D'_\eta=D'\otimes_{S'}\kappa(\eta)=H_\eta.
\]
Moreover, by EGA~IV\(_3\), 9.7.7 and 9.3.3, we may suppose
that \(D'\) has geometrically reduced fibers. In case
\textup{(a)}, another application of EGA~IV\(_3\), 9.7.7 and
9.3.3 allows us to suppose that \(D'\) has connected fibers,
hence geometrically connected fibers
\textup{(cf.\ \hyperref[prop:sga3-VIA-2.1.1]%
{\SGARefLink{prop:sga3-VIA-2.1.1}{VI\(_A\), 2.1.1}})}. In case \textup{(b)}, we may suppose,
by \hyperref[lem:sga3-VIB-10.4]{10.4}, that \(D'\) is
invariant in \(G'\).

We also saw, in the course of the proof of
\hyperref[prop:sga3-VIB-7.8]{7.8}, that there is an integer
\(n\) such that
\phantomsection\label{form:sga3-VIB-10.12-generic-image}%
\[
  \nu_\eta^n
  \left((A_\eta\times_{\kappa(\eta)}B_\eta)^n\right)
  =D'_\eta,
\]
where \(\nu_\eta\) and \(\nu_\eta^n\) are defined as in
\hyperref[def:sga3-VIB-7.2.2-a]{7.2.2\textup{(a)}} and
\hyperref[concl:sga3-VIB-7.1-ii]{7.1\textup{(ii)}}.
The same formulas define morphisms
\phantomsection\label{form:sga3-VIB-10.12-nu-prime}%
\[
  \nu':A'\times_{S'}B'\longrightarrow G',
  \qquad
  {\nu'}^n:(A'\times_{S'}B')^n\longrightarrow G',
\]
and
\[
  {\nu'}^n\otimes_{S'}\kappa(\eta)=\nu_\eta^n.
\]

It follows from \hyperref[par:sga3-VIB-10.1]{10.1} that we
may choose \(S'\) so that the morphism \({\nu'}^n\) is flat
and factors through \(D'\), and so that the resulting
morphism
\[
  {\nu'}^n:(A'\times_{S'}B')^n\longrightarrow D'
\]
is surjective. By \hyperref[lem:sga3-VIB-7.5]{7.5}, the
morphism%
\footnote{\label{note:sga3-VIB-98}Editor's note: \(n+1\) has
been corrected below to \(2n\).}
\phantomsection\label{form:sga3-VIB-10.12-covering-map}%
\[
  {\nu'}^{2n}
  =\mu\circ\bigl({\nu'}^n\times_{S'}{\nu'}^n\bigr):
  (A'\times_{S'}B')^{2n}\longrightarrow D'
\]
is a covering for the \textup{(fppf)} topology. Therefore,
by \hyperref[prop:sga3-VIB-7.6]{7.6}, \(D'\) represents the
\textup{(fppf)} sheaf associated with the presheaf of
commutators of \(A'\) and \(B'\) in \(G'\).

Furthermore, for every \(s\in S'\), \({\nu'}^n\) induces a
surjective morphism
\[
  \nu_s^n:
  (A_s\times_{\kappa(s)}B_s)^n\longrightarrow D'_s.
\]%
\footnote{\label{note:sga3-VIB-99}Editor's note: the argument
that follows has been expanded from the original.}
Thus \(D'_s\) is a closed subgroup of \(G_s\) containing
\[
  \nu_s(A_s\times_{\kappa(s)}B_s),
\]
and hence containing \((A_s,B_s)\). Since \(\nu_s^n\) is
surjective, \(D'_s=(A_s,B_s)\); and by
\hyperref[prop:sga3-VIB-7.6]{7.6}, it represents the
\textup{(fppf)} sheaf of commutators of \(A_s\) and \(B_s\)
in \(G_s\).

\medskip
\phantomsection\label{cor:sga3-VIB-10.12.1}%
\noindent\textbf{Corollary 10.12.1.---}%
\footnote{\label{note:sga3-VIB-100}Editor's note: this
corollary, which is used in the proof of
\hyperref[prop:sga3-VIB-8.4]{8.4}, has been added.}
\begingroup\itshape
Let \(S\) be an integral scheme with generic point \(\eta\),
and let \(G\) be an \(S\)-group of finite presentation with
smooth fibers. Put
\[
  K_\eta^0=G_\eta,
  \qquad
  K_\eta^p=(K_\eta^{p-1},K_\eta^{p-1})
  \quad
  \textup{(respectively, }
    K_\eta^p=(G_\eta,K_\eta^{p-1})\textup{)}
\]
for every \(p\in\mathbb N^*\).%
\begingroup
\renewcommand{\thefootnote}{\fnsymbol{footnote}}%
\footnote[1]{\label{note:sga3-VIB-10.12.1-Geta}
Translator's note: the French re-edition prints \(G\) in
the second recursion. Since \(K_\eta^{p-1}\) is a
\(\kappa(\eta)\)-group subscheme of \(G_\eta\), the
commutator must be formed with \(G_\eta\).}%
\endgroup
Fix \(n\in\mathbb N^*\).
Then there are a nonempty open subset \(S'\) of \(S\) and a
group subscheme \(D\), invariant in \(G|_{S'}\), of finite
presentation and with smooth fibers, such that
\[
  D_s=K_s^n
\]
for every \(s\in S'\).
\par\endgroup

This follows from \hyperref[prop:sga3-VIB-10.12]{10.12}, by
induction on \(n\).

\medskip
\phantomsection\label{cor:sga3-VIB-10.13}%
\noindent\textbf{Corollary 10.13.---}
\begingroup\itshape
Let \(S\) be an integral scheme with generic point \(\eta\).
Let \(G\) be an \(S\)-group and let \(H\) be an invariant
group subscheme of \(G\). Suppose that \(G\) and \(H\) are of
finite presentation over \(S\) and have smooth generic
fibers.%
\footnote{\label{note:sga3-VIB-101}Editor's note: the
original assumed that \(G\) and \(H\) were of finite
presentation and that \(H\) had smooth fibers; the hypothesis
has been modified in order to apply
\hyperref[prop:sga3-VIB-10.12]{10.12}. The proof has also
been expanded.}

If
\[
  (H_\eta,H_\eta)=H_\eta
  \qquad
  \textup{(respectively, }
    (G_\eta,H_\eta)=H_\eta\textup{)},
\]
then there is a nonempty open subset \(S'\) of \(S\) such
that, for every \(s\in S'\),
\[
  (H_s,H_s)=H_s
  \qquad
  \textup{(respectively, }
    (G_s,H_s)=H_s\textup{)}.
\]
\par\endgroup

Indeed, the proof of \hyperref[prop:sga3-VIB-10.12]{10.12}
gives a nonempty open subset \(S'\) of \(S\) and an
\(S'\)-group subscheme \(D\) of \(G|_{S'}\), of finite
presentation and with smooth fibers, such that, for every
\(s\in S'\),
\[
  D_s=(H_s,H_s)
  \qquad
  \textup{(respectively, }
    D_s=(G_s,H_s)\textup{)}.
\]
On the other hand, since \(D_\eta=H_\eta\) and \(D\) and
\(H\) are of finite presentation over \(S'\), EGA~IV\(_3\),
8.8.2.5 gives a nonempty open subset \(S''\) of \(S'\) such
that
\[
  D|_{S''}=H|_{S''}.
\]
Consequently, for every \(s\in S''\),
\[
  H_s=D_s=(H_s,H_s)
  \qquad
  \textup{(respectively, }
    H_s=D_s=(G_s,H_s)\textup{)}.
\]

% END INCLUDED SGA3_English_current_buildable_source_20260823/source/inputs/VI/tex_reference_v2/components/80_expose_VIB_cor10111_through_cor1013_r2_en.tex

% BEGIN INCLUDED SGA3_English_current_buildable_source_20260823/source/inputs/VI/tex_reference_v2/components/81_expose_VIB_par1014_through_prop1016_en.tex
% SGA 3, Expose VI_B, Loop 1 English reconstruction.
% Scope: paragraph 10.14, Definition 10.15, Proposition 10.16,
% and the complete proof of Proposition 10.16, including editor's
% notes 102--103.
% Authority: Polo--Gille Exp6B-13oct24.pdf, component-local pages
% 74--75 (combined-reader pages 415--416; re-edition headers
% 404--405; original-margin printed pages 394--395).
% Section 11 begins on local page 76 and is excluded.
% OCR and Reinhold English Markdown are comparison witnesses only.
% This unit contains no diagrams.

\medskip
\phantomsection\label{par:sga3-VIB-10.14}%
\noindent\textbf{10.14.---}
Statements \hyperref[lem:sga3-VIB-10.2]{10.2} and
\hyperref[lem:sga3-VIB-10.3]{10.3}, concerning the category
of \(S\)-groups of finite presentation, extend to the category
of pairs consisting of an \(S\)-group of finite presentation
and an \(S\)-scheme of finite presentation with operator group
\(G\). More precisely, in the situation recalled at the
beginning of \hyperref[par:sga3-VIB-10.0]{10.0}:

\begin{enumerate}
\item[\textup{(i)}]\phantomsection\label{concl:sga3-VIB-10.14-i}
  let \(j\in I\), and let \(G_j\) and \(G'_j\) be two
  \(S_j\)-groups of finite presentation; let \(H_j\)
  \textup{(respectively, \(H'_j\))} be an \(S_j\)-scheme of
  finite presentation with operator group \(G_j\)
  \textup{(respectively, \(G'_j\))}. For \(i\in I\),
  \(i\geq j\), put
  \[
    G_i=G_j\times_{S_j}S_i,
    \qquad
    G=G_j\times_{S_j}S,
  \]
  and define \(G'_i\), \(G'\), \(H_i\), \(H\), \(H'_i\),
  and \(H'\) similarly. Denote by
  \[
    \operatorname{Dihom}_{S\text{-}\mathrm{gr.}}
      \bigl((G,H),(G',H')\bigr)
  \]
  the set of di-morphisms of \(S\)-groups and \(S\)-schemes
  with operator groups from the pair \((G,H)\) to the pair
  \((G',H')\). Then the canonical map
  \phantomsection\label{form:sga3-VIB-10.14-Dihom-limit}%
  \[
    \varinjlim_{i\geq j}
      \operatorname{Dihom}_{S_i\text{-}\mathrm{gr.}}
        \bigl((G_i,H_i),(G'_i,H'_i)\bigr)
      \longrightarrow
      \operatorname{Dihom}_{S\text{-}\mathrm{gr.}}
        \bigl((G,H),(G',H')\bigr)
  \]
  is bijective.

\item[\textup{(ii)}]\phantomsection\label{concl:sga3-VIB-10.14-ii}
  let \(G\) be an \(S\)-group of finite presentation and
  \(H\) an \(S\)-scheme of finite presentation with operator
  group \(G\). Then there are an index \(j\in I\), an
  \(S_j\)-group \(G_j\) of finite presentation, an
  \(S_j\)-scheme \(H_j\) of finite presentation with operator
  group \(G_j\), and a di-isomorphism of \(S\)-groups and
  \(S\)-schemes with operator groups from
  \[
    \bigl(G_j\times_{S_j}S,\ H_j\times_{S_j}S\bigr)
  \]
  onto \((G,H)\).%
  \begingroup
  \renewcommand{\thefootnote}{\fnsymbol{footnote}}%
  \footnote[1]{\label{note:sga3-VIB-10.14-GH}
  Translator's note: the French re-edition prints \((G,X)\)
  here. Since this assertion begins with the pair \((G,H)\)
  and \(X\) is not defined in \SGARefLink{par:sga3-VIB-10.14}{10.14}, \(H\) is restored.}%
  \endgroup
\end{enumerate}

\medskip
\phantomsection\label{def:sga3-VIB-10.15}%
\noindent\textbf{Definition 10.15.---}%
\footnote{\label{note:sga3-VIB-102}Editor's note: the original
has been corrected by distinguishing the notions of a
``formally homogeneous'' object and a ``homogeneous'' object;
see IV, \S\S \SGARefLink{sga3:iv:subsection:5:1}{5.1} and \SGARefLink{sga3:iv:subsection:6:7}{6.7}.}
\emph{Let \(\mathcal T\) be a topology on
\((\mathbf{Sch}/S)\) coarser than the canonical topology.
Given an \(S\)-group scheme \(G\) and an \(S\)-scheme \(X\)
with operator group \(G\), one says that \(X\) is a
\emph{formally homogeneous space under \(G\)}
\textup{(relative to the topology \(\mathcal T\))} if the
morphism}
\phantomsection\label{form:sga3-VIB-10.15-Phi}%
\[
  \Phi:G\times_S X\longrightarrow X\times_S X,
  \qquad
  (g,x)\longmapsto(gx,x),
\]
\emph{defined for every \(S'\to S\), \(g\in G(S')\), and
\(x\in X(S')\), is an epimorphism in the category of sheaves
for the topology \(\mathcal T\); equivalently, \(\Phi\) is a
covering for the topology \(\mathcal T\)
\textup{(cf.\
\hyperref[sga3:iv:proposition:4:4:3]{IV, 4.4.3})}.}

\emph{One says that \(X\) is a \emph{homogeneous space} if it
is formally homogeneous and, in addition, the morphism
\(X\to S\) is also a covering for the topology \(\mathcal T\).}

\emph{In particular, one says that \(X\) is a
\emph{formally principal homogeneous space under \(G\)} if
\(\Phi\) is an isomorphism, and that \(X\) is a
\emph{principal homogeneous bundle} \textup{(or a
\(G\)-torsor)} if \(\Phi\) is an isomorphism and, in addition,
the morphism \(X\to S\) is a covering for the topology
\(\mathcal T\) \textup{(cf.\
\hyperref[sga3:iv:definition:5:1:5]{IV, 5.1.5} and
\hyperref[sga3:iv:proposition:5:1:6:item:ii:01]{5.1.6(ii)})}.}

\medskip
\phantomsection\label{prop:sga3-VIB-10.16}%
\noindent\textbf{Proposition 10.16.---}
\emph{Retain the situation considered at the beginning of
\hyperref[par:sga3-VIB-10.0]{10.0}. Let \(j\in I\), let
\(G_j\) be an \(S_j\)-group, and let \(X_j\) be an
\(S_j\)-scheme with operator group \(G_j\). Suppose that
\(G_j\) and \(X_j\) are of finite presentation over \(S_j\).}

\emph{For}
\[
  X=X_j\times_{S_j}S
\]
\emph{to be a homogeneous space \textup{(respectively, a
principal homogeneous bundle)} under}
\[
  G=G_j\times_{S_j}S
\]
\emph{for the \textup{(fppf)} topology, it is necessary and
sufficient that there be an index \(i\geq j\) such that}
\[
  X_i=X_j\times_{S_j}S_i
\]
\emph{is a homogeneous space \textup{(respectively, a
principal homogeneous bundle)} under}
\[
  G_i=G_j\times_{S_j}S_i.
\]

Taking into account \hyperref[par:sga3-VIB-10.14]{10.14} and
EGA~IV\(_3\), 8.8.2, 8.8.3, and 8.10.5, the assertion follows
from the property of covering morphisms for the
\textup{(fppf)} topology recalled in
\hyperref[par:sga3-VIB-10.1]{10.1}.%
\footnote{\label{note:sga3-VIB-103}Editor's note: for
limit-passage properties of torsors expressed in terms of
\v{C}ech cohomology, see also SGA~4, VII, 5.7 and its
corollaries. This is developed in detail in \textup{[Ma07]}.}

% END INCLUDED SGA3_English_current_buildable_source_20260823/source/inputs/VI/tex_reference_v2/components/81_expose_VIB_par1014_through_prop1016_en.tex

% BEGIN INCLUDED SGA3_English_current_buildable_source_20260823/source/inputs/VI/tex_reference_v2/components/82_expose_VIB_section11_through_prop1131_loop2_r1_en.tex
% SGA 3, Expose VI_B, Loop 1 English reconstruction.
% Scope: Section 11 heading, paragraph 11.0, Notation,
% Lemma 11.1, Corollary 11.2, Definition 11.3, and
% Proposition 11.3.1, including editor's notes 104--107.
% Authority: Polo--Gille Exp6B-13oct24.pdf, component-local
% pages 76--78 (combined-reader pages 417--419; re-edition
% headers 406--408; original-margin printed pages 396--398).
% Paragraph 11.4 begins on local page 78 and is excluded.
% OCR and Reinhold English Markdown are comparison witnesses only.
% Loop 2 replaces both source-bound commutative-square crops with
% faithful native LaTeX reconstructions; both are closed.

\section*{11. Affine group schemes}
\phantomsection\label{sec:sga3-VIB-11}%
\addcontentsline{toc}{section}{11. Affine group schemes}

\medskip
\phantomsection\label{par:sga3-VIB-11.0}%
\noindent\textbf{11.0. Reminders.---}%
\footnote{\label{note:sga3-VIB-104}Editor's note: this
paragraph of reminders has been added.}
Let \(q:X\to S\) be a quasi-compact and quasi-separated
morphism of schemes \textup{(cf.\ EGA~IV\(_1\), 1.1 and
1.2)}, and let \(\mathcal F\) be a quasi-coherent
\(\mathcal O_X\)-module. Recall that \(q_*(\mathcal F)\) is a
quasi-coherent \(\mathcal O_S\)-module
\textup{(EGA~I, 9.2.1)}. Moreover, by EGA~III, 1.4.15
\textup{(completed by EGA~IV\(_1\), 1.7.21)}, one has
\hyperref[concl:sga3-VIB-11.0-c]{\textup{(c)}} below, and the
proof in \emph{loc.\ cit.} also gives
\hyperref[concl:sga3-VIB-11.0-a]{\textup{(a)}} and
\hyperref[concl:sga3-VIB-11.0-b]{\textup{(b)}}:
\begin{enumerate}
\item[\textup{(a)}]\phantomsection\label{concl:sga3-VIB-11.0-a}
  If \(\mathcal F\) is a filtered direct limit of
  quasi-coherent submodules \(\mathcal F_\alpha\), then
  \[
    q_*(\mathcal F)
      =\varinjlim_\alpha q_*(\mathcal F_\alpha).
  \]

\item[\textup{\SGAInlineRefTarget{par:sga3-VIB-11.0:item:b}(b)}]\phantomsection\label{concl:sga3-VIB-11.0-b}
  If \(\mathcal E\) is a flat \(\mathcal O_S\)-module, the
  canonical morphism
  \phantomsection\label{form:sga3-VIB-11.0-projection}%
  \[
    \mathcal E\otimes_{\mathcal O_S}q_*(\mathcal O_X)
      \longrightarrow q_*q^*(\mathcal E)
  \]
  is an isomorphism.

\item[\textup{(c)}]\phantomsection\label{concl:sga3-VIB-11.0-c}
  Let \(p:S'\to S\) be a flat morphism, let
  \(q':X'\to S'\) be obtained from \(q\) by base change, and
  let \(\mathcal F'\) be the inverse image of \(\mathcal F\)
  on \(X'\). Then the canonical morphism
  \phantomsection\label{form:sga3-VIB-11.0-flat-base-change}%
  \[
    p^*q_*(\mathcal F)\longrightarrow q'_*(\mathcal F')
  \]
  is an isomorphism.
\end{enumerate}

Indeed, let \(U=\operatorname{Spec}(A)\) be an arbitrary
affine open subset of \(S\). By hypothesis, \(q^{-1}(U)\) is
the union of affine open subsets
\[
  V_i=\operatorname{Spec}(B_i),\qquad i=1,\ldots,n,
\]
and every intersection \(V_i\cap V_j\) is the union of
finitely many affine open subsets
\[
  W_{ijk}=\operatorname{Spec}(C_{ijk}).
\]
Then
\(\Gamma(U,q_*(\mathcal F))
  =\Gamma(q^{-1}(U),\mathcal F)\)
is the kernel of the morphism
\phantomsection\label{form:sga3-VIB-11.0-Cech-F}%
\[
  \bigoplus_{i=1}^{n}\Gamma(V_i,\mathcal F)
    \longrightarrow
  \bigoplus_{i,j,k}\Gamma(W_{ijk},\mathcal F).
\]
Assertion~\hyperref[concl:sga3-VIB-11.0-a]{\textup{(a)}}
follows, since every term above commutes with filtered direct
limits \textup{(the \(V_i\) and \(W_{ijk}\) being affine,
hence quasi-compact)}.

Let us prove
\hyperref[concl:sga3-VIB-11.0-b]{\textup{(b)}}. The module
\(E=\Gamma(U,\mathcal E)\) is flat over \(A\), and
\(\Gamma(U,q_*q^*(\mathcal E))\) is the kernel \(K(E)\) of
the morphism
\phantomsection\label{form:sga3-VIB-11.0-Cech-E}%
\[
  \bigoplus_{i=1}^{n}B_i\otimes_AE
    \longrightarrow
  \bigoplus_{i,j,k}C_{ijk}\otimes_AE.
\]
Since \(E\) is flat over \(A\), this kernel is identified
with
\[
  K(A)\otimes_AE
    =\mathcal O_X(q^{-1}(U))\otimes_AE.
\]
Finally, if \(U'\) is an arbitrary affine open subset of
\(S'\) lying over \(U\), then
\(A'=\mathcal O(U')\) is a flat \(A\)-algebra, and the same
argument gives%
\begingroup
\renewcommand{\thefootnote}{\fnsymbol{footnote}}%
\footnote[1]{\label{note:sga3-VIB-11.0-U-prime}
Translator's note: the French re-edition prints
\(q'^{-1}(U)\) in the right-hand side below. Since
\(q':X'\to S'\) and the affine open subset just introduced
is \(U'\subset S'\), the inverse image is
\(q'^{-1}(U')\).}%
\endgroup
\phantomsection\label{form:sga3-VIB-11.0-local-base-change}%
\[
  \mathcal F(q^{-1}(U))\otimes_AA'
    \xrightarrow{\ \sim\ }
  \mathcal F'(q'^{-1}(U')).
\]

\medskip
\phantomsection\label{not:sga3-VIB-11-A}%
\noindent\textbf{Notation.---}
Let \(S\) be a scheme, let \(X\) be an \(S\)-scheme, and let
\(f:X\to S\) be the structural morphism. We shall put
\phantomsection\label{form:sga3-VIB-11-A-X}%
\[
  \mathcal A(X)=f_*(\mathcal O_X).
\]

\medskip
\phantomsection\label{lem:sga3-VIB-11.1}%
\noindent\textbf{Lemma 11.1.---}
\emph{Let \(X\) and \(Y\) be two quasi-compact and
quasi-separated \(S\)-schemes, and let \(f:X\to S\) and
\(g:Y\to S\) be their structural morphisms. Then the
canonical homomorphism}
\phantomsection\label{form:sga3-VIB-11.1-phi}%
\[
  \varphi:
  \mathcal A(X)\otimes_{\mathcal O_S}\mathcal A(Y)
    \longrightarrow
  \mathcal A(X\times_SY)
\]
\emph{is an isomorphism in each of the following cases:}%
\footnote{\label{note:sga3-VIB-105}Editor's note: if \(k\) is
a field and \(X\) is an infinite disjoint union of copies of
\(S=\operatorname{Spec}k\) \textup{(so that \(X\) is not
quasi-compact)}, then \(\mathcal A(X)=k^X\), and the
canonical morphism
\[
  k^X\otimes_k k^X\longrightarrow k^{X\times X}
\]
is not surjective.}
\begin{enumerate}
\item[\textup{(a)}]\phantomsection\label{concl:sga3-VIB-11.1-a}
  \emph{\(f\) and \(g\) are affine;}

\item[\textup{(b)}]\phantomsection\label{concl:sga3-VIB-11.1-b}
  \emph{\(g\) \textup{(or \(f\))} is flat and affine;}

\item[\textup{(c)}]\phantomsection\label{concl:sga3-VIB-11.1-c}
  \emph{\(g\) is flat and \(f_*(\mathcal O_X)\) is a flat
  \(\mathcal O_S\)-module.}
\end{enumerate}

In case~\hyperref[concl:sga3-VIB-11.1-b]{\textup{(b)}}, we
shall suppose that it is \(g\) that is flat and affine. Put
\[
  S'=\operatorname{Spec}\mathcal A(X),\qquad
  Y'=Y\times_SS',\qquad
  g'=g\times_SS',
\]
and write \(v:S'\to S\) for the structural morphism:
\begin{center}
  \phantomsection\label{diag:sga3-VIB-11.1-base-change}%
  \begin{tikzcd}[row sep=large,column sep=large]
    Y' \arrow[r] \arrow[d,"g'"'] &
    Y \arrow[d,"g"] \\
    \operatorname{Spec}\mathcal A(X)=S'
      \arrow[r,"v"'] &
    S
  \end{tikzcd}
\end{center}
In cases~\hyperref[concl:sga3-VIB-11.1-a]{\textup{(a)}} and
\hyperref[concl:sga3-VIB-11.1-b]{\textup{(b)}}, the morphism
\(g\) is affine, and hence EGA~II, 1.5.2 gives
\phantomsection\label{form:sga3-VIB-11.1-eq1}%
\[
  \tag{1}
  g'_*(\mathcal O_{Y'})
    =v^*g_*(\mathcal O_Y)
    =\mathcal A(Y)\otimes_{\mathcal O_S}\mathcal O_{S'}.
\]
The same equality holds in
case~\hyperref[concl:sga3-VIB-11.1-c]{\textup{(c)}}, by
\hyperref[concl:sga3-VIB-11.0-c]{11.0\textup{(c)}}, since
\(S'\) is flat over \(S\) and \(g\) is quasi-compact and
quasi-separated.

On the other hand, by EGA~II, 1.2.7, the morphism
\(f:X\to S\) factors through \(v\) by means of a morphism
\(p:X\to S'\), and
\[
  p_*(\mathcal O_X)=\mathcal O_{S'},
  \qquad
  X\times_SY=X\times_{S'}Y'.
\]
Since \(f\) is quasi-separated, so is \(p\)
\textup{(EGA~IV\(_1\), 1.2.2)}; and since \(f\) is
quasi-compact and \(v\) is quasi-separated, \(p\) is also
quasi-compact \textup{(EGA~IV\(_1\), 1.2.4)}. Consider the
cartesian square
\begin{center}
  \phantomsection\label{diag:sga3-VIB-11.1-cartesian}%
  \begin{tikzcd}[row sep=large,column sep=large]
    X\times_SY \arrow[r,"p'"] \arrow[d] &
    Y' \arrow[d,"g'"] \\
    X \arrow[r,"p"'] &
    S'
  \end{tikzcd}
\end{center}
In cases~\hyperref[concl:sga3-VIB-11.1-b]{\textup{(b)}} and
\hyperref[concl:sga3-VIB-11.1-c]{\textup{(c)}}, the scheme
\(Y\) is flat over \(S\), and hence \(Y'\) is flat over
\(S'\). Applying
\hyperref[concl:sga3-VIB-11.0-c]{11.0\textup{(c)}} once
again, we obtain
\phantomsection\label{form:sga3-VIB-11.1-eq2}%
\[
  \tag{2}
  p'_*(\mathcal O_{X\times_SY})
    =g'^*p_*(\mathcal O_X)
    =g'^*(\mathcal O_{S'})
    =\mathcal O_{Y'}.
\]
The same equality holds in
case~\hyperref[concl:sga3-VIB-11.1-a]{\textup{(a)}}, because
\(p\) and \(p'\) are then isomorphisms.

Finally, since \(v\) is affine, EGA~II, 1.4.7 gives
\[
  v_*(\mathcal F\otimes_{\mathcal O_S}\mathcal O_{S'})
    =\mathcal F\otimes_{\mathcal O_S}\mathcal A(X)
\]
for every quasi-coherent \(\mathcal O_S\)-module
\(\mathcal F\). Combined with
\hyperref[form:sga3-VIB-11.1-eq2]{\textup{(2)}} and
\hyperref[form:sga3-VIB-11.1-eq1]{\textup{(1)}}, this gives
\[
\begin{aligned}
  \mathcal A(X\times_SY)
    &=v_*g'_*p'_*(\mathcal O_{X\times_SY})\\
    &\overset{\textup{(2)}}{=}v_*g'_*(\mathcal O_{Y'})
     \overset{\textup{(1)}}{=}
       v_*(\mathcal A(Y)\otimes_{\mathcal O_S}\mathcal O_{S'})\\
    &=\mathcal A(Y)\otimes_{\mathcal O_S}\mathcal A(X).
\end{aligned}
\]

\medskip
\phantomsection\label{cor:sga3-VIB-11.2}%
\noindent\textbf{Corollary 11.2.---}
\emph{The functor}
\[
  X\longmapsto\operatorname{Spec}\mathcal A(X),
\]
\emph{from the full subcategory of
\((\mathbf{Sch}/S)\) formed by the \(S\)-schemes \(X\) that
are flat, quasi-compact, and quasi-separated over \(S\), and
for which \(\mathcal A(X)\) is a flat
\(\mathcal O_S\)-module, to the category of flat affine
\(S\)-schemes, commutes with finite products and therefore
takes \(S\)-groups to \(S\)-groups.}

\medskip
\phantomsection\label{def:sga3-VIB-11.3}%
\noindent\textbf{Definition 11.3.---}
\emph{Given a flat \(S\)-group \(G\), quasi-compact and
quasi-separated over \(S\), such that \(\mathcal A(G)\) is
flat over \(\mathcal O_S\),}%
\footnote{\label{note:sga3-VIB-106}Editor's note: if \(S\) is
a regular locally Noetherian scheme of dimension at most
\(2\), and \(X\) is an \(S\)-scheme that is flat,
quasi-compact, and quasi-separated, then \(\mathcal A(X)\) is
a flat \(\mathcal O_S\)-module; cf.\ \textup{[Ray70a]},
VII, \SGARefLink{sga3:VIIB:3.2:prop}{3.2}.}
\emph{we write \(G_{\mathrm{af}}\), and call the
\(S\)-group}
\[
  G_{\mathrm{af}}=\operatorname{Spec}\mathcal A(G)
\]
\emph{the \emph{affine envelope} of \(G\).}

\medskip
\phantomsection\label{prop:sga3-VIB-11.3.1}%
\noindent\textbf{Proposition 11.3.1.---}%
\footnote{\label{note:sga3-VIB-107}Editor's note: this
proposition has been added; see also the additional
\hyperref[sec:sga3-VIB-12]{Section~12} below for a study of the morphism
\(G\to G_{\mathrm{af}}\) and of its kernel.}
\emph{The canonical morphism}
\[
  \tau_G:G\longrightarrow G_{\mathrm{af}}
\]
\emph{is a morphism of \(S\)-groups. Moreover, it has the
following universal property:}
\begin{enumerate}
\item[\textup{(i)}]\phantomsection\label{concl:sga3-VIB-11.3.1-i}
  \emph{For every morphism of \(S\)-schemes
  \(\varphi:G\to H\), where \(H\) is affine over \(S\), there
  is a unique morphism of \(S\)-schemes
  \(\varphi':G_{\mathrm{af}}\to H\) such that}
  \[
    \varphi=\varphi'\circ\tau_G.
  \]

\item[\textup{(ii)}]\phantomsection\label{concl:sga3-VIB-11.3.1-ii}
  \emph{If, moreover, \(H\) is an \(S\)-group and
  \(\varphi\) is a morphism of \(S\)-groups, then the same is
  true of \(\varphi'\).}
\end{enumerate}

% END INCLUDED SGA3_English_current_buildable_source_20260823/source/inputs/VI/tex_reference_v2/components/82_expose_VIB_section11_through_prop1131_loop2_r1_en.tex

% BEGIN INCLUDED SGA3_English_current_buildable_source_20260823/source/inputs/VI/tex_reference_v2/components/83_expose_VIB_par114_through_rem117bis_loop2_r1_en.tex
% SGA 3, Expose VI_B, Loop 1 English reconstruction.
% Scope: paragraph 11.4, Proposition 11.5, paragraph 11.6,
% Remark 11.6.A, Proposition 11.6.1, Lemma 11.7, and
% Remark 11.7.bis, including editor's notes 108--111.
% Authority: Polo--Gille Exp6B-13oct24.pdf, component-local
% pages 78--82 (combined-reader pages 419--423; re-edition
% headers 408--412; original-margin printed pages 397--401).
% Paragraph 11.4 begins on local page 78, still within printed page 397;
% printed page 398 begins there only at Proposition 11.5. Definition
% 11.8.0 begins on local page 82, still within printed page 401, and is
% excluded. Printed page 402 begins on local page 85.
% OCR and Reinhold English Markdown are comparison witnesses only.
% Loop 2 replaces all seven source-bound commutative-diagram crops
% with faithful native LaTeX reconstructions; all are closed.

\medskip
\phantomsection\label{par:sga3-VIB-11.4}%
\noindent\textbf{11.4.---}%
\footnote{\label{note:sga3-VIB-108}Editor's note: in
\SGARefLink{par:sga3-VIB-11.4}{11.4}--\SGARefLink{par:sga3-VIB-11.6}{11.6}, we have considered
\(\underline{\operatorname{Hom}}_{\mathcal O_S}
  (\mathrm W(\mathcal E),\mathrm W(\mathcal F))\)
instead of
\(\underline{\operatorname{Hom}}_{\mathcal O_S}
  (\mathrm V(\mathcal F),\mathrm V(\mathcal E))\)
\textup{(cf.\ I, \SGARefLink{sga3:i:proposition:4-6-3}{4.6.3})}, and simplified the original by
taking account of the addition \hyperref[concl:sga3-VIB-11.0-b]{11.0}\textup{(b)}.}
Let \(\mathcal E\) and \(\mathcal F\) be two
quasi-coherent \(\mathcal O_S\)-modules. Consider the
\(S\)-functor
\[
  \underline{\operatorname{Hom}}_{\mathcal O_S}
    \bigl(\mathrm W(\mathcal E),\mathrm W(\mathcal F)\bigr)
\]
\textup{(cf.\ I, \SGARefLink{sga3:i:definition:3-1-4}{3.1.4})}, that is, for every \(S\)-scheme
\(f:X\to S\),
\phantomsection\label{form:sga3-VIB-11.4-HomX}%
\[
  \underline{\operatorname{Hom}}_{\mathcal O_S}
    \bigl(\mathrm W(\mathcal E),\mathrm W(\mathcal F)\bigr)(X)
  =
  \operatorname{Hom}_{\mathcal O_X}
    \bigl(\mathrm W(\mathcal E)_X,\mathrm W(\mathcal F)_X\bigr).
\]
Moreover, by I, \SGARefLink{sga3:i:proposition:4-6-2}{4.6.2}, one has
\(\mathrm W(\mathcal E)_X=\mathrm W(f^*(\mathcal E))\)
\textup{(and similarly for \(\mathcal F\))}, and
\[
  \operatorname{Hom}_{\mathcal O_X}
    \bigl(\mathrm W(f^*(\mathcal E)),
          \mathrm W(f^*(\mathcal F))\bigr)
  =
  \operatorname{Hom}_{\mathcal O_X}
    \bigl(f^*(\mathcal E),f^*(\mathcal F)\bigr).
\]
We therefore obtain \textup{(using the adjunction formula
for the last equality)}
\phantomsection\label{form:sga3-VIB-11.4-adjunction}%
\[
  \tag{\(\dagger\)}
  \underline{\operatorname{Hom}}_{\mathcal O_S}
    \bigl(\mathrm W(\mathcal E),\mathrm W(\mathcal F)\bigr)(X)
  =
  \operatorname{Hom}_{\mathcal O_X}
    \bigl(f^*(\mathcal E),f^*(\mathcal F)\bigr)
  =
  \operatorname{Hom}_{\mathcal O_S}
    \bigl(\mathcal E,f_*f^*(\mathcal F)\bigr).
\]

\medskip
\phantomsection\label{prop:sga3-VIB-11.5}%
\noindent\textbf{Proposition 11.5.---}
\emph{Let \(X\) be an \(S\)-scheme that is quasi-compact and
quasi-separated over \(S\), let \(f:X\to S\) be the
structural morphism, and let \(\mathcal E\) and
\(\mathcal E'\) be two quasi-coherent
\(\mathcal O_S\)-modules. Suppose that one of the following
two conditions holds:}
\begin{enumerate}
\item[\textup{a)}]\phantomsection\label{concl:sga3-VIB-11.5-a}
  \emph{\(f\) is affine;}

\item[\textup{b)}]\phantomsection\label{concl:sga3-VIB-11.5-b}
  \emph{\(\mathcal E\) is flat over \(\mathcal O_S\).}
\end{enumerate}
\emph{Then the canonical morphism}
\phantomsection\label{form:sga3-VIB-11.5-projection}%
\[
  \mathcal E\otimes_{\mathcal O_S}\mathcal A(X)
    \longrightarrow f_*f^*(\mathcal E)
\]
\emph{is an isomorphism, and hence}
\phantomsection\label{form:sga3-VIB-11.5-Hom}%
\[
  \underline{\operatorname{Hom}}_{\mathcal O_S}
    \bigl(\mathrm W(\mathcal E'),\mathrm W(\mathcal E)\bigr)(X)
  =
  \operatorname{Hom}_{\mathcal O_S}
    \bigl(\mathcal E',
          \mathcal E\otimes_{\mathcal O_S}\mathcal A(X)\bigr).
\]

Indeed, the second assertion follows from
\hyperref[par:sga3-VIB-11.4]{11.4} and the first; the latter
follows from EGA~II, 1.4.7 in
case~\hyperref[concl:sga3-VIB-11.5-a]{\textup{(a)}}, and from
\hyperref[concl:sga3-VIB-11.0-b]{11.0\textup{(b)}} in
case~\hyperref[concl:sga3-VIB-11.5-b]{\textup{(b)}}.% 
\footnote{\label{note:sga3-VIB-109}Editor's note: we have
simplified the original, which used the isomorphism
\[
  \underline{\operatorname{Hom}}_{\mathcal O_S}
    \bigl(\mathrm W(\mathcal E'),\mathrm W(\mathcal E)\bigr)
  \simeq
  \underline{\operatorname{Hom}}_{\mathcal O_S}
    \bigl(\mathrm V(\mathcal E),\mathrm V(\mathcal E')\bigr),
\]
then the inclusion of the right-hand term in
\[
  \underline{\operatorname{Hom}}_S
    \bigl(\mathrm V(\mathcal E),\mathrm V(\mathcal E')\bigr)
  =
  \operatorname{Hom}_{\mathcal O_S}
    \bigl(\mathcal E',
          f_*f^*(\operatorname{Sym}(\mathcal E))\bigr),
\]
and applied EGA~III, 4.1.15 to
\(\mathrm V(\mathcal E)=
  \operatorname{Spec}(\operatorname{Sym}(\mathcal E))\)
to deduce \SGARefLink{par:sga3-VIB-11.0}{11.0}\textup{\SGARefLink{par:sga3-VIB-11.0:item:b}{(b)}}.}

\medskip
\phantomsection\label{par:sga3-VIB-11.6}%
\noindent\textbf{11.6.---}%
\footnote{\label{note:sga3-VIB-110}Editor's note: we have
expanded \SGARefLink{par:sga3-VIB-11.6}{11.6} and displayed the results obtained in the form
of \hyperref[prop:sga3-VIB-11.6.1]{Proposition 11.6.1}.}
Let \(G\) be an \(S\)-group and let \(\mathcal E\) be a
quasi-coherent \(\mathcal O_S\)-module. Giving \(\mathcal E\)
a \(G\)-\(\mathcal O_S\)-module structure \textup{(that is,
an \(\mathcal O_S\)-linear action of \(G\) on
\(\mathrm W(\mathcal E)\), cf.\ I, \SGARefLink{sga3:i:definition:4-7-1}{4.7.1})} is equivalent to
giving a morphism of \(S\)-functors in monoids
\phantomsection\label{form:sga3-VIB-11.6-rho}%
\[
  \rho:G\longrightarrow
  \underline{\operatorname{End}}_{\mathcal O_S}
    \bigl(\mathrm W(\mathcal E)\bigr);
\]
indeed, such a \(\rho\) necessarily takes \(G\) into
\(\underline{\operatorname{Aut}}_{\mathcal O_S}
  (\mathrm W(\mathcal E))\).

By \hyperref[par:sga3-VIB-11.4]{11.4}, giving a morphism of
\(S\)-functors
\[
  \rho:G\longrightarrow
  \underline{\operatorname{End}}_{\mathcal O_S}
    \bigl(\mathrm W(\mathcal E)\bigr)
\]
is equivalent to giving an element \(\theta\) of
\[
  \operatorname{Hom}_{\mathcal O_G}
    \bigl(f^*(\mathcal E),f^*(\mathcal E)\bigr),
\]
which corresponds by adjunction to an element \(\delta\) of
\[
  \operatorname{Hom}_{\mathcal O_S}
    \bigl(\mathcal E,f_*f^*(\mathcal E)\bigr),
\]
where \(f\) denotes the projection \(G\to S\).

Let \(m:G\times_SG\to G\) be multiplication, let
\(\delta_G\) be the morphism
\[
  \mathcal O_G\longrightarrow
  m_*(\mathcal O_{G\times_SG}),
\]
and let \(\varphi:G\times_SG\to S\) be the projection
\textup{(which equals \(f\circ m\))}. It is convenient to
write \(\boxtimes\) for the ``external'' tensor product. We
thus obtain a morphism
\phantomsection\label{form:sga3-VIB-11.6-id-deltaG}%
\[
  \operatorname{id}_{\mathcal E}\boxtimes\delta_G:
  f^*(\mathcal E)
    =\mathcal E\boxtimes_{\mathcal O_S}\mathcal O_G
  \longrightarrow
  \mathcal E\boxtimes_{\mathcal O_S}
    m_*(\mathcal O_{G\times_SG}).
\]
By abuse of notation, we again write
\(\operatorname{id}_{\mathcal E}\boxtimes\delta_G\) for the
composite of this morphism with the canonical morphism
\[
  \mathcal E\boxtimes_{\mathcal O_S}
    m_*(\mathcal O_{G\times_SG})
  \longrightarrow
  m_*m^*f^*(\mathcal E)
  =m_*\varphi^*(\mathcal E).
\]

On the other hand, let \(h:G\to S\) denote a second copy of
\(f:G\to S\), and consider the following commutative diagram,
in which \(p,q\) denote the two projections:
\begin{center}
  \phantomsection\label{diag:sga3-VIB-11.6-projections}%
  \begin{tikzcd}[row sep=large,column sep=large]
    G\times_SG
      \arrow[r,"q"]
      \arrow[d,"p"']
      \arrow[dr,"\varphi" description] &
    G \arrow[d,"h"] \\
    G \arrow[r,"f"'] &
    S
  \end{tikzcd}
\end{center}
Writing \(\delta\) again for the morphism
\(\mathcal E\to h_*h^*(\mathcal E)\), we obtain the morphism
\phantomsection\label{form:sga3-VIB-11.6-delta-id}%
\[
  \delta\boxtimes\operatorname{id}_{\mathcal O_G}:
  f^*(\mathcal E)
    =\mathcal E\boxtimes_{\mathcal O_S}\mathcal O_G
  \longrightarrow
  h_*h^*(\mathcal E)\boxtimes_{\mathcal O_S}\mathcal O_G
    =f^*h_*h^*(\mathcal E).
\]
By abuse of notation, we again write
\(\delta\boxtimes\operatorname{id}_{\mathcal O_G}\) for the
composite of this morphism with the canonical morphism
\[
  f^*h_*h^*(\mathcal E)\longrightarrow
  p_*\varphi^*(\mathcal E).
\]

The condition that \(\rho\) be compatible with multiplication
is then equivalent to saying that, for every open subset
\(U\) of \(S\), the following diagram is commutative:
\begin{center}
  \phantomsection\label{diag:sga3-VIB-11.6-1}%
  \makebox[0pt][r]{\textup{(1)}\qquad}%
  {\scriptsize
  \begin{tikzcd}[row sep=large,column sep=large]
    \Gamma(U,\mathcal E)
      \arrow[r,"\delta"]
      \arrow[d,"\delta"'] &
    \Gamma\!\left(
      U\times_SG,\,
      \mathcal E\boxtimes_{\mathcal O_S}\mathcal O_G
    \right)
      \arrow[d,"{\operatorname{id}_{\mathcal E}\boxtimes\delta_G}"] \\
    \Gamma\!\left(
      U\times_SG,\,
      \mathcal E\boxtimes_{\mathcal O_S}\mathcal O_G
    \right)
      \arrow[r,"{\delta\boxtimes\operatorname{id}_{\mathcal O_G}}"'] &
    \Gamma\!\left(
      U\times_SG\times_SG,\,
      \mathcal E\boxtimes_{\mathcal O_S}\mathcal O_G
        \boxtimes_{\mathcal O_S}\mathcal O_G
    \right)
  \end{tikzcd}}
\end{center}
Furthermore, the unit section \(\varepsilon:S\to G\) induces
a morphism \(u\) from \(\mathcal O_G\) to
\[
  \varepsilon_*\varepsilon^*(\mathcal O_G)
  =\varepsilon_*(\mathcal O_S),
\]
and the condition that \(\rho\) preserve unit elements is
equivalent to the commutativity of the diagram
\begin{center}
  \phantomsection\label{diag:sga3-VIB-11.6-2}%
  \makebox[0pt][r]{\textup{(2)}\qquad}%
  {\small
  \begin{tikzcd}[row sep=large,column sep=large]
    \Gamma(U,\mathcal E)
      \arrow[rr,"\delta"]
      \arrow[dr,"\sim"'] &
    {} &
    \Gamma\!\left(
      U\times_SG,\,
      \mathcal E\boxtimes_{\mathcal O_S}\mathcal O_G
    \right)
      \arrow[dl,"{\operatorname{id}_{\mathcal E}\boxtimes u}"] \\
    {} &
    \Gamma\!\left(
      U\times_SS,\,
      \mathcal E\boxtimes_{\mathcal O_S}\mathcal O_S
    \right)\mathpunct{.} &
    {}
  \end{tikzcd}}
\end{center}

Thus, giving \(\mathcal E\) a
\(G\)-\(\mathcal O_S\)-module structure is equivalent to
giving a morphism of \(\mathcal O_S\)-modules
\[
  \delta:\mathcal E\longrightarrow f_*f^*(\mathcal E)
\]
satisfying conditions
\hyperref[diag:sga3-VIB-11.6-1]{\textup{(1)}} and
\hyperref[diag:sga3-VIB-11.6-2]{\textup{(2)}} above. In that
case, the morphism
\[
  \theta:f^*(\mathcal E)\longrightarrow f^*(\mathcal E)
\]
deduced from \(\delta\) by adjunction is an isomorphism,
because it corresponds to the isomorphism
\[
  G\times_S\mathrm W(\mathcal E)
    \xrightarrow{\ \sim\ }
  G\times_S\mathrm W(\mathcal E)
\]
defined set-theoretically by
\((g,x)\longmapsto(g,gx)\); see also I, \SGARefLink{sga3:i:remark:6-5-4}{6.5.4}.

Now suppose that \(G\) is flat, quasi-compact, and
quasi-separated over \(S\), and that \(\mathcal A(G)\) is a
flat \(\mathcal O_S\)-module. Then, by
\hyperref[concl:sga3-VIB-11.1-c]{11.1\textup{(c)}}, the
canonical morphism
\[
  \mathcal A(G)\otimes_{\mathcal O_S}\mathcal A(G)
    \longrightarrow
  \mathcal A(G\times_SG)
\]
is an isomorphism, and we write \(\Delta\) for the morphism
\[
  \delta_G:\mathcal A(G)\longrightarrow
  \mathcal A(G)\otimes_{\mathcal O_S}\mathcal A(G).
\]

If, moreover,
\[
  \mathcal E\otimes_{\mathcal O_S}\mathcal A(G)
    =f_*f^*(\mathcal E),
\]
as is the case by \hyperref[prop:sga3-VIB-11.5]{11.5} if
\(G\to S\) is affine or if \(\mathcal E\) is flat over
\(\mathcal O_S\), then conditions
\hyperref[diag:sga3-VIB-11.6-1]{\textup{(1)}} and
\hyperref[diag:sga3-VIB-11.6-2]{\textup{(2)}} are equivalent
to the following conditions, which express that
\[
  \delta:\mathcal E\longrightarrow
  \mathcal E\otimes_{\mathcal O_S}\mathcal A(G)
\]
makes \(\mathcal E\) a right \(\mathcal A(G)\)-comodule
\textup{(cf.\ I, \SGARefLink{sga3:i:proposition:4-7-2}{4.7.2})}:

\smallskip
\phantomsection\label{cond:sga3-VIB-11.6-CM1}%
\noindent\textup{(CM 1)}
Put \(\mathcal A=\mathcal A(G)\). The following diagram is
commutative:
\begin{center}
  \phantomsection\label{diag:sga3-VIB-11.6-CM1}%
  \begin{tikzcd}[row sep=large,column sep=large]
    \mathcal E
      \arrow[r,"\delta"]
      \arrow[d,"\delta"'] &
    \mathcal E\otimes\mathcal A
      \arrow[d,"{\operatorname{id}_{\mathcal E}\otimes\Delta}"] \\
    \mathcal E\otimes\mathcal A
      \arrow[r,"{\delta\otimes\operatorname{id}_{\mathcal A}}"'] &
    \mathcal E\otimes\mathcal A\otimes\mathcal A
  \end{tikzcd}
\end{center}

\smallskip
\phantomsection\label{cond:sga3-VIB-11.6-CM2}%
\noindent\textup{(CM 2)}
Write
\(\eta:\mathcal A\to\mathcal O_S\) for the morphism
\(\mathcal A(\varepsilon)\). The following diagram is
commutative:
\begin{center}
  \phantomsection\label{diag:sga3-VIB-11.6-CM2}%
  \begin{tikzcd}[row sep=large,column sep=large]
    \mathcal E
      \arrow[rr,"\delta"]
      \arrow[dr,"\sim"'] &
    {} &
    \mathcal E\otimes\mathcal A
      \arrow[dl,"{\operatorname{id}_{\mathcal E}\otimes\eta}"] \\
    {} &
    \mathcal E\otimes\mathcal O_S\mathpunct{.} &
    {}
  \end{tikzcd}
\end{center}

\medskip
\phantomsection\label{rem:sga3-VIB-11.6.A}%
\noindent\textbf{Remark 11.6.A.---}
Recall that \(\mathrm V(\mathcal E)\) denotes the vector
fibration over \(S\) representing the functor
\(\mathrm V(\mathcal E)\); that is, for every \(S'\to S\),
\[
  \mathrm V(\mathcal E)(S')
  =
  \operatorname{Hom}_{\mathcal O_{S'}}
    \bigl(\mathcal E_{S'},\mathcal O_{S'}\bigr).
\]
Since I, \SGARefLink{sga3:i:proposition:4-6-2}{4.6.2} gives an anti-isomorphism of \(S\)-functors
in monoids
\[
  \underline{\operatorname{End}}_{\mathcal O_S}
    \bigl(\mathrm W(\mathcal E)\bigr)
  \simeq
  \underline{\operatorname{End}}_{\mathcal O_S}
    \bigl(\mathrm V(\mathcal E)\bigr),
\]
if \(\mathcal E\) is a left
\(G\)-\(\mathcal O_S\)-module, then there is a right action
\[
  \mu:\mathrm V(\mathcal E)\times_SG
    \longrightarrow\mathrm V(\mathcal E)
\]
of \(G\) on \(\mathrm V(\mathcal E)\), defined
set-theoretically by
\[
  (\phi g)(x)=\phi(gx)
\]
for every \(g\in G(S')\),
\(x\in\Gamma(S',\mathcal E_{S'})\), and
\(\phi\in
  \operatorname{Hom}_{\mathcal O_{S'}}
    (\mathcal E_{S'},\mathcal O_{S'})\).
We therefore have the commutative diagrams
\begin{center}
  \phantomsection\label{diag:sga3-VIB-11.6.A-actions}%
  {\small
  \begin{tabular}{c@{\qquad}c}
    \begin{tikzcd}[row sep=large,column sep=large]
      \mathrm V(\mathcal E)\times_SG\times_SG
        \arrow[r,"{\mu\times\operatorname{id}_G}"]
        \arrow[d,"{\operatorname{id}_{\mathrm V(\mathcal E)}\times m}"'] &
      \mathrm V(\mathcal E)\times_SG
        \arrow[d,"\mu"] \\
      \mathrm V(\mathcal E)\times_SG
        \arrow[r,"\mu"'] &
      \mathrm V(\mathcal E)
    \end{tikzcd}
    &
    \begin{tikzcd}[row sep=large,column sep=large]
      \mathrm V(\mathcal E) &
      \mathrm V(\mathcal E)\times_SG
        \arrow[l,"\mu"'] \\
      {} &
      \mathrm V(\mathcal E)\times_SS\mathpunct{.}
        \arrow[u,"{\operatorname{id}_{\mathrm V(\mathcal E)}\times\varepsilon}"]
        \arrow[ul,"\sim"']
    \end{tikzcd}
  \end{tabular}}
\end{center}
When \(G\) is flat, quasi-compact, and quasi-separated over
\(S\), when \(\mathcal A(G)\) is a flat
\(\mathcal O_S\)-module, and when one of the conditions of
\hyperref[prop:sga3-VIB-11.5]{11.5} is satisfied, one
similarly recovers conditions
\hyperref[cond:sga3-VIB-11.6-CM1]{\textup{(CM 1)}} and
\hyperref[cond:sga3-VIB-11.6-CM2]{\textup{(CM 2)}}.

Consequently, we have proved:

\medskip
\phantomsection\label{prop:sga3-VIB-11.6.1}%
\noindent\textbf{Proposition 11.6.1.---}
\emph{Let \(G\) be an \(S\)-group that is flat,
quasi-compact, and quasi-separated over \(S\), such that
\(\mathcal A(G)\) is a flat \(\mathcal O_S\)-module, and let
\(\mathcal E\) be a quasi-coherent
\(\mathcal O_S\)-module.}
\begin{enumerate}
\item[\textup{(i)}]\phantomsection\label{concl:sga3-VIB-11.6.1-i}
  \emph{Giving an \(\mathcal A(G)\)-comodule structure}
  \[
    \delta:\mathcal E\longrightarrow
      \mathcal E\otimes_{\mathcal O_S}\mathcal A(G)
  \]
  \emph{is equivalent to giving
  \(\mathcal E\) a
  \(G_{\mathrm{af}}\)-\(\mathcal O_S\)-module structure
  \textup{(that is, an \(\mathcal O_S\)-linear action of
  \(G_{\mathrm{af}}\) on \(\mathcal E\))}. By composition
  with the morphism of \(S\)-groups
  \(G\to G_{\mathrm{af}}\), this defines a
  \(G\)-\(\mathcal O_S\)-module structure on
  \(\mathcal E\).}

\item[\textup{(ii)}]\phantomsection\label{concl:sga3-VIB-11.6.1-ii}
  \emph{If, moreover, \(\mathcal E\) is flat, every
  \(\mathcal O_S\)-linear action of \(G\) on
  \(\mathcal E\) factors through \(G_{\mathrm{af}}\) and
  corresponds to a unique \(\mathcal A(G)\)-comodule
  structure on \(\mathcal E\).}
\end{enumerate}

\medskip
\phantomsection\label{lem:sga3-VIB-11.7}%
\noindent\textbf{Lemma 11.7.---}
\emph{Let \(G\) be an \(S\)-group that is flat,
quasi-compact, and quasi-separated over \(S\), such that
\(\mathcal A=\mathcal A(G)\) is a flat
\(\mathcal O_S\)-module. Let \(\mathcal E\) be a
quasi-coherent \(\mathcal O_S\)-module, let}
\[
  \delta:\mathcal E\longrightarrow\mathcal E\otimes\mathcal A
\]
\emph{be an \(\mathcal A\)-comodule structure, and let}
\[
  \rho:G\longrightarrow
  \underline{\operatorname{Aut}}_{\mathcal O_S}
    \bigl(\mathrm W(\mathcal E)\bigr)
\]
\emph{be the associated action of \(G\) on \(\mathcal E\).
Let \(\mathcal E_0\) be a quasi-coherent
\(\mathcal O_S\)-submodule of \(\mathcal E\) such that the
restriction \(\delta_0\) of \(\delta\) to
\(\mathcal E_0\)}%
\begingroup
\renewcommand{\thefootnote}{\fnsymbol{footnote}}%
\footnote[1]{\label{note:sga3-VIB-11.7-delta-source}
Translator's note: the French source writes ``the restriction
\(\delta_0\) of \(\mu\)'' here. The surrounding notation and
the displayed square show that \(\delta\) is intended.}%
\endgroup
\emph{ factors through}
\[
  \mathcal E_0\longrightarrow\mathcal E_0\otimes\mathcal A;
\]
\emph{that is, suppose that the following diagram is
commutative:}
\begin{center}
  \phantomsection\label{diag:sga3-VIB-11.7-stability}%
  \begin{tikzcd}[row sep=large,column sep=large]
    \mathcal E_0
      \arrow[r,hook]
      \arrow[d,"\delta_0"'] &
    \mathcal E
      \arrow[d,"\delta"] \\
    \mathcal E_0\otimes\mathcal A
      \arrow[r,hook] &
    \mathcal E\otimes\mathcal A\mathpunct{.}
  \end{tikzcd}
\end{center}
\emph{\textup{(N.B.\ The morphism
\(\mathcal E_0\otimes\mathcal A
  \to\mathcal E\otimes\mathcal A\)
is injective because \(\mathcal A\) is flat over
\(\mathcal O_S\).)}
Then \(\delta_0\) makes \(\mathcal E_0\) an
\(\mathcal A(G)\)-comodule, and hence defines an action
\(\rho_0\) of \(G\) on \(\mathcal E_0\). This is called the
action on \(\mathcal E_0\) induced by \(\rho\), and we say
that \(\mathcal E_0\) is stable under \(\rho\).}

This follows immediately from the definitions and from
\hyperref[par:sga3-VIB-11.6]{11.6}. Notice, however, that in
general the canonical map
\[
  \mathrm W(\mathcal E_0)\longrightarrow\mathrm W(\mathcal E)
\]
is not a monomorphism.

\medskip
\phantomsection\label{rem:sga3-VIB-11.7.bis}%
\noindent\textbf{Remark 11.7.bis.---}%
\footnote{\label{note:sga3-VIB-111}Editor's note: we have
added this remark, which generalizes \SGARefLink{lem:sga3-VIB-11.7}{11.7} and will be useful
in \SGARefLink{prop:sga3-VIB-11.10.bis}{11.10.b}is.}
Let \(G\) be a flat \(S\)-group and let \(\mathcal E\) be a
quasi-coherent \(G\)-\(\mathcal O_S\)-module. Write \(f\) for
the morphism \(G\to S\) and \(\delta\) for the morphism
\[
  \mathcal E\longrightarrow f_*f^*(\mathcal E)
\]
defined in \hyperref[par:sga3-VIB-11.6]{11.6}. Let
\(\mathcal E_0\) be a quasi-coherent
\(\mathcal O_S\)-submodule of \(\mathcal E\). Since \(f\) is
flat, \(f_*f^*(\mathcal E_0)\) is an
\(\mathcal O_S\)-submodule of \(f_*f^*(\mathcal E)\), and the
same holds for \(\varphi=f\times f\). Consequently, if the
restriction \(\delta_0\) of \(\delta\) to \(\mathcal E_0\)
factors through \(f_*f^*(\mathcal E_0)\), then it makes
\(\mathcal E_0\) a \(G\)-\(\mathcal O_S\)-module. In this
case, we say that \(\mathcal E_0\) is a \(G\)-stable
submodule of \(\mathcal E\).

% END INCLUDED SGA3_English_current_buildable_source_20260823/source/inputs/VI/tex_reference_v2/components/83_expose_VIB_par114_through_rem117bis_loop2_r1_en.tex

% BEGIN INCLUDED SGA3_English_current_buildable_source_20260823/source/inputs/VI/tex_reference_v2/components/84_expose_VIB_def1180_through_cor1110_loop2_r1_en.tex
% SGA 3, Expose VI_B, Loop 1 English reconstruction.
% Loop 2 R1: the source-bound coassociativity crop is replaced by
% a faithful native tikz-cd diagram; prose, mathematics, and labels
% otherwise remain identical to the reviewed Loop-1 component.
% Scope: Definition 11.8.0 through Corollary 11.10, including
% the transition sentence immediately following Corollary 11.10.
% Authority: Polo--Gille Exp6B-13oct24.pdf, component-local
% pages 82--86 (combined-reader pages 423--427; re-edition
% headers 412--416). Proposition 11.10.bis begins on local
% page 87 / combined page 428 / header 417 and is excluded.
% Original-print margin markers do not advance one-for-one with
% re-edition pages: the body enters under marker 401 inherited
% from local page 81 and visibly crosses marker 402 on local 85.
% OCR and Reinhold English Markdown are comparison witnesses only.
% Loop 2 replaces the source-bound coassociativity crop with a
% faithful native LaTeX reconstruction; the replacement is closed.

\medskip
\phantomsection\label{def:sga3-VIB-11.8.0}%
\noindent\textbf{Definition 11.8.0.---}%
\footnote{\label{note:sga3-VIB-112}Editor's note: the
statements \SGARefLink{lem:sga3-VIB-11.8}{11.8} and \SGARefLink{prop:sga3-VIB-11.9}{11.9} concern only the notion of a
comodule over a coalgebra. We have therefore introduced
\hyperref[def:sga3-VIB-11.8.0]{Definition 11.8.0} and reformulated \hyperref[lem:sga3-VIB-11.8]{11.8} and \hyperref[prop:sga3-VIB-11.9]{11.9}
accordingly.}
Let \(S\) be a scheme. An \(\mathcal O_S\)-\emph{coalgebra}
is an \(\mathcal O_S\)-module \(\mathcal C\), together with
two morphisms of \(\mathcal O_S\)-modules
\[
  \Delta:\mathcal C\longrightarrow\mathcal C\otimes\mathcal C,
  \qquad
  \varepsilon:\mathcal C\longrightarrow\mathcal O_S,
\]
satisfying the following two axioms \textup{(cf.\ I, \SGARefLink{sga3:i:subsection:4-2}{4.2})}:
\begin{enumerate}
\item[\textup{(CO 1)}]\phantomsection\label{ax:sga3-VIB-11.8.0-CO1}
  \(\Delta\) is coassociative: the following diagram is
  commutative:
  \begin{center}
    \phantomsection\label{diag:sga3-VIB-11.8.0-coassociativity}%
    \begin{tikzcd}[row sep=large,column sep=large]
      & \mathcal C\otimes\mathcal C
        \arrow[dr,"\mathrm{id}\otimes\Delta"] & \\
      \mathcal C
        \arrow[ur,"\Delta"]
        \arrow[dr,"\Delta"'] & &
      \mathcal C\otimes\mathcal C\otimes\mathcal C \\
      & \mathcal C\otimes\mathcal C
        \arrow[ur,"\Delta\otimes\mathrm{id}"'] &
    \end{tikzcd}
  \end{center}

\item[\textup{(CO 2)}]\phantomsection\label{ax:sga3-VIB-11.8.0-CO2}
  \(\varepsilon\) is a counit; that is, the following two
  composites are the identity:
  \phantomsection\label{form:sga3-VIB-11.8.0-counit-right}%
  \[
    \mathcal C\xrightarrow{\Delta}
    \mathcal C\otimes\mathcal C
    \xrightarrow{\mathrm{id}\otimes\varepsilon}
    \mathcal C\otimes\mathcal O_S
    \xrightarrow{\sim}\mathcal C,
  \]
  \phantomsection\label{form:sga3-VIB-11.8.0-counit-left}%
  \[
    \mathcal C\xrightarrow{\Delta}
    \mathcal C\otimes\mathcal C
    \xrightarrow{\varepsilon\otimes\mathrm{id}}
    \mathcal O_S\otimes\mathcal C
    \xrightarrow{\sim}\mathcal C.
  \]
\end{enumerate}

A \emph{\(\mathcal C\)-comodule \textup{(on the right)}} is
an \(\mathcal O_S\)-module \(\mathcal E\), together with a
morphism of \(\mathcal O_S\)-modules
\[
  \mu:\mathcal E\longrightarrow\mathcal E\otimes\mathcal C
\]
satisfying axioms \textup{(CM 1)} and \textup{(CM 2)}
of \hyperref[par:sga3-VIB-11.6]{11.6}.

We say that \(\mathcal C\) \textup{(respectively,
\(\mathcal E\))} is a \emph{quasi-coherent coalgebra}
\textup{(respectively, a \emph{quasi-coherent comodule})} if
it is a quasi-coherent \(\mathcal O_S\)-module.

Let \(A\) be a commutative ring and let \(C\) be an
\(A\)-coalgebra. Then
\[
  C^\vee=\operatorname{Hom}_A(C,A)
\]
is an \(A\)-algebra. We denote by
\[
  \operatorname{ev}:C\otimes_A C^\vee\longrightarrow A
\]
the natural evaluation map.

\medskip
\phantomsection\label{lem:sga3-VIB-11.8}%
\noindent\textbf{Lemma 11.8.---}
\emph{Let \(C\) be an \(A\)-coalgebra, \(V\) a
\(C\)-comodule, and \(M\) an \(A\)-submodule of \(V\).
Suppose that \(C\) is a projective \(A\)-module.}%
\footnote{\label{note:sga3-VIB-113}Editor's note: in the
original, \(C\) is assumed to be a free \(A\)-module. The
generalization to the case where \(C\) is a projective
\(A\)-module, pointed out by J.-P.\ Serre, is mentioned in
\hyperref[rem:sga3-VIB-11.10.1]{Remark 11.10.1}. We have incorporated this generalization
here and in \hyperref[prop:sga3-VIB-11.9]{11.9}, and have expanded the proofs
accordingly.}
\emph{Let \(c(M)\) be the image of the morphism of
\(A\)-modules}
\phantomsection\label{form:sga3-VIB-11.8-theta}%
\[
  \theta:
  M\otimes_A C^\vee
  \xrightarrow{\ \mu\otimes\mathrm{id}\ }
  V\otimes_A C\otimes_A C^\vee
  \xrightarrow{\ \mathrm{id}\otimes\operatorname{ev}\ }
  V.
\]
\emph{Then \(c(M)\) is the smallest subcomodule of \(V\)
containing \(M\), and it is a finitely generated
\(A\)-module if \(M\) is. We call \(c(M)\) the
\emph{subcomodule generated by \(M\)}.}

\emph{Moreover, for every homomorphism of rings
\(A\to A'\), if \(M'\) denotes the image of
\(M\otimes_AA'\) in \(V'=V\otimes_AA'\), then \(c(M')\) is
the image of \(c(M)\otimes_AA'\) in \(V'\). Thus the
formation of \(c(M)\) commutes with base change.}

First, \(M\subset c(M)\) by \textup{(CM 2)}. If \(N\) is a
subcomodule of \(V\) containing \(M\), then
\(\mu(M)\subset N\otimes C\), and hence \(c(M)\subset N\).

By hypothesis, \(C\) is a direct summand of a free
\(A\)-module \(L\) with basis \((e_i)_{i\in I}\). Let
\(\varphi_i\) denote the restriction to \(C\) of the linear
form \(e_i^*\), defined by
\[
  e_i^*(e_j)=\delta_{ij}.
\]
Let \(x\in M\). We may write
\phantomsection\label{form:sga3-VIB-11.8-eq1}%
\[
  \tag{1}
  \mu(x)=\sum_{i\in J}x_i\otimes e_i,
\]
where \(x_i\in V\) and \(J\) is a finite subset of \(I\).
Then
\[
  x_i=\theta(x\otimes\varphi_i)\in c(Ax),
\]
and consequently
\[
  c(Ax)=\sum_{i\in J}Ax_i.
\]
Since \(C\) is a direct summand of \(L\), say
\(L=C\oplus R\), whence
\[
  V\otimes L=(V\otimes C)\oplus(V\otimes R),
\]
we obtain
\phantomsection\label{form:sga3-VIB-11.8-intersection}%
\[
  (c(Ax)\otimes L)\cap(V\otimes C)=c(Ax)\otimes C.
\]
Consequently, \(\mu(x)\) may also be written in the form
\phantomsection\label{form:sga3-VIB-11.8-eq2}%
\[
  \tag{2}
  \mu(x)=\sum_{j\in J}x_j\otimes b_j,
\]
with \(b_j\in C\). Write
\[
  \Delta(b_j)=\sum_{i\in I}b_{ij}\otimes e_i,
  \qquad b_{ij}\in C.
\]
Applying \(\mu\otimes\mathrm{id}\) to
\hyperref[form:sga3-VIB-11.8-eq1]{\textup{(1)}} and
\(\mathrm{id}\otimes\Delta\) to
\hyperref[form:sga3-VIB-11.8-eq2]{\textup{(2)}}, and using
axiom \textup{(CM 1)}, we obtain, for every \(i\in J\),
\phantomsection\label{form:sga3-VIB-11.8-coaction-generators}%
\[
  \mu(x_i)=\sum_{j\in J}x_j\otimes b_{ij}
    \in c(Ax)\otimes C.
\]
This proves that \(c(M)\) is a subcomodule of \(V\), and
therefore that it is the smallest subcomodule of \(V\)
containing \(M\).

It is clear that \(c(M)\) is a finitely generated
\(A\)-module whenever \(M\) is: if
\[
  M=Ax_1+\cdots+Ax_n,
  \qquad
  \mu(x_k)=\sum_i x_{ik}\otimes e_i,
\]
then the \(x_{ik}\), for \(k=1,\ldots,n\) and \(i\) ranging
over a finite subset of \(I\), generate \(c(M)\).

Finally, let \(A\to A'\) be a homomorphism of rings, and let
\(M'\) be the image of \(M\otimes_AA'\) in
\(V'=V\otimes_AA'\). Then \(c(M')\), respectively the image
of \(c(M)\otimes_AA'\) in \(V'\), is the image of the
morphism \(\theta'\) below, respectively of the composite
\(\theta'\circ\tau\):
\phantomsection\label{form:sga3-VIB-11.8-base-change-theta}%
\[
  M\otimes_AA'\otimes_A C^\vee
  \xrightarrow{\ \tau\ }
  M\otimes_A\operatorname{Hom}_A(C,A')
  \xrightarrow{\ \theta'\ }V'.
\]
These two morphisms have the same image. Indeed, let
\(\psi\in\operatorname{Hom}_A(C,A')\) and \(x\in M\), and
write
\[
  \mu(x)=\sum_{i\in J}x_i\otimes e_i.
\]
% correction SGA-SRC-000656: project the basis elements of L to C before applying psi.
Then
\[
  \theta'(x\otimes\psi)
    =\sum_{i\in J}\psi(p(e_i))x_i
\]
is the image under \(\theta'\circ\tau\) of the element
\[
  \sum_{i\in J}x\otimes\psi(p(e_i))\otimes\varphi_i
    \in M\otimes_AA'\otimes_A C^\vee.
\]
This proves the lemma.

We also have the following proposition.

\medskip
\phantomsection\label{prop:sga3-VIB-11.8-bis}%
\noindent\textbf{Proposition 11.8.bis.---}%
\footnote{\label{note:sga3-VIB-114}Editor's note: this
proposition, taken from \textup{[Se68]}, Section 1.5,
Proposition 2, has been added.}
\emph{Let \(A\) be a Noetherian ring, let \(C\) be an
\(A\)-coalgebra flat over \(A\), let \(V\) be a
\(C\)-comodule, and let \(M\) be a finitely generated
\(A\)-submodule of \(V\). Then there is a subcomodule \(W\)
of \(V\), finitely generated over \(A\), that contains
\(M\).}

Indeed, since \(M\) is finitely generated, so is
\(\Delta_V(M)\). Hence there is a finitely generated
\(A\)-submodule \(M'\) of \(V\) such that
\[
  \Delta_V(M)\subset M'\otimes_A C.
\]
Let \(\pi:V\to V/M'\) be the projection, put
\[
  \overline{\Delta}_V
    =(\pi\otimes\mathrm{id}_C)\Delta_V,
\]
and set
\phantomsection\label{form:sga3-VIB-11.8-bis-W}%
\[
  W=\{x\in V\mid\Delta_V(x)\in M'\otimes_A C\}
    =\ker\overline{\Delta}_V.
\]
This is an \(A\)-submodule of \(V\) that contains \(M\) and
is contained in \(M'\), because
\[
  x=(\mathrm{id}_V\otimes\varepsilon)\Delta_V(x).
\]
It is therefore finitely generated over \(A\). Moreover,
\[
  (\overline{\Delta}_V\otimes\mathrm{id}_C)\Delta_V
    =(\pi\otimes\Delta_C)\Delta_V
\]
vanishes on \(W\); in other words, \(\Delta_V(W)\) is
contained in the kernel \(K\) of
\(\overline{\Delta}_V\otimes\mathrm{id}_C\). Since \(C\) is
flat over \(A\), one has
\[
  K=W\otimes_A C.
\]
Thus \(W\) is a subcomodule of \(V\).

\medskip
\phantomsection\label{lem:sga3-VIB-11.8.1}%
\noindent\textbf{Lemma 11.8.1.---}%
\footnote{\label{note:sga3-VIB-115}Editor's note: this lemma,
which is used in the proof of \SGARefLink{prop:sga3-VIB-11.9}{11.9}, has been added.}
\emph{Let \(C\) be an \(A\)-coalgebra, let \(V\) be a
\(C\)-comodule, let \(M\) be an \(A\)-submodule of \(V\),
and let \(f:A\to A'\) be a faithfully flat homomorphism of
rings. Suppose that \(C'=C\otimes_AA'\) is a projective
\(A'\)-module.}
\begin{enumerate}
\item[\textup{(i)}]\phantomsection\label{concl:sga3-VIB-11.8.1-i}
  \emph{There is a smallest subcomodule \(t(M)\) of \(V\)
  containing \(M\), and \(t(M)\) is a finitely generated
  \(A\)-module if \(M\) is. Moreover, the formation of
  \(t(M)\) commutes with base change.}

\item[\textup{(ii)}]\phantomsection\label{concl:sga3-VIB-11.8.1-ii}
  \emph{More precisely, \(C\) is a projective
  \(A\)-module, and \(t(M)=c(M)\).}
\end{enumerate}

\emph{Proof.}
For \hyperref[concl:sga3-VIB-11.8.1-ii]{\textup{(ii)}},
\textup{[RG71]} \textup{(see \hyperref[prop:sga3-VIB-11.8.2]{Proposition 11.8.2} below)}
shows that \(C\) is a projective \(A\)-module. We may
therefore apply \hyperref[lem:sga3-VIB-11.8]{Lemma 11.8}: \(c(M)\) is the smallest
subcomodule of \(V\) containing \(M\), it is a finitely
generated \(A\)-module if \(M\) is, and its formation
commutes with base change.

To avoid an anachronism \textup{(\([RG71]\) postdates
SGA 3)}, let us sketch a direct proof of
\hyperref[concl:sga3-VIB-11.8.1-i]{\textup{(i)}}. Since
\(A\to A'\) is flat,
\[
  M'=M\otimes_AA'
\]
is an \(A'\)-submodule of \(V'=V\otimes_AA'\). Since
\(C'\) is a projective \(A'\)-module, \(c(M')\) is the
smallest subcomodule of \(V'\) containing \(M'\). Denote by
\(V'\otimes_AA'\) and \(A'\otimes_AV'\) the two
\(A'\otimes_AA'\)-comodule structures on
\[
  V''=V'\otimes_{A'}(A'\otimes_AA')
\]
obtained by the two base changes
\[
  A'\rightrightarrows A'\otimes_AA',
  \qquad
  a'\longmapsto a'\otimes1,
  \quad
  a'\longmapsto1\otimes a'.
\]
The \(A'\)-comodule \(V'\) carries an isomorphism of
\(A'\otimes_AA'\)-comodules
\phantomsection\label{form:sga3-VIB-11.8.1-descent-phi}%
\[
  \varphi:
  V'\otimes_AA'\xrightarrow{\ \sim\ }A'\otimes_AV',
  \qquad
  (x\otimes a')\otimes b'
    \longmapsto b'\otimes(x\otimes a'),
\]
which is a descent datum, that is,
\[
  \varphi_{31}=\varphi_{32}\circ\varphi_{21}.
\]

Since \(M'=M\otimes_AA'\), the isomorphism \(\varphi\)
takes \(M'\otimes_AA'\) onto \(A'\otimes_AM'\), and hence
takes \(c(M'\otimes_AA')\) onto \(c(A'\otimes_AM')\).
Because the formation of \(c(M')\) commutes with base
change,
\[
  c(M'\otimes_AA')=c(M')\otimes_AA',
  \qquad
  c(A'\otimes_AM')=A'\otimes_Ac(M').
\]
Thus
\phantomsection\label{form:sga3-VIB-11.8.1-descent-subcomodule}%
\[
  \varphi\bigl(c(M')\otimes_AA'\bigr)
    =A'\otimes_Ac(M'),
\]
and \(\varphi\) equips \(c(M')\) with a descent datum. By
fpqc descent, there is a unique subcomodule \(t(M)\) of
\(V\) such that
\[
  c(M')=t(M)\otimes_AA'.
\]
It contains \(M\), since \(t(M)\otimes_AA'\) contains
\(M'\). Moreover, if \(N\) is a subcomodule of \(V\)
containing \(M\), then \(N\) contains \(t(M)\), because
\(N\otimes_AA'\) contains
\[
  c(M')=t(M)\otimes_AA'.
\]
Hence \(t(M)\) is the smallest subcomodule of \(V\)
containing \(M\).

Finally, let \(A\to B\) be a homomorphism of rings. Put
\[
  B'=B\otimes_AA',
\]
and let \(M_B\), respectively \(M'_{B'}\), be the image of
\(M\otimes_AB\) in \(V_B=V\otimes_AB\), respectively of
\(M'\otimes_{A'}B'\) in
\(V'\otimes_{A'}B'=V\otimes_AB'\). Then
\[
  M_B\otimes_BB'=M'_{B'}.
\]
On the one hand, the preceding construction, applied to
\(C_B\) and \(B\to B'\), gives
\phantomsection\label{form:sga3-VIB-11.8.1-base-change-first}%
\[
  c(M_B\otimes_BB')
    =t(M_B)\otimes_BB'
    =t(M_B)\otimes_AA'.
\]
On the other hand, because the formation of \(c(M')\)
commutes with base change, \(c(M'_{B'})\) is the image in
\[
  V'\otimes_{A'}B'=V\otimes_AB\otimes_AA'
\]
of
\phantomsection\label{form:sga3-VIB-11.8.1-base-change-second}%
\[
  c(M')\otimes_{A'}B'
    =t(M)\otimes_AB\otimes_AA'.
\]
It follows that \(t(M_B)\) is the image of
\(t(M)\otimes_AB\) in \(V_B\).

\medskip
\phantomsection\label{prop:sga3-VIB-11.8.2}%
\noindent\textbf{Proposition 11.8.2 \textup{(Gruson--Raynaud)}.---}
\emph{Let \(f:A\to A'\) be a faithfully flat homomorphism.
Then \(f\) ``descends projectivity'': if \(M\) is an
\(A\)-module and \(M\otimes_AA'\) is a projective
\(A'\)-module, then \(M\) is a projective \(A\)-module.}

Indeed, by \textup{[RG71]}, II, \SGARefLink{lem:sga3-VIA-2.5.1}{2.5.1}, \(f\) descends the
Mittag--Leffler condition, and therefore, by
\emph{loc.\ cit.}, II, 3.1.3, \(f\) descends
projectivity.%
\footnote{\label{note:sga3-VIB-116}Editor's note: let us
also point out that assertion II, \SGARefLink{prop:sga3-VIA-2.5.2}{2.5.2} of \emph{loc.\ cit.},
which is more general than II, \SGARefLink{lem:sga3-VIA-2.5.1}{2.5.1}, is corrected in
\textup{[Gr73]}; this does not affect the case of faithfully
flat homomorphisms.}

\medskip
\phantomsection\label{prop:sga3-VIB-11.9}%
\noindent\textbf{Proposition 11.9.---}%
\footnote{\label{note:sga3-VIB-117}Editor's note: as
indicated in editor's \hyperref[note:sga3-VIB-112]{note 112}, we have rewritten the
statement for an \(\mathcal O_S\)-coalgebra \(\mathcal C\)
\textup{(rather than for an \(S\)-group \(G\) satisfying the
hypotheses in \hyperref[cor:sga3-VIB-11.10]{Corollary 11.10})}. We have also expanded the
proof \textup{(the original merely said: ``\(\cdots\) the
proposition follows from \hyperref[lem:sga3-VIB-11.8]{Lemma 11.8}'')}.}
\emph{Let \(\mathcal C\) be an \(\mathcal O_S\)-coalgebra,
let \(\mathcal E\) be a \(\mathcal C\)-comodule, and let
\(\mathcal F\) be an \(\mathcal O_S\)-submodule of
\(\mathcal E\), all of them quasi-coherent. Suppose that
\(S\) is covered by affine open subsets
\(U_\alpha=\operatorname{Spec}A_\alpha\), and that for each
\(\alpha\) there is a faithfully flat homomorphism of rings
\(A_\alpha\to A'_\alpha\) such that}
\[
  \Gamma(U_\alpha,\mathcal C)\otimes_{A_\alpha}A'_\alpha
\]
\emph{is a projective \(A'_\alpha\)-module.}%
\textsuperscript{\textup{(*)}}
\begingroup
\renewcommand{\thefootnote}{*}%
\footnotetext{This holds, for example, when
\(\mathcal C=\mathcal A(G)\), where \(G\) is a reductive
\(S\)-group, as we shall see in \SGARefLink{sga3:xxii:cor:5.7.8}{Exposé XXII, 5.7.8}.}
\endgroup

\emph{Then there is a smallest quasi-coherent subcomodule
\(t(\mathcal F)\) of \(\mathcal E\) containing
\(\mathcal F\), and \(t(\mathcal F)\) is a finitely
generated \(\mathcal O_S\)-module if \(\mathcal F\) is.
Moreover, for every base change \(S'\to S\), if
\(\mathcal F'\) denotes the image of
\(\mathcal F\otimes_{\mathcal O_S}\mathcal O_{S'}\) in}
\[
  \mathcal E'
    =\mathcal E\otimes_{\mathcal O_S}\mathcal O_{S'},
\]
\emph{then \(t(\mathcal F')\) is the image of
\(t(\mathcal F)\otimes_{\mathcal O_S}\mathcal O_{S'}\) in
\(\mathcal E'\). In other words, the formation of
\(t(\mathcal F)\) commutes with base change.}

\emph{Proof.}
For each \(\alpha\), let \(T_\alpha\) be the
\(A_\alpha\)-module
\[
  T_\alpha=t\bigl(\Gamma(U_\alpha,\mathcal F)\bigr).
\]
By \hyperref[concl:sga3-VIB-11.8.1-i]{11.8.1}\textup{\SGARefLink{concl:sga3-VIB-11.8.1-i}{(i)}}, the associated
\(\mathcal O_{U_\alpha}\)-module
\(\widetilde T_\alpha\) is the smallest quasi-coherent
subcomodule of \(\mathcal E|_{U_\alpha}\) containing
\(\mathcal F|_{U_\alpha}\); and it is a finitely generated
\(\mathcal O_{U_\alpha}\)-module if
\(\mathcal F|_{U_\alpha}\) is.

For every \(\alpha,\beta\), put
\[
  U_{\alpha\beta}=U_\alpha\cap U_\beta.
\]
Since the construction of \(t(M)\) commutes with base
change, there are canonical isomorphisms of
\(\mathcal O_{U_{\alpha\beta}}\)-modules
\phantomsection\label{form:sga3-VIB-11.9-gluing}%
\[
  \varphi_{\alpha\beta}:
  \widetilde T_\beta
    \otimes_{A_\beta}\mathcal O_{U_{\alpha\beta}}
  \xrightarrow{\ \sim\ }
  \widetilde T_\alpha
    \otimes_{A_\alpha}\mathcal O_{U_{\alpha\beta}}.
\]
They satisfy the cocycle condition
\[
  \varphi'_{\alpha\gamma}
    =\varphi'_{\alpha\beta}\circ\varphi'_{\beta\gamma},
\]
where \(\varphi'_{\alpha\gamma}\), respectively the other
maps, denotes the restriction of
\(\varphi_{\alpha\gamma}\), respectively the other maps, to
\(U_\alpha\cap U_\beta\cap U_\gamma\).

Consequently, the \(\widetilde T_\alpha\) glue to a
quasi-coherent subcomodule \(t(\mathcal F)\) of
\(\mathcal E\) containing \(\mathcal F\). We leave it to the
reader to verify that \(t(\mathcal F)\) is the smallest
quasi-coherent subcomodule of \(\mathcal E\) containing
\(\mathcal F\), and that its formation commutes with base
change.

\medskip
\phantomsection\label{def:sga3-VIB-11.9.1}%
\noindent\textbf{Definition 11.9.1.---}%
\footnote{\label{note:sga3-VIB-118}Editor's note: this
definition, taken from the bottom of page 82 of
\textup{[RG71]}, has been added. Thus the hypothesis in
\hyperref[prop:sga3-VIB-11.9]{Proposition 11.9} is that the coalgebra \(\mathcal C\) be a
locally projective \(\mathcal O_S\)-module, and we have used
this terminology in \hyperref[cor:sga3-VIB-11.10]{Corollary 11.10}.}
\emph{Let \(S\) be a scheme and let \(\mathcal P\) be a
quasi-coherent \(\mathcal O_S\)-module. The following
conditions are equivalent:}
\begin{enumerate}
\item[\textup{(i)}]\phantomsection\label{concl:sga3-VIB-11.9.1-i}
  \emph{For every affine open subset \(U\) of \(S\), the
  module \(\Gamma(U,\mathcal P)\) is a projective
  \(\mathcal O_S(U)\)-module.}

\item[\textup{(ii)}]\phantomsection\label{concl:sga3-VIB-11.9.1-ii}
  \emph{There is a covering \((U_\alpha)\) of \(S\) by
  affine open subsets such that every
  \(\Gamma(U_\alpha,\mathcal P)\) is a projective
  \(\mathcal O_S(U_\alpha)\)-module.}

\item[\textup{\SGAInlineRefTarget{def:sga3-VIB-11.9.1:item:iii}(iii)}]\phantomsection\label{concl:sga3-VIB-11.9.1-iii}
  \emph{There is a covering \((U_\alpha)\) of \(S\) by
  affine open subsets, and faithfully flat homomorphisms of
  rings}
  \[
    A_\alpha=\mathcal O_S(U_\alpha)\longrightarrow A'_\alpha,
  \]
  \emph{such that, for every \(\alpha\),}
  \[
    \Gamma(U_\alpha,\mathcal P)
      \otimes_{A_\alpha}A'_\alpha
  \]
  \emph{is a projective \(A'_\alpha\)-module.}
\end{enumerate}

It is clear that
\[
  \textup{(i)}\Longrightarrow\textup{(ii)}
  \Longrightarrow\textup{(iii)}.
\]
Conversely, if \textup{\SGARefLink{def:sga3-VIB-11.9.1:item:iii}{(iii)}} holds, \hyperref[prop:sga3-VIB-11.8.2]{11.8.2} implies that
every \(\Gamma(U_\alpha,\mathcal P)\) is a projective
\(\mathcal O_S(U_\alpha)\)-module, and hence
\textup{(ii)} holds.

Finally, suppose that \textup{\SGARefLink{concl:sga3-VIB-11.9.1-ii}{(ii)}} holds, and let
\[
  V=\operatorname{Spec}A
\]
be an arbitrary affine open subset. It is covered by
finitely many affine open subsets
\[
  V_1,\ldots,V_n,
  \qquad
  V_i=\operatorname{Spec}A_i,
\]
each contained in at least one \(V\cap U_\alpha\), so that
\(\Gamma(V_i,\mathcal P)\) is a projective
\(A_i\)-module. Put
\[
  A'=A_1\times\cdots\times A_n.
\]
Then \(A\to A'\) is faithfully flat and
\[
  \Gamma(V,\mathcal P)\otimes_AA'
\]
is a projective \(A'\)-module. Therefore \SGARefLink{prop:sga3-VIB-11.8.2}{11.8.2} implies that
\(\Gamma(V,\mathcal P)\) is a projective \(A\)-module.

When these equivalent conditions hold, we say that
\(\mathcal P\) is a \emph{locally projective}
\(\mathcal O_S\)-module.

\medskip
\phantomsection\label{cor:sga3-VIB-11.10}%
\noindent\textbf{Corollary 11.10.---}
\emph{Let \(S\) be a quasi-compact and quasi-separated
scheme, let \(G\) be an \(S\)-group, and let \(\rho\) be a
linear action of \(G\) on a quasi-coherent
\(\mathcal O_S\)-module \(\mathcal E\). Suppose that:}
\begin{enumerate}
\item[\textup{(i)}]\phantomsection\label{hyp:sga3-VIB-11.10-i}
  \emph{\(G\) satisfies one of the following conditions:}
  \begin{enumerate}
  \item[\textup{(a)}]\phantomsection\label{hyp:sga3-VIB-11.10-i-a}
    \emph{\(G\) is affine and flat over \(S\);}

  \item[\textup{(b)}]\phantomsection\label{hyp:sga3-VIB-11.10-i-b}
    \emph{\(G\) is flat, quasi-compact, and quasi-separated
    over \(S\), and \(\mathcal E\) is flat;}
  \end{enumerate}

\item[\textup{(ii)}]\phantomsection\label{hyp:sga3-VIB-11.10-ii}
  \emph{\(\mathcal A(G)\) is a locally projective
  \(\mathcal O_S\)-module.}
\end{enumerate}
\emph{Then \(\mathcal E\) is the direct limit of an
increasing filtered family of finitely generated
quasi-coherent \(\mathcal O_S\)-submodules of
\(\mathcal E\) that are stable under \(G\).}

By hypothesis
\hyperref[hyp:sga3-VIB-11.10-i]{\textup{(i)}} and \hyperref[prop:sga3-VIB-11.6.1]{11.6.1},
\(\mathcal E\) is endowed with an
\(\mathcal A(G)\)-comodule structure. On the other hand,
since \(S\) is quasi-compact and quasi-separated,
\(\mathcal E\) is the direct limit of its finitely
generated quasi-coherent submodules
\textup{(EGA I, 9.4.9 and EGA IV\(_1\), 1.7.7)}.
Consequently, the corollary follows from \hyperref[prop:sga3-VIB-11.9]{Proposition 11.9},
applied to the coalgebra \(\mathcal A(G)\).

We also have the following proposition.

% END INCLUDED SGA3_English_current_buildable_source_20260823/source/inputs/VI/tex_reference_v2/components/84_expose_VIB_def1180_through_cor1110_loop2_r1_en.tex

% BEGIN INCLUDED SGA3_English_current_buildable_source_20260823/source/inputs/VI/tex_reference_v2/components/85_expose_VIB_prop1110bis_through_rem11111_loop2_r1_en.tex
% SGA 3, Expose VI_B, Loop 1 English reconstruction.
% R2: independent visual review restored the intersentence
% space after the local symbolic translator's footnote.
% R3: independent PDF-link replay assigned distinct anchors
% to the eleven enumerated conclusion targets.
% Scope: Proposition 11.10bis, Remark 11.10.1,
% Proposition 11.11, and Remark 11.11.1, including
% editor's notes 119--124.
% Authority: Polo--Gille Exp6B-13oct24.pdf, component-local
% pages 87--90 (combined-reader pages 428--431; re-edition
% headers 417--420; original-margin printed pages 402--405).
% Lemma 11.12 begins on local page 90 and is excluded.
% OCR is locator-only; Reinhold e7a259f English Markdown is
% comparison-only and is credited in the component report.
% Loop 2 replaces all three source-bound commutative-diagram crops
% with faithful native LaTeX reconstructions; all are closed.

\medskip
\phantomsection\label{prop:sga3-VIB-11.10.bis}%
\noindent\textbf{Proposition 11.10.bis.---}%
\footnote{\label{note:sga3-VIB-119}Editor's note: this
proposition has been added; it is a special case of
\textup{[Th87]}, \SGARefLink{prop:sga3-VIB-1.4}{1.4}--\SGARefLink{lem:sga3-VIB-1.5}{1.5}. The author there refers to an
argument of Deligne \textup{(cf.\ \textup{[Kn71]},
III, Th.\ \SGARefLink{par:sga3-VIB-1.1}{1.1})}; one may also note its similarity with
Serre's argument \textup{(\textup{[Se68]}, Prop.\ 2)}
recalled in \hyperref[prop:sga3-VIB-11.8-bis]{11.8.bis}.}
\emph{Let \(S\) be a Noetherian scheme, let \(G\) be an
\(S\)-group that is flat, quasi-compact, and quasi-separated
over \(S\), let \(\mathcal E\) be a quasi-coherent
\(G\)-\(\mathcal O_S\)-module, and let \(\mathcal M\) be a
coherent \(\mathcal O_S\)-submodule of \(\mathcal E\). Then
\(\mathcal M\) is contained in a coherent
\(\mathcal O_S\)-submodule stable under \(G\).}

\emph{Proof.}
Write \(f\) for the morphism \(G\to S\), and \(\tau\) for
the adjunction morphism
\[
  \mathcal E\longrightarrow f_*f^*(\mathcal E).
\]
By \hyperref[par:sga3-VIB-11.6]{11.6}, the
\(G\)-\(\mathcal O_S\)-module structure on \(\mathcal E\)
is given by an automorphism \(\theta\) of the
\(\mathcal O_G\)-module \(f^*(\mathcal E)\), such that the
morphism
\[
  \delta=\theta\circ\tau
\]
satisfies conditions
\hyperref[diag:sga3-VIB-11.6-1]{\textup{(1)}} and
\hyperref[diag:sga3-VIB-11.6-2]{\textup{(2)}} of
\hyperref[par:sga3-VIB-11.6]{11.6}. (The situation considered
in \textup{[Th87]} is more general: the author considers a
\(G\)-scheme \(X\) and a \(G\)-equivariant
\(\mathcal O_X\)-module \(\mathcal E\)
\textup{(cf.\ \SGARefLink{sga3:i:section:6}{Expos\'e I, Section 6})}; here \(X=S\), endowed
with the trivial action of \(G\).)

Since \(S\) is Noetherian, \(\mathcal E\) is the filtered
direct limit of its coherent submodules
\(\mathcal F_\alpha\) \textup{(cf.\ EGA~I, 9.4.9)}. Then
\(f^*(\mathcal E)\) is the filtered direct limit of the
\(f^*(\mathcal F_\alpha)\), which are submodules of
\(f^*(\mathcal E)\) because \(f\) is flat. Moreover, since
\(f\) is quasi-compact and quasi-separated, by
\hyperref[par:sga3-VIB-11.0]{11.0} the
\(\mathcal O_S\)-module \(f_*f^*(\mathcal E)\) is
quasi-coherent and is the filtered direct limit of the
quasi-coherent submodules \(f_*f^*(\mathcal F_\alpha)\).
Consequently, \(\mathcal E\) is the filtered direct limit of
the quasi-coherent \(\mathcal O_S\)-submodules
\(\mathcal E_\alpha\), where \(\mathcal E_\alpha\) denotes
the inverse image under \(\delta\) of
\(f_*f^*(\mathcal F_\alpha)\); equivalently, it is the kernel
of the composite
\phantomsection\label{form:sga3-VIB-11.10bis-delta-alpha}%
\[
  \overline{\delta}_\alpha:\quad
  \mathcal E\xrightarrow{\ \delta\ }f_*f^*(\mathcal E)
    \longrightarrow f_*f^*(\mathcal E/\mathcal F_\alpha).
\]
Since \(\mathcal M\) is coherent and \(S\) is Noetherian,
every increasing sequence of submodules of \(\mathcal M\)
is stationary; hence \(\mathcal M\) is contained in some
\(\mathcal E_\alpha\). We shall show that every
\(\mathcal E_\alpha\) is coherent and \(G\)-stable.

Let
\[
  u:f^*(\mathcal E)\longrightarrow\varepsilon_*(\mathcal E)
\]
be the morphism corresponding by adjunction to the identity
morphism from \(\mathcal E\) to
\(f_*\varepsilon_*(\mathcal E)=\mathcal E\). Then the
composite
\phantomsection\label{form:sga3-VIB-11.10bis-adjunction-retract}%
\[
  \mathcal E\xrightarrow{\ \delta\ }f_*f^*(\mathcal E)
    \xrightarrow{\ f_*(u)\ }f_*\varepsilon_*(\mathcal E)
    =\mathcal E
\]
is the identity%
\footnote{\label{source-note:sga3-VIB-11.10bis-delta-arrows}The
source PDF labels the first arrow in this composite and the lower
horizontal arrow in the following diagram by \(\tau\). They must be
\(\delta=\theta\circ\tau\): condition 11.6\textup{(2)} gives
\(f_*(u)\circ\delta=\operatorname{id}_{\mathcal E}\), and the definition
of \(\mathcal E_\alpha\) makes the square commute.}
\textup{(cf.\ \hyperref[diag:sga3-VIB-11.6-2]{11.6\textup{(2)}})}.
We have the commutative diagram
\begin{center}
  \phantomsection\label{diag:sga3-VIB-11.10bis-coherence}%
  \begin{tikzcd}[row sep=large,column sep=large]
    \mathcal E_\alpha
      \arrow[r,"\delta"]
      \arrow[d,"i_\alpha"'] &
    f_*f^*(\mathcal F_\alpha)
      \arrow[r,"f_*(u)"]
      \arrow[d] &
    \mathcal F_\alpha
      \arrow[d,"j_\alpha"] \\
    \mathcal E
      \arrow[r,"\delta"'] &
    f_*f^*(\mathcal E)
      \arrow[r,"f_*(u)"'] &
    \mathcal E\mathpunct{,}
  \end{tikzcd}
\end{center}
where \(i_\alpha\) \textup{(respectively \(j_\alpha\))}
denotes the inclusion of \(\mathcal E_\alpha\)
\textup{(respectively \(\mathcal F_\alpha\))} in
\(\mathcal E\). It follows that \(i_\alpha\) factors through
\(j_\alpha\); thus \(\mathcal E_\alpha\) is a submodule of
\(\mathcal F_\alpha\), and is therefore coherent.

Finally, let us show that \(\mathcal E_\alpha\) is
\(G\)-stable
\textup{(cf.\ \hyperref[rem:sga3-VIB-11.7.bis]{11.7.bis})}.
Let \(h:G\to S\) denote a second copy of \(f:G\to S\), and
consider the following commutative diagram:
\begin{center}
  \phantomsection\label{diag:sga3-VIB-11.10bis-projections}%
  \begin{tikzcd}[row sep=large,column sep=large]
    G\times_S G
      \arrow[r,"q"]
      \arrow[d,"p"']
      \arrow[dr,"\varphi" description] &
    G \arrow[d,"h"] \\
    G \arrow[r,"f"'] &
    S\mathpunct{.}
  \end{tikzcd}
\end{center}
The exact sequence
\phantomsection\label{seq:sga3-VIB-11.10bis-first-exact}%
\[
  0\longrightarrow\mathcal E_\alpha
    \longrightarrow\mathcal E
    \xrightarrow{\ \overline{\delta}_\alpha\ }
    h_*h^*(\mathcal E/\mathcal F_\alpha)
\]
therefore gives, since \(f\) is flat, the exact sequence
\phantomsection\label{seq:sga3-VIB-11.10bis-dagger}%
\[
  \tag{\(\dagger\)}
  0\longrightarrow f_*f^*(\mathcal E_\alpha)
    \longrightarrow f_*f^*(\mathcal E)
    \xrightarrow{\ f_*f^*(\overline{\delta}_\alpha)\ }
    f_*f^*h_*h^*(\mathcal E/\mathcal F_\alpha).
\]
Moreover, since \(h\) is quasi-compact and quasi-separated
and \(f\) is flat, the canonical morphism
\phantomsection\label{form:sga3-VIB-11.10bis-flat-base-change}%
\[
  f^*h_*h^*(\mathcal E/\mathcal F_\alpha)
    \longrightarrow
  p_*q^*h^*(\mathcal E/\mathcal F_\alpha)
    =p_*\varphi^*(\mathcal E/\mathcal F_\alpha)
\]
is an isomorphism. Thus the right-hand term in
\hyperref[seq:sga3-VIB-11.10bis-dagger]{\textup{(\(\dagger\))}}
is \(\varphi_*\varphi^*(\mathcal E/\mathcal F_\alpha)\).
Using the notation of \hyperref[par:sga3-VIB-11.6]{11.6}
and writing \(\pi_\alpha\) for the projection
\(\mathcal E\to\mathcal E/\mathcal F_\alpha\), we obtain the
following commutative diagram, whose bottom row is exact:
\begin{center}
  \phantomsection\label{diag:sga3-VIB-11.10bis-stability}%
  {\small
  \begin{tikzcd}[row sep=large,column sep=large]
    {} &
    \mathcal E_\alpha
      \arrow[r,"\delta"]
      \arrow[d,"\delta"'] &
    f_*f^*(\mathcal E)
      \arrow[d,"{f_*(f^*(\pi_\alpha)\boxtimes\delta_G)}"] \\
    f_*f^*(\mathcal E_\alpha)
      \arrow[r] &
    f_*f^*(\mathcal E)
      \arrow[r,"{f_*f^*(\overline{\delta}_\alpha)}"'] &
    \varphi_*\varphi^*(\mathcal E/\mathcal F_\alpha)
  \end{tikzcd}}
\end{center}
Since
\[
  f_*\bigl(f^*(\pi_\alpha)\boxtimes\delta_G\bigr)\circ\delta
\]
vanishes on \(\mathcal E_\alpha\), it follows that
\(\delta\) carries \(\mathcal E_\alpha\) into
\(f_*f^*(\mathcal E_\alpha)\); in other words,
\(\mathcal E_\alpha\) is \(G\)-stable. This proves the
proposition.

\medskip
\phantomsection\label{rem:sga3-VIB-11.10.1}%
\noindent\textbf{Remark 11.10.1.---}%
\footnote{\label{note:sga3-VIB-120}Editor's note: the first
part of the original \hyperref[rem:sga3-VIB-11.10.1]{Remark 11.10.1} has been incorporated
into \hyperref[lem:sga3-VIB-11.8]{11.8} and
\hyperref[prop:sga3-VIB-11.9]{11.9}
\textup{(with ``free'' replaced by ``projective'')}; the
counterexamples below correct its second part.}
In \hyperref[lem:sga3-VIB-11.8]{11.8}, it is not enough to
assume that \(C\) is a flat \(A\)-module, even when \(A\)
is a principal ideal domain. Indeed, we have the following
counterexamples, pointed out to us independently by
O.\ Gabber and J.-P.\ Serre.

\smallskip
\phantomsection\label{rem:sga3-VIB-11.10.1-a}%
\noindent\textup{(a)}
Let \((A,\mathfrak m)\) be a discrete valuation ring, let
\(K\) be its field of fractions, and let \(G\) be the
\(A\)-group obtained by ``extension by zero'' from the
\(K\)-group \((\mathbb Z/2\mathbb Z)_K\). The constant
group \((\mathbb Z/2\mathbb Z)_A\), and hence also its
subgroup \(G\), acts on the free \(A\)-module \(V\) with
basis \(v_1,v_2\) by interchanging \(v_1\) and \(v_2\).
Then \(Av_1\) is not a \(G\)-submodule of \(V\), but it is
the intersection of the \(G\)-submodules
\[
  Av_1+\mathfrak m^n v_2,\qquad n\in\mathbb N^*.
\]
Thus there is no smallest \(G\)-submodule of \(V\)
containing \(v_1\).

\smallskip
\phantomsection\label{rem:sga3-VIB-11.10.1-b}%
\noindent\textup{(b)}
Let \(A\) be an integral domain distinct from its field of
fractions \(K\), and let \(G\) be the flat affine
\(A\)-group corresponding to the Hopf algebra
\phantomsection\label{form:sga3-VIB-11.10.1-Hopf-example}%
\[
  \mathcal A(G)=\{P\in K[T]\mid P(0)\in A\}.
\]
Its comultiplication, counit, and antipode are respectively
defined by
\[
  \Delta(T)=T\otimes1+1\otimes T,\qquad
  \varepsilon(T)=0,\qquad
  \tau(T)=-T.
\]
\textup{(Thus \(G\otimes_AK=\mathbb G_{a,K}\).)}

Let \(V=Av_1\oplus Av_2\) be free over \(A\), and let
\(u\) be the endomorphism of \(V\) defined by
\[
  u(v_1)=v_2,\qquad u(v_2)=0,
\]
so that \(u^2=0\). Then \(V\) has the
\(\mathcal A(G)\)-comodule structure
\phantomsection\label{form:sga3-VIB-11.10.1-coaction}%
\[
  \mu(m)=m\otimes1+u(m)\otimes T.
\]
The \(G\)-submodules of \(V\) containing \(v_1\) are
precisely the \(A\)-submodules \(Av_1\oplus Iv_2\), where
\(I\) is a nonzero ideal of \(A\). Their intersection is
\(Av_1\), which is not a \(G\)-submodule. Hence there is no
smallest \(G\)-submodule of \(V\) containing \(v_1\).
Moreover,
\[
  C=A\oplus K\cdot T
\]
is a subcoalgebra of \(\mathcal A(G)\), flat over \(A\),
and the coaction
\[
  \mu:V\longrightarrow V\otimes\mathcal A(G)
\]
factors through \(V\otimes C\). Thus the ``very simple''
coalgebra \(C\) also gives a counterexample.

Finally, the preceding two examples are special cases of
the following construction. Let \(A\) be an integral
domain distinct from its field of fractions \(K\), and let
\(B\) be a Hopf \(A\)-algebra, free over \(A\). Write
\[
  \varepsilon:B\longrightarrow A,\qquad
  I=\ker(\varepsilon)
\]
for the augmentation and the augmentation ideal. Since
\(B=A\cdot1\oplus I\), one readily sees that
\phantomsection\label{form:sga3-VIB-11.10.1-B-prime}%
\[
  B'=\{b\in B\otimes_AK\mid \varepsilon(b)\in A\}
\]
is a Hopf subalgebra of \(B_K\). Suppose that \(V\) is a
\(B\)-comodule, free with basis
\((v_1,\ldots,v_n)\) as an \(A\)-module, and that
\(\mu_V(v_1)\neq v_1\otimes1\). Then \(Av_1\) is not a
subcomodule of \(V\), but it is the intersection of the
subcomodules
\[
  Av_1+J\cdot V
\]
as \(J\) ranges over the nonzero ideals of \(A\).%
\begingroup
\renewcommand{\thefootnote}{\fnsymbol{footnote}}%
\footnote[1]{\label{note:sga3-VIB-11.10.1-ideal-letter}
Translator's note: the French re-edition uses \(I\) both
for the fixed augmentation ideal above and for the ideals
ranging in this final intersection. We use \(J\) for the
latter to keep the two roles distinct.}%
\endgroup\ 
Thus there is no smallest subcomodule of \(V\) containing
\(v_1\).

\medskip
\phantomsection\label{prop:sga3-VIB-11.11}%
\noindent\textbf{Proposition 11.11.---}
\emph{Let \(G\) be an algebraic group over a field \(k\).
The following conditions are equivalent:}
\begin{enumerate}
\item[\textup{(i)}]\phantomsection\label{concl:sga3-VIB-11.11-i}
  \emph{\(G\) is affine.}
\item[\textup{(ii)}]\phantomsection\label{concl:sga3-VIB-11.11-ii}
  \emph{\(G\) is quasi-affine.}
\item[\textup{(iii)}]\phantomsection\label{concl:sga3-VIB-11.11-iii}
  \emph{\(G\) acts faithfully on a quasi-affine
  \(k\)-scheme \(X\).}
\item[\textup{(iv)}]\phantomsection\label{concl:sga3-VIB-11.11-iv}
  \emph{\(G\) acts linearly and faithfully on a
  \(k\)-vector space \textup{(not necessarily
  finite-dimensional)}.}
\item[\textup{(v)}]\phantomsection\label{concl:sga3-VIB-11.11-v}
  \emph{\(G\) is isomorphic to a closed subgroup of
  \(\operatorname{GL}(n)_k\).}
\end{enumerate}

\emph{Proof.}
The implication
\hyperref[concl:sga3-VIB-11.11-i]{\textup{(i)}}
\(\Rightarrow\)
\hyperref[concl:sga3-VIB-11.11-ii]{\textup{(ii)}} is
immediate, and
\hyperref[concl:sga3-VIB-11.11-ii]{\textup{(ii)}}
\(\Rightarrow\)
\hyperref[concl:sga3-VIB-11.11-iii]{\textup{(iii)}} because
\(G\) acts faithfully on itself by translations.

Suppose that \(G\) acts faithfully on the right on a
quasi-affine \(k\)-scheme \(X\).%
\footnote{\label{note:sga3-VIB-121}Editor's note: the
original has been expanded in what follows. We have also
chosen a right action of \(G\) on \(X\), so as to obtain a
left linear action of \(G\) on \(\mathcal O(X)\).}
Since \(X\) is quasi-affine, it is separated and
quasi-compact.%
\footnote{\label{note:sga3-VIB-122}Editor's note: recall
that a \(k\)-scheme \(X\) is called quasi-affine if it is
isomorphic to a quasi-compact open subscheme of an affine
\(k\)-scheme \textup{(EGA~II, 5.1.1)}.}
Similarly, \(G\) is separated
\textup{(\hyperref[sga3:VIA:0.3]{VI\(_A\), 0.3})} and
quasi-compact \textup{(since it is of finite type over
\(k\))}. Hence, by
\hyperref[concl:sga3-VIB-11.1-c]{11.1\textup{(c)}}, there
are canonical isomorphisms
\phantomsection\label{form:sga3-VIB-11.11-functions-products}%
\[
  \mathcal O(X\times G)
    =\mathcal O(X)\otimes\mathcal O(G),
  \qquad
  \mathcal O(X\times G\times G)
    =\mathcal O(X)\otimes\mathcal O(G)\otimes\mathcal O(G).
\]
It follows that the morphism
\[
  \mu:\mathcal O(X)\longrightarrow
    \mathcal O(X)\otimes\mathcal O(G)
\]
induced by \(X\times G\to X\) gives
\(\mathcal O(X)\) the structure of a right
\(\mathcal O(G)\)-comodule; that is, \(G\) acts linearly on
the left on the \(k\)-algebra \(\mathcal O(X)\).
Consequently, \(G\) also acts on the right on the affine
envelope
\[
  X_{\mathrm{af}}=\operatorname{Spec}\mathcal O(X)
\]
of \(X\), and the canonical morphism
\(X\to X_{\mathrm{af}}\) is \(G\)-equivariant.

Since \(X\) is quasi-affine, the morphism
\(X\to X_{\mathrm{af}}\) is an open immersion
\textup{(EGA~II, 5.1.2)}. Hence, a fortiori, \(G\) acts
faithfully on \(X_{\mathrm{af}}\), and therefore the left
linear action of \(G\) on the \(k\)-algebra
\(\mathcal O(X)=\mathcal O(X_{\mathrm{af}})\) is faithful.
This proves
\hyperref[concl:sga3-VIB-11.11-iii]{\textup{(iii)}}
\(\Rightarrow\)
\hyperref[concl:sga3-VIB-11.11-iv]{\textup{(iv)}}.

Now suppose that \(G\) acts faithfully on a \(k\)-vector
space \(V\). By \hyperref[cor:sga3-VIB-11.10]{11.10},
\(V\) is the direct limit of finite-dimensional vector
subspaces \(V_i\) stable under \(G\). Let \(K_i\) be the
kernel of the induced action of \(G\) on \(V_i\); that is,
the kernel of \(G\to\operatorname{Aut}(V_i)\). Then \(K_i\)
is a closed subscheme of \(G\), and the faithfulness of the
action of \(G\) means that the intersection of the \(K_i\)
is the identity subgroup of \(G\). Since \(G\) is
Noetherian, one of the \(K_i\) is already the identity
subgroup. Thus \(G\to\operatorname{Aut}(V_i)\) is a
monomorphism, and hence, by
\hyperref[cor:sga3-VIB-1.4.2]{1.4.2}, a closed immersion.
This proves
\hyperref[concl:sga3-VIB-11.11-iv]{\textup{(iv)}}
\(\Rightarrow\)
\hyperref[concl:sga3-VIB-11.11-v]{\textup{(v)}}. Finally,
\hyperref[concl:sga3-VIB-11.11-v]{\textup{(v)}}
\(\Rightarrow\)
\hyperref[concl:sga3-VIB-11.11-i]{\textup{(i)}} is
immediate, and the proposition follows.

\medskip
\phantomsection\label{rem:sga3-VIB-11.11.1}%
\noindent\textbf{Remark 11.11.1.---}
Proposition \hyperref[prop:sga3-VIB-11.11]{11.11} admits
the following generalization. Let \(S\) be a regular,
locally Noetherian scheme of dimension at most \(1\), and
let \(G\) be a group scheme that is flat, quasi-compact, and
quasi-separated over \(S\). \textup{(In this case,
\(\mathcal A(G)\) is a torsion-free
\(\mathcal O_S\)-module, and hence is flat.)}

\begin{enumerate}
\item[\textup{(a)}]\phantomsection\label{concl:sga3-VIB-11.11.1-a}
  The following conditions are equivalent:%
  \footnote{\label{note:sga3-VIB-123}Editor's note: the
  equivalence of these conditions is proved in the added
  \hyperref[sec:sga3-VIB-12]{Section 12}.}
  \begin{enumerate}
  \item[\textup{(i)}]\phantomsection\label{concl:sga3-VIB-11.11.1-a-i}
    \(G\) is affine over \(S\).
  \item[\textup{(ii)}]\phantomsection\label{concl:sga3-VIB-11.11.1-a-ii}
    \(G\) is quasi-affine over \(S\).
  \item[\textup{(iii)}]\phantomsection\label{concl:sga3-VIB-11.11.1-a-iii}
    \(G\) acts faithfully on a quasi-affine flat
    \(S\)-scheme.
  \item[\textup{(iv)}]\phantomsection\label{concl:sga3-VIB-11.11.1-a-iv}
    \(G\) acts linearly and faithfully on a flat
    quasi-coherent \(\mathcal O_S\)-module.
  \end{enumerate}

\item[\textup{(b)}]\phantomsection\label{concl:sga3-VIB-11.11.1-b}
  If, moreover, \(G\) is of finite type over \(S\) and
  \(S\) is Noetherian, these conditions imply that \(G\)
  is isomorphic to a closed subgroup of
  \(\operatorname{Aut}(\mathcal E)\), where
  \(\mathcal E\) is a finite-rank locally free
  \(\mathcal O_S\)-module.%
  \footnote{\label{note:sga3-VIB-124}Editor's note: this is
  proved, with various generalizations, in the added
  \hyperref[sec:sga3-VIB-13]{Section 13}.}
\end{enumerate}

% END INCLUDED SGA3_English_current_buildable_source_20260823/source/inputs/VI/tex_reference_v2/components/85_expose_VIB_prop1110bis_through_rem11111_loop2_r1_en.tex

% BEGIN INCLUDED SGA3_English_current_buildable_source_20260823/source/inputs/VI/tex_reference_v2/components/86_expose_VIB_lemma1112_through_thm1116_r3_en.tex
% SGA 3, Expose VI_B, Loop 1 English reconstruction.
% R3: independent convention-v2 PDF replay assigned distinct
% precise anchors to the two custom-labelled conclusions 11.15(i)/(ii).
% Scope: Lemma 11.12 through Theorem 11.16, including the
% complete proof of Theorem 11.16 and editor's notes 125--129.
% Authority: Polo--Gille Exp6B-13oct24.pdf, component-local
% pages 90--94 (combined-reader pages 431--435; re-edition
% headers 420--424; original-print margin markers 405--409).
% Lemma 11.17.0 begins on local page 94 / combined page 435 /
% header 424 under original-print marker 409 and is excluded.
% OCR is locator-only. Reinhold e7a259f English Markdown is
% comparison-only and is credited in the component report.
% No source diagram occurs in this range.

\medskip
\phantomsection\label{lem:sga3-VIB-11.12}%
\noindent\textbf{Lemma 11.12.---}
\emph{Let \(k\) be a field and let \(G\) be an affine
\(k\)-group. Put \(A=\mathcal A(G)\). Given \(x\in A\),
there is a finitely generated \(k\)-subalgebra \(B\) of
\(A\) such that}
\[
  x\in B,\qquad
  \Delta(B)\subset B\otimes_kB,\qquad
  u(B)\subset B,
\]
\emph{where \(u\) denotes the involution of \(A\)
corresponding to the inversion morphism of \(G\).}%
\footnote{\label{note:sga3-VIB-125}Editor's note: on the
one hand, this is generalized in \hyperref[sec:sga3-VIB-13]{Section 13} to the case
where \(G\) is affine and flat over a regular base of
dimension at most \(2\). On the other hand, the original
has been expanded and corrected in what follows.}

We may suppose that \(x\neq0\). Then
\(\Delta(x)\neq0\), since
\[
  (\varepsilon\otimes\mathrm{id})\Delta(x)=x,
\]
where \(\varepsilon\) denotes the augmentation
\textup{(counit)} of \(A\). Write
\phantomsection\label{form:sga3-VIB-11.12-eq1}%
\[
  \tag{1}
  \Delta(x)=\sum_{j=1}^{n}e_j\otimes a_j,
  \qquad n\ \text{minimal}.
\]
The \(e_j\) \textup{(respectively the \(a_j\))} are then
linearly independent. Complete
\((e_1,\ldots,e_n)\) to a basis
\((e_i)_{i\in I}\) of \(A\), and, for
\[
  j\in J=\{1,\ldots,n\},
\]
write
\[
  \Delta(e_j)=\sum_{i\in I}e_i\otimes b_{ij}.
\]
Applying \(\Delta\otimes\mathrm{id}\) and
\(\mathrm{id}\otimes\Delta\) to
\hyperref[form:sga3-VIB-11.12-eq1]{\textup{(1)}}, the
axiom
\hyperref[ax:sga3-VIB-11.8.0-CO1]{\textup{(CO 1)}} of
\hyperref[def:sga3-VIB-11.8.0]{11.8.0}
\textup{(see also \textup{(HA 1)} in I, \SGARefLink{sga3:i:subsection:4-2}{4.2})} gives
\[
  \sum_{j\in J}e_j\otimes\Delta(a_j)
  =
  \sum_{\ell\in J}\Delta(e_\ell)\otimes a_\ell
  =
  \sum_{i\in I}e_i\otimes
    \left(\sum_{\ell\in J}b_{i\ell}\otimes a_\ell\right).
\]
Since the \(e_i\) are linearly independent, it follows that
\phantomsection\label{form:sga3-VIB-11.12-eq2}%
\[
  \tag{2}
  \Delta(a_j)=\sum_{\ell\in J}b_{j\ell}\otimes a_\ell
  \qquad (j\in J).
\]

Let \(B\) be the finitely generated \(k\)-subalgebra of
\(A\) generated by the \(b_{ij}\) and the \(u(b_{ij})\),
for \(i,j\in J\). It is clear that \(u(B)=B\).

Applying \(\Delta\otimes\mathrm{id}\) and
\(\mathrm{id}\otimes\Delta\) to
\hyperref[form:sga3-VIB-11.12-eq2]{\textup{(2)}}, we obtain
once again from \textup{(CO 1)} the equalities
\[
  \sum_{\ell\in J}\Delta(b_{j\ell})\otimes a_\ell
  =
  \sum_{i\in J}b_{ji}\otimes\Delta(a_i)
  =
  \sum_{i,\ell\in J}b_{ji}\otimes b_{i\ell}\otimes a_\ell.
\]
Since the \(a_\ell\) are linearly independent, this yields
\phantomsection\label{form:sga3-VIB-11.12-eq3}%
\[
  \tag{3}
  \Delta(b_{j\ell})
    =\sum_{i\in J}b_{ji}\otimes b_{i\ell}
  \qquad (j,\ell\in J).
\]
Moreover,
\[
  \Delta\circ u=(u\otimes u)\circ v\circ\Delta,
  \qquad v(a\otimes b)=b\otimes a,
\]
and hence
\phantomsection\label{form:sga3-VIB-11.12-eq4}%
\[
  \tag{4}
  \Delta\bigl(u(b_{j\ell})\bigr)
    =\sum_{i\in J}u(b_{i\ell})\otimes u(b_{ji})
  \qquad (j,\ell\in J).
\]
Since \(\Delta\) is an algebra homomorphism,
\hyperref[form:sga3-VIB-11.12-eq3]{\textup{(3)}} and
\hyperref[form:sga3-VIB-11.12-eq4]{\textup{(4)}} imply
\[
  \Delta(B)\subset B\otimes_kB.
\]
Finally, axiom
\hyperref[ax:sga3-VIB-11.8.0-CO2]{\textup{(CO 2)}} of
\hyperref[def:sga3-VIB-11.8.0]{11.8.0}
\textup{(see also \textup{(HA 2)} in I, \SGARefLink{sga3:i:subsection:4-2}{4.2})}, applied to
\hyperref[form:sga3-VIB-11.12-eq2]{\textup{(2)}} and
\hyperref[form:sga3-VIB-11.12-eq1]{\textup{(1)}}, gives%
\begingroup
\renewcommand{\thefootnote}{\fnsymbol{footnote}}%
\footnote[1]{\label{note:sga3-VIB-11.12-counit-source-reading}
Source-reading note: the French re-edition prints
\(a_j=\sum_{i\in I}\varepsilon(a_j)b_{ij}\).
Applying \(\mathrm{id}\otimes\varepsilon\) to
\hyperref[form:sga3-VIB-11.12-eq2]{\textup{(2)}} gives the
first identity displayed here. The Reinhold comparison
lineage repeats the printed formula.}
\endgroup
\phantomsection\label{form:sga3-VIB-11.12-counit-identities}%
\[
  a_j=\sum_{i\in J}b_{ji}\varepsilon(a_i),
  \qquad
  x=\sum_{j\in J}\varepsilon(e_j)a_j.
\]
Thus \(x\in B\), proving the lemma.

\medskip
\phantomsection\label{prop:sga3-VIB-11.13}%
\noindent\textbf{Proposition 11.13.---}
\emph{Let \(k\) be a field and let \(G\) be an affine
\(k\)-group with algebra \(A\). Then \(G\) is the
projective limit of an increasing filtered system of
affine \(k\)-groups of finite type whose transition
morphisms are faithfully flat.}

If \(B\) and \(B'\) are two finitely generated
subalgebras of \(A\) stable under \(\Delta\) and \(u\), the
subalgebra generated by \(B\) and \(B'\) has the same
properties. Hence, by
\hyperref[lem:sga3-VIB-11.12]{Lemma 11.12}, \(A\) is the
direct limit of an increasing filtered family
\((B_i)_{i\in I}\) of finitely generated subalgebras stable
under \(\Delta\) and \(u\). Each \(B_i\), equipped with the
restriction of \(u\) and with the morphism
\[
  B_i\longrightarrow B_i\otimes_kB_i
\]
induced by \(\Delta\), is a Hopf algebra. Thus, by I, \SGARefLink{sga3:i:subsection:4-2}{4.2},
it is the algebra of an affine \(k\)-group \(G_i\) of
finite type over \(k\). Finally, since
\[
  A=\varinjlim_i B_i,
\]
we have
\[
  G=\varprojlim_i G_i
\]
\textup{(cf.\ EGA IV\(_3\), 8.2.3)}. The transition
morphisms are faithfully flat by the following lemma.

\medskip
\phantomsection\label{lem:sga3-VIB-11.14}%
\noindent\textbf{Lemma 11.14.---}
\emph{Let \(k\) be a field, let
\(u:G\to H\) be a morphism between affine \(k\)-groups of
finite type, and let}
\[
  u^\natural:B\longrightarrow A
\]
\emph{be the corresponding morphism of \(k\)-algebras.}%
\footnote{\label{note:sga3-VIB-126}Editor's note: we write
\(u^\natural\) \textup{(instead of \(u^\circ\))} for the
morphism \(B\to A\) corresponding to
\(u:\operatorname{Spec}(A)\to\operatorname{Spec}(B)\).}
\emph{Then \(u\) is faithfully flat if and only if
\(u^\natural\) is injective.}

The condition is plainly necessary
\textup{(cf.\ EGA 0\(_{\mathrm I}\), 6.6.1)}. Conversely,
suppose that \(u^\natural\) is injective, and put
\[
  N=\operatorname{Ker}(u).
\]
By
\hyperref[thm:sga3-VIA-3.3.2]{VI\(_A\), 3.3.2} and
\hyperref[prop:sga3-VIA-5.4.1]{VI\(_A\), 5.4.1},
\(G/N\) is a \(k\)-group of finite type and \(u\) factors
as
\[
  G\xrightarrow{\ p\ }G/N\xrightarrow{\ v\ }H,
\]
where \(p\) is faithfully flat and \(v\) is a closed
immersion. Since \(H\) is affine, \(G/N\) is affine, and
\[
  v^\natural:B\longrightarrow\mathcal O(G/N)
\]
is surjective \textup{(cf.\ EGA I, 4.2.3)}. But
\[
  u^\natural=p^\natural\circ v^\natural
\]
is injective, so \(v^\natural\) is also injective. Hence
\(v^\natural\), and therefore \(v\), is an isomorphism.
Since \(p\) is faithfully flat, so is \(u\).

\medskip
\phantomsection\label{def:sga3-VIB-11.15}%
\noindent\textbf{Definition 11.15.---}%
\footnote{\label{note:sga3-VIB-127}Editor's note: we have
added the hypothesis that \(G\) be quasi-compact and have
spelled out the equivalence of conditions
\textup{(i)} and \textup{\SGARefLink{concl:sga3-VIB-11.15-i}{(ii)}}.}
\emph{Let \(k\) be a field, let \(G\) be a quasi-compact
\(k\)-group, and let \(V\) be a \(k\)-vector space endowed
with a \(k\)-linear action of \(G\), hence, by
\hyperref[prop:sga3-VIB-11.6.1]{11.6.1}, with an
\(\mathcal A(G)\)-comodule structure}
\[
  \delta:V\longrightarrow V\otimes\mathcal A(G).
\]
\emph{Let \(v\in V\) be nonzero. The following conditions
are equivalent:}
\begin{enumerate}
\item[\textup{(i)}]\phantomsection\label{concl:sga3-VIB-11.15-i}
  \emph{There is a \(\lambda\in\mathcal A(G)\),
  necessarily unique, such that}
  \[
    \delta(v)=v\otimes\lambda.
  \]

\item[\textup{(ii)}]\phantomsection\label{concl:sga3-VIB-11.15-ii}
  \emph{For every \(k\)-algebra \(R\) and every
  \(g\in G(R)\), one has \(g\cdot v\in Rv\); that is,
  there is a necessarily unique \(f\in R\) such that,
  in \(V\otimes R\),}
  \[
    g\cdot v=v\otimes f.
  \]
\end{enumerate}

It is clear that
\hyperref[concl:sga3-VIB-11.15-i]{\textup{(i)}}
implies
\hyperref[concl:sga3-VIB-11.15-ii]{\textup{(ii)}}.
Conversely, if \textup{\SGARefLink{concl:sga3-VIB-11.15-i}{(ii)}} holds, apply it with
\[
  R=\mathcal A(G),
  \qquad
  g=\operatorname{id}_{\mathcal A(G)}.
\]
There is then a unique \(\lambda\in\mathcal A(G)\) such
that \(\delta(v)=v\otimes\lambda\).

When \(v\) satisfies these conditions, we call \(v\) a
\emph{semi-invariant vector under \(G\)}, and
\(\lambda\) the \emph{weight} of \(v\). We also say that
\(v\) is a \emph{semi-invariant of weight \(\lambda\)}.

Let \(\Delta\) denote the comultiplication of
\(\mathcal A(G)\). The equality
\phantomsection\label{form:sga3-VIB-11.15-group-like}%
\[
  v\otimes\lambda\otimes\lambda
  =
  (\delta\otimes\mathrm{id})(\delta(v))
  =
  (\mathrm{id}\otimes\Delta)(\delta(v))
  =
  v\otimes\Delta(\lambda)
\]
implies
\[
  \Delta(\lambda)=\lambda\otimes\lambda.
\]
Consequently, \(\lambda\) defines a morphism of Hopf
algebras
\[
  \mathcal A(\mathbb G_{m,k})
    =k[T,T^{-1}]
    \longrightarrow\mathcal A(G),
  \qquad T\longmapsto\lambda,
\]
and hence a morphism of \(k\)-groups
\[
  \lambda:G\longrightarrow\mathbb G_{m,k}.
\]
Thus \(\lambda\) is a \emph{character} of \(G\), called the
\emph{character associated with the semi-invariant vector
\(v\)}.

\medskip
\phantomsection\label{lem:sga3-VIB-11.16.0}%
\noindent\textbf{Lemma 11.16.0.---}%
\footnote{\label{note:sga3-VIB-128}Editor's note: this
lemma, taken from \textup{[DG70]}, Section II.2, \SGARefLink{prop:sga3-VIB-3.5}{3.5}, has
been added.}
\emph{Let \(k\) be a field, let \(H\) be an affine
\(k\)-group, let \(V\) be an \(H\)-module of dimension
\(n\), and let \(U\) be a \(d\)-dimensional vector
subspace of \(V\). Consider the line}
\[
  D=\bigwedge\nolimits^dU
    \subset\bigwedge\nolimits^dV.
\]
\emph{Then \(U\) is stable under \(H\) if and only if
\(D\) is stable under \(H\).}

Necessity is clear. To prove sufficiency, we may assume
\(d<n\). Let \((e_1,\ldots,e_d)\) be a basis of \(U\), and
complete it to a basis \((e_1,\ldots,e_n)\) of \(V\).
For every \(k\)-algebra \(R\), the module
\[
  V_R=V\otimes R
\]
is free over \(R\), and
\phantomsection\label{form:sga3-VIB-11.16.0-UR}%
\[
  U_R
  =
  \{\,v\in V_R\mid
    v\wedge(e_1\wedge\cdots\wedge e_d)=0\,\}.
\]
Indeed, for \(i>d\), the elements
\[
  e_i\wedge e_1\wedge\cdots\wedge e_d
\]
are linearly independent in
\(\bigwedge_R^{d+1}V_R\). Since \(H(R)\) acts on
\(\bigwedge_R^\bullet V_R\) by
\[
  h(x_1\wedge\cdots\wedge x_s)
  =
  h(x_1)\wedge\cdots\wedge h(x_s),
\]
it follows that, if \(H(R)\) stabilizes
\[
  D_R=R(e_1\wedge\cdots\wedge e_d),
\]
then it also stabilizes \(U_R\).

\medskip
\phantomsection\label{thm:sga3-VIB-11.16}%
\noindent\textbf{Theorem 11.16 \textup{(Chevalley)}.---}
\emph{Let \(k\) be a field, let \(G\) be an affine
algebraic \(k\)-group, and let \(H\) be a closed group
subscheme of \(G\).}%
\footnote{\label{note:sga3-VIB-129}Editor's note: first,
recall that every group subscheme of \(G\) is closed
\textup{(\hyperref[cor:sga3-VIB-1.4.2]{1.4.2})}. Second,
we have stated the result in its usual form, namely that
``\(H\) is the stabilizer of a line in a representation of
\(G\),'' while retaining the original formulation in terms
of a sequence \(a_1,\ldots,a_n\) of semi-invariants in
\(\mathcal A(G)\).}
\emph{Then there are a finite-dimensional \(G\)-module
\(V\) and a line \(D\) in \(V\) such that}
\[
  H=\operatorname{Norm}_G(D);
\]
\emph{that is, for every \(k\)-algebra \(R\),}
\phantomsection\label{form:sga3-VIB-11.16-stabilizer}%
\[
  H(R)=\{\,g\in G(R)\mid g(D_R)=D_R\,\}.
\]

\emph{Equivalently, there are finitely many elements
\(a_1,\ldots,a_n\in\mathcal A(G)\), all semi-invariant of
the same weight \(\lambda\) for the ``right'' action of
\(H\); that is,}
\[
  a_i(gh)=\lambda(h)a_i(g)
\]
\emph{for every \(k\)-algebra \(R\),
\(g\in G(R)\), and \(h\in H(R)\), such that \(H\) is the
largest closed group subscheme of \(G\) under which the
\(a_i\) are semi-invariant.}

\emph{Proof.}
Let \(\Delta\) and \(\varepsilon\) denote, respectively,
the comultiplication and the augmentation of
\[
  A=\mathcal A(G).
\]
There is an ideal \(I\) of \(A\), contained in
\(\operatorname{Ker}\varepsilon\), such that
\[
  H=\operatorname{Spec}(A/I),
  \qquad
  \Delta(I)\subset I\otimes A+A\otimes I.
\]
Put \(B=A/I\), and let \(\pi:A\to B\) be the projection.
Consider the left action of \(H\) on \(A\) given by
\[
  (h\phi)(g)=\phi(gh).
\]
The corresponding \(B\)-comodule structure is
\phantomsection\label{form:sga3-VIB-11.16-H-coaction-A}%
\[
  \overline{\Delta}:\quad
  A\xrightarrow{\ \Delta\ }A\otimes A
    \xrightarrow{\ \mathrm{id}_A\otimes\pi\ }A\otimes B.
\]
Then \(I\) is an \(H\)-submodule of \(A\), since
\[
  \overline{\Delta}(I)\subset I\otimes B.
\]

The \(k\)-algebra \(A\) is finitely generated, and hence
Noetherian, so \(I\) has a finite system of generators
\((x_1,\ldots,x_r)\). By
\hyperref[lem:sga3-VIB-11.8]{11.8}, the \(x_i\) are
contained in a finite-dimensional \(G\)-submodule \(V\) of
\(A\). Then
\[
  W=V\cap I
\]
is a finite-dimensional \(H\)-module; write
\[
  d=\dim_kW.
\]
Since \(V\) contains every \(x_i\), \(W\) generates the
ideal \(I\).

Put
\[
  E=\bigwedge\nolimits^dV,
\]
choose a basis \((w_1,\ldots,w_d)\) of \(W\), and choose a
basis \((e_0,\ldots,e_n)\) of \(E\) containing
\[
  e_0=w_1\wedge\cdots\wedge w_d.
\]
The action of \(G\) on \(V\) canonically induces an action
of \(G\) on \(E\), hence an \(A\)-comodule structure
\[
  \rho:E\longrightarrow E\otimes A.
\]
For \(j=0,\ldots,n\), write%
\begingroup
\renewcommand{\thefootnote}{\fnsymbol{footnote}}%
\footnote[2]{\label{note:sga3-VIB-11.16-rho-index-source-reading}
Source-reading note: the French re-edition and the Reinhold
comparison lineage print the lower summation bound as
\(i=1\), although the chosen basis begins with \(e_0\) and
the ensuing identity sums from \(0\). The lower bound is
therefore \(i=0\) here.}
\endgroup
\phantomsection\label{form:sga3-VIB-11.16-rho-matrix}%
\[
  \rho(e_j)=\sum_{i=0}^{n}e_i\otimes b_{ij}.
\]
As in the proof of
\hyperref[lem:sga3-VIB-11.12]{11.12}, one has
\phantomsection\label{form:sga3-VIB-11.16-star}%
\[
  \tag{\(*\)}
  \Delta(b_{ij})
    =\sum_{\ell=0}^{n}b_{i\ell}\otimes b_{\ell j}.
\]

Put \(a_i=b_{i0}\). Equivalently, if
\((e_0^*,\ldots,e_n^*)\) is the dual basis, then the \(a_i\)
are the matrix coefficients
\[
  g\longmapsto e_i^*(ge_0).
\]
The action of \(H\) on \(E\) corresponds to
\phantomsection\label{form:sga3-VIB-11.16-H-coaction-E}%
\[
  \overline{\rho}:\quad
  E\xrightarrow{\ \rho\ }E\otimes A
    \xrightarrow{\ \mathrm{id}_E\otimes\pi\ }E\otimes B.
\]
Since \(W\) is stable under \(H\), the vector \(e_0\) is
semi-invariant under \(H\). Hence
\[
  \overline{\rho}(e_0)=e_0\otimes\pi(a_0),
  \qquad
  a_i=b_{i0}\in I\quad(1\leq i\leq n).
\]
Substitution in
\hyperref[form:sga3-VIB-11.16-star]{\textup{(\(*\))}} gives
\[
  \overline{\Delta}(a_i)
    =a_i\otimes\pi(a_0)
  \qquad(0\leq i\leq n).
\]
Thus the \(a_i\) are semi-invariant under \(H\), all of
weight \(\pi(a_0)\). Moreover, after replacing \(E\), if
necessary, by the \(G\)-submodule generated by \(e_0\)
\textup{(cf.\ \hyperref[lem:sga3-VIB-11.8]{11.8})}, we may
suppose that the \(a_i\) are linearly independent.

Conversely, let
\[
  H'=\operatorname{Spec}B',
  \qquad
  B'=A/I',
\]
be a closed group subscheme of \(G\) under which every
\(a_i\) is semi-invariant, of weight
\(\lambda_i\in B'\). This includes, in particular, the
case in which \(e_0\) is semi-invariant under \(H'\) of
weight \(\lambda'\). We shall prove that \(H'\subset H\).%
\begingroup
\renewcommand{\thefootnote}{\fnsymbol{footnote}}%
\footnote[3]{\label{note:sga3-VIB-11.16-maximality-source-reading}
Source-reading note: the French re-edition says at this
point ``let us show that \(H'=H\),'' as does the Reinhold
comparison lineage. The maximality assertion requires
\(H'\subset H\), and that is exactly the conclusion proved
below.}
\endgroup\ 
Let
\(\pi':A\to B'\) be the projection. The hypothesis gives
\phantomsection\label{form:sga3-VIB-11.16-converse-coefficients}%
\[
  a_i\otimes\lambda_i
  =
  (\mathrm{id}_A\otimes\pi')\Delta(b_{i0})
  =
  \sum_{\ell=0}^{n}b_{i\ell}\otimes\pi'(a_\ell).
\]
It follows that
\[
  \lambda_i=\pi'(a_0),
  \qquad
  a_\ell\in I'\quad(1\leq\ell\leq n),
\]
and hence \(e_0\) is semi-invariant under \(H'\). By
\hyperref[lem:sga3-VIB-11.16.0]{Lemma 11.16.0}, \(W\) is
stable under \(H'\). Since \(W\) generates \(I\), the ideal
\(I\) is also stable under \(H'\), and therefore
\[
  \Delta(I)\subset I\otimes A+A\otimes I'.
\]
Finally,
\[
  I\subset\operatorname{Ker}\varepsilon,
  \qquad
  (\varepsilon\otimes\mathrm{id}_A)\circ\Delta
    =\mathrm{id}_A,
\]
so \(I\subset I'\), whence \(H'\subset H\). This completes
the proof.

% END INCLUDED SGA3_English_current_buildable_source_20260823/source/inputs/VI/tex_reference_v2/components/86_expose_VIB_lemma1112_through_thm1116_r3_en.tex

% BEGIN INCLUDED SGA3_English_current_buildable_source_20260823/source/inputs/VI/tex_reference_v2/components/87_expose_VIB_lemma11170_through_cor11184_loop2_r1_en.tex
% SGA 3, Expose VI_B, Loop 1 English reconstruction.
% Loop 2 R1: two source-bound crops are replaced by faithful native
% tikz-cd diagrams; prose, mathematics, and labels otherwise remain
% identical to the independently reviewed Loop-1 component.
% Scope: Lemma 11.17.0 through Corollary 11.18.4, including
% the complete proof of Theorem 11.17, Lemma 11.17.1,
% Subsection 11.18, Lemma 11.18.1, Theorem 11.18.2,
% Remark 11.18.3, Corollary 11.18.4, and editor's notes
% 130--134.
% Authority: Polo--Gille Exp6B-13oct24.pdf, component-local
% pages 94--99 (combined-reader pages 435--440; re-edition
% headers 424--429; original-print margin markers 409--411).
% Remarks 11.18.5 begin later on local page 99 / combined
% page 440 / header 429 and are excluded.
% OCR is locator-only. Reinhold e7a259f English Markdown is
% comparison-only and is credited in the component report.
% Loop 2 replaces both source-bound commutative-diagram crops with
% faithful native LaTeX reconstructions; both are closed.

\medskip
\phantomsection\label{lem:sga3-VIB-11.17.0}%
\noindent\textbf{Lemma 11.17.0.---}%
\footnote{\label{note:sga3-VIB-130}Editor's note: we have
added this lemma, drawn from the proof of \SGARefLink{par:sga3-VIB-5.6}{Theorem 5.6} of
\textup{[DG70]}, Section III.3. By an abuse of notation, the
same letter \(E\) denotes a \(k\)-vector space and the
\(k\)-module scheme
\(W(E)=\operatorname{Spec}S(E^*)\).}
\emph{Let \(k\) be an algebraically closed field, let \(G\)
be a reduced affine algebraic \(k\)-group, let \(V\) be a
finite-dimensional \(G\)-module over \(k\), and let
\(v\in V\). Let \(E\) be the vector subspace of \(V\)
spanned by the vectors \(gv\), for \(g\in G(k)\). Then
\(E\) is the smallest \(G\)-submodule of \(V\) containing
\(v\). Consequently, the morphism}
\[
  G\times E\longrightarrow V,\qquad (g,x)\longmapsto gx,
\]
\emph{factors through \(E\).}

\emph{Proof.}
By \hyperref[lem:sga3-VIB-11.8]{11.8}, there is a smallest
\(G\)-submodule \(U\) of \(V\) containing \(v\). Indeed,
if
\[
  \mu:V\longrightarrow V\otimes\mathcal A(G)
\]
denotes the comodule structure and
\[
  \mu(v)=\sum_{i=1}^{n}v_i\otimes a_i
\]
with the \(a_i\) linearly independent, then
\[
  U=\operatorname{Span}_k(v_1,\ldots,v_n).
\]
It is clear that \(U\) contains \(E\), and that the
morphism \(G\times U\to V\) factors through \(U\).

Conversely, the inverse image of \(E\) under the morphism
\[
  \mu_v:G\longrightarrow V,\qquad g\longmapsto gv,
\]
is a closed subset of \(G\) containing all the rational
points. Those points are dense in \(G\), since \(G\) is of
finite type over \(k\)
\textup{(cf.\ EGA IV\(_3\), 10.4.8)}. Hence
\[
  \mu_v^{-1}(E)=G.
\]
Since \(G\) is reduced, \(\mu_v\) therefore factors through
\(E\), whence \(E=U\).

\medskip
\phantomsection\label{thm:sga3-VIB-11.17}%
\noindent\textbf{Theorem 11.17 \textup{(Chevalley)}.---}
\emph{Let \(k\) be a field, let \(G\) be an affine
\(k\)-group \textup{(not necessarily of finite type)}, and
let \(N\) be a closed group subscheme of \(G\), invariant
in \(G\). Then the \textup{(fpqc)} sheaf quotient \(G/N\)
is representable by an affine \(k\)-group.}%
\footnote{\label{note:sga3-VIB-131}Editor's note: for
another proof of this theorem that does not use the results
of Expos\'e VI\(_A\), see \textup{[Ta72]}, \SGARefLink{prop:sga3-VIB-5.2}{Theorem 5.2}
\textup{(see also \hyperref[rem:sga3-VIB-11.18.5]{Remark 11.18.5})}.}

\emph{Proof.}
Suppose first that \(G\) is of finite type. By
\hyperref[thm:sga3-VIA-3.2]{VI\(_A\), 3.2} and
\hyperref[par:sga3-VIA-5.2]{5.2}, the \textup{(fpqc)}
sheaf quotient \(G/N\) is representable by a \(k\)-group
\(Q\); it remains to prove that \(Q\) is affine. We proceed
in several steps, first supposing that \(k\) is
algebraically closed.%
\footnote{\label{note:sga3-VIB-132}Editor's note: in what
follows, we have expanded the original and corrected the
erroneous assertion
\((G/N)_{\mathrm{red}}=G_{\mathrm{red}}/N_{\mathrm{red}}\),
drawing on \textup{[DG70]}, Section III.3, \SGARefLink{par:sga3-VIB-5.6}{5.6}.}

\smallskip
\phantomsection\label{case:sga3-VIB-11.17-a}%
\noindent\textup{(a)}
Suppose, in addition, that \(G\) is reduced and connected,
and that \(N\) is reduced. By
\hyperref[thm:sga3-VIB-11.16]{11.16}, there are a
finite-dimensional \(G\)-module \(V\) over \(k\) and a
line \(D=ke_0\) such that
\[
  \operatorname{Norm}_G(D)=N.
\]
In particular, \(N\) acts on \(D\) through a character
\[
  \chi:N\longrightarrow\mathbb G_{m,k}.
\]

Fix \(h\in N(k)\). For every \(g\in G(k)\), one has
\[
  hge_0=g(g^{-1}hg)e_0
       =\chi(g^{-1}hg)\,ge_0.
\]
Thus \(\chi(g^{-1}hg)\) is an eigenvalue of \(h\). The
continuous map
\[
  \varphi:G(k)\longrightarrow k,\qquad
  g\longmapsto\chi(g^{-1}hg),
\]
therefore takes only finitely many values. Since \(G(k)\)
is irreducible \textup{(it is dense in \(G\))}, it follows
that
\[
  \varphi(g)=\varphi(e)=\chi(h)
  \qquad\text{for every }g\in G(k).
\]
Consequently,
\phantomsection\label{form:sga3-VIB-11.17-character-invariance}%
\[
  \chi(g^{-1}hg)=\chi(h)
  \qquad
  \text{for all }g\in G(k),\ h\in N(k).
\]

Let \(E\) be the vector subspace of \(V\) spanned by the
vectors \(ge_0\), for \(g\in G(k)\). By
\hyperref[lem:sga3-VIB-11.17.0]{Lemma 11.17.0}, this is the
\(G\)-submodule of \(V\) generated by \(e_0\).

By the preceding identity, the two morphisms
\[
  N\times E\longrightarrow E,\qquad
  (h,x)\longmapsto hx
  \quad\text{and}\quad
  (h,x)\longmapsto\chi(h)x,
\]
agree on
\[
  (N\times E)(k)=N(k)\times E(k),
\]
which is dense in \(N\times E\). Since \(N\times E\) is
reduced \textup{(and \(E\) is separated)}, the two
morphisms are equal. Thus \(N\) acts on \(E\) by scalar
multiplications.

It follows that \(N\) is contained in the kernel \(K\) of
the morphism
\[
  \rho:G\longrightarrow
    \operatorname{GL}\bigl(\operatorname{End}_k(E)\bigr)
\]
defined, for every \(k\)-algebra \(R\), by
\[
  \rho(g)(u)=gug^{-1}
  \qquad
  \bigl(g\in G(R),\
  u\in\operatorname{End}_R(E\otimes_kR)\bigr).
\]
Conversely, if \(g\in K(R)\), then
\[
  g(Re_0)=Re_0,
\]
and hence \(g\in N(R)\). Thus \(N=K\). By
\hyperref[prop:sga3-VIA-5.4.1]{VI\(_A\), 5.4.1}, the
morphism
\[
  G/N\longrightarrow
  \operatorname{GL}\bigl(\operatorname{End}_k(E)\bigr)
\]
is a closed immersion. Therefore \(G/N\) is affine.

\smallskip
\phantomsection\label{case:sga3-VIB-11.17-b}%
\noindent\textup{(b)}
Suppose now that \(G\) and \(N\) are reduced, without
assuming that \(G\) is connected. Put
\[
  N'=N\cap G^0.
\]
Then \(G^0/N'\) is affine by \textup{(a)}. Moreover,
\(NG^0\) is an invariant subgroup of \(G\), and
\(G/NG^0\), being a quotient of the finite constant group
\(G/G^0\)
\textup{(cf.\ \hyperref[prop:sga3-VIA-5.5.1]
{VI\(_A\), 5.5.1})}, is itself a finite constant group.
The scheme \(G/N\) is therefore the disjoint union of the
fibers of
\[
  G/N\longrightarrow G/NG^0.
\]
By \hyperref[par:sga3-VIA-5.3.3]{VI\(_A\), 5.3.3}, all
these fibers are isomorphic to \(NG^0/N\), hence to
\(G^0/N'\). Thus \(G/N\) is affine.

\smallskip
\phantomsection\label{case:sga3-VIB-11.17-c}%
\noindent\textup{(c)}
Suppose that \(G\) is reduced, with \(N\) arbitrary. The
morphism
\[
  G\times N\longrightarrow N,\qquad
  (g,h)\longmapsto ghg^{-1},
\]
induces a morphism
\[
  (G\times N)_{\mathrm{red}}\longrightarrow N_{\mathrm{red}}.
\]
Since \(G\) is reduced and \(k\) is algebraically closed,
\[
  (G\times N)_{\mathrm{red}}=G\times N_{\mathrm{red}},
\]
so \(N_{\mathrm{red}}\) is an invariant subgroup of \(G\).
\textup{(This fails when \(G\) is not reduced; see
\hyperref[sga3:VIA:0.2]{VI\(_A\), 0.2}.)}

Hence, by \textup{\SGARefLink{case:sga3-VIB-11.17-b}{(b)}},%
\begingroup
\renewcommand{\thefootnote}{\fnsymbol{footnote}}%
\footnote[1]{\label{note:sga3-VIB-11.17-case-c-source-reading}
Source-reading note: the French re-edition prints
``by \textup{(a)}'' here, and the Reinhold comparison
lineage repeats that reference. At this stage \(G\) need
not be connected; the applicable preceding case is
\textup{\SGARefLink{case:sga3-VIB-11.17-b}{(b)}}.}
\endgroup\ 
\[
  G'=G/N_{\mathrm{red}}
\]
is affine. By
\hyperref[prop:sga3-VIA-5.6.1]{VI\(_A\), 5.6.1},
\[
  N'=N/N_{\mathrm{red}}
\]
is a finite \(k\)-group. Hence, by
\hyperref[thm:sga3-V-4.1]{Theorem 4.1 of Expos\'e V},
\[
  G/N=G'/N'
\]
is affine.

\smallskip
\phantomsection\label{case:sga3-VIB-11.17-d}%
\noindent\textup{(d)}
For arbitrary \(G\) and \(N\), base change of the
equivalence relation
\[
  G\times N\rightrightarrows G
\]
along \(G_{\mathrm{red}}\to G\) gives
\phantomsection\label{form:sga3-VIB-11.17-reduced-relation}%
\[
  G_{\mathrm{red}}\times N'
    \rightrightarrows G_{\mathrm{red}},
  \qquad
  N'=N\cap G_{\mathrm{red}}.
\]
Since the underlying spaces are the same \textup{(and the
quotient is the quotient ringed space)}, the morphism
\[
  G_{\mathrm{red}}/N'\longrightarrow G/N
\]
is a homeomorphism. The quotient
\(G_{\mathrm{red}}/N'\) is reduced, because
\[
  p:G_{\mathrm{red}}\longrightarrow G_{\mathrm{red}}/N'
\]
is faithfully flat. It follows that
\[
  (G/N)_{\mathrm{red}}
    \simeq G_{\mathrm{red}}/N'.
\]
The scheme on the right is affine by \textup{\SGARefLink{case:sga3-VIB-11.17-c}{(c)}}. Since
\(G/N\) is of finite type over \(k\)
\textup{(cf.\
\hyperref[thm:sga3-VIA-3.3.2]{VI\(_A\), 3.3.2})},
EGA I, 5.1.10 implies that \(G/N\) is affine.

Finally, let \(k\) be arbitrary and let \(\overline{k}\)
be an algebraic closure. By
\hyperref[concl:sga3-VIB-9.2-v]{9.2\textup{(v)}}, one has
\[
  (G\otimes_k\overline{k})/
    (N\otimes_k\overline{k})
  \simeq
  (G/N)\otimes_k\overline{k}.
\]
The first scheme is affine, hence so is the second.
Therefore \(G/N\) is affine by \textup{(fpqc)} descent
\textup{(cf.\ EGA IV\(_2\), 2.7.1)}. This proves
\hyperref[thm:sga3-VIB-11.17]{Theorem 11.17} when \(G\) is of finite type. To pass to the
general case, we need the following lemma.%
\footnote{\label{note:sga3-VIB-133}Editor's note: we have
added this lemma, taken from \textup{[DG70]}, Section III.3,
7.1.}

\medskip
\phantomsection\label{lem:sga3-VIB-11.17.1}%
\noindent\textbf{Lemma 11.17.1.---}
\emph{Let \((C_i\to A_i)_{i\in I}\) be a filtered direct
system of ring homomorphisms, all faithfully flat. Then}
\[
  A=\varinjlim_i A_i
\]
\emph{is faithfully flat over}
\[
  C=\varinjlim_i C_i.
\]

\emph{Proof.}
By \textup{[BAC]}, \SGARefLink{sec:sga3-VIB-3}{Section I.3}, Proposition 9, a ring
homomorphism \(B\to B'\) is faithfully flat if and only if
it is injective and \(B'/B\) is a flat \(B\)-module. Since
each \(C_i\to A_i\) is faithfully flat, there are exact
sequences
\phantomsection\label{seq:sga3-VIB-11.17.1-stage}%
\[
  0\longrightarrow C_i\longrightarrow A_i
   \longrightarrow A_i/C_i\longrightarrow0,
\]
and \(A_i/C_i\) is a flat \(C_i\)-module. Thus
\[
  (A_i/C_i)\otimes_{C_i}C
\]
is a flat \(C\)-module. Since direct limits are exact and
commute with tensor products, we obtain an exact sequence
\phantomsection\label{seq:sga3-VIB-11.17.1-limit}%
\[
  0\longrightarrow C\longrightarrow A
   \longrightarrow A/C\longrightarrow0
\]
and an isomorphism
\phantomsection\label{form:sga3-VIB-11.17.1-limit-quotient}%
\[
  \varinjlim_i\bigl((A_i/C_i)\otimes_{C_i}C\bigr)
    =(A/C)\otimes_CC=A/C.
\]
It follows that \(A/C\) is a flat \(C\)-module. Hence
\(A\) is faithfully flat over \(C\).

We now return to the proof of \hyperref[thm:sga3-VIB-11.17]{Theorem 11.17} in the general
case, where \(G\) is not assumed to be of finite type. Put
\[
  \mathcal A(G)=A,
  \qquad
  \mathcal A(N)=B=A/J.
\]
By \hyperref[prop:sga3-VIB-11.13]{11.13}, \(A\) is the
direct limit of an increasing filtered family
\((A_i)_{i\in I}\) of finitely generated sub-Hopf
algebras. Thus
\[
  G=\varprojlim_iG_i,
  \qquad
  G_i=\operatorname{Spec}(A_i),
\]
where the \(G_i\) are affine algebraic \(k\)-groups. Let
\(\Delta\) and \(\tau\) denote, respectively, the
comultiplication and antipode of \(A\), and put
\[
  \Delta^2=(\Delta\otimes\mathrm{id}_A)\circ\Delta.
\]

For each \(i\),
\[
  B_i=A_i/(J\cap A_i)
\]
is a quotient Hopf algebra of \(A_i\), so
\[
  N_i=\operatorname{Spec}(B_i)
\]
is a closed subgroup of \(G_i\). Since \(N\) is invariant
in \(G\), the morphism
\[
  G\times N\longrightarrow G,\qquad
  (g,n)\longmapsto gng^{-1},
\]
factors through \(N\). Equivalently, the pair \((A,J)\)
satisfies
\phantomsection\label{form:sga3-VIB-11.17-conjugation-coideal}%
\[
  \bigl(m_{13}\circ(\Delta\otimes\tau)\circ\Delta\bigr)(J)
    \subset A\otimes_kJ,
\]
where
\[
  m_{13}(a_1\otimes a_2\otimes a_3)
    =a_1a_3\otimes a_2.
\]
It follows that \((A_i,A_i\cap J)\) satisfies the analogous
property, so \(N_i\) is invariant in \(G_i\). Moreover,
\[
  \varinjlim_iB_i=B,
  \qquad
  \varprojlim_iN_i=N.
\]

By the finite-type case already proved, each
\textup{(fpqc)} sheaf quotient \(G_i/N_i\) is represented
by an affine \(k\)-group
\[
  Q_i=\operatorname{Spec}(C_i).
\]
Put
\[
  C=\varinjlim_iC_i.
\]
The \(Q_i\) form a filtered projective system of affine
\(k\)-groups, and its projective limit
\[
  Q=\varprojlim_iQ_i
\]
is the \(k\)-group \(\operatorname{Spec}(C)\)
\textup{(cf.\ EGA IV\(_3\), 8.2.3)}. We have an exact
sequence of \(k\)-groups
\phantomsection\label{seq:sga3-VIB-11.17-Q}%
\[
  1\longrightarrow N\longrightarrow G\longrightarrow Q.
\]

We claim that \(Q\) represents the \textup{(fpqc)} sheaf
quotient of \(G\) by \(N\). It is enough to verify that
\[
  G\longrightarrow Q
\]
is a covering for the \textup{(fpqc)} topology
\textup{(cf.\ IV, \SGARefLink{sga3:iv:scholium:3:3:2:1}{3.3.2.1} and 5.1.7.1)}. Every morphism
\[
  G_i\longrightarrow Q_i
\]
is faithfully flat
\textup{(cf.\ \hyperref[concl:sga3-VIB-9.2-xi]
{9.2\textup{(xi)}})}; equivalently, \(A_i\) is faithfully
flat over \(C_i\). Since
\[
  A=\varinjlim_iA_i,
  \qquad
  C=\varinjlim_iC_i,
\]
\hyperref[lem:sga3-VIB-11.17.1]{Lemma 11.17.1} shows that \(A\) is faithfully flat over
\(C\). Thus \(G\to Q\) is faithfully flat. It is also
affine, hence quasi-compact, and is therefore an
\textup{(fpqc)} covering. This completes the proof of
\hyperref[thm:sga3-VIB-11.17]{Theorem 11.17}.

\medskip
\phantomsection\label{par:sga3-VIB-11.18}%
\noindent\textbf{11.18. Complements.---}%
\footnote{\label{note:sga3-VIB-134}Editor's note: we have
added this subsection.}
\hyperref[thm:sga3-VIB-11.17]{Theorem 11.17} and its proof also give the following
results, taken from \textup{[DG70]}, Section III.3.7. Let
\(k\) be a field. We begin with a lemma
\textup{(cf.\ \textup{[An73]}, 2.3.3.2)} that will be used
later \textup{(see \hyperref[prop:sga3-VIB-12.10]{12.10})}.

\medskip
\phantomsection\label{lem:sga3-VIB-11.18.1}%
\noindent\textbf{Lemma 11.18.1.---}
\emph{Let \(u:G\to G'\) be a morphism of \(k\)-groups and
put \(N=\operatorname{Ker}(u)\). Suppose that \(G'\) is
affine and that \(G\) is of finite type.}
\begin{enumerate}
\item[\textup{\SGAInlineRefTarget{lem:sga3-VIB-11.18.1:item:i}(i)}]\phantomsection
  \label{concl:sga3-VIB-11.18.1-i}%
  \emph{The morphism
  \(\tau:G/N\to G'\) is a closed immersion. In particular,
  \(G/N\) is affine.}

\item[\textup{(ii)}]\phantomsection
  \label{concl:sga3-VIB-11.18.1-ii}%
  \emph{If, moreover, the morphism
  \(u^\natural:\mathcal O(G')\to\mathcal O(G)\) is
  injective, then \(\tau\) is an isomorphism. Thus \(G'\)
  is of finite type and \(u\) is faithfully flat.}
\end{enumerate}

Indeed, by \hyperref[prop:sga3-VIB-11.13]{11.13},
\(G'\) is the projective limit of a filtered system of
affine algebraic \(k\)-groups \(G'_i\). Let
\[
  u_i:G\longrightarrow G'\longrightarrow G'_i
\]
be the composite, and let \(N_i\) be its kernel. The
\(N_i\) form a filtered decreasing system of closed
subgroups of \(G\), whose intersection is \(N\). Since
\(G\) is Noetherian, there is an index \(i\) such that
\[
  N=N_i.
\]
The groups \(G\) and \(G'_i\) are of finite type.
Therefore, by
\hyperref[thm:sga3-VIA-3.2]{VI\(_A\), 3.2} and
\hyperref[prop:sga3-VIA-5.4.1]{5.4.1}, the quotient
\(G/N\) is a \(k\)-group of finite type, and \(u_i\) is
the composite of the faithfully flat projection
\[
  p:G\longrightarrow G/N
\]
and a closed immersion
\[
  \tau_i:G/N\hookrightarrow G'_i.
\]

Consider the commutative diagram
\begin{center}
  \phantomsection\label{diag:sga3-VIB-11.18.1-closed-immersion}%
  \begin{tikzcd}[row sep=large,column sep=large]
    G \arrow[r,"u"] \arrow[d,"p"'] &
    G' \arrow[d,"q_i"] \\
    G/N \arrow[r,"\tau_i"'] \arrow[ur,"\tau"] &
    G'_i\mathpunct{.}
  \end{tikzcd}
\end{center}
Since
\[
  q_i\circ\tau=\tau_i
\]
is a closed immersion and \(q_i\) is separated
\textup{(\(G'\) is separated; cf.\
\hyperref[sga3:VIA:0.3]{VI\(_A\), 0.3})}, \(\tau\) is a
closed immersion
\textup{(cf.\ EGA I, 5.4.4)}. Hence \(G/N\) is affine and
\[
  \tau^\natural:\mathcal O(G')\longrightarrow
    \mathcal O(G/N)
\]
is surjective, proving \textup{\SGARefLink{lem:sga3-VIB-11.18.1:item:i}{(i)}}.

If \(u^\natural:\mathcal O(G')\to\mathcal O(G)\) is also
injective, then so is \(\tau^\natural\). It follows that
\(\tau^\natural\) is an isomorphism, and hence so is
\(\tau\), since \(G'\) and \(G/N\) are affine. This proves
\textup{\SGARefLink{concl:sga3-VIB-11.18.1-i}{(ii)}}.

\medskip
\phantomsection\label{thm:sga3-VIB-11.18.2}%
\noindent\textbf{Theorem 11.18.2.---}
\emph{Let \(G\) and \(G'\) be affine \(k\)-groups, with
algebras \(A\) and \(A'\), and let
\(u:G\to G'\) be a morphism of \(k\)-groups. Put
\(N=\operatorname{Ker}(u)\), and let
\(\varphi:A'\to A\) be the morphism induced by \(u\).}
\begin{enumerate}
\item[\textup{\SGAInlineRefTarget{thm:sga3-VIB-11.18.2:item:i}(i)}]\phantomsection
  \label{concl:sga3-VIB-11.18.2-i}%
  \emph{If \(\varphi\) is injective, then \(u\) is
  faithfully flat and identifies \(G'\) with \(G/N\).}

\item[\textup{\SGAInlineRefTarget{thm:sga3-VIB-11.18.2:item:ii}(ii)}]\phantomsection
  \label{concl:sga3-VIB-11.18.2-ii}%
  \emph{One has
  \(G/N=\operatorname{Spec}(B)\), where
  \(B=\varphi(A')\). Thus \(u\) is the composite of the
  faithfully flat morphism \(G\to G/N\), corresponding to
  the inclusion \(B\hookrightarrow A\), and the closed
  immersion \(G/N\hookrightarrow G'\), corresponding to
  the surjection \(A'\to B\). Moreover, \(N\) is defined
  in \(G\) by the ideal \(AB_+\), where \(B_+\) denotes
  the augmentation ideal of \(B\).}

\item[\textup{\SGAInlineRefTarget{thm:sga3-VIB-11.18.2:item:iii}(iii)}]\phantomsection
  \label{concl:sga3-VIB-11.18.2-iii}%
  \emph{In particular, if \(u\) is a monomorphism, it is
  a closed immersion.}
\end{enumerate}

\emph{Proof.}
\textup{(i)}
Suppose that \(\varphi\) is injective and identify \(A'\)
with a sub-Hopf algebra of \(A\). By
\hyperref[prop:sga3-VIB-11.13]{11.13}, \(A\) is the
filtered union of finitely generated sub-Hopf algebras
\[
  A_i=\mathcal O(G_i).
\]
Put
\[
  A'_i=A'\cap A_i,
  \qquad
  G'_i=\operatorname{Spec}(A'_i),
\]
and let \(N_i\) be the kernel of the morphism
\[
  G_i\longrightarrow G'_i
\]
induced by \(A'_i\hookrightarrow A_i\). By \hyperref[lem:sga3-VIB-11.18.1]{Lemma 11.18.1},
\[
  G_i/N_i\simeq G'_i.
\]
Thus, for each \(i\), there is an exact sequence
\phantomsection\label{seq:sga3-VIB-11.18.2-stage}%
\[
  1\longrightarrow N_i\longrightarrow G_i
   \xrightarrow{\ p_i\ }G'_i\longrightarrow1,
\]
where \(p_i\) is faithfully flat. Since
\[
  G=\varprojlim_iG_i,
  \qquad
  G'=\varprojlim_iG'_i,
\]
we obtain an exact sequence
\begingroup
\renewcommand{\thefootnote}{\fnsymbol{footnote}}%
\footnote[2]{\label{note:sga3-VIB-11.18.2-limit-target-source-reading}
Source-reading note: the French re-edition prints \(G/N\)
as the final term of this limiting exact sequence, and the
Reinhold comparison lineage repeats it. At this point the
argument is proving that \(G'\) represents \(G/N\); taking
the projective limits of the preceding sequences gives
\(G'\) as the final term.}
\endgroup
\phantomsection\label{seq:sga3-VIB-11.18.2-limit}%
\[
  1\longrightarrow N\longrightarrow G
   \xrightarrow{\ p\ }G',
\]
where \(N=\varprojlim_iN_i\). By \hyperref[lem:sga3-VIB-11.17.1]{Lemma 11.17.1}, \(p\) is
faithfully flat \textup{(and affine)}. Therefore \(G'\)
represents the \textup{(fpqc)} sheaf quotient of \(G\) by
\(N\), proving \textup{\SGARefLink{thm:sga3-VIB-11.18.2:item:i}{(i)}}.

In general, \(B=\varphi(A')\) is a sub-Hopf algebra of
\(A\). Let
\[
  H=\operatorname{Spec}(B),
\]
and let \(N'\) be the kernel of the morphism \(G\to H\)
induced by \(B\hookrightarrow A\). By \textup{\SGARefLink{thm:sga3-VIB-11.18.2:item:i}{(i)}}, \(H\)
is identified with \(G/N'\), and \(u\) is the composite
of the projection \(G\to G/N'\) and the closed immersion
\[
  G/N'\hookrightarrow G'
\]
induced by the surjection \(A'\to B\). Hence \(N'=N\).
Moreover, by
\hyperref[concl:sga3-VIB-9.2-ii]{9.2\textup{(ii)}}, the square
\begin{center}
  \phantomsection\label{diag:sga3-VIB-11.18.2-cartesian}%
  \begin{tikzcd}[row sep=large,column sep=large]
    N \arrow[r] \arrow[d] &
    G \arrow[d,"p"] \\
    \operatorname{Spec}(k) \arrow[r,"\varepsilon"'] &
    H
  \end{tikzcd}
\end{center}
is cartesian. Here \(\varepsilon\) is the unit section of
\(H\), corresponding to the augmentation \(B\to k\).\footnote{Canon
source note (SGA-SRC-000862): the re-edition prints \(G'\) in the lower
right corner and calls \(\varepsilon\) its unit. The augmentation
\(B\to k\) instead defines the unit of \(H=\operatorname{Spec}(B)=G/N\).}
It follows that \(N\) is defined in \(G\) by the ideal
\(AB_+\). This proves \textup{\SGARefLink{thm:sga3-VIB-11.18.2:item:ii}{(ii)}}, and \textup{\SGARefLink{thm:sga3-VIB-11.18.2:item:iii}{(iii)}}
follows.

\medskip
\phantomsection\label{rem:sga3-VIB-11.18.3}%
\noindent\textbf{Remark 11.18.3.---}
Let \(G\) be an affine \(k\)-group and \(N\) an invariant
\(k\)-subgroup. Since
\[
  p:G\longrightarrow G/N
\]
is faithfully flat and quasi-compact,
\(\mathcal O(G/N)\) is, by IV, \SGARefLink{sga3:iv:remark:3:3:3:2}{3.3.3.2}, the subalgebra of
\(\mathcal O(G)\) formed by the functions \(\phi\) that
are right \(N\)-invariant; that is,
\[
  \phi(gh)=\phi(g)
\]
for every \(k\)-scheme \(S\), every \(g\in G(S)\), and
every \(h\in N(S)\). Let \(J\) be the ideal of
\[
  A=\mathcal O(G)
\]
defining \(N\). The preceding condition is equivalent to
\phantomsection\label{form:sga3-VIB-11.18.3-invariants}%
\[
  \Delta(\phi)-\phi\otimes1
    \in\mathcal O(G)\otimes J,
\]
where \(\Delta\) is the comultiplication of \(A\).

The preceding theorem may therefore be reformulated in
terms of Hopf algebras as follows.

\medskip
\phantomsection\label{cor:sga3-VIB-11.18.4}%
\noindent\textbf{Corollary 11.18.4.---}
\emph{Let \(k\) be a field, let \(A\) be a commutative
\(k\)-Hopf algebra, and put \(G=\operatorname{Spec}(A)\).}
\begin{enumerate}
\item[\textup{\SGAInlineRefTarget{cor:sga3-VIB-11.18.4:item:i}(i)}]\phantomsection
  \label{concl:sga3-VIB-11.18.4-i}%
  \emph{If \(B\) is a sub-Hopf algebra of \(A\), then
  \(A\) is faithfully flat over \(B\).}

\item[\textup{(ii)}]\phantomsection
  \label{concl:sga3-VIB-11.18.4-ii}%
  \emph{The map}
  \[
    N\longmapsto\mathcal O(G/N)
  \]
  \emph{is a bijection from the set of invariant
  subgroups of \(G\) to the set of sub-Hopf algebras of
  \(A\). Its inverse is}
  \[
    B\longmapsto\operatorname{Spec}(A/AB_+).
  \]
  \emph{Moreover, if \(J\) is the ideal of \(A\) defining
  \(N\), then}
  \phantomsection\label{form:sga3-VIB-11.18.4-invariant-ring}%
  \[
    \mathcal O(G/N)
      =
    \{\,x\in A\mid
      \Delta(x)-x\otimes1\in A\otimes J\,\}.
  \]
\end{enumerate}

% END INCLUDED SGA3_English_current_buildable_source_20260823/source/inputs/VI/tex_reference_v2/components/87_expose_VIB_lemma11170_through_cor11184_loop2_r1_en.tex

% BEGIN INCLUDED SGA3_English_current_buildable_source_20260823/source/inputs/VI/tex_reference_v2/components/88_expose_VIB_rem11185_through_section12_end_loop2_r1_en.tex
% SGA 3, Expose VI_B, Loop 1 English reconstruction.
% Loop 2 R1: the source-bound base-change crop is replaced by a
% faithful native tikz-cd square; prose, mathematics, and labels
% otherwise remain identical to the reviewed Loop-1 component.
% Scope: Remarks 11.18.5 and all of Section 12, through
% Remarks 12.10.2, including editor's notes 135--138.
% Authority: Polo--Gille Exp6B-13oct24.pdf, component-local
% pages 99--105 (combined-reader pages 440--446; re-edition
% headers 429--435). Opening Remarks 11.18.5 remain under
% original-print marker 411. Editor-added Section 12 has no
% original-print pagination.
% Section 13 begins on local page 105 / combined page 446 /
% header 435 and is excluded; it too is editor-added.
% OCR is locator-only. Reinhold e7a259f English Markdown is
% comparison-only and is credited in the component report.
% Loop 2 replaces the source-bound cartesian-square crop in the proof
% of Theorem 12.2 with a faithful native LaTeX reconstruction; it is closed.

\medskip
\phantomsection\label{rem:sga3-VIB-11.18.5}%
\noindent\textbf{Remarks 11.18.5.---}
\begin{enumerate}
\item[\textup{(a)}]\phantomsection\label{rem:sga3-VIB-11.18.5-a}%
  One consequence of the preceding theorem is that the category
  of commutative affine \(k\)-groups is abelian. For this, as well
  as for other results on affine \(k\)-groups, we refer to
  \textup{[DG70]}, \S~III.3, \SGARefLink{prop:sga3-VIB-7.4}{7.4}--\SGARefLink{prop:sga3-VIB-7.8}{7.8}.

\item[\textup{(b)}]\phantomsection\label{rem:sga3-VIB-11.18.5-b}%
  Finally, let us point out that M.~Takeuchi gave another proof
  of the results of \hyperref[thm:sga3-VIB-11.17]{11.17}--\hyperref[cor:sga3-VIB-11.18.4]{11.18.4}; see \textup{[Ta72]}, \S~5.
  He moreover strengthened \hyperref[concl:sga3-VIB-11.18.4-i]{11.18.4}\textup{\SGARefLink{cor:sga3-VIB-11.18.4:item:i}{(i)}} above by showing
  that \(A\) is even a projective \(B\)-module; see
  \textup{[Ta79]}, \SGARefLink{sec:sga3-VIB-5}{Theorem~5} \textup{(see also [MW94],
  Theorem~3.6)}.
\end{enumerate}

\section*{12. Complements on \(G_{\mathrm{af}}\) and the
  ``anti-affine'' groups}
\phantomsection\label{sec:sga3-VIB-12}%
\addcontentsline{toc}{section}{12. Complements on Gaf and the
  ``anti-affine'' groups}

\footnote{\label{note:sga3-VIB-135}Editor's note: this
section has been added.}
We begin with the following lemma, which extends
\hyperref[lem:sga3-VIB-11.18.1]{11.18.1} to the case where
\(G\) is not assumed to be of finite type.%
\footnote{\label{note:sga3-VIB-136}Editor's note: this
lemma was communicated to us by M.~Raynaud; it will be used
in the proof of Proposition~\hyperref[prop:sga3-VIB-12.9]{12.9}.}

\medskip
\phantomsection\label{lem:sga3-VIB-12.1}%
\noindent\textbf{Lemma 12.1.---}
\emph{Let \(k\) be a field and let
\(u:G\to H\) be a monomorphism of \(k\)-groups, with \(H\)
affine. Suppose that \(u\) is quasi-compact. Then \(u\) is a
closed immersion.}

By \textup{(fpqc)} descent, we may suppose that \(k\) is
algebraically closed. By
\hyperref[prop:sga3-VIA-6.4]{VI\(_A\), 6.4}, the closed image
\(I\) of \(u\) is a closed group subscheme of \(H\), hence is
still affine. Replacing \(H\) by \(I\), we may therefore suppose
that \(u\) is schematically dominant. Since \(k\) is
algebraically closed, \(H'=H_{\mathrm{red}}\) is a group
subscheme of \(H\). Put \(G'=G\times_HH'\). The morphism
\[
  u':G'\longrightarrow H'
\]
deduced from \(u\) by base change is a quasi-compact
monomorphism and is dominant \textup{(the underlying continuous
map being the same for \(u\) and \(u'\))}. Hence, by
\hyperref[cor:sga3-VIA-6.2]{VI\(_A\), 6.2}, \(u'\) is faithfully
flat. It is therefore a quasi-compact faithfully flat
monomorphism, and thus an isomorphism \textup{(cf.\ IV, \SGARefLink{sga3:iv:remark:1:14}{1.14})}.

Consequently \(u:G\to H\) is a homeomorphism, and hence is
affine by \hyperref[lem:sga3-VIB-2.9.1]{2.9.1}. Thus \(G\) is
affine, and \(u\) is a closed immersion by
\hyperref[thm:sga3-VIB-11.18.2]{11.18.2\textup{(iii)}}.

\medskip
\phantomsection\label{thm:sga3-VIB-12.2}%
\noindent\textbf{Theorem 12.2.---}
\emph{Let \(G\) be an algebraic \(k\)-group. Denote by
\(\rho:G\to G_{\mathrm{af}}\) the canonical morphism and by
\(N\) its kernel.}
\begin{enumerate}
\item[\textup{(i)}]\phantomsection\label{thm:sga3-VIB-12.2-i}%
  \emph{The canonical morphism
  \(G/N\to G_{\mathrm{af}}\) is an isomorphism. Consequently
  \(G_{\mathrm{af}}\) is an affine algebraic group and
  \(\rho\) is faithfully flat.}

\item[\textup{(ii)}]\phantomsection\label{thm:sga3-VIB-12.2-ii}%
  \emph{There is a canonical isomorphism
  \((G/N)_{\mathrm{af}}=G_{\mathrm{af}}\).}

\item[\textup{(iii)}]\phantomsection\label{thm:sga3-VIB-12.2-iii}%
  \emph{\(N\) is a characteristic subgroup of \(G\).}

\item[\textup{(iv)}]\phantomsection\label{thm:sga3-VIB-12.2-iv}%
  \emph{\(\mathcal O(N)=k\).}

\item[\textup{(v)}]\phantomsection\label{thm:sga3-VIB-12.2-v}%
  \emph{\(N\) is smooth, connected, and commutative.}
\end{enumerate}

Assertion~\hyperref[thm:sga3-VIB-12.2-i]{\textup{(i)}} is a
special case of
\hyperref[lem:sga3-VIB-11.18.1]{11.18.1}, and
\hyperref[thm:sga3-VIB-12.2-ii]{\textup{(ii)}} follows from the
universal properties of \(G_{\mathrm{af}}\) and
\((G/N)_{\mathrm{af}}\).

Let us prove
\hyperref[thm:sga3-VIB-12.2-iii]{\textup{(iii)}}. For an
arbitrary \(k\)-scheme
\(\pi:S\to\operatorname{Spec}k\), consider the cartesian square
\begin{center}
  \phantomsection\label{diag:sga3-VIB-12.2-base-change}%
  \begin{tikzcd}[row sep=large,column sep=large]
    G_S \arrow[r] \arrow[d,"q"'] &
    G \arrow[d,"p"] \\
    S \arrow[r,"\pi"] &
    \operatorname{Spec}k\mathpunct{,}
  \end{tikzcd}
\end{center}
Since \(p\) is quasi-compact and separated and \(\pi\) is flat,
EGA~III, 1.4.15 and EGA~IV\(_1\), 1.7.21 give
\phantomsection\label{form:sga3-VIB-12.2-base-change-global-sections}%
\[
  q_*(\mathcal O_{G_S})
   =\pi^*(\mathcal O(G))
   =\pi^*(\mathcal O(G_{\mathrm{af}})).
\]
Hence EGA~II, 1.5.2 gives
\[
  (G_S)_{\mathrm{af}}=(G_{\mathrm{af}})_S.
\]
Thus \(N_S\), being the kernel of the canonical morphism
\(G_S\to(G_S)_{\mathrm{af}}\), is invariant under every
automorphism of \(G_S\); in other words, \(N\) is a
characteristic subgroup of \(G\).

To prove \hyperref[thm:sga3-VIB-12.2-iv]{\textup{(iv)}}, put
\[
  N'=\operatorname{Ker}(N\longrightarrow N_{\mathrm{af}}).
\]
By \hyperref[thm:sga3-VIB-12.2-ii]{\textup{(ii)}}, this is an
invariant subgroup of \(G\). Since \(N\) is algebraic
\textup{(being a closed subgroup of \(G\))},
\hyperref[thm:sga3-VIB-12.2-i]{\textup{(i)}} gives
\(N/N'\simeq N_{\mathrm{af}}\). Moreover, by VI\(_A\), \hyperref[thm:sga3-VIA-3.2]{3.2}
and \SGARefLink{par:sga3-VIA-5.3.2}{5.3.2}, there is an isomorphism of \(k\)-groups
\phantomsection\label{form:sga3-VIB-12.2-double-quotient}%
\[
  (G/N')/N_{\mathrm{af}}
    \simeq (G/N')/(N/N')
    \simeq G/N.
\]
Since \(N_{\mathrm{af}}\) is affine, the projection
\(G/N'\to G/N\) is affine by
\hyperref[prop:sga3-VIB-9.2]{9.2\textup{(vii)}}; and since
\(G/N=G_{\mathrm{af}}\) is affine, \(G/N'\) is affine as well.
The universal property of \(G_{\mathrm{af}}\) therefore makes
the projection
\[
  p':G\longrightarrow G/N'
\]
factor through \(G_{\mathrm{af}}=G/N\). Hence \(N\subset N'\),
and consequently \(N=N'\). Thus \(N_{\mathrm{af}}\) is the
trivial group, and \(\mathcal O(N)=k\).

Finally,
\hyperref[thm:sga3-VIB-12.2-v]{\textup{(v)}} follows from the
next lemma.

\medskip
\phantomsection\label{lem:sga3-VIB-12.3}%
\noindent\textbf{Lemma 12.3.---}
\emph{Let \(k\) be a field and let \(N\) be an algebraic
\(k\)-group such that \(\mathcal O(N)=k\). Then \(N\) is
smooth, connected, and commutative.}

Indeed, we may suppose that \(k\) is algebraically closed.
Then
\[
  H=N_{\mathrm{red}}^0
\]
is a \(k\)-subgroup of \(N\), and the quotient \(k\)-scheme
\(X=N/H\) is finite \textup{(hence affine)} over \(k\), by
\hyperref[prop:sga3-VIA-5.5.1]{VI\(_A\), 5.5.1} and
\hyperref[prop:sga3-VIA-5.6.1]{5.6.1}. Since
\(p:N\to X\) is faithfully flat,
\[
  \mathcal O(X)\subset\mathcal O(N)=k.
\]
It follows that \(N=N_{\mathrm{red}}^0\), and hence that \(N\)
is smooth \textup{(VI\(_A\), \hyperref[prop:sga3-VIA-1.3.1]{1.3.1})} and connected.

Let \(Z\) be the center of \(N\). By
\hyperref[cor:sga3-VIB-6.2.6]{6.2.6}, \(N/Z\) is affine, and
the preceding argument gives \(\mathcal O(N/Z)=k\). Hence
\(N=Z\). This proves the lemma and completes the proof of
\hyperref[thm:sga3-VIB-12.2]{12.2}.

We also record, without proof, the following theorem. Recall
that an \emph{abelian variety} over a field \(k\) is a proper,
smooth, connected \(k\)-group scheme.

\medskip
\phantomsection\label{thm:sga3-VIB-12.4}%
\noindent\textbf{Theorem 12.4 \textup{(Chevalley)}.---}
\emph{Let \(k\) be a perfect field and let \(G\) be a smooth,
connected algebraic \(k\)-group. Then there exists a smooth,
connected affine \(k\)-subgroup \(L\), invariant in \(G\),
such that \(G/L\) is an abelian variety. Moreover, \(L\) is
unique and its formation commutes with extension of the base
field.}

\medskip
\phantomsection\label{rem:sga3-VIB-12.5}%
\noindent\textbf{Remarks 12.5.---}
\begin{enumerate}
\item[\textup{(1)}]\phantomsection\label{rem:sga3-VIB-12.5-1}%
  This theorem was announced in 1953 by C.~Chevalley, who
  published his proof in 1960 \textup{([Ch60])}. Meanwhile,
  other proofs were obtained independently by I.~Barsotti and
  M.~Rosenlicht \textup{([Ba55, Ro56])}; see \textup{[Se99]}
  for historical comments.

\item[\textup{(2)}]\phantomsection\label{rem:sga3-VIB-12.5-2}%
  A modern version \textup{(that is, in the language of
  schemes)} of Chevalley's proof was given by B.~Conrad
  \textup{([Co02])}. \textup{(In \emph{loc.\ cit.},
  ``algebraic group'' means a smooth connected \(k\)-group
  scheme.)}

\item[\textup{(3)}]\phantomsection\label{rem:sga3-VIB-12.5-3}%
  A modern version of Rosenlicht's proof, on the other hand,
  was given by Ngô B.-C.\ in a course at Orsay in 2005--2006.

\item[\textup{(4)}]\phantomsection\label{rem:sga3-VIB-12.5-4}%
  If the hypothesis that \(k\) be perfect is omitted, there
  still exists a smallest invariant connected affine subgroup
  \(L\), not necessarily smooth, such that \(G/L\) is an
  abelian variety \textup{([BLR], \S~\SGARefLink{prop:sga3-VIB-9.2}{9.2}, Theorem~1)}.

\item[\textup{(5)}]\phantomsection\label{rem:sga3-VIB-12.5-5}%
  One may also omit the hypothesis that \(G\) be smooth over
  \(k\). Indeed, by VII\(_A\), \SGARefLink{prop:sga3-VIB-8.3}{8.3}, there is an integer
  \(n\geq1\) such that the quotient
  \[
    G'=G/(\operatorname{Fr}^nG)
  \]
  is smooth. Then \(G'\) contains a subgroup \(L'\) as in
  \hyperref[rem:sga3-VIB-12.5-4]{\textup{(4)}} above, and the
  inverse image of \(L'\) in \(G\) again has the same
  properties. Thus, for every connected algebraic group \(G\)
  over a field \(k\), there is an exact sequence
  \phantomsection\label{form:sga3-VIB-12.5-connected-exact-sequence}%
  \[
    1\longrightarrow H\longrightarrow G
      \longrightarrow A\longrightarrow1,
  \]
  where \(H\) is an affine \(k\)-group and \(A\) is an abelian
  \(k\)-variety. Moreover, by \textup{[Per76]},
  Corollary~4.2.9, such an exact sequence exists for every
  connected \(k\)-group \(G\), not necessarily algebraic.

\item[\textup{(6)}]\phantomsection\label{rem:sga3-VIB-12.5-6}%
  Let \(k\) be an algebraically closed field and let \(G\) be
  the semidirect product of an elliptic curve \(E\) by the
  constant \(k\)-group \(\{\pm1\}_k\), for the action defined
  by \((-1)\cdot x=-x\). In this case, if \(L\) is a closed
  invariant subgroup of \(G\) such that \(G/L\) is connected,
  then \(L=G\).
\end{enumerate}

\medskip
\phantomsection\label{rem:sga3-VIB-12.6}%
\noindent\textbf{Remark 12.6.---}
Following \textup{[Br09]}, a \(k\)-group \(N\) will be called
\emph{anti-affine} if \(\mathcal O(N)=k\). By
\hyperref[lem:sga3-VIB-12.3]{12.3} and
\hyperref[thm:sga3-VIB-12.4]{12.4}, if \(k\) is perfect, every
anti-affine algebraic \(k\)-group is an extension of an
abelian variety by a smooth, connected, commutative affine
algebraic \(k\)-group. For the precise structure of
anti-affine algebraic groups over a perfect field, and various
consequences, see the recent articles of M.~Brion and
C.\ \& F.~Sancho de Salas \textup{([Br09, SS09])}.

To conclude this section, we shall prove two results due to
M.~Raynaud, the first being
\hyperref[rem:sga3-VIB-11.11.1]{Remark 11.11.1}, and the
second \SGARefLink{prop:sga3-VIB-2.1}{Proposition~2.1} of Exposé~XVII, Appendix~III. We shall
need the following lemma,%
\footnote{\label{note:sga3-VIB-137}Editor's note: this is a
version improved by O.~Gabber of a statement communicated by
M.~Raynaud.}
which improves, for a complete discrete valuation ring \(R\),
the flatness criteria given in \textup{[BAC]}, \S~III.5.

\medskip
\phantomsection\label{lem:sga3-VIB-12.7}%
\noindent\textbf{Lemma 12.7.---}
\emph{Let \(R\) be a discrete valuation ring, \(K\) its
fraction field, and \(\pi\) a uniformizer. Let \(A\) be a flat
\(R\)-algebra and let \(M\) be an \(A\)-module flat over \(R\).
Suppose that}
\begin{enumerate}
\item[\textup{(i)}]\phantomsection\label{lem:sga3-VIB-12.7-i}%
  \emph{\(M/\pi M\) is flat over
  \(\overline A=A/\pi A\),}

\item[\textup{(ii)}]\phantomsection\label{lem:sga3-VIB-12.7-ii}%
  \emph{\(M\otimes_RK\) is flat over \(A\otimes_RK\).}
\end{enumerate}
\emph{Then \(M\) is flat over \(A\).}

By the flatness criterion in the nilpotent case
\textup{(cf.\ [BAC], III \S~\SGARefLink{prop:sga3-VIB-5.2}{5.2}, Theorem~1)},
\[
  M/\pi^nM
\]
is flat over \(A/\pi^nA\) for every
\(n\in\mathbb N^*\). It then follows from \textup{[RG71]},
II, Lemma~1.4.2.1, that \(M\) is flat over \(A\). For the
reader's convenience, let us briefly indicate the proof. Put
\[
  s=\pi^n,\qquad P=M/sM.
\]
Since \(P\) is flat over \(A/sA\), and \(A/sA\) has
projective dimension \(1\) over \(A\), the spectral sequence
of composite functors shows that \(P\) has
Tor-dimension at most \(1\) over \(A\). Since \(M\) is flat
over \(R\), and thus has no \(\pi\)-torsion, \(M_K/M\) is the
direct limit of the \(A\)-modules
\[
  \pi^{-n}M/M\simeq M/\pi^nM.
\]
Consequently \(M_K/M\) also has Tor-dimension at most \(1\).
From the exact sequence
\phantomsection\label{form:sga3-VIB-12.7-MK-exact-sequence}%
\[
  0\longrightarrow M\longrightarrow M_K
    \longrightarrow M_K/M\longrightarrow0
\]
and the hypothesis that \(M_K\) is flat over \(A_K\), hence
over \(A\), it follows that \(M\) is flat.

For completeness, let us also give the following simpler
proof, pointed out by O.~Gabber. Let \(I\) be a finitely
generated ideal of \(A\). We must show that
\[
  u:M\otimes_AI\longrightarrow M
\]
is injective. By
\hyperref[lem:sga3-VIB-12.7-ii]{\textup{(ii)}},
\(u\otimes_RK\) is injective, so \(\operatorname{Ker}(u)\)
is an \(R\)-module of \(\pi\)-torsion. It therefore suffices
to show that the \(\pi\)-torsion submodule of
\(M\otimes_AI\) is zero. This submodule is a quotient of
\(\operatorname{Tor}_1^A(M,I/\pi I)\), as is seen by
tensoring with \(M\) the exact sequence
\phantomsection\label{form:sga3-VIB-12.7-I-exact-sequence}%
\[
  0\longrightarrow I\xrightarrow{\ \pi\ }I
    \longrightarrow I/\pi I\longrightarrow0.
\]
On the other hand, since \(M\) has no \(\pi\)-torsion
\textup{(being flat over \(R\))}, one has
\[
  \operatorname{Tor}_i^A(M,\overline A)=0
  \qquad(i\geq1).
\]
Consequently, if \(P_\bullet\) is a projective resolution of
the \(A\)-module \(M\), then
\[
  P_\bullet\otimes_A\overline A
\]
is a projective resolution of the \(\overline A\)-module
\(\overline M=M/\pi M\). Thus, for every
\(\overline A\)-module \(Q\),
\[
  \operatorname{Tor}_i^A(M,Q)
    =\operatorname{Tor}_i^{\overline A}(\overline M,Q),
\]
which vanishes for \(i\geq1\), since \(\overline M\) is flat
over \(\overline A\). Hence
\(\operatorname{Tor}_1^A(M,I/\pi I)=0\), proving the lemma.

\medskip
\phantomsection\label{rem:sga3-VIB-12.8}%
\noindent\textbf{Remark 12.8.---}
Let \(S\) be a regular locally noetherian scheme of dimension
\(1\), and let \(X\) be an \(S\)-scheme that is flat,
quasi-separated, and quasi-compact over \(S\). Then
\(\mathcal A(X)\) is a flat \(\mathcal O_S\)-module. Indeed,
we may suppose that \(S\) is local. Let \(s\) be its closed
point, let \(i:X_s\hookrightarrow X\) be the inclusion, and
let \(\pi\) be a uniformizer of \(R=\mathcal O(S)\). Since
\(X\) is flat over \(S\), there is an exact sequence of
sheaves
\phantomsection\label{form:sga3-VIB-12.8-sheaf-exact-sequence}%
\[
  0\longrightarrow\mathcal O_X
    \xrightarrow{\ \pi\ }\mathcal O_X
    \longrightarrow i_*(\mathcal O_{X_s})
    \longrightarrow0.
\]
Taking global sections shows that \(\mathcal A(X)\) has no
\(\pi\)-torsion as an \(R\)-module, and hence is flat.%
\footnote{\label{note:sga3-VIB-138}Editor's note: this is
true more generally when \(S\) is regular, locally noetherian,
and of dimension at most \(2\); see \textup{[Ray70a]},
\SGARefLink{sga3:VIIB:3.2:prop}{VII~3.2}.}
Moreover, the morphism
\[
  \mathcal A((X_{\mathrm{af}})_s)
    =\mathcal A(X)/\pi\mathcal A(X)
    \longrightarrow\mathcal A(X_s),
\]
induced by \(X\to X_{\mathrm{af}}\), is injective.

We can now prove the following proposition
\textup{(cf.\ \hyperref[rem:sga3-VIB-11.11.1]{Remark 11.11.1})}.

\medskip
\phantomsection\label{prop:sga3-VIB-12.9}%
\noindent\textbf{Proposition 12.9.---}
\emph{Let \(S\) be a regular locally noetherian scheme of
dimension at most \(1\), and let \(G\) be an \(S\)-group
scheme that is flat, quasi-separated, and quasi-compact over
\(S\). Then the following conditions are equivalent:}
\begin{enumerate}
\item[\textup{(i)}]\phantomsection\label{prop:sga3-VIB-12.9-i}%
  \emph{\(G\) is affine over \(S\).}

\item[\textup{(ii)}]\phantomsection\label{prop:sga3-VIB-12.9-ii}%
  \emph{\(G\) is quasi-affine over \(S\).}

\item[\textup{(iii)}]\phantomsection\label{prop:sga3-VIB-12.9-iii}%
  \emph{\(G\) acts faithfully on a quasi-affine and flat
  \(S\)-scheme \(X\).}

\item[\textup{(iv)}]\phantomsection\label{prop:sga3-VIB-12.9-iv}%
  \emph{\(G\) acts linearly and faithfully on a
  quasi-coherent module \(\mathcal E\) that is flat over
  \(S\).}

\item[\textup{(v)}]\phantomsection\label{prop:sga3-VIB-12.9-v}%
  \emph{The morphism \(\rho:G\to G_{\mathrm{af}}\) is a
  monomorphism.}
\end{enumerate}

The implication
\hyperref[prop:sga3-VIB-12.9-i]{\textup{(i)}}\(\Rightarrow\)
\hyperref[prop:sga3-VIB-12.9-ii]{\textup{(ii)}} is clear, as
is
\hyperref[prop:sga3-VIB-12.9-ii]{\textup{(ii)}}\(\Rightarrow\)
\hyperref[prop:sga3-VIB-12.9-iii]{\textup{(iii)}} \textup{(take
\(X=G\))}.

Suppose
\hyperref[prop:sga3-VIB-12.9-iii]{\textup{(iii)}} holds.
Since \(\mathcal A(G)\) and \(\mathcal A(X)\) are flat
\(\mathcal O_S\)-modules, the argument of \hyperref[prop:sga3-VIB-11.11]{11.11} shows that
\(G\) acts faithfully on the right on \(X_{\mathrm{af}}\),
and hence acts linearly and faithfully on the left on the
flat quasi-coherent \(\mathcal O_S\)-module \(\mathcal A(X)\).

If \hyperref[prop:sga3-VIB-12.9-iv]{\textup{(iv)}} holds,
\hyperref[prop:sga3-VIB-11.6.1]{11.6.1\textup{(ii)}} shows
that the monomorphism
\[
  G\longrightarrow\operatorname{Aut}_{\mathcal O_S}(\mathcal E)
\]
factors through \(G_{\mathrm{af}}\). Hence
\(G\to G_{\mathrm{af}}\) is a monomorphism.

Finally, suppose
\hyperref[prop:sga3-VIB-12.9-v]{\textup{(v)}} holds, and let
us prove that \(\rho:G\to G_{\mathrm{af}}\) is an
isomorphism. Replacing \(S\) by one of its connected
components, we may suppose that \(S\) is irreducible, with
generic point \(\eta\). Since the formation of
\(G_{\mathrm{af}}\) commutes with flat base change,
\[
  (G_{\mathrm{af}})_\eta=(G_\eta)_{\mathrm{af}}.
\]
Therefore \(G_\eta\to(G_\eta)_{\mathrm{af}}\) is a
monomorphism, hence a closed immersion by
\hyperref[lem:sga3-VIB-12.1]{12.1}, and consequently an
isomorphism because
\[
  \mathcal O((G_\eta)_{\mathrm{af}})=\mathcal O(G_\eta).
\]
If \(S=\operatorname{Spec}(\kappa(\eta))\), the proof is
complete; we may therefore suppose that \(\dim S=1\).

Let \(s\) be a closed point of \(S\). We shall show that
\[
  G_s\longrightarrow(G_{\mathrm{af}})_s
\]
is an isomorphism and that \(\rho\) is flat at every point
of \(G_s\). We may suppose that \(S\) is local, with closed
point \(s\). The base-change morphism
\[
  \rho_s:G_s\longrightarrow(G_{\mathrm{af}})_s
\]
is a monomorphism, hence a closed immersion by
\hyperref[lem:sga3-VIB-12.1]{12.1}. Thus the morphism
\[
  \mathcal O((G_{\mathrm{af}})_s)
    \longrightarrow\mathcal O(G_s)
\]
induced by \(\rho_s\) is surjective. By the preceding remark
it is also injective, and hence is an isomorphism. In
particular, \(\rho\) is surjective.

It follows from \hyperref[lem:sga3-VIB-12.7]{Lemma 12.7} that
\(\rho:G\to G_{\mathrm{af}}\) is faithfully flat. Since \(G\)
is quasi-compact over \(S\) and \(G_{\mathrm{af}}\) is
separated over \(S\), \(\rho\) is also quasi-compact
\textup{(cf.\ EGA~I, 6.6.4)}. Therefore \(\rho\) is a
quasi-compact faithfully flat monomorphism, and is thus an
isomorphism. This proves the proposition.

Finally, we prove Proposition~2.1 of Exposé~XVII,
Appendix~III. Taking \textup{[Per76]}, Corollary~4.2.5 into
account, we have replaced ``of finite type'' in the
hypotheses by ``quasi-compact and quasi-separated''.
\textup{(If \(G\) is assumed to be of finite type,
\hyperref[lem:sga3-VIB-12.1]{12.1} may be used in place of
\emph{loc.\ cit.})}

\medskip
\phantomsection\label{prop:sga3-VIB-12.10}%
\noindent\textbf{Proposition 12.10.---}
\emph{Let \(S\) be a regular locally noetherian scheme of
dimension at most \(1\), and let \(G\) be a flat,
quasi-compact, and quasi-separated \(S\)-group scheme.}
\begin{enumerate}
\item[\textup{(i)}]\phantomsection\label{prop:sga3-VIB-12.10-i}%
  \emph{The canonical morphism
  \(\rho:G\to G_{\mathrm{af}}\) is faithfully flat and
  quasi-compact. Consequently, \(G_{\mathrm{af}}\) represents
  the \textup{(fpqc)} sheaf quotient of \(G\) by
  \(N=\operatorname{Ker}(\rho)\).}

\item[\textup{(ii)}]\phantomsection\label{prop:sga3-VIB-12.10-ii}%
  \emph{If moreover \(G\) is of finite type over \(S\), then
  \(\rho\) is of finite presentation, \(G_{\mathrm{af}}\)
  represents the \textup{(fppf)} sheaf quotient of \(G\) by
  \(N\), and \(G_{\mathrm{af}}\) is of finite type over \(S\).}

\item[\textup{(iii)}]\phantomsection\label{prop:sga3-VIB-12.10-iii}%
  \emph{Suppose moreover that \(G_\eta\) is affine for every
  maximal point \(\eta\) of \(S\). Then \(N\) is an étale
  \(S\)-group and is the unit group if \(G\) is separated
  over \(S\).}
\end{enumerate}

Since \(G_{\mathrm{af}}\) is affine, and therefore separated
over \(S\), the morphism \(\rho\) is quasi-compact
\textup{(cf.\ EGA~I, 6.6.4)}, and its kernel
\[
  N=\operatorname{Ker}(\rho)
\]
is a closed subgroup of \(G\). Replacing \(S\) by one of its
connected components, we may suppose that \(S\) is
irreducible, with generic point \(\eta\).

To prove
\hyperref[prop:sga3-VIB-12.10-i]{\textup{(i)}} and
\hyperref[prop:sga3-VIB-12.10-ii]{\textup{(ii)}}, it is
enough to show that \(\rho\) is faithfully flat. Indeed,
\SGARefLink{sga3:iv:corollary:5:1:7:1}{Exposé~IV, 5.1.7.1} then shows that \(G_{\mathrm{af}}\)
represents the \textup{(fpqc)} sheaf quotient of \(G\) by
\(N\). If \(G\) is moreover of finite type over \(S\), it is
of finite presentation because \(S\) is locally noetherian;
by \hyperref[prop:sga3-VIB-9.2]{9.2\textup{(xiii)}},
\(\rho\) is then of finite presentation, as is
\(G_{\mathrm{af}}\to S\), and \(\rho\) is therefore a
covering for the \textup{(fppf)} topology.

We may thus suppose that
\[
  S=\operatorname{Spec}(R),
\]
where \(R\) is a discrete valuation ring if \(\dim S=1\), and
is the field \(\kappa(\eta)\) if \(\dim S=0\). Let \(s\) be
the closed point of \(S\). Since formation of
\(G_{\mathrm{af}}\) commutes with flat base change, the
canonical morphism \(G_\eta\to(G_\eta)_{\mathrm{af}}\) is
identified with
\[
  \rho_\eta:G_\eta\longrightarrow(G_{\mathrm{af}})_\eta.
\]
Since
\[
  \mathcal O((G_{\mathrm{af}})_\eta)
    =\mathcal O((G_\eta)_{\mathrm{af}})
    =\mathcal O(G_\eta),
\]
\(\rho_\eta\) is faithfully flat by \textup{[Per76]}, 4.2.5
\textup{(see also the addition VI\(_A\), \hyperref[cor:sga3-VIA-6.6]{6.6})}. If
\(\dim S=1\), the same argument shows that \(\rho_s\) is
faithfully flat, because the morphism
\[
  \mathcal O((G_{\mathrm{af}})_s)
    \longrightarrow\mathcal O(G_s)
\]
is injective by \hyperref[rem:sga3-VIB-12.8]{Remark 12.8}.
Lemma~\hyperref[lem:sga3-VIB-12.7]{12.7} now shows that
\(\rho\) is faithfully flat. This proves
\hyperref[prop:sga3-VIB-12.10-i]{\textup{(i)}} and
\hyperref[prop:sga3-VIB-12.10-ii]{\textup{(ii)}}. In
particular, \(N\) is flat over \(S\).

Suppose now that \(G_\eta\) is affine. Since \(\rho_\eta\)
coincides with the canonical morphism
\[
  G_\eta\longrightarrow(G_\eta)_{\mathrm{af}},
\]
its kernel \(N_\eta\) is the unit group. Let us prove that
\(N\) is étale over \(S\). Since \(N\) is flat over \(S\),
it remains to show that \(N_s\) is étale over \(\kappa(s)\)
for every point \(s\neq\eta\) of \(S\). There is nothing to
prove when \(S=\operatorname{Spec}(\kappa(\eta))\), so we may
suppose that \(S=\operatorname{Spec}(R)\), where \(R\) is a
discrete valuation ring. Let \(s\) be the closed point of
\(S\), let \(K\) be the fraction field of \(R\), and let
\(\pi\) be a uniformizer.

Let \(x\in N_s\), and let \(U_x\) be an affine open
neighborhood of \(x\) in \(N\). Since \(N\) is flat over
\(S\), the intersection \(U_x\cap N_\eta\) is nonempty, and
hence equals \(\{\epsilon(\eta)\}\), where \(\epsilon\)
denotes the unit section. Thus
\[
  A=\mathcal O(U_x)
\]
is a flat \(R\)-algebra such that \(A_K=K\) and
\(\pi^{-1}\notin A\), since \(\pi\) belongs to the maximal
ideal of \(\mathcal O_{U_x,x}\). It follows that \(A=R\), and
hence the projection \(U_x\to S\) is an isomorphism. Thus
\(N\) is étale over \(S\). If \(G\) is also separated over
\(S\), the inverse isomorphism \(S\to U_x\) is the unit
section, since the two coincide on the dense open subset
\(\{\eta\}\) of \(S=\operatorname{Spec}(R)\). Hence \(N\) is
the unit group. This proves the proposition.

In particular, we obtain the following corollary, two other
proofs of which may be found in \textup{[An73]},
\SGARefLink{prop:sga3-VIA-2.3.1}{Proposition~2.3.1}, and \textup{[PY06]}, \SGARefLink{def:sga3-VIB-3.1}{Proposition~3.1}.

\medskip
\phantomsection\label{cor:sga3-VIB-12.10.1}%
\noindent\textbf{Corollary 12.10.1.---}
\emph{Let \(R\) be a discrete valuation ring, \(K\) its
fraction field, and \(G\) a separated, flat \(R\)-group
scheme of finite type over \(R\). If \(G_K\) is affine, then
\(G\) is affine.}

\medskip
\phantomsection\label{rem:sga3-VIB-12.10.2}%
\noindent\textbf{Remarks 12.10.2.---}
\begin{enumerate}
\item[\textup{(a)}]\phantomsection\label{rem:sga3-VIB-12.10.2-a}%
  O.~Gabber has pointed out to us examples in which \(G\) is
  a flat group of finite type over a discrete valuation ring,
  its generic fiber is an abelian variety, and the kernel
  \(N\) of \(G\to G_{\mathrm{af}}\) is not smooth.

\item[\textup{(b)}]\phantomsection\label{rem:sga3-VIB-12.10.2-b}%
  M.~Raynaud, on the other hand, gave an example, with \(S\)
  the affine plane of dimension \(2\) over a field \(k\), of
  a smooth quasi-affine \(S\)-group scheme with affine,
  connected fibers that is not affine over \(S\); see
  \textup{[Ray70a]}, \S~VII.3, p.~116.
\end{enumerate}

% END INCLUDED SGA3_English_current_buildable_source_20260823/source/inputs/VI/tex_reference_v2/components/88_expose_VIB_rem11185_through_section12_end_loop2_r1_en.tex

% BEGIN INCLUDED SGA3_English_current_buildable_source_20260823/source/inputs/VI/tex_reference_v2/components/89_expose_VIB_section13_through_bibliography_loop2_r1_en.tex
% SGA 3, Expose VI_B, Loop 1 English reconstruction.
% Loop 2 R1: the two source-bound crops are replaced by faithful
% native tikz-cd diagrams; prose, mathematics, labels, and the prior
% integration-path normalization otherwise remain unchanged.
% R2 integration successor: source-reviewed prose, mathematics, labels,
% and locators are unchanged; only the two diagram paths are rebound to
% the cumulative reader's shared figures directory.
% Scope: Section 13 in full and the complete VI_B bibliography,
% including editor's notes 139--140.
% Authority: Polo--Gille Exp6B-13oct24.pdf, component-local
% pages 105--111 (combined-reader pages 446--452; re-edition
% headers 435--441; original-print margin markers 420--426).
% Component-local page 112 is the terminal blank and is excluded.
% OCR is locator-only. Reinhold e7a259f English Markdown is a
% credited comparison witness, not source authority.
% The projective system in Example 13.3 and cartesian square in
% Proposition 13.7 are exact source-bound Loop-1 PNG crops.

\section*{13. Flat affine groups over a regular base of
  dimension at most \(2\)}
\phantomsection\label{sec:sga3-VIB-13}%
\addcontentsline{toc}{section}{13. Flat affine groups over a
  regular base of dimension at most 2}

Let us begin by observing\footnote{\label{note:sga3-VIB-139}%
Editor's note: this section has been added.} that the familiar
argument showing that every affine algebraic group over a field
\(k\) is linear, as well as
\hyperref[lem:sga3-VIB-11.12]{Lemma 11.12}, extends to a group
scheme \(G\) that is affine, flat, and of finite type over a
Noetherian regular base scheme \(S\) of dimension at most \(2\).
For \(S\) of dimension at most \(1\), this proves
\hyperref[concl:sga3-VIB-11.11.1-b]{Remark 11.11.1(b)}. The
extension to \(\dim S=2\) rests on the following lemma,
communicated to us by O.~Gabber.

\medskip
\phantomsection\label{lem:sga3-VIB-13.1}%
\noindent\textbf{Lemma 13.1.---}
\emph{Let \(S\) be a normal Noetherian scheme, let
\(\mathcal A\) be a flat quasi-coherent
\(\mathcal O_S\)-module, let \(\mathcal F\) be a finite-type
\(\mathcal O_S\)-submodule of \(\mathcal A\), let
\(\mathcal F^{**}\) be its bidual, and let \(U\) be the
flatness locus of \(\mathcal F\), that is, the set of points
\(s\in S\) such that \(\mathcal F_s\) is a flat
\(\mathcal O_{S,s}\)-module.}
\begin{enumerate}
\item[\textup{(i)}]\phantomsection\label{concl:sga3-VIB-13.1-i}
  \emph{\(U\) is open in \(S\), and}
  \[
    \mathcal F^{**}=j_*j^*(\mathcal F),
  \]
  \emph{where \(j:U\hookrightarrow S\) denotes the
  inclusion.}%
  \begingroup
  \renewcommand{\thefootnote}{\fnsymbol{footnote}}%
  \footnote[1]{\label{note:sga3-VIB-13.1-i-U-ambient}
  Source-reading note: the French re-edition prints
  \(j:U\hookrightarrow X\) in the statement. The ambient
  scheme throughout the lemma is \(S\), and its proof explicitly
  defines \(j:U\hookrightarrow S\); the latter is used here.
  The Reinhold comparison lineage retains the printed \(X\).}%
  \endgroup

\item[\textup{(ii)}]\phantomsection\label{concl:sga3-VIB-13.1-ii}
  \emph{For every quasi-coherent \(\mathcal O_U\)-module
  \(\mathcal E\), the canonical morphism}
  \phantomsection\label{form:sga3-VIB-13.1-projection}%
  \[
    j_*(\mathcal E)\otimes_{\mathcal O_S}\mathcal A
      \longrightarrow
    j_*\bigl(\mathcal E\otimes_{\mathcal O_U}j^*(\mathcal A)\bigr)
  \]
  \emph{is an isomorphism.}

\item[\textup{(iii)}]\phantomsection\label{concl:sga3-VIB-13.1-iii}
  \emph{In particular,}
  \[
    \mathcal V=j_*j^*(\mathcal F)
      \subset \mathcal A=j_*j^*(\mathcal A),
  \]
  \emph{and the canonical morphism}
  \phantomsection\label{form:sga3-VIB-13.1-tensor}%
  \[
    \mathcal V\otimes_{\mathcal O_S}\mathcal A
      \longrightarrow
    j_*j^*\bigl(\mathcal F\otimes_{\mathcal O_S}\mathcal A\bigr)
  \]
  \emph{is an isomorphism.}
\end{enumerate}

\emph{Proof.}
Replacing \(S\) by one of its connected components, we may
suppose that \(S\) is integral. By EGA~IV\(_2\), 2.1.12, the
flatness locus of \(\mathcal F\), namely the set of points
\(s\in S\) such that \(\mathcal F_s\) is a flat
\(\mathcal O_{S,s}\)-module, is an open subset \(U\) of \(S\);
write \(j:U\hookrightarrow S\) for its inclusion. Since
\(\mathcal A_s\) is flat, hence torsion-free, the same is true
of \(\mathcal F_s\). Thus \(U\) contains every point of
codimension at most \(1\). Consequently,
by \hyperref[bib:sga3-VIB-BAC]{\textup{[BAC]}}, VII, \S~4.2,
Corollary to Theorem~1,
\[
  \mathcal F^{**}=j_*j^*(\mathcal F),
\]
and hence there is a monomorphism
\[
  \mathcal F^{**}\longrightarrow j_*j^*(\mathcal A).
\]

The proof of \hyperref[concl:sga3-VIB-13.1-ii]{\textup{(ii)}}
is analogous to that of EGA~III, 1.4.15, recalled in
\hyperref[par:sga3-VIB-11.0]{11.0}. Since \(S\) is normal,
the morphism
\[
  \mathcal O_S\longrightarrow j_*j^*(\mathcal O_S)
\]
is an isomorphism \textup{(cf.\ EGA~IV\(_2\), 5.8.6 and
5.10.5)}. Applying
\hyperref[concl:sga3-VIB-13.1-ii]{\textup{(ii)}} to
\(\mathcal E=\mathcal O_U\) therefore gives
\[
  \mathcal A=j_*j^*(\mathcal A).
\]
Finally, since
\[
  j^*(\mathcal F\otimes_{\mathcal O_S}\mathcal A)
  =
  j^*(\mathcal F)\otimes_{\mathcal O_U}j^*(\mathcal A),
\]
the last assertion of
\hyperref[concl:sga3-VIB-13.1-iii]{\textup{(iii)}} follows
from \hyperref[concl:sga3-VIB-13.1-ii]{\textup{(ii)}} applied
to \(\mathcal E=j^*(\mathcal F)\). This proves the lemma.

Recall also that a finitely generated \(R\)-module \(M\) is
called \emph{reflexive} if the canonical morphism from \(M\)
to its bidual \(M^{**}\) is an isomorphism. If \(R\) is a
Noetherian regular ring of dimension at most \(2\), this
implies that \(M\) is projective. Indeed, for every finitely
generated \(R\)-module \(N\), choose a resolution
\[
  L_1\longrightarrow L_0\longrightarrow N\longrightarrow0,
\]
where \(L_0\) and \(L_1\) are finite free \(R\)-modules.
There is then an exact sequence
\phantomsection\label{seq:sga3-VIB-13.1-dual-resolution}%
\[
  0\longrightarrow N^*
    \longrightarrow L_0^*
    \longrightarrow L_1^*
    \longrightarrow Q
    \longrightarrow0,
\]
where \(Q\) is the cokernel of \(L_0^*\to L_1^*\). Since
\(R\) has homological dimension at most \(2\)
\textup{(\hyperref[bib:sga3-VIB-BAC]{[BAC]}, X, \S~\SGARefLink{prop:sga3-VIB-4.2}{4.2},
Corollary~1 to Theorem~1)}, it follows that \(N^*\) is
projective.

\medskip
\phantomsection\label{prop:sga3-VIB-13.2}%
\noindent\textbf{Proposition 13.2.---}
\emph{Let \(S\) be a Noetherian regular scheme of dimension at
most \(2\), let \(G\) be an affine flat \(S\)-group, and let
\(\mathcal A(G)\) be its affine algebra.}
\begin{enumerate}
\item[\textup{(i)}]\phantomsection\label{concl:sga3-VIB-13.2-i}
  \emph{If \(G\) is of finite type over \(S\), it is
  isomorphic to a closed subgroup of}
  \[
    H=\operatorname{Aut}_{\mathcal O_S}(\mathcal V)
  \]
  \emph{for some locally free \(\mathcal O_S\)-module
  \(\mathcal V\) of finite rank. If, moreover, \(S\) is
  affine, one may take
  \(\mathcal V=\mathcal O_S^{\oplus d}\) for some \(d\), so
  that \(H=\operatorname{GL}_{d,S}\).}

\item[\textup{(ii)}]\phantomsection\label{concl:sga3-VIB-13.2-ii}
  \emph{\(\mathcal A(G)\) is the filtered direct limit of
  flat finite-type \(\mathcal O_S\)-sub-Hopf algebras.}
\end{enumerate}

\emph{Proof.}
Let \(\mathcal B\) be a finite-type
\(\mathcal O_S\)-subalgebra of \(\mathcal A(G)\). Since every
coherent module on an open subset of \(S\) extends to a
coherent module on \(S\) \textup{(EGA~I, 9.4.5)}, there is a
coherent \(\mathcal O_S\)-submodule \(\mathcal M\) of
\(\mathcal B\) that generates \(\mathcal B\) as an
\(\mathcal O_S\)-algebra \textup{(\emph{loc.\ cit.}, 9.6.5)}.
By \hyperref[prop:sga3-VIB-11.10.bis]{11.10.bis},
\(\mathcal M\) is contained in a coherent \(G\)-stable
\(\mathcal O_S\)-submodule \(\mathcal F\).
\textup{(N.B. Since \(G\) is affine over \(S\), the proof of
\emph{loc.\ cit.} is simpler here: one may replace
\(f_*f^*(\mathcal E)\) by
\(\mathcal E\otimes_{\mathcal O_S}\mathcal A(G)\), and so
on.)}

Let \(j:U\hookrightarrow S\) be the inclusion, where \(U\)
is the flatness locus of \(\mathcal F\). By
\hyperref[lem:sga3-VIB-13.1]{Lemma 13.1} and the preceding
observation,
\[
  \mathcal V=j_*j^*(\mathcal F)
\]
is a locally free \(\mathcal O_S\)-submodule of
\(\mathcal A(G)\), and the canonical morphism
\phantomsection\label{form:sga3-VIB-13.2-V-stability}%
\[
  \mathcal V\otimes_{\mathcal O_S}\mathcal A(G)
    \longrightarrow
  j_*j^*\bigl(\mathcal F\otimes_{\mathcal O_S}\mathcal A(G)\bigr)
\]
is an isomorphism. Thus \(\mathcal V\) is a \(G\)-submodule
of \(\mathcal A(G)\). The action of \(G\) on \(\mathcal V\)
induces a morphism of affine \(S\)-groups
\[
  \rho_{\mathcal V}:
  G\longrightarrow
  H=\operatorname{Aut}_{\mathcal O_S}(\mathcal V)
\]
and therefore a morphism of \(\mathcal O_S\)-Hopf algebras
\[
  \varphi_{\mathcal V}:\mathcal A(H)\longrightarrow\mathcal A(G).
\]
Write \(\mathcal A_{\mathcal V}\) for the image of
\(\varphi_{\mathcal V}\). It is the affine algebra of a
closed subgroup \(G_{\mathcal V}\) of \(H\), the closed image
of \(\rho_{\mathcal V}\). We shall show that
\(\mathcal A_{\mathcal V}\) contains \(\mathcal B\).

The question is local on \(S\), so we may suppose that
\(S=\operatorname{Spec}(R)\) and that
\[
  V=\Gamma(S,\mathcal V)
\]
is a free \(R\)-module with basis \(v_1,\ldots,v_n\).
Then \(H\simeq\operatorname{GL}_{n,R}\), and
\(\mathcal A(H)\) is generated as an \(R\)-algebra by the
``matrix coefficients'' \(c_{ij}\) and by \(d^{-1}\), where
\(d\) denotes the determinant. Let \(\Delta\)
\textup{(respectively, \(\varepsilon\))} denote the
comultiplication \textup{(respectively, augmentation)} of
\(A=\Gamma(S,\mathcal A(G))\). For \(j=1,\ldots,n\), write
\phantomsection\label{form:sga3-VIB-13.2-matrix-coefficients}%
\[
  \Delta(v_j)=\sum_{i=1}^{n}v_i\otimes a_{ij}.
\]
Then \(a_{ij}=\varphi(c_{ij})\) belongs to
\(\operatorname{Im}(\varphi)\). Since \(V\) is an
\(R\)-submodule of \(A\), the identity
\[
  (\varepsilon\otimes\mathrm{id}_A)\circ\Delta
    =\mathrm{id}_A
\]
gives
\[
  v_j=\sum_i\varepsilon(v_i)a_{ij}
    \in\operatorname{Im}(\varphi).
\]
As \(V\) contains a set of generators of
\(B=\Gamma(S,\mathcal B)\), it follows that
\[
  B\subset\operatorname{Im}(\varphi).
\]

If \(G\) is of finite type over \(S\), take
\(\mathcal B=\mathcal A(G)\). Then \(\varphi\) is
surjective, so the morphism
\[
  G\longrightarrow
  H=\operatorname{Aut}_{\mathcal O_S}(\mathcal V)
\]
is a closed immersion.

If \(S\) is also affine, there is a locally free
\(\mathcal O_S\)-module \(\mathcal V'\) of finite rank such
that
\[
  \mathcal V\oplus\mathcal V'
    =\mathcal O_S^{\oplus d}.
\]
Regard \(\mathcal V'\) as a trivial \(G\)-module and replace
\(\mathcal V\) by \(\mathcal O_S^{\oplus d}\). This realizes
\(G\) as a closed subgroup of \(\operatorname{GL}_{d,S}\).

Finally, return to an arbitrary affine flat \(S\)-group
\(G\). By EGA~I, 9.4.9, \(\mathcal A(G)\) is the union of
its coherent \(\mathcal O_S\)-submodules \(\mathcal M\),
hence also of the \(\mathcal O_S\)-sub-Hopf algebras
\(\mathcal A_{\mathcal V}\) constructed above. This proves
\hyperref[concl:sga3-VIB-13.2-ii]{\textup{(ii)}}.

\medskip
\phantomsection\label{ex:sga3-VIB-13.3}%
\noindent\textbf{Example 13.3.---}
\emph{Let \(R\) be a discrete valuation ring with uniformizer
\(\pi\) and fraction field \(K\). Consider the filtered
projective system of \(R\)-groups}
\begin{center}
  \phantomsection\label{diag:sga3-VIB-13.3-projective-system}%
  \begin{tikzcd}[column sep=large]
    \cdots \arrow[r] &
    \mathbb G_{a,R} \arrow[r,"\times\pi"] &
    \mathbb G_{a,R} \arrow[r,"\times\pi"] &
    \mathbb G_{a,R}
  \end{tikzcd}
\end{center}
\emph{corresponding to the direct system}
\[
  R[X_0]\longrightarrow R[X_1]\longrightarrow R[X_2]
    \longrightarrow\cdots,
\]
\emph{whose transition morphisms are given by
\(X_n=\pi X_{n+1}\). Its projective limit \(G\) is a flat
affine \(R\)-group scheme, not of finite type, with trivial
special fiber and generic fiber \(\mathbb G_{a,K}\). The
affine \(R\)-algebra \(\mathcal A(G)\) is the subring of
\(K[X]\) consisting of the polynomials whose constant
coefficient lies in \(R\).
\textup{(N.B. We already encountered this example in
\hyperref[rem:sga3-VIB-11.10.1]{Remark 11.10.1}.)}
Furthermore, \(G\) represents the functor assigning to every
\(R\)-algebra \(B\) the set of sequences
\((x_n)_{n\in\mathbb N}\) of elements of \(B\) such that}
\[
  x_n=\pi x_{n+1}\qquad\text{for every \(n\)}.
\]
\emph{\textup{(In particular, every \(x_n\) is indefinitely
\(\pi\)-divisible.)}}

Now let \(S\) be a Noetherian scheme such that every coherent
\(\mathcal O_S\)-module is a quotient of a locally free
\(\mathcal O_S\)-module of finite type. This holds, for
example, when \(S\) is a separated Noetherian regular scheme;
see SGA~6, Expos\'e~II, 2.1.1 and 2.2.7. \textup{(One can
show that every Noetherian regular scheme of dimension at
most \(1\) also has this property. It does not hold, however,
when \(S\) is the affine plane over a field \(k\) with the
origin doubled; see \emph{loc.\ cit.}, 2.2.7.2.)}

\medskip
\phantomsection\label{def:sga3-VIB-13.4}%
\noindent\textbf{Definition 13.4.---}
\emph{Let \(G\) be an affine flat \(S\)-group. Following
R.~W.~Thomason \textup{(\hyperref[bib:sga3-VIB-Th87]{[Th87]},
\SGARefLink{prop:sga3-VIB-2.1}{2.1})}, we say that the pair \((G,S)\) has the
\emph{equivariant resolution property}, or satisfies
\textup{(RE)}, if, for every coherent
\(G\)-\(\mathcal O_S\)-module \(\mathcal F\), there are a
locally free \(G\)-\(\mathcal O_S\)-module \(\mathcal E\) of
finite rank and a \(G\)-equivariant epimorphism}
\[
  \mathcal E\longrightarrow\mathcal F.
\]

In \emph{loc.\ cit.}, \SGARefLink{def:sga3-VIB-3.1}{Theorem~3.1}, Thomason proves the result
below under the hypothesis that \(G\) be essentially free
over \(S\) \textup{(see the remark below)}. Gabber pointed
out to us the simpler proof given here, which does not use
that hypothesis.

\medskip
\phantomsection\label{prop:sga3-VIB-13.5}%
\noindent\textbf{Proposition 13.5.---}
\emph{Let \(S\) be a Noetherian scheme, and let \(G\) be an
affine \(S\)-group, flat and of finite type, with affine
algebra \(\mathcal A(G)\). Suppose that \((G,S)\) satisfies
\textup{(RE)}. Then:}
\begin{enumerate}
\item[\textup{(i)}]\phantomsection\label{concl:sga3-VIB-13.5-i}
  \emph{\(G\) is isomorphic to a closed subgroup of}
  \[
    H=\operatorname{Aut}_{\mathcal O_S}(\mathcal V)
  \]
  \emph{for some locally free \(\mathcal O_S\)-module
  \(\mathcal V\) of finite rank.}

\item[\textup{(ii)}]\phantomsection\label{concl:sga3-VIB-13.5-ii}
  \emph{If, moreover, \(S\) is affine, one may take
  \(\mathcal V=\mathcal O_S^{\oplus d}\) for some \(d\), so
  that \(H=\operatorname{GL}_{d,S}\).}%
  \begingroup
  \renewcommand{\thefootnote}{\fnsymbol{footnote}}%
  \footnote[2]{\label{note:sga3-VIB-13.5-ii-rank-letter}
  Source-reading note: the French statement prints ``for some
  \(n\)'' after \(\mathcal O_S^{\oplus d}\), while its conclusion
  is \(H=\operatorname{GL}_{d,S}\). The same rank is denoted by
  \(d\) here. The Reinhold comparison lineage likewise uses
  \(d\).}%
  \endgroup
\end{enumerate}

The proof is analogous to that of
\hyperref[prop:sga3-VIB-13.2]{13.2}. As there, there is a
coherent \(G\)-stable \(\mathcal O_S\)-submodule
\(\mathcal F\) that generates
\(\mathcal A=\mathcal A(G)\) as an
\(\mathcal O_S\)-algebra. Replacing \(S\) by one of its
connected components, we may suppose that \(S\) is
connected. By \textup{(RE)}, there are a locally free
\(G\)-\(\mathcal O_S\)-module \(\mathcal E\) of rank \(n\)
and an epimorphism of \(\mathcal A\)-comodules
\[
  \pi:\mathcal E\longrightarrow\mathcal F.
\]
Put \(H=\operatorname{Aut}_{\mathcal O_S}(\mathcal E)\).
This is an \(S\)-group scheme locally isomorphic to
\(\operatorname{GL}_{n,S}\). The action of \(G\) on
\(\mathcal E\) induces a morphism of affine \(S\)-groups
\[
  \rho:G\longrightarrow H,
\]
corresponding to a morphism of \(\mathcal O_S\)-Hopf algebras
\[
  \varphi:\mathcal A(H)\longrightarrow\mathcal A(G).
\]
We shall show that \(\rho\) is a closed immersion.

The question is local on \(S\), so suppose that
\(S=\operatorname{Spec}(R)\) and that
\[
  V=\Gamma(S,\mathcal E)
\]
is a free \(R\)-module with basis \(v_1,\ldots,v_n\). Then
\(H\simeq\operatorname{GL}_{n,R}\), and
\[
  B=\Gamma(S,\mathcal A(H))
\]
is generated as an \(R\)-algebra by the matrix coefficients
\(c_{ij}\) and by \(d^{-1}\), where \(d\) denotes the
determinant. Let \(\Delta\)
\textup{(respectively, \(\varepsilon\))} be the
comultiplication \textup{(respectively, augmentation)} of
\[
  A=\Gamma(S,\mathcal A(G)),
\]
let
\[
  \mu:V\longrightarrow V\otimes_R A
\]
be the \(A\)-comodule structure on \(V\), and put
\[
  F=\Gamma(S,\mathcal F).
\]
For \(j=1,\ldots,n\), write
\phantomsection\label{form:sga3-VIB-13.5-mu-matrix}%
\[
  \mu(v_j)=\sum_{i=1}^{n}v_i\otimes a_{ij}.
\]
Then \(\varphi(c_{ij})=a_{ij}\). Since
\(\pi:V\to F\) is a morphism of \(A\)-comodules, one has
\phantomsection\label{form:sga3-VIB-13.5-comodule-map}%
\[
  \sum_{i=1}^{n}\pi(v_i)\otimes a_{ij}
    =
  (\pi\otimes_R\mathrm{id}_A)(\mu(v_j))
    =
  \Delta(\pi(v_j)),
\]
and therefore
\phantomsection\label{form:sga3-VIB-13.5-coefficient-recovery}%
\[
  \begin{aligned}
  \pi(v_j)
    &=(\varepsilon\otimes_R\mathrm{id}_A)\Delta(\pi(v_j))\\
    &=\sum_{i=1}^{n}\varepsilon(\pi(v_i))a_{ij}\\
    &=\varphi\left(
        \sum_{i=1}^{n}\varepsilon(\pi(v_i))c_{ij}
      \right).
  \end{aligned}
\]
Thus \(\varphi(B)\) contains \(\pi(V)=F\), which generates
\(A\) as an \(R\)-algebra. Hence \(\varphi\) is surjective,
and \(\rho\) is a closed immersion. This proves
\hyperref[concl:sga3-VIB-13.5-i]{\textup{(i)}}.

If \(S\) is affine, there is a locally free
\(\mathcal O_S\)-module \(\mathcal V'\) of finite rank such
that
\[
  \mathcal V\oplus\mathcal V'
    =\mathcal O_S^{\oplus d}.
\]
Regard \(\mathcal V'\) as a trivial \(G\)-module and replace
\(\mathcal V\) by \(\mathcal O_S^{\oplus d}\). This realizes
\(G\) as a closed subgroup of \(\operatorname{GL}_{d,S}\)
and completes the proof.

\medskip
\phantomsection\label{rem:sga3-VIB-13.6}%
\noindent\textbf{Remarks 13.6.---}
\begin{enumerate}
\item[\textup{(a)}]\phantomsection\label{rem:sga3-VIB-13.6-a}
  For completeness, let us briefly sketch Thomason's
  argument \textup{(\hyperref[bib:sga3-VIB-Th87]{[Th87]},
  \SGARefLink{def:sga3-VIB-3.1}{Theorem~3.1})}, retaining the preceding notation. The
  \(G\)-equivariant epimorphism
  \[
    \pi:\mathcal E\longrightarrow\mathcal F
  \]
  induces a closed immersion
  \[
    \tau:G\longrightarrow V(\mathcal E)
  \]
  such that
  \[
    \tau(g'g)=\tau(g')\cdot g.
  \]
  \textup{(N.B. \(G\) acts on the right on
  \(V(\mathcal E)\).)}
  There is an isomorphism
  \[
    H\simeq
    \operatorname{Aut}_{\mathcal O_S}
      \bigl(V(\mathcal E)\bigr)^{\mathrm{op}}
  \]
  compatible with the left action of \(G\) on
  \(\mathcal E\) and its right action on \(V(\mathcal E)\).
  Let
  \[
    N=\operatorname{Transp}^{\mathrm{str}}_H
      \bigl(\tau(G),\tau(G)\bigr)
  \]
  be the strict transporter. If \(G\) is essentially free
  over \(S\), it follows from
  \hyperref[ex:sga3-VIB-6.2.4-e]{6.2.4(e)} that \(N\) is a
  closed subgroup scheme of \(H\), hence affine over \(S\).
  Moreover, \(\rho\) factors through a morphism of
  \(S\)-groups
  \[
    \rho':G\longrightarrow N.
  \]
  For every \(S'\to S\) and \(h\in N(S')\), put
  \[
    \pi(h)=\tau(1)\cdot h,
  \]
  where \(1\) is the unit element of \(G(S')\). This defines
  a morphism of \(S\)-schemes
  \[
    \pi:N\longrightarrow\tau(G)
  \]
  that is a retraction of \(\rho'\), after identifying
  \(G\) with \(\tau(G)\). Since \(N\) is separated over
  \(S\), it follows that \(\rho'\) is a closed immersion.

\item[\textup{(b)}]\phantomsection\label{rem:sga3-VIB-13.6-b}
  The proof of \hyperref[bib:sga3-VIB-Th87]{[Th87]},
  \SGARefLink{def:sga3-VIB-3.1}{Theorem~3.1}, appears to require the hypothesis that \(G\)
  be essentially free over \(S\), although that hypothesis
  is not stated there \textup{(the author invokes instead
  the fact that \(H\) is essentially free)}. It is,
  however, satisfied when \(G\) is reductive
  \textup{(cf.\ \SGARefLink{sga3:xxii:cor:5.7.8}{Expos\'e~XXII, 5.7.8})}, and hence in all
  cases considered in \emph{loc.\ cit.}, \SGARefLink{rem:sga3-VIB-3.2}{Corollary~3.2}.
  In particular, Thomason proves in \emph{loc.\ cit.}, \SGARefLink{prop:sga3-VIB-2.5}{2.5},
  that if \(S\) is separated, Noetherian, and regular of
  dimension at most \(2\), if \(G\) is affine and of finite
  presentation, and if \(\mathcal A(G)\) is a locally
  projective \(\mathcal O_S\)-module, then \((G,S)\)
  satisfies \textup{(RE)}. By
  \hyperref[prop:sga3-VIB-13.5]{13.5}, this gives
  \hyperref[prop:sga3-VIB-13.2]{13.2} under a slightly more
  restrictive hypothesis.
\end{enumerate}

To conclude, let us point out that the proof of
\hyperref[bib:sga3-VIB-Th87]{[Th87]}, \SGARefLink{prop:sga3-VIB-2.5}{2.5}, can be simplified
slightly as follows. \textup{(For brevity, we work in the
affine situation \(S=\operatorname{Spec}(R)\).)}

\medskip
\phantomsection\label{prop:sga3-VIB-13.7}%
\noindent\textbf{Proposition 13.7.---}
\emph{Let \(R\) be a Noetherian regular ring of dimension at
most \(2\), let \(C\) be an \(R\)-coalgebra that is
projective as an \(R\)-module, and let \(F\) be a
\(C\)-comodule of finite type over \(R\). Then \(F\) is a
quotient of a \(C\)-comodule \(V\) that is projective and of
finite type over \(R\).}

\emph{Proof.}
Replacing \(\operatorname{Spec}(R)\) by one of its connected
components, we may suppose that \(R\) is integral, with
fraction field \(K\). Let \(\Delta\)
\textup{(respectively, \(\varepsilon\))} denote the
comultiplication \textup{(respectively, augmentation)} of
\(C\), and let
\[
  \rho:F\longrightarrow F\otimes C
\]
be the comodule structure on \(F\). Let
\[
  \pi:W\longrightarrow F
\]
be a surjective morphism, with \(W\) a finite free
\(R\)-module. Give \(W\otimes C\) the comodule structure
defined by \(\mathrm{id}_W\otimes\Delta\), and do the same for
\(F\otimes C\). Then
\(\rho:F\to F\otimes C\) is a morphism of \(C\)-comodules
with retraction \(\mathrm{id}_F\otimes\varepsilon\).

Let \(W'\) be the \(C\)-comodule defined by the following
cartesian square:
\begin{center}
  \phantomsection\label{diag:sga3-VIB-13.7-cartesian-square}%
  \begin{tikzcd}[row sep=large,column sep=large]
    W' \arrow[r] \arrow[d] &
    F \arrow[d,"\rho"] \\
    W\otimes C \arrow[r,"\pi\otimes\mathrm{id}_C"] &
    F\otimes C
  \end{tikzcd}
\end{center}
Thus \(W'\) identifies with the kernel of the morphism
\[
  W\otimes C\longrightarrow(F\otimes C)/\rho(F),
\]
and the projection
\[
  \pi':W'\longrightarrow F,
  \qquad
  x\longmapsto(\pi\otimes\varepsilon)(x),
\]
is surjective. Since \(F\) is a finitely generated
\(R\)-module, there is a \(C\)-subcomodule \(V'\) of \(W'\),
finitely generated over \(R\), such that
\[
  \pi'(V')=F.
\]
Since \(W\otimes C\) is \(R\)-torsion-free, so is \(V'\).
Replacing \(F\) by \(V'\), we may therefore suppose from
the outset that \(F\) is torsion-free.

Apply the preceding construction to this new \(F\), obtaining
\(V'\) as above. Consider the subcomodule \(V\) that is the
kernel of the morphism
\phantomsection\label{form:sga3-VIB-13.7-E-quotient}%
\[
  W\otimes C\longrightarrow
  E=\frac{W\otimes C\otimes K}{V'\otimes K}.
\]
Then \(V\) contains \(V'\), and
\[
  Q=(W\otimes C)/V
\]
is an \(R\)-submodule of the \(K\)-vector space \(E\). Put
\[
  Q'=E/Q.
\]
Since \(W\otimes C\) and \(E\) are flat \(R\)-modules, for
every \(R\)-module \(N\) one obtains
\phantomsection\label{form:sga3-VIB-13.7-Tor}%
\[
  \operatorname{Tor}_1^R(V,N)
    \simeq
  \operatorname{Tor}_2^R(Q,N)
    \simeq
  \operatorname{Tor}_3^R(Q',N).
\]
Because \(R\) is regular of dimension at most \(2\), the
right-hand term vanishes. Hence \(V\) is a flat
\(R\)-module.

It remains to prove that \(V\) is finitely generated. Put
\[
  M=W\otimes C.
\]
By construction,
\[
  V/V'\simeq(M/V')_{\mathrm{tors}},
\]
the \(R\)-torsion submodule of \(M/V'\). Since \(M\) is a
projective \(R\)-module, there is a projective \(R\)-module
\(P\) such that
\[
  M\oplus P=L
\]
is free. Then
\[
  (M/V')_{\mathrm{tors}}
    \simeq
  (L/V')_{\mathrm{tors}}.
\]
Since \(V'\) is finitely generated, there is a direct
summand
\[
  L'\simeq R^n
\]
of \(L\) containing \(V'\). Hence
\[
  (L/V')_{\mathrm{tors}}
    \simeq
  (L'/V')_{\mathrm{tors}},
\]
and the latter is finitely generated because \(L'/V'\) is.
Thus \(V\) is a finitely generated flat \(R\)-module, hence
a finitely generated projective \(R\)-module, since \(R\) is
Noetherian. This proves \hyperref[prop:sga3-VIB-13.7]{Proposition~13.7}.

\section*{Bibliography}
\phantomsection\label{bib:sga3-VIB}%
\addcontentsline{toc}{section}{Bibliography}
\footnote{\label{note:sga3-VIB-140}Editor's note: additional
references cited in this Expos\'e.}

\begin{description}
\item[{[An73]}]\phantomsection\label{bib:sga3-VIB-An73}%
  S.~Anantharaman,
  \emph{Sch\'emas en groupes, espaces homog\`enes et espaces
  alg\'ebriques sur une base de dimension 1},
  M\'em. Soc. Math. France \textbf{33} (1973), 5--79.
\item[{[Ba55]}]\phantomsection\label{bib:sga3-VIB-Ba55}%
  I.~Barsotti,
  \emph{Un teorema di struttura per le variet\`a gruppali},
  Atti Acad. Naz. Lincei Rend. Cl. Sci. Fis. Mat. Nat.
  \textbf{18} (1955), 43--50.
\item[{[BLR]}]\phantomsection\label{bib:sga3-VIB-BLR}%
  S.~Bosch, W.~L\"utkebohmert, and M.~Raynaud,
  \emph{N\'eron Models}, Springer-Verlag, 1990.
\item[{[BAC]}]\phantomsection\label{bib:sga3-VIB-BAC}%
  N.~Bourbaki,
  \emph{Alg\`ebre commutative}, Chapters I--IV, V--VII,
  and X, Masson, 1985 and 1998.
\item[{[BTop]}]\phantomsection\label{bib:sga3-VIB-BTop}%
  N.~Bourbaki,
  \emph{Topologie g\'en\'erale}, Chapters I--IV, Hermann,
  1971.
\item[{[Br09]}]\phantomsection\label{bib:sga3-VIB-Br09}%
  M.~Brion,
  \emph{Anti-affine algebraic groups},
  J. Algebra \textbf{321} (2009), no.~3, 934--952.
\item[{[Ch60]}]\phantomsection\label{bib:sga3-VIB-Ch60}%
  C.~Chevalley,
  \emph{Une d\'emonstration d'un th\'eor\`eme sur les groupes
  alg\'ebriques},
  J. Maths. Pure Appl. \textbf{39} (1960), 307--317.
\item[{[Co02]}]\phantomsection\label{bib:sga3-VIB-Co02}%
  B.~Conrad,
  \emph{A modern proof of Chevalley's theorem on algebraic
  groups},
  J. Ramanujan Math. Soc. \textbf{17} (2002), 1--18.
\item[{[DG70]}]\phantomsection\label{bib:sga3-VIB-DG70}%
  M.~Demazure and P.~Gabriel,
  \emph{Groupes alg\'ebriques}, Masson \& North-Holland,
  1970.
\item[{[Gr73]}]\phantomsection\label{bib:sga3-VIB-Gr73}%
  L.~Gruson,
  \emph{Dimension homologique des modules plats sur un anneau
  noeth\'erien},
  Symposia Mathematica \textbf{XI} (1973), 243--254.
\item[{[Kn71]}]\phantomsection\label{bib:sga3-VIB-Kn71}%
  D.~Knutson,
  \emph{Algebraic spaces},
  Lecture Notes in Mathematics \textbf{203},
  Springer-Verlag, 1971.
\item[{[Ma07]}]\phantomsection\label{bib:sga3-VIB-Ma07}%
  B.~Margaux,
  \emph{Passage to the Limit in Non-Abelian \v{C}ech
  Cohomology},
  J. Lie Theory \textbf{17} (2007), 591--596.
\item[{[MW94]}]\phantomsection\label{bib:sga3-VIB-MW94}%
  A.~Masuoka and D.~Wigner,
  \emph{Faithful flatness of Hopf algebras},
  J. Algebra \textbf{170} (1994), 156--184.
\item[{[Per76]}]\phantomsection\label{bib:sga3-VIB-Per76}%
  D.~Perrin,
  \emph{Approximation des sch\'emas en groupes,
  quasi-compacts sur un corps},
  Bull. Soc. Math. France \textbf{104} (1976), 323--335.
\item[{[Pes66]}]\phantomsection\label{bib:sga3-VIB-Pes66}%
  C.~Peskine,
  \emph{Une g\'en\'eralisation du ``Main Theorem'' de Zariski},
  Bull. Sci. Math. \textbf{90} (1966), 119--127.
\item[{[PY06]}]\phantomsection\label{bib:sga3-VIB-PY06}%
  G.~Prasad and J.-K.~Yu,
  \emph{On quasi-reductive group schemes},
  J. Algebraic Geom. \textbf{15} (2006), 507--549.
\item[{[Ray70a]}]\phantomsection\label{bib:sga3-VIB-Ray70a}%
  M.~Raynaud,
  \emph{Faisceaux amples sur les sch\'emas en groupes et les
  espaces homog\`enes},
  Lecture Notes in Mathematics \textbf{119},
  Springer-Verlag, 1970.
\item[{[Ray70b]}]\phantomsection\label{bib:sga3-VIB-Ray70b}%
  M.~Raynaud,
  \emph{Anneaux locaux hens\'eliens},
  Lecture Notes in Mathematics \textbf{169},
  Springer-Verlag, 1970.
\item[{[RG71]}]\phantomsection\label{bib:sga3-VIB-RG71}%
  M.~Raynaud and L.~Gruson,
  \emph{Crit\`eres de platitude et de projectivit\'e},
  Invent. Math. \textbf{13} (1971), 1--89.
\item[{[Ro56]}]\phantomsection\label{bib:sga3-VIB-Ro56}%
  M.~Rosenlicht,
  \emph{Some basic theorems on algebraic groups},
  Amer. J. Math. \textbf{78} (1956), 401--443.
\item[{[SS09]}]\phantomsection\label{bib:sga3-VIB-SS09}%
  C.~Sancho de Salas and F.~Sancho de Salas,
  \emph{Principal bundles, quasi-abelian varieties and
  structure of algebraic groups},
  J. Algebra \textbf{322} (2009), no.~8, 2751--2772.
\item[{[Se68]}]\phantomsection\label{bib:sga3-VIB-Se68}%
  J.-P.~Serre,
  \emph{Groupes de Grothendieck des sch\'emas en groupes
  r\'eductifs d\'eploy\'es},
  Publ. Math. I.H.\'E.S. \textbf{34} (1968), 37--52.
\item[{[Se99]}]\phantomsection\label{bib:sga3-VIB-Se99}%
  C.~S.~Seshadri,
  \emph{Chevalley: some reminiscences},
  Transform. Groups \textbf{4} (1999), nos.~2--3, 119--125.
\item[{[Ta72]}]\phantomsection\label{bib:sga3-VIB-Ta72}%
  M.~Takeuchi,
  \emph{A correspondence between Hopf ideals and sub-Hopf
  algebras},
  Manuscripta Math. \textbf{7} (1972), 251--270.
\item[{[Ta79]}]\phantomsection\label{bib:sga3-VIB-Ta79}%
  M.~Takeuchi,
  \emph{Relative Hopf modules---Equivalences and freeness
  criteria},
  J. Algebra \textbf{60} (1979), 452--471.
\item[{[Th87]}]\phantomsection\label{bib:sga3-VIB-Th87}%
  R.~W.~Thomason,
  \emph{Equivariant resolution, linearization, and Hilbert's
  fourteenth problem over arbitrary base schemes},
  Adv. Math. \textbf{65} (1987), 16--34.
\end{description}

% End of Expose VI_B. Component-local page 112 is a terminal blank.

% END INCLUDED SGA3_English_current_buildable_source_20260823/source/inputs/VI/tex_reference_v2/components/89_expose_VIB_section13_through_bibliography_loop2_r1_en.tex



% END INCLUDED SGA3_English_current_buildable_source_20260823/source/inputs/VI/body.tex


\SGAStartExpose{VII}
% BEGIN INCLUDED SGA3_English_current_buildable_source_20260823/source/inputs/VII/body.tex
% Mechanically extracted after \mainmatter from the frozen standalone master.

% BEGIN INCLUDED SGA3_English_current_buildable_source_20260823/source/inputs/VII/tex/components/00_expose_VIIA_title_through_par123_en.tex
% SGA 3, Expose VII_A, English reconstruction -- Loop 1.
% Scope: authority local pages 1--5, from the title through the end of 1.2.3,
% stopping immediately before 1.3 on local page 5.
% Combined-reader pages: 454--458. Re-edition printed pages: 443--447.
% Original-edition marginal locators within this span: 412--415.
% Authority: Polo--Gille Exp7A-13oct24.pdf,
% SHA-256 D46DB8C79D0385643B6EED4B96900613248375F9B73B4202478A4F57FD3E8522.
% Jacob Reinhold's CC BY 4.0 Markdown was consulted as comparison lineage.

\chapter*{Expos\'e \texorpdfstring{VII\textsubscript{A}}{VII-A}\\
Infinitesimal Study of Group Schemes}
\phantomsection
\label{sga3:VIIA:title}
\addcontentsline{toc}{chapter}{Expos\'e VII-A. Infinitesimal Study of Group Schemes}
\markboth{EXPOS\'E VII\textsubscript{A}. INFINITESIMAL STUDY OF GROUP SCHEMES}%
  {EXPOS\'E VII\textsubscript{A}. INFINITESIMAL STUDY OF GROUP SCHEMES}

\begin{center}
  \textit{by P. Gabriel}
\end{center}

% Authority: Exp7A local p. 1; combined-reader p. 454; printed p. 443.
In Expos\'e II we restricted ourselves to the study of first-order
differential invariants, and we did not address certain phenomena special
to characteristic \(p>0\) or to characteristic zero. Our purpose in
Part~A of this expos\'e is to fill that gap.

Moreover, the infinitesimal study of a group scheme to arbitrary order is
related to that of the associated formal group; the purpose of the second
part of this expos\'e is to present the first definitions and properties
concerning formal groups.

\renewcommand{\thefootnote}{\fnsymbol{footnote}}
\setcounter{footnote}{0}
\section*{A) Differential Operators and Restricted
  \texorpdfstring{\(p\)}{p}-Lie Algebras\footnote{Part~A of the present
  expos\'e was not treated seriously in the oral seminars.}}
\phantomsection
\label{sga3:VIIA:A}
\addcontentsline{toc}{section}{A) Differential Operators and Restricted p-Lie Algebras}
\renewcommand{\thefootnote}{\arabic{footnote}}
\setcounter{footnote}{0}

\section*{1. Differential Operators}
\phantomsection
\label{sga3:VIIA:1}
\addcontentsline{toc}{section}{1. Differential Operators}

In this section, as well as in \SGARefLink{sga3:VIIA:2}{Sections~2} and~3, \(S\) denotes a fixed
scheme, and all products under consideration are cartesian products in
the category of \(S\)-schemes.\footnote{In particular, if \(X\) and \(Y\)
are two \(S\)-schemes, \(X\times_S Y\) is denoted simply by \(X\times Y\).
For the material in \SGARefLink{sga3:VIIA:1}{Sections~1} and~2, one may also consult \emph{Groupes
alg\'ebriques} by Demazure and Gabriel, II.4, nos.~5--6; see also
Jantzen, \emph{Representations of Algebraic Groups}, I.7.}
If \(X\) is an \(S\)-scheme, we write \(p_{X/S}\), \(p_X\), or simply
\(p\) for the structural morphism from \(X\) to \(S\).

\medskip
\noindent\textbf{1.1.}
\phantomsection
\label{sga3:VIIA:1.1}
Let \(u:Y\to X\) be a morphism of \(S\)-schemes, and endow the direct
image \(u_*(\mathcal O_Y)\) of the structure sheaf of \(Y\) with the
\(\mathcal O_X\)-module structure induced by \(u\). The sheaf
\[
  \mathcal H=
  \mathcal Hom_{p_X^{-1}(\mathcal O_S)}
    \bigl(\mathcal O_X,u_*(\mathcal O_Y)\bigr)
\]
of homomorphisms of \(p_X^{-1}(\mathcal O_S)\)-modules from
\(\mathcal O_X\) to \(u_*(\mathcal O_Y)\) is therefore naturally an
\(\mathcal O_X\)-bimodule. If \(U\) is an open subset of \(X\), and
\(f\) and \(d\) are sections over \(U\) of \(\mathcal O_X\) and
\(\mathcal H\), respectively, then \(fd\) and \(df\) are the morphisms
\[
  g\longmapsto f\,d(g)
  \qquad\text{and}\qquad
  g\longmapsto d(fg)
\]
from \(\mathcal O_X\) to \(u_*(\mathcal O_Y)\). From now on we write
\((\operatorname{ad}f)(d)\) in place of \(fd-df\).

% Authority: Exp7A local p. 2; combined-reader p. 455; printed p. 444.
\medskip
\noindent\textbf{Definition 1.1.1.}
\phantomsection
\label{sga3:VIIA:1.1.1}
\textit{An \(S\)-deviation of order \(\leq n\) is a pair
\(D=(u,d)\), consisting of a morphism of \(S\)-schemes \(u:Y\to X\)
and a morphism of \(p_X^{-1}(\mathcal O_S)\)-modules
\[
  d:\mathcal O_X\longrightarrow u_*(\mathcal O_Y),
\]
such that, for every open subset \(U\) of \(X\) and every sequence of
\(n+1\) sections \(f_0,\ldots,f_n\in\mathcal O_X(U)\), one has}
\begin{equation}
  (\operatorname{ad}f_0)(\operatorname{ad}f_1)\cdots
  (\operatorname{ad}f_n)(d)=0.
  \tag{\(\ast_n\)}
  \label{sga3:VIIA:eq:star-n}
\end{equation}
\footnote{Equivalently, for every \(x\in X\) and
\(f_0,\ldots,f_n,g\in\mathcal O_{X,x}\),
\[
  (\operatorname{ad}f_0)\cdots(\operatorname{ad}f_n)(d_x)(g)=0.
\]
Recall also that the adjunction isomorphism
\[
 \theta:
 \operatorname{Hom}_{p_X^{-1}(\mathcal O_S)}
   \bigl(\mathcal O_X,u_*(\mathcal O_Y)\bigr)
 \xrightarrow{\sim}
 \operatorname{Hom}_{p_Y^{-1}(\mathcal O_S)}
   \bigl(u^{-1}\mathcal O_X,\mathcal O_Y\bigr)
\]
sends \(d\) to \(d'=\varepsilon\circ u^{-1}(d)\), where
\(\varepsilon:u^{-1}u_*(\mathcal O_Y)\to\mathcal O_Y\) is canonical.
Conversely,
\(\theta^{-1}(d')=u_*(d')\circ\eta\), with
\(\eta:\mathcal O_X\to u_*u^{-1}\mathcal O_X\) canonical. Hence \(d\)
satisfies \((\ast_n)\) if and only if \(d'\) satisfies the analogous
identity for all local sections of \(u^{-1}\mathcal O_X\).}

In this case we also say that \(d\) is an \(S\)-deviation of \(u\) of
order \(\leq n\). In particular, an \(S\)-deviation of \(u\) of order
\(\leq0\) is a morphism of \(\mathcal O_X\)-modules from
\(\mathcal O_X\) to \(u_*(\mathcal O_Y)\), that is, an element of
\(\Gamma(Y,\mathcal O_Y)\).

\medskip
\noindent\textbf{Definition 1.1.2.}
\phantomsection
\label{sga3:VIIA:1.1.2}
\textit{A morphism of \(p_X^{-1}(\mathcal O_S)\)-modules
\[
  d:\mathcal O_X\longrightarrow u_*(\mathcal O_Y)
\]
is an \(S\)-deviation of \(u\) if, for every point \(y\in Y\), there
exist an open neighborhood \(U\) of \(u(y)\) in \(X\) and an open
neighborhood \(V\) of \(y\) in \(Y\) satisfying:
\begin{enumerate}
\item[\textup{a)}] \(u(V)\subset U\);
\item[\textup{b)}] if \(v:V\to U\) is induced by \(u\), there is an
integer \(n\) such that the morphism
\(\mathcal O_U\to v_*(\mathcal O_V)\) induced by \(d\) is an
\(S\)-deviation of \(v\) of order \(\leq n\).
\end{enumerate}}\footnote{If \(X\) and \(u\) are quasi-compact, every
\(S\)-deviation of \(u\) is therefore of order \(\leq n\) for some
integer \(n\).}

If \(d\) is an \(S\)-deviation of \(u\), we also call the pair
\(D=(u,d)\) an \(S\)-deviation, and we may write
\[
  Y\xrightarrow{D}X
  \qquad\text{or}\qquad
  Y\mathop{\longrightarrow}^{d}_{u}X.
\]
When \(d\) is the homomorphism of algebras
\[
  u^\natural:\mathcal O_X\longrightarrow u_*(\mathcal O_Y)
\]
corresponding to the morphism \(u:Y\to X\), we also write \(u\) in
place of \(D\).

\medskip
\noindent\textbf{Remarks 1.1.3.}\footnote{These remarks were added by
the editors because they are useful in \SGARefLink{sga3:VIIA:1.3}{1.3}, \SGARefLink{sga3:VIIA:1.4}{1.4}, and \SGARefLink{sga3:VIIA:2.1}{2.1}.}
\phantomsection
\label{sga3:VIIA:1.1.3}
Let \(\operatorname{Dev}(u)\) (respectively,
\(\operatorname{Dev}_{\leq n}(u)\)) be the set of \(S\)-deviations of
\(u\) (respectively, of order \(\leq n\)). It has a natural
\(\mathcal O_Y(Y)\)-module structure: if
\(\lambda\in\mathcal O_Y(Y)\), then \(\lambda d\) is the deviation
that sends \(f\) to \(\lambda d(f)\) for every local section \(f\) of
\(\mathcal O_X\).

For every open subset \(V\) of \(Y\), set
\(\mathcal{Dev}(u)(V)=\operatorname{Dev}(u|_V)\). Thus
\(\mathcal{Dev}(u)(V)\) is the set of
\begin{align*}
 d_V&\in
 \operatorname{Hom}_{p_X^{-1}(\mathcal O_S)}
   \bigl(\mathcal O_X,(u|_V)_*(\mathcal O_V)\bigr)\\
 &\simeq
 \operatorname{Hom}_{p_Y^{-1}(\mathcal O_S)}
   \bigl((u|_V)^{-1}\mathcal O_X,\mathcal O_V\bigr)\\
 &\simeq
 \operatorname{Hom}_{p_Y^{-1}(\mathcal O_S)}
   \bigl(u^{-1}\mathcal O_X,\mathcal O_Y\bigr)(V)
\end{align*}
such that, for every open \(U\subset X\), the map
\[
 d_V(U):\mathcal O_X(U)\longrightarrow
 \mathcal O_Y\bigl(u^{-1}(U)\cap V\bigr)
\]
satisfies \((\ast_n)\). This defines a presheaf of
\(\mathcal O_Y\)-modules on \(Y\), and it is easily seen to be a sheaf
(more precisely, a subsheaf of
\(\mathcal Hom_{p_Y^{-1}(\mathcal O_S)}
  (u^{-1}\mathcal O_X,\mathcal O_Y)\)).

% Authority: Exp7A local p. 3; combined-reader p. 456; printed p. 445.
\medskip
\noindent\textbf{1.2.}
\phantomsection
\label{sga3:VIIA:1.2}
Consider two \(S\)-deviations \(D=(u,d)\) and \(E=(v,e)\):
\[
  Z\mathop{\longrightarrow}^{e}_{v}Y
  \mathop{\longrightarrow}^{d}_{u}X.
\]
As \(U\) ranges over the open subsets of \(X\), the composites
\[
 \Gamma(U,\mathcal O_X)
 \xrightarrow{d(U)}
 \Gamma(u^{-1}U,\mathcal O_Y)
 \xrightarrow{e(u^{-1}U)}
 \Gamma(v^{-1}u^{-1}U,\mathcal O_Z)
\]
define an \(S\)-deviation of \(uv\), denoted by \(de\). If \(d\) has
order \(\leq m\) and \(e\) has order \(\leq n\), then \(de\) has order
\(\leq m+n\). We also write
\begin{equation}
  D\circ E=(uv,de)
  \tag{\(\dagger\)}
  \label{sga3:VIIA:eq:composition}
\end{equation}
\footnote{With this notation, \(de\) means the composite ``\(d\)
followed by \(e\).''}
and call \(D\circ E\), or \(DE\), the \emph{composite
\(S\)-deviation}. When \(d=u^\natural\), so that \(D=u\) under the
convention of \SGARefLink{sga3:VIIA:1.1}{1.1}, we also say that \(DE\) is the \emph{image of \(E\)
under \(u\)}.

The map \((D,E)\mapsto D\circ E\) allows us to speak from now on of the
\emph{category of \(S\)-deviations}, whose objects are \(S\)-schemes
and whose morphisms are \(S\)-deviations.\footnote{Often one considers
only the \(S\)-deviations of \(\operatorname{id}_X\); they form the
algebra of \(S\)-differential operators on \(X\), cf.~\SGARefLink{sga3:VIIA:1.4}{1.4} below. The
more general framework provides a convenient functorial language for
statements such as the following: if \(G\) is an \(S\)-group, the
algebra of left-translation-invariant \(S\)-differential operators on
\(G\) is isomorphic to the algebra of \(S\)-deviations of the unit
section \(\varepsilon:S\to G\); cf.~\SGARefLink{sga3:VIIA:2.1}{2.1} and~\SGARefLink{sga3:VIIA:2.4}{2.4}.}

\medskip
\noindent\textbf{Definition 1.2.0.}\footnote{The editors expanded this
paragraph and assigned the number \SGARefLink{sga3:VIIA:1.2.0}{1.2.0} to this definition and 1.2.1
to the lemma that follows.}
\phantomsection
\label{sga3:VIIA:1.2.0}
\textit{Let \(w:Z\to X\) be an \(S\)-morphism. An
\(S\)-derivation of \(w\), or an \(S\)-derivation of \(\mathcal O_X\)
into \(w_*(\mathcal O_Z)\), is a morphism of
\(p^{-1}(\mathcal O_S)\)-modules
\[
  d:\mathcal O_X\longrightarrow w_*(\mathcal O_Z)
\]
such that, for every open \(U\subset X\) and
\(f,g\in\mathcal O_X(U)\),
\[
  d(fg)=w^\natural(f)d(g)+w^\natural(g)d(f).
\]
Then \(d\) is a deviation of \(w\) of order \(\leq1\), vanishing on
the unit section of \(\mathcal O_X\). We denote the set of
\(S\)-derivations of \(w\) by \(\operatorname{Der}_S(w)\); it is an
\(\mathcal O(Z)\)-module.}

With the notation of \SGARefLink{sga3:VIIA:1.2}{1.2}, take
\[
  Y=I_Z=\operatorname{Spec}\mathcal O_Z[t],\qquad t^2=0,
\]
let \(v\) be the zero section \(\tau:Z\to I_Z\), defined by the
homomorphism \(\mathcal O_Z[t]\to\mathcal O_Z\) sending \(t\) to zero,
and let \(e\) be the morphism of \(\mathcal O_Z\)-modules
\(\sigma:\mathcal O_Z[t]\to\mathcal O_Z\) defined by
\(\sigma(1)=0\) and \(\sigma(t)=1\), which it is convenient to denote
by \(\partial_t\).\footnote{The editors added this notation
\(\partial_t\).}

If \(u:I_Z\to X\) is a morphism satisfying \(w=u\circ\tau\), then
\(\sigma\circ u^\natural\) is an \(S\)-derivation from
\(\mathcal O_X\) to \(w_*(\mathcal O_Z)\). Conversely, an
\(S\)-derivation \(d\) determines the morphism \(u:I_Z\to X\) which
equals \(w\) on the underlying spaces and satisfies
\[
  u^\natural(f)=w^\natural(f)+d(f)t
\]
for every local section \(f\) of \(\mathcal O_X\).

% Authority: Exp7A local p. 4; combined-reader p. 457; printed p. 446.
\medskip
\noindent\textbf{Lemma 1.2.1.}
\phantomsection
\label{sga3:VIIA:1.2.1}
\textit{Let \(E=(\tau,\partial_t)\) be the deviation of
\(\tau:Z\to I_Z\) defined above. For every \(S\)-morphism \(w:Z\to X\),
the map \(u\mapsto u\circ E\) is a bijection between the
\(S\)-morphisms \(u:I_Z\to X\) such that \(u\circ\tau=w\), and the
\(S\)-derivations of \(w\).}

\medskip
\noindent\textbf{1.2.2.}
\phantomsection
\label{sga3:VIIA:1.2.2}
Let \(d\) be an \(S\)-deviation of \(u:Y\to X\). On the one hand, \(d\)
is plainly an \(S'\)-deviation of \(u\) for every morphism \(s:S\to S'\).

On the other hand, let \(t:T\to S\) be a morphism with target \(S\),
let \(u_T:Y_T\to X_T\) be obtained from \(u\) by base change, and let
\(t_Y:Y_T\to Y\) and \(t_X:X_T\to X\) be the canonical projections.
There exists a unique \(T\)-deviation of \(u_T\), denoted \(d_T\) or
\(d\times T\), which satisfies
\[
  t_Xd_T=dt_Y
\]
in the sense of \eqref{sga3:VIIA:eq:composition}. Thus, for every open
\(U\subset X\), the following diagram is commutative:
\begin{center}
  \phantomsection
  \label{sga3:VIIA:1.2.2-base-change-square}
  % SGA3-VII-DIAG-001
  % BEGIN INCLUDED SGA3_English_current_buildable_source_20260823/source/inputs/VII/native_diagrams/exp7/viia_00_04/diag_001_base_change_square.tex
\begin{tikzcd}[column sep=large,row sep=large]
\mathcal O(U)
  \arrow[r,"t_X^{\natural}"]
  \arrow[d,"d(U)"']
&
\mathcal O(U\times T)
  \arrow[d,"d_T(U\times T)"]
\\
\mathcal O(u^{-1}U)
  \arrow[r,"t_Y^{\natural}"]
&
\mathcal O(u^{-1}U\times T)\,.
\end{tikzcd}

% END INCLUDED SGA3_English_current_buildable_source_20260823/source/inputs/VII/native_diagrams/exp7/viia_00_04/diag_001_base_change_square.tex

\end{center}
\footnote{Explicitly, if \(V\) is an affine open subset of \(S\), and
\(U\) (respectively, \(U'\)) is an affine open subset of \(X\)
(respectively, \(T\)) over \(V\), then
\[
 \mathcal O_{X\times T}(U\times U')
 =\mathcal O_X(U)\otimes_{\mathcal O_S(V)}\mathcal O_T(U'),
\]
and \(d_T(U\times U')\) is the composite
\[
 \mathcal O_X(U)\otimes_{\mathcal O_S(V)}\mathcal O_T(U')
 \xrightarrow{d(U)\otimes\mathrm{id}}
 \mathcal O_Y(u^{-1}U)\otimes_{\mathcal O_S(V)}\mathcal O_T(U')
 \longrightarrow
 \mathcal O_{Y\times T}(u^{-1}U\times U').
\]
The author left to the reader the verification that \(d_T\) is
well-defined, and the editors do the same.}
If \(D=(u,d)\), we also write \(D_T=(u_T,d_T)\), and say that \(d_T\)
and \(D_T\) are obtained from \(d\) and \(D\) by base change.

\medskip
\noindent\textbf{1.2.3.}
\phantomsection
\label{sga3:VIIA:1.2.3}
For example, let \(u:Y\to X\) and \(v:Z\to T\) be two \(S\)-morphisms,
and let \(d\) and \(e\) be \(S\)-deviations of \(u\) and \(v\).
There is a commutative diagram
\begin{center}
  \phantomsection
  \label{sga3:VIIA:1.2.3-product-scheme-square}
  % SGA3-VII-DIAG-002
  % BEGIN INCLUDED SGA3_English_current_buildable_source_20260823/source/inputs/VII/native_diagrams/exp7/viia_00_04/diag_002_product_scheme_square.tex
\begin{tikzcd}[column sep=large,row sep=large]
X\times T
&
Y\times T
  \arrow[l,"u_T"']
\\
X\times Z
  \arrow[u,"v_X"]
&
Y\times Z
  \arrow[l,"u_Z"']
  \arrow[u,"v_Y"']
  \arrow[ul,"u\times v"']
\end{tikzcd}

% END INCLUDED SGA3_English_current_buildable_source_20260823/source/inputs/VII/native_diagrams/exp7/viia_00_04/diag_002_product_scheme_square.tex

\end{center}
and we denote by \(d\times e\), the \emph{product of \(d\) and \(e\)},
the \(S\)-deviation of \(u\times v\) equal to
\[
  d_Te_Y=e_Xd_Z
\]
under the convention of \eqref{sga3:VIIA:eq:composition}. More
explicitly, for every open \(U\subset X\times T\), set
\[
  W=v_Y^{-1}u_T^{-1}U=u_Z^{-1}v_X^{-1}U.
\]
Then the following diagram is commutative:
\begin{center}
  \phantomsection
  \label{sga3:VIIA:1.2.3-product-deviation-square}
  % SGA3-VII-DIAG-003
  % BEGIN INCLUDED SGA3_English_current_buildable_source_20260823/source/inputs/VII/native_diagrams/exp7/viia_00_04/diag_003_product_deviation_square.tex
\begin{tikzcd}[column sep=huge,row sep=large]
\mathcal O(U)
  \arrow[r,"d_T(U)"]
  \arrow[d,"e_X(U)"']
  \arrow[dr,dashed,"(d\times e)(U)" description]
&
\mathcal O(u_T^{-1}U)
  \arrow[d,"e_Y(u_T^{-1}U)"]
\\
\mathcal O(v_X^{-1}U)
  \arrow[r,"d_Z(v_X^{-1}U)"]
&
\mathcal O(W)\,.
\end{tikzcd}

% END INCLUDED SGA3_English_current_buildable_source_20260823/source/inputs/VII/native_diagrams/exp7/viia_00_04/diag_003_product_deviation_square.tex

\end{center}
If \(D=(u,d)\) and \(E=(v,e)\), we also write
\[
  D\times E=(u\times v,d\times e).
\]
\footnote{Source note: the French sentence at this point prints
\(E=(v,d)\); the established notation is \(E=(v,e)\).}

% Next cursor: Exp7A local p. 5, combined-reader p. 458,
% re-edition printed p. 447, at the heading 1.3.

% END INCLUDED SGA3_English_current_buildable_source_20260823/source/inputs/VII/tex/components/00_expose_VIIA_title_through_par123_en.tex

% BEGIN INCLUDED SGA3_English_current_buildable_source_20260823/source/inputs/VII/tex/components/01_expose_VIIA_par13_through_par15_en.tex
% SGA 3, Expose VII_A, English reconstruction -- Loop 1.
% Scope: authority local pages 5--8, from 1.3 through the end of 1.5,
% stopping before the Section 2 heading on local page 9.
% Combined-reader pages: 458--461. Re-edition printed pages: 447--450.
% Original-edition marginal locators within this span: 415--418.
% Authority: Polo--Gille Exp7A-13oct24.pdf,
% SHA-256 D46DB8C79D0385643B6EED4B96900613248375F9B73B4202478A4F57FD3E8522.
% Jacob Reinhold's CC BY 4.0 Markdown was consulted as comparison lineage.

% Authority resumes on Exp7A local p. 5; combined p. 458; printed p. 447.
\medskip
\noindent\textbf{1.3.}\footnote{The editors expanded the original in this
paragraph; see also their \SGARefLink{sga3:VIIA:nde:2}{note~2} to \SGARefLink{sga3:VIIA:1.1.1}{Definition~1.1.1}.}
\phantomsection
\label{sga3:VIIA:1.3}
Let \(u:Y\to X\) be a morphism of \(S\)-schemes. Recall that the
adjunction isomorphism
\[
 \operatorname{Hom}_{p_X^{-1}(\mathcal O_S)}
   \bigl(\mathcal O_X,u_*(\mathcal O_Y)\bigr)
 \xrightarrow{\sim}
 \operatorname{Hom}_{p_Y^{-1}(\mathcal O_S)}
   \bigl(u^{-1}\mathcal O_X,\mathcal O_Y\bigr)
\]
sends a morphism of \(p^{-1}(\mathcal O_S)\)-modules
\[
  d:\mathcal O_X\longrightarrow u_*(\mathcal O_Y)
\]
to \(d'=\varepsilon\circ u^{-1}(d)\), where
\(\varepsilon:u^{-1}u_*(\mathcal O_Y)\to\mathcal O_Y\) is the
canonical morphism.

Write \(\mathcal J_u\) (respectively, \(\mathcal I_u\)) for the kernel
of the homomorphism of algebras
\[
 u^\natural:\mathcal O_X\longrightarrow u_*(\mathcal O_Y)
 \quad
 \left(\text{respectively, }
 u^{\natural\prime}:u^{-1}\mathcal O_X\longrightarrow\mathcal O_Y
 \right),
\]
and let \(d:\mathcal O_X\to u_*(\mathcal O_Y)\) be a morphism of
\(p^{-1}(\mathcal O_S)\)-modules. If \(U\subset X\) is open and
\(f_0,\ldots,f_n,g\in\mathcal O_X(U)\), induction on \(n\) shows that
condition \((\ast_n)\) is equivalent to
\begin{equation}
  0=
  \sum_{I\subset\{0,\ldots,n\}}
  (-1)^{|I|}
  u^\natural\bigl(f_{\{0,\ldots,n\}\setminus I}\bigr)d(f_Ig),
  \tag{\(\ast\ast_n\)}
  \label{sga3:VIIA:eq:starstar-n}
\end{equation}
where \(f_I\) denotes the product of the \(f_i\) for \(i\in I\).
Consequently, if \(d\) satisfies \((\ast_n)\), then \(d\) vanishes on
the ideal \(\mathcal J_u^{\,n+1}\).

Now suppose that \(Y=S\). Then \(u:S\to X\) is a section of
\(p:X\to S\), hence an immersion. On the one hand,
\[
 \varepsilon:u^{-1}u_*\mathcal O_S\longrightarrow\mathcal O_S
\]
is an isomorphism, and therefore \(u^{-1}\mathcal J_u=\mathcal I_u\).
On the other hand, there is an isomorphism
\begin{equation}
  u^{-1}\mathcal O_X\simeq\mathcal O_S\oplus\mathcal I_u.
  \tag{\(\star\)}
  \label{sga3:VIIA:eq:section-splitting}
\end{equation}
Suppose that \(d\) vanishes on \(\mathcal J_u^{\,n+1}\). Then
\(d'=\varepsilon\circ u^{-1}(d)\) vanishes on
\(\mathcal I_u^{\,n+1}\), and therefore satisfies the analogues
\((\ast\ast'_n)\) and \((\ast'_n)\) of
\((\ast\ast_n)\) and \((\ast_n)\) when
\[
 f_0,\ldots,f_n\in\mathcal I_u\bigl(u^{-1}(U)\bigr).
\]
Moreover,
\[
  (\operatorname{ad}a)(\varphi)=0
\]
for every \(a\in\mathcal O_S(u^{-1}U)\) and every morphism of
\(\mathcal O_{u^{-1}U}\)-modules
\(\varphi:u^{-1}\mathcal O_U\to\mathcal O_{u^{-1}U}\).
The splitting \eqref{sga3:VIIA:eq:section-splitting} therefore implies
that \(d'\) satisfies \((\ast'_n)\), and hence that \(d\) satisfies
\((\ast_n)\).

\medskip
\noindent\textbf{Lemma.}
\phantomsection
\label{sga3:VIIA:1.3-section-lemma}
\textit{If \(u:S\to X\) is a section of \(p:X\to S\), then \(d\) is an
\(S\)-deviation of \(u\) of order \(\leq n\) if and only if \(d'\)
vanishes on \(\mathcal I_u^{\,n+1}\).}

This interpretation generalizes as follows. Let \(u:Y\to X\) be any
\(S\)-morphism, and let \(\Gamma_u:Y\to Y\times X\) be its graph, with
components \(\operatorname{id}_Y\) and \(u\). Every \(S\)-deviation
\(d\) of \(u\) of order \(\leq n\) gives, by composition,
\[
  Y\xrightarrow{\mathrm{diag}}Y\times Y
   \mathop{\longrightarrow}^{d_Y}_{u_Y}Y\times X,
\]
a \(Y\)-deviation of \(\Gamma_u\) of order \(\leq n\), denoted
\(\Gamma_d\), the graph of \(d\).

Conversely, a \(Y\)-deviation \(e\) of \(\Gamma_u\) gives the composite
\(S\)-deviation
\[
 e_X=\operatorname{pr}_2\circ e:
 \qquad
 Y\mathop{\longrightarrow}^{e}_{\Gamma_u}Y\times X
  \xrightarrow{\operatorname{pr}_2}X.
\]
One has immediately \((\Gamma_d)_X=d\). The identity
\(\Gamma_{e_X}=e\) follows from the \(\mathcal O_Y\)-linearity of
\(e\).\footnote{If \(\lambda\) and \(f\) are local sections of
\(\mathcal O_Y\) and \(\mathcal O_X\), then
\[
 (\Gamma_{e_X})(\lambda\otimes f)
 =\lambda\,e(1\otimes f)
 =e(\lambda\otimes f).
\]
The French note prints \(g\) in the two right-hand tensor expressions;
the variables just introduced require \(f\).}
Thus there is an isomorphism of \(\mathcal O_Y(Y)\)-modules
\[
\left\{
\begin{array}{c}
S\text{-deviations of }u\\[-2pt]
\text{of order }\leq n
\end{array}
\right\}
\xrightarrow{\sim}
\left\{
\begin{array}{c}
Y\text{-deviations of }\Gamma_u\\[-2pt]
\text{of order }\leq n
\end{array}
\right\},
\qquad d\longmapsto\Gamma_d.
\]
Moreover, \(d\) is an \(S\)-derivation of \(u\) if and only if
\(\Gamma_d\) is a \(Y\)-derivation of \(\Gamma_u\).

% Authority: Exp7A local p. 6; combined p. 459; printed p. 448.
Let \(\mathcal I_{\Gamma_u}\) denote the kernel of the homomorphism of
algebras
\[
  \Gamma_u^{-1}\mathcal O_{Y\times X}\longrightarrow\mathcal O_Y
\]
corresponding to \(\Gamma_u\). The preceding lemma yields:

\medskip
\noindent\textbf{Proposition.}
\phantomsection
\label{sga3:VIIA:1.3-graph-proposition}
\textit{Let \(u:Y\to X\) be an \(S\)-morphism and
\(\Gamma_u:Y\to Y\times X\) its graph. The \(S\)-deviations of \(u\)
of order \(\leq n\) identify with the \(Y\)-deviations of
\(\Gamma_u\) of order \(\leq n\), and these are in bijection with}
\[
 \operatorname{Hom}_{\mathcal O_Y}
 \left(
  \frac{\Gamma_u^{-1}\mathcal O_{Y\times X}}
       {\mathcal I_{\Gamma_u}^{\,n+1}},
  \mathcal O_Y
 \right).
\]

\medskip
\noindent\textbf{1.3.1.}\footnote{This paragraph was added by the
editors.}
\phantomsection
\label{sga3:VIIA:1.3.1}
Return to the case in which \(u:S\to X\) is a section of \(p:X\to S\).
The homomorphism
\[
 \varphi:u^{-1}\mathcal O_X\longrightarrow\mathcal O_S
\]
has a section, denoted simply by \(g\mapsto g\cdot1\). Thus, with the
notation of \SGARefLink{sga3:VIIA:1.3}{1.3},
\[
  u^{-1}\mathcal O_X\simeq\mathcal O_S\oplus\mathcal I_u,
\]
and, for every section \(f\) of \(u^{-1}\mathcal O_X\),
\[
  f-\varphi(f)\cdot1\in\mathcal I_u.
\]

Let \(d\) be an \(S\)-deviation of \(u\) of order \(\leq1\), and let
\(d':u^{-1}\mathcal O_X\to\mathcal O_S\) be the
\(\mathcal O_S\)-morphism corresponding to \(d\). If \(a,b\) are
sections of \(u^{-1}\mathcal O_X\), then
\begin{align*}
0
 &=d'\bigl((a-\varphi(a)\cdot1)(b-\varphi(b)\cdot1)\bigr)\\
 &=d'(ab)-\varphi(a)d'(b)-\varphi(b)d'(a)+\varphi(ab)d'(1).
\end{align*}
Consequently, \(d\) is an \(S\)-derivation of \(u\) if and only if
\(d'(1)=0\).

% Authority: Exp7A local p. 7; combined p. 460; printed p. 449.
\medskip
\noindent\textbf{Lemma.}
\phantomsection
\label{sga3:VIIA:1.3.1-derivation-lemma}
\textit{The \(S\)-derivations of \(u\) are precisely the
\(S\)-deviations of \(u\) of order \(1\) which vanish on the unit
section of \(\mathcal O_X\). They correspond to the
\(\mathcal O_S(S)\)-module}
\[
  \operatorname{Hom}_{\mathcal O_S}
  \bigl(\mathcal I_u/\mathcal I_u^2,\mathcal O_S\bigr),
\]
\textit{and there is an isomorphism of \(\mathcal O_S(S)\)-modules}
\[
 \operatorname{Dev}_{\leq1}(u)
 \simeq
 \mathcal O_S(S)\oplus\operatorname{Der}_S(u).
\]

Returning to the general case gives:

\medskip
\noindent\textbf{Corollary.}
\phantomsection
\label{sga3:VIIA:1.3.1-graph-corollary}
\textit{Let \(u:Y\to X\) be an \(S\)-morphism and
\(\Gamma_u:Y\to Y\times X\) its graph. There is a canonical
isomorphism of \(\mathcal O_Y(Y)\)-modules}
\[
 \operatorname{Der}_S(u)
 \simeq
 \operatorname{Der}_Y(\Gamma_u)
 \simeq
 \operatorname{Hom}_{\mathcal O_Y}
 \bigl(\mathcal I_{\Gamma_u}/\mathcal I_{\Gamma_u}^2,\mathcal O_Y\bigr).
\]

\medskip
\noindent\textbf{Definition 1.4.}
\phantomsection
\label{sga3:VIIA:1.4}
\textit{Let \(X\) be an \(S\)-scheme. An \(S\)-differential operator
(respectively, an \(S\)-differential operator of order \(\leq n\)) on
\(X\) is an \(S\)-deviation (respectively, an \(S\)-deviation of order
\(\leq n\)) of the identity morphism of \(X\).}

By \SGARefLink{sga3:VIIA:1.1}{1.1}, an \(S\)-differential operator of order \(\leq n\) is therefore
an endomorphism of the \(p^{-1}(\mathcal O_S)\)-module
\(\mathcal O_X\) satisfying the identities \((\ast_n)\) of \SGARefLink{sga3:VIIA:1.1}{1.1}. We
write \(\operatorname{Dif}^n_{X/S}\) for the
\(\Gamma(\mathcal O_S)\)-module\footnote{In this expos\'e the ring
\(\Gamma(S,\mathcal O_S)=\mathcal O_S(S)\) is denoted by
\(\Gamma(\mathcal O_S)\).}
of \(S\)-differential operators of order \(\leq n\), and
\(\operatorname{Dif}_{X/S}\) for the module of all
\(S\)-differential operators.

As observed in \SGARefLink{sga3:VIIA:1.2}{1.2}, \(S\)-deviations of \(\operatorname{id}_X\) can be
composed. This makes \(\operatorname{Dif}_{X/S}\) a
\(\Gamma(\mathcal O_S)\)-algebra, called the \emph{algebra of
differential operators of \(X/S\)}.

Similarly, for every open subset \(V\subset X\), set
\[
 \mathcal{Dif}_{X/S}(V)
 =\operatorname{Dif}_{V/S}
 =\operatorname{Dev}(\operatorname{id}_V).
\]
By \SGARefLink{sga3:VIIA:1.1.3}{1.1.3} this defines a sheaf of \(\mathcal O_X\)-modules, called the
\emph{sheaf of \(S\)-differential operators on \(X\)}.\footnote{The
editors modified the original here. It mentioned the sheaf
\(U\mapsto\operatorname{Dif}_{X_U/U}\), where \(U\) ranges over open
subsets of \(S\); that sheaf is the direct image of
\(\mathcal{Dif}_{X/S}\) under \(p_X:X\to S\).}

\medskip
\noindent\textbf{1.4.1.}
\phantomsection
\label{sga3:VIIA:1.4.1}
As in \SGARefLink{sga3:VIIB:1.3.prop}{1.3}, differential operators on \(X/S\) may be interpreted using
the graph of the identity morphism of \(X\), namely the diagonal
\[
  \Delta=\Delta_{X/S}:X\longrightarrow X\times X.
\]
Endow \(\mathcal O_{X\times X}\) with the
\(\operatorname{pr}_1^{-1}(\mathcal O_X)\)-algebra structure defined
by \(\operatorname{pr}_1\). Then
\(\Delta^{-1}\mathcal O_{X\times X}\) is an algebra over
\[
  \mathcal O_X=\Delta^{-1}\operatorname{pr}_1^{-1}\mathcal O_X.
\]
Let \(\mathcal I_{X/S}\) be the kernel of the homomorphism
\[
 \Delta^{-1}\mathcal O_{X\times X}
 \xrightarrow{m}
 \mathcal O_X
\]
adjoint to
\(\mathcal O_{X\times X}\to\Delta_*(\mathcal O_X)\), and set
\[
  \mathcal P^m_{X/S}
  =
  \Delta^{-1}\mathcal O_{X\times X}/\mathcal I_{X/S}^{\,m+1}.
\]
If \(V\subset S\) and \(U\subset X\) are affine open subsets with
\(U\) above \(V\), and if
\[
  k=\Gamma(V,\mathcal O_S),
  \qquad
  A=\Gamma(U,\mathcal O_X),
\]
then
\[
  \Gamma(U,\mathcal P^m_{X/S})
  =(A\otimes_kA)/I^{m+1},
\]
where \(I\) is generated by the elements \(a\otimes1-1\otimes a\),
for \(a\in A\).

% Authority: Exp7A local p. 8; combined p. 461; printed p. 450.
By \SGARefLink{sga3:VIIB:1.3.prop}{1.3} there is an isomorphism of \(\mathcal O_X(X)\)-modules
\[
 j_X:\operatorname{Dif}^m_{X/S}
 \xrightarrow{\sim}
 \operatorname{Hom}_{\mathcal O_X}
   (\mathcal P^m_{X/S},\mathcal O_X).
\]
It is given explicitly as follows: if
\(d\in\operatorname{Dif}^m_{X/S}\) and a section \(c\) of
\(\mathcal P^m_{X/S}\) over \(U\) has the form
\[
  c=a\otimes b+I^{m+1},
\]
then \(j_X(d)(c)=a\,d(b)\).\footnote{Under this isomorphism, the
\(X\)-derivations of \(\Delta_{X/S}\) correspond, by \SGARefLink{sga3:VIIA:1.3.1}{1.3.1}, to the
\(S\)-derivations of \(\operatorname{id}_X\), that is, to the
\(p^{-1}(\mathcal O_S)\)-derivations of \(\mathcal O_X\).}

\medskip
\noindent\textbf{1.4.2.}
\phantomsection
\label{sga3:VIIA:1.4.2}
Let \(d\) be a differential operator and \(u\) a section of \(X\) over
\(S\). The \emph{value of \(d\) at \(u\)} is the composite
\(S\)-deviation
\[
  S\xrightarrow{u}X
  \mathop{\longrightarrow}^{d}_{\operatorname{id}_X}X.
\]
By \SGARefLink{sga3:VIIB:1.3.prop}{1.3} and \SGARefLink{sga3:VIIA:1.4.1}{1.4.1}, if \(d\) has order \(\leq m\), then \(du\)
(respectively, \(d\)) is canonically associated with a morphism of
\(\mathcal O_S\)-modules
\[
 d':u^{-1}\mathcal O_X/\mathcal I_u^{\,m+1}\longrightarrow\mathcal O_S
\]
(respectively, a morphism of \(\mathcal O_X\)-modules
\(d'':\mathcal P^m_{X/S}\to\mathcal O_X\)).

The morphism \(d'\) may be recovered from \(d''\) as follows. The square
\begin{center}
  \phantomsection
  \label{sga3:VIIA:1.4.2-value-cartesian-square}
  % SGA3-VII-DIAG-004
  % BEGIN INCLUDED SGA3_English_current_buildable_source_20260823/source/inputs/VII/native_diagrams/exp7/viia_00_04/diag_004_value_cartesian_square.tex
\begin{tikzcd}[column sep=huge,row sep=large]
X\simeq S\times X
  \arrow[r,"u\times X"]
  \arrow[d,"p"']
&
X\times X
  \arrow[d,"\operatorname{pr}_1"]
\\
S
  \arrow[r,"u"]
&
X
\end{tikzcd}

% END INCLUDED SGA3_English_current_buildable_source_20260823/source/inputs/VII/native_diagrams/exp7/viia_00_04/diag_004_value_cartesian_square.tex

\end{center}
is cartesian. Hence \(X\) identifies with
\(S\times_X(X\times X)\), \(u\) with \(S\times_X\Delta\), and
\(u^*(\mathcal P^m_{X/S})\) with
\(u^{-1}\mathcal O_X/\mathcal I_u^{\,m+1}\). Under these
identifications, \(u^*(d'')\) is the morphism
\[
  u^{-1}\mathcal O_X/\mathcal I_u^{\,m+1}\longrightarrow\mathcal O_S
\]
which is precisely \(d'\).

\medskip
\noindent\textbf{1.5.}
\phantomsection
\label{sga3:VIIA:1.5}
Set, as usual,
\[
  I_S=\operatorname{Spec}\mathcal O_S[T]/(T^2).
\]
Let \(\tau:S\to I_S\) be the zero section, and let \(\sigma\) be the
canonical deviation of \(\tau\) defined in \SGARefLink{sga3:VIIA:1.2.0}{1.2.0}: it is the
\(\mathcal O_S\)-module homomorphism which vanishes on the unit section
of \(\mathcal O_S[T]/(T^2)\) and sends the class \(t\) of \(T\) modulo
\(T^2\) to the unit section of \(\mathcal O_S\).

Let \(X\) be an \(S\)-scheme. Every \(I_S\)-automorphism \(u\) of
\(I_S\times X\) inducing the identity on \(X\) determines, by
composition, a differential operator \(D_u\) on \(X\):
\begin{center}
  \phantomsection
  \label{sga3:VIIA:1.5-infinitesimal-automorphism-sequence}
  % SGA3-VII-DIAG-005
  % BEGIN INCLUDED SGA3_English_current_buildable_source_20260823/source/inputs/VII/native_diagrams/exp7/viia_00_04/diag_005_infinitesimal_automorphism_sequence.tex
\begin{tikzcd}[column sep=large]
X\simeq S\times X
  \arrow[r,"\sigma\times X"]
&
I_S\times X
  \arrow[r,"u"]
&
I_S\times X
  \arrow[r,"\operatorname{pr}_2"]
&
X\,.
\end{tikzcd}

% END INCLUDED SGA3_English_current_buildable_source_20260823/source/inputs/VII/native_diagrams/exp7/viia_00_04/diag_005_infinitesimal_automorphism_sequence.tex

\end{center}
By \SGARefLink{sga3:ii:subsection:3-14}{Expos\'e II, 3.14}, the map \(u\mapsto D_u\) is an isomorphism from
the \(\Gamma(\mathcal O_S)\)-Lie algebra
\[
  \operatorname{Lie}(\operatorname{Aut}X)
  :=\operatorname{Lie}(\operatorname{Aut}X)(S)
\]
onto the \(\Gamma(\mathcal O_S)\)-Lie algebra of
\(p^{-1}(\mathcal O_S)\)-derivations of \(\mathcal O_X\). Its inverse
sends a derivation \(D\) to the automorphism of \(I_S\times X\)
corresponding to
\[
  a+bt\longmapsto a+(Da+b)t
\]
on \(\mathcal O_X[T]/(T^2)\).

% Next cursor: Exp7A local p. 9, combined-reader p. 462,
% re-edition printed p. 451, at the Section 2 heading.

% END INCLUDED SGA3_English_current_buildable_source_20260823/source/inputs/VII/tex/components/01_expose_VIIA_par13_through_par15_en.tex

% BEGIN INCLUDED SGA3_English_current_buildable_source_20260823/source/inputs/VII/tex/components/02_expose_VIIA_section2_invariant_differential_operators_en.tex
% SGA 3, Expose VII_A, English reconstruction -- Loop 1.
% Scope: authority local pages 9--14, all of Section 2,
% stopping before the Section 3 heading on local page 15.
% Combined-reader pages: 462--467. Re-edition printed pages: 451--456.
% Original-edition marginal locators within this span: 419--425.
% Authority: Polo--Gille Exp7A-13oct24.pdf,
% SHA-256 D46DB8C79D0385643B6EED4B96900613248375F9B73B4202478A4F57FD3E8522.
% Jacob Reinhold's CC BY 4.0 Markdown was consulted as comparison lineage.

\SGAInlineRefTarget{sga3:VIIA:nde:2}\section*{2. Invariant Differential Operators on Group Schemes}
\phantomsection
\label{sga3:VIIA:2}
\addcontentsline{toc}{section}{2. Invariant Differential Operators on Group Schemes}

% Authority: Exp7A local p. 9; combined p. 462; printed p. 451.
\medskip
\noindent\textbf{2.1.}
\phantomsection
\label{sga3:VIIA:2.1}
Let \(G\) be an \(S\)-group scheme; we denote its unit section by
\(\varepsilon\), or by \(\varepsilon_G:S\to G\).

\medskip
\noindent\textbf{Definition.}
\textit{Let \(U(G)\) be the \(\Gamma(\mathcal O_S)\)-module of
\(S\)-deviations of \(\varepsilon_G\) (or \(S\)-deviations of the
origin), as in \SGARefLink{sga3:VIIA:1.1}{1.1}.}

\textit{If \(d,e\in U(G)\), then \(d\times e\) is an \(S\)-deviation of
\[
  \varepsilon\times\varepsilon:
  S\simeq S\times S\longrightarrow G\times G.
\]
The image of \(d\times e\) under the multiplication morphism
\(m:G\times G\to G\), as in \SGARefLink{sga3:VIIA:1.2}{1.2}, is called the \emph{product} of
\(d\) and \(e\) and is denoted \(d\cdot e\).}

\textit{The \(\Gamma(\mathcal O_S)\)-module \(U(G)\) is thereby made
into an associative \(\Gamma(\mathcal O_S)\)-algebra with unit
\(\varepsilon_G\). We call \(U(G)\) the \emph{infinitesimal algebra}
of \(G\).}\footnote{This is now usually called the algebra of
distributions (at the origin) of \(G\); see Demazure--Gabriel,
II.4, \SGARefLink{sga3:VIIA:6.1}{6.1}, and Jantzen, I.7.7.}

As \(T\) ranges over schemes above \(S\), the infinitesimal algebra
\(U(G_T)\) of the \(T\)-group \(G\times T\) plainly varies
contravariantly with \(T\); thus one may speak of the
\emph{infinitesimal-algebra functor}.

Restricting \(T\) to the open subsets of \(S\) gives a presheaf
\(T\mapsto U(G_T)\) of \(\mathcal O_S\)-algebras. By \SGARefLink{sga3:VIIA:1.1.3}{1.1.3} it is a
sheaf. We denote it by \(\mathcal U(G)\) and call it the
\emph{sheaf of infinitesimal algebras of \(G\)}.

The algebra \(U(G)\) is also covariantly functorial in \(G\). Indeed,
if \(u:G\to H\) is a homomorphism of \(S\)-groups and \(d\) is an
\(S\)-deviation of \(\varepsilon_G\), its image under \(u\) is the
element
\[
  U(u)(d)=u d
\]
of \(U(H)\). The resulting map \(U(u):U(G)\to U(H)\) is a
homomorphism of \(\Gamma(\mathcal O_S)\)-algebras. Likewise one obtains
a homomorphism
\(\mathcal U(u):\mathcal U(G)\to\mathcal U(H)\).

\medskip
\noindent\textbf{2.2.}
\phantomsection
\label{sga3:VIIA:2.2}
Let \(d\in U(G)\), that is, let \(d\) be an \(S\)-deviation of the
origin of \(G\). The base change of \(d\) is the \(S\)-deviation
\[
  d\times G
\]
of
\(\varepsilon\times G:G\simeq S\times G\to G\times G\), as in
\SGARefLink{sga3:VIIA:1.2.2}{1.2.2}. Its image under multiplication \(m:G\times G\to G\) is an
\(S\)-deviation of
\[
  m\circ(\varepsilon\times\operatorname{id}_G)
  =\operatorname{id}_G,
\]
hence an element of \(\operatorname{Dif}_{G/S}\), which we denote
by \(d^G\).

The map \(d\mapsto d^G\) is plainly
\(\Gamma(\mathcal O_S)\)-linear. The following commutative diagram
shows that
\[
  (e\cdot d)^G=d^G\cdot e^G.
\]
\begin{center}
  \phantomsection
  \label{sga3:VIIA:2.2-right-translation-associativity}
  % SGA3-VII-DIAG-006
  % BEGIN INCLUDED SGA3_English_current_buildable_source_20260823/source/inputs/VII/native_diagrams/exp7/viia_00_04/diag_006_right_translation_associativity.tex
\begin{tikzcd}[column sep=small,row sep=large]
&&
G\times G\times G
  \arrow[rr,"m\times G"]
  \arrow[dr,"G\times m"]
&&
G\times G
  \arrow[dd,"m"]
\\
&
G\times G
  \arrow[ur,"G\times d\times G"]
  \arrow[dr,"m"]
&&
G\times G
  \arrow[dr,"m"]
\\
G
  \arrow[ur,shift left=.7ex,"e\times G"]
  \arrow[ur,shift right=.7ex,"\varepsilon\times G"']
  \arrow[rr,shift left=.55ex,"e^G"]
  \arrow[rr,shift right=.55ex,"\operatorname{id}_G"']
&&
G
  \arrow[ur,shift left=.7ex,"d\times G"]
  \arrow[ur,shift right=.7ex,"\varepsilon\times G"']
  \arrow[rr,shift left=.55ex,"d^G"]
  \arrow[rr,shift right=.55ex,"\operatorname{id}_G"']
&&
G
\end{tikzcd}

% END INCLUDED SGA3_English_current_buildable_source_20260823/source/inputs/VII/native_diagrams/exp7/viia_00_04/diag_006_right_translation_associativity.tex

\end{center}
\footnote{The editors corrected the original diagram by replacing
\(d\times G\times G\) with \(G\times d\times G\). The composite on
the left side of the upper triangle is then
\((e\times d)\times G\), and \(d\mapsto d^G\) becomes an
anti-isomorphism from \(U(G)\) to the right-invariant differential
operators (\SGARefLink{sga3:VIIA:2.3}{2.3}--\SGARefLink{sga3:VIIA:2.4}{2.4}). If instead \({}^Gd\) is defined as the image of
\(G\times d\) under \(m\), one similarly obtains an isomorphism onto
the left-invariant differential operators. \SGARefLink{sga3:VIIA:2.4}{Sections 2.4} and \SGARefLink{sga3:VIIA:2.5}{2.5} have
been corrected accordingly.}

The two lower triangles commute by the definitions of \(d^G\) and
\(e^G\). The composite \(S\)-deviation of \(e\times G\) and
\(G\times d\times G\) is \((e\times d)\times G\), by \SGARefLink{sga3:VIIA:1.2.2}{1.2.2}. Its image
under \(m\times G\) is \((e\cdot d)\times G\), and the image of the
latter under \(m\) is therefore \((e\cdot d)^G\).

Thus one obtains an antihomomorphism
\[
  U(G)\longrightarrow\operatorname{Dif}_{G/S}
\]
of \(\Gamma(\mathcal O_S)\)-algebras, called \emph{right
translation}.\footnote{It would be preferable to call this a
\emph{left action}. For example, if \(d\) is an \(S\)-derivation of
the origin, then by \SGARefLink{sga3:VIIA:1.2.1}{1.2.1} it is the composite of
\((\tau,\partial_t):S\to I_S\) with a morphism \(x:I_S\to G\)
satisfying \(x\circ\tau=\varepsilon\). The derivation \(d^G\) of
\(\mathcal O_G\) sends a local section \(\varphi\) to the section
\(g\mapsto\partial_t\varphi(xg)\). With this terminology one could say
that the left action commutes with right translations.}

If \(\mathcal{Dif}_{G/S}\) denotes the sheaf of \(S\)-differential
operators on \(G\), as in \SGARefLink{sga3:VIIA:1.4}{1.4}, and \(p:G\to S\) is the structural
morphism, one likewise defines a right translation
\[
  \mathcal U(G)\longrightarrow p_*(\mathcal{Dif}_{G/S}).
\]

\medskip
\noindent\textbf{2.3.}
\phantomsection
\label{sga3:VIIA:2.3}
We now characterize the differential operators on \(G/S\) of the
form \(d^G\). Let \(g:S\to G\) be a section of the structural
morphism and let \(g_G\) denote right translation by \(g\), namely
\[
  g_G:
  G\simeq G\times S
  \xrightarrow{G\times g}G\times G
  \xrightarrow{m}G.
\]
For every differential operator \(D\) on \(G/S\), the composite
\[
  g_G^{-1}Dg_G
\]
is again an \(S\)-deviation of \(\operatorname{id}_G\), hence an
element of \(\operatorname{Dif}_{G/S}\). We write
\[
  D^g=g_G^{-1}Dg_G.
\]
We call \(D\) \emph{right-invariant} if for every base change
\(t:T\to S\) and every section \(g:T\to G\times T\), one has
\[
  (D_T)^g=D_T.
\]

\medskip
\noindent\textbf{Lemma.}
\phantomsection
\label{sga3:VIIA:2.3-invariance-lemma}
\textit{For every differential operator \(D\) on \(G/S\), the
following assertions are equivalent, where \(m\) is the
multiplication morphism of \(G\):}
\begin{enumerate}
\item[\textit{(i)}] \textit{\(D\) is right-invariant.}
\item[\textit{(ii)}] \textit{The following two deviations of \(m\)
are equal:}
\[
  Dm=m(D\times G).
\]
\end{enumerate}

\noindent\((ii)\Rightarrow(i)\).
Since condition \SGARefLink{sga3:vii:paragraph:5-1:equation:ii}{(ii)} is stable under base change, it suffices to show
that it implies \(D^g=D\) for every section \(g:S\to G\). Let
\[
  h=G\times g:G\simeq G\times S\longrightarrow G\times G,
\]
so that \(m\circ h=g_G\). The equality \(D^g=D\) is equivalent to
\(g_G\circ D=D\circ g_G\), and this follows from the commutative
diagram
\begin{center}
  \phantomsection
  \label{sga3:VIIA:2.3-ii-implies-i}
  % SGA3-VII-DIAG-007
  % BEGIN INCLUDED SGA3_English_current_buildable_source_20260823/source/inputs/VII/native_diagrams/exp7/viia_00_04/diag_007_invariance_implication_square.tex
{\small
\begin{tikzcd}[column sep=large,row sep=large]
G
  \arrow[d,shift right=.55ex,"D"']
  \arrow[d,shift left=.55ex,"\operatorname{id}_G"]
&
G\times G
  \arrow[l,"m"']
  \arrow[d,shift right=.55ex,"D\times G"']
  \arrow[d,shift left=.55ex,"\operatorname{id}_{(G\times G)}"]
&
G\,.
  \arrow[l,"h"']
  \arrow[d,shift right=.55ex,"D"']
  \arrow[d,shift left=.55ex,"\operatorname{id}_G"]
\\
G
&
G\times G
  \arrow[l,"m"']
&
G
  \arrow[l,"h"']
\end{tikzcd}}

% END INCLUDED SGA3_English_current_buildable_source_20260823/source/inputs/VII/native_diagrams/exp7/viia_00_04/diag_007_invariance_implication_square.tex

\end{center}

\noindent\((i)\Rightarrow(ii)\).
Take \(t:T\to S\) to be the structural morphism \(p:G\to S\), and
take the section \(g:T\to G\times T\) to be the diagonal
\(\Delta:G\to G\times G\). Right translation
\[
  \Delta_{G\times G}:G\times G\longrightarrow G\times G
\]
then has components \(m\) and \(\operatorname{pr}_2\). The equality
\((D_G)^\Delta=D_G\) is equivalent to commutativity of the first
square in
\begin{center}
  \phantomsection
  \label{sga3:VIIA:2.3-i-implies-ii}
  % SGA3-VII-DIAG-008
  % BEGIN INCLUDED SGA3_English_current_buildable_source_20260823/source/inputs/VII/native_diagrams/exp7/viia_00_04/diag_008_delta_translation_square.tex
{\small
\begin{tikzcd}[column sep=large,row sep=large]
G\times G
  \arrow[r,"\Delta_{G\times G}"]
  \arrow[d,shift right=.55ex,"D_G"']
  \arrow[d,shift left=.55ex,"\operatorname{id}_{G\times G}"]
&
G\times G
  \arrow[r,"\operatorname{pr}_1"]
  \arrow[d,shift right=.55ex,"D_G"']
  \arrow[d,shift left=.55ex,"\operatorname{id}_{G\times G}"]
&
G\,.
  \arrow[d,shift right=.55ex,"D"']
  \arrow[d,shift left=.55ex,"\operatorname{id}_G"]
\\
G\times G
  \arrow[r,"\Delta_{G\times G}"]
&
G\times G
  \arrow[r,"\operatorname{pr}_1"]
&
G
\end{tikzcd}}

% END INCLUDED SGA3_English_current_buildable_source_20260823/source/inputs/VII/native_diagrams/exp7/viia_00_04/diag_008_delta_translation_square.tex

\end{center}
because \(m=\operatorname{pr}_1\circ\Delta_{G\times G}\). This proves
\SGARefLink{sga3:vii:paragraph:5-1:equation:ii}{(ii)}.

Now consider \(d\in U(G)\). Every square in the following diagram is
commutative:
\begin{center}
  \phantomsection
  \label{sga3:VIIA:2.3-origin-deviation-invariant}
  % SGA3-VII-DIAG-009
  % BEGIN INCLUDED SGA3_English_current_buildable_source_20260823/source/inputs/VII/native_diagrams/exp7/viia_00_04/diag_009_invariant_operator_rectangle.tex
{\scriptsize
\begin{tikzcd}[column sep=small,row sep=large]
G\times G
  \arrow[r,equal]
  \arrow[d,"m"']
&
S\times G\times G
  \arrow[r,shift left=.55ex,"d\times G\times G"]
  \arrow[r,shift right=.55ex,"\varepsilon\times G\times G"']
  \arrow[d,"S\times m"']
&
G\times G\times G
  \arrow[r,"m\times G"]
  \arrow[d,"G\times m"]
&
G\times G
  \arrow[d,"m"]
\\
G
  \arrow[r,equal]
&
S\times G
  \arrow[r,shift left=.55ex,"d\times G"]
  \arrow[r,shift right=.55ex,"\varepsilon\times G"']
&
G\times G
  \arrow[r,"m"]
&
G
\end{tikzcd}}

% END INCLUDED SGA3_English_current_buildable_source_20260823/source/inputs/VII/native_diagrams/exp7/viia_00_04/diag_009_invariant_operator_rectangle.tex

\end{center}
Since
\[
  m\circ(d\times G)=d^G
  \qquad\text{and}\qquad
  (m\times G)\circ(d\times G\times G)=d^G\times G,
\]
one also has
\[
  d^G\circ m=m\circ(d^G\times G).
\]
Thus, for every \(S\)-deviation \(d\) of the origin, \(d^G\) is a
right-invariant differential operator.

\medskip
\noindent\textbf{Theorem 2.4.}
\phantomsection
\label{sga3:VIIA:2.4}
\begin{enumerate}
\item[\textit{\SGAInlineRefTarget{sga3:VIIA:2.4:item:i}(i)}]
\textit{The map \(d\mapsto d^G\) is an anti-isomorphism of the
infinitesimal algebra \(U(G)\) onto the subalgebra
\(\operatorname{Dif}^{G}_{G/S}\) of \(\operatorname{Dif}_{G/S}\)
formed by the right-invariant differential operators.}
\item[\textit{\SGAInlineRefTarget{sga3:VIIA:2.4:item:ii}(ii)}]
\textit{Likewise, the map \(d\mapsto{}^Gd\) is an isomorphism from
\(U(G)\) onto the subalgebra of \(\operatorname{Dif}_{G/S}\) formed
by the left-invariant differential operators.}
\end{enumerate}
\footnote{The editors corrected ``isomorphism'' to
``anti-isomorphism'' and added assertion \SGARefLink{sga3:vii:paragraph:5-1:equation:ii}{(ii)}; see the preceding
editorial note.}

Let \(D\) be an arbitrary differential operator on \(G/S\), and let
\(D_0\) denote its \emph{value at the origin}, that is, the composite
deviation
\[
  S\xrightarrow{\varepsilon}G
   \mathop{\longrightarrow}^{D}_{\operatorname{id}_G}G.
\]
The right-invariant differential operator \((D_0)^G\) is obtained by
the composite
\begin{center}
  \phantomsection
  \label{sga3:VIIA:2.4-D0-composite}
  % SGA3-VII-DIAG-010
  % BEGIN INCLUDED SGA3_English_current_buildable_source_20260823/source/inputs/VII/native_diagrams/exp7/viia_00_04/diag_010_D0_composition_sequence.tex
{\small
\begin{tikzcd}[column sep=normal]
G\simeq S\times G
  \arrow[r,"\varepsilon\times G"]
&
G\times G
  \arrow[r,shift left=.55ex,"D\times G"]
  \arrow[r,shift right=.55ex,"\operatorname{id}_{G\times G}"']
&
G\times G
  \arrow[r,"m"]
&
G\,.
\end{tikzcd}}

% END INCLUDED SGA3_English_current_buildable_source_20260823/source/inputs/VII/native_diagrams/exp7/viia_00_04/diag_010_D0_composition_sequence.tex

\end{center}
If \(D\) is right-invariant, then \(Dm=m(D\times G)\), and hence
\[
  D=Dm(\varepsilon\times G)
   =m(D\times G)(\varepsilon\times G)
   =(D_0)^G.
\]
Thus \(d\mapsto d^G\) is surjective.

Conversely, for every \(S\)-deviation \(d\) of the origin one has the
commutative square
\begin{center}
  \phantomsection
  \label{sga3:VIIA:2.4-origin-value-square}
  % SGA3-VII-DIAG-011
  % BEGIN INCLUDED SGA3_English_current_buildable_source_20260823/source/inputs/VII/native_diagrams/exp7/viia_00_04/diag_011_origin_value_square.tex
\begin{tikzcd}[column sep=huge,row sep=large]
G\times G
&
G
  \arrow[l,"d\times G"']
\\
G\times S\simeq G
  \arrow[u,"G\times\varepsilon"]
&
S
  \arrow[l,"d"']
  \arrow[u,"\varepsilon"']
\end{tikzcd}

% END INCLUDED SGA3_English_current_buildable_source_20260823/source/inputs/VII/native_diagrams/exp7/viia_00_04/diag_011_origin_value_square.tex

\end{center}
from which
\[
  d=m(G\times\varepsilon)d
   =m(d\times G)\varepsilon
   =(d^G)_0.
\]
\textit{A fortiori}, \(d\mapsto d^G\) is injective. This proves the
theorem.

As \(S\) varies, \SGARefLink{sga3:VIIA:2.4}{Theorem~2.4} also shows that right translation
\[
  \mathcal U(G)\longrightarrow p_*(\mathcal{Dif}_{G/S})
\]
is an anti-isomorphism from \(\mathcal U(G)\) onto the sheaf of
\(\mathcal O_S\)-algebras
\[
  p_*(\mathcal{Dif}_{G/S})^G,
\]
which assigns to every open \(V\subset S\) the algebra
\(\operatorname{Dif}^{G_V}_{G_V/V}\).

\medskip
\noindent\textbf{Remark 2.4.1.}
\phantomsection
\label{sga3:VIIA:2.4.1}
Consider the commutative diagram
\begin{center}
  \phantomsection
  \label{sga3:VIIA:2.4.1-eta-cartesian-square}
  % SGA3-VII-DIAG-012
  % BEGIN INCLUDED SGA3_English_current_buildable_source_20260823/source/inputs/VII/native_diagrams/exp7/viia_00_04/diag_012_eta_cartesian_square.tex
\begin{tikzcd}[column sep=huge,row sep=large]
G
  \arrow[d,shift right=.55ex,"p"']
&
G\times G
  \arrow[l,"\eta"']
  \arrow[d,shift right=.55ex,"\operatorname{pr}_1"']
\\
S
  \arrow[u,shift right=.55ex,"\varepsilon"']
&
G,
  \arrow[l,"p"']
  \arrow[u,shift right=.55ex,"\Delta"']
\end{tikzcd}

% END INCLUDED SGA3_English_current_buildable_source_20260823/source/inputs/VII/native_diagrams/exp7/viia_00_04/diag_012_eta_cartesian_square.tex

\end{center}
where \(\eta\) denotes the morphism
\((x,y)\mapsto yx^{-1}\).\footnote{Thus \(G\) acts on itself on the
left by right translations.}
It induces morphisms
\[
  \eta':\eta^{-1}(\mathcal O_G)\longrightarrow\mathcal O_{G\times G}
\]
and
\[
  \Delta^{-1}(\eta'):
  p^{-1}\varepsilon^{-1}(\mathcal O_G)
  \longrightarrow\Delta^{-1}(\mathcal O_{G\times G}).
\]

For every integer \(n\geq1\), set
\[
  \mathfrak p^n_{G/S}
  =
  \varepsilon^{-1}(\mathcal O_G)/
  \mathcal I_\varepsilon^{\,n+1}.
\]
\footnote{The editors corrected the original here: the relevant
cartesian square is formed by \(\varepsilon,\eta,\Delta,p\), not by
\(p,p,\eta,\operatorname{pr}_1\).}
Since the square formed by \(\varepsilon,\eta,\Delta,p\) is cartesian,
\(\Delta^{-1}(\eta')\) induces an isomorphism of
\(\mathcal O_G\)-modules
\[
  p^*(\mathfrak p^n_{G/S})
  \xrightarrow{\sim}
  \mathcal P^n_{G/S}.
\]
The differential operators on \(G/S\) of order \(\leq n\) therefore
correspond bijectively to the homomorphisms of
\(\mathcal O_G\)-modules
\[
  p^*(\mathfrak p^n_{G/S})\longrightarrow\mathcal O_G,
\]
or equivalently to the homomorphisms of \(\mathcal O_S\)-modules
\[
  \mathfrak p^n_{G/S}\longrightarrow p_*(\mathcal O_G).
\]
Under this bijection the right-invariant differential operators
correspond to the composites
\begin{center}
  \phantomsection
  \label{sga3:VIIA:2.4.1-invariant-factorization}
  % SGA3-VII-DIAG-013
  % BEGIN INCLUDED SGA3_English_current_buildable_source_20260823/source/inputs/VII/native_diagrams/exp7/viia_00_04/diag_013_invariant_factorization_sequence.tex
\(
\mathfrak p^n_{G/S}
  \longrightarrow
\mathcal O_S
  \xrightarrow{\mathrm{can.}}
p_*(\mathcal O_G)\,.
\)

% END INCLUDED SGA3_English_current_buildable_source_20260823/source/inputs/VII/native_diagrams/exp7/viia_00_04/diag_013_invariant_factorization_sequence.tex

\end{center}
This recovers the isomorphism of \SGARefLink{sga3:VIIA:2.4}{Theorem~2.4}.

\medskip
\noindent\textbf{2.5.}
\phantomsection
\label{sga3:VIIA:2.5}
\footnote{The editors reordered this paragraph, beginning with the
definition of \(\alpha:\operatorname{Lie}(G)\to U(G)\), and corrected
the original as indicated in the note to \SGARefLink{sga3:VIIA:2.2}{2.2}.}
Let \(\operatorname{Lie}(G)\) be the Lie algebra of \(G\).\footnote{In
this expos\'e, for an \(S\)-group scheme \(G\) (respectively, an
\(S\)-scheme \(X\)), \(\operatorname{Lie}(G)\) (respectively,
\(\operatorname{Lie}(\operatorname{Aut}X)\)) denotes, in the notation
of Expos\'e II,
\(\operatorname{Lie}(G/S)(S)\) (respectively,
\(\operatorname{Lie}(\operatorname{Aut}_S(X)/S)(S)\)). It is a
\(\Gamma(\mathcal O_S)\)-Lie algebra by II, \SGARefLink{sga3:ii:subsection:4-11}{4.11} and \SGARefLink{sga3:ii:subsection:3-14}{3.14}.}
We shall define a homomorphism of
\(\Gamma(\mathcal O_S)\)-Lie algebras
\[
  \alpha:\operatorname{Lie}(G)\longrightarrow U(G).
\]

Let \(s:S\to I_S\) be the zero section and let \(\sigma\) be the
deviation of \(s\) defined in \SGARefLink{sga3:VIIA:1.2.0}{1.2.0}. Recall from II, \SGARefLink{sga3:ii:subsection:4-1}{4.1} that
\(\operatorname{Lie}(G)\) is the set of morphisms \(x:I_S\to G\)
such that \(x\circ s=\varepsilon_G\). The composite
\begin{center}
  \phantomsection
  \label{sga3:VIIA:2.5-sigma-x}
  % SGA3-VII-DIAG-014
  % BEGIN INCLUDED SGA3_English_current_buildable_source_20260823/source/inputs/VII/native_diagrams/exp7/viia_00_04/diag_014_sigma_x_sequence.tex
\begin{tikzcd}[column sep=large]
S
  \arrow[r,shift left=.55ex,"\sigma"]
  \arrow[r,shift right=.55ex,"s"']
&
I_S
  \arrow[r,"x"]
&
G
\end{tikzcd}

% END INCLUDED SGA3_English_current_buildable_source_20260823/source/inputs/VII/native_diagrams/exp7/viia_00_04/diag_014_sigma_x_sequence.tex

\end{center}
is an \(S\)-deviation of \(\varepsilon_G\), hence an element of
\(U(G)\); in the notation of \SGARefLink{sga3:VIIA:1.2}{1.2} it is denoted \(\sigma x\).
Moreover, by \SGARefLink{sga3:VIIA:1.2.1}{1.2.1} the map
\[
  \alpha:x\longmapsto\sigma x
\]
is an isomorphism of \(\mathcal O_S(S)\)-modules from
\(\operatorname{Lie}(G)\) onto the submodule
\(\operatorname{Der}(\varepsilon_G)\subset U(G)\) formed by the
\(S\)-derivations of \(\varepsilon_G\). We show that \(\alpha\) is a
homomorphism of Lie algebras.\footnote{See also \SGARefLink{sga3:ii:subsection:4-11}{Expos\'e II, 4.11}.}

Let
\[
  \rho':U(G)\longrightarrow\operatorname{Dif}_{G/S}
\]
be the algebra homomorphism which sends an \(S\)-deviation \(d\) of
\(\varepsilon_G\) to the left-invariant differential operator
\({}^Gd\), as in \SGARefLink{sga3:VIIA:2.2}{2.2}.

Let
\[
  \rho:G\longrightarrow\operatorname{Aut}_S(G)
\]
be the homomorphism of group functors which sends an \(S\)-morphism
\(g:T\to G\) to right translation by \(g\) on \(G_T\), namely
\begin{center}
  \phantomsection
  \label{sga3:VIIA:2.5-rho-right-translation}
  % SGA3-VII-DIAG-015
  % BEGIN INCLUDED SGA3_English_current_buildable_source_20260823/source/inputs/VII/native_diagrams/exp7/viia_00_04/diag_015_rho_translation_sequence.tex
{\small
\begin{tikzcd}[column sep=normal]
G_T\simeq T\times_TG_T
  \arrow[r,"G_T\times g"]
&
G_T\times_TG_T
  \arrow[r,"m_T"]
&
G_T\,.
\end{tikzcd}}

% END INCLUDED SGA3_English_current_buildable_source_20260823/source/inputs/VII/native_diagrams/exp7/viia_00_04/diag_015_rho_translation_sequence.tex

\end{center}
Recall also from \SGARefLink{sga3:VIIA:1.5}{1.5} and II, \SGARefLink{sga3:ii:subsection:3-14}{3.14} that
\[
 \operatorname{Lie}(\operatorname{Aut}G)
 =
 \operatorname{Lie}(\operatorname{Aut}_S(G)/S)(S)
\]
identifies with the infinitesimal automorphisms of \(G\), that is, the
automorphisms of \(I_S\times G\) which induce the identity on \(G\).

Since \(\rho\) is a monomorphism, so is
\[
 \operatorname{Lie}(\rho):
 \operatorname{Lie}(G/S)
 \longrightarrow
 \operatorname{Lie}(\operatorname{Aut}_S(G)/S)
\]
(see, for example, Expos\'e II, editorial note~50). Consequently
\[
  \operatorname{Lie}(\rho):
  \operatorname{Lie}(G)\longrightarrow
  \operatorname{Lie}(\operatorname{Aut}G)
\]
is injective.

On the other hand, by \SGARefLink{sga3:VIIA:1.5}{1.5} the map \(\beta\), which assigns to every
infinitesimal automorphism \(u\) of \(G\) the differential operator
\(D_u\),
\begin{center}
  \phantomsection
  \label{sga3:VIIA:2.5-beta-sequence}
  % SGA3-VII-DIAG-016
  % BEGIN INCLUDED SGA3_English_current_buildable_source_20260823/source/inputs/VII/native_diagrams/exp7/viia_00_04/diag_016_beta_automorphism_sequence.tex
{\small
\begin{tikzcd}[column sep=normal]
G\simeq S\times G
  \arrow[r,"\sigma\times G"]
&
I_S\times G
  \arrow[r,"u"]
&
I_S\times G
  \arrow[r,"\operatorname{pr}_2"]
&
G
\end{tikzcd}}

% END INCLUDED SGA3_English_current_buildable_source_20260823/source/inputs/VII/native_diagrams/exp7/viia_00_04/diag_016_beta_automorphism_sequence.tex

\end{center}
is an isomorphism from
\(\operatorname{Lie}(\operatorname{Aut}G)\) onto the Lie subalgebra
of \(\operatorname{Dif}_{G/S}\) formed by the
\(p^{-1}(\mathcal O_S)\)-derivations of \(\mathcal O_G\).

For every \(x\in\operatorname{Lie}(G)\), the following commutative
square determines the image of \(x\) under \(\operatorname{Lie}(\rho)\):
\begin{center}
  \phantomsection
  \label{sga3:VIIA:2.5-lie-rho-square}
  % SGA3-VII-DIAG-017
  % BEGIN INCLUDED SGA3_English_current_buildable_source_20260823/source/inputs/VII/native_diagrams/exp7/viia_00_04/diag_017_lie_rho_square.tex
\begin{tikzcd}[column sep=huge,row sep=large]
I_S\times G
  \arrow[r,"\operatorname{Lie}(\rho)(x)"]
  \arrow[d,"x\times G"']
&
I_S\times G
  \arrow[d,"\operatorname{pr}_2"]
\\
G\times G
  \arrow[r,"m"]
&
G\,.
\end{tikzcd}

% END INCLUDED SGA3_English_current_buildable_source_20260823/source/inputs/VII/native_diagrams/exp7/viia_00_04/diag_017_lie_rho_square.tex

\end{center}
It follows from this diagram that the image of
\(\operatorname{Lie}(\rho)(x)\) under \(\beta\) is the composite
deviation
\begin{center}
  \phantomsection
  \label{sga3:VIIA:2.5-composite-deviation}
  % SGA3-VII-DIAG-018
  % BEGIN INCLUDED SGA3_English_current_buildable_source_20260823/source/inputs/VII/native_diagrams/exp7/viia_00_04/diag_018_composite_deviation_sequence.tex
{\small
\begin{tikzcd}[column sep=normal]
G\simeq S\times G
  \arrow[r,shift left=.55ex,"\sigma\times G"]
  \arrow[r,shift right=.55ex,"s\times G"']
&
I_S\times G
  \arrow[r,"G\times x"]
&
G\times G
  \arrow[r,"m"]
&
G
\end{tikzcd}}

% END INCLUDED SGA3_English_current_buildable_source_20260823/source/inputs/VII/native_diagrams/exp7/viia_00_04/diag_018_composite_deviation_sequence.tex

\end{center}
which, by the editorial note to \SGARefLink{sga3:VIIA:2.2}{2.2}, is precisely
\[
  {}^G(\sigma x)=\rho'(\alpha(x)).
\]
Thus one obtains the commutative diagram
\begin{center}
  \phantomsection
  \label{sga3:VIIA:2.5-lie-algebra-square}
  % SGA3-VII-DIAG-019
  % BEGIN INCLUDED SGA3_English_current_buildable_source_20260823/source/inputs/VII/native_diagrams/exp7/viia_00_04/diag_019_lie_algebra_square.tex
\begin{tikzcd}[column sep=huge,row sep=large]
\operatorname{Lie}(G)
  \arrow[r,hook,"\operatorname{Lie}(\rho)"]
  \arrow[d,hook,"\alpha"']
&
\operatorname{Lie}(\operatorname{Aut}G)
  \arrow[d,hook,"\beta"]
\\
U(G)
  \arrow[r,hook,"\rho'"]
&
\operatorname{Dif}_{G/S}
\end{tikzcd}

% END INCLUDED SGA3_English_current_buildable_source_20260823/source/inputs/VII/native_diagrams/exp7/viia_00_04/diag_019_lie_algebra_square.tex

\end{center}
in which \(\operatorname{Lie}(\rho)\), \(\beta\), and \(\rho'\) are
homomorphisms of Lie algebras. Since \(\rho'\) is injective,
\(\alpha\) is also a homomorphism of Lie algebras. Hence:

\medskip
\noindent\textbf{Proposition.}
\phantomsection
\label{sga3:VIIA:2.5-proposition}
\textit{The map \(\alpha\) is an isomorphism of
\(\mathcal O_S(S)\)-Lie algebras from \(\operatorname{Lie}(G)\) onto
the Lie algebra of \(S\)-derivations of \(\varepsilon_G\); the latter
is itself isomorphic, via \(\operatorname{Lie}(\rho)\), to the Lie
algebra of left-invariant \(S\)-derivations of \(G\).}
\footnote{There exist Lie algebras \(\mathfrak g\) over a ring \(A\)
for which the canonical map \(\mathfrak g\to U(\mathfrak g)\) is not
injective; see Bourbaki, \emph{Lie Groups and Lie Algebras}, I.2,
Exercise~9. The result above shows, since \(\alpha\) factors through
\[
  \operatorname{Lie}(G)\longrightarrow
  U(\operatorname{Lie}(G))\longrightarrow U(G),
\]
that this cannot happen for ``algebraic'' Lie algebras, namely those
of the form \(\operatorname{Lie}(G)\) for an \(A\)-group scheme \(G\).}

% Next cursor: Exp7A local p. 15, combined-reader p. 468,
% re-edition printed p. 457, at the Section 3 heading.

% END INCLUDED SGA3_English_current_buildable_source_20260823/source/inputs/VII/tex/components/02_expose_VIIA_section2_invariant_differential_operators_en.tex

% BEGIN INCLUDED SGA3_English_current_buildable_source_20260823/source/inputs/VII/tex/components/03_expose_VIIA_section3_coalgebras_cartier_en.tex
% SGA 3, Expose VII_A, English reconstruction -- Loop 1.
% Scope: authority local pages 15--21, all of Section 3,
% stopping before the Section 4 heading on local page 22.
% Combined-reader pages: 468--474. Re-edition printed pages: 457--463.
% Original-edition marginal locators within this span: 426--433.
% Authority: Polo--Gille Exp7A-13oct24.pdf,
% SHA-256 D46DB8C79D0385643B6EED4B96900613248375F9B73B4202478A4F57FD3E8522.
% Jacob Reinhold's CC BY 4.0 Markdown was consulted as comparison lineage.

\section*{3. Coalgebras and Cartier Duality}
\phantomsection
\label{sga3:VIIA:3}
\addcontentsline{toc}{section}{3. Coalgebras and Cartier Duality}

% Authority: Exp7A local p. 15; combined p. 468; printed p. 457.
\medskip
\noindent\textbf{3.1.}
\phantomsection
\label{sga3:VIIA:3.1}
Let \(S\) be a scheme (or, more generally, a ringed space). An
\(\mathcal O_S\)-\emph{coalgebra}\footnote{The term \emph{cogebra} is
also used; see Bourbaki, \emph{Algebra}, III, \S 11.1. Notice also
that in this expos\'e, as well as in Expos\'e VII\(_B\), we work in
the category of cocommutative coalgebras (that is, coalgebras
satisfying condition~\SGARefLink{sga3:vii:paragraph:5-1:equation:i}{(i)}). This is crucial for defining the product
and the notion of a coalgebra in groups; see \SGARefLink{sga3:VIIA:3.1.0}{3.1.0} and \SGARefLink{sga3:VIIA:3.2}{3.2}.}
is a pair \((\mathcal U,\Delta_{\mathcal U})\) consisting of an
\(\mathcal O_S\)-module \(\mathcal U\) and a morphism of
\(\mathcal O_S\)-modules
\[
  \Delta_{\mathcal U}:\mathcal U\longrightarrow
  \mathcal U\otimes_{\mathcal O_S}\mathcal U,
\]
called the \emph{diagonal morphism} or \emph{comultiplication}, such
that:

\smallskip
\noindent
\SGARefLink{sga3:vii:paragraph:5-1:equation:i}{(i)} \(\sigma\circ\Delta_{\mathcal U}=\Delta_{\mathcal U}\), where
\(\sigma(a\otimes b)=b\otimes a\).

\smallskip
\noindent
\SGARefLink{sga3:vii:paragraph:5-1:equation:ii}{(ii)} The square
\begin{center}
  \phantomsection
  \label{sga3:VIIA:3.1-coassociativity-square}
  % SGA3-VII-DIAG-020
  % BEGIN INCLUDED SGA3_English_current_buildable_source_20260823/source/inputs/VII/native_diagrams/exp7/viia_00_04/diag_020_coassociativity_square.tex
\begin{tikzcd}[column sep=huge,row sep=large]
\mathcal U
  \arrow[r,"\Delta_{\mathcal U}"]
  \arrow[d,"\Delta_{\mathcal U}"']
&
\mathcal U\otimes_{\mathcal O_S}\mathcal U
  \arrow[d,"\operatorname{id}_{\mathcal U}\otimes\Delta_{\mathcal U}"]
\\
\mathcal U\otimes_{\mathcal O_S}\mathcal U
  \arrow[r,"\Delta_{\mathcal U}\otimes\operatorname{id}_{\mathcal U}"]
&
\mathcal U\otimes_{\mathcal O_S}\mathcal U\otimes_{\mathcal O_S}\mathcal U
\end{tikzcd}

% END INCLUDED SGA3_English_current_buildable_source_20260823/source/inputs/VII/native_diagrams/exp7/viia_00_04/diag_020_coassociativity_square.tex

\end{center}
is commutative.

\smallskip
\noindent
(iii) There exists a morphism of \(\mathcal O_S\)-modules
\[
  \varepsilon_{\mathcal U}:\mathcal U\longrightarrow\mathcal O_S,
\]
called the \emph{augmentation}, such that the two composites
\begin{center}
  \phantomsection
  \label{sga3:VIIA:3.1-augmentation-sequences}
  % SGA3-VII-DIAG-021
  % BEGIN INCLUDED SGA3_English_current_buildable_source_20260823/source/inputs/VII/native_diagrams/exp7/viia_00_04/diag_021_augmentation_sequences.tex
\(
\begin{gathered}
\mathcal U
  \xrightarrow{\Delta_{\mathcal U}}
\mathcal U\otimes_{\mathcal O_S}\mathcal U
  \xrightarrow{\operatorname{id}_{\mathcal U}\otimes\varepsilon_{\mathcal U}}
\mathcal U\otimes_{\mathcal O_S}\mathcal O_S
\simeq\mathcal U
\\[0.6em]
\mathcal U
  \xrightarrow{\Delta_{\mathcal U}}
\mathcal U\otimes_{\mathcal O_S}\mathcal U
  \xrightarrow{\varepsilon_{\mathcal U}\otimes\operatorname{id}_{\mathcal U}}
\mathcal O_S\otimes_{\mathcal O_S}\mathcal U
\simeq\mathcal U
\end{gathered}
\)

% END INCLUDED SGA3_English_current_buildable_source_20260823/source/inputs/VII/native_diagrams/exp7/viia_00_04/diag_021_augmentation_sequences.tex

\end{center}
are the identity morphism of \(\mathcal U\).

If \(\varepsilon_{\mathcal U}\) and
\(\varepsilon'_{\mathcal U}\) are two augmentations, then
\[
  \varepsilon_{\mathcal U}
  =(\varepsilon_{\mathcal U}\otimes\varepsilon'_{\mathcal U})
    \circ\Delta_{\mathcal U}
  =\varepsilon'_{\mathcal U};
\]
the augmentation is therefore uniquely determined by~(iii).

If \((\mathcal U,\Delta_{\mathcal U})\) and
\((\mathcal V,\Delta_{\mathcal V})\) are two
\(\mathcal O_S\)-coalgebras, a \emph{morphism} from the first to the
second is a morphism of \(\mathcal O_S\)-modules
\(f:\mathcal U\to\mathcal V\) for which the diagrams
\begin{center}
  \phantomsection
  \label{sga3:VIIA:3.1-coalgebra-morphism-diagrams}
  % SGA3-VII-DIAG-022
  % BEGIN INCLUDED SGA3_English_current_buildable_source_20260823/source/inputs/VII/native_diagrams/exp7/viia_00_04/diag_022_coalgebra_morphism_diagrams.tex
\begin{minipage}{0.47\linewidth}
\centering
\begin{tikzcd}[column sep=large,row sep=large]
\mathcal U
  \arrow[r,"f"]
  \arrow[d,"\Delta_{\mathcal U}"]
&
\mathcal V
  \arrow[d,"\Delta_{\mathcal V}"]
\\
\mathcal U\otimes\mathcal U
  \arrow[r,"f\otimes f"]
&
\mathcal V\otimes\mathcal V
\end{tikzcd}
\end{minipage}
\hfill
\text{and}
\hfill
\begin{minipage}{0.37\linewidth}
\centering
\begin{tikzcd}[column sep=large,row sep=large]
\mathcal U
  \arrow[rr,"f"]
  \arrow[dr,"\varepsilon_{\mathcal U}"']
&&
\mathcal V
  \arrow[dl,"\varepsilon_{\mathcal V}"]
\\
&
\mathcal O_S
\end{tikzcd}
\end{minipage}

% END INCLUDED SGA3_English_current_buildable_source_20260823/source/inputs/VII/native_diagrams/exp7/viia_00_04/diag_022_coalgebra_morphism_diagrams.tex

\end{center}
are commutative. Morphisms of coalgebras compose like morphisms of
\(\mathcal O_S\)-modules, so we may speak of the
\emph{category of \(\mathcal O_S\)-coalgebras}.

% Authority: Exp7A local p. 16; combined p. 469; printed p. 458.
\medskip
\noindent\textbf{3.1.0.}
\phantomsection
\label{sga3:VIIA:3.1.0}
\footnote{The editors added the numbering \SGARefLink{sga3:VIIA:3.1.0}{3.1.0} for later
references.}
This category has finite products: its final object is the
\(\mathcal O_S\)-module \(\mathcal O_S\), with the identity as
comultiplication. The product of two coalgebras
\((\mathcal U,\Delta_{\mathcal U})\) and
\((\mathcal V,\Delta_{\mathcal V})\) is the tensor product
\(\mathcal U\otimes_{\mathcal O_S}\mathcal V\), whose
comultiplication is the composite
\begin{center}
  \phantomsection
  \label{sga3:VIIA:3.1.0-product-comultiplication}
  % SGA3-VII-DIAG-023
  % BEGIN INCLUDED SGA3_English_current_buildable_source_20260823/source/inputs/VII/native_diagrams/exp7/viia_00_04/diag_023_product_comultiplication_sequence.tex
\resizebox{0.98\linewidth}{!}{%
\(
\mathcal U\otimes\mathcal V
\xrightarrow{\Delta_{\mathcal U}\otimes\Delta_{\mathcal V}}
\mathcal U\otimes\mathcal U\otimes\mathcal V\otimes\mathcal V
\xrightarrow{\operatorname{id}_{\mathcal U}\otimes\sigma\otimes\operatorname{id}_{\mathcal V}}
\mathcal U\otimes\mathcal V\otimes\mathcal U\otimes\mathcal V
\)
}

% END INCLUDED SGA3_English_current_buildable_source_20260823/source/inputs/VII/native_diagrams/exp7/viia_00_04/diag_023_product_comultiplication_sequence.tex

\end{center}
where \(\sigma(a\otimes b)=b\otimes a\). The canonical projections
of \(\mathcal U\otimes\mathcal V\) onto the factors
\(\mathcal U\) and \(\mathcal V\) are
\(\operatorname{id}_{\mathcal U}\otimes\varepsilon_{\mathcal V}\)
and
\(\varepsilon_{\mathcal U}\otimes\operatorname{id}_{\mathcal V}\),
respectively.\footnote{The preceding assertion was added by the
editors. Recall also that, to show that
\(\mathcal U\otimes\mathcal V\) is indeed the product of
\(\mathcal U\) and \(\mathcal V\) in the category of cocommutative
\(\mathcal O_S\)-cogebras, one argues as follows. Given an arbitrary
\(\mathcal O_S\)-cogebra \(\mathcal E\) and cogebra morphisms
\(f:\mathcal E\to\mathcal U\) and
\(g:\mathcal E\to\mathcal V\), any cogebra morphism
\(\varphi:\mathcal E\to\mathcal U\otimes\mathcal V\) satisfying
\(\operatorname{pr}_{\mathcal U}\circ\varphi=f\) and
\(\operatorname{pr}_{\mathcal V}\circ\varphi=g\) must equal
\((f\otimes g)\circ\Delta_{\mathcal E}\); and this is a cogebra
morphism if and only if it also equals
\((g\otimes f)\circ\Delta_{\mathcal E}\).}
The ``diagonal morphism''
\(\mathcal U\to\mathcal U\otimes\mathcal U\), corresponding to the
pair
\((\operatorname{id}_{\mathcal U},
  \operatorname{id}_{\mathcal U})\),
is precisely the comultiplication \(\Delta_{\mathcal U}\).

\medskip
\noindent\textbf{3.1.1.}
\phantomsection
\label{sga3:VIIA:3.1.1}
Let \(\mathcal A\) be a commutative \(\mathcal O_S\)-algebra which,
as an \(\mathcal O_S\)-module, is locally free and of finite type.
Put
\[
  \mathcal A^{*}
  =\mathcal Hom_{\mathcal O_S\text{-Mod.}}
    (\mathcal A,\mathcal O_S).
\]
The canonical morphism
\[
  \varphi:\mathcal A^{*}\otimes_{\mathcal O_S}\mathcal A^{*}
  \longrightarrow
  (\mathcal A\otimes_{\mathcal O_S}\mathcal A)^{*}
\]
is invertible. If
\(m:\mathcal A\otimes\mathcal A\to\mathcal A\) is the morphism
defining the multiplication of \(\mathcal A\), composition gives a
diagonal morphism
\begin{center}
  \phantomsection
  \label{sga3:VIIA:3.1.1-dual-coalgebra-sequence}
  % SGA3-VII-DIAG-024
  % BEGIN INCLUDED SGA3_English_current_buildable_source_20260823/source/inputs/VII/native_diagrams/exp7/viia_00_04/diag_024_dual_coalgebra_sequence.tex
\(
\Delta_{\mathcal A^*}:\quad
\mathcal A^*
\xrightarrow{m^*}
(\mathcal A\otimes\mathcal A)^*
\xrightarrow{\varphi^{-1}}
\mathcal A^*\otimes\mathcal A^*\,.
\)

% END INCLUDED SGA3_English_current_buildable_source_20260823/source/inputs/VII/native_diagrams/exp7/viia_00_04/diag_024_dual_coalgebra_sequence.tex

\end{center}
This diagonal morphism plainly makes \(\mathcal A^{*}\) into an
\(\mathcal O_S\)-coalgebra whose augmentation is the transpose of
the morphism \(\mathcal O_S\to\mathcal A\) defined by the unit
section of \(\mathcal A\). Moreover, it is clear that:

\begin{quote}
\itshape
The functor \(\mathcal A\mapsto\mathcal A^{*}\) is an
anti-equivalence from the category of \(\mathcal O_S\)-algebras
which are locally free and of finite type as
\(\mathcal O_S\)-modules onto the category of
\(\mathcal O_S\)-coalgebras which are locally free and of finite
type as \(\mathcal O_S\)-modules.
\end{quote}

\medskip
\noindent\textbf{3.1.2.}
\phantomsection
\label{sga3:VIIA:3.1.2}
To every \(\mathcal O_S\)-coalgebra \(\mathcal U\) there is
canonically associated an \(S\)-functor
\begin{center}
  \phantomsection
  \label{sga3:VIIA:3.1.2-spec-star-functor}
  % SGA3-VII-DIAG-025
  % BEGIN INCLUDED SGA3_English_current_buildable_source_20260823/source/inputs/VII/native_diagrams/exp7/viia_00_04/diag_025_spec_star_functor_sequence.tex
\(
\operatorname{Spec}^{*}\mathcal U:
(\mathbf{Sch}/S)^{\circ}\longrightarrow(\mathbf{Ens})
\,.
\)

% END INCLUDED SGA3_English_current_buildable_source_20260823/source/inputs/VII/native_diagrams/exp7/viia_00_04/diag_025_spec_star_functor_sequence.tex

\end{center}
Indeed, for every \(S\)-scheme \(q:T\to S\),
\[
  q^{*}(\mathcal U\otimes_{\mathcal O_S}\mathcal U)
  \simeq
  q^{*}(\mathcal U)\otimes_{\mathcal O_T}q^{*}(\mathcal U),
\]
so \(q^{*}(\Delta_{\mathcal U})\) makes
\(\mathcal U_T=q^{*}(\mathcal U)\) into an
\(\mathcal O_T\)-coalgebra. Thus, by definition and with an obvious
abuse of notation,\footnote{For
\(x\otimes y\in
 \Gamma(T,\mathcal U_T)\otimes_{\mathcal O(T)}
 \Gamma(T,\mathcal U_T)\),
its image in
\(\Gamma(T,\mathcal U_T\otimes_{\mathcal O_T}\mathcal U_T)\)
is again denoted \(x\otimes y\).}
\[
  (\operatorname{Spec}^{*}\mathcal U)(T)
  =
  \left\{
    x\in\Gamma(T,\mathcal U_T)
    \ \middle|\
    \varepsilon_{\mathcal U_T}(x)=1,\quad
    \Delta_{\mathcal U_T}(x)=x\otimes x
  \right\}.
\]

% Authority: Exp7A local p. 17; combined p. 470; printed p. 459.
Sections \(x\) of \(\mathcal U_T\) plainly correspond to morphisms
of \(\mathcal O_T\)-modules
\(\xi:\mathcal O_T\to\mathcal U_T\). The conditions
\(\varepsilon(x)=1\) and \(\Delta(x)=x\otimes x\) say exactly that
\(\xi\) is a morphism of coalgebras. Hence also
\[
  (\operatorname{Spec}^{*}\mathcal U)(T)
  =
  \operatorname{Hom}_{\mathcal O_T\text{-coalg.}}
    (\mathcal O_T,\mathcal U_T).
\]
In particular:

\medskip
\noindent\textbf{Proposition 3.1.2.1.}
\phantomsection
\label{sga3:VIIA:3.1.2.1}
\footnote{The editors added the numbering \SGARefLink{sga3:VIIA:3.1.2.1}{3.1.2.1} for later
references. Notice also that the \(S\)-functor
\(\operatorname{Spec}^{*}\mathcal U\) is a sheaf for the Zariski
topology; it is even an fpqc sheaf if \(\mathcal U\) is a
quasi-coherent \(\mathcal O_S\)-module.}
\textit{Let \(\mathcal A\) be a commutative
\(\mathcal O_S\)-algebra which is locally free of finite type as an
\(\mathcal O_S\)-module. Then the \(S\)-functor
\(\operatorname{Spec}^{*}\mathcal A^{*}\) is represented by
\(\operatorname{Spec}\mathcal A\).}

Indeed, for every \(S\)-scheme \(T\), there are canonical
isomorphisms
\[
  \begin{aligned}
  (\operatorname{Spec}^{*}\mathcal A^{*})(T)
  &=
  \operatorname{Hom}_{\mathcal O_T\text{-coalg.}}
    (\mathcal O_T,\mathcal A_T^{*})                                      \\
  &\simeq
  \operatorname{Hom}_{\mathcal O_T\text{-alg.}}
    (\mathcal A_T,\mathcal O_T)
  \simeq
  (\operatorname{Spec}\mathcal A)(T).
  \end{aligned}
\]

\medskip
\noindent\textbf{3.2.}
\phantomsection
\label{sga3:VIIA:3.2}
An \(\mathcal O_S\)-\emph{coalgebra in groups}, that is, a group in
the category of \(\mathcal O_S\)-coalgebras, consists of an
\(\mathcal O_S\)-coalgebra
\((\mathcal U,\Delta_{\mathcal U})\) together with three morphisms
of \(\mathcal O_S\)-coalgebras
\[
  m_{\mathcal U}:\mathcal U\otimes\mathcal U\longrightarrow\mathcal U,
  \qquad
  \eta_{\mathcal U}:\mathcal O_S\longrightarrow\mathcal U,
  \qquad
  c_{\mathcal U}:\mathcal U\longrightarrow\mathcal U,
\]
satisfying conditions \((\mathrm{ii})^{*}\),
\((\mathrm{iii})^{*}\), and~(vi) below. The requirement that
\(m_{\mathcal U}\) be a morphism of coalgebras is expressed by the
commutativity of diagrams~(iv) and~(v):
\begin{center}
  \phantomsection
  \label{sga3:VIIA:3.2-compatibility-iv-v}
  % SGA3-VII-DIAG-026
  % BEGIN INCLUDED SGA3_English_current_buildable_source_20260823/source/inputs/VII/native_diagrams/exp7/viia_00_04/diag_026_group_coalgebra_compatibility_iv_v.tex
{\scriptsize
\begin{tikzcd}[column sep=small,row sep=normal]
\text{\rm(iv)}
&
\mathcal U\otimes\mathcal U
  \arrow[rr,"m_{\mathcal U}"]
  \arrow[d,"\Delta_{\mathcal U}\otimes\Delta_{\mathcal U}"']
&&
\mathcal U
  \arrow[dd,"\Delta_{\mathcal U}"]
\\
&
\mathcal U\otimes\mathcal U\otimes\mathcal U\otimes\mathcal U
  \arrow[d,"\operatorname{id}_{\mathcal U}\otimes\sigma\otimes\operatorname{id}_{\mathcal U}"']
&&
\\
&
\mathcal U\otimes\mathcal U\otimes\mathcal U\otimes\mathcal U
  \arrow[rr,"m_{\mathcal U}\otimes m_{\mathcal U}"]
&&
\mathcal U\otimes\mathcal U
\\[0.5em]
\text{\rm(v)}
&
\mathcal U\otimes\mathcal U
  \arrow[rr,"m_{\mathcal U}"]
  \arrow[dr,"\varepsilon_{\mathcal U}\otimes\varepsilon_{\mathcal U}"']
&&
\mathcal U
  \arrow[dl,"\varepsilon_{\mathcal U}"]
\\
&&
\mathcal O_S
&
\end{tikzcd}}

% END INCLUDED SGA3_English_current_buildable_source_20260823/source/inputs/VII/native_diagrams/exp7/viia_00_04/diag_026_group_coalgebra_compatibility_iv_v.tex

\end{center}

\noindent
\((\mathrm{ii})^{*}\) The square
\begin{center}
  \phantomsection
  \label{sga3:VIIA:3.2-associativity-ii-star}
  % SGA3-VII-DIAG-027
  % BEGIN INCLUDED SGA3_English_current_buildable_source_20260823/source/inputs/VII/native_diagrams/exp7/viia_00_04/diag_027_group_coalgebra_associativity_ii_star.tex
\begin{tikzcd}[column sep=huge,row sep=large]
\mathcal U\otimes\mathcal U\otimes\mathcal U
  \arrow[r,"\operatorname{id}_{\mathcal U}\otimes m_{\mathcal U}"]
  \arrow[d,"m_{\mathcal U}\otimes\operatorname{id}_{\mathcal U}"']
&
\mathcal U\otimes\mathcal U
  \arrow[d,"m_{\mathcal U}"]
\\
\mathcal U\otimes\mathcal U
  \arrow[r,"m_{\mathcal U}"]
&
\mathcal U
\end{tikzcd}

% END INCLUDED SGA3_English_current_buildable_source_20260823/source/inputs/VII/native_diagrams/exp7/viia_00_04/diag_027_group_coalgebra_associativity_ii_star.tex

\end{center}
is commutative.

% Authority: Exp7A local p. 18; combined p. 471; printed p. 460.
\noindent
\((\mathrm{iii})^{*}\) The two composites
\begin{center}
  \phantomsection
  \label{sga3:VIIA:3.2-unit-sequences-iii-star}
  % SGA3-VII-DIAG-028
  % BEGIN INCLUDED SGA3_English_current_buildable_source_20260823/source/inputs/VII/native_diagrams/exp7/viia_00_04/diag_028_group_coalgebra_unit_sequences_iii_star.tex
\(
\begin{gathered}
\mathcal U\simeq\mathcal U\otimes\mathcal O_S
\xrightarrow{\operatorname{id}_{\mathcal U}\otimes\eta_{\mathcal U}}
\mathcal U\otimes\mathcal U
\xrightarrow{m_{\mathcal U}}
\mathcal U
\\[0.7em]
\mathcal U\simeq\mathcal O_S\otimes\mathcal U
\xrightarrow{\eta_{\mathcal U}\otimes\operatorname{id}_{\mathcal U}}
\mathcal U\otimes\mathcal U
\xrightarrow{m_{\mathcal U}}
\mathcal U.
\end{gathered}
\)

% END INCLUDED SGA3_English_current_buildable_source_20260823/source/inputs/VII/native_diagrams/exp7/viia_00_04/diag_028_group_coalgebra_unit_sequences_iii_star.tex

\end{center}
are the identity morphism of \(\mathcal U\).

\smallskip
\noindent
(vi) The composite
\begin{center}
  \phantomsection
  \label{sga3:VIIA:3.2-antipode-sequence-vi}
  % SGA3-VII-DIAG-029
  % BEGIN INCLUDED SGA3_English_current_buildable_source_20260823/source/inputs/VII/native_diagrams/exp7/viia_00_04/diag_029_group_coalgebra_antipode_sequence_vi.tex
\(
\mathcal U
\xrightarrow{\Delta_{\mathcal U}}
\mathcal U\otimes\mathcal U
\xrightarrow{c_{\mathcal U}\otimes\operatorname{id}_{\mathcal U}}
\mathcal U\otimes\mathcal U
\xrightarrow{m_{\mathcal U}}
\mathcal U
\,.
\)

% END INCLUDED SGA3_English_current_buildable_source_20260823/source/inputs/VII/native_diagrams/exp7/viia_00_04/diag_029_group_coalgebra_antipode_sequence_vi.tex

\end{center}
is equal to \(\eta_{\mathcal U}\circ\varepsilon_{\mathcal U}\).

\medskip
\noindent\textbf{3.2.1.}
\phantomsection
\label{sga3:VIIA:3.2.1}
The morphisms \(\eta_{\mathcal U}\) and \(c_{\mathcal U}\) are
uniquely determined by \(m_{\mathcal U}\). Conditions
\((\mathrm{ii})^{*}\) and \((\mathrm{iii})^{*}\) simply say that
\(m_{\mathcal U}\) makes \(\mathcal U\) into an
\(\mathcal O_S\)-algebra whose unit section is the image under
\(\eta_{\mathcal U}\) of the unit section of \(\mathcal O_S\).
Condition~(iv) also says that the diagonal morphism
\(\Delta_{\mathcal U}\) is compatible with multiplication. Indeed,
\[
  \Delta_{\mathcal U}:\mathcal U\longrightarrow
  \mathcal U\otimes\mathcal U
\]
must be a homomorphism of coalgebras in groups, which also implies
the commutativity of the triangle
\begin{center}
  \phantomsection
  \label{sga3:VIIA:3.2.1-unit-triangle-v-star}
  % SGA3-VII-DIAG-030
  % BEGIN INCLUDED SGA3_English_current_buildable_source_20260823/source/inputs/VII/native_diagrams/exp7/viia_00_04/diag_030_group_coalgebra_unit_triangle_v_star.tex
\(\textup{(v)}^{*}\qquad\)
\begin{tikzcd}[column sep=huge,row sep=large,baseline=(current bounding box.center)]
&
\mathcal O_S
  \arrow[dl,"\eta_{\mathcal U}"']
  \arrow[dr,"\eta_{\mathcal U}\otimes\eta_{\mathcal U}"]
&
\\
\mathcal U
  \arrow[rr,"\Delta_{\mathcal U}"']
&& 
\mathcal U\otimes\mathcal U\,.
\end{tikzcd}

% END INCLUDED SGA3_English_current_buildable_source_20260823/source/inputs/VII/native_diagrams/exp7/viia_00_04/diag_030_group_coalgebra_unit_triangle_v_star.tex

\end{center}

On the other hand, as in every category, the antipode
\(c_{\mathcal U}\) is an isomorphism from \(\mathcal U\) onto the
``opposite'' group object;\footnote{That is, the object endowed with
the multiplication
\(m'_{\mathcal U}=m_{\mathcal U}\circ\sigma\).}
in particular, \(c_{\mathcal U}\) induces an algebra isomorphism
from \(\mathcal U\) onto the opposite algebra
\(\mathcal U^{\circ}\).

\medskip
\noindent\textbf{3.2.2.}
\phantomsection
\label{sga3:VIIA:3.2.2}
Since the functor
\(\mathcal U\mapsto\operatorname{Spec}^{*}\mathcal U\) commutes
with finite products, it carries a coalgebra in groups to a
group-valued \(S\)-functor. Indeed, for every \(S\)-scheme \(T\),
the elements
\[
  x\in\Gamma(T,\mathcal U_T)
\]
belonging to
\((\operatorname{Spec}^{*}\mathcal U)(T)\) form a group under the
multiplication of the algebra \(\Gamma(T,\mathcal U_T)\), and the
inverse of \(x\) is \(c_{\mathcal U}(x)\). By \SGARefLink{sga3:VIIA:3.1.2.1}{3.1.2.1}:

\medskip
\noindent\textbf{Scholium 3.2.2.1.}
\phantomsection
\label{sga3:VIIA:3.2.2.1}
\footnote{The editors added this scholium, which is implicit in the
original.}
\textit{Let \(\mathcal U\) be an \(\mathcal O_S\)-coalgebra in
groups, finite and locally free as an \(\mathcal O_S\)-module. Then
the group-valued \(S\)-functor
\(\operatorname{Spec}^{*}\mathcal U\) is represented by the finite
locally free \(S\)-group \(\operatorname{Spec}\mathcal U^{*}\).}

\medskip
\noindent\textbf{Remark 3.2.2.2.}
\phantomsection
\label{sga3:VIIA:3.2.2.2}
Let \(\mathcal L\) be an \(\mathcal O_S\)-Lie algebra and
\(\mathcal U(\mathcal L)\) its enveloping algebra, that is, the
sheaf on \(S\) associated with the presheaf assigning to every open
subset \(V\) the enveloping algebra
\[
  U\bigl(\Gamma(V,\mathcal L)\bigr)
\]
of the Lie algebra \(\Gamma(V,\mathcal L)\).

Every homomorphism from \(\mathcal L\) to the Lie algebra underlying
an associative \(\mathcal O_S\)-algebra factors uniquely through
the canonical morphism
\(\mathcal L\to\mathcal U(\mathcal L)\). Besides the functoriality
of \(\mathcal U(\mathcal L)\) in \(\mathcal L\), this universal
property also shows that the enveloping algebra of a product of Lie
algebras is identified with the tensor product of their enveloping
algebras.

% Authority: Exp7A local p. 19; combined p. 472; printed p. 461.
In particular, the diagonal morphism
\(\delta:\mathcal L\to\mathcal L\times\mathcal L\) induces an
algebra homomorphism
\[
  \Delta:\mathcal U(\mathcal L)\longrightarrow
  \mathcal U(\mathcal L\times\mathcal L)
  \simeq
  \mathcal U(\mathcal L)\otimes\mathcal U(\mathcal L).
\]
The zero morphism \(\mathcal L\to0\) induces a homomorphism
\[
  \varepsilon:\mathcal U(\mathcal L)\longrightarrow
  \mathcal U(0)\simeq\mathcal O_S.
\]
The isomorphism \(x\mapsto-x\) from \(\mathcal L\) onto the opposite
Lie algebra \(\mathcal L^{\circ}\) induces an anti-isomorphism \(c\)
of the algebra \(\mathcal U(\mathcal L)\). One then checks readily
that the multiplication \(m\) of \(\mathcal U(\mathcal L)\) makes
\((\mathcal U(\mathcal L),\Delta)\) into an
\(\mathcal O_S\)-coalgebra in groups with augmentation
\(\varepsilon\) and antipode \(c\).\footnote{The group-valued
\(S\)-functor \(\operatorname{Spec}^{*}\mathcal U(\mathcal L)\) is
not representable in general. Later, in \SGARefLink{sga3:VIIA:5.5}{5.5}, it will be shown that
if \(S\) is a scheme of characteristic \(p\), if \(\mathcal L\) is
finite locally free over \(\mathcal O_S\), and if
\(\mathcal U_p(\mathcal L)\) is its restricted enveloping algebra
(see \SGARefLink{sga3:VIIA:5.3}{5.3}), then
\(\operatorname{Spec}^{*}\mathcal U_p(\mathcal L)\) is represented
by a finite locally free \(S\)-group.}

\medskip
\noindent\textbf{3.2.3.}
\phantomsection
\label{sga3:VIIA:3.2.3}
\footnote{The editors expanded this paragraph.}
Let \(\mathcal U\) be an \(\mathcal O_S\)-coalgebra in groups. We
shall see that the group-valued \(S\)-functor
\[
  G=\operatorname{Spec}^{*}\mathcal U
\]
is \emph{very good}, in the sense of \SGARefLink{sga3:ii:definition:4-6}{Expos\'e II, 4.6} and \SGARefLink{sga3:ii:definition:4-10}{4.10}.

Let \(\mathcal M\) be a free \(\mathcal O_S\)-module of rank \(r\),
and let \(T\to S\) be an \(S\)-scheme. Since
\[
  I_T(\mathcal M)=
  \operatorname{Spec}(\mathcal O_T\oplus\mathcal M_T),
\]
the morphism \(\pi:I_T(\mathcal M)\to T\) is affine, and
\[
  \pi_{*}\bigl(\mathcal U_{I_T(\mathcal M)}\bigr)
  =
  \mathcal U_T\otimes_{\mathcal O_T}
  \pi_{*}\bigl(\mathcal O_{I_T(\mathcal M)}\bigr)
  =
  \mathcal U_T\otimes_{\mathcal O_T}
  (\mathcal O_T\oplus\mathcal M_T).
\]
Consequently,
\begin{equation}
  \Gamma\bigl(I_T(\mathcal M),\mathcal U_{I_T(\mathcal M)}\bigr)
  \simeq
  \Gamma(T,\mathcal U_T)\otimes_{\mathcal O(T)}
  \bigl(\mathcal O(T)\oplus\Gamma(T,\mathcal M_T)\bigr).
  \tag{1}
  \label{sga3:VIIA:3.2.3-eq1}
\end{equation}

Let \((d_1,\ldots,d_r)\) be a basis of \(\mathcal M\). An element
\[
  u_0+\sum_i u_i d_i
  \in
  \Gamma\bigl(I_T(\mathcal M),\mathcal U_{I_T(\mathcal M)}\bigr)
\]
belongs to \(G(I_T(\mathcal M))\) if and only if
\[
  1
  =\varepsilon\left(u_0+\sum_i u_i d_i\right)
  =\varepsilon(u_0)+\sum_i\varepsilon(u_i)d_i
\]
and
\[
  \left(u_0+\sum_i u_i d_i\right)
  \otimes
  \left(u_0+\sum_i u_i d_i\right)
  =
  \Delta\left(u_0+\sum_i u_i d_i\right)
  =
  \Delta(u_0)+\sum_i\Delta(u_i)d_i.
\]
Equivalently,
\begin{equation}
  \left\{
  \begin{aligned}
    \varepsilon(u_0)&=1,&
    \Delta u_0&=u_0\otimes u_0
      &&\bigl(\text{that is, }u_0\in G(T)\bigr),\\
    \varepsilon(u_i)&=0,&
    \Delta(u_i)&=u_i\otimes u_0+u_0\otimes u_i
      &&(i=1,\ldots,r).
  \end{aligned}
  \right.
  \tag{2}
  \label{sga3:VIIA:3.2.3-eq2}
\end{equation}

Moreover, the morphism
\(G(I_T(\mathcal M))\to G(T)\) corresponding to the zero section of
\(I_T(\mathcal M)\to T\) is given by
\[
  u_0+\sum_i u_i d_i\longmapsto u_0.
\]
Combining this with (1) and (2), one sees that if \(\mathcal N\) is
a second free \(\mathcal O_S\)-module of finite rank, then the
diagram of sets
\begin{center}
  \phantomsection
  \label{sga3:VIIA:3.2.3-condition-E-square}
  % SGA3-VII-DIAG-031
  % BEGIN INCLUDED SGA3_English_current_buildable_source_20260823/source/inputs/VII/native_diagrams/exp7/viia_00_04/diag_031_condition_E_cartesian_square.tex
\begin{tikzcd}[column sep=huge,row sep=large]
G(I_T(\mathcal M\oplus\mathcal N))
  \arrow[r]
  \arrow[d]
&
G(I_T(\mathcal N))
  \arrow[d]
\\
G(I_T(\mathcal M))
  \arrow[r]
&
G(T)
\end{tikzcd}

% END INCLUDED SGA3_English_current_buildable_source_20260823/source/inputs/VII/native_diagrams/exp7/viia_00_04/diag_031_condition_E_cartesian_square.tex

\end{center}
is cartesian; that is, \(G\) satisfies condition~(E) of \SGARefLink{sga3:ii:definition:3-5}{Expos\'e
II, 3.5}.

% Authority: Exp7A local p. 20; combined p. 473; printed p. 462.
Write
\[
  \operatorname{Prim}\Gamma(T,\mathcal U_T)
\]
for the \(\mathcal O(T)\)-submodule of \(\Gamma(T,\mathcal U_T)\)
formed by the \emph{primitive elements}, that is, the elements \(u\)
satisfying, with the abuse of notation noted in \SGARefLink{sga3:VIIA:3.1.2}{3.1.2},
\[
  \Delta u=u\otimes1+1\otimes u,
  \qquad
  \varepsilon(u)=0.
\]
\footnote{The second condition follows from the first, since the
first gives
\(u=(\operatorname{id}\otimes\varepsilon)\Delta(u)
=u+\varepsilon(u)\), and hence \(\varepsilon(u)=0\).}

Since \((\operatorname{Lie}G)(T)\) is the set of elements
\(u_0+ud\in G(I_T)\) above the unit element \(u_0=1\) of \(G(T)\),
there is an isomorphism of \(\mathcal O(T)\)-modules, functorial in
\(T\),\footnote{The \(\mathcal O_S\)-module structure on
\(\operatorname{Lie}G\) is defined in \SGARefLink{sga3:ii:proposition:3-6}{Expos\'e II, Proposition
3.6}.}
\[
  (\operatorname{Lie}G)(T)
  \simeq
  \operatorname{Prim}\Gamma(T,\mathcal U_T).
\]
On the other hand, \SGARefLink{sga3:VIIA:3.2.3-eq1}{(1)} yields
\[
  \operatorname{Prim}
    \Gamma\bigl(I_T(\mathcal M),
      \mathcal U_{I_T(\mathcal M)}\bigr)
  \simeq
  \operatorname{Prim}\Gamma(T,\mathcal U_T)
  \otimes_{\mathcal O(T)}\mathcal O(I_T(\mathcal M)).
\]
It follows that the natural morphism of
\(\mathcal O(I_T(\mathcal M))\)-modules
\[
  (\operatorname{Lie}G)(T)
  \otimes_{\mathcal O(T)}\mathcal O(I_T(\mathcal M))
  \longrightarrow
  (\operatorname{Lie}G)(I_T(\mathcal M))
\]
is an isomorphism; thus \(\operatorname{Lie}G\) is a good
\(\mathcal O_S\)-module (see Expos\'e II, Definition 4.4).

Therefore \(G\) is a good group-valued \(S\)-functor (see \SGARefLink{sga3:ii:definition:4-6}{Expos\'e
II, Definition 4.6}), and by \SGARefLink{sga3:ii:definition:4-7-2}{Expos\'e II, 4.7.2},
\(\operatorname{Lie}G\) carries an \(\mathcal O_S\)-bilinear Lie
bracket satisfying the Jacobi identity. It remains to show that
\(G\) is very good, that is, that the bracket on
\((\operatorname{Lie}G)(T)\) satisfies
\([u,u]=0\) for every \(u\in(\operatorname{Lie}G)(T)\); see \SGARefLink{sga3:ii:definition:4-10}{Expos\'e
II, 4.10}.

Let \(u,v\) be two elements of \((\operatorname{Lie}G)(T)\), that
is, two primitive elements of \(\Gamma(T,\mathcal U_T)\). Put
\[
  I=\operatorname{Spec}\mathcal O_S[d]/(d^2),
  \qquad
  I'=\operatorname{Spec}\mathcal O_S[d']/(d'^2).
\]
Since the composition law on \(G(I\times I')\) is induced by the
multiplication in \(\mathcal U_{I\times I'}\), one has in
\(G(I\times I')\)
\[
  \begin{aligned}
  &(1+ud)(1+vd')(1+ud)^{-1}(1+vd')^{-1}\\
  &\qquad
    =(1+ud)(1+vd')(1-ud)(1-vd')\\
  &\qquad
    =1+(uv-vu)dd'.
  \end{aligned}
\]
\begin{quote}
\small
\phantomsection
\label{sga3:VIIA:3.2.3-source-note-missing-prime}
\emph{Source note.} The re-edition prints \(1-vd\), without a prime,
in the last factor of the middle expression. Since that factor is
the inverse of \(1+vd'\), and since the displayed result is
proportional to \(dd'\), the English reconstruction makes the
required reading \(1-vd'\) explicit. Reinhold's comparison Markdown
reproduces the printed missing prime.
\end{quote}
By the description of the bracket \([u,v]\) given before
\SGARefLink{sga3:ii:proposition:4-8}{Proposition 4.8} of Expos\'e II, it follows that
\[
  [u,v]=uv-vu,
\]
where the right-hand side is the commutator of \(u\) and \(v\) in
the algebra \(\Gamma(T,\mathcal U_T)\). Hence \([u,u]=0\), and we
have proved:

\medskip
\noindent\textbf{Proposition.}
\phantomsection
\label{sga3:VIIA:3.2.3-primitive-elements-proposition}
\textit{Let \(\mathcal U\) be an \(\mathcal O_S\)-coalgebra in
groups. The group-valued \(S\)-functor
\(G=\operatorname{Spec}^{*}\mathcal U\) is very good, and there is
an isomorphism
\[
  \operatorname{Lie}G
  \simeq
  \operatorname{Prim}\mathcal W(\mathcal U)
\]
of \(\mathcal O_S\)-Lie algebras, where
\(\operatorname{Prim}\mathcal W(\mathcal U)\) denotes the functor
which assigns to every \(T\to S\) the \(\mathcal O(T)\)-Lie algebra
of primitive elements of
\(\mathcal W(\mathcal U)(T)=\Gamma(T,\mathcal U_T)\).}
\footnote{The editors added this proposition to summarize the
preceding discussion.}

% Authority: Exp7A local p. 21; combined p. 474; printed p. 463.
\medskip
\noindent\textbf{3.3.}
\phantomsection
\label{sga3:VIIA:3.3}
Suppose finally that \(\mathcal U\) is a
\emph{coalgebra in commutative groups}, that is, that the triangle
\begin{center}
  \phantomsection
  \label{sga3:VIIA:3.3-commutative-triangle-i-star}
  % SGA3-VII-DIAG-032
  % BEGIN INCLUDED SGA3_English_current_buildable_source_20260823/source/inputs/VII/native_diagrams/exp7/viia_00_04/diag_032_commutative_group_coalgebra_triangle_i_star.tex
\(\textup{(i)}^{*}\qquad\)
\begin{tikzcd}[column sep=huge,row sep=large,baseline=(current bounding box.center)]
\mathcal U\otimes\mathcal U
  \arrow[rr,"\sigma"]
  \arrow[dr,"m_{\mathcal U}"']
&&
\mathcal U\otimes\mathcal U
  \arrow[dl,"m_{\mathcal U}"]
\\
&
\mathcal U
&
\end{tikzcd}

% END INCLUDED SGA3_English_current_buildable_source_20260823/source/inputs/VII/native_diagrams/exp7/viia_00_04/diag_032_commutative_group_coalgebra_triangle_i_star.tex

\end{center}
is commutative, or equivalently that \(m_{\mathcal U}\) makes
\(\mathcal U\) into a commutative \(\mathcal O_S\)-algebra.
Conditions \SGARefLink{sga3:vii:paragraph:5-1:equation:i}{(i)}, \SGARefLink{sga3:vii:paragraph:5-1:equation:ii}{(ii)}, (iii), (iv), (v), (vi),
\((\mathrm{i})^{*}\), \((\mathrm{ii})^{*}\),
\((\mathrm{iii})^{*}\), and \((\mathrm{v})^{*}\) then also say that
\(\mathcal U\) is a cogroup in the category of commutative
\(\mathcal O_S\)-algebras. Thus, if in addition \(\mathcal U\) is a
quasi-coherent \(\mathcal O_S\)-module, the affine \(S\)-scheme
\(\operatorname{Spec}\mathcal U\) is a commutative \(S\)-group
scheme.

In this case the diagonal morphism \(\Delta'\) of
\(\mathcal O_S[T,T^{-1}]\) sends \(T\) to \(T\otimes T\).
Consequently, morphisms of \(S\)-groups from
\(\operatorname{Spec}\mathcal U\) to
\(\mathbb G_{m,S}\) (Expos\'e I, 4.3.2) correspond bijectively to
morphisms of unital \(\mathcal O_S\)-algebras
\[
  \varphi:\mathcal O_S[T,T^{-1}]\longrightarrow\mathcal U
\]
such that
\[
  (\varphi\otimes\varphi)\circ\Delta'
  =\Delta_{\mathcal U}\circ\varphi.
\]
In this situation,
\(\varepsilon_{\mathcal U}\circ\varphi\) is the neutral element of
\(\mathbb G_{m,S}(S)\), that is, the augmentation
\(\varepsilon'\). Such a morphism \(\varphi\) is determined by the
image \(\varphi(T)\), which must be an invertible element
\(x\in\mathcal U\) satisfying
\[
  \Delta_{\mathcal U}x=x\otimes x,
  \qquad
  \varepsilon_{\mathcal U}(x)=\varepsilon'(T)=1.
\]
Therefore
\[
  \operatorname{Hom}_{S\text{-gr.}}
  (\operatorname{Spec}\mathcal U,\mathbb G_{m,S})
  \simeq
  (\operatorname{Spec}^{*}\mathcal U)(S).
\]
Since this formula remains valid after arbitrary base change, it
gives
\[
  \operatorname{Spec}^{*}\mathcal U
  \simeq
  \operatorname{Hom}_{S\text{-gr.}}
  (\operatorname{Spec}\mathcal U,\mathbb G_{m,S}).
\]
We have therefore proved:

\medskip
\noindent\textbf{Proposition 3.3.0.}
\phantomsection
\label{sga3:VIIA:3.3.0}
\textit{If \(\mathcal U\) is an \(\mathcal O_S\)-coalgebra in
commutative groups, quasi-coherent as an \(\mathcal O_S\)-module,
then the affine \(S\)-scheme
\(G=\operatorname{Spec}\mathcal U\) is a commutative \(S\)-group
scheme, and there is an isomorphism of group-valued \(S\)-functors
\[
  \operatorname{Spec}^{*}\mathcal U
  \simeq
  \operatorname{Hom}_{S\text{-gr.}}(G,\mathbb G_{m,S}).
\]}

If, moreover, \(\mathcal U\) is locally free of finite type as an
\(\mathcal O_S\)-module, then by \SGARefLink{sga3:VIIA:3.1.2.1}{3.1.2.1} the group-valued
\(S\)-functor \(\operatorname{Spec}^{*}\mathcal U\) is represented
by \(\operatorname{Spec}\mathcal U^{*}\). Hence:

\medskip
\noindent\textbf{Proposition 3.3.1 (Cartier duality).}
\phantomsection
\label{sga3:VIIA:3.3.1}
\textit{The functor
\[
  \mathscr A(G)\longmapsto
  \mathscr A(G)^{*}
  =
  \mathcal Hom_{\mathcal O_S\text{-Mod.}}
  (\mathscr A(G),\mathcal O_S)
\]
induces a duality\textsuperscript{\(\ast\)} on the category of
commutative finite locally free \(S\)-group schemes; it associates
to \(G\) the \(S\)-group
\(\operatorname{Hom}_{S\text{-gr.}}(G,\mathbb G_{m,S})\).}

\begingroup
\renewcommand{\thefootnote}{\fnsymbol{footnote}}
\footnotetext[1]{A duality of a category \(\mathscr C\) is a pair
\((D,\varphi)\) consisting of a contravariant functor
\(D:\mathscr C\to\mathscr C\) and a natural isomorphism
\(\varphi:\operatorname{Id}_{\mathscr C}\to DD\), such that the
isomorphisms
\(\varphi D:D\to DDD\) and
\(D\varphi^{-1}:DDD\to D\) are inverse to one another. The editors
corrected \(D\varphi\) to \(D\varphi^{-1}\).}
\endgroup

% Next cursor: Exp7A local p. 22, combined-reader p. 475,
% re-edition printed p. 464, at the Section 4 heading.

% END INCLUDED SGA3_English_current_buildable_source_20260823/source/inputs/VII/tex/components/03_expose_VIIA_section3_coalgebras_cartier_en.tex

% BEGIN INCLUDED SGA3_English_current_buildable_source_20260823/source/inputs/VII/tex/components/04_expose_VIIA_section4_frobeniuseries_40_through_412_en.tex
% SGA 3, Expose VII_A, English reconstruction -- Loop 1.
% Scope: authority local pages 22--24, Section 4 opening through 4.1.2,
% stopping before 4.1.3 on local page 25.
% Combined-reader pages: 475--477. Re-edition printed pages: 464--466.
% Original-edition marginal locators within this span: 434--436.
% Authority: Polo--Gille Exp7A-13oct24.pdf,
% SHA-256 D46DB8C79D0385643B6EED4B96900613248375F9B73B4202478A4F57FD3E8522.
% Jacob Reinhold's CC BY 4.0 Markdown was consulted as comparison lineage.

\section*{4. ``Frobeniuseries''}
\phantomsection
\label{sga3:VIIA:4}
\addcontentsline{toc}{section}{4. ``Frobeniuseries''}

% Authority: Exp7A local p. 22; combined p. 475; printed p. 464.
Fix a prime number \(p\), and let
\((\mathbf{Sch}/\mathbb F_p)\) be the category of schemes of
characteristic \(p\), that is, schemes over the prime field
\(\mathbb F_p\). Following the general conventions of this seminar,
we identify \((\mathbf{Sch}/\mathbb F_p)\) with a subcategory of
\((\mathbf{Sch}/\mathbb F_p)^\wedge\) by means of the functor \(h\)
of \SGARefLink{sga3:i:subsection:1-1}{Expos\'e I, 1.1}. We likewise use the isomorphism
\[
  \operatorname{Hom}(h_X,F)\xrightarrow{\sim}F(X)
\]
defined there to identify these sets whenever \(X\) is an
\(\mathbb F_p\)-scheme and \(F\) is an object of
\((\mathbf{Sch}/\mathbb F_p)^\wedge\).

\medskip
\noindent\textbf{Notation 4.0.}
\phantomsection
\label{sga3:VIIA:4.0}
\footnote{The editors added the number \SGARefLink{sga3:VIIA:4.0}{4.0} for later references.}
\textit{If \(T\) is an \(\mathbb F_p\)-scheme, a \(T\)-functor is a
morphism \(q:F\to T\) in
\((\mathbf{Sch}/\mathbb F_p)^\wedge\) with target \(T\). For every
\(T\)-scheme \(r:X\to T\), the set of \(T\)-morphisms \(X\to F\),
that is, of \(\mathbb F_p\)-morphisms \(s:X\to F\) satisfying
\(q\circ s=r\), is denoted}
\[
  q(r),\qquad q(X/T),\qquad F(r),\qquad\text{or}\qquad F(X/T).
\]
\textit{It may also be denoted \(F(X)\) when no confusion with
\(\operatorname{Hom}(h_X,F)\) can result.}

\medskip
\noindent\textbf{4.1.}
\phantomsection
\label{sga3:VIIA:4.1}
For every scheme \(S\) of characteristic \(p\), let
\(\operatorname{fr}(S)\), or simply \(\operatorname{fr}\), denote the
endomorphism of \(S\) which is the identity on the underlying
topological space and which sends a section \(x\) of \(\mathcal O_S\)
over an open subset to \(x^p\).

The assignment
\[
  \operatorname{fr}:S\longmapsto\operatorname{fr}(S)
\]
is an endomorphism of the identity functor of
\((\mathbf{Sch}/\mathbb F_p)\).\footnote{Thus, for every morphism
\(f:Y\to X\) of \(\mathbb F_p\)-schemes, the following square
commutes:
\par\noindent\parbox{\linewidth}{%
  \centering
  % SGA3-VII-DIAG-033
  % BEGIN INCLUDED SGA3_English_current_buildable_source_20260823/source/inputs/VII/native_diagrams/exp7/viia_00_04/diag_033_absolute_frobenius_naturality_square.tex
\begin{tikzcd}[column sep=large,row sep=large]
Y
  \arrow[r,"f"]
  \arrow[d,"\operatorname{fr}(Y)"']
&
X
  \arrow[d,"\operatorname{fr}(X)"]
\\
Y
  \arrow[r,"f"]
&
X\,.
\end{tikzcd}

% END INCLUDED SGA3_English_current_buildable_source_20260823/source/inputs/VII/native_diagrams/exp7/viia_00_04/diag_033_absolute_frobenius_naturality_square.tex

  \par\vspace*{\baselineskip}
  \vspace*{\baselineskip}
}\par}
This has the following consequences. If \(E\) is an
\(\mathbb F_p\)-functor, the assignment which sends every
\(\mathbb F_p\)-scheme \(S\) to the endomorphism
\(E(\operatorname{fr}(S))\) of \(E(S)\) is a functorial endomorphism
of \(E\), denoted \(\operatorname{fr}(E)\) or simply
\(\operatorname{fr}\). This notation is compatible with the
identification of \((\mathbf{Sch}/\mathbb F_p)\) with a subcategory
of \((\mathbf{Sch}/\mathbb F_p)^\wedge\). Moreover,
\(E\mapsto\operatorname{fr}(E)\) is itself an endomorphism of the
identity functor of \((\mathbf{Sch}/\mathbb F_p)^\wedge\), again
denoted \(\operatorname{fr}\).\footnote{The morphism
\(\operatorname{fr}(X)\) is called the \emph{absolute} Frobenius
morphism of \(X\), to distinguish it from the relative Frobenius
morphism \(\operatorname{Fr}(X/S)\) introduced below.}

For every \(\mathbb F_p\)-scheme \(S\) and every \(S\)-functor
\(q:X\to S\), let \(X^{(p/S)}\), or \(X^{(p)}\), denote the pullback
of \(X\) along \(\operatorname{fr}(S)\):
\begin{center}
  \phantomsection
  \label{sga3:VIIA:4.1-frobenius-basechange-square}
  % SGA3-VII-DIAG-034
  % BEGIN INCLUDED SGA3_English_current_buildable_source_20260823/source/inputs/VII/native_diagrams/exp7/viia_00_04/diag_034_frobenius_basechange_square.tex
\begin{tikzcd}[column sep=large,row sep=large]
X^{(p/S)}
  \arrow[r,"\operatorname{pr}_X"]
  \arrow[d]
&
X
  \arrow[d,"q"]
\\
S
  \arrow[r,"\operatorname{fr}(S)"]
&
S\,.
\end{tikzcd}

% END INCLUDED SGA3_English_current_buildable_source_20260823/source/inputs/VII/native_diagrams/exp7/viia_00_04/diag_034_frobenius_basechange_square.tex

\end{center}

% Authority: Exp7A local p. 23; combined p. 476; printed p. 465.
The commutative square
\begin{center}
  \phantomsection
  \label{sga3:VIIA:4.1-absolute-frobenius-square}
  % SGA3-VII-DIAG-035
  % BEGIN INCLUDED SGA3_English_current_buildable_source_20260823/source/inputs/VII/native_diagrams/exp7/viia_00_04/diag_035_absolute_frobenius_square.tex
\begin{tikzcd}[column sep=large,row sep=large]
X
  \arrow[r,"\operatorname{fr}(X)"]
  \arrow[d,"q"']
&
X
  \arrow[d,"q"]
\\
S
  \arrow[r,"\operatorname{fr}(S)"]
&
S
\end{tikzcd}

% END INCLUDED SGA3_English_current_buildable_source_20260823/source/inputs/VII/native_diagrams/exp7/viia_00_04/diag_035_absolute_frobenius_square.tex

\end{center}
induces an \(S\)-morphism
\[
  \operatorname{Fr}(X/S):X\longrightarrow X^{(p/S)}
\]
such that
\[
  \operatorname{fr}(X)
  =
  \operatorname{pr}_X\circ\operatorname{Fr}(X/S):
\]
\begin{center}
  \phantomsection
  \label{sga3:VIIA:4.1-relative-frobenius-factorization}
  % SGA3-VII-DIAG-036
  % BEGIN INCLUDED SGA3_English_current_buildable_source_20260823/source/inputs/VII/native_diagrams/exp7/viia_00_04/diag_036_relative_frobenius_factorization.tex
\begin{tikzcd}[column sep=large,row sep=large]
X
  \arrow[drr,bend left=24,"\operatorname{fr}(X)"]
  \arrow[dr,"\operatorname{Fr}(X/S)"]
  \arrow[ddr,"q"']
\\
&
X^{(p/S)}
  \arrow[r,"\operatorname{pr}_X"]
  \arrow[d]
&
X
  \arrow[d,"q"]
\\
&
S
  \arrow[r,"\operatorname{fr}(S)"]
&
S\,.
\end{tikzcd}

% END INCLUDED SGA3_English_current_buildable_source_20260823/source/inputs/VII/native_diagrams/exp7/viia_00_04/diag_036_relative_frobenius_factorization.tex

\end{center}
We call \(\operatorname{Fr}(X/S)\) the \emph{Frobenius morphism of
\(X\) relative to \(S\)}. Plainly the assignment
\[
  \operatorname{Fr}:X\longmapsto\operatorname{Fr}(X/S)
\]
is a functorial homomorphism.

\footnote{The editors expanded the original in the remainder of this
paragraph.}
Let \(r:T\to S\) be an \(S\)-scheme. For every
\[
  \varphi\in X(r)=\operatorname{Hom}_S(T,X),
\]
one has the commutative diagram
\begin{center}
  \phantomsection
  \label{sga3:VIIA:4.1-point-phi-diagram}
  % SGA3-VII-DIAG-037
  % BEGIN INCLUDED SGA3_English_current_buildable_source_20260823/source/inputs/VII/native_diagrams/exp7/viia_00_04/diag_037_point_phi_relative_frobenius.tex
{\small
\begin{tikzcd}[column sep=normal,row sep=large]
X
  \arrow[r,"\operatorname{Fr}(X/S)"]
  \arrow[dr,"q" description]
&
X^{(p/S)}
  \arrow[r,"\operatorname{pr}_X"]
  \arrow[d,"q^{(p/S)}"]
&
X
  \arrow[d,"q"]
\\
T
  \arrow[u,"\varphi"]
  \arrow[r,"r"]
&
S
  \arrow[r,"\operatorname{fr}(S)"]
&
S\,.
\end{tikzcd}}

% END INCLUDED SGA3_English_current_buildable_source_20260823/source/inputs/VII/native_diagrams/exp7/viia_00_04/diag_037_point_phi_relative_frobenius.tex

\end{center}
By the definition of \(X^{(p/S)}\) as a fiber product,
\(\operatorname{pr}_X\) induces a bijection
\[
 X^{(p/S)}(r)
 =
 \operatorname{Hom}_S(T,X^{(p/S)})
 \xrightarrow{\sim}
 \operatorname{Hom}_S(T,X)
 =
 X(\operatorname{fr}(S)\circ r).
\]
On the other hand,
\[
  r\circ\operatorname{fr}(T)
  =
  \operatorname{fr}(S)\circ r,
\]
because \(\operatorname{fr}\) is an endomorphism of the identity
functor. Hence
\[
  \operatorname{Fr}(X/S)(r):
  X(r)\longrightarrow X^{(p/S)}(r)
\]
is characterized by commutativity of
\begin{center}
  \phantomsection
  \label{sga3:VIIA:4.1-pointwise-frobenius-square}
  % SGA3-VII-DIAG-038
  % BEGIN INCLUDED SGA3_English_current_buildable_source_20260823/source/inputs/VII/native_diagrams/exp7/viia_00_04/diag_038_pointwise_frobenius_square.tex
\begin{tikzcd}[column sep=huge,row sep=large]
X(r)
  \arrow[r,"\operatorname{Fr}(X/S)(r)"]
  \arrow[d,"X(\operatorname{fr}(T))"']
&
X^{(p/S)}(r)
  \arrow[d,"\wr"]
\\
X(r\circ\operatorname{fr}(T))
  \arrow[r,equal]
&
X(\operatorname{fr}(S)\circ r)\,.
\end{tikzcd}

% END INCLUDED SGA3_English_current_buildable_source_20260823/source/inputs/VII/native_diagrams/exp7/viia_00_04/diag_038_pointwise_frobenius_square.tex

\end{center}

For example, if \(X\) is the subscheme of \(S\) defined by a
quasi-coherent ideal \(\mathcal I\), then \(X^{(p)}\) is the subscheme
defined by the ideal \(\mathcal I^{\{p\}}\) generated by the \(p\)-th
powers of local sections of \(\mathcal I\). Moreover,
\(\operatorname{Fr}(X/S)\) is the canonical immersion
\[
  \operatorname{Spec}(\mathcal O_S/\mathcal I)
  \longrightarrow
  \operatorname{Spec}(\mathcal O_S/\mathcal I^{\{p\}}).
\]
\footnote{The French source prints
\(\operatorname{Spec}(\mathcal O_X/\mathcal I)\) for the source of
this immersion. Since \(X\subset S\) is defined by \(\mathcal I\), the
intended source is
\(\operatorname{Spec}(\mathcal O_S/\mathcal I)=X\), as used here.}

\medskip
\noindent\textbf{4.1.1.}
\phantomsection
\label{sga3:VIIA:4.1.1}
\footnote{The editors expanded the original in this paragraph.}
Let \(t:T\to S\) be a base change and put
\[
  X_T=X\times_{q,t}T.
\]
Consider the pullback of \(X_T\) along \(\operatorname{fr}(T)\):
\begin{center}
  \phantomsection
  \label{sga3:VIIA:4.1.1-basechange-rectangle}
  % SGA3-VII-DIAG-039
  % BEGIN INCLUDED SGA3_English_current_buildable_source_20260823/source/inputs/VII/native_diagrams/exp7/viia_00_04/diag_039_basechange_rectangle.tex
\begin{tikzcd}[column sep=large,row sep=large]
(X_T)^{(p/T)}
  \arrow[r]
  \arrow[d]
&
X_T
  \arrow[r]
  \arrow[d]
&
X
  \arrow[d,"q"]
\\
T
  \arrow[r,"\operatorname{fr}(T)"]
&
T
  \arrow[r,"t"]
&
S\,.
\end{tikzcd}

% END INCLUDED SGA3_English_current_buildable_source_20260823/source/inputs/VII/native_diagrams/exp7/viia_00_04/diag_039_basechange_rectangle.tex

\end{center}
Since
\[
  t\circ\operatorname{fr}(T)
  =
  \operatorname{fr}(S)\circ t,
\]
\((X_T)^{(p/T)}\) identifies with the pullback of \(X^{(p/S)}\)
along \(t\). Equivalently, there is a canonical isomorphism
\[
  X_T^{(p/T)}
  \xrightarrow{\sim}
  \bigl(X^{(p/S)}\bigr)_T.
\]
Under this identification,
\(\operatorname{Fr}(X_T/T)\) identifies with the pullback
\(\operatorname{Fr}(X/S)_T\).

\medskip
\noindent\textbf{4.1.1.1.}
\phantomsection
\label{sga3:VIIA:4.1.1.1}
In particular, if \(S=\operatorname{Spec}\mathbb F_p\), then
\(X^{(p/S)}=X\) and
\(\operatorname{Fr}(X/S)=\operatorname{fr}(X)\). Consequently,
\(X_T^{(p/T)}\) identifies with \(X_T\), and
\(\operatorname{Fr}(X_T/T)\) with
\(\operatorname{fr}(X)_T\).

For example, if \(E\) is a set and \(E_T\) the constant \(T\)-scheme
of type \(E\), then
\[
  E_T^{(p/T)}\simeq E_T,
  \qquad
  \operatorname{Fr}(E_T/T)\simeq\operatorname{id}_{E_T}.
\]

\medskip
\noindent\textbf{4.1.2.}
\phantomsection
\label{sga3:VIIA:4.1.2}
The functor \(X\mapsto X^{(p/S)}\) plainly commutes with products.
It therefore sends an \(S\)-group \(G\) to an \(S\)-group
\(G^{(p/S)}\). Since \(\operatorname{Fr}\) is functorial,
\[
  \operatorname{Fr}(G/S):G\longrightarrow G^{(p/S)}
\]
is a homomorphism of \(S\)-groups. We denote its kernel by
\(\operatorname{Fr}G\).

If \(r:T\to S\) is an \(S\)-scheme, the
\hyperref[sga3:VIIA:4.1-pointwise-frobenius-square]{pointwise
Frobenius diagram} shows that the value
of \(\operatorname{Fr}G\) at \(r\) is the kernel of
\[
  G(\operatorname{fr}(T)):
  G(r)\longrightarrow G(r\circ\operatorname{fr}(T)).
\]

When \(T=I_R\) is the scheme of dual numbers over an \(S\)-scheme
\(R\), the absolute Frobenius of \(I_R\) factors as
\begin{center}
  \phantomsection
  \label{sga3:VIIA:4.1.2-dual-numbers-factorization}
  % SGA3-VII-DIAG-040
  % BEGIN INCLUDED SGA3_English_current_buildable_source_20260823/source/inputs/VII/native_diagrams/exp7/viia_00_04/diag_040_dual_numbers_frobenius_factorization.tex
\begin{tikzcd}[column sep=large]
I_R
  \arrow[r,"\mathrm{can.}"]
&
R
  \arrow[r,"\operatorname{fr}(R)"]
&
R
  \arrow[r,"s"]
&
I_R,
\end{tikzcd}

% END INCLUDED SGA3_English_current_buildable_source_20260823/source/inputs/VII/native_diagrams/exp7/viia_00_04/diag_040_dual_numbers_frobenius_factorization.tex

\end{center}
where \(s\) is the zero section. It follows that
\((\operatorname{Fr}G)(I_R)\) contains the kernel
\(\operatorname{Lie}(G/S)(R)\) of
\[
  G(s):G(I_R)\longrightarrow G(R),
\]
and hence
\[
  \operatorname{Lie}(G/S)
  =
  \operatorname{Lie}(\operatorname{Fr}G/S).
\]

% Next cursor: Exp7A local p. 25, combined-reader p. 478,
% re-edition printed p. 467, at 4.1.3.

% END INCLUDED SGA3_English_current_buildable_source_20260823/source/inputs/VII/tex/components/04_expose_VIIA_section4_frobeniuseries_40_through_412_en.tex

% BEGIN INCLUDED SGA3_English_current_buildable_source_20260823/source/inputs/VII/tex/components/05_expose_VIIA_par413_through_par414_en.tex
% SGA 3, Expose VII_A, English reconstruction -- Loop 1.
% Scope: authority local pages 25--26, paragraphs 4.1.3--4.1.4,
% stopping before 4.2 on local page 27.
% Combined-reader pages: 478--479. Re-edition printed pages: 467--468.
% Original-edition marginal locators within this span: 437--439.
% Authority: Polo--Gille Exp7A-13oct24.pdf,
% SHA-256 D46DB8C79D0385643B6EED4B96900613248375F9B73B4202478A4F57FD3E8522.
% Jacob Reinhold's CC BY 4.0 Markdown was consulted as comparison lineage.

\medskip
\noindent\textbf{4.1.3.}
\phantomsection
\label{sga3:VIIA:4.1.3}

% Authority: Exp7A local p. 25; combined p. 478; printed p. 467.
More generally, for every \(S\)-functor \(X\), define the
\(S\)-functor \(X^{(p^n)}\) recursively in \(n\) by
\[
  X^{(p)}=X^{(p/S)}
  \qquad\text{and}\qquad
  X^{(p^n)}
  =
  \bigl(X^{(p^{n-1})}\bigr)^{(p)}.
\]
Likewise, \(\operatorname{Fr}^n(X/S)\), or simply
\(\operatorname{Fr}^n\), denotes the composite functorial
homomorphism
\begin{center}
  \phantomsection
  \label{sga3:VIIA:4.1.3-iterated-relative-frobenius-chain}
  % SGA3-VII-DIAG-041
  % BEGIN INCLUDED SGA3_English_current_buildable_source_20260823/source/inputs/VII/native_diagrams/exp7/viia_05_14/diagram_041.tex
\(\displaystyle
X \xrightarrow{\operatorname{Fr}(X/S)} X^{(p)}
  \xrightarrow{\operatorname{Fr}(X^{(p)}/S)} X^{(p^2)}
  \longrightarrow \cdots \longrightarrow X^{(p^{n-1})}
  \xrightarrow{\operatorname{Fr}(X^{(p^{n-1})}/S)} X^{(p^n)} .
\)

% END INCLUDED SGA3_English_current_buildable_source_20260823/source/inputs/VII/native_diagrams/exp7/viia_05_14/diagram_041.tex

\end{center}
By \SGARefLink{sga3:VIIA:4.1.1}{4.1.1},
\(\operatorname{Fr}(X^{(p)}/S)\) coincides with
\(\operatorname{Fr}(X/S)^{(p)}\); in other words, the following
diagram commutes:
\begin{center}
  \phantomsection
  \label{sga3:VIIA:4.1.3-iterated-frobenius-naturality-square}
  % SGA3-VII-DIAG-042
  % BEGIN INCLUDED SGA3_English_current_buildable_source_20260823/source/inputs/VII/native_diagrams/exp7/viia_05_14/diagram_042.tex
\begin{tikzcd}[row sep=3.2em,column sep=4.2em]
X^{(p)} \arrow[r] \arrow[d,swap,"\operatorname{Fr}(X^{(p)}/S)"]
  & X \arrow[d,"\operatorname{Fr}(X/S)"] \\
X^{(p^2)} \arrow[r]
  & X^{(p)} .
\end{tikzcd}

% END INCLUDED SGA3_English_current_buildable_source_20260823/source/inputs/VII/native_diagrams/exp7/viia_05_14/diagram_042.tex

\end{center}

If \(G\) is a group \(S\)-functor, then \(G^{(p^n)}\) is again a
group \(S\)-functor and
\(\operatorname{Fr}^n(G/S)\) is a homomorphism of group
\(S\)-functors.

\medskip
\noindent\textbf{Definition.}
\phantomsection
\label{sga3:VIIA:4.1.3-height-definition}
\textit{We denote the kernel of \(\operatorname{Fr}^n(G/S)\) by
\(\operatorname{Fr}^nG\), and say that \(G\) has height
\(\leq n\) if \(\operatorname{Fr}^n(G/S)\) is zero, that is, if
\(\operatorname{Fr}^nG=G\).}

\medskip
\noindent\textbf{Lemma.}
\phantomsection
\label{sga3:VIIA:4.1.3-characteristic-kernel-lemma}
\textit{The group subfunctor \(\operatorname{Fr}^nG\) of \(G\) is
characteristic: for every \(S\)-scheme \(T\), every endomorphism
\(\varphi\) of the group \(T\)-functor \(G_T\) induces an
endomorphism of \((\operatorname{Fr}^nG)_T\).}

\begin{proof}
By \SGARefLink{sga3:VIIA:4.1.1}{4.1.1}, the constructions of \(G^{(p^n)}\) and
\(\operatorname{Fr}^n(G/S)\) commute with base change. We may
therefore suppose that \(T=S\). The assertion then follows from the
fact that \(\operatorname{Fr}^n(G/S)\) is a functorial
homomorphism.
\end{proof}

\medskip
\noindent\textbf{4.1.4. Examples.}
\phantomsection
\label{sga3:VIIA:4.1.4}

\smallskip
\noindent\textit{a)}
Let \(M\) first be an ``abstract'' abelian group, and let
\[
  G=D_S(M)
\]
be the diagonalizable group of type \(M\) (I, \SGARefLink{sga3:i:subsection:4-4}{4.4}). For every
\(S\)-scheme \(T\),
\[
  G(T)
  =
  \operatorname{Hom}_{(\mathbf{Ab})}
    \bigl(M,\Gamma(T,\mathcal O_T)^\times\bigr).
\]
Since \(G\) is the pullback of the diagonalizable group \(D(M)\)
over \(\mathbb F_p\), \(G^{(p)}\) identifies with \(G\), and
\(\operatorname{Fr}(G/S)(T)\) identifies with the endomorphism
\(x\mapsto x^p\) of \(G(T)\) (\SGARefLink{sga3:VIIA:4.1.1}{4.1.1}). In particular, when
\(M=\mathbb Z\), one has \(D_S(M)=\mathbb G_{m,S}\). Hence:

\begin{quote}
\(\operatorname{Fr}\mathbb G_{m,S}\) is the \(S\)-group
\(\boldsymbol\mu_{p,S}\) which associates to every \(S\)-scheme
\(T\) the group of \(p\)-th roots of unity in
\(\Gamma(T,\mathcal O_T)^\times\).
\end{quote}

\smallskip
\noindent\textit{b)}
Now let \(S\) be a scheme of characteristic \(p\), and let
\(\mathcal E\) be a sheaf of modules on \(S\). By I, \SGARefLink{sga3:i:proposition:4-6-2}{4.6.2} there is
a canonical isomorphism
\[
  \mathcal W(\mathcal E)^{(p)}
  \xrightarrow{\sim}
  \mathcal W(\mathcal E^{(p)}),
\]
where \(\mathcal E^{(p)}\) is the pullback of \(\mathcal E\) along
\(\operatorname{fr}(S)\). For every \(S\)-scheme
\(\pi:T\to S\), the map
\(\operatorname{Fr}(\mathcal W(\mathcal E)/S)(\pi)\) is determined,
by \SGARefLink{sga3:VIIA:4.1}{4.1}~\((\dagger)\), by the commutative triangle
\begin{center}
  \phantomsection
  \label{sga3:VIIA:4.1.4-vector-group-frobenius-triangle}
  % SGA3-VII-DIAG-043
  % BEGIN INCLUDED SGA3_English_current_buildable_source_20260823/source/inputs/VII/native_diagrams/exp7/viia_05_14/diagram_043.tex
\begin{tikzcd}[row sep=3.6em,column sep=2.6em]
\Gamma\!\left(T,\pi^*\operatorname{fr}(S)^*\mathcal E\right)
  \arrow[rr,"\sim","{\mathrm{can.}}"']
  && \Gamma\!\left(T,\operatorname{fr}(T)^*\pi^*\mathcal E\right) \\
& \Gamma(T,\pi^*\mathcal E)
  \arrow[ul,"\operatorname{Fr}(\mathcal W(\mathcal E)/S)(\pi)"]
  \arrow[ur,"f'"']
& ,
\end{tikzcd}

% END INCLUDED SGA3_English_current_buildable_source_20260823/source/inputs/VII/native_diagrams/exp7/viia_05_14/diagram_043.tex

\end{center}
where \(f'\) is the map induced by \(\operatorname{fr}(T)\).

% Authority: Exp7A local p. 26; combined p. 479; printed p. 468.
In particular, if \(\mathcal E=\mathcal O_S\), then
\(\mathcal W(\mathcal E)\) identifies with the additive group
\(\mathbb G_{a,S}\). In this case
\(\mathcal E^{(p)}=\mathcal E=\mathcal O_S\), and the Frobenius
morphism \(\operatorname{Fr}(\mathbb G_{a,S}/S)\) sends
\(x\in\Gamma(T,\mathcal O_T)\) to \(x^p\). Thus:

\begin{quote}
\(\operatorname{Fr}\mathbb G_{a,S}\) is the \(S\)-group
\(\boldsymbol\alpha_{p,S}\) which associates to every
\(S\)-scheme \(T\) the group
\[
  \bigl\{x\in\Gamma(T,\mathcal O_T)\mid x^p=0\bigr\}.
\]
\end{quote}

\smallskip
\noindent\textit{c)}
Similarly, for every quasi-coherent \(\mathcal O_S\)-algebra
\(\mathcal A\), the scheme
\((\operatorname{Spec}\mathcal A)^{(p)}\) identifies with
\(\operatorname{Spec}\mathcal A^{(p)}\), the spectrum of the
pullback of \(\mathcal A\) along \(\operatorname{fr}(S)\). If
\(\pi\) denotes the endomorphism \(x\mapsto x^p\) of the sheaf of
rings \(\mathcal O_S\), then
\[
  \mathcal A^{(p)}
  =
  \mathcal A\otimes_\pi\mathcal O_S,
  \tag{44}\label{sga3:vii:paragraph:4-1-4:equation:44}
\]
and \(\operatorname{Fr}((\operatorname{Spec}\mathcal A)/S)\) is
induced by the morphism of \(\mathcal O_S\)-algebras
\[
  \mathcal A\otimes_\pi\mathcal O_S
  \longrightarrow
  \mathcal A,
  \qquad
  a\otimes_\pi x\longmapsto a^p x.
\]
\footnote{\(\mathcal A\otimes_\pi\mathcal O_S\) denotes the
\(\mathcal O_S\)-algebra obtained by extension of scalars along
\(\pi:\mathcal O_S\to\mathcal O_S\); thus
\(ax\otimes_\pi1=a\otimes_\pi x^p\) and
\(x\cdot(a\otimes_\pi1)=a\otimes_\pi x\).}

Finally, for every quasi-coherent \(\mathcal O_S\)-module
\(\mathcal E\), there are canonical isomorphisms
\[
  \mathcal V(\mathcal E)^{(p)}
  \xrightarrow{\sim}
  \mathcal V(\mathcal E^{(p)})
  \qquad\text{and}\qquad
  \mathcal S(\mathcal E)^{(p)}
  \xrightarrow{\sim}
  \mathcal S(\mathcal E^{(p)}),
\]
where \(\mathcal S(\mathcal E)\) denotes the symmetric algebra of
the \(\mathcal O_S\)-module \(\mathcal E\).

\smallskip
\noindent\textit{d)}
Let \(\mathcal U\) be an \(\mathcal O_S\)-coalgebra (\SGARefLink{sga3:VIIB:3.1:lemma}{3.1}), and let
\(S\) be a scheme of characteristic \(p\). If
\(\mathcal U^{(p/S)}\), or \(\mathcal U^{(p)}\), denotes the
pullback of the coalgebra \(\mathcal U\) along
\(\operatorname{fr}(S)\), then, as in \textit{b)}, there is a
canonical isomorphism
\[
  \bigl(\operatorname{Spec}^{*}\mathcal U\bigr)^{(p)}
  \xrightarrow{\sim}
  \operatorname{Spec}^{*}\mathcal U^{(p)}.
\]
If \(\mathcal U\) is a group coalgebra, then for an \(S\)-scheme
\(T\), the value of
\(\operatorname{Fr}(\operatorname{Spec}^{*}\mathcal U)\), that is,
of the kernel of the Frobenius morphism
\[
  \operatorname{Spec}^{*}\mathcal U
  \longrightarrow
  \bigl(\operatorname{Spec}^{*}\mathcal U\bigr)^{(p)},
\]
is the set of elements \(\gamma\) of
\[
  \bigl(\operatorname{Spec}^{*}\mathcal U\bigr)(T)
  =
  \left\{
    x\in\Gamma(T,\mathcal U_T)
    \,\middle|\,
    \varepsilon_{\mathcal U_T}(x)=1,\ 
    \Delta_{\mathcal U_T}x=x\otimes x
  \right\}
\]
such that the image in
\(\Gamma(T,\mathcal U_T\otimes_{\operatorname{fr}(T)}\mathcal O_T)\)
of the element
\[
  \gamma\otimes_{\operatorname{fr}(T)}1
  \in
  \Gamma(T,\mathcal U_T)\otimes_{\operatorname{fr}(T)}
    \mathcal O(T)
\]
is equal to \(1\).

% Next cursor: Exp7A local p. 27, combined-reader p. 480,
% re-edition printed p. 469, at 4.2.

% END INCLUDED SGA3_English_current_buildable_source_20260823/source/inputs/VII/tex/components/05_expose_VIIA_par413_through_par414_en.tex

% BEGIN INCLUDED SGA3_English_current_buildable_source_20260823/source/inputs/VII/tex/components/06_expose_VIIA_section42_symmetric_frobenius_en.tex
% SGA 3, Expose VII_A, English reconstruction -- Loop 1.
% Scope: authority local pages 27--28, paragraphs 4.2--4.2.1,
% stopping before 4.3 on local page 29.
% Combined-reader pages: 480--481. Re-edition printed pages: 469--470.
% Original-edition marginal locators within this span: 440--441.
% Authority: Polo--Gille Exp7A-13oct24.pdf,
% SHA-256 D46DB8C79D0385643B6EED4B96900613248375F9B73B4202478A4F57FD3E8522.
% Jacob Reinhold's CC BY 4.0 Markdown was consulted as comparison lineage.

\medskip
\noindent\textbf{4.2.}
\phantomsection
\label{sga3:VIIA:4.2}
\footnote{For the material in \SGARefLink{sga3:VIIA:4.2}{4.2} and \SGARefLink{sga3:VIIA:4.3}{4.3}, see also
\textup{[DG70]}, IV, \S3, nos.~4--6.}

% Authority: Exp7A local p. 27; combined p. 480; printed p. 469.
We now consider a construction closely related to the preceding
one. Let \(S\) be a scheme of characteristic \(p\), let \(X\) be
an \(S\)-scheme, and let \(X_S^p\) be the product, in
\((\mathbf{Sch}/S)\), of \(p\) copies of \(X\).

Denote by \(U^p(X)\) the open subscheme of \(X_S^p\) which is the
union of the products \(U_S^p\) as \(U\) ranges over the affine
open subsets of \(X\). Thus a point \(x\) of \(X_S^p\) belongs to
\(U^p(X)\) if and only if its projections
\(\operatorname{pr}_i x\) on the factors of \(X_S^p\) belong to a
common affine open subset of \(X\). For example, if every finite
subset of \(X\) is contained in an affine open subset, then
\[
  U^p(X)=X_S^p.
\]

The symmetric group \(\mathfrak S_p\) on \(p\) letters acts on
\(X_S^p\) by permuting the factors and leaves \(U^p(X)\) stable.
We call the quotient of \(X_S^p\) by \(\mathfrak S_p\) in the
category of ringed spaces the \emph{\(p\)-fold symmetric product
of \(X\)}, and denote it by \(\Sigma^p X\). Let
\[
  q(X),\ \text{or simply }q:
  X_S^p\longrightarrow\Sigma^p X
\]
be the canonical projection.

The morphism \(q\) maps \(U^p(X)\) onto an open subset
\(V^p(X)\) of the symmetric product, which may be described as
follows (cf.\ V, \SGARefLink{thm:sga3-V-4.1}{4.1}). The structure sheaf of \(\Sigma^p X\)
induces a scheme structure on \(V^p(X)\), and the morphism
\[
  q'(X):U^p(X)\longrightarrow V^p(X)
\]
induced by \(q(X)\) is affine, indeed integral. As \(U\) ranges
over the affine open subsets of \(X\) which map into a varying
affine open \(V\) of \(S\), the schemes \(\Sigma^p U\) form an
affine open covering of \(V^p(X)\). If \(R\) is the affine algebra
of \(V\), and \(A\) that of \(U\), the affine algebra of
\(\Sigma^p U\) is the subalgebra \(\Sigma^p A\) of
\(\bigotimes_R^p A\) consisting of the symmetric tensors.

Now consider the diagonal morphism
\[
  \delta:X\longrightarrow U^p(X).
\]
If \(V=\operatorname{Spec}R\) is an affine open subset of \(S\)
and \(U=\operatorname{Spec}A\) is an affine open subset of \(X\)
lying over \(V\), the restriction of \(\delta\) to \(U\) is
defined by the algebra homomorphism
\[
  \eta:\bigotimes_R^p A\longrightarrow A,
  \qquad
  a_1\otimes\cdots\otimes a_p
  \longmapsto
  a_1a_2\cdots a_p.
\]
Consequently, if \(N\) denotes the symmetrization operator,
\[
\begin{aligned}
  \eta\bigl(N(a_1\otimes\cdots\otimes a_p)\bigr)
  &=
  \eta\left(
    \sum_{\sigma\in\mathfrak S_p}
      a_{\sigma(1)}\otimes\cdots\otimes a_{\sigma(p)}
  \right) \\
  &=
  p!\,a_1\cdots a_p
  =
  0.
\end{aligned}
\]
In other words, \(\eta\) vanishes on the subspace
\[
  N\!\left(\bigotimes_R^p A\right)
  \subseteq
  \Sigma^p A
\]
consisting of symmetrized tensors. Moreover, if \(f\) is a
symmetric tensor, then plainly \(N(fa)=fN(a)\). Hence
\(N(\bigotimes_R^p A)\) is an ideal of \(\Sigma^p A\). We put
\[
  U^{[p/S]}
  =
  \operatorname{Spec}\left(
    \Sigma^p A
    \big/
    N\!\left(\bigotimes_R^p A\right)
  \right).
\]
This is a closed subscheme of
\(\Sigma^p(U)=V^p(U)\). As \(U\) ranges over the affine open
subsets of \(X\) which map into a varying affine open \(V\) of
\(S\), the schemes \(U^{[p/S]}\) glue to a closed subscheme of
\(V^p(X)\), denoted \(X^{[p/S]}\).

If \(i(X)\) denotes the inclusion of \(X^{[p/S]}\) in \(V^p(X)\),
the preceding calculation shows that
\(q'(X)\circ\delta\) factors through \(X^{[p/S]}\). We therefore
obtain a morphism
\[
  F^{[p]}(X/S):X\longrightarrow X^{[p/S]}.
\]
\footnote{In the original, this morphism, respectively the relative
Frobenius morphism, was denoted \(F'\), respectively \(F\).}
The factorization is displayed by the commutative diagram
\begin{center}
  \phantomsection
  \label{sga3:VIIA:4.2-pfold-symmetric-factorization}
  % SGA3-VII-DIAG-044
  % BEGIN INCLUDED SGA3_English_current_buildable_source_20260823/source/inputs/VII/native_diagrams/exp7/viia_05_14/diagram_044.tex
\begin{tikzcd}[row sep=3.2em,column sep=1.8em]
X_S^p \arrow[dd,swap,"q(X)"]
  & \supset
  & \mathrm U^p(X) \arrow[dd,swap,"q'(X)"]
  && X \arrow[ll,swap,"\delta(X)"] \arrow[dd,"F^{[p]}(X/S)"] \\
\\[-1.2em]
\Sigma^p(X)
  & \supset
  & \mathrm V^p(X)
  && X^{[p/S]} . \arrow[ll,swap,"i(X)"]
\end{tikzcd}

% END INCLUDED SGA3_English_current_buildable_source_20260823/source/inputs/VII/native_diagrams/exp7/viia_05_14/diagram_044.tex

\end{center}
It is clear that \(X^{[p/S]}\) is functorial in \(X\), and that
\[
  F^{[p]}:X\longmapsto F^{[p]}(X/S)
\]
is a functorial homomorphism.

\medskip
\noindent\textbf{4.2.1.}
\phantomsection
\label{sga3:VIIA:4.2.1}

% Authority: Exp7A local p. 28; combined p. 481; printed p. 470.
The schemes \(X^{[p/S]}\) and \(X^{(p/S)}\) are naturally related.
Let \(V\) be an affine open subset of \(S\), with affine algebra
\(R\), and let \(U\) be an affine open subset of \(X\) lying over
\(V\), with affine algebra \(A\). If \(\pi\) denotes the
endomorphism \(x\mapsto x^p\) of \(R\), the affine algebra of
\(U^{(p/S)}\) is \(A\otimes_\pi R\). The map
\[
  a\otimes_\pi\lambda
  \longmapsto
  \left(
    \lambda a\otimes\cdots\otimes a
    \bmod
    N\!\left(\bigotimes_R^p A\right)
  \right)
\]
is a homomorphism of \(R\)-algebras
\[
  A\otimes_\pi R
  \longrightarrow
  \Sigma^p A\big/
    N\!\left(\bigotimes_R^p A\right).
\]
It induces a morphism
\[
  \varphi(U):U^{[p/S]}\longrightarrow U^{(p/S)}
\]
satisfying
\[
  \varphi(U)\circ F^{[p]}(U/S)=\operatorname{Fr}(U/S).
\]
Gluing these morphisms gives the commutative triangle
\begin{center}
  \phantomsection
  \label{sga3:VIIA:4.2.1-phi-frobenius-triangle}
  % SGA3-VII-DIAG-045
  % BEGIN INCLUDED SGA3_English_current_buildable_source_20260823/source/inputs/VII/native_diagrams/exp7/viia_05_14/diagram_045.tex
\begin{tikzcd}[row sep=3.6em,column sep=4.4em]
& X
  \arrow[dl,swap,"F^{[p]}(X/S)"]
  \arrow[dr,"\operatorname{Fr}(X/S)"] & \\
X^{[p/S]} \arrow[rr,"\varphi(X)"]
  && X^{(p/S)} .
\end{tikzcd}

% END INCLUDED SGA3_English_current_buildable_source_20260823/source/inputs/VII/native_diagrams/exp7/viia_05_14/diagram_045.tex

\end{center}

For example, if \(X\) is the subscheme of \(S\) defined by a
quasi-coherent ideal \(\mathcal I\), then
\(F^{[p]}(X/S)\) identifies with the identity morphism of \(X\).
Thus \(\varphi(X)\) is the canonical immersion
\[
  \operatorname{Spec}(\mathcal O_S/\mathcal I)
  \longrightarrow
  \operatorname{Spec}(\mathcal O_S/\mathcal I^{\{p\}}).
\]
This shows that \(\varphi(X)\) need not be an isomorphism.

On the other hand, when \(M\) is a free \(R\)-module, the map
\[
\begin{aligned}
  M\otimes_\pi R
  &\longrightarrow
  \Sigma^p M\big/
    N\!\left(\bigotimes_R^p M\right), \\
  m\otimes_\pi\lambda
  &\longmapsto
  \left(
    \lambda m\otimes\cdots\otimes m
    \bmod
    N\!\left(\bigotimes_R^p M\right)
  \right)
\end{aligned}
\]
is plainly bijective. It remains bijective when \(M\) is
\(R\)-flat, because every flat module is a filtered direct limit of
free modules.\footnote{D.~Lazard,
\emph{Sur les modules plats}, C.\ R.\ Acad.\ Sci.\ Paris
\textbf{258} (1964), 6313--6316. See also
D.~Lazard, Bull.\ Soc.\ Math.\ France \textbf{97} (1969), 81--128,
or \textup{[BAlg]}, X, \S1.6, Th.~1.}
It follows that
\[
  \boxed{\;
  \varphi(X):X^{[p/S]}\longrightarrow X^{(p/S)}
  \text{ is an isomorphism if \(X\) is flat over \(S\).}
  \;}
\]

% Next cursor: Exp7A local p. 29, combined-reader p. 482,
% re-edition printed p. 471, at 4.3.

% END INCLUDED SGA3_English_current_buildable_source_20260823/source/inputs/VII/tex/components/06_expose_VIIA_section42_symmetric_frobenius_en.tex

% BEGIN INCLUDED SGA3_English_current_buildable_source_20260823/source/inputs/VII/tex/components/07_expose_VIIA_section43_verschiebung_en.tex
% SGA 3, Expose VII_A, English reconstruction -- Loop 1.
% Scope: authority local pages 29--31, paragraphs 4.3--4.3.3,
% stopping before Section 5 on local page 31.
% Combined-reader pages: 482--484. Re-edition printed pages: 471--473.
% Original-edition marginal locators within this span: 442--444.
% Authority: Polo--Gille Exp7A-13oct24.pdf,
% SHA-256 D46DB8C79D0385643B6EED4B96900613248375F9B73B4202478A4F57FD3E8522.
% Jacob Reinhold's CC BY 4.0 Markdown was consulted as comparison lineage.

\medskip
\noindent\textbf{4.3.}
\phantomsection
\label{sga3:VIIA:4.3}

% Authority: Exp7A local p. 29; combined p. 482; printed p. 471.
Finally, let \(G\) be an \(S\)-group scheme in abelian groups. The
composite morphism
\[
  \mu(G):
  U^p(G)\lhook\joinrel\longrightarrow G_S^p
  \longrightarrow G,
\]
defined by multiplication, factors through \(V^p(G)\). Thus there
is a morphism
\[
  \nu(G):V^p(G)\longrightarrow G
\]
such that
\[
  \nu(G)\circ q'(G)=\mu(G).
\]
The resulting commutative diagram is
\begin{center}
  \phantomsection
  \label{sga3:VIIA:4.3-verschiebung-master-diagram}
  % SGA3-VII-DIAG-046
  % BEGIN INCLUDED SGA3_English_current_buildable_source_20260823/source/inputs/VII/native_diagrams/exp7/viia_05_14/diagram_046.tex
\begin{tikzpicture}[
  baseline=(current bounding box.center),
  >=Stealth,
  every node/.style={font=\normalsize}
]
  \node (gleft) at (-3.3,2.1) {$G$};
  \node (upg)   at (0,2.1) {$\mathrm U^p(G)$};
  \node (gright) at (3.3,2.1) {$G$};
  \node (vpg)   at (-1.65,0.35) {$\mathrm V^p(G)$};
  \node (gbr)   at (1.65,0.35) {$G^{[p/S]}$};
  \node (gps)   at (0,-1.65) {$G^{(p/S)}$};

  \draw[->] (upg) -- node[above] {$\mu(G)$} (gleft);
  \draw[->] (gright) -- node[above] {$\delta(G)$} (upg);
  \draw[->] (vpg) -- node[pos=.52,above,fill=white,inner sep=1pt]
    {$\nu(G)$} (gleft);
  \draw[->] (upg) -- node[pos=.53,right,fill=white,inner sep=1pt]
    {$q'(G)$} (vpg);
  \draw[->] (gbr) -- node[below] {$i(G)$} (vpg);
  \draw[->] (gright) -- node[pos=.5,left,fill=white,inner sep=1pt]
    {$F^{[p]}(G/S)$} (gbr);
  \draw[->] (gbr) -- node[pos=.55,anchor=west,xshift=2pt,fill=white,inner sep=1pt]
    {$\varphi(G)$} (gps);
  \draw[->] (gright) to[out=-92,in=18]
    node[right] {$\operatorname{Fr}(G/S)$} (gps);
  \draw[->,dotted] (gps) to[out=172,in=-87]
    node[left] {$\operatorname{Ver}(G/S)$} (gleft);
  \node at (3.55,-1.65) {$.$};
\end{tikzpicture}

% END INCLUDED SGA3_English_current_buildable_source_20260823/source/inputs/VII/native_diagrams/exp7/viia_05_14/diagram_046.tex

\end{center}

When \(G\) is flat over \(S\), \(\varphi(G)\) is an isomorphism,
and we may define the morphism called the \emph{Verschiebung}
\[
  \operatorname{Ver}(G/S):
  G^{(p/S)}
  \longrightarrow
  G
\]
by
\[
  \operatorname{Ver}(G/S)
  =
  \nu(G)\circ i(G)\circ\varphi(G)^{-1}.
\]
As \(G\) ranges over the flat commutative \(S\)-group schemes, the
assignment
\[
  \operatorname{Ver}:G\longmapsto\operatorname{Ver}(G/S)
\]
is plainly a functorial homomorphism. Consequently,
\(\operatorname{Ver}(G/S)\) is a group homomorphism.

For every \(S\)-scheme \(T\), the composite
\[
  G(T)
  \xrightarrow{\delta(G)(T)}
  U^p(G)(T)
  \xrightarrow{\mu(G)(T)}
  G(T)
\]
sends \(x\in G(T)\) to \(p\cdot x\). We may therefore write
\(p\cdot\operatorname{id}_G\) in place of
\(\mu(G)\circ\delta(G)\), obtaining the classical formula
\[
  \tag{*}\label{sga3:vii:paragraph:4-3:equation:star}
  \operatorname{Ver}(G/S)\circ\operatorname{Fr}(G/S)
  =
  p\cdot\operatorname{id}_G.
\]

\medskip
\noindent\textbf{Examples 4.3.1.}
\phantomsection
\label{sga3:VIIA:4.3.1}

\smallskip
\noindent\textit{a)}
If \(G\) is a constant \(S\)-scheme in abelian groups, then
\(\operatorname{Fr}(G/S)\) identifies with the identity morphism
of \(G\) (cf.\ \SGARefLink{sga3:VIIA:4.1.1.1}{4.1.1.1}). Hence
\[
  \operatorname{Ver}(G/S)
  =
  p\cdot\operatorname{id}_G.
\]

\smallskip
\noindent\textit{b)}
If \(G\) is the diagonalizable \(S\)-group of type \(M\), then
\(\operatorname{Fr}(G/S)=p\cdot\operatorname{id}_G\) by
\SGARefLink{sga3:VIIA:4.1.4}{4.1.4}~\textit{a)}. It follows at once that
\(\operatorname{Ver}(G/S)\) is the identity morphism of \(G\).

\smallskip
\noindent\textit{c)}
If \(\mathcal E\) is a flat \(\mathcal O_S\)-module and
\(G=\mathcal V(\mathcal E)\), then
\(\operatorname{Ver}(G/S)\) is zero, as is
\(p\cdot\operatorname{id}_G\). In Expos\'e VII\textsubscript{B} we
shall see that a commutative algebraic group \(G\) over a field
\(k\) is ``unipotent'' if and only if, for some \(n\), the
composite homomorphism
\begin{center}
  \phantomsection
  \label{sga3:VIIA:4.3.1-iterated-verschiebung-chain}
  % SGA3-VII-DIAG-047
  % BEGIN INCLUDED SGA3_English_current_buildable_source_20260823/source/inputs/VII/native_diagrams/exp7/viia_05_14/diagram_047.tex
\(\displaystyle
G^{(p^n)}
  \xrightarrow{\operatorname{Ver}(G^{(p^{n-1})}/S)}
G^{(p^{n-1})}
  \longrightarrow \cdots \longrightarrow
G^{(p)}
  \xrightarrow{\operatorname{Ver}(G/S)}
G
\)

% END INCLUDED SGA3_English_current_buildable_source_20260823/source/inputs/VII/native_diagrams/exp7/viia_05_14/diagram_047.tex

\end{center}
is zero. Here
\[
  G^{(p^n)}
  =
  \bigl(G^{(p^{n-1})}\bigr)^{(p)}
\]
as in \SGARefLink{sga3:VIIA:4.1.3}{4.1.3}.\footnote{This criterion is not stated explicitly in
Expos\'e VII\textsubscript{B}; see \textup{[DG70]}, IV, \S3,
Prop.~4.11.}

\medskip
\noindent\textbf{4.3.2.}
\phantomsection
\label{sga3:VIIA:4.3.2}

% Authority: Exp7A local p. 30; combined p. 483; printed p. 472.
Since
\(\operatorname{Ver}:G\mapsto\operatorname{Ver}(G/S)\) is a
functorial homomorphism on the category of flat commutative
\(S\)-group schemes, the square
\begin{center}
  \phantomsection
  \label{sga3:VIIA:4.3.2-frobenius-verschiebung-square}
  % SGA3-VII-DIAG-048
  % BEGIN INCLUDED SGA3_English_current_buildable_source_20260823/source/inputs/VII/native_diagrams/exp7/viia_05_14/diagram_048.tex
\begin{tikzcd}[row sep=3.3em,column sep=5.2em]
G^{(p)} \arrow[r,"\operatorname{Ver}(G/S)"]
  \arrow[d,swap,"\operatorname{Fr}(G/S)^{(p)}"]
  & G \arrow[d,"\operatorname{Fr}(G/S)"] \\
G^{(p^2)} \arrow[r,"\operatorname{Ver}(G^{(p)}/S)"]
  & G^{(p)}
\end{tikzcd}

% END INCLUDED SGA3_English_current_buildable_source_20260823/source/inputs/VII/native_diagrams/exp7/viia_05_14/diagram_048.tex

\end{center}
commutes. Here \(\operatorname{Fr}(G/S)^{(p)}\) denotes the
pullback of \(\operatorname{Fr}(G/S)\) along the base change
\(\operatorname{fr}(S)\). By \SGARefLink{sga3:VIIA:4.1.1}{4.1.1},
\[
  \operatorname{Fr}(G/S)^{(p)}
  =
  \operatorname{Fr}(G^{(p)}/S).
\]
Applying formula \((*)\) of \SGARefLink{sga3:VIIA:4.3}{4.3} to \(G^{(p)}\), we obtain
\[
  \tag{**}\label{sga3:vii:paragraph:4-3-2:equation:double-star}
  \operatorname{Fr}(G/S)\circ\operatorname{Ver}(G/S)
  =
  \operatorname{Ver}(G^{(p)}/S)
    \circ\operatorname{Fr}(G^{(p)}/S)
  =
  p\cdot\operatorname{id}_{G^{(p)}}.
\]

\medskip
\noindent\textbf{4.3.3.}
\phantomsection
\label{sga3:VIIA:4.3.3}

Suppose finally that \(G\) is a commutative \(S\)-group which is
finite and locally free. Let \(\mathcal A\) be its affine
\(\mathcal O_S\)-algebra, and let \(\pi\) be the endomorphism of
the sheaf of rings \(\mathcal O_S\) which sends a section \(x\) to
\(x^p\).\footnote{The editors changed the order of presentation,
introducing first the objects which occur in the diagram below.}
Denote by \(\Sigma^p\mathcal A\) the subalgebra of
\(\bigotimes_{\mathcal O_S}^p\mathcal A\) consisting of the
sections invariant under the symmetric-group action, and by
\[
  i(\mathcal A):
  \Sigma^p\mathcal A
  \longrightarrow
  \bigotimes_{\mathcal O_S}^p\mathcal A
\]
its inclusion.

Let
\[
  \Delta^p(\mathcal A):
  \mathcal A
  \longrightarrow
  \bigotimes_{\mathcal O_S}^p\mathcal A
\]
be the morphism obtained by iterating the diagonal morphism of the
coalgebra \(\mathcal A\). It corresponds to the multiplication
morphism \(U^p(G)=G_S^p\to G\). By the beginning of \SGARefLink{sga3:VIIA:4.3}{4.3},
\(\Delta^p(\mathcal A)\) factors through
\(\Sigma^p\mathcal A\); that is, it induces a morphism
\[
  a(\mathcal A):
  \mathcal A\longrightarrow\Sigma^p\mathcal A
\]
such that
\[
  i(\mathcal A)\circ a(\mathcal A)
  =
  \Delta^p(\mathcal A).
\]

On the other hand, let \(S^p(\mathcal A)\) be the homogeneous
component of degree \(p\) in the symmetric algebra of
\(\mathcal A\), and let
\[
  q(\mathcal A):
  \bigotimes_{\mathcal O_S}^p\mathcal A
  \longrightarrow
  S^p(\mathcal A)
\]
be the canonical projection. The multiplication
\[
  m^p(\mathcal A):
  \bigotimes_{\mathcal O_S}^p\mathcal A
  \longrightarrow
  \mathcal A
\]
factors through \(S^p(\mathcal A)\); equivalently, it induces a map
\[
  b(\mathcal A):
  S^p(\mathcal A)\longrightarrow\mathcal A
\]
such that
\[
  b(\mathcal A)\circ q(\mathcal A)
  =
  m^p(\mathcal A).
\]

Since \(\Sigma^p\mathcal A\) is the affine algebra of
\(V^p(G)\), the composite
\(i(G)\circ\varphi(G)^{-1}\) induces, again by the beginning of
\SGARefLink{sga3:VIIA:4.3}{4.3}, an algebra homomorphism
\[
  r(\mathcal A):
  \Sigma^p\mathcal A
  \longrightarrow
  \mathcal A\otimes_\pi\mathcal O_S.
\]
It vanishes on sections of the form
\[
  \sum_{\sigma\in\mathfrak S_p}
    a_{\sigma(1)}\otimes\cdots\otimes a_{\sigma(p)}
\]
and sends \(a\otimes\cdots\otimes a\) to
\(a\otimes_\pi1\). Similarly, \(j(\mathcal A)\) is the morphism
of \(\mathcal O_S\)-modules
\[
  a\otimes_\pi1
  \longmapsto
  q(a\otimes\cdots\otimes a).
\]

% Authority: Exp7A local p. 31; combined p. 484; printed p. 473.
These maps form the commutative diagram
\begin{center}
  \phantomsection
  \label{sga3:VIIA:4.3.3-cartier-duality-algebra-diagram}
  % SGA3-VII-DIAG-049
  % BEGIN INCLUDED SGA3_English_current_buildable_source_20260823/source/inputs/VII/native_diagrams/exp7/viia_05_14/diagram_049.tex
\begin{tikzpicture}[
  baseline=(current bounding box.center),
  >=Stealth,
  every node/.style={font=\normalsize}
]
  \node (atl) at (-3.2,2.0) {$\mathcal A$};
  \node (atm) at (0,2.0) {$\displaystyle\bigotimes_{\mathcal O_S}^{p}\mathcal A$};
  \node (atr) at (3.2,2.0) {$\mathcal A$};
  \node (aml) at (-1.65,0.15) {$\Sigma^p\mathcal A$};
  \node (amr) at (1.65,0.15) {$\mathrm S^p(\mathcal A)$};
  \node (ab)  at (0,-1.75) {$\mathcal A\otimes_{\pi}\mathcal O_S$};
  \node at (-4.85,0.15) {$({\mathcal A})$};

  \draw[->] (atl) -- node[above] {$\Delta^p(\mathcal A)$} (atm);
  \draw[->] (atm) -- node[above] {$m^p(\mathcal A)$} (atr);
  \draw[->] (atl) -- node[midway,left,fill=white,inner sep=1pt]
    {$a(\mathcal A)$} (aml);
  \draw[->] (aml) -- node[midway,left,fill=white,inner sep=1pt]
    {$i(\mathcal A)$} (atm);
  \draw[->] (atm) -- node[midway,right,fill=white,inner sep=1pt]
    {$q(\mathcal A)$} (amr);
  \draw[->] (amr) -- node[midway,right,fill=white,inner sep=1pt]
    {$b(\mathcal A)$} (atr);
  \draw[->] (aml) -- node[midway,left,fill=white,inner sep=1pt]
    {$r(\mathcal A)$} (ab);
  \draw[->] (ab) -- node[midway,right,fill=white,inner sep=1pt]
    {$j(\mathcal A)$} (amr);
  \node at (3.55,-1.75) {$.$};
\end{tikzpicture}

% END INCLUDED SGA3_English_current_buildable_source_20260823/source/inputs/VII/native_diagrams/exp7/viia_05_14/diagram_049.tex

\end{center}
The composite \(r(\mathcal A)\circ a(\mathcal A)\) corresponds to
the Verschiebung morphism \(\operatorname{Ver}(G/S)\), while
\(b(\mathcal A)\circ j(\mathcal A)\) corresponds to the Frobenius
morphism \(\operatorname{Fr}(G/S)\).

The preceding commutative diagram \((\mathcal A)\) is self-dual.
Indeed, let \(D\) be the functor which associates to each
\(\mathcal O_S\)-module \(\mathcal M\) its dual
\[
  D\mathcal M
  =
  \operatorname{Hom}_{\mathcal O_S}
    (\mathcal M,\mathcal O_S).
\]
The image of diagram \((\mathcal A)\) under \(D\) is precisely
diagram \((D\mathcal A)\). The morphisms
\[
  Dr(\mathcal A),\quad
  Da(\mathcal A),\quad
  Dj(\mathcal A),\quad
  Db(\mathcal A)
\]
identify respectively with
\[
  j(D\mathcal A),\quad
  b(D\mathcal A),\quad
  r(D\mathcal A),\quad
  a(D\mathcal A).
\]
By \SGARefLink{sga3:VIIA:3.3.1}{3.3.1}, it follows that:

\begin{quote}
\textit{In the category of commutative \(S\)-groups which are
finite and locally free, Cartier duality interchanges the
Frobenius morphism and the Verschiebung.}
\end{quote}
\footnote{See also \textup{[DG70]}, IV, \S3, 4.9.}

% Next cursor: Exp7A local p. 31, combined-reader p. 484,
% re-edition printed p. 473, at Section 5.

% END INCLUDED SGA3_English_current_buildable_source_20260823/source/inputs/VII/tex/components/07_expose_VIIA_section43_verschiebung_en.tex

% BEGIN INCLUDED SGA3_English_current_buildable_source_20260823/source/inputs/VII/tex/components/08_expose_VIIA_section5_restricted_p_lie_51_52_en.tex
% SGA 3, Expose VII_A, English reconstruction -- Loop 1.
% Scope: authority local pages 31--33, Section 5 opening through 5.2.1,
% stopping before 5.3 on local page 33.
% Combined-reader pages: 484--486. Re-edition printed pages: 473--475.
% Original-edition marginal locators within this span: 445--448.
% Authority: Polo--Gille Exp7A-13oct24.pdf,
% SHA-256 D46DB8C79D0385643B6EED4B96900613248375F9B73B4202478A4F57FD3E8522.
% Jacob Reinhold's CC BY 4.0 Markdown was consulted as comparison lineage.

\section*{5. Restricted \texorpdfstring{\(p\)}{p}-Lie Algebras}
\phantomsection
\label{sga3:VIIA:5}
\addcontentsline{toc}{section}{5. Restricted p-Lie Algebras}

We first recall some results from the \emph{S\'eminaire Sophus
Lie}.\footnote{See P.~Cartier,
\emph{Exemples d'hyperalg\`ebres}, S\'eminaire Sophus Lie
1955/56, Exp.~3.}

\medskip
\noindent\textbf{5.1.}
\phantomsection
\label{sga3:VIIA:5.1}

% Authority: Exp7A local p. 31; combined p. 484; printed p. 473.
Let \(p\) be a prime number, let \(R\) be a commutative ring of
characteristic \(p\), and let \(A\) be an associative
\(R\)-algebra, not necessarily commutative. For \(a,b\in A\), put
\[
  [a,b]=ab-ba,
  \qquad
  a\cdot b=L_a(b)=R_b(a).
\]
Then
\[
\begin{aligned}
  (\operatorname{ad}x^p)(y)
  &=[x^p,y] \\
  &=(L_x^p-R_x^p)(y) \\
  &=(L_x-R_x)^p(y) \\
  &=(\operatorname{ad}x)^p(y),
\end{aligned}
\]
which gives Jacobson's first formula:
\[
  \tag{i}\label{sga3:vii:paragraph:5-1:equation:i}
  \operatorname{ad}(x^p)
  =
  (\operatorname{ad}x)^p.
\]

If \(a_1,\ldots,a_p\) are arbitrary elements of \(A\), and \(N\)
is the symmetrization operator (cf.\ \SGARefLink{sga3:VIIA:4.2}{4.2}), then
\begin{equation}
  \tag{*}\label{sga3:vii:paragraph:5-1:equation:star}
\begin{aligned}
  N(a_1\otimes\cdots\otimes a_p)
  &=
  \sum_{\sigma\in\mathfrak S_p}
    a_{\sigma(1)}\cdots a_{\sigma(p)} \\
  &=
  \sum_{\tau\in\mathfrak S_{p-1}}
  [a_{\tau(1)},
    [a_{\tau(2)},
      [\cdots,
        [a_{\tau(p-1)},a_p]
      \cdots]
    ]
  ].
\end{aligned}
\end{equation}
Here \(\sigma\) ranges over the permutations of \(p\) letters and
\(\tau\) over those of \(p-1\) letters.

% Authority: Exp7A local p. 32; combined p. 485; printed p. 474.
Indeed, the final expression is equal to
\[
  \sum_{\tau}
  \sum_{r=0}^{p-1}
  \sum_{i_1<\cdots<i_r}
  (-1)^s
  a_{\tau(i_1)}\cdots a_{\tau(i_r)}
  a_p
  a_{\tau(j_s)}\cdots a_{\tau(j_1)},
\]
where \(\tau\) ranges over the permutations of \(p-1\) letters;
\(i_1,\ldots,i_r\) ranges over the strictly increasing sequences
in \([1,p-1]\); and \(j_1,\ldots,j_s\) is the strictly increasing
sequence of those integers of \([1,p-1]\) which are different
from \(i_1,\ldots,i_r\). For fixed \(r\), the sum of the terms
\[
  (-1)^s
  a_{\tau(i_1)}\cdots a_{\tau(i_r)}
  a_p
  a_{\tau(j_s)}\cdots a_{\tau(j_1)}
\]
is plainly
\[
  (-1)^s\binom{p-1}{s}
  \sum_{\rho}
  a_{\rho(1)}\cdots a_{\rho(r)}
  a_p
  a_{\rho(r+1)}\cdots a_{\rho(p-1)},
\]
where \(\rho\) ranges over the permutations of \(p-1\) letters.
But
\[
  (-1)^s\binom{p-1}{s}=1
  \quad\text{in }\mathbb F_p,
\]
because in \(\mathbb F_p[x]\),
\[
  (x-1)^p=x^p-1=(x-1)(x^{p-1}+\cdots+1),
\]
and hence
\[
  (x-1)^{p-1}=x^{p-1}+\cdots+1.
\]
This proves \((*)\).

On the other hand, if \(x_0,x_1\in A\), then
\[
  (x_0+x_1)^p
  =
  x_0^p+x_1^p+
  \sum_z x_{z(1)}x_{z(2)}\cdots x_{z(p)},
\]
where \(z\) ranges over the nonconstant maps
\([1,p]\to\{0,1\}\). It follows that
\[
\begin{aligned}
  (x_0+x_1)^p
  ={}&
  x_0^p+x_1^p \\
  &+
  \sum_{0<r<p}
  \frac{1}{r!(p-r)!}\,
  N(
    \underbrace{x_0,\ldots,x_0}_{r},
    \underbrace{x_1,\ldots,x_1}_{p-r}
  ).
\end{aligned}
\]
\footnote{The editors inserted the following explanation from
\textup{[DG70]}, III.7, \SGARefLink{sga3:VIIB:3.2:prop}{3.2}.}
By \((*)\),
\[
\begin{aligned}
  &N(
    \underbrace{x_0,\ldots,x_0}_{r},
    \underbrace{x_1,\ldots,x_1}_{p-r}
  ) \\
  &\qquad =
  r!(p-1-r)!
  \sum_t
  [x_{t(1)},
    [x_{t(2)},
      [\cdots,
        [x_{t(p-1)},x_1]
      \cdots]
    ]
  ],
\end{aligned}
\]
where \(t\) ranges over the maps
\([1,p-1]\to\{0,1\}\) which take the value \(0\) exactly \(r\)
times. We obtain Jacobson's second formula:
\begin{equation}
  \tag{ii}\label{sga3:vii:paragraph:5-1:equation:ii}
\begin{aligned}
  (x_0+x_1)^p
  ={}&
  x_0^p+x_1^p \\
  &-
  \sum_{0<r<p}\sum_t\frac1r
  [x_{t(1)},
    [x_{t(2)},
      [\cdots,
        [x_{t(p-1)},x_1]
      \cdots]
    ]
  ],
\end{aligned}
\end{equation}
where \(t\) has the same meaning.

\medskip
\noindent\textbf{5.2.}
\phantomsection
\label{sga3:VIIA:5.2}

Let now \(\mathfrak g\) be an \(R\)-Lie algebra. A map
\[
  x\longmapsto x^{(p)}
\]
from \(\mathfrak g\) to itself is said to make \(\mathfrak g\) a
\emph{restricted \(p\)-Lie algebra over \(R\)} if the following
conditions hold:
\[
\begin{array}{ll}
  (0) &
  (\lambda x)^{(p)}
  =
  \lambda^p x^{(p)}
  \quad
  (\lambda\in R,\ x\in\mathfrak g), \\[0.4em]
  (i) &
  \operatorname{ad}x^{(p)}
  =
  (\operatorname{ad}x)^p
  \quad
  (x\in\mathfrak g),
\end{array}
\]
and
\[
\begin{aligned}
  (ii)\quad
  (x_0+x_1)^{(p)}
  ={}&
  x_0^{(p)}+x_1^{(p)} \\
  &-
  \sum_{0<r<p}\sum_t\frac1r
  [x_{t(1)},
    [x_{t(2)},
      [\cdots,
        [x_{t(p-1)},x_1]
      \cdots]
    ]
  ],
\end{aligned}
\]
where \(t\) ranges over the maps
\([1,p-1]\to\{0,1\}\) which take the value \(0\) exactly \(r\)
times, and \(x_0,x_1\in\mathfrak g\). The map
\(x\mapsto x^{(p)}\) is called the \emph{symbolic \(p\)-th power}.

% Authority: Exp7A local p. 33; combined p. 486; printed p. 475.
For example, if \(A\) is an associative \(R\)-algebra, \SGARefLink{sga3:VIIB:5.1:proposition}{5.1} shows
that the \(R\)-module underlying \(A\), endowed with
\[
  [x,y]=xy-yx,
  \qquad
  x^{(p)}=x^p,
\]
is a \(p\)-Lie algebra, denoted \(A_{\mathrm{Lie}}\). We call it
the \(p\)-Lie algebra underlying \(A\).

We shall chiefly consider restricted \(p\)-Lie subalgebras of
algebras of the form \(A_{\mathrm{Lie}}\). Here is one example.
Let \(S\) be a scheme of characteristic \(p>0\), and let \(X\) be
an \(S\)-scheme. Recall that a derivation of \(X\) over \(S\) is
an endomorphism \(D\) of the sheaf of abelian groups
\(\mathcal O_X\) such that
\[
  D(\lambda s)=\lambda D(s),
  \qquad
  D(st)=(Ds)t+s(Dt),
\]
whenever \(\lambda\), respectively \(s,t\), are sections of
\(\mathcal O_S\), respectively \(\mathcal O_X\), on open subsets
for which the formulas make sense. Leibniz's formula
\[
  D^n(st)
  =
  \sum_{i=0}^n
  \binom ni
  (D^is)(D^{n-i}t)
\]
shows that \(D^p\) is again a derivation of \(X\) over \(S\), since
\[
  \binom pi\equiv0\pmod p
  \qquad(i\ne0,p).
\]
It follows that:

\begin{quote}
\textit{The algebra \(\operatorname{Der}_{X/S}\) of derivations of
\(X\) over \(S\) is a restricted \(p\)-Lie subalgebra of the
\(\Gamma(S,\mathcal O_S)\)-algebra of differential operators of
\(X\) over \(S\).}
\end{quote}

\medskip
\noindent\textbf{5.2.1.}
\phantomsection
\label{sga3:VIIA:5.2.1}
If \(\mathfrak g\) and \(\mathfrak h\) are two restricted
\(p\)-Lie algebras, a \emph{homomorphism}
\[
  h:\mathfrak g\longrightarrow\mathfrak h
\]
is an \(R\)-linear map such that, for \(x,y\in\mathfrak g\),
\[
  h([x,y])=[h(x),h(y)],
  \qquad
  h(x^{(p)})=h(x)^{(p)}.
\]
The composite of two homomorphisms is again a homomorphism, so we
may speak of the \emph{category of restricted \(p\)-Lie algebras
over \(R\)}.

If \((X,\mathcal R)\) is a ringed space, we say that an
\(\mathcal R\)-module \(\mathfrak g\) is endowed with the
structure of a restricted \(p\)-Lie algebra over \(\mathcal R\) if,
for every open subset \(U\), the module
\(\Gamma(U,\mathfrak g)\) is a restricted \(p\)-Lie algebra over
\(\Gamma(U,\mathcal R)\), and if all restriction maps are
homomorphisms.

% Next cursor: Exp7A local p. 33, combined-reader p. 486,
% re-edition printed p. 475, at 5.3.

% END INCLUDED SGA3_English_current_buildable_source_20260823/source/inputs/VII/tex/components/08_expose_VIIA_section5_restricted_p_lie_51_52_en.tex

% BEGIN INCLUDED SGA3_English_current_buildable_source_20260823/source/inputs/VII/tex/components/09_expose_VIIA_section53_restricted_enveloping_en.tex
% SGA 3, Expose VII_A, English reconstruction -- Loop 1.
% Scope: authority local pages 33--37, paragraphs 5.3--5.3.4,
% stopping before 5.4 on local page 37.
% Combined-reader pages: 486--490. Re-edition printed pages: 475--479.
% Original-edition marginal locators within this span: 448--452.
% Authority: Polo--Gille Exp7A-13oct24.pdf,
% SHA-256 D46DB8C79D0385643B6EED4B96900613248375F9B73B4202478A4F57FD3E8522.
% Jacob Reinhold's CC BY 4.0 Markdown was consulted as comparison lineage.

\medskip
\noindent\textbf{5.3.}
\phantomsection
\label{sga3:VIIA:5.3}

% Authority: Exp7A local p. 33; combined p. 486; printed p. 475.
We are now interested in the left adjoint of the functor
\[
  A\longmapsto A_{\mathrm{Lie}}
\]
of \SGARefLink{sga3:VIIB:5.2:theorem}{5.2}. Let \(\mathfrak g\) be a restricted \(p\)-Lie algebra
over a ring \(R\) of characteristic \(p\), let
\(U(\mathfrak g)\) be the enveloping algebra of the Lie algebra
underlying \(\mathfrak g\) (cf.\ \textup{[BLie]}, I, \S2.1), and
let
\[
  i_{\mathfrak g}:\mathfrak g\longrightarrow U(\mathfrak g)
\]
(or simply \(i\)) be the canonical map.

Let \(A\) be a unital associative \(R\)-algebra. For every Lie
algebra homomorphism
\[
  \varphi:\mathfrak g\longrightarrow A_{\mathrm{Lie}},
\]
there is a unique homomorphism of unital \(R\)-algebras
\[
  \psi:U(\mathfrak g)\longrightarrow A
\]
such that \(\psi\circ i=\varphi\). Moreover, \(\varphi\) is a
homomorphism of restricted \(p\)-Lie algebras if and only if
\(\psi\) vanishes on all the elements
\[
  i(x)^p-i\bigl(x^{(p)}\bigr),
  \qquad x\in\mathfrak g.
\]

\noindent\textbf{Definition.}
The quotient of \(U(\mathfrak g)\) by the two-sided ideal
generated by the elements
\[
  i(x)^p-i\bigl(x^{(p)}\bigr)
\]
is denoted \(U_p^R(\mathfrak g)\), or simply
\(U_p(\mathfrak g)\). We denote by \(j_{\mathfrak g}\), or simply
\(j\), the composite
\[
  \mathfrak g
  \xrightarrow{i}
  U(\mathfrak g)
  \longrightarrow
  U_p(\mathfrak g).
\]
The algebra \(U_p(\mathfrak g)\) is called the
\emph{restricted enveloping algebra} of \(\mathfrak g\).

The preceding discussion gives the following result.

\medskip
\noindent\textbf{Proposition.}
\textit{For every unital associative \(R\)-algebra \(A\) and
every homomorphism of restricted \(p\)-Lie algebras
\[
  \varphi:\mathfrak g\longrightarrow A_{\mathrm{Lie}},
\]
there is a unique homomorphism of unital algebras
\[
  \psi:U_p(\mathfrak g)\longrightarrow A
\]
such that \(\psi\circ j=\varphi\). In other words, the functor
\(\mathfrak g\mapsto U_p(\mathfrak g)\) is left adjoint to the
forgetful functor \(A\mapsto A_{\mathrm{Lie}}\).}

\medskip
\noindent\textbf{5.3.1.}
\phantomsection
\label{sga3:VIIA:5.3.1}

% Authority: Exp7A local p. 34; combined p. 487; printed p. 476.
With the notation of \SGARefLink{sga3:VIIA:5.3}{5.3}, put
\[
  \beta(x)=i(x)^p-i\bigl(x^{(p)}\bigr).
\]
For every \(y\in\mathfrak g\), \SGARefLink{sga3:VIIA:5.1}{formulas 5.1}~\SGARefLink{sga3:vii:paragraph:5-1:equation:i}{(i)} and \SGARefLink{sga3:VIIB:5.2:theorem}{5.2}~\SGARefLink{sga3:vii:paragraph:5-1:equation:i}{(i)}
give
\[
\begin{aligned}
  \beta(x)i(y)
  &=
  i(y)\beta(x)+[\beta(x),i(y)]\\
  &=
  i(y)\beta(x)
  +(\operatorname{ad}i(x))^p i(y)
  -i\bigl((\operatorname{ad}x)^p y\bigr)\\
  &=
  i(y)\beta(x).
\end{aligned}
\]
Thus \(\beta(x)\) belongs to the center of \(U(\mathfrak g)\);
in particular, the left ideal generated by the elements
\(\beta(x)\) is already two-sided.

On the other hand,
\[
  \beta(\lambda x)=\lambda^p\beta(x)
  \qquad(\lambda\in R),
\]
and \SGARefLink{sga3:VIIA:5.1}{formulas 5.1}~\SGARefLink{sga3:vii:paragraph:5-1:equation:ii}{(ii)} and \SGARefLink{sga3:VIIA:5.2}{5.2}~\SGARefLink{sga3:vii:paragraph:5-1:equation:ii}{(ii)} give
\[
  \beta(x+y)=\beta(x)+\beta(y)
  \qquad(x,y\in\mathfrak g).
\]
Consequently, if \((x_\alpha)\) is a family of generators of the
\(R\)-module \(\mathfrak g\), the left ideal generated by all
\(\beta(x)\) is already generated by the elements
\(\beta(x_\alpha)\).

\medskip
\noindent\textbf{5.3.2. Proposition.}\footnote{In this paragraph,
the editors changed the order, stating the result first and then
giving the detailed proof.}
\phantomsection
\label{sga3:VIIA:5.3.2}
\textit{Let \(\mathfrak g\) be an \(R\)-Lie algebra whose
underlying \(R\)-module is free with basis \((x_\alpha)\). Then
the restricted \(p\)-Lie algebra structures on \(\mathfrak g\)
correspond bijectively to the families \((y_\alpha)\) of elements
of \(\mathfrak g\) such that
\[
  \operatorname{ad}y_\alpha
  =
  (\operatorname{ad}x_\alpha)^p.
\]}

Indeed, suppose that \(\mathfrak g\) is endowed with a restricted
\(p\)-Lie algebra structure \(x\mapsto x^{(p)}\). By \SGARefLink{sga3:VIIB:5.2:theorem}{5.2}~\SGARefLink{sga3:vii:paragraph:5-1:equation:i}{(i)},
\SGARefLink{sga3:VIIB:0}{(0)}, and \SGARefLink{sga3:vii:paragraph:5-1:equation:ii}{(ii)}, the elements
\[
  y_\alpha=x_\alpha^{(p)}
\]
satisfy
\[
  \operatorname{ad}y_\alpha
  =
  (\operatorname{ad}x_\alpha)^p,
\]
and they determine the restricted \(p\)-Lie algebra structure.

We prove the converse. Since \(\mathfrak g\) is a free
\(R\)-module, the canonical map
\[
  i:\mathfrak g\longrightarrow U(\mathfrak g)
\]
is injective by the Poincar\'e--Birkhoff--Witt theorem
(cf.\ \textup{[BLie]}, I, \S2.7). We may therefore identify
\(\mathfrak g\) with an \(R\)-submodule of \(U(\mathfrak g)\).
Suppose that \((y_\alpha)\) is a family of elements of
\(\mathfrak g\) such that
\[
  \operatorname{ad}y_\alpha
  =
  (\operatorname{ad}x_\alpha)^p.
\]
Let
\[
  \pi:R\longrightarrow R,
  \qquad
  r\longmapsto r^p,
\]
and let \(\mathfrak g\otimes_\pi R\) be the \(R\)-Lie algebra
obtained by extension of scalars along \(\pi\).\footnote{Thus
\(xr\otimes_\pi1=x\otimes_\pi r^p\) and
\(r\cdot(x\otimes_\pi1)=x\otimes_\pi r\), for
\(x\in\mathfrak g\) and \(r\in R\).}

There is an \(R\)-linear map
\[
  \gamma:\mathfrak g\otimes_\pi R
  \longrightarrow
  U(\mathfrak g)
\]
which sends \(x_\alpha\otimes_\pi1\) to
\(x_\alpha^p-y_\alpha\). Moreover, for every
\(x\in\mathfrak g\),
\[
  (\operatorname{ad}x_\alpha^p)(x)
  =
  (\operatorname{ad}x_\alpha)^p(x)
  =
  (\operatorname{ad}y_\alpha)(x).
\]
It follows that \(\gamma\) maps
\(\mathfrak g\otimes_\pi R\) into the center of
\(U(\mathfrak g)\). For \(x\in\mathfrak g\), put
\[
  x^{(p)}
  =
  x^p-\gamma(x\otimes_\pi1).
\]

% Authority: Exp7A local p. 35; combined p. 488; printed p. 477.
Then \(x_\alpha^{(p)}=y_\alpha\) for every \(\alpha\). If
\[
  x=\sum_\alpha\lambda_\alpha x_\alpha,
\]
\SGARefLink{sga3:VIIA:5.1}{formula 5.1}~\SGARefLink{sga3:vii:paragraph:5-1:equation:ii}{(ii)}, applied inductively to the number of indices
\(\alpha\) for which \(\lambda_\alpha\ne0\), gives
\[
  x^p-\sum_\alpha\lambda_\alpha^p x_\alpha^p
  \in\mathfrak g.
\]
Denote this element by \(z\). Then
\[
  x^{(p)}
  =
  \sum_\alpha\lambda_\alpha^p y_\alpha+z,
\]
and hence \(x^{(p)}\in\mathfrak g\).

It is clear that
\[
  (\lambda x)^{(p)}
  =
  \lambda^p x^{(p)}.
\]
Since \(\gamma(x\otimes_\pi1)\) is central, we have
\(\operatorname{ad}x^{(p)}=\operatorname{ad}x^p\). Jacobson's
first formula, \SGARefLink{sga3:VIIA:5.1}{5.1}~\SGARefLink{sga3:vii:paragraph:5-1:equation:i}{(i)}, therefore yields
\[
  \operatorname{ad}x^{(p)}
  =
  (\operatorname{ad}x)^p.
\]
Finally, Jacobson's second formula, \SGARefLink{sga3:VIIA:5.1}{5.1}~\SGARefLink{sga3:vii:paragraph:5-1:equation:ii}{(ii)}, shows that the map
\(x\mapsto x^{(p)}\) satisfies condition \SGARefLink{sga3:vii:paragraph:5-1:equation:ii}{(ii)} of \SGARefLink{sga3:VIIB:5.2:theorem}{5.2}. It thus
makes \(\mathfrak g\) into a restricted \(p\)-Lie algebra. This
proves the proposition.

\medskip
\noindent\textbf{5.3.3. Proposition.}
\phantomsection
\label{sga3:VIIA:5.3.3}
\textit{Let \(\mathfrak g\) be a restricted \(p\)-Lie algebra
over \(R\) whose underlying module is free with basis
\((x_\alpha)\). Then
\[
  j:\mathfrak g\longrightarrow U_p(\mathfrak g)
\]
is injective. If \(z_\alpha=j(x_\alpha)\), then
\(U_p(\mathfrak g)\) has as a basis the monomials
\[
  \prod_\alpha z_\alpha^{n_\alpha},
  \qquad
  0\le n_\alpha<p.
\]
Here the \(n_\alpha\) are zero except for finitely many
\(\alpha\); the basis is assumed to be totally ordered, and the
products are taken in increasing order.}

Identify \(\mathfrak g\) with a submodule of the enveloping
algebra \(U(\mathfrak g)\) by means of the canonical map \(i\).
For every family
\[
  n=(n_\alpha)
\]
of nonnegative integers, all but finitely many zero, put
\[
  |n|=\sum_\alpha n_\alpha,
  \qquad
  x^n=\prod_\alpha x_\alpha^{n_\alpha}.
\]
Write
\[
  n_\alpha=m_\alpha+p\ell_\alpha,
  \qquad
  0\le m_\alpha<p,
\]
and put
\[
  T_n
  =
  \prod_\alpha
  x_\alpha^{m_\alpha}
  \beta(x_\alpha)^{\ell_\alpha},
\]
where
\[
  \beta(x)=x^p-x^{(p)}
\]
is the map \(\mathfrak g\to U(\mathfrak g)\) defined in \SGARefLink{sga3:VIIA:5.3.1}{5.3.1}.

For \(r\in\mathbb N\), let \(U^r\) be the \(R\)-submodule of
\(U(\mathfrak g)\) generated by the \(x^n\) for which
\(|n|\le r\). Since the graded ring
\[
  \bigoplus_r U^r/U^{r-1}
\]
is commutative (cf.\ \textup{[BLie]}, I, \S2.6), for every
\(n\) we have
\[
  T_n-\prod_\alpha x_\alpha^{n_\alpha}
  \in U^{|n|-1}.
\]
For every \(s\in\mathbb N\), the \(x^n\) with \(|n|=s\) form a
basis of \(U^s/U^{s-1}\) by the Poincar\'e--Birkhoff--Witt
theorem (\emph{loc.\ cit.}, \S2.7); the same is therefore true
of the \(T_n\) with \(|n|=s\).

As \(s=|n|\) varies, the \(T_n\) consequently form a basis of
\(U(\mathfrak g)\). The kernel \(J\) of the canonical map
\[
  U(\mathfrak g)\longrightarrow U_p(\mathfrak g)
\]
is the left ideal of \(U(\mathfrak g)\) generated by the central
elements \(\beta(x_\alpha)\), by \SGARefLink{sga3:VIIA:5.3.1}{5.3.1}. The \(T_n\) for which
\(\ell=(\ell_\alpha)\ne0\) therefore form a basis of \(J\), while
the \(T_n\) for which \(n_\alpha<p\) for every \(\alpha\) form a
basis of
\[
  U_p(\mathfrak g)=U(\mathfrak g)/J.
\]

\medskip
\noindent\textbf{5.3.3 bis.}
\phantomsection
\label{sga3:VIIA:5.3.3bis}

% Authority: Exp7A local p. 36; combined p. 489; printed p. 478.
Let \(\mathfrak g\) be a restricted \(p\)-Lie algebra over \(R\),
and let \(f:R\to R'\) be an extension of the base ring. We claim
that the \(R'\)-module
\[
  R'\otimes_R\mathfrak g
\]
admits one and only one restricted \(p\)-Lie algebra structure
such that
\[
  \tag{*}\label{sga3:vii:paragraph:5-3-3-bis:equation:star}
  [\lambda\otimes x,\mu\otimes y]
  =
  \lambda\mu\otimes[x,y],
  \qquad
  (\lambda\otimes x)^{(p)}
  =
  \lambda^p\otimes x^{(p)}.
\]
In particular, it will follow that the functor
\[
  \mathfrak g\longmapsto R'\otimes_R\mathfrak g
\]
is left adjoint to restriction of scalars from \(R'\) to \(R\).

The uniqueness of the structure defined by \((*)\) is clear; we
prove its existence. If \(\mathfrak g\) is free with basis
\((x_\alpha)\), \SGARefLink{sga3:VIIA:5.3.2}{Proposition 5.3.2} gives a unique restricted
\(p\)-Lie algebra structure on the Lie algebra
\(R'\otimes_R\mathfrak g\) such that
\[
  (1\otimes x_\alpha)^{(p)}
  =
  1\otimes x_\alpha^{(p)}.
\]
This is the structure sought.

Suppose now that \(\mathfrak g\) is arbitrary. There is a
restricted \(p\)-Lie algebra \(L_0\), free as an \(R\)-module,
together with a surjective homomorphism
\[
  q_0:L_0\longrightarrow\mathfrak g.
\]
For example, we may take
\[
  L_0=R\otimes_{\mathbb F_p}\mathfrak g
\]
and let \(q_0\) be the homomorphism
\[
  \lambda\otimes x\longmapsto\lambda x;
\]
here \(\mathfrak g\) is free over the prime field
\(\mathbb F_p\). The kernel of \(q_0\) is a \(p\)-ideal of
\(L_0\), that is, an ideal of the Lie algebra \(L_0\) stable
under \(x\mapsto x^{(p)}\). Hence there is also a restricted
\(p\)-Lie algebra \(L_1\), free as an \(R\)-module, and a
homomorphism \(q_1:L_1\to L_0\) whose image is
\(\operatorname{Ker}q_0\). This gives the exact sequence
\begin{center}
  \phantomsection
  \label{sga3:VIIA:5.3.3bis-exact-sequence}
  % SGA3-VII-DIAG-050
  % BEGIN INCLUDED SGA3_English_current_buildable_source_20260823/source/inputs/VII/native_diagrams/exp7/viia_05_14/diagram_050.tex
\(\displaystyle
L_1 \xrightarrow{q_1} L_0 \xrightarrow{q_0} \mathfrak g \longrightarrow 0 .
\)

% END INCLUDED SGA3_English_current_buildable_source_20260823/source/inputs/VII/native_diagrams/exp7/viia_05_14/diagram_050.tex

\end{center}
and therefore an exact sequence of \(R'\)-Lie algebras
\begin{center}
  \phantomsection
  \label{sga3:VIIA:5.3.3bis-basechange-exact-sequence}
  % SGA3-VII-DIAG-051
  % BEGIN INCLUDED SGA3_English_current_buildable_source_20260823/source/inputs/VII/native_diagrams/exp7/viia_05_14/diagram_051.tex
\(\displaystyle
R'\otimes_R L_1
  \xrightarrow{R'\otimes_R q_1}
R'\otimes_R L_0
  \xrightarrow{R'\otimes_R q_0}
R'\otimes_R\mathfrak g
  \longrightarrow 0 .
\)

% END INCLUDED SGA3_English_current_buildable_source_20260823/source/inputs/VII/native_diagrams/exp7/viia_05_14/diagram_051.tex

\end{center}

Since \(R'\otimes_Rq_1\) is manifestly a homomorphism of
restricted \(p\)-Lie algebras, the kernel of
\(R'\otimes_Rq_0\) is a \(p\)-ideal. Thus the symbolic \(p\)-th
power operation on \(R'\otimes_RL_0\) induces, by passage to the
quotient, a map from \(R'\otimes_R\mathfrak g\) to itself (use
formula \SGARefLink{sga3:vii:paragraph:5-1:equation:ii}{(ii)} of \SGARefLink{sga3:VIIB:5.2:theorem}{5.2}). This map equips \(R'\otimes_R\mathfrak g\)
with the required restricted \(p\)-Lie algebra structure.

\medskip
\noindent\textbf{5.3.4.}
\phantomsection
\label{sga3:VIIA:5.3.4}

The canonical map
\[
  j_{\mathfrak g}:\mathfrak g\longrightarrow U_p(\mathfrak g)
\]
induces, for every extension \(R\to R'\) of the base ring, a
homomorphism
\[
  R'\otimes_Rj_{\mathfrak g}:
  R'\otimes_R\mathfrak g
  \longrightarrow
  R'\otimes_RU_p(\mathfrak g).
\]
It therefore induces a homomorphism
\[
  h:
  U_p(R'\otimes_R\mathfrak g)
  \longrightarrow
  R'\otimes_RU_p(\mathfrak g)
\]
such that
\[
  h\circ j_{R'\otimes_R\mathfrak g}
  =
  R'\otimes_Rj_{\mathfrak g}.
\]
The universal properties of \(R'\otimes_R\mathfrak g\) and of
the restricted enveloping algebra show immediately that \(h\) is
an isomorphism. We may therefore identify
\[
  U_p(R'\otimes_R\mathfrak g)
  \simeq
  R'\otimes_RU_p(\mathfrak g).
\]

% Authority: Exp7A local p. 37; combined p. 490; printed p. 479.
In particular, let \(r\in R\), and let \(R'=R_r\) be the
localized ring. Then
\[
  \mathfrak g_r=R_r\otimes_R\mathfrak g
\]
has a canonical restricted \(p\)-Lie algebra structure over
\(R_r\). Consequently, the sheaf \(\widetilde{\mathfrak g}\) on
\(\operatorname{Spec}R\) is a quasi-coherent restricted
\(p\)-Lie algebra on \(\operatorname{Spec}R\). Moreover,
\[
  U_p^{R_r}(\mathfrak g_r)
  \simeq
  U_p^R(\mathfrak g)_r,
\]
so the sheaf associated with the presheaf
\[
  V\longmapsto
  U_p\bigl(\Gamma(V,\widetilde{\mathfrak g})\bigr)
\]
is quasi-coherent.

\noindent\textbf{Definition.}
More generally, let \(S\) be a scheme of characteristic \(p\),
and let \(\mathcal G\) be a quasi-coherent restricted \(p\)-Lie
algebra over \(\mathcal O_S\). The sheaf associated with the
presheaf
\[
  V\longmapsto U_p\bigl(\Gamma(V,\mathcal G)\bigr)
\]
is quasi-coherent. It is denoted
\(\mathcal U_p(\mathcal G)\) and is called the
\emph{restricted enveloping algebra} of \(\mathcal G\). If \(V\)
is affine, then
\[
  U_p\bigl(\Gamma(V,\mathcal G)\bigr)
  \simeq
  \Gamma\bigl(V,\mathcal U_p(\mathcal G)\bigr).
\]

% Next cursor: Exp7A local p. 37, combined-reader p. 490,
% re-edition printed p. 479, at 5.4.

% END INCLUDED SGA3_English_current_buildable_source_20260823/source/inputs/VII/tex/components/09_expose_VIIA_section53_restricted_enveloping_en.tex

% BEGIN INCLUDED SGA3_English_current_buildable_source_20260823/source/inputs/VII/tex/components/10_expose_VIIA_section54_55_coalgebra_group_en.tex
% SGA 3, Expose VII_A, English reconstruction -- Loop 1.
% Scope: authority local pages 37--40, paragraphs 5.4--5.5.3,
% stopping before Section 6 on local page 40.
% Combined-reader pages: 490--493. Re-edition printed pages: 479--482.
% Original-edition marginal locators within this span: 452--454.
% Authority: Polo--Gille Exp7A-13oct24.pdf,
% SHA-256 D46DB8C79D0385643B6EED4B96900613248375F9B73B4202478A4F57FD3E8522.
% Jacob Reinhold's CC BY 4.0 Markdown was consulted as comparison lineage.

\medskip
\noindent\textbf{5.4.}
\phantomsection
\label{sga3:VIIA:5.4}

% Authority: Exp7A local p. 37; combined p. 490; printed p. 479.
The universal property of \(U_p(\mathfrak g)\) implies that
\(U_p(\mathfrak g)\) is functorial in \(\mathfrak g\). Every
homomorphism of restricted \(p\)-Lie algebras
\[
  \varphi:\mathfrak g\longrightarrow\mathfrak h
\]
induces a unique homomorphism of unital algebras \(U_p(\varphi)\)
such that
\[
  j_{\mathfrak h}\circ\varphi
  =
  U_p(\varphi)\circ j_{\mathfrak g}.
\]
Here are some examples.

\smallskip
\noindent\textit{a)}
If \(\mathfrak h=0\), then \(U_p(\mathfrak h)\) identifies with the
base ring, and \(U_p(\varphi)\) is an algebra homomorphism
\[
  \varepsilon_{\mathfrak g}:
  U_p(\mathfrak g)\longrightarrow R
\]
called the \emph{augmentation}.

\smallskip
\noindent\textit{b)}
Let now \(\mathfrak h=\mathfrak g^\circ\), the Lie algebra opposite
to \(\mathfrak g\). Thus \(\mathfrak g^\circ\) has the same
underlying module and the same symbolic \(p\)-th power as
\(\mathfrak g\), while its bracket is the negative of the bracket
of \(\mathfrak g\). We may identify
\(U_p(\mathfrak g^\circ)\) with the algebra opposite to
\(U_p(\mathfrak g)\). Moreover, the isomorphism
\[
  x\longmapsto-x
\]
from \(\mathfrak g\) onto \(\mathfrak g^\circ\) induces an
isomorphism
\[
  c_{\mathfrak g}:
  U_p(\mathfrak g)
  \longrightarrow
  U_p(\mathfrak g^\circ)
  \simeq
  U_p(\mathfrak g)^\circ.
\]
We call \(c_{\mathfrak g}\) the \emph{antipode} of
\(U_p(\mathfrak g)\).

\smallskip
\noindent\textit{c)}
Finally, let \(\mathfrak f\) and \(\mathfrak g\) be two restricted
\(p\)-Lie algebras, and let
\(\mathfrak h=\mathfrak f\times\mathfrak g\) be their product.
Its underlying \(R\)-module is the direct product, and
\[
  [(x,y),(x',y')]
  =
  ([x,x'],[y,y']),
  \qquad
  (x,y)^{(p)}
  =
  (x^{(p)},y^{(p)}).
\]
If
\[
  h_1:\mathfrak f\to\mathfrak k,
  \qquad
  h_2:\mathfrak g\to\mathfrak k
\]
are restricted \(p\)-Lie algebra homomorphisms satisfying
\[
  [h_1(x),h_2(y)]=0
  \quad
  (x\in\mathfrak f,\ y\in\mathfrak g),
\]
then
\[
  h_1+h_2:(x,y)\longmapsto h_1(x)+h_2(y)
\]
is a restricted \(p\)-Lie algebra homomorphism. Conversely, every
homomorphism \(\mathfrak f\times\mathfrak g\to\mathfrak k\) is of
this form. This characterizes \(\mathfrak f\times\mathfrak g\) by
a universal property.

For example, the maps
\[
  h_1:x\longmapsto j_{\mathfrak f}(x)\otimes1,
  \qquad
  h_2:y\longmapsto1\otimes j_{\mathfrak g}(y)
\]
induce a homomorphism \(h_1+h_2\) from
\(\mathfrak f\times\mathfrak g\) to the restricted \(p\)-Lie
algebra underlying
\[
  U_p(\mathfrak f)\otimes U_p(\mathfrak g).
\]
The universal properties of \(\mathfrak f\times\mathfrak g\) and
of the restricted enveloping algebras show that this homomorphism
extends to an isomorphism
\[
  \varphi:
  U_p(\mathfrak f\times\mathfrak g)
  \xrightarrow{\sim}
  U_p(\mathfrak f)\otimes U_p(\mathfrak g).
\]

\medskip
\noindent\textbf{Definition.}
\phantomsection
\label{sga3:VIIA:5.4-comultiplication-definition}
If \(\mathfrak f=\mathfrak g\), the diagonal map
\[
  \delta:x\longmapsto(x,x)
\]
from \(\mathfrak g\) to
\(\mathfrak g\times\mathfrak g\) induces a homomorphism from
\(U_p(\mathfrak g)\) to
\(U_p(\mathfrak g\times\mathfrak g)\). We denote by
\[
  \Delta_{\mathfrak g}:
  U_p(\mathfrak g)
  \longrightarrow
  U_p(\mathfrak g)\otimes U_p(\mathfrak g)
\]
the composite of this homomorphism with \(\varphi\).
\footnote{Thus
\(\Delta_{\mathfrak g}(x)=x\otimes1+1\otimes x\) for every
\(x\in\mathfrak g\); in particular,
\(\Delta_{\mathfrak g}\) is cocommutative.}

It is then easy to see that \(\Delta_{\mathfrak g}\), together
with the multiplication of \(U_p(\mathfrak g)\), makes
\(U_p(\mathfrak g)\) an \(R\)-coalgebra in groups (cf.\ \SGARefLink{sga3:VIIB:3.2:prop}{3.2}), with
augmentation \(\varepsilon_{\mathfrak g}\) and antipode
\(c_{\mathfrak g}\).

\medskip
\noindent\textbf{5.5.}
\phantomsection
\label{sga3:VIIA:5.5}
\footnote{The editors reorganized \S5.4.1 of the original into this
\S5.5. \SGARefLink{sga3:VIIA:5.5.1}{Proposition 5.5.1} combines results from \SGARefLink{sga3:VIIA:5}{Section~5} and
\SGARefLink{sga3:VIIA:3.2.3}{Proposition~3.2.3} and incorporates the original \SGARefLink{sga3:VIIA:7.3}{Lemma~7.3}; the
proof of Proposition 5.5.3(ii) expands the proof of
\((i)\Rightarrow(ii)\) in the original \SGARefLink{sga3:VIIA:7.4}{Theorem~7.4}.}

% Authority: Exp7A local p. 38; combined p. 491; printed p. 480.
Let \(S\) be a scheme of characteristic \(p\). First let
\(\mathcal U\) be an \(\mathcal O_S\)-coalgebra in groups and put
\[
  G=\operatorname{Spec}^{*}\mathcal U.
\]
We saw in \SGARefLink{sga3:VIIA:3.2.3}{3.2.3} that, for every \(T\to S\),
\((\operatorname{Lie}G)(T)\) is the Lie subalgebra of
\(\Gamma(T,\mathcal U_T)\) consisting of primitive elements. If
\(x\) is primitive, then
\[
  \Delta(x^p)
  =
  x^p\otimes1+1\otimes x^p,
\]
since \(\binom pi\equiv0\pmod p\) for \(0<i<p\). Thus \(x^p\) is
again primitive. By \SGARefLink{sga3:VIIB:5.1:proposition}{5.1} and \SGARefLink{sga3:VIIB:5.2:theorem}{5.2}, the map \(x\mapsto x^p\) makes
\((\operatorname{Lie}G)(T)\) into an
\(\mathcal O(T)\)-restricted \(p\)-Lie algebra.

Now let \(\mathcal L\) be a quasi-coherent restricted
\(p\)-Lie algebra over \(\mathcal O_S\). As \(V\) ranges over the
open subsets of \(S\), the group-coalgebra structures defined
above on
\[
  U_p\bigl(\Gamma(V,\mathcal L)\bigr)
\]
induce on the associated sheaf
\(\mathcal U_p(\mathcal L)\) a structure of
\(\mathcal O_S\)-coalgebra in groups. Moreover, for every
\(S\)-scheme \(T\), there is an isomorphism
\[
  \mathcal U_p(\mathcal L_T)
  \xrightarrow{\sim}
  \mathcal U_p(\mathcal L)_T.
\]

Let \(\operatorname{Prim}\mathcal U_p(\mathcal L)\) denote the
subpresheaf which assigns to every open subset \(V\) the primitive
elements of \(\mathcal U_p(\mathcal L)(V)\); this is readily seen
to be a sheaf. As \(V\) ranges over the open subsets of \(S\), the
composites
\begin{center}
  \phantomsection
  \label{sga3:VIIA:5.5-primitive-sheaf-sequence}
  % SGA3-VII-DIAG-052
  % BEGIN INCLUDED SGA3_English_current_buildable_source_20260823/source/inputs/VII/native_diagrams/exp7/viia_05_14/diagram_052.tex
\(\displaystyle
\Gamma(V,\mathcal L)
  \xrightarrow{j}
\operatorname{Prim}\mathrm U_p\!\bigl(\Gamma(V,\mathcal L)\bigr)
  \longrightarrow
\operatorname{Prim}\mathcal U_p(\mathcal L)(V)
\)

% END INCLUDED SGA3_English_current_buildable_source_20260823/source/inputs/VII/native_diagrams/exp7/viia_05_14/diagram_052.tex

\end{center}
define a morphism
\[
  \mathcal L
  \longrightarrow
  \operatorname{Prim}\mathcal U_p(\mathcal L),
\]
again denoted \(j\) or \(j_{\mathcal L}\). It in turn defines a
morphism of \(\mathcal O_S\)-restricted \(p\)-Lie algebras
\[
  \mathcal W(\mathcal L)
  \longrightarrow
  \operatorname{Prim}
    \mathcal W\bigl(\mathcal U_p(\mathcal L)\bigr)
\]
(cf.\ \SGARefLink{sga3:VIIA:3.2.3}{3.2.3}).

\medskip
\noindent\textbf{Proposition 5.5.1.}
\phantomsection
\label{sga3:VIIA:5.5.1}
\textit{Let \(\mathcal L\) be an
\(\mathcal O_S\)-restricted \(p\)-Lie algebra, locally free as an
\(\mathcal O_S\)-module. Then \(j_{\mathcal L}\) induces an
isomorphism of \(\mathcal O_S\)-restricted \(p\)-Lie algebras}
\[
  \mathcal W(\mathcal L)
  \xrightarrow{\sim}
  \operatorname{Prim}
    \mathcal W\bigl(\mathcal U_p(\mathcal L)\bigr).
\]

\begin{proof}
Let \(T\) be an \(S\)-scheme. In view of the identification
\[
  \mathcal U_p(\mathcal L_T)
  =
  \mathcal U_p(\mathcal L)_T,
\]
we must show that
\[
  \Gamma(T,\mathcal L_T)
  \longrightarrow
  \operatorname{Prim}
    \Gamma\bigl(T,\mathcal U_p(\mathcal L_T)\bigr)
\]
is bijective. Replacing \(S\) by \(T\), we reduce to \(T=S\); it
then suffices to prove that the morphism of sheaves
\[
  j_{\mathcal L}:
  \mathcal L
  \longrightarrow
  \operatorname{Prim}\mathcal U_p(\mathcal L)
\]
is an isomorphism. The question is local on \(S\), so suppose that
\(S=\operatorname{Spec}R\) and that \(\mathcal L\) is the sheaf
associated with an \(R\)-restricted \(p\)-Lie algebra \(L\) with a
totally ordered basis \((x_\alpha)\).

As in \SGARefLink{sga3:VIIA:5.3.3}{5.3.3}, let \(z_\alpha\) be the image of \(x_\alpha\) in
\(U=U_p(L)\). For a family
\(n=(n_\alpha)\), with \(0\leq n_\alpha\leq p-1\) and only
finitely many nonzero terms, put
\[
  z^{(n)}
  =
  \prod_\alpha\frac{z_\alpha^{n_\alpha}}{n_\alpha!},
\]
the product being taken in increasing order. Since
\[
  \Delta(z_\alpha)
  =
  z_\alpha\otimes1+1\otimes z_\alpha,
\]
one obtains
\[
  \Delta(z^{(n)})
  =
  \sum_r z^{(n-r)}\otimes z^{(r)},
\]
where the sum is over the finite set of \(r\) satisfying
\[
  0\leq r_\alpha\leq n_\alpha
  \quad\text{for every }\alpha.
\]
The \(z^{(n)}\), respectively the
\(z^{(n)}\otimes z^{(m)}\), form a basis of \(U\), respectively
of \(U\otimes U\). It follows that
\[
  \Delta(u)=u\otimes1+1\otimes u
\]
holds for \(u\in U\) if and only if \(u\) is a linear combination
of the \(z_\alpha\). This proves the proposition.
\end{proof}

\medskip
\noindent\textbf{Remark 5.5.2.}
\phantomsection
\label{sga3:VIIA:5.5.2}
Recall from \SGARefLink{sga3:VIIA:3.2.2}{3.2.2} and \SGARefLink{sga3:VIIA:3.2.3}{3.2.3} that the group \(S\)-functor
\[
  G=\operatorname{Spec}^{*}\mathcal U_p(\mathcal L)
\]
is very good and that
\[
  \operatorname{Lie}G
  =
  \operatorname{Prim}
    \mathcal W\bigl(\mathcal U_p(\mathcal L)\bigr).
\]
Thus \SGARefLink{sga3:VIIA:5.5.1}{Proposition 5.5.1} says that \(j_{\mathcal L}\) induces an
isomorphism
\[
  \mathcal W(\mathcal L)
  \xrightarrow{\sim}
  \operatorname{Lie}G.
\]

If, moreover, \(\mathcal L\) is locally free of finite rank, then
\(\mathcal U_p(\mathcal L)\) is finite locally free over
\(\mathcal O_S\) by \SGARefLink{sga3:VIIA:5.3.3}{5.3.3}. Hence
\(\operatorname{Spec}^{*}\mathcal U_p(\mathcal L)\) is represented
by the \(S\)-group
\[
  \mathfrak G_p(\mathcal L)
  =
  \operatorname{Spec}\mathcal U_p(\mathcal L)^{*}
\]
(cf.\ \SGARefLink{sga3:VIIA:3.2.2.1}{3.2.2.1}). We obtain the following more precise statement.

\medskip
\noindent\textbf{Proposition 5.5.3.}
\phantomsection
\label{sga3:VIIA:5.5.3}
\textit{Let \(\mathcal L\) be an
\(\mathcal O_S\)-restricted \(p\)-Lie algebra, locally free of
finite rank as an \(\mathcal O_S\)-module. Put}
\[
  \mathcal A=\mathcal U_p(\mathcal L)^{*},
  \qquad
  G=\mathfrak G_p(\mathcal L)=\operatorname{Spec}\mathcal A.
\]
\textit{Then:}

\smallskip
\noindent
\textit{\SGAInlineRefTarget{sga3:VIIA:5.5.3:item:i}(i) \(j_{\mathcal L}\) induces an isomorphism}
\[
  \mathcal W(\mathcal L)
  \xrightarrow{\sim}
  \operatorname{Lie}G
\]
\textit{of \(\mathcal O_S\)-restricted \(p\)-Lie algebras.}

\smallskip
\noindent
\textit{\SGAInlineRefTarget{sga3:VIIA:5.5.3:item:ii}(ii) Let \(\mathcal I\) be the augmentation ideal of
\(\mathcal A\) and put
\(\omega_G=\mathcal I/\mathcal I^2\) (cf.\ II, \SGARefLink{sga3:ii:remark:4-11-4}{4.11.4}). Then}
\[
  \omega_G
  \simeq
  \mathcal L^{*}
  =
  \operatorname{Hom}_{\mathcal O_S}
    (\mathcal L,\mathcal O_S).
\]
\textit{In particular, \(\omega_G\) is locally free of finite rank,
and \(\omega_{G/S}^{*}\simeq\mathcal L\).}

\begin{proof}
Part~\SGARefLink{sga3:vii:paragraph:5-1:equation:i}{(i)} follows from \SGARefLink{sga3:VIIA:5.5.2}{Remark~5.5.2}. We prove~\SGARefLink{sga3:vii:paragraph:5-1:equation:ii}{(ii)}. Let
\(\eta_{\mathcal U}\) and \(\varepsilon_{\mathcal U}\) denote the
unit and augmentation of
\(\mathcal U=\mathcal U_p(\mathcal L)\), and let
\(\eta_{\mathcal A}\) and \(\varepsilon_{\mathcal A}\) denote
those of \(\mathcal A\). Put
\[
  \mathcal J=\ker\varepsilon_{\mathcal U}.
\]
Then
\begin{equation}
  \tag{1}\label{sga3:vii:proposition:5-5-3:equation:1}
  \mathcal U
  =
  \eta_{\mathcal U}(\mathcal O_S)\oplus\mathcal J.
\end{equation}
Let \(\delta\) be the morphism defined by the diagram
\begin{center}
  \phantomsection
  \label{sga3:VIIA:5.5.3-augmentation-kernel-delta-square}
  % SGA3-VII-DIAG-053
  % BEGIN INCLUDED SGA3_English_current_buildable_source_20260823/source/inputs/VII/native_diagrams/exp7/viia_05_14/diagram_053.tex
\begin{tikzcd}[row sep=3.2em,column sep=4.2em]
\mathcal J \arrow[r,"\delta"] \arrow[d,swap,"\tau"]
  & \mathcal J\otimes\mathcal J \\
\mathcal U \arrow[r,"\Delta"]
  & \mathcal U\otimes\mathcal U \arrow[u,swap,"\pi"]
\end{tikzcd}

% END INCLUDED SGA3_English_current_buildable_source_20260823/source/inputs/VII/native_diagrams/exp7/viia_05_14/diagram_053.tex

\end{center}
where \(\tau\) and \(\pi\) are the inclusion and projection
deduced from decomposition~\SGARefLink{sga3:vii:proposition:5-5-3:equation:1}{(1)}. There is an exact sequence
\begin{center}
  \phantomsection
  \label{sga3:VIIA:5.5.3-primitive-exact-sequence}
  % SGA3-VII-DIAG-054
  % BEGIN INCLUDED SGA3_English_current_buildable_source_20260823/source/inputs/VII/native_diagrams/exp7/viia_05_14/diagram_054.tex
\(\displaystyle
(\ast)\qquad
0 \longrightarrow \mathcal L^*
  \xrightarrow{j_{\mathcal L}}
\mathcal J
  \xrightarrow{\delta}
\mathcal J\otimes\mathcal J .
\)

% END INCLUDED SGA3_English_current_buildable_source_20260823/source/inputs/VII/native_diagrams/exp7/viia_05_14/diagram_054.tex

\end{center}

By \SGARefLink{sga3:VIIA:5.3.3}{5.3.3}, the \(\mathcal O_S\)-module
\(\mathcal J/\mathcal L\) is locally free, and by \SGARefLink{sga3:VIIA:5.5.1}{5.5.1} the
preceding sequence remains exact after every base change.
Consequently, \(\delta\) induces an isomorphism from
\(\mathcal J/\mathcal L\) onto a locally direct summand \(Q\) of
\(\mathcal J\otimes\mathcal J\). By duality, the exact sequence
therefore becomes
\begin{center}
  \phantomsection
  \label{sga3:VIIA:5.5.3-dual-exact-sequence}
  % SGA3-VII-DIAG-055
  % BEGIN INCLUDED SGA3_English_current_buildable_source_20260823/source/inputs/VII/native_diagrams/exp7/viia_05_14/diagram_055.tex
\(\displaystyle
(\ast\ast)\qquad
0 \longleftarrow \mathcal L
  \xleftarrow{{}^{t}\!j_{\mathcal L}}
\mathcal I
  \xleftarrow{{}^{t}\!\delta}
\mathcal I\otimes\mathcal I .
\)

% END INCLUDED SGA3_English_current_buildable_source_20260823/source/inputs/VII/native_diagrams/exp7/viia_05_14/diagram_055.tex

\end{center}

% Authority: Exp7A local p. 40; combined p. 493; printed p. 482.
Decomposition~\SGARefLink{sga3:vii:proposition:5-5-3:equation:1}{(1)} corresponds by duality to
\begin{equation}
  \tag{2}\label{sga3:vii:proposition:5-5-3:equation:2}
  \mathcal A
  =
  \eta_{\mathcal A}(\mathcal O_S)\oplus\mathcal I.
\end{equation}
The transpose of \(\Delta\) is the multiplication
\[
  m_{\mathcal A}:
  \mathcal A\otimes\mathcal A
  \longrightarrow
  \mathcal A.
\]
Since \(\mathcal I\) is an ideal of \(\mathcal A\),
\(m_{\mathcal A}\) sends
\(\mathcal I\otimes\mathcal I\) into \(\mathcal I\). More
precisely, decomposition~\SGARefLink{sga3:vii:proposition:5-5-3:equation:2}{(2)} gives the commutative square
\begin{center}
  \phantomsection
  \label{sga3:VIIA:5.5.3-augmentation-ideal-multiplication-square}
  % SGA3-VII-DIAG-056
  % BEGIN INCLUDED SGA3_English_current_buildable_source_20260823/source/inputs/VII/native_diagrams/exp7/viia_05_14/diagram_056.tex
\begin{tikzcd}[row sep=3.1em,column sep=4.2em]
\mathcal I
  & \mathcal I\otimes\mathcal I \arrow[l,swap,"m'"] \arrow[d,"{}^{t}\!\pi"] \\
\mathcal A \arrow[u,"{}^{t}\!\tau"]
  & \mathcal A\otimes\mathcal A . \arrow[l,swap,"m_{\mathcal A}"]
\end{tikzcd}

% END INCLUDED SGA3_English_current_buildable_source_20260823/source/inputs/VII/native_diagrams/exp7/viia_05_14/diagram_056.tex

\end{center}
This shows that the restriction \(m'\) of \(m_{\mathcal A}\) to
\(\mathcal I\otimes\mathcal I\) is the transpose of \(\delta\).
The dual exact sequence therefore yields
\[
  \mathcal I/\mathcal I^2
  \simeq
  \mathcal L^{*},
\]
which proves the proposition.
\end{proof}

% Next cursor: Exp7A local p. 40, combined-reader p. 493,
% re-edition printed p. 482, at Section 6.

% END INCLUDED SGA3_English_current_buildable_source_20260823/source/inputs/VII/tex/components/10_expose_VIIA_section54_55_coalgebra_group_en.tex

% BEGIN INCLUDED SGA3_English_current_buildable_source_20260823/source/inputs/VII/tex/components/11_expose_VIIA_section6_p_lie_group_scheme_en.tex
% SGA 3, Expose VII_A, English reconstruction -- Loop 1.
% Scope: authority local pages 40--44, all of Section 6,
% stopping before Section 7 on local page 45.
% Combined-reader pages: 493--497. Re-edition printed pages: 482--486.
% Original-edition marginal locators within this span: 454--460.
% Authority: Polo--Gille Exp7A-13oct24.pdf,
% SHA-256 D46DB8C79D0385643B6EED4B96900613248375F9B73B4202478A4F57FD3E8522.
% Jacob Reinhold's CC BY 4.0 Markdown was consulted as comparison lineage.

\section*{6. The \(p\)-Lie Algebra of an \(S\)-Group Scheme}
\phantomsection
\label{sga3:VIIA:6}
\addcontentsline{toc}{section}{6. The p-Lie Algebra of an S-Group Scheme}

% Authority: Exp7A local p. 40; combined p. 493; printed p. 482.
Let \(S\) be a scheme of characteristic \(p>0\). In \SGARefLink{sga3:VIIA:5.5}{5.5} we
associated with every quasi-coherent
\(\mathcal O_S\)-\(p\)-Lie algebra \(\mathcal L\) the group
\(S\)-functor
\[
  \mathfrak G_p(\mathcal L)
  =
  \operatorname{Spec}^{*}\mathcal U_p(\mathcal L).
\]
We shall now see that, for every \(S\)-group scheme \(G\), the
\(\mathcal O_S\)-Lie algebra \(\operatorname{Lie}(G/S)\) defined
in II, \SGARefLink{sga3:ii:subsection:4-11}{4.11} is naturally endowed with the structure of an
\(\mathcal O_S\)-\(p\)-Lie algebra.

\medskip
\noindent\textbf{6.1.}
\phantomsection
\label{sga3:VIIA:6.1}

First identify \(\operatorname{Lie}(G/S)(S)\) and
\(\operatorname{Lie}(\operatorname{Aut}G/S)(S)\), respectively,
with Lie subalgebras of \(\mathrm U(G)\) and
\(\operatorname{Dif}_{G/S}\) by means of the injections
\(\alpha\) and \(\beta\) of \SGARefLink{sga3:VIIB:2.5:proposition}{2.5}. Thus
\(\operatorname{Lie}(\operatorname{Aut}G/S)(S)\) is identified
with the \(\Gamma(S,\mathcal O_S)\)-Lie algebra of
\(S\)-derivations of \(\mathcal O_G\). By \SGARefLink{sga3:VIIB:5.2:theorem}{5.2}, the latter is a
restricted \(p\)-Lie subalgebra of
\(\operatorname{Dif}_{G/S}\).

On the other hand, the image of
\[
  L=\operatorname{Lie}(G/S)(S)
\]
under the injective algebra homomorphism
\[
  \ell:\mathrm U(G)\longrightarrow\operatorname{Dif}_{G/S},
  \qquad
  d\longmapsto {}_Gd,
\]
consists of the left-invariant derivations (cf.\ \SGARefLink{sga3:VIIA:2.2}{2.2},
editor's note~(17), \SGARefLink{sga3:VIIA:2.4}{2.4}, and \SGARefLink{sga3:VIIB:2.5:proposition}{2.5}). If \(x\in L\), then
\(\ell(x)^p=\ell(x^p)\) by the cited passages. Since
\(\ell(x)^p\) is again a derivation, \(x^p\) belongs to
\(\operatorname{Lie}(G/S)(S)\). Hence
\begin{center}
  \textit{\(\operatorname{Lie}(G/S)(S)\) is a restricted
  \(p\)-Lie subalgebra of the infinitesimal algebra
  \(\mathrm U(G)\).}
\end{center}
\footnote{One may also prove directly, without passing through
\(\operatorname{Dif}_{G/S}\), that the Lie algebra of
derivations of \(G\) at the origin (isomorphic to
\(\operatorname{Lie}(G/S)(S)\) by \SGARefLink{sga3:VIIB:2.5:proposition}{2.5}) is stable under taking
\(p\)-th powers in \(\mathrm U(G)\). This is done in \SGARefLink{sga3:VIIA:6.2}{6.2} below.}

\medskip
\noindent\textbf{6.1.1.}
\phantomsection
\label{sga3:VIIA:6.1.1}

% Authority: Exp7A local p. 41; combined p. 494; printed p. 483.
Let \(\varphi:G\to H\) be a homomorphism of \(S\)-group schemes.
The homomorphisms
\(\operatorname{Lie}(\varphi/S)(S)\) and \(\mathrm U(\varphi)\)
are compatible with the identifications of
\(\operatorname{Lie}(G/S)(S)\) and
\(\operatorname{Lie}(H/S)(S)\) with restricted \(p\)-Lie
subalgebras of \(\mathrm U(G)\) and \(\mathrm U(H)\). Since
\(\mathrm U(\varphi)\) is an algebra homomorphism,
\(\operatorname{Lie}(\varphi/S)(S)\) is a homomorphism of
restricted \(p\)-Lie algebras.

Likewise, if \(s:T\to S\) is a base change, the map from
\(\operatorname{Lie}(G/S)(S)\) to
\(\operatorname{Lie}(G/S)(T)\) induced by \(s\) is a
homomorphism of restricted \(p\)-Lie algebras. Equivalently, the
functor \(\operatorname{Lie}(G/S)\) is endowed with a structure
of \(\mathcal O_S\)-\(p\)-Lie algebra. In particular, as \(T\)
ranges over the open subsets of \(S\), the sheaf
\(\operatorname{Lie}(G/S)\) is endowed with a structure of
\(\mathcal O_S\)-\(p\)-Lie algebra.

\medskip
\noindent\textbf{6.2.}
\phantomsection
\label{sga3:VIIA:6.2}

Following an idea of Demazure, we shall generalize the preceding
construction to certain group \(S\)-functors which need not be
representable. We first give another definition of the symbolic
\(p\)-th power in the Lie algebra of an \(S\)-group scheme \(G\).

Let \(D\) be a derivation of \(G\) at the
origin.\footnote{The editors expanded the argument in the
original text in what follows.} By \SGARefLink{sga3:VIIA:1.2.1}{1.2.1}, \(D\) is the composite
of the \(S\)-derivation
\[
  \delta=(\tau,\partial_t)
\]
of the zero section \(\tau:S\to I_S\), and a morphism
\(x:I_S\to G\) such that \(x\circ\tau=\varepsilon\), that is,
\(x\in\operatorname{Lie}(G/S)(S)\). By the definition in \SGARefLink{sga3:VIIA:2.1}{2.1},
\(D^p\) is the following composite deviation:
\begin{center}
  \phantomsection
  \label{sga3:VIIA:6.2-demazure-pfold-composite}
  % SGA3-VII-DIAG-057
  % BEGIN INCLUDED SGA3_English_current_buildable_source_20260823/source/inputs/VII/native_diagrams/exp7/viia_05_14/diagram_057.tex
\resizebox{0.98\linewidth}{!}{%
  \(\displaystyle
  S \simeq
  \underbrace{S\times S\times\cdots\times S}_{p}
    \xrightarrow{\delta\times\cdots\times\delta}
  I_S\times\cdots\times I_S
    \xrightarrow{x\times\cdots\times x}
  G\times\cdots\times G
    \xrightarrow{m^{(p)}} G
  \)%
}

% END INCLUDED SGA3_English_current_buildable_source_20260823/source/inputs/VII/native_diagrams/exp7/viia_05_14/diagram_057.tex

\end{center}
Here \(m^{(p)}\) is induced by the multiplication
\(m:G\times G\to G\). Since
\(I_S\times\cdots\times I_S\) is affine over \(S\), with affine
algebra
\[
  \mathcal B
  =
  \mathcal O_S[d_1,\ldots,d_p]/
  (d_1^2,\ldots,d_p^2),
\]
the deviation \(\delta\times\cdots\times\delta\) is defined by
the homomorphism of \(\mathcal O_S\)-modules
\[
  \phi:\mathcal B\longrightarrow\mathcal O_S
\]
which sends the monomial \(d_1d_2\cdots d_p\) to \(1\), and all
the other monomials \(d_{i_1}\cdots d_{i_r}\), for
\(0\le r<p\), to \(0\).

Let \(\operatorname{pr}_i\) denote the projection of
\(I_S^p\) onto its \(i\)-th factor, and let \(x_i\) be the image
of \(x\) in \(G(I_S^p)\) under \(G(\operatorname{pr}_i)\). Then
\[
  m^{(p)}\circ(x\times\cdots\times x)
\]
is precisely the product \(x_1x_2\cdots x_p\). Consequently,
\(D^p\) is also the composite deviation
\begin{center}
  \phantomsection
  \label{sga3:VIIA:6.2-demazure-product-composite}
  % SGA3-VII-DIAG-058
  % BEGIN INCLUDED SGA3_English_current_buildable_source_20260823/source/inputs/VII/native_diagrams/exp7/viia_05_14/diagram_058.tex
\(\displaystyle
S \xrightarrow{\delta\times\cdots\times\delta}
I_S\times\cdots\times I_S
\xrightarrow{x_1x_2\cdots x_p}
G .
\)

% END INCLUDED SGA3_English_current_buildable_source_20260823/source/inputs/VII/native_diagrams/exp7/viia_05_14/diagram_058.tex

\end{center}

This description gives another proof that \(D^p\) is a
derivation of \(G\) at the origin. Since \(G\) is a very good
group (II, \SGARefLink{sga3:ii:subsection:4-11}{4.11}), the images \(G(\operatorname{pr}_1)(x)\) and
\(G(\operatorname{pr}_2)(x)\) of \(x\) in
\(G(I_S\times I_S)\) commute. It follows that the elements
\(x_i\in G(I_S^p)\) commute pairwise. Hence, for every
permutation \(\gamma\) of the factors of \(I_S^p\),
\[
  (x_1\cdots x_p)\circ\gamma=x_1\cdots x_p.
\]
Thus \(x_1\cdots x_p\) factors through the canonical projection
from \(I_S^p\) to the symmetric product \(\Sigma^p I_S\)
(cf.\ \SGARefLink{sga3:VIIA:4.2}{4.2}).

The affine algebra of \(\Sigma^p I_S\) is the subalgebra
\(\mathcal A\) of \(\mathcal B\) having as an
\(\mathcal O_S\)-basis the elementary symmetric functions
\[
  1=\sigma_0,\sigma_1,\ldots,\sigma_p
\]
in \(d_1,\ldots,d_p\). Let
\(\kappa:\mathcal A\hookrightarrow\mathcal B\) be the
inclusion, and let
\[
  \pi:\mathcal A\longrightarrow\mathcal O_S[t]/(t^2)
\]
be the homomorphism of \(\mathcal O_S\)-algebras which kills
\(\sigma_1,\ldots,\sigma_{p-1}\) and sends \(\sigma_p\) to
\(t\). Then
\[
  \phi\circ\kappa=\partial_t\circ\pi,
\]
where \(\partial_t:\mathcal O_S[t]/(t^2)\to\mathcal O_S\)
kills \(1\) and sends \(t\) to \(1\). If
\[
  i:I_S\hookrightarrow\Sigma^p I_S
\]
denotes the closed immersion defined by \(\pi\), we obtain the
commutative diagram
\begin{center}
  \phantomsection
  \label{sga3:VIIA:6.2-symbolic-power-commutative-diagram}
  % SGA3-VII-DIAG-059
  % BEGIN INCLUDED SGA3_English_current_buildable_source_20260823/source/inputs/VII/native_diagrams/exp7/viia_05_14/diagram_059.tex
\begin{tikzcd}[row sep=3.5em,column sep=4.3em]
S
  \arrow[r,swap,"\delta\times\cdots\times\delta"]
  \arrow[d,swap,"\delta"]
  \arrow[rr,bend left=23,"D^p"]
  & I_S^p \arrow[r,swap,"x_1\cdots x_p"] \arrow[d,"\mathrm{can.}"]
  & G \arrow[d,equal] \\
I_S
  \arrow[r,"i"]
  \arrow[rr,bend right=23,swap,"y"]
  & \Sigma^p I_S \arrow[r]
  & G
\end{tikzcd}

% END INCLUDED SGA3_English_current_buildable_source_20260823/source/inputs/VII/native_diagrams/exp7/viia_05_14/diagram_059.tex

\end{center}
which shows that \(D^p\) has the form \(y\circ\delta\), and is
therefore indeed a derivation of \(G\) at the origin.

\medskip
\noindent\textbf{6.3.}
\phantomsection
\label{sga3:VIIA:6.3}

% Authority: Exp7A local p. 42; combined p. 495; printed p. 484.
Let \(\mathfrak S_p\) be the symmetric group of degree \(p\),
and let \(I_S^p\times\mathfrak S_p\) denote the disjoint union
of a family of copies of \(I_S^p\) indexed by
\(\mathfrak S_p\). Let
\[
  \pi:I_S^p\times\mathfrak S_p\longrightarrow I_S^p
\]
be the canonical projection, and let
\[
  \mu:I_S^p\times\mathfrak S_p\longrightarrow I_S^p
\]
be the morphism defining the action of \(\mathfrak S_p\) on
\(I_S^p\): if \(\rho\in\mathfrak S_p\), the restriction of
\(\mu\) to \(I_S^p\times\rho\) has
\(\operatorname{pr}_{\rho(j)}\) as its \(j\)-th component.

\noindent\textbf{Definition.}
A functor
\[
  X:(\operatorname{Sch}/S)^\circ\longrightarrow(\operatorname{Ens})
\]
satisfies \emph{condition (F)} if:
\begin{enumerate}
  \item[\textit{a)}] \(X\) takes finite disjoint unions to direct
  products;
  \item[\textit{b)}] for every \(S\)-scheme \(T\), the following
  sequence is exact:
\end{enumerate}
\begin{center}
  \phantomsection
  \label{sga3:VIIA:6.3-conditionF-equalizer}
  % SGA3-VII-DIAG-060
  % BEGIN INCLUDED SGA3_English_current_buildable_source_20260823/source/inputs/VII/native_diagrams/exp7/viia_05_14/diagram_060.tex
\resizebox{0.98\linewidth}{!}{%
\begin{tikzcd}[ampersand replacement=\&,column sep=4.4em]
X(T\times\Sigma^p I_S)
  \arrow[r]
  \& X(T\times I_S^p)
    \arrow[r,shift left=.8ex,"X(\operatorname{id}_T\times\pi)"]
    \arrow[r,shift right=.8ex,swap,"X(\operatorname{id}_T\times\mu)"]
  \& X(T\times I_S^p\times\mathfrak S_p)\, .
\end{tikzcd}%
}

% END INCLUDED SGA3_English_current_buildable_source_20260823/source/inputs/VII/native_diagrams/exp7/viia_05_14/diagram_060.tex

\end{center}

Every \(S\)-scheme satisfies (F). If \(\mathcal F\) is an
\(\mathcal O_S\)-module, then \(\mathcal W(\mathcal F)\)
satisfies (F). Every projective limit of functors satisfying
(F) again satisfies (F). Finally, if \(Y\) satisfies (F) and
\(X\) is an arbitrary \(S\)-functor, then
\(\operatorname{Hom}_S(X,Y)\) satisfies (F).

Let \(X\) be a very good group (II, \SGARefLink{sga3:ii:definition:4-10}{4.10}) satisfying condition
(F). Let \(x:I_S\to X\) be a morphism extending the unit section
of \(X\), and use the notation of \SGARefLink{sga3:VIIA:6.2}{6.2}. As above,
\[
  x_1\cdots x_p:I_S^p\longrightarrow X
\]
factors through \(\Sigma^p I_S\):
\begin{center}
  \phantomsection
  \label{sga3:VIIA:6.3-symbolic-power-factorization}
  % SGA3-VII-DIAG-061
  % BEGIN INCLUDED SGA3_English_current_buildable_source_20260823/source/inputs/VII/native_diagrams/exp7/viia_05_14/diagram_061.tex
\begin{tikzcd}[row sep=3.5em,column sep=2.1em]
I_S^p \arrow[rr,"x_1\cdots x_p"] \arrow[dr,swap,"\mathrm{can.}"]
  && X \\
& \Sigma^p I_S \arrow[ur,swap,"\Sigma^p(x)"]
&
\end{tikzcd}

% END INCLUDED SGA3_English_current_buildable_source_20260823/source/inputs/VII/native_diagrams/exp7/viia_05_14/diagram_061.tex

\end{center}
and therefore defines, by composition, a morphism
\begin{center}
  \phantomsection
  \label{sga3:VIIA:6.3-symbolic-pth-power-composite}
  % SGA3-VII-DIAG-062
  % BEGIN INCLUDED SGA3_English_current_buildable_source_20260823/source/inputs/VII/native_diagrams/exp7/viia_05_14/diagram_062.tex
\(\displaystyle
x^{(p)}:\quad
I_S \xrightarrow{i} \Sigma^p I_S
  \xrightarrow{\Sigma^p(x)} X
\)

% END INCLUDED SGA3_English_current_buildable_source_20260823/source/inputs/VII/native_diagrams/exp7/viia_05_14/diagram_062.tex

\end{center}
which we call the \emph{symbolic \(p\)-th power} of \(x\).

% Authority: Exp7A local p. 43; combined p. 496; printed p. 485.
The endomorphism
\[
  x\longmapsto x^{(p)}
\]
of \(\operatorname{Lie}(X/S)(S)\) is plainly compatible with
base change and functorial in \(X\). It would be interesting to
know for which \(X\) this endomorphism makes
\(\operatorname{Lie}(X/S)(S)\) into a restricted \(p\)-Lie
algebra.

\medskip
\noindent\textbf{6.4.}
\phantomsection
\label{sga3:VIIA:6.4}

The preceding definition of the symbolic \(p\)-th power is
particularly well suited to computation. Here are some
examples.

\medskip
\noindent\textbf{6.4.1.}
\phantomsection
\label{sga3:VIIA:6.4.1}

Let \(M\) be an ``abstract'' abelian group, and let \(D_S(M)\)
be the diagonalizable \(S\)-group of type \(M\) (I, \SGARefLink{sga3:i:paragraph:4-4-2}{4.4.2}). For
every \(S\)-scheme \(T\),
\[
  D_S(M)(T)
  =
  \operatorname{Hom}_{(\mathrm{Ab})}
  \bigl(M,\mathcal O(T)^\times\bigr).
\]
Let \(x\in\operatorname{Lie}(D_S(M)/S)(S)\), that is, an
abelian-group homomorphism
\[
  x:M\longrightarrow
  \Gamma(S,\mathcal O_S+d\mathcal O_S)^\times
\]
of the form
\[
  m\longmapsto1+d\xi(m),
  \qquad
  \xi\in
  \operatorname{Hom}_{(\mathrm{Ab})}(M,\mathcal O(S)).
\]
With the notation of \SGARefLink{sga3:VIIA:6.2}{6.2} and \SGARefLink{sga3:VIIA:6.3}{6.3}, the product
\(x_1\cdots x_p\) associates with \(m\in M\) the expression
\[
  (1+d_1\xi(m))\cdots(1+d_p\xi(m)),
\]
that is,
\[
  1+\sigma_1\xi(m)+\sigma_2\xi(m)^2+\cdots+
  \sigma_p\xi(m)^p.
\]
This expression belongs to \(\mathcal O(\Sigma^p I_S)\).
Projecting it to \(\mathcal O(S)[d]/(d^2)\), killing
\(\sigma_1,\ldots,\sigma_{p-1}\), and sending \(\sigma_p\) to
\(d\), we see that \(x^{(p)}\) is the homomorphism
\[
  m\longmapsto1+d\xi(m)^p.
\]
Thus, under the identification
\[
  \operatorname{Lie}(D_S(M)/S)(S)
  \simeq
  \operatorname{Hom}_{(\mathrm{Ab})}(M,\mathcal O(S))
\]
of II, \SGARefLink{sga3:ii:proposition:5-1}{5.1}, the symbolic \(p\)-th power sends \(\xi\) to
\[
  \xi^{(p)}:m\longmapsto\xi(m)^p.
\]

\medskip
\noindent\textbf{6.4.2.}
\phantomsection
\label{sga3:VIIA:6.4.2}

Let \(\mathcal F\) be an \(\mathcal O_S\)-module, and let
\(G=\mathcal W(\mathcal F)\) be the group \(S\)-functor in
abelian groups (cf.\ I, \SGARefLink{sga3:i:subsection:4-6}{4.6}). Let
\[
  y\in\mathcal W(\mathcal F)(S)=\Gamma(S,\mathcal F),
\]
and let \(y'\) be its image in
\(\mathcal W(\mathcal F)(I_S)\) under
\(\mathcal W(\mathcal F)(I_S\to S)\).

By II, \SGARefLink{sga3:ii:examples:4-4-2}{4.4.2} and \SGARefLink{sga3:ii:scholium:4-5-1}{4.5.1}, the map
\[
  y\longmapsto dy'
\]
is an isomorphism of \(\mathcal O(S)\)-modules from
\(\mathcal W(\mathcal F)(S)\) onto
\(\operatorname{Lie}(\mathcal W(\mathcal F)/S)(S)\). Put
\(x=dy'\). Then the element \(x_i\) of \SGARefLink{sga3:VIIA:6.2}{6.2} is
\(d_i y''\), where \(y''\) denotes the canonical image of
\(y'\)\footnote{The editors corrected \(x\) to \(y'\).} in
\(\mathcal W(\mathcal F)(I_S^p)\). Hence
\[
  x_1+\cdots+x_p
  =
  (d_1+\cdots+d_p)y''
  =
  \sigma_1y'',
\]
and this belongs to
\(\mathcal W(\mathcal F)(\Sigma^p I_S)\). Since the
homomorphism
\[
  \mathcal O(\Sigma^p I_S)\longrightarrow\mathcal O(I_S)
\]
defining the morphism \(i\) kills \(\sigma_1\), we have
\(x^{(p)}=0\). Therefore
\begin{center}
  \textit{For every \(\mathcal O_S\)-module \(\mathcal F\), the
  operation \(x\mapsto x^{(p)}\) on
  \(\operatorname{Lie}\mathcal W(\mathcal F)\) is zero.}
\end{center}

\medskip
\noindent\textbf{6.4.3.}
\phantomsection
\label{sga3:VIIA:6.4.3}

% Authority: Exp7A local p. 44; combined p. 497; printed p. 486.
Let \(X\) be an \(S\)-scheme, let
\(G=\operatorname{Aut}_S X\), and let \(D\) be an
\(S\)-derivation of the structure sheaf \(\mathcal O_X\). By
\SGARefLink{sga3:VIIA:1.5}{1.5}, \(D\) may be identified with an \(I_S\)-automorphism \(x\)
of \(X_{I_S}\) inducing the identity on \(X\). It can be
described as follows. If \(f\) is a section of
\(\mathcal O_X[d]/(d^2)\) of the form \(a+bd\), put
\[
  D_{I_S}f=Da+(Db)d.
\]
Thus \(D_{I_S}\) is obtained from \(D\) by the base change
\(I_S\to S\), and the automorphism of \(X_{I_S}\) is associated
with the endomorphism
\[
  f\longmapsto f+(D_{I_S}f)d
  =
  a+(D(a)+b)d
\]
of \(\mathcal O_X[d]/(d^2)\).

Similarly, let \(D_{I_S^p}\) be the differential operator on
\(X_{I_S^p}\) obtained from \(D\) by the base change
\(I_S^p\to S\). With the notation of \SGARefLink{sga3:VIIA:6.2}{6.2}, the automorphism
\(x_i\) of \(X_{I_S^p}\) is associated with the endomorphism
\[
  f\longmapsto f+(D_{I_S^p}f)d_i
\]
of
\[
  \mathcal O_X[d_1,\ldots,d_p]/(d_1^2,\ldots,d_p^2).
\]
The product \(x_1\cdots x_p\) is therefore associated with
\[
  (1+d_1D_{I_S^p})(1+d_2D_{I_S^p})\cdots
  (1+d_pD_{I_S^p}),
\]
that is, with
\[
  1+\sigma_1D_{I_S^p}
  +\sigma_2D_{I_S^p}^{\,2}
  +\cdots+
  \sigma_pD_{I_S^p}^{\,p}.
\]
The coefficient of \(\sigma_p\) is \(D_{I_S^p}^{\,p}\).
Consequently, the Lie-algebra isomorphism
\[
  \operatorname{Der}_S(\mathcal O_X)
  \xrightarrow{\ \sim\ }
  \operatorname{Lie}(\operatorname{Aut}_S X),
  \qquad
  D\longmapsto x
\]
(cf.\ \SGARefLink{sga3:VIIA:1.5}{1.5}) is also an isomorphism of restricted \(p\)-Lie
algebras.

\medskip
\noindent\textbf{6.4.4.}
\phantomsection
\label{sga3:VIIA:6.4.4}

The same method shows that, for every \(\mathcal O_S\)-module
\(\mathcal F\), the Lie-algebra isomorphism
\[
  \operatorname{Lie}
  \bigl(
    \operatorname{Aut}_{\mathcal O_S\text{-mod.}}
    \mathcal W(\mathcal F)/S
  \bigr)(S)
  \xrightarrow{\ \sim\ }
  \bigl(
    \operatorname{End}_{\mathcal O_S\text{-mod.}}
    \mathcal W(\mathcal F)
  \bigr)(S)
\]
(cf.\ II, \SGARefLink{sga3:ii:proposition:4-8}{4.8}) is also an isomorphism of restricted
\(p\)-Lie algebras.

\medskip
\noindent\textbf{6.4.5.}\footnote{The editors supplied the
number \SGARefLink{sga3:VIIA:6.4.5}{6.4.5} and expanded the original text.}
\phantomsection
\label{sga3:VIIA:6.4.5}

Let \(\mathcal U\) be an \(\mathcal O_S\)-coalgebra in groups,
and let
\[
  G=\operatorname{Spec}^{*}\mathcal U
\]
be the group \(S\)-functor; suppose that \(G\) is representable.
For every \(T\to S\), \SGARefLink{sga3:VIIA:5.5}{5.5} and \SGARefLink{sga3:VIIA:6.1.1}{6.1.1} define two restricted
\(p\)-Lie algebra structures on
\[
  L(T)=\operatorname{Lie}(G)(T).
\]
There is a commutative diagram
\begin{center}
  \phantomsection
  \label{sga3:VIIA:6.4.5-group-coalgebra-comparison}
  % SGA3-VII-DIAG-063
  % BEGIN INCLUDED SGA3_English_current_buildable_source_20260823/source/inputs/VII/native_diagrams/exp7/viia_05_14/diagram_063.tex
\begin{tikzcd}[row sep=3.4em,column sep=4.3em]
L(T) \arrow[r,hook,"\tau"] \arrow[d,hook,swap,"\alpha"] \arrow[dr,"i"]
  & \Gamma(T,\mathcal U_T) \\
\mathrm U(G_T)
  & \mathrm U(L(T)) \arrow[l,swap,"\varphi"] \arrow[u,swap,"\psi"]
\end{tikzcd}

% END INCLUDED SGA3_English_current_buildable_source_20260823/source/inputs/VII/native_diagrams/exp7/viia_05_14/diagram_063.tex

\end{center}
in which \(\mathrm U(L(T))\) is the enveloping algebra of
\(L(T)\), and \(\varphi,\psi\) are the algebra homomorphisms
induced by \(\alpha,\tau\). It follows that the two restricted
\(p\)-Lie algebra structures coincide: if
\(x\in L(T)\) is identified with its image in
\(\mathrm U(G_T)\), respectively in
\(\Gamma(T,\mathcal U_T)\), then \(x^{(p)}\) is the image under
\(\varphi\), respectively under \(\psi\), of the element
\(x^p\in\mathrm U(L(T))\).

% Next cursor: Exp7A local p. 45, combined-reader p. 498,
% re-edition printed p. 487, at Section 7.

% END INCLUDED SGA3_English_current_buildable_source_20260823/source/inputs/VII/tex/components/11_expose_VIIA_section6_p_lie_group_scheme_en.tex

% BEGIN INCLUDED SGA3_English_current_buildable_source_20260823/source/inputs/VII/tex/components/12_expose_VIIA_section7_radicial_height1_en.tex
% SGA 3, Expose VII_A, English reconstruction -- Loop 1.
% Scope: authority local pages 45--52, all of Section 7,
% including Remark 7.5 and stopping before Section 8.
% Combined-reader pages: 498--505. Re-edition printed pages: 487--494.
% Original-edition marginal locators within this span: 461--469.
% Authority: Polo--Gille Exp7A-13oct24.pdf,
% SHA-256 D46DB8C79D0385643B6EED4B96900613248375F9B73B4202478A4F57FD3E8522.
% Jacob Reinhold's CC BY 4.0 Markdown was consulted as comparison lineage.

\section*{7. Radicial Groups of Height 1}
\phantomsection
\label{sga3:VIIA:7}
\addcontentsline{toc}{section}{7. Radicial Groups of Height 1}

% Authority: Exp7A local p. 45; combined p. 498; printed p. 487.
Let \(S\) be a scheme of characteristic
\(p>0\).\footnote{For the results of this section, one may also
consult [DG70], \S~II.7, nos.~3--4.}
We shall say that an \(\mathcal O_S\)-algebra \(\mathcal A\)
(respectively, an \(\mathcal O_S\)-\(p\)-Lie algebra
\(\mathcal L\)) is \emph{finite locally free} if its underlying
\(\mathcal O_S\)-module is locally free and of finite type. If
\(\mathcal L\) is a finite locally free
\(\mathcal O_S\)-\(p\)-Lie algebra, we know (cf.\ \SGARefLink{sga3:VIIA:5.5.2}{5.5.2}) that the
group \(S\)-functor
\[
  \operatorname{Spec}^{*}\mathcal U_p(\mathcal L)
\]
is represented by a finite locally free \(S\)-group scheme
\(\mathfrak G_p(\mathcal L)\). We shall see that this
\(S\)-group scheme solves a universal problem (\SGARefLink{sga3:VIIA:7.2}{7.2}), and we
shall characterize the \(S\)-group schemes of the form
\(\mathfrak G_p(\mathcal L)\) (\SGARefLink{sga3:VIIA:7.4}{7.4}).

\medskip
\noindent\textbf{Definition 7.0.}\footnote{The editors added
this definition (cf.\ [DG70], \S~II.4, \SGARefLink{sga3:VIIA:7.1}{7.1}), which will be used
in \SGARefLink{sga3:VIIA:7.2.1}{7.2.1}.}
\phantomsection
\label{sga3:VIIA:7.0}
\textit{Let \(H=\operatorname{Spec}\mathcal A\) be a finite
locally free \(S\)-group scheme. We say that \(H\) is
\emph{infinitesimal} if its unit section
\(\varepsilon_H:S\to H\) is a homeomorphism; equivalently, the
augmentation ideal of \(\mathcal A\) is locally nilpotent.}

\medskip
\noindent\textbf{7.1.}\footnote{The editors simplified the
original, taking into account the additions made in \SGARefLink{sga3:VIIA:5.5}{5.5}.}
\phantomsection
\label{sga3:VIIA:7.1}

Let \(\mathcal L\) be a finite locally free
\(\mathcal O_S\)-\(p\)-Lie algebra, and let
\(\mathfrak G_p(\mathcal L)\) be the affine \(S\)-group
\(\operatorname{Spec}\mathcal U_p(\mathcal L)\). By \SGARefLink{sga3:VIIA:5.5}{5.5},
\(\mathcal L\) is identified with
\(\operatorname{Lie}\mathfrak G_p(\mathcal L)\).

Now let \(G\) be a very good group \(S\)-functor satisfying
condition (F) of \SGARefLink{sga3:VIIA:6.3}{6.3}, and let
\[
  \varphi:\mathfrak G_p(\mathcal L)\longrightarrow G
\]
be a morphism of group \(S\)-functors. By \SGARefLink{sga3:VIIA:6.3}{6.3}, the morphism of
\(\mathcal O_S\)-Lie algebras
\[
  \operatorname{Lie}\varphi:
  \operatorname{Lie}\mathfrak G_p(\mathcal L)
  \longrightarrow
  \operatorname{Lie}G
\]
is compatible with the symbolic \(p\)-th power. If
\(\operatorname{Hom}_p(\mathcal L,\operatorname{Lie}G)\)
denotes the set of morphisms of \(\mathcal O_S\)-Lie algebras
compatible with the symbolic \(p\)-th power, we therefore have a
map
\[
  \operatorname{Lie}:
  \operatorname{Hom}_{S\text{-Gr.}}
    \bigl(\mathfrak G_p(\mathcal L),G\bigr)
  \longrightarrow
  \operatorname{Hom}_p
    \bigl(\mathcal L,\operatorname{Lie}G\bigr),
  \qquad
  \varphi\longmapsto\operatorname{Lie}\varphi.
\]

\medskip
\noindent\textbf{7.2. Theorem.}
\phantomsection
\label{sga3:VIIA:7.2}

\textit{If \(\mathcal L\) is a finite locally free
\(\mathcal O_S\)-\(p\)-Lie algebra, the map}
\[
  \operatorname{Hom}_{S\text{-gr.}}
    \bigl(\mathfrak G_p(\mathcal L),G\bigr)
  \longrightarrow
  \operatorname{Hom}_p
    \bigl(\mathcal L,\operatorname{Lie}G\bigr)
\]
\textit{is bijective in each of the following cases:}
\begin{enumerate}
  \item[\textit{(i)}]
  \textit{\(G\) is an \(S\)-group scheme;}
  \item[\textit{\SGAInlineRefTarget{sga3:VIIA:7.2:item:ii}(ii)}]
  \textit{\(G\) is of the form \(\operatorname{Aut}_S X\), where
  \(X\) is an \(S\)-scheme;}
  \item[\textit{(iii)}]
  \textit{\(G\) is of one of the forms
  \(\mathcal W(\mathcal F)\) or
  \(\operatorname{Aut}_{\mathcal O_S\text{-mod.}}
  \mathcal W(\mathcal F)\), where \(\mathcal F\) is a
  quasi-coherent \(\mathcal O_S\)-module.}
\end{enumerate}

The proof of the theorem rests on the following lemma.

\medskip
\noindent\textbf{Lemma.}
\phantomsection
\label{sga3:VIIA:7.2-frobenius-lemma}

\textit{If \(\mathcal L\) is a finite locally free
\(\mathcal O_S\)-\(p\)-Lie algebra, the \(S\)-group
\[
  G=\mathfrak G_p(\mathcal L)
\]
is annihilated by the Frobenius morphism
\(\operatorname{Fr}:G\to G^{(p)}\). In particular, \(G\) is
infinitesimal.}

% Authority: Exp7A local p. 46; combined p. 499; printed p. 488.
Indeed, let \(\mathcal U\) be the restricted enveloping algebra
of \(\mathcal L\), let \(\mathcal A=\mathcal U^{*}\) be the
affine algebra of \(G\), and let
\(\mathcal I=\ker\varepsilon_{\mathcal A}\) be the augmentation
ideal of \(\mathcal A\).\footnote{The editors supplied details
in what follows. For another proof, see [DG70], \S~II.7, 3.9.}
We have
\begin{equation}\SGAInlineEquationTarget{sga3:viia:paragraph:7-2:equation:1}
  \mathcal A
  =
  \mathcal I\oplus\eta_{\mathcal A}(\mathcal O_S),
  \tag{1}
  \label{sga3:VIIA:7.2-lemma-decomposition-A}
\end{equation}
where \(\eta_{\mathcal A}\) denotes the unit section of
\(\mathcal A\). Since \(\varepsilon_{\mathcal A}\)
(respectively, \(\eta_{\mathcal A}\)) is the transpose of
\(\eta_{\mathcal U}\) (respectively,
\(\varepsilon_{\mathcal U}\)), this decomposition corresponds
by duality to
\begin{equation}\SGAInlineEquationTarget{sga3:viia:paragraph:7-2:equation:2}
  \mathcal U
  =
  \mathcal J\oplus\eta_{\mathcal U}(\mathcal O_S),
  \tag{2}
  \label{sga3:VIIA:7.2-lemma-decomposition-U}
\end{equation}
where \(\mathcal J\) is the augmentation ideal of
\(\mathcal U\); thus \(\mathcal J=\mathcal I^{*}\).

Let \(\pi\) denote the endomorphism
\(x\mapsto x^p\) of \(\mathcal O_S\). We must show that
\(\operatorname{Fr}:G\to G^{(p)}\) factors through the unit
section of \(G^{(p)}\). By \SGARefLink{sga3:VIIA:4.1.4}{4.1.4}~(c), this is equivalent to
saying that the morphism
\[
  \Phi:
  \mathcal I\otimes_{\pi}\mathcal O_S
  \longrightarrow
  \mathcal A,
  \qquad
  a\otimes_{\pi}x\longmapsto a^p x,
\]
is zero. Since \(\mathcal A\) is finite locally free over
\(\mathcal O_S\), it suffices to show that the transpose
\({}^{t}\Phi\) is zero.

Now \(\Phi\) is precisely the following composite:
\begin{center}
  \phantomsection
  \label{sga3:VIIA:7.2-frobenius-phi-composite}
  % SGA3-VII-DIAG-064
  % BEGIN INCLUDED SGA3_English_current_buildable_source_20260823/source/inputs/VII/native_diagrams/exp7/viia_05_14/diagram_064.tex
\(\displaystyle
\mathcal J\otimes_{\pi}\mathcal O_S
  \xrightarrow{\tau}
\mathcal A\otimes_{\pi}\mathcal O_S
  \xrightarrow{j(\mathcal A)}
\mathrm S^p\mathcal A
  \xrightarrow{b(\mathcal A)}
\mathcal A ,
\)

% END INCLUDED SGA3_English_current_buildable_source_20260823/source/inputs/VII/native_diagrams/exp7/viia_05_14/diagram_064.tex

\end{center}
Here \(\tau\) is induced by the inclusion
\(\mathcal I\hookrightarrow\mathcal A\), while
\(b(\mathcal A)\) and \(j(\mathcal A)\) are defined as in
\SGARefLink{sga3:VIIA:4.3.3}{4.3.3}: \(b(\mathcal A)\) is induced by multiplication in
\(\mathcal A\), and \(j(\mathcal A)\) sends
\(a\otimes_{\pi}1\) to the image of
\(a\otimes\cdots\otimes a\) in
\(\mathrm S^p\mathcal A\). The
\(\mathcal O_S\)-dual of \(\mathrm S^p\mathcal A\) is the
submodule \(\Sigma^p\mathcal U\) of
\(\bigotimes^p\mathcal U\) consisting of the sections invariant
under the symmetric group of degree \(p\). Hence
\({}^{t}\Phi\) is the following composite:
\begin{center}
  \phantomsection
  \label{sga3:VIIA:7.2-transposed-phi-composite}
  % SGA3-VII-DIAG-065
  % BEGIN INCLUDED SGA3_English_current_buildable_source_20260823/source/inputs/VII/native_diagrams/exp7/viia_05_14/diagram_065.tex
\(\displaystyle
\mathcal U
  \xrightarrow{a(\mathcal U)}
\Sigma^p\mathcal U
  \xrightarrow{r(\mathcal U)}
\mathcal U\otimes_{\pi}\mathcal O_S
  \xrightarrow{q}
\mathcal J\otimes_{\pi}\mathcal O_S ,
\)

% END INCLUDED SGA3_English_current_buildable_source_20260823/source/inputs/VII/native_diagrams/exp7/viia_05_14/diagram_065.tex

\end{center}
Here \(q\) is induced by the projection
\(\mathcal U\to\mathcal J\) with kernel
\(\eta_{\mathcal U}(\mathcal O_S)\),
\(a(\mathcal U)\) is induced by the comultiplication of
\(\mathcal U\), and \(r(\mathcal U)\) vanishes on symmetrized
tensors and sends a section
\(x\otimes\cdots\otimes x\) to
\(x\otimes_{\pi}1\) (cf.\ \SGARefLink{sga3:VIIA:4.3.3}{4.3.3}).

It is clear that
\({}^{t}\Phi\circ\eta_{\mathcal U}=0\); by (2), it remains to
show that \({}^{t}\Phi\) annihilates the augmentation ideal
\(\mathcal J\). Since \({}^{t}\Phi\) is an algebra morphism and
\(\mathcal J\) is generated by \(\mathcal L\), identified with
its image in \(\mathcal U\), it suffices to show that
\({}^{t}\Phi(x)=0\) for every section \(x\) of
\(\mathcal L\). But
\[
  -a(\mathcal U)(x)=(p-1)!\,a(\mathcal U)(x)
\]
is the symmetrization of
\(x\otimes1\otimes\cdots\otimes1\), so its image under
\(r(\mathcal U)\) is zero. This proves the first assertion of
the lemma.

The second follows. Every local section of \(\mathcal I\) has
zero \(p\)-th power, and \(\mathcal I\) is an
\(\mathcal O_S\)-module of finite type; hence \(\mathcal I\) is
locally nilpotent. Explicitly, if \(V\) is an affine open subset
of \(S\) such that
\[
  I=\Gamma(V,\mathcal I)
\]
is generated by \(r\) elements, then
\[
  I^{r(p-1)+1}=0.
\]
Thus \(G_{\mathrm{red}}=S_{\mathrm{red}}\), and the unit section
\(\varepsilon_G:S\to G\) is a homeomorphism.

\medskip
\noindent\textbf{7.2.1.}\footnote{The editors supplied details
in what follows.}
\phantomsection
\label{sga3:VIIA:7.2.1}

We first prove assertion \SGARefLink{sga3:vii:paragraph:5-1:equation:ii}{(ii)} of \SGARefLink{sga3:VIIA:7.2}{Theorem 7.2}. Let
\(\pi:X\to S\) be an \(S\)-scheme, and first consider an
arbitrary infinitesimal \(S\)-group \(H\).

% Authority: Exp7A local p. 47; combined p. 500; printed p. 489.
Morphisms \(\varphi:H\to\operatorname{Aut}X\) correspond
bijectively to left actions
\(\mu:H\times X\to X\) of \(H\) on \(X\). For such an action,
if \(\varepsilon\) is the unit section of \(H\), the composite
\begin{center}
  \phantomsection
  \label{sga3:VIIA:7.2.1-action-unit-composite}
  % SGA3-VII-DIAG-066
  % BEGIN INCLUDED SGA3_English_current_buildable_source_20260823/source/inputs/VII/native_diagrams/exp7/viia_05_14/diagram_066.tex
\(\displaystyle
X \simeq S\times X
  \xrightarrow{\varepsilon\times X}
H\times X
  \xrightarrow{\mu}
X
\)

% END INCLUDED SGA3_English_current_buildable_source_20260823/source/inputs/VII/native_diagrams/exp7/viia_05_14/diagram_066.tex

\end{center}
must be the identity. Since
\((H\times X)_{\mathrm{red}}\) is identified with
\(X_{\mathrm{red}}\), the morphism \(\mu\) must induce the
identity on the associated reduced schemes. In particular,
\(\mu\) induces an action of \(H\) on every open subset of
\(X\). Thus, for every open \(U\subset X\) affine over \(S\), we
obtain a morphism of unital associative algebras
\[
  \mathcal A(U)
  \longrightarrow
  \mathcal A(H)\otimes\mathcal A(U)
\]
which makes \(\mathcal A(U)\) a left
\(\mathcal A(H)\)-comodule, in such a way that the restriction
maps
\[
  \mathcal A(U)\longrightarrow\mathcal A(U'),
  \qquad U'\subset U,
\]
are comodule morphisms. Conversely, every datum of this kind
comes from a unique left action \(\mu:H\times X\to X\).

We also have the following lemma.

\medskip
\noindent\textbf{Lemma.}
\phantomsection
\label{sga3:VIIA:7.2.1-action-representation-lemma}

\textit{Let \(X=\operatorname{Spec}\mathcal C\) be an affine
\(S\)-scheme, let \(H=\operatorname{Spec}\mathcal A\) be an
infinitesimal \(S\)-group, and put
\[
  \mathcal U
  =
  \mathcal A^{*}
  =
  \operatorname{Hom}_{\mathcal O_S}
    (\mathcal A,\mathcal O_S).
\]
Left actions of \(H\) on \(X\) correspond bijectively to
representations of the algebra \(\mathcal U\) in the
\(\mathcal O_S\)-module \(\mathcal C\) such that}
\begin{align*}
  \textit{(a)}\quad&
  u(1_{\mathcal C})
  =
  \varepsilon(u)\,1_{\mathcal C},
  \\
  \textit{(b)}\quad&
  u(xy)
  =
  \sum_i v_i(x)w_i(y)
  \quad\textit{if}\quad
  \Delta u
  =
  \sum_i v_i\otimes w_i.
\end{align*}
\textit{Here \(u\) is an arbitrary section of \(\mathcal U\)
over an affine open subset \(V\) of \(S\), while \(x\) and \(y\)
are sections of \(\mathcal C\) over \(V\). The symbol
\(1_{\mathcal C}\) denotes the unit section of \(\mathcal C\),
and \(\varepsilon\) and \(\Delta\) denote the augmentation and
the diagonal morphism of \(\mathcal U\), respectively.}

Indeed, a left action \(\mu\) of \(H\) on \(X\) is defined by a
morphism of unital associative algebras
\[
  \lambda:\mathcal C\longrightarrow\mathcal A\otimes\mathcal C
\]
which makes \(\mathcal C\) a left
\(\mathcal A\)-comodule. Let \(\alpha\) be the composite
\begin{center}
  \phantomsection
  \label{sga3:VIIA:7.2.1-action-contraction-composite}
  % SGA3-VII-DIAG-067
  % BEGIN INCLUDED SGA3_English_current_buildable_source_20260823/source/inputs/VII/native_diagrams/exp7/viia_05_14/diagram_067.tex
\resizebox{0.98\linewidth}{!}{%
  \(\displaystyle
  \mathcal U\otimes_{\mathcal O_S}\mathcal C
    \xrightarrow{\mathcal U\otimes\lambda}
  \mathcal U\otimes_{\mathcal O_S}\mathcal A\otimes_{\mathcal O_S}\mathcal C
    \xrightarrow{\gamma\otimes\mathcal C}
  \mathcal O_S\otimes_{\mathcal O_S}\mathcal C
    \simeq \mathcal C
  \)%
}

% END INCLUDED SGA3_English_current_buildable_source_20260823/source/inputs/VII/native_diagrams/exp7/viia_05_14/diagram_067.tex

\end{center}
where \(\gamma\) is the ``contraction'' from
\(\mathcal A^{*}\otimes_{\mathcal O_S}\mathcal A\) to
\(\mathcal O_S\). Since \(\mathcal A\) is finite locally free
over \(\mathcal O_S\), the map
\[
  \lambda
  \longmapsto
  (\gamma\otimes\mathcal C)(\mathcal U\otimes\lambda)
\]
is a bijection from
\[
  \operatorname{Hom}_{\mathcal O_S}
    (\mathcal C,\mathcal A\otimes\mathcal C)
  \quad\text{onto}\quad
  \operatorname{Hom}_{\mathcal O_S}
    (\mathcal U\otimes\mathcal C,\mathcal C).
\]
Moreover, by duality, the condition that \(\lambda\) define a
left \(\mathcal A\)-comodule structure (respectively, be a
morphism of unital associative algebras) is equivalent to the
condition that \(\alpha\) be a representation of
\(\mathcal U\) in \(\mathcal C\) (respectively, satisfy
conditions (a) and (b)). This proves the lemma.

For any representation of \(\mathcal U\) in the
\(\mathcal O_S\)-module \(\mathcal C\), it is also clear that
the sections \(u\) of \(\mathcal U\) satisfying conditions (a)
and (b) form a subalgebra of \(\mathcal U\).

In the particular case
\[
  H=\mathfrak G_p(\mathcal L)
\]
which concerns us, these conditions will therefore hold for
every section \(u\) of
\(\mathcal U=\mathcal U_p(\mathcal L)\) if they hold for the
sections \(u\) of \(\mathcal L\), identified with its image in
\(\mathcal U_p(\mathcal L)\).

% Authority: Exp7A local p. 48; combined p. 501; printed p. 490.
If \(u\) is a section of \(\mathcal L\), conditions (a) and (b)
simply mean
\[
  u(1_{\mathcal C})=0,
  \qquad
  u(xy)=u(x)y+xu(y);
\]
that is, \(\alpha(u)\) is an \(\mathcal O_S\)-derivation of
\(\mathcal C=\mathcal A(X)\). Assertion \SGARefLink{sga3:VIIA:7.2:item:ii}{(ii)} of \SGARefLink{sga3:VIIA:7.2}{7.2} follows.
Indeed, every morphism
\[
  \varphi:\mathfrak G_p=\mathfrak G_p(\mathcal L)
  \longrightarrow\operatorname{Aut}X
\]
defines a morphism of \(p\)-Lie algebras
\(\operatorname{Lie}\varphi\) from
\(\mathcal L=\operatorname{Lie}\mathfrak G_p\) to
\[
  \pi_*
  \bigl(
    \operatorname{Der}_{\mathcal O_S}\mathcal O_X
  \bigr),
\]
and conversely every datum of this kind comes, by what precedes,
from a unique action
\(\mu:\mathfrak G_p(\mathcal L)\times X\to X\).

\medskip
\noindent\textbf{7.2.2.}
\phantomsection
\label{sga3:VIIA:7.2.2}

We now show how assertion \SGARefLink{sga3:vii:paragraph:5-1:equation:i}{(i)} of \SGARefLink{sga3:VIIA:7.2}{Theorem 7.2} follows from \SGARefLink{sga3:vii:paragraph:5-1:equation:ii}{(ii)}.
Let \(G\) be an \(S\)-group scheme. If \(T\) is an
\(S\)-scheme and \(x\in G(T)\), let
\(\ell_x^T\) (respectively, \(r_x^T\)) denote the left
(respectively, right) translation of \(G_T\) defined by \(x\).
The maps
\[
  \ell^T:x\longmapsto\ell_x^T
\]
therefore determine a homomorphism
\(\ell:G\to\operatorname{Aut}G\).

On the other hand, let \(f\) be a \(T\)-automorphism of \(G_T\),
and define
\[
  {}^x f=(r_x^T)^{-1}fr_x^T;
\]
thus, for every \(T'\to T\) and \(g\in G(T')\),
\[
  ({}^x f)(g)=f(gx)x^{-1}.
\]
In this way \(G\) acts on the left on the group \(S\)-functor
\(\operatorname{Aut}G\), and hence also on the functors
\[
  T\longmapsto
  \operatorname{Hom}_{T\text{-Gr.}}
    \bigl(\mathfrak G_p(\mathcal L_T),
          \operatorname{Aut}G_T\bigr)
\]
and
\[
  T\longmapsto
  \operatorname{Hom}_p
    \bigl(
      \mathcal L_T,
      \operatorname{Lie}(\operatorname{Aut}G_T/T)
    \bigr).
\]

The morphism
\(\ell^T:G_T\to\operatorname{Aut}G_T\) identifies \(G_T\) with
the group of automorphisms of the \(T\)-scheme \(G_T\) which
commute with right translations. Its derived morphism
\(\operatorname{Lie}(\ell^T)\) identifies
\(\operatorname{Lie}G_T\) with the \(p\)-Lie algebra of
\(\mathcal O_T\)-derivations of \(\mathcal O_{G_T}\) which
commute with right translations (cf.\ II, \SGARefLink{sga3:ii:scholium:4-11-1}{4.11.1}). They
therefore induce the commutative squares
\begin{center}
  \phantomsection
  \label{sga3:VIIA:7.2.2-lie-invariant-commutative-squares}
  % SGA3-VII-DIAG-068
  % BEGIN INCLUDED SGA3_English_current_buildable_source_20260823/source/inputs/VII/native_diagrams/exp7/viia_05_14/diagram_068.tex
\resizebox{0.98\linewidth}{!}{%
\begin{tikzcd}[ampersand replacement=\&,row sep=3.5em,column sep=4.5em]
\operatorname{Hom}_{T\text{-}\mathrm{Gr.}}
  \bigl(\mathfrak G_p(\mathcal L_T),G_T\bigr)
  \arrow[r,"\operatorname{Lie}"]
  \arrow[d,swap,"\ell_T"]
  \&
\operatorname{Hom}_p
  \bigl(\mathcal L_T,\operatorname{Lie}(G_T/T)\bigr)
  \arrow[d,"\operatorname{Lie}\ell_T"]
  \\
\operatorname{Hom}_{T\text{-}\mathrm{Gr.}}
  \bigl(\mathfrak G_p(\mathcal L_T),\underline{\operatorname{Aut}}\,G_T\bigr)
  \arrow[r,"\operatorname{Lie}"]
  \&
\operatorname{Hom}_p
  \bigl(\mathcal L_T,
    \operatorname{Lie}(\underline{\operatorname{Aut}}\,G_T/T)\bigr) .
\end{tikzcd}%
}

% END INCLUDED SGA3_English_current_buildable_source_20260823/source/inputs/VII/native_diagrams/exp7/viia_05_14/diagram_068.tex

\end{center}
The images of the two vertical arrows are the subfunctors
formed by the invariants under the action of the \(S\)-group
\(G\). Since the lower horizontal arrow is invertible by \SGARefLink{sga3:VIIA:7.2.1}{7.2.1}
and is compatible with the action of \(G\), the upper
horizontal arrow is invertible as well. This proves \SGARefLink{sga3:VIIA:7.2}{7.2}~\SGARefLink{sga3:vii:paragraph:5-1:equation:i}{(i)}.

\medskip
\noindent\textbf{7.2.3.}
\phantomsection
\label{sga3:VIIA:7.2.3}

Consider the case in which
\(
  G=
  \operatorname{Aut}_{\mathcal O_S\text{-mod.}}
    \mathcal W(\mathcal F)
\).\footnote{In what follows, the editors expanded and
simplified the original, taking VIB, \SGARefLink{prop:sga3-VIB-11.6.1}{11.6.1} into account.}
Put
\[
  \mathcal U=\mathcal U_p(\mathcal L),
  \qquad
  \mathcal A=\mathcal U^{*},
  \qquad
  H=\mathfrak G_p(\mathcal L)=\operatorname{Spec}\mathcal A.
\]
Since \(H\) is affine over \(S\), VIB, \SGARefLink{prop:sga3-VIB-11.6.1}{11.6.1} shows that a
morphism of \(S\)-groups from \(H\) to
\(\operatorname{Aut}_{\mathcal O_S\text{-mod.}}
\mathcal W(\mathcal F)\) is the same as a right
\(\mathcal A\)-comodule structure
\[
  \mu:\mathcal F\longrightarrow\mathcal F\otimes\mathcal A.
\]
Since \(\mathcal A\) is finite locally free over
\(\mathcal O_S\), this is in turn equivalent to giving a
representation
\[
  \alpha:
  \mathcal U
  \longrightarrow
  \operatorname{End}_{\mathcal O_S}(\mathcal F)
\]
of \(\mathcal U\) in \(\mathcal F\). Finally, by the universal
property of \(\mathcal U=\mathcal U_p(\mathcal L)\), giving such
a morphism \(\alpha\) is equivalent to giving its restriction
\(\rho\) to \(\mathcal L\), identified with its image in
\(\mathcal U\), and \(\rho\) is a morphism of \(p\)-Lie
algebras from \(\mathcal L\) to
\(\operatorname{End}_{\mathcal O_S}(\mathcal F)\).

% Authority: Exp7A local p. 49; combined p. 502; printed p. 491.
Finally, consider the case \(G=\mathcal W(\mathcal F)\), with
the preceding notation.\footnote{The editors added what
follows. The original said that ``the case of
\(\mathcal W(\mathcal F)\) is analogous.''}
To give a morphism of \(S\)-functors
\(\varphi:H\to\mathcal W(\mathcal F)\) is first of all
equivalent to giving an element \(\theta\) of
\(\Gamma(H,\mathcal F\otimes\mathcal O_H)\). Since \(H\) is
finite locally free over \(S\), we have
\[
  \Gamma(H,\mathcal F\otimes\mathcal O_H)
  =
  \Gamma(S,\mathcal F\otimes\mathcal A)
  =
  \operatorname{Hom}_{\mathcal O_S}(\mathcal U,\mathcal F).
\]
The condition that \(\varphi\) be a group morphism means that
\(\theta\), regarded as a morphism of \(\mathcal O_S\)-modules
\(\mathcal U\to\mathcal F\), vanishes on
\(1_{\mathcal U}\) and on \(\mathcal J^2\), where
\(\mathcal J\) is the augmentation ideal of \(\mathcal U\).
Consequently,
\begin{equation}
  \operatorname{Hom}_{S\text{-gr.}}
    \bigl(H,\mathcal W(\mathcal F)\bigr)
  =
  \operatorname{Hom}_{\mathcal O_S}
    \bigl(\mathcal J/\mathcal J^2,\mathcal F\bigr).
  \tag{1}
  \label{sga3:VIIA:7.2.3-hom-group-additive}
\end{equation}

Next consider the quasi-coherent sheaf
\([\mathcal L,\mathcal L]\), the image of the morphism
\[
  \mathcal L\otimes\mathcal L
  \longrightarrow
  \mathcal L,
  \qquad
  x\otimes y\longmapsto[x,y].
\]
For every affine open subset \(V\) of \(S\), we have
\[
  [\mathcal L,\mathcal L](V)
  =
  [\mathcal L(V),\mathcal L(V)].
\]
There is an exact sequence
\begin{center}
  \phantomsection
  \label{sga3:VIIA:7.2.3-augmentation-quotient-exact-sequence}
  % SGA3-VII-DIAG-069
  % BEGIN INCLUDED SGA3_English_current_buildable_source_20260823/source/inputs/VII/native_diagrams/exp7/viia_05_14/diagram_069.tex
\(\displaystyle
(\dagger)\qquad
0 \longrightarrow [\mathcal L,\mathcal L]
  \longrightarrow \mathcal L
  \xrightarrow{\pi}
\mathcal J/\mathcal J^2
  \longrightarrow 0 ,
\)

% END INCLUDED SGA3_English_current_buildable_source_20260823/source/inputs/VII/native_diagrams/exp7/viia_05_14/diagram_069.tex

\end{center}
in which \(\pi\) is the composite of the inclusion
\(\mathcal L\hookrightarrow\mathcal J\) and the projection
\(\mathcal J\to\mathcal J/\mathcal J^2\).

Indeed, the question is local on \(S\), so we may suppose that
\(S\) is affine with ring \(R\), and that
\[
  L=\mathcal L(S)
\]
is free with basis \((x_1,\ldots,x_r)\). Identify \(L\) with
its image in \(U=U_p(L)\), and let \(K\) be the sub-\(R\)-module
of \(U\) which is the direct sum of \([L,L]\) and the submodule
with basis the monomials
\[
  x_1^{n_1}\cdots x_r^{n_r}
  \quad\text{such that}\quad
  n_1+\cdots+n_r\geq2.
\]
One checks that \(K\) is a two-sided ideal of \(U\). Since
\(K\) is contained in \(J^2\), where \(J\) is the augmentation
ideal of \(U\), and contains all the products \(x_ix_j\), which
generate \(J^2\), it follows that \(K=J^2\). Hence
\[
  J^2\cap L=[L,L],
\]
and the displayed sequence is exact.

By \SGARefLink{sga3:VIIA:6.4.2}{6.4.2},
\(\operatorname{Lie}\mathcal W(\mathcal F)\) is simply
\(\mathcal W(\mathcal F)\), with zero Lie bracket and zero
symbolic \(p\)-th power. From this and the preceding discussion,
we obtain
\begin{equation}
  \operatorname{Hom}_p(\mathcal L,\mathcal F)
  =
  \operatorname{Hom}_{\mathcal O_S}
    \bigl(\mathcal L/[\mathcal L,\mathcal L],\mathcal F\bigr)
  =
  \operatorname{Hom}_{\mathcal O_S}
    \bigl(\mathcal J/\mathcal J^2,\mathcal F\bigr).
  \tag{2}
  \label{sga3:VIIA:7.2.3-hom-p-additive}
\end{equation}
Together with \SGARefLink{sga3:VIIA:7.2.3-hom-group-additive}{(1)}, this completes the proof of \SGARefLink{sga3:VIIA:7.2}{Theorem 7.2}.

\medskip
\noindent\textbf{7.3. Lemma.}
\phantomsection
\label{sga3:VIIA:7.3}

\textit{If \(\mathcal L\) is a finite locally free
\(\mathcal O_S\)-\(p\)-Lie algebra, the morphism}
\[
  j_{\mathcal L}:
  \mathcal L
  \longrightarrow
  \operatorname{Lie}\mathfrak G_p(\mathcal L)
\]
\textit{of \SGARefLink{sga3:VIIA:5.5}{5.5} is invertible.}\footnote{To avoid changing the
numbering, the editors retained statement \SGARefLink{sga3:VIIA:7.3}{7.3}, although they
included it, together with its proof, in \SGARefLink{sga3:VIIA:5.5.1}{5.5.1}.}

For the proof, see \SGARefLink{sga3:VIIA:5.5.1}{5.5.1}.

\medskip
\noindent\textbf{7.4.}
\phantomsection
\label{sga3:VIIA:7.4}

To conclude this section, we give a characterization of the
\(S\)-group schemes of the form
\(\mathfrak G_p(\mathcal L)\), where \(\mathcal L\) is a finite
locally free \(\mathcal O_S\)-\(p\)-Lie algebra.

Let \(G\) be an \(S\)-group scheme, let \(\varepsilon_G\) be its
unit section, and let \(\mathcal I'\) be the kernel of the
morphism
\[
  \varepsilon_G^{-1}(\mathcal O_G)
  \longrightarrow
  \mathcal O_S
\]
corresponding to \(\varepsilon_G\). By \SGARefLink{sga3:VIIB:2.5:proposition}{2.5} and \SGARefLink{sga3:VIIA:1.3.1}{1.3.1}, the image
of \(\operatorname{Lie}(G/S)(S)\) in \(\mathrm U(G)\) is
identified with the morphisms of \(\mathcal O_S\)-modules from
\(\varepsilon_G^{-1}(\mathcal O_G)\) to \(\mathcal O_S\) which
vanish on the unit section of
\(\varepsilon_G^{-1}(\mathcal O_G)\) and on
\(\mathcal I'^{\,2}\).

% Authority: Exp7A local p. 50; combined p. 503; printed p. 492.
We thereby recover the canonical isomorphism
\[
  \operatorname{Lie}(G/S)(S)
  \xrightarrow{\sim}
  \operatorname{Hom}_{\mathcal O_S}
    (\mathcal I'/\mathcal I'^{\,2},\mathcal O_S)
\]
of II, \SGARefLink{sga3:VIIA:3.3}{3.3} and \SGARefLink{sga3:ii:remark:4-11-4}{4.11.4}.\footnote{If \(G\) is affine over \(S\)
and \(\mathcal I\) denotes the augmentation ideal of
\(\mathcal A(G)\), then
\(\mathcal I'/\mathcal I'^{\,2}\) is identified with
\(\varepsilon_G^{*}(\mathcal I/\mathcal I^2)\); cf.\ loc. cit.}
As in the cited passage, put
\[
  \omega_{G/S}=\mathcal I'/\mathcal I'^{\,2}.
\]
Thus the sheaf \(\operatorname{Lie}(G/S)\) is identified with
\[
  \omega_{G/S}^{*}
  =
  \operatorname{Hom}_{\mathcal O_S}
    (\omega_{G/S},\mathcal O_S).
\]
Moreover,\footnote{The editors added the following sentence.}
if \(G=\mathfrak G_p(\mathcal L)\), where
\(\mathcal L\) is a finite locally free
\(\mathcal O_S\)-\(p\)-Lie algebra, then \SGARefLink{sga3:VIIA:5.5.3}{5.5.3} gives
\[
  \omega_{G/S}=\mathcal L^{*}.
\]

\medskip
\noindent\textbf{Theorem.}
\phantomsection
\label{sga3:VIIA:7.4-theorem}

\textit{If \(G\) is a group scheme over a scheme \(S\) of
characteristic \(p>0\), the following assertions are
equivalent:}
\begin{enumerate}
  \item[\textit{\SGAInlineRefTarget{sga3:VIIA:7.4-theorem:item:i}(i)}]
  \textit{There exists a finite locally free
  \(\mathcal O_S\)-\(p\)-Lie algebra \(\mathcal L\) such that
  \(G\simeq\mathfrak G_p(\mathcal L)\).}

  \item[\textit{(i\(^{\prime}\))}]
  \textit{The \(\mathcal O_S\)-\(p\)-Lie algebra
  \(\operatorname{Lie}(G)\) is finite locally free, and
  \(G\simeq\mathfrak G_p(\operatorname{Lie}(G))\).}

  \item[\textit{(ii)}]
  \textit{\(G\) is affine over \(S\), the
  \(\mathcal O_S\)-module \(\omega_{G/S}\) is locally free of
  finite type, and the affine algebra of \(G\) is locally
  isomorphic to the quotient of the symmetric algebra
  \(\mathrm S_{\mathcal O_S}(\omega_{G/S})\) by the ideal
  generated by the \(p\)-th powers of the sections of
  \(\omega_{G/S}\).}

  \item[\textit{\SGAInlineRefTarget{sga3:VIIA:7.4-theorem:item:iii}(iii)}]
  \textit{\(G\) is locally of finite presentation over \(S\),
  is of height \(\leq1\), and \(\omega_{G/S}\) is locally free.}

  \item[\textit{(iii\(^{\prime}\))}]
  \textit{\(G\) is locally of finite type over \(S\), is of
  height \(\leq1\), and \(\omega_{G/S}\) is locally free.}

  \item[\textit{\SGAInlineRefTarget{sga3:VIIA:7.4-theorem:item:iv}(iv)}]
  \textit{\(G\) is locally of finite presentation and flat over
  \(S\), and is of height \(\leq1\).}\footnote{The editors added
  assertion (i\(^{\prime}\)), implicit in the original, and
  assertions (iii\(^{\prime}\)) and \SGARefLink{sga3:VIIA:7.4-theorem:item:iv}{(iv)}, pointed out by
  O.~Gabber. Assertion \SGARefLink{sga3:VIIA:7.4-theorem:item:iv}{(iv)} takes up a footnote in the original,
  which stated: ``The condition on \(\omega_{G/S}\) is in fact
  unnecessary, as is easily seen by reducing to the case where
  \(S\) is local with residue field \(k\), and applying the
  theorem to \(G_k\).'' As Gabber pointed out, this is incorrect
  without a flatness hypothesis. Indeed, let \(A\) be an
  Artinian local ring of characteristic \(p>0\), let \(J\) be a
  proper nonzero ideal of \(A\), and let \(H\) be the subgroup
  \[
    \operatorname{Spec}A[x]/(x^p,Jx)
  \]
  of \(\alpha_{p,A}\); that is, for every \(A\)-algebra \(R\),
  \[
    H(R)=\{x\in R\mid x^p=0\ \text{and}\ Jx=0\}.
  \]
  Then \(H\) is not flat over \(A\), and hence is not of the
  form \(\mathfrak G_p(\mathcal L)\) with \(\mathcal L\) a
  finite-rank free \(p\)-Lie algebra over \(A\).}
\end{enumerate}

\medskip
\noindent\textbf{7.4.1.}
\phantomsection
\label{sga3:VIIA:7.4.1}

The equivalence \SGARefLink{sga3:vii:paragraph:5-1:equation:i}{(i)} \(\Leftrightarrow\)
(i\(^{\prime}\)) follows from \SGARefLink{sga3:VIIA:5.5.3}{5.5.3}~\SGARefLink{sga3:vii:paragraph:5-1:equation:i}{(i)}, and the implications
\SGARefLink{sga3:VIIA:5.5.3:item:ii}{(ii)} \(\Rightarrow\) \SGARefLink{sga3:VIIA:7.4-theorem:item:iii}{(iii)} \(\Rightarrow\)
(iii\(^{\prime}\)) are clear. We also have
\SGARefLink{sga3:VIIA:5.5.3:item:i}{(i)} \(\Rightarrow\) \SGARefLink{sga3:VIIA:7.4-theorem:item:iv}{(iv)}, since
\(\mathfrak G_p(\mathcal L)\) is finite locally free and of
height \(\leq1\), by \SGARefLink{sga3:VIIA:5.5.2}{5.5.2} and the lemma in \SGARefLink{sga3:VIIA:7.2}{7.2}. We show that
\SGARefLink{sga3:vii:paragraph:5-1:equation:i}{(i)} implies \SGARefLink{sga3:vii:paragraph:5-1:equation:ii}{(ii)}.

Let \(\mathcal I\) be the augmentation ideal of
\[
  \mathcal A=\mathcal U_p(\mathcal L)^{*}.
\]
We saw in \SGARefLink{sga3:VIIA:5.5.3}{5.5.3}~\SGARefLink{sga3:vii:paragraph:5-1:equation:ii}{(ii)} that
\(\omega_{G/S}=\mathcal I/\mathcal I^2\) is identified with
\(\mathcal L^{*}\), and hence is finite locally free.

Now suppose that \(S\) is affine. There is a section
\[
  \sigma:\omega_{G/S}\longrightarrow\mathcal I
\]
of the projection
\(\mathcal I\to\mathcal I/\mathcal I^2\). It induces an algebra
morphism
\[
  \sigma':
  \mathrm S_{\mathcal O_S}(\omega_{G/S})
  \longrightarrow
  \mathcal A,
\]
and, by the lemma in \SGARefLink{sga3:VIIA:7.2}{7.2}, \(\sigma'\) factors through a morphism
\[
  \varphi:
  \mathrm S_{\mathcal O_S}(\omega_{G/S})/\mathcal K
  \longrightarrow
  \mathcal A,
\]
where \(\mathcal K\) denotes the ideal generated by the
\(p\)-th powers of sections of \(\omega_{G/S}\). Filter
\(\mathcal A\) by powers of \(\mathcal I\), and filter
\(\mathrm S_{\mathcal O_S}(\omega_{G/S})/\mathcal K\) by
powers of the ideal generated by \(\omega_{G/S}\). The morphism
\(\varphi\) plainly induces an epimorphism on the associated
graded modules. Hence \(\varphi\) is an epimorphism between
locally free \(\mathcal O_S\)-modules of the same rank
(cf.\ \SGARefLink{sga3:VIIA:5.3.3}{5.3.3}), and is therefore an isomorphism. This proves that
\SGARefLink{sga3:vii:paragraph:5-1:equation:i}{(i)} implies \SGARefLink{sga3:vii:paragraph:5-1:equation:ii}{(ii)}.

\medskip
\noindent\textbf{7.4.2.}
\phantomsection
\label{sga3:VIIA:7.4.2}

% Authority: Exp7A local p. 51; combined p. 504; printed p. 493.
Suppose now that \(G\) is of height \(\leq1\) and locally of
finite presentation over \(S\).\footnote{The editors supplied
details and simplified the original in what follows.}
The Frobenius morphism
\(\operatorname{Fr}:G\to G^{(p)}\) is integral and factors
through the unit section of \(G^{(p)}\); consequently, \(G\) is
integral, and hence affine, over \(S\). Put
\(\mathcal A=\mathcal A(G)\). Since \(G\) is locally of finite
presentation over \(S\), it follows that \(G\) is finite and of
finite presentation over \(S\), and hence that
\(\mathcal A(G)\) is an \(\mathcal O_S\)-module of finite
presentation (cf.\ EGA IV\(_1\), 1.4.7).

Let \(\mathcal I\) be the augmentation ideal of
\(\mathcal A\). Since
\[
  \mathcal A
  =
  \eta_{\mathcal A}(\mathcal O_S)\oplus\mathcal I,
\]
where \(\eta_{\mathcal A}\) is the unit section of
\(\mathcal A\), the module \(\mathcal I\) is of finite
presentation over \(\mathcal O_S\), as is
\(\omega_G=\mathcal I/\mathcal I^2\). If instead \(G\) is
assumed to be of height \(\leq1\) and locally of finite type
over \(S\), the same argument shows that
\(\mathcal A(G)\), \(\mathcal I\), and
\(\omega_G=\mathcal I/\mathcal I^2\) are
\(\mathcal O_S\)-modules of finite type.

Thus, under hypothesis (iii\(^{\prime}\)),
\(\omega_{G/S}\) is finite locally free over
\(\mathcal O_S\), as is
\[
  \mathcal L
  =
  \operatorname{Lie}(G/S)
  =
  \omega_{G/S}^{*}.
\]
Put
\[
  \mathcal B=\mathcal U_p(\mathcal L)^{*},
  \qquad
  H=\mathfrak G_p(\mathcal L)=\operatorname{Spec}\mathcal B.
\]
By \SGARefLink{sga3:VIIA:7.2}{Theorem 7.2}, the identity map of \(\mathcal L\) corresponds
to a group morphism from
\(H=\mathfrak G_p(\mathcal L)\) to \(G\), and hence to a
morphism of \(\mathcal O_S\)-algebras
\[
  \theta:\mathcal A\longrightarrow\mathcal B.
\]
By definition, \(\theta\) induces an isomorphism from
\(\omega_{G/S}\) to \(\omega_{H/S}\); we must show that
\(\theta\) itself is an isomorphism.

We may restrict to the case in which \(S\) is affine. There is
then a section \(\tau\) of the projection
\(\mathcal I\to\omega_{G/S}\), inducing an algebra morphism
\[
  \tau':
  \mathrm S_{\mathcal O_S}(\omega_{G/S})
  \longrightarrow
  \mathcal A.
\]
Every local section of \(\mathcal I\) has zero \(p\)-th power,
because \(\operatorname{Fr}:G\to G^{(p)}\) factors through the
unit section of \(G^{(p)}\). Thus \(\tau'\) induces a morphism
of \(\mathcal O_S\)-algebras \(\psi\) fitting into the
commutative diagram
\begin{center}
  \phantomsection
  \label{sga3:VIIA:7.4.2-height-one-characterization-triangle}
  % SGA3-VII-DIAG-070
  % BEGIN INCLUDED SGA3_English_current_buildable_source_20260823/source/inputs/VII/native_diagrams/exp7/viia_05_14/diagram_070.tex
\begin{tikzcd}[row sep=3.5em,column sep=4.6em]
\mathrm S_{\mathcal O_S}(\omega_{G/S})/\mathcal K
  \arrow[r,"\psi"]
  \arrow[dr,swap,"\varphi"]
  & \mathcal A \arrow[d,"\theta"] \\
& \mathcal B
\end{tikzcd}

% END INCLUDED SGA3_English_current_buildable_source_20260823/source/inputs/VII/native_diagrams/exp7/viia_05_14/diagram_070.tex

\end{center}
where \(\mathcal K\) is the ideal generated by the \(p\)-th
powers of sections of \(\omega_{G/S}\). As in \SGARefLink{sga3:VIIA:7.4.1}{7.4.1},
\(\psi\) is an epimorphism of \(\mathcal O_S\)-modules. On the
other hand, \SGARefLink{sga3:VIIA:7.4.1}{7.4.1} shows that
\(\varphi=\theta\circ\psi\) is an isomorphism. Hence
\(\theta\) is an isomorphism as well. This proves
(iii\(^{\prime}\)) \(\Rightarrow\) \SGARefLink{sga3:VIIA:7.4-theorem:item:i}{(i)}.

\medskip
\noindent\textbf{7.4.3.}\footnote{The editors added this
paragraph to prove \SGARefLink{sga3:VIIA:7.4-theorem:item:iv}{(iv)} \(\Rightarrow\) \SGARefLink{sga3:VIIA:7.4-theorem:item:iii}{(iii)}; cf.\ editor's
note~\SGARefLink{sga3:VIIB:nde:71}{(71)}.}
\phantomsection
\label{sga3:VIIA:7.4.3}

We finally prove that \SGARefLink{sga3:VIIA:7.4-theorem:item:iv}{(iv)} implies \SGARefLink{sga3:VIIA:7.4-theorem:item:iii}{(iii)}. It suffices to show
that \(\omega_{G/S}\) is locally free, so we may suppose that
\(S\) is affine with ring \(R\). As observed at the beginning
of \SGARefLink{sga3:VIIA:7.4.2}{7.4.2}, hypothesis \SGARefLink{sga3:VIIA:7.4-theorem:item:iv}{(iv)} then implies that
\(G=\operatorname{Spec}A\), where \(A\) is an \(R\)-algebra
which is an \(R\)-module of finite presentation, as is
\[
  \omega_{G/S}=I/I^2,
\]
where \(I\) is the augmentation ideal of \(A\). Since \(G\) is
also assumed flat over \(S\), the algebra \(A\) is a finite
locally free \(R\)-module (cf.\ [BAC] II, \S~\SGARefLink{sga3:VIIB:5.2:theorem}{5.2}, Theorem~1 and
Corollary~2). By the cited result, to show that
\(\omega_{G/S}\) is locally free of finite rank, it suffices to
show that
\[
  (\omega_{G/S})_{\mathfrak m}
\]
is flat for every maximal ideal \(\mathfrak m\) of \(R\).
Thus we may suppose that \(R\) is local and that \(A\) is free
of rank \(n+1\); then \(I\) is free of rank \(n\). Let
\(\mathfrak m\) be the maximal ideal of \(R\), and put
\[
  k=R/\mathfrak m.
\]

% Authority: Exp7A local p. 52; combined p. 505; printed p. 494.
Let \(I_k\) be the augmentation ideal of \(A_k\), and let \(r\)
be the dimension of
\[
  \omega_{G_k/k}=I_k/I_k^2.
\]
Choose a basis \((e_1,\ldots,e_n)\) of \(I_k\) such that
\((e_{r+1},\ldots,e_n)\) is a basis of \(I_k^2\), and choose
elements \(x_1,\ldots,x_n\) of \(I\) lifting the \(e_i\).
Nakayama's lemma shows that \((x_1,\ldots,x_n)\) is a basis of
\(I\) over \(R\).

Let \(N\) be the sub-\(R\)-module of \(I\) with basis
\((x_1,\ldots,x_r)\), and let \(B\) be the quotient of the
symmetric algebra of \(N\) by the ideal generated by the
elements \(x^p\), for \(x\in N\). Since every element of \(I\)
has zero \(p\)-th power, we obtain a morphism of \(R\)-algebras
\[
  \psi:B\longrightarrow A.
\]
By \SGARefLink{sga3:VIIA:7.4.2}{7.4.2}, \(\psi\otimes k\) is an isomorphism. It follows that
\(\operatorname{Coker}\psi=0\), and, if
\(K=\ker\psi\), the morphism
\[
  \tau:K\otimes k\longrightarrow B\otimes k
\]
is zero. But \(\psi\) is surjective and \(A\) is flat over
\(R\), so \(\tau\) is also injective. Hence
\(K\otimes k=0\).

Since \(A\) is an \(R\)-module of finite presentation, \(K\) is
an \(R\)-module of finite type (cf.\ [BAC] I, \S~\SGARefLink{sga3:VIIB:2.8}{2.8},
Lemma~9). Nakayama's lemma now gives \(K=0\). Thus \(\psi\) is
an isomorphism of \(R\)-algebras. Since \(\psi^{-1}(I)\)
contains the augmentation ideal \(J\) of \(B\), it follows that
\(\psi^{-1}(I)=J\). Therefore \(\psi^{-1}\) induces an
isomorphism of \(R\)-modules
\[
  I/I^2
  \xrightarrow{\sim}
  J/J^2=N.
\]
This proves that \(\omega_{G/S}\) is finite locally free, and
hence proves \SGARefLink{sga3:VIIA:7.4-theorem:item:iv}{(iv)} \(\Rightarrow\) \SGARefLink{sga3:VIIA:7.4-theorem:item:iii}{(iii)}. The proof of
\SGARefLink{sga3:VIIA:7.4}{Theorem 7.4} is complete.

\medskip
\noindent\textbf{Remark 7.5.}\footnote{The editors added this
remark.}
\phantomsection
\label{sga3:VIIA:7.5}

It follows immediately from \SGARefLink{sga3:VIIA:7.2}{Theorems 7.2} and \SGARefLink{sga3:VIIA:7.4}{7.4} that the
functors
\[
  G\longmapsto\operatorname{Lie}(G),
  \qquad
  \mathcal L\longmapsto\mathfrak G_p(\mathcal L)
\]
induce quasi-inverse equivalences between the category of
\(S\)-groups which are locally of finite presentation, flat,
and of height \(\leq1\), and the full subcategory of the
category of \(\mathcal O_S\)-\(p\)-Lie algebras consisting of
the finite locally free \(\mathcal O_S\)-\(p\)-Lie algebras.

% Next cursor: Exp7A local p. 52, combined-reader p. 505,
% re-edition printed p. 494, original marginal locator 470,
% at the heading "8. The Case of a Base Field".

% END INCLUDED SGA3_English_current_buildable_source_20260823/source/inputs/VII/tex/components/12_expose_VIIA_section7_radicial_height1_en.tex

% BEGIN INCLUDED SGA3_English_current_buildable_source_20260823/source/inputs/VII/tex/components/13_expose_VIIA_section8_base_field_81_83_en.tex
% SGA 3, Expose VII_A, English reconstruction -- Loop 1.
% Scope: authority local pages 52--55, Section 8 opening through 8.3.1,
% stopping before 8.4 on local page 56.
% Combined-reader pages: 505--508. Re-edition printed pages: 494--497.
% Original-edition marginal locators within this span: 470--472.
% Authority: Polo--Gille Exp7A-13oct24.pdf,
% SHA-256 D46DB8C79D0385643B6EED4B96900613248375F9B73B4202478A4F57FD3E8522.
% Jacob Reinhold's CC BY 4.0 Markdown was consulted as comparison lineage.

\section*{8. The Case of a Base Field}
\phantomsection
\label{sga3:VIIA:8}
\addcontentsline{toc}{section}{8. The Case of a Base Field}

\medskip
\noindent\textbf{8.1.}
\phantomsection
\label{sga3:VIIA:8.1}

% Authority: Exp7A local p. 52; combined p. 505; printed p. 494.
We now summarize the results obtained when \(S\) is the spectrum
of a field \(k\) of characteristic \(p>0\). We call a
\(k\)-group scheme \emph{algebraic} if its underlying scheme is
of finite type over \(k\). \SGARefLink{sga3:VIIA:7.2}{Theorem 7.2} then gives the following
result.\footnote{The editors supplied the numbering \SGARefLink{sga3:VIIA:8.1.1}{8.1.1}
through \SGARefLink{sga3:VIIA:8.1.3}{8.1.3} in order to highlight the results stated here.}

\medskip
\noindent\textbf{Theorem 8.1.1.}
\phantomsection
\label{sga3:VIIA:8.1.1}
\textit{The functor \(\mathfrak G_p\), which associates with
every finite-dimensional restricted \(p\)-Lie algebra
\(\mathcal L\) over \(k\) the \(k\)-group
\(\mathfrak G_p(\mathcal L)\), is left adjoint to the functor
which associates \(\operatorname{Lie}(G)\) with every algebraic
\(k\)-group \(G\).}

Combining this with \SGARefLink{sga3:VIIA:7.4}{Theorem 7.4} gives:

\medskip
\noindent\textbf{Theorem 8.1.2.}
\phantomsection
\label{sga3:VIIA:8.1.2}
\textit{The functors
\[
  \mathfrak G_p:
  \mathcal L\longmapsto\mathfrak G_p(\mathcal L),
  \qquad
  G\longmapsto\operatorname{Lie}(G),
\]
induce quasi-inverse equivalences between the category of
finite-dimensional restricted \(p\)-Lie algebras over \(k\)
and the category of algebraic \(k\)-groups of height at most
\(1\).}

% Authority: Exp7A local p. 53; combined p. 506; printed p. 495.
Since \(\mathfrak G_p\) is a left adjoint, it commutes with
inductive limits,\footnote{The editors expanded the argument
which follows, as well as the proof of \SGARefLink{sga3:VIIA:8.1.3}{Corollary 8.1.3}.} and in
particular with cokernels. Let
\[
  \varphi:G\longrightarrow H,
  \qquad
  \varphi':G'\longrightarrow H
\]
be two morphisms between algebraic \(k\)-groups of height at
most \(1\). The fiber product \(G\times_HG'\) is again an
algebraic \(k\)-group of height at most \(1\), because the
Frobenius morphism
\[
  \operatorname{Fr}:G\longrightarrow G^{(p)}
\]
commutes with fiber products. Thus the inclusion of the
category of algebraic \(k\)-groups of height at most \(1\) into
the category of all algebraic \(k\)-groups commutes with fiber
products, and hence in particular with kernels. We obtain:

\medskip
\noindent\textbf{Corollary 8.1.3.}
\phantomsection
\label{sga3:VIIA:8.1.3}
\textit{The functor \(\mathfrak G_p\) is exact in the following
sense. Let
\[
  \pi:\mathcal L_1\longrightarrow\mathcal L_2
\]
be a surjective homomorphism between finite-dimensional
restricted \(p\)-Lie algebras over \(k\), and let \(i\) be the
inclusion of \(\mathcal L_0=\ker\pi\) into \(\mathcal L_1\).
Then the following sequence of algebraic \(k\)-groups is exact:}
\begin{center}
  \phantomsection
  \label{sga3:VIIA:8.1.3-Gp-exact-sequence}
  % SGA3-VII-DIAG-071
  % BEGIN INCLUDED SGA3_English_current_buildable_source_20260823/source/inputs/VII/native_diagrams/exp7/viia_05_14/diagram_071.tex
\(\displaystyle
1 \longrightarrow \mathfrak G_p(\mathcal L_0)
  \xrightarrow{\mathfrak G_p(i)}
\mathfrak G_p(\mathcal L_1)
  \xrightarrow{\mathfrak G_p(\pi)}
\mathfrak G_p(\mathcal L_2)
  \longrightarrow 1 .
\)

% END INCLUDED SGA3_English_current_buildable_source_20260823/source/inputs/VII/native_diagrams/exp7/viia_05_14/diagram_071.tex

\end{center}
\footnote{Moreover, by VIA, \SGARefLink{thm:sga3-VIA-3.2}{3.2},
\(\mathfrak G_p(\mathcal L_2)\) represents the (fppf) quotient
sheaf of \(\mathfrak G_p(\mathcal L_1)\) by
\(\mathfrak G_p(\mathcal L_0)\).}

Indeed, by the preceding discussion, \(\mathfrak G_p(i)\)
induces an isomorphism from
\(\mathfrak G_p(\mathcal L_0)\) onto
\(\ker(\mathfrak G_p(\pi))\). This kernel is the same whether
it is taken in the category of all algebraic \(k\)-groups or in
the subcategory of those of height at most \(1\). Furthermore,
\[
  \mathfrak G_p(\pi):
  \mathfrak G_p(\mathcal L_1)
  \longrightarrow
  \mathfrak G_p(\mathcal L_2)
\]
identifies \(\mathfrak G_p(\mathcal L_2)\) with the quotient of
\(\mathfrak G_p(\mathcal L_1)\) by
\(\mathfrak G_p(\mathcal L_0)\) in the category of algebraic
\(k\)-groups.

\medskip
\noindent\textbf{Remark 8.1.4.}\footnote{The editors added this
remark, communicated by O.~Gabber, for use in \SGARefLink{sga3:VIIA:8.3.1}{8.3.1}.}
\phantomsection
\label{sga3:VIIA:8.1.4}

Let \(\varphi:G\to H\) be a morphism of \(k\)-groups and put
\(K=\ker(\varphi)\). Suppose that \(\varphi\) is covering for
the (fpqc) topology. This is in particular the case if it is
covering for a coarser topology, such as the (fppf) topology.
Then \(\varphi\) is a \(K\)-torsor over \(H\) (cf.\ IV,
\SGARefLink{sga3:iv:corollary:5:1:7:1}{5.1.7.1}). Moreover, by IV, \SGARefLink{sga3:iv:proposition:6:3:1}{6.3.1}, there exists a covering of
\(H\) by affine open subsets \(S_i\), and for each \(i\) an
affine faithfully flat morphism \(T_i\to S_i\) factoring through
\(\varphi\). Now \(G\times_HT_i\) is \(T_i\)-isomorphic to
\(K\times T_i\), hence faithfully flat over \(T_i\). By (fpqc)
descent, \(G\times_HS_i\to S_i\) is faithfully flat, and
therefore \(\varphi\) is faithfully flat.

Conversely, if \(\varphi\) is faithfully flat and quasi-compact,
respectively faithfully flat and locally of finite
presentation, then it is covering for the (fpqc), respectively
(fppf), topology (cf.\ IV, \SGARefLink{sga3:iv:proposition:6:3:1}{6.3.1}). Finally, a morphism of
sheaves is covering if and only if it is an epimorphism
(cf.\ IV, \SGARefLink{sga3:iv:proposition:4:4:3}{4.4.3}). In particular, a quasi-compact morphism of
\(k\)-groups is faithfully flat if and only if it is an
epimorphism of (fpqc) sheaves.

\medskip
\noindent\textbf{8.2. Proposition.}
\phantomsection
\label{sga3:VIIA:8.2}
\textit{Consider an exact sequence\footnote{Thus \(\pi\) is
faithfully flat and \(\tau\) identifies \(G'\) with
\(\ker\pi\), so \(G''\) represents the (fppf) quotient sheaf of
\(G\) by \(G'\); cf.\ VIA, \SGARefLink{thm:sga3-VIA-3.2}{3.2} and \SGARefLink{sga3:VIIB:5.2:theorem}{5.2}.} of algebraic groups
over a field \(k\) of characteristic \(p>0\):}
\begin{center}
  \phantomsection
  \label{sga3:VIIA:8.2-algebraic-groups-exact-sequence}
  % SGA3-VII-DIAG-072
  % BEGIN INCLUDED SGA3_English_current_buildable_source_20260823/source/inputs/VII/native_diagrams/exp7/viia_05_14/diagram_072.tex
\(\displaystyle
1 \longrightarrow G'
  \xrightarrow{\tau} G
  \xrightarrow{\pi} G''
  \longrightarrow 1 .
\)

% END INCLUDED SGA3_English_current_buildable_source_20260823/source/inputs/VII/native_diagrams/exp7/viia_05_14/diagram_072.tex

\end{center}
\textit{Consider the following assertions:
\begin{enumerate}
  \item[\textup{(i)}] \(\pi\) is smooth;
  \item[\textup{(ii)}] \(G'\) is smooth;
  \item[\textup{(iii)}] for every integer \(n>0\), the sequence
  induced by \(\tau\) and \(\pi\),
\end{enumerate}}
\begin{center}
  \phantomsection
  \label{sga3:VIIA:8.2-frobenius-kernel-exact-sequence}
  % SGA3-VII-DIAG-073
  % BEGIN INCLUDED SGA3_English_current_buildable_source_20260823/source/inputs/VII/native_diagrams/exp7/viia_05_14/diagram_073.tex
\(\displaystyle
1 \longrightarrow \operatorname{Fr}^nG'
  \longrightarrow \operatorname{Fr}^nG
  \longrightarrow \operatorname{Fr}^nG''
  \longrightarrow 1 .
\)

% END INCLUDED SGA3_English_current_buildable_source_20260823/source/inputs/VII/native_diagrams/exp7/viia_05_14/diagram_073.tex

\end{center}
\textit{\noindent is exact;
\begin{enumerate}
  \item[\textup{\SGAInlineRefTarget{sga3:VIIA:8.2:item:iv}(iv)}] the morphism
  \[
    \operatorname{Fr}\pi:
    \operatorname{Fr}G\longrightarrow\operatorname{Fr}G''
  \]
  is an epimorphism of (fppf) sheaves;
  \item[\textup{\SGAInlineRefTarget{sga3:VIIA:8.2:item:v}(v)}] the morphism
  \[
    \operatorname{Lie}(\pi):
    \operatorname{Lie}(G)\longrightarrow\operatorname{Lie}(G'')
  \]
  is surjective.
\end{enumerate}
Then
\[
  \textup{(i)}\Longleftrightarrow\textup{(ii)}
  \Longrightarrow\textup{(iii)}
  \Longrightarrow\textup{(iv)}
  \Longleftrightarrow\textup{(v)},
\]
and all five assertions are equivalent when \(G\) is smooth
over \(k\).}

% Authority: Exp7A local p. 54; combined p. 507; printed p. 496.
Indeed, \SGARefLink{sga3:vii:paragraph:5-1:equation:i}{(i)} is equivalent to \SGARefLink{sga3:vii:paragraph:5-1:equation:ii}{(ii)} by VIB, \SGARefLink{prop:sga3-VIB-9.2}{9.2}~\SGARefLink{concl:sga3-VIB-9.2-vii}{(vii)}, and it is
clear that \SGARefLink{concl:sga3-VIB-9.2-iii}{(iii)} implies \SGARefLink{sga3:VIIA:8.2:item:iv}{(iv)}. The equivalence of \SGARefLink{sga3:VIIA:8.2:item:iv}{(iv)} and \SGARefLink{sga3:VIIA:8.2:item:v}{(v)}
follows from \SGARefLink{sga3:VIIA:8.1.3}{8.1.3}.

The implication \SGARefLink{concl:sga3-VIB-9.2-ii}{(ii)} \(\Rightarrow\) \SGARefLink{concl:sga3-VIB-9.2-iii}{(iii)} follows from the
diagram
\begin{center}
  \phantomsection
  \label{sga3:VIIA:8.2-frobenius-exact-rows-diagram}
  % SGA3-VII-DIAG-074
  % BEGIN INCLUDED SGA3_English_current_buildable_source_20260823/source/inputs/VII/native_diagrams/exp7/viia_05_14/diagram_074.tex
\resizebox{0.98\linewidth}{!}{%
\begin{tikzcd}[ampersand replacement=\&,row sep=3.6em,column sep=3.5em]
1 \arrow[r]
  \& G' \arrow[r,"\tau"] \arrow[d,swap,"\operatorname{Fr}^n(G'/k)"]
  \& G \arrow[r,"\pi"] \arrow[d,swap,"\operatorname{Fr}^n(G/k)"]
  \& G'' \arrow[r] \arrow[d,swap,"\operatorname{Fr}^n(G''/k)"]
  \& 1 \\
1 \arrow[r]
  \& G'^{(p^n)} \arrow[r,"\tau^{(p^n)}"']
  \& G^{(p^n)} \arrow[r,"\pi^{(p^n)}"']
  \& G''^{(p^n)} \arrow[r]
  \& 1
\end{tikzcd}%
}

% END INCLUDED SGA3_English_current_buildable_source_20260823/source/inputs/VII/native_diagrams/exp7/viia_05_14/diagram_074.tex

\end{center}
whose two rows are exact. Since
\(\operatorname{Fr}^n(G'/k)\) is an epimorphism of (fppf)
sheaves by \SGARefLink{sga3:VIIA:8.3.1}{Corollary 8.3.1} below, \(\pi\) induces an
epimorphism from \(\operatorname{Fr}^nG\) onto
\(\operatorname{Fr}^nG''\); here one uses the snake lemma for
sheaves of groups, not necessarily commutative.

Likewise, suppose that \(G\) is smooth over \(k\). Then
\(\operatorname{Fr}(G/k)\) is an epimorphism. If
\(\operatorname{Fr}\pi\) is also an epimorphism, the same snake
lemma, applied to the preceding diagram with \(n=1\), shows
that \(\operatorname{Fr}(G'/k)\) is an epimorphism. Hence \(G'\)
is smooth over \(k\) by \SGARefLink{sga3:VIIA:8.3.1}{8.3.1} below.

\medskip
\noindent\textbf{8.3. Proposition.}\footnote{The editors
replaced ``algebraic'' by ``locally of finite type.''}
\phantomsection
\label{sga3:VIIA:8.3}
\textit{Let \(G\) be a group locally of finite type over a field
\(k\) of characteristic \(p>0\). There exists an integer
\(n_0\) such that
\[
  G/(\operatorname{Fr}^nG)
\]
is smooth over \(k\) for every \(n\ge n_0\).}

The construction of \(G/(\operatorname{Fr}^nG)\) commutes with
extension of the base field (\SGARefLink{sga3:VIIA:4.1.1}{4.1.1} and VIA, \SGARefLink{thm:sga3-VIA-3.3.2}{3.3.2}), so we may
suppose that \(k\) is perfect. Then \(G_{\mathrm{red}}\) is a
\(k\)-group locally of finite type (cf.\ VIA, \SGARefLink{sga3:VIA:0.2}{0.2}), and, putting
\[
  H=G_{\mathrm{red}}\backslash G,
\]
we have the following commutative exact diagram:
\begin{center}
  \phantomsection
  \label{sga3:VIIA:8.3-reduced-quotient-frobenius-diagram}
  % SGA3-VII-DIAG-075
  % BEGIN INCLUDED SGA3_English_current_buildable_source_20260823/source/inputs/VII/native_diagrams/exp7/viia_05_14/diagram_075.tex
\resizebox{0.94\linewidth}{!}{%
\begin{tikzcd}[ampersand replacement=\&,row sep=3.8em,column sep=3.7em]
1 \arrow[r]
  \& G_{\mathrm{red}} \arrow[r] \arrow[d,"\operatorname{Fr}^n(G_{\mathrm{red}}/k)"]
  \& G \arrow[r] \arrow[d,"\operatorname{Fr}^n(G/k)"]
  \& H \arrow[d,"\operatorname{Fr}^n(H/k)"] \\
1 \arrow[r]
  \& G_{\mathrm{red}}^{(p^n)} \arrow[r]
  \& G^{(p^n)} \arrow[r]
  \& H^{(p^n)}\, .
\end{tikzcd}%
}

% END INCLUDED SGA3_English_current_buildable_source_20260823/source/inputs/VII/native_diagrams/exp7/viia_05_14/diagram_075.tex

\end{center}

% Authority: Exp7A local p. 55; combined p. 508; printed p. 497.
The scheme \(H\) is the spectrum of a finite local \(k\)-algebra
with residue field \(k\) (cf.\ VIA, \SGARefLink{prop:sga3-VIA-5.6.1}{5.6.1}). There is therefore
an integer \(n_0\) such that, for every \(n\ge n_0\),
\(\operatorname{Fr}^n(H/k)\) factors through the unit section of
\(H^{(p^n)}\). It follows that, for \(n\ge n_0\),
\(\operatorname{Fr}^n(G/k)\) factors through
\(G_{\mathrm{red}}^{(p^n)}\), and VIA, \SGARefLink{prop:sga3-VIA-5.4.1}{5.4.1} gives the
commutative diagram
\begin{center}
  \phantomsection
  \label{sga3:VIIA:8.3-frobenius-factorization-triangle}
  % SGA3-VII-DIAG-076
  % BEGIN INCLUDED SGA3_English_current_buildable_source_20260823/source/inputs/VII/native_diagrams/exp7/viia_05_14/diagram_076.tex
\begin{tikzcd}[row sep=3.6em,column sep=4.6em]
G \arrow[r,"\operatorname{Fr}^n(G/k)"] \arrow[d,swap,"\pi"]
  & G_{\mathrm{red}}^{(p^n)} \\
G/(\operatorname{Fr}^nG) \arrow[ur,swap,"i"]
\end{tikzcd}

% END INCLUDED SGA3_English_current_buildable_source_20260823/source/inputs/VII/native_diagrams/exp7/viia_05_14/diagram_076.tex

\end{center}
where \(i\) is a closed immersion and \(\pi\) is the canonical
projection. Moreover, \(i\) induces a homeomorphism of the
underlying topological spaces, so it is an isomorphism. Since
\(k\) is perfect, \(G_{\mathrm{red}}^{(p^n)}\) is smooth over
\(k\) (VIA, \SGARefLink{prop:sga3-VIA-1.3.1}{1.3.1}). Thus \(G/(\operatorname{Fr}^nG)\) is smooth
over \(k\) for every \(n\ge n_0\).

\medskip
\noindent\textbf{8.3.1. Corollary.}
\phantomsection
\label{sga3:VIIA:8.3.1}
\textit{Let \(G\) be a group locally of finite type over a field
\(k\) of characteristic \(p>0\), and let \(n\ge1\).\footnote{The
editors made explicit the equivalence of \SGARefLink{sga3:vii:paragraph:5-1:equation:ii}{(ii)} and \SGARefLink{sga3:VIIA:8.3.1:item:iii}{(iii)} and
expanded the proof.} The following conditions are equivalent:
\begin{enumerate}
  \item[\textup{(i)}] \(G\) is smooth over \(k\);
  \item[\textup{(ii)}]
  \(\operatorname{Fr}^n(G/k):G\to G^{(p^n)}\) is an
  epimorphism of (fppf) sheaves;
  \item[\textup{\SGAInlineRefTarget{sga3:VIIA:8.3.1:item:iii}(iii)}]
  \(\operatorname{Fr}^n(G/k):G\to G^{(p^n)}\) is faithfully
  flat.
\end{enumerate}}

Since \(G\) is locally of finite type over \(k\),
\(\operatorname{Fr}^n(G/k)\) is of finite presentation, so the
equivalence of \SGARefLink{sga3:vii:paragraph:5-1:equation:ii}{(ii)} and \SGARefLink{sga3:VIIA:8.3.1:item:iii}{(iii)} follows from \SGARefLink{sga3:VIIA:8.1.4}{8.1.4}. Suppose that
\(G\) is smooth over \(k\), and hence reduced. Since
\(\operatorname{Fr}^n(G/k)\) is surjective, it is faithfully
flat (cf.\ VIA, \SGARefLink{cor:sga3-VIA-6.2}{6.2} or VIB, \SGARefLink{prop:sga3-VIB-1.3}{1.3}).

Conversely, suppose that \(\operatorname{Fr}^n(G/k)\) is
faithfully flat. The morphism
\[
  \operatorname{Fr}^n(G^{(p^n)}/k)
\]
is obtained from \(\operatorname{Fr}^n(G/k)\) by base change
(cf.\ \SGARefLink{sga3:VIIA:4.1.3}{4.1.3}), so it too is faithfully flat, as is the composite
\[
  \operatorname{Fr}^{2n}(G/k):
  G\longrightarrow G^{(p^n)}
  \longrightarrow G^{(p^{2n})}.
\]
It follows that, for every \(m\in\mathbb N\),
\[
  \operatorname{Fr}^{mn}(G/k):
  G\longrightarrow G^{(p^{mn})}
\]
is faithfully flat. Hence it induces an isomorphism
\[
  G/(\operatorname{Fr}^{mn}G)
  \simeq
  G^{(p^{mn})}
\]
(cf.\ VIA, \SGARefLink{prop:sga3-VIA-5.4.1}{5.4.1}). By \SGARefLink{sga3:VIIA:8.3}{Proposition 8.3},
\(G^{(p^{mn})}\) is smooth over \(k\) for \(m\) sufficiently
large, and therefore \(G\) is smooth as well by (fpqc) descent
(cf.\ EGA IV\(_4\), 17.7.1).

% Next cursor: Exp7A local p. 56, combined-reader p. 509,
% re-edition printed p. 498, at 8.4.

% END INCLUDED SGA3_English_current_buildable_source_20260823/source/inputs/VII/tex/components/13_expose_VIIA_section8_base_field_81_83_en.tex

% BEGIN INCLUDED SGA3_English_current_buildable_source_20260823/source/inputs/VII/tex/components/14_expose_VIIA_section84_85_bibliography_en.tex
% SGA 3, Expose VII_A, English reconstruction -- Loop 1.
% Scope: authority local pages 56--60, paragraphs 8.4--8.5.3 and
% the terminal VIIA bibliography.
% Combined-reader pages: 509--513. Re-edition printed pages: 498--502.
% Original-edition marginal locators within this span: 473--476.
% Authority: Polo--Gille Exp7A-13oct24.pdf,
% SHA-256 D46DB8C79D0385643B6EED4B96900613248375F9B73B4202478A4F57FD3E8522.
% Jacob Reinhold's CC BY 4.0 Markdown was consulted as comparison lineage.

\medskip
\noindent\textbf{8.4.}
\phantomsection
\label{sga3:VIIA:8.4}

% Authority: Exp7A local p. 56; combined p. 509; printed p. 498.
In the two statements which conclude this Expos\'e, we return to
the case of a field \(k\) of arbitrary characteristic.

If \(k\) has characteristic \(0\), let \(n\ge1\). If \(k\) has
characteristic \(p>0\), let \(n\ge1\) be prime to \(p\). In both
cases we simply say that \(n\) is \emph{prime to the
characteristic of \(k\)}. If \(G\) is a group scheme over \(k\),
we denote by
\[
  n_G:G\longrightarrow G
\]
the morphism of \(k\)-schemes which sends
\(x\in G(T)\) to \(x^n\in G(T)\), for every \(k\)-scheme \(T\).

\medskip
\noindent\textbf{Proposition.}
\phantomsection
\label{sga3:VIIA:8.4-proposition}
\textit{Let \(G\) be an algebraic group over a field \(k\), and
let \(n\) be prime to the characteristic of \(k\). Then
\[
  n_G:G\longrightarrow G
\]
is an \'etale morphism.}

By VIB, \SGARefLink{prop:sga3-VIB-1.3}{1.3}, it suffices to show that \(n_G\) is \'etale at the
origin.\footnote{The editors changed ``\'etale at the origin''
to ``\'etale'' in the statement and added the preceding
sentence.} Let \(A\) be the local ring of \(G\) at the origin
and \(I\) its maximal ideal. By II, \SGARefLink{sga3:ii:corollary:3-9-4}{3.9.4}, the map
\[
  \operatorname{Lie}(n_G):
  \operatorname{Lie}(G)\longrightarrow\operatorname{Lie}(G)
\]
induced by \(n_G\) is multiplication by \(n\). It is therefore
an isomorphism, as is the endomorphism induced by \(n_G\) on
\[
  I/I^2=\operatorname{Lie}(G)^*.
\]

Suppose first that \(\operatorname{char}k=0\). The group \(G\)
is smooth over \(k\) (VIB, 1.6.1; see also VIIB, 3.3.1), and the
canonical homomorphism
\[
  \operatorname{Sym}(I/I^2)\longrightarrow
  \operatorname{gr}_I(A)
\]
is an isomorphism, where \(\operatorname{gr}_I(A)\) is the
graded ring associated with the \(I\)-adic filtration. Hence
\(n_G\) induces an automorphism of \(\operatorname{gr}_I(A)\),
and therefore also of the completion \(\widehat A\). Thus
\(n_G\) is \'etale at the origin (cf.\ EGA IV\(_4\), 17.6.3).

Suppose next that \(\operatorname{char}k=p>0\) and \(G\) has
height at most \(1\). Then \(A\) is isomorphic to the quotient
of the symmetric algebra on
\[
  \omega_{G/k}=I/I^2
\]
by the ideal generated by the \(p\)-th powers of the elements
of \(\omega_{G/k}\) (cf.\ 7.4). The same argument as in
characteristic \(0\) shows that \(n_G\) induces an automorphism
of \(A\).

Suppose that \(G\) has height at most \(r\), and that the result
has been proved for groups of height at most \(r-1\). Let
\(B,A,A'\) be the affine algebras of
\(\operatorname{Fr}G,G\), and
\[
  G'=\operatorname{Fr}G\backslash G,
\]
and let \(n_B,n_A,n_{A'}\) be the endomorphisms induced by
\(n_{\operatorname{Fr}G},n_G,n_{G'}\), respectively. Put
\[
  I'=I\cap A',
\]
the maximal ideal of \(A'\). The cartesian square
\begin{center}
  \phantomsection
  \label{sga3:VIIA:8.4-height-induction-cartesian-square}
  % SGA3-VII-DIAG-077
  % BEGIN INCLUDED SGA3_English_current_buildable_source_20260823/source/inputs/VII/native_diagrams/exp7/viia_05_14/diagram_077.tex
\begin{tikzcd}[row sep=3.4em,column sep=4.2em]
\operatorname{Fr}G \arrow[r] \arrow[d]
  & G \arrow[d] \\
e \arrow[r]
  & G'
\end{tikzcd}

% END INCLUDED SGA3_English_current_buildable_source_20260823/source/inputs/VII/native_diagrams/exp7/viia_05_14/diagram_077.tex

\end{center}
gives
\[
  B=A/I'A.
\]
\footnote{The editors expanded the original argument in what
follows.}
The endomorphism \(n_{A'}\), respectively \(n_B\), is the
endomorphism induced by \(n_A\) on \(A'\), respectively on
\(B\). By VIA, \SGARefLink{thm:sga3-VIA-3.2}{3.2}, \(A\) is a faithfully flat
\(A'\)-module. Since \(A'\) is an Artinian local ring, because
\(G'\) is an algebraic \(k\)-group of height at most \(r-1\),
\(A\) is a free \(A'\)-module. The restriction of \(n_A\) to
\(A'\) is \(n_{A'}\), which is an isomorphism by the induction
hypothesis. Nakayama's lemma therefore shows that \(n_A\) is an
automorphism as soon as the endomorphism it induces on
\(A/I'A\) is an automorphism.

% Authority: Exp7A local p. 57; combined p. 510; printed p. 499.
That endomorphism is \(n_B\), which is an automorphism because
\(\operatorname{Fr}G\) has height at most \(1\). Hence \(n_A\)
is an automorphism.

Finally, let \(G\) be any algebraic group over a field of
characteristic \(p>0\). The preceding argument shows that
\(n_G\) induces automorphisms of the \(k\)-schemes
\(\operatorname{Fr}^rG\). These schemes are affine over \(k\),
and their algebras are the quotients of the local algebra \(A\)
by the ideal
\[
  I^{\{p^r\}}
\]
generated by the \(p^r\)-th powers of the elements of \(I\).
Since \(n_G\) induces automorphisms of all the algebras
\(A/I^{\{p^r\}}\), passage to the projective limit shows that it
induces an automorphism of \(\widehat A\). Thus \(n_G\) is
\'etale at the origin (EGA IV\(_4\), 17.6.3).

\medskip
\noindent\textbf{8.5. Proposition.}
\phantomsection
\label{sga3:VIIA:8.5}
\textit{Let \(G\) be a finite algebraic group of rank \(n\) over
the field \(k\). Then
\[
  n_G:G\longrightarrow G
\]
is the zero morphism of \(G\).}

Combining 8.4 and 8.5 immediately gives the following
corollary.

\medskip
\noindent\textbf{Corollary 8.5.1.}
\phantomsection
\label{sga3:VIIA:8.5.1}
\textit{Let \(G\) be a finite algebraic group of rank \(n\) over
the field \(k\). If \(n\) is prime to the characteristic of
\(k\), then \(G\) is \'etale over \(k\).}
\footnote{The editors added this corollary, which the original
indicated implicitly by ``cf.\ \SGARefLink{sga3:VIIA:8.4}{8.4}.'' For another proof which
does not use \SGARefLink{sga3:VIIA:8.5}{8.5}, see \textup{[TO70]}, \SGARefLink{sga3:VIIA:5}{Lemma~5}.}

We now prove \SGARefLink{sga3:VIIA:8.5}{8.5}. Let \(H\) be a normal subgroup of \(G\), of
rank \(m\) over \(k\), and let
\[
  \lambda:H\times G\longrightarrow G
\]
be induced by the multiplication of \(G\). With the notation of
VIA, 3.2, there is a cartesian square
\begin{center}
  \phantomsection
  \label{sga3:VIIA:8.5-normal-subgroup-cartesian-square}
  % SGA3-VII-DIAG-078
  % BEGIN INCLUDED SGA3_English_current_buildable_source_20260823/source/inputs/VII/native_diagrams/exp7/viia_05_14/diagram_078.tex
\begin{tikzcd}[row sep=3.4em,column sep=4.2em]
H\times G \arrow[r,"\lambda"] \arrow[d,swap,"\operatorname{pr}_2"]
  & G \arrow[d,"\pi"] \\
G \arrow[r,"\pi"]
  & H\backslash G
\end{tikzcd}

% END INCLUDED SGA3_English_current_buildable_source_20260823/source/inputs/VII/native_diagrams/exp7/viia_05_14/diagram_078.tex

\end{center}
Since \(\pi:G\to H\backslash G\) is faithfully flat and
quasi-compact (VIA, 3.2), while
\(\operatorname{pr}_2\) is locally free of rank \(m\), EGA
IV\(_2\), 2.5.2 shows that \(G\to H\backslash G\) is locally
free of rank \(m\). If
\[
  r=\operatorname{rk}_k(H\backslash G),
\]
then
\[
  n=\operatorname{rk}_k(G)=rm.
\]

On the other hand, for every \(k\)-scheme \(T\) there is an
exact sequence of abstract groups
\[
  1\longrightarrow H(T)\longrightarrow G(T)
  \longrightarrow(H\backslash G)(T).
\]
It is therefore clear that \(n_G\) is zero whenever \(m_H\) and
\(r_{H\backslash G}\) are zero. Taking \(H=G^0\), the identity
component, the quotient \(G^0\backslash G\) is \'etale
(cf.\ VIA, \SGARefLink{prop:sga3-VIA-5.5.1}{5.5.1}). We may consequently assume that \(G\) is
\'etale over \(k\), or that it is infinitesimal (cf.\ \SGARefLink{sga3:VIIA:7.0}{7.0}).

If \(G\) is \'etale, extension of the base field reduces us to
the case where \(k\) is algebraically closed. Then \(G\) is a
constant group (cf.\ I, \SGARefLink{sga3:i:subsection:4-1}{4.1}), and the assertion is classical.

If \(G\) is nonzero and infinitesimal, then \(k\) necessarily
has characteristic \(p>0\) (cf.\ VIB, \SGARefLink{cor:sga3-VIB-1.6.1}{1.6.1} or VIIB, \SGARefLink{sga3:VIIA:3.3.1}{3.3.1}).
The subgroups \(\operatorname{Fr}^rG\) form a composition series
of \(G\) whose successive quotients have height at most \(1\).

% Authority: Exp7A local p. 58; combined p. 511; printed p. 500.
We are thus reduced to the case where \(G\) has height at most
\(1\). Let \(A\), respectively \(L\), be the affine algebra,
respectively the Lie algebra, of \(G\), and put
\[
  U=U_p(L).
\]
By \SGARefLink{sga3:VIIA:7.4}{7.4}, \(G=\mathfrak G_p(L)\), and hence \(A=U^*\). Therefore,
if \(\dim_kL=r\), the rank of \(G\) over \(k\) is \(p^r\)
(cf.\ \SGARefLink{sga3:VIIA:5.3.3}{5.3.3}). We shall study the morphism
\[
  p_G:G\longrightarrow G
\]
defined by taking the \(p\)-th power. It induces an
endomorphism \(p_A\) of \(A\), and, by duality, an endomorphism
\(p_U\) of \(U\).

Let \(I\) be the augmentation ideal of \(A\). We shall prove
\[
  \tag{*}\label{sga3:vii:corollary:8-5-1:equation:star}
  p_A(I)\subset I^p.
\]
Assuming this, one has
\[
  p_A^r(I)\subset I^{p^r}.
\]
Moreover,
\[
  I^{r(p-1)+1}=0,
\]
because \(I\) is generated by \(r\) elements whose \(p\)-th
powers vanish. Since \(p^r>r(p-1)\), we obtain
\(p_A^r(I)=0\), and therefore \(p_G^r\) is the zero morphism.

For every integer \(s\ge1\), let
\[
  m_A^{s-1}:A^{\otimes s}\longrightarrow A,
  \qquad
  \Delta_U^{s-1}:U\longrightarrow U^{\otimes s},
\]
denote the maps induced by the multiplication \(m_A\) of \(A\)
and the comultiplication \(\Delta_U\) of \(U\), respectively.
Then \(p_A\) is the composite
\[
  A\xrightarrow{\Delta_A^{p-1}}
  A^{\otimes p}
  \xrightarrow{m_A^{p-1}}A.
\]
Since the transpose of \(m_A\), respectively of \(\Delta_A\),
is \(\Delta_U\), respectively \(m_U\), the endomorphism
\[
  p_U={}^tp_A
\]
is the composite
\[
  U\xrightarrow{\Delta_U^{p-1}}
  U^{\otimes p}
  \xrightarrow{m_U^{p-1}}U.
\]

Let \(J\) be the augmentation ideal of \(U\).\footnote{The
editors expanded the original argument in this paragraph.}
Then
\[
  U=k1_U\oplus J,
\]
and we write \(\pi:U\to J\) for the projection with kernel
\(k1_U\). For \(s\ge1\), let \((I^s)^\perp\) be the orthogonal
of \(I^s\) in \(A^*=U\). Thus \((I^s)^\perp\) consists of the
elements \(u\in U\) for which the composite
\[
  I^{\otimes s}
  \xrightarrow{m_A^{s-1}}I
  \xrightarrow{u}k
\]
is zero. Since the transpose of \(m_A\) is \(\Delta_U\),
\((I^s)^\perp\) is the vector subspace \(P_{s-1}\) consisting
of those \(u\in U\) for which \(\Delta_U^{s-1}(u)\) vanishes on
\(I^{\otimes s}\). If
\[
  \overline\Delta_U^{s-1}
  =
  \pi^{\otimes s}\circ\Delta_U^{s-1},
\]
then
\[
  (I^s)^\perp
  =
  P_{s-1}
  =
  \ker\overline\Delta_U^{s-1}
\]
(see also VIIB, \SGARefLink{sga3:VIIB:1.3.6}{1.3.6}).

To prove \((*)\), it is enough to show that the transpose
\(p_U={}^tp_A\) maps \(P_{p-1}\) into
\[
  I^\perp=k1_U.
\]
Since \(p_U(1_U)=1_U\), it suffices to prove
\[
  \tag{**}\label{sga3:vii:corollary:8-5-1:equation:double-star}
  p_U(P_{p-1}^+)=0,
  \qquad
  P_{p-1}^+=J\cap P_{p-1}.
\]
An easy induction on \(s\) shows that \(P_{s-1}^+\) is spanned
by the products
\[
  x_1\cdots x_t,
  \qquad
  1\le t\le s-1,\quad x_i\in L
\]
(cf.\ VIIB, \SGARefLink{sga3:VIIA:4.3}{4.3}).

% Authority: Exp7A local p. 59; combined p. 512; printed p. 501.
For \(x_1,\ldots,x_t\in L\), one has
\[
\begin{aligned}
  p_U(x_1x_2\cdots x_t)
  &=
  m_U^{p-1}
  \left(
    \prod_{j=1}^t
    \sum_{i=1}^p
    1\otimes\cdots\otimes
    \underset{\text{\scriptsize position \(i\)}}{x_j}
    \otimes\cdots\otimes1
  \right).
\end{aligned}
\]
The expression inside parentheses is a sum of \(p^t\) terms
\(x_h\), indexed by the maps
\[
  h:\{1,\ldots,t\}\longrightarrow\{1,\ldots,p\}.
\]
Such a map determines an ordered partition
\(\mathcal P_h\) of \(\{1,\ldots,t\}\) into at most \(p\)
parts.\footnote{The editors expanded the original argument and
replaced the notion of a preorder by the equivalent notion of
an ordered partition.} Indeed, let
\[
  i_1<\cdots<i_q
\]
be the elements in the image of \(h\), put
\[
  I_a=h^{-1}(i_a)
  \quad(1\le a\le q),
  \qquad
  x_{I_a}=\prod_{j\in I_a}x_j,
\]
where the product is taken in increasing order. The map \(h\)
corresponds to the \(p\)-tensor
\[
  1\otimes\cdots\otimes x_{I_1}\otimes\cdots\otimes
  x_{I_q}\otimes\cdots\otimes1,
\]
where \(x_{I_a}\) occupies position \(i_a\). Its image under the
iterated multiplication is the product
\[
  x_{I_1}\cdots x_{I_q}.
\]
This depends only on the ordered partition
\[
  \mathcal P=(I_1,\ldots,I_q);
\]
denote it by \(x_{\mathcal P}\). For a fixed \(\mathcal P\),
the same term is obtained for every choice of
\(i_1<\cdots<i_q\) in \(\{1,\ldots,p\}\). Consequently,
\[
  p_U(x_1x_2\cdots x_t)
  =
  \sum_{\mathcal P}
  \binom{p}{n(\mathcal P)}x_{\mathcal P},
\]
where \(\mathcal P\) ranges over the ordered partitions of
\(\{1,\ldots,t\}\) into at most \(p\) parts and
\(n(\mathcal P)\) is the number of parts. Here
\[
  1\le n(\mathcal P)\le\min(t,p).
\]
If \(t<p\), all the coefficients
\(\binom{p}{n(\mathcal P)}\) vanish in \(k\), and hence
\[
  p_U(x_1\cdots x_t)=0.
\]
Thus \(p_U\) vanishes on \(P_{p-1}^+\). This proves \((**)\),
then \((*)\), and completes the proof of \SGARefLink{sga3:VIIA:8.5}{8.5}.

\medskip
\noindent\textbf{Corollary 8.5.2.}\footnote{The editors added
this corollary, which is noted in \SGARefLink{sga3:viii:remark:7:3:1}{Expos\'e VIII, Remark 7.3.1}.}
\phantomsection
\label{sga3:VIIA:8.5.2}
\textit{Let \(S\) be a reduced scheme and \(G\) a finite locally
free \(S\)-group of rank \(n\). Then
\[
  n_G:G\longrightarrow G
\]
is the zero morphism of \(G\).}

Let \(S'\) be the disjoint union of the
\(\operatorname{Spec}\mathcal O_{S,\eta}\), where \(\eta\)
ranges over the maximal points of \(S\). Since \(S\) is
reduced, the morphism \(S'\to S\) is schematically dominant.
The same is therefore true of
\[
  f:G_{S'}\longrightarrow G,
\]
because \(G\) is finite locally free over \(S\) (cf.\ EGA
IV\(_3\), 11.10.5). Since \(G\to S\) is affine, and hence
separated, the coincidence locus of \(n_G\) and the zero
morphism is a closed subscheme of \(G\). By 8.5 it contains the
image of \(f\), and therefore equals \(G\). Thus \(n_G\) is the
zero morphism.

\medskip
\noindent\textbf{Remark 8.5.3.}
\phantomsection
\label{sga3:VIIA:8.5.3}

A theorem of P.~Deligne (see \textup{[TO70]}, p.~4) also says
that if \(G\) is a commutative finite locally free \(S\)-group
of rank \(n\) over an arbitrary base \(S\), then \(n_G=0\).

\chapter*{Bibliography}
\phantomsection
\label{sga3:VIIA:bibliography}
\addcontentsline{toc}{chapter}{Bibliography}

% Authority: Exp7A local p. 60; combined p. 513; printed p. 502.
\begin{description}
  \item[\textup{[BAlg]}] N.~Bourbaki,
  \emph{Alg\`ebre}, Chapters I--III, Hermann, 1974;
  Chapter X, Masson, 1980.

  \item[\textup{[BAC]}] N.~Bourbaki,
  \emph{Alg\`ebre commutative}, Chapters I--IV, Masson, 1985.

  \item[\textup{[BLie]}] N.~Bourbaki,
  \emph{Groupes et alg\`ebres de Lie}, Chapter I,
  Hermann, 1971.

  \item[\textup{[DG70]}] M.~Demazure and P.~Gabriel,
  \emph{Groupes alg\'ebriques}, Masson \& North-Holland, 1970.

  \item[\textup{[Ja03]}] J.~C.~Jantzen,
  \emph{Representations of Algebraic Groups},
  Academic Press, 1987; 2nd ed., American Mathematical Society,
  2003.

  \item[\textup{[TO70]}] J.~Tate and F.~Oort,
  ``Group schemes of prime order,''
  \emph{Annales scientifiques de l'\'Ecole Normale
  Sup\'erieure} \textbf{3} (1970), 1--21.
\end{description}

% End of VIIA. Next cursor: Exp7B local p. 1, combined-reader
% p. 514, re-edition printed p. 503, at the title of VII_B.

% END INCLUDED SGA3_English_current_buildable_source_20260823/source/inputs/VII/tex/components/14_expose_VIIA_section84_85_bibliography_en.tex

% BEGIN INCLUDED SGA3_English_current_buildable_source_20260823/source/inputs/VII/tex/components/15_expose_VIIB_title_section01_pseudocompact_rings_en.tex
% SGA 3, Expose VII_B, English reconstruction -- Loop 1.
% Scope: authority local pages 1--3, title through 0.1.3,
% stopping immediately before 0.2 on local page 3.
% Combined-reader pages: 514--516. Re-edition printed pages: 503--505.
% Original-edition marginal locators within this span: 477--480.
% Authority: Polo--Gille Exp7B-13oct24.pdf,
% SHA-256 D303690E0195654BBE7C3DAA5D89EA23E24CD709A3575217234AC19AE8836694.
% Jacob Reinhold's CC BY 4.0 Markdown was consulted as comparison lineage.

\chapter*{Expos\'e \texorpdfstring{VII\textsubscript{B}}{VII-B}\\
Infinitesimal Study of Group Schemes}
\phantomsection
\label{sga3:VIIB:title}
\addcontentsline{toc}{chapter}{Expos\'e VII-B. Infinitesimal Study of Group Schemes}
\markboth{EXPOS\'E VII\textsubscript{B}. INFINITESIMAL STUDY OF GROUP SCHEMES}%
  {EXPOS\'E VII\textsubscript{B}. INFINITESIMAL STUDY OF GROUP SCHEMES}

\begin{center}
  \textit{by P. Gabriel}
\end{center}

\setcounter{footnote}{-1}
\section*{B) Formal Groups}
\phantomsection
\label{sga3:VIIB:B}
\addcontentsline{toc}{section}{B) Formal Groups}

% Authority: Exp7B local p. 1; combined-reader p. 514; printed p. 503.
\noindent\footnote{Editors' note: Version of 13 October 2024.}
The study of formal groups is usually extremely simple. If this does not
appear clearly in the pages that follow, the responsibility lies with an
arithmetician who claims to know formal groups over ``something other than
fields.''\footnote{The interest of formal groups over a complete
Noetherian local ring appears, for example, in the work of Lubin and Tate
(cf.\ \textup{[LT65]}). The study of formal groups over an arbitrary base,
and of lifting and deformation questions---in particular for
Barsotti--Tate groups (``\(p\)-divisible groups'')---has given rise to an
abundant literature; see, for example, \textup{[LT66, Ta67, Gr74, Me72,
La75, Fo77, Il85, Br00]}. In particular, the results of the present
expos\'e are largely taken up again in Chapter~I of \textup{[Fo77]}.}
We have therefore developed, for formal groups ``locally free over
projective limits of Artinian rings,'' the generalities ordinarily stated
for formal groups defined over a field. For a more detailed study of the
latter, we refer to the 1964/65 Heidelberg--Strasbourg algebraic geometry
seminar.\footnote{The editors found only such a seminar dated 1965/66 and
entitled ``Linear Algebraic Groups,'' in which the notion of a formal group
does not appear; see instead \textup{[De72]}.}

\section*{0. Review of Pseudocompact Rings and Modules}
\phantomsection
\label{sga3:VIIB:0}
\addcontentsline{toc}{section}{0. Review of Pseudocompact Rings and Modules}

This section contains a few technical preliminaries. We recall and
complete here several results from \textup{[CA]} (\emph{Des cat\'egories
ab\'eliennes}, Bull.\ Soc.\ Math.\ France \textbf{90} (1962)).

\medskip
\noindent\textbf{0.1.}
\phantomsection
\label{sga3:VIIB:0.1}
A \emph{left pseudocompact ring} is a separated, complete topological
ring with identity which has a basis of neighborhoods of \(0\) consisting
of left ideals \(\mathfrak l\) of finite colength, that is,
\[
  \operatorname{length}_A(A/\mathfrak l)<+\infty.
\]
We shall assume here that \(A\) is commutative, so that no distinction
need be made ``between left and right.''

In particular, the quotients \(A/\mathfrak l\) are Artinian rings, and
\(A\) identifies with the topological projective limit of these rings,
each endowed with the discrete topology.

A complete Noetherian local ring \((A,\mathfrak m)\) is plainly
pseudocompact.\footnote{When it is endowed with the
\(\mathfrak m\)-adic topology.}

% Authority: Exp7B local p. 2; combined-reader p. 515; printed p. 504.
\medskip
\noindent\textbf{0.1.1.}
\phantomsection
\label{sga3:VIIB:0.1.1}
Every closed ideal \(I\) of \(A\) is the intersection of the open ideals
that contain it.\footnote{Indeed, if \(x\notin I\), there is an open ideal
\(\mathfrak l\) such that
\((x+\mathfrak l)\cap I=\varnothing\); then \(I+\mathfrak l\) is an open
ideal that does not contain \(x\). In what follows, the editors have also
made explicit the fact that every ``closed maximal'' ideal is maximal and
open.}
Every closed maximal ideal is therefore open. Moreover, if
\(\mathfrak l\) is an open ideal of \(A\), the maximal ideals of
\(A/\mathfrak l\) correspond bijectively to the maximal ideals
\(\mathfrak m\) of \(A\) containing \(\mathfrak l\); the latter are thus
open and closed. Consequently, every closed maximal ideal is an open
maximal ideal (and hence closed), the converse being clear. We denote the
set of these ideals by \(\Upsilon(A)\).

If \(\mathfrak l\) is an open ideal of \(A\) and
\(\mathfrak m\in\Upsilon(A)\), the localization
\((A/\mathfrak l)_{\mathfrak m}\) is a local ring when
\(\mathfrak m\) contains \(\mathfrak l\), and is zero otherwise. Since
\(A/\mathfrak l\) is Artinian, it is a direct product of finitely many
local rings; thus
\[
  A/\mathfrak l
  \simeq
  \prod_{\mathfrak m\in\Upsilon(A)}
    (A/\mathfrak l)_{\mathfrak m}.
\]
It follows that there are ``canonical'' isomorphisms
\[
  A
  \simeq \varprojlim_{\mathfrak l}(A/\mathfrak l)
  \simeq \varprojlim_{\mathfrak l}
    \prod_{\mathfrak m}(A/\mathfrak l)_{\mathfrak m}
  \simeq \prod_{\mathfrak m}
    \varprojlim_{\mathfrak l}(A/\mathfrak l)_{\mathfrak m}
  \simeq \prod_{\mathfrak m}A_{\mathfrak m},
\]
where
\[
  A_{\mathfrak m}
  =
  \varprojlim_{\mathfrak l}(A/\mathfrak l)_{\mathfrak m}.
\]
This \emph{local component} \(A_{\mathfrak m}\) is a filtered
projective limit of Artinian local rings endowed with the discrete
topology. It is therefore a local ring which is pseudocompact for the
projective-limit topology.\footnote{Notice that \(A_{\mathfrak m}\) is
the localization of \(A\) at the maximal ideal \(\mathfrak m\). Indeed,
the identity element \(e_{\mathfrak m}\) of \(A_{\mathfrak m}\) is an
idempotent of \(A\) not belonging to \(\mathfrak m\). Since
\(A=A_{\mathfrak m}\times B\), where
\(B=A(1-e_{\mathfrak m})\), the ring \(A_{\mathfrak m}\) identifies with
the localization \(A_{e_{\mathfrak m}}\), and hence also with
\(S^{-1}A\), where \(S=A\setminus\mathfrak m\). Furthermore,
\(e_{\mathfrak m}(1-e_{\mathfrak m})=0\), so \(\mathfrak m\) contains
\(1-e_{\mathfrak m}\), hence also \(B\); thus \(\mathfrak m\) identifies
with \(\mathfrak mA_{\mathfrak m}\times B\).}

\medskip
\noindent\textbf{0.1.2.}
\phantomsection
\label{sga3:VIIB:0.1.2}
Let \(r(A)\) be the intersection of the open maximal ideals of \(A\),
that is, the cartesian product of the ideals
\(\mathfrak mA_{\mathfrak m}\) under the identification
\[
  A=\prod_{\mathfrak m}A_{\mathfrak m}.
\]
For every open ideal \(\mathfrak l\) of \(A\), the image of \(r(A)\) in
\(A/\mathfrak l\) is contained in the radical of \(A/\mathfrak l\).
Some power of this image is therefore zero, so
\(r(A)^n\subset\mathfrak l\) for all sufficiently large \(n\). Thus the
sequence \(r(A)^n\) tends to \(0\).

The same is true of \(x^n\) when \(x\in r(A)\). In other words, every
element of \(r(A)\) is topologically nilpotent, and the converse is
clear. It follows that the sequence with general term
\[
  1+x+\cdots+x^n
\]
converges, and converges to \(1/(1-x)\), whenever \(x\in r(A)\). This
shows that \(r(A)\) is the Jacobson radical of \(A\), that is, the
intersection of all maximal ideals of \(A\) (cf.\ Bourbaki,
\emph{Alg\`ebre}, Chap.~8, \S6, Th.~1).\footnote{Indeed, let
\(x\in r(A)\). If \(\mathfrak m\) were a maximal ideal not containing
\(x\), there would be \(y\in A\) such that
\(1-xy\in\mathfrak m\). But \(xy\in r(A)\), so \(1-xy\) is invertible,
a contradiction. Notice the following consequence: if
\(\Upsilon(A)=\{\mathfrak m_1,\ldots,\mathfrak m_r\}\) is finite, then
the \(\mathfrak m_i\) are all the maximal ideals of \(A\).}

\medskip
\noindent\textbf{Remarks.}\footnote{The editors added these remarks in
order to compare the definition of the formal spectrum
\(\operatorname{Spf}(A)\) given in 1.1 with those in EGA~I, 10.1.2 and
10.4.2.}

\noindent\textit{a)} If \(\mathfrak p\) is an open prime ideal of \(A\),
then, since \(A/\mathfrak p\) is Artinian, \(\mathfrak p\) is maximal.
Consequently, \(\Upsilon(A)\) is the set of open prime ideals of \(A\).

\smallskip
\noindent\textit{b)} Each \(\mathfrak mA_{\mathfrak m}\) is an ideal of
definition of \(A_{\mathfrak m}\), that is, an open ideal \(I\) for
which \(I^n\) tends to \(0\) (cf.\ EGA~0\textsubscript{I}, 7.1.2).
Consequently,
\[
  \operatorname{Spec}
    (A_{\mathfrak m}/\mathfrak mA_{\mathfrak m}),
\]
endowed with the topological ring \(A_{\mathfrak m}\), is an affine
formal scheme in the sense of EGA~I, 10.1.2.

\smallskip
\noindent\textit{c)} The topological ring \(A\) is admissible in the
sense of EGA~0\textsubscript{I}, 7.1.2 if and only if \(r(A)\) is open
(and hence an ideal of definition), and this holds if and only if
\(\Upsilon(A)\) is finite. In that case the affine formal scheme
\[
  \operatorname{Spf}(A)
  =
  \operatorname{Spec}(A/r(A))
\]
(cf.\ EGA~I, 10.1.2) has \(\Upsilon(A)\), endowed with the discrete
topology, as its underlying space; the sections of its structure sheaf
over a subset \(E\subset\Upsilon(A)\) form the ring
\[
  \prod_{\mathfrak m\in E}A_{\mathfrak m}.
\]

\smallskip
\noindent\textit{d)} Let \(A\) be an arbitrary pseudocompact ring. In
\SGARefLink{sga3:VIIB:1.1}{1.1} below, the space \(\Upsilon(A)\) is endowed with the discrete
topology and with the sheaf of rings whose ring of sections over any
subset \(E\) is
\[
  \prod_{\mathfrak m\in E}A_{\mathfrak m}.
\]
By \textit{b)}, every point then has an open neighborhood which is an
affine formal scheme; this therefore defines a formal scheme (EGA~I,
10.4.2), denoted by \(\operatorname{Spf}(A)\). For this formal scheme to
be affine, it must be quasi-compact, so \(\Upsilon(A)\) must be finite;
in that case \(\operatorname{Spf}(A)\) agrees with the definition of
EGA~I, 10.1.2.

% Authority: Exp7B local p. 3; combined-reader p. 516; printed p. 505.
\medskip
\noindent\textbf{0.1.3.}
\phantomsection
\label{sga3:VIIB:0.1.3}
If \(A\) and \(B\) are pseudocompact rings, a \emph{homomorphism} from
\(A\) to \(B\) is, by definition, a continuous map compatible with
addition, multiplication, and the identity elements. Such a
homomorphism sends a topologically nilpotent element of \(A\) to a
topologically nilpotent element of \(B\); it therefore maps the radical
\(r(A)\) into the radical \(r(B)\).

% Next cursor: Exp7B local p. 3, combined-reader p. 516,
% re-edition printed p. 505, at 0.2.

% END INCLUDED SGA3_English_current_buildable_source_20260823/source/inputs/VII/tex/components/15_expose_VIIB_title_section01_pseudocompact_rings_en.tex

% BEGIN INCLUDED SGA3_English_current_buildable_source_20260823/source/inputs/VII/tex/components/16_expose_VIIB_section02_statements_en.tex
% SGA 3, Expose VII_B, English reconstruction -- Loop 1.
% Scope: authority local pages 3--4, 0.2 through Corollary 0.2.D,
% stopping immediately before the proof of Corollary 0.2.D.
% Combined-reader pages: 516--517. Re-edition printed pages: 505--506.
% Original-edition marginal locator in this span: 480.
% Authority: Polo--Gille Exp7B-13oct24.pdf,
% SHA-256 D303690E0195654BBE7C3DAA5D89EA23E24CD709A3575217234AC19AE8836694.
% Jacob Reinhold's CC BY 4.0 Markdown was consulted as comparison lineage.

\medskip
\noindent\textbf{0.2.}
\phantomsection
\label{sga3:VIIB:0.2}
Let \(A\) be a commutative pseudocompact ring. A
\emph{pseudocompact \(A\)-module} \(M\) is a separated, complete
topological \(A\)-module having a basis of neighborhoods of \(0\)
consisting of submodules \(M'\) for which \(M/M'\) has finite length
over \(A\).

If \(M\) and \(N\) are pseudocompact \(A\)-modules, a
\emph{morphism} from \(M\) to \(N\) is, by definition, a continuous
\(A\)-linear map. We write
\[
  \operatorname{Hom}_c(M,N)
\]
for the group of these morphisms.

\medskip
\noindent\SGAInlineRefTarget{sga3:vii:proposition:0-2-b}\textbf{Proposition 0.2.B.}\footnote{The editors have stated
the results of this paragraph separately in the following proposition,
and have then indicated the steps of the proof; cf.\ \textup{[CA]},
\S~IV.3, or \textup{[DG70]}, \S~V.2.}
\phantomsection
\label{sga3:VIIB:0.2.B}
{\itshape
\begin{enumerate}
  \item[\textup{(i)}]
  The pseudocompact \(A\)-modules form an abelian category, denoted
  \(\operatorname{PC}(A)\). In particular, for every morphism
  \(f:M\to N\), the submodule \(\operatorname{Im}(f)\) is complete and
  hence closed in \(N\).

  \item[\textup{(ii)}]
  The pseudocompact \(A\)-modules of finite length, whose topology is
  therefore discrete, form a system of cogenerators of
  \(\operatorname{PC}(A)\).

  \item[\textup{\SGAInlineRefTarget{sga3:vii:proposition:0-2-b:item:iii}(iii)}]
  Infinite products and filtered projective limits are exact; in other
  words, \(\operatorname{PC}(A)\) satisfies axiom
  \((\mathrm{AB}5^{*})\).\footnote{Cf.\ \textup{[Gr57]}, I, \S~\SGARefLink{sga3:VIIA:1.5}{1.5}
  and Proposition~1.8. One may also consult, for example,
  \textup{[Po73]}, \S~\SGARefLink{sga3:VIIB:2.8}{2.8}, or \textup{[We95]}, Appendix~A.4.}
\end{enumerate}}

For the reader's convenience, we briefly indicate the steps of the
proof. First we use the following lemma (\textup{[CA]}, \S~IV.3,
Lemma~1; for the proof, see \textup{[BEns]}, III, \S~\SGARefLink{sga3:VIIA:7.4}{7.4}, Theorem~1
and Example~2).

\medskip
\noindent\textbf{Lemma 0.2.C.}
\phantomsection
\label{sga3:VIIB:0.2.C}
\textit{Let \(B\) be a ring, let \(I\) be a filtered ordered set, and
let \((M_i)\) and \((N_i)\) be two projective systems of left
\(B\)-modules indexed by \(I\). Let
\[
  (s_i):(M_i)\longrightarrow(N_i)
\]
be a morphism of projective systems such that \(s_i\) is surjective
and has Artinian kernel for every \(i\). Then the map}
\[
  \varprojlim_i s_i:
  \varprojlim_i M_i
  \longrightarrow
  \varprojlim_i N_i
\]
\textit{is surjective.}

\medskip
\noindent\textbf{Corollary 0.2.D}
\textit{(\textup{[CA]}, \S~IV.3, Propositions~10 and~11).}
\phantomsection
\label{sga3:VIIB:0.2.D}
{\itshape Let \(M\) be a pseudocompact \(A\)-module.
\begin{enumerate}
  \item[\textup{(i)}]
  If \(K\) is a closed submodule of \(M\), then \(M/K\), endowed with
  the quotient topology, is a pseudocompact \(A\)-module.

  \item[\textup{(ii)}]
  Let \((M_i)\) be a filtered decreasing family of closed submodules
  of \(M\).
  \begin{enumerate}
    \item[\textup{\SGAInlineRefTarget{sga3:VIIB:0.2.D:item:a}(a)}]
    The canonical map
    \[
      M\longrightarrow\varprojlim_i M/M_i
    \]
    is surjective and has kernel \(\bigcap_i M_i\).

    \item[\textup{\SGAInlineRefTarget{sga3:VIIB:0.2.D:item:b}(b)}]
    For every closed submodule \(N\) of \(M\), one has
    \[
      N+\bigcap_i M_i
      =
      \bigcap_i(N+M_i).
    \]
  \end{enumerate}
\end{enumerate}}

% Next cursor: Exp7B local p. 4, combined-reader p. 517,
% re-edition printed p. 506, at the proof of Corollary 0.2.D.

% END INCLUDED SGA3_English_current_buildable_source_20260823/source/inputs/VII/tex/components/16_expose_VIIB_section02_statements_en.tex

% BEGIN INCLUDED SGA3_English_current_buildable_source_20260823/source/inputs/VII/tex/components/17_expose_VIIB_section02D_proof_en.tex
% SGA 3, Expose VII_B, English reconstruction -- Loop 1.
% Scope: authority local pages 4--6, proof of Corollary 0.2.D
% through the end of the proof of Proposition 0.2.B.
% Combined-reader pages: 517--519. Re-edition printed pages: 506--508.
% Authority: Polo--Gille Exp7B-13oct24.pdf,
% SHA-256 D303690E0195654BBE7C3DAA5D89EA23E24CD709A3575217234AC19AE8836694.
% Jacob Reinhold's CC BY 4.0 Markdown was consulted as comparison lineage.

\medskip
\noindent\textit{Proof.}
\phantomsection
\label{sga3:VIIB:0.2.D-proof}
Let \((L_j)\) be the filtered decreasing family of open submodules of
\(M\). We endow \(M/K\) with the quotient topology; that is, a basis
of neighborhoods of \(0\) is formed by the open submodules
\((K+L_j)/K\). Since \(K\) is closed, it is equal to the intersection
of the \(K+L_j\), and hence the map
\[
  \tau:M/K\longrightarrow\varprojlim_j M/(K+L_j)
\]
is injective. It is also open, since the right-hand side is the
inverse limit of the discrete modules \(M/(K+L_j)\). Moreover, for
each \(j\), the map
\[
  t_j:M/L_j\longrightarrow M/(K+L_j)
\]
is surjective and has Artinian kernel. Thus, by the preceding lemma,
the map \(t\) in the following commutative diagram is surjective:
\begin{center}
  \phantomsection
  \label{sga3:VIIB:0.2.D-quotient-completion-square}
  % SGA3-VII-DIAG-079
  % BEGIN INCLUDED SGA3_English_current_buildable_source_20260823/source/inputs/VII/native_diagrams/exp7/viib_15_95/diag_079.tex
% SGA3-VII-DIAG-079
% Authority: Exp7B-13oct24.pdf, local p. 4.
% Crop: figures/Exp7B_localp004_cor02D_quotient_completion_square.png
% Locator: tex/components/17_expose_VIIB_section02D_proof_en.tex, Loop 1 line 32.
\begin{tikzcd}[row sep=2.8em,column sep=4.6em]
M \arrow[r,"p","{\sim}"'] \arrow[d]
  & \varprojlim_j M/L_j \arrow[d,"t"] \\
M/K \arrow[r,"\tau"]
  & \varprojlim_j M/(K+L_j).
\end{tikzcd}

% END INCLUDED SGA3_English_current_buildable_source_20260823/source/inputs/VII/native_diagrams/exp7/viib_15_95/diag_079.tex

\end{center}
Since \(p\) is an isomorphism because \(M\) is complete, it follows
that \(\tau\) is surjective, and hence is an isomorphism. This proves
\textup{\SGARefLink{sga3:vii:paragraph:5-1:equation:i}{(i)}}.

% Authority: Exp7B local p. 5; combined-reader p. 518;
% re-edition printed p. 507.
\phantomsection
\label{sga3:VIIB:0.2.D-proof-ii-a}
Let us prove \textup{\SGARefLink{sga3:vii:paragraph:5-1:equation:ii}{(ii)}(a)}. By the foregoing, for every \(i\) there
is an isomorphism
\[
  M/M_i\xrightarrow{\sim}\varprojlim_j M/(M_i+L_j),
\]
so the two horizontal arrows in the following commutative diagram are
isomorphisms:
\begin{center}
  \phantomsection
  \label{sga3:VIIB:0.2.D-filtered-family-square}
  % SGA3-VII-DIAG-080
  % BEGIN INCLUDED SGA3_English_current_buildable_source_20260823/source/inputs/VII/native_diagrams/exp7/viib_15_95/diag_080.tex
% SGA3-VII-DIAG-080
% Authority: Exp7B-13oct24.pdf, local p. 5.
% Crop: figures/Exp7B_localp005_cor02D_filtered_family_square.png
% Locator: tex/components/17_expose_VIIB_section02D_proof_en.tex, Loop 1 line 53.
\begin{tikzcd}[row sep=2.8em,column sep=4.6em]
M \arrow[r,"p","{\sim}"'] \arrow[d,"g"]
  & \varprojlim_j M/L_j \arrow[d,"s"] \\
\varprojlim_i M/M_i \arrow[r,"{\sim}"']
  & \varprojlim_{i,j} M/(M_i+L_j).
\end{tikzcd}

% END INCLUDED SGA3_English_current_buildable_source_20260823/source/inputs/VII/native_diagrams/exp7/viib_15_95/diag_080.tex

\end{center}
Moreover, for each \(j\), the family of submodules
\((M_i+L_j)/L_j\) has a least element, since \(M/L_j\) is Artinian.
Consequently, the morphism
\[
  s_j:M/L_j\longrightarrow\varprojlim_i M/(L_j+M_i)
\]
is surjective; therefore, by the preceding lemma, \(s\) is
surjective. It follows that \(g\) is surjective. Finally, the kernel
of \(g\) is the inverse limit of the \(M_i\), that is, their
intersection. This proves \textup{\SGARefLink{sga3:VIIB:0.2.D:item:a}{(a)}}.

\phantomsection
\label{sga3:VIIB:0.2.D-proof-ii-b}
Let us deduce \textup{(b)} from this. Since \(N\) is a closed
submodule, and hence separated and complete, it is a pseudocompact
module for the topology induced by that of \(M\). Therefore, by
\textup{(a)}, the morphisms \(f\) and \(g\) in the following
commutative exact diagram are surjective:
\begin{center}
  \phantomsection
  \label{sga3:VIIB:0.2.D-snake-exact-diagram}
  % SGA3-VII-DIAG-081
  % BEGIN INCLUDED SGA3_English_current_buildable_source_20260823/source/inputs/VII/native_diagrams/exp7/viib_15_95/diag_081.tex
% SGA3-VII-DIAG-081
% Authority: Exp7B-13oct24.pdf, local p. 5.
% Crop: figures/Exp7B_localp005_cor02D_snake_exact_diagram.png
% Locator: tex/components/17_expose_VIIB_section02D_proof_en.tex, Loop 1 line 77.
\begin{tikzcd}[row sep=2.8em,column sep=2.15em]
0 \arrow[r] & N \arrow[r] \arrow[d,"f"]
  & M \arrow[r] \arrow[d,"g"]
  & M/N \arrow[r] \arrow[d,"h"]
  & 0 \\
0 \arrow[r] & \varprojlim_i N/(N\cap M_i) \arrow[r]
  & \varprojlim_i M/M_i \arrow[r]
  & \varprojlim_i M/(N+M_i)\, . & {}
\end{tikzcd}

% END INCLUDED SGA3_English_current_buildable_source_20260823/source/inputs/VII/native_diagrams/exp7/viib_15_95/diag_081.tex

\end{center}
The ``snake lemma'' then shows that the sequence
\[
  0\longrightarrow\ker f\longrightarrow\ker g
  \longrightarrow\ker h\longrightarrow 0
\]
is exact, and the equality
\[
  N+\bigcap_i M_i=\bigcap_i(N+M_i)
\]
follows.

\phantomsection
\label{sga3:VIIB:0.2.B-proof}
We can now prove \SGARefLink{sga3:vii:proposition:0-2-b}{Proposition~0.2.B}. Let \(f:M\to N\) be a morphism of
pseudocompact \(A\)-modules. Then \(K=\ker(f)\) is a closed submodule
of \(M\), and hence is separated and complete; thus \(K\), with the
topology induced by that of \(M\), is a pseudocompact module. By
\SGARefLink{sga3:VIIB:0.2.D}{0.2.D}\,\textup{\SGARefLink{sga3:vii:paragraph:5-1:equation:i}{(i)}}, the module \(M/K\), endowed with the quotient
topology, is pseudocompact.

Let us show that the continuous bijective morphism
\[
  M/K\longrightarrow\operatorname{Im}(f)
\]
is bicontinuous. After identifying \(M/K\) with
\(\operatorname{Im}(f)\), it is enough to show that the topology
\(T\) induced by that of \(N\) is finer than the quotient topology
\(Q\).\footnote{The French re-edition says here that \(Q\) is finer
than \(T\); the ensuing argument proves the reverse inclusion of open
sets, which is the one needed for bicontinuity.} Let \((L_j)\),
respectively \((N_i)\), be the filtered
decreasing families of open submodules of \(M\), respectively \(N\),
and put
\[
  N_i'=N_i\cap\operatorname{Im}(f).
\]
Let \(P=(K+L_j)/K\) be a submodule of \(M/K\) open for \(Q\). Since
\(M/(K+L_j)\) is Artinian, the family \(N_i'+P\) has a least element
\(N_{i_0}'+P\). The \(N_i'\) are open, and hence closed, for \(T\),
and therefore also closed for \(Q\). It follows from
\SGARefLink{sga3:VIIB:0.2.D}{0.2.D}\,\textup{\SGARefLink{sga3:vii:paragraph:5-1:equation:ii}{(ii)}\SGARefLink{sga3:VIIB:0.2.D:item:b}{(b)}} that
\[
  N_{i_0}'+P
  =\bigcap_i(N_i'+P)
  =P+\bigcap_i N_i'
  =P,
\]
whence \(N_{i_0}'\subset P\). Thus \(P\) is open for \(T\), and
\(M/K\to\operatorname{Im}(f)\) is consequently an isomorphism of
pseudocompact modules.

In particular, \(\operatorname{Im}(f)\) is complete for \(T\), and
hence closed in \(N\). By \SGARefLink{sga3:VIIB:0.2.D}{0.2.D}\,\textup{\SGARefLink{sga3:vii:paragraph:5-1:equation:i}{(i)}} once more,
\(\operatorname{Coker}(f)\), endowed with the quotient topology, is
pseudocompact. This proves that \(\operatorname{PC}(A)\) is an
abelian category.

Arbitrary inverse limits also exist in \(\operatorname{PC}(A)\): if
\((M_i)\) is an inverse system of pseudocompact modules, its inverse
limit has as
% Authority: Exp7B local p. 6; combined-reader p. 519;
% re-edition printed p. 508.
underlying module the inverse limit of the underlying
modules, endowed with the inverse-limit topology. Moreover, if one
has a family of exact sequences in \(\operatorname{PC}(A)\),
\[
  0\longrightarrow K_i\longrightarrow M_i
  \longrightarrow Q_i\longrightarrow 0,
\]
then the sequence
\[
  0\longrightarrow\prod_i K_i\longrightarrow\prod_i M_i
  \longrightarrow\prod_i Q_i\longrightarrow 0
\]
is exact. Assertion \textup{\SGARefLink{sga3:vii:proposition:0-2-b:item:iii}{(iii)}} of \SGARefLink{sga3:vii:proposition:0-2-b}{0.2.B} follows, since in every
abelian category in which arbitrary products exist, conditions
\textup{\SGARefLink{sga3:VIIB:0.2.D:item:a}{(a)}} and \textup{\SGARefLink{sga3:VIIB:0.2.D:item:b}{(b)}} of \SGARefLink{sga3:VIIB:0.2.D}{0.2.D} are equivalent, and are
equivalent to the exactness of filtered inverse limits
(\textup{cf.}\ [Mi65], III, \SGARefLink{sga3:iii:subsection:1-2}{1.2}--1.9, or [Po73], Chap.~2,
Th.~8.6).

Finally, every pseudocompact module \(M\) is a submodule of the
product
\[
  \prod_L M/L,
\]
where \(L\) ranges over the open submodules of \(M\). Thus the objects
of finite length form a system of cogenerators of
\(\operatorname{PC}(A)\). (Moreover, every object of length \(n\) is
isomorphic to a quotient \(A^n/L\), where \(L\) is an open submodule
of \(A^n\) of colength \(n\); these quotients therefore form a
\emph{set} of cogenerators.) This completes the proof of \SGARefLink{sga3:vii:proposition:0-2-b}{0.2.B}.

% Next cursor: Exp7B local p. 6, combined-reader p. 519,
% re-edition printed p. 508, at the editors' note (10) and the paragraph
% beginning ``Let (Ab) be the category of abelian groups''. The revised
% authority inserts 0.2.E--0.2.G here, before 0.2.1.

% END INCLUDED SGA3_English_current_buildable_source_20260823/source/inputs/VII/tex/components/17_expose_VIIB_section02D_proof_en.tex

% BEGIN INCLUDED SGA3_English_current_buildable_source_20260823/source/inputs/VII/tex/components/18_expose_VIIB_sections021_022_duality_en.tex
% SGA 3, Expose VII_B, English reconstruction -- Loop 1.
% Scope: authority local pages 6--8, immediately after the proof of 0.2.B
% through the end of 0.2.2, stopping immediately before 0.3.
% Combined-reader pages: 519--521. Re-edition printed pages: 508--510.
% Original-edition marginal locators in this span: 481--482.
% Authority: Polo--Gille Exp7B-13oct24.pdf,
% SHA-256 D303690E0195654BBE7C3DAA5D89EA23E24CD709A3575217234AC19AE8836694.
% Jacob Reinhold's CC BY 4.0 Markdown was consulted as comparison lineage.

\medskip
Let \((\mathrm{Ab})\)\footnote{\label{sga3:VIIB:nde:10}The editors
inserted \SGARefLink{sga3:VIIB:0.2.E}{Proposition 0.2.E} and \SGARefLink{sga3:VIIB:0.2.F}{Corollary 0.2.F} into this paragraph;
they will be useful in \SGARefLink{sga3:VIIB:0.2.2}{0.2.2}. In the original, these results appeared
in \SGARefLink{sga3:VIIB:0.3}{0.3}.} be the category of abelian groups and let
\(\operatorname{LF}(A)\) be the full subcategory of
\(\operatorname{PC}(A)\) formed by the objects of finite length. For
every object \(M\) of \(\operatorname{PC}(A)\), let \(h_c^M\) denote
the functor
\[
  \operatorname{LF}(A)\longrightarrow(\mathrm{Ab}),
  \qquad
  N\longmapsto\operatorname{Hom}_c(M,N).
\]
By \textup{[CA]}, \S~II.4, Theorem~1, Lemma~4, and Corollary~1, we have
the following results.\footnote{\label{sga3:VIIB:nde:11}Indeed,
\(\operatorname{PC}(A)\) has a set of Artinian cogenerators, filtered
projective limits are exact in it, and \(\operatorname{LF}(A)\) is the
subcategory of Artinian objects. The dual category
\(\operatorname{PC}(A)^\circ\) therefore has a set of Noetherian
generators, and filtered inductive limits are exact in it. By the proof
of \textup{[CA]}, \S~II.4, Theorem~1, the functor
\(M\mapsto\operatorname{Hom}_c(M,-)\) is an anti-equivalence from
\(\operatorname{PC}(A)\) to
\(\operatorname{Lex}(\operatorname{LF}(A),(\mathrm{Ab}))\).
Likewise, Lemma~4 and Corollary~1 of \textit{loc.\ cit.} state results
``dual'' to those of \SGARefLink{sga3:VIIB:0.2.F}{Corollary~0.2.F}. For a ``direct'' proof, see
\textup{[DG70]}, \S~V.2, Theorem~3.1, Lemma~3.5, and
Corollaries~3.3 and~3.4.}

\medskip
\noindent\textbf{Proposition 0.2.E.}
\phantomsection
\label{sga3:VIIB:0.2.E}
\textit{The functor \(M\mapsto h_c^M\) is an anti-equivalence from
\(\operatorname{PC}(A)\) to the category}
\[
  \operatorname{Lex}\bigl(\operatorname{LF}(A),(\mathrm{Ab})\bigr)
\]
\textit{of left-exact functors
\(\operatorname{LF}(A)\longrightarrow(\mathrm{Ab})\).}

\medskip
\noindent\textbf{Corollary 0.2.F.}
\phantomsection
\label{sga3:VIIB:0.2.F}
{\itshape
\begin{enumerate}
  \item[\textup{(i)}]
  An object \(P\) of \(\operatorname{PC}(A)\) is projective if and
  only if the functor \(h_c^P\) is exact; that is, if and only if the
  functor \(\operatorname{Hom}_c(P,-)\) is exact on
  \(\operatorname{LF}(A)\).

  \item[\textup{\SGAInlineRefTarget{sga3:VIIB:0.2.F:item:ii}(ii)}]
  Let \((M_i)\) be a filtered projective
  system\footnote{\label{sga3:VIIB:nde:12}Since every infinite product
  is a filtered projective limit of finite products, every additive
  functor that commutes with filtered projective limits also commutes
  with infinite products.} of objects of \(\operatorname{PC}(A)\).
  For every object \(N\) of \(\operatorname{LF}(A)\), there is an
  isomorphism, functorial in \(N\),
  \[
    \operatorname{Hom}_c\!\left(\varprojlim_i M_i,N\right)
    \simeq
    \varinjlim_i\operatorname{Hom}_c(M_i,N).
  \]

  \item[\textup{\SGAInlineRefTarget{sga3:VIIB:0.2.F:item:iii}(iii)}]
  Every filtered projective limit and every
  product\footnotemark[\value{footnote}] of projective objects of
  \(\operatorname{PC}(A)\) is a projective object of
  \(\operatorname{PC}(A)\).
\end{enumerate}}

Finally, \SGARefLink{sga3:VIIB:0.2.F}{Corollary~0.2.F} gives the following result.

\medskip
\noindent\textbf{Corollary 0.2.G.}
\phantomsection
\label{sga3:VIIB:0.2.G}
\textit{Let \((M_i)_{i\in I}\) be a family of objects of
\(\operatorname{PC}(A)\). Then \(\prod_{i\in I}M_i\) is projective if
and only if each \(M_i\) is projective.}

Indeed, for every \(N\in\operatorname{Ob}\operatorname{LF}(A)\), one
has
\[
  \operatorname{Hom}_c\!\left(\prod_i M_i,N\right)
  \simeq
  \bigoplus_i\operatorname{Hom}_c(M_i,N).
\]

\medskip
\noindent\textbf{0.2.1.}
\phantomsection
\label{sga3:VIIB:0.2.1}
Each local component \(A_{\mathfrak m}\) of \(A\) is a direct factor
of \(A\), hence a projective object of \(\operatorname{PC}(A)\)
(\(A\) is manifestly projective). Moreover,
\[
  S_{\mathfrak m}
  =
  A_{\mathfrak m}/\mathfrak mA_{\mathfrak m}
\]
is the unique simple quotient of \(A_{\mathfrak m}\), so
\(A_{\mathfrak m}\) is indecomposable. On the other hand, every simple
object of \(\operatorname{PC}(A)\) is isomorphic to a unique
\(S_{\mathfrak m}\). By \textup{[CA]}, IV, \S~3, Corollary~1 of
Theorem~3,\footnote{\label{sga3:VIIB:nde:13}This refers to the
``dual'' statements established in \textit{loc.\ cit.}, \S~2,
Theorem~2, and \S~1, Proposition~2. For a ``direct'' proof, see
\textup{[DG70]}, V, \S~2, Theorem~4.5 and Example~4.6(b).} we
therefore have:

\begin{samepage}
\medskip
\noindent\textbf{Proposition.}
\phantomsection
\label{sga3:VIIB:0.2.1:projectives}
{\itshape
\begin{enumerate}
  \item[\textup{(i)}]
  Every projective object of \(\operatorname{PC}(A)\) is a direct
  product of indecomposable projective objects, uniquely determined up
  to isomorphism.

  \item[\textup{\SGAInlineRefTarget{sga3:VIIB:0.2.1:item:ii}(ii)}]
  Every indecomposable projective object is isomorphic to a unique
  \(A_{\mathfrak m}\), where
  \(\mathfrak m\in\Upsilon(A)\).
\end{enumerate}}
\end{samepage}

\medskip
\noindent\textbf{Definition.}
\phantomsection
\label{sga3:VIIB:0.2.1:topologically-free}
A pseudocompact \(A\)-module \(M\) is said to be
\emph{topologically free} if it is isomorphic to the product of a
family \((A_i)\) of copies of \(A\).

In this case, a family \((m_i)\) of elements of \(M\) is called a
\emph{pseudobasis} of \(M\) if the \(A\)-linear maps
\(A_i\longrightarrow M\) sending the unit element of \(A_i\) to
\(m_i\) extend to an isomorphism
\[
  \prod_i A_i\xrightarrow{\ \sim\ }M.
\]

\medskip
\noindent\textbf{0.2.2.}
\phantomsection
\label{sga3:VIIB:0.2.2}
\footnote{\label{sga3:VIIB:nde:14}The editors have spelled out the
results of this paragraph, which play an important role below; cf.\
\SGARefLink{sga3:VIIB:1.2.3}{1.2.3}, \SGARefLink{sga3:VIIB:1.3.5}{1.3.5}, \SGARefLink{sga3:VIIB:2.2.1}{2.2.1}, and so on.}
If \(M\) is a pseudocompact \(A\)-module, we write \(M^\dagger\) for
the \(A\)-module, without topology,
\[
  \operatorname{Hom}_c(M,A).
\]

Conversely, if \(N\) is an \(A\)-module, we write
\[
  N^*=\operatorname{Hom}_A(N,A)
\]
for its dual, endowed with the topology of pointwise convergence. Thus
a basis of neighborhoods of \(0\) in \(N^*\) is formed by the
submodules
\[
  \mathcal V(x,\mathfrak l)
  =
  \{f\in N^*\mid f(x)\in\mathfrak l\},
\]
where \(x\in N\) and \(\mathfrak l\) is an open ideal of \(A\).

This makes \(N^*\) a pseudocompact \(A\)-module. Indeed, first suppose
that \(N=A\). Then \(N^*=A\), endowed with its topology as a
pseudocompact ring. If \(N\) is a free module \(A^{(I)}\), then
\(N^*\) is the product \(A^I\), endowed with the product topology.
Furthermore, for every morphism
\(\varphi:N_1\longrightarrow N_2\), the transpose
\[
  {}^t\!\varphi:N_2^*\longrightarrow N_1^*
\]
is continuous, because the inverse image under \({}^t\!\varphi\) of
a submodule \(\mathcal V(x_1,\mathfrak l)\) of \(N_1^*\) is precisely
the submodule \(\mathcal V(\varphi(x_1),\mathfrak l)\) of \(N_2^*\).
For arbitrary \(N\), choose a presentation
\[
  A^{(J)}
  \xrightarrow{\ \varphi\ }
  A^{(I)}
  \xrightarrow{\ \pi\ }
  N
  \longrightarrow
  0.
\]
Then \(N^*\) is the kernel of the continuous morphism
\[
  {}^t\!\varphi:A^I\longrightarrow A^J,
\]
and hence \(N^*\) is pseudocompact.

When \(A\) is Artinian (in which case one may take
\(\mathfrak l=0\) above), \SGARefLink{sga3:VIIB:0.2.F}{Corollary~0.2.F} gives:

\medskip
\noindent\textbf{Proposition.}
\phantomsection
\label{sga3:VIIB:0.2.2:artinian-duality}
\textit{When \(A\) is Artinian, the functors}
\[
  P\longmapsto P^\dagger
  \qquad\textit{and}\qquad
  Q\longmapsto Q^*,
\]
\textit{where \(P\), respectively \(Q\), is a projective object of
\(\operatorname{PC}(A)\), respectively a projective \(A\)-module,
establish an anti-equivalence between the category of projective
pseudocompact \(A\)-modules and the category of projective
\(A\)-modules.}\footnote{\label{sga3:VIIB:nde:15}Recall also that,
over an Artinian ring, a module is projective if and only if it is
flat; see, for example, Expos\'e VI\textsubscript{B}, editors'
note~(54), or \SGARefLink{sga3:VIIB:0.3.7}{0.3.7} below.}

In particular, when \(A\) is a field \(k\), the functor
\[
  P\longmapsto P^\dagger
\]
is an anti-equivalence between the category of all pseudocompact
\(k\)-modules (also called \emph{linearly compact \(k\)-vector
spaces}) and the category of \(k\)-vector spaces.\footnote{
\label{sga3:VIIB:nde:16}The direct sum in \(\operatorname{PC}(k)\) of
a family \((V_i)_{i\in I}\) of linearly compact \(k\)-vector spaces
is \(\left(\prod_{i\in I}V_i^\dagger\right)^*\).}

% Next cursor: Exp7B local p. 8, combined-reader p. 521,
% re-edition printed p. 510, immediately before 0.3.

% END INCLUDED SGA3_English_current_buildable_source_20260823/source/inputs/VII/tex/components/18_expose_VIIB_sections021_022_duality_en.tex

% BEGIN INCLUDED SGA3_English_current_buildable_source_20260823/source/inputs/VII/tex/components/19_expose_VIIB_sections03_034_completed_tensor_en.tex
% SGA 3, Expose VII_B, English reconstruction -- Loop 1.
% Scope: authority local pages 8--11, 0.3 through 0.3.4 inclusive,
% stopping immediately before 0.3.5 on local page 11.
% Combined-reader pages: 521--524. Re-edition printed pages: 510--513.
% Original-edition marginal locators within this span: 482--484.
% Authority: Polo--Gille Exp7B-13oct24.pdf,
% SHA-256 D303690E0195654BBE7C3DAA5D89EA23E24CD709A3575217234AC19AE8836694.
% Jacob Reinhold's CC BY 4.0 Markdown was consulted as comparison lineage.
% No genuine diagram occurs in this span; all displayed mathematics is native LaTeX.

\medskip
\noindent\textbf{0.3.}\footnote{Editors' note: This paragraph has been
modified to take account of the additions made in \SGARefLink{sga3:VIIB:0.2}{0.2}.}
\phantomsection
\label{sga3:VIIB:0.3}
Let \(L\) and \(M\) be two pseudocompact \(A\)-modules. The functor
\[
  \operatorname{LF}(A)\longrightarrow(\mathrm{Ab}),
  \qquad
  N\longmapsto\operatorname{Bil}_{c}(L\times M,N),
\]
where \(\operatorname{Bil}_{c}(L\times M,N)\) denotes the abelian group
of continuous \(A\)-bilinear maps from \(L\times M\) to \(N\), is left
exact.

By \SGARefLink{sga3:VIIB:0.2.E}{0.2.E}, there therefore exists a pseudocompact \(A\)-module
\(L\widehat{\otimes}_{A}M\), unique up to unique isomorphism, representing
this functor; that is, for every object \(N\) of
\(\operatorname{LF}(A)\), there is a functorial isomorphism
\[
  \operatorname{Hom}_{c}
    \bigl(L\widehat{\otimes}_{A}M,N\bigr)
  \simeq
  \operatorname{Bil}_{c}(L\times M,N).
\]

Moreover, \(L\widehat{\otimes}_{A}M\) identifies with the projective
limit
\[
  P=\widehat{L\otimes_{A}M}
\]
of the discrete \(A\)-modules
\[
  (L/L')\otimes_{A}(M/M'),
\]
where \(L'\) and \(M'\) range respectively over the open submodules of
\(L\) and \(M\).

Indeed, let
\(\varphi:L\times M\to N\) be a continuous bilinear map into a discrete
\(A\)-module \(N\) of finite length. By \SGARefLink{sga3:VIIB:0.3.1}{Lemma~0.3.1} below, there exist
open submodules \(L'\) and \(M'\) of \(L\) and \(M\) such that
\[
  \varphi(L'\times M)=\varphi(L\times M')=\{0\}.
\]
This means that the map
\[
  \overline{\varphi}:L\otimes_{A}M\longrightarrow N
\]
induced by \(\varphi\) has the form \(\varphi'\circ p\), where \(p\) is
the canonical projection from \(L\otimes_{A}M\) onto
\((L/L')\otimes_{A}(M/M')\). If \(\widehat{\varphi}\) denotes the
composite
\[
  P
  \longrightarrow
  (L/L')\otimes_{A}(M/M')
  \xrightarrow{\ \varphi'\ }
  N,
\]
then the map \(\varphi\mapsto\widehat{\varphi}\) is a functorial
bijection from \(\operatorname{Bil}_{c}(L\times M,N)\) onto
\(\operatorname{Hom}_{c}(P,N)\). Hence
\[
  P\simeq L\widehat{\otimes}_{A}M.
\]

Thus the pseudocompact module \(L\widehat{\otimes}_{A}M\) is the
separated completion of \(L\otimes_{A}M\) for the linear topology
defined by the kernels of the canonical projections
\[
  L\otimes_{A}M
  \longrightarrow
  (L/L')\otimes_{A}(M/M'),
\]
and it will be called the \emph{completed tensor product} of \(L\) and
\(M\).

If \(x\in L\) and \(y\in M\), the image of \(x\otimes_{A}y\) in
\(L\widehat{\otimes}_{A}M\) will be denoted by
\(x\widehat{\otimes}_{A}y\).

\medskip
\noindent\textbf{Lemma 0.3.1.}
\phantomsection
\label{sga3:VIIB:0.3.1}
\textit{Let \(L\), \(M\), and \(N\) be pseudocompact \(A\)-modules,
with \(N\) of finite length. If
\(\varphi:L\times M\to N\) is a continuous \(A\)-bilinear map, there
exist open submodules \(L'\) and \(M'\) of \(L\) and \(M\) such that}
\[
  \varphi(L'\times M)=\varphi(L\times M')=\{0\}.
\]

\noindent\textit{Proof.}
The inverse image \(\varphi^{-1}(0)\) is an open neighborhood of
\((0,0)\), and therefore contains an open set of the form
\(L_{1}\times M_{1}\), where \(L_{1}\) and \(M_{1}\) are open
submodules of \(L\) and \(M\). Since \(L/L_{1}\) has finite length,
there are elements \(x_{1},\ldots,x_{r}\) of \(L\) such that
\[
  L_{1}+Ax_{1}+\cdots+Ax_{r}=L.
\]
If \(M'\subset M_{1}\) is sufficiently small, then
\(\varphi(x_i,M')=0\) for every \(i\), since the map
\(y\mapsto\varphi(x_i,y)\) is continuous. It follows that
\(\varphi(L,M')=\{0\}\). Likewise,
\(\varphi(L',M)=\{0\}\) when \(L'\) is sufficiently small.
\hfill\(\square\)

\medskip
\noindent\textbf{Corollary 0.3.1.1.}\footnote{Editors' note: This
corollary has been added.}
\phantomsection
\label{sga3:VIIB:0.3.1.1}
\textit{Let \(M\) be a pseudocompact \(A\)-module.}
\begin{enumerate}
  \item[\textup{(i)}]
  \textit{For every open submodule \(M'\), there exists an open ideal
  \(\mathfrak l\) of \(A\) such that}
  \[
    \overline{\mathfrak lM}\subset M'.
  \]

  \item[\textup{(ii)}]
  \textit{Consequently,}
  \[
    M\simeq
    \varprojlim_{\mathfrak l}
      M/\overline{\mathfrak lM},
  \]
  \textit{where \(\mathfrak l\) ranges over the filtered projective
  system of open ideals of \(A\), and every
  \(M/\overline{\mathfrak lM}\) is endowed with the quotient topology,
  which makes it a pseudocompact module; cf.\ \SGARefLink{sga3:VIIB:0.2.D}{0.2.D}.}
\end{enumerate}

\noindent\textit{Proof.}
Consider the map
\[
  \varphi:A\times M\longrightarrow M/M',
  \qquad
  (a,m)\longmapsto am+M'.
\]
By \SGARefLink{sga3:VIIB:0.3.1}{0.3.1}, there is an open ideal \(\mathfrak l\) of \(A\) such that
\(\mathfrak lM\subset M'\). Since \(M'\) is also closed, it contains
\(\overline{\mathfrak lM}\). Since the intersection of the open
submodules of \(M\) is zero, it follows that
\[
  \bigcap_{\mathfrak l}\overline{\mathfrak lM}=(0).
\]
On the other hand, by \SGARefLink{sga3:VIIB:0.2.D}{0.2.D} the map
\[
  \phi:M\longrightarrow
  \varprojlim_{\mathfrak l}M/\overline{\mathfrak lM}
\]
is surjective. By the proof of \SGARefLink{sga3:vii:proposition:0-2-b}{0.2.B}, \(\phi\) therefore induces an
isomorphism
\[
  M/\operatorname{Ker}(\phi)
  \xrightarrow{\ \sim\ }
  \varprojlim_{\mathfrak l}M/\overline{\mathfrak lM}.
\]
But we have just seen that
\[
  \operatorname{Ker}(\phi)
  =
  \bigcap_{\mathfrak l}\overline{\mathfrak lM}
  =
  (0).
\]
\hfill\(\square\)

\medskip
\noindent\textbf{Remark 0.3.1.2.}\footnote{Editors' note: This remark
has been added.}
\phantomsection
\label{sga3:VIIB:0.3.1.2}
The completed tensor product satisfies the usual associativity
condition: if \(L\), \(M\), and \(P\) are pseudocompact
\(A\)-modules, there is a canonical isomorphism
\[
  (L\widehat{\otimes}_{A}M)\widehat{\otimes}_{A}P
  \simeq
  L\widehat{\otimes}_{A}(M\widehat{\otimes}_{A}P).
\]
Indeed, both objects represent the functor which associates to every
object \(N\) of \(\operatorname{LF}(A)\) the abelian group of
continuous \(A\)-trilinear maps from \(L\times M\times P\) to \(N\).

\medskip
\noindent\textbf{0.3.2.}
\phantomsection
\label{sga3:VIIB:0.3.2}
Let
\[
  L'\xrightarrow{\,f\,}L\xrightarrow{\,g\,}L''\longrightarrow0
\]
be an exact sequence, and let \(M\) be an object of
\(\operatorname{PC}(A)\). It is clear that for every object \(N\) of
\(\operatorname{LF}(A)\), the induced sequences
\[
\begin{array}{ccccccc}
0 &\longrightarrow&
  \operatorname{Bil}_{c}(L''\times M,N)
&\longrightarrow&
  \operatorname{Bil}_{c}(L\times M,N)
&\longrightarrow&
  \operatorname{Bil}_{c}(L'\times M,N)
\\[0.35em]
&&\Big\Vert&&\Big\Vert&&\Big\Vert
\\[-0.1em]
0 &\longrightarrow&
  \operatorname{Hom}_{c}(L''\widehat{\otimes}_{A}M,N)
&\longrightarrow&
  \operatorname{Hom}_{c}(L\widehat{\otimes}_{A}M,N)
&\longrightarrow&
  \operatorname{Hom}_{c}(L'\widehat{\otimes}_{A}M,N)
\end{array}
\]
are exact. By \SGARefLink{sga3:VIIB:0.2.E}{0.2.E}, this is equivalent to saying that the sequence
\begin{equation}
  L'\widehat{\otimes}_{A}M
  \xrightarrow{\,f\widehat{\otimes}M\,}
  L\widehat{\otimes}_{A}M
  \xrightarrow{\,g\widehat{\otimes}M\,}
  L''\widehat{\otimes}_{A}M
  \longrightarrow0
  \tag{\(*\)}
  \label{sga3:VIIB:0.3.2.star}
\end{equation}
is exact. Hence:

\medskip
\noindent\textbf{Corollary.}
\phantomsection
\label{sga3:VIIB:0.3.2.cor}
\textit{For every pseudocompact \(A\)-module \(M\), the functor
\(-\widehat{\otimes}_{A}M\) is right exact.}

In particular, take \(L=A\), let \(f\) be the inclusion of a closed
ideal \(\mathfrak a\) into \(A\), and let \(g\) be the canonical
projection of \(A\) onto \(A/\mathfrak a\). We may then identify
\(A\widehat{\otimes}_{A}M\) with \(M\) by means of the map
\[
  x\widehat{\otimes}_{A}m\longmapsto xm.
\]
Since the image of
\(\mathfrak a\widehat{\otimes}_{A}M\) is closed in \(M\)
(cf.\ \SGARefLink{sga3:vii:proposition:0-2-b}{0.2.B}), and the image of
\(\mathfrak a\otimes_{A}M\) is dense in
\(\mathfrak a\widehat{\otimes}_{A}M\), the image of
\(f\widehat{\otimes}M\) is precisely the closure
\(\overline{\mathfrak aM}\) of \(\mathfrak aM\) in \(M\).
The exact sequence \(\textup{(\(*\))}\) therefore gives the
isomorphism
\[
  (A/\mathfrak a)\widehat{\otimes}_{A}M
  \xrightarrow{\ \sim\ }
  M/\overline{\mathfrak aM}.
\]

\medskip
\noindent\textbf{0.3.3. Nakayama's Lemma.}
\phantomsection
\label{sga3:VIIB:0.3.3}
\textit{Let \(A\) be a pseudocompact ring, \(M\) a pseudocompact
\(A\)-module, and \(\mathfrak a\) an ideal of \(A\) contained in the
radical \(r(A)\). Then}
\[
  \overline{\mathfrak aM}=M
  \qquad\Longrightarrow\qquad
  M=0.
\]

\noindent\textit{Proof.}
Suppose that
\(\overline{\mathfrak aM}=M\).\footnote{Editors' note: Replacing
\(\mathfrak a\) by its closure if necessary, we may assume that
\(\mathfrak a\) is closed.}
Let \(M'\) be an open submodule of \(M\), and put \(M''=M/M'\).
Since \(M''\) is discrete, \(\mathfrak aM''\) is closed in \(M''\),
and hence equals \(\overline{\mathfrak aM''}\). By \SGARefLink{sga3:VIIB:0.3.2}{0.3.2}, the
canonical map
\[
  M/\overline{\mathfrak aM}
  \longrightarrow
  M''/\overline{\mathfrak aM''}
\]
is surjective, so
\[
  \mathfrak aM''
  =
  \overline{\mathfrak aM''}
  =
  M''.
\]
Since \(M''\) is a finitely generated \(A\)-module and
\(\mathfrak a\subset r(A)\), the usual Nakayama lemma implies that
\(M''=0\). Consequently, every open submodule \(M'\) of \(M\) is
equal to \(M\), and therefore \(M=0\).\footnote{Editors' note: Indeed,
\(M\) is the projective limit of the quotients \(M/M'=0\).}
\hfill\(\square\)

\medskip
\noindent\textbf{0.3.4.}
\phantomsection
\label{sga3:VIIB:0.3.4}
The usual consequences of Nakayama's lemma follow.

\medskip
\noindent\textbf{Corollary.}
\phantomsection
\label{sga3:VIIB:0.3.4.cor}
\textit{Let \(\mathfrak a\) be a closed ideal contained in \(r(A)\),
and let \(f:M\to N\) be a morphism of pseudocompact \(A\)-modules.}
\begin{enumerate}
  \item[\textup{(i)}]
  \textit{The map \(f\) is surjective if the induced map}
  \[
    f':
    M/\overline{\mathfrak aM}
    \longrightarrow
    N/\overline{\mathfrak aN}
  \]
  \textit{is surjective.}\footnote{Editors' note: Consequently,
  every pseudocompact \(A\)-module is a quotient of a topologically
  free \(A\)-module. One first reduces to the case in which \(A\) is
  local; cf.\ the proof of \SGARefLink{sga3:VIIB:0.3.7}{0.3.7}.}

  \item[\textup{(ii)}]
  \textit{If \(N\) is projective, then \(f\) is an isomorphism whenever
  \(f'\) is an isomorphism.}
\end{enumerate}

\noindent\textit{Proof.}
Part \textup{\SGARefLink{sga3:vii:paragraph:5-1:equation:i}{(i)}} follows by applying \SGARefLink{sga3:VIIB:0.3.3}{Lemma~0.3.3} to
\(\operatorname{Coker}(f)\). For \textup{\SGARefLink{sga3:vii:paragraph:5-1:equation:ii}{(ii)}}, suppose that \(f'\) is
an isomorphism. By \textup{\SGARefLink{sga3:vii:paragraph:5-1:equation:i}{(i)}}, \(f\) is then surjective, and hence
has a section. Apply \SGARefLink{sga3:VIIB:0.3.3}{Lemma~0.3.3} to \(\operatorname{Ker}(f)\).
\hfill\(\square\)\par\smallskip

When \(A\) is local with maximal ideal \(\mathfrak m\), one may also
deduce from \SGARefLink{sga3:VIIB:0.3.3}{0.3.3} the following exchange theorem.

\medskip
\noindent\textbf{Theorem.}
\phantomsection
\label{sga3:VIIB:0.3.4.exchange}
\textit{Let \(A\) be a local pseudocompact ring, let \(\mathfrak m\)
be its maximal ideal, let \(M\) be a topologically free \(A\)-module
with pseudobasis \((m_i)_{i\in I}\) \textup{(\SGARefLink{sga3:VIIB:0.2.1}{0.2.1})}, and let \(N\)
be a closed direct summand of \(M\). There is a pseudobasis of \(M\)
consisting of elements of \(N\) together with some of the \(m_i\).}

\noindent\textit{Proof.}
This is clear when \(A\) is a field: use the duality of \SGARefLink{sga3:VIIB:0.2.2}{0.2.2} and apply
the usual exchange theorem. Consequently,
\[
  N/\overline{\mathfrak mN}
\]
has a complementary topologically free \(A/\mathfrak m\)-module with
pseudobasis \((\overline{m}_i)_{i\in J}\), where
\(\overline{m}_i\) is the image of \(m_i\) in
\[
  M/\overline{\mathfrak mM},
\]
and \(J\) is a subset of \(I\).\footnote{Editors' note: We have
corrected \(N/\mathfrak mN\) to
\(N/\overline{\mathfrak mN}\), and likewise for
\(M/\overline{\mathfrak mM}\) below.}
If
\[
  P=\prod_{i\in J}Am_i,
\]
then the canonical map
\[
  N\oplus P\longrightarrow M
\]
is ``bijective modulo \(\mathfrak m\).'' It is therefore bijective by
what precedes. For another proof, see \textup{[CA]}, \S~IV.2,
Proposition~8.
\hfill\(\square\)

% Next cursor: Exp7B local p. 11, combined-reader p. 524,
% re-edition printed p. 513, immediately before 0.3.5.

% END INCLUDED SGA3_English_current_buildable_source_20260823/source/inputs/VII/tex/components/19_expose_VIIB_sections03_034_completed_tensor_en.tex

% BEGIN INCLUDED SGA3_English_current_buildable_source_20260823/source/inputs/VII/tex/components/20_expose_VIIB_sections035_038_tensor_projectives_en.tex
% SGA 3, Expose VII_B, English reconstruction -- Loop 1.
% Scope: authority local pages 11--14, 0.3.5 through 0.3.8,
% stopping immediately before 0.4 on local page 14.
% Combined-reader pages: 524--527. Re-edition printed pages: 513--516.
% Original-edition marginal locators within this span: 485--488.
% Authority: Polo--Gille Exp7B-13oct24.pdf,
% SHA-256 D303690E0195654BBE7C3DAA5D89EA23E24CD709A3575217234AC19AE8836694.
% Jacob Reinhold's CC BY 4.0 Markdown was consulted as comparison lineage.

\medskip
\noindent\textbf{0.3.5.}
\phantomsection
\label{sga3:VIIB:0.3.5}
Consider three pseudocompact \(A\)-modules \(L,M,N\), where \(N\) has
finite length. Give the \(A\)-module
\(\operatorname{Hom}_c(M,N)\) the discrete topology. Every element
\[
  \psi\in
  \operatorname{Hom}_c
    \bigl(L,\operatorname{Hom}_c(M,N)\bigr)
\]
defines a continuous bilinear map
\[
  \psi':L\times M\longrightarrow N,
  \qquad
  (\ell,m)\longmapsto\psi(\ell)(m).
\]
This gives a natural isomorphism
\begin{equation}
  \operatorname{Hom}_c
    \bigl(L,\operatorname{Hom}_c(M,N)\bigr)
  \xrightarrow{\sim}
  \operatorname{Hom}_c
    \bigl(L\widehat{\otimes}_A M,N\bigr).
  \tag{1}
  \label{sga3:VIIB:0.3.5-eq1}
\end{equation}
This is another characterization of
\(\operatorname{Hom}_c(L\widehat{\otimes}_A M,N)\), which we now use
when \(M\) is the projective limit of a filtered projective system
\((M_i)\) of pseudocompact \(A\)-modules. By (1) and \SGARefLink{sga3:VIIB:0.2.F}{0.2.F}~\SGARefLink{sga3:vii:paragraph:5-1:equation:ii}{(ii)}, there
are natural isomorphisms
\begin{align}
  \operatorname{Hom}_c
    \left(L\widehat{\otimes}_A\varprojlim_iM_i,N\right)
  &\simeq
  \operatorname{Hom}_c
    \left(L,\operatorname{Hom}_c(\varprojlim_iM_i,N)\right)
  \notag\\
  &\simeq
  \operatorname{Hom}_c
    \left(L,\varinjlim_i\operatorname{Hom}_c(M_i,N)\right).
  \tag{2}
  \label{sga3:VIIB:0.3.5-eq2}
\end{align}

Moreover, because
\(\varinjlim_i\operatorname{Hom}_c(M_i,N)\) is discrete, every
continuous map from \(L\) into this module factors through a
finite-length quotient of \(L\). Hence the natural map
\begin{equation}
  \varinjlim_i
  \operatorname{Hom}_c
    \bigl(L,\operatorname{Hom}_c(M_i,N)\bigr)
  \xrightarrow{\sim}
  \operatorname{Hom}_c
    \left(L,\varinjlim_i\operatorname{Hom}_c(M_i,N)\right)
  \tag{3}
  \label{sga3:VIIB:0.3.5-eq3}
\end{equation}
is an isomorphism. Finally, (1) and \SGARefLink{sga3:VIIB:0.2.F}{0.2.F}~\SGARefLink{sga3:vii:paragraph:5-1:equation:ii}{(ii)} again give natural
isomorphisms
\begin{align}
  \varinjlim_i
  \operatorname{Hom}_c
    \bigl(L,\operatorname{Hom}_c(M_i,N)\bigr)
  &\simeq
  \varinjlim_i
  \operatorname{Hom}_c
    \bigl(L\widehat{\otimes}_A M_i,N\bigr)
  \notag\\
  &\simeq
  \operatorname{Hom}_c
    \left(\varprojlim_i
      (L\widehat{\otimes}_A M_i),N\right).
  \tag{4}
  \label{sga3:VIIB:0.3.5-eq4}
\end{align}
Combining \SGARefLink{sga3:VIIB:0.3.5-eq2}{(2)}, \SGARefLink{sga3:VIIB:0.3.5-eq3}{(3)}, and \SGARefLink{sga3:VIIB:0.3.5-eq4}{(4)} yields the following result.

\medskip
\noindent\textbf{Proposition.}
\phantomsection
\label{sga3:VIIB:0.3.5-proposition}
{\itshape Let \((M_i)\) be a filtered projective system of objects of
\(\operatorname{PC}(A)\), let
\(L\in\operatorname{PC}(A)\), and let
\(N\in\operatorname{LF}(A)\). There is an isomorphism, functorial in
\(N\),
\[
  \operatorname{Hom}_c
    \left(L\widehat{\otimes}_A\varprojlim_iM_i,N\right)
  \simeq
  \operatorname{Hom}_c
    \left(\varprojlim_i(L\widehat{\otimes}_A M_i),N\right),
\]
and consequently an isomorphism
\[
  L\widehat{\otimes}_A\varprojlim_iM_i
  \simeq
  \varprojlim_i(L\widehat{\otimes}_A M_i).
\]
Thus the completed tensor product commutes with filtered projective
limits.\footnote{This also recovers \SGARefLink{sga3:VIIB:0.3.1.1}{0.3.1.1}:
\[
  L=L\widehat{\otimes}_A A
  =
  L\widehat{\otimes}_A
    \varprojlim_{\mathfrak l}(A/\mathfrak l)
  \simeq
  \varprojlim_{\mathfrak l}L/\mathfrak lL,
\]
where \(\mathfrak l\) ranges over the open ideals of \(A\).}}

% Authority: Exp7B local p. 12; combined-reader p. 525; printed p. 514.
\medskip
\noindent\textbf{0.3.6.}
\phantomsection
\label{sga3:VIIB:0.3.6}
In particular,\footnote{Cf.\ editor's note~(12).} the completed tensor
product commutes with arbitrary products. For example, since \(A\) is
the product of its local components \(A_{\mathfrak m}\) (\SGARefLink{sga3:VIIB:0.1.1}{0.1.1}), every
pseudocompact \(A\)-module
\[
  M\simeq A\widehat{\otimes}_A M
\]
identifies with the product of the modules
\[
  M_{\mathfrak m}
  =
  A_{\mathfrak m}\widehat{\otimes}_A M,
\]
which are called the local components of \(M\).

Likewise, let \(M\) and \(N\) be pseudocompact \(A\)-modules. Recall
from \SGARefLink{sga3:VIIB:0.2.2}{0.2.2} that
\[
  M^\dagger=\operatorname{Hom}_c(M,A).
\]
Consider the map
\[
  \varphi:
  M^\dagger\otimes_A N^\dagger
  \longrightarrow
  (M\widehat{\otimes}_A N)^\dagger
\]
which sends \(f\otimes g\) to the map
\[
  m\widehat{\otimes}n\longmapsto f(m)g(n)
\]
from \(M\widehat{\otimes}_A N\) to \(A\). The map \(\varphi\) is
bijective when \(M\simeq A\).

\medskip
\noindent\textbf{Corollary.}
\phantomsection
\label{sga3:VIIB:0.3.6-corollary}
\textit{If \(A\) is Artinian, then \(\varphi\) is an isomorphism
whenever \(M\) is topologically free, or more generally projective.}

Indeed, with \(N\) fixed, each of the functors
\[
  M\longmapsto(M\widehat{\otimes}_A N)^\dagger
  \qquad\text{and}\qquad
  M\longmapsto M^\dagger\otimes_A N^\dagger
\]
takes arbitrary products to direct sums, by the preceding result and
\SGARefLink{sga3:VIIB:0.2.F}{0.2.F}.

\medskip
\noindent\textbf{Remark 0.3.6.A.}\footnote{The editors added this
remark, which will be useful in \SGARefLink{sga3:VIIB:2.3.1}{2.3.1}.}
\phantomsection
\label{sga3:VIIB:0.3.6.A}
\textit{The same use of \SGARefLink{sga3:VIIB:0.2.F}{0.2.F} gives the following result. Let \(A\) be
an Artinian ring, let \(M,Q\in\operatorname{PC}(A)\), and let
\(N\in\operatorname{LF}(A)\). If \(Q\) is projective, then there are
natural isomorphisms}
\[
  \operatorname{Hom}_c(M,Q)
  \xrightarrow{\sim}
  \operatorname{Hom}_A(Q^\dagger,M^\dagger),
  \qquad
  Q^\dagger\otimes_A N
  \xrightarrow{\sim}
  \operatorname{Hom}_c(Q,N).
\]

\medskip
\noindent\textbf{0.3.7.}
\phantomsection
\label{sga3:VIIB:0.3.7}
For every \(\mathfrak m\in\Upsilon(A)\), the functor
\[
  M\longmapsto A_{\mathfrak m}\widehat{\otimes}_A M
\]
is plainly exact. Since every projective pseudocompact \(A\)-module
\(P\) is a product of modules of the form \(A_{\mathfrak m}\), it
follows that
\[
  M\longmapsto P\widehat{\otimes}_A M
\]
is exact whenever \(P\) is projective. The converse is also true.

\medskip
\noindent\textbf{Proposition.}
\phantomsection
\label{sga3:VIIB:0.3.7-proposition}
{\itshape Let \(A\) be a pseudocompact ring and \(P\) a pseudocompact
\(A\)-module. The following conditions are equivalent:
\begin{enumerate}
  \item[\textup{(i)}]
  \(P\) is a projective object of \(\operatorname{PC}(A)\);
  \item[\textup{(ii)}]
  every local component \(P_{\mathfrak m}\) is a topologically free
  \(A_{\mathfrak m}\)-module;
  \item[\textup{(iii)}]
  the functor \(M\mapsto P\widehat{\otimes}_A M\) is exact.
\end{enumerate}}

The equivalence of \SGARefLink{sga3:vii:paragraph:5-1:equation:i}{(i)} and \SGARefLink{sga3:vii:paragraph:5-1:equation:ii}{(ii)} follows from \SGARefLink{sga3:VIIB:0.2.F}{0.2.F}~\SGARefLink{sga3:VIIB:0.2.F:item:iii}{(iii)} and \SGARefLink{sga3:VIIB:0.2.1}{0.2.1},
and we have just proved \SGARefLink{sga3:VIIB:0.2.1:item:ii}{(ii)}~\(\Rightarrow\)~\SGARefLink{sga3:VIIB:0.2.F:item:iii}{(iii)}. Suppose that the
functor \(M\mapsto P\widehat{\otimes}_A M\) is exact. Since
\(P\widehat{\otimes}_A M\) is the product of its local components,
\[
  (P\widehat{\otimes}_A M)_{\mathfrak m}
  \simeq
  P_{\mathfrak m}
    \widehat{\otimes}_{A_{\mathfrak m}}M_{\mathfrak m},
\]
we reduce to the case in which \(A\) is local. We shall prove that
\(P\) is then topologically free.

Let \(\mathfrak m\) be the maximal ideal of \(A\). The quotient
\(P/\overline{\mathfrak m}P\) is a linearly compact vector space over
\(A/\mathfrak m\), and hence a product of copies of \(A/\mathfrak m\)
(\SGARefLink{sga3:VIIB:0.2.2}{0.2.2}). Therefore there is a family \((A_i)_{i\in I}\) of copies of
\(A\) and an isomorphism
\[
  \phi:
  \prod_{i\in I}(A_i/\mathfrak m)
  \xrightarrow{\sim}
  P/\overline{\mathfrak m}P.
\]
Since \(\prod_{i\in I}A_i\) is projective, there is a commutative
square
\begin{center}
  \phantomsection
  \label{sga3:VIIB:0.3.7-projective-lift-square}
  % SGA3-VII-DIAG-082
  % BEGIN INCLUDED SGA3_English_current_buildable_source_20260823/source/inputs/VII/native_diagrams/exp7/viib_15_95/diag_082.tex
% SGA3-VII-DIAG-082
% Authority: Exp7B-13oct24.pdf, local p. 13.
% Crop: figures/Exp7B_localp013_projective_lift_square.png
% Locator: tex/components/20_expose_VIIB_sections035_038_tensor_projectives_en.tex, Loop 1 line 260.
\begin{tikzcd}[row sep=3em,column sep=4.8em]
\displaystyle\prod_{i\in I}A_i \arrow[r,"\psi"] \arrow[d,"p"']
  & P \arrow[d,"q"] \\
\displaystyle\prod_{i\in I}(A_i/\mathfrak m)
  \arrow[r,"\varphi","{\sim}"']
  & P/\mathfrak mP
\end{tikzcd}

% END INCLUDED SGA3_English_current_buildable_source_20260823/source/inputs/VII/native_diagrams/exp7/viib_15_95/diag_082.tex

\end{center}
where \(p\) and \(q\) are the canonical projections. Apply Nakayama's
lemma to \(\operatorname{Coker}\psi\). Since
\[
  (A/\mathfrak m)\widehat{\otimes}_A\psi
\]
is precisely \(\phi\), the map \(\psi\) is
surjective.\footnote{This proves the result announced in editor's
note~(22): every pseudocompact \(A\)-module \(M\) is a quotient of a
topologically free \(A\)-module. Since products are exact in
\(\operatorname{PC}(A)\), one reduces to the local case treated here.}

% Authority: Exp7B local p. 13; combined-reader p. 526; printed p. 515.
Put
\[
  B=\prod_{i\in I}A_i,
  \qquad
  N=\ker\psi.
\]
There is then the following commutative exact diagram:
\begin{center}
  \phantomsection
  \label{sga3:VIIB:0.3.7-nakayama-exact-diagram}
  % SGA3-VII-DIAG-083
  % BEGIN INCLUDED SGA3_English_current_buildable_source_20260823/source/inputs/VII/native_diagrams/exp7/viib_15_95/diag_083.tex
% SGA3-VII-DIAG-083
% Authority: Exp7B-13oct24.pdf, local p. 13.
% Crop: figures/Exp7B_localp013_nakayama_exact_diagram.png
% Locator: tex/components/20_expose_VIIB_sections035_038_tensor_projectives_en.tex, Loop 1 line 285.
\begin{tikzcd}[row sep=2.05em,column sep=1.65em]
{} & {} & {} &
  0
    \arrow[d]
    \arrow[ddlll,dashed,rounded corners=10pt,
      to path={
        -- ([xshift=6em]\tikztostart.east)
        |- ([yshift=.8em]\tikztotarget.north)
        -- (\tikztotarget)
      }]
  & {} \\
{} & \mathfrak m\widehat\otimes_A N \arrow[r] \arrow[d]
  & \mathfrak m\widehat\otimes_A B \arrow[r] \arrow[d]
  & \mathfrak m\widehat\otimes_A P \arrow[r] \arrow[d]
  & 0 \\
0
  \arrow[r]
  \arrow[dr,dashed,out=-90,in=180,looseness=1.35]
  & A\widehat\otimes_A N \arrow[r] \arrow[d]
  & A\widehat\otimes_A B \arrow[r,"\psi"] \arrow[d]
  & A\widehat\otimes_A P \arrow[r] \arrow[d]
  & 0 \\
{} & (A/\mathfrak m)\widehat\otimes_A N \arrow[r]
  & (A/\mathfrak m)\widehat\otimes_A B \arrow[r,"\varphi"]
  & (A/\mathfrak m)\widehat\otimes_A P \arrow[r]
  & 0 \qquad .
\end{tikzcd}

% END INCLUDED SGA3_English_current_buildable_source_20260823/source/inputs/VII/native_diagrams/exp7/viib_15_95/diag_083.tex

\end{center}
The snake lemma, applied to the first two rows, shows that in the
bottom row the morphism
\[
  (A/\mathfrak m)\widehat{\otimes}_A N
  \longrightarrow
  (A/\mathfrak m)\widehat{\otimes}_A B
\]
is a monomorphism. But \(\phi\) is an isomorphism, so
\((A/\mathfrak m)\widehat{\otimes}_A N=0\). Hence \(N=0\) by \SGARefLink{sga3:VIIB:0.3.3}{0.3.3},
and \(\psi\) is an isomorphism.\footnote{A completely analogous proof,
using the ``nilpotent Nakayama lemma,'' shows that if \(A\) is Artinian
and \(P\) is a flat \(A\)-module, then every local component of \(P\)
is a free \(A\)-module (cf.\ \textup{[BAC]}, II, \S~\SGARefLink{sga3:VIIB:3.2:prop}{3.2}, \SGARefLink{sga3:VIIB:0.3.5-eq2}{Corollary~2}
to Proposition~5).}

\medskip
\noindent\textbf{0.3.8. Corollary.}
\phantomsection
\label{sga3:VIIB:0.3.8}
\textit{Let \(A\) be a complete Noetherian local ring and \(P\) a
pseudocompact \(A\)-module. Then \(P\) is topologically free if and
only if \(P\) is flat over \(A\).}

Indeed, the canonical map
\[
  M\otimes_A P\longrightarrow M\widehat{\otimes}_A P
\]
is bijective when \(M=A\), and hence also when \(M\) is Noetherian:
take a finite presentation of \(M\) and use the right exactness of
both the tensor product and the completed tensor product.

Now \(P\) is flat if and only if
\[
  M\longmapsto M\otimes_A P
\]
is exact on Noetherian modules. Likewise, the proof of \SGARefLink{sga3:VIIB:0.3.7}{Proposition
0.3.7} shows that \(P\) is topologically free provided the sequence
\[
  0\longrightarrow
  \mathfrak m\widehat{\otimes}_A P
  \longrightarrow
  A\widehat{\otimes}_A P
  \longrightarrow
  (A/\mathfrak m)\widehat{\otimes}_A P
  \longrightarrow0
\]
is exact. Thus \(P\) is topologically free if and only if the functor
\[
  M\longmapsto M\widehat{\otimes}_A P
\]
is exact on Noetherian modules. The corollary follows from the equality
\[
  M\otimes_A P=M\widehat{\otimes}_A P
\]
established above.

% Next cursor: Exp7B local p. 14, combined-reader p. 527,
% re-edition printed p. 516, at 0.4.

% END INCLUDED SGA3_English_current_buildable_source_20260823/source/inputs/VII/tex/components/20_expose_VIIB_sections035_038_tensor_projectives_en.tex

% BEGIN INCLUDED SGA3_English_current_buildable_source_20260823/source/inputs/VII/tex/components/21_expose_VIIB_section04_profinite_algebras_en.tex
% SGA 3, Expose VII_B, English reconstruction -- Loop 1.
% Scope: authority local pages 14--15, all of 0.4 through 0.4.2,
% stopping immediately before 0.5 on local page 16.
% Combined-reader pages: 527--528. Re-edition printed pages: 516--517.
% Original-edition marginal locators within this span: 488--490.
% Authority: Polo--Gille Exp7B-13oct24.pdf,
% SHA-256 D303690E0195654BBE7C3DAA5D89EA23E24CD709A3575217234AC19AE8836694.
% Jacob Reinhold's CC BY 4.0 Markdown was consulted as comparison lineage.

\section*{0.4. Profinite Algebras}
\phantomsection
\label{sga3:VIIB:0.4}
\addcontentsline{toc}{subsection}{\SGARefLink{sga3:VIIB:0.4}{0.4}. Profinite Algebras}

Let \(k\) be a pseudocompact ring. A \emph{topological
\(k\)-algebra} is a commutative topological ring \(A\), together with
a morphism of topological rings \(k\to A\). The algebra \(A\) is called
a \emph{profinite \(k\)-algebra} if its underlying topological
\(k\)-module is pseudocompact.

Let \(\mathfrak l\) be an open \(k\)-submodule of such an \(A\). The
composite map
\[
  \varphi:
  A\times A
  \xrightarrow{\mathrm{mult}}A
  \longrightarrow A/\mathfrak l
\]
is continuous. \SGARefLink{sga3:VIIB:0.3.1}{Lemma 0.3.1} therefore gives an open \(k\)-submodule
\(\mathfrak n\) of \(A\) such that
\(\varphi(A\times\mathfrak n)=0\). Thus \(\mathfrak l\) contains the
open ideal \(A\mathfrak n\), and \(A\) is a pseudocompact ring.

Likewise, let \(M\) be a topological \(A\)-module whose underlying
\(k\)-module is pseudocompact. If \(M'\) is an open \(k\)-submodule of
\(M\), applying \SGARefLink{sga3:VIIB:0.3.1}{Lemma 0.3.1} to
\[
  A\times M
  \xrightarrow{\mathrm{mult}}M
  \longrightarrow M/M'
\]
shows that \(M'\) contains an open \(A\)-submodule. Hence \(M\) is also
a pseudocompact \(A\)-module.\footnote{The editors added the following
lemma; cf.\ editor's note~(36) in \SGARefLink{sga3:VIIB:0.5}{Proposition 0.5}.}
Conversely:

\medskip
\noindent\textbf{Lemma.}
\phantomsection
\label{sga3:VIIB:0.4-restriction-lemma}
\textit{Let \(A\) be a profinite \(k\)-algebra and \(M\) a
pseudocompact \(A\)-module. The \(k\)-module \(M|_k\), obtained by
restriction of scalars, is pseudocompact.}

Indeed, every finite-length pseudocompact \(A\)-module is isomorphic
to a quotient
\[
  A^n/L,
\]
where \(L\) is an open submodule of \(A^n\), and is therefore a
pseudocompact \(k\)-module. Since \(M|_k\) is a projective limit of
modules of this kind, it is a pseudocompact \(k\)-module.

\medskip
\noindent\textbf{0.4.1.}
\phantomsection
\label{sga3:VIIB:0.4.1}
If \(A\) and \(B\) are profinite \(k\)-algebras, a \emph{morphism}
from \(A\) to \(B\) is, by definition, a continuous homomorphism of
\(k\)-algebras. We denote the category of profinite \(k\)-algebras by
\[
  \operatorname{Alp}/k.
\]
It plainly has projective limits: the algebra underlying a projective
limit is the projective limit of the underlying algebras, endowed with
the projective-limit topology.

It also has finite direct limits.\footnote{In \SGARefLink{sga3:VIIB:0.4.2}{0.4.2} we shall see that
it has arbitrary direct limits.} For example, if
\[
  f:A\longrightarrow B,
  \qquad
  g:A\longrightarrow C
\]
are morphisms of profinite \(k\)-algebras, the amalgamated sum of
\(B\) and \(C\) over \(A\) has
\[
  B\widehat{\otimes}_A C
\]
as its underlying topological \(A\)-module. Here \SGARefLink{sga3:VIIB:0.4}{0.4} makes \(B\) and
\(C\) pseudocompact \(A\)-modules via \(f\) and \(g\), and the
multiplication is characterized by
\[
  (b\widehat{\otimes}c)
  (b'\widehat{\otimes}c')
  =
  bb'\widehat{\otimes}cc'
\]
for \(b,b'\in B\) and \(c,c'\in C\).

% Authority: Exp7B local p. 15; combined-reader p. 528; printed p. 517.
\medskip
\noindent\textbf{Definition 0.4.2.}
\phantomsection
\label{sga3:VIIB:0.4.2}
\textit{A profinite \(k\)-algebra \(C\) is said to have
\emph{finite length} if its underlying \(k\)-module has finite length,
and is therefore discrete. We denote by \(\operatorname{Alf}/k\) the
full subcategory of \(\operatorname{Alp}/k\) formed by the
\(k\)-algebras of finite length.}\footnote{Caution: \(k\) is an object
of \(\operatorname{Alf}/k\) only when \(k\) is Artinian; cf.\ \SGARefLink{sga3:VIIB:1.2.2}{1.2.2}
below.}

For every profinite \(k\)-algebra \(A\), let \(h_A\) be the functor
\[
  h_A:\operatorname{Alf}/k\longrightarrow(\mathbf{Sets}),
  \qquad
  C\longmapsto
  \operatorname{Hom}_{\operatorname{Alp}/k}(A,C).
\]
The functor \(h_A\) is left exact,\footnote{That is, it commutes with
finite projective limits.} and the canonical projections
\(A\to A/\mathfrak l\), where \(\mathfrak l\) ranges over the open
ideals of \(A\), induce for every \(C\in\operatorname{Alf}/k\) an
isomorphism
\[
  \varinjlim_{\mathfrak l}
  \operatorname{Hom}_{\operatorname{Alf}/k}(A/\mathfrak l,C)
  \xrightarrow{\sim}
  \operatorname{Hom}_{\operatorname{Alp}/k}(A,C),
\]
functorial in \(C\). In other words, \(h_A\) is the direct limit of the
representable functors \(h_{A/\mathfrak l}\):
\begin{equation}
  h_A\simeq\varinjlim_{\mathfrak l}h_{A/\mathfrak l}.
  \tag{\(\ast\)}
  \label{sga3:VIIB:0.4.2-star}
\end{equation}

If \(B\) is another profinite \(k\)-algebra, the general properties of
the bifunctor \(\operatorname{Hom}\), together with
\[
  \operatorname{Hom}(h_C,h_B)=h_B(C)
\]
for \(C\) of finite length, give isomorphisms
\begin{align*}
  \operatorname{Hom}_{\operatorname{Alp}/k}(B,A)
  &\simeq
  \varprojlim_{\mathfrak l}
  \operatorname{Hom}_{\operatorname{Alp}/k}(B,A/\mathfrak l)
  \\
  &\simeq
  \varprojlim_{\mathfrak l}
  \operatorname{Hom}(h_{A/\mathfrak l},h_B)
  \\
  &\simeq
  \operatorname{Hom}
    \left(\varinjlim_{\mathfrak l}h_{A/\mathfrak l},h_B\right).
\end{align*}
Together with \eqref{sga3:VIIB:0.4.2-star}, this proves that the
contravariant functor \(A\mapsto h_A\) is fully faithful. In fact:

\medskip
\noindent\textbf{Proposition.}
\phantomsection
\label{sga3:VIIB:0.4.2-proposition}
\textit{The functor \(A\mapsto h_A\) is an anti-equivalence from
\(\operatorname{Alp}/k\) onto the category of left-exact functors from
\(\operatorname{Alf}/k\) to \((\mathbf{Sets})\).}

By what precedes, it remains only to show that every left-exact functor
\[
  F:\operatorname{Alf}/k\longrightarrow(\mathbf{Sets})
\]
is isomorphic to a functor \(h_A\). The algebra \(A\) may be constructed
as follows (cf.\ TDTE~II, \S~3).

Because \(F\) is left exact, for every finite-length \(k\)-algebra \(C\)
and every \(\xi\in F(C)\), there is a smallest subalgebra \(C'\) of
\(C\) such that \(\xi\) belongs to the image of \(F(C')\) in \(F(C)\).
If \(C'=C\), the pair \((C,\xi)\) is called \emph{minimal}.

The minimal pairs form a category when a morphism from \((C,\xi)\) to
\((D,\eta)\) is defined to be a homomorphism
\(\phi:C\to D\) such that
\[
  (F\phi)(\xi)=\eta.
\]
Such a \(\phi\) is necessarily surjective, and the category of minimal
pairs is left filtered. Moreover, one may restrict to pairs
\((C,\xi)\) for which \(C\) belongs to a set containing one
finite-length \(k\)-algebra of every isomorphism
type.\footnote{Indeed, every discrete \(k\)-module of length \(n\) is
isomorphic to \(k^n/L\), where \(L\) is an open submodule of \(k^n\) of
colength \(n\). One may then consider the set of isomorphism classes of
profinite \(k\)-algebra structures on each such module.}
The functor \((C,\xi)\mapsto C\), from the category of minimal pairs to
the category of profinite \(k\)-algebras, therefore has a projective
limit; take \(A\) to be this limit.

\medskip
\noindent\textbf{Corollary.}
\phantomsection
\label{sga3:VIIB:0.4.2-corollary}
\textit{The category \(\operatorname{Alp}/k\) has arbitrary direct
limits.}

Indeed, the category of left-exact functors from
\(\operatorname{Alf}/k\) to \((\mathbf{Sets})\) has projective limits,
computed pointwise: if \((F_i)\) is a projective system of such
functors, then for every \(C\in\operatorname{Alf}/k\),
\[
  (\varprojlim_iF_i)(C)=\varprojlim_iF_i(C).
\]
\footnote{If \((A_i)_{i\in I}\) is a direct system in
\(\operatorname{Alp}/k\), its direct limit in \(\operatorname{Alp}/k\)
is the separated completion of the \(k\)-algebra \(B\), the direct
limit of the underlying \(k\)-algebras, for the topology defined by the
ideals \(I\) such that \(I\cap A_i\) is open in \(A_i\) for every \(i\)
and \(\operatorname{length}_k(B/I)<\infty\). Notice, for example, that
if \(K/k\) is an algebraic extension of fields of infinite degree and
\((k_i)\) is the filtered direct system of finite-degree subextensions,
then the direct limit of the system \((k_i)\) in
\(\operatorname{Alp}/k\) is the zero ring.}

% Next cursor: Exp7B local p. 16, combined-reader p. 529,
% re-edition printed p. 518, at 0.5.

% END INCLUDED SGA3_English_current_buildable_source_20260823/source/inputs/VII/tex/components/21_expose_VIIB_section04_profinite_algebras_en.tex

% BEGIN INCLUDED SGA3_English_current_buildable_source_20260823/source/inputs/VII/tex/components/22_expose_VIIB_section05_base_change_en.tex
% SGA 3, Expose VII_B, English reconstruction -- Loop 1.
% Scope: authority local page 16, all of 0.5,
% stopping immediately before Section 1 on local page 17.
% Combined-reader page: 529. Re-edition printed page: 518.
% Original-edition marginal locator: 491.
% Authority: Polo--Gille Exp7B-13oct24.pdf,
% SHA-256 D303690E0195654BBE7C3DAA5D89EA23E24CD709A3575217234AC19AE8836694.
% Jacob Reinhold's CC BY 4.0 Markdown was consulted as comparison lineage.

\section*{0.5. Base Change}
\phantomsection
\label{sga3:VIIB:0.5}
\addcontentsline{toc}{subsection}{\SGARefLink{sga3:VIIB:0.5}{0.5}. Base Change}

Let
\[
  \phi:k\longrightarrow\ell
\]
be a homomorphism of pseudocompact rings (\SGARefLink{sga3:VIIB:0.1.3}{0.1.3}). The construction of
\SGARefLink{sga3:VIIB:0.3}{0.3} admits the following generalization.\footnote{The editors expanded
the details of this paragraph.}

\medskip
\noindent\textbf{Definition 0.5.A.}
\phantomsection
\label{sga3:VIIB:0.5.A}
{\itshape For every \(M\in\operatorname{PC}(k)\) and
\(N\in\operatorname{PC}(\ell)\), let
\[
  M\widehat{\otimes}_kN
\]
be the separated completion of \(M\otimes_kN\) for the linear topology
defined by the kernels of the projections
\[
  M\otimes_kN
  \longrightarrow
  (M/M')\otimes_k(N/N'),
\]
where \(M'\), respectively \(N'\), ranges over the open submodules of
\(M\), respectively \(N\). Then \(M\widehat{\otimes}_kN\) is a
pseudocompact \(\ell\)-module. If \(m\in M\) and \(x\in N\), the
canonical image of \(m\otimes_kx\) in
\(M\widehat{\otimes}_kN\) is denoted by
\[
  m\widehat{\otimes}_kx.
\]}

This applies in particular to \(N=\ell\). In that case
\[
  M\widehat{\otimes}_k\ell
\]
is called the pseudocompact \(\ell\)-module obtained from \(M\) by the
base change \(k\to\ell\).

\medskip
\noindent\textbf{Remarks 0.5.B.}
\phantomsection
\label{sga3:VIIB:0.5.B}

\noindent\textit{\SGARefLink{sga3:vii:paragraph:5-1:equation:i}{(i)}}
In such a base change, \(\ell\) need not be a profinite
\(k\)-algebra. The typical example is that in which \(k\) is a field
and \(\ell=K\) is an arbitrary field extension of \(k\).

\smallskip
\noindent\textit{\SGARefLink{sga3:vii:paragraph:5-1:equation:ii}{(ii)}}
If, however, the \(k\)-module underlying \(N\) is pseudocompact---for
example, if \(\ell\) is a profinite \(k\)-algebra---then \SGARefLink{sga3:VIIB:0.4}{0.4} shows
that every open \(k\)-submodule of \(N\) contains an open
\(\ell\)-submodule. In that case \(M\widehat{\otimes}_kN\) agrees with
the completed tensor product of \SGARefLink{sga3:VIIB:0.3}{0.3} for the pseudocompact \(k\)-modules
\(M\) and \(N\), so the notation is unambiguous.

In every case, the \(k\)-module \(N|_k\) obtained by restriction of
scalars is a topological \(k\)-module: the map
\[
  k\times N\longrightarrow N,
  \qquad
  (t,n)\longmapsto\phi(t)n
\]
is continuous. We write
\(\operatorname{Hom}_c(M,N|_k)\) for the abelian group of continuous
\(k\)-linear maps from \(M\) to \(N|_k\).

\medskip
\noindent\textbf{Proposition 0.5.C.}
\phantomsection
\label{sga3:VIIB:0.5.C}
\textit{For every \(M\in\operatorname{PC}(k)\) and
\(N\in\operatorname{PC}(\ell)\), there is a functorial isomorphism}
\[
  \operatorname{Hom}_{\operatorname{PC}(\ell)}
    (M\widehat{\otimes}_k\ell,N)
  \simeq
  \operatorname{Hom}_c(M,N|_k).
\]
\footnote{Consequently, when \(\ell\) is a profinite \(k\)-algebra,
the functor \(M\mapsto M\widehat{\otimes}_k\ell\) is left adjoint to
the restriction-of-scalars functor.}

Indeed, let \(u:M\to N|_k\) be a continuous homomorphism. The map
\[
  u':M\times\ell\longrightarrow N,
  \qquad
  (m,\lambda)\longmapsto\lambda u(m)
\]
is continuous and ``bilinear,'' meaning \(k\)-linear in the first
factor and \(\ell\)-linear in the second. If \(N'\) is an open
\(\ell\)-submodule of \(N\), the argument of \SGARefLink{sga3:VIIB:0.3.1}{Lemma 0.3.1} gives an open
submodule \(M'\subset M\) and an open ideal
\(\ell'\subset\ell\) such that
\[
  u'(M'\times\ell)\subset N',
  \qquad
  u'(M\times\ell')\subset N'.
\]
Thus \(u\) induces a continuous \(\ell\)-linear map
\[
  U:M\widehat{\otimes}_k\ell\longrightarrow N
\]
satisfying
\[
  U(m\widehat{\otimes}_k\lambda)=\lambda u(m).
\]
Conversely, every morphism
\[
  f:M\widehat{\otimes}_k\ell\longrightarrow N
\]
gives the morphism
\[
  f':M\longrightarrow N|_k,
  \qquad
  m\longmapsto f(m\widehat{\otimes}_k1).
\]

As in \SGARefLink{sga3:VIIB:0.3.2}{0.3.2}, \SGARefLink{sga3:VIIB:0.3.5}{0.3.5}, and \SGARefLink{sga3:VIIB:0.3.1.2}{0.3.1.2}, this yields:

\medskip
\noindent\textbf{Corollary 0.5.D.}
\phantomsection
\label{sga3:VIIB:0.5.D}
{\itshape The functor
\[
  \operatorname{PC}(k)\longrightarrow\operatorname{PC}(\ell),
  \qquad
  M\longmapsto M\widehat{\otimes}_k\ell
\]
is right exact and commutes with filtered projective limits. Thus, if
\((M_i)\) is a filtered projective system in
\(\operatorname{PC}(k)\), there is a canonical isomorphism
\[
  (\varprojlim_iM_i)\widehat{\otimes}_k\ell
  \simeq
  \varprojlim_i(M_i\widehat{\otimes}_k\ell).
\]
Moreover, if \(M\) and \(N\) are pseudocompact \(k\)-modules, there is
a canonical isomorphism
\[
  (M\widehat{\otimes}_kN)\widehat{\otimes}_k\ell
  \simeq
  (M\widehat{\otimes}_k\ell)
    \widehat{\otimes}_{\ell}
  (N\widehat{\otimes}_k\ell).
\]}

\medskip
\noindent\textbf{Definition 0.5.E.}
\phantomsection
\label{sga3:VIIB:0.5.E}
{\itshape Finally, if \(A\) is a profinite \(k\)-algebra, there is a
unique profinite \(\ell\)-algebra structure on
\(A\widehat{\otimes}_k\ell\) such that
\[
  (a\widehat{\otimes}_k\lambda)
  (b\widehat{\otimes}_k\mu)
  =
  ab\widehat{\otimes}_k(\lambda\mu)
\]
for \(a,b\in A\) and \(\lambda,\mu\in\ell\). The profinite
\(\ell\)-algebra \(A\widehat{\otimes}_k\ell\) is said to be obtained
from \(A\) by extension of scalars, or base change, along
\(k\to\ell\).}

% Next cursor: Exp7B local p. 17, combined-reader p. 530,
% re-edition printed p. 519, at Section 1.

% END INCLUDED SGA3_English_current_buildable_source_20260823/source/inputs/VII/tex/components/22_expose_VIIB_section05_base_change_en.tex

% BEGIN INCLUDED SGA3_English_current_buildable_source_20260823/source/inputs/VII/tex/components/23_expose_VIIB_sections11_122_bicoalgebras_en.tex
% SGA 3, Expose VII_B, English reconstruction -- Loop 1.
% Scope: authority local pages 17--20, Section 1 through the end of 1.2.2,
% stopping immediately before 1.2.3 on local page 20.
% Combined-reader pages: 530--533. Re-edition printed pages: 519--522.
% Original-edition marginal locators in this span: 492--495.
% Editors' notes in scope: 37--44.
% Authority: Polo--Gille Exp7B-13oct24.pdf,
% SHA-256 D303690E0195654BBE7C3DAA5D89EA23E24CD709A3575217234AC19AE8836694.
% Jacob Reinhold's CC BY 4.0 Markdown was consulted as comparison lineage.
% The genuine double-arrow diagram in 1.2.C is a source-bound PDF crop;
% all other displayed mathematics is native LaTeX.

% Authority: Exp7B local p. 17; combined-reader p. 530; printed p. 519.
\section*{1. Formal Varieties over a Pseudocompact Ring}
\phantomsection
\label{sga3:VIIB:1}
\addcontentsline{toc}{section}{1. Formal Varieties over a Pseudocompact Ring}

\medskip
\noindent\textbf{1.1.}
\phantomsection
\label{sga3:VIIB:1.1}
To every pseudocompact ring \(A\) one may associate a formal scheme
(EGA~I, 10.4.2) as follows. Its underlying topological space is the
set \(\Upsilon(A)\) of open prime ideals of \(A\), which are therefore
maximal, endowed with the discrete topology. For every subset
\(E\subset\Upsilon(A)\), the space of sections of the structure sheaf
over \(E\) is the Cartesian product
\[
  \prod_{\mathfrak m\in E}A_{\mathfrak m}.
\]
% Authority: Exp7B local p. 18; combined-reader p. 531; printed p. 520.
The formal scheme so obtained is denoted by \(\operatorname{Spf}(A)\),
the \emph{formal spectrum} of \(A\).\footnote{\label{sga3:VIIB:nde:37}%
Cf.\ \SGARefLink{sga3:VIIB:0.1.2}{Remarks 0.1.2}.}

If \(A\) and \(B\) are two pseudocompact rings, a \emph{morphism} from
\(\operatorname{Spf}(B)\) to \(\operatorname{Spf}(A)\) consists of a
map
\[
  f:\Upsilon(B)\longrightarrow\Upsilon(A)
\]
and a family of continuous homomorphisms
\[
  f_y^\natural:A_{f(y)}\longrightarrow B_y,
  \qquad y\in\Upsilon(B).
\]
Such a morphism induces a continuous homomorphism
\[
  f^\natural:
  A=\prod_{x\in\Upsilon(A)}A_x
  \longrightarrow
  B=\prod_{y\in\Upsilon(B)}B_y.
\]
The converse is true:

\medskip
\noindent\textbf{Proposition.}
\phantomsection
\label{sga3:VIIB:1.1-spf-fully-faithful}
\textit{The contravariant functor
\(A\mapsto\operatorname{Spf}(A)\) is fully faithful.}

Indeed, let \(\phi:A\to B\) be a continuous homomorphism. The inverse
image \(\phi^{-1}(\mathfrak n)\) of an open maximal ideal of \(B\) is
an open prime ideal of \(A\), and hence is maximal in \(A\). We
therefore obtain a map
\[
  \mathfrak n\longmapsto\phi^{-1}(\mathfrak n)
\]
from \(\Upsilon(B)\) to \(\Upsilon(A)\), and \(\phi\) induces a
continuous homomorphism
\[
  A_{\phi^{-1}(\mathfrak n)}\longrightarrow B_{\mathfrak n}.
\]
Thus \(\phi\) induces a morphism of formal schemes
\[
  \operatorname{Spf}(\phi):
  \operatorname{Spf}(B)\longrightarrow\operatorname{Spf}(A).
\]
It is readily verified that
\[
  \operatorname{Spf}(\phi)^\natural=\phi,
  \qquad
  \operatorname{Spf}(f^\natural)=f
\]
for every morphism
\(f:\operatorname{Spf}(B)\to\operatorname{Spf}(A)\), which proves the
proposition.

Although we shall speak here only of formal schemes of the form
\(\operatorname{Spf}(A)\), we shall use the language of formal schemes
rather than that of pseudocompact rings, so as to ground our assertions
in geometric intuition.

\medskip
\noindent\textbf{1.2.}
\phantomsection
\label{sga3:VIIB:1.2}
Let \(k\) be a pseudocompact ring.

\medskip
\noindent\textbf{Definition 1.2.A.}\footnote{%
\label{sga3:VIIB:nde:38}The editors added the numbering \SGARefLink{sga3:VIIB:1.2.A}{1.2.A}--\SGARefLink{sga3:VIIB:1.2.E}{1.2.E}
for later reference.}
\phantomsection
\label{sga3:VIIB:1.2.A}
A \emph{formal variety over \(k\)} is a formal scheme \(X\)
over \(\operatorname{Spf}(k)\) which is isomorphic to a formal
\(k\)-scheme \(\operatorname{Spf}(A)\) for some profinite
\(k\)-algebra \(A\). The algebra \(A\) is then isomorphic to the
\emph{affine algebra} of \(X\), that is, the algebra of sections of
the structure sheaf \(\mathcal O_X\) of \(X\).

We denote by \(\operatorname{Vaf}/k\) the full subcategory of the
category of formal schemes over \(\operatorname{Spf}(k)\) whose
objects are the formal \(k\)-varieties.\footnote{%
\label{sga3:VIIB:nde:39}By the definition of morphisms in \SGARefLink{sga3:VIIB:1.1}{1.1}, every
\(X\in\operatorname{Ob}(\operatorname{Vaf}/k)\) is the direct sum in
\(\operatorname{Vaf}/k\) of the one-point formal varieties
\(\operatorname{Spf}(\mathcal O_{X,x})\), for \(x\in X\).}

By \SGARefLink{sga3:VIIB:1.1}{1.1}, the functor \(A\mapsto\operatorname{Spf}(A)\) is an
anti-equivalence from \(\operatorname{Alp}/k\) (\SGARefLink{sga3:VIIB:0.4.1}{0.4.1}) onto
\(\operatorname{Vaf}/k\). Hence, by the corollary of \SGARefLink{sga3:VIIB:0.4.2}{0.4.2}:

\medskip
\noindent\textbf{Proposition 1.2.B.}
\phantomsection
\label{sga3:VIIB:1.2.B}
\textit{The category \(\operatorname{Vaf}/k\) has projective and
direct limits.}\footnote{\label{sga3:VIIB:nde:40}In particular,
cokernels exist in \(\operatorname{Vaf}/k\); see below. Moreover, a
filtered projective limit in \(\operatorname{Vaf}/k\), all of whose
transition morphisms are surjective, can be empty; cf.\ editor's
note~(34).}

For example, let \(X\in\operatorname{Ob}(\operatorname{Vaf}/k)\), let
\[
  f:Y\longrightarrow X,
  \qquad
  g:Z\longrightarrow X
\]
be two formal \(k\)-varieties over \(X\), and let \(A,B,C\) be the
affine algebras of \(X,Y,Z\), respectively.
% Authority: Exp7B local p. 19; combined-reader p. 532; printed p. 521.
By \SGARefLink{sga3:VIIB:0.4.1}{0.4.1}, the affine algebra of the fiber product \(Y\times_X Z\)
identifies with \(B\widehat{\otimes}_A C\),\footnote{%
\label{sga3:VIIB:nde:41}The formal variety \(Y\times_X Z\)
is the direct sum, over \(x\in X\), of the formal varieties
\(B_x\widehat{\otimes}_{\mathcal O_{X,x}}C_x\),
where \(B_x\) is the product of the \(\mathcal O_{Y,y}\) for
\(y\in Y\) such that \(f(y)=x\), and \(C_x\) is defined analogously.
This will be used in \SGARefLink{sga3:VIIB:2.5.1:corollary}{2.5.1}.}
Consequently, the inclusion of \(\operatorname{Vaf}/k\) into the
category of all formal \(k\)-schemes commutes with finite projective
limits (cf.\ EGA~I, 10.7).

Direct limits of formal \(k\)-varieties correspond to projective
limits of their affine algebras.

\medskip
\noindent\textbf{Example 1.2.C (Cokernels).}
\phantomsection
\label{sga3:VIIB:1.2.C}
Let \(d,e:X\rightrightarrows Y\) be a double arrow in
\(\operatorname{Vaf}/k\). The affine algebra of
\(\operatorname{Coker}(d,e)\) is isomorphic to the kernel of the
homomorphisms induced on the affine algebras of \(X\) and \(Y\).
There is also the following construction of
\(\operatorname{Coker}(d,e)\). Its underlying topological space is the
cokernel of the underlying spaces.\footnote{%
\label{sga3:VIIB:nde:42}That is, it is the quotient set of \(Y\) by
the identifications \(d(x)=e(x)\), for every \(x\in X\), endowed with
the quotient topology, which here is discrete.} Let \(p\) be the
canonical projection from the set underlying \(Y\) to this cokernel,
and let \(z\) be a point of the cokernel. The local algebra of
\(\operatorname{Coker}(d,e)\) at \(z\) is the kernel of the double
arrow
\begin{center}
  \phantomsection
  \label{sga3:VIIB:1.2.C-cokernel-double-arrow}
  % SGA3-VII-DIAG-084
  % BEGIN INCLUDED SGA3_English_current_buildable_source_20260823/source/inputs/VII/native_diagrams/exp7/viib_15_95/diag_084.tex
% SGA3-VII-DIAG-084
% Authority: Exp7B-13oct24.pdf, local p. 19.
% Crop: figures/Exp7B_localp019_example12C_cokernel_double_arrow.png
% Locator: tex/components/23_expose_VIIB_sections11_122_bicoalgebras_en.tex, Loop 1 line 183.
\(
d^\natural,e^\natural:\quad
\prod_{p(y)=z}\mathcal O_{Y,y}
\mathrel{\substack{\longrightarrow\\[-.65ex]\longrightarrow}}
\prod_{q(x)=z}\mathcal O_{X,x}
\)

% END INCLUDED SGA3_English_current_buildable_source_20260823/source/inputs/VII/native_diagrams/exp7/viib_15_95/diag_084.tex

\end{center}
where
\[
  q=p\circ d=p\circ e,
\]
and \(d^\natural\) and \(e^\natural\) are induced by the
homomorphisms \(d_x^\natural\) and \(e_x^\natural\), in the notation
of \SGARefLink{sga3:VIIB:1.1}{1.1}.

\medskip
\noindent\textbf{Definition 1.2.D.}
\phantomsection
\label{sga3:VIIB:1.2.D}
Let \(\phi:k\to\ell\) be a homomorphism of pseudocompact
rings and let \(X\) be a formal \(k\)-variety with affine algebra
\(A\). The formal scheme
\[
  X\times_{\operatorname{Spf}(k)}\operatorname{Spf}(\ell),
\]
obtained by base change, is a formal \(\ell\)-variety. We also denote
it by
\[
  X\widehat{\otimes}_k\ell.
\]
Its affine algebra is the completed tensor product
\(A\widehat{\otimes}_k\ell\) (cf.\ 0.5 and EGA~I, \S~10).

\medskip
\noindent\textbf{Remark 1.2.E.}
\phantomsection
\label{sga3:VIIB:1.2.E}
Since every formal variety over \(k\) decomposes into formal
varieties over the local components of \(k\), \emph{in certain proofs
we shall assume that \(k\) is a local pseudocompact ring}.

We now give some examples while fixing our terminology.

\medskip
\noindent\textbf{1.2.1.}
\phantomsection
\label{sga3:VIIB:1.2.1}
A \(k\)-\emph{functor} is, by definition, a covariant functor from
\(\operatorname{Alf}/k\) to \((\mathbf{Sets})\). By \SGARefLink{sga3:VIIB:1.1}{1.1} and \SGARefLink{sga3:VIIB:0.4.2}{0.4.2},
we may identify
\[
  \operatorname{Vaf}/k
  \simeq
  (\operatorname{Alp}/k)^\circ
\]
with a full subcategory of the category of \(k\)-functors by
associating to every formal \(k\)-variety \(X\) the functor
\[
  \operatorname{Alf}/k\longrightarrow(\mathbf{Sets}),
  \qquad
  C\longmapsto
  X(C)=
  \operatorname{Hom}_{\operatorname{Vaf}/k}
    \bigl(\operatorname{Spf}(C),X\bigr).
\]

% Authority: Exp7B local p. 20; combined-reader p. 533; printed p. 522.
We shall later encounter \(k\)-functors \(\mathbf F\) which associate
to every object \(C\) of \(\operatorname{Alf}/k\) a \(C\)-module
\(\mathbf F(C)\), and to every morphism
\(\phi:C\to D\) in \(\operatorname{Alf}/k\) a \(k\)-linear map
\[
  \mathbf F(\phi):\mathbf F(C)\longrightarrow\mathbf F(D)
\]
such that, for \(x\in\mathbf F(C)\) and \(\lambda\in C\),
\[
  \mathbf F(\phi)(\lambda x)
  =
  \phi(\lambda)\mathbf F(\phi)(x).
\]
By \SGARefLink{sga3:i:definition:3-1}{Expos\'e I, 3.1}, such an \(\mathbf F\) is endowed with the structure
of an \(\mathbf O_k\)-module, where \(\mathbf O_k\) denotes the
\(k\)-functor in rings which assigns to every object \(C\) of
\(\operatorname{Alf}/k\) its underlying ring.

\medskip
\noindent\textbf{Definitions.}
\phantomsection
\label{sga3:VIIB:1.2.1-definitions}

\noindent\textit{\SGARefLink{sga3:vii:paragraph:5-1:equation:i}{(i)}}
An \(\mathbf O_k\)-module \(\mathbf F\) is called
\emph{admissible} if every morphism \(\phi:C\to D\) in
\(\operatorname{Alf}/k\) induces a bijection from
\(D\otimes_C\mathbf F(C)\) onto \(\mathbf F(D)\).\footnote{%
\label{sga3:VIIB:nde:43}The editors introduced this
terminology; cf.\ editor's note~(47) in \SGARefLink{sga3:VIIB:1.2.3}{1.2.3}.}

\smallskip
\noindent\textit{\SGARefLink{sga3:vii:paragraph:5-1:equation:ii}{(ii)}}
The module \(\mathbf F\) is called \emph{flat} if it is admissible and
if, for every object \(C\) of \(\operatorname{Alf}/k\), the
\(C\)-module \(\mathbf F(C)\) is flat.\footnote{%
\label{sga3:VIIB:nde:44}It is therefore projective, since \(C\) is
Artinian.}

For example, if \(M\) is a \(k\)-module, not necessarily
pseudocompact, we denote by \(\mathbf W(M)\), as in \SGARefLink{sga3:i:definition:4-6-1}{Expos\'e I, 4.6.1},
the \(\mathbf O_k\)-module
\[
  C\longmapsto C\otimes_kM.
\]
Then \(\mathbf W(M)\) is flat whenever \(M\) is flat over \(k\).

Moreover, the functor \(M\mapsto\mathbf W(M)\) has as right adjoint
the functor which assigns to every \(\mathbf O_k\)-module
\(\mathbf F\) the \(k\)-module
\[
  \varprojlim_{\mathfrak l}\mathbf F(k/\mathfrak l),
\]
where \(\mathfrak l\) ranges over the open ideals of \(k\).

\medskip
\noindent\textbf{1.2.2.}
\phantomsection
\label{sga3:VIIB:1.2.2}
In what follows, an \(\mathbf O_k\)-module will always be denoted by a
boldface letter such as \(\mathbf F\). When \(k\) is Artinian, we shall
simply write \(F\) instead of \(\mathbf F(k)\). In that case, it is
clear that the functor
\[
  \mathbf F\longmapsto F
\]
induces an equivalence from the category of flat
\(\mathbf O_k\)-modules onto the category of flat \(k\)-modules! The
purpose of this terminology is solely to allow us to reason ``as if
\(k\) were always Artinian.''

In accordance with Expos\'e I, \S~\SGARefLink{sga3:i:definition:3-1}{3.1}, we shall use analogous
conventions for other algebraic structures. Thus an
\(\mathbf O_k\)-algebra, respectively an \(\mathbf O_k\)-coalgebra,
an \(\mathbf O_k\)-Lie algebra, or an \(\mathbf O_k\)-\(p\)-Lie
algebra, consists of an \(\mathbf O_k\)-module \(\mathbf M\) together
with, for every object \(C\) of \(\operatorname{Alf}/k\), a structure
of \(C\)-algebra, respectively \(C\)-coalgebra, \(C\)-Lie algebra, or
\(C\)-\(p\)-Lie algebra, on \(\mathbf M(C)\). We further require that
for every morphism \(\phi:C\to D\) in \(\operatorname{Alf}/k\), the
map
\[
  D\otimes_C\mathbf M(C)\longrightarrow\mathbf M(D)
\]
induced by \(\mathbf M(\phi)\) be a homomorphism of \(D\)-algebras,
respectively of \(D\)-coalgebras, and so on.

Finally, if \(\mathbf F\) and \(\mathbf G\) are two
\(\mathbf O_k\)-modules, then \(\mathbf F\otimes_k\mathbf G\) denotes
the \(\mathbf O_k\)-module
\[
  C\longmapsto\mathbf F(C)\otimes_C\mathbf G(C).
\]

% Next cursor: Exp7B local p. 20, combined-reader p. 533,
% re-edition printed p. 522, at 1.2.3 and editor's note (45),
% beginning ``We begin with the following lemma.''

% END INCLUDED SGA3_English_current_buildable_source_20260823/source/inputs/VII/tex/components/23_expose_VIIB_sections11_122_bicoalgebras_en.tex

% BEGIN INCLUDED SGA3_English_current_buildable_source_20260823/source/inputs/VII/tex/components/24_expose_VIIB_section123_anticoalgebras_en.tex
% SGA 3, Expose VII_B, English reconstruction -- Loop 1.
% Scope: authority local pages 20--23, all of 1.2.3,
% stopping immediately before 1.2.4 on local page 23.
% Combined-reader pages: 533--536. Re-edition printed pages: 522--525.
% Original-edition marginal locators in this span: 495--497.
% Authority: Polo--Gille Exp7B-13oct24.pdf,
% SHA-256 D303690E0195654BBE7C3DAA5D89EA23E24CD709A3575217234AC19AE8836694.
% Jacob Reinhold's CC BY 4.0 Markdown was consulted as comparison lineage.
% No genuine diagram occurs in this span; all displayed mathematics is native LaTeX.

\medskip
\noindent\textbf{1.2.3.}\footnote{\label{sga3:VIIB:nde:45}The editors
added the following lemma and introduced the numbering \SGARefLink{sga3:VIIB:1.2.3.A}{1.2.3.A} through
\SGARefLink{sga3:VIIB:1.2.3.F}{1.2.3.F} for later references.}
\phantomsection
\label{sga3:VIIB:1.2.3}
We begin with the following lemma.

\medskip
\noindent\textbf{Lemma 1.2.3.A.}
\phantomsection
\label{sga3:VIIB:1.2.3.A}
{\itshape Let \(k\longrightarrow K\) be a morphism of pseudocompact
rings, let \(B\) be a pseudocompact \(K\)-module, and let \(M\) be a
projective \(k\)-module, without topology. There is a canonical
isomorphism of pseudocompact \(K\)-modules
\begin{equation}
  \theta:
  \operatorname{Hom}_k(M,k)\widehat{\otimes}_k B
  \xrightarrow{\ \sim\ }
  \operatorname{Hom}_k(M,B).
  \tag{2}
  \label{sga3:VIIB:1.2.3:eq2}
\end{equation}}

Here \(M^*=\operatorname{Hom}_k(M,k)\) is endowed with the topology
defined in \SGARefLink{sga3:VIIB:0.2.2}{0.2.2}, which makes it a pseudocompact \(k\)-module. The
topology on \(\operatorname{Hom}_k(M,B)\) is defined analogously: a
basis of neighborhoods of \(0\) is formed by the following
\(K\)-submodules, where \(x\in M\) and \(B'\) is an open
\(K\)-submodule of \(B\):
\[
  \mathscr H(x,B')
  =
  \{f\in\operatorname{Hom}_k(M,B)\mid f(x)\in B'\}.
\]
Finally, \(M^*\widehat{\otimes}_kB\) is the pseudocompact
\(K\)-module obtained from \(M^*\) by base change; cf.\ \SGARefLink{sga3:VIIB:0.5.A}{0.5.A}.

With these conventions, \(\theta\) is plainly an isomorphism when
\(M=k\). Moreover, both sides of \textup{(2)}, regarded as functors in
\(M\), carry direct sums to products; in particular, they commute
with finite direct sums. It follows that \textup{\SGARefLink{sga3:VIIB:1.2.3:eq2}{(2)}} is an
isomorphism first when \(M\) is a free \(k\)-module, and then when
\(M\) is projective.

\medskip
\noindent\textbf{Definition 1.2.3.B.}
\phantomsection
\label{sga3:VIIB:1.2.3.B}
Let \(N\) be a pseudocompact \(k\)-module. We write
\(\mathbf V_k^f(N)\), or
\(\mathbf N^\dagger\),\footnote{\label{sga3:VIIB:nde:46}In fact, the
second notation will not be used; see editors' note~\SGARefLink{sga3:VIIB:nde:47}{(47)}.} for the
\(\mathbf O_k\)-module which assigns to every
\(C\in\operatorname{Ob}(\operatorname{Alf}/k)\) the \(C\)-module
\[
  (C\widehat{\otimes}_kN)^\dagger,
\]
the dual of \(C\widehat{\otimes}_kN\) in the sense of \SGARefLink{sga3:VIIB:0.2.2}{0.2.2}; that is,
the \(C\)-module
\[
  \operatorname{Hom}_c(C\widehat{\otimes}_kN,C)
  \simeq
  \operatorname{Hom}_c(N,C|_k),
\]
where \(C|_k\) denotes the \(k\)-module obtained from \(C\) by
restriction of scalars. This \(\mathbf O_k\)-module
\(\mathbf V_k^f(N)\) is called the \emph{dual} of \(N\).

If \(k_{\mathfrak m}\) is a local component of \(k\), then for every
object \(C\) of \(\operatorname{Alf}/k\),
\[
  \operatorname{Hom}_c(k_{\mathfrak m},C|_k)
  =
  C_{\mathfrak m}
  =
  C\otimes_k k_{\mathfrak m}
\]
is a projective \(C\)-module. Moreover,
\[
  D\otimes_C C_{\mathfrak m}=D_{\mathfrak m}
\]
for every morphism \(C\longrightarrow D\) in
\(\operatorname{Alf}/k\). Thus
\(\mathbf V_k^f(k_{\mathfrak m})\) is a flat
\(\mathbf O_k\)-module; cf.\ \SGARefLink{sga3:VIIB:1.2.1}{1.2.1}. Since every projective
pseudocompact \(k\)-module is a product of modules
\(k_{\mathfrak m}\) (\SGARefLink{sga3:VIIB:0.2.1}{Proposition~0.2.1}), \SGARefLink{sga3:VIIB:0.2.F}{Corollary~0.2.F} gives:

\medskip
\noindent\textbf{Corollary 1.2.3.C.}
\phantomsection
\label{sga3:VIIB:1.2.3.C}
\textit{The \(\mathbf O_k\)-module \(\mathbf V_k^f(N)\) is flat
whenever \(N\) is a projective pseudocompact \(k\)-module.}

\medskip
\noindent\textbf{Definition 1.2.3.D.}
\phantomsection
\label{sga3:VIIB:1.2.3.D}
Conversely, let \(\mathbf M\) be an admissible
\(\mathbf O_k\)-module (cf.\ \SGARefLink{sga3:VIIB:1.2.1}{1.2.1}).\footnote{
\label{sga3:VIIB:nde:47}The editors do not understand the definition
that follows unless \(\mathbf M\) is assumed to be admissible, which
is why this hypothesis has been added. They also write
\(\Gamma^*(\mathbf M)\), rather than \(\mathbf M^*\), to avoid
confusion between \(\mathbf M^*\), the module dual to the functor
\(\mathbf M\), and \(\mathbf N^\dagger\), the functor dual to the
module \(N\); cf.\ editors' note~(46). Finally, the editors have
spelled out the construction of \(\Gamma^*(\mathbf M)\).}
We call the following pseudocompact \(k\)-module the \emph{dual} of
\(\mathbf M\), and denote it by \(\Gamma^*(\mathbf M)\). As
\(\mathfrak l\) ranges over the open ideals of \(k\), endow each
\(k/\mathfrak l\)-module
\[
  \mathbf M(k/\mathfrak l)^*
  =
  \operatorname{Hom}_{k/\mathfrak l}
    \bigl(\mathbf M(k/\mathfrak l),k/\mathfrak l\bigr)
\]
with the topology described in \SGARefLink{sga3:VIIB:0.2.2}{0.2.2}; it is then a pseudocompact
\(k\)-module. Since
\[
  \mathbf M(k/\mathfrak l)
  \simeq
  \mathbf M(k/\mathfrak l')
    \otimes_{k/\mathfrak l'}(k/\mathfrak l)
  \qquad
  (\mathfrak l\supset\mathfrak l'),
\]
there are transition morphisms
\[
  \begin{split}
  \operatorname{Hom}_{k/\mathfrak l'}
    \bigl(\mathbf M(k/\mathfrak l'),k/\mathfrak l'\bigr)
  &\longrightarrow
  \operatorname{Hom}_{k/\mathfrak l'}
    \bigl(\mathbf M(k/\mathfrak l'),k/\mathfrak l\bigr)\\
  &=
  \operatorname{Hom}_{k/\mathfrak l}
    \bigl(\mathbf M(k/\mathfrak l),k/\mathfrak l\bigr).
  \end{split}
\]
By definition, \(\Gamma^*(\mathbf M)\) is the projective limit of this
filtered projective system of pseudocompact \(k\)-modules.

From now on, suppose in addition that \(\mathbf M\) is a flat
\(\mathbf O_k\)-module (cf.\ \SGARefLink{sga3:VIIB:1.2.1}{1.2.1}). Each
\(\mathbf M(k/\mathfrak l')\) is then a projective
\(k/\mathfrak l'\)-module, so the transition morphisms above are
surjective. Consequently the projections
\[
  \Gamma^*(\mathbf M)
  \longrightarrow
  \mathbf M(k/\mathfrak l)^*
\]
are also surjective, since filtered projective limits are exact in
\(\operatorname{PC}(k)\).

Let \(\tau:k\longrightarrow K\) be a morphism of pseudocompact rings.
We write \(\mathbf M\otimes_kK\), or simply \(\mathbf M_K\), for the
\(\mathbf O_K\)-functor defined as follows. If \(C\) is a
finite-length \(K\)-algebra, then the kernel of
\(k\longrightarrow C\) is an open ideal \(\mathfrak l\), and we set
\[
  \mathbf M_K(C)
  =
  \mathbf M(k/\mathfrak l)\otimes_k C.
\]
The same value is obtained from
\(\mathbf M(k/\mathfrak l')\otimes_kC\) for every open ideal
\(\mathfrak l'\subset\mathfrak l\). We then define
\(\Gamma_K^*(\mathbf M_K)\) as the projective limit, with \(I\)
ranging over the open ideals of \(K\), of the pseudocompact
\(K\)-modules
\[
  \operatorname{Hom}_{K/I}\bigl(\mathbf M_K(K/I),K/I\bigr)
  =
  \operatorname{Hom}_{k/\mathfrak l}
    \bigl(\mathbf M(k/\mathfrak l),K/I\bigr),
\]
where on the right \(\mathfrak l\) is any open ideal of \(k\) such
that \(\tau(\mathfrak l)\subset I\).

By \SGARefLink{sga3:VIIB:1.2.3.A}{Lemma~1.2.3.A}, the right-hand term identifies with
\[
  \mathbf M(k/\mathfrak l)^*
  \widehat{\otimes}_k(K/I).
\]
Since the projections
\(\Gamma^*(\mathbf M)\longrightarrow\mathbf M(k/\mathfrak l)^*\)
are surjective, the projective limit of these latter modules is
precisely the pseudocompact \(K\)-module
\[
  \Gamma^*(\mathbf M)\widehat{\otimes}_kK
\]
(cf.\ \SGARefLink{sga3:VIIB:0.5.A}{0.5.A}). We have therefore proved that, for every flat
\(\mathbf O_k\)-module \(\mathbf M\), formation of
\(\Gamma^*(\mathbf M)\) commutes with extension of the base; that is,
\begin{equation}
  \Gamma_K^*(\mathbf M\otimes_kK)
  \simeq
  \Gamma^*(\mathbf M)\widehat{\otimes}_kK.
  \tag{\(\star\)}
  \label{sga3:VIIB:1.2.3:star}
\end{equation}

\medskip
\noindent\textbf{Proposition 1.2.3.E.}
\phantomsection
\label{sga3:VIIB:1.2.3.E}
{\itshape
\begin{enumerate}
  \item[\textup{(i)}]
  The functors
  \[
    N\longmapsto\mathbf V_k^f(N),
    \qquad
    \mathbf M\longmapsto\Gamma^*(\mathbf M)
  \]
  define an anti-equivalence between the category of projective
  pseudocompact \(k\)-modules and the category of flat
  \(\mathbf O_k\)-modules.\footnote{\label{sga3:VIIB:nde:48}The
  editors added assertion \textup{\SGARefLink{sga3:vii:paragraph:5-1:equation:ii}{(ii)}} below and the proof that
  follows. The original said: ``If \(\mathbf M\) is a flat
  \(\mathbf O_k\)-module, it is clear that \(\mathbf M\) is the dual
  of \(\mathbf M^*\); hence the functors \(N\mapsto N^*\) and
  \(\mathbf M\mapsto\mathbf M^*\) define an
  anti-equivalence \ldots''}

  \item[\textup{(ii)}]
  Moreover, if \(k\longrightarrow K\) is a morphism of pseudocompact
  rings, the preceding anti-equivalence ``commutes with base change''
  in the following sense: if
  \(N\simeq\Gamma^*(\mathbf M)\), then
  \[
    N\widehat{\otimes}_kK
    \simeq
    \Gamma_K^*(\mathbf M\otimes_kK).
  \]
\end{enumerate}}

\noindent\textit{Proof.}
First, there is a natural transformation
\[
  \phi_N:
  N\longrightarrow\Gamma^*\bigl(\mathbf V_k^f(N)\bigr).
\]
The functor \(\Gamma^*\circ\mathbf V_k^f\) commutes with products by
\SGARefLink{sga3:VIIB:0.3.5}{0.3.5} and \SGARefLink{sga3:VIIB:0.2.F}{0.2.F}, so it is enough to verify that \(\phi_N\) is an
isomorphism when \(N\) is a local component \(k_{\mathfrak m}\) of
\(k\). In this case, for every open ideal \(\mathfrak l\) of \(k\)
contained in \(\mathfrak m\), the natural morphism
\[
  (k/\mathfrak l)_{\mathfrak m}
  \longrightarrow
  \operatorname{Hom}_{k/\mathfrak l}
  \left(
    \operatorname{Hom}_c
      (k_{\mathfrak m}\widehat{\otimes}_k k/\mathfrak l,
       k/\mathfrak l),
    k/\mathfrak l
  \right)
\]
is an isomorphism, which proves the assertion.

Conversely, let \(\mathbf M\) be a flat \(\mathbf O_k\)-module. We
show that
\[
  \Gamma^*(\mathbf M)
  =
  \varprojlim_{\mathfrak l}\mathbf M(k/\mathfrak l)^*
\]
is a projective object of \(\operatorname{PC}(k)\). Let
\(N\longrightarrow N'\) be a surjective morphism between objects of
\(\operatorname{LF}(k)\). By \SGARefLink{sga3:VIIB:0.2.F}{0.2.F}\,\textup{\SGARefLink{sga3:vii:paragraph:5-1:equation:i}{(i)}} and
\textup{\SGARefLink{sga3:vii:paragraph:5-1:equation:ii}{(ii)}}, it is enough to prove that the natural map
\[
  \varinjlim_{\mathfrak l}
    \operatorname{Hom}_c(\mathbf M(k/\mathfrak l)^*,N)
  \longrightarrow
  \varinjlim_{\mathfrak l}
    \operatorname{Hom}_c(\mathbf M(k/\mathfrak l)^*,N')
\]
is surjective. This is clear: \(N\) and \(N'\) are
\(k/\mathfrak l_0\)-modules for some open ideal
\(\mathfrak l_0\), and if
\(\mathfrak l\subset\mathfrak l_0\), every morphism
\(\mathbf M(k/\mathfrak l)^*\longrightarrow N'\) lifts to a morphism
\(\mathbf M(k/\mathfrak l)^*\longrightarrow N\), since
\(\mathbf M(k/\mathfrak l)^*\) is a projective object of
\(\operatorname{PC}(k/\mathfrak l)\).

There is a morphism
\[
  \psi:
  \mathbf M\longrightarrow
  \mathbf V_k^f\bigl(\Gamma^*(\mathbf M)\bigr)
\]
of functors from \(\operatorname{Alf}/k\) to
\((\mathbf{Sets})\). Let \(B\) be an object of
\(\operatorname{Alf}/k\). We claim that
\begin{equation}
  \begin{split}
  \psi(B):\mathbf M(B)
  &\longrightarrow
  \mathbf V_k^f\bigl(\Gamma^*(\mathbf M)\bigr)(B)\\
  &=
  \operatorname{Hom}_c
  \left(
    \varprojlim_{\mathfrak l}
      \bigl(\mathbf M(k/\mathfrak l)^*
        \widehat{\otimes}_kB\bigr),
    B
  \right)
  \end{split}
  \tag{1}
  \label{sga3:VIIB:1.2.3:eq1}
\end{equation}
is a bijection. In the equality above we used the fact that
\(\widehat{\otimes}\) commutes with filtered projective limits. Let
\(\mathfrak l_0\) be an open ideal of \(k\) contained in the kernel
of \(k\longrightarrow B\). By \SGARefLink{sga3:VIIB:1.2.3.A}{Lemma~1.2.3.A}, for every
\(\mathfrak l\subset\mathfrak l_0\), there is a canonical
isomorphism of pseudocompact \(B\)-modules
\[
  \mathbf M(k/\mathfrak l)^*\widehat{\otimes}_kB
  \xrightarrow{\ \sim\ }
  \operatorname{Hom}_{k/\mathfrak l}
    \bigl(\mathbf M(k/\mathfrak l),B\bigr).
\]
Since
\[
  \mathbf M(B)
  =
  \mathbf M(k/\mathfrak l)
    \otimes_{k/\mathfrak l}B,
\]
the right-hand side is
\(\operatorname{Hom}_B(\mathbf M(B),B)\). Thus the projective system
in \textup{\SGARefLink{sga3:VIIB:1.2.3:eq1}{(1)}} is constant for
\(\mathfrak l\subset\mathfrak l_0\), and \textup{(1)} reduces to the
canonical morphism
\[
  \mathbf M(B)
  \longrightarrow
  \operatorname{Hom}_c
  \bigl(\operatorname{Hom}_B(\mathbf M(B),B),B\bigr).
\]
This is an isomorphism by \SGARefLink{sga3:VIIB:0.2.2}{0.2.2}, since \(B\) is Artinian and
\(\mathbf M(B)\) is a projective \(B\)-module. This proves assertion
\textup{\SGARefLink{sga3:vii:paragraph:5-1:equation:i}{(i)}} of the proposition; assertion \textup{\SGARefLink{sga3:vii:paragraph:5-1:equation:ii}{(ii)}} follows
from the isomorphism \textup{(\(\star\))} established above.
\hfill\(\square\)

\medskip
\noindent\textbf{Remark 1.2.3.F.}
\phantomsection
\label{sga3:VIIB:1.2.3.F}
Let us return to the statement of the proposition and suppose, in
addition, that \(N\) is a topologically free pseudocompact
\(k\)-module. In this case one can choose, ``coherently,'' bases for
the \(C\)-modules \(\mathbf V_k^f(N)(C)\).

Indeed, let \((n_i)\) be a pseudobasis of \(N\) (\SGARefLink{sga3:VIIB:0.2.1}{0.2.1}), and let
\(n_i^C\) be the canonical image of \(n_i\) in
\(C\widehat{\otimes}_kN\), for
\(C\in\operatorname{Alf}/k\). Define
\(\delta_i^C\in(C\widehat{\otimes}_kN)^\dagger\) by
\[
  \delta_i^C(n_i^C)=1,
  \qquad
  \delta_i^C(n_j^C)=0\quad(i\ne j).
\]
Then \((\delta_i^C)\) is a basis of
\(\mathbf V_k^f(N)(C)\). Moreover, for every morphism
\(\varphi:C\longrightarrow D\) in \(\operatorname{Alf}/k\), the map
\(\mathbf V_k^f(N)(\varphi)\) sends \(\delta_i^C\) to
\(\delta_i^D\).

% Next cursor: Exp7B local p. 23, combined-reader p. 536,
% re-edition printed p. 525, immediately before 1.2.4.

% END INCLUDED SGA3_English_current_buildable_source_20260823/source/inputs/VII/tex/components/24_expose_VIIB_section123_anticoalgebras_en.tex

% BEGIN INCLUDED SGA3_English_current_buildable_source_20260823/source/inputs/VII/tex/components/25_expose_VIIB_sections124_126_tensor_coalgebra_en.tex
% SGA 3, Expose VII_B, English reconstruction -- Loop 1.
% Scope: authority local pages 23--26, 1.2.4 through 1.2.6 in full,
% stopping immediately before 1.3 on local page 26.
% Combined-reader pages: 536--539. Re-edition printed pages: 525--528.
% Original-edition marginal locators within this span: 497--498.
% Authority: Polo--Gille Exp7B-13oct24.pdf,
% SHA-256 D303690E0195654BBE7C3DAA5D89EA23E24CD709A3575217234AC19AE8836694.
% Jacob Reinhold's CC BY 4.0 Markdown was consulted as comparison lineage.
% No genuine diagram occurs in this span; all displayed mathematics is native LaTeX.

\medskip
\noindent\SGAInlineRefTarget{sga3:vii:paragraph:1-2-4}\textbf{1.2.4.}\footnote{Editors' note: In this paragraph the
order has been changed: \(\widehat S_k(E)\) is defined first, and it is
then stated that \(\mathbf V^f_k(E)\) represents
\(\mathbf V^f_k(E)\).}
\phantomsection
\label{sga3:VIIB:1.2.4}
For every pseudocompact \(k\)-module \(E\), let
\(\widehat S_k(E)\) denote the \emph{completed symmetric algebra} of
\(E\), defined as follows. Let \(T_k(E)\) be the direct sum of the
pseudocompact \(k\)-modules
\[
  \widehat{\bigotimes}_{k}^{\,n}E
  =
  E\widehat{\otimes}_{k}\cdots\widehat{\otimes}_{k}E
  \qquad
  \left(
    n\geq0;\quad
    \widehat{\bigotimes}_{k}^{\,0}E=k
  \right).
\]
Give \(T_k(E)\) its usual graded \(k\)-algebra multiplication. Let
\(S_k(E)\) be the quotient of \(T_k(E)\) by the homogeneous ideal
whose \(n\)-th component is the closed \(k\)-submodule of
\(\widehat{\bigotimes}_{k}^{\,n}E\) generated by the elements
\[
\begin{aligned}
&x_1\widehat{\otimes}\cdots\widehat{\otimes}x_i
 \widehat{\otimes}x_{i+1}\widehat{\otimes}\cdots
 \widehat{\otimes}x_n
\\[-0.1em]
&\qquad{}-
x_1\widehat{\otimes}\cdots\widehat{\otimes}x_{i+1}
 \widehat{\otimes}x_i\widehat{\otimes}\cdots
 \widehat{\otimes}x_n.
\end{aligned}
\]
If this quotient is denoted by \(S_k^n(E)\), then \(S_k(E)\) is a
graded \(k\)-algebra whose \(n\)-th component is \(S_k^n(E)\).

Endow \(S_k(E)\) with the linear topology defined by the set
\(\Upsilon\) of ideals \(\mathfrak l\) such that
\(S_k(E)/\mathfrak l\) is a \(k\)-module of finite length and
\[
  S_k^n(E)\cap\mathfrak l
\]
is an open submodule of \(S_k^n(E)\) for every \(n\). The profinite
algebra \(\widehat S_k(E)\) is the separated completion of \(S_k(E)\)
for this topology.\footnote{Editors' note: For example, let \(k\) be
a field, let \(E\) be a pseudocompact \(k\)-vector space, and put
\(V=\operatorname{Hom}_c(E,k)\); then \(E\simeq V^*\)
\textup{(cf.\ \SGARefLink{sga3:VIIB:0.2.2}{0.2.2})}. For every finite-dimensional subspace \(W\) of
\(V\), let \(\mathcal F(W)\) be the set of closed points of the
\(k\)-scheme
\(\mathbf V(W)=\operatorname{Spec}S_k(W^*)\), and let
\(\mathcal F(V)=\varinjlim_W\mathcal F(W)\). Then
\(\widehat S_k(E)\) is the product, over \(x\in\mathcal F(V)\), of
the pseudocompact \(k\)-algebras
\[
  \widehat S_k(E)_x
  =
  \varprojlim_W\widehat{\mathcal O}_{W,x},
\]
where \(W\) ranges over the finite-dimensional subspaces of \(V\)
such that \(x\in\mathcal F(W)\), and
\(\widehat{\mathcal O}_{W,x}\) is the separated completion of the
local ring \(\mathcal O_{\mathbf V(W),x}\) for the topology defined
by the ideals of finite codimension \textup{(which here coincides
with the \(\mathfrak m\)-adic topology)}. If \(k\) is algebraically
closed and \((v_i)_{i\in I}\) is a basis of \(V\), so that \(E\) has
a pseudobasis \((e_i)_{i\in I}\), then every local component
\(\widehat S_k(E)_x\) is isomorphic to the formal power-series ring
\(k[[e_i,\ i\in I]]\), endowed with the topology defined by the
ideals of finite codimension.}

We write
\[
  \mathbf V^f_k(E)
  =
  \operatorname{Spf}\bigl(\widehat S_k(E)\bigr)
\]
for the corresponding formal \(k\)-variety. It represents the
\(k\)-functor \(\mathbf V^f_k(E)\); that is, for every object \(C\)
of \(\operatorname{Alf}/k\), there is a canonical isomorphism
\[
\begin{aligned}
  \mathbf V^f_k(E)(C)
  &=
  \operatorname{Hom}_c(E,C|_k)
  \xrightarrow{\ \sim\ }
  \operatorname{Hom}_{\operatorname{Alp}/k}
    \bigl(\widehat S_k(E),C\bigr)
\\
  &=
  \operatorname{Hom}_{\operatorname{Vaf}/k}
    \bigl(\operatorname{Spf}(C),\mathbf V^f_k(E)\bigr).
\end{aligned}
\]
In what follows, we identify the formal variety
\(\mathbf V^f_k(E)\) with the \(k\)-functor \(\mathbf V^f_k(E)\).

\medskip
\noindent\textbf{1.2.5.}
\phantomsection
\label{sga3:VIIB:1.2.5}
By \SGARefLink{sga3:vii:paragraph:1-2-4}{1.2.4}, the zero morphism from \(E\) to \(k\) corresponds to a
morphism of profinite algebras
\[
  \pi:\widehat S_k(E)\longrightarrow k.
\]
The morphism \(\pi\) induces the zero map on \(S_k^n(E)\) for
\(n\geq1\), and defines a section of the structure morphism
\[
  \mathbf V^f_k(E)\longrightarrow\operatorname{Spf}(k).
\]

We denote by \(\mathbf V^{f,0}_k(E)\) the formal variety whose points
are the images of the points of \(\operatorname{Spf}(k)\) under the
section \(\operatorname{Spf}(\pi)\), and whose local algebras at
those points are the same as those of \(\mathbf V^f_k(E)\).
\footnote{Editors' note: The remainder of this paragraph has been
spelled out in detail.}

The affine algebra of \(\mathbf V^{f,0}_k(E)\) is thus the separated
completion of \(S_k(E)\) for the topology defined by the ideals
\(\mathfrak l\in\Upsilon\) \textup{(cf.\ \SGARefLink{sga3:vii:paragraph:1-2-4}{1.2.4})} which contain
\(S_k^n(E)\) for all sufficiently large \(n\). Hence it is the
infinite product
\[
  k[[E]]
  =
  k\times E\times S_k^2(E)\times S_k^3(E)\times\cdots.
\]

Let \(C\) be an object of \(\operatorname{Alf}/k\), let
\[
  u\in\mathbf V^f_k(E)(C)=\operatorname{Hom}_c(E,C),
\]
and let
\[
  \phi:\widehat S_k(E)\longrightarrow C
\]
be the corresponding morphism of profinite \(k\)-algebras. Then
\(\operatorname{Ker}(\phi)\) is an open ideal---that is, \(\phi\)
belongs to \(\mathbf V^{f,0}_k(E)(C)\)---if and only if
\(\operatorname{Ker}(\phi)\) contains \(S_k^n(E)\) for all
sufficiently large \(n\), or equivalently if and only if \(u(E)\) is
contained in the radical \(r(C)\) of \(C\). Consequently, for every
object \(C\) of \(\operatorname{Alf}/k\), there are canonical
isomorphisms
\[
  \mathbf V^{f,0}_k(E)(C)
  \simeq
  \operatorname{Hom}_{\operatorname{Alp}/k}(k[[E]],C)
  \simeq
  \operatorname{Hom}_c(E,r(C)).
\]

\medskip
\noindent\textbf{1.2.6.}
\phantomsection
\label{sga3:VIIB:1.2.6}
A formal \(k\)-variety \(V\) is said to be \emph{of finite length} if
its affine algebra has finite length. Similarly, if \(S\) is a scheme,
an \(S\)-scheme \(X\) is said to be \emph{of finite length} if \(X\)
is finite over \(S\) and the direct image of \(\mathcal O_X\) on
\(S\) is an \(\mathcal O_S\)-module of finite length.\footnote{Editors'
note: The following sentence has been added.}
Thus giving an \(S\)-scheme \(X\) of finite length is the same as
giving a finite set
\[
  \{s_1,\ldots,s_n\}
\]
of closed points of \(S\), together with an
\(\mathcal O_{S,s_i}\)-algebra of finite length at each such point.

It follows that the \(S\)-schemes of finite length identify with the
formal varieties of finite length over the following formal scheme
\(\widehat S\). The underlying topological space of \(\widehat S\)
is the set of closed points of \(S\), endowed with the discrete
topology. If \(s\) is such a closed point, the stalk
\(\mathcal O_{\widehat S,s}\) of the structure sheaf is the separated
completion
\[
  \widehat{\mathcal O}_{S,s}
\]
of \(\mathcal O_{S,s}\) for the linear topology defined by the ideals
of finite colength. Thus
\[
  \widehat S
  =
  \operatorname{Spf}\mathcal A(\widehat S),
\]
where \(\mathcal A(\widehat S)\) is the product of the
\(\widehat{\mathcal O}_{S,s}\), as \(s\) ranges over the closed points
of \(S\), endowed with the product topology.

\medskip
\noindent\textbf{Definition.}
\phantomsection
\label{sga3:VIIB:1.2.6.definition}
\textit{If \(X\) is an \(S\)-scheme, let
\(\widehat X/\widehat S\) denote the formal variety over
\[
  k=\mathcal A(\widehat S)
\]
defined as follows. Its underlying topological space consists of the
points \(x\in X\) such that
\[
  [\kappa(x):\kappa(s)]<\infty,
\]
where \(s\) is the image of \(x\) in \(S\). The local ring
\(\mathcal O_{\widehat X/\widehat S,x}\) is the separated completion
of \(\mathcal O_{X,x}\) for the linear topology defined by the ideals
\(I\) of \(\mathcal O_{X,x}\) for which
\(\mathcal O_{X,x}/I\) has finite length as an
\(\mathcal O_{S,s}\)-module. Since
\([\kappa(x):\kappa(s)]<\infty\), this is equivalent to requiring
that \(\mathcal O_{X,x}/I\) have finite length as an
\(\mathcal O_{X,x}\)-module.}

Let \(\operatorname{Vaf}^{\ell f}/\widehat S\) be the category of
formal varieties of finite length over \(\widehat S\), identified
with the category of \(S\)-schemes of finite length. By \SGARefLink{sga3:VIIB:1.1}{1.1} and
\SGARefLink{sga3:VIIB:1.2.1}{1.2.1}, the category \(\operatorname{Vaf}/\widehat S\) of formal
varieties over \(\widehat S\) is equivalent to the category of
left-exact contravariant functors
\[
  \operatorname{Vaf}^{\ell f}/\widehat S
  \longrightarrow
  (\mathbf{Sets}).
\]
In particular, for every \(S\)-scheme \(X\), the functor
\[
  T\longmapsto
  \operatorname{Hom}_{(\operatorname{Sch}/S)}(T,X)
\]
from \(\operatorname{Vaf}^{\ell f}/\widehat S\) to
\((\mathbf{Sets})\) is left exact and therefore corresponds to a
formal variety over \(\widehat S\). This variety is precisely
\(\widehat X/\widehat S\).\footnote{Editors' note: The following
proposition has been expanded by adding the fact that the functor
\(X\mapsto\widehat X/\widehat S\) commutes with finite projective
limits. In the original this appeared in the proof of \SGARefLink{sga3:VIIB:1.3.4}{1.3.4}; the
proof given here is more direct than the original one. This result
will also be used in \SGARefLink{sga3:VIIB:2}{Section~2} and in \SGARefLink{sga3:VIIA:3.3.1}{3.3.1}.}

\medskip
\noindent\textbf{Proposition.}
\phantomsection
\label{sga3:VIIB:1.2.6.proposition}
\textit{For every \(S\)-scheme \(X\), the functors}
\[
  \operatorname{Hom}_{\operatorname{Vaf}/\widehat S}
    (-,\widehat X/\widehat S)
  \qquad\textit{and}\qquad
  \operatorname{Hom}_{(\operatorname{Sch}/S)}(-,X)
\]
\textit{have the same restriction to
\(\operatorname{Vaf}^{\ell f}/\widehat S\). This defines a functor}
\[
  X\longmapsto\widehat X/\widehat S
\]
\textit{from \((\operatorname{Sch}/S)\) to
\(\operatorname{Vaf}/\widehat S\), and this functor commutes with
finite projective limits.}

\noindent\textit{Proof.}
It is immediate that the formal variety \(\widehat X/\widehat S\)
has the required property and that the correspondence
\(X\mapsto\widehat X/\widehat S\) is functorial. It remains to prove
the second assertion.

Put
\[
  \mathscr S=(\operatorname{Sch}/S),
  \qquad
  \mathscr V=\operatorname{Vaf}/\widehat S,
\]
and write \(\widehat X\) for \(\widehat X/\widehat S\).
By \SGARefLink{sga3:VIIB:1.2.B}{1.2.B}, the category \(\mathscr V\) has arbitrary projective
limits. Let \((X_i)_{i\in I}\) be a projective system in
\(\mathscr S\), and suppose that
\[
  X=\varprojlim_iX_i
\]
exists in \(\mathscr S\), as it does when \(I\) is finite. The functor
which associates with every object \(Y\) of \(\mathscr S\) the
functor
\[
  h_Y=\operatorname{Hom}_{\mathscr S}(-,Y)
\]
and, respectively, with every object \(V\) of \(\mathscr V\) the
functor
\[
  h_V=\operatorname{Hom}_{\mathscr V}(-,V)
\]
commutes with projective limits. Hence, for every \(S\)-scheme \(T\)
of finite length, there are functorial isomorphisms
\[
\begin{aligned}
  \operatorname{Hom}_{\mathscr S}(T,X)
  &\simeq
  \varprojlim_i\operatorname{Hom}_{\mathscr S}(T,X_i)
\\
  &\simeq
  \varprojlim_i\operatorname{Hom}_{\mathscr V}(T,\widehat X_i)
  \simeq
  \operatorname{Hom}_{\mathscr V}
    \left(T,\varprojlim_i\widehat X_i\right).
\end{aligned}
\]
Consequently, the functor \(X\mapsto\widehat X/\widehat S\) commutes
with projective limits whenever they exist in
\((\operatorname{Sch}/S)\), and in particular with finite
projective limits.
\hfill\(\square\)

% Next cursor: Exp7B local p. 26, combined-reader p. 539,
% re-edition printed p. 528, immediately before 1.3.

% END INCLUDED SGA3_English_current_buildable_source_20260823/source/inputs/VII/tex/components/25_expose_VIIB_sections124_126_tensor_coalgebra_en.tex

% BEGIN INCLUDED SGA3_English_current_buildable_source_20260823/source/inputs/VII/tex/components/26_expose_VIIB_sections13_134_topological_flatness_en.tex
% SGA 3, Expose VII_B, English reconstruction -- Loop 1.
% Scope: authority local pages 26--29, Section 1.3 through the end of 1.3.4,
% stopping immediately before 1.3.5.
% Combined-reader pages: 539--542. Re-edition printed pages: 528--531.
% Original-edition marginal locators: 499--503.
% Authority: Polo--Gille Exp7B-13oct24.pdf,
% SHA-256 D303690E0195654BBE7C3DAA5D89EA23E24CD709A3575217234AC19AE8836694.
% Jacob Reinhold's CC BY 4.0 Markdown was consulted as comparison lineage.

\subsection*{1.3}
\phantomsection
\label{sga3:VIIB:1.3}
\addcontentsline{toc}{subsection}{1.3. Monomorphisms and topological flatness}

\noindent\textbf{Proposition 1.3.}
\phantomsection
\label{sga3:VIIB:1.3.prop}
{\itshape Let \(f:X\to Y\) be a morphism of formal varieties over
\(k\), let \(A\) and \(B\) be the affine algebras of \(X\) and \(Y\),
and let \(g:B\to A\) be the morphism induced by \(f\). Then \(f\) is a
monomorphism in \(\operatorname{Vaf}/k\) if and only if \(g\) is
surjective.}

By \SGARefLink{sga3:VIIB:1.1}{1.1},\footnote{The editors expanded the beginning of the proof.}
the functor \(A\mapsto\operatorname{Spf}(A)\) is an anti-equivalence
from \(\operatorname{Alp}/k\) to \(\operatorname{Vaf}/k\).
Consequently, \(f\) is a monomorphism if and only if \(g\) is an
epimorphism, which it certainly is when \(g\) is surjective.

Conversely, suppose that \(g\) is an epimorphism. For every
\(\mathfrak n\in\Upsilon(B)\), put
\[
  A_{\mathfrak n}
  =
  A\widehat{\otimes}_B B_{\mathfrak n}.
\]
By \SGARefLink{sga3:VIIB:0.4}{0.4}, \(A\) is a pseudocompact \(B\)-module, and hence is the product
of the \(A_{\mathfrak n}\) \textup{(\SGARefLink{sga3:VIIB:0.3.6}{0.3.6})}. The morphism \(g\) is
therefore the product of the morphisms
\[
  g_{\mathfrak n}:B_{\mathfrak n}\longrightarrow A_{\mathfrak n}
\]
obtained from \(g\) by base change. These remain epimorphisms. We are
thus reduced to the case in which \(B\) is local, with maximal ideal
\(\mathfrak n\); set \(K=B/\mathfrak n\).

By Nakayama's \SGARefLink{sga3:VIIB:0.3.3}{lemma 0.3.3}, it is enough to prove that
\[
  g\widehat{\otimes}_B K
\]
is surjective. This map is obtained from \(g\) by base change and is
therefore an epimorphism in \(\operatorname{Alp}/K\). We may thus
suppose that \(B=K\) is a field. Now \(f\) is a monomorphism if and
only if the diagonal
\[
  X\longrightarrow X\times_YX
\]
is an isomorphism; equivalently, the homomorphism
\[
  A\widehat{\otimes}_K A\longrightarrow A,
  \qquad
  x\widehat{\otimes}_Kx'\longmapsto xx',
\]
is an isomorphism. Since \(K\) is a field, this implies \(A=K\).

\medskip
\noindent\textbf{Remark.}
\phantomsection
\label{sga3:VIIB:1.3.remark}
{\itshape It follows from the proposition that every monomorphism
\(f:X\to Y\) of formal varieties is an isomorphism from \(X\) onto a
necessarily closed formal subvariety of \(Y\).}\footnote{The editors
added this remark.}

\subsubsection*{1.3.1}
\phantomsection
\label{sga3:VIIB:1.3.1}
\addcontentsline{toc}{subsubsection}{\SGARefLink{sga3:VIIA:1.3.1}{1.3.1}. Effective epimorphisms}

The preceding proposition implies, in particular, that every
monomorphism \(f:X\to Y\) in \(\operatorname{Vaf}/k\) is effective
\textup{(cf.\ \SGARefLink{sga3:iv:definition:1:3}{IV 1.3})}.\footnote{In the opposite category
\((\operatorname{Vaf}/k)^\circ=\operatorname{Alp}/k\), the
corresponding map \(g:B\to A\) is an effective epimorphism: since it is
surjective, it induces an isomorphism of profinite \(k\)-algebras
\(B/I\xrightarrow{\sim}A\), where \(I=\ker g\). Hence every morphism
\(\phi:B\to C\) in \(\operatorname{Alp}/k\) that vanishes on \(I\)
descends to a morphism \(\bar\phi:A\to C\) such that
\(\phi=\bar\phi\circ g\).}
The same is not true of epimorphisms, as one sees by slightly modifying
the counterexample of Expos\'e V, \S~2.c).\footnote{Let \(k\) be a
field, \(X=\operatorname{Spf}(k[[T]])\), and
\(Y=\operatorname{Spf}(B)\), where \(B\) is the \(k\)-subalgebra of
\(A=k[[T]]\) topologically generated by \(T^3\) and \(T^4\); thus
\(B\) consists of the series \(\sum a_nT^n\) with
\(a_n=0\) for \(n=1,2,5\). Then \(X\to Y\) is an epimorphism that is
not effective: the cokernel of
\(X\times_YX\rightrightarrows X\) is
\(\operatorname{Spf}(B')\), where
\[
  B'=\{a\in A\mid
       a\widehat{\otimes}1=1\widehat{\otimes}a\},
\]
and \(B'\) contains \(T^5\).}
This is why we now consider a useful class of effective epimorphisms.

\medskip
\noindent\textbf{Lemma.}
\phantomsection
\label{sga3:VIIB:1.3.1.flatness-criterion}
\footnote{The editors added this lemma, which explains the term
``topologically flat.''}
{\itshape Let \(f:X\to Y\) be a morphism of formal \(k\)-varieties,
let \(A\) and \(B\) be the affine algebras of \(X\) and \(Y\), and let
\(f^\natural:B\to A\) be the morphism induced by \(f\). The following
conditions are equivalent:}
\begin{enumerate}
  \item[\textup{(i)}]
  {\itshape For every \(x\in X\), the homomorphism
  \[
    f_x^\natural:
    \mathcal O_{Y,f(x)}
    \longrightarrow
    \mathcal O_{X,x}
  \]
  makes \(\mathcal O_{X,x}\) a topologically free pseudocompact
  \(\mathcal O_{Y,f(x)}\)-module.}

  \item[\textup{(ii)}]
  {\itshape For every \(y\in Y\), the local component
  \[
    A_y=\prod_{f(x)=y}\mathcal O_{X,x}
  \]
  is a topologically free pseudocompact \(B_y\)-module.}

  \item[\textup{(iii)}]
  {\itshape The map \(f^\natural:B\to A\) makes \(A\) a projective
  pseudocompact \(B\)-module.}

  \item[\textup{(iv)}]
  {\itshape The functor
  \[
    \operatorname{PC}(B)\longrightarrow\operatorname{PC}(A),
    \qquad
    M\longmapsto M\widehat{\otimes}_B A,
  \]
  is exact.}
\end{enumerate}
{\itshape When these conditions hold, \(f\) is said to be
\emph{topologically flat}.}

The implications
\[
  \textup{(i)}\Longrightarrow\textup{(ii)}
  \Longleftrightarrow\textup{(iii)}
  \Longleftrightarrow\textup{(iv)}
\]
follow from \SGARefLink{sga3:VIIB:0.2.F}{0.2.F}\,\textup{\SGARefLink{sga3:VIIB:0.2.F:item:iii}{(iii)}} and \SGARefLink{sga3:VIIB:0.3.7}{0.3.7}. Conversely, suppose that
\textup{\SGARefLink{sga3:VIIB:0.2.F:item:ii}{(ii)}} holds, and let \(x\in X\), \(y=f(x)\). Since
\(\mathcal O_{X,x}\) is a direct factor of \(A_y\), it is a projective
pseudocompact \(B_y\)-module, and therefore is topologically free by
\SGARefLink{sga3:VIIB:0.2.1}{0.2.1}, because \(B_y\) is local.

A morphism \(f:X\to Y\) of formal \(k\)-varieties is called
\emph{surjective} if it induces a surjection of the underlying sets.

\medskip
\noindent\textbf{Proposition.}
\phantomsection
\label{sga3:VIIB:1.3.1.effective-epi}
{\itshape Let \(f:X\to Y\) be a surjective, topologically flat
morphism of formal \(k\)-varieties. Then \(f\) is an effective
epimorphism \textup{(cf.\ \SGARefLink{sga3:iv:definition:1:3}{IV 1.3})}.}

Let \(A\) and \(B\) be the affine algebras of \(X\) and \(Y\), and let
\(g:B\to A\) be the morphism induced by \(f\). We must prove that \(Y\)
identifies with the cokernel of
\[
  X\times_YX\rightrightarrows X;
\]
that is, for every \(\mathfrak n\in\Upsilon(B)\),
\(B_{\mathfrak n}\) must identify with the subring of
\[
  A_{\mathfrak n}
  =
  A\widehat{\otimes}_B B_{\mathfrak n}
\]
formed by the elements \(a\) such that
\[
  a\widehat{\otimes}1=1\widehat{\otimes}a.
\]

We may suppose that \(B\) is local, with maximal ideal
\(\mathfrak n\). The hypothesis then says that \(g\) makes \(A\) a
nonzero topologically free \(B\)-module. By Nakayama's \SGARefLink{sga3:VIIB:0.3.3}{lemma 0.3.3},
\(A/\overline{\mathfrak nA}\ne0\), and hence the induced morphism
\[
  g':B/\mathfrak n\longrightarrow
     A/\overline{\mathfrak nA}
\]
is injective. By \SGARefLink{sga3:VIIB:1.3.2}{Lemma 1.3.2} below, \(B\) is a direct factor of \(A\)
as a \(B\)-module; write \(A=B\oplus A'\). It follows that \(B\)
identifies with the set of \(a\in A\) such that
\[
  a\widehat{\otimes}_B1
  =
  1\widehat{\otimes}_B a.
\]

\subsubsection*{1.3.2}
\phantomsection
\label{sga3:VIIB:1.3.2}

\noindent\textbf{Lemma 1.3.2.}
{\itshape Let \(B\) be a local pseudocompact ring, let
\(\mathfrak n\) be its maximal ideal, let \(M,N\) be projective
pseudocompact \(B\)-modules, and let \(g:M\to N\) be a morphism. If
\[
  (B/\mathfrak n)\widehat{\otimes}_B g
\]
is injective, then \(g\) is an isomorphism from \(M\) onto a direct
factor of \(N\).}

Indeed, suppose
\[
  g'=(B/\mathfrak n)\widehat{\otimes}_B g
\]
is injective. Since \(B/\mathfrak n\) is a field, \(g'\) has a
retraction \(r'\). Let \(p\) and \(q\) be the canonical projections
from \(M\) and \(N\) onto
\[
  M/\overline{\mathfrak nM}
  \qquad\text{and}\qquad
  N/\overline{\mathfrak nN}.
\]
Since \(N\) is projective, there exists \(r:N\to M\) such that
\[
  p\circ r=r'\circ q;
\]
thus \(r'\) is obtained from \(r\) by passage to the quotient. Since
\(r'\circ g'\) is an isomorphism, \SGARefLink{sga3:VIIB:0.3.4}{0.3.4} shows that the same is true of
\(r\circ g\), because \(M\) is projective. If \(s\) denotes the
inverse of \(r\circ g\), then \(s\circ r\) is a retraction of \(g\).

\subsubsection*{1.3.3}
\phantomsection
\label{sga3:VIIB:1.3.3}

\noindent\textbf{Proposition 1.3.3.}
{\itshape Let \(f:X\to Y\) and \(g:Y\to Z\) be morphisms of formal
\(k\)-varieties.}
\begin{enumerate}
  \item[\textup{(i)}]
  {\itshape If \(f\) and \(g\) are topologically flat, so is
  \(g\circ f\).}

  \item[\textup{\SGAInlineRefTarget{sga3:VIIB:1.3.3:item:ii}(ii)}]
  {\itshape If \(f\) and \(g\circ f\) are topologically flat and
  \(f\) is surjective, then \(g\) is topologically flat.}

  \item[\textup{\SGAInlineRefTarget{sga3:VIIB:1.3.3:item:iii}(iii)}]
  {\itshape If \(f\) is topologically flat, then so is
  \(f\times_YY'\) after every base change \(Y'\to Y\).}
\end{enumerate}

Assertions \textup{\SGARefLink{sga3:vii:paragraph:5-1:equation:i}{(i)}} and \textup{\SGARefLink{sga3:VIIB:1.3.3:item:iii}{(iii)}} are clear. To prove
\textup{\SGARefLink{sga3:VIIB:1.3.3:item:ii}{(ii)}}, let \(A,B,C\) be the affine algebras of \(X,Y,Z\), and
let \(f':B\to A\), \(g':C\to B\) be the homomorphisms induced by
\(f,g\). Since \(g\circ f\) is topologically flat,
\(f'\circ g'\) makes \(A\) a projective pseudocompact \(C\)-module.
Likewise, \(f'\) makes \(A\) a faithful projective pseudocompact
\(B\)-module. Thus, as \(P\) ranges over pseudocompact \(C\)-modules
and \(N\) over pseudocompact \(B\)-modules, the functors
\[
  P\longmapsto P\widehat{\otimes}_C A,
  \qquad
  N\longmapsto N\widehat{\otimes}_B A
\]
are exact; the second is also faithful. Hence the functor
\[
  P\longmapsto P\widehat{\otimes}_C B
\]
is exact, and \SGARefLink{sga3:VIIB:0.3.7}{0.3.7} implies that \(B\) is a projective
pseudocompact \(C\)-module.

\subsubsection*{1.3.4}
\phantomsection
\label{sga3:VIIB:1.3.4}

\noindent\textbf{Proposition 1.3.4.}
{\itshape Let \(S\) be a scheme, let \(Y\) be a locally Noetherian
\(S\)-scheme, and let \(f:X\to Y\) be a faithfully flat
\(S\)-morphism locally of finite type. Thus \(f\) is an effective
epimorphism; equivalently, the following sequence is exact
\textup{(cf.\ \SGARefLink{sga3:iv:proposition:6:3:1}{IV 6.3.1}\,\textup{\SGARefLink{sga3:iv:proposition:6:3:1:item:iv:01}{(iv)}} and \SGARefLink{sga3:iv:definition:1:3}{IV 1.3})}:}

\begin{center}
  % SGA3-VII-DIAG-085
  % BEGIN INCLUDED SGA3_English_current_buildable_source_20260823/source/inputs/VII/native_diagrams/exp7/viib_15_95/diag_085.tex
% SGA3-VII-DIAG-085
% Authority: Exp7B-13oct24.pdf, local p. 28.
% Crop: figures/Exp7B_localp028_effective_epimorphism_sequence.png
% Locator: tex/components/26_expose_VIIB_sections13_134_topological_flatness_en.tex, Loop 1 line 298.
\(
(\ast)\qquad
X\times_YX
\mathrel{\substack{\xrightarrow{\ \operatorname{pr}_1\ }\\[-.7ex]
                   \xrightarrow[\ \operatorname{pr}_2\ ]{}}}
X\xrightarrow{\ f\ }Y.
\)

% END INCLUDED SGA3_English_current_buildable_source_20260823/source/inputs/VII/native_diagrams/exp7/viib_15_95/diag_085.tex

\end{center}

{\itshape Then the morphism of formal
\(\widehat S\)-varieties
\[
  \widehat f:
  \widehat X/\widehat S
  \longrightarrow
  \widehat Y/\widehat S
\]
\textup{(cf.\ \SGARefLink{sga3:VIIB:1.2.6}{1.2.6})} is surjective and topologically flat, and the
following sequence, obtained from the preceding one, is exact:}

\begin{center}
  % SGA3-VII-DIAG-086
  % BEGIN INCLUDED SGA3_English_current_buildable_source_20260823/source/inputs/VII/native_diagrams/exp7/viib_15_95/diag_086.tex
% SGA3-VII-DIAG-086
% Authority: Exp7B-13oct24.pdf, local p. 28.
% Crop: figures/Exp7B_localp028_completion_exact_sequence.png
% Locator: tex/components/26_expose_VIIB_sections13_134_topological_flatness_en.tex, Loop 1 line 314.
\(
\widehat{(\ast)}\qquad
\widehat{X\times_YX}/\widehat S
\mathrel{\substack{\xrightarrow{\ \widehat{\operatorname{pr}}_1\ }\\[-.7ex]
                   \xrightarrow[\ \widehat{\operatorname{pr}}_2\ ]{}}}
\widehat X/\widehat S
\xrightarrow{\ \widehat f\ }
\widehat Y/\widehat S.
\)

% END INCLUDED SGA3_English_current_buildable_source_20260823/source/inputs/VII/native_diagrams/exp7/viib_15_95/diag_086.tex

\end{center}

Let \(y\) be a point of \(Y\), with image \(s\) in \(S\), such that
\(\kappa(y)\) is a finite extension of the residue field
\(\kappa(s)\). Since \(f\) is surjective and locally of finite type,
\(f^{-1}(y)\) is nonempty and locally of finite type over
\(\kappa(y)\). Its closed points are precisely the points of
\(\widehat X/\widehat S\) mapping to \(y\). This proves that
\(\widehat f\) is surjective.

Let \(x\) be a closed point of
\(f^{-1}(y)\).\footnote{The editors expanded the argument that follows
and then used the material added in \SGARefLink{sga3:VIIB:1.2.6}{1.2.6}.} Since \(Y\) is locally
Noetherian and \(f\) is locally of finite type, the local rings
\(\mathcal O_{Y,y}\) and \(\mathcal O_{X,x}\) are Noetherian. The
powers of their maximal ideals therefore have finite colength, and the
local rings of \(\widehat Y/\widehat S\) at \(y\) and of
\(\widehat X/\widehat S\) at \(x\) are the completions
\[
  \widehat{\mathcal O}_{Y,y}
  \qquad\text{and}\qquad
  \widehat{\mathcal O}_{X,x}
\]
for their adic topologies. Since \(f\) is flat,
\(\widehat{\mathcal O}_{X,x}\) is flat over
\(\widehat{\mathcal O}_{Y,y}\) by SGA~1, IV 5.8. By \SGARefLink{sga3:VIIB:0.3.8}{0.3.8} it is
therefore a topologically free
\(\widehat{\mathcal O}_{Y,y}\)-module. Thus \(\widehat f\) is
topologically flat.

\SGARefLink{sga3:VIIB:1.3.1}{Proposition 1.3.1} now shows that \(\widehat f\) is an effective
epimorphism. In other words, the following sequence is exact, where
\(\widehat X\) is written for \(\widehat X/\widehat S\):

\begin{center}
  % SGA3-VII-DIAG-087
  % BEGIN INCLUDED SGA3_English_current_buildable_source_20260823/source/inputs/VII/native_diagrams/exp7/viib_15_95/diag_087.tex
% SGA3-VII-DIAG-087
% Authority: Exp7B-13oct24.pdf, local p. 29.
% Crop: figures/Exp7B_localp029_completed_effective_epimorphism_sequence.png
% Locator: tex/components/26_expose_VIIB_sections13_134_topological_flatness_en.tex, Loop 1 line 351.
\(
\widehat X\times_{\widehat Y}\widehat X
\mathrel{\substack{\xrightarrow{\ \widehat{\operatorname{pr}}_1\ }\\[-.7ex]
                   \xrightarrow[\ \widehat{\operatorname{pr}}_2\ ]{}}}
\widehat X\xrightarrow{\ \widehat f\ }\widehat Y.
\)

% END INCLUDED SGA3_English_current_buildable_source_20260823/source/inputs/VII/native_diagrams/exp7/viib_15_95/diag_087.tex

\end{center}

Moreover, \SGARefLink{sga3:VIIB:1.2.6}{1.2.6} gives a natural isomorphism, compatible in particular
with the projections onto \(\widehat X\):

\begin{center}
  % SGA3-VII-DIAG-088
  % BEGIN INCLUDED SGA3_English_current_buildable_source_20260823/source/inputs/VII/native_diagrams/exp7/viib_15_95/diag_088.tex
% SGA3-VII-DIAG-088
% Authority: Exp7B-13oct24.pdf, local p. 29.
% Crop: figures/Exp7B_localp029_completion_fiber_product_isomorphism.png
% Locator: tex/components/26_expose_VIIB_sections13_134_topological_flatness_en.tex, Loop 1 line 359.
\(
\widehat{X\times_YX}
\simeq
\widehat X\times_{\widehat Y}\widehat X.
\)

% END INCLUDED SGA3_English_current_buildable_source_20260823/source/inputs/VII/native_diagrams/exp7/viib_15_95/diag_088.tex

\end{center}

Consequently, the completed sequence displayed in the proposition is
exact.

% Next cursor: Exp7B local p. 29, combined-reader p. 542,
% re-edition printed p. 531, immediately before 1.3.5.

% END INCLUDED SGA3_English_current_buildable_source_20260823/source/inputs/VII/tex/components/26_expose_VIIB_sections13_134_topological_flatness_en.tex

% BEGIN INCLUDED SGA3_English_current_buildable_source_20260823/source/inputs/VII/tex/components/27_expose_VIIB_section135_topological_flatness_duality_en.tex
% SGA 3, Expose VII_B, English reconstruction -- Loop 1.
% Scope: Section 1.3.5 in full, beginning on authority local page 29
% and stopping immediately before 1.3.6 on local page 31.
% Combined-reader pages: 542--544. Re-edition printed pages: 531--533.
% The translated body occupies local pages 29--30; local page 31 fixes
% the terminal boundary. Original-edition marginal locators: 502--503.
% Editorial notes in scope: 60--61.
% Authority: Polo--Gille Exp7B-13oct24.pdf,
% SHA-256 D303690E0195654BBE7C3DAA5D89EA23E24CD709A3575217234AC19AE8836694.
% Jacob Reinhold's CC BY 4.0 Markdown was consulted as comparison lineage.
% No genuine diagram occurs in this span; all displayed mathematics is
% reconstructed in native LaTeX.

\subsubsection*{1.3.5}
\phantomsection
\label{sga3:VIIB:1.3.5}

Let \(k\) be a pseudocompact ring. A formal variety \(X\) over \(k\)
is said to be \emph{topologically flat} if its affine algebra \(A\) is
a \emph{projective} pseudocompact \(k\)-module, that is, if the
structure morphism
\[
  X\longrightarrow\operatorname{Spf}(k)
\]
is topologically flat.

Let us first note\footnote{\label{sga3:VIIB:nde:60}The editors added
the following lemma, used implicitly in the original; they also
introduced the numbering \SGARefLink{sga3:VIIB:1.3.5.A}{1.3.5.A}--\SGARefLink{sga3:VIIB:1.3.5.D}{1.3.5.D} for later references.} that
\SGARefLink{sga3:VIIB:0.2.2}{0.2.2} and \SGARefLink{sga3:VIIB:0.3.6}{0.3.6} imply the following result, analogous to
VII\textsubscript{A}, \SGARefLink{sga3:VIIA:3.1.1}{3.1.1}.

\medskip
\noindent\textbf{Lemma 1.3.5.A.}
\phantomsection
\label{sga3:VIIB:1.3.5.A}
{\itshape Suppose that \(k\) is Artinian. The functors
\[
  A\longmapsto A^\dagger=\operatorname{Hom}_c(A,k),
  \qquad
  C\longmapsto C^*=\operatorname{Hom}_k(C,k)
\]
define an anti-equivalence between the category of topologically flat
profinite \(k\)-algebras and the category of flat
\(k\)-coalgebras.}

Indeed, if \(A\) is a topologically flat profinite \(k\)-algebra,
then, by \SGARefLink{sga3:VIIB:0.3.6}{0.3.6}, there is an isomorphism of \(k\)-modules
\[
  A^\dagger\otimes_k A^\dagger
  \xrightarrow{\ \sim\ }
  \bigl(A\widehat{\otimes}A\bigr)^\dagger.
\]
Thus the multiplication
\[
  A\widehat{\otimes}A\longrightarrow A
\]
induces by duality a \(k\)-coalgebra structure on \(A^\dagger\). The
rest then follows from \SGARefLink{sga3:VIIB:0.2.2}{Proposition 0.2.2}.

Let us return to the case of an arbitrary pseudocompact ring \(k\).

\medskip
\noindent\textbf{Definition 1.3.5.B.}
\phantomsection
\label{sga3:VIIB:1.3.5.B}
To every formal \(k\)-variety \(X\) whose affine algebra \(A\) is a
projective pseudocompact \(k\)-module, one associates a flat
\(\mathbf O_k\)-coalgebra \(\mathbf H(X)\), defined as follows.

Denote by \(\mathbf H(X)\) the \(\mathbf O_k\)-module
\(\mathbf V_k^f(A)\), the ``dual of \(A\)''. It is a flat
\(\mathbf O_k\)-module, since the pseudocompact \(k\)-module
underlying \(A\) is projective \textup{(cf.\ \SGARefLink{sga3:VIIB:1.2.3.C}{1.2.3.C})}. Moreover, by
% Authority local p. 30; combined-reader p. 543; printed p. 532.
\SGARefLink{sga3:VIIB:0.3.6}{0.3.6}, there is an isomorphism
\[
  \mathbf V_k^f\bigl(A\widehat{\otimes}A\bigr)
  \simeq
  \mathbf V_k^f(A)\otimes\mathbf V_k^f(A).
\]
Consequently, the multiplication of \(A\) induces by transposition a
diagonal morphism
\[
  \mathbf H(X)=\mathbf V_k^f(A)
  \longrightarrow
  \mathbf V_k^f\bigl(A\widehat{\otimes}A\bigr)
  =\mathbf H(X)\otimes\mathbf H(X),
\]
which makes \(\mathbf H(X)\) a flat \(\mathbf O_k\)-coalgebra. We
shall say that \(\mathbf H(X)\) is the \emph{coalgebra of \(X\) over
\(\mathbf O_k\)}.

\medskip
\noindent\textbf{Definition 1.3.5.C.}
\phantomsection
\label{sga3:VIIB:1.3.5.C}
Conversely, to every \(\mathbf O_k\)-coalgebra \(\mathbf C\) one can
associate a \(k\)-functor \textup{(cf.\ \SGARefLink{sga3:VIIA:1.2.1}{1.2.1})}
\(\operatorname{Spf}^*(\mathbf C)\), defined as follows. For every
object \(B\) of \(\operatorname{Alf}/k\), set, with the notation of
VII\textsubscript{A} \SGARefLink{sga3:VIIA:3.1}{3.1},
\begin{equation}
\begin{aligned}
  \operatorname{Spf}^*(\mathbf C)(B)
  &=
  \operatorname{Hom}_{B\text{-coalg.}}\bigl(B,\mathbf C(B)\bigr)
  \\
  &=
  \bigl\{x\in\mathbf C(B)\mid
    \varepsilon_{\mathbf C(B)}(x)=1
    \ \text{and}\
    \Delta_{\mathbf C(B)}(x)=x\otimes x\bigr\}.
\end{aligned}
\tag{1}
\label{sga3:VIIB:1.3.5.C:1}
\end{equation}

Suppose in addition\footnote{\label{sga3:VIIB:nde:61}The editors
supplied the details that follow.} that the \(\mathbf O_k\)-module
underlying \(\mathbf C\) is admissible \textup{(cf.\ \SGARefLink{sga3:VIIA:1.2.1}{1.2.1})}, and put
\begin{equation}
  A
  =
  \Gamma^*(\mathbf C)
  =
  \varprojlim_{\mathfrak l}\mathbf C(k/\mathfrak l)^*.
\tag{2}
\label{sga3:VIIB:1.3.5.C:2}
\end{equation}
The algebra structures on the modules
\(\mathbf C(k/\mathfrak l)^*\) endow \(A\) with the structure of a
profinite \(k\)-algebra. For every object \(B\) of
\(\operatorname{Alf}/k\), one has
\begin{equation}
\begin{aligned}
  \operatorname{Hom}_{\operatorname{Vaf}/k}
    \bigl(\operatorname{Spf}(B),\operatorname{Spf}(A)\bigr)
  &=
  \operatorname{Hom}_{\operatorname{Alp}/k}(A,B)
  \\
  &=
  \operatorname{Hom}_{\operatorname{Alp}/B}
    \bigl(A\widehat{\otimes}B,B\bigr).
\end{aligned}
\tag{3}
\label{sga3:VIIB:1.3.5.C:3}
\end{equation}

Suppose finally that \(\mathbf C\) is a flat \(\mathbf O_k\)-module.
Then, by \SGARefLink{sga3:VIIB:1.2.3.E}{1.2.3.E},
\(A=\Gamma^*(\mathbf C)\) is a projective pseudocompact
\(k\)-module. Moreover, we saw in the proof of \emph{loc.\ cit.} that
if \(\mathfrak l_0\) is an open ideal of \(k\) contained in the kernel
of \(k\to B\), there are isomorphisms
\begin{equation}
\begin{aligned}
  A\widehat{\otimes}B
  &=
  \varprojlim_{\mathfrak l}
    \mathbf C(k/\mathfrak l)^*\widehat{\otimes}_kB
  \\
  &\simeq
  \operatorname{Hom}_{k/\mathfrak l_0}
    \bigl(\mathbf C(k/\mathfrak l_0),B\bigr)
  \\
  &\simeq
  \operatorname{Hom}_B\bigl(\mathbf C(B),B\bigr).
\end{aligned}
\tag{4}
\label{sga3:VIIB:1.3.5.C:4}
\end{equation}
We shall denote the term on the right by \(\mathbf C(B)^*\). Finally,
\SGARefLink{sga3:VIIB:1.3.5.A}{Lemma 1.3.5.A}, applied to the Artinian ring \(B\), gives a natural
isomorphism
\begin{equation}
  \operatorname{Hom}_{B\text{-coalg.}}\bigl(B,\mathbf C(B)\bigr)
  \xrightarrow{\ \sim\ }
  \operatorname{Hom}_{\operatorname{Alp}/B}
    \bigl(\mathbf C(B)^*,B\bigr).
\tag{5}
\label{sga3:VIIB:1.3.5.C:5}
\end{equation}
Combining \textup{\SGARefLink{sga3:VIIB:1.3.5.C:1}{(1)}}, \textup{\SGARefLink{sga3:VIIB:1.3.5.C:5}{(5)}}, \textup{\SGARefLink{sga3:VIIB:1.3.5.C:4}{(4)}}, \textup{\SGARefLink{sga3:VIIB:1.3.5.C:3}{(3)}},
and \textup{(2)}, we obtain, when \(\mathbf C\) is a flat
\(\mathbf O_k\)-coalgebra, an isomorphism of functors
\begin{equation}
  \operatorname{Spf}^*(\mathbf C)
  \simeq
  \operatorname{Spf}(A)
  =
  \operatorname{Spf}\bigl(\Gamma^*(\mathbf C)\bigr).
\tag{\(\star\)}
\label{sga3:VIIB:1.3.5.C:star}
\end{equation}

Consequently, if \(\mathcal A(X)\) denotes the affine algebra of a
formal \(k\)-variety \(X\), then, taking \SGARefLink{sga3:VIIB:1.2.3.E}{1.2.3.E} into account, we
obtain the following.

\medskip
\noindent\textbf{Proposition 1.3.5.D.}
\phantomsection
\label{sga3:VIIB:1.3.5.D}
{\itshape
\begin{enumerate}
  \item[\textup{(i)}]
  The functors
  \[
    X\longmapsto
    \mathbf H(X)=\mathbf V_k^f\bigl(\mathcal A(X)\bigr),
    \qquad
    \mathbf C\longmapsto
    \operatorname{Spf}^*(\mathbf C)
    =
    \operatorname{Spf}\bigl(\Gamma^*(\mathbf C)\bigr)
  \]
  define an equivalence between the category of topologically flat
  formal \(k\)-varieties and the category of flat
  \(\mathbf O_k\)-coalgebras.

  \item[\textup{(ii)}]
  Moreover, this equivalence ``commutes with base change'': if
  \(k\to K\) is a morphism of pseudocompact rings, then
  \(X\widehat{\otimes}_kK\) corresponds to
  \(\mathbf H(X)\otimes_kK\).
\end{enumerate}}

% Next cursor: Exp7B local p. 31, combined-reader p. 544,
% re-edition printed p. 533, immediately before 1.3.6.

% END INCLUDED SGA3_English_current_buildable_source_20260823/source/inputs/VII/tex/components/27_expose_VIIB_section135_topological_flatness_duality_en.tex

% BEGIN INCLUDED SGA3_English_current_buildable_source_20260823/source/inputs/VII/tex/components/28_expose_VIIB_section136_filtered_flat_coalgebras_en.tex
% SGA 3, Expose VII_B, English reconstruction -- Loop 1.
% Scope: all of Section 1.3.6 on authority local pages 31--34,
% stopping immediately before Section 1.4.
% Combined-reader pages: 544--547. Re-edition printed pages: 533--536.
% Original-edition marginal locators in this span: 504--506.
% Authority: Polo--Gille Exp7B-13oct24.pdf,
% SHA-256 D303690E0195654BBE7C3DAA5D89EA23E24CD709A3575217234AC19AE8836694.
% Jacob Reinhold's CC BY 4.0 Markdown was consulted as comparison lineage.
% The sole genuine diagram is retained as an exact authority-PDF crop in Loop 1;
% all other displayed mathematics is reconstructed in native LaTeX.

\subsubsection*{1.3.6}
\phantomsection
\label{sga3:VIIB:1.3.6}
\addcontentsline{toc}{subsubsection}{\SGARefLink{sga3:VIIB:1.3.6}{1.3.6}. Filtered flat coalgebras}

In the remainder of this Expos\'e, we shall several times define a
topologically flat formal \(k\)-variety \(X\) by exhibiting the
coalgebra \(\mathbf H(X)\). We shall then have to interpret, by means
of \(\mathbf H(X)\), certain geometric properties of \(X\). Here is
one example of this situation: suppose that a section \(\sigma\) of
the structure morphism \(X\to\operatorname{Spf}(k)\) is given, and
ask under what condition \(\sigma\) induces an isomorphism of the
underlying topological spaces.\footnote{\label{sga3:VIIB:nde:62}In
\SGARefLink{sga3:VIIB:1.3.6.A}{Lemma 1.3.6.A} below, the editors supplied detailed proofs of
parts \textup{\SGARefLink{sga3:vii:paragraph:5-1:equation:i}{(i)}} and \textup{\SGARefLink{sga3:vii:paragraph:5-1:equation:ii}{(ii)}} and added part \textup{\SGARefLink{sga3:VIIB:1.3.6.A:item:iii}{(iii)}}.}

\textit{To begin, suppose that \(k\) is Artinian.} Let
\((H,\Delta,\varepsilon)\) be a flat \(k\)-coalgebra, put
\(H^+=\ker(\varepsilon)\), and let \(A=H^*\) be the profinite
\(k\)-algebra dual to \(H\). Suppose that a morphism of
\(k\)-coalgebras \(k\to H\) is given; that is, an element
\(\phi\in H\) such that
\[
  \varepsilon(\phi)=1,
  \qquad
  \Delta(\phi)=\phi\otimes\phi.
\]
On the one hand, \(\phi\) defines a continuous algebra homomorphism
\(\Phi:A\to k\), and hence a section
\[
  \sigma:\operatorname{Spf}(k)\longrightarrow\operatorname{Spf}(A)
\]
of the projection
\(\operatorname{Spf}(A)\to\operatorname{Spf}(k)\).

On the other hand, define \(k\)-submodules of \(H\) by setting
\(H_0=k\phi\) and, for \(n\geq1\),
\[
  H_n
  =
  \{x\in H\mid
    \Delta(x)-x\otimes\phi\in H_{n-1}\otimes H^+\}.
\]
This formula also applies when \(n=0\) if \(H_{-1}=0\). Induction on
\(n\) shows that \(H_{n-1}\subset H_n\). We call
\[
  H_0\subset H_1\subset\cdots
\]
the \emph{filtration of \(H\) defined by \(\phi\)}.

\medskip
\noindent\textbf{Remark.}
\phantomsection
\label{sga3:VIIB:1.3.6.remark}
Since
\[
  \Delta(H_n)
  \subset
  H_n\otimes H_0\oplus H_{n-1}\otimes H^+,
\]
we have \(\Delta(H_n)\subset H_n\otimes H\). Since \(\Delta\) is
cocommutative---that is,
\(\tau\circ\Delta=\Delta\), where
\(\tau(a\otimes b)=b\otimes a\)---we likewise have
\(\Delta(H_n)\subset H\otimes H_n\). When \(H/H_n\) is flat over
\(k\), it follows that \(H_n\) is a subcoalgebra of \(H\)
\textup{(see also \SGARefLink{sga3:VIIB:1.3.6.A}{1.3.6.A}\,\textup{\SGARefLink{sga3:VIIB:1.3.6.A:item:iii}{(iii)}} below)}. In general,
however, the map
\[
  \Delta:H_n\longrightarrow H_n\otimes H
\]
does not factor through \(H_n\otimes H_n\).\footnote{
\label{sga3:VIIB:nde:63}For example, let \(k_0\) be a field, let
\(k=k_0[T]/(T^n)\), where \(n\geq4\), and let
\(H=k\phi\oplus kx\), with \(\varepsilon(x)=0\) and
\[
  \Delta(x)
  =
  x\otimes\phi+\phi\otimes x+t\,x\otimes x,
\]
where \(t\) is the image of \(T\) in \(k\). Then
\(H_i=k\phi\oplus t^{\,n-i}x\) for \(i=0,\ldots,n\), but for
\(2\leq i\leq n-2\), the element
\(\Delta(t^{\,n-i}x)\) does not belong to the image of
\(H_i\otimes H_i\) in \(H\otimes H\).}

\medskip
\noindent\textbf{Lemma 1.3.6.A.}
\phantomsection
\label{sga3:VIIB:1.3.6.A}
{\itshape Let \(k\) be an Artinian ring, let \(H\) be a flat
\(k\)-coalgebra, and let \(A=H^*\) be its dual profinite
\(k\)-algebra. Let \(\phi\in H\) satisfy
\(\varepsilon(\phi)=1\) and
\(\Delta(\phi)=\phi\otimes\phi\). Let
\(\Phi:A\to k\) be the continuous algebra homomorphism, let
\(\sigma:\operatorname{Spf}(k)\to\operatorname{Spf}(A)\) be the
section of
\(\operatorname{Spf}(A)\to\operatorname{Spf}(k)\), and let
\((H_n)\) be the filtration of \(H\) defined by \(\phi\). Put
\(I=\ker\Phi\).
\begin{enumerate}
  \item[\textup{(i)}]
  For every \(n\geq1\), \(H_{n-1}\) is the orthogonal in \(H\) of
  the closure \(\overline{I^n}\) of \(I^n\).

  \item[\textup{(ii)}]
  Consequently, \(\sigma\) induces a bijection of the underlying
  sets if and only if
  \[
    H=\bigcup_n H_n.
  \]
  \footnote{\label{sga3:VIIB:nde:64}In this case the coalgebra
  \(H\) is said to be \emph{connected}; cf.\ Addition \SGARefLink{sga3:VIIB:2.9}{2.9}.}

  \item[\textup{\SGAInlineRefTarget{sga3:VIIB:1.3.6.A:item:iii}(iii)}]
  If, in addition, every \(H/H_n\) is flat over \(k\), then for every
  \(n\geq0\),
  \begin{equation}
    \Delta(H_n)
    \subset
    \sum_{i=0}^n H_i\otimes H_{n-i}.
    \tag{*}
    \label{sga3:VIIB:1.3.6.A:star}
  \end{equation}
  In particular, every \(H_n\) is then a subcoalgebra of \(H\).
\end{enumerate}}

\textit{Proof.}
For every \(x\in H\), the map
\[
  A\longrightarrow k,\qquad f\longmapsto f(x),
\]
is continuous. Thus, if \(I^n\) is annihilated by \(x\), so is its
closure \(\overline{I^n}\). For \(n\geq1\), put
\[
  (\overline{I^n})^\perp
  =
  \{x\in H\mid
    f(x)=0\text{ for every }f\in\overline{I^n}\}.
\]

Suppose first that
\[
  \sigma:\mathfrak m\longmapsto\Phi^{-1}(\mathfrak m)
\]
is a bijection from \(\operatorname{Spf}(k)\) onto
\(\operatorname{Spf}(A)\). Since \(I\) is contained in the
intersection of the ideals \(\Phi^{-1}(\mathfrak m)\), it follows
from \SGARefLink{sga3:VIIB:0.1.2}{0.1.2} that the sequence of ideals
\((\overline{I^n})\) tends to zero. Let \(x\in H\). The submodule
\[
  J(x)=\{f\in A\mid f(x)=0\}
\]
is open and closed in \(A\); hence it contains
\(\overline{I^n}\) for all sufficiently large \(n\). In other words,
\(x\in(\overline{I^n})^\perp\) for all sufficiently large \(n\).

Conversely, suppose that
\[
  H=\bigcup_n(\overline{I^n})^\perp.
\]
Let \(\mathfrak p\) be an open prime ideal of \(A\). By the
definition of the topology on \(A\) \textup{(\SGARefLink{sga3:VIIB:0.2.2}{0.2.2})},
\(\mathfrak p\) contains an open \(k\)-submodule of the form
\[
  \mathscr V(x_1,\ldots,x_s)
  =
  \{f\in A\mid f(x_1)=\cdots=f(x_s)=0\}.
\]
By hypothesis, there is an integer \(n\) such that
\[
  x_1,\ldots,x_s\in(\overline{I^n})^\perp;
\]
therefore \(\overline{I^n}\subset\mathfrak p\). Moreover, since
\(k\) is Artinian,
\[
  \operatorname{Spf}(k)=\{\mathfrak m_1,\ldots,\mathfrak m_r\}
\]
is finite, and there is an integer \(t\geq1\) such that
\[
  \left(\prod_i\mathfrak m_i\right)^t=0.
\]
It follows that
\[
  \prod_i\Phi^{-1}(\mathfrak m_i)^t\subset I.
\]
Thus \(\mathfrak p\) contains the product of the ideals
\(\Phi^{-1}(\mathfrak m_i)^{tn}\). Since \(\mathfrak p\) is prime,
it contains, and hence equals, one of the
\(\Phi^{-1}(\mathfrak m_i)\). We have therefore proved that
\footnotemark
\begin{equation}
  \sigma\text{ is a bijection}
  \quad\Longleftrightarrow\quad
  H=\bigcup_n(\overline{I^n})^\perp.
  \tag{65}
  \label{sga3:VIIB:1.3.6:eq65}
\end{equation}
\footnotetext{\label{sga3:VIIB:nde:65}This remains valid without
assuming that \(H\) is cocommutative, with \(k\) still a commutative
Artinian ring. In that case a basis of neighborhoods of zero in
\(A=H^*\) is given by the two-sided ideals \(J\) such that \(A/J\)
has finite length over \(k\), and the preceding proof shows that
\(H=\bigcup_n(\overline{I^n})^\perp\) if and only if the
\(\Phi^{-1}(\mathfrak m_i)\) are the only open prime ideals of
\(A\).}

We also have \(H=k\phi\oplus H^+\); let
\(\pi:H\to H^+\) be the projection with kernel \(k\phi\). For
\(n\geq0\), let
\[
  \Delta^n:H\longrightarrow H^{\otimes(n+1)}
\]
be the iterated comultiplication, let
\(\overline{\Delta}^{\,n}\) be the composite of \(\Delta^n\) with
the projection
\[
  \pi^{\otimes(n+1)}:
  H^{\otimes(n+1)}
  \longrightarrow
  (H^+)^{\otimes(n+1)},
\]
and put
\[
  H_n'
  =
  \ker(\overline{\Delta}^{\,n})
  =
  \left\{
    x\in H\ \middle|\
    \Delta^n(x)\in
    \sum_{i=0}^n
      H^{\otimes(n-i)}\otimes H_0\otimes H^{\otimes i}
  \right\}.
\]
We set \(\overline{\Delta}^{\,0}=\operatorname{id}_H\), so that
\(H_0'=H_0\). A straightforward induction on \(n\) gives
\begin{equation}
  H_n
  \subset
  H_n'
  \subset
  (\overline{I^{n+1}})^\perp.
  \tag{*}
  \label{sga3:VIIB:1.3.6.A:star-prime}
\end{equation}

Up to this point we have not used the hypothesis that \(H\) is flat
over \(k\). Assume now that \(H\) is flat, and hence projective, over
\(k\). Then \(A^\dagger=H\) by \SGARefLink{sga3:VIIB:0.2.2}{0.2.2}. We shall prove that
\[
  H_n=(\overline{I^{n+1}})^\perp.
\]
This is clear for \(n=0\). Suppose it has been proved for \(n<r\).
Then \(H_{r-1}\) is the kernel of the homomorphism
\[
  H\longrightarrow(\overline{I^r})^\dagger.
\]
Since \(H^+\) is flat, the homomorphism
\[
  (H/H_{r-1})\otimes H^+
  \longrightarrow
  (\overline{I^r})^\dagger\otimes H^+
\]
is injective. On the other hand, the hypothesis implies that \(I\)
is a projective pseudocompact \(k\)-module, since it is a direct
factor of \(A=H^*\). Thus, by \SGARefLink{sga3:VIIB:0.3.6}{0.3.6},
\[
  \bigl(\overline{I^r}\widehat{\otimes}I\bigr)^\dagger
  \simeq
  (\overline{I^r})^\dagger\otimes I^\dagger
  =
  (\overline{I^r})^\dagger\otimes H^+.
\]
The exact sequence
\[
  \overline{I^r}\widehat{\otimes}I
  \longrightarrow A
  \longrightarrow A/\overline{I^{r+1}}
  \longrightarrow0
\]
therefore gives, by duality, the exact sequence
\begin{equation}
  (\overline{I^r})^\dagger\otimes H^+
  \xleftarrow{\ \delta\ }
  H
  \longleftarrow
  (A/\overline{I^{r+1}})^\dagger
  \longleftarrow0,
  \tag{1}
  \label{sga3:VIIB:1.3.6.A:eq1}
\end{equation}
where \(\delta\) is obtained by composing \(\Delta_H\) with the
projection
\begin{equation}
  H\otimes H
  \longrightarrow
  H\otimes H^+
  \longrightarrow
  (H/H_{r-1})\otimes H^+
  \lhook\joinrel\longrightarrow
  (\overline{I^r})^\dagger\otimes H^+.
  \tag{2}
  \label{sga3:VIIB:1.3.6.A:eq2}
\end{equation}
For every \(u\in H\), the projection of \(\Delta(u)\) onto
\(H\otimes H^+\) is \(\Delta(u)-u\otimes\phi\). Equations
\textup{\SGARefLink{sga3:VIIB:1.3.6.A:eq1}{(1)}} and \textup{\SGARefLink{sga3:VIIB:1.3.6.A:eq2}{(2)}} show that if
\[
  u\in
  (A/\overline{I^{r+1}})^\dagger
  =
  (\overline{I^{r+1}})^\perp,
\]
then \(\Delta(u)-u\otimes\phi\) belongs to the kernel of
\[
  H\otimes H^+
  \longrightarrow
  (H/H_{r-1})\otimes H^+,
\]
namely \(H_{r-1}\otimes H^+\). Thus \(u\in H_r\). This completes the
proof of \textup{\SGARefLink{sga3:vii:paragraph:5-1:equation:i}{(i)}} and \textup{\SGARefLink{sga3:vii:paragraph:5-1:equation:ii}{(ii)}}, and also proves that
\[
  H_n=\ker(\overline{\Delta}^{\,n}).
\]

We now prove \textup{\SGARefLink{sga3:VIIB:1.3.6.A:item:iii}{(iii)}}. For \(i\geq0\), put
\[
  H_i^+=H_i\cap H^+.
\]
Let \(n\geq1\). For every \(x\in H_n\), the element
\(\overline x=x-\varepsilon(x)\phi\) belongs to \(H_n^+\), and
\[
  \Delta(x)
  =
  \varepsilon(x)\phi\otimes\phi
  +\overline x\otimes\phi
  +\phi\otimes\overline x
  +\overline{\Delta}(\overline x).
\]
It is therefore enough to prove that
\[
  \overline{\Delta}(H_n^+)
  \subset
  \sum_{i=1}^{n-1}H_i^+\otimes H_{n-i}^+.
\]
For \(i=0,\ldots,n-1\), the map
\(\overline{\Delta}^{\,n}\) factors as follows:
\begin{center}
  \phantomsection
  \label{sga3:VIIB:1.3.6.A:iterated-reduced-comultiplication}
  % SGA3-VII-DIAG-089
  % BEGIN INCLUDED SGA3_English_current_buildable_source_20260823/source/inputs/VII/native_diagrams/exp7/viib_15_95/diag_089.tex
% SGA3-VII-DIAG-089
% Authority: Exp7B-13oct24.pdf, local p. 33.
% Crop: figures/Exp7B_localp033_iterated_reduced_comultiplication_diagram.png
% Locator: tex/components/28_expose_VIIB_section136_filtered_flat_coalgebras_en.tex, Loop 1 line 364.
\begin{tikzcd}[row sep=3.4em,column sep=3.8em]
H^+ \arrow[r,"\overline\Delta"]
  & H^+\otimes H^+
      \arrow[r,"{\overline\Delta^{\,i}\otimes
                    \overline\Delta^{\,n-i-1}}"]
      \arrow[d]
  & (H^+)^{\otimes(i+1)}\otimes(H^+)^{\otimes(n-i)} \\
{}
  & \displaystyle
    \frac{H^+}{H_i^+}\otimes\frac{H^+}{H_{n-i-1}^+}
      \arrow[r,"f"]
  & \displaystyle
    \frac{H^+}{H_i^+}\otimes(H^+)^{\otimes(n-i)}
      \arrow[u,"g"]
      \, .
\end{tikzcd}

% END INCLUDED SGA3_English_current_buildable_source_20260823/source/inputs/VII/native_diagrams/exp7/viib_15_95/diag_089.tex

\end{center}
Moreover, since \(H^+/H_i^+\) and
\((H^+)^{\otimes(n-i)}\) are flat, the maps \(f\) and \(g\) in this
diagram are injective. It follows that
\[
  \overline{\Delta}(H_n^+)
  \subset
  H_i^+\otimes H^+
  +
  H^+\otimes H_{n-i-1}^+
\]
for every \(i=0,\ldots,n-1\). Part \textup{\SGARefLink{sga3:VIIB:1.3.6.A:item:iii}{(iii)}} now follows from
the next lemma, applied with \(M=H^+\), \(E_0=0\), and
\(E_i=H_{i-1}^+\) for \(i\geq1\).

\medskip
\noindent\textbf{Lemma 1.3.6.B.}
\phantomsection
\label{sga3:VIIB:1.3.6.B}
{\itshape Let \(k\) be a ring, and let
\[
  0=E_0\subset E_1\subset\cdots\subset E_n\subset M
\]
be \(k\)-modules. Suppose that \(M/E_i\) is flat for every \(i\).
Then
\footnote{The French source prints \(M\otimes E_i\), rather than
\(M\otimes E_{n-i}\), in the second summand on the left. The proof
immediately below, in particular its use of \(M/E_{n-i+1}\), requires
the latter indexing, which is used here.}
\[
  \bigcap_{i=0}^n
    \bigl(E_i\otimes M+M\otimes E_{n-i}\bigr)
  =
  \sum_{i=1}^n E_i\otimes E_{n-i+1}.
\]}

\textit{Proof.}
Denote the left-hand and right-hand sides by \(K\) and \(S\),
respectively. Clearly \(S\subset K\); we prove the reverse inclusion.
For \(i=0,\ldots,n\), put
\[
  K_i=K\cap(E_i\otimes M).
\]
For \(i=1,\ldots,n\), both \(M/E_{n-i+1}\) and
\(E_i/E_{i-1}\) are flat. Consequently, the homomorphism
\(\tau_i\) below is injective, and the composite
\[
  (E_i/E_{i-1})\otimes M
  \longrightarrow
  (E_i/E_{i-1})\otimes(M/E_{n-i+1})
  \xhookrightarrow{\ \tau_i\ }
  (M/E_{i-1})\otimes(M/E_{n-i+1})
\]
has kernel
\[
  (E_i/E_{i-1})\otimes E_{n-i+1}.
\]
The image of \(K_i\) in
\((E_i\otimes M)/(E_{i-1}\otimes M)\) is both contained in and
contains this kernel. Hence
\[
  K_i=K_{i-1}+E_i\otimes E_{n-i+1},
\]
which proves the lemma.

To conclude this paragraph, return to an arbitrary pseudocompact ring
\(k\).\footnote{\label{sga3:VIIB:nde:66}The editors added what
follows; the original treated only the case in which \(k\) is
Artinian.} Let
\((\mathbf H,\Delta,\varepsilon)\) be a flat
\(\mathbf O_k\)-coalgebra, put
\(\mathbf H^+=\ker(\varepsilon)\), and let
\[
  A=\Gamma^*(\mathbf H)
\]
be the dual profinite \(k\)-algebra. Put
\(X=\operatorname{Spf}(A)\), so that
\(\mathbf H=\mathbf H(X)\) \textup{(cf.\ \SGARefLink{sga3:VIIB:1.3.5}{1.3.5})}. Suppose that a
morphism of \(\mathbf O_k\)-coalgebras
\[
  \phi:\mathbf O_k\longrightarrow\mathbf H
\]
is given. It defines a continuous homomorphism of \(k\)-algebras
\(A\to k\), and hence a section \(\sigma\) of the structure morphism
\(X\to\operatorname{Spf}(k)\).

For every object \(B\) of \(\operatorname{Alf}/k\), put
\[
  \mathbf H_0(B)=\phi(B)=B\phi_B,
  \qquad
  \phi_B=\phi(1_B)\in\mathbf H(B).
\]
For \(n\geq1\), define \(\mathbf O_k\)-submodules
\(\mathbf H_n\) of \(\mathbf H\) by
\[
  \mathbf H_n(B)
  =
  \left\{
    u\in\mathbf H(B)\ \middle|\
    \Delta(u)-u\otimes\phi_B
    \in
    \mathbf H_{n-1}(B)\otimes\mathbf H^+(B)
  \right\}.
\]
This gives a filtration
\[
  \mathbf H_0\subset\mathbf H_1\subset\cdots
\]
of \(\mathbf H(X)\). The preceding result gives:

\medskip
\noindent\textbf{Proposition.}
\phantomsection
\label{sga3:VIIB:1.3.6.proposition}
{\itshape The section \(\sigma\) induces an isomorphism of the
underlying topological spaces if and only if
\[
  \mathbf H(X)=\bigcup_n\mathbf H_n.
\]}

% Next cursor: Exp7B local p. 34, combined-reader p. 547,
% re-edition printed p. 536, immediately before Section 1.4.

% END INCLUDED SGA3_English_current_buildable_source_20260823/source/inputs/VII/tex/components/28_expose_VIIB_section136_filtered_flat_coalgebras_en.tex

% BEGIN INCLUDED SGA3_English_current_buildable_source_20260823/source/inputs/VII/tex/components/29_expose_VIIB_section14_descent_exactness_en.tex
% SGA 3, Expose VII_B, English reconstruction -- Loop 1.
% Scope: authority local pages 34--37, Theorem 1.4 and 1.4.1--1.4.3,
% stopping immediately before 1.5 on local page 37.
% Combined-reader pages: 547--550. Re-edition printed pages: 536--539.
% Original-edition marginal locators: 506--509.
% Authority: Polo--Gille Exp7B-13oct24.pdf,
% SHA-256 D303690E0195654BBE7C3DAA5D89EA23E24CD709A3575217234AC19AE8836694.
% Jacob Reinhold's CC BY 4.0 Markdown was consulted as comparison lineage.

\subsection*{1.4}
\phantomsection
\label{sga3:VIIB:1.4}
\addcontentsline{toc}{subsection}{1.4. Descent and exactness}

\noindent\textbf{Theorem 1.4.}
\phantomsection
\label{sga3:VIIB:1.4.thm}
{\itshape Let \(k\) be a pseudocompact ring and let
\[
  d_0,d_1:X_1\rightrightarrows X
\]
be an equivalence pair in \(\operatorname{Vaf}/k\)
\textup{(cf.\ Expos\'e V, \S~2.b)} such that \(d_1\) is
topologically flat.}
\begin{enumerate}
  \item[\textup{(i)}]
  {\itshape The canonical projection from \(X\) onto
  \(X/X_1=\operatorname{Coker}(d_0,d_1)\) is surjective and
  topologically flat, and the morphism
  \[
    X_1\longrightarrow X\times_{X/X_1}X
  \]
  with components \(d_0\) and \(d_1\) is an isomorphism.}

  \item[\textup{(ii)}]
  {\itshape If \(X\) is topologically flat over \(k\), then so is
  \(X/X_1\).}
\end{enumerate}

First note that \textup{\SGARefLink{sga3:vii:paragraph:5-1:equation:ii}{(ii)}} follows from \textup{\SGARefLink{sga3:vii:paragraph:5-1:equation:i}{(i)}} by
\SGARefLink{sga3:VIIB:1.3.3}{1.3.3}\,\textup{\SGARefLink{sga3:vii:paragraph:5-1:equation:ii}{(ii)}}. The proof of \textup{\SGARefLink{sga3:vii:paragraph:5-1:equation:i}{(i)}} occupies
\SGARefLink{sga3:VIIA:1.4.1}{1.4.1}, \SGARefLink{sga3:VIIA:1.4.2}{1.4.2}, and \SGARefLink{sga3:VIIB:1.4.3}{1.4.3}.

\subsubsection*{1.4.1}
\phantomsection
\label{sga3:VIIB:1.4.1}

We first show that we may reduce to the case in which \(X\) has only
one point. Since we are dealing with an equivalence pair, one sees,
as in Expos\'e V, \S~3.e), that an equivalence relation is defined on
the underlying set of \(X\) by declaring two points \(x,y\) to be
equivalent if there exists a point \(z\in X_1\) such that
\[
  d_0(z)=y
  \qquad\text{and}\qquad
  d_1(z)=x.
\]
We may plainly suppose, without loss of generality, that \(X\)
contains a single equivalence class for this relation; in other words,
that \(X/X_1\) has only one point \textup{(see the construction of
\(X/X_1\) given in \SGARefLink{sga3:VIIB:1.2}{1.2})}.

In this case, let \(x\) be any point of \(X\), and let \(U\) be the
formal variety having \(x\) as its only point and the same local ring
as \(X\) at \(x\). As in Expos\'e V, \S~6, one then sees that the
equivalence relation induced by \((d_0,d_1)\) on \(U\) still satisfies
the hypotheses of the theorem and that it is enough to prove the
theorem for this latter equivalence relation \textup{(\(U\) is a
``quasi-section'')}.

Let us briefly recall the principle of the proof given in
Expos\'e V, \S~6. Put
\[
  V=d_0^{-1}(U)=U\times_{i,d_0}X_1,
\]
where \(i\) is the inclusion of \(U\) into \(X\), and let \(u\) and
\(v\) be the morphisms with source \(V\) induced respectively by
\(d_0\) and \(d_1\):
\begin{center}
  \phantomsection
  \label{sga3:VIIB:1.4.1:quasisection-span}
  % SGA3-VII-DIAG-090
  % BEGIN INCLUDED SGA3_English_current_buildable_source_20260823/source/inputs/VII/native_diagrams/exp7/viib_15_95/diag_090.tex
% SGA3-VII-DIAG-090
% Authority: Exp7B-13oct24.pdf, local p. 35.
% Crop: figures/Exp7B_localp035_quasisection_span.png
% Locator: tex/components/29_expose_VIIB_section14_descent_exactness_en.tex, Loop 1 line 82.
\(
X\xleftarrow{\ v\ }V\xrightarrow{\ u\ }U.
\)

% END INCLUDED SGA3_English_current_buildable_source_20260823/source/inputs/VII/native_diagrams/exp7/viib_15_95/diag_090.tex

\end{center}
It is clear that \(u\) and \(v\) are topologically flat and that
\(u\) is surjective; since \(X\) contains a single equivalence class,
\(v\) is surjective. If \((v_0,v_1)\) is the inverse image of the
equivalence pair \((d_0,d_1)\) under \(v\) \textup{(cf.\ V, 3.a))},
it follows from V, 3.c) and 3.d) that \(X/X_1\) and the quotient of
\(U\) by the equivalence relation induced by \((d_0,d_1)\) both
identify with \(\operatorname{Coker}(v_0,v_1)\). As in the proof of
V, \SGARefLink{lem:sga3-V-6.1}{6.1}, one then sees that if the conclusion of \SGARefLink{sga3:VIIB:1.4.thm}{Theorem~1.4} holds for
\(U\), it also holds for \(X\).

\subsubsection*{1.4.2}
\phantomsection
\label{sga3:VIIB:1.4.2}

We are thus reduced to the case in which \(X\) has only one
point.\footnote{\label{sga3:VIIB:nde:67}Editors' note: Thus one may
suppose that \(k\) is local.}
Consider the following commutative diagram
\textup{(cf.\ V, \S~1, \textup{(0,1,2)})}:
\begin{center}
  \phantomsection
  \label{sga3:VIIB:1.4.2:equivalence-groupoid-diagram}
  % SGA3-VII-DIAG-091
  % BEGIN INCLUDED SGA3_English_current_buildable_source_20260823/source/inputs/VII/native_diagrams/exp7/viib_15_95/diag_091.tex
% SGA3-VII-DIAG-091
% Authority: Exp7B-13oct24.pdf, local p. 35.
% Crop: figures/Exp7B_localp035_equivalence_groupoid_diagram.png
% Locator: tex/components/29_expose_VIIB_section14_descent_exactness_en.tex, Loop 1 line 107.
\begin{tikzcd}[row sep=3.3em,column sep=3.9em]
X_2
  \arrow[r,shift left=.6ex,"d'_1"]
  \arrow[r,shift right=.6ex,"d'_0"']
  \arrow[d,"d'_2"']
  & X_1 \arrow[r,"d_0"] \arrow[d,"d_1"]
  & X \arrow[d] \\
X_1
  \arrow[r,shift left=.6ex,"d_1"]
  \arrow[r,shift right=.6ex,"d_0"']
  & X \arrow[r]
  & X/X_1
\end{tikzcd}

% END INCLUDED SGA3_English_current_buildable_source_20260823/source/inputs/VII/native_diagrams/exp7/viib_15_95/diag_091.tex

\end{center}
Here
\[
  X_2=X_1\times_{d_1,d_0}X_1,
\]
and \(d'_0,d'_1,d'_2\) are respectively the
morphisms\footnote{\label{sga3:VIIB:nde:68}Editors' note: Cf.\
Expos\'e V, \S~2.b).}
\[
  (x,y,z)\longmapsto(x,y),\qquad
  (x,y,z)\longmapsto(x,z),\qquad
  (x,y,z)\longmapsto(y,z).
\]

If \(B,A,A_1,A_2\) denote respectively the affine algebras of
\(X/X_1,X,X_1,X_2\), the preceding diagram induces the commutative
diagram
\begin{center}
  \phantomsection
  \label{sga3:VIIB:1.4.2:affine-algebra-diagram}
  % SGA3-VII-DIAG-092
  % BEGIN INCLUDED SGA3_English_current_buildable_source_20260823/source/inputs/VII/native_diagrams/exp7/viib_15_95/diag_092.tex
% SGA3-VII-DIAG-092
% Authority: Exp7B-13oct24.pdf, local p. 35.
% Crop: figures/Exp7B_localp035_affine_algebra_diagram.png
% Locator: tex/components/29_expose_VIIB_section14_descent_exactness_en.tex, Loop 1 line 129.
\begin{tikzcd}[row sep=3.3em,column sep=4.1em]
A_2
  & A_1
      \arrow[l,shift left=.6ex,"j_0"]
      \arrow[l,shift right=.6ex,"j_1"']
  & A \arrow[l,"i_0"'] \\
A_1 \arrow[u,"j_2"]
  & A
      \arrow[l,shift left=.6ex,"i_0"]
      \arrow[l,shift right=.6ex,"i_1"']
      \arrow[u,"i_1"']
  & B \arrow[l,"i"'] \arrow[u,"i"]
\end{tikzcd}

% END INCLUDED SGA3_English_current_buildable_source_20260823/source/inputs/VII/native_diagrams/exp7/viib_15_95/diag_092.tex

\end{center}
in which both rows are exact and the squares determined by
\((i_0,j_0)\) and \((i_1,j_1)\) are cocartesian. Since the morphism
\[
  X_1\longrightarrow X\times X
\]
with components \(d_0\) and \(d_1\) is a monomorphism by hypothesis,
the morphism
\[
  A\widehat{\otimes}_k A\longrightarrow A_1
\]
with components \(i_0\) and \(i_1\) is surjective by
\SGARefLink{sga3:VIIB:1.3.prop}{Proposition~1.3}.

\begin{samepage}
This means that \(i_1\) makes \(A_1\) a pseudocompact \(A\)-module,
topologically free by hypothesis, generated by \(i_0(A)\). Since
\(A\) is local, \SGARefLink{sga3:VIIB:1.4.3}{Lemma~1.4.3} below gives a topologically free
\(k\)-module \(V\) and a morphism of pseudocompact \(k\)-modules
\(f:V\to A\) such that the morphism
\[
  \alpha_1:
  A\widehat{\otimes}_kV\longrightarrow A_1,
  \qquad
  a\widehat{\otimes}v
  \longmapsto
  i_1(a)\cdot i_0\bigl(f(v)\bigr)
\]
is an isomorphism.
\end{samepage}
\begin{samepage}
Let
\[
  \alpha:B\widehat{\otimes}_kV\longrightarrow A
  \qquad\text{and}\qquad
  \alpha_2:A_1\widehat{\otimes}_kV\longrightarrow A_2
\]
be the morphisms
\[
  b\widehat{\otimes}v
  \longmapsto
  i(b)\cdot f(v)
  \qquad\text{and}\qquad
  a_1\widehat{\otimes}v
  \longmapsto
  j_2(a_1)\cdot j_0i_0\bigl(f(v)\bigr).
\]
\end{samepage}
In the following commutative diagram, both rows are exact and the two
left-hand squares are cocartesian. Since \(\alpha_1\) is an
isomorphism, the same is true of \(\alpha_2\), and hence of
\(\alpha\).
\begin{center}
  \phantomsection
  \label{sga3:VIIB:1.4.2:alpha-cocartesian-diagram}
  % SGA3-VII-DIAG-093
  % BEGIN INCLUDED SGA3_English_current_buildable_source_20260823/source/inputs/VII/native_diagrams/exp7/viib_15_95/diag_093.tex
% SGA3-VII-DIAG-093
% Authority: Exp7B-13oct24.pdf, local p. 36.
% Crop: figures/Exp7B_localp036_alpha_cocartesian_diagram.png
% Locator: tex/components/29_expose_VIIB_section14_descent_exactness_en.tex, Loop 1 line 186.
\begin{tikzcd}[row sep=3.4em,column sep=3.5em]
A_2
  & A_1
      \arrow[l,shift left=.6ex,"j_0"]
      \arrow[l,shift right=.6ex,"j_1"']
  & A \arrow[l,"i_0"'] \\
A_1\widehat\otimes_k V \arrow[u,"\alpha_2"]
  & A\widehat\otimes_k V
      \arrow[l,shift left=.6ex,"{i_0\widehat\otimes V}"]
      \arrow[l,shift right=.6ex,"{i_1\widehat\otimes V}"']
      \arrow[u,"\alpha_1"']
  & B\widehat\otimes_k V
      \arrow[l,"{i\widehat\otimes V}"']
      \arrow[u,"\alpha"']
      \, .
\end{tikzcd}

% END INCLUDED SGA3_English_current_buildable_source_20260823/source/inputs/VII/native_diagrams/exp7/viib_15_95/diag_093.tex

\end{center}

This shows, on the one hand, that \(A\) is topologically free over
\(B\), with pseudobasis \(f(V)\) \textup{(cf.\ \SGARefLink{sga3:VIIB:0.2.1}{0.2.1})}, and that a
pseudobasis of \(A_1\) over \(A\), where \(A_1\) is regarded as an
\(A\)-module through \(i_1\), is obtained by taking the image under
\(i_0\) of \(f(V)\). It follows that the morphism
\[
  A\widehat{\otimes}_B A\longrightarrow A_1
\]
with components \(i_1\) and \(i_0\) is an isomorphism:
\[
  A\widehat{\otimes}_B A
  \simeq
  A\widehat{\otimes}_B B\widehat{\otimes}_kV
  \simeq
  A\widehat{\otimes}_kV
  \simeq
  A_1.
\]
This proves \SGARefLink{sga3:VIIB:1.4.thm}{Theorem~1.4}, modulo the following lemma.

\subsubsection*{1.4.3}
\phantomsection
\label{sga3:VIIB:1.4.3}

\begin{samepage}
\noindent\textbf{Lemma 1.4.3.}
{\itshape Let \(k\) be a pseudocompact ring, let \(A\) be a local
profinite \(k\)-algebra, let \(A_1\) be a topologically free
\(A\)-module, and let
\[
  i_0:M\longrightarrow A_1
\]
be a morphism of pseudocompact \(k\)-modules. Suppose that the map
\[
  A\widehat{\otimes}_kM\longrightarrow A_1,
  \qquad
  a\widehat{\otimes}m
  \longmapsto
  a\cdot i_0(m),
\]
is surjective. Then there exist a topologically free \(k\)-module
\(V\) and a morphism of pseudocompact \(k\)-modules
\[
  f:V\longrightarrow M
\]
such that the map
\[
  A\widehat{\otimes}_kV\longrightarrow A_1,
  \qquad
  a\widehat{\otimes}v
  \longmapsto
  a\cdot i_0\bigl(f(v)\bigr),
\]
is bijective.}
\end{samepage}

Since every pseudocompact \(k\)-module is a quotient of a
topologically free \(k\)-module \textup{(cf.\ Editors' note
\textup{(27)})}, we may suppose, without loss of generality, that
\(M\) is topologically free. Thus write
\[
  M=\prod_{i\in I}M_i,
\]
where each \(M_i\) is a copy of \(k\). In this case,
\[
  A\widehat{\otimes}_kM
  =
  \prod_{i\in I}
  A\widehat{\otimes}_kM_i.
\]
Since the map
\[
  a\widehat{\otimes}m\longmapsto a\cdot i_0(m)
\]
is surjective and \(A_1\) is projective, its kernel is a direct
factor of \(A\widehat{\otimes}_kM\). As \(A\) is local, the exchange
theorem \textup{(\SGARefLink{sga3:VIIB:0.3.4}{0.3.4})} shows that this kernel has as a complement a
partial product
\[
  \prod_{i\in J}A\widehat{\otimes}_kM_i,
\]
where \(J\) is a subset of \(I\). One may therefore take
\[
  V=\prod_{i\in J}M_i.
\]

% Next cursor: Exp7B local p. 37, combined-reader p. 550,
% re-edition printed p. 539, immediately before 1.5.

% END INCLUDED SGA3_English_current_buildable_source_20260823/source/inputs/VII/tex/components/29_expose_VIIB_section14_descent_exactness_en.tex

% BEGIN INCLUDED SGA3_English_current_buildable_source_20260823/source/inputs/VII/tex/components/30_expose_VIIB_section15_flat_topology_en.tex
% SGA 3, Expose VII_B, English reconstruction -- Loop 1.
% Scope: authority local page 37, Section 1.5 in full, stopping before 1.6.
% Combined-reader page: 550. Re-edition printed page: 539.
% Original-edition marginal locator in this span: 510.
% Authority: Polo--Gille Exp7B-13oct24.pdf,
% SHA-256 D303690E0195654BBE7C3DAA5D89EA23E24CD709A3575217234AC19AE8836694.
% Jacob Reinhold's CC BY 4.0 Markdown was consulted as comparison lineage.

\subsection*{1.5}
\phantomsection
\label{sga3:VIIB:1.5}
\addcontentsline{toc}{subsection}{1.5. The flat topology}

Let \(k\) be a pseudocompact ring.

\medskip
\noindent\textbf{Definition.}
\phantomsection
\label{sga3:VIIB:1.5.definition}
We shall say that a family of morphisms
\[
  f_i:X_i\longrightarrow X
\]
in \(\operatorname{Vaf}/k\) is a \emph{surjective topologically flat
family} if the morphism
\[
  \coprod_i X_i\longrightarrow X
\]
induced by the \(f_i\) is surjective and topologically flat. This
means that every \(f_i\) is topologically flat and that every point of
\(X\) lies in the image of at least one \(X_i\).

It follows from \SGARefLink{sga3:VIIB:1.3.3}{1.3.3} that surjective topologically flat families
define a pretopology on \(\operatorname{Vaf}/k\) \textup{(IV,
\SGARefLink{sga3:iv:definition:4:2:5}{4.2.5})}. The corresponding topology will be called the \emph{flat
topology} on \(\operatorname{Vaf}/k\).

By IV, \SGARefLink{sga3:iv:corollary:4:3:5}{4.3.5}, a functor
\[
  F:(\operatorname{Vaf}/k)^{\mathrm{op}}\longrightarrow
  (\mathbf{Sets})
\]
is a sheaf for the flat topology if and only if \(F\) carries every
direct sum to a direct product and the sequence
\begin{center}
  \phantomsection
  \label{sga3:VIIB:1.5:flat-sheaf-sequence}
  % SGA3-VII-DIAG-094
  % BEGIN INCLUDED SGA3_English_current_buildable_source_20260823/source/inputs/VII/native_diagrams/exp7/viib_15_95/diag_094.tex
% SGA3-VII-DIAG-094
% Authority: Exp7B-13oct24.pdf, local p. 37.
% Crop: figures/Exp7B_localp037_flat_sheaf_equalizer_sequence.png
% Locator: tex/components/30_expose_VIIB_section15_flat_topology_en.tex, Loop 1 line 48.
\(
F(Y)\xrightarrow{\ F(f)\ }F(X)
\mathrel{\substack{\xrightarrow{\ F(\operatorname{pr}_1)\ }\\[-.7ex]
                   \xrightarrow[\ F(\operatorname{pr}_2)\ ]{}}}
F(X\times_YX)
\)

% END INCLUDED SGA3_English_current_buildable_source_20260823/source/inputs/VII/native_diagrams/exp7/viib_15_95/diag_094.tex

\end{center}
is exact for every surjective topologically flat morphism
\(f:X\to Y\).

By IV, \SGARefLink{sga3:iv:subsection:4:5}{4.5},\footnote{\label{sga3:VIIB:nde:69}The editors modified the
original in what follows. In particular, they replaced the statement:
``if \(d_0,d_1:X_1\rightrightarrows X\) is an equivalence relation
such that \(d_1\) is topologically flat, formation of the quotient
commutes with \(h\),'' by the theorem below.}
\SGARefLink{sga3:VIIB:1.3.1}{Proposition~1.3.1} implies that the flat topology is coarser than the
canonical topology. In other words, for every object \(X\) of
\(\operatorname{Vaf}/k\), the functor
\[
  h_X:T\longmapsto
  \operatorname{Hom}_{\operatorname{Vaf}/k}(T,X)
\]
is a sheaf for the flat topology. In what follows, as usual
\textup{(cf.\ Exposé I)}, we identify \(X\) with \(h_X\).

By IV, \SGARefLink{sga3:iv:lemma:4:6:5}{4.6.5}, \SGARefLink{sga3:VIIB:1.4.thm}{Theorem~1.4} may be reformulated as follows.

\medskip
\noindent\textbf{Theorem.}
\phantomsection
\label{sga3:VIIB:1.5.theorem}
\textit{Let \(k\) be a pseudocompact ring, let
\(d_0,d_1:X_1\rightrightarrows X\) be an equivalence pair in
\(\operatorname{Vaf}/k\), and let \(X/X_1\) be the quotient formal
variety \textup{(that is,
\(\operatorname{Coker}(d_0,d_1)\), cf.\ \SGARefLink{sga3:VIIB:1.2}{1.2})}. If \(d_1\) is
topologically flat, then \(X/X_1\) represents the quotient sheaf for
the flat topology.}

% END INCLUDED SGA3_English_current_buildable_source_20260823/source/inputs/VII/tex/components/30_expose_VIIB_section15_flat_topology_en.tex

% BEGIN INCLUDED SGA3_English_current_buildable_source_20260823/source/inputs/VII/tex/components/31_expose_VIIB_sections16_16C_formally_etale_en.tex
% SGA 3, Expose VII_B, English reconstruction -- Loop 1.
% Scope: authority local pages 37--39, Section 1.6 opening through 1.6.C,
% stopping immediately before Lemma 1.6.D.
% Combined-reader pages: 550--552. Re-edition printed pages: 539--541.
% Original-edition marginal locators: 510--511.
% Authority: Polo--Gille Exp7B-13oct24.pdf,
% SHA-256 D303690E0195654BBE7C3DAA5D89EA23E24CD709A3575217234AC19AE8836694.
% Jacob Reinhold's CC BY 4.0 Markdown was consulted as comparison lineage.
% No genuine diagram occurs in this span; all mathematics is native LaTeX.

\subsection*{1.6}
\phantomsection
\label{sga3:VIIB:1.6}
\addcontentsline{toc}{subsection}{\SGARefLink{sga3:VIIB:1.6}{1.6}. Étale formal varieties}

To complete these generalities on formal varieties, it remains to
define briefly the formal varieties that are étale over \(k\).
\footnote{\label{sga3:VIIB:nde:70}The original continued as follows:
``A formal variety \(X\) over \(k\) is said to be étale if the
diagonal morphism
\(\Delta_X:X\to X\times X\) is a local isomorphism; that is, if
\(\Delta_X\) induces an isomorphism from
\(\mathcal O_{X\times X,\Delta_X(x)}\) onto
\(\mathcal O_{X,x}\) for every point \(x\) of \(X\). It follows
readily from SGA~I that this formulation is equivalent to each of the
following two: the formal variety \(X\) is topologically flat and,
for every point \(x\in X\) projecting to
\(s\in\operatorname{Spf}(k)\),
\(\mathcal O_{X,x}\widehat\otimes_k\kappa(s)\) is a finite separable
extension of the residue field \(\kappa(s)\); or, if \(A\) denotes the
affine algebra of \(X\), the local components \textup{(\SGARefLink{sga3:VIIB:0.1}{0.1})} of
\(A\widehat\otimes_k(k/\mathfrak l)\) are finite étale algebras over
\(k/\mathfrak l\), for every open ideal \(\mathfrak l\) of \(k\).''
In what follows, the editors corrected the omission of the flatness
hypothesis from the first condition above and supplied details for
the equivalence of the conditions.}

\medskip
\noindent\textbf{Recollections \SGARefLink{sga3:VIIB:1.6.A}{1.6.A}.}
\phantomsection
\label{sga3:VIIB:1.6.A}
\textit{(i)} Recall first \textup{(cf.\ EGA \(0_{\mathrm{IV}}\),
19.10.2)} that a topological \(k\)-algebra \(A\) is said to be
\emph{formally étale over \(k\)}---for the given topologies on \(k\)
and \(A\), respectively for the discrete topologies---if, for every
discrete topological \(k\)-algebra \(C\) \textup{(not necessarily
Artinian)} and every nilpotent ideal \(J\) of \(C\), every continuous
morphism of \(k\)-algebras
\[
  A\longrightarrow C/J
\]
lifts uniquely to a continuous morphism \(A\to C\), where \(A\) is
endowed with the given topology, respectively with the discrete
topology.

This property is plainly preserved under base change. Thus, for
every morphism \(k\to k'\) of pseudocompact rings,
\[
  A\widehat\otimes_k k'
\]
is formally étale over \(k'\). It is also enough to verify the
lifting condition for square-zero ideals \(J\); cf.\ EGA
\(\mathrm{IV}_4\), 17.1.2~(ii).

The algebra \(A\) is said to be \emph{étale over \(k\)} if it is
formally étale over \(k\) for the discrete topologies and is, in
addition, a \(k\)-algebra of finite presentation; cf.\ EGA
\(\mathrm{IV}_4\), 17.3.2~(ii). In what follows, when \(k\) is a
pseudocompact ring and \(A\) a profinite \(k\)-algebra, ``formally
étale'' refers to the given topologies unless the contrary is stated.

\smallskip
\noindent\textit{(ii)} Recall also that if \(F\in k[T]\) is a
separable monic polynomial of degree \(d\geq1\), meaning that the
ideal generated by \(F\) and its derivative \(F'\) is all of
\(k[T]\), then
\[
  B=k[T]/(F),
\]
which is free of rank \(d\) over \(k\) and is endowed with the product
topology, is formally étale over \(k\).

Indeed, let \(C\) be a discrete \(k\)-algebra, so that the kernel
\(\mathfrak l\) of \(k\to C\) is an open ideal of \(k\); let \(J\)
be a square-zero ideal of \(C\); and let
\(\phi:B\to C/J\) be a continuous morphism of \(k\)-algebras. Since
\(B\) is a free \(k\)-module of finite rank,
\(\mathfrak lB\) is an open ideal of \(B\), and hence every lift
\(\Phi:B\to C\) of \(\phi\) is automatically continuous.

Let \(t\) be the image of \(T\) in \(B\), and choose any lift
\(u_0\in C\) of \(\phi(t)\). Then \(F(u_0)\in J\), since
\(\phi(t)\) is a root of \(F\). There exist \(G,H\in k[T]\) such that
\[
  GF+HF'=1.
\]
Consequently,
\[
  H(u_0)F'(u_0)=1-G(u_0)F(u_0).
\]
The right-hand side is invertible because \(F(u_0)\) is square-zero,
and therefore \(F'(u_0)\) is invertible. We seek \(h\in J\) such that
\(u=u_0+h\) is a root of \(F\). This is equivalent to
\[
  0=F(u)=F(u_0)+F'(u_0)h,
\]
whose unique solution is
\[
  h=-F'(u_0)^{-1}F(u_0)\in J.
\]
The same proof, with no continuity assumptions on the morphisms
\(k\to C\), \(\phi\), and \(\Phi\), also shows that \(B\) is an
étale \(k\)-algebra.

\smallskip
\noindent\textit{(iii)} Finally, if
\[
  A=A_1\times\cdots\times A_n
\]
is a finite product, then \(A\) is formally étale over \(k\) if and
only if every \(A_i\) is.\footnote{\label{sga3:VIIB:nde:71}The
editors point out that the proof given in EGA
\(0_{\mathrm{IV}}\), 19.3.5~(v), is erroneous.}
It is enough to prove this for \(n=2\). Let \(e=1_{A_1}\) be the
idempotent such that \(A_1=Ae\) and \(A_2=A(1-e)\), and suppose that
a continuous morphism
\[
  \phi:A\longrightarrow C/J
\]
is given, where \(J\) is square-zero. The polynomial
\(F=X^2-X\) is separable, since
\[
  F'=2X-1
  \qquad\text{and}\qquad
  (F')^2-4F=1.
\]
Therefore the idempotent \(\phi(e)\) of \(C/J\) lifts uniquely to an
idempotent \(f\) of \(C\). Hence
\[
  C=Cf\oplus C(1-f),
\]
and giving a lift \(\Phi\) is equivalent to giving two morphisms
\[
  \Phi_1:A_1\longrightarrow Cf,
  \qquad
  \Phi_2:A_2\longrightarrow C(1-f),
\]
lifting the respective restrictions of \(\phi\). The same argument
shows that if \(e\) is an idempotent of \(k\) and \(A=Ae\), then
\(A\) is formally étale over \(k\) if and only if it is formally
étale over the local component \(k_e\), identified with \(ke\).

Now let \(X\) be a formal \(k\)-variety with affine algebra \(A\).
If \(x\) is a point of \(X\), that is, an open maximal ideal
\(\mathfrak m\) of \(A\), and \(s\) is its image in
\(\operatorname{Spf}(k)\), we denote by \(k_{\mathfrak m}\), or
\(k_s\), the local component of \(k\) corresponding to \(s\).

\medskip
\noindent\textbf{Definition 1.6.B.}
\phantomsection
\label{sga3:VIIB:1.6.B}
\textit{The following conditions are equivalent:}
\begin{enumerate}
\item[\textup{(a)}] \(A\) is formally étale over \(k\).
\item[\textup{(b)}] For every \(\mathfrak m\in\Upsilon(A)\),
  \(A_{\mathfrak m}\) is formally étale over \(k\), or equivalently
  over \(k_{\mathfrak m}\).
\item[\textup{(c)}] For every open ideal \(\mathfrak l\) of \(k\),
  \[
    A\widehat\otimes_k(k/\mathfrak l)
    =
    A/A\mathfrak l
  \]
  is formally étale over \(k/\mathfrak l\).
\item[\textup{(d)}] For every \(\mathfrak m\in\Upsilon(A)\) and every
  open ideal \(\mathfrak l\) of \(k_{\mathfrak m}\),
  \(A_{\mathfrak m}/A_{\mathfrak m}\mathfrak l\) is formally étale
  over \(k_{\mathfrak m}/\mathfrak l\).
\end{enumerate}

We say that \(X\) is \emph{étale over \(k\)} if these conditions
hold, and write
\[
  \operatorname{Vaf}^{\mathrm{\acute{e}t}}/k
\]
for the full subcategory of \(\operatorname{Vaf}/k\) consisting of
formal varieties that are étale over \(k\).
\footnote{\label{sga3:VIIB:nde:72}The category
\(\operatorname{Vaf}^{\mathrm{\acute{e}t}}/k\) is stable under finite
projective limits.}

Indeed, suppose that
\[
  \phi:A\longrightarrow C/J
\]
is a continuous morphism of \(k\)-algebras, where \(C\) is a
discrete \(k\)-algebra and \(J\) is square-zero. Then
\(I=\operatorname{Ker}(\phi)\) is an open ideal of \(A\), so \(A/I\)
is Artinian. Thus \(I\) is contained in only finitely many open
maximal ideals
\[
  \mathfrak m_1,\ldots,\mathfrak m_r.
\]
It therefore contains the product of the components
\(A_{\mathfrak m}\) for
\(\mathfrak m\notin\{\mathfrak m_1,\ldots,\mathfrak m_r\}\). This
product is \(A(1-e)\), where \(e\) is the idempotent of \(A\) for
which
\[
  Ae=\prod_{i=1}^{r}A_{\mathfrak m_i}.
\]
Hence \(\phi(e)=1_{C/J}\), and giving a continuous lift of \(\phi\)
is equivalent to giving a continuous lift of the induced morphism
\[
  Ae\simeq A/A(1-e)\longrightarrow C/J.
\]

On the other hand, by editors' note~(24),
\[
  A\simeq\varprojlim_{\mathfrak l}A/A\mathfrak l.
\]
The equivalence of the stated conditions follows readily from these
observations and the preceding recollections.

\medskip
\noindent\textbf{Definitions 1.6.C.}
\phantomsection
\label{sga3:VIIB:1.6.C}
Put
\[
  \kappa(k)
  =
  \prod_{s\in\operatorname{Spf}(k)}\kappa(s),
\]
endowed with the product topology. Thus the formal variety
\(\operatorname{Spf}(\kappa(k))\) is the direct sum of the schemes
\(\operatorname{Spec}\kappa(s)\), for
\(s\in\operatorname{Spf}(k)\). We write \(S_{\kappa(k)}\) for the
scheme that is the direct sum of the same
\(\operatorname{Spec}\kappa(s)\).

For every formal variety \(X\) over \(k\), put
\[
  X_\kappa
  =
  X\widehat\otimes_k\kappa(k).
\]
This is the formal variety over \(\kappa(k)\) obtained by base
change. It has the same points as \(X\), and if
\(x\in X\) projects to \(s\in\operatorname{Spf}(k)\), then
\[
  \mathcal O_{X_\kappa,x}
  =
  \mathcal O_{X,x}\widehat\otimes_k\kappa(s).
\]
The base-change functor
\[
  \operatorname{Vaf}/k
  \longrightarrow
  \operatorname{Vaf}/\kappa(k)
\]
sends \(\operatorname{Vaf}^{\mathrm{\acute{e}t}}/k\) into
\(\operatorname{Vaf}^{\mathrm{\acute{e}t}}/\kappa(k)\), by \SGARefLink{sga3:VIIB:1.6.A}{1.6.A}~\SGARefLink{sga3:vii:paragraph:5-1:equation:i}{(i)}.

% END INCLUDED SGA3_English_current_buildable_source_20260823/source/inputs/VII/tex/components/31_expose_VIIB_sections16_16C_formally_etale_en.tex

% BEGIN INCLUDED SGA3_English_current_buildable_source_20260823/source/inputs/VII/tex/components/32_expose_VIIB_section16D_residue_base_change_en.tex
% SGA 3, Expose VII_B, English reconstruction -- Loop 1.
% Scope: authority local pages 39--40, Lemma 1.6.D in full,
% stopping immediately before Proposition 1.6.E.
% Combined-reader pages: 552--553. Re-edition printed pages: 541--542.
% Original-edition marginal locator: 511.
% Authority: Polo--Gille Exp7B-13oct24.pdf,
% SHA-256 D303690E0195654BBE7C3DAA5D89EA23E24CD709A3575217234AC19AE8836694.
% Jacob Reinhold's CC BY 4.0 Markdown was consulted as comparison lineage.
% No genuine diagram occurs in this span; all mathematics is native LaTeX.

One then has the following result \textup{(cf.\ SGA~1, \SGARefLink{sga3:i:remark:6-2}{I 6.2})}.

\medskip
\noindent\textbf{Lemma 1.6.D.}
\phantomsection
\label{sga3:VIIB:1.6.D}
\textit{For every
\[
  Y\in\operatorname{Ob}(\operatorname{Vaf}^{\mathrm{\acute{e}t}}/k)
  \qquad\text{and}\qquad
  X\in\operatorname{Ob}(\operatorname{Vaf}/k),
\]
the canonical map
\[
  \operatorname{Hom}_{\operatorname{Vaf}/k}(X,Y)
  \longrightarrow
  \operatorname{Hom}_{\operatorname{Vaf}/\kappa(k)}
    (X_\kappa,Y_\kappa)
\]
is bijective. In particular, the functor
\[
  \operatorname{Vaf}^{\mathrm{\acute{e}t}}/k
  \longrightarrow
  \operatorname{Vaf}^{\mathrm{\acute{e}t}}/\kappa(k)
\]
is fully faithful; below it will be shown to be an equivalence.}

Indeed, let \(A\), respectively \(B\), be the affine algebra of
\(X\), respectively \(Y\), and let \(r\) be the radical of \(k\).
Suppose that a morphism
\[
  B\widehat\otimes_k\kappa(k)
  \longrightarrow
  A\widehat\otimes_k\kappa(k)
\]
is given, or equivalently a morphism
\[
  \phi:B\longrightarrow A/rA.
\]

For every open ideal \(\mathfrak l\) of \(k\), there is an integer
\(n\geq1\) such that
\[
  r^n\subset\mathfrak l.
\]
Because the multiplication map
\[
  m^{n-1}:A^n\longrightarrow A
\]
is continuous, it follows that
\[
  (rA)^n\subset\mathfrak lA.
\]
Thus \(rA/\mathfrak lA\) is a nilpotent ideal of
\(A/\mathfrak lA\). Consequently, \(\phi\) lifts uniquely to a
morphism
\[
  \phi_{\mathfrak l}:B\longrightarrow A/\mathfrak lA.
\]
By uniqueness, these morphisms form a projective system and therefore
give a continuous morphism
\[
  \Phi:
  B\longrightarrow
  \varprojlim_{\mathfrak l}A/\mathfrak lA
  =
  A.
\]
Moreover, \(\Phi\) is unique: if \(\Phi'\) is another lift of
\(\phi\), then \(\Phi'\) and \(\Phi\) agree modulo
\(\mathfrak lA\) for every \(\mathfrak l\), and hence are equal.

% END INCLUDED SGA3_English_current_buildable_source_20260823/source/inputs/VII/tex/components/32_expose_VIIB_section16D_residue_base_change_en.tex

% BEGIN INCLUDED SGA3_English_current_buildable_source_20260823/source/inputs/VII/tex/components/33_expose_VIIB_section16E_etale_criteria_en.tex
% SGA 3, Expose VII_B, English reconstruction -- Loop 1.
% Scope: authority local pages 40--41, Proposition 1.6.E in full,
% stopping before editor's note 73 and Proposition 1.6.F.
% Combined-reader pages: 553--554. Re-edition printed pages: 542--543.
% Original-edition marginal locator: 511.
% Authority: Polo--Gille Exp7B-13oct24.pdf,
% SHA-256 D303690E0195654BBE7C3DAA5D89EA23E24CD709A3575217234AC19AE8836694.
% Jacob Reinhold's CC BY 4.0 Markdown was consulted as comparison lineage.
% No genuine diagram occurs in this span; all mathematics is native LaTeX.

\medskip
\noindent\textbf{Proposition 1.6.E.}
\phantomsection
\label{sga3:VIIB:1.6.E}
\textit{(a) Let \(X\) be a formal variety over \(k\), with affine
algebra \(A\). The following conditions are equivalent:
\begin{enumerate}
\item[\textup{(i)}] \(X\) is étale over \(k\).
\item[\textup{\SGAInlineRefTarget{sga3:VIIB:1.6.E:item:ii}(ii)}] \(X\) is topologically flat over \(k\), and the
  diagonal morphism
  \[
    \Delta:X\longrightarrow X\times X
  \]
  is a local isomorphism; that is, for every point \(x\) of \(X\),
  \(\Delta_X\) induces an isomorphism
  \[
    \mathcal O_{X\times X,\Delta(x)}
    \xrightarrow{\ \sim\ }
    \mathcal O_{X,x}.
  \]
\item[\textup{\SGAInlineRefTarget{sga3:VIIB:1.6.E:item:iii}(iii)}] For every \(x\in X\) projecting to
  \(s\in\operatorname{Spf}(k)\), there is an isomorphism
  \[
    \mathcal O_{X,x}\simeq k_s[T]/(F),
  \]
  where \(F\in k_s[T]\) is a separable monic polynomial
  \textup{(cf.\ \SGARefLink{sga3:VIIB:1.6.A}{1.6.A}~\SGARefLink{sga3:VIIB:1.6.E:item:ii}{(ii)})}.
\item[\textup{\SGAInlineRefTarget{sga3:VIIB:1.6.E:item:iv}(iv)}] \(X\) is topologically flat over \(k\), and, for
  every \(x\in X\) projecting to \(s\in\operatorname{Spf}(k)\),
  \[
    \mathcal O_{X,x}\widehat\otimes_k\kappa(s)
  \]
  is a finite separable extension of \(\kappa(s)\).
\item[\textup{\SGAInlineRefTarget{sga3:VIIB:1.6.E:item:v}(v)}] For every open ideal \(\mathfrak l\) of \(k\),
  each local component of
  \[
    A\widehat\otimes_k(k/\mathfrak l)
  \]
  is a finite étale algebra over the Artinian ring
  \(k/\mathfrak l\).
\end{enumerate}
(b) The category
\(\operatorname{Vaf}^{\mathrm{\acute{e}t}}/\kappa(k)\) identifies
with the category of étale schemes over \(S_{\kappa(k)}\)
\textup{(cf.\ \SGARefLink{sga3:VIIB:1.6.C}{1.6.C})}, and the functor \(X\mapsto X_\kappa\) induces
an equivalence
\[
  \operatorname{Vaf}^{\mathrm{\acute{e}t}}/k
  \xrightarrow{\ \sim\ }
  \operatorname{Vaf}^{\mathrm{\acute{e}t}}/\kappa(k)
\]
\textup{(cf.\ SGA~1, \SGARefLink{sga3:i:definition:6-1}{I 6.1})}.}

\medskip
\noindent\textit{Proof.}
Let \(I\) be the kernel of the multiplication morphism
\[
  m:A\widehat\otimes_kA\longrightarrow A.
\]
Suppose first that \(X\) is étale over \(k\), so that \(A\) is
formally étale over \(k\). By EGA \(0_{\mathrm{IV}}\), 20.7.4, the
separated completion of \(I/I^2\), for the quotient topology induced
from \(I\), is zero; hence
\[
  I=I^2.
\]
For every \(x\in X\), the ideal \(I\) is contained in the maximal
ideal \(\mathfrak m_{\Delta(x)}\) of
\(A\widehat\otimes_kA\). Thus its localization
\(I_{\mathfrak m_{\Delta(x)}}\) lies in the maximal ideal of
\(\mathcal O_{X\times X,\Delta(x)}\). Nakayama's \SGARefLink{sga3:VIIB:0.3}{lemma 0.3} therefore
gives
\[
  I_{\mathfrak m_{\Delta(x)}}=0.
\]
Consequently, \(\Delta:X\to X\times X\) is a local isomorphism.

Suppose next that \(\Delta\) is a local isomorphism and that \(k\) is
a field \(\kappa\). We show that every \(\mathcal O_{X,x}\) is a
finite-dimensional étale \(\kappa\)-algebra. Replacing \(X\) by
\(\operatorname{Spf}(\mathcal O_{X,x})\), we may assume that
\[
  A=\mathcal O_{X,x}
\]
is local. We proceed as in the proof of EGA
\(\mathrm{IV}_4\), 17.4.1, \textup{(b)}\(\Rightarrow\)\textup{(d'')}.
Let \(K\) be a finite normal extension of \(\kappa\) containing the
residue field \(\kappa(x)\), and put
\[
  B=A\otimes_\kappa K
   =A\widehat\otimes_\kappa K,
  \qquad
  X_K=\operatorname{Spf}(B)
      =X\widehat\otimes_\kappa K.
\]
The two tensor products agree because \([K:\kappa]<\infty\).

The morphism
\[
  \Delta_K:
  X_K\longrightarrow X_K\times_KX_K
\]
is again a local isomorphism. Hence, for every \(y\in X_K\),
multiplication induces an isomorphism
\[
  \left(
    B_y\widehat\otimes_KB_y
  \right)_{\mathfrak m_{\Delta_K(y)}}
  \xrightarrow{\ \sim\ }
  B_y.
\]
The residue field of \(B_y\) is \(K\); cf., for example,
\(\mathrm{VI}_A\), \SGARefLink{sga3:VIIA:1.1.1}{1.1.1}, editors' note~(11). Moreover,
\[
  C=B_y\widehat\otimes_KB_y
\]
is already local. Indeed,
\[
  \mathfrak n
  =
  \mathfrak m_y\widehat\otimes_KB_y
  +
  B_y\widehat\otimes_K\mathfrak m_y
\]
consists of topologically nilpotent elements and therefore lies in
the radical of \(C\), while
\[
  C/\mathfrak n
  =
  K\widehat\otimes_KK
  =
  K
\]
is a field. It follows that multiplication itself is an isomorphism
\[
  B_y\widehat\otimes_KB_y
  \xrightarrow{\ \sim\ }
  B_y.
\]
Choosing a pseudobasis of \(B_y\) over \(K\) that contains the unit
element \(1\), we deduce that \(B_y=K\). Since \(B\) is finite over
\(A\), the set \(X_K\) is finite. Thus
\[
  B=A\otimes_\kappa K
\]
is a product of finitely many copies of \(K\), and \(A\) is a finite
étale \(\kappa\)-algebra.

We have therefore shown that, over a field \(\kappa\), every profinite
\(\kappa\)-algebra \(A\) that is étale over \(\kappa\) is a product
of finite separable extensions \(K_i/\kappa\), endowed with the
product topology. Consequently,
\(\operatorname{Spf}(A)\) is the direct sum of the schemes
\[
  \operatorname{Spf}(K_i)=\operatorname{Spec}(K_i),
\]
and \(\operatorname{Vaf}^{\mathrm{\acute{e}t}}/\kappa\) identifies
with the category of étale \(\kappa\)-schemes.

The preceding argument proves
\textup{\SGARefLink{sga3:VIIB:1.6.E:item:ii}{(ii)}}\(\Rightarrow\)\textup{\SGARefLink{sga3:VIIB:1.6.E:item:iv}{(iv)}}, because
\(\mathcal O_{X,x}\widehat\otimes_k\kappa(s)\) is formally étale over
\(\kappa(s)\), and also proves part~\textup{(b)}. Indeed, return to
an arbitrary pseudocompact ring \(k\). The category
\(\operatorname{Vaf}^{\mathrm{\acute{e}t}}/\kappa(k)\) has just been
identified with the category of étale schemes over
\(S_{\kappa(k)}\). By \SGARefLink{sga3:VIIB:1.6.D}{1.6.D}, to show that
\[
  \operatorname{Vaf}^{\mathrm{\acute{e}t}}/k
  \longrightarrow
  \operatorname{Vaf}^{\mathrm{\acute{e}t}}/\kappa(k)
\]
is an equivalence, it remains only to prove essential surjectivity.

Let \(s\in\operatorname{Spf}(k)\), and let \(K/\kappa(s)\) be a
finite separable extension. Choose a primitive element
\(\xi\in K\), let \(n=[K:\kappa(s)]\), and choose a monic polynomial
\[
  F\in k_s[T]
\]
of degree \(n\) whose image
\(\overline F\in\kappa(s)[T]\) is the minimal polynomial of \(\xi\).
Since \(\overline F'\) is invertible in
\(\kappa(s)[T]/(\overline F)\), Nakayama's lemma shows that \(F'\) is
invertible in \(k_s[T]/(F)\). Thus \(F\) is separable, and \SGARefLink{sga3:VIIB:1.6.A}{1.6.A}~\SGARefLink{sga3:VIIB:1.6.E:item:ii}{(ii)}
shows that
\[
  A=k_s[T]/(F)
\]
is an étale \(k_s\)-algebra satisfying
\[
  A\widehat\otimes_k\kappa(s)\simeq K.
\]
It follows that every local profinite étale \(k_s\)-algebra is free
of finite rank over \(k_s\), and therefore is topologically free over
\(k\). Hence every profinite étale \(k\)-algebra is topologically
free over \(k\).

Condition~\textup{\SGARefLink{sga3:VIIB:1.6.E:item:v}{(v)}} implies condition~\textup{(d)} of \SGARefLink{sga3:VIIB:1.6.B}{1.6.B}, and
therefore implies~\textup{\SGARefLink{sga3:vii:paragraph:5-1:equation:i}{(i)}}. We have thus proved that
\textup{\SGARefLink{sga3:vii:paragraph:5-1:equation:i}{(i)}}, \textup{\SGARefLink{sga3:VIIB:1.6.E:item:iii}{(iii)}}, and \textup{\SGARefLink{sga3:VIIB:1.6.E:item:v}{(v)}} are equivalent and
imply~\textup{\SGARefLink{sga3:vii:paragraph:5-1:equation:ii}{(ii)}}, which in turn implies~\textup{\SGARefLink{sga3:VIIB:1.6.E:item:iv}{(iv)}}.

Finally, let \(A\) be a profinite \(k\)-algebra satisfying
\textup{\SGARefLink{sga3:VIIB:1.6.E:item:iv}{(iv)}}. We show that \(A\) is formally étale over \(k\). We
may suppose that \(A\) and \(k\) are local, and write \(\kappa\) for
the residue field of \(k\). By hypothesis,
\[
  K=A\widehat\otimes_k\kappa
\]
is a finite separable extension of \(\kappa\), say of degree \(n\).
By the preceding argument and \SGARefLink{sga3:VIIB:1.6.D}{Lemma~1.6.D}, there is a \(k\)-algebra
\(B\), free of rank \(n\) and formally étale over \(k\), together
with a continuous morphism
\[
  \phi:B\longrightarrow A
\]
such that
\[
  \phi\widehat\otimes_k\kappa
\]
is an isomorphism. Since \(A\) is topologically flat over \(k\), this
implies
\[
  \operatorname{Coker}(\phi)\widehat\otimes_k\kappa=0
  \qquad\text{and}\qquad
  \operatorname{Ker}(\phi)\widehat\otimes_k\kappa=0.
\]
Nakayama's \SGARefLink{sga3:VIIB:0.3}{lemma 0.3} gives
\[
  \operatorname{Coker}(\phi)=0
  =
  \operatorname{Ker}(\phi).
\]
Thus \(\phi\) is an isomorphism \textup{(cf.\ the proof of \SGARefLink{sga3:vii:proposition:0-2-b}{0.2.B})},
which completes the proof of \SGARefLink{sga3:VIIB:1.6.E}{1.6.E}.

% END INCLUDED SGA3_English_current_buildable_source_20260823/source/inputs/VII/tex/components/33_expose_VIIB_section16E_etale_criteria_en.tex

% BEGIN INCLUDED SGA3_English_current_buildable_source_20260823/source/inputs/VII/tex/components/34_expose_VIIB_section16F_etale_reflection_en.tex
% SGA 3, Expose VII_B, English reconstruction -- Loop 1.
% Scope: authority local pages 41--42, editor's note 73 and Proposition
% 1.6.F in full, stopping immediately before 1.6.1.
% Combined-reader pages: 554--555. Re-edition printed pages: 543--544.
% Original-edition marginal locator: 511.
% Authority: Polo--Gille Exp7B-13oct24.pdf,
% SHA-256 D303690E0195654BBE7C3DAA5D89EA23E24CD709A3575217234AC19AE8836694.
% Jacob Reinhold's CC BY 4.0 Markdown was consulted as comparison lineage.

Let \(X\) be a formal variety over \(k\). If \(x\in X\) projects to
\(s\in\operatorname{Spf}(k)\), then the residue field \(\kappa(x)\)
is a finite extension of \(\kappa(s)\). We denote by
\(\kappa_e(x)\) the separable closure of \(\kappa(s)\) in
\(\kappa(x)\).\footnote{\label{sga3:VIIB:nde:73}In what follows, the
editors used the preceding additions to give details of the
construction of the functor \(X\mapsto X_e\) and to prove that it
commutes with finite products; this is used in \SGARefLink{sga3:VIIB:2.5.1:corollary}{2.5.1}.}

\medskip
\noindent\textbf{Proposition 1.6.F.}
\phantomsection
\label{sga3:VIIB:1.6.F}
{\itshape \SGARefLink{sga3:vii:paragraph:5-1:equation:i}{(i)} The inclusion
\[
  \operatorname{Vaf}^{\mathrm{\acute{e}t}}/k
  \longrightarrow
  \operatorname{Vaf}/k
\]
has a left adjoint \(X\mapsto X_e\). The formal variety \(X_e\) has
the same points as \(X\). For each \(x\in X\) projecting to
\(s\in\operatorname{Spf}(k)\), choose a primitive element
\(\xi\) of \(\kappa_e(x)/\kappa(s)\), let \(n\) be its degree, choose
any lift \(u\in\mathcal O_{X,x}\) of \(\xi\), and let
\(F\in k[T]\) be a monic polynomial of degree \(n\) that annihilates
\(u\). Then put
\[
  \mathcal O_{X_e,x}=k[T]/(F).
\]
For every
\(Y\in\operatorname{Ob}(\operatorname{Vaf}^{\mathrm{\acute{e}t}}/k)\),
the canonical maps in the following diagram are bijective:
\begin{center}
  \phantomsection
  \label{sga3:VIIB:1.6.F:hom-square}
  % SGA3-VII-DIAG-095
  % BEGIN INCLUDED SGA3_English_current_buildable_source_20260823/source/inputs/VII/native_diagrams/exp7/viib_15_95/diag_095.tex
% SGA3-VII-DIAG-095
% Authority: Exp7B-13oct24.pdf, local p. 42.
% Crop: figures/Exp7B_localp042_etale_reflection_hom_square.png
% Locator: tex/components/34_expose_VIIB_section16F_etale_reflection_en.tex, Loop 1 line 45.
\begin{tikzcd}[row sep=2.8em,column sep=3.5em]
\operatorname{Hom}_{\operatorname{Vaf}/k}(X,Y)
  \arrow[r,"{\sim}"]
  & \operatorname{Hom}_{\operatorname{Vaf}/\kappa(k)}
      (X_\kappa,Y_\kappa)
      \arrow[d,"\wr"] \\
\operatorname{Hom}_{\operatorname{Vaf}/k}(X_e,Y)
  \arrow[r,"{\sim}"]
  & \operatorname{Hom}_{\operatorname{Vaf}/\kappa(k)}
      ((X_e)_\kappa,Y_\kappa)\, ,
\end{tikzcd}

% END INCLUDED SGA3_English_current_buildable_source_20260823/source/inputs/VII/native_diagrams/exp7/viib_15_95/diag_095.tex

\end{center}
The vertical arrow is induced by the inclusions
\[
  \kappa_e(x)\lhook\joinrel\longrightarrow\kappa(x),
  \qquad x\in X.
\]
In particular, this defines a morphism
\[
  p:X\longrightarrow X_e,
\]
and every morphism of formal \(k\)-varieties
\(\phi:X\to Y\), with \(Y\) étale over \(k\), factors uniquely
through \(p\).

\SGARefLink{sga3:vii:paragraph:5-1:equation:ii}{(ii)} The functor \(X\mapsto X_e\) commutes with finite products.}

\medskip
\noindent\textit{Proof.}
Assertion~\textup{\SGARefLink{sga3:vii:paragraph:5-1:equation:i}{(i)}} follows from \SGARefLink{sga3:VIIB:1.6.E}{1.6.E}~\textup{(b)} and \SGARefLink{sga3:VIIB:1.6.D}{1.6.D}.
We prove~\textup{\SGARefLink{sga3:vii:paragraph:5-1:equation:ii}{(ii)}}. By the equivalence of categories in
\SGARefLink{sga3:VIIB:1.6.E}{1.6.E}~\textup{(b)}, we may suppose that \(k\) is a field. In that
case, if \(X\) is a semilocal formal \(k\)-variety, so that its
affine algebra \(A\) is semilocal, the affine algebra of \(X_e\) is
the largest subalgebra of \(A\) that is étale over \(k\); denote it
by \(A_e\).

It is therefore enough to show that, if \(K\) and \(L\) are two
finite extensions of \(k\), then
\[
  K_e\otimes_kL_e
  =
  (K\otimes_kL)_e.
\]
Let \(p\) be the characteristic exponent of \(k\), so that \(p=1\)
if \(\operatorname{char}(k)=0\), and
\(p=\operatorname{char}(k)\) otherwise. For every
\(x\in K\otimes_kL\), there is an integer \(n\geq1\) such that
\[
  x^{p^n}\in K_e\otimes_kL_e.
\]
Thus \(K\otimes_kL\) is radicial over \(K_e\otimes_kL_e\), and its
largest étale \(k\)-subalgebra is precisely
\(K_e\otimes_kL_e\). Hence
\[
  K_e\otimes_kL_e=(K\otimes_kL)_e.
\]

Keeping this notation, observe that
\(\mathcal O_{X_e,x}\) is not necessarily a subalgebra of
\(\mathcal O_{X,x}\). It is a subalgebra when \(X\) is
topologically flat over \(k\), by the proposition that follows.

% END INCLUDED SGA3_English_current_buildable_source_20260823/source/inputs/VII/tex/components/34_expose_VIIB_section16F_etale_reflection_en.tex

% BEGIN INCLUDED SGA3_English_current_buildable_source_20260823/source/inputs/VII/tex/components/35_expose_VIIB_sections161_162_flat_etale_reflection_en.tex
% SGA 3, Expose VII_B, English reconstruction -- Loop 1.
% Scope: authority local pages 42--43, Section 1.6.1, its concluding
% proposition, and Remark 1.6.2 in full, stopping immediately before
% Section 2 on local page 43.
% Combined-reader pages: 555--556. Re-edition printed pages: 544--545.
% Original-edition marginal locator within this span: 511.
% Editorial notes in scope: 74--75.
% Authority: Polo--Gille Exp7B-13oct24.pdf,
% SHA-256 D303690E0195654BBE7C3DAA5D89EA23E24CD709A3575217234AC19AE8836694.
% Jacob Reinhold's CC BY 4.0 Markdown was consulted as comparison lineage.
% The three genuine diagrams are retained as exact PDF-derived Loop 1
% source crops; all ordinary mathematics is native LaTeX.

\subsubsection*{1.6.1}
\phantomsection
\label{sga3:VIIB:1.6.1}

Let \(Y\) be an étale formal \(k\)-variety and let
\(f:X\to Y\) be a morphism in \(\operatorname{Vaf}/k\). Then the
following square is Cartesian, where \(\Gamma_f\) is the graph
morphism \(X\to X\times Y\), with components
\(\operatorname{id}_X\) and \(f\):

\begin{center}
  \phantomsection
  \label{sga3:VIIB:1.6.1:cartesian-square}
  % SGA3-VII-DIAG-096
  % BEGIN INCLUDED SGA3_English_current_buildable_source_20260823/source/inputs/VII/native_diagrams/exp7/viib_15_95/diag_096.tex
% SGA3-VII-DIAG-096
% Authority: Exp7B-13oct24.pdf, local p. 42.
% Crop: figures/Exp7B_localp042_section161_cartesian_square.png
% Locator: tex/components/35_expose_VIIB_sections161_162_flat_etale_reflection_en.tex, Loop 1 line 27.
\begin{tikzcd}[row sep=2.8em,column sep=3.7em]
X \arrow[r,"\Gamma_f"] \arrow[d,"f"']
  & X\times Y \arrow[d,"{f\times\operatorname{id}_Y}"] \\
Y \arrow[r,"\Delta_Y"]
  & Y\times Y\, .
\end{tikzcd}

% END INCLUDED SGA3_English_current_buildable_source_20260823/source/inputs/VII/native_diagrams/exp7/viib_15_95/diag_096.tex

\end{center}

It follows that \(\Gamma_f\) is a local isomorphism, and hence that
\[
  f=\operatorname{pr}_Y\circ\Gamma_f
\]
is topologically flat if \(\operatorname{pr}_Y\) is; this holds, for
example, if \(X\) is topologically flat over \(k\).

Conversely, since \(Y\to k\) is topologically flat, \(X\to k\) will
also be topologically flat if \(f\) is \textup{(cf.\ \SGARefLink{sga3:VIIB:1.3.3}{1.3.3})}. Taking
in particular for \(f\) the canonical morphism
\[
  p:X\longrightarrow X_e
\]
of \SGARefLink{sga3:VIIB:1.6}{1.6}, we obtain the following.

\medskip
\noindent\textbf{Proposition.}
\phantomsection
\label{sga3:VIIB:1.6.1.prop}
{\itshape Let \(X\) be a formal variety over \(k\). The morphism
\(X\to X_e\) is topologically flat if and only if \(X\) is
topologically flat over \(k\).}

% Authority local p. 43; combined-reader p. 556; printed p. 545.
\medskip
\noindent\textbf{Remark 1.6.2.}\footnote{%
\label{sga3:VIIB:nde:74}The editors inserted this remark here; in the
original it appeared in \SGARefLink{sga3:VIIB:2.5.2:proposition}{2.5.2}.}
\phantomsection
\label{sga3:VIIB:1.6.2}
When \(k\) is a \emph{perfect} field, the functor
\[
  X\longmapsto X_e
\]
is also right adjoint to the inclusion
\[
  \operatorname{Vaf}^{\mathrm{\acute{e}t}}/k
  \lhook\joinrel\longrightarrow
  \operatorname{Vaf}/k.
\]
Indeed, in this case
\[
  \mathcal O_{X_e,x}=\kappa(x)
\]
for every \(x\in X\), and the canonical projections
\[
  \mathcal O_{X,x}\longrightarrow\kappa(x)
\]
define a morphism \(s:X_e\to X\), which is a section of
\(p:X\to X_e\). For every morphism \(f:X\to Y\), the following
diagrams are commutative:

\begin{center}
  \phantomsection
  \label{sga3:VIIB:1.6.2:local-ring-square}
  % SGA3-VII-DIAG-097
  % BEGIN INCLUDED SGA3_English_current_buildable_source_20260823/source/inputs/VII/native_diagrams/exp7/viib_15_95/diag_097.tex
% SGA3-VII-DIAG-097
% Authority: Exp7B-13oct24.pdf, local p. 43.
% Crop: figures/Exp7B_localp043_remark162_local_ring_square.png
% Locator: tex/components/35_expose_VIIB_sections161_162_flat_etale_reflection_en.tex, Loop 1 line 86.
\begin{tikzcd}[row sep=2.6em,column sep=3.5em]
\mathcal O_{Y,f(x)} \arrow[r] \arrow[d]
  & \mathcal O_{X,x} \arrow[d] \\
\kappa(f(x)) \arrow[r]
  & \kappa(x)
\end{tikzcd}

% END INCLUDED SGA3_English_current_buildable_source_20260823/source/inputs/VII/native_diagrams/exp7/viib_15_95/diag_097.tex

  \hspace{.07\textwidth}
  \phantomsection
  \label{sga3:VIIB:1.6.2:functoriality-square}
  % SGA3-VII-DIAG-098
  % BEGIN INCLUDED SGA3_English_current_buildable_source_20260823/source/inputs/VII/native_diagrams/exp7/viib_15_95/diag_098.tex
% SGA3-VII-DIAG-098
% Authority: Exp7B-13oct24.pdf, local p. 43.
% Crop: figures/Exp7B_localp043_remark162_functoriality_square.png
% Locator: tex/components/35_expose_VIIB_sections161_162_flat_etale_reflection_en.tex, Loop 1 line 91.
\begin{tikzcd}[row sep=2.6em,column sep=3.2em]
X \arrow[r,"f"] & Y \\
X_e \arrow[r,"f_e"] \arrow[u,"s_X"]
  & Y_e \arrow[u,"s_Y"']
\end{tikzcd}

% END INCLUDED SGA3_English_current_buildable_source_20260823/source/inputs/VII/native_diagrams/exp7/viib_15_95/diag_098.tex

\end{center}

Thus \(s\) is functorial in \(X\), and it follows that
\(X\mapsto X_e\) is right adjoint to the inclusion
\[
  \operatorname{Vaf}^{\mathrm{\acute{e}t}}/k
  \lhook\joinrel\longrightarrow
  \operatorname{Vaf}/k.
\]
Consequently, being a right adjoint, the functor \(X\mapsto X_e\)
commutes with projective limits whenever they exist in
\(\operatorname{Vaf}^{\mathrm{\acute{e}t}}/k\),\footnote{%
\label{sga3:VIIB:nde:75}This is the case for finite projective
limits.} and hence, in particular, with finite products.

\noindent
(This can also be verified directly: for every formal
\(k\)-variety \(X\), with affine algebra \(A\), the affine algebra of
\(X_e\) is the quotient of \(A\) by its radical \(r(A)\); and, since
\(k\) is perfect, the quotient of
\(\mathcal O_{X,x}\widehat{\otimes}\mathcal O_{Y,y}\) by its radical
is the algebra \(\kappa(x)\otimes_k\kappa(y)\), because the latter is
semisimple.)

% Next cursor: Exp7B local p. 43, combined-reader p. 556,
% re-edition printed p. 545, immediately before Section 2.

% END INCLUDED SGA3_English_current_buildable_source_20260823/source/inputs/VII/tex/components/35_expose_VIIB_sections161_162_flat_etale_reflection_en.tex

% BEGIN INCLUDED SGA3_English_current_buildable_source_20260823/source/inputs/VII/tex/components/36_expose_VIIB_section21_formal_group_cogroups_en.tex
% SGA 3, Expose VII_B, English reconstruction -- Loop 1.
% Scope: authority local pages 43--44, Section 2 heading and 2.1 in full,
% stopping immediately before 2.2.
% Combined-reader pages: 556--557. Re-edition printed pages: 545--546.
% Original-edition marginal locators: 512--513.
% Authority: Polo--Gille Exp7B-13oct24.pdf,
% SHA-256 D303690E0195654BBE7C3DAA5D89EA23E24CD709A3575217234AC19AE8836694.
% Jacob Reinhold's CC BY 4.0 Markdown was consulted as comparison lineage.

\section*{2. Generalities on Formal Groups}
\phantomsection
\label{sga3:VIIB:2}
\addcontentsline{toc}{section}{2. Generalities on Formal Groups}

\subsection*{2.1}
\phantomsection
\label{sga3:VIIB:2.1}
\addcontentsline{toc}{subsection}{2.1. Formal groups and cogroups}

Let \(k\) be a pseudocompact ring and let \(G\) be a formal
\(k\)-group, that is, a group object in the category
\(\operatorname{Vaf}/k\) of formal varieties over \(k\). Let \(A\)
be the affine algebra of \(G\). The composition law of \(G\) plainly
defines a diagonal morphism, namely a homomorphism of profinite
\(k\)-algebras
\[
  \Delta_A:A\longrightarrow A\widehat\otimes_kA.
\]
This homomorphism has the following properties.

\begin{enumerate}
\item[\textup{\SGAInlineRefTarget{sga3:VIIB:2.1:item:i}(i)}] The diagram
\begin{center}
  \phantomsection
  \label{sga3:VIIB:2.1:coassociativity-square}
  % SGA3-VII-DIAG-099
  % BEGIN INCLUDED SGA3_English_current_buildable_source_20260823/source/inputs/VII/native_diagrams/exp7/viib_15_95/diag_099.tex
% SGA3-VII-DIAG-099
% Authority: Exp7B-13oct24.pdf, local p. 43.
% Crop: figures/Exp7B_localp043_coassociativity_square.png
% Locator: tex/components/36_expose_VIIB_section21_formal_group_cogroups_en.tex, Loop 1 line 36.
\begin{tikzcd}[row sep=3em,column sep=4.7em]
A \arrow[r,"\Delta_A"] \arrow[d,"\Delta_A"']
  & A\widehat\otimes_k A
      \arrow[d,"{\Delta_A\widehat\otimes\operatorname{id}_A}"] \\
A\widehat\otimes_k A
  \arrow[r,"{\operatorname{id}_A\widehat\otimes\Delta_A}"]
  & A\widehat\otimes_k A\widehat\otimes_k A
\end{tikzcd}

% END INCLUDED SGA3_English_current_buildable_source_20260823/source/inputs/VII/native_diagrams/exp7/viib_15_95/diag_099.tex

\end{center}
is commutative.

\item[\textup{(ii)}] There is a necessarily unique
\emph{augmentation}, that is, a homomorphism of profinite
\(k\)-algebras
\[
  \varepsilon_A:A\longrightarrow k,
\]
such that the two composites displayed below are the identity maps of
\(A\):
\begin{center}
  \phantomsection
  \label{sga3:VIIB:2.1:counit-compositions}
  % SGA3-VII-DIAG-100
  % BEGIN INCLUDED SGA3_English_current_buildable_source_20260823/source/inputs/VII/native_diagrams/exp7/viib_15_95/diag_100.tex
% SGA3-VII-DIAG-100
% Authority: Exp7B-13oct24.pdf, local p. 43.
% Crop: figures/Exp7B_localp043_counit_compositions.png
% Locator: tex/components/36_expose_VIIB_section21_formal_group_cogroups_en.tex, Loop 1 line 52.
\(
\begin{gathered}
A\xrightarrow{\ \Delta_A\ }A\widehat\otimes_k A
 \xrightarrow{\ \varepsilon_A\widehat\otimes\operatorname{id}_A\ }
 k\widehat\otimes_k A\simeq A
\\[.7ex]
\text{and}\qquad
A\xrightarrow{\ \Delta_A\ }A\widehat\otimes_k A
 \xrightarrow{\ \operatorname{id}_A\widehat\otimes\varepsilon_A\ }
 A\widehat\otimes_k k\simeq A
\end{gathered}
\)

% END INCLUDED SGA3_English_current_buildable_source_20260823/source/inputs/VII/native_diagrams/exp7/viib_15_95/diag_100.tex

\end{center}

\item[\textup{\SGAInlineRefTarget{sga3:VIIB:2.1:item:iii}(iii)}] There is a necessarily unique \emph{antipode},
that is, a homomorphism of profinite \(k\)-algebras
\[
  c_A:A\longrightarrow A,
\]
such that the composite
\begin{center}
  \phantomsection
  \label{sga3:VIIB:2.1:antipode-composition}
  % SGA3-VII-DIAG-101
  % BEGIN INCLUDED SGA3_English_current_buildable_source_20260823/source/inputs/VII/native_diagrams/exp7/viib_15_95/diag_101.tex
% SGA3-VII-DIAG-101
% Authority: Exp7B-13oct24.pdf, local p. 44.
% Crop: figures/Exp7B_localp044_antipode_composition.png
% Locator: tex/components/36_expose_VIIB_section21_formal_group_cogroups_en.tex, Loop 1 line 65.
\(
A\xrightarrow{\ \Delta_A\ }A\widehat\otimes_k A
\xrightarrow{\ c_A\widehat\otimes\operatorname{id}_A\ }
A\widehat\otimes_k A
\xrightarrow{\ m_A\ }A
\)

% END INCLUDED SGA3_English_current_buildable_source_20260823/source/inputs/VII/native_diagrams/exp7/viib_15_95/diag_101.tex

\end{center}
equals \(\eta_A\circ\varepsilon_A\). Here \(m_A\) is the linear map
that sends \(a\widehat\otimes b\) to \(ab\), and
\(\eta_A:k\to A\) sends \(\lambda\) to \(\lambda1_A\).
\end{enumerate}

Conversely, data
\[
  (\Delta_A,\varepsilon_A,c_A)
\]
satisfying \textup{\SGARefLink{sga3:VIIB:2.1:item:i}{(i)}}--\textup{\SGARefLink{sga3:VIIB:2.1:item:iii}{(iii)}} endow \(G\) with the
structure of a formal \(k\)-group.\footnote{
\label{sga3:VIIB:nde:76}One also says that \(A\) is a \emph{cogroup}
in the category of profinite \(k\)-algebras. The editors expanded the
material that follows, in particular by making explicit what a
morphism of formal groups \(K\to G\) means; cf.\ \SGARefLink{sga3:VIIB:2.3.1}{Proposition~2.3.1}.}
Explicitly, for every profinite \(k\)-algebra \(B\), the set
\[
  \operatorname{Hom}_c(A,B)
\]
of continuous morphisms of \(k\)-modules \(\phi:A\to B\) carries a
group structure, functorial in \(B\), defined by
\[
  \phi\cdot\phi'
  =
  m_B\circ
  (\phi\widehat\otimes\phi')
  \circ\Delta_A.
\]
The identity element is \(\eta_B\circ\varepsilon_A\), where \(m_B\)
is the multiplication of \(B\) and
\(\eta_B:k\to B\) sends \(\lambda\) to \(\lambda1_B\); the inverse
of \(\phi\) is \(\phi\circ c_A\). The set
\[
  \operatorname{Hom}_{\operatorname{Alp}/k}(A,B)
\]
of continuous morphisms of \(k\)-algebras \(A\to B\) is a subgroup,
because \(B\) is commutative.

\medskip
\noindent\textbf{Definition.}
\phantomsection
\label{sga3:VIIB:2.1.definition}
A \emph{morphism of formal \(k\)-groups}
\(\theta:K\to G\) is, by definition, a morphism of formal
\(k\)-varieties that respects the group structures. Let \(B\),
respectively \(A\), be the affine algebra of \(K\), respectively
\(G\), and let \(f:A\to B\) be the morphism corresponding to
\(\theta\). The condition is equivalently that \(f\) be compatible
with comultiplication:
\[
  (f\widehat\otimes f)\circ\Delta_A
  =
  \Delta_B\circ f.
\]
The identities
\[
  \varepsilon_B\circ f=\varepsilon_A,
  \qquad
  c_B\circ f=f\circ c_A
\]
then follow automatically. We write \(\operatorname{Grf}/k\) for
the category of formal \(k\)-groups.

\medskip
\noindent\textbf{Notation.}
\phantomsection
\label{sga3:VIIB:2.1.notation}
In what follows, the ideal
\[
  I_A=\operatorname{Ker}(\varepsilon_A)
\]
will be called the \emph{augmentation ideal} of \(A\), and we write
\[
  \omega_{G/k}=I_A/\overline{I_A^2}
\]
for the pseudocompact \(k\)-module obtained by quotienting \(I_A\)
by the closed ideal generated by the products \(xy\), with
\(x,y\in I_A\).

% END INCLUDED SGA3_English_current_buildable_source_20260823/source/inputs/VII/tex/components/36_expose_VIIB_section21_formal_group_cogroups_en.tex

% BEGIN INCLUDED SGA3_English_current_buildable_source_20260823/source/inputs/VII/tex/components/37_expose_VIIB_section22_bialgebras_en.tex
% SGA 3, Expose VII_B, English reconstruction -- Loop 1.
% Scope: authority local pages 44--45, Section 2.2 through the paragraph
% preceding 2.2.1.
% Combined-reader pages: 557--558. Re-edition printed pages: 546--547.
% Original-edition marginal locators: 513--514.
% Authority: Polo--Gille Exp7B-13oct24.pdf,
% SHA-256 D303690E0195654BBE7C3DAA5D89EA23E24CD709A3575217234AC19AE8836694.
% Jacob Reinhold's CC BY 4.0 Markdown was consulted as comparison lineage.
% No genuine diagram occurs in this span; all mathematics is native LaTeX.

\subsection*{2.2}
\phantomsection
\label{sga3:VIIB:2.2}
\addcontentsline{toc}{subsection}{\SGARefLink{sga3:VIIA:2.2}{2.2}. Bialgebras over \(\mathbf O_k\)}

Let \(\mathbf H\) be a group object in the category of coalgebras
over \(\mathbf O_k\). Thus, for every object \(C\) of
\(\operatorname{Alf}/k\), the object \(\mathbf H(C)\) carries the
structure of a \(C\)-coalgebra in groups \textup{(cf.\
VII\textsubscript{A} \SGARefLink{sga3:VIIA:3.2}{3.2})}. Following Manin, we shall say
\emph{bialgebra}\footnote{\label{sga3:VIIB:nde:77}The term
``bigebra'' is also used \textup{(cf.\ [BAlg], III \S~11.4)}. Recall
\textup{(cf.\ VII\textsubscript{A}, \SGARefLink{sga3:VIIB:3.1:lemma}{3.1}, editors' note~(26))} that
all bialgebras considered here are assumed to be cocommutative and
equipped with an antipode; they are therefore cocommutative Hopf
algebras.} instead of ``coalgebra in groups.'' Moreover, for every
morphism
\[
  \phi:C\longrightarrow D
\]
in \(\operatorname{Alf}/k\), the map
\[
  D\otimes_C\mathbf H(C)\longrightarrow\mathbf H(D)
\]
is a homomorphism of \(D\)-bialgebras.

\medskip
\noindent\textbf{Definition.}
\phantomsection
\label{sga3:VIIB:2.2.definition}
These properties will be summarized by saying that \(\mathbf H\) is
a \emph{bialgebra over \(\mathbf O_k\)}.

The functor
\[
  \mathbf H\longmapsto\operatorname{Spf}^*(\mathbf H)
\]
of \SGARefLink{sga3:VIIB:1.3.5}{1.3.5} plainly commutes with finite products. It therefore sends a
bialgebra over \(\mathbf O_k\) to a \(k\)-group functor, that is, a
covariant functor from \(\operatorname{Alf}/k\) to the category of
groups.

Indeed, for every \(k\)-algebra \(C\) of finite length, the elements
of
\[
\begin{aligned}
  \operatorname{Spf}^*(\mathbf H)(C)
  &=
  \operatorname{Spf}^*(\mathbf H(C))
  \\
  &=
  \bigl\{
    x\in\mathbf H(C)
    \mid
    \varepsilon(x)=1
    \ \text{and}\
    \Delta(x)=x\otimes x
  \bigr\}
\end{aligned}
\]
form a group under multiplication in the algebra \(\mathbf H(C)\);
cf.\ VII\textsubscript{A} \SGARefLink{sga3:VIIA:3.2.2}{3.2.2}. Notice also that
\(\Delta(x)=x\otimes x\) implies
\[
  \varepsilon(x)=\varepsilon(x)^2.
\]
Thus \(\varepsilon(x)=1\) whenever \(C\) is local and
\(x\neq0\).\footnote{\label{sga3:VIIB:nde:78}Since
\(\Delta(x)=x\otimes x\), one has
\(x=\varepsilon(x)x\); hence \(\varepsilon(x)=0\) can occur only
when \(x=0\).}

% END INCLUDED SGA3_English_current_buildable_source_20260823/source/inputs/VII/tex/components/37_expose_VIIB_section22_bialgebras_en.tex

% BEGIN INCLUDED SGA3_English_current_buildable_source_20260823/source/inputs/VII/tex/components/38_expose_VIIB_section221_flat_bialgebra_equivalence_en.tex
% SGA 3, Expose VII_B, English reconstruction -- Loop 1.
% Scope: authority local pages 45--47, Sections 2.2.1 and 2.2.2 in full,
% stopping immediately before Section 2.3.
% Combined-reader pages: 558--560. Re-edition printed pages: 547--549.
% Original-edition marginal locators: 514--515.
% Editorial notes in scope: 79--84.
% Authority: Polo--Gille Exp7B-13oct24.pdf,
% SHA-256 D303690E0195654BBE7C3DAA5D89EA23E24CD709A3575217234AC19AE8836694.
% Jacob Reinhold's CC BY 4.0 Markdown was consulted as comparison lineage.
% No genuine diagram occurs in this span; all mathematics is native LaTeX.

\subsubsection*{2.2.1}
\phantomsection
\label{sga3:VIIB:2.2.1}

A bialgebra \(\mathbf H\) over \(\mathbf O_k\) is said to be
\emph{flat} if its underlying module is flat
\textup{(cf.\ \SGARefLink{sga3:VIIB:1.2.1}{1.2.1})}.\footnote{\label{sga3:VIIB:nde:79}The editors
have expanded the original in what follows.} If \(\mathbf H\) is
flat, then, by \SGARefLink{sga3:VIIB:1.3.5}{1.3.5},
\[
  A=\Gamma^*(\mathbf H)
\]
is a topologically flat profinite \(k\)-algebra, and
\(\operatorname{Spf}^*(\mathbf H)\), as a functor from
\[
  (\operatorname{Alf}/k)^{\mathrm{op}}
  =
  \operatorname{Vaf}/k
\]
to the category of sets, is isomorphic to the functor
\[
  \operatorname{Spf}(A):
  C\longmapsto
  \operatorname{Hom}_{\operatorname{Vaf}/k}
    \bigl(\operatorname{Spf}(C),\operatorname{Spf}(A)\bigr).
\]
The group structure on \(\operatorname{Spf}^*(\mathbf H)\) therefore
endows
\[
  \mathscr G(\mathbf H)=\operatorname{Spf}(A)
\]
with the structure of a formal group, described explicitly as
follows.

For every object \(C\) of \(\operatorname{Alf}/k\), the \(C\)-module
underlying \(\mathbf H(C)\) is projective. \SGARefLink{sga3:VIIB:1.2.3.A}{Lemma 1.2.3.A} therefore
gives, by induction on \(n\), natural isomorphisms
\[
\begin{aligned}
  \bigl(\mathbf H(C)^*\bigr)^{\widehat\otimes(n+1)}
  &\simeq
  \operatorname{Hom}_C
    \bigl(
      \mathbf H(C),
      \bigl(\mathbf H(C)^*\bigr)^{\widehat\otimes n}
    \bigr)
  \\
  &\simeq
  \operatorname{Hom}_C
    \bigl(
      \mathbf H(C),
      \bigl(\mathbf H(C)^{\otimes n}\bigr)^*
    \bigr)
  \\
  &\simeq
  \bigl(\mathbf H(C)^{\otimes(n+1)}\bigr)^*.
\end{aligned}
\]
For \(n=1,2\), it follows that the \(C\)-algebra structure on
\(\mathbf H(C)\) endows \(\mathbf H(C)^*\) with a diagonal map
satisfying conditions \SGARefLink{sga3:VIIB:2.1}{2.1}\textup{\SGARefLink{sga3:vii:paragraph:5-1:equation:i}{(i)}}--\textup{\SGARefLink{sga3:VIIB:2.1:item:iii}{(iii)}}, functorially
in \(C\).

Consequently,
\[
  A
  =
  \Gamma^*(\mathbf H)
  =
  \varprojlim_{\mathfrak l}\mathbf H(k/\mathfrak l)^*
\]
is endowed with the structure of a cogroup in
\(\operatorname{Alp}/k\), which defines the announced formal-group
structure on \(\mathscr G(\mathbf H)\).

Conversely, let \(G\) be a topologically flat formal \(k\)-group with
affine algebra \(A\), and write
\[
  \mathbf H(G)=\mathbf V_k^f(A)
\]
for its \(\mathbf O_k\)-coalgebra \textup{(cf.\ \SGARefLink{sga3:VIIB:1.2.3}{1.2.3})}. The
diagonal morphism
\[
  \Delta_A:A\longrightarrow A\widehat\otimes_kA
\]
then induces, for every finite-length \(k\)-algebra \(C\), a
\(C\)-linear map
\[
  \mathbf V_k^f(A)(C)
  \otimes_C
  \mathbf V_k^f(A)(C)
  \longrightarrow
  \mathbf V_k^f(A)(C),
\]
which makes the coalgebra \(\mathbf V_k^f(A)(C)\) into a
\(C\)-bialgebra. We shall call \(\mathbf H(G)\) the
\emph{covariant bialgebra} of the formal group \(G\).\footnote{%
\label{sga3:VIIB:nde:80}The editors added the adjective
``covariant'' to make clear that the functor
\(G\mapsto\mathbf H(G)\) is covariant; this terminology is used in
[Di73], I \S~2.14.} Thus, by \SGARefLink{sga3:VIIB:1.3.5.D}{Proposition 1.3.5.D}, we obtain the
following.

\medskip
\noindent\textbf{Proposition.}
\phantomsection
\label{sga3:VIIB:2.2.1.prop}
{\itshape
\begin{enumerate}
\item[\textup{(i)}]
The functors
\[
  G\longmapsto\mathbf H(G),
  \qquad
  \mathbf H\longmapsto\mathscr G(\mathbf H)
\]
define an equivalence between the category of topologically flat
formal \(k\)-groups and the category of flat
\(\mathbf O_k\)-bialgebras.\footnote{\label{sga3:VIIB:nde:81}Recall
that in this Exposé, ``bialgebra'' means ``cocommutative Hopf
algebra.''}

\item[\textup{(ii)}]
This equivalence ``commutes with base change'': if \(k\to K\) is a
morphism of pseudocompact rings, then
\[
  \mathbf H\bigl(G\widehat\otimes_kK\bigr)
  =
  \mathbf H(G)\otimes_kK
\]
and
\[
  \mathscr G\bigl(\mathbf H\otimes_kK\bigr)
  =
  \mathscr G(\mathbf H)\widehat\otimes_kK.
\]
\end{enumerate}}

When \(k\) is an Artinian ring and \(G\) is a topologically flat
formal \(k\)-group, the functor \(\mathbf H(G)\) is plainly
determined by its value
\[
  \mathbf H(G)=\mathbf H(G)(k)
\]
at \(k\). We shall also call \(\mathbf H(G)\) the \emph{(covariant)
bialgebra} of \(G\).\footnote{\label{sga3:VIIB:nde:82}The editors
have introduced here the notation
\(\mathcal H^{\mathrm{cocom}}_{\mathrm{flat}/k}\), which will be
useful below.} Let \(\mathcal H_k\) denote the category of
\(k\)-Hopf algebras, and let
\[
  \mathcal H^{\mathrm{cocom}}_{\mathrm{flat}/k}
\]
be its full subcategory formed by the \(k\)-Hopf algebras that are
flat over \(k\) and cocommutative. We then obtain the following.

\medskip
\noindent\textbf{Corollary.}
\phantomsection
\label{sga3:VIIB:2.2.1.cor}
{\itshape Let \(k\) be an Artinian ring.
\begin{enumerate}
\item[\textup{(i)}]
The functors
\[
  G\longmapsto\mathbf H(G),
  \qquad
  \mathbf H\longmapsto
  \mathscr G(\mathbf H)=\operatorname{Spf}^*(\mathbf H)
\]
define an equivalence between the category of topologically flat
formal \(k\)-groups and
\(\mathcal H^{\mathrm{cocom}}_{\mathrm{flat}/k}\).

\item[\textup{(ii)}]
This equivalence ``commutes with base change'': if \(k\to K\) is a
morphism of Artinian rings, then
\[
  \mathbf H\bigl(G\widehat\otimes_kK\bigr)
  =
  \mathbf H(G)\otimes_kK
\]
and
\[
  \mathscr G\bigl(\mathbf H\otimes_kK\bigr)
  =
  \mathscr G(\mathbf H)\widehat\otimes_kK.
\]
\end{enumerate}}

On the other hand, let
\[
  \mathcal H^{\mathrm{com}}_{\mathrm{flat}/k}
\]
denote the full subcategory of \(\mathcal H_k\) formed by the
\(k\)-Hopf algebras that are flat over \(k\) and commutative. Recall
that the functor
\[
  T\longmapsto\mathcal O(T)
\]
is an anti-equivalence from the category of affine \(k\)-group
schemes onto \(\mathcal H^{\mathrm{com}}_{\mathrm{flat}/k}\).

\subsubsection*{2.2.2}
\phantomsection
\label{sga3:VIIB:2.2.2}

Suppose, for simplicity, that \(k\) is Artinian, and let \(G\) be a
topologically flat formal \(k\)-group.\footnote{%
\label{sga3:VIIB:nde:83}The editors have expanded the original in
what follows in order to introduce the notation \(D'(G)\) and
\(D(K)\); cf.\ [Ca62], \S~14.} Then \(G\) is commutative if and only
if its affine algebra \(\mathcal A(G)\) has a cocommutative
comultiplication, which is equivalent to saying that the bialgebra
\(\mathbf H(G)\) has a commutative multiplication. In that case,
\(\mathbf H(G)\) is a commutative and cocommutative Hopf algebra,
flat over \(k\). Thus, if
\[
  D'(G)=\operatorname{Spec}\mathbf H(G),
\]
then \(D'(G)\) is a commutative \(k\)-group scheme, affine and flat
over \(k\).

Conversely, if \(T\) is such a \(k\)-group scheme, its affine algebra
\(\mathcal O(T)\) is a commutative group object in the category of
cocommutative \(k\)-coalgebras flat over \(k\). Hence \SGARefLink{sga3:VIIB:1.3.5.C}{1.3.5.C} gives
a topologically flat formal \(k\)-group \(D(T)\), defined by
\[
  D(T)
  =
  \operatorname{Spf}^*\mathcal O(T)
  =
  \operatorname{Spf}(\mathcal A),
  \qquad
  \mathcal A=\mathcal O(T)^*.
\]
By \SGARefLink{sga3:VIIB:1.3.5.D}{1.3.5.D}, there are canonical isomorphisms
\[
  G=\operatorname{Spf}^*\mathbf H(G),
  \qquad
  \mathbf H(D(T))=\mathcal O(T).
\]
We therefore obtain canonical isomorphisms
\[
\begin{aligned}
  D(D'(G))
  &=
  \operatorname{Spf}^*\mathcal O(D'(G))
  =
  \operatorname{Spf}^*\mathbf H(G)
  =
  G,
  \\
  D'(D(T))
  &=
  \operatorname{Spec}\mathbf H(D(T))
  =
  \operatorname{Spec}\mathcal O(T)
  =
  T.
\end{aligned}
\]
Moreover, if \(k\text{-}\operatorname{Gr.}\) denotes the category of
\(k\)-group schemes, \SGARefLink{sga3:VIIB:2.2.1}{Corollary 2.2.1} gives functorial isomorphisms
\[
  \operatorname{Hom}_{\operatorname{Grf}/k}\bigl(G,D(T)\bigr)
  \simeq
  \operatorname{Hom}_{\mathcal H_k}
    \bigl(\mathbf H(G),\mathcal O(T)\bigr)
  \simeq
  \operatorname{Hom}_{k\text{-}\operatorname{Gr.}}
    \bigl(T,D'(G)\bigr).
\]
This proves the following.

\medskip
\noindent\textbf{Proposition (Cartier duality).}
\phantomsection
\label{sga3:VIIB:2.2.2.prop}
{\itshape Let \(k\) be an Artinian ring.
\begin{enumerate}
\item[\textup{(i)}]
The functors
\[
  G\longmapsto D'(G)=\operatorname{Spec}\mathbf H(G),
  \qquad
  T\longmapsto D(T)=\operatorname{Spf}^*\mathcal O(T)
\]
induce an anti-equivalence between the category
\(\mathcal{FC}_k\) of commutative topologically flat formal
\(k\)-groups and the category \(\mathcal{AC}_k\) of commutative
affine flat \(k\)-group schemes. In other words, \(G\) and \(D'(G)\),
respectively \(T\) and \(D(T)\), are related by the equalities
\footnote{\label{sga3:VIIB:nde:84}Let
\(\mathcal{FC}_k^f\), respectively \(\mathcal{AC}_k^f\), denote the
full subcategory of \(\mathcal{FC}_k\), respectively
\(\mathcal{AC}_k\), consisting of the objects \(G\), respectively
\(T\), for which \(\mathbf H(G)\), respectively \(\mathcal O(T)\),
is a finite \(k\)-module \textup{(and hence finite locally free)}.
Then \(\mathcal{FC}_k^f\) and \(\mathcal{AC}_k^f\) both have the
same objects as the category \(\mathcal C\) of commutative and
cocommutative Hopf \(k\)-algebras that are finite and flat over
\(k\). The correspondence \(G\mapsto\mathbf H(G)\), respectively
\(T\mapsto\mathcal O(T)\), is covariant, respectively
contravariant. Thus one recovers the ``Cartier duality'' of the
category \(\mathcal{AC}_k\), already seen in
VII\textsubscript{A} \SGARefLink{sga3:VIIA:3.3.1}{3.3.1}.}
\[
  \mathbf H(G)=\mathcal O(D'(G)),
  \qquad
  \mathcal O(T)=\mathbf H(D(T)).
\]

\item[\textup{(ii)}]
This anti-equivalence ``commutes with base change'': if \(k\to K\)
is a morphism of Artinian rings, then
\[
  D'\bigl(G\widehat\otimes_kK\bigr)
  =
  D'(G)\otimes_kK
\]
and
\[
  D\bigl(T\otimes_kK\bigr)
  =
  D(T)\widehat\otimes_kK.
\]

\item[\textup{(iii)}]
In particular, if \(k\) is a field, one obtains an
anti-equivalence between the category of commutative formal
\(k\)-groups and the category of commutative affine \(k\)-group
schemes, and this anti-equivalence commutes with extension of the
base field.
\end{enumerate}}

% Next cursor: Exp7B local p. 47, combined-reader p. 560,
% re-edition printed p. 549, immediately before Section 2.3.

% END INCLUDED SGA3_English_current_buildable_source_20260823/source/inputs/VII/tex/components/38_expose_VIIB_section221_flat_bialgebra_equivalence_en.tex

% BEGIN INCLUDED SGA3_English_current_buildable_source_20260823/source/inputs/VII/tex/components/39_expose_VIIB_section23_covariant_algebra_en.tex
% SGA 3, Expose VII_B, English reconstruction -- Loop 1.
% Scope: authority local page 47, Section 2.3 through the paragraph
% preceding Proposition 2.3.1.
% Combined-reader page: 560. Re-edition printed page: 549.
% Original-edition marginal locator: 516.
% Authority: Polo--Gille Exp7B-13oct24.pdf,
% SHA-256 D303690E0195654BBE7C3DAA5D89EA23E24CD709A3575217234AC19AE8836694.
% Jacob Reinhold's CC BY 4.0 Markdown was consulted as comparison lineage.
% No genuine diagram occurs in this span; all mathematics is native LaTeX.

\subsection*{2.3}
\phantomsection
\label{sga3:VIIB:2.3}
\addcontentsline{toc}{subsection}{2.3. The covariant algebra}

Let \(G\) now be an arbitrary formal \(k\)-group,\footnote{
\label{sga3:VIIB:nde:85}That is, \(G\) is not necessarily
topologically flat over \(k\).} with affine algebra \(A\). We
continue to write
\[
  \mathbf H(G)=\mathbf V_k^f(A)
\]
for the \(\mathbf O_k\)-module dual to \(A\), and denote by
\[
  \phi_G:
  \mathbf H(G)\otimes_k\mathbf H(G)
  \longrightarrow
  \mathbf H(G\times G)
\]
the functorial homomorphism induced, for every object \(C\) of
\(\operatorname{Alf}/k\), by the natural map of \SGARefLink{sga3:VIIB:0.3.6}{0.3.6}:
\[
\begin{aligned}
  &(A\widehat\otimes_kC)^\dagger
  \otimes_C
  (A\widehat\otimes_kC)^\dagger
  \\
  &\hspace{3em}\longrightarrow
  \left(
    (A\widehat\otimes_kC)
    \widehat\otimes_C
    (A\widehat\otimes_kC)
  \right)^\dagger.
\end{aligned}
\]

If
\[
  m:G\times G\longrightarrow G
\]
is the multiplication of \(G\), the composite
\[
  \mathbf H(G)\otimes_k\mathbf H(G)
  \xrightarrow{\ \phi_G\ }
  \mathbf H(G\times G)
  \xrightarrow{\ \mathbf H(m)\ }
  \mathbf H(G)
\]
endows \(\mathbf H(G)\) with an algebra structure over
\(\mathbf O_k\). For every
\(C\in\operatorname{Alf}/k\), the unit element of
\[
  \mathbf H(G)(C)
  =
  (A\widehat\otimes_kC)^\dagger
\]
is the augmentation of \(A\widehat\otimes_kC\); cf.\ \SGARefLink{sga3:VIIA:2.1}{2.1}.
\footnote{\label{sga3:VIIB:nde:86}The editors changed the order of
the sentences that follow.}

If \(G\) is not topologically flat over \(k\), the morphism
\(\phi_G\) need not be an isomorphism. Consequently, the morphism
\[
  \delta_G:
  \mathbf H(G)\longrightarrow\mathbf H(G\times G)
\]
induced by the diagonal \(x\mapsto(x,x)\) of \(G\) need not factor
through
\[
  \mathbf H(G)\otimes_k\mathbf H(G).
\]
Thus \(\mathbf H(G)\) need not be a bialgebra over
\(\mathbf O_k\). For this reason, in the general case we simply call
\(\mathbf H(G)\) the \emph{covariant algebra} of the formal group
\(G\).

When \(G\) is topologically flat over \(k\), of course, \(\phi_G\) is
an isomorphism, and one recovers the \(\mathbf O_k\)-bialgebra
structure on \(\mathbf H(G)\) defined in \SGARefLink{sga3:VIIB:2.2.1}{2.2.1}.

% END INCLUDED SGA3_English_current_buildable_source_20260823/source/inputs/VII/tex/components/39_expose_VIIB_section23_covariant_algebra_en.tex

% BEGIN INCLUDED SGA3_English_current_buildable_source_20260823/source/inputs/VII/tex/components/40_expose_VIIB_proposition231_covariant_algebra_morphisms_en.tex
% SGA 3, Expose VII_B, English reconstruction -- Loop 1.
% Scope: authority local pages 48--49, Proposition 2.3.1 and its proof
% in full, stopping immediately before 2.3.2.
% Combined-reader pages: 561--562. Re-edition printed pages: 550--551.
% Original-edition marginal locators: 517--518.
% Editorial note in scope: 87.
% Authority: Polo--Gille Exp7B-13oct24.pdf,
% SHA-256 D303690E0195654BBE7C3DAA5D89EA23E24CD709A3575217234AC19AE8836694.
% Jacob Reinhold's CC BY 4.0 Markdown was consulted as comparison lineage.
% The three genuine diagrams are exact PDF-derived Loop 1 source crops;
% all ordinary mathematics is native LaTeX.

\medskip
\noindent\textbf{Proposition 2.3.1.}
\phantomsection
\label{sga3:VIIB:2.3.1}
{\itshape Let \(K\) and \(G\) be two formal \(k\)-groups, with
affine algebras \(B\) and \(A\), respectively. Suppose that \(K\) is
topologically flat over \(k\). Then there is a canonical bijection
from
\[
  \operatorname{Hom}_{\operatorname{Grf}/k}(K,G)
\]
onto the set of homomorphisms of unital \(\mathbf O_k\)-algebras
\[
  h:\mathbf H(K)\longrightarrow\mathbf H(G)
\]
for which the diagram
\begin{center}
  \phantomsection
  \label{sga3:VIIB:2.3.1:compatibility-diagram}
  % SGA3-VII-DIAG-102
  % BEGIN INCLUDED SGA3_English_current_buildable_source_20260823/source/inputs/VII/native_diagrams/exp7/viib_15_95/diag_102.tex
% SGA3-VII-DIAG-102
% Authority: Exp7B-13oct24.pdf, local p. 48.
% Crop: figures/Exp7B_localp048_proposition231_compatibility_diagram.png
% Locator: tex/components/40_expose_VIIB_proposition231_covariant_algebra_morphisms_en.tex, Loop 1 line 32.
\begin{tikzcd}[row sep=1.25em,column sep=4em]
\mathbf H(K)\otimes_k\mathbf H(K)
  \arrow[r,"h\otimes h"]
  & \mathbf H(G)\otimes_k\mathbf H(G)
      \arrow[dr,"\varphi_G"] & {} \\
{} & {} & \mathbf H(G\times G) \\
\mathbf H(K) \arrow[uu,"\Delta_{\mathbf H(K)}"]
  \arrow[r,"h"]
  & \mathbf H(G) \arrow[ur,"\delta_G"']
  & {}
\end{tikzcd}

% END INCLUDED SGA3_English_current_buildable_source_20260823/source/inputs/VII/native_diagrams/exp7/viib_15_95/diag_102.tex

\end{center}
is commutative.}

Since \(K\) is topologically flat, \(\mathbf H(K)\) is endowed with
a bialgebra structure \textup{(cf.\ \SGARefLink{sga3:VIIB:2.2}{2.2})}, and
\(\Delta_{\mathbf H(K)}\) is its diagonal morphism. In other words,
with the notation of \SGARefLink{sga3:VIIB:2.3}{2.3},
\[
  \Delta_{\mathbf H(K)}=\phi_K^{-1}\circ\delta_K.
\]
When \(G\) too is topologically flat over \(k\), the proposition
follows from the equivalence of categories established in \SGARefLink{sga3:VIIB:2.2.1}{2.2.1}.

In the general case, we may suppose that \(k\) is Artinian and argue
with the algebras
\[
  \mathbf H(K)=B^\dagger,
  \qquad
  \mathbf H(G)=A^\dagger.
\]
Let \(\operatorname{Hom}_c(A,B)\) denote the set of continuous
\(k\)-linear maps from \(A\) to \(B\), and let
\(\operatorname{Hom}_k(B^\dagger,A^\dagger)\) denote the set of
\(k\)-linear maps from \(B^\dagger\) to \(A^\dagger\).
\footnote{\label{sga3:VIIB:nde:87}The editors supplied the details of
the remainder of the proof.}

By \SGARefLink{sga3:VIIB:0.3.6.A}{0.3.6.A}, if \(M\) and \(P\) are pseudocompact \(k\)-modules and
\(P\) is projective, the canonical map
\[
  \operatorname{Hom}_c(M,P)
  \longrightarrow
  \operatorname{Hom}_k(P^\dagger,M^\dagger),
  \qquad
  f\longmapsto{}^t f,
\]
where \({}^t f\) denotes the transpose of \(f\), is bijective. We
shall apply this either with
\[
  M=A\widehat\otimes A,\quad P=B,
\]
or with
\[
  M=A,\quad P=B\widehat\otimes B.
\]

Let \(f\in\operatorname{Hom}_c(A,B)\). Consider the following
diagrams. The squares marked \textup{\SGARefLink{sga3:VIIB:0}{(0)}} are commutative, and the
two unnamed vertical arrows in the second diagram are
\({}^t(f\widehat\otimes f)\).
\begin{center}
  \phantomsection
  \label{sga3:VIIB:2.3.1:affine-diagram}
  % SGA3-VII-DIAG-103
  % BEGIN INCLUDED SGA3_English_current_buildable_source_20260823/source/inputs/VII/native_diagrams/exp7/viib_15_95/diag_103.tex
% SGA3-VII-DIAG-103
% Authority: Exp7B-13oct24.pdf, local p. 48.
% Crop: figures/Exp7B_localp048_proposition231_affine_diagram.png
% Locator: tex/components/40_expose_VIIB_proposition231_covariant_algebra_morphisms_en.tex, Loop 1 line 87.
\begin{tikzcd}[row sep=3.1em,column sep=3.5em]
A\widehat\otimes A
  \arrow[r,"m_A"] \arrow[d,"{f\widehat\otimes f}"']
  \arrow[dr,phantom,"\textup{(1)}" description]
  & A \arrow[r,"\Delta_A"] \arrow[d,"f"]
      \arrow[dr,phantom,"\textup{(2)}" description]
  & A\widehat\otimes A \arrow[d,"{f\widehat\otimes f}"] \\
B\widehat\otimes B \arrow[r,"m_B"]
  & B \arrow[r,"\Delta_B"]
  & B\widehat\otimes B
\end{tikzcd}

% END INCLUDED SGA3_English_current_buildable_source_20260823/source/inputs/VII/native_diagrams/exp7/viib_15_95/diag_103.tex


  \medskip
  \phantomsection
  \label{sga3:VIIB:2.3.1:dual-diagram}
  % SGA3-VII-DIAG-104
  % BEGIN INCLUDED SGA3_English_current_buildable_source_20260823/source/inputs/VII/native_diagrams/exp7/viib_15_95/diag_104.tex
% SGA3-VII-DIAG-104
% Authority: Exp7B-13oct24.pdf, local p. 48.
% Crop: figures/Exp7B_localp048_proposition231_dual_diagram.png
% Locator: tex/components/40_expose_VIIB_proposition231_covariant_algebra_morphisms_en.tex, Loop 1 line 93.
\begin{tikzcd}[row sep=3.4em,column sep=2.15em]
A^\dagger\otimes A^\dagger
  \arrow[r,"\varphi_G"]
  \arrow[dr,phantom,"\textup{(0)}" description]
  & (A\widehat\otimes A)^\dagger
  & A^\dagger
      \arrow[l,"{\delta_G=m_A^t}"']
      \arrow[dr,phantom,"\textup{(2')}" description]
  & (A\widehat\otimes A)^\dagger
      \arrow[l,"\Delta_A^t"']
      \arrow[dr,phantom,"\textup{(0)}" description]
  & A^\dagger\otimes A^\dagger \arrow[l,"\varphi_G"'] \\
B^\dagger\otimes B^\dagger
  \arrow[r,equals] \arrow[u,"{{}^tf\otimes{}^tf}"]
  & (B\widehat\otimes B)^\dagger
      \arrow[u]
  & B^\dagger
      \arrow[l,"{\delta_K=m_B^t}"]
      \arrow[u,"{{}^tf}"']
      \arrow[ul,phantom,"\textup{(1')}" description]
  & (B\widehat\otimes B)^\dagger
      \arrow[l,"\Delta_B^t"]
      \arrow[u]
  & B^\dagger\otimes B^\dagger
      \arrow[l,equals] \arrow[u,"{{}^tf\otimes{}^tf}"']
\end{tikzcd}

% END INCLUDED SGA3_English_current_buildable_source_20260823/source/inputs/VII/native_diagrams/exp7/viib_15_95/diag_104.tex

\end{center}

If \(f:A\to B\) corresponds to a morphism of formal groups
\(K\to G\), then the squares \textup{(1)} and \textup{(2)} are
commutative and
\[
  \varepsilon_B\circ f=\varepsilon_A.
\]
Consequently, the squares \textup{(1\(^{\prime}\))} and
\textup{(2\(^{\prime}\))} are commutative, and \({}^t f\) sends the
unit of \(B^\dagger=\mathbf H(K)\) to that of
\(A^\dagger=\mathbf H(G)\). Thus \({}^t f\) is a homomorphism of
unital \(k\)-algebras
\[
  \mathbf H(K)\longrightarrow\mathbf H(G)
\]
for which the diagram in the proposition is commutative.

Conversely, suppose that \({}^t f\) satisfies these conditions.
Then
\[
  \varepsilon_B\circ f=\varepsilon_A
\]
and the squares \textup{(1\(^{\prime}\))} and
\textup{(2\(^{\prime}\))} are commutative. For
\[
  M=A\widehat\otimes A,\quad P=B
  \qquad
  \text{(respectively \(M=A,\ P=B\widehat\otimes B\))},
\]
the map \(g\mapsto{}^t g\) is injective. Hence the squares
\textup{(1)} and \textup{(2)} are commutative, so \(f\) is
compatible with the multiplications and diagonal morphisms of \(A\)
and \(B\).

It remains to prove that \(f(1_A)=1_B\). What precedes gives
\[
  \varepsilon_B f(1)=1,
  \qquad
  \Delta_B f(1)=f(1)\widehat\otimes f(1),
  \qquad
  f(1)\cdot f(1)=f(1).
\]
By \SGARefLink{sga3:VIIA:2.1}{2.1}\textup{(iii)}, the first two conditions imply that \(f(1)\)
has inverse \(c_Bf(1)\) in \(B\). The equality
\(f(1)\cdot f(1)=f(1)\) therefore implies \(f(1)=1\). Thus
\[
  f:A\longrightarrow B
\]
is a morphism in \(\operatorname{Alp}/k\), compatible with the
comultiplications of \(A\) and \(B\).

% END INCLUDED SGA3_English_current_buildable_source_20260823/source/inputs/VII/tex/components/40_expose_VIIB_proposition231_covariant_algebra_morphisms_en.tex

% BEGIN INCLUDED SGA3_English_current_buildable_source_20260823/source/inputs/VII/tex/components/41_expose_VIIB_section232_cartier_universal_property_en.tex
% SGA 3, Expose VII_B, English reconstruction -- Loop 1.
% Scope: authority local pages 49--51, Section 2.3.2 in full, stopping
% immediately before Section 2.4.
% Combined-reader pages: 562--564. Re-edition printed pages: 551--553.
% Original-edition marginal locators: 518--519.
% Editorial notes in scope: 88--89.
% Authority: Polo--Gille Exp7B-13oct24.pdf,
% SHA-256 D303690E0195654BBE7C3DAA5D89EA23E24CD709A3575217234AC19AE8836694.
% Jacob Reinhold's CC BY 4.0 Markdown was consulted as comparison lineage.
% The beta diagram is an exact PDF-derived Loop 1 source crop; all
% ordinary mathematics is native LaTeX.

\subsection*{2.3.2}
\phantomsection
\label{sga3:VIIB:2.3.2}
\addcontentsline{toc}{subsection}{\SGARefLink{sga3:VIIB:2.3.2:prop}{2.3.2}}

Let us now suppose, for simplicity, that the ring \(k\) is Artinian.
When \(G\) is topologically flat over \(k\), the algebra
\[
  \mathbf H(G)=\mathbf H(G)(k)
\]
may be characterized by a universal property due to Cartier. Recall
\textup{(cf.\ \SGARefLink{sga3:VIIA:1.2.1}{1.2.1})} that, if \(U\) is a \(k\)-module, we write
\(\mathbf W(U)\) for the functor that assigns to every finite-length
\(k\)-algebra \(C\) the \(C\)-module
\[
  U\otimes_k C.
\]
\footnote{\label{sga3:VIIB:nde:88}The editors added the preceding
sentence and, in what follows, wrote \(\mathbf W(U)^\times\) in
place of \(U_m\). Recall also \textup{(cf.\ \SGARefLink{sga3:VIIA:1.2.1}{1.2.1})} that a
``\(k\)-functor'' is a covariant functor from
\(\operatorname{Alf}/k\) to \((\mathbf{Sets})\).}
If \(U\) is a \(k\)-algebra \textup{(associative and with a unit
element)}, then so is \(U\otimes_k C\). We shall write
\(\mathbf W(U)^\times\) for the \(k\)-group functor that assigns to
every \(C\in\operatorname{Ob}(\operatorname{Alf}/k)\) the
multiplicative group of invertible elements of the algebra
\(U\otimes_k C\):
\[
  \mathbf W(U)^\times(C)=(U\otimes_k C)^\times.
\]

Moreover, identify \(G\) with the \(k\)-group functor
\[
  C\longmapsto
  \operatorname{Hom}_{\operatorname{Vaf}/k}
    \bigl(\operatorname{Spf}(C),G\bigr),
\]
and write
\[
  \operatorname{Hom}_{k\text{-}\operatorname{Gr.}}
    \bigl(G,\mathbf W(U)^\times\bigr)
\]
for the set of homomorphisms of \(k\)-group functors from \(G\) to
\(\mathbf W(U)^\times\). We have the following.

\medskip
\noindent\textbf{Proposition.}
\phantomsection
\label{sga3:VIIB:2.3.2:prop}
{\itshape Let \(k\) be an Artinian ring. For every formal group \(G\)
topologically flat over \(k\), and every \(k\)-algebra \(U\), there
is a canonical isomorphism
\[
  \operatorname{Hom}_{k\text{-}\operatorname{Gr.}}
    \bigl(G,\mathbf W(U)^\times\bigr)
  \xrightarrow{\ \sim\ }
  \operatorname{Hom}_{k\text{-}\operatorname{Alg.}}
    \bigl(\mathbf H(G),U\bigr).
\]}

Let \(A\) be the affine algebra of \(G\). By hypothesis, \(A\) is a
projective object of \(\operatorname{PC}(k)\), and
\[
  \mathbf H(G)=A^\dagger.
\]
For every object \(P\) of \(\operatorname{PC}(k)\), write
\[
  U\widehat\otimes_k P
\]
for the inverse limit of the \(k\)-modules
\[
  U\otimes_k(P/N),
\]
where \(N\) runs through the open submodules of \(P\). There are
linear maps
\[
  U\otimes_k(P/N)
  \longrightarrow
  \operatorname{Hom}_k\bigl((P/N)^*,U\bigr)
\]
that send \(u\otimes x\) to the \(k\)-linear map
\[
  f\longmapsto f(x)u
\]
and form a filtered inverse system. Passing to the inverse limit
therefore gives a morphism
\begin{equation}
  U\widehat\otimes_kP
  \xrightarrow{\ \psi_P\ }
  \operatorname{Hom}_k(P^\dagger,U).
  \tag{1}
  \label{sga3:VIIB:2.3.2:eq1}
\end{equation}

When \(P=k\), the morphism \(\psi_k\) is plainly an isomorphism.
Moreover, the two sides of \textup{\SGARefLink{sga3:VIIB:2.3.2:eq1}{(1)}}, considered as functors of
\(P\), commute with infinite products \textup{(every product being
a filtered inverse limit of finite products)}. Thus \textup{\SGARefLink{sga3:VIIB:2.3.2:eq1}{(1)}}
is an isomorphism when \(P\) is a product of copies of \(k\), and
then when \(P\) is a projective object of \(\operatorname{PC}(k)\)
\textup{(the two sides of \textup{\SGARefLink{sga3:VIIB:2.3.2:eq1}{(1)}} commute with finite direct
sums)}.

Now let
\[
  \operatorname{Hom}_{\mathcal F}\bigl(G,\mathbf W(U)\bigr)
\]
denote the set of morphisms of \(k\)-functors from \(G\) to
\(\mathbf W(U)\). Since
\[
  G=\operatorname{Spf}(A)
   =\varinjlim_{\mathfrak l}\operatorname{Spf}(A/\mathfrak l),
\]
where \(\mathfrak l\) runs through the open ideals of \(A\), there
are canonical isomorphisms
\[
\begin{aligned}
  \operatorname{Hom}_{\mathcal F}\bigl(G,\mathbf W(U)\bigr)
  &=
  \varprojlim_{\mathfrak l}
  \operatorname{Hom}_{\mathcal F}
    \bigl(\operatorname{Spf}(A/\mathfrak l),\mathbf W(U)\bigr)
  \\
  &=
  \varprojlim_{\mathfrak l}
  U\otimes_k(A/\mathfrak l)
  =
  U\widehat\otimes_kA.
\end{aligned}
\]
The preceding argument therefore gives a canonical isomorphism
\begin{equation}
  \operatorname{Hom}_{\mathcal F}\bigl(G,\mathbf W(U)\bigr)
  =
  U\widehat\otimes_kA
  \xrightarrow{\ \psi_A\ }
  \operatorname{Hom}_k\bigl(\mathbf H(G),U\bigr).
  \tag{2}
  \label{sga3:VIIB:2.3.2:eq2}
\end{equation}

For every finite-length \(k\)-algebra \(C\), multiplication makes
\(U\otimes_kC\) a monoid with a unit element, and every homomorphism
of monoids with unit
\[
  G(C)\longrightarrow U\otimes_kC
\]
is necessarily a group homomorphism
\[
  G(C)\longrightarrow(U\otimes_kC)^\times.
\]
It follows that
\[
  \operatorname{Hom}_{k\text{-}\operatorname{Gr.}}
    \bigl(G,\mathbf W(U)^\times\bigr)
\]
is the part of
\(\operatorname{Hom}_{\mathcal F}(G,\mathbf W(U))\) consisting of
morphisms of \(k\)-functors into monoids with a unit element.

It remains to show that these morphisms correspond to the
\(k\)-linear maps from \(\mathbf H(G)\) to \(U\) that preserve
multiplication and unit elements.
\footnote{\label{sga3:VIIB:nde:89}The original stated that this
``follows from the functoriality of \(\psi_A\).'' The editors gave
the details that follow.}
To simplify the notation, write
\[
  H=\mathbf H(G)=A^\dagger
\]
and write \(\widehat\otimes\) instead of
\(\widehat\otimes_k\). Denote by
\[
  \Delta_A,\ m_A,\ \varepsilon_A
  \qquad\text{and}\qquad
  \Delta_H,\ m_H,\ \varepsilon_H
\]
the comultiplication, multiplication, and augmentation of \(A\)
and \(H\), respectively. Let
\[
  \theta\in
  \operatorname{Hom}_{\mathcal F}\bigl(G,\mathbf W(U)\bigr),
\]
let \(\gamma\) be its image in \(U\widehat\otimes A\), and let
\[
  \phi:H\longrightarrow U
\]
be the associated \(k\)-linear map. The morphism \(\theta\) sends
the unit section \(s\in G(k)\) to an element \(u\) of \(U\). Since
\(s\) corresponds to the augmentation
\[
  \varepsilon:A\longrightarrow k,
\]
which is the unit element \(1_H\) of \(H\), we have
\[
  \theta(s)=1_U
  \quad\Longleftrightarrow\quad
  \phi(1_H)=1_U.
\]

The morphism
\[
  \theta\circ m_G:G\times G\longrightarrow\mathbf W(U)
\]
corresponds to the element
\[
  (\operatorname{id}_U\widehat\otimes\Delta_A)(\gamma),
\]
which, by duality, corresponds to the map
\[
  \phi\circ m_H:H\otimes H\longrightarrow U.
\]

On the other hand, the morphism
\[
  \theta\circ\operatorname{pr}_1:
  G\times G\longrightarrow\mathbf W(U)
\]
corresponds to the element
\[
  \gamma\widehat\otimes1_A
  \in
  U\widehat\otimes A\widehat\otimes A,
\]
which, by duality, corresponds to
\[
  \phi\circ(\operatorname{id}_H\otimes\varepsilon_H):
  H\otimes H\longrightarrow U.
\]
Similarly, \(\theta\circ\operatorname{pr}_2\) corresponds to the
element
\[
  \tau(\gamma\widehat\otimes1_A)
  \in
  U\widehat\otimes A\widehat\otimes A,
\]
where
\[
  \tau(u\widehat\otimes a\widehat\otimes b)
  =
  u\widehat\otimes b\widehat\otimes a,
\]
and this element corresponds by duality to
\[
  \phi\circ(\varepsilon_H\otimes\operatorname{id}_H):
  H\otimes H\longrightarrow U.
\]

Finally, the multiplication map
\[
  \mu=m_{U\widehat\otimes A\widehat\otimes A}
\]
given by
\[
\begin{aligned}
  &(U\widehat\otimes A\widehat\otimes A)
   \times
   (U\widehat\otimes A\widehat\otimes A)
  \longrightarrow
  U\widehat\otimes A\widehat\otimes A,
  \\
  &\bigl(
    u\widehat\otimes a_1\widehat\otimes a_2,\,
    u'\widehat\otimes a_3\widehat\otimes a_4
  \bigr)
  \longmapsto
  uu'\widehat\otimes a_1a_3\widehat\otimes a_2a_4
\end{aligned}
\]
may be viewed as the composite of the endomorphism
\(\sigma_{23}\) of
\[
  (U\otimes U)\widehat\otimes A^{\widehat\otimes4},
\]
which ``exchanges the factors \(a_2\) and \(a_3\),'' and the map
\[
  (U\otimes U)
  \widehat\otimes A^{\widehat\otimes2}
  \widehat\otimes A^{\widehat\otimes2}
  \xrightarrow{\ m_U\widehat\otimes m_A\widehat\otimes m_A\ }
  U\widehat\otimes A\widehat\otimes A.
\]

It follows that the map
\[
  \mu\circ
  \bigl(\theta\circ\operatorname{pr}_1,\,
        \theta\circ\operatorname{pr}_2\bigr):
  G\times G\longrightarrow\mathbf W(U)
\]
corresponds to the composite map \(\beta\) below:
\begin{center}
  \phantomsection
  \label{sga3:VIIB:2.3.2:beta-diagram}
  % SGA3-VII-DIAG-105
  % BEGIN INCLUDED SGA3_English_current_buildable_source_20260823/source/inputs/VII/native_diagrams/exp7/viib_15_95/diag_105.tex
% SGA3-VII-DIAG-105
% Authority: Exp7B-13oct24.pdf, local p. 51.
% Crop: figures/Exp7B_localp051_cartier_universal_beta_diagram.png
% Locator: tex/components/41_expose_VIIB_section232_cartier_universal_property_en.tex, Loop 1 line 308.
\begin{tikzcd}[row sep=2.7em,column sep=3em]
{} & H\otimes H
     \arrow[r,"\beta"]
     \arrow[dl,"{\operatorname{id}_H\otimes\sigma_{23}\otimes
                   \operatorname{id}_H}"']
   & U \\
H\otimes H\otimes H\otimes H
  \arrow[dr,"{(\operatorname{id}_H\otimes\varepsilon_H)
              \otimes(\varepsilon_H\otimes\operatorname{id}_H)}"']
  & {} & U\otimes U \arrow[u,"m_U"'] \\
{} & H\otimes H \arrow[ur,"\phi\otimes\phi"'] & {.}
\end{tikzcd}

% END INCLUDED SGA3_English_current_buildable_source_20260823/source/inputs/VII/native_diagrams/exp7/viib_15_95/diag_105.tex

\end{center}

Finally, \(\theta\) is compatible with the laws of \(G\) and
\(\mathbf W(U)^\times\) if and only if
\[
  \theta\circ m_G
  =
  \mu\circ
  \bigl(\theta\circ\operatorname{pr}_1,\,
        \theta\circ\operatorname{pr}_2\bigr),
\]
\noindent\textit{Source note.}\enspace
The French re-edition prints \(\phi\circ\operatorname{pr}_i\) in
these two ordered pairs. Since
\(\operatorname{pr}_i:G\times G\to G\) whereas
\(\phi:H\to U\), Canon record \texttt{SGA-SRC-000847} restores the
type-correct \(\theta\circ\operatorname{pr}_i\).

which, by the preceding argument, is equivalent to
\[
  \phi\circ m_H=\beta.
\]
But plainly
\[
  (\phi\circ m_H)(x\otimes y)=\phi(xy),
\]
whereas
\[
  \beta(x\otimes y)=\phi(x)\phi(y).
\]

% END INCLUDED SGA3_English_current_buildable_source_20260823/source/inputs/VII/tex/components/41_expose_VIIB_section232_cartier_universal_property_en.tex

% BEGIN INCLUDED SGA3_English_current_buildable_source_20260823/source/inputs/VII/tex/components/42_expose_VIIB_section24_formal_group_quotients_en.tex
% SGA 3, Expose VII_B, English reconstruction -- Loop 1.
% Scope: authority local pages 51--52, Section 2.4 in full, including
% its theorem and corollary, stopping immediately before 2.4.A.
% Combined-reader pages: 564--565. Re-edition printed pages: 553--554.
% Original-edition marginal locator in scope: 520.
% Editorial notes in scope: 90--91.
% Authority: Polo--Gille Exp7B-13oct24.pdf,
% SHA-256 D303690E0195654BBE7C3DAA5D89EA23E24CD709A3575217234AC19AE8836694.
% Jacob Reinhold's CC BY 4.0 Markdown was consulted as comparison lineage.
% The three genuine diagrams are exact PDF-derived Loop 1 source crops;
% all ordinary mathematics is native LaTeX.

\subsection*{2.4}
\phantomsection
\label{sga3:VIIB:2.4}
\addcontentsline{toc}{subsection}{2.4. Quotients of formal groups}

We now return to an arbitrary pseudocompact ring \(k\), in order to
apply to formal groups the results of \SGARefLink{sga3:VIIB:1.4.thm}{1.4}--\SGARefLink{sga3:VIIB:1.5}{1.5} on passage to the
quotient by a topologically flat equivalence relation.
\footnote{\label{sga3:VIIB:nde:90}The editors supplied the details of
the original argument that follow.}

Let \(u:H\to G\) be a monomorphism of formal \(k\)-groups, let
\(\mu:G\times G\to G\) be the multiplication morphism of \(G\), and
let \(\lambda\) be the composite
\begin{center}
  \phantomsection
  \label{sga3:VIIB:2.4:lambda-composite}
  % SGA3-VII-DIAG-106
  % BEGIN INCLUDED SGA3_English_current_buildable_source_20260823/source/inputs/VII/native_diagrams/exp7/viib_15_95/diag_106.tex
% SGA3-VII-DIAG-106
% Authority: Exp7B-13oct24.pdf, local p. 51.
% Crop: figures/Exp7B_localp051_quotient_action_lambda_composite.png
% Locator: tex/components/42_expose_VIIB_section24_formal_group_quotients_en.tex, Loop 1 line 30.
\(
\lambda:\quad
G\times H\xrightarrow{\ \operatorname{id}_G\times u\ }
G\times G\xrightarrow{\ \mu\ }G.
\)

% END INCLUDED SGA3_English_current_buildable_source_20260823/source/inputs/VII/native_diagrams/exp7/viib_15_95/diag_106.tex

\end{center}
Since \(u\) is a monomorphism, the pair
\begin{center}
  \phantomsection
  \label{sga3:VIIB:2.4:equivalence-pair}
  % SGA3-VII-DIAG-107
  % BEGIN INCLUDED SGA3_English_current_buildable_source_20260823/source/inputs/VII/native_diagrams/exp7/viib_15_95/diag_107.tex
% SGA3-VII-DIAG-107
% Authority: Exp7B-13oct24.pdf, local p. 51.
% Crop: figures/Exp7B_localp051_quotient_equivalence_pair.png
% Locator: tex/components/42_expose_VIIB_section24_formal_group_quotients_en.tex, Loop 1 line 37.
\begin{tikzpicture}[baseline=(lhs.base)]
  \node (lhs) {$G\times H$};
  \node[right=2.5cm of lhs] (rhs) {$G$};
  \draw[-{Straight Barb[length=4pt,width=4pt]}]
    ([yshift=2.5pt]lhs.east) --
    node[above] {$\operatorname{pr}_1$}
    ([yshift=2.5pt]rhs.west);
  \draw[-{Straight Barb[length=4pt,width=4pt]}]
    ([yshift=-2.5pt]lhs.east) --
    node[below] {$\lambda$}
    ([yshift=-2.5pt]rhs.west);
\end{tikzpicture}

% END INCLUDED SGA3_English_current_buildable_source_20260823/source/inputs/VII/native_diagrams/exp7/viib_15_95/diag_107.tex

\end{center}
is an equivalence pair in \(\operatorname{Vaf}/k\)
\textup{(cf.\ V, 2.b)}. Recall \textup{(cf.\ \SGARefLink{sga3:VIIB:1.2.C}{1.2.C})} that the
cokernel \(G/H\) of this pair is defined as follows.

Let \(\mathcal O(G)\) and \(\mathcal O(H)\) be the affine algebras of
\(G\) and \(H\), let
\[
  \Delta:\mathcal O(G)
  \longrightarrow
  \mathcal O(G)\widehat\otimes_k\mathcal O(G)
\]
be the diagonal morphism of \(\mathcal O(G)\), and let \(I\) be the
kernel of
\[
  u^\natural:\mathcal O(G)\longrightarrow\mathcal O(H).
\]
By \SGARefLink{sga3:VIIB:1.3.prop}{Proposition 1.3}, the morphism \(u^\natural\) induces an
isomorphism
\[
  \mathcal O(G)/I\xrightarrow{\ \sim\ }\mathcal O(H).
\]
The affine algebra \(\mathcal O(G/H)\) of \(G/H\) is then the kernel
of the pair
\begin{center}
  \phantomsection
  \label{sga3:VIIB:2.4:affine-equalizer-pair}
  % SGA3-VII-DIAG-108
  % BEGIN INCLUDED SGA3_English_current_buildable_source_20260823/source/inputs/VII/native_diagrams/exp7/viib_15_95/diag_108.tex
% SGA3-VII-DIAG-108
% Authority: Exp7B-13oct24.pdf, local p. 51.
% Crop: figures/Exp7B_localp051_quotient_affine_equalizer_pair.png
% Locator: tex/components/42_expose_VIIB_section24_formal_group_quotients_en.tex, Loop 1 line 66.
\begin{tikzpicture}[baseline=(lhs.base)]
  \node (lhs) {$\mathcal O(G)$};
  \node[right=3cm of lhs] (rhs)
    {$\mathcal O(G)\widehat\otimes_k\mathcal O(H),$};
  \draw[-{Straight Barb[length=4pt,width=4pt]}]
    ([yshift=2.5pt]lhs.east) --
    node[above] {$\tau_1$}
    ([yshift=2.5pt]rhs.west);
  \draw[-{Straight Barb[length=4pt,width=4pt]}]
    ([yshift=-2.5pt]lhs.east) --
    node[below] {$(\operatorname{id}\widehat\otimes u^\natural)\Delta$}
    ([yshift=-2.5pt]rhs.west);
\end{tikzpicture}

% END INCLUDED SGA3_English_current_buildable_source_20260823/source/inputs/VII/native_diagrams/exp7/viib_15_95/diag_108.tex

\end{center}
where \(\tau_1(x)=x\widehat\otimes1\). Thus
\[
  \mathcal O(G/H)
  =
  \bigl\{
    x\in\mathcal O(G)
    \mathrel{\big|}
    \Delta(x)-x\widehat\otimes1
    \in\mathcal O(G)\widehat\otimes I
  \bigr\}.
\]

If, moreover, \(H\) is topologically flat over \(k\), then
\(\operatorname{pr}_1\) is topologically flat, and \SGARefLink{sga3:VIIB:1.4.thm}{Theorem 1.4} gives
the following result.

\medskip
\noindent\textbf{Theorem.}
\phantomsection
\label{sga3:VIIB:2.4:theorem}
{\itshape Let \(u:H\to G\) be a monomorphism of formal
\(k\)-groups. Suppose that \(H\) is topologically flat over \(k\).
Then the projection
\[
  p:G\longrightarrow G/H
\]
is surjective and topologically flat, there is an isomorphism
\[
  G\times H
  \xrightarrow{\ \sim\ }
  G\times_{G/H}G,
  \tag{*}\label{sga3:VIIB:2.4:equation:quotient-fiber-product-star}
\]
and \(G/H\) represents the quotient sheaf for the flat topology.

Consequently, \(G/H\) has a canonical structure of an object on which
\(G\) acts, and \(p:G\to G/H\) is an equivariant morphism. If, in
addition, \(u\) identifies \(H\) with a normal subgroup of \(G\),
then \(G/H\) has a canonical structure of formal \(k\)-group for
which \(p:G\to G/H\) is a morphism of formal \(k\)-groups, and \(H\)
is the kernel of \(p\).}

Indeed, the first assertion follows from \SGARefLink{sga3:VIIB:1.4.thm}{1.4}, and the other two from
\SGARefLink{sga3:iv:corollary:5:2:2}{Exposé IV, Corollaries 5.2.2} and \SGARefLink{sga3:VIIB:5.2.4}{5.2.4}.

\medskip
\noindent\textbf{Corollary.}
\phantomsection
\label{sga3:VIIB:2.4:corollary}
\footnote{\label{sga3:VIIB:nde:91}The editors made this corollary
explicit; it is used, for example, in \SGARefLink{sga3:VIIB:5.2.1:section-proposition}{5.2.1} and \SGARefLink{sga3:VIIB:5.2.3}{5.2.3}.}
{\itshape Let \(G\) be a formal \(k\)-group and \(H\) a formal
subgroup of \(G\). Let \(A\), \(A/J\), and \(B\) be the affine
algebras of \(G\), \(H\), and \(G/H\), respectively; let \(I_A\) be
the augmentation ideal of \(A\), and put \(I_B=B\cap I_A\). Suppose
that \(H\) is topologically flat over \(k\). Then
\[
  J=AI_B,
\]
the closed ideal generated by \(I_B\).}

Indeed, the projection \(B\to B/I_B\) corresponds to the unit section
\[
  e:\operatorname{Spf}(k)\longrightarrow G/H
\]
of \(G/H\). By \textup{(*)}, \(H\) is identified with the fiber
product
\[
  \operatorname{Spf}(k)\times_{G/H}G.
\]
Its affine algebra \(A/J\) is therefore identified with
\[
  (B/I_B)\widehat\otimes_BA
  \simeq
  A/AI_B.
\]

% END INCLUDED SGA3_English_current_buildable_source_20260823/source/inputs/VII/tex/components/42_expose_VIIB_section24_formal_group_quotients_en.tex

% BEGIN INCLUDED SGA3_English_current_buildable_source_20260823/source/inputs/VII/tex/components/43_expose_VIIB_section24A_semidirect_products_en.tex
% SGA 3, Expose VII_B, English reconstruction -- Loop 1.
% Scope: authority local pages 52--53, Section 2.4.A in full,
% stopping immediately before Section 2.4.B.
% Combined-reader pages: 565--566. Re-edition printed pages: 554--555.
% Editorial note in scope: 92 (shared by 2.4.A and 2.4.B).
% Authority: Polo--Gille Exp7B-13oct24.pdf,
% SHA-256 D303690E0195654BBE7C3DAA5D89EA23E24CD709A3575217234AC19AE8836694.
% Jacob Reinhold's CC BY 4.0 Markdown was consulted as comparison lineage.
% The two genuine diagrams are exact PDF-derived Loop 1 source crops;
% all ordinary mathematics is native LaTeX.

\subsubsection*{2.4.A}
\phantomsection
\label{sga3:VIIB:2.4.A}

Let \(G\) and \(Q\) be formal \(k\)-groups that are topologically
flat. Suppose that there are homomorphisms
\[
  \sigma:Q\longrightarrow G
  \qquad\text{and}\qquad
  \pi:G\longrightarrow Q
\]
such that
\[
  \pi\circ\sigma=\operatorname{id}_Q.
\]
\footnote{\label{sga3:VIIB:nde:92}The editors added paragraphs
2.4.A and \SGARefLink{sga3:VIIB:2.4.B}{2.4.B}, which will be useful in \SGARefLink{sga3:VIIB:5.1.3}{5.1.3} and \SGARefLink{sga3:VIIB:2.9}{2.9},
respectively.}
In particular, \(\sigma\) is a monomorphism, so \(Q\) is a formal
subgroup of \(G\) \textup{(cf.\ \SGARefLink{sga3:VIIB:1.3}{Remark 1.3})}. Let
\[
  N=\operatorname{Ker}(\pi)
\]
and let \(\sigma'\) be the inclusion \(N\hookrightarrow G\).
Then \(G\) is the semidirect product of \(N\) by \(Q\)
\textup{(cf.\ I, \SGARefLink{sga3:i:proposition:2-3-8}{2.3.8})}. In other words, for every
\[
  B\in\operatorname{Ob}(\operatorname{Alf}/k),
\]
identify \(N(B)\) and \(Q(B)\) with their images in \(G(B)\) under
\(\sigma'\) and \(\sigma\). Then \(N(B)\) is a normal subgroup of
\(G(B)\), and the map
\begin{equation}
  \mu:N(B)\times Q(B)\longrightarrow G(B),
  \qquad
  (x,q)\longmapsto xq
  \tag{1}
  \label{sga3:VIIB:2.4.A:eq1}
\end{equation}
is bijective.

The morphism of formal \(k\)-varieties
\begin{equation}
  \theta:G(B)\longrightarrow N(B),
  \qquad
  g\longmapsto g\cdot\sigma\pi(g^{-1})
  \tag{2}
  \label{sga3:VIIB:2.4.A:eq2}
\end{equation}
is then a retraction of \(\sigma'\), and the inverse of \(\mu\) is
\begin{equation}
  g\longmapsto\bigl(\theta(g),\pi(g)\bigr).
  \tag{3}
  \label{sga3:VIIB:2.4.A:eq3}
\end{equation}
If \(p:G\to G/Q\) denotes the quotient projection, then
\[
  p\circ\sigma':N\longrightarrow G/Q
\]
is an isomorphism of formal \(k\)-varieties.

\smallskip
\noindent\emph{Source note.}
The Polo--Gille text prints
\(\theta\circ\sigma':N\to G/Q\) here. Since
\(\theta:G\to N\), that composite is ill-typed; the isomorphism
with the displayed source and target is the quotient composite
\(p\circ\sigma'\).

\smallskip
In particular, \(N\) is topologically flat over \(k\), by \SGARefLink{sga3:VIIB:1.4.thm}{1.4}
\textup{\SGARefLink{sga3:vii:paragraph:5-1:equation:ii}{(ii)}}. Let
\[
  \alpha:Q\times N\longrightarrow N
\]
and
\[
  \beta:G\times N=N\times Q\times N\longrightarrow N
\]
be the maps defined on points by
\[
  \alpha(q,y)=qyq^{-1},
  \qquad
  \beta(x,q,y)=x\alpha(q,y).
\]
Then the following diagram is commutative:
\begin{center}
  \phantomsection
  \label{sga3:VIIB:2.4.A:eq4-diagram}
  % SGA3-VII-DIAG-109
  % BEGIN INCLUDED SGA3_English_current_buildable_source_20260823/source/inputs/VII/native_diagrams/exp7/viib_15_95/diag_109.tex
% SGA3-VII-DIAG-109
% Authority: Exp7B-13oct24.pdf, local p. 53.
% Crop: figures/Exp7B_localp053_semidirect_product_diagram4.png
% Locator: tex/components/43_expose_VIIB_section24A_semidirect_products_en.tex, Loop 1 line 101.
\begin{tikzcd}[row sep=3.2em,column sep=3.5em]
{} \arrow[d,phantom,"\textup{(4)}" description]
  & G\times G
  \arrow[r,"{\operatorname{id}\times\theta}"]
  \arrow[d,"m_G"']
  & G\times N
      \arrow[r,"{\theta\times\operatorname{id}}"]
      \arrow[d,"\beta"]
  & N\times N \arrow[dl,"m_N"] \\
{} & G \arrow[r,"\theta"] & N & {.}
\end{tikzcd}

% END INCLUDED SGA3_English_current_buildable_source_20260823/source/inputs/VII/native_diagrams/exp7/viib_15_95/diag_109.tex

\end{center}

This can be expressed as follows in terms of the affine algebras
\(A\), \(A_0\), and \(A'\) of \(G\), \(Q\), and \(N\)
\textup{(cf.\ \SGARefLink{sga3:VIIB:5.1.3}{5.1.3} below)}. Let
\[
  \rho':A\longrightarrow A',
  \qquad
  \rho:A\longrightarrow A_0,
  \qquad
  \tau:A_0\longrightarrow A
\]
be the homomorphisms of \(k\)-bialgebras corresponding to
\(\sigma'\), \(\sigma\), and \(\pi\), and put
\[
  I=\operatorname{Ker}(\rho).
\]
By the preceding corollary, \(A'\) is identified with
\[
  A\big/\overline{A\,\tau(J_0)},
\]
where \(J_0\) is the augmentation ideal of \(A_0\).

On the other hand, let \(B\) be the affine algebra of \(G/Q\).
It is the kernel of the pair of morphisms
\phantomsection
\label{sga3:VIIB:2.4.A:affine-equalizer-pair}
\begin{equation*}
  A
  \mathrel{\overset{\tau_1}
    {\underset{(\operatorname{id}_A\widehat\otimes\rho)
      \circ\Delta_A}{\rightrightarrows}}}
  A\widehat\otimes_kA_0,
\end{equation*}
where \(\tau_1(x)=x\widehat\otimes1\). Thus
\[
  B
  =
  \bigl\{
    x\in A
    \mathrel{\big|}
    \Delta_A(x)-x\widehat\otimes1
    \in A\widehat\otimes I
  \bigr\}.
\]
Let
\[
  \gamma=\theta^\natural:A'\longrightarrow A
\]
be the corresponding continuous homomorphism of \(k\)-algebras.
It is a section of \(\rho'\), and it is an isomorphism of profinite
\(k\)-algebras from \(A'\) onto \(B\). Moreover,
\[
  \tau=\pi^\natural
\]
identifies \(A_0\) with a sub-bialgebra of \(A\), namely the affine
algebra of the quotient \(N\backslash G\).

It follows from equations~\eqref{sga3:VIIB:2.4.A:eq1}
and~\eqref{sga3:VIIB:2.4.A:eq3} that there is an
isomorphism of profinite \(k\)-algebras
\begin{equation}
  A'\widehat\otimes_kA_0
  \xrightarrow{\ \sim\ }
  A,
  \qquad
  a'\widehat\otimes a_0
  \longmapsto
  \gamma(a')\tau(a_0),
  \tag{*}
  \label{sga3:VIIB:2.4.A:star}
\end{equation}
whose inverse is the map
\[
  a\longmapsto
  (\rho'\widehat\otimes\rho)\Delta_A(a).
\]

Finally, identify \(A'\) with its image in \(A\) under \(\gamma\),
so that the projection \(A\to A'\) is then \(\gamma\rho'\).
Writing \(\Delta_{A'}\) for the comultiplication of \(A'\), one
deduces from \hyperref[sga3:VIIB:2.4.A:eq4-diagram]{\textup{(4)}}
that
\[
  \Delta_A(A')\subset A\widehat\otimes A'
\]
and that the following diagram is commutative:
\begin{center}
  \phantomsection
  \label{sga3:VIIB:2.4.A:Aprime-comultiplication-square}
  % SGA3-VII-DIAG-110
  % BEGIN INCLUDED SGA3_English_current_buildable_source_20260823/source/inputs/VII/native_diagrams/exp7/viib_15_95/diag_110.tex
% SGA3-VII-DIAG-110
% Authority: Exp7B-13oct24.pdf, local p. 53.
% Crop: figures/Exp7B_localp053_Aprime_comultiplication_square.png
% Locator: tex/components/43_expose_VIIB_section24A_semidirect_products_en.tex, Loop 1 line 193.
\begin{tikzcd}[row sep=3em,column sep=3.7em]
A' \arrow[r,"\Delta_A"] \arrow[d,equals]
  & A\widehat\otimes A' \arrow[d,"{\gamma\rho'\widehat\otimes
                                      \operatorname{id}}" right] \\
A' \arrow[r,"\Delta_{A'}"]
  & A'\widehat\otimes A'
\end{tikzcd}

% END INCLUDED SGA3_English_current_buildable_source_20260823/source/inputs/VII/native_diagrams/exp7/viib_15_95/diag_110.tex

\end{center}
Consequently, one also has
\[
  \gamma\rho'\circ c_A=c_{A'},
\]
where \(c_A\) and \(c_{A'}\) are the antipodes of \(A\) and \(A'\),
respectively. On the other hand, writing \(m_A\) for multiplication
in \(A\), equation~\eqref{sga3:VIIB:2.4.A:eq2} gives, for every
\(a\in A\),
\[
  \gamma\rho'(a)
  =
  \bigl(
    m_A
    \circ
    (\operatorname{id}_A\widehat\otimes\tau\rho c_A)
    \circ
    \Delta_A
  \bigr)(a).
\]

% END INCLUDED SGA3_English_current_buildable_source_20260823/source/inputs/VII/tex/components/43_expose_VIIB_section24A_semidirect_products_en.tex

% BEGIN INCLUDED SGA3_English_current_buildable_source_20260823/source/inputs/VII/tex/components/44_expose_VIIB_section24B_covariant_semidirect_products_en.tex
% SGA 3, Expose VII_B, English reconstruction -- Loop 1.
% Scope: authority local pages 53--54, Section 2.4.B in full,
% stopping immediately before Proposition 2.4.1.
% Combined-reader pages: 566--567. Re-edition printed pages: 555--556.
% Authority: Polo--Gille Exp7B-13oct24.pdf,
% SHA-256 D303690E0195654BBE7C3DAA5D89EA23E24CD709A3575217234AC19AE8836694.
% Jacob Reinhold's CC BY 4.0 Markdown was consulted as comparison lineage.
% Editorial note 92 is attached to 2.4.A and explicitly covers both 2.4.A
% and 2.4.B; it is not duplicated here.
% The one genuine diagram is an exact PDF-derived Loop 1 source crop;
% all ordinary mathematics is native LaTeX.

\subsubsection*{2.4.B}
\phantomsection
\label{sga3:VIIB:2.4.B}

Suppose, for simplicity, that \(k\) is Artinian. The preceding
discussion can then be expressed more simply in terms of the
covariant bialgebras of \(G\), \(Q\), and \(N\). Indeed, since
\[
  G=N\times Q
\]
as formal \(k\)-varieties, one has
\[
  \mathbf H(G)
  =
  \mathbf H(N)\otimes_k\mathbf H(Q)
\]
as \(k\)-coalgebras. Moreover, since multiplication in \(G\) is
given by
\[
  (x,q)\cdot(x',q')
  =
  \bigl(x\alpha(q,x'),qq'\bigr),
  \qquad
  \alpha(q,x')=qx'q^{-1},
\]
multiplication in \(\mathbf H(G)\) is given, for
\(x\in\mathbf H(N)\) and \(q\in\mathbf H(Q)\), by
\[
  (x\otimes q)\cdot(x'\otimes q')
  =
  x\phi(q,x')\otimes qq',
\]
where
\[
  \phi:
  \mathbf H(Q)\otimes_k\mathbf H(N)
  \longrightarrow
  \mathbf H(N)
\]
is the morphism of \(k\)-coalgebras induced by \(\alpha\).

The morphism \(\alpha\) is the following composite. Here
\(\delta_G\) and \(m_G\) denote, respectively, the diagonal and
multiplication morphisms of \(G\), \(c_Q\) is inversion in \(Q\),
and
\[
  v(q\otimes q'\otimes x)=q\otimes x\otimes q'.
\]
\begin{center}
  \phantomsection
  \label{sga3:VIIB:2.4.B:conjugation-composite}
  % SGA3-VII-DIAG-111
  % BEGIN INCLUDED SGA3_English_current_buildable_source_20260823/source/inputs/VII/native_diagrams/exp7/viib_15_95/diag_111.tex
% SGA3-VII-DIAG-111
% Authority: Exp7B-13oct24.pdf, local p. 54.
% Crop: figures/Exp7B_localp054_semidirect_conjugation_composite.png
% Locator: tex/components/44_expose_VIIB_section24B_covariant_semidirect_products_en.tex, Loop 1 line 64.
\(
Q\times N
\xrightarrow{\ v\circ(\delta_G\times\operatorname{id})\ }
Q\times N\times Q
\xrightarrow{\ \operatorname{id}\times\operatorname{id}\times c_Q\ }
Q\times N\times Q
\xrightarrow{\ m_G\ }G,
\)

% END INCLUDED SGA3_English_current_buildable_source_20260823/source/inputs/VII/native_diagrams/exp7/viib_15_95/diag_111.tex

\end{center}
Writing \(c_Q\) also for the antipode of \(\mathbf H(Q)\), we obtain
\[
  \phi(q\otimes x')
  =
  \sum_i q_i x'c_Q(q_i')
  \quad\text{if}\quad
  \Delta_{\mathbf H(Q)}(q)
  =
  \sum_i q_i\otimes q_i'.
  \tag{\(\star\)}\label{sga3:VIIB:2.4.B:equation:covariant-antipode-star}
\]

In particular, let \(M\) be an abstract group and suppose that
\(\mathbf H(Q)\) is the group algebra \(kM\), so that
\[
  Q=\operatorname{Spf}^{*}kM
\]
is the constant formal \(k\)-group \(M_k\). Then for every
\(\gamma\in M\) and \(x'\in\mathbf H(N)\),
\[
  \phi(\gamma\otimes x')
  =
  \gamma x'\gamma^{-1}.
\]
This defines an action of \(M\) on \(\mathbf H(N)\) by Hopf-algebra
automorphisms.

% END INCLUDED SGA3_English_current_buildable_source_20260823/source/inputs/VII/tex/components/44_expose_VIIB_section24B_covariant_semidirect_products_en.tex

% BEGIN INCLUDED SGA3_English_current_buildable_source_20260823/source/inputs/VII/tex/components/45_expose_VIIB_proposition241_kernel_quotient_image_en.tex
% SGA 3, Expose VII_B, English reconstruction -- Loop 1.
% Scope: authority local pages 54--55, Proposition 2.4.1, its proof,
% and its corollary in full, stopping immediately before 2.4.2.
% Combined-reader pages: 567--568. Re-edition printed pages: 556--557.
% Original-edition marginal locator in scope: 521.
% Editorial note in scope: 93.
% Authority: Polo--Gille Exp7B-13oct24.pdf,
% SHA-256 D303690E0195654BBE7C3DAA5D89EA23E24CD709A3575217234AC19AE8836694.
% Jacob Reinhold's CC BY 4.0 Markdown was consulted as comparison lineage.
% The one genuine diagram is an exact PDF-derived Loop 1 source crop.

\medskip
\noindent\textbf{Proposition 2.4.1.}
\phantomsection
\label{sga3:VIIB:2.4.1}
{\itshape Let \(f:G\to K\) be a morphism of formal \(k\)-groups. If
\[
  H=\ker(f)
\]
is topologically flat over \(k\), the homomorphism
\[
  f':G/H\longrightarrow K
\]
induced by \(f\) is a monomorphism.}

This is a consequence of the results of Exposé IV,\footnote{%
\label{sga3:VIIB:nde:93}See \SGARefLink{sga3:iv:remark:5:2:6}{Exposé IV, 5.2.6}.} but we give a direct
proof. Let \(T\) be a formal variety of finite length over \(k\), and
let \(t\in(G/H)(T)\) be such that \(f'\circ t\) is the identity
element of \(K(T)\). We must show that \(t\) is the identity element
of \((G/H)(T)\). Let
\[
  p:G\longrightarrow G/H
\]
be the projection, and let
\[
  X=T\times_{G/H}G.
\]

By \SGARefLink{sga3:VIIB:2.4}{2.4}, \(p\) is surjective and topologically flat. The same is
therefore true of
\[
  \operatorname{pr}_1:X\longrightarrow T.
\]
By \SGARefLink{sga3:VIIB:1.3.1}{Proposition 1.3.1}, \(\operatorname{pr}_1\) is consequently an
epimorphism, so it is enough to show that
\(t\circ\operatorname{pr}_1\) is the identity element of
\((G/H)(X)\).

Let \(\operatorname{pr}_2:X\to G\) be the second projection. We have
\[
  t\circ\operatorname{pr}_1
  =
  p\circ\operatorname{pr}_2,
\]
and hence
\[
  1
  =
  f'\circ t\circ\operatorname{pr}_1
  =
  f'\circ p\circ\operatorname{pr}_2
  =
  f\circ\operatorname{pr}_2.
\]
The exact sequence
\begin{center}
  \phantomsection
  \label{sga3:VIIB:2.4.1:exact-sequence}
  % SGA3-VII-DIAG-112
  % BEGIN INCLUDED SGA3_English_current_buildable_source_20260823/source/inputs/VII/native_diagrams/exp7/viib_15_95/diag_112.tex
% SGA3-VII-DIAG-112
% Authority: Exp7B-13oct24.pdf, local p. 54.
% Crop: figures/Exp7B_localp054_proposition241_exact_sequence.png
% Locator: tex/components/45_expose_VIIB_proposition241_kernel_quotient_image_en.tex, Loop 1 line 70.
\(
1\longrightarrow H\longrightarrow G\xrightarrow{\ f\ }K
\)

% END INCLUDED SGA3_English_current_buildable_source_20260823/source/inputs/VII/native_diagrams/exp7/viib_15_95/diag_112.tex

\end{center}
therefore shows that \(\operatorname{pr}_2\) factors through \(H\).
Thus \(p\circ\operatorname{pr}_2\) is the zero morphism. Since
\[
  p\circ\operatorname{pr}_2
  =
  t\circ\operatorname{pr}_1,
\]
this proves the proposition.

The following corollary follows. Write
\(\mathcal O(G)\), \(\mathcal O(K)\), and \(\mathcal O(G/H)\) for
the affine algebras of \(G\), \(K\), and \(G/H\), respectively. We
saw in \SGARefLink{sga3:VIIB:2.4}{2.4} that \(p\) induces an injection
\[
  \mathcal O(G/H)\lhook\joinrel\longrightarrow\mathcal O(G).
\]
Moreover, since \(f':G/H\to K\) is a monomorphism, \SGARefLink{sga3:VIIB:1.3.prop}{Proposition 1.3}
shows that
\[
  \mathcal O(K)\longrightarrow\mathcal O(G/H)
\]
is surjective. Hence:

\medskip
\noindent\textbf{Corollary.}
\phantomsection
\label{sga3:VIIB:2.4.1:corollary}
{\itshape Let \(f:G\to K\) be a morphism of formal \(k\)-groups and
let \(H=\ker(f)\). If \(H\) is topologically flat over \(k\), then
\[
  \mathcal O(G/H)
\]
is the image of \(\mathcal O(K)\) in \(\mathcal O(G)\).}

% END INCLUDED SGA3_English_current_buildable_source_20260823/source/inputs/VII/tex/components/45_expose_VIIB_proposition241_kernel_quotient_image_en.tex

% BEGIN INCLUDED SGA3_English_current_buildable_source_20260823/source/inputs/VII/tex/components/46_expose_VIIB_section242_abelian_categories_en.tex
% SGA 3, Expose VII_B, English reconstruction -- Loop 1.
% Scope: authority local page 55, Section 2.4.2 in full, stopping
% immediately before Section 2.5.
% Combined-reader page: 568. Re-edition printed page: 557.
% Original-edition marginal locator in scope: 522.
% Editorial note in scope: 94.
% Authority: Polo--Gille Exp7B-13oct24.pdf,
% SHA-256 D303690E0195654BBE7C3DAA5D89EA23E24CD709A3575217234AC19AE8836694.
% Jacob Reinhold's CC BY 4.0 Markdown was consulted as comparison lineage.
% No genuine diagram occurs in this span.

\subsubsection*{2.4.2}
\phantomsection
\label{sga3:VIIB:2.4.2}

Retain the preceding notation and suppose that \(H\) and \(G\) are
topologically flat over \(k\). Then \(G\) is topologically flat both
over \(k\) and over \(G/H\); hence, by \SGARefLink{sga3:VIIB:1.3.3}{1.3.3}, \(G/H\) is
topologically flat over \(k\). Consequently, by \SGARefLink{sga3:VIIB:2.4}{2.4}, the canonical
projection
\[
  q:K\longrightarrow\operatorname{Coker}(f')
\]
is topologically flat, and \(f'\) is an isomorphism from \(G/H\) onto
\(\ker(q)\). It is also clear that
\[
  \operatorname{Coker}(f)
  =
  \operatorname{Coker}(f').
\]
Thus, under the hypothesis that \(G\) and \(H=\ker(f)\) are
topologically flat over \(k\), we have obtained an isomorphism between
\(\ker(q)\), the image of \(f\), and
\[
  G/\ker(f),
\]
the coimage of \(f\). This proves the following theorem.

\medskip
\noindent\textbf{Theorem.}
\phantomsection
\label{sga3:VIIB:2.4.2:theorem}
{\itshape Let \(k\) be a field. The commutative formal \(k\)-groups
form an abelian category.}

\medskip
\noindent\textbf{Corollary.}
\phantomsection
\label{sga3:VIIB:2.4.2:corollary}
{\itshape Let \(k\) be a field. The commutative affine
\(k\)-group schemes form an abelian category.}
\footnote{\label{sga3:VIIB:nde:94}See also the remarks following
Exposé VI\textsubscript{A}, \SGARefLink{cor:sga3-VIA-5.4.3}{5.4.3}.}

This follows from the theorem and the equivalence of categories in
\SGARefLink{sga3:VIIB:2.2.2}{2.2.2}.

% END INCLUDED SGA3_English_current_buildable_source_20260823/source/inputs/VII/tex/components/46_expose_VIIB_section242_abelian_categories_en.tex

% BEGIN INCLUDED SGA3_English_current_buildable_source_20260823/source/inputs/VII/tex/components/47_expose_VIIB_section25_etale_formal_groups_en.tex
% SGA 3, Expose VII_B, English reconstruction -- Loop 1.
% Scope: authority local pages 55--56, Section 2.5 through its main
% proposition and Remark 2.5.A, stopping immediately before 2.5.1.
% Combined-reader pages: 568--569. Re-edition printed pages: 557--558.
% Original-edition marginal locator in scope: 522.
% Editorial notes in scope: 95--96.
% Authority: Polo--Gille Exp7B-13oct24.pdf,
% SHA-256 D303690E0195654BBE7C3DAA5D89EA23E24CD709A3575217234AC19AE8836694.
% Jacob Reinhold's CC BY 4.0 Markdown was consulted as comparison lineage.
% No genuine diagram occurs in this span.

\subsection*{2.5}
\phantomsection
\label{sga3:VIIB:2.5}
\addcontentsline{toc}{subsection}{2.5. Étale formal groups}

A formal \(k\)-group is said to be \emph{étale} if its underlying
formal variety is étale \textup{(cf.\ \SGARefLink{sga3:VIIB:1.6}{1.6})}. Such formal groups have
a very simple structure. Suppose that \(k\) is local, let \(\kappa\)
be its residue field, let \(\kappa_s\) be a separable closure of
\(\kappa\), and let
\[
  \Gamma=\operatorname{Gal}(\kappa_s/\kappa)
\]
be the topological Galois group. A \emph{\(\Gamma\)-set} means a set
\(E\) with a continuous action of \(\Gamma\); equivalently, the
stabilizer of every element \(x\in E\) is an open subgroup of
\(\Gamma\).

For every formal \(k\)-variety \(X\), put
\[
  X(\kappa_s)
  =
  \varinjlim_{\ell}X(\ell),
\]
where \(\ell\) ranges over the finite extensions of \(\kappa\)
contained in \(\kappa_s\).\footnote{%
\label{sga3:VIIB:nde:95}If
\([\kappa_s:\kappa]=\infty\), then \(\kappa_s\) is not a profinite
\(\kappa\)-algebra, so the preceding notation is an abuse. The
editors also supplied the details that follow.}
The group \(\Gamma\) acts continuously on every \(X(\ell)\), and
hence on \(X(\kappa_s)\). Moreover, if
\[
  X_\kappa=X\widehat\otimes_k\kappa
  \qquad\text{\textup{(cf.\ 1.6.C)}},
\]
then \(X(\ell)=X_\kappa(\ell)\) for every \(\ell\), and therefore
\[
  X(\kappa_s)=X_\kappa(\kappa_s).
\]

Now suppose that \(X\) is étale over \(k\). Then \(X_\kappa\) is the
formal \(\kappa\)-variety that is the disjoint union of the
\(\operatorname{Spec}\kappa(x)\), for \(x\in X\). If
\(X'_\kappa\) denotes the \(\kappa\)-scheme given by the same
disjoint union, then
\[
\begin{aligned}
  X_\kappa(\kappa_s)
  &=
  X'_\kappa(\kappa_s)
  \\
  &=
  \operatorname{Hom}_{\operatorname{Sch}/\kappa}
  \bigl(\operatorname{Spec}\kappa_s,X'_\kappa\bigr).
\end{aligned}
\]

Let \(\mathcal C\) be the full subcategory
\[
  \operatorname{Sch}^{\mathrm{\acute{e}t}}/\kappa
  \subset
  \operatorname{Sch}/\kappa
\]
of étale \(\kappa\)-schemes. It is known that
\[
  X'\longmapsto X'(\kappa_s)
\]
is an equivalence from \(\mathcal C\) onto the category
\(\mathcal C'\) of \(\Gamma\)-sets
\textup{(cf.\ SGA 1, V, §§7--8, or \([DG70]\), §I.4, 6.4)}.
It therefore induces an equivalence between the category of group
objects in \(\mathcal C\) and the category of group objects in
\(\mathcal C'\). A group object in \(\mathcal C'\) is simply an
abstract group endowed with a continuous action of \(\Gamma\) by
group automorphisms; we call this a \emph{\(\Gamma\)-group}.

Taking into account the equivalences of \SGARefLink{sga3:VIIB:1.6.E}{1.6.E},
\[
  \operatorname{Vaf}^{\mathrm{\acute{e}t}}/k
  \xrightarrow{\ \sim\ }
  \operatorname{Vaf}^{\mathrm{\acute{e}t}}/\kappa
  \xrightarrow{\ \sim\ }
  \operatorname{Sch}^{\mathrm{\acute{e}t}}/\kappa,
\]
we obtain the following.

\medskip
\noindent\textbf{Proposition.}
\phantomsection
\label{sga3:VIIB:2.5:proposition}
{\itshape Let \(k\) be a local pseudocompact ring, \(\kappa\) its
residue field, \(\kappa_s\) a separable closure of \(\kappa\), and
\(\Gamma=\operatorname{Gal}(\kappa_s/\kappa)\).
\begin{enumerate}
\item[\textup{(i)}]
The functor \(X\mapsto X(\kappa_s)\) is an equivalence from the
category of étale formal \(k\)-varieties onto the category of
\(\Gamma\)-sets.

\item[\textup{(ii)}]
It induces an equivalence from the category of étale formal
\(k\)-groups onto the category of \(\Gamma\)-groups.
\end{enumerate}}

\medskip
\noindent\textbf{Remark 2.5.A.}
\phantomsection
\label{sga3:VIIB:2.5.A}
\footnote{\label{sga3:VIIB:nde:96}The editors added this remark.}
Let \(k\) be a field, let \(k_s\) be a separable closure of \(k\),
let \(G\) be an étale formal \(k\)-group, and put
\[
  M=G(k_s).
\]
Choose a set \(X\) of representatives for the orbits of
\(\Gamma=\operatorname{Gal}(k_s/k)\) in \(M\). For each \(x\in X\),
let \(\Gamma_x\) be its stabilizer, an open subgroup of finite index
in \(\Gamma\), and put
\[
  L_x=k_s^{\Gamma_x}.
\]
Then \(L_x/k\) has degree \([\Gamma:\Gamma_x]\)
\textup{(see, for example, \([BAlg]\), V, §10.10)}.

By the preceding equivalence of categories, the affine algebra
\(\mathcal A(G)\) of \(G\) is the product of the \(L_x\), endowed
with the product topology. The \(L_x\) are therefore exactly the
simple quotients
\[
  \mathcal A(G)/\mathfrak m,
\]
where \(\mathfrak m\) ranges over the open maximal ideals of
\(\mathcal A(G)\). Since these quotients correspond by duality to
the simple subcoalgebras of \(\mathbf H(G)\), the coalgebra
\(\mathbf H(G)\) is pointed---that is,
\[
  \dim_kC=1
\]
for every simple subcoalgebra \(C\)---if and only if \(L_x=k\) for
every \(x\). In that case, \(\mathcal A(G)\) is the topological
algebra \(k^M\), so \(\mathbf H(G)\) is the group algebra \(kM\), and
\[
  M
  =
  \bigl\{
    x\in\mathbf H(G)
    \mathrel{\big|}
    \varepsilon(x)=1,\ \Delta(x)=x\otimes x
  \bigr\}
  =
  G(k).
\]

% END INCLUDED SGA3_English_current_buildable_source_20260823/source/inputs/VII/tex/components/47_expose_VIIB_section25_etale_formal_groups_en.tex

% BEGIN INCLUDED SGA3_English_current_buildable_source_20260823/source/inputs/VII/tex/components/48_expose_VIIB_section251_infinitesimal_identity_component_en.tex
% SGA 3, Expose VII_B, English reconstruction -- Loop 1.
% Scope: authority local pages 56--57, Section 2.5.1 in full,
% stopping immediately before 2.5.2.
% Combined-reader pages: 569--570. Re-edition printed pages: 558--559.
% Original-edition marginal locator in scope: 523.
% Editorial notes in scope: 97--99.
% Authority: Polo--Gille Exp7B-13oct24.pdf,
% SHA-256 D303690E0195654BBE7C3DAA5D89EA23E24CD709A3575217234AC19AE8836694.
% Jacob Reinhold's CC BY 4.0 Markdown was consulted as comparison lineage.
% The three genuine diagrams are exact PDF-derived Loop 1 source crops.

\subsubsection*{2.5.1}
\phantomsection
\label{sga3:VIIB:2.5.1}

Suppose again that \(k\) is an arbitrary pseudocompact ring. Since
the functor
\[
  X\longmapsto X_e
\]
of \SGARefLink{sga3:VIIB:1.6.F}{1.6.F} commutes with finite products, it sends every formal group
\(G\) to an étale formal group \(G_e\). The morphism
\(p:X\to X_e\) of that result is functorial in \(X\), so in this
case
\[
  p:G\longrightarrow G_e
\]
is a morphism of formal groups.

Consider its kernel \(\ker(p)\).\footnote{%
\label{sga3:VIIB:nde:97}The editors supplied the details of the
original argument that follow.} It is the fiber product of the
diagram
\begin{center}
  \phantomsection
  \label{sga3:VIIB:2.5.1:kernel-fiber-square}
  % SGA3-VII-DIAG-113
  % BEGIN INCLUDED SGA3_English_current_buildable_source_20260823/source/inputs/VII/native_diagrams/exp7/viib_15_95/diag_113.tex
% SGA3-VII-DIAG-113
% Authority: Exp7B-13oct24.pdf, local p. 56.
% Crop: figures/Exp7B_localp056_kernel_etale_fiber_square.png
% Locator: tex/components/48_expose_VIIB_section251_infinitesimal_identity_component_en.tex, Loop 1 line 37.
\begin{tikzcd}[row sep=2.5em,column sep=3.8em]
{} & G \arrow[d,"p"] \\
\operatorname{Spf}(k) \arrow[r,"\varepsilon"]
  & G_e
\end{tikzcd}

% END INCLUDED SGA3_English_current_buildable_source_20260823/source/inputs/VII/native_diagrams/exp7/viib_15_95/diag_113.tex

\end{center}
where \(\varepsilon\) is the identity section of \(G_e\). Since
\(p\) induces a bijection on the underlying sets, it follows
\textup{(cf.\ \SGARefLink{sga3:VIIB:1.2}{1.2}, editorial note 41)} that the underlying set of
\(\ker(p)\) is the image of \(\operatorname{Spf}(k)\) under
\(\varepsilon\), and that, for every point \(s\) of
\(\operatorname{Spf}(k)\), the local algebra of \(\ker(p)\) at
\(\varepsilon(s)\) is
\[
  \mathcal O_{G,\varepsilon(s)}
  \widehat\otimes_{\mathcal O_{G_e,\varepsilon(s)}}
  \kappa(s).
\]
At each point \(\varepsilon(s)\), the residue field of
\(\mathcal O_{G,\varepsilon(s)}\) is \(\kappa(s)\), and hence
\[
  \mathcal O_{G_e,\varepsilon(s)}=\kappa(s),
  \qquad
  \mathcal O_{\ker(p),\varepsilon(s)}
  =
  \mathcal O_{G,\varepsilon(s)}.
\]

For these reasons, we call \(\ker(p)\) the \emph{infinitesimal
neighborhood of the identity} of \(G\), and write
\[
  \ker(p)=G_0.
\]
We have thus obtained the exact sequence
\begin{center}
  \phantomsection
  \label{sga3:VIIB:2.5.1:infinitesimal-exact-sequence}
  % SGA3-VII-DIAG-114
  % BEGIN INCLUDED SGA3_English_current_buildable_source_20260823/source/inputs/VII/native_diagrams/exp7/viib_15_95/diag_114.tex
% SGA3-VII-DIAG-114
% Authority: Exp7B-13oct24.pdf, local p. 57.
% Crop: figures/Exp7B_localp057_infinitesimal_exact_sequence.png
% Locator: tex/components/48_expose_VIIB_section251_infinitesimal_identity_component_en.tex, Loop 1 line 71.
\(
1\longrightarrow G_0\xrightarrow{\ \mathrm{incl.}\ }G
\xrightarrow{\ p\ }G_e.
\)

% END INCLUDED SGA3_English_current_buildable_source_20260823/source/inputs/VII/native_diagrams/exp7/viib_15_95/diag_114.tex

\end{center}

In what follows, we say that \(G\) is \emph{infinitesimal} if
\[
  G=G_0.
\]
\footnote{\label{sga3:VIIB:nde:98}The editors added the sentence that
follows.}
This is equivalent to saying that, for every
\(s\in\operatorname{Spf}(k)\), or equivalently for every continuous
morphism \(k\to\kappa\), where \(\kappa\) is a field with the
discrete topology, the identity is the unique element of
\(G(\kappa(s))\), respectively of \(G(\kappa)\).

Suppose in addition that \(G\) is topologically flat over
\(k\).\footnote{\label{sga3:VIIB:nde:99}The editors supplied the
details that follow.} By \SGARefLink{sga3:VIIB:1.6.1}{1.6.1}, the morphism
\[
  p:G\longrightarrow G_e
\]
is topologically flat. Since it is bijective, \SGARefLink{sga3:VIIB:1.3.1}{Proposition 1.3.1}
shows that it is an effective epimorphism. Consequently, \(G_e\) is
identified with the quotient \(G/G_0\). We have proved:

\medskip
\noindent\textbf{Corollary.}
\phantomsection
\label{sga3:VIIB:2.5.1:corollary}
{\itshape Let \(G\) be a formal group topologically flat over \(k\).
Then \(G_e\) is identified with the quotient \(G/G_0\); in other
words, there is an exact sequence of formal groups}
\begin{center}
  \phantomsection
  \label{sga3:VIIB:2.5.1:short-exact-sequence}
  % SGA3-VII-DIAG-115
  % BEGIN INCLUDED SGA3_English_current_buildable_source_20260823/source/inputs/VII/native_diagrams/exp7/viib_15_95/diag_115.tex
% SGA3-VII-DIAG-115
% Authority: Exp7B-13oct24.pdf, local p. 57.
% Crop: figures/Exp7B_localp057_infinitesimal_etale_short_exact_sequence.png
% Locator: tex/components/48_expose_VIIB_section251_infinitesimal_identity_component_en.tex, Loop 1 line 107.
\(
1\longrightarrow G_0\xrightarrow{\ \mathrm{incl.}\ }G
\xrightarrow{\ p\ }G_e\longrightarrow1.
\)

% END INCLUDED SGA3_English_current_buildable_source_20260823/source/inputs/VII/native_diagrams/exp7/viib_15_95/diag_115.tex

\end{center}

% END INCLUDED SGA3_English_current_buildable_source_20260823/source/inputs/VII/tex/components/48_expose_VIIB_section251_infinitesimal_identity_component_en.tex

% BEGIN INCLUDED SGA3_English_current_buildable_source_20260823/source/inputs/VII/tex/components/49_expose_VIIB_section252_perfect_field_splitting_en.tex
% SGA 3, Expose VII_B, English reconstruction -- Loop 1.
% Scope: authority local pages 57--58, Section 2.5.2 and
% Remark 2.5.2.A in full, stopping immediately before 2.5.3.
% Combined-reader pages: 570--571. Re-edition printed pages: 559--560.
% Original-edition marginal locator in scope: 524.
% Editorial notes in scope: 100--102.
% Authority: Polo--Gille Exp7B-13oct24.pdf,
% SHA-256 D303690E0195654BBE7C3DAA5D89EA23E24CD709A3575217234AC19AE8836694.
% Jacob Reinhold's CC BY 4.0 Markdown was consulted as comparison lineage.
% The two genuine diagrams are exact PDF-derived Loop 1 source crops.

\subsubsection*{2.5.2}
\phantomsection
\label{sga3:VIIB:2.5.2}

Suppose that \(k\) is a perfect field. In this case, we saw in \SGARefLink{sga3:VIIB:1.6.2}{1.6.2}
that the morphism
\[
  p:X\longmapsto X_e
\]
has a section
\[
  s:X_e\longrightarrow X
\]
that depends functorially on \(X\). Consequently, when \(X\) is a
formal group, this section is a morphism of formal groups. We thus
obtain:

\medskip
\noindent\textbf{Proposition.}
\phantomsection
\label{sga3:VIIB:2.5.2:proposition}
{\itshape When \(k\) is a perfect field, every formal \(k\)-group
\(G\) is canonically isomorphic to the semidirect product of an
infinitesimal group \(G_0\) and an étale group \(G_e\) acting on
\(G_0\).}\footnote{%
\label{sga3:VIIB:nde:100}The same argument also shows that, for an
arbitrary field \(k\), if all the residue fields of \(G\) are equal
to \(k\), then \(G_e\) is the constant \(k\)-group \((M)_k\), where
\[
  M=G(k)=\operatorname{Spf}^{*}\mathbf H(G),
\]
and \(\mathbf H(G)\) is the semidirect product of
\(\mathbf H(G_0)\) by the group algebra \(kM\); see the addition in
\SGARefLink{sga3:VIIB:2.9}{2.9} below.}

If, moreover, \(G\) is commutative, then \(G\) is the product of
\(G_0\) and \(G_e\). By \SGARefLink{sga3:VIIB:2.2.2}{2.2.2}, this canonical decomposition of
commutative formal \(k\)-groups corresponds to an analogous
decomposition of affine commutative \(k\)-group schemes; see \SGARefLink{sga3:VIIB:2.5.3}{2.5.3}
below.

\medskip
\noindent\textbf{Remark 2.5.2.A.}
\phantomsection
\label{sga3:VIIB:2.5.2.A}
\footnote{\label{sga3:VIIB:nde:101}The editors added this remark.}
When \(k\) is not perfect, the exact sequence
\[
  1\longrightarrow G_0\longrightarrow G\longrightarrow G_e
  \longrightarrow 1
\]
need not split. Consider the following example, taken from
\([DG70]\), §III.6, 8.6. Let \(k\) be a nonperfect field of
characteristic \(p>0\), choose
\[
  \lambda\in k\setminus k^p,
\]
and let \(L_\lambda\) be the abelian restricted \(p\)-Lie algebra
with basis \((x,y)\), whose symbolic \(p\)-th power
\textup{(cf.\ VII\textsubscript{A}, \SGARefLink{sga3:VIIA:5.2}{5.2})} is given by
\[
  x^{(p)}=x,
  \qquad
  y^{(p)}=\lambda x.
\]
Let
\[
  U_p(L_\lambda)
  =
  k[x,y]/(x^p-x,\ y^p-\lambda x)
\]
be the restricted enveloping algebra of \(L_\lambda\)
\textup{(cf.\ VII\textsubscript{A}, \SGARefLink{sga3:VIIA:5.3}{5.3})}, and let \(G_\lambda\)
be the commutative formal \(k\)-group whose affine algebra is
\(U_p(L_\lambda)\).

Then \(G_\lambda\) is a nonsplit extension of the constant étale
\(k\)-group \((\mathbb Z/p\mathbb Z)_k\) by the infinitesimal
\(k\)-group
\[
  \alpha_{p,k}=\operatorname{Spf}k[x]/(x^p).
\]
\phantomsection
\label{sga3:VIIB:2.5.2.A:alpha-p-source-note}
\noindent\textit{Source note.}\enspace
The French re-edition prints
\(\operatorname{Spf}k[x]/(x^p-x)\) here. That algebra is étale,
whereas the asserted kernel is the standard infinitesimal group
\(\alpha_{p,k}=\operatorname{Spf}k[x]/(x^p)\). Indeed, the kernel
fiber \(x=0\) in \(U_p(L_\lambda)\) has coordinate algebra
\(k[y]/(y^p)\). The English text uses the standard infinitesimal
reading while preserving the printed reading in this note.
Thus there is a nonsplit exact sequence of commutative formal
\(k\)-groups:\footnote{%
\label{sga3:VIIB:nde:102}The groups
\(\alpha_{p,k}\), \((\mathbb Z/p\mathbb Z)_k\), and \(G_\lambda\)
are also finite flat algebraic \(k\)-groups; see editorial note 84
in \SGARefLink{sga3:VIIB:2.2.2}{2.2.2}.}
\begin{center}
  \phantomsection
  \label{sga3:VIIB:2.5.2.A:nonsplit-formal-sequence}
  % SGA3-VII-DIAG-116
  % BEGIN INCLUDED SGA3_English_current_buildable_source_20260823/source/inputs/VII/native_diagrams/exp7/viib_15_95/diag_116.tex
% SGA3-VII-DIAG-116
% Authority: Exp7B-13oct24.pdf, local p. 58.
% Crop: figures/Exp7B_localp058_nonsplit_formal_group_exact_sequence.png
% Locator: tex/components/49_expose_VIIB_section252_perfect_field_splitting_en.tex, Loop 1 line 113.
\(
0\longrightarrow\boldsymbol{\alpha}_{p,k}
\longrightarrow G_\lambda
\longrightarrow(\mathbb Z/p\mathbb Z)_k
\longrightarrow0.
\)

% END INCLUDED SGA3_English_current_buildable_source_20260823/source/inputs/VII/native_diagrams/exp7/viib_15_95/diag_116.tex

\end{center}

By duality, this corresponds to a nonsplit exact sequence of
commutative algebraic \(k\)-groups:
\begin{center}
  \phantomsection
  \label{sga3:VIIB:2.5.2.A:nonsplit-dual-sequence}
  % SGA3-VII-DIAG-117
  % BEGIN INCLUDED SGA3_English_current_buildable_source_20260823/source/inputs/VII/native_diagrams/exp7/viib_15_95/diag_117.tex
% SGA3-VII-DIAG-117
% Authority: Exp7B-13oct24.pdf, local p. 58.
% Crop: figures/Exp7B_localp058_nonsplit_dual_algebraic_group_exact_sequence.png
% Locator: tex/components/49_expose_VIIB_section252_perfect_field_splitting_en.tex, Loop 1 line 122.
\(
0\longrightarrow\boldsymbol{\mu}_{p,k}
\longrightarrow D'(G_\lambda)
\longrightarrow\boldsymbol{\alpha}_{p,k}
\longrightarrow0
\)

% END INCLUDED SGA3_English_current_buildable_source_20260823/source/inputs/VII/native_diagrams/exp7/viib_15_95/diag_117.tex

\end{center}
where
\[
  \mu_{p,k}=\operatorname{Spec}k[t]/(t^p-1).
\]
In this way one obtains all extensions of \(\alpha_{p,k}\) by
\(\mu_{p,k}\); see \([DG70]\), §III.6, 8.6.

% END INCLUDED SGA3_English_current_buildable_source_20260823/source/inputs/VII/tex/components/49_expose_VIIB_section252_perfect_field_splitting_en.tex

% BEGIN INCLUDED SGA3_English_current_buildable_source_20260823/source/inputs/VII/tex/components/50_expose_VIIB_section253_unipotent_multiplicative_type_en.tex
% SGA 3, Expose VII_B, English reconstruction -- Loop 1.
% Scope: authority local pages 58--59, Section 2.5.3 in full,
% stopping immediately before 2.6.
% Combined-reader pages: 571--572. Re-edition printed pages: 560--561.
% Original-edition marginal locator in scope: 524.
% Editorial notes in scope: 103--106.
% Authority: Polo--Gille Exp7B-13oct24.pdf,
% SHA-256 D303690E0195654BBE7C3DAA5D89EA23E24CD709A3575217234AC19AE8836694.
% Jacob Reinhold's CC BY 4.0 Markdown was consulted as comparison lineage.
% No genuine diagram occurs in this span.

\subsubsection*{2.5.3}
\phantomsection
\label{sga3:VIIB:2.5.3}

Let \(k\) be a field.\footnote{\label{sga3:VIIB:nde:103}The editors
introduced the numbering \SGARefLink{sga3:VIIB:2.5.3}{2.5.3} and \SGARefLink{sga3:VIIB:2.5.3.A}{2.5.3.A}--\SGARefLink{sga3:VIIB:2.5.3.C}{2.5.3.C}.}

\medskip
\noindent\textbf{Definition 2.5.3.A.}
\phantomsection
\label{sga3:VIIB:2.5.3.A}
A commutative \(k\)-group scheme is said to be \emph{unipotent} if
it is isomorphic to
\[
  \operatorname{Spec}\mathbf H(G),
\]
where \(G\) is an infinitesimal commutative formal
\(k\)-group.\footnote{%
\label{sga3:VIIB:nde:104}This agrees with the ``usual'' notion of a
unipotent \(k\)-group scheme; cf.\ the addition in \SGARefLink{sga3:VIIB:2.8}{2.8} at the end of
\SGARefLink{sga3:VIIB:2}{Section 2}.}

On the other hand, following \SGARefLink{sga3:ix:definition:1:1}{Exposé IX, Definition
1.1},\footnote{\label{sga3:VIIB:nde:105}The original stated:
``Let us say that a commutative \(k\)-group scheme is of
multiplicative type if it is isomorphic to
\(\operatorname{Spec}\mathbf H(G)\), where \(G\) is an étale
commutative formal \(k\)-group.'' The editors have given here the
``usual'' definition, taken from \SGARefLink{sga3:ix:definition:1:1}{Exposé IX, 1.1}, and proved its
equivalence with the preceding condition; see also \([DG70]\),
§IV.1, \SGARefLink{sga3:VIIB:2.2}{Theorem 2.2}.}
a \(k\)-group scheme \(H\) is said to be \emph{of multiplicative
type} if there exists a scheme \(S\), faithfully flat and
quasi-compact over \(\operatorname{Spec}k\), such that \(H_S\) is a
diagonalizable \(S\)-group; that is,
\[
  H_S\simeq D_S(M)
\]
for some ``abstract'' abelian group \(M\). By fpqc descent, this
implies that \(H\) is affine and commutative. Moreover, since \(S\)
may be replaced by the residue field of one of its points, \(H\) is
of multiplicative type if and only if there exists an extension
\(K\) of \(k\) such that \(H_K\) is a diagonalizable \(K\)-group.

\medskip
\noindent\textbf{Proposition 2.5.3.B.}
\phantomsection
\label{sga3:VIIB:2.5.3.B}
{\itshape Let \(T\) be an affine commutative \(k\)-group scheme, and
let \(k_s\) be a separable closure of \(k\).
\begin{enumerate}
\item[\textup{(i)}]
The scheme \(T\) is of multiplicative type if and only if its dual
\(D(T)\) is an étale commutative formal \(k\)-group.

\item[\textup{(ii)}]
Consequently, \(T\) is of multiplicative type if and only if
\[
  T\otimes_k k_s
\]
is diagonalizable.
\end{enumerate}}

\textit{Proof.}
Let \(\mathcal A\) denote the affine algebra of \(D(T)\), and let
\(\mathcal A_0\) be its local component at the identity. Then
\(D(T)\) is étale over \(k\) if and only if
\[
  \mathcal A_0=k.
\]
If \(K\) is an extension of \(k\), then
\(\mathcal A_0\widehat\otimes_k K\) is local, being an inverse limit
of Artinian local rings, and hence coincides with the local
component
\[
  \bigl(\mathcal A\widehat\otimes_k K\bigr)_0.
\]
Moreover, since the formation of \(D(T)\) commutes with base change
\textup{(cf.\ \SGARefLink{sga3:VIIB:2.2.2}{2.2.2})}, one also has
\[
  \mathcal A\widehat\otimes_k K
  \simeq
  \mathcal A\bigl(D(T_K)\bigr).
\]
\noindent\begin{minipage}{\linewidth}
\phantomsection
\label{sga3:VIIB:2.5.3.B:base-change-source-note}
\textit{Source note.}\enspace
The French re-edition prints
\(\mathcal A\widehat\otimes_k K\simeq D(T_K)\). The left-hand side
is an algebra and the right-hand side a formal group. The English
therefore uses the type-correct intended reading above and preserves
the printed formula here.
\end{minipage}

Suppose that \(T\) is of multiplicative type. Then there exists an
extension \(K\) of \(k\) such that
\[
  \mathcal O(T)\otimes_k K
\]
is isomorphic, as a Hopf \(K\)-algebra, to the group algebra \(KM\)
for some abelian group \(M\). It follows that
\(\mathcal A\widehat\otimes_k K\) is isomorphic to the algebra
\(K^M\), endowed with the product topology. Thus
\[
  K
  =
  \bigl(\mathcal A\widehat\otimes_k K\bigr)_0
  =
  \mathcal A_0\widehat\otimes_k K.
\]
This implies that \(\mathcal A_0=k\), and hence that \(D(T)\) is
étale.

Conversely, suppose that \(D(T)\) is étale. Then
\[
  D(T)\widehat\otimes_k k_s
  =
  D(T_{k_s})
\]
is the constant \(k_s\)-group associated with a group \(M\).
Therefore its covariant bialgebra is the group algebra \(k_sM\)
\textup{(cf.\ \SGARefLink{sga3:VIIB:2.5.A}{Remark 2.5.A})}, and consequently
\[
  \mathcal O(T)\otimes_k k_s=k_sM.
\]
This proves that \(T\) is of multiplicative type and is split by the
extension \(k\to k_s\). The proposition follows.

By \SGARefLink{sga3:VIIB:2.2.2}{2.2.2}, \SGARefLink{sga3:VIIB:2.5.1:corollary}{2.5.1}, and \SGARefLink{sga3:VIIB:2.5}{2.5}, we obtain:

\medskip
\noindent\textbf{Corollary 2.5.3.C.}
\phantomsection
\label{sga3:VIIB:2.5.3.C}
{\itshape Let \(k\) be a field, and let \(G\) be an affine
commutative \(k\)-group scheme.
\begin{enumerate}
\item[\textup{(i)}]
The group \(G\) contains a subgroup \(G_m\) of multiplicative type
such that \(G/G_m\) is unipotent.

\item[\textup{(ii)}]
When \(k\) is perfect, there is in addition a canonical retraction
of \(G\) onto \(G_m\), so that \(G\) is the product of a unipotent
group and a group of multiplicative type.

\item[\textup{\SGAInlineRefTarget{sga3:VIIB:2.5.3.C:item:iii}(iii)}]
\footnote{\label{sga3:VIIB:nde:106}The editors added item
\textup{\SGARefLink{sga3:VIIB:2.5.3.C:item:iii}{(iii)}}, a consequence of \SGARefLink{sga3:VIIB:2.5}{Proposition 2.5}.}
Let \(k_s\) be a separable closure of \(k\), and put
\[
  \Gamma=\operatorname{Gal}(k_s/k).
\]
The category of \(k\)-group schemes of multiplicative type is
anti-equivalent to the category of continuous \(\Gamma\)-modules.
\end{enumerate}}

% END INCLUDED SGA3_English_current_buildable_source_20260823/source/inputs/VII/tex/components/50_expose_VIIB_section253_unipotent_multiplicative_type_en.tex

% BEGIN INCLUDED SGA3_English_current_buildable_source_20260823/source/inputs/VII/tex/components/51_expose_VIIB_section26_lie_algebra_opening_en.tex
% SGA 3, Expose VII_B, English reconstruction -- Loop 1.
% Scope: authority local pages 59--61, Section 2.6 through Section
% 2.6.1 in full, beginning after Corollary 2.5.3.C and stopping
% immediately before 2.6.2.
% Combined-reader pages: 572--574. Re-edition printed pages: 561--563.
% Original-edition marginal locators in scope: 525--527.
% Editorial notes in scope: 107--109.
% Authority: Polo--Gille Exp7B-13oct24.pdf,
% SHA-256 D303690E0195654BBE7C3DAA5D89EA23E24CD709A3575217234AC19AE8836694.
% Jacob Reinhold's CC BY 4.0 Markdown was consulted as comparison
% lineage only and never as authority.
% The one genuine diagram is an exact PDF-derived Loop 1 source crop;
% all other mathematics is native LaTeX.

\subsection*{2.6}
\phantomsection
\label{sga3:VIIB:2.6}
\addcontentsline{toc}{subsection}{\SGARefLink{sga3:VIIB:2.6}{2.6}. Lie algebras of formal groups}

We shall now study infinitesimal formal groups, to which the sections
that follow are devoted. In this study, Lie algebras play a fundamental
role.

Suppose first that the base ring \(k\) is Artinian, and let \(G\) be a
formal group over \(k\). The Lie algebra of \(G\) admits three different
interpretations, all of which we shall use:

\medskip
\noindent\textbf{a)}
Let
\[
  D=k[d]/(d^2)
\]
be the algebra of dual numbers over \(k\), and let
\[
  \delta:D\longrightarrow k
\]
be the homomorphism that annihilates \(d\). For every formal group
\(G\) over \(k\), \(\operatorname{Lie}(G)\) is the kernel of
\(G(\delta)\). Thus, by definition, there is an exact sequence of
groups
\begin{center}
  \phantomsection
  \label{sga3:VIIB:2.6:dual-numbers-exact-sequence}
  % SGA3-VII-DIAG-118
  % BEGIN INCLUDED SGA3_English_current_buildable_source_20260823/source/inputs/VII/native_diagrams/exp7/viib_15_95/diag_118.tex
% SGA3-VII-DIAG-118
% Authority: Exp7B-13oct24.pdf, local p. 59.
% Crop: figures/Exp7B_localp059_lie_dual_numbers_exact_sequence.png
% Locator: tex/components/51_expose_VIIB_section26_lie_algebra_opening_en.tex, Loop 1 line 45.
\(
1\longrightarrow\operatorname{Lie}(G)
\longrightarrow G(D)
\xrightarrow{\ G(\delta)\ }G(k)
\longrightarrow1.
\)

% END INCLUDED SGA3_English_current_buildable_source_20260823/source/inputs/VII/native_diagrams/exp7/viib_15_95/diag_118.tex

\end{center}

\medskip
\noindent\textbf{b)}
Let \(A\) be the affine algebra of \(G\), and let
\[
  I_A=\ker\varepsilon_A
\]
be its augmentation ideal. The elements of \(G(D)\) are the morphisms
of profinite \(k\)-algebras
\[
  f:A\longrightarrow D.
\]
The condition \(G(\delta)(f)=1\) is equivalent to
\[
  \delta\circ f=\varepsilon_A.
\]
Since
\[
  x-\varepsilon_A(x)\cdot1_A\in I_A
  \qquad\text{for every \(x\in A\),}
\]
this is equivalent to \(f(I_A)\subset k\cdot d\). Hence
\(\ker(f)\) contains \(I_A^2\), and therefore also its closure
\(\overline{I_A^2}\). Thus \(f\) induces a continuous linear map
\[
  f':
  I_A/\overline{I_A^2}
  =
  \omega_{G/k}
  \longrightarrow k
\]
such that, for every \(x\in A\),
\[
  f(x)
  =
  \varepsilon_A(x)\cdot1_D+f'(\bar x)\cdot d,
\]
where \(\bar x\) denotes the image of
\(x-\varepsilon_A(x)\cdot1_A\) in
\(I_A/\overline{I_A^2}\). It follows that the map
\[
  f\longmapsto f'
\]
defines a bijection from \(\operatorname{Lie}(G)\) onto the
topological dual
\[
  \omega_{G/k}^{\dagger}
\]
of \(\omega_{G/k}\) \textup{(cf.\ \SGARefLink{sga3:VIIB:0.2.2}{0.2.2})}.

This bijection respects the group structures. Indeed, let \(f\) and
\(g\) be two elements of \(\operatorname{Lie}(G)\). Their product
\(f\cdot g\) is the composite \(h\circ\Delta_A\), where
\[
  h:A\widehat\otimes A\longrightarrow D
\]
is defined by
\[
  h(b\widehat\otimes b')=f(b)g(b').
\]
If \(a\in I_A\), then, by \SGARefLink{sga3:VIIA:2.1}{2.1}\textup{\SGARefLink{sga3:vii:paragraph:5-1:equation:ii}{(ii)}},
\[
  \Delta_A(a)-a\widehat\otimes1-1\widehat\otimes a
  \in
  I_A\widehat\otimes I_A.
\]
Consequently,
\[
  (f\cdot g)(a)=f(a)+g(a)
\]
\textup{(cf.\ also II, \SGARefLink{sga3:ii:lemma:3-10}{3.10})}.

Henceforth we identify
\[
  \operatorname{Lie}(G)=\omega_{G/k}^{\dagger}
\]
by means of the bijection \(f\mapsto f'\) just described. The group
\(\operatorname{Lie}(G)\) is thereby endowed with a \(k\)-module
structure.

\medskip
\noindent\SGAInlineRefTarget{sga3:VIIB:2.6:item:c}\textbf{c)}
Let \(A^\dagger\) and \(D^\dagger\) be the \(k\)-modules dual to
\(A\) and \(D\), and let
\[
  \{1_D^\dagger,d^\dagger\}
\]
be the basis dual to the basis \(\{1_D,d\}\) of \(D\) over \(k\);
thus \(1_D^\dagger=\delta\).\footnote{The French authority prints
\(1_D^*\) and \(d^*\) in the two occurrences immediately below,
after defining \(1_D^\dagger\) and \(d^\dagger\) as the dual basis.
The consistent \(\dagger\)-notation is used here.}
Since \(D\) is free of finite rank over \(k\), the canonical map
\[
  \operatorname{Hom}_c(A,D)
  \longrightarrow
  \operatorname{Hom}_k(D^\dagger,A^\dagger),
  \qquad
  f\longmapsto{}^t f,
\]
is bijective. On the other hand, \({}^t f\) is determined by the
values
\[
  {}^t f(1_D^\dagger)
  \qquad\text{and}\qquad
  {}^t f(d^\dagger)=x.
\]
The condition \(G(\delta)(f)=1\) is equivalent to
\[
  {}^t f(1_D^\dagger)=\varepsilon_A.
\]
Moreover, one sees readily that \(f\) is compatible with
multiplication if and only if \textup{(cf.\ \SGARefLink{sga3:VIIB:2.3}{2.3})}
\[
  \tag{*}
  \label{sga3:VIIB:2.6:primitive-relation}
  \delta_G(x)=\phi_G(x\otimes1+1\otimes x).
\]

Finally, it is clear that a continuous linear map
\[
  f:A\longrightarrow D
\]
that is compatible with multiplication and satisfies
\(\delta\circ f=\varepsilon_A\) sends the unit element of \(A\) to
that of \(D\).\footnote{\label{sga3:VIIB:nde:107}Indeed,
\(\delta\circ f=\varepsilon_A\) implies
\(f(1)=1+\lambda d\), with \(\lambda\in k\), and then
\[
  f(1)=f(1)^2=1+2\lambda d
\]
gives \(\lambda=0\).}
The map \(f\mapsto x\) therefore allows us to identify
\(\operatorname{Lie}(G)\) with the set of ``primitive elements'' of
\(\mathbf H(G)\), that is, the elements \(x\in\mathbf H(G)\)
satisfying relation \textup{(*)}. If \(x\) and \(y\) are two such
elements, then
\[
\begin{aligned}
  \delta_G(xy)
  &=
  \delta_G(x)\delta_G(y)
  \\
  &=
  \phi_G\bigl(
    (x\otimes1+1\otimes x)
    (y\otimes1+1\otimes y)
  \bigr)
  \\
  &=
  \phi_G\bigl(
    xy\otimes1+x\otimes y+y\otimes x+1\otimes xy
  \bigr),
\end{aligned}
\]
whence
\[
  \delta_G(xy-yx)
  =
  \phi_G\bigl(
    (xy-yx)\otimes1+1\otimes(xy-yx)
  \bigr).
\]
This shows that the \(k\)-module \(\operatorname{Lie}(G)\) is
identified with a Lie subalgebra of \(\mathbf H(G)\); we call
\(\operatorname{Lie}(G)\) the \emph{Lie algebra} of
\(G\).\footnote{\label{sga3:VIIB:nde:108}Comparing with
VII\textsubscript{A}, \SGARefLink{sga3:VIIB:2.5:proposition}{2.5}, we see that if \(K\) is a group scheme of
finite type over \(k\) and \(G\) is the formal completion of \(K\) at
the identity---that is, if \(\mathcal A(G)\) is the completion of the
local ring \(\mathcal O_{K,e}\) for the \(\mathfrak m\)-adic topology,
where
\[
  \mathfrak m=\ker\bigl(
    \varepsilon:\mathcal O_{K,e}\longrightarrow k
  \bigr)
\]
---then \(\mathbf H(G)\) is identified with the algebra of
distributions \(U(K)\), and \(\operatorname{Lie}(G)\) with
\(\operatorname{Lie}(K)\). The hypothesis that \(K\) is of finite
type over \(k\) is used to ensure that
\(\mathfrak m/\mathfrak m^2\) is a \(k\)-module of finite length, and
hence discrete, so that its topological dual coincides with its
ordinary dual. In particular, when \(G\) is a finite formal
\(k\)-group---that is, when \(\mathcal A(G)\) is a finite
\(k\)-module---in which case \(G\) may also be regarded as the
\(k\)-group scheme \(\operatorname{Spec}\mathcal A(G)\), the two
definitions of \(\operatorname{Lie}(G)\) coincide.}

\subsubsection*{2.6.1}
\phantomsection
\label{sga3:VIIB:2.6.1}

Let \(k\) now be an arbitrary pseudocompact ring, and let \(G\) be a
formal group over \(k\). We call the functor \(\operatorname{Lie}(G)\)
that assigns to every object \(C\) of \(\operatorname{Alf}/k\) the
\(C\)-Lie algebra of the \(C\)-formal group
\[
  G'=G\widehat\otimes_k C
\]
the \emph{\(\mathbf O_k\)-Lie algebra of \(G\)}.\footnote{%
\label{sga3:VIIB:nde:109}The editors expanded the original in what
follows.}
Put
\[
  A'=A\widehat\otimes_k C.
\]
Since \(I_A\) is a direct factor of \(A\), we have
\[
  I_{A'}
  =
  I_A\widehat\otimes_k C
  =
  I_A\widehat\otimes_A A'.
\]
Moreover,
\[
  \omega_{G/k}=I_A\widehat\otimes_A k,
  \qquad
  \omega_{G'/C}=I_{A'}\widehat\otimes_{A'}C.
\]
It follows that
\[
  \omega_{G'/C}
  =
  \omega_{G/k}\widehat\otimes_A C,
\]
and hence
\[
\begin{aligned}
  \operatorname{Lie}(G)(C)
  &=
  \operatorname{Hom}_c\bigl(
    \omega_{G/k}\widehat\otimes_A C,C
  \bigr),
  \\
  \text{i.e.}\qquad
  \operatorname{Lie}(G)
  &=
  \mathbf V_k^f(\omega_{G/k}).
\end{aligned}
\]
With the notation of \SGARefLink{sga3:VIIB:1.2.3.B}{1.2.3.B}, \SGARefLink{sga3:VIIB:1.2.3.E}{Proposition 1.2.3.E} therefore shows
that \(\operatorname{Lie}(G)\) is flat over \(\mathbf O_k\) if and
only if \(\omega_{G/k}\) is a projective pseudocompact \(k\)-module.

% END INCLUDED SGA3_English_current_buildable_source_20260823/source/inputs/VII/tex/components/51_expose_VIIB_section26_lie_algebra_opening_en.tex

% BEGIN INCLUDED SGA3_English_current_buildable_source_20260823/source/inputs/VII/tex/components/52_expose_VIIB_section262_enveloping_formal_group_en.tex
% SGA 3, Expose VII_B, English reconstruction -- Loop 1.
% Scope: Section 2.6.2 in full on authority local pages 61--62,
% stopping immediately before Section 2.6.3.
% Combined-reader pages: 574--575. Re-edition printed pages: 563--564.
% Original-edition marginal locator in scope: 528.
% Editorial note in scope: 110.
% Authority: Polo--Gille Exp7B-13oct24.pdf,
% SHA-256 D303690E0195654BBE7C3DAA5D89EA23E24CD709A3575217234AC19AE8836694.
% Jacob Reinhold's CC BY 4.0 Markdown was consulted as comparison lineage.
% No genuine diagram occurs in this component; all mathematics is native LaTeX.

\subsubsection*{2.6.2}
\phantomsection
\label{sga3:VIIB:2.6.2}

Conversely, every Lie algebra \(\mathbf L\) over \(\mathbf O_k\)
defines a \(k\)-group functor. Indeed, let \(U(\mathbf L)\) denote
the functor assigning to every object \(C\) of
\(\operatorname{Alf}/k\) the enveloping algebra
\[
  U\bigl(\mathbf L(C)\bigr)
\]
of the \(C\)-Lie algebra \(\mathbf L(C)\). By
VII\textsubscript{A}, \SGARefLink{sga3:VIIA:3.2.2}{3.2.2}, \(U(\mathbf L)\) is a bialgebra over
\(\mathbf O_k\), and hence, by \SGARefLink{sga3:VIIA:2.2}{2.2}, determines the \(k\)-group
functor
\[
  \operatorname{Spf}^{*}U(\mathbf L),
\]
which from now on we denote by \(\mathscr G(\mathbf L)\). Thus
\(\mathscr G(\mathbf L)(C)\) is the group of elements
\(z\in U(\mathbf L(C))\) having augmentation \(1\) and satisfying
\[
  \Delta_{U(\mathbf L(C))}(z)=z\otimes z.
\]

When \(\mathbf L\) is flat over \(\mathbf O_k\), we have the
following proposition.

\medskip
\noindent\textbf{Proposition.}\footnote{%
\label{sga3:VIIB:nde:110}The editors set off parts \textup{\SGARefLink{sga3:vii:paragraph:5-1:equation:i}{(i)}} and
\textup{\SGARefLink{sga3:vii:paragraph:5-1:equation:ii}{(ii)}} and added part \textup{(iii)}, which will be used in
\SGARefLink{sga3:VIIB:2.6.3:prop}{2.6.3} and \SGARefLink{sga3:VIIB:3.3.2}{3.3.2}.}
\phantomsection
\label{sga3:VIIB:2.6.2:prop}
{\itshape Let \(\mathbf L\) be a flat Lie algebra over
\(\mathbf O_k\).
\begin{enumerate}
\item[\textup{\SGAInlineRefTarget{sga3:VIIB:2.6.2:prop:item:i}(i)}]
\(\mathscr G(\mathbf L)\) is a formal group topologically flat over
\(k\), whose \(\mathbf O_k\)-bialgebra is \(U(\mathbf L)\).

\item[\textup{(ii)}]
\(\mathscr G(\mathbf L)\) is infinitesimal.

\item[\textup{\SGAInlineRefTarget{sga3:VIIB:2.6.2:prop:item:iii}(iii)}]
For every object \(C\) of \(\operatorname{Alf}/k\),
\(\operatorname{Lie}(\mathscr G(\mathbf L))(C)\) is identified with
the set
\[
  \operatorname{Prim}U(\mathbf L(C))
  =
  \left\{
    x\in U(\mathbf L(C))
    \,\middle|\,
    \varepsilon(x)=0,\quad
    \Delta(x)=x\otimes1+1\otimes x
  \right\}
\]
of primitive elements of \(U(\mathbf L(C))\). In particular, there
is a natural morphism of Lie algebras over \(\mathbf O_k\),
\[
  \tau_{\mathbf L}:
  \mathbf L\longrightarrow
  \operatorname{Lie}\bigl(\mathscr G(\mathbf L)\bigr).
\]
\end{enumerate}}

\textit{Proof.}
The hypothesis that \(\mathbf L\) is flat over \(\mathbf O_k\)
means that, for every morphism \(C\to D\) in
\(\operatorname{Alf}/k\),
\[
  \mathbf L(D)=\mathbf L(C)\otimes_C D,
\]
and that, for every local component \(C'\) of \(C\), the
\(C'\)-module \(\mathbf L(C')\) is free. The first condition,
together with the universal properties of tensor products and
enveloping algebras, gives
\[
  U\bigl(\mathbf L(D)\bigr)
  =
  U\bigl(\mathbf L(C)\bigr)\otimes_C D.
\]
The second condition and the Poincaré--Birkhoff--Witt theorem
\textup{(cf.\ Bourbaki, \emph{Groupes et algèbres de Lie},
I, 2.7)} show that \(U(\mathbf L(C'))\) is a free \(C'\)-module.
Consequently, the bialgebra \(U(\mathbf L)\) is flat over
\(\mathbf O_k\), proving \textup{\SGARefLink{sga3:VIIB:2.6.2:prop:item:i}{(i)}}.

To prove \textup{\SGARefLink{sga3:vii:paragraph:5-1:equation:ii}{(ii)}} and \textup{\SGARefLink{sga3:VIIB:2.6.2:prop:item:iii}{(iii)}}, we may suppose that
\(k\) is Artinian. Put
\[
  L=\mathbf L(k),\qquad
  U=U(L),\qquad
  U_0=k\cdot1_U,
\]
and let \(U^+\) be the two-sided ideal of \(U\) generated by the
image of \(L\). For every \(n\geq1\), put
\[
  U_n
  =
  \left\{
    u\in U
    \,\middle|\,
    \Delta_U(u)-u\otimes1\in U_{n-1}\otimes U^+
  \right\}.
\]
By \SGARefLink{sga3:VIIB:1.3.6}{1.3.6}, it is enough to show that \(U\) is the union of the
\(U_n\). Identify \(L\) with its canonical image in \(U\). Then
\(L\subset U_1\). If \(x_1,\ldots,x_n\in L\), then
\[
  \Delta_U(x_1\cdots x_n)
  =
  (x_1\otimes1+1\otimes x_1)\cdots
  (x_n\otimes1+1\otimes x_n).
\]
Induction on \(n\) therefore shows that
\(x_1\cdots x_n\in U_n\), and hence
\[
  U=\bigcup_n U_n.
\]
This proves \textup{\SGARefLink{sga3:vii:paragraph:5-1:equation:ii}{(ii)}}.

Finally, let
\[
  D=k[d]/(d^2)
\]
be the algebra of dual numbers over \(k\). By hypothesis,
\[
  \mathbf L(D)\simeq L\otimes_kD,
  \qquad
  U(\mathbf L(D))\simeq U\otimes_kD.
\]
It follows that
\(\operatorname{Lie}(\mathscr G(\mathbf L))(k)\) is identified with
the set of elements
\[
  z=1+xd\in U\oplus Ud
  \qquad (x\in U)
\]
such that \(\varepsilon(z)=1\) and \(\Delta(z)=z\otimes z\).
These conditions are equivalent to
\[
  \varepsilon(x)=0,
  \qquad
  \Delta(x)=x\otimes1+1\otimes x,
\]
that is, to \(x\in\operatorname{Prim}U\). In particular,
\[
  \tau_{\mathbf L}:x\longmapsto1+dx
\]
is a morphism of Lie algebras over \(\mathbf O_k\) from
\(\mathbf L\) to
\(\operatorname{Lie}(\mathscr G(\mathbf L))\).


% END INCLUDED SGA3_English_current_buildable_source_20260823/source/inputs/VII/tex/components/52_expose_VIIB_section262_enveloping_formal_group_en.tex

% BEGIN INCLUDED SGA3_English_current_buildable_source_20260823/source/inputs/VII/tex/components/53_expose_VIIB_section263_universal_property_en.tex
% SGA 3, Expose VII_B, English reconstruction -- Loop 1.
% Scope: Section 2.6.3 in full on authority local page 62,
% stopping immediately before Section 2.7.
% Combined-reader page: 575. Re-edition printed page: 564.
% Original-edition marginal locator in scope: 529.
% Editorial note in scope: 111.
% Authority: Polo--Gille Exp7B-13oct24.pdf,
% SHA-256 D303690E0195654BBE7C3DAA5D89EA23E24CD709A3575217234AC19AE8836694.
% Jacob Reinhold's CC BY 4.0 Markdown was consulted as comparison lineage.
% The genuine commutative diagram is an exact PDF-derived Loop 1 source crop.

\subsubsection*{2.6.3}
\phantomsection
\label{sga3:VIIB:2.6.3}

If \(\mathbf L\) is a flat Lie algebra over \(\mathbf O_k\), the
formal group \(\mathscr G(\mathbf L)\) can be characterized by a
universal property.\footnote{%
\label{sga3:VIIB:nde:111}The editors modified what follows to take
advantage of the addition made in \SGARefLink{sga3:VIIB:2.6.2:prop}{2.6.2}.}
Indeed, every morphism
\[
  \phi:\mathscr G(\mathbf L)\longrightarrow G
\]
to a formal group \(G\) induces a morphism
\[
  \operatorname{Lie}(\phi):
  \operatorname{Lie}\bigl(\mathscr G(\mathbf L)\bigr)
  \longrightarrow
  \operatorname{Lie}(G).
\]
Composing it with
\[
  \tau_{\mathbf L}:
  \mathbf L\longrightarrow
  \operatorname{Lie}\bigl(\mathscr G(\mathbf L)\bigr)
\]
\textup{(cf.\ \SGARefLink{sga3:VIIB:2.6.2:prop}{2.6.2})} gives a morphism
\[
  \phi':\mathbf L\longrightarrow\operatorname{Lie}(G).
\]
We then have:

\medskip
\noindent\textbf{Proposition.}
\phantomsection
\label{sga3:VIIB:2.6.3:prop}
{\itshape Let \(G\) be a formal \(k\)-group and let \(\mathbf L\)
be a flat Lie algebra over \(\mathbf O_k\). The map
\(\phi\mapsto\phi'\) defined above is a bijection
\[
  \operatorname{Hom}_{\operatorname{Grf}/k}
    \bigl(\mathscr G(\mathbf L),G\bigr)
  \xrightarrow{\ \sim\ }
  \operatorname{Hom}_{\operatorname{Lie}}
    \bigl(\mathbf L,\operatorname{Lie}(G)\bigr),
\]
where the right-hand side is the set of morphisms of Lie algebras
over \(\mathbf O_k\) from \(\mathbf L\) to
\(\operatorname{Lie}(G)\).}

\textit{Proof.}
We may immediately reduce to the case in which \(k\) is Artinian.
Put \(L=\mathbf L(k)\). By \SGARefLink{sga3:VIIB:2.3.1}{2.3.1},
\[
  \operatorname{Hom}_{\operatorname{Grf}/k}
    \bigl(\mathscr G(\mathbf L),G\bigr)
\]
is in bijection with the set of unital algebra morphisms
\[
  h:U(L)\longrightarrow\mathbf H(G)
\]
for which the following diagram is commutative:
\begin{center}
  \phantomsection
  \label{sga3:VIIB:2.6.3:universal-property-diagram}
  % SGA3-VII-DIAG-119
  % BEGIN INCLUDED SGA3_English_current_buildable_source_20260823/source/inputs/VII/native_diagrams/exp7/viib_15_95/diag_119.tex
% SGA3-VII-DIAG-119
% Authority: Exp7B-13oct24.pdf, local p. 62.
% Crop: figures/Exp7B_localp062_universal_property_commutative_diagram.png
% Locator: tex/components/53_expose_VIIB_section263_universal_property_en.tex, Loop 1 line 77.
\begin{tikzcd}[row sep=1.2em,column sep=4.2em]
U(L) \arrow[r,"h"] \arrow[dd,"\Delta_{U(L)}"']
  & \mathbf H(G) \arrow[dr,"\delta_G"] & {} \\
{} & {} & \mathbf H(G\times G) \\
U(L)\otimes U(L) \arrow[r,"h\otimes h"]
  & \mathbf H(G)\otimes\mathbf H(G)
      \arrow[ur,"\varphi_G"']
  & {.}
\end{tikzcd}

% END INCLUDED SGA3_English_current_buildable_source_20260823/source/inputs/VII/native_diagrams/exp7/viib_15_95/diag_119.tex

\end{center}
Now \(h\) is determined by its restriction to \(L\), which is a
morphism from \(L\) to the Lie algebra underlying
\(\mathbf H(G)\). Commutativity of the diagram says precisely that
\(h\) sends \(L\) into the subset of \(\mathbf H(G)\) formed by
the elements \(x\) satisfying
\[
  \delta_G(x)
  =
  \varphi_G(x\otimes1+1\otimes x).
\]
This subset is exactly \(\operatorname{Lie}(G)\); cf.\ \SGARefLink{sga3:VIIB:2.6:item:c}{2.6(c)}.
The proposition follows.

% END INCLUDED SGA3_English_current_buildable_source_20260823/source/inputs/VII/tex/components/53_expose_VIIB_section263_universal_property_en.tex

% BEGIN INCLUDED SGA3_English_current_buildable_source_20260823/source/inputs/VII/tex/components/54_expose_VIIB_section27_infinitesimal_formal_monoids_en.tex
% SGA 3, Expose VII_B, English reconstruction -- Loop 1.
% Scope: Section 2.7 in full on authority local pages 62--63,
% stopping immediately before Section 2.8.
% Combined-reader pages: 575--576. Re-edition printed pages: 564--565.
% Original-edition marginal locator in scope: 530.
% Editorial notes in scope: 112--114.
% Authority: Polo--Gille Exp7B-13oct24.pdf,
% SHA-256 D303690E0195654BBE7C3DAA5D89EA23E24CD709A3575217234AC19AE8836694.
% Jacob Reinhold's CC BY 4.0 Markdown was consulted as comparison lineage.
% No genuine diagram occurs in this component; all mathematics is native LaTeX.
% The printed c_U/c_n and x/u inconsistencies are recorded in the source note
% below; the global antipode and the indexed induction are normalized there.

\subsection*{2.7}
\phantomsection
\label{sga3:VIIB:2.7}
\addcontentsline{toc}{subsection}{\SGARefLink{sga3:VIIB:2.7:prop}{2.7}. Infinitesimal formal monoids}

We conclude these generalities with a statement that goes back to
S.~Lie and that will be used in \SGARefLink{sga3:VIIB:5.1}{5.1}. A \emph{formal monoid} over
\(k\) is, by definition, a pair \((M,m)\) consisting of a formal
variety \(M\) and a morphism
\[
  m:M\times M\longrightarrow M
\]
such that \(m(C)\) makes \(M(C)\) a monoid for every object \(C\)
of \(\operatorname{Alf}/k\).\footnote{%
\label{sga3:VIIB:nde:112}Here and in what follows, the editors have
written ``monoid'' instead of ``monoid with identity element''
\textup{(recall that a monoid is, by definition, a set equipped
with an associative composition law and possessing an identity
element)}.}
In particular, the ``unit section,'' which assigns to every object
\(C\) the identity element of \(M(C)\), defines a section
\[
  \varepsilon_M:\operatorname{Spf}(k)\longrightarrow M
\]
of the canonical projection
\(M\to\operatorname{Spf}(k)\). We shall say that the formal
monoid \(M\) is \emph{infinitesimal} if \(\varepsilon_M\) induces
a bijection of the underlying sets.

\medskip
\noindent\textbf{Proposition.}
\phantomsection
\label{sga3:VIIB:2.7:prop}
{\itshape Every topologically flat infinitesimal formal
\(k\)-monoid \(M\) is a formal \(k\)-group.}\footnote{%
\label{sga3:VIIB:nde:113}By the equivalence of categories in
\SGARefLink{sga3:VIIB:1.3.5.D}{1.3.5.D}, a monoid in the category of topologically flat formal
\(k\)-varieties ``is the same thing'' as a monoid in the category
of flat cocommutative \(\mathbf O_k\)-coalgebras, that is, a
cocommutative \(\mathbf O_k\)-bialgebra \textup{(in the usual
sense, hence not necessarily equipped with an antipode)}.
Moreover, by \SGARefLink{sga3:VIIB:1.3.6}{1.3.6}, the hypothesis that \(M\) be infinitesimal is
equivalent to saying that the corresponding bialgebra is
connected. Thus, if \(k\) is an Artinian ring, the proposition is
equivalent to the following statement: every connected
cocommutative \(k\)-bialgebra that is flat over \(k\) is a
\(k\)-Hopf algebra, that is, it possesses an antipode. In fact,
the cocommutativity hypothesis is superfluous; cf.\ the proof.}

\textit{Proof.}
We must show that \(M(C)\) is a group for every object \(C\) of
\(\operatorname{Alf}/k\). We therefore reduce immediately to the
case in which \(k\) is Artinian. Let
\[
  U=\mathbf H(M)
\]
be the coalgebra of \(M\) \textup{(\SGARefLink{sga3:VIIB:1.3.5}{1.3.5})}. The multiplication
\[
  m:M\times M\longrightarrow M
\]
induces a morphism of coalgebras
\[
  m_U:U\otimes U\longrightarrow U
\]
that makes \(U\) an associative algebra over \(k\). The identity
element of this algebra is the image of the identity element of
\(k\) under the map \(k\to U\) induced by the unit section
\(\varepsilon_M\) of \(M\). Likewise, the projection
\[
  M\longrightarrow\operatorname{Spf}(k)
\]
induces a homomorphism
\[
  \varepsilon_U:U\longrightarrow k;
\]
we write \(U^+\) for the kernel of \(\varepsilon_U\).

We must show that there is an antipode, that is, a morphism of
coalgebras
\[
  c_U:U\longrightarrow U
\]
such that, for every \(u\in U\),
\begin{equation}
  \bigl(m_U\circ(c_U\otimes\operatorname{id}_U)
    \circ\Delta_U\bigr)(u)
  =
  \varepsilon_U(u)\,1_U.
  \tag{*}
  \label{sga3:VIIB:2.7:antipode-left-identity}
\end{equation}
\phantomsection
\label{sga3:VIIB:2.7:antipode-source-note}
\noindent\textit{Source note.}\enspace
In \textup{(*)}, the French re-edition prints \(c_n\) in place of
\(c_U\), applies the left-hand side to \(x\), and retains \(u\) on
the right after quantifying over \(u\). At the end of the proof it
again prints
\(m_U\circ(c_n\otimes\operatorname{id}_U)\circ\Delta_U=\eta\)
while calling this the right-inverse identity. The global identity
above therefore uses \(c_U\) and the single variable \(u\); below,
\(c_n\) denotes the coherent map at the \(n\)-th stage of the
induction, and the final right-inverse identity has the tensor
factors in the order required by that assertion.

Let \((U_n)\) be the filtration of \(U\) defined in \SGARefLink{sga3:VIIB:1.3.6}{1.3.6}, and put
\[
  U_n^+=U^+\cap U_n.
\]
Since \(M\) is infinitesimal, \(U^+\) is the union of the
submodules \(U_n^+\).\footnote{%
\label{sga3:VIIB:nde:114}The original continued as follows:
``one then shows easily, by induction on \(n\), that there exists
one and only one linear map \(c_n:U_n\to U_n\) such that the
composite
\[
  m_U\circ(c_n\otimes\operatorname{id}_U)\circ\Delta_U:
  U_n^+\longrightarrow U
\]
is zero.'' The editors have supplied the details of the proof,
which rests on that of \SGARefLink{sga3:VIIB:1.3.6.A}{Lemma~1.3.6.A}.}
Set
\[
  c_0(1)=1
\]
and
\[
  c_1(x)=-x
  \qquad\text{if }x\in U_1^+,
\]
that is, if \(x\) is primitive. Suppose that
\[
  c_{n-1}:U_{n-1}\longrightarrow U
\]
has been constructed so as to satisfy \textup{(*)}, with
\(c_{n-1}\) in place of \(c_U\), for every \(x\in U_{n-1}\), and
let \(x\in U_n^+\). By the proof of \SGARefLink{sga3:VIIB:1.3.6.A}{Lemma~1.3.6.A},
\[
  \Delta_U(x)-x\otimes1
  \in
  U_{n-1}\otimes U^+;
\]
this is where the hypothesis that \(U\) is flat over \(k\) enters.
We may therefore write
\[
  \Delta_U(x)
  =
  x\otimes1+\sum_i y_i\otimes z_i,
  \qquad
  y_i\in U_{n-1},
\]
and then set
\[
  c_n(x)
  =
  -\sum_i c_{n-1}(y_i)z_i.
\]
Together with \(c_n(1)=1\), this prescription extends
\(c_{n-1}\) coherently to \(U_n\). In this way we obtain a
\(k\)-linear map
\[
  c:U\longrightarrow U
\]
that is the left inverse of \(\operatorname{id}_U\) for the monoid
law on \(\operatorname{End}_k(U)\) defined by
\[
  f\mathbin{\cdot}g
  =
  m_U\circ(f\otimes g)\circ\Delta_U;
\]
the identity element is the map
\[
  \eta:u\longmapsto\varepsilon_U(u)\,1_U.
\]
It follows that \(c\) is uniquely determined and is also the right
inverse of \(\operatorname{id}_U\); that is,
\[
  m_U\circ(\operatorname{id}_U\otimes c)\circ\Delta_U
  =
  \eta,
\]
without any assumption that \(U\) be cocommutative. Thus
\(c_U:=c\) is the required antipode, and \(M\) is a formal
\(k\)-group.

% Next cursor: Section 2.8 on authority local page 63,
% combined-reader page 576, re-edition printed page 565.

% END INCLUDED SGA3_English_current_buildable_source_20260823/source/inputs/VII/tex/components/54_expose_VIIB_section27_infinitesimal_formal_monoids_en.tex

% BEGIN INCLUDED SGA3_English_current_buildable_source_20260823/source/inputs/VII/tex/components/55_expose_VIIB_section28_unipotent_group_schemes_en.tex
% SGA 3, Expose VII_B, English reconstruction -- Loop 1.
% Scope: Section 2.8 in full on authority local pages 63--65,
% stopping immediately before Section 2.9.
% Combined-reader pages: 576--578. Re-edition printed pages: 565--567.
% Editorial notes in scope: 115--117.
% Authority: Polo--Gille Exp7B-13oct24.pdf,
% SHA-256 D303690E0195654BBE7C3DAA5D89EA23E24CD709A3575217234AC19AE8836694.
% Jacob Reinhold's CC BY 4.0 Markdown was consulted as comparison lineage.
% No genuine diagram occurs in this component; all mathematics is native LaTeX.

\subsection*{2.8. Unipotent group schemes over a field}
\phantomsection
\label{sga3:VIIB:2.8}
\addcontentsline{toc}{subsection}{\SGARefLink{sga3:VIIB:2.8}{2.8}. Unipotent group schemes over a field}
\footnote{\label{sga3:VIIB:nde:115}The editors added this subsection.}

Let \(k\) be a field. Recall that a \(k\)-group scheme \(G\) is
called \emph{unipotent} if it satisfies the following two conditions
\textup{(cf.\ [DG70], \S~IV.2, \SGARefLink{sga3:VIIB:2.5:proposition}{Proposition 2.5})}:
\begin{enumerate}
\item[\textup{(a)}]
\(G\) is affine.

\item[\textup{(b)}]
Every simple \(\mathcal O(G)\)-comodule is trivial. Explicitly, if
\[
  \rho:V\longrightarrow V\otimes_k\mathcal O(G)
\]
is an \(\mathcal O(G)\)-comodule structure on a nonzero
\(k\)-vector space \(V\), and there is no nonzero proper subspace
\(W\subset V\) such that
\[
  \rho(W)\subset W\otimes_k\mathcal O(G),
\]
then \(V\) is one-dimensional and
\[
  \rho(v)=v\otimes1
  \qquad\text{for every }v\in V.
\]
\end{enumerate}
According to the cited result, when \(G\) is of finite type over
\(k\), this is equivalent to the definition given in
Exposé XVII, \S~1: \(G\) has a finite composition series whose
successive quotients are isomorphic to \(k\)-subgroups of
\(\mathbb G_{a,k}\).

For every affine \(k\)-group scheme \(G\), the comultiplication of
\(\mathcal O(G)\) endows the pseudocompact \(k\)-module
\[
  A=\mathcal O(G)^*
\]
with the structure of a profinite \(k\)-algebra, not necessarily
commutative; its unit \(1_A\) is the augmentation
\[
  \varepsilon:\mathcal O(G)\longrightarrow k.
\]
Set
\[
  I=\{f\in A\mid f(1_{\mathcal O(G)})=0\}.
\]
Then \(I\) is a two-sided ideal of \(A\), and
\[
  A=k\cdot1_A\oplus I;
\]
cf.\ \SGARefLink{sga3:VIIB:1.3.6}{1.3.6}.

Let \(V\) be a finite-codimensional subspace of \(A\), and consider
the continuous \(k\)-bilinear map
\[
  \varphi:A\times A\longrightarrow A/V,
  \qquad
  (a,b)\longmapsto ab+V.
\]
By \SGARefLink{sga3:VIIB:0.3.1}{Lemma 0.3.1}, there are finite-codimensional subspaces
\(L_1,L_2\subset A\) such that \(V\) contains \(AL_2\) and
\(L_1A\). Thus
\[
  L=L_1\cap L_2
\]
is finite-codimensional, and \(V\) contains the two-sided ideal
\(ALA\), which is itself finite-codimensional. Hence the
finite-codimensional two-sided ideals form a neighborhood basis of
zero in \(A\).

It follows that an \(\mathcal O(G)\)-comodule is the same thing as
a continuous \(A\)-module: an \(A\)-module \(V\), equipped with the
discrete topology, for which the map
\[
  A\times V\longrightarrow V
\]
is continuous. Every such module is the union of
finite-dimensional \(A\)-submodules \(V_i\), each of which is a
module over a finite-dimensional quotient \(k\)-algebra \(A_i\) of
\(A\).

Consequently, a simple continuous \(A\)-module \(M\) is
finite-dimensional over \(k\) and is a faithful simple module over
the finite-dimensional algebra \(A/\operatorname{Ann}(M)\). The
latter is therefore a finite-dimensional simple \(k\)-algebra, so
\(\operatorname{Ann}(M)\) is an open maximal ideal of \(A\).

Conversely, let \(P\) be an open prime ideal of
\(A\).\footnote{\label{sga3:VIIB:nde:116}Recall that a two-sided
ideal \(P\) of a ring \(A\) is called prime if, in \(A/P\), the
product of any two nonzero two-sided ideals is nonzero.}
Then \(A/P\) is a finite-dimensional \(k\)-algebra in which the
zero ideal is prime, and hence is a finite-dimensional simple
\(k\)-algebra. Up to isomorphism, there is therefore a unique
simple continuous \(A\)-module with annihilator \(P\). Thus
\[
  M\longmapsto\operatorname{Ann}(M)
\]
is a bijection from the isomorphism classes of simple continuous
\(A\)-modules to the open prime ideals of \(A\).

In particular, the \(A\)-module \(A/I\), which is one-dimensional
over \(k\), is called the \emph{trivial module}. It corresponds to
the trivial one-dimensional \(\mathcal O(G)\)-comodule \(V\), for
which
\[
  \rho(v)=v\otimes1_{\mathcal O(G)}
  \qquad\text{for every }v\in V.
\]
We have proved:

\medskip
\noindent\textbf{Proposition 2.8.1.}
\phantomsection
\label{sga3:VIIB:2.8.1}
{\itshape Let \(k\) be a field, let \(G\) be an affine
\(k\)-group scheme, and put \(A=\mathcal O(G)^*\).
\begin{enumerate}
\item[\textup{(i)}]
\(G\) is unipotent if and only if \(I\) is the unique open prime
ideal of \(A\).\footnote{\label{sga3:VIIB:nde:117}Equivalently,
\(k\cdot1_{\mathcal O(G)}\) is the unique simple
\(k\)-subcoalgebra of \(\mathcal O(G)\); see, for example,
[Ab80], 3.1.4. The editors also note a misprint in the cited work,
p.~130, line~4: \(M\simeq C^*/\operatorname{ann}M\) should be
replaced by
\(C^*/\operatorname{ann}M\simeq\operatorname{End}_k(M)\).}

\item[\textup{(ii)}]
In particular, suppose that \(G\) is commutative, so that
\[
  \mathcal O(G)=\mathbf H(D(G)),
\]
where \(D(G)=\operatorname{Spf}(A)\) is the Cartier dual of \(G\).
Then \(G\) is unipotent if and only if \(D(G)\) is infinitesimal.
\end{enumerate}}


% END INCLUDED SGA3_English_current_buildable_source_20260823/source/inputs/VII/tex/components/55_expose_VIIB_section28_unipotent_group_schemes_en.tex

% BEGIN INCLUDED SGA3_English_current_buildable_source_20260823/source/inputs/VII/tex/components/56_expose_VIIB_section29_pointed_cocommutative_hopf_algebras_en.tex
% SGA 3, Expose VII_B, English reconstruction -- Loop 1.
% Scope: Section 2.9 in full, including Proposition 2.9.1 and
% Remarks 2.9.2, on authority local pages 65--66.
% Stops immediately before Section 3.
% Combined-reader pages: 578--579. Re-edition printed pages: 567--568.
% Editorial note in scope: 118.
% Authority: Polo--Gille Exp7B-13oct24.pdf,
% SHA-256 D303690E0195654BBE7C3DAA5D89EA23E24CD709A3575217234AC19AE8836694.
% Jacob Reinhold's CC BY 4.0 Markdown was consulted as comparison lineage.
% No genuine diagram occurs in this component; all mathematics is native LaTeX.

\subsection*{2.9. Pointed cocommutative Hopf algebras over a field}
\phantomsection
\label{sga3:VIIB:2.9}
\addcontentsline{toc}{subsection}{\SGARefLink{sga3:VIIB:2.9}{2.9}. Pointed cocommutative Hopf algebras over a field}
\footnote{\label{sga3:VIIB:nde:118}The editors added this subsection
to express Proposition 2.5.2, together with the variant noted in
editorial note 100, in the language of cocommutative Hopf
algebras.}

Let \(k\) be a field, let \(C\) be a \(k\)-coalgebra, and endow
\[
  A=C^*
\]
with the structure of a profinite \(k\)-algebra, not necessarily
commutative, described in \SGARefLink{sga3:VIIB:2.8}{2.8}. By \SGARefLink{sga3:VIIB:0.2.2}{0.2.2},
\[
  C=A^\dagger=\operatorname{Hom}_c(A,k).
\]
The map
\[
  D\longmapsto
  D^\perp=\{f\in A\mid f(D)=0\}
\]
is therefore a bijection from the subcoalgebras of \(C\) to the
closed two-sided ideals of \(A\) \textup{(which we shall simply
call ideals below)}. Its inverse is
\[
  I\longmapsto
  I^\perp
  =
  \{x\in C=A^\dagger\mid x(I)=0\}.
\]
Since every maximal closed ideal is a maximal open ideal
\textup{(cf.\ \SGARefLink{sga3:VIIB:0.2.1}{0.2.1})}, every subcoalgebra contains a simple
subcoalgebra, necessarily finite-dimensional.

Recall that a subcoalgebra \(D\subset C\) is called
\emph{irreducible} if it contains a unique simple subcoalgebra
\(S_0\). Equivalently,
\[
  \mathfrak m_0=S_0^\perp
\]
is the unique open maximal ideal containing \(D^\perp\), or,
equivalently,
\[
  D^\perp+\mathfrak m=A
  \qquad
  \text{for every open maximal ideal }
  \mathfrak m\ne\mathfrak m_0.
\]
This applies in particular when \(D=S_0\).

Let \(\Sigma_0\) be the sum of all irreducible subcoalgebras
\(C_i\) containing \(S_0\). It is plainly a subcoalgebra, and it is
irreducible because, for every \(\mathfrak m\ne\mathfrak m_0\),
\SGARefLink{sga3:VIIB:0.2.D}{0.2.D} gives
\[
  \mathfrak m+\bigcap_i C_i^\perp
  =
  \bigcap_i\bigl(\mathfrak m+C_i^\perp\bigr)
  =
  A.
\]
The coalgebra \(\Sigma_0\) is called the \emph{irreducible
component of \(C\) corresponding to \(S_0\)}.
\textit{The French prints \(C_0\) in this last phrase; the
preceding definition makes \(S_0\) the intended simple
subcoalgebra.}

A coalgebra \(C\) is called \emph{pointed} if every simple
\(k\)-subcoalgebra of \(C\) is one-dimensional. Equivalently,
\[
  A/\mathfrak m=k
\]
for every open maximal ideal \(\mathfrak m\) of \(A\). Recall also
that \(C\) is called \emph{connected} if it is irreducible and
pointed. If \(C\) is a bialgebra, then it is connected if and only
if it is irreducible, since \(k\cdot1_C\) is a simple
subcoalgebra.

Suppose from now on that \(C\) is cocommutative. Then \(A\) is
commutative and hence is the product of its local components
\[
  A_{\mathfrak m},
  \qquad
  \mathfrak m\in\Upsilon(A);
\]
cf.\ \SGARefLink{sga3:VIIB:0.1.1}{0.1.1}. Put
\[
  S_{\mathfrak m}
  =
  \mathfrak m^\perp
  \simeq
  (A/\mathfrak m)^*,
\]
and let \(\Sigma_{\mathfrak m}\) be its irreducible component.
This component can be described explicitly. Let
\[
  J_{\mathfrak m}
  =
  \ker\bigl(A\longrightarrow A_{\mathfrak m}\bigr).
\]
Then \(J_{\mathfrak m}\subset\mathfrak m\), and it is the smallest
closed ideal \(I\subset\mathfrak m\) satisfying
\[
  I+\mathfrak m'=A
  \qquad
  \text{for every open maximal ideal }
  \mathfrak m'\ne\mathfrak m.
\]
Indeed, any such \(I\) contains \(A_{\mathfrak m'}\) for every
\(\mathfrak m'\ne\mathfrak m\), and hence contains
\(J_{\mathfrak m}\). Since
\[
  A=J_{\mathfrak m}\oplus A_{\mathfrak m},
\]
we obtain
\[
  \Sigma_{\mathfrak m}
  =
  J_{\mathfrak m}^{\perp}
  \simeq
  A_{\mathfrak m}^{\dagger}.
\]
We can now prove:

\medskip
\noindent\textbf{Proposition 2.9.1.}
\phantomsection
\label{sga3:VIIB:2.9.1}
{\itshape Let \(k\) be a field.
\begin{enumerate}
\item[\textup{\SGAInlineRefTarget{sga3:VIIB:2.9.1:item:i}(i)}]
Let \(G\) be a formal \(k\)-group all of whose residue fields are
equal to \(k\). Then \(G_e\) is the constant \(k\)-group \(M_k\),
where
\[
  M
  =
  G(k)
  =
  \left\{
    x\in\mathbf H(G)
    \,\middle|\,
    \varepsilon(x)=1,\quad
    \Delta(x)=x\otimes x
  \right\},
\]
and \(\mathbf H(G)\) is the semidirect product of
\(\mathbf H(G_0)\) by the group algebra \(kM\)
\textup{(cf.\ \SGARefLink{sga3:VIIB:2.4.B}{2.4.B})}.

\item[\textup{(ii)}]
Equivalently, let \(H\) be a pointed cocommutative Hopf
\(k\)-algebra. Then \(H\) is the semidirect product of the
irreducible component \(H_0\) of the unit \(1_H\) by \(kM\), where
\[
  M
  =
  \{x\in H\mid \varepsilon(x)=1,\ \Delta(x)=x\otimes x\}.
\]
\end{enumerate}}

\textit{Proof.}
We first prove \textup{\SGARefLink{sga3:VIIB:2.9.1:item:i}{(i)}}. Since every residue field of \(G\) is
equal to \(k\), the projection
\[
  \pi:G\longrightarrow G_e
\]
has a section \(s:G_e\to G\), defined at every point \(g\in G\) by
\[
  \mathcal O_{G,g}\longrightarrow\kappa(g)=k.
\]
Moreover, for all \(g,h\in G\), the algebra
\[
  \mathcal O_{G,g}
  \widehat\otimes_k
  \mathcal O_{G,h}
\]
is local with residue field \(k\). Thus \(s\times s\) is a section
of
\[
  \pi\times\pi:
  G\times G\longrightarrow
  (G\times G)_e=G_e\times G_e.
\]
Since \(\pi\) is a group morphism,
\[
  \pi\circ m_G\circ(s\times s)
  =
  m_{G_e}
  =
  \pi\circ s\circ m_{G_e}.
\]
As \(\pi\) is an epimorphism, it follows that \(s\) is a group
morphism. Hence
\[
  G=G_0\rtimes G_e,
\]
and \(\mathbf H(G)\) is the semidirect product of
\(\mathbf H(G_0)\) by \(\mathbf H(G_e)\).

Since every residue field of \(G\) is \(k\),
\(\mathbf H(G_e)\) is the group algebra \(kM\), where
\[
  M=G_e(k);
\]
cf.\ \SGARefLink{sga3:VIIB:2.5.A}{2.5.A}. Finally, because \(G_0\) is infinitesimal, the map
\[
  G(k)\longrightarrow G_e(k)
\]
is injective. It is also surjective, since it has the section
induced by \(s\). Thus \(M=G(k)\), proving \textup{\SGARefLink{sga3:vii:paragraph:5-1:equation:i}{(i)}}.

For \textup{\SGARefLink{sga3:vii:paragraph:5-1:equation:ii}{(ii)}}, it remains only to show that
\[
  H_0=\mathbf H(G_0).
\]
The unit \(1_H\) of \(H\) is the augmentation
\[
  \varepsilon_{\mathcal A}:\mathcal A\longrightarrow k,
\]
which is the identity section of \(G(k)\). Therefore the local
component of the affine algebra \(\mathcal A(G)\) corresponding to
\(H_0\) is precisely \(\mathcal A(G_0)\). By the description above,
\[
  H_0
  =
  \mathcal A(G_0)^\dagger
  =
  \mathbf H(G_0).
\]
This proves the proposition.

\medskip
\noindent\textbf{Remarks 2.9.2.}
\phantomsection
\label{sga3:VIIB:2.9.2}
\begin{enumerate}
\item[\textup{(a)}]
The preceding proposition, which is implicit in \SGARefLink{sga3:VIIB:2.5.2:proposition}{2.5.2}, was obtained
independently by B.\ Kostant \textup{(cf.\ [Sw69], Preface)}.
Combined with Cartier's \SGARefLink{sga3:VIIB:3.3:cartier-theorem}{Theorem 3.3} below
\textup{(cf.\ [Ca62], \S~12, Theorem 3)}, which was also obtained
by Kostant \textup{(cf.\ [Sw69], loc.\ cit.)}, this result is often
called the \emph{Cartier--Gabriel--Kostant theorem}.

\item[\textup{(b)}]
In the form \textup{\SGARefLink{sga3:vii:paragraph:5-1:equation:ii}{(ii)}}, \SGARefLink{sga3:VIIB:2.9.1}{Proposition 2.9.1} was extended by
R.\ G.\ Heyneman and M.\ E.\ Sweedler to the case in which \(H\)
is assumed to be pointed and the direct sum of its irreducible
components, but is not necessarily cocommutative; see [HS69],
Theorem 3.5.8.
\end{enumerate}

% END INCLUDED SGA3_English_current_buildable_source_20260823/source/inputs/VII/tex/components/56_expose_VIIB_section29_pointed_cocommutative_hopf_algebras_en.tex

% BEGIN INCLUDED SGA3_English_current_buildable_source_20260823/source/inputs/VII/tex/components/57_expose_VIIB_section3_and_31_primitives_en.tex
% SGA 3, Expose VII_B, English reconstruction -- Loop 1.
% Scope: Section 3 heading, opening convention, and Section 3.1 in full,
% on authority local pages 66--67, stopping immediately before Section 3.2.
% Combined-reader pages: 579--580. Re-edition printed pages: 568--569.
% Original-edition marginal locator in scope: 531.
% Authority: Polo--Gille Exp7B-13oct24.pdf,
% SHA-256 D303690E0195654BBE7C3DAA5D89EA23E24CD709A3575217234AC19AE8836694.
% Jacob Reinhold's CC BY 4.0 Markdown was consulted as comparison lineage.
% No genuine diagram occurs in this component; all mathematics is native LaTeX.

\section*{3. Phenomena specific to characteristic \(0\)}
\phantomsection
\label{sga3:VIIB:3}
\addcontentsline{toc}{section}{3. Phenomena specific to characteristic 0}

Throughout \SGARefLink{sga3:VIIB:3}{Section~3}, we assume that the pseudocompact ring \(k\)
contains the field of rational numbers \(\mathbb Q\).

\subsection*{3.1. Lemma}
\phantomsection
\label{sga3:VIIB:3.1}

\noindent\textbf{Lemma 3.1.}
\phantomsection
\label{sga3:VIIB:3.1:lemma}
{\itshape Let \(C\) be a commutative unital
\(\mathbb Q\)-algebra, and let \(L\) be a Lie algebra over \(C\)
whose underlying \(C\)-module is free. Then the canonical map
\[
  L\longrightarrow U(L)
\]
identifies \(L\) with the set of primitive elements of \(U(L)\).}

\textit{Proof.}
Identify \(L\) with its canonical image in \(U(L)\). Let \(I\) be
a totally ordered set, and let
\[
  (x_i)_{i\in I}
\]
be a basis of \(L\) indexed by \(I\). Denote by
\(\mathbb N^{(I)}\) the set of families
\[
  n=(n_i)_{i\in I}
\]
of nonnegative integers for which \(n_i=0\) except possibly at
finitely many indices
\[
  i_1<i_2<\cdots<i_s;
\]
these indices depend on \(n\). Put
\[
  x^n
  =
  x_{i_1}^{\,n_{i_1}}
  x_{i_2}^{\,n_{i_2}}
  \cdots
  x_{i_s}^{\,n_{i_s}},
  \qquad
  n!
  =
  (n_{i_1}!)(n_{i_2}!)\cdots(n_{i_s}!).
\]
By the Poincaré--Birkhoff--Witt theorem, the elements \(x^n\)
form a basis of \(U(L)\). Moreover,
\begin{equation}
  \Delta_{U(L)}
  \left(\frac{x^n}{n!}\right)
  =
  \sum_{0\leq m\leq n}
  \frac{x^m}{m!}
  \otimes
  \frac{x^{\,n-m}}{(n-m)!}.
  \tag{*}
  \label{sga3:VIIB:3.1:pbw-comultiplication}
\end{equation}
Here the sum ranges over all \(m\in\mathbb N^{(I)}\) such that
\(0\leq m\leq n\), that is,
\[
  0\leq m_i\leq n_i
  \qquad\text{for every }i.
\]
It follows immediately that
\[
  \Delta_{U(L)}(u)=u\otimes1+1\otimes u
\]
if and only if \(u\) is a linear combination of the \(x_i\).
This proves the lemma.

% END INCLUDED SGA3_English_current_buildable_source_20260823/source/inputs/VII/tex/components/57_expose_VIIB_section3_and_31_primitives_en.tex

% BEGIN INCLUDED SGA3_English_current_buildable_source_20260823/source/inputs/VII/tex/components/58_expose_VIIB_section32_exponential_logarithm_en.tex
% SGA 3, Expose VII_B, English reconstruction -- Loop 1.
% Scope: Section 3.2 from its opening through Proposition 3.2,
% on authority local pages 67--68, stopping immediately before 3.2.1.
% Combined-reader pages: 580--581. Re-edition printed pages: 569--570.
% Original-edition marginal locators in scope: 532--534.
% Editorial notes in scope: 119--123.
% Authority: Polo--Gille Exp7B-13oct24.pdf,
% SHA-256 D303690E0195654BBE7C3DAA5D89EA23E24CD709A3575217234AC19AE8836694.
% Jacob Reinhold's CC BY 4.0 Markdown was consulted as comparison
% lineage only and never as authority.
% No genuine diagram occurs in this component; all mathematics is native LaTeX.

\subsection*{3.2}
\phantomsection
\label{sga3:VIIB:3.2}

Suppose now that \(C\) is Artinian, with radical \(\mathfrak r\).
For every \(C\)-algebra \(U\) \textup{(associative and unital)}, the
ideal \(\mathfrak rU\) therefore consists of nilpotent elements. If
\(x\in\mathfrak rU\), put
\[
  \exp_U(x)
  =
  1+x+\frac{x^2}{2!}+\cdots.
\]
\footnote{\label{sga3:VIIB:nde:119}If
\(x,x'\in\mathfrak rU\) commute, then
\[
  \exp_U(x+x')=\exp_U(x)\exp_U(x').
\]}
This gives a bijection from \(\mathfrak rU\) onto
\(1+\mathfrak rU\); its inverse sends \(1+y\in1+\mathfrak rU\) to
\[
  \log_U(1+y)
  =
  y-\frac{y^2}{2}+\frac{y^3}{3}-\cdots.
\]
Moreover, the map \(\exp_U\) is plainly functorial in
\(U\).\footnote{\label{sga3:VIIB:nde:120}That is, for every
homomorphism \(\phi:U\to V\) of \(C\)-algebras,
\[
  \phi\bigl(\exp_U(x)\bigr)=\exp_V\bigl(\phi(x)\bigr).
\]}

Still assuming \(C\) Artinian, suppose that \(U\) is endowed with a
bialgebra structure over \(C\) \textup{(cf.\ \SGARefLink{sga3:VIIA:2.2}{2.2})}. For every
primitive element \(x\in\mathfrak rU\)
\textup{(cf.\ VII\textsubscript{A}, \SGARefLink{sga3:VIIA:3.2.3}{3.2.3})}, we then have
\[
\begin{aligned}
  \Delta_U\bigl(\exp_U(x)\bigr)
  &=
  \exp_{U\otimes_CU}\bigl(\Delta_U(x)\bigr)
  \\
  &=
  \exp_{U\otimes_CU}(x\otimes1+1\otimes x)
  \\
  &=
  \exp_{U\otimes_CU}(x\otimes1)\,
  \exp_{U\otimes_CU}(1\otimes x)
  \\
  &=
  \bigl(\exp_U(x)\otimes1\bigr)
  \bigl(1\otimes\exp_U(x)\bigr)
  \\
  &=
  \exp_U(x)\otimes\exp_U(x).
\end{aligned}
\]
Thus the bijection \(\exp_U\) sends a primitive element of
\(\mathfrak rU\) to an element \(z\in1+\mathfrak rU\) satisfying
\[
  \Delta_U(z)=z\otimes z.
\]
Conversely, if \(z\) satisfies these conditions, then, on putting
\(x=\log_U(z)\), the preceding calculation shows that
\[
  \exp_{U\otimes_CU}(x\otimes1+1\otimes x)
  =
  z\otimes z
  =
  \Delta_U\bigl(\exp_U(x)\bigr)
  =
  \exp_{U\otimes_CU}\bigl(\Delta_U(x)\bigr);
\]
hence
\[
  x\otimes1+1\otimes x=\Delta_U(x).
\]
\footnote{\label{sga3:VIIB:nde:121}The editors supplied the
details of the preceding argument and added the sentence that
follows.}
Furthermore, if \(z\in1+\mathfrak rU\) satisfies
\(\Delta_U(z)=z\otimes z\), then
\[
  \varepsilon_U(z)^2=\varepsilon_U(z).
\]
Since \(\varepsilon_U(z)\) is invertible \textup{(because \(z\) is
invertible, and \(\mathfrak r\) is nilpotent)}, it follows that
\(\varepsilon_U(z)=1\).

In particular, let \(L\) be a Lie algebra flat over \(C\), take for
\(U\) the enveloping algebra \(U(L)\) of \(L\) over \(C\), and
identify \(L\) with its canonical image in \(U\). By \SGARefLink{sga3:VIIB:3.1:lemma}{Lemma~3.1},
\(L\) is the set of primitive elements of \(U\): indeed, on every
local component of \(C\), \(L\) is a free module. Consider the
\(C\)-formal group
\[
  \mathscr G(L)=\operatorname{Spf}^{*}U(L),
\]
whose covariant bialgebra is \(U=U(L)\)
\textup{(cf.\ \SGARefLink{sga3:VIIB:2.6.2:prop}{2.6.2})}. Let \(\mathfrak m\) be a maximal ideal of
\(C\), and put
\[
  \overline C=C/\mathfrak m,
  \qquad
  \overline U=U\otimes_C\overline C.
\]
Since \(\mathscr G(L)\) is infinitesimal \textup{(loc.\ cit.)}, the
unit element of \(\overline U\) is the only element
\(\overline z\in\overline U\) such
that\footnote{\label{sga3:VIIB:nde:122}The editors added the
condition \(\varepsilon_{\overline U}(\overline z)=1\), which was
omitted in the original.}
\[
  \varepsilon_{\overline U}(\overline z)=1,
  \qquad
  \Delta_{\overline U}(\overline z)
  =
  \overline z\otimes\overline z.
\]
It follows that every \(z\in U\) satisfying
\[
  \varepsilon_U(z)=1,
  \qquad
  \Delta_U(z)=z\otimes z
\]
necessarily belongs to \(1+\mathfrak rU\).

Finally, since
\[
  L\cap\mathfrak rU
  =
  \mathfrak rL
  =
  L\otimes_C\mathfrak r,
\]
we conclude that \(\exp_U\) defines a bijection from
\(L\otimes_C\mathfrak r\) onto the group \(\mathscr G(L)(C)\).
We summarize this as follows.

\medskip
\noindent\textbf{Proposition 3.2.}
\phantomsection
\label{sga3:VIIB:3.2:prop}
{\itshape Let \(k\) be a pseudocompact ring containing
\(\mathbb Q\), and let \(\mathbf L\) be a flat Lie algebra over
\(\mathbf O_k\).
\begin{enumerate}
\item[\textup{(i)}]
For every object \(C\) of \(\operatorname{Alf}/k\), let
\(\mathfrak r(C)\) denote its radical. Then the map
\[
  \exp_{U(\mathbf L(C))}:
  \mathbf L(C)\otimes_C\mathfrak r(C)
  \longrightarrow
  \mathscr G(\mathbf L)(C)
\]
is bijective and functorial in \(C\) and \(\mathbf L\).

\item[\textup{\SGAInlineRefTarget{sga3:VIIB:3.2:prop:item:ii}(ii)}]
The natural morphism
\[
  \mathbf L
  \longrightarrow
  \operatorname{Lie}\bigl(\mathscr G(\mathbf L)\bigr)
\]
\textup{(cf.\ \SGARefLink{sga3:VIIB:2.6.2:prop}{2.6.2})} is an isomorphism of Lie algebras over
\(\mathbf O_k\).
\end{enumerate}}
\footnote{\label{sga3:VIIB:nde:123}The editors added part
\textup{\SGARefLink{sga3:vii:paragraph:5-1:equation:ii}{(ii)}}, which follows immediately from \SGARefLink{sga3:VIIB:3.1:lemma}{3.1}.}


% END INCLUDED SGA3_English_current_buildable_source_20260823/source/inputs/VII/tex/components/58_expose_VIIB_section32_exponential_logarithm_en.tex

% BEGIN INCLUDED SGA3_English_current_buildable_source_20260823/source/inputs/VII/tex/components/59_expose_VIIB_section321_campbell_hausdorff_en.tex
% SGA 3, Expose VII_B, English reconstruction -- Loop 1.
% Scope: Section 3.2.1 in full on authority local pages 68--69,
% stopping immediately before Section 3.3.
% Combined-reader pages: 581--582. Re-edition printed pages: 570--571.
% Original-edition marginal locator in scope: 534.
% Editorial note in scope: 124.
% Authority: Polo--Gille Exp7B-13oct24.pdf,
% SHA-256 D303690E0195654BBE7C3DAA5D89EA23E24CD709A3575217234AC19AE8836694.
% Jacob Reinhold's CC BY 4.0 Markdown was consulted as comparison
% lineage only and never as authority.
% No genuine diagram occurs in this component; all mathematics is native LaTeX.

\subsubsection*{3.2.1}
\phantomsection
\label{sga3:VIIB:3.2.1}

The bijection \(\exp_{U(L(C))}\) allows us, by transport of
structure, to define a group law on the set
\[
  L(C)\otimes_C\mathfrak r(C),
\]
which, as in \SGARefLink{sga3:VIIB:3.2:prop}{3.2}, we identify with a subset of \(U(L(C))\). If
\(x\) and \(y\) are two elements of
\(L(C)\otimes_C\mathfrak r(C)\), this law is given by
\phantomsection
\label{sga3:VIIB:3.2.1:transported-product}
\[
\begin{aligned}
  x\mathbin{\cdot}y
  &=
  \log\bigl((\exp x)(\exp y)\bigr)
  =
  \log\left(
    1+\sum_{p+q>0}\frac{x^p y^q}{p!\,q!}
  \right)
  \\
  &=
  \sum_{m\geq1}
  \sum_{p_i+q_i>0}
  \frac{(-1)^{m-1}}{m}
  \frac{x^{p_1}y^{q_1}}{p_1!\,q_1!}
  \cdots
  \frac{x^{p_m}y^{q_m}}{p_m!\,q_m!}
  \\
  &=
  \sum_{\ell\geq1}P_\ell(x,y).
\end{aligned}
\]
Here \(P_\ell(x,y)\) denotes the sum of the monomials of total
degree \(\ell\) in \(x\) and \(y\). For example,
\[
  P_1(x,y)=x+y,
\]
and
\[
\begin{aligned}
  P_2(x,y)
  &=
  \underbrace{
    \frac{x^2}{2}+xy+\frac{y^2}{2}
  }_{m=1}
  -
  \underbrace{
    \frac12\bigl(x^2+xy+yx+y^2\bigr)
  }_{m=2}
  \\
  &=
  \frac12(xy-yx)
  =
  \frac12[x,y].
\end{aligned}
\]
Moreover, \(P_3(x,y)\) is the sum of the following three terms:
\[
\begin{aligned}
  P_3(x,y)
  ={}&
  \underbrace{
    \frac{x^3}{6}
    +\frac{x^2y}{2}
    +\frac{xy^2}{2}
    +\frac{y^3}{6}
  }_{m=1}
  \\
  &-
  \underbrace{
    \frac12\left(
      x^3+\frac32x^2y+\frac12yx^2+xyx+yxy
      +\frac12y^2x+\frac32xy^2+y^3
    \right)
  }_{m=2}
  \\
  &+
  \underbrace{
    \frac13\left(
      x^3+x^2y+yx^2+xyx+yxy+y^2x+xy^2+y^3
    \right)
  }_{m=3}.
\end{aligned}
\]
Consequently,
\phantomsection
\label{sga3:VIIB:3.2.1:P3}
\[
  P_3(x,y)
  =
  \frac1{12}
  \bigl(
    x^2y+yx^2-2xyx-2yxy+y^2x+xy^2
  \bigr)
  =
  \frac1{12}
  \bigl(
    [[y,x],x]+[y,[y,x]]
  \bigr).
\]

More generally, one can prove the Campbell--Hausdorff
formula:\footnote{%
\label{sga3:VIIB:nde:124}The editors corrected the formula given
in the original, which was erroneous, and added the reference
\textup{[BLie]}.}
\phantomsection
\label{sga3:VIIB:3.2.1:campbell-hausdorff}
\[
\begin{aligned}
  P_\ell(x,y)
  ={}&
  \sum_{m=1}^{\ell}
  \frac{(-1)^{m-1}}{m\cdot\ell}
  \sum_{\substack{
    p_1,\ldots,p_m\\
    q_1,\ldots,q_{m-1}
  }}
  \left(
    \prod_{i=1}^{m-1}
    \frac{(\operatorname{ad}x)^{p_i}}{p_i!}
    \frac{(\operatorname{ad}y)^{q_i}}{q_i!}
  \right)
  \frac{(\operatorname{ad}x)^{p_m}}{p_m!}(y)
  \\
  &+
  \sum_{m=1}^{\ell}
  \frac{(-1)^{m-1}}{m\cdot\ell}
  \sum_{\substack{
    p_1,\ldots,p_{m-1}\\
    q_1,\ldots,q_{m-1}
  }}
  \left(
    \prod_{i=1}^{m-1}
    \frac{(\operatorname{ad}x)^{p_i}}{p_i!}
    \frac{(\operatorname{ad}y)^{q_i}}{q_i!}
  \right)(x).
\end{aligned}
\]
Here the \(p_j,q_i\in\mathbb N\) satisfy
\[
  p_i+q_i\geq1
  \quad(1\leq i\leq m-1),
  \qquad
  p_m+\sum_{i=1}^{m-1}(p_i+q_i)=\ell-1.
\]
In other words, in the sums above every nonzero ``Lie monomial''
has total degree \(\ell\). For a proof, see N.~Jacobson,
\emph{Lie Algebras} (Interscience, 1962), \S V.5, or
\textup{[BLie]}, II, \S6.4, Theorem~2.

% END INCLUDED SGA3_English_current_buildable_source_20260823/source/inputs/VII/tex/components/59_expose_VIIB_section321_campbell_hausdorff_en.tex

% BEGIN INCLUDED SGA3_English_current_buildable_source_20260823/source/inputs/VII/tex/components/60_expose_VIIB_section33_cartier_theorem_en.tex
% SGA 3, Expose VII_B, English reconstruction -- Loop 1.
% Scope: Section 3.3 through the end of Cartier's theorem and proof,
% on authority local pages 69--71; stops immediately before Section 3.3.1.
% Combined-reader pages: 582--584. Re-edition printed pages: 571--573.
% Original-edition marginal locators in scope: 535--538.
% Editorial notes in scope: 125--128.
% Authority: Polo--Gille Exp7B-13oct24.pdf,
% SHA-256 D303690E0195654BBE7C3DAA5D89EA23E24CD709A3575217234AC19AE8836694.
% Jacob Reinhold's CC BY 4.0 Markdown was consulted as comparison lineage.
% No genuine diagram occurs in this component; all mathematics is native LaTeX.

\subsection*{3.3}
\phantomsection
\label{sga3:VIIB:3.3}

If \(G\) is a formal \(k\)-group with affine algebra \(A\), recall
that \(I_A\) denotes the augmentation ideal of \(A\), and that
\[
  \omega_{G/k}
  =
  I_A/I_A^2
  \simeq
  I_A\widehat\otimes_Ak
\]
is a pseudocompact \(k\)-module.

\medskip
\noindent\textbf{Theorem 3.3 \textup{(Cartier)}.}
\phantomsection
\label{sga3:VIIB:3.3:cartier-theorem}
{\itshape Let \(k\) be a pseudocompact ring containing
\(\mathbb Q\),\footnote{\label{sga3:VIIB:nde:125}The editors
removed the hypothesis that \(k\) be local; the proof reduces to
that case.} and let \(G\) be a formal \(k\)-group. The following
assertions are equivalent:
\begin{enumerate}
\item[\textup{\SGAInlineRefTarget{sga3:VIIB:3.3:cartier-theorem:item:i}(i)}] There exists an \(\mathcal O_k\)-Lie algebra
\(\mathbf L\), flat over \(\mathcal O_k\), such that \(G\) is
isomorphic to \(\mathscr G(\mathbf L)\). In this case
\(\mathbf L=\operatorname{Lie}(G)\) by \SGARefLink{sga3:VIIB:3.2:prop}{3.2}.

\item[\textup{\SGAInlineRefTarget{sga3:VIIB:3.3:cartier-theorem:item:ii}(ii)}] There exists a projective pseudocompact
\(k\)-module \(\omega\) such that the formal variety underlying
\(G\) is isomorphic to
\[
  \mathbf V_k^{f,0}(\omega)
\]
\textup{(cf.\ \SGARefLink{sga3:VIIB:1.2.5}{1.2.5})}, whose affine algebra is \(k[[\omega]]\).
In this case \(\omega\simeq\omega_{G/k}\).

\item[\textup{\SGAInlineRefTarget{sga3:VIIB:3.3:cartier-theorem:item:iii}(iii)}] The group \(G\) is infinitesimal, and
\(\omega_{G/k}\) is a projective pseudocompact \(k\)-module.

\item[\textup{\SGAInlineRefTarget{sga3:VIIB:3.3:cartier-theorem:item:iv}(iv)}] The group \(G\) is infinitesimal and
topologically flat over \(k\).
\end{enumerate}}

\noindent\textit{Proof of \textup{\SGARefLink{sga3:vii:paragraph:5-1:equation:i}{(i)}}
\(\boldsymbol{\Rightarrow}\) \textup{\SGARefLink{sga3:vii:paragraph:5-1:equation:ii}{(ii)}}.}
Let
\[
  \omega=\Gamma^*(\mathbf L)
\]
be the pseudocompact \(k\)-module dual to \(\mathbf L\)
\textup{(cf.\ \SGARefLink{sga3:VIIB:1.2.3.D}{1.2.3.D})}. For every \(C\in\operatorname{Alf}/k\),
we must construct, functorially in \(C\), an isomorphism
\[
  \mathbf V_k^{f,0}(\omega)(C)
  \xrightarrow{\ \sim\ }
  \mathscr G(\mathbf L)(C).
\]
By \SGARefLink{sga3:VIIB:1.2.5}{1.2.5}, the left-hand side identifies with
\[
  \operatorname{Hom}_c\bigl(\omega,\mathfrak r(C)\bigr),
\]
where \(\mathfrak r(C)\) is the radical of \(C\). This set is
naturally isomorphic to the set
\[
  \operatorname{Hom}_c\bigl(
    \omega\widehat\otimes_kC,\mathfrak r(C)
  \bigr)
\]
of continuous \(C\)-linear maps. Since
\(\omega\widehat\otimes_kC\) is a projective pseudocompact
\(C\)-module, the canonical map
\[
  \bigl(\omega\widehat\otimes_kC\bigr)^\dagger
  \otimes_C\mathfrak r(C)
  \longrightarrow
  \operatorname{Hom}_c\bigl(
    \omega\widehat\otimes_kC,\mathfrak r(C)
  \bigr)
\]
is bijective \textup{(cf.\ \SGARefLink{sga3:VIIB:0.3.6.A}{0.3.6.A})}. By \SGARefLink{sga3:VIIB:1.2.3.E}{1.2.3.E},
\[
  \mathbf L(C)
  \simeq
  \mathbf V_k^f\bigl(\Gamma^*(\mathbf L)\bigr)(C)
  =
  \bigl(\omega\widehat\otimes_kC\bigr)^\dagger.
\]
Thus \(\mathbf V_k^{f,0}(\omega)(C)\) is canonically isomorphic
to
\[
  \mathbf L(C)\otimes_C\mathfrak r(C),
\]
which \SGARefLink{sga3:VIIB:3.2:prop}{Proposition~3.2} identifies canonically with
\(\mathscr G(\mathbf L)(C)\). This proves
\textup{\SGARefLink{sga3:vii:paragraph:5-1:equation:i}{(i)}} \(\Rightarrow\) \textup{\SGARefLink{sga3:VIIB:3.2:prop:item:ii}{(ii)}}.

\medskip
\noindent\textit{Proof of \textup{\SGARefLink{sga3:vii:paragraph:5-1:equation:ii}{(ii)}}
\(\boldsymbol{\Rightarrow}\) \textup{\SGARefLink{sga3:VIIB:3.3:cartier-theorem:item:iii}{(iii)}}.}
Let \(\omega\) be a projective object of \(\operatorname{PC}(k)\),
and let
\[
  h:k[[\omega]]\xrightarrow{\ \sim\ }A
\]
be an isomorphism onto the affine algebra of \(G\). The composite
\[
  \varepsilon_A\circ h:k[[\omega]]\longrightarrow k
\]
is determined by its restriction
\(\lambda:\omega\to\mathfrak r(k)\). The map
\[
  x\longmapsto x-\lambda(x)
\]
from \(\omega\) to the radical of \(k[[\omega]]\) extends, by the
universal property of \(k[[\omega]]\), to an endomorphism
\(\ell_\lambda\) of \(k[[\omega]]\). Since
\[
  \ell_\lambda\circ\ell_{-\lambda}
  =
  \ell_{-\lambda}\circ\ell_\lambda
  =
  \operatorname{id},
\]
the map \(\ell_\lambda\) is an automorphism. Replacing \(h\) by
\(h\circ\ell_\lambda\), we may assume that
\(\varepsilon_A\circ h\) vanishes on the closed ideal \(I\) of
\(k[[\omega]]\) generated by \(\omega\). Then \(h\) induces an
isomorphism
\[
  I/I^2\xrightarrow{\ \sim\ }I_A/I_A^2.
\]
Since \(I/I^2\simeq\omega\), we obtain
\[
  \omega_{G/k}\simeq\omega,
\]
so \(\omega_{G/k}\) is projective. The formal variety
\(\mathbf V_k^{f,0}(\omega)\), and hence \(G\), is infinitesimal.

\medskip
\noindent\textit{Proof of \textup{\SGARefLink{sga3:VIIB:3.3:cartier-theorem:item:iii}{(iii)}}
\(\boldsymbol{\Rightarrow}\) \textup{\SGARefLink{sga3:vii:paragraph:5-1:equation:i}{(i)}}.}
Assume that \(G\) is infinitesimal and \(\omega_{G/k}\) is
projective. Let \(\mathbf L\) be the \(\mathcal O_k\)-Lie algebra
of \(G\). By \SGARefLink{sga3:VIIB:2.6.1}{2.6.1}, its underlying \(\mathcal O_k\)-module is
\[
  \mathbf V_k^f(\omega_{G/k}).
\]
Consequently, \(\mathbf L\) is flat over \(\mathcal O_k\), and
\SGARefLink{sga3:VIIB:1.2.3.E}{Proposition~1.2.3.E} gives
\[
  \Gamma^*(\mathbf L)\simeq\omega_{G/k}.
\]
The proof of \textup{\SGARefLink{sga3:vii:paragraph:5-1:equation:i}{(i)}} \(\Rightarrow\) \textup{\SGARefLink{sga3:vii:paragraph:5-1:equation:ii}{(ii)}} therefore
identifies the affine algebra of \(\mathscr G(\mathbf L)\) with
\(k[[\omega_{G/k}]]\). By \SGARefLink{sga3:VIIB:2.6.3:prop}{2.6.3}, the identity morphism of
\(\mathbf L\) corresponds canonically to a morphism
\[
  \mathscr G(\mathbf L)\longrightarrow G,
\]
and hence to a continuous \(k\)-algebra morphism
\[
  \varphi:A\longrightarrow k[[\omega_{G/k}]].
\]
Let \(I\) be the closed ideal generated by \(\omega_{G/k}\).
Filter \(k[[\omega_{G/k}]]\) by the closures of the powers \(I^n\),
and \(A\) by the closures of the powers \(I_A^n\). By definition,
\(\varphi\) induces the identity on
\[
  \omega_{G/k}=I_A/I_A^2=I/I^2.
\]

Because \(\omega_{G/k}\) is projective, the canonical projection
\[
  I_A\longrightarrow\omega_{G/k}
\]
has a section \(\tau\). By the universal property of
\(k[[\omega_{G/k}]]\), this section defines a continuous algebra
morphism
\[
  \psi:k[[\omega_{G/k}]]\longrightarrow A.
\]
The composite \(\varphi\circ\psi\) induces the identity on
\(\omega_{G/k}\), and hence on the associated graded algebra of
\(k[[\omega_{G/k}]]\). It follows from
\textup{[CA], \S V.7, Lemma~1} that \(\varphi\circ\psi\) is an
isomorphism.\footnote{\label{sga3:VIIB:nde:126}Compare \SGARefLink{sga3:VIIB:5.1.5}{5.1.5}
below, and also \textup{[BAC], III, \S \SGARefLink{sga3:VIIB:2.8}{2.8}, Theorem~1 and its
corollaries}. The editors supplied the details of the original
argument that follow.}

Moreover, \(\psi\) induces an isomorphism
\[
  I/I^2\xrightarrow{\ \sim\ }I_A/I_A^2,
\]
and hence a surjection on associated graded algebras. Since \(G\)
is radicial, \(I_A\) is contained in the radical of \(A\), so
\[
  \bigcap_n I_A^n=0.
\]
The cited lemma therefore shows that \(\psi\) is surjective. Since
\(\varphi\circ\psi\) is an isomorphism and \(\psi\) is surjective,
both \(\psi\) and \(\varphi\) are isomorphisms. This proves
\textup{\SGARefLink{sga3:VIIB:3.3:cartier-theorem:item:iii}{(iii)}} \(\Rightarrow\) \textup{\SGARefLink{sga3:VIIB:3.3:cartier-theorem:item:i}{(i)}}.

\medskip
Finally, \textup{\SGARefLink{sga3:vii:paragraph:5-1:equation:i}{(i)}} or \textup{\SGARefLink{sga3:vii:paragraph:5-1:equation:ii}{(ii)}} plainly implies
\textup{\SGARefLink{sga3:VIIB:3.3:cartier-theorem:item:iv}{(iv)}}, so it remains to prove
\textup{\SGARefLink{sga3:VIIB:3.3:cartier-theorem:item:iv}{(iv)}} \(\Rightarrow\) \textup{\SGARefLink{sga3:VIIB:3.3:cartier-theorem:item:ii}{(ii)}}. We may assume that
\(k\) is local, with residue field \(k_0\). Put
\[
  G_0=G\widehat\otimes_kk_0,
  \qquad
  \omega=\omega_{G/k},
  \qquad
  \omega_0=\omega\widehat\otimes_kk_0.
\]
\footnote{\label{sga3:VIIB:nde:127}The editors supplied the
details of the original argument that follow.}
The pseudocompact \(k_0\)-module \(\omega_0\) has a pseudobasis
\((y_\lambda)_{\lambda\in\Lambda}\). Let
\[
  \omega'=\prod_{\lambda\in\Lambda}k_\lambda
\]
be the topologically free \(k\)-module with that index set, and
lift each \(y_\lambda\) to an element \(x_\lambda\in\omega\).
This gives a continuous \(k\)-linear map
\[
  f:\omega'\longrightarrow\omega
\]
such that
\[
  f_0=f\widehat\otimes_kk_0
\]
is invertible.\footnote{\label{sga3:VIIB:nde:128}One does not know
a priori that \(\omega\) is a projective pseudocompact
\(k\)-module; this will follow from the argument. Compare the
proof of \textup{\SGARefLink{sga3:VIIB:3.3:cartier-theorem:item:iv}{(iv)}} \(\Rightarrow\) \textup{\SGARefLink{sga3:VIIB:3.3:cartier-theorem:item:iii}{(iii)}} in
VIIA~\SGARefLink{sga3:VIIA:7.4}{7.4}.}
Since \(\omega'\) is projective, \(f\) lifts to a continuous
\(k\)-linear map
\[
  g:\omega'\longrightarrow I_A
\]
with \(\pi\circ g=f\), where
\(\pi:I_A\to\omega\) is the projection. The universal property
of \(k[[\omega']]\) gives a morphism of topological algebras
\[
  \phi:k[[\omega']]\longrightarrow A.
\]

Now
\[
  \omega'\widehat\otimes_kk_0
  \simeq
  \omega\widehat\otimes_kk_0
  =
  \omega_{G_0/k_0},
\]
and consequently
\[
  k[[\omega']]\widehat\otimes_kk_0
  \simeq
  k_0[[\omega_{G_0/k_0}]].
\]
The pair \(k_0,G_0\) satisfies \textup{\SGARefLink{sga3:VIIB:3.3:cartier-theorem:item:iii}{(iii)}}, and the proof of
\textup{\SGARefLink{sga3:VIIB:3.3:cartier-theorem:item:iii}{(iii)}} \(\Rightarrow\) \textup{\SGARefLink{sga3:vii:paragraph:5-1:equation:i}{(i)}} shows that
\[
  \phi_0=\phi\widehat\otimes_kk_0
\]
is invertible. Both \(k[[\omega']]\) and \(A\) are projective
pseudocompact \(k\)-modules, so \SGARefLink{sga3:VIIB:0.3.4}{0.3.4} implies that \(\phi\) itself
is invertible. In particular, if \(I\) is the augmentation ideal
of \(k[[\omega']]\), then \(\phi\) induces an isomorphism
\[
  \omega'=I/I^2
  \xrightarrow{\ \sim\ }
  I_A/I_A^2=\omega.
\]
This completes the proof.

% END INCLUDED SGA3_English_current_buildable_source_20260823/source/inputs/VII/tex/components/60_expose_VIIB_section33_cartier_theorem_en.tex

% BEGIN INCLUDED SGA3_English_current_buildable_source_20260823/source/inputs/VII/tex/components/61_expose_VIIB_section331_smoothness_corollary_en.tex
% SGA 3, Expose VII_B, English reconstruction -- Loop 1.
% Scope: Corollary 3.3.1 in full on authority local pages 71--72,
% stopping immediately before Corollary 3.3.2.
% Combined-reader pages: 584--585. Re-edition printed pages: 573--574.
% Original-edition marginal locator in scope: 538.
% Editorial notes in scope: 129--130.
% Authority: Polo--Gille Exp7B-13oct24.pdf,
% SHA-256 D303690E0195654BBE7C3DAA5D89EA23E24CD709A3575217234AC19AE8836694.
% Jacob Reinhold's CC BY 4.0 Markdown was consulted as comparison lineage.
% No genuine diagram occurs in this component; all mathematics is native LaTeX.

\subsubsection*{3.3.1. Corollary}
\phantomsection
\label{sga3:VIIB:3.3.1}

{\itshape Let \(S\) be a locally Noetherian scheme over
\(\mathbb Q\), and let \(G\) be an \(S\)-group scheme that is flat
and locally of finite type.\footnote{%
\label{sga3:VIIB:nde:129}The editors added the flatness hypothesis,
which was omitted from the original.}
If \(\omega_{G/S}\) is a locally free \(\mathcal O_S\)-module,
then \(G\) is smooth over \(S\) at every point of the unit
section.}

\textit{Proof.}
Let \(x\) be a point of the unit section, let \(s\) be its image
in \(S\), and denote the local rings at \(x\) and \(s\) by
\(\mathcal O_x\) and \(\mathcal O_s\), respectively.\footnote{%
\label{sga3:VIIB:nde:130}The editors supplied the details of the
argument that follows.}
By hypothesis, \(\mathcal O_s\) and \(\mathcal O_x\) are
Noetherian local rings. The maximal-ideal-adic topology on each
of them therefore agrees with the ``pseudocompact'' topology
defined by the ideals of finite colength. Write
\[
  \widehat{\mathcal O}_x
  \qquad\text{and}\qquad
  \widehat{\mathcal O}_s
\]
for their completions in this topology. By EGA~IV\textsubscript{4},
17.5.3, \(G\) is smooth over \(S\) at \(x\) if and only if
\(\widehat{\mathcal O}_x\) is formally smooth over
\(\widehat{\mathcal O}_s\), both algebras being endowed with their
maximal-ideal-adic topologies. It is therefore enough to prove this
last assertion.

The rings \(\widehat{\mathcal O}_x\) and
\(\widehat{\mathcal O}_s\) are the local rings at \(x\) and \(s\)
of the formal varieties \(\widehat G/\widehat S\) and
\(\widehat S\) defined in \SGARefLink{sga3:VIIB:1.2.6}{1.2.6}. Put
\[
  k=\widehat{\mathcal O}_s,
  \qquad
  H=\operatorname{Spf}\bigl(\widehat{\mathcal O}_x\bigr).
\]
Then \(H\) is an infinitesimal formal \(k\)-variety. Since the
formation of \(\widehat G/\widehat S\) commutes with products
\textup{(loc.\ cit.)}, \(H\) is in fact an infinitesimal formal
\(k\)-group.

Let \(I\) be the augmentation ideal of \(\mathcal O_x\).
By hypothesis,
\[
  \omega_{G/S,x}=I/I^2
\]
is a free \(\mathcal O_s\)-module of finite rank \(n\). Since
\(\mathcal O_s\) is Noetherian, the module \(\omega_{H/k}\),
which is the completion of \(\omega_{G/S,x}\), is identified with
\[
  \omega_{G/S,x}
  \otimes_{\mathcal O_s}
  \widehat{\mathcal O}_s.
\]
It is therefore a topologically free \(k\)-module of rank \(n\).
By the implication \textup{(iv)}\(\Rightarrow\)\textup{\SGARefLink{sga3:vii:paragraph:5-1:equation:ii}{(ii)}} of
\SGARefLink{sga3:VIIB:3.3:cartier-theorem}{Theorem~3.3},
\[
  \widehat{\mathcal O}_x
  \simeq
  k[[\omega_{H/k}]]
  \simeq
  k[[t_1,\ldots,t_n]].
\]
Finally, the latter formal power-series algebra is formally smooth
over \(k\), by EGA~0\textsubscript{IV}, 19.3.3. This proves the
corollary.

Thus we recover a result obtained independently for group schemes
locally of finite presentation over an arbitrary scheme \(S\);
cf.\ VI\textsubscript{B}, \SGARefLink{prop:sga3-VIB-1.6}{1.6}.

% END INCLUDED SGA3_English_current_buildable_source_20260823/source/inputs/VII/tex/components/61_expose_VIIB_section331_smoothness_corollary_en.tex

% BEGIN INCLUDED SGA3_English_current_buildable_source_20260823/source/inputs/VII/tex/components/62_expose_VIIB_section332_lie_formal_group_equivalence_en.tex
% SGA 3, Expose VII_B, English reconstruction -- Loop 1.
% Scope: Section 3.3.2 in full on authority local page 72;
% stops immediately before Section 3.3.3.
% Combined-reader page: 585. Re-edition printed page: 574.
% Original-edition marginal locator in scope: 539.
% Editorial note in scope: 131.
% Authority: Polo--Gille Exp7B-13oct24.pdf,
% SHA-256 D303690E0195654BBE7C3DAA5D89EA23E24CD709A3575217234AC19AE8836694.
% Jacob Reinhold's CC BY 4.0 Markdown was consulted as comparison lineage.
% No genuine diagram occurs in this component; all mathematics is native LaTeX.

\subsubsection*{3.3.2 Corollary}
\phantomsection
\label{sga3:VIIB:3.3.2}

\noindent\textbf{Corollary 3.3.2.}
{\itshape Let \(k\) be a field of characteristic \(0\). The
functor
\[
  \mathbf L\longmapsto\mathscr G(\mathbf L)
\]
is an equivalence from the category of \(k\)-Lie algebras to the
category of infinitesimal formal \(k\)-groups.}
\footnote{\label{sga3:VIIB:nde:131}By \SGARefLink{sga3:VIIB:2.7:prop}{2.7}, \SGARefLink{sga3:VIIB:2.2.1}{2.2.1}, and \SGARefLink{sga3:VIIB:1.3.6}{1.3.6}, an
infinitesimal formal \(k\)-group \(G\) is ``the same thing'' as a
connected cocommutative \(k\)-bialgebra \(H\)
\textup{(cf.\ \SGARefLink{sga3:VIIB:2.9}{2.9})}. The corollary is therefore equivalent to the
theorem below, obtained independently by Kostant
\textup{(cf.\ \SGARefLink{sga3:VIIB:2.9.2}{2.9.2})}. J.~W.~Milnor and J.~C.~Moore had earlier
proved it in \textup{[MM65]} under the additional hypothesis that
\(H\) be generated as an algebra by its primitive elements.
Although published in 1965, their text had circulated before 1960
\textup{(see the review [Ja65])}; the result is consequently often
called the Cartier--Kostant--Milnor--Moore theorem.

\smallskip
\noindent\textbf{Cartier--Kostant--Milnor--Moore Theorem.}
Let \(k\) be a field of characteristic \(0\). The functors
\[
  \mathbf L\longmapsto U(\mathbf L),
  \qquad
  H\longmapsto\operatorname{Prim}(H),
\]
define an equivalence between the category of \(k\)-Lie algebras
and the category of connected cocommutative \(k\)-bialgebras.

\smallskip
If \(k\) is an Artinian ring containing \(\mathbb Q\), then
\SGARefLink{sga3:VIIB:3.2:prop}{3.2}\textup{\SGARefLink{sga3:vii:paragraph:5-1:equation:ii}{(ii)}} and \SGARefLink{sga3:VIIB:3.3:cartier-theorem}{3.3}, combined with \SGARefLink{sga3:VIIB:2.7:prop}{2.7}, \SGARefLink{sga3:VIIB:2.2.1}{2.2.1}, and \SGARefLink{sga3:VIIB:1.3.6}{1.3.6},
likewise show that
\[
  \mathbf L\longmapsto\mathscr G(\mathbf L)
\]
\textup{(respectively, \(\mathbf L\mapsto U(\mathbf L)\))}
is an equivalence from the category of \(k\)-Lie algebras flat
over \(k\) to the category of topologically flat infinitesimal
formal \(k\)-groups \textup{(respectively, to the category of
connected cocommutative \(k\)-bialgebras flat over \(k\))}.}

Indeed, as \(G\) ranges over the infinitesimal formal
\(k\)-groups, \SGARefLink{sga3:VIIB:3.3:cartier-theorem}{Theorem~3.3} identifies the functor
\[
  G\longmapsto\mathscr G\bigl(\operatorname{Lie}G\bigr)
\]
with the identity functor. Likewise, \SGARefLink{sga3:VIIB:3.2:prop}{3.2}\textup{\SGARefLink{sga3:vii:paragraph:5-1:equation:ii}{(ii)}} identifies
\[
  \mathbf L
  \longmapsto
  \operatorname{Lie}\bigl(\mathscr G(\mathbf L)\bigr)
  =
  \operatorname{Prim}U(\mathbf L)
\]
with the identity functor on \(k\)-Lie algebras.

% END INCLUDED SGA3_English_current_buildable_source_20260823/source/inputs/VII/tex/components/62_expose_VIIB_section332_lie_formal_group_equivalence_en.tex

% BEGIN INCLUDED SGA3_English_current_buildable_source_20260823/source/inputs/VII/tex/components/63_expose_VIIB_section333_gamma_lie_algebras_en.tex
% SGA 3, Expose VII_B, English reconstruction -- Loop 1.
% Scope: Section 3.3.3 in full on authority local pages 72--73,
% stopping immediately before Section 4.
% Combined-reader pages: 585--586. Re-edition printed pages: 574--575.
% Original-edition marginal locators in scope: 539--540.
% Editorial notes in scope: 132--133.
% Authority: Polo--Gille Exp7B-13oct24.pdf,
% SHA-256 D303690E0195654BBE7C3DAA5D89EA23E24CD709A3575217234AC19AE8836694.
% Jacob Reinhold's CC BY 4.0 Markdown was consulted as comparison lineage.
% The genuine fiber-product diagram is an exact PDF-derived Loop 1 crop.

\subsubsection*{3.3.3}
\phantomsection
\label{sga3:VIIB:3.3.3}

Continue to suppose that \(k\) is a field of characteristic \(0\).
Let \(\overline k\) be an algebraic closure of \(k\), and let
\(\Gamma\) be the topological Galois group of \(\overline k\) over
\(k\).

For every formal \(k\)-group \(G\), denote by
\(\mathbf{Aut}_k(G)\) the \(k\)-group functor assigning to every
finite-dimensional commutative \(k\)-algebra \(C\) the group of
automorphisms of the formal \(C\)-group
\[
  G\widehat\otimes_k C.
\]
\footnote{\label{sga3:VIIB:nde:132}The original stated:
``This \(k\)-functor is manifestly left exact
\textup{(\(H\) is topologically flat over \(k\)!)}.'' The editors
supplied the details that follow.}
Since \(k\) is a field, \(G\) is topologically flat over \(k\).
Equivalently, its affine algebra
\[
  A=\mathcal A(G)
\]
is topologically flat over \(k\), or its covariant bialgebra
\[
  H=\mathbf H(G)
\]
is flat over \(k\). The functor \(\mathbf{Aut}_k(G)\) is therefore
a formal \(k\)-group.

Indeed, \(\mathbf{Aut}_k(G)\) is the fiber product of the following
diagram \textup{(cf.\ \SGARefLink{sga3:i:remark:1-7-3}{Exposé~I, 1.7.3})}:
\begin{center}
  \phantomsection
  \label{sga3:VIIB:3.3.3:automorphism-fiber-product-diagram}
  % SGA3-VII-DIAG-120
  % BEGIN INCLUDED SGA3_English_current_buildable_source_20260823/source/inputs/VII/native_diagrams/exp7/viib_15_95/diag_120.tex
% SGA3-VII-DIAG-120
% Authority: Exp7B-13oct24.pdf, local p. 72.
% Crop: figures/Exp7B_localp072_automorphism_fiber_product_diagram.png
% Locator: tex/components/63_expose_VIIB_section333_gamma_lie_algebras_en.tex, Loop 1 line 49.
\begin{tikzcd}[row sep=3.2em,column sep=3.8em]
{} &
\underline{\operatorname{End}}_k(G)
  \times\underline{\operatorname{End}}_k(G) \arrow[d] \\
\operatorname{Spf}(k) \arrow[r] &
\underline{\operatorname{End}}_k(G)
  \times\underline{\operatorname{End}}_k(G)
\end{tikzcd}

% END INCLUDED SGA3_English_current_buildable_source_20260823/source/inputs/VII/native_diagrams/exp7/viib_15_95/diag_120.tex

\end{center}
The vertical arrow sends
\[
  (\phi,\psi)\longmapsto(\phi\circ\psi,\psi\circ\phi),
\]
whereas the horizontal arrow is the unit section. It is therefore
enough to show that the \(k\)-functor \(\mathbf{End}_k(G)\) is
represented by an object of \(\operatorname{Vaf}/k\); by \SGARefLink{sga3:VIIB:1.1}{1.1} and
\SGARefLink{sga3:VIIB:0.4.2}{0.4.2}, this amounts to showing that it is left exact, or,
equivalently, that it commutes with fiber products of
\(k\)-algebras.

To give an element of \(\mathbf{End}_k(G)(C)\) is equivalent to
giving a \(k\)-algebra morphism
\[
  \phi:H\longrightarrow H\otimes_k C
\]
that also respects the coalgebra structure, that is, for which the
usual compatibility diagrams commute. Since \(H\) is flat over
\(k\), the functor
\[
  H\otimes_k -
\]
commutes with fiber products of \(k\)-algebras. Thus
\(\mathbf{End}_k(G)\) is left exact and can be treated as a formal
\(k\)-group.

If \(L\) is a \(k\)-Lie algebra and
\(G=\mathscr G(L)\), \SGARefLink{sga3:VIIB:3.3:cartier-theorem}{Theorem~3.3} shows that
\(\mathbf{Aut}_k(G)\) is isomorphic to the \(k\)-group functor
\(\mathbf{Aut}_k(L)\) which assigns to each finite-dimensional
\(k\)-algebra \(C\) the group of automorphisms of the
\(C\)-Lie algebra
\[
  C\otimes_k L.
\]

Now let \(G\) be an arbitrary formal \(k\)-group. We saw in \SGARefLink{sga3:VIIB:2.5.2:proposition}{2.5.2}
that \(G\) decomposes canonically as a semidirect product of an
étale formal group \(G_e\) and an infinitesimal formal group
\(G_0\). This semidirect product is determined by the data of
\[
  L=\operatorname{Lie}(G)=\operatorname{Lie}(G_0),
\]
the \(\Gamma\)-group
\[
  M=G_e(\overline k)=G(\overline k),
\]
and a morphism of formal \(k\)-groups
\[
  \Phi:
  G_e\longrightarrow
  \mathbf{Aut}_k(G_0)
  \simeq
  \mathbf{Aut}_k(L).
\]
Such a morphism sends \(G_e\) into the étale part
\(\mathbf{Aut}_k(L)_e\) of \(\mathbf{Aut}_k(L)\)
\textup{(cf.\ \SGARefLink{sga3:VIIB:2.5.1:corollary}{2.5.1})}. It is consequently determined by a morphism
of \(\Gamma\)-groups
\[
  \phi=\Phi(\overline k):
  M\longrightarrow
  \bigl(\mathbf{Aut}_kL\bigr)(\overline k)
  =
  \operatorname{Aut}_{\overline k}
    \bigl(L\otimes_k\overline k\bigr).
\]

Let \(\Gamma\) act on \(L\otimes_k\overline k\) by
\[
  \gamma(\ell\otimes t)=\ell\otimes\gamma(t).
\]
Then \(\phi\) makes \(M\) act on
\(L\otimes_k\overline k\) by automorphisms of
\(\overline k\)-Lie algebras in such a way that
\[
  \phi\bigl(\gamma(m)\bigr)
  =
  \gamma\circ\phi(m)\circ\gamma^{-1}.
\]
Equivalently,
\[
  \gamma(m)(\ell\otimes t)
  =
  \gamma\!\left(
    m\bigl(\ell\otimes\gamma^{-1}(t)\bigr)
  \right)
\]
for every \(m\in M\), \(\gamma\in\Gamma\), and
\(\ell\otimes t\in L\otimes_k\overline k\). We express this last
condition by saying that \(M\) acts on
\(L\otimes_k\overline k\) \emph{compatibly with \(\Gamma\)}.

The situation can be summarized in a ``highbrow'' statement. A
\emph{\(\Gamma\)-Lie algebra over \(k\)} is a triple
\[
  (L,M,\phi)
\]
consisting of a \(k\)-Lie algebra \(L\), a \(\Gamma\)-group \(M\),
and an action \(\phi\) of \(M\) on
\(L\otimes_k\overline k\) that is compatible with the actions of
\(\Gamma\) on \(M\) and on \(L\otimes_k\overline k\).

If \((L,M,\phi)\) and \((L',M',\phi')\) are two such
\(\Gamma\)-Lie algebras, a morphism from the first to the second is
a pair \((f,\theta)\), where
\[
  f:L\longrightarrow L'
\]
is a morphism of \(k\)-Lie algebras and
\[
  \theta:M\longrightarrow M'
\]
is a morphism of \(\Gamma\)-groups, such that
\[
  (f\otimes1)(m\cdot x)
  =
  \theta(m)\cdot(f\otimes1)(x)
\]
for every \(x\in L\otimes_k\overline k\) and \(m\in M\).
We have therefore proved:

\medskip
\noindent\textbf{Theorem.}
\phantomsection
\label{sga3:VIIB:3.3.3:gamma-lie-equivalence}
{\itshape Let \(k\) be a field of characteristic \(0\). The functor
that assigns to a formal \(k\)-group \(G\) the triple
\[
  \bigl(
    \operatorname{Lie}(G),
    G(\overline k),
    \Phi(\overline k)
  \bigr)
\]
is an equivalence from the category of formal \(k\)-groups to the
category of \(\Gamma\)-Lie algebras.}\footnote{%
\label{sga3:VIIB:nde:133}In particular, when
\(k=\overline k\), this recovers the
Cartier--Gabriel--Kostant theorem mentioned in \SGARefLink{sga3:VIIB:2.9.2}{2.9.2}.}

% END INCLUDED SGA3_English_current_buildable_source_20260823/source/inputs/VII/tex/components/63_expose_VIIB_section333_gamma_lie_algebras_en.tex

% BEGIN INCLUDED SGA3_English_current_buildable_source_20260823/source/inputs/VII/tex/components/64_expose_VIIB_section4_and_41_relative_frobenius_en.tex
% SGA 3, Expose VII_B, English reconstruction -- Loop 1.
% Scope: Section 4 heading and opening convention, followed by Section 4.1
% in full, on authority local pages 73--74; stops before Section 4.1.1.
% Combined-reader pages: 586--587. Re-edition printed pages: 575--576.
% Original-edition marginal locator in scope: 541.
% Authority: Polo--Gille Exp7B-13oct24.pdf,
% SHA-256 D303690E0195654BBE7C3DAA5D89EA23E24CD709A3575217234AC19AE8836694.
% Jacob Reinhold's CC BY 4.0 Markdown was consulted as comparison lineage.
% No genuine diagram occurs in this component; all mathematics is native LaTeX.

\section*{4. Phenomena specific to characteristic \(p>0\)}
\phantomsection
\label{sga3:VIIB:4}
\addcontentsline{toc}{section}{4. Phenomena specific to characteristic p > 0}

Throughout \SGARefLink{sga3:VIIB:4}{Section~4}, \(p\) denotes a prime number, \(k\) a
pseudocompact ring of characteristic \(p\), and
\[
  \pi:k\longrightarrow k,\qquad x\longmapsto x^p,
\]
the continuous Frobenius endomorphism of \(k\).

\subsection*{4.1}
\phantomsection
\label{sga3:VIIB:4.1}

Let \(X\) be a formal \(k\)-variety with affine algebra \(A\).
We write \(X^{(p/k)}\), or simply \(X^{(p)}\), for the formal
\(k\)-variety
\[
  X\widehat\otimes_{\pi}k
\]
obtained by the base change \(\pi:k\to k\)
\textup{(cf.\ \SGARefLink{sga3:VIIB:1.2.D}{1.2.D})}. Its affine algebra is the completed tensor
product
\[
  A\widehat\otimes_{\pi}k.
\]
The continuous morphism
\[
  A\widehat\otimes_{\pi}k\longrightarrow A,
  \qquad
  a\widehat\otimes_{\pi}\lambda
  \longmapsto
  a^p\lambda,
\]
therefore induces a morphism of formal \(k\)-varieties
\[
  \operatorname{Fr}(X/k):
  X\longrightarrow X^{(p/k)}.
\]
We call \(\operatorname{Fr}(X/k)\) the \emph{Frobenius morphism of
\(X\) relative to \(k\)}, and below we shall often write
\(\operatorname{Fr}\) in place of \(\operatorname{Fr}(X/k)\).


% END INCLUDED SGA3_English_current_buildable_source_20260823/source/inputs/VII/tex/components/64_expose_VIIB_section4_and_41_relative_frobenius_en.tex

% BEGIN INCLUDED SGA3_English_current_buildable_source_20260823/source/inputs/VII/tex/components/65_expose_VIIB_section411_formal_completion_frobenius_en.tex
% SGA 3, Expose VII_B, English reconstruction -- Loop 1.
% Scope: Section 4.1.1 in full on authority local page 74,
% stopping immediately before Section 4.1.2.
% Combined-reader page: 587. Re-edition printed page: 576.
% Original-edition marginal locator in scope: 542.
% Editorial note in scope: 134.
% Authority: Polo--Gille Exp7B-13oct24.pdf,
% SHA-256 D303690E0195654BBE7C3DAA5D89EA23E24CD709A3575217234AC19AE8836694.
% Jacob Reinhold's CC BY 4.0 Markdown was consulted as comparison lineage.
% The residue-field arrow diagram is an exact PDF-derived Loop 1 crop.

\subsubsection*{4.1.1}
\phantomsection
\label{sga3:VIIB:4.1.1}
\footnote{\label{sga3:VIIB:nde:134}The editors corrected the
original, which gave the inclusion
\[
  (\widehat X/\widehat S)^{(p)}
  \subset
  \widehat{X^{(p)}}/\widehat S
\]
instead of the reverse inclusion. They also note that this
paragraph is not used below.}

Let \(S\) be a scheme over the prime field \(\mathbb F_p\), and
let \(\eta:X\to S\) be an \(S\)-scheme. Write
\(\operatorname{fr}(S)\) for the ``absolute'' Frobenius morphism
of \(S\): it induces the identity on the underlying topological
space and sends every section \(f\) of the structure sheaf to
\(f^p\); cf.\ VII\textsubscript{A}, \SGARefLink{sga3:VIIA:4.1}{4.1}. Let
\[
  X^{(p)}=X\times_{\operatorname{fr}(S)}S
\]
be the corresponding \(S\)-scheme
\textup{(VII\textsubscript{A}, loc.\ cit.)}; thus its structural
morphism to \(S\) is \(\operatorname{fr}(S)\circ\eta\).

Let \(\widehat S\) be the formal scheme whose underlying
topological space is the set of closed points of \(S\), endowed
with the discrete topology, and whose local ring at a closed point
\(s\) is the separated completion
\[
  \mathcal O_{\widehat S,s}
  =
  \widehat{\mathcal O}_{S,s}
\]
of \(\mathcal O_{S,s}\) for the linear topology defined by ideals
of finite colength \textup{(cf.\ \SGARefLink{sga3:VIIB:1.2.6}{1.2.6})}. Its affine algebra is
therefore
\[
  k=\mathcal A(\widehat S)
  =
  \prod_{s\in S_{\mathrm{cl}}}
  \widehat{\mathcal O}_{S,s}.
\]

Recall \textup{(loc.\ cit.)} that
\(\widehat X/\widehat S\) denotes the formal \(k\)-variety
consisting of those points \(x\in X\) for which
\[
  [\kappa(x):\kappa(s)]<\infty,
\]
where \(s\) is the image of \(x\) in \(S\). Its local ring at
\(x\) is the separated completion
\[
  \mathcal O_{\widehat X/\widehat S,x}
  =
  \widehat{\mathcal O}_{X,x}
\]
of \(\mathcal O_{X,x}\) for the linear topology defined by those
ideals \(I\) such that \(\mathcal O_{X,x}/I\) has finite length as
an \(\mathcal O_{X,x}\)-module, and hence also as an
\(\mathcal O_{S,s}\)-module. By base change we can therefore form
the formal \(k\)-variety
\[
  (\widehat X/\widehat S)^{(p)}
  =
  (\widehat X/\widehat S)
  \widehat\otimes_{\pi}k.
\]

We can also consider the formal \(k\)-variety
\(\widehat{X^{(p)}}/\widehat S\). Its underlying set consists of
the points \(x\in X^{(p)}=X\) such that, if \(s\) denotes the image
of \(x\) in \(S\), the composite
\begin{center}
  \phantomsection
  \label{sga3:VIIB:4.1.1:residue-field-frobenius-diagram}
  % SGA3-VII-DIAG-121
  % BEGIN INCLUDED SGA3_English_current_buildable_source_20260823/source/inputs/VII/native_diagrams/exp7/viib_15_95/diag_121.tex
% SGA3-VII-DIAG-121
% Authority: Exp7B-13oct24.pdf, local p. 74.
% Crop: figures/Exp7B_localp074_residue_field_frobenius_diagram.png
% Locator: tex/components/65_expose_VIIB_section411_formal_completion_frobenius_en.tex, Loop 1 line 89.
\(
\kappa(s)\xrightarrow{\ \operatorname{fr}\ }\kappa(s)
\xrightarrow{\ \eta^\natural\ }\kappa(x)
\)

% END INCLUDED SGA3_English_current_buildable_source_20260823/source/inputs/VII/native_diagrams/exp7/viib_15_95/diag_121.tex

\end{center}
makes \(\kappa(x)\) a finite extension of \(\kappa(s)\). In that
case the same is true of
\[
  \eta^\natural:\kappa(s)\longrightarrow\kappa(x),
\]
so \(x\) is a point of \(\widehat X/\widehat S\), and
\[
  \mathcal O_{\widehat{X^{(p)}}/\widehat S,x}
  =
  \widehat{\mathcal O}_{X,x}
  \widehat\otimes_{\pi}
  \widehat{\mathcal O}_{S,s}
  =
  \widehat{\mathcal O}_{X^{(p)},x}.
\]
The second equality follows because the construction
\(X\mapsto\widehat X/\widehat S\) commutes with products
\textup{(cf.\ \SGARefLink{sga3:VIIB:1.2.6}{1.2.6})}.

It follows that
\[
  \widehat{X^{(p)}}/\widehat S
  \ \subset\
  (\widehat X/\widehat S)^{(p)}
\]
as a formal subvariety. Equality holds if and only if, for every
closed point \(s\) of \(S\), the field \(\kappa(s)\) has finite
degree over \(\kappa(s)^p\). This is the case, for example, if
\(S\) is locally of finite type over a field \(\kappa\) such that
\[
  [\kappa:\kappa^p]<\infty.
\]

% END INCLUDED SGA3_English_current_buildable_source_20260823/source/inputs/VII/tex/components/65_expose_VIIB_section411_formal_completion_frobenius_en.tex

% BEGIN INCLUDED SGA3_English_current_buildable_source_20260823/source/inputs/VII/tex/components/66_expose_VIIB_section412_frobenius_kernels_height_en.tex
% SGA 3, Expose VII_B, English reconstruction -- Loop 1.
% Scope: Section 4.1.2 in full on authority local page 75;
% stops immediately before Section 4.2.
% Combined-reader page: 588. Re-edition printed page: 577.
% Authority: Polo--Gille Exp7B-13oct24.pdf,
% SHA-256 D303690E0195654BBE7C3DAA5D89EA23E24CD709A3575217234AC19AE8836694.
% Jacob Reinhold's CC BY 4.0 Markdown was consulted as comparison lineage.
% No genuine diagram occurs in this component; all mathematics is native LaTeX.

\subsubsection*{4.1.2}
\phantomsection
\label{sga3:VIIB:4.1.2}

Let \(G\) be a formal \(k\)-group. Since the functor
\[
  X\longmapsto X^{(p)}
  =
  X\widehat\otimes_\pi k
\]
commutes with finite products, it carries formal \(k\)-groups to
formal \(k\)-groups. Moreover, \(\operatorname{Fr}(X/k)\) is
functorial in \(X\), so
\[
  \operatorname{Fr}:G\longrightarrow G^{(p)}
\]
is a homomorphism of formal \(k\)-groups.

Define recursively
\[
  G^{(p^{\,n+1})}
  =
  \bigl(G^{(p^n)}\bigr)^{(p)}.
\]
Then the same is true of the composite
\[
\begin{aligned}
  \operatorname{Fr}^{n}
  =
  \operatorname{Fr}^{n}(G/k):
  G
  &\xrightarrow{\ \operatorname{Fr}\ }
  G^{(p)}
  \xrightarrow{\ \operatorname{Fr}\ }
  G^{(p^2)}
  \xrightarrow{\ \operatorname{Fr}\ }
  \cdots
  \xrightarrow{\ \operatorname{Fr}\ }
  G^{(p^n)}.
\end{aligned}
\]

Let \(A\) be the affine algebra of \(G\), and let \(I_A\) be its
augmentation ideal.

\medskip
\noindent\textbf{Definitions.}
\begin{enumerate}
\item[\textup{(a)}] The kernel of \(\operatorname{Fr}^{n}\) is
denoted
\[
  \operatorname{Fr}_nG.
\]
It follows directly from the definitions that
\(\operatorname{Fr}_nG\) is infinitesimal and has affine algebra
\[
  A/I_A^{\{p^n\}},
\]
where \(I_A^{\{p^n\}}\) is the closed ideal of \(A\) generated by
the \(p^n\)-th powers of the elements of \(I_A\).

\item[\textup{(b)}] The group \(G\) is said to have
\emph{height at most \(n\)} if
\[
  I_A^{\{p^n\}}=0;
\]
equivalently,
\[
  \operatorname{Fr}_nG=G.
\]
\end{enumerate}

% END INCLUDED SGA3_English_current_buildable_source_20260823/source/inputs/VII/tex/components/66_expose_VIIB_section412_frobenius_kernels_height_en.tex

% BEGIN INCLUDED SGA3_English_current_buildable_source_20260823/source/inputs/VII/tex/components/67_expose_VIIB_sections42_421_restricted_enveloping_group_en.tex
% SGA 3, Expose VII_B, English reconstruction -- Loop 1.
% Scope: Sections 4.2 and 4.2.1 in full on authority local page 75,
% stopping immediately before Section 4.2.2.
% Combined-reader page: 588. Re-edition printed page: 577.
% Original-edition marginal locator in scope: 543.
% Authority: Polo--Gille Exp7B-13oct24.pdf,
% SHA-256 D303690E0195654BBE7C3DAA5D89EA23E24CD709A3575217234AC19AE8836694.
% Jacob Reinhold's CC BY 4.0 Markdown was consulted as comparison lineage.
% No genuine diagram occurs in this component; all mathematics is native LaTeX.

\subsection*{4.2}
\phantomsection
\label{sga3:VIIB:4.2}

It is clear that the Lie algebra \(\operatorname{Lie}(G)\) of a
formal \(k\)-group \(G\) is a \(p\)-Lie subalgebra of
\(\mathbf H(G)\) \textup{(cf.\ \SGARefLink{sga3:VIIA:2.3}{2.3})}. Indeed, we may immediately
reduce to the case in which \(k\) is Artinian. Then
\(\operatorname{Lie}(G)\) is the subset of \(\mathbf H(G)\)
consisting of the elements \(x\) such that
\[
  \varphi_G(x\otimes1+1\otimes x)=\Delta_G(x),
\]
with the notation of \SGARefLink{sga3:VIIB:2.3}{2.3} and \SGARefLink{sga3:VIIB:2.6}{2.6}\textup{(c)}. Hence
\[
\begin{aligned}
  \varphi_G(x^p\otimes1+1\otimes x^p)
  &=
  \varphi_G\bigl((x\otimes1+1\otimes x)^p\bigr)
  \\
  &=
  \varphi_G(x\otimes1+1\otimes x)^p
  \\
  &=
  \Delta_G(x)^p
  =
  \Delta_G(x^p).
\end{aligned}
\]

\subsubsection*{4.2.1}
\phantomsection
\label{sga3:VIIB:4.2.1}

Conversely, every \(p\)-Lie algebra \(\mathbf L\) over
\(\mathbf O_k\) defines a \(k\)-group functor. Let
\(U_p(\mathbf L)\) denote the functor assigning to every object
\(C\) of \(\operatorname{Alf}/k\) the restricted enveloping
algebra
\[
  U_p\bigl(\mathbf L(C)\bigr)
\]
of the \(p\)-Lie algebra \(\mathbf L(C)\) over \(C\)
\textup{(VII\textsubscript{A}, \SGARefLink{sga3:VIIA:5.3}{5.3})}. By
VII\textsubscript{A}, \SGARefLink{sga3:vii:corollary:5-4}{5.4}, \(U_p(\mathbf L)\) is a bialgebra over
\(\mathbf O_k\), and by \SGARefLink{sga3:VIIA:2.2}{2.2} it therefore determines the
\(k\)-group functor
\[
  \operatorname{Spf}^{*}U_p(\mathbf L),
\]
which we henceforth denote by \(\mathscr G_p(\mathbf L)\).

Thus \(\mathscr G_p(\mathbf L)(C)\) is the group of elements
\[
  z\in U_p\bigl(\mathbf L(C)\bigr)
\]
having augmentation \(1\) and satisfying
\[
  \Delta_{U_p(\mathbf L(C))}(z)=z\otimes z.
\]


% END INCLUDED SGA3_English_current_buildable_source_20260823/source/inputs/VII/tex/components/67_expose_VIIB_sections42_421_restricted_enveloping_group_en.tex

% BEGIN INCLUDED SGA3_English_current_buildable_source_20260823/source/inputs/VII/tex/components/68_expose_VIIB_section422_restricted_universal_property_en.tex
% SGA 3, Expose VII_B, English reconstruction -- Loop 1.
% Scope: Section 4.2.2 in full on authority local pages 75--76,
% stopping immediately before Section 4.3.
% Combined-reader pages: 588--589. Re-edition printed pages: 577--578.
% Original-edition marginal locator in scope: 544.
% Editorial notes in scope: 135--136.
% Authority: Polo--Gille Exp7B-13oct24.pdf,
% SHA-256 D303690E0195654BBE7C3DAA5D89EA23E24CD709A3575217234AC19AE8836694.
% Jacob Reinhold's CC BY 4.0 Markdown was consulted as comparison lineage.
% The genuine universal-property diagram is an exact PDF-derived Loop 1 crop.

\subsubsection*{4.2.2}
\phantomsection
\label{sga3:VIIB:4.2.2}

Suppose that \(\mathbf L\) is a flat \(p\)-Lie algebra over
\(\mathbf O_k\). Taking VII\textsubscript{A}, \SGARefLink{sga3:VIIA:5.3.3}{5.3.3} into account,
the argument of part \textup{\SGARefLink{sga3:vii:paragraph:5-1:equation:i}{(i)}} of \SGARefLink{sga3:VIIB:2.6.2:prop}{Proposition~2.6.2} shows that
\(U_p(\mathbf L)\) is flat over \(\mathbf O_k\). Consequently,
\(\mathscr G_p(\mathbf L)\) is a formal \(k\)-group topologically
flat over \(k\), with \(U_p(\mathbf L)\) as its covariant
\(\mathbf O_k\)-bialgebra.

\footnote{\label{sga3:VIIB:nde:135}The editors supplied the
details that follow.}
By the proof of \SGARefLink{sga3:VIIB:2.6.2:prop}{2.6.2}\textup{\SGARefLink{sga3:VIIB:2.6.2:prop:item:iii}{(iii)}}, for every finite-length
\(k\)-algebra \(C\), the Lie algebra
\(\operatorname{Lie}(\mathscr G_p(\mathbf L))(C)\) is the set of
primitive elements of
\[
  U_p(\mathbf L)(C)
  =
  U_p\bigl(\mathbf L(C)\bigr);
\]
see also VII\textsubscript{A}, \SGARefLink{sga3:VIIA:3.2.3}{3.2.3}. Moreover, by
VII\textsubscript{A}, \SGARefLink{sga3:VIIA:5.5.3}{5.5.3}, the canonical morphism
\(\mathbf L\to U_p(\mathbf L)\) induces an isomorphism of
\(p\)-Lie algebras
\[
  \tau_{p,\mathbf L}:
  \mathbf L
  \xrightarrow{\ \sim\ }
  \operatorname{Lie}\bigl(\mathscr G_p(\mathbf L)\bigr).
\]
Compare this with \SGARefLink{sga3:VIIB:3.1:lemma}{3.1} in characteristic \(0\).

The formal group \(\mathscr G_p(\mathbf L)\) can be characterized
by a universal property. Every morphism
\[
  h:\mathscr G_p(\mathbf L)\longrightarrow G
\]
to a formal group \(G\) induces a morphism
\[
  \operatorname{Lie}(h):
  \operatorname{Lie}\bigl(\mathscr G_p(\mathbf L)\bigr)
  \longrightarrow
  \operatorname{Lie}(G).
\]
Composing this morphism with \(\tau_{p,\mathbf L}\) gives
\[
  h':\mathbf L\longrightarrow\operatorname{Lie}(G).
\]

\medskip
\noindent\textbf{Proposition.}
\phantomsection
\label{sga3:VIIB:4.2.2:universal-property}
{\itshape Let \(k\) be a pseudocompact ring of characteristic
\(p>0\), let \(G\) be a formal \(k\)-group, and let
\(\mathbf L\) be a flat \(p\)-Lie algebra over \(\mathbf O_k\).
The map \(h\mapsto h'\) defined above is a bijection
\[
  \operatorname{Hom}_{\operatorname{Grf}/k}
    \bigl(\mathscr G_p(\mathbf L),G\bigr)
  \xrightarrow{\ \sim\ }
  \operatorname{Hom}_{p\text{-}\operatorname{Lie}}
    \bigl(\mathbf L,\operatorname{Lie}(G)\bigr),
\]
where the right-hand side is the set of morphisms of \(p\)-Lie
algebras from \(\mathbf L\) to \(\operatorname{Lie}(G)\).}

\textit{Proof.}\footnote{\label{sga3:VIIB:nde:136}The proof is
identical to that of \SGARefLink{sga3:VIIB:2.6.3:prop}{2.6.3}.}
We may immediately reduce to the case in which \(k\) is Artinian.
Put
\[
  L=\mathbf L(k).
\]
By \SGARefLink{sga3:VIIB:2.3.1}{2.3.1},
\[
  \operatorname{Hom}_{\operatorname{Grf}/k}
    \bigl(\mathscr G_p(\mathbf L),G\bigr)
\]
is in bijection with the set of unital algebra morphisms
\[
  h:U_p(L)\longrightarrow\mathbf H(G)
\]
for which the following diagram commutes:
\begin{center}
  \phantomsection
  \label{sga3:VIIB:4.2.2:restricted-universal-diagram}
  % SGA3-VII-DIAG-122
  % BEGIN INCLUDED SGA3_English_current_buildable_source_20260823/source/inputs/VII/native_diagrams/exp7/viib_15_95/diag_122.tex
% SGA3-VII-DIAG-122
% Authority: Exp7B-13oct24.pdf, local p. 76.
% Crop: figures/Exp7B_localp076_restricted_enveloping_universal_diagram.png
% Locator: tex/components/68_expose_VIIB_section422_restricted_universal_property_en.tex, Loop 1 line 102.
\begin{tikzcd}[row sep=1.2em,column sep=3.9em]
U_p(L) \arrow[r,"h"] \arrow[dd,"\Delta_{U_p(L)}"']
  & \mathbf H(G) \arrow[dr,"\delta_G"] & {} \\
{} & {} & \mathbf H(G\times G) \\
U_p(L)\otimes U_p(L) \arrow[r,"h\otimes h"]
  & \mathbf H(G)\otimes\mathbf H(G)
      \arrow[ur,"\varphi_G"']
  & {.}
\end{tikzcd}

% END INCLUDED SGA3_English_current_buildable_source_20260823/source/inputs/VII/native_diagrams/exp7/viib_15_95/diag_122.tex

\end{center}
The morphism \(h\) is determined by its restriction to \(L\),
which is a morphism of \(p\)-Lie algebras from \(L\) to the
underlying \(p\)-Lie algebra of \(\mathbf H(G)\). The commutativity
of the diagram says precisely that \(h\) sends \(L\) into the
subset of \(\mathbf H(G)\) consisting of those \(x\) such that
\[
  \delta_G(x)
  =
  \varphi_G(x\otimes1+1\otimes x).
\]
This subset is \(\operatorname{Lie}(G)\), by \SGARefLink{sga3:VIIB:2.6}{2.6}\textup{(c)}.
The proposition follows.

% END INCLUDED SGA3_English_current_buildable_source_20260823/source/inputs/VII/tex/components/68_expose_VIIB_section422_restricted_universal_property_en.tex

% BEGIN INCLUDED SGA3_English_current_buildable_source_20260823/source/inputs/VII/tex/components/69_expose_VIIB_section43_restricted_pbw_filtration_en.tex
% SGA 3, Expose VII_B, English reconstruction -- Loop 1.
% Scope: Section 4.3 opening in full on authority local pages 76--77,
% stopping immediately before Section 4.3.1.
% Combined-reader pages: 589--590. Re-edition printed pages: 578--579.
% Original-edition marginal locator in scope: 545.
% Authority: Polo--Gille Exp7B-13oct24.pdf,
% SHA-256 D303690E0195654BBE7C3DAA5D89EA23E24CD709A3575217234AC19AE8836694.
% Jacob Reinhold's CC BY 4.0 Markdown was consulted as comparison lineage.
% No genuine diagram occurs in this component; all mathematics is native LaTeX.

\subsection*{4.3}
\phantomsection
\label{sga3:VIIB:4.3}

We now study the formal \(k\)-group \(\mathscr G_p(\mathbf L)\) in
greater detail when \(\mathbf L\) is a flat \(p\)-Lie algebra over
\(\mathbf O_k\).

For this purpose, first let \(C\) be a ring of characteristic
\(p\), and let \(L\) be a \(p\)-Lie algebra over \(C\) whose
underlying module is free with basis
\[
  (x_i)_{i\in I}.
\]
Let \(\mathcal P\) be the set of families
\[
  n=(n_i)_{i\in I}
\]
of nonnegative integers such that
\[
  0\leq n_i<p
\]
and \(n_i=0\) except possibly for finitely many indices. Endow
\(I\) with a total order, and let
\[
  i_1<i_2<\cdots<i_r
\]
be the indices for which \(n_i\ne0\). Put
\[
  |n|=n_{i_1}+\cdots+n_{i_r},
\]
and
\[
  x^n
  =
  x_{i_1}^{\,n_{i_1}}
  x_{i_2}^{\,n_{i_2}}
  \cdots
  x_{i_r}^{\,n_{i_r}},
  \qquad
  n!
  =
  (n_{i_1}!)\cdots(n_{i_r}!).
\]

By VII\textsubscript{A}, \SGARefLink{sga3:VIIA:5.3.3}{5.3.3}, the monomials \(x^n/n!\) form a
basis of \(U_p(L)\), and
\begin{equation}
  \Delta\left(\frac{x^n}{n!}\right)
  =
  \sum_{m\leq n}
  \frac{x^m}{m!}
  \otimes
  \frac{x^{\,n-m}}{(n-m)!}.
  \tag{*}
  \label{sga3:VIIB:4.3:restricted-pbw-comultiplication}
\end{equation}
The sum ranges over all \(m\in\mathcal P\) such that \(m\leq n\),
that is,
\[
  m_i\leq n_i
  \qquad\text{for every }i\in I.
\]

Formula \textup{(*)} gives an immediate description of the natural
filtration of \(U_p(L)\) \textup{(cf.\ \SGARefLink{sga3:VIIB:1.3.6}{1.3.6})}. Put
\[
  U=U_p(L),
\]
let \(U^+\) be the two-sided ideal generated by \(L\), and put
\[
  U_0=C\cdot1_U.
\]
As in \SGARefLink{sga3:VIIB:1.3.6}{1.3.6}, define a filtration of \(U\) by
\(C\)-subcoalgebras by setting, for \(n\geq1\),
\[
  U_n
  =
  \left\{
    u\in U
    \,\middle|\,
    \Delta_U(u)-u\otimes1
    \in U_{n-1}\otimes U^+
  \right\}.
\]
Formula \textup{(*)} then shows that \(U_n\) is the free
\(C\)-module with basis the monomials \(x^m\) satisfying
\[
  |m|\leq n.
\]


% END INCLUDED SGA3_English_current_buildable_source_20260823/source/inputs/VII/tex/components/69_expose_VIIB_section43_restricted_pbw_filtration_en.tex

% BEGIN INCLUDED SGA3_English_current_buildable_source_20260823/source/inputs/VII/tex/components/70_expose_VIIB_section431_group_like_elements_en.tex
% SGA 3, Expose VII_B, English reconstruction -- Loop 1.
% Scope: Section 4.3.1 in full on authority local page 77,
% stopping immediately before Section 4.3.2.
% Combined-reader page: 590. Re-edition printed page: 579.
% Original-edition marginal locators in scope: 546--547.
% Editorial note in scope: 137.
% Authority: Polo--Gille Exp7B-13oct24.pdf,
% SHA-256 D303690E0195654BBE7C3DAA5D89EA23E24CD709A3575217234AC19AE8836694.
% Jacob Reinhold's CC BY 4.0 Markdown was consulted as comparison lineage.
% No genuine diagram occurs in this component; all mathematics is native LaTeX.

\subsubsection*{4.3.1}
\phantomsection
\label{sga3:VIIB:4.3.1}

With the notation of \SGARefLink{sga3:VIIB:4.3}{4.3}, let us determine the elements
\(\xi\in U=U_p(L)\) such that
\[
  \varepsilon_U(\xi)=1,
  \qquad
  \Delta_U(\xi)=\xi\otimes\xi.
\]
Every element \(\xi\in U\) has a unique expression
\[
  \xi
  =
  \sum_{n\in\mathcal P}
  \xi_n\frac{x^n}{n!},
  \qquad
  \xi_n\in C.
\]
The condition \(\varepsilon_U(\xi)=1\) gives
\[
  \xi_0=1,
\]
where \(0\) denotes the element of \(\mathcal P\) all of whose
components vanish. The condition
\(\Delta_{U_p(L)}(\xi)=\xi\otimes\xi\) gives
\[
  \sum_{\substack{n\in\mathcal P\\m\leq n}}
  \xi_n
  \frac{x^m}{m!}
  \otimes
  \frac{x^{\,n-m}}{(n-m)!}
  =
  \sum_{q,r\in\mathcal P}
  \xi_q\xi_r
  \frac{x^q}{q!}
  \otimes
  \frac{x^r}{r!}.
\]
Equivalently,
\[
  \xi_{q+r}=\xi_q\xi_r
  \quad\text{if }q_i+r_i<p\text{ for every }i,
\]
and
\[
  \xi_q\xi_r=0
  \quad\text{otherwise}.
\]

Let \(\xi_i\) denote the coefficient \(\xi_n\) for which
\(n_i=1\) and \(n_j=0\) whenever \(j\ne i\). The preceding
conditions say precisely that
\[
  \xi_n=\prod_i\xi_i^{\,n_i}
  \quad(n\in\mathcal P),
  \qquad
  \xi_i^p=0
  \quad(i\in I).
\]
We obtain:

\medskip
\noindent\textbf{Proposition.}
\phantomsection
\label{sga3:VIIB:4.3.1:group-like-parametrization}
{\itshape Let \(k\) be a local pseudocompact ring\footnote{%
\label{sga3:VIIB:nde:137}The editors added the hypothesis that
\(k\) be local, so that every topologically flat pseudocompact
\(k\)-module is topologically free; see the proof.}
of characteristic \(p>0\), let \(\mathbf L\) be a flat \(p\)-Lie
algebra over \(\mathbf O_k\), and let \(C\) be an object of
\(\operatorname{Alf}/k\). Write
\[
  \sqrt[p]{0_C}
  =
  \{x\in C\mid x^p=0\}.
\]
There is a bijection, ``functorial in \(C\),''
\[
  \mathbf L(C)\otimes_C\sqrt[p]{0_C}
  \xrightarrow{\ \sim\ }
  \mathscr G_p(\mathbf L)(C).
\]}

\textit{Proof.}
By \SGARefLink{sga3:VIIB:1.2.3.F}{Remark~1.2.3.F}, one may choose a basis
\[
  ({}^C x_i)_{i\in I}
\]
of \(\mathbf L(C)\) over \(C\) in such a way that, for every
morphism \(\phi:C\to D\) in \(\operatorname{Alf}/k\), the map
\(\mathbf L(\phi)\) sends \({}^C x_i\) to \({}^D x_i\).
The preceding calculation shows that
\[
  \sum_i({}^C x_i)\otimes\xi_i
  \longmapsto
  \sum_{n\in\mathcal P}
  \left(\prod_i\xi_i^{\,n_i}\right)
  \frac{({}^C x)^n}{n!}
\]
is a functorial bijection from
\(\mathbf L(C)\otimes_C\sqrt[p]{0_C}\) onto
\(\mathscr G_p(\mathbf L)(C)\).


% END INCLUDED SGA3_English_current_buildable_source_20260823/source/inputs/VII/tex/components/70_expose_VIIB_section431_group_like_elements_en.tex

% BEGIN INCLUDED SGA3_English_current_buildable_source_20260823/source/inputs/VII/tex/components/71_expose_VIIB_section432_no_functorial_campbell_hausdorff_en.tex
% SGA 3, Expose VII_B, English reconstruction -- Loop 1.
% Scope: Section 4.3.2 in full on authority local pages 78--79,
% stopping immediately before Section 4.4.
% Combined-reader pages: 591--592. Re-edition printed pages: 580--581.
% Original-edition marginal locator in scope: 548.
% Editorial note in scope: 138.
% Authority: Polo--Gille Exp7B-13oct24.pdf,
% SHA-256 D303690E0195654BBE7C3DAA5D89EA23E24CD709A3575217234AC19AE8836694.
% Jacob Reinhold's CC BY 4.0 Markdown was consulted as comparison lineage.
% The genuine naturality square is an exact PDF-derived Loop 1 crop.

\subsubsection*{4.3.2}
\phantomsection
\label{sga3:VIIB:4.3.2}

``There is no Campbell--Hausdorff formula in characteristic
\(p>0\).'' Let us explain. The functorial isomorphism of \SGARefLink{sga3:VIIA:4.3.1}{4.3.1}
depends on the choice of the bases
\(({}^C x_i)_{i\in I}\). Unlike the situation in \SGARefLink{sga3:VIIB:3.2:prop}{3.2}, in
characteristic \(0\), there is in general no bijection
\[
  \mathbf L(C)\otimes_C\sqrt[p]{0_C}
  \longrightarrow
  \mathscr G_p(\mathbf L)(C)
\]
that is ``functorial both in \(C\) and in \(\mathbf L\).''
The following argument shows, for example, that no such
isomorphism exists when \(k\) is a field containing all
\((p^2-1)\)-st roots of unity.

Choose \(\mathbf L\) so that, for every
\(C\in\operatorname{Alf}/k\), the algebra \(\mathbf L(C)\) is the
abelian \(p\)-Lie algebra over \(C\) generated by an element \(x\)
subject to
\[
  x^{(p^2)}=0.
\]
Thus
\[
  \mathbf L(C)
  =
  Cx\oplus Cx^{(p)},
  \qquad
  U_p\bigl(\mathbf L(C)\bigr)
  \simeq
  C[x]/(x^{p^2}).
\]
\phantomsection
\label{sga3:VIIB:4.3.2:source-notation-note}
\noindent\textit{Source note.}\enspace
The French prints \(k[x]/(x^{p^2})\) in the last display, although
the enveloping algebra here is formed over \(C\). A little later it
also identifies \(\mathscr G_p(\mathbf L)(C)\), whose elements are
group-like, with \(\operatorname{Prim}U_p(\mathbf L(C))\).
The target uses the coefficient ring \(C\) and retains the
group-like set defined in \SGARefLink{sga3:VIIA:4.2.1}{4.2.1}; primitive elements instead
describe its Lie algebra in 4.2.2.

We shall show that the only functorial morphism
\[
  \chi_C:
  \mathbf L(C)\otimes_C\sqrt[p]{0_C}
  \longrightarrow
  U_p\bigl(\mathbf L(C)\bigr)
\]
which is compatible with the endomorphisms of \(\mathbf L\) and
sends \(\mathbf L(C)\otimes_C\sqrt[p]{0_C}\) into
\(\mathscr G_p(\mathbf L)(C)\) is the constant morphism with value
\(1\).

Let
\[
  \psi_C:
  \mathbf L(C)\otimes_C\sqrt[p]{0_C}
  \xrightarrow{\ \sim\ }
  \mathscr G_p(\mathbf L)(C)
\]
be the bijection of \SGARefLink{sga3:VIIA:4.3.1}{4.3.1}, namely
\[
  x\otimes a+x^{(p)}\otimes b
  \longmapsto
  \sum_{0\leq i,j<p}
  a^i b^j x^{\,i+pj}.
\]
Composing \(\chi_C\) with \(\psi_C^{-1}\) gives a morphism,
functorial in \(C\),
\[
\begin{aligned}
  \theta_C:
  \mathbf L(C)\otimes_C\sqrt[p]{0_C}
  &\longrightarrow
  \mathbf L(C)\otimes_C\sqrt[p]{0_C},
  \\
  x\otimes a+x^{(p)}\otimes b
  &\longmapsto
  x\otimes P(a,b)+x^{(p)}\otimes Q(a,b).
\end{aligned}
\]
Functoriality in \(C\) implies that \(P(a,b)\) and \(Q(a,b)\) are
the values at \((a,b)\) of two polynomials
\[
  P,Q\in k[X,Y]
\]
which may be assumed to have degree less than \(p\) in each of
\(X\) and \(Y\).\footnote{\label{sga3:VIIB:nde:138}The editors
expanded the original in the argument that follows.}

Let \(\alpha\in k\), and let \(\ell_\alpha\) be the endomorphism of
the \(p\)-Lie algebra \(\mathbf L\) sending
\[
  x\longmapsto\alpha x,
  \qquad
  x^{(p)}\longmapsto\alpha^p x^{(p)}.
\]
Then \(U_p(\ell_\alpha)\) sends \(x\) to \(\alpha x\), and hence
sends \(x^i\) to \(\alpha^i x^i\) for \(i<p^2\). The following
square commutes:
\begin{center}
  \phantomsection
  \label{sga3:VIIB:4.3.2:naturality-endomorphism-square}
  % SGA3-VII-DIAG-123
  % BEGIN INCLUDED SGA3_English_current_buildable_source_20260823/source/inputs/VII/native_diagrams/exp7/viib_15_95/diag_123.tex
% SGA3-VII-DIAG-123
% Authority: Exp7B-13oct24.pdf, local p. 78.
% Crop: figures/Exp7B_localp078_naturality_endomorphism_square.png
% Locator: tex/components/71_expose_VIIB_section432_no_functorial_campbell_hausdorff_en.tex, Loop 1 line 121.
\begin{tikzcd}[row sep=3.2em,column sep=4.8em]
\mathbf L(C)\otimes_C\sqrt[p]{0_C}
  \arrow[r,"\psi_C"]
  \arrow[d,"{\ell_\alpha(C)\otimes_C\operatorname{id}}"]
  & U_p(\mathbf L(C))
      \arrow[d,"{U_p(\ell_\alpha)(C)}"] \\
\mathbf L(C)\otimes_C\sqrt[p]{0_C}
  \arrow[r,"\psi_C"]
  & U_p(\mathbf L(C))
\end{tikzcd}

% END INCLUDED SGA3_English_current_buildable_source_20260823/source/inputs/VII/native_diagrams/exp7/viib_15_95/diag_123.tex

\end{center}
The hypothesis
\[
  \chi_C\circ(\ell_\alpha(C)\otimes\operatorname{id})
  =
  U_p(\ell_\alpha)(C)\circ\chi_C
\]
therefore gives
\[
  P(\alpha a,\alpha^p b)=\alpha P(a,b),
  \qquad
  Q(\alpha a,\alpha^p b)=\alpha^p Q(a,b).
\]

Write
\[
  P=\sum_{i,j<p}\lambda_{ij}X^iY^j,
  \qquad
  Q=\sum_{i,j<p}\mu_{ij}X^iY^j,
\]
and take
\[
  C=k[X,Y]/(X^p,Y^p).
\]
If \(\lambda_{ij}\ne0\), then
\[
  \alpha^{\,i-1+pj}=1;
\]
if \(\mu_{ij}\ne0\), then
\[
  \alpha^{\,i+p(j-1)}=1.
\]
Taking \(\alpha\) to be a primitive root of unity of order
\(p^2-1\), we conclude that
\[
  P=\lambda X,
  \qquad
  Q=\mu Y
\]
for some \(\lambda,\mu\in k\). Consequently,
\[
  \chi_C(x\otimes a+x^{(p)}\otimes b)
  =
  \sum_{0\leq i,j<p}
  (\lambda a)^i(\mu b)^j x^{\,i+pj}.
\]

Finally, consider the endomorphism \(f\) of \(\mathbf L\) sending
\[
  x\longmapsto x+x^{(p)}.
\]
Take \(C=k[a]/(a^3)\) and compare the coefficients of \(x^p\) and
\(x^{p+1}\) in
\[
  (\chi_C\circ f(C))(x\otimes a)
  \qquad\text{and}\qquad
  \bigl(U_p(f)(C)\circ\chi_C\bigr)(x\otimes a).
\]
This gives
\[
  \lambda=\mu,
  \qquad
  \lambda^2a^2=0,
\]
and hence \(\lambda=\mu=0\), as claimed.

% END INCLUDED SGA3_English_current_buildable_source_20260823/source/inputs/VII/tex/components/71_expose_VIIB_section432_no_functorial_campbell_hausdorff_en.tex

% BEGIN INCLUDED SGA3_English_current_buildable_source_20260823/source/inputs/VII/tex/components/72_expose_VIIB_sections44_441_height_one_groups_en.tex
% SGA 3, Expose VII_B, English reconstruction -- Loop 1.
% Scope: Theorem 4.4, Remark 4.4.A, and Corollary 4.4.1 in full
% on authority local pages 79--80; stops immediately before Section 4.4.2.
% Combined-reader pages: 592--593. Re-edition printed pages: 581--582.
% Original-edition marginal locators in scope: 549--550.
% Editorial notes in scope: 139--140.
% Authority: Polo--Gille Exp7B-13oct24.pdf,
% SHA-256 D303690E0195654BBE7C3DAA5D89EA23E24CD709A3575217234AC19AE8836694.
% Jacob Reinhold's CC BY 4.0 Markdown was consulted as comparison lineage.
% No genuine diagram occurs in this component; all mathematics is native LaTeX.

\subsection*{\SGAInlineRefTarget{sga3:viib:subsection:4-4}4.4 Theorem}
\phantomsection
\label{sga3:VIIB:4.4}

\noindent\textbf{Theorem 4.4.}
{\itshape Let \(k\) be a pseudocompact ring of characteristic
\(p>0\), and let \(G\) be a formal \(k\)-group. The following
assertions are equivalent:
\begin{enumerate}
\item[\textup{(i)}] There exists a \(p\)-Lie algebra \(\mathbf L\),
flat over \(\mathcal O_k\), such that \(G\) is isomorphic to
\(\mathscr G_p(\mathbf L)\). In this case,
\(\mathbf L=\operatorname{Lie}(G)\) by \SGARefLink{sga3:VIIB:4.2.2}{4.2.2}.

\item[\textup{(ii)}] There exists a projective pseudocompact
\(k\)-module \(\omega\) such that the affine algebra of \(G\) is
isomorphic to the quotient of \(k[[\omega]]\) by the closed ideal
generated by the elements \(x^p\), \(x\in\omega\). In this case,
\(\omega\simeq\omega_{G/k}\).

\item[\textup{\SGAInlineRefTarget{sga3:viib:subsection:4-4:item:iii}(iii)}] The group \(G\) has height at most \(1\), and
\(\omega_{G/k}\) is a projective pseudocompact \(k\)-module.

\item[\textup{(iv)}] The group \(G\) has height at most \(1\), and
is topologically flat over \(k\).
\end{enumerate}}

\noindent\textit{Proof of \textup{\SGARefLink{sga3:vii:paragraph:5-1:equation:i}{(i)}} \(\boldsymbol{\Rightarrow}\)
\textup{\SGARefLink{sga3:vii:paragraph:5-1:equation:ii}{(ii)}}.}
Set
\[
  \omega=\Gamma^*(\mathbf L)
\]
\textup{(cf.\ \SGARefLink{sga3:VIIB:1.2.3.D}{1.2.3.D})}, and let \(A\) be the quotient of
\(k[[\omega]]\) by the closed ideal generated by the elements
\(x^p\), \(x\in\omega\). Every continuous morphism
\[
  h:A\longrightarrow C,
  \qquad C\in\operatorname{Alf}/k,
\]
is determined by its restriction \(h'\) to \(\omega\), and this
restriction maps \(\omega\) into \(\sqrt[p]{0_C}\). Thus there is a
canonical bijection from
\[
  \operatorname{Hom}_{\operatorname{Alp}/k}(A,C)
\]
onto the set
\[
  \operatorname{Hom}_{c}\bigl(\omega,\sqrt[p]{0_C}\bigr)
\]
of continuous \(k\)-linear maps from \(\omega\) into
\(\sqrt[p]{0_C}\). The latter set is canonically isomorphic to
\[
  \mathbf L(C)\otimes_C\sqrt[p]{0_C}
\]
\textup{(see the proof of \SGARefLink{sga3:VIIB:3.3:cartier-theorem}{3.3})}. The implication
\textup{\SGARefLink{sga3:vii:paragraph:5-1:equation:i}{(i)}} \(\Rightarrow\) \textup{\SGARefLink{sga3:vii:paragraph:5-1:equation:ii}{(ii)}} therefore follows from
the functorial bijection
\[
  \mathbf L(C)\otimes_C\sqrt[p]{0_C}
  \xrightarrow{\ \sim\ }
  \mathscr G_p(\mathbf L)(C)
\]
established in \SGARefLink{sga3:VIIA:4.3.1}{4.3.1}.

For the remaining implications, consult the analogous proofs of
\SGARefLink{sga3:VIIB:3.3:cartier-theorem}{Theorems~3.3} and VIIA~\SGARefLink{sga3:VIIA:7.4.1}{7.4.1}.

\medskip
\noindent\textbf{Remark 4.4.A.}
\phantomsection
\label{sga3:VIIB:4.4.A}
Let \(G\) be an infinitesimal formal \(k\)-group, with affine
algebra \(A\), such that
\[
  \omega_{G/k}=I_A/I_A^2
\]
is a projective pseudocompact \(k\)-module. Then the projection
\[
  I_A\longrightarrow\omega_{G/k}
\]
has a section
\[
  \tau:\omega_{G/k}\longrightarrow I_A,
\]
and \(\tau\) induces a surjective continuous morphism of algebras
\[
  \psi:k[[\omega_{G/k}]]\longrightarrow A;
\]
compare the proof of the implication
\textup{\SGARefLink{sga3:viib:subsection:4-4:item:iii}{(iii)}} \(\Rightarrow\) \textup{\SGARefLink{sga3:vii:paragraph:5-1:equation:i}{(i)}} in \SGARefLink{sga3:VIIB:3.3:cartier-theorem}{3.3}.\footnote{%
\label{sga3:VIIB:nde:139}The editors added this remark, which is
used in \SGARefLink{sga3:VIIB:4.4.2}{4.4.2}.}

\subsubsection*{4.4.1 Corollary}
\phantomsection
\label{sga3:VIIB:4.4.1}

\noindent\textbf{Corollary 4.4.1.}
{\itshape If \(k\) is a field of characteristic \(p>0\), the functor
\[
  \mathbf L\longmapsto\mathscr G_p(\mathbf L)
\]
is an equivalence from the category of \(p\)-Lie algebras over
\(k\) to the category of formal \(k\)-groups of height at most
\(1\).}

Indeed, as \(G\) ranges over the formal groups of height at most
\(1\), \SGARefLink{sga3:viib:subsection:4-4}{Theorem~4.4} identifies the functor
\[
  G\longmapsto\mathscr G_p\bigl(\operatorname{Lie}G\bigr)
\]
with the identity functor. Likewise, \SGARefLink{sga3:VIIB:4.2.2}{4.2.2} identifies the functor
\[
  \mathbf L\longmapsto
  \operatorname{Lie}\mathscr G_p(\mathbf L)
\]
with the identity functor \textup{(see also VIIA~\SGARefLink{sga3:VIIA:5.5.3}{5.5.3})}.\footnote{%
\label{sga3:VIIB:nde:140}If \(k\) is an Artinian ring of
characteristic \(p>0\), the same proof gives an equivalence between
the category of \(p\)-Lie algebras flat over \(k\) and the category
of formal \(k\)-groups of height at most \(1\) that are
topologically flat over \(k\).}

% END INCLUDED SGA3_English_current_buildable_source_20260823/source/inputs/VII/tex/components/72_expose_VIIB_sections44_441_height_one_groups_en.tex

% BEGIN INCLUDED SGA3_English_current_buildable_source_20260823/source/inputs/VII/tex/components/73_expose_VIIB_section442_finite_height_filtration_en.tex
% SGA 3, Expose VII_B, English reconstruction -- Loop 1.
% Scope: Section 4.4.2 in full on authority local page 80; stops
% immediately before Section 5.
% Combined-reader page: 593. Re-edition printed page: 582.
% Editorial note in scope: 141.
% Authority: Polo--Gille Exp7B-13oct24.pdf,
% SHA-256 D303690E0195654BBE7C3DAA5D89EA23E24CD709A3575217234AC19AE8836694.
% Jacob Reinhold's CC BY 4.0 Markdown was consulted as comparison lineage.
% The genuine Frobenius/height-reduction diagram is an exact
% PDF-derived Loop 1 crop.

\subsubsection*{4.4.2}
\phantomsection
\label{sga3:VIIB:4.4.2}

Continue to assume that \(k\) is a field of characteristic \(p\).
Let \(G\) be an infinitesimal formal \(k\)-group, with affine
algebra \(A\). Since \(G\) is infinitesimal, every open ideal of
\(A\) contains
\[
  I_A^{\{p^n\}}
\]
for all sufficiently large \(n\). Thus \(A\) is the projective
limit of the affine algebras
\[
  A/I_A^{\{p^n\}}
\]
of the groups \(\operatorname{Fr}_nG\)
\textup{(cf.\ \SGARefLink{sga3:VIIA:4.1.2}{4.1.2})}. Equivalently, every infinitesimal formal
\(k\)-group is an inductive limit of formal \(k\)-groups of finite
height.

Suppose that \(G\) has height at most \(n\), and write
\[
  H=G/\operatorname{Fr}G.
\]
\footnote{\label{sga3:VIIB:nde:141}The editors supplied the details
of the original argument in what follows.}
By \SGARefLink{sga3:VIIB:2.4}{2.4} and \SGARefLink{sga3:VIIB:2.4.1}{2.4.1}, the morphism
\[
  \operatorname{Fr}:G\longrightarrow G^{(p)}
\]
factors as an epimorphism
\[
  \pi:G\longrightarrow H
\]
followed by a monomorphism
\[
  i:H\longrightarrow G^{(p)}.
\]
We therefore have the following commutative diagram:
\begin{center}
  \phantomsection
  \label{sga3:VIIB:4.4.2:height-reduction-frobenius-diagram}
  % SGA3-VII-DIAG-124
  % BEGIN INCLUDED SGA3_English_current_buildable_source_20260823/source/inputs/VII/native_diagrams/exp7/viib_15_95/diag_124.tex
% SGA3-VII-DIAG-124
% Authority: Exp7B-13oct24.pdf, local p. 80.
% Crop: figures/Exp7B_localp080_height_reduction_frobenius_diagram.png
% Locator: tex/components/73_expose_VIIB_section442_finite_height_filtration_en.tex, Loop 1 line 55.
\begin{tikzcd}[row sep=2.8em,column sep=3em]
G \arrow[r,"\pi"] \arrow[dr,"\operatorname{Fr}"']
  & H \arrow[r,"\operatorname{Fr}^{\,n-1}"] \arrow[d,"i"]
  & H^{(p^{\,n-1})}
      \arrow[d,"i^{(p^{\,n-1})}" right] \\
{} & G^{(p)} \arrow[r,"\operatorname{Fr}^{\,n-1}"]
  & G^{(p^n)}
\end{tikzcd}

% END INCLUDED SGA3_English_current_buildable_source_20260823/source/inputs/VII/native_diagrams/exp7/viib_15_95/diag_124.tex

\end{center}
Since the functor \(X\mapsto X^{(p)}\) commutes with fiber
products, the morphism
\[
  i^{(p^{\,n-1})}
\]
is again a monomorphism. Moreover,
\(\operatorname{Fr}^{n}(G/k)\) is zero and \(\pi\) is an
epimorphism. Hence
\[
  \operatorname{Fr}^{n-1}\bigl(G^{(p)}/k\bigr)\circ i
  =
  i^{(p^{\,n-1})}\circ
  \operatorname{Fr}^{n-1}(H/k)
\]
is zero. Because \(i^{(p^{\,n-1})}\) is a monomorphism,
\[
  \operatorname{Fr}^{n-1}(H/k)=0,
\]
so \(H\) has height at most \(n-1\).

It follows that every formal \(k\)-group of finite height has a
composition series whose quotients have height at most \(1\), and
therefore can be described by \(p\)-Lie algebras over \(k\).

Finally, the affine algebra \(A\) of \(G\) is a quotient of
\[
  k[[\omega_{G/k}]]
\]
by \SGARefLink{sga3:VIIB:4.4.A}{Remark~4.4.A}. If \(\omega_{G/k}\) is finite-dimensional over
\(k\), all the algebras
\[
  A/I_A^{\{p^n\}}
\]
are finite-dimensional \(k\)-vector spaces. Consequently, every
infinitesimal formal group \(G\) over a field of characteristic
\(p>0\) for which \(\omega_{G/k}\) is finite-dimensional over
\(k\) is an inductive limit of finite formal groups
\textup{(that is, groups of finite length; cf.\ \SGARefLink{sga3:VIIB:1.2.6}{1.2.6})}.

% END INCLUDED SGA3_English_current_buildable_source_20260823/source/inputs/VII/tex/components/73_expose_VIIB_section442_finite_height_filtration_en.tex

% BEGIN INCLUDED SGA3_English_current_buildable_source_20260823/source/inputs/VII/tex/components/74_expose_VIIB_section5_and_50_homogeneous_spaces_hopf_en.tex
% SGA 3, Expose VII_B, English reconstruction -- Loop 1.
% Scope: Section 5 heading and Section 5.0 through Theorem 5.0.1
% and its bibliographic discussion, on authority local pages 80--82;
% stops immediately before Section 5.1.
% Combined-reader pages: 593--595. Re-edition printed pages: 582--584.
% Original-edition marginal locator in scope: 551.
% Editorial note in scope: 142.
% Authority: Polo--Gille Exp7B-13oct24.pdf,
% SHA-256 D303690E0195654BBE7C3DAA5D89EA23E24CD709A3575217234AC19AE8836694.
% Jacob Reinhold's CC BY 4.0 Markdown was consulted as comparison lineage.
% No genuine diagram occurs in this component; all mathematics is native LaTeX.

\section*{5. Homogeneous spaces of infinitesimal formal groups over a field}
\phantomsection
\label{sga3:VIIB:5}
\addcontentsline{toc}{section}{5. Homogeneous spaces of infinitesimal formal groups over a field}

\subsection*{5.0}
\phantomsection
\label{sga3:VIIB:5.0}

For simplicity, suppose that \(k\) is a field.\footnote{%
\label{sga3:VIIB:nde:142}The editors added subsection 5.0 in order
to express \SGARefLink{sga3:VIIB:5.1:proposition}{Proposition~5.1} below in the language of cocommutative
Hopf algebras and to cite results subsequently obtained in this
direction.}
Let \(G\) be a formal \(k\)-group with affine algebra
\[
  A=\mathscr A(G).
\]
Let \(A^+\) be the augmentation ideal of \(A\). For \(a\in A\),
write
\[
  \overline a=a-\varepsilon_A(a)1_A
\]
for its projection onto \(A^+\). Let \(F\) be a formal subgroup of
\(G\), defined by a closed ideal \(J\) of \(A\), and let
\[
  \pi:A\longrightarrow A/J
\]
be the projection. Write
\[
  \overline\pi:A\longrightarrow A^+/J
\]
for the composite of the projections
\[
  A\longrightarrow A^+\longrightarrow A^+/J.
\]
For every \(a\in A\), the projection of \(\Delta_A(a)\) onto
\(A\widehat\otimes A^+\) is
\[
  \Delta_A(a)-a\widehat\otimes1.
\]

By \SGARefLink{sga3:VIIB:2.4:theorem}{Theorem~2.4}, the quotient formal \(k\)-variety \(X=G/F\)
exists. Its affine algebra is
\[
\begin{aligned}
  B
  &=
  \operatorname{Eq}\!\left(
    A
    \mathrel{\mathop{\rightrightarrows}^{\tau_1}_{%
      (\operatorname{id}_A\widehat\otimes\pi)\Delta_A}}
    A\widehat\otimes_k(A/J)
  \right)
  \\
  &=
  \left\{
    a\in A\ \middle|\
    (\operatorname{id}_A\widehat\otimes\pi)\Delta_A(a)
    =
    a\widehat\otimes\pi(1)
  \right\}
  \\
  &=
  \ker\!\left(
    (\operatorname{id}_A\widehat\otimes\overline\pi)\Delta_A
  \right),
\end{aligned}
\]
where
\[
  \tau_1(a)=a\widehat\otimes\pi(1).
\]

Equivalently, \(B\) is the subalgebra of \(A\) consisting of the
elements \(\varphi\) such that
\[
  \varphi(gh)=\varphi(g)
\]
for every \(C\in\operatorname{Alf}/k\),
\(g\in G(C)\), and \(h\in F(C)\). For
\(g,g'\in G(C)\) and \(h\in F(C)\), one has
\[
  \varphi(gg'h)=\varphi(gg').
\]
Thus \(\Delta_A(\varphi)\) belongs to the kernel of
\[
  \operatorname{id}_A\widehat\otimes_k
  \bigl((\operatorname{id}_A\widehat\otimes\overline\pi)
        \Delta_A\bigr),
\]
which equals \(A\widehat\otimes_kB\), since \(A\) is
topologically flat over \(k\). Consequently,
\[
  \Delta_A(B)\subset A\widehat\otimes_kB;
\]
that is, the closed subalgebra \(B\) is also a left coideal.

Conversely, \(B\) determines \(F\), because \SGARefLink{sga3:VIIB:2.4.1:corollary}{Corollary~2.4.1} gives
\[
  J=AB^+,
  \qquad
  B^+=B\cap A^+;
\]
in other words, \(J\) is the closed ideal generated by \(B^+\).
We therefore obtain an injective map
\[
  \mathcal Q:\mathcal F\longrightarrow\mathcal B
\]
from the set \(\mathcal F\) of formal subgroups of \(G\) to the set
\(\mathcal B\) of closed subalgebras of \(A\) that are left
coideals. \SGARefLink{sga3:VIIB:5.1:proposition}{Proposition~5.1} below shows that \(\mathcal Q\) is
bijective when \(G\) is infinitesimal. In fact, it is bijective for
every formal \(k\)-group \(G\).

Recall from \SGARefLink{sga3:VIIB:2.2.1}{2.2.1} that the functor
\[
  G\longmapsto\mathbf H(G)
\]
is an equivalence from formal \(k\)-groups to cocommutative
\(k\)-Hopf algebras. If \(F\) is defined by the closed ideal
\(J\subset A\), then the Hopf subalgebra
\[
  \mathbf H(F)\subset H=\mathbf H(G)
\]
is the orthogonal complement of \(J\) for the duality
\[
  A=H^*,
  \qquad
  H=A^\dagger
\]
\textup{(cf.\ \SGARefLink{sga3:VIIB:0.2.2}{0.2.2})}. Conversely, if \(B\) is a closed subalgebra
of \(A\) that is also a left coideal, then
\[
  I=B^\perp
\]
is a coideal and a left ideal of \(H\); explicitly,
\[
  \Delta_H(I)\subset I\otimes H+H\otimes I,
  \qquad
  \varepsilon_H(I)=0.
\]

Let \(\mathcal H\) be the set of Hopf subalgebras of \(H\), and
let \(\mathcal I\) be the set of left ideals of \(H\) that are
also coideals. For \(I\in\mathcal I\), denote by
\[
  \pi_I:H\longrightarrow H/I
\]
the projection, and by
\[
  \overline\pi_I:H\longrightarrow H^+/I
\]
the composite
\[
  H\longrightarrow H^+\longrightarrow H^+/I,
\]
where \(H^+\) is the augmentation ideal of \(H\).

Let \(K\) be a Hopf subalgebra of \(H\), and put
\[
  K^+=K\cap H^+.
\]
If \(F\) is the formal subgroup corresponding to \(K\), then
\[
  J=K^\perp
\]
and \(A^+/J\) identifies with the dual of \(K^+\). Since
\(B=\mathcal Q(F)\) is the kernel of
\[
  A
  \xrightarrow{\ \Delta_A\ }
  A\widehat\otimes_kA
  \xrightarrow{\ \operatorname{id}\widehat\otimes\overline\pi\ }
  A\widehat\otimes_k(A^+/J),
\]
duality identifies \(\mathcal Q\) with the map
\[
  \Phi:\mathcal H\longrightarrow\mathcal I,
  \qquad
  K\longmapsto HK^+.
\]
Indeed, \(HK^+\) is the image of
\[
  H\otimes_kK^+
  \xrightarrow{\ \operatorname{id}\otimes\operatorname{can}\ }
  H\otimes_kH
  \xrightarrow{\ m_H\ }
  H,
\]
and is both a left ideal and a coideal.

Similarly, the map that assigns \(AB^+\) to
\(B\in\mathcal B\) corresponds by duality to
\[
  \Psi:\mathcal I\longrightarrow\mathcal H,
\]
where
\[
  \Psi(I)
  =
  \left\{
    x\in H\ \middle|\
    (\operatorname{id}_H\otimes\pi_I)\Delta_H(x)
    =
    x\otimes\pi_I(1)
  \right\}.
\]

\medskip
\noindent\textbf{Theorem 5.0.1.}
\phantomsection
\label{sga3:VIIB:5.0.1}
{\itshape Let \(k\) be a field and \(H\) a cocommutative Hopf
algebra. The maps \(\Phi\) and \(\Psi\) above are mutually inverse
bijections between the Hopf subalgebras of \(H\) and the left
ideals of \(H\) that are coideals.}

This theorem was first proved by K.~Newman
\textup{([Ne75], \SGARefLink{sga3:VIIB:4.1}{Theorem~4.1}, where the word
``cocommutative'' was omitted)}. Newman's proof uses the
Cartier--Gabriel--Kostant theorem \textup{(cf.\ \SGARefLink{sga3:VIIB:2.9}{2.9})} to reduce to
the case in which \(H\) is connected, and then uses the existence
of a ``Sweedler basis'' in that case \textup{([Sw67],
Theorem~3)}, a result dual to the Dieudonn\'e--Cartier
\SGARefLink{sga3:VIIB:5.2.2}{Theorem~5.2.2} below. A second, shorter proof was given by
H.~J.~Schneider \textup{([Sch90], Theorem~4.15)}. A.~Masuoka
subsequently generalized the result to the case in which only the
coradical \(H_0\) of \(H\), the sum of its simple subcoalgebras, is
assumed commutative \textup{([Ma91], Theorem~1.3\,(3))}.

Finally, for a commutative \(k\)-Hopf algebra corresponding to an
affine \(k\)-group scheme \(G\), one cannot expect an analogue of
\SGARefLink{sga3:VIIB:5.0.1}{Theorem~5.0.1} without additional hypotheses: for a
\(k\)-subgroup \(F\subset G\), the quotient \(G/F\) need not be
affine. Nevertheless, M.~Takeuchi established analogous
bijections. In \textup{[Tak72], \SGARefLink{sga3:VIIB:4.3}{Theorem~4.3}}, invariant
\(k\)-subgroups \(F\) correspond to subalgebras
\(B\subset\mathcal O(G)\) satisfying
\[
  \Delta(B)\subset A\otimes B
\]
and stable under the antipode. In \textup{[Tak79], Theorem~3},
the \(k\)-subgroups \(F\) for which \(G/F\) is affine correspond to
such subalgebras \(B\) for which \(B\to A\) is faithfully flat.

% END INCLUDED SGA3_English_current_buildable_source_20260823/source/inputs/VII/tex/components/74_expose_VIIB_section5_and_50_homogeneous_spaces_hopf_en.tex

% BEGIN INCLUDED SGA3_English_current_buildable_source_20260823/source/inputs/VII/tex/components/75_expose_VIIB_section51_quotient_criterion_en.tex
% SGA 3, Expose VII_B, English reconstruction -- Loop 1.
% Scope: Section 5.1 opening, proposition, and proof reduction in full
% on authority local pages 82--83; stops immediately before Section 5.1.1.
% Combined-reader pages: 595--596. Re-edition printed pages: 584--585.
% Editorial notes in scope: 143--144.
% Authority: Polo--Gille Exp7B-13oct24.pdf,
% SHA-256 D303690E0195654BBE7C3DAA5D89EA23E24CD709A3575217234AC19AE8836694.
% Jacob Reinhold's CC BY 4.0 Markdown was consulted as comparison lineage.
% No genuine diagram occurs in this component; exact sequences are native LaTeX.

\subsection*{5.1}
\phantomsection
\label{sga3:VIIB:5.1}

Let \(k\) be a pseudocompact ring.\footnote{%
\label{sga3:VIIB:nde:143}The original assumes in \SGARefLink{sga3:VIIB:5.1:proposition}{5.1} that \(k\)
is a field. In fact, that hypothesis can be replaced by flatness
hypotheses; the editors modified \SGARefLink{sga3:VIIB:5.1}{5.1}--\SGARefLink{sga3:VIIB:5.1.5}{5.1.5} accordingly.}
Let \(G\) be a topologically flat infinitesimal formal
\(k\)-group, let \(A\) be its affine algebra, and let
\(B\subset A\) be a closed subalgebra. Put
\[
  X=\operatorname{Spf}(B),
\]
and let
\[
  r:G\longrightarrow X
\]
be the epimorphism induced by the inclusion \(B\subset A\). We
shall determine when \(r\) makes \(X\) the quotient of \(G\) on
the right by a subgroup \(H\)
\textup{(cf.\ \SGARefLink{sga3:VIIB:2.4}{2.4})}.\footnote{%
\label{sga3:VIIB:nde:144}The editors replaced ``left'' by
``right'' and modified the statement of \SGARefLink{sga3:VIIB:5.1:proposition}{Proposition~5.1} so as to
display both the equivalent conditions \textup{\SGARefLink{sga3:vii:paragraph:5-1:equation:i}{(i)}}, \textup{\SGARefLink{sga3:vii:paragraph:5-1:equation:ii}{(ii)}}
and the conclusion \(\operatorname{Spf}(B)\simeq G/H\) more
clearly.}

\medskip
\noindent\textbf{Proposition 5.1.}
\phantomsection
\label{sga3:VIIB:5.1:proposition}
{\itshape Let \(G\) be a topologically flat infinitesimal formal
\(k\)-group, let \(A\) be its affine algebra, let \(I_A\) be the
augmentation ideal of \(A\), and let \(B\) be a closed subalgebra
of \(A\). Put
\[
  I_B=B\cap I_A,
  \qquad
  J_B=AI_B.
\]
Assume that \(A\), and also
\[
  J_B^n/J_B^{n+1}
\]
for every \(n\geq0\), are topologically flat over \(k\). Then the
following two assertions are equivalent:
\begin{enumerate}
\item[\textup{(i)}] For every \(x\in B\),
\[
  \Delta_A(x)-x\widehat\otimes1
  \in
  A\widehat\otimes_kJ_B.
\]

\item[\textup{\SGAInlineRefTarget{sga3:VIIB:5.1:proposition:item:ii}(ii)}] Let
\[
  \pi:A\longrightarrow A/J_B
\]
be the projection, and put
\[
  \tau_1(a)=a\widehat\otimes1.
\]
The sequence
\[
  B\longrightarrow A
  \mathrel{\mathop{\rightrightarrows}^{\tau_1}_{%
    (\operatorname{id}_A\widehat\otimes\pi)\Delta_A}}
  A\widehat\otimes_k(A/J_B)
  \tag{*}\label{sga3:vii:proposition:5-1:equation:star}
\]
is exact. Equivalently, \(B\) is the set of all \(x\in A\) for
which
\[
  \Delta_A(x)-x\widehat\otimes1
  \in
  A\widehat\otimes_kJ_B.
\]
\end{enumerate}
In this case,
\[
  H=\operatorname{Spf}(A/J_B)
\]
is a formal subgroup of \(G\), and the sequence
\[
  G\times H
  \mathrel{\mathop{\rightrightarrows}^{\operatorname{pr}_1}_{\lambda}}
  G
  \longrightarrow
  \operatorname{Spf}(B)
  \tag{**}\label{sga3:vii:proposition:5-1:equation:double-star}
\]
is exact, where \(\lambda\) is the restriction of the
multiplication of \(G\) to \(G\times H\). Thus
\[
  \operatorname{Spf}(B)\simeq G/H.
\]}

\phantomsection
\label{sga3:VIIB:5.1:source-arrow-note}
\noindent\textit{Source note.}\enspace
In \textup{(*)}, the French PDF labels the lower arrow
\((\pi\widehat\otimes\operatorname{id}_A)\Delta_A\). The displayed
equivalent condition immediately before and after the sequence,
as well as the stated right-quotient convention, require
\((\operatorname{id}_A\widehat\otimes\pi)\Delta_A\), which is used
here.

\medskip
\noindent\textit{Proof.}
Put
\[
  \overline A=A/J_B,
  \qquad
  H=\operatorname{Spf}(\overline A).
\]
Then \(H\) is a formal subvariety of \(G\). Since
\(J_B\subset I_A\), the augmentation \(\varepsilon_A\) induces a
continuous \(k\)-algebra morphism
\[
  \overline\varepsilon:\overline A\longrightarrow k.
\]

If \textup{\SGARefLink{sga3:vii:paragraph:5-1:equation:i}{(i)}} holds, then
\[
  \Delta_A(I_B)
  \subset
  I_B\widehat\otimes1
  +
  A\widehat\otimes J_B.
\]
Therefore \(\Delta_A\) induces, after passage to the quotient, a
diagonal morphism
\[
  \overline\Delta:\overline A
  \longrightarrow
  \overline A\widehat\otimes_k\overline A.
\]
The maps \(\overline\Delta\) and \(\overline\varepsilon\) make
\(H\) a formal submonoid of \(G\). Since \(G\) is infinitesimal,
so is \(H\); \SGARefLink{sga3:VIIB:2.7:prop}{Proposition~2.7} then shows that \(H\) is a formal
subgroup of \(G\). By the definition of \(G/H\)
\textup{(cf.\ \SGARefLink{sga3:VIIB:2.4}{2.4})}, assertion \textup{\SGARefLink{sga3:vii:paragraph:5-1:equation:ii}{(ii)}} implies the final
assertion of the proposition.

It is clear that \textup{\SGARefLink{sga3:vii:paragraph:5-1:equation:ii}{(ii)}} implies \textup{\SGARefLink{sga3:vii:paragraph:5-1:equation:i}{(i)}}. The proof of
the converse occupies \SGARefLink{sga3:VIIB:5.1.1}{5.1.1}--\SGARefLink{sga3:VIIB:5.1.5}{5.1.5}.

% END INCLUDED SGA3_English_current_buildable_source_20260823/source/inputs/VII/tex/components/75_expose_VIIB_section51_quotient_criterion_en.tex

% BEGIN INCLUDED SGA3_English_current_buildable_source_20260823/source/inputs/VII/tex/components/76_expose_VIIB_section511_cogroups_of_pairs_en.tex
% SGA 3, Expose VII_B, English reconstruction -- Loop 1.
% Scope: Section 5.1.1 in full on authority local pages 83--84;
% stops immediately before Section 5.1.2.
% Combined-reader pages: 596--597. Re-edition printed pages: 585--586.
% Original-edition marginal locator in scope: 552.
% Editorial note in scope: 145.
% Authority: Polo--Gille Exp7B-13oct24.pdf,
% SHA-256 D303690E0195654BBE7C3DAA5D89EA23E24CD709A3575217234AC19AE8836694.
% Jacob Reinhold's CC BY 4.0 Markdown was consulted as comparison lineage.
% No genuine diagram occurs in this component; all mathematics is native LaTeX.

\subsubsection*{5.1.1}
\phantomsection
\label{sga3:VIIB:5.1.1}

Consider the following category \(\mathcal C\). An object of
\(\mathcal C\) is a pair \((A,J)\) consisting of a profinite
\(k\)-algebra \(A\) and a closed ideal \(J\subset A\). A morphism
\[
  \psi:(A,J)\longrightarrow(A',J')
\]
in \(\mathcal C\) is a continuous \(k\)-algebra homomorphism
\[
  A\longrightarrow A'
\]
that maps \(J\) into \(J'\).

The assignment
\[
  (A,J)
  \longmapsto
  \bigl(\operatorname{Spf}(A/J),\operatorname{Spf}(A)\bigr)
\]
is an antiequivalence from \(\mathcal C\) to the category of pairs
\((Z,Y)\) consisting of a formal \(k\)-variety \(Y\) and a formal
subvariety \(Z\subset Y\). A morphism
\[
  \varphi:(Z,Y)\longrightarrow(Z',Y')
\]
is a morphism of formal \(k\)-varieties
\[
  Y\longrightarrow Y'
\]
that maps \(Z\) into \(Z'\).

A cogroup structure on an object \((A,J)\) of \(\mathcal C\)
consists of a formal-group structure on
\(\operatorname{Spf}(A)\) satisfying the following conditions
\textup{(with the notation of \SGARefLink{sga3:VIIB:2.1}{2.1})}:
\[
\begin{aligned}
  \textup{(1)}\quad&
  \Delta_A(J)
  \subset
  J\widehat\otimes_kA
  +
  A\widehat\otimes_kJ,
  \\
  \textup{(2)}\quad&
  \varepsilon_A(J)=0,
  \\
  \textup{(3)}\quad&
  c_A(J)\subset J.
\end{aligned}
\]
These conditions also say that
\[
  H=\operatorname{Spf}(A/J)
\]
is a formal subgroup of
\[
  G=\operatorname{Spf}(A).
\]
\footnote{\label{sga3:VIIB:nde:145}If \((A',J')\) is a second
cogroup of \(\mathcal C\), corresponding to a pair
\(H'\subset G'\) of formal \(k\)-groups, then giving a cogroup
morphism
\[
  (A',J')\longrightarrow(A,J)
\]
is equivalent to giving a morphism of formal \(k\)-groups
\[
  G\longrightarrow G'
\]
that maps \(H\) into \(H'\).}

Suppose in addition that \(A\) is local; equivalently,
\(\operatorname{Spf}(A)\) is an infinitesimal formal group. If
\(J\ne A\), then conditions \textup{(2)} and \textup{(3)} follow
from \textup{(1)}. Indeed, a proper closed ideal \(J\) of \(A\)
is contained in the augmentation ideal \(I_A\), so
\textup{(2)} holds. Moreover,
\[
  M=\operatorname{Spf}(A/J)
\]
is a formal submonoid of \(G\). Since \(G\) is infinitesimal,
\SGARefLink{sga3:VIIB:2.7:prop}{Proposition~2.7} shows that \(M\) is a formal subgroup of \(G\);
hence \textup{(3)} also holds.

% END INCLUDED SGA3_English_current_buildable_source_20260823/source/inputs/VII/tex/components/76_expose_VIIB_section511_cogroups_of_pairs_en.tex

% BEGIN INCLUDED SGA3_English_current_buildable_source_20260823/source/inputs/VII/tex/components/77_expose_VIIB_section512_graded_profinite_algebras_lemma_en.tex
% SGA 3, Expose VII_B, English reconstruction -- Loop 1.
% Scope: Section 5.1.2 from its opening through the filtration lemma
% statement on authority local page 84; the proof follows in component 78.
% Combined-reader page: 597. Re-edition printed page: 586.
% Original-edition marginal locators in scope: 553--554.
% Editorial notes in scope: 146--147.
% Authority: Polo--Gille Exp7B-13oct24.pdf,
% SHA-256 D303690E0195654BBE7C3DAA5D89EA23E24CD709A3575217234AC19AE8836694.
% Jacob Reinhold's CC BY 4.0 Markdown was consulted as comparison lineage.
% No genuine diagram occurs in this component; all mathematics is native LaTeX.

\subsubsection*{5.1.2}
\phantomsection
\label{sga3:VIIB:5.1.2}

Denote by \(\mathbf{Alpg}/k\) the category of \emph{graded
profinite \(k\)-algebras}. An object of this category consists of
a sequence
\[
  A_0,A_1,\ldots,A_n,\ldots
\]
of pseudocompact \(k\)-modules, together with a profinite algebra
structure on
\[
  \prod_{n\geq0}A_n
\]
such that
\[
  A_nA_m\subset A_{m+n}.
\]
Here \(A_n\) is identified with a subset of
\(\prod_{i\geq0}A_i\) by means of the canonical injection. A
morphism
\[
  \psi:(A_n)\longrightarrow(B_n)
\]
is a sequence of continuous linear maps
\[
  \psi_n:A_n\longrightarrow B_n
\]
such that
\[
  \psi_{m+n}(aa')=\psi_m(a)\psi_n(a')
\]
for \(a\in A_m\) and \(a'\in A_n\).

\medskip
\noindent\textit{Definitions.}
Two graded profinite \(k\)-algebras \((A_n)\) and \((B_n)\)
clearly have a coproduct\footnote{%
\label{sga3:VIIB:nde:146}The editors replaced ``direct sum'' by
``coproduct''.}
in \(\mathbf{Alpg}/k\). Its \(n\)-th component is
\[
  \prod_{i=0}^{n}
  A_i\widehat\otimes_k B_{n-i},
\]
the product of the pseudocompact \(k\)-modules
\[
  A_i\widehat\otimes_k B_{n-i}.
\]
This coproduct will be denoted by
\[
  (A_n)\widehat\otimes_k(B_n).
\]

A cogroup structure on an object \((A_n)\) of
\(\mathbf{Alpg}/k\) is therefore the datum of continuous
\(k\)-linear maps
\[
  \Delta_n:A_n\longrightarrow
  \prod_{i=0}^{n}A_i\widehat\otimes_k A_{n-i},
  \qquad
  \varepsilon:A_0\longrightarrow k,
\]
which induce on \(\prod_{n\geq0}A_n\), upon setting
\(\varepsilon(A_i)=0\) for \(i\geq1\), a cogroup structure in
\(\mathbf{Alp}/k\).

Finally, for every object \((A,J)\) of \(\mathcal C\), denote by
\[
  \operatorname{Gr}_J(A)
\]
the graded profinite algebra associated with the filtration of
\(A\) by the closures \(\overline{J^n}\) of the powers of \(J\).
Thus
\[
  \operatorname{Gr}_J(A)_n
  =
  \overline{J^n}/\overline{J^{n+1}},
\]
and the multiplication of \(\operatorname{Gr}_J(A)\) is induced
by that of \(A\).

\medskip
\noindent\textbf{Lemma.}
\phantomsection
\label{sga3:VIIB:5.1.2:filtration-lemma}
\footnote{\label{sga3:VIIB:nde:147}In the original, the lemma is
stated when \(k\) is a field, and in that case the proof is left
to the reader.}
{\itshape Let \(U,V\) be two pseudocompact \(k\)-modules, with
\(U\) topologically flat, and let
\[
  U=U_0\supset U_1\supset\cdots,
  \qquad
  V=V_0\supset V_1\supset\cdots
\]
be two decreasing sequences of closed \(k\)-submodules. Filter
the completed tensor product
\[
  W=U\widehat\otimes_kV
\]
by means of the closed submodules
\[
  W_n
  =
  U_n\widehat\otimes_kV_0
  +
  U_{n-1}\widehat\otimes_kV_1
  +\cdots+
  U_0\widehat\otimes_kV_n.
\]
Suppose that every \(U_i/U_{i+1}\) is topologically flat over
\(k\), so that \(U/U_n\), and hence also \(U_n\), is
topologically flat for every \(n\). Then, for every \(n\), there
is an isomorphism
\[
  W_n/W_{n+1}
  \simeq
  \bigoplus_{i+j=n}
  (U_i/U_{i+1})\widehat\otimes_k(V_j/V_{j+1}).
\]}


% END INCLUDED SGA3_English_current_buildable_source_20260823/source/inputs/VII/tex/components/77_expose_VIIB_section512_graded_profinite_algebras_lemma_en.tex

% BEGIN INCLUDED SGA3_English_current_buildable_source_20260823/source/inputs/VII/tex/components/78_expose_VIIB_section512_filtration_lemma_proof_corollary_en.tex
% SGA 3, Expose VII_B, English reconstruction -- Loop 1.
% Scope: Proof and corollary completing Section 5.1.2 on authority
% local pages 84--85; stops immediately before Section 5.1.3.
% Combined-reader pages: 597--598. Re-edition printed pages: 586--587.
% Original-edition marginal locator in scope: 554.
% Editorial note in scope: 148.
% Authority: Polo--Gille Exp7B-13oct24.pdf,
% SHA-256 D303690E0195654BBE7C3DAA5D89EA23E24CD709A3575217234AC19AE8836694.
% Jacob Reinhold's CC BY 4.0 Markdown was consulted as comparison lineage.
% The genuine exact-rows proof diagram is an exact PDF-derived Loop 1 crop.

\noindent\textit{Proof.}
For every \(i,j\geq0\), put
\[
  W_{i,j}=U_i\widehat\otimes V_j,
  \qquad
  \overline W_{i,j}
  =
  (U_i/U_{i+1})\widehat\otimes_k(V_j/V_{j+1}).
\]
We prove by induction on \(n\) that the natural map
\[
  \pi_n:W_n\longrightarrow
  \bigoplus_{i+j=n}\overline W_{i,j}
\]
is surjective and that the inclusion
\[
  W_{n+1}\subset\ker(\pi_n)
\]
is an equality. For \(n=0\), the projection
\[
  \pi_0:
  U_0\widehat\otimes_kV_0
  \longrightarrow
  (U_0/U_1)\widehat\otimes_k(V_0/V_1)
\]
is surjective. Since \(U_0\), \(U_0/U_1\), and hence \(U_1\),
are topologically flat over \(k\), we have
\[
  \ker(\pi_0)
  =
  U_0\widehat\otimes_kV_1
  +
  U_1\widehat\otimes_kV_0,
\]
and moreover
\[
  \bigl(U_0\widehat\otimes_kV_1\bigr)
  \cap
  \bigl(U_1\widehat\otimes_kV_0\bigr)
  =
  U_1\widehat\otimes_kV_1.
\]

Suppose now that \(n>0\) and that the result has been proved for
\(n-1\). Put
\[
  M_0=U_0\widehat\otimes_kV_n,
  \qquad
  S_0=\sum_{i=1}^{n}U_i\widehat\otimes_kV_{n-i}.
\]
We have \(S_0\subset U_1\widehat\otimes_kV_0\). Applying the
preceding argument to \(V_0\supset V_n\) in place of
\(V_0\supset V_1\), we obtain
\[
  M_0\cap S_0
  \subset
  \bigl(U_0\widehat\otimes_kV_n\bigr)
  \cap
  \bigl(U_1\widehat\otimes_kV\bigr)
  =
  U_1\widehat\otimes_kV_n.
\]
It follows that
\[
  M_0\cap S_0=U_1\widehat\otimes_kV_n.
\]
Since \(W_n=M_0+S_0\), we obtain a commutative diagram with exact
rows, where
\[
  U'_i=U_{i+1},
  \qquad
  W'_{i,n-1-i}=W_{i+1,n-1-i}
  \quad(0\leq i\leq n-1):
\]
\begin{center}
  \phantomsection
  \label{sga3:VIIB:5.1.2:filtration-exact-rows-diagram}
  % SGA3-VII-DIAG-125
  % BEGIN INCLUDED SGA3_English_current_buildable_source_20260823/source/inputs/VII/native_diagrams/exp7/viib_15_95/diag_125.tex
% SGA3-VII-DIAG-125
% Authority: Exp7B-13oct24.pdf, local p. 85.
% Crop: figures/Exp7B_localp085_section512_filtration_exact_rows_diagram.png
% Locator: tex/components/78_expose_VIIB_section512_filtration_lemma_proof_corollary_en.tex, Loop 1 line 89.
\begin{tikzcd}[row sep=3em,column sep=1.8em]
0 \arrow[r] & S_0 \arrow[r] \arrow[d,"\pi'_{n-1}"']
  & W_n \arrow[r] \arrow[d,"\pi_n"']
  & (U_0/U_1)\widehat\otimes_k V_n
      \arrow[r] \arrow[d,"p"']
  & 0 \\
0 \arrow[r]
  & \displaystyle\bigoplus_{i=0}^{n-1}\overline W'_{i,n-1-i}
      \arrow[r]
  & \displaystyle\bigoplus_{i=0}^{n}\overline W_{i,n-i}
      \arrow[r]
  & \overline W_{0,n} \arrow[r]
  & 0
\end{tikzcd}

% END INCLUDED SGA3_English_current_buildable_source_20260823/source/inputs/VII/native_diagrams/exp7/viib_15_95/diag_125.tex

\end{center}

\phantomsection
\label{sga3:VIIB:5.1.2:diagram-source-notation-note}
\noindent\textit{Source note.}\enspace
The French diagram's lower middle term is printed
\(\bigoplus_{i=0}^{n}\overline W_{i,n-1}\); the natural map
\(\pi_n\) defined immediately above instead has target
\(\bigoplus_{i+j=n}\overline W_{i,j}\), which would give
\(\bigoplus_{i=0}^{n}\overline W_{i,n-i}\). Canon record
\texttt{SGA-SRC-000840} accepts this complementary index, which is
used in the diagram above while the printed reading is retained here.

Then \(p\) is surjective, with kernel
\[
  (U_0/U_1)\widehat\otimes_kV_{n+1}.
\]
Moreover, by the induction hypothesis applied to the sequence
\((U'_i)\), the map \(\pi'_{n-1}\) is surjective, with kernel
\[
  W'_n
  =
  \sum_{i=1}^{n}W_{i,n+1-i}.
\]
It follows that \(\pi_n\) is surjective and that
\[
  W_{n+1}\subset\ker(\pi_n)
\]
is an equality. This proves the lemma.

Let us return to an object \((A,J)\) of \(\mathcal C\). By
\SGARefLink{sga3:VIIB:0.2.G}{0.2.G}, the hypothesis that every
\[
  \overline{J^n}/\overline{J^{n+1}}
\]
is topologically flat over \(k\) is equivalent to saying that
\(\operatorname{Gr}_J(A)\) is topologically flat over \(k\).

\medskip
\noindent\textbf{Corollary.}
\phantomsection
\label{sga3:VIIB:5.1.2:corollary}
{\itshape Let \(\mathcal P\) be the full subcategory of
\(\mathcal C\) consisting of the objects \((A,J)\) such that
\(A\) and \(\operatorname{Gr}_J(A)\) are topologically flat over
\(k\). Then the functor
\[
  \mathcal P\longrightarrow\mathbf{Alpg}/k,
  \qquad
  (A,J)\longmapsto\operatorname{Gr}_J(A),
\]
commutes with finite coproducts and hence carries a cogroup of
\(\mathcal P\) to a cogroup of \(\mathbf{Alpg}/k\).

In particular, if \(k\) is a field, then for every cogroup
\((A,J)\) of \(\mathcal C\), the algebra
\(\operatorname{Gr}_J(A)\) is a cogroup of
\(\mathbf{Alpg}/k\).}
\footnote{\label{sga3:VIIB:nde:148}Thus, to a pair
\(H\subset G\) of formal \(k\)-groups one associates the
``formal completion \(\widehat G_H\) of \(G\) along \(H\)'',
which is a formal \(k\)-group. Moreover, we shall see in
\SGARefLink{sga3:VIIB:5.1.3}{5.1.3}--\SGARefLink{sga3:VIIB:5.1.4}{5.1.4} that the inclusion
\(\sigma:H\hookrightarrow\widehat G_H\) has a retraction
\(\pi:\widehat G_H\to H\), and that the formal \(k\)-group
\(N=\ker(\pi)\) identifies, as a formal variety, with the
completion of the homogeneous space \(G/H\) along the unit
section. This will be useful in \SGARefLink{sga3:VIIB:5.2.2}{5.2.2}.}

% END INCLUDED SGA3_English_current_buildable_source_20260823/source/inputs/VII/tex/components/78_expose_VIIB_section512_filtration_lemma_proof_corollary_en.tex

% BEGIN INCLUDED SGA3_English_current_buildable_source_20260823/source/inputs/VII/tex/components/79_expose_VIIB_section513_graded_cogroup_splitting_en.tex
% SGA 3, Expose VII_B, English reconstruction -- Loop 1.
% Scope: Section 5.1.3 in full on authority local pages 85--86;
% stops immediately before Section 5.1.4.
% Combined-reader pages: 598--599. Re-edition printed pages: 587--588.
% Editorial note 148 occurs immediately before this component and is
% preserved in component 78; no editorial note begins inside this component.
% Authority: Polo--Gille Exp7B-13oct24.pdf,
% SHA-256 D303690E0195654BBE7C3DAA5D89EA23E24CD709A3575217234AC19AE8836694.
% Jacob Reinhold's CC BY 4.0 Markdown was consulted as comparison lineage.
% The genuine kernel-pair diagram is an exact PDF-derived Loop 1 crop.

\subsubsection*{5.1.3}
\phantomsection
\label{sga3:VIIB:5.1.3}

Identify every profinite \(k\)-algebra \(\Gamma\) with the graded
profinite \(k\)-algebra \((\Gamma_n)_{n\geq0}\) such that
\[
  \Gamma_0=\Gamma,
  \qquad
  \Gamma_n=0\quad\text{if }n>0.
\]
In particular, if \((A_n)_{n\geq0}\) is a graded profinite
\(k\)-algebra, we shall regard \(A_0\), without distinction, as
either a profinite \(k\)-algebra or a graded profinite
\(k\)-algebra. Denote by
\[
  \rho:(A_n)\longrightarrow A_0
\]
the morphism of \(\mathbf{Alpg}/k\) such that
\[
  \rho_0=\operatorname{id}_{A_0},
  \qquad
  \rho_n=0\quad\text{if }n>0.
\]
Similarly,
\[
  \tau:A_0\longrightarrow(A_n)
\]
will denote the section of \(\rho\) such that
\[
  \tau_0=\operatorname{id}_{A_0},
  \qquad
  \tau_n=0\quad\text{if }n>0.
\]

Every cogroup structure on \((A_n)_{n\in\mathbb N}\) induces a
cogroup structure on \(A_0\) for which \(\rho\) and \(\tau\) are
cogroup homomorphisms. In this case, let \(I_0\) be the
augmentation ideal of \(A_0\), and put
\[
  A'_n=A_n/I_0A_n
  \qquad(n\geq0),
\]
so that \(A'_0=k\).

Then \((A'_n)_{n\in\mathbb N}\) is a cogroup in
\(\mathbf{Alpg}/k\). Indeed, since
\[
  A_0=I_0\oplus k\cdot1,
\]
the direct summand
\[
  I_0\widehat\otimes_kA_n\simeq I_0A_n
\]
of \(A_n\) is defined for every \(n\). Since \(\tau\) is a
section of \(\rho\), it follows from \SGARefLink{sga3:VIIB:2.4.A}{2.4.A} that the cogroup
\((A_n)_{n\in\mathbb N}\) is the ``semidirect coproduct'' of
\(A_0\) and the cogroup \((A'_n)_{n\in\mathbb N}\).

More precisely, \((A'_n)_{n\in\mathbb N}\) is isomorphic, as an
object of \(\mathbf{Alpg}/k\), to the kernel of the following
pair:
\begin{center}
  \phantomsection
  \label{sga3:VIIB:5.1.3:graded-cogroup-kernel-pair}
  % SGA3-VII-DIAG-126
  % BEGIN INCLUDED SGA3_English_current_buildable_source_20260823/source/inputs/VII/native_diagrams/exp7/viib_15_95/diag_126.tex
% SGA3-VII-DIAG-126
% Authority: Exp7B-13oct24.pdf, local p. 86.
% Crop: figures/Exp7B_localp086_graded_cogroup_kernel_pair.png
% Locator: tex/components/79_expose_VIIB_section513_graded_cogroup_splitting_en.tex, Loop 1 line 77.
\begin{tikzpicture}[baseline=(lhs.base)]
  \node (lhs) {$(A_n)$};
  \node[right=2.6cm of lhs] (rhs)
    {$(A_n)\widehat\otimes_k A_0$};
  \draw[-{Straight Barb[length=4pt,width=4pt]}]
    ([yshift=2.5pt]lhs.east) --
    node[above] {$\tau_1$}
    ([yshift=2.5pt]rhs.west);
  \draw[-{Straight Barb[length=4pt,width=4pt]}]
    ([yshift=-2.5pt]lhs.east) --
    node[below] {$(\operatorname{id}\widehat\otimes\rho)\Delta$}
    ([yshift=-2.5pt]rhs.west);
\end{tikzpicture}

% END INCLUDED SGA3_English_current_buildable_source_20260823/source/inputs/VII/native_diagrams/exp7/viib_15_95/diag_126.tex

\end{center}
Here
\[
  \Delta:(A_n)\longrightarrow(A_n)\widehat\otimes(A_n)
\]
is the comultiplication of \((A_n)\), and
\[
  \tau_1(x)=x\widehat\otimes1.
\]
Identifying \(A'_n\) with its image in \(A_n\), the map
\[
  A'_n\widehat\otimes_kA_0\longrightarrow A_n,
  \qquad
  a'_n\widehat\otimes a_0\longmapsto a'_na_0,
\]
is an isomorphism in \(\mathbf{Alpg}/k\).

This is not an isomorphism of cogroups. Rather,
\[
  \Delta(A')\subset A\widehat\otimes A',
\]
and
\[
  \bigl(\gamma'_\rho\widehat\otimes\operatorname{id}\bigr)
  \circ\Delta\big|_{A'}
  =
  \Delta',
\]
where \(\Delta'\) is the comultiplication of \(A'\) and
\(\gamma'_\rho:A\to A'\) is the projection
\textup{(cf.\ \SGARefLink{sga3:VIIB:2.4.A}{2.4.A})}.

% END INCLUDED SGA3_English_current_buildable_source_20260823/source/inputs/VII/tex/components/79_expose_VIIB_section513_graded_cogroup_splitting_en.tex

% BEGIN INCLUDED SGA3_English_current_buildable_source_20260823/source/inputs/VII/tex/components/80_expose_VIIB_section514_graded_generation_en.tex
% SGA 3, Expose VII_B, English reconstruction -- Loop 1.
% Scope: Section 5.1.4 from its opening through formula (3) on authority
% local pages 86--87; the application to Proposition 5.1 follows in
% component 81.
% Combined-reader pages: 599--600. Re-edition printed pages: 588--589.
% Original-edition marginal locator in scope: 555.
% Editorial note in scope: 149.
% Authority: Polo--Gille Exp7B-13oct24.pdf,
% SHA-256 D303690E0195654BBE7C3DAA5D89EA23E24CD709A3575217234AC19AE8836694.
% Jacob Reinhold's CC BY 4.0 Markdown was consulted as comparison lineage.
% The genuine generation diagram is an exact PDF-derived Loop 1 crop.

\subsubsection*{5.1.4}
\phantomsection
\label{sga3:VIIB:5.1.4}

Let \((A,J)\) be an object of \(\mathcal C\), and let
\[
  (A_n)=\operatorname{Gr}_J(A)
\]
be the associated object of \(\mathbf{Alpg}/k\); thus
\[
  A_n=\overline{J^n}/\overline{J^{n+1}}
  \qquad(n\geq0).
\]
It is clear that the algebra
\[
  \mathscr A=\prod_{n\geq0}A_n
\]
is generated by \(A_0\) and \(A_1\). In other words, for
\(n\geq1\), the multiplication map
\[
  A_1\widehat\otimes_{A_0}A_1
  \widehat\otimes_{A_0}\cdots
  \widehat\otimes_{A_0}A_1
  \longrightarrow A_n
  \tag{1}
  \label{sga3:VIIB:5.1.4:eq1}
\]
is surjective.

Suppose further that \((A,J)\) is a cogroup of \(\mathcal C\)
and that \(A\) and the quotients
\[
  \overline{J^n}/\overline{J^{n+1}}
\]
are topologically flat over \(k\). By \SGARefLink{sga3:VIIB:5.1.2:corollary}{Corollary~5.1.2},
\(\operatorname{Gr}_J(A)\) is then a cogroup of
\(\mathbf{Alpg}/k\). Hence, by \SGARefLink{sga3:VIIB:5.1.3}{5.1.3}, if one puts
\[
\begin{split}
  A'_n
  =
  \biggl\{
    x\in\overline{J^n}/\overline{J^{n+1}}
    \ \biggm|\
    \Delta(x)-x\widehat\otimes1
    \in
    \bigoplus_{i=1}^{n}
    &\bigl(
      \overline{J^{\,n-i}}/\overline{J^{\,n-i+1}}
    \bigr)
    \widehat\otimes_k
    \\
    &\bigl(
      \overline{J^i}/\overline{J^{i+1}}
    \bigr)
  \biggr\},
\end{split}
\tag{2}
\label{sga3:VIIB:5.1.4:eq2}
\]
then the map
\[
  A'_n\widehat\otimes_kA_0\longrightarrow A_n,
  \qquad
  a'_n\widehat\otimes a_0\longmapsto a'_na_0,
\]
is an isomorphism in \(\mathbf{Alpg}/k\).\footnote{%
\label{sga3:VIIB:nde:149}Let \(G\), \(H\), and
\(\widehat G_H\) be the formal \(k\)-groups corresponding to
\(A\), \(A_0=A/J\), and \(\operatorname{Gr}_J(A)\),
respectively. Then
\(\tau:A_0\hookrightarrow\operatorname{Gr}_J(A)\) corresponds
to a retraction
\(\pi:\widehat G_H\to H\) of the inclusion
\(H\hookrightarrow\widehat G_H\), and the preceding assertion
means that \(\widehat G_H\) is the semidirect product of
\(N=\ker(\pi)\) by \(H\).}

\phantomsection
\label{sga3:VIIB:5.1.4:eq2-source-tensor-note}
\noindent\textit{Source note.}\enspace
In the displayed element-wise description of this map, the
French PDF prints
\(a'_n\widehat\otimes a'_n\mapsto a'_na_0\).
The domain and the product require
\(a'_n\widehat\otimes a_0\mapsto a'_na_0\), which is used here.

Let \(I_0\) denote the augmentation ideal of \(A_0\). It follows
from \textup{\SGARefLink{sga3:VIIB:5.1.4:eq1}{(1)}} and the following commutative diagram, in which
\[
  A_1'{}^{\widehat\otimes n}
  =
  \underbrace{
    A'_1\widehat\otimes_k\cdots\widehat\otimes_kA'_1
  }_{n\text{ factors}},
\]
that:
\begin{center}
  \phantomsection
  \label{sga3:VIIB:5.1.4:degree-one-generation-diagram}
  % SGA3-VII-DIAG-127
  % BEGIN INCLUDED SGA3_English_current_buildable_source_20260823/source/inputs/VII/native_diagrams/exp7/viib_15_95/diag_127.tex
% SGA3-VII-DIAG-127
% Authority: Exp7B-13oct24.pdf, local p. 87.
% Crop: figures/Exp7B_localp087_degree_one_generation_diagram.png
% Locator: tex/components/80_expose_VIIB_section514_graded_generation_en.tex, Loop 1 line 113.
\begin{tikzcd}[row sep=3em,column sep=2.1em]
A_1\widehat\otimes_{A_0}\cdots\widehat\otimes_{A_0}A_1
  \arrow[d,"m"']
  & A_1'{}^{\widehat\otimes n}\widehat\otimes_k A_0
      \arrow[l,"{\sim}"']
      \arrow[d,"{m'\widehat\otimes\operatorname{id}}"']
  & \bigl(A_1'{}^{\widehat\otimes n}\widehat\otimes_k I_0\bigr)
      \oplus A_1'{}^{\widehat\otimes n}
      \arrow[l,"{\sim}"']
      \arrow[d,"{(m'\widehat\otimes\operatorname{id})\oplus m'}"] \\
A_n
  & A'_n\widehat\otimes_k A_0 \arrow[l,"{\sim}"']
  & (A'_n\widehat\otimes_k I_0)\oplus A'_n
      \arrow[l,"{\sim}"']
\end{tikzcd}

% END INCLUDED SGA3_English_current_buildable_source_20260823/source/inputs/VII/native_diagrams/exp7/viib_15_95/diag_127.tex

\end{center}
the multiplication map
\[
  m':
  A'_1\widehat\otimes_k\cdots\widehat\otimes_kA'_1
  \longrightarrow A'_n
  \tag{3}
  \label{sga3:VIIB:5.1.4:eq3}
\]
is surjective. In other words, the profinite \(k\)-algebra
\[
  \mathscr A'=\prod_{n\geq0}A'_n
\]
is generated by its degree-one terms.

% END INCLUDED SGA3_English_current_buildable_source_20260823/source/inputs/VII/tex/components/80_expose_VIIB_section514_graded_generation_en.tex

% BEGIN INCLUDED SGA3_English_current_buildable_source_20260823/source/inputs/VII/tex/components/81_expose_VIIB_section514_proposition51_completion_en.tex
% SGA 3, Expose VII_B, English reconstruction -- Loop 1.
% Scope: The remainder of Section 5.1.4, completing the proof of
% Proposition 5.1 and its supplement, on authority local pages 87--88;
% stops immediately before Section 5.1.5.
% Combined-reader pages: 600--601. Re-edition printed pages: 589--590.
% Original-edition marginal locator in scope: 556.
% Editorial notes in scope: 150--151.
% Authority: Polo--Gille Exp7B-13oct24.pdf,
% SHA-256 D303690E0195654BBE7C3DAA5D89EA23E24CD709A3575217234AC19AE8836694.
% Jacob Reinhold's CC BY 4.0 Markdown was consulted as comparison lineage.
% The genuine kernel-pair and exact commutative diagrams are exact
% PDF-derived Loop 1 crops.

Let us now return to hypothesis \textup{\SGARefLink{sga3:vii:paragraph:5-1:equation:i}{(i)}} of
\SGARefLink{sga3:VIIB:5.1:proposition}{Proposition~5.1}.\footnote{%
\label{sga3:VIIB:nde:150}The editors expanded the original in
what follows and moved to the end the ``supplement''
\(\overline{I_B^n}=B\cap\overline{J^n}\), which is not needed
to prove \SGARefLink{sga3:VIIB:5.1:proposition}{Proposition~5.1}.}
Let \(G\) be a topologically flat infinitesimal formal
\(k\)-group, \(A\) its affine algebra, \(I_A\) the augmentation
ideal of \(A\), \(B\) a closed subalgebra of \(A\), and
\[
  J=AI_B,
  \qquad
  I_B=B\cap I_A.
\]
Let \(H\) be the formal subgroup \(\operatorname{Spf}(A/J)\),
and let
\[
  \pi:A\longrightarrow A/J
\]
be the projection. Suppose that
\[
  A/\overline{J^n}
\]
is topologically flat over \(k\) for every \(n\geq1\), and that
\(B\) is contained in the kernel \(\widetilde B\) of the
following pair:
\begin{center}
  \phantomsection
  \label{sga3:VIIB:5.1.4:quotient-kernel-pair}
  % SGA3-VII-DIAG-128
  % BEGIN INCLUDED SGA3_English_current_buildable_source_20260823/source/inputs/VII/native_diagrams/exp7/viib_15_95/diag_128.tex
% SGA3-VII-DIAG-128
% Authority: Exp7B-13oct24.pdf, local p. 87.
% Crop: figures/Exp7B_localp087_quotient_kernel_pair.png
% Locator: tex/components/81_expose_VIIB_section514_proposition51_completion_en.tex, Loop 1 line 43.
\begin{tikzpicture}[baseline=(lhs.base)]
  \node (lhs) {$A$};
  \node[right=3cm of lhs] (rhs)
    {$A\widehat\otimes_k(A/J).$};
  \draw[-{Straight Barb[length=4pt,width=4pt]}]
    ([yshift=2.5pt]lhs.east) --
    node[above] {$\tau_1$}
    ([yshift=2.5pt]rhs.west);
  \draw[-{Straight Barb[length=4pt,width=4pt]}]
    ([yshift=-2.5pt]lhs.east) --
    node[below] {$(\operatorname{id}_A\widehat\otimes_k\pi)\Delta_A$}
    ([yshift=-2.5pt]rhs.west);
\end{tikzpicture}

% END INCLUDED SGA3_English_current_buildable_source_20260823/source/inputs/VII/native_diagrams/exp7/viib_15_95/diag_128.tex

\end{center}

Put
\[
  (A_n)=\operatorname{Gr}_J(A),
\]
let
\[
  I_0=I_A/J
\]
be the augmentation ideal of \(A_0=A/J\), and define
\((A'_n)_{n\in\mathbb N}\) as in \textup{(2)} above. Let
\((B_n)_{n\in\mathbb N}\), respectively
\((\widetilde B_n)_{n\in\mathbb N}\), be the object of
\(\mathbf{Alpg}/k\) associated with the filtration of \(B\),
respectively of \(\widetilde B\), induced by that of \(A\); these
filtrations are defined by the ideals
\[
  B\cap\overline{J^n},
  \qquad\text{respectively}\qquad
  \widetilde B\cap\overline{J^n}.
\]
It is clear that
\[
  B_n\subset\widetilde B_n\subset A'_n
  \qquad\text{for every }n,
\]
and that
\[
  \mathscr B=\prod_{n\geq0}B_n
  \subset
  \widetilde{\mathscr B}
  =
  \prod_{n\geq0}\widetilde B_n
\]
are subalgebras of
\[
  \mathscr A'=\prod_{n\geq0}A'_n.
\]

On the other hand, \(J\), respectively \(\overline{J^2}\), is
the image in \(A\) of
\[
  I_B\widehat\otimes_kA,
  \qquad\text{respectively}\qquad
  \overline{I_B^2}\widehat\otimes_kA.
\]
Consequently, the map
\[
  (I_B/\overline{I_B^2})\widehat\otimes_kA
  \longrightarrow
  J/\overline{J^2}
  =
  A_1
  \simeq
  A'_1\widehat\otimes_kA_0
\]
is surjective, and it factors through
\[
  (I_B/\overline{I_B^2})\widehat\otimes_kA_0.
\]
Since
\[
  A_0=k\cdot1\oplus I_0,
\]
the following commutative exact diagram shows that \(A'_1\) is
the image of \(I_B/\overline{I_B^2}\):
\begin{center}
  \phantomsection
  \label{sga3:VIIB:5.1.4:degree-one-image-diagram}
  % SGA3-VII-DIAG-129
  % BEGIN INCLUDED SGA3_English_current_buildable_source_20260823/source/inputs/VII/native_diagrams/exp7/viib_15_95/diag_129.tex
% SGA3-VII-DIAG-129
% Authority: Exp7B-13oct24.pdf, local p. 88.
% Crop: figures/Exp7B_localp088_degree_one_image_diagram.png
% Locator: tex/components/81_expose_VIIB_section514_proposition51_completion_en.tex, Loop 1 line 115.
\begin{tikzcd}[row sep=2.7em,column sep=3em]
(I_B/\overline{I_B^2})\widehat\otimes_k A_0 \arrow[d]
  & \bigl((I_B/\overline{I_B^2})\widehat\otimes_k I_0\bigr)
      \oplus(I_B/\overline{I_B^2})
      \arrow[l,"{\sim}"'] \arrow[d] \\
J/\overline{J^2} \arrow[d]
  & (A'_1\widehat\otimes_k I_0)\oplus A'_1
      \arrow[l,"{\sim}"'] \arrow[d] \\
0 & 0
\end{tikzcd}

% END INCLUDED SGA3_English_current_buildable_source_20260823/source/inputs/VII/native_diagrams/exp7/viib_15_95/diag_129.tex

\end{center}
Thus
\[
  B_1=A'_1.
\]
Since \(\mathscr A'\) is generated by \(A'_1\), it follows that,
for every \(n\geq1\),
\[
  A'_n\subset B_n\subset\widetilde B_n\subset A'_n,
\]
and hence
\[
  B_n=\widetilde B_n=A'_n.
\]
\footnote{\label{sga3:VIIB:nde:151}With the notation of
editorial note~149, this implies that the formal completion of
\(G/H\) along the unit section, whose affine algebra is
\(\prod_nB_n\), is isomorphic as a formal variety to the formal
\(k\)-group \(N\).}

Finally, because \(G\) is infinitesimal,
\[
  \mathfrak r(A)=I_A,
\]
and hence the ideals \(\overline{I_A^n}\) tend to zero
\textup{(cf.\ \SGARefLink{sga3:VIIB:0.1.2}{0.1.2})}. A fortiori, the ideals
\(\overline{J^n}\) tend to zero, so the induced filtrations on
\(\widetilde B\) and \(B\) are separated. Moreover, since \(B\)
is a closed subalgebra of \(A\), it is complete for the topology
defined by the ideals \(B\cap\overline{J^n}\). It therefore
follows from \textup{[CA]}, \S~V.7, Lemma~1
\textup{(see also \SGARefLink{sga3:VIIB:5.1.5}{5.1.5} below)} that
\[
  B=\widetilde B.
\]
This completes the proof of \SGARefLink{sga3:VIIB:5.1:proposition}{Proposition~5.1}.

There is also the following supplement. For every \(n\),
\[
  \overline{J^n}=A\overline{I_B^n}
\]
is the image in \(A\) both of
\[
  A\widehat\otimes_B\overline{I_B^n}
\]
and of
\[
  A\widehat\otimes_B(B\cap\overline{J^n}).
\]
By hypothesis, the affine algebra \(A/J\) of the formal subgroup
\(H\) is topologically flat over \(k\). Thus, by \SGARefLink{sga3:VIIB:2.4:theorem}{Theorem~2.4},
the morphism
\[
  G=\operatorname{Spf}(A)
  \longrightarrow
  G/H=\operatorname{Spf}(B)
\]
is surjective and topologically flat. Consequently,
\[
  A\widehat\otimes_B\overline{I_B^n}
  =
  \overline{J^n}
  =
  A\widehat\otimes_B(B\cap\overline{J^n}),
\]
which implies
\[
  \overline{I_B^n}=B\cap\overline{J^n}
  \qquad\text{for every }n.
\]
This also follows from the surjectivity of the maps
\[
  \overline{I_B^n}/\overline{I_B^{n+1}}
  \longrightarrow
  A'_n
  =
  (B\cap\overline{J^n})/
  (B\cap\overline{J^{n+1}})
\]
and from \SGARefLink{sga3:VIIB:5.1.5}{5.1.5}\textup{\SGARefLink{sga3:vii:paragraph:5-1:equation:ii}{(ii)}} below, applied to
\[
  B'_n=\overline{I_B^n},
  \qquad
  B_n=B\cap\overline{J^n}.
\]

\phantomsection
\label{sga3:VIIB:5.1.4:supplement-index-source-note}
\noindent\textit{Source note.}\enspace
In the preceding \(n\)-indexed display, the French PDF prints
\(A'_i\). The surrounding numerators and denominators, and the
application to \(B'_n\) and \(B_n\), require \(A'_n\), which is
used here.

% END INCLUDED SGA3_English_current_buildable_source_20260823/source/inputs/VII/tex/components/81_expose_VIIB_section514_proposition51_completion_en.tex

% BEGIN INCLUDED SGA3_English_current_buildable_source_20260823/source/inputs/VII/tex/components/81b_expose_VIIB_section515_filtered_groups_lemma_en.tex
% SGA 3, Expose VII_B, English reconstruction -- Loop 1.
% Scope: Section 5.1.5 in full on authority local pages 88--89;
% stops immediately before Section 5.2.
% Combined-reader pages: 601--602. Re-edition printed pages: 590--591.
% Editorial note in scope: 152.
% Authority: Polo--Gille Exp7B-13oct24.pdf,
% SHA-256 D303690E0195654BBE7C3DAA5D89EA23E24CD709A3575217234AC19AE8836694.
% Jacob Reinhold's CC BY 4.0 Markdown was consulted as comparison lineage.
% No genuine diagram occurs in this component; all mathematics is native LaTeX.

\subsubsection*{5.1.5}
\phantomsection
\label{sga3:VIIB:5.1.5}

\medskip
\noindent\textbf{Lemma 5.1.5.}
\phantomsection
\label{sga3:VIIB:5.1.5:filtered-groups-lemma}
\footnote{\label{sga3:VIIB:nde:152}The original stated only
part~\textup{\SGARefLink{sga3:VIIB:5.1.5:item:ii}{(ii)}}. For the reader's convenience, the editors
stated in part~\textup{\SGARefLink{sga3:VIIB:5.1.5:item:i}{(i)}} Lemma~1 of \textup{[CA]},
\S~V.7.}
{\itshape
\begin{enumerate}
\item[\textup{\SGAInlineRefTarget{sga3:VIIB:5.1.5:item:i}(i)}]
Let \(M\) and \(N\) be two abelian groups filtered by decreasing
sequences of subgroups
\[
  (M_n)_{n\in\mathbb Z},
  \qquad
  (N_n)_{n\in\mathbb Z}.
\]
Suppose that
\[
  \bigcup_{n\in\mathbb Z}M_n=M,
  \qquad
  \bigcup_{n\in\mathbb Z}N_n=N,
\]
that
\[
  \bigcap_{n\in\mathbb Z}M_n=0,
  \qquad
  \bigcap_{n\in\mathbb Z}N_n=0,
\]
and that \(M\) is complete for the topology defined by the
\(M_n\). Let
\[
  f:M\longrightarrow N
\]
be a homomorphism of filtered groups.

\begin{enumerate}
\item[\textup{a)}]
If \(f\) induces a surjection on the associated graded groups,
then \(f\) is surjective and \(N\) is complete for the topology
defined by the \(N_n\).

\item[\textup{b)}]
If \(f\) induces an injection on the associated graded groups,
then \(f\) is injective.
\end{enumerate}

\item[\textup{\SGAInlineRefTarget{sga3:VIIB:5.1.5:item:ii}(ii)}]
Let \(B\) be an abelian group, and let
\[
  B=B'_0\supset B'_1\supset\cdots,
  \qquad
  B=B_0\supset B_1\supset\cdots
\]
be two separated filtrations of \(B\) by subgroups such that
\[
  B'_n\subset B_n
  \qquad\text{for every }n.
\]
Suppose that \(B\) is complete for the topology defined by the
filtration \((B'_n)\). If the map
\[
  B'_i/B'_{i+1}
  \longrightarrow
  B_i/B_{i+1}
\]
is surjective for every \(i\), then
\[
  B'_n=B_n
  \qquad\text{for every }n.
\]
\end{enumerate}
}

\noindent\textit{Proof.}
Part~\textup{\SGARefLink{sga3:VIIB:5.1.5:item:i}{(i)}} is Lemma~1 of \textup{[CA]}, \S~V.7
\textup{(see also [BAC], III, \S~\SGARefLink{sga3:VIIB:2.8}{2.8})}. Part~\textup{\SGARefLink{sga3:VIIB:5.1.5:item:ii}{(ii)}}
follows by taking
\[
  M=B'_n\supset B'_{n+1}\supset\cdots,
  \qquad
  N=B_n\supset B_{n+1}\supset\cdots.
\]

% END INCLUDED SGA3_English_current_buildable_source_20260823/source/inputs/VII/tex/components/81b_expose_VIIB_section515_filtered_groups_lemma_en.tex

% BEGIN INCLUDED SGA3_English_current_buildable_source_20260823/source/inputs/VII/tex/components/82_expose_VIIB_section52_height_theorem_en.tex
% SGA 3, Expose VII_B, English reconstruction -- Loop 1.
% Scope: Section 5.2 opening, notation, and the Dieudonne--Cartier theorem
% on authority local pages 89--90; stops before Section 5.2.1.
% Combined-reader pages: 602--603. Re-edition printed pages: 591--592.
% Original-edition marginal locators in scope: 557--558.
% Editorial notes in scope: 153--154.
% Authority: Polo--Gille Exp7B-13oct24.pdf,
% SHA-256 D303690E0195654BBE7C3DAA5D89EA23E24CD709A3575217234AC19AE8836694.
% Jacob Reinhold's CC BY 4.0 Markdown was consulted as comparison lineage.
% No genuine diagram occurs in this component; all mathematics is native LaTeX.

\subsection*{5.2}
\phantomsection
\label{sga3:VIIB:5.2}

Throughout the rest of \SGARefLink{sga3:VIIB:5}{Section~5}, \(k\) denotes a perfect field
of characteristic \(p>0\).

Put
\[
  \overline{\mathbb N}=\mathbb N\cup\{\infty\}.
\]
If \(B\) is a profinite \(k\)-algebra and
\(r\in\overline{\mathbb N}\), denote by
\[
  \bigl((x^{p^r})\bigr)_{x\in\mathfrak r(B)}
\]
the closed ideal of \(B\) generated by the elements \(x^{p^r}\),
where \(x\) ranges over the radical \(\mathfrak r(B)\) of \(B\).
If \(r=\infty\), use the same notation with the convention that
\[
  \bigl((x^{p^\infty})\bigr)_{x\in\mathfrak r(B)}
\]
is the zero ideal. In both cases, \(B_r\) denotes the quotient
\[
  B_r
  =
  B\big/\bigl((x^{p^r})\bigr)_{x\in\mathfrak r(B)}.
\]

We say that \(B\) is \emph{of height at most \(r\)} if
\[
  \bigl((x^{p^r})\bigr)_{x\in\mathfrak r(B)}
\]
is the zero ideal. If this holds and \(r\) is finite, we say that
\(B\) is \emph{of finite height}.

Consider in particular the case
\[
  B=k[[\omega]],
\]
where \(\omega\) is a pseudocompact \(k\)-vector space
\textup{(cf.\ \SGARefLink{sga3:VIIB:1.2.5}{1.2.5})}.\footnote{%
\label{sga3:VIIB:nde:153}In the original, the author uses
``linearly compact vector space,'' which is equivalent to
``pseudocompact vector space'' \textup{(cf.\ [BAC],
\S III.2, Exercises 15\,a), 19\,a), and 20\,d))}. The editors
preferred to retain the terminology ``pseudocompact,'' used up to
this point.}
We then say that \(B\) is a \emph{formal power series algebra}
and that \(B_r\) is a \emph{truncated formal power series algebra}
\(\bigl(r\in\overline{\mathbb N}\bigr)\); thus, by convention,
\[
  B=B_\infty
\]
is also called ``truncated.'' If \(B=k[[\omega]]\), we also write
\[
  \bigl((x^{p^r})\bigr)_{x\in\omega}
\]
instead of
\[
  \bigl((x^{p^r})\bigr)_{x\in\mathfrak r(B)}.
\]

\medskip
\noindent\textbf{Notation.}
Let \(\omega\) be a pseudocompact \(k\)-vector space filtered by
an increasing sequence of closed vector subspaces
\[
  0=\omega_0\subset\omega_1\subset\omega_2
  \subset\omega_3\subset\cdots.
\]
\begin{enumerate}
\item[\textup{\SGAInlineRefTarget{sga3:VIIB:5.2:item:a}(a)}]
The closed ideal of \(k[[\omega]]\) generated by the elements
\(x^{p^r}\), where \(r\) ranges over \(\mathbb N\) and
\(x\) ranges over \(\omega_r\), will be denoted by
\[
  \bigl((x^{p^r})\bigr)_{x\in\omega_r}.
\]

\item[\textup{(b)}]
For \(r\in\mathbb N\), denote by \({}^r\omega\) the filtered
pseudocompact vector space defined by
\[
  {}^r\omega_i=\omega_i\quad\text{if }i<r,
  \qquad
  {}^r\omega_i=\omega\quad\text{if }i\geq r.
\]
\end{enumerate}

\phantomsection
\label{sga3:VIIB:5.2:notation-source-note}
\noindent\textit{Source note.}\enspace
In Notation \textup{\SGARefLink{sga3:VIIB:5.2:item:a}{(a)}}, the French PDF prints the subscript
\(r\in\omega_r\). The preceding defining sentence and every
subsequent use require \(x\in\omega_r\), which is used here.

\medskip
\noindent\textbf{Theorem (Dieudonn\'e--Cartier).}
\phantomsection
\label{sga3:VIIB:5.2:theorem}
{\itshape Let
\[
  H\longrightarrow G
\]
be a monomorphism of infinitesimal formal groups over a perfect
field \(k\) of characteristic \(p>0\). Let \(B\) be the affine
algebra of the homogeneous space \(G/H\), and suppose that one of
the following three conditions holds:\footnote{%
\label{sga3:VIIB:nde:154}The editors replaced \(H\backslash G\)
by \(G/H\), and did the same in the proof; they also added
condition \textup{\SGARefLink{sga3:VIIB:5.2:theorem:item:iii}{(iii)}}.}
\begin{enumerate}
\item[\textup{(i)}]
\(B\) has finite height \textup{(in particular, this holds if
\(G\) has finite height)}.

\item[\textup{\SGAInlineRefTarget{sga3:VIIB:5.2:theorem:item:ii}(ii)}]
\(B\) is a complete Noetherian local ring.

\item[\textup{\SGAInlineRefTarget{sga3:VIIB:5.2:theorem:item:iii}(iii)}]
\(B\) is a reduced ring.
\end{enumerate}
Then \(B\) is isomorphic to the completed tensor product of a
finite family of truncated formal power series algebras.}

The proof of this theorem occupies \SGARefLink{sga3:VIIB:5.2.1:section-proposition}{5.2.1}--\SGARefLink{sga3:VIIB:5.2.5}{5.2.5}.

% END INCLUDED SGA3_English_current_buildable_source_20260823/source/inputs/VII/tex/components/82_expose_VIIB_section52_height_theorem_en.tex

% BEGIN INCLUDED SGA3_English_current_buildable_source_20260823/source/inputs/VII/tex/components/83_expose_VIIB_section521_reduction_proposition_en.tex
% SGA 3, Expose VII_B, English reconstruction -- Loop 1.
% Scope: Section 5.2.1 opening and reduction proposition in full
% on authority local page 90; stops immediately before Section 5.2.1.A.
% Combined-reader page: 603. Re-edition printed page: 592.
% Original-edition marginal locator in scope: 558.
% Authority: Polo--Gille Exp7B-13oct24.pdf,
% SHA-256 D303690E0195654BBE7C3DAA5D89EA23E24CD709A3575217234AC19AE8836694.
% Jacob Reinhold's CC BY 4.0 Markdown was consulted as comparison lineage.
% No genuine diagram occurs in this component; all mathematics is native LaTeX.

\subsubsection*{5.2.1}
\phantomsection
\label{sga3:VIIB:5.2.1}

Let \(A\) be the affine algebra of \(G\), let \(I_A\) be its
augmentation ideal, and put
\[
  I=B\cap I_A.
\]
By \SGARefLink{sga3:VIIB:2.4}{2.4},
\[
  H=\operatorname{Spf}\bigl(A/\overline{AI}\bigr)
\]
and
\[
  B
  =
  \left\{
    x\in A\ \middle|\
    \Delta(x)-x\widehat\otimes1
    \in A\widehat\otimes\overline{AI}
  \right\}.
\]
\phantomsection
\label{sga3:VIIB:5.2.1:tensor-order-source-note}
\noindent\textit{Source note.}\enspace
The French PDF prints
\(\Delta(x)-1\widehat\otimes x\) here. The right-quotient
convention fixed in \SGARefLink{sga3:VIIB:5.0}{5.0} and \SGARefLink{sga3:VIIB:5.1:proposition}{Proposition~5.1}, together with the
defining condition
\((\operatorname{id}_A\widehat\otimes\pi)\Delta(x)
=x\widehat\otimes1\), requires
\(\Delta(x)-x\widehat\otimes1\), which is used above.

Put
\[
  \omega=I/\overline{I^2}.
\]
For each \(r\), denote by \(I_r\) the closed subideal of \(I\)
consisting of the elements \(x\) such that
\[
  x^{p^r}=0,
\]
and denote by \(\omega_r\) the canonical image of \(I_r\) in
\(\omega\). We shall prove the following proposition.

\medskip
\noindent\textbf{Proposition.}
\phantomsection
\label{sga3:VIIB:5.2.1:section-proposition}
{\itshape If there exists a continuous section
\[
  \sigma:\omega\longrightarrow I
\]
of the projection \(I\to I/\overline{I^2}\) such that
\[
  \sigma(\omega_r)\subset I_r
\]
for every \(r\), then \(B\) is isomorphic to
\[
  k[[\omega]]
  \big/
  \bigl((x^{p^r})\bigr)_{x\in\omega_r}.
\]}

Indeed, such a section extends to a continuous morphism
\[
  k[[\omega]]\longrightarrow B,
\]
which factors through
\[
  B'
  =
  k[[\omega]]
  \big/
  \bigl((x^{p^r})\bigr)_{x\in\omega_r}.
\]
In \SGARefLink{sga3:VIIB:5.2.2}{5.2.2}--\SGARefLink{sga3:VIIB:5.2.5}{5.2.5} we prove that the resulting morphism
\[
  \varphi:
  B'
  =
  k[[\omega]]
  \big/
  \bigl((x^{p^r})\bigr)_{x\in\omega_r}
  \longrightarrow B
\]
is an isomorphism.

% END INCLUDED SGA3_English_current_buildable_source_20260823/source/inputs/VII/tex/components/83_expose_VIIB_section521_reduction_proposition_en.tex

% BEGIN INCLUDED SGA3_English_current_buildable_source_20260823/source/inputs/VII/tex/components/84_expose_VIIB_section521A_pseudobasis_splittings_en.tex
% SGA 3, Expose VII_B, English reconstruction -- Loop 1.
% Scope: Section 5.2.1.A in full on authority local pages 90--91;
% stops immediately before Section 5.2.1.B.
% Combined-reader pages: 603--604. Re-edition printed pages: 592--593.
% Editorial note in scope: 155.
% Authority: Polo--Gille Exp7B-13oct24.pdf,
% SHA-256 D303690E0195654BBE7C3DAA5D89EA23E24CD709A3575217234AC19AE8836694.
% Jacob Reinhold's CC BY 4.0 Markdown was consulted as comparison lineage.
% The two genuine diagrams are exact PDF-derived Loop-1 PNG crops.

\paragraph{\SGARefLink{sga3:VIIB:5.2.1.A}{5.2.1.A}.}
\phantomsection
\label{sga3:VIIB:5.2.1.A}
\footnote{\label{sga3:VIIB:nde:155}The editors added
\SGARefLink{sga3:VIIB:5.2.1.A}{5.2.1.A} and \SGARefLink{sga3:VIIB:5.2.1.B}{5.2.1.B}.}
For every \(r\in\mathbb N\), put
\[
  \mathcal I_r=\overline{I^2}+I_r,
\]
so that
\[
  \mathcal I_r/\overline{I^2}\simeq\omega_r.
\]
There is then a commutative diagram with exact rows:
\begin{center}
  \phantomsection
  \label{sga3:VIIB:5.2.1.A:exact-rows}
  % SGA3-VII-DIAG-130
  % BEGIN INCLUDED SGA3_English_current_buildable_source_20260823/source/inputs/VII/native_diagrams/exp7/viib_15_95/diag_130.tex
% SGA3-VII-DIAG-130
% Authority: Exp7B-13oct24.pdf, local p. 90.
% Crop: figures/Exp7B_localp090_521A_exact_rows.png
% Locator: tex/components/84_expose_VIIB_section521A_pseudobasis_splittings_en.tex, Loop 1 line 28.
\begin{tikzcd}[row sep=3em,column sep=2.6em]
0 \arrow[r] & \mathcal I_{r-1} \arrow[r] \arrow[d]
  & \mathcal I_r \arrow[r] \arrow[d]
  & \mathcal I_r/\mathcal I_{r-1}
      \arrow[r] \arrow[d,"\wr"]
  & 0 \\
0 \arrow[r] & \omega_{r-1} \arrow[r]
  & \omega_r \arrow[r]
  & \omega_r/\omega_{r-1} \arrow[r]
  & 0
\end{tikzcd}

% END INCLUDED SGA3_English_current_buildable_source_20260823/source/inputs/VII/native_diagrams/exp7/viib_15_95/diag_130.tex

\end{center}

Since \(k\) is a field, the rows split. Thus a pseudobasis
\(\mathscr B_{r-1}\) of \(\mathcal I_{r-1}\) can be completed
to a pseudobasis
\[
  \mathscr B_{r-1}\cup\mathscr B'_r
\]
of \(\mathcal I_r\). The closed subspace \(S_r\) with
pseudobasis \(\mathscr B'_r\) is then a complement of
\(\mathcal I_{r-1}\) in \(\mathcal I_r\), and the projection
\[
  \pi:I\longrightarrow\omega
\]
induces an isomorphism from \(S_r\) onto a complement
\(\omega'_r\) of \(\omega_{r-1}\) in \(\omega_r\).

Let \(\mathcal I_\infty\) be the closed ideal
\[
  \mathcal I_\infty
  =
  \overline{\bigcup_r\mathcal I_r}.
\]
It likewise has a complement \(S_\infty\) in \(I\), and
\(\pi\) induces an isomorphism from
\[
  \mathcal I_\infty/\overline{I^2}
\]
onto the closure
\[
  \omega_\infty
  =
  \overline{\bigcup_r\omega_r}
\]
of the union of the \(\omega_r\), and an isomorphism from
\(S_\infty\) onto a complement \(\omega'_\infty\) of
\(\omega_\infty\) in \(\omega\). Denote the latter isomorphism by
\[
  \eta:S_\infty\xrightarrow{\sim}\omega'_\infty.
\]
We thus obtain continuous linear maps:
\begin{center}
  \phantomsection
  \label{sga3:VIIB:5.2.1.A:splitting-maps}
  % SGA3-VII-DIAG-131
  % BEGIN INCLUDED SGA3_English_current_buildable_source_20260823/source/inputs/VII/native_diagrams/exp7/viib_15_95/diag_131.tex
% SGA3-VII-DIAG-131
% Authority: Exp7B-13oct24.pdf, local p. 90.
% Crop: figures/Exp7B_localp090_521A_splitting_maps.png
% Locator: tex/components/84_expose_VIIB_section521A_pseudobasis_splittings_en.tex, Loop 1 line 74.
\begin{tikzcd}[row sep=3.2em,column sep=3em]
\overline{I^2}\times S_\infty\times\widehat{\bigoplus}_r S_r
  \arrow[r,"\phi"]
  \arrow[d,"\eta" left]
  \arrow[d,phantom,"{\times\theta}" right]
  & I \\
\omega=\omega'_\infty\times\omega_\infty & {}
\end{tikzcd}

% END INCLUDED SGA3_English_current_buildable_source_20260823/source/inputs/VII/native_diagrams/exp7/viib_15_95/diag_131.tex

\end{center}

Here
\[
  \widehat{\bigoplus}_r S_r
\]
is the direct sum of the \(S_r\) in \(\operatorname{PC}(k)\),
that is,
\[
  \left(\prod_r S_r^\dagger\right)^*
\]
\textup{(cf.\ editorial note \textup{(16)} to \SGARefLink{sga3:VIIB:0.2.2}{0.2.2})}, and
\[
  \theta:
  \widehat{\bigoplus}_r S_r
  \longrightarrow\omega_\infty
\]
is induced by the maps
\[
  S_r\xrightarrow{\sim}\omega'_r
  \lhook\joinrel\longrightarrow\omega_\infty.
\]
It follows that a sufficient condition \textup{(but not a
necessary one, as seen below)} for obtaining a section
\[
  \sigma:\omega\longrightarrow I
\]
of the desired kind is that \(\theta\) be an isomorphism.
By duality \textup{(cf.\ \SGARefLink{sga3:VIIB:0.2.2}{0.2.2})}, this is equivalent to saying
that the linear map
\[
  \omega_\infty^\dagger
  \longrightarrow
  \prod_r S_r^\dagger
\]
is bijective.

% END INCLUDED SGA3_English_current_buildable_source_20260823/source/inputs/VII/tex/components/84_expose_VIIB_section521A_pseudobasis_splittings_en.tex

% BEGIN INCLUDED SGA3_English_current_buildable_source_20260823/source/inputs/VII/tex/components/85_expose_VIIB_section521B_split_filtered_spaces_en.tex
% SGA 3, Expose VII_B, English reconstruction -- Loop 1.
% Scope: Section 5.2.1.B in full on authority local page 91;
% stops immediately before Section 5.2.1.C.
% Combined-reader page: 604. Re-edition printed page: 593.
% Editorial note in scope: 156.
% Authority: Polo--Gille Exp7B-13oct24.pdf,
% SHA-256 D303690E0195654BBE7C3DAA5D89EA23E24CD709A3575217234AC19AE8836694.
% Jacob Reinhold's CC BY 4.0 Markdown was consulted as comparison lineage.
% No genuine diagram occurs in this component; all mathematics is native LaTeX.

\paragraph{\SGARefLink{sga3:VIIB:5.2.1.B}{5.2.1.B}.}
\phantomsection
\label{sga3:VIIB:5.2.1.B}

As above, put
\[
  \omega_\infty
  =
  \overline{\bigcup_{n\in\mathbb N}\omega_n}.
\]
A second case in which a section
\[
  \sigma:\omega\longrightarrow I
\]
of the desired kind exists is the case in which
\(\omega_\infty\) has a pseudobasis \(\mathscr B_\infty\) that
is the union, for \(n\in\mathbb N^*\), of pseudobases of the
spaces
\[
  \omega_n/\omega_{n-1}.
\]
One can then complete it by a pseudobasis
\(\mathscr B'_\infty\) of \(\omega/\omega_\infty\), obtaining a
pseudobasis of \(\omega\) compatible with the filtration.

Put
\[
  V=\omega_\infty^\dagger,
\]
and let \(V_n\) be the orthogonal complement of \(\omega_n\) in
\(V\). The preceding condition is equivalent to saying that, in
the category of ordinary \(k\)-vector spaces, the separated
decreasing filtration
\[
  V=V_0\supset V_1\supset V_2\supset\cdots
\]
splits: namely, \(V\) is the direct sum, over
\(n\in\mathbb N\), of subspaces \(F_n\) such that
\[
  F_n\simeq V_n/V_{n+1}.
\]

This need not hold. For example, let
\[
  V=\mathscr S=k^{\mathbb N}
\]
be the space of sequences of elements of \(k\), and let
\(\mathscr S_n\) be the subspace of sequences
\((u_i)\) such that \(u_i=0\) for \(i<n\). Then
\[
  \dim\mathscr S_n/\mathscr S_{n+1}=1,
\]
but \(\mathscr S\) is not isomorphic to the direct sum of the
\(\mathscr S_n/\mathscr S_{n+1}\), because \(\mathscr S\) is not
countable-dimensional. In this example, however, \(\mathscr S\)
is the product of the \(\mathscr S_n/\mathscr S_{n+1}\)
\textup{(cf.\ \SGARefLink{sga3:VIIB:5.2.1.A}{5.2.1.A})}.

The splitting does hold when \(V\) is
countable-dimensional.\footnote{%
\label{sga3:VIIB:nde:156}This paragraph arose from discussions
with J.-M. Fontaine and E. Bouscaren; in particular, Bouscaren
indicated the proof that follows.}
Indeed, let
\[
  (e_n)_{n\in\mathbb N}
\]
be a basis of \(V\). We shall construct inductively an increasing
function
\[
  g:\mathbb N\longrightarrow\mathbb N
\]
and subspaces \(F_i\), for \(i=0,\ldots,g(n)\), such that
\[
  F_i\simeq V_i/V_{i+1}
\]
and
\[
  F_{\leq g(n)}
  =
  \bigoplus_{i=0}^{g(n)}F_i
\]
is a complement of \(V_{g(n)+1}\) containing
\[
  e_0,\ldots,e_n.
\]
It will then follow that
\[
  V=\bigoplus_{i\geq0}F_i.
\]

Let \(n+1\in\mathbb N\), and assume the assertion already proved
for \(n\) \textup{(it is vacuous for \(n=-1\))}. If
\[
  e_{n+1}\in F_{\leq g(n)},
\]
put \(g(n+1)=g(n)\). Otherwise write
\[
  e_{n+1}=f+x,
  \qquad
  f\in F_{\leq g(n)},
  \qquad
  0\ne x\in V_{g(n)+1}.
\]
Let \(j\) be the least integer such that
\[
  x\in V_j\setminus V_{j+1}.
\]
For
\[
  i=g(n)+1,\ldots,j,
\]
choose a complement \(F_i\) of \(V_{i+1}\) in \(V_i\), choosing
\(F_j\) so that \(x\in F_j\). Then
\[
  e_{n+1}\in F_{\leq j};
\]
put \(g(n+1)=j\).

\phantomsection
\label{sga3:VIIB:5.2.1.B:induction-source-note}
\noindent\textit{Source note.}\enspace
At the last induction step, the French PDF prints
\(e_{n+1}\in F_j\). From
\(e_{n+1}=f+x\), with \(f\in F_{\leq g(n)}\), the argument
instead gives \(x\in F_j\) and
\(e_{n+1}\in F_{\leq j}\); these are the relations used here.

% END INCLUDED SGA3_English_current_buildable_source_20260823/source/inputs/VII/tex/components/85_expose_VIIB_section521B_split_filtered_spaces_en.tex

% BEGIN INCLUDED SGA3_English_current_buildable_source_20260823/source/inputs/VII/tex/components/86_expose_VIIB_section521C_stationary_filtration_en.tex
% SGA 3, Expose VII_B, English reconstruction -- Loop 1.
% Scope: Section 5.2.1.C in full on authority local pages 91--92;
% stops immediately before Section 5.2.2.
% Combined-reader pages: 604--605. Re-edition printed pages: 593--594.
% Editorial note in scope: 157.
% Authority: Polo--Gille Exp7B-13oct24.pdf,
% SHA-256 D303690E0195654BBE7C3DAA5D89EA23E24CD709A3575217234AC19AE8836694.
% Jacob Reinhold's CC BY 4.0 Markdown was consulted as comparison lineage.
% No genuine diagram occurs in this component; all mathematics is native LaTeX.

\paragraph{\SGARefLink{sga3:VIIB:5.2.1.C}{5.2.1.C}.}
\phantomsection
\label{sga3:VIIB:5.2.1.C}
\footnote{\label{sga3:VIIB:nde:157}Here the editors return to the
original, shortened to take account of the preceding additions.}
In particular, the two preceding conditions \textup{(\SGARefLink{sga3:VIIB:5.2.1.A}{5.2.1.A} and
\SGARefLink{sga3:VIIB:5.2.1.B}{5.2.1.B})} hold when the filtration of \(\omega\) is stationary,
that is, when there exists an integer \(n_0\) such that
\[
  \omega_n=\omega_{n_0}
  \qquad
  \text{for }n_0\leq n<+\infty.
\]
In this case,
\[
  k[[\omega]]
  \big/
  \bigl((x^{p^r})\bigr)_{x\in\omega_r}
\]
is isomorphic to the completed tensor product
\[
\begin{aligned}
  &
  \frac{k[[\omega_1]]}
       {\bigl((x^p)\bigr)_{x\in\omega_1}}
  \widehat\otimes_k
  \frac{k[[\omega'_2]]}
       {\bigl((x^{p^2})\bigr)_{x\in\omega'_2}}
  \widehat\otimes_k\cdots
  \\
  &\hspace{3.2em}
  {}\widehat\otimes_k
  \frac{k[[\omega'_{n_0}]]}
       {\bigl((x^{p^{n_0}})\bigr)_{x\in\omega'_{n_0}}}
  \widehat\otimes_k
  k[[\omega'_\infty]],
\end{aligned}
\]
where \(\omega'_n\) is a complement of \(\omega_{n-1}\) in
\(\omega_n\), and \(\omega'_\infty\) is a complement of
\[
  \omega_\infty=\omega_{n_0}
\]
in \(\omega\).

The filtration of \(\omega\) is plainly stationary in case
\textup{\SGARefLink{sga3:vii:paragraph:5-1:equation:i}{(i)}}, where \(\omega_r=\omega\) for all sufficiently
large \(r\); in case \textup{\SGARefLink{sga3:vii:paragraph:5-1:equation:ii}{(ii)}}, where \(\omega\) is
finite-dimensional; and also in case \textup{(iii)}, where
\[
  I_r=0
\]
for every \(r\). In the last case, \(B\) is isomorphic to the
formal power series algebra \(k[[\omega]]\). The preceding remarks
therefore imply the theorem, subject to \SGARefLink{sga3:VIIB:5.2.2}{5.2.2}--\SGARefLink{sga3:VIIB:5.2.5}{5.2.5} below.

% END INCLUDED SGA3_English_current_buildable_source_20260823/source/inputs/VII/tex/components/86_expose_VIIB_section521C_stationary_filtration_en.tex

% BEGIN INCLUDED SGA3_English_current_buildable_source_20260823/source/inputs/VII/tex/components/87_expose_VIIB_section522_height_one_case_en.tex
% SGA 3, Expose VII_B, English reconstruction -- Loop 1.
% Scope: Section 5.2.2 in full on authority local page 92;
% stops immediately before Section 5.2.3.
% Combined-reader page: 605. Re-edition printed page: 594.
% Original-edition marginal locator in scope: 559.
% Authority: Polo--Gille Exp7B-13oct24.pdf,
% SHA-256 D303690E0195654BBE7C3DAA5D89EA23E24CD709A3575217234AC19AE8836694.
% Jacob Reinhold's CC BY 4.0 Markdown was consulted as comparison lineage.
% No genuine diagram occurs in this component; all mathematics is native LaTeX.

\subsubsection*{5.2.2}
\phantomsection
\label{sga3:VIIB:5.2.2}

Suppose first that \(B\) has height at most \(1\); that is,
\[
  x^p=0
  \qquad\text{for every }x\in I.
\]
By \SGARefLink{sga3:VIIB:5.1.4}{5.1.4}, the graded algebra
\[
  \operatorname{Gr}_I(B)
\]
associated with the filtration
\[
  I\supset I^2\supset I^3\supset\cdots
\]
of \(B\) is endowed with a cogroup structure in the category
\(\mathbf{Alpg}/k\). In other words, the profinite algebra
\[
  C=\prod_{n\geq0}\operatorname{Gr}_I(B)_n
\]
is the affine algebra of a formal \(k\)-group \(N\).
Clearly,
\[
  \omega_{N/k}=I/I^2,
\]
and \(N\) is infinitesimal of height at most \(1\).

By \SGARefLink{sga3:viib:subsection:4-4}{4.4}, the identity map of \(\omega_{N/k}\) therefore induces
an isomorphism
\[
  B'
  =
  k[[\omega_{N/k}]]
  \big/
  \bigl((x^p)\bigr)_{x\in\omega_{N/k}}
  \xrightarrow{\ \sim\ }
  C.
\]
In particular, the map \(\phi\) of \SGARefLink{sga3:VIIB:5.2.1:section-proposition}{5.2.1} induces an isomorphism
between the associated graded algebras of \(B'\) and \(B\), when
both are filtered by the powers of their augmentation ideals.
It follows from \textup{[CA]}, \(\S\)V.7, Lemma~1
\textup{(see also \SGARefLink{sga3:VIIB:5.1.5}{5.1.5})} that \(\phi\) is an isomorphism.

% END INCLUDED SGA3_English_current_buildable_source_20260823/source/inputs/VII/tex/components/87_expose_VIIB_section522_height_one_case_en.tex

% BEGIN INCLUDED SGA3_English_current_buildable_source_20260823/source/inputs/VII/tex/components/88_expose_VIIB_section523_finite_height_induction_en.tex
% SGA 3, Expose VII_B, English reconstruction -- Loop 1.
% Scope: Section 5.2.3 in full on authority local pages 92--93;
% stops immediately before Section 5.2.4.
% Combined-reader pages: 605--606. Re-edition printed pages: 594--595.
% Original-edition marginal locator in scope: 560.
% Editorial notes in scope: 158--159.
% Authority: Polo--Gille Exp7B-13oct24.pdf,
% SHA-256 D303690E0195654BBE7C3DAA5D89EA23E24CD709A3575217234AC19AE8836694.
% Jacob Reinhold's CC BY 4.0 Markdown was consulted as comparison lineage.
% Three genuine diagrams are preserved as two exact PDF-derived PNGs
% for Loop 1 and remain due for native reconstruction in Loop 2.

\subsubsection*{5.2.3}
\phantomsection
\label{sga3:VIIB:5.2.3}

Suppose now that \(B\) has finite height at most \(r\).
Let
\[
  \pi:k\xrightarrow{\ \sim\ }k,
  \qquad
  x\longmapsto x^p,
\]
be the Frobenius automorphism. The linear map
\[
  B\widehat\otimes_{\pi}k\longrightarrow B,
  \qquad
  b\widehat\otimes_{\pi}x
  \longmapsto
  b^px=(bx^{1/p})^p,
\]
has closed image equal to the closed subalgebra
\[
  B^p=\{\,b^p\mid b\in B\,\}
\]
of \(B\). Put
\[
  J=B^p\cap I=B^p\cap I_A.
\]

\footnote{\label{sga3:VIIB:nde:158}The editors added what follows.}
Let \(G_1\) be the kernel of the Frobenius morphism
\[
  \operatorname{Fr}:G\longrightarrow G^{(p)},
\]
and let \(HG_1\) be the formal subgroup of \(G\) that is the
inverse image of the formal subgroup \(H^{(p)}\subset G^{(p)}\).
Then \(HG_1\) is defined by the closed ideal generated by the
\(p\)-th powers of the elements of \(AI\); this ideal is \(AJ\).
Moreover, the formation of \(G/H\) commutes with base change,
because \(G/H\) represents the quotient sheaf for the flat
topology \textup{(cf.\ \SGARefLink{sga3:VIIB:2.4}{2.4})}. Hence
\[
  (G/H)^{(p)}=G^{(p)}/H^{(p)},
\]
and we have the following commutative diagrams:

\begin{figure}[H]
  \centering
  % SGA3-VII-DIAG-132
  % BEGIN INCLUDED SGA3_English_current_buildable_source_20260823/source/inputs/VII/native_diagrams/exp7/viib_15_95/diag_132.tex
% SGA3-VII-DIAG-132
% Authority: Exp7B-13oct24.pdf, local p. 92.
% Crop: figures/Exp7B_localp092_frobenius_quotient_dual_diagrams.png
% Locator: tex/components/88_expose_VIIB_section523_finite_height_induction_en.tex, Loop 1 line 60.
\begin{tikzcd}[row sep=2.5em,column sep=2.8em]
G \arrow[r,"\operatorname{Fr}"] \arrow[d]
  & G^{(p)} \arrow[d] \\
G/H \arrow[r,"\operatorname{Fr}"']
  & G^{(p)}/H^{(p)}
\end{tikzcd}
\qquad
\begin{tikzcd}[row sep=2.5em,column sep=2.8em]
A
  & A\widehat\otimes_\pi k
      \arrow[l,"{a\widehat\otimes1\mapsto a^p}"'] \\
B \arrow[u]
  & B\widehat\otimes_\pi k
      \arrow[l,"{b\widehat\otimes1\mapsto b^p}"']
      \arrow[u]
\end{tikzcd}

% END INCLUDED SGA3_English_current_buildable_source_20260823/source/inputs/VII/native_diagrams/exp7/viib_15_95/diag_132.tex

\end{figure}

It follows that \(B^p\) is the affine algebra of the quotient
\(G/HG_1\).\footnote{\label{sga3:VIIB:nde:159}As indicated in the
original, this also follows from \SGARefLink{sga3:VIIB:5.1:proposition}{Proposition~5.1}. The editors
preferred to give the preceding argument, which does not use the
implication \textup{\SGARefLink{sga3:vii:paragraph:5-1:equation:i}{(i)}} \(\Rightarrow\) \textup{\SGARefLink{sga3:VIIB:5.1:proposition:item:ii}{(ii)}} of that
proposition.}
Temporarily denote by \(C\) the affine algebra of the quotient
\(HG_1/H\). Since the formation of \(G/H\) commutes with base
change, the square

\begin{figure}[H]
  \centering
  % SGA3-VII-DIAG-133
  % BEGIN INCLUDED SGA3_English_current_buildable_source_20260823/source/inputs/VII/native_diagrams/exp7/viib_15_95/diag_133.tex
% SGA3-VII-DIAG-133
% Authority: Exp7B-13oct24.pdf, local p. 92.
% Crop: figures/Exp7B_localp092_cartesian_HG1_square.png
% Locator: tex/components/88_expose_VIIB_section523_finite_height_induction_en.tex, Loop 1 line 76.
\begin{tikzcd}[row sep=2.8em,column sep=3.5em]
HG_1 \arrow[r] \arrow[d]
  & G \arrow[d] \\
HG_1/H \arrow[r]
  & G/H
\end{tikzcd}

% END INCLUDED SGA3_English_current_buildable_source_20260823/source/inputs/VII/native_diagrams/exp7/viib_15_95/diag_133.tex

\end{figure}

\noindent
is Cartesian. Consequently,
\[
  A\widehat\otimes_BC
  \simeq
  A/AJ
  =
  A\widehat\otimes_B(B/BJ).
\]
Since \(A\) is topologically free over \(B\)
\textup{(by 2.4, because \(A\) and \(B\) are local)}, the natural
morphism
\[
  B/BJ\longrightarrow C
\]
is an isomorphism. Thus \(B/BJ\) is the affine algebra of
\(HG_1/H\), and of course
\[
  B/BJ=B_1
\]
has height at most \(1\), since
\[
  J=\bigl((x^p)\bigr)_{x\in\mathfrak r(B)}.
\]

Let
\[
  B'
  =
  k[[\omega]]
  \big/
  \bigl((x^{p^r})\bigr)_{x\in\omega_r},
\]
let \(\phi:B'\to B\) be the morphism introduced in \SGARefLink{sga3:VIIB:5.2.1:section-proposition}{5.2.1}, let
\[
  B'^p=\{\,x^p\mid x\in B'\,\},
\]
and let \(J'\) be the augmentation ideal of \(B'^p\). We then
have the commutative diagram

\begin{figure}[H]
  \centering
  % SGA3-VII-DIAG-134
  % BEGIN INCLUDED SGA3_English_current_buildable_source_20260823/source/inputs/VII/native_diagrams/exp7/viib_15_95/diag_134.tex
% SGA3-VII-DIAG-134
% Authority: Exp7B-13oct24.pdf, local p. 93.
% Crop: figures/Exp7B_localp093_phi_reduction_square.png
% Locator: tex/components/88_expose_VIIB_section523_finite_height_induction_en.tex, Loop 1 line 122.
\begin{tikzcd}[row sep=2.8em,column sep=3.7em]
B' \arrow[r,"\psi"] \arrow[d]
  & B \arrow[d] \\
B'_1=B'/\overline{B'J'}
  \arrow[r,"\phi_1"]
  & B_1=B/\overline{BJ}
\end{tikzcd}

% END INCLUDED SGA3_English_current_buildable_source_20260823/source/inputs/VII/native_diagrams/exp7/viib_15_95/diag_134.tex

\end{figure}

\noindent
and, by \SGARefLink{sga3:VIIB:5.2.2}{5.2.2}, \(\phi_1\) is an isomorphism.

On the other hand, by 2.4, \(A\) is topologically flat over
\[
  B=\mathscr A(G/H)
  \qquad\text{and over}\qquad
  B^p=\mathscr A(G/HG_1).
\]
It follows from \SGARefLink{sga3:VIIB:1.3.3}{1.3.3} that \(B\) is topologically flat over
\(B^p\). Moreover, \SGARefLink{sga3:VIIB:5.2.4}{5.2.4} below shows that the morphism
\[
  B'^p\longrightarrow B^p
\]
induced by \(\phi\) is an isomorphism. We may therefore apply
\SGARefLink{sga3:VIIB:0.3.4}{0.3.4} to the pseudocompact ring \(B'^p=B^p\) and the
pseudocompact \(B^p\)-modules \(M=B'\), \(N=B\). By what
precedes,
\[
  \phi_1
  =
  \phi\widehat\otimes_{B^p}k
\]
is an isomorphism, and it follows that \(\phi\) itself is an
isomorphism. This proves \SGARefLink{sga3:VIIB:5.2:theorem}{5.2} when \(B\) has finite height, subject
only to \SGARefLink{sga3:VIIB:5.2.4}{5.2.4} below.

% END INCLUDED SGA3_English_current_buildable_source_20260823/source/inputs/VII/tex/components/88_expose_VIIB_section523_finite_height_induction_en.tex

% BEGIN INCLUDED SGA3_English_current_buildable_source_20260823/source/inputs/VII/tex/components/89_expose_VIIB_section524_frobenius_twist_induction_en.tex
% SGA 3, Expose VII_B, English reconstruction -- Loop 1.
% Scope: Section 5.2.4 in full on authority local pages 93--94;
% stops immediately before Remark 5.2.4.A.
% Combined-reader pages: 606--607. Re-edition printed pages: 595--596.
% Original-edition marginal locator in scope: 561.
% Editorial note in scope: 160.
% Authority: Polo--Gille Exp7B-13oct24.pdf,
% SHA-256 D303690E0195654BBE7C3DAA5D89EA23E24CD709A3575217234AC19AE8836694.
% Jacob Reinhold's CC BY 4.0 Markdown was consulted as comparison lineage.
% The genuine exact-rows diagram is preserved as an exact PDF-derived
% PNG for Loop 1 and remains due for native reconstruction in Loop 2.

\subsubsection*{5.2.4}
\phantomsection
\label{sga3:VIIB:5.2.4}

For every pseudocompact vector space \(V\) over \(k\), denote by
\[
  {}_{\pi}V
  =
  V\widehat\otimes_{\pi}k
\]
the space obtained from \(V\) by extension of scalars along
\(x\mapsto x^p\).\footnote{\label{sga3:VIIB:nde:160}Since \(k\)
is perfect, one may identify \({}_{\pi}V\) with the abelian group
\(V\) on which \(k\) acts by
\(\lambda\cdot v=\lambda^{1/p}v\).}
We then have a commutative diagram with exact rows:

\begin{figure}[H]
  \centering
  % SGA3-VII-DIAG-135
  % BEGIN INCLUDED SGA3_English_current_buildable_source_20260823/source/inputs/VII/native_diagrams/exp7/viib_15_95/diag_135.tex
% SGA3-VII-DIAG-135
% Authority: Exp7B-13oct24.pdf, local p. 93.
% Crop: figures/Exp7B_localp093_pi_exact_rows_diagram.png
% Locator: tex/components/89_expose_VIIB_section524_frobenius_twist_induction_en.tex, Loop 1 line 32.
\begin{tikzcd}[row sep=2.8em,column sep=2.6em]
0 \arrow[r]
  & {}_\pi\overline{I^2} \arrow[r,"\alpha"] \arrow[d,"u"']
  & {}_\pi I \arrow[r,"\beta"] \arrow[d,"v"']
  & {}_\pi\omega \arrow[r] \arrow[d,"w"']
  & 0 \\
0 \arrow[r]
  & \overline{J^2} \arrow[r,"\gamma"]
  & J \arrow[r,"\delta"]
  & \overline\omega \arrow[r]
  & 0\, ,
\end{tikzcd}

% END INCLUDED SGA3_English_current_buildable_source_20260823/source/inputs/VII/native_diagrams/exp7/viib_15_95/diag_135.tex

\end{figure}

\noindent
Here
\[
  \overline\omega=J/J^2,
\]
and the maps \(u,v,w\) are induced by the linear map
\[
  {}_{\pi}B\longrightarrow B^p,
  \qquad
  x\widehat\otimes a\longmapsto x^pa.
\]
Since \(u\) and \(v\) are surjective, so is \(w\), and its kernel
is the image \({}_{\pi}\omega_1\) of
\[
  {}_{\pi}I_1=\ker(v).
\]

Put
\[
  J_n=\{\,x\in J\mid x^{p^n}=0\,\},
  \qquad
  \overline\omega_n=\delta(J_n).
\]
Then, for every \(n\geq0\),
\[
  J_n=v({}_{\pi}I_{n+1}),
  \qquad
  \overline\omega_n=w({}_{\pi}\omega_{n+1}).
\]

\phantomsection
\label{sga3:VIIB:5.2.4:delta-source-note}
\noindent\textit{Source note.}\enspace
The French PDF prints
\(\overline\omega_n=\delta(I_n)\), although the displayed exact
row defines \(\delta\) on \(J\), and the immediately preceding
definition introduces \(J_n\subset J\). The type-correct
\(\delta(J_n)\) is used above; all subsequent identities agree
with this reading.

The section
\[
  {}_{\pi}\sigma:{}_{\pi}\omega\longrightarrow{}_{\pi}I
\]
induced by the section \(\sigma\) of \SGARefLink{sga3:VIIB:5.2.1:section-proposition}{5.2.1} therefore descends to
a section
\[
  \tau:\overline\omega\longrightarrow J
\]
compatible with the filtrations of \(J\) and
\(\overline\omega\). Since \(B^p\) has height at most \(r-1\),
the induction hypothesis shows that this section induces an
isomorphism
\[
  \psi:
  B''
  =
  k[[\overline\omega]]
  \big/
  \bigl((x^{p^n})\bigr)_{x\in\overline\omega_n}
  \xrightarrow{\ \sim\ }
  B^p.
\]
But \(B''\) identifies with \(B'^p\), and \(\psi\) identifies
with the morphism \(B'^p\to B^p\) induced by \(\phi\). Thus the
theorem is proved when \(B\) has finite height.

% END INCLUDED SGA3_English_current_buildable_source_20260823/source/inputs/VII/tex/components/89_expose_VIIB_section524_frobenius_twist_induction_en.tex

% BEGIN INCLUDED SGA3_English_current_buildable_source_20260823/source/inputs/VII/tex/components/90_expose_VIIB_remark524A_finite_height_filtration_en.tex
% SGA 3, Expose VII_B, English reconstruction -- Loop 1.
% Scope: Remark 5.2.4.A in full on authority local page 94;
% stops immediately before Section 5.2.5.
% Combined-reader page: 607. Re-edition printed page: 596.
% Editorial note in scope: 161.
% Authority: Polo--Gille Exp7B-13oct24.pdf,
% SHA-256 D303690E0195654BBE7C3DAA5D89EA23E24CD709A3575217234AC19AE8836694.
% Jacob Reinhold's CC BY 4.0 Markdown was consulted as comparison lineage.
% No genuine diagram occurs in this component; all mathematics is native LaTeX.

\medskip
\noindent\textbf{Remark 5.2.4.A.}
\phantomsection
\label{sga3:VIIB:5.2.4.A}
\footnote{\label{sga3:VIIB:nde:161}The editors added this remark,
which is used in \SGARefLink{sga3:VIIB:5.2.5}{5.2.5}.}
Suppose that \(B\) has height at most \(r+1\), with
\(r\in\mathbb N^*\). Then \(I_{r+1}=I\), and there is an
isomorphism
\[
\tag{1}\label{sga3:vii:remark:5-2-4-a:equation:1}
  B
  \simeq
  \frac{k[[S_1]]}{\bigl((x_1^p)\bigr)_{x_1\in S_1}}
  \widehat\otimes_k\cdots\widehat\otimes_k
  \frac{k[[S_r]]}{\bigl((x_r^{p^r})\bigr)_{x_r\in S_r}}
  \widehat\otimes_k
  \frac{k[[S_{r+1}]]}
       {\bigl((x_{r+1}^{p^{r+1}})\bigr)_{x_{r+1}\in S_{r+1}}},
\]
where each \(S_n\) is a complement of
\[
  I^2+I_{n-1}
\]
in \(I^2+I_n\). Thus \(\omega=I/I^2\) identifies with
\[
  \prod_{i=1}^{r+1}S_i,
\]
and it is immediate that, for \(n=1,\ldots,r+1\), the image
\(\omega_n\) of \(I_n\) in \(\omega\) identifies with
\[
  \prod_{i=1}^{n}S_i.
\]

This has the following consequence. Let
\[
  B_r=B/\check J_r,
  \qquad
  \check J_r
  =
  \bigl((x^{p^r})\bigr)_{x\in\mathfrak r(B)},
\]
and let
\[
  \mathfrak m=I/\check J_r
\]
be the augmentation ideal of \(B_r\). Since
\(\check J_r\subset I^2\), the space
\[
  \omega(r)=\mathfrak m/\mathfrak m^2
\]
identifies with \(\omega\). For \(n=1,\ldots,r\), denote by
\(\omega(r)_n\) the image in \(\omega(r)\) of
\[
  \mathfrak m_n
  =
  \{\,y\in\mathfrak m\mid y^{p^n}=0\,\}.
\]
It is also the image in \(\omega\) of
\[
  \{\,x\in I\mid x^{p^n}\in\check J_r\,\};
\]
hence \(\omega(r)_n\) contains \(\omega_n\).

On the other hand, isomorphism \textup{\SGARefLink{sga3:vii:remark:5-2-4-a:equation:1}{(1)}} gives
\[
  \check J_r
  =
  \bigl((x_{r+1}^{p^r})\bigr)_{x_{r+1}\in S_{r+1}},
\]
and therefore
\[
\tag{2}\label{sga3:vii:remark:5-2-4-a:equation:2}
  B_r
  \simeq
  \frac{k[[S_1]]}{\bigl((x_1^p)\bigr)_{x_1\in S_1}}
  \widehat\otimes_k\cdots\widehat\otimes_k
  \frac{k[[S_r]]}{\bigl((x_r^{p^r})\bigr)_{x_r\in S_r}}
  \widehat\otimes_k
  \frac{k[[S_{r+1}]]}
       {\bigl((x_{r+1}^{p^r})\bigr)_{x_{r+1}\in S_{r+1}}}.
\]
It follows that \(\omega(r)_n\) identifies with
\(\prod_{i=1}^{n}S_i\) for \(n=1,\ldots,r-1\). Thus the
inclusion
\[
  \omega_n\subset\omega(r)_n
\]
is an equality for \(n=1,\ldots,r-1\).

% END INCLUDED SGA3_English_current_buildable_source_20260823/source/inputs/VII/tex/components/90_expose_VIIB_remark524A_finite_height_filtration_en.tex

% BEGIN INCLUDED SGA3_English_current_buildable_source_20260823/source/inputs/VII/tex/components/91_expose_VIIB_section525_infinite_height_limit_en.tex
% SGA 3, Expose VII_B, English reconstruction -- Loop 1.
% Scope: Section 5.2.5 in full on authority local pages 94--95;
% stops immediately before Section 5.2.6.
% Combined-reader pages: 607--608. Re-edition printed pages: 596--597.
% Original-edition marginal locator in scope: 562.
% Editorial notes in scope: 162--163.
% Authority: Polo--Gille Exp7B-13oct24.pdf,
% SHA-256 D303690E0195654BBE7C3DAA5D89EA23E24CD709A3575217234AC19AE8836694.
% Jacob Reinhold's CC BY 4.0 Markdown was consulted as comparison lineage.
% No genuine diagram occurs in this component; all mathematics is native LaTeX.

\subsubsection*{5.2.5}
\phantomsection
\label{sga3:VIIB:5.2.5}

It remains to consider the case in which \(B\) has infinite
height and the projection \(I\to\omega\) admits a section
\(\sigma\) compatible with the filtrations of \(\omega\) and
\(I\). Consider the morphism induced by \(\sigma\):
\[
  \phi:
  B'
  =
  k[[\omega]]
  \big/
  \bigl((x^{p^n})\bigr)_{x\in\omega_n}
  \longrightarrow B.
\]
It is enough to prove that, for every \(r\in\overline{\mathbb N}\),
the map
\[
  \phi_r:B'_r\longrightarrow B_r
\]
induced by \(\phi\) is an isomorphism.\footnote{%
\label{sga3:VIIB:nde:162}Indeed, \(B'\), respectively \(B\), is
the inverse limit of the \(B'_r\), respectively the \(B_r\).
Moreover, the editors modified the original in what follows in
order to take account of the addition made in \SGARefLink{sga3:VIIB:5.2.3}{5.2.3}.}

For every \(r\in\mathbb N^*\), let \(G_r\) be the kernel of the
iterated Frobenius morphism
\[
  G\longrightarrow G^{(p^r)},
\]
and let \(HG_r\) be the formal subgroup of \(G\) that is the
inverse image of \(H^{(p^r)}\subset G^{(p^r)}\). Thus \(HG_r\)
is defined by the closed ideal generated by the \(p^r\)-th
powers of the elements of \(AI\), namely
\[
  A\check J_r,
  \qquad
  \check J_r=\{\,x^{p^r}\mid x\in I\,\}.
\]
Exactly as in \SGARefLink{sga3:VIIB:5.2.3}{5.2.3}, it follows that
\[
  B_r=B/B\check J_r
\]
is the affine algebra of \(HG_r/H\); in particular, it has
height at most \(r\).

Let
\[
  \mathfrak m(r)=I/B\check J_r
\]
be the augmentation ideal of \(B_r\). The canonical projection
\(B\to B_r\) induces an isomorphism
\[
  \omega=I/I^2
  \xrightarrow{\ \sim\ }
  \omega(r)
  =
  \mathfrak m(r)/\mathfrak m(r)^2,
\]
which we use to identify the two spaces. Let \(\omega(r)_n\) be
the image in \(\omega(r)\) of the closed ideal
\[
  \mathfrak m(r)_n
  =
  \{\,y\in\mathfrak m(r)\mid y^{p^n}=0\,\}.
\]
It is also the image in \(\omega\) of the closed ideal
\[
  I(r)_n
  =
  \{\,x\in I\mid x^{p^n}\in B\check J_r\,\}.
\]
\footnote{\label{sga3:VIIB:nde:163}The editors expanded the
original in what follows.}

Clearly,
\[
  \omega(r)_n=\omega
  \qquad\text{if }n\geq r.
\]
We shall prove that
\[
  \omega(r)_n=\omega_n
  \qquad\text{if }n<r.
\]
For every \(r,n\), the sequence
\[
  0
  \longrightarrow
  I(r)_n\cap I^2
  \longrightarrow
  I(r)_n
  \longrightarrow
  \omega(r)_n
  \longrightarrow0
\]
is exact. For fixed \(n\),
\[
  \bigcap_r I(r)_n=I_n,
\]
because \(\bigcap_rB\check J_r=0\). Filtered inverse limits are
exact in \(\operatorname{PC}(k)\) \textup{(cf.\ \SGARefLink{sga3:VIIB:0.2}{0.2})}, so for
every \(n\),
\[
\tag{*}\label{sga3:vii:remark:5-2-4-a:equation:star}
  \omega_n=\bigcap_r\omega(r)_n.
\]
On the other hand, \SGARefLink{sga3:VIIB:5.2.4.A}{Remark~5.2.4.A} gives
\[
  \omega(r)_n=\omega(r+1)_n
  \qquad\text{if }n<r.
\]
Together with \textup{(*)}, this implies
\[
  \omega(r)_n=\omega_n
  \qquad\text{if }n<r.
\]

Consequently, the vector space \(\omega\), filtered by the
subspaces \(\bigl(\omega(r)_n\bigr)_{n\geq0}\), is precisely
\({}^{r}\omega\) in the notation of \SGARefLink{sga3:VIIB:5.2:theorem}{5.2}. The composite
\[
  \sigma(r):
  \omega\xrightarrow{\ \sigma\ }I
  \longrightarrow\mathfrak m(r)
\]
is therefore compatible with the filtrations
\(\bigl(\omega(r)_n\bigr)\) and
\(\bigl(\mathfrak m(r)_n\bigr)\). Since
\[
  k[[{}^{r}\omega]]
  \big/
  \bigl((x^{p^n})\bigr)_{x\in{}^{r}\omega_n}
\]
is exactly \(B'_r\), and \(\phi_r:B'_r\to B_r\) is induced by
\(\sigma(r)\), the result already proved for finite-height
algebras shows that \(\phi_r\) is an isomorphism.

% END INCLUDED SGA3_English_current_buildable_source_20260823/source/inputs/VII/tex/components/91_expose_VIIB_section525_infinite_height_limit_en.tex

% BEGIN INCLUDED SGA3_English_current_buildable_source_20260823/source/inputs/VII/tex/components/92_expose_VIIB_section526_infinite_completed_tensor_product_en.tex
% SGA 3, Expose VII_B, English reconstruction -- Loop 1.
% Scope: Definition 5.2.6 in full on authority local page 95;
% stops immediately before Section 5.3.
% Combined-reader page: 608. Re-edition printed page: 597.
% Editorial note in scope: 164.
% Authority: Polo--Gille Exp7B-13oct24.pdf,
% SHA-256 D303690E0195654BBE7C3DAA5D89EA23E24CD709A3575217234AC19AE8836694.
% Jacob Reinhold's CC BY 4.0 Markdown was consulted as comparison lineage.
% No genuine diagram occurs in this component; all mathematics is native LaTeX.

\subsubsection*{5.2.6}
\phantomsection
\label{sga3:VIIB:5.2.6}

\noindent\textbf{Definition 5.2.6.}
\footnote{\label{sga3:VIIB:nde:164}The editors added this number
in order to define the infinite tensor products used in \SGARefLink{sga3:VIIB:5.3:item:a}{5.3(a)}.}
Let \((A_\lambda)_{\lambda\in\Lambda}\) be a family of profinite
\(k\)-algebras, each endowed with an augmentation
\[
  \varepsilon_\lambda:A_\lambda\longrightarrow k.
\]
\textit{(This holds in particular if every \(A_\lambda\) is
local with residue field \(k\).)}
Define the infinite completed tensor product
\[
  \mathcal A
  =
  \mathop{\widehat{\bigotimes}}\nolimits^{c}_{\lambda\in\Lambda}
  A_\lambda
\]
to be the inverse limit in \(\operatorname{Alp}/k\) of
\[
  \mathcal A_F
  =
  \mathop{\widehat{\bigotimes}}\nolimits^{c}_{\lambda\in F}
  A_\lambda,
\]
where \(F\) ranges through the finite subsets of \(\Lambda\).
For \(F'=F\cup\{\lambda\}\), the transition morphism
\[
  \mathcal A_{F'}\longrightarrow\mathcal A_F
\]
is
\[
  \operatorname{id}\widehat\otimes\varepsilon_\lambda.
\]

In particular, suppose \(\Lambda=\mathbb N^*\), and write
\[
  X_n=\operatorname{Spf}(A_n).
\]
Then \(\operatorname{Spf}(\mathcal A)\) represents the functor
that assigns to every \(C\in\operatorname{Alf}/k\) the set of
\emph{finite} sequences in
\(\prod_{n\geq1}X_n(C)\); that is, the sequences
\[
  (x_1,x_2,\ldots)\in\prod_{n\geq1}X_n(C)
\]
such that
\[
  x_n=\varepsilon_n
\]
for all sufficiently large \(n\). Here, by abuse of notation,
\(\varepsilon_n\) denotes the composite of
\[
  \varepsilon_n:A_n\longrightarrow k
\]
with the structural morphism \(k\to C\).

If, in addition, every \(A_n\) is a quotient of an algebra
\(k[[\omega'_n]]\), one may write \(0\) for the unique morphism
\[
  A_n\longrightarrow C
\]
that vanishes on \(\omega'_n\). The preceding description then
becomes the set of sequences for which \(x_n=0\) for all
sufficiently large \(n\).

% END INCLUDED SGA3_English_current_buildable_source_20260823/source/inputs/VII/tex/components/92_expose_VIIB_section526_infinite_completed_tensor_product_en.tex

% BEGIN INCLUDED SGA3_English_current_buildable_source_20260823/source/inputs/VII/tex/components/93_expose_VIIB_sections53_54_stationary_algebras_corollary_en.tex
% SGA 3, Expose VII_B, English reconstruction -- Loop 1.
% Scope: Sections 5.3 and 5.4 in full on authority local pages 95--97;
% stops immediately before Section 5.5.
% Combined-reader pages: 608--610. Re-edition printed pages: 597--599.
% Original-edition marginal locator in scope: 563.
% Editorial notes in scope: 165--168.
% Authority: Polo--Gille Exp7B-13oct24.pdf,
% SHA-256 D303690E0195654BBE7C3DAA5D89EA23E24CD709A3575217234AC19AE8836694.
% Jacob Reinhold's CC BY 4.0 Markdown was consulted as comparison lineage.
% No genuine diagram occurs in this component; all mathematics is native LaTeX.

\subsection*{5.3. Remarks}
\phantomsection
\label{sga3:VIIB:5.3}
\addcontentsline{toc}{subsection}{5.3. Remarks}

\noindent\SGAInlineRefTarget{sga3:VIIB:5.3:item:a}\textup{(a)}
Call a profinite \(k\)-algebra \emph{stationary} if it is the
completed tensor product of a family of truncated formal power
series algebras.\footnote{\label{sga3:VIIB:nde:165}The author
doubtless had in mind a completed tensor product
\[
  \mathcal A
  =
  \mathop{\widehat{\bigotimes}}\nolimits^{c}_{n\in\mathbb N^*}
  \frac{k[[\omega'_n]]}
       {\bigl((x^{p^n})\bigr)_{x\in\omega'_n}},
\]
where the \(\omega'_n\) are arbitrary pseudocompact vector
spaces. Then \(\omega=\omega_{\mathcal A/k}\) identifies with
\(\prod_n\omega'_n\), and the filtration is
\(\omega_n=\prod_{i=1}^{n}\omega'_i\); this is the case of
\SGARefLink{sga3:VIIB:5.2.1.B}{5.2.1.B}. For this reason these algebras would be better called
\emph{stable}, rather than \emph{stationary}; cf.\ \textup{[Di73]},
II, \(\S\)\SGARefLink{sga3:VIIB:2.9}{2.9}, p.~75. If, for example,
\(\dim\omega'_n=1\) for every \(n\in\mathbb N^*\), then
\(\mathcal A\) represents the functor that assigns to
\(C\in\operatorname{Alf}/k\) the sequences
\((x_n)_{n\in\mathbb N^*}\) of elements of \(C\) such that
\(x_n^{p^n}=0\) and \(x_n=0\) for all sufficiently large \(n\).
This case, where \(\omega=\prod_{n\in\mathbb N^*}\omega'_n\),
corresponds in the dual setting of connected cocommutative Hopf
algebras to the case studied by M.~E.~Sweedler; cf.\
\textup{[Sw67]}, Theorem~3.}

Let \(G\) be an infinitesimal formal \(k\)-group and let \(B\)
be the affine algebra of a homogeneous space of \(G\). It follows
from \SGARefLink{sga3:VIIB:5.2:theorem}{Theorem~5.2} that, for every \(r\in\mathbb N\), the algebra
\[
  B
  \big/
  \bigl((x^{p^r})\bigr)_{x\in\mathfrak r(B)}
\]
is stationary. In particular, \(B\) is an inverse limit of
stationary algebras.\footnote{\label{sga3:VIIB:nde:166}Such an
inverse limit need not itself be a stable profinite
\(k\)-algebra in the sense of the preceding editor's note. For
example, let \(S\) be the ordinary \(k\)-vector space of
sequences \((u_1,u_2,\ldots)\) of elements of \(k\), and put
\(\omega=S^*\). Then \(\omega\) is the direct sum in
\(\operatorname{PC}(k)\) of copies \(kx_n\) of \(k\), for
\(n\in\mathbb N^*\), where \(x_n(u)=u_n\). Thus this is the case
of \SGARefLink{sga3:VIIB:5.2.1.A}{5.2.1.A}. The algebra
\[
  \mathcal A
  =
  k[[\omega]]
  \big/
  \bigl((x_n^{p^n})\bigr)_{n\in\mathbb N^*}
\]
has \(\omega_{\mathcal A/k}=\omega\) and
\(\omega_n=\prod_{i=1}^{n}kx_i\), but it is not stable:
\(\operatorname{Spf}(\mathcal A)\) represents the functor that
assigns to every \(C\in\operatorname{Alf}/k\) the infinite
sequences \((x_n)_{n\in\mathbb N^*}\) of elements of \(C\) such
that \(x_n^{p^n}=0\) for every \(n\).}

\medskip
\noindent\textup{(b)}
I do not know whether, with the notation of \SGARefLink{sga3:VIIB:5.2.1:section-proposition}{5.2.1}, \(A\) and
\(B\) can be chosen so that there exists no section
\[
  \sigma:\omega\longrightarrow I
\]
compatible with the filtrations.\footnote{%
\label{sga3:VIIB:nde:167}The editors do not know either, outside
the cases considered in \SGARefLink{sga3:VIIB:5.2.1.A}{5.2.1.A} and \SGARefLink{sga3:VIIB:5.2.1.B}{5.2.1.B}.}
Nevertheless, \(\omega\) may be any pseudocompact vector space
filtered by an increasing sequence of closed subspaces. Indeed,
given
\[
  \omega_1\subset\omega_2\subset\cdots\subset\omega,
\]
one can define on
\[
  B=A
  =
  k[[\omega]]
  \big/
  \bigl((x^{p^r})\bigr)_{x\in\omega_r}
\]
a diagonal morphism
\[
  \Delta_A:A\longrightarrow A\widehat\otimes_kA
\]
satisfying conditions \textup{\SGARefLink{sga3:vii:paragraph:5-1:equation:i}{(i)}}, \textup{\SGARefLink{sga3:vii:paragraph:5-1:equation:ii}{(ii)}}, and
\textup{(iii)} of \SGARefLink{sga3:VIIB:2.1}{2.1}. It suffices to put
\[
  \Delta_A(y)
  =
  y\widehat\otimes1+1\widehat\otimes y
\]
when \(y\) is the image in \(A\) of an element of \(\omega\).

\subsection*{5.4. Corollary}
\phantomsection
\label{sga3:VIIB:5.4}
\addcontentsline{toc}{subsection}{5.4. Corollary}

\noindent\SGAInlineRefTarget{sga3:vii:corollary:5-4}\textbf{Corollary 5.4.}
{\itshape Let \(G\) be an algebraic group over a perfect field
\(k\) of characteristic \(p>0\), let \(H\) be an algebraic
subgroup of \(G\), let \(e\) be the image of the identity element
of \(G\) in \(G/H\), and let \(A\) be the local ring of \(G/H\)
at \(e\). Then \(\widehat A\) is isomorphic to an algebra of the
form
\[
  k[[X_1,\ldots,X_r,\ldots,X_s]]
  \big/
  (X_1^{p^{n_1}},\ldots,X_r^{p^{n_r}}).
\]}

\medskip
\noindent\textit{Proof.}
Consider the infinitesimal formal groups
\[
  \widehat G
  =
  \operatorname{Spf}(\widehat{\mathcal O}_{G,e}),
  \qquad
  \widehat H
  =
  \operatorname{Spf}(\widehat{\mathcal O}_{H,e}).
\]
By \SGARefLink{sga3:VIIB:1.3.4}{1.3.4}, the completion \(\widehat A\) of
\[
  A=\mathcal O_{G/H,e}
\]
is isomorphic to
\[
  \mathscr A(\widehat G/\widehat H),
\]
and the corollary follows from \SGARefLink{sga3:VIIB:5.2:theorem}{Theorem~5.2}\textup{\SGARefLink{sga3:VIIB:5.2:theorem:item:ii}{(ii)}}.\footnote{%
\label{sga3:VIIB:nde:168}See also \textup{[DG70]},
\(\S\)III.3, Theorem~6.1.}

% END INCLUDED SGA3_English_current_buildable_source_20260823/source/inputs/VII/tex/components/93_expose_VIIB_sections53_54_stationary_algebras_corollary_en.tex

% BEGIN INCLUDED SGA3_English_current_buildable_source_20260823/source/inputs/VII/tex/components/94_expose_VIIB_section55_complements_en.tex
% SGA 3, Expose VII_B, English reconstruction -- Loop 1.
% Scope: Section 5.5 in full, from its opening definitions through
% Remarks 5.5.2(b), on authority local pages 96--97.
% Combined-reader pages: 609--610. Re-edition printed pages: 598--599.
% Editorial note in scope: 169.
% Authority: Polo--Gille Exp7B-13oct24.pdf,
% SHA-256 D303690E0195654BBE7C3DAA5D89EA23E24CD709A3575217234AC19AE8836694.
% Jacob Reinhold's CC BY 4.0 Markdown was consulted as comparison lineage.
% No genuine diagram occurs in this component; the exact sequence is
% reconstructed as native LaTeX.

\subsection*{5.5. Complements}
\phantomsection
\label{sga3:VIIB:5.5}
\addcontentsline{toc}{subsection}{5.5. Complements}
\footnote{\label{sga3:VIIB:nde:169}The editors added this
subsection in order to give some consequences of \SGARefLink{sga3:vii:corollary:5-4}{5.4} mentioned
in Expos\'es III and VIA.}

Recall the following definitions. A Noetherian local ring \(A\)
is called a \emph{complete intersection} if its completion
\(\widehat A\) is the quotient of a complete regular Noetherian
local ring \(B\) by an ideal \(I\) generated by a regular
sequence of elements of \(B\) (cf.\ EGA IV\(_4\), 19.3.1).

On the other hand, let
\[
  \tau:Y\hookrightarrow X
\]
be a closed immersion of schemes. If \(y\in Y\), then \(\tau\)
is called a \emph{regular immersion at the point \(y\)} if the
kernel of
\[
  \mathcal O_{X,y}\longrightarrow\mathcal O_{Y,y}
\]
is generated by a regular sequence. If, moreover, \(X\) is
locally Noetherian and \(\tau\) is a regular immersion at every
point, then \(\tau\) is called a \emph{regular immersion}
(cf.\ \textit{loc.\ cit.}, Proposition~16.9.10 and
Definition~16.9.2).

\medskip
\noindent\textbf{Corollary 5.5.1.}
\phantomsection
\label{sga3:VIIB:5.5.1}
\textit{If \(G\) is an algebraic group over a field \(k\), then
the local ring \(\mathcal O_{G,e}\) is a complete intersection.}

\medskip
\noindent\textit{Proof.}
By EGA IV\(_4\), 19.3.4, we may assume that \(k\) is
algebraically closed. If \(\operatorname{char}(k)=0\), then
\(G\) is already known to be smooth (cf.\ \SGARefLink{sga3:VIIA:3.3.1}{3.3.1} or VIB, \SGARefLink{cor:sga3-VIB-1.6.1}{1.6.1});
hence \(\widehat{\mathcal O}_{G,e}\) is a \(k\)-algebra of
formal power series, by EGA IV\(_4\), 17.5.3
\((\mathrm d^{\prime\prime})\). If
\(\operatorname{char}(k)=p>0\), then \SGARefLink{sga3:vii:corollary:5-4}{5.4}, applied with
\(H=\{e\}\), shows that \(\widehat{\mathcal O}_{G,e}\) is a
complete intersection.

\medskip
\noindent\textbf{Remarks 5.5.2.}
\phantomsection
\label{sga3:VIIB:5.5.2}
\textit{Let \(k\) be a field, let \(G\) be a smooth algebraic
\(k\)-group, and let \(H\) be a closed subgroup of \(G\).}

\smallskip
\noindent\textit{(a)}
We saw in \SGARefLink{sga3:iii:proposition:4-15}{Expos\'e III, 4.15}, that the immersion
\[
  H\hookrightarrow G
\]
is regular. This can also be deduced from \SGARefLink{sga3:vii:corollary:5-4}{5.4} as follows. As in
\textit{loc.\ cit.}, we may assume that \(k\) is algebraically
closed, and it is enough to show that the kernel \(I\) of
\[
  \mathcal O_{G,e}\longrightarrow\mathcal O_{H,e}
\]
is generated by a regular sequence. Put
\[
  A=\mathcal O_{G,e},
  \qquad
  \widehat G=\operatorname{Spf}(\widehat A),
  \qquad
  \widehat H
  =
  \operatorname{Spf}(\widehat{\mathcal O}_{H,e}).
\]
Since \(A\) is Noetherian, there is an exact sequence
\[
  0
  \longrightarrow I\otimes_A\widehat A
  \longrightarrow \widehat A
  \xrightarrow{\ \pi\ }
  \mathscr A(\widehat H)
  \longrightarrow 0,
\]
and \(\widehat A\) is faithfully flat over \(A\). Therefore, by
EGA IV\(_4\), 16.9.10(ii) and 19.1.5(ii), it is enough to show
that the kernel
\[
  \widehat I=I\otimes_A\widehat A
\]
of \(\pi\) is generated by a regular sequence of elements of
\(\widehat A\).

Since \(G\) is smooth, \(\widehat A\) is reduced. By \SGARefLink{sga3:vii:corollary:5-4}{5.4}, the
subalgebra
\[
  B=\mathscr A(\widehat G/\widehat H)
\]
is therefore isomorphic to a formal power series algebra
\[
  k[[x_1,\ldots,x_n]].
\]
Consequently, the identity section of
\(\widehat G/\widehat H\) is defined in \(B\) by the regular
sequence \((x_1,\ldots,x_n)\). Since \(\widehat A\) is
Noetherian, the ideal \(J\) of \(\widehat A\) generated by
\(x_1,\ldots,x_n\) is closed, and hence equals \(\widehat I\),
by Corollary~1.4. Moreover, \(\widehat A\) is topologically
flat, and hence flat over \(B\) (cf.\ \SGARefLink{sga3:VIIB:0.3.8}{0.3.8}). Thus
\((x_1,\ldots,x_n)\) is a regular sequence in \(\widehat A\),
by EGA IV\(_4\), 19.1.5(ii).

\smallskip
\noindent\textit{(b)}
The following more precise result can also be deduced from
\SGARefLink{sga3:VIIB:5.2:theorem:item:ii}{5.2(ii)}. Suppose that \(k\) is perfect. Applying \SGARefLink{sga3:VIIB:5.2:theorem:item:ii}{5.2(ii)} to the
algebra
\[
  C=\mathscr A(H),
\]
there are a basis
\[
  (y_1,\ldots,y_{r+s})
\]
of \(\omega_H\) and integers
\[
  1\leq n_1\leq\cdots\leq n_r
\]
such that \(\mathscr A(H)\) is isomorphic to the quotient of
\[
  k[[y_1,\ldots,y_{r+s}]]
\]
by the ideal generated by
\[
  y_i^{p^{n_i}},
  \qquad i=1,\ldots,r.
\]
Lift the \(y_i\) to elements \(x_i\) of \(\omega_G\), and extend
\((x_1,\ldots,x_{r+s})\) to a basis
\((x_1,\ldots,x_n)\) of \(\omega_G\). Since
\(\mathscr A(G)\) is reduced, the morphism
\[
  k[[x_1,\ldots,x_n]]
  \longrightarrow
  \mathscr A(G)
\]
is an isomorphism, by \SGARefLink{sga3:VIIB:5.2:theorem:item:iii}{5.2(iii)}.

Thus there exists a ``coordinate system''
\((x_1,\ldots,x_n)\) on \(G\), that is, an isomorphism
\[
  \mathscr A(G)
  \simeq
  k[[x_1,\ldots,x_n]],
\]
such that \(H\) is defined by the equations
\[
  x_i^{p^{n_i}}=0
  \quad (i=1,\ldots,r),
  \qquad
  x_i=0
  \quad (i>r+s).
\]
This is ``Dieudonn\'e's theorem''; compare \textup{[Di55]},
\(\S\)19, Theorem~6, and \textup{[Di73]}, II, \(\S\)\SGARefLink{sga3:VIIB:3.2:prop}{3.2},
Proposition~3 and the material preceding it.

% Next cursor: the terminal VIIB bibliography on authority local
% page 97, combined-reader page 610, re-edition printed page 599.

% END INCLUDED SGA3_English_current_buildable_source_20260823/source/inputs/VII/tex/components/94_expose_VIIB_section55_complements_en.tex

% BEGIN INCLUDED SGA3_English_current_buildable_source_20260823/source/inputs/VII/tex/components/95_expose_VIIB_bibliography_en.tex
% SGA 3, Expose VII_B, English reconstruction -- Loop 1.
% Scope: the complete terminal bibliography, from [CA] through [We95],
% on authority local pages 97--99.
% Combined-reader pages: 610--612. Re-edition printed pages: 599--601.
% Editorial note in scope: 170.
% Authority: Polo--Gille Exp7B-13oct24.pdf,
% SHA-256 D303690E0195654BBE7C3DAA5D89EA23E24CD709A3575217234AC19AE8836694.
% Jacob Reinhold's CC BY 4.0 Markdown was consulted as comparison lineage.
% Bibliographic titles, author spellings, journal abbreviations, and
% publication data follow the Polo--Gille authority.

\chapter*{Bibliography}
\phantomsection
\label{sga3:VIIB:bibliography}
\addcontentsline{toc}{chapter}{Bibliography}

\begin{description}
  \item[\textup{[CA]}] P.~Gabriel,
  \emph{Des cat\'egories ab\'eliennes},
  Bull.\ Soc.\ Math.\ France \textbf{90} (1962), 323--448.
  \footnote{\label{sga3:VIIB:nde:170}The editors added to this
  reference, which appeared in the original, the references that
  follow.}

  \item[\textup{[Ab80]}] E.~Abe,
  \emph{Hopf algebras}, Cambridge Univ.\ Press, 1980.

  \item[\textup{[Br00]}] C.~Breuil,
  \emph{Groupes \(p\)-divisibles, groupes finis et modules
  filtr\'es}, Ann.\ of Math.\ \textbf{152} (2000),
  n\textsuperscript{o}~2, 489--549.

  \item[\textup{[BAlg]}] N.~Bourbaki,
  \emph{Alg\`ebre}, Chap.\ I--III et IV--VII, Hermann, 1970,
  et Masson, 1981.

  \item[\textup{[BAC]}] N.~Bourbaki,
  \emph{Alg\`ebre commutative}, Chap.\ I--IV, Masson, 1985.

  \item[\textup{[BEns]}] N.~Bourbaki,
  \emph{Th\'eorie des ensembles}, Chap.\ I--IV, Hermann, 1970.

  \item[\textup{[BLie]}] N.~Bourbaki,
  \emph{Groupes et alg\`ebres de Lie}, Chap.\ II--III,
  Hermann, 1970.

  \item[\textup{[Ca62]}] P.~Cartier,
  \emph{Groupes alg\'ebriques et groupes formels}, pp.~87--111
  in: \emph{Colloque sur la th\'eorie des groupes alg\'ebriques}
  (Bruxelles, 1962), Gauthier-Villars, 1962.

  \item[\textup{[DG70]}] M.~Demazure, P.~Gabriel,
  \emph{Groupes alg\'ebriques}, Masson \& North-Holland, 1970.

  \item[\textup{[De72]}] M.~Demazure,
  \emph{Lectures on \(p\)-divisible groups},
  Lect.\ Notes Math.\ \textbf{302}, Springer-Verlag, 1972.

  \item[\textup{[Di55]}] J.~Dieudonn\'e,
  \emph{Groupes de Lie et hyperalg\`ebres de Lie sur un corps
  de caract\'eristique \(p>0\)}, Comm.\ Math.\ Helv.\
  \textbf{28} (1955), 87--118.

  \item[\textup{[Di73]}] J.~Dieudonn\'e,
  \emph{Introduction to the theory of formal groups},
  Marcel Dekker, 1973.

  \item[\textup{[Fo77]}] J.-M.~Fontaine,
  \emph{Groupes \(p\)-divisibles sur les corps locaux},
  Ast\'erisque \textbf{47--48}, Soc.\ Math.\ France, 1977.

  \item[\textup{[Gr57]}] A.~Grothendieck,
  \emph{Sur quelques points d'alg\`ebre homologique},
  T\^ohoku Math.\ J.\ \textbf{9} (1957), 119--221.

  \item[\textup{[Gr74]}] A.~Grothendieck,
  \emph{Groupes de Barsotti-Tate et cristaux de Dieudonn\'e},
  Presses Univ.\ Montr\'eal, 1974.

  \item[\textup{[HS69]}] R.~G.~Heyneman, M.~E.~Sweedler,
  \emph{Affine Hopf Algebras I}, J.\ Algebra \textbf{13}
  (1969), 194--241.

  \item[\textup{[Il85]}] L.~Illusie,
  \emph{D\'eformations de groupes de Barsotti-Tate, d'apr\`es
  A.~Grothendieck}, pp.~151--198 in:
  \emph{S\'eminaire sur les pinceaux arithm\'etiques: la
  conjecture de Mordell}, Ast\'erisque \textbf{127},
  Soc.\ Math.\ France, 1985.

  \item[\textup{[Ja65]}] I.~M.~James,
  review of \textup{[MM65]} in Math.\ Reviews, MR0174052.

  \item[\textup{[La75]}] M.~Lazard,
  \emph{Commutative formal groups},
  Lect.\ Notes Math.\ \textbf{443}, Springer-Verlag, 1975.

  \item[\textup{[LT65]}] J.~Lubin, J.~Tate,
  \emph{Formal complex multiplication in local fields},
  Ann.\ of Math.\ \textbf{81} (1965), 380--387.

  \item[\textup{[LT66]}] J.~Lubin, J.~Tate,
  \emph{Formal moduli for one-parameter formal Lie groups},
  Bull.\ Soc.\ Math.\ France \textbf{94} (1966), 49--60.

  \item[\textup{[Ma91]}] A.~Masuoka,
  \emph{On Hopf algebras with cocommutative coradicals},
  J.\ Algebra \textbf{144} (1991), 451--466.

  \item[\textup{[Me72]}] W.~Messing,
  \emph{The crystals associated with Barsotti-Tate Groups:
  with applications to abelian schemes},
  Lect.\ Notes Math.\ \textbf{264}, Springer-Verlag, 1972.

  \item[\textup{[MM65]}] J.~W.~Milnor, J.~C.~Moore,
  \emph{On the structure of Hopf algebras},
  Ann.\ of Math.\ \textbf{81} (1965), 211--264.

  \item[\textup{[Mi65]}] B.~Mitchell,
  \emph{Theory of categories}, Academic Press, 1965.

  \item[\textup{[Ne75]}] K.~Newman,
  \emph{A correspondence between bi-ideals and sub-Hopf
  algebras in cocommutative Hopf algebras},
  J.\ Algebra \textbf{36} (1975), 1--15.

  \item[\textup{[Po73]}] N.~Popescu,
  \emph{Abelian categories with applications to rings and
  modules}, Academic Press, 1973.

  \item[\textup{[Sch90]}] H.~J.~Schneider,
  \emph{Principal homogeneous spaces for arbitrary Hopf
  algebras}, Israel J.\ Math.\ \textbf{72} (1990),
  n\textsuperscript{os}~1--2, 167--195.

  \item[\textup{[Sw67]}] M.~Sweedler,
  \emph{Hopf algebras with one grouplike element},
  Trans.\ Amer.\ Math.\ Soc.\ \textbf{127} (1967),
  n\textsuperscript{o}~3, 515--526.

  \item[\textup{[Sw69]}] M.~Sweedler,
  \emph{Hopf algebras}, Benjamin, 1969.

  \item[\textup{[Tak72]}] M.~Takeuchi,
  \emph{A correspondence between Hopf ideals and sub-Hopf
  algebras}, Manuscripta Math.\ \textbf{7} (1972), 251--270.

  \item[\textup{[Tak79]}] M.~Takeuchi,
  \emph{Relative Hopf modules--Equivalence and freeness
  criteria}, J.\ Algebra \textbf{60} (1979), 452--471.

  \item[\textup{[Ta67]}] J.~Tate,
  \emph{\(p\)-divisible groups}, pp.~158--183 in:
  \emph{Proc.\ Conf.\ Local Fields} (Driebergen, 1966)
  (\'ed.\ T.~A.~Springer), Springer-Verlag, 1967.

  \item[\textup{[We95]}] C.~A.~Weibel,
  \emph{An introduction to homological algebra},
  Cambridge Univ.\ Press, 1995.
\end{description}

% End of Expose VII_B and of Expose VII.
% The next excluded cursor is Expose VIII local p. 1, combined-reader
% p. 617. Tome I index pages 613--616 are not part of Expose VII.

% END INCLUDED SGA3_English_current_buildable_source_20260823/source/inputs/VII/tex/components/95_expose_VIIB_bibliography_en.tex



% END INCLUDED SGA3_English_current_buildable_source_20260823/source/inputs/VII/body.tex


\clearpage
\chapter*{Index --- Tome I}
\addcontentsline{toc}{chapter}{Index --- Tome I}
\markboth{Index --- Tome I}{Index --- Tome I}
\noindent\emph{Page numbers refer to the Polo--Gille re-edition.}
\medskip
\begin{multicols}{2}
\small
\setlength{\parskip}{1.2pt}
\interlinepenalty=10000
\SGAInputRoot{TOMEI_INDEX}
% BEGIN INCLUDED SGA3_English_current_buildable_source_20260823/source/inputs/TOMEI_INDEX/body.tex
% Exact R15 special component.
% BEGIN INCLUDED SGA3_English_current_buildable_source_20260823/source/inputs/TOMEI_INDEX/index_tomeI/tex/components/00_tomeI_index_entries_en.tex
\Idx{Abelian variety}{428}
\Idx{Adjoint representation}{70, 90}
\Idx{Admissible composition law}{129}
\Idx{Affine groups}{20, 403}
\Idx{\(\mathbf{Alf}_{/k}\) (category of finite-length \(k\)-algebras)}{515}
\Idx{Hopf algebra}{21, 427, 546, 565, 582}
\Idx{Lie algebra}{}
\Idx{\quad of a formal group}{559}
\Idx{\quad of a group scheme}{89, 90}
\Idx{Lie algebra of the centralizer (normalizer)}{}
\Idx{\quad of a subgroup}{94}
\Idx{\quad of a submodule}{97}
\Idx{Truncated formal power-series algebra}{589}
\Idx{Distribution algebra of \(G\)}{449}
\Idx{Restricted enveloping algebra}{473, 476, 575}
\Idx{Infinitesimal algebra of \(G\)}{449}
\Idx{\(\mathbf{Alp}_{/k}\) (category of profinite \(k\)-algebras)}{514}
\Idx{\(\alpha_{p,k}\)}{99}
\Idx{Anti-affine \(k\)-groups}{429}
\Idx{\(a(P)\) (sheaf associated with the presheaf \(P\))}{209}
\Idx{\(\operatorname{Aut}(X)\)}{9, 12}
\Idx{Cocommutative bialgebra}{544}
\Idx{Bigebra (see also bialgebra)}{20}
\Idx{Good \(\mathcal O_S\)-module}{81}
\Idx{Good group functor}{83}
\Idx{Locally Boolean space}{363}
\Idx{Cantor space}{364}
\Idx{Characteristic subgroup}{377}
\Idx{Category of commutative \(k\)-groups}{}
\Idx{\quad affine}{555}
\Idx{\quad algebraic}{323}
\Idx{\quad affine algebraic}{323}
\Idx{\quad quasi-compact}{327}
\Idx{\(\widehat{\mathcal C}\)}{1}
\Idx{\(\widehat{\mathcal C}_{/S}\)}{3}
\Idx{\(\widehat{\mathcal C}^{\,d}_{/S}\), \(\widehat{\mathcal C}^{\,b}_{hS}\)}{3}
\Idx{Central subgroup}{17}
\Idx{Centralizer}{15, 366}
\Idx{\(\underline{\operatorname{Centr}}_G(X)\)}{15}
\Idx{\(\underline{\operatorname{Centr}}(G)\)}{17}
\Idx{Chevalley theorems}{420, 422, 428}
\Idx{Coalgebra \(H(X)\) of a formal variety \(X\)}{530}
\Idx{Cocommutative coalgebras}{455}
\Idx{Coalgebras in groups}{457, 477}
\Idx{Cogebra (see also coalgebra)}{409}
\Idx{Cohen--Macaulay}{294}
\Idx{Commutators}{}
\Idx{\quad sheaf of commutators}{388}
\Idx{\quad commutator subgroup}{386}
\Idx{Comodules}{28, 409}
\Idx{Connected component \(G^0\)}{299, 308, 345}
\Idx{Local components \(A_m\), \(M_m\)}{502, 511}
\Idx{Condition (E)}{59}
\Idx{Geometrically connected}{297}
\Idx{Conormal sheaf}{46, 147}
\Idx{Cokernel of a pair of arrows}{178, 186, 250, 310, 492, 519}
\Idx{Constant objects}{10}
\Idx{\quad in \((\mathrm{Sch})\)}{20}
\Idx{Equivalence pair}{254}
\Idx{Covering morphism or family}{200}
\Idx{Covariant bialgebra or algebra}{545, 547}
\Idx{Sieves}{196}
\Idx{Bracket on \(\operatorname{Lie}(G/S)\)}{85}
\Idx{\(c(X,Y,f)\)}{152}
\Idx{\(\Delta^{(1)}_{X/S}\)}{105}
\Idx{\(S\)-derivation of an \(S\)-morphism \(Y\to X\)}{443}
\Idx{Invariant derivations}{89}

\Idx{Descent}{}
\Idx{\quad descent datum}{182}
\Idx{\quad descent morphism}{183}
\Idx{Descent of projectivity}{413}
\Idx{Effective-descent morphism}{183}
\Idx{\(S\)-deviations of order \(\leq n\)}{441}
\Idx{Fiber dimension}{350}
\Idx{\(D_{\mathcal O_S}(M)\)}{52}
\Idx{Cartier duality}{461, 547}
\Idx{\(d(Y_0,Y)\)}{148}
\Idx{Affine envelope \(G_{\mathrm{af}}\)}{405, 427}
\Idx{Epimorphisms}{}
\Idx{\quad effective}{178}
\Idx{\quad universally effective}{178}
\Idx{\quad universal}{177}
\Idx{Equivariant objects and modules}{38}
\Idx{Quotient ringed space}{250}
\Idx{Tangent space at \(u\)}{56}
\Idx{Homogeneous spaces of formal groups}{582, 590}
\Idx{Essentially free}{367}
\Idx{Étale}{}
\Idx{\quad formally étale algebra}{539}
\Idx{\quad formal group}{556}
\Idx{\quad formal variety}{539}
\Idx{Extension by zero}{416}
\Idx{Sheaf associated with a presheaf}{208}
\Idx{Relative sheaf}{222}
\Idx{Quotient sheaves}{217, 259, 272, 277, 280, 283, 286, 287, 393, 492, 493, 537, 552}
\Idx{Vector bundle \(V(F)\)}{25}
\Idx{Principal homogeneous bundle}{231}
\Idx{Tangent bundle}{55}
\Idx{Faithful action}{17}
\Idx{Functor \(\operatorname{Aut}(X)\)}{9, 12}
\Idx{Ring functor \(\mathcal O\) on \((\mathrm{Sch})\)}{22}
\Idx{Functor \(\operatorname{Hom}(X,Y)\)}{7}
\Idx{Functor \(\operatorname{Hom}_{Z/S}(X,Y)\)}{50}
\Idx{Functor \(\prod_{Z/S}Y\)}{51, 368, 373}
\Idx{Functor \(X\mapsto X_e\)}{542}
\Idx{Module functors \(V,W\)}{24}
\Idx{Formal scheme}{503, 517}
\Idx{Formal variety}{518}
\Idx{Formally étale algebra}{539}
\Idx{Formally homogeneous space}{246}
\Idx{Formally principal homogeneous}{102, 230}
\Idx{Frobenius morphism}{}
\Idx{\quad absolute, \(\operatorname{fr}(S)\)}{461, 574}
\Idx{\quad relative, \(\operatorname{Fr}(X/S)\)}{461, 573}
\Idx{\(\Gamma^*(M)\)}{522}
\Idx{\(\mathbb G_a\)}{21}
\Idx{\(\mathbb G_m\)}{22}
\Idx{\(G\)-\(\mathcal O\)-modules}{19}
\Idx{\(G\)-\(\mathcal O_S\)-modules}{27, 43}
\Idx{\(G_{\mathrm{red}}\) over a perfect field}{291}
\Idx{\(G_{\mathrm{red}}\) over a nonperfect field}{297}
\Idx{Group in a category}{12, 13}
\Idx{Object with an operator group}{15}
\Idx{Étale group \(G/G^0\)}{323}
\Idx{Smooth group over a field}{296}
\Idx{Groups with connected fibers}{353}
\Idx{Affine groups over a Dedekind scheme \(S\)}{430, 431}
\Idx{Diagonalizable groups}{23}
\Idx{\quad cohomology of}{37}
\Idx{Quasi-compact groups over a field \(k\)}{326}
\Idx{Groupoids}{251}
\Idx{\(H\)-set}{62}
\Idx{\(h_A\), the functor \(\operatorname{Hom}_{\mathbf{Alp}/k}(A,-)\)}{515}
\Idx{Height \(\leq 1\), for \(S\)-groups and formal groups}{489, 492, 579}
\Idx{Height \(\leq n\), for an \(S\)-group or formal group}{464, 575, 590}
\Idx{\(H^\vee_0(S,P)\)}{205}
\Idx{\(h^c_M\), the functor \(\operatorname{Hom}_c(M,-)\)}{506}
\Idx{Homogeneous space}{246}
\Idx{Crossed homomorphisms}{74}
\Idx{\(\operatorname{Hom}(X,Y)\)}{7}
\Idx{\(\operatorname{Hom}_{\mathcal O}(F,F')\)}{19}
\Idx{\(\operatorname{Hom}_{Z/S}(X,Y)\)}{50}
\Idx{\(\operatorname{Hom}_{(Z/S)\text{-gr.}}(X,Y)\)}{74}
\Idx{\(H(X)\)}{530}
\Idx{\(h_X\)}{1}
\Idx{Regular immersion}{149, 161, 597}
\Idx{Infinitesimal \(S\)-group or formal group}{484, 557, 569}
\Idx{Infinitesimal}{}
\Idx{\quad automorphisms}{98}
\Idx{\quad endomorphisms}{68, 69}
\Idx{Complete intersection}{597}
\Idx{Locally complete intersection}{149, 160}
\Idx{Invariant subgroup}{17}
\Idx{Invariants, subobject of}{15, 28}
\Idx{\(I_S(M)\)}{52}
\Idx{Jacobson formulas}{471, 472}
\Idx{\(\operatorname{Ker}(f)\)}{17}
\Idx{Koszul complex}{149}
\Idx{\(\mathbf{LF}(A)\), category of finite-length \(A\)-modules}{506}
\Idx{Freely acting group}{188}
\Idx{\(\operatorname{Lie}(X/S)\)}{63, 70}
\Idx{\(\operatorname{Lie}(X/S,M)\)}{63, 70}
\Idx{\(\operatorname{Lie}(f)\), the derived morphism of \(f\)}{70}
\Idx{\(\operatorname{Lie}_0(X/S,M)\)}{77}
\Idx{\(\operatorname{Lie}(G/S)\)}{89, 90}

\Idx{\(\operatorname{Lie}(\mathbb G_{m,S}/S)\), \(\operatorname{Lie}(D_S(M)/S)\)}{91}
\Idx{\(\operatorname{Lie}({}_nG)\), the derived morphism of \(g\mapsto g^n\)}{64}
\Idx{Inverse limits of group schemes}{396}
\Idx{Linearity of affine algebraic groups}{}
\Idx{\quad flat over regular \(S\) of dimension \(\leq2\)}{433}
\Idx{\quad over a field}{417}
\Idx{Smoothness of \(G/\operatorname{Fr}^nG\)}{493}
\Idx{Smoothness of \(G\) over a characteristic-zero field \(k\)}{336, 571}
\Idx{\(\mathrm{LP}\)}{205}
\Idx{\(L^u_{X/S}\)}{56}
\Idx{\(L^u_{X/S}(M)\)}{56}
\Idx{\(L_X\)}{104}
\Idx{\(L^0_X\), \(L'_X\)}{114}
\Idx{\((M)\)-effectivity}{193}
\Idx{\(M^\dagger,N^*\)}{507}
\Idx{Monomorphisms and closed immersions}{301, 327, 335}
\Idx{The morphism \(g\mapsto g^n\) is étale}{495}
\Idx{The morphism \(g\mapsto g^n\) is zero}{496}
\Idx{\(\mu_n\), the \(n\)-th roots of unity}{23}
\Idx{Noether isomorphism theorems}{235}
\Idx{Scheme of dual numbers over \(S\)}{52}
\Idx{Normalizer}{15, 366}
\Idx{Norm}{}
\Idx{\quad of an invertible sheaf}{267}
\Idx{\quad of a finite locally free \(A\)-algebra}{261}
\Idx{\(\underline{\operatorname{Norm}}_G(X)\)}{15}
\Idx{Kernel}{17}
\Idx{\(N_{Y/X}\), conormal sheaf of \(Y\) in \(X\)}{46, 147}
\Idx{\(\Omega^1_{X/S}\)}{54, 89, 105}
\Idx{\(\omega^1_{X/S}\)}{89}
\Idx{\(\mathcal O\)-modules}{18}
\Idx{Flat \(\mathcal O_k\)-modules}{520}
\Idx{Differential operators}{441}
\Idx{Invariant differential operators on \(G\)}{449}
\Idx{\(p\)-Lie algebra}{470}
\Idx{\quad of an \(S\)-group}{480}
\Idx{\quad of a formal group}{575}
\Idx{\(\mathbf{PC}(A)\), category of pseudocompact \(A\)-modules}{503}
\Idx{\(\prod_{T/S}F\), cohomology of}{124}
\Idx{\(\prod_{Z/S}Y\)}{51}
\Idx{Pointed cogebra or bigebra}{556, 565}
\Idx{Geometrically pointwise irreducible}{306}
\Idx{Category of presheaves}{1}
\Idx{Separated presheaves}{202}
\Idx{Pre-equivalence relation}{249, 253}
\Idx{Pretopology}{201}
\Idx{Primitive elements}{459, 561, 566}
\Idx{Completed tensor product \(\widehat{\otimes}\)}{508, 509, 511, 516, 519}
\Idx{Profinite \(k\)-algebra}{514}
\Idx{Pseudobasis}{507}
\Idx{Pseudocompact}{}
\Idx{\quad ring}{501}
\Idx{\quad module}{503}
\Idx{Squarable morphism}{179}
\Idx{Quasi-sections}{268}
\Idx{Quasi-separated scheme}{307}
\Idx{Schemes quasi-separated over \(S\)}{363, 403, 409}
\Idx{Quotient by a groupoid}{}
\Idx{\quad finite and flat}{259}
\Idx{\quad flat, not necessarily proper}{277}
\Idx{\quad proper and flat}{272}
\Idx{Quotients \(G/H\) over a local Artinian ring \(A\)}{311, 315}
\Idx{Quotients in \(\mathbf{Vaf}_{/k}\)}{534, 537}
\Idx{Quotients by a group scheme}{282}
\Idx{\(r(A)\), the Jacobson radical of \(A\)}{502}
\Idx{Refinements}{199}
\Idx{Equivalence relation}{186}
\Idx{\quad effective}{191}
\Idx{\quad universally effective}{191}
\Idx{Representability}{}
\Idx{\quad of \(\prod_{Z/S}Y\)}{368, 373}
\Idx{\quad of centralizers}{371, 376}
\Idx{\quad of normalizers}{371, 376}
\Idx{Representable functor}{2}
\Idx{Solvable or nilpotent groups}{390}
\Idx{Weil restriction of scalars}{51}
\Idx{Retrocompact}{307}
\Idx{\(S\)-\(H\)-functor}{62}
\Idx{\((\mathrm{Sch})\), \((\mathrm{Sch})/S\), \((\mathrm{Sch}_{/S})\)}{20}
\Idx{Schematically dense}{374}
\Idx{Schematically dominant}{325, 326, 378, 427}
\Idx{Semidirect product}{16, 17, 552}
\Idx{Separated (every group over a field is)}{292}
\Idx{\(\widehat S_k(E)\), completed symmetric algebra}{523}
\Idx{Subgroup generated by \(f:X\to G\)}{378}
\Idx{\(\operatorname{Spec}^*(U)\)}{456}
\Idx{\(\operatorname{Spf}^*(C)\)}{530}
\Idx{\(\underline{\operatorname{Stab}}_G(x)\)}{15}
\Idx{Stabilizer}{15, 44}
\Idx{Strictly rational point}{293}
\Idx{Topology}{}
\Idx{\quad \((\mathrm{\acute etfg})\)}{244}
\Idx{\quad \((\mathrm{fpqc})\), \((\mathrm{fppf})\), \((\mathrm{\acute et})\), \((\mathrm{\acute etf})\)}{241}
\Idx{\quad canonical}{204}
\Idx{\quad chaotic (coarse)}{211, 245}
\Idx{\quad Zariski}{236}
\Idx{\quad flat}{537, 552}

\Idx{Topologies}{199}
\Idx{Topologically free}{507}
\Idx{Topologically nilpotent}{502}
\Idx{Topologically flat morphism}{527}
\Idx{Topologically flat formal variety}{529}
\Idx{Torsor}{231}
\Idx{Transporter}{366}
\Idx{Strict transporter}{366}
\Idx{Very good group functor}{88, 89, 459}
\Idx{\(T_{X/S}\)}{55}
\Idx{\(T_{X/S}(M)\)}{55}
\Idx{\(T'_{X/S}(M)\)}{77}
\Idx{Group of multiplicative type over \(k\)}{558}
\Idx{Unipotent \(k\)-group}{558, 563}
\Idx{\(\Upsilon(A)\), maximal open ideals of \(A\)}{502}
\Idx{\(\mathbf{Vaf}_{/k}\), category of formal varieties over \(k\)}{518}
\Idx{\(\mathbf{Vaf}_{/S}^{f\flat}\)}{525}
\Idx{\(\mathbf{Vaf}_{/k}^{\acute et}\)}{539}
\Idx{Abelian variety}{428}
\Idx{Formal variety \(\operatorname{Spf}^*(C)\) of a coalgebra \(C\)}{530}
\Idx{Verschiebung}{468}
\Idx{\(V^f_k(N)\)}{520}
\Idx{\(V(\Omega^1_{X/S})\)}{55}
\Idx{\(V^f_k(E)\)}{524}
\Idx{\(V^{f,0}_k(E)\)}{524}
\Idx{\(X\), underlying set of a scheme \(X\)}{249}
\Idx{\(X^+=\prod_{S_J/S}X_J\)}{102}
\Idx{\(X_e\), the étale formal variety associated with \(X\)}{542}
\Idx{\(\widehat X_{/S}\), the formal variety associated with the \(S\)-scheme \(X\)}{525}
\Idx{Yoneda lemma}{1}

% END INCLUDED SGA3_English_current_buildable_source_20260823/source/inputs/TOMEI_INDEX/index_tomeI/tex/components/00_tomeI_index_entries_en.tex


% END INCLUDED SGA3_English_current_buildable_source_20260823/source/inputs/TOMEI_INDEX/body.tex

\end{multicols}

\newgeometry{margin=30mm}
\SGATagsLeft

\SGAStartExpose{VIII}
% BEGIN INCLUDED SGA3_English_current_buildable_source_20260823/source/inputs/VIII/body.tex
% Mechanically extracted after \mainmatter from the frozen standalone master.

% BEGIN INCLUDED SGA3_English_current_buildable_source_20260823/source/inputs/VIII/tex/components/01_expose_VIII_opening_through_corollary_1_3.tex
% SGA 3, Exposé VIII, English reconstruction.
% Scope: opening and §1 through the proof of Corollary 1.3.
% Authority: Polo--Gille Exp8-8nov09.pdf, component-local pp. 1--4 (partial),
% complete-reader pp. 617--620 (partial).
% Jacob Reinhold's CC BY 4.0 Markdown lineage was used only as a phrasing
% comparison. The Polo--Gille PDF controls text, formulas, numbering, and notes.

\setcounter{footnote}{-1}
\SGARefTarget{sga3:viii:chapter:viii}\chapter*{Expos\'e VIII\\Diagonalizable Groups}
\addcontentsline{toc}{chapter}{Expos\'e VIII. Diagonalizable Groups}
\markboth{EXPOS\'E VIII. DIAGONALIZABLE GROUPS}%
  {EXPOS\'E VIII. DIAGONALIZABLE GROUPS}

\begin{center}
  \textit{by A. Grothendieck}\footnote{Editor's note: version \SGARefLink{sga3:viii:definition:1:1}{1.1} of
  November 8, 2009. Additions were made in \SGARefLink{sga3:viii:theorem:1:2}{1.2}, \SGARefLink{sga3:viii:corollary:1:4}{1.4}, \SGARefLink{sga3:viii:paragraph:1-7}{1.7}, \SGARefLink{sga3:viii:theorem:3:1}{3.1}, \SGARefLink{sga3:viii:corollary:3:4}{3.4}, \SGARefLink{sga3:viii:remark:4:5:1}{4.5.1},
  \SGARefLink{sga3:viii:theorem:6:4}{6.4}, and \SGARefLink{sga3:viii:remark:6:8}{6.8}; 1.5.1 and \SGARefLink{sga3:viii:section:7}{Section~7} remain to be revised.}
\end{center}

\SGARefTarget{sga3:viii:section:1}\section*{1. Biduality}
\addcontentsline{toc}{section}{1. Biduality}

% Source page: Exp8 component-local p. 1; complete-reader p. 617.
Let \(\mathscr C\) be a category, which as usual we identify with a full
subcategory of
\[
  \widehat{\mathscr C}
  =\operatorname{Hom}(\mathscr C^{\circ},\mathbf{Set})
\]
(cf. Expos\'e~I). Let \(I\) be a commutative group functor, that is, an
object of \(\widehat{\mathscr C}\) endowed with a commutative group structure
(cf. I, \SGARefLink{sga3:i:subsection:2-1}{2.1}).\footnote{Editor's note: the original has been expanded in what
follows.} For every \(X\in\operatorname{Ob}(\widehat{\mathscr C})\), the
object \(\operatorname{Hom}(X,I)\) carries the commutative group structure
induced by that of \(I\). For every group \(G\) in
\(\widehat{\mathscr C}\), put
\[
  D(G)=\operatorname{Hom}_{\mathrm{gr}}(G,I),
\]
the subobject of \(\operatorname{Hom}(G,I)\) defined, for every
\(S\in\operatorname{Ob}(\mathscr C)\), by
\SGARefTarget{sga3:viii:eq:1:x}
\begin{equation}
  \tag{x}
  D(G)(S)=\operatorname{Hom}_{S\text{-}\mathrm{gr}}(G_S,I_S),
\end{equation}
where \(G_S=G\times S\) and \(I_S=I\times S\) are regarded as
\(S\)-groups, that is, as groups in \(\widehat{\mathscr C}_{/S}\). Thus
\(D(G)\) is a \(\widehat{\mathscr C}\)-subgroup of
\(\operatorname{Hom}(G,I)\). In this way we obtain a contravariant functor
\(D\) from the category of \(\widehat{\mathscr C}\)-groups to the category
of commutative \(\widehat{\mathscr C}\)-groups.

The right-hand side of \(\SGARefLink{sga3:viii:eq:1:x}{(x)}\) may also
be interpreted as the subset of \(\operatorname{Hom}(G\times S,I)\)
consisting of morphisms \(G\times S\to I\) that are ``multiplicative in the
first argument \(G\).'' Moreover, the preceding formulas remain valid when
\(S\) is an arbitrary object of \(\widehat{\mathscr C}\), not necessarily
one arising from \(\mathscr C\).

Now take for \(S\) a group in \(\widehat{\mathscr C}\), denoted \(G'\). In
the left-hand side \(\operatorname{Hom}(G',D(G))\) of
\(\SGARefLink{sga3:viii:eq:1:x}{(x)}\), consider the subset
\(\operatorname{Hom}_{\mathrm{gr}}(G',D(G))\) consisting of morphisms that
respect the group structures of \(G'\) and \(D(G)\).

% Source page: Exp8 component-local p. 2; complete-reader p. 618.
It corresponds to the subset of \(\operatorname{Hom}(G\times G',I)\)
consisting of morphisms that are multiplicative in each of the two
arguments---which one may call \emph{bilinear morphisms from
\(G\times G'\) to \(I\)}, or \emph{pairings of \(G\) and \(G'\) with
values in \(I\)}. We therefore have
\SGARefTarget{sga3:viii:eq:1:xx}
\begin{equation}
  \tag{xx}
  \operatorname{Hom}_{\mathrm{gr}}(G',D(G))
  \xrightarrow{\ \sim\ }
  \operatorname{Hom}_{\mathrm{bil}}(G\times G',I),
\end{equation}
functorially in the pair \((G,G')\). Since the right-hand side is symmetric
in \(G\) and \(G'\), this gives a functorial bijection
\SGARefTarget{sga3:viii:eq:1:xxx}
\begin{equation}
  \tag{xxx}
  \operatorname{Hom}_{\mathrm{gr}}(G',D(G))
  \xrightarrow{\ \sim\ }
  \operatorname{Hom}_{\mathrm{gr}}(G,D(G')).
\end{equation}
In other words, giving a group homomorphism \(G'\to D(G)\) ``amounts to the
same thing'' as giving a group homomorphism \(G\to D(G')\): both amount to
giving a pairing \(G\times G'\to I\).

Apply this with \(G'=D(G)\) and with the identity homomorphism
\(G'\to D(G)\). We obtain a canonical homomorphism
\SGARefTarget{sga3:viii:eq:1:xxxx}
\begin{equation}
  \tag{xxxx}
  G\longrightarrow D(D(G)).
\end{equation}

\medskip
\noindent\SGARefTarget{sga3:viii:definition:1:0}\textbf{Definition 1.0.}
\footnote{Editor's note: the numbering \SGARefLink{sga3:viii:definition:1:0}{1.0}, \SGARefLink{sga3:viii:proposition:1:0:1}{1.0.1}, and so on has been added
to make the definitions and statements contained here explicit.}
\textit{The group \(G\) is called \emph{reflexive} (relative to \(I\)) if
\SGARefLink{sga3:viii:eq:1:xxxx}{the preceding homomorphism} is an isomorphism. This implies, in particular,
that \(G\) is commutative.}

We therefore have:

\medskip
\noindent\SGARefTarget{sga3:viii:proposition:1:0:1}\textbf{Proposition 1.0.1.}
\textit{The functor \(D\) induces an anti-equivalence from the category of
reflexive \(\widehat{\mathscr C}\)-groups to itself.}

In particular, if \(G\) and \(H\) are reflexive groups, \(D\) induces an
isomorphism
\[
  \operatorname{Hom}_{\mathrm{gr}}(G,H)
  \xrightarrow{\ \sim\ }
  \operatorname{Hom}_{\mathrm{gr}}(D(H),D(G)).
\]
It is enough, in fact, that \(H\) be reflexive, as follows from
\(\SGARefLink{sga3:viii:eq:1:xxx}{(xxx)}\).

\medskip
\noindent\SGARefTarget{sga3:viii:definition:1:0:2}\textbf{Definition 1.0.2.}
\textit{As usual, a \(\mathscr C\)-group \(G\) is called \emph{reflexive} if
it is reflexive as a \(\widehat{\mathscr C}\)-group, regardless of whether
\(D(G)\) is representable.}

Thus \(D\) gives an anti-equivalence from the category of reflexive
\(\mathscr C\)-groups \(G\) for which \(D(G)\) is representable to itself.

\medskip
\noindent\SGARefTarget{sga3:viii:remark:1:0:3}\textbf{Remark 1.0.3.}
To conclude these generalities, note that formation of the dual \(D(G)\)
commutes with base extension, which therefore carries reflexive groups to
reflexive groups.

We shall henceforth work in the following setting. Let
\(\mathscr C=(\mathbf{Sch})_{/S}\) be the category of preschemes over \(S\),
and let \(I=\mathbf G_{m,S}\) be the ``multiplicative group over \(S\)''
(cf. Expos\'e~I). For every ordinary group \(M\), consider the constant
\(S\)-group \(M_S\). For every group prescheme \(J\) over \(S\), there is a
canonical isomorphism, functorial in \(M\) and \(J\) and compatible with base
extension,
\[
  \operatorname{Hom}_{S\text{-}\mathrm{gr}}(M_S,J)
  =\operatorname{Hom}_{\mathrm{gr}}(M,J(S)).
\]

% Source page: Exp8 component-local p. 3; complete-reader p. 619.
Applying this with \(J=I=\mathbf G_{m,S}\) and with a variable \(S'\) over
\(S\), we obtain a functorial isomorphism
\SGARefTarget{sga3:viii:eq:1:0:4}
\begin{equation}
  \tag{1.0.4}
  D(M_S)(S')
  \xrightarrow{\ \sim\ }
  \operatorname{Hom}_{\mathrm{gr}}(M,\mathbf G_m(S')).
\end{equation}
We thus recover the functor already considered in I, \SGARefLink{sga3:i:subsection:4-4}{4.4}, also denoted
\(D_S(M)\), which is representable when \(M\) is commutative, since
\[
  D_S(M)=D(M_S)=\operatorname{Spec}\mathscr O_S(M),
\]
where \(\mathscr O_S(M)\) denotes the group algebra of \(M\) over
\(\mathscr O_S\). In general, moreover, \(D(M)\) is unchanged when \(M\)
is replaced by its abelianization, so no information is lost by assuming
that \(M\) is commutative.

\medskip
\noindent\SGARefTarget{sga3:viii:definition:1:1}\textbf{Definition 1.1.}
\textit{A group prescheme \(G\) over \(S\) is called \emph{diagonalizable}
if it is isomorphic to a scheme of the form}
\[
  D_S(M)=D(M_S)
  =\operatorname{Hom}_{S\text{-}\mathrm{gr}}(M_S,\mathbf G_m)
\]
\textit{for a suitable commutative group \(M\). It is called
\emph{locally diagonalizable} if every point of \(S\) has an open
neighborhood \(U\) such that \(G|_U\) is diagonalizable.}

\medskip
\noindent\SGARefTarget{sga3:viii:theorem:1:2}\textbf{Theorem 1.2.}
\textit{Let \(\Gamma\) be a constant commutative group scheme over \(S\);
that is, suppose that it is isomorphic to a group scheme of the form \(M_S\),
where \(M\) is an ordinary commutative group. Then \(\Gamma\) is reflexive:
the canonical homomorphism}
\[
  \Gamma\longrightarrow D(D(\Gamma))
\]
\textit{is an isomorphism.\footnote{Editor's note: the sentence that follows
has been added.} Hence the diagonalizable group \(D(M_S)\) is also
reflexive.}

In view of the definitions, this follows from \SGARefLink{sga3:viii:corollary:1:3}{the next statement}, applied to
a prescheme \(S'\) over \(S\).

\medskip
\noindent\SGARefTarget{sga3:viii:corollary:1:3}\textbf{Corollary 1.3.}
\textit{Let \(G=D(M_S)\). Every homomorphism of \(S\)-groups}
\[
  u:G\longrightarrow\mathbf G_{m,S}
\]
\textit{is defined by a unique section of \(M_S\) over \(S\), equivalently
by a unique locally constant map from \(S\) to \(M\).}

\emph{Proof.}
By definition,
\[
  \mathbf G_{m,S}
  =\operatorname{GL}(1)_S
  =\operatorname{Aut}_{\mathscr O_S\text{-}\mathrm{mod}}(\mathscr O_S).
\]
Thus giving a group homomorphism \(G\to\mathbf G_{m,S}\) is equivalent to
giving \(\mathscr O_S\) the structure of a
\(G\)-\(\mathscr O_S\)-module compatible with its natural
\(\mathscr O_S\)-module structure (cf. I, \SGARefLink{sga3:i:subsection:4-7}{4.7}). By I, \SGARefLink{sga3:i:proposition:4-7-3}{4.7.3}, this is also
equivalent to giving an \(M\)-grading on \(\mathscr O_S\), that is, a direct
sum decomposition
\[
  \mathscr O_S=\bigoplus_{m\in M}\mathscr L_m.
\]
It is well known that a direct summand of a locally free module of finite
type is locally free of finite type. Hence, near each point of \(S\), every
\(\mathscr L_m\) is either zero or free of rank one; in the latter case it
identifies with \(\mathscr O_S\) on that neighborhood. Let \(S_m\) be the
open subset of \(S\) on which the latter alternative occurs. The direct-sum
decomposition shows that the \(S_m\) cover \(S\) and are pairwise disjoint.

% Source page: Exp8 component-local p. 4 (opening); complete-reader p. 620.
Consequently, giving a group homomorphism \(G\to\mathbf G_{m,S}\) is
equivalent to giving a decomposition of \(S\) as a union of pairwise
disjoint open subsets \(S_m\), \(m\in M\), hence to giving a locally
constant map from \(S\) to \(M\). This proves \SGARefLink{sga3:viii:corollary:1:3}{Corollary~1.3} and therefore
\SGARefLink{sga3:viii:theorem:1:2}{Theorem~1.2}. \qed

% END INCLUDED SGA3_English_current_buildable_source_20260823/source/inputs/VIII/tex/components/01_expose_VIII_opening_through_corollary_1_3.tex

% BEGIN INCLUDED SGA3_English_current_buildable_source_20260823/source/inputs/VIII/tex/components/02_expose_VIII_corollary_1_4_through_remark_1_7_4.tex
% SGA 3, Exposé VIII, English reconstruction.
% Scope: Corollary 1.4 through Remark 1.7.4, ending immediately before §2.
% Authority: Polo--Gille Exp8-8nov09.pdf, component-local pp. 4--7 (partial),
% complete-reader pp. 620--623 (partial).
% The authority controls; comparison English is phrasing assistance only.

% Source page: Exp8 component-local p. 4; complete-reader p. 620.
\medskip
\noindent\SGARefTarget{sga3:viii:corollary:1:4}\textbf{Corollary 1.4.}
\textit{A diagonalizable group is reflexive, and therefore so is a locally
diagonalizable group. If \(M\) and \(N\) are ordinary commutative groups,
the natural homomorphism}
\SGARefTarget{sga3:viii:corollary:1:4:display:natural-isomorphism:01}\[
  \operatorname{Hom}_{S\text{-}\mathrm{gr}}(M_S,N_S)
  \longrightarrow
  \operatorname{Hom}_{S\text{-}\mathrm{gr}}(D(N_S),D(M_S))
\]
\textit{is bijective.}

\footnote{Editor's note: the discussion that follows has been expanded.}
\SGARefTarget{sga3:viii:corollary:1:4:discussion:01}Since \SGARefLink{sga3:viii:corollary:1:4:display:natural-isomorphism:01}{the preceding isomorphism} is compatible with base extension, it yields
an isomorphism of \(S\)-functors in groups
\SGARefTarget{sga3:viii:eq:1:4:1}
\begin{equation}
  \tag{1}
  \operatorname{\underline{Hom}}_{S\text{-}\mathrm{gr}}(M_S,N_S)
  \xrightarrow{\ \sim\ }
  \operatorname{\underline{Hom}}_{S\text{-}\mathrm{gr}}
    (D(N_S),D(M_S)).
\end{equation}
For every \(S\)-prescheme \(T\),
\[
  \operatorname{\underline{Hom}}_{S\text{-}\mathrm{gr}}(M_S,N_S)(T)
  =
  \operatorname{Hom}_{\mathrm{gr}}\bigl(M,\Gamma(N_T/T)\bigr),
\]
and by I, \SGARefLink{sga3:i:subsection:1:8}{1.8}, \(\Gamma(N_T/T)\) is the abelian group of locally constant
maps \(T\to N\). On the other hand, let
\(\operatorname{Hom}_{\mathrm{gr}}(M,N)_S\) be the constant \(S\)-group
associated with the ordinary abelian group
\(\operatorname{Hom}_{\mathrm{gr}}(M,N)\). There is an evident
homomorphism of \(S\)-functors in commutative groups
\SGARefTarget{sga3:viii:eq:1:4:2}
\begin{equation}
  \tag{2}
  \operatorname{Hom}_{\mathrm{gr}}(M,N)_S
  \xrightarrow{\ \theta\ }
  \operatorname{\underline{Hom}}_{S\text{-}\mathrm{gr}}(M_S,N_S),
\end{equation}
which is always a monomorphism and is an isomorphism if \(M\) is of finite
type.\footnote{Editor's note: write \(F(M)\) and \(G(M)\) for the left- and
right-hand sides of \SGARefLink{sga3:viii:eq:1:4:2}{(2)}, respectively. If \(M=M_1\oplus M_2\), there are
canonical isomorphisms
\(F(M)=F(M_1)\oplus F(M_2)\) and
\(G(M)=G(M_1)\oplus G(M_2)\). It is therefore enough to check that
\(\theta\) is an isomorphism for \(M=\mathbf Z/r\mathbf Z\), \(r\geq0\).
In that case \(F(M)=({}_rN)_S\), where \({}_rN\) is the kernel of
\(r\operatorname{id}_N\), and for every \(T\to S\) the homomorphism
\[
  F(M)(T)=\Gamma({}_rN_T/T)
  \longrightarrow
  G(M)(T)={}_r\Gamma(N_T/T)
\]
is bijective.}

Part \SGARefLink{sga3:viii:corollary:1:5:item:a}{(a)} of \SGARefLink{sga3:viii:corollary:1:5}{the following corollary} follows from \SGARefLink{sga3:viii:corollary:1:4:discussion:01}{the preceding discussion};
part \SGARefLink{sga3:viii:corollary:1:5:item:b}{(b)} follows from it by the descent results recalled in \SGARefLink{sga3:viii:paragraph:1-7}{1.7}.
\footnote{Editor's note: the original stated after \SGARefLink{sga3:viii:corollary:1:5}{1.5} that, more generally,
if \(G\) and \(H\) are locally diagonalizable and \(H\) is of finite type,
then \(\operatorname{\underline{Hom}}_{S\text{-}\mathrm{gr}}(G,H)\) is
representable. \SGARefLink{sga3:viii:corollary:1:5:item:b}{That assertion} has been included in \SGARefLink{sga3:viii:corollary:1:5}{1.5} and its proof made
explicit in \SGARefLink{sga3:viii:paragraph:1-7}{1.7}.}

\medskip
\noindent\SGARefTarget{sga3:viii:corollary:1:5}\textbf{Corollary 1.5.}
\begin{enumerate}
  \item[(a)]\SGARefTarget{sga3:viii:corollary:1:5:item:a}
  \textit{Let \(M\) and \(N\) be ordinary commutative groups, with \(M\) of
  finite type. Then there is an isomorphism}
  \[
    \operatorname{Hom}_{\mathrm{gr}}(M,N)_S
    \xrightarrow{\ \sim\ }
    \operatorname{\underline{Hom}}_{S\text{-}\mathrm{gr}}
      (D(N_S),D(M_S));
  \]
  \textit{consequently,
  \(\operatorname{\underline{Hom}}_{S\text{-}\mathrm{gr}}
  (D(N_S),D(M_S))\) is representable.}

  \item[(b)]\SGARefTarget{sga3:viii:corollary:1:5:item:b}
  \textit{More generally, if \(G\) and \(H\) are locally diagonalizable and
  \(H\) is of finite type, then
  \(\operatorname{\underline{Hom}}_{S\text{-}\mathrm{gr}}(G,H)\) is
  representable.}
\end{enumerate}

\footnote{Editor's note: a statement 1.5.1 treating the functor
\(\operatorname{\underline{Isom}}_{S\text{-}\mathrm{gr}}(G,H)\), considered
in X, \SGARefLink{sga3:x:corollary:5-10}{5.10} and 5.11, should be added.}
Corollary~\SGARefLink{sga3:viii:corollary:1:5}{1.5} gives:

\medskip
\noindent\SGARefTarget{sga3:viii:corollary:1:6}\textbf{Corollary 1.6.}
\textit{Under the hypotheses of \SGARefLink{sga3:viii:corollary:1:5}{1.5}, if \(S\) is connected, then}
\[
  \operatorname{Hom}_{S\text{-}\mathrm{gr}}(D_S(N),D_S(M))
  \xrightarrow{\ \sim\ }
  \operatorname{Hom}_{\mathrm{gr}}(M,N)
\]
\textit{and}
\[
  \operatorname{Isom}_{S\text{-}\mathrm{gr}}(D_S(N),D_S(M))
  \xrightarrow{\ \sim\ }
  \operatorname{Isom}_{\mathrm{gr}}(M,N).
\]

\medskip
\noindent\SGAInlineRefTarget{sga3:viii:paragraph:1-7}\SGARefTarget{sga3:viii:subsection:1:7}\textbf{1.7. Descent of
representability.}
\footnote{Editor's note: \SGARefLink{sga3:viii:paragraph:1-7}{this paragraph} has been added to make explicit the
``local on \(S\)'' character of the representability of certain sheaves on
\(S\), which is used repeatedly below and was used implicitly in the
original.}
We recall here several descent results that will be used frequently below.

\medskip
\noindent\SGARefTarget{sga3:viii:scholium:1:7:1}\textbf{Scholium 1.7.1.}
\textit{Let \(S\) be a prescheme and let \(X,Y,T\) be \(S\)-preschemes. If
\((T_i)\) is an open covering of \(T\), and if
\(T_{ij}=T_i\cap T_j=T_i\times_TT_j\), then, because giving a morphism of
\(T\)-preschemes \(X_T\to Y_T\) is local on \(T\), there is an exact
sequence of sets}
\SGARefTarget{sga3:viii:eq:1:7:1}
\begin{equation}
  \tag{1}
  \operatorname{Hom}_T(X_T,Y_T)
  \longrightarrow
  \prod_i\operatorname{Hom}_{T_i}(X_{T_i},Y_{T_i})
  \rightrightarrows
  \prod_{i,j}\operatorname{Hom}_{T_{ij}}(X_{T_{ij}},Y_{T_{ij}}).
\end{equation}
\textit{Thus \(\operatorname{\underline{Hom}}_S(X,Y)\) is a local
\(S\)-functor, that is, a sheaf on \((\mathbf{Sch})_{/S}\) for the Zariski
topology.}

More generally, by IV, \SGARefLink{sga3:iv:corollary:4:5:13}{4.5.13},
\(\operatorname{\underline{Hom}}_S(X,Y)\) is a sheaf on
\((\mathbf{Sch})_{/S}\) for every topology coarser than the canonical
topology, for example the fpqc topology. If \(G\) and \(H\) are group
preschemes over \(S\), it follows that the subfunctor
\(\operatorname{\underline{Hom}}_{S\text{-}\mathrm{gr}}(G,H)\) is an fpqc
sheaf, hence a fortiori a local functor.

\medskip
\noindent\SGARefTarget{sga3:viii:lemma:1:7:2}\textbf{Lemma 1.7.2.}
\footnote{Editor's note: see also Remark XI, \SGARefLink{sga3:xi:remark:3:4}{3.4}.}
\textit{Let \(F\) be a local \(S\)-functor.}
\begin{enumerate}
  \item[(i)]\SGARefTarget{sga3:viii:lemma:1:7:2:item:i}
  \textit{Suppose that there is an open covering \((S_i)\) of \(S\) such
  that the restriction \(F_i=F\times_SS_i\) of \(F\) to each \(S_i\) is
  representable by an \(S_i\)-prescheme \(X_i\). Then \(F\) is
  representable by an \(S\)-prescheme \(X\).}

  \item[(ii)]\SGARefTarget{sga3:viii:lemma:1:7:2:item:ii}
  \textit{Suppose that \(F\) is an fpqc sheaf and that there is a faithfully
  flat quasi-compact morphism \(S'\to S\) such that
  \(F'=F\times_SS'\) is representable by an \(S'\)-prescheme \(X'\). Then
  \(X'\) carries a descent datum relative to \(S'\to S\) (cf. IV, \SGARefLink{sga3:iv:definition:2:1}{2.1}).
  If this descent datum is effective---as it is when \(X'\) is affine over
  \(S'\)---then \(F\) is representable by an \(S\)-prescheme \(X\).}
\end{enumerate}

\emph{Proof.}
\SGARefLink{sga3:viii:lemma:1:7:2:item:i}{(i)} The hypothesis implies that \(X_i\times_SS_j\) and
\(X_j\times_SS_i\) both represent the restriction of \(F\) to
\(S_{ij}=S_i\times_SS_j\). The Yoneda lemma therefore gives a unique
isomorphism of \(S_{ij}\)-preschemes
\[
  c_{ji}:X_i\times_SS_j\xrightarrow{\ \sim\ }X_j\times_SS_i.
\]
Over \(S_{ijk}=S_i\times_SS_j\times_SS_k\), we obtain the commutative
diagram of isomorphisms
\SGARefTarget{sga3:viii:diagram:1:7:2:i}
\[
\begin{tikzcd}[column sep=small,row sep=large]
X_i\times_SS_j\times_SS_k
  \arrow[r,"c_{ji}\times\operatorname{id}_{S_k}"]
  \arrow[d,equal]
&
X_j\times_SS_i\times_SS_k
  \arrow[r,equal]
&
X_j\times_SS_k\times_SS_i
  \arrow[d,"c_{kj}\times\operatorname{id}_{S_i}"]
\\
X_i\times_SS_k\times_SS_j
  \arrow[r,"c_{ki}\times\operatorname{id}_{S_j}"]
&
X_k\times_SS_i\times_SS_j
  \arrow[r,equal]
&
X_k\times_SS_j\times_SS_i .
\end{tikzcd}
\]
Indeed, all of its objects represent the restriction of \(F\) to
\(S_{ijk}\). Thus the \(c_{ji}\) satisfy the usual cocycle relation
\(c_{kj}\circ c_{ji}=c_{ki}\).

It follows that the \(X_i\) glue to an \(S\)-prescheme \(X\) with
\(X\times_SS_i=X_i\) for every \(i\). For every \(Y\) over \(S_i\),
\SGARefTarget{sga3:viii:eq:1:7:2:star}
\begin{equation}
  \tag{\(\ast\)}
  F(Y)=F_i(Y)
  =\operatorname{Hom}_{S_i}(Y,X\times_SS_i)
  =\operatorname{Hom}_S(Y,X)=h_X(Y).
\end{equation}
For an arbitrary \(Y\to S\), the \(Y_i=Y\times_SS_i\) form an open covering
of \(Y\); put \(Y_{ij}=Y_i\times_YY_j=Y\times_SS_{ij}\). Since \(F\) is
local by hypothesis and \(h_X\) is local because the Zariski topology is
coarser than the canonical topology, both \(F(Y)\) and \(h_X(Y)\), in view
of \SGARefLink{sga3:viii:eq:1:7:2:star}{\((\ast)\)}, identify with the
kernels of the respective double arrows
\[
  \prod_iF(Y_i)\rightrightarrows\prod_{i,j}F(Y_{ij}),
  \qquad
  \prod_ih_X(Y_i)\rightrightarrows\prod_{i,j}h_X(Y_{ij}).
\]
This proves \SGARefLink{sga3:viii:lemma:1:7:2:item:i}{(i)}.

\SGARefLink{sga3:viii:lemma:1:7:2:item:ii}{(ii)} Put \(S''=S'\times_SS'\), and let \(S''_1\) and \(S''_2\) denote
\(S''\) regarded as an \(S'\)-prescheme through the first and second
projections. The functors \(F''_1=F'\times_{S'}S''_1\) and
\(F''_2=F'\times_{S'}S''_2\) are represented respectively by
\(X''_1=X'\times_{S'}S''_1\) and \(X''_2=X'\times_{S'}S''_2\). Since
\[
  F''_1=F\times_SS''=F''_2,
\]
there is a unique \(S''\)-isomorphism
\[
  c:X''_1\xrightarrow{\ \sim\ }X''_2.
\]
Let \(q_i\), respectively \(p_{ji}\), be the projection from
\(S'''=S'\times_SS'\times_SS'\) onto its \(i\)-th factor, respectively
onto its \(i\)-th and \(j\)-th factors. Let
\(X'''_i=X'\times_{S'}S'''_i\), where \(S'''_i\) means \(S'''\) viewed as
an \(S'\)-prescheme through \(q_i\), and let
\(p_{ji}^{*}(c):X'''_i\xrightarrow{\sim}X'''_j\) be obtained from \(c\)
by base change. We obtain the diagram
\SGARefTarget{sga3:viii:diagram:1:7:2:ii}
\[
\begin{tikzcd}[column sep=large,row sep=large]
X'''_1
  \arrow[r,"p_{21}^{*}(c)"]
  \arrow[dr,"p_{31}^{*}(c)"']
&
X'''_2
  \arrow[d,"p_{32}^{*}(c)"]
\\
&X'''_3 .
\end{tikzcd}
\]
All three objects represent the restriction of \(F\) to \(S'''\), so the
diagram commutes:
\[
  p_{32}^{*}(c)\circ p_{21}^{*}(c)=p_{31}^{*}(c).
\]
Thus \(c\) is a descent datum on \(X'\) relative to \(S'\to S\)
(cf. IV, \SGARefLink{sga3:iv:definition:2:1}{2.1}).

Suppose that this descent datum is effective, so there is an \(S\)-prescheme
\(X\) such that \(X'\simeq X\times_SS'\). This is the case when \(X'\) is
affine over \(S'\), by SGA~1, VIII, \SGARefLink{sga3:viii:proposition:2:1}{2.1}.\footnote{Editor's note: for another
effectivity criterion, see X, \SGARefLink{sga3:x:lemma:5-4}{5.4}--\SGARefLink{sga3:viii:corollary:5:6}{5.6}.} For every \(Y\to S'\), one then has
\SGARefTarget{sga3:viii:eq:1:7:2:double-star}
\begin{equation}
  \tag{\(\ast\ast\)}
  F(Y)=F'(Y)
  =\operatorname{Hom}_{S'}(Y,X\times_SS')
  =\operatorname{Hom}_S(Y,X)=h_X(Y).
\end{equation}
For arbitrary \(Y\to S\), put \(Y'=Y\times_SS'\) and
\(Y''=Y'\times_YY'\simeq Y\times_SS''\). Like \(S'\to S\), the morphism
\(Y'\to Y\) is faithfully flat and quasi-compact, hence an
\(\mathcal M\)-effective epimorphism, where \(\mathcal M\) is the family of
faithfully flat quasi-compact morphisms. Thus the equivalence relation
\[
  Y'\times_YY'\rightrightarrows Y'
\]
is \(\mathcal M\)-effective and has quotient \(Y\). Since \(F\) is an fpqc
sheaf by hypothesis, while \(h_X\) is one because the fpqc topology is
coarser than the canonical topology, both \(F(Y)\) and \(h_X(Y)\), in view
of \SGARefLink{sga3:viii:eq:1:7:2:double-star}{\((\ast\ast)\)}, identify
with the respective kernels of
\[
  F(Y')\rightrightarrows F(Y'\times_YY'),
  \qquad
  h_X(Y')\rightrightarrows h_X(Y'\times_YY').
\]
This proves \SGARefLink{sga3:viii:lemma:1:7:2:item:ii}{(ii)}. \qed

\medskip
\noindent\SGARefTarget{sga3:viii:corollary:1:7:3}\textbf{Corollary 1.7.3.}
\textit{Let \(F\) be an fpqc sheaf on \((\mathbf{Sch})_{/S}\). Suppose
there is an open covering \((S_i)\) of \(S\), and for each \(i\) a
faithfully flat quasi-compact morphism \(S'_i\to S_i\), such that
\(F'_i=F\times_SS'_i\) is representable by an \(S'_i\)-prescheme \(X'_i\)
affine over \(S'_i\). Then \(F\) is representable by an \(S\)-prescheme
\(X\) affine over \(S\), with \(X\times_SS'_i=X'_i\) for every \(i\). If,
moreover, every \(X'_i\to S'_i\) is a closed immersion (respectively, a
finite \'etale morphism), then so is \(X\to S\).}

The first assertion follows from
\SGARefLink{sga3:viii:lemma:1:7:2}{1.7.2}. For the second, it is enough to
check that every morphism \(X\times_SS_i\to S_i\) is a closed immersion
(respectively, finite and \'etale); this follows from EGA~IV\(_2\), 2.7.1
(respectively, together with IV\(_4\), 17.7.3).

\medskip
\noindent\SGARefTarget{sga3:viii:remark:1:7:4}\textbf{Remark 1.7.4.}
Assertion \SGARefLink{sga3:viii:corollary:1:5:item:b}{1.5(b)} follows,
as announced, from \SGARefLink{sga3:viii:scholium:1:7:1}{1.7.1} and
\SGARefLink{sga3:viii:lemma:1:7:2:item:i}{1.7.2(i)}.

% END INCLUDED SGA3_English_current_buildable_source_20260823/source/inputs/VIII/tex/components/02_expose_VIII_corollary_1_4_through_remark_1_7_4.tex

% BEGIN INCLUDED SGA3_English_current_buildable_source_20260823/source/inputs/VIII/tex/components/delegated_sections2_4/01_expose_VIII_s2_s4_pp007_013.tex
% SGA 3, Expose VIII, English reconstruction.
% Scope: Sections 2 and 3 in full, and Section 4 from its heading through
% the word "Moreover," in the proof of Corollary 4.7.
% Authority: Polo--Gille Exp8-8nov09.pdf, standalone pages 7--13
% (combined-reader pages 623--629).
% Jacob Reinhold's CC BY 4.0 Markdown lineage was used only as a phrasing
% comparison. The Polo--Gille PDF controls prose, formulas, numbering,
% notes, references, and diagrams.
%
% Start cursor: "2. Proprietes schematiques des groupes diagonalisables"
% and "Elles sont resumees dans la" (standalone p. 7).
% End cursor: "De plus," in the proof of Corollary 4.7
% (standalone p. 13); the sentence continues on standalone p. 14.

\setcounter{footnote}{10}

% Source page: Exp8 standalone p. 7; combined-reader p. 623.
\SGARefTarget{sga3:viii:section:2}\section*{2. Scheme-theoretic Properties of Diagonalizable Groups}
\addcontentsline{toc}{section}{2. Scheme-theoretic Properties of Diagonalizable Groups}

These properties are summarized in \SGARefLink{sga3:viii:proposition:2:1}{the following proposition}.

% Source page: Exp8 standalone p. 8; combined-reader p. 624.
\medskip
\noindent\SGARefTarget{sga3:viii:proposition:2:1}\textbf{Proposition 2.1.}
\textit{Let \(S\) be a nonempty prescheme, let \(M\) be an ordinary
commutative group, and let \(G=D(M_S)\) be the diagonalizable \(S\)-group
defined by \(M\). Then:}
\begin{enumerate}
  \item[(a)]\SGARefTarget{sga3:viii:proposition:2:1:item:a:01}
  \textit{\(G\) is faithfully flat over \(S\) and affine over \(S\)
  (a fortiori, quasi-compact over \(S\));}

  \item[(b)]\SGARefTarget{sga3:viii:proposition:2:1:item:b:01}
  \textit{\(M\) is of finite type if and only if \(G\) is of finite type
  over \(S\), if and only if \(G\) is of finite presentation over \(S\);}

  \item[(c)]\SGARefTarget{sga3:viii:proposition:2:1:item:c:01}
  \textit{\(M\) is finite if and only if \(G\) is finite over \(S\), if
  and only if \(G\) is of finite type over \(S\) and is annihilated by an
  integer \(n>0\). In that case,}
  \[
    \deg(G/S)=\operatorname{Card}(M);
  \]

  \item[\((c')\)]\SGARefTarget{sga3:viii:proposition:2:1:item:cprime:01}
  \textit{\(M\) is a torsion group if and only if \(G\) is integral over
  \(S\);}

  \item[(d)]\SGARefTarget{sga3:viii:proposition:2:1:item:d:01}
  \textit{\(M=0\) if and only if \(G\) is the unit \(S\)-group;}

  \item[(e)]\SGARefTarget{sga3:viii:proposition:2:1:item:e:01}
  \textit{\(M\) is of finite type and the order of its torsion subgroup is
  prime to the residue characteristics of \(S\) if and only if \(G\) is
  smooth over \(S\).}
\end{enumerate}

The verification of \SGARefLink{sga3:viii:proposition:2:1:item:a:01}{(a)}
through \SGARefLink{sga3:viii:proposition:2:1:item:d:01}{(d)} is immediate
and is left to the reader. Let us prove
\SGARefLink{sga3:viii:proposition:2:1:item:e:01}{(e)}. If \(G\) is smooth
over \(S\), it is locally of finite presentation over \(S\), hence of
finite presentation over \(S\), since it is affine over \(S\); thus \(M\)
is of finite type. We may therefore assume from the outset that \(M\) is
of finite type, and hence that \(G\) is of finite presentation over
\(S\). Then\footnote{Editor's note: this uses the fact that \(G\) is flat
over \(S\), by \SGARefLink{sga3:viii:proposition:2:1:item:a:01}{(a)}.}
\(G\) is smooth over \(S\) if and only if its geometric fibers are
smooth, reducing us to the case in which \(S\) is the spectrum of an
algebraically closed field \(k\).

Write \(M=T\times L\), where \(T\) is the torsion subgroup and \(L\) is
free, with \(L\simeq\mathbb Z^r\). Then
\[
  D(M)=D(T)\times D(L),
  \qquad
  D(L)\simeq\mathbb G_m^r,
\]
and \(D(L)\) is smooth over \(k\). Thus \(G=D(M)\) is smooth over \(k\)
if and only if \(D(T)\) is smooth. Since \(D(T)\) is finite over \(k\) of
degree equal to the order \(n\) of \(T\), this means that \(D(T)(k)\) has
\(n\) elements. Now \(T\) is isomorphic to a sum of groups
\(\mathbb Z/n_i\mathbb Z\), where \(n=\prod_i n_i\); consequently
\[
  D(T)=\prod_iD(\mathbb Z/n_i\mathbb Z)
      =\prod_i\mu_{n_i},
\]
where \(\mu_{n_i}\) is the group scheme of \(n_i\)-th roots of unity.
Hence
\[
  \operatorname{card}\bigl(D(T)(k)\bigr)
  =\prod_i\operatorname{card}\bigl(\mu_{n_i}(k)\bigr).
\]
Here \(\operatorname{card}(\mu_{n_i}(k))\), the number of \(n_i\)-th
roots of unity in \(k\), is at most \(n_i\), and equality holds if and
only if \(n_i\) is prime to the characteristic \(p\) of \(k\).
\emph{[The source prints \(\mu_i(k)\) in both occurrences.]}
It follows
that
\[
  \operatorname{card}\bigl(D(T)(k)\bigr)=n,
  \qquad n=\prod_i n_i,
\]
if and only if every \(n_i\) is prime to \(p\), that is, if and only if
\(n\) is prime to \(p\).
\hfill\(\square\)

\SGARefTarget{sga3:viii:section:3}\section*{3. Exactness Properties of the Functor
\texorpdfstring{\(D_S\)}{D\_S}}
\addcontentsline{toc}{section}{3. Exactness Properties of the Functor
\texorpdfstring{\(D_S\)}{D\_S}}

\medskip
\noindent\SGARefTarget{sga3:viii:theorem:3:1}\textbf{Theorem 3.1.}
\textit{Let \(S\) be a prescheme, and let}
\SGARefTarget{sga3:viii:theorem:3:1:diagram:01}
\[
  0\longrightarrow M'
   \xrightarrow{u}M
   \xrightarrow{v}M''
   \longrightarrow0
\]
\textit{be an exact sequence of ordinary commutative groups. Consider the
sequence of transposed homomorphisms}
\SGARefTarget{sga3:viii:theorem:3:1:diagram:02}
\[
  0\longrightarrow D_S(M'')
   \xrightarrow{v^t}D_S(M)
   \xrightarrow{u^t}D_S(M')
   \longrightarrow0.
\]
\begin{enumerate}
  \item[(i)]\SGARefTarget{sga3:viii:theorem:3:1:item:i:01}
  \textit{The homomorphism \(v^t\) induces an isomorphism from
  \(D_S(M'')\) onto the kernel of \(u^t\), and \(u^t\) is faithfully flat
  and quasi-compact.}

% Source page: Exp8 standalone p. 9; combined-reader p. 625.
  \item[(ii)]\SGARefTarget{sga3:viii:theorem:3:1:item:ii:01}%
  \footnote{Editor's note: what follows has been added, and the proof has
  been expanded accordingly.}
  \textit{\(D_S(M')\) represents the quotient sheaf
  \(D_S(M)/D_S(M'')\) for the \((\mathrm{fpqc})\) topology.}
\end{enumerate}

Let \(\mathcal M\) denote the family of faithfully flat quasi-compact
morphisms. First,
\SGARefLink{sga3:viii:theorem:3:1:item:ii:01}{(ii)} follows from
\SGARefLink{sga3:viii:theorem:3:1:item:i:01}{(i)} (cf. IV,
\SGARefLink{sga3:iv:corollary:4:6:5:1}{4.6.5.1}). Indeed, the equivalence
relation on \(D_S(M)\) defined by \(u^t\) is the same as that defined by
the subgroup
\[
  \operatorname{Ker}(u^t)=D_S(M'').
\]
Since \(u^t\in\mathcal M\), this equivalence relation is
\(\mathcal M\)-effective (cf. IV,
\SGARefLink{sga3:iv:scholium:3:3:2:1}{3.3.2.1}); therefore \(D_S(M')\)
represents the quotient sheaf for the \((\mathrm{fpqc})\) topology (cf.
IV, \SGARefLink{sga3:iv:lemma:4:6:5}{4.6.5}).

The first assertion of
\SGARefLink{sga3:viii:theorem:3:1:item:i:01}{(i)} is an immediate
consequence of the definition of the functors \(D_S(-)\). More generally,
every exact sequence
\SGARefTarget{sga3:viii:theorem:3:1:diagram:03}
\[
  M'\longrightarrow M\longrightarrow M''\longrightarrow0
\]
with no zero on the left gives rise to an exact transposed sequence
\SGARefTarget{sga3:viii:theorem:3:1:diagram:04}
\[
  0\longrightarrow D_S(M'')
   \longrightarrow D_S(M)
   \longrightarrow D_S(M').
\]
(This remains valid in the more general setting considered at the
beginning of \(\SGARefLink{sga3:viii:section:1}{\S1}\).) On the other hand, because \(D_S(M)\) and
\(D_S(M')\) are affine over \(S\), the morphism \(u^t\) is necessarily
affine, and a fortiori quasi-compact, whatever the homomorphism
\(u:M'\to M\). The second assertion of
\SGARefLink{sga3:viii:theorem:3:1:item:i:01}{(i)} will therefore follow
from \SGARefLink{sga3:viii:corollary:3:2:item:a:01}{part (a)} of \SGARefLink{sga3:viii:corollary:3:2}{the
following corollary}.

\medskip
\noindent\SGARefTarget{sga3:viii:corollary:3:2}\textbf{Corollary 3.2.}
\textit{Let \(S\) be a nonempty prescheme, let \(u:M'\to M\) be a
homomorphism of ordinary commutative groups, and let
\(u^t:G\to G'\) be the transposed homomorphism. Then:}
\begin{enumerate}
  \item[(a)]\SGARefTarget{sga3:viii:corollary:3:2:item:a:01}
  \textit{\(u\) is a monomorphism if and only if \(u^t\) is faithfully
  flat;}

  \item[(b)]\SGARefTarget{sga3:viii:corollary:3:2:item:b:01}
  \textit{\(u\) is an epimorphism if and only if \(u^t\) is a
  monomorphism; in that case \(u^t\) is in fact a closed immersion.}
\end{enumerate}

To prove \SGARefLink{sga3:viii:corollary:3:2:item:a:01}{(a)}, observe
that if \(u\) is a monomorphism, then \(\mathcal O_S(M)\), as a module
over \(\mathcal O_S(M')\), has a nonempty basis: namely, the system of
sections defined by any system of representatives of \(M\) modulo
\(M'\). It is therefore faithfully flat. Conversely, if
\(\mathcal O_S(M)\) is faithfully flat over \(\mathcal O_S(M')\), then
\[
  u^t:\mathcal O_S(M')\longrightarrow\mathcal O_S(M)
\]
is injective, which, when \(S\ne\varnothing\), implies that
\(u:M'\to M\) is injective.

To prove \SGARefLink{sga3:viii:corollary:3:2:item:b:01}{(b)}, observe
that if \(u\) is an epimorphism, then
\(\mathcal O_S(M')\to\mathcal O_S(M)\) is surjective, so \(u^t\) is a
closed immersion and, a fortiori, a monomorphism. Conversely, if \(u^t\)
is a monomorphism, then \(\operatorname{Ker}u^t\) is the unit group.
Putting \(M''=\operatorname{Coker}u\), we have seen that
\[
  \operatorname{Ker}u^t\simeq D_S(M'').
\]
Thus \SGARefLink{sga3:viii:proposition:2:1:item:d:01}{2.1(d)} gives
\(M''=0\), and \(u\) is an epimorphism.

The usual argument now gives \SGARefLink{sga3:viii:corollary:3:3}{the following consequence} of
\SGARefLink{sga3:viii:theorem:3:1}{3.1}.

\medskip
\noindent\SGARefTarget{sga3:viii:corollary:3:3}\textbf{Corollary 3.3.}
\textit{Let}
\SGARefTarget{sga3:viii:corollary:3:3:diagram:00}
\[
  M'\xrightarrow{u}M\xrightarrow{v}M''
\]
\textit{be an exact sequence of ordinary commutative groups, and consider
the transposed sequence}
\SGARefTarget{sga3:viii:corollary:3:3:diagram:01}
\[
  G''\xrightarrow{v^t}G\xrightarrow{u^t}G'.
\]
\textit{Then \(v^t\) induces a faithfully flat quasi-compact morphism
from \(G''\) to \(\operatorname{Ker}u^t\), and this kernel is a
diagonalizable group isomorphic to}
\[
  D_S(v(M))=D_S(\operatorname{Coker}u).
\]

% Source page: Exp8 standalone p. 10; combined-reader p. 626.
\medskip
\noindent\SGARefTarget{sga3:viii:corollary:3:4}\textbf{Corollary 3.4.}
\textit{Let \(S\) be a prescheme, let \(u:G\to H\) be a homomorphism of
\(S\)-group preschemes that are locally diagonalizable, with \(H\) of
finite type over \(S\), and put \(G'=\operatorname{Ker}u\). Then:}
\begin{enumerate}
  \item[(a)]\SGARefTarget{sga3:viii:corollary:3:4:item:a:01}
  \textit{\(G'\) is locally diagonalizable, and is of finite type over
  \(S\) if \(G\) is;}

  \item[(b)]\SGARefTarget{sga3:viii:corollary:3:4:item:b:01}
  \textit{the quotient \(G/G'\) ``exists'': more precisely, the
  equivalence relation on \(G\) defined by \(G'\) is
  \(\mathcal M\)-effective, where \(\mathcal M\) is the family of
  faithfully flat quasi-compact morphisms (cf. IV,
  \SGARefLink{sga3:iv:subsection:3:4}{3.4}). Moreover, \(G/G'\) is
  locally diagonalizable and of finite type over \(S\);}

  \item[(c)]\SGARefTarget{sga3:viii:corollary:3:4:item:c:01}
  \textit{the homomorphism \(u:G\to H\) factors uniquely as}
  \SGARefTarget{sga3:viii:corollary:3:4:diagram:01}
  \[
    G\xrightarrow{v}G/G'\xrightarrow{w}H,
  \]
  \textit{where \(v\) is the canonical homomorphism, and is therefore
  faithfully flat and quasi-compact. Moreover, \(w\) is a closed
  immersion, and a fortiori a monomorphism.}
\end{enumerate}
\textit{Finally, the quotient}
\[
  H'=H/\operatorname{Im}w
    =\operatorname{Coker}w
    =\operatorname{Coker}u
\]
\textit{exists: more precisely, the equivalence relation on \(H\)
defined by \(G/G'\) is \(\mathcal M\)-effective, and \(H'\) is of finite
type over \(S\).}

The first assertion of
\SGARefLink{sga3:viii:corollary:3:4:item:c:01}{(c)} follows from
\SGARefLink{sga3:viii:corollary:3:4:item:b:01}{(b)}, by the definition
of the quotient \(G/G'\) (cf. IV,
\SGARefLink{sga3:iv:lemma:3:2:3}{3.2.3}).
\footnote{Editor's note: the original has been expanded in what follows.}
Let us show that the \((\mathrm{fpqc})\) quotient sheaf
\(\widetilde G/\widetilde{G'}\) is representable. This can be checked
locally on \(S\), as can all the other assertions; we may therefore
suppose that \(G\) and \(H\) are diagonalizable, of the form \(D_S(M)\)
and \(D_S(N)\).

Since \(H\) is of finite type over \(S\), the group \(N\) is of finite
type by \SGARefLink{sga3:viii:proposition:2:1:item:b:01}{2.1(b)}.
Therefore, by \SGARefLink{sga3:viii:corollary:1:5}{1.5}, \(u\) is defined by a homomorphism \(u':N\to M\).
It follows from \SGARefLink{sga3:viii:theorem:3:1}{3.1} and
\SGARefLink{sga3:viii:corollary:3:2}{3.2} that
\[
  G'\simeq D_S(\operatorname{Coker}u')
\]
and that \(\widetilde G/\widetilde{G'}\) is represented by
\(D_S(\operatorname{Im}u')\). Moreover, on considering the exact
sequence
\SGARefTarget{sga3:viii:corollary:3:4:diagram:02}
\[
  0\longrightarrow\operatorname{Ker}u'
   \longrightarrow N
   \xrightarrow{w'}\operatorname{Im}u'
   \longrightarrow0,
\]
we find that \(w\) is a closed immersion and that
\[
  H'=H/\operatorname{Im}w=D_S(\operatorname{Ker}u').
\]
The latter is of finite type over \(S\), since \(N\), and hence
\(\operatorname{Ker}u'\), is of finite type.

\medskip
\noindent\SGARefTarget{sga3:viii:remarks:3:4}\textbf{Remarks.}
The existence result for quotients in
\SGARefLink{sga3:viii:corollary:3:4}{3.4} will be substantially
generalized in \(\SGARefLink{sga3:viii:section:5}{\S5}\).

On the other hand, note that in \SGARefLink{sga3:viii:section:3}{the present section} and \SGARefLink{sga3:viii:section:2}{the preceding
one}, the hypothesis \(S\ne\varnothing\) was used only to ensure the
validity of certain converses, allowing properties of the corresponding
ordinary groups to be deduced from certain hypotheses on diagonalizable
\(S\)-groups. The results in the ``forward'' direction hold without any
restriction on \(S\), and the proofs given here apply in general.

\medskip
\noindent\SGARefTarget{sga3:viii:corollary:3:5}\textbf{Corollary 3.5.}
\textit{Let \(G\) be a diagonalizable group prescheme over \(S\), and let
\(n\ne0\) be an integer. Then the subgroup \({}_nG\) of \(G\), the
kernel of the homomorphism}
\[
  n\mathbin{\cdot}\operatorname{id}_G:G\longrightarrow G,
\]
\textit{is integral over \(S\), and is finite over \(S\) if \(G\) is of
finite type over \(S\).}

% Source page: Exp8 standalone p. 11; combined-reader p. 627.
Indeed, if \(G=D_S(M)\), then
\[
  {}_nG=D_S(M/nM)
\]
by \SGARefLink{sga3:viii:theorem:3:1}{3.1}, and the result follows from
\SGARefLink{sga3:viii:proposition:2:1:item:b:01}{2.1(b)},
\SGARefLink{sga3:viii:proposition:2:1:item:c:01}{(c)}, and
\SGARefLink{sga3:viii:proposition:2:1:item:cprime:01}{\((c')\)}.

\SGARefTarget{sga3:viii:section:4}\section*{4. Torsors under a Diagonalizable Group}
\addcontentsline{toc}{section}{4. Torsors under a Diagonalizable Group}

\SGARefTarget{sga3:viii:section:4:problem:01}Let \(S\) be a prescheme and let \(G=D_S(M)\) be a diagonalizable group
over \(S\). We shall determine the \(G\)-torsors (or principal
homogeneous \(G\)-bundles) over \(S\) for the faithfully flat
quasi-compact topology; cf. Expos\'e IV,
\SGARefLink{sga3:iv:subsection:5:1}{5.1}. Recall that a prescheme \(P\)
over \(S\), with \(S\)-operator group \(G\), is called a \emph{torsor} or
a \emph{principal homogeneous bundle} if every point of \(S\) has an
open neighborhood \(U\) and a faithfully flat quasi-compact morphism
\(S'\to U\) such that
\[
  P'=P\times_SS'
\]
is isomorphic, as an object with operators, to
\[
  G'=G\times_SS',
\]
where \(G'\) acts on itself by right translations. Since \(G\) is affine
over \(S\), SGA 1, VIII 5.6 implies that \(P\) is necessarily affine over
\(S\). Also, because \(G\) itself is faithfully flat and quasi-compact
over \(S\), \(P\) is principal homogeneous under \(G\) if and only if it
is formally principal homogeneous and is, in addition, faithfully flat
and quasi-compact over \(S\) (cf. IV,
\SGARefLink{sga3:iv:proposition:5:1:6}{5.1.6}).

Recall also (\SGARefLink{sga3:i:proposition:4-7-3}{Expos\'e I, 4.7.3}) that giving an \(S\)-prescheme \(P\), affine
over \(S\), with \(S\)-operator group \(G=D_S(M)\), amounts to giving a
commutative quasi-coherent algebra graded by \(M\) over \(S\). In other
words, it amounts to giving a quasi-coherent algebra \(\mathscr A\) over
\(S\), together with a direct-sum decomposition as a module
\[
  \mathscr A=\coprod_{m\in M}\mathscr A_m,
\]
such that
\[
  \mathscr A_m\mathbin{\cdot}\mathscr A_{m'}
  \subseteq\mathscr A_{m+m'}
  \qquad\text{for }m,m'\in M.
\]
With this understood, the answer to \SGARefLink{sga3:viii:section:4:problem:01}{the preceding problem} is as follows.

\medskip
\noindent\SGARefTarget{sga3:viii:proposition:4:1}\textbf{Proposition 4.1.}
\textit{Let \(P\) be the prescheme with \(S\)-operator group
\(G=D_S(M)\) defined by the \(M\)-graded algebra \(\mathscr A\). Then
\(P\) is a principal homogeneous bundle under \(G\) if and only if
\(\mathscr A\) satisfies the following conditions:}
\begin{enumerate}
  \item[(a)]\SGARefTarget{sga3:viii:proposition:4:1:item:a:01}
  \textit{for every \(m\in M\), \(\mathscr A_m\) is an invertible module
  over \(S\);}

  \item[(b)]\SGARefTarget{sga3:viii:proposition:4:1:item:b:01}
  \textit{for \(m,m'\in M\), the homomorphism}
  \[
    \mathscr A_m\otimes_{\mathcal O_S}\mathscr A_{m'}
    \longrightarrow\mathscr A_{m+m'}
  \]
  \textit{induced by multiplication in \(\mathscr A\) is an
  isomorphism.}
\end{enumerate}

The necessity of \SGARefLink{sga3:viii:proposition:4:1}{these conditions} follows immediately by descent, since
they hold when \(P\) is the trivial principal homogeneous bundle, that
is, when \(\mathscr A=\mathcal O_S(M)\). For sufficiency, note that
\SGARefLink{sga3:viii:proposition:4:1:item:a:01}{(a)} already implies
that \(P\) is faithfully flat over \(S\); it is in any event
quasi-compact over \(S\), since it is affine over \(S\). It therefore
remains to check that \(P\) is formally principal homogeneous under
\(G\), that is, that the familiar morphism
\SGARefTarget{sga3:viii:proposition:4:1:diagram:01}
\[
  P\times_SG\longrightarrow P\times_SP
\]
is an isomorphism. On affine algebras \SGARefLink{sga3:viii:proposition:4:1:diagram:01}{this morphism} is given explicitly
by the homomorphism
\SGARefTarget{sga3:viii:proposition:4:1:diagram:02}
\[
  \mathscr A\otimes\mathscr A
  \longrightarrow
  \mathscr A(M)=\mathscr A\otimes\mathcal O_S(M).
\]

% Source page: Exp8 standalone p. 12; combined-reader p. 628.
In bidegree \((m,n)\), where \(m,n\in M\), it is given by
\[
  x_m\otimes y_n\longmapsto x_my_n\otimes e_n.
\]
At the level of degrees, \SGARefLink{sga3:viii:proposition:4:1:diagram:02}{this homomorphism} is compatible with the
homomorphism \(M\times M\to M\times M\) given by
\[
  (m,n)\longmapsto(m+n,n),
\]
which is an isomorphism. This shows that
\SGARefLink{sga3:viii:proposition:4:1:item:b:01}{(b)} expresses precisely,
and independently of \SGARefLink{sga3:viii:proposition:4:1:item:a:01}{(a)},
that \(P\) is formally principal homogeneous, and proves
\SGARefLink{sga3:viii:proposition:4:1}{4.1}.

Faithfully flat descent also gives:

\medskip
\noindent\SGARefTarget{sga3:viii:corollary:4:2}\textbf{Corollary 4.2.}
\textit{The conditions of \SGARefLink{sga3:viii:proposition:4:1}{4.1}
imply that the homomorphism}
\[
  \mathcal O_S\longrightarrow\mathscr A_0
\]
\textit{is an isomorphism.}

For example, if \(M=\mathbb Z\), then, under the conditions of
\SGARefLink{sga3:viii:proposition:4:1}{4.1}, the algebra \(\mathscr A\)
is essentially determined by \(\mathscr A_1=\mathscr L\), namely
\[
  \mathscr A\simeq\coprod_{n\in\mathbb Z}\mathscr L^{\otimes n}
\]
as graded algebras. Thus we recover the familiar result:

\medskip
\noindent\SGARefTarget{sga3:viii:corollary:4:3}\textbf{Corollary 4.3.}
\textit{The category of principal homogeneous bundles \(P\) over \(S\)
under \(\mathbb G_{m,S}\) is equivalent to the category of invertible
modules \(\mathscr L\) over \(S\), where the morphisms in both categories
are the isomorphisms of the structures in question. Quasi-inverse
functors are obtained by assigning to \(P\) the degree-\(1\) component of
its \(\mathbb Z\)-graded affine algebra, and by assigning to
\(\mathscr L\) the spectrum of the \(\mathbb Z\)-graded algebra}
\[
  \coprod_{n\in\mathbb Z}\mathscr L^{\otimes n}.
\]

In particular:

\medskip
\noindent\SGARefTarget{sga3:viii:corollary:4:4}\textbf{Corollary 4.4.}
\textit{The group of isomorphism classes of principal homogeneous
bundles over \(S\) under \(\mathbb G_{m,S}\) is isomorphic to the group
\(\operatorname{Pic}(S)\) of isomorphism classes of invertible modules
over \(S\), that is, to}
\[
  H^1(S,\mathcal O_S^\times).
\]

Since \(\mathbb G_{m,S}\) is the scheme of automorphisms of the module
\(\mathcal O_S\), Corollary \SGARefLink{sga3:viii:corollary:4:4}{4.4} is
equivalent to \SGARefLink{sga3:viii:corollary:4:5}{the following statement}, one of the forms of Hilbert's
``\SGARefLink{sga3:viii:corollary:4:5}{Theorem 90}.''

\medskip
\noindent\SGARefTarget{sga3:viii:corollary:4:5}\textbf{Corollary 4.5.}
\textit{Every principal homogeneous bundle over \(S\) under
\(\mathbb G_m\) is locally trivial for the Zariski topology.}

\medskip
\noindent\SGARefTarget{sga3:viii:remark:4:5:1}\textbf{Remark 4.5.1.}
\SGARefLink{sga3:viii:corollary:4:5}{The preceding statement} is no longer true in general for a group such as
\(\mu_n\), or for a ``twisted form'' of \(\mathbb G_m\). For example,
the unique twisted form of \(\mathbb G_m\) over the field
\(\mathbb R\) of real numbers has a first cohomology group equal to
\(\mathbb Z/2\mathbb Z\).

% Source page: Exp8 standalone p. 13; combined-reader p. 629.
\footnote{Editor's note: what follows has been added.}
Indeed, let \(\mathbb S^1\) be the kernel of the norm morphism
\[
  N:\prod_{\mathbb C/\mathbb R}\mathbb G_{m,\mathbb C}
  \longrightarrow\mathbb G_{m,\mathbb R};
\]
it is a \(\mathbb C/\mathbb R\)-twisted form of
\(\mathbb G_{m,\mathbb R}\). The equation \(N(z)=-1\) in
\(\prod_{\mathbb C/\mathbb R}(\mathbb G_{m,\mathbb C})\) defines an
\(\mathbb S^1\)-torsor \(X\) over \(\operatorname{Spec}(\mathbb R)\),
locally trivial for the \(\acute{e}\)tale topology, but not trivial, since
\(X(\mathbb R)=\varnothing\). Let us show that
\[
  H^1_{\mathrm{\acute et}}(\mathbb R,\mathbb S^1)
  \simeq\mathbb Z/2\mathbb Z.
\]
There is an exact sequence of smooth commutative algebraic
\(\mathbb R\)-groups
\SGARefTarget{sga3:viii:remark:4:5:1:diagram:01}
\[
  1\longrightarrow\mathbb S^1
  \longrightarrow\prod_{\mathbb C/\mathbb R}\mathbb G_{m,\mathbb C}
  \longrightarrow\mathbb G_{m,\mathbb R}
  \longrightarrow1,
\]
which gives rise to a long exact sequence in \(\acute{e}\)tale cohomology
(or in Galois cohomology):
\SGARefTarget{sga3:viii:remark:4:5:1:diagram:02}
\[
\begin{aligned}
  0\longrightarrow\mathbb S^1(\mathbb R)
  &\longrightarrow\mathbb C^\times
   \xrightarrow{N}\mathbb R^\times
   \longrightarrow H^1_{\mathrm{\acute et}}(\mathbb R,\mathbb S^1)\\
  &\longrightarrow
   H^1_{\mathrm{\acute et}}\!
   \left(\mathbb R,
     \prod_{\mathbb C/\mathbb R}\mathbb G_{m,\mathbb C}\right)
   \longrightarrow\cdots .
\end{aligned}
\]
Now (see, for example, XXIV, \SGARefLink{sga3:xxiv:proposition:8-4}{8.4}),
\[
  H^1_{\mathrm{\acute et}}\!
  \left(\mathbb R,
    \prod_{\mathbb C/\mathbb R}\mathbb G_{m,\mathbb C}\right)
  \simeq
  H^1_{\mathrm{\acute et}}(\mathbb C,\mathbb G_{m,\mathbb C}),
\]
and the latter group is zero by
\SGARefLink{sga3:viii:corollary:4:5}{4.5} (or, in this case, because
\(\mathbb C\) is algebraically closed). We therefore obtain an
isomorphism
\[
  H^1_{\mathrm{\acute et}}(\mathbb R,\mathbb S^1)
  \simeq\mathbb R^\times/N(\mathbb C^\times)
  \simeq\{\pm1\}.
\]

We shall need \SGARefLink{sga3:viii:proposition:4:6}{the following result} in \SGARefLink{sga3:viii:section:5}{the next section}.

\medskip
\noindent\SGARefTarget{sga3:viii:proposition:4:6}\textbf{Proposition 4.6.}
\textit{Under the conditions of
\SGARefLink{sga3:viii:proposition:4:1}{4.1}, conditions
\SGARefLink{sga3:viii:proposition:4:1:item:a:01}{(a)} and
\SGARefLink{sga3:viii:proposition:4:1:item:b:01}{(b)} are equivalent to
the following conditions:}
\begin{enumerate}
  \item[\((a')\)]\SGARefTarget{sga3:viii:proposition:4:6:item:aprime:01}
  \textit{the homomorphism
  \(\mathcal O_S\to\mathscr A_0\) is an isomorphism;}

  \item[\((b')\)]\SGARefTarget{sga3:viii:proposition:4:6:item:bprime:01}
  \textit{for every \(m\in M\) (it is enough to require this for the
  members of a system of generators of \(M\)), one has}
  \[
    \mathscr A_m\mathbin{\cdot}\mathscr A_{-m}=\mathscr A_0.
  \]
\end{enumerate}

The necessity is clear in view of
\SGARefLink{sga3:viii:corollary:4:2}{4.2}.\footnote{Editor's note:
``\SGARefLink{sga3:viii:proposition:4:1}{4.1}'' has been corrected to ``\SGARefLink{sga3:viii:corollary:4:2}{4.2}.''} We are thus reduced to proving
\SGARefLink{sga3:viii:corollary:4:7}{the following corollary}.

\medskip
\noindent\SGARefTarget{sga3:viii:corollary:4:7}\textbf{Corollary 4.7.}
\textit{Let}
\[
  A=\coprod_{n\in\mathbb Z}A_n
\]
\textit{be a \(\mathbb Z\)-graded ring such that}
\[
  A_1\mathbin{\cdot}A_{-1}=A_0.
\]
\textit{Then the \(A_n\) are invertible \(A_0\)-modules, and, for
\(n,n'\in\mathbb Z\), the homomorphism}
\[
  A_n\otimes_{A_0}A_{n'}\longrightarrow A_{n+n'}
\]
\textit{induced by multiplication in \(A\) is an isomorphism.}

By hypothesis there exist \(f_i\in A_1\) and \(g_i\in A_{-1}\) such that
\SGARefTarget{sga3:viii:corollary:4:7:eq:star:01}
\[
  \makebox[\displaywidth][s]{%
    \(\text{(*)}\)\hfill
    \(\displaystyle\sum_i f_ig_i=1.\)\hfill
    \(\phantom{\text{(*)}}\)%
  }
\]
Since the assertion to be proved is local on
\(\operatorname{Spec}(A_0)\),\footnote{Editor's note: the following
passage has been expanded.} and since
\(\SGARefLink{sga3:viii:corollary:4:7:eq:star:01}{(*)}\) shows that
\(\operatorname{Spec}(A_0)\) is covered by the affine open subsets
\(D(f_ig_i)\), we are reduced to the case in which there is an element
\(f\in A_1\) that is invertible in \(A\). Then, for every
\(n\in\mathbb Z\), \(f^n\) is an invertible element of \(A_n\), and
therefore defines an isomorphism
\[
  A_0\xrightarrow{\ \sim\ }A_n,
  \qquad h\longmapsto f^nh.
\]
Moreover,

% END INCLUDED SGA3_English_current_buildable_source_20260823/source/inputs/VIII/tex/components/delegated_sections2_4/01_expose_VIII_s2_s4_pp007_013.tex

% BEGIN INCLUDED SGA3_English_current_buildable_source_20260823/source/inputs/VIII/tex/components/03_expose_VIII_section4_continuation_p14.tex
% SGA 3, Expose VIII, English reconstruction.
% Scope: continuation of the proof of Corollary 4.7 and completion of
% Proposition 4.6, at the top of authority standalone page 14
% (combined-reader page 630). Stops before the Section 5 heading.
% Authority: Polo--Gille Exp8-8nov09.pdf.
% Authority standalone p. 14; combined-reader p. 630.
this shows that one obtains an isomorphism
\[
  A_0[t,t^{-1}]\longrightarrow A
\]
of graded \(A_0\)-algebras by sending \(t\) to \(f\), which completes
the proof of \SGARefLink{sga3:viii:corollary:4:7}{4.7}.

Under the hypotheses of
\SGARefLink{sga3:viii:proposition:4:6}{4.6},
\SGARefLink{sga3:viii:corollary:4:7}{4.7} already shows that the
\(\mathscr A_m\), for \(m\in M\), are invertible. To prove condition
\SGARefLink{sga3:viii:proposition:4:1:item:b:01}{4.1(b)}, we may
therefore suppose that \(\mathscr A_m\) and \(\mathscr A_{m'}\) have
bases \(f_m\) and \(f_{m'}\), with inverses
\[
  f_m^{-1}\in\Gamma(\mathscr A_{-m})
  \qquad\text{and}\qquad
  f_{m'}^{-1}\in\Gamma(\mathscr A_{-m'}).
\]
Multiplication by
\[
  f_m^{-1}f_{m'}^{-1}
  \in\Gamma(\mathscr A_{-m-m'})
\]
then defines a homomorphism
\[
  \mathscr A_{m+m'}\longrightarrow
  \mathscr A_0\simeq\mathcal O_S
\]
that sends the image of \(f_m\otimes f_{m'}\) to the unit section of
\(\mathcal O_S\). In the diagram
\SGARefTarget{sga3:viii:proposition:4:6:diagram:01}
\[
\begin{tikzcd}[column sep=large]
  \mathscr A_m\otimes\mathscr A_{m'}
    \arrow[r,"u"]
    \arrow[rr,bend left=32,"w"]
  & \mathscr A_{m+m'}
    \arrow[r,"v"]
  & \mathscr A_0\simeq\mathcal O_S ,
\end{tikzcd}
\]
the maps \(w\) and \(v\) are therefore epimorphisms of invertible
sheaves, hence isomorphisms; consequently \(u\) is an isomorphism.
\hfill\(\square\)

% END INCLUDED SGA3_English_current_buildable_source_20260823/source/inputs/VIII/tex/components/03_expose_VIII_section4_continuation_p14.tex

% BEGIN INCLUDED SGA3_English_current_buildable_source_20260823/source/inputs/VIII/tex/components/delegated_sections5_6/05_expose_VIII_section5_en.tex
% SGA 3, Expose VIII, English reconstruction.
% Scope: all of Section 5, stopping immediately before the Section 6 heading.
% Authority: Polo--Gille Exp8-8nov09.pdf, standalone pages 14--17
% (combined-reader pages 630--633).
% The Jacob Reinhold Markdown lineage was used only as a phrasing comparison.
% The Polo--Gille PDF controls prose, formulas, numbering, references, and notes.

\setcounter{footnote}{16}

% Authority standalone p. 14; original-edition pagination marker 15.
\SGARefTarget{sga3:viii:section:5}\section*{5. Quotient of an Affine Scheme by a Diagonalizable Group Acting Freely}
\addcontentsline{toc}{section}{5. Quotient of an Affine Scheme by a Diagonalizable Group Acting Freely}

\noindent\footnote{Editor's note: We have added the following sentence.}
We denote by \(\mathcal M\) the set of faithfully flat quasi-compact
morphisms, and recall that torsors are understood with respect to the
\((\mathrm{fpqc})\) topology.

\medskip
\noindent\SGARefTarget{sga3:viii:theorem:5:1}\textbf{Theorem 5.1.}
\begin{itshape}
Let \(S\) be a prescheme, \(M\) an ordinary commutative group,
\(G=D_S(M)\) the diagonalizable group over \(S\) that it defines, and
\(P\) an \(S\)-prescheme affine over \(S\), on which \(G\) acts freely on
the right.

Then the equivalence relation defined by \(G\) on \(P\) is
\(\mathcal M\)-effective (cf. IV,
\SGARefLink{sga3:iv:subsection:3:4}{3.4}); that is, the quotient
\(X=P/G\) exists and \(P\) is a torsor over \(X\) under
\(G_X=D_X(M)\). Moreover, \(P/G\) is affine over \(S\). More precisely,
if \(P\) is defined by the \(M\)-graded algebra \(\mathcal A\), then
\(P/G\) is isomorphic to \(\operatorname{Spec}(\mathcal A_0)\), where
\(\mathcal A_0=\mathcal A^G\) is the degree-zero component of
\(\mathcal A\).
\end{itshape}

\smallskip
\noindent\textit{Proof.}
Put \(X=\operatorname{Spec}(\mathcal A_0)\). The inclusion
\(\mathcal A_0\to\mathcal A\) gives a natural morphism \(P\to X\), which
is invariant under the action of \(G\). Thus \(P\) becomes an
\(X\)-prescheme, affine over \(X\), with operator group
\(G_X=D_X(M)\); the hypothesis that \(G\) acts freely on \(P/S\) implies
that \(G_X\) acts freely on \(P/X\). It remains only to show that \(P\)
is a principal homogeneous bundle under \(G_X\), using the fact that
\(\mathcal B_0=\mathcal O_X\), where \(\mathcal B\) is the \(M\)-graded
algebra over \(X\) defining \(P/X\).

We may therefore suppose that \(X=S\) and that \(S\) is affine. Then
\(P\) is affine, defined by an \(M\)-graded ring \(A\); write \(A_m\) for
its homogeneous part of degree \(m\), so that
\(S=\operatorname{Spec}(A_0)\). In view of
\SGARefLink{sga3:viii:proposition:4:6}{4.6}, it remains to verify
\[
  A_m\cdot A_{-m}=A_0
  \qquad\text{for every \(m\in M\).}
  \tag{x}\label{sga3:viii:theorem:5:1:eq:x:01}
\]
A direct calculation also shows that
\SGARefLink{sga3:viii:theorem:5:1:eq:x:01}{(x)} is equivalent to saying
that
\[
  P\times_SG\longrightarrow P\times_SP
\]
is a closed immersion, rather than merely a monomorphism (as follows
from the hypothesis that \(G\) acts freely); equivalently, the
homomorphism of affine rings

% Authority standalone p. 15; original-edition pagination marker 16.
\[
  \theta:A\otimes_{A_0}A\longrightarrow A(M)
\]
is surjective.\footnote{Editor's note: Indeed, for every \(b\in A\) and
\(a_m\in A_m\), one has
\(\theta(b\otimes a_m)=ba_m\otimes e_m\); hence the surjectivity of
\(\theta\) is equivalent to \SGARefLink{sga3:viii:theorem:5:1:eq:x:01}{(x)}.}
This proves \SGARefLink{sga3:viii:theorem:5:1}{5.1} if one assumes that
the equivalence relation defined by \(G\) on \(P\) is closed. We shall
show, however, that this hypothesis already follows from the fact that
\(G\) acts freely. This is also implicit in
\SGARefLink{sga3:viii:theorem:5:1}{Theorem 5.1}: the scheme
\(G\times_SP\) must be isomorphic to \(P\times_XP\), which is closed in
\(P\times_SP\) because \(X\) is affine, hence separated, over \(S\).

Let \(R=P\times_SG\). The hypothesis that \(G\) acts freely, that is,
that \(R\to P\times_SP\) is a monomorphism, says that the diagonal
morphism
\[
  R\longrightarrow
  R\times_{(P\times_SP)}R=R'
\]
is an isomorphism. We have \(R=\operatorname{Spec}(A(M))\) and
\[
  R'=\operatorname{Spec}\bigl(A(M\times M)/K\bigr),
\]
\footnote{Editor's note: We have denoted by \(K\) the ideal denoted by
\(J\) in the original, in order to distinguish it from the ideals
\(J_p\) of \(A_0\) occurring in
\SGARefLink{sga3:viii:theorem:5:1:eq:xxx:01}{(xxx)}. We have also made
explicit below that the relations
\SGARefLink{sga3:viii:theorem:5:1:eq:xx:01}{(xx)} are equivalent to the
equality \(K=K'\); see what follows.}
where \(K\) is the ideal generated by the elements
\[
  x_m(e_{m,0}-e_{0,m}),
  \qquad m\in M,\quad x_m\in A_m.
\]
Let
\[
  \phi:A(M\times M)\longrightarrow A(M)
\]
be the surjective ring homomorphism defined by
\[
  x e_{m,n}\longmapsto x e_{m+n}
  \qquad (m,n\in M,\ x\in A),
\]
where the \(e_m\), respectively
\(e_{m,n}=e_m\otimes e_n\), are the elements of the canonical basis of
\(A(M)\), respectively \(A(M\times M)\). The diagonal morphism
\(R\to R'\) then corresponds to the homomorphism
\[
  \bar\phi:A(M\times M)/K\longrightarrow A(M)
\]
obtained by passing to the quotient by \(K\). The kernel of \(\phi\) is
the ideal \(K'\) generated by
\[
  d_m=e_{m,0}-e_{0,m}.
\]
\emph{[The source prints \(\ker(\bar\phi)=K'\); after quotienting by
\(K\), the corresponding kernel is \(K'/K\).]}
We have \(K\subseteq K'\), and the hypothesis that \(G\) acts freely on
\(P\), that is, that \(\bar\phi\) is an isomorphism, is equivalent to
\(K'=K\), expressed by the relations
\[
  d_m\in K=\sum_p A(M\times M)\,A_p\,d_p
  \qquad\text{for every \(m\in M\).}
  \tag{xx}\label{sga3:viii:theorem:5:1:eq:xx:01}
\]

Using the natural trigrading of \(A(M\times M)\), and the fact that the
first degree of \(d_m\) is zero, this says that \(d_m\) can be written as
a sum of elements of the form
\[
  f e_{r,s}(e_{p,0}-e_{0,p}),
  \qquad f\in A_{-p}\cdot A_p.
\]
Since the total degree of \(d_m\) is \(m\), it is enough to use terms
satisfying
\[
  r+s+p=m.
\]

% Authority standalone p. 16; original-edition pagination marker 17.
Thus, for every \(m\in M\), there must be an expression
\[
  \left\{
  \begin{aligned}
    d_m=e_{m,0}-e_{0,m}
      &=\sum_{r,s}\lambda_{r,s}
        (e_{m-s,s}-e_{r,m-r}),\\
    \lambda_{r,s}&\in J_p=A_p\cdot A_{-p}\subseteq A_0,
      \qquad p=m-(r+s).
  \end{aligned}
  \right.
  \tag{xxx}\label{sga3:viii:theorem:5:1:eq:xxx:01}
\]
We must deduce \SGARefLink{sga3:viii:theorem:5:1:eq:x:01}{(x)}, that is,
the relations
\[
  J_n=A_0
  \qquad\text{for every \(n\in M\).}
  \tag{xxxx}\label{sga3:viii:theorem:5:1:eq:xxxx:01}
\]
For this it is enough to establish \SGARefLink{sga3:viii:theorem:5:1:eq:xxxx:01}{the same relation} modulo every
maximal ideal of \(A_0\). Since all the hypotheses are preserved under
such a reduction, we may assume that \(A_0\) is a field.

\medskip
\noindent\SGARefTarget{sga3:viii:lemma:5:2}\textbf{Lemma 5.2.}
\textit{Under the preceding conditions, with \(A_0\) a field, if
\(M\ne0\), there exists \(p\in M-\{0\}\) such that \(J_p=A_0\).}

Indeed, otherwise the sum on the right-hand side of
\SGARefLink{sga3:viii:theorem:5:1:eq:xxx:01}{(xxx)} would vanish for
every \(m\in M\), which is absurd.

% Original-edition pagination marker 18.
\medskip
\noindent\SGARefTarget{sga3:viii:corollary:5:3}\textbf{Corollary 5.3.}
\textit{Under the preceding conditions, but without assuming that
\(A_0\) is a field, there are finitely many elements
\(p_i\in M-\{0\}\) such that}
\[
  \sum_iJ_{p_i}=A_0.
\]

Indeed, apply \SGARefLink{sga3:viii:lemma:5:2}{5.2} to the situations
obtained from \(A/A_0\) by reduction modulo the maximal ideals of
\(A_0\).\footnote{Editor's note: It follows from
\SGARefLink{sga3:viii:lemma:5:2}{5.2} that
\(\sum_{p\ne0}J_p=A_0\); hence \(1\) can be written as a finite sum
\(\sum_i x_i\), with \(x_i\in J_{p_i}\).}

\medskip
\noindent\SGARefTarget{sga3:viii:corollary:5:4}\textbf{Corollary 5.4.}
\textit{Suppose again that \(A_0\) is a field. For every subgroup
\(N\) of \(M\) with \(N\ne M\), there exists \(p\in M-N\) such that
\(J_p=A_0\).}

Indeed, let \(M'=M/N\), and consider the \(M'\)-graded ring \(A'\), whose
underlying ring is \(A\) and whose grading is given by
\[
  A'_{m'}=\coprod_{m\in h^{-1}(m')}A_m,
\]
where \(h:M\to M'=M/N\) is the canonical homomorphism. Geometrically,
this construction amounts to considering the subgroup
\(G'=D_S(M')\) of \(G\), and the structure on \(P\) of a scheme with
operator group \(G'\) induced by the action of \(G\). It is then clear
that \(G'\) acts freely on \(P\)
\emph{[the source prints \(P'\)]}; that is, the pair \((M',A')\)
satisfies the hypotheses of \SGARefLink{sga3:viii:corollary:5:3}{5.3}.
Thus
\[
  1=\sum_i f_i g_i,
  \qquad
  f_i\in A'_{m_i'},\quad
  g_i\in A'_{-m_i'},\quad
  m_i'\in M'-\{0\}.
\]
Taking the components of the right-hand side along \(A_0\), and using
the fact that \(A_0\) is a field, immediately gives
\SGARefLink{sga3:viii:corollary:5:4}{5.4}.

Now observe that
\[
  J_pJ_q\subseteq J_{p+q}
  \qquad\text{and}\qquad
  J_p=J_{-p}.
\]

% Authority standalone p. 17; original-edition pagination marker 19.
Thus, if \(N\) denotes the set of \(m\in M\) such that \(J_m=A_0\)
\emph{[the source prints \(J_p=A_0\)]}, then
\(N\) is a subgroup of \(M\). By
\SGARefLink{sga3:viii:corollary:5:4}{5.4}, it equals \(M\). This
completes the proof of \SGARefLink{sga3:viii:theorem:5:1}{Theorem 5.1}.

As noted in the course of the proof,
\SGARefLink{sga3:viii:theorem:5:1}{Theorem 5.1} implies:

\medskip
\noindent\SGARefTarget{sga3:viii:corollary:5:5}\textbf{Corollary 5.5.}
\textit{Under the hypotheses of
\SGARefLink{sga3:viii:theorem:5:1}{5.1}, the graph morphism}
\[
  P\times_SG\longrightarrow P\times_SP
\]
\textit{is a closed immersion.}

It follows immediately that:

\medskip
\noindent\SGARefTarget{sga3:viii:corollary:5:6}\textbf{Corollary 5.6.}
\textit{Let \(\sigma\) be a section of \(P\) over \(S\). Then the
morphism \(g\mapsto\sigma\cdot g\) from \(G\) to \(P\) defined by
\(\sigma\) is a closed immersion.}

\medskip
\noindent\SGARefTarget{sga3:viii:corollary:5:7}\textbf{Corollary 5.7.}
\textit{Let \(G\) and \(H\) be two \(S\)-groups, with \(G\)
diagonalizable and \(H\) affine over \(S\), and let \(u:G\to H\) be a
homomorphism of \(S\)-groups that is a monomorphism. Then \(u\) is a
closed immersion, \(H/G=X\) exists, \(H\) is a principal homogeneous
bundle over \(X\) under \(G_X\), and \(X\) is affine over \(S\).}

\medskip
\noindent\SGARefTarget{sga3:viii:corollary:5:8}\textbf{Corollary 5.8.}
\textit{Under the hypotheses of
\SGARefLink{sga3:viii:theorem:5:1}{5.1}, if \(P\) is of finite type
(respectively, of finite presentation) over \(S\), then the same is
true of \(X=P/G\).}

Indeed, the hypothesis implies that the fibers of \(G_X\) are of finite
type. Hence \(G_X\) is of finite presentation over \(X\), by
\SGARefLink{sga3:viii:proposition:2:1:item:b:01}{2.1(b)}; since \(P\) is a torsor
under \(G_X\), it is therefore of finite presentation over
\(X\).\footnote{Editor's note: Cf. EGA II, 2.7.1(vi).}
It is also faithfully flat over \(X\), so the conclusion follows from
Expos\'e V, Proposition \SGARefLink{sga3:v:proposition:9:1}{9.1}.

% END INCLUDED SGA3_English_current_buildable_source_20260823/source/inputs/VIII/tex/components/delegated_sections5_6/05_expose_VIII_section5_en.tex

% BEGIN INCLUDED SGA3_English_current_buildable_source_20260823/source/inputs/VIII/tex/components/delegated_sections5_6/06_expose_VIII_section6_en.tex
% SGA 3, Expose VIII, English reconstruction.
% Scope: all of Section 6, stopping immediately before the Section 7 heading.
% Authority: Polo--Gille Exp8-8nov09.pdf, standalone pages 17--22
% (combined-reader pages 633--638).
% The Jacob Reinhold Markdown lineage was used only as a phrasing comparison.
% The Polo--Gille PDF controls prose, formulas, numbering, references, and notes.

% Authority standalone p. 17; original-edition pagination marker 20.
\SGARefTarget{sga3:viii:section:6}\section*{6. Essentially Free Morphisms and Representability of Certain Functors of the Form \(\prod_{Y/S}Z/Y\)\textsuperscript{(*)}}
\addcontentsline{toc}{section}{6. Essentially Free Morphisms and Representability of Certain Functors}

\begingroup
\renewcommand{\thefootnote}{(*)}
\footnotetext[0]{\SGARefLink{sga3:viii:section:6}{This section} is independent of the theory of diagonalizable
groups; its natural place would be in Expos\'e VI\textsubscript{B}.}
\endgroup

\noindent\SGARefTarget{sga3:viii:definition:6:1}\textbf{Definition 6.1.}
Let \(f:X\to S\) be a morphism of preschemes. We say that \(f\) is
\emph{essentially free}, or that \(X\) is \emph{essentially free over
\(S\)}, if there are an affine open covering \((S_i)\) of \(S\), for
each \(i\) an \(S_i\)-prescheme \(S_i'\) that is affine and faithfully
flat over \(S_i\), and a covering \((X'_{ij})_j\) of
\(X_i'=X\times_SS_i'\) by affine opens \(X'_{ij}\), such that, for every
\((i,j)\), the ring of \(X'_{ij}\) is a free module over the ring of
\(S_i'\).\footnote{Editor's note: ``Free'' may be replaced by
``projective''; cf. \SGARefLink{sga3:viii:remark:6:8}{6.8} below. This
notion should also be compared with that of an \(S\)-prescheme that is
flat and pure, introduced and developed in \SGARefLink{sga3:viii:bib:rg71}{[RG71]}; see in particular
\emph{loc. cit.}, Part I, 3.3.12, and Part II, 3.1.4.1.}

% Authority standalone p. 18.
\medskip
\noindent\SGARefTarget{sga3:viii:proposition:6:2}\textbf{Proposition 6.2.}
\begin{enumerate}
  \item[(a)]\SGARefTarget{sga3:viii:proposition:6:2:item:a:01}
  \textit{If \(X\) is essentially free over \(S\), then it is flat over
  \(S\). The converse holds if \(S\) is Artinian.}

  \item[(b)]\SGARefTarget{sga3:viii:proposition:6:2:item:b:01}
  \textit{If \(S\) is the spectrum of a field, every \(S\)-prescheme is
  essentially free over \(S\).}

  \item[(c)]\SGARefTarget{sga3:viii:proposition:6:2:item:c:01}
  \textit{If \(X\) is essentially free over \(S\), and if \(S'\to S\)
  is a base-change morphism, then \(X'=X\times_SS'\) is essentially
  free over \(S'\). The converse holds if \(S'\to S\) is faithfully
  flat and quasi-compact.}
\end{enumerate}

The proof is immediate, using for the converse in \SGARefLink{sga3:viii:proposition:6:2:item:a:01}{(a)} the fact that a
flat module over an Artinian local ring is free.\footnote{Editor's note:
Indeed, let \((A,\mathfrak m)\) be an Artinian local ring, \(k\) its
residue field, \(M\) an arbitrary \(A\)-module, and
\((x_i)_{i\in I}\) elements of \(M\) whose images form a basis of
\(M/\mathfrak mM\) over \(k\). Let \(F\) be the free \(A\)-module with
basis \((e_i)_{i\in I}\), and let \(\phi:F\to M\) be the
\(A\)-morphism defined by \(\phi(e_i)=x_i\). Then
\(Q=\operatorname{Coker}\phi\) satisfies \(Q=\mathfrak mQ\), whence
\(Q=0\), since \(\mathfrak m\) is nilpotent. Suppose moreover that
\(M\) is flat over \(A\). Then \(K=\operatorname{Ker}\phi\) satisfies
\(K\otimes_Ak=0\), that is, \(K=\mathfrak mK\), whence \(K=0\).}

\medskip
\noindent\SGARefTarget{sga3:viii:proposition:6:3}\textbf{Proposition 6.3.}
\textit{Let \(H\) be a diagonalizable \(S\)-group prescheme
(more generally, one that becomes diagonalizable after a suitable
faithfully flat quasi-compact extension of every affine open of \(S\);
that is, \(H\) is ``of multiplicative type,'' cf. IX, \SGARefLink{sga3:ix:definition:1:1}{1.1}). Then \(H\)
is essentially free over \(S\).}

Indeed, if \(H\) is diagonalizable, it is affine over \(S\) and is
defined by an algebra that is a free \(\mathcal O_S\)-module.

The introduction of \SGARefLink{sga3:viii:definition:6:1}{Definition 6.1}
is justified by \SGARefLink{sga3:viii:theorem:6:4}{the following theorem}.

% Original-edition pagination marker 21.
\medskip
\noindent\SGARefTarget{sga3:viii:theorem:6:4}\textbf{Theorem 6.4.}
\textit{Let \(S\) be a prescheme, \(Z\) an \(S\)-prescheme essentially
free over \(S\), and \(Y\) a closed subprescheme of \(Z\). Consider the
functor}\footnote{Editor's note: Cf. II.1, where \SGARefLink{sga3:viii:theorem:6:4:display:functor:01}{this functor} is denoted
\(\prod_{Z/S}Y\).}
\SGARefTarget{sga3:viii:theorem:6:4:display:functor:01}
\[
  \begin{aligned}
  F=\prod_{Z/S}Y/Z:\;(\mathbf{Sch}/S)^\circ
    &\longrightarrow\mathbf{Set},\\
  F(S')=\Gamma(Y_{S'}/Z_{S'})
    &=
    \begin{cases}
      \varnothing,& Z_{S'}\ne Y_{S'},\\
      \{\operatorname{id}_{Z_{S'}}\},& Z_{S'}=Y_{S'}.
    \end{cases}
  \end{aligned}
\]
\textit{\SGARefLink{sga3:viii:theorem:6:4:display:functor:01}{This functor} is representable by a closed subprescheme of
\(S\).}\footnote{Editor's note: We have corrected \(Z\) to \(S\).}

\smallskip
\noindent\footnote{Editor's note: We have expanded the original in what
follows; see also \SGARefLink{sga3:viii:corollary:1:7:3}{1.7.3}.}
First note that \(F\) is a sheaf for the \((\mathrm{fpqc})\) topology.
Since \(F(S')\) is either \(\varnothing\) or \(\{\mathrm{pt}\}\) for
every \(S'\), it is enough to check the following assertion: if
\((S_i)\) is an open covering of \(S\) (respectively, if
\(S'\to S\) is faithfully flat and quasi-compact), and every
\(Y_{S_i}\to Z_{S_i}\) is an isomorphism (respectively,
\(Y_{S'}\to Z_{S'}\) is an isomorphism), then \(Y\to Z\) is an
isomorphism. This is clear in the first case and follows from SGA 1,
VIII, 5.4, or EGA IV\textsubscript{2}, 2.7.1, in the second.

Moreover, by SGA 1, VIII, 2.1 and 5.5, faithfully flat quasi-compact
morphisms are morphisms of effective descent for the fibered category
of arrows that are closed immersions. Thus, with the notation of
\SGARefLink{sga3:viii:definition:6:1}{6.1}, we may reduce to the case
\(S=S_i'\).

Let \((Z_j)\) be a covering of \(Z\) by affine opens such that
\(\mathcal O(Z_j)\) is a free module over \(A=\mathcal O(S)\), put
\(Y_j=Y\cap Z_j\), and let
\[
  F_j:(\mathbf{Sch}/S)^\circ\longrightarrow\mathbf{Set}
\]
be the functor defined from \((Z_j,Y_j)\) as \(F\) is defined from
\((Z,Y)\). It is a subfunctor of the terminal functor, and clearly
\[
  F=\bigcap_jF_j.
\]

% Authority standalone p. 19.
It therefore suffices to prove that every \(F_j\) is representable by a
closed subscheme \(T_j\) of \(S\), for then \(F\) is represented by the
closed subscheme \(T=\bigcap_jT_j\). We may consequently suppose that
\(Z\) is affine as well, say \(Z=\operatorname{Spec}(B)\), where \(B\)
is a free \(A\)-module. Let \(J\) be a subset of \(B\) defining the
closed subprescheme \(Y\) of \(Z\), and let \(K\) be the ideal of \(A\)
generated by the \(u_i(J)\subseteq A\), where the \(u_i:B\to A\) are
the coordinate forms for the chosen basis. It is immediate that
\[
  T=\mathcal V(K)=\operatorname{Spec}(A/K)
\]
has the required property. This completes the proof.

\medskip
\noindent\SGARefTarget{sga3:viii:examples:6:5}\textbf{Examples 6.5.}
We give some important examples of functors that reduce to functors
\(\prod_{Z/S}Y/Z\) of the type considered in
\SGARefLink{sga3:viii:theorem:6:4}{6.4}, and for which criteria of
representability will be useful below. Here \(S\) denotes a prescheme,
and \(X,Y,Z,\ldots\) denote preschemes over \(S\).

\smallskip
\noindent\SGARefTarget{sga3:viii:examples:6:5:item:a:01}\textit{(a)}
Suppose given an \(S\)-morphism
\[
  q:X\longrightarrow\underline{\operatorname{Hom}}_S(Y,Z),
  \tag{x}\label{sga3:viii:examples:6:5:eq:x:01}
\]
(``\(X\) acts on \(Y\) with values in \(Z\)''), that is, a morphism
\[
  r:X\times_SY\longrightarrow Z.
  \tag{xx}\label{sga3:viii:examples:6:5:eq:xx:01}
\]
Let \(Z'\) be a subprescheme of \(Z\). It gives a monomorphism
\[
  \underline{\operatorname{Hom}}_S(Y,Z')
  \longrightarrow
  \underline{\operatorname{Hom}}_S(Y,Z),
\]
making the first functor a subfunctor of the second. Let \(X'\) be the
inverse image of this subfunctor under
\SGARefLink{sga3:viii:examples:6:5:eq:x:01}{(x)}. Thus \(X'\) is the
subfunctor of \(X\) for which \(X'(T)\) is the set of \(x\in X(T)\)
such that
\[
  q(x):Y_T\longrightarrow Z_T
\]
factors through \(Z_T'\). This functor has the following description.
Put \(P=X\times_SY\), and let \(P'\) be the inverse image of \(Z'\)
under \(r:P\to Z\). Then there is an evident isomorphism
\[
  X'\simeq\prod_{P/X}P'/P.
  \tag{xxx}\label{sga3:viii:examples:6:5:eq:xxx:01}
\]
It follows that \textit{if \(Y\) is essentially free over \(S\) and
\(Z'\) is closed in \(Z\), then the subfunctor \(X'\) of \(X\) is
representable by a closed subprescheme of \(X\).}

\smallskip
\noindent\SGARefTarget{sga3:viii:examples:6:5:item:b:01}\textit{(b)}
Suppose given two actions of \(X\) on \(Y\) with values in \(Z\), that
is, two morphisms
\[
  q_1,q_2:X\rightrightarrows
  \underline{\operatorname{Hom}}_S(Y,Z),
\]
and put \(X'=\operatorname{Ker}(q_1,q_2)\). This is the subfunctor of
\(X\) for which \(X'(T)\) is the set of \(x\in X(T)\) such that the two
morphisms
\[
  q_1(x),q_2(x):Y_T\rightrightarrows Z_T
\]
are equal. Giving \(q_1,q_2\) is equivalent to giving a morphism
\[
  q:X\longrightarrow
  \underline{\operatorname{Hom}}_S(Y,Z\times_SZ),
\]
or, equivalently, a morphism
\[
  r:X\times_SY\longrightarrow Z\times_SZ.
\]
Put \(U=Z\times_SZ\), and let \(U'\) be the diagonal subprescheme of
\(Z\times_SZ\). Then \(X'\) is the inverse image under \(q\) of the
subfunctor
\[
  \underline{\operatorname{Hom}}_S(Y,U')
  \longrightarrow
  \underline{\operatorname{Hom}}_S(Y,U).
\]
It can therefore be written in the form
\SGARefLink{sga3:viii:examples:6:5:eq:xxx:01}{(xxx)}, with
\(P=X\times_SY\) and \(P'\) the inverse image of the diagonal under
\(r\), that is, the kernel of
\[
  X\times_SY
  \mathop{\rightrightarrows}^{r_1}_{r_2}
  Z.
\]
Thus we are in the situation of \SGARefLink{sga3:viii:examples:6:5:item:a:01}{(a)}.

% Authority standalone p. 20; original-edition pagination marker 24.
Consequently, \textit{if \(Y\) is essentially free over \(S\) and \(Z\)
is separated over \(S\), then the subfunctor \(X'\) of \(X\) is
representable by a closed subprescheme of \(X\).}

\smallskip
\noindent\SGARefTarget{sga3:viii:examples:6:5:item:c:01}\textit{(c)}
Suppose given a morphism
\SGARefTarget{sga3:viii:examples:6:5:item:c:display:q:01}\[
  q:X\longrightarrow\underline{\operatorname{Hom}}_S(Y,Y);
\]
that is, ``\(X\) acts on \(Y\).'' Let \(X'\) be the ``kernel'' of \SGARefLink{sga3:viii:examples:6:5:item:c:display:q:01}{this
morphism}, namely the subfunctor of \(X\) for which \(X'(T)\) is the set
of \(x\in X(T)\) such that \(q(x):Y_T\to Y_T\) is the identity. This
functor falls under \SGARefLink{sga3:viii:examples:6:5:item:b:01}{(b)}, as one sees by introducing a second morphism
\[
  q':X\longrightarrow\underline{\operatorname{Hom}}_S(Y,Y)
\]
that ``makes \(X\) act trivially on \(Y\).'' Hence, \textit{if \(Y\) is
essentially free and separated over \(S\), then the kernel subfunctor of
\(q\) is representable by a closed subprescheme of \(X\).}

\smallskip
\noindent\SGARefTarget{sga3:viii:examples:6:5:item:d:01}\textit{(d)}
Under the hypotheses of \SGARefLink{sga3:viii:examples:6:5:item:c:01}{(c)}, consider the subfunctor \(Y'\) of \(Y\) of
``invariants under \(X\).'' Thus \(Y'(T)\) is the set of \(y\in Y(T)\)
such that the corresponding morphism \(\bar q(y):X_T\to Y_T\) is the
``\(T\)-constant morphism with value \(y\).'' Introducing \(q'\) as in
\SGARefLink{sga3:viii:examples:6:5:item:c:01}{(c)}, and writing \(\bar q\) and \(\bar q'\) for the corresponding morphisms,
\[
  \bar q,\bar q':Y\rightrightarrows
  \underline{\operatorname{Hom}}_S(X,Y),
\]
one sees that \(Y'=\operatorname{Ker}(\bar q,\bar q')\). It therefore again
falls under \SGARefLink{sga3:viii:examples:6:5:item:b:01}{(b)}, with the roles of \(X\) and \(Y\) reversed and \(Z=Y\).
Consequently, \textit{if \(X\) is essentially free over \(S\) and \(Y\)
is separated over \(S\), then the subfunctor \(Y'\) of \(Y\) of
invariants under \(X\) is representable by a closed subprescheme of
\(Y\).}

\smallskip
\noindent\SGARefTarget{sga3:viii:examples:6:5:item:e:01}\textit{(e)}
Constructions of the type described in \SGARefLink{sga3:viii:examples:6:5}{the preceding examples} occur
especially often in group theory. Thus, when a group prescheme \(G\)
over \(S\) acts on the \(S\)-prescheme \(X\),
\[
  q:G\longrightarrow\underline{\operatorname{Aut}}_S(X),
\]
the kernel of \(q\), the ``subgroup of \(G\) acting trivially,'' is a
closed subprescheme of \(G\) provided that \(X\) is essentially free
and separated over \(S\), by \SGARefLink{sga3:viii:examples:6:5:item:c:01}{(c)}. The subobject \(X^G\) of invariants
is a closed subprescheme of \(X\), provided that \(G\) is essentially
free over \(S\) and \(X\) is separated over
\(S\),\footnote{Editor's note: We have corrected \(S_X\) to \(S\).}
by \SGARefLink{sga3:viii:examples:6:5:item:d:01}{(d)}.

Let \(Y,Z\) be subpreschemes of \(X\), and consider the subfunctor
\(\underline{\operatorname{Transp}}_G(Y,Z)\) of \(G\), the
``transporter of \(Y\) into \(Z\).'' Its points with values in an
\(S\)-prescheme \(T\) are the \(g\in G(T)\) such that the corresponding
automorphism of \(X_T\) satisfies \(g(Y_T)\subseteq Z_T\), that is,
induces a morphism \(Y_T\to X_T\) factoring through \(Z_T\). Hence,
\textit{if \(Y\) is essentially free over \(S\), and \(Z\) is closed in
\(X\), then \(\underline{\operatorname{Transp}}_G(Y,Z)\) is a closed
subprescheme of \(G\)}, by \SGARefLink{sga3:viii:examples:6:5:item:a:01}{(a)}.

One may also consider the \emph{strict transporter} of \(Y\) into
\(Z\),\footnote{Editor's note: It is denoted
\(\underline{\operatorname{Transp}}^{\mathrm{str}}_G(Y,Z)\).}
whose \(T\)-valued points are the \(g\in G(T)\) such that
\(g(Y_T)=Z_T\). It is precisely
\[
  \underline{\operatorname{Transp}}_G(Y,Z)
  \cap
  \sigma\bigl(\underline{\operatorname{Transp}}_G(Z,Y)\bigr),
\]
where \(\sigma\) is the symmetry of \(G\). Consequently, if \(Y\) and
\(Z\) are essentially free over \(S\) and closed in \(X\), then the
strict transporter of \(Y\) into \(Z\) is a closed subprescheme of
\(G\).

% Authority standalone p. 21; original-edition pagination marker 25.
An important case is \(X=G\), with \(G\) acting on itself by inner
automorphisms. If \(H\) is a subprescheme of \(G\), the strict
transporter of \(H\) into itself is also called the \emph{normalizer of
\(H\) in \(G\)}, and is denoted
\(\underline{\operatorname{Norm}}_G H\). Hence, \textit{if \(H\) is a
closed subgroup prescheme of \(G\), essentially free over \(S\), then
\(\underline{\operatorname{Norm}}_G H\) is representable by a closed
subgroup prescheme of \(G\).}

Finally, let \(Z\) be a subprescheme of \(G\). Its \emph{centralizer}
\(\underline{\operatorname{Centr}}_G(Z)\) in \(G\) is the group
subfunctor of \(G\) defined by the procedure in \SGARefLink{sga3:viii:examples:6:5:item:d:01}{(d)}, regarding ``\(Z\)
acts on \(G\)'' through the action induced by that of \(G\). Thus,
\textit{if \(Z\) is essentially free over \(S\) and \(G\) is separated
over \(S\), then \(\underline{\operatorname{Centr}}_G(Z)\) is a closed
subgroup prescheme of \(G\).} In particular, \textit{if \(G\) is
essentially free and separated over \(S\), then the center \(C\) of
\(G\), which is none other than
\(\underline{\operatorname{Centr}}_G(G)\), is a closed subgroup
prescheme of \(G\).}

When \(S\) is the spectrum of a field,
\SGARefLink{sga3:viii:proposition:6:2:item:b:01}{6.2(b)}
\emph{[the source prints 6.3(b)]} shows that in examples
\SGARefLink{sga3:viii:examples:6:5:item:a:01}{(a)}--\SGARefLink{sga3:viii:examples:6:5:item:e:01}{(e)} above the ``essentially free'' hypotheses are automatically
satisfied; only separation hypotheses remain. Recalling that a
group prescheme over a field is necessarily separated, one obtains,
for example:

\medskip
\noindent\SGARefTarget{sga3:viii:corollary:6:7}\textbf{Corollary 6.7.}
\footnote{Editor's note: We have retained the numbering of the original;
there is no 6.6.}
\textit{Let \(G\) be a group prescheme over a field \(k\). Then:}
\begin{list}{--}{\setlength{\leftmargin}{1.8em}\setlength{\itemsep}{0.35em}}
  \item\SGARefTarget{sga3:viii:corollary:6:7:item:1:01}
  \textit{For every subprescheme \(Z\) of \(G\), the centralizer of
  \(Z\) in \(G\) is a closed subgroup prescheme of \(G\). In particular,
  this holds for the center
  \(\underline{\operatorname{Centr}}_G(G)\) of \(G\).}

  \item\SGARefTarget{sga3:viii:corollary:6:7:item:2:01}
  \textit{More generally, if \(u,v:X\to G\) are morphisms of
  preschemes, then \(\underline{\operatorname{Transp}}_G(u,v)\) is
  representable by a closed subprescheme of \(G\).}

  \item\SGARefTarget{sga3:viii:corollary:6:7:item:3:01}
  \textit{For two subpreschemes \(Y,Z\) of \(G\), with \(Z\) closed,
  \(\underline{\operatorname{Transp}}_G(Y,Z)\) is a closed subprescheme
  of \(G\). If \(Y\) is also closed, \SGARefLink{sga3:viii:corollary:6:7:item:3:01}{the same conclusion} holds for
  \(\underline{\operatorname{Transp}}^{\mathrm{str}}_G(Y,Z)\).}

  \item\SGARefTarget{sga3:viii:corollary:6:7:item:4:01}
  \textit{For every subgroup prescheme}\footnote{Editor's note: Indeed,
  over a field \(k\), every subgroup prescheme of \(G\) is closed; cf.
  Expos\'e VI\textsubscript{A}, 0.6.1.}
  \textit{\(H\) of \(G\), the normalizer
\(\underline{\operatorname{Norm}}_G(H)\) is a closed subgroup
prescheme of \(G\).}
\end{list}

\medskip
\noindent\SGARefTarget{sga3:viii:remark:6:8}\textbf{Remark 6.8.}
\footnote{Editor's note: We have expanded the original in what follows;
in particular, we have added \SGARefLink{sga3:viii:lemma:6:8:1}{Lemma 6.8.1}.}
Let \(A\) be a commutative ring, let \(M\) be an \(A\)-module, and put
\[
  M^\vee=\operatorname{Hom}_A(M,A).
\]
Endow the \(A\)-module \(\operatorname{End}_A(M)\) with the topology of
discrete pointwise convergence. Thus a basis of neighborhoods of zero
is formed by the following \(A\)-submodules, where \(n\in\mathbb N\)
and \(m_1,\ldots,m_n\in M\):
\SGARefTarget{sga3:viii:remark:6:8:display:neighborhood:01}
\[
  K(m_1,\ldots,m_n)
  =
  \{u\in\operatorname{End}_A(M)\mid
  u(m_i)=0\ \text{for \(i=1,\ldots,n\)}\}.
\]
We say that \(M\) is a \emph{quasi-free \(A\)-module} if the image of
the canonical morphism
\[
  \Theta:M\otimes_AM^\vee\longrightarrow\operatorname{End}_A(M)
\]
contains \(\operatorname{id}_M\) in its closure; that is, if \SGARefLink{sga3:viii:remark:6:8:eq:star:01}{the
following condition} holds:

% Authority standalone p. 22.
\[
  \left\{
  \begin{aligned}
    &\text{for all \(m_1,\ldots,m_n\in M\), there exist
      \(x_1,\ldots,x_r\in M\) and \(f_1,\ldots,f_r\in M^\vee\)}\\
    &\text{such that }\quad
    m_i
    =
    \Theta\!\left(\sum_{s=1}^r x_s\otimes f_s\right)(m_i)
    =
    \sum_{s=1}^r f_s(m_i)x_s
    \quad\text{for \(i=1,\ldots,n\).}
  \end{aligned}
  \right.
  \tag{*}\label{sga3:viii:remark:6:8:eq:star:01}
\]
\noindent(In this case \(\operatorname{Im}\Theta\) is dense in
\(\operatorname{End}_A(M)\), since for every
\(u\in\operatorname{End}_A(M)\) one has
\(
  u(m_i)
  =
  \sum_{s=1}^r f_s(m_i)u(x_s)
  =
  \Theta\!\left(\sum_{s=1}^r u(x_s)\otimes f_s\right)(m_i).
\))

First note that \SGARefLink{sga3:viii:remark:6:8:eq:star:01}{this property} is stable under base change. Indeed, let
\(\phi:A\to A'\) be a ring homomorphism, put \(M'=M\otimes_AA'\), and
let \(m_1',\ldots,m_n'\in M'\). Write
\[
  m_i'=\sum_jm_{ij}\otimes b_{ij},
  \qquad m_{ij}\in M,\quad b_{ij}\in A'.
\]
By hypothesis there exist \(x_1,\ldots,x_r\in M\) and
\(f_1,\ldots,f_r\in\operatorname{Hom}_A(M,A)\) such that
\[
  m_{ij}=\sum_sx_sf_s(m_{ij})
  \qquad\text{for all \(i,j\).}
\]
Write \(\phi\circ f_s\) for the image of \(f_s\) in
\[
  \operatorname{Hom}_A(M,A')
  =
  \operatorname{Hom}_{A'}(M',A').
\]
Then, for \(i=1,\ldots,n\),
\SGARefTarget{sga3:viii:remark:6:8:display:base-change:01}
\[
  \begin{aligned}
  \left(\sum_sx_s\otimes\phi\circ f_s\right)(m_i')
    &=
    \sum_{s,j}x_s\otimes\phi(f_s(m_{ij}))\,b_{ij}\\
    &=
    \sum_j\left(\sum_sx_sf_s(m_{ij})\right)\otimes b_{ij}\\
    &=m_i'.
  \end{aligned}
\]
This proves that \(M'\) is quasi-free over \(A'\).

Also, every projective \(A\)-module \(P\) is quasi-free. Indeed, there
are \(A\)-morphisms
\[
  A^{(I)}\xrightarrow{\ \pi\ }P
  \xrightarrow{\ \tau\ }A^{(I)}
  \qquad\text{such that}\qquad
  \pi\circ\tau=\operatorname{id}_P.
\]
Let \((e_i)\) be the canonical basis of \(A^{(I)}\), and put
\(f_i=e_i^*\circ\tau\), a linear form on \(P\). Given
\(m_1,\ldots,m_n\in P\), there is a finite subset \(J\) of \(I\) such
that
\[
  m_k=\sum_{i\in J}f_i(m_k)\pi(e_i)
  \qquad\text{for \(k=1,\ldots,n\).}
\]

Theorem \SGARefLink{sga3:viii:theorem:6:4}{6.4} remains valid if, in the
statement of \SGARefLink{sga3:viii:definition:6:1}{Definition 6.1},
``free'' is replaced by ``projective'' or, more generally, by
``quasi-free.'' Proceeding as in the proof of
\SGARefLink{sga3:viii:theorem:6:4}{6.4}, one reduces to \SGARefLink{sga3:viii:lemma:6:8:1}{the following
lemma}.

\medskip
\noindent\SGARefTarget{sga3:viii:lemma:6:8:1}\textbf{Lemma 6.8.1.}
\textit{Let \(M\) be a quasi-free \(A\)-module, \(N\) a submodule, and
\(F\) the covariant functor}
\[
  F:(\text{\(A\)-algebras})\longrightarrow\mathbf{Set}
\]
\textit{such that \(F(B)=\{\mathrm{pt}\}\) if
\(M_B=(M/N)_B\), and \(F(B)=\varnothing\) otherwise. Then there is an
ideal \(K\) of \(A\) such that \(F(B)=\{\mathrm{pt}\}\) if and only if
the morphism \(A\to B\) factors through \(A/K\); that is, there is an
isomorphism, functorial in \(B\),}
\SGARefTarget{sga3:viii:lemma:6:8:1:display:representing:01}
\[
  F(B)\simeq\operatorname{Hom}_{A\text{-alg.}}(A/K,B).
\]

\smallskip
\noindent\textit{Proof.}
Let \((n_j)\) be a system of generators of \(N\), and let \(F_j\) be the
subfunctor of the terminal functor \(e\), where
\(e(B)=\{\mathrm{pt}\}\) for every \(B\), corresponding to the
submodule \(An_j\). We have \(F(B)=\{\mathrm{pt}\}\) if and only if the
image of every \(n_j\) in \(M_B\) is zero. Hence \(F\) is the
intersection of the functors \(F_j\), reducing us to the case in which
\(N\) is generated by one element \(n\).

Let \(B\) be an \(A\)-algebra, with structure map \(\phi:A\to B\).
If the image \(n\otimes1\) of
\(n\) in \(M_B\) is zero, then for every
\(f\in M^\vee=\operatorname{Hom}_A(M,A)\),
\[
  0=f(n)\otimes1=\phi(f(n)).
\]
Conversely, because \(M\) is quasi-free, there exist
\(x_1,\ldots,x_r\in M\) and \(f_1,\ldots,f_r\in M^\vee\) such that
\[
  n=\sum_sf_s(n)x_s,
\]
whence
\[
  n\otimes1=\sum_sx_s\otimes\phi(f_s(n)).
\]
It follows that \(n\otimes1=0\) if and only if \(\phi\) factors through
\(A/K\), where \(K\) is the ideal generated by the \(f_s(n)\), for
\(s=1,\ldots,r\). This proves the lemma.

% END INCLUDED SGA3_English_current_buildable_source_20260823/source/inputs/VIII/tex/components/delegated_sections5_6/06_expose_VIII_section6_en.tex

% BEGIN INCLUDED SGA3_English_current_buildable_source_20260823/source/inputs/VIII/tex/components/05_expose_VIII_section7_part1.tex
% SGA 3, Exposé VIII, English reconstruction.
% Scope: §7 opening through Corollary 7.4.
% Authority: Polo--Gille Exp8-8nov09.pdf, component-local pp. 23--26,
% complete-reader pp. 639--642.
% The current Reinhold English is comparison material only; its own version
% notice says §7 requires revision. Every formula and statement here was
% checked against the Polo--Gille authority.

\setcounter{footnote}{31}
\SGARefTarget{sga3:viii:section:7}\section*{7. Appendix: On Monomorphisms of
Group Preschemes}
\addcontentsline{toc}{section}{7. Appendix: On Monomorphisms of Group Preschemes}

% Source page: Exp8 component-local p. 23; complete-reader p. 639.
The result proved in \SGARefLink{sga3:viii:section:7}{this section} is not needed in the remainder of the
seminar, except in X, \SGARefLink{sga3:x:theorem:8-8}{8.8} and Expos\'es XV and XVI. It depends essentially
on the existence theorem for quotient groups in Expos\'e VI\(_A\).
\footnote{Editor's note: cf. Expos\'e VI\(_A\), \SGARefLink{sga3:viii:corollary:3:2}{Theorems 3.2} and \SGARefLink{thm:sga3-VIA-3.3.2}{3.3.2}.}

\SGARefLink{sga3:viii:corollary:5:7}{Corollary 5.7} raises the question of when a monomorphism
\(u:G\to H\) of \(S\)-groups can be asserted to be an immersion, or even a
closed immersion.

We saw in VI\(_B\), \SGARefLink{cor:sga3-VIB-1.4.2}{1.4.2} that this is so when \(S\) is the spectrum of a
field \(k\), provided that \(G\) is of finite type over \(k\) and \(H\) is
locally of finite type over \(k\). It follows readily that the same result
remains valid when \(S\) is merely assumed Artinian.\footnote{Editor's note:
the additions made in Expos\'e VI\(_B\) should be taken into account.}

On the other hand, examples of bijective monomorphisms that are not
immersions are easy to give, for instance over the affine line over a field
or over the spectrum of a discrete valuation ring. Take
\[
  H=(\mathbf Z/2\mathbf Z)_S,\qquad G=G_1\times_SG_2,
\]
where \(G_1\) is the open subgroup of \(H\) complementary to the closed
point \(x\ne0\) in the fiber \(H_s\simeq\mathbf Z/2\mathbf Z\), with \(s\)
a fixed closed point of \(S\), and \(G_2\) is the closed subscheme of \(H\)
which is the union of the subscheme given by the identity section and the
reduced closed subscheme defined by \(x\). One checks immediately that
\(G_2\) is stable under multiplication in \(H\), and hence is a group
scheme. The immersions \(G_i\to H\), \(i=1,2\), then define a homomorphism
of \(S\)-groups
\[
  G=G_1\times_SG_2\longrightarrow H
\]
which is visibly a bijective monomorphism, and even a local immersion, but
is not an immersion. Indeed, \(G\) and \(H\) are reduced; \(G\) has three
disjoint irreducible components, whereas \(H\) has only two.

Notice also that \(G_1\to H\) is an example of an open subgroup of \(H\)
that is not closed, in contrast to the situation for algebraic groups over
a field. The theory of degenerations of elliptic curves gives further
examples of this phenomenon in which \(H\) is smooth over \(S\) of relative
dimension one and \(G\) has connected fibers.

One might hope\begingroup
\renewcommand{\thefootnote}{(\(\ast\))}%
\footnote[1]{In fact, M. Raynaud constructed a counterexample with \(G\)
smooth and with connected fibers; cf. XVI, \SGARefLink{sga3:xvi:subsection:1:1:item:c}{1.1(c)}. Without the connected
fiber hypothesis, one may take
\(S=\operatorname{Spec}(\mathbf Z_2)\), let
\(G=(\mathbf Z/2\mathbf Z)_S\) with its nonidentity closed point removed,
and let \(H=(\mu_2)_S\).}%
\endgroup
that once \(G\) is assumed flat over \(S\), and, say, \(G\) and \(H\) are
of finite presentation over \(S\), every monomorphism \(u:G\to H\) of
\(S\)-groups is automatically an immersion. We shall prove a result of this
kind under additional hypotheses.

First, we may assume that \(S\) is affine and, by the finite-presentation
hypothesis on \(G\) and \(H\), that \(S\) is Noetherian: one reduces to the
spectrum of a ring of finite type over \(\mathbf Z\).\footnote{Editor's
note: cf. EGA~IV\(_3\), §8 and Expos\'e VI\(_B\), §10.} Then \(G\) and
\(H\) are Noetherian. Saying that \(u:G\to H\) is a closed immersion
(respectively, an immersion) is then equivalent to saying that \(u\) is a
monomorphism---which holds by hypothesis---and is proper (respectively,
proper at every point of \(u(G)\)).\footnote{Editor's note: cf.
EGA~IV\(_3\), 8.11.5 and 15.7.1.}

% Source page: Exp8 component-local p. 24; complete-reader p. 640.
The valuative criterion for properness shows that it is enough to check,
after every base change \(S'\to S\) with \(S'\) the spectrum of a discrete
valuation ring, complete if desired, that the morphism
\(u':G'\to H'\) has the corresponding properness property. The case of
local properness was omitted from EGA~II, 7.3 and was to appear as an
erratum in EGA~IV.\begingroup
\renewcommand{\thefootnote}{(\(\ast\))}%
\footnote[0]{Cf. EGA~IV\(_3\), 15.7.}%
\endgroup
We are therefore reduced to the case in which \(S\) itself is the spectrum
of a complete discrete valuation ring, subject to the proviso that any
additional hypotheses imposed on \(G,H,u\) be stable under base change.

Let \(s\), respectively \(s'\), be the closed, respectively generic, point
of \(S\). The induced homomorphisms on the fibers,
\[
  u_s:G_s\longrightarrow H_s,
  \qquad
  u_{s'}:G_{s'}\longrightarrow H_{s'},
\]
are closed immersions because they are monomorphisms of group schemes of
finite type over fields (VI\(_B\), \SGARefLink{cor:sga3-VIB-1.4.2}{1.4.2}). We may therefore identify
\(G_{s'}\) with a closed subscheme of \(H_{s'}\). We shall use \SGARefLink{sga3:viii:lemma:7:1}{the
following result}.

\medskip
\noindent\SGARefTarget{sga3:viii:lemma:7:1}\textbf{Lemma 7.1.}
\textit{Let \(S\) be the spectrum of a discrete valuation ring, with generic
point \(s'\). Let \(H\) be an \(S\)-prescheme and let \(L'\) be a closed
subprescheme of the generic fiber \(H_{s'}\), hence also a subprescheme of
\(H\). Its schematic closure \(L\) in \(H\)---the smallest closed
subprescheme of \(H\) containing \(L'\), cf. EGA~I, 9.5---exists. It is
also the unique closed subprescheme of \(H\), flat over \(S\), whose generic
fiber is \(L'\). Formation of \(L\) is functorial in the pair \((H,L')\)
and commutes with fiber products over \(S\). In particular, if \(H\) is an
\(S\)-group and \(L'\) is a subgroup of \(H_{s'}\), then \(L\) is a
subgroup of \(H\).}

The proof is immediate and is left to the reader.\begingroup
\renewcommand{\thefootnote}{(\(\ast\ast\))}%
\footnote[0]{Cf. EGA~IV\(_2\), 2.8.}%
\endgroup
Apply this to the pair \((H,G_{s'})\). The monomorphism \(u:G\to H\)
factors as
\[
  G\xrightarrow{\,u'\,}L\longrightarrow H,
\]
where \(L\) is a subgroup of \(H\), a closed subprescheme flat over \(S\),
and \(u':G\to L\) induces an isomorphism on generic fibers. Thus \(u:G\to H\)
is an immersion (respectively, a closed immersion) if and only if
\(u':G\to L\) is one. We are reduced to the case in which \(H\) is flat over
\(S\) and \(u_{s'}:G_{s'}\to H_{s'}\) is an isomorphism, again subject
to the requirement that any additional hypotheses remain valid when \(H\)
is replaced by a closed subgroup prescheme.

Since \(H\) is then the schematic closure of \(H_{s'}\), if \(u\) is an
immersion its image \(u(G)\) is a subprescheme\footnote{Editor's note: the
word ``closed'' has been removed.} of \(H\) containing \(H_{s'}\). Its
schematic closure is therefore \(H\), so \(u(G)\) is an open subprescheme of
\(H\). We must consequently prove that \(u\) is an open immersion, or an
isomorphism if the desired conclusion is that it is a closed immersion.

% Source page: Exp8 component-local p. 25; complete-reader p. 641.
Because \(G\) and \(H\) are flat over \(S\), this is equivalent to the
induced maps on the fibers being open immersions, respectively
isomorphisms (cf. SGA~1, I, \SGARefLink{sga3:viii:corollary:5:7}{5.7}). It already holds for \(u_{s'}\), so it
remains to prove that \(u_s\) is an open immersion, respectively an
isomorphism.

Since \(G\) and \(H\) are flat over \(S\), their fiber dimensions are
constant (VI\(_B\), \SGARefLink{sga3:viii:corollary:4:3}{4.3}). Their generic fibers agree, hence these dimensions
are the same for \(G\) and \(H\). Thus \(u_s:G_s\to H_s\) is a
monomorphism of algebraic groups of the same dimension. It follows that
\(G_s\) is open in \(H_s\), and set-theoretically is a union of connected
components of \(H_s\). Indeed, after passing to an algebraically closed
base field and reducing \(G\) and \(H\), both are smooth over the field and
the assertion is immediate.

The complement \(H_s-G_s\) is closed in \(H_s\), hence in \(H\). Its
complement \(H'\) in \(H\) is open and is stable under the group law. Upon
replacing \(H\) by \(H'\), we may therefore reduce, with the usual proviso,
to the case in which \(u\) is bijective, that is, \(u_s:G_s\to H_s\) is
bijective. In every case we are reduced to proving that \(u\), equivalently
\(u_s\), is an isomorphism, possibly under this bijectivity hypothesis.

Assume first that \(u\) is bijective. If \(H_s\) is reduced, then \(u_s\) is
an isomorphism, because \(G_s\) identifies with a closed subscheme of
\(H_s\) having the same underlying set. In particular, if
\(k=\kappa(s)\) has characteristic zero, every algebraic group over \(k\)
is reduced by Cartier's theorem (cf. VI\(_B\), 1.6.1; VII\(_B\), 3.3.1;
or EGA~IV\(_4\), 16.12.2 and 17.12.5). We have proved:

\medskip
\noindent\SGARefTarget{sga3:viii:proposition:7:2}\textbf{Proposition 7.2.}
\textit{Let \(u:G\to H\) be a homomorphism of group preschemes of finite
presentation over \(S\). If \(u\) is a monomorphism, \(G\) is flat over
\(S\), and every residue field of \(S\) has characteristic zero, then \(u\)
is an immersion.}

Suppose now that \(\kappa(s)\) has characteristic \(p>0\), and restrict to
the case in which \(G\) is commutative. Under the preceding reductions,
\(H\) is also commutative, because it is the schematic closure of
\(H_{s'}\simeq G_{s'}\). For \(n\geq0\), set
\[
  S_n=\operatorname{Spec}(V/\mathfrak m^{n+1}),\qquad
  G_n=G\times_SS_n,\qquad H_n=H\times_SS_n,
\]
where \(V\) is the valuation ring defining \(S\) and \(\mathfrak m\) its
maximal ideal. For \(m>0\), denote by \({}_mG\) and \({}_mH\) the kernels
of multiplication by \(m\) in \(G\) and \(H\). Define similarly
\[
  {}_m(G_n)=({}_mG)_n,
\]
which we abbreviate to \({}_mG_n\), and likewise for \(H\).

By VI\(_A\), \SGARefLink{sga3:viii:corollary:3:2}{3.2}, the quotients
\[
  Q_n=H_n/G_n
\]
exist. Each \(Q_n\) is a commutative group scheme over \(S_n\), flat over
\(S_n\), and \(H_n\to Q_n\) is faithfully flat with kernel \(G_n\).
Formation of the quotient commutes with base extension,\footnote{Editor's
note: this should be made explicit in Expos\'es V and VI\(_A\).} so
\[
  Q_n\simeq Q_m\times_{S_m}S_n\qquad(m\geq n).
\]
In particular, \(Q_n\times_{S_n}S_0=Q_0=H_0/G_0\). Since \(G_0\) and
\(H_0\) have the same underlying set, \(Q_0\) is set-theoretically a single
point: it is a purely infinitesimal group.

% Source page: Exp8 component-local p. 26; complete-reader p. 642.
It follows that every \(Q_n\) is finite and flat over \(S_n\). Let it be
defined by an algebra \(C_n\) over \(V_n=V/\mathfrak m^{n+1}\), finite free
as a \(V_n\)-module. For \(m\geq n\),
\[
  C_n=C_m\otimes_{V_m}V_n,
\]
compatibly also with the comultiplication maps. Hence
\[
  C=\varprojlim C_n
  \quad\text{is a finite free module over}\quad
  V=\varprojlim V_n,
\]
and the comultiplications on the \(C_n\) induce one on \(C\). Thus
\(Q=\operatorname{Spec}(C)\) is a finite flat group scheme over \(S\)
satisfying
\[
  Q\times_SS_n=Q_n
\]
for every \(n\).

\medskip
\noindent\SGARefTarget{sga3:viii:lemma:7:3}\textbf{Lemma 7.3.}
\footnote{Editor's note: the hypothesis that \(Q\) be commutative has been
removed here, and \SGARefLink{sga3:viii:remark:7:3:1}{7.3.1} modified accordingly. If \(Q\) is not commutative,
the \(n\)-th power morphism need not be a group homomorphism.}
\textit{Let \(K\) be a field and let \(Q\) be a finite group scheme over
\(K\), of degree \(n\). Then the \(n\)-th power morphism on \(Q\) is zero.}

See VII\(_A\), \SGARefLink{sga3:VIIA:8.5}{8.5}.

\medskip
\noindent\SGARefTarget{sga3:viii:remark:7:3:1}\textbf{Remark 7.3.1.}
The statement of Lemma~\SGARefLink{sga3:viii:lemma:7:3}{7.3} still makes
sense for a finite locally free group scheme \(Q/S\) over an arbitrary base
prescheme \(S\). It would be interesting to find a proof in that generality.

Lemma~\SGARefLink{sga3:viii:lemma:7:3}{7.3} (that is, VII\(_A\), 8.5)
does show that the proposed statement holds whenever \(S\) is reduced: apply
it to the fibers of \(Q\) at the maximal points of \(S\), equivalently the
generic points of its irreducible components.\footnote{Editor's note: this
should be added to Expos\'e VII\(_A\). On the other hand, the statement is
true for every \(S\) when \(Q\) is commutative; cf. Deligne's theorem in
\SGARefLink{sga3:viii:bib:to70}{[TO70]}, p.~4.}

In the preceding situation, where \(S\) is the spectrum of a discrete
valuation ring and \(Q\) is commutative, it follows that
\(n\operatorname{id}_Q=0\). Here \(n\) is a power
\(p^{\nu_0}\) of the residue characteristic. Thus:

\medskip
\noindent\SGARefTarget{sga3:viii:corollary:7:4}\textbf{Corollary 7.4.}
\textit{With \(Q\) and the \(Q_n\) as above, suppose that their common
degree is \(p^{\nu_0}\). Then}
\[
  p^{\nu_0}\operatorname{id}_Q=0,
  \qquad
  p^{\nu_0}\operatorname{id}_{Q_n}=0
  \quad\textit{for every \(n\)}.
\]

% END INCLUDED SGA3_English_current_buildable_source_20260823/source/inputs/VIII/tex/components/05_expose_VIII_section7_part1.tex

% BEGIN INCLUDED SGA3_English_current_buildable_source_20260823/source/inputs/VIII/tex/components/06_expose_VIII_section7_part2_and_bibliography.tex
% SGA 3, Exposé VIII, English reconstruction.
% Scope: Corollary 7.5 through §7.13 and terminal bibliography.
% Authority: Polo--Gille Exp8-8nov09.pdf, component-local pp. 26--30,
% complete-reader pp. 642--646. Hard stop before Exposé IX.
% The Polo--Gille authority controls every statement, formula, note, and
% diagram; the comparison translation was consulted only for phrasing.

% Source page: Exp8 component-local p. 26; complete-reader p. 642.
\medskip
\noindent\SGARefTarget{sga3:viii:corollary:7:5}\textbf{Corollary 7.5.}
\textit{Suppose in addition that \(G_0\) is smooth over \(k\) and that
\(p\operatorname{id}_{G_0}\) is flat. Equivalently, by the structure theory
of algebraic groups over an algebraically closed field \(\overline{k}\),
\(G_{\overline{k}}\) contains no subgroup isomorphic to the additive group.
Then, for every \(\nu\geq\nu_0\) and every \(n>0\), the group
\({}_{p^\nu}H_n\) is flat over \(S_n\).}

Indeed, \(p^{\nu_0}\operatorname{id}_{Q_n}=0\) implies that
\(p^\nu\operatorname{id}_{H_n}\) factors through \(G_n\), giving the
commutative diagram
\SGARefTarget{sga3:viii:diagram:7:5}
\[
\begin{tikzcd}[column sep=huge,row sep=large]
H_n
  \arrow[r,"p^\nu\operatorname{id}_{H_n}"]
  \arrow[dr,"v_n"]
&
H_n
\\
G_n
  \arrow[u,"u_n"]
  \arrow[r,"p^\nu\operatorname{id}_{G_n}"]
&
G_n
  \arrow[u,"u_n"']
.
\end{tikzcd}
\]
We claim that \(v_n\) is flat. Since \(H_n\) and \(G_n\) are flat over
\(S_n\), it is enough to check that \(v_0\) is flat (SGA~1, IV, 5.9).
Thus we may take \(n=0\), so that \(S_n=\operatorname{Spec}(k)\). Since
\(p\operatorname{id}_{G_0}\), and hence
\(p^\nu\operatorname{id}_{G_0}\), is flat, its image is an induced open
subgroup \(G'_0\) of \(G_0\). The map \(u_0:G_0\to H_0\) is surjective, so
\(v_0\) takes its values in \(G'_0\) and may be regarded as a homomorphism
to \(G'_0\). Its composite with \(u_0\) is an epimorphism; hence \(v_0\)
itself is an epimorphism, and therefore a flat homomorphism into \(G'_0\),
and thus into \(G_0\). Consequently \(v_n\) is flat and
% The authority prints _{p^\nu}H over S here; the immediately preceding
% definition of v_n and the statement of 7.5 require _{p^\nu}H_n over S_n.
\(\ker v_n={}_{p^\nu}H_n\)%
\begingroup
\renewcommand{\thefootnote}{\(\dagger\)}
\footnote[0]{English editor's note: the source prints
\({}_{p^\nu}H\) over \(S\) here and \({}_{p^\nu}H\) over \(S_n\) in
the proof of \SGARefLink{sga3:viii:lemma:7:7}{Lemma 7.7}. The subscripts \(n\) are supplied in both places
because those arguments concern the base-changed family over \(S_n\).}
\endgroup
is flat over \(S_n\). \qed

\medskip
\noindent\SGARefTarget{sga3:viii:remark:7:5:after:01}\textbf{Remark.}
We did not explicitly use the assumption that \(G_0\) is smooth over \(k\).
It is, however, an easy consequence of the flatness of
\(p\operatorname{id}_{G_0}\)
\emph{[the source prints \(p\operatorname{id}_{G}\)]}; this is why the
condition was stated explicitly
in Corollary~\SGARefLink{sga3:viii:corollary:7:5}{7.5}.

\medskip
\noindent\SGARefTarget{sga3:viii:lemma:7:6}\textbf{Lemma 7.6.}
\textit{Let \(u:G\to H\) be a surjective monomorphism of commutative
algebraic groups over a field \(k\) of characteristic \(p>0\), and put
\(Q=H/G\), a purely infinitesimal group. There is an integer \(\nu_1\) such
that, for every \(\nu\geq\nu_1\), the sequence}
\[
  0\longrightarrow{}_{p^\nu}G
  \longrightarrow{}_{p^\nu}H
  \longrightarrow Q\longrightarrow0
\]
\textit{is exact.}

It is enough to establish exactness at \(Q\), which in turn amounts to the
injectivity of the homomorphism
\SGARefTarget{sga3:viii:eq:7:6:x}
\begin{equation}
  \tag{x}
  \mathscr O_{Q,e}\longrightarrow
  \mathscr O_{{}_{p^\nu}H,e}
\end{equation}
of local rings at the identity elements. Recall that \(Q\) is
set-theoretically supported at the single point \(e\). There is a natural
homomorphism
\SGARefTarget{sga3:viii:eq:7:6:xx}
\begin{equation}
  \tag{xx}
  \mathscr O_{{}_{p^\nu}H,e}
  \longrightarrow
  \mathscr O_{H,e}/\mathfrak m^{p^\nu},
\end{equation}
where \(\mathfrak m\) is the augmentation ideal, equivalently the maximal
ideal, of \(\mathscr O_{H,e}\). Indeed,
\(p^\nu\operatorname{id}_H\) vanishes on the kernel of the iterated
Frobenius homomorphism \(F^\nu\). The composite of
\SGARefLink{sga3:viii:eq:7:6:x}{(x)} and
\SGARefLink{sga3:viii:eq:7:6:xx}{(xx)} is the
natural composite
\SGARefTarget{sga3:viii:eq:7:6:xxx}
\begin{equation}
  \tag{xxx}
  \mathscr O_{Q,e}\longrightarrow
  \mathscr O_{H,e}\longrightarrow
  \mathscr O_{H,e}/\mathfrak m^{p^\nu}.
\end{equation}
Now \(\mathscr O_{Q,e}\to\mathscr O_{H,e}\) is injective, because
\(H\to Q\) is an epimorphism and therefore flat. Moreover,
\(\mathscr O_{Q,e}\) is Artinian, while
\(\bigcap_\nu\mathfrak m^{p^\nu}=0\) in \(\mathscr O_{H,e}\), so the
intersection of the corresponding ideals in \(\mathscr O_{Q,e}\) is also
zero. One of them is therefore zero. Thus
\SGARefLink{sga3:viii:eq:7:6:xxx}{(xxx)} is
injective, and a fortiori so is
\SGARefLink{sga3:viii:eq:7:6:x}{(x)}.

% Source page: Exp8 component-local p. 28; complete-reader p. 644.
\medskip
\noindent\SGARefTarget{sga3:viii:lemma:7:7}\textbf{Lemma 7.7.}
\textit{Under the hypotheses of \SGARefLink{sga3:viii:corollary:7:5}{7.5}, there is an integer \(\nu_1\) such
that, for every \(\nu\geq\nu_1\) and every \(n\), the sequence of
\(S_n\)-groups}
\[
  0\longrightarrow{}_{p^\nu}G_n
  \longrightarrow{}_{p^\nu}H_n
  \xrightarrow{\,w_n\,}Q_n
  \longrightarrow0
\]
\textit{is exact. More precisely, \(w_n\) is faithfully flat and its kernel
is \({}_{p^\nu}G_n\).}

Take \(\nu_1\) as in Lemma~\SGARefLink{sga3:viii:lemma:7:6}{7.6}, applied
to \(u_0:G_0\to H_0\), and require \(\nu_1\geq\nu_0\), where
\(p^{\nu_0}=\operatorname{rank}Q_0\). It remains only to check that \(w_n\)
is faithfully flat. By Corollary~\SGARefLink{sga3:viii:corollary:7:5}{7.5},
% The authority prints _{p^\nu}H over S_n here; the base-changed family
% under discussion is _{p^\nu}H_n over S_n.
\({}_{p^\nu}H_n\) is flat over \(S_n\), as is \(Q_n\); it is therefore
enough to check that
\[
  w_0:{}_{p^\nu}H_0\longrightarrow Q_0
\]
is faithfully flat, equivalently an epimorphism. This follows from the
choice of \(\nu_1\).

\medskip
\noindent\SGARefTarget{sga3:viii:corollary:7:8}\textbf{Corollary 7.8.}
\textit{Suppose in addition that \(H\) is separated over \(S\), or more
generally that \({}_{p^\nu}H\) is separated over \(S\) for every \(\nu\),
and that \({}_{p^\nu}G\) is finite over \(S\) for every \(\nu\). Then, for
\(\nu\geq\nu_1\), the group schemes \({}_{p^\nu}G\) and
\({}_{p^\nu}H\) are finite and flat over \(S\), and there is an exact
sequence}
\SGARefTarget{sga3:viii:corollary:7:8:display:exact-sequence:01}\[
  0\longrightarrow{}_{p^\nu}G
  \longrightarrow{}_{p^\nu}H
  \xrightarrow{\,w\,}Q
  \longrightarrow0.
\]

Because \(u:G\to H\) is surjective, the induced morphism
\({}_{p^\nu}G\to{}_{p^\nu}H\) is also surjective. The source is finite over
\(S\) and the target is separated over \(S\), hence the target is finite
over \(S\). To prove that \({}_{p^\nu}G\) and \({}_{p^\nu}H\) are flat over
\(S\), it is enough to prove this after every base change to \(S_n\), where
it is contained in Corollary~\SGARefLink{sga3:viii:corollary:7:5}{7.5}.
Finally, \SGARefLink{sga3:viii:corollary:7:8:display:exact-sequence:01}{the exact sequence above} is obtained from the exact sequences in
Lemma~\SGARefLink{sga3:viii:lemma:7:7}{7.7} as \(n\) varies.

Taking generic fibers in the sequence of
Corollary~\SGARefLink{sga3:viii:corollary:7:8}{7.8}, and recalling that
\(u_{s'}:G_{s'}\to H_{s'}\) is an isomorphism, gives that \(Q_{s'}\) is
the unit group. Since \(Q\) is flat over \(S\), \(Q\) itself is the unit
group; therefore so is \(Q_s\), and \(u_s:G_s\to H_s\) is an isomorphism.
We have proved:

\medskip
\noindent\SGARefTarget{sga3:viii:proposition:7:9}\textbf{Proposition 7.9.}
\textit{Let \(u:G\to H\) be a homomorphism of group preschemes of finite
presentation over a prescheme \(S\). Assume:}
\begin{enumerate}
  \item[(a)]\SGARefTarget{sga3:viii:proposition:7:9:item:a}
  \textit{\(u\) is a monomorphism;}
  \item[(b)]\SGARefTarget{sga3:viii:proposition:7:9:item:b}
  \textit{\(G\) is flat over \(S\);}
  \item[(c)]\SGARefTarget{sga3:viii:proposition:7:9:item:c}
  \textit{for every \(s\in S\) whose residue field has characteristic
  \(p>0\), let \(S'=\operatorname{Spec}(\mathscr O_{S,s})\), and let
  \(u':G'\to H'\) be obtained from \(u\) by base change to \(S'\). Then
  \(G'\) is commutative, the special fiber \(G'_0=G_s\) is smooth over
  \(\kappa(s)\), and, for every \(\nu>0\), the group
  \({}_{p^\nu}G'\) is finite over \(S'\) while \({}_{p^\nu}H'\) is
  separated over \(S'\).}
\end{enumerate}
\textit{Under these assumptions, \(u\) is an immersion.}

Indeed, the finiteness of \({}_pG'\) over \(S'\) implies the finiteness of
its special fiber over \(\kappa(s)\). Hence
\(G'_0\otimes_{\kappa(s)}\overline{\kappa(s)}\) contains no subgroup
isomorphic to the additive group, and the hypotheses of
Corollary~\SGARefLink{sga3:viii:corollary:7:5}{7.5} apply.

% Source page: Exp8 component-local p. 29; complete-reader p. 645.
\medskip
\noindent\SGARefTarget{sga3:viii:remark:7:10}\textbf{Remark 7.10.}
Examples of M. Raynaud (XVI, 1.1(a),\SGARefLink{sga3:viii:proposition:7:9:item:b}{(b)}) show that in condition \SGARefLink{sga3:viii:proposition:7:9:item:c}{(c)} one
cannot omit either the finiteness of the \({}_{p^\nu}G'\) over \(S'\) or
the separatedness of the \({}_{p^\nu}H'\) over \(S'\).

We now seek conditions under which \(u\) is a closed immersion. Retain the
hypotheses preceding Proposition~\SGARefLink{sga3:viii:proposition:7:9}{7.9},
but no longer suppose that \(u\) is surjective; assume only that it is an
isomorphism on generic fibers. We already know that \(u:G\to H\) is an open
immersion. The same is therefore true of \(u_0:G_0\to H_0\), whose image
contains the identity component \((H_0)^0\). Since \(G_0\) is smooth over
\(k\) and has no additive component after extending to \(\overline{k}\),
the same is true of \(H_0\). We use:

\medskip
\noindent\SGARefTarget{sga3:viii:lemma:7:11}\textbf{Lemma 7.11.}
\textit{Let \(H\) be a commutative algebraic group over a field \(k\), and
suppose that \(H\otimes_k\overline{k}\) contains no subgroup isomorphic to
\(\mathbf G_a\). If \(n=\deg(H/H^0)\), then}
\[
  {}_nH\longrightarrow H/H^0
\]
\textit{is surjective.}

We may assume that \(k\) is algebraically closed, in which case this follows
from the familiar fact that \(H^0(k)\) is a divisible group.

Return to \(u:G\to H\), and suppose that the \({}_nG\), \(n>0\), are proper
over \(S\), while the \({}_nH\) are separated over \(S\). The \({}_nH\)
are also flat over \(S\). Indeed, it is enough to check flatness at the
points above \(s\), and hence to check that
\(n\operatorname{id}_{H_0}\) is flat. As noted above, this is equivalent to
\((H_0)^0\) being smooth\footnote{Editor's note: the terminology has been
updated, replacing ``simple'' by ``smooth''; see, for example, A.
Grothendieck's note (1) in SGA~1, Expos\'e II.} over \(k\) and having no
\(\mathbf G_a\)-component over \(\overline{k}\). Since
\((H_0)^0=(G_0)^0\), this follows from the corresponding hypothesis on
\(G_0\).

Because \(H\) is separated over \(S\), so are the \({}_nH\). The morphism
\({}_nG\to{}_nH\) is proper and therefore has closed image. Its image
contains the generic fiber of \({}_nH\), because
\(u_{s'}:G_{s'}\to H_{s'}\) is an isomorphism. Since \({}_nH\) is flat
over \(S\), it is the schematic closure of its generic fiber. Hence
\({}_nG\to{}_nH\), and therefore \({}_nG_0\to{}_nH_0\), is surjective.
Together with \(G_0\supset(H_0)^0\), Lemma
\SGARefLink{sga3:viii:lemma:7:11}{7.11} now implies that
\(G_0\to H_0\) is surjective. Hence \(u\) is surjective. We obtain:

\medskip
\noindent\SGARefTarget{sga3:viii:proposition:7:12}\textbf{Proposition 7.12.}
\textit{With the notation of \SGARefLink{sga3:viii:proposition:7:9}{7.9}, suppose that conditions \SGARefLink{sga3:viii:proposition:7:9:item:a}{(a)} and \SGARefLink{sga3:viii:proposition:7:9:item:b}{(b)} hold,
together with the stronger condition (c\(^{\prime}\)) obtained from \SGARefLink{sga3:viii:proposition:7:9:item:c}{(c)} by
requiring, for every integer \(n>0\), not only those of the form \(p^\nu\),
that \({}_nG'\) be finite over \(S'\) and \({}_nH'\) separated over \(S'\).
Then \(u\) is a closed immersion.}

\medskip
\noindent\SGARefTarget{sga3:viii:remark:7:13}\textbf{Remark 7.13.}
\begin{enumerate}
  \item[(a)] The separatedness hypothesis on the \({}_nH\) easily implies
  that \(H\) itself is separated over \(S\). This is also formally contained
  in Proposition~\SGARefLink{sga3:viii:proposition:7:12}{7.12}, by taking
  \(G\) to be the unit \(S\)-group. If \(S\) is locally Noetherian, in
  Proposition~\SGARefLink{sga3:viii:proposition:7:9}{7.9} it is enough to
  assume that \(H\) is locally of finite type over \(S\), rather than of
  finite type; the proof applies verbatim. For
  Proposition~\SGARefLink{sga3:viii:proposition:7:12}{7.12}, one also
  assumes that the fibers of \(H\) are of finite type.

  % Source page: Exp8 component-local p. 30; complete-reader p. 646.
  \item[(b)]\SGAInlineRefTarget{sga3:viii:remark:7-13:item:b} Using Proposition~\SGARefLink{sga3:viii:proposition:7:12}{7.12},
  it is not difficult to prove that if \(u:G\to H\) is a monomorphism of
  group preschemes of finite presentation over \(S\), \(G\) is
  diagonalizable---or, more generally, of multiplicative type---and \(H\)
  is separated over \(S\), then \(u\) is a closed immersion. This is a
  satisfactory generalization of the first conclusion in \SGARefLink{sga3:viii:corollary:5:7}{5.7}. If \(G\) is
  smooth over \(S\), apply \SGARefLink{sga3:viii:proposition:7:12}{Proposition~7.12} directly; the general case
  reduces readily to this one. If \(H\) is not assumed separated over \(S\),
  one can still show that \(u\) is an immersion; for a torus \(G\), this
  also follows from the material that comes next.

  \item[(c)] If \(G\) is assumed to have connected fibers in
  Proposition~\SGARefLink{sga3:viii:proposition:7:9}{7.9}, the
  separatedness of \({}_{p^\nu}H'\) over \(S'\) may be omitted from
  condition \SGARefLink{sga3:viii:proposition:7:9:item:c}{(c)}. Indeed, after the reductions in the proof, \(H\) may be
  assumed flat over \(S\) and \(u\) bijective. Thus \(H\) has connected
  fibers, and Raynaud's theorem (VI\(_B\), 5.5) implies that \(H\) is
  separated over \(S\).\footnote{Editor's note: the attribution, which
  appeared in SGAD 1965, should be made explicit in Expos\'e VI\(_B\), §5,
  together with the changes between SGAD and Lecture Notes 151.}
\end{enumerate}

\SGARefTarget{sga3:viii:bibliography}\section*{Bibliography}
\addcontentsline{toc}{section}{Bibliography}
\footnote{Editor's note: additional references cited in this expos\'e.}

\begin{description}
  \item[\SGARefTarget{sga3:viii:bib:rg71}{[RG71]}]
  M. Raynaud and L. Gruson, \textit{Crit\`eres de platitude et de
  projectivit\'e}, Invent. Math. \textbf{13} (1971), 1--89.

  \item[\SGARefTarget{sga3:viii:bib:to70}{[TO70]}]
  J. Tate and F. Oort, \textit{Group schemes of prime order}, Ann. Sci.
  \'Ec. Norm. Sup\'er. (4) \textbf{3} (1970), 1--21.
\end{description}

% HARD STOP: Exposé VIII ends here. Exposé IX begins at complete-reader p. 647
% and is outside this frozen allocation.

% END INCLUDED SGA3_English_current_buildable_source_20260823/source/inputs/VIII/tex/components/06_expose_VIII_section7_part2_and_bibliography.tex



% END INCLUDED SGA3_English_current_buildable_source_20260823/source/inputs/VIII/body.tex


\SGAStartExpose{IX}
% BEGIN INCLUDED SGA3_English_current_buildable_source_20260823/source/inputs/IX/body.tex
% Mechanically extracted after \mainmatter from the frozen standalone master.

% BEGIN INCLUDED SGA3_English_current_buildable_source_20260823/source/inputs/IX/tex/components/01_expose_IX_opening_and_section_1.tex
% SGA 3, Exposé IX, English reconstruction.
% Scope: title and §1 in full.
% Authority: Polo--Gille Exp9-8nov09.pdf, component-local pp. 1--3,
% complete-reader pp. 647--649, printed pp. 37--40.
% Jacob Reinhold's CC BY 4.0 Markdown lineage was used only as a phrasing
% comparison. The Polo--Gille PDF controls text, formulas, numbering, and notes.

\setcounter{footnote}{-1}
\SGARefTarget{sga3:ix:chapter:ix}\chapter*{Exposé IX\\
Groups of Multiplicative Type:\\
Homomorphisms into a Group Scheme}
\addcontentsline{toc}{chapter}{Exposé IX. Groups of Multiplicative Type:
Homomorphisms into a Group Scheme}
\markboth{EXPOSÉ IX. GROUPS OF MULTIPLICATIVE TYPE I (HOMOMORPHISMS)}
  {EXPOSÉ IX. GROUPS OF MULTIPLICATIVE TYPE I (HOMOMORPHISMS)}

\begin{center}
  \textit{by A. Grothendieck}\footnote{Editor's note: version 1.0 of
  November 8, 2009. Additions were made in the proof of \SGARefLink{sga3:ix:theorem:3:6:bis}{3.6 bis}, in
  \SGARefLink{sga3:ix:lemma:4:4}{4.4}--\SGARefLink{sga3:ix:theorem:4:7}{7}, \SGARefLink{sga3:ix:lemma:5:0}{5.0}--\SGARefLink{sga3:ix:theorem:5:6}{6}, \SGARefLink{sga3:ix:lemma:6:1}{6.1}, \SGARefLink{sga3:ix:theorem:7:1}{7.1}, and \SGARefLink{sga3:ix:proposition:8:2}{8.2}. Sections \SGARefLink{sga3:ix:proposition:8:1}{8.1} and \SGARefLink{sga3:ix:lemma:4:5}{4.5} remain
  to be revised.}
\end{center}

\SGARefTarget{sga3:ix:section:1}\section*{1. Definitions}
\addcontentsline{toc}{section}{1. Definitions}

% Source page: Exp9 component-local p. 1; complete-reader p. 647; printed p. 37.
\medskip
\noindent\SGARefTarget{sga3:ix:definition:1:1}\textbf{Definition 1.1.} ---
\textit{Let \(S\) be a prescheme and let \(G\) be a group prescheme over
\(S\). The group \(G\) is said to be \emph{of multiplicative type} if it
is locally diagonalizable for the faithfully flat quasi-compact topology
(cf.\ \SGARefLink{sga3:iv:subsection:6:3}{IV 6.3}); that is, if for every \(s\in S\), there are an open
neighborhood \(U\) of \(s\) and a faithfully flat quasi-compact morphism
\(S'\to U\) such that
\[
  G'=G\times_S S'
\]
is a diagonalizable \(S'\)-group (\SGARefLink{sga3:i:subsection:4-4}{I 4.4} and \SGARefLink{sga3:viii:definition:1:1}{VIII 1.1}).}

\textit{The group \(G\) is said to be \emph{of quasi-isotrivial
multiplicative type} if it is locally diagonalizable even for the étale
topology (\SGARefLink{sga3:iv:subsection:6:3}{IV 6.3}); that is, if in the preceding definition one can take
\(S'\to U\) to be étale and surjective. Equivalently---as one sees by
taking the disjoint union of the schemes \(S'\) attached to the various
\(s\in S\)---there is an étale surjective morphism \(S'\to S\) such that
\(G'=G\times_S S'\) is a locally diagonalizable \(S'\)-group.}

\textit{If one can moreover choose \(S'\to S\) to be finite, étale, and
surjective, then \(G\) is said to be \emph{of isotrivial multiplicative
type}.}

\textit{Finally, \(G\) is said to be a group \emph{of locally trivial
multiplicative type} (respectively, \emph{of locally isotrivial
multiplicative type}) if every \(s\in S\) has an open neighborhood \(U\)
such that \(G\times_S U\) is a diagonalizable \(U\)-group (respectively,
is of isotrivial multiplicative type; in other words, there is a finite
étale surjective morphism \(S'\to U\) such that \(G\times_S S'\) is a
diagonalizable \(S'\)-group).}

\medskip
\noindent\SGARefTarget{sga3:ix:comments:1:2}\textbf{Comments 1.2.} ---
The five preceding notions are all obtained from the notion of a
diagonalizable group by localization, with respect to five different
``topologies'' associated with \((\mathbf{Sch})\).

It is understood in general that, when the word ``locally'' is not
qualified by naming the topology under consideration, the Zariski
topology is meant, in accordance with established terminology. In the
terminology introduced here, ``of locally trivial multiplicative type''
is equivalent to ``locally diagonalizable'' in \SGARefLink{sga3:viii:definition:1:1}{VIII 1.1}; likewise,
``trivial'' translates as ``diagonalizable.'' The five notions introduced
above fit into the following implication diagram, which follows from the
corresponding relations among the topologies involved:

% Source page: Exp9 component-local p. 2; complete-reader p. 648; printed pp. 38--39.
\SGARefTarget{sga3:ix:diagram:1:2}
\begin{center}
\begin{tikzpicture}[>=stealth]
  \node (general) at (0,4.0) {general multiplicative type};
  \node (quasi) at (0,3.0) {quasi-isotrivial};
  \node (lociso) at (0,2.0) {locally isotrivial};
  \node (iso) at (-2.7,1.0) {isotrivial};
  \node (loctriv) at (2.7,1.0) {locally trivial};
  \node (triv) at (0,0) {trivial};
  \draw[double,->] (quasi) -- (general);
  \draw[double,->] (lociso) -- (quasi);
  \draw[double,->] (iso) -- (lociso);
  \draw[double,->] (loctriv) -- (lociso);
  \draw[double,->] (triv) -- (iso);
  \draw[double,->] (triv) -- (loctriv);
\end{tikzpicture}
\end{center}

For practical purposes, let us point out at once that every group of
multiplicative type that we shall encounter will be quasi-isotrivial. Thus
we shall see in the next exposé (\SGARefLink{sga3:x:corollary:4-5}{X 4.5}) that, when \(G\) is of finite type
over \(S\), it is automatically quasi-isotrivial; we shall nevertheless
give examples in which it is not locally isotrivial. We shall likewise see
that \(G\) may be locally trivial without being isotrivial, and hence
\emph{a fortiori} without being trivial. This readily implies that the
inclusions in \SGARefLink{sga3:ix:diagram:1:2}{the preceding diagram} are strict.

On the other hand, we shall also see there (\SGARefLink{sga3:x:theorem:5-16}{X 5.16}) that, when \(S\) is
locally noetherian and normal---or, more generally, geometrically
unibranch---every group of multiplicative type and finite type over \(S\)
is necessarily isotrivial; moreover, it is trivial as soon as it is
locally trivial (or even merely trivial over a dense open subset). This
explains why most groups of multiplicative type encountered in practice
will doubtless be isotrivial, all the more so because we shall later see
that the maximal tori of semisimple group schemes are automatically
isotrivial.

\medskip
\noindent\SGARefTarget{sga3:ix:definition:1:3}\textbf{Definition 1.3.} ---
\textit{Let \(S\) and \(G\) be as in
\SGARefLink{sga3:ix:definition:1:1}{1.1}. The group \(G\) is called a
\emph{torus} if, for the faithfully flat quasi-compact topology, it is
locally isomorphic to a group of the form \(\mathbf G_m^r\), where
\(r\) is an integer with \(r\geq 0\).}

With the notation of \SGARefLink{sga3:ix:definition:1:1}{1.1}, this says
that \(S'\) may be chosen so that \(G'\) is isomorphic to a group of the
form \((\mathbf G_{m,S'})^r\). The integer \(r\) depends on \(s\in S\):
it is the dimension of the fiber
\[
  G_s=G\otimes_S\kappa(s).
\]
It is a locally constant function of \(s\), as is seen at once. This
observation should be generalized as follows.

% Source page: Exp9 component-local p. 3; complete-reader p. 649; printed pp. 39--40.
\medskip
\noindent\SGARefTarget{sga3:ix:definition:1:4}\textbf{Definition 1.4.} ---
\textit{Let \(G\) be a diagonalizable group scheme over a field \(k\).
Thus \(G\) is isomorphic to a group \(D_k(M)\), where \(M\) is an ordinary
commutative group determined up to isomorphism by this condition---more
precisely,
\[
  M\simeq\operatorname{Hom}_{k\text{-}\mathrm{gr}}
    (G,\mathbf G_{m,k})
\]
(\SGARefLink{sga3:viii:corollary:1:3}{VIII 1.3}). The isomorphism class of \(M\) is called the \emph{type} of
the diagonalizable group \(G\); it is plainly invariant under extension
of the base field.}

\textit{Now let \(G\) be a group scheme of multiplicative type over an
arbitrary prescheme \(S\). For every \(s\in S\), there is an extension
\(k\) of \(\kappa(s)\) such that
\(G\times_S\operatorname{Spec}(k)\) is a diagonalizable
\(k\)-group.\footnote{Editor's note: indeed, by hypothesis there are an
open neighborhood \(U\) of \(s\) and a faithfully flat quasi-compact
morphism \(S'\to U\) such that \(G_{S'}\) is diagonalizable. For every
\(s'\in S'\) above \(s\), the field \(\kappa(s')\) is an extension of
\(\kappa(s)\), and
\(G_{s'}=G\times_S\operatorname{Spec}(\kappa(s'))\) is diagonalizable.}
Its type is independent of the chosen extension by the preceding
observation and is called the \emph{type of \(G\) at \(s\)}, or the
\emph{type of \(G_s\)}. In particular, if \(G\) itself is diagonalizable,
and hence isomorphic to a group of the form \(D_S(M)\), then for every
\(s\in S\) the type of \(G\) at \(s\) is the isomorphism class of \(M\).}

\textit{More generally, if \(G\) is a group of multiplicative type over
\(S\) and \(M\) is an ordinary commutative group, then \(G\) is said to
be \emph{of type \(M\)} if it is of type \(M\) at every point \(s\in S\);
equivalently, if \(G\) is locally isomorphic to \(D_S(M)\) for the
faithfully flat quasi-compact topology.}

\medskip
\noindent\SGARefTarget{sga3:ix:remark:1:4:1}\textbf{Remark 1.4.1.}
\footnote{Editor's note: the number \SGARefLink{sga3:ix:remark:1:4:1}{1.4.1} has been added to the remarks
that follow, for use in later references.}
\SGARefTarget{sga3:ix:remark:1:4:1:item:a}\textit{a) For a group \(G\) of multiplicative type over \(S\), the
function
\[
  s\longmapsto \text{the type of }G_s
\]
is locally constant on \(S\). Indeed, with the notation of
\SGARefLink{sga3:ix:definition:1:1}{1.1}, if \(G'\) is of type \(M\),
then \(G\) is of type \(M\) over the neighborhood \(U\) of \(s\).
Consequently \(S\) has a canonical decomposition as a disjoint union of
preschemes \(S_i\) such that, for every \(i\), \(G_{S_i}\) is of type
\(M_i\), where the commutative groups \(M_i\) are pairwise
nonisomorphic.}

\SGARefTarget{sga3:ix:remark:1:4:1:item:b}\textit{b) In particular, if \(S\) is connected, the type of the fibers
of \(G\) is constant; that is, there is a commutative group \(M\) such
that \(G\) is of type \(M\). Finally, if \(G\) is a torus, the type of
\(G\) at \(s\) is determined by the integer \(\dim G_s\): indeed,
\(G_s\) is of type \(\mathbb Z^n\), where \(n=\dim G_s\).}

\medskip
\noindent\SGARefTarget{sga3:ix:remark:1:5}\textbf{Remark 1.5.} ---
\SGARefTarget{sga3:ix:remark:1:5:item:a}\textit{a) It follows immediately from Definitions
\SGARefLink{sga3:ix:definition:1:1}{1.1} and
\SGARefLink{sga3:ix:definition:1:3}{1.3} that these notions are
``compatible with base extension.'' Thus, if \(G\) is a group prescheme
over \(S\) and \(S'\to S\) is a base-change morphism, then whenever \(G\)
is of multiplicative type (respectively, of isotrivial multiplicative
type, and so on), the same is true of \(G'\).}

\SGARefTarget{sga3:ix:remark:1:5:item:b}\textit{b) If, moreover, \(S'\to S\) is faithfully flat and
quasi-compact, then \(G'\) being of multiplicative type, respectively a
torus, implies the same for \(G\). If \(S'\to S\) is also étale
(respectively, finite étale), then \(G'\) being quasi-isotrivial
(respectively, isotrivial) implies the same for \(G\).}

\SGARefTarget{sga3:ix:remark:1:5:item:c}\textit{c) Finally, return to an arbitrary base-change morphism
\(f:S'\to S\). If \(s'\in S'\) and \(s=f(s')\), then the type of \(G'\)
at \(s'\) is the same as that of \(G\) at \(s\), since
\[
  G'_{s'}=G_s\otimes_{\kappa(s)}\kappa(s').
\]}

% END INCLUDED SGA3_English_current_buildable_source_20260823/source/inputs/IX/tex/components/01_expose_IX_opening_and_section_1.tex

% BEGIN INCLUDED SGA3_English_current_buildable_source_20260823/source/inputs/IX/tex/components/02_expose_IX_sections_2_3.tex
% SGA 3, Expose IX, English reconstruction.
% Scope: Sections 2 and 3 in full.
% Authority: Polo--Gille Exp9-8nov09.pdf, component-local pp. 4--10
% (complete-reader pp. 650--656), stopping immediately before Section 4.
% Jacob Reinhold's CC BY 4.0 Markdown lineage was used only as a phrasing
% comparison. The Polo--Gille PDF controls text, formulas, numbering, notes,
% and references.

\setcounter{footnote}{2}

% Source page: Exp9 component-local p. 4; complete-reader p. 650;
% original printed pp. 41--42.
\SGARefTarget{sga3:ix:section:2}\section*{2. Extension of Certain Properties of Diagonalizable Groups to Groups of Multiplicative Type}
\addcontentsline{toc}{section}{2. Extension of Certain Properties of Diagonalizable Groups to Groups of Multiplicative Type}

The extensions in question are essentially trivial consequences of the
results of the preceding Expos\'e, in view of
\SGARefLink{sga3:ix:definition:1:1}{the definitions in 1.1} and the
``local'' nature (for the faithfully flat and quasi-compact topology) of
the results under consideration.

\medskip
\noindent\SGARefTarget{sga3:ix:definition:2:0}\textbf{Definition 2.0.}
\footnote{Editor's note: This definition has been added; it occurs in
\SGARefLink{sga3:ix:proposition:2:3}{Propositions~2.3} and
\SGARefLink{sga3:ix:proposition:2:7}{2.7}.}
\textit{We denote by \(\mathcal M\) the set of faithfully flat and
quasi-compact morphisms.}

\medskip
\noindent\SGARefTarget{sga3:ix:proposition:2:1}\textbf{Proposition 2.1.}
\textit{Let \(G\) be a group of multiplicative type over a prescheme
\(S\). Then:}
\begin{enumerate}
  \item[(a)]\SGARefTarget{sga3:ix:proposition:2:1:item:a:01}
  \textit{\(G\) is faithfully flat over \(S\), and affine over \(S\)
  (a fortiori, quasi-compact over \(S\));}

  \item[(b)]\SGARefTarget{sga3:ix:proposition:2:1:item:b:01}
  \textit{\(G\) is of finite type over \(S\) if and only if \(G\) is of
  finite presentation over \(S\), if and only if, for every \(s\in S\),
  the type of \(G\) at \(s\) is given by a commutative group of finite
  type;}

  \item[(c)]\SGARefTarget{sga3:ix:proposition:2:1:item:c:01}
  \textit{\(G\) is finite over \(S\) if and only if, for every
  \(s\in S\), the type of \(G\) at \(s\) is given by a finite commutative
  group, if and only if (when \(S\) is quasi-compact) \(G\) is of finite
  type over \(S\) and is annihilated by an integer \(n>0\);}

  \item[\((c')\)]\SGARefTarget{sga3:ix:proposition:2:1:item:cprime:01}
  \textit{\(G\) is integral over \(S\) if and only if, for every
  \(s\in S\), the type of \(G\) at \(s\) is given by a torsion
  commutative group;}

  \item[(d)]\SGARefTarget{sga3:ix:proposition:2:1:item:d:01}
  \textit{\(G\) is the unit group if and only if, for every \(s\in S\),
  the type of \(G\) at \(s\) is given by the zero commutative group;}

  \item[(e)]\SGARefTarget{sga3:ix:proposition:2:1:item:e:01}
  \textit{\(G\) is smooth over \(S\) if and only if, for every
  \(s\in S\), the type of \(G\) at \(s\) is given by a commutative group
  of finite type whose torsion subgroup has order prime to the
  characteristic of \(\kappa(s)\).}
\end{enumerate}

These statements follow from \SGARefLink{sga3:viii:proposition:2:1}{VIII~2.1}, taking into account that the
properties in question descend under faithfully flat and quasi-compact
morphisms (cf. SGA~1, VIII, or EGA~IV\textsubscript{2}, \S2).

Using \SGARefLink{sga3:viii:corollary:3:5}{VIII~3.5}, one likewise obtains:

\medskip
\noindent\SGARefTarget{sga3:ix:proposition:2:2}\textbf{Proposition 2.2.}
\textit{Let \(G\) be a group of multiplicative type and of finite type over
\(S\). For every integer \(n\ne0\), the kernel
\({}_nG=\ker(n\cdot\operatorname{id}_G)\) is a group of multiplicative
type, finite over \(S\).}

\medskip
\noindent\SGARefTarget{sga3:ix:proposition:2:3}\textbf{Proposition 2.3.}
\textit{Let \(G\) be a group of multiplicative type over the prescheme
\(S\), acting freely on the right on an \(S\)-prescheme \(X\) that is
affine over \(S\). Then:}
\begin{enumerate}
  \item[(a)]\SGARefTarget{sga3:ix:proposition:2:3:item:a:01}
  \textit{the equivalence relation defined by \(G\) in \(X\) is
  \(\mathcal M\)-effective (\SGARefLink{sga3:iv:subsection:3:4}{IV~3.4}), and \(Y=X/G\) is affine over \(S\);}

  \item[(b)]\SGARefTarget{sga3:ix:proposition:2:3:item:b:01}
  \textit{if, moreover, \(X\) is of finite presentation (respectively, of
  finite type) over \(S\), then the same is true of \(Y\).}
\end{enumerate}

\SGARefLink{sga3:ix:proposition:2:3:item:a:01}{The first assertion} follows from \SGARefLink{sga3:viii:theorem:5:1}{VIII~5.1}, which treats the case in which
\(G\) is diagonalizable, and from \SGARefLink{sga3:iv:proposition:3:5:2}{IV~3.5.2}, which allows one to reduce to
that case. Indeed, faithfully flat and quasi-compact morphisms
\(S'\to S\) are effective descent morphisms for the fibered category of
affine schemes over schemes: if \(Y'\) is affine over \(S'\) and is
equipped with a descent datum relative to \(S'\to S\), then this datum is
effective, that is, \(Y'\) comes from a scheme \(Y\) affine over \(S\)
(cf. SGA~1, \SGARefLink{sga3:viii:proposition:2:1}{VIII~2.1}).

% Source page: Exp9 component-local p. 5; complete-reader p. 651;
% original printed p. 43.
For \SGARefLink{sga3:ix:proposition:2:3:item:b:01}{the second assertion}, one is likewise reduced to the diagonalizable
case of \SGARefLink{sga3:viii:corollary:5:8}{VIII~5.8}, since the finiteness conditions under consideration
descend under faithfully flat and quasi-compact morphisms (SGA~1,
VIII~3.3 and 3.6\footnote{Editor's note: See also V, 9.1, or
EGA~IV\textsubscript{2}, 2.7.1.}). Proceeding as in VIII,
Corollaries~5.5--5.7, one deduces from
\SGARefLink{sga3:ix:proposition:2:3}{2.3}:

\medskip
\noindent\SGARefTarget{sga3:ix:corollary:2:4}\textbf{Corollary 2.4.}
\textit{Under the hypotheses of \SGARefLink{sga3:ix:proposition:2:3}{2.3},
the graph morphism}
\SGARefTarget{sga3:ix:corollary:2:4:diagram:01}
\[
  X\times_S G\longrightarrow X\times_S X
\]
\textit{is a closed immersion. For every section \(\sigma\) of \(X\)
over \(S\), the corresponding morphism}
\[
  G\longrightarrow X,\qquad g\longmapsto\sigma\cdot g,
\]
\textit{is a closed immersion.}

In particular:

\medskip
\noindent\SGARefTarget{sga3:ix:corollary:2:5}\textbf{Corollary 2.5.}
\textit{Let \(u:G\to H\) be a monomorphism of \(S\)-preschemes in groups,
with \(G\) of multiplicative type and \(H\) affine over \(S\). Then \(u\)
is a closed immersion, \(H/G=Y\) exists and is affine over \(S\), and
\(H\) is a principal homogeneous bundle over \(Y\) under \(G_Y\).}

\medskip
\noindent\SGARefTarget{sga3:ix:remark:2:6}\textbf{Remark 2.6.}
Let \(u:G\to H\) be a monomorphism of \(S\)-preschemes in groups, where
\(G\) is of multiplicative type and of finite type over \(S\), and \(H\)
is of finite presentation and separated over \(S\). One can then show
that \(u\) is a closed immersion by using \SGARefLink{sga3:viii:proposition:7:12}{VIII~7.12} (see
Remark~\SGARefLink{sga3:viii:remark:7-13:item:b}{VIII~7.13(b)}).

\medskip
\noindent\SGARefTarget{sga3:ix:proposition:2:7}\textbf{Proposition 2.7.}
\textit{Let \(u:G\to H\) be a homomorphism of \(S\)-groups of
multiplicative type, with \(H\) of finite type over \(S\). Then:}
\begin{enumerate}
  \item[(i)]\SGARefTarget{sga3:ix:proposition:2:7:item:i:01}
  \textit{\(G_0=\ker u\) is an \(S\)-group of multiplicative type; the
  equivalence relation defined by \(G_0\) in \(G\) is
  \(\mathcal M\)-effective, and hence (\SGARefLink{sga3:iv:subsection:3:4}{IV~3.4}) \(u\) factors as}
  \SGARefTarget{sga3:ix:proposition:2:7:diagram:01}
  \[
    G\longrightarrow G/G_0=I\longrightarrow H,
  \]
  \textit{where \(I\to H\) is a monomorphism\footnote{Editor's note:
  Indeed, even a closed immersion.} of \(S\)-groups. The group \(I\) is of
  multiplicative type, the equivalence relation in \(H\) defined by \(I\)
  is \(\mathcal M\)-effective, and consequently \(H_0=H/I\) exists and is
  also of multiplicative type;}

  \item[(ii)]\SGARefTarget{sga3:ix:proposition:2:7:item:ii:01}
  \textit{\(H_0\) and \(I\) are of finite type, and so is \(G_0\) when
  \(G\) is.}
\end{enumerate}

\emph{Proof.}
Proceeding as for \SGARefLink{sga3:ix:proposition:2:3}{2.3}, one is
reduced to the case in which \(G\) and \(H\) are diagonalizable; then
\SGARefLink{sga3:ix:proposition:2:7}{2.7} reduces to \SGARefLink{sga3:viii:corollary:3:4}{VIII~3.4}.

We record the following corollaries.

% Source page: Exp9 component-local p. 6; complete-reader p. 652;
% original printed p. 44.
\medskip
\noindent\SGARefTarget{sga3:ix:corollary:2:8}\textbf{Corollary 2.8.}
\begin{enumerate}
  \item[(a)]\SGARefTarget{sga3:ix:corollary:2:8:item:a:01}
  \textit{Let \(S\) be a prescheme. The category of \(S\)-groups of
  multiplicative type and of finite type is an abelian category.}

  \item[(b)]\SGARefTarget{sga3:ix:corollary:2:8:item:b:01}
  \textit{Let \(u:G\to H\) be a homomorphism in this category. It is a
  monomorphism (respectively, an epimorphism, respectively, an
  isomorphism) in this category if and only if \(u\) is a monomorphism of
  preschemes (respectively, \(u\) is faithfully flat, respectively, \(u\)
  is an isomorphism of preschemes).}
\end{enumerate}

It is enough to note that the product of two \(S\)-groups of
multiplicative type is again of multiplicative type (which is immediate),
and is plainly of finite type over \(S\) if both factors are. The rest of
\SGARefLink{sga3:ix:corollary:2:8}{2.8} follows immediately from
\SGARefLink{sga3:ix:proposition:2:7}{2.7}; the details are left to the
reader.

\medskip
\noindent\SGARefTarget{sga3:ix:corollary:2:9}\textbf{Corollary 2.9.}
\textit{Let \(u:G\to H\) be a homomorphism of \(S\)-groups of
multiplicative type and of finite type. Let \(U\) be the set of points
\(s\in S\) for which \(u_s:G_s\to H_s\) is a monomorphism
(respectively, faithfully flat, respectively, an isomorphism). Then
\(U\) is both open and closed, and the induced homomorphism
\(u_U:G_U\to H_U\) is a monomorphism (respectively, faithfully flat,
respectively, an isomorphism).}

Let \(P\) (respectively, \(Q\)) be the kernel (respectively, cokernel) of
\(u\). By \SGARefLink{sga3:ix:proposition:2:7}{2.7}, \(Q\) exists and
\(P\) and \(Q\) are of multiplicative type; moreover, formation of \(P\)
and \(Q\) commutes with every base change \(S'\to S\), in particular with
the formation of fibers. On the other hand, \(u\) is a monomorphism
(respectively, faithfully flat, respectively, an isomorphism) if and only
if \(P=0\) (respectively, \(Q=0\), respectively, \(P=Q=0\)). These
remarks reduce us to checking the following assertion: if \(R\) is an
\(S\)-group of multiplicative type, then the set \(U\) of points
\(s\in S\) for which \(R_s=0\) is open and closed, and \(R_U=0\). This is
contained in \SGARefLink{sga3:ix:remark:1:4:1:item:a}{Remark~1.4.1(a)}.

\medskip
\noindent\SGARefTarget{sga3:ix:corollary:2:10}\textbf{Corollary 2.10.}
\textit{Let \(G\) be an \(S\)-group of multiplicative type and of finite
type, and let \(H\) and \(H'\) be two subgroups of multiplicative type.}
\begin{enumerate}
  \item[(a)]\SGARefTarget{sga3:ix:corollary:2:10:item:a:01}
  \textit{\(H''=H\cap H'\) is a subgroup of \(G\) of multiplicative
  type.}

  \item[(b)]\SGARefTarget{sga3:ix:corollary:2:10:item:b:01}
  \textit{Let \(U\) be the set of \(s\in S\) for which
  \(H_s\subseteq H'_s\) (respectively, \(H_s=H'_s\)). Then \(U\) is open
  and closed, and \(H_U\subseteq H'_U\) (respectively, \(H_U=H'_U\)).}
\end{enumerate}

Here, of course, the intersection sign in \(H\cap H'\) denotes the
intersection in the functorial sense, namely \(H\times_GH'\), which is
plainly an \(S\)-subgroup of \(G\). Likewise, the inclusion sign
(respectively, the equality sign) denotes the order relation
(respectively, equality) between subfunctors of \(H\), not inclusion
(respectively, equality) of the underlying sets.

Applying \SGARefLink{sga3:ix:proposition:2:7}{2.7} first to the inclusion
morphism \(H\to G\), one finds that \(H\) is of finite type; the same
holds for \(H'\). It then follows from
\SGARefLink{sga3:ix:corollary:2:8}{2.8} that \(H\cap H'\) is of
multiplicative type and of finite type. (Note that the canonical functor
from the category considered in \SGARefLink{sga3:ix:corollary:2:8}{2.8}
to the category \((\mathbf{Sch})/S\) commutes with finite inverse limits.)

Formation of \(H\cap H'\) commutes with base extension, in particular
with formation of the fibers. Moreover,
\(H\subseteq H'\) (respectively, \(H=H'\)) is equivalent to
\(H''=H\) (respectively, to \(H''=H\) and \(H''=H'\)). Consider the
canonical homomorphisms \(H''\to H\) and \(H''\to H'\). Then \(U\) is the
set of points \(s\in S\) for which the induced homomorphism
\(H''_s\to H_s\) is an isomorphism (respectively, for which
\(H''_s\to H_s\) and \(H''_s\to H'_s\) are isomorphisms).

% Source page: Exp9 component-local p. 7; complete-reader p. 653;
% original printed pp. 45--46.
By \SGARefLink{sga3:ix:corollary:2:9}{2.9}, it follows that \(U\) is open
and closed, and that \(H''_U\to H_U\) (respectively,
\(H''_U\to H_U\) and \(H''_U\to H'_U\)) are isomorphisms. This gives the
desired conclusion.

\medskip
\noindent\SGARefTarget{sga3:ix:proposition:2:11}\textbf{Proposition 2.11.}
\textit{Let \(S\) be a prescheme, let \(G\) be an \(S\)-group of
multiplicative type and of finite type over \(S\), let \(H\) be a subgroup
of \(G\) of multiplicative type, and put \(K=G/H\), a quotient group of
\(G\) of multiplicative type.}
\begin{enumerate}
  \item[(i)]\SGARefTarget{sga3:ix:proposition:2:11:item:i:01}
  \textit{Suppose \(G\) is trivial, that is, diagonalizable. There is a
  partition of \(S\) into open and closed subsets \(S_i\) such that, for
  every \(i\), \(H_{S_i}\) and \(K_{S_i}\) are diagonalizable. In
  particular, if \(S\) is connected, \(H\) and \(K\) are
  diagonalizable.}

  \item[(ii)]\SGARefTarget{sga3:ix:proposition:2:11:item:ii:01}
  \textit{\SGARefLink{sga3:ix:proposition:2:11:item:i:01}{The same assertion} holds with ``diagonalizable'' replaced by
  ``isotrivial,'' provided that \(S\) is connected, locally connected, or
  quasi-compact.}

  \item[(iii)]\SGARefTarget{sga3:ix:proposition:2:11:item:iii:01}
  \textit{If \(G\) is locally trivial (respectively, locally isotrivial,
  respectively, quasi-isotrivial), then the same is true of \(K\) and
  \(H\).}
\end{enumerate}

\emph{Proof.}
\SGARefLink{sga3:ix:proposition:2:11:item:i:01}{(i)} By hypothesis, \(G=D_S(M)\), where \(M\) is a commutative group of
finite type. For each quotient group \(M_i\) of \(M\), let
\(H_i=D_S(M_i)\) be the corresponding diagonalizable subgroup of \(G\)
(\SGARefLink{sga3:viii:theorem:3:1}{VIII~3.1}). Let \(S_i\) be the set of points \(s\in S\) for which
\(H_s=(H_i)_s\). By \SGARefLink{sga3:ix:corollary:2:10}{2.10}, \(S_i\)
is open and closed, and \(H_{S_i}=(H_i)_{S_i}\); thus \(H_{S_i}\) is
diagonalizable, and hence so is \(K_{S_i}\) (\SGARefLink{sga3:viii:theorem:3:1}{VIII~3.1}). The \(S_i\) are
plainly pairwise disjoint, and we claim that they cover \(S\). Indeed, for
every \(s\), the group \(H_s\), being a subgroup of the diagonalizable
group \(G_s\), is diagonalizable (cf. \SGARefLink{sga3:ix:proposition:8:1}{8.1} below\footnote{Editor's note:
The erroneous reference to IV~4.7.5 has been corrected. Note that
\SGARefLink{sga3:ix:section:8}{Section~8} of the present Expos\'e is independent of Sections~\SGARefLink{sga3:ix:section:3}{3}--\SGARefLink{sga3:ix:section:7}{7}.}).
Hence \(H_s\) has the form \(D_{\kappa(s)}(M_i)\), by \SGARefLink{sga3:viii:corollary:1:5}{VIII~1.5} and
\SGARefLink{sga3:ix:theorem:3:2:item:b}{3.2(b)}. Restricting to the family of nonempty \(S_i\), assertion \SGARefLink{sga3:ix:proposition:2:11:item:i:01}{(i)}
follows.

\SGARefLink{sga3:ix:proposition:2:11:item:ii:01}{(ii)} By hypothesis, there exists a finite surjective \'etale morphism
\(S'\to S\) such that \(G_{S'}\) is diagonalizable. Every point of \(S'\)
therefore has an open and closed neighborhood \(U'\) such that \(H_{U'}\)
and \(K_{U'}\) are diagonalizable. The image \(U\) of \(U'\) in \(S\) is
open and closed, and \(S'\to S\) still induces a finite surjective
\'etale
morphism \(U'\to U\). Thus every point \(s\in S\) has an open and closed
neighborhood \(U\) such that \(H_U\) and \(K_U\) are isotrivial.
Assertion \SGARefLink{sga3:ix:proposition:2:11:item:ii:01}{(ii)} follows at once.

\SGARefLink{sga3:ix:proposition:2:11:item:iii:01}{(iii)} This follows at once from \SGARefLink{sga3:ix:proposition:2:11:item:i:01}{(i)}, \SGARefLink{sga3:ix:proposition:2:11:item:ii:01}{(ii)}, and the definitions. Moreover,
the quasi-isotrivial case will also follow from the more general fact
``finite type \(\Longrightarrow\) quasi-isotrivial,'' announced in
\SGARefLink{sga3:ix:comments:1:2}{1.2}
(cf. \SGARefLink{sga3:x:corollary:4-5}{X~4.5}).

\SGARefTarget{sga3:ix:section:3}\section*{3. Infinitesimal Properties: Lifting and Conjugation Theorems}
\addcontentsline{toc}{section}{3. Infinitesimal Properties: Lifting and Conjugation Theorems}

The fundamental infinitesimal properties of groups of multiplicative type
follow from \SGARefLink{sga3:ix:theorem:3:1}{the following theorem}.

\medskip
\noindent\SGARefTarget{sga3:ix:theorem:3:1}\textbf{Theorem 3.1.}
\textit{Let \(S\) be an affine scheme, let \(H\) be an \(S\)-group of
multiplicative type, and let \(\mathscr F\) be a quasi-coherent sheaf on
\(S\) on which \(H\) acts (\SGARefLink{sga3:i:subsection:4-7}{I~4.7}). Then}
\[
  \mathrm H^i(H,\mathscr F)=0\qquad\text{for }i>0,
\]
\textit{where \(\mathrm H^i\) denotes the Hochschild cohomology studied
in Expos\'e~I, \S5.}

% Source page: Exp9 component-local p. 8; complete-reader p. 654;
% original printed pp. 47--48.
Indeed, by loc.\ cit., \SGARefLink{sga3:ix:corollary:5:3}{5.3}, Hochschild cohomology is expressed, in terms
of the affine rings \(A\) and \(B\) of \(S\) and \(G\), and the
\(A\)-module \(M\) defining \(\mathscr F\), as the cohomology of a complex
of \(A\)-modules \(C^\bullet(H,M)\), whose formation commutes with every
base change \(A\to A'\). Consequently, for a base change \(A\to A'\) with
\(A'\) flat over \(A\), one has
\[
  \mathrm H^i(H',\mathscr F')
  =\mathrm H^i(H,\mathscr F)\otimes_AA'.
\]
If \(A'\) is moreover faithfully flat over \(A\), it is therefore enough,
in order to prove that \(\mathrm H^i(H,\mathscr F)\) vanishes, to prove
that \(\mathrm H^i(H',\mathscr F')\) vanishes. This observation reduces
us to checking \SGARefLink{sga3:ix:theorem:3:1}{3.1} in the case where
\(G\) is diagonalizable; in that case it was proved in \SGARefLink{sga3:i:theorem:5-3-3}{I~5.3.3}.

Using the results of Expos\'e~III, we shall derive from
\SGARefLink{sga3:ix:theorem:3:1}{3.1} several consequences of a geometric
nature.

\medskip
\noindent\SGARefTarget{sga3:ix:theorem:3:2}\textbf{Theorem 3.2.}
\textit{Let \(S\) be an affine scheme, let \(S_0\) be an affine subscheme
defined by an ideal \(J\), let \(H\) be a group of multiplicative type
over \(S\), and let \(G\) be an arbitrary group prescheme over \(S\).
Let \(u,v:H\to G\) be two group homomorphisms, and let
\(u_0,v_0:H_0\to G_0\) be the homomorphisms obtained from \(u,v\) by the
base change \(S_0\to S\). If \(J^2=0\) and \(u_0=v_0\), then there exists
\(g\in G(S)\) such that}
\[
  v=\operatorname{int}(g)\circ u
  \qquad\textit{and}\qquad
  g_0=e.
\]

This follows from \SGARefLink{sga3:iii:theorem:2-1:item:ii}{III~2.1(ii)} and
\SGARefLink{sga3:ix:theorem:3:1}{3.1} (the vanishing of
\(\mathrm H^1\)).

\medskip
\noindent\SGARefTarget{sga3:ix:corollary:3:3}\textbf{Corollary 3.3.}
\textit{Under the preliminary hypotheses of
\SGARefLink{sga3:ix:theorem:3:2}{3.2}, suppose in addition that \(G\) is
smooth over \(S\), but assume only that \(J\) is nilpotent (not
necessarily square-zero). If there exists \(g_0\in G_0(S_0)\) such that
\(v_0=\operatorname{int}(g_0)\circ u_0\), then \(g_0\) lifts to an element
\(g\in G(S)\) such that \(v=\operatorname{int}(g)\circ u\).}

By induction on an integer \(n\) such that \(J^n=0\), one is reduced to
the case in which \(J\) is square-zero. Since \(G\) is smooth over \(S\),
one may lift \(g_0\) to an element \(g'\) of \(G(S)\). Replacing \(v\) by
\(v'=\operatorname{int}(g')\circ u\), one is then reduced to the
situation of \SGARefLink{sga3:ix:theorem:3:2}{3.2}.

\medskip
\noindent\SGARefTarget{sga3:ix:corollary:3:4}\textbf{Corollary 3.4.}
\textit{Under the preliminary hypotheses of
\SGARefLink{sga3:ix:theorem:3:2}{3.2}, with \(J\) nilpotent, suppose in
addition that \(u\) is central (for example, that \(G\) is commutative).
Then \(u_0=v_0\) implies \(u=v\).}

Indeed, the reduction to the case in which \(J\) is square-zero is again
immediate, and then it is enough to apply
\SGARefLink{sga3:ix:theorem:3:2}{3.2}. In particular:

\medskip
\noindent\SGARefTarget{sga3:ix:corollary:3:5}\textbf{Corollary 3.5.}
\textit{Let \(S,S_0,H,G\) be as in
\SGARefLink{sga3:ix:theorem:3:2}{3.2}, with \(J\) nilpotent. Let
\(u:H\to G\) be a homomorphism of \(S\)-groups such that
\(u_0:H_0\to G_0\) is the unit homomorphism. Then \(u\) is the unit
homomorphism.}

\medskip
\noindent\SGARefTarget{sga3:ix:theorem:3:6}\textbf{Theorem 3.6.}
\textit{Let \(S\) be an affine scheme, let \(S_0\) be a subscheme defined
by a nilpotent ideal \(J\), let \(H\) be a group of multiplicative type
over \(S\), let \(G\) be a group prescheme smooth over \(S\), and let
\(u_0:H_0\to G_0\) be a homomorphism of the \(S_0\)-groups obtained by
base change \(S_0\to S\). Then there exists a homomorphism
\(u:H\to G\) of \(S\)-groups that lifts \(u_0\). Moreover, by
\SGARefLink{sga3:ix:corollary:3:3}{3.3}, any two such liftings \(u,u'\)
are conjugate by an element of \(G(S)\) whose reduction in \(G(S_0)\) is
the identity element.}

% Source page: Exp9 component-local p. 9; complete-reader p. 655;
% original printed p. 49.
This follows from \SGARefLink{sga3:iii:theorem:2-1}{III~2.1} and \SGARefLink{sga3:ix:proposition:2:3}{2.3} and from
\SGARefLink{sga3:ix:theorem:3:1}{3.1}: the vanishing of
\(\mathrm H^2\) gives the existence of the lifting of \(u_0\), and the
vanishing of \(\mathrm H^1\) gives uniqueness up to conjugation.

One may prove in the same way the following variants of
\SGARefLink{sga3:ix:theorem:3:2}{3.2}--%
\SGARefLink{sga3:ix:theorem:3:6}{3.6}. In fact, these variants imply the
preceding statements when applied to the graph subgroups of the
homomorphisms considered there.

\medskip
\noindent\SGARefTarget{sga3:ix:theorem:3:2:bis}\textbf{Theorem 3.2 bis.}
\textit{Let \(S\) be an affine scheme, let \(S_0\) be a subscheme defined
by an ideal \(J\), let \(G\) be an \(S\)-group prescheme, and let \(H,H'\)
be subgroup preschemes of multiplicative type. Let \(G_0,H_0,H'_0\) be
the groups over \(S_0\) obtained by base change \(S_0\to S\). If
\(J^2=0\) and \(H_0=H'_0\), then there exists \(g\in G(S)\) such that}
\[
  H'=\operatorname{int}(g)(H)
  \qquad\textit{and}\qquad
  g_0=e.
\]

\medskip
\noindent\SGARefTarget{sga3:ix:corollary:3:3:bis}\textbf{Corollary 3.3 bis.}
\textit{Under the preliminary hypotheses of
\SGARefLink{sga3:ix:theorem:3:2:bis}{3.2 bis}, suppose in addition that
\(G\) is smooth over \(S\), and that \(J\) is nilpotent (not necessarily
square-zero). If \(g_0\in G_0(S_0)\) satisfies
\(H'_0=\operatorname{int}(g_0)(H_0)\), then \(g_0\) lifts to an element
\(g\in G(S)\) such that \(H'=\operatorname{int}(g)(H)\).}

\medskip
\noindent\SGARefTarget{sga3:ix:corollary:3:4:bis}\textbf{Corollary 3.4 bis.}
\textit{Under the preliminary hypotheses of
\SGARefLink{sga3:ix:theorem:3:2:bis}{3.2 bis}, suppose that \(J\) is
nilpotent and that \(H\) is central (for example, that \(G\) is
commutative). Then \(H_0=H'_0\) implies \(H=H'\).}

\medskip
\noindent\SGARefTarget{sga3:ix:theorem:3:6:bis}\textbf{Theorem 3.6 bis.}
\textit{Let \(S\) be an affine scheme, let \(S_0\) be a subscheme defined
by a nilpotent ideal \(J\), let \(G\) be an \(S\)-group prescheme smooth
over \(S\), and let \(H_0\) be a subgroup prescheme of
\(G_0=G\times_SS_0\), of multiplicative type. Then:}
\begin{enumerate}
  \item[(a)]\SGARefTarget{sga3:ix:theorem:3:6:bis:item:a:01}
  \textit{there exists a subgroup prescheme \(H\) of \(G\), flat over
  \(S\), such that \(H\times_SS_0=H_0\);}

  \item[\SGAInlineRefTarget{sga3:ix:theorem:3:2:item:b}(b)]\SGARefTarget{sga3:ix:theorem:3:6:bis:item:b:01}
  \textit{such an \(H\) is necessarily of multiplicative type;}

  \item[(c)]\SGARefTarget{sga3:ix:theorem:3:6:bis:item:c:01}
  \textit{finally, by
  \SGARefLink{sga3:ix:theorem:3:2:bis}{3.2 bis}, any two such liftings
  \(H,H'\) of \(H_0\) are conjugate by an element \(g\in G(S)\) whose
  reduction in \(G(S_0)\) is the identity element.}
\end{enumerate}

Let us indicate how one can prove
\SGARefLink{sga3:ix:theorem:3:2:bis}{3.2 bis} (of which
\SGARefLink{sga3:ix:corollary:3:3:bis}{3.3 bis} and
\SGARefLink{sga3:ix:corollary:3:4:bis}{3.4 bis} are immediate
consequences) and assertion
\SGARefLink{sga3:ix:theorem:3:6:bis:item:a:01}{(a)} of
\SGARefLink{sga3:ix:theorem:3:6:bis}{3.6 bis}. The latter follows from
\SGARefLink{sga3:ix:theorem:3:1}{3.1} (the vanishing of
\(\mathrm H^2\)) and \SGARefLink{sga3:iii:corollary:4-37:item:ii}{III~4.37(ii)}. Similarly,
\SGARefLink{sga3:ix:theorem:3:2:bis}{3.2 bis} follows from
\SGARefLink{sga3:ix:theorem:3:1}{3.1} (the vanishing of
\(\mathrm H^1\)) and \SGARefLink{sga3:iii:corollary:4-37:item:i}{III~4.37(i)}, at least when \(G\) and \(H\) are
smooth over \(S\).

Using the results of the following Expos\'e, however, one can very simply
deduce \SGARefLink{sga3:ix:theorem:3:2:bis}{3.2 bis} and
\SGARefLink{sga3:ix:theorem:3:6:bis}{3.6 bis}, in the general form stated
here, from \SGARefLink{sga3:ix:theorem:3:2}{3.2} and
\SGARefLink{sga3:ix:theorem:3:6}{3.6}, respectively. For
\SGARefLink{sga3:ix:theorem:3:2:bis}{3.2 bis}, it is enough to note that,
by \SGARefLink{sga3:x:proposition:2:1}{X~2.1}, \(H\) and \(H'\) are isomorphic, which reduces us to
\SGARefLink{sga3:ix:theorem:3:2}{3.2}.

For \SGARefLink{sga3:ix:theorem:3:6:bis}{3.6 bis}, note that \(H_0\) is
necessarily of finite type over \(S_0\):\footnote{Editor's note: The
original said, ``by \SGARefLink{sga3:ix:proposition:2:1:item:b:01}{2.1(b)}, since its fibers are.'' The editors did not
understand this argument and have replaced it by the argument that
follows.} indeed, as a subgroup prescheme of \(G_0\), which is smooth
over \(S_0\), the group \(H_0\) is locally of finite type over \(S_0\).
But \(H_0\) is affine over \(S_0\) by
\SGARefLink{sga3:ix:proposition:2:1:item:a:01}{2.1(a)}, and is therefore
of finite type over \(S_0\). By \SGARefLink{sga3:x:corollary:4-5}{X~4.5}, \(H_0\) is consequently
quasi-isotrivial, and hence comes, by \SGARefLink{sga3:x:proposition:2:1}{X~2.1}, from a group \(H\) of
multiplicative type over \(S\).\footnote{Editor's note: The sentence that
follows has been expanded.} By
\SGARefLink{sga3:ix:theorem:3:6}{3.6}, there exists a homomorphism of
\(S\)-groups \(u:H\to G\) lifting the immersion \(H_0\hookrightarrow
G_0\). Since \(H\) and \(H_0\) (respectively, \(G\) and \(G_0\)) have the
same underlying topological space, and since, for every \(h\in H\), the
morphism

% Source page: Exp9 component-local p. 10; complete-reader p. 656;
% original printed p. 50.
\[
  \mathscr O_{G,u(h)}\longrightarrow\mathscr O_{H,h}
\]
is surjective (because it is so after reduction modulo the nilpotent ideal
\(J\)), the morphism \(u\) is also an immersion (cf. EGA~I, 4.2.2).

Finally, every lifting \(H\) of \(H_0\) to a flat subgroup of \(G\) is
necessarily of multiplicative type by \SGARefLink{sga3:x:corollary:2:3}{X~2.3},\footnote{Editor's note:
\SGARefLink{sga3:x:corollary:2:3}{X~2.3} depends essentially on \SGARefLink{sga3:ix:theorem:3:6}{Theorem~3.6} of the present Expos\'e.} which
also proves assertion
\SGARefLink{sga3:ix:theorem:3:6:bis:item:b:01}{(b)} of
\SGARefLink{sga3:ix:theorem:3:6:bis}{3.6 bis}. The reader may verify that
the results \SGARefLink{sga3:ix:theorem:3:2:bis}{3.2 bis}--%
\SGARefLink{sga3:ix:theorem:3:6:bis}{3.6 bis} are not used in
Expos\'e~X,
so there is no vicious circle.

\medskip
\noindent\SGARefTarget{sga3:ix:proposition:3:8}\textbf{Proposition 3.8.}
\footnote{Editor's note: The numbering of the original has been retained;
there is no 3.7.}
\textit{Let \(S\) be a prescheme, let \(S_0\) be a closed subscheme
defined by a locally nilpotent ideal \(J\), let \(G\) be an \(S\)-group
of multiplicative type, and let \(X\) be an \(S\)-prescheme locally of
finite presentation. Suppose that \(G\) acts on \(X\) in such a way that
\(G_0=G\times_SS_0\) acts trivially on \(X_0=X\times_SS_0\). Then \(G\)
acts trivially on \(X\).}

The proof is the proof of \SGARefLink{sga3:iii:corollary:2-4:item:b}{III~2.4(b)}. One may also reduce to loc.\ cit. by
observing that \SGARefLink{sga3:ix:proposition:3:8}{3.8} immediately
reduces to the case in which \(G\) is diagonalizable, which is the case
treated there.

\medskip
\noindent\SGARefTarget{sga3:ix:corollary:3:9}\textbf{Corollary 3.9.}
\textit{Let \(S,S_0\) be as above, and let \(u:G\to H\) be a homomorphism
of \(S\)-groups, with \(G\) of multiplicative type and \(H\) locally of
finite presentation over \(S\). Suppose that the homomorphism
\(u_0:G_0\to H_0\) obtained by base change \(S_0\to S\) is central. Then
\(u\) is central.}

Indeed, it is enough to apply
\SGARefLink{sga3:ix:proposition:3:8}{3.8} with \(X=H\), letting \(G\) act
on \(H\) by
\[
  (g,h)\longmapsto\operatorname{int}(u(g))\cdot h
  =u(g)h\,u(g)^{-1}.
\]

% Hard stop: immediately before Section 4, Exp9 component-local p. 10
% (complete-reader p. 656; original printed p. 50).

% END INCLUDED SGA3_English_current_buildable_source_20260823/source/inputs/IX/tex/components/02_expose_IX_sections_2_3.tex

% BEGIN INCLUDED SGA3_English_current_buildable_source_20260823/source/inputs/IX/tex/components/03_expose_IX_sections_4_5.tex
% SGA 3, Expose IX, English reconstruction.
% Scope: Sections 4 and 5 in full.
% Authority: Polo--Gille Exp9-8nov09.pdf, component-local pp. 10--22
% (partial), complete-reader pp. 656--668 (partial), printed pp. 50--63.
% Start cursor: "4. Le theoreme de densite" on component-local p. 10.
% End cursor: immediately before
% "6. Monomorphismes des groupes de type multiplicatif, et factorisation"
% on component-local p. 22.
% Jacob Reinhold's CC BY 4.0 Markdown lineage was used only as a phrasing
% comparison. The Polo--Gille PDF controls prose, formulas, numbering,
% notes, references, and diagrams.

\setcounter{footnote}{10}

% Source page: Exp9 component-local p. 10; complete-reader p. 656.
\SGARefTarget{sga3:ix:section:4}\section*{4. The Density Theorem}
\addcontentsline{toc}{section}{4. The Density Theorem}

This theorem (cf. \SGARefLink{sga3:ix:theorem:4:7}{4.7 below})%
\begingroup
\renewcommand{\thefootnote}{\(\ast\)}%
\footnote{All the results from
\SGARefLink{sga3:ix:definition:4:1}{4.1} through
\SGARefLink{sga3:ix:corollary:4:6}{4.6} are contained in EGA
IV\(_3\), 11.10, to which we refer the reader for a systematic account of
the notion of schematic density.}%
\addtocounter{footnote}{-1}%
\endgroup,
together with the theorem\footnote{Editor's note: the ``algebraicity''
theorem \SGARefLink{sga3:ix:theorem:7:1}{7.1}.} of \SGARefLink{sga3:ix:section:7}{Section~7}, will be the
essential tool in the present expos\'e and the two following ones for
passing from the infinitesimal properties of groups of multiplicative type
developed above to their ``finite'' properties.

\medskip
\noindent\SGARefTarget{sga3:ix:definition:4:1}\textbf{Definition 4.1.}
\textit{Let \(X\) be a prescheme. A family
\((Z_i)_{i\in I}\) of subpreschemes of \(X\) is said to be
\emph{schematically dense} if, for every open subset \(U\) of \(X\) and
every closed subprescheme \(U_0\) of \(U\) containing all the
\(Z_i\cap U\), one has \(U_0=U\).}

\textit{A subprescheme \(Z\) of \(X\) is said to be
\emph{schematically dense in \(X\)} if the one-member family consisting
of \(Z\) is schematically dense.}

It is immediate (cf. EGA IV\(_3\), 11.10.1) that the definition is
equivalent to saying that, for every open subset \(U\) of \(X\), every
section \(f\) of \(\mathscr O_U\) that vanishes on each \(Z_i\cap U\)
is zero. Equivalently, the intersection of the kernels of the canonical
homomorphisms
\SGARefTarget{sga3:ix:definition:4:1:display:01}
\[
  u_i:\mathscr O_X\longrightarrow (v_i)_*(\mathscr O_{Z_i})
\]
is zero, where \(v_i:Z_i\to X\) is the canonical immersion.%
\footnote{Editor's note: all the results on schematic density stated in
EGA IV\(_3\), \S11.10 follow from the results on ``separating families of
homomorphisms'' proved in \emph{loc.\ cit.}, \S11.9, to which some of the
editor's notes below will refer.}
When \(X\) lies over a prescheme \(S\), this is also equivalent to the
following assertion: for every open subset \(U\) of \(X\), every
\(S\)-prescheme \(Y\) separated over \(S\), and every pair of morphisms
\(u,v:U\to Y\) that coincide on the \(Z_i\cap U\), one has \(u=v\).
Indeed, \(u=v\) is equivalent to \(U_0=U\), where \(U_0\) is the inverse
image of the diagonal in \(Y\times_S Y\) under the \(S\)-morphism
\(U\to Y\times_S Y\) defined by \((u,v)\). The diagonal is a closed
subprescheme of \(Y\times_S Y\), so \(U_0\) is a closed subprescheme of
\(U\) containing all the \(Z_i\cap U\). Thus schematic density gives
\(U_0=U\), hence \(u=v\). Conversely, take
\(Y=\operatorname{Spec}\mathscr O_S[T]\).

In the terminology introduced at the end of EGA I, 9.5, saying that the
subprescheme \(Z\) is schematically dense in \(X\) also means that \(X\)
is the \emph{closure subprescheme} of \(Z\) in \(X\).

% Source page: Exp9 component-local p. 11; complete-reader p. 657.
\medskip
\noindent\SGARefTarget{sga3:ix:lemma:4:2}\textbf{Lemma 4.2.}
\textit{Let \(X\) be a flat \(S\)-prescheme, and let
\((Z_i)_{i\in I}\) be a family of subpreschemes of \(X\), flat over
\(S\). Let \(S_0\) be a subprescheme of \(S\) defined by a nilpotent
ideal \(\mathscr J\), and suppose that the modules
\(\mathscr J^n/\mathscr J^{n+1}\) are locally free over \(S_0\).
Let \(X_0,Z_{i0}\) be obtained from \(X,Z_i\) by the base change
\(S_0\to S\). If \((Z_{i0})_{i\in I}\) is schematically dense in
\(X_0\), then \((Z_i)_{i\in I}\) is schematically dense in \(X\).}

\emph{Proof.}
Suppose \(\mathscr J^{n+1}=0\), where \(n\geq0\), and argue by induction
on \(n\); the assertion is trivial for \(n=0\). Define
\(S_m,X_m,Z_{im}\) by reduction modulo \(\mathscr J^{m+1}\), as usual.
The induction hypothesis already implies that
\((Z_{i,n-1})_{i\in I}\) is schematically dense in \(X_{n-1}\).
After replacing \(X\) by an induced open subset, it is enough to prove
that every section \(f\) of \(\mathscr O_X\) that vanishes on all the
\(Z_i\) is zero.

The section \(f_{n-1}\) of
\[
  \mathscr O_{X_{n-1}}
  =\mathscr O_X/\mathscr J^n\mathscr O_X
\]
defined by \(f\) vanishes on the \(Z_{i,n-1}\), hence is zero.
Consequently, \(f\) is a section of
\(\mathscr J^n\mathscr O_X\). Since \(X\) is flat over \(S\), there is
an isomorphism
\SGARefTarget{sga3:ix:lemma:4:2:display:01}
\[
  \mathscr J^n\mathscr O_X
  \xleftarrow{\ \sim\ }
  \mathscr E\otimes_{\mathscr O_{S_0}}\mathscr O_{X_0},
  \qquad
  \mathscr E=\mathscr J^n=\mathscr J^n/\mathscr J^{n+1}.
\]
Similarly, since each \(Z_i\) is flat over \(S\), the restriction
\(f_i\) of \(f\) to \(Z_i\) may be regarded as a section of
\SGARefTarget{sga3:ix:lemma:4:2:display:02}
\[
  \mathscr J^n\mathscr O_{Z_i}
  \xleftarrow{\ \sim\ }
  \mathscr E\otimes_{\mathscr O_{S_0}}\mathscr O_{Z_{i0}}.
\]
By hypothesis, all the \(f_i\) are zero. The module \(\mathscr E\) is
locally free by assumption, and therefore so is
\[
  \mathscr F
  =\mathscr E\otimes_{\mathscr O_{S_0}}\mathscr O_{X_0}.
\]
Thus \(f\) is a section of the locally free module \(\mathscr F\) on
\(X_0\) whose restriction to every \(Z_{i0}\) is zero. Since the family
\((Z_{i0})_{i\in I}\) is schematically dense in \(X_0\), it follows at
once that \(f=0\). \qed

% Source page: Exp9 component-local p. 12; complete-reader p. 658.
\medskip
\noindent\SGARefTarget{sga3:ix:corollary:4:3}\textbf{Corollary 4.3.}
\textit{Let \(X\) be a locally Noetherian \(S\)-prescheme, flat over
\(S\), and let \((Z_i)_{i\in I}\) be a family of subpreschemes of \(X\),
flat over \(S\). Suppose that, for every \(s\in S\), the family of fibers
\((Z_{is})_{i\in I}\) is schematically dense in \(X_s\). Then
\((Z_i)_{i\in I}\) is schematically dense in \(X\).}

\emph{Proof.}
After replacing \(X\) by an induced open subset, it is enough to prove
that every section \(f\) of \(\mathscr O_X\) whose restriction to all
the \(Z_i\) is zero is itself zero.%
\footnote{Editor's note: the original has been expanded in what follows.}
For this, it suffices to show that the image of \(f\) in
\(\mathscr O_{X,x}\) is zero for every \(x\in X\). Let
\(\mathfrak m_x\) be the maximal ideal of \(\mathscr O_{X,x}\), and let
\(\mathfrak m_s\) be that of \(\mathscr O_{S,s}\), where \(s\) is the
image of \(x\) in \(S\). Since \(\mathscr O_{X,x}\) is Noetherian,
\[
  \bigcap_{n\in\mathbb N}\mathfrak m_x^n=0,
  \qquad\text{and hence, a fortiori,}\qquad
  \bigcap_{n\in\mathbb N}\mathscr O_{X,x}\mathfrak m_s^n=0.
\]
It is therefore enough to show, for every integer \(n\), that the section
induced by \(f\) on
\[
  X\times_S
  \operatorname{Spec}\bigl(\mathscr O_{S,s}/\mathfrak m_s^{n+1}\bigr)
\]
is zero. By hypothesis,
\[
  Z_{is}=Z_i\otimes_{\mathscr O_{S,s}}\kappa(s)
\]
is a schematically dense family in
\[
  X_s=X\otimes_{\mathscr O_{S,s}}\kappa(s).
\]
We are thus reduced to the case in which \(S\) is the spectrum of a
local ring whose maximal ideal \(\mathscr J\) is nilpotent,
\(S_0=\operatorname{Spec}\kappa(s)\), and the hypotheses of
\SGARefLink{sga3:ix:lemma:4:2}{Lemma~4.2} hold. The conclusion follows.
\qed

\medskip
\noindent\SGARefTarget{sga3:ix:lemma:4:4}\textbf{Lemma 4.4.}
\textit{Let \(S\) be a locally Noetherian prescheme, let \(X\) be a
locally Noetherian \(S\)-prescheme, flat over \(S\), and let
\((Z_i)_{i\in I}\) be a family of subpreschemes of \(X\), flat over
\(S\). Suppose that, for every \(s\in S\), the family of fibers
\((Z_{is})_{i\in I}\) is schematically dense in \(X_s\).
Then, for every base change \(S'\to S\), with \(S'\) not necessarily
locally Noetherian, the family \((Z_i')_{i\in I}\) is schematically
dense in \(X'\). In other words, \((Z_i)_{i\in I}\) is universally
schematically dense in \(X\) relative to \(S\) (cf. EGA IV\(_3\),
Definition 11.10.8).}

\emph{Proof.}%
\footnote{Editor's note: for another proof, reducing to the case in which
\(S'\) is locally Noetherian and using
\SGARefLink{sga3:ix:corollary:4:3}{4.3} and
\SGARefLink{sga3:ix:lemma:4:5}{4.5} as here, see EGA
IV\(_3\), 11.9.16 and 11.9.12. In the last line of the proof of 11.9.16,
replace 11.9.5 by 11.9.12.}
Let \(f'\) be a section of \(\mathscr O_{X'}\) over an open subset
\(U'\) of \(X'\), vanishing on the \(Z_i'\); we must show that \(f'=0\).
Fix \(s'\in S'\). It is enough to prove that \(f'\) vanishes at every
point of \(U'\) above \(s'\). If \(s\) is the image of \(s'\) in \(S\),
then, after replacing \(S,S'\) by the spectra of their local rings at
\(s,s'\), we may suppose that \(S,S'\) are local and that \(s,s'\) are
their closed points. We may also suppose that \(X\) is affine, so that
\SGARefTarget{sga3:ix:lemma:4:4:display:01}
\[
  \begin{gathered}
    S=\operatorname{Spec}(A),\qquad
    S'=\operatorname{Spec}(A'),\\
    X=\operatorname{Spec}(B),\qquad
    X'=\operatorname{Spec}(B'),\qquad
    B'=B\otimes_A A',
  \end{gathered}
\]
where \(A,A'\) are local, \(A\to A'\) is a local homomorphism, and
\(A,B\) are Noetherian. We may further suppose that \(U'\) is affine,
of the form \(X'_{g'}\), with \(g'\in B'\).

For every \(A\)-subalgebra \(A''\) of \(A'\), put
\[
  S''=\operatorname{Spec}(A'')
\]
and let \(X'',Z_i''\) be obtained from \(X,Z_i\) by the base change
\(S''\to S\). We obtain the diagram
\SGARefTarget{sga3:ix:lemma:4:4:diagram:01}
\[
\begin{tikzcd}[column sep=large,row sep=large]
  % The source visibly prints the n-subscripts in the upper middle and
  % upper right nodes; they are retained despite the surrounding prose.
  Z^i \arrow[d] & Z_n^{i''} \arrow[l]\arrow[d]
    & Z_n^{i'} \arrow[l]\arrow[d] \\
  X \arrow[d] & X'' \arrow[l]\arrow[d] & X' \arrow[l]\arrow[d] \\
  S & S'' \arrow[l] & S' \arrow[l].
\end{tikzcd}
\]

% Source page: Exp9 component-local p. 13; complete-reader p. 659.
In what follows, we restrict to those \(A''\) that are localizations of
\(A\)-subalgebras \(A_1''\) of \(A'\), of finite type over \(A\), at
\(\mathfrak m'\cap A_1''\), where \(\mathfrak m'\) is the maximal ideal
of \(A'\). Thus the homomorphisms \(A\to A''\to A'\) are local. These
\(A''\) form an increasing filtered family of subalgebras whose union,
and hence whose direct limit, is \(A'\). Likewise,
\[
  B'=\varinjlim_{A''}\bigl(B\otimes_A A''\bigr).
\]
There therefore exist such an \(A''\) and an element
\(g''\in B''=B\otimes_A A''\) whose image in \(B'\) is \(g'\).

By \SGARefLink{sga3:ix:lemma:4:5}{Lemma~4.5 below}, the property that
\((Z_{is})_{i\in I}\) be schematically dense in \(X_s\) is preserved by
every base change \(S''\to S\).%
\footnote{Editor's note: the preceding sentence has been added, and the
original has been expanded in what follows.}
Consequently, because each \(S''\) is locally Noetherian, replacing
\(S\) by a suitable \(S''\) and \(X,Z_i\) by \(X'',Z_i''\) preserves
the hypotheses of \SGARefLink{sga3:ix:lemma:4:4}{Lemma~4.4}.

We may therefore suppose that \(g'\) comes from an element \(g\in B\).
After replacing \(X\) by the open subset \(X_g=U\), we may suppose that
\(U'=X'\). In the same way, we may suppose that \(f'\) comes from a
section \(f\) of \(\mathscr O_X\) over \(X\).

Let \(Y=V(f)\) be the subscheme of \(X\) defined by \(f\), and put
\[
  Y_i=Z_i\times_XY.
\]
Thus \(Y_i\) is the closed subscheme \(V(f_i)\) of \(Z_i\), where \(f_i\)
is the section of \(\mathscr O_{Z_i}\) induced by \(f\). Denoting by
\(Y',Y_i'\) the \(S'\)-schemes obtained by base change, we have
\SGARefTarget{sga3:ix:lemma:4:4:display:02}
\[
  Y'=V(f'),
  \qquad
  Y_i'=Z_i'\times_{X'}Y'=V(f_i'),
\]
and analogous relations hold for \(Y'',Y_i''\). The assertion that
\(f'\) vanishes on the \(Z_i'\) is equivalent to
\[
  Y_i'=Z_i'\qquad\text{for every }i.
\]
\footnote{Editor's note: the original added ``at least at every point
\(x'\) of \(X'\) above \(s'\),'' but this restriction seems unnecessary.}

Let \(\mathfrak m\) be the maximal ideal of \(A\). For every
\(n\geq0\), introduce the closed subscheme
\[
  S_n=\operatorname{Spec}(A/\mathfrak m^{n+1})
\]
of \(S\), and let \(X_n,Y_n,Z_n^i,Y_n^i\) be obtained from
\(X,Y,Z_i,Y_i\) by the base change \(S_n\to S\). More generally, for
every \(S\)-prescheme \(P\), write \(P_n=P\times_SS_n\).

For every \(n\) and \(i\in I\), consider the functor
\SGARefTarget{sga3:ix:lemma:4:4:display:03}
\[
  F_n^i
  =\prod_{Z_n^i/S_n}Y_n^i/Z_n^i:
  (\mathbf{Sch})_{/S_n}^{\circ}\longrightarrow\mathbf{Set},
\]
defined, for every \(P\) over \(S_n\), by
\SGARefTarget{sga3:ix:lemma:4:4:display:04}
\[
F_n^i(P)
=\Gamma\bigl((Y_n^i)_P/(Z_n^i)_P\bigr)
=
\begin{cases}
  \varnothing,&(Y_n^i)_P\ne(Z_n^i)_P,\\
  \{\operatorname{id}\},&(Y_n^i)_P=(Z_n^i)_P.
\end{cases}
\]
Since \(S_n\) is local Artinian and \(Z_n^i\) is flat, hence essentially
free, over \(S_n\), \SGARefLink{sga3:viii:theorem:6:4}{VIII, 6.4} shows
that each \(F_n^i\) is represented by a closed subprescheme \(T_n^i\)
of \(S_n\).

% Source page: Exp9 component-local p. 14; complete-reader p. 660.
The completion \(\widehat A\) of the local ring \(A\) is Noetherian
because \(A\) is, and it is the projective limit of the rings \(A_n\).%
\footnote{Editor's note: the original has been expanded and corrected to
show that the projective limit of the rings
\(\mathscr O(T_n^i)\) defines a closed formal subscheme of
\(\operatorname{Spf}(\widehat A)\).}
Let \(K_n^i\) be the ideal of
\[
  A_n=A/\mathfrak m^{n+1}
\]
defining \(T_n^i\), and let \(K^i=\varprojlim_nK_n^i\); this is an ideal
of \(\widehat A\).

For fixed \(i\), integers \(m\geq n\), and any \(S_n\)-prescheme \(P\),
\[
  (Y_n^i)_P
  =Y_i\times_SS_n\times_{S_n}P
  =Y_i\times_SP
  =(Y_m^i)_P,
\]
and similarly \((Z_n^i)_P=(Z_m^i)_P\). It follows that \(F_n^i\) is the
restriction \(F_m^i\times_{S_m}S_n\) of \(F_m^i\) to
\((\mathbf{Sch})_{/S_n}\), and therefore
\[
  T_n^i=T_m^i\times_{S_m}S_n.
\]
We thus have a commutative diagram with exact rows:
\SGARefTarget{sga3:ix:lemma:4:4:diagram:02}
\[
\begin{tikzcd}[column sep=large,row sep=large]
  0 \arrow[r] & K_m^i \arrow[r]\arrow[d,two heads]
    & A_m \arrow[r]\arrow[d,two heads]
    & \mathscr O(T_m^i) \arrow[r]\arrow[d,two heads] & 0\\
  0 \arrow[r] & K_n^i \arrow[r]
    & A_n \arrow[r]
    & \mathscr O(T_n^i) \arrow[r] & 0
\end{tikzcd}
\]
All the vertical arrows are surjective. It follows, first, that the
projective system \((K_n^i)_{n\in\mathbb N}\) satisfies the
Mittag--Leffler condition, so that
\[
  \varprojlim_n\mathscr O(T_n^i)
  \simeq\widehat A/K^i
\]
as topological rings (cf. EGA \(0_{\mathrm{III}}\), \S13.2). Second, the
map \(K^i\to K_n^i\) is surjective (cf. \SGARefLink{sga3:ix:bib:bens}{[BEns]}, III, \S7.4,
Proposition~5). Hence
\[
  K_n^i\simeq
  \frac{K^i+\mathfrak m^{n+1}\widehat A}
       {\mathfrak m^{n+1}\widehat A},
  \qquad
  \mathscr O(T_n^i)\simeq
  (\widehat A/K^i)\otimes_A(A/\mathfrak m^{n+1}).
\]

In other words, \((T_n^i)_n\) is an inductive system of affine Artinian
schemes, and its inductive limit
\[
  T^i=\varinjlim_nT_n^i
\]
is a closed formal subscheme of
\[
  \widehat S=\operatorname{Spf}(\widehat A)
\]
(cf. EGA I, \S10), whose reduction modulo \(\mathfrak m^{n+1}\) is
\(T_n^i\).

Let \(T\) be the closed formal subscheme of \(\widehat S\) given by the
intersection of the \(T^i\), that is, defined by
\[
  K=\sum_iK^i.
\]
Since \(\widehat A\) is Noetherian, there is a finite subset \(J\) of
\(I\) such that
\[
  K=\sum_{i\in J}K^i,
  \qquad\text{equivalently}\qquad
  T=\bigcap_{i\in J}T^i.
\]
For every \(n\), therefore,
\[
  K_n=\sum_{i\in J}K_n^i,
\]
where \(K_n\) is the image of \(K\) in \(A_n\).

Recall that \(f_i\) is the image of
\(f\in B=\mathscr O(X)\) in \(\mathscr O(Z_i)\), and that \(f_i'\) is
its image in
\[
  \mathscr O(Z_i')
  =\mathscr O(Z_i)\otimes_AA'.
\]
The relation \(Y_i'=Z_i'\) is equivalent to \(f_i'=0\). Now
\(\mathscr O(Z_i)\otimes_AA'\) is the direct limit of the subalgebras
\(\mathscr O(Z_i)\otimes_AA''\), for \(A''\) as above. Consequently,
there is an \(A''\) for which \(Y_i''=Z_i''\). A priori \(A''\) depends
on \(i\), but, because \(J\) is finite, one \(A''\) can be chosen that
works for every \(i\in J\). Put
\[
  S''=\operatorname{Spec}(A''),
  \qquad
  S''_{(n)}=S''\times_SS_n.
\]
\footnote{Editor's note: the original has been expanded and corrected,
distinguishing \(S''_{(n)}\) from the subscheme
\(S_n''=S_n''(s'')\) introduced below.}

% Source page: Exp9 component-local p. 15; complete-reader p. 661.
For every \(i\in J\), the equality
\[
  Y_i''\times_{S''}S''_{(n)}
  =(Y_n^i)_{S''_{(n)}}
  =Z_i''\times_{S''}S''_{(n)}
  =(Z_n^i)_{S''_{(n)}}
\]
shows, by the definition of \(T_n^i\), that \(S''_{(n)}\to S_n\)
factors through \(T_n^i\). It therefore factors through
\[
  T_n=\bigcap_{i\in J}T_n^i,
\]
and hence through \(T_n^i\) for every \(i\in I\). Let \(f''\) be the
image of \(f\) in \(B''=B\otimes_AA''\), and let \(f''_{(n)}\) be its
image in
\[
  B''_{(n)}=B''/\mathfrak m^{n+1}B''.
\]
The preceding factorization means that, for every \(n\), the section
\(f''_{(n)}\) vanishes on all the \(Z_i\times_SS''_{(n)}\).

Because \(A\to A''\to A'\) are local homomorphisms, the image \(s''\)
of \(s'\) in \(S''\) is the closed point of \(S''\), corresponding to
the maximal ideal \(\mathfrak m''\) of \(A''\), and its image in \(S\)
is \(s\). \emph{[The source prints \(B''\) here; since
\(s''\in S''=\operatorname{Spec}(A'')\), the intended ring is
\(A''\).]} Fix \(n\), and now let \(S_n''\) denote the \(n\)-th
infinitesimal neighborhood of \(s''\) in \(S''\):
\[
  S_n''=\operatorname{Spec}(A''/\mathfrak m''^{\,n+1}).
\]
Put
\[
  B_n''=B''/\mathfrak m''^{\,n+1}B'',
  \qquad
  Z_{n}^{i''}=Z_i\times_SS_n''.
\]
\emph{[The source appends
\(Z_n^{i''}=\operatorname{Spec}(B_n'')\); the displayed equality would
identify each \(Z_n^{i''}\) with \(X_n''\), so the base-change
definition is retained here.]}
Then the image \(f_n''\) of \(f''_{(n)}\) in \(B_n''\) vanishes on every
\(Z_n^{i''}\). By \SGARefLink{sga3:ix:lemma:4:5}{Lemma~4.5 below}, the
family of fibers
\[
  Z_0^{i'''}
  =Z_n^{i''}\times_{S_n''}\kappa(s'')
  =Z_{is}\times_{\kappa(s)}\kappa(s'')
\]
is schematically dense in
\[
  X_0'''
  =X_n''\times_{S_n''}\kappa(s'')
  =X_s\times_{\kappa(s)}\kappa(s'').
\]
Thus \(S_n'',X_n'',(Z_n^{i''})_{i\in I}\), and
\[
  S_0'''=\operatorname{Spec}\kappa(s'')
\]
satisfy the hypotheses of \SGARefLink{sga3:ix:lemma:4:2}{Lemma~4.2}.
It follows that \(f_n''=0\), that is,
\[
  f''\in\mathfrak m''^{\,n+1}B''
  \qquad\text{for every }n.
\]

Since \(B\) is Noetherian and \(B''=B\otimes_AA''\) is a localization
of a finite-type \(B\)-algebra, \(B''\) is Noetherian. Its local rings
are therefore separated for their usual topology. Hence \(f''\) is zero
at every point \(x''\) of \(X''\) for which
\(\mathfrak m''B''\) is contained in the maximal ideal of
\(\mathscr O_{X'',x''}\), that is, at every point of \(X''\) above
\(s''\). A fortiori, \(f'\) is zero at every point of \(X'\) above
\(s'\). \qed

\medskip
\noindent\SGARefTarget{sga3:ix:lemma:4:5:0}\textbf{Lemma 4.5.0.}%
\footnote{Editor's note: this lemma has been inserted here because it is
used in the proofs of \SGARefLink{sga3:ix:lemma:4:5}{4.5} and
\SGARefLink{sga3:ix:theorem:4:7}{4.7}.}
\textit{Let \(X\) be an \(S\)-prescheme, let \((Z_i)_{i\in I}\) be a
family of subpreschemes of \(X\), and let \(S'\to S\) be faithfully
flat. If \((Z_i')_{i\in I}\) is schematically dense in \(X'\), then
\((Z_i)_{i\in I}\) is schematically dense in \(X\).}

This follows from EGA IV\(_3\), 11.10.5(i) and 11.9.10(i). It remains to
give the proof of \SGARefLink{sga3:ix:lemma:4:5}{the following lemma}.

\medskip
\noindent\SGARefTarget{sga3:ix:lemma:4:5}\textbf{Lemma 4.5.}
\textit{Let \(X\) be a locally Noetherian prescheme}%
\footnote{Editor's note: this hypothesis is in fact superfluous; cf. EGA
IV\(_3\), 11.10.6 and 11.9.13.}
\textit{ over a field \(k\), let \((Z_i)_{i\in I}\) be a family of
subpreschemes of \(X\), let \(k'/k\) be a field extension, and let
\(X',Z_i'\) be obtained from \(X,Z_i\) by base change. Then
\((Z_i)_{i\in I}\) is schematically dense in \(X\) if and only if
\((Z_i')_{i\in I}\) is schematically dense in \(X'\).}

\emph{Proof.}
The ``if'' direction follows from
\SGARefLink{sga3:ix:lemma:4:5:0}{Lemma~4.5.0}; let us prove the
converse.%
\footnote{Editor's note: the original has been expanded in what follows.}
First, we may suppose \(X=\operatorname{Spec}(B)\) affine. For every
\(i\), let \((Z_{ij})_{j\in J_i}\) be an affine open covering of \(Z_i\).
Replacing \((Z_i)_{i\in I}\) by
\((Z_{ij})_{(i,j)\in J}\), where
\[
  J=\coprod_{i\in I}J_i,
\]
we may also suppose that every \(Z_i\) is affine.

Let
\[
  g\in B'=B\otimes_kk',
  \qquad
  t'\in B_g'
\]
be a section of \(\mathscr O_{X'}\) on the affine open
\(U'=X_g'\) whose restriction to each \(Z_i'\cap U'\) is zero. Thus its
image in every
\[
  \bigl(\mathscr O(Z_i)\otimes_kk'\bigr)_g
\]
is zero. There exists a finite-type \(k\)-subalgebra \(A\) of \(k'\)
such that
\[
  g\in B_A=B\otimes_kA,
  \qquad
  t'\in(B_A)_g.
\]
The map
\[
  \mathscr O(Z_i)\otimes_kA
  \longrightarrow
  \mathscr O(Z_i)\otimes_kk'
\]
is injective, and therefore so is
\[
  \bigl(\mathscr O(Z_i)\otimes_kA\bigr)_g
  \longrightarrow
  \bigl(\mathscr O(Z_i)\otimes_kk'\bigr)_g.
\]
The hypothesis consequently implies that the image of \(t'\) in every
\(\bigl(\mathscr O(Z_i)\otimes_kA\bigr)_g\) is zero. It remains to show
that \((Z_{iA})_{i\in I}\) is schematically dense in \(X_A\), which
follows from EGA IV\(_3\), 11.9.10(b).%
\footnote{Editor's note: the original continued: ``Using
\SGARefLink{sga3:ix:corollary:4:3}{4.3} for the
\(X_{A'}\), where \(A'\) is a quotient of \(A\) by a power of a maximal
ideal, and Hilbert's Nullstellensatz, one is reduced to the case where
\(k'\) is a finite extension of \(k\).'' This reduction could be
expanded and that argument pursued, since the case of a finite extension
\(k'/k\) is somewhat simpler than the more general case treated in EGA
IV\(_3\), 11.9.10(b).}
\qed

Let us record in passing \SGARefLink{sga3:ix:corollary:4:6}{the following result}, which will be used in a
later expos\'e.%
\footnote{Editor's note: this reference remains to be specified.}

% Source page: Exp9 component-local p. 16; complete-reader p. 662.
\medskip
\noindent\SGARefTarget{sga3:ix:corollary:4:6}\textbf{Corollary 4.6.}
\textit{Let \(X\) be a flat \(S\)-prescheme and \(U\) an open subset of
\(X\). Suppose that \(U_s\) is schematically dense in \(X_s\) for every
\(s\in S\), and that either \(X\) is locally Noetherian or \(X\) is
locally of finite presentation over \(S\). Then \(U\) is schematically
dense in \(X\).}

\emph{Proof.}
For simplicity, suppose \(U\) is retrocompact in \(X\) (cf. EGA
IV\(_3\), 11.10.10 and 11.9.17 for the general case). The locally
Noetherian case is included only for reference; it is contained in
\SGARefLink{sga3:ix:corollary:4:3}{4.3}.

In the second case, we may plainly suppose \(S\) and \(X\) affine. Then
\(X\) and \(U\) are of finite presentation over
\(S=\operatorname{Spec}(A)\). The ring \(A\) is the direct limit of its
finite-type \(\mathbb Z\)-subalgebras \(A_i\). The ``patented
procedure'' already used (cf. EGA IV\(_3\), 8.8.2 and 8.10.5(iii))
shows that there are an index \(i\), an affine scheme \(X_i\) over
\[
  S_i=\operatorname{Spec}(A_i),
\]
and an open subset \(U_i\) of \(X_i\), from which \(X,U\) are obtained
by the base change \(S\to S_i\). Let \(E_i\) be the subset of \(S_i\)
consisting of those \(s\in S_i\) for which \((U_i)_s\) is schematically
dense in \((X_i)_s\).

By EGA IV\(_2\), 5.9.9 and 5.10.2, \(E_i\) is the set of points
\(s\in S_i\) for which \((U_i)_s\) contains the set
\[
  \operatorname{Ass}\mathscr O_{(X_i)_s}
\]
of points associated with the structure sheaf of \((X_i)_s\). By EGA
IV\(_3\), 11.9.17.1, this is a constructible condition; hence \(E_i\)
is constructible.%
\footnote{Editor's note: the original has been expanded in this paragraph
and the next.}

By \SGARefLink{sga3:ix:lemma:4:5}{4.5}, the inverse image of \(E_i\)
under \(S_j\to S_i\) is \(E_j\), and its inverse image under
\(S\to S_i\) is the set \(E\) of those \(s\in S\) for which \(U_s\)
is schematically dense in \(X_s\).%
\footnotemark[24]
By hypothesis, \(E=S\), which is also the inverse image of \(S_i\) under
every \(S\to S_i\). EGA IV\(_3\), 8.3.11 therefore implies that there
is a \(j\geq i\) such that \(E_j=S_j\); that is, \((U_j)_s\) is
schematically dense in \((X_j)_s\) for every \(s\in S_j\).

Apply \SGARefLink{sga3:ix:lemma:4:4}{4.4} to
\((S_j,X_j,U_j)\) and the base change \(S\to S_j\). It follows that
\(U\) is schematically dense in \(X\). Notice that here
\SGARefLink{sga3:ix:lemma:4:4}{4.4} is used only
for a family with a finite index set (in fact, a singleton); in this
case its proof simplifies considerably. \qed

% Source page: Exp9 component-local p. 17; complete-reader p. 663.
\medskip
\noindent\SGARefTarget{sga3:ix:theorem:4:7}\textbf{Theorem 4.7.}
\textit{Let \(S\) be a prescheme, and let \(G\) be a group of
multiplicative type and of finite type over \(S\). For every integer
\(n>0\), let \({}_nG\) be the kernel of \(n\cdot\operatorname{id}_G\).
Then the family of subpreschemes \(({}_nG)_{n>0}\) is schematically
dense in \(G\) in the sense of
\SGARefLink{sga3:ix:definition:4:1}{Definition~4.1}.}

\emph{Proof.}
We distinguish two cases.

\smallskip
\SGARefTarget{sga3:ix:theorem:4:7:proof:item:a}\noindent\textit{(a) \(S\) locally Noetherian.}
By \SGARefLink{sga3:ix:corollary:4:3}{4.3}, we are reduced to the case
in which \(S\) is the spectrum of a field \(k\). By
\SGARefLink{sga3:ix:lemma:4:5:0}{4.5.0}, we may suppose that \(k\) is
algebraically closed and that \(G\) is diagonalizable, say
\[
  G=D_k(M),
\]
where \(M\) is a commutative group of finite type. Then
\[
  M\simeq\Gamma\times\mathbb Z^r
\]
with \(\Gamma\) finite, and consequently
\[
  G\simeq G_1\times G_2,
  \qquad
  G_1=D(\Gamma),
  \qquad
  G_2=\mathbf G_m^r.
\]
For \(n\) sufficiently divisible---namely, for \(n\) a multiple of the
order of \(\Gamma\)---one has
\[
  {}_nG=G_1\times{}_nG_2,
  \qquad\text{since}\qquad
  {}_nG_1=G_1.
\]
Applying \SGARefLink{sga3:ix:corollary:4:3}{4.3} again to the projection
\(G\to G_1\), we are reduced to
\(G_2\), hence to the case \(G=\mathbf G_m^r\). Since \(G\) is then
reduced, schematic density of \(({}_nG)_{n>0}\) is equivalent to
density of \(\bigcup_n{}_nG\) in the ordinary topology. As
\[
  {}_nG=({}_n\mathbf G_m)^r,
\]
we are reduced to \(G=\mathbf G_m\), an irreducible group of dimension
one. The assertion now follows because
\(\bigcup_n{}_nG\), the set of roots of unity in \(k\), is infinite.

\smallskip
\SGARefTarget{sga3:ix:theorem:4:7:proof:item:b}\noindent\textit{(b) The general case.}
For every \(s\in S\), there are an open neighborhood \(U\) of \(s\)
and a faithfully flat quasi-compact morphism \(S'\to U\) such that
\(G'=G_{S'}\) is diagonalizable, say \(G'=D_{S'}(M)\). By
\SGARefLink{sga3:ix:lemma:4:5:0}{4.5.0} we may
reduce to \(G'\), and therefore suppose that \(G\) itself is
diagonalizable. It is obtained by base change
\[
  S\longrightarrow\operatorname{Spec}(\mathbb Z)
\]
from the absolute group \(H=D_{\mathbb Z}(M)\). Part~\SGARefLink{sga3:ix:theorem:4:7:proof:item:a}{(a)} shows that,
for every \(s\in\operatorname{Spec}(\mathbb Z)\), the family
\(({}_nH_s)_{n>0}\) is schematically dense in \(H_s\). It remains only
to apply \SGARefLink{sga3:ix:lemma:4:4}{4.4}.%
\footnote{Editor's note: \SGARefLink{sga3:ix:theorem:4:7}{this result} will be needed in X, \SGARefLink{sga3:x:lemma:4-3}{4.3} for a
non-locally-Noetherian base, namely
\(S=\operatorname{Spec}(\widehat A\otimes_A\widehat A)\), where
\(\widehat A\) is the completion of the Noetherian local ring \(A\).}
\qed

\medskip
\noindent\SGARefTarget{sga3:ix:corollary:4:8}\textbf{Corollary 4.8.}
\begin{enumerate}
  \item[(a)]\SGARefTarget{sga3:ix:corollary:4:8:item:a}
  \textit{Let \(u,v:H\to G\) be two homomorphisms of \(S\)-preschemes
  in groups, with \(H\) of multiplicative type}%
  \footnote{Editor's note: ``diagonalizable'' has been replaced by ``of
  multiplicative type.''}
  \textit{ and of finite type. Suppose that, for every integer \(n>0\),
  the restrictions of \(u\) and \(v\) to \({}_nH\) are equal. Then
  \(u=v\).}

  \item[(b)]\SGARefTarget{sga3:ix:corollary:4:8:item:b}
  \textit{Let \(H_1,H_2\) be two subgroups of multiplicative type and
  of finite type of a prescheme in groups \(G\). If
  \({}_nH_1={}_nH_2\) for every integer \(n>0\), then \(H_1=H_2\).}
\end{enumerate}

\emph{Proof.}
Part~\SGARefLink{sga3:ix:corollary:4:8:item:a}{(a)} follows from part~\SGARefLink{sga3:ix:corollary:4:8:item:b}{(b)} by considering the subgroups \(H_1,H_2\)
of \(H\times G\) that are the graphs of \(u,v\). For part~\SGARefLink{sga3:ix:corollary:4:8:item:b}{(b)}, put
\[
  H=H_1\cap H_2.
\]
This is a subgroup prescheme of \(H_i\), for \(i=1,2\), and we must show
that it equals \(H_i\). The hypothesis says precisely that \(H\)
contains all the \({}_nH_i\). We are therefore reduced to \SGARefLink{sga3:ix:corollary:4:9}{the next
corollary}. \qed

\medskip
\noindent\SGARefTarget{sga3:ix:corollary:4:9}\textbf{Corollary 4.9.}
\textit{Let \(G\) be an \(S\)-group of multiplicative type and of finite
type, and let \(H\) be a subgroup prescheme containing \({}_nG\) for
every \(n>0\). Then \(H=G\).}

\emph{Proof.}
By \SGARefLink{sga3:ix:theorem:4:7}{4.7}, it is enough to prove that
\(H\) is closed, or equivalently that it equals \(G\) set-theoretically.
This reduces us to the case in which \(S\) is the spectrum of a field.
Every subgroup prescheme of \(G\) is then closed, by VI\(_B\), \SGARefLink{cor:sga3-VIB-1.4.2}{1.4.2}.
\qed

% Source page: Exp9 component-local p. 18; complete-reader p. 664.
\medskip
\noindent\SGARefTarget{sga3:ix:remark:4:10}\textbf{Remark 4.10.}
Under the hypotheses of \SGARefLink{sga3:ix:theorem:4:7}{4.7}, let
\(m>0\) be an integer with the following properties: for every
\(s\in S\), \(m\) is not a power of the characteristic of \(\kappa(s)\);
and if \(G_s\) has type \(M\), every prime divisor of the torsion
subgroup of \(M\) divides \(m\). The second condition is automatically
satisfied if \(G\) is a torus. The proof shows that in
\SGARefLink{sga3:ix:theorem:4:7}{Theorem~4.7} and
\SGARefLink{sga3:ix:corollary:4:8}{Corollary~4.8} and
\SGARefLink{sga3:ix:corollary:4:9}{Corollary~4.9} one may restrict to
the subgroups
\[
  {}_{m^r}G,\qquad r>0.
\]

Moreover, if \(G\) is smooth over \(S\), then for every \(s\in S\) there
are an open neighborhood \(U\) of \(s\) and an integer \(m>0\), prime
to every residue characteristic on \(U\), that satisfy the preceding
conditions for \(G_U\). For instance, take \(m\) to be the order of the
torsion subgroup of the type of \(G\) at \(s\) (cf.
\SGARefLink{sga3:ix:remark:1:4:1:item:a}{1.4.1(a)} and
\SGARefLink{sga3:ix:proposition:2:1:item:e:01}{2.1(e)}). After restriction to
\(U\), the groups \({}_{m^r}G\) are then
finite and \'etale over \(S\), and their family is schematically dense.
In certain applications---in particular those involving
\SGARefLink{sga3:ix:theorem:3:2}{Theorem~3.2} and
\SGARefLink{sga3:ix:theorem:3:2:bis}{Theorem~3.2~bis} (cf. XI, 6),
though not those involving \SGARefLink{sga3:ix:theorem:3:6}{Theorem~3.6}
and \SGARefLink{sga3:ix:theorem:3:6:bis}{Theorem~3.6~bis}---this
observation allows one to avoid
\SGARefLink{sga3:ix:theorem:3:1}{Theorem~3.1} and its use of Hochschild
cohomology.

\SGARefTarget{sga3:ix:section:5}\section*{5. Central Homomorphisms of Groups of Multiplicative Type}
\addcontentsline{toc}{section}{5. Central Homomorphisms of Groups of Multiplicative Type}

\medskip
\noindent\SGARefTarget{sga3:ix:lemma:5:0}\textbf{Lemma 5.0.}%
\footnote{Editor's note: this lemma has been made explicit because it is
used several times below; it was implicit in the original.}
\textit{Let \((A,\mathfrak m)\) be a Noetherian local ring, put
\(S=\operatorname{Spec}(A)\), and let \(H\) be a finite scheme over
\(S\), say \(H=\operatorname{Spec}(B)\), where \(B\) is finite over
\(A\). Let \(K\) be a subscheme of \(H\) such that reduction modulo
\(\mathfrak m^{n+1}\) gives \(K_n=H_n\) for every \(n\). Then \(K=H\).}

\emph{Proof.}
Let \(s\) be the closed point of \(S\). First, \(K\) is a closed
subscheme of \(H\). Indeed, a priori it is a closed subscheme of an open
subset \(U\) of \(H\). But \(K\), and hence \(U\), contains the fiber
\[
  K_s=H_s.
\]
Since \(H\to S\) is finite, and therefore closed, the complement of
\(U\) is empty; thus \(U=H\). Hence \(K\) is defined by an ideal \(I\)
of \(B\). The hypothesis implies
\[
  I\subset\mathfrak m^nB
  \qquad\text{for every }n.
\]
Since \(B\) is a finite \(A\)-module, it is separated for the
\(\mathfrak m\)-adic topology. Therefore \(I=0\). \qed

\medskip
\noindent\SGARefTarget{sga3:ix:theorem:5:1}\textbf{Theorem 5.1.}
\textit{Let \(u,v:H\to G\) be two homomorphisms of \(S\)-preschemes in
groups, where \(H\) is of multiplicative type and of finite type.
Suppose either that \(S\) is locally Noetherian or that \(G\) is of
finite presentation over \(S\). Let \(s\in S\) satisfy \(u_s=v_s\), and
suppose that \(u_s\) is central. Then there is an open neighborhood
\(U\) of \(s\) such that \(u_U=v_U\).}

\emph{Proof.}
We distinguish the two cases.

\smallskip
\SGARefTarget{sga3:ix:theorem:5:1:proof:item:a}\noindent\textit{(a) \(S\) locally Noetherian.}%
\footnote{Editor's note: the original has been expanded in what follows.}
Let
\[
  K=\operatorname{Ker}(u,v)
\]
be the inverse image of the diagonal of \(G\times_SG\) under
\((u,v):H\to G\times_SG\). It is a subgroup prescheme of \(H\). We seek
an open \(U\) for which \(K_U=H_U\). Because \(S\) is locally
Noetherian and \(H\) is of finite type over \(S\), the prescheme \(H\)
is locally Noetherian. Hence the immersion \(K\hookrightarrow H\) is of
finite type (cf. EGA I, 6.3.5), so \(K\) is of finite type, and therefore
of finite presentation, over \(S\). By EGA IV\(_3\), 8.8.2.4, it is
enough to prove that
\[
  K_{S_0}=H_{S_0},
  \qquad
  S_0=\operatorname{Spec}(\mathscr O_{S,s}).
\]
We may thus suppose that \(S\) is local with closed point \(s\).
Replacing the Noetherian local ring \(A\) by its completion
\(\widehat A\), if desired, gives a faithfully flat quasi-compact base
change \(\widehat S\to S\), and allows us to suppose that \(A\) is
complete.%
\footnote{Editor's note: this is because \(K=H\) if
\(K_{\widehat S}=H_{\widehat S}\); cf. SGA~1, VIII, \SGARefLink{sga3:viii:corollary:5:4}{5.4}. The argument
below does not use completeness of \(A\).}

By \SGARefLink{sga3:ix:corollary:3:4}{3.4}, one has \(u_n=v_n\) for every
\(n\), where the subscript denotes reduction modulo
\(\mathfrak m^{n+1}\), with \(\mathfrak m\) the maximal ideal of \(A\).
For every \(m>0\), let
\[
  {}_mu,{}_mv:{}_mH\longrightarrow G
\]
be the homomorphisms induced by \(u,v\). Then
\[
  ({}_mu)_n=({}_mv)_n
  \qquad\text{for every }n.
\]
The group \({}_mH\) is finite over \(S\), by \SGARefLink{sga3:ix:proposition:2:2}{2.2}, so
\SGARefLink{sga3:ix:lemma:5:0}{5.0} gives \({}_mu={}_mv\). This holds
for every \(m\), and therefore
\SGARefLink{sga3:ix:theorem:4:7}{4.7} gives \(u=v\).

\smallskip
\SGARefTarget{sga3:ix:theorem:5:1:proof:item:b}\noindent\textit{(b) \(G\) of finite presentation.}%
\footnote{Editor's note: the original has been expanded in what follows.}
The group \(H\) is also of finite presentation over \(S\). By EGA
IV\(_3\), 8.8.2, we may again suppose that \(S\) is local with closed
point \(s\), and it suffices to prove \(u=v\). Let \(f:S'\to S\) be a
faithfully flat quasi-compact morphism. Write
\[
  f_H:H'\to H,
  \qquad
  f_G:G'\to G,
  \qquad
  u',v':H'\to G'
\]
for the resulting morphisms. If \(u'=v'\), then
\[
  u\circ f_H=f_G\circ u'=f_G\circ v'=v\circ f_H,
\]
and \(u=v\) because \(f_H\) is an epimorphism. After a faithfully flat
quasi-compact extension of the base, we may therefore suppose that
\(H\) is diagonalizable:
\[
  H=D_S(M)
\]
for a commutative group \(M\) of finite type.

As in the proof of \SGARefLink{sga3:ix:corollary:4:6}{4.6}, introduce
the increasing filtered family of \(\mathbb Z\)-subalgebras \(A_i\) of
\(A=\mathscr O_{S,s}\) that are of finite type, and put
\[
  S_i=\operatorname{Spec}(A_i).
\]
\footnotemark[30]
For every \(i\), the group \(H=D_S(M)\) comes from
\[
  H_i=D_{S_i}(M).
\]
Since \(G\) is of finite presentation over \(S\), EGA IV\(_3\), 8.8.2
(see also VI\(_B\), 10.2 and 10.3) gives an index \(i\), an
\(S_i\)-group prescheme \(G_i\) of finite presentation, and
homomorphisms
\[
  u_i,v_i:H_i\longrightarrow G_i
\]
whose base changes are \(u,v\). Let \(s_i\) be the image of \(s\) in
\(S_i\), and define
\[
  \rho_i:H_i\times_{S_i}G_i\longrightarrow G_i,
  \qquad
  \rho_i(h,g)=u_i(h)g\,u_i(h)^{-1}.
\]
Because \(u_s\) is central,
\[
  \rho_s=\rho_i\times_{\kappa(s_i)}\kappa(s)
\]
is the second projection. The extension
\(\kappa(s_i)\to\kappa(s)\) is faithfully flat, so \(\rho_i\) is also
the second projection over \(s_i\); hence \((u_i)_{s_i}\) is central.
Likewise \(u_s=v_s\) implies
\[
  (u_i)_{s_i}=(v_i)_{s_i}.
\]
Part~\SGARefLink{sga3:ix:theorem:5:1:proof:item:a}{(a)}, applied over \(S_i\), now gives the conclusion. \qed

\medskip
\noindent\SGARefTarget{sga3:ix:corollary:5:2}\textbf{Corollary 5.2.}
\textit{Let \(u:H\to G\) be a homomorphism of \(S\)-preschemes in
groups, where \(H\) is of multiplicative type and of finite type, and
suppose either that \(S\) is locally Noetherian or that \(G\) is of
finite presentation over \(S\). If \(s\in S\) and \(u_s\) is the unit
homomorphism, then there is an open neighborhood \(U\) of \(s\) on
which \(u_U\) is the unit homomorphism.}

\medskip
\noindent\SGARefTarget{sga3:ix:corollary:5:3}\textbf{Corollary 5.3.}
\textit{Let \(u,H,G\) be as in
\SGARefLink{sga3:ix:corollary:5:2}{5.2}, and suppose in addition that \(G\)
is separated over \(S\). Let \(U\) be the set of points \(s\in S\) for
which \(u_s\) is the unit homomorphism. Then \(U\) is open and closed
in \(S\), and \(u_U:H_U\to G_U\) is the unit homomorphism.}

\emph{Proof.}
Only closedness remains to be proved. Put \(K=\operatorname{Ker}(u)\).
Since \(G\) is separated over \(S\), \(K\) is a closed subprescheme of
\(H\), and \(U\) is the set of points \(s\in S\) for which \(K_s=H_s\).
It follows readily that \(U\) is closed, for example by
\SGARefLink{sga3:viii:theorem:6:4}{VIII, 6.4}, since \(H\) is
essentially free over \(S\) by VIII, \SGARefLink{sga3:viii:proposition:6:3}{6.3}. Alternatively, we may suppose
that \(S\) is reduced and that \(U\) is dense in \(S\), hence
schematically dense. Then \(H_U\) is schematically dense in \(H\)
because \(H\) is flat over \(S\),%
\footnote{Editor's note: and of finite presentation over \(S\); cf. EGA
IV\(_3\), 11.10.5(ii) and 11.9.10(ii).}
and \(u\) and the unit homomorphism, being equal on \(H_U\), are equal
on \(H\). \qed

% Source page: Exp9 component-local p. 20; complete-reader p. 666.
\medskip
\noindent\SGARefTarget{sga3:ix:corollary:5:4}\textbf{Corollary 5.4.}
\textit{Let \(S\) be a prescheme; let \(H,G\) be \(S\)-groups, with
\(H\) of multiplicative type and of finite type and \(G\) separated
and of finite presentation over \(S\); and let
\(\pi:S'\to S\) be a}%
\footnote{Editor's note: ``covering morphism for the faithfully flat
quasi-compact topology'' has been replaced by ``universal effective
epimorphism,'' so that the result applies to the morphism \(K\to S\) in
\SGARefLink{sga3:ix:corollary:5:5}{5.5}; the beginning of the proof has
been modified accordingly.}
\textit{ universal effective epimorphism---for example, a faithfully
flat quasi-compact morphism or a morphism admitting a section (cf. IV,
\SGARefLink{sga3:iv:proposition:1:12}{1.12})---whose fibers are geometrically connected.}%
\footnote{Editor's note: this holds, for example, when
\(S=\operatorname{Spec}(k)\) and \(S'=\operatorname{Spec}(k')\), where
\(k\) is a field and \(k'/k\) is radicial; cf. X, \SGARefLink{sga3:x:proposition:1:4}{Proposition~1.4}.}

\textit{Let \(u':H'\to G'\) be a central homomorphism of the
\(S'\)-groups obtained from \(H,G\) by base change. Then there is a
unique homomorphism \(u:H\to G\) of \(S\)-groups such that
\(u\times_SS'=u'\). If \(S'/S\) admits a section \(g\), then \(u\) is
the morphism obtained from \(u'\) by the base change
\(g:S\to S'\).}

\emph{Proof.}
Because \(\pi\) is a universal effective epimorphism, so is
\(H'\to H\); uniqueness follows as at the beginning of the proof of
\SGARefLink{sga3:ix:theorem:5:1:proof:item:b}{5.1(b)}. If \(\pi\) has a section
\(g\), then
\[
  u'=\pi^*u
  \quad\Longrightarrow\quad
  u=g^*\pi^*u=g^*u'.
\]

For existence, IV, \SGARefLink{sga3:iv:proposition:2:3}{2.3} reduces us to showing that the two homomorphisms
\[
  u_1'',u_2'':H''\longrightarrow G''
\]
of \(S''\)-groups obtained from \(u'\) by the two base changes
\[
  \operatorname{pr}_1,\operatorname{pr}_2:
  S''=S'\times_SS'\longrightarrow S'
\]
are equal. Their pullbacks along the diagonal \(S'\to S''\) are equal,
since both are \(u'\).%
\footnote{Editor's note: the preceding sentence has been added.}
Because \(u_1'',u_2''\) are central, we may apply
\SGARefLink{sga3:ix:corollary:5:3}{5.3} to
\[
  u_1''(u_2'')^{-1}.
\]
There is therefore an open-and-closed subset \(U\) of \(S''\),
containing the diagonal, over which \(u_1''\) and \(u_2''\) coincide.
The fibers of \(S'/S\) are geometrically connected, hence so are those
of \(S''/S\), and in particular they are connected. Since \(U\)
contains the diagonals of these fibers, it contains every fiber and
therefore equals \(S''\). This proves existence. \qed

\medskip
\noindent\SGARefTarget{sga3:ix:corollary:5:5}\textbf{Corollary 5.5.}
\textit{Let \(S\) be a prescheme, let \(K\) be an \(S\)-group prescheme
locally of finite type}%
\footnote{Editor's note: the hypothesis ``locally of finite type'' has
been added to agree with VI\(_A\), \SGARefLink{sga3:ix:corollary:2:4}{2.4}. In fact, over a field every
connected group scheme is geometrically connected.}
\textit{ and with connected fibers, and let \(H\) be an invariant
subgroup of \(K\), of multiplicative type and of finite type. Then
\(H\) is a central subgroup of \(K\).}

\emph{Proof.}%
\footnote{Editor's note: the original has been expanded in what follows.}
First, \(\pi:K\to S\) has geometrically connected fibers, because for a
group scheme locally of finite type over a field, connectedness implies
geometric connectedness (cf. VI\(_A\), \SGARefLink{sga3:ix:corollary:2:4}{2.4}). Apply
\SGARefLink{sga3:ix:corollary:5:4}{5.4} with \(G=H\) and \(S'=K\) to
the homomorphism of \(K\)-groups
\[
  u':H_K\longrightarrow H_K,
  \qquad
  (h,k)\longmapsto(khk^{-1},k).
\]
It is central because \(H\) is commutative. Its pullback along the unit
section \(\epsilon:S\to K\) is the identity homomorphism of \(H\).
\emph{[The source prints ``of \(K\)'' here; the displayed endomorphism
of \(H_K\) shows that \(H\) is intended.]}
\SGARefLink{sga3:ix:corollary:5:4}{Corollary~5.4} therefore implies
that \(u'\) itself is the identity on
\(H_K\). Thus \(H\) is central in \(K\). \qed

Let us state the variants of the preceding results for central subgroups
of multiplicative type. Arguing as above, and using
\SGARefLink{sga3:ix:theorem:3:2:bis}{3.2~bis}, gives:

% Source page: Exp9 component-local p. 21; complete-reader p. 667.
\medskip
\noindent\SGARefTarget{sga3:ix:theorem:5:1:bis}\textbf{Theorem 5.1 bis.}
\textit{Let \(G\) be an \(S\)-group prescheme, and let \(H_1,H_2\) be
two subgroup preschemes of multiplicative type and of finite type, with
\((H_1)_s\) central. Suppose either that \(S\) is locally Noetherian or
that \(G\) is of finite presentation over \(S\). For every \(s\in S\)
such that \((H_1)_s=(H_2)_s\), there is an open neighborhood \(U\) of
\(s\) for which \(H_1|_U=H_2|_U\).}

\medskip
\noindent\SGARefTarget{sga3:ix:corollary:5:3:bis}\textbf{Corollary 5.3 bis.}
\textit{Under the preceding hypotheses, suppose moreover that \(G\) is
separated over \(S\). Let \(U\) be the set of points \(s\in S\) for
which \((H_1)_s=(H_2)_s\). Then \(U\) is open and closed in \(S\), and
\(H_1|_U=H_2|_U\).}

\medskip
\noindent\SGARefTarget{sga3:ix:corollary:5:4:bis}\textbf{Corollary 5.4 bis.}
\textit{Let \(S\) be a prescheme, let \(G\) be a group prescheme
separated and of finite presentation over \(S\), let \(S'\to S\) be a
covering morphism for the faithfully flat quasi-compact topology with
geometrically connected fibers, and let \(H'\) be a subgroup of
multiplicative type and of finite type of \(G'=G\times_SS'\).
Then there is a unique subgroup \(H\) of \(G\) such that
\(H'=H\times_SS'\). The subgroup \(H\) is of multiplicative type and
of finite type over \(S\), by \SGARefLink{sga3:ix:definition:1:1}{1.1} and \SGARefLink{sga3:ix:proposition:2:1:item:b:01}{2.1(b)}. If \(S'\to S\) admits a
section \(g\), then \(H=g^*H'\).}

\medskip
\noindent\SGARefTarget{sga3:ix:theorem:5:6}\textbf{Theorem 5.6.}
\textit{Let \(S\) be a prescheme and let \(u:H\to G\) be a homomorphism
of \(S\)-preschemes in groups, where \(H\) is of multiplicative type and
of finite type and \(G\) is of finite presentation over \(S\).}
\begin{enumerate}
  \item[(a)]\SGARefTarget{sga3:ix:theorem:5:6:item:a}
  \textit{Suppose that \(G\) has connected fibers. If \(s\in S\) and
  \(u_s:H_s\to G_s\) is central, then there is an open neighborhood
  \(U\) of \(s\) such that the homomorphism \(H|_U\to G|_U\) induced by
  \(u\) is central.}

  \item[(b)]\SGARefTarget{sga3:ix:theorem:5:6:item:b}
  \textit{If \(u_s:H_s\to G_s\) is central for every \(s\in S\), then
  \(u\) is central.}
\end{enumerate}

\emph{Proof.}
We first prove \SGARefLink{sga3:ix:theorem:5:6:item:a}{(a)}.%
\footnote{Editor's note: the original has been expanded in the proof of
\SGARefLink{sga3:ix:theorem:5:6:item:a}{(a)}.}
Arguing exactly as in the proof of
\SGARefLink{sga3:ix:theorem:5:1:proof:item:b}{5.1(b)}, we may suppose first that
\(S\) is local with closed point \(s\), then that
\(H=D_S(M)\) is diagonalizable, and finally that there exist a
finite-type \(\mathbb Z\)-subalgebra \(A_i\) of
\(A=\mathscr O_{S,s}\) and a homomorphism
\[
  u_i:D_{S_i}(M)\longrightarrow G_i
\]
whose base change is \(u\) and for which \((u_i)_{s_i}\) is central,
where \(s_i\) is the image of \(s\) in \(S_i\). We are thereby reduced
to the case in which \(S\) is Noetherian local with closed point \(s\),
and we must prove that \(u\) is central.

% Source page: Exp9 component-local p. 22; complete-reader p. 668.
Let \(K=\operatorname{Ker}(v,w)\), where
\[
  v,w:H\times_SG\rightrightarrows G
\]
are defined by
\SGARefTarget{sga3:ix:theorem:5:6:display:01}
\[
  v(h,g)=g,
  \qquad
  w(h,g)=\operatorname{int}(u(h))\cdot g
        =u(h)g\,u(h)^{-1}.
\]
Then \(K\) is a subgroup over \(G\) of
\[
  H_G=H\times_SG,
\]
and we must prove that \(K=H_G\). By
\SGARefLink{sga3:ix:corollary:4:9}{4.9}, it is enough to show that \(K\)
contains
\[
  {}_m(H_G)=({}_mH)\times_SG
\]
for every \(m>0\). We are thus reduced to the case \(H={}_mH\), so that
\(H\) is finite over \(S\).

Let \(e\) be the unit element of the fiber \(G_s\), and put
\[
  S^0=\operatorname{Spec}(\mathscr O_{G,e}),
  \qquad
  S_n^0=S^0\times_SS_n,
  \qquad
  S_n=\operatorname{Spec}(A/\mathfrak m^{n+1}).
\]
The scheme \(S^0\) is Noetherian local because \(G\) is of finite
presentation over the Noetherian scheme \(S\). The base change
\[
  K_{S^0}=K\times_GS^0
\]
is a subscheme of
\[
  H_{S^0}=H\times_SS^0,
\]
and \SGARefLink{sga3:ix:corollary:3:9}{3.9} gives
\[
  K_{S_n^0}=H_{S_n^0}
  \qquad\text{for every }n.
\]
\footnote{Editor's note: the original has been expanded and corrected,
replacing ``the unit section of \((G_H)_n\) over \(H_n\)'' by
\(H_{S_n^0}\).}
Since \(H_{S^0}\) is finite over \(S^0\),
\SGARefLink{sga3:ix:lemma:5:0}{5.0}, applied to the Noetherian local
ring \(\mathscr O_{G,e}\), gives
\[
  K_{S^0}=H_{S^0}.
\]
\emph{[The source prints the undefined local ring \(B\) here; the
Noetherian local ring \(\mathscr O_{G,e}\) just introduced is intended.]}

On the other hand, \(H\times_SG\) is Noetherian, being of finite
presentation over the Noetherian scheme \(S\). Thus the immersion
\[
  K\hookrightarrow H\times_SG
\]
is of finite type (cf. EGA I, 6.3.5), and \(K\) is of finite
presentation over \(G\). The equality
\[
  K_{S^0}=H_G\times_GS^0
\]
therefore implies, by EGA IV\(_3\), 8.8.2.4, that there is an open
neighborhood \(W\) of \(e\) in \(G\) such that
\[
  K\times_GW
  =H_G\times_GW
  =H\times_SW.
\]
Hence \(K\) contains the open neighborhood
\[
  V=H\times_SW
\]
of the unit section of \(G_H=G\times_SH\) over \(H\).

For every \(t\in H\), the fiber \(G_t\) is an algebraic group over
\(\kappa(t)\), hence Cohen--Macaulay by VI\(_A\), 1.1.1 and therefore
has no embedded components. It is also connected, hence irreducible by
VI\(_A\), 2.4, so its generic point is its unique associated point.
EGA IV\(_2\), 3.1.8 therefore shows that \(V_t\) is schematically dense
in \(G_t\). By \SGARefLink{sga3:ix:corollary:4:3}{4.3}, \(V\) is
schematically dense in \(G_H\).

Moreover, since \(K\) is an \(H\)-subgroup of \(G_H\), it induces on
each fiber \(G_h\) a subgroup \(K_h\). This subgroup contains an open
neighborhood of the unit element; because \(G_h\) is connected, it
follows that \(K_h=G_h\). Thus \(K=G_H\) set-theoretically. Finally,
\(K\) is a closed subprescheme of \(G_H\) containing the schematically
dense subprescheme \(V\), and hence \(K=G_H\). This proves \SGARefLink{sga3:ix:theorem:5:6:item:a}{(a)}.

Part~\SGARefLink{sga3:ix:theorem:5:6:item:b}{(b)} follows directly from
\SGARefLink{sga3:ix:corollary:3:9}{3.9}: to verify that a subprescheme
\(K\) of \(P=G\times_SH\) equals \(P\), it is enough to verify, after
every base change \(S'\to S\) in which \(S'\) is the spectrum of an
Artinian quotient of a local ring of \(S\), that \(K'=P'\). \qed

% End cursor: before Section 6 on Exp9 component-local p. 22,
% complete-reader p. 668.

% END INCLUDED SGA3_English_current_buildable_source_20260823/source/inputs/IX/tex/components/03_expose_IX_sections_4_5.tex

% BEGIN INCLUDED SGA3_English_current_buildable_source_20260823/source/inputs/IX/tex/components/04_expose_IX_sections_6_7.tex
% SGA 3, Expose IX, English reconstruction.
% Scope: Sections 6 and 7 in full.
% Authority: Polo--Gille Exp9-8nov09.pdf, component-local pp. 22 (partial)--30
% (partial), complete-reader pp. 668 (partial)--676 (partial),
% printed pp. 64--73 and the opening of printed p. 74.
% Start cursor: "6. Monomorphismes des groupes de type multiplicatif..."
% on component-local p. 22.
% End cursor: immediately before
% "8. Sous-groupes, groupes quotients et extensions..." on
% component-local p. 30.
% Jacob Reinhold's CC BY 4.0 Markdown lineage was used only as a phrasing
% comparison. The Polo--Gille PDF controls prose, formulas, numbering,
% notes, references, and diagrams.

\setcounter{footnote}{38}

% Source page: Exp9 component-local p. 22; complete-reader p. 668;
% printed p. 64.
\SGARefTarget{sga3:ix:section:6}
\section*{6. Monomorphisms of Groups of Multiplicative Type, and
Canonical Factorization of a Homomorphism from Such a Group}
\addcontentsline{toc}{section}{6. Monomorphisms of Groups of
Multiplicative Type, and Canonical Factorization}

\medskip
\noindent\SGARefTarget{sga3:ix:lemma:6:1}\textbf{Lemma 6.1.}
\textit{Let \(S\) be a quasi-compact prescheme, and let \(G\) be a
commutative \(S\)-group prescheme, of finite presentation and
quasi-finite over \(S\). Then there is an integer \(n>0\) such that
\[
  n\cdot\operatorname{id}_G=0,
  \qquad\text{that is,}\qquad
  G={}_nG.
\]}

\emph{Proof.}
If \(S\) is the spectrum of a field \(k\), then \(G\) is finite over
\(k\), and by VII\textsubscript{A} \SGARefLink{sga3:VIIA:8.5}{8.5} it is enough to take
\[
  n=\deg(G/k).
\]
Suppose now that \(S=\operatorname{Spec}(A)\), where \(A\) is local
Artinian. Let \(k\) be the residue field of \(A\), put
\[
  G_0=G\otimes_A k,
  \qquad
  n_0=\deg(G_0/k),
\]
and distinguish two cases.

\smallskip
\SGARefTarget{sga3:ix:lemma:6:1:proof:item:a}\noindent\textit{(a) \(k\) has characteristic zero.}
Then \(G_0\) is separable over \(k\), and hence \(G\) is unramified over
\(S\). The unit section of \(G\) is therefore an open immersion, so
\[
  {}_nG=\operatorname{Ker}(n\cdot\operatorname{id}_G)
\]
is an open subscheme of \(G\). For it to equal \(G\), it is consequently
enough that it do so set-theoretically; thus one may take \(n=n_0\).

\smallskip
\SGARefTarget{sga3:ix:lemma:6:1:proof:item:b}\noindent\textit{(b) \(k\) has characteristic \(p>0\).}
Let \(\mathfrak m\) be the maximal ideal of \(A\), and put
\[
  S_r=\operatorname{Spec}(A/\mathfrak m^{r+1})
  \qquad (r\geq0).
\]
Choose an integer \(m\) such that \(\mathfrak m^{m+1}=0\). We claim that
one may take \(n=p^mn_0\). Indeed, arguing by induction on \(m\) and
putting \(n'=p^{m-1}n_0\), we may suppose that
\(n'\cdot\operatorname{id}_G\) induces the zero endomorphism on
\[
  G_{m-1}=G\times_S S_{m-1}.
\]
\footnote{Editor's note: the original has been expanded in what follows,
with the results of Expos\'e III that are used made explicit.}
Let \(E\) be the set of endomorphisms \(u\) of \(G\) having this
property, and let \(K\) be the kernel of the morphism of group functors
\[
  G\longrightarrow \prod_{S_{m-1}/S}G
\]
(cf.\ \SGARefLink{sga3:iii:remark:0-1-2}{III 0.1.2}). Then
\[
  E=\operatorname{Hom}_{S\text{-}\mathrm{gr}}(G,K);
\]
in particular, the abelian group law on \(E\) is induced by that on
\(K\). By \SGARefLink{sga3:iii:corollary:0-9}{III 0.9}, \(K\) is the \(S\)-group functor which, to every
\(g:T\to S\), assigns the \(k\)-vector space
\SGARefTarget{sga3:ix:lemma:6:1:display:01}
\[
  \operatorname{Hom}_{\mathscr O_{T_0}}
  \bigl(
    g_0^*(\Omega^1_{G_0/S_0}),
    \mathfrak m^m\mathscr O_T
  \bigr).
\]
Thus \(pu=0\) for every \(u\in E\). It follows that
\[
  pn'\cdot\operatorname{id}_G=0,
  \qquad\text{that is,}\qquad
  p^mn_0\cdot\operatorname{id}_G=0.
\]

Suppose now that \(S\) is Noetherian; this is the case to which the
general assertion reduces by the usual reduction to the Noetherian
case.\footnote{Editor's note: since \(S\) is quasi-compact, it is covered
by finitely many affine open subsets, so one reduces to
\(S=\operatorname{Spec}(A)\). Since \(G\) is of finite presentation over
\(S\), EGA IV\(_3\), 8.8.2 may then be applied.}
It is enough to combine the foregoing with \SGARefLink{sga3:ix:lemma:6:2}{the following lemma}.

\medskip
\noindent\SGARefTarget{sga3:ix:lemma:6:2}\textbf{Lemma 6.2.}
\textit{Let \(X\) be a quasi-finite prescheme over a Noetherian
prescheme \(S\), and let \((X_i)_{i\in I}\) be an increasing filtered
family of subpreschemes having the following property:
\[
\left\{
\begin{aligned}
 &\text{for every \(s\in S\) and every \(n\geq0\), putting}\\
 &S_{s,n}
   =\operatorname{Spec}
      \bigl(\mathscr O_{S,s}/\mathfrak m^{n+1}\bigr),
   \text{ there is an \(i\in I\) such that}\\
 &X_i\times_S S_{s,n}=X\times_S S_{s,n}.
\end{aligned}
\right.
\]
Then there is an \(i\in I\) such that \(X_i=X\).}

\emph{Proof.}
Since \(S\) is Noetherian, there is a maximal open subset \(U\) on which
\[
  X|_U=X_i|_U
\]
for all sufficiently large \(i\). We shall prove that \(U=S\). In other
words, if \(U\neq S\), we shall find an open subset \(U'\) strictly
larger than \(U\), and an \(i\), such that
\[
  X|_{U'}=X_i|_{U'}.
\]
Localizing at a maximal point \(s\) of \(S-U\), we reduce to the case in
which \(S\) is local with closed point \(s\), and
\[
  U=S-\{s\}.
\]
Indeed, put \(S'=\operatorname{Spec}(\mathscr O_{S,s})\). If
\[
  X\times_S S'=X_i\times_S S'
\]
for some \(i\),\footnote{Editor's note: \(X-\{s\}\) has been corrected
above to \(S-\{s\}\). Recall also that a quasi-finite morphism is assumed
to be of finite type; cf.\ EGA II, 6.2.3, and EGA III\(_2\),
Err\textsubscript{III} 20 for the definition of ``locally
quasi-finite.'' Thus here, since \(S\) is Noetherian, \(X\) is of finite
presentation over \(S\); every immersion \(X_i\hookrightarrow X\) is of
finite type by EGA I, 6.3.5, so \(X_i\) is of finite presentation over
\(S\), and EGA IV\(_3\), 8.8.2 applies.}
there is an open neighborhood \(V\) of \(s\) such that
\[
  X|_V=X_i|_V.
\]
Taking \(i\) large enough that \(X_i|_U=X|_U\), we then have
\[
  X|_W=X_i|_W,
  \qquad W=U\cup V.
\]

For all sufficiently large \(i\), the equality \(X|_U=X_i|_U\) shows
that \(X_i\) is a closed subprescheme of \(X\), defined by an ideal
\(\mathscr I^{(i)}\) whose support is contained in
\[
  X_s=X_0=X\times_S S_0,
\]
where
\[
  S_n=\operatorname{Spec}(A/\mathfrak m^{n+1}),
  \qquad S=\operatorname{Spec}(A).
\]
Since \(X_0\) is quasi-finite over \(S_0\), it is a finite closed subset
of the Noetherian prescheme \(X\). Hence \(\mathscr I^{(i)}\) is a
module of finite length. There is therefore an integer \(n\geq0\) such
that
\[
  \mathscr I^{(i)}
  \cap\mathfrak m^{n+1}\mathscr O_X=0.
\]
On the other hand, by the hypothesis of the lemma, after increasing
\(i\) if necessary we may suppose that the image of
\(\mathscr I^{(i)}\) in
\[
  \mathscr O_{X_n}
  =\mathscr O_X/\mathfrak m^{n+1}\mathscr O_X
\]
is zero. It follows that \(\mathscr I^{(i)}=0\), and therefore
\(X_i=X\). \qed

\medskip
\noindent\SGARefTarget{sga3:ix:lemma:6:3}\textbf{Lemma 6.3.}
\textit{Let \(S\) be a prescheme, let \(K\) be a group prescheme over
\(S\), locally of finite presentation, and let \(s\in S\). Suppose that
\(K_s\) is quasi-finite (respectively, unramified) over \(\kappa(s)\) at
the unit element. Then there is an open neighborhood \(U\) of \(s\) such
that \(K|_U\) is locally quasi-finite (respectively, unramified) over
\(U\).}%
\begingroup
\renewcommand{\thefootnote}{\(\ast\)}%
\footnote{Cf.\ VI\textsubscript{B} \SGARefLink{concl:sga3-VIB-2.5-ii}{2.5(ii)} for a more general
statement.}%
\addtocounter{footnote}{-1}%
\endgroup

\emph{Proof.}
Let \(V\) be the set of points \(x\in K\) having the following
property: if \(t\) is the image of \(x\) in \(S\), then \(K_t\) is
quasi-finite (respectively, unramified) over \(\kappa(t)\) at \(x\);
that is, \(x\) is isolated in \(K_t\) (respectively, it is isolated and
its local ring in \(K_t\) is a separable extension of \(\kappa(t)\)).
The set \(V\) is open because \(K\) is locally of finite presentation
over \(S\).\footnote{Editor's note: cf.\ EGA IV\(_3\), 13.1.4 and EGA
IV\(_4\), 17.4.1. In addition, \(t\), rather than \(s\), has been used
for the image of \(x\).}
If \(\epsilon\) denotes the unit section of \(K\), then
\(\epsilon^{-1}(V)\) is open. By hypothesis it contains \(s\), and is
therefore an open neighborhood \(U\) of \(s\). This \(U\) has the
required property. Indeed, if \(t\in U\), then \(K_t\), being a group
locally of finite type over \(\kappa(t)\), is locally quasi-finite
(respectively, unramified) over \(\kappa(t)\) because it has that
property at \(\epsilon(t)\); cf.\ VI\textsubscript{B} \SGARefLink{prop:sga3-VIB-1.3}{1.3}. \qed

Combining \SGARefLink{sga3:ix:lemma:6:1}{6.1} and
\SGARefLink{sga3:ix:lemma:6:3}{6.3}, one obtains:

\medskip
\noindent\SGARefTarget{sga3:ix:theorem:6:4}\textbf{Theorem 6.4.}
\textit{Let \(S\) be a prescheme, let \(H\) be an \(S\)-group of
multiplicative type and of finite type, and let \(K\) be a closed group
subprescheme of \(H\), of finite presentation over \(S\). Let \(s\in S\)
be such that \(K_s\) is finite. Then there are an open neighborhood
\(U\) of \(s\) and an integer \(n>0\) such that
\[
  K|_U\subset {}_nH|_U.
\]
In particular, since \({}_nH\) is finite over \(S\), the group
\(K|_U\) is finite over \(U\).}

Nakayama's lemma now gives:

\medskip
\noindent\SGARefTarget{sga3:ix:corollary:6:5}\textbf{Corollary 6.5.}
\textit{With the preceding notation, suppose that \(K_s\) is the unit
group. Then there is an open neighborhood \(U\) of \(s\) such that
\(K|_U\) is the unit group.}

\medskip
\noindent\SGARefTarget{sga3:ix:corollary:6:6}\textbf{Corollary 6.6.}
\textit{Let \(u:H\to G\) be a homomorphism of \(S\)-group preschemes,
where \(H\) is of multiplicative type and of finite type and \(G\) is
separated over \(S\). Suppose in addition that either \(S\) is locally
Noetherian or \(G\) is locally of finite presentation over
\(S\).\footnote{Editor's note: it is enough to suppose that \(G\) is
locally of finite type over \(S\). By EGA IV\(_4\), 1.4.3(v), this
implies, since \(H\to S\) is of finite presentation, that \(H\to G\) is
locally of finite presentation; the same is then true of \(K\to S\),
which is obtained from it by base change.}
If \(s\in S\) and
\[
  u_s:H_s\longrightarrow G_s
\]
is a monomorphism, then there is an open neighborhood \(U\) of \(s\)
such that \(u\) induces a monomorphism
\[
  H|_U\longrightarrow G|_U.
\]}

\emph{Proof.}
Put \(K=\operatorname{Ker}(u)\). The hypothesis on \(u_s\) says that
\(K_s\) is the unit group, and the desired conclusion says that
\(K|_U\) is the unit group for a suitable \(U\). Since \(G\) is
separated over \(S\), the group \(K\) is a closed subgroup of \(H\).
If \(S\) is not assumed locally Noetherian, but \(G\) is assumed locally
of finite presentation over \(S\), then \(K\) is locally of finite
presentation over \(S\).\footnotemark[43] It is therefore of finite
presentation over \(S\), since it is separated over \(S\) (as \(H\)
is) and quasi-compact over \(S\) (being closed in \(H\), which is
quasi-compact over \(S\)). Corollary
\SGARefLink{sga3:ix:corollary:6:5}{6.5} now applies. \qed

\medskip
\noindent\SGARefTarget{sga3:ix:remark:6:6:1}\textbf{Remark 6.6.1.}%
\footnote{Editor's note: the number \SGARefLink{sga3:ix:remark:6:6:1}{6.6.1} has been added for use in
later references.}
\textit{Under the hypotheses of
\SGARefLink{sga3:ix:corollary:6:6}{6.6}, note that if, in addition,
\(G\) is affine over \(S\) (respectively, of finite presentation over
\(S\)), then
\[
  H|_U\longrightarrow G|_U
\]
is even a closed immersion, as was observed in \SGARefLink{sga3:ix:corollary:2:5}{2.5} (respectively,
\SGARefLink{sga3:ix:remark:2:6}{2.6}).}

\medskip
\noindent\SGARefTarget{sga3:ix:corollary:6:7}\textbf{Corollary 6.7.}
\textit{Let \(u:H\to G\) be a homomorphism of \(S\)-group preschemes,
where \(H\) is of multiplicative type and of finite type. Suppose that
either \(S\) is locally Noetherian or \(G\) is locally of finite
presentation over \(S\). If every homomorphism on fibers
\[
  u_s:H_s\longrightarrow G_s
\]
is a monomorphism, then \(u\) is a monomorphism.}

\emph{Proof.}
The argument is the same as in
\SGARefLink{sga3:ix:corollary:6:6}{6.6}. Here the hypothesis that \(G\)
be separated over \(S\) is unnecessary in order to ensure that
\(K=\operatorname{Ker}(u)\) is closed in \(H\), since the assumption
that all the \(u_s\) are monomorphisms implies that \(K\) is
set-theoretically reduced to the unit section. \qed

\medskip
\noindent\SGARefTarget{sga3:ix:theorem:6:8}\textbf{Theorem 6.8.}
\begin{itshape}
Let \(u:H\to G\) be a homomorphism of \(S\)-group preschemes,
where \(H\) is of multiplicative type and of finite type and \(G\) is
separated over \(S\). Suppose in addition that either \(S\) is locally
Noetherian or \(G\) is of finite presentation over \(S\). Then
\[
  K=\operatorname{Ker}(u)
\]
is a subgroup of \(H\) of multiplicative type and of finite type, and
\(u\) factors as
\SGARefTarget{sga3:ix:theorem:6:8:diagram:01}
\[
\begin{tikzcd}[column sep=large]
  H
    \arrow[r,"u'"]
  & H/K
    \arrow[r,"u''"]
  & G,
\end{tikzcd}
\]
where \(H/K\) is of multiplicative type and of finite type,
\(u'\) is the canonical homomorphism and is faithfully flat (and
affine), and \(u''\) is a monomorphism.
\end{itshape}

\emph{N.B.}
As noted in \SGARefLink{sga3:ix:remark:6:6:1}{6.6.1}, the homomorphism
\(u''\) is a closed immersion if \(G\) is affine over \(S\), or if
\(G\) is of finite presentation over \(S\).

\emph{Proof.}
It is enough to prove that \(K\) is of multiplicative type; the rest
then follows from \SGARefLink{sga3:ix:proposition:2:7}{2.7} and \SGARefLink{sga3:iv:remark:5:2:6}{IV 5.2.6}.

Suppose first that \(G\) is of finite presentation over \(S\). Since
this hypothesis is stable under base change, in proving that \(K\) is
of multiplicative type we may reduce to the case in which \(H\) is
diagonalizable:
\[
  H=D_S(M),
\]
where \(M\) is a commutative group of finite type. Let \(s\in S\). The
group \(K_s\) is a closed subgroup of
\[
  H_s=D_{\kappa(s)}(M).
\]
By \SGARefLink{sga3:ix:proposition:8:1}{8.1},\footnote{Editor's note:
the erroneous reference to IV 4.7.5 has been corrected. Note that
\SGARefLink{sga3:ix:section:8}{Section 8} of the present expos\'e is independent of Sections \SGARefLink{sga3:ix:section:2}{2} through
\SGARefLink{sga3:ix:section:7}{7}.}
it is therefore of the form \(D_{\kappa(s)}(N)\), where \(N\) is a
quotient of \(M\). Put
\[
  K'=D_S(N).
\]
Then \(K'\) is a subgroup of \(H\) of multiplicative type. Let
\[
  v':K'\longrightarrow G
\]
be the homomorphism induced by \(u\).\footnote{Editor's note:
``induced by \(G\)'' has been corrected to ``induced by \(u\).''}
By construction, \(v'_s\) is the unit homomorphism. Corollary
\SGARefLink{sga3:ix:corollary:5:2}{5.2} therefore gives an open
neighborhood \(U\) of \(s\) such that
\[
  v'|_U:K'|_U\longrightarrow G|_U
\]
is the unit homomorphism. Replacing \(S\) by \(U\), we may suppose that
\(v'\) itself is the unit homomorphism, and hence that \(u\) factors as
\SGARefTarget{sga3:ix:theorem:6:8:diagram:02}
\[
\begin{tikzcd}[column sep=large]
  H
    \arrow[r,"u'"]
  & H/K'
    \arrow[r,"u''"]
  & G.
\end{tikzcd}
\]
The homomorphism \(u''_s\) is obtained from \(u_s\) by factoring through
\[
  H_s\longrightarrow H_s/\operatorname{Ker}(u_s),
\]
and is therefore a monomorphism by \SGARefLink{sga3:iv:remark:5:2:6}{IV 5.2.6}. By
\SGARefLink{sga3:ix:corollary:6:6}{6.6}, after restricting to a suitable
open neighborhood \(U\) of \(s\), the homomorphism
\[
  u''|_U:(H/K')|_U\longrightarrow G|_U
\]
is a monomorphism. Thus, after this restriction,
\[
  \operatorname{Ker}(u)
  =\operatorname{Ker}(u')
  =K',
\]
which proves that \(\operatorname{Ker}(u)\) is of multiplicative type.

The same proof applies if, instead of assuming \(G\) of finite
presentation over \(S\), one assumes \(S\) locally Noetherian, at least
when \(H\) is diagonalizable. Without this additional hypothesis on
\(H\), one must show that there is a covering base change
\[
  S'\longrightarrow S,
\]
with \(S'\) locally Noetherian, which trivializes \(H\). This will be
proved in the next expos\'e, \SGARefLink{sga3:x:corollary:4-6}{X 4.6}. \qed

% Source page: Exp9 component-local p. 26; complete-reader p. 672;
% printed pp. 68--69.
\SGARefTarget{sga3:ix:section:7}
\section*{7. Algebraicity of Formal Homomorphisms into an Affine Group}
\addcontentsline{toc}{section}{7. Algebraicity of Formal Homomorphisms
into an Affine Group}

\medskip
\noindent\SGARefTarget{sga3:ix:theorem:7:1}\textbf{Theorem 7.1.}
\textit{Let \(A\) be a Noetherian ring equipped with an ideal \(I\) for
which \(A\) is separated and complete in the \(I\)-adic topology. Put
\[
  S=\operatorname{Spec}(A),
  \qquad
  S_n=\operatorname{Spec}(A/I^{n+1}),
\]
and let \(H,G\) be \(S\)-group preschemes, with \(H\) of isotrivial
multiplicative type and \(G\) affine. Then the canonical map
\SGARefTarget{sga3:ix:theorem:7:1:equation:x}
\begin{equation}
  \tag{x}
  \theta:
  \operatorname{Hom}_{S\text{-}\mathrm{gr}}(H,G)
  \longrightarrow
  \varprojlim_n
  \operatorname{Hom}_{S_n\text{-}\mathrm{gr}}(H_n,G_n)
\end{equation}
is bijective, where \(H_n,G_n\) are the \(S_n\)-groups obtained from
\(H,G\) by the base change \(S_n\to S\).}

\emph{Proof.}
Suppose first that \(H\) is diagonalizable, hence of the form
\[
  H=\operatorname{Spec}A(M),
\]
where \(B=A(M)\) is the algebra of the commutative group \(M\) with
coefficients in \(A\). We may also write
\[
  G=\operatorname{Spec}(C),
\]
where \(C\) is an \(A\)-algebra equipped with a diagonal map satisfying
the usual axioms. Homomorphisms of \(S\)-groups \(H\to G\) correspond
bijectively to homomorphisms of \(A\)-algebras
\[
  \varphi:C\longrightarrow B
\]
compatible with the diagonal maps; that is, such that, for every
\(f\in C\),
\SGARefTarget{sga3:ix:theorem:7:1:display:01}
\[
  \Delta_H(\varphi(f))
  =(\varphi\otimes\varphi)(\Delta_G(f)).
\]
There is an analogous description of the homomorphisms of \(S_n\)-groups
\(H_n\to G_n\), given by homomorphisms of \(A_n\)-algebras
\[
  \varphi_n:C_n\longrightarrow B_n,
\]
where
\[
  A_n=A/I^{n+1},
  \qquad
  B_n=B\otimes_AA_n,
  \qquad
  C_n=C\otimes_AA_n.
\]
\emph{[The source prints \(C=C\otimes_AA_n\) in the last equality;
\(C_n\) is intended.]}
Put
\SGARefTarget{sga3:ix:theorem:7:1:display:02}
\[
  \widehat B=\varprojlim_n B_n,
  \qquad
  \widehat C=\varprojlim_n C_n.
\]
Then
\[
  \widehat B=\varprojlim_n A_n(M)
\]
identifies with a submodule of the product \(A^M\), namely the submodule
of families \((a_m)_{m\in M}\) of elements of \(A\) that tend to zero
for the \(I\)-adic topology along the filter of complements of finite
subsets of \(M\). In particular, the canonical homomorphism
\[
  B\longrightarrow\widehat B
\]
is injective.\footnote{Editor's note: the argument that follows has
been expanded.}
The projective system
\[
  \theta(\varphi)=(\varphi_n)_{n\geq0}
\]
defines a homomorphism
\[
  \widehat\varphi:\widehat C\longrightarrow\widehat B
\]
such that \SGARefLink{sga3:ix:theorem:7:1:diagram:01}{the following diagram} is commutative:
\SGARefTarget{sga3:ix:theorem:7:1:diagram:01}
\[
\begin{tikzcd}[column sep=large,row sep=large]
  C
    \arrow[r,"\varphi"]
    \arrow[d]
  & B
    \arrow[d]
  \\
  \widehat C
    \arrow[r,"\widehat\varphi"]
  & \widehat B.
\end{tikzcd}
\]

We prove that \(\theta\) is surjective. Let
\((\varphi_n)_{n\geq0}\) be a projective system; we must show that it is
obtained by reduction from some \(\varphi\) in the left-hand side of
\SGARefLink{sga3:ix:theorem:7:1:equation:x}{\((x)\)}. A priori, the system \((\varphi_n)\) defines a homomorphism
on the completed algebras,
\[
  \widehat\varphi:\widehat C\longrightarrow\widehat B.
\]
It is enough to prove that the composite
\[
  \Phi:
  C\longrightarrow\widehat C
   \xrightarrow{\widehat\varphi}\widehat B
\]
sends \(C\) into \(B\). Indeed, in that case it gives a homomorphism of
\(A\)-algebras
\SGARefTarget{sga3:ix:theorem:7:1:display:phi:01}\[
  \varphi:C\longrightarrow B
\]
whose reductions are the \(\varphi_n\). \SGARefLink{sga3:ix:theorem:7:1:display:phi:01}{This homomorphism} is immediately
seen to be compatible with the diagonal maps, since the \(\varphi_n\)
are so, and since
\[
  B\otimes_AB\longrightarrow\widehat{B\otimes_AB}
\]
is injective, as follows from the same argument after replacing \(M\)
by \(M\times M\).

Passing to the completions, the diagonal homomorphism \(\Delta_H\)
defines
\SGARefTarget{sga3:ix:theorem:7:1:display:03}
\[
  \widehat\Delta_H:
  \widehat B\longrightarrow\widehat{B\otimes_AB}.
\]
Taking the projective limit of the analogous diagrams defined by the
\(\varphi_n\), one obtains the commutative diagram
\SGARefTarget{sga3:ix:theorem:7:1:diagram:02}
\[
\begin{tikzcd}[column sep=huge,row sep=large]
  C
    \arrow[r,"\Phi"]
    \arrow[d,"\Delta_G"']
  & \widehat B
    \arrow[d,"\widehat\Delta_H"]
  \\
  C\otimes_A C
    \arrow[r,"\Psi"]
  & \widehat{B\otimes_A B},
\end{tikzcd}
\]
where \(\Psi\) is the composite
\SGARefTarget{sga3:ix:theorem:7:1:display:04}
\[
  C\otimes_A C
  \xrightarrow{\ \Phi\otimes\Phi\ }
  \widehat B\otimes_A\widehat B
  \longrightarrow
  \widehat{B\otimes_A B};
\]
the last arrow is the evident canonical homomorphism. Consequently, for
every \(f\in C\), the element \(\Phi(f)\in\widehat B\) has the property
that its image under \(\widehat\Delta_H\) is a ``decomposable'' element
of \(\widehat{B\otimes_A B}\), namely an element in the image of
\(\widehat B\otimes_A\widehat B\).

Let \((e_m)_{m\in M}\) be the canonical basis of \(A(M)\), and let
\((e_{m,m'})\) be that of
\[
  A(M\times M)=A(M)\otimes_AA(M).
\]
Since
\[
  \Delta_H(e_m)=e_m\otimes e_m=e_{m,m}
  \qquad (m\in M),
\]
it remains only to apply \SGARefLink{sga3:ix:lemma:7:2}{the following lemma}.

\medskip
\noindent\SGARefTarget{sga3:ix:lemma:7:2}\textbf{Lemma 7.2.}
\begin{itshape}
Let \(A\) be a Noetherian ring, let \(M\) be a set, and let
\((a_{m,m'})\) be a family of elements of \(A\), indexed by
\(M\times M\), satisfying:
\begin{enumerate}
  \item[(i)] \SGARefTarget{sga3:ix:lemma:7:2:item:i}
  \(a_{m,m'}=0\) if \(m\neq m'\); that is, the support of the family is
  contained in the diagonal of \(M\times M\);

  \item[(ii)] \SGARefTarget{sga3:ix:lemma:7:2:item:ii}
  one has
  \SGARefTarget{sga3:ix:lemma:7:2:display:01}
  \[
    a_{m,m'}=\sum_i b_m^ic_{m'}^i,
  \]
  where the \(b^i,c^i\) are finitely many elements of \(A^M\); in other
  words, \((a_{m,m'})\) belongs to the image of the canonical
  homomorphism
  \[
    A^M\otimes_AA^M\longrightarrow A^{M\times M}.
  \]
\end{enumerate}
Then the support of \((a_{m,m'})\) is finite.
\end{itshape}

\emph{Proof.}
By \SGARefLink{sga3:ix:lemma:7:2:item:i}{(i)}, the family \((a_{m,m'})\) is determined by the elements
\[
  a_m=a_{m,m}.
\]
For \(x=(x_n)_{n\in M}\in A(M)\), put
\[
  (u\cdot x)_m=\sum_{m'}a_{m,m'}x_{m'}.
\]
Thus \((a_{m,m'})\) is viewed as the matrix of a homomorphism
\[
  u:A(M)\longrightarrow A^M.
\]
Condition \SGARefLink{sga3:ix:lemma:7:2:item:i}{(i)} gives simply
\[
  (u\cdot x)_m=a_mx_m.
\]
On the other hand, condition \SGARefLink{sga3:ix:lemma:7:2:item:ii}{(ii)} gives
\[
  u\cdot x\in\sum_i A\cdot b^i.
\]
Hence \(u(A(M))\) lies in a finitely generated \(A\)-module. If
\((e_m)_{m\in M}\) denotes the canonical basis of
\[
  A(M)\subset A^M,
\]
then all the elements \(a_me_m\) lie in a finitely generated
\(A\)-module. Since \(A\) is Noetherian, the module generated by these
elements is itself finitely generated. Because the \(e_m\) are linearly
independent, all but finitely many of the \(a_m\) are zero. This proves
the lemma, and hence Theorem \SGARefLink{sga3:ix:theorem:7:1}{7.1} when
\(H\) is diagonalizable. \qed

We now prove the general case of
\SGARefLink{sga3:ix:theorem:7:1}{7.1}, in which \(H\) is only assumed
isotrivial. There is then a finite \'etale surjective morphism
\[
  S'\longrightarrow S
\]
such that
\[
  H'=H\times_SS'
\]
is diagonalizable. We shall use only the fact that \(S'\to S\) is finite
and a covering for the faithfully flat quasi-compact topology, or simply
for the canonical topology on \((\mathbf{Sch})\). Thus ``\'etale'' could
be replaced by ``flat.''

Put
\[
  S''=S'\times_SS',
\]
and similarly introduce
\[
  S'_n
  \quad\text{and}\quad
  S''_n
  =S''\times_SS_n
  =S'_n\times_{S_n}S'_n,
\]
together with \(H',G',H'',G'',H'_n,G'_n,H''_n,G''_n\), obtained from
\(H,G\) by the evident base changes. The groups \(H'\) and \(H''\) are
diagonalizable. Moreover, \(S'\), and hence \(S''\), is affine. If
\[
  S'=\operatorname{Spec}(A'),
  \qquad
  S''=\operatorname{Spec}(A''),
\]
then \(A'\) and \(A''\) are separated and complete for the topologies
defined by \(IA'\) and \(IA''\), respectively, because both are finite
over \(A\).

Since \(S'\to S\) and \(S'_n\to S_n\) are covering morphisms, we obtain
a commutative diagram of maps of sets whose two rows are exact:
\SGARefTarget{sga3:ix:theorem:7:1:diagram:03}
\[
\begin{tikzcd}[
  column sep=small,
  row sep=large,
  cells={nodes={font=\small}}
]
  \operatorname{Hom}_{S\text{-}\mathrm{gr}}(H,G)
    \arrow[r]
    \arrow[d,"u"]
  & \operatorname{Hom}_{S'\text{-}\mathrm{gr}}(H',G')
    \arrow[r,shift left=.6ex]
    \arrow[r,shift right=.6ex]
    \arrow[d,"u'"]
  & \operatorname{Hom}_{S''\text{-}\mathrm{gr}}(H'',G'')
    \arrow[d,"u''"]
  \\
  \varprojlim_n
    \operatorname{Hom}_{S_n\text{-}\mathrm{gr}}(H_n,G_n)
    \arrow[r]
  & \varprojlim_n
    \operatorname{Hom}_{S'_n\text{-}\mathrm{gr}}(H'_n,G'_n)
    \arrow[r,shift left=.6ex]
    \arrow[r,shift right=.6ex]
  & \varprojlim_n
    \operatorname{Hom}_{S''_n\text{-}\mathrm{gr}}(H''_n,G''_n).
\end{tikzcd}
\]
The second row is exact because it is the projective limit of the exact
diagrams associated with the various \(S'_n\to S_n\). By the
diagonalizable case already proved, \(u'\) and \(u''\) are bijective.
Thus \(u\) is bijective as well, completing the proof of
\SGARefLink{sga3:ix:theorem:7:1}{7.1}. \qed

\medskip
\noindent\SGARefTarget{sga3:ix:corollary:7:3}\textbf{Corollary 7.3.}
\textit{Under the hypotheses of
\SGARefLink{sga3:ix:theorem:7:1}{7.1}, suppose in addition that \(G\)
is smooth over \(S\), and let
\[
  u_0:H_0\longrightarrow G_0
\]
be a homomorphism of \(S_0\)-groups. Then there is a homomorphism of
\(S\)-groups
\[
  u:H\longrightarrow G
\]
lifting \(u_0\). Any two such liftings \(u,u'\) are conjugate by an
element \(g\in G(S)\) whose reduction is the unit element of
\(G_0(S_0)\).}

\emph{Proof.}
This follows by combining
\SGARefLink{sga3:ix:theorem:3:6}{3.6} and
\SGARefLink{sga3:ix:theorem:7:1}{7.1}. To construct \(u\), construct
the \(u_n\) successively, as is possible by \SGARefLink{sga3:ix:theorem:3:6}{3.6}; \SGARefLink{sga3:ix:theorem:7:1}{Theorem 7.1} then shows
that the system \((u_n)\) comes from an \(S\)-homomorphism \(u\). Given
two liftings \(u,u'\), construct successively elements \(g_n\) such
that
\[
  g_0=1,
  \qquad
  g_n\text{ is the reduction of }g_{n+1},
  \qquad
  u'_n=\operatorname{int}(g_n)u_n.
\]
Again, this is possible by \SGARefLink{sga3:ix:theorem:3:6}{3.6}. Since \(A\) is separated and complete,
the \(g_n\) come from an element \(g\in G(S)\). Finally, the injectivity
in \SGARefLink{sga3:ix:theorem:7:1}{7.1} gives
\[
  u'=\operatorname{int}(g)u.
\]
\qed

\medskip
\noindent\SGARefTarget{sga3:ix:remark:7:4}\textbf{Remark 7.4.}
\textit{Compare \SGARefLink{sga3:ix:theorem:7:1}{7.1} with EGA III,
5.4.1, which implies that the statement of \SGARefLink{sga3:ix:theorem:7:1}{7.1} remains true if, instead
of assuming \(H\) of multiplicative type and \(G\) affine, one assumes
that \(H\) is proper over \(S\), and that \(G\) is separated and locally
of finite type over \(S\). A statement such as \SGARefLink{sga3:ix:theorem:7:1}{7.1} without a properness
hypothesis is rather exceptional, and should be understood here as one
aspect of the great ``rigidity'' of the structure of a group of
multiplicative type.}

\textit{The analogous statement with
\[
  G=H=\mathbf G_a
\]
is false in general. Indeed, take \(A\) of characteristic \(p>0\), and
define the \(u_n\), by reduction modulo \(I^{n+1}\), from an additive
formal power series
\SGARefTarget{sga3:ix:remark:7:4:display:01}
\[
  u(T)=\sum_{n\geq0}a_nT^{p^n},
\]
where the \(a_n\in A\) tend to zero for the \(I\)-adic topology, but not
for the discrete topology.}

\textit{The statement of \SGARefLink{sga3:ix:theorem:7:1}{7.1} also becomes false if the hypothesis that
\(G\) be affine is omitted, even when \(H=\mathbf G_m\). An example,
with \(A\) a complete discrete valuation ring, is obtained from an
elliptic curve over the fraction field \(K\) of \(A\) whose reduction,
in the N\'eron--Kodaira theory, say, is the group \(\mathbf G_m\) over the
residue field \(k\).\footnote{Editor's note: such an example could be
made explicit.} One then has a smooth commutative group scheme \(G\)
over \(S\), with special fiber \(\mathbf G_{m,k}\). By \SGARefLink{sga3:ix:theorem:3:6}{3.6} this permits
the construction of a projective system of isomorphisms
\[
  u_n:H_n\xrightarrow{\ \sim\ }G_n,
  \qquad H=\mathbf G_{m,A},
\]
while the generic fiber of \(G\) is an abelian variety. Consequently,
there is no homomorphism of \(S\)-groups \(H\to G\) other than the zero
homomorphism.}

% Hard stop: before Section 8 on Exp9 component-local p. 30,
% complete-reader p. 676, printed p. 74.

% END INCLUDED SGA3_English_current_buildable_source_20260823/source/inputs/IX/tex/components/04_expose_IX_sections_6_7.tex

% BEGIN INCLUDED SGA3_English_current_buildable_source_20260823/source/inputs/IX/tex/components/05_expose_IX_section_8_and_bibliography.tex
% SGA 3, Exposé IX, English reconstruction.
% Scope: §8 in full, terminal bibliography, and editor notes 49--55.
% Authority: Polo--Gille Exp9-8nov09.pdf, component-local pp. 30 (partial)--32,
% complete-reader pp. 676 (partial)--678, printed pp. 74--76.
% Jacob Reinhold's CC BY 4.0 Markdown lineage was used only as a phrasing
% comparison. The Polo--Gille PDF controls text, formulas, numbering, and notes.

\setcounter{footnote}{48}
\SGARefTarget{sga3:ix:section:8}
\section*{8. Subgroups, Quotient Groups, and Extensions of Groups of
Multiplicative Type over a Field\textsuperscript{*}}
\addcontentsline{toc}{section}{8. Subgroups, Quotient Groups, and Extensions
of Groups of Multiplicative Type over a Field}
\begingroup
\renewcommand{\thefootnote}{*}
\footnotetext{Added in July 1969. This afterthought section is independent
of Sections \SGARefLink{sga3:ix:section:3}{3} through \SGARefLink{sga3:ix:section:7}{7}.}
\endgroup

% Source page: Exp9 component-local p. 30; complete-reader p. 676; printed p. 74.
\medskip
\noindent\SGARefTarget{sga3:ix:scholium:8:0}\textbf{Scholium 8.0.}
\footnote{Editor's note: this scholium has been added.} ---
\textit{Let \(k\) be a field and let \(H\) be an affine \(k\)-group
scheme. Write \(k[H]\) for its affine algebra and
\[
  X(H)=\operatorname{Hom}_{k\text{-}\mathrm{gr}}
    (H,\mathbf G_{m,k})
\]
for the group of characters of \(H\). By the independence lemma for
characters, \(X(H)\) is a linearly independent subset of \(k[H]\).
Consequently, \(H\) is diagonalizable if and only if \(X(H)\) spans
\(k[H]\) as a \(k\)-vector space.}

\medskip
\noindent\SGARefTarget{sga3:ix:proposition:8:1}\textbf{Proposition 8.1.} ---
\textit{Let \(G\) be a diagonalizable group scheme (respectively, a group
scheme of multiplicative type) over a field \(k\).\footnote{Editor's note:
the words ``over a field \(k\),'' implicit in the original, have been
added. In addition, ``algebraic group'' has been replaced by ``group
scheme,'' since \(G\) is not assumed to be of finite type.} Then, for
every subgroup scheme \(H\) of \(G\), the schemes \(H\) and \(G/H\) are
diagonalizable (respectively, of multiplicative type).}

Definition \SGARefLink{sga3:ix:definition:1:1}{1.1} immediately reduces
the assertion to the case without the parenthetical qualification. An
easy passage to the limit, using VI\textsubscript{B} \SGARefLink{prop:sga3-VIB-11.13}{11.13} and \SGARefLink{sga3:viii:theorem:3:1}{VIII 3.1},
reduces us to the case in which \(G\) is of finite type over
\(k\).\footnote{Editor's note: this ``passage to the limit'' should be
spelled out. It uses \SGARefLink{sga3:ix:scholium:8:0}{8.0} and the
equality
\[
  G/H=\varprojlim_i G_i/H_i
\]
(cf.\ VI\textsubscript{B}, proof of \SGARefLink{thm:sga3-VIB-11.17}{11.17}). To see that
\(H=\varprojlim_i(H\cap G_i)\), does one use the fact that \(H\) is
closed in \(G\)? Using this fact, one can give a direct proof.}
By VI\textsubscript{B} \SGARefLink{thm:sga3-VIB-11.16}{11.16},\footnote{Editor's note: VI\textsubscript{B}
\SGARefLink{thm:sga3-VIB-11.16}{11.16} should doubtless be rewritten in the usual, more agreeable form.
In particular, a single \(f_i\) suffices below.} one can find a finite
family of nonzero elements
\SGARefTarget{sga3:ix:eq:8:1:1}
\begin{equation}
  \tag{8.1.1}
  f_i=\sum_m a_{im}m,\qquad a_{im}\in k,
\end{equation}
of the affine ring \(k(M)\) of \(G=D_k(M)\), such that the points of
\(H\), with values in an arbitrary \(k\)-algebra \(k'\), are precisely
the points \(g\in G(k')\) for which
\SGARefTarget{sga3:ix:eq:8:1:2}
\begin{equation}
  \tag{8.1.2}
  \tau_g f_i=\lambda_i(g)f_i,
  \qquad \lambda_i(g)\in k',
\end{equation}
where \(\tau_g\) denotes translation by \(g\). But
\SGARefTarget{sga3:ix:eq:8:1:3}
\begin{equation}
  \tag{8.1.3}
  \tau_g f_i=\sum_m a_{im}\chi_m(g)m,
\end{equation}
where \(\chi_m:G\to\mathbf G_{m,k}\) is the character associated with
\(m\). Thus \SGARefLink{sga3:ix:eq:8:1:2}{(8.1.2)} is equivalent to
\SGARefTarget{sga3:ix:eq:8:1:4}
\begin{equation}
  \tag{8.1.4}
  \chi_m(g)=\chi_{m'}(g)
  \qquad\text{if }m,m'\in Z_i,
\end{equation}
where \(Z_i\) denotes the set of \(m\in M\) for which \(a_{im}\neq0\).
\SGARefLink{sga3:ix:eq:8:1:4}{The same relation} may be written
\[
  \chi_{m'-m}(g)=1
  \qquad\text{if }m,m'\in Z_i.
\]

% Source page: Exp9 component-local p. 31; complete-reader p. 677; printed p. 75.
Let \(N\) be the subgroup of \(M\) generated by all the elements
\(m'-m\), where \(i\) varies and \(m,m'\in Z_i\). It follows from the
definition of \(D_k(M)\simeq G\) that \(H\) identifies with
\(D_k(M/N)\). Hence \SGARefLink{sga3:viii:theorem:3:1}{VIII 3.1} identifies \(G/H\) with \(D_k(N)\).
\hfill\(\square\)

\medskip
\noindent\SGARefTarget{sga3:ix:proposition:8:2}\textbf{Proposition 8.2.}
\footnote{Editor's note: both the statement and its proof have been
expanded.} ---
\textit{Let \(k\) be a field; let \(H\) and \(K\) be \(k\)-group schemes
of multiplicative type and finite type; and let \(G\) be a \(k\)-group
scheme fitting into an exact sequence
\SGARefTarget{sga3:ix:eq:8:2:0}
\[
  1\longrightarrow H\longrightarrow G\longrightarrow K\longrightarrow 1
\]
(which implies that \(G\) is of finite type over \(k\)).}

\SGARefTarget{sga3:ix:proposition:8:2:item:i}\textit{(i) If \(G\) is commutative or \(K\) is connected, then \(G\) is
of multiplicative type.}

\SGARefTarget{sga3:ix:proposition:8:2:item:ii}\textit{(ii) If \(K\) and \(H\) are diagonalizable and \(K\) is a torus,
then \(G\) is diagonalizable.}

For the proof of \SGARefLink{sga3:ix:proposition:8:2:item:i}{(i)}, the reader is referred to \SGARefLink{sga3:xvii:proposition:7:1:1}{XVII 7.1.1}, of which
\SGARefLink{sga3:ix:proposition:8:2:item:i}{8.2(i)} is a special case. The case
of a field is treated in part \SGARefLink{sga3:xvii:proposition:7:1:1:item:i}{(i)} of the proof of \SGARefLink{sga3:xvii:proposition:7:1:1}{XVII 7.1.1}, without
using the results of the later exposés.

For \SGARefLink{sga3:ix:proposition:8:2:item:ii}{(ii)}, now suppose that \(K\) and \(H\) are diagonalizable:
\[
  K\simeq D_k(M)
  \qquad\text{and}\qquad
  H\simeq D_k(N),
\]
and suppose that \(G\) is of multiplicative type, as follows from \SGARefLink{sga3:ix:proposition:8:2:item:i}{(i)}
when \(K\) is a torus. By \SGARefLink{sga3:x:proposition:1:4}{X 1.4}, \(G\) is isotrivial over \(k\); that is,
there is a finite separable extension \(k'/k\) such that
\[
  G'=G\times_k k'
\]
is diagonalizable. Hence \(G'=D_{k'}(E)\) for a commutative group \(E\),
and there is an exact sequence
\SGARefTarget{sga3:ix:eq:8:2:1}
\begin{equation}
  \tag{1}
  0\longrightarrow D_{k'}(N)\longrightarrow D_{k'}(E)
    \longrightarrow D_{k'}(M)\longrightarrow 0.
\end{equation}
By \SGARefLink{sga3:viii:theorem:3:1}{VIII 3.1} and \SGARefLink{sga3:ix:theorem:3:2}{3.2}, \(M\) is therefore a subgroup of \(E\), and one has
an exact sequence
\SGARefTarget{sga3:ix:eq:8:2:2}
\begin{equation}
  \tag{2}
  0\longrightarrow M\longrightarrow E\longrightarrow N\longrightarrow 0.
\end{equation}

For a given extension \(E_0\) of \(N\) by \(M\), consider the
diagonalizable \(k\)-group \(G_0=D_k(E_0)\). The \(k\)-group functor
\(A\) of automorphisms of the extension
\SGARefTarget{sga3:ix:eq:8:2:3}
\begin{equation}
  \tag{3}
  1\longrightarrow H\longrightarrow G_0\longrightarrow K\longrightarrow 1,
\end{equation}
that is, the group subfunctor of
\(\operatorname{Aut}_{k\text{-}\mathrm{gr}}(G_0)\) whose points on a
\(k\)-prescheme \(T\) are the
\[
  \phi\in\operatorname{Aut}_{T\text{-}\mathrm{gr}}(G_T)
\]
inducing the identity on \(H_T\) and \(K_T\), identifies with
\(\operatorname{Hom}_{k\text{-}\mathrm{gr}}(K,H)\). By \SGARefLink{sga3:viii:corollary:1:5}{VIII 1.5}, this is
the constant \(k\)-group with value
\[
  L=\operatorname{Hom}_{\mathrm{gr}}(N,M).
\]

It follows that the extensions \(G\) of \(K\) by \(H\) which become
isomorphic to \SGARefLink{sga3:ix:eq:8:2:3}{(3)} over a separable closure
\(k_s\) of \(k\) are classified in the same way as the \(k\)-torsors,
for the étale topology, under the constant group \(L_k\). Put
\(\Gamma=\operatorname{Gal}(k_s/k)\). Since \(L\) is a trivial
\(\Gamma\)-module, these torsors are classified by
\SGARefTarget{sga3:ix:eq:8:2:4}
\[
  H^1_{\mathrm{\acute et}}(k,L_k)
  =H^1(\Gamma,L)
  =\operatorname{Hom}_{\mathrm{gr.\,top.}}(\Gamma,L).
\]

% Source page: Exp9 component-local p. 32; complete-reader p. 678; printed p. 76.
If \(K\) is moreover a torus, then \(M\) is torsion-free, and hence so
is \(L\). Therefore
\[
  \operatorname{Hom}_{\mathrm{gr.\,top.}}(\Gamma,L)=0,
\]
and every extension of \(K\) by \(H\) is already diagonalizable.

\medskip
\noindent\SGARefTarget{sga3:ix:remark:8:3}\textbf{Remark 8.3.} ---
\textit{As already noted in the proof of
\SGARefLink{sga3:ix:proposition:8:2}{8.2}, \SGARefLink{sga3:ix:proposition:8:2:item:i}{the first assertion} is stated
and proved over an arbitrary base in \SGARefLink{sga3:xvii:proposition:7:1:1}{XVII 7.1.1}. \SGARefLink{sga3:ix:proposition:8:2:item:ii}{The second assertion}
generalizes, with essentially the same proof, to the case of an integral
normal base scheme---or, more generally, a geometrically unibranch
one---by using \SGARefLink{sga3:x:lemma:5-13}{X 5.13} below.}

\textit{Finally, \SGARefLink{sga3:ix:proposition:8:1}{8.1} may also be
regarded as a corollary of the distinctly less trivial result \SGARefLink{sga3:ix:theorem:6:8}{6.8}, which
is likewise valid over an arbitrary base scheme.}\footnote{Editor's note:
the proof of \SGARefLink{sga3:ix:theorem:6:8}{6.8} nevertheless uses
\SGARefLink{sga3:ix:proposition:8:1}{8.1}.}

\SGARefTarget{sga3:ix:bibliography}
\section*{Bibliography}
\addcontentsline{toc}{section}{Bibliography}
\footnote{Editor's note: additional references cited in this exposé.}

\noindent \SGARefTarget{sga3:ix:bib:bens}[BEns] N. Bourbaki, \textit{Théorie des ensembles},
Chapters I--IV, Hermann, 1970.

% END INCLUDED SGA3_English_current_buildable_source_20260823/source/inputs/IX/tex/components/05_expose_IX_section_8_and_bibliography.tex



% END INCLUDED SGA3_English_current_buildable_source_20260823/source/inputs/IX/body.tex


\SGAStartExpose{X}
% BEGIN INCLUDED SGA3_English_current_buildable_source_20260823/source/inputs/X/body.tex
% Mechanically extracted after \mainmatter from the frozen standalone master.

\markboth{Exposé X. Groups of Multiplicative Type}{Exposé X. Groups of Multiplicative Type}
% BEGIN INCLUDED SGA3_English_current_buildable_source_20260823/source/inputs/X/tex/components/01_expose_X_sections1_3.tex
% SGA 3, Exposé X, Loop-1 English reconstruction.
% Scope: title and Sections 1--3, through Corollary 3.9.
% Authority: Polo--Gille Expo10-8nov09.pdf, local pages 1--11
% (combined-reader pages 679--689; printed pages 63--73, with marginal
% pagination 77--92). OCR and the Reinhold English Markdown were used
% only as comparison witnesses.

\setcounter{footnote}{0}
\SGARefTarget{sga3:x:chapter:x}\chapter*{Exposé X\\
Characterization and Classification\\
of Groups of Multiplicative Type}
\addcontentsline{toc}{chapter}{Exposé X. Characterization and
Classification of Groups of Multiplicative Type}
\markboth{EXPOSÉ X. GROUPS OF MULTIPLICATIVE TYPE}
  {EXPOSÉ X. GROUPS OF MULTIPLICATIVE TYPE}

\begin{center}
  \textit{by A. Grothendieck}
\end{center}
\footnotetext{xy version of November 6, 2009: Addenda placed in
\SGARefLink{sga3:x:section:9}{Section~9}; reviewed through \SGARefLink{sga3:x:theorem:8-8}{8.8}.}

\SGARefTarget{sga3:x:section:1}\section*{1. Classification of Isotrivial
Groups. Case of a Base Field}
\addcontentsline{toc}{section}{1. Classification of Isotrivial Groups.
Case of a Base Field}

Recall (\SGARefLink{sga3:ix:definition:1:1}{IX~1.1}) that a group of multiplicative type \(H\) over the
prescheme \(S\) is said to be \emph{isotrivial} if there exists a finite
étale surjective morphism \(S'\to S\) such that
\(H'=H\times_S S'\) is diagonalizable. When \(S\) is connected,
\(S'\) decomposes as a finite disjoint sum of connected components
\(S'_i\), and hence, after replacing \(S'\) by one of the \(S'_i\), we
may suppose that \(S'\) is connected. Finally, \(S'\) can be dominated
by a connected finite étale \(S'_1\) that is Galois, that is, a
principal homogeneous bundle over \(S\) with group \(\Gamma_S\), where
\(\Gamma\) is an ordinary finite group (cf.\ SGA~1, V, nos.~7 and~3
when \(S\) is locally noetherian, and EGA~IV\textsubscript{4}, 18.2.9
in general).\footnote{Editor's note: More precisely, this follows from
the ``axiomatic conditions of a Galois theory,'' cf.\ SGA~1, V,
no.~4~g); when \(S\) is locally noetherian, the axioms are verified in
\emph{loc.\ cit.}, nos.~7 and~3, especially \SGARefLink{sga3:x:theorem:3:7}{3.7}, which rests on
SGA~1, I, 10.9. The latter result is proved without noetherian
hypotheses in EGA~IV\textsubscript{4}, 18.2.9.}
We suppose \(S'\) chosen in this way and propose to determine the
groups of multiplicative type \(H\) over \(S\) that are
``trivialized'' by \(S'\), that is, those for which
\(H'=H\times_S S'\) is diagonalizable.\footnote{Editor's note: In what
follows, an \(S\)-group of multiplicative type will be called
``split'' if it is ``trivial'' in the sense of \SGARefLink{sga3:ix:comments:1:2}{IX~1.2}, that is, if it
is a diagonalizable \(S\)-group.}
By descent theory (cf.\ SGA~1, VIII, \SGARefLink{sga3:viii:corollary:5:4}{5.4}), the category of such \(H\)
is equivalent to the category of diagonalizable groups \(H'\) over
\(S'\), equipped with actions of \(\Gamma\) on \(H'\) compatible with
the actions of \(\Gamma\) on \(S'\). (N.B. Since the groups under
consideration are affine over the base, every descent datum here is
effective; cf.\ SGA~1, VIII, \SGARefLink{sga3:viii:proposition:2:1}{2.1}.) Since \(S'\) is connected, the
contravariant functor
\[
  M\longmapsto D_{S'}(M)
\]
is an anti-equivalence from the category of ordinary commutative
groups to the category of diagonalizable groups over \(S'\)
(cf.\ \SGARefLink{sga3:viii:corollary:1:6}{VIII~1.6}); a quasi-inverse is\footnote{Editor's note: \(S\) has
been corrected to \(S'\).}
\[
  H\longmapsto
  \operatorname{Hom}_{S'\text{-}\mathrm{gr}}(H,\mathbb G_{m,S'}).
\]

\medskip
\noindent\SGARefTarget{sga3:x:proposition:1:1}%
\textbf{Proposition 1.1.} ---
\textit{Let \(S\) be a connected prescheme and let \(S'\) be a
connected principal cover of \(S\) with finite group \(\Gamma\). Then
the category of groups of multiplicative type over \(S\) split by
\(S'\)\footnote{Editor's note: Here and below, ``splittés'' has been
replaced in the French text by ``déployés,'' and ``splitte'' by
``déploie,'' etc. Both are rendered here by ``split.''} is
anti-equivalent to the category of \(\Gamma\)-modules, that is, of
ordinary commutative groups \(M\) equipped with a homomorphism
\[
  \Gamma\longrightarrow\operatorname{Aut}_{\mathrm{gr}}(M).
\]}

The following standard consequence results.

\medskip
\noindent\SGARefTarget{sga3:x:corollary:1:2}%
\textbf{Corollary 1.2.} ---
\textit{Let \(S\) be a connected prescheme, and let
\(\xi:\operatorname{Spec}(\Omega)\to S\) be a ``geometric point'' of
\(S\), that is, a morphism into \(S\) from the spectrum of an
algebraically closed field \(\Omega\). Consider the fundamental group
of \(S\) at \(\xi\) (cf.\ SGA~1, V, no.~7):
\[
  \pi=\pi_1(S,\xi).
\]
Then the category of isotrivial groups of multiplicative type \(H\)
over \(S\) is anti-equivalent to the category of ``Galois modules''
under \(\pi\), that is, of \(\pi\)-modules \(M\) such that the
stabilizer in \(\pi\) of every element of \(M\) is an open subgroup.}

\textit{Under this correspondence, the group associated with an
isotrivial group of multiplicative type \(H\) is
\[
  M=\operatorname{Hom}_{\Omega\text{-}\mathrm{gr}}
       (H_\xi,\mathbb G_{m,\Omega}),
\]
where \(H_\xi\) is the fiber of \(H\) at \(\xi\); this group is
naturally endowed with an action of \(\pi_1(S,\xi)\).}

\medskip
\noindent\SGARefTarget{sga3:x:remark:1:3}%
\textbf{Remark 1.3.} ---
We shall see below (cf.\ \SGARefLink{sga3:x:theorem:5-16}{5.16}) that if \(S\) is normal, or more
generally geometrically unibranch, then every group of multiplicative
type and of finite type over \(S\) is necessarily isotrivial.
Consequently, the classification principle of \SGARefLink{sga3:x:corollary:1:2}{1.2} applies to groups of
multiplicative type and of finite type over \(S\); they correspond to
Galois \(\pi\)-modules that are of finite type over \(\mathbb Z\). For
the moment, let us confine ourselves to the following result.

\medskip
\noindent\SGARefTarget{sga3:x:proposition:1:4}%
\textbf{Proposition 1.4.} ---
\textit{Let \(k\) be a field and \(H\) a group of multiplicative type
and of finite type over \(k\). Then \(H\) is isotrivial; that is, there
exists a finite separable extension \(k'\) of \(k\) that splits \(H\).}

\textit{Consequently, by \SGARefLink{sga3:x:corollary:1:2}{1.2}, if \(\pi\) is the topological Galois
group of an algebraic closure \(\Omega\) of \(k\), the category of
groups \(H\) of multiplicative type and of finite type over \(k\) is
anti-equivalent to the category of Galois modules \(M\) under \(\pi\)
that are of finite type as \(\mathbb Z\)-modules.}

Since \(H\) is of finite type over \(k\), the ``principle of finite
extension'' (cf.\ EGA~IV\textsubscript{3}, 9.1.4) first gives a finite
extension \(k'\) of \(k\) that splits \(H\). Recall the argument. By
hypothesis, there exist a diagonalizable group \(G\) of finite type
over \(k\), a faithfully flat \(S'\) over
\(S=\operatorname{Spec}(k)\), and an isomorphism of \(S'\)-groups
\[
  H_{S'}\simeq G_{S'}.
\]
After replacing \(S'\) by the residue field of one of its points, we
may suppose that \(S'\) is the spectrum of an extension \(K\) of
\(k\). The field \(K\) is the filtered direct limit of its finitely
generated \(k\)-subalgebras \(A_i\); it follows readily (cf.\
EGA~IV\textsubscript{3}, 8.8.2.4) that the given isomorphism \(u\)
comes, for sufficiently large \(i\), from an \(A_i\)-isomorphism
\[
  u_i:H_{A_i}\simeq G_{A_i}.
\]
By the Nullstellensatz, \(A_i\) has a quotient ring \(k'\) that is a
finite extension of \(k\). Thus \(k'\) splits \(H\).

Now \(k'\) is a radicial extension of a separable extension \(k'_s\)
of \(k\). By \SGARefLink{sga3:ix:corollary:5:4}{IX~5.4}, the isomorphism
\[
  u':H_{k'}\simeq G_{k'}
\]
comes from an isomorphism \(H_{k'_s}\simeq G_{k'_s}\). Thus \(k'_s\)
splits \(H\), proving \SGARefLink{sga3:x:proposition:1:4}{1.4}.

\medskip
\noindent\SGARefTarget{sga3:x:remark:1:5}%
\textbf{Remark 1.5.} ---
Statement \SGARefLink{sga3:x:corollary:1:2}{1.2} gives, in particular, a characterization of the
isotrivial tori over \(S\) of relative dimension \(n\): setting
\(\pi=\pi_1(S,\xi)\), they correspond to classes, up to
``equivalence,'' of representations
\[
  \pi\longrightarrow\operatorname{GL}(n,\mathbb Z)
\]
whose kernel is an open subgroup of \(\pi\).

\medskip
\noindent\SGARefTarget{sga3:x:paragraph:1:6}\textbf{1.6.}
Even when \(S\) is an algebraic curve, there may exist tori over \(S\)
(of relative dimension \(2\)) that are not locally isotrivial (and
\emph{a fortiori} not isotrivial); there may also exist locally
trivial tori that are not isotrivial. (Observe, however, that such
phenomena can occur only when \(S\) is not normal, as already noted in
\SGARefLink{sga3:x:remark:1:3}{1.3}.) For example, let \(S\) be an irreducible algebraic curve (over
an algebraically closed field, to fix ideas) with an ordinary double
point \(a\), let \(S'\) be its normalization, and let \(b\) and \(c\)
be the two points of \(S'\) above \(a\). A connected principal
homogeneous bundle \(P\) over \(S\), with structural group
\(\mathbb Z\), is obtained by attaching infinitely many copies of
\(S'\) according to the diagram
\begin{center}
  % BEGIN INCLUDED SGA3_English_current_buildable_source_20260823/source/inputs/X/native_diagrams/exp10/001_normalization_gluing_chain.tex
\begin{tikzpicture}[
  x=1cm,
  y=1cm,
  line cap=round,
  line join=round,
  every node/.style={inner sep=1pt}
]
  \coordinate (b1) at (2.8,0.55);
  \coordinate (c1) at (4.4,-0.55);
  \coordinate (b2) at (6.0,0.55);
  \coordinate (c2) at (7.6,-0.55);
  \coordinate (b3) at (9.2,0.55);
  \coordinate (c3) at (10.8,-0.55);
  \coordinate (b4) at (12.4,0.55);

  \draw[line width=.45pt]
    (1.75,1.27) -- (b1) -- (c1) -- (b2) -- (c2) --
    (b3) -- (c3) -- (b4) -- (13.45,-0.17);
  \draw[line width=.45pt] (1.75,-0.72) -- (b1) -- (3.52,1.27);
  \draw[line width=.45pt] (3.68,-1.18) -- (c1) -- (5.12,-1.18);
  \draw[line width=.45pt] (5.28,1.27) -- (b2) -- (6.72,1.27);
  \draw[line width=.45pt] (6.88,-1.18) -- (c2) -- (8.32,-1.18);
  \draw[line width=.45pt] (8.48,1.27) -- (b3) -- (9.92,1.27);
  \draw[line width=.45pt] (10.08,-1.18) -- (c3) -- (11.52,-1.18);
  \draw[line width=.45pt] (11.68,1.27) -- (b4) -- (13.12,1.27);

  \node[below=1pt] at (b1) {$b$};
  \node[below=1pt] at (b2) {$b$};
  \node[below=1pt] at (b3) {$b$};
  \node[below=1pt] at (b4) {$b$};
  \node[below=1pt] at (c1) {$c$};
  \node[below=1pt] at (c2) {$c$};
  \node[below=1pt] at (c3) {$c$};

  \node at (0.65,0) {$\cdots\!\cdots$};
  \node at (14.55,0) {$\cdots\!\cdots$};
\end{tikzpicture}

% END INCLUDED SGA3_English_current_buildable_source_20260823/source/inputs/X/native_diagrams/exp10/001_normalization_gluing_chain.tex

\end{center}
(N.B. This is a principal bundle for the étale topology.) There is a
homomorphism
\[
  \varphi:\mathbb Z\longrightarrow\operatorname{GL}(2,\mathbb Z),
  \qquad
  \varphi(n)=
  \begin{pmatrix}
    1 & n\\
    0 & 1
  \end{pmatrix},
\]
which makes it possible to construct a torus \(T\) over \(S\), of
relative dimension \(2\), from the trivial torus \(\mathbb G_m^2\)
over \(P\) and the descent datum on it deduced from \(\varphi\).
(N.B. The projection \(P\to S\) is a covering for the étale topology
and \emph{a fortiori} for the canonical topology of
\((\mathbf{Sch})\); the descent datum in question is necessarily
effective because \(\mathbb G_{m,P}^2\) is affine over \(P\).) It is
not difficult to prove that \(T\) is not isotrivial in a neighborhood
of \(a\)\footnote{Editor's note: See \SGARefLink{sga3:x:example:7-3}{7.3} below.} (although it is
trivial on \(S-\{a\}\)).

A variant of this construction is obtained by taking for \(S\) a
curve with two irreducible components \(S_1\) and \(S_2\) meeting in
two points \(a'\) and \(a''\). This makes it possible to construct a
connected, locally trivial principal homogeneous bundle \(P\) over
\(S\) with group \(\mathbb Z_S\), and hence an associated torus \(T\)
that is locally trivial but not isotrivial.

\SGARefTarget{sga3:x:section:2}\section*{2. Infinitesimal Variations of
Structure}
\addcontentsline{toc}{section}{2. Infinitesimal Variations of Structure}

\footnote{Editor's note: This ``recall'' has been added for later
references.}
Let us begin by recalling the following result (cf.\ SGA~1, I, 8.3
when \(S\) is locally noetherian, and EGA~IV\textsubscript{4}, 18.1.2
in general).

\medskip
\noindent\SGARefTarget{sga3:x:recall:2:0}\textbf{Recall 2.0.} ---
\textit{Let \(S\) be a prescheme and \(S_0\) a subprescheme with the
same underlying set. Then the functor
\[
  X\longmapsto X_0=X\times_S S_0
\]
is an equivalence between the category of étale preschemes over \(S\)
and the analogous category over \(S_0\).}

\medskip
\noindent\SGARefTarget{sga3:x:proposition:2:1}%
\textbf{Proposition 2.1.} ---
\textit{Let \(S\) be a prescheme and \(S_0\) a subprescheme with the
same underlying set (that is, defined by a nilideal \(\mathscr I\)).
Then the functor
\[
  H\longmapsto H_0=H\times_S S_0
\]
from the category of group preschemes of multiplicative type over
\(S\) to the analogous category over \(S_0\) is fully faithful.}

\textit{Moreover, it induces an equivalence between the category of
quasi-isotrivial groups of multiplicative type over \(S\) and the
analogous category over \(S_0\).}

We first prove full faithfulness, that is, if \(H\) and \(G\) over
\(S\) are of multiplicative type, then
\[
  \operatorname{Hom}_{S\text{-}\mathrm{gr}}(H,G)
  \longrightarrow
  \operatorname{Hom}_{S_0\text{-}\mathrm{gr}}(H_0,G_0)
\]
is bijective. The question is local on \(S\), so we may suppose that
\(S\) is affine. There then exists a faithfully flat quasi-compact
morphism \(S'\to S\) that splits \(H\) and \(G\). Let
\(S''=S'\times_S S'\); write \(H',G'\), respectively \(H'',G''\), for
the groups obtained from \(H,G\) by the base changes \(S'\to S\),
respectively \(S''\to S\). Define \(S'_0\) and \(S''_0\) similarly;
the latter is also isomorphic to \(S'_0\times_{S_0}S'_0\). We obtain a
commutative diagram with exact rows:
\begin{center}
  % BEGIN INCLUDED SGA3_English_current_buildable_source_20260823/source/inputs/X/native_diagrams/exp10/002_descent_hom_exact_rows.tex
{\small
\begin{tikzcd}[column sep=2.1em, row sep=3.4em]
  \operatorname{Hom}_{S\text{-}\mathrm{gr.}}(H,G)
    \arrow[r]
    \arrow[d,"u"]
  &
  \operatorname{Hom}_{S'\text{-}\mathrm{gr.}}(H',G')
    \arrow[r,shift left=.65ex]
    \arrow[r,shift right=.65ex]
    \arrow[d,"u'"]
  &
  \operatorname{Hom}_{S''\text{-}\mathrm{gr.}}(H'',G'')
    \arrow[d,"u''"]
  \\
  \operatorname{Hom}_{S_0\text{-}\mathrm{gr.}}(H_0,G_0)
    \arrow[r]
  &
  \operatorname{Hom}_{S'_0\text{-}\mathrm{gr.}}(H'_0,G'_0)
    \arrow[r,shift left=.65ex]
    \arrow[r,shift right=.65ex]
  &
  \operatorname{Hom}_{S''_0\text{-}\mathrm{gr.}}(H''_0,G''_0),
\end{tikzcd}
}

% END INCLUDED SGA3_English_current_buildable_source_20260823/source/inputs/X/native_diagrams/exp10/002_descent_hom_exact_rows.tex

\end{center}
Thus, to prove that \(u\) is bijective, it is enough to prove that
\(u'\) and \(u''\) are bijective. This reduces us to the case where
\(H\) and \(G\) are diagonalizable, hence of the form \(D_S(M)\) and
\(D_S(N)\), where \(M\) and \(N\) are ordinary commutative groups.
Likewise,
\[
  H_0=D_{S_0}(M),\qquad G_0=D_{S_0}(N).
\]
We then have a commutative diagram\footnote{Editor's note: The
original has been corrected by interchanging \(M\) and \(N\) in the
right-hand terms.}
\begin{center}
  % BEGIN INCLUDED SGA3_English_current_buildable_source_20260823/source/inputs/X/native_diagrams/exp10/003_diagonalizable_hom_square.tex
{\small
\begin{tikzcd}[column sep=3.2em, row sep=3.5em]
  \operatorname{Hom}_{S\text{-}\mathrm{gr.}}(N_S,M_S)
    \arrow[r,"\sim"]
    \arrow[d]
  &
  \operatorname{Hom}_{S\text{-}\mathrm{gr.}}\!\bigl(D_S(M),D_S(N)\bigr)
    \arrow[d]
  \\
  \operatorname{Hom}_{S_0\text{-}\mathrm{gr.}}(N_{S_0},M_{S_0})
    \arrow[r,"\sim"]
  &
  \operatorname{Hom}_{S_0\text{-}\mathrm{gr.}}\!\bigl(D_{S_0}(M),D_{S_0}(N)\bigr),
\end{tikzcd}
}

% END INCLUDED SGA3_English_current_buildable_source_20260823/source/inputs/X/native_diagrams/exp10/003_diagonalizable_hom_square.tex

\end{center}
whose horizontal arrows are isomorphisms by \SGARefLink{sga3:viii:corollary:1:4}{VIII~1.4}. We are therefore
reduced to proving that the homomorphism
\begin{equation}
  \operatorname{Hom}_{S\text{-}\mathrm{gr}}(N_S,M_S)
  \longrightarrow
  \operatorname{Hom}_{S_0\text{-}\mathrm{gr}}(N_{S_0},M_{S_0})
  \tag{x}\label{sga3:x:proposition:2-1:equation:x}
\end{equation}
is bijective; in other words, the functor
\(M_S\mapsto M_{S_0}\), from constant commutative group preschemes
over \(S\) to constant commutative group preschemes over \(S_0\), is
fully faithful. Now \SGARefLink{sga3:x:proposition:2-1:equation:x}{(x)} is also identified with the natural map
\[
  \operatorname{Hom}_{\mathrm{gr}}(N,\Gamma(M_S))
  \longrightarrow
  \operatorname{Hom}_{\mathrm{gr}}(N,\Gamma(M_{S_0}))
\]
deduced from \(\Gamma(M_S)\to\Gamma(M_{S_0})\). The latter map is
plainly bijective, since \(\Gamma(M_S)\), the set of locally constant
maps from \(S\) to \(M\), depends only on the underlying topological
space of \(S\). This proves the desired assertion.

To prove the second assertion of \SGARefLink{sga3:x:proposition:2:1}{2.1}, it remains to show that every
quasi-isotrivial group \(H_0\) of multiplicative type over \(S_0\)
comes from a quasi-isotrivial group \(H\) of multiplicative type over
\(S\). Let \(S'_0\to S_0\) be an étale surjective morphism that splits
\(H_0\).

By \SGARefLink{sga3:x:recall:2:0}{2.0}, there exist an étale morphism \(S'\to S\) and an
\(S_0\)-isomorphism
\[
  S'\times_S S_0\simeq S'_0.
\]
Thus we may suppose that \(S'_0\) is obtained from \(S'\) by
reduction. Since \(H'_0\) is diagonalizable, it is immediately seen
to be isomorphic to the group obtained by the base change
\(S'_0\to S'\) from a diagonalizable group \(H'\) over \(S'\). (N.B.
If \(H'_0=D_{S'_0}(M)\), take \(H'=D_{S'}(M)\).) As usual, set
\[
  S''=S'\times_S S',
  \qquad
  S'''=S'\times_S S'\times_S S',
\]
and similarly define \(S''_0\) and \(S'''_0\) by base change from
\(S\) to \(S_0\). They are also isomorphic to the fibered square and
fibered cube, respectively, of \(S'_0\) over \(S_0\). Applying the
full faithfulness already proved to \((S'',S''_0)\) and
\((S''',S'''_0)\), we see that the natural descent datum on \(H'_0\)
relative to \(S'_0\to S_0\) (cf.\ \SGARefLink{sga3:iv:definition:2:1}{IV~2.1}) comes from a uniquely
determined descent datum on \(H'\) relative to \(S'\to S\). This
descent datum is effective because \(H'\) is affine over \(S'\)
(SGA~1, VIII, \SGARefLink{sga3:viii:proposition:2:1}{2.1}). Hence there exists an \(S\)-group \(H\) such that
\[
  H\times_S S'=H'=D_{S'}(M),
\]
and \(H\) is therefore a quasi-isotrivial group of multiplicative
type.

Using the full-faithfulness result for \((S',S'_0)\), one now checks
readily that the given isomorphism between \(H'_0\) and
\(H'\times_{S'}S'_0\) comes from an isomorphism between \(H_0\) and
\(H\times_S S_0\). (For a more formal presentation of results of this
kind, see Giraud's article in preparation on descent theory.%
\begingroup
\renewcommand{\thefootnote}{\(\ast\)}%
\footnote{Cf.\ J.~Giraud, ``Méthode de la descente,'' \emph{Bull.\
Soc.\ Math.\ France}, Mémoire~2 (1964).}%
\endgroup
\addtocounter{footnote}{-1})

\medskip
\noindent\SGARefTarget{sga3:x:corollary:2:2}%
\textbf{Corollary 2.2.} ---
\textit{Let \(H\) be an \(S\)-group of multiplicative type, and set
\(H_0=H\times_S S_0\). Then \(H\) is quasi-isotrivial (respectively,
locally isotrivial, isotrivial, locally trivial, or trivial) if and
only if \(H_0\) has the corresponding property.}

The ``only if'' direction is immediate. The ``if'' direction has
already been proved in the quasi-isotrivial case: by full
faithfulness, it is enough to know that every quasi-isotrivial group
over \(S_0\) lifts to a quasi-isotrivial group over \(S\). The same
argument applies to ``trivial.'' For ``isotrivial,'' repeat the
argument establishing the second assertion of \SGARefLink{sga3:x:proposition:2:1}{2.1}, but take
\(S'_0\to S_0\) to be finite étale and surjective. The locally
isotrivial and locally trivial cases follow at once from the
isotrivial and trivial cases.

When \(\mathscr I\) is nilpotent, \SGARefLink{sga3:x:corollary:2:2}{2.2} can be generalized slightly
without assuming \emph{a priori} that \(H\) is of multiplicative
type.\footnote{Editor's note: Recall that \SGARefLink{sga3:x:corollary:2:3}{2.3} is used in proving
Theorem~\SGARefLink{sga3:ix:theorem:3:6}{IX~3.6}~bis.}

\medskip
\noindent\SGARefTarget{sga3:x:corollary:2:3}%
\textbf{Corollary 2.3.} ---
\textit{Suppose that the ideal \(\mathscr I\) defining \(S_0\) is
locally nilpotent. Let \(H\) be a group prescheme over \(S\), flat
over \(S\), and set \(H_0=H\times_S S_0\). Then \(H\) is of
quasi-isotrivial multiplicative type if and only if \(H_0\) is so.}

Suppose that \(H_0\) is of quasi-isotrivial multiplicative type. We
prove that the same is true of \(H\). The question is local for the
étale topology, and by \SGARefLink{sga3:x:recall:2:0}{2.0} the category of étale preschemes over \(S\)
is equivalent to the category of étale preschemes over \(S_0\) by the
functor \(S'\mapsto S'\times_S S_0\). We are therefore immediately
reduced to the case where \(H_0\) is diagonalizable, hence isomorphic
to \(D_{S_0}(M)\). Let \(G=D_S(M)\). We then have an isomorphism
\[
  u_0:H_0\xrightarrow{\sim}G_0.
\]
I claim that \(u_0\) comes from a unique homomorphism \(u:H\to G\),
which is necessarily an isomorphism because \(u_0\) is one, \(H\) and
\(G\) are flat over \(S\), and \(\mathscr I\) is locally
nilpotent.\footnote{Editor's note: Since \(u_0\) is an isomorphism, it
is enough to see that for every \(h\in H\), the morphism
\(\varphi:\mathscr O_{G,u(h)}\to\mathscr O_{H,h}\) is bijective. It is
bijective after reduction modulo
\(\mathscr I=\mathscr I_s\), where \(s\) is the image of \(h\) in
\(S\); hence its cokernel \(C\) satisfies \(C=\mathscr IC\), so
\(C=0\) because \(\mathscr I\) is nilpotent. Since
\(\mathscr O_{H,h}\) is flat
over \(\mathscr O_{S,s}\), the kernel \(K\) of \(\varphi\) likewise
satisfies \(K=\mathscr IK\), and hence \(K=0\).}
This will prove \SGARefLink{sga3:x:corollary:2:3}{2.3}. We have (cf.\ VIII~1~(xxx))
\[
  \operatorname{Hom}_{S\text{-}\mathrm{gr}}(H,G)
  \simeq
  \operatorname{Hom}_{S\text{-}\mathrm{gr}}
  \bigl(M_S,
    \underline{\operatorname{Hom}}_{S\text{-}\mathrm{gr}}
      (H,\mathbb G_{m,S})\bigr),
\]
and the right-hand side is also identified with
\[
  \operatorname{Hom}_{\mathrm{gr}}
  \bigl(M,
    \operatorname{Hom}_{S\text{-}\mathrm{gr}}
      (H,\mathbb G_{m,S})\bigr).
\]
Thus the homomorphism
\begin{equation}
  \operatorname{Hom}_{S\text{-}\mathrm{gr}}(H,G)
  \longrightarrow
  \operatorname{Hom}_{S_0\text{-}\mathrm{gr}}(H_0,G_0)
  \tag{x}\label{sga3:x:corollary:2-3:equation:x}
\end{equation}
is isomorphic to the homomorphism
\[
  \operatorname{Hom}_{\mathrm{gr}}
  \bigl(M,\operatorname{Hom}_{S\text{-}\mathrm{gr}}
    (H,\mathbb G_{m,S})\bigr)
  \longrightarrow
  \operatorname{Hom}_{\mathrm{gr}}
  \bigl(M,\operatorname{Hom}_{S_0\text{-}\mathrm{gr}}
    (H_0,\mathbb G_{m,S_0})\bigr)
\]
deduced from the restriction homomorphism
\begin{equation}
  \operatorname{Hom}_{S\text{-}\mathrm{gr}}(H,\mathbb G_{m,S})
  \longrightarrow
  \operatorname{Hom}_{S_0\text{-}\mathrm{gr}}
    (H_0,\mathbb G_{m,S_0}).
  \tag{xx}\label{sga3:x:corollary:2-3:equation:xx}
\end{equation}
It is therefore enough to prove that \SGARefLink{sga3:x:corollary:2-3:equation:xx}{(xx)} is bijective. The question
is local on \(S\), so we may suppose that \(S\) is affine. Since
\(\mathbb G_{m,S}\) is commutative and smooth over \(S\), \SGARefLink{sga3:ix:theorem:3:6}{IX~3.6}
applies and completes the proof.

\medskip
\noindent\SGARefTarget{sga3:x:corollary:2:4}%
\textbf{Corollary 2.4.} ---
\textit{Let \(A\) be an artinian local ring with residue field \(k\),
\(S=\operatorname{Spec}(A)\), and \(S_0=\operatorname{Spec}(k)\).}
\begin{enumerate}
\item[\SGAInlineRefTarget{sga3:x:corollary:2:4:item:i}(i)]\textit{\footnote{Editor's note: The statement has been
rewritten, retaining only the nontrivial implication in order to
highlight it.} Let \(H\) be an \(S\)-group prescheme, flat and locally
of finite type, such that \(H_0=H\times_S S_0\) is of multiplicative
type. Then \(H\) is of multiplicative type, of finite type, and
isotrivial. In particular, every \(S\)-group prescheme \(H\) of
multiplicative type and of finite type is isotrivial.}
\item[\SGAInlineRefTarget{sga3:x:corollary:2:4:item:ii}(ii)]\textit{The functor \(H\mapsto H_0\) is an equivalence
between the category of groups of multiplicative type and of finite
type over \(A\) and the analogous category over \(k\).}
\end{enumerate}

Indeed, let \(H\) be as in \SGARefLink{sga3:x:corollary:2:4:item:i}{(i)}. Then \(H_0\) is of multiplicative type
and locally of finite type, hence of finite type, and therefore
isotrivial by \SGARefLink{sga3:x:proposition:1:4}{1.4}. By \SGARefLink{sga3:x:corollary:2:3}{2.3} and \SGARefLink{sga3:x:corollary:2:2}{2.2}, \(H\) is consequently of
multiplicative type (and hence of finite type) and isotrivial.
Assertion \SGARefLink{sga3:x:corollary:2:4:item:ii}{(ii)} now follows from \SGARefLink{sga3:x:proposition:2:1}{2.1}.

\medskip
\noindent\SGARefTarget{sga3:x:corollary:2:5}%
\textbf{Corollary 2.5.} ---
\textit{Let \(S'\to S\) be a faithfully flat radicial morphism.}
\begin{enumerate}
\item[(i)]\textit{The functor
\[
  H\longmapsto H'=H\times_S S'
\]
from the category of groups of multiplicative type over \(S\) to the
analogous category over \(S'\) is fully faithful. Moreover, it induces
an equivalence between the subcategories of quasi-isotrivial groups
of multiplicative type.}
\item[(ii)]\textit{If \(H\) is of multiplicative type, then \(H\) is
quasi-isotrivial (respectively, locally isotrivial, isotrivial,
locally trivial, or trivial) if and only if \(H'\) has the
corresponding property.}
\end{enumerate}

Indeed, let
\[
  S''=S'\times_S S',
  \qquad
  S'''=S'\times_S S'\times_S S'.
\]
The assumption that \(S'\to S\) is radicial implies that the diagonal
immersions \(S'\to S''\) and \(S'\to S'''\) are surjective, so base
change by either immersion is governed by \SGARefLink{sga3:x:proposition:2:1}{2.1} and \SGARefLink{sga3:x:corollary:2:2}{2.2}. Taking into
account that \(S'\to S\) is a morphism of effective descent for the
fibered category of groups of multiplicative type over preschemes
(because it is faithfully flat and quasi-compact\footnote{Editor's
note: Spell out this point \ldots}), the assertion follows formally
from \SGARefLink{sga3:x:proposition:2:1}{2.1} and \SGARefLink{sga3:x:corollary:2:2}{2.2} (cf.\ the already cited article of Giraud for a
formal argument).

\SGARefTarget{sga3:x:section:3}\section*{3. Finite Variations of
Structure: Complete Base Ring}
\addcontentsline{toc}{section}{3. Finite Variations of Structure:
Complete Base Ring}

\medskip
\noindent\SGARefTarget{sga3:x:lemma:3:1}\textbf{Lemma 3.1.} ---
\textit{Let \(A\) be a noetherian ring equipped with an ideal \(I\)
such that \(A\) is separated and complete for the \(I\)-adic topology;
set \(S=\operatorname{Spec}(A)\) and
\(S_0=\operatorname{Spec}(A/I)\). Let \(G\) and \(H\) be two
\(S\)-group preschemes, with \(G\) of multiplicative type and
isotrivial, and with \(H\) affine over \(S\), flat over \(S\) at the
points of \(H_0=H\times_S S_0\), where \(H_0\) is of
quasi-isotrivial multiplicative type. Then the natural map
\[
  \operatorname{Hom}_{S\text{-}\mathrm{gr}}(G,H)
  \longrightarrow
  \operatorname{Hom}_{S_0\text{-}\mathrm{gr}}(G_0,H_0)
\]
is bijective.}

For every integer \(n\geq 0\), let
\[
  S_n=\operatorname{Spec}(A/I^{n+1}),
\]
and let \(G_n,H_n\) be obtained from \(G,H\) by the base change
\(S_n\to S\). Since \(G/S\) is of isotrivial multiplicative type and
\(H\) is affine over \(S\), \SGARefLink{sga3:ix:theorem:7:1}{IX~7.1} shows that the natural
homomorphism
\[
  \operatorname{Hom}_{S\text{-}\mathrm{gr}}(G,H)
  \longrightarrow
  \varprojlim_n
  \operatorname{Hom}_{S_n\text{-}\mathrm{gr}}(G_n,H_n)
\]
is bijective. On the other hand, \SGARefLink{sga3:x:corollary:2:3}{2.3} shows that the \(H_n\) are of
quasi-isotrivial multiplicative type, and \SGARefLink{sga3:x:proposition:2:1}{2.1} shows that all
transition homomorphisms in the inverse system
\[
  \bigl(
    \operatorname{Hom}_{S_n\text{-}\mathrm{gr}}(G_n,H_n)
  \bigr)_n
\]
are isomorphisms. \SGARefLink{sga3:x:lemma:3:1}{Lemma~3.1} follows at once.

\medskip
\noindent\SGARefTarget{sga3:x:theorem:3:2}%
\textbf{Theorem 3.2.} ---
\textit{Let \(A\) be a noetherian ring equipped with an ideal \(I\)
such that \(A\) is separated and complete for the \(I\)-adic topology,
and set \(S=\operatorname{Spec}(A)\) and
\(S_0=\operatorname{Spec}(A/I)\). Then:}
\begin{enumerate}
\item[\SGAInlineRefTarget{sga3:x:theorem:3:2:item:i}(i)]\textit{The functor
\[
  H\longmapsto H_0=H\times_S S_0
\]
is an equivalence between the category of isotrivial groups of
multiplicative type over \(S\) and the analogous category over
\(S_0\).}
\item[\SGAInlineRefTarget{sga3:x:theorem:3:2:item:ii}(ii)]\textit{Let \(H\) be an \(S\)-group of multiplicative type
and of finite type. Then \(H\) is isotrivial if and only if \(H_0\)
is isotrivial.}
\end{enumerate}

\textit{Proof.}
For \SGARefLink{sga3:x:theorem:3:2:item:i}{(i)}, one may either repeat the proof of \SGARefLink{sga3:x:proposition:2:1}{2.1} or use \SGARefLink{sga3:x:corollary:1:2}{1.2}. In either
case one uses the fact that the functor
\[
  S'\longmapsto S'_0=S'\times_S S_0
\]
from the category of finite étale schemes over \(S\) to the category
of finite étale schemes over \(S_0\) is an equivalence (SGA~1, I,
\SGARefLink{sga3:x:lemma:8-4}{8.4}). Equivalently, after reducing to the case where \(S\) is
connected (that is, \(S_0\) is connected) and choosing a geometric
point \(\xi\) of \(S_0\), the canonical homomorphism
\[
  \pi_1(S_0,\xi)\longrightarrow\pi_1(S,\xi)
\]
is an isomorphism.

We prove \SGARefLink{sga3:x:theorem:3:2:item:ii}{(ii)}, namely that if \(H_0\) is isotrivial, then \(H\) is
isotrivial. By \SGARefLink{sga3:x:theorem:3:2:item:i}{(i)}, there exist an isotrivial group \(G\) of
multiplicative type over \(S\) and an \(S_0\)-isomorphism
\[
  u_0:G_0\xrightarrow{\sim}H_0.
\]
\footnote{Editor's note: The original has been expanded in what
follows.}
Since \(H\) is of finite type, so are \(H_0\) and \(G_0\). Therefore,
by \SGARefLink{sga3:ix:proposition:2:1}{IX~2.1}~b), the type of \(G\) at every point of \(S\) is a
finitely generated abelian group, and hence \(G\) is of finite type
over \(S\). Moreover, by \SGARefLink{sga3:x:lemma:3:1}{3.1}, \(u_0\) comes from a homomorphism of
\(S\)-groups
\[
  u:G\longrightarrow H.
\]
Finally, \(G\) and \(H\) are groups of multiplicative type and of
finite type over \(S\), and \(u_0\) is an isomorphism. It follows from
\SGARefLink{sga3:ix:corollary:2:9}{IX~2.9} that \(u\) is an isomorphism, since every neighborhood of
\(S_0\) in \(S\) is equal to \(S\).

\medskip
\noindent\SGARefTarget{sga3:x:corollary:3:3}%
\textbf{Corollary 3.3.} ---
\textit{Let \(A\) be a complete noetherian local ring with residue
field \(k\).}
\begin{enumerate}
\item[\SGAInlineRefTarget{sga3:x:corollary:3:3:item:i}(i)]\textit{Every group of multiplicative type and of finite
type over \(A\) is isotrivial.}
\item[(ii)]\textit{The functor
\[
  H\longmapsto H\otimes_A k=H_0
\]
is an equivalence between the categories of groups of multiplicative
type and of finite type over \(A\) and over \(k\).}
\end{enumerate}

\footnotemark[12]
First, \SGARefLink{sga3:x:corollary:3:3:item:i}{(i)} follows from \SGARefLink{sga3:x:theorem:3:2}{3.2}~\SGARefLink{sga3:x:theorem:3:2:item:ii}{(ii)} and \SGARefLink{sga3:x:proposition:1:4}{1.4}. Then \SGARefLink{sga3:x:theorem:3:2:item:ii}{(ii)} follows from
\SGARefLink{sga3:x:theorem:3:2}{3.2}~\SGARefLink{sga3:x:theorem:3:2:item:i}{(i)}, taking into account that \(H\) is of finite type if and only
if \(H_0\) is of finite type (cf.\ the proof of \SGARefLink{sga3:x:theorem:3:2}{3.2}~\SGARefLink{sga3:x:theorem:3:2:item:ii}{(ii)}).

\medskip
\noindent\SGARefTarget{sga3:x:remark:3:3:1}%
\textbf{Remark 3.3.1.} ---
By \SGARefLink{sga3:x:proposition:1:4}{1.4}, statement \SGARefLink{sga3:x:corollary:3:3}{3.3} gives a complete classification of groups of
multiplicative type and of finite type over \(A\) in terms of the
topological Galois group of an algebraic closure \(\Omega\) of \(k\).

\medskip
\noindent\SGARefTarget{sga3:x:remarks:3:4}%
\textbf{Remarks 3.4.} ---
Under the hypotheses of \SGARefLink{sga3:x:theorem:3:2}{3.2} (that is, \(A\) noetherian, separated and
complete for the \(I\)-adic topology, with no further hypothesis on
\(A/I\)), it will follow from no.~5 that the functor
\[
  H\longmapsto H_0
\]
from the category of groups of multiplicative type and of finite type
over \(S\) to the analogous category over \(S_0\) is still fully
faithful, without any isotriviality hypothesis.\footnote{Editor's
note: Specify this reference, possibly adding a corollary in no.~5
\ldots}

It is not, however, essentially surjective in general. Indeed, there
may exist a group \(H_0\) of multiplicative type and of finite type
over \(S_0\), even locally trivial if desired (but not isotrivial),
that does not arise by reduction from a group \(H\) of multiplicative
type over \(S\).

To see this, take either example from \SGARefLink{sga3:x:paragraph:1:6}{1.6} of a non-isotrivial group of
multiplicative type over a nonnormal curve. The curve may plainly be
taken affine, say \(S_0\), and embedded in affine \(2\)-space; it is
therefore defined by an ideal \(J\) in \(k[X,Y]\). Take \(A\) to be
the completion of \(k[X,Y]\) for the \(J\)-adic topology. Then \(A\)
is a regular ring of dimension \(2\), and \emph{a fortiori} normal.
We shall see in \SGARefLink{sga3:x:theorem:5-16}{5.16} that every group of multiplicative type and of
finite type over \(S=\operatorname{Spec}(A)\) is consequently
isotrivial. Thus \(H_0\), being of finite type and non-isotrivial,
cannot come from a group \(H\) of multiplicative type over \(S\),
since such an \(H\) would necessarily be of finite type and hence
isotrivial.

\medskip
\noindent\SGARefTarget{sga3:x:lemma:3:5}\textbf{Lemma 3.5.}%
\begingroup
\renewcommand{\thefootnote}{\(\ast\)}%
\footnote{Cf.\ also \(\mathrm{VI}_{\mathrm B}\), \SGARefLink{sga3:x:corollary:2:5}{2.5}, for more
systematic developments of this kind.}%
\endgroup
\addtocounter{footnote}{-1}
---
\textit{Let \(S\) be a prescheme and
\(u:G\to H\) a homomorphism of \(S\)-group preschemes that are
locally of finite presentation and flat over \(S\). Let \(U\) be the
set of \(s\in S\) such that the homomorphism induced on the fibers,
\[
  u_s:G_s\longrightarrow H_s,
\]
is flat (respectively, smooth, unramified, étale, or quasi-finite).
Then \(U\) is open, and the restriction
\[
  u|_U:G|_U\longrightarrow H|_U
\]
is flat (respectively, smooth, unramified, étale, or quasi-finite).}

Let \(V\) be the set of points at which \(u\) is flat (respectively,
\(\ldots\)). It is known that \(V\) is open (cf.\ SGA~1, I through IV
in the locally noetherian case, and EGA~IV in
general\footnote{Editor's note: Specify these references. Since
\(G,H\) are locally of finite presentation, \(u\) is locally of
finite presentation (IV\textsubscript{1}, 1.4.3~(v)); one then applies
IV\textsubscript{3}, 11.3.10 and 13.1.4 for flat and quasi-finite,
and IV\textsubscript{4}, 17.4.1, 17.5.1, and 17.6.1 for unramified,
smooth, and étale; compare \(\mathrm{VI}_{\mathrm B}\), \SGARefLink{sga3:x:corollary:2:5}{2.5}.}), and
for \(x\in G\) above \(s\in S\), one has \(x\in V\) if and only if
\(u_s:G_s\to H_s\) is flat (respectively, \(\ldots\)) at \(x\)
(same reference; in the flat, smooth, and étale cases, the flatness of
\(G\) and \(H\) over \(S\) is used here). Since \(u_s\) is a
homomorphism of locally algebraic groups, this is also equivalent to
\(u_s\) being flat (respectively, \(\ldots\)) everywhere
(\(\mathrm{VI}_{\mathrm B}\), \SGARefLink{sga3:x:remark:1:3}{1.3}), that is, to \(s\in U\).
Therefore \(U\) is the inverse image of the open set \(V\) under the
unit section of \(G\), and hence is open. Moreover, \(V=G|_U\), so
\(G|_U\to H|_U\) is flat (respectively, \(\ldots\)), completing the
proof. (N.B. In the unramified and quasi-finite cases, the flatness
hypothesis on \(G\) and \(H\) is unnecessary.)

\medskip
\noindent\SGARefTarget{sga3:x:lemma:3:6}\textbf{Lemma 3.6.} ---
\textit{Let \(H\) be a commutative algebraic group over a field \(k\)
that admits an open subgroup \(G\) of multiplicative type. Then the
family of subschemes \({}_nH\), \(n>0\), is schematically dense in
\(H\). In particular, if \({}_nH={}_nG\) for every \(n>0\), then
\(H=G\).}

Here \({}_nH\) denotes the kernel of
\(n\cdot\operatorname{id}_H\).\footnote{Editor's note: The preceding
sentence has been added.}
Let \(\bar{k}\) be an algebraic closure of \(k\). It is enough to show
that the family \(({}_nH_{\bar{k}})_{n>0}\) is schematically dense in
\(H_{\bar{k}}\), for then \(({}_nH)_{n>0}\) is schematically dense in
\(H\) (cf.\ \SGARefLink{sga3:ix:lemma:4:5}{IX~4.5}). We may therefore suppose that \(k\) is
algebraically closed, so that \(G=D_k(M)\), where \(M\) is an ordinary
finitely generated commutative group.

Let
\[
  M_0=M/\operatorname{Torsion}(M).
\]
Then \(T=D_k(M_0)\) is the largest torus contained in \(G\), and
\(H/T\) is finite. Hence \(H(k)/T(k)\) is annihilated by an integer
\(\nu>0\). There exist finitely many \(g_i\in H(k)\) such that
\[
  H=\coprod_i g_i\cdot G,
\]
and \(g_i^\nu\in T(k)\). Since \(k\) is algebraically closed, the map
\(\nu\cdot\operatorname{id}_T\) is surjective on
\[
  T(k)\simeq(k^\times)^d.
\]
After replacing \(g_i\) by \(g_it_i^{-1}\), where \(t_i\in T(k)\)
satisfies \(t_i^\nu=g_i^\nu\), we may therefore suppose that
\(g_i^\nu=1\). If \(n\) is now a multiple of \(\nu\), then
\[
  {}_nH\supset g_i\cdot{}_nG.
\]
As \(n>0\) varies, the family \({}_nG\) is schematically dense in
\(G\) by \SGARefLink{sga3:ix:theorem:4:7}{IX~4.7}, and the conclusion of \SGARefLink{sga3:x:lemma:3:6}{3.6} follows.

\medskip
\noindent\SGARefTarget{sga3:x:theorem:3:7}%
\textbf{Theorem 3.7.} ---
\textit{Let \(A\) be a noetherian ring equipped with an ideal \(I\)
such that \(A\) is separated and complete for the \(I\)-adic topology;
set \(S=\operatorname{Spec}(A)\) and
\(S_0=\operatorname{Spec}(A/I)\). Let \(H\) be an affine
\(S\)-group scheme of finite type, flat over \(S\) at the points of
\(H_0\), such that \(H_0=H\times_S S_0\) is of multiplicative type
and isotrivial. Then there exists an open and closed subgroup \(G\) of
\(H\), of isotrivial multiplicative type (and of finite type), such
that \(G_0=H_0\).}

By \SGARefLink{sga3:x:theorem:3:2}{3.2}~\SGARefLink{sga3:x:theorem:3:2:item:i}{(i)}, there exist an isotrivial group \(G\) of multiplicative
type over \(S\) and an isomorphism
\[
  u_0:G_0\xrightarrow{\sim}H_0.
\]
By \SGARefLink{sga3:x:lemma:3:1}{3.1}, \(u_0\) comes from a unique homomorphism of \(S\)-groups
\[
  u:G\longrightarrow H.
\]
Using \SGARefLink{sga3:ix:corollary:6:6}{IX~6.6}, we see that \(u\) is a monomorphism: if \(U\) is the set
of \(s\in S\) such that \(u_s:G_s\to H_s\) is a monomorphism, then
\(U\) is an open neighborhood of \(S_0\), hence equals \(S\), and
\(G|_U\to H|_U\) is a monomorphism. By \SGARefLink{sga3:ix:corollary:2:5}{IX~2.5}, \(u\) is even a closed
immersion.

Thus \(G\) is of finite type, and therefore of finite presentation,
over \(S\).\footnote{Editor's note: The preceding sentence has been
added.}
Applying \SGARefLink{sga3:x:lemma:3:5}{Lemma~3.5} in the étale case, there is an open neighborhood
\(U\) of \(S_0\), hence \(U=S\), such that \(G|_U\to H|_U\) is étale.
Thus \(u\) is étale and therefore an open immersion, since it is an
étale monomorphism.\footnote{Editor's note: Cf.\
EGA~IV\textsubscript{4}, 17.9.1.}
This completes the proof.

\medskip
\noindent\SGARefTarget{sga3:x:corollary:3:8}%
\textbf{Corollary 3.8.} ---
\textit{Let \(A\) be a noetherian ring equipped with an ideal \(I\)
such that \(A\) is separated and complete for the \(I\)-adic topology;
set \(S=\operatorname{Spec}(A)\) and
\(S_0=\operatorname{Spec}(A/I)\). Let \(H\) be an \(S\)-group
prescheme of finite type, affine and flat over \(S\).}

\textit{Then \(H\) is of multiplicative type and isotrivial if and
only if \(H_0\) is so and \(H\) satisfies either one of the following
conditions \SGARefLink{sga3:x:corollary:3:8:item:a}{(a)} and \SGARefLink{sga3:x:corollary:3:8:item:b}{(b)}, which are therefore equivalent under the
preceding hypotheses:}
\begin{enumerate}
\item[\SGAInlineRefTarget{sga3:x:corollary:3:8:item:a}(a)]\textit{The fibers of \(H\) are of multiplicative type and
have constant type on each connected component of \(S\).}
\item[\SGAInlineRefTarget{sga3:x:corollary:3:8:item:b}(b)]\textit{\(H\) is commutative and the \({}_nH\), \(n>0\),
are finite over \(S\).}
\end{enumerate}
\textit{Both of these conditions are implied by the following one:}
\begin{enumerate}
\item[\SGAInlineRefTarget{sga3:x:corollary:3:8:item:c}(c)]\textit{The fibers of \(H\) are connected.}
\end{enumerate}

Of course, if \(H\) is of multiplicative type (and isotrivial), then
\SGARefLink{sga3:x:corollary:3:8:item:a}{(a)} and \SGARefLink{sga3:x:corollary:3:8:item:b}{(b)} hold by \SGARefLink{sga3:ix:remark:1:4:1}{IX~1.4.1}~a) and \SGARefLink{sga3:x:proposition:2:1}{2.1}~c). Thus only the ``if''
direction requires proof. We use the subgroup \(G\) provided by \SGARefLink{sga3:x:theorem:3:7}{3.7}.
When \SGARefLink{sga3:x:corollary:3:8:item:c}{(c)} holds, plainly \(G=H\), and there is nothing to prove.

In case \SGARefLink{sga3:x:corollary:3:8:item:b}{(b)}, the immersion \(u:G\to H\) induces an open immersion
\[
  {}_nu:{}_nG\longrightarrow{}_nH
\]
and an isomorphism
\[
  ({}_nG)_0\xrightarrow{\sim}({}_nH)_0.
\]
Since \({}_nH\) is finite over \(S\), it follows at once that
\({}_nu\) is an isomorphism: the complement of its image is finite
over \(S\) and reduces to the empty set. By \SGARefLink{sga3:x:lemma:3:6}{3.6}, the induced
morphisms on fibers
\[
  u_s:G_s\longrightarrow H_s
\]
are isomorphisms. Hence \(u\) is surjective and therefore an
isomorphism.

Finally, in case \SGARefLink{sga3:x:corollary:3:8:item:a}{(a)}, we may suppose that \(S\) is
connected.\footnote{Editor's note: The original said ``that is,
\(S_0\) connected.'' Since \(A\) is separated for the \(I\)-adic
topology, \(S_0\) meets every connected component of \(S\); and since
\(A\) is complete, the connected components of \(S_0\) and \(S\) are
in bijection.}
It follows that, for every \(s\in S\),
\[
  u_s:G_s\longrightarrow H_s
\]
is a monomorphism of algebraic groups of multiplicative type and of
the same type over \(\kappa(s)\).\footnote{Editor's note: This is
because \(G_s\) and \(H_s\) have the same type for every \(s\in S_0\),
and hence for every \(s\).}
Such a homomorphism is necessarily an isomorphism, which again
completes the proof.\footnote{Editor's note: Indeed, \(u\) will again
be a surjective open immersion.}
After extending the base field, we may suppose that both groups over
\(\kappa(s)\) are diagonalizable. The assertion then follows from
\SGARefLink{sga3:viii:corollary:3:2}{VIII~3.2}~b) and the fact that a surjective homomorphism between
isomorphic finitely generated \(\mathbb Z\)-modules
\[
  M\longrightarrow N
\]
is necessarily bijective.

\medskip
\noindent\SGARefTarget{sga3:x:corollary:3:9}%
\textbf{Corollary 3.9.} ---
\textit{Let \(A\) be a noetherian ring equipped with an ideal \(I\)
such that \(A\) is separated and complete for the \(I\)-adic topology,
set \(S=\operatorname{Spec}(A)\), and let \(H\) be an \(S\)-group
prescheme of finite type, affine and flat over \(S\), with connected
fibers. Then \(H\) is an isotrivial torus if and only if
\(H_0=H\times_S\operatorname{Spec}(A/I)\) is an isotrivial torus.}

% END INCLUDED SGA3_English_current_buildable_source_20260823/source/inputs/X/tex/components/01_expose_X_sections1_3.tex

% BEGIN INCLUDED SGA3_English_current_buildable_source_20260823/source/inputs/X/tex/components/02_expose_X_sections4_6.tex
\section*{\SGAInlineRefTarget{sga3:x:section:4}4. Case of an arbitrary base. Quasi-isotriviality theorem}
\phantomsection
\addcontentsline{toc}{section}{4. Case of an arbitrary base. Quasi-isotriviality theorem}

Let \(A\) be a local ring. Recall, following Nagata, that \(A\) is said to be
\emph{henselian} if every algebra \(B\) finite over \(A\) is a product of local
algebras \(B_i\) finite over \(A\).

\par\medskip\noindent
\textbf{\SGAInlineRefTarget{sga3:x:paragraph:4-0}Recall 4.0.}\footnote{\textit{Editors' note.} This recall has been
spelled out, and \SGARefLink{sga3:x:remark:4-0-1}{Remark 4.0.1} has been added.}\enspace
\emph{Let \(A\) be a henselian local ring, \(k\) its residue field,
\(S=\operatorname{Spec}(A)\), \(S_0=\operatorname{Spec}(k)\), and \(\xi\) a
geometric point of \(S_0\). Then the functor}
\[
  X\longmapsto X_0=X\times_S S_0
\]
\emph{is an equivalence between the category of étale coverings of \(S\) and
the analogous category over \(S_0\) (cf.\ EGA IV\(_4\), \S~18.5).
Consequently (cf.\ SGA 1, V),}
\[
  \pi_1(S_0,\xi)=\pi_1(S,\xi).
\]

Suppose in addition that \(A\) is noetherian, and write \(A'\) for its
completion and \(S'=\operatorname{Spec}(A')\). Then \(A'\) is a complete
noetherian local ring, hence henselian (\emph{loc.\ cit.}, 18.5.14), and the
functor
\[
  X\longmapsto X'=X\times_S S'
\]
from the category of étale coverings of \(S\) to the analogous category over
\(S'\) fits into a commutative diagram
\[
  % BEGIN INCLUDED SGA3_English_current_buildable_source_20260823/source/inputs/X/native_diagrams/exp10/004_etale_cover_equivalence_triangle.tex
{\small
\begin{tikzcd}[column sep=2.5em, row sep=2.8em]
  \mathrm{\acute{E}tCov}(S)
    \arrow[rr]
    \arrow[dr,"\sim"']
  &&
  \mathrm{\acute{E}tCov}(S')
    \arrow[dl,"\sim"]
  \\
  &
  \mathrm{\acute{E}tCov}(S_0)
  &
\end{tikzcd}
}

% END INCLUDED SGA3_English_current_buildable_source_20260823/source/inputs/X/native_diagrams/exp10/004_etale_cover_equivalence_triangle.tex

\]
and is therefore also an equivalence of categories, whence
\(\pi_1(S_0,\xi)=\pi_1(S',\xi)\).

\par\medskip\noindent
\textbf{\SGAInlineRefTarget{sga3:x:remark:4-0-1}Remark 4.0.1.}\enspace
Since \(S\) is connected (\(A\) being local), it follows from \SGARefLink{sga3:x:corollary:1:2}{1.2} that the
category of isotrivial groups of multiplicative type over \(S\) is equivalent
to the analogous category over \(S_0\) (and also over \(S'\) if \(A\) is
noetherian).

\par\medskip\noindent
\textbf{\SGAInlineRefTarget{sga3:x:lemma:4-1}Lemma 4.1.}\enspace
\begingroup\itshape
Let \(A\) be a henselian local ring with residue field \(k\),
\(S=\operatorname{Spec}(A)\), and \(S_0=\operatorname{Spec}(k)\).

\smallskip
\noindent
\(\mathrm{(i)}\) The functor
\[
  H\longmapsto H_0=H\times_S S_0
\]
is an equivalence between the category of finite groups of multiplicative type
over \(S\) and the analogous category over \(S_0\).

\smallskip
\noindent
\(\mathrm{(ii)}\) If, in addition, \(A\) is noetherian, writing \(A'\) for its
completion and \(S'=\operatorname{Spec}(A')\), the functor
\[
  H\longmapsto H'=H\times_S S'
\]
is an equivalence between the category of finite groups of multiplicative type
over \(S\) and the analogous category over \(S'\).
\par\endgroup

As in \SGARefLink{sga3:x:paragraph:4-0}{4.0}, the second assertion follows from the first; let us prove the first.
We already know that the functor under consideration is essentially
surjective: every group of multiplicative type \(H_0\) over
\(S_0=\operatorname{Spec}(k)\), being finite and hence of finite type over
\(k\), is isotrivial by \SGARefLink{sga3:x:proposition:1:4}{1.4}, and therefore comes from an isotrivial group of
multiplicative type over \(S\), by \SGARefLink{sga3:x:remark:4-0-1}{4.0.1}.

It remains to prove full faithfulness, that is, that for two finite groups
\(G,H\) of multiplicative type over \(S\), the following map is bijective:
\[
  \operatorname{Hom}_{S\text{-gr.}}(G,H)
  \longrightarrow
  \operatorname{Hom}_{S_0\text{-gr.}}(G_0,H_0).
\]
In other words,\footnote{\textit{Editors' note.} What follows has been added.}
if \(F\) denotes the \(S\)-functor
\(\operatorname{Hom}_{S\text{-gr.}}(G,H)\), the natural map
\[
  \operatorname{Hom}_S(S,F)
  \longrightarrow
  \operatorname{Hom}_{S_0}(S_0,F_0)
\]
induced by the base change \(S_0\to S\) is bijective. In view of Recall \SGARefLink{sga3:x:paragraph:4-0}{4.0},
it is enough to prove the following lemma.

\par\medskip\noindent
\textbf{\SGAInlineRefTarget{sga3:x:lemma:4-2}Lemma 4.2.}\enspace
\emph{Let \(G,H\) be two finite groups of multiplicative type over a
prescheme \(S\). Then
\(F=\operatorname{Hom}_{S\text{-gr.}}(G,H)\) is representable by a finite
étale scheme over \(S\).}

\footnote{\textit{Editors' note.} The original has been expanded in what
follows.}
Let \(f:S'\to S\) be a faithfully flat quasi-compact morphism such that
\(G'\) and \(H'\) are diagonalizable. It is enough to show that \(F_{S'}\) is
represented by a prescheme \(X'\) étale and finite (hence affine) over \(S'\).
Indeed, the descent datum on \(X'\) relative to \(f\) (cf.\ VIII 1.7.2) is then
effective (by SGA 1, VIII 2.1), so there is an \(S\)-prescheme \(X\) such that
\(X\times_S S'=X'\), representing \(F\), and \(X\) is étale and finite over
\(S\) (cf.\ SGA 1, V 5.7 and EGA IV\(_4\), 17.7.3 (ii)).

We may therefore suppose that \(G=D_S(M)\) and \(H=D_S(N)\), where \(M\) and
\(N\) are finite commutative groups (cf.\ \SGARefLink{sga3:viii:proposition:2:1}{VIII 2.1} c)). Then
\(K=\operatorname{Hom}_{\mathrm{gr.}}(N,M)\) is a finite abelian group and,
by \SGARefLink{sga3:viii:corollary:1:5}{VIII 1.5}, there is an isomorphism
\[
  \operatorname{Hom}_{S\text{-gr.}}(G,H)\simeq K_S,
\]
which completes the proof of \SGARefLink{sga3:x:lemma:4-2}{4.2} and hence of \SGARefLink{sga3:x:lemma:4-1}{4.1}.

\par\medskip\noindent
\textbf{\SGAInlineRefTarget{sga3:x:proposition:4-3-0}Proposition 4.3.0.}\footnote{\textit{Editors' note.} This recall,
which is used in the proof of 4.3, has been added (the reference to
EGA IV\(_4\), 18.5.11 was not entirely satisfactory).}\enspace
\begingroup\itshape
Let \(S\) be a henselian local scheme, \(s\) its closed point,
\(g:X\to S\) a morphism locally of finite type, and \(x\) an isolated point
of the fiber \(X_s\).

\smallskip
\noindent
\(\mathrm{(i)}\) Then \(\mathcal O_{X,x}\) is finite over
\(\mathcal O_{S,s}\). (In particular, if the residue-field extension
\(\kappa(x)/\kappa(s)\) is trivial, then
\(\mathcal O_{S,s}\to\mathcal O_{X,x}\) is surjective, by Nakayama's lemma.)

\smallskip
\noindent
\(\mathrm{(ii)}\) If \(g\) is also separated, then
\(X'=\operatorname{Spec}(\mathcal O_{X,x})\) is an open and closed subscheme
of \(X\); that is, there is a decomposition into a disjoint sum
\(X=X'\amalg X''\).
\par\endgroup

\par\smallskip\noindent
\emph{Proof.}
By the local form of Zariski's Main Theorem found in [Pes66], or in
[Ray70, Ch.~IV, Th.~1], \(x\) has an affine open neighborhood
\(U=\operatorname{Spec}(B)\), of finite type and quasi-finite over
\(A=\mathcal O_{S,s}\), and there is an open immersion
\(U\hookrightarrow Y=\operatorname{Spec}(C)\), where \(C\) is a finite
\(A\)-algebra. (N.B. To obtain this conclusion one may also use Chevalley's
semicontinuity theorem (EGA IV\(_3\), 13.1.4), followed by the form of
Zariski's Main Theorem given in \emph{loc.\ cit.}, 8.12.8.)

Since \(A\) is henselian, \(Y\) is the disjoint sum of local schemes
\(Y_1,\ldots,Y_n\), each finite over \(S\), and the points of \(Y\) above
\(s\) are their closed points \(y_1,\ldots,y_n\). Thus \(x=y_i\) for some
\(i\), and
\(\mathcal O_{X,x}=\mathcal O_{U,x}=C_i\) is finite over \(A\). Moreover,
\(X'=\operatorname{Spec}(C_i)\) is an open subscheme of \(U\), hence of
\(X\).

Suppose in addition that \(g\) is separated. Since \(X'\to S\) is finite
(\(C_i\) being finite over \(A\)), so is the immersion \(X'\to X\)
(cf.\ EGA II, 6.1.5 (v)); therefore \(X'\) is also closed in \(X\).
\hfill\(\square\)

\par\medskip\noindent
\textbf{\SGAInlineRefTarget{sga3:x:lemma:4-3}Lemma 4.3.}\enspace
\begingroup\itshape
Let \(A\) be a noetherian henselian local ring, \(A'\) its completion,
\(S=\operatorname{Spec}(A)\), \(S'=\operatorname{Spec}(A')\), and \(s\) the
closed point of \(S\). Let \(G,H\) be two \(S\)-preschemes in groups, with
\(G\) of multiplicative type and of finite type over \(S\), and with \(H\)
locally of finite type and separated over \(S\), \(H_s\) of multiplicative
type, and \(H\) flat over \(S\) at the points of \(H_s\).

Let \(G',H'\) be obtained from \(G,H\) by the base change \(S'\to S\). Then
the following natural map is bijective:
\par\endgroup
\[
  \operatorname{Hom}_{S\text{-gr.}}(G,H)
  \xrightarrow{\ \sim\ }
  \operatorname{Hom}_{S'\text{-gr.}}(G',H').
\]

Since \(S'\) is faithfully flat and quasi-compact over \(S\), SGA 1,
VIII 5.2 shows that this map identifies the first member with the subset of
the second formed by the \(u':G'\to H'\) for which the two inverse images
\(u''_1,u''_2:G''\to H''\) of \(u'\) on
\(S''=S'\times_S S'\), under the projections
\(\operatorname{pr}_1,\operatorname{pr}_2:S''\to S'\), are equal.

It therefore remains to prove that \(u''_1=u''_2\) for every homomorphism of
\(S'\)-groups \(u':G'\to H'\). By the density theorem \SGARefLink{sga3:ix:corollary:4:8}{IX 4.8}, it is enough
to prove that \(u''_1\) and \(u''_2\) agree on \({}_nG''\) for every integer
\(n>0\). (N.B. Here the density theorem is needed essentially in a case
where the base prescheme, namely \(S''\), is not locally noetherian.)

Replacing \(G\) by \({}_nG\), this reduces us to the case in which
\(n\cdot\operatorname{id}_G=0\) for some integer \(n>0\), and hence \(G\)
is finite over \(S\). Likewise let \({}_nH\) be the kernel of the \(n\)-th
power morphism \(\varphi_n\) on \(H\). (N.B. We have not assumed \(H\)
commutative, so \(\varphi_n\), respectively \({}_nH\), need not be a group
homomorphism, respectively a subgroup of \(H\).) Since \({}_nH\) is defined
as the fiber product of \(\varphi_n:H\to H\) and the unit section
\(\varepsilon:S\to H\), its formation commutes with every base change
\(T\to S\):
\[
  ({}_nH)_T={}_n(H_T).
\]

\footnote{\textit{Editors' note.} The original has been expanded in what
follows, on the basis of indications by M.\ Raynaud.}
We use a right subscript \(m\) for reduction modulo
\(\mathfrak m^{m+1}\), where \(\mathfrak m\) is the maximal ideal of \(A\).
Then \(H_m\) is flat over \(S_m\) (because \(H\) is flat over \(S\) at the
points of \(H_0\)); hence, by \SGARefLink{sga3:x:corollary:2:4}{2.4}, \(H_m\) is of multiplicative type and of
finite type, since \(H_0\) is. Thus every
\(({}_nH)_m={}_n(H_m)\) is a subgroup of \(H_m\), of multiplicative type,
finite and flat over \(S_m\).

In particular, \(({}_nH)_0={}_n(H_0)\) is finite over \(S_0\). Since \(S\)
is henselian local and \({}_nH\), like \(H\), is locally of finite type and
separated over \(S\), \SGARefLink{sga3:x:proposition:4-3-0}{Proposition 4.3.0} shows that the local rings of
\({}_nH\) at the points of \(({}_nH)_0\) are finite over \(S\), and that
there is a decomposition as a disjoint sum of open subsets
\[
  \tag{*}\label{sga3:x:lemma:4-3:equation:star}
  {}_nH={}_nH^+\amalg{}_nH^-,
\]
where \(Z={}_nH^+\) is finite over \(S\), while \({}_nH^-\) lies above
\(S-\{s\}\).

Observe that for every finite \(S\)-scheme \(Y\) (such as \(G\)), every
\(S\)-morphism \(Y\to{}_nH\) factors through \(Z\), and that the formation
of the decomposition \((*)\) commutes with the base change \(S'\to S\),
where \(S'=\operatorname{Spec}(A')\) and \(A'\) is the completion of \(A\).

Then \(Z'=Z\times_S S'\) is finite over \(S'\), as is
\(P'=Z'\times_{S'}Z'\). Let \(\nu\) denote the restriction to \(P'\) of
the multiplication on \(H'\), and let \(\sigma\) be the automorphism of
\(P'\) interchanging the two factors. Since \(P'\) is finite over \(S'\),
while \(H'\) is separated and locally of finite type over \(S'\), EGA II
5.4.3 and IV\(_1\) 1.1.3 show that \(Y=\nu(P')\) is a closed subscheme of
\(H'\), universally closed and quasi-compact, hence of finite type and
therefore proper over \(S'\). Moreover, \(Y\to S'\) has finite fibers
(because \(P'\to S'\) does). Since \(S'\) is noetherian, \(Y\to S'\) is
therefore finite (cf.\ EGA III, 4.4.2). Since \(Z'\subseteq Y\) and
\(Z'_m=Y_m\) for every \(m\), Lemma IX 5.0 gives \(Z'=Y\); thus \(Z'\) is
a subgroup of \(H'\). Similarly, the kernel
\(K=\operatorname{Ker}(\nu,\nu\circ\sigma)\) is a closed subscheme of
\(P'\) such that \(K_m=P'_m\) for every \(m\) (because
\(Z'_m={}_nH_m\) is commutative). Hence \(K=P'\), that is, \(Z'\) is a
commutative subgroup of \(H'\).

On the other hand, every reduction \(Z'_m={}_nH_m\) is flat over \(S_m\).
The ``local criterion for flatness'' (cf.\ EGA \(0_{\mathrm{III}}\),
10.2.2, or [BAC], III \S~5, Example 1 and Th.~1) therefore shows that
\(Z'\) is flat over \(S'\). Moreover,
\(Z'_0={}_nH_0\) is a finite group of multiplicative type over \(S_0\),
and hence isotrivial by \SGARefLink{sga3:x:proposition:1:4}{1.4}. It follows from \SGARefLink{sga3:x:corollary:3:8}{3.8} b) that \(Z'\) is of
multiplicative type (and isotrivial) over \(S'\). Since \(S'\to S\) is
faithfully flat and quasi-compact, the multiplication
\(Z\times_S Z\to H\) factors through \(Z\) and makes \(Z\) a subgroup of
\(H\), finite over \(S\) and of multiplicative type (because \(Z'\) is).

Finally, by the observations made above concerning \((*)\), \(Z'\) is the
``finite part'' of \({}_nH'\), and therefore \(u':G'\to{}_nH'\) takes its
values in \(Z'\). Since \(Z\), as well as \(G={}_nG\), is finite over \(S\)
and of multiplicative type, \SGARefLink{sga3:x:lemma:4-1}{Lemma 4.1} shows that \(u'\) comes from a
\(u:G\to Z\), and hence \(u''_1=u''_2\). This completes the proof of \SGARefLink{sga3:x:lemma:4-3}{4.3}.

\par\medskip\noindent
\textbf{\SGAInlineRefTarget{sga3:x:lemma:4-4-0}Lemma 4.4.0.}\footnote{\textit{Editors' note.} This lemma, used
several times below, has been added.}\enspace
\begingroup\itshape
Let \(A\) be a ring, \(S=\operatorname{Spec}(A)\), and let \(H\) be
an \(S\)-group scheme of multiplicative type and of finite type,
quasi-isotrivial (respectively isotrivial). Let \((A_i)\) be a filtered
family of subrings of \(A\) whose inductive limit is \(A\), and put
\(S_i=\operatorname{Spec}(A_i)\).

Then there are an index \(i\) and an \(S_i\)-group scheme \(H_i\), of
multiplicative type and of finite type, quasi-isotrivial (respectively
isotrivial), such that \(H=H_i\times_{S_i}S\).
\par\endgroup

\par\medskip\noindent
\textbf{\SGAInlineRefTarget{sga3:x:theorem:4-4}Theorem 4.4.}\enspace
\begingroup\itshape
Let \(S\) be a prescheme, let \(H\) be an affine \(S\)-prescheme in
groups of finite presentation over \(S\), and let \(s\in S\). Suppose:

\smallskip
\noindent
\(\alpha)\) \(H\) is flat over \(S\) at the points of \(H_s\);

\smallskip
\noindent
\(\beta)\) \(H_s\) is of multiplicative type.

\smallskip
Then there exist an étale morphism \(S'\to S\), a point \(s'\in S'\) above
\(s\) such that \(\kappa(s')/\kappa(s)\) is trivial, and an open and closed
subgroup \(G'\) of \(H'=H\times_S S'\), of multiplicative type and finite
type, and isotrivial, such that \(G'_{s'}=H'_{s'}\).
\par\endgroup

\footnote{\textit{Editors' note.} The reductions that follow have been
spelled out (the original said: ``One reduces at once to the case where
\(S\) is the spectrum of a noetherian local ring \(A\), and where \(s\) is
the closed point of \(S\).'')}
\noindent\textbf{a)}
Write provisionally \(T=\operatorname{Spec}(\mathcal O_{S,s})\), and let us
show, by ``descent of the conclusions'', that one may reduce to the case
\(S=T\). Suppose indeed that we have found an étale morphism \(g:T'\to T\),
a point \(s'\in T'\), and a subgroup \(G'\) of \(H'=H_{T'}\) satisfying the
conditions of the statement. Replacing \(T'\) by an affine open
neighborhood of \(s'\), we may assume that \(T'\) is of finite presentation
over \(T\).

Since \(G'\) is isotrivial, there is a finite étale surjective morphism
\(f:\widetilde T\to T'\) such that
\(\widetilde G=G'\times_{T'}\widetilde T\) is isomorphic to
\(D_{\widetilde T}(M)\), where \(M\) is a finitely generated abelian group.
Since \(T'\), \(H\), \(\widetilde T\), and \(G'\) are of finite presentation
over \(T=\operatorname{Spec}(\mathcal O_{S,s})\), and
\(\mathcal O_{S,s}\) is the inductive limit of the subalgebras
\(A_i=\mathcal O_S(S_i)\), where \(S_i\) runs through the affine open
neighborhoods of \(s\) in \(S\), EGA IV\(_3\), 8.8.2 (and Exposé
VI\(_{\mathrm B}\), 10.2 and 10.3) gives an index \(i\), \(S_i\)-preschemes
\(S'_i,\widetilde T_i\) of finite presentation (respectively an
\(S_i\)-prescheme in groups \(G'_i\) of finite presentation), and morphisms
\[
  g_i:S'_i\to S_i,\qquad f_i:\widetilde T_i\to S'_i,
  \qquad
  u_i:G'_i\to H\times_S S'_i,
\]
from which \(g,f\), respectively \(u:G'\to H'\), are obtained by base
change. Taking \(i\) large enough, \(g_i\) is étale, \(f_i\) is finite
étale surjective, and \(u_i\) is an open and closed immersion
(cf.\ EGA IV, 8.10.5 and 17.7.8).

Now \(\widetilde G\) comes from both
\(\widetilde G_i=G'_i\times_{S'_i}\widetilde T_i\) and
\(D_{\widetilde T_i}(M)\). By EGA IV\(_3\), 8.8.2 (i) (and
VI\(_{\mathrm B}\), 10.2), there is an index \(j\geq i\) such that
\(\widetilde G_j\simeq D_{\widetilde T_j}(M)\), so \(G'_j\) is an
isotrivial group of multiplicative type. Let \(s'_j\) be the image of
\(s'\) in \(S'_j\). Then the étale morphism
\(S'_j\to S_j\hookrightarrow S\), the point \(s'_j\), and the open and
closed subgroup \(G'_j\) of \(H\times_S S'_j\) satisfy the conditions of
the statement. Thus we may assume that \(S=\operatorname{Spec}(A)\) is
local, with closed point \(s\).

\par\smallskip\noindent
\textbf{b)}
The ring \(A\) is the inductive limit of local rings \(A_i\) which are
localizations of finitely generated \(\mathbb Z\)-algebras; put
\(S_i=\operatorname{Spec}(A_i)\). Let us show that hypotheses \(\alpha\)
and \(\beta\) ``descend'' to some \(A_i\). Since \(H\to S\) is of finite
presentation, the set of \(x\in H\) at which \(H\) is flat over \(S\) is
an open set \(W\). It contains \(H_s\) by hypothesis, and hence contains a
quasi-compact open set \(V\) containing \(H_s\) (since the affine scheme
\(H_s\) can be covered by finitely many affine open subsets contained in
\(W\)). The open immersion \(\tau:V\hookrightarrow H\) is of finite
presentation, so \(V\to S\) is also of finite presentation.

By EGA IV\(_3\), 8.8.2 (and Exposé VI\(_{\mathrm B}\), 10.2 and 10.3),
there are an index \(i\), an \(S_i\)-prescheme \(V_i\) of finite
presentation (respectively an \(S_i\)-prescheme in groups \(H_i\) of
finite presentation), and an \(S_i\)-morphism \(\tau_i:V_i\to H_i\), from
which \(V,H,\tau\) are obtained by base change \(S\to S_i\). Taking \(i\)
large enough, \(\tau_i\) is an open immersion and \(V_i\) is flat over
\(S_i\), by EGA IV\(_3\), 8.10.5 and 11.2.6. Let \(s_i\) be the image of
\(s\) in \(S_i\). Since the open immersion
\((V_i)_{s_i}\hookrightarrow(H_i)_{s_i}\), after base change
\(\kappa(s_i)\to\kappa(s)\), gives the equality \(V_s=H_s\), one already
has \((V_i)_{s_i}=(H_i)_{s_i}\). Thus \(H_i\) is flat over \(S_i\) at the
points of \((H_i)_{s_i}\). Finally, \((H_i)_{s_i}\) is of multiplicative
type, because
\[
  H_s=(H_i)_{s_i}\otimes_{\kappa(s_i)}\kappa(s)
\]
is. Hence \((S_i,H_i,s_i)\) satisfies the hypotheses of \SGARefLink{sga3:x:theorem:4-4}{4.4}, and if the
desired assertion holds for this triple, it also holds for \((S,H,s)\) by
base change. We are reduced to the case in which \(A\) is local and
noetherian. We now distinguish two cases.

\par\smallskip\noindent
\(\mathbf{1^\circ}\) \emph{\(A\) is local noetherian and henselian.}
Let \(\widehat A\) be its completion,
\(\widehat S=\operatorname{Spec}(\widehat A)\), and
\(\widehat H=H\times_S\widehat S\). Applying \SGARefLink{sga3:x:theorem:3:7}{Theorem 3.7}, we find an
isotrivial \(\widehat S\)-group \(\widehat G\) of multiplicative type and
of finite type, together with a homomorphism
\[
  \widehat u:\widehat G\longrightarrow\widehat H
\]
which is both an open and a closed immersion, and which induces an
isomorphism
\[
  \widehat u_0:\widehat G_0\xrightarrow{\ \sim\ }\widehat H_0.
\]
By \SGARefLink{sga3:x:remark:4-0-1}{Remark 4.0.1}, base change along \(\widehat S\to S\) induces an
equivalence between the categories of isotrivial groups of multiplicative
type over \(S\) and over \(\widehat S\). In particular, \(\widehat G\)
comes from an isotrivial \(S\)-group \(G\) of multiplicative type and of
finite type. By \SGARefLink{sga3:x:lemma:4-3}{4.3}, \(\widehat u\) comes from a homomorphism
\[
  u:G\longrightarrow H.
\]
Moreover, \(u\) is an open and closed immersion and induces an isomorphism
\(u_0:G_0\to H_0\), since this is true after the faithfully flat
quasi-compact base change \(\widehat S\to S\). This proves \SGARefLink{sga3:x:theorem:4-4}{4.4} in this case
(taking \(S'=S\) and \(s'=s\) in its conclusion).

\par\smallskip\noindent
\(\mathbf{2^\circ}\) \emph{\(A\) is local noetherian.}
The reduction to case \(1^\circ\) is immediate, by applying that case to
the henselization \(A^h\) of \(A\). More precisely, one sees easily (using
SGA 1, I \S~5%
\begingroup
\renewcommand{\thefootnote}{\(\ast\)}
\footnote{Or EGA IV\(_4\), \S~18.6.}
\addtocounter{footnote}{-1}
\endgroup)
that the local rings \(\mathcal O_{S',s'}\) of étale \(S\)-preschemes
\(S'\), equipped with a point \(s'\) above \(s\) for which
\(\kappa(s')/\kappa(s)\) is trivial, form a filtered inductive system whose
inductive limit is a noetherian henselian local ring \(A^h\) (called the
henselization of \(A\)); for details of this construction, due to Nagata in
the normal case, see SGA 4, VIII \S~4. The standard arguments of
EGA IV\(_3\), \S~8 then allow one, as in part a), to deduce from the known
result over the inductive limit \(A^h\) of the \(\mathcal O_{S',s'}\) an
analogous result over one of the pairs \((S',s')\). This proves \SGARefLink{sga3:x:theorem:4-4}{4.4}.

\par\medskip\noindent
\textbf{\SGAInlineRefTarget{sga3:x:corollary:4-5}Corollary 4.5.}\enspace
\emph{Let \(S\) be a prescheme and \(H\) an \(S\)-prescheme in groups of
multiplicative type and of finite type. Then \(H\) is quasi-isotrivial;
that is, \(H\) is split by an étale surjective morphism \(S'\to S\).}

Indeed, let \(s\in S\). By \SGARefLink{sga3:x:theorem:4-4}{4.4}, there are an étale morphism \(S'\to S\), a
point \(s'\in S'\) above \(s\) with \(\kappa(s)=\kappa(s')\), and an
isotrivial subgroup \(G'\) of \(H'\), of multiplicative type and of finite
type, such that \(G'_{s'}=H'_{s'}\). Since \(G'\) and \(H'\) are of
multiplicative type and finite type, \SGARefLink{sga3:ix:corollary:2:9}{IX 2.9} gives an open neighborhood
\(U'\) of \(s'\) such that \(G'|_{U'}=H'|_{U'}\).

\footnote{\textit{Editors' note.} The original has been expanded in what
precedes and what follows.}
Suppose in addition that \(S\) is henselian local, with closed point \(s\).
By EGA IV\(_4\), 18.5.11 b), there is a section \(\sigma\) of \(S'\to S\)
such that \(\sigma(s)=s'\). (N.B. One can see directly that
\(B=\mathcal O_{S',s'}\) equals \(A=\mathcal O_{S,s}\) as follows. By
4.3.0 (i), \(B\simeq A/I\). Since \(B\) is a flat \(A\)-algebra of finite
presentation, \(I\) is a finitely generated ideal (cf.\ EGA IV\(_1\),
1.4.7), and \(I=I\mathfrak m\), where \(\mathfrak m\) is the maximal ideal
of \(A\); hence \(I=0\).) Thus \(H\) is already isotrivial. We obtain:

\par\medskip\noindent
\textbf{\SGAInlineRefTarget{sga3:x:corollary:4-6}Corollary 4.6.}\enspace
\begingroup\itshape
Let \(A\) be a henselian local ring, \(k\) its residue field, and
\(\pi\) the topological Galois group of an algebraic closure of \(k\).

\smallskip
\noindent
\(\mathrm{(i)}\) Every group of multiplicative type and of finite type over
\(S=\operatorname{Spec}(A)\) is isotrivial.

\smallskip
\noindent
\(\mathrm{(ii)}\) The category of these groups over \(S\) is equivalent,
through the functor \(H\mapsto H\otimes_A k\), to the analogous category
over \(S_0=\operatorname{Spec}(k)\). It is therefore anti-equivalent, by
\SGARefLink{sga3:x:proposition:1:4}{1.4}, to the category of Galois modules under \(\pi\) which are finitely
generated as \(\mathbb Z\)-modules.
\par\endgroup

\par\medskip\noindent
\textbf{\SGAInlineRefTarget{sga3:x:corollary:4-7}Corollary 4.7.}\enspace
\begingroup\itshape
Under the hypotheses of \SGARefLink{sga3:x:theorem:4-4}{4.4}, suppose in addition that one of the
following conditions holds:

\smallskip
\noindent
a) for every generization \(t\) of \(s\), \(H_t\) is of multiplicative
type and has the same type as \(H_s\);

\smallskip
\noindent
b) \(H\) is commutative and the \({}_nH\), \(n>0\), are finite over \(S\)
in a neighborhood of \(s\);

\smallskip
\noindent
c) for every generization \(t\) of \(s\), the fiber \(H_t\) is connected.

\smallskip
Then there is an open neighborhood \(U\) of \(s\) such that \(H|_U\) is of
multiplicative type.
\par\endgroup

Since an étale morphism is open, it is enough to show, with the notation of
the conclusion of \SGARefLink{sga3:x:theorem:4-4}{4.4}, that there is an open neighborhood \(U'\) of \(s'\)
such that \(G'|_{U'}=H'|_{U'}\). Put
\(S''=\operatorname{Spec}(\mathcal O_{S',s'})\). Since \(G'\) and \(H'\)
are of finite presentation over \(S'\), EGA IV\(_3\), 8.8.2 reduces the
matter to showing that \(G''=H''\). We may therefore assume that \(S=S''\);
then conditions a), b), c) become those of the following lemma, whose proof
is the same as that of \SGARefLink{sga3:x:corollary:3:8}{3.8}.

\par\medskip\noindent
\textbf{\SGAInlineRefTarget{sga3:x:lemma:4-7-1}Lemma 4.7.1.}\enspace
\begingroup\itshape
Let \(S\) be a local scheme with closed point \(s\), let \(H\) be an
\(S\)-group scheme of finite type, and let \(G\) be an open and closed
subgroup of \(H\), of multiplicative type, such that \(G_s=H_s\). Suppose
in addition that one of the following conditions holds:

\smallskip
\noindent
a) the fibers of \(H\) are of multiplicative type and have the same type as
\(H_s\);

\smallskip
\noindent
b) \(H\) is commutative and the \({}_nH\), \(n>0\), are finite over \(S\);

\smallskip
\noindent
c) the fibers of \(H\) are connected.

\smallskip
Then \(G=H\).
\par\endgroup

\par\medskip\noindent
\textbf{\SGAInlineRefTarget{sga3:x:corollary:4-8}Corollary 4.8.}\enspace
\begingroup\itshape
Let \(S\) be a prescheme and let \(H\) be an affine \(S\)-prescheme
in groups, flat and of finite presentation over \(S\), with fibers of
multiplicative type.

For \(H\) to be of multiplicative type, it is necessary and sufficient that
it satisfy one of the following two conditions (which are equivalent under
the preceding hypotheses):

\smallskip
\noindent
\SGAInlineRefTarget{sga3:x:corollary:4-8:item:a}a) the type of \(H_s\) (cf.\ \SGARefLink{sga3:ix:definition:1:4}{IX 1.4}) is a locally constant function of
\(s\in S\);

\smallskip
\noindent
\SGAInlineRefTarget{sga3:x:corollary:4-8:item:b}b) \(H\) is commutative and the \({}_nH\), \(n>0\), are finite over \(S\).

\smallskip
Moreover, a) and b) are implied by:

\smallskip
\noindent
c) the fibers of \(H\) are connected.
\par\endgroup

In particular:

\par\medskip\noindent
\textbf{\SGAInlineRefTarget{sga3:x:corollary:4-9}Corollary 4.9.}\enspace
\begingroup\itshape
Let \(S\) be a prescheme and \(H\) an \(S\)-prescheme in groups,
flat and of finite presentation over \(S\). Suppose in addition that \(H\)
is affine over \(S\) and has connected fibers.\footnote{\textit{Editors'
note.} This result is generalized in \SGARefLink{sga3:x:theorem:8-1}{8.1}: it is enough to suppose that the
fibers of \(H\) are affine and connected.} If \(s\in S\) is such that
\(H_s\) is a torus, there is an open neighborhood \(U\) of \(s\) such that
\(H|_U\) is a torus.

In particular, if all the fibers of \(H\) are tori, \(H\) is a torus.
\par\endgroup

\section*{\SGAInlineRefTarget{sga3:x:section:5}5. Scheme of homomorphisms from one group of multiplicative type to
another. Twisted constant groups and groups of multiplicative type}
\phantomsection
\addcontentsline{toc}{section}{5. Scheme of homomorphisms from one group of
multiplicative type to another. Twisted constant groups and groups of
multiplicative type}

\par\medskip\noindent
\textbf{\SGAInlineRefTarget{sga3:x:definition:5-1}Definition 5.1.}\enspace
\begingroup\itshape
a) Let \(S\) be a prescheme and \(R\) a prescheme in groups over
\(S\). We say that \(R\) is a \emph{twisted constant group} over \(S\) if
it is locally isomorphic, for the faithfully flat quasi-compact topology,
to a constant group scheme, that is, one of the form \(M_S\) for an
ordinary group \(M\).

\smallskip
\noindent
b) A twisted constant group \(R\) over \(S\) is said to be
\emph{quasi-isotrivial}, respectively \emph{isotrivial}, respectively
\emph{locally isotrivial}, respectively \emph{locally trivial},
respectively \emph{trivial}, if in the preceding definition one can replace
the faithfully flat quasi-compact topology by the étale topology,
respectively the global finite étale topology, respectively the finite
étale topology, respectively the Zariski topology, respectively the
coarsest (or ``chaotic'') topology; cf.\ \SGARefLink{sga3:ix:definition:1:1}{IX 1.1} and \SGARefLink{sga3:x:corollary:1:2}{1.2}, and \SGARefLink{sga3:iv:subsection:6:6}{IV 6.6}.

Thus, to say that \(R\) is quasi-isotrivial (respectively isotrivial) means
that there is an étale surjective (respectively also finite) morphism
\(S'\to S\) such that \(R'=R\times_S S'\) is a constant group over \(S'\);
to say that \(R\) is trivial means that \(R\) is a constant group.

\smallskip
\noindent
c) As in \SGARefLink{sga3:viii:corollary:1:4}{VIII 1.4}, one defines the \emph{type} of a twisted constant group
\(R\) over \(S\) at a point \(s\in S\). It is an isomorphism class of
ordinary groups and, as \(s\) varies, is a locally constant function of
\(s\), hence a constant function if \(S\) is connected. We also say that
\(R\) is ``of type \(M\)'' if all the fibers of \(R\) are of type \(M\).

One should beware that \(R\) is quasi-compact over \(S\) only if it is
finite over \(S\), that is, only if its type at every \(s\in S\) is a
finite group.\footnote{\textit{Editors' note.} See \SGARefLink{sga3:x:lemma:5-12}{Lemma 5.12} below.}

\smallskip
\noindent
d) The case of greatest interest to us is that in which \(R\) is
commutative. We then say that \(R\) is \emph{finitely generated} if its
type at every point \(s\in S\) is given by a finitely generated
\(\mathbb Z\)-module. This notion must not be confused with the scheme-
theoretic notion ``\(R\) of finite type over \(S\)'' (cf.\ above).
\par\endgroup

\par\medskip\noindent
\textbf{\SGAInlineRefTarget{sga3:x:remark:5-2}Remark 5.2.}\enspace
We shall also have to consider \(S\)-preschemes \(X\) which are locally
isomorphic, for the faithfully flat quasi-compact topology, to constant
preschemes, independently of any group structure. We then call \(X\) a
\emph{twisted constant bundle} over \(S\), and extend to these preschemes
the terminology introduced in \SGARefLink{sga3:x:definition:5-1}{5.1}. Of course, when \(X\) is endowed with an
\(S\)-group structure, the meaning of the expressions ``twisted constant'',
``isotrivial'', and so on changes according as one does or does not take the
group structure over \(S\) into account. The same applies when \(X\)
carries any other kind of algebraic structure (for example, that of a
Galois principal bundle, to be considered in the next section).

\par\medskip\noindent
\textbf{\SGAInlineRefTarget{sga3:x:proposition:5-3}Proposition 5.3.}\enspace
\begingroup\itshape
Let \(R\) be a commutative twisted constant group over \(S\).

\smallskip
\noindent
\(\mathrm{(i)}\) The functor \(H=D_S(R)\) (cf.\ VIII 1) is representable
and is a group of multiplicative type over \(S\).

\smallskip
\noindent
\(\mathrm{(ii)}\) For every \(s\in S\), the type of \(R\) at \(s\) equals
the type of \(H\) at \(s\).

\smallskip
\noindent
\(\mathrm{(iii)}\) The group \(R\) is quasi-isotrivial (respectively
isotrivial, respectively trivial, respectively locally isotrivial,
respectively locally trivial) if and only if \(H\) is so.
\par\endgroup

\par\medskip\noindent
\textbf{\SGAInlineRefTarget{sga3:x:remark:5-3-1}Remark 5.3.1.}\footnote{\textit{Editors' note.} This remark, which
appeared after the proof in the original, has been moved here.}\enspace
One may be concerned that in \SGARefLink{sga3:x:proposition:5-3}{5.3} the word ``type'' is being used in two
different senses, according as it refers to \(R\) or to \(H\). Fortunately,
when an \(S\)-prescheme in groups \(G\) is both a twisted constant group
and a group of multiplicative type, its type in either sense is the same,
thanks to the fact that an ordinary finite commutative group is isomorphic
to its dual.

\par\smallskip\noindent
\emph{Proof.}
Covering families for the faithfully flat quasi-compact topology are
families of effective descent for the fibered category of affine
preschemes in groups over variable base preschemes (SGA 1, VIII 2.1).
Consequently \(H\) is representable (and affine over \(S\)), since this is
true ``locally''\footnote{\textit{Editors' note.} Cf.\ VIII \S~1.7.}
(when \(R\) is constant, \(H\) is a diagonalizable group).

That \(H\) is of multiplicative type is then immediate from the definition,
as is the equality of the types of \(R\) and \(H\) at \(s\in S\). Finally,
since
\[
  H_{S'}=D_{S'}(R_{S'}),
\]
the last assertion reduces to the ``trivial'' case, namely to checking that
\(R\) is trivial if and only if \(H\) is trivial; this follows at once from
the biduality theorem \SGARefLink{sga3:viii:theorem:1:2}{VIII 1.2}. \hfill\(\square\)

To complete the description of the correspondence between twisted constant
groups and groups of multiplicative type, one must start with a group of
multiplicative type \(H\) and study \(R=D_S(H)\). If \(R\) is
representable, it is plainly a twisted constant group and
\(H\simeq D_S(R)\). In other words:

\par\medskip\noindent
\textbf{\SGAInlineRefTarget{sga3:x:scholium:5-4-0}Scholium 5.4.0.}\enspace
\emph{The functor \(R\mapsto D_S(R)\) is an
anti-equivalence\footnote{\textit{Editors' note.} ``Equivalence'' has been
corrected to ``anti-equivalence''.} between the category of twisted
constant groups over \(S\) and the category of groups of multiplicative
type \(H\) over \(S\) for which \(D_S(H)\) is representable.}

I do not know whether this condition on \(H\) is always satisfied. We shall
see, however, that it is satisfied when \(H\) is quasi-isotrivial, and in
particular when \(H\) is of finite type.

\par\medskip\noindent
\textbf{\SGAInlineRefTarget{sga3:x:lemma:5-4}Lemma 5.4.}\enspace
\begingroup\itshape
Let \(S'\to S\) be a faithfully flat morphism locally of finite
presentation, and let \(X'\) be an \(S'\)-prescheme separated, locally of
finite presentation, and locally quasi-finite over \(S'\).

Then every descent datum on \(X'\) relative to \(S'\to S\) is effective:
there are an \(S\)-prescheme \(X\) and an \(S'\)-isomorphism
\(X\times_S S'\simeq X'\) compatible with the descent datum.
\par\endgroup

We have already observed that the hypothesis on \(S'\to S\) implies that it
is a covering morphism for the faithfully flat quasi-compact topology
(\SGARefLink{sga3:iv:proposition:6:3:1:item:iv:01}{IV 6.3.1 (iv)}), and hence \emph{a fortiori} a morphism of effective
descent.

When \(X'\to S'\) is quasi-compact, and hence of finite presentation and
quasi-finite, it is a quasi-affine morphism (cf.\ SGA 1, VIII
6.2\footnote{\textit{Editors' note.} When \(S'\) is locally noetherian;
EGA IV\(_3\), 8.11.2 in general.}), and effectivity follows in this case
from SGA 1, VIII 7.9. In the general case, one reduces at once to the case
where \(S\) and \(S'\) are affine. Cover \(X'\) by affine open subsets
\(U'_i\), and let \(V'_i\) be the saturation of \(U'_i\) under the
equivalence relation on \(X'\) defined by the descent datum:
\[
  V'_i=q_2\bigl(q_1^{-1}(U'_i)\bigr).
\]
Here \(q_1,q_2\) are the two projections to \(X'\) of
\[
  X''_1=X'\times_{S'}(S'',p_1)
  \quad\text{and}\quad
  X''_2=X'\times_{S'}(S'',p_2),
\]
respectively: \(q_1=\operatorname{pr}_1\), and \(q_2\) is obtained from the
first projection of \(X''_2\) by means of the given descent isomorphism
\(X''_1\simeq X''_2\). Since \(S'\to S\) is faithfully flat,
quasi-compact, and locally of finite presentation, the same is true of
\[
  p_1:S''=S'\times_S S'\longrightarrow S',
\]
and hence of \(q_1,q_2\). These are consequently open morphisms
(SGA 1, IV 6.6\footnote{\textit{Editors' note.} When \(S'\) is locally
noetherian; EGA IV\(_2\), 2.4.6 in general.}). Therefore \(V'_i\) is an
open quasi-compact part of \(X'\). By the case already proved, the induced
descent data on the \(V'_i\) are effective; hence so is the descent datum on
\(X'\) (SGA 1, VIII 7.2).

\par\medskip\noindent
\textbf{\SGAInlineRefTarget{sga3:x:corollary:5-5}Corollary 5.5.}\enspace
\emph{A faithfully flat morphism of finite presentation \(S'\to S\) is a
morphism of effective descent for the fibered category of twisted constant
groups over variable base preschemes.}

Indeed, this amounts to the effectivity assertion of \SGARefLink{sga3:x:lemma:5-4}{5.4} when \(X'\) is a
constant \(S'\)-prescheme.

\par\medskip\noindent
\textbf{\SGAInlineRefTarget{sga3:x:theorem:5-6}Theorem 5.6.}\enspace
\begingroup\itshape
Let \(S\) be a prescheme, and let \(G,H\) be two quasi-isotrivial
\(S\)-preschemes in groups of multiplicative type, with
\(G\)\footnote{\textit{Editors' note.} \(H\) has been corrected to \(G\).}
of finite type.

Then \(\operatorname{Hom}_{S\text{-gr.}}(H,G)\) is representable and is a
quasi-isotrivial twisted constant group over
\(S\);\footnote{\textit{Editors' note.} What follows has been added.} for
every \(s\in S\), if the type of \(G\) at \(s\) is \(M\) and the type of
\(H\) at \(s\) is \(N\), then the type of
\(\operatorname{Hom}_{S\text{-gr.}}(H,G)\) is
\(\operatorname{Hom}_{\mathrm{gr.}}(M,N)\).
\par\endgroup

The proof proceeds as in \SGARefLink{sga3:x:lemma:4-2}{4.2}, using the fact that the assertion is known
(\SGARefLink{sga3:viii:corollary:1:5}{VIII 1.5}) when \(G\) and \(H\) are trivial. The required effectivity
criterion is provided by \SGARefLink{sga3:x:corollary:5-5}{5.5}, applied to an étale surjective morphism
\(S'\to S\).

Taking \(G=\mathbb G_{m,S}\), in particular, gives:

\par\medskip\noindent
\textbf{\SGAInlineRefTarget{sga3:x:corollary:5-7}Corollary 5.7.}\enspace
\begingroup\itshape
\(\mathrm{(i)}\) Let \(H\) be a quasi-isotrivial \(S\)-group of
multiplicative type. Then the \(S\)-group \(D_S(H)\) is representable and
is a quasi-isotrivial twisted constant group over \(S\).

\smallskip
\noindent
\(\mathrm{(ii)}\) The functors \(R\mapsto D_S(R)\) and
\(H\mapsto D_S(H)\) are mutually quasi-inverse anti-equivalences between
the category of quasi-isotrivial twisted constant \(S\)-groups and the
category of quasi-isotrivial \(S\)-groups of multiplicative type.

\smallskip
\noindent
\(\mathrm{(iii)}\) These functors induce anti-equivalences between the
subcategories of isotrivial, respectively locally isotrivial, respectively
locally trivial, respectively trivial groups.
\par\endgroup

The last assertion is included only for reference, since it is already
contained in \SGARefLink{sga3:x:proposition:5-3}{5.3}.

Moreover, by \SGARefLink{sga3:x:corollary:4-5}{4.5} every \(S\)-group of multiplicative type and of finite type
is quasi-isotrivial. \SGARefLink{sga3:x:theorem:5-6}{Theorem 5.6} therefore gives:

\par\medskip\noindent
\textbf{\SGAInlineRefTarget{sga3:x:corollary:5-8}Corollary 5.8.}\enspace
\emph{Let \(S\) be a prescheme, and let \(G,H\) be two \(S\)-preschemes in
groups of multiplicative type and of finite type. Then
\(\operatorname{Hom}_{S\text{-gr.}}(H,G)\) is representable and is a
finitely generated quasi-isotrivial twisted constant \(S\)-group.}

Also, in \SGARefLink{sga3:x:proposition:5-3}{5.3}, \(R\) is finitely generated if and only if
\(H=D_S(R)\) is of finite type (\SGARefLink{sga3:ix:proposition:2:1}{IX 2.1} b)). By \SGARefLink{sga3:x:corollary:4-5}{4.5}, \(H\) is then
quasi-isotrivial, and hence so is \(R\). Thus:

\par\medskip\noindent
\textbf{\SGAInlineRefTarget{sga3:x:corollary:5-9}Corollary 5.9.}\enspace
\emph{The functors of \SGARefLink{sga3:x:corollary:5-7}{5.7} induce mutually quasi-inverse anti-equivalences
between the category of \(S\)-groups \(H\) of multiplicative type and of
finite type and the category of finitely generated twisted constant
\(S\)-groups \(R\). Moreover, every such group \(R\) is quasi-isotrivial.}

\par\medskip\noindent
\textbf{\SGAInlineRefTarget{sga3:x:corollary:5-10}Corollary 5.10.}\enspace
\begingroup\itshape
Let \(H,G\) be two \(S\)-preschemes in groups of multiplicative type
and of finite type.

\smallskip
\noindent
\(\mathrm{(i)}\)
\(\operatorname{Isom}_{S\text{-gr.}}(H,G)\) is representable by an open and
closed subprescheme of \(\operatorname{Hom}_{S\text{-gr.}}(H,G)\), and is
a twisted constant \(S\)-prescheme.

\smallskip
\noindent
\(\mathrm{(ii)}\) In particular,
\(\operatorname{Aut}_{S\text{-gr.}}(H)\) is representable and is a twisted
constant \(S\)-group, in general noncommutative.
\par\endgroup

This follows from \SGARefLink{sga3:x:corollary:5-8}{5.8} and \SGARefLink{sga3:viii:corollary:1:6}{VIII
1.6}.\footnote{\textit{Editors' note.} This reference should be corrected by
treating the functor \(\operatorname{Isom}_{S\text{-gr.}}(H,G)\) in an
addition 1.5.1 to VIII \ldots}

\par\medskip\noindent
\textbf{\SGAInlineRefTarget{sga3:x:paragraph:5-11-0}Recall 5.11.0.}\footnote{\textit{Editors' note.} This recall, used
several times below, has been added. See also EGA I, 6.1.9; note, however,
that the proof in \emph{loc.\ cit.} seems unnecessarily complicated.}\enspace
Recall that if \(X\) is a locally noetherian prescheme, its connected
components (which are always closed) are open. Indeed, let \(C\) be a
connected component of \(X\), let \(x\in C\), and let \(U\) be a noetherian
open neighborhood of \(x\). The prescheme \(U\) has only finitely many
connected components, so the connected component \(U_x\) of \(x\) in
\(U\) is open in \(U\), hence in \(X\). Since \(U_x\subseteq C\), this
shows that \(C\) is open in \(X\).

\par\medskip\noindent
\textbf{\SGAInlineRefTarget{sga3:x:proposition:5-11}Proposition 5.11.}\enspace
\begingroup\itshape
Let \(S\) be a prescheme, let \(R\) be a commutative twisted constant
group over \(S\), and let \(H=D_S(R)\) be the group of multiplicative type
that it defines. Consider the following conditions:

\smallskip
\noindent
\(\mathrm{(i)}\) \(H\) is isotrivial (that is, \(R\) is isotrivial);

\smallskip
\noindent
\SGAInlineRefTarget{sga3:x:proposition:5-11:item:ii}\(\mathrm{(ii)}\) \(R\) is a union of open and closed subpreschemes \(R_i\)
which are quasi-compact over \(S\) (and then necessarily finite over \(S\));

\smallskip
\noindent
\(\mathrm{(iii)}\) the connected components of \(R\) are finite over \(S\).

\smallskip
\noindent
a) Then \(\mathrm{(i)}\Rightarrow\mathrm{(ii)}
\Rightarrow\mathrm{(iii)}\).\footnote{\textit{Editors' note.} Compare the
examples of \SGARefLink{sga3:x:paragraph:1:6}{1.6} \ldots}

\smallskip
\noindent
b) If \(S\) is locally noetherian, then
\(\mathrm{(i)}\Leftrightarrow\mathrm{(ii)}
\Leftrightarrow\mathrm{(iii)}\).

\smallskip
\noindent
c) Finally, if \(R\) is finitely generated (that is, if \(H\) is of finite
type over \(S\)), then \(\mathrm{(i)}\Leftrightarrow\mathrm{(ii)}\), at
least when \(S\) is quasi-compact or when its connected components are
open.
\par\endgroup

First decompose \(S\) as a prescheme sum of preschemes \(S_i\) on each of
which \(H\) has constant type (cf.\ \SGARefLink{sga3:ix:remark:1:4:1}{IX 1.4.1}). We are reduced to the case
where \(H\), and therefore \(R\), has constant type \(M\). We shall use:

\par\medskip\noindent
\textbf{\SGAInlineRefTarget{sga3:x:lemma:5-12}Lemma 5.12.}\enspace
\emph{Let \(S\) be a prescheme and \(R\) a twisted constant prescheme over
\(S\). Every closed subprescheme \(Z\) of \(R\) which is quasi-compact
over \(S\) is finite over \(S\).}

Indeed, we reduce to the case where \(R\) is constant, hence of the form
\(I_S\) for a set \(I\). It is the filtered increasing union of the
\(J_S\), where \(J\) runs through the finite subsets of \(I\). We may also
assume that \(S\) is affine. Then \(Z\) is quasi-compact, and is therefore
contained in some \(J_S\). Since \(J_S\) is finite over \(S\), so is \(Z\).

\SGARefLink{sga3:x:lemma:5-12}{Lemma 5.12} already proves the parenthetical assertion in
\SGARefLink{sga3:x:proposition:5-11:item:ii}{5.11 (ii)}.\footnote{\textit{Editors' note.} The original has been expanded
in what follows.} The implication
\(\mathrm{(ii)}\Rightarrow\mathrm{(iii)}\) is then clear, since the
connected components of \(R\) are closed. If \(S\) is locally noetherian,
they are open and closed by \SGARefLink{sga3:x:paragraph:5-11-0}{5.11.0} (\(R\) being étale and hence locally of
finite type over \(S\)); this gives
\(\mathrm{(iii)}\Rightarrow\mathrm{(ii)}\) in that case.

Let us prove \(\mathrm{(i)}\Rightarrow\mathrm{(ii)}\). Choose a finite
étale surjective morphism \(S'\to S\) which splits \(H\), and hence \(R\)
(cf.\ \SGARefLink{sga3:x:proposition:5-3}{5.3}), so that
\[
  R'\simeq M_{S'}=\coprod_{m\in M}R'_m,
\]
where the \(R'_m\) are disjoint open subsets of \(R'\), each
\(S'\)-isomorphic to \(S'\). Let \(g:R'\to R\) be the projection. It is a
finite étale surjective morphism, and hence both open and closed. The
subsets \(R_m=g(R'_m)\) are therefore open and closed parts of \(R\), and
are plainly quasi-compact over \(S\), since the \(R'_m\) are.

Finally, suppose that \(H\) is of finite type over \(S\), and let us prove
\(\mathrm{(ii)}\Rightarrow\mathrm{(i)}\). If the connected components of
\(S\) are open, the matter reduces at once to the case where \(S\) is
connected; thus we may assume that \(S\) is quasi-compact or connected.
Since \(M\) is finitely generated, write
\[
  M=\mathbb Z^r\times N,
\]
where \(r\geq 0\) and \(N\) is a finite abelian group. Let
\(G=D_S(M)\), and consider the preschemes
\[
  P=\operatorname{Isom}_{S\text{-gr.}}(H,G)
  \subset
  \operatorname{Hom}_{S\text{-gr.}}(H,G)=Q
\]
(cf.\ \SGARefLink{sga3:x:corollary:5-8}{5.8} and \SGARefLink{sga3:x:corollary:5-10}{5.10}). There are isomorphisms
\begin{align*}
  Q
  &\simeq \operatorname{Hom}_{S\text{-gr.}}(M_S,R)\\
  &\simeq
    \operatorname{Hom}_{S\text{-gr.}}(\mathbb Z^r_S,R)
    \times
    \operatorname{Hom}_{S\text{-gr.}}(N_S,R)\\
  &\simeq R^r\times E,
\end{align*}
where
\[
  E=\operatorname{Hom}_{S\text{-gr.}}(N_S,R)
\]
is finite over \(S\).\footnote{\textit{Editors' note.} Indeed, it is a
twisted constant group of type \(\operatorname{End}_{\mathrm{gr.}}(N)\);
cf.\ \SGARefLink{sga3:x:theorem:5-6}{5.6} and \SGARefLink{sga3:x:corollary:5-8}{5.8}.} It follows that \(Q\) is a union of open and closed
subpreschemes \(Q_i\) finite over \(S\). Hence \(P\) is a union of the open
and closed subpreschemes \(P_i=P\cap Q_i\), finite over \(S\). Since they
are étale over \(S\), their images in \(S\) are open and closed parts
\(S_i\), and these cover \(S\). If \(S\) is connected or quasi-compact,
there is a finite set of indices \(i\) for which the \(S_i\) cover \(S\).
Let \(S'\) be the union of the corresponding \(P_i\). Then
\(S'\to S\) is finite étale surjective. Putting
\[
  P'=P\times_S S'
     =\operatorname{Isom}_{S'\text{-gr.}}(H',G'),
\]
we see that \(P'\) has a section over \(S'\); hence there is an isomorphism
of \(S'\)-groups
\[
  H'=H\times_S S'
  \xrightarrow{\ \sim\ }
  G'=G\times_S S'=D_{S'}(M).
\]
Thus \(S'\) splits \(H\), completing the proof of
\SGARefLink{sga3:x:proposition:5-11}{5.11}.\footnote{\textit{Editors' note.} This is modulo verifying
\(\mathrm{(ii)}\Rightarrow\mathrm{(i)}\) when \(S\) is locally noetherian
and \(R\) is not finitely generated \ldots}

\par\medskip\noindent
\textbf{\SGAInlineRefTarget{sga3:x:remark:5-11-bis}Remark 5.11 bis.}\enspace
If \(H\) is of finite type over \(S\) and of constant type, one can also
give the following isotriviality criterion, with no restriction on \(S\):
\(H=D_S(R)\) is isotrivial if and only if \(R\) is the union of a sequence
of open and closed parts finite over \(S\).\footnote{\textit{Editors'
note.} To spell this out: since \(M\) is a finitely generated
\(\mathbb Z\)-module, it is countable, and the proof of
\SGARefLink{sga3:x:proposition:5-11}{5.11} \(\mathrm{(i)}\Rightarrow\mathrm{(ii)}\) shows that \(R\) is the
union of a countable family of open and closed parts finite over \(S\).
The converse would have to be checked \ldots}

\par\medskip\noindent
\textbf{\SGAInlineRefTarget{sga3:x:lemma:5-13}Lemma 5.13.}\enspace
\emph{Let \(S\) be a locally noetherian connected prescheme, let \(P\) be
a quasi-isotrivial twisted constant \(S\)-prescheme, and let \(Z\) be an
open and closed part of \(P\). If \(Z_s\) is finite for some \(s\in S\),
then \(Z\) is finite over \(S\).}

First, do not assume \(S\) connected. Let \(U\) be the set of \(s\in S\)
for which \(Z_s\) is finite. We shall prove that \(U\) is both open and
closed, and that \(Z|_U\) is finite over \(U\). This statement is
essentially equivalent to \SGARefLink{sga3:x:lemma:5-13}{5.13}, but has the advantage of being local in
\(S\) for the étale topology, say. We are therefore reduced to the case
where \(P\) is constant, of the form \(I_S\) for a set \(I\). (N.B. The
locally noetherian hypothesis is preserved by étale base change; this is
where the quasi-isotriviality of \(P\) over \(S\) is used.)

We may also suppose that \(S\) is connected, since its connected components
are open (\(S\) being locally noetherian). Then necessarily \(Z=J_S\) for
a subset \(J\subseteq I\), and consequently \(U=\varnothing\) or \(U=S\)
according as \(J\) is infinite or finite. This proves the assertion.

\par\medskip\noindent
\textbf{\SGAInlineRefTarget{sga3:x:paragraph:5-14-0}Recalls 5.14.0.}\footnote{\textit{Editors' note.} The original
treated in \SGARefLink{sga3:x:corollary:5-14}{5.14} the case where \(S\) is locally noetherian and normal, and
observed in Remark 5.15 that the argument applies more generally when \(S\)
is merely geometrically unibranch rather than normal. The statement of \SGARefLink{sga3:x:corollary:5-14}{5.14}
(and also \SGARefLink{sga3:x:theorem:5-16}{5.16}) has accordingly been modified, and these ``recalls'', drawn
from EGA IV\(_4\), 18.10.6 and 18.10.7, have been added; they show that the
proof of 5.14 applies verbatim to the geometrically unibranch case.}\enspace
Let \(S\) be a locally noetherian prescheme and let \(\widetilde S\) be the
normalization of \(S_{\mathrm{red}}\) (cf.\ EGA II, 6.3.8). Recall that
\(S\) is said to be \emph{geometrically unibranch} (cf.\ EGA
\(0_{\mathrm{IV}}\), \S~23.2 and IV\(_2\), \S~6.15) if the canonical
morphism \(\widetilde S\to S\) is radicial, and hence a universal
homeomorphism. In particular, the connected components of \(S\) are
irreducible.

Suppose now that \(S\) is connected, and hence irreducible. Let \(\eta\) be
its generic point and let \(f:P\to S\) be a flat, locally quasi-finite
morphism. Let \(P_i\) be the irreducible components of \(P\), with generic
points \(\xi_i\). Since \(P\) is flat over \(S\), every \(\xi_i\) lies
above \(\eta\) (cf.\ EGA IV\(_2\), 2.3.4), and
\[
  (P_i)_\eta=P_i\cap P_\eta
\]
is the closure of \(\xi_i\) in \(P_\eta\). Since the fiber \(P_\eta\) is
discrete, one has \((P_i)_\eta=\{\xi_i\}\). This applies in particular when
\(f\) is étale. In that case \(P\) is also locally noetherian and
geometrically unibranch (cf.\ EGA IV\(_4\), 17.5.7), so its irreducible
components are its connected components and are open (and closed).

\par\medskip\noindent
\textbf{\SGAInlineRefTarget{sga3:x:corollary:5-14}Corollary 5.14.}\enspace
\emph{Let \(S\) be a locally noetherian geometrically unibranch prescheme,
and let \(P\) be a quasi-isotrivial twisted constant \(S\)-prescheme. Then
the connected components of \(P\) are finite over \(S\).}

We may plainly suppose that \(S\) is connected and hence irreducible, and
write \(\eta\) for its generic point. By the preceding recalls, every
connected component \(P_i\) of \(P\) is open and closed and meets the fiber
\(P_\eta\) in a single point. \SGARefLink{sga3:x:lemma:5-13}{Lemma 5.13} therefore applies and shows that
every \(P_i\) is finite over \(S\).

\par\medskip\noindent
\textbf{\SGAInlineRefTarget{sga3:x:theorem:5-16}Theorem 5.16.}\footnote{\textit{Editors' note.} Remark 5.15,
rendered obsolete by the addition of \SGARefLink{sga3:x:paragraph:5-14-0}{5.14.0} (cf.\ the preceding editors'
note), has been deleted, and in \SGARefLink{sga3:x:theorem:5-16}{5.16} ``normal'' has been replaced by
``geometrically unibranch''.}\enspace
\emph{Let \(S\) be a locally noetherian geometrically unibranch prescheme.
Then every \(S\)-group of multiplicative type and of finite type is
isotrivial.}

Indeed, we may suppose \(S\) connected, and hence \(H\) of constant type
\(M\). It is enough to apply \SGARefLink{sga3:x:corollary:5-14}{5.14} to
\[
  P=\operatorname{Isom}_{S\text{-gr.}}(H,G),
  \qquad G=D_S(M),
\]
and then argue as in the proof of
\SGARefLink{sga3:x:proposition:5-11}{5.11} \(\mathrm{(ii)}\Rightarrow\mathrm{(i)}\). Alternatively, one may
apply \SGARefLink{sga3:x:corollary:5-14}{5.14} to \(P=R=D_S(H)\) (cf.\ \SGARefLink{sga3:x:corollary:5-9}{5.9}), and then use \SGARefLink{sga3:x:proposition:5-11}{5.11}.

\section*{\SGAInlineRefTarget{sga3:x:section:6}6. Infinite Galois principal bundles and the enlarged fundamental group}
\phantomsection
\addcontentsline{toc}{section}{6. Infinite Galois principal bundles and the
enlarged fundamental group}

(The results of the present section and of the next will not be used again
in the remainder of this Seminar.)

Let \(S\) be a prescheme. We propose to determine the principal homogeneous
bundles \(P\) over \(S\) whose structural group is of the form \(G_S\), the
constant \(S\)-group defined by an ordinary group \(G\), not necessarily
finite. We shall also call them \emph{Galois principal bundles over \(S\)
with group \(G\)}. We use ``principal bundle'' in the sense of the
faithfully flat quasi-compact topology (cf.\ Exposé IV, Def.~\SGARefLink{sga3:iv:definition:5:1:5}{5.1.5}). For
such a \(P\), however, the structural morphism \(P\to S\) is necessarily
étale and surjective, hence a covering for the étale topology; consequently
\(P\) is also locally trivial for the étale topology (cf.\ IV,
Prop.~\SGARefLink{sga3:iv:proposition:5:1:6}{5.1.6}).\footnote{\textit{Editors' note.} Notice that \(P\) is assumed
to be a prescheme. In the \emph{a priori} more general case of an fpqc sheaf
\(P\) which is a \(G_S\)-torsor, is \(P\) necessarily representable?}

We assume that \(S\) is a sum of connected preschemes, that is, that its
connected components are open; this reduces us at once to the case in which
\(S\) is connected. Choose a ``geometric point'' \(\xi\) of \(S\), that
is, an \(S\)-scheme \(\xi\) which is the spectrum of an algebraically
closed field \(\Omega=\kappa(\xi)\). For every Galois principal bundle
\(P\) over \(S\) with group \(G\), \(P_\xi\) is a Galois principal bundle
over the algebraically closed field \(\kappa(\xi)\), and is therefore
trivial. We refine the initial problem by asking for the category of Galois
principal bundles over \(S\) \emph{pointed} above \(\xi\), that is, endowed
with an \(S\)-morphism \(\xi\to P\), or equivalently with a trivialization
of \(P_\xi\). For fixed \(G\), let
\(\bar\pi^{\,1}(S,\xi;G)\) denote the set of classes of such bundles, up to
an isomorphism respecting the base point. Then the set
\(\bar\pi^{\,1}(S;G)\) of isomorphism classes of Galois principal bundles
over \(S\) with group \(G\), without a specified base point, is the set of
orbits of \(G\) in \(\bar\pi^{\,1}(S,\xi;G)\). Here \(G\) acts naturally
by translating the marked point of a pointed Galois principal bundle:
\[
  \bar\pi^{\,1}(S;G)=\bar\pi^{\,1}(S,\xi;G)/G.
\]

Let \(S'\to S\) be any morphism which is a morphism of universal effective
descent for the fibered category of twisted constant preschemes over a
variable base (for example, \(S'\to S\) faithfully flat and locally of
finite presentation, cf.\ 5.4; for other examples, cf.\ SGA 1 IX). We
propose to determine the subsets
\[
  \bar\pi^{\,1}(S'/S,\xi;G)
  \quad\text{and}\quad
  \bar\pi^{\,1}(S'/S;G)
\]
of the preceding sets, formed by the Galois principal bundles over \(S\)
which become trivial over \(S'\), or, as one says, are ``split'' by \(S'\).
In fact, we shall determine the category of Galois principal bundles \(P\)
over \(S\) split by \(S'\). Of course,
\[
  \bar\pi^{\,1}(S,\xi;G)
  =
  \varinjlim_{S'}\bar\pi^{\,1}(S'/S,\xi;G),
\]
where \(S'\) runs through a cofinal system among the \(S'/S\) which cover
for the étale topology (for example, if \(S\) is quasi-compact, among the
quasi-compact \(S'\) over \(S\) whose structural morphism is étale and
surjective). Likewise, the category of Galois principal bundles over \(S\)
is the inductive limit of the subcategories defined by the \(S'\), formed
by the bundles split over \(S'\).

By the assumption on \(S'/S\), the category of Galois principal bundles
over \(S\) split by \(S'\) is equivalent to the category of trivial Galois
principal bundles over \(S'\), hence of the form \(G_{S'}\), with \(G\)
acting by right translations, equipped with descent data relative to
\(S'\to S\). In terms of the corresponding trivial bundle \(P'\) over
\(S'\) and its descent datum, a base point on a Galois principal bundle
\(P\) over \(S\) split by \(S'\) amounts to a trivialization of
\(P'\times_{S'}S'_\xi\) compatible with the induced descent datum relative
to \(S'_\xi\to\xi\), where \(S'_\xi=S'\times_S\xi\); equivalently, it is a
section \(\sigma\) of \(P'_\xi\) over \(S'_\xi\) compatible with the
descent datum.

It is therefore useful, for an arbitrary \(S\)-prescheme \(S'\), without
now assuming that \(S'\to S\) is a morphism of universal effective descent
for the fibered category of twisted constant bundles, to define
\(\bar\pi^{\,1}(S'/S;G)\) and
\(\bar\pi^{\,1}(S'/S,\xi;G)\) as the sets of isomorphism classes of the
structures with descent data just specified. These functors in \(G\) then
have a very simple simplicial description in terms of the fibered square
and cube \(S''\) and \(S'''\) of \(S'\) over \(S\); we shall sketch it
below (cf.\ \SGARefLink{sga3:x:paragraph:6-3}{6.3}).

The important conclusion is the following.

\par\medskip\noindent
\textbf{\SGAInlineRefTarget{sga3:x:proposition:6-1}Proposition 6.1.}\enspace
\begingroup\itshape
Suppose that the connected components of \(S'\) and \(S''\) are open, and,
for example, that the quotient of \(\pi_0(S')\) by the equivalence relation
induced by the two projections
\(\pi_0(S'')\to\pi_0(S')\) is a single point.\footnote{\textit{Editors'
note.} The hypothesis that the connected components of \(S'\) be open, as
well as the second hypothesis, has been added. The latter says that the
simplicial set \(K_\bullet=\pi_0(S_\bullet)\) defined in \SGARefLink{sga3:x:paragraph:6-3}{6.3} is connected.
In fact a weaker hypothesis suffices: the cone \(\widetilde K_\bullet\) of
a certain morphism of simplicial sets
\(k_\bullet\to K_\bullet\) need only be connected (cf.\ \emph{loc.\ cit.}).}

\smallskip
\noindent
\(\mathrm{(i)}\) The functor
\[
  G\longmapsto\bar\pi^{\,1}(S'/S,\xi;G)
\]
from the category of groups to the category of sets is representable by a
group, denoted \(\pi_1(S'/S,\xi)\) and called the \emph{fundamental group
of \(S\) at \(\xi\) relative to \(S'\to S\)}. Thus there is a functorial
bijection
\[
  \bar\pi^{\,1}(S'/S,\xi;G)
  \simeq
  \operatorname{Hom}_{\mathrm{gr.}}\bigl(\pi_1(S'/S,\xi),G\bigr).
\]

\smallskip
\noindent
\(\mathrm{(ii)}\) This group has a set of generators in bijection with
\(\pi_0(S'')\), and is described in terms of those generators by relations
in bijection with the elements of
\(\pi_0(S''')\).\footnote{\textit{Editors' note.} The superfluous
hypothesis that the connected components of \(S'''\) be open has been
deleted. On the other hand, the later description (cf.\ \ldots) gives the
natural set of generators
\(\pi_0(S'')\amalg\pi_0(S'_\xi)\); the relations between these generators
arising from the \(2\)-cells then reduce it to \(\pi_0(S'')\). See, for
example, [Kan58], \S~19 for a finer description.} In particular,
\(\pi_1(S'/S,\xi)\) is finitely generated (respectively finitely presented)
if \(\pi_0(S'')\) is finite (respectively if both \(\pi_0(S'')\) and
\(\pi_0(S''')\) are finite).

\smallskip
\noindent
\(\mathrm{(iii)}\) The category of Galois principal bundles over \(S\)
split by \(S'\), with base point above \(\xi\), is equivalent to the
category of ordinary groups \(G\) equipped with a homomorphism
\(\pi_1(S'/S,\xi)\to G\).
\par\endgroup

The proof is given below; cf.\ \ldots

\par\medskip\noindent
\textbf{\SGAInlineRefTarget{sga3:x:paragraph:6-2}6.2.}\enspace
When \(S\) is a connected locally noetherian prescheme, every étale
prescheme \(S'\) over \(S\) is locally noetherian and therefore has open
connected components. It follows from the preceding
discussion\footnote{\textit{Editors' note.} This deduction could be
spelled out: \(I\) is in bijection with a cofinal set of morphisms
\(S'\to S\) covering for the étale topology \ldots} that the functor
\[
  G\longmapsto\bar\pi^{\,1}(S,\xi;G)
\]
from ordinary groups to sets is \emph{strictly pro-representable}
(cf.\ \emph{Séminaire Bourbaki}, February 1960, No.~195,
\S\S~A.2 and A.3). Thus there is a projective system
\[
  \Pi=\Pi_1(S;\xi)=(\pi_i)_{i\in I}
\]
of ordinary groups over an increasing filtered index set \(I\), which is
``strict'', meaning that its transition morphisms
\(\pi_j\to\pi_i\) are surjective, together with an isomorphism of functors
in \(G\)
\[
  \bar\pi^{\,1}(S,\xi;G)
  \simeq
  \varinjlim_i\operatorname{Hom}_{\mathrm{gr.}}(\pi_i,G).
\]
The second member is also denoted
\(\operatorname{Hom}_{\mathrm{pro\text{-}gr.}}(\Pi,G)\)
(cf.\ \emph{loc.\ cit.}).

If the projective limit
\[
  \pi=\varprojlim_i\pi_i
\]
is ``large enough'', more precisely if the canonical homomorphisms
\(\pi\to\pi_i\) are surjective, it is natural to equip \(\pi\) with the
projective-limit topology of the discrete topologies on the \(\pi_i\).
The preceding isomorphism can then also be written
\[
  \bar\pi^{\,1}(S,\xi;G)
  \simeq
  \operatorname{Hom}_{\mathrm{gr.\,top.}}(\pi,G),
\]
where the right-hand side is the set of homomorphisms of topological
groups, \(G\) being given the discrete topology.

The hypothesis just formulated on the projective system \(\Pi\) holds, as
is well known, when the \(\pi_i\) are finite groups
(cf.\ [BEns], III \S~\SGARefLink{sga3:x:remark:7-4}{7.4}, Th.~1). This last condition plainly also means
that every Galois principal bundle over \(S\) is isotrivial, that is, split
by a finite étale surjective morphism. This is the case when \(S\) is
geometrically unibranch (for example normal), as follows at once from
\SGARefLink{sga3:x:corollary:5-14}{5.14}.\footnote{\textit{Editors' note.} \SGARefLink{sga3:x:lemma:5-13}{5.13} has been corrected to \SGARefLink{sga3:x:corollary:5-14}{5.14}.}
When the \(\pi_i\) are finite, \(\pi\) also coincides with the fundamental
group \(\pi_1(S,\xi)\) introduced in SGA 1, V.

Thus, in the favorable case in which the maps \(\pi\to\pi_i\) are
surjective, one might call \(\pi\) the \emph{enlarged fundamental group of
\(S\) at \(\xi\)}. Outside this favorable case, \(\pi\) itself is of
little interest, and the role of the usual fundamental group is played by
the projective system \(\Pi\) itself, which we call the \emph{enlarged
fundamental pro-group of \(S\) at \(\xi\)}. (Any better terminological
suggestion than ``enlarged'' is welcome!\footnote{\textit{Editors' note.}
Some authors speak of the ``true'' fundamental group.}) Knowledge of this
pro-group is more precise than knowledge of the usual fundamental group
\(\pi_1(S,\xi)\) of SGA 1, V; more precisely, the latter is the projective
limit of the projective system formed by the finite quotients of the
\(\pi_i\).

\par\medskip\noindent
\textbf{\SGAInlineRefTarget{sga3:x:paragraph:6-3}6.3.}\enspace
Let us indicate briefly the ``calculation'' of \(\pi_1(S'/S,\xi)\). Let
\(S_i\) be the \((i+1)\)-st fibered power of \(S'\) over \(S\), so that
\(S_0=S'\), \(S_1=S''\), and so on. The evident simplicial operations
between the \(S_i\) make \((S_i)_{i\in\mathbb N}\) a simplicial object of
\((\mathrm{Sch})/S\).

Applying to this simplicial object the functor ``set of connected
components''
\[
  \pi_0:(\mathrm{Sch})/S\longrightarrow(\mathrm{Ens}),
\]
we obtain a simplicial set \(K_\bullet=(K_i)_{i\in\mathbb N}\), with
\[
  K_i=\pi_0(S_i).
\]

Likewise, the \(S_{i,\xi}\), the \((i+1)\)-st fibered powers of
\(S'_\xi\) over \(\xi\), form a simplicial object of
\((\mathrm{Sch})/\xi\), hence of \((\mathrm{Sch})/S\), together with a
natural morphism of simplicial objects into \((S_i)_{i\in\mathbb N}\).
Thus we obtain a simplicial set \(k_\bullet\), with
\[
  k_i=\pi_0(S_{i,\xi}),
\]
and a canonical morphism
\[
  k_\bullet\longrightarrow K_\bullet.
\]
Taking the cone of this morphism (cf.\ \SGARefLink{sga3:x:definition:9-5-1}{9.5.1}), we form a new simplicial set
\[
  \widetilde K_\bullet
  =
  \operatorname{Cone}(k_\bullet\longrightarrow K_\bullet).
\]

\footnote{\textit{Editors' note.} The original, which takes
\(\widetilde K_\bullet=K_\bullet/k_\bullet\), has been corrected: on the
one hand \(k_\bullet\) is already contractible (cf.\ \SGARefLink{sga3:x:proposition:9-7-1}{9.7.1}), and on the
other hand \(k_\bullet\to K_\bullet\) is not in general injective. It is an
epimorphism if \(S'/S\) is finite étale (cf.\ \SGARefLink{sga3:x:proposition:9-7-3}{9.7.3}).}
In this way we obtain a ``pointed simplicial set'' \(\widetilde K_\bullet\),
that is, a simplicial set equipped with a morphism
\[
  \widetilde\xi:e_\bullet\longrightarrow\widetilde K_\bullet,
\]
where \(e_\bullet\) is the final simplicial set. We may form its familiar
combinatorial invariants
\[
  \pi_0(\widetilde K_\bullet,\widetilde\xi)
  \quad\text{and}\quad
  \pi_1(\widetilde K_\bullet,\widetilde\xi),
\]
whose constructions involve only the components of degree at most \(1\)
and at most \(2\), respectively. These invariants are defined without any
restriction on \(S\) or \(S'\). When the connected components of \(S_0\)
and \(S_1\) are open and \(\widetilde K_\bullet\) is
connected,\footnote{\textit{Editors' note.} The hypothesis that
\(\widetilde K_\bullet\) be connected has been added; for the proof, see
the addendum below (\SGARefLink{sga3:x:section:9}{Section 9}).} one checks without difficulty that
\(\pi_1(\widetilde K_\bullet,\widetilde\xi)\) represents the functor
\(G\mapsto\bar\pi^{\,1}(S'/S,\xi;G)\); that is,
\[
  \pi_1(S'/S;\xi)
  \simeq
  \pi_1(\widetilde K_\bullet,\widetilde\xi).
\]

We also note that if \(S'\to S\) is ``universally submersive''
(cf.\ SGA 1, IV 2.1) and the connected components of \(S'\) are open, then
the simplicial set \(K_\bullet\), and hence also \(\widetilde K_\bullet\),
is connected.\footnote{\textit{Editors' note.} The sequel has been
modified to take account of the correction made above; cf.\ Editors' note
(53). (The original read:
``\(\pi_0(\widetilde K_\bullet,\widetilde\xi)\) is also canonically
isomorphic to the set \(\pi_0(S,\xi)\) of connected components of \(S\),
pointed by the connected component of \(\xi\) in \(S\).'')}

\par\medskip\noindent
\textbf{\SGAInlineRefTarget{sga3:x:example:6-4}Examples 6.4.}\enspace
It remains to give examples of enlarged fundamental groups. Return to the
examples of \SGARefLink{sga3:x:paragraph:1:6}{1.6}. Let \(k\) be an algebraically closed field; let \(C_1\) be
a complete rational curve over \(k\) with exactly one singular point, that
point being an ordinary double point; and let \(C_2\) be a curve which is
the union of two irreducible components, each isomorphic to
\(\mathbb P^1_k\), meeting in exactly two points, which are ordinary double
points of \(C_2\). In either case, the enlarged fundamental group of the
curve is isomorphic to \(\mathbb Z\).\footnote{\textit{Editors' note.} This
could be spelled out. First, by \SGARefLink{sga3:x:corollary:5-14}{5.14}, the enlarged fundamental group of the
projective line \(\mathbb P^1_k\) is zero; that is,
\(\mathbb P^1_k\) is ``truly'' simply connected. One would then have to
extend Corollary 5.4 of SGA 1, IX, and the discussion following it, to the
enlarged fundamental group and to the category of principal bundles not
necessarily finite. Let \(\Gamma_1,\Gamma_2\) be two copies of
\(\mathbb P^1_k\), and let \(a_i,b_i\) be two distinct points of
\(\Gamma_i\). Let \(C''_2\) be the disjoint union of \(\Gamma_1\) and
\(\Gamma_2\), and let \(C'_2\) be the curve obtained by identifying
\(a_1\) with \(a_2\); then \(C_2\) is obtained from \(C'_2\) by also
identifying \(b_1\) with \(b_2\). The discussion following
\emph{loc.\ cit.}, extended to the enlarged case, shows that the enlarged
fundamental group of \(C'_2\), respectively \(C_2\), is zero, respectively
\(\mathbb Z\).}

More generally, one should take up, simplifying and correcting them, the
results of SGA 1, IX 5 in the setting of the enlarged fundamental group.
The examples of \emph{loc.\ cit.}, 5.5 would give as many examples of
enlarged fundamental pro-groups which are not profinite. Thus, if in the
first example the ordinary double point were replaced by a double point
with \(n\) distinct branches,\footnote{\textit{Editors' note.} This means
an \(n\)-fold point with \(n\) distinct tangents, for example the curve
\(X^n-Y^n=X^{n+1}\).} the enlarged fundamental group would be the free
(discrete!) group on \(n-1\) generators.

% END INCLUDED SGA3_English_current_buildable_source_20260823/source/inputs/X/tex/components/02_expose_X_sections4_6.tex

% BEGIN INCLUDED SGA3_English_current_buildable_source_20260823/source/inputs/X/tex/components/03_expose_X_sections7_9_bibliography.tex
\section*{\SGAInlineRefTarget{sga3:x:section:7}7. Classification of twisted constant preschemes and of groups of
multiplicative type of finite type in terms of the enlarged fundamental group}
\addcontentsline{toc}{section}{7. Classification in terms of the enlarged
fundamental group}

\setcounter{footnote}{57}

\subsection*{\SGAInlineRefTarget{sga3:x:subsection:7-0}7.0.}

Let \(S\) be a prescheme, which we continue to assume locally noetherian, in
order to ensure that \(S\), and certain preschemes over \(S\) that we shall
consider (notably those étale over \(S\), and more generally those locally of
finite type over \(S\)), are locally connected.

\par\medskip\noindent\SGAInlineRefTarget{sga3:x:proposition:7-0-1}\textbf{Proposition 7.0.1.---}\enspace
\emph{Every twisted constant prescheme \(X\) over \(S\) that is locally
trivial for the \((\mathrm{fppf})\) topology (that is, that is split by a
faithfully flat morphism \(S'\to S\) locally of finite presentation) is
quasi-isotrivial (that is, one may even choose \(S'\to S\) étale and
surjective).}

Indeed, we may suppose that \(S\) is connected, and hence that \(X\) is of
type \(I\), where \(I\) is a fixed set. Thus
\[
  X'=X\times_S S'
\]
is isomorphic to \(I_{S'}\), so \(I_{S'}\) is endowed with descent data
relative to \(S'\to S\); in other words, there is an isomorphism
\[
  I_{S''}\xrightarrow{\sim}I_{S''}
\]
satisfying the usual transitivity condition. Now
\(S''=S'\times_S S'\) is locally noetherian and hence locally connected.
It follows that the automorphisms of \(I_{S''}\) correspond to the sections
of \(G_{S''}\), where \(G=\operatorname{Aut}(I)\) is the group of
permutations of \(I\).

In this way one obtains descent data on \(G_{S'}\) (regarded as a trivial
Galois principal bundle) relative to \(S'\to S\). By \SGARefLink{sga3:x:lemma:5-4}{5.4} these descent data
are effective, and therefore give a Galois principal bundle \(P\) over \(S\)
with group \(G\). By construction, \(P\) represents the functor
\(\operatorname{Isom}_S(I_S,X)\) in the category of preschemes over \(S\)
that are locally noetherian. Consequently, the étale surjective base change
\(P\to S\) splits \(X\), so \(X\) is indeed quasi-isotrivial.

\par\medskip\noindent\SGAInlineRefTarget{sga3:x:remark:7-0-2}\textbf{Remark 7.0.2.---}\enspace
One should bear in mind that, even when \(S\) is the spectrum of a field, it
is not true in general that every twisted constant bundle over \(S\) is
quasi-isotrivial. For example, one may take \(X\) to be the disjoint union
of a sequence of schemes of the form \(\operatorname{Spec}(k_i)\), where
the \(k_i\) are separable extensions of \(k\) of strictly increasing
degrees.

The proof just given shows at the same time that the classification of
quasi-isotrivial twisted constant bundles \(X\) over \(S\) of type \(I\) is
equivalent to that of Galois principal bundles over \(S\) with group
\(G=\operatorname{Aut}(I)\). Indeed, this is an equivalence of categories.

It can be put into a more convenient form, as in SGA~1, V. To this end,
suppose that \(S\) is connected and equipped with a geometric point
\(\xi\). The enlarged fundamental pro-group
\(\Pi=\Pi_1(S,\xi)\) is then defined. For every quasi-isotrivial twisted
constant bundle \(X\) over \(S\), let \(I=X(\xi)\) be its set-theoretic
fiber at \(\xi\). Thus \(X\) is of type \(I\), and is consequently
associated, as above, with a Galois principal bundle
\[
  P=\operatorname{Isom}_S(I_S,X)
\]
over \(S\), with group \(G=\operatorname{Aut}(I)\).

By the definition of \(\Pi\), one therefore obtains a canonical
homomorphism from \(\Pi\) to \(G\), that is, from one of the \(\Pi_i\) to
\(G\). Since \(G\) is the group of permutations of \(I=X(\xi)\), this says
that \(\Pi\) ``acts continuously on \(I=X(\xi)\)'', meaning that the
\(\pi_i\), for \(i\) sufficiently large, act on \(I\) compatibly with the
transition morphisms.

We leave it to the reader to verify that every \(S\)-morphism \(X\to Y\)
between quasi-isotrivial twisted constant bundles over \(S\) induces a map
\(X(\xi)\to Y(\xi)\) compatible with the actions of \(\Pi\), and that the
resulting functor is an equivalence of categories:

\par\medskip\noindent\SGAInlineRefTarget{sga3:x:proposition:7-0-3}\textbf{Proposition 7.0.3.---}\enspace
\emph{Let \(S\) be a locally noetherian connected prescheme, let \(\xi\) be
a geometric point of \(S\), and let
\(\Pi=\Pi_1(S,\xi)\) be the enlarged fundamental pro-group of \(S\) at
\(\xi\). Then the functor}
\[
  X\longmapsto X(\xi)
\]
\emph{is an equivalence between the category of quasi-isotrivial twisted
constant bundles over \(S\) and the category of sets on which \(\Pi\) acts
continuously.}

This functor is compatible with finite sums and finite inverse limits. It
follows, for example, that quasi-isotrivial twisted constant groups (or
rings, and so forth) over \(S\) correspond to ordinary groups
(respectively rings, and so forth) on which the pro-group \(\Pi\) acts
continuously. In particular:

\par\medskip\noindent\SGAInlineRefTarget{sga3:x:corollary:7-0-4}\textbf{Corollary 7.0.4.---}\enspace
\emph{The category of quasi-isotrivial twisted constant commutative groups
over \(S\) is equivalent to the category of ``\(\Pi\)-modules'', that is,
commutative groups \(M\) on which \(\Pi\) acts continuously.}

Using \SGARefLink{sga3:x:corollary:5-7}{5.7}, we now obtain:

\par\medskip\noindent\SGAInlineRefTarget{sga3:x:theorem:7-1}\textbf{Theorem 7.1.---}\enspace
\emph{Let \(S\) be a locally noetherian connected prescheme, let \(\xi\) be
a geometric point of \(S\), and let
\(\Pi=\Pi_1(S,\xi)\) be the enlarged fundamental pro-group of \(S\) at
\(\xi\). Then the functor}
\[
  G\longmapsto
  \operatorname{Hom}_{\kappa(\xi)\text{-}\mathrm{gr}}
    (G_\xi,\mathbf{G}_{m,\xi})
\]
\emph{induces an anti-equivalence between the category of quasi-isotrivial
groups of multiplicative type over \(S\) and the category of
\(\Pi\)-modules.}

Using \SGARefLink{sga3:x:corollary:4-5}{4.5}, we obtain:

\par\medskip\noindent\SGAInlineRefTarget{sga3:x:corollary:7-2}\textbf{Corollary 7.2.---}\enspace
\emph{The preceding functor induces an anti-equivalence between the
category of groups of multiplicative type and of finite type over \(S\),
and the category of \(\Pi\)-modules that are of finite type over
\(\mathbb{Z}\).}

\par\medskip\noindent\SGAInlineRefTarget{sga3:x:example:7-3}\textbf{Example 7.3.---}\enspace
Take for \(S\), for example, a complete rational curve over an
algebraically closed field, having exactly one multiple point with
\(n+1\) distinct branches. By \SGARefLink{sga3:x:example:6-4}{6.4}, the enlarged fundamental group
\(\Pi(S,\xi)\) is a free group on \(n\) generators. Thus, by \SGARefLink{sga3:x:corollary:7-2}{7.2}, the
classification of tori of relative dimension \(m\) over \(S\) is
equivalent to the classification of systems of \(n\) endomorphisms
\(A_1,\ldots,A_n\) of the \(\mathbb{Z}\)-module \(\mathbb{Z}^m\), up to
automorphism of \(\mathbb{Z}^m\). Except when \(n\leq 1\) or \(m\leq 1\),
an explicit classification of such systems seems hopeless. One can at
least define a multitude of nontrivial invariants for such a system, such
as the characteristic polynomials of the \(A_i\).\footnote{\textit{Editors'
note.} In particular, this shows that the tori of relative dimension \(2\)
considered in \SGARefLink{sga3:x:paragraph:1:6}{1.6} are non-isotrivial.}

\par\medskip\noindent\SGAInlineRefTarget{sga3:x:remark:7-4}\textbf{Remark 7.4.---}\enspace
Without any hypothesis on \(S\), it remains true that, for a given ordinary
commutative group \(M\) of finite type over \(\mathbb{Z}\), the category of
groups of multiplicative type of type \(M\) over \(S\) is anti-equivalent
to the category of Galois principal bundles over \(S\) with group
\(G=\operatorname{Aut}_{\mathrm{gr}}(M)\). This follows easily from \SGARefLink{sga3:x:corollary:5-9}{5.9}
and \SGARefLink{sga3:x:corollary:5-10}{5.10}.

\par\medskip\noindent\SGAInlineRefTarget{sga3:x:remark:7-5}\textbf{Remark 7.5.---}\enspace
The theory of the fundamental pro-group sketched in the present two
sections is more advantageously formulated in the setting of general
sites. In that form it applies equally well, for example, to ordinary
topological spaces, and gives a satisfactory theory at least for a locally
connected topological space (not necessarily locally simply connected).
Here too it seems that one cannot make do with a fundamental group, and
that a pro-group is needed. Finally, once one has at one's disposal the
language of topologies and descent (which really lies at the heart of these
questions), the account sketched here is also technically simpler than
that of SGA~1, V, and should therefore in principle replace it.

\section*{\SGAInlineRefTarget{sga3:x:section:8}8. Appendix: Elimination of certain affineness hypotheses}
\addcontentsline{toc}{section}{8. Appendix: Elimination of certain affineness
hypotheses}

Our aim is to prove the following generalization of 4.9.

\par\medskip\noindent\SGAInlineRefTarget{sga3:x:theorem:8-1}\textbf{Theorem 8.1.---}\enspace
\emph{Let \(S\) be a prescheme, and let \(H\) be a flat \(S\)-prescheme in
groups of finite presentation, with connected affine
fibers.\footnote{\textit{Editors' note.} Observe that the hypotheses imply
that \(H\) is separated over \(S\), by VI\textsubscript{B}, \SGARefLink{thm:sga3-VIB-5.3}{Theorem~5.3} or
\SGARefLink{sga3:x:corollary:5-5}{Corollary~5.5}.} Let \(s\in S\), and suppose that \(H_s\) is a torus. Then
there exists an open neighborhood \(U\) of \(s\) such that \(H|_U\) is a
torus.}

We immediately obtain:

\par\medskip\noindent\SGAInlineRefTarget{sga3:x:corollary:8-2}\textbf{Corollary 8.2.---}\enspace
\emph{Let \(H\) be an \(S\)-prescheme in groups, flat and of finite
presentation over \(S\). Then \(H\) is a torus if and only if its fibers
are tori.}

\par\medskip\noindent\SGAInlineRefTarget{sga3:x:remark:8-3}\textbf{Remark 8.3.---}\enspace
Even when \(S\) is the spectrum of a discrete valuation ring, the
hypothesis in 8.1 that the fibers of \(H\) (here, the generic fiber) be
affine cannot be dropped: there are examples of smooth groups over \(S\)
whose generic fiber is an elliptic curve and whose special fiber is
\(\mathbf{G}_m\).

\medskip
\noindent\emph{Proof of 8.1.}
We may clearly suppose that \(S=\operatorname{Spec}(A)\) is affine. By the
standard procedure (cf.\ EGA~IV\textsubscript{2}, 8.8.2), this reduces us
to the case in which \(S\) is moreover noetherian. We begin by proving 8.2
in this case.

By 4.9, we are reduced to proving that \(H\) is affine over \(S\). We may
therefore suppose\footnote{\textit{Editors' note.} By
EGA~IV\textsubscript{2}, 8.8.2 once again.} that \(A\) is noetherian and
local. Since its completion \(\widehat A\) is faithfully flat over \(A\),
descent reduces us to the case in which \(A\) is a complete noetherian
local ring. Let \(S'\) be the normalization of \(S_{\mathrm{red}}\).
Nagata's theorem says that it is finite over \(S\)
(EGA~\(0_{\mathrm{IV}}\), 23.1.5). Moreover, \(S'\to S\) is surjective,
so
\[
  H'=H\times_S S'\longrightarrow H
\]
is finite and surjective. Hence, to prove that \(H\) is affine, it suffices
to show that \(H'\) is affine (EGA~II, 6.7.1).

Replacing \(S\) by a connected component of \(S'\), we are reduced to the
case in which \(A\) is a noetherian (semilocal), normal, integral
ring.\footnote{\textit{Editors' note.} The original has been expanded in
what follows.} Moreover, after possibly replacing \(A\) by its
normalization in a finite separable extension of its field of fractions,
we may suppose that the generic fiber \(H_\eta\) of \(H\) is
diagonalizable, that is, that there is an isomorphism
\[
  u_\eta:H_\eta\xrightarrow{\sim}T_\eta,
  \qquad T=\mathbf{G}_{m,S}^{\,r}.
\]
We now have:

\par\medskip\noindent\SGAInlineRefTarget{sga3:x:lemma:8-4}\textbf{Lemma 8.4.---}\enspace
\emph{Let \(S\) be a locally noetherian normal irreducible prescheme with
generic point \(\eta\). Let \(H\) be a smooth prescheme in groups over
\(S\), with connected fibers, let \(T\) be a prescheme in groups of
multiplicative type and of finite type over \(S\), and let
\(u_\eta:H_\eta\to T_\eta\) be a homomorphism of algebraic groups over
\(\kappa(\eta)\). Then \(u_\eta\) extends to a homomorphism of groups
\(u:H\to T\).}

After possibly replacing \(S\) by its normalization in a finite separable
extension of its function field, we may suppose that \(T\) is
diagonalizable (which is already the case in the application we have in
mind). Then \(T\) is a closed subgroup of a group of the form
\(\mathbf{G}_{m,S}^{\,r}\), reducing us to the case
\(T=\mathbf{G}_{m,S}\).

It remains only to prove that \(u_\eta\), regarded as a rational map from
\(H\) to \(T=\mathbf{G}_{m,S}\), is everywhere defined (for then the
morphism \(u:H\to T\) extending it is necessarily a group homomorphism).
We may regard \(u_\eta\) as an invertible section \(f\) of the structure
sheaf of \(H_\eta\), and must show that it extends to an invertible section
of the structure sheaf of \(H\). Since \(H\) is smooth over the normal
prescheme \(S\), it is normal (SGA~1, II 3.1). It therefore suffices to
find a closed subset \(Z\) of \(H\), of codimension at least \(2\), such
that \(f\) extends to an invertible section of the structure sheaf of
\(H-Z\). By localizing at the codimension-one points of \(S\), this at
once reduces us to the case in which \(S\) is the spectrum of a discrete
valuation ring \(A\).

Let \(t\) be a uniformizer of \(A\), and let \(t'\) be the section of
\(\mathcal O_H\) that it defines, so that the special fiber \(H_0\) equals
\(V(t')\). By hypothesis, \(H_0\) is smooth and connected over the residue
field \(k\). The function \(f\) is then a rational function on \(H\) with
neither zeros nor poles in \(H-H_0\). Since \(H_0=V(t')\) is an
irreducible divisor, there is an integer \(n\in\mathbb{Z}\) such that
\[
  t'^{\,n}f=t^n f
\]
has neither zeros nor poles, and hence is an invertible section of
\(\mathcal O_H\). It therefore defines a morphism
\(v:H\to\mathbf{G}_{m,S}\). Since
\(v_\eta=t'^{\,n}u_\eta\), and \(u_\eta\) is a group homomorphism, \(v\)
sends the identity section of \(H\) to a section of
\(\mathbf{G}_{m,S}\) whose value at the generic point of \(S\) is \(t^n\).
Because this is a section of \(\mathbf{G}_{m,S}\), \(t^n\) must be a unit;
thus \(n=0\). Hence \(v\) extends \(u_\eta\), completing the proof of \SGARefLink{sga3:x:lemma:8-4}{8.4}.

Applying this lemma in the present situation, we obtain a group
homomorphism
\[
  u:H\longrightarrow T=\mathbf{G}_{m,S}^{\,r}
\]
that induces an isomorphism on the generic fibers. We shall prove that
\(u\) is an isomorphism.

\par\medskip\noindent\SGAInlineRefTarget{sga3:x:lemma:8-5}\textbf{Lemma 8.5.---}\enspace
\emph{For every integer \(n>0\) prime to the residual characteristic of
\(A\), the group
\[
  {}_nH=\ker(n\cdot\operatorname{id}_H)
\]
is finite over \(S\).}

If \(n\) is prime to the residual characteristic of \(A\), then it is
prime to all the residual characteristics of the points of \(S\).
Consequently, \(n\cdot\operatorname{id}_H\) induces an étale morphism on
every fiber of \(H\), and hence is itself étale (SGA~1, I 5.9); its kernel
\({}_nH\) is therefore étale over \(S\). On the other hand, \({}_nH\) is
separated over \(S\), because \(H\) is.
\begingroup
\renewcommand{\thefootnote}{\(\ast\)}
\footnote{By Raynaud's theorem, VI\textsubscript{B}, 5.3.}
\endgroup
\addtocounter{footnote}{-1}%
\footnote{\textit{Editors' note.} This point should be expanded in
connection with the modifications in VI\textsubscript{B}, \S5.}
Moreover, all its fibers have the same rank \(n^r\), since the fibers of
\(H\) are tori, all of the same dimension \(r\) (as \(H\) is smooth over
\(S\)). We conclude that \({}_nH\) is finite over \(S\) (SGA~1, I 10.9,
or EGA~IV\textsubscript{4}, 18.2.9).

By the preceding lemma, \(u({}_nH)\) is therefore a closed subset of
\(T\). On the generic fiber \(T_\eta\), this subset is identical to
\({}_n(T_\eta)\), so it contains the closure of that subset, namely
\({}_nT\). For every \(s\in S\), the groups \({}_n(T_s)\), as \(n\)
varies, are dense in \(T_s\); since \(u_s(H_s)\) is a closed subset
(VI\textsubscript{B}, 1.4.2) containing them, we have
\(u_s(H_s)=T_s\). Thus \(u\) is surjective. A surjective homomorphism
between tori of the same dimension over a field is
flat,\footnote{\textit{Editors' note.} This follows, for example, from
VIII 3.2(a). More generally, if \(f:G\to H\) is a surjective morphism of
algebraic groups over a field \(k\), and \(H\) is reduced, then \(f\) is
flat (cf.\ VI\textsubscript{B}, 1.3).} and hence \(u\) induces a flat
morphism on every fiber. Therefore \(u\) is flat (SGA~1, I 5.9).
Consequently \(K=\ker(u)\) is flat over
\(S\),\footnote{\textit{Editors' note.} The original has been expanded in
what follows.} and hence equals the closure of its generic fiber
\(K_\eta\). By construction, \(K_\eta\) is the identity group. Since
\(K\) is separated over \(S\) (because \(H\) is), its identity section is
closed, and it follows that \(K\) is the identity group. Thus \(u\) is a
monomorphism; since we have seen that it is faithfully flat, it is an
isomorphism (cf.\ SGA~1, I 5.1, or EGA~IV\textsubscript{4}, 17.9.1).
This proves that \(H\) is a torus, and completes the proof of
8.2.\footnote{\textit{Editors' note.} At least in the case considered up
to this point, namely \(S\) locally noetherian.}

\par\medskip\noindent\SGAInlineRefTarget{sga3:x:remark:8-5-1}\textbf{Remark 8.5.1.---}\enspace
Instead of invoking 8.5, one can also invoke Zariski's ``Main Theorem'',
which directly implies that \(u\) is an open immersion and hence an
isomorphism.\footnote{\textit{Editors' note.} The editors did not
understand this remark, since they did not see why \(u\) would
\emph{a priori} have finite and surjective fibers.}

To prove \SGARefLink{sga3:x:theorem:8-1}{8.1}, the quasi-isotriviality theorem reduces us to the case in
which \(S\) is local and \(s\) is its closed
point,\footnote{\textit{Editors' note.} Expand this reduction.} and it
remains to prove that, under the other hypotheses, \(H\) is then a torus.
By \SGARefLink{sga3:x:corollary:8-2}{8.2}, already proved, we are reduced to showing that the fibers of \(H\)
are tori. We may suppose that \(S\) is the spectrum of a complete
noetherian local ring \(A\). By \SGARefLink{sga3:x:corollary:3:3}{3.3}, for every \(n>0\) there exists a
group of multiplicative type \(T_n\), finite over \(S\), and an
isomorphism
\[
  (T_n)_s\simeq {}_n(H_s),
\]
where \(s\) is the closed point of \(S\). Arguing as in \SGARefLink{sga3:x:lemma:3:1}{3.1} and using the
fact that \(T_n\) is finite over
\(S\),\footnote{\textit{Editors' note.} Expand this point.} we see that
the preceding isomorphism comes from a homomorphism
\(u_n:T_n\to H\), which is moreover uniquely determined. (Pass to the
limit over the artinian quotients of \(A\).)

Moreover, by the uniqueness properties, when \(m\) is a multiple of \(n\),
the homomorphism \(u_n:T_n\to H\) is obtained from
\(u_m:T_m\to H\) by restriction to
\({}_n(T_m)\simeq T_n\). It follows from \SGARefLink{sga3:ix:corollary:6:6}{IX 6.6} that the \(u_n\) are
monomorphisms, so the \(T_n\) are subgroups of \(H\), and for \(m\) a
multiple of \(n\) one has
\[
  T_n={}_n(T_m).
\]

Thus, for every \(t\in S\), the groups \((T_n)_t\) are subgroups of
\(H_t\), of type \((\mathbb{Z}/n\mathbb{Z})^r\), where
\(r=\dim H_s=\dim H_t\), and when \(m\) is a multiple of \(n\) one has
\[
  (T_n)_t={}_n(T_m)_t.
\]
The fact that \(H_t\) is a torus now follows from:

\par\medskip\noindent\SGAInlineRefTarget{sga3:x:lemma:8-6}\textbf{Lemma 8.6.---}\enspace
\emph{Let \(H\) be a smooth affine algebraic group over an algebraically
closed field \(k\), and let \((T_n)_{n>0}\) be a family of subgroups of
multiplicative type of \(H\), such that for every integer \(n>0\), \(T_n\)
is of type \((\mathbb{Z}/n\mathbb{Z})^r\), and for every multiple \(m\) of
\(n\), one has \(T_n={}_n(T_m)\). Under these conditions, \(H\) contains
a torus of dimension at least \(r\) containing the \(T_n\). Consequently,
if \(H\) is connected and has dimension at most \(r\), then \(H\) is a
torus of dimension \(r\).}

This is an exercise on affine algebraic groups, which we shall handle by
repeated reference to \emph{Bible}. We restrict attention to the \(T_n\)
for \(n\) prime to the characteristic. Let \(K\) be the closure of the
union of the \(T_n\) in \(H\), endowed with the induced reduced structure.
Standard arguments then show that \(K\) is a commutative algebraic
subgroup of \(H\). By \emph{Bible}, \S4.5, Theorem~4, \(K\) is isomorphic
to a product \(K_u\times K_s\), with \(K_u\) ``unipotent'' and \(K_s\)
diagonalizable. Every diagonalizable subgroup of \(K\) is contained in
\(K_s\), so the \(T_n\) are subgroups of \(K_s\), and therefore
\(K=K_s\). Write \(K=D(M)\), where \(M\) is an ordinary commutative group
of finite type over \(\mathbb{Z}\). The inclusion \(T_n\subseteq K\)
means that \(M\) has a quotient group isomorphic to
\((\mathbb{Z}/n\mathbb{Z})^r\). Since this holds for every integer \(n\)
prime to the characteristic of \(k\) (powers of a single fixed prime
would suffice), \(M\) has rank at least \(r\). Thus \(K\) contains a torus
\(T\) of dimension \(r\). When \(H\) is connected of dimension \(r\), it
follows that \(T=H\), completing the proof of \SGARefLink{sga3:x:lemma:8-6}{8.6}. \SGARefLink{sga3:x:theorem:8-1}{Theorem~8.1} is now
proved.

\par\medskip\noindent\SGAInlineRefTarget{sga3:x:remark:8-7}\textbf{Remarks 8.7.---}\enspace
Using \SGARefLink{sga3:x:theorem:8-1}{8.1}, it should not be difficult to give analogous generalizations of
\SGARefLink{sga3:x:corollary:4-7}{4.7} and \SGARefLink{sga3:x:corollary:4-8}{4.8}. A more interesting problem would be to study the situation of
\SGARefLink{sga3:x:theorem:8-1}{8.1} after dropping the hypothesis that the fibers of \(H\) be affine. One
can show that there is then an open neighborhood \(U\) of \(s\) such that
\(H|_U\) is commutative and, for every \(t\in U\), the geometric fiber
\(H_t\) is an extension of an abelian variety by a
torus.\footnote{\textit{Editors' note.} Give a reference here?} Of course,
in questions of this kind one may restrict to the case in which \(S\) is
the spectrum of a discrete valuation ring, \(s\) its closed point, and
\(t\) its generic point.

This result can be generalized as follows. For every smooth connected
algebraic group \(G\) over an algebraically closed field \(k\), a
well-known theorem of Chevalley says that \(G\) is, in a unique way, an
extension of an abelian variety \(A\) by a smooth connected affine group
\(V\). Define the \emph{abelian rank} (respectively, \emph{reductive rank},
\emph{nilpotent rank}, and \emph{semisimple dimension}) of \(G\), denoted
\(\rho_{\mathrm{ab}}(G)\) (respectively,
\(\rho_r(G)\), \(\rho_n(G)\), and \(d_s(G)\)), to be the dimension of
\(A\) (respectively, the dimension of the maximal tori of \(G\), the
dimension of the Cartan subgroups\footnote{\textit{Editors' note.} Recall
the definition of ``Cartan subgroup''.} of \(G\), and the dimension of the
quotient of \(G\), or equivalently of \(V\), by its radical); see
\emph{Bible} for all these notions. Also define the \emph{unipotent rank}
by
\[
  \rho_u(G)=\rho_u(V)
  =\rho_n(G)-\rho_r(G)-\rho_{\mathrm{ab}}(G).
\]
When the ground field is not algebraically closed, the same names and
notations designate the corresponding invariants of \(G_{\bar{k}}\),
where \(\bar{k}\) is an algebraic closure of \(k\).

Now let \(G\) be a smooth group scheme over the spectrum \(S\) of a
discrete valuation ring. Let \(\rho_{\mathrm{ab}}\), and so on
(respectively, \(\rho'_{\mathrm{ab}}\), and so on), denote the invariants
associated with the special fiber (respectively, the generic fiber). Then
one has the inequalities
\[
  \left\{
  \begin{aligned}
    \rho_{\mathrm{ab}}&\leq \rho'_{\mathrm{ab}},\\
    \rho_r+\rho_{\mathrm{ab}}&\leq
      \rho'_r+\rho'_{\mathrm{ab}},\\
    d_s&\leq d'_s,
  \end{aligned}
  \right.
  \qquad
  \left\{
  \begin{aligned}
    \rho_n&\geq \rho'_n,\\
    \rho_u&\geq \rho'_u.
  \end{aligned}
  \right.
\]
Equivalently, if \(G\) is smooth of finite type over an arbitrary base
\(S\), the functions
\[
  s\longmapsto \rho_{\mathrm{ab}}(s),\qquad
  s\longmapsto \rho_{\mathrm{ab}}(s)+\rho_r(s),\qquad
  s\longmapsto d_s(s)
\]
are lower semicontinuous, while the functions
\[
  s\longmapsto \rho_n(s),\qquad
  s\longmapsto \rho_u(s)
\]
are upper semicontinuous.\footnote{\textit{Editors' note.} Give a
reference for these results?}

The same results probably remain valid without assuming \(G\) smooth over
\(S\), but assuming only that it is flat and of finite presentation over
\(S\), provided that, for an algebraic group \(G\) over an algebraically
closed field \(k\), \(\rho_{\mathrm{ab}}(G)\), and so on, are understood
to mean the corresponding invariants of \(G_{\mathrm{red}}^0\).

In this Seminar we present some results of this kind when \(G\) is affine
over \(S\), or more generally has affine fibers. In this case we shall
verify the semicontinuity properties for \(\rho_r\) and \(\rho_n\), hence
for \(\rho_u=\rho_n-\rho_r\), and the continuity of \(\rho_s\) in a
neighborhood of a point \(s\) whose fiber is a reductive
group.\footnotemark[71]

When \(G\) is already assumed commutative, \SGARefLink{sga3:x:corollary:8-2}{8.2} can be generalized as
follows.

\par\medskip\noindent\SGAInlineRefTarget{sga3:x:theorem:8-8}\textbf{Theorem 8.8.---}\enspace
\footnote{\textit{Editors' note.} G.~Prasad and J.-K.~Yu generalized this
result (subject to some additional hypotheses) without assuming \(G\)
commutative, and with ``torus'' replaced by ``reductive group'' in
hypothesis~(a) and in the conclusion; cf.\ \([{\rm PY06}]\),
Theorem~6.2.}%
\emph{Let \(G\) be a commutative \(S\)-prescheme in groups, flat and of
finite presentation over \(S\), with affine connected fibers. Let
\(s\in S\), and suppose that:
\begin{enumerate}
\renewcommand{\labelenumi}{\alph{enumi})}
\item if \(\bar{k}\) denotes an algebraic closure of \(\kappa(s)\), then
  \((G_{\bar{k}})_{\mathrm{red}}\) is a torus;
\item there is a generization \(t\) of \(s\) such that \(G_t\) is smooth
  over \(\kappa(t)\).
\end{enumerate}
Under these conditions, there exists an open neighborhood \(U\) of \(s\)
such that \(G|_U\) is a torus.}

\emph{(Observe moreover that, if one assumes only that \(G_t\) is affine,
respectively connected, for every generization \(t\) of \(s\), it follows
easily that the same conclusion holds for \(t\) in an open neighborhood
of \(s\).)}

\medskip
\noindent\emph{Proof of \SGARefLink{sga3:x:theorem:8-8}{8.8}.}
It suffices to prove that \(G_s\) is smooth over \(\kappa(s)\). Indeed,
since \(G\) is flat and of finite presentation over \(S\), it then follows
that \(G\) is smooth over \(S\) above a neighborhood of \(s\)
(cf.\ \SGARefLink{sga3:x:lemma:3:5}{3.5}), and we are then in the situation of \SGARefLink{sga3:x:theorem:8-1}{8.1}.

To prove that \(G_s\) is smooth over \(\kappa(s)\), the usual procedure
reduces us to the case in which \(S\) is affine and noetherian. Choose a
morphism \(S'\to S\), where \(S'\) is the spectrum of a discrete
valuation ring, sending the closed point to \(s\) and the generic point
\(\eta\) to \(t\). We are thereby reduced to the case in which \(S\) is
itself the spectrum of a discrete valuation ring; we may moreover suppose
that it is complete and has algebraically closed residue field, and that
\(s\) and \(t=\eta\) are respectively the closed point and the generic
point of \(S\).

Thus \(G\) is flat, separated, and of finite type over \(S\); the generic
fiber \(G_\eta\) is smooth and connected, and the special fiber \(G_0\)
is such that
\[
  T_0=(G_0)_{\mathrm{red}}
\]
is a torus. Let \(m>0\) be an integer prime to the residual
characteristic at \(s\), and hence also to that at \(\eta\). Then
\(m\cdot\operatorname{id}_G\) is étale on every fiber
(VII\textsubscript{A}, 8.4), and is therefore étale because \(G\) is flat
over \(S\) (SGA~1, I 5.9). Its kernel \({}_mG\) is thus étale over \(S\);
since \(G\) is separated
\begingroup
\renewcommand{\thefootnote}{\(\ast\)}
\footnote{By Raynaud's theorem, VI\textsubscript{B}, 5.3.}
\endgroup
\addtocounter{footnote}{-1}%
\footnote{\textit{Editors' note.} See Editors' note (62) in 8.5.}
and of finite type over \(S\), the same holds for \({}_mG\). Since
\[
  ({}_mG)_0={}_m(T_0),
\]
its degree is \(m^r\), where
\[
  r=\dim T_0=\dim G_0=\dim G.
\]
It follows that the rank of
\(({}_mG)_\eta={}_m(G_\eta)\) is at least \(m^r\), which already proves,
using \SGARefLink{sga3:x:lemma:8-6}{8.6}, that \(G_\eta\) is a torus of dimension \(r\), since the two
fibers of \({}_mG\) have the same rank, and hence, as in \SGARefLink{sga3:x:lemma:8-5}{8.5}, that
\({}_mG\) is finite over \(S\).\footnote{\textit{Editors' note.} Revise
the preceding sentence.}

Since \(S\) is complete with algebraically closed residue field, the
finite étale cover \({}_mG\) splits completely. Thus, through every point
of \({}_mG_0\) there passes a section of \(G\) over \(S\); in particular,
the set of points of \(G_0\) through which there passes a section of \(G\)
over \(S\) is everywhere dense. We now use the following lemma.

\par\medskip\noindent\SGAInlineRefTarget{sga3:x:lemma:8-9}\textbf{Lemma 8.9.---}\enspace
\emph{Let \(S\) be a locally noetherian regular prescheme of dimension
\(1\), and let \(G\) be an \(S\)-prescheme in groups, flat and locally of
finite type, such that \(G_\eta\) is smooth over \(\kappa(\eta)\) for
every maximal point \(\eta\) of \(S\) (which implies that \(G\) is
reduced). Suppose that the normalization \(X\) of \(G\) is finite over
\(G\) (this is the case, by Nagata, when \(S\) is the spectrum of a
complete discrete valuation ring; cf.\ EGA~IV\textsubscript{2}, 7.7.4),
and let \(G'\) be the open subset of \(X\) formed by the points at which
\(X\) is smooth over \(S\). With these notations:
\begin{enumerate}
\renewcommand{\labelenumi}{\alph{enumi})}
\item If the projection \(G'\to S\) is surjective, then \(G'\) carries a
  unique structure of an \(S\)-prescheme in groups for which the
  canonical morphism \(G'\to G\) is a group homomorphism.
\item Suppose that, for every closed point \(s\) of \(S\), the set of
  points of \(G_s^0\) through which an étale quasi-section passes is dense
  in \(G_s^0\) for the Zariski topology. Then \(X\) is regular, the
  hypotheses of \textup{(a)} hold, and the map
  \[
    \Gamma(G'/S)\longrightarrow\Gamma(G/S)
  \]
  is bijective.
\end{enumerate}}

In the case that concerns us, this lemma applies and gives a group
homomorphism \(u:G'\to G\), where \(G'\) is smooth over \(S\) and
\[
  u_\eta:G'_\eta\xrightarrow{\sim}G_\eta
\]
is an isomorphism. After possibly replacing \(G'\) by an open subgroup,
we may suppose that \(G'_0\) is connected. Since
\(u_0:G'_0\to G_0\) induces a surjective morphism
\(G'_0\to T_0\), where \(T_0\) is a torus of the same dimension as
\(G'_0\), it follows easily that \(G'_0\) is a torus (for example, by
using \SGARefLink{sga3:x:lemma:8-6}{8.6}, or by referring to \emph{Bible}, \S7.3, Theorem~3(a)).
Consequently, by \SGARefLink{sga3:x:corollary:8-2}{8.2}, \(G'\) is a torus. By \SGARefLink{sga3:ix:theorem:6:8}{IX 6.8},
\(\ker u\) is then a subgroup of multiplicative type of \(G'\); since its
generic fiber is the identity group, it is itself the identity group, and
therefore \(u\) is a monomorphism. It follows from \SGARefLink{sga3:viii:proposition:7:9}{VIII 7.9} that \(u\) is
an immersion. Since it is surjective and \(G\) is reduced, \(u\) is an
isomorphism, completing the proof of \SGARefLink{sga3:x:theorem:8-8}{8.8}.

It remains to prove \SGARefLink{sga3:x:lemma:8-9}{8.9}. For part \SGARefLink{sga3:viii:proposition:7:9:item:a}{(a)}, the uniqueness of the group law on
\(G'\) for which \(u:G'\to G\) is a group homomorphism is clear, since
the group law on the generic fiber of \(G'\) is known (we may suppose
\(S\) irreducible). For existence, we reduce easily to the case in which
\(S\) is local, the spectrum of a discrete valuation ring \(A\). By
uniqueness, and because integral closure commutes with an étale extension
of the base, we may make étale base extensions on \(S\). This reduces us
to the case in which \(A\) is ``strictly local'', that is, henselian with
separably closed residue field. The same reduction applies to part (b);
under the hypothesis made there, ``étale quasi-section'' can now be
replaced by ``section''.

Since \(G'\) is smooth over \(S\) and \(X\) is normal, the prescheme
\(G'\times_S X\) is normal (being smooth over the normal prescheme \(X\)).
Thus the composite morphism
\[
  G'\times_S X\longrightarrow G\times_S G\longrightarrow G
\]
factors as
\[
  p:G'\times_S X\longrightarrow X.
\]
We shall prove that the restriction of this morphism to the open subset
\(G'\times_S G'\) of \(G'\times_S X\) induces a morphism
\[
  G'\times_S G'\longrightarrow G'.
\]

It is therefore enough to show that \(G'_0\times_k G'_0\) maps into the
open subset \(G'_0\) of \(X_0\). It suffices to check that, for every
\(k\)-valued point \(g'_0\) of \(G'_0\), the morphism
\[
  \tag{\(\times\)}\label{sga3:x:lemma:8-9:equation:times}
  h'_0\longmapsto p_0(g'_0,h'_0)
\]
from \(G'_0\) to \(X_0\) takes its values in \(G'_0\). Since \(G'\) is
smooth over \(S\) and \(S\) is henselian, every such \(g'_0\) is induced
by a section \(g'\) of \(G'\) over \(S\). The morphism \((\times)\) is
then induced by the morphism
\[
  h'\longmapsto p(g',h')
\]
from \(G'\) to \(X\), which itself is obtained by transport of structure
from the automorphism \(h\mapsto g\cdot h\) of \(G\), namely left
translation by the section \(g\) of \(G\) that is the image of \(g'\).
Thus \(h'\mapsto p(g',h')\) is itself an automorphism of \(X\), and hence
maps \(G'\) into \(G'\). This proves the assertion.

It remains to prove that the composition law thereby obtained on \(G'\)
is a group law. Associativity follows at once from associativity on the
generic fiber (which is isomorphic to that of \(G\)). Moreover, the
inversion automorphism of the \(S\)-prescheme \(G\) induces an automorphism
of \(X\), which therefore preserves \(G'\) and induces an automorphism
\(\sigma\) of \(G'\). One verifies that \(\sigma\) has the properties of
an inverse for the composition law on \(G'\). Indeed, this again amounts
to verifying the commutativity of certain diagrams involving fiber
powers of \(G'\) over \(S\); since these are smooth over \(S\), it
suffices to check commutativity on the generic fiber, where it is clear.
This proves part (a) of 8.9.

We prove part (b). Let \(Z'\) be the set of \(x\in X\) for which
\(\mathcal O_{X,x}\) is not regular. It is closed by a theorem of Nagata
(EGA~IV\textsubscript{2}, 6.12.6), and is clearly contained in \(X_0\).
Let \(Z\) be its image in \(G\), which is therefore a closed subset of
\(G_0\). Then \(Z\) is nowhere dense in \(G_0\), that is, it contains no
maximal point \(y\) of \(G_0\). Indeed, \(G_0\) is defined in \(G\) by an
equation \(t=0\), where \(t\) is a uniformizer of the discrete valuation
ring \(A\). By the Hauptidealsatz, \(\mathcal O_{G,y}\) has dimension
\(1\). Hence, for every \(x\in X\) above \(y\),
\(\mathcal O_{X,x}\) has dimension \(1\), and is therefore a discrete
valuation ring, since \(X\) is normal and hence regular in codimension
\(1\).

On the other hand, for every section \(g\) of \(G\) over \(S\), the set
\(Z'\) is plainly stable under the automorphism of \(X\) obtained by
transport of structure from left translation by \(g\) on \(G\). Thus
\(Z\) is stable under the left translation on \(G_0\) defined by \(g_0\).
By hypothesis, the set of these \(g_0\) is dense in \(G_0^0\). Since
\(Z\) is stable under translation by these \(g_0\) and is a nowhere-dense
closed subset, it follows at once that \(Z=\varnothing\), whence
\(Z'=\varnothing\), so \(X\) is regular. Consequently \(X\) is smooth
over \(S\) at every point through which a section passes. Every section
of \(G\) over \(S\) lifts uniquely to a section of \(X\) over \(S\), and
hence to a section of \(G'\) over \(S\). This applies in particular to the
identity section, proving that the image of \(G'\) in \(S\) is all of
\(S\), that is, that the hypothesis of (a) holds. This completes the
proof of \SGARefLink{sga3:x:lemma:8-9}{8.9}, and hence of \SGARefLink{sga3:x:theorem:8-8}{8.8}.

\section*{\SGAInlineRefTarget{sga3:x:section:9}9. Addenda}
\addcontentsline{toc}{section}{9. Addenda}
\setcounter{section}{9}
\setcounter{subsection}{0}
\footnote{\textit{Editors' note.} This additional section was written by
Fabrice Orgogozo (following indications by Ofer Gabber).}

\subsection{Simplicial sets, topoi, groupoids, and topological spaces}

\par\medskip\noindent\SGAInlineRefTarget{sga3:x:notation:9-1-1}\textbf{Notations 9.1.1.---}\enspace
\emph{(1) Let \(E\) be a simplicial set. The following objects can be
associated with it:
\begin{enumerate}
\renewcommand{\labelenumi}{(\arabic{enumi})}
\setcounter{enumi}{1}
\item a topos \(T=E^{\sim}\), obtained from the topoi naturally
  associated with the sets \(E_i\), by the procedure described in
  \([{\rm Del74}]\), 6.3.1 (see also \([{\rm Ill72}]\), VI.5.2 and
  SGA~4, VI.7);
\item a groupoid \(G=\Pi(E)\), whose objects are the elements of \(E_0\)
  (the ``vertices'') and whose arrows are defined in \([{\rm GZ67}]\),
  II.7;
\item a topological space \(X=|E|\), a cellular complex called the
  ``geometric realization'' (or ``topological realization'')
  (cf.\ \emph{loc.\ cit.}, III.1).
\end{enumerate}}

Observe that a sheaf on \(T\) is nothing other than a simplicial set over
\(E\).

\subsection{Locally constant sheaves; descent data}

\par\medskip\noindent\SGAInlineRefTarget{sga3:x:definition:9-2-1}\textbf{Definitions 9.2.1.---}\enspace
A \emph{locally constant sheaf} on
\begin{enumerate}
\renewcommand{\labelenumi}{(\arabic{enumi})}
\item a simplicial set \(E\) is a morphism of simplicial sets
  \(E'\to E\) such that, for every \(i\in\mathbb{N}\) and every
  \(e\in E_i\), the face operators \(d\) induce isomorphisms
  \[
    E'_e\xrightarrow{\sim}E'_{d(e)}
  \]
  between the fibers (cf.\ \([{\rm AM69}]\), \S10);
\item a topos \(T\) is an object \(F\) of \(T\) such that there exist an
  epimorphism \(U\to 1\) and an isomorphism
  \[
    F\times U\simeq f^*I\times U,
  \]
  where \(I\) is a set and \(f:T\to\mathbf{Ens}\) is the final morphism
  (cf.\ SGA~4, IX.2);
\item a groupoid \(G\) is a presheaf on \(G\), that is, a contravariant
  functor from \(G\) to the category of sets
  (cf.\ \([{\rm GZ67}]\), Appendix I.1.2);
\item a topological space \(X\) is a sheaf of sets on \(X\) that is
  locally constant in the usual sense.
\end{enumerate}
Finally:
\begin{enumerate}
\renewcommand{\labelenumi}{(\arabic{enumi})}
\setcounter{enumi}{4}
\item a \emph{descent datum} on a simplicial set \(E\) consists of a
  sheaf \(F\) on \(E_0^{\sim}\) (that is, a set-valued function on
  \(E_0\)) and an isomorphism
  \[
    \alpha:p_1^*F\xrightarrow{\sim}p_2^*F,
  \]
  where \(p_1,p_2:E_1^{\sim}\to E_0^{\sim}\) are the morphisms induced by
  the faces, satisfying the usual cocycle relation
  (cf.\ \SGARefLink{sga3:iv:definition:2:1}{Exposé~IV, 2.1}%
  \begingroup
  \renewcommand{\thefootnote}{1}
  \footnote{PP: I have replaced ``FGA, \emph{Technique de descente} I,
  \SGARefLink{sga3:i:subsection:1-6}{1.6}'' by ``\SGARefLink{sga3:iv:definition:2:1}{Exposé IV, 2.1}''.}
  \endgroup
  \addtocounter{footnote}{-1}%
  \ and below).
\end{enumerate}

The morphisms between these five kinds of objects, as well as the
associated inverse-image functors, are defined in the evident way.

\subsection{Some equivalences of categories}

Let \(E\) be a simplicial set, and let \(T\), \(G\), and \(X\) be,
respectively, the associated topos, groupoid, and topological space.

\par\medskip\noindent\SGAInlineRefTarget{sga3:x:proposition:9-3-1}\textbf{Proposition 9.3.1.---}\enspace
\emph{The categories of locally constant sheaves on \(E\), \(T\), \(G\),
and \(X\), as well as the category of descent data on \(E\), are
equivalent.}

\medskip
\noindent\emph{Sketch of proof.}
Denote by \SGARefLink{sga3:x:proposition:9-3-1:item:1:03}{(1)} through \SGARefLink{sga3:x:proposition:9-3-1:item:5}{(5)} the categories of objects defined in the
preceding paragraph.

\begin{itemize}
\item[\((1)\Longleftrightarrow(5)\).]
  This is a special case of \([{\rm AM69}]\), 10.6 (see also
  \([{\rm GZ67}]\), Appendix \SGARefLink{sga3:x:corollary:2:3}{I.2.3}, \([{\rm Fri82}]\), ?.5.6, or
  \([{\rm Ill72}]\), VI.8.1.6). An equivalence of categories is given by
  the functor that associates with the object \(E'\to E\) the pair
  \((E'_0,\alpha)\), where \(E'_0\) is regarded as a sheaf on \(E_0\)
  and \(\alpha\) is the unique isomorphism
  \[
    p_1^*E'_0\xrightarrow{\sim}p_2^*E'_0
  \]
  whose fiber at each \(y\in E_1\), with images denoted \(x_1\) and
  \(x_2\) under the two projections, is given by the isomorphisms
  \[
    (E'_0)_{x_1}\xleftarrow{\sim}(E'_1)_y
      \xrightarrow{\sim}(E'_0)_{x_2}.
  \]
\item[\(\SGAInlineRefTarget{sga3:x:proposition:9-3-1:item:5}(5)\Longleftrightarrow(3)\).]
  This is evident: one of the two relations defining the morphisms in the
  groupoid \(G\) associated with \(E\) is a cocycle relation.
\item[\((1)\Longleftrightarrow(4)\).]
  Cf.\ \([{\rm GZ67}]\), Appendix I.3.2.1.
\item[\(\SGAInlineRefTarget{sga3:x:proposition:9-3-1:item:1:03}(1)\Longleftrightarrow(2)\).]
  One must show that an object over \(E\) is a locally constant sheaf in
  the simplicial sense if and only if it is one as a sheaf on
  \(T=E^{\sim}\). This follows from the fact that the locally constant
  sheaves \(F_\bullet\) on \(T\) are precisely the cartesian sheaves,
  that is, those for which the arrows
  \[
    ([n]\to[m])^*F_m\longrightarrow F_n
  \]
  are isomorphisms. This last assertion is a special case of a general
  fact about simplicial topoi, together with the fact that every sheaf on
  \(E_0^{\sim}\) is locally constant. (See also \([{\rm Ill72}]\),
  VI.8.1.6.)
\end{itemize}
\hfill Q.E.D.

For every group \(H\), the preceding equivalences of categories induce
equivalences between the categories of \(H\)-torsors. Here these are
locally constant sheaves equipped with an action of \(H\), whose fibers
are isomorphic to \(H\) acting on itself by translation.

\subsection{Fundamental groups and groupoids}

It follows from the preceding equivalences of categories that, for every
group \(H\) and every simplicial set \(E\), the sets of isomorphism
classes of \(H\)-torsors
\[
  H^1(E,H),\qquad
  H^1(E^{\sim},H),\qquad
  \pi_0\bigl(\operatorname{Hom}(\Pi(E),BH)\bigr),\qquad
  H^1(|E|,H)
\]
are naturally in bijection. Recall that \(BH\) denotes the one-object
groupoid associated with \(H\), and that \(\pi_0\) is the functor that
associates with a category the set of isomorphism classes of its objects.

Likewise, if \(e\) is a point of \(E\), the preceding equivalences induce
bijections between the sets of isomorphism classes of \(H\)-torsors
trivialized over \(e\), denoted respectively by
\[
  H^1(E\mathbin{\mathrm{rel}}e,H),\qquad
  H^1(E^{\sim}\mathbin{\mathrm{rel}}e^{\sim},H),\qquad
  \operatorname{Hom}\bigl(\pi_1(\Pi(E),e),H\bigr),\qquad
  H^1(|E|\mathbin{\mathrm{rel}}|e|,H).
\]
Recall that
\[
  \pi_1(\Pi(E),e)=\operatorname{Isom}_{\Pi(E)}(e,e).
\]
For variable \(H\), these functors are represented, in the connected
case, by a group denoted \(\pi_1(E,e)\). The group \(\pi_1(E,e)\) is
isomorphic to \(\pi_1(\Pi(E),e)\) and to \(\pi_1(|E|,|e|)\), and hence
also to the fundamental group of a simplicial set as defined by Kan
(cf., for example, \([{\rm May67}]\), 16.1, or \([{\rm Ill72}]\),
I.2.1.1). Recall also that \(H^1(E,H)\) is in turn isomorphic to the set
\(H^1(\pi_1(E,e),H)\) of homomorphisms to \(H\) modulo conjugation, also
denoted
\[
  \operatorname{Hom\,ext}(\pi_1(E,e),H).
\]

\subsection{Cones}

\par\medskip\noindent\SGAInlineRefTarget{sga3:x:definition:9-5-1}\textbf{Definitions 9.5.1.---}\enspace
\emph{\SGAInlineRefTarget{sga3:x:definition:9-5-1:item:1}(1) Let \(f:E'\to E\) be a morphism of simplicial sets. Recall
(cf., for example, \([{\rm Del74}]\), 6.3.1) that \(C(f)\) denotes the
pointed simplicial set whose underlying set in degree \(n\geq0\) is}
\[
  E_n\amalg\left(\coprod_{i<n}E'_i\right)\amalg\star,
\]
\emph{where \(\star\) is a singleton. We leave it to the reader to define
the simplicial maps; the faces in degrees at most two are recalled below.
(See also \([{\rm GZ67}]\), VI.2, for a pointed variant.) The category of
locally constant sheaves on \(C(f)\) is equivalent to the category of
locally constant sheaves on \(E\) equipped with a trivialization of the
inverse image on \(E'\). The simplicial set \(C(f)\) is naturally pointed
by \(\star\to C(f)\).}

\emph{(2) Let \(f:T'\to T\) be a morphism of topoi. Denote by \(C(f)\) the
topos whose objects are quintuples}
\[
  \bigl(F,F',A,\alpha:f^*F\to F',
                 \beta:A\to\Gamma(T',F')\bigr),
\]
\emph{where \(F\) (respectively \(F'\)) is an object of \(T\)
(respectively \(T'\)), \(A\) is a set, and \(\alpha\) (respectively
\(\beta\)) is a morphism in \(T'\) (respectively in \(\mathbf{Ens}\)).
(See also \([{\rm Del80}]\), 4.3.4 and \([{\rm Ill72}]\), III.4 for a
variant of this construction.) The category of locally constant sheaves
on \(C(f)\) is equivalent to the category of locally constant sheaves on
\(T\) equipped with a trivialization of the inverse image on \(T'\). The
topos \(C(f)\) is naturally pointed by the fiber functor that sends the
quintuple above to the set \(A\).}

\emph{(3) Let \(f:G'\to G\) be a morphism of groupoids. Denote by \(C(f)\)
the colimit of the diagram}
\[
  G\longleftarrow G'\longrightarrow B1,
\]
\emph{where \(B1\) is the one-object category (one object, one arrow).
The category of locally constant sheaves on \(C(f)\) is equivalent to the
category of locally constant sheaves on \(G\) equipped with a
trivialization of the inverse image on \(G'\).}

\emph{(4) Let \(f:X'\to X\) be a morphism of topological spaces. Denote by
\(C(f)\) the colimit of the diagram}
\[
  \star\longleftarrow X'\times\{0\}
  \longrightarrow X'\times[0,1]
  \longleftarrow X'\times\{1\}
  \longrightarrow X.
\]
\emph{The category of locally constant sheaves on \(C(f)\) is equivalent
to the category of locally constant sheaves on \(X\) equipped with a
trivialization of the inverse image on \(X'\). The topological space
\(C(f)\) is naturally pointed by \(\star\to C(f)\).}

\subsubsection*{9.5.2. Descent data on the cone of a simplicial map}

Let \(f:E'\to E\) be a morphism of simplicial sets, and let \(C(f)\) be
its cone (cf.\ above, \SGARefLink{sga3:x:definition:9-5-1:item:1}{(1)}). Use \(p\) (respectively \(q,r\)) to denote the
face maps of \(E'\) (respectively \(E,C(f)\)). Before stating the
proposition below, we describe the faces \(r\) in degrees at most two in
terms of \(p\) and \(q\). By convention, \(p_{ij}\) (respectively \(p_i\))
is the face map
\[
  E'_2=E'_{\{1,2,3\}}\longrightarrow
  E'_1=E'_{\{1,2\}}
  \quad
  \left(\text{respectively }
  E'_1=E'_{\{1,2\}}\longrightarrow E'_0=E'_{\{1\}}\right)
\]
corresponding to the increasing map
\[
  \{1,2\}\longrightarrow\{1,2,3\}
  \quad
  \left(\text{respectively }
  \{1\}\longrightarrow\{1,2\}\right)
\]
whose image is \(\{i,j\}\) (respectively
\(\{i\}\)).%
\begingroup
\renewcommand{\thefootnote}{2}
\footnote{This convention departs from the present simplicial convention
but is closer to the notation used by Grothendieck in his theory of
descent.}
\endgroup
\addtocounter{footnote}{-1}%
The same convention is adopted for the faces of \(E\) and \(C(f)\).

The morphism
\[
  r_1:C(f)_1=E_1\amalg E'_0\amalg\star
  \longrightarrow C(f)_0=E_0\amalg\star
\]
(respectively \(r_2\)) is:
\begin{itemize}
\item \(q_1:E_1\to E_0\) (respectively \(q_2:E_1\to E_0\)) on \(E_1\);
\item \(E'_0\to\star\) (respectively \(f_0:E'_0\to E_0\)) on \(E'_0\);
\item \(\star\to\star\) on \(\star\).
\end{itemize}
Likewise, the morphism
\[
  r_{21}:C(f)_2=E_2\amalg E'_1\amalg E'_0\amalg\star
  \longrightarrow C(f)_1=E_1\amalg E'_0\amalg\star
\]
(respectively \(r_{32},r_{31}\)) is:
\begin{itemize}
\item \(q_{21}:E_2\to E_1\) (respectively \(q_{32},q_{31}\)) on \(E_2\);
\item \(p_1:E'_1\to E'_0\) (respectively
  \(f_1:E'_1\to E_1\), \(p_2:E'_1\to E'_0\)) on \(E'_1\);
\item \(E'_0\to\star\) (respectively
  \(\operatorname{Id}:E'_0\to E'_0\),
  \(\operatorname{Id}:E'_0\to E'_0\)) on \(E'_0\).
\end{itemize}

Let \(H\) be a descent datum on \(C(f)\), that is, a sheaf \(H_0\) on
\[
  C(f)_0=E_0\amalg\star
\]
equipped with an isomorphism
\[
  \gamma:r_1^*H_0\xrightarrow{\sim}r_2^*H_0
\]
between its inverse images on
\[
  C(f)_1=E_1\amalg E'_0\amalg\star,
\]
satisfying the cocycle relation
\[
  r_{31}^*\gamma=r_{32}^*\gamma\circ r_{21}^*\gamma.
\]

The restriction of \(\gamma\) to \(E_1\) (respectively \(E'_0\)) is a
descent datum
\[
  \beta:q_1^*G_0\xrightarrow{\sim}q_2^*G_0,
\]
where \(G_0\) is the restriction of \(H_0\) to \(E_0\) (respectively, a
trivialization
\[
  t:H_0(\star)\xrightarrow{\sim}f^*G_0=:F_0,
\]
where \(H_0(\star)\) is the constant sheaf with stalk \(H_0(\star)\)).
The restriction of the cocycle relation satisfied by \(\gamma\) to
\(E_2\) (respectively \(E'_1\)) is the cocycle relation for \(\beta\)
(respectively,
\[
  p_2^*(t)=f_1^*(\beta)\circ p_1^*(t)).
\]
This latter relation says that the trivialization \(t\) is compatible
with the descent datum \(f_1^*(\beta)\) induced by \(\beta\) on \(F_0\).
It follows that:

\par\medskip\noindent\SGAInlineRefTarget{sga3:x:proposition:9-5-3}\textbf{Proposition 9.5.3.---}\enspace
\emph{Let \(f:E'\to E\) be a morphism of simplicial sets, and let \(C(f)\)
be its cone, pointed by \(\star\). The category of descent data on
\(C(f)\) trivialized over \(\star\) is equivalent to the category of
descent data on \(E\) equipped with a trivialization of the descent datum
induced on \(E'\).}

\subsection{Representability of the functor
\(\overline{\pi}^{\,1}(S'/S,\xi;{-})\)}

The notation is that of page 90, and from now on we suppose that the
connected components of
\[
  S_0=S'
  \qquad\text{and}\qquad
  S_1=S'\times_S S'
\]
are open. Under this hypothesis we have the following explicit
combinatorial description, which follows immediately from the preceding
recall.

\par\medskip\noindent\SGAInlineRefTarget{sga3:x:proposition:9-6-1}\textbf{Proposition 9.6.1.---}\enspace
\emph{For every group \(G\), the set
\(\overline{\pi}^{\,1}(S'/S,\xi;G)\) of isomorphism classes of descent
data equipped with a trivialization over \(\xi\) is naturally in
bijection with the relative nonabelian cohomology set}
\[
  H^1(K\mathbin{\mathrm{rel}}k,G)
  :=H^1(\widetilde K\mathbin{\mathrm{rel}}\star,G).
\]

It follows that this functor is representable by a group, denoted
\(\pi_1(K,k)\), the \emph{relative fundamental group}, as soon as the
simplicial set
\[
  \widetilde K=C(k\to K)
\]
is connected. One verifies immediately that this holds if and only if
the map \(\pi_0(k)\to\pi_0(K)\) is surjective. It is enough, in
particular, that the simplicial set \(K\) be connected. As indicated in
the text, this is the case when the scheme \(S\) is connected and the
morphism \(S'\to S\) is universally submersive.

\subsection{Contractibility of \(k\); counterexample to the injectivity of
\(k_\bullet\to K_\bullet\)}

The following two propositions make Editors' note (53) more precise.

\par\medskip\noindent\SGAInlineRefTarget{sga3:x:proposition:9-7-1}\textbf{Proposition 9.7.1.---}\enspace
\emph{Let \(\kappa\) be an algebraically closed field, and let \(X\) be a
connected \(\kappa\)-scheme. For every integer \(i\geq0\), let \(X_i\)
denote the \((i+1)\)-st fiber power of \(X\) over \(\kappa\). For every
integer \(i\geq0\), the canonical map}
\[
  \pi_0(X_i)\longrightarrow\pi_0(X)^{i+1}
\]
\emph{is a bijection. In particular, the simplicial set
\(\pi_0(X_\bullet)\) is contractible.}

It suffices to prove the following lemma, which follows by passage to the
limit from the Künneth formula. [Expand?]

\par\medskip\noindent\SGAInlineRefTarget{sga3:x:lemma:9-7-2}\textbf{Lemma 9.7.2.---}\enspace
\emph{Let \(\kappa\) be an algebraically closed field, and let \(X\) and
\(Y\) be two connected \(\kappa\)-schemes. Then the fiber product
\(X\times_\kappa Y\) is connected.}

\par\medskip\noindent\SGAInlineRefTarget{sga3:x:proposition:9-7-3}\textbf{Proposition 9.7.3.---}\enspace
\emph{Let \(X\) be a connected scheme, let \(X'\to X\) be a finite étale
morphism, and let \(\bar{x}\) be a geometric point of \(X\) localized at
a point \(x\). The canonical morphism}
\[
  \pi_0(X'_{\bar{x}})\longrightarrow\pi_0(X')
\]
\emph{is surjective.}

Indeed, the image of every connected component of \(X'\) is open and
closed in \(X\), and therefore meets the fiber \(X'_x\). Thus
\[
  \pi_0(X'_x)\longrightarrow\pi_0(X')
\]
is surjective, as is
\[
  \pi_0(X'_{\bar{x}})\longrightarrow\pi_0(X'_x).
\]

\section*{Bibliography}
\addcontentsline{toc}{section}{Bibliography}
\footnote{\textit{Editors' note.} Additional references cited in this
Exposé; the references \([{\rm AM69}]\) and following concern the
Addenda.}

\begin{description}
\item[\({[{\rm BAC}]}\)]
  N.~Bourbaki, \emph{Algèbre commutative}, Chap.~I--IV, Hermann, 1961;
  Masson, 1985; Springer-Verlag, 2006.
\item[\({[{\rm BEns}]}\)]
  N.~Bourbaki, \emph{Théorie des ensembles}, Chap.~I--IV, Hermann, 1970.
\item[\({[{\rm Pes66}]}\)]
  C.~Peskine, \emph{Une généralisation du ``Main Theorem'' de Zariski},
  Bull.\ Sci.\ Math.\ \textbf{90} (1966), 119--127.
\item[\({[{\rm PY06}]}\)]
  G.~Prasad and J.-K.~Yu, \emph{On quasi-reductive group schemes},
  J.\ Alg.\ Geom.\ \textbf{15} (2006), 507--549.
\item[\({[{\rm Ray70}]}\)]
  M.~Raynaud, \emph{Anneaux locaux henséliens}, Lect.\ Notes Maths.\
  \textbf{169}, Springer-Verlag, 1970.
\item[\({[{\rm AM69}]}\)]
  M.~Artin and B.~Mazur, \emph{Etale homotopy}, Lect.\ Notes Maths.\
  \textbf{100}, Springer, 1969.
\item[\({[{\rm Del74}]}\)]
  P.~Deligne, \emph{Théorie de Hodge III}, Publ.\ Math.\ IHÉS
  \textbf{44} (1974), 5--77.
\item[\({[{\rm Del80}]}\)]
  P.~Deligne, \emph{La conjecture de Weil II}, Publ.\ Math.\ IHÉS
  \textbf{52} (1980), 137--252.
\item[\({[{\rm Fri82}]}\)]
  E.~M.~Friedlander, \emph{Etale homotopy of simplicial schemes},
  Princeton Univ.\ Press, 1982.
\item[\({[{\rm GZ67}]}\)]
  P.~Gabriel and M.~Zisman, \emph{Calculus of fractions and homotopy
  theory}, Springer, 1967.
\item[\({[{\rm Ill72}]}\)]
  L.~Illusie, \emph{Complexe cotangent et déformations I \& II},
  Lect.\ Notes Maths.\ \textbf{239} \& \textbf{283}, Springer,
  1971 \& 1972.
\item[\({[{\rm Kan58}]}\)]
  D.~Kan, \emph{A combinatorial definition of homotopy groups},
  Ann.\ of Math.\ \textbf{67} (1958), 282--312.
\item[\({[{\rm May67}]}\)]
  J.~P.~May, \emph{Simplicial objects in algebraic topology},
  Univ.\ of Chicago Press, 1967.
\end{description}

% END INCLUDED SGA3_English_current_buildable_source_20260823/source/inputs/X/tex/components/03_expose_X_sections7_9_bibliography.tex



% END INCLUDED SGA3_English_current_buildable_source_20260823/source/inputs/X/body.tex


\SGAStartExpose{XI}
% BEGIN INCLUDED SGA3_English_current_buildable_source_20260823/source/inputs/XI/body.tex
% Mechanically extracted after \mainmatter from the frozen standalone master.

% BEGIN INCLUDED SGA3_English_current_buildable_source_20260823/source/inputs/XI/tex/components/01_expose_XI_opening_through_remark_1_5.tex
% SGA 3, Exposé XI, English reconstruction.
% Scope: title, Introduction, and §1 through Remark 1.5.
% Authority: Polo--Gille Expo11.pdf, component-local pp. 1--4,
% complete-reader pp. 723--726, printed pp. 127--132 and 86--88.
% Jacob Reinhold's CC BY 4.0 Markdown lineage was used only as a phrasing
% comparison. The Polo--Gille PDF controls text, formulas, numbering, and notes.

\setcounter{footnote}{-1}
\SGARefTarget{sga3:xi:chapter:xi}\chapter*{Exposé XI\\
Representability Criteria\\
Applications to Multiplicative-Type Subgroups\\
of Affine Group Schemes}
\addcontentsline{toc}{chapter}{Exposé XI. Representability Criteria and
Multiplicative-Type Subgroups}
\markboth{EXPOSÉ XI. REPRESENTABILITY CRITERIA}
  {EXPOSÉ XI. REPRESENTABILITY CRITERIA}

\begin{center}
  \textit{by A. Grothendieck}\footnote{Editor's note: xy version of
  December 5, 2008.}
\end{center}

\SGARefTarget{sga3:xi:section:0}\section*{0. Introduction}
\addcontentsline{toc}{section}{0. Introduction}

% Source page: Expo11 local p. 1; combined-reader p. 723; printed pp. 127--128.
As we have already seen in Exposé X, nos.~4 and 5, the representability
of certain functors---for example, functors of the form
\(\operatorname{Hom}_S(X,Y)\) and various related constructions---plays
an important role in many questions concerning group preschemes.

Among the results that are particularly useful in this theory, besides
the representability of quotients studied in Exposés V and
\(\mathrm{VI}_{\mathrm B}\) and in Exposé VIII, no.~5, let us mention the
representability of functors of the form
\[
  \prod_{X/S} Y/X ,
\]
where \(Y\) is a subobject of \(X\).  This was studied in a very
elementary case in Exposé VIII, no.~6; variants will be given in no.~\SGARefLink{sga3:xi:section:6}{6}
of the present exposé.  These results provide the representability of
various centralizers, normalizers, and transporters.

Less elementary representability criteria, using results that will
appear in EGA VI, are indicated in \SGARefLink{sga3:xi:remark:6:12}{6.12} and in Exposés XV and XVI.
There we shall give a criterion for the representability of quotients
\(G/H\) in cases not covered by the preceding exposés; this criterion
was not developed in the oral lectures.

Our principal purpose in the present exposé is to prove Theorems \SGARefLink{sga3:xi:theorem:4:1}{4.1}
and \SGARefLink{sga3:xi:corollary:4:2}{4.2}, which provide a typical example of a nonprojective
construction technique, close to the technique that will be developed
in EGA VI.  Since the oral lecture and the drafting of the present
text, it has moreover become apparent that the affine hypotheses in
\SGARefLink{sga3:xi:theorem:4:1}{4.1} and \SGARefLink{sga3:xi:corollary:4:2}{4.2} can be removed to a large extent (cf.\ Exposé XV), and
that the main part of the theory developed in the following exposé can
dispense with \SGARefLink{sga3:xi:theorem:4:1}{4.1} and \SGARefLink{sga3:xi:corollary:4:2}{4.2}.  Finally, it might be interesting to prove
analogues of these results for a general reductive group prescheme
(for example, a semisimple one) in place of a group of multiplicative
type; in that case \SGARefLink{sga3:xi:theorem:4:1}{4.1} and \SGARefLink{sga3:xi:corollary:4:2}{4.2} would doubtless be the key result in the
proof.

\SGARefTarget{sga3:xi:section:1}\section*{1. Review of Smooth, Étale,
and Unramified Morphisms}
\addcontentsline{toc}{section}{1. Review of Smooth, Étale, and
Unramified Morphisms}

% Source page: Expo11 local p. 2; combined-reader p. 724; printed pp. 86 and 129.
The reader is referred to EGA IV, §§17 and 18, and, pending its
publication, to SGA 1, Exposés I, II, and III.  Certain noetherian
hypotheses used there, which are inconvenient in applications, should
however be replaced by finite-presentation hypotheses.

\medskip
\noindent\SGARefTarget{sga3:xi:definition:1:1}\textbf{Definition 1.1.} ---
\textit{Let \(S\) be a prescheme and let
\[
  F:(\mathbf{Sch})^\circ/S\longrightarrow(\mathbf{Ens})
\]
be a functor.  The functor \(F\) is said to be \emph{formally smooth}
(respectively, \emph{formally unramified}\footnote{The source notes
that one would now rather say ``net'' in place of ``unramified.''},
respectively, \emph{formally étale}) if, for every \(S\)-prescheme
\(S'\) that is affine in the absolute sense and every subscheme
\(S'_0\subset S'\) defined by a nilpotent ideal \(\mathscr J\), the map
\[
  F(S')\longrightarrow F(S'_0)
\]
is surjective (respectively, injective, respectively, bijective).}

\textit{An \(S\)-prescheme \(X\) is said to be formally smooth over
\(S\) (respectively, formally unramified over \(S\), respectively,
formally étale over \(S\)) if the corresponding functor has that
property.  It is said to be smooth over \(S\) (respectively,
unramified over \(S\), respectively, étale over \(S\)) if it satisfies
the preceding condition and is, in addition, locally of finite
presentation over \(S\).}

The usefulness of these definitions for \(X\) lies, on the one hand,
in their remarkably simple expression in terms of the functor
represented by \(X\) (and in practice \(X\) is often given as the
object over \(S\) representing an explicitly described functor), and,
on the other hand, in the remarkable local structural properties that
they express.  We recall those properties in the following statements;
for their proofs, see the references just cited.

\medskip
\noindent\SGARefTarget{sga3:xi:proposition:1:2}\textbf{Proposition 1.2.} ---
\textit{Let \(X\) be a prescheme locally of finite presentation over
\(S\).  Then:}

\begin{enumerate}
\item[(i)] \SGARefTarget{sga3:xi:proposition:1:2:item:i} \textit{\(X\) is smooth over \(S\) if and only if it is flat
over \(S\) and its geometric fibers
\[
  X\otimes_S\operatorname{Spec}\kappa(s)
\]
are regular schemes.  More generally, \(X\) is smooth over \(S\) in a
neighborhood of \(x\in X\) (one then says that \(X\) is smooth over
\(S\) at \(x\)) if and only if \(X\) is flat over \(S\) at \(x\) and
\(X_s\) is smooth over \(\kappa(s)\) at \(x\); that is,
\[
  X_s\otimes_{\kappa(s)}\overline{\kappa(s)}
\]
is regular at the points (or merely at one point) lying over \(x\).}

\item[(ii)] \SGARefTarget{sga3:xi:proposition:1:2:item:ii} \textit{Suppose that \(S\), and hence \(X\), is locally
noetherian.  Let \(x\in X\), and let \(s\in S\) be its image.  The
smoothness of \(X\) over \(S\) at \(x\) can then be detected from the
local homomorphism
\[
  A=\mathscr O_{S,s}\longrightarrow B=\mathscr O_{X,x}
\]
of noetherian local rings (or even from the local homomorphism
\(\widehat A\to\widehat B\) of their completions), by the following
characteristic property: \(B\) is flat over \(A\), and
\(B\otimes_A k\) is geometrically regular over \(k\), the residue
field of \(A\).  Equivalently, for every finite extension \(k'\) of
\(k\),
\[
  (B\otimes_A k)\otimes_k k'=B\otimes_A k'
\]
is a regular semilocal ring.  When the residue-field extension
\(k(B)/k(A)\) is trivial, these conditions are also equivalent to the
following one: \(\widehat B\) is isomorphic, as an
\(\widehat A\)-algebra, to an algebra of formal power series
\[
  \widehat A[[t_1,\ldots,t_n]].
\]}
\end{enumerate}

Thus, from the ``formal'' point of view, the structure of \(X\) over
\(S\) is that of the standard affine space \(S[t_1,\ldots,t_n]\) over
\(S\).

% Source page: Expo11 local p. 3; combined-reader p. 725; printed pp. 87 and 130--131.
\medskip
\noindent\SGARefTarget{sga3:xi:proposition:1:3}\textbf{Proposition 1.3.} ---
\textit{Let \(X\) be a prescheme locally of finite presentation over
\(S\).  Then:}

\begin{enumerate}
\item[(i)] \SGARefTarget{sga3:xi:proposition:1:3:item:i} \textit{The following conditions are equivalent:}
  \begin{enumerate}
  \item[(a)] \SGARefTarget{sga3:xi:proposition:1:3:item:i:item:a} \textit{\(X\) is unramified over \(S\).}
  \item[(b)] \SGARefTarget{sga3:xi:proposition:1:3:item:i:item:b} \textit{The diagonal morphism
  \(X\to X\times_S X\) is an open immersion.}
  \item[(c)] \SGARefTarget{sga3:xi:proposition:1:3:item:i:item:c} \textit{\(\Omega^1_{X/S}=0\).}
  \item[(d)] \SGARefTarget{sga3:xi:proposition:1:3:item:i:item:d} \textit{The geometric fibers of \(X/S\) are discrete and
  reduced; that is, they are isomorphic to disjoint sums of copies of
  the base field.}
  \item[(e)] \SGARefTarget{sga3:xi:proposition:1:3:item:i:item:e} \textit{For every \(s\in S\), the fiber
  \[
    X\otimes_S\operatorname{Spec}\kappa(s)=X_s
  \]
  is unramified over \(\kappa(s)\), or equivalently is isomorphic to a
  disjoint sum of spectra of finite separable extensions of
  \(\kappa(s)\).}
  \end{enumerate}

\item[(ii)] \SGARefTarget{sga3:xi:proposition:1:3:item:ii} \textit{There are analogous pointwise conditions for
\(X\) to be unramified over \(S\) at a given point \(x\), that is, in
a neighborhood of that point.  For example, it is necessary and
sufficient that \(\mathscr O_{X_s,x}\) be a finite separable extension
of \(\kappa(s)\).  In terms of the local homomorphism
\[
  A=\mathscr O_{S,s}\longrightarrow B=\mathscr O_{X,x},
\]
this says that \(B\otimes_A k\) is a finite separable extension of
the residue field \(k\) of \(A\); equivalently,
\[
  \mathfrak r(A)B=\mathfrak r(B)
  \quad\text{and}\quad
  k(B)/k(A)\ \text{is finite separable}.
\]
When \(A\), and hence \(B\), is noetherian and
\(k(A)\xrightarrow{\sim}k(B)\), this also means that
\(\widehat A\to\widehat B\) is surjective.}
\end{enumerate}

Thus, from the ``formal'' point of view, saying that \(X\) is
unramified over \(S\) means that \(X\) is essentially a subprescheme
of \(S\).  Observe also that conditions \SGARefLink{sga3:xi:proposition:1:3:item:i:item:d}{(d)} and \SGARefLink{sga3:xi:proposition:1:3:item:i:item:e}{(e)} are expressed
solely in terms of the fibers of \(X/S\).

\medskip
\noindent\SGARefTarget{sga3:xi:proposition:1:4}\textbf{Proposition 1.4.} ---
\textit{Let \(X\) be a prescheme locally of finite presentation over
\(S\).  Then \(X\) is étale over \(S\) if and only if it is smooth and
unramified over \(S\) (which is immediate from the definition), so the
criteria in \SGARefLink{sga3:xi:proposition:1:2}{1.2} and \SGARefLink{sga3:xi:proposition:1:3}{1.3} apply.  In particular:}

\begin{enumerate}
\item[(i)] \SGARefTarget{sga3:xi:proposition:1:4:item:i} \textit{\(X\) is étale over \(S\) if and only if it is flat
and unramified over \(S\).  There is an analogous local criterion for
\(X\) to be étale over \(S\) at a given point \(x\).}

\item[(ii)] \SGARefTarget{sga3:xi:proposition:1:4:item:ii} \textit{Suppose in addition that \(S\), and hence \(X\),
is locally noetherian.  Whether \(X\) is étale over \(S\) at \(x\) can
then be detected from the local homomorphism
\[
  A=\mathscr O_{S,s}\longrightarrow B=\mathscr O_{X,x}
\]
\footnote{The source prints \(\mathscr O_{X,s}\); the local ring at the
point of \(X\) is \(\mathscr O_{X,x}\).}
(and even from \(\widehat A\to\widehat B\)) by the following
characteristic property: \(B\) is flat over \(A\), and \(B\otimes_A k\)
is a finite separable extension of \(k\), the residue
field\footnote{Editor's note: ``residue extension'' was replaced by
``residue field.''} of \(A\).  When
\(k(A)\xrightarrow{\sim}k(B)\), this condition simply says that
\(\widehat A\to\widehat B\) is an isomorphism.}
\end{enumerate}

Thus, from the ``formal'' point of view, saying that \(X\) is étale
over \(S\) simply means that \(X\) is locally isomorphic to \(S\).

\medskip
\noindent\SGARefTarget{sga3:xi:remark:1:5}\textbf{Remark 1.5.} ---
When \(X\) is locally of finite presentation over \(S\) and a point
\(x\in X\) is given, the fact that \(X\) is smooth (respectively,
unramified, respectively, étale) over \(S\) at \(x\), that is, in a
neighborhood of \(x\), can also be detected from the functor
\[
  X:(\mathbf{Sch})^\circ/S\longrightarrow(\mathbf{Ens})
\]
by the following property.  For every \(S\)-scheme \(S'\) that is the
spectrum of a local ring, every subscheme \(S'_0\subset S'\) defined
by a nilpotent ideal, and every \(S\)-morphism
\(u_0:S'_0\to X\) sending the closed point \(s'\) of \(S'\) to \(x\),
there exists at least one (respectively, at most one, respectively,
exactly one) \(S\)-morphism \(u:S'\to X\) extending \(u_0\).

This statement shows in particular that, in \SGARefLink{sga3:xi:definition:1:1}{Definition 1.1}, one may
restrict to local schemes \(S'\), provided the functor under
consideration is representable by a prescheme locally of finite
presentation over \(S\).  Moreover, when \(S\) is locally noetherian,
one may even restrict the preceding pointwise criterion to local
artinian schemes \(S'\), and the same restriction may therefore be
made in \SGARefLink{sga3:xi:definition:1:1}{Definition 1.1}, provided that \(S\) is locally noetherian and
that the functor in question is represented by a prescheme locally of
finite type over \(S\).

% Exact continuation: Remark 1.6, Expo11 local p. 4 / combined p. 726 /
% printed p. 132.

% END INCLUDED SGA3_English_current_buildable_source_20260823/source/inputs/XI/tex/components/01_expose_XI_opening_through_remark_1_5.tex

% BEGIN INCLUDED SGA3_English_current_buildable_source_20260823/source/inputs/XI/tex/components/02_expose_XI_remark_1_6.tex
% SGA 3, Exposé XI, English reconstruction.
% Scope: Remark 1.6 only, stopping immediately before Remark 1.7.
% Authority: Polo--Gille Expo11.pdf, component-local pp. 4--5,
% complete-reader pp. 726--727, printed pp. 132--133 and 88--89.
% Jacob Reinhold's CC BY 4.0 Markdown lineage was used only as a phrasing
% comparison. The Polo--Gille PDF controls text, formulas, and diagrams.

\medskip
\noindent\SGARefTarget{sga3:xi:remark:1:6}\textbf{Remark 1.6.} ---
Of course, given a morphism \(f:X\to Y\) of preschemes, one says that
\(f\) is smooth (respectively, unramified, respectively, étale) if
\(f\) makes \(X\) a smooth (respectively, unramified, respectively,
étale) \(Y\)-prescheme.  When \(X\) and \(Y\) are \(S\)-preschemes and
\(f\) is an \(S\)-morphism of finite presentation, these properties of
\(f\) can be expressed immediately in terms of the morphism of
functors
\[
  F,G:(\mathbf{Sch})^\circ/S\longrightarrow(\mathbf{Ens})
\]
defined by \(f\).  Namely, \(f\) is smooth (respectively, unramified,
respectively, étale) if, for every \(S\)-prescheme \(S'\) that is
affine in the absolute sense and every subscheme \(S'_0\subset S'\)
defined by a nilpotent ideal \(\mathscr J\), the map
\[
  F(S')\longrightarrow
  F(S'_0)\times_{G(S'_0)}G(S')
\]
deduced from the commutative square

\SGARefTarget{sga3:xi:diagram:1:6:functor-square}
\[
\begin{tikzcd}[column sep=huge,row sep=large]
  F(S') \arrow[r] \arrow[d] & G(S') \arrow[d] \\
  F(S'_0) \arrow[r] & G(S'_0)
\end{tikzcd}
\]
is surjective (respectively, injective, respectively, bijective).
Equivalently, for every commutative square of morphisms of functors
over \(S\), where \(S'\) and \(S'_0\) are as above and
\(i:S'_0\to S'\) is the canonical immersion,

\SGARefTarget{sga3:xi:diagram:1:6:lifting-square}
\[
\begin{tikzcd}[column sep=huge,row sep=large]
  \text{\rm(Q)} & S'_0 \arrow[r,"i"] \arrow[d,"u_0"] &
    S' \arrow[d,"v"'] \\
  & F \arrow[r,"f"] & G ,
\end{tikzcd}
\]
there exists at least one (respectively, at most one, respectively,
exactly one) morphism

\SGARefTarget{sga3:xi:diagram:1:6:lift}
\[
\begin{tikzcd}[column sep=huge,row sep=large]
  & S' \arrow[dl,"u"'] \\
  F &
\end{tikzcd}
\]
that makes the two corresponding triangles commute:
\[
  u_0=u\circ i
  \qquad\text{and}\qquad
  v=f\circ u.
\]

This property for a homomorphism of functors \(f:F\to G\) over \(S\)
retains its meaning even when the functors under consideration are
not representable.  If it holds, the homomorphism \(f:F\to G\) is
said to be \emph{formally smooth} (respectively,
\emph{formally unramified}, respectively, \emph{formally étale}).
Notice that this property depends only on the homomorphism of the
underlying functors
\[
  (\mathbf{Sch})^\circ\longrightarrow(\mathbf{Ens})
\]
defined by \(F\) and \(G\) (cf.\ \SGARefLink{sga3:i:proposition:1:4:1}{Exposé I, 1.4.1}), and not on the
structure morphisms \(F\to S\) and \(G\to S\).  An equivalent
formulation of the preceding definition, somewhat easier to use in
applications, is the following: a morphism of functors \(F\to G\)
over \(S\) is formally smooth (respectively, formally unramified,
respectively, formally étale) if, for every \(S'\) over \(S\) and
every morphism \(S'\to G\), the functor
\[
  F'=F\times_G S'
\]
over \(S'\) is formally smooth (respectively, formally unramified,
respectively, formally étale).

% END INCLUDED SGA3_English_current_buildable_source_20260823/source/inputs/XI/tex/components/02_expose_XI_remark_1_6.tex

% BEGIN INCLUDED SGA3_English_current_buildable_source_20260823/source/inputs/XI/tex/components/03_expose_XI_remarks_1_7_1_8.tex
% SGA 3, Exposé XI, English reconstruction.
% Scope: Remarks 1.7 and 1.8 only, beginning after Remark 1.6 and
% stopping immediately before Proposition 1.9.
% Authority: Polo--Gille Expo11.pdf, component-local pp. 5--6,
% complete-reader pp. 727--728, source-margin pp. 133--135.
% Jacob Reinhold's CC BY 4.0 Markdown lineage was used only as a phrasing
% comparison. The Polo--Gille PDF controls text, formulas, and diagrams.

\medskip
\noindent\SGARefTarget{sga3:xi:remark:1:7}\textbf{Remark 1.7.} ---
Under the hypotheses of \SGARefLink{sga3:xi:remark:1:6}{1.6}, suppose that we are given a ``point of
\(F\) with values in a field \(k\),'' that is, an element
\(x\in F(\operatorname{Spec}k)\), or equivalently a morphism
\[
  x:\operatorname{Spec}k\longrightarrow F.
\]
We likewise say that \(f\) is \emph{formally smooth} (respectively,
\emph{formally unramified}, respectively, \emph{formally étale})
\emph{at \(x\)} if the condition in the preceding definition is
satisfied whenever \(S'\) is a local scheme and the morphism
\(S'_0\to F\) in diagram \(\mathrm{(Q)}\) above is ``compatible with
the marked points'' in the following sense.  Let \(k'\) denote the
residue field of \(S'\) and \(S'_0\) at their closed point.  Then the
diagram

\SGARefTarget{sga3:xi:diagram:1:7:marked-point-compatibility}
\[
\begin{tikzcd}[column sep=huge,row sep=large]
  \operatorname{Spec}(k') \arrow[d,"j"'] &
    \operatorname{Spec}(k) \arrow[d,"x"] \\
  S'_0 \arrow[r,"u_0"] & F
\end{tikzcd}
\]
in which \(j:\operatorname{Spec}(k')\to S'_0\) is the canonical
immersion, must admit a completion to a commutative diagram

\SGARefTarget{sga3:xi:diagram:1:7:common-field-extension}
\[
\begin{tikzcd}[column sep=huge,row sep=large]
  & \operatorname{Spec}(k'') \arrow[dl] \arrow[dr] & \\
  \operatorname{Spec}(k') \arrow[d] &&
    \operatorname{Spec}(k) \arrow[d] \\
  S'_0 \arrow[rr] && F ,
\end{tikzcd}
\]
where \(k''\) is a field.\footnote{The source says that \(k''\) is
``the spectrum of a field''; the displayed object
\(\operatorname{Spec}(k'')\) makes clear that \(k''\) itself is the
field.}

When \(F\) and \(G\) are represented by \(S\)-preschemes \(X\) and
\(Y\), and the \(S\)-morphism \(f:X\to Y\) is locally of finite
presentation, this condition says precisely, by \SGARefLink{sga3:xi:remark:1:5}{1.5}, that
\(f:X\to Y\) is smooth at the point of \(X\) that is the image of
\(\operatorname{Spec}(k)\) under
\(x:\operatorname{Spec}(k)\to X\).

\medskip
\noindent\SGARefTarget{sga3:xi:remark:1:8}\textbf{Remark 1.8.} ---
When the condition in the preceding definition is satisfied after
restricting to local artinian schemes \(S'\), we say that
\(F\to G\) is \emph{infinitesimally smooth} (respectively,
\emph{infinitesimally unramified}, respectively,
\emph{infinitesimally étale}) \emph{at \(x\)}.  We say that
\(F\to G\) is infinitesimally smooth (respectively,
infinitesimally unramified, respectively, infinitesimally étale) if
it has the corresponding property at every point \(x\); equivalently,
if the condition considered in \SGARefLink{sga3:xi:remark:1:6}{1.6} is satisfied whenever \(S'\) is
local artinian.

This variant of the preceding notions is technically useful.  Since
it is weaker, it is often easier to verify, while it is nevertheless
frequently sufficient to imply the stronger condition; this is the
case, for example, when \(F\to G\) is a morphism locally of finite
presentation and \(G\) is represented by a locally noetherian
prescheme.

% Exact continuation: Proposition 1.9, Expo11 local p. 6 /
% combined-reader p. 728 / source-margin p. 135.

% END INCLUDED SGA3_English_current_buildable_source_20260823/source/inputs/XI/tex/components/03_expose_XI_remarks_1_7_1_8.tex

% BEGIN INCLUDED SGA3_English_current_buildable_source_20260823/source/inputs/XI/tex/components/04_expose_XI_proposition_1_9_through_corollary_1_11.tex
% SGA 3, Exposé XI, English reconstruction.
% Scope: Proposition 1.9 and its explanatory discussion, Proposition 1.10
% with proof, and Corollary 1.11 with proof; stopping immediately before
% Section 2.
% Authority: Polo--Gille Expo11.pdf, component-local pp. 6--7,
% complete-reader pp. 728--729, printed pp. 135--137.
% Jacob Reinhold's CC BY 4.0 Markdown lineage was used only as a phrasing
% comparison. The Polo--Gille PDF controls text and formulas.

\medskip
\noindent\SGARefTarget{sga3:xi:proposition:1:9}\textbf{Proposition 1.9.} ---
\textit{Let \(f:X\to S\) be a smooth morphism of preschemes.  Then
\(\Omega^1_{X/S}\) is a locally free module of finite type on \(X\),
and its rank at a point \(x\in X\) is equal to the dimension of the
fiber \(X_s\), where \(s=f(x)\), in a neighborhood of \(x\).}

This dimension is called the \emph{relative dimension of \(X\) over
\(S\) at \(x\)}.  Notice that it is zero (when \(f\) is smooth at
\(x\)) if and only if \(f\) is étale at \(x\).  It can also be
computed directly in terms of the functor \(F\) represented by \(X\),
as follows.  Let \(\xi\) be a point of \(F\) with values in a field
\(k\), ``localized at \(x\),'' that is, a morphism
\[
  \operatorname{Spec}(k)\longrightarrow F=X
\]
whose image is \(x\).  Consider the algebra
\[
  D(k)=k[t]/(t^2)
\]
of dual numbers over \(k\), regarded as an \(S\)-prescheme, and the
map
\SGARefTarget{sga3:xi:display:1:9:dual-number-map}
\[
  F(D(k))\longrightarrow F(k)
\]
induced by the augmentation \(D(k)\to k\).  Finally, let
\(F(D(k),\xi)\) be the inverse image of \(\xi\in F(k)\) under this
map.  This set is naturally endowed with the structure of a vector
space over \(k\).  In fact, it is the vector-space dual of
\[
  \Omega^1_{X/S}\otimes_{\mathscr O_{X,x}}k,
\]
and its dimension is the relative dimension of \(X\) over \(S\) at
\(x\).

To describe the vector-space law on \(F(D(k),\xi)\) explicitly, it is
useful to introduce, more generally and as in Exposé II, for every
vector space \(V\) over \(k\), the algebra
\[
  D_k(V)=k+V,
\]
where \(V\) is an ideal of square zero, and to regard
\(F(D_k(V),\xi)\), the inverse image of \(\xi\) under
\[
  F(D_k(V))\longrightarrow F(k),
\]
as a covariant functor of \(V\) with values in
\((\mathbf{Ens})\).  It is then enough that this functor commute with
products of two factors.  This means that \(F\) carries certain
amalgamated sums of a very special kind to fiber products; compare
Exposé II.  This condition always holds when \(F\) is representable.
It follows that the sets \(F(D_k(V),\xi)\), and in particular
\(F(D(k),\xi)\), are endowed with vector-space structures over \(k\).
Thus one can define the relative dimension of \(F\) over \(X\) at
the ``point'' \(\xi\) under conditions appreciably more general than
the representability of \(F\).

In the present Exposé, the fact that certain functors to be described
explicitly are representable by preschemes smooth over \(S\) will be
used chiefly through \SGARefLink{sga3:xi:proposition:1:10}{the following result}.  It will be our technical
means of passing from constructions over the completion of the local
ring \(A\) of a point \(s\) of a noetherian scheme \(S\) to a local
ring \(A'\) over \(A\) that is étale over \(A\), and hence, in
particular, of passing from \(s\) to neighboring points.

\medskip
\noindent\SGARefTarget{sga3:xi:proposition:1:10}\textbf{Proposition 1.10.} ---
\textit{Let \(f:X\to S\) be a smooth morphism of preschemes, let
\(s\in S\), and let \(x\in X\) lie over \(s\), with \(\kappa(x)\) a
finite separable extension of \(\kappa(s)\).  Then there exists a
subscheme \(S'\subset X\), étale over \(S\), passing through
\(x\).\footnote{The French re-edition prints ``a subscheme \(S'\) of
\(S\), étale over \(S\), passing through \(x\).''  Since \(x\in X\),
and both the equivalent formulation and the construction immediately
below use \(S'\to X\) and sections of \(\mathscr O_X\), the intended
reading is \(S'\subset X\).}  Thus one can find an étale morphism
\(S'\to S\), a point \(s'\in S'\) above \(s\) whose residue-field
extension is \(\kappa(x)/\kappa(s)\), and an \(S\)-morphism
\(S'\to X\) carrying \(s'\) to \(x\).}

To construct \(S'\), simply choose sections
\(g_1,\ldots,g_n\) of \(\mathscr O_X\) on a neighborhood of \(x\)
whose images at \(x\) form a regular system of parameters for the
local ring of the fiber \(X_s\) at \(x\).  The subscheme \(S'\)
defined by the \(g_i\) is then étale over \(S\) at \(x\), and, after
shrinking \(S'\), it is étale over \(S\).

We shall use
\SGARefLink{sga3:xi:proposition:1:10}{1.10} when
\(\kappa(x)=\kappa(s)\), that is, when \(x\) is rational over
\(\kappa(s)\), or equivalently when \(x\) may be regarded as a
section of \(X_s\) over \(\operatorname{Spec}\kappa(s)\).  In that
case, \SGARefLink{sga3:xi:proposition:1:10}{1.10} has the character
of an extension theorem for sections, after an étale extension of
the base.  It takes a particularly simple form in the following
special case.

\medskip
\noindent\SGARefTarget{sga3:xi:corollary:1:11}\textbf{Corollary 1.11.} ---
\textit{Under the hypotheses of
\SGARefLink{sga3:xi:proposition:1:10}{1.10}, suppose that \(S\) is
the spectrum of a henselian local ring and that
\(\kappa(x)=\kappa(s)\).  Then there exists a section of \(X\) over
\(S\) passing through \(x\).  It is uniquely determined if \(X\) is
in fact étale over \(S\) at \(x\).}

Indeed, since \(S\) is henselian, under the hypotheses of
\SGARefLink{sga3:xi:proposition:1:10}{1.10} the prescheme \(S'\)
contains an open subscheme that is finite over \(S\) and whose fiber
at \(s\) is reduced to \(x\).  Since this open subscheme is étale
over \(S\), it is isomorphic to \(S\), which proves the assertion.
Notice that when \(S\) is the spectrum of a complete local ring,
\SGARefLink{sga3:xi:proposition:1:10}{1.10} or
\SGARefLink{sga3:xi:corollary:1:11}{1.11} is more or less equivalent
to the classical ``Hensel's lemma,'' and one sometimes refers to the
result by that name.

% Exact continuation: Section 2, Expo11 local p. 7 /
% complete-reader p. 729 / printed p. 137.

% END INCLUDED SGA3_English_current_buildable_source_20260823/source/inputs/XI/tex/components/04_expose_XI_proposition_1_9_through_corollary_1_11.tex

% BEGIN INCLUDED SGA3_English_current_buildable_source_20260823/source/inputs/XI/tex/components/05_expose_XI_section_2.tex
% SGA 3, Exposé XI, English reconstruction.
% Scope: all of §2, from its heading through Proposition 2.5 and its proof;
% stopping immediately before §3.
% Authority: Polo--Gille Expo11.pdf, component-local pp. 7--9,
% complete-reader pp. 729--731, source-margin pp. 137--141.
% Jacob Reinhold's CC BY 4.0 Markdown lineage was used only as a phrasing
% comparison. The Polo--Gille PDF controls text, formulas, and numbering.

\SGARefTarget{sga3:xi:section:2}
\section*{2. Examples of Formally Smooth Functors from the Theory of
Groups of Multiplicative Type}
\addcontentsline{toc}{section}{2. Examples of Formally Smooth Functors
from the Theory of Groups of Multiplicative Type}

We shall interpret, in the language introduced in the preceding
number, the results stated in Exposé IX, no.~3, concerning
infinitesimal extensions of a homomorphism from a group of
multiplicative type.  Those results are consequences of the
vanishing of the Hochschild cohomology of such a group, established
in Exposé I.

\medskip
\noindent\SGARefTarget{sga3:xi:proposition:2:1}\textbf{Proposition 2.1.} ---
\textit{Let \(S\) be a prescheme, let \(H\) be a group of
multiplicative type over \(S\), and let \(G\) be an \(S\)-group
prescheme smooth over \(S\).  Consider the functor over \(S\)}
\SGARefTarget{sga3:xi:display:2:1:hom-functor}
\[
  M_H=\operatorname{Hom}_{S\text{-}\mathrm{gr}}(H,G).
\]
\textit{Its value at an \(S'\) over \(S\) is the set of
homomorphisms of \(S'\)-groups from \(H_{S'}\) to \(G_{S'}\).
\SGARefLink{sga3:xi:display:2:1:hom-functor}{This functor} is formally smooth over \(S\).}

Cf.\ \SGARefLink{sga3:ix:theorem:3:6}{Exposé IX, 3.6}.  More generally:

\medskip
\noindent\SGARefTarget{sga3:xi:corollary:2:2}\textbf{Corollary 2.2.} ---
\textit{Let \(S\) and \(G\) be as above, and consider a homomorphism
\(u:H_1\to H_2\) of group schemes of multiplicative type over \(S\).
With the preceding notation this induces a morphism of functors over
\(S\)}
\SGARefTarget{sga3:xi:display:2:2:restriction-map}
\[
  M_{H_2}\longrightarrow M_{H_1},
  \qquad
  M_{H_i}=\operatorname{Hom}_{S\text{-}\mathrm{gr}}(H_i,G),
\]
\textit{given by \(w\mapsto w\circ u\).  \SGARefLink{sga3:xi:display:2:2:restriction-map}{This morphism} is formally
smooth.}

Indeed, by the definitions in
\SGARefLink{sga3:xi:remark:1:6}{1.6}, this is equivalent to the
following assertion.  Suppose that \(S\) is affine and that
\(S_0\subset S\) is a subscheme defined by a nilpotent ideal.  Let
\SGARefTarget{sga3:xi:display:2:2:v}
\[
  v:H_1\longrightarrow G
\]
be a homomorphism of \(S\)-groups, and let
\SGARefTarget{sga3:xi:display:2:2:w-zero}
\[
  w_0:(H_2)_{S_0}\longrightarrow G_{S_0}
\]
be a homomorphism of \(S_0\)-groups such that
\[
  w_0u_{S_0}=v_{S_0}.
\]
Then there exists a homomorphism of \(S\)-groups
\SGARefTarget{sga3:xi:display:2:2:w}
\[
  w:H_2\longrightarrow G
\]
extending \(w_0\) and satisfying \(wu=v\).

To see this, first extend \(w_0\) to a homomorphism of \(S\)-groups
\[
  w':H_2\longrightarrow G,
\]
which is possible by
\SGARefLink{sga3:xi:proposition:2:1}{2.1}.  Set
\[
  v'=w'u:H_1\longrightarrow G.
\]
By the hypothesis on \(w_0\), one has \(v'_{S_0}=v_{S_0}\).
Consequently, by Exposé VIII, 3.6, there exists an element
\(g\in G(S)\) whose image in \(G(S_0)\) is the identity element and
such that
\[
  v=\operatorname{int}(g)v'.
\]
It follows that
\[
  v=\bigl(\operatorname{int}(g)w'\bigr)u,
\]
and it is therefore enough to take
\[
  w=\operatorname{int}(g)w'.
\]

\medskip
\noindent\SGARefTarget{sga3:xi:corollary:2:3}\textbf{Corollary 2.3.} ---
\textit{With the notation of
\SGARefLink{sga3:xi:proposition:2:1}{2.1}, regard \(M_H=M\) as a
functor with operator group \(G\), where \(G\) acts by
\[
  (v,g)\longmapsto\operatorname{int}(g)\circ v.
\]
Then the corresponding morphism}
\SGARefTarget{sga3:xi:display:2:3:action-map}
\[
  R:G\times_S M\longrightarrow M\times_S M,
\]
\textit{defined by
\[
  (g,v)\longmapsto\bigl(\operatorname{int}(g)\circ v,v\bigr),
\]
is formally smooth.}

After a base change \(S'\to S\), this is equivalent to the following
assertion.

\medskip
\noindent\SGARefTarget{sga3:xi:corollary:2:4}\textbf{Corollary 2.4.} ---
\textit{With the notation of
\SGARefLink{sga3:xi:proposition:2:1}{2.1}, let
\(v_1,v_2:H\to G\) be two morphisms of \(S\)-groups, and let
\(\operatorname{Transp}(v_1,v_2)\) be the subfunctor of \(G\)
consisting of the \(g\) such that
\[
  \operatorname{int}(g)\circ v_1=v_2.
\]
Then this functor is formally smooth over \(S\).  In particular, if
\(v_1=v_2=v\), the functor \(\operatorname{Centr}(v)\), the subgroup
of \(G\) consisting of the \(h\) such that
\(\operatorname{int}(h)v=v\), is formally smooth over \(S\).}

\emph{N.B.} The pair \((v_1,v_2)\) may be regarded as a section over
\(S\) of the second factor in the codomain of the morphism \(R\) in
\SGARefLink{sga3:xi:corollary:2:3}{2.3}, and
\(\operatorname{Transp}(v_1,v_2)\) as the inverse-image functor of
that section under \(R\).

The assertion in
\SGARefLink{sga3:xi:corollary:2:4}{2.4} is itself equivalent to the
following.  Suppose that \(S\) is affine and that \(S_0\subset S\) is
a subscheme defined by a nilpotent ideal.  If
\(g_0\in G(S_0)\) satisfies
\[
  \operatorname{int}(g_0)(v_1)_{S_0}=(v_2)_{S_0},
\]
then \(g_0\) extends to an element \(g\in G(S)\) such that
\[
  \operatorname{int}(g)v_1=v_2.
\]
To prove this, first extend \(g_0\) to a section \(g'\) of \(G\) over
\(S\), which is possible because \(G\) is smooth over \(S\), and put
\[
  v'_2=\operatorname{int}(g')v_1.
\]
The morphisms \(v_2\) and \(v'_2\) have the same restriction over
\(S_0\).  Hence, by \SGARefLink{sga3:ix:theorem:3:6}{Exposé IX, 3.6}, there exists
\(g''\in G(S)\) inducing the identity section over \(S_0\) and such
that
\[
  v_2=\operatorname{int}(g'')v'_2.
\]
Therefore
\[
  v_2
  =\operatorname{int}(g'')\operatorname{int}(g')v_1
  =\operatorname{int}(g''g')v_1,
\]
and it is enough to take \(g=g''g'\).

\medskip
\noindent\SGARefTarget{sga3:xi:proposition:2:1:bis}\textbf{Proposition 2.1 bis.} ---
\textit{Let \(S\) be a prescheme and let \(G\) be an \(S\)-group
prescheme smooth over \(S\).  Consider the functor}
\SGARefTarget{sga3:xi:display:2:1:bis:functor}
\[
  M:(\mathbf{Sch})^\circ/S\longrightarrow(\mathbf{Ens}),
\]
\textit{defined by}
\SGARefTarget{sga3:xi:display:2:1:bis:value}
\[
  M(S')
  =\left\{
    \begin{array}{c}
      \text{multiplicative-type subgroups}\\[-2pt]
      \text{of \(G_{S'}\)}
    \end{array}
  \right\}.
\]
\textit{Then \(M\) is formally smooth over \(S\).}

Cf.\ \SGARefLink{sga3:ix:theorem:3:6:bis}{Exposé IX, 3.6 bis}.

\medskip
\noindent\SGARefTarget{sga3:xi:corollary:2:2:bis}\textbf{Corollary 2.2 bis.} ---
\textit{Let \(n\) be an integer, and consider the morphism of
functors}
\SGARefTarget{sga3:xi:display:2:2:bis:phi}
\[
  \varphi_n:M\longrightarrow M
\]
\textit{defined by}
\SGARefTarget{sga3:xi:display:2:2:bis:kernel}
\[
  \varphi_n(H)={}_nH=\ker\!\bigl(n\cdot\operatorname{id}_H\bigr).
\]
\textit{Then \(\varphi_n\) is formally smooth.  For every integer
\(p>0\), let \(M_p\) denote the subfunctor of \(M\) for which
\(M_p(S')\) is the set of multiplicative-type subgroups \(H\) of
\(G_{S'}\) satisfying \({}_pH=H\).  Then the morphism induced by
\(\varphi_n\)}
\SGARefTarget{sga3:xi:display:2:2:bis:restricted-map}
\[
  M_{np}\longrightarrow M_n
\]
\textit{is formally smooth.}

The second assertion is trivially contained in the first and is
included only for the convenience of a later reference.  The proof
of the first assertion is analogous to that of
\SGARefLink{sga3:xi:corollary:2:2}{2.2}, now invoking \SGARefLink{sga3:ix:theorem:3:6:bis}{Exposé IX,
3.6 bis}.

\medskip
\noindent\SGARefTarget{sga3:xi:corollary:2:3:bis}\textbf{Corollary 2.3 bis.} ---
\textit{With the notation of
\SGARefLink{sga3:xi:proposition:2:1:bis}{2.1 bis}, regard \(M\) as a
functor with operator group \(G\), where \(G\) acts by
\[
  (g,H)\longmapsto\operatorname{int}(g)(H).
\]
Then the corresponding morphism}
\SGARefTarget{sga3:xi:display:2:3:bis:action-map}
\[
  R:G\times_S M\longrightarrow M\times_S M,
\]
\textit{defined by
\[
  (g,v)\longmapsto\bigl(\operatorname{int}(g)v,v\bigr),
\]
is formally smooth.}

This is equivalent to the following assertion.

\medskip
\noindent\SGARefTarget{sga3:xi:corollary:2:4:bis}\textbf{Corollary 2.4 bis.} ---
\textit{With the notation of
\SGARefLink{sga3:xi:proposition:2:1:bis}{2.1 bis}, let \(H_1\) and
\(H_2\) be two multiplicative-type subgroups of \(G\), and let
\(\operatorname{Transp}_G(H_1,H_2)\) be the subfunctor of \(G\)
consisting of the \(g\) such that
\[
  \operatorname{int}(g)H_1=H_2.
\]
Then this functor is formally smooth over \(S\).  In particular, if
\(H_1=H_2=H\), the subfunctor
\(\operatorname{Norm}_G(H)\) of \(G\), the normalizer of \(H\), is
formally smooth over \(S\).}

The proof is analogous to that of
\SGARefLink{sga3:xi:corollary:2:4}{2.4}, again invoking \SGARefLink{sga3:ix:theorem:3:6:bis}{Exposé IX,
3.6 bis}.

\medskip
\noindent\SGARefTarget{sga3:xi:proposition:2:5}\textbf{Proposition 2.5.} ---
\textit{Let \(S\) be a prescheme, let \(G\) be an \(S\)-group
prescheme smooth over \(S\), let \(K\) be an \(S\)-group
subprescheme that is smooth over \(S\) or of multiplicative type,
let \(H\) be an \(S\)-group prescheme of multiplicative type, and let
\(u:H\to G\) be a homomorphism of \(S\)-groups.  Denote by
\(\operatorname{Transp}_G(u,K)\) the subfunctor of \(G\) whose value
at an \(S'\) over \(S\) consists of the elements \(g\in G(S')\) for
which}
\SGARefTarget{sga3:xi:display:2:5:transporter-condition}
\[
  \operatorname{int}(g)u_{S'}:
  H_{S'}\longrightarrow G_{S'}
\]
\textit{factors through \(K_{S'}\).  Then this functor is formally
smooth over \(S\).}

The proof is analogous to that of
\SGARefLink{sga3:xi:corollary:2:2}{2.2}, invoking \SGARefLink{sga3:ix:theorem:3:6}{Exposé IX, 3.6} and
\SGARefLink{sga3:x:proposition:2:1}{Exposé X, 2.1}, the latter in the case where \(K\) is of
multiplicative type.

% Exact continuation: Section 3, Expo11 local p. 10 /
% complete-reader p. 732 / source-margin p. 141.

% END INCLUDED SGA3_English_current_buildable_source_20260823/source/inputs/XI/tex/components/05_expose_XI_section_2.tex

% BEGIN INCLUDED SGA3_English_current_buildable_source_20260823/source/inputs/XI/tex/components/06_expose_XI_section_3_opening_through_remark_3_6.tex
% SGA 3, Exposé XI, English reconstruction.
% Scope: Section 3 heading through the proof following Remark 3.6,
% stopping immediately before Proposition 3.7.
% Authority: Polo--Gille Expo11.pdf, component-local pp. 10--13,
% complete-reader pp. 732--735, source-margin pp. 142--147.
% Jacob Reinhold's CC BY 4.0 Markdown lineage was used only as a phrasing
% comparison. The Polo--Gille PDF controls text, formulas, and notes.

\SGARefTarget{sga3:xi:section:3}\section*{3. Auxiliary Representability
Results}
\addcontentsline{toc}{section}{3. Auxiliary Representability Results}

\medskip
\noindent\SGARefTarget{sga3:xi:proposition:3:1}\textbf{Proposition 3.1.} ---
\textit{Let \(F\to S\) be a functor over the prescheme \(S\).  The
following conditions are equivalent:}

\begin{enumerate}
\item[(i)] \SGARefTarget{sga3:xi:proposition:3:1:item:i}
\textit{\(F\) is representable and \(F\to S\) is an open immersion
(one also says simply that \(F\to S\) is an open immersion).}

\item[(ii)] \SGARefTarget{sga3:xi:proposition:3:1:item:ii}
\textit{\(F\) is a sheaf for the faithfully flat quasi-compact
topology, ``commutes with filtered direct limits of rings,'' and
\(F\to S\) is a monomorphism.  Moreover, the following condition
holds: for every local prescheme \(S'\) over \(S\), with residue field
\(k\), and every \(S\)-morphism
\[
  \operatorname{Spec}(k)=S'_0\longrightarrow F,
\]
there exists an \(S\)-morphism \(S'\to F\) extending it.\footnote{The
French re-edition prints \(S'\to S\).  Since this morphism must extend
the displayed morphism \(S'_0\to F\), its intended target is \(F\).}}

\item[(iii)] \SGARefTarget{sga3:xi:proposition:3:1:item:iii}
\textit{Suppose that \(S\) is locally noetherian.  Then \(F\) is a
sheaf for the faithfully flat quasi-compact topology, commutes with
filtered direct limits of rings, commutes with adic inverse limits of
local artinian rings, and \(F\to S\) is a monomorphism and is
infinitesimally étale
(cf.\ \SGARefLink{sga3:xi:remark:1:8}{1.8}).}
\end{enumerate}

Let us first make two points of terminology precise.

\medskip
\noindent\SGARefTarget{sga3:xi:remark:3:2}\textbf{Remark 3.2.} ---
A functor\footnote{Editor's note: here and in \SGARefLink{sga3:xi:remark:3:3}{Remark 3.3}, the functor
is contravariant.} \(F\) over \(S\) is said to \emph{commute with
direct limits of rings} (understood to mean filtered direct limits) if
the following holds.  Let \((S'_i)_{i\in I}\) be any filtered inverse
system of affine \(S\)-schemes lying over an affine open subscheme of
\(S\), and let \(A'_i\) be their rings.  Then the natural homomorphism

\SGARefTarget{sga3:xi:display:3:2:direct-limit}
\[
  \tag{*}
  \varinjlim_i F(S'_i)\longrightarrow F(S'),
  \qquad
  S'=\operatorname{Spec}(A'),
  \quad
  A'=\varinjlim_i A'_i
\]
is bijective.

Notice that \(S'\) is precisely the inverse limit of the \(S'_i\) in
the category of preschemes, and even in the category of all ringed
spaces.  Thus the condition under consideration has the character of
a right-exactness condition---commutation with certain direct limits
in \((\mathbf{Sch})^\circ/S\)---just as does the condition of being a
sheaf for a topology.

Care must be taken that this condition is essentially relative: it
involves the morphism \(F\to S\), not merely the functor
\[
  F:(\mathbf{Sch})^\circ\longrightarrow(\mathbf{Ens}).
\]
More precisely, in \SGARefLink{sga3:xi:display:3:2:direct-limit}{\((*)\)}, \(F(S'_i)\) and \(F(S')\) mean
\(\operatorname{Hom}_S(S'_i,F)\) and
\(\operatorname{Hom}_S(S',F)\).  Consequently, when \(F\) is
representable, the condition says that \(F\) is locally of finite
presentation over \(S\).  We used several times in the preceding two
exposés the fact that a functor represented by an \(S\)-prescheme
locally of finite presentation commutes with direct limits of rings.

\medskip
\noindent\SGARefTarget{sga3:xi:remark:3:3}\textbf{Remark 3.3.} ---
A functor \(F\) over \(S\) is said to \emph{commute with adic inverse
limits of local artinian rings} if the following condition holds.
For every \(S'\) over \(S\) that is the spectrum of a complete
noetherian local ring \(A'\), set
\[
  S'_n=\operatorname{Spec}
  \bigl(A'/\operatorname{rad}(A')^{\,n+1}\bigr).
\]
Then the natural map

\SGARefTarget{sga3:xi:display:3:3:adic-limit}
\[
  \tag{**}
  F(S')\longrightarrow\varprojlim_n F(S'_n)
\]
is required to be bijective.

This condition has the character of a left-exactness condition and
is satisfied whenever \(F\) is representable.  In contrast to
commutation with direct limits of rings, it is intrinsic to \(F\) as
an object of \(\widehat{\mathbf{Sch}}\); that is, it does not involve
the morphism \(F\to S\).

\medskip
\noindent\SGARefTarget{sga3:xi:remark:3:4}\textbf{Remark 3.4.} ---
Let \(F\) be a functor over \(S\) that is a sheaf for the Zariski
topology, or, as one also says, is ``of local nature.''  It is enough
for this that \(F\) be a sheaf for a finer topology, such as the
faithfully flat quasi-compact topology.  Let \((S_i)\) be an open
covering of \(S\).  A straightforward gluing argument shows that
\(F\) is representable if and only if each
\[
  F_i=F\times_S S_i
\]
is representable; this allows one, for example, to reduce to the case
where \(S\) is affine.

Suppose that the functor \(F\), of local nature, commutes with direct
limits of rings.  Then \(F\) is representable if and only if its
restriction to the category of preschemes locally of finite
presentation over \(S\) is representable.  Necessity was noted in
\SGARefLink{sga3:xi:remark:3:2}{3.2}.  Sufficiency amounts to the
following assertion: if \(X\) is a prescheme locally of finite
presentation over \(S\), and \(X\to F\) is a morphism such that, for
every \(S'\) locally of finite presentation over \(S\), the induced
map

\SGARefTarget{sga3:xi:display:3:4:hom-comparison}
\[
  \operatorname{Hom}_S(S',X)\longrightarrow
  \operatorname{Hom}_S(S',F)
\]
is bijective, then \(X\to F\) is an isomorphism.  This follows readily
from the fact that \(X\) and \(F\) are both functors of local nature
that commute with direct limits of rings.

\medskip
\noindent\textit{Proof of \SGARefLink{sga3:xi:proposition:3:1}{Proposition 3.1}.}
The implications \SGARefLink{sga3:xi:proposition:3:1:item:i}{(i)} \(\Rightarrow\) \SGARefLink{sga3:xi:proposition:3:1:item:ii}{(ii)} and
\SGARefLink{sga3:xi:proposition:3:1:item:i}{(i)} \(\Rightarrow\) \SGARefLink{sga3:xi:proposition:3:1:item:iii}{(iii)} are immediate.  We prove the converse
implications.

Assume \SGARefLink{sga3:xi:proposition:3:1:item:ii}{(ii)}.  Let \(U\) be the set of points \(s\in S\) for which the
canonical monomorphism
\[
  \operatorname{Spec}\kappa(s)\longrightarrow S
\]
factors through \(F\).  By the last condition in \SGARefLink{sga3:xi:proposition:3:1:item:ii}{(ii)}, for every
\(s\in U\), the canonical map
\(\operatorname{Spec}\mathscr O_{S,s}\to S\), equivalently its
factorization through \(U\), factors through \(F\).  The local ring
\(\mathscr O_{S,s}\) is the direct limit of the rings of affine
neighborhoods of \(s\).  Since \(F\) commutes with direct limits of
rings, for every \(s\in U\)\footnote{The French re-edition prints
\(s\in S\) at this occurrence; the argument and the conclusion
\(U^s\subset U\) require \(s\in U\).} there is an open neighborhood
\(U^s\) such that the canonical immersion \(U^s\to S\) factors
through \(F\).  It follows that \(U^s\subset U\), and hence that
\(U\) is open.

Because \(F\to S\) is a monomorphism and \(F\) is of local nature,
the \(S\)-morphisms \(U^s\to F\) agree on the intersections
\(U^s\cap U^{s'}\), for \(s,s'\in U\), and therefore glue to an
\(S\)-morphism \(U\to F\).  It remains to prove that this morphism is
an isomorphism, or equivalently that every \(S\)-morphism \(S'\to F\)
factors uniquely through \(U\), where \(S'\) is an \(S\)-prescheme.
Since both \(F\to S\) and \(U\to S\) are monomorphisms, it is
equivalent to say that the structure morphism \(S'\to S\) factors
through \(U\).

By locality, we are reduced to the case in which \(S'\) is the
spectrum of a field and hence has a single point \(s'\).  Let \(s\)
be the point of \(S\) below \(s'\).  We claim that the given
\(S\)-morphism \(S'\to F\) factors through
\[
  \operatorname{Spec}\kappa(s)=S_0\longrightarrow F.
\]
This implies \(s\in U\) and proves the assertion.  Since
\(S'\to S_0\) is an fpqc covering and \(F\) is a sheaf for the fpqc
topology, it is enough that the two composites

\SGARefTarget{sga3:xi:display:3:1:fpqc-descent}
\[
  S''=S'\times_{S_0}S'
  \rightrightarrows S'\longrightarrow F
\]
coincide.  They do coincide because \(F\to S\) is a monomorphism.
This proves \SGARefLink{sga3:xi:proposition:3:1:item:ii}{(ii)} \(\Rightarrow\) \SGARefLink{sga3:xi:proposition:3:1:item:i}{(i)}.

Now assume \SGARefLink{sga3:xi:proposition:3:1:item:iii}{(iii)}, with \(S\) locally noetherian.  It is enough to
prove the last condition of \SGARefLink{sga3:xi:proposition:3:1:item:ii}{(ii)}; by the preceding proof, it is
enough to do so when
\[
  S'=\operatorname{Spec}\mathscr O_{S,s}
\]
for some \(s\in S\).  Put
\[
  A=\mathscr O_{S,s},
  \qquad
  A_n=A/\mathfrak m^{\,n+1},
  \qquad
  S_n=\operatorname{Spec}(A_n).
\]
Because \(F\to S\) is infinitesimally smooth, the given morphism
\(S'_0\to F\) extends to morphisms \(S_n\to F\).  Since \(F\to S\)
is a monomorphism, these morphisms form an element of
\(\varprojlim_n F(S_n)\).  Since \(F\) commutes with adic inverse
limits of local artinian rings, they arise from a morphism
\[
  \operatorname{Spec}(\widehat A)=\widehat S'\longrightarrow F.
\]
Finally, \(F\to S\) is a monomorphism, \(F\) is an fpqc sheaf, and
\(\widehat S'\to S'\) is an fpqc covering.  The last morphism
therefore descends to \(S'\to F\), completing the proof of
\SGARefLink{sga3:xi:proposition:3:1:item:iii}{(iii)} \(\Rightarrow\) \SGARefLink{sga3:xi:proposition:3:1:item:ii}{(ii)}. \hfill\(\square\)

% Source page: Expo11 local p. 12; combined-reader p. 734;
% source-margin p. 145.
\medskip
\noindent\SGARefTarget{sga3:xi:proposition:3:5}\textbf{Proposition 3.5.} ---
\textit{Let \(S\) be a locally noetherian prescheme, let \(F\to S\)
be a functor over \(S\), and let
\((X_i,u_i)_{i\in I}\) be a family of \(S\)-morphisms
\(u_i:X_i\to F\), where the \(X_i\) are preschemes locally of finite
type over \(S\).  Suppose that the following conditions hold:}

\begin{enumerate}
\item[(a)] \SGARefTarget{sga3:xi:proposition:3:5:item:a}
\textit{\(F\) is a sheaf for the faithfully flat quasi-compact
topology, commutes with direct limits of rings, and commutes with adic
inverse limits of local artinian rings.}

\item[(b)] \SGARefTarget{sga3:xi:proposition:3:5:item:b}
\textit{The morphisms \(u_i:X_i\to F\) are monomorphisms and are
infinitesimally étale
(cf.\ \SGARefLink{sga3:xi:remark:1:8}{1.8}).}

\item[(c)] \SGARefTarget{sga3:xi:proposition:3:5:item:c}
\textit{The family \((u_i)\) is ``set-theoretically surjective.''}
\end{enumerate}

\textit{Under these conditions, \(F\) is representable by a
prescheme locally of finite type over \(S\).  Moreover, the \(u_i\)
are open immersions and the family \((X_i)\) is an open covering of
\(F\).}

\medskip
\noindent\SGARefTarget{sga3:xi:remark:3:6}\textbf{Remark 3.6.} ---
Proceeding as at the end of
\SGARefLink{sga3:xi:remark:1:7}{Remark 1.7}, one defines, for every
functor
\[
  F:(\mathbf{Sch})^\circ\longrightarrow(\mathbf{Ens}),
\]
an ``underlying set'' \(\operatorname{ens}(F)\): it is the quotient
of the set of field-valued points of \(F\) by the equivalence relation
described in \SGARefLink{sga3:xi:remark:1:7}{1.7}.  When \(F\) is represented by \(X\), this recovers
the underlying set of \(X\).

Plainly, \(\operatorname{ens}(F)\) depends functorially on \(F\).
Thus, if \(G\to F\) is a morphism of functors, we say that it is
\emph{set-theoretically surjective} if the induced map
\[
  \operatorname{ens}(G)\longrightarrow\operatorname{ens}(F)
\]
is surjective.  Equivalently, every point of \(F\) with values in a
field \(k\) ``comes from'' a point of \(G\) with values in a suitable
extension of \(k\).  This definition extends immediately to a family
of morphisms \(G_i\to F\), and thereby specifies the meaning of
condition \SGARefLink{sga3:xi:proposition:3:5:item:c}{(c)}.

\medskip
\noindent\textit{Proof of \SGARefLink{sga3:xi:proposition:3:5}{Proposition 3.5}.}
For \((i,j)\in I\times I\), put

\SGARefTarget{sga3:xi:display:3:5:fiber-product}
\[
  X_{i,j}=X_i\times_F X_j
\]
and consider the projections

\SGARefTarget{sga3:xi:display:3:5:projections}
\[
  v_{i,j}:X_{i,j}\longrightarrow X_i,
  \qquad
  w_{i,j}:X_{i,j}\longrightarrow X_j.
\]
We claim that these morphisms are representable by open immersions.
Apply criterion
\SGARefLink{sga3:xi:proposition:3:1:item:iii}{3.1(iii)}.
The functor \(X_{i,j}\) satisfies the three exactness conditions:
it is a sheaf, it commutes with direct limits of rings, and it
commutes with adic inverse limits of local artinian rings.  Indeed,
\(F\), \(X_i\), and \(X_j\) satisfy these conditions, and the
conditions are stable under finite inverse limits, in particular
under fiber products.  Since every \(X_i\to F\) is a monomorphism,
both projections are monomorphisms, each being obtained from one of
these maps by base change.  Finally, being infinitesimally étale is
also preserved by base change.  Thus the hypotheses of \SGARefLink{sga3:xi:proposition:3:1:item:iii}{3.1(iii)} are
satisfied.

We may now use the \(X_i\), \(X_{i,j}\), \(v_{i,j}\), and \(w_{i,j}\)
to construct in the usual way an \(S\)-prescheme \(X\), in which the
\(X_i\) are identified with open subschemes, the \(X_{i,j}\) with the
intersections \(X_i\cap X_j\), and \(v_{i,j}\), \(w_{i,j}\) with the
canonical immersions.

Put\footnote{The French re-edition labels this coproduct \(X\) and
later writes \(R=X\times_F X\), while simultaneously using
\(R\rightrightarrows Y\) and \(X=Y/R\).  The coherent notation,
adopted here, is \(Y=\coprod_iX_i\) and \(R=Y\times_FY\).}

\SGARefTarget{sga3:xi:display:3:5:coproduct}
\[
  Y=\coprod_i X_i.
\]
Then \(X\) is also the quotient of \(Y\) by the equivalence relation

\SGARefTarget{sga3:xi:display:3:5:relation}
\[
  R=\coprod_{i,j}X_{i,j}
  \rightrightarrows Y,
\]
whose two projections \(v,w\) are defined by the \(v_{i,j}\) and
\(w_{i,j}\), respectively.

More precisely, because \(F\) is an fpqc sheaf, the morphisms
\(u_i:X_i\to F\) induce a morphism \(u:Y\to F\), and

\SGARefTarget{sga3:xi:display:3:5:kernel-relation}
\[
  R=Y\times_FY
\]
is exactly the equivalence relation defined by \(u\).  The quotient
\(X=Y/R\) is also a quotient in the category of fpqc sheaves, and
even in the category of Zariski sheaves; this follows directly from
the definitions of quotient and sheaf.  Consequently, \(u\) factors
uniquely through a morphism

\SGARefTarget{sga3:xi:display:3:5:induced-map}
\[
  \overline u:X\longrightarrow F,
\]
and \SGARefLink{sga3:xi:display:3:5:induced-map}{this morphism} is a monomorphism.

It remains to show that \(\overline u\) is an isomorphism.  Since
\(F\) is of local nature, we may assume that \(S\) is affine.  Since
\(F\) also commutes with direct limits of rings, it is enough, by
\SGARefLink{sga3:xi:remark:3:4}{3.4}, to verify that, for every
affine \(T\) of finite type over \(S\), every morphism \(T\to F\)
factors through \(X\).  Consider

\SGARefTarget{sga3:xi:display:3:5:base-change}
\[
  G=X\times_FT\longrightarrow T.
\]
It is enough to prove that \SGARefLink{sga3:xi:display:3:5:base-change}{this morphism} is an isomorphism.  Since
\(T\) is noetherian, the argument above shows that it is an open
immersion.  Here one uses that \(X\to F\) is infinitesimally étale,
which follows immediately because the induced morphisms
\(u_i:X_i\to F\) are infinitesimally étale.  By hypothesis,
\(X\to F\) is set-theoretically surjective, and this condition is
stable under base change.  Hence \(G\to T\) is set-theoretically
surjective.  Being also an open immersion, it is an isomorphism.
\hfill\(\square\)

% Exact continuation: Proposition 3.7, Expo11 local p. 13 /
% complete-reader p. 735 / source-margin p. 147.

% END INCLUDED SGA3_English_current_buildable_source_20260823/source/inputs/XI/tex/components/06_expose_XI_section_3_opening_through_remark_3_6.tex

% BEGIN INCLUDED SGA3_English_current_buildable_source_20260823/source/inputs/XI/tex/components/07_expose_XI_proposition_3_7_through_corollary_3_11.tex
% SGA 3, Exposé XI, English reconstruction.
% Scope: Proposition 3.7 through Corollary 3.11 inclusive, stopping
% immediately before Proposition 3.12.
% Authority: Polo--Gille Expo11.pdf, component-local pp. 13--17,
% complete-reader pp. 735--739, source-margin pp. 147--153.
% Jacob Reinhold's CC BY 4.0 Markdown lineage was used only as a phrasing
% comparison. The Polo--Gille PDF controls text, formulas, and diagrams.

\medskip
\noindent\SGARefTarget{sga3:xi:proposition:3:7}\textbf{Proposition 3.7.} ---
\textit{Let \(S\) be a locally noetherian prescheme, let \(I\) be a
directed partially ordered index set, let \((T_i)_{i\in I}\) be a
projective system of \(S\)-preschemes locally of finite type, let
\[
  T=\varprojlim_{i\in I}T_i
\]
be the projective-limit functor, let \(F\) be a functor over \(S\),
and let \(u:F\to T\) be an \(S\)-morphism.  Suppose that the following
conditions hold:}

\begin{enumerate}
\item[(a)] \SGARefTarget{sga3:xi:proposition:3:7:item:a}%
\textit{\(F\) is a sheaf for the faithfully flat quasi-compact
topology, commutes with inductive limits of rings, and commutes with
adic projective limits of local artinian rings.}

\item[(b)] \SGARefTarget{sga3:xi:proposition:3:7:item:b}%
\textit{The morphism \(u:F\to T\) is a monomorphism.}

\item[\(b'\)] \SGARefTarget{sga3:xi:proposition:3:7:item:b-prime}%
\textit{The morphism \(u:F\to T\) is infinitesimally étale.}

\item[(c)] \SGARefTarget{sga3:xi:proposition:3:7:item:c}%
\textit{For every point \(\xi\) of \(F\) with values in the spectrum
of a field \(k\), let
\(\xi_i\in T_i(\operatorname{Spec}k)\) be its image and let \(t_i\)
be the corresponding point of \(T_i\).  There exists \(i\in I\) such
that, for every \(j\geq i\), the transition morphism
\[
  p_{i,j}:T_j\longrightarrow T_i
\]
is étale at \(t_j\).}

\item[(d)] \SGARefTarget{sga3:xi:proposition:3:7:item:d}%
\textit{For every prescheme \(X\) locally of finite type over \(S\)
and every \(S\)-morphism \(X\to F\), the set of points \(x\in X\) at
which this morphism is infinitesimally étale is open.}
\end{enumerate}

\textit{Under these conditions, \(F\) is representable by a
prescheme locally of finite type over \(S\).}

Let us note at once that, in the case to be considered in the \SGARefLink{sga3:xi:section:4}{next
section}, conditions \SGARefLink{sga3:xi:proposition:3:7:item:c}{(c)} and \SGARefLink{sga3:xi:proposition:3:7:item:d}{(d)} will be verified by means of \SGARefLink{sga3:xi:corollary:3:8}{the
following corollary}.

\medskip
\noindent\SGARefTarget{sga3:xi:corollary:3:8}\textbf{Corollary 3.8.} ---
\textit{Assuming \SGARefLink{sga3:xi:proposition:3:7:item:a}{(a)}, \SGARefLink{sga3:xi:proposition:3:7:item:b}{(b)}, and \SGARefLink{sga3:xi:proposition:3:7:item:b-prime}{\((b')\)}, conditions \SGARefLink{sga3:xi:proposition:3:7:item:c}{(c)} and \SGARefLink{sga3:xi:proposition:3:7:item:d}{(d)} are
implied by the following conditions:}

\begin{enumerate}
\item[\(c'\)] \SGARefTarget{sga3:xi:corollary:3:8:item:c-prime}%
\textit{The \(T_i\) are smooth over \(S\), and the transition
morphisms \(p_{i,j}:T_j\to T_i\) are smooth.}

\item[\(d'\)] \SGARefTarget{sga3:xi:corollary:3:8:item:d-prime}%
\textit{For every point \(\xi\) of \(F\) with values in the spectrum
of a field \(k\), let \(t_i(\xi)\) be the point of \(T_i\) defined by
\(\xi\), let \(d_i(\xi)\) be the relative dimension of \(T_i\) over
\(S\) at \(t_i(\xi)\), and put
\[
  d(\xi)=\sup_i d_i(\xi).
\]
Then:}
  \begin{enumerate}
  \item[\(1^\circ\)] \SGARefTarget{sga3:xi:corollary:3:8:item:d-prime:1}%
  \textit{for every \(\xi\) as above,
  \(d(\xi)<+\infty\); and}

  \item[\(2^\circ\)] \SGARefTarget{sga3:xi:corollary:3:8:item:d-prime:2}%
  \textit{for every prescheme \(X\) locally of finite type over
  \(S\) and every \(S\)-monomorphism \(v:X\to F\), the function
  \[
    x\longmapsto d(\xi_x)
  \]
  on \(X\) is locally constant, where, for \(x\in X\), \(\xi_x\)
  denotes the point of \(F\) with values in \(\kappa(x)\) induced by
  \(v\).}
  \end{enumerate}
\end{enumerate}

\smallskip
\noindent\emph{Proof of Proposition
\SGARefLink{sga3:xi:proposition:3:7}{3.7}.}
Work under condition \SGARefLink{sga3:xi:proposition:3:7:item:c}{(c)}.  Let \(t\) be the element of
\(\operatorname{ens}(F)\) defined by \(\xi\) (cf.\ \SGARefLink{sga3:xi:remark:3:6}{3.6}), and put
\SGARefTarget{sga3:xi:display:3:7:local-ring}
\[
  \mathscr O_t=\varinjlim_i\mathscr O_{T_i,t_i}.
\]
Using condition \SGARefLink{sga3:xi:proposition:3:7:item:c}{(c)}, one sees readily that \(\mathscr O_t\) is a
noetherian local ring (EGA \(0_{\mathrm{IV}}\), 10.3.1.3).  Its
residue field \(\kappa(t)\) is the inductive limit of the residue
fields \(\kappa(t_i)\), and \(k\) is an extension of it.  If
\(\eta\) denotes \(\operatorname{Spec}\kappa(t)\), and \(\xi\)
denotes \(\operatorname{Spec}k\), there is a commutative diagram

\SGARefTarget{sga3:xi:diagram:3:7:residue-field-square}
\[
\begin{tikzcd}[column sep=huge,row sep=large]
  \xi \arrow[r] \arrow[d] & F \arrow[d] \\
  \eta \arrow[r] & T ,
\end{tikzcd}
\]
whose definition is evident and is left to the reader.  Since
\(\xi\to\eta\) is a covering for the fpqc topology, \(F\) is a sheaf
for that topology by \SGARefLink{sga3:xi:proposition:3:7:item:a}{(a)}, and \(F\to T\) is a monomorphism by \SGARefLink{sga3:xi:proposition:3:7:item:b}{(b)}, the
morphism \(\eta\to T\) factors uniquely through a morphism
\(\eta\to F\).  We now record \SGARefLink{sga3:xi:lemma:3:9}{the following lemma}.

\medskip
\noindent\SGARefTarget{sga3:xi:lemma:3:9}\textbf{Lemma 3.9.} ---
\textit{Under the hypotheses of
\SGARefLink{sga3:xi:proposition:3:7}{3.7}, let
\[
  T'=\operatorname{Spec}(A)
  \quad\text{and}\quad
  T'_0=\operatorname{Spec}(k(A))
\]
with \(T'\) a noetherian local scheme, and let
\(i:T'_0\to T'\) be the canonical immersion.  Suppose that a
commutative square of morphisms is given:}

\SGARefTarget{sga3:xi:diagram:3:9:local-square}
\[
\begin{tikzcd}[column sep=huge,row sep=large]
  T'_0 \arrow[r] \arrow[d,"i"'] & F \arrow[d] \\
  T' \arrow[r] & T .
\end{tikzcd}
\]
\textit{Then there exists a unique \(T\)-morphism \(T'\to F\).}

The proof is the proof of \(\SGARefLink{sga3:xi:proposition:3:1}{3.1}\), \((\mathrm{iii})\Rightarrow
(\mathrm{ii})\), in which the representability of the \(S\) occurring
in that statement was not used.  One uses the facts that \(F\) is a
sheaf for the fpqc topology, that it commutes with adic projective
limits of local artinian rings, here the rings
\[
  A/\operatorname{rad}(A)^{n+1},
\]
and that \(F\to T\) is both a monomorphism and infinitesimally étale.

We apply Lemma \SGARefLink{sga3:xi:lemma:3:9}{3.9} with
\[
  T'=\operatorname{Spec}(\mathscr O_t),
  \qquad
  T'_0=\operatorname{Spec}(\kappa(t)).
\]
We have just constructed \(T'_0\to F\), and the canonical
homomorphisms
\[
  \mathscr O_{T_i,t_i}\longrightarrow\mathscr O_t
\]
define a projective system of morphisms
\(\operatorname{Spec}(\mathscr O_t)\to T_i\), hence a canonical
morphism \(\operatorname{Spec}(\mathscr O_t)\to T\).
\SGARefTarget{sga3:xi:diagram:3:9:application-square}%
The corresponding square \(\mathrm{(x)}\) plainly commutes: by the
definition of \(T\), it is enough to check this after replacing \(T\)
by \(T_i\).  We therefore obtain a unique \(T\)-morphism

\SGARefTarget{sga3:xi:display:3:9:v-prime}
\[
  v':T'=\operatorname{Spec}(\mathscr O_t)\longrightarrow F.
\]
Since \(F\) commutes with inductive limits of rings, \SGARefLink{sga3:xi:display:3:9:v-prime}{this morphism}
factors, for sufficiently large \(i\), through

\SGARefTarget{sga3:xi:display:3:9:v-i-prime}
\[
  v'_i:T'_i=\operatorname{Spec}(\mathscr O_{T_i,t_i})
  \longrightarrow F.
\]
For such an \(i\), one has
\SGARefTarget{sga3:xi:display:3:9:local-isomorphism}
\[
  T'\xrightarrow{\sim}T'_i,
  \qquad\text{that is,}\qquad
  \mathscr O_{T_i,t_i}\xrightarrow{\sim}\mathscr O_t.
\]
Indeed, \(T'\to T'_i\) is faithfully flat and quasi-compact, hence an
effective epimorphism, so it is enough to show that it is a
monomorphism.  Suppose that two morphisms \(R\rightrightarrows T'\),
with \(R\) representable, have the same composite with
\(T'\to T'_i\).  Their composites with \(T'\to T\), which factors
through \(T'\to T'_i\), are then equal.  Hence their composites with
\(T'\to T_j\), and therefore with \(T'\to T'_j\), are equal for
every \(j\).  The two original morphisms are equal because \(T'\) is
the projective limit of the \(T'_j\) in the category
\((\mathbf{Sch})\).

Consider the diagram

\SGARefTarget{sga3:xi:diagram:3:9:factorization-triangle}
\[
\begin{tikzcd}[column sep=huge,row sep=large]
  T' \arrow[r,"v'"] \arrow[d] & F \arrow[d] \\
  T'_i \arrow[r] \arrow[ur,"v'_i"'] & T_i ,
\end{tikzcd}
\]
where the unlabeled arrows are the evident ones.  The square and the
upper triangle commute.  Since \(T'\to T'_i\) is an epimorphism, the
lower triangle also commutes.

The ring \(\mathscr O_{T_i,t_i}\) is the inductive limit of the rings
of affine open neighborhoods of \(t_i\) in \(T_i\).  Since \(F\)
commutes with inductive limits of rings, \(v'_i\) comes from a
morphism
\SGARefTarget{sga3:xi:display:3:7:v-t}
\[
  v_t:U_t\longrightarrow F
\]
defined on an open neighborhood \(U_t\) of \(t_i\) in \(T_i\).
After shrinking this neighborhood if necessary, \(v_t\) is a
\(T_i\)-morphism.  Indeed, \(T_i\), being locally of finite type over
\(S\), also commutes with inductive limits of rings.  The composite of
\(v_t\) with \(F\to T_i\) is a monomorphism, and therefore \(v_t\) is
a monomorphism.

We claim that \(v_t\) is infinitesimally étale at \(t_i\).  Since
\(F\to T\) is infinitesimally étale, this is equivalent to the
composite
\[
  T'_i\xrightarrow{v'_i}F\longrightarrow T
\]
being infinitesimally étale at \(t_i\), or to the induced morphism
\(T'_i\to T\) being infinitesimally étale.  Since
\(T'\xrightarrow{\sim}T'_i\), it is in turn equivalent to
\(T'\to T\) being infinitesimally étale.  This is immediate, because
that morphism is the projective limit of the infinitesimally étale
morphisms \(T'_j\to T_j\).

Apply condition \SGARefLink{sga3:xi:proposition:3:7:item:d}{(d)}, which has not yet been used.  It follows that
\(v_t\) is infinitesimally étale in a neighborhood of \(t_i\).
After replacing \(U_t\) by a smaller open, we may therefore suppose
that \(v_t\) is a monomorphism that is infinitesimally étale.

As \(t\) varies in \(\operatorname{ens}(F)\), the family of morphisms
\(v_t\) satisfies \SGARefLink{sga3:xi:proposition:3:5}{Proposition 3.5}, which gives the conclusion of
Proposition \SGARefLink{sga3:xi:proposition:3:7}{3.7}.

\smallskip
\noindent\emph{Proof of Corollary
\SGARefLink{sga3:xi:corollary:3:8}{3.8}.}
Assume conditions \SGARefLink{sga3:xi:corollary:3:8:item:c-prime}{\((c')\)} and \SGARefLink{sga3:xi:corollary:3:8:item:d-prime}{\((d')\)}.  With the notation of
\SGARefLink{sga3:xi:corollary:3:8:item:d-prime}{\((d')\)}, \(d_i(\xi)\) is constant for all sufficiently large \(i\).
Consequently, the relative dimension of \(T_j\) over \(T_i\) at
\(t_j\), which is
\[
  d_j(\xi)-d_i(\xi),
\]
is zero.  Thus \(T_j\to T_i\) is étale at \(t_j\), which proves
condition \SGARefLink{sga3:xi:proposition:3:7:item:c}{(c)}.

The preceding proof now gives, for each
\(t\in\operatorname{ens}(F)\), an index \(i\), an open neighborhood
\(U_t\) of \(t_i\), and a morphism \(U_t\to F\) that is a
monomorphism and is infinitesimally étale at \(t_i\).  It remains to
show that \SGARefLink{sga3:xi:display:3:7:v-t}{this morphism} is infinitesimally étale throughout a
neighborhood of \(t_i\).  In condition \SGARefLink{sga3:xi:corollary:3:8:item:d-prime}{\((d')\)}, \SGARefLink{sga3:xi:corollary:3:8:item:d-prime:2}{\(2^\circ\)}, take
\(X=U_t\).  Replacing \(U_t\), if necessary, by the connected
component of \(t_i\) in \(U_t\), we may suppose that

\SGARefTarget{sga3:xi:display:3:8:dimension-constant}
\[
  d(\xi_x)=d(\xi_{t_i})
  \qquad\text{for every }x\in U_t.
\]
Now \(d(\xi_{t_i})\) is the relative dimension of \(T_i\) over \(S\)
at \(t_i\).  Since that relative dimension remains constant on
\(U_t\), we also have, for every \(x\in U_t\),

\SGARefTarget{sga3:xi:display:3:8:relative-dimension-star}
\[
  \tag{*}
  d(\xi_x)=\text{the relative dimension of \(U_t\) over \(S\) at \(x\)}.
\]

On the other hand, using the notation of the proof of
\SGARefLink{sga3:xi:proposition:3:7}{3.7}, consider any local
noetherian prescheme
\[
  R'=\operatorname{Spec}(A),
  \qquad
  R'_0=\operatorname{Spec}(k(A)).
\]
Every morphism \(R'\to F\) whose induced point \(R'_0\to F\)
represents \(t\in\operatorname{ens}(F)\), that is, sends the closed
point of \(R'\) to \(t\), factors uniquely through \(T'\).  The proof
is the proof of \SGARefLink{sga3:xi:proposition:3:1}{3.1} or of Lemma
\SGARefLink{sga3:xi:lemma:3:9}{3.9}: one uses that \(T'\to F\) is a
monomorphism infinitesimally étale at \(t\), and that \(T'\) is an
fpqc sheaf commuting with adic projective limits.

Apply this result to the morphism
\[
  \operatorname{Spec}(\mathscr O_{U_t,x})=R'\longrightarrow F
\]
induced by \(v_t\), and to the point
\(t_x\in\operatorname{ens}(F)\) that is the image of the closed point
of \(R'\), or equivalently the image of \(x\) under \(v_t\).  We
obtain a factorization

\SGARefTarget{sga3:xi:display:3:8:local-factorization}
\[
  \operatorname{Spec}(\mathscr O_{U_t,x})
  \longrightarrow
  \operatorname{Spec}(\mathscr O_{t_x})
  \longrightarrow F.
\]
The second arrow is infinitesimally étale at \(t_x\).  To prove that
the composite is infinitesimally étale at \(x\), it is therefore
enough to prove the same for the first arrow at \(x\).

By \SGARefLink{sga3:xi:diagram:3:9:application-square}{\(\mathrm{(x)}\)}
above and formula
\SGARefLink{sga3:xi:display:3:8:relative-dimension-star}{\((*)\)},
the first arrow is induced by a local \(S\)-homomorphism between the
localizations, at \(x\) and \(t_x\), of the smooth \(S\)-preschemes
\(U_t\) and \(U_{t_x}\), which have the same relative dimension
\(d\) over \(S\) at those points.  By EGA I, 6.5.1, this local
homomorphism is induced by an \(S\)-morphism

\SGARefTarget{sga3:xi:display:3:8:w}
\[
  w:U\longrightarrow V,
\]
where \(U\) is an open neighborhood of \(x\) in \(U_t\) and
\(V=U_{t_x}\).  It remains to prove that \(w\) is étale at \(x\).
Moreover, \(U\) and \(V\) carry monomorphisms to \(F\), and, after
shrinking \(U\) once more, \(w\) is an \(F\)-morphism.  Hence \(w\)
is a monomorphism.  It is now enough to prove \SGARefLink{sga3:xi:lemma:3:10}{the following lemma}.

\medskip
\noindent\SGARefTarget{sga3:xi:lemma:3:10}\textbf{Lemma 3.10.} ---
\textit{Let \(U\) and \(V\) be two smooth \(S\)-preschemes of the
same relative dimension \(d\) over \(S\), and let \(w:U\to V\) be an
\(S\)-morphism that is a monomorphism.  Then \(w\) is an open
immersion, and hence is étale.}

By SGA 1, I, 5.7, we reduce to the case in which \(S\) is the
spectrum of a field, which may be assumed algebraically closed.  By
SGA 1, I, 5.1, it is enough to prove that \(w\) is étale, and it is
enough to do so at the closed points of \(U\).  Let \(x\) be such a
point and put \(y=w(x)\).  Choose a regular system of parameters
\(f_1,\ldots,f_d\) of \(\mathscr O_{V,y}\).  Then
\[
  A=
  \mathscr O_{U,x}\big/\sum_i f_i\mathscr O_{U,x}
\]
is the trivial extension of
\[
  k=
  \mathscr O_{V,y}\big/\sum_i f_i\mathscr O_{V,y};
\]
that is, \(A\simeq k\).  Indeed, since \(w\) is a monomorphism, so is
the structural morphism
\[
  \operatorname{Spec}(A)\longrightarrow\operatorname{Spec}(k)
\]
obtained from it by base change.  Since \(\mathscr O_{U,x}\) is a
regular local ring of dimension \(d\), the \(f_i\) form a regular
system of parameters of that ring.  It follows immediately that
\(w\) is étale at \(x\).  This proves Lemma
\SGARefLink{sga3:xi:lemma:3:10}{3.10}, and hence Corollary
\SGARefLink{sga3:xi:corollary:3:8}{3.8}.

\medskip
\noindent\SGARefTarget{sga3:xi:corollary:3:11}\textbf{Corollary 3.11.} ---
\textit{Under the hypotheses of
\SGARefLink{sga3:xi:corollary:3:8}{3.8}, for every quasi-compact open
subfunctor \(U\) of \(F\) that is separated over \(S\), there exists
\(i\in I\) such that, for every \(j\geq i\), the morphism
\[
  u_j|_U:U\longrightarrow T_j
\]
is an open immersion.  In particular, if the \(T_i\) are quasi-affine
over \(S\), then every open subfunctor \(U\) of \(F\) that is
quasi-compact over \(S\), that is, of finite type over \(S\), is
quasi-affine over \(S\).}

The proof of Proposition
\SGARefLink{sga3:xi:proposition:3:7}{3.7} shows that for every
\(t\in\operatorname{ens}(F)\) there exists \(i\in I\) such that
\(u_i:F\to T_i\) is a local isomorphism at \(t\), and then \(u_j\)
is a local isomorphism at that point for every \(j\geq i\).  By
quasi-compactness, \(i\) may be chosen independently of \(x\in U\).
It remains to prove that, for sufficiently large \(i\),
\[
  u_i|_U:U\longrightarrow T_i
\]
is a monomorphism.

Since \(U\to T\) is a monomorphism, the intersection of the
equivalence relations
\[
  U\times_{T_i}U\ \subset\ U\times_S U
\]
is the diagonal.  The prescheme \(U\times_S U\) is noetherian, and
the \(U\times_{T_i}U\) are closed subpreschemes.  It follows that one
of them is already the diagonal; equivalently, \(u_i|_U\) is a
monomorphism.  This proves the first assertion of Corollary
\SGARefLink{sga3:xi:corollary:3:11}{3.11}, and the second follows
immediately.

% Exact continuation: Proposition 3.12, Expo11 local p. 17 /
% complete-reader p. 739 / source-margin p. 153.

% END INCLUDED SGA3_English_current_buildable_source_20260823/source/inputs/XI/tex/components/07_expose_XI_proposition_3_7_through_corollary_3_11.tex

% BEGIN INCLUDED SGA3_English_current_buildable_source_20260823/source/inputs/XI/tex/components/08_expose_XI_proposition_3_12_and_remark_3_13.tex
% SGA 3, Exposé XI, English reconstruction.
% Scope: Proposition 3.12 and Remark 3.13, including the complete direct
% proof, stopping immediately before Section 4.
% Authority: Polo--Gille Expo11.pdf, component-local pp. 17--20,
% complete-reader pp. 739--742, printed pp. 154--156.
% Jacob Reinhold's CC BY 4.0 Markdown lineage was used only as a phrasing
% comparison. The Polo--Gille PDF controls text and formulas.

\medskip
\noindent\SGARefTarget{sga3:xi:proposition:3:12}\textbf{Proposition 3.12.} ---
\textit{Let \(S\) be a prescheme and \(G\) an affine \(S\)-group
prescheme.}

\begin{enumerate}
\item[\SGARefTarget{sga3:xi:proposition:3:12:item:a}\textit{(a)}]
\textit{Let
\[
  F:(\mathbf{Sch})^\circ/S\longrightarrow(\mathbf{Ens})
\]
be the functor such that, for every \(T\) over \(S\),}
\SGARefTarget{sga3:xi:display:3:12:finite-multiplicative-subgroups}
\[
  F(T)=
  \left\{
  \begin{array}{c}
  \text{multiplicative-type subgroups of \(G_T\)}\\
  \text{that are finite over \(T\)}
  \end{array}
  \right\}.
\]
\textit{Suppose that \(S\) is locally noetherian or that \(G\) is of
finite presentation over \(S\).  Then \(F\) is representable and
affine over \(S\).  If \(G\) is of finite presentation over \(S\),
then \(F\) is locally of finite presentation over \(S\).}

\item[\SGARefTarget{sga3:xi:proposition:3:12:item:b}\textit{(b)}]
\textit{Let \(H\) be a finite \(S\)-group prescheme of multiplicative
type.  Then
\(\operatorname{Hom}_{S\text{-}\mathrm{gr}}(H,G)\) is representable
and affine over \(S\).  If \(G\) is of finite type (respectively, of
finite presentation) over \(S\), then the same is true of
\(\operatorname{Hom}_{S\text{-}\mathrm{gr}}(H,G)\).}
\end{enumerate}

\medskip
\noindent\SGARefTarget{sga3:xi:remark:3:13}\textbf{Remark 3.13.} ---
Apart from the assertion that
\(\operatorname{Hom}_{S\text{-}\mathrm{gr}}\) is affine, and from the
case in which \(G\) is of finite presentation over \(S\), which is
the case that will suffice for us,
\SGARefLink{sga3:xi:proposition:3:12}{3.12} is an immediate
consequence of the theory of Hilbert schemes (A.~Grothendieck,
\textit{Techniques de construction et théorèmes d'existence en
Géométrie Algébrique: IV. Les Schémas de Hilbert}, Séminaire
Bourbaki, May 1961, no.~221).  It is enough, in fact, that \(G\) be
quasi-projective over \(S\).  In case \textit{\SGARefLink{sga3:xi:proposition:3:12:item:a}{(a)}} one can also
represent the larger functor
\[
  F_0(T)=
  \left\{
  \begin{array}{c}
  \text{group subpreschemes of \(G_T\) that are flat, proper,}\\
  \text{and of finite presentation over \(T\)}
  \end{array}
  \right\}.
\]
The canonical monomorphism \(F\to F_0\) is an ``open immersion,'' as
follows from the criterion in \SGARefLink{sga3:xi:proposition:3:1}{3.1} and \SGARefLink{sga3:x:corollary:4-7}{Exposé X, 4.7}~\textit{(b)}.
Thus the representability of \(F_0\) implies that \(F\) is
representable by an open subprescheme of \(F_0\).  In case \textit{\SGARefLink{sga3:xi:proposition:3:12:item:b}{(b)}}
it is enough to assume that \(H\) is projective and of finite
presentation over \(S\).  In both cases one obtains a functor locally
of finite presentation over \(S\).

In the present Exposé, \SGARefLink{sga3:xi:proposition:3:12}{3.12} is
only a technical lemma used to prove a key result in the \SGARefLink{sga3:xi:section:4}{next
section}.  We shall therefore sketch an easy direct proof that does
not use Hilbert schemes.

We first prove \textit{\SGARefLink{sga3:xi:proposition:3:12:item:b}{(b)}}.  It is enough to prove that
\(\operatorname{Hom}_S(H,G)\) is representable, independently of any
group structure on \(H\) and \(G\), and that it has the additional
properties asserted for
\(\operatorname{Hom}_{S\text{-}\mathrm{gr}}(H,G)\).  Indeed, applying
the same result to \(\operatorname{Hom}_S(H\times_S H,G)\), the
subfunctor
\(\operatorname{Hom}_{S\text{-}\mathrm{gr}}(H,G)\) of
\(\operatorname{Hom}_S(H,G)\) can be expressed by a finite inverse
limit---in fact, by fiber products---involving \(G\),
\(\operatorname{Hom}_S(H,G)\), and
\(\operatorname{Hom}_S(H\times_S H,G)\).  We leave this expression
to the reader.

Write
\[
  H=\operatorname{Spec}(B),
\]
where \(B\) is a sheaf of algebras on \(S\) that is locally free as a
sheaf of modules; this is the only hypothesis on \(H\) that the
argument uses.  Since the representability question is local on
\(S\), we may suppose that \(S\) is affine, with ring \(A\).  We may
also write \(G=\operatorname{Spec}(C)\), where \(C\) is an
\(A\)-algebra.

When \(G=S[t]=G_0\), with \(t\) an indeterminate, the Hom functor is
simply
\[
  T\longmapsto\Gamma(T,B_T).
\]
Since \(B\) is locally free, this functor is represented by the
vector bundle \(V(B^\vee)\), where \(B^\vee\) is the dual module
sheaf.  When \(G=S[(t_i)_{i\in I}]=G_0^I\), for a not necessarily
finite family of indeterminates, it follows that\footnote{The source
display prints
\(\operatorname{Hom}_S(H,G)^I=V(B^\vee)^I=V(B^{\vee(I)})\).
Since the preceding sentence sets \(G=G_0^I\), the displayed equality
requires the explicit factor
\(\operatorname{Hom}_S(H,G)=\operatorname{Hom}_S(H,G_0)^I\).}
\SGARefTarget{sga3:xi:display:3:13:hom-affine-space}
\[
  \operatorname{Hom}_S(H,G)
  =\operatorname{Hom}_S(H,G_0)^I
  =V(B^\vee)^I
  =V\bigl(B^{\vee(I)}\bigr).
\]

In general \(G\) is isomorphic to a closed subscheme of a scheme of
the form \(S[(t_i)_{i\in I}]\); equivalently, \(C\) is a quotient of
an \(A\)-algebra \(A[(t_i)_{i\in I}]\).  Let \((F_j)\) be a family of
generators of the ideal by which one takes the quotient.  After
covering \(S\) by smaller affine open subsets, we may assume that
\(B\) is free of rank \(n\).  Choose a basis
\((e_k)_{1\leq k\leq n}\) of \(B\).  Writing
\[
  x_i=\sum_{k=1}^{n}x_{ik}e_k
\]
with indeterminate coefficients \(x_{ik}\) in an unspecified
\(A\)-algebra \(A'\), expand \(F_j((x_i))\) in this basis.  For every
\(F_j\) this gives \(n\) polynomials
\[
  F_{j,k}\qquad(1\leq k\leq n)
\]
in the variables
\((x_{i,k})_{i\in I,\,1\leq k\leq n}\), with coefficients in \(A\).
It follows at once that \(\operatorname{Hom}_S(H,G)\) is represented
by the spectrum of the quotient of the polynomial ring
\(A[(x_{i,k})]\) by the ideal generated by the \(F_{j,k}\).  This
proves \textit{\SGARefLink{sga3:xi:proposition:3:12:item:b}{(b)}}.

We now prove \textit{\SGARefLink{sga3:xi:proposition:3:12:item:a}{(a)}}.  For every ordinary finite commutative
group \(M\), let \(F_M\) be the subfunctor of \(F\) obtained by
restricting to those subgroups of \(G_T\) that are of multiplicative
type and of type \(M\).  A gluing argument like the one used in \SGARefLink{sga3:xi:proposition:3:5}{3.5}
(which we should have stated after decomposing the problem a little
further) shows that it is enough to prove that each \(F_M\) is
representable.  Then \(F\) is represented by the disjoint sum of the
\(F_M\), with \(M\) ranging over the isomorphism classes of finite
commutative groups.  Indeed, \(F\) is the corresponding sum in the
category of sheaves.

Fix \(M\) and write \(F\) in place of \(F_M\).  Put
\(H=D_S(M)\), and let
\[
  F'=\operatorname{Imm}_{S\text{-}\mathrm{gr}}(H,G)
  \subset\operatorname{Hom}_{S\text{-}\mathrm{gr}}(H,G)
\]
be the subfunctor whose value at \(T\) is the set of homomorphisms of
\(T\)-groups \(H_T\to G_T\) that are closed immersions.  By
\textit{\SGARefLink{sga3:xi:proposition:3:12:item:b}{(b)}},
\(\operatorname{Hom}_{S\text{-}\mathrm{gr}}(H,G)\) is represented by
an affine \(S\)-prescheme.  Using \SGARefLink{sga3:ix:theorem:6:8}{Exposé IX, 6.8}, one sees immediately
that \(F'\) is represented by an open-and-closed subprescheme of it;
in particular, \(F'\) is affine over \(S\).

Consider the canonical morphism \(F'\to F\), which assigns to every
monomorphism \(H_T\to G_T\) its image group.  Directly from the
definitions, \(F'\) is thereby a principal homogeneous bundle over
\(F\), in the category of fpqc sheaves, under
\[
  \Gamma_F=\Gamma_S\times_S F,
  \qquad
  \Gamma=\operatorname{Aut}_{\mathrm{gr}}(M).
\]
It follows by descent that \(F'\to F\) is representable: for every
morphism \(T\to F\) with \(T\) representable, the prescheme
\(T\times_F F'\) is representable over \(T\), in fact by a principal
Galois bundle under \(\Gamma\).

By \SGARefLink{sga3:iv:theorem:4:6:6}{Exposé IV, 4.6.6}, \(F\) is therefore representable if and only if
the quotient \(F'/F''\), where \(F''=F'\times_F F'\), exists in
\((\mathbf{Sch})\) and is universally effective for faithfully flat
quasi-compact morphisms.  Equivalently, the quotient \(F'/\Gamma\)
must exist and be universally effective for such morphisms.  Since
\(F'\) is affine over \(S\), \SGARefLink{thm:sga3-V-4.1}{Exposé V, 4.1} shows that this condition
is satisfied.  This proves the representability of \(F\) in
\textit{\SGARefLink{sga3:xi:proposition:3:12:item:a}{(a)}}.

Finally, suppose that \(G\) is of finite presentation over \(S\).
The additional assertion follows immediately from the preceding
proof: by \textit{\SGARefLink{sga3:xi:proposition:3:12:item:b}{(b)}}, \(F'\) is then locally of finite presentation
over \(S\), so that Exposé V may be applied.\footnote{Editor's note:
the source marks this final reference to Exposé V as needing
verification or further specification.}

% Exact continuation: Section 4, Expo11 local p. 20 /
% complete-reader p. 742 / printed p. 157.

% END INCLUDED SGA3_English_current_buildable_source_20260823/source/inputs/XI/tex/components/08_expose_XI_proposition_3_12_and_remark_3_13.tex

% BEGIN INCLUDED SGA3_English_current_buildable_source_20260823/source/inputs/XI/tex/components/09_expose_XI_section_4.tex
% SGA 3, Exposé XI, English reconstruction.
% Scope: all of §4, from its heading and Theorem 4.1 through Remark 4.6;
% stopping immediately before §5.
% Authority: Polo--Gille Expo11.pdf, component-local pp. 20--25,
% complete-reader pp. 742--747, source-margin pp. 157--164.
% Jacob Reinhold's CC BY 4.0 Markdown lineage was used only as a phrasing
% comparison. The Polo--Gille PDF controls text, formulas, and numbering.

\SGARefTarget{sga3:xi:section:4}
\section*{4. The Scheme of Multiplicative-Type Subgroups\\
of a Smooth Affine Group}
\addcontentsline{toc}{section}{4. The Scheme of Multiplicative-Type
Subgroups of a Smooth Affine Group}

The principal result of the present Exposé is \SGARefLink{sga3:xi:theorem:4:1}{the following}.

\medskip
\noindent\SGARefTarget{sga3:xi:theorem:4:1}\textbf{Theorem 4.1.} ---
\textit{Let \(S\) be a prescheme, let \(G\) be an \(S\)-group
prescheme smooth and affine over \(S\), and let}
\[
  F\colon(\mathbf{Sch})^\circ/S\longrightarrow(\mathbf{Ens})
\]
\textit{be the functor defined by}
\SGARefTarget{sga3:xi:display:4:1:functor}
\[
  F(T)=\{\text{multiplicative-type subgroups of \(G_T\)}\}.
\]
\textit{Then \(F\) is representable, smooth, and separated over
\(S\).}

Let us record at once \SGARefLink{sga3:xi:corollary:4:2}{the following variant}.

\medskip
\noindent\SGARefTarget{sga3:xi:corollary:4:2}\textbf{Corollary 4.2.} ---
\textit{Let \(G\) and \(H\) be two \(S\)-group preschemes, with
\(G\) smooth and affine over \(S\), and \(H\) of multiplicative type
and of finite type.  Then
\(\operatorname{Hom}_{S\text{-}\mathrm{gr}}(H,G)\) is representable,
smooth, and separated over \(S\).}

Indeed, when \(H\) is smooth over \(S\), so is \(H\times_S G\), and
one may apply \SGARefLink{sga3:xi:theorem:4:1}{4.1} to the latter.
By considering graph subgroups,
\(\operatorname{Hom}_{S\text{-}\mathrm{gr}}(H,G)\) becomes a
subfunctor \(F'\) of the functor \(F\) of multiplicative-type
subgroups of \(H\times_S G\): namely, \(F'(T)\) is the set of
multiplicative-type subgroups \(K\) of
\[
  (H\times_S G)_T=H_T\times_TG_T
\]
such that the homomorphism induced by the first projection,
\SGARefTarget{sga3:xi:display:4:2:first-projection}
\[
  K\longrightarrow H_T,
\]
is an isomorphism.  By \SGARefLink{sga3:ix:corollary:2:9}{Exposé IX, 2.9}, the canonical morphism
\(F'\to F\) is an open immersion.  Since \(F\) is representable,
the same is true of \(F'\), which is represented by an open
subprescheme of \(F\); and since \(F\) is separated and smooth over
\(S\), so is \(F'\).

When \(H\) is finite over \(S\), it is enough to apply
\SGARefLink{sga3:xi:proposition:3:12:item:b}{3.12(b)} to obtain the
representability of
\(\operatorname{Hom}_{S\text{-}\mathrm{gr}}(H,G)\).  When
\(H=H_1\times_SH_2\), with \(H_1\) smooth over \(S\) and \(H_2\)
finite over \(S\), then
\(\operatorname{Hom}_{T\text{-}\mathrm{gr}}(H_T,G_T)\) is identified
with the subset of
\SGARefTarget{sga3:xi:display:4:2:product-hom}
\[
 \operatorname{Hom}_{T\text{-}\mathrm{gr}}\bigl((H_1)_T,G_T\bigr)
 \times
 \operatorname{Hom}_{T\text{-}\mathrm{gr}}\bigl((H_2)_T,G_T\bigr)
\]
formed by the pairs \((u_1,u_2)\) such that
\(u_1\mathbin{\cdot}u_2=u_2\mathbin{\cdot}u_1\).  It follows that
\(\operatorname{Hom}_{S\text{-}\mathrm{gr}}(H,G)\) is represented
by a closed subprescheme of
\SGARefTarget{sga3:xi:display:4:2:representing-product}
\[
 \operatorname{Hom}_{S\text{-}\mathrm{gr}}(H_1,G)
 \times_S
 \operatorname{Hom}_{S\text{-}\mathrm{gr}}(H_2,G)=X.
\]
Indeed, apply \SGARefLink{sga3:viii:examples:6:5:item:b:01}{Exposé VIII, 6.5(b)}, with \(Y=H\), \(Z=G\), and with
\(q_1\) and \(q_2\) defined respectively by
\[
  (u_1,u_2)\longmapsto u_1\mathbin{\cdot}u_2
  \qquad\text{and}\qquad
  (u_1,u_2)\longmapsto u_2\mathbin{\cdot}u_1.
\]

In the general case, the question is local on \(S\) for the Zariski
topology, so we may suppose that \(S\) is affine and that \(H\) has
constant type over \(S\).  Since \(H\) is quasi-isotrivial
(\SGARefLink{sga3:x:corollary:4-5}{Exposé X, 4.5}), there is a surjective étale morphism \(S'\to S\),
with \(S'\) affine, that splits \(H\), that is, such that
\(H'=H_{S'}\) is diagonalizable.  \SGARefLink{sga3:xi:corollary:4:2}{The preceding result} then applies,
because a diagonalizable group is the product of a diagonalizable
torus and a finite diagonalizable group.  It remains to see that the
descent datum obtained on the \(S'\)-prescheme
\[
  \operatorname{Hom}_{S'\text{-}\mathrm{gr}}(H',G')=X'
\]
is effective.  This follows from the argument of \SGARefLink{sga3:x:lemma:5-4}{Exposé X, 5.4},
noting that, there, the hypothesis that \(X'\to S'\) be separated
and locally quasi-finite was used only to ensure that every open
part of \(X'\) quasi-compact over \(S'\) was quasi-affine over
\(S'\).  That property still holds here, as follows readily from
the fact that it holds for the functor \(F\) of
\SGARefLink{sga3:xi:theorem:4:1}{4.1} (as will be seen during the
proof of \SGARefLink{sga3:xi:theorem:4:1}{4.1}).  This argument---or
any similar variant of this small dévissage---establishes the
representability of
\(\operatorname{Hom}_{S\text{-}\mathrm{gr}}(H,G)\), and at the same
time shows that it is separated over \(S\).  It is locally of finite
presentation over \(S\), as follows, for example, from the fact
already noted in \SGARefLink{sga3:xi:remark:3:2}{3.2} that this
functor ``commutes with filtered direct limits of rings.''  Finally,
the functor is formally smooth over \(S\) by
\SGARefLink{sga3:xi:proposition:2:1}{2.1}, and is therefore smooth
over \(S\).

\medskip
\noindent\SGARefTarget{sga3:xi:remark:4:3}\textbf{Remark 4.3.} ---
We have deduced \SGARefLink{sga3:xi:corollary:4:2}{4.2} from
\SGARefLink{sga3:xi:theorem:4:1}{4.1}; this is genuinely immediate
only when \(H\) is also smooth over \(S\).  To make the deduction
without contortions in the general case, one would need the
representability assertion in
\SGARefLink{sga3:xi:theorem:4:1}{4.1} without assuming that \(G\) is
smooth over \(S\), assuming only that it is affine and of finite
presentation over \(S\).  Of course, \(F\) would then no longer be
smooth over \(S\) in general.  There is little doubt that
\SGARefLink{sga3:xi:theorem:4:1}{4.1} remains true under these more
general hypotheses, but the proof appears likely to be more delicate,
because one can no longer invoke
\SGARefLink{sga3:xi:corollary:3:8}{3.8}.\footnote{This is indeed
proved for \(G\) flat and quasi-affine over \(S\), with connected
fibres, provided one restricts to the central subtori of \(G\)
(\SGARefLink{sga3:xv:theorem:8-8}{Exposé XV, 8.8}).}\footnote{Editor's note: when \(H\) is smooth (not
necessarily affine) over a normal locally noetherian base \(S\),
M.~Raynaud showed that the largest representable open subfunctor of
the functor of subtori of \(H\) is a disjoint sum of opens smooth and
affine over \(S\).  This is Theorem IX.9.26 in \emph{Faisceaux amples
sur les schémas en groupes et sur les espaces homogènes}, Lecture
Notes in Mathematics 119 (1970).}

Let us nevertheless point out that when \(G\) is a closed subgroup
of an affine smooth group \(G_0\) over \(S\), the functor \(F\) of
multiplicative-type subgroups of \(G\) is represented by a closed
subprescheme of the prescheme representing the analogous functor
\(F_0\) for \(G_0\), which is covered by
\SGARefLink{sga3:xi:theorem:4:1}{4.1}; this follows readily by
applying \SGARefLink{sga3:viii:theorem:6:4}{Exposé VIII, 6.4}.  This also raises the question whether an
affine \(S\)-group scheme \(G\), affine and of finite presentation
over \(S\), is isomorphic to a subgroup scheme of
\(\operatorname{GL}(n)_S\) for some suitable \(n\).  This is true
when \(S\) is the spectrum of a field (cf.\ Exposé
\(\mathrm{VI}_{\mathrm B}\), 11.11), but is unfortunately false in
general, even for tori; see
\SGARefLink{sga3:xi:remark:4:6}{4.6}.  Finally, one could also prove
\SGARefLink{sga3:xi:corollary:4:2}{4.2} directly by exactly the same
method as \SGARefLink{sga3:xi:theorem:4:1}{4.1}.

We now prove \SGARefLink{sga3:xi:theorem:4:1}{4.1}.  Since the
functor \(F\) is evidently local in nature, we may suppose that
\(S\) is affine, say \(S=\operatorname{Spec}(A)\), where \(A\) is a
ring.  Regarding \(A\) as the filtered direct limit of its subrings
of finite type over \(\mathbb Z\), and noting that \(G\) descends
from a smooth affine group over one such subring (EGA IV, 8), we are
reduced to the case in which \(S\) is noetherian.

For every integer \(n>0\), let \(T_n\) be the functor defined like
\(F\), but restricted to multiplicative-type subgroups \(H\) of
\(G_T\) such that \(n\cdot\operatorname{id}_H=0\), that is, such
that \({}_nH=H\).  Order the set \(I\) of positive integers by
divisibility.  When \(m\) is a multiple of \(n\), define
\SGARefTarget{sga3:xi:display:4:1:p-n-m}
\[
  p_{n,m}\colon T_m\longrightarrow T_n
\]
by
\[
  p_{n,m}(H)={}_nH=\operatorname{Ker}
  \bigl(n\cdot\operatorname{id}_H\bigr).
\]
In this way the \(T_n\) form an inverse system of functors over
\(S\).  By
\SGARefLink{sga3:xi:proposition:3:12:item:a}{3.12(a)}, the functors
\(T_n\) are representable, affine, and of finite type over \(S\).
Similarly define morphisms
\SGARefTarget{sga3:xi:display:4:1:u-n}
\[
  u_n\colon F\longrightarrow T_n
\]
by \(u_n(H)={}_nH\).  Thus one obtains a morphism
\SGARefTarget{sga3:xi:display:4:1:u-inverse-limit}
\[
  u\colon F\longrightarrow T=\varprojlim_n T_n,
\]
where the inverse limit is taken in the category of functors over
\(S\).  Since the \(T_n\) are representable and affine over \(S\),
so is \(T\): it is the spectrum of the filtered direct limit of the
quasi-coherent algebras on \(S\) defining the \(T_n\).  Of course,
\(T\) is not in general of finite type over \(S\).

We shall apply \SGARefLink{sga3:xi:proposition:3:7}{3.7}.  It is
enough to verify conditions \textit{\SGARefLink{sga3:xi:proposition:3:7:item:a}{(a)}}--\textit{\SGARefLink{sga3:xi:proposition:3:7:item:d}{(d)}} there; these
will imply that \(F\) is represented by a prescheme locally of finite
type over \(S\).  It then follows from
\SGARefLink{sga3:xi:proposition:2:1:bis}{2.1 bis} that \(F\) is even
smooth over \(S\).  Moreover, since \(F\) is a subfunctor of \(T\),
which is affine over \(S\), \(F\) is separated over \(S\): it is
separated over \(T\), and \(T\) is separated over \(S\).  We can
also prove at once the supplementary assertion used above: every
open \(U\subset F\) that is quasi-compact over \(S\) is quasi-affine
over \(S\).  This follows from
\SGARefLink{sga3:xi:corollary:3:11}{3.11} and the fact that the
\(T_n\) are affine over \(S\).

Let us now verify the conditions of
\SGARefLink{sga3:xi:proposition:3:7}{3.7}.

\smallskip
\noindent\SGARefTarget{sga3:xi:theorem:4:1:condition:a}\textit{(a)}
\(F\) is a sheaf for the faithfully flat quasi-compact topology, by
the descent theory of SGA~1, Exposé VIII, which applies here because
groups of multiplicative type over \(T\) are affine over
\(T\).\footnote{The printed source says ``affine over \(S\)'' here;
the descent assertion concerns groups over \(T\), so the required
base is \(T\).} It
commutes with filtered direct limits of rings by the standard
machinery of EGA IV, 8.  Let us show that it commutes with adic
inverse limits of local artinian rings.  For the analogous functor
considered in \SGARefLink{sga3:xi:corollary:4:2}{4.2}, this is
precisely \SGARefLink{sga3:ix:theorem:7:1}{Exposé IX, 7.1}, in the special case where \(A\) is a
complete noetherian local ring endowed with an ideal of definition
for its usual topology.  Notice that this is exactly where the
hypothesis that \(G\) be affine enters essentially.  In the present
case we are reduced to \SGARefLink{sga3:xi:lemma:4:4}{the following lemma}.

\medskip
\noindent\SGARefTarget{sga3:xi:lemma:4:4}\textbf{Lemma 4.4.} ---
\textit{Let \(A\) be a complete noetherian local ring with maximal
ideal \(\mathfrak m\), let \(G\) be a group scheme affine over
\(S=\operatorname{Spec}(A)\), and, for every integer \(n\geq0\), put}
\SGARefTarget{sga3:xi:display:4:4:adic-thickenings}
\[
  S_n=\operatorname{Spec}\bigl(A/\mathfrak m^{\,n+1}\bigr),
  \qquad
  G_n=G\times_SS_n.
\]
\textit{For every \(n\), let \(H_n\) be a multiplicative-type
subgroup of finite type of \(G_n\), and suppose that, whenever
\(m\geq n\), \(H_n\) is obtained from \(H_m\) by reduction.  Then
there exists a unique multiplicative-type subgroup \(H\) of \(G\)
whose reductions are the \(H_n\).}

By \SGARefLink{sga3:x:theorem:3:2}{Exposé X, 3.2}, there exists a multiplicative-type group \(H\) over
\(S\), necessarily of finite type and isotrivial and determined up
to unique isomorphism, together with an isomorphism
\[
  H\times_SS_0\simeq H_0.
\]
Here one uses the fact that \(H_0\) is isotrivial, since it is of
finite type over a field; see \SGARefLink{sga3:x:proposition:1:4}{Exposé X, 1.4}.  By \SGARefLink{sga3:x:proposition:2:1}{Exposé X, 2.1}, for
every \(n\) the isomorphism \(H\times_SS_0\simeq H_0\) lifts to a
unique isomorphism \(H\times_SS_n\simeq H_n\).  \SGARefLink{sga3:ix:theorem:7:1}{Exposé IX, 7.1},
already cited, now shows that the homomorphisms \(H_n\to G_n\) arise
from a unique homomorphism of \(S\)-groups
\SGARefTarget{sga3:xi:display:4:4:u}
\[
  u\colon H\longrightarrow G.
\]
By \SGARefLink{sga3:ix:corollary:6:6}{Exposé IX, 6.6}, \SGARefLink{sga3:xi:display:4:4:u}{this homomorphism} is a monomorphism because
\(u_0\colon H_0\to G_0\) is one.  This proves
\SGARefLink{sga3:xi:lemma:4:4}{4.4}.

\smallskip
\noindent\SGARefTarget{sga3:xi:theorem:4:1:condition:b}\textit{(b)}
The morphism \(u\colon F\to T\) is a monomorphism.  This follows from
the density theorem, \SGARefLink{sga3:ix:theorem:4:7}{Exposé IX, 4.7}, in the form of \SGARefLink{sga3:ix:corollary:4:8:item:b}{Corollary
4.8(b)}.  It is essential for the present application that this
result is available over a base that need not be noetherian.  Since
the functor \(T\) does not commute with filtered direct limits of
rings, one cannot reduce to the noetherian case a priori.

\smallskip
\noindent\SGARefTarget{sga3:xi:theorem:4:1:condition:b-prime}
\textit{(\(b'\))} The morphism \(u\colon F\to T\) is infinitesimally
étale; in other words, we have \SGARefLink{sga3:xi:lemma:4:5}{the following lemma}.

\medskip
\noindent\SGARefTarget{sga3:xi:lemma:4:5}\textbf{Lemma 4.5.} ---
\textit{Let \(A\) be an artinian local ring with residue field
\(k\), let \(S=\operatorname{Spec}(A)\), let
\(I\subset\operatorname{rad}(A)\) be an ideal, and put
\(S_0=\operatorname{Spec}(A/I)\).  Let \(G\) be an \(S\)-group
prescheme smooth over \(S\), and put \(G_0=G\times_SS_0\).  Suppose
given a multiplicative-type subgroup \(H_0\) of \(G_0\) and, for
every integer \(n>0\), a multiplicative-type subgroup \(H(n)\) of
\(G\), in such a way that}
\begin{enumerate}
\item[\SGARefTarget{sga3:xi:lemma:4:5:item:1}\(1^\circ\)]
\textit{for every multiple \(m\) of \(n\),
\(H(n)={}_nH(m)\); and}
\item[\SGARefTarget{sga3:xi:lemma:4:5:item:2}\(2^\circ\)]
\textit{\(H(n)_0={}_nH_0\).}
\end{enumerate}
\textit{Then there exists a unique multiplicative-type subgroup
\(H\) of \(G\) such that \({}_nH=H(n)\) for every \(n\).}

Uniqueness is already contained in \textit{\SGARefLink{sga3:xi:theorem:4:1:condition:b}{(b)}}.  For existence, an
immediate induction reduces us to the case in which
\(I\mathfrak m=0\), where \(\mathfrak m\) is the maximal ideal of
\(A\).\footnote{The printed source abruptly writes
\(J\mathfrak m=0\) here and uses \(J\) throughout the ensuing
argument, although only \(I\) was introduced in the statement.  We
retain \(I\), the same ideal, throughout.} Put
\[
  k=A/\mathfrak m,\qquad
  S_0=\operatorname{Spec}(k),\qquad
  G_0=G\times_SS_0,
\]
and let \(\mathfrak g_0\) and \(\mathfrak h_0\) be the Lie algebras
of \(G_0\) and \(H_0\), respectively.  There is a canonical
isomorphism of groups
\SGARefTarget{sga3:xi:display:4:5:infinitesimal-kernel}
\[
  \mathfrak g_0\otimes_kI
  \simeq
  \operatorname{Ker}\bigl(G(S)\longrightarrow G(S_0)\bigr);
\]
see Exposé III.  By \SGARefLink{sga3:ix:theorem:3:6:bis}{Exposé IX, 3.6 bis} and \SGARefLink{sga3:xi:proposition:3:7}{3.7}, there exists a
multiplicative-type subgroup \(H\) of \(G\) reducing to \(H_0\), and
such an \(H\) is determined, up to an inner automorphism by an
element of
\[
  N=\operatorname{Ker}\bigl(G(S)\longrightarrow G(S_0)\bigr).
\]
Thus the set \(P\) of liftings \(H\) of \(H_0\) is a principal
homogeneous set under \(N/M\), where
\[
  M=\{g\in N\mid \operatorname{int}(g)\mathbin{\cdot}H=H\}.
\]

One sees readily (Exposé III) that this subgroup is precisely the
vector subspace of \(\mathfrak g_0\otimes_kI\) consisting of the
invariants under \(H_0\), where \(H_0\) acts on
\(\mathfrak g_0\otimes_kI\) through the representation induced by
the adjoint representation of \(G_0\).  Likewise, the set \(P(n)\)
of liftings of \(H(n)_0={}_nH_0\) is a principal homogeneous set
under \(N/M(n)\), where \(M(n)\) is the vector subspace of
\(\mathfrak g_0\otimes_kI=N\) consisting of the invariants under
\({}_nH_0\).  The density theorem, \SGARefLink{sga3:ix:theorem:4:7}{Exposé IX, 4.7}, shows readily
that \(M=M(n)\) for sufficiently large \(n\), where ``large'' is
understood with respect to the divisibility order on the positive
integers.  Consequently, the natural map
\[
  H\longmapsto{}_nH
\]
from \(P\) to \(P(n)\), which is compatible with the actions of
\(N\), and hence with the homomorphism
\[
  N/M\longrightarrow N/M(n)
\]
between the acting groups, is bijective for sufficiently large
\(n\).  The conclusion of
\SGARefLink{sga3:xi:lemma:4:5}{4.5} follows immediately.

To verify conditions \textit{\SGARefLink{sga3:xi:proposition:3:7:item:c}{(c)}} and \textit{\SGARefLink{sga3:xi:proposition:3:7:item:d}{(d)}} of
\SGARefLink{sga3:xi:proposition:3:7}{3.7}, we use
\SGARefLink{sga3:xi:corollary:3:8}{3.8}, which reduces us to checking
\textit{(\(c'\))} and \textit{(\(d'\))} below.

\smallskip
\noindent\SGARefTarget{sga3:xi:theorem:4:1:condition:c-prime}
\textit{(\(c'\))} The \(T_n\) are smooth over \(S\), and the
transition morphisms \(p_{n,m}\colon T_m\to T_n\) are smooth.

This is precisely
\SGARefLink{sga3:xi:corollary:2:2:bis}{2.2 bis}.

\smallskip
\noindent\SGARefTarget{sga3:xi:theorem:4:1:condition:d-prime}
\textit{(\(d'\))} With the notation of
\SGARefLink{sga3:xi:corollary:3:8}{3.8}, a point \(\xi\) of \(F\)
with values in a field \(k\) over \(S\) is just a
multiplicative-type subgroup \(H_0\) of \(G_0=G_k\).  Repeating the
argument in \SGARefLink{sga3:xi:theorem:4:1:condition:b-prime}{\textit{(\(b'\))}} above, we see that the integer
\(d(\xi)\) considered in
\SGARefLink{sga3:xi:corollary:3:8}{3.8} is
\SGARefTarget{sga3:xi:display:4:1:d-xi-quotient}
\[
  d(\xi)=\dim_k\bigl(\mathfrak g_0/\mathfrak g_0^{H_0}\bigr),
\]
where \(\mathfrak g_0\) is the Lie algebra of \(G_0\) and
\(\mathfrak g_0^{H_0}\) is the vector subspace of invariants under
\(H_0\).  This is a finite integer, so condition
\SGARefLink{sga3:xi:corollary:3:8:item:d-prime:1}
{\(\textit{(d')}1^\circ\textit{)}\)} of
\SGARefLink{sga3:xi:corollary:3:8}{3.8} holds.

With the notation of
\SGARefLink{sga3:xi:corollary:3:8:item:d-prime:2}
{\(\textit{(d')}2^\circ\textit{)}\)}, giving a morphism \(v\colon
X\to F\) amounts to giving a multiplicative-type subgroup \(H\) of
\(G_X\).  For \(x\in X\), the integer \(d(\xi_x)\)\footnote{The printed source gives only
\(\bigl(\mathfrak g\otimes\kappa(x)\bigr)^{H_x}\) here and later
\(\operatorname{rank}_x\mathfrak g^H\), omitting the quotient.  That
reading contradicts both \SGARefLink{sga3:xi:display:4:1:d-xi-quotient}{the immediately preceding formula}
\(d(\xi)=\dim(\mathfrak g_0/\mathfrak g_0^{H_0})\) and the lifting
torsor \(N/M\).  The quotient is restored in both places.} is then
the dimension of
\SGARefTarget{sga3:xi:display:4:1:d-xi-x-invariants}
\[
  \bigl(\mathfrak g\otimes\kappa(x)\bigr)
  \big/
  \bigl(\mathfrak g\otimes\kappa(x)\bigr)^{H_x}.
\]
where \(\mathfrak g\) is the sheaf of Lie algebras of \(G_X\).  It
is locally free of finite type over \(X\), because \(G_X\) is smooth
over \(X\), and \(\mathfrak g\otimes\kappa(x)\) is the Lie algebra of
the fibre \(G_x\) of \(G_X\) at \(x\); here \(H_x\) is the fibre of
\(H\) at \(x\).  Now \(\mathfrak g\) is a locally free module on
\(X\)\footnote{The printed source says ``over \(S\)'' here, although
\(\mathfrak g\) is the Lie-algebra sheaf of \(G_X\); the base is
therefore \(X\).} on which the multiplicative-type group \(H\) acts, and the
invariant subfunctor \(\mathfrak g^H\) is therefore given by a
locally direct-summand, hence locally free, subsheaf of
\(\mathfrak g\).  By descent, one reduces to the case in which
\(H\) is diagonalizable, and then applies \SGARefLink{sga3:i:proposition:4-7-3}{Exposé I, 4.7.3}, noting
that the invariant subsheaf corresponds to the degree-zero
component.  Consequently,
\[
  d(\xi_x)=\operatorname{rank}_x
  \bigl(\mathfrak g/\mathfrak g^H\bigr),
\]
so this is a locally constant function of \(x\).  This proves
condition \SGARefLink{sga3:xi:theorem:4:1:condition:d-prime}{\textit{(\(d'\))}}.

We have thus verified the conditions of
\SGARefLink{sga3:xi:proposition:3:7}{3.7}, completing the proof of
\SGARefLink{sga3:xi:theorem:4:1}{4.1}.

\medskip
\noindent\SGARefTarget{sga3:xi:remark:4:6}\textbf{Remark 4.6.} ---
When \(G=\operatorname{GL}(n)_S\), one can give a much simpler and
more explicit direct proof of
\SGARefLink{sga3:xi:theorem:4:1}{4.1}, using \SGARefLink{sga3:i:proposition:4-7-3}{Exposé I, 4.7.3}.  The
proof moreover shows that, in this case, the moduli scheme is a
disjoint sum of affine \(S\)-schemes.  Proceeding as indicated in
\SGARefLink{sga3:xi:remark:4:3}{4.3}, one obtains the same result
whenever \(G\) is a closed subgroup of a group of the form
\(\operatorname{GL}(n)_S\).

One must not, however, suppose that the preschemes representing the
functors in \SGARefLink{sga3:xi:theorem:4:1}{4.1} or
\SGARefLink{sga3:xi:corollary:4:2}{4.2} are always disjoint sums of
affine \(S\)-schemes.  For example, let \(H\) be a
multiplicative-type group of finite type over a locally noetherian
prescheme \(S\).  By \SGARefLink{sga3:x:proposition:5-11}{Exposé X, 5.11}, \(H\) is isotrivial if and only
if
\[
  \operatorname{Hom}_{S\text{-}\mathrm{gr}}(H,\mathbb G_m)
\]
is a disjoint sum of preschemes affine over \(S\).  But, as noted in
\SGARefLink{sga3:x:paragraph:1:6}{Exposé X, 1.6}, there may exist nonisotrivial tori \(H\) of relative
dimension \(2\).  For such a torus,
\(\operatorname{Hom}_{S\text{-}\mathrm{gr}}(H,\mathbb G_m)\) is not
a disjoint sum of \(S\)-preschemes affine over \(S\), and likewise
the ``dual'' twisted constant group
\[
  \operatorname{Hom}_{S\text{-}\mathrm{gr}}(\mathbb G_m,H)
\]
is not such a sum.  Indeed, if two commutative twisted constant
groups \(R,R'\) of finite presentation are dual to one another,
\[
  R'=\operatorname{Hom}_{S\text{-}\mathrm{gr}}(R,\mathbb Z_S),
\]
then one readily sees that one is isotrivial if and only if the
other is.

This last point also shows that such an \(H\) is not isomorphic to a
subgroup of a group of the form \(\operatorname{GL}(n)_S\).  More
precisely, one can show that a multiplicative-type group of finite
type over a connected locally noetherian \(S\) is isomorphic to a
subgroup of a group of the form
\(\operatorname{Aut}_{\mathrm{Mod}}(E)\), where \(E\) is a locally
free module of finite type over \(S\), if and only if it is
isotrivial.  Finally, taking \(G=H\times\mathbb G_m\) in the two
preceding examples gives an example in which the prescheme
representing the functor \(F\) of
\SGARefLink{sga3:xi:theorem:4:1}{4.1} is not a disjoint sum of
affine \(S\)-preschemes; if desired, \(G\) may be taken to be a
torus of relative dimension \(3\).

% Exact continuation: Section 5, Expo11 local p. 25 /
% complete-reader p. 747 / source-margin p. 165.

% END INCLUDED SGA3_English_current_buildable_source_20260823/source/inputs/XI/tex/components/09_expose_XI_section_4.tex

% BEGIN INCLUDED SGA3_English_current_buildable_source_20260823/source/inputs/XI/tex/components/10_expose_XI_section_5.tex
% SGA 3, Exposé XI, English reconstruction.
% Scope: all of Section 5, from its heading through Remark 5.11,
% stopping immediately before the heading of Section 6.
% Authority: Polo--Gille Expo11.pdf, component-local pp. 25--29,
% complete-reader pp. 747--751, printed pp. 165--171.
% Jacob Reinhold's CC BY 4.0 Markdown lineage was used only as a phrasing
% comparison. The Polo--Gille PDF controls text, formulas, and notes.

\SGARefTarget{sga3:xi:section:5}\section*{5. First Corollaries of the
Representability Theorem}
\addcontentsline{toc}{section}{5. First Corollaries of the
Representability Theorem}

Let \(H\) and \(G\) be as in \SGARefLink{sga3:xi:corollary:4:2}{4.2}.  Put
\SGARefTarget{sga3:xi:display:5:hom-scheme}
\[
  M=\operatorname{Hom}_{S\text{-}\mathrm{gr}}(H,G),
\]
which, by \SGARefLink{sga3:xi:corollary:4:2}{4.2}, is a smooth \(S\)-prescheme separated over \(S\).
Notice that \(G\) acts on the functor
\(M=\operatorname{Hom}_{S\text{-}\mathrm{gr}}(H,G)\) by
\SGARefTarget{sga3:xi:display:5:conjugation-action}
\[
  (g,u)\longmapsto\operatorname{int}(g)\circ u.
\]
Hence there is a canonical morphism
\SGARefTarget{sga3:xi:display:5:x}
\[
  \tag{x}
  G\times_S M\longrightarrow M\times_S M,
\]
whose components are \SGARefLink{sga3:xi:display:5:conjugation-action}{the preceding morphism}
\(G\times_S M\to M\) and the second projection
\(G\times_S M\to M\).

\medskip
\noindent\SGARefTarget{sga3:xi:corollary:5:1}\textbf{Corollary 5.1.} ---
\textit{The preceding morphism \(\mathrm{\SGARefLink{sga3:xi:display:5:x}{(x)}}\) is smooth.}

This follows from \SGARefLink{sga3:xi:corollary:4:2}{4.2} and \SGARefLink{sga3:xi:corollary:2:3}{2.3}.  This statement is equivalent to the
following one.

\medskip
\noindent\SGARefTarget{sga3:xi:corollary:5:2}\textbf{Corollary 5.2.} ---
\textit{Let \(u_1,u_2:H\rightrightarrows G\) be two homomorphisms of
\(S\)-groups.  Then the subfunctor
\(\operatorname{Transp}(u_1,u_2)\) of \(G\) (cf.\ \SGARefLink{sga3:xi:corollary:2:4}{2.4}) is
representable by a closed subprescheme of \(G\), smooth over \(S\).}

It remains only to prove that
\(\operatorname{Transp}(u_1,u_2)\to G\) is indeed a closed
immersion.  This follows because it is the equalizer of a pair of
morphisms \(G\rightrightarrows M\), namely
\[
  g\longmapsto\operatorname{int}(g)u_1
  \qquad\text{and}\qquad
  g\longmapsto u_2,
\]
the latter being the ``constant'' morphism, and because \(M\) is
separated over \(S\).  In particular:

\medskip
\noindent\SGARefTarget{sga3:xi:corollary:5:3}\textbf{Corollary 5.3.} ---
\textit{Let \(u:H\to G\) be a morphism of \(S\)-groups.  Then
\[
  \operatorname{Centr}_G(u)=\operatorname{Transp}(u,u)
\]
is representable by a closed group subprescheme of \(G\), smooth
over \(S\).  Moreover, the quotient of \(G\) by
\(\operatorname{Centr}_G(u)\) is representable by an open
subprescheme of \(M\).}

It remains to prove the last assertion.  The morphism
\SGARefTarget{sga3:xi:display:5:orbit-map}
\[
  g\longmapsto\operatorname{int}(g)\circ u
\]
from \(G\) to \(M\) is smooth and of finite type by \SGARefLink{sga3:xi:corollary:5:2}{5.2}, hence is an
open morphism (EGA IV, 6.6).  Let \(U\) denote its image, equipped
with the structure induced from \(M\).  The induced morphism
\(G\to U\) is smooth, surjective, and of finite type, and is
therefore a covering for the faithfully flat quasi-compact
topology.  Moreover, it is clear that \SGARefLink{sga3:xi:display:5:orbit-map}{the preceding morphism}
\(G\to M\) makes \(G\) a formally principal homogeneous sheaf under
\(\operatorname{Centr}_G(u)_M\).  It follows that the sheaf
\(G/\operatorname{Centr}_G(u)\) is represented by \(U\).

\medskip
\noindent\SGARefTarget{sga3:xi:corollary:5:4}\textbf{Corollary 5.4.} ---
\textit{Let \(u_1,u_2:H\rightrightarrows G\) be two homomorphisms of
\(S\)-groups.  The following conditions are equivalent:}

\begin{enumerate}
\item[(i)] \SGARefTarget{sga3:xi:corollary:5:4:item:i}
\textit{\(u_1\) and \(u_2\) are conjugate locally for the étale
topology.}

\item[(i bis)] \SGARefTarget{sga3:xi:corollary:5:4:item:i-bis}
\textit{\(u_1\) and \(u_2\) are conjugate locally for the faithfully
flat quasi-compact topology.}

\item[(ii)] \SGARefTarget{sga3:xi:corollary:5:4:item:ii}
\textit{For every \(s\in S\), let \(\overline{s}\) denote the
spectrum of an algebraic closure of \(\kappa(s)\).  The morphisms
\[
  u_{1,\overline{s}},u_{2,\overline{s}}:
  H_{\overline{s}}\rightrightarrows G_{\overline{s}}
\]
are conjugate by an element of
\(G\bigl(\kappa(\overline{s})\bigr)\).}

\item[(ii bis)] \SGARefTarget{sga3:xi:corollary:5:4:item:ii-bis}
\textit{The structural morphism
\(\operatorname{Transp}(u_1,u_2)\to S\) is surjective.}

\item[(iii)] \SGARefTarget{sga3:xi:corollary:5:4:item:iii}
\textit{\(\operatorname{Transp}(u_1,u_2)\) is a torsor under the
action of the smooth \(S\)-group prescheme of finite type
\(\operatorname{Centr}_G(u_1)\).}
\end{enumerate}

The implications \SGARefLink{sga3:xi:corollary:5:4:item:i}{(i)} \(\Rightarrow\) \SGARefLink{sga3:xi:corollary:5:4:item:i-bis}{(i bis)} and
\SGARefLink{sga3:xi:corollary:5:4:item:ii}{(ii)} \(\Rightarrow\) \SGARefLink{sga3:xi:corollary:5:4:item:ii-bis}{(ii bis)} are immediate, the latter by the
Nullstellensatz.  The implication \SGARefLink{sga3:xi:corollary:5:4:item:i-bis}{(i bis)} \(\Rightarrow\) \SGARefLink{sga3:xi:corollary:5:4:item:ii}{(ii)}
follows from the ``principle of finite extension'' (EGA IV, 9).
Moreover, \SGARefLink{sga3:xi:corollary:5:4:item:ii-bis}{(ii bis)} \(\Rightarrow\) \SGARefLink{sga3:xi:corollary:5:4:item:iii}{(iii)}: since
\(\operatorname{Transp}(u_1,u_2)\) is smooth over \(S\), it is flat
over \(S\), and since it is of finite type, it is quasi-compact over
\(S\).  Thus it is faithfully flat and quasi-compact over \(S\) if
and only if its structural morphism is surjective.  On the other
hand, it is formally principal homogeneous under
\(\operatorname{Centr}_G(u_1)\), which is faithfully flat and
quasi-compact over \(S\).  The last condition is therefore also
equivalent to saying that \(\operatorname{Transp}(u_1,u_2)\) is a
torsor under \(\operatorname{Centr}_G(u_1)\), in the faithfully flat
quasi-compact topology.  Finally, \SGARefLink{sga3:xi:corollary:5:4:item:iii}{(iii)} \(\Rightarrow\) \SGARefLink{sga3:xi:corollary:5:4:item:i}{(i)} follows
from the smoothness of \(\operatorname{Transp}(u_1,u_2)\) over \(S\)
and from Hensel's lemma in the form \SGARefLink{sga3:xi:proposition:1:10}{1.10}.

\medskip
\noindent\SGARefTarget{sga3:xi:remark:5:5}\textbf{Remark 5.5.} ---
Fix \(u_1=u:H\to G\).  Consider the functor
\SGARefTarget{sga3:xi:display:5:5:functor}
\[
  (\mathbf{Sch})^\circ/S\longrightarrow(\mathbf{Ens})
\]
that assigns to every \(T\) over \(S\) the set of homomorphisms of
\(T\)-groups \(u_2:H_T\to G_T\) that are locally conjugate, for the
étale topology, to \(u_T:H_T\to G_T\).  \SGARefLink{sga3:xi:display:5:5:functor}{This functor} is represented
precisely by the open subprescheme of \(M\), isomorphic to
\(G/\operatorname{Centr}_G(u)\), considered in \SGARefLink{sga3:xi:corollary:5:3}{5.3}.

Let us sketch the variants of the preceding results obtained by
applying \SGARefLink{sga3:xi:theorem:4:1}{4.1} instead of \SGARefLink{sga3:xi:corollary:4:2}{4.2}.  Thus let \(G\) be a group prescheme
smooth and affine over \(S\), and now let \(M\) denote the smooth
\(S\)-prescheme separated over \(S\) that represents the functor
considered in \SGARefLink{sga3:xi:theorem:4:1}{4.1}.  Again \(G\) acts on \(M\):
\SGARefTarget{sga3:xi:display:5:subgroup-conjugation-action}
\[
  (g,H)\longmapsto\operatorname{int}(g)(H).
\]
Hence, as above, there is a morphism
\SGARefTarget{sga3:xi:display:5:x-bis}
\[
  \tag{x bis}
  G\times_S M\longrightarrow M\times_S M.
\]

\medskip
\noindent\SGARefTarget{sga3:xi:corollary:5:1:bis}\textbf{Corollary 5.1 bis.} ---
\textit{\SGARefLink{sga3:xi:display:5:x-bis}{The preceding morphism} is smooth.}

This follows from \SGARefLink{sga3:xi:theorem:4:1}{4.1} and \SGARefLink{sga3:xi:corollary:2:3:bis}{2.3 bis}.  We again conclude:

\medskip
\noindent\SGARefTarget{sga3:xi:corollary:5:2:bis}\textbf{Corollary 5.2 bis.} ---
\textit{Let \(H_1,H_2\) be two multiplicative-type subgroups of
\(G\).  Then the subfunctor
\(\operatorname{Transp}_G(H_1,H_2)\) of \(G\) is representable by a
closed subprescheme of \(G\), smooth over \(S\).}

In particular:

\medskip
\noindent\SGARefTarget{sga3:xi:corollary:5:3:bis}\textbf{Corollary 5.3 bis.} ---
\textit{Let \(H\) be a multiplicative-type subgroup of \(G\).  Then
the group subfunctor \(\operatorname{Norm}_G(H)\) of \(G\) is
representable by a closed subprescheme of \(G\), smooth over \(S\).
Moreover, the quotient \(G/\operatorname{Norm}_G(H)\) is
representable by an open subprescheme of \(M\).}

\medskip
\noindent\SGARefTarget{sga3:xi:corollary:5:4:bis}\textbf{Corollary 5.4 bis.} ---
\textit{Let \(H_1,H_2\) be two multiplicative-type subgroups of
\(G\).  The following conditions are equivalent:}

\begin{enumerate}
\item[(i)] \SGARefTarget{sga3:xi:corollary:5:4:bis:item:i}
\textit{\(H_1\) and \(H_2\) are conjugate locally for the étale
topology.}

\item[(i bis)] \SGARefTarget{sga3:xi:corollary:5:4:bis:item:i-bis}
\textit{\(H_1\) and \(H_2\) are conjugate locally for the faithfully
flat quasi-compact topology.}

\item[(ii)] \SGARefTarget{sga3:xi:corollary:5:4:bis:item:ii}
\textit{For every \(s\in S\), let \(\overline{s}\) denote the
spectrum of an algebraic closure of \(\kappa(s)\).  The subgroups
\(H_{1,\overline{s}}\) and \(H_{2,\overline{s}}\) of
\(G_{\overline{s}}\) are conjugate by an element of
\(G\bigl(\kappa(\overline{s})\bigr)\).}

\item[(ii bis)] \SGARefTarget{sga3:xi:corollary:5:4:bis:item:ii-bis}
\textit{The structural morphism
\(\operatorname{Transp}(H_1,H_2)\to S\) is surjective.}

\item[(iii)] \SGARefTarget{sga3:xi:corollary:5:4:bis:item:iii}
\textit{\(\operatorname{Transp}(H_1,H_2)\) is a principal
homogeneous bundle under the action of the smooth \(S\)-group
prescheme of finite type \(\operatorname{Norm}_G(H_1)\).}
\end{enumerate}

\medskip
\noindent\SGARefTarget{sga3:xi:remark:5:5:bis}\textbf{Remark 5.5 bis.} ---
\SGARefLink{sga3:xi:remark:5:5}{Remark 5.5} carries over to the present case as well.

\medskip
\noindent\SGARefTarget{sga3:xi:remark:5:6}\textbf{Remark 5.6.} ---
To prove \SGARefLink{sga3:xi:corollary:5:2}{5.2}, and hence also the first assertion of \SGARefLink{sga3:xi:corollary:5:3}{5.3} and the
assertions of \SGARefLink{sga3:xi:corollary:5:4}{5.4}, the reference to \SGARefLink{sga3:xi:corollary:4:2}{4.2} may be replaced by a
reference to VIII, \SGARefLink{sga3:viii:theorem:6:4}{6.4}, whose proof is much easier.  This also shows
that the hypothesis that \(G\) be affine over \(S\) is unnecessary
there.  Moreover, in \SGARefLink{sga3:xi:section:6}{Section 6} we shall show how a variant of this
method extends these results to certain groups \(H\) more general
than groups of multiplicative type.  The same observations apply to
the variants \SGARefLink{sga3:xi:corollary:5:2:bis}{5.2 bis} and following.  By contrast, the result \SGARefLink{sga3:xi:theorem:5:8}{5.8}
below makes essential use of all the hypotheses imposed---in
particular, that \(G\) be affine and smooth over \(S\), and that
\(H\) be of multiplicative type---and the lecturer knows no proof
other than one using the representability theorems \SGARefLink{sga3:xi:theorem:4:1}{4.1} or
\SGARefLink{sga3:xi:corollary:4:2}{4.2}.\footnote{The situation has changed since this text was written;
cf.\ Exposés XV and XIX, no.~6.}

\medskip
\noindent\SGARefTarget{sga3:xi:paragraph:5:7}\textbf{5.7.}
Since the morphism \(\mathrm{\SGARefLink{sga3:xi:display:5:x}{(x)}}\), respectively
\SGARefLink{sga3:xi:display:5:x-bis}{\(\mathrm{(x\ bis)}\)}, is smooth and hence open, its image is open.
Let \(R\) be this image, equipped with the structure induced from
\(M\times_S M\).  One checks readily that \(R\) is an equivalence
relation on \(M\), namely the relation made explicit in \SGARefLink{sga3:xi:corollary:5:4}{5.4},
respectively \SGARefLink{sga3:xi:corollary:5:4:bis}{5.4 bis}.  It would be interesting to know whether the
quotient sheaf \(M/R\), which is formally étale over \(S\), is
representable; if so, it is represented by a prescheme étale over
\(S\).  This is the case, for example, when \(S\) is the spectrum of
a field.  Notice also that \(R\) need not be closed in
\(M\times_S M\), even when \(G\) is quasi-finite over \(S\).  Thus,
when \(M/R\) is representable, it may be nonseparated over \(S\).

\medskip
\noindent\SGARefTarget{sga3:xi:theorem:5:8}\textbf{Theorem 5.8.} ---
\textit{Let \(S\) be a prescheme, let \(G\) be a smooth affine
\(S\)-prescheme, and let \(s\in S\).  Then:}

\begin{enumerate}
\item[(a)] \SGARefTarget{sga3:xi:theorem:5:8:item:a}
\textit{For every multiplicative-type subgroup \(H_0\) of \(G_s\),
there exist an étale morphism \(S'\to S\), a point \(s'\in S'\)
lying over \(s\) such that the residue-field extension
\(\kappa(s')/\kappa(s)\) is trivial, and a multiplicative-type
subgroup \(H'\) of \(G'=G\times_S S'\), such that}
\SGARefTarget{sga3:xi:display:5:8:item:a:fiber}
\[
  H'_{s'}=H_0\otimes_{\kappa(s)}\kappa(s').
\]

\item[(b)] \SGARefTarget{sga3:xi:theorem:5:8:item:b}
\textit{For every group homomorphism \(u_0:H_s\to G_s\), where
\(H\) is a multiplicative-type \(S\)-group prescheme of finite
type, there exist an étale morphism \(S'\to S\), a point \(s'\in S'\)
lying over \(s\) such that \(\kappa(s')/\kappa(s)\) is trivial, and
a group homomorphism
\[
  u':H'\longrightarrow G',
  \qquad
  H'=H\times_S S',
  \quad
  G'=G\times_S S',
\]
such that}
\SGARefTarget{sga3:xi:display:5:8:item:b:fiber}
\[
  u'_{s'}:H'_{s'}\longrightarrow G'_{s'}
  \quad\text{equals}\quad
  u_0\otimes_{\kappa(s)}\kappa(s').
\]
\end{enumerate}

This follows from \SGARefLink{sga3:xi:theorem:4:1}{4.1}, respectively \SGARefLink{sga3:xi:corollary:4:2}{4.2}, and from Hensel's lemma in
the form \SGARefLink{sga3:xi:proposition:1:10}{1.10}.

\medskip
\noindent\SGARefTarget{sga3:xi:proposition:5:9}\textbf{Proposition 5.9.} ---
\textit{Let \(S\) be a prescheme, let \(G\) be a smooth affine
\(S\)-prescheme, and let \(H\) be a multiplicative-type subgroup of
\(G\).  Consider}
\[
  \operatorname{Centr}_G(H)\lhook\joinrel\longrightarrow
  \operatorname{Norm}_G(H),
\]
\textit{which, by \SGARefLink{sga3:xi:corollary:5:3}{5.3} and \SGARefLink{sga3:xi:corollary:5:3:bis}{5.3 bis}, are closed group subpreschemes of
\(G\), smooth over \(S\).  The first is then an open and closed
subprescheme of the second, and the quotient sheaf}
\SGARefTarget{sga3:xi:display:5:9:weyl-quotient}
\[
  W_G(H)=
  \operatorname{Norm}_G(H)\big/\operatorname{Centr}_G(H)
\]
\textit{is represented by an open group subprescheme of
\(\operatorname{Aut}_{S\text{-}\mathrm{gr}}(H)\).  It is therefore
a quasi-finite, étale, and separated \(S\)-group prescheme.}

Indeed, consider the evident homomorphism
\SGARefTarget{sga3:xi:display:5:9:theta}
\[
  \theta:\operatorname{Norm}_G(H)\longrightarrow
  \operatorname{Aut}_{S\text{-}\mathrm{gr}}(H),
\]
whose kernel is, by definition, \(\operatorname{Centr}_G(H)\).
Since \(\operatorname{Aut}_{S\text{-}\mathrm{gr}}(H)\) is
represented by an étale \(S\)-group prescheme separated over \(S\)
(X, \SGARefLink{sga3:x:corollary:5-10}{5.10}), its unit section is an open and closed immersion.  Its
inverse image under \(\theta\) is therefore an open and closed
subgroup of \(\operatorname{Norm}_G(H)\).

Moreover, \(\theta\) is smooth.  This follows formally from the
definitions and from the facts that \(\operatorname{Norm}_G(H)\) is
smooth over \(S\) and that
\(\operatorname{Aut}_{S\text{-}\mathrm{gr}}(H)\) is étale over
\(S\).  As in \SGARefLink{sga3:xi:corollary:5:3}{5.3}, the image of \(\theta\) is consequently an open
subprescheme \(U\) of
\(\operatorname{Aut}_{S\text{-}\mathrm{gr}}(H)\), and, with the
induced structure, \(U\) represents the quotient sheaf
\(\operatorname{Norm}_G(H)/\operatorname{Centr}_G(H)\).  The latter
is therefore étale and separated over \(S\).  It is also quasi-finite
over \(S\), since it is quasi-compact over \(S\), being the image of
\(\operatorname{Norm}_G(H)\), which is quasi-compact.  This proves
\SGARefLink{sga3:xi:proposition:5:9}{5.9}.

\medskip
\noindent\SGARefTarget{sga3:xi:corollary:5:10}\textbf{Corollary 5.10.} ---
\textit{For every \(s\in S\), put}
\SGARefTarget{sga3:xi:display:5:10:w}
\[
  w(s)=
  \operatorname{rank}
  \bigl(
    \operatorname{Norm}_{G_s}(H_s)/
    \operatorname{Centr}_{G_s}(H_s)
  \bigr).
\]
\textit{This is also the index of
\(\operatorname{Centr}_{G(\overline{k})}(H(\overline{k}))\) in
\(\operatorname{Norm}_{G(\overline{k})}(H(\overline{k}))\), where
\(\overline{k}\) is an algebraic closure of \(\kappa(s)\).  Then the
function \(s\mapsto w(s)\) is lower semicontinuous.  It is constant
in a neighborhood of \(s\) if and only if \(W_G(H)\) is finite over
\(S\) in a neighborhood of \(s\).}

Indeed, for every \(S\)-prescheme \(W\) that is étale, of finite
type, and separated over \(S\), the function
\[
  s\longmapsto
  \operatorname{rank}[W_s:\kappa(s)]
\]
is lower semicontinuous, and it is constant in a neighborhood of
\(s\) if and only if \(W\) is finite over \(S\) in a neighborhood of
\(s\).  This fact was noted in SGA 1, I, 10.9; its proof, which
presents no difficulty, is given in EGA IV.\footnote{EGA IV, 15.5.1
and 18.10.7.}

\medskip
\noindent\SGARefTarget{sga3:xi:remark:5:11}\textbf{Remark 5.11.}%
\footnote{Editor's note: For a generalization to the nonaffine case,
see M.~Raynaud, \textit{Faisceaux amples sur les schémas en groupes
et les espaces homogènes}, Lecture Notes in Mathematics 119 (1970),
IX.2.8.} ---
Let \(G\) be a group prescheme affine and smooth over \(S\), and let
\(H\) be a multiplicative-type subgroup.  Then \SGARefLink{sga3:xi:corollary:5:3}{5.3} and \SGARefLink{sga3:xi:corollary:5:3:bis}{5.3 bis} imply
that the quotients
\SGARefTarget{sga3:xi:display:5:11:quotients}
\[
  G/\operatorname{Centr}_G(H)
  \qquad\text{and}\qquad
  G/\operatorname{Norm}_G(H)
\]
are preschemes smooth and quasi-affine over \(S\).  Indeed, in both
cases the moduli prescheme \(M\) in which the quotient is embedded
has the property that every open \(U\subset M\) quasi-compact over
\(S\) is quasi-affine over \(S\), by \SGARefLink{sga3:xi:corollary:3:11}{3.11}.

We shall moreover see in the next Exposé\footnote{Editor's note:
Cf.\ Remark XII.5.7.  We mention here the following generalization.
Without assuming \(G\) affine, if \(H\) is a torus and \(S\) is
locally noetherian, \(U\) is affine and smooth by a theorem of
Raynaud (\textit{loc.\ cit.}, IX.2.6 and IX.2.8).  Indeed,
\(G/\operatorname{Norm}_G(H)\) is an open subprescheme of \(M\),
and the largest representable open subfunctor of \(M\) is a disjoint
sum of open subpreschemes smooth and affine over \(S\).} that for
every such \(U\), the schematic closure \(\overline{U}\) of \(U\) in
\(M\) is even affine over \(S\).  Thus the quotients under
consideration are affine over \(S\), provided that the open
subprescheme \(U\) of the corresponding moduli prescheme---of group
homomorphisms from \(H\) to \(G\), respectively of
multiplicative-type subgroups of \(G\)---is closed in \(M\).  This
is the case, for example, when \(H\) is a ``maximal torus'' in the
second case, or when \(S\) is the spectrum of a field; cf.\
XII.5.4.\footnote{Editor's note: This assertion assumes that \(G\)
has locally constant reductive rank.}  I do not know whether the
quotients under consideration are affine over \(S\) in general.

% Exact continuation: Section 6 heading, Expo11 local p. 29 /
% complete-reader p. 751 / printed p. 171.

% END INCLUDED SGA3_English_current_buildable_source_20260823/source/inputs/XI/tex/components/10_expose_XI_section_5.tex

% BEGIN INCLUDED SGA3_English_current_buildable_source_20260823/source/inputs/XI/tex/components/11_expose_XI_section_6_through_end.tex
% SGA 3, Exposé XI, English reconstruction.
% Scope: all of §6, from its heading and Proposition 6.1 through
% terminal Remark 6.12 and its Murre/Bourbaki footnote; stopping
% immediately before Exposé XII.
% Authority: Polo--Gille Expo11.pdf, component-local pp. 29--34,
% complete-reader pp. 751--756, source-margin pp. 170--179.
% Jacob Reinhold's CC BY 4.0 Markdown lineage was used only as a phrasing
% comparison. The Polo--Gille PDF controls text, formulas, and numbering.

\SGARefTarget{sga3:xi:section:6}
\section*{6. A Rigidity Property for Homomorphisms\\
of Certain Group Schemes, and the Representability\\
of Certain Transporters\footnote{\SGARefLink{sga3:xi:section:6}{The
present section} does not use the results of Sections \SGARefLink{sga3:xi:section:3}{3}, \SGARefLink{sga3:xi:section:4}{4}, and \SGARefLink{sga3:xi:section:5}{5}; its
natural place would be in Exposé VI\textsubscript{B}.}}
\addcontentsline{toc}{section}{6. A Rigidity Property for Homomorphisms
of Certain Group Schemes, and the Representability of Certain
Transporters}

\medskip
\noindent\SGARefTarget{sga3:xi:proposition:6:1}\textbf{Proposition 6.1.}
---
\textit{Let \(S\) be a locally noetherian prescheme, let \(H\) be a
commutative \(S\)-group prescheme of finite type over \(S\), and let
\(E\) be a set of positive integers stable under multiplication.
Suppose that the following conditions hold:}

\begin{enumerate}
\item[(a)] \SGARefTarget{sga3:xi:proposition:6:1:item:a}%
\textit{for every \(n\in E\), the subgroup}
\SGARefTarget{sga3:xi:display:6:1:n-torsion}
\[
  {}_nH=\operatorname{Ker}\bigl(n\cdot\operatorname{id}_H\bigr)
\]
\textit{is finite and flat over \(S\);}

\item[(b)] \SGARefTarget{sga3:xi:proposition:6:1:item:b}%
\textit{for every \(s\in S\), the family of the
\({}_nH_s\), \(n\in E\), is schematically dense in the fibre
\(H_s\).\footnote{The printed French text has \(H\) here; the
fibrewise condition and its use in the proof require \(H_s\).}}
\end{enumerate}

\textit{Let \(Y\) be a closed subprescheme of \(H\).  Then the
subfunctor \(\prod_{H/S}(Y/H)\) of \(S\) (compare \SGARefLink{sga3:viii:section:6}{Exposé VIII, 6}) is
representable by a closed subprescheme \(T\) of \(S\).  If \(S\) is
noetherian, there exists \(n\in E\) such that\footnote{The printed
formula writes the right-hand quotient as
\((Y\cap{}_nH)/S\).  Since the product functor is taken over
\({}_nH/S\), and as the proof requires, the quotient is
\((Y\cap{}_nH)/{}_nH\).}}
\SGARefTarget{sga3:xi:display:6:1:cross}
\[
  \tag{x}
  \prod_{H/S}(Y/H)
  =
  \prod_{{}_nH/S}\bigl((Y\cap{}_nH)/{}_nH\bigr).
\]

Indeed, by \SGARefLink{sga3:viii:theorem:6:4}{Exposé VIII, 6.4}, condition \SGARefLink{sga3:xi:proposition:6:1:item:a}{(a)} implies that \({}_nH\) is
finite and flat, and hence a fortiori ``essentially free'' over \(S\).
Thus the right-hand side of \SGARefLink{sga3:xi:display:6:1:cross}{\((x)\)} is representable by a closed
subprescheme \(T_n\) of \(S\).\footnote{The printed French text has
``\(T_n\) of \(T\)'' here; the next sentence makes the \(T_n\) closed
subpreschemes of \(S\), and \(T\) is defined only afterward as their
stable value.}  Order \(E\) by divisibility.  The
\(T_n\) then form a decreasing family of closed subpreschemes of
\(S\); consequently, when \(S\) is noetherian---as we may now
suppose---this family is stationary for sufficiently large \(n\).
Let \(T\) be the stable value of \(T_n\).  We claim that \(T\)
represents the left-hand side of \SGARefLink{sga3:xi:display:6:1:cross}{\((x)\)}, which will complete
the proof.

It is enough to prove the following.  Let \(S'\) be a prescheme over
\(S\) such that
\SGARefTarget{sga3:xi:display:6:1:base-change-condition}
\[
  (Y\cap{}_nH)_{S'}=({}_nH)_{S'}
  \qquad\text{for every \(n\in E\),}
\]
or, equivalently, such that
\(Y_{S'}\supset({}_nH)_{S'}\) for every \(n\in E\).  Then
\(Y_{S'}=H_{S'}\).  This follows at once from \SGARefLink{sga3:ix:lemma:4:4}{Exposé IX, 4.4}:
in view of conditions \SGARefLink{sga3:xi:proposition:6:1:item:a}{(a)} and \SGARefLink{sga3:xi:proposition:6:1:item:b}{(b)}, the family of the
\(({}_nH)_{S'}\) is schematically dense in \(H_{S'}\).

\medskip
\noindent\SGARefTarget{sga3:xi:theorem:6:2}\textbf{Theorem 6.2.} ---
\textit{Let \(S\) be a prescheme, let \(H,E\) be as in
\SGARefLink{sga3:xi:proposition:6:1}{6.1} and satisfy conditions
\SGARefLink{sga3:xi:proposition:6:1:item:a}{(a)} and
\SGARefLink{sga3:xi:proposition:6:1:item:b}{(b)}, let
\(u:H\to G\) be a homomorphism of \(S\)-groups from \(H\) to an
\(S\)-group prescheme locally of finite type over \(S\), and let
\(K\) be a closed subgroup prescheme of \(G\).\footnote{The printed
French text has \(H\) here; the proof immediately uses \(K\subset G\),
as required by the transporter.}  Consider the
subfunctor \(\operatorname{Transp}_G(u,K)\) of \(G\) (cf.\ \SGARefLink{sga3:xi:proposition:2:5}{2.5}).
It is representable by a closed subprescheme of \(G\).  If \(G\) is
noetherian---for example, if \(S\) is noetherian and \(G\) is of
finite type over \(S\)---there exists \(n\in E\) such that}
\SGARefTarget{sga3:xi:display:6:2:transporter-stabilization}
\[
  \operatorname{Transp}_G(u,K)
  =
  \operatorname{Transp}_G\bigl(u|_{{}_nH},K\bigr).
\]
\textit{Suppose, finally, that \(G\) is smooth over \(S\), that
\(K\) is smooth over \(S\) or is of multiplicative type, and that
\(H,E\) satisfy the following condition, stronger than \SGARefLink{sga3:xi:proposition:6:1:item:a}{(a)}:}

\begin{enumerate}
\item[\(a'\)] \SGARefTarget{sga3:xi:theorem:6:2:item:a-prime}%
\textit{for every \(n\in E\), the subgroup \({}_nH\) of \(H\) is
of multiplicative type.}
\end{enumerate}

\textit{Then \(\operatorname{Transp}_G(u,K)\) is smooth over \(S\).}

Consider the \(G\)-group
\SGARefTarget{sga3:xi:display:6:2:hg}
\[
  H_G=H\times_SG,
\]
which plainly satisfies the conditions of
\SGARefLink{sga3:xi:proposition:6:1}{6.1}, with \(S\) replaced by
\(G\) and \(H\) by \(H_G\).  Let \(Y\) be the closed subprescheme of
\(H_G\) that is the inverse image of \(K\subset G\) under the map
\SGARefTarget{sga3:xi:display:6:2:conjugation-map}
\[
  H_G=H\times_SG\longrightarrow G,\qquad
  (h,g)\longmapsto\operatorname{int}(g)\cdot u(h).
\]
Then
\SGARefTarget{sga3:xi:display:6:2:product-identification}
\[
  \operatorname{Transp}_G(u,K)
  =
  \prod_{H_G/G}(Y/H_G)
\]
(compare \SGARefLink{sga3:viii:examples:6:5:item:e:01}{Exposé VIII, Example 6.5(e)}).  The first assertions
therefore follow from \SGARefLink{sga3:xi:proposition:6:1}{6.1}.
Moreover, for every quasi-compact open \(U\subset G\), there exists
\(n\in E\) for which the restrictions to \(U\) of the two functors
\SGARefTarget{sga3:xi:display:6:2:local-stabilization}
\[
  \operatorname{Transp}_G(u,K)
  \qquad\text{and}\qquad
  \operatorname{Transp}_G\bigl(u|_{{}_nH},K\bigr)
\]
coincide.  To verify the last assertion of
\SGARefLink{sga3:xi:theorem:6:2}{6.2}, we may therefore replace
\(H\) by some \({}_nH\), after which it is enough to apply \SGARefLink{sga3:xi:proposition:2:5}{2.5},
since \({}_nH\) is assumed to be of multiplicative type over \(S\).

\medskip
\noindent\SGARefTarget{sga3:xi:remark:6:3}\textbf{Remark 6.3.} ---
The preceding proof uses only the very elementary result of \SGARefLink{sga3:viii:theorem:6:4}{Exposé
VIII, 6.4}, together, for its last part, with \SGARefLink{sga3:xi:proposition:2:5}{2.5}; when \(K\) is
smooth over \(S\), the latter is the infinitesimal result of \SGARefLink{sga3:ix:theorem:3:6}{Exposé
IX, 3.6}, and hence the vanishing of the cohomology of
multiplicative-type groups.  In the most important cases (cf.\
\SGARefLink{sga3:xi:unit:6:7}{6.7}), the integers \(n\in E\) may
even be taken prime to the residue characteristics of \(S\), that
is, invertible in \(\mathscr O_S\).  The \({}_nH\) are then finite
étale over \(S\), and the cohomological result invoked is practically
trivial.  Thus, in that situation,
\SGARefLink{sga3:xi:theorem:6:2}{6.2} is independent of the theory of
multiplicative-type groups.

\medskip
\noindent\SGARefTarget{sga3:xi:unit:6:4}\textbf{6.4.} ---
As usual, one sees that \SGARefLink{sga3:xi:proposition:6:1}{6.1}
extends to the following situation.  Let \(S'\) be an \(S\)-prescheme
(not necessarily locally noetherian), put \(H'=H_{S'}\), and let
\(Y'\) be a closed subprescheme of \(H'\), provided that
\(Y'\to H'\) is of finite presentation, that is, that the ideal
defining \(Y'\) is of finite type.  Then
\(\prod_{H'/S'}(Y'/H')\) is representable by a closed subprescheme
\(T'\) of \(S'\), the morphism \(T'\to S'\) is of finite
presentation, and, if \(S'\) is quasi-compact, there exists \(n\in E\)
for which the analogue of \SGARefLink{sga3:xi:display:6:1:cross}{\((x)\)} holds.  It follows as well
that the first assertion of \SGARefLink{sga3:xi:theorem:6:2}{6.2}
remains valid without assuming that \(G\) is locally of finite type
over \(S\), provided that the immersion \(K\to G\) is of finite
presentation.

\medskip
\noindent\SGARefTarget{sga3:xi:unit:6:5}\textbf{6.5.} ---
As announced in \SGARefLink{sga3:xi:remark:5:6}{5.6}, Theorem
\SGARefLink{sga3:xi:theorem:6:2}{6.2} makes it possible to extend to
group preschemes satisfying \SGARefLink{sga3:xi:theorem:6:2:item:a-prime}{\((a')\)} and \SGARefLink{sga3:xi:proposition:6:1:item:b}{(b)} above certain results
proved by another method, under more restrictive hypotheses, for
groups of multiplicative type.  This applies to \SGARefLink{sga3:xi:corollary:5:2}{5.2}, the beginning
of \SGARefLink{sga3:xi:corollary:5:3}{5.3}, and \SGARefLink{sga3:xi:corollary:5:4}{5.4} and their corresponding ``bis'' variants, and also
to \SGARefLink{sga3:xi:proposition:5:9}{5.9} and \SGARefLink{sga3:xi:corollary:5:10}{5.10}.  It also applies to the results of
\SGARefLink{sga3:ix:section:5}{Exposé IX, Sections 5} and 6, with the exception of \SGARefLink{sga3:ix:theorem:6:8}{Exposé IX, 6.8}.
The latter assertion is already false for a homomorphism of abelian
schemes \(u:H\to G\) over the spectrum \(S\) of a discrete valuation
ring of residue characteristic \(p>0\): it can happen that
\(\operatorname{Ker}u\) has the trivial group as its generic fibre
and, as its special fibre, a radicial group that is not reduced to
the trivial group.

\medskip
\noindent\SGARefTarget{sga3:xi:unit:6:6}\textbf{6.6.} ---
We have just given an example of a rigidity property of
multiplicative-type groups that is not shared by abelian schemes.
Another, extremely important, example is that the existence theorem
for infinitesimal extensions of homomorphisms in \SGARefLink{sga3:ix:theorem:3:6}{Exposé IX, 3.6} no
longer holds when \(H\) is an abelian scheme.  Thus a nonzero abelian
scheme over a field admits nontrivial infinitesimal variations, in
contrast to a multiplicative-type group.  This is the infinitesimal
aspect of the fact that there is a theory of moduli---one that,
moreover, is far from complete---for abelian varieties, whereas the
theory of moduli of multiplicative-type groups is empty.

Another, ``global,'' aspect of this infinitesimal difference is the
following.  Let \(H\) be an abelian scheme over a locally noetherian
prescheme \(S\), and let \(G\) be a commutative \(S\)-group prescheme
locally of finite type over \(S\).  One can show that
\(\operatorname{Hom}_{S\text{-}\mathrm{gr}}(H,G)\) is representable
by a prescheme locally of finite type over \(S\).  Unlike the
corresponding prescheme when \(H\) is of multiplicative type,
however, it is not étale over \(S\), but only unramified.  Thus, for
example, if \(S\) is the spectrum of a complete discrete valuation
ring and \(H,G\) are abelian schemes over \(S\), there may exist
homomorphisms \(H_0\to G_0\) on the special fibres that do not arise
``by specialization'' from a homomorphism on the generic fibres.

\medskip
\noindent\SGARefTarget{sga3:xi:unit:6:7}\textbf{6.7.} ---
Theorem \SGARefLink{sga3:xi:theorem:6:2}{6.2} applies whenever \(H\)
is an abelian scheme over \(S\), or, more generally, an extension of
such a scheme by a torus.  Indeed, the question is local on \(S\), so
one may suppose that there is a prime number \(\ell\) prime to the
residue characteristics of \(S\).  Taking \(E\) to be the set of
powers of \(\ell\) then suffices to satisfy conditions \SGARefLink{sga3:xi:theorem:6:2:item:a-prime}{\((a')\)} and
\SGARefLink{sga3:xi:proposition:6:1:item:b}{(b)}.

An argument analogous to that of
\SGARefLink{sga3:xi:proposition:6:1}{6.1} gives \SGARefLink{sga3:xi:theorem:6:8}{the following result}.

\medskip
\noindent\SGARefTarget{sga3:xi:theorem:6:8}\textbf{Theorem 6.8.} ---
\textit{Let \(X\) be a prescheme smooth over \(S\), with nonempty
connected geometric fibres.  For every closed subprescheme \(Y\) of
\(X\), the functor \(\prod_{X/S}(Y/X)\) is representable by a closed
subprescheme \(T\) of \(S\).  If \(Y\) is of finite presentation
over \(X\), then \(T\) is of finite presentation over \(S\).}

Since \(f:X\to S\) is faithfully flat and locally of finite
presentation, it is a covering for the fpqc topology.  On the other
hand, \(T=\prod_{X/S}(Y/X)\) is plainly a subsheaf of \(S\) for the
fpqc topology.  Hence both the representability of \(T\) by a closed
subprescheme of \(S\) and the question whether \(T\) is of finite
presentation over \(S\) are local on \(S\) for the fpqc topology.
After the base change \(S'\to S\) with \(S'=X\), we are reduced to
the case in which \(X\) admits a section \(e\) over \(S\).  We may
also suppose that \(S\) is affine, and hence in particular
quasi-compact.  We then have:

\medskip
\noindent\SGARefTarget{sga3:xi:corollary:6:9}\textbf{Corollary 6.9.}
---
\textit{Under the hypotheses of
\SGARefLink{sga3:xi:theorem:6:8}{6.8}, suppose that \(S\) is
quasi-compact and that \(X\) admits a section \(e\) over \(S\).  For
every integer \(n\geq0\), let \(X_n\) be the order-\(n\)
infinitesimal neighbourhood of the section \(e\) in \(X\).  Suppose
that \(Y\) is of finite presentation over \(X\).  Then there exists
an integer \(n\geq0\) such that}
\SGARefTarget{sga3:xi:display:6:9:stabilization}
\[
  \prod_{X/S}(Y/X)
  =
  \prod_{X_n/S}(Y_n/X_n),
  \qquad Y_n=Y\cap X_n.
\]

This corollary implies
\SGARefLink{sga3:xi:theorem:6:8}{6.8} when \(Y\) is of finite
presentation over \(X\), by \SGARefLink{sga3:viii:theorem:6:4}{Exposé VIII, 6.4}.  Indeed, since \(X\)
is smooth over \(S\), \(X_n\) is finite and locally free over \(S\),
and a fortiori ``essentially free'' over \(S\) in the sense of
\SGARefLink{sga3:viii:definition:6:1}{Exposé VIII, 6.1}.  Thus \(\prod_{X_n/S}(Y_n/X_n)\) is representable
by a closed subprescheme of \(S\).  Furthermore, either the proof of
the cited result or reduction to the noetherian case shows at once
that this closed subprescheme is of finite presentation over \(S\).

Let us first prove \SGARefLink{sga3:xi:corollary:6:9}{6.9}, and hence
\SGARefLink{sga3:xi:theorem:6:8}{6.8}, when \(S\) is noetherian.  Put
\SGARefTarget{sga3:xi:display:6:9:tn}
\[
  T_n=\prod_{X_n/S}(Y_n/X_n).
\]
The \(T_n\) form a decreasing sequence of closed subpreschemes of
\(S\), and this sequence is stationary because \(S\) is noetherian.
Let
\SGARefTarget{sga3:xi:display:6:9:r}
\[
  R=\bigcap_{n\geq0}\prod_{X_n/S}(Y_n/X_n)
\]
be their common value for large \(n\).  Plainly \(T\subset R\), and
it remains to show that \(R\subset T\).  After the base change
\(R\to S\), we are reduced to the case \(R=S\), that is,
\(Y_n=X_n\) for every \(n\), or equivalently
\(Y\supset X_n\) for every \(n\).  We must then prove \(T=S\), that
is, \(Y=X\).

The inclusions \(Y\supset X_n\) for all \(n\), together with the fact
that \(X\) is locally noetherian, imply that in a neighbourhood of
each point of \(e(S)\), \(Y\) is an induced open subprescheme of
\(X\).  Hence there is an induced open subprescheme \(U\) of \(X\)
such that \(e(S)\subset U\subset Y\).  By \SGARefLink{sga3:ix:corollary:4:3}{Exposé IX, 4.3}, since the
fibres of \(X/S\) are integral, \(U\) is schematically dense in
\(X\).  As \(Y\) is a closed subprescheme containing \(U\), it
follows that \(Y=X\).  This proves
\SGARefLink{sga3:xi:corollary:6:9}{6.9}, and hence
\SGARefLink{sga3:xi:theorem:6:8}{6.8}, in the noetherian case.

The general case is reduced to the preceding one.  For every
\(s\in S\), there is an affine open neighbourhood \(U\) of \(s\)
and an affine open neighbourhood \(V\) of \(e(s)\) such that
\(f(V)\subset U\).  Then \(f(V)\) is an open neighbourhood of \(s\)
contained in \(U\).  Choose an affine open neighbourhood \(S'\) of
\(s\) contained in \(e^{-1}(V)\cap f(V)\), and put
\(X'=V\cap f^{-1}(S')\).  Then \(X'\) and \(S'\) are affine open
subpreschemes of \(X\) and \(S\), respectively, and \(X'/S'\)
admits a section.  By the local nature of
\SGARefLink{sga3:xi:theorem:6:8}{6.8} and
\SGARefLink{sga3:xi:corollary:6:9}{6.9}, we may suppose that
\(S=S'\).

We claim that
\SGARefTarget{sga3:xi:display:6:9:open-reduction}
\[
  \prod_{X/S}(Y/X)=\prod_{X'/S'}(Y'/X'),
  \qquad Y'=Y\cap X'.
\]
Indeed, by \SGARefLink{sga3:ix:corollary:4:6}{Exposé IX, 4.6}, \(X'\) is schematically dense in \(X\),
at least when \(X\) is quasi-separated over \(S\), so that \(X'\)
is retrocompact in \(X\).  In fact, one can readily prove that
\SGARefLink{sga3:ix:corollary:4:6}{Exposé IX, 4.6} remains valid without the retrocompactness
hypothesis.  Likewise, after every base change \(S_1\to S\),
\(X'\times_SS_1\) is schematically dense in \(X\times_SS_1\), and
the claimed equality follows immediately.  We are thus reduced to
the case \(X=X'\), so we may suppose that both \(S\) and \(X\) are
affine.

Now let \(X=\operatorname{Spec}(B)\), and let \(J\) be the ideal of
\(B\) defining \(Y\).  The ideal \(J\) is the filtered colimit of
its finite-type subideals.  Consequently, \(Y\) is the intersection
of closed subpreschemes \(Y_i\) of \(X\) that are of finite
presentation over \(S\), and
\SGARefTarget{sga3:xi:display:6:9:intersection}
\[
  \prod_{X/S}(Y/X)
  =
  \bigcap_i\prod_{X/S}(Y_i/X).
\]
To prove \SGARefLink{sga3:xi:theorem:6:8}{6.8}, it is therefore
enough to treat the case in which \(Y\) is of finite presentation
over \(X\).  It remains to prove
\SGARefLink{sga3:xi:corollary:6:9}{6.9} with \(S\) and \(X\)
affine.  But \(X\) and \(Y\) over \(S\) then arise by base change
\(S\to S_0\) from an analogous situation \(X_0,Y_0\) over a
noetherian prescheme \(S_0\).  This reduces us to the already treated
case in which \(S\) is noetherian, and completes the proof of
\SGARefLink{sga3:xi:theorem:6:8}{6.8} and
\SGARefLink{sga3:xi:corollary:6:9}{6.9}.

\medskip
\noindent\SGARefTarget{sga3:xi:corollary:6:10}\textbf{Corollary 6.10.}
---
\textit{Let \(X\) be an \(S\)-group prescheme, smooth and of finite
presentation, with connected fibres.  Let \(Y\) be a group
prescheme of finite presentation over \(S\), and let \(i:Y\to X\)
be a monomorphism of \(S\)-group preschemes, so that \(Y\) is a
subgroup of \(X\).  Then \(\prod_{X/S}(Y/X)\) is representable by a
closed subprescheme of \(S\) of finite presentation.  If \(S\) is
quasi-compact, let \(X_n\) denote, for every \(n\geq0\), the
order-\(n\) infinitesimal neighbourhood of the unit section of
\(X\), and put \(Y_n=X_n\cap Y\).  Then, for all sufficiently large
\(n\),}
\SGARefTarget{sga3:xi:display:6:10:stabilization}
\[
  \prod_{X/S}(Y/X)=\prod_{X_n/S}(Y_n/X_n).
\]

The proof is essentially that of
\SGARefLink{sga3:xi:corollary:6:9}{6.9}.  The unit sections of \(X\)
and \(Y\) induce bijective immersions \(S\to X_n\) and
\(S\to Y_n\), and hence isomorphisms
\SGARefTarget{sga3:xi:display:6:10:reductions}
\[
  S_{\mathrm{red}}\xrightarrow{\sim}(X_n)_{\mathrm{red}},
  \qquad
  S_{\mathrm{red}}\xrightarrow{\sim}(Y_n)_{\mathrm{red}}.
\]
It follows that the injection \(Y_n\to X_n\) is proper.  Since it is
a monomorphism of finite presentation, it is therefore a closed
immersion.  \SGARefLink{sga3:viii:theorem:6:4}{Exposé VIII, 6.4}, already used above, now shows that
\(\prod_{X_n/S}(Y_n/X_n)\) is representable by a closed
subprescheme of \(S\) of finite presentation over \(S\).  It remains
to prove the last assertion of
\SGARefLink{sga3:xi:corollary:6:10}{6.10} under the additional
assumption that \(S\) is affine.

We reduce once more to the case in which \(S\) is noetherian.  With
the notation of the proof of
\SGARefLink{sga3:xi:corollary:6:9}{6.9}, it is enough to prove
\(R=T\), or, equivalently, that \(Y\supset X_n\) for every \(n\)
implies \(Y=X\).  This hypothesis implies that \(i:Y\to X\) is
étale at the points of the unit section of \(Y\) over \(S\), and
therefore that \(Y\) is smooth over \(S\) at those points.  Hence the
open subprescheme \(Y'\) consisting of the points where \(Y\) is
smooth over \(S\) is an induced open subgroup of \(Y\).  By \SGARefLink{sga3:x:lemma:3:5}{Exposé
X, 3.5}, the morphism \(Y'\to X\) is an étale monomorphism and hence
an open immersion.  Since the fibres of \(X\) are connected and
every open subgroup of an algebraic group is also closed, this open
immersion is surjective, and therefore an isomorphism.  Thus
\(Y'=X\), and a fortiori \(Y=X\), completing the proof of
\SGARefLink{sga3:xi:corollary:6:10}{6.10}.

Arguing as in \SGARefLink{sga3:viii:examples:6:5}{Exposé VIII, 6.5}, we deduce:

\medskip
\noindent\SGARefTarget{sga3:xi:corollary:6:11}\textbf{Corollary 6.11.}
---
\textit{Let \(G\) be an \(S\)-group prescheme locally of finite type
and quasi-separated over \(S\), let \(H\) be a group prescheme
smooth and of finite presentation over \(S\), with connected
fibres, and let \(i:H\to G\) be a monomorphism of \(S\)-groups, so
that \(H\) is a subgroup of \(G\).  Then:}

\begin{enumerate}
\item[(a)] \SGARefTarget{sga3:xi:corollary:6:11:item:a}%
\textit{\(\operatorname{Centr}_G(H)\) and
\(\operatorname{Norm}_G(H)\) are representable by closed
subpreschemes of \(G\) of finite presentation over \(G\).  Likewise,
for every monomorphism \(j:K\to G\) of finite presentation between
\(S\)-group preschemes, \(\operatorname{Transp}_G(H,K)\) is
representable by a closed subprescheme of \(G\) of finite
presentation over \(G\).}

\item[(b)] \SGARefTarget{sga3:xi:corollary:6:11:item:b}%
\textit{If \(G\) is quasi-compact, then in each of the cases in
\SGARefLink{sga3:xi:corollary:6:11:item:a}{(a)} there exists an integer \(n\geq0\) such that, if \(H_n\)
denotes the order-\(n\) infinitesimal neighbourhood of the unit
section of \(H\) (and \(K_n\) is defined similarly in the transporter
case),}
\SGARefTarget{sga3:xi:display:6:11:stabilization}
\[
  \begin{aligned}
  \operatorname{Centr}_G(H)
    &=\operatorname{Centr}_G(H_n),\\
  \operatorname{Norm}_G(H)
    &=\operatorname{Norm}_G(H_n),\\
  \operatorname{Transp}_G(H,K)
    &=\operatorname{Transp}_G(H_n,K)
     =\operatorname{Transp}_G(H_n,K_n).
  \end{aligned}
\]
\end{enumerate}

Apply \SGARefLink{sga3:xi:corollary:6:10}{6.10} over the base
prescheme \(G\) to the group prescheme
\SGARefTarget{sga3:xi:display:6:11:hg}
\[
  X=H_G=H\times_SG.
\]
For the centralizer, take as \(Y\) the inverse image of the diagonal
subgroup of \((G\times_SG)_G\) under a suitable homomorphism of
\(G\)-groups from \(X\) to \((G\times_SG)_G\).  For the transporter,
take the inverse image of \(K_G\) under a suitable homomorphism of
\(G\)-groups from \(X\) to \(G_G\).  The normalizer case reduces to
the transporter case by taking \(K=H\).  The hypothesis on \(G\)
ensures that \(H\to G\) is of finite presentation, and hence that
\(Y\to X\) is of finite presentation.  In the centralizer case, the
same hypothesis ensures that the diagonal subgroup of
\(G\times_SG\) is of finite presentation over \(G\times_SG\);
again it follows that \(Y\) is of finite presentation over \(X\).

\medskip
\noindent\SGARefTarget{sga3:xi:remark:6:12}\textbf{Remark 6.12.} ---
One can prove, using rather delicate results of EGA
VI,\footnote{Cf. also J.\,P. Murre, ``Representation of unramified
functors. Applications,'' \emph{Séminaire Bourbaki}, no.\ 294
(May 1965), Theorem 3, p.\ 13.} that if \(X\) is a prescheme flat
and of finite presentation over \(S\), and is either proper over
\(S\) or has nonempty connected geometric fibres, then for every
closed subprescheme \(Y\) of \(X\) of finite presentation over
\(S\), the functor \(\prod_{X/S}(Y/X)\) is representable by a closed
subprescheme of \(S\) of finite presentation over \(S\).  Likewise,
if \(X\) is an \(S\)-group prescheme flat and of finite presentation
with connected fibres, and \(i:Y\to X\) is a monomorphism of
\(S\)-groups, with \(Y\) an \(S\)-group prescheme of finite
presentation, then \(\prod_{X/S}(Y/X)\) is representable by a closed
subprescheme of \(S\) of finite presentation over \(S\).  In
particular, \SGARefLink{sga3:xi:corollary:6:11:item:a}{6.11(a)}
remains valid if the hypothesis ``\(H\) smooth over \(S\)'' is
replaced by ``\(H\) flat over \(S\).''

% Exact terminal boundary: end of Exposé XI, Expo11 local p. 34 /
% complete-reader p. 756 / source-margin p. 179. Hard stop before
% Exposé XII, complete-reader p. 757.

% END INCLUDED SGA3_English_current_buildable_source_20260823/source/inputs/XI/tex/components/11_expose_XI_section_6_through_end.tex



% END INCLUDED SGA3_English_current_buildable_source_20260823/source/inputs/XI/body.tex


\SGAStartExpose{XII}
\chapter*{Expos\'e XII\\
  Maximal Tori, the Weyl Group, Cartan Subgroups, and the Reductive Center}
\addcontentsline{toc}{chapter}{Expos\'e XII. Maximal Tori, the Weyl Group,
  Cartan Subgroups, and the Reductive Center}
\begin{center}
  \textit{by A. Grothendieck}
\end{center}
% BEGIN INCLUDED SGA3_English_current_buildable_source_20260823/source/inputs/XII/body.tex
% Mechanically extracted after \mainmatter from the frozen standalone master.

\markboth{Exposé XII. Maximal Tori and the Weyl Group}
  {Exposé XII. Maximal Tori and the Weyl Group}
% BEGIN INCLUDED SGA3_English_current_buildable_source_20260823/source/inputs/XII/tex/components/00_expose_XII_sections1_3_en.tex
% Authority: Expo12.pdf, local p. 1 / combined pp. 757--758.
% This component will continue through Section 3.

\SGARefTarget{sga3:xii:expose}

\setcounter{footnote}{-1}
\footnotetext{Version xy of 2 December 2008.}

From the present Exposé onwards, in contrast to the preceding Exposés,
we shall use the known results on the structure of smooth algebraic
groups over an algebraically closed field \(k\), and above all Borel's
theory of affine algebraic groups, expounded in the 1956/57 Chevalley
Seminar, \emph{Classification of Algebraic Lie Groups}.\footnote{Editor's
note: this seminar has since been published, in a form revised by
P.~Cartier, as volume~3 of Chevalley's collected works (Springer,
2005).} Following the usage in the theory of algebraic groups, we shall
refer to that seminar by the siglum \emph{BIBLE}. In the following
Exposés we shall need, in particular, the results of \emph{BIBLE}~4,
5, 6, and~7 (the number following \emph{BIBLE} is the number of the
Exposé). Moreover, it appears that scheme theory brings no notable
simplification to Borel's theory as it is presented in \emph{BIBLE}.
For this reason it did not seem useful to reproduce that theory in the
present Seminar: our purpose here is to deduce from results known over
algebraically closed fields analogous results valid over an arbitrary
base prescheme. (The situation will be different for Chevalley's
structure theory of semisimple groups, which, it appears, is more
advantageously treated \emph{ab ovo} over an arbitrary base
prescheme.)

In the oral Exposés (which we followed in Nos.~1 to~4), we restricted
ourselves to affine group preschemes over \(S\), relying essentially on
the representability theorems of Exposé~XI, No.~4. In the present notes
(cf. Nos.~6 to~8), we show quickly how the affineness hypotheses can be
eliminated by a simpler method that does not use the most delicate
results of Exposé~XI. For further generalizations of the results
contained in the present Exposé, see also Exposés~XV and~XVI. (It is
clear that the contents of Exposés~XI, XII, XV, and~XVI ought to be
completely rewritten.)

\section{Maximal tori}
\SGARefTarget{sga3:xii:section:1}

\paragraph{\textbf{1.0.}}
\SGARefTarget{sga3:xii:1.0}
Let first \(G\) be an algebraic group over an algebraically closed
field \(k\). A \emph{maximal torus of \(G\)} is an algebraic subgroup
\(T\) of \(G\) which is a torus (here, since \(k\) is algebraically
closed, this means that it is isomorphic to a group of the form
\(\mathbf G_m^r\)) and which is maximal for this property. Notice that,
since \(k\) is perfect, \(G_{\mathrm{red}}\) is a subgroup of \(G\),
smooth over \(k\), and hence is essentially an algebraic group in the
sense of \emph{BIBLE}. Moreover, every reduced subgroup of \(G\) (such
as a torus) is automatically a subgroup of \(G_{\mathrm{red}}\).
Consequently the maximal tori of \(G\) are precisely the maximal tori
of \(G_{\mathrm{red}}\). When \(G\) is affine, and therefore
\(G_{\mathrm{red}}\) is affine, a fundamental theorem of Borel says
that any two maximal tori of \(G\) are conjugate by an element of
\[
  G(k)=G_{\mathrm{red}}(k)
\]
(\emph{BIBLE}~6, Theorem~4(c)); in particular, they have the same
dimension. We call the common dimension of the maximal tori of \(G\)
the \emph{reductive rank of \(G\)}. The assumption that \(G\) be affine
is in fact unnecessary. This follows from a theorem of Chevalley
asserting that every connected smooth algebraic group over \(k\) is an
extension of an abelian variety by an affine group; cf. No.~6. In
Nos.~1 to~4 we shall nevertheless most often restrict ourselves to
affine group preschemes over the base.

Let now \(G\) be a smooth algebraic group over \(k\), and let \(T\) be
a maximal torus of \(G\). The centralizer \(C(T)\) of \(T\) in \(G\)
will be called the \emph{Cartan subgroup of \(G\) associated with
\(T\)}. For us it is a group subprescheme defined by \SGARefLink{sga3:viii:corollary:6:7}{VIII~6.7}. Since
\(G\) is smooth over \(k\), however, its Cartan subgroup \(C\) is also
smooth over \(k\), by \SGARefLink{sga3:xi:corollary:5:3}{XI~5.3}. Thus \(C\) is in this case the unique
group subprescheme of \(G\), smooth over \(k\) (that is, an algebraic
subgroup of \(G\) in the sense of \emph{BIBLE}), for which \(C(k)\) is
the subgroup of \(G(k)\) centralizing \(T(k)\); in other words, it is
essentially the centralizer in the sense of \emph{BIBLE}.

By the conjugacy theorem just cited, the Cartan subgroups associated
with the various maximal tori are conjugate to one another and
therefore have the same dimension. We call their common dimension the
\emph{nilpotent rank of \(G\)}; it is equal to that of
\(G_{\mathrm{red}}\). If \(\rho_r(G)\) and \(\rho_n(G)\) denote the
reductive and nilpotent ranks of \(G\), then
\[
  \rho_r(G)\leq \rho_n(G),
\]
and the difference
\[
  \rho_u(G)=\rho_n(G)-\rho_r(G)=\dim(C/T)
\]
might be called, when \(G\) is affine, the \emph{unipotent rank of
\(G\)}. If \(G\) is smooth, affine, and connected, then \(C\) is a
connected nilpotent algebraic group (\emph{BIBLE}~6,
Theorem~6(a),(c)). Hence, by \emph{BIBLE}~6, Theorem~2, it is
isomorphic to the product
\[
  C_s\times C_u,
\]
where \(C_s=T\) is the original maximal torus (and \emph{a fortiori}
a maximal torus of \(C\)), while \(C_u\) is a smooth unipotent
subgroup, that is, a successive extension of groups isomorphic to
\(\mathbf G_a\) (\emph{BIBLE}~6, Theorem~1, Corollary~1, and
\emph{BIBLE}~7, Theorem~4). In this case one also has
\[
  \rho_u(G)=\dim C_u.
\]

\paragraph{\textbf{Remark 1.1.}}
\SGARefTarget{sga3:xii:1.1}
Besides the three notions of rank just specified for an affine
algebraic group, two others are useful. The first is the
\emph{semisimple rank} \(\rho_s(G)\), defined as the reductive rank of
the quotient \(G/R\), where \(R\) is the radical of \(G\). The second
is the \emph{infinitesimal rank} \(\rho_i(G)\), defined as the
nilpotent rank of the Lie algebra of \(G\), which will be defined and
studied in the following Exposé. We shall not use them in the present
Exposé. We record only the inequalities
\[
  \rho_s\leq \rho_r\leq \rho_n\leq \rho_i.
\]

\paragraph{\textbf{Lemma 1.2.}}
\SGARefTarget{sga3:xii:1.2}
\emph{Let \(G\) be an algebraic group over the algebraically closed
field \(k\), let \(T\) be an algebraic subgroup of \(G\), and let
\(k'\) be an algebraically closed extension of \(k\). Let \(G'\) and
\(T'\) be the groups obtained from \(G\) and \(T\) by base extension.
Then \(T\) is a maximal torus of \(G\) if and only if \(T'\) is a
maximal torus of \(G'\).}

This is an immediate consequence of the ``principle of finite
extension'' (EGA~IV, 9.1.1).

\paragraph{\textbf{Definition 1.3.}}
\SGARefTarget{sga3:xii:1.3}
Let \(S\) be a prescheme, let \(G\) be an \(S\)-group prescheme of
finite type, and let \(T\) be a group subprescheme of \(G\). One says
that \(T\) is a \emph{maximal torus of \(G\)} if:
\begin{enumerate}
\renewcommand{\labelenumi}{\alph{enumi})}
\item \(T\) is a torus (\SGARefLink{sga3:ix:definition:1:3}{IX~1.3}); and
\item for every \(s\in S\), if \(\bar s\) denotes the spectrum of an
  algebraic closure of \(\kappa(s)\), then \(T_{\bar s}\) is a
  maximal torus of \(G_{\bar s}\).
\end{enumerate}

\paragraph{\textbf{Remarks 1.4.}}
\SGARefTarget{sga3:xii:1.4}
It follows from \SGARefLink{sga3:xii:1.2}{1.2} that, when \(S\) is the spectrum of an
algebraically closed field, this recovers the usual definition, and
that property~\SGARefLink{sga3:xii:1.3}{1.3} is stable under arbitrary base change. Notice that,
by \SGARefLink{sga3:x:theorem:8-8}{X~8.8}, condition~(a) in \SGARefLink{sga3:xii:1.3}{Definition~1.3} may be replaced by:
\begin{quote}
\textit{a\(^{\prime}\)) \(T\) is of finite presentation and flat over
\(S\).}
\end{quote}

One should take care: a maximal torus in the sense of \SGARefLink{sga3:xii:1.3}{Definition~1.3}
is indeed maximal among the subtori of \(G\), as follows immediately
from \SGARefLink{sga3:ix:corollary:2:9}{IX~2.9}, but the converse cannot hold---say when the base \(S\) is
connected---unless \(G\) actually admits a maximal torus in the sense
of \SGARefLink{sga3:xii:1.3}{1.3}. This is not true in general, even when
\(S=\operatorname{Spec}(\mathbb Z)\) and \(G\) is ``semisimple''; see
also \SGARefLink{sga3:xii:1.6}{1.6}. We shall see in Exposé~XIV, however, that the converse is
true when \(S\) is artinian, or when \(S\) is a local scheme and \(G\)
is ``reductive'': in that case every torus of \(G\) is contained in a
maximal torus.

\paragraph{\textbf{Definition 1.5.}}
\SGARefTarget{sga3:xii:1.5}
Let \(G\) be an algebraic group over a field \(k\). The
\emph{reductive rank} (respectively the \emph{nilpotent rank}, the
\emph{unipotent rank}, and so on) of \(G\) is the reductive rank
(respectively, and so on) of \(G_{\bar k}\), where \(\bar k\) is an
algebraic closure of \(k\).

By \SGARefLink{sga3:xii:1.2}{1.2} and the compatibility of the formation of
\(\operatorname{Centr}_G(T)\) with base extension, the notions of rank
introduced in \SGARefLink{sga3:xii:1.5}{1.5} are invariant under extension of the base field.
For algebraically closed \(k\), they coincide with the notions
introduced at the beginning of the present number.

\paragraph{\textbf{Remark 1.6.}}
\SGARefTarget{sga3:xii:1.6}
It is not difficult to construct a smooth affine group scheme \(G\)
over the spectrum \(S\) of a discrete valuation ring whose generic
fiber is isomorphic to \(\mathbf G_m\) and whose special fiber is
isomorphic to \(\mathbf G_a\). Such a \(G\) contains no torus other
than the trivial torus \(T\), reduced to the identity subgroup, and
\(T\) is plainly not a maximal torus. More precisely, in the special
fiber \(G_0=\mathbf G_{a,k}\), the torus \(T_0\) is indeed maximal,
whereas in the generic fiber \(G_1=\mathbf G_{m,K}\), the torus
\(T_1\) is no longer maximal; here \(k\) is the residue field and
\(K\) the fraction field of the valuation ring. This example also
shows that the reductive rank of \(G_s\), for \(s\in S\), is not a
continuous function of \(s\). Nevertheless, one has the following
results.

\paragraph{\textbf{Theorem 1.7.}}
\SGARefTarget{sga3:xii:1.7}
\emph{Let \(S\) be a prescheme and let \(G\) be a smooth affine
\(S\)-group prescheme. For every \(s\in S\), set
\[
  \rho_r(s)=\rho_r(G_s),\qquad
  \rho_n(s)=\rho_n(G_s),
\]
the reductive and nilpotent ranks of \(G_s\) as in \SGARefLink{sga3:xii:1.5}{1.5}. Then:}

\begin{enumerate}
\renewcommand{\labelenumi}{\alph{enumi})}
\item\SGAInlineRefTarget{sga3:xii:1.7:item:a} \emph{The function \(\rho_r\) on \(S\) is lower
  semicontinuous, while \(\rho_n\) is upper semicontinuous. Hence
  \(\rho_u=\rho_n-\rho_r\) is upper semicontinuous.}

\item\SGAInlineRefTarget{sga3:xii:1.7:item:b} \emph{The following conditions, each stable under arbitrary
  base change, are equivalent:}
  \begin{enumerate}
  \renewcommand{\labelenumii}{(\roman{enumii})}
  \item \emph{the function \(\rho_r\) on \(S\) is locally constant;}
  \item \emph{\(G\) admits a maximal torus locally for the étale
    topology;}
  \item[\textup{(ii bis)}] \emph{\(G\) admits a maximal torus locally
    for the faithfully flat quasi-compact topology.}
  \end{enumerate}

\item \emph{Let \(T_1,T_2\) be two maximal tori of \(G\), which in
  particular places us under the conditions of (b). Then \(T_1\) and
  \(T_2\) are locally conjugate for the étale topology: there exists a
  surjective étale morphism \(S'\to S\) such that the subgroups
  \((T_1)_{S'}\) and \((T_2)_{S'}\) of \(G_{S'}\) are conjugate by a
  section of \(G_{S'}\) over \(S'\).}

\item \emph{Under the conditions of (b), that is, when \(\rho_r\) is
  locally constant, \(\rho_n\) is locally constant as well, and hence
  so is \(\rho_u=\rho_n-\rho_r\).}
\end{enumerate}

\emph{Proof.}
a) Observe that for every morphism \(S'\to S\), if
\(G'=G\times_S S'\), then the functions \(\rho'_r\), and so on, on
\(S'\), defined in terms of \(G'\) as \(\rho_r\), and so on, are
defined in terms of \(G\), are obtained simply by composing the latter
functions with \(S'\to S\). When \(S'\to S\) is faithfully flat and
quasi-compact, it follows that \(\rho'\) is continuous, respectively
upper semicontinuous, respectively lower semicontinuous, if and only
if \(\rho\) is, since the Zariski topology on \(S\) is the quotient of
that on \(S'\) (SGA~1, \SGARefLink{sga3:viii:corollary:4:3}{VIII~4.3}). Consequently, the assertions in (a)
are local for the faithfully flat quasi-compact topology.

Let \(s\in S\). We wish to show that the set \(U\) of \(t\in S\) such
that
\[
  \rho_r(t)\geq \rho_r(s)
  \qquad
  \bigl(\text{respectively }\rho_n(t)\leq \rho_n(s)\bigr)
\]
is a neighborhood of \(s\). By the principle of finite extension,
there exists a finite extension \(k\) of \(\kappa(s)\) such that
\(G_k\) admits a maximal torus. There then exist an open neighborhood
\(U\) of \(s\) and a finite, surjective, locally free morphism
\(S'\to U\) such that the fiber \(S'_s\) is
\(\kappa(s)\)-isomorphic to \(\operatorname{Spec}(k)\) (cf. EGA~III,
10.3.2, where the noetherian hypothesis is plainly unnecessary).
Since \(S'\to S\) is an open morphism (SGA~1, \SGARefLink{sga3:iv:subsection:6:6}{IV~6.6}), we are reduced
to the case \(S'=S\), that is, to the case in which \(G_s\) admits a
maximal torus \(T_s\). Moreover, by \SGARefLink{sga3:xi:theorem:5:8:item:a}{XI~5.8(a)}, after possibly replacing
\(S\) once more by an étale \(S'\) over \(S\), equipped with a point
\(s'\) above \(s\), we may suppose that \(T_s\) is the fiber at \(s\)
of a subtorus \(T\) of \(G\). Then, for every \(t\in S\),
\[
  \rho_r(t)=\rho_r(G_t)
  \geq \dim T_t
  =\dim T_s
  =\rho_r(G_s)
  =\rho_r(s),
\]
which proves that \(\rho_r\) is lower semicontinuous.

On the other hand, by \SGARefLink{sga3:xi:corollary:5:3}{XI~5.3}, the functor
\[
  \operatorname{Centr}_G(T)
\]
is represented by
\[
  C=\operatorname{Centr}_G(T),\footnote{Editor's note: the text has
  been modified here to introduce the prescheme
  \(C=\operatorname{Centr}_G(T)\) representing the functor
  \(\operatorname{Centr}_G(T)\).}
\]
a closed group subprescheme of \(G\), smooth over \(S\). There is
therefore a neighborhood \(U\) of \(s\) such that \(t\in U\) implies
\[
  \dim C_t=\dim C_s=\rho_n(s).
\]
The upper semicontinuity of \(\rho_n\) now follows from the relation
\[
  \rho_n(t)\leq \dim C_t
  \qquad\text{for every }t\in S,
\]
which is itself contained in the following purely geometric lemma.

\paragraph{\textbf{Lemma 1.8.}}
\SGARefTarget{sga3:xii:1.8}
\emph{Let \(k\) be a field, let \(G\) be a smooth affine algebraic
group over \(k\), let \(T\) be a torus in \(G\), and let \(C\) be its
centralizer. Then}
\[
  \rho_n(G)\leq \dim C.
\]

Indeed, we may suppose \(k\) algebraically closed, so that \(T\) is
contained in a maximal torus \(T'\). If \(C'\) is the centralizer of
\(T'\), then \(C'\subset C\), and hence
\(\dim C'\leq \dim C\). This proves the lemma.

b) If \(\rho_r\) is locally constant, then for every torus \(T\) in
\(G\) and every \(s\in S\), if \(T_s\) is a maximal torus in \(G_s\),
there exists an open neighborhood \(U\) of \(s\) such that
\(T|_U\) is a maximal torus in \(G|_U\). Applying the argument of (a),
we see that (i) implies (ii bis). Conversely, (ii bis) implies (i),
because when \(G\) admits a maximal torus \(T\), it is clear that
\[
  \rho_r(s)=\dim T_s
\]
is a locally constant function of \(s\); and we observed at the
beginning of the proof of (a) that continuity of \(\rho_r\) is local
for the faithfully flat topology. Thus (i) is equivalent to
(ii bis), while (ii) plainly implies (ii bis); it remains to prove
the converse implication (i) \(\Rightarrow\) (ii).

For this purpose, introduce the functor \(F\) of \SGARefLink{sga3:xi:theorem:4:1}{XI~4.1},\footnote{Editor's
note: this is the functor of multiplicative-type subgroups of \(G\).}
which is a smooth separated prescheme over \(S\), and consider the
subfunctor \(\mathcal T\) of \(F\) whose value on an \(S'\) over \(S\)
is the set of maximal tori in \(G_{S'}\). By the preceding observation
that, under (i), a torus in \(G_{S'}\) which is maximal in the fiber
over a point \(s\in S'\) is maximal over an open neighborhood of
\(s\), the functor \(\mathcal T\) is represented by an open
subprescheme of \(F\), and is therefore smooth and separated over
\(S\). Since the structural morphism \(\mathcal T\to S\) is plainly
surjective, it admits a section locally for the étale topology by
\SGARefLink{sga3:xi:proposition:1:10}{XI~1.10}. This proves (i) \(\Rightarrow\) (ii).

c) This is an immediate consequence of \SGARefLink{sga3:xi:corollary:5:4:bis}{XI~5.4 bis}, taking account of
Borel's conjugacy theorem recalled at the beginning of this number.

d) By the observation at the beginning of the proof of (a), and in
view of (b), we may suppose that \(G\) admits a maximal torus \(T\).
If \(C\) is its centralizer, then \(C\) is representable and smooth
over \(S\) by \SGARefLink{sga3:xi:corollary:5:3}{XI~5.3}, so the function
\[
  s\longmapsto \rho_n(s)=\dim C_s
\]
is indeed locally constant.

This completes the proof of \SGARefLink{sga3:xii:1.7}{1.7}. We shall refer to the conditions in
\SGARefLink{sga3:xii:1.7:item:b}{1.7(b)} by saying that, in this case, \(G\) is \emph{of locally
constant reductive rank}. We record the following consequence.

\paragraph{\textbf{Corollary 1.9.}}
\SGARefTarget{sga3:xii:1.9}
\emph{Let \(G\) be as in \SGARefLink{sga3:xii:1.7}{1.7}, and let \(s\in S\) be such that
\(\rho_u(s)=0\), that is,
\(\rho_r(s)=\rho_n(s)\) (equivalently, the Cartan subgroups of
\(G_{\bar k}\) are tori, where \(\bar k\) is an algebraic closure of
\(\kappa(s)\)). Then there exists an open neighborhood \(U\) of \(s\)
on which \(\rho_r\) and \(\rho_n\) are constant. In particular, for
every \(t\in U\), the unipotent rank \(\rho_u(t)\) of \(G_t\) is
zero.}

Indeed, this follows immediately from \SGARefLink{sga3:xii:1.7:item:a}{1.7(a)} and the inequality
\(\rho_r(t)\leq \rho_n(t)\) for every \(t\in S\).

Notice also that, in the course of proving (b), we proved the
following result.

\paragraph{\textbf{Corollary 1.10.}}
\SGARefTarget{sga3:xii:1.10}
\emph{Let \(G\) be as in \SGARefLink{sga3:xii:1.7}{1.7}, and suppose that \(G\) has locally
constant reductive rank. Consider the functor
\[
  \mathcal T:(\mathbf{Sch})^\circ\longrightarrow(\mathbf{Ens})
\]
such that, for every \(S'\) over \(S\),
\[
  \mathcal T(S')=
  \{\text{maximal tori of }G_{S'}\}.
\]
Then \(\mathcal T\) is represented by a smooth, separated prescheme of
finite type over \(S\).}

It remains only to verify that \(\mathcal T\) is of finite type over
\(S\). Since the question is local for the faithfully flat
quasi-compact topology, we may suppose that \(G\) admits a maximal
torus \(T\). By \SGARefLink{sga3:xi:corollary:5:3:bis}{XI~5.3 bis} and the bis version of \SGARefLink{sga3:xi:remark:5:5}{XI~5.5},
\(\operatorname{Norm}_G(T)\) and
\(G/\operatorname{Norm}_G(T)\) are represented by preschemes
\(\operatorname{Norm}_G(T)\) and
\(G/\operatorname{Norm}_G(T)\), and \(\mathcal T\) is isomorphic to
\[
  G/\operatorname{Norm}_G(T).\footnote{Editor's note: the text has
  been modified here to introduce the preschemes
  \(\operatorname{Norm}_G(T)\) and
  \(G/\operatorname{Norm}_G(T)\).}
\]
The morphism \(G\to\mathcal T\) defined by
\(g\mapsto\operatorname{int}(g)(T)\) is surjective. Since \(G\) is
quasi-compact over \(S\), so is \(\mathcal T\), which completes the
proof.

\paragraph{\textbf{Remarks 1.11.}}
\SGARefTarget{sga3:xii:1.11}
a) The prescheme \(\mathcal T\) of \SGARefLink{sga3:xii:1.10}{1.10} will properly be called the
\emph{prescheme of maximal tori of \(G\)}. We shall see in No.~5 that
it is in fact affine over \(S\). This can be checked directly when
\(G\) is isomorphic to a closed group subprescheme of a prescheme of
the form \(\operatorname{GL}(n)\), using \SGARefLink{sga3:xi:remark:4:6}{XI~4.6} and observing that, by
the argument in the proof of \SGARefLink{sga3:xii:1.7:item:b}{1.7(b)}, \(\mathcal T\) identifies with a
subprescheme both open and closed in the prescheme \(F\) representing
the multiplicative-type subgroups of \(G\).

b) It is possible to prove \SGARefLink{sga3:xii:1.7}{1.7}, and hence \SGARefLink{sga3:xii:1.9}{1.9}, without using the
delicate results of Exposé~XI, using only the elementary results
\SGARefLink{sga3:xi:proposition:3:12}{XI~3.12} and \SGARefLink{sga3:xi:theorem:6:2}{XI~6.2} and working solely with multiplicative-type groups
finite over the base (morally, the groups \({}_nT\), where \(T\) is a
maximal torus); compare No.~7. The same observation applies to the
proof of \SGARefLink{sga3:xii:1.10}{1.10}.

We end this number with examples in which there is a unique maximal
torus.

\paragraph{\textbf{Proposition 1.12.}}
\SGARefTarget{sga3:xii:1.12}
\emph{Let \(G\) be an \(S\)-group prescheme of multiplicative type
and finite type. Then \(G\) admits a unique maximal torus, and every
torus in \(G\) is contained in this maximal torus.}

Uniqueness plainly follows from the latter assertion, which
characterizes the maximal torus as the largest subtorus of \(G\).
Uniqueness also makes existence a question local for the faithfully
flat quasi-compact topology. We may therefore suppose that \(G\) is
diagonalizable, say \(G=D_S(M)\), where \(M\) is a finitely generated
commutative group. Let \(M_0\) be the quotient of \(M\) by its torsion
subgroup. We claim that the torus
\[
  T=D_S(M_0)
\]
in \(G\) is a maximal torus and the largest subtorus. Indeed, a
subtorus \(T'\) of \(G\) is locally diagonalizable for the fpqc
topology. Thus, in order to prove \(T'\subset T\), we may suppose that
\(T'\) is diagonalizable, and hence of the form \(D_S(N)\), where
\(N\) is a free quotient of \(M\) (\SGARefLink{sga3:viii:corollary:1:4}{VIII~1.4} and \SGARefLink{sga3:viii:corollary:3:2:item:b:01}{3.2(b)}). It follows
that \(N\) is a quotient of \(M_0\), and therefore that
\(T'\subset T\). Since the construction of \(T=D_S(M_0)\) is
compatible with arbitrary base extension, this proves at the same
time that \(T\) is a maximal torus of \(G\), and completes the proof.
When \(G\) is smooth over \(S\), \SGARefLink{sga3:xii:1.12}{Proposition~1.12} has the following
generalization.

\paragraph{\textbf{Proposition 1.13.}}
\SGARefTarget{sga3:xii:1.13}
\emph{Let \(G\) be an \(S\)-group prescheme of finite presentation
over \(S\). Suppose that, locally for the faithfully flat
quasi-compact topology, \(G\) admits a central maximal torus. Then
\(G\) admits globally a unique maximal torus,\footnote{Which is
central!} and it is the largest subtorus of \(G\).}

Here again, uniqueness follows trivially from the last assertion and
makes existence a local question for the fpqc topology. We may
therefore suppose that \(G\) admits a central maximal torus \(T\). We
shall show that every torus \(R\) of \(G\) is contained in \(T\). This
follows from the next lemma.

\paragraph{\textbf{Lemma 1.14.}}
\SGARefTarget{sga3:xii:1.14}
\emph{Let \(G\) be an \(S\)-group prescheme of finite presentation
over \(S\), let \(T\) be a maximal torus of \(G\), let \(R\) be a
subtorus of \(G\), and suppose that \(T\) and \(R\) commute. Then
\(R\subset T\).}

Indeed, since \(R\) and \(T\) commute, the morphism
\[
  R\times_S T\longrightarrow G,\qquad (r,t)\longmapsto r\cdot t,
\]
is a group homomorphism. By \SGARefLink{sga3:ix:theorem:6:8}{IX~6.8} it therefore has an image subgroup
\(T'\) in \(G\), which is a multiplicative-type group quotient of
\(R\times_S T\), hence a torus, and plainly contains \(T\). Since
\(T\) is a maximal torus, \(T'=T\) by \SGARefLink{sga3:xii:1.4}{1.4}, and therefore
\(R\subset T\). This proves \SGARefLink{sga3:xii:1.14}{1.14} and hence \SGARefLink{sga3:xii:1.13}{1.13}. In particular, using
\SGARefLink{sga3:xii:1.7:item:b}{1.7(b)}, we obtain:

\paragraph{\textbf{Corollary 1.15.}}
\SGARefTarget{sga3:xii:1.15}
\emph{Let \(G\) be a commutative, smooth, affine \(S\)-group
prescheme of locally constant reductive rank. Then \(G\) admits a
unique maximal torus, and this torus contains every subtorus of
\(G\).}

\paragraph{\textbf{Corollary 1.16.}}
\SGARefTarget{sga3:xii:1.16}
\emph{Let \(G\) be a smooth affine \(S\)-group prescheme. Suppose
that, for every \(s\in S\), if \(\bar s\) denotes the spectrum of an
algebraic closure \(\bar k\) of \(\kappa(s)\), the geometric fiber
\(G_{\bar s}\) is a connected nilpotent algebraic group (in the sense
of \emph{BIBLE}, that is, the group \(G(\bar k)\) of its
\(\bar k\)-valued points is nilpotent). Suppose moreover that \(G\)
has locally constant reductive rank. Then \(G\) admits a unique
maximal torus \(T\); moreover, \(T\) is central and is the largest
subtorus of \(G\).}

By 1.7(b), \(G\) admits a maximal torus locally for the fpqc topology,
and by \SGARefLink{sga3:xii:1.13}{1.13} we are reduced to proving that a maximal torus of \(G\) is
central. By \SGARefLink{sga3:ix:theorem:5:6:item:b}{IX~5.6(b)}, we are reduced to the case in which \(S\) is
the spectrum of a field, which we may suppose algebraically closed.
Then \(T(k)\) lies in the center of \(G(k)\) by \emph{BIBLE}~6,
\SGARefLink{sga3:xii:section:2}{Theorem~2}. This implies that \(T\) lies in the center of \(G\), because
\(\operatorname{Centr}_G(T)\) is a closed subscheme of \(G\)
(VIII~6.6) containing all the points of \(G(k)\), and is therefore
equal to \(G\) by the Nullstellensatz, since \(G\) is reduced.

Finally, for later reference, we record the following elementary
proposition, which we have already used implicitly.

\paragraph{\textbf{Proposition 1.17.}}
\SGARefTarget{sga3:xii:1.17}
\emph{Let \(G\supset H\supset T\) be group preschemes of finite type
over \(S\). The following conditions are equivalent:}
\begin{enumerate}
\renewcommand{\labelenumi}{(\roman{enumi})}
\item \emph{\(T\) is a maximal torus of \(G\);}
\item \emph{\(T\) is a maximal torus of \(H\), and for every
  \(s\in S\) one has equality of reductive ranks
  \[
    \rho_r(G_s)=\rho_r(H_s).
  \]}
\end{enumerate}

\section{The Weyl group}
\SGARefTarget{sga3:xii:section:2}

Let first \(k\) be an algebraically closed field, and let \(G\) be a
smooth affine algebraic group over \(k\). If \(T\) is a maximal torus,
\(C\) its centralizer, and \(N\) its normalizer, then by \SGARefLink{sga3:xi:proposition:5:9}{XI~5.9} these
are smooth closed subgroups of \(G\), and \(C\) is an open subgroup of
\(N\). Thus
\[
  W=N/C
\]
is a finite étale group over \(k\), and is consequently determined by
the group \(W(k)\) of its \(k\)-valued points:
\[
  W=W(k)_k
\]
(the constant group defined by the ordinary finite group \(W(k)\)).
The finite group \(W(k)\) will be called the \emph{geometric Weyl
group}, or simply the \emph{Weyl group} when no confusion is likely,
of \(G\) relative to \(T\). By Borel's conjugacy theorem, the Weyl
groups relative to the various maximal tori are isomorphic to one
another; this is why one sometimes speaks of ``the'' Weyl group of
\(G\), without specifying a maximal torus. Since the formation of
\(C\), \(N\), and \(N/C\) commutes with arbitrary base extension, if
\(k'\) is an algebraically closed extension of \(k\), the geometric
Weyl group of \(G_{k'}\) relative to \(T_{k'}\) is canonically
isomorphic to that of \(G\) relative to \(T\). Consequently, ``the''
geometric Weyl group of \(G\), properly speaking an isomorphism class
of ordinary finite groups, coincides with that of \(G_{k'}\).

This permits us, when \(G\) is a smooth affine algebraic group over an
arbitrary field \(k\), to speak of the \emph{geometric Weyl group of
\(G\)} as the isomorphism class of the geometric Weyl group of
\(G_{k'}\), where \(k'\) is any algebraically closed extension of
\(k\). When \(G\) admits a maximal torus \(T\), we can form
\[
  C=\operatorname{Centr}_G(T),\qquad
  N=\operatorname{Norm}_G(T),\qquad
  W=N/C
\]
as above. The latter is a finite étale group over \(k\), called the
\emph{Weyl group of \(G\) relative to \(T\)}. The geometric Weyl group
is then simply the class of the group of points of \(W\) with values
in any algebraically closed extension \(k'\) of \(k\). In general,
knowledge of the geometric Weyl group \(W(k')\) plainly no longer
suffices to reconstruct the algebraic group \(W\): one must also know
the action on \(W(k')\) of the Galois group of the separable algebraic
closure of \(k\) in \(k'\).

Finally, let \(G\) be a group prescheme over an arbitrary base \(S\),
smooth and affine over \(S\), and let \(T\) be a maximal torus of
\(G\). By \SGARefLink{sga3:xi:proposition:5:9}{XI~5.9} we may again form the group
\[
  W(T)=N(T)/C(T),
\]
which is an étale, separated, quasi-finite \(S\)-group prescheme. Its
geometric fibers, relative to algebraic closures of the residue
fields \(\kappa(s)\), \(s\in S\), are the geometric Weyl groups of the
fibers \(G_s\). Thus \SGARefLink{sga3:xi:corollary:5:10}{XI~5.10} gives information about the variation of
these groups with \(s\in S\). We can make this information more
precise and generalize it as follows.

\paragraph{\textbf{Theorem 2.1.}}
\SGARefTarget{sga3:xii:2.1}
\emph{Let \(S\) be a prescheme and let \(G\) be a smooth affine
\(S\)-group prescheme. For every \(s\in S\), let \(w(s)\) be the
geometric Weyl group of \(G_s\), hence an isomorphism class of finite
groups. On the set \(E\) of such classes, introduce the following
preorder: \(w\leq w'\) if and only if \(w\) and \(w'\) can be
represented by finite groups \(W\) and \(W'\) such that \(W\) is
isomorphic to a quotient of a subgroup of \(W'\). With this
convention:}
\begin{enumerate}
\renewcommand{\labelenumi}{\alph{enumi})}
\item\SGAInlineRefTarget{sga3:xii:2.1:item:a} \emph{The function \(s\mapsto w(s)\) from \(S\) to \(E\) is
  lower semicontinuous.}
\item\SGAInlineRefTarget{sga3:xii:2.1:item:b} \emph{Suppose that \(G\) has locally constant reductive rank.
  Then the following conditions are equivalent:}
  \begin{enumerate}
  \renewcommand{\labelenumii}{(\roman{enumii})}
  \item \emph{the function \(s\mapsto w(s)\) is locally constant;}
  \item[\textup{(i bis)}] \emph{the function
    \(s\mapsto\operatorname{card}w(s)\) is locally constant;}
  \item \emph{locally for the étale topology (or merely for the fpqc
    topology), there exists a maximal torus \(T\) such that \(W(T)\)
    is finite over its base;}
  \item[\textup{(ii bis)}] \emph{for every \(S'\) over \(S\) and
    every maximal torus \(T\) of \(G_{S'}\), the associated Weyl group
    \(W(T)\) is finite over \(S'\).}
  \end{enumerate}
\end{enumerate}

\emph{Proof.}
a) Arguing as in \SGARefLink{sga3:xii:1.7:item:a}{1.7(a)}, in order to prove that for every \(s\in S\)
there exists an open neighborhood \(U\) of \(s\) such that
\(t\in U\) implies \(w(t)\geq w(s)\), we are reduced to the case in
which there is a torus \(R\) in \(G\) such that \(R_s\) is a maximal
torus in \(G_s\). Put
\[
  W(R)=\operatorname{Norm}_G(R)/\operatorname{Centr}_G(R).
\]
As in \SGARefLink{sga3:xi:proposition:5:9}{XI~5.9}, this is an étale, separated, quasi-finite group
prescheme over \(S\). For \(t\in S\), let \(w'(t)\in E\) be its
geometric fiber at \(t\). Since \(R_s\) is a maximal torus in \(G_s\),
and the formation of \(\operatorname{Norm}\),
\(\operatorname{Centr}\), and
\(\operatorname{Norm}/\operatorname{Centr}\) commutes with arbitrary
base extension, in particular with passage to fibers, one has
\[
  w(s)=w'(s).
\]
We claim moreover that for \(t\) near \(s\),
\[
  w(t)\geq w'(t)
  \qquad\text{and}\qquad
  w'(t)\geq w'(s),
\]
which will suffice to prove (a). These two inequalities are contained
in the following two lemmas.

\paragraph{\textbf{Lemma 2.2.}}
\SGARefTarget{sga3:xii:2.2}
\emph{Let \(S\) be a prescheme, and let \(W\) be an étale,
separated, quasi-finite \(S\)-group prescheme. For every \(s\in S\),
let \(f(s)\) be the class of the geometric fiber of \(W\) at \(s\),
an element of the ordered set \(E\) of isomorphism classes of finite
groups introduced in \SGARefLink{sga3:xii:2.1}{2.1}. Then the function
\(f:S\to E\) is lower semicontinuous. It is constant in a
neighborhood of \(s\in S\) if and only if the function
\(s\mapsto\operatorname{card}f(s)\) is constant there, and this holds
if and only if \(W\) is finite over \(S\) above some open
neighborhood \(U\) of \(s\).}

This result refines, in group-theoretic terms, the result used in the
proof of \SGARefLink{sga3:xi:corollary:5:10}{XI~5.10}. We give only a sketch of the proof, which is entirely
standard. As usual, we are reduced to the case in which \(S\) is
affine and noetherian. The function \(f\) is immediately seen to be
constructible (EGA \(0_{\mathrm{III}}\), 9.3.1 and 9.3.2, together
with the sorites of EGA~IV, 9), and EGA
\(0_{\mathrm{III}}\), 9.3.4 reduces semicontinuity to proving that if
\(t\) is a generalization of \(s\), then
\[
  f(t)\geq f(s).
\]
By EGA~II, 7.1.7, this reduces us to the case in which \(S\) is the
spectrum of a discrete valuation ring, which may be supposed complete
with algebraically closed residue field. Let \(s\) be the closed point
of \(S\). Since \(W\) is étale and separated over \(S\), it contains
a subscheme \(W'\), both open and closed and finite over \(S\), such
that \(W'_s=W_s\) (EGA~II, 6.2.6); it is immediate here that \(W'\)
is a subgroup of \(W\). Moreover, since \(W'\) is finite étale over
\(S=\operatorname{Spec}(V)\), with \(V\) complete and with
algebraically closed residue field, it is a constant group, hence of
the form \(A_S\), where
\[
  A=W'(\kappa(s))=W(\kappa(s))
\]
has class \(f(s)\). If \(B\) is the geometric fiber of \(W\) at the
generic point \(t\) of \(S\), there is therefore a canonical
monomorphism \(A\to B\), proving \(f(s)\leq f(t)\). Notice that this
argument actually proves semicontinuity for an order on \(E\) finer
than the one specified in \SGARefLink{sga3:xii:2.1}{2.1}.

The equivalence between constancy of \(f\) at \(s\) and constancy of
\(s\mapsto\operatorname{card}f(s)\) follows from the fact that, for
\(w,w'\in E\),
\[
  w\leq w'
  \quad\text{and}\quad
  \operatorname{card}w=\operatorname{card}w'
  \quad\Longrightarrow\quad
  w=w'.
\]
The equivalence of this condition with finiteness of \(W\) over a
neighborhood of \(s\) is independent of the group structure on
\(W\); it was noted after \SGARefLink{sga3:xi:corollary:5:10}{XI~5.10}, and also follows readily from the
preceding argument using the valuative criterion of properness
(EGA~II, 7.3.8).

\paragraph{\textbf{Lemma 2.3.}}
\SGARefTarget{sga3:xii:2.3}
\emph{Let \(G\) be a smooth affine algebraic group over an
algebraically closed field \(k\), and let \(R\subset T\) be two
subtori. Let \(W(R)\) and \(W(T)\) be the two associated finite groups
of \SGARefLink{sga3:xi:proposition:5:9}{XI~5.9}, obtained as the quotient of the normalizer by the
centralizer. Then \(W(R)\) is isomorphic to a quotient of a subgroup
of \(W(T)\).}

Indeed, consider the diagram
\[
\begin{tikzcd}[column sep=large,row sep=large]
T \arrow[r, hook] &
C(T) \arrow[r, hook] \arrow[d, hook] &
C(R) \arrow[d, hook] \\
& N(T)\cap N(R) \arrow[r, hook] \arrow[d, hook] &
N(R) \\
& N(T) &
\end{tikzcd}
\]
Then \((N(T)\cap N(R))/C(T)\) is a subgroup of
\[
  W(T)=N(T)/C(T),
\]
and there is an evident homomorphism
\[
  (N(T)\cap N(R))/C(T)
  \longrightarrow
  W(R)=N(R)/C(R).
\]
Everything reduces to proving that this last homomorphism is
surjective, that is, that for every \(k\)-valued point \(g\) of
\(N(R)\), there exists a \(k\)-valued point \(c\) of \(C(R)\) such
that \(cg\) normalizes \(T\), or equivalently such that
\[
  \operatorname{int}(c)\bigl(\operatorname{int}(g)T\bigr)=T.
\]
For this, it suffices to observe that
\(\operatorname{int}(g)T\) is a torus of \(N(R)\), hence of \(C(R)\),
which is an open subgroup of \(N(R)\). The tori \(T\) and
\(\operatorname{int}(g)T\) are maximal tori of \(C(R)\), since they
are maximal in \(G\), and the assertion follows from Borel's conjugacy
theorem.

This proves \SGARefLink{sga3:xii:2.3}{2.3}, and hence \SGARefLink{sga3:xii:2.1:item:a}{2.1(a)}.

b) We have already observed that (i) and (i bis) are plainly
equivalent. They imply (ii bis), by the converse part of \SGARefLink{sga3:xii:2.2}{2.2} or,
equivalently, \SGARefLink{sga3:xi:corollary:5:10}{XI~5.10}. Condition (ii bis) implies (ii) by \SGARefLink{sga3:xii:1.7:item:b}{1.7(b)}.
Finally, (ii) implies (i bis): as in \SGARefLink{sga3:xii:1.7:item:a}{1.7(a)}, condition (i) is local
for the fpqc topology, so we may suppose that \(G\) admits a maximal
torus \(T\) such that \(W(T)\) is finite over \(S\), and \SGARefLink{sga3:xi:corollary:5:10}{XI~5.10}
again gives the conclusion. This completes the proof of \SGARefLink{sga3:xii:2.1}{2.1}.

\section{Cartan subgroups}
\SGARefTarget{sga3:xii:section:3}

\paragraph{\textbf{Definition 3.1.}}
\SGARefTarget{sga3:xii:3.1}
Let \(G\) be a smooth group prescheme of finite type over a prescheme
\(S\). A \emph{Cartan subgroup of \(G\)} is a group subprescheme
\(C\) of \(G\), smooth over \(S\), such that for every \(s\in S\), if
\(\bar s\) denotes the spectrum of an algebraic closure of
\(\kappa(s)\), then \(C_{\bar s}\) is a Cartan subgroup, in the sense
of No.~1, of \(G_{\bar s}\).

It is immediate that if \(C\) is a Cartan subgroup of \(G\), then for
every \(S'\) over \(S\), \(C_{S'}\) is a Cartan subgroup of
\(G_{S'}\). If \(S\) is the spectrum of an algebraically closed field,
this recovers the notion recalled in No.~1. Finally, the property that
a group subprescheme \(C\) of \(G\) be a Cartan subgroup is plainly
local for the faithfully flat quasi-compact topology.

\paragraph{\textbf{Theorem 3.2.}}
\SGARefTarget{sga3:xii:3.2}
\emph{Let \(G\) over \(S\) be as in \SGARefLink{sga3:xii:3.1}{3.1}, and suppose that \(G\) is
affine over \(S\) and has locally constant reductive rank. Then the
map
\[
  T\longmapsto\operatorname{Centr}_G(T)
\]
induces a bijection from the set of maximal tori of \(G\) onto the set
of Cartan subgroups\footnote{The proof given here in fact proves
\SGARefLink{sga3:xii:3.2}{Theorem~3.2} only for closed Cartan subgroups of \(G\). However, \SGARefLink{sga3:xii:7.1:item:a}{7.1(a)}
establishes \SGARefLink{sga3:xii:3.2}{Theorem~3.2} as stated, and implies that the Cartan
subgroups of \(G\) are closed.} of \(G\). If \(C\) corresponds to
\(T\), then \(T\) is the unique maximal torus of \(C\).}

If \(T\) is a maximal torus of \(G\), the functor
\(\operatorname{Centr}_G(T)\) is represented, by \SGARefLink{sga3:xi:corollary:5:3}{XI~5.3} or \SGARefLink{sga3:xii:6.2}{6.2}, by
a closed smooth subprescheme
\[
  C=\operatorname{Centr}_G(T)
\]
of \(G\), and the definitions show that \(C\) is a Cartan subgroup
of \(G\). Moreover, \(T\) is plainly a maximal torus of \(C\), and,
being central, it is the unique maximal torus of \(C\) by \SGARefLink{sga3:xii:1.13}{1.13}.
Thus the map \(T\mapsto\operatorname{Centr}_G(T)\) from the set of
maximal tori to the set of Cartan subgroups is injective; it remains
to prove surjectivity.

Let \(C\) be a Cartan subgroup of \(G\). We shall prove that
\(C=\operatorname{Centr}_G(T)\) for a maximal torus \(T\) of \(G\).
It is enough to find a maximal torus \(T\) of \(C\), for then \(T\)
is a maximal torus of \(G\): for every \(s\in S\), the groups
\(G_s\) and \(C_s\) have the same reductive rank. By \SGARefLink{sga3:ix:theorem:5:6:item:b}{IX~5.6(b)},
\(T\) lies in the center of \(C\), so
\[
  C\subset C'=\operatorname{Centr}_G(T).
\]
The smooth subgroup \(C\) of the smooth group \(C'\) over \(S\)
coincides with \(C'\) fiber by fiber; hence \(C=C'\).

Since \(G\) has locally constant reductive rank, so does \(C\).
Therefore \(C\) admits a maximal torus locally for the étale topology
by \SGARefLink{sga3:xii:1.7:item:b}{1.7(b)}; that torus is central by the preceding argument, and \SGARefLink{sga3:xii:1.13}{1.13}
then shows that \(C\) itself admits a maximal torus. This completes
the proof.

\paragraph{\textbf{Corollary 3.3.}}
\SGARefTarget{sga3:xii:3.3}
\emph{Let \(G\) be a smooth affine group prescheme over \(S\), of
locally constant reductive rank, and let
\[
  \mathcal C:(\mathrm{Sch})^\circ_{/S}\longrightarrow(\mathrm{Ens})
\]
be the functor defined by
\[
  \mathcal C(S')=
  \{\text{Cartan subgroups of }G_{S'}\}.
\]
Then \(\mathcal C\) is isomorphic to the functor \(\mathcal T\) of
\SGARefLink{sga3:xii:1.10}{1.10}, and is therefore represented by the same smooth separated
prescheme of finite type over \(S\).}

\paragraph{\textbf{Corollary 3.4.}}
\SGARefTarget{sga3:xii:3.4}
\emph{Under the hypotheses of \SGARefLink{sga3:xii:3.2}{3.2}, if
\(C=\operatorname{Centr}_G(T)\), then
\[
  \operatorname{Norm}_G(C)=\operatorname{Norm}_G(T).
\]}

\paragraph{\textbf{Remark 3.5.}}
\SGARefTarget{sga3:xii:3.5}
When \(G\) does not have locally constant reductive rank, it may
nevertheless admit Cartan subgroups. For example, if \(G\) has
connected nilpotent fibers, then \(G\) is a Cartan subgroup of itself,
although it need not have locally constant reductive rank; see \SGARefLink{sga3:xii:1.6}{1.6}.
In Exposé XV we develop the theory of Cartan subgroups without
assuming either that \(G\) is affine over \(S\) or that its reductive
rank is locally constant, using the theory of regular elements from
the next Exposé. See also Nos.~6 and 7 below for the removal of the
affineness hypothesis.

% END INCLUDED SGA3_English_current_buildable_source_20260823/source/inputs/XII/tex/components/00_expose_XII_sections1_3_en.tex

% BEGIN INCLUDED SGA3_English_current_buildable_source_20260823/source/inputs/XII/tex/components/01_expose_XII_sections4_6_en.tex
% Exposé XII Sections 4--6.

\section{The reductive center}
\SGARefTarget{sga3:xii:section:4}

\paragraph{\textbf{Definition 4.1.}}
\SGARefTarget{sga3:xii:4.1}
Let \(S\) be a prescheme, let \(G\) be an \(S\)-group of finite
presentation over \(S\) with affine fibers, and let \(Z\) be a group
subprescheme of \(G\). We say that \(Z\) is a \emph{reductive center}
of \(G\) if:

\begin{enumerate}
\item[(i)] \(Z\) is central and of multiplicative type;
\item[(ii)] for every base change \(S'\to S\) and every central
homomorphism
\[
  u:H\longrightarrow G_{S'},
\]
where \(H\) is a group of multiplicative type and of finite type over
\(S'\), the homomorphism \(u\) factors through \(Z_{S'}\).
\end{enumerate}

(For a variant of this notion when the fibers of \(G\) are not
assumed affine, see \SGARefLink{sga3:xii:8.6}{8.6}.)

Such a \(Z\) is necessarily of finite type over \(S\), since its
fibers are, and is therefore uniquely determined as the largest
central subgroup of multiplicative type in \(G\). It is easy to give
examples in which \(G\), smooth and affine over \(S\), admits a
largest central subgroup \(Z\) of multiplicative type but \(Z\) is
not a reductive center---that is, \(G\) admits no reductive center;
see, for example, \SGARefLink{sga3:xii:1.6}{1.6}. It follows readily, however, from \SGARefLink{sga3:ix:theorem:6:8}{IX~6.8} that
a subgroup \(Z\) of \(G\) is a reductive center if and only if it is
a largest central subgroup of multiplicative type and retains this
property after every base change. See also \SGARefLink{sga3:xii:4.3}{4.3} below.

It is clear from \SGARefLink{sga3:xii:4.1}{4.1} that if \(Z\) is the reductive center of \(G\),
then \(Z_{S'}\) is the reductive center of \(G_{S'}\) after every base
change \(S'\to S\). This fact and the uniqueness of the reductive
center, together with faithfully flat quasi-compact descent
(SGA~1, VIII), give:

\paragraph{\textbf{Proposition 4.2.}}
\SGARefTarget{sga3:xii:4.2}
\emph{Let \(G\) be an \(S\)-group prescheme of finite presentation
over \(S\), with affine fibers. If \(G\) admits a reductive center
locally for the faithfully flat quasi-compact topology, then it
admits a reductive center. For \(Z\) to be a reductive center of
\(G\), it is necessary and sufficient that it be so locally for the
fpqc topology.}

We also record:

\paragraph{\textbf{Proposition 4.3.}}
\SGARefTarget{sga3:xii:4.3}
\emph{Let \(G\) be an \(S\)-group prescheme of finite presentation
with affine fibers, and let \(Z\) be a group subprescheme. For \(Z\)
to be a reductive center of \(G\), it is necessary and sufficient
that it be of multiplicative type and that, for every \(s\in S\),
\(Z_s\) be a reductive center of \(G_s\).}

Indeed, \SGARefLink{sga3:ix:theorem:5:6:item:b}{IX~5.6(b)} first shows that \(Z\) is central. Since the
conditions under consideration are stable under base change, it
remains to prove that every central homomorphism
\[
  u:H\longrightarrow G,
\]
with \(H\) of multiplicative type and of finite type, factors through
\(Z\). Since \(Z\) is central, \(u\) and the canonical immersion
\(v:Z\to G\) define a group homomorphism
\[
  w:H\times_S Z\longrightarrow G.
\]
By \SGARefLink{sga3:ix:theorem:6:8}{IX~6.8}, this homomorphism has a group image \(K\), which is a
subgroup of multiplicative type in \(G\). Everything reduces to
showing that \(K=Z\). This holds fiber by fiber by the hypothesis on
\(Z\), and \SGARefLink{sga3:ix:theorem:5:1:bis}{IX~5.1 bis} now applies, completing the proof of \SGARefLink{sga3:xii:4.3}{4.3}.

In criterion \SGARefLink{sga3:xii:4.3}{4.3}, the hypothesis that \(Z_s\) be the reductive center
of \(G_s\) for every \(s\in S\) is in fact purely geometric: it is
enough to verify it after passing to an algebraic closure of
\(\kappa(s)\), as follows from the second assertion of \SGARefLink{sga3:xii:4.2}{4.2}.

\paragraph{\textbf{Theorem 4.4.}}
\SGARefTarget{sga3:xii:4.4}
\emph{Let \(G\) be an affine algebraic group over a field \(k\).
Then \(G\) admits a reductive center.}

Since the center of \(G\) is represented by a closed subgroup of
\(G\) by \SGARefLink{sga3:viii:corollary:6:7}{VIII~6.7}, we are immediately reduced to the case in which
\(G\) is commutative. Moreover, by \SGARefLink{sga3:xii:4.2}{4.2} we may suppose that \(k\) is
algebraically closed. We shall see in Exposé XVII that in this case
\(G\) can be written as a product \(Z\times U\), where \(Z\) is of
multiplicative type and \(U\) is ``unipotent'', from which it follows
at once that \(Z\) is a reductive center of \(G\). We give here a
proof independent of the general structure theorem for commutative
algebraic groups just indicated.

By \SGARefLink{sga3:ix:theorem:6:8}{IX~6.8}, the set of subgroups \(H\) of multiplicative type in
\(G\) is filtered upward. We show that it has a largest element.

When \(G\) is smooth over \(k\), the structure theorem
BIBLE~4, Th.~4 gives
\[
  G=G_s\times G_u,
\]
where \(G_s\) is of multiplicative type and \(G_u\) is ``unipotent'':
here this means that \(G_u\) has a composition series whose factors
are subgroups of \(\mathbb G_a\), smooth if desired. Indeed,
\(G_u/G_u^0\) is a unipotent group by BIBLE~4, Corollary to Th.~3,
hence a \(p\)-group, where \(p\) is the characteristic exponent, by
BIBLE~4, Prop.~4. On the other hand, \(G_u^0\) has a composition
series with smooth connected one-dimensional quotients by BIBLE~6,
Th.~1, Cor.~1, and these quotients are isomorphic to \(\mathbb G_a\)
by BIBLE~7, Th.~4. Every homomorphism from a group of multiplicative
type to \(\mathbb G_a\), and therefore also to \(G_u\), is trivial,
as follows from the lemma below. Thus every subgroup of
multiplicative type in \(G\) is contained in \(G_s\).

In the general case, consider the subgroup
\[
  G_0=G_{\mathrm{red}}
\]
of \(G\). It is smooth over \(k\), and hence by the preceding case
has a largest subgroup \(Z_0\) of multiplicative type. The subgroups
of multiplicative type in \(G\) containing \(Z_0\) correspond to
the subgroups of multiplicative type in \(G'=G/Z_0\). We have an
exact sequence
\[
  0\longrightarrow G_0/Z_0
   \longrightarrow G/Z_0
   \longrightarrow G/G_0
   \longrightarrow 0,
\]
where \(G_0/Z_0\) has no subgroup of multiplicative type other than
the unit group. Since every subgroup of a group of multiplicative
type is again of multiplicative type by IX~8, every subgroup \(H\)
of multiplicative type in \(G/Z_0\) satisfies
\[
  H\cap(G_0/Z_0)=0;
\]
thus \(H\to G/G_0\) is injective. The algebraic group \(G/G_0\) is
radicial, hence finite over \(k\), so \(H\) is itself radicial and
its rank is bounded above by that of \(G/G_0\). Consequently the
family of subgroups of multiplicative type in \(G\) has a largest
element, say \(Z\).

We claim that \(G/Z\) has no subgroup of multiplicative type other
than the unit subgroup. Indeed, a commutative extension of two
algebraic groups of multiplicative type is of multiplicative type
(\SGARefLink{sga3:xvii:proposition:7:1:1}{7.1.1}). If \(Z'\) is the inverse image in \(G\) of a subgroup of
multiplicative type in \(G/Z\), then \(Z'\) is of multiplicative
type. The maximality of \(Z\) gives \(Z'=Z\), and hence \(Z'/Z=0\).

The ``principle of finite extension'' now readily shows that for
every extension \(K\) of \(k\),
\[
  (G/Z)_K=G_K/Z_K
\]
likewise has no subgroup of multiplicative type other than the unit
group.

We can now prove that \(Z\) is a reductive center of \(G\). Let
\[
  u:H\longrightarrow G_S
\]
be a homomorphism of \(S\)-groups, where \(S\) is a prescheme over
\(k\) and \(H\) is of multiplicative type and of finite type over
\(S\). We must show that \(u\) factors through \(Z_S\), or
equivalently that the composite
\[
  H\longrightarrow G_S\longrightarrow
  (G/Z)_S=G_S/Z_S
\]
is trivial. Put \(U=G/Z\). Every homomorphism
\(H\to U_S\) is trivial: by \SGARefLink{sga3:ix:corollary:5:2}{IX~5.2} one reduces to the case where
\(S\) is the spectrum of a field \(K\), and \SGARefLink{sga3:ix:theorem:6:8}{IX~6.8} then reduces the
claim to the fact that \(U_K\) has no nontrivial subgroup of
multiplicative type. This proves \SGARefLink{sga3:xii:4.4}{4.4}, modulo the following lemma.

\paragraph{\textbf{Lemma 4.4.1.}}
\SGARefTarget{sga3:xii:4.4.1}
\emph{Let \(H\) be an \(S\)-group prescheme of multiplicative type.
Then every homomorphism of \(S\)-groups
\[
  u:H\longrightarrow\mathbb G_a
\]
is trivial.}

Regard \(E=\mathcal O_S^2\) as an extension with submodule
\(E'\simeq\mathcal O_S\) and quotient \(E''\simeq\mathcal O_S\).
Then \(\mathbb G_a\) identifies with the prescheme of automorphisms
of this extension. A homomorphism \(u:H\to\mathbb G_a\) therefore
amounts to an \(H\)-module structure on \(E\) preserving the
extension, with \(E'\) stable under \(H\) and with the induced
actions on \(E'\) and \(E''\) trivial. By \SGARefLink{sga3:i:proposition:4-7-3}{I~4.7.3}, \(H\) acts
trivially on \(E\); hence \(u\) is trivial.

\paragraph{\textbf{Remark 4.4.2.}}
\SGARefTarget{sga3:xii:4.4.2}
If \(G\) is a nonzero abelian variety over \(k\), then \(G\) admits
no reductive center in the sense of \SGARefLink{sga3:xii:4.1}{4.1} with the restriction
``\(G\) affine'' omitted. Indeed, if \(n\) is prime to the
characteristic, then \({}_nG\) is étale over \(k\), of order prime
to the characteristic, and hence of multiplicative type. The family
of the \({}_nG\) is schematically dense in \(G\). A reductive center,
if it existed, would therefore equal \(G\), which is impossible
because \(G\) is not of multiplicative type. This explains the
requirement in \SGARefLink{sga3:xii:4.1}{4.1} that \(G\) have affine fibers.

\paragraph{\textbf{Lemma 4.5.}}
\SGARefTarget{sga3:xii:4.5}
\emph{Let \(S\) be a prescheme, let \(G\) be a smooth affine
\(S\)-group prescheme with connected fibers, let \(T\) be a maximal
torus of \(G\), and let \(u:H\to G\) be a central homomorphism, where
\(H\) is an \(S\)-group prescheme of multiplicative type and of
finite type. Then \(u\) factors through \(T\).}

Let \(C=\operatorname{Centr}_G(T)\). By \SGARefLink{sga3:xi:corollary:5:3}{XI~5.3} it is a closed smooth
group subprescheme of \(G\), hence affine over \(S\). Since \(T\) is
central in \(C\), it is invariant, and the quotient group
\[
  U=C/T
\]
is represented by \SGARefLink{sga3:viii:theorem:5:1}{VIII~5.1}. Because \(u\) is central it factors
through \(C\), and it remains only to show that the composite
\[
  H\longrightarrow C\longrightarrow U=C/T
\]
is trivial. By \SGARefLink{sga3:ix:corollary:5:2}{IX~5.2} we reduce to the case where \(S\) is the
spectrum of a field, which we may suppose algebraically closed.
Then BIBLE~6, Th.~2 shows that \(U\) is a connected algebraic group,
smooth over \(k\), which is ``unipotent''; as recalled above, it has
a composition series with quotients isomorphic to \(\mathbb G_a\).
Every homomorphism from the finite-type group \(H\) of multiplicative
type to \(U\) is therefore trivial. This proves \SGARefLink{sga3:xii:4.5}{4.5}.

\paragraph{\textbf{Corollary 4.6.}}
\SGARefTarget{sga3:xii:4.6}
\emph{Let \(G\) be as in \SGARefLink{sga3:xii:4.5}{4.5}. If \(G\) admits a reductive center,
then that center is contained in every maximal torus of \(G\).}

\paragraph{\textbf{Theorem 4.7.}}
\SGARefTarget{sga3:xii:4.7}
\emph{Let \(S\) be a prescheme and let \(G\) be an \(S\)-group
prescheme, smooth and affine over \(S\), with connected fibers.}

\begin{enumerate}
\item[\textup{\SGAInlineRefTarget{sga3:xii:4.7:item:a}(a)}] \emph{For every \(s\in S\), let \(z(s)\) denote
the type of the reductive center of \(G_s\), which is defined by
\SGARefLink{sga3:xii:4.4}{4.4}. Order the set \(E\) of isomorphism classes of finite-type
\(\mathbb Z\)-modules by declaring the class of \(M\) to be greater
than that of \(M'\) when \(M'\) is isomorphic to a quotient of \(M\).
Then the map
\[
  S\longrightarrow E,\qquad s\longmapsto z(s),
\]
is lower semicontinuous.}

\item[\textup{\SGAInlineRefTarget{sga3:xii:4.7:item:b}\SGAInlineRefTarget{sga3:xii:theorem:4-7:item:b}(b)}] \emph{The group \(G\) admits a reductive center
\(Z\) if and only if the preceding function \(z:S\to E\) is locally
constant. If this holds, then \(G/Z\) is representable
\textup{(cf.\ \SGARefLink{sga3:viii:theorem:5:1}{VIII~5.1})}, and its reductive center is the unit
subgroup.}

\item[\textup{\SGAInlineRefTarget{sga3:xii:4.7:item:c}\SGAInlineRefTarget{sga3:xii:theorem:4-7:item:c}(c)}] \emph{Suppose that the reductive rank of \(G\)
is locally constant \textup{(cf.\ \SGARefLink{sga3:xii:1.7:item:b}{1.7(b)})}. Then \(G\) admits a
reductive center \(Z\). If \(G/Z\) is representable---for example,
if \(G\) is affine over \(S\)---then the maximal tori \(T\) of \(G\)
\textup{(respectively, the Cartan subgroups \(C\) of \(G\))} and the
maximal tori \(T'\) of \(G'=G/Z\) \textup{(respectively, the Cartan
subgroups \(C'\) of \(G'\))} are in bijective correspondence:
\[
  T'=T/Z,\qquad C'=C/Z,
\]
and conversely
\[
  T=\varphi^{-1}(T'),\qquad C=\varphi^{-1}(C'),
\]
where \(\varphi:G\to G'\) is the canonical homomorphism.}

\item[\textup{\SGAInlineRefTarget{sga3:xii:4.7:item:d}(d)}] \emph{Let \(T\) be a maximal torus of \(G\), let
\(\mathfrak g\) denote the Lie algebra of \(G\), and consider the
homomorphism
\[
  \theta:T\longrightarrow\operatorname{GL}(\mathfrak g)
\]
induced by the adjoint representation of \(G\) \textup{(II~4)}.
Then \(\ker\theta\) is a reductive center of \(G\).}
\end{enumerate}

\emph{Proof.}
\textup{\SGARefLink{sga3:xii:4.7:item:a}{(a)}} and \textup{\SGARefLink{sga3:xii:4.7:item:b}{(b)}}. The proof of \textup{\SGARefLink{sga3:xii:4.7:item:a}{(a)}} and of the
first assertion of \textup{\SGARefLink{sga3:xii:4.7:item:b}{(b)}} is essentially identical to that of
\SGARefLink{sga3:xii:1.7:item:a}{1.7(a)} and \SGARefLink{sga3:xii:4.7:item:b}{(b)}, and we shall not repeat the argument. We point out
only that in \textup{\SGARefLink{sga3:xii:4.7:item:a}{(a)}} one must appeal to \SGARefLink{sga3:ix:theorem:5:6}{IX~5.6}, using the fact
that \(G\) has connected fibers.

We prove the second assertion of \textup{\SGARefLink{sga3:ix:theorem:5:6:item:b}{(b)}}. Suppose that \(Z\) is
a reductive center of \(G\), and put \(G'=G/Z\). By \SGARefLink{sga3:viii:theorem:5:1}{VIII~5.1} the
quotient is representable. It is affine over \(S\), of finite
presentation over \(S\) by \SGARefLink{sga3:viii:corollary:5:8}{VIII~5.8}, and even smooth over \(S\):
this may be checked on the fibers because \(G'\) is flat and of
finite presentation over \(S\), and VI\textsubscript{B}~\SGARefLink{concl:sga3-VIB-9.2-xii}{9.2(xii)}
states that a quotient of a smooth algebraic group over a field is
smooth. Since \(G\) has connected fibers, so does \(G'\). Thus \(G'\)
satisfies the same hypotheses as \(G\).

By \SGARefLink{sga3:xii:4.3}{4.3} we may reduce to the case in which \(S\) is the spectrum of a
field. Let \(Z'\) be a central subgroup of multiplicative type in
\(G'\); it remains to prove that \(Z'\) is the unit subgroup. Let
\(Z_1\) be the inverse image of \(Z'\) in \(G\), and consider the
action of \(Z_1\) on \(G\) by inner automorphisms. Since \(Z\) is
central, it acts trivially, and hence \(Z_1\) acts through an action
of \(Z'\) on \(G\). Moreover \(Z'\) acts trivially both on \(Z\) and
on \(G/Z=G'\). Since \(Z\) is central, the action of \(Z'\) on \(G\)
therefore comes from a group homomorphism
\[
  u:Z'\longrightarrow
    \operatorname{Hom}_{S\text{-}\mathrm{gr}}(G',Z),
\]
characterized by
\[
  r(z')\cdot g
    =g\cdot u(z')\bigl(\varphi(g)\bigr),
\]
where \(r:Z'\to\operatorname{Aut}_{S\text{-}\mathrm{gr}}(G)\) is
the representation just considered and \(\varphi:G\to G'\) is the
canonical homomorphism. Giving \(u\) is equivalent to giving a group
homomorphism
\[
  v:G'\longrightarrow
    \operatorname{Hom}_{S\text{-}\mathrm{gr}}(Z',Z).
\]
By \SGARefLink{sga3:x:corollary:5-8}{X~5.8} the group on the right is representable and étale over
\(S\). Its unit section is therefore an open immersion, so
\(\ker v\) is an open subgroup of \(G'\). Since \(G'\) has connected
fibers, \(\ker v=G'\). Thus \(v=0\), hence \(u=0\), and \(Z'\), and
therefore \(Z_1\), acts trivially on \(G\). Consequently \(Z_1\) is
central in \(G\). It is a commutative extension of the group \(Z'\)
of multiplicative type by the group \(Z\) of multiplicative type,
and hence, over a field, is itself of multiplicative type by \SGARefLink{sga3:xvii:proposition:7:1:1}{7.1.1}.
Being central in \(G\), it is contained in the reductive center
\(Z\). Thus \(Z_1=Z\), whence \(Z'\) is the unit subgroup. This
proves \textup{\SGARefLink{sga3:xii:4.7:item:b}{(b)}}.

\textup{\SGARefLink{sga3:xii:4.7:item:d}{(d)}} Put \(Z=\ker\theta\), a subgroup of multiplicative type
in \(T\), for example by \SGARefLink{sga3:ix:theorem:6:8}{IX~6.8}. By \SGARefLink{sga3:xii:4.5}{4.5}, every central homomorphism
\(u:H\to G\), with \(H\) of multiplicative type and of finite type,
factors through \(T\), and therefore through \(Z\). Since formation
of \(Z\) commutes with every base change, it remains to show that
\(Z\) is central, or equivalently that its centralizer \(C\) is
equal to \(G\). By \SGARefLink{sga3:xi:corollary:5:3}{XI~5.3}, \(C\) is a closed smooth subgroup of
\(G\). Since \(Z\) acts trivially on \(\mathfrak g\), we have
\[
  \operatorname{Lie}(C)=\mathfrak g.
\]
It follows that \(C=G\). Indeed, the immersion \(C\to G\) is étale,
as may be checked fiberwise: it is an unramified homomorphism of
smooth algebraic groups of the same dimension
\textup{(VI\textsubscript{B}~\SGARefLink{sga3:xii:1.3}{1.3})}, so \SGARefLink{sga3:x:lemma:3:5}{X~3.5} applies. Hence
\(C\to G\) is an étale closed immersion, and therefore an open
immersion by SGA~1, \SGARefLink{sga3:i:subsection:5-1}{I~5.1}. Since it is both open and closed and \(G\)
has connected fibers, it is an isomorphism.

\textup{\SGARefLink{sga3:xii:4.7:item:c}{(c)}} By 1.7(b), \(G\) admits a maximal torus locally for the
fpqc topology. \SGARefLink{sga3:xii:4.2}{Proposition~4.2} and part \textup{\SGARefLink{sga3:xii:corollary:6-7:equation:d}{(d)}} therefore show
that \(G\) admits a reductive center \(Z\). Part \textup{\SGARefLink{sga3:xii:corollary:6-7:equation:d}{(d)}} also
shows that every maximal torus \(T\) of \(G\) contains \(Z\).
Clearly \(T'=T/Z\) is a torus. To prove that it is a maximal torus
of \(G'\), \SGARefLink{sga3:xii:1.3}{Definition~1.3} reduces us to the case in which \(S\) is
the spectrum of an algebraically closed field; the assertion then
follows from \emph{BIBLE}~7, Th.~3\SGARefLink{sga3:xii:4.7:item:a}{(a)}.

Conversely, let \(T'\) be a maximal torus of \(G'\), and set
\(T=\varphi^{-1}(T')\). We prove that \(T\) is a maximal torus of
\(G\). The question is local for the fpqc topology, so we may
suppose that \(G\) already admits a maximal torus \(T_0\). By what
precedes, \(T'_0=T_0/Z\) is a maximal torus of \(G'\). By the
conjugacy theorem 1.7(c), \(T'\) and \(T'_0\) are locally conjugate
for the fpqc topology. Since \(\varphi:G\to G'\) is a covering for
this topology, sections of \(G'\) lift locally to sections of \(G\);
thus \(T\) and \(T_0\) are likewise locally conjugate. Hence \(T\)
is a maximal torus. The argument for Cartan subgroups is identical.
This completes the proof.

Here is a useful reformulation of \textup{\SGARefLink{sga3:xii:corollary:6-7:equation:d}{(d)}} when \(T\) is
diagonalizable, say \(T=D_S(M)\), where \(M\) is a finite-type free
\(\mathbb Z\)-module. By \SGARefLink{sga3:i:proposition:4-7-3}{I~4.7.3}, the \(T\)-module \(\mathfrak g\)
decomposes into the direct sum of its \(T\)-submodules
\(\mathfrak g_m\), \(m\in M\):
\[
  \mathfrak g=\bigoplus_{m\in M}\mathfrak g_m.
\]
These submodules are necessarily locally free. Suppose that for
every \(m\in M\) the rank of \(\mathfrak g_m\) is constant, as is
the case in particular when \(S\) is connected. Then the set
\[
  R=\{m\in M\mid\mathfrak g_m\ne0\}
\]
of ``roots'' is finite. We obtain:

\paragraph{\textbf{Corollary 4.8.}}
\SGARefTarget{sga3:xii:4.8}
\emph{Under the preceding hypotheses and with the preceding
notation, the reductive center of \(G\) is the intersection of the
kernels in \(T\) of the root characters \(m\in R\). Thus there is
an isomorphism
\[
  Z\simeq D_S(N),
\]
where \(N\) is the quotient of \(M\) by the subgroup generated by
\(R\).}

\paragraph{\textbf{Corollary 4.9.}}
\SGARefTarget{sga3:xii:4.9}
\emph{Let \(G\) be an algebraic group over an algebraically closed
field \(k\). Suppose that \(G\) is smooth, connected, and affine
over \(k\), that its reductive center is the unit subgroup, and that
its Lie algebra \(\mathfrak g\) is nilpotent. Then \(G\) is
``unipotent,'' that is, it admits a composition series whose
quotients are isomorphic to \(\mathbb G_a\).}

By \emph{BIBLE}~6, Th.~4, Cor.~3, it is enough to prove that a
maximal torus \(T\) of \(G\) is the unit subgroup, or equivalently
that its Lie algebra \(\mathfrak t\) is zero. The \(T\)-module
\(\mathfrak g\) decomposes according to the characters of \(T\) by
\SGARefLink{sga3:i:proposition:4-7-3}{I~4.7.3}, so for every \(t\in\mathfrak t\), the endomorphism
\(\operatorname{ad}_{\mathfrak g}(t)\) of \(\mathfrak g\) is
semisimple. If \(\mathfrak g\) is nilpotent, this endomorphism is
also nilpotent, hence zero. By \SGARefLink{sga3:xii:4.7:item:d}{4.7(d)}, the hypothesis on the
reductive center makes the homomorphism
\[
  T\longrightarrow\operatorname{GL}(\mathfrak g)
\]
a monomorphism, and it therefore induces an injective map on Lie
algebras. By \SGARefLink{sga3:ii:proposition:4-5}{II~4.5}, this says that
\(\operatorname{ad}_{\mathfrak g}(t)=0\) implies \(t=0\).
Consequently \(\mathfrak t=0\), as required.

\paragraph{\textbf{Proposition 4.10.}}
\SGARefTarget{sga3:xii:4.10}
\emph{Let \(G\) be a smooth connected affine algebraic group over
an algebraically closed field \(k\). Then the reductive center of
\(G\) is the intersection of the maximal tori of \(G\).}

The intersection is meant schematically, or equivalently in the
sense of subfunctors of \(G\): it is the largest closed
subprescheme of \(G\) contained in every maximal torus. Since \(G\)
is noetherian, it is already the intersection of a suitable finite
family of maximal tori.

Let \(Z\) denote this intersection. It is a closed subgroup of a
torus, hence is of multiplicative type, and by \SGARefLink{sga3:xii:4.5}{4.5} it contains the
reductive center of \(G\). To prove equality it remains to show that
\(Z\) is central. Since \(G\) is connected, \SGARefLink{sga3:ix:corollary:5:5}{IX~5.5} reduces this to
showing that \(Z\) is invariant. By construction \(Z\) is invariant
under every inner automorphism \(\operatorname{int}(g)\), with
\(g\in G(k)\). Its normalizer \(N\), represented by \SGARefLink{sga3:viii:corollary:6:7}{VIII~6.7}, is
therefore a closed subgroup of \(G\) containing all rational points
of \(G\). Since \(G\) is reduced, \(N=G\), which proves the result.

\paragraph{\textbf{Proposition 4.11.}}
\SGARefTarget{sga3:xii:4.11}
\emph{Let \(S\) be a prescheme and let \(G\) be a smooth affine
\(S\)-group prescheme with connected fibers and unipotent rank zero.
Then \(G\) has locally constant reductive rank by \SGARefLink{sga3:xii:1.9}{1.9}, and hence the
``prescheme of maximal tori of \(G\),'' denoted by
\(\mathcal T\), is defined by \SGARefLink{sga3:xii:1.10}{1.10} and is smooth, separated, and of
finite type over \(S\). Let \(G\) act on \(\mathcal T\) through
inner automorphisms, giving a homomorphism of group functors
\[
  u:G\longrightarrow\operatorname{Aut}_S(\mathcal T).
\]
Under these hypotheses the following three subfunctors of \(G\) are
identical:
\begin{enumerate}
\item[\textup{(i)}] the reductive center \(Z\) of \(G\);
\item[\textup{(ii)}] the center \(Z'\) of \(G\);
\item[\textup{(iii)}] the kernel \(Z''=\ker u\) of the preceding
homomorphism.
\end{enumerate}
In particular, the center of \(G\) is represented by a subgroup of
multiplicative type in \(G\).}

\emph{Proof.}
Clearly \(Z\subset Z'\subset Z''\). It remains to prove
\(Z''\subset Z\), or, since the hypotheses are stable under base
change, that every section \(g\) of \(G\) over \(S\) which acts
trivially on \(\mathcal T\) is a section of \(Z\). Put \(G'=G/Z\).
By \SGARefLink{sga3:xii:4.7:item:b}{4.7(b)} and \SGARefLink{sga3:xii:4.7:item:c}{(c)}, the prescheme \(\mathcal T'\) of maximal tori of
\(G'\) is canonically isomorphic to \(\mathcal T\). We are therefore
reduced to the case \(G=G'\), in which the reductive center \(Z\) is
the unit subgroup. Notice that \SGARefLink{sga3:xii:4.7:item:c}{4.7(c)} also shows that the unipotent
rank of \(G'\) equals that of \(G\), and hence is zero.

We must prove that \(g\) is the unit section of \(G\). The usual
reduction process reduces us first to the case in which \(S\) is
noetherian, and then to the case in which \(S\) is local artinian:
when \(S\) is locally noetherian, it is enough to check the assertion
after every base change \(S'\to S\) with \(S'\) local artinian. In
this case \(Z''=\ker u\) is representable by \SGARefLink{sga3:viii:proposition:6:2:item:a:01}{VIII~6.2(a)} and \SGARefLink{sga3:viii:examples:6:5:item:c:01}{6.5(c)};
alternatively, representability follows from \SGARefLink{sga3:xi:theorem:6:8}{XI~6.8}. By Nakayama's
lemma it is enough to show that the fiber \(Z''_0\) is the unit
subgroup. We may therefore suppose that \(S\) is the spectrum of a
field \(k\), which we may take to be algebraically closed.

The subgroup \(Z''\) lies in the stabilizer of every point of
\(\mathcal T(k)\), that is, in the normalizer \(N(T)\) of every
maximal torus \(T\) of \(G\). Since the unipotent rank of \(G\) is
zero, \SGARefLink{sga3:xi:proposition:5:9}{XI~5.9} says that \(T\) is an open subgroup of \(N(T)\). Thus
\[
  \operatorname{Lie}N(T)=\operatorname{Lie}T,
\]
and \(\operatorname{Lie}Z''\) is contained in
\(\operatorname{Lie}T\) for every maximal torus \(T\). \SGARefLink{sga3:xii:4.10}{Proposition
4.10} shows that the intersection of the Lie algebras of all maximal
tori is the Lie algebra of the reductive center \(Z\), which is zero
in the present reduction. Hence \(\operatorname{Lie}Z''=0\), so
\(Z''\) is étale over \(k\).

Moreover, \(Z''\) is invariant in \(G\). Since \(G\) is connected,
it follows that \(Z''\) is central in \(G\). Thus it is contained in
\[
  T=\operatorname{Centr}_G(T)
\]
for every maximal torus \(T\), and therefore in their intersection.
By \SGARefLink{sga3:xii:4.10}{4.10} this intersection is \(Z=0\), completing the proof.

\section[Application to the prescheme of subgroups of multiplicative type]
{Application to the prescheme of subgroups of multiplicative
type\texorpdfstring{\protect\footnotemark}{}}
\footnotetext{For a generalization of the results of this section to
the case in which \(G\) is not assumed affine over \(S\), see
M.~Raynaud, ``Ample sheaves on group schemes and homogeneous
spaces,'' Chapter~IX.}
\SGARefTarget{sga3:xii:5}

\paragraph{\textbf{Theorem 5.1.}}
\SGARefTarget{sga3:xii:5.1}
\emph{Let \(S\) be a prescheme, let \(G\) be a smooth affine
\(S\)-group prescheme, and let \(M\) be the ``prescheme of subgroups
of multiplicative type in \(G\),'' representing the functor
described in \SGARefLink{sga3:xi:theorem:4:1}{XI~4.1}. For every integer \(n>0\), let \(T_n\) be the
subfunctor of \(M\) for which \(T_n(S')\) is the set of subgroup
preschemes \(H\) of multiplicative type in \(G_{S'}\) satisfying
\[
  n\operatorname{id}_H=0.
\]
Thus \(T_n\) is representable and affine over \(S\) by \SGARefLink{sga3:xi:proposition:3:12:item:a}{XI~3.12(a)}.
Let
\[
  u_n:M\longrightarrow T_n
\]
be the morphism defined by \(u_n(H)={}_nH\), where
\({}_nH=\ker(n\operatorname{id}_H)\). Then:}

\begin{enumerate}
\item[\textup{\SGAInlineRefTarget{sga3:xii:theorem:5-1:item:a}\SGAInlineRefTarget{sga3:xii:5.1:item:a}(a)}] \emph{Every subprescheme \(U\) of \(M\) of
finite type over \(S\) is contained in a closed subprescheme of
finite type over \(S\). Every closed subprescheme \(X\) of \(M\) of
finite type over \(S\) is affine over \(S\).}

\item[\textup{\SGAInlineRefTarget{sga3:xii:5.1:item:b}(b)}] \emph{Suppose that \(S\) is quasi-compact, and
let \(X\) be a closed subprescheme of \(M\) of finite type over
\(S\). Then there is an integer \(n>0\) such that, for every
multiple \(m\) of \(n\), the induced morphism
\[
  u_m|_X:X\longrightarrow T_m
\]
is a closed immersion.}
\end{enumerate}

\emph{Proof.}
\textup{(a)} For the first assertion, take for \(X\) the schematic
closure of \(U\) \textup{(EGA~I, 9.5.1 and 9.5.3)}. This closure is
defined because the immersion \(i:U\to M\) is quasi-compact:
\(U\) is of finite type over \(S\), and \(M\) is separated over
\(S\) by \SGARefLink{sga3:xi:theorem:4:1}{XI~4.1}. It remains to prove that such an \(X\) is affine
over \(S\), which will prove the second assertion at the same time.
The question is local on \(S\), so we may suppose that \(S\) is
affine. Since \(U\) is of finite type over \(S\), it is contained
in a quasi-compact open of \(M\). We are therefore reduced to the
case in which \(U\) itself is an open subprescheme of finite type
over \(S\), hence quasi-compact. Notice that a closed subscheme of a
scheme affine over \(S\) is affine over \(S\).

It is enough to show that such a \(U\) is contained in a closed
affine subprescheme of \(M\). In this form the usual reduction
process reduces us immediately to the case in which \(S\) is
noetherian. The same argument applies to \textup{\SGARefLink{sga3:xii:5.1:item:b}{(b)}}, reducing it
to the case in which \(S\) is affine noetherian.

Recall the inverse limit \(T\) of the \(T_n\) used in the proof of
\SGARefLink{sga3:xi:theorem:4:1}{XI~4.1}. It is an affine prescheme over \(S\), though not of finite
type, and its local rings are noetherian, as was observed at the
beginning of the proof of IX~3.7. More explicitly, if
\(t=(t_n)\in T\), then the transition morphisms \(T_m\to T_n\) are
smooth and the dimensions of the \(T_m\) are bounded. Hence there
is an \(n_0\) such that \(T_n\to T_{n_0}\) is étale at \(t_n\) for
every multiple \(n\) of \(n_0\). The proof just cited, or \SGARefLink{sga3:xi:corollary:3:11}{XI~3.11},
shows that the canonical morphism
\[
  u:M\longrightarrow T
\]
is an immersion and induces isomorphisms on local rings. It need not
be an open immersion or a quasi-compact morphism.

Let \(U\) be a quasi-compact open of \(M\). We shall show that its
closure in \(T\) is contained in \(M\). If \(X\) is the schematic
closure of \(U\) in \(M\), this will make
\[
  u|_X:X\longrightarrow T
\]
a closed immersion. Then \(X\) is affine, proving \textup{\SGARefLink{sga3:xii:5.1:item:a}{(a)}}.
Since \(U\) is noetherian it has only finitely many irreducible
components. Every \(t\in T\) in the closure of \(U\) is a
specialization of a point \(x\in U\). The local ring of \(t\) in
\(T\) is noetherian, and so is the local ring \(A\) of \(t\) in
\(\overline{\{x\}}\), endowed with the induced reduced structure.

The canonical morphism from the noetherian integral local scheme
\[
  S'=\operatorname{Spec}(A)
\]
to \(T\) sends its closed point to \(t\) and its generic point to
\(x\in M\). We must show that \(t\in M\). Replacing \(A\), if
necessary, by the quotient by a minimal prime of a suitable complete
local \(A\)-algebra flat over \(A\) \textup{(EGA
0\textsubscript{III}, 10.3.1)}, we may suppose that \(A\) is
complete and has algebraically closed residue field. One could even
reduce further to a discrete valuation ring by EGA~II, 7.1.7.
Renaming \(S'\) as \(S\), we are reduced to the following lemma.

\paragraph{\textbf{Lemma 5.2.}}
\SGARefTarget{sga3:xii:5.2}
\emph{Let \(S\) be the spectrum of a complete noetherian integral
local ring with algebraically closed residue field. Let \(G\) be a
smooth affine \(S\)-group prescheme, let
\(\bigl(H(n)\bigr)_{n>0}\) be a system of subgroups of
multiplicative type in \(G\), indexed by the positive integers, and
let \(\eta\) be the generic point of \(S\). Suppose that:}

\begin{enumerate}
\item[\textup{(a)}] \emph{if \(m\) is a multiple of \(n\), then
\[
  H(n)={}_nH(m);
\]}
\item[\textup{\SGAInlineRefTarget{sga3:xii:5.2:item:b}(b)}] \emph{there is a subgroup \(H_\eta\) of
multiplicative type in \(G_\eta\) such that
\[
  H(n)_\eta={}_n(H_\eta)
\]
for every \(n>0\).}
\end{enumerate}

\emph{Then there is a subgroup \(H\) of multiplicative type in
\(G\) such that \(H(n)={}_nH\) for every \(n>0\).}

\emph{Proof.}
For every \(n\), let
\[
  C_n=\operatorname{Centr}_G\bigl(H(n)\bigr).
\]
This is a closed subgroup prescheme of \(G\), smooth over \(S\), by
\SGARefLink{sga3:xi:corollary:5:3}{XI~5.3}. The positive integers being ordered by divisibility, the
\(C_n\) form a decreasing family of closed subpreschemes of \(G\).
Since \(G\) is noetherian, this family is stationary. Let \(C\) be
the common value of the \(C_n\) for sufficiently divisible \(n\).
Then \(C\) satisfies the same conditions as \(G\), and all the
\(H(n)\) are in fact subgroups of \(C\). Replacing \(G\) by \(C\),
we may therefore suppose that all the \(H(n)\) are central subgroups
of \(G\).

Let \(s\) be the closed point of \(S\), and let \(Z_s\) be the
reductive center of \(G_s\), defined in \SGARefLink{sga3:xii:4.4}{4.4}. By the definition in
\SGARefLink{sga3:xii:4.1}{4.1}, \(Z_s\) contains every \(H(n)_s\). Since \(A\) is complete,
\(Z_s\) lifts to a subgroup prescheme \(Z\) of multiplicative type
in \(G\), by \SGARefLink{sga3:xi:theorem:5:8}{XI~5.8}. We claim that \(Z\) contains every \(H(n)\).
Indeed, \(H(n)\) is central and hence commutes with \(Z\) in \(G\).
The resulting homomorphism
\[
  Z\times H(n)\longrightarrow G
\]
has, by \SGARefLink{sga3:ix:theorem:6:8}{IX~6.8}, an image \(K\) which is a subgroup of multiplicative
type in \(G\) containing \(Z\). It remains only to prove \(K=Z\);
this follows from \(K_s=Z_s\) and \SGARefLink{sga3:ix:theorem:5:1}{IX~5.1}~bis.

We are thus reduced to proving the analogue of the lemma with \(G\)
replaced by a group \(Z\) of multiplicative type and of finite type
over \(S\), not necessarily smooth over \(S\). Since \(A\) is
complete with algebraically closed residue field, \(Z\) is
diagonalizable by \SGARefLink{sga3:x:corollary:3:3}{X~3.3} and \SGARefLink{sga3:xii:1.4}{1.4}, hence is of the form
\(D_S(M)\), where \(M\) is a finitely generated commutative group.
Every subgroup of multiplicative type \(H\) of \(Z\) is therefore
diagonalizable by \SGARefLink{sga3:ix:proposition:2:11:item:i:01}{IX~2.11(i)}, and so is of the form \(D_S(N)\),
where \(N\) is a quotient group of \(M\), by \SGARefLink{sga3:viii:corollary:3:2}{VIII~3.2}. Thus the
\(H(n)\) correspond to quotient groups \(M(n)\) of \(M\), and the
existence of a subgroup \(H\) such that \(H(n)={}_nH\) for every
\(n\) is equivalent to the existence of a quotient group \(N\) of
\(M\) such that
\[
  M(n)=N/nN
\]
for every \(n\). This follows at once from the existence of the
subgroup \(H_\eta\) of \(G_\eta\) such that
\(H(n)_\eta={}_n(H_\eta)\) for every \(n\). This proves \SGARefLink{sga3:xii:5.2}{Lemma~5.2}
and hence \SGARefLink{sga3:xii:theorem:5-1:item:a}{Theorem~5.1(a)}.

\textup{\SGARefLink{sga3:xii:5.2:item:b}{(b)}} We already know that
\[
  u|_X:X\longrightarrow T
\]
is a closed immersion. It follows readily that, for all sufficiently
large \(m\), the composite
\[
  u_m|_X:X\longrightarrow T_m
\]
is also a closed immersion. Since we shall not need this fact below,
we omit the details of the proof.

\paragraph{\textbf{Corollary 5.3.}}
\SGARefTarget{sga3:xii:5.3}
\emph{With the notation of \SGARefLink{sga3:xii:5.1}{5.1}, let \(U\) be a subprescheme of
\(M\), both open and closed, and of finite type over \(S\). Then
\(U\), with the structure induced by \(M\), is affine over \(S\).
If \(S\) is quasi-compact, there is an integer \(n>0\) such that,
for every multiple \(m\) of \(n\), the induced morphism
\[
  u_m|_U:U\longrightarrow T_m
\]
is both an open and a closed immersion; equivalently, it is an
isomorphism onto an open-and-closed subprescheme of \(T_m\), endowed
with the induced structure.}

\emph{Proof.}
The first assertion follows from \SGARefLink{sga3:xii:5.1:item:a}{5.1(a)}, and the second from \SGARefLink{sga3:xii:5.1:item:b}{5.1(b)},
since \(u_m:M\to T_m\) is smooth by \SGARefLink{sga3:xi:corollary:2:2}{XI~2.2}~bis and a smooth
immersion, that is, an étale immersion, is an open immersion.

\paragraph{\textbf{Corollary 5.4.}}
\SGARefTarget{sga3:xii:5.4}
\emph{Let \(S\) be a prescheme, let \(G\) be a smooth affine
\(S\)-group prescheme whose reductive rank is locally constant
\textup{(\SGARefLink{sga3:xii:1.7:item:b}{1.7(b)})}, and let \(\mathcal T\) be the ``prescheme of
maximal tori of \(G\)'' \textup{(\SGARefLink{sga3:xii:1.10}{1.10})}. Then \(\mathcal T\) is
smooth and affine over \(S\). If \(T\) is a maximal torus of \(G\)
and \(N(T)\) its normalizer, then \(G/N(T)\) \textup{(\SGARefLink{sga3:xi:corollary:5:3}{XI~5.3}~bis)}
is affine over \(S\). The same is true of \(G/C(T)\), where \(C(T)\)
is the centralizer of \(T\), provided that
\[
  W(T)=N(T)/C(T)
\]
is finite over \(S\); cf.\ \SGARefLink{sga3:xii:2.1:item:b}{2.1(b)}.}

\emph{Proof.}
The second assertion is contained in the first, since the conjugacy
theorem identifies \(G/N(T)\) with \(\mathcal T\), by \SGARefLink{sga3:xi:remark:5:5}{XI~5.5}~bis.
For the first assertion, observe that by construction
\(\mathcal T\) is isomorphic to an open-and-closed subprescheme of
\(M\). Indeed, one may suppose that the reductive rank of \(G\) is
constant and equal to \(r\); then \(\mathcal T\) is the subprescheme
of \(M\) corresponding to the subtori of relative dimension \(r\),
that is, the largest subprescheme of \(M\) over which the universal
subgroup of multiplicative type
\[
  H\subset G_M
\]
is a torus of relative dimension \(r\). \SGARefLink{sga3:xii:5.3}{Corollary~5.3} therefore
applies. Finally, \(G/C(T)\) is finite over
\[
  G/N(T)\simeq \bigl(G/C(T)\bigr)/W(T),
\]
and hence is affine because \(G/N(T)\) is affine.

\paragraph{\textbf{Corollary 5.5.}}
\SGARefTarget{sga3:xii:5.5}
\emph{Let \(G\) be a smooth affine algebraic group over a field
\(k\). Then the scheme \(M\) of subgroups of multiplicative type in
\(G\) is a disjoint union of affine \(k\)-schemes. For every
subgroup \(H\) of multiplicative type in \(G\), if \(C\) and \(N\)
denote respectively its centralizer and its normalizer, then the
quotients \(G/C\) and \(G/N\) are affine.}

\emph{Proof.}
By \SGARefLink{sga3:xi:corollary:5:1}{XI~5.1}~bis, the saturation under the action of \(G\) on \(M\) of
every finite closed subset of \(M\) is open. Indeed, one reduces to
the case in which \(k\) is algebraically closed, and hence to the
orbit of a \(k\)-rational point \(x\); then, by the cited result, the
morphism
\[
  G\longrightarrow M,\qquad g\longmapsto g\cdot x
\]
is smooth and therefore open, so its image is open.

Let \(U\) be the union of the equivalence classes of the closed
points of \(M\) under the action of \(G\). Then \(U\) is open and
contains every closed point of \(M\); by the Nullstellensatz,
\(U=M\). Thus \(M\) is a disjoint union of open subpreschemes, which
are necessarily also closed. Hence \(M\) is the sum prescheme of
subpreschemes \(M_i\), each of which is the \(G\)-orbit of a closed
point. Each \(M_i\) is quasi-compact, hence of finite type, and is
therefore affine by \SGARefLink{sga3:xii:5.3}{5.3}.

If \(H\) is a subgroup of multiplicative type in \(G\), it
corresponds to a \(k\)-rational point \(x\) of \(M\), and \(G/N\)
identifies with the orbit of \(x\) under \(G\), by \SGARefLink{sga3:xi:remark:5:5}{XI~5.5}~bis. Thus
\(G/N\) is affine. Since \(C\) is an open subgroup of \(N\), by
\SGARefLink{sga3:xi:proposition:5:9}{XI~5.9}, the quotient \(G/C\) is finite over \(G/N\); indeed,
\[
  G/N\simeq (G/C)/(N/C).
\]
Consequently \(G/C\) is affine as well.

This proof also gives the following result.

\paragraph{\textbf{Corollary 5.6.}}
\SGARefTarget{sga3:xii:5.6}
\emph{Under the hypotheses of \SGARefLink{sga3:xii:5.5}{5.5} on \(G\) and \(H\), the
subprescheme \(U\) of \(M\) consisting of the subgroups of
multiplicative type in \(G\) which are locally conjugate to \(H\)
is an open-and-closed subprescheme of \(M\).}

In picturesque language, every subgroup \(H'\) of multiplicative
type in \(G\) which is a limit of subgroups conjugate to \(H\) is
itself conjugate to \(H\).

\paragraph{\textbf{Remarks 5.7.}}
\SGARefTarget{sga3:xii:5.7}
Let \(S\) be a prescheme, let \(G\) be a smooth affine \(S\)-group
prescheme, let \(H\) be an \(S\)-prescheme of multiplicative type
and of finite type, and put
\[
  M=\operatorname{Hom}_{S\text{-gr}}(H,G).
\]
By \SGARefLink{sga3:xi:corollary:4:2}{XI~4.2}, \(M\) is representable, smooth over \(S\), and separated
over \(S\). One can prove for \(M\) a result entirely analogous to
\SGARefLink{sga3:xii:5.1}{5.1}, either by reducing to \SGARefLink{sga3:xii:5.1}{5.1} through an argument analogous to that
of \SGARefLink{sga3:xi:corollary:4:2}{XI~4.2}, or directly by an argument modeled on the proof of \SGARefLink{sga3:xii:5.1}{5.1}.
The corresponding variants of \SGARefLink{sga3:xii:5.3}{5.3}, \SGARefLink{sga3:xii:5.5}{5.5}, and \SGARefLink{sga3:xii:5.6}{5.6} follow; we leave
their formulation to the reader.

\section[Maximal tori and Cartan subgroups]
{Maximal tori and Cartan subgroups of not necessarily affine
algebraic groups \textup{(algebraically closed base field)}}
\SGARefTarget{sga3:xii:6}

\paragraph{\textbf{Lemma 6.1.}}
\SGARefTarget{sga3:xii:6.1}
\emph{Let \(G\) be a connected algebraic group over a field \(k\),
and let \(Z\) be an algebraic subgroup of finite index in the center
of \(G\). Then \(G/Z\) is affine.}

\emph{Proof.}
We may suppose that \(Z\) is the center of \(G\), since a scheme
finite over an affine scheme is affine. Consider the vector spaces
\[
  P^n=P^n(G)=
  \mathcal O_{G,e}/\mathfrak m_{G,e}^{\,n+1}
  \qquad (n\geq 0),
\]
where \(\mathfrak m_{G,e}\) is the maximal ideal at the identity.
The group \(G\) acts on the \(P^n(G)\) through the adjoint
representation. Let \(Z_n\) be the kernel of the corresponding
homomorphism
\[
  \operatorname{ad}_n:G\longrightarrow \operatorname{GL}(P^n).
\]
Using the connectedness of \(G\), one readily verifies that \(Z\)
is the intersection of the \(Z_n\). Since \(G\) is noetherian, this
intersection is equal to one of the \(Z_n\). Passing to the quotient,
\(\operatorname{ad}_n\) defines a monomorphism
\[
  G/Z_n\longrightarrow \operatorname{GL}(P^n),
\]
which is a closed immersion. Hence \(G/Z_n\) is affine, and so is
\(G/Z\).

\paragraph{\textbf{Lemma 6.2.}}
\SGARefTarget{sga3:xii:6.2}
\emph{Let \(G\) be a smooth algebraic group over a field \(k\), let
\(Z\) be a central algebraic subgroup, put \(G'=G/Z\), and let
\[
  u:G\longrightarrow G'
\]
be the canonical homomorphism. Let \(T\) be a subgroup of
multiplicative type in \(G\), let \(T'=u(T)\), let \(C\) be the
centralizer of \(T\) in \(G\), and let \(C'\) be the centralizer of
\(T'\) in \(G'\). Then
\[
  C'^{\,0}\subset u(C).
\]}

\emph{Proof.}
We may suppose that \(k\) is algebraically closed. Put
\[
  C_1=\bigl(u^{-1}(C'^{\,0})\bigr)_{\mathrm{red}}.
\]
It is enough to prove that \(C_1\subset C\). Indeed, \(u(C_1)\) is
connected and of finite index in \(C'^{\,0}\), hence is equal to it
set-theoretically and therefore scheme-theoretically, since \(C'\),
and consequently \(C'^{\,0}\), is smooth.

The commutator morphism
\[
  (g,t)\longmapsto gtg^{-1}t^{-1}
\]
from \(G\times G\) to \(G\) induces a morphism
\[
  \varphi:C_1\times T\longrightarrow Z_1,
  \qquad Z_1=Z_{\mathrm{red}}.
\]
Indeed, the source is reduced, and for \(g\in C_1(k)\) and
\(t\in T(k)\), the element \(gtg^{-1}t^{-1}\) belongs to \(Z(k)\)
because \(g\) centralizes \(T\) modulo \(Z\).

A calculation on \(k\)-rational points shows that
\(\varphi(g,t)\) is additive in \(g\) and in \(t\), hence
``bilinear.'' We claim that this homomorphism is identically zero;
this proves \(C_1\subset C\). By the density theorem for \(T\), we
reduce to the case in which \(T\) is finite, so that
\(n\operatorname{id}_T=0\) for some integer \(n>0\). The morphism
\(\varphi\) is then defined by a group homomorphism
\[
  C_1\longrightarrow \operatorname{Hom}_{k\text{-gr}}(T,Z_1).
\]
Since \(Z_1\) is commutative and smooth, the group on the right is
represented by an étale algebraic group over \(k\). Since \(C_1\)
is connected, every group homomorphism from \(C_1\) to that group
is zero.

\paragraph{\textbf{Corollary 6.3.}}
\SGARefTarget{sga3:xii:6.3}
\emph{Under the preceding hypotheses, suppose that \(C'\) is
connected. Then
\[
  C'=u(C),
  \qquad
  C=u^{-1}(C').
\]}

\emph{Proof.}
In this case \(C'=C'^{\,0}\). Moreover, \(C\) plainly contains
\(Z\), and hence
\[
  C=u^{-1}\bigl(u(C)\bigr).
\]

\paragraph{\textbf{Lemma 6.4.}}
\SGARefTarget{sga3:xii:6.4}
\emph{Let \(G\) be an algebraic group over a field \(k\), and let
\(Z\) be a central algebraic subgroup such that
\[
  G/Z=G'
\]
is a torus. Then \(G\) is commutative. If \(k\) is algebraically
closed, there is a torus \(T\) in \(G\) such that \(u(T)=G'\), where
\[
  u:G\longrightarrow G/Z=G'
\]
is the canonical homomorphism.}

\emph{Proof.}
We may suppose that \(k\) is algebraically closed. Consider again
the commutator morphism
\[
  G\times G\longrightarrow G,\qquad
  (g,h)\longmapsto ghg^{-1}h^{-1}.
\]
Since \(Z\) is central, this morphism factors through a morphism
\(G'\times G'\to G\). Since \(G'\) is commutative, the latter takes
its values in \(Z\), and even in \(Z_{\mathrm{red}}\), because
\(G'\times G'\) is reduced. As above, the resulting morphism
\[
  \varphi:G'\times G'\longrightarrow Z_{\mathrm{red}}
\]
is bilinear and hence zero. Thus \(G\) is commutative.

To find a torus \(T\) in \(G\) with \(u(T)=G'\), we may, after
replacing \(G\) by \(G_{\mathrm{red}}^{\,0}\), suppose that \(G\)
is smooth and connected. Replacing \(Z\) by
\(Z_{\mathrm{red}}^{\,0}\), we may also suppose that \(Z\) is
smooth and connected. By a well-known theorem of
Chevalley,\footnote{Editor's note: references should be supplied
here.} \(Z\) is an extension of an abelian variety by an affine
group. We are therefore reduced to the two cases in which \(Z\) is
affine or \(Z\) is an abelian variety. In the first case \(G\) is
affine and the result is classical \textup{(BIBLE~7, Theorem~3(a))}.

Suppose then that \(Z\) is an abelian variety. Every homomorphism
from the additive group \(\mathbb G_a\) to the torus \(G'=R\), or
to the abelian variety \(Z\), is trivial. The same is therefore true
of every homomorphism from \(\mathbb G_a\) to \(G\); thus \(G\)
contains no subgroup isomorphic to \(\mathbb G_a\). Chevalley's
structure theorem gives an exact sequence
\[
  0\longrightarrow T\longrightarrow G\longrightarrow A
  \longrightarrow 0, \tag{*}\label{sga3:xii:lemma:6-4:equation:star}
\]
where \(A\) is an abelian variety and \(T\) is a smooth connected
affine group. Since \(T\) is commutative and contains no additive
subgroup, \(T\) is a torus; it is plainly the unique maximal torus
of \(G\). It remains to prove that every epimorphism
\[
  u:G\longrightarrow R
\]
to a torus \(R\) satisfies \(u(T)=R\).

Put
\[
  \operatorname{Hom}_{\mathrm{gr}}(T,\mathbb G_m)=M,
  \qquad
  \operatorname{Hom}_{\mathrm{gr}}(R,\mathbb G_m)=P.
\]
These are finitely generated free \(\mathbb Z\)-modules, and
\[
  T=D_k(M),\qquad R=D_k(P).
\]
Let
\[
  B=\operatorname{Ext}^{1}_{k\text{-gr}}(A,\mathbb G_m);
\]
this is also the set of \(k\)-rational points of the abelian variety
dual to \(A\). We have
\[
  \operatorname{Ext}^{1}(A,T)
    =\operatorname{Hom}_{\mathrm{gr}}(M,B),
  \qquad
  \operatorname{Ext}^{1}(A,R)
    =\operatorname{Hom}_{\mathrm{gr}}(P,B). \tag{**}\label{sga3:xii:lemma:6-4:equation:double-star}
\]
In particular, the extension \(G\) of \(A\) by \(T\) is determined
by a homomorphism
\[
  \theta:M\longrightarrow B.
\]

The exact sequence \textup{(*)} also gives
\[
  0\longrightarrow \operatorname{Hom}(A,R)
  \longrightarrow \operatorname{Hom}(G,R)
  \longrightarrow \operatorname{Hom}(T,R)
  \longrightarrow \operatorname{Ext}^{1}(A,R).
\]
Moreover, \(\operatorname{Hom}(A,R)=0\), and the last map identifies
with the homomorphism
\[
  \operatorname{Hom}(P,M)\longrightarrow
  \operatorname{Hom}(P,B)
\]
induced by \(\theta:M\to B\). Setting
\[
  N=\ker\theta,
\]
we obtain canonical bijections
\[
  \operatorname{Hom}(G,R)
  \simeq \operatorname{Hom}(P,N)
  \simeq \operatorname{Hom}(S,R),
  \qquad S=D_k(N).
\]
This has the following immediate geometric description.

\paragraph{\textbf{Lemma 6.5.}}
\SGARefTarget{sga3:xii:6.5}
\emph{Let \(G\) be an extension of an abelian variety \(A\) by a
torus \(T=D_k(M)\), defined by a homomorphism
\[
  \theta:M\longrightarrow
  B=\operatorname{Ext}^{1}(A,\mathbb G_m),
\]
over an algebraically closed field \(k\). Put \(N=\ker\theta\), and
let \(S=D_k(N)\) be the corresponding torus, isomorphic to \(T/U\),
where \(U=D_k(M/N)\). Consider the extension \(H=G/U\) of \(A\) by
\(S\). This extension splits,\footnote{Editor's note: the extension
is split.} so there is a unique projection from \(H\) to \(S\), and
hence a unique homomorphism
\[
  v:G\longrightarrow S
\]
extending the canonical homomorphism \(T\to S\). For every torus
\(R\) and every homomorphism \(u:G\to R\), there is a unique
homomorphism \(u':S\to R\) such that \(u=u'v\). In particular,
\[
  \operatorname{Im}(u)=\operatorname{Im}(u')=u(T).
\]}

It follows in particular that if \(u\) is an epimorphism, then its
restriction to \(T\) is also an epimorphism. This completes the
proof of \SGARefLink{sga3:xii:6.4}{6.4}.

\paragraph{\textbf{Theorem 6.6.}}
\SGARefTarget{sga3:xii:6.6}
\emph{Let \(G\) be a smooth connected algebraic group over an
algebraically closed field \(k\).}

\begin{enumerate}
\item[\textup{\SGAInlineRefTarget{sga3:xii:theorem:6-6:item:a}\SGAInlineRefTarget{sga3:xii:6.6:item:a}(a)}] \emph{The maximal tori of \(G\) are conjugate.}

\item[\textup{\SGAInlineRefTarget{sga3:xii:6.6:item:b}\SGAInlineRefTarget{sga3:xii:theorem:6-6:item:b}(b)}] \emph{If \(T\) is a torus in \(G\), then its
centralizer is smooth and connected.}

\item[\textup{\SGAInlineRefTarget{sga3:xii:6.6:item:c}\SGAInlineRefTarget{sga3:xii:theorem:6-6:item:c}(c)}] \emph{The assignment
\[
  T\longmapsto \operatorname{Centr}_G(T)
\]
gives a bijection between the maximal tori \(T\) of \(G\) and the
Cartan subgroups of \(G\) defined in No.~1. An algebraic subgroup
\(C\) of \(G\) is a Cartan subgroup if and only if it is smooth,
nilpotent, and set-theoretically equal to its connected normalizer;
in that case it is even scheme-theoretically equal to its connected
normalizer. Moreover,
\[
  C=\operatorname{Centr}_G(T),
\]
where \(T\) is the unique maximal torus of \(C\), and
\[
  \operatorname{Norm}_G(C)=\operatorname{Norm}_G(T).
\]}

\item[\textup{\SGAInlineRefTarget{sga3:xii:6.6:item:d}(d)}] \emph{Let \(v:G\to H\) be an epimorphism of
smooth connected algebraic groups. The maximal tori, respectively
the Cartan subgroups, of \(H\) are the images of the maximal tori,
respectively the Cartan subgroups, of \(G\). If \(T\) is a maximal
torus of \(G\) and \(C\) its centralizer, then \(v(C)\) is the
centralizer of \(v(T)\).}

\item[\textup{\SGAInlineRefTarget{sga3:xii:theorem:6-6:item:e}\SGAInlineRefTarget{sga3:xii:6.6:item:e}(e)}] \emph{Under the hypotheses of \textup{\SGARefLink{sga3:xii:corollary:6-7:equation:d}{(d)}},
suppose that \(\ker v\) is a central subgroup of \(G\). Then
\[
  T\longmapsto v(T),
  \qquad
  C\longmapsto v(C)
\]
give bijections between the maximal tori, respectively the Cartan
subgroups, of \(G\) and those of \(H\). The Cartan subgroups of \(G\)
contain the center of \(G\), and hence \(\ker v\); they are precisely
the groups \(v^{-1}(C')\), where \(C'\) is a Cartan subgroup of
\(H\).}
\end{enumerate}

We also record the following immediate consequence of
\SGARefLink{sga3:xii:theorem:6-6:item:c}{Theorem~6.6(c)}.

\paragraph{\textbf{Corollary 6.7.}}
\SGARefTarget{sga3:xii:6.7}
\emph{The group \(G\) is nilpotent, that is, \(G(k)\) is nilpotent,
if and only if every torus in \(G\) is central, or equivalently if
\(G\) has a unique maximal torus. In that case this torus is the
largest subtorus of \(G\).}

\begin{proof}[Proof of \SGARefLink{sga3:xii:6.6}{Theorem 6.6}]
Let \(Z\) be a central algebraic subgroup of \(G\), put
\(G'=G/Z\), and let
\[
  u:G\longrightarrow G'
\]
be the canonical homomorphism. Then \(G'\) is smooth and connected.
If \(T'\) is a subtorus of \(G'\), it follows from \SGARefLink{sga3:xii:6.4}{6.4} that
\(u^{-1}(T')\) is commutative and that \(T'\) is the image of a
subtorus of \(u^{-1}(T')\), hence of a subtorus of \(G\). Since
\(u^{-1}(T')\) is commutative, it plainly has a largest subtorus
\(T\): the sum of two subtori is a third subtorus containing both.
Thus
\[
  u(T)=T'.
\]
It follows immediately that, for every maximal torus \(T\) of
\(G\), its image \(T'=u(T)\) is a maximal torus of \(G'\), and that
\[
  T\longmapsto u(T)
\]
is a bijection between the maximal tori of \(G\) and those of
\(G'\).

Now take
\[
  Z=\operatorname{Centr}(G)_{\mathrm{red}}^{0}.
\]
Then \(G'\) is affine by \SGARefLink{sga3:xii:6.1}{6.1}. Its maximal tori are conjugate, and
therefore so are those of \(G\), which proves \textup{\SGARefLink{sga3:xii:6.6:item:a}{(a)}}.
Moreover, \(G\) is nilpotent, respectively has a unique maximal
torus, if and only if \(G'\) has the same property. Since \(G'\) is
affine, these two conditions on \(G'\) are equivalent
\((\mathrm{BIBLE}\ 6,\ \mathrm{Th.}\ 4,\ \mathrm{Cor.}\ 2)\);
hence they are also equivalent for \(G\). If \(G\) has a unique
maximal torus \(T\), then \(T\) is invariant and therefore central;
and because every torus of \(G\) is contained in a maximal torus,
every torus is central. Conversely, if every torus is central, then
so are the maximal tori, and the conjugacy theorem
\textup{\SGARefLink{sga3:xii:6.6:item:a}{(a)}} shows that there is only one. This proves \SGARefLink{sga3:xii:6.7}{6.7}.

Let \(T\) be any torus of \(G\), let \(T'=u(T)\), and put
\[
  C'=\operatorname{Centr}_{G'}(T').
\]
The group \(C'\) is connected
\((\mathrm{BIBLE}\ 6,\ \mathrm{Th.}\ 6\textup{(a)})\).
By \SGARefLink{sga3:xii:6.3}{6.3} the centralizer \(C\) of \(T\) is
\(u^{-1}(C')\), and is therefore connected because \(Z\) is
connected. This proves \textup{\SGARefLink{sga3:xii:6.6:item:b}{(b)}}. If \(T\) is maximal, then
\(T'\) is maximal; hence \(C'\) is nilpotent, and so is \(C\),
which is a central extension of \(C'\). Moreover, \(T\) is a
maximal torus of \(C\), and by \SGARefLink{sga3:xii:6.7}{6.7} it is the unique maximal torus
of \(C\). Consequently
\[
  T\longmapsto\operatorname{Centr}_G(T)
\]
is a bijection from the maximal tori of \(G\) to the Cartan
subgroups of \(G\). We also have
\[
  \operatorname{Centr}_G(T)=C
  \subset \operatorname{Norm}_G(C)
  \subset \operatorname{Norm}_G(T).
\]
The centralizer of \(T\) has finite index in its normalizer
\((\text{cf. XI~5.9})\); the argument there remains valid without
the affineness hypothesis, because it uses only the representability
of the two functors in question, as noted in \SGARefLink{sga3:xi:unit:6:5}{XI~6.5}. Since \(C\) is
smooth and connected, it follows that
\[
  C=\operatorname{Norm}_G(C)^{0}.
\]
Furthermore, the bijection between maximal tori and Cartan
subgroups shows that the normalizers of \(C\) and \(T\) have the
same \(k\)-valued points. Since the latter normalizer is smooth,
\[
  \operatorname{Norm}_G(T)=\operatorname{Norm}_G(C).
\]

To finish the proof of \textup{\SGARefLink{sga3:xii:6.6:item:c}{(c)}}, it remains to show that a
smooth, connected, nilpotent subgroup \(C\) of \(G\) which has
finite index in its normalizer is a Cartan subgroup. Since \(Z\) is
central, the normalizer of \(C\) contains \(Z\); and since \(Z\) is
smooth and connected, it follows that \(Z\subset C\). Hence
\[
  C=u^{-1}(C'),\qquad C'=u(C),
\]
and
\[
  \operatorname{Norm}_G(C)
  =u^{-1}\bigl(\operatorname{Norm}_{G'}(C')\bigr).
\]
Thus \(C'\) is connected and nilpotent and has finite index in its
normalizer. By Theorem~1 of \(\mathrm{BIBLE}\ 7\), it is a Cartan
subgroup of \(G'\), and it follows at once that \(C\) is a Cartan
subgroup of \(G\).

We next prove \textup{\SGARefLink{sga3:xii:6.6:item:e}{(e)}}. We are in the situation considered at
the beginning of the proof, with \(\ker u=Z\) and \(G'=H\), and we
have already seen that \(T\mapsto u(T)\) is a bijection between the
maximal tori of \(G\) and those of \(G'\). Using the bijection just
proved between maximal tori and Cartan subgroups, we obtain a
bijection between Cartan subgroups \(C\) of \(G\) and Cartan
subgroups \(C'\) of \(G'\): to
\[
  C=\operatorname{Centr}_G(T)
\]
we associate
\[
  C'=\operatorname{Centr}_{G'}(T'),
  \qquad T'=u(T).
\]
Since \(C'\) is connected by \textup{\SGARefLink{sga3:xii:6.6:item:b}{(b)}}, \SGARefLink{sga3:xii:6.3}{6.3} gives
\[
  C'=u(C),\qquad C=u^{-1}(C'),
\]
which proves \textup{\SGARefLink{sga3:xii:6.6:item:e}{(e)}}.

It remains to prove \textup{\SGARefLink{sga3:xii:corollary:6-7:equation:d}{(d)}}. Let \(T\) be a maximal torus of
\(G\). We first prove that \(v(T)\) is a maximal torus of \(H\);
by the conjugacy theorem \textup{\SGARefLink{sga3:xii:6.6:item:a}{(a)}}, this will imply that all
maximal tori of \(H\) arise in this way. Let \(R\) be a torus of
\(H\) containing \(v(T)\). Replacing \(H\) by \(R\) and \(G\) by
\(v^{-1}(R)_{\mathrm{red}}^{0}\), we may suppose that \(R=H\), so
that \(H\) is a torus; it remains to prove \(v(T)=H\). Put again
\[
  Z=\operatorname{Centr}(G)_{\mathrm{red}}^{0},
  \qquad G'=G/Z,\qquad H'=H/v(Z).
\]
We have a commutative diagram with exact rows
\begin{equation}
\tag{D}\label{sga3:xii:corollary:6-7:equation:d}
\begin{tikzcd}[column sep=large,row sep=large]
e \arrow[r] &
Z \arrow[r] \arrow[d,"v''"'] &
G \arrow[r,"u"] \arrow[d,"v"'] &
G' \arrow[r] \arrow[d,"v'"'] &
e \\
e \arrow[r] &
v(Z) \arrow[r] &
H \arrow[r,"u'"'] &
H' \arrow[r] &
e .
\end{tikzcd}
\end{equation}
We already know that \(u(T)\) is a maximal torus of \(G'\).
Because \(G'\) is affine, the epimorphism
\[
  v':G'\longrightarrow H'
\]
is an epimorphism of smooth connected affine groups. Therefore
\(v'(u(T))\) is a maximal torus of \(H'\)
\((\mathrm{BIBLE}\ 7,\ \mathrm{Th.}\ 3\textup{(a)})\), and hence
equals \(H'\). Thus \(u'(v(T))=H'\), so in order to prove
\(v(T)=H\) it is enough to show that \(v(T)\supset v(Z)\). Now
\(v(Z)\) is a smooth connected subgroup of \(H\), hence a torus,
and
\[
  v'':Z\longrightarrow v(Z)
\]
is an epimorphism. By \SGARefLink{sga3:xii:6.4}{6.4},
\[
  v(Z)=v''(S)
\]
for a torus \(S\) of \(Z\). Since \(S\) is central in \(G\), it is
contained in the maximal torus \(T\), and therefore
\(v(Z)\subset v(T)\). This proves \textup{\SGARefLink{sga3:xii:6.6:item:d}{(d)}} for maximal tori.

In view of the bijection between maximal tori and Cartan subgroups,
it remains to prove that, if \(T\) is a maximal torus of \(G\) and
\(C\) its centralizer, then \(v(C)\) is the centralizer of \(v(T)\).
Return to diagram \textup{(D)}, now without assuming that \(H\) is
a torus, and put
\[
  T'=u(T),\qquad C'=u(C).
\]
By \textup{\SGARefLink{sga3:xii:6.6:item:e}{(e)}}, \(C'\) is the centralizer of the maximal torus
\(T'\). Since \(G'\), and hence \(H'\), is affine,
\(v'(C')\) is the centralizer of the maximal torus \(v'(T')\) of
\(H'\)
\((\mathrm{BIBLE}\ 7,\ \mathrm{Th.}\ 3\textup{(a)})\).
Equivalently, \(u'(v(C))\) is the centralizer of \(u'(v(T))\).
Because \(C\) contains \(Z\), the group \(v(C)\) contains \(v(Z)\).
It is therefore the inverse image under \(u'\) of \(u'(v(C))\),
that is, of the centralizer of \(u'(v(T))\). By
\textup{\SGARefLink{sga3:xii:6.6:item:e}{(e)}}, applied to \(u':H\to H'\) and the maximal torus
\(v(T)\) of \(H\), this inverse image is precisely the centralizer
of \(v(T)\). This completes the proof of \SGARefLink{sga3:xii:6.6}{6.6}.
\end{proof}

% END INCLUDED SGA3_English_current_buildable_source_20260823/source/inputs/XII/tex/components/01_expose_XII_sections4_6_en.tex

% BEGIN INCLUDED SGA3_English_current_buildable_source_20260823/source/inputs/XII/tex/components/02_expose_XII_sections7_9_bibliography_en.tex
% Exposé XII Sections 7--9 and bibliography.

\section{Application to smooth group preschemes not necessarily affine}
\SGARefTarget{sga3:xii:section:7}

\paragraph{\textbf{Theorem 7.1.}}
\SGARefTarget{sga3:xii:7.1}
\emph{Let \(S\) be a prescheme and let \(G\) be an \(S\)-group
prescheme that is smooth, separated, and of finite type over \(S\).
Suppose that \(G\) admits a maximal torus locally for the faithfully
flat quasi-compact topology. Then:}

\begin{enumerate}
\item[\emph{\SGAInlineRefTarget{sga3:xii:7.1:item:a}(a)}] \emph{The map}
\[
  T\longmapsto \operatorname{Centr}_G(T)
\]
\emph{induces a bijection from the set of maximal tori of \(G\)
onto the set of Cartan subgroups of \(G\). If \(C\) corresponds to
\(T\), then \(T\) is the unique maximal torus of \(C\).}

\item[\emph{\SGAInlineRefTarget{sga3:xii:7.1:item:b}(b)}] \emph{Let \(T,T'\) be two maximal tori of \(G\),
and let \(C,C'\) be the corresponding Cartan subgroups. Then}
\[
  \operatorname{Transp}_G(T,T')
  =\operatorname{Transp}_G(C,C')
  =\operatorname{Transp}_G(T,C').
\]
\emph{The first two terms are also identical to the strict
transporters. Moreover, the functor in question is representable by
a closed subprescheme of \(G\), smooth over \(S\). The tori \(T,T'\)
and the Cartan subgroups \(C,C'\) are locally conjugate for the étale
topology.}

\item[\emph{\SGAInlineRefTarget{sga3:xii:7.1:item:c}(c)}] \emph{A maximal torus of \(G\), and a Cartan
subgroup of \(G\), exist locally for the étale topology.}

\item[\emph{\SGAInlineRefTarget{sga3:xii:theorem:7-1:item:d}\SGAInlineRefTarget{sga3:xii:7.1:item:d}(d)}] \emph{Suppose that every finite subset of a fiber
\(G_s\) is contained in an affine open of \(G\)---for example, that
\(G\) is quasi-projective over \(S\), or that \(S\) is artinian.
Then the functor}
\[
  \mathcal T:(\mathrm{Sch})^\circ/S\longrightarrow(\mathrm{Sets})
\]
\emph{defined in \SGARefLink{sga3:xii:1.10}{1.10}, the functor of maximal tori of \(G\), is
isomorphic to the functor}
\[
  \mathcal C:(\mathrm{Sch})^\circ/S\longrightarrow(\mathrm{Sets})
\]
\emph{of Cartan subgroups of \(G\), and is representable by a smooth,
separated prescheme of finite type over \(S\). This prescheme is
quasi-projective over \(S\) when \(G\) is, and is affine over \(S\)
when \(G\) is affine over \(S\), or when \(S\) is artinian.}

\item[\emph{\SGAInlineRefTarget{sga3:xii:7.1:item:e}(e)}] \emph{Let \(u:G\to G'\) be a morphism of
\(S\)-group preschemes, where \(G'\) is smooth, separated, and of
finite type over \(S\), and suppose that}
\[
  u_s(G_s)=G'_s
\]
\emph{for every \(s\in S\), that is, that \(u_s\) is faithfully
flat.\footnote{The editors did not locate the printed reference in
Exposé VI\(_B\); see M.\ Demazure and P.\ Gabriel,
\emph{Groupes algébriques}, I, Masson (1970),
Proposition II.5.1\SGARefLink{sga3:xii:7.1:item:c}{(c)}.} For every maximal torus \(T\) of \(G\),
\(u(T)\) is a maximal torus of \(G'\). If \(G'\) has connected
fibers, then for every Cartan subgroup \(C\) of \(G\), \(u(C)\) is a
Cartan subgroup of \(G'\); and if \(C\) is the centralizer of \(T\),
then \(u(C)\) is the centralizer of \(u(T)\). In either case the
induced morphism \(T\to u(T)\), respectively \(C\to u(C)\), is
faithfully flat.}

\item[\emph{\SGAInlineRefTarget{sga3:xii:7.1:item:f}(f)}] \emph{Under the hypotheses of \textup{\SGARefLink{sga3:xii:7.1:item:e}{(e)}}, suppose
in addition that \(\ker u\) is a central subgroup of \(G\). Then}
\[
  T\longmapsto u(T)
\]
\emph{is a bijection from the set of maximal tori of \(G\) onto the
set of maximal tori of \(G'\). If \(G'\) has connected fibers, then}
\[
  C\longmapsto u(C)
\]
\emph{is likewise a bijection from the set of Cartan subgroups of
\(G\) onto the set of Cartan subgroups of \(G'\); if \(C'=u(C)\),
then \(C=u^{-1}(C')\).}
\end{enumerate}

\paragraph{\textbf{Remarks 7.2.}}
\SGARefTarget{sga3:xii:7.2}
We shall see in Exposé XV that the conclusion of
\SGARefLink{sga3:xii:7.1}{Theorem~7.1}\textup{\SGARefLink{sga3:xii:7.1:item:d}{(d)}} remains valid without the restrictive
condition on the finite subsets of the \(G_s\). We shall also prove
there the conclusions concerning Cartan subgroups alone that occur
in \textup{\SGARefLink{sga3:xii:7.1:item:b}{(b)}}, \textup{\SGARefLink{sga3:xii:7.1:item:c}{(c)}}, \textup{\SGARefLink{sga3:xii:7.1:item:d}{(d)}}, and \textup{\SGARefLink{sga3:xii:7.1:item:e}{(e)}}, under
the weaker hypothesis that \(G\) admit a Cartan subgroup locally for
the faithfully flat quasi-compact topology, but not necessarily a
maximal torus. For that purpose we shall moreover need to use \SGARefLink{sga3:xii:7.1}{7.1}
when \(S\) is artinian. The method of Exposé XV also considerably
simplifies the proof of \textup{\SGARefLink{sga3:xii:7.1:item:c}{(c)}} and \textup{\SGARefLink{sga3:xii:7.1:item:d}{(d)}} when \(S\) is
not artinian.

\begin{proof}[Proof of \SGARefLink{sga3:xii:7.1}{Theorem 7.1}]
\emph{\SGAInlineRefTarget{sga3:xii:7.2:item:a}(a)} Proceeding as in \SGARefLink{sga3:xii:3.2}{3.2}, one sees that for every maximal
torus \(T\) of \(G\),
\[
  C=\operatorname{Centr}_G(T)
\]
is indeed a Cartan subgroup of \(G\), and that \(T\) is determined
by \(C\) as the unique maximal torus of \(C\). It remains only to
show that every Cartan subgroup \(C\) of \(G\) is defined by a
maximal torus of \(G\), or equivalently that \(C\) admits a central
maximal torus.

The question is still local for the fpqc topology, so we may suppose
that \(G\) admits a maximal torus \(T\). By \SGARefLink{sga3:xi:theorem:6:2}{XI~6.2},
\[
  \operatorname{Transp}_G(T,C)
\]
is representable by a closed subprescheme \(S'\) of \(G\), smooth
over \(S\); moreover \(S'\to S\) is surjective by \SGARefLink{sga3:xii:6.6:item:c}{6.6(c)}. After the
base change \(S'\to S\), we may therefore suppose that there is a
section \(g\) of \(G\) such that
\[
  \operatorname{ad}(g)\cdot T\subset C.
\]
Then \(\operatorname{ad}(g)\cdot T\) is a maximal torus of \(C\).
It is central fiber by fiber by \SGARefLink{sga3:xii:6.6:item:c}{6.6(c)}, and hence central by
\SGARefLink{sga3:ix:theorem:5:6:item:b}{IX~5.6(b)}. This proves \textup{\SGARefLink{sga3:xii:7.2:item:a}{(a)}}.

\emph{\SGAInlineRefTarget{sga3:xii:7.2:item:b}(b)} Let \(T,T'\) be two maximal tori of \(G\), and \(C,C'\)
their corresponding Cartan subgroups. The following conditions are
equivalent:
\[
\begin{array}{ccccc}
(1)\ T\subset T',&
(2)\ T=T',&
(3)\ T\subset C',&
(4)\ C\subset C',&
(5)\ C=C'.
\end{array}
\]
This follows immediately from \textup{\SGARefLink{sga3:xii:7.2:item:a}{(a)}}. Applying the same result
after arbitrary base change gives the asserted identities among the
transporters and strict transporters.

We have already observed in \textup{\SGARefLink{sga3:xii:7.2:item:a}{(a)}} that
\(\operatorname{Transp}_G(T,C')\) is representable by a closed
subprescheme of \(G\), smooth over \(S\), whose structural morphism
is surjective. By Hensel's lemma this prescheme has a section locally
for the étale topology. Thus \(T,T'\), on the one hand, and \(C,C'\),
on the other, are locally conjugate for the étale topology. This
proves \textup{\SGARefLink{sga3:xii:7.2:item:b}{(b)}}.

\emph{\SGAInlineRefTarget{sga3:xii:7.2:item:c}(c)} Suppose first that \(S\) is artinian and local. If \(G\)
admits a maximal torus \(T\), the conjugacy theorem just proved shows
that the functor \(\mathcal T\) of maximal tori of \(G\) is
represented by the homogeneous space \(G/N\), where \(N\) is the
normalizer of \(T\) in \(G\); it is smooth by \textup{\SGARefLink{sga3:xii:7.2:item:b}{(b)}}.
Moreover, as observed in VI\(_A\)~\SGARefLink{par:sga3-VIA-3.2.1}{3.2.1}, every finite subset of
\(\mathcal T\) is contained in an affine open, because \(\mathcal T\)
is a homogeneous space under \(G\).

In general there is a finite extension \(k'/k\) of the residue field
of \(S\) such that \(G_{k'}\) has a maximal torus. The extension
\(k'/k\) is obtained by a finite flat base change \(S'\to S\), and
the maximal torus of \(G_{k'}\) lifts to a maximal torus of
\(G_{S'}\) by \SGARefLink{sga3:xi:proposition:2:1:bis}{XI~2.1 bis}. Thus \(\mathcal T_{S'}\) is represented by
a smooth, separated prescheme of finite type over \(S'\), every
finite subset of which is contained in an affine open. Its natural
descent datum is effective. Hence \(\mathcal T\) is representable,
and descent shows that it is smooth, separated, and of finite type
over \(S\). By smoothness and Hensel's lemma, \(\mathcal T\) has a
section locally for the étale topology. This proves \textup{\SGARefLink{sga3:xii:7.2:item:c}{(c)}} and
\textup{\SGARefLink{sga3:xii:7.1:item:d}{(d)}} when \(S\) is artinian. We shall prove below that in
this case \(\mathcal T\) is in fact affine over \(S\).

Now let \(S\) be arbitrary. To prove \textup{\SGARefLink{sga3:xii:7.2:item:c}{(c)}} and \textup{\SGARefLink{sga3:xii:7.1:item:d}{(d)}},
which are local assertions on \(S\) for the Zariski topology, we may
suppose that \(S\) is affine, that there is a prime \(\ell\) prime
to all residue characteristics of \(S\), and that the reductive rank
of the fibers of \(G\) is constant, say equal to \(r\). The latter
rank is visibly locally constant because of the assumed fpqc-local
existence of a maximal torus.

Choose a faithfully flat quasi-compact morphism \(S'\to S\), with
\(S'\) affine, such that \(G_{S'}\) admits a maximal torus \(T_0\),
and let \(C_0\) be its centralizer. I claim that there is a suitable
power \(n\) of \(\ell\) such that
\[
  C_0=\operatorname{Centr}_{G_{S'}}({}_nT_0).
\]
Indeed, this reduces immediately to the case where \(S'\) is
noetherian, which was treated in \SGARefLink{sga3:xi:theorem:6:2}{XI~6.2}. Fix such an \(n\), and put
\[
  M=(\mathbb Z/n\mathbb Z)^r,
  \qquad
  P=\operatorname{Hom}_{S\text{-gr}}(M_S,G).
\]
Then \(P\) is represented by a closed subprescheme of finite
presentation of \(({}_nG)^r\), where
\[
  {}_nG=\operatorname{Hom}_{S\text{-gr}}
    ((\mathbb Z/n\mathbb Z)_S,G).
\]
This is also the kernel of the \(n\)-th-power morphism on \(G\), and
is therefore represented by a closed subprescheme of finite
presentation of \(G\). Moreover, \(P\) is smooth over \(S\) by
\SGARefLink{sga3:xi:proposition:2:1}{XI~2.1}.

Let \(s\in S\). From what precedes, there is a finite separable
extension \(k''/\kappa(s)\) for which \(G_{k''}\) has a maximal
torus \(T''\). Since \({}_nT''\) is étale over \(k''\), we may,
after replacing \(k''\) by a finite separable extension, suppose
that
\[
  {}_nT''\simeq M_{k''}.
\]
Choose an étale morphism \(S''\to S\) and a point \(s''\in S''\)
above \(s\), inducing the residue-field extension
\(\kappa(s'')\simeq k''\). We obtain a section of
\(P_{S''}\otimes\kappa(s'')\) over
\(\operatorname{Spec}\kappa(s'')\). By Hensel's lemma, after
replacing the pointed prescheme \((S'',s'')\) by an étale
neighborhood if necessary, this section extends to a section of
\(P_{S''}\). Equivalently, there is a homomorphism
\[
  M_{S''}\longrightarrow G_{S''}
\]
inducing the given isomorphism
\[
  M_{\kappa(s'')}\simeq{}_nT''.
\]
By \SGARefLink{sga3:ix:theorem:6:4}{IX~6.4}---which here reduces to Nakayama's lemma---this
homomorphism is a closed immersion above an
open neighborhood of \(s''\), which we may take to be all of
\(S''\). Let \(H''\) be its image, and let \(C''\) be its
centralizer in \(G_{S''}\); by \SGARefLink{sga3:xi:theorem:6:2}{XI~6.2}, \(C''\) is a closed smooth
subprescheme of \(G_{S''}\).

\paragraph{\textbf{Lemma 7.3.}}
\SGARefTarget{sga3:xii:7.3}
\emph{Under the preceding hypotheses on \(\ell,n,r\), for every
prescheme \(S''\) over \(S\) and every maximal torus \(T''\) of
\(G_{S''}\), one has}
\[
  \operatorname{Centr}_{G_{S''}}(T'')
  =
  \operatorname{Centr}_{G_{S''}}({}_nT'').
\]

Indeed, after faithfully flat descent from \(S'\times_S S''\), we
reduce to the case in which \(G_{S''}\) admits a maximal torus
\(T''_1\) for which the displayed equality holds. By
\textup{(b)}, \(T''\) and \(T''_1\) are locally conjugate for the
étale topology, so the same equality holds for \(T''\).

Apply this result over \(\operatorname{Spec}\kappa(s'')\). It shows
that the fiber \(C''_{s''}\) is a Cartan subgroup of \(G_{s''}\).
Now, since \(G\) is smooth over \(S\), the union of the identity
components of the fibers \(G_s\) is an open subprescheme of \(G\);
this is the general EGA IV result for smooth morphisms recorded by
the editors without a precise reference. This open is visibly
stable under the group law and hence is a subgroup, denoted \(G^0\).
It satisfies the same preliminary hypotheses as \(G\), and maximal
tori of \(G\) correspond bijectively to maximal tori of \(G^0\).
Thus, for the proof of \textup{(c)} and \textup{(d)}, we may and do
suppose that \(G\) has connected fibers. Its Cartan subgroups then
also have connected fibers by \SGARefLink{sga3:xii:6.6:item:a}{6.6(a)}.

I claim that \(C''{}^0\) is a Cartan subgroup of \(G_{S''}\) above
an open neighborhood of \(s''\). This will complete the proof of
\textup{\SGARefLink{sga3:xii:7.2:item:c}{(c)}}. The claim follows from the next lemma.

\paragraph{\textbf{Lemma 7.4.}}
\SGARefTarget{sga3:xii:7.4}
\emph{Under the hypotheses of \SGARefLink{sga3:xii:7.1}{7.1}, suppose that \(G\) has connected
fibers. Let \(C\) be a closed group subprescheme of \(G\), smooth
over \(S\), and let \(s\in S\) be such that \(C_s\) is a Cartan
subgroup of \(G_s\). Then \(C^0\) is a Cartan subgroup of \(G\) over
an open neighborhood of \(s\).}

We reduce easily to the case where \(S\) is local and \(s\) is its
closed point, and to proving there that \(C^0\) is a Cartan subgroup
of \(G\). By flat descent we may further suppose that \(G\) admits a
maximal torus \(T\). By \SGARefLink{sga3:xi:theorem:6:2}{XI~6.2},
\[
  \operatorname{Transp}_G(T,C)
\]
is represented by a closed subprescheme of \(G\), smooth over \(S\).
The fiber of this transporter at \(s\) is nonempty by the hypothesis
on \(C_s\) and the conjugacy \SGARefLink{sga3:xii:6.6:item:a}{theorem 6.6(a)}. Faithfully flat descent
therefore reduces us to the case where the transporter has a section
over \(S\), hence to the case where \(C\) contains a maximal torus
\(T\) of \(G\). Then \(T\) is central in \(C^0\) by \SGARefLink{sga3:ix:theorem:5:6:item:a}{IX~5.6(a)}, so
\[
  C^0\subset \operatorname{Centr}_G(T).
\]
This is an inclusion of smooth group preschemes over \(S\), with
connected fibers and the same relative dimension---namely, the
dimension of their common fiber at \(s\). It is therefore an
equality, proving the lemma and completing \textup{\SGARefLink{sga3:xii:7.2:item:c}{(c)}}.

\emph{(d)} Retain the preceding notation and hypotheses on
\(\ell,n,r\), and the assumption that \(G\) has connected fibers.
Let
\[
  Q:(\mathrm{Sch})^\circ/S\longrightarrow(\mathrm{Sets})
\]
be the functor defined by
\[
  Q(S')=
  \left\{
  \begin{array}{c}
  \text{subgroups of multiplicative type in \(G_{S'}\)}\\
  \text{of type \(M=(\mathbb Z/n\mathbb Z)^r\)}
  \end{array}
  \right\}
\]
in the sense of \SGARefLink{sga3:ix:definition:1:4}{IX~1.4}. The rule
\[
  T\longmapsto{}_nT
\]
defines a morphism
\[
  \varphi:\mathcal T\longrightarrow Q,
\]
which is a monomorphism by \SGARefLink{sga3:xii:7.3}{Lemma~7.3}.

I claim that \(Q\) is representable by a separated prescheme of
finite presentation over \(S\). As noted in the proof of
\SGARefLink{sga3:xi:proposition:3:12:item:a}{XI~3.12(a)}, there is an isomorphism
\[
  Q\simeq P_0/\Gamma,
\]
where \(P_0\) is the open and closed subprescheme of
\[
  P=\operatorname{Hom}_{S\text{-gr}}(M_S,G)
\]
corresponding to monomorphisms \(M_{S'}\to G_{S'}\), as in \SGARefLink{sga3:ix:theorem:6:8}{IX~6.8},
and
\[
  \Gamma=(\operatorname{Aut}_{\mathrm{gr}}(M))_S.
\]
The hypothesis that every finite subset of a fiber \(G_s\) lies in
an affine open is stable under passage to a closed subprescheme and
under fiber products. It is therefore inherited by \(({}_nG)^r\),
then by \(P\), and then by \(P_0\). Hence \(Q\) is represented by a
prescheme of finite presentation over \(S\).\footnote{The editors
left the printed reference to Exposé V unspecified.} The same
argument shows that \(Q\) is quasi-projective, respectively affine,
over \(S\) when \(G\) is. In every case \(Q\) is separated over
\(S\).

\paragraph{\textbf{Lemma 7.5.}}
\SGARefTarget{sga3:xii:7.5}
\emph{The preceding morphism
\(\varphi:\mathcal T\to Q\) is representable by an open immersion.}

In other words, let \(H\) be a subgroup of multiplicative type in
\(G\), of type \(M\). We must prove that there is an open subset
\(U\subset S\) with the following property: for every \(S'\to S\),
the group \(H_{S'}\) is of the form \({}_nT_{S'}\), for a suitable
maximal torus \(T_{S'}\) of \(G_{S'}\), if and only if
\(S'\to S\) factors through \(U\).

We may suppose that the nilpotent rank of the fibers of \(G\) is
constant, say equal to \(r'\); it is locally constant by
\textup{\SGARefLink{sga3:xii:7.2:item:b}{(b)}}. Put
\[
  C=\operatorname{Centr}_G(H),
\]
a closed group subprescheme smooth over \(S\) by \SGARefLink{sga3:xi:theorem:6:2}{XI~6.2}. Replacing
\(S\) by the open and closed subset where \(C\) has relative
dimension \(r'\), we may suppose that this is its relative dimension
everywhere. Then \(H={}_nT\) for a maximal torus \(T\) of \(G\) if
and only if \(C^0\) admits a central maximal torus \(T\), everywhere
of relative dimension \(r\), and \(H={}_nT\). This gives another
description of the subfunctor \(U\).

By flat descent, we may suppose that \(G\) admits a maximal torus
\(T_1\). Consider the subfunctor
\[
  R=
  \operatorname{Transp}_G
  \bigl(T_1,\operatorname{Centr}_{C^0}(C^0)\bigr)
  \subset
  \operatorname{Transp}_G(T_1,C).
\]
The transporter on the right is represented by a smooth prescheme
over \(S\) by \SGARefLink{sga3:xi:theorem:6:2}{XI~6.2}. The one on the left is represented by the
induced open subprescheme, as follows immediately from \SGARefLink{sga3:ix:theorem:5:6:item:a}{IX~5.6(a)};
in particular it too is smooth over \(S\). Its structural morphism
is therefore open. Replacing \(S\) by its image, equipped with the
induced structure, we may suppose that this morphism is surjective.
Applying \SGARefLink{sga3:xii:1.13}{1.13} to \(C^0\), we find that \(C^0\) admits a central
maximal torus \(T\), because it admits one fpqc-locally. Its relative
dimension is that of \(T_1\), namely \(r\). The remaining condition
on \(H\) is \(H={}_nT\), and by \SGARefLink{sga3:ix:corollary:2:10}{IX~2.10} this again cuts out a suitable
open and closed subset of \(S\). This proves the lemma.

\SGARefLink{sga3:xii:7.5}{Lemma~7.5} shows that \(\mathcal T\) is represented by a separated
prescheme locally of finite presentation over \(S\), and in fact of
finite presentation. For the latter assertion, one may either
inspect the proof of \SGARefLink{sga3:xii:7.5}{Lemma~7.5} to see that \(\varphi\) is a
quasi-compact open immersion, or descend fpqc-locally to the case
where \(G\) has a maximal torus \(T\), when
\[
  \mathcal T\simeq G/\operatorname{Norm}_G(T).
\]
This expression, or \SGARefLink{sga3:xi:proposition:2:1}{XI~2.1}, also shows that \(\mathcal T\) is smooth
over \(S\). If \(G\) is quasi-projective over \(S\), then so are
\(Q\) and \(\mathcal T\). If \(G\) is affine over \(S\), the
affineness of \(\mathcal T\) over \(S\) was proved in 5.4. It is
unknown here whether, without the hypothesis that \(G/S\) is
affine, \(n\) can be chosen so that the open immersion of Lemma~7.5
is also closed.

It remains to show that \(\mathcal T\) is affine over \(S\) when
\(S\) is artinian. By EGA I~6.1.7, we reduce to the case where
\(S\) is the spectrum of a field, which we may suppose algebraically
closed. In view of \textup{(f)}, proved below, it suffices to prove
the assertion for
\[
  G/\operatorname{Centr}(G).
\]
This quotient is affine by 6.1, so the preceding affine case
applies. This proves \textup{(d)}.

\emph{\SGAInlineRefTarget{sga3:xii:7.5:item:e}(e)} By IX~6.8, there is a subtorus \(T'\) of \(G'\) such
that \(u\) induces a faithfully flat morphism
\[
  T\longrightarrow T'.
\]
This characterizes \(T'\) as the subsheaf \(u(T)\) of \(G'\). Let
\(C\) and \(C'\) be the centralizers of \(T\) and \(T'\),
respectively. We prove that \(C\to C'\) is flat, and faithfully flat
when \(G'\) has connected fibers. Since \(C\) and \(C'\) are flat
and of finite presentation over \(S\), SGA~1, I~5.9 reduces us to a
base field, which we may suppose algebraically closed. Replacing
\(G,G'\) by their identity components does not change \(C^0,C'{}^0\),
so we may suppose \(G,G'\) connected. The assertion then follows
from 6.6(d), together with \textup{(a)}.

\emph{(f)} By \textup{(a)} and \textup{(e)}, it is enough to prove
the assertion about Cartan subgroups. A Cartan subgroup \(C\) of
\(G\), being the centralizer of a maximal torus, contains the center
of \(G\), and therefore contains \(\ker u\). Thus
\[
  C=u^{-1}(C'),
  \qquad C'=u(C),
\]
where \(C'\) is the Cartan subgroup of \(G'\) supplied by
\textup{\SGARefLink{sga3:xii:7.5:item:e}{(e)}}. Hence \(C\mapsto u(C)\) is injective. To prove
surjectivity, it suffices to show that \(u^{-1}(C')\) is a Cartan
subgroup of \(G\) for every Cartan subgroup \(C'\) of \(G'\).

The question is fpqc-local. We may suppose that \(G\) admits a
Cartan subgroup \(C_1\). Then \(C'_1=u(C_1)\) is a Cartan subgroup
of \(G'\), and is fpqc-locally conjugate to \(C'\) by \textup{(b)}.
Since
\[
  u^{-1}(C'_1)=C_1
\]
is a Cartan subgroup of \(G\), the same is true of \(u^{-1}(C')\).
This proves \textup{(f)} and completes the proof of the theorem.
\end{proof}

Alternatively, to prove that \(u^{-1}(C')\) is a Cartan subgroup,
observe that it is flat over \(S\), since \(u\) is; SGA~1, I~5.9
then reduces the assertion, by definition, to the case of a base
field, where 6.6(e) applies.

\paragraph{\textbf{Corollary 7.6.}}
\SGARefTarget{sga3:xii:7.6}
\emph{Under the hypotheses of 7.1\textup{(e)}, \(G'\) admits a
maximal torus locally for the étale topology, and therefore satisfies
the preliminary hypotheses imposed on \(G\). If moreover \(\ker u\)
is central, then the functors \(\mathcal T_G,\mathcal C_G\) of
maximal tori and Cartan subgroups of \(G\) are isomorphic to the
analogous functors \(\mathcal T_{G'},\mathcal C_{G'}\) for \(G'\).
Consequently, whenever these functors are representable, they are
represented by isomorphic \(S\)-preschemes.}

\paragraph{\textbf{Remark 7.7.}}
\SGARefTarget{sga3:xii:7.7}
\emph{(a)} In contrast with the case where \(G\) is affine over
\(S\)---in fact it is enough that \(G\) have affine fibers, as will
be seen in Exposé XVI---local constancy of the reductive rank of
\(G\) does not imply that \(G\) admits a maximal torus fpqc-locally.
For example, let \(S\) be the spectrum of a discrete valuation ring,
and let \(G_1,G_2\) be smooth separated group schemes of finite type
over \(S\), such that the generic fiber of \(G_1\) is an elliptic
curve and its special fiber is \(\mathbb G_m\), while the generic
fiber of \(G_2\) is \(\mathbb G_m\) and its special fiber is
\(\mathbb G_a\). Put
\[
  G=G_1\times_S G_2.
\]
Both fibers of \(G\) have reductive rank \(1\), but \(G\) plainly
does not admit a maximal torus fpqc-locally.

It is, on the other hand, very plausible that the following
condition on a smooth separated group of finite type over a
prescheme \(S\) is sufficient for the étale-local existence of a
maximal torus: the reductive rank and the abelian rank of the fibers
of \(G\) are locally constant functions.

\emph{(b)} In the proof of \SGARefLink{sga3:xii:7.1}{7.1}, particularly in \textup{\SGARefLink{sga3:xii:7.1:item:a}{(a)}}, we
used \SGARefLink{sga3:xi:theorem:6:2}{XI~6.2} in cases where \(S\) was not assumed locally noetherian.
In the instances at hand, however, the reduction to affine
noetherian \(S\) is immediate.

Here is a variant of \SGARefLink{sga3:xii:7.1:item:b}{7.1(b)}.

\paragraph{\textbf{Proposition 7.8.}}
\SGARefTarget{sga3:xii:7.8}
\emph{Let \(G\) be a smooth \(S\)-group prescheme of finite
presentation with connected fibers, let \(H\) be a smooth
\(S\)-group prescheme, and let}
\[
  i:H\longrightarrow G
\]
\emph{be a monomorphism of \(S\)-groups, making \(H\) an
\(S\)-subgroup of \(G\). For every Cartan subgroup
\(C\) of \(G\),\footnote{The proof only requires \(C\) to be a
smooth subgroup of \(G\) whose geometric fibers are connected
centralizers of subgroups of maximal tori.} the transporter
\(\operatorname{Transp}_G(C,H)\) is represented by a closed
subprescheme of \(G\), smooth over \(S\).}

The representability of the transporter by a closed subprescheme of
\(G\), of finite presentation, follows from \SGARefLink{sga3:xi:corollary:6:11}{XI~6.11}, because \(C\)
has connected fibers when \(G\) does, by \SGARefLink{sga3:xii:6.6:item:b}{6.6(b)}. To prove smoothness,
the standard procedure reduces first to affine noetherian \(S\),
then to artinian local \(S\), and finally, by descent, to an
algebraically closed residue field. In that case \(C\) has a maximal
torus \(T\), which is also a maximal torus of \(G\). Unless the
reductive and nilpotent ranks of the fiber \(H_0\) equal those of
\(G_0\), the transporter is empty. Under this equality one sees,
using the connectedness of \(C\) and the smoothness of the
centralizer in \(H\) of a maximal torus of \(H\), that
\[
  \operatorname{Transp}_G(C,H)
  =
  \operatorname{Transp}_G(T,H).
\]
The right-hand side is smooth by \SGARefLink{sga3:xi:proposition:2:5}{XI~2.5}, and therefore so is the
left-hand side.

The same argument proves the more general assertion
\SGARefLink{sga3:xii:7.9:item:b:02}{7.9(b)} below.

\paragraph{\textbf{Proposition 7.9.}}
\SGARefTarget{sga3:xii:7.9}
\emph{Let \(G\) and \(i:H\to G\) be as in \SGARefLink{sga3:xii:7.8}{7.8}. Suppose moreover
that, for every \(s\in S\), \(H_s\) is connected and has the same
reductive rank and the same nilpotent rank as \(G_s\); equivalently,
\(H_s\) contains a Cartan subgroup of \(G_s\). Then:}

\begin{enumerate}
\item[\emph{(a)}] \emph{\(\operatorname{Norm}_G(H)\) is represented
by a closed subprescheme of \(G\), of finite presentation over
\(S\), and the canonical monomorphism}
\[
  H\longrightarrow\operatorname{Norm}_G(H)
\]
\emph{is an open immersion. Consequently \(i\) is an immersion and}
\[
  H=\operatorname{Norm}_G(H)^0.
\]

\item[\emph{(b)}] \emph{For every Cartan subgroup \(C\) of \(G\),
\(\operatorname{Transp}_G(C,H)\) is a closed subprescheme of \(G\),
smooth over \(S\), with surjective structural morphism. If \(C\) is
the centralizer of a maximal torus \(T\) of \(G\), then moreover}
\[
  \operatorname{Transp}_G(T,H)
  =
  \operatorname{Transp}_G(C,H).
\]

\item[\emph{(c)}] \emph{Let \(C\) be a group subprescheme of \(H\).
Then \(C\) is a Cartan subgroup of \(H\) if and only if it is a
Cartan subgroup of \(G\).}

\item[\emph{(d)}] \emph{If \(G\) admits a Cartan subgroup,
respectively a maximal torus, locally for the étale topology or for
the fpqc topology, then so does \(H\).}
\end{enumerate}

\begin{proof}
\emph{\SGAInlineRefTarget{sga3:xii:7.9:item:a:02}(a)} The representability of \(\operatorname{Norm}_G(H)\) by
a closed subprescheme of \(G\), of finite presentation over \(S\),
is contained in \SGARefLink{sga3:xi:corollary:6:11}{XI~6.11}. Since \(H\) is smooth, hence flat and
locally of finite presentation over \(S\), VI\(_B\)~\SGARefLink{cor:sga3-VIB-2.6}{2.6} reduces the
assertion that
\[
  H\longrightarrow\operatorname{Norm}_G(H)
\]
is an open immersion to the fibers. We therefore reduce to the case
where \(S=\operatorname{Spec}k\), with \(k\) algebraically closed.

The assertion then reduces to showing that the corresponding map on
Lie algebras is an isomorphism,\footnote{The editors did not locate
the printed reference in Exposé VI\(_B\); see M.\ Demazure and
P.\ Gabriel, \emph{Groupes algébriques}, I, Masson (1970),
Corollary II.5.6.} or equivalently that
\[
  (\mathfrak g/\mathfrak h)^H=0,
\]
where \(\mathfrak g,\mathfrak h\) are the Lie algebras of \(G,H\)
and the superscript denotes invariants; compare II~5.2.3(i). By
hypothesis \(H\) contains a Cartan subgroup \(C\) of \(G\), the
centralizer of a maximal torus \(T\) of \(G\). It therefore suffices
to prove
\[
  (\mathfrak g/\mathfrak h)^T=0.
\]
By complete reducibility of representations of \(T\)
(\SGARefLink{sga3:i:proposition:4-7-3}{I~4.7.3}), this follows from
\[
  (\mathfrak g/\mathfrak c)^T=0,
  \qquad \mathfrak c=\operatorname{Lie}(C).
\]
The latter equality is equivalent to
\[
  \mathfrak g^T=\mathfrak c.
\]
It says that the centralizer \(C\) and the normalizer of \(T\) have
the same Lie algebra, and follows from the fact that \(C\) is an
open subgroup of that normalizer by \SGARefLink{sga3:xi:proposition:5:9}{XI~5.9}. This proves
\textup{\SGARefLink{sga3:xii:7.9:item:a:02}{(a)}}.

\emph{\SGAInlineRefTarget{sga3:xii:7.9:item:b:02}(b)} was proved in \SGARefLink{sga3:xii:7.8}{7.8}.

\emph{\SGAInlineRefTarget{sga3:xii:7.9:item:c:02}(c)} Since \(H\) is a subprescheme of \(G\), the assertion
reduces immediately to an algebraically closed base field, where it
follows directly from the hypothesis on \(H\).

\emph{\SGAInlineRefTarget{sga3:xii:7.9:item:d:02}(d)} follows from \textup{\SGARefLink{sga3:xii:7.9:item:b:02}{(b)}}, \textup{\SGARefLink{sga3:xii:7.9:item:c:02}{(c)}}, and the Hensel
lemma \SGARefLink{sga3:xi:proposition:1:10}{XI~1.10}.
\end{proof}

\paragraph{\textbf{Corollary 7.10.}}
\SGARefTarget{sga3:xii:7.10}
\emph{Let \(G,H\) be as in \SGARefLink{sga3:xii:7.9}{7.9}, and suppose that for every
algebraically closed field \(k\) over \(S\), every element of
\(G_k(k)\) normalizing \(H_k\) belongs to \(H_k(k)\). Then \(H\) is
a closed subprescheme of \(G\) and is its own normalizer.}

This follows immediately from \SGARefLink{sga3:xii:7.9:item:a:02}{7.9(a)}. We shall apply it in
particular to Borel subgroups, and more generally to parabolic
subgroups, of \(G\).

\paragraph{\textbf{Corollary 7.11.}}
\SGARefTarget{sga3:xii:7.11}
\emph{Let \(G,H\) be as in \SGARefLink{sga3:xii:7.9}{7.9}. Then \(H\) contains every group
subprescheme \(Z\) of \(G\) which is central in \(G\), flat, and of
finite presentation over \(S\).}

As usual, we reduce first to affine noetherian \(S\), then to
artinian \(S\); the hypotheses of \SGARefLink{sga3:xii:7.1}{7.1} then hold. By \SGARefLink{sga3:xii:7.9:item:c:02}{7.9(c)}, we may
suppose that \(H\) contains a Cartan subgroup \(C\) of \(G\), and it
remains to show \(C\supset Z\). Since \(S\) is artinian, the quotient
\[
  G'=G/Z
\]
is represented by a group prescheme of finite type over \(S\); the
canonical morphism \(u:G\to G'\) is faithfully flat and has kernel
\(Z\) by VI\(_A\)~\SGARefLink{sga3:xii:3.2}{3.2}. The group \(G'\) is smooth over \(S\).
\SGARefLink{sga3:xii:7.1:item:f}{Theorem~7.1(f)} therefore shows that \(C=u^{-1}(C')\) for a Cartan
subgroup \(C'\) of \(G'\), and hence that \(C\) contains \(Z\).

\paragraph{\textbf{Corollary 7.12.}}
\SGARefTarget{sga3:xii:7.12}
\emph{Let \(G,G'\) be smooth \(S\)-group preschemes of finite
presentation, and let \(u:G\to G'\) be a faithfully flat group
homomorphism, meaning that \(u_s\) is faithfully flat for every
\(s\in S\). Suppose that \(\ker u\) is central and that \(G\) has
connected fibers. Then}
\[
  H'\longmapsto H=u^{-1}(H')
\]
\emph{defines a bijection between:}

\begin{itemize}
\item \emph{the group subpreschemes \(H'\) of \(G'\), smooth and of
finite presentation over \(S\), having at every \(s\in S\) the same
reductive rank and the same nilpotent rank as \(G'\); and}
\item \emph{the group subpreschemes \(H\) of \(G\), smooth and of
finite presentation over \(S\), having at every \(s\in S\) the same
reductive rank and the same nilpotent rank as \(G\).}
\end{itemize}

\emph{Moreover, \(H\) has connected fibers if and only if \(H'\)
does.}

Let \(H\) be a group subprescheme of \(G\) with the stated
properties. By \SGARefLink{sga3:xii:7.11}{7.11} it contains \(Z=\ker u\), so faithfully flat
descent gives
\[
  H=u^{-1}(H')
\]
for a uniquely determined group subprescheme \(H'\) of \(G'\).
Using \SGARefLink{sga3:xii:6.6:item:e}{6.6(e)}, one checks immediately that \(H'\) has the asserted
properties. If \(H\) has connected fibers, so does \(H'=H/Z\).

Conversely, start with \(H'\subset G'\) having the stated
properties, and put \(H=u^{-1}(H')\). In view of \SGARefLink{sga3:xii:6.6:item:e}{6.6(e)}, it remains
to show that \(H\) is smooth over \(S\), and that it has connected
fibers when \(H'\) does. Since \(u\) is flat and \(H'\) is flat over
\(S\), \(H\) is already flat over \(S\). It therefore suffices to
check smoothness, respectively connectedness, on geometric fibers,
reducing to an algebraically closed base field.

Then \(H'\) contains a Cartan subgroup \(C'\) of \(G'\), and \(H\)
contains its inverse image \(C\). By \SGARefLink{sga3:xii:6.6:item:e}{6.6(e)}, \(C\) is a Cartan
subgroup of \(G\), hence is smooth and connected. The conclusion now
follows from the next lemma.

\paragraph{\textbf{Lemma 7.13.}}
\SGARefTarget{sga3:xii:7.13}
\emph{Let \(G,G'\) be group preschemes flat and of finite
presentation over \(S\), let \(u:G\to G'\) be a faithfully flat
group homomorphism, and let \(C'\) be a group subprescheme of \(G'\),
of finite presentation over \(S\), such that
\(C=u^{-1}(C')\) is smooth over \(S\), respectively has connected
fibers. For every group subprescheme \(H'\) of \(G'\), of finite
presentation over \(S\), containing \(C'\), and smooth over \(S\),
respectively with connected fibers, its inverse image
\(H=u^{-1}(H')\) is smooth over \(S\), respectively has connected
fibers.}

As noted above, this reduces at once to the case where \(S\) is the
spectrum of a field. The group \(H\) is then a principal \(C\)-bundle
over \(H/C\), by VI\(_B\)~9, while
\[
  H/C\simeq H'/C'
\]
by Exposé IV. Since \(H'\) is smooth, respectively connected, so is
\(H'/C'\), hence so is \(H/C\). The same property holds for \(C\) by
hypothesis and therefore for the bundle \(H\). This proves the
lemma and \SGARefLink{sga3:xii:7.12}{Corollary~7.12}.

\section{Semisimple elements; unions and intersections of maximal
tori in group schemes not necessarily affine}
\SGARefTarget{sga3:xii:section:8}

Throughout this section, \(G\) denotes an \(S\)-group prescheme that
is smooth and of finite presentation over \(S\), with connected
fibers.

Suppose first that \(S\) is the spectrum of an algebraically closed
field \(k\). When \(G\) is affine, the notion of a semisimple element
of \(G(k)\) was defined in \emph{BIBLE}~4, No.~4; one checks at once
that this notion is invariant under an algebraically closed extension
\(k'/k\) of the base field. Moreover, by \emph{BIBLE}~6,
Theorem~5(c), an element \(g\in G(k)\) is semisimple if and only if
it is contained in a maximal torus of \(G\).

When \(G\) is no longer assumed affine, Chevalley's theorem gives a
canonical expression of \(G\) as an extension of an abelian variety
by a smooth connected affine algebraic group:
\begin{equation}
\tag{*}\label{sga3:xii:lemma:7-13:equation:star}
  e\longrightarrow V\longrightarrow G\longrightarrow A
  \longrightarrow e.
\end{equation}
We shall say that \(g\in G(k)\) is \emph{semisimple} if it is a
semisimple element of \(V(k)\). Since the maximal tori of \(V\) are
precisely the maximal tori of \(G\), this is equivalent to saying
that \(g\) belongs to a maximal torus of \(G\). This notion is again
invariant under algebraically closed extensions of the base field.

Suppose next that \(S=\operatorname{Spec}k\) for an arbitrary field
\(k\), and let \(g\in G\). After choosing an algebraically closed
extension \(K\) of \(\kappa(g)\), the point \(g\) is the image of a
geometric point \(g'\) of \(G\) with values in an algebraically
closed extension \(K\) of \(k\). We call \(g\) semisimple when \(g'\)
is semisimple. What precedes shows that this is independent of the
choice of \(g'\). If \(k'/k\) is any extension, the set of
semisimple elements of \(G_{k'}\) is the inverse image of the set
\(G^{\mathrm{ss}}\) of semisimple elements of \(G\).

Finally, for arbitrary \(S\), a point \(g\in G\) is called
semisimple when it is semisimple in its fiber \(G_s\). Thus, for
every morphism \(S'\to S\), the subset \(G^{\mathrm{ss}}_{S'}\) is
the inverse image of \(G^{\mathrm{ss}}\). Suppose that the functor
\(\mathcal T\) of maximal tori defined in \SGARefLink{sga3:xii:1.10}{1.10} is representable by a
prescheme of finite presentation over \(S\). This is so, for
example, if \(G\) admits a maximal torus locally for the fpqc
topology, by \SGARefLink{sga3:xii:7.1:item:d}{7.1(d)}, at least when \(G\) is quasi-projective over
\(S\). Let \(T\) be the canonical maximal torus of
\(G_{\mathcal T}\), the ``universal maximal torus of \(G\)'', and
consider the morphism
\[
  u:T\longrightarrow G
\]
induced by the projection \(G_{\mathcal T}\to G\). It follows
immediately from the definition that \(G^{\mathrm{ss}}\) is the
image of this morphism. We obtain:

\paragraph{\textbf{Proposition 8.1.}}
\SGARefTarget{sga3:xii:8.1}
\emph{The set \(G^{\mathrm{ss}}\) of semisimple elements of \(G\)
is locally constructible. It is therefore constructible when \(S\),
and hence \(G\), is quasi-compact and quasi-separated.}

As usual, we reduce to affine noetherian \(S\). By the standard
noetherian criterion for constructibility (EGA \(0_{\mathrm{III}}\),
9.2.3), we may further suppose that \(S\) is integral and need only
show that there is a nonempty open \(U\subset S\) such that
\(G^{\mathrm{ss}}|_U\) is constructible. After shrinking \(U\), and
replacing it if necessary by a finite covering, we may suppose that
\(G\) is separated over \(S\) and contains a maximal torus. By
7.1(d), \(\mathcal T\) is then represented by a prescheme of finite
presentation over \(S\); the same is true of its universal torus,
whose image in \(G\) is therefore constructible.

Suppose again that \(k\) is a field, and let \(H\) be the subgroup
of \(G\) generated by the preceding morphism (VI\(_B\)~\SGARefLink{sga3:xii:1.2}{1.2}). It is
a smooth group subscheme of \(G\), connected because the universal
torus is connected, and its formation commutes with every extension
of the base field, as does the formation of \(\mathcal T\).

When \(k\) is algebraically closed, \(H\) is also the algebraic
subgroup of \(G\) generated by the maximal tori of \(G\), or,
equivalently, by all the tori of \(G\). It is likewise the smallest
algebraic subgroup of \(G\) containing all semisimple elements of
\(G(k)\). These characterizations remain valid whenever \(k\) is
infinite, because, as will be proved in Exposé XIV, the set of
\(k\)-rational points of \(\mathcal T\) is dense in \(\mathcal T\).

The subgroup \(H\) is invariant in \(G\). Indeed, it is enough to
check this after passing to an algebraically closed field; since
\(H\) and \(G\) are then smooth, it suffices to check stability
under the inner automorphisms \(\operatorname{int}(g)\),
\(g\in G(k)\), which is immediate. One can even show that \(H\) is
characteristic in \(G\), that is, stable under
\(\operatorname{Aut}_{k\text{-}\mathrm{gr}}(G)\).

It follows from \SGARefLink{sga3:xii:6.6:item:d}{6.6(d)} that \(G/H\) has reductive rank zero. More
precisely, if \(K\) is an invariant algebraic subgroup of \(G\),
then, over an algebraically closed field, the maximal tori of \(G/K\)
are the direct images of the maximal tori of \(G\). Hence \(G/K\)
has reductive rank zero if and only if \(K\) contains all maximal
tori of \(G\), or, over an arbitrary field, if and only if
\(K\supset H\). Thus \(H\) is the smallest invariant algebraic
subgroup of \(G\) for which \(G/H\) has reductive rank zero. It is
also the smallest algebraic subgroup of \(G\) having the same
reductive rank as \(G\).

Finally, \(H\) is affine. This may again be checked after passing to
an algebraic closure. In the exact sequence \((*)\), every maximal
torus of \(G\) is contained in \(V\); consequently \(H\subset V\),
and \(V\) is affine. We have proved:

\paragraph{\textbf{Proposition 8.2.}}
\SGARefTarget{sga3:xii:8.2}
\emph{Let \(G\) be a smooth connected algebraic group over a field
\(k\), and let \(\bar k\) be an algebraic closure of \(k\). There
exists an algebraic subgroup \(H\) of \(G\) such that
\(H_{\bar k}\) is the algebraic subgroup of \(G_{\bar k}\) generated
by the maximal tori, or equivalently by all the tori, of
\(G_{\bar k}\). The group \(H\) is also characterized as the
smallest algebraic subgroup of \(G\) having the same reductive rank
as \(G\), or as the smallest invariant algebraic subgroup of \(G\)
such that \(G/H\) has reductive rank zero. It is a smooth, connected,
invariant, affine subgroup of \(G\), and its formation commutes with
every extension of the base field.}

The characterization in terms of \(G/H\) gives the following
consequences.

\paragraph{\textbf{Corollary 8.3.}}
\SGARefTarget{sga3:xii:8.3}
\emph{Let \(G\) be a smooth connected algebraic group over an
algebraically closed field \(k\). Then \(G\) has reductive rank zero
---equivalently, in the notation of \SGARefLink{sga3:xii:8.2}{8.2}, \(H=(e)\)---if and only if
\(G\) is an extension of an abelian variety by a smooth connected
unipotent algebraic group, that is, a successive extension of groups
isomorphic to the additive group \(\mathbb G_a\).}

Indeed, Chevalley's exact sequence \((*)\) reduces the assertion to
the fact that a smooth connected affine group \(V\) has reductive
rank zero if and only if it is unipotent; this is contained in
\emph{BIBLE}~\SGARefLink{sga3:xii:6.4}{6.4}, Theorem~4, \SGARefLink{sga3:xii:bibliography}{Corollary~3}. Thus, in the notation of
\SGARefLink{sga3:xii:8.2}{8.2}, \(H\) is the smallest invariant algebraic subgroup of \(G\)
such that \((G/H)_{\bar k}\) is an extension of an abelian variety
by a smooth connected affine unipotent algebraic group.

\paragraph{\textbf{Corollary 8.4.}}
\SGARefTarget{sga3:xii:8.4}
\emph{In the notation of \SGARefLink{sga3:xii:8.2}{8.2}, \(G=H\), or equivalently
\(G_{\bar k}\) is generated by its maximal tori, if and only if
\(G\) is affine and every homomorphism
\(G_{\bar k}\to\mathbb G_a\) is trivial.}

\paragraph{\textbf{Remarks 8.5.}}
\SGARefTarget{sga3:xii:8.5}
\begin{enumerate}
\item[\emph{\SGAInlineRefTarget{sga3:xii:8.5:item:a}(a)}] Let \(V\) be the largest smooth connected affine
algebraic subgroup of \(G_{\bar k}\), so that \(G_{\bar k}\) is an
extension of an abelian variety by \(V\). It is well known, by
Rosenlicht,\footnote{The editors ask that references be supplied
here.} that when \(k\) is not perfect, \(V\) is not in general
defined over \(k\): there need not exist an algebraic subgroup
\(V_0\) of \(G\) such that \((V_0)_{\bar k}=V\). If, however, \(V\)
is generated by its maximal tori---equivalently, if \(V\) has no
quotient isomorphic to \(\mathbb G_{a,\bar k}\)---then such a group
does exist, namely the group \(H\) of \SGARefLink{sga3:xii:8.2}{8.2}. Thus, in this rationality
question as in many others (compare Exposé XIV, No.~6), all the
difficulties arise from unipotent groups, hence from the additive
group, whereas the presence of sufficiently many groups of
multiplicative type makes descent work.

\item[\emph{(b)}] In the notation of \SGARefLink{sga3:xii:8.2}{8.2}, let \(H'\) be the inverse
image in \(G\) of the commutator subgroup of \(G/H\). Then \(H'\) is
the smallest invariant algebraic subgroup of \(G\) such that
\((G/H')_{\bar k}\) is a commutative extension of an abelian variety
by a smooth connected affine unipotent algebraic group. It is again
a smooth connected affine subgroup of \(G\).

Assume that the characteristic \(p\) is positive, and let \(H''\)
be the inverse image in \(G\) of \(p^n(G/H')\) for \(n\) sufficiently
large. Then \(H''/H'\) is an abelian variety, and \(H''\) is the
smallest invariant algebraic subgroup of \(G\) such that \(G/H''\)
is a commutative smooth connected affine unipotent algebraic group.
It follows that \(V=H_{\bar k}\), in the notation of \textup{\SGARefLink{sga3:xii:8.5:item:a}{(a)}},
if and only if \(H''=G\), or equivalently if and only if every
homomorphism \(G_{\bar k}\to\mathbb G_a\) is trivial.
\end{enumerate}

We now generalize to smooth groups with connected fibers the notion
of reductive center developed in \SGARefLink{sga3:xii:section:4}{Section~4}, following the idea of
\SGARefLink{sga3:xii:4.10}{4.10}. Let \(\mathcal Z\) be the subfunctor of \(G\) defined by
\[
\begin{split}
  \mathcal Z(S')=\{\, &g'\in G_{S'}(S'):\text{ for every }
  S''\to S'\text{ and every maximal torus }T''\subset G_{S''},\\
  &\text{the pullback }g''\text{ of }g'\text{ is a section of }
  T''\text{ over }S''\,\}.
\end{split}
\]
Let \(\mathcal T(S')\) be the set of maximal tori of \(G_{S'}\), and
let \(\widetilde{\mathcal T}(S')\) be the set of pairs \((T',g')\)
where \(T'\) is a maximal torus of \(G_{S'}\) and \(g'\) is a section
of \(T'\) over \(S'\). Then \(\widetilde{\mathcal T}\) is a
subfunctor of
\[
  \mathcal T\times_S G=\mathcal T_G,
\]
and in this notation one may also write
\[
  \mathcal Z=
  \prod_{\mathcal T_G/G}
  \widetilde{\mathcal T}/\mathcal T_G.
\]
Using \SGARefLink{sga3:xi:theorem:6:8}{XI~6.8} and \SGARefLink{sga3:xii:7.1:item:d}{7.1(d)}, which under suitable hypotheses represents
\(\mathcal T\) by a smooth \(S\)-prescheme, one could deduce
representability of \(\mathcal Z\) under corresponding hypotheses.
We shall instead obtain it more directly below.

\paragraph{\textbf{Definition 8.6.}}
\SGARefTarget{sga3:xii:8.6}
\emph{Let \(G\) be an \(S\)-group prescheme, smooth and of finite
presentation with connected fibers. We say that \(G\) admits a
\emph{reductive center} if the preceding functor \(\mathcal Z\),
which is evidently a subgroup of \(G\), is representable by a group
of multiplicative type. In that case \(\mathcal Z\) is called the
reductive center of \(G\).}

If \(Z\) is a reductive center of \(G\), then \(Z_{S'}\) is a
reductive center of \(G_{S'}\) for every base change \(S'\to S\).
The existence of a reductive center is consequently local for the
fpqc topology. The terminology is justified because \(\mathcal Z\)
is invariant under
\[
  \operatorname{Aut}_S(G),
\]
and hence is an invariant subgroup of \(G\); by \SGARefLink{sga3:ix:corollary:5:5}{IX~5.5}, it is
central.

\SGARefLink{sga3:xii:4.5}{Lemma~4.5} is replaced here by the following statement.

\paragraph{\textbf{Lemma 8.7.}}
\SGARefTarget{sga3:xii:8.7}
\emph{Let \(u:H\to G\) be a central homomorphism, where \(H\) is of
multiplicative type and of finite type over \(S\). Suppose that for
every algebraically closed field \(k\) over \(S\), the subgroup
\(u_k(H_k)\) is contained in the largest smooth connected affine
subgroup of \(G_k\). Then \(u\) factors through every maximal torus
\(T\) of \(G\). In particular, \(u\) factors through the subfunctor
\(\mathcal Z\) defined above.}

An easy variant of \SGARefLink{sga3:ix:theorem:5:1:bis}{IX~5.1 bis}, with equality replaced by inclusion,
reduces the assertion to the case in which \(S\) is the spectrum of
an algebraically closed field: first reduce to affine noetherian
\(S\), then to artinian \(S\), and use \SGARefLink{sga3:ix:theorem:3:6}{IX~3.6}. Since \(T\) is
contained in the largest smooth connected affine subgroup \(V\) of
\(G\), we are reduced to \(G=V\), hence to the affine case proved
in \SGARefLink{sga3:xii:4.5}{4.5}.

\paragraph{\textbf{Proposition 8.8.}}
\SGARefTarget{sga3:xii:8.8}
\emph{Let \(G\) be an \(S\)-group prescheme, smooth and of finite
presentation with connected fibers.}

\begin{enumerate}
\item[\emph{(a)}] \emph{If \(S=\operatorname{Spec}k\), then \(G\)
admits a reductive center \(Z\). When \(G\) is an extension of an
abelian variety by a smooth connected affine algebraic group \(V\)
---for example, when \(k\) is algebraically closed---then \(Z\) is
also the reductive center of \(V\), and it is the largest central
subgroup of multiplicative type of \(V\).}

\item[\emph{(b)}] \emph{Let \(Z\) be a group subprescheme of
multiplicative type in \(G\). Then \(Z\) is a reductive center of
\(G\) if and only if \(Z_s\) is a reductive center of \(G_s\) for
every \(s\in S\). It is then the largest subgroup of multiplicative
type \(K\) in \(G\) such that \(K_s\) is contained in the reductive
center of \(G_s\) for every \(s\in S\). More generally, let
\(u:H\to G\) be a homomorphism with \(H\) of multiplicative type and
of finite type over \(S\), and suppose that for every algebraically
closed field \(k\) over \(S\), \(u_k\) factors through the largest
smooth connected affine subgroup of \(G_k\). Then \(u\) factors
through \(Z\), and in particular is central.}

\item[\emph{(c)}] \emph{If \(G\) admits a maximal torus locally for
the fpqc topology, then \(G\) admits a reductive center.}

\item[\emph{(d)}] \emph{If \(T\) is a maximal torus of \(G\), then}
\[
  T\cap\operatorname{Centr}(G)
  =\ker\bigl(T\longrightarrow\operatorname{GL}(\mathfrak g)\bigr),
\]
\emph{and this group is a reductive center of \(G\).}

\item[\emph{(e)}] \emph{Let \(Z\) be a reductive center of \(G\),
and suppose that \(G'=G/Z\) is representable, for example when \(S\)
is artinian. Then the unit subgroup is a reductive center of \(G'\),
and}
\[
  T'\longmapsto T=u^{-1}(T')
\]
\emph{is a bijection from the set of maximal tori of \(G'\) to the
set of maximal tori of \(G\).}
\end{enumerate}

\begin{proof}
\emph{\SGAInlineRefTarget{sga3:xii:8.8:item:a:02}(a)} Let \(S=\operatorname{Spec}k\). To prove existence, we
may pass to an algebraic closure of \(k\), and thus suppose that
\(G\) is an extension of an abelian variety by a smooth connected
affine algebraic group \(V\). For every \(S'\) over \(k\), the
maximal tori of \(G_{S'}\) are those of \(V_{S'}\): this follows
from \SGARefLink{sga3:ix:corollary:5:2}{IX~5.2} and from the fact that a homomorphism from a torus to an
abelian variety over an algebraically closed field is trivial.
Consequently the functor \(\mathcal Z\) defined using \(G\) is the
same as the one defined using \(V\), and we reduce to affine \(G\).

The functor \(\mathcal T\) is representable and is necessarily
``essentially free'' over \(k\) (\SGARefLink{sga3:viii:definition:6:1}{VIII~6.1}); hence
\(\mathcal T_G\) is essentially free over \(G\). Since the universal
torus is a closed subscheme of
\(\mathcal T_G=G_{\mathcal T}\) (for example by \SGARefLink{sga3:viii:corollary:5:7}{VIII~5.7}),
\SGARefLink{sga3:viii:theorem:6:4}{VIII~6.4} shows that
\[
  \mathcal Z=
  \prod_{\mathcal T_G/G}
  \widetilde{\mathcal T}/\mathcal T_G
\]
is represented by a closed subscheme of \(G\). It is therefore a
group subscheme. It is of multiplicative type: after passing to an
algebraic closure, \(G\) admits a maximal torus \(T\), and by
definition \(Z\subset T\); a subgroup of a group of multiplicative
type is of multiplicative type (IX~8). Thus \(Z\) is a reductive
center. \SGARefLink{sga3:xii:8.7}{Lemma~8.7} shows that it is the largest central subgroup of
multiplicative type of \(V\).

\emph{\SGAInlineRefTarget{sga3:xii:8.8:item:b:02}(b)} Only sufficiency requires proof. Suppose that \(Z\) is
of multiplicative type and that \(Z_s\) is the reductive center of
\(G_s\) for every \(s\in S\). \SGARefLink{sga3:xii:8.7}{Lemma~8.7} first shows that \(Z\) is
contained in every maximal torus of \(G\), and this remains true
after every base change. It remains to show that if \(g\) is a
section of \(G\) over \(S\) whose pullback to every \(S'\) over
\(S\) is a section of every maximal torus of \(G_{S'}\), then \(g\)
is a section of \(Z\).

As usual, since the fpqc sheaf
\[
  \mathcal Z_0=
  \prod_{\mathcal T_G/G}
  \widetilde{\mathcal T}/\mathcal T_G
\]
commutes with filtered direct limits of rings, we reduce successively
to affine, noetherian, and artinian \(S\). By \SGARefLink{sga3:xii:7.1:item:d}{7.1(d)},
\(\mathcal T\) is then represented by a smooth prescheme over \(S\).
Arguing as in \textup{\SGARefLink{sga3:xii:8.8:item:a:02}{(a)}}, \(\mathcal Z_0\) is represented by a
closed subscheme \(Z_0\) of \(G\). We already know \(Z\subset Z_0\);
by hypothesis \(Z_k=(Z_0)_k\) over the residue field \(k\), and
\(Z\) is flat over \(S\). It follows that \(Z=Z_0\), proving that
\(Z\) is the reductive center. The remaining assertions of
\textup{\SGARefLink{sga3:xii:8.8:item:b:02}{(b)}} follow from \SGARefLink{sga3:xii:8.7}{8.7}.

\emph{\SGAInlineRefTarget{sga3:xii:8.8:item:c:02}(c)} follows immediately from \textup{\SGARefLink{sga3:xii:8.8:item:d:02}{(d)}}.

\emph{\SGAInlineRefTarget{sga3:xii:8.8:item:d:02}(d)} Here \(\mathfrak g\) denotes the Lie algebra of \(G\),
and \(T\to\operatorname{GL}(\mathfrak g)\) is induced by the adjoint
representation. Trivially,
\[
  T\cap\operatorname{Centr}(G)
  \subset
  \ker\bigl(T\longrightarrow\operatorname{GL}(\mathfrak g)\bigr).
\]
The proof of the reverse inclusion is the same as in \SGARefLink{sga3:xii:4.7:item:d}{4.7(d)}. Put
\[
  Z=\ker\bigl(T\longrightarrow\operatorname{GL}(\mathfrak g)\bigr).
\]
Then \(Z\) is of multiplicative type (\SGARefLink{sga3:ix:theorem:6:8}{IX~6.8}). By \textup{\SGARefLink{sga3:xii:8.8:item:b:02}{(b)}}, we
may check that it is a reductive center over a field. Let \(Z_0\)
be the reductive center supplied by \textup{\SGARefLink{sga3:xii:8.8:item:a:02}{(a)}}. Plainly
\[
  Z_0\subset T\cap\operatorname{Centr}(G)=Z.
\]
Conversely, \(Z\) is a central subgroup of multiplicative type
contained in the smooth connected affine group \(T\), so
\textup{\SGARefLink{sga3:xii:8.8:item:a:02}{(a)}} gives \(Z\subset Z_0\). Thus \(Z=Z_0\).

\emph{\SGAInlineRefTarget{sga3:xii:8.8:item:e:02}(e)} Under the hypotheses of \SGARefLink{sga3:xii:7.1:item:f}{7.1(f)}, put \(Z=\ker u\), and
suppose that for every algebraically closed field \(k\) over \(S\),
\(Z_k\) is contained in a maximal torus of \(G_k\). One checks
readily that
\[
  T'\longmapsto u^{-1}(T')
\]
is a bijection from the maximal tori of \(G'\) to those of \(G\).
Applying this observation to the situation of \textup{\SGARefLink{sga3:xii:8.8:item:e:02}{(e)}} proves
the assertion.
\end{proof}

\paragraph{\textbf{Remarks 8.9.}}
\SGARefTarget{sga3:xii:8.9}
\begin{enumerate}
\item[\emph{(a)}] The proof of \SGARefLink{sga3:xii:8.8}{8.8} is independent of the results of
Section~4; in particular, the proof of \SGARefLink{sga3:xii:8.8:item:a:02}{8.8(a)} does not use \SGARefLink{sga3:xii:4.4}{4.4},
whose proof is somewhat laborious.

\item[\emph{(b)}] The subfunctor \(\mathcal Z\) considered in \SGARefLink{sga3:xii:8.6}{8.6}
is always central, whether or not it is representable.\ It may
therefore be tempting to call it the reductive center of \(G\) in
all cases.

\item[\emph{(c)}] One may likewise generalize \SGARefLink{sga3:xii:4.9}{4.9}. If \(G\) is a
smooth connected algebraic group over an algebraically closed
field, then \(G\) is an extension of an abelian variety by a smooth
connected unipotent algebraic group---equivalently, \(G\) has
reductive rank zero, by \SGARefLink{sga3:xii:8.3}{8.3}---if and only if the reductive center of
\(G\) is the unit group and the Lie algebra of \(G\) is nilpotent.
\end{enumerate}

\section{Complement: actions of group schemes and fixed points}
\SGARefTarget{sga3:xii:section:9}

The purpose of this section, added in January 2008, is to study fixed
points under the action of a group scheme.

\subsection{Representability of the fixed-point functor}
\SGARefTarget{sga3:xii:9.1}

Let \(S\) be a scheme, let \(G\) be an \(S\)-group scheme, and
suppose that \(G\) acts on an \(S\)-scheme \(X\). Define the
subfunctor \(X^G\) of fixed points of \(X\) under \(G\) by requiring,
for every \(S\)-scheme \(T\), that \(X^G(T)\) be the subset of
\(X(T)\) consisting of the points fixed under \(G\) (VIII~6\SGARefLink{sga3:xii:8.8:item:e:02}{(e)}).

Recall the notion of an \(S\)-pure scheme introduced in
Raynaud--Gruson, \S3.3, and in Raynaud. Let \(G\) be a flat
\(S\)-group scheme of finite presentation. Since the fibers of \(G\)
have no embedded components, purity takes the following form:
\(G\) is \emph{\(S\)-pure} if, for every strictly henselian local
\(S\)-scheme \(T\) with closed point \(t\), every generic point
\(\widetilde x\) of a fiber of \(G\times_S T\) specializes to a
point of \(G_t\); equivalently, the closure of \(\widetilde x\) in
\(G\times_S T\) meets \(G_t\).

\paragraph{\textbf{Examples.}}
\begin{enumerate}
\item[\textup{(1)}] If \(G\) is quasi-finite and separated over
\(S\), then \(G\) is \(S\)-pure if and only if it is finite over
\(S\).

\item[\textup{(2)}] If \(S=\operatorname{Spec}R\) and
\(G=\operatorname{Spec}A\), then \(G\) is \(S\)-pure if and only if
\(A\) is a projective \(R\)-module (Raynaud--Gruson, 3.3.5). In
particular, a diagonalizable \(G\) is \(S\)-pure.

\item[\textup{(3)}] Every group scheme of multiplicative type is
\(S\)-pure, since purity is local for the étale topology on \(S\).

\item[\textup{(4)}] The scheme \(G\) is \(S\)-pure if it is proper
over \(S\), if its fibers are irreducible, or if \(S\) is
semilocal artinian.

\item[\textup{(5)}] In particular, every reductive \(S\)-group
scheme is \(S\)-pure.
\end{enumerate}

\paragraph{\textbf{Proposition 9.2.}}
\SGARefTarget{sga3:xii:9.2}
\emph{Let \(G\) be an \(S\)-group scheme of finite presentation,
flat and \(S\)-pure over \(S\), acting on an \(S\)-scheme \(X\) of
finite presentation. Then the fixed-point functor \(X^G\) is
represented by a closed subscheme of \(X\), of finite presentation
over \(S\).}

Indeed, let \(a\in X(S)\), and let \(H\) be the group subscheme of
\(G\) fixing \(a\). The point \(a\) is fixed under \(G\) if and only
if \(H=G\). The coincidence functor of \(H\) and \(G\)---whose value
on an \(S\)-scheme \(T\) is a singleton when
\[
  H\times_S T\longrightarrow G\times_S T
\]
is bijective and is empty otherwise---is represented by a
closed subscheme of \(S\), defined by a finite-type sheaf of ideals
(Raynaud--Gruson, 4.1.1). Apply this result with \(S=X\) and with
\(a\) the universal point of \(X\), that is, the identity element
of
\[
  X(X)=\operatorname{Hom}_S(X,X).
\]

\setcounter{subsection}{2}
\subsection{The infinitesimal obstruction}
\SGARefTarget{sga3:xii:9.3}

Suppose that \(G\) is a flat \(S\)-group scheme of finite
presentation acting on the left on a smooth \(S\)-prescheme \(X\).
We shall study the formal smoothness of \(X^G\).

Let \(S_0\) be a closed subscheme of \(S\), defined by a
quasi-coherent sheaf of ideals \(\mathcal I\) with
\(\mathcal I^2=0\), and write
\[
  j:S_0\longrightarrow S,\qquad
  G_0=G\times_S S_0,\qquad
  X_0=X\times_S S_0.
\]
Let \(\epsilon_0\in X^G(S_0)\) be a point that lifts to \(X(S)\).
We shall study the obstruction to lifting \(\epsilon_0\) to
\(X^G(S)\).

Since \(X\) is smooth, the set of lifts of \(\epsilon_0\) to \(X(S)\)
is a trivial principal homogeneous space under the abelian group
\[
  \operatorname{Hom}_{\mathcal O_{S_0}}
  \bigl(\epsilon_0^*\Omega^1_{X_0/S_0},\mathcal I\bigr)
\]
(\SGARefLink{sga3:iii:proposition:0-2}{III~0.2}, \SGARefLink{sga3:iii:remark:0-3}{0.3}). Put
\[
  L_0=
  \operatorname{Hom}_{\mathcal O_{S_0}}
  \bigl(\epsilon_0^*\Omega^1_{X_0/S_0},\mathcal I\bigr).
\]
This is an \(\mathcal O_{S_0}\)-module. Since \(\epsilon_0\) is
fixed under \(G\), it is naturally a
\(G_0\)-\(\mathcal O_{S_0}\)-module. Let
\[
  \rho_0:G_0\longrightarrow\operatorname{Aut}_{S_0}(L_0),
  \qquad
  \rho:G\longrightarrow\operatorname{Aut}_S(j_*L_0)
\]
be the associated representations.

\paragraph{\textbf{Lemma 9.4.}}
\SGARefTarget{sga3:xii:9.4}
\emph{There is a canonically defined class}
\[
  c(\epsilon_0)\in H^1(G,j_*L_0)
  \simeq H^1(G_0,L_0)
\]
\emph{such that \(\epsilon_0\) lifts to \(X^G(S)\) if and only if
\(c(\epsilon_0)=0\).}

The displayed adjunction is the one in \SGARefLink{sga3:iii:lemma:1-1-2}{III~1.1.2}. We work on the
small fppf site of \(S\). For every \(S\)-scheme \(T\), flat and of
finite presentation, write
\[
  G_T=G\times_S T,\qquad X_T=X\times_S T.
\]
Consider the sheaf \(\mathcal A\) on this site defined by
\[
  \mathcal A(T)=
  \{\,\text{lifts of }\epsilon_0\times_S T\text{ to }X(T)\,\}.
\]
Since the formation of \(L_0\) commutes with flat base change,
\(\mathcal A(T)\) is a trivial principal homogeneous space under
\[
  H^0(T,j_*L_0).
\]
Moreover, because \(\epsilon_0\) is fixed, every \(g\in G(T)\) acts
by affine automorphisms on \(\mathcal A(T)\), compatibly with the
action of \(G\) on \(j_*L_0\):
\[
  g(a_T+v)=g(a_T)+\rho(g)(v),
  \qquad v\in H^0(T,j_*L_0).
\]
Since \(G\) is flat and of finite type, we may take \(T=G\) and
\(g\) equal to the universal point \(\operatorname{id}_G\). We
thereby obtain an action of the fppf sheaf \(G\) on \(\mathcal A\).

Let \(a\) be a lift of \(\epsilon_0\) to \(X(S)\), and let
\(g^\sharp\) denote the universal point of \(G\). Define
\[
  v^\sharp\in H^0(G,j_*L_0)
\]
by the relation
\[
  \rho(g^\sharp)(v^\sharp)=g^\sharp\cdot a_G-a_G.
\]
For every \(S\)-prescheme \(Y\), set
\[
  z(Y):G(Y)\longrightarrow H^0(Y,j_*L_0),
  \qquad g\longmapsto v^\sharp(g).
\]
This defines a \(1\)-cocycle \(z\in Z^1(G,j_*L_0)\) (Exposé I,
\S5). Its class \(c(\epsilon_0)\) does not depend on the choice of
\(a\). In particular, if \(\epsilon_0\) lifts to \(X^G(S)\), then
\(c(\epsilon_0)=0\).

Conversely, if \(c(\epsilon_0)=0\), there is a section
\[
  w\in H^0(S,j_*L_0)
\]
such that, for every \(S\)-prescheme \(Y\) and every \(g\in G(Y)\),
\[
  z(Y)(g)=g\cdot w_Y-w_Y.
\]
Apply this with \(Y=G\) and \(g=g^\sharp\). It follows that
\(a-w\in X^G(S)\). Thus \(\epsilon_0\) lifts to \(X^G(S)\) if and
only if \(c(\epsilon_0)=0\).

\paragraph{\textbf{Remark 9.5.}}
\SGARefTarget{sga3:xii:9.5}
When \(S\) is affine and \(G\) is affine, flat, and of finite type
over \(S\), this obstruction can be refined to a class
\[
  \widetilde c(\epsilon_0)\in H^1_{G_0}(X_0,L_0),
\]
where the right-hand side is the \(G_0\)-equivariant cohomology
group defined by Wewers (Appendix~C). In this case it is unnecessary
to assume that \(\epsilon_0\) lifts to \(X(S)\).

\setcounter{subsection}{5}
\subsection{Smoothness of fixed points}
\SGARefTarget{sga3:xii:9.6}

\paragraph{\textbf{Theorem 9.7.}}
\SGARefTarget{sga3:xii:9.7}
\emph{Let \(G\) be a flat \(S\)-group scheme of finite type over a
noetherian base \(S\), acting on a smooth \(S\)-scheme \(X\).
Suppose that \(G\) is \(S\)-pure and that, for every geometric point
\(s\) of \(S\) with algebraically closed residue field,}
\[
  H^1\bigl(G_s,(\Omega^1_{X_s})^*\bigr)=0.
\]
\emph{Then the fixed-point functor \(X^G\) is represented by a
closed subscheme of \(X\), smooth over \(S\).}

By 9.2, \(X^G\) is represented by a closed subscheme of \(X\), of
finite presentation over \(S\). We may suppose
\(S=\operatorname{Spec}A\) with \(A\) local. By the smoothness
criterion of SGA~1, III~3.1(iii bis), it is enough to verify formal
smoothness for
\[
  S'=\operatorname{Spec}A'
\]
and its closed subscheme
\[
  S''=\operatorname{Spec}(A'/I'),
\]
where \(A'\) is an artinian local ring and \(I'\) is an ideal of
square zero. By dévissage, we may assume that \(I'\) is annihilated
by the maximal ideal of \(A'\). If \(k\) is the residue field of
\(A'\), then \(I'\) is therefore a \(k\)-vector space.

Let \(\epsilon''\in X^G(S'')\) be a point that we wish to lift to
\(X^G(S')\), and put
\[
\begin{gathered}
  X'=X\times_S S',\qquad G'=G\times_S S',\\
  X''=X\times_S S'',\qquad G''=G\times_S S''.
\end{gathered}
\]
Since \(S'\) is affine and \(X'\) is smooth over \(S'\),
\(\epsilon''\) lifts to \(X(S')\). By \SGARefLink{sga3:xii:9.4}{Lemma~9.4}, the obstruction to
lifting \(\epsilon''\) to \(X^G(S')\) is a class
\[
  c(\epsilon'')\in H^1(G'',L''),
\]
where
\[
  L''=
  \operatorname{Hom}_{\mathcal O_{S''}}
  \bigl((\epsilon'')^*\Omega^1_{X''/S''},I'\bigr).
\]
Because \(L''\) is a \(k\)-vector space, there is a canonical
isomorphism
\[
  H^1(G'',L'')
  \simeq
  H^1(G''\times_{A'/I'}k,L'').
\]
Let \(\bar k\) be an algebraic closure of \(k\). By \SGARefLink{sga3:ix:theorem:3:1}{IX~3.1} in the
affine case, and by \SGARefLink{sga3:xii:9.11}{Lemma~9.11} in general, there is an isomorphism
\[
  H^1(G''\times_{A'/I'}k,L'')\otimes_k\bar k
  \xrightarrow{\ \sim\ }
  H^1(G''\times_{A'/I'}\bar k,L''\otimes_k\bar k).
\]
The representation \(L''\otimes_k\bar k\) is a direct sum of copies
of
\[
  (\Omega^1_{X_{\bar k}})^*.
\]
The group on the right therefore vanishes by hypothesis. It follows
that \(H^1(G'',L'')=0\), so the obstruction vanishes and \(X^G\) is
formally smooth.

\paragraph{\textbf{Corollary 9.8.}}
\SGARefTarget{sga3:xii:9.8}
\emph{Let \(G\) be a flat \(S\)-group scheme of finite presentation
acting on a smooth \(S\)-scheme \(X\) of finite presentation.
Suppose that \(G\) admits a composition series whose factors are of
one of the following types:}
\begin{enumerate}
\item[\emph{(1)}] \emph{an abelian \(S\)-scheme, that is, a smooth
\(S\)-group scheme whose fibers are abelian varieties;}
\item[\emph{(2)}] \emph{an \(S\)-group scheme of multiplicative
type;}
\item[\emph{(3)}] \emph{a finite étale \(S\)-group scheme of degree
invertible on \(S\);}
\item[\emph{(4)}] \emph{a reductive \(S\)-group scheme, when \(S\)
is a \(\mathbb Q\)-scheme.}
\end{enumerate}
\emph{Then the fixed-point functor \(X^G\) is represented by a
closed subscheme of \(X\), of finite presentation and smooth over
\(S\).}

Indeed, suppose that
\[
  1\longrightarrow G'\longrightarrow G\longrightarrow G''
  \longrightarrow 1
\]
is an exact sequence of \(S\)-groups satisfying the hypotheses.
The group \(G''\) acts on \(X^{G'}\), and \(X^G\) is the fixed-point
functor of \(X^{G'}\) under \(G''\). We are therefore reduced to the
case of a group of one of the four types listed above. By passage to
the limit, we may suppose \(S\) noetherian. We verify the
cohomological criterion of \SGARefLink{sga3:xii:9.7}{Theorem~9.7} for a geometric fiber \(G_s\)
and a finite-dimensional linear representation \(V\) of \(G_s\).

For an abelian scheme, every morphism from \(G_s^n\) to \(V\) is
constant; consequently
\[
  H^i(G_s,V)=0\qquad(i>0).
\]
For a group scheme of multiplicative type, \(G_s\) is
diagonalizable, and \SGARefLink{sga3:i:theorem:5-3-3}{I~5.3.3} again gives
\[
  H^i(G_s,V)=0\qquad(i>0).
\]

In the remaining two cases the assertion follows from semisimplicity
of the linear representations of \(G_s\). If \(G_s\) is finite
étale of degree invertible in \(\kappa(s)\), it is a constant finite
group of invertible order, and every linear representation of
\(G_s\) is semisimple. In the reductive case,
\(\kappa(s)\) is algebraically closed of characteristic zero, and
every linear representation of \(G_s\) is semisimple
(Tauvel--Yu, Theorem~27.3.3).

\paragraph{\textbf{Remark 9.9.}}
\SGARefTarget{sga3:xii:9.9}
This applies in particular to an action of \(G/S\) on a smooth
\(S\)-group scheme \(H\). If \(G\) is of multiplicative type and
acts by inner automorphisms on \(H\), with \(H\) smooth and affine
over \(S\), one recovers \SGARefLink{sga3:xi:corollary:5:3}{XI~5.3}: the centralizer of \(G\) in \(H\)
is smooth.

\paragraph{\textbf{Remark 9.10.}}
\SGARefTarget{sga3:xii:9.10}
When \(S=\operatorname{Spec}k\) and \(G\) is a linearly reductive
algebraic group over an algebraically closed field \(k\), the result
is due independently to Fogarty (Theorem~5.4) and Iversen
(Proposition~1.3).

The proof of \SGARefLink{sga3:xii:9.7}{Theorem~9.7} used the following lemma, which is well
known for affine group schemes (\SGARefLink{sga3:ix:theorem:3:1}{IX~3.1}).

\paragraph{\textbf{Lemma 9.11.}}
\SGARefTarget{sga3:xii:9.11}
\emph{Let \(G\) be a group scheme over an affine scheme
\(S=\operatorname{Spec}A\), and let
\(f:S'=\operatorname{Spec}A'\to S\) be flat. Put
\(G'=G\times_S S'\). Then, for every quasi-coherent
\(G\)-\(\mathcal O_S\)-module \(M\), there are canonical
isomorphisms}
\[
  H^i(G,M)\otimes_A A'
  \xrightarrow{\ \sim\ }
  H^i(G',M\otimes_A A')
  \qquad(i\geq 0).
\]

For \(n\geq1\), the degree-\(n\) Hochschild cochains are determined
by their value at
\[
  (\operatorname{id},\ldots,\operatorname{id})
  \in G(G)\times\cdots\times G(G).
\]
Thus there is an isomorphism of \(A\)-modules
\[
  C^n(G,M)
  \simeq
  \Gamma\bigl(G^n,M\otimes_A\mathcal O_{G^n}\bigr).
\]
The group \(H^n(G,M)\) is the \(n\)-th cohomology group of the
complex
\[
  M\longrightarrow
  \Gamma(G,M\otimes_A\mathcal O_G)\longrightarrow
  \Gamma(G^2,M\otimes_A\mathcal O_{G^2})\longrightarrow
  \Gamma(G^3,M\otimes_A\mathcal O_{G^3})\longrightarrow\cdots.
\]
Since \(A'\) is flat over \(A\), EGA IV\(_1\), 1.7.21 gives
\[
\begin{split}
  \Gamma\bigl(G^n,M\otimes_A\mathcal O_{G^n}\bigr)\otimes_A A'
  \xrightarrow{\ \sim\ }{}&
  \Gamma\bigl((G')^n,\\
  &(M\otimes_A A')\otimes_{A'}
  \mathcal O_{(G')^n}\bigr).
\end{split}
\]
Taking cohomology proves the assertion.

\begin{thebibliography}{Ray72}
\SGARefTarget{sga3:xii:bibliography}
\addcontentsline{toc}{section}{Bibliography}

\item[]\emph{Editor's note: the following are additional
references cited in this Exposé.}

\bibitem[Fo73]{sga3-xii-Fo73}
J.\ Fogarty,
\emph{Fixed point schemes},
Amer.\ J.\ Math.\ \textbf{95} (1973), 35--51.

\bibitem[RG71]{sga3-xii-RG71}
M.\ Raynaud and L.\ Gruson,
\emph{Critères de platitude et de projectivité},
Invent.\ Math.\ \textbf{13} (1971), 1--89.

\bibitem[Iv72]{sga3-xii-Iv72}
B.\ Iversen,
\emph{A fixed point formula for action of tori on algebraic
varieties},
Invent.\ Math.\ \textbf{16} (1972), 229--236.

\bibitem[Ray72]{sga3-xii-Ray72}
M.\ Raynaud,
\emph{Flat modules in algebraic geometry},
Compositio Math.\ \textbf{24} (1972), 11--31.

\bibitem[TY06]{sga3-xii-TY06}
P.\ Tauvel and R.\ W.\ T.\ Yu,
\emph{Lie algebras and algebraic groups},
Springer-Verlag, 2006.

\bibitem[We05]{sga3-xii-We05}
S.\ Wewers,
\emph{Formal deformation of curves with group scheme action},
Ann.\ Inst.\ Fourier (Grenoble) \textbf{55} (2005), 1105--1165.

\end{thebibliography}

% END INCLUDED SGA3_English_current_buildable_source_20260823/source/inputs/XII/tex/components/02_expose_XII_sections7_9_bibliography_en.tex



% END INCLUDED SGA3_English_current_buildable_source_20260823/source/inputs/XII/body.tex


\SGAStartExpose{XIII}
% BEGIN INCLUDED SGA3_English_current_buildable_source_20260823/source/inputs/XIII/body.tex
% Mechanically extracted after \mainmatter from the frozen standalone master.

\markboth{Exposé XIII. Regular Elements}{Exposé XIII. Regular Elements}
% BEGIN INCLUDED SGA3_English_current_buildable_source_20260823/source/inputs/XIII/source/tex/components/01_expose_XIII_opening.tex
% Authority: Expo13.pdf, local pp. 1--3 / combined pp. 805--807.
% Stops before: “Nous supposons par la suite que S est le spectre...”

\chapter*{Exposé XIII\\Regular Elements of Algebraic Groups and Lie Algebras}
\addcontentsline{toc}{chapter}{Exposé XIII. Regular Elements of Algebraic Groups and Lie Algebras}
\SGARefTarget{sga3:xiii:expose}

\begin{center}
\textit{by A. Grothendieck}
\end{center}

\section{An auxiliary lemma on varieties with operators}
\SGARefTarget{sga3:xiii:section:1}

Let \(S\) be a prescheme, let \(G\) be an \(S\)-group prescheme acting
on the left on an \(S\)-prescheme \(V\), let \(W\) be a closed
\(S\)-subprescheme of \(V\), and let \(N\) be its stabilizer in \(G\).
Thus \(N\) is the subgroup of \(G\) whose points with values in an
\(S'\) over \(S\) are the elements \(g\in G(S')\) such that
\[
  g\cdot W_{S'}=W_{S'}.
\]
We endow \((\mathrm{Sch}/S)\) with the faithfully flat
quasi-compact topology and identify \(G\), \(V\), and \(W\) with the
corresponding sheaves.\footnote{The French source prints \(M\) in this
list; the surrounding sentence and all subsequent uses require \(W\).}
We shall therefore reason in the category of sheaves on
\((\mathrm{Sch}/S)\); throughout \SGARefLink{sga3:xiii:section:1}{this section}, ``locally'' refers to
the topology just specified on \((\mathrm{Sch}/S)\).

Since \(N\) is a sheaf, consider the quotient sheaf \(G/N\). It is
immediately seen to be isomorphic to the following functor: to every
\(S'\) over \(S\), it assigns the set of subsheaves \(W'\) of
\(V_{S'}\) that are locally conjugate to \(W_{S'}\) under \(G\).
Let \(X\) be the subsheaf of \(G/N\times_S V\) whose value on an
\(S'\) over \(S\) is the set of pairs \((W',v)\), where \(W'\) is as
above and \(v\) is a section of \(W'\) over \(S'\) (hence also a
section of \(V_{S'}\) over \(S'\)). Let \(Z\) be the inverse image of
\(X\) in \(G\times_S V\). We then have \SGARefLink{sga3:xiii:equation:x}{the cartesian diagram}
\SGARefTarget{sga3:xiii:diagram:x}
\[
\tag{x}\label{sga3:xiii:equation:x}
\begin{tikzcd}[column sep=large,row sep=large]
Z \arrow[r] \arrow[d,hook] & X \arrow[d,hook,"i"] \\
G\times_S V \arrow[r] & G/N\times_S V ,
\end{tikzcd}
\]
where \(i\) is the canonical immersion, while the lower horizontal
arrow is induced by the canonical morphism \(G\to G/N\).

This morphism sends \(g\in G(S')\) to the subsheaf
\(g\cdot W_{S'}\) of \(V_{S'}\). Consequently, \(Z(S')\) is the set
of pairs \((g,v)\in G(S')\times V(S')\) such that
\(v\in (g\cdot W_{S'})(S')\). It follows that \(Z\) is isomorphic to
the sheaf \(G\times_S W\), via the isomorphism
\[
  G\times_S W \xrightarrow{\sim} Z,\qquad
  (g,w)\longmapsto(g,g\cdot w).
\]
Thus \SGARefLink{sga3:xiii:equation:x}{the preceding cartesian diagram} yields \SGARefLink{sga3:xiii:equation:xx}{the cartesian diagram}
\SGARefTarget{sga3:xiii:diagram:xx}
\[
\tag{xx}\label{sga3:xiii:equation:xx}
\begin{tikzcd}[column sep=large,row sep=large]
G\times_S W \arrow[r,"q"] \arrow[d,"\lambda"'] &
  X \arrow[d,"i"] \\
G\times_S V \arrow[r] & G/N\times_S V
\end{tikzcd}
\]
where
\[
  \lambda(g,w)=(g,g\cdot w),\qquad
  q(g,w)=(\bar g,g\cdot w),
\]
where \(\bar g\) denotes the image of \(g\) under the canonical map
\(G(S')\to(G/N)(S')\).

Diagram \(\SGARefLink{sga3:xiii:equation:x}{\mathrm{(x)}}\) also shows that \(Z\to X\) makes \(Z\) a
principal bundle over \(X\), with \(N\) acting on the right by
\[
  (g,v)\cdot n=(gn,v).
\]
Accordingly, in diagram \(\SGARefLink{sga3:xiii:equation:xx}{\mathrm{(xx)}}\), the morphism
\(q:G\times_S W\to X\) makes \(G\times_S W\) a principal bundle over
\(X\), with right \(N\)-action
\[
  (g,w)\cdot n=(gn,n^{-1}\cdot w).
\]

The principal morphisms just introduced are summarized by
\SGARefTarget{sga3:xiii:diagram:d}
\[
\tag{D}\label{sga3:xiii:equation:d}
\begin{tikzcd}[column sep=large,row sep=large]
& G\times_S W \arrow[d,"q"] \arrow[dr,"\varphi"] & \\
G/N &
X \arrow[l,"p"] \arrow[r,"\psi"] \arrow[d,hook,"i"] &
V \\
& G/N\times_S V
  \arrow[ul,dashed,"\operatorname{pr}_1"]
  \arrow[ur,dashed,"\operatorname{pr}_2"'] &
\end{tikzcd}
\]
where
\[
  \psi=\operatorname{pr}_2\circ i,\qquad
  \varphi=\psi\circ q,\qquad
  \varphi(g,w)=g\cdot w.
\]

If \(v\) is a section of \(V\) over \(S\), let \(X_v\) be its inverse
image in \(X\) under \(\psi\). For every \(S'\) over \(S\),
\(X_v(S')\) is the set of subsheaves \(W'\) of \(V_{S'}\) that are
locally conjugate to \(W_{S'}\) under \(G\) and contain the section
\(v_{S'}\) of \(V_{S'}\). The inverse image of \(v\) in
\(G\times_S W\) under \(\varphi\) is isomorphic to the subsheaf
\(M_v\) of \(G\) defined by
\[
  M_v(S')=
  \bigl\{g\in G(S')\mid v_{S'}\in(g\cdot W_{S'})(S')\bigr\}
  =
  \bigl\{g\in G(S')\mid g^{-1}v_{S'}\in W(S')\bigr\}.
\]
If \(v\) is a section of \(W\), and not merely of \(V\), then
\(M_v\) evidently contains \(N\).

None of these explicit descriptions used the representability of
\(G\), \(V\), or \(W\) (nor even the fact that the site in question is
defined in terms of preschemes). Suppose now, however, that \(N\) is
representable, faithfully flat, and quasi-compact over \(S\), and that
\(G/N\) is representable. When \(S\) is the spectrum of a field and
\(G\) is of finite type over that field, this hypothesis is necessarily
satisfied (VIII~6 and VI\textsubscript{B}.11.18). Using the cartesian
diagram \(\SGARefLink{sga3:xiii:equation:xx}{\mathrm{(xx)}}\), faithfully flat quasi-compact descent, and
the fact that \(Z\to G\times_S V\) is a closed immersion, one sees that
\(X\) is representable: it is obtained by descending the closed
subprescheme \(Z\) of \(G\times_S V\) along the faithfully flat
quasi-compact morphism
\[
  G\times_S V\longrightarrow G/N\times_S V.
\]
Thus diagram \(\SGARefLink{sga3:xiii:equation:d}{\mathrm{(D)}}\) is a diagram of morphisms of
preschemes over \(S\).

% END INCLUDED SGA3_English_current_buildable_source_20260823/source/inputs/XIII/source/tex/components/01_expose_XIII_opening.tex

% BEGIN INCLUDED SGA3_English_current_buildable_source_20260823/source/inputs/XIII/source/tex/components/02_expose_XIII_section1_remainder_and_section2.tex
% Authority: Expo13.pdf, from “Nous supposons par la suite que S est le
% spectre d’un corps k...” through the end of §2, immediately before §3.
% PDF local pp. 3--14 (requested boundary: local p. 3 / combined p. 807
% through local p. 14); printed pp. 251--269.

We henceforth assume that \(S\) is the spectrum of a field \(k\), and
that \(G\), \(V\), and \(W\) are of finite type over \(k\). Let
\(\mathfrak n\) be the Lie algebra of \(N\). Then
\(\dim N\leq\operatorname{rank}_k\mathfrak n\), with equality if and
only if \(N\) is smooth over \(k\) (Exposé VI\footnote{Editor's note in
the French source: since this reference could not be identified, see
Theorem II.5.2.1 of M.~Demazure and P.~Gabriel,
\emph{Groupes algébriques I}, Masson (1970).}). Let \(a\in W(k)\), and
consider the subscheme \(M_a\) of \(G\) defined above; it contains
\(N\) and is isomorphic to \(\varphi^{-1}(a)\). We denote by
\(\mathfrak m_a\) its Zariski tangent space at the identity element
\(e\) of \(N\). Thus\footnote{The French source omits
\(\operatorname{rank}_k\) before \(\mathfrak m_a\) in \SGARefLink{sga3:xiii:equation:1}{(1)}; it is
restored here, as required by the surrounding text and \SGARefLink{sga3:xiii:lemma:1:1}{Lemma 1.1}.}
\SGARefTarget{sga3:xiii:equation:1}
\[
\tag{1}
  \mathfrak n\subset\mathfrak m_a,\qquad
  \dim N\leq\operatorname{rank}_k\mathfrak n
  \leq\operatorname{rank}_k\mathfrak m_a.
\]

\medskip
\noindent\SGARefTarget{sga3:xiii:lemma:1:1}\textbf{Lemma 1.1.} ---
\textit{With the preceding notation:}

\begin{enumerate}
\item[\textit{a)}] \SGARefTarget{sga3:xiii:lemma:1:1:item:a} Consider the following conditions:
  \begin{enumerate}
  \item[(i)] \SGARefTarget{sga3:xiii:lemma:1:1:item:i}
    \(\mathfrak n=\mathfrak m_a\), and \(N\) is smooth over \(k\).
  \item[(i bis)] \SGARefTarget{sga3:xiii:lemma:1:1:item:i-bis}
    \(\dim N=\operatorname{rank}_k\mathfrak m_a\).
  \item[(ii)] \SGARefTarget{sga3:xiii:lemma:1:1:item:ii}
    The morphism \(\psi:X\to V\) is unramified at
    \((\bar e,a)\).
  \item[(iii)] \SGARefTarget{sga3:xiii:lemma:1:1:item:iii}
    \(M_a\) and \(N\) coincide in a neighborhood of \(e\).
  \end{enumerate}
  Then
  \[
    \SGARefLink{sga3:xiii:lemma:1:1:item:i}{\mathrm{(i)}}\Longleftrightarrow\SGARefLink{sga3:xiii:lemma:1:1:item:i-bis}{\mathrm{(i\ bis)}}
    \Longrightarrow\SGARefLink{sga3:xiii:lemma:1:1:item:ii}{\mathrm{(ii)}}
    \Longleftrightarrow\SGARefLink{sga3:xiii:lemma:1:1:item:iii}{\mathrm{(iii)}}.
  \]
\item[\textit{b)}] \SGARefTarget{sga3:xiii:lemma:1:1:item:b}
  Suppose that \(\varphi:G\times_S W\to V\) is smooth at \((e,a)\).
  Then \(M_a\) is smooth over \(k\) at \(e\), and \(\psi\) is smooth
  at \((\bar e,a)\).
\end{enumerate}

\noindent\textit{Proof.}
\SGARefTarget{sga3:xiii:lemma:1:1:proof:item:a}
\textit{a)} The equivalence of \SGARefLink{sga3:xiii:lemma:1:1:item:i}{(i)} and \SGARefLink{sga3:xiii:lemma:1:1:item:i-bis}{(i bis)} follows immediately
from \SGARefLink{sga3:xiii:equation:1}{(1)} and from the fact noted above that \(N\) is smooth over \(k\)
if and only if
\(\dim N=\operatorname{rank}_k\mathfrak n\). Consider the inclusion
morphism \(N\to M_a\). It is well known\footnote{Editor's note in the
French source: see, for example, EGA IV\textsubscript{4},
Theorem 17.11.1(d).} that if \(N\) is smooth over \(k\) at \(e\) and
the tangent map at \(e\) is surjective, then \(N\to M_a\) is smooth
at \(e\); since it is an immersion, it is therefore an isomorphism at
\(e\). This proves that \SGARefLink{sga3:xiii:lemma:1:1:item:i}{(i)} implies \SGARefLink{sga3:xiii:lemma:1:1:item:iii}{(iii)}.

To prove the equivalence of \SGARefLink{sga3:xiii:lemma:1:1:item:ii}{(ii)} and \SGARefLink{sga3:xiii:lemma:1:1:item:iii}{(iii)}, let
\(X_a=\psi^{-1}(a)\), as above. From the isomorphisms
\[
  M_a\simeq\varphi^{-1}(a)\simeq q^{-1}(X_a)
\]
we obtain a morphism \(p_a:M_a\to X_a\) making \(M_a\) a principal
homogeneous bundle with group \(N_{X_a}\). Consider \SGARefLink{sga3:xiii:lemma:1:1:diagram:01}{the diagram}
\SGARefTarget{sga3:xiii:lemma:1:1:diagram:01}
\[
\begin{tikzcd}[column sep=large,row sep=large]
M_a \arrow[d,"p_a"'] &
N \arrow[l,"j_a"'] \arrow[d] \\
X_a \arrow[r,"j'_a"] &
S=\operatorname{Spec}(k).
\end{tikzcd}
\]
Here \(\operatorname{Spec}(k)\to X_a\)\footnote{\SGARefLink{sga3:xiii:lemma:1:1:diagram:01}{The printed diagram}
directs \(j'_a:X_a\to S\), whereas the immediately following prose
defines the point morphism \(S\to X_a\) and calls it \(j'_a\). Both
\SGARefLink{sga3:xiii:lemma:1:1:diagram:01}{the printed diagram direction} and the verbal definition are retained
here. The source also prints \(X_n\) here and \(X_{a,b}\) in the two
local rings below. They are corrected to \(X_a\) and
\(X_{a,\bar a}\), in accordance with the fiber notation and the point
\(\bar a\).} is defined by the point
\(\bar a=(\bar e,a)\) of \(X\), and \(j_a:N\to M_a\) is the canonical
immersion. To say that \(\psi\) is unramified at \((\bar e,a)\) means
that \(j'_a\) is an open immersion, or equivalently that it induces an
isomorphism
\SGARefTarget{sga3:xiii:lemma:1:1:display:residue-field-map}
\[
  S\xrightarrow{\sim}\operatorname{Spec}
  \bigl(\mathcal O_{X_a,\bar a}\bigr).
\]
Since \(p_a\) is flat, this is equivalent to saying that the morphism
obtained from \SGARefLink{sga3:xiii:lemma:1:1:display:residue-field-map}{the preceding one} by the base change
\[
  \operatorname{Spec}\bigl(\mathcal O_{M_a,e}\bigr)
  \longrightarrow
  \operatorname{Spec}\bigl(\mathcal O_{X_a,\bar a}\bigr)
\]
is an isomorphism. The resulting morphism is precisely
\[
  \operatorname{Spec}\bigl(\mathcal O_{N,e}\bigr)
  \longrightarrow
  \operatorname{Spec}\bigl(\mathcal O_{M_a,e}\bigr),
\]
which proves the equivalence of \SGARefLink{sga3:xiii:lemma:1:1:item:ii}{(ii)} and \SGARefLink{sga3:xiii:lemma:1:1:item:iii}{(iii)}.

Notice moreover that the proof shows that \SGARefLink{sga3:xiii:lemma:1:1:item:ii}{(ii)} and \SGARefLink{sga3:xiii:lemma:1:1:item:iii}{(iii)} imply \SGARefLink{sga3:xiii:lemma:1:1:item:iii-bis}{the
following condition}, apparently stronger than \SGARefLink{sga3:xiii:lemma:1:1:item:iii}{(iii)}:
\begin{enumerate}
\item[(iii bis)] \SGARefTarget{sga3:xiii:lemma:1:1:item:iii-bis}
  \(N\) is an open and closed subscheme of \(M_a\).
\end{enumerate}

\SGARefTarget{sga3:xiii:lemma:1:1:proof:item:b}
\textit{b)} The first assertion follows from the isomorphism
\(M_a\simeq\varphi^{-1}(a)\), and the second from the fact that \(q\)
is flat and \(q(e,a)=(\bar e,a)\).

\section{Density theorem and theory of the regular points of \(G\)}
\SGARefTarget{sga3:xiii:section:2}

We apply the constructions and notation of \SGARefLink{sga3:xiii:section:1}{the preceding section} in
the following situation. Let \(G\) be a smooth connected algebraic
group over \(k\); take \(V=G\), with \(G\) acting on \(V\) by inner
automorphisms; and let \(W=H\), where \(H\) is a smooth connected
algebraic subgroup of \(G\). We denote the Lie algebras of \(G\) and
\(H\) by \(\mathfrak g\) and \(\mathfrak h\), the normalizer of \(H\)
in \(G\) by \(N\), and its Lie algebra by \(\mathfrak n\). For
\(a\in G(k)\), we again denote by \(M_a\), as in \SGARefLink{sga3:xiii:section:1}{Section 1}, the
inverse of the transporter of \(a\) into \(H\). Thus, if
\(a\in H(k)\), then \(N\subset M_a\); in that case
\(\mathfrak m_a\) denotes the Zariski tangent space of \(M_a\) at the
identity element \(e\) of \(G\). Note that
\[
  \mathfrak h\subset\mathfrak n\subset\mathfrak m_a
  \subset\mathfrak g
  \qquad\text{for }a\in H(k).
\]
We shall use \SGARefLink{sga3:xiii:lemma:2:0}{the following lemma}.

\medskip
\noindent\SGARefTarget{sga3:xiii:lemma:2:0}\textbf{Lemma 2.0.} ---
\textit{The equality \(H=N^0\) holds if and only if}
\SGARefTarget{sga3:xiii:lemma:2:0:display:invariants-zero}
\[
  (\mathfrak g/\mathfrak h)^H=0,
\]
\textit{where the left-hand side denotes the subspace of invariants
for the action of \(H\) induced by the adjoint representation. When
\SGARefLink{sga3:xiii:lemma:2:0:display:invariants-zero}{this condition} holds, \(N\) is smooth and \(\dim X=\dim G\). In all
cases, \(\dim X\leq\dim G\), and equality holds if and only if \(H\)
has finite index in \(N\).}

Indeed, by \SGARefLink{sga3:ii:theorem:5-2-3:item:i}{II~5.2.3(i)}, \(\mathfrak n\) is the inverse image of
\((\mathfrak g/\mathfrak h)^H\) under
\(\mathfrak g\to\mathfrak g/\mathfrak h\). Hence
\((\mathfrak g/\mathfrak h)^H=0\) is equivalent to
\(\mathfrak h=\mathfrak n\), which in turn is equivalent to
\(H=N^0\), since \(H\) is a smooth connected algebraic subgroup of
the algebraic group \(N\) (cf. VI.2). This also shows that \(N\) is
smooth. Furthermore,
\[
  \dim X=(\dim G-\dim N)+\dim H
        =\dim G-(\dim N-\dim H),
\]
and therefore
\[
  \dim X\leq\dim G,
\]
with equality if and only if \(\dim H=\dim N\), that is, if and only
if \(H\) has finite index in \(N\). This holds in particular when
\(H=N^0\), and proves the lemma.

\medskip
\noindent\SGARefTarget{sga3:xiii:theorem:2:1}\textbf{Theorem 2.1.} ---
\textit{Let \(G\) be a smooth connected algebraic group over an
algebraically closed field \(k\), let \(H\) be a smooth connected
algebraic subgroup, and let \(N\) be its normalizer. Write
\(\mathfrak g,\mathfrak h,\mathfrak n\) for their Lie algebras, let
\(X=G\times^N H\) be the scheme introduced in \SGARefLink{sga3:xiii:section:1}{Section 1} (fibered over
\(G/N\) with typical fiber \(H\)), and let \(\psi:X\to G\) be the
canonical morphism. Its image is also the image of
\(\varphi:G\times H\to G\), where}
\[
  \varphi(g,h)=\operatorname{int}(g)h=ghg^{-1}.
\]
\textit{The following conditions are equivalent:}

\begin{enumerate}
\item[(i)] \SGARefTarget{sga3:xiii:theorem:2:1:item:i}
  \(H\) contains a Cartan subgroup \(C\) of \(G\) (XII~1).
\item[(i bis)] \SGARefTarget{sga3:xiii:theorem:2:1:item:i-bis}
  \(H\) has the same reductive rank and the same nilpotent rank
  (XII~1) as \(G\).
\item[(ii)] \SGARefTarget{sga3:xiii:theorem:2:1:item:ii}
  \(H\) contains a maximal torus \(T\) of \(G\), and
  \((\mathfrak g/\mathfrak h)^T=0\).
\item[(iii)] \SGARefTarget{sga3:xiii:theorem:2:1:item:iii}
  The set of conjugates of \(H\) containing a given maximal torus is
  finite and nonempty, and \(H\) has finite index in its normalizer.
\item[(iv)] \SGARefTarget{sga3:xiii:theorem:2:1:item:iv}
  There exists \(a\in H(k)\) that belongs to only finitely many
  conjugates of \(H\) (or merely such that \(\psi^{-1}(a)\) has an
  isolated point), and \(H\) has finite index in its normalizer.
\item[(iv bis)] \SGARefTarget{sga3:xiii:theorem:2:1:item:iv-bis}
  The morphism \(\psi:X\to G\) is generically quasi-finite (that is,
  there is a dense open subset of \(X\) on which \(\psi\) is
  quasi-finite), and \(H\) has finite index in its normalizer.
\item[(v)] \SGARefTarget{sga3:xiii:theorem:2:1:item:v}
  There is a dense open subset \(U\) of \(G\) such that, for every
  \(x\in U(k)\), the set of conjugates of \(H\) containing \(x\) is
  finite and nonempty; equivalently, \(\psi:X\to G\) is dominant and
  generically quasi-finite.
\item[(vi)] \SGARefTarget{sga3:xiii:theorem:2:1:item:vi}
  There is a dense open subset \(U\) of \(G\) such that every
  \(x\in U(k)\) is contained in a conjugate of \(H\); equivalently,
  \(\psi:X\to G\) is dominant.
\item[(vii)] \SGARefTarget{sga3:xiii:theorem:2:1:item:vii}
  There exists \(a\in H(k)\) such that the fixed subspace of
  \(\mathfrak g/\mathfrak h\) under
  \(\operatorname{ad}_{\mathfrak g/\mathfrak h}(a)\) is zero.
\end{enumerate}

\textit{Moreover, these conditions imply that \(H\) is its own
connected normalizer: \(N\) is smooth and
\(\dim H=\dim N\). They also imply that \(\psi:X\to G\) is
generically étale.}

\noindent\textit{Proof.}
By \SGARefLink{sga3:xiii:lemma:2:0}{2.0}, \(\dim X\leq\dim G\), with equality if and only if
\(\dim H=\dim N\), or equivalently if and only if \(H\) has finite
index in \(N\). The inequality \(\dim X\leq\dim G\) shows that
\(\psi\) is dominant if and only if it is dominant and generically
quasi-finite, and also if and only if it is generically quasi-finite
and \(\dim X=\dim G\). Since the last equality is equivalent, by
\SGARefLink{sga3:xiii:lemma:2:0}{2.0}, to \(H\) having finite index in \(N\), conditions \SGARefLink{sga3:xiii:theorem:2:1:item:vi}{(vi)}, \SGARefLink{sga3:xiii:theorem:2:1:item:v}{(v)}, and
\SGARefLink{sga3:xiii:theorem:2:1:item:iv-bis}{(iv bis)} are equivalent. The equivalence of \SGARefLink{sga3:xiii:theorem:2:1:item:iv}{(iv)} and \SGARefLink{sga3:xiii:theorem:2:1:item:iv-bis}{(iv bis)} is
immediate.

The equivalence of \SGARefLink{sga3:xiii:theorem:2:1:item:i}{(i)} and \SGARefLink{sga3:xiii:theorem:2:1:item:i-bis}{(i bis)} follows directly from the
definitions and is left to the reader. If \(H\) contains a Cartan
subgroup \(C\) of \(G\), it contains the maximal torus \(T\) of
\(C\), which is a maximal torus of \(G\). Since \(C\) is the
centralizer of \(T\), its Lie algebra \(\mathfrak c\) is
\[
  \mathfrak c=\mathfrak g^T
\]
(\SGARefLink{sga3:ii:theorem:5-2-3:item:ii}{II~5.2.3(ii)}). Hence
\(\mathfrak g^T=\mathfrak c\subset\mathfrak h\), and \SGARefLink{sga3:i:proposition:4-7-3}{I~4.7.3} gives
the relation\footnote{The French source has two evident typographical
slips in this sentence, writing \(\mathfrak n\) for \(\mathfrak c\)
and reversing the inclusion involving \(\mathfrak h\) and
\(\mathfrak g^T\). The mathematically required relation is restored
here.}
\SGARefTarget{sga3:xiii:theorem:2:1:equation:x}
\[
\tag{x}
  (\mathfrak g/\mathfrak h)^T=0.
\]
Conversely, suppose that \(H\) contains the maximal torus \(T\) and
that \SGARefLink{sga3:xiii:theorem:2:1:equation:x}{(x)} holds, so that \(\mathfrak h\supset\mathfrak c\). We claim
that \(H\) contains the centralizer \(C\) of \(T\), which will prove
the equivalence of \SGARefLink{sga3:xiii:theorem:2:1:item:i}{(i)} and \SGARefLink{sga3:xiii:theorem:2:1:item:ii}{(ii)}. This follows from \SGARefLink{sga3:xiii:lemma:2:1:1}{the next lemma}.

\medskip
\noindent\SGARefTarget{sga3:xiii:lemma:2:1:1}\textbf{Lemma 2.1.1.} ---
\textit{Let \(G\) be a smooth algebraic group over a field \(k\), let
\(T\) be a subgroup of multiplicative type of \(G\), let \(C\) be its
connected centralizer (equal to the centralizer of \(T\) when \(G\)
is connected and \(T\) is a torus, by \SGARefLink{sga3:xii:6.6:item:b}{XII~6.6(b)}), and let \(H\) be
a smooth subgroup of \(G\) containing \(T\). Then \(H\) contains
\(C\) if and only if its Lie algebra \(\mathfrak h\) contains the Lie
algebra \(\mathfrak c\) of \(C\).}

Indeed, \(\operatorname{Centr}_G(T)\) is smooth over \(k\) by \SGARefLink{sga3:xi:corollary:2:4}{XI~2.4},
and hence \(C\) is smooth over \(k\). Likewise,
\(\operatorname{Centr}_H(T)\) is smooth, and
\[
  \operatorname{Centr}_H(T)=\operatorname{Centr}_G(T)\cap H
\]
has Lie algebra \(\mathfrak h\cap\mathfrak c\). The hypothesis
therefore implies that this smooth subgroup has the same Lie algebra
as the smooth group \(\operatorname{Centr}_G(T)\), so it contains the
connected component \(C\) of the latter. Thus \(H\) contains \(C\).
\hfill\(\square\)

We now prove the equivalence of \SGARefLink{sga3:xiii:theorem:2:1:item:i}{(i)} and \SGARefLink{sga3:xiii:theorem:2:1:item:iii}{(iii)}. If \(H\) contains the
maximal torus \(T\) of \(G\), we must show that \(H\supset C\)
(equivalently, that \SGARefLink{sga3:xiii:theorem:2:1:equation:x}{(x)} holds) if and only if \(H\) has finite index
in its normalizer and the set of conjugates of \(H\) containing \(T\)
is finite. If \(H\supset C\), then \SGARefLink{sga3:xiii:theorem:2:1:equation:x}{(x)} implies, a fortiori,
\SGARefTarget{sga3:xiii:theorem:2:1:equation:xx}
\[
\tag{xx}
  (\mathfrak g/\mathfrak h)^H=0.
\]
The Lie algebra \(\mathfrak n\) is the inverse image of the left-hand
side under the canonical homomorphism
\(\mathfrak g\to\mathfrak g/\mathfrak h\) (II~5.2.3(i)).
Consequently, by \SGARefLink{sga3:xiii:lemma:2:0}{2.0}, \SGARefLink{sga3:xiii:theorem:2:1:equation:xx}{(xx)} says that \(H=N^0\); in particular, \(H\)
has finite index in its normalizer. Consider \SGARefLink{sga3:xiii:theorem:2:1:diagram:01}{the diagram of
subgroups}
\SGARefTarget{sga3:xiii:theorem:2:1:diagram:01}
\[
\begin{tikzcd}[column sep=large,row sep=large]
T \arrow[r] &
N(T)\cap H \arrow[r] \arrow[d] &
H \arrow[d] \\
{} &
N(T)\cap N(H) \arrow[r] \arrow[d] &
N(H) \\
{} &
N(T). &
\end{tikzcd}
\]
By the conjugacy theorem for maximal tori in \(H\) (\SGARefLink{sga3:xii:6.6:item:a}{XII~6.6(a)}),
every conjugate of \(H\) containing \(T\) is obtained from \(H\) by
conjugation by an element of \(N(T)(k)\). Thus the set of such
conjugates is in bijection with
\[
  \bigl(N(T)/(N(T)\cap N(H))\bigr)(k).
\]
Since \(H\supset C\), we have \(N(T)\cap H\supset C\), so the
preceding set is a quotient of \((N(T)/C)(k)\), which is finite.
This proves \SGARefLink{sga3:xiii:theorem:2:1:item:i}{(i)} \(\Rightarrow\) \SGARefLink{sga3:xiii:theorem:2:1:item:iii}{(iii)}.

Conversely, suppose \SGARefLink{sga3:xiii:theorem:2:1:item:iii}{(iii)}, so that
\(N(T)/(N(T)\cap N(H))\) and \(N(H)/H\) are finite. The conjugacy
theorem in \(H\) shows that the homomorphism
\[
  (N(T)\cap N(H))/(N(T)\cap H)\longrightarrow N(H)/H
\]
induced by \SGARefLink{sga3:xiii:theorem:2:1:diagram:01}{the preceding diagram} is bijective on \(k\)-valued points
(indeed, it is an isomorphism). Since the latter quotient is finite,
so is the former. Hence \(N(T)\cap H\) has finite index in \(N(T)\)
and therefore contains \(C=N(T)^0\). Thus \(H\supset C\), and \SGARefLink{sga3:xiii:theorem:2:1:item:i}{(i)},
\SGARefLink{sga3:xiii:theorem:2:1:item:i-bis}{(i bis)}, \SGARefLink{sga3:xiii:theorem:2:1:item:ii}{(ii)}, and \SGARefLink{sga3:xiii:theorem:2:1:item:iii}{(iii)} are equivalent.

We next prove \SGARefLink{sga3:xiii:theorem:2:1:item:ii}{(ii)} \(\Rightarrow\) \SGARefLink{sga3:xiii:theorem:2:1:item:vii}{(vii)}. Both conditions are
invariant under an extension \(k\to k'\) with \(k'\) algebraically
closed. We may therefore assume that \(k\) has infinite
transcendence degree over its prime subfield. There is then an
element \(a\in T(k)\) such that the subgroup of \(T(k)\) generated by
\(a\) is Zariski dense in \(T\). It follows that
\[
  (\mathfrak g/\mathfrak h)^T
  =(\mathfrak g/\mathfrak h)^{\operatorname{ad}(a)}.
\]
The left-hand side is zero by hypothesis, which proves \SGARefLink{sga3:xiii:theorem:2:1:item:vii}{(vii)}.

The implication \SGARefLink{sga3:xiii:theorem:2:1:item:vii}{(vii)} \(\Rightarrow\) \SGARefLink{sga3:xiii:theorem:2:1:item:vi}{(vi)} is contained in \SGARefLink{sga3:xiii:corollary:2:2}{the
following result}, which sharpens \SGARefLink{sga3:xiii:theorem:2:1}{2.1}.

\medskip
\noindent\SGARefTarget{sga3:xiii:corollary:2:2}\textbf{Corollary 2.2.} ---
\textit{Let \(G\) be a smooth algebraic group over a field \(k\), let
\(H\) be a smooth algebraic subgroup, let \(N\) be its normalizer in
\(G\), and let}
\[
  \varphi:G\times H\longrightarrow G,\qquad
  \varphi(g,h)=\operatorname{ad}(g)h=ghg^{-1}.
\]
\textit{Let \(\psi:X=G\times^N H\to G\) be the morphism obtained from
\(\varphi\) by passage to the quotient (cf. \SGARefLink{sga3:xiii:section:1}{Section 1}), and let
\(a\in H(k)\). The following conditions are equivalent:}

\begin{enumerate}
\item[(i)] \SGARefTarget{sga3:xiii:corollary:2:2:item:i}
  \(\varphi\) is smooth at \((e,a)\).
\item[(ii)] \SGARefTarget{sga3:xiii:corollary:2:2:item:ii}
  \(\psi\) is étale at \((\bar e,a)\), and \(N\) is smooth over \(k\).
\item[(iii)] \SGARefTarget{sga3:xiii:corollary:2:2:item:iii}
  \((\mathfrak g/\mathfrak h)^{\operatorname{ad}(a)}=0\), where
  \(\mathfrak g\) and \(\mathfrak h\) are the Lie algebras of \(G\)
  and \(H\).
\end{enumerate}

\textit{These conditions imply \(H^0=N^0\).}

The smoothness at a \(k\)-rational point of \(\varphi\), a morphism
between smooth \(k\)-preschemes, is equivalent to the surjectivity of
the tangent map at that point. Under the usual identifications of
the tangent spaces of \(G\) and \(H\) with their Lie algebras, a
direct calculation gives
\SGARefTarget{sga3:xiii:corollary:2:2:display:tangent-map}
\[
  d\varphi(\xi,\eta)
  =\bigl(\operatorname{id}-\operatorname{ad}(a)\bigr)\xi+\eta,
\]
as a map from \(\mathfrak g\times\mathfrak h\) to \(\mathfrak g\).
Its surjectivity is therefore equivalent to the surjectivity of
\(\operatorname{id}-\operatorname{ad}(a)\) on
\(\mathfrak g/\mathfrak h\), which is precisely \SGARefLink{sga3:xiii:corollary:2:2:item:iii}{(iii)}.

Condition \SGARefLink{sga3:xiii:corollary:2:2:item:iii}{(iii)} implies, a fortiori,
\((\mathfrak g/\mathfrak h)^H=0\), and hence \(H^0=N^0\) by \SGARefLink{sga3:xiii:lemma:2:0}{2.0}
(the connectedness hypothesis made at the beginning of the section
is not needed there). It follows that \(N\) is smooth and that
\(\dim H=\dim N\), whence \(\dim X=\dim G\). Since
\(q:G\times H\to X\) is flat and
\(\psi\circ q=\varphi\)\footnote{The source reverses this composition,
printing \(\psi=\varphi\circ q\). The displayed morphisms and the
definition of \(\varphi\) require \(\varphi=\psi\circ q\).} is smooth
at \((e,a)\), the morphism \(\psi\)
is smooth at \(q(e,a)=(\bar e,a)\), and hence is étale there by
dimension. We have proved
\(\SGARefLink{sga3:xiii:corollary:2:2:item:i}{\mathrm{(i)}}\Longleftrightarrow\SGARefLink{sga3:xiii:corollary:2:2:item:iii}{\mathrm{(iii)}}
\Longrightarrow\SGARefLink{sga3:xiii:corollary:2:2:item:ii}{\mathrm{(ii)}}\). Conversely, \SGARefLink{sga3:xiii:corollary:2:2:item:ii}{(ii)} implies \SGARefLink{sga3:xiii:corollary:2:2:item:i}{(i)},
because the smoothness of \(N\) implies that of \(q\).
\hfill\(\square\)

In particular, under condition \SGARefLink{sga3:xiii:theorem:2:1:item:vii}{(vii)} of \SGARefLink{sga3:xiii:theorem:2:1}{2.1} the morphism
\(\varphi\) is smooth at \((e,a)\), so its image contains a nonempty
open subset of \(G\). Thus \(\psi\) is dominant, proving
\SGARefLink{sga3:xiii:theorem:2:1:item:vii}{(vii)} \(\Rightarrow\) \SGARefLink{sga3:xiii:theorem:2:1:item:vi}{(vi)}.

It remains to prove \SGARefLink{sga3:xiii:theorem:2:1:item:vi}{(vi)} \(\Rightarrow\) \SGARefLink{sga3:xiii:theorem:2:1:item:i}{(i)}. Suppose first that
\(G\) is affine. Let \(U\) be a nonempty open subset of \(G\) such
that every \(x\in U(k)\) is contained in a conjugate of \(H\), and
let \(C\) be a Cartan subgroup of \(G\). Applying the already proved
implication \SGARefLink{sga3:xiii:theorem:2:1:item:i}{(i)} \(\Rightarrow\) \SGARefLink{sga3:xiii:theorem:2:1:item:iv}{(iv)} with \(C\) in place of \(H\)
(Borel's ``density theorem''), we find a conjugate of \(C\) meeting
\(U\). We may therefore assume that \(U\cap C\neq\varnothing\);
equivalently, there is a nonempty open subset \(V\) of \(C\) such
that every \(x\in V(k)\) belongs to a conjugate of \(H\).

Write
\SGARefTarget{sga3:xiii:theorem:2:1:display:cartan-decomposition}
\[
  C=T\cdot C_u,
\]
where \(T\) is the maximal torus of \(C\), hence a maximal torus of
\(G\), and \(C_u\) is the unipotent part of \(C\); \(T\) lies in the
center of \(C\) (\emph{BIBLE} 6, Theorem 2). We may again assume that
\(k\) has infinite transcendence degree over its prime subfield.
Choose \(t\in T(k)\) in the projection of \(V\) onto \(T\), a
nonempty open subset of \(T\), so that
\(tC_u\cap V\neq\varnothing\), and choose it so that it generates a
Zariski-dense subgroup of \(T\). Every algebraic subgroup of \(G\)
containing a product \(tu\), with \(t\in T(k)\) and
\(u\in C_u(k)\), contains both factors (\emph{BIBLE} 4,
Theorem 3). Taking \(tu\in V(k)\), we conclude that a conjugate of
\(H\) contains \(t\), and therefore contains \(T\). After replacing
\(H\) by that conjugate, we may assume
\[
  T\subset H.
\]

Let \(W\) be the inverse image in \(C_u\) of \(V\) under
\(u\mapsto tu\). For every \(x\in(T\cdot W)(k)\), there is then a conjugate
of \(H\) containing both \(T\) and \(x\). As observed above, every
such conjugate is of the form \(\operatorname{int}(g)H\) for some
\(g\in N(T)(k)\). Consider
\SGARefTarget{sga3:xiii:theorem:2:1:display:f}
\[
  f:N(T)\times H\longrightarrow G,\qquad
  f(g,h)=\operatorname{int}(g)h=ghg^{-1}.
\]
The image of \(f\) contains \(TW\). Since \(C\) has finite index in
\(N(T)\), the latter is a finite union of translates \(Cg_i\), with
\(g_i\in N(T)(k)\). It follows that some dense open subset \(V'\) of
\(C=T C_u\) is contained in the image under \(f\) of
\((Cg_i)\times H\). Replacing \(H\) by
\(\operatorname{int}(g_i)H\), we may assume \(g_i=e\), so that
\[
  f(C\times H)\supset V'.
\]
Thus for every \(u\in V'(k)\) there are \(v\in C(k)\) and
\(h\in H(k)\) such that
\[
  v^{-1}hv=u,\qquad\text{hence}\qquad vuv^{-1}\in H(k).
\]
Set
\[
  C'=C\cap H=\operatorname{Centr}_H(T).
\]
Since \(vuv^{-1}\) also belongs to \(C\), it lies in \(C'\), and
hence \(u\in\operatorname{int}(v^{-1})C'\). Therefore the union of
the conjugates of \(C'\) in \(C\), by elements of \(C(k)\), is
dense. As in the preceding argument for the pair \((G,H)\), now
applied to \((C,C')\), this implies that \(C'\) has finite index in
its normalizer in \(C\). By the lemma of \emph{BIBLE} 7, \SGARefLink{sga3:xiii:section:1}{Section 1},
we obtain \(C'=C\), and consequently \(C\subset H\).\footnote{The
source prints \(C=H\); only the inclusion \(C\subset H\) follows and
is used.} This proves
\SGARefLink{sga3:xiii:theorem:2:1:item:vi}{(vi)} \(\Rightarrow\) \SGARefLink{sga3:xiii:theorem:2:1:item:i}{(i)} when \(G\) is affine.

In the general case we argue by induction on \(n=\dim G\), the
assertion being trivial for \(n=0\). Let \(Z\) be the center of
\(G\). There are two cases.

\begin{enumerate}
\item[\(1^\circ\)] \SGARefTarget{sga3:xiii:theorem:2:1:proof:case:1} If \(\dim(Z\cap H)>0\), set
  \(G'=G/(Z\cap H)\). Then \(\dim G'<n\), and condition \SGARefLink{sga3:xiii:theorem:2:1:item:vi}{(vi)} for
  \(H\) implies \SGARefLink{sga3:xiii:theorem:2:1:item:vi}{the same condition} for its image \(H'\) in \(G'\).
  By induction, \(H'\) contains a Cartan subgroup \(C'\) of \(G'\).
  Hence \(H\) contains the inverse image \(C\) of \(C'\), which is a
  Cartan subgroup by \SGARefLink{sga3:xii:theorem:6-6:item:e}{XII~6.6(e)}.

\item[\(2^\circ\)] \SGARefTarget{sga3:xiii:theorem:2:1:proof:case:2} If \(\dim(Z\cap H)=0\), the canonical morphism
  \(H\to G/Z\) is finite. Since \(G/Z\) is affine by \SGARefLink{sga3:xii:6.1}{XII~6.1}, \(H\)
  is affine. Every homomorphism from \(H\) to an abelian variety is
  therefore zero (indeed, every morphism of preschemes from \(H\) to
  an abelian variety is constant). This follows from the fact that a
  smooth connected affine algebraic group over an algebraically
  closed field is a rational variety, or simply from the fact that
  it is the union of its Borel subgroups (\emph{BIBLE} 6,
  Theorem 5(b)); the structure theorems \emph{BIBLE} 6.2 and 6.3
  show readily that a smooth connected solvable affine group is a
  rational variety.

  Apply Chevalley's structure theorem to \(G\): the group \(G\) is
  an extension of an abelian variety \(A\) by a smooth affine group.
  Since \(H\) is affine, its image in \(A\) is zero. On the other
  hand, that image must be all of \(A\), because the union of its
  conjugates in \(A\) is dense. Hence \(A=0\), so \(G\) is affine,
  reducing us to the case already proved.
\end{enumerate}

This completes the proof of \SGARefLink{sga3:xiii:theorem:2:1}{2.1}.

\medskip
\noindent\SGARefTarget{sga3:xiii:corollary:2:3}\textbf{Corollary 2.3.} ---
\textit{Assume that the equivalent conditions of \SGARefLink{sga3:xiii:theorem:2:1}{2.1} hold.}

\begin{enumerate}
\item[\textit{a)}] \SGARefTarget{sga3:xiii:corollary:2:3:item:a}
  Let \(k(X)\) and \(k(G)\) be the fields of rational functions of
  \(X\) and \(G\), respectively. Then \(k(X)\) is a finite separable
  extension of \(k(G)\). Denote its degree by \(d\).

\item[\textit{b)}] \SGARefTarget{sga3:xiii:corollary:2:3:item:b}
  Let \(T\) be a maximal torus of \(G\) contained in \(H\), which
  exists by \SGARefLink{sga3:xiii:theorem:2:1:item:ii}{2.1(ii)}, and let \(C\) be the corresponding Cartan
  subgroup of \(G\). Then \(C\subset H\). Moreover,
  \(\operatorname{Norm}_G(T)\) is a smooth subgroup of \(G\), and
  \[
    \operatorname{Norm}_G(T)\cap\operatorname{Norm}_G(H)
    =
    \operatorname{Norm}_{\operatorname{Norm}_G(H)}(T)
  \]
  is a smooth subgroup of it, of finite index equal to the degree
  \(d\) defined in \SGARefLink{sga3:xiii:corollary:2:3:item:a}{(a)}. The number of conjugates of \(H\) containing
  a given maximal torus, or a given Cartan subgroup, is \(d\).

\item[\textit{c)}] \SGARefTarget{sga3:xiii:corollary:2:3:item:c}
  Let \(U\) be the largest open subset of \(G\) such that
  \(\psi:X\to G\) induces a finite étale morphism
  \[
    \psi^{-1}(U)\longrightarrow U.
  \]
  Then \(U\) is dense. For \(g\in G(k)\), the following are
  equivalent:
  \begin{enumerate}
  \item \SGARefTarget{sga3:xiii:corollary:2:3:item:c:condition:i} \(g\in U(k)\);
  \item \SGARefTarget{sga3:xiii:corollary:2:3:item:c:condition:ii} exactly \(d\) conjugates of \(H\) contain \(g\);
  \item \SGARefTarget{sga3:xiii:corollary:2:3:item:c:condition:iii} there are at least \(d\) distinct conjugates \(H_i\) of \(H\)
    containing \(g\), and for every \(i\),
    \[
      (\mathfrak g/\mathfrak h_i)^{\operatorname{ad}(g)}=0,
      \qquad \mathfrak h_i=\operatorname{Lie}(H_i).
    \]
  \end{enumerate}
\end{enumerate}

Assertion \SGARefLink{sga3:xiii:corollary:2:3:item:a}{(a)} follows from the generic étaleness of \(\psi\), stated
at the end of \SGARefLink{sga3:xiii:theorem:2:1}{2.1}. This also shows that the open set \(U\) in \SGARefLink{sga3:xiii:corollary:2:3:item:c}{(c)} is
nonempty, hence dense, and gives the two stated characterizations of
its \(k\)-points: here one uses that \(\psi\) is separated, that
\(X\) is integral, that \(G\) is integral and normal (SGA 1,
I~10.11), and that
\((\mathfrak g/\mathfrak h_i)^{\operatorname{ad}(g)}=0\) means that
\(\psi\) is étale at the point \(x_i\in\psi^{-1}(g)\) corresponding
to \(H_i\).

If \(H\) contains the maximal torus \(T\) of \(G\), the centralizers
of \(T\) in \(H\) and \(G\) have the same dimension and are smooth
and connected (\SGARefLink{sga3:xii:theorem:6-6:item:b}{XII~6.6(b)}); they are therefore equal, proving that
\(C\subset H\).\footnote{The source prints the tautological
\(C\subset G\); equality of the two centralizers gives
\(C\subset H\).} The normalizer of \(T\) in any smooth group
containing it is smooth (\SGARefLink{sga3:xi:corollary:2:4:bis}{XI~2.4 bis}), so
\(\operatorname{Norm}_G(T)\) and
\(\operatorname{Norm}_{\operatorname{Norm}_G(H)}(T)\) are smooth.
Recall also that \(N=\operatorname{Norm}_G(H)\) is smooth by the
last assertion of \SGARefLink{sga3:xiii:theorem:2:1}{2.1}. Furthermore,
\(\operatorname{Norm}_N(T)\) contains \(C\), which has finite index
in \(\operatorname{Norm}_G(T)\); hence
\(\operatorname{Norm}_N(T)\) itself has finite index in
\(\operatorname{Norm}_G(T)\).

The conjugacy theorem for maximal tori in \(H\) shows that this index
is the number of conjugates of \(H\) containing \(T\), equivalently
the number containing \(C\). The union of the conjugates of \(C\) in
\(G\) is dense, by \SGARefLink{sga3:xiii:theorem:2:1:item:i}{2.1(i)} \(\Rightarrow\) \SGARefLink{sga3:xiii:theorem:2:1:item:vi}{(vi)} applied to \(C\) in
place of \(H\), and \(U\) is stable under inner automorphisms. Hence
\(C\cap U\neq\varnothing\). Arguing as in the proof of
\SGARefLink{sga3:xiii:theorem:2:1:item:vi}{2.1(vi)} \(\Rightarrow\) \SGARefLink{sga3:xiii:theorem:2:1:item:ii}{(ii)}, and making a harmless extension of the
base field if necessary, one finds
\(g\in(C\cap U)(k)\) such that every conjugate of \(H\) containing
\(g\) contains \(T\), and hence \(C\). Thus the conjugates of \(H\)
containing \(C\) are precisely those containing \(g\); since
\(g\in U(k)\), their number is \(d\). This proves \SGARefLink{sga3:xiii:corollary:2:3:item:b}{(b)}.

In fact, the argument shows that the set of conjugates of \(H\)
containing \(T\) is a homogeneous set under the group of rational
points of
\[
  W_G(T)=
  \operatorname{Norm}_G(T)/\operatorname{Centr}_G(T).
\]
In particular,
\[
  d\leq\text{the order of the Weyl group of \(G\)}.
\]

\medskip
\noindent\SGARefTarget{sga3:xiii:corollary:2:4}\textbf{Corollary 2.4.} ---
\textit{With the notation of \SGARefLink{sga3:xiii:theorem:2:1}{2.1}, the following conditions are
equivalent:}

\begin{enumerate}
\item[(i)] \SGARefTarget{sga3:xiii:corollary:2:4:item:i}
  \(\psi:X\to G\) is birational.
\item[(ii)] \SGARefTarget{sga3:xiii:corollary:2:4:item:ii}
  Exactly one conjugate of \(H\) contains a given Cartan subgroup of
  \(G\).
\item[(iii)] \SGARefTarget{sga3:xiii:corollary:2:4:item:iii}
  \(H\) contains a Cartan subgroup \(C\) of \(G\), and
  \(\operatorname{Norm}_G(H)\supset\operatorname{Norm}_G(C)\).
\item[(iv)] \SGARefTarget{sga3:xiii:corollary:2:4:item:iv}
  There is a nonempty open subset \(V\) of \(G\) such that every
  \(g\in V(k)\) belongs to exactly one conjugate of \(H\).
\end{enumerate}

This follows immediately from \SGARefLink{sga3:xiii:theorem:2:1}{2.1} and \SGARefLink{sga3:xiii:corollary:2:3}{2.3}.

\medskip
\noindent\SGARefTarget{sga3:xiii:corollary:2:5}\textbf{Corollary 2.5.} ---
\textit{Assume that the conditions of \SGARefLink{sga3:xiii:corollary:2:4}{2.4} hold, and let
\(g\in G(k)\). The following conditions are equivalent:}

\begin{enumerate}
\item[(i)] \SGARefTarget{sga3:xiii:corollary:2:5:item:i}
  \(g\in U(k)\), where \(U\) is defined in \SGARefLink{sga3:xiii:corollary:2:3:item:c}{2.3(c)}; that is, \(g\) is
  contained in exactly one conjugate of \(H\).
\item[(ii)] \SGARefTarget{sga3:xiii:corollary:2:5:item:ii}
  The set of conjugates of \(H\) containing \(g\) is finite and
  nonempty.
\item[(iii)] \SGARefTarget{sga3:xiii:corollary:2:5:item:iii}
  The scheme \(\psi^{-1}(g)\), the ``scheme of conjugates of \(H\)
  containing \(g\),'' has an isolated point.
\item[(iv)] \SGARefTarget{sga3:xiii:corollary:2:5:item:iv}
  There is a conjugate \(H'\) of \(H\) containing \(g\) such that
  \[
    (\mathfrak g/\mathfrak h')^{\operatorname{ad}(g)}=0,
    \qquad \mathfrak h'=\operatorname{Lie}(H').
  \]
\end{enumerate}

\textit{Finally, \(U\) is also the largest open subset of \(G\) on
which \(\psi\) induces an isomorphism}
\[
  \psi^{-1}(U)\xrightarrow{\sim}U.
\]

The equivalence of \SGARefLink{sga3:xiii:corollary:2:5:item:i}{(i)}, \SGARefLink{sga3:xiii:corollary:2:5:item:ii}{(ii)}, and \SGARefLink{sga3:xiii:corollary:2:5:item:iii}{(iii)}, together with the last
assertion, follows from the ``Main Theorem''\footnote{Editor's note in
the French source: Zariski's Main Theorem.} applied to the
birational morphism \(\psi:X\to G\), since \(G\) is normal. Their
equivalence with \SGARefLink{sga3:xiii:corollary:2:5:item:iv}{(iv)} follows immediately from \SGARefLink{sga3:xiii:corollary:2:2}{Corollary~2.2}, which
characterizes the points of \(X\) at which \(\psi\) is
étale.\footnote{The source cites the final assertion of \SGARefLink{sga3:xiii:theorem:2:1}{2.1} here;
the pointwise étaleness criterion is \SGARefLink{sga3:xiii:corollary:2:2}{Corollary 2.2}.}

\medskip
\noindent\SGARefTarget{sga3:xiii:theorem:2:6}\textbf{Theorem 2.6.} ---
\textit{Let \(G\) be a smooth connected algebraic group over an
algebraically closed field \(k\), let \(C\) be a Cartan subgroup
associated with a maximal torus \(T\), and put}
\[
  N=\operatorname{Norm}_G(C)=\operatorname{Norm}_G(T)
  \qquad\text{(cf. XII~8.4).}
\]
\textit{Let \(X=G\times^N C\), where \(N\) acts on the left factor
\(G\) by right translations and on the right factor \(C\) by inner
automorphisms, and let \(\psi:X\to G\) be the canonical morphism.}

\begin{enumerate}
\item[\textit{a)}] \SGARefTarget{sga3:xiii:theorem:2:6:item:a}
  \(\psi\) is birational.

\item[\textit{b)}] \SGARefTarget{sga3:xiii:theorem:2:6:item:b}
  Let \(U\) be the largest open subset of \(G\) such that \(\psi\)
  induces an isomorphism
  \[
    \psi^{-1}(U)\xrightarrow{\sim}U
  \]
  (cf. \SGARefLink{sga3:xiii:corollary:2:5}{2.5}), and let
  \SGARefTarget{sga3:xiii:theorem:2:6:display:rho}
  \[
    \rho=\rho_\nu(G)=\dim C
  \]
  be the nilpotent rank of \(G\). For every \(g\in G(k)\), the
  multiplicity of the eigenvalue \(1\) in the action of
  \(\operatorname{ad}(g)\) on \(\mathfrak g\) is at least \(\rho\);
  it equals \(\rho\) if and only if \(g\in U(k)\).
\end{enumerate}

\noindent\textit{Proof.}
\SGARefLink{sga3:xiii:theorem:2:1:item:i}{Condition 2.1(i)} holds, so \SGARefLink{sga3:xiii:corollary:2:4:item:iii}{2.4(iii)} \(\Rightarrow\) \SGARefLink{sga3:xiii:corollary:2:4:item:i}{(i)} proves \SGARefLink{sga3:xiii:theorem:2:6:item:a}{(a)}.
In \emph{BIBLE} 7, when \(G\) is affine, the points of \(U(k)\) are
called the \emph{regular points} of \(G(k)\). We shall use the same
terminology and call \(U\) the \emph{open subset of regular points}
of \(G\). Notice that the proof in \emph{BIBLE} that this set is
itself open is incorrect; \SGARefLink{sga3:xiii:expose}{the present Exposé} establishes its
openness under more general hypotheses.

To prove \SGARefLink{sga3:xiii:theorem:2:6:item:b}{(b)}, for each \(g\in G(k)\) introduce \SGARefLink{sga3:xiii:theorem:2:6:display:characteristic-polynomial}{the characteristic
polynomial}
\SGARefTarget{sga3:xiii:theorem:2:6:display:characteristic-polynomial}
\[
  P(\operatorname{ad}(g),t)
  =t^n+c_1(g)t^{n-1}+\cdots+c_n(g).
\]
Replacing \(k\) by an arbitrary \(k\)-algebra shows immediately that
the functions \(c_i(g)\) arise from uniquely determined sections
\[
  c_i\in\Gamma(G,\mathcal O_G).
\]

Suppose that \(g\in G(k)\) belongs to a Cartan subgroup, which we may
take to be \(C\); this includes the case of a regular element. By
\SGARefLink{sga3:xiii:corollary:2:5:item:iv}{2.5(iv)},
\[
  (\mathfrak g/\mathfrak c)^{\operatorname{ad}(g)}=0
\]
if and only if \(g\) is regular, where
\(\mathfrak c=\operatorname{Lie}(C)\). Since \(C\) is nilpotent,
\(\operatorname{ad}_{\mathfrak c}(g)\) has only the eigenvalue \(1\).
It follows that the multiplicity of the eigenvalue \(1\) in
\(\operatorname{ad}_{\mathfrak g}(g)\) is at least \(\rho\), and is
exactly \(\dim C=\rho\) if and only if \(g\) is regular.

In particular, \SGARefLink{sga3:xiii:theorem:2:6:display:characteristic-polynomial}{the characteristic polynomial above} is divisible by
\((t-1)^\rho\). Divisibility by \((t-1)^\rho\) is expressed by
linear relations, with integer coefficients, among the coefficients
of the polynomial. These relations hold for \(g\in U(k)\), where
\(U\) is dense; since \(G\) is reduced, they therefore hold
everywhere. Thus\footnote{The French source prints \(c_0\) in place
of the omitted tail \(+\cdots+c_n\) on the left-hand side of
\SGARefLink{sga3:xiii:theorem:2:6:equation:dagger}{\((\dagger)\)}. \SGARefLink{sga3:xiii:theorem:2:6:display:characteristic-polynomial}{The characteristic polynomial displayed immediately
above} determines the correction made here.}
\SGARefTarget{sga3:xiii:theorem:2:6:equation:dagger}
\[
\tag{\(\dagger\)}
  t^n+c_1t^{n-1}+\cdots+c_n
  =(t-1)^\rho
    \bigl(t^{n-\rho}+b_1t^{n-\rho-1}+\cdots+b_{n-\rho}\bigr)
\]
in the polynomial ring over \(\Gamma(G,\mathcal O_G)\). In
particular, for every field extension \(K/k\) and every
\(g\in G(K)\), the eigenvalue \(1\) of \(\operatorname{ad}(g)\) has
multiplicity at least \(\rho\).

We have already seen that equality holds when \(g\) is regular; let
us prove the converse. Suppose first that \(G\) is affine, and write
\[
  g=g_sg_u
\]
as the product of its semisimple and unipotent parts
(\emph{BIBLE} 4, Section 4). Then
\[
  \operatorname{ad}(g)
  =\operatorname{ad}(g_s)\operatorname{ad}(g_u)
\]
is the corresponding decomposition of \(\operatorname{ad}(g)\)
(loc. cit., corollary to Theorem 3). Hence
\(\operatorname{ad}(g)\) and \(\operatorname{ad}(g_s)\) have the same
eigenvalues, counted with multiplicity; in particular, the
eigenvalue \(1\) has the same multiplicity in both.

By \emph{BIBLE} 7, Theorem 2, Corollary 1, the element \(g\) is
regular if and only if \(g_s\) is regular. We may therefore assume
that \(g=g_s\) is semisimple. It is then contained in a maximal torus
by \emph{BIBLE} 6, Theorem 5(c), and hence in a Cartan subgroup, a
case already treated. This proves \SGARefLink{sga3:xiii:theorem:2:6:item:b}{(b)} when \(G\) is affine.

In general, let \(Z=\operatorname{Centr}(G)_{\mathrm{red}}\). By
\SGARefLink{sga3:xii:6.6:item:e}{XII~6.6(e)}, the Cartan subgroups of \(G\) are the inverse images of
those of \(G'=G/Z\). Thus \(g\) is regular in \(G\) if and only if
its image \(g'\) is regular in \(G'\). Since \(Z\) is smooth, the Lie
algebra of \(G'\) is
\(\mathfrak g'=\mathfrak g/\mathfrak z\), where
\(\mathfrak z=\operatorname{Lie}(Z)\), and
\(\operatorname{ad}(g')\) is the endomorphism induced by
\(\operatorname{ad}(g)\) on \(\mathfrak g/\mathfrak z\).
Consequently, the multiplicity of the eigenvalue \(1\) in
\(\operatorname{ad}(g)\) is
\[
  d=\dim Z
\]
plus its multiplicity in \(\operatorname{ad}(g')\). The former
therefore equals the nilpotent rank of \(G\) if and only if the latter
equals the nilpotent rank of \(G'\). We are reduced to \(G'\), which
is affine by \SGARefLink{sga3:xii:6.1}{XII~6.1}, and that case has already been proved. This
completes the proof of \SGARefLink{sga3:xiii:theorem:2:6}{2.6}.

\medskip
\noindent\SGARefTarget{sga3:xiii:corollary:2:7}\textbf{Corollary 2.7.} ---
\textit{With the notation of the preceding proof,\footnote{Editor's
note in the French source: cf. \SGARefLink{sga3:xiii:theorem:2:6:equation:dagger}{\((\dagger)\)} above.} let}
\[
  b=1+b_1+\cdots+b_{n-\rho}\in\Gamma(G,\mathcal O_G).
\]
\textit{Then the open subset of regular points of \(G\) is}
\[
  U=G_b,
\]
\textit{the set of points of \(G\) at which \(b\) is invertible. In
particular, \(U\) is an affine open subset when \(G\) is affine.}

\medskip
\noindent\SGARefTarget{sga3:xiii:corollary:2:8}\textbf{Corollary 2.8.} ---
\textit{Let \(H\) be a smooth connected algebraic subgroup of \(G\)
containing a Cartan subgroup of \(G\).}

\begin{enumerate}
\item[\textit{a)}] \SGARefTarget{sga3:xiii:corollary:2:8:item:a}
  Let \(C\) be an algebraic subgroup of \(H\). Then \(C\) is a
  Cartan subgroup of \(H\) if and only if it is a Cartan subgroup of
  \(G\).

\item[\textit{b)}] \SGARefTarget{sga3:xiii:corollary:2:8:item:b}
  Let \(g\in G(k)\), and let \(d\) be the integer introduced in \SGARefLink{sga3:xiii:corollary:2:3}{2.3}.
  Then \(g\) is a regular point of \(G\) if and only if there are
  exactly \(d\) conjugates \(H_i\) of \(H\) containing \(g\), and
  \(g\) is a regular element of every \(H_i\). Equivalently, there
  are at most \(d\) conjugates of \(H\) containing \(g\), and \(g\)
  is regular in one of them. When these conditions hold, if \(C\) is
  the unique Cartan subgroup of \(G\) containing \(g\), then the
  conjugates of \(H\) containing \(g\) are precisely the conjugates
  of \(H\) containing \(C\).

\item[\textit{c)}] \SGARefTarget{sga3:xiii:corollary:2:8:item:c}
  Let \(g\in H(k)\). Then \(g\) is regular in \(G\) if and only if it
  is regular in \(H\) and
  \[
    (\mathfrak g/\mathfrak h)^{\operatorname{ad}(g)}=0.
  \]
\end{enumerate}

\noindent\textit{Proof.}
\SGARefTarget{sga3:xiii:corollary:2:8:proof:item:a}
\textit{a)} Under either hypothesis on \(C\), its unique maximal
torus \(T\) is a maximal torus of both \(H\) and \(G\), since \(H\)
has the same reductive rank as \(G\). The smooth connected groups
\(\operatorname{Centr}_H(T)\subset\operatorname{Centr}_G(T)\) have
the same dimension, and therefore are equal. Hence \(C\) is equal to
one of these two centralizers if and only if it is equal to the
other, proving \SGARefLink{sga3:xiii:corollary:2:8:item:a}{(a)}.

\SGARefTarget{sga3:xiii:corollary:2:8:proof:item:b}
\textit{b)} Suppose first that \(g\) is regular in \(G\), and let
\(C\) be the unique Cartan subgroup of \(G\) containing \(g\). By
\SGARefLink{sga3:xiii:corollary:2:3:item:b}{2.3(b)}, exactly \(d\) conjugates \(H_i\) of \(H\) contain \(C\).
Since
\[
  (\mathfrak g/\mathfrak c)^{\operatorname{ad}(g)}=0,
\]
the endomorphism \(\operatorname{ad}(g)\) has no eigenvalue \(1\) on
\(\mathfrak g/\mathfrak c\), and a fortiori
\[
  (\mathfrak g/\mathfrak h_i)^{\operatorname{ad}(g)}=0.
\]
By \SGARefLink{sga3:xiii:corollary:2:3:item:c}{2.3(c)}, exactly \(d\) conjugates of \(H\) contain \(g\), namely
the \(H_i\). A Cartan subgroup of such an \(H_i\) containing \(g\) is
a Cartan subgroup of \(G\) by \SGARefLink{sga3:xiii:corollary:2:8:item:a}{(a)}, and hence is \(C\). Thus \(g\) is
regular in every \(H_i\).

Conversely, suppose that at most \(d\) conjugates \(H_i\) of \(H\)
contain \(g\), and that \(g\) is regular in one of them, which we may
take to be \(H\). Since \(g\) is regular in \(H\), it belongs to a
unique Cartan subgroup \(C\) of \(H\); by \SGARefLink{sga3:xiii:corollary:2:8:item:a}{(a)}, \(C\) is a Cartan
subgroup of \(G\). Let \(C'\) be any Cartan subgroup of \(G\)
containing \(g\). By \SGARefLink{sga3:xiii:corollary:2:3:item:b}{2.3(b)}, exactly \(d\) conjugates of \(H\)
contain \(C'\). They all contain \(g\), so they are precisely the
\(H_i\); in particular, \(H\) contains \(C'\). Thus \(C\) and \(C'\)
are two Cartan subgroups of \(H\) containing the same regular element
\(g\) of \(H\), and therefore \(C=C'\). Hence \(g\) is regular in
\(G\).\footnote{The source prints ``regular in \(C\)'' at this
conclusion; the argument proves regularity in \(G\).}
\hfill\(\square\)

\SGARefTarget{sga3:xiii:corollary:2:8:proof:item:c}
\textit{c)} For an endomorphism \(u\) of a finite-dimensional vector
space, let \(\nu(u)\) denote the nullity\footnote{Editor's note in the
French source: that is, the dimension of the nullspace of
\(\operatorname{id}-u\).} of \(\operatorname{id}-u\). Then
\SGARefTarget{sga3:xiii:corollary:2:8:display:nullity}
\[
  \nu\bigl(\operatorname{ad}(g)_{\mathfrak g}\bigr)
  =
  \nu\bigl(\operatorname{ad}(g)_{\mathfrak h}\bigr)
  +
  \nu\bigl(\operatorname{ad}(g)_{\mathfrak g/\mathfrak h}\bigr).
\]
The two terms on the right are respectively at least the nilpotent
rank of \(H\), which equals the nilpotent rank \(\rho\) of \(G\) by
\SGARefLink{sga3:xiii:corollary:2:8:item:a}{(a)}, and at least \(0\). Hence
\(\nu(\operatorname{ad}(g)_{\mathfrak g})=\rho\) if and only if
\[
  \nu\bigl(\operatorname{ad}(g)_{\mathfrak h}\bigr)=\rho,
  \qquad
  \nu\bigl(\operatorname{ad}(g)_{\mathfrak g/\mathfrak h}\bigr)=0.
\]
This says precisely that \(g\) is regular in \(G\) if and only if it
is regular in \(H\) and
\(\operatorname{ad}(g)_{\mathfrak g/\mathfrak h}\) has no nonzero
invariants. \hfill\(\square\)

\medskip
\noindent\SGARefTarget{sga3:xiii:remark:2:9}\textbf{Remarks 2.9.} ---
In the statement of \SGARefLink{sga3:xiii:theorem:2:1}{2.1}, condition \SGARefLink{sga3:xiii:theorem:2:1:item:iii}{(iii)} cannot be weakened by
assuming only that \(H\) contains a maximal torus and has finite
index in its normalizer, even if one also requires that normalizer to
be smooth, that is, that \(H=N^0\), and even when \(G\) is affine and
solvable. An example is the group \(G\) of matrices of the form
\SGARefTarget{sga3:xiii:remark:2:9:display:matrix}
\[
  g=
  \begin{pmatrix}
    t & a & c\\
    0 & 1 & b\\
    0 & 0 & 1
  \end{pmatrix},
\]
and the subgroup \(H\) formed by \SGARefLink{sga3:xiii:remark:2:9:display:matrix}{the matrices above} for which
\(b=c=0\). The Cartan subgroup of \(G\) here consists of the matrices
\(g\) for which \(a=c=0\).

\medskip
\noindent\SGARefTarget{sga3:xiii:remark:2:10}\textbf{Remarks 2.10.} ---
Let \(G\) be a smooth algebraic group over \(k\), and let \(H\) be a
smooth algebraic subgroup; we no longer assume that \(G\) and \(H\)
are connected. Suppose that \(H^0\) contains a Cartan subgroup of
\(G^0\). Then
\[
  (\mathfrak g/\mathfrak h)^H
  \subset(\mathfrak g/\mathfrak h)^{H^0}=0,
\]
so \(H^0=N^0\), where \(N\) is the normalizer of \(H\); in
particular, \(N\) is smooth.

Nevertheless, even with \(G\) connected, examples are easily
constructed in which \(H\) has a connected component \(H_i\) such
that
\[
  (\mathfrak g/\mathfrak h)^{\operatorname{ad}(h)}\neq0
  \qquad\text{for every }h\in H_i(k).
\]
Equivalently, the morphism
\[
  G\times H_i\longrightarrow G,\qquad
  (g,h)\longmapsto\operatorname{ad}(g)\cdot h
\]
is nowhere étale, or even quasi-finite. One may take, for example,
\(H\) to be the normalizer of a maximal torus in
\(\operatorname{SL}(2)_k\).

Similarly, even if the image \(H''\) of \(H\) in the finite group
\(G''=G/G^0\) is all of \(G''\), the union of the conjugates of \(H\)
in \(G\) need not be dense. For example, take \(G\) to be the
semidirect product of \(\mathbb Z/2\mathbb Z\) with
\(G^0=\operatorname{SL}(2)_k\), on which it acts by ``symmetry,'' and
take \(H\) to be the semidirect product
\(\mathbb Z/2\mathbb Z\cdot T\), where \(T\) is a maximal torus of
\(G^0\).

On the other hand, if one does not assume in advance that \(H^0\)
contains a Cartan subgroup of \(G^0\), but assumes that the union of
the conjugates of \(H\) in \(G\) is dense, then \(H^0\) necessarily
contains a Cartan subgroup of \(G^0\). To see this, one may plainly
assume that \(G\) is connected and repeat the proof of
\SGARefLink{sga3:xiii:theorem:2:1:item:vi}{2.1(vi)} \(\Rightarrow\) \SGARefLink{sga3:xiii:theorem:2:1:item:i}{(i)}, which does not require \(H\) to be
connected.

% END INCLUDED SGA3_English_current_buildable_source_20260823/source/inputs/XIII/source/tex/components/02_expose_XIII_section1_remainder_and_section2.tex

% BEGIN INCLUDED SGA3_English_current_buildable_source_20260823/source/inputs/XIII/source/tex/components/03_expose_XIII_sections3_4.tex
% Authority: Expo13.pdf, local pp. 14--22 / combined pp. 818--826
% (printed pp. 180--188), from the heading of section 3 through the
% end of Proposition 4.8, stopping immediately before the heading of section 5.
% Authority SHA-256:
% 24735CAFA57291A71C603712C0588358977856D5971503E6E73F3F2E40A70798

\makeatletter
\@ifundefined{remarks}{%
  \theoremstyle{remark}
  \newtheorem{remarks}[theorem]{Remarks}%
}{}
\makeatother

\section{The case of an arbitrary base prescheme}
\SGARefTarget{sga3:xiii:section:3}

Suppose first that the base is a field \(k\), not necessarily
algebraically closed. Since conditions \SGARefLink{sga3:xiii:theorem:2:1}{2.1} \SGARefLink{sga3:xiii:theorem:2:1:item:i-bis}{(i bis)}, \SGARefLink{sga3:xiii:theorem:2:1:item:iv-bis}{(iv bis)}, \SGARefLink{sga3:xiii:theorem:2:1:item:v}{(v)},
and \SGARefLink{sga3:xiii:theorem:2:1:item:vi}{(vi)} are invariant under extension of the base field, passage to
the algebraic closure \(\overline{k}\) of \(k\) shows that they are
equivalent to one another, and equivalent to the assertion that
\(H_{\overline{k}}\) contains a Cartan subgroup of
\(G_{\overline{k}}\). When \SGARefLink{sga3:xiii:theorem:2:1:item:i}{this condition} is satisfied, it is still
true (with the notation of \SGARefLink{sga3:xiii:corollary:2:3}{2.3}) that \(k(X)\) is a finite separable
extension of \(k(G)\), of degree \(d\) independent of every extension
of the base. If \(U\) is the largest open subset of \(G\) such that
\(\psi\) induces a finite étale morphism
\[
  \psi^{-1}(U)\longrightarrow U,
\]
then the formation of \(U\) commutes with extension of the base field.
If \(\psi\) is birational, then \(U\) is also the largest open subset
of \(G\) such that \(\psi\) induces an isomorphism
\[
  \psi^{-1}(U)\xrightarrow{\sim}U;
\]
and if \(g\in U(k)\), there is one and only one subgroup \(H'\) of
\(G\), conjugate to \(H\) over \(\overline{k}\), such that
\(H'\ni g\).

A point \(g\in G(k)\) is called \emph{regular} if it is regular as an
element of
\[
  G(\overline{k})=G_{\overline{k}}(\overline{k}).
\]
More generally, the construction of \SGARefLink{sga3:xiii:corollary:2:7}{2.7} gives an open subset of \(G\),
whose formation commutes with every extension of the base field,
called the \emph{open subset of regular points} of \(G\). It is also
characterized by the following property: for every algebraically
closed extension \(K\) of \(k\) and every \(g\in G(K)\), the point
\(g\) is regular in \(G_K\) if and only if \(g\in U(K)\). If
\(g\in U(k)\)\footnote{Editor's note: \(G(k)\) has been corrected to
\(U(k)\).}, then \(g\) is contained in one and only one Cartan
subgroup of \(G\), as we shall prove below under more general
hypotheses.

Let \(G\) be a smooth, separated group prescheme of finite type, with
connected fibers over the prescheme \(S\), and consider the functor
\SGARefTarget{sga3:xiii:section:3:display:cartan-functor}
\[
  \mathcal C:(\mathrm{Sch})^\circ_{/S}\longrightarrow(\mathrm{Ens})
\]
defined by
\[
  \mathcal C(S')
  =
  \left\{
    \begin{array}{c}
      \text{Cartan subgroups (XII~3.1) of \(G_{S'}\)}
    \end{array}
  \right\}.
\]
Suppose that \SGARefLink{sga3:xiii:section:3:display:cartan-functor}{this functor} is representable by a prescheme smooth over
\(S\). In Exposé XV we give a necessary and sufficient condition for
this to hold, but we already know that the hypothesis is satisfied if
\(G\) is affine over \(S\) with locally constant reductive rank
(\SGARefLink{sga3:xii:3.3}{XII~3.3}), or, more generally, if \(G\) locally admits a maximal torus
for the fpqc topology (\SGARefLink{sga3:xii:7.1}{XII~7.1} a)); this is the case, for example, if
\(S\) is the spectrum of a field.

Let \(X\) be the Cartan subgroup of the \(\mathcal C\)-group
prescheme \(G_{\mathcal C}\), the ``universal Cartan subgroup'' of
\(G\). Thus, as a prescheme over \(S\), \(X\) represents the functor
\[
  X(S')
  =
  \left\{
    \begin{array}{c}
      (C,g)\ \bigm|\ C\text{ is a Cartan subgroup of }G_{S'},\\
      g\text{ is a section of }C\text{ over }S'
    \end{array}
  \right\}.
\]
Consider the canonical projection morphism \((C,g)\mapsto g\),
\[
  \psi:X\longrightarrow G.
\]
We then have \SGARefLink{sga3:xiii:theorem:3:1}{the following theorem}.

\begin{theorem}
\SGARefTarget{sga3:xiii:theorem:3:1}
Under the preceding hypotheses on \(G\), and with the preceding
notation, let \(U\) be the set of \(g\in G\) such that \(g\) is a
regular element of its fiber \(G_s\). Then \(U\) is open, and it is
also the largest open subset \(U\) of \(G\) such that \(\psi\)
induces an isomorphism
\[
  \psi^{-1}(U)\xrightarrow{\sim}U.
\]
\end{theorem}

Let us first prove that \(U\) is open. By the representability of
\(\mathcal C\) as a prescheme smooth over \(S\), and since its
structure morphism is plainly surjective, it follows at once that
\(G\) locally admits a Cartan subgroup for the étale topology and
that the nilpotent rank of the fibers of \(G\) is locally constant.
The same holds for the dimension of the fibers of \(G\). After
localizing on \(S\), we may therefore suppose that these two integers
are constant, say \(\rho\) and \(n\). Consider the Killing polynomial
\[
  P_G(t)=t^n+c_1t^{n-1}+\cdots+c_n\in A[t],
  \qquad A=\Gamma(G,\mathcal O_G).
\]
The restriction of this polynomial to the fibers \(G_s\) of \(G\),
and in particular to the fibers at the maximal points of \(S\), is
divisible by \((t-1)^\rho\). This says that certain linear
combinations of the \(c_i\), with integer coefficients, vanish on the
fibers \(G_s\). When \(S\) is reduced, as we may assume in order to
prove that \(U\) is open, it follows that these combinations
themselves vanish. Hence the Killing polynomial itself is divisible
by \((t-1)^\rho\), say\footnote{The source prints the undefined
exponent \(n-r-1\). It is corrected to \(n-\rho-1\), as required by
the degree \(n-\rho\) of the second factor.}
\[
  P_G(t)
  =
  (t-1)^\rho
  \bigl(t^{n-\rho}+b_1t^{n-\rho-1}+\cdots+b_{n-\rho}\bigr).
\]
Let \(b\) be the sum of the coefficients
\(b_0=1,b_1,\ldots,b_{n-\rho}\) of the second factor. By \SGARefLink{sga3:xiii:corollary:2:7}{2.7}, applied
to the fibers of \(G\), we have
\[
  U=G_b,
\]
which proves that \(U\) is open.

To prove that \(\psi^{-1}(U)\to U\) is an isomorphism, SGA~1,
I~5.7 reduces us to checking this fiber by fiber. We are thereby
reduced to the case of a base field, which may be assumed
algebraically closed. There is then a Cartan subgroup \(C\) of \(G\).
If \(N\) is its normalizer, the conjugacy theorem \SGARefLink{sga3:xii:7.1}{XII~7.1} a) and b)
identifies \(\mathcal C\) with \(G/N\), and the morphism
\(\psi:X\to G\) considered here is precisely the one defined in
\SGARefLink{sga3:xiii:section:2}{Section~2}. The conclusion follows from \SGARefLink{sga3:xiii:theorem:2:6:item:b}{2.6 b)}. The same argument also
shows that \(U\) is the largest open subset of \(G\) for which
\(\psi\) induces an isomorphism \(\psi^{-1}(U)\to U\).

\begin{corollary}
\SGARefTarget{sga3:xiii:corollary:3:2}
Under the hypotheses of \SGARefLink{sga3:xiii:theorem:3:1}{3.1}, let \(g\) be a regular section of \(G\),
that is, suppose that \(g(s)\) is a regular point of \(G_s\) for every
\(s\in S\). Then there is one and only one Cartan subgroup \(C\) of
\(G\) such that \(g\) is a section of \(C\).
\end{corollary}

Indeed, the hypothesis on \(g\) says that \(g\) is a section of
\(U\), while the conclusion says that it admits a unique lift to a
section of \(X\). This is simply another way of saying that
\(\psi^{-1}(U)\to U\) is an isomorphism.

Observe now that the open subset \(\psi^{-1}(U)\) of the Cartan
subgroup \(X\) of \(G_{\mathcal C}\) is precisely the open subset of
\(X\) consisting of the points regular in \(G_{\mathcal C}\), by
which we mean regular in their fibers. We thus obtain a natural
``fibration'' of the dense open subset \(U\) of regular points of
\(G\) over the prescheme \(\mathcal C\); its fibers are dense open
subsets of Cartan subgroups of the fibers of \(G_{\mathcal C}\),
namely their open subsets of points regular in \(G\). For example, we
obtain \SGARefLink{sga3:xiii:corollary:3:3}{the following result}, which will be considerably sharpened in
the next Exposé.

\begin{corollary}
\SGARefTarget{sga3:xiii:corollary:3:3}
Let \(G\) be a smooth connected algebraic group over the field \(k\),
and let \(\mathcal T\) be the scheme of maximal tori of \(G\), which
is isomorphic to the scheme of Cartan subgroups of \(G\). Then the
function field \(k(G)\) of \(G\) is isomorphic to the function field
of a smooth connected affine nilpotent algebraic group \(C\) over the
function field \(k(\mathcal T)\) of \(\mathcal T\), namely
\[
  C=\text{``the generic Cartan subgroup of \(G\)''}.
\]
If \(G\) is affine and has unipotent rank zero, that is, if the Cartan
subgroups of \(G_{\overline k}\)\footnote{The source prints the
undefined base field \(\overline K\); it is corrected to
\(\overline k\).} are tori, then \(k(G)\) is a
unirational extension of \(k(\mathcal T)\).
\end{corollary}

Of course, by the generic Cartan subgroup of \(G\) we mean, by an
abuse of language, the Cartan subgroup of \(G_{k(\mathcal T)}\) that
is the generic fiber of \(X\) over \(\mathcal T\). It remains only to
prove the last assertion of \SGARefLink{sga3:xiii:corollary:3:3}{3.3}, which is contained in \SGARefLink{sga3:xiii:lemma:3:4}{the following
well-known result}, due to Chevalley.

\begin{lemma}
\SGARefTarget{sga3:xiii:lemma:3:4}
Let \(k\) be a field, \(T\) a torus over \(k\), and \(k(T)\) the
field of rational functions on \(T\). Then \(k(T)\) is a
unirational extension of \(k\), that is, it is contained in a purely
transcendental extension of \(k\).
\end{lemma}

Indeed, let \(k'\) be a finite separable extension of \(k\) that
splits\footnote{Editor's note: ``splits'' (?), or: ``over which \(T\)
is split''.} \(T\) (\SGARefLink{sga3:x:proposition:1:4}{X~1.4}). Then \(T\otimes_k k'\) is a rational
variety: it contains a dense open subset isomorphic to a dense open
subset of affine \(n\)-space \(\mathbb A^n_{k'}\). Hence
\[
  T'
  =
  \prod_{\operatorname{Spec}(k')/\operatorname{Spec}(k)}
  T_{k'}/\operatorname{Spec}(k')
\]
is a rational variety, since it has a dense open subset isomorphic to
a dense open subset of
\[
  \prod_{\operatorname{Spec}(k')/\operatorname{Spec}(k)}
  \mathbb A^n_{k'},
\]
which is isomorphic to affine \(mn\)-space over \(k\), where
\(m=[k':k]\). Consider the norm homomorphism from \(T'\) to \(T\),
which is defined whenever \(T\) is a commutative group scheme over
\(k\). The composite
\[
  T\longrightarrow T'\longrightarrow T
\]
is the \(m\)-th power map on \(T\), and is therefore dominant.
Consequently \(T'\to T\) is dominant, proving that \(T\) is
unirational.

Return to the hypotheses of \SGARefLink{sga3:xiii:theorem:3:1}{3.1}, but suppose in addition that \(G\)
locally admits a maximal torus for the fpqc topology (\SGARefLink{sga3:xii:7.1}{XII~7.1}). Let
\(Y\) be the maximal torus of the Cartan subgroup \(X\) of
\(G_{\mathcal C}\). The morphism \(\psi:X\to G\) then induces a
morphism \(Y\to G\) whose set-theoretic image consists of the
semisimple elements of the fibers of \(G\) (XII~8). Finally, \SGARefLink{sga3:xiii:theorem:3:1}{3.1}
shows that the restriction of \(\psi\) to the open subset
\(Y^{\mathrm{reg}}\) of regular points of \(Y\) induces a closed
immersion
\[
  Y^{\mathrm{reg}}\longrightarrow U=G^{\mathrm{reg}}.
\]
Making explicit the functorial meaning of
\(Z=Y^{\mathrm{reg}}\), we obtain \SGARefLink{sga3:xiii:corollary:3:5}{the following}.

\begin{corollary}
\SGARefTarget{sga3:xiii:corollary:3:5}
Let \(G\) be a smooth, separated \(S\)-group prescheme of finite
type, with connected fibers over the prescheme \(S\), and suppose
that \(G\) locally admits a maximal torus for the fpqc topology. Let
\[
  Z:(\mathrm{Sch})^\circ_{/S}\longrightarrow(\mathrm{Ens})
\]
be the functor defined by
\[
  Z(S')
  =
  \left\{
    \begin{array}{c}
      \text{regular sections of \(G_{S'}\) over \(S'\) that are}\\
      \text{contained in a maximal torus of \(G_{S'}\)}
    \end{array}
  \right\}.
\]
Then \(Z\) is representable by a closed subprescheme, smooth over
\(S\), of the open subset \(U=G^{\mathrm{reg}}\) of \(G\) introduced
in \SGARefLink{sga3:xiii:theorem:3:1}{3.1}.
\end{corollary}

To conclude, let us record \SGARefLink{sga3:xiii:corollary:3:6}{the following result}, which sharpens the
density theorem \SGARefLink{sga3:xiii:theorem:2:1}{2.1} \SGARefLink{sga3:xiii:theorem:2:1:item:i}{(i)} \(\Rightarrow\) \SGARefLink{sga3:xiii:theorem:2:1:item:vi}{(vi)}.

\begin{corollary}
\SGARefTarget{sga3:xiii:corollary:3:6}
Under the hypotheses of \SGARefLink{sga3:xiii:corollary:3:5}{3.5}, let \(C\) be a Cartan subgroup of \(G\),
and consider the morphism
\[
  \varphi:Z\times C\longrightarrow G
\]
defined by
\[
  \varphi(g,h)=\operatorname{ad}(g)h=ghg^{-1}.
\]
Then \(\varphi\) is dominant.
\end{corollary}

It plainly suffices to prove the assertion fiber by fiber, which
reduces us to the case where \(S\) is the spectrum of an
algebraically closed field \(k\). Let \(T\) be the maximal torus of
\(C\), let \(t_0\in T(k)\) be regular in \(G\), and let
\(c_0\in C(k)\) be regular in \(G\). Consider
\(\varphi^{-1}(\varphi(t_0,c_0))\). Its \(k\)-rational points are the
pairs \((t,c)\), with \(t\in Z(k)\) and \(c\in C(k)\), such that
\[
  \operatorname{ad}(t)c=\operatorname{ad}(t_0)c_0,
  \qquad\text{that is,}\qquad
  c=\operatorname{ad}(t^{-1}t_0)c_0.
\]
They are therefore in bijection with the \(t\in Z(k)\) such that
\(\operatorname{ad}(t^{-1}t_0)c_0\in C\), or equivalently, since
\(c_0\) is regular, such that \(t^{-1}t_0\in N\), where \(N\) is the
normalizer of \(C\); in other words, \(t\in N\). Restricting to those
\(t\in Z(k)\) that belong to \(C(k)\) gives an open and closed part
of this fiber. We have therefore found a connected component of
\(\varphi^{-1}(\varphi(t_0,c_0))\) isomorphic to \(T\). The preceding
set-theoretic argument indeed gives an isomorphism of schemes: one
simply replaces \(k\)-valued points by points with values in an
arbitrary \(k\)-prescheme.

It follows that the generic fiber of \(\varphi\) has dimension at
most \(\dim T\), and therefore
\[
  \dim\operatorname{Im}(\varphi)
  \geq
  \dim(Z\times C)-\dim T
  =
  \dim Z+\dim C-\dim T.
\]
But
\[
  \dim Z
  =
  \dim Y
  =
  \dim\mathcal C+\dim T
  =
  \dim G-\dim C+\dim T.
\]
Thus \(\dim\operatorname{Im}(\varphi)\geq\dim G\), and \(\varphi\)
is dominant.
\par\hfill\textsc{q.e.d.}

\begin{remarks}
\SGARefTarget{sga3:xiii:remark:3:7}
The argument shows in addition that the connected component at
\((t_0,c_0)\) of the fiber
\(\varphi^{-1}(\varphi(t_0,c_0))\) is isomorphic to \(T\). In
particular, it is smooth over \(k\) and has the same dimension as the
generic fiber. It follows that \(\varphi\) is in fact smooth at
\((t_0,c_0)\), something that one should also be able to verify by
calculating the tangent map. Consequently, under the hypotheses of
\SGARefLink{sga3:xiii:corollary:3:6}{3.6}, the induced morphism
\[
  Z\times_S C^{\mathrm{reg}}\longrightarrow G^{\mathrm{reg}},
  \qquad C^{\mathrm{reg}}=C\cap G^{\mathrm{reg}},
\]
is smooth. Likewise, the analogous morphism
\[
  Z\times T^{\mathrm{reg}}\longrightarrow Z,
\]
where \(T\) is a maximal torus of \(G\), is smooth. More generally,
for every smooth connected normal algebraic subgroup \(H\) of \(C\)
containing an element \(c_0\in G(k)\) that is regular in \(G\), the
image of \(Z\times H\to G\) is dense in the image of
\(G\times H\to G\).
\end{remarks}

\section{Lie algebras over a field: rank, regular elements, Cartan subalgebras}
\SGARefTarget{sga3:xiii:section:4}

In the remainder of \SGARefLink{sga3:xiii:expose}{this Exposé} we take up the theory developed by
Chevalley in his book \emph{Théorie des Groupes de Lie III}
(Act.\ Sc.\ Ind.\ 1226, Paris 1955); the methods of schemes allow us
to dispense with the hypothesis of characteristic zero. We begin by
recalling in \SGARefLink{sga3:xiii:section:4}{this section} several well-known notions and results.

Let \(\mathfrak g\) be a Lie algebra over a ring \(k\). For every
\(a\in\mathfrak g\), denote by \(\operatorname{ad}(a)\) the
endomorphism
\[
  \operatorname{ad}(a)\cdot x=[a,x]
\]
of \(\mathfrak g\); it is a derivation of the Lie algebra
\(\mathfrak g\). For every derivation \(D\) of \(\mathfrak g\), the
nil-space of \(D\), that is, the union of the kernels of the iterates
of \(D\), is a Lie subalgebra of \(\mathfrak g\), as follows from the
Leibniz formula
\[
  D^m([x,y])
  =
  \sum_{0\leq p\leq m}
  \binom{m}{p}[D^p x,D^{m-p}y].
\]
We set
\[
  \operatorname{Nil}(a,\mathfrak g)
  =
  \text{nil-space of }\operatorname{ad}(a)
  =
  \bigcup_{m\geq0}\operatorname{Ker}\operatorname{ad}(a)^m.
\]
When no confusion can arise, we write simply
\(\operatorname{Nil}(a)\), and call it the \emph{nil-space of \(a\)
in \(\mathfrak g\)}.

\begin{proposition}
\SGARefTarget{sga3:xiii:proposition:4:1}
For every \(a\in\mathfrak g\), its nil-space
\(\operatorname{Nil}(a,\mathfrak g)\) is a Lie subalgebra of
\(\mathfrak g\), equal to its own normalizer.
\end{proposition}

It remains only to prove that it is its own normalizer. In other
words, every element of
\(\mathfrak g/\operatorname{Nil}(a,\mathfrak g)\) annihilated by the
adjoint representation of \(\operatorname{Nil}(a,\mathfrak g)\) on
this quotient must be zero. This is immediate, since every element
of the quotient annihilated by \(\operatorname{ad}(a)\)\footnote{The
source prints \(\operatorname{Ad}(a)\); because
\(a\in\mathfrak g\), the endomorphism here is
\(\operatorname{ad}(a)\).} is zero.

For the remainder of \SGARefLink{sga3:xiii:section:4}{this section}, suppose that \(k\) is a field and
that \(\mathfrak g\) is finite-dimensional over \(k\). We denote by
\(W(\mathfrak g)\) the scheme over \(k\) defined by \(\mathfrak g\);
its points in a \(k\)-algebra \(A\) are the elements of
\(\mathfrak g\otimes_k A\). If \(a\in\mathfrak g\), the
characteristic polynomial of \(\operatorname{ad}(a)\) is also called
the \emph{characteristic polynomial}, or the \emph{Killing
polynomial}, of \(a\) in \(\mathfrak g\):
\[
  P_{\mathfrak g}(a,t)
  =
  t^n+c_1(a)t^{n-1}+\cdots+c_n(a),
\]
where \(n=\operatorname{rank}_k\mathfrak g\) and \(c_i(a)\in k\).
Taking the same polynomial for
\(a\in\mathfrak g\otimes_k A\)\footnote{Editor's note:
\(W\otimes_k A\) has been corrected to
\(\mathfrak g\otimes_k A\).}, where \(A\) is any \(k\)-algebra, we
see that the \(c_i(a)\) arise from uniquely determined sections
\(c_i\) of the structure sheaf of \(W(\mathfrak g)\), that is, from
elements of the symmetric algebra
\[
  A=\operatorname{Sym}_k(\mathfrak g^\vee),
\]
where \(\mathfrak g^\vee\) is the dual of the \(k\)-module
\(\mathfrak g\). When \(k\) is infinite, the \(c_i\) are determined
by the corresponding polynomial functions \(\mathfrak g\to k\), but
this is no longer true when \(k\) is finite.

Let \(r\) be the largest integer such that the Killing polynomial
\[
  P_{\mathfrak g}(t)
  =
  t^n+c_1t^{n-1}+\cdots+c_n
  \in A[t]
\]
is divisible by \(t^r\); thus
\[
  P_{\mathfrak g}(t)
  =
  t^n+c_1t^{n-1}+\cdots+c_{n-r}t^r,
  \qquad c_{n-r}\neq0.
\]
The integer \(r\) is called the \emph{nilpotent rank} of the Lie
algebra \(\mathfrak g\). It is invariant under extension of the base
field.

\begin{proposition}
\SGARefTarget{sga3:xiii:proposition:4:2}
Let \(r\) be the nilpotent rank of \(\mathfrak g\), and let
\(a\in\mathfrak g\). Then
\[
  \operatorname{rank}_k\operatorname{Nil}(a,\mathfrak g)\geq r,
\]
with equality if and only if
\[
  c_{n-r}(a)\neq0.
\]
In this case \(\operatorname{Nil}(a,\mathfrak g)\) is a nilpotent
Lie algebra. We shall see the converse in \SGARefLink{sga3:xiii:corollary:5:7:part:b}{5.7 b)} when
\(\mathfrak g\) is the Lie algebra of an algebraic group \(G\)
smooth over \(k\).
\end{proposition}

The first assertion is immediate, since by definition
\[
  \operatorname{rank}_k\operatorname{Nil}(a,\mathfrak g)
  =
  \text{the multiplicity of the zero root in }P_{\mathfrak g}(a,t).
\]
Suppose now that \(c_{n-r}(a)\neq0\). We prove that
\(\operatorname{Nil}(a)\) is nilpotent, which is equivalent to saying
that \(\operatorname{ad}_{\operatorname{Nil}(a)}(x)\) is a nilpotent
endomorphism for every \(x\in\operatorname{Nil}(a)\). We may suppose
that \(k\) is algebraically closed. Since
\[
  \operatorname{ad}(a)_{\mathfrak g/\operatorname{Nil}(a)}
\]
is injective, there is a nonempty open subset
\(U\subset W(\operatorname{Nil}(a))\) such that, for every
\(x\in U(k)\), the endomorphism
\(\operatorname{ad}(x)_{\mathfrak g/\operatorname{Nil}(a)}\) is
injective and hence
\(\operatorname{Nil}(x)\subset\operatorname{Nil}(a)\). We may further
suppose that \(U\) is contained in the open subset on which
\(c_{n-r}\) does not vanish; that open subset is nonempty because
\(c_{n-r}(a)\neq0\). Then
\(\operatorname{Nil}(x)\) and \(\operatorname{Nil}(a)\) have the same
dimension, so
\[
  \operatorname{Nil}(x)=\operatorname{Nil}(a).
\]
It follows that
\(\operatorname{ad}(x)_{\operatorname{Nil}(a)}\) is nilpotent for
every \(x\in U(k)\). By the principle of extension of algebraic
identities, this remains true for every
\(x\in\operatorname{Nil}(a)\), and hence \(\operatorname{Nil}(a)\) is
nilpotent.

An element \(a\in\mathfrak g\) is called \emph{regular} if
\(c_{n-r}(a)\neq0\), that is, if
\[
  \operatorname{rank}_k\operatorname{Nil}(a,\mathfrak g)=r.
\]
When \(k\) is infinite, this also says that
\(\operatorname{rank}_k\operatorname{Nil}(a,\mathfrak g)\) is as
small as possible as \(a\) varies in \(\mathfrak g\). In every case,
the notion of a regular element of \(\mathfrak g\) is invariant under
extension of the base field. The set of regular points of
\(W(\mathfrak g)\), meaning those that arise from regular points with
values in a suitable extension of \(k\), is open, since it is
\[
  W(\mathfrak g)_{c_{n-r}},
\]
the set of points at which \(c_{n-r}\) is invertible.

\begin{corollary}
\SGARefTarget{sga3:xiii:corollary:4:3}
Let \(a\) be a regular element of \(\mathfrak g\), and let
\(\mathfrak h\) be a Lie subalgebra of \(\mathfrak g\) containing
\(a\). Then \(\mathfrak h\) is nilpotent if and only if
\[
  \mathfrak h\subset\operatorname{Nil}(a,\mathfrak g).
\]
In particular, \(\operatorname{Nil}(a,\mathfrak g)\) is a maximal
nilpotent subalgebra of \(\mathfrak g\).
\end{corollary}

Since \(\operatorname{Nil}(a)\) is nilpotent, the inclusion
\(\mathfrak h\subset\operatorname{Nil}(a)\) implies that
\(\mathfrak h\) is nilpotent. Conversely, if \(\mathfrak h\) is
nilpotent, it is contained in the nil-space of its element \(a\);
that is, \(\mathfrak h\subset\operatorname{Nil}(a)\).

\begin{proposition}
\SGARefTarget{sga3:xiii:proposition:4:4}
Suppose that \(k\) is infinite. Let \(\mathfrak d\) be a Lie
subalgebra of \(\mathfrak g\). Consider the following conditions:
\begin{enumerate}
\item[(i)]\SGARefTarget{sga3:xiii:proposition:4:4:item:i}
\(\mathfrak d\) is maximal nilpotent and contains a regular element
of \(\mathfrak g\).
\item[(i bis)]\SGARefTarget{sga3:xiii:proposition:4:4:item:i-bis}
\(\mathfrak d=\operatorname{Nil}(a,\mathfrak g)\) for some regular
element \(a\) of \(\mathfrak g\).
\item[(ii)]\SGARefTarget{sga3:xiii:proposition:4:4:item:ii}
\(\mathfrak d\) is nilpotent and
\(\mathfrak d=\operatorname{Nil}(a,\mathfrak g)\) for some
\(a\in\mathfrak g\).
\item[(ii bis)]\SGARefTarget{sga3:xiii:proposition:4:4:item:ii-bis}
\(\mathfrak d\) is nilpotent, and there exists \(a\in\mathfrak d\)
such that \(\operatorname{ad}(a)_{\mathfrak g/\mathfrak d}\) is
injective.
\item[(iii)]\SGARefTarget{sga3:xiii:proposition:4:4:item:iii}
\(\mathfrak d\) is nilpotent and equal to its own normalizer.
\end{enumerate}
Then
\[
  \SGARefLink{sga3:xiii:proposition:4:4:item:i}{\text{\rm(i)}}
  \Longleftrightarrow
  \text{\rm\SGARefLink{sga3:xiii:proposition:4:4:item:i-bis}{(i bis)}}
  \Longrightarrow
  \SGARefLink{sga3:xiii:proposition:4:4:item:ii}{\text{\rm(ii)}}
  \Longleftrightarrow
  \text{\rm\SGARefLink{sga3:xiii:proposition:4:4:item:ii-bis}{(ii bis)}}
  \Longleftrightarrow
  \SGARefLink{sga3:xiii:proposition:4:4:item:iii}{\text{\rm(iii)}}.
\]
We shall see in \SGARefLink{sga3:xiii:corollary:5:7:part:a}{5.7 a)} that if \(\mathfrak g\) is the Lie algebra of
a smooth algebraic group, then all the preceding conditions are
equivalent.
\end{proposition}

The equivalence of \SGARefLink{sga3:xiii:proposition:4:4:item:i}{(i)} and \SGARefLink{sga3:xiii:proposition:4:4:item:i-bis}{(i bis)} follows immediately from \SGARefLink{sga3:xiii:corollary:4:3}{4.3}, and
these conditions plainly imply \SGARefLink{sga3:xiii:proposition:4:4:item:ii}{(ii)}. The equivalence of \SGARefLink{sga3:xiii:proposition:4:4:item:ii}{(ii)} and
\SGARefLink{sga3:xiii:proposition:4:4:item:ii-bis}{(ii bis)} is also immediate, as is \SGARefLink{sga3:xiii:proposition:4:4:item:ii-bis}{(ii bis)} \(\Rightarrow\) \SGARefLink{sga3:xiii:proposition:4:4:item:iii}{(iii)}
(cf. \SGARefLink{sga3:xiii:proposition:4:1}{4.1}). It remains to prove \SGARefLink{sga3:xiii:proposition:4:4:item:iii}{(iii)} \(\Rightarrow\) \SGARefLink{sga3:xiii:proposition:4:4:item:ii-bis}{(ii bis)}. This
is the only implication that uses the infinitude of \(k\), and it
follows at once from \SGARefLink{sga3:xiii:lemma:4:5}{the next lemma}.

\begin{lemma}
\SGARefTarget{sga3:xiii:lemma:4:5}
Let \(\mathfrak d\) be a nilpotent Lie algebra over an infinite field
\(k\), acting on a finite-dimensional vector space \(V\). Suppose
that the endomorphism \(u(x)\) is not injective for every
\(x\in\mathfrak d\). Then there is a nonzero element \(v\in V\)
annihilated by \(\mathfrak d\).
\end{lemma}

We may suppose that \(k\) is algebraically closed and that
\(\mathfrak d\) is finite-dimensional. Then \(V\) is known to be a
direct sum of finitely many nonzero stable vector subspaces \(V_i\),
\(1\leq i\leq n\), such that, for every \(i\) and every
\(x\in\mathfrak d\), the restriction \(u(x)|_{V_i}\) has a single
eigenvalue \(\lambda_i(x)\) (cf. Bourbaki, \emph{Groupes et Algèbres
de Lie}, Chap.~I, \(\S\)4, Exercise~22). Let \(c_i(x)\) be the
constant term of the characteristic polynomial of
\(u(x)|_{V_i}\), so that
\[
  \lambda_i(x)=0\quad\Longleftrightarrow\quad c_i(x)=0.
\]
Each \(c_i\) is a polynomial function on \(\mathfrak d\), and the
hypothesis says that \(\mathfrak d\) is the union of their zero sets.
Thus one of the \(c_i\) is identically zero. Replacing \(V\) by
\(V_i\), we are reduced to the case in which all \(u(x)\),
\(x\in\mathfrak d\), are nilpotent. Engel's theorem (Bourbaki,
loc.\ cit., Th.~1) then implies that some nonzero \(v\in V\) is
annihilated by \(\mathfrak d\).
\par\hfill\textsc{q.e.d.}

It is easy to see, still assuming that \(k\) is infinite, that
conditions \SGARefLink{sga3:xiii:proposition:4:4:item:i}{(i)} and \SGARefLink{sga3:xiii:proposition:4:4:item:i-bis}{(i bis)} of \SGARefLink{sga3:xiii:proposition:4:4}{4.4} are invariant under every extension
of the base field. When they hold, \(\mathfrak d\) is called a
\emph{Cartan subalgebra} of \(\mathfrak g\). In general, with \(k\)
not necessarily infinite, \(\mathfrak d\) is called a Cartan
subalgebra of \(\mathfrak g\) if it becomes one after one, and hence
every, extension \(k\to k'\) with \(k'\) infinite. It follows in
particular that \(\mathfrak d\) is nilpotent and equal to its own
normalizer.

\begin{proposition}
\SGARefTarget{sga3:xiii:proposition:4:6}
\SGARefTarget{sga3:xiii:proposition:4:6:item:a}
\textup{a)} Let \(a\in\mathfrak g\). If \(a\) is regular, then it is
contained in one and only one Cartan subalgebra of \(\mathfrak g\).
We shall see the converse in \SGARefLink{sga3:xiii:theorem:6:1:part:d}{6.1 d)}, when \(k\) is algebraically
closed and \(\mathfrak g\) is the Lie algebra of a smooth algebraic
group.

\SGARefTarget{sga3:xiii:proposition:4:6:item:b}
\textup{b)} Let \(\mathfrak d\) be a Cartan subalgebra of
\(\mathfrak g\), and let \(a\in\mathfrak d\). Then \(a\) is regular
in \(\mathfrak g\) if and only if
\(\operatorname{ad}(a)_{\mathfrak g/\mathfrak d}\) is injective.
\end{proposition}

For \SGARefLink{sga3:xiii:proposition:4:6:item:a}{a)}, if \(a\) is regular, then
\(\operatorname{Nil}(a,\mathfrak g)\) is a Cartan subalgebra of
\(\mathfrak g\), since this is true after an infinite extension
\(k'\) of \(k\). It follows at once from \SGARefLink{sga3:xiii:corollary:4:3}{4.3} that every Cartan
subalgebra of \(\mathfrak g\) containing \(a\) is equal to this one.
For \SGARefLink{sga3:xiii:proposition:4:6:item:b}{b)}, observe that the nullity\footnote{Editor's note: that is, the
dimension of the nil-space.} of
\(\operatorname{ad}(a)_{\mathfrak g}\) is the sum of the nullities of
\(\operatorname{ad}(a)_{\mathfrak d}\) and
\(\operatorname{ad}(a)_{\mathfrak g/\mathfrak d}\). Since the first
is \(r\), their sum is \(r\) if and only if
\(\operatorname{ad}(a)_{\mathfrak g/\mathfrak d}\) is injective.
\par\hfill\textsc{q.e.d.}

\begin{corollary}
\SGARefTarget{sga3:xiii:corollary:4:7}
Let \(a\) be a regular element of \(\mathfrak g\), let
\(\mathfrak d\) be a Cartan subalgebra of \(\mathfrak g\) containing
\(a\), let \(A\) be a \(k\)-algebra, and let
\(\mathfrak g_A\) and \(\mathfrak d_A\subset\mathfrak g_A\) be the
\(A\)-Lie algebras obtained from \(\mathfrak g\) and \(\mathfrak d\)
by base change; write \(a_A\) for the image of \(a\) in
\(\mathfrak g_A\). Let \(u\) be an automorphism of
\(\mathfrak g_A\). Then
\[
  u(\mathfrak d_A)=\mathfrak d_A
\]
if and only if
\[
  u(a_A)\in\mathfrak d_A.
\]
\end{corollary}

The condition is plainly necessary; let us prove that it is
sufficient. Suppose it holds. Then
\[
  \mathfrak d'=u^{-1}(\mathfrak d_A)
\]
is a Lie subalgebra containing \(a_A\), and every \(b\in\mathfrak d'\)
is such that \(\operatorname{ad}(b)_{\mathfrak d'}\) is
nilpotent.\footnote{The source prints
\(\mathfrak d'=u(\mathfrak d_A)\) and
\(\operatorname{ad}(b)_{\mathfrak d}\). Under the hypothesis
\(u(a_A)\in\mathfrak d_A\), the inverse image
\(u^{-1}(\mathfrak d_A)\) is the subalgebra that contains \(a_A\),
and the nilpotence assertion concerns its restriction
\(\operatorname{ad}(b)_{\mathfrak d'}\).}
Indeed, \(\mathfrak d'\) is isomorphic to \(\mathfrak d_A\), which
has this property, as follows immediately from the definition of
``nilpotent'' in Bourbaki, \emph{Groupes et Algèbres de Lie},
Chap.~I, \(\S\)4, Def.~1. Taking \(b=a_A\), we see that
\(\operatorname{Nil}(b,\mathfrak g_A)\) contains \(\mathfrak d'\).
On the other hand,
\[
  \operatorname{Nil}(b,\mathfrak g_A)=\mathfrak d_A.
\]
Moreover, \(\mathfrak d'\) is locally a direct summand of
\(\mathfrak g_A\), since \(\mathfrak d\) is, and hence is locally a
direct summand of \(\mathfrak d_A\). It is a projective module of the
same rank \(r\) as \(\mathfrak d_A\), and therefore equals
\(\mathfrak d_A\).
\par\hfill\textsc{q.e.d.}

\begin{proposition}
\SGARefTarget{sga3:xiii:proposition:4:8}
Let \(\mathfrak h\) be a Lie subalgebra of \(\mathfrak g\).

\SGARefTarget{sga3:xiii:proposition:4:8:item:a}
\textup{a)} If \(k\) is infinite, the following conditions are
equivalent:
\begin{enumerate}
\item[(i)]\SGARefTarget{sga3:xiii:proposition:4:8:item:a:i}
\(\mathfrak h\) contains a Cartan subalgebra \(\mathfrak d\) of
\(\mathfrak g\).
\item[(ii)]\SGARefTarget{sga3:xiii:proposition:4:8:item:a:ii}
\(\mathfrak h\) contains a regular element \(a\) of \(\mathfrak g\),
and an element \(b\) such that
\(\operatorname{ad}(b)_{\mathfrak g/\mathfrak h}\) is injective.
\item[(iii)]\SGARefTarget{sga3:xiii:proposition:4:8:item:a:iii}
\(\mathfrak h\) has the same nilpotent rank as \(\mathfrak g\), and
contains a regular element of \(\mathfrak g\).
\end{enumerate}
These conditions are invariant under extension of the base field
\(k\).

\SGARefTarget{sga3:xiii:proposition:4:8:item:b}
\textup{b)} Suppose that these conditions hold after a suitable
infinite extension \(k'\) of \(k\). For \(a\in\mathfrak h\), the
element \(a\) is regular in \(\mathfrak g\) if and only if it is
regular in \(\mathfrak h\) and
\(\operatorname{ad}(a)_{\mathfrak g/\mathfrak h}\) is injective;
equivalently, if and only if
\[
  \operatorname{Nil}(a,\mathfrak h)
  =
  \operatorname{Nil}(a,\mathfrak g)
\]
and this common Lie algebra is a Cartan subalgebra of
\(\mathfrak h\).

\SGARefTarget{sga3:xiii:proposition:4:8:item:c}
\textup{c)} Under the hypotheses of \SGARefLink{sga3:xiii:proposition:4:8:item:b}{b)}, let \(\mathfrak d\) be a Lie
subalgebra of \(\mathfrak h\). In order that \(\mathfrak d\) be a
Cartan subalgebra of \(\mathfrak h\), it is sufficient that it be a
Cartan subalgebra of \(\mathfrak g\). We shall see in \SGARefLink{sga3:xiii:corollary:5:8}{5.8} that \SGARefLink{sga3:xiii:proposition:4:8:item:c}{this
condition} is also necessary when \(\mathfrak g\) is the Lie algebra
of a smooth algebraic group \(G\), and \(\mathfrak h\) is the Lie
algebra of a smooth algebraic subgroup \(H\) of \(G\).
\end{proposition}

Conditions \SGARefLink{sga3:xiii:proposition:4:8:item:a:ii}{(ii)} and \SGARefLink{sga3:xiii:proposition:4:8:item:a:iii}{(iii)} of \SGARefLink{sga3:xiii:proposition:4:8:item:a}{a)} are visibly invariant under extension
of the base field \(k\), assumed infinite. In \SGARefLink{sga3:xiii:proposition:4:8:item:b}{b)} and \SGARefLink{sga3:xiii:proposition:4:8:item:c}{c)}, we may
suppose that \(k\) is infinite, and shall do so. If \(\mathfrak h\)
contains the Cartan subalgebra
\[
  \mathfrak d=\operatorname{Nil}(a,\mathfrak g),
\]
then \(a\) is a regular element of \(\mathfrak g\) such that
\(\operatorname{ad}(a)_{\mathfrak g/\mathfrak h}\) is injective.
Thus \SGARefLink{sga3:xiii:proposition:4:8:item:a:i}{(i)} \(\Rightarrow\) \SGARefLink{sga3:xiii:proposition:4:8:item:a:ii}{(ii)}. Conversely, if \SGARefLink{sga3:xiii:proposition:4:8:item:a:ii}{(ii)} holds, then a
``sufficiently general'' element \(a\) of \(\mathfrak h\)
simultaneously satisfies both conditions in \SGARefLink{sga3:xiii:proposition:4:8:item:a:ii}{(ii)}. Hence
\(\operatorname{Nil}(a,\mathfrak g)\) is a Cartan subalgebra of
\(\mathfrak g\) contained in \(\mathfrak h\), proving \SGARefLink{sga3:xiii:proposition:4:8:item:a:i}{(i)}.
Conditions \SGARefLink{sga3:xiii:proposition:4:8:item:a:i}{(i)} and \SGARefLink{sga3:xiii:proposition:4:8:item:a:ii}{(ii)} are therefore equivalent.

Suppose they hold, and let \(a\) vary in \(\mathfrak h\). Then
\SGARefTarget{sga3:xiii:equation:4:8:x}
\[
\tag{x}
  \operatorname{rank}_k\operatorname{Nil}(a,\mathfrak g)
  =
  \operatorname{rank}_k\operatorname{Nil}(a,\mathfrak h)
  +
  \operatorname{rank}_k
  \operatorname{Nil}\bigl(
    \operatorname{ad}(a)_{\mathfrak g/\mathfrak h}
  \bigr).
\]
The two terms on the right are respectively at least \(r'\), the
nilpotent rank of \(\mathfrak h\), and at least \(0\); moreover, both
equalities are attained\footnote{Editor's note: ``the equality'' has
been replaced by ``the equalities''.} for a sufficiently general
element of \(\mathfrak h\). We also have
\[
  \operatorname{rank}_k\operatorname{Nil}(a,\mathfrak g)\geq r,
\]
where \(r\) is the nilpotent rank of \(\mathfrak g\). Equality is
attained for a sufficiently general element of \(\mathfrak h\), and
holds if and only if \(a\) is regular in \(\mathfrak g\).

It follows that \(r=r'\), and that \(a\) is regular if and only if
the two terms on the right-hand side of \(\SGARefLink{sga3:xiii:equation:4:8:x}{\mathrm{(x)}}\) are
respectively \(r'\) and \(0\); that is, if and only if \(a\) is
regular in \(\mathfrak h\) and
\(\operatorname{ad}(a)_{\mathfrak g/\mathfrak h}\) is injective.
This proves \SGARefLink{sga3:xiii:proposition:4:8:item:b}{b)}, and \SGARefLink{sga3:xiii:proposition:4:8:item:c}{c)} follows immediately by taking an element
\(a\in\mathfrak d\) regular in \(\mathfrak g\), so that
\(\operatorname{Nil}(a,\mathfrak g)=\mathfrak d\).

\SGARefLink{sga3:xiii:proposition:4:8:item:b}{The preceding result} also proves \SGARefLink{sga3:xiii:proposition:4:8:item:a:i}{(i)} \(\Rightarrow\) \SGARefLink{sga3:xiii:proposition:4:8:item:a:iii}{(iii)}. Finally,
\SGARefLink{sga3:xiii:proposition:4:8:item:a:iii}{(iii)} \(\Rightarrow\) \SGARefLink{sga3:xiii:proposition:4:8:item:a:i}{(i)}: under \SGARefLink{sga3:xiii:proposition:4:8:item:a:iii}{(iii)}, a sufficiently general
\(a\in\mathfrak h\) is regular both in \(\mathfrak h\) and in
\(\mathfrak g\). Thus
\[
  \operatorname{Nil}(a,\mathfrak h)
  \subset
  \operatorname{Nil}(a,\mathfrak g)
\]
are respectively Cartan subalgebras of \(\mathfrak h\) and
\(\mathfrak g\). Since they have the same rank over \(k\), they are
equal, proving \SGARefLink{sga3:xiii:proposition:4:8:item:a:i}{(i)}. This completes the proof of \SGARefLink{sga3:xiii:proposition:4:8}{4.8}.

% END INCLUDED SGA3_English_current_buildable_source_20260823/source/inputs/XIII/source/tex/components/03_expose_XIII_sections3_4.tex

% BEGIN INCLUDED SGA3_English_current_buildable_source_20260823/source/inputs/XIII/source/tex/components/04_expose_XIII_sections5_6.tex
% Authority: Expo13.pdf, local pp. 22--29 / combined pp. 826--833.
% SHA-256: 24735CAFA57291A71C603712C0588358977856D5971503E6E73F3F2E40A70798.
% The blank local p. 30 / combined p. 834 is excluded; Exposé XIV begins
% on combined p. 835.

\section{The case of the Lie algebra of a smooth algebraic group:
the density theorem}
\SGARefTarget{sga3:xiii:section:5}

Let \(G\) be a smooth algebraic group over the field \(k\), and let
\(\mathfrak g\) be its Lie algebra. Let \(\mathfrak h\) be a Lie
subalgebra of \(\mathfrak g\). Let \(G\) act on \(W(\mathfrak g)\)
through the adjoint representation, and consider the subscheme
\(W(\mathfrak h)\). The construction of \SGARefLink{sga3:xiii:section:1}{Section~1} leads us to introduce
\[
  N=\operatorname{Norm}_G(\mathfrak h)
   =\operatorname{Norm}_G\bigl(W(\mathfrak h)\bigr),
\]
which is an algebraic subgroup of \(G\) (not necessarily smooth), and
\[
  \mathfrak n=\operatorname{Lie}(N),
\]
together with the scheme
\[
  X=G\times^N W(\mathfrak h),
\]
the quotient of \(\widetilde X=G\times W(\mathfrak h)\) by the right
action of \(N\) given by
\[
  (g,x)\cdot n=\bigl(gn,\operatorname{Ad}(n^{-1})x\bigr).
\]
We consider the canonical morphisms
\SGARefTarget{sga3:xiii:section:5:diagram:canonical-morphisms}
\[
\begin{tikzcd}[column sep=large,row sep=large]
G\times W(\mathfrak h)
  \arrow[d,"q"']
  \arrow[dr,"\varphi"] & \\
X \arrow[r,"\psi"] & W(\mathfrak g).
\end{tikzcd}
\]

\begin{theorem}
\SGARefTarget{sga3:xiii:theorem:5:1}
With the preceding notation, suppose that \(k\) is infinite. Consider
the following conditions:
\begin{enumerate}
  \item[(i)]\SGARefTarget{sga3:xiii:theorem:5:1:condition:i}
  \(\mathfrak h\) contains a Cartan subalgebra \(\mathfrak d\) of
  \(\mathfrak g\).

  \item[(ii)]\SGARefTarget{sga3:xiii:theorem:5:1:condition:ii}
  There exists \(a\in\mathfrak h\) such that
  \(\operatorname{ad}(a)_{\mathfrak g/\mathfrak h}\) is injective.

  \item[(iii)]\SGARefTarget{sga3:xiii:theorem:5:1:condition:iii}
  The morphism
  \(\varphi:G\times W(\mathfrak h)\to W(\mathfrak g)\) is generically
  smooth.

  \item[(iv)]\SGARefTarget{sga3:xiii:theorem:5:1:condition:iv}
  \SGARefLink{sga3:xiii:theorem:5:1:condition:iii}{The preceding morphism} \(\varphi\) (or, equivalently,
  \(\psi:X\to W(\mathfrak g)\)) is dominant, and \(\mathfrak h\) has
  the same nilpotent rank as \(\mathfrak g\).

  \item[(v)]\SGARefTarget{sga3:xiii:theorem:5:1:condition:v}
  The morphism \(\psi:X\to W(\mathfrak g)\) is generically smooth and
  \(\mathfrak h=\mathfrak n\).

  \item[(vi)]\SGARefTarget{sga3:xiii:theorem:5:1:condition:vi}
  The morphism \(\psi:X\to W(\mathfrak g)\) is dominant, and
  \(\mathfrak h=\mathfrak n\).

  \item[(vii)]\SGARefTarget{sga3:xiii:theorem:5:1:condition:vii}
  The morphism \(\psi:X\to W(\mathfrak g)\) is dominant.

  \item[(viii)]\SGARefTarget{sga3:xiii:theorem:5:1:condition:viii}
  The morphism \(\psi:X\to W(\mathfrak g)\) is dominant, and \(N\) is
  smooth.

  \item[(ix)]\SGARefTarget{sga3:xiii:theorem:5:1:condition:ix}
  The morphism \(\psi:X\to W(\mathfrak g)\) is dominant,
  \(\mathfrak h=\mathfrak n\), and \(\mathfrak h\) is the Lie algebra
  of a smooth algebraic subgroup \(H\) of \(G\).

  \item[(x)]\SGARefTarget{sga3:xiii:theorem:5:1:condition:x}
  The morphism \(\psi:X\to W(\mathfrak g)\) is generically
  quasi-finite, and \(\mathfrak n=\mathfrak h\).

  \item[(xi)]\SGARefTarget{sga3:xiii:theorem:5:1:condition:xi}
  The morphism \(\psi:X\to W(\mathfrak g)\) is generically étale, and
  \(\mathfrak n=\mathfrak h\).

  \item[(xii)]\SGARefTarget{sga3:xiii:theorem:5:1:condition:xii}
  There exists a smooth algebraic subgroup \(H\) of \(G\), with Lie
  algebra \(\mathfrak h\), such that \(\mathfrak h\) contains a Cartan
  subalgebra \(\mathfrak d\) of \(\mathfrak g\).

  \item[(xiii)]\SGARefTarget{sga3:xiii:theorem:5:1:condition:xiii}
  There exists \(a\in\mathfrak h\) such that \(N\) and the transporter
  \(M_a\) of \(a\) into \(\mathfrak h\) (cf.\ \SGARefLink{sga3:xiii:section:1}{Section~1}) coincide in a
  neighborhood of \(e\), and \(\mathfrak n=\mathfrak h\).
\end{enumerate}
One then has \SGARefLink{sga3:xiii:theorem:5:1:diagram:implications}{the following diagram} of implications:
\SGARefTarget{sga3:xiii:theorem:5:1:diagram:implications}
\[
\begin{tikzcd}[column sep=1.45em,row sep=3.2em]
\SGARefLink{sga3:xiii:theorem:5:1:condition:xi}{(\mathrm{xi})}
  \arrow[rr,Leftrightarrow] &&
\SGARefLink{sga3:xiii:theorem:5:1:condition:xii}{(\mathrm{xii})}
  \arrow[r,Leftrightarrow] &
\SGARefLink{sga3:xiii:theorem:5:1:condition:xiii}{(\mathrm{xiii})}
  \arrow[r,Rightarrow] &
\SGARefLink{sga3:xiii:theorem:5:1:condition:viii}{(\mathrm{viii})}
  \arrow[r,Leftrightarrow] &
\SGARefLink{sga3:xiii:theorem:5:1:condition:ix}{(\mathrm{ix})}
  \arrow[r,Leftrightarrow] &
\SGARefLink{sga3:xiii:theorem:5:1:condition:x}{(\mathrm{x})} \\
\SGARefLink{sga3:xiii:theorem:5:1:condition:i}{(\mathrm{i})}
  \arrow[r,Leftrightarrow] &
\SGARefLink{sga3:xiii:theorem:5:1:condition:ii}{(\mathrm{ii})}
  \arrow[r,Leftrightarrow] &
\SGARefLink{sga3:xiii:theorem:5:1:condition:iii}{(\mathrm{iii})}
  \arrow[u,Leftarrow]
  \arrow[r,Leftrightarrow] &
\SGARefLink{sga3:xiii:theorem:5:1:condition:iv}{(\mathrm{iv})}
  \arrow[r,Rightarrow] &
\SGARefLink{sga3:xiii:theorem:5:1:condition:v}{(\mathrm{v})}
  \arrow[r,Rightarrow] &
\SGARefLink{sga3:xiii:theorem:5:1:condition:vi}{(\mathrm{vi})}
  \arrow[u,Leftarrow]
  \arrow[r,Rightarrow] &
\SGARefLink{sga3:xiii:theorem:5:1:condition:vii}{(\mathrm{vii})}.
\end{tikzcd}
\]
When \(k\) has characteristic zero, all the conditions under
consideration are equivalent. Finally,
\[
  \SGARefLink{sga3:xiii:theorem:5:1:condition:xi}{(\mathrm{xi})}
  \Longleftrightarrow
  [\SGARefLink{sga3:xiii:theorem:5:1:condition:i}{(\mathrm{i})}\ \text{and}\ \SGARefLink{sga3:xiii:theorem:5:1:condition:viii}{(\mathrm{viii})}]
  \Longleftrightarrow
  [\SGARefLink{sga3:xiii:theorem:5:1:condition:v}{(\mathrm{v})}\ \text{and}\ \SGARefLink{sga3:xiii:theorem:5:1:condition:viii}{(\mathrm{viii})}].
\]
\end{theorem}

Let us first record the trivial implications
\[
  \SGARefLink{sga3:xiii:theorem:5:1:condition:v}{(\mathrm{v})}\Rightarrow\SGARefLink{sga3:xiii:theorem:5:1:condition:vi}{(\mathrm{vi})}\Rightarrow\SGARefLink{sga3:xiii:theorem:5:1:condition:vii}{(\mathrm{vii})},
  \qquad
  \SGARefLink{sga3:xiii:theorem:5:1:condition:ix}{(\mathrm{ix})}\Rightarrow\SGARefLink{sga3:xiii:theorem:5:1:condition:vi}{(\mathrm{vi})},
  \qquad
  \SGARefLink{sga3:xiii:theorem:5:1:condition:xi}{(\mathrm{xi})}
  \Longleftrightarrow
  [\SGARefLink{sga3:xiii:theorem:5:1:condition:x}{(\mathrm{x})}\ \text{and}\ \SGARefLink{sga3:xiii:theorem:5:1:condition:v}{(\mathrm{v})}].
\]

We prove the equivalence of conditions \SGARefLink{sga3:xiii:theorem:5:1:condition:i}{(i)} through \SGARefLink{sga3:xiii:theorem:5:1:condition:iv}{(iv)}, and that they
imply \SGARefLink{sga3:xiii:theorem:5:1:condition:v}{(v)}. The implication \SGARefLink{sga3:xiii:theorem:5:1:condition:i}{(i)}\(\Rightarrow\)\SGARefLink{sga3:xiii:theorem:5:1:condition:ii}{(ii)} is immediate. On the
other hand, when \(k\) is algebraically closed, condition \SGARefLink{sga3:xiii:theorem:5:1:condition:iii}{(iii)} says
that there is a \(k\)-rational point of \(G\times W(\mathfrak h)\) at
which the tangent map of \(\varphi\) is surjective. It is clear that,
after translating by a suitable element of \(G(k)\), this point may be
taken to have the form \((e,a)\), with \(a\in\mathfrak h\). It follows
that when \(k\) is infinite, though not necessarily algebraically
closed, \SGARefLink{sga3:xiii:theorem:5:1:condition:iii}{this condition}---which is plainly sufficient for generic
smoothness---is also necessary.

The tangent map is readily computed. Upon identifying the tangent
space of \(W(\mathfrak h)\) at \(a\) with \(\mathfrak h\), it is the
map
\[
  (\xi,x)\longmapsto[\xi,a]+x
\]
from \(\mathfrak g\times\mathfrak h\) to \(\mathfrak g\). Its
surjectivity is equivalent to the surjectivity of
\(\operatorname{ad}(a)_{\mathfrak g/\mathfrak h}\), or, what amounts
to the same thing, to its injectivity. This proves the equivalence of
\SGARefLink{sga3:xiii:theorem:5:1:condition:ii}{(ii)} and \SGARefLink{sga3:xiii:theorem:5:1:condition:iii}{(iii)}. Moreover, \SGARefLink{sga3:xiii:theorem:5:1:condition:ii}{(ii)} plainly implies
\(\mathfrak h=\mathfrak n\), and \SGARefLink{sga3:xiii:theorem:5:1:condition:iii}{(iii)} implies that \(\psi\) is
generically smooth: if \(\varphi\) is smooth at a point \(u\), then,
since \(q\) is flat, \(\psi\) is smooth at \(q(u)\). Thus \SGARefLink{sga3:xiii:theorem:5:1:condition:ii}{(ii)} and
\SGARefLink{sga3:xiii:theorem:5:1:condition:iii}{(iii)} imply \SGARefLink{sga3:xiii:theorem:5:1:condition:v}{(v)}.

We now show that they imply \SGARefLink{sga3:xiii:theorem:5:1:condition:i}{(i)}. Since \(\psi\) is dominant and the set
of regular points of \(\mathfrak g\) is dense and open, \(\mathfrak h\)
contains regular elements of \(\mathfrak g\). Hence a sufficiently
general element \(b\) of \(\mathfrak h\) is regular in \(\mathfrak g\)
and makes \(\operatorname{ad}(b)_{\mathfrak g/\mathfrak h}\)
injective. Consequently
\(\operatorname{Nil}(b,\mathfrak g)\subset\mathfrak h\), so
\(\mathfrak h\) contains the Cartan subalgebra
\(\mathfrak d=\operatorname{Nil}(b,\mathfrak g)\). Conditions \SGARefLink{sga3:xiii:theorem:5:1:condition:i}{(i)},
\SGARefLink{sga3:xiii:theorem:5:1:condition:ii}{(ii)}, and \SGARefLink{sga3:xiii:theorem:5:1:condition:iii}{(iii)} are therefore equivalent. Finally,
\SGARefLink{sga3:xiii:theorem:5:1:condition:i}{(i)}\(\Longleftrightarrow\)\SGARefLink{sga3:xiii:theorem:5:1:condition:iv}{(iv)}: we have already observed that if
\(\mathfrak h\) contains a Cartan subalgebra, then it has the same
nilpotent rank as \(\mathfrak g\) (\SGARefLink{sga3:xiii:proposition:4:6}{4.6}), so \SGARefLink{sga3:xiii:theorem:5:1:condition:i}{(i)} implies \SGARefLink{sga3:xiii:theorem:5:1:condition:iv}{(iv)}.
Conversely, if \SGARefLink{sga3:xiii:theorem:5:1:condition:iv}{(iv)} holds, then \(\mathfrak h\) contains a regular
element of \(\mathfrak g\), and since it has the same rank as
\(\mathfrak g\), it contains a Cartan subalgebra by \SGARefLink{sga3:xiii:proposition:4:6}{4.6}.

We next prove the equivalence of conditions \SGARefLink{sga3:xiii:theorem:5:1:condition:viii}{(viii)} through \SGARefLink{sga3:xiii:theorem:5:1:condition:x}{(x)}. We
first note \SGARefLink{sga3:xiii:lemma:5:2}{the following facts}.

\begin{lemma}
\SGARefTarget{sga3:xiii:lemma:5:2}
\begin{enumerate}
  \item[a)]\SGARefTarget{sga3:xiii:lemma:5:2:part:a}
  If \(N\) is smooth, then
  \[
    \dim X\leq\dim G,
  \]
  with equality if and only if \(\mathfrak h=\mathfrak n\).

  \item[b)]\SGARefTarget{sga3:xiii:lemma:5:2:part:b}
  If \(\mathfrak h=\mathfrak n\), then
  \[
    \dim X\geq\dim G,
  \]
  with equality if and only if \(N\) is smooth.
\end{enumerate}
\end{lemma}

These assertions follow at once from \SGARefLink{sga3:xiii:lemma:5:2:display:dimension-formula}{the formula}
\SGARefTarget{sga3:xiii:lemma:5:2:display:dimension-formula}
\[
  \dim X=\dim G-\dim N+\dim_k\mathfrak h
\]
and from the fact that
\(\dim N=\dim_k\mathfrak n\) is equivalent to the smoothness of \(N\).

Now \SGARefLink{sga3:xiii:theorem:5:1:condition:viii}{(viii)} implies \SGARefLink{sga3:xiii:theorem:5:1:condition:x}{(x)}: by dominance, \SGARefLink{sga3:xiii:theorem:5:1:condition:viii}{(viii)} gives
\(\dim X\geq\dim G\), so \SGARefLink{sga3:xiii:lemma:5:2:part:a}{Lemma~5.2(a)} gives equality and
\(\mathfrak h=\mathfrak n\), hence \SGARefLink{sga3:xiii:theorem:5:1:condition:x}{(x)}. The same argument in reverse,
using \SGARefLink{sga3:xiii:lemma:5:2:part:b}{Lemma~5.2(b)}, shows that \SGARefLink{sga3:xiii:theorem:5:1:condition:x}{(x)} implies \SGARefLink{sga3:xiii:theorem:5:1:condition:viii}{(viii)}. Moreover, \SGARefLink{sga3:xiii:theorem:5:1:condition:viii}{(viii)}
implies \SGARefLink{sga3:xiii:theorem:5:1:condition:ix}{(ix)}, since it implies \(\mathfrak h=\mathfrak n\), so
\(\mathfrak h\) is the Lie algebra of the smooth algebraic subgroup
\(N\) of \(G\). Conversely, \SGARefLink{sga3:xiii:theorem:5:1:condition:ix}{(ix)} implies \SGARefLink{sga3:xiii:theorem:5:1:condition:viii}{(viii)}: since \(H\)
normalizes its Lie algebra \(\mathfrak h\), it is contained in \(N\);
and since \(H\) is smooth and has the same Lie algebra as \(N\), it
follows that \(N\) is smooth.

Finally, we prove the equivalence of \SGARefLink{sga3:xiii:theorem:5:1:condition:xi}{(xi)}, \SGARefLink{sga3:xiii:theorem:5:1:condition:xii}{(xii)}, and \SGARefLink{sga3:xiii:theorem:5:1:condition:xiii}{(xiii)}, and show
that they imply \SGARefLink{sga3:xiii:theorem:5:1:condition:iii}{(iii)}; this will complete \SGARefLink{sga3:xiii:theorem:5:1:diagram:implications}{the implication diagram}.
We have \SGARefLink{sga3:xiii:theorem:5:1:condition:xi}{(xi)}\(\Longleftrightarrow\)\SGARefLink{sga3:xiii:theorem:5:1:condition:xiii}{(xiii)}. Indeed, if
\(\mathfrak n=\mathfrak h\), then \SGARefLink{sga3:xiii:lemma:5:2:part:b}{Lemma~5.2(b)} gives
\(\dim X\geq\dim G\). Since \(W(\mathfrak g)\) is normal, condition
\SGARefLink{sga3:xiii:theorem:5:1:condition:xi}{(xi)} is therefore equivalent to \(\psi\) being generically
unramified, and this in turn is equivalent to \SGARefLink{sga3:xiii:theorem:5:1:condition:xiii}{(xiii)} by
\SGARefLink{sga3:xiii:lemma:1:1}{1.1}, \SGARefLink{sga3:xiii:theorem:5:1:condition:ii}{(ii)}\(\Longleftrightarrow\)\SGARefLink{sga3:xiii:theorem:5:1:condition:iii}{(iii)}, arguing as above in the proof
of \SGARefLink{sga3:xiii:theorem:5:1:condition:ii}{(ii)}\(\Longleftrightarrow\)\SGARefLink{sga3:xiii:theorem:5:1:condition:iii}{(iii)}.

Since \SGARefLink{sga3:xiii:theorem:5:1:condition:xi}{(xi)}\(\Rightarrow\)\SGARefLink{sga3:xiii:theorem:5:1:condition:x}{(x)}\(\Rightarrow\)\SGARefLink{sga3:xiii:theorem:5:1:condition:viii}{(viii)}, what we have
already proved shows that \SGARefLink{sga3:xiii:theorem:5:1:condition:xi}{(xi)} implies that \(N\) is smooth. Thus
\(q:G\times W(\mathfrak h)\to X\) is smooth, and the composite
\(\varphi=\psi\circ q\) is generically smooth; this is \SGARefLink{sga3:xiii:theorem:5:1:condition:iii}{(iii)}. Since
\SGARefLink{sga3:xiii:theorem:5:1:condition:iii}{(iii)} implies \SGARefLink{sga3:xiii:theorem:5:1:condition:i}{(i)}, it also follows that
\SGARefLink{sga3:xiii:theorem:5:1:condition:xi}{(xi)}\(\Rightarrow\)\SGARefLink{sga3:xiii:theorem:5:1:condition:xii}{(xii)}.\footnote{Editor's note: because \(H=N\)
will do.} Finally, \SGARefLink{sga3:xiii:theorem:5:1:condition:xii}{(xii)} implies \SGARefLink{sga3:xiii:theorem:5:1:condition:xi}{(xi)}. Indeed, \SGARefLink{sga3:xiii:theorem:5:1:condition:xii}{(xii)} plainly implies
\SGARefLink{sga3:xiii:theorem:5:1:condition:i}{(i)}, and we have seen that
\SGARefLink{sga3:xiii:theorem:5:1:condition:i}{(i)}\(\Rightarrow\)\SGARefLink{sga3:xiii:theorem:5:1:condition:iii}{(iii)}\(\Rightarrow\)\SGARefLink{sga3:xiii:theorem:5:1:condition:v}{(v)}; hence \SGARefLink{sga3:xiii:theorem:5:1:condition:xii}{(xii)} implies \SGARefLink{sga3:xiii:theorem:5:1:condition:v}{(v)}.
It also implies \SGARefLink{sga3:xiii:theorem:5:1:condition:ix}{(ix)}, and \SGARefLink{sga3:xiii:theorem:5:1:condition:ix}{(ix)} implies \SGARefLink{sga3:xiii:theorem:5:1:condition:x}{(x)}. Thus \SGARefLink{sga3:xiii:theorem:5:1:condition:xii}{(xii)} implies both
\SGARefLink{sga3:xiii:theorem:5:1:condition:v}{(v)} and \SGARefLink{sga3:xiii:theorem:5:1:condition:x}{(x)}, and therefore \SGARefLink{sga3:xiii:theorem:5:1:condition:xi}{(xi)}, since
\[
  \text{generically étale}
  =
  \text{generically smooth}
  +
  \text{generically quasi-finite}.
\]

If \(k\) has characteristic zero, then \SGARefLink{sga3:xiii:theorem:5:1:condition:vii}{(vii)} implies \SGARefLink{sga3:xiii:theorem:5:1:condition:viii}{(viii)}, because
Cartier's theorem makes \(N\) automatically smooth
(VI\textsubscript{B}~1.6.1). Moreover, [\SGARefLink{sga3:xiii:theorem:5:1:condition:viii}{(viii)} and \SGARefLink{sga3:xiii:theorem:5:1:condition:x}{(x)}] implies \SGARefLink{sga3:xiii:theorem:5:1:condition:xi}{(xi)},
because in characteristic zero, for a morphism of integral
preschemes,
\[
  \text{generically étale}
  =
  \text{dominant and generically quasi-finite}.
\]
Thus in this case all conditions \SGARefLink{sga3:xiii:theorem:5:1:condition:i}{(i)} through \SGARefLink{sga3:xiii:theorem:5:1:condition:xiii}{(xiii)} are equivalent.

\begin{corollary}
\SGARefTarget{sga3:xiii:corollary:5:3}
Under the equivalent conditions \SGARefLink{sga3:xiii:theorem:5:1:condition:viii}{(viii)} through \SGARefLink{sga3:xiii:theorem:5:1:condition:x}{(x)}, there exists a
unique smooth connected algebraic subgroup \(H\) of \(G\) whose Lie
algebra is \(\mathfrak h\), and
\[
  \operatorname{Norm}_G(H)
  =\operatorname{Norm}_G(\mathfrak h)
  =N,
  \qquad
  H=N^0.
\]
\end{corollary}

Indeed, \(H=N^0\) has the required properties. Conversely, if \(H\)
has these properties, then \(H\subset N\), since \(H\) normalizes its
Lie algebra \(\mathfrak h\). This is an inclusion of smooth groups
with the same Lie algebra, and \(H\) is connected, so \(H=N^0\). To
prove
\(\operatorname{Norm}_G(H)=\operatorname{Norm}_G(\mathfrak h)\), we
may assume that \(k\) is algebraically closed. What we have just proved
then shows immediately that the two groups have the same \(k\)-points.
On the other hand, the inclusions
\[
  H\subset\operatorname{Norm}_G(H)\subset N
\]
show that \(\operatorname{Norm}_G(H)\) and \(N\) have the same Lie
algebra; hence they are equal.

\begin{corollary}
\SGARefTarget{sga3:xiii:corollary:5:4}
Under the equivalent conditions \SGARefLink{sga3:xiii:theorem:5:1:condition:i}{(i)} through \SGARefLink{sga3:xiii:theorem:5:1:condition:iv}{(iv)}, let
\(a\in\mathfrak h\). Then the following conditions are equivalent, and
they hold whenever \(a\) is regular in \(\mathfrak g\):
\begin{enumerate}
  \item[(i)]\SGARefTarget{sga3:xiii:corollary:5:4:condition:i}
  \(\varphi\) is smooth at \((e,a)\).

  \item[(ii)]\SGARefTarget{sga3:xiii:corollary:5:4:condition:ii}
  \(M_a=\operatorname{Transp}_G(a,\mathfrak h)\) is smooth at \(e\),
  and
  \(\dim_e(M_a)=\operatorname{rank}_k\mathfrak h\).

  \item[(iii)]\SGARefTarget{sga3:xiii:corollary:5:4:condition:iii}
  \(\operatorname{ad}(a)_{\mathfrak g/\mathfrak h}\) is injective
  (equivalently, bijective).
\end{enumerate}
Under the equivalent conditions \SGARefLink{sga3:xiii:theorem:5:1:condition:xi}{(xi)} through \SGARefLink{sga3:xiii:theorem:5:1:condition:xiii}{(xiii)}, let \(H\) be the
algebraic subgroup of \(G\) considered in \SGARefLink{sga3:xiii:corollary:5:3}{Corollary~5.3}. The preceding
conditions are also equivalent to:
\begin{enumerate}
  \item[(iv)]\SGARefTarget{sga3:xiii:corollary:5:4:condition:iv}
  \(\psi\) is étale at \((\overline e,a)\).

  \item[(v)]\SGARefTarget{sga3:xiii:corollary:5:4:condition:v}
  If \(M_a^0\) denotes the connected component of \(e\) in the
  transporter \(M_a\) of \(a\) into \(\mathfrak h\), endowed with the
  structure induced by \(M_a\), then
  \[
    H=M_a^0.
  \]
\end{enumerate}
\end{corollary}

Plainly \SGARefLink{sga3:xiii:corollary:5:4:condition:i}{(i)} implies \SGARefLink{sga3:xiii:corollary:5:4:condition:ii}{(ii)}, since \(M_a\) is isomorphic to the fiber
\(\varphi^{-1}(a)\), with \(e\) corresponding to \((e,a)\).
Conversely, \SGARefLink{sga3:xiii:corollary:5:4:condition:ii}{(ii)} implies \SGARefLink{sga3:xiii:corollary:5:4:condition:i}{(i)}: condition \SGARefLink{sga3:xiii:corollary:5:4:condition:ii}{(ii)} says that \(\varphi\) is
``equidimensional'' at \((e,a)\), in the sense that the dimension of
the fiber through this point is the dimension of the generic fiber.
Since \(G\times W(\mathfrak h)\) and \(W(\mathfrak g)\) are regular,
this implies that \(\varphi\) is flat at \((e,a)\), hence smooth there
because its fiber is smooth at that point. The equivalence of \SGARefLink{sga3:xiii:corollary:5:4:condition:i}{(i)} and
\SGARefLink{sga3:xiii:corollary:5:4:condition:iii}{(iii)} was obtained in the proof of \SGARefLink{sga3:xiii:theorem:5:1}{Theorem~5.1} by the elementary
calculation of the tangent map. We also saw in \SGARefLink{sga3:xiii:proposition:4:8:item:b}{4.8(b)} that regularity
of \(a\) in \(\mathfrak g\) implies \SGARefLink{sga3:xiii:corollary:5:4:condition:iii}{(iii)}.

Under conditions \SGARefLink{sga3:xiii:theorem:5:1:condition:xi}{(xi)} through \SGARefLink{sga3:xiii:theorem:5:1:condition:xiii}{(xiii)}, the morphism
\(q:G\times W(\mathfrak h)\to X\) is smooth, since \(N\) is smooth.
Thus \SGARefLink{sga3:xiii:corollary:5:4:condition:i}{(i)} is equivalent to the smoothness of \(\psi\) at
\((\overline e,a)\); and since \(\psi\) is generically étale, this is
equivalent to \SGARefLink{sga3:xiii:corollary:5:4:condition:iv}{(iv)}. Finally, as was noted at the end of the proof of
\SGARefLink{sga3:xiii:lemma:1:1}{1.1}, condition \SGARefLink{sga3:xiii:corollary:5:4:condition:iv}{(iv)} implies that \(N\) is the prescheme induced on
\(M_a\) by an open and closed subset of \(M_a\), which gives \SGARefLink{sga3:xiii:corollary:5:4:condition:v}{(v)}.
Conversely, \SGARefLink{sga3:xiii:corollary:5:4:condition:v}{(v)} trivially implies \SGARefLink{sga3:xiii:corollary:5:4:condition:ii}{(ii)}. Alternatively, \SGARefLink{sga3:xiii:corollary:5:4:condition:v}{(v)} implies
\SGARefLink{sga3:xiii:corollary:5:4:condition:iv}{(iv)} by \SGARefLink{sga3:xiii:lemma:1:1}{1.1}, since \(\psi\) is dominant and \(W(\mathfrak g)\) is
normal, and ``unramified'' is equivalent here to ``étale''. This
completes the proof of \SGARefLink{sga3:xiii:corollary:5:4}{Corollary~5.4}.

\begin{corollary}
\SGARefTarget{sga3:xiii:corollary:5:5}
Let \(G\) be a smooth algebraic group over a field \(k\), and let
\(H\) be a smooth algebraic subgroup such that its Lie algebra
\(\mathfrak h\) contains, at least after a suitable extension of
\(k\), a Cartan subalgebra of the Lie algebra \(\mathfrak g\) of
\(G\). Let \(K\) be a connected algebraic subgroup of \(G\), not
necessarily smooth, with Lie algebra \(\mathfrak k\), and suppose that
\(\mathfrak k\) contains a regular element \(a\) of \(\mathfrak g\),
again at least after a suitable extension of \(k\). Then \(H\)
contains \(K\) if and only if \(\mathfrak h\) contains
\(\mathfrak k\).
\end{corollary}

Indeed, \SGARefLink{sga3:xiii:corollary:5:4}{Corollary~5.4}, \SGARefLink{sga3:xiii:corollary:5:4:condition:iii}{(iii)}\(\Rightarrow\)\SGARefLink{sga3:xiii:corollary:5:4:condition:v}{(v)}, applied to the
connected subgroup supplied by \SGARefLink{sga3:xiii:corollary:5:3}{Corollary~5.3}, gives
\(H^0=M_a^0\).\footnote{The French source prints \(H=M_a^0\) here. In the stated
generality the argument gives \(H^0=M_a^0\), which is all that is
needed.}
This identity is invariant under extension of the base field, so it
remains valid without assuming that the base field is infinite.
Moreover, \(\mathfrak k\subset\mathfrak h\) plainly implies
\(K\subset M_a\). Since \(K\) is connected,
\(K\subset M_a^0\subset H\). \hfill\(\square\)

\begin{corollary}
\SGARefTarget{sga3:xiii:corollary:5:6}
Let \(G\) and \(H\) be as in \SGARefLink{sga3:xiii:corollary:5:5}{Corollary~5.5}, and let \(K\) be an
algebraic subgroup of \(G\). Suppose that \(H\) is connected and \(K\)
is smooth. Then \(K\) contains \(H\) if and only if
\(\mathfrak k\) contains \(\mathfrak h\).
\end{corollary}

Indeed, if \(\mathfrak k\supset\mathfrak h\), then \(K\) satisfies,
with respect to \(H\), the hypothesis imposed on \(H\) in
\SGARefLink{sga3:xiii:corollary:5:5}{Corollary~5.5}. Meanwhile \(H\) plainly satisfies, with respect to
\(K\), the condition imposed on \(K\) there. The conclusion follows
from \SGARefLink{sga3:xiii:corollary:5:5}{Corollary~5.5}.

\begin{corollary}
\SGARefTarget{sga3:xiii:corollary:5:7}
Let \(\mathfrak g\) be the Lie algebra of a smooth algebraic group
\(G\) over a field \(k\). Then:
\begin{enumerate}
  \item[a)]\SGARefTarget{sga3:xiii:corollary:5:7:part:a}
  A Lie subalgebra \(\mathfrak d\) of \(\mathfrak g\) is a Cartan
  subalgebra if and only if it is nilpotent and equal to its own
  normalizer.

  \item[b)]\SGARefTarget{sga3:xiii:corollary:5:7:part:b}
  An element \(a\in\mathfrak g\) is regular if and only if
  \(\operatorname{Nil}(a,\mathfrak g)\) is nilpotent.
\end{enumerate}
\end{corollary}

After extending the base field, we may assume that \(k\) is infinite.
In view of \SGARefLink{sga3:xiii:proposition:4:4}{4.4}, to prove \SGARefLink{sga3:xiii:corollary:5:7:part:a}{(a)} it remains only to show that if
\(\mathfrak d\) is nilpotent and contains an element \(a\) for which
\(\operatorname{ad}(a)_{\mathfrak g/\mathfrak d}\) is injective, then
\(\mathfrak d\) is a Cartan subalgebra. By \SGARefLink{sga3:xiii:theorem:5:1}{Theorem~5.1},
\SGARefLink{sga3:xiii:theorem:5:1:condition:ii}{(ii)}\(\Rightarrow\)\SGARefLink{sga3:xiii:theorem:5:1:condition:i}{(i)}, \(\mathfrak d\) contains a Cartan subalgebra
\[
  \mathfrak d'=\operatorname{Nil}(a,\mathfrak g).
\]
Since \(\mathfrak d\) is nilpotent, \SGARefLink{sga3:xiii:corollary:4:3}{4.3} now gives
\(\mathfrak d=\mathfrak d'\). For \SGARefLink{sga3:xiii:corollary:5:7:part:b}{(b)}, one notes that
\(\operatorname{Nil}(a,\mathfrak g)\) is a Cartan subalgebra of
\(\mathfrak g\) by \SGARefLink{sga3:xiii:corollary:5:7:part:a}{(a)}, and therefore \(a\) is regular.

\begin{corollary}
\SGARefTarget{sga3:xiii:corollary:5:8}
Let \(G\) be a smooth algebraic group over a field \(k\), and let
\(H\) be a smooth algebraic subgroup, with Lie algebras
\(\mathfrak g\) and \(\mathfrak h\), respectively. Suppose that after
a suitable extension of the base field, \(\mathfrak h\) contains a
Cartan subalgebra of \(\mathfrak g\). Let \(\mathfrak d\) be a Lie
subalgebra of \(\mathfrak h\). Then \(\mathfrak d\) is a Cartan
subalgebra of \(\mathfrak h\) if and only if it is a Cartan subalgebra
of \(\mathfrak g\).
\end{corollary}

In view of \SGARefLink{sga3:xiii:proposition:4:8:item:c}{4.8(c)}, it remains to show that if \(\mathfrak d\) is a
Cartan subalgebra of \(\mathfrak h\), then it is one of
\(\mathfrak g\). It is enough to show that \(\mathfrak d\) contains
an element \(a\) regular in \(\mathfrak g\), and we may assume that
\(k\) is algebraically closed. But there is a dense open subset of
\(\mathfrak h\) consisting of points regular in \(\mathfrak g\), so
the assertion follows from \SGARefLink{sga3:xiii:theorem:5:1}{Theorem~5.1},
\SGARefLink{sga3:xiii:theorem:5:1:condition:i}{(i)}\(\Rightarrow\)\SGARefLink{sga3:xiii:theorem:5:1:condition:vii}{(vii)}, applied to
\((\mathfrak h,\mathfrak d)\).

\section{Cartan subalgebras and subgroups of type \texorpdfstring{\((C)\)}{(C)}
associated with a smooth algebraic group}
\SGARefTarget{sga3:xiii:section:6}

For simplicity, in \SGARefLink{sga3:xiii:theorem:6:1}{the following theorem} we restrict ourselves to an
algebraically closed base field; the case of an arbitrary base
prescheme will be treated in the next Exposé.

\begin{theorem}
\SGARefTarget{sga3:xiii:theorem:6:1}
Let \(G\) be a smooth algebraic group over an algebraically closed
field \(k\), and let \(\mathfrak g\) be its Lie algebra. Then:
\begin{enumerate}
  \item[a)]\SGARefTarget{sga3:xiii:theorem:6:1:part:a}
  The Cartan subalgebras of \(\mathfrak g\) are conjugate.

  \item[b)]\SGARefTarget{sga3:xiii:theorem:6:1:part:b}
  Let \(\mathfrak d\) be a Cartan subalgebra of \(\mathfrak g\).
  Its normalizer \(N\) in \(G\) is smooth, and \(D=N^0\) is the
  unique smooth connected subgroup of \(G\) whose Lie algebra is
  \(\mathfrak d\). One has
  \[
    \operatorname{Norm}_G(\mathfrak d)
    =\operatorname{Norm}_G(D)
    =N,
    \qquad\text{and hence}\qquad
    D=\operatorname{Norm}_G(D)^0.
  \]

  \item[c)]\SGARefTarget{sga3:xiii:theorem:6:1:part:c}
  With \(\mathfrak d\) as in \SGARefLink{sga3:xiii:theorem:6:1:part:b}{(b)}, set, as in \SGARefLink{sga3:xiii:section:5}{Section~5},
  \[
    X=G\times^N W(\mathfrak d),
  \]
  and consider the canonical morphism
  \[
    \psi:X\longrightarrow W(\mathfrak g).
  \]
  Its fiber over \(a\in\mathfrak g\) has as its \(k\)-valued points
  the set of Cartan subalgebras of \(\mathfrak g\) containing \(a\).
  Then \(\psi\) is birational.

  \item[d)]\SGARefTarget{sga3:xiii:theorem:6:1:part:d}
  With the notation of \SGARefLink{sga3:xiii:theorem:6:1:part:c}{(c)}, let \(U\) be the largest open subset of
  \(W(\mathfrak g)\) for which \(\psi\) induces an isomorphism
  \[
    \psi^{-1}(U)\xrightarrow{\sim}U.
  \]
  Then, for \(a\in\mathfrak g\), the following conditions are
  equivalent:
  \begin{enumerate}
    \item[(i)]\SGARefTarget{sga3:xiii:theorem:6:1:part:d:condition:i}
    \(a\in U(k)\).

    \item[(ii)]\SGARefTarget{sga3:xiii:theorem:6:1:part:d:condition:ii}
    The element \(a\) belongs to one and only one Cartan subalgebra of
    \(\mathfrak g\).

    \item[(iii)]\SGARefTarget{sga3:xiii:theorem:6:1:part:d:condition:iii}
    The set of Cartan subalgebras of \(\mathfrak g\) containing \(a\)
    is finite and nonempty.

    \item[(iv)]\SGARefTarget{sga3:xiii:theorem:6:1:part:d:condition:iv}
    The fiber \(\psi^{-1}(a)\) has an isolated point.

    \item[(v)]\SGARefTarget{sga3:xiii:theorem:6:1:part:d:condition:v}
    If \(a\in\mathfrak d\), the morphism \(\psi\) is étale---or merely
    quasi-finite---at the point \((\overline e,a)\).

    \item[(vi)]\SGARefTarget{sga3:xiii:theorem:6:1:part:d:condition:vi}
    The element \(a\) is regular in \(\mathfrak g\).

    \item[(vii)]\SGARefTarget{sga3:xiii:theorem:6:1:part:d:condition:vii}
    With the notation of \SGARefLink{sga3:xiii:corollary:5:4:condition:iii}{Corollary~5.4(iii)}, one has \(N=M_a\).
  \end{enumerate}
\end{enumerate}
\end{theorem}

\begin{proof}
Apply \SGARefLink{sga3:xiii:theorem:5:1}{Theorem~5.1} with the Cartan subalgebra \(\mathfrak d\) of
\(\mathfrak g\) in place of \(\mathfrak h\). We show that the strongest
conditions, \SGARefLink{sga3:xiii:theorem:5:1:condition:xi}{(xi)} through \SGARefLink{sga3:xiii:theorem:5:1:condition:xiii}{(xiii)}, hold. This is clear either in the
form \SGARefLink{sga3:xiii:theorem:5:1:condition:xiii}{(xiii)}, by \SGARefLink{sga3:xiii:corollary:4:7}{4.7}---which says that if a regular element
\(a\in\mathfrak g\) lies in \(\mathfrak d\), then \(N=M_a\)---or in
the form
\[
  \SGARefLink{sga3:xiii:theorem:5:1:condition:xi}{(\mathrm{xi})}=\SGARefLink{sga3:xiii:theorem:5:1:condition:x}{(\mathrm{x})}+\SGARefLink{sga3:xiii:theorem:5:1:condition:i}{(\mathrm{i})}.
\]
Indeed, \SGARefLink{sga3:xiii:theorem:5:1:condition:i}{(i)} is automatic, while \SGARefLink{sga3:xiii:theorem:5:1:condition:x}{(x)} follows from the fact that a
regular point of \(\mathfrak g\) belongs to a unique Cartan subalgebra
of \(\mathfrak g\) (\SGARefLink{sga3:xiii:proposition:4:6:item:a}{4.6(a)}), and therefore, a fortiori, to a unique
conjugate of \(\mathfrak d\).

Part \SGARefLink{sga3:xiii:theorem:6:1:part:b}{(b)} now follows from \SGARefLink{sga3:xiii:corollary:5:3}{Corollary~5.3}. Part \SGARefLink{sga3:xiii:theorem:6:1:part:c}{(c)} follows because
\(\psi\) is generically étale and because a sufficiently general
point of \(\mathfrak g\)---more precisely, a regular point---belongs
to a unique Cartan subalgebra of \(\mathfrak g\). Under these
conditions, the equivalence of conditions \SGARefLink{sga3:xiii:theorem:6:1:part:d:condition:i}{(i)} through \SGARefLink{sga3:xiii:theorem:6:1:part:d:condition:v}{(v)} on \(a\)
follows immediately from Zariski's Main Theorem, applied to the
separated birational morphism \(\psi\), since \(W(\mathfrak g)\) is
normal and \(X\) is integral. The equivalence of \SGARefLink{sga3:xiii:theorem:6:1:part:d:condition:v}{(v)} and \SGARefLink{sga3:xiii:theorem:6:1:part:d:condition:vi}{(vi)} is the
special case of \SGARefLink{sga3:xiii:corollary:5:4}{Corollary~5.4},
\SGARefLink{sga3:xiii:theorem:6:1:part:d:condition:iii}{(iii)}\(\Longleftrightarrow\)\SGARefLink{sga3:xiii:theorem:6:1:part:d:condition:iv}{(iv)}, obtained after carrying \(a\) into
\(\mathfrak d\) by a suitable element \(g\in G(k)\), in view of
\SGARefLink{sga3:xiii:proposition:4:6:item:b}{4.6(b)}. Moreover, \SGARefLink{sga3:xiii:corollary:5:4}{Corollary~5.4} shows that \SGARefLink{sga3:xiii:theorem:6:1:part:d:condition:v}{(v)} and \SGARefLink{sga3:xiii:theorem:6:1:part:d:condition:vi}{(vi)} are also
equivalent to
\[
  N\supset M_a^0.
\]
By \SGARefLink{sga3:xiii:corollary:4:7}{4.7} this even implies
\(N=M_a\).\footnote{The French source prints \(N=M\) at this point; condition
\SGARefLink{sga3:xiii:theorem:6:1:part:d:condition:vii}{(vii)} and the surrounding notation require \(N=M_a\).}
This proves \SGARefLink{sga3:xiii:theorem:6:1:part:d}{(d)}.

Of course, \SGARefLink{sga3:xiii:theorem:6:1:part:b}{(b)}, \SGARefLink{sga3:xiii:theorem:6:1:part:c}{(c)}, and the equivalence of \SGARefLink{sga3:xiii:theorem:6:1:part:d:condition:i}{(i)}, \SGARefLink{sga3:xiii:theorem:6:1:part:d:condition:iv}{(iv)}, \SGARefLink{sga3:xiii:theorem:6:1:part:d:condition:v}{(v)}, \SGARefLink{sga3:xiii:theorem:6:1:part:d:condition:vi}{(vi)}, and
\SGARefLink{sga3:xiii:theorem:6:1:part:d:condition:vii}{(vii)} remain valid over an arbitrary base field. We now prove \SGARefLink{sga3:xiii:theorem:6:1:part:a}{(a)},
using the assumption that \(k\) is algebraically closed. By
\SGARefLink{sga3:xiii:theorem:5:1}{Theorem~5.1}, \SGARefLink{sga3:xiii:theorem:5:1:condition:i}{(i)}\(\Rightarrow\)\SGARefLink{sga3:xiii:theorem:5:1:condition:vii}{(vii)}, the morphism
\(\psi:X\to W(\mathfrak g)\) is dominant. Hence there is a dense open
subset \(V\) of \(W(\mathfrak g)\) such that every \(a\in V(k)\) is
the image of a point of \(X(k)\), or equivalently, belongs to a
conjugate of \(\mathfrak d\). Apply \SGARefLink{sga3:xiii:theorem:5:1:diagram:implications}{the same result} to a second Cartan
subalgebra \(\mathfrak d'\) of \(\mathfrak g\). We may choose \(V\)
so that every element of \(V(k)\) is conjugate both to an element of
\(\mathfrak d\) and to an element of \(\mathfrak d'\). Choose a
regular element of \(V(k)\). Some conjugate \(\mathfrak d''\) of
\(\mathfrak d'\) then contains an element of \(\mathfrak d\) regular
in \(\mathfrak g\), and \SGARefLink{sga3:xiii:proposition:4:6:item:a}{4.6(a)} gives
\(\mathfrak d''=\mathfrak d\). This proves \SGARefLink{sga3:xiii:theorem:6:1:part:a}{(a)} and completes the proof.
\end{proof}

\begin{definition}
\SGARefTarget{sga3:xiii:definition:6:2}
Let \(G\) be a smooth algebraic group over a field \(k\). A
\emph{subgroup of type \((C)\)} of \(G\) is a smooth connected
algebraic subgroup whose Lie algebra is a Cartan subalgebra of
\(\mathfrak g\). The \emph{infinitesimal rank} of \(G\) is the
nilpotent rank of its Lie algebra \(\mathfrak g\); it is equal to the
dimension of every subgroup of type \((C)\) of \(G\).
\end{definition}

\SGARefLink{sga3:xiii:theorem:6:1:part:b}{Theorem~6.1(b)} immediately gives:

\begin{corollary}
\SGARefTarget{sga3:xiii:corollary:6:3}
Under the hypotheses of \SGARefLink{sga3:xiii:definition:6:2}{Definition~6.2}, the map
\[
  D\longmapsto\mathfrak d=\operatorname{Lie}(D)
\]
establishes a bijection between the subgroups of type \((C)\) of
\(G\) and the Cartan subalgebras \(\mathfrak d\) of \(\mathfrak g\).
If \(D\) is a subgroup of type \((C)\) of \(G\), then \(D\) is its
own connected normalizer:
\[
  D=\operatorname{Norm}_G(D)^0.
\]
\end{corollary}

Combining \SGARefLink{sga3:xiii:theorem:6:1:part:a}{Theorem~6.1(a)} and \SGARefLink{sga3:xiii:corollary:6:3}{Corollary~6.3} gives:

\begin{corollary}
\SGARefTarget{sga3:xiii:corollary:6:4}
If \(k\) is algebraically closed, the subgroups of type \((C)\) of
\(G\) are conjugate to one another.
\end{corollary}

\begin{corollary}
\SGARefTarget{sga3:xiii:corollary:6:5}
Let \(G\) be a smooth algebraic group over an algebraically closed
field \(k\), and let \(H\) be a smooth algebraic subgroup of \(G\),
with Lie algebras \(\mathfrak g\) and \(\mathfrak h\), respectively.
Then \(\mathfrak h\) contains a Cartan subalgebra of \(\mathfrak g\)
if and only if \(H\) contains a subgroup \(D\) of type \((C)\) of
\(G\).
\end{corollary}

The condition is plainly sufficient. It is also necessary, since by
\SGARefLink{sga3:xiii:corollary:5:5}{Corollary~5.5},
\[
  H\supset D
  \quad\Longleftrightarrow\quad
  \mathfrak h\supset\mathfrak d.
\]

\refstepcounter{theorem}
\par\medskip
\noindent\SGARefTarget{sga3:xiii:remarks:6:6}%
\textbf{Remarks \thetheorem.}
\begin{enumerate}
  \item[a)]\SGARefTarget{sga3:xiii:remarks:6:6:part:a}
  The connected subgroups \(H\) of \(G\) described in
  \SGARefLink{sga3:xiii:corollary:6:5}{Corollary~6.5} correspond bijectively to the Lie subalgebras
  \(\mathfrak h\) of \(\mathfrak g\) satisfying the strongest
  condition, \SGARefLink{sga3:xiii:theorem:5:1:condition:xii}{(xii)}, of \SGARefLink{sga3:xiii:theorem:5:1}{Theorem~5.1}. By \SGARefLink{sga3:xiii:corollary:5:5}{Corollary~5.5}, inclusion
  relations among such subgroups can already be read from their Lie
  algebras.

  \item[b)]\SGARefTarget{sga3:xiii:remarks:6:6:part:b}
  Continue to suppose that \(k\) is algebraically closed, and let
  \(D\) be a subgroup of type \((C)\) of \(G\). It is then easy to
  show that \(D\) contains a Cartan subgroup \(C\) of \(G\). Indeed,
  put \(V=\mathfrak g/\mathfrak d\). For a general element
  \(a\in\mathfrak d\), the endomorphism
  \(\operatorname{ad}(a)\) of \(V\) is injective; it suffices to take
  \(a\) regular in \(\mathfrak g\). From this one readily deduces that,
  for a general \(g\in D(k)\), the action of
  \(\operatorname{Ad}(g)\) on \(V\) has no nonzero invariant vector.
  Thus \SGARefLink{sga3:xiii:theorem:2:1}{2.1}, \SGARefLink{sga3:xiii:theorem:2:1:item:vii}{(vii)}\(\Rightarrow\)\SGARefLink{sga3:xiii:theorem:2:1:item:i}{(i)}, applies. In fact, in the next
  Exposé we shall prove the more precise result that every Cartan
  subgroup \(C\) of \(G\) is contained in one and only one subgroup
  of type \((C)\) of \(G\).

  \item[c)]\SGARefTarget{sga3:xiii:remarks:6:6:part:c}
  In general one must not confuse Cartan subgroups with subgroups of
  type \((C)\). The subgroups of type \((C)\) of \(G\) are Cartan
  subgroups if and only if they are nilpotent, since Cartan subgroups
  are maximal nilpotent subgroups. It may happen that \(\mathfrak g\)
  is nilpotent although \(G\) is not; an example is
  \(\operatorname{SL}(2)_k\) when \(k\) has characteristic \(2\).
  Then the Cartan subalgebras of \(\mathfrak g\) are all equal to
  \(\mathfrak g\). Thus, if \(G\) is connected, \(G\) is a subgroup
  of type \((C)\) of itself, but it is not a Cartan subgroup of
  itself.

  On the other hand, we shall see in Exposé~XIV that if \(G\) is an
  adjoint semisimple group, its subgroups of type \((C)\) are precisely
  its Cartan subgroups. The same conclusion holds in characteristic
  zero for every smooth algebraic group \(G\) over \(k\). This follows
  at once from the fact that, in characteristic zero, a connected
  algebraic group is nilpotent if and only if its Lie algebra is
  nilpotent. In turn, this follows from the fact that in
  characteristic zero the center \(Z\) of a connected algebraic group
  \(G\) has Lie algebra equal to the center of \(\mathfrak g\).
  Indeed, one obtains a priori the space of invariants in
  \(\mathfrak g\) under the adjoint representation of \(G\). Since
  \(G\) is connected and \(k\) has characteristic zero, these are also
  the infinitesimal invariants of the adjoint representation of
  \(\mathfrak g\), namely the center of \(\mathfrak g\). Finally,
  Cartier's theorem (cf.\ VI\textsubscript{B}~1.6.1) shows that \(Z\)
  is smooth.
\end{enumerate}

% END INCLUDED SGA3_English_current_buildable_source_20260823/source/inputs/XIII/source/tex/components/04_expose_XIII_sections5_6.tex



% END INCLUDED SGA3_English_current_buildable_source_20260823/source/inputs/XIII/body.tex


\SGAStartExpose{XIV}
\chapter*{Expos\'e XIV\\Regular Elements: Continuation\\
Applications to Algebraic Groups}
\addcontentsline{toc}{chapter}{Expos\'e XIV. Regular Elements: Continuation}
\begin{center}
  \textit{by A. Grothendieck, with an appendix by J.-P. Serre}
\end{center}
% BEGIN INCLUDED SGA3_English_current_buildable_source_20260823/source/inputs/XIV/body.tex
% Mechanically extracted after \mainmatter from the frozen standalone master.

\markboth{Exposé XIV. Regular Elements: Continuation}
  {Exposé XIV. Regular Elements: Continuation}
% BEGIN INCLUDED SGA3_English_current_buildable_source_20260823/source/inputs/XIV/tex/components/00_expose_XIV_section1_opening_en.tex
% Authority: Expo14.pdf, local pp. 1--3 / combined pp. 835--837.
% Continuous English currently runs through the statement of Lemma 2.1.

\SGARefTarget{sga3:xiv:expose}

\setcounter{footnote}{0}
\footnotetext{Version xy of 1 December 2008.}

\section{Construction of Cartan subgroups and maximal tori for a
smooth algebraic group}
\SGARefTarget{sga3:xiv:section:1}

\paragraph{\textbf{Theorem 1.1.}}
\SGARefTarget{sga3:xiv:1.1}
\emph{Let \(G\) be a smooth algebraic group over a field \(k\).
Then \(G\) admits a maximal torus \(T\), and hence a Cartan subgroup
\(C=\Centr_G(T)\).}

By \SGARefLink{sga3:xii:3.2}{XII~3.2}, finding a maximal torus \(T\) of \(G\) is equivalent to
finding a Cartan subgroup \(C\) of \(G\). Moreover, since the maximal
tori of \(G\) are those of \(G^0\), we may suppose that \(G\) is
connected. We distinguish two cases.

\medskip
\noindent\textup{1\textdegree)} \emph{The field \(k\) is finite.}
Let \(\mathcal T\) be the scheme of maximal tori of \(G\)
(\SGARefLink{sga3:xii:1.10}{XII~1.10}); it is smooth over \(k\). The group \(G\) acts on
\(\mathcal T\) through inner automorphisms, and the conjugacy theorem
(\SGARefLink{sga3:xii:6.6:item:a}{XII~6.6(a)}) says that any two \(\bar k\)-rational points of
\(\mathcal T_{\bar k}\) are conjugate under
\(G_{\bar k}(\bar k)\). Since \(\mathcal T\) is smooth over \(k\),
\(\mathcal T_{\bar k}\) is smooth over \(\bar k\); it follows that
\(\mathcal T_{\bar k}\) is isomorphic to
\[
  G_{\bar k}/N,
\]
where \(N\) is the stabilizer of an element of
\(\mathcal T_{\bar k}(\bar k)\), that is, the normalizer of a
maximal torus \(\bar T\) of \(G_{\bar k}\). Consequently
\(\mathcal T\) is a homogeneous space under the action of \(G\).
A well-known theorem of Lang
(\emph{Amer.\ J.\ Math.} \textbf{78} (1956), 555--563) says that
every homogeneous space under a smooth connected algebraic group
over a finite field \(k\) has a rational point. In particular,
\(\mathcal T\) has a rational point; equivalently, \(G\) has a
maximal torus \(T\).

\medskip
\noindent\textup{2\textdegree)} \emph{The field \(k\) is infinite.}
We shall use the following lemma.

\paragraph{\textbf{Lemma 1.2.}}
\SGARefTarget{sga3:xiv:1.2}
\emph{Let \(G\) be a smooth algebraic group over a field \(k\).
Then \(G\) admits a subgroup of type \textup{(C)}
\textup{(\SGARefLink{sga3:xiii:definition:6:2}{XIII~6.2})}.}

By \SGARefLink{sga3:xiii:corollary:6:3}{XIII~6.3}, this is equivalent to saying that \(\mathfrak g\)
contains a Cartan subalgebra \(\mathfrak d\). When \(k\) is infinite
this is immediate: \(\mathfrak g\) contains a regular element \(a\),
and one takes
\[
  \mathfrak d=\Nil(a,\mathfrak g).
\]
The case of finite \(k\) is treated exactly as in the proof of
\textup{1\textdegree)} above, but it requires first constructing the
scheme \(\mathcal D\) of Cartan subalgebras of \(\mathfrak g\) and
proving that this scheme is smooth over \(k\); this will be done
below in \SGARefLink{sga3:xiv:2.16}{2.16}.

To prove \SGARefLink{sga3:xiv:1.1}{1.1} in case \textup{2\textdegree)}, which is the case now
under consideration, it is therefore enough to know \SGARefLink{sga3:xiv:1.2}{1.2} when \(k\)
is infinite. We also record the following result for later use.

\paragraph{\textbf{Lemma 1.3.}}
\SGARefTarget{sga3:xiv:1.3}
\emph{Let \(G\) be a smooth connected affine algebraic group whose
reductive center \textup{(\SGARefLink{sga3:xii:4.1}{XII~4.1} and \SGARefLink{sga3:xiv:4.4}{4.4})} is reduced to the identity
group. Then \(G\) is nilpotent if and only if its Lie algebra
\(\mathfrak g\) is nilpotent.}

This is contained in \SGARefLink{sga3:xii:4.9}{XII~4.9}.

We can now give a procedure for constructing Cartan subgroups of
\(G\). The same procedure applies when \(k\) is finite, provided
\SGARefLink{sga3:xiv:1.2}{Lemma~1.2} is granted in that case. Suppose first that \(G\) is
affine. We argue by induction on \(n=\dim G\), the assertion being
trivial when \(n=0\). Suppose therefore that \(n>0\), and that the
assertion has been proved in dimensions \(n'<n\). Let \(Z\) be the
reductive center of \(G\), and let
\[
  u\colon G\longrightarrow G'=G/Z
\]
be the canonical homomorphism. By \SGARefLink{sga3:xii:theorem:4-7:item:c}{XII~4.7(c)}, the rule
\[
  C'\longmapsto C=u^{-1}(C')
\]
is a bijection from the Cartan subgroups of \(G'\) to those of
\(G\). Replacing \(G\) by \(G'\), we may thus suppose that the
reductive center \(Z\) of \(G\) is reduced to the identity element
\textup{(the same is true for the reductive center of \(G'\), by
\SGARefLink{sga3:xii:theorem:4-7:item:b}{XII~4.7(b)})}.

By \SGARefLink{sga3:xiv:1.2}{Lemma~1.2},\footnote{The source prints ``\SGARefLink{sga3:xiv:1.3}{1.3}'' here; the cited
existence result is \SGARefLink{sga3:xiv:1.2}{Lemma~1.2}.} \(G\) admits a subgroup \(D\) of
type \textup{(C)}. Over the algebraic closure \(\bar k\) of \(k\),
the group \(D_{\bar k}\) contains a Cartan subgroup of
\(G_{\bar k}\) \textup{(\SGARefLink{sga3:xiii:remarks:6:6:part:b}{XIII~6.6(b)})}; hence every Cartan subgroup
of \(D\) is a Cartan subgroup of \(G\) \textup{(\SGARefLink{sga3:xiii:corollary:2:8:item:a}{XIII~2.8(a)})}.
We are therefore reduced to finding a Cartan subgroup of \(D\).

If \(\dim D=\dim G\), that is, if \(D=G\), then the Lie algebra of
\(G\) is a Cartan subalgebra of itself and is therefore nilpotent.
\SGARefLink{sga3:xiv:1.3}{Lemma~1.3} then shows that \(G\) is nilpotent, so \(G\) is a Cartan
subgroup of itself. If \(\dim D<\dim G\), the induction hypothesis
gives a Cartan subgroup of \(D\), which is consequently a Cartan
subgroup of \(G\). This proves the assertion when \(G\) is affine.

In the general case, let \(Z\) be the center of \(G\). Then
\(G/Z=G'\) is affine \textup{(\SGARefLink{sga3:xii:6.1}{XII~6.1})}, and the inverse image in
\(G\) of every Cartan subgroup \(C'\) of \(G'\) is a Cartan
subgroup of \(G\) \textup{(\SGARefLink{sga3:xii:6.6:item:e}{XII~6.6(e)})}. We are thus reduced to
finding a Cartan subgroup of the smooth affine algebraic group
\(G'\), which is the case already treated.

\paragraph{\textbf{Corollary 1.4.}}
\SGARefTarget{sga3:xiv:1.4}
\emph{Let \(G\) be a smooth group scheme of finite type over an
Artinian scheme \(S\). Then \(G\) admits a maximal torus \(T\), and
hence a Cartan subgroup \(C=\Centr_G(T)\). Every torus in \(G\) is
contained in a maximal torus.}

Indeed, we may suppose that \(S\) is local, with residue field
\(k\). By \SGARefLink{sga3:xiv:1.1}{1.1}, the group
\[
  G_0=G\times_S\operatorname{Spec}(k)
\]
admits a maximal torus \(T_0\). By \SGARefLink{sga3:ix:theorem:3:6}{IX~3.6}~bis and \SGARefLink{sga3:x:corollary:2:3}{X~2.3}, \(T_0\)
comes from a torus \(T\) of \(G\), which is plainly maximal. The
last assertion follows by applying the preceding result to the
centralizer of the given torus, which is smooth over \(k\)
\textup{(\SGARefLink{sga3:xi:corollary:2:4}{XI~2.4})}.

\paragraph{\textbf{Remarks 1.5.}}
\SGARefTarget{sga3:xiv:1.5}
\textup{(a)} We shall give below, in \SGARefLink{sga3:xiv:3.20}{3.20}, \SGARefLink{sga3:xiv:3.21}{3.21}, and Exposé XV,
variants of \SGARefLink{sga3:xiv:1.4}{1.4} in which \(S\) is not assumed Artinian.

\medskip
\noindent\textup{\SGAInlineRefTarget{sga3:xiv:remark:1-5:item:b}(b)} I do not know whether every algebraic group
\(G\), not necessarily smooth, over a field \(k\) admits a maximal
torus. The question arises only in characteristic \(p>0\). Applying
\SGARefLink{sga3:xiv:1.1}{1.1} to a smooth quotient group of the form \(G'=G/I\), where \(I\)
is a suitable radicial subgroup of \(G\) \textup{(for example, the
kernel of a sufficiently high power of the Frobenius homomorphism)},
and then taking in \(G\) the inverse image of a maximal torus of
\(G'\), reduces the question to the case in which
\((G_{\bar k})_{\mathrm{red}}\) is a torus; here, as always,
\(\bar k\) denotes the algebraic closure of \(k\). The answer is
easily seen to be affirmative when \(G\) is commutative, or more
generally nilpotent: in that case \(G\) has a unique maximal torus,
which may for example be constructed by descent from the maximal
torus of \(G_{\bar k}\).\begingroup
\renewcommand{\thefootnote}{*}\footnote{M. Raynaud gave a negative
answer to the question raised here; see XVII, \SGARefLink{sga3:xvii:subsection:5-9:item:c}{Example~5.9(c)}.}
\endgroup

\medskip
\noindent\textup{(c)} When \(G\) is affine and either \(k\) is
perfect or \(G\) is solvable, \SGARefLink{sga3:xiv:1.1}{1.1} was already known, by a theorem of
Rosenlicht; his proof is quite different from the one given here.

\medskip
\noindent\textup{(d)} When \(k\) is infinite, \SGARefLink{sga3:xiv:1.1}{1.1} follows from the
much more precise result, proved below in \SGARefLink{sga3:xiv:6.1}{6.1}, that the scheme
\(\mathcal T\) of maximal tori of \(G\) is a rational variety. The
method essentially combines the proof of \SGARefLink{sga3:xiv:1.1}{1.1} with an explicit
description of the structure of the scheme \(\mathcal D\) of Cartan
subalgebras of \(\mathfrak g\). To obtain the desired result, we
must first extend to an arbitrary base prescheme certain results of
XIII~4--6 \textup{(this is the purpose of the next two sections)},
and then refine the preceding construction proving \SGARefLink{sga3:xiv:1.1}{1.1}, using the
fact that every Cartan subgroup of \(G\) is contained in one and
only one subgroup of type \textup{(C)} in \(G\) \textup{(Nos.~4
and 5)}.

\section{Lie algebras over an arbitrary prescheme: regular sections
and Cartan subalgebras}
\SGARefTarget{sga3:xiv:section:2}

\paragraph{\textbf{Lemma 2.1.}}
\SGARefTarget{sga3:xiv:2.1}
\emph{Let \(A\) be a ring, let \(\mathfrak d\) be a Lie algebra over
\(A\), and for each \(s\in\operatorname{Spec}(A)\) let
\[
  \mathfrak d(s)=\mathfrak d\otimes_A k(s)
\]
be the corresponding Lie algebra over the residue field \(k(s)\).
Suppose that the \(A\)-module \(\mathfrak d\) is of finite
presentation. The following conditions are equivalent:
\begin{enumerate}
\item[\textup{(i)}] for every \(s\in\operatorname{Spec}(A)\), the
  Lie algebra \(\mathfrak d(s)\) is nilpotent;
\item[\textup{(ii)}] for every \(x\in\mathfrak d\), the endomorphism
  \(\operatorname{ad}(x)\) is nilpotent;
\item[\textup{(iii)}] there is an integer \(N\geq 0\) such that,
  for every sequence \(x_1,\ldots,x_N\) of elements of
  \(\mathfrak d\), one has
  \[
    \operatorname{ad}(x_1)\operatorname{ad}(x_2)\cdots
    \operatorname{ad}(x_N)=0.
  \]
\end{enumerate}}

% END INCLUDED SGA3_English_current_buildable_source_20260823/source/inputs/XIV/tex/components/00_expose_XIV_section1_opening_en.tex

% BEGIN INCLUDED SGA3_English_current_buildable_source_20260823/source/inputs/XIV/tex/components/01_expose_XIV_lemma21_proof_through_definition25_en.tex
% Authority: Expo14.pdf, local pp. 3--4 / combined pp. 837--838.
% Continuous English ends immediately before Proposition 2.6.

When \(A\) is a field, the equivalence of \textup{(i)} and
\textup{(iii)} is the definition of nilpotence, while the equivalence
of \textup{(ii)} and \textup{(iii)} is the familiar consequence of
Engel's theorem \textup{(Bourbaki, \emph{Groupes et algèbres de Lie},
Chapter I, \S4, No.~2)}. In the general case,
\textup{(iii)} plainly implies \textup{(ii)}, and \textup{(ii)}
implies \textup{(i)} by the preceding result and the fact that
\textup{(ii)} is preserved under passage to quotients and under
localization. It remains to prove that \textup{(i)} implies
\textup{(iii)}.

Suppose first that \(A\) is local Artinian, with maximal ideal
\(\mathfrak m\). Choose \(n>0\) such that \(\mathfrak m^n=0\), and
choose \(N\) such that condition \textup{(iii)} holds for
\[
  \mathfrak d(s)=\mathfrak d\otimes_A A/\mathfrak m.
\]
Taking \(N'=nN\), one sees at once that \(N'\) satisfies
\textup{(iii)} for \(\mathfrak d\).

If \(A\) is Noetherian, there are finitely many points
\(s_i\in\operatorname{Spec}(A)\) and ideals of definition
\(\mathfrak q_i\subset A_{s_i}\) such that the natural map
\[
  \mathfrak d\longrightarrow
  \prod_i\bigl(\mathfrak d\otimes_A A_i\bigr),
  \qquad A_i=A_{s_i}/\mathfrak q_i,
\]
is injective. By the preceding case, for each \(i\) there is an
integer \(N_i\) satisfying \textup{(iii)} for the Lie algebra
\(\mathfrak d\otimes_A A_i\). Taking \(N\) to be the largest of the
\(N_i\), condition \textup{(iii)} follows for \(\mathfrak d\).
Finally, the general case reduces to the Noetherian case by the
usual procedure described in EGA~IV~8.

\paragraph{\textbf{Definition 2.2.}}
\SGARefTarget{sga3:xiv:2.2}
Let \(S\) be a prescheme and let \(\mathfrak d\) be a
quasi-coherent Lie algebra over \(S\) whose underlying module is of
finite presentation. We say that \(\mathfrak d\) is
\emph{locally nilpotent} if, for every \(s\in S\), the Lie algebra
\(\mathfrak d(s)\) over \(k(s)\) is nilpotent. We say that
\(\mathfrak d\) is \emph{strictly locally nilpotent} if it is
locally free and, on every open subset \(U\) of \(S\) over which it
has constant rank \(r\), its Killing polynomial reduces to
\[
  P_{\mathfrak d}(t)=t^r.
\]

Of course, when \(\mathfrak d\) is a Lie algebra over \(S\) that is
locally free as a module, its Killing polynomial is defined as a
polynomial
\[
  P_{\mathfrak d}\in A[t],
  \qquad
  A=\Gamma\!\left(\operatorname{Sym}_{\mathcal O_S}(\mathfrak d)\right)
   \simeq \Gamma\!\left(W(\mathfrak d)\right),
\]
where \(A\) is the ring of sections of the structure sheaf of the
vector bundle \(W(\mathfrak d)\) defined by \(\mathfrak d\).

The two notions introduced in \SGARefLink{sga3:xiv:2.2}{2.2} are plainly stable under base
change and local for the fpqc topology. We record:

\paragraph{\textbf{Proposition 2.3.}}
\SGARefTarget{sga3:xiv:2.3}
\emph{If \(\mathfrak d\) is strictly locally nilpotent, then it is
locally nilpotent. The converse holds when \(\mathfrak d\) is
locally free and \(S\) is reduced.}

The proof is immediate.

\paragraph{\textbf{Definition 2.4.}}
\SGARefTarget{sga3:xiv:2.4}
Let \(S\) be a prescheme, let \(\mathfrak g\) be a Lie algebra over
\(S\) whose underlying module is locally free of finite type, and
let \(\mathfrak d\) be a Lie subalgebra of \(\mathfrak g\). We call
\(\mathfrak d\) a \emph{Cartan subalgebra} if it satisfies the
following conditions:
\begin{enumerate}
\item[\textup{\SGAInlineRefTarget{sga3:xiv:2.4:item:i}(i)}] \(\mathfrak d\) is locally a direct-factor
  submodule \textup{(and hence is itself locally free of finite
  type)};
\item[\textup{(ii)}] for every \(s\in S\), the Lie algebra
  \(\mathfrak d(s)\) is a Cartan subalgebra of
  \(\mathfrak g(s)\).
\end{enumerate}

\paragraph{\textbf{Definition 2.5.}}
\SGARefTarget{sga3:xiv:2.5}
Let \(S\) and \(\mathfrak g\) be as in \SGARefLink{sga3:xiv:2.4}{2.4}. A section \(a\) of
\(\mathfrak g\) is called \emph{quasi-regular} if, for every
\(s\in S\), the element \(a(s)\in\mathfrak g(s)\) is regular in the
Lie algebra \(\mathfrak g(s)\) over \(k(s)\). We call \(a\) a
\emph{regular section} if it is quasi-regular and is contained in a
Cartan subalgebra of \(\mathfrak g\).

The notions in \SGARefLink{sga3:xiv:2.4}{2.4} and \SGARefLink{sga3:xiv:2.5}{2.5} are again stable under base change and
local for the fpqc topology. Only the last assertion, and only for
the notion of a regular section, requires proof. It follows from the
fact that the Cartan subalgebra containing a given regular section
is uniquely determined. More precisely:

% END INCLUDED SGA3_English_current_buildable_source_20260823/source/inputs/XIV/tex/components/01_expose_XIV_lemma21_proof_through_definition25_en.tex

% BEGIN INCLUDED SGA3_English_current_buildable_source_20260823/source/inputs/XIV/tex/components/02_expose_XIV_proposition26_corollary27_en.tex
% Authority: Expo14.pdf, local p. 5 / combined p. 839.
% Continuous English ends immediately before Corollary 2.8.

\paragraph{\textbf{Proposition 2.6.}}
\SGARefTarget{sga3:xiv:2.6}
\emph{Let \(S\) and \(\mathfrak g\) be as in \SGARefLink{sga3:xiv:2.4}{2.4}, and let \(a\) be a
quasi-regular section of \(\mathfrak g\). There is at most one Cartan
subalgebra of \(\mathfrak g\) containing \(a\). Such a Cartan
subalgebra exists---that is, \(a\) is a regular section---if and only
if \(a\) satisfies the following condition:
\[
  \mathfrak d=\Nil(a,\mathfrak g)
\]
is locally a direct-factor submodule of \(\mathfrak g\), and
\(\operatorname{ad}(a)\) induces an automorphism of
\(\mathfrak g/\mathfrak d\). In that case \(\mathfrak d\) is the
unique Cartan subalgebra of \(\mathfrak g\) containing \(a\).}

Suppose first that \(a\) is contained in a Cartan subalgebra
\(\mathfrak d\) of \(\mathfrak g\). Then
\(\operatorname{ad}(a)_{\mathfrak g/\mathfrak d}\) is bijective on
every fiber. Since \(\mathfrak g/\mathfrak d\) is locally free of
finite type, it follows that this endomorphism is an automorphism of
\(\mathfrak g/\mathfrak d\). On the other hand, by \SGARefLink{sga3:xiv:2.1}{2.1} the
endomorphism \(\operatorname{ad}(a)_{\mathfrak d}\) is locally
nilpotent. Hence
\[
  \mathfrak d=\Nil(a,\mathfrak g),
\]
which proves both the uniqueness of \(\mathfrak d\) and the
necessity of the stated regularity criterion.

For sufficiency, observe that the hypothesis on \(a\) implies that
the formation of \(\Nil(a,\mathfrak g)\) commutes with base
extension, and in particular with passage to fibers. Thus the fibers
\(\mathfrak d(s)\) of
\(\mathfrak d=\Nil(a,\mathfrak g)\) are Cartan subalgebras of the
\(\mathfrak g(s)\). Moreover, \(\mathfrak d\) is a Lie subalgebra of
\(\mathfrak g\) by \SGARefLink{sga3:xiii:proposition:4:1}{XIII~4.1}. It is therefore a Cartan subalgebra.

\paragraph{\textbf{Corollary 2.7.}}
\SGARefTarget{sga3:xiv:2.7}
\emph{Under the hypotheses of \SGARefLink{sga3:xiv:2.4}{2.4}, let \(\mathfrak d\) be a Cartan
subalgebra of \(\mathfrak g\), let \(a\) be a section of
\(\mathfrak d\) that is regular in \(\mathfrak g\), and let \(u\) be
an automorphism of \(\mathfrak g\). Then \(u\) preserves
\(\mathfrak d\) if and only if \(u(a)\) is a section of
\(\mathfrak d\).}

Indeed, by transport of structure \(u(a)\) is then a regular section
of \(\mathfrak g\), contained in both Cartan subalgebras
\(\mathfrak d\) and \(u(\mathfrak d)\); \SGARefLink{sga3:xiv:2.6}{Proposition~2.6} therefore
shows that these two subalgebras are equal.

% END INCLUDED SGA3_English_current_buildable_source_20260823/source/inputs/XIV/tex/components/02_expose_XIV_proposition26_corollary27_en.tex

% BEGIN INCLUDED SGA3_English_current_buildable_source_20260823/source/inputs/XIV/tex/components/03_expose_XIV_corollary28_and_nilpotent_rank_en.tex
% Authority: Expo14.pdf, local p. 5 / combined p. 839.
% Continuous English ends immediately before Proposition 2.9.

\paragraph{\textbf{Corollary 2.8.}}
\SGARefTarget{sga3:xiv:2.8}
\emph{Under the hypotheses of \SGARefLink{sga3:xiv:2.4}{2.4}, let \(\mathfrak d\) be a Cartan
subalgebra of \(\mathfrak g\). For every \(s\in S\) such that
\(k(s)\) is infinite, there exist an open neighborhood \(V\) of
\(s\) and a regular section \(a\) of \(\mathfrak g\) over \(V\)
such that
\[
  \mathfrak d|_V=\Nil\!\left(a,\mathfrak g|_V\right);
\]
equivalently, \(a\) is a section of \(\mathfrak d|_V\).}

Indeed, since \(k(s)\) is infinite, there is a regular element
\(a_0\) of \(\mathfrak g(s)\) contained in \(\mathfrak d(s)\).
After replacing \(S\) by an open neighborhood \(U\) of \(s\), we
may extend \(a_0\) to a section \(a\) of \(\mathfrak d\) over
\(U\). The endomorphism induced by \(\operatorname{ad}(a)\) on the
fiber
\[
  \mathfrak g(s)/\mathfrak d(s)
\]
is an automorphism. Shrinking \(U\) to a smaller open neighborhood
\(V\), we may therefore suppose that
\(\operatorname{ad}(a)_{\mathfrak g/\mathfrak d}\) is an
automorphism. It follows that \(a\) is a quasi-regular section of
\(\mathfrak g|_V\) satisfying the required condition.

Let \(\mathfrak g\) be as in \SGARefLink{sga3:xiv:2.4}{2.4}. Inspection of its Killing
polynomial immediately shows that the function
\[
  s\longmapsto\text{the nilpotent rank of }\mathfrak g(s)
\]
on \(S\) is upper semicontinuous. We shall be especially interested
in the case in which this function is continuous, that is, locally
constant. Here are several variants of this property.

% END INCLUDED SGA3_English_current_buildable_source_20260823/source/inputs/XIV/tex/components/03_expose_XIV_corollary28_and_nilpotent_rank_en.tex

% BEGIN INCLUDED SGA3_English_current_buildable_source_20260823/source/inputs/XIV/tex/components/04_expose_XIV_proposition29_through_corollary211_en.tex
% Authority: Expo14.pdf, local pp. 5--6 / combined pp. 839--840.
% Continuous English ends immediately before Remark 2.12.

\paragraph{\textbf{Proposition 2.9.}}
\SGARefTarget{sga3:xiv:2.9}
{\itshape Let \(S,\mathfrak g\) be as in \SGARefLink{sga3:xiv:2.4}{2.4}. Consider the following
conditions:

\smallskip
\noindent
\textup{\((\mathrm C_0)\)}
The nilpotent rank of \(\mathfrak g(s)\), for \(s\in S\), is a
locally constant function of \(s\).

\smallskip
\noindent
\textup{\((\mathrm C_1)\)}
Locally for the fpqc topology, there exists a Cartan subalgebra of
\(\mathfrak g\).

\smallskip
\noindent
\textup{\((\mathrm C'_1)\)}
The same as \((\mathrm C_1)\), with ``fpqc topology'' replaced by
``étale topology''.

\smallskip
\noindent
\textup{\((\mathrm C_2)\)}
Condition \((\mathrm C_0)\) holds, and for every \(S'\) over \(S\),
every quasi-regular section of
\[
  \mathfrak g_{S'}=\mathfrak g\otimes_S S'
\]
is regular.

\smallskip
\noindent
\textup{\((\mathrm C_3)\)}
Every \(s\in S\) has an open neighborhood \(V\) on which the
Killing polynomial of \(\mathfrak g\) has the form
\[
  t^r\bigl(t^{n-r}+c_1t^{n-r-1}+\cdots+c_{n-r}\bigr),
\]
where, for every \(s\in V\),
\[
  c_{n-r}(s)\in
  \operatorname{Sym}\!\bigl(\mathfrak g(s)^\vee\bigr)
\]
is nonzero.

\smallskip
With this notation, one has the implications
\[
  (\mathrm C_3)\Longrightarrow(\mathrm C_2)
  \Longrightarrow(\mathrm C'_1)\Longrightarrow(\mathrm C_1)
  \Longrightarrow(\mathrm C_0),
\]
and, when \(S\) is reduced, the five conditions above are
equivalent.}

Notice also that these conditions are manifestly stable under base
change and are local for the fpqc topology.

The implications
\[
  (\mathrm C'_1)\Longrightarrow(\mathrm C_1)
  \Longrightarrow(\mathrm C_0)
\]
are trivial, and the implication
\((\mathrm C_0)\Longrightarrow(\mathrm C_3)\) when \(S\) is
reduced is immediate. We note moreover:

\paragraph{\textbf{Corollary 2.10.}}
\SGARefTarget{sga3:xiv:2.10}
\emph{Suppose that condition \((\mathrm C_0)\) holds. Let \(U\) be
the set of elements of \(W(\mathfrak g)\) that are regular in their
fiber. Then \(U\) is open. In particular, for every section \(a\)
of \(\mathfrak g\) over \(S\), the set \(V\) of \(s\in S\) for
which \(a(s)\in\mathfrak g(s)\) is regular is open.}

Indeed, the first assertion follows from the second, applied to
\(\mathfrak g_{S'}\) for every base change \(S'\) over \(S\). For
the second, we may suppose that \(S\) is reduced; then
\((\mathrm C_3)\) holds, and it is enough to inspect the Killing
polynomial of \(a\) in \(\mathfrak g\).

The implication
\((\mathrm C_2)\Longrightarrow(\mathrm C'_1)\) now follows easily
from part~b) of the following more precise statement.

\paragraph{\textbf{Corollary 2.11.}}
\SGARefTarget{sga3:xiv:2.11}
{\itshape Suppose that condition \((\mathrm C_2)\) holds. Then:

\smallskip
\noindent
\textup{a)}
For every \(s\in S\) and every Cartan subalgebra
\(\mathfrak d_0\) of \(\mathfrak g(s)\) such that
\(\mathfrak d_0\) contains a regular element of \(\mathfrak g(s)\)
\textup{(a condition automatically satisfied when \(k(s)\) is
infinite)}, there exist an open neighborhood \(V\) of \(s\) and a
Cartan subalgebra \(\mathfrak d\) of \(\mathfrak g|_V\) whose fiber
at \(s\) is \(\mathfrak d_0\). If \(S_1\) is a subscheme of \(S\)
containing \(s\), and if \(\mathfrak d_0\) has already been
extended to a Cartan subalgebra \(\mathfrak d_1\) of
\(\mathfrak g\otimes_S S_1\), then one can find an open
neighborhood \(V\) of \(s\) in \(S\) and a Cartan subalgebra
\(\mathfrak d\) of \(\mathfrak g|_V\) such that
\[
  \mathfrak d\otimes_V(S_1\cap V)
  =\mathfrak d_1|_{S_1\cap V}.
\]

\smallskip
\noindent
\textup{b)}
For every \(s\in S\) such that \(\mathfrak g(s)\) contains a
regular element \textup{(a condition automatically satisfied when
\(k(s)\) is infinite)}, there exist an open neighborhood \(V\) of
\(s\) and a Cartan subalgebra \(\mathfrak d\) of
\(\mathfrak g|_V\).}

Part~b) follows from part~a) by taking
\[
  \mathfrak d_0=\Nil\!\left(a_0,\mathfrak g(s)\right),
\]
where \(a_0\) is a regular element of \(\mathfrak g(s)\).
To prove part~a), in its second formulation, choose a regular
element \(a_0\) of \(\mathfrak g(s)\) contained in
\(\mathfrak d_0\). Extend it near \(s\) in \(S_1\) to a section of
\(\mathfrak d_1\), and extend that section in turn near \(s\) to a
section of \(\mathfrak g\). By \SGARefLink{sga3:xiv:2.10}{2.10} this section is quasi-regular
on an open neighborhood \(V\) of \(s\), hence regular by
\((\mathrm C_2)\). Therefore, by \SGARefLink{sga3:xiv:2.6}{2.6},
\[
  \mathfrak d=\Nil\!\left(a,\mathfrak g|_V\right)
\]
has the required properties.

% END INCLUDED SGA3_English_current_buildable_source_20260823/source/inputs/XIV/tex/components/04_expose_XIV_proposition29_through_corollary211_en.tex

% BEGIN INCLUDED SGA3_English_current_buildable_source_20260823/source/inputs/XIV/tex/components/05_expose_XIV_remark212_lemma213_en.tex
% Authority: Expo14.pdf, local p. 7 / combined p. 841.
% Continuous English ends after the proof of Lemma 2.13, hence of 2.9.

\paragraph{\textbf{Remark 2.12.}}
\SGARefTarget{sga3:xiv:2.12}
I do not know whether the assertions of \SGARefLink{sga3:xiv:2.11}{2.11}~a) and b) remain valid
without the hypothesis that regular points exist, when \(k(s)\) is
finite.

It remains to prove
\[
  (\mathrm C_3)\Longrightarrow(\mathrm C_2).
\]
We also record the following equivalent form of
\((\mathrm C_3)\):

\smallskip
\noindent
\textup{\((\mathrm C'_3)\)}
Condition \((\mathrm C_0)\) holds, that is, the nilpotent rank of
\(\mathfrak g(s)\), for \(s\in S\), is locally constant; and on
every open subset \(V\) of \(S\) on which this rank has the value
\(r\), the Killing polynomial of \(\mathfrak g|_V\) is divisible
by \(t^r\).

\smallskip
We must show that this condition implies that every quasi-regular
section \(a\) of \(\mathfrak g\) is regular. In view of \SGARefLink{sga3:xiv:2.6}{2.6}, this
follows from the implication
\(\textup{(iv)}\Longrightarrow\textup{(iii)}\) in the following
lemma, applied to the endomorphism
\(\operatorname{ad}(a)\) of \(\mathfrak g\).

\paragraph{\textbf{Lemma 2.13.}}
\SGARefTarget{sga3:xiv:2.13}
{\itshape Let \(A\) be a ring, let \(M\) be a projective
\(A\)-module of finite type, and let \(u\) be an endomorphism of
\(M\). The following conditions are equivalent:

\smallskip
\noindent
\textup{(i)}
\(M\) is the direct sum of two \(u\)-stable submodules \(M'\) and
\(M''\), such that \(u|_{M'}\) is nilpotent and \(u|_{M''}\) is an
automorphism of \(M''\).

\smallskip
\noindent
\textup{(ii)}
There exists an integer \(n>0\) such that
\[
  \operatorname{Im}u^n+\operatorname{Ker}u^n=M.
\]

\smallskip
\noindent
\textup{(iii)}
The nilspace
\[
  \Nil(u)=\bigcup_{n>0}\operatorname{Ker}u^n
\]
is a direct factor of \(M\), and
\[
  M=\Nil(u)+u(M).
\]

\smallskip
These conditions are implied by the following one, and are
equivalent to it when \(A\) is reduced:

\smallskip
\noindent
\textup{(iv)}
Locally on \(\operatorname{Spec}(A)\), for the Zariski topology, the
characteristic polynomial \(P_u(t)\) of \(u\) can be written in the
form
\[
  t^r\bigl(t^{n-r}+c_1t^{n-r-1}+\cdots+c_{n-r}\bigr),
\]
where \(c_{n-r}\) is invertible.}

The equivalence of \textup{(i)}, \textup{(ii)}, and \textup{(iii)}
is immediate and is recalled for reference. When \(A\) is reduced,
\textup{(i)} implies \textup{(iv)} because a nilpotent
endomorphism of a projective module of rank \(r\) then has
characteristic polynomial \(t^r\). On the other hand, in all cases
the constant term of the characteristic polynomial of an
automorphism of a projective module of finite type is, up to sign,
the determinant of \(u\), locally on \(\operatorname{Spec}(A)\);
it is therefore an invertible element of \(A\).

Finally, to prove
\(\textup{(iv)}\Longrightarrow\textup{(i)}\), regard \(M\) as a
module over the polynomial ring \(A[t]\), with \(t\) acting as
\(u\). The familiar identity
\[
  P(u)=0
\]
shows that \(M\) is annihilated by \(P A[t]\), and hence may be
regarded as a module over \(A[t]/P A[t]\). Write
\[
  P=t^rQ,
\]
where the constant term of \(Q\) is invertible. Then
\[
  P A[t]=t^rA[t]\cap Q A[t],
\]
so that \(A[t]/P A[t]\) decomposes as the product of the rings
\[
  A[t]/t^rA[t]
  \qquad\text{and}\qquad
  A[t]/Q A[t].
\]
This gives a corresponding decomposition of \(M\) into the direct
sum of two \(A[t]\)-modules, or equivalently two \(u\)-stable
\(A\)-submodules \(M'\) and \(M''\); this is precisely the
decomposition of \textup{(i)}.

This completes the proof of \SGARefLink{sga3:xiv:2.13}{2.13}, and hence that of \SGARefLink{sga3:xiv:2.9}{2.9}.

% END INCLUDED SGA3_English_current_buildable_source_20260823/source/inputs/XIV/tex/components/05_expose_XIV_remark212_lemma213_en.tex

% BEGIN INCLUDED SGA3_English_current_buildable_source_20260823/source/inputs/XIV/tex/components/06_expose_XIV_corollary214_remark215_en.tex
% Authority: Expo14.pdf, local p. 8 / combined p. 842.
% Continuous English ends immediately after Remark 2.15(c).

\paragraph{\textbf{Corollary 2.14.}}
\SGARefTarget{sga3:xiv:2.14}
\emph{When condition \((\mathrm C_3)\) of \SGARefLink{sga3:xiv:2.9}{2.9} holds, the Cartan
subalgebras of \(\mathfrak g\) are strictly nilpotent
\textup{(\SGARefLink{sga3:xiv:2.2}{2.2})}.}

The proof is immediate.

\paragraph{\textbf{Remark 2.15.}}
\SGARefTarget{sga3:xiv:2.15}
\noindent
\textup{a)}
One can prove a converse to \SGARefLink{sga3:xiv:2.14}{2.14}: condition \((\mathrm C_3)\) is
equivalent to the assertion that, for every \(S'\) over \(S\),
every quasi-regular section of \(\mathfrak g_{S'}\) is regular and
every Cartan subalgebra of \(\mathfrak g_{S'}\) is strictly
nilpotent; equivalently, every quasi-regular section of
\(\mathfrak g_{S'}\) is contained in a strictly nilpotent Cartan
subalgebra of \(\mathfrak g_{S'}\).

\smallskip
\noindent
\textup{b)}
Unlike the other conditions \((\mathrm C_0)\) through
\((\mathrm C_2)\), condition \((\mathrm C_3)\) is of
``infinitesimal nature.'' More precisely, when \(S\) is locally
noetherian, \(\mathfrak g\) satisfies \((\mathrm C_3)\) if and only
if, for every base change \(S'\to S\) with \(S'\) local Artinian,
\(\mathfrak g_{S'}\) satisfies the same condition. If desired, one
may take \(S'\) to be the spectrum of an Artinian quotient of a
local ring of \(S\). Similarly, when \((\mathrm C_0)\) holds, the
condition for a section of \(\mathfrak g\) to be regular is of
infinitesimal nature.

\smallskip
\noindent
\textup{c)}
When \(\mathfrak g\) is the Lie algebra of a smooth group
prescheme of finite presentation over \(S\), we shall see that the
conditions
\[
  (\mathrm C_0),\quad(\mathrm C_1),\quad
  (\mathrm C'_1),\quad(\mathrm C_2)
\]
on \(\mathfrak g\) are equivalent \textup{(\SGARefLink{sga3:xiv:5.2}{5.2}~a)}. I do not know
what happens in general, except that even when \(S\) is local
Artinian, \((\mathrm C_0)\) does not imply \((\mathrm C_1)\).
However, even when \(\mathfrak g\) comes from a group \(G\) and
\(S\) is local Artinian, it is not true in general that
\((\mathrm C_2)\) implies \((\mathrm C_3)\), since
\(\mathfrak g\) may be nilpotent without being strictly
nilpotent.

An example is obtained from a smooth affine group scheme \(G\)
over the spectrum \(S\) of a discrete valuation ring, such that
the Lie algebra of the generic fiber is not nilpotent---for
example, the generic fiber may be an adjoint semisimple
group---whereas that of the special fiber is nilpotent---for
example, the special fiber may be a vector group. Then, for
sufficiently large \(n\), the Lie algebra of
\[
  G_n=G\times_S S_n
\]
is not strictly nilpotent, although it is nilpotent.

% END INCLUDED SGA3_English_current_buildable_source_20260823/source/inputs/XIV/tex/components/06_expose_XIV_corollary214_remark215_en.tex

% BEGIN INCLUDED SGA3_English_current_buildable_source_20260823/source/inputs/XIV/tex/components/07_expose_XIV_cartans_functor_theorem216_statement_en.tex
% Authority: Expo14.pdf, local pp. 8--9 / combined pp. 842--843.
% Continuous English ends immediately before the proof of Theorem 2.16.

We continue under the hypotheses of \SGARefLink{sga3:xiv:2.4}{2.4}. Let
\[
  \mathcal D\colon
  \bigl((\mathrm{Sch})_{/S}\bigr)^\circ
  \longrightarrow \mathrm{Sets}
\]
be the functor defined by
\[
  \mathcal D(S')
  =\left\{
    \text{Cartan subalgebras of }\mathfrak g_{S'}
   \right\}.
\]
Introduce also the functor
\[
  X(S')
  =\left\{
    (\mathfrak d,a)\ \middle|\
    \begin{array}{l}
      \mathfrak d\text{ is a Cartan subalgebra of }\mathfrak g_{S'},\\
      a\text{ is a section of }\mathfrak d
    \end{array}
   \right\}.
\]
The two projections
\[
  (\mathfrak d,a)\longmapsto\mathfrak d
  \qquad\text{and}\qquad
  (\mathfrak d,a)\longmapsto a
\]
give morphisms
\[
  p\colon X\longrightarrow\mathcal D
  \qquad\text{and}\qquad
  \psi\colon X\longrightarrow W(\mathfrak g).
\]

When \((\mathrm C_0)\) holds, let \(U\) also denote the open subset
of regular points of \(W(\mathfrak g)\) \textup{(\SGARefLink{sga3:xiv:2.10}{2.10})}. Condition
\((\mathrm C_2)\) is then equivalent to saying that the morphism
\[
  \psi^{-1}(U)\longrightarrow U
\]
induced by \(\psi\) is an isomorphism; a priori it is a
monomorphism by \SGARefLink{sga3:xiv:2.6}{2.6}.

It is immediate that
\[
  p\colon X\longrightarrow\mathcal D
\]
is representable by a projection of vector bundles. More
explicitly, for every \(S\)-morphism \(S'\to\mathcal D\)
corresponding to a Cartan subalgebra \(\mathfrak d\) of
\(\mathfrak g_{S'}\), the fiber product
\[
  X\times_{\mathcal D}S'
\]
is represented by the vector bundle \(W(\mathfrak d)\) over
\(S'\). Thus, if \(\mathcal D\) is representable, then so is \(X\).
We have:

\paragraph{\textbf{Theorem 2.16.}}
\SGARefTarget{sga3:xiv:2.16}
{\itshape Let \(S\) be a prescheme and let \(\mathfrak g\) be a Lie
algebra over \(S\) whose underlying \(\mathcal O_S\)-module is
locally free of finite type. Suppose that condition
\((\mathrm C_0)\) of \SGARefLink{sga3:xiv:2.9}{2.9} holds.

\smallskip
\noindent
\textup{a)}
The functor \(\mathcal D\) of Cartan subalgebras of
\(\mathfrak g\), defined above, is represented by a
quasi-projective prescheme of finite presentation over \(S\).
The same holds for the functor \(X\).

\smallskip
\noindent
\textup{b)}
If condition \((\mathrm C_2)\) of \SGARefLink{sga3:xiv:2.9}{2.9} holds, then
\(\mathcal D\) and \(X\) are smooth over \(S\), and the morphism
\[
  \psi^{-1}(U)\longrightarrow U
\]
induced by \(\psi\) is an isomorphism.

\smallskip
\noindent
\textup{c)}
Continue to suppose that \((\mathrm C_2)\) holds. Let \(s\in S\),
and let \(\mathfrak d_0\) be a Cartan subalgebra of
\(\mathfrak g(s)\), corresponding to a \(k(s)\)-rational point
\(d\) of \(\mathcal D\). Suppose that \(\mathfrak d_0\) contains a
regular element of \(\mathfrak g(s)\), a condition automatically
satisfied when \(k(s)\) is infinite. If \(r\) is the infinitesimal
rank of \(\mathfrak g(s)\) and \(n\) its rank over \(k(s)\), then
there exists an open neighborhood \(V\) of \(d\) in
\(\mathcal D\) that is \(S\)-isomorphic to an open subset \(V'\)
of
\[
  S[t_1,\ldots,t_{n-r}].
\]}

% END INCLUDED SGA3_English_current_buildable_source_20260823/source/inputs/XIV/tex/components/07_expose_XIV_cartans_functor_theorem216_statement_en.tex

% BEGIN INCLUDED SGA3_English_current_buildable_source_20260823/source/inputs/XIV/tex/components/08_expose_XIV_theorem216_proof_corollary217_remark218_en.tex
% Authority: Expo14.pdf, local pp. 9--11 / combined pp. 843--845.
% Continuous English ends at the close of Section 2, before Section 3.

\paragraph{\textit{Proof.}}
We may suppose that \(\mathfrak g\) has constant rank \(n\) and
constant nilpotent rank \(r\). The assertions concerning \(X\) in
\textup{a)} and \textup{b)} follow immediately from those concerning
\(\mathcal D\) and from the fact that \(X\) is a vector bundle over
\(\mathcal D\), defined by a locally free module; they were included
only for completeness.

\medskip
\noindent\textup{a)}
The functor \(\mathcal D\) is a subfunctor of
\(\operatorname{Grass}_{n-r}(\mathfrak g)\), whose value at \(S'\)
is the set of locally free quotient modules of rank \(n-r\) of
\(\mathfrak g_{S'}\). It is well known that the latter functor is
represented by a projective prescheme smooth over \(S\); see, for
example, Cartan Seminar 1960/61, Exposé~12, Nos.~2 and~3, whose
constructions carry over unchanged to preschemes.\begingroup
\renewcommand{\thefootnote}{*}\footnote{See also EGA~I, second
edition \textup{(announced as forthcoming from North-Holland
Publishing Company)}.}
\endgroup\footnote{Editors' note: See \S \SGARefLink{sga3:xiv:1.3}{I.1.3} of M.~Demazure and
P.~Gabriel, \emph{Groupes algébriques}, Masson \textup{(1970)}.}

We are therefore reduced to the following relative problem. Given
a locally free quotient module of rank \(n-r\) of \(\mathfrak g\),
or, equivalently, a locally free submodule \(\mathfrak d\) of rank
\(r\) that is locally a direct factor, represent the functor
\[
 F(S')=
 \begin{cases}
   \varnothing,
     &\text{if \(\mathfrak d_{S'}\) is not a Cartan subalgebra of
       \(\mathfrak g_{S'}\),}\\
   \{\varnothing\},
     &\text{otherwise.}
 \end{cases}
\]
We shall show that \(F\) is represented by a subprescheme of finite
presentation over \(S\). This will show that
\[
  \mathcal D\longrightarrow
  \operatorname{Grass}_{n-r}(\mathfrak g)
\]
is represented by an immersion of finite presentation and will
complete the proof of \textup{a)}.

First express the condition that \(\mathfrak d_{S'}\) be a Lie
subalgebra of \(\mathfrak g_{S'}\). This condition says precisely
that \(S'\to S\) factors through a certain closed subprescheme
\(S_1\) of finite presentation over \(S\); locally on \(S\), its
equations are obtained at once from a basis of \(\mathfrak g\)
adapted to \(\mathfrak d\). We may therefore suppose that
\(S=S_1\).

Next we must express that \(\mathfrak d_{S'}\) contains,
fpqc-locally, a quasi-regular section of \(\mathfrak g_{S'}\). For
this purpose consider
\[
  V=W(\mathfrak d)\cap U,
\]
where \(U\) is the open subset of regular points of
\(W(\mathfrak g)\) \textup{(\SGARefLink{sga3:xiv:2.10}{2.10})}. Since the structure morphism
\(V\to S\) is smooth and quasi-compact, its image \(S_1\) is an
open subset of \(S\), and the immersion \(S_1\to S\) is
quasi-compact, hence of finite presentation. The condition on
\(S'\) is exactly that \(S'\to S\) factor through \(S_1\). We are
thus reduced to the case \(S=S_1\), and, by descent, to the case in
which \(\mathfrak d\) has a section \(a\) that is quasi-regular in
\(\mathfrak g\).

It remains to express that, for the induced section \(a_{S'}\), the
endomorphism
\[
  \operatorname{ad}(a_{S'})
    _{\mathfrak g_{S'}/\mathfrak d_{S'}}
\]
is bijective. This again says that \(S'\to S\) factors through an
open subprescheme of finite presentation, namely \(S_D\), where
\[
  D=\det\!\left(
    \operatorname{ad}(a)_{\mathfrak g/\mathfrak d}
  \right).
\]
But \(\mathfrak d|_{S_D}\) is then a Cartan subalgebra of
\(\mathfrak g|_{S_D}\). Consequently \(S_D\) represents \(F\),
which proves \textup{a)}.

\medskip
\noindent\textup{b)}
This follows immediately from \SGARefLink{sga3:xiv:2.11}{2.11}\textup{(a)} and \SGARefLink{sga3:xi:remark:1:5}{XI~1.5}. It is
also, of course, a consequence of the more precise assertion
\textup{c)}.

\medskip
\noindent\textup{c)}
Let \(a_0\) be a regular point of \(\mathfrak g(s)\) contained in
\(\mathfrak d_0\), and extend it to a section \(a\) of
\(\mathfrak g\) over an open neighborhood of \(s\). We may plainly
suppose that this neighborhood is all of \(S\). Choose a complement
\(M_0\) of \(\mathfrak d_0\) in the vector space
\(\mathfrak g(s)\). After again replacing \(S\) by an open
neighborhood of \(s\), there is a submodule \(M\) of
\(\mathfrak g\), locally a direct factor, such that
\(M(s)=M_0\); we again suppose that the neighborhood is \(S\).

Let \(\mathcal V\) be the subfunctor of \(\mathcal D\) for which
\(\mathcal V(S')\) consists of the Cartan subalgebras
\(\mathfrak d'\) of \(\mathfrak g_{S'}\) satisfying:
\begin{enumerate}
\item[\textup{1\textdegree)}]
  \(\mathfrak d'\) is a complement of \(M_{S'}\);
\item[\textup{2\textdegree)}]
  the unique section in
  \((a_{S'}+M_{S'})\cap\mathfrak d'\) is a regular section of
  \(\mathfrak g_{S'}\).
\end{enumerate}
Condition \textup{1\textdegree)} defines an open \(V_1\) of
\(\mathcal D\), induced by the open of
\(\operatorname{Grass}_{n-r}(\mathfrak g)\) given by the same
condition. By \SGARefLink{sga3:xiv:2.10}{2.10} and \((\mathrm C_2)\), the conjunction of
\textup{1\textdegree)} and \textup{2\textdegree)} defines an open
of \(V_1\). Thus \(\mathcal V\) is represented by an open
subprescheme \(V\) of \(\mathcal D\), plainly containing the point
\(d\).

Now let \(\mathcal V'\) be the subfunctor of \(W(M)\) defined by
\[
 \mathcal V'(S')
 =
 \left\{
   u'\in\Gamma(S',M_{S'})
   \ \middle|\
   \begin{array}{l}
   \textup{(i) }a_{S'}+u'\text{ is a regular section of }
     \mathfrak g_{S'},\\
   \textup{(ii) the unique Cartan subalgebra }\mathfrak d'
     \text{ of }\mathfrak g_{S'}\\
   \quad\text{containing }a_{S'}+u'
     \text{ is a complement of }M_{S'}
   \end{array}
 \right\}.
\]
Condition \textup{(i)} defines an open subprescheme \(V'_1\) of
\(W(M)\): it is the inverse image of the open \(U\) of regular
points of \(W(\mathfrak g)\) \textup{(\SGARefLink{sga3:xiv:2.10}{2.10})} under the translation
morphism \(m\mapsto a+m\). The conjunction of \textup{(i)} and
\textup{(ii)} defines an open \(V'\) of \(V'_1\), namely the
inverse image of \(V\) under the evident morphism
\[
  V'_1\longrightarrow\mathcal D
\]
that associates with \(u'\) the unique Cartan subalgebra
\(\mathfrak d'\) containing \(a_{S'}+u'\). Its restriction gives
a morphism
\[
  V'\longrightarrow V,
\]
which is plainly an isomorphism. This proves \textup{c)}.

\paragraph{\textbf{Corollary 2.17.}}
\SGARefTarget{sga3:xiv:2.17}
{\itshape Let \(\mathfrak g\) be a finite-dimensional Lie algebra
over a field \(k\). Then the scheme \(\mathcal D\) of Cartan
subalgebras of \(\mathfrak g\) \textup{(\SGARefLink{sga3:xiv:2.16}{2.16}\textup{(a)})} is
quasi-projective, smooth, and irreducible. If \(\mathfrak g\)
contains a regular element \textup{(for example, if \(k\) is
infinite)}, then \(\mathcal D\) is a rational variety; that is, its
function field is a purely transcendental extension of \(k\).}

Indeed, \(\mathcal D\) is irreducible because there is a surjective
morphism
\[
  \psi^{-1}(U)\longrightarrow\mathcal D,
\]
and \(\psi^{-1}(U)\) is irreducible, being isomorphic to the open
subset \(U\) of \(W(\mathfrak g)\). The assertion concerning the
function field follows immediately from \SGARefLink{sga3:xiv:2.16}{2.16}\textup{(c)}.

\paragraph{\textbf{Remarks 2.18.}}
\SGARefTarget{sga3:xiv:2.18}
I do not know whether this conclusion remains true when \(k\) is
finite without assuming that \(\mathfrak g\) contains a regular
point; compare \SGARefLink{sga3:xiv:2.12}{2.12}. It can be proved when \(\mathfrak g\) is the
Lie algebra of an algebraic group \(G\) smooth over \(k\), at least
when \(G/\operatorname{radical}(G)\) is an ``adjoint'' semisimple
group, by using a result of Chevalley cited below \textup{(see the
Appendix)}. It is plausible that the result remains valid without
restriction on \(G\); for this it would suffice that the cited
result of Chevalley be proved for every semisimple algebraic group,
not necessarily an adjoint group.

% END INCLUDED SGA3_English_current_buildable_source_20260823/source/inputs/XIV/tex/components/08_expose_XIV_theorem216_proof_corollary217_remark218_en.tex

% BEGIN INCLUDED SGA3_English_current_buildable_source_20260823/source/inputs/XIV/tex/components/09_expose_XIV_theorem31_en.tex
% Authority: Expo14.pdf, local pp. 11--12 / combined pp. 845--846.
% Continuous English ends after the proof of Theorem 3.1.

\section{Subgroups of type \textup{(C)} of group preschemes over an
arbitrary prescheme}
\SGARefTarget{sga3:xiv:section:3}

\paragraph{\textbf{Theorem 3.1.}}
\SGARefTarget{sga3:xiv:3.1}
{\itshape Let \(S\) be a prescheme, let \(G\) be a smooth
\(S\)-group prescheme, and let \(\mathfrak g\) be its Lie algebra,
which is a locally free module of finite type over \(S\). Let
\(\mathfrak h\) be a Lie subalgebra of \(\mathfrak g\) that, as a
module, is locally a direct factor in \(\mathfrak g\), and suppose
that, for every \(s\in S\), the geometric fiber
\(\mathfrak h_{\bar s}\) contains a Cartan subalgebra of
\(\mathfrak g_{\bar s}\). Let \(a\) be a quasi-regular section of
\(\mathfrak g\) \textup{(\SGARefLink{sga3:xiv:2.5}{2.5})}. Then
\[
  M_a=\operatorname{Transp}_G(a,\mathfrak h),
\]
the subfunctor of \(G\) whose \(S'\)-valued points are the
\(g\in G(S')\) such that
\[
  \operatorname{ad}(g)\cdot a_{S'}
    \in\Gamma(S',\mathfrak h_{S'}),
\]
is represented by a closed subprescheme of \(G\), smooth over
\(S\), whose structure morphism to \(S\) is surjective.}
\footnote{Editors' note: This transporter is denoted
\(\operatorname{Transp}_G(a,\mathfrak h)\).}

\paragraph{\textit{Proof.}}
Consider the canonical morphism
\[
  \varphi\colon G\times_S W(\mathfrak h)\longrightarrow
  W(\mathfrak g),
  \qquad
  (g,x)\longmapsto\operatorname{ad}(g)\cdot x.
\]
Then \(M_a\) is \(S\)-isomorphic to \(\varphi^{-1}(a)\), the inverse
image under \(\varphi\) of \(a\), viewed as a section of
\(W(\mathfrak g)\) over \(S\). To prove that \(M_a\) is smooth, it
is enough to show that \(\varphi\) is smooth at the points of
\(G\times_S W(\mathfrak h)\) lying over \(\operatorname{Im}(a)\).
More generally, \(\varphi\) is smooth at every point lying over a
point of \(W(\mathfrak g)\) that is regular in its fiber
\(W(\mathfrak g(s))\) over \(S\).

Indeed, the source and target of \(\varphi\) are smooth over \(S\),
and hence are flat and locally of finite presentation over \(S\).
We may therefore check the assertion fiber by fiber. This reduces
us to the case in which \(S\) is the spectrum of an algebraically
closed field \(k\), \(G\) is a locally algebraic group over \(k\),
\(\mathfrak h\) is a Lie subalgebra of its Lie algebra
\(\mathfrak g\) containing a Cartan subalgebra of \(\mathfrak g\),
and \(a\) is a regular point of \(\mathfrak g\).

Since \(\varphi\) is a \(G\)-morphism, we may suppose that the point
of \(G\times W(\mathfrak h)\) under consideration is of the form
\((e,a)\). We may also suppose that \(G\) is connected, and hence
of finite type over \(k\). The assertion is then precisely
\SGARefLink{sga3:xiii:corollary:5:4}{XIII~5.4}.

It is also immediate that \(M_a\) is a closed subprescheme of \(G\)
of finite presentation over \(S\): it is the inverse image of
\(W(\mathfrak h)\) under the morphism
\[
  G\longrightarrow W(\mathfrak g),
  \qquad
  g\longmapsto\operatorname{ad}(g)\cdot a.
\]
Finally, the surjectivity of \(M_a\to S\) likewise reduces to the
case of an algebraically closed base field. By hypothesis,
\(\mathfrak h\) then contains a Cartan subalgebra
\(\mathfrak d'\). By the conjugacy theorem \SGARefLink{sga3:xiii:theorem:6:1}{XIII~6.1}\textup{(a)},
\(\mathfrak d'\) is conjugate to
\[
  \mathfrak d=\Nil(a,\mathfrak g).
\]
It follows that \(\mathfrak h\) contains a conjugate of \(a\).
This completes the proof.

% END INCLUDED SGA3_English_current_buildable_source_20260823/source/inputs/XIV/tex/components/09_expose_XIV_theorem31_en.tex

% BEGIN INCLUDED SGA3_English_current_buildable_source_20260823/source/inputs/XIV/tex/components/10_expose_XIV_corollary32_lemma33_en.tex
% Authority: Expo14.pdf, local p. 12 / combined p. 846.
% Continuous English ends immediately before Corollary 3.4.

\paragraph{\textbf{Corollary 3.2.}}
\SGARefTarget{sga3:xiv:3.2}
{\itshape Let \(G\) and \(\mathfrak g\) be as in \SGARefLink{sga3:xiv:3.1}{3.1}, with \(G\)
of finite type over \(S\), and let \(\mathfrak k\) and
\(\mathfrak h\) be two Lie subalgebras of \(\mathfrak g\) that are
locally direct factors as modules. Suppose that one of the following
two hypotheses holds:

\smallskip
\noindent
\textup{a)}
For every \(s\in S\), the geometric fiber
\(\mathfrak k_{\bar s}\) is nilpotent and contains a regular element
of \(\mathfrak g_{\bar s}\), while the geometric fiber
\(\mathfrak h_{\bar s}\) contains a Cartan subalgebra of
\(\mathfrak g_{\bar s}\).

\smallskip
\noindent
\textup{b)}
\(\mathfrak k\) is a Cartan subalgebra of \(\mathfrak g\).

\smallskip
\noindent
Then \(\operatorname{Transp}_G(\mathfrak k,\mathfrak h)\) is a
closed subprescheme of \(G\), smooth over \(S\). Moreover, in case
\textup{a)}, its structure morphism to \(S\) is surjective.}

The representability of the transporter by a closed subprescheme of
\(G\) of finite presentation is immediate and is left to the
reader. We first prove smoothness in case \textup{a)}. Suppose first
that \(\mathfrak k\) has a section \(a\) that is quasi-regular in
\(\mathfrak g\). It is then enough to apply \SGARefLink{sga3:xiv:3.1}{3.1} and the following
lemma.

\paragraph{\textbf{Lemma 3.3.}}
\SGARefTarget{sga3:xiv:3.3}
{\itshape Under the hypotheses of \SGARefLink{sga3:xiv:3.2}{3.2}\textup{(a)}, if \(a\) is a
section of \(\mathfrak k\) that is quasi-regular in
\(\mathfrak g\), then
\[
  \operatorname{Transp}_G(\mathfrak k,\mathfrak h)
  =
  \operatorname{Transp}_G(a,\mathfrak h).
\]}

Indeed, by the definitions it is enough to show that, if \(a\) is
also a section of \(\mathfrak h\), then
\(\mathfrak k\subset\mathfrak h\). Since \(\mathfrak k\) is
locally nilpotent by hypothesis, \SGARefLink{sga3:xiv:2.1}{2.1} gives
\[
  \mathfrak k\subset\Nil(a,\mathfrak g).
\]
On the other hand,
\[
  \Nil(a,\mathfrak g)\subset\mathfrak h,
\]
because
\[
  \operatorname{ad}(a)_{\mathfrak g/\mathfrak h}
\]
is injective; this can be checked fiber by fiber by \SGARefLink{sga3:xiii:proposition:4:8:item:b}{XIII~4.8(b)}.
The lemma follows.

In the general case, the usual standard argument first reduces us
to the case where \(S\) is affine noetherian, and then to the case
where \(S\) is local artinian, smoothness being an infinitesimal
property. By flat descent we may further suppose that the residue
field is infinite. The fiber of \(\mathfrak k\) then contains an
element regular in the corresponding fiber of \(\mathfrak g\).
Lifting this element to a section of \(\mathfrak k\) reduces us to
the preceding case.

Thus the transporter is smooth in case \textup{a)}. Its structure
morphism is surjective because this may be checked after reducing to
the case where \(S\) is the spectrum of an algebraically closed
field; \(\mathfrak k\) then contains a regular point of
\(\mathfrak g\), and \SGARefLink{sga3:xiv:3.3}{3.3} and \SGARefLink{sga3:xiv:3.1}{3.1} apply.

For case \textup{b)}, the definition of smoothness \textup{(\SGARefLink{sga3:xi:definition:1:1}{XI~1.1})}
reduces the assertion to the following lifting statement. Let \(S\)
be affine, let \(S_0\) be a closed subscheme defined by a
quasi-coherent nilpotent ideal \(J\), and let
\(g_0\in G(S_0)\) transport \(\mathfrak k_0\) into
\(\mathfrak h_0\). Then \(g_0\) lifts to an element \(g\in G(S)\)
that transports \(\mathfrak k\) into \(\mathfrak h\). The
hypothesis on \(g_0\) places this lifting problem under case
\textup{a)}, already proved. This completes the proof.

Of course, if in \SGARefLink{sga3:xiv:3.2}{3.2}\textup{(b)} the subalgebra \(\mathfrak h\)
satisfies the stronger hypothesis of \SGARefLink{sga3:xiv:3.1}{3.1}, then, and only then, the
structure morphism
\[
  \operatorname{Transp}_G(\mathfrak k,\mathfrak h)
  \longrightarrow S
\]
is surjective. Using Hensel's lemma \SGARefLink{sga3:xi:proposition:1:10}{XI~1.10}, one obtains from \SGARefLink{sga3:xiv:3.1}{3.1}
and \SGARefLink{sga3:xiv:3.2}{3.2} the following corollary.

% END INCLUDED SGA3_English_current_buildable_source_20260823/source/inputs/XIV/tex/components/10_expose_XIV_corollary32_lemma33_en.tex

% BEGIN INCLUDED SGA3_English_current_buildable_source_20260823/source/inputs/XIV/tex/components/11_expose_XIV_corollary34_through_definition38_en.tex
% Authority: Expo14.pdf, local pp. 12--13 / combined pp. 846--847.
% Continuous English ends immediately before Theorem 3.9.

\paragraph{\textbf{Corollary 3.4.}}
\SGARefTarget{sga3:xiv:3.4}
{\itshape Under the hypotheses of \SGARefLink{sga3:xiv:3.1}{3.1} for \(G\) and
\(\mathfrak h\), and assuming that \(G\) is of finite type over
\(S\), the following assertions hold:

\smallskip
\noindent
\textup{a)}
For every quasi-regular section \(a\) of \(\mathfrak g\), étale
locally there is a conjugate of \(a\) that is a section of
\(\mathfrak h\).

\smallskip
\noindent
\textup{b)}
For every Cartan subalgebra \(\mathfrak d\) of \(\mathfrak g\),
étale locally \(\mathfrak d\) is conjugate to a subalgebra of
\(\mathfrak h\).}

In particular, when \(\mathfrak h\) is itself a Cartan subalgebra
of \(\mathfrak g\), one obtains the following.

\paragraph{\textbf{Corollary 3.5.}}
\SGARefTarget{sga3:xiv:3.5}
{\itshape Let \(G\) be a smooth \(S\)-prescheme in groups of finite
type, let \(\mathfrak g\) be its Lie algebra, and let
\(\mathfrak d\) and \(\mathfrak d'\) be two Cartan subalgebras of
\(\mathfrak g\). Then
\(\operatorname{Transp}_G(\mathfrak d,\mathfrak d')\) is identical
to the strict transporter of \(\mathfrak d\) into
\(\mathfrak d'\), and is a closed subprescheme of \(G\), smooth
over \(S\), whose structure morphism is surjective. Étale locally,
\(\mathfrak d\) and \(\mathfrak d'\) are conjugate.}

Here the transporter is identical to the strict transporter simply
because \(\mathfrak d\) and \(\mathfrak d'\) are locally direct
factors in \(\mathfrak g\) and have the same rank at every point.
Thus \SGARefLink{sga3:xiv:3.5}{3.5} is a special case of \SGARefLink{sga3:xiv:3.4}{3.4}. Taking
\(\mathfrak d=\mathfrak d'\) gives the following.

\paragraph{\textbf{Corollary 3.6.}}
\SGARefTarget{sga3:xiv:3.6}
{\itshape Let \(G\) be as in \SGARefLink{sga3:xiv:3.5}{3.5}, and let \(\mathfrak d\) be a
Cartan subalgebra of \(\mathfrak g\). Then
\(\operatorname{Norm}_G(\mathfrak d)\) is a closed subprescheme in
groups of \(G\), smooth over \(S\), and its Lie algebra is
\(\mathfrak d\).}

Indeed, the last assertion says that \(\mathfrak d\) is its own
normalizer in \(\mathfrak g\), and this follows at once because it
is true fiber by fiber.

\paragraph{\textbf{Corollary 3.7.}}
\SGARefTarget{sga3:xiv:3.7}
{\itshape Let \(G\) and \(\mathfrak g\) be as in \SGARefLink{sga3:xiv:3.5}{3.5}. Then the
conditions \((\mathrm C_2)\), \((\mathrm C'_1)\), and
\((\mathrm C_1)\) of \SGARefLink{sga3:xiv:2.9}{2.9} are equivalent. In other words, if
\(\mathfrak g\) admits a Cartan subalgebra fpqc locally, then every
quasi-regular section of \(\mathfrak g\) is regular.}

Indeed, let \(a\) be a quasi-regular section; we prove that it is
regular. Since the question is fpqc local, we may suppose that
\(\mathfrak g\) admits a Cartan subalgebra \(\mathfrak d\). By
\SGARefLink{sga3:xiv:3.4}{3.4}\textup{(a)}, étale locally \(a\) is conjugate to a section of
\(\mathfrak d\). We are thereby reduced to the case in which
\(a\) is a section of \(\mathfrak d\), where the conclusion follows
immediately from the definition.

\paragraph{\textbf{Definition 3.8.}}
\SGARefTarget{sga3:xiv:3.8}
{\itshape Let \(G\) be a smooth prescheme in groups over a
prescheme \(S\). A \emph{subgroup of type \((\mathrm C)\)} of
\(G\) is a subprescheme in groups \(D\subset G\), smooth over
\(S\) and with connected fibers, such that
\[
  \mathfrak d=\operatorname{Lie}(D)
\]
is a Cartan subalgebra \textup{(\SGARefLink{sga3:xiv:2.4}{2.4})} of
\(\mathfrak g=\operatorname{Lie}(G)\). Equivalently, for every
\(s\in S\), the fiber \(D_s\) is a subgroup of type
\((\mathrm C)\) of the algebraic group \(G_s\)
\textup{(\SGARefLink{sga3:xiii:definition:6:2}{XIII~6.2})}.}

% END INCLUDED SGA3_English_current_buildable_source_20260823/source/inputs/XIV/tex/components/11_expose_XIV_corollary34_through_definition38_en.tex

% BEGIN INCLUDED SGA3_English_current_buildable_source_20260823/source/inputs/XIV/tex/components/12_expose_XIV_theorem39_through_corollary311_en.tex
% Authority: Expo14.pdf, local pp. 13--14 / combined pp. 847--848.
% Continuous English ends immediately before Proposition 3.12.

\paragraph{\textbf{Theorem 3.9.}}
\SGARefTarget{sga3:xiv:3.9}
{\itshape Let \(G\) be a smooth \(S\)-prescheme in groups of finite
presentation, and let \(\mathfrak g\) be its Lie algebra. Then:

\smallskip
\noindent
\textup{a)}
The map
\[
  D\longmapsto\mathfrak d=\operatorname{Lie}(D)
\]
establishes a bijective correspondence between the subgroups of
type \((\mathrm C)\) of \(G\) and the Cartan subalgebras of
\(\mathfrak g\).

\smallskip
\noindent
\textup{b)}
If \(D\) and \(\mathfrak d\) correspond, then
\[
  \operatorname{Norm}_G(D)=\operatorname{Norm}_G(\mathfrak d).
\]
This is a closed subprescheme of \(G\), smooth over \(S\), and
\[
  D=\operatorname{Norm}_G(D)^0
   =\operatorname{Norm}_G(\mathfrak d)^0.
\]

\smallskip
\noindent
\textup{c)}
Any two subgroups \(D\) and \(D'\) of type \((\mathrm C)\) in
\(G\) are conjugate étale locally.}

\emph{Proof.}
Let \(D\) be a subgroup of type \((\mathrm C)\) of \(G\), and let
\(\mathfrak d\) be its Lie algebra. Then
\[
  D\subset\operatorname{Norm}_G(\mathfrak d).
\]
By \SGARefLink{sga3:xiv:3.8}{Definition 3.8} and \SGARefLink{sga3:xiv:3.6}{Corollary 3.6}, this is an inclusion of smooth
preschemes in groups over \(S\) that induces an isomorphism on Lie
algebras. Since \(D\) has connected fibers, it follows that
\[
  D=\operatorname{Norm}_G(\mathfrak d)^0.
\]
Thus the map in \textup{a)} is injective. We prove that it is
surjective. Let \(\mathfrak d\) be a Cartan subalgebra of
\(\mathfrak g\). By \SGARefLink{sga3:xiv:3.6}{3.6},
\[
  N=\operatorname{Norm}_G(\mathfrak d)
\]
is a closed subprescheme in groups of \(G\), smooth over \(S\), and
has Lie algebra \(\mathfrak d\). Since \(G\) is of finite
presentation over \(S\), so is \(N\). Consequently, as noted in
XII after \SGARefLink{sga3:xii:7.3}{7.3}, the union of the connected components of the fibers
of \(N\) is the underlying set of an open subprescheme in groups
\(N^\circ\subset N\). Evidently \(N^\circ\) is a subgroup of type
\((\mathrm C)\) of \(G\) with Lie algebra \(\mathfrak d\). This
proves \textup{a)}. The first assertion of \textup{b)} follows at
once, and the formula
\[
  D=\operatorname{Norm}_G(\mathfrak d)^0
\]
has already been proved. Finally, \textup{c)} follows from
\textup{a)} and \SGARefLink{sga3:xiv:3.5}{3.5}.

\paragraph{\textbf{Corollary 3.10.}}
\SGARefTarget{sga3:xiv:3.10}
{\itshape Suppose that \(\mathfrak g\) admits a Cartan subalgebra
fpqc locally, or equivalently by \SGARefLink{sga3:xiv:3.9}{3.9}\textup{(a)}, that \(G\) admits
a subgroup of type \((\mathrm C)\) fpqc locally. Consider the
functor
\[
  \mathcal D\colon(\mathrm{Sch})^\circ/S\longrightarrow
  (\mathrm{Sets}),
\]
where \(\mathcal D(S')\) is the set of subgroups of type
\((\mathrm C)\) of \(G_{S'}\). This functor is representable by a
quasi-projective prescheme, smooth over \(S\), with connected
geometric fibers. If \(S=\operatorname{Spec}(k)\), so that \(G\)
is a smooth algebraic group over \(k\), and if \(\mathfrak g\)
admits a regular point---a condition automatically satisfied when
\(k\) is infinite---then \(\mathcal D\) is a rational variety over
\(k\).}

Indeed, by \SGARefLink{sga3:xiv:3.9}{3.9}\textup{(a)} the functor \(\mathcal D\) is
canonically isomorphic to the functor considered in \SGARefLink{sga3:xiv:2.16}{2.16}. On the
other hand, by \SGARefLink{sga3:xiv:3.7}{3.7} condition \((\mathrm C_2)\) is satisfied.
Therefore \SGARefLink{sga3:xiv:3.10}{3.10} follows from \SGARefLink{sga3:xiv:2.16}{2.16} and \SGARefLink{sga3:xiv:2.17}{2.17}.

\paragraph{\textbf{Corollary 3.11.}}
\SGARefTarget{sga3:xiv:3.11}
{\itshape Let \(D\) be a subgroup of type \((\mathrm C)\) of
\(G\), and let \(N\) be its normalizer in \(G\). Then the quotient
sheaf \(G/N\) is canonically isomorphic to the functor
\(\mathcal D\) of \SGARefLink{sga3:xiv:3.10}{3.10}. It is therefore representable by a
quasi-projective prescheme, smooth over \(S\), with connected
geometric fibers.}

% END INCLUDED SGA3_English_current_buildable_source_20260823/source/inputs/XIV/tex/components/12_expose_XIV_theorem39_through_corollary311_en.tex

% BEGIN INCLUDED SGA3_English_current_buildable_source_20260823/source/inputs/XIV/tex/components/13_expose_XIV_proposition312_through_corollary315_en.tex
% Authority: Expo14.pdf, local pp. 14--15 / combined pp. 848--849.
% Continuous English ends immediately before Corollary 3.16.

\paragraph{\textbf{Proposition 3.12.}}
\SGARefTarget{sga3:xiv:3.12}
{\itshape Let \(G\) be a smooth \(S\)-prescheme in groups of finite
presentation over \(S\). Let \(H\) and \(K\) be two subpreschemes
in groups, smooth and of finite presentation over \(S\), with
\(K\) having connected fibers, and let
\(\mathfrak g,\mathfrak h,\mathfrak k\) be the corresponding Lie
algebras. Suppose that one of the following conditions holds on
their geometric fibers:

\smallskip
\noindent
\textup{a)}
For every \(s\in S\), \(\mathfrak k_{\bar s}\) contains a regular
element of \(\mathfrak g_{\bar s}\), and
\(\mathfrak h_{\bar s}\) contains a Cartan subalgebra of
\(\mathfrak g_{\bar s}\).

\smallskip
\noindent
\textup{b)}
For every \(s\in S\), \(\mathfrak k_{\bar s}\) contains a Cartan
subalgebra of \(\mathfrak g_{\bar s}\).

\smallskip
\noindent
Under these conditions, \(H\supset K\) if and only if
\(\mathfrak h\supset\mathfrak k\).}

It remains only to prove that
\(\mathfrak h\supset\mathfrak k\) implies \(H\supset K\). In case
\textup{b)}, the inclusion
\(\mathfrak h\supset\mathfrak k\) places us under the hypotheses of
case \textup{a)}, so it is enough to prove \textup{a)}. Proceeding
as in \SGARefLink{sga3:xiv:3.2}{3.2}\textup{(a)}, by reduction to the case where \(S\) is
local artinian, we reduce to the case in which there exists a
section \(a\) of \(\mathfrak k\) that is quasi-regular in
\(\mathfrak g\). Arguing then as in \SGARefLink{sga3:xiii:corollary:5:5}{XIII~5.5} reduces us to the
following statement.

\paragraph{\textbf{Corollary 3.13.}}
\SGARefTarget{sga3:xiv:3.13}
{\itshape Let \(G\) be a smooth \(S\)-prescheme in groups of finite
presentation over \(S\), let \(H\subset G\) be a subprescheme in
groups smooth and of finite presentation over \(S\), and let
\(\mathfrak g\) and \(\mathfrak h\) be their Lie algebras. Let
\(a\) be a section of \(\mathfrak h\) that is quasi-regular in
\(\mathfrak g\). Suppose that for every \(s\in S\), the geometric
fiber \(\mathfrak h_{\bar s}\) contains a Cartan subalgebra of
\(\mathfrak g_{\bar s}\). Put
\[
  M_a=\operatorname{Transp}_G(a,\mathfrak h).
\]
This is a closed subprescheme of \(G\), smooth over \(S\)
\textup{(cf.\ \SGARefLink{sga3:xiv:3.1}{3.1})}. Hence \(M_a^0\), the union in the fibers of
\(M_a\) of the connected components containing the identity, is
an open subset of \(M_a\), endowed with the structure induced from
\(M_a\). Then
\[
  H^0=M_a^0.
\]}

Clearly \(H\subset M_a\), and therefore
\(H^0\subset M_a^0\). Since this is an inclusion of preschemes
smooth over \(S\), it is enough to check equality on the fibers.
This reduces us to the case where \(S\) is the spectrum of an
algebraically closed field, which was treated in \SGARefLink{sga3:xiii:corollary:5:4}{XIII~5.4}.

\paragraph{\textbf{Corollary 3.14.}}
\SGARefTarget{sga3:xiv:3.14}
{\itshape Let \(G\) be a smooth \(S\)-prescheme in groups of finite
presentation, let \(D\) be a subgroup of type \((\mathrm C)\) of
\(G\), and let \(H\subset G\) be a subprescheme in groups smooth
and of finite presentation over \(S\). If \(\mathfrak d\) and
\(\mathfrak h\) are their Lie algebras, then
\[
  \operatorname{Transp}_G(D,H)
  =
  \operatorname{Transp}_G(\mathfrak d,\mathfrak h),
\]
and this functor is representable by a closed subprescheme of
\(G\), smooth over \(S\).}

Indeed, the identity between the two transporters follows from
\SGARefLink{sga3:xiv:3.12}{3.12}\textup{(b)}, after which \SGARefLink{sga3:xiv:3.2}{3.2} applies.

\paragraph{\textbf{Corollary 3.15.}}
\SGARefTarget{sga3:xiv:3.15}
{\itshape Let \(G\) and \(H\) be as in \SGARefLink{sga3:xiv:3.14}{3.14}, and suppose that for
every \(s\in S\), the geometric fiber \(\mathfrak h_{\bar s}\)
contains a Cartan subalgebra of \(\mathfrak g_{\bar s}\). Suppose
also that \(\mathfrak g\) admits a Cartan subalgebra fpqc locally.
Then, étale locally, \(H\) contains a subgroup of type
\((\mathrm C)\) of \(G\).}

By \SGARefLink{sga3:xiv:3.7}{3.7} and \SGARefLink{sga3:xiv:3.9}{3.9}\textup{(a)}, \(G\) admits a subgroup of type
\((\mathrm C)\) étale locally. We may therefore suppose that \(G\)
admits such a subgroup \(D\). The hypothesis on \(\mathfrak h\)
then also says that the structure morphism of the transporter in
\SGARefLink{sga3:xiv:3.14}{3.14} is surjective, by the conjugacy theorem \SGARefLink{sga3:xiii:theorem:6:1}{XIII~6.1}\textup{(a)}.
The conclusion follows from Hensel's lemma \SGARefLink{sga3:xi:proposition:1:10}{XI~1.10}.

% END INCLUDED SGA3_English_current_buildable_source_20260823/source/inputs/XIV/tex/components/13_expose_XIV_proposition312_through_corollary315_en.tex

% BEGIN INCLUDED SGA3_English_current_buildable_source_20260823/source/inputs/XIV/tex/components/14_expose_XIV_corollary316_remarks317_en.tex
% Authority: Expo14.pdf, local pp. 15--16 / combined pp. 849--850.
% Continuous English ends immediately before Theorem 3.18.

\paragraph{\textbf{Corollary 3.16.}}
\SGARefTarget{sga3:xiv:3.16}
{\itshape Let \(G,H,K\) be as in \SGARefLink{sga3:xiv:3.12}{3.12}\textup{(a)}, and suppose in
addition that every geometric fiber \(\mathfrak k_{\bar s}\) is
nilpotent, that is, that \(\mathfrak k\) is locally nilpotent. Then
\[
  \operatorname{Transp}_G(K,H)
  =
  \operatorname{Transp}_G(\mathfrak k,\mathfrak h),
\]
and this functor is representable by a closed subprescheme of
\(G\), smooth over \(S\), whose structure morphism to \(S\) is
surjective. Étale locally, \(H\) contains a subgroup conjugate to
\(K\).}

The identity between the two transporters is again contained in
\SGARefLink{sga3:xiv:3.12}{3.12}\textup{(a)}. The assertion about its structure is precisely
\SGARefLink{sga3:xiv:3.2}{3.2}\textup{(a)}, and the last assertion of \SGARefLink{sga3:xiv:3.16}{3.16} then follows from
Hensel's lemma.

\paragraph{\textbf{Remarks 3.17.}}
\SGARefTarget{sga3:xiv:3.17}
\textup{a)}
In \SGARefLink{sga3:xiv:3.12}{3.12} and \SGARefLink{sga3:xiv:3.16}{3.16}, the hypothesis that \(K\) be smooth over \(S\)
may be replaced by the weaker hypothesis that the sheaf
\[
  e_K^*(\Omega^1_{K/S})=\omega^1_K
\]
of relative \(1\)-differentials along the identity section be
locally free. In this form, \SGARefLink{sga3:xiv:3.12}{3.12} includes \SGARefLink{sga3:xiii:corollary:5:5}{XIII~5.5}.

\smallskip
\noindent
\textup{b)}
Let \(G,\mathfrak g,\mathfrak h\) be as in \SGARefLink{sga3:xiv:3.1}{3.1}, with \(G\) of
finite presentation over \(S\). Then
\[
  N=\operatorname{Norm}_G(\mathfrak h)
\]
need not be smooth over \(S\) along the identity section.
Equivalently, there need not exist a subprescheme in groups
\(H\subset G\), smooth over \(S\), whose Lie algebra is
\(\mathfrak h\), even when \(S\) is the spectrum of a field. When
such an \(H\) exists, taking \(H\) with connected fibers gives
\[
  N=\operatorname{Norm}_G(H);
\]
I do not know whether \(N\) is then smooth over \(S\). In this
question one may plainly reduce to the case where \(S\) is local
artinian.

% END INCLUDED SGA3_English_current_buildable_source_20260823/source/inputs/XIV/tex/components/14_expose_XIV_corollary316_remarks317_en.tex

% BEGIN INCLUDED SGA3_English_current_buildable_source_20260823/source/inputs/XIV/tex/components/15_expose_XIV_theorem318_corollary319_en.tex
% Authority: Expo14.pdf, local pp. 16--17 / combined pp. 850--851.
% Continuous English ends immediately before Corollary 3.20.

To conclude this section, we examine the case in which \(G\) is
``semisimple'' over \(S\).

\paragraph{\textbf{Theorem 3.18.}}
\SGARefTarget{sga3:xiv:3.18}
{\itshape Let \(S\) be a prescheme and let \(G\) be a smooth
\(S\)-prescheme in groups whose geometric fibers are semisimple
``adjoint'' groups, that is, semisimple groups whose reductive
center \textup{(\SGARefLink{sga3:xii:4.1}{XII~4.1} and \SGARefLink{sga3:xiv:4.4}{4.4})} is the identity group. Then the
subgroups of type \((\mathrm C)\) of \(G\) are precisely its
maximal tori, and hence also precisely its Cartan subgroups
\textup{(\SGARefLink{sga3:xii:3.1}{XII~3.1})}.}

In view of the definitions, it is enough to treat the case where
\(S\) is the spectrum of an algebraically closed field. We must
then show that, for a maximal torus \(T\) of \(G\), its Lie algebra
\(\mathfrak t\) is a Cartan subalgebra of \(\mathfrak g\). Put
\(r=\dim T=\operatorname{rank}\mathfrak t\). Taking into account
the inequality
\[
  \operatorname{nil\,rank}(\mathfrak g)
  \geq
  \operatorname{nil\,rank}(G)
  =
  \dim T
  =
  \operatorname{rank}\mathfrak t
  =
  r,
\]
it is enough to find \(x\in\mathfrak t\) such that
\[
  \dim\Nil(x,\mathfrak g)=r.
\]
Since \(\mathfrak t\) is abelian, and therefore nilpotent, this is
equivalent to finding \(a\in\mathfrak t\) for which
\[
  \operatorname{ad}(a)_{\mathfrak g/\mathfrak t}
\]
is injective \textup{(\SGARefLink{sga3:xiii:corollary:5:7}{XIII~5.7}\textup{(a)})}.

Consider the characters \(\alpha\) of \(T\) occurring in the
representation of \(T\) induced by the adjoint representation of
\(G\). The structure theory of the semisimple group \(G\)
\textup{\emph{(Bible)}, 13, \SGARefLink{sga3:xiv:7:lemma1}{Theorem~1}\textup{(a)}, and Theorem~3,
\SGARefLink{sga3:xiv:7:lemma2}{Corollary~2}}, or more precisely the ``big cell'' of \(G\), shows
that the eigenspace of \(\mathfrak g\) for the trivial character is
exactly \(\mathfrak t\). The other eigenspaces are
one-dimensional, and their characters are precisely the roots of
\(G\) relative to \(T\). Here the big cell is built from \(T\) and
the subgroups \(P_\alpha\), each isomorphic to the additive group
\(\mathbf G_a\), stable under \(T\), and corresponding to the root
characters.

The reductive center of \(G\) is the intersection of the kernels
of the characters of \(T\) occurring in the adjoint representation
\textup{(\SGARefLink{sga3:xii:4.8}{XII~4.8})}. Thus the assertion that \(G\) is adjoint says
that its roots generate the lattice
\[
  M=\operatorname{Hom}(T,\mathbf G_m).
\]
A familiar lemma from root theory says that every root belongs to
a system of simple roots, hence to a basis of the group generated
by the roots, and therefore to a basis of the dual lattice of
\(T\).\footnote{The re-edition replaces \(\widehat T\) by \(T\)
at this point.} We obtain the following.

\paragraph{\textbf{Corollary 3.19.}}
\SGARefTarget{sga3:xiv:3.19}
{\itshape Let \(G\) be an adjoint semisimple algebraic group over
an algebraically closed field \(k\), and let \(T\) be a maximal
torus of \(G\). For every root
\[
  \alpha\colon T\longrightarrow\mathbf G_m
\]
of \(G\) relative to \(T\), the corresponding homomorphism
\[
  \alpha'\colon\mathfrak t\longrightarrow k
\]
is nonzero.}

This assertion is essentially equivalent to \SGARefLink{sga3:xiv:3.18}{Theorem 3.18}. Indeed,
for \(t\in\mathfrak t\), the endomorphism
\(\operatorname{ad}(t)\) is semisimple, and its eigenvalues on
\(\mathfrak g/\mathfrak t\) are precisely the numbers
\(\alpha'(t)\). Hence
\(\operatorname{ad}(t)_{\mathfrak g/\mathfrak t}\) is injective if
and only if all the \(\alpha'(t)\) are nonzero; and such a
\(t\in\mathfrak t\) exists if and only if every \(\alpha'\) is
nonzero.

% END INCLUDED SGA3_English_current_buildable_source_20260823/source/inputs/XIV/tex/components/15_expose_XIV_theorem318_corollary319_en.tex

% BEGIN INCLUDED SGA3_English_current_buildable_source_20260823/source/inputs/XIV/tex/components/16_expose_XIV_corollary320_remarks321_en.tex
% Authority: Expo14.pdf, local pp. 17--19 / combined pp. 851--853.
% This closes Section 3; continuous English next begins at Section 4.

\paragraph{\textbf{Corollary 3.20.}}
\SGARefTarget{sga3:xiv:3.20}
{\itshape Let \(S\) be a prescheme and let \(G\) be a smooth
\(S\)-prescheme in groups of finite presentation over \(S\), whose
geometric fibers are connected reductive algebraic groups, that is,
extensions of a semisimple group by a torus. For every \(s\in S\)
there is an open neighborhood \(U\) of \(s\) such that
\(G|_U\) admits a maximal torus.}\footnote{This result, as well as
\SGARefLink{sga3:xiv:2.11}{2.11} on which it rests, immediately generalizes to the case in
which \(s\) is replaced by a finite subset of \(S\) contained in
an affine open.}

We shall see in Exposé XVI that the hypothesis on \(G\) implies
that \(G\) is affine over \(S\) and has locally constant reductive
rank. Hence, by \SGARefLink{sga3:xii:4.7}{XII~4.7}\textup{\SGARefLink{sga3:xii:4.7:item:c}{(c)}}, \(G\) admits a reductive center
\(Z\), the quotient
\[
  G'=G/Z
\]
is a smooth affine group over \(S\) whose reductive center is the
identity group, and the maximal tori of \(G\) correspond
bijectively to those of \(G'\). Moreover, the geometric fibers of
\(G'\) are plainly connected semisimple groups, and are adjoint by
definition because their reductive center is trivial.

It is therefore enough to treat the case where \(G\) is
semisimple and adjoint. By \SGARefLink{sga3:xiv:3.18}{3.18}, the problem reduces to finding an
open neighborhood \(U\) of \(s\) and a subgroup of type
\((\mathrm C)\) of \(G|_U\), or equivalently by \SGARefLink{sga3:xiv:3.9}{3.9}\textup{(a)}, a
Cartan subalgebra of \(\mathfrak g|_U\). This is possible when
\(k(s)\) is infinite: by \SGARefLink{sga3:xiv:3.7}{3.7}, \(\mathfrak g\) satisfies condition
\((\mathrm C_2)\) of \SGARefLink{sga3:xiv:2.9}{2.9}, and \SGARefLink{sga3:xiv:2.11}{2.11}\textup{(b)} applies.

In fact, \SGARefLink{sga3:xiv:3.20}{3.20} remains true without assuming that \(k(s)\) is
infinite. By the preceding argument, it is enough to know that, for
every adjoint semisimple group \(G\) over a finite field \(k\), the
Lie algebra \(\mathfrak g\) contains a regular element. Chevalley
proved this assertion using the properties of a Coxeter element of
the Weyl group; see the appendix below by J.-P. Serre.

\paragraph{\textbf{Remarks 3.21.}}
\SGARefTarget{sga3:xiv:3.21}
\textup{a)}
Assertion \SGARefLink{sga3:xiv:3.20}{3.20} remains valid, with essentially the same proof, if
\(G\) is replaced by a closed subprescheme in groups \(H\), smooth
over \(S\), whose rank is everywhere the same as that of \(G\)
\textup{(for example, a ``parabolic subgroup'' of \(G\))}, provided
that \(k(s)\) is infinite. I do not know whether the infinitude
hypothesis is again superfluous. It can at least be removed when
\(H\) is parabolic: constructing the radical \(U\) of \(H\) and
the quotient \(H/U\) reduces the question to the semisimple case.
Unfortunately, the method of regular elements seems powerless
here. One readily constructs examples---for instance, the
projective group with its ``standard'' Borel subgroup---in which
the Lie algebra \(\mathfrak h\) of \(H\) contains no regular
element.

\smallskip
\noindent
\textup{b)}
When \(k(s)\) is infinite, the proof of \SGARefLink{sga3:xiv:3.20}{3.20} in fact gives the
following sharper statement, using \SGARefLink{sga3:xiv:3.9}{3.9}\textup{(b)}: every maximal
torus \(T_0\) of \(G_s\) is induced by a maximal torus \(T\) over
an open neighborhood of \(s\). I do not know whether this remains
true without the infinitude hypothesis. The evident difficulty is
that the Lie algebra \(\mathfrak t_0\) of \(T_0\) need not contain
a regular element of the Lie algebra \(\mathfrak g_0\) of the
fiber \(G_s\).

An affirmative answer would imply the following assertion, which
is known only when the residue field is infinite or when \(A\) is
separated and complete. Let \(A\) be a local ring with residue
field \(k\), and let \(M\) be an ``Azumaya algebra'' over \(A\),
meaning an algebra that is a finite free \(A\)-module and whose
reduction
\[
  M_0=M\otimes_A k
\]
is a central simple algebra over \(k\). Let \(D_0\subset M_0\) be
a commutative subalgebra separable over \(k\), with
\[
  [M_0:k]=[D_0:k]^2.
\]
Then there should exist a commutative subalgebra \(D\subset M\)
that is a direct summand as an \(A\)-module and satisfies
\[
  D\otimes_A k=D_0.
\]

The datum of \(M\) is equivalent to that of a principal homogeneous
bundle under the projective group \(\operatorname{PGL}(n)_A\), and
hence to an inner twisted form \(G\) of \(\operatorname{PGL}(n)\).
The maximal tori of \(G\) correspond bijectively to the commutative
subalgebras \(D\subset M\) that are étale of rank \(n\) over \(A\).

\smallskip
\noindent
\textup{c)}
Applying \SGARefLink{sga3:xiv:3.20}{3.20} to the centralizer of a subtorus \(Q\) of a reductive
\(S\)-prescheme in groups \(G\), one sees that every such \(Q\) is
contained, Zariski locally, in a maximal torus of \(G\).

% END INCLUDED SGA3_English_current_buildable_source_20260823/source/inputs/XIV/tex/components/16_expose_XIV_corollary320_remarks321_en.tex

% BEGIN INCLUDED SGA3_English_current_buildable_source_20260823/source/inputs/XIV/tex/components/17_expose_XIV_section4_definition41_through_definition45_en.tex
% Authority: Expo14.pdf, local pp. 18--19 / combined pp. 852--853.
% Exact next cursor: Remarks 4.6 on local p. 19 / combined p. 853.

\section{A Digression on Borel Subgroups}
\SGARefTarget{sga3:xiv:4}

\paragraph{\textbf{Definition 4.1.}}
\SGARefTarget{sga3:xiv:4.1}
{\itshape Let \(G\) be a smooth algebraic group over an algebraically
closed field. A \emph{Borel subgroup} of \(G\) is a smooth, solvable,
connected algebraic subgroup that is maximal with these properties.}

When \(G\) is affine, this agrees with the terminology of
\textup{\emph{(Bible)}, 6, \SGARefLink{sga3:xiv:7:lemma1}{Definition~1}}. Observe at once that if
\(Z\) is a connected smooth subgroup of \(G\) contained in its center
\textup{(or, more generally, if \(Z\) is solvable and normal)}, then,
for every Borel subgroup \(B\) of \(G\), the image \(BZ\) of
\(B\times Z\) under the morphism
\[
  (b,z)\longmapsto bz
\]
from \(B\times Z\) to \(G\) is a smooth, solvable, connected subgroup
of \(G\) containing \(B\). It is therefore equal to \(B\). Thus \(B\)
contains \(Z\), and consequently \(B\) is the inverse image of an
algebraic subgroup \(B'\) of \(G'=G/Z\); it is immediate that \(B'\)
is a Borel subgroup of \(G'\).

Taking
\[
  Z=\Centr(G^0)_{\mathrm{red}}^0,
\]
the quotient \(G'=G/Z\) is affine by \SGARefLink{sga3:xii:6.1}{XII~6.1}. Hence the Borel
subgroups of \(G'\) are conjugate, and \(G'/B'\) is a projective
variety for every such \(B'\)
\textup{\emph{(Bible)}, 6, \SGARefLink{sga3:xiv:7:lemma4}{Theorem~4}}. Consequently:

\paragraph{\textbf{Proposition 4.2.}}
\SGARefTarget{sga3:xiv:4.2}
{\itshape Let \(G\) be as in \SGARefLink{sga3:xiv:4.1}{4.1}. Then the Borel subgroups of \(G\)
are conjugate. If \(B\) is a Borel subgroup, then \(G/B\) is a
projective variety. The maximal tori of \(B\) are maximal tori of
\(G\); and, when \(G\) is connected, the Cartan subgroups of \(B\)
are Cartan subgroups of \(G\).}

It remains to prove the last assertion, and we may plainly suppose
that \(G=G^0\). For Cartan subgroups, it follows from the analogous
assertion in \(G'\)
\textup{\emph{(Bible)}, 6, \SGARefLink{sga3:xiv:7:lemma4}{Theorem~4}, \SGARefLink{sga3:xiv:7:lemma4}{Corollary~4}}
and \SGARefLink{sga3:xii:6.6}{XII~6.6}\textup{\SGARefLink{sga3:xii:6.6:item:e}{(e)}}. The assertion for maximal tori follows from
this one, since by \SGARefLink{sga3:xii:6.6}{XII~6.6}\textup{\SGARefLink{sga3:xii:6.6:item:c}{(c)}} the maximal tori of a smooth
algebraic group are precisely the maximal tori of its Cartan
subgroups.

\paragraph{\textbf{Corollary 4.3.}}
\SGARefTarget{sga3:xiv:4.3}
{\itshape Suppose that \(G\) is connected. Then every element of
\(G\) is contained in a Borel subgroup of \(G\).}

This reduces to the same assertion in \(G'\), where it is well known
\textup{\emph{(Bible)}, 6, Theorem~5\textup{(d)}}.

\paragraph{\textbf{Corollary 4.4.}}
\SGARefTarget{sga3:xiv:4.4}
{\itshape Let \(B\) be a Borel subgroup of \(G\), let \(C\) be a
Cartan subgroup of \(B\), and let \(N\) be its normalizer in \(G\).
Then
\[
  N\cap B=C.
\]}

Indeed,
\[
  N\cap B=\operatorname{Norm}_B(C),
\]
so it remains to show that, when \(G\) is connected and solvable, a
Cartan subgroup \(C\) is its own connected normalizer. With the
preceding notation, \(C\) is the inverse image of a Cartan subgroup
\(C'\) of \(G'\), and we are therefore reduced to the case in which
\(G\) is affine. The normalizer of a Cartan subgroup is smooth
\textup{(\(C\) being its own connected normalizer; see, for example,
\SGARefLink{sga3:xii:6.6}{XII~6.6}\textup{\SGARefLink{sga3:xii:6.6:item:c}{(c)}})}. It is thus enough to check that \(C\) and
\(N\) have the same \(k\)-valued points, which is exactly
\textup{\emph{(Bible)}, Theorem~6\textup{\SGARefLink{sga3:xii:6.6:item:d}{(d)}}}.

\paragraph{\textbf{Definition 4.5.}}
\SGARefTarget{sga3:xiv:4.5}
{\itshape Let \(G\) be a smooth prescheme in groups of finite
presentation over a prescheme \(S\). A \emph{Borel subgroup} of
\(G\) is a smooth subprescheme in groups \(B\) of \(G\), of finite
presentation, such that for every \(s\in S\) the geometric fiber
\(B_{\bar s}\) is a Borel subgroup of \(G_{\bar s}\).}

As is immediately verified, this notion is stable under base change
and is local for the fpqc topology. Indeed, if \(k'\) is an
algebraically closed extension of an algebraically closed field
\(k\), then an algebraic subgroup \(B\) of a smooth algebraic group
\(G\) over \(k\) is a Borel subgroup of \(G\) if and only if
\(B_{k'}\) is a Borel subgroup of \(G_{k'}\).

It follows from the definition that, if \(G\) is a smooth algebraic
group over an arbitrary field \(k\) and \(B\) is a Borel subgroup
of \(G\), then \(G/B\) is a projective variety; every maximal torus
\(T\) of \(B\) is a maximal torus of \(G\); and the normalizer of
\(T\) in \(B\) equals its centralizer \(C\), which is a Cartan
subgroup of \(G\) when \(G\) is connected.

% END INCLUDED SGA3_English_current_buildable_source_20260823/source/inputs/XIV/tex/components/17_expose_XIV_section4_definition41_through_definition45_en.tex

% BEGIN INCLUDED SGA3_English_current_buildable_source_20260823/source/inputs/XIV/tex/components/18_expose_XIV_remarks46_through_remarks481_en.tex
% Authority: Expo14.pdf, local pp. 19--20 / combined pp. 853--854.
% The source has no item numbered 4.7.
% Exact next cursor: the paragraph preceding Proposition 4.9,
% local p. 20 / combined p. 854.

\paragraph{\textbf{Remarks 4.6.}}
\SGARefTarget{sga3:xiv:4.6}
Unfortunately, it is no longer true in general that two Borel
subgroups of \(G\) are locally conjugate for the fpqc topology.
This can fail even when \(G\) is affine over \(S\) and \(S\) is
the spectrum of the algebra of dual numbers over an algebraically
closed field \(k\).

As a consequence of this regrettable fact, let us point out that
if \(G\) is a smooth, affine, connected algebraic group over a
nonperfect field \(k\), it is in general impossible to define
naturally a homogeneous space \(D\) under \(G\) that plays the
role of a flag variety, that is, of the variety of Borel subgroups
of \(G\). Over the algebraic closure \(\bar k\) of \(k\), such a
space would have to be isomorphic to \(G_{\bar k}/B\), where \(B\)
is a Borel subgroup of \(G_{\bar k}\).

Indeed, suppose that the quotient of \(G_{\bar k}\) by its radical
\(\bar R\) is an adjoint semisimple group. The kernel of
\[
  G_{\bar k}\longrightarrow
  \operatorname{Aut}_{\bar k}(\bar D)
\]
is then the radical of \(G_{\bar k}\). Consequently, if \(\bar D\)
came from a homogeneous space \(D\) under \(G\), the radical
\(\bar R\) would come from a subgroup \(R\) of \(G\). But examples
are easily constructed for which \(G_{\bar k}/\bar R\) is adjoint
while \(\bar R\) is not defined over \(k\).

Under these conditions, the functor
\[
  \mathfrak B\colon(\operatorname{Sch}/S)^\circ
    \longrightarrow \operatorname{Ens},
  \qquad
  \mathfrak B(S')=
  \{\text{Borel subgroups of }G_{S'}\},
\]
is not representable by a smooth \(S\)-prescheme. From the
infinitesimal point of view \textup{(III, \S3)}, the failure of the
conjugacy theorem is expressed by the fact that, if \(B\) is a
Borel subgroup of the smooth algebraic group \(G\), the cohomology
group
\[
  H^1(B,\mathfrak g/\mathfrak b)
\]
may be nonzero.\footnote{Here \(\mathfrak b\) denotes the Lie
algebra of \(B\).}

We shall see, on the other hand, in a later exposé that when \(G\)
is semisimple, or more generally reductive, such unpleasant
phenomena do not occur. These phenomena, together no doubt with
the absence of good existence theorems, explain why Borel
subgroups play only a relatively unobtrusive role in the
scheme-theoretic study of general group schemes, whereas they
dominate the theory of semisimple group schemes in the later
exposés.

\paragraph{\textbf{Proposition 4.8.}}
\SGARefTarget{sga3:xiv:4.8}
\footnote{There is no item numbered 4.7 in the source.}
{\itshape Let \(G\) be a smooth \(S\)-prescheme in groups of finite
presentation with connected fibers, and let \(B\) be a Borel
subgroup of \(G\). Then \(B\) is its own normalizer and is a closed
subprescheme of \(G\).}

By \SGARefLink{sga3:xii:7.10}{XII~7.10}, it is enough to prove that, over an algebraically
closed field \(k\), every element of \(G(k)\) that normalizes \(B\)
belongs to \(B(k)\). When \(G\) is affine this is a fundamental
result of Chevalley
\textup{\emph{(Bible)}, 9, \SGARefLink{sga3:xiv:7:lemma1}{Theorem~1}}; the general case reduces
immediately to it by the reduction already used in \SGARefLink{sga3:xiv:4.2}{4.2}.

\paragraph{\textbf{Remarks 4.8.1.}}
\SGARefTarget{sga3:xiv:4.8.1}
\SGARefLink{sga3:xiv:4.5}{Definition~4.5} may be generalized by also introducing the notion
of a \emph{parabolic subgroup} of \(G\). This means a subprescheme
in groups \(P\) of \(G\), smooth and of finite presentation over
\(S\), such that for every \(s\in S\) the geometric fiber
\(P_{\bar s}\) is a parabolic subgroup of \(G_{\bar s}\), that is,
contains a Borel subgroup of \(G_{\bar s}\).

\SGARefLink{sga3:xiv:4.8}{Proposition~4.8} extends to a parabolic subgroup \(P\) of \(G\), by
the same reduction to the known set-theoretic assertion. We note
the following consequence \textup{(compare Exposé XVI)}. \SGAInlineRefTarget{sga3:xiv:paragraph:4-8-bis}If \(P\)
is a parabolic subgroup of \(G\), then \(G/P\) is representable by
a quasi-projective prescheme of finite presentation over \(S\);
here \(G\) is assumed to have connected fibers. Moreover, \(G/P\)
is plainly smooth over \(S\), and its geometric fibers are
connected and proper. It follows readily from EGA~III, 5.5.1 that
\[
  D=G/P
\]
is proper, hence projective, over \(S\). If its relative dimension
is \(n\), it is known that the inverse of the invertible sheaf
\(\Omega^n_{D/S}\) induces ample sheaves on the geometric fibers of
\(D/S\); hence, by EGA~III, 4.7.1,
\[
  \bigl(\Omega^n_{D/S}\bigr)^{-1}
\]
is ample on \(D\) relative to \(S\).

% END INCLUDED SGA3_English_current_buildable_source_20260823/source/inputs/XIV/tex/components/18_expose_XIV_remarks46_through_remarks481_en.tex

% BEGIN INCLUDED SGA3_English_current_buildable_source_20260823/source/inputs/XIV/tex/components/19_expose_XIV_proposition49_corollary410_en.tex
% Authority: Expo14.pdf, local pp. 20--21 / combined pp. 854--855.
% Exact next cursor: Theorem 4.11, local p. 21 / combined p. 855.

It is readily seen, by reduction to the affine case and to an
algebraically closed base field, that if
\[
  u\colon G\longrightarrow G'
\]
is an epimorphism of smooth algebraic groups, then \(u(B)=B'\) is a
Borel subgroup of \(G'\) for every Borel subgroup \(B\) of \(G\).
We now consider the case in which this gives a bijective
correspondence between the Borel subgroups of \(G\) and those of
\(G'\).

\paragraph{\textbf{Proposition 4.9.}}
\SGARefTarget{sga3:xiv:4.9}
{\itshape Let \(G\) and \(G'\) be smooth \(S\)-preschemes in groups
of finite presentation with connected fibers, and let
\[
  u\colon G\longrightarrow G'
\]
be a faithfully flat, that is, surjective, homomorphism of
groups.\footnote{The editors ask whether a reference should be
added here.} Put \(N=\ker u\), and suppose that one of the following
two conditions holds:

\smallskip
\noindent
\textup{a)} \(N\) is central in \(G\).

\smallskip
\noindent
\textup{b)} \(S=\operatorname{Spec}(k)\) for a field \(k\), and, if
\(\bar k\) is an algebraic closure of \(k\), then
\[
  N_{\bar k}=\ker(u_{\bar k})
\]
is contained in the radical of \(G_{\bar k}\), that is, in its
largest smooth, connected, solvable, normal subgroup.

\smallskip
Then the map
\[
  B'\longmapsto u^{-1}(B')
\]
induces a bijection from the set of Borel subgroups of \(G'\) onto
the set of Borel subgroups of \(G\).}

Case \textup{b)} follows immediately from the correspondence between
algebraic subgroups of \(G'\) and algebraic subgroups of \(G\)
containing \(N\), together with the fact that, over an algebraically
closed field, the Borel subgroups of \(G\) contain its radical. The
latter fact follows at once from the argument preceding \SGARefLink{sga3:xiv:4.2}{4.2}.

We prove case \textup{a)}. By \SGARefLink{sga3:xii:7.12}{XII~7.12}, the map
\[
  H'\longmapsto H=u^{-1}(H')
\]
gives a bijection between the subpreschemes in groups \(H'\) of
\(G'\) that are smooth and of finite presentation over \(S\), have
connected fibers, and at every \(s\in S\) have the same reductive
rank and the same nilpotent rank as \(G'\), and the subpreschemes
in groups \(H\) of \(G\) having the analogous properties. Borel
subgroups of \(G'\) and of \(G\) have these properties.

It remains to show that, when \(H'\) and \(H\) correspond, \(H'\)
is a Borel subgroup of \(G'\) if and only if \(H\) is a Borel
subgroup of \(G\). By definition, this question reduces to the case
where \(S\) is the spectrum of an algebraically closed field. Since
\(N\) is central in \(G\), and hence in \(H\), the group \(H\) is
solvable if and only if \(H'\) is solvable. Finally, the
correspondence between algebraic subgroups of \(G'\) and algebraic
subgroups of \(G\) containing \(N\) shows immediately that \(H'\)
has the maximality property of \SGARefLink{sga3:xiv:4.1}{Definition~4.1} if and only if \(H\)
does. This completes the proof.

\paragraph{\textbf{Corollary 4.10.}}
\SGARefTarget{sga3:xiv:4.10}
{\itshape With the notation of \SGARefLink{sga3:xiv:4.9}{4.9},\footnote{The source prints
``\SGARefLink{sga3:xiv:4.9}{4.7}''; the intended reference is \SGARefLink{sga3:xiv:4.9}{4.9}.} if \(B'\) and \(B\) are
corresponding Borel subgroups of \(G'\) and \(G\), then
\[
  \mathfrak b=\operatorname{Lie}(u)^{-1}(\mathfrak b'),
\]
where \(\mathfrak g,\mathfrak g',\mathfrak b,\mathfrak b'\) are
the Lie algebras of \(G,G',B,B'\), respectively, and
\[
  \operatorname{Lie}(u)\colon\mathfrak g\longrightarrow\mathfrak g'
\]
is the homomorphism induced by \(u\).}

This follows immediately from the definitions and the identity
\[
  u^{-1}(B')=B.
\]

We can now prove the principal result of this section.

% END INCLUDED SGA3_English_current_buildable_source_20260823/source/inputs/XIV/tex/components/19_expose_XIV_proposition49_corollary410_en.tex

% BEGIN INCLUDED SGA3_English_current_buildable_source_20260823/source/inputs/XIV/tex/components/20_expose_XIV_theorem411_corollary412_en.tex
% Authority: Expo14.pdf, local pp. 21--22 / combined pp. 855--856.
% This closes Section 4.
% Exact next cursor: Section 5 and Theorem 5.1,
% local p. 22 / combined p. 856.

\paragraph{\textbf{Theorem 4.11.}}
\SGARefTarget{sga3:xiv:4.11}
{\itshape Let \(G\) be a smooth algebraic group over an
algebraically closed field \(k\), and let \(\mathfrak g\) be its
Lie algebra. Then \(\mathfrak g\) is the union of the Lie algebras
\(\mathfrak b\) of the Borel subgroups \(B\) of \(G\).}

We may plainly suppose that \(G\) is connected. Let \(R\) be the
radical of \(G\), and put \(G'=G/R\). By \SGARefLink{sga3:xiv:4.9}{4.9}\textup{(b)} and \SGARefLink{sga3:xiv:4.10}{4.10},
it is enough to prove the theorem for \(G'\); thus we may suppose
that \(G\) is semisimple.

Let \(Z\) be the center of \(G\), which is its reductive center, and
put \(G'=G/Z\). The same argument, now using \SGARefLink{sga3:xiv:4.9}{4.9}\textup{(a)},
reduces the theorem to \(G'\). We may therefore suppose that \(G\)
is semisimple and adjoint.

Let \(B\) be a Borel subgroup of \(G\), let \(T\) be a maximal torus
of \(B\), hence of \(G\), and let \(\mathfrak b\) and
\(\mathfrak t\) be their Lie algebras. By \SGARefLink{sga3:xiv:3.18}{3.18}, \(T\) is a subgroup
of type \((\mathrm C)\) of \(G\), that is, \(\mathfrak t\) is a
Cartan subalgebra of \(\mathfrak g\). The union of the conjugates
of \(\mathfrak t\) is consequently dense in \(\mathfrak g\)
\textup{(\SGARefLink{sga3:xiii:theorem:5:1}{XIII~5.1}, \textup{(i)}\(\Rightarrow\)\textup{(vii)})}.
\emph{A fortiori}, the union of the conjugates of \(\mathfrak b\)
is dense in \(\mathfrak g\).

Let \(X\) be the closed subscheme of
\[
  G/B\times W(\mathfrak g)
\]
whose \(k\)-valued points are the pairs \((gB,x)\) such that
\[
  x\in\operatorname{Ad}(g)\mathfrak b
\]
\textup{(compare XIII~1)}. The morphism
\[
  \psi\colon X\longrightarrow W(\mathfrak g)
\]
induced by the second projection is proper because \(G/B\) is
proper over \(k\). We have just seen that it is dominant; it is
therefore surjective. This proves \SGARefLink{sga3:xiv:4.11}{4.11}.

The only result of this section that we shall use in the remainder
of this exposé is the following corollary.

\paragraph{\textbf{Corollary 4.12.}}
\SGARefTarget{sga3:xiv:4.12}
{\itshape Let \(k\) be an infinite field, let \(G\) be a smooth
algebraic group over \(k\), and let \(T\) be a maximal torus of
\(G\). Write \(\mathfrak g\) and \(\mathfrak t\) for their Lie
algebras. Let
\[
  u\colon G\longrightarrow\operatorname{GL}(V)
\]
be a linear representation, where \(V\) is finite-dimensional over
\(k\), and let
\[
  u'\colon\mathfrak g\longrightarrow\mathfrak{gl}(V)
\]
be the induced representation of Lie algebras. Then the minimum of
the nullity of \(u'(x)\), for \(x\in\mathfrak g\), is attained at
an element \(x\in\mathfrak t\).}

We reduce immediately to the case where \(k\) is algebraically
closed. We may plainly suppose that \(G\) is connected, and, after
dividing \(G\) by \((\ker u)_{\mathrm{red}}\), that \(G\) is
affine. Using \SGARefLink{sga3:xiv:4.11}{4.11} and the conjugacy theorem for the maximal tori
of \(G\), we reduce to the further case in which \(G\) is solvable.

Then \(G\) is a semidirect product \(T\cdot U\), where \(U\) is the
unipotent part of \(G\), a smooth connected unipotent group
\textup{\emph{(Bible)}, 6, Theorem~3}. Thus, as a vector space,
\[
  \mathfrak g=\mathfrak t\oplus\mathfrak u,
  \qquad
  \mathfrak u=\operatorname{Lie}(U).
\]
By the Lie--Kolchin theorem
\textup{\emph{(Bible)}, 6, \SGARefLink{sga3:xiv:7:lemma1}{Theorem~1}}, the representation space
\(V\) has a composition series by stable subspaces \(V_i\), with
one-dimensional quotients
\[
  W_i=V_i/V_{i+1}.
\]
For every \(i\) there is therefore an induced representation
\[
  u_i\colon G\longrightarrow
  \operatorname{GL}(W_i)=\mathbb G_m
\]
and a corresponding homomorphism of Lie algebras
\[
  u_i'\colon\mathfrak g\longrightarrow k.
\]
For every \(x\in\mathfrak g\), the nullity of \(u'(x)\) is the
number of indices \(i\) for which \(u_i'(x)=0\).

Since \(U\) is unipotent, every \(u_i\) is trivial on \(U\), and
hence every \(u_i'\) vanishes on \(\mathfrak u\). If
\[
  x=t+v,\qquad t\in\mathfrak t,\quad v\in\mathfrak u,
\]
then \(u_i'(x)=u_i'(t)\) for every \(i\). Thus \(u'(x)\) and
\(u'(t)\) have the same nullity, and \SGARefLink{sga3:xiv:4.12}{4.12} follows.

% END INCLUDED SGA3_English_current_buildable_source_20260823/source/inputs/XIV/tex/components/20_expose_XIV_theorem411_corollary412_en.tex

% BEGIN INCLUDED SGA3_English_current_buildable_source_20260823/source/inputs/XIV/tex/components/21_expose_XIV_section5_theorem51_through_lemma53_en.tex
% Authority: Expo14.pdf, local pp. 22--24 / combined pp. 856--858.
% Exact next cursor: Corollary 5.4, local p. 24 / combined p. 858.

\section{Relations Between Cartan Subgroups and Cartan Subalgebras}
\SGARefTarget{sga3:xiv:5}

Applying \SGARefLink{sga3:xiv:4.12}{4.12} to the adjoint representation of \(G\), we obtain:

\paragraph{\textbf{Theorem 5.1.}}
\SGARefTarget{sga3:xiv:5.1}
{\itshape Let \(G\) be a smooth algebraic group over an infinite
field, let \(T\) be a maximal torus of \(G\), and let
\(\mathfrak g\supset\mathfrak t\) be their Lie algebras. Then
\(\mathfrak t\) contains a regular element of \(\mathfrak g\).}

\paragraph{\textbf{Corollary 5.2.}}
\SGARefTarget{sga3:xiv:5.2}
{\itshape Let \(G\) be a smooth \(S\)-prescheme in groups of finite
presentation, and let \(\mathfrak g\) be its Lie algebra.

\smallskip
\noindent
\textup{a)} Conditions \((\mathrm C_0)\) through \((\mathrm C_2)\)
of \SGARefLink{sga3:xiv:2.9}{2.9} for \(\mathfrak g\) are equivalent. In particular, if the
infinitesimal rank of the fibers of \(\mathfrak g\) at the points
of \(S\) is locally constant, then, étale locally,
\(\mathfrak g\) admits a Cartan subalgebra and hence, by
\SGARefLink{sga3:xiv:3.9}{3.9}\textup{(a)}, \(G\) admits a subgroup of type \((\mathrm C)\).

\smallskip
\noindent
\textup{b)} Let \(H\) be a subprescheme in groups of \(G\), smooth
and of finite presentation over \(S\), whose fibers are connected
and at every \(s\in S\) have the same reductive rank as \(G\);
for example, \(H\) may be a maximal torus or a Cartan subgroup of
\(G\). Let \(D\) be a subgroup of type \((\mathrm C)\) of \(G\).
Then
\[
  H\subset D\quad\Longleftrightarrow\quad
  \mathfrak h\subset\mathfrak d.
\]

\smallskip
\noindent
\textup{c)} Suppose that condition \((\mathrm C_0)\) holds, that is,
the infinitesimal rank of \(G\) is locally constant. Let \(H\) be a
subprescheme in groups of \(G\), smooth and of finite presentation
over \(S\), with connected nilpotent fibers, and having at every
point the same reductive rank as \(G\); for example, \(H\) may be a
maximal torus or a Cartan subgroup of \(G\). Then \(H\) is contained
in a unique subgroup \(D\) of type \((\mathrm C)\) of \(G\).

\smallskip
\noindent
\textup{d)} If \(G\) admits a Cartan subgroup fpqc locally, then so
does every subgroup \(D\) of type \((\mathrm C)\) of \(G\).}

\smallskip
\noindent
\emph{Proof.}
\textup{a)}
Assume \((\mathrm C_0)\), and let us prove \((\mathrm C_2)\), that
is, that every quasi-regular section \(a\) of \(\mathfrak g\) is
regular. As usual, we reduce first to \(S\) affine noetherian and
then, since the question is infinitesimal by \SGARefLink{sga3:xiv:2.15}{2.15}\textup{(b)}, to
\(S\) local Artinian. In this last case \((\mathrm C_0)\) is
automatic. We may moreover suppose that the residue field \(k\)
of \(S\) is infinite.

By \SGARefLink{sga3:xiv:3.7}{3.7}, it is enough to show that \(\mathfrak g\) admits a Cartan
subalgebra. Let \(T\) be a maximal torus of \(G\). By \SGARefLink{sga3:xiv:5.1}{5.1} there is
a quasi-regular element \(t\) in the Lie algebra \(\mathfrak t\)
of \(T\). We shall show that \(t\) is regular.

Consider the linear representation of \(T\) on \(\mathfrak g\)
induced by the adjoint representation of \(G\). There is a finite
set of characters \((u_i)_{i\in I}\) of \(T\) and a direct-sum
decomposition
\[
  \mathfrak g=\bigoplus_{i\in I}\mathfrak g_i
\]
into \(T\)-stable submodules such that \(T\) acts on
\(\mathfrak g_i\) through \(u_i\)
\textup{(I, \S4.7.3)}. Let
\[
  u_i'\colon\mathfrak t\longrightarrow A
\]
be the homomorphism of Lie algebras induced by
\(u_i\colon T\to\mathbb G_m\), where \(A\) is the ring of \(S\).
Let \(u_{i0}\) and \(u_{i0}'\) be their reductions under
\(A\to k\). Put
\[
  I'=\{i\in I\mid u_{i0}'\ne0\},
  \qquad I''=I\setminus I'.
\]
The regularity of \(t_0\) is expressed by
\[
  u_{i0}'(t_0)\ne0\qquad(i\in I'),
\]
and therefore \(u_i'(t)\) is invertible for \(i\in I'\).

The Cartan subalgebra of \(\mathfrak g_0\) defined by \(t_0\), that
is, the kernel of the semisimple endomorphism
\(\operatorname{ad}(t_0)\), is
\[
  \bigoplus_{i\in I''}(\mathfrak g_i)_0.
\]
Set
\[
  \mathfrak d=\bigoplus_{i\in I''}\mathfrak g_i.
\]
Then \(\operatorname{ad}(t)\) is nilpotent on \(\mathfrak d\), and
is an automorphism of
\[
  \mathfrak g/\mathfrak d
  \simeq\bigoplus_{i\in I'}\mathfrak g_i.
\]
By \SGARefLink{sga3:xiv:2.6}{2.6}, \(t\) is regular.

\smallskip
\noindent
\textup{b)}
This follows from \SGARefLink{sga3:xiv:3.12}{3.12}\textup{(a)}. By \SGARefLink{sga3:xiv:5.1}{5.1}, \(H\) satisfies the
required hypothesis that every geometric fiber of \(\mathfrak h\)
contains a regular element of the corresponding fiber of
\(\mathfrak g\).

\smallskip
\noindent
\textup{c)}
By \textup{b)}, it is enough to show that \(\mathfrak h\) is
contained in a unique Cartan subalgebra of \(\mathfrak g\). By
\textup{a)}, \(\mathfrak g\) satisfies \((\mathrm C_2)\). We are
therefore reduced to the following lemma.

\paragraph{\textbf{Lemma 5.3.}}
\SGARefTarget{sga3:xiv:5.3}
{\itshape Let \(\mathfrak g\) be a Lie algebra over \(S\), locally
free of finite type as a module, and satisfying condition
\((\mathrm C_2)\) of \SGARefLink{sga3:xiv:2.9}{2.9}. Let \(\mathfrak h\) be a Lie subalgebra
of \(\mathfrak g\) such that \(\mathfrak h\) is locally a direct
summand, is locally nilpotent, and, for every \(s\in S\), the
geometric fiber \(\mathfrak h_{\bar s}\) contains a regular element
of \(\mathfrak g_{\bar s}\). Then \(\mathfrak h\) is contained in
a unique Cartan subalgebra of \(\mathfrak g\).}

In the situation at hand, \(\mathfrak h\) has these properties: it
is locally nilpotent because \(H\) has nilpotent fibers, and the
condition on regular elements follows from \SGARefLink{sga3:xiv:5.1}{5.1}.

Since \SGARefLink{sga3:xiv:5.3}{5.3} is fpqc local, it is enough, at a point \(s\in S\) with
\(\kappa(s)\) infinite, to find an open neighborhood \(U\) of \(s\)
on which existence and uniqueness hold after every base change
\(S'\to S\) factoring through \(U\). Choose a regular element of
\(\mathfrak g\otimes\kappa(s)\) lying in
\(\mathfrak h\otimes\kappa(s)\), and extend it to a section \(a\)
of \(\mathfrak h\) over an open neighborhood \(U\) of \(s\).
Condition \((\mathrm C_2)\) and \SGARefLink{sga3:xiv:2.10}{2.10} allow us to shrink \(U\) so
that \(a\) is regular. We may thus suppose \(U=S\).

Every Cartan subalgebra of \(\mathfrak g\) containing
\(\mathfrak h\) contains \(a\), and hence equals
\[
  \mathfrak d=\Nil(a,\mathfrak g)
\]
by \SGARefLink{sga3:xiv:2.6}{2.6}; this proves uniqueness. Conversely,
\(\mathfrak h\subset\mathfrak d\) because \(\mathfrak h\) is
locally nilpotent, which proves existence.

\smallskip
\noindent
\textup{d)}
This is a special case of \SGARefLink{sga3:xii:7.9}{XII~7.9}\textup{\SGARefLink{sga3:xii:7.9:item:d:02}{(d)}}.

% END INCLUDED SGA3_English_current_buildable_source_20260823/source/inputs/XIV/tex/components/21_expose_XIV_section5_theorem51_through_lemma53_en.tex

% BEGIN INCLUDED SGA3_English_current_buildable_source_20260823/source/inputs/XIV/tex/components/22_expose_XIV_corollary54_through_corollary56_en.tex
% Authority: Expo14.pdf, local pp. 24--25 / combined pp. 858--859.
% This closes Section 5.
% Exact next cursor: Section 6 and Theorem 6.1,
% local p. 25 / combined p. 859.

\paragraph{\textbf{Corollary 5.4.}}
\SGARefTarget{sga3:xiv:5.4}
{\itshape Let \(G\) be a smooth connected algebraic group over an
algebraically closed field \(k\). Let \(H\) be a connected
algebraic subgroup such that \(\mathfrak h\) contains a Cartan
subalgebra of \(\mathfrak g\), or equivalently such that \(H\)
contains a subgroup of type \((\mathrm C)\) of \(G\). Then the
number of conjugates of \(H\) containing a regular element of
\(G(k)\) equals the number of conjugates of \(\mathfrak h\)
containing a regular element of \(\mathfrak g\).}

Indeed, let \(g\) be a regular element of \(G(k)\), and let \(C\)
be the unique Cartan subgroup of \(G\) containing \(g\)
\textup{(XIII~2)}. The conjugates of \(H\) that contain \(g\) are
exactly those that contain \(C\)
\textup{(\SGARefLink{sga3:xiii:corollary:2:8}{XIII~2.8}\textup{(b)})}.

Similarly, let \(x\) be a regular element of \(\mathfrak g\). If a
conjugate \(\mathfrak h'\) of \(\mathfrak h\) contains \(x\), then
\[
  \operatorname{ad}(x)_{\mathfrak g/\mathfrak h'}
\]
is injective by \SGARefLink{sga3:xiii:corollary:5:4}{XIII~5.4}, and therefore
\[
  \mathfrak h'\supset\Nil(x,\mathfrak g)=\mathfrak d.
\]
Thus the number of conjugates of \(\mathfrak h\) containing \(x\)
equals the number containing the Cartan subalgebra
\(\mathfrak d\).

The subalgebra \(\mathfrak d\) is the Lie algebra of a subgroup
\(D\) of type \((\mathrm C)\) of \(G\). We may suppose that
\(C\subset D\). It remains to show that
\[
  H\supset C\quad\Longleftrightarrow\quad
  \mathfrak h\supset\mathfrak d.
\]
By \SGARefLink{sga3:xiii:corollary:5:5}{XIII~5.5}, the inclusion
\(\mathfrak h\supset\mathfrak d\) implies \(H\supset D\), and
\emph{a fortiori} \(H\supset C\). Conversely, by \SGARefLink{sga3:xiv:5.1}{5.1} the Lie
algebra \(\mathfrak t\) of \(C\) contains a regular element \(x\)
of \(\mathfrak g\). Hence \(H\supset C\) implies
\(x\in\mathfrak h\), and, as observed above, this implies
\(\mathfrak h\supset\mathfrak d\). This completes the proof.

Let \(G\) be as in \SGARefLink{sga3:xiv:5.2}{5.2}, and suppose that the infinitesimal rank of
its fibers is locally constant, that is, condition
\((\mathrm C_0)\) holds. By \SGARefLink{sga3:xiv:5.2}{5.2}\textup{(c)}, there is a morphism of
functors on \((\operatorname{Sch}/S)^\circ\),
\[
  \mathcal C\longrightarrow\mathcal D,
\]
where
\[
  \begin{aligned}
  \mathcal C(S')&=
    \{\text{Cartan subgroups of }G_{S'}\},\\
  \mathcal D(S')&=
    \{\text{subgroups of type \((\mathrm C)\) of }G_{S'}\}.
  \end{aligned}
\]
By \SGARefLink{sga3:xiv:3.10}{3.10} and \SGARefLink{sga3:xiv:5.2}{5.2}\textup{(a)}, the functor \(\mathcal D\) is
represented by a smooth quasi-projective prescheme over \(S\).
Let \(D\) be the universal subgroup of type \((\mathrm C)\) of
\(G_{\mathcal D}\). One may then define the functor
\[
  \mathcal C_D\colon
  (\operatorname{Sch}/\mathcal D)^\circ
  \longrightarrow\operatorname{Ens}
\]
for the \(\mathcal D\)-group \(D\), exactly as \(\mathcal C\) was
defined for the \(S\)-group \(G\).

\paragraph{\textbf{Proposition 5.5.}}
\SGARefTarget{sga3:xiv:5.5}
{\itshape Under the preceding hypotheses, regard \(\mathcal C\) as
a functor over \(\mathcal D\). Then \(\mathcal C\) is
\(\mathcal D\)-isomorphic to the functor \(\mathcal C_D\) of
Cartan subgroups of \(D\).}

This follows immediately from the next corollary.

\paragraph{\textbf{Corollary 5.6.}}
\SGARefTarget{sga3:xiv:5.6}
{\itshape Let \(G\) be a smooth \(S\)-prescheme in groups of finite
presentation, and let \(D\) be a subgroup of type \((\mathrm C)\)
of \(G\). There is a bijective correspondence between the Cartan
subgroups of \(G\) contained in \(D\) and the Cartan subgroups of
\(D\). More precisely, for a subgroup \(H\subset D\), \(H\) is a
Cartan subgroup of \(G\) if and only if it is a Cartan subgroup of
\(D\).}

This is a special case of \SGARefLink{sga3:xii:7.9}{XII~7.9}\textup{\SGARefLink{sga3:xii:7.9:item:c:02}{(c)}}, taking into account
that, over an algebraically closed field, every subgroup of type
\((\mathrm C)\) of \(G\) contains a Cartan subgroup of \(G\).

For the next section, the principal result obtained here is
\SGARefLink{sga3:xiv:5.2}{5.2}\textup{(c)} with \(H\) a Cartan subgroup of \(G\). It permits
the formulation of \SGARefLink{sga3:xiv:5.5}{5.5} and thus gives a useful dévissage of
\(\mathcal C\).

% END INCLUDED SGA3_English_current_buildable_source_20260823/source/inputs/XIV/tex/components/22_expose_XIV_corollary54_through_corollary56_en.tex

% BEGIN INCLUDED SGA3_English_current_buildable_source_20260823/source/inputs/XIV/tex/components/23_expose_XIV_section6_theorem61_en.tex
% Authority: Expo14.pdf, local pp. 25--27 / combined pp. 859--861.
% This begins Section 6 and contains Theorem 6.1.
% Exact next cursor: Corollary 6.2, local p. 27 / combined p. 861.

\section{Applications to the Structure of Algebraic Groups}
\SGARefTarget{sga3:xiv:6}

\paragraph{\textbf{Theorem 6.1.}}
\SGARefTarget{sga3:xiv:6.1}
{\itshape Let \(G\) be a smooth algebraic group over a field
\(k\). Consider the scheme \(\mathcal T\) of maximal tori of
\(G\), isomorphic to the scheme \(\mathcal C\) of Cartan
subgroups of \(G\). This is a homogeneous space under \(G^0\) and
a smooth affine connected algebraic scheme
\textup{(\SGARefLink{sga3:xii:7.1}{XII~7.1}\textup{\SGARefLink{sga3:xii:7.1:item:d}{(d)}})}. Then \(\mathcal C\) is a rational
variety: its field of rational functions is a purely
transcendental extension of \(k\).}

We first prove the theorem when \(k\) is infinite. We may plainly
suppose that \(G\) is connected, since \(\mathcal T\), and hence
\(\mathcal C\), is unchanged when \(G\) is replaced by \(G^0\).
Moreover, by \SGARefLink{sga3:xii:7.6}{XII~7.6}, \(\mathcal C\) is unchanged when \(G\) is
divided by a central subgroup. We may therefore first divide by
the center of \(G\) and suppose that \(G\) is affine
\textup{(\SGARefLink{sga3:xii:6.1}{XII~6.1})}, and then divide by its reductive center
\textup{(\SGARefLink{sga3:xii:4.1}{XII~4.1} and \SGARefLink{sga3:xiv:4.4}{4.4})} and suppose that the reductive center
of \(G\) is trivial \textup{(\SGARefLink{sga3:xii:4.7}{XII~4.7}\textup{\SGARefLink{sga3:xii:4.7:item:b}{(b)}})}.

We proceed by induction on \(n=\dim G\), assuming the result for
dimensions \(n'<n\). If \(G\) is nilpotent, then \(\mathcal C\)
consists of a single \(k\)-rational point and the assertion is
trivial. Otherwise, the Lie algebra of \(G\) is not nilpotent
\textup{(\SGARefLink{sga3:xiv:1.3}{1.3})}. Its Cartan subalgebras, and therefore the
subgroups of type \((\mathrm C)\) of \(G\), have some dimension
\(n'<n\). Consider the morphism
\[
  \mathcal C\longrightarrow\mathcal D
\]
of \SGARefLink{sga3:xiv:5.5}{5.5}. By \SGARefLink{sga3:xiv:3.10}{3.10}, since the infinite-field Lie algebra
\(\mathfrak g\) contains a regular element, \(\mathcal D\) is a
rational variety. Thus its function field \(K\) is purely
transcendental over \(k\).

Let \(x\) be the generic point of \(\mathcal D\). By \SGARefLink{sga3:xiv:5.5}{5.5}, the
fiber of \(\mathcal C\) at \(x\) is the scheme of Cartan
subgroups of a smooth connected algebraic group \(D_x\) over
\(K=\kappa(x)\), namely the generic subgroup of type
\((\mathrm C)\) of \(G\). The function field \(L\) of
\(\mathcal C\) is therefore the function field of
\(\mathcal C_{D_x}\). Since \(\dim D_x=n'<n\), the induction
hypothesis makes \(L\) purely transcendental over \(K\), and
hence, by transitivity, purely transcendental over \(k\).

Suppose now that \(k\) is finite. A different argument is
required. We may again suppose that \(G\) is affine and
connected. Since \(k\) is perfect, the radical of
\(G_{\bar k}\) descends to a subgroup \(R\) of \(G\).
Suppose first that \(R\ne G\), so that \(G\) is not solvable, and
let
\[
  u\colon G\longrightarrow G'=G/R
\]
be the canonical morphism. By \SGARefLink{sga3:xii:7.1}{XII~7.1}\textup{\SGARefLink{sga3:xii:7.1:item:e}{(e)}}, it induces
\[
  v\colon\mathcal C_G\longrightarrow\mathcal C_{G'},
  \qquad C\longmapsto u(C).
\]
Let \(x\) be the generic point of \(\mathcal C_{G'}\), and put
\(K=\kappa(x)\). The fiber \(v^{-1}(x)\) is the scheme of Cartan
subgroups of \(G_K\) whose image in \(G'_K\) is the generic
Cartan subgroup \(C'_x\) of \(G'\). By \SGARefLink{sga3:xii:7.9}{XII~7.9}\textup{\SGARefLink{sga3:xii:7.9:item:c:02}{(c)}}, it is
also the scheme of Cartan subgroups of
\[
  H=u_K^{-1}(C'_x).
\]
The field \(K\) is infinite. The infinite-field case just proved
therefore shows that the function field \(L\) of
\(\mathcal C_G\), which is also that of \(\mathcal C_H\), is
purely transcendental over \(K\).\footnote{The French source
refers here to the ``already proved part of \SGARefLink{sga3:xiv:6.1}{7.1}''; the argument
is the infinite-field part of \SGARefLink{sga3:xiv:6.1}{Theorem 6.1}.}

It remains to prove that \(K\) is purely transcendental over
\(k\), and we are thus reduced to the case in which \(G\) is
semisimple. Dividing by its reductive center, we may further
suppose that \(G\) is adjoint. Then \SGARefLink{sga3:xiv:3.18}{3.18} gives
\(\mathcal C_G\simeq\mathcal D_G\), and by \SGARefLink{sga3:xiv:3.10}{3.10} it is enough to
show that \(\mathfrak g\) has a regular point. As noted in \SGARefLink{sga3:xiv:3.20}{3.20},
this is an unpublished result of Chevalley.\footnote{See the
Appendix by J.-P. Serre.}

It remains only to treat the case in which \(k\) is finite and
\(G\) is connected, affine, and solvable. In fact, the next
result holds over an arbitrary field.

% END INCLUDED SGA3_English_current_buildable_source_20260823/source/inputs/XIV/tex/components/23_expose_XIV_section6_theorem61_en.tex

% BEGIN INCLUDED SGA3_English_current_buildable_source_20260823/source/inputs/XIV/tex/components/24_expose_XIV_corollary62_lemma63_en.tex
% Authority: Expo14.pdf, local pp. 27--28 / combined pp. 861--862.
% This completes the proof of Theorem 6.1.
% Exact next cursor: Corollary 6.4, local p. 28 / combined p. 862.

\paragraph{\textbf{Corollary 6.2.}}
\SGARefTarget{sga3:xiv:6.2}
{\itshape Let \(G\) be a smooth solvable algebraic group over a
field \(k\). Then the variety \(\mathcal C\) of Cartan subgroups
of \(G\) is isomorphic to an affine space
\(\operatorname{Spec}k[t_1,\ldots,t_N]\).}

We may suppose that \(G\) is connected and affine. Let
\(G_\infty\) be the terminal group in the descending central
series
\[
  G_0=G,\qquad G_{i+1}=[G,G_i].
\]
Thus \(G_\infty\) is the smallest invariant algebraic subgroup
of \(G\) for which \(G/G_\infty\) is nilpotent. Let \(C\) be a
Cartan subgroup of \(G\), which exists by \SGARefLink{sga3:xiv:1.1}{1.1}. Its image in
\(G/G_\infty\) is a Cartan subgroup and hence is all of
\(G/G_\infty\). Consequently the morphism
\[
  G_\infty\longrightarrow\mathcal C,\qquad
  g\longmapsto\operatorname{Ad}(g)C
\]
is an epimorphism and identifies
\[
  \mathcal C\simeq G_\infty/(N\cap G_\infty),
\]
where \(N\) is the normalizer of \(C\) in \(G\). In fact \(N=C\),
as recalled in \SGARefLink{sga3:xiv:4.4}{4.4}, but that is immaterial here.

The group \(U=G_\infty\) is smooth, connected, and unipotent:
over \(\bar k\) it lies in the unipotent part of \(G\), by the
known structure of smooth affine solvable groups
\textup{\emph{Bible}~6, Theorem~3}. If \(k\) is perfect, as in
the finite-field case needed for \SGARefLink{sga3:xiv:6.1}{6.1}, it follows easily that
\(U\) is \(k\)-unipotent: it has a composition series by
algebraic subgroups \(U_i\) with
\(U_i/U_{i+1}\simeq\mathbb G_a\).

Rosenlicht proved that a group of the form \(G_\infty\) above is
\(k\)-unipotent without any restriction on \(k\)
\textup{(M.~Rosenlicht, \emph{Questions of rationality for
solvable algebraic groups over non perfect fields}, Annali di
Matematica, 1963, pp.~97--120, \SGARefLink{sga3:xiv:7:lemma4}{Theorem~4}, \SGARefLink{sga3:xiv:7:lemma2}{Corollary~2})}. We admit
this substantially more delicate result and apply the following
lemma.

\paragraph{\textbf{Lemma 6.3.}}
\SGARefTarget{sga3:xiv:6.3}
{\itshape Let \(U\) be a smooth connected algebraic group over a
field \(k\), and let \(X=U/V\) be a homogeneous space under \(U\)
having a \(k\)-rational point. If \(U\) is \(k\)-unipotent, then,
as a \(k\)-scheme, \(X\) is isomorphic to an affine space
\(\operatorname{Spec}k[t_1,\ldots,t_N]\).}

Choose a composition series
\[
  U=U_0\supset U_1\supset\cdots\supset U_n=(e)
\]
by smooth connected normal subgroups of \(U\), with
\(U_i/U_{i+1}\simeq\mathbb G_a\). The groups
\[
  K_i=U_iV
\]
are algebraic subgroups of \(U\). They need not be smooth or
connected when \(V\) is not, but this does not matter, and
\(K_{i+1}\) is normal in \(K_i\). The canonical epimorphism
\[
  U_i/U_{i+1}\longrightarrow K_i/K_{i+1}
\]
shows that \(K_i/K_{i+1}\) is either the unit group or isomorphic
to \(\mathbb G_a/H_i\), where \(H_i\) is a finite subgroup of
\(\mathbb G_a\). Rosenlicht's \SGARefLink{sga3:xiv:7:lemma2}{Theorem~2} in the cited paper gives
\[
  K_i/K_{i+1}\simeq\mathbb G_a;
\]
this last assertion is immediate when \(k\) is perfect.

Put \(X_i=U/K_i\). We prove inductively that every \(X_i\) is an
affine space. If \(K_i/K_{i+1}\) is trivial, then
\(X_i=X_{i+1}\). Otherwise \(X_{i+1}\) is a principal
\(\mathbb G_a\)-bundle over \(X_i\). Since \(X_i\) is affine,
the bundle is trivial, so
\[
  X_{i+1}\simeq X_i\times\mathbb G_a.
\]
The induction proves the lemma. It proves \SGARefLink{sga3:xiv:6.2}{6.2} and thereby
completes the proof of \SGARefLink{sga3:xiv:6.1}{6.1}.

% END INCLUDED SGA3_English_current_buildable_source_20260823/source/inputs/XIV/tex/components/24_expose_XIV_corollary62_lemma63_en.tex

% BEGIN INCLUDED SGA3_English_current_buildable_source_20260823/source/inputs/XIV/tex/components/25_expose_XIV_corollary64_through_corollary610_en.tex
% Authority: Expo14.pdf, local pp. 28--30 / combined pp. 862--864.
% Exact next cursor: Remarks 6.11, local p. 30 / combined p. 864.

\paragraph{\textbf{Corollary 6.4.}}
\SGARefTarget{sga3:xiv:6.4}
{\itshape Let \(G\) be a smooth algebraic group over an infinite
field \(k\). Then the set of \(k\)-rational points of
\(\mathcal C\), in the notation of \SGARefLink{sga3:xiv:6.1}{6.1}, is Zariski dense. The
union of the Cartan subgroups of \(G\) is dense in \(G\).}

The first assertion holds for every rational variety over an
infinite field and is, for us, the most important arithmetic
consequence of the rationality result. The second follows from
the first and the density theorem \SGARefLink{sga3:xiii:theorem:2:1}{XIII~2.1}.

\paragraph{\textbf{Corollary 6.5.}}
\SGARefTarget{sga3:xiv:6.5}
{\itshape Let \(G\) be a smooth connected algebraic group over
\(k\). Then the variety \(Z\) of regular semisimple points of
\(G\) \textup{(\SGARefLink{sga3:xiii:corollary:3:5}{XIII~3.5})} is unirational. In particular, if
\(k\) is infinite, its \(k\)-rational points are dense in \(Z\).}

Indeed, \(Z\) is an open subscheme of a torus over
\(\mathcal C\). Its function field \(L\) is therefore the
function field of a torus over the function field \(K\) of
\(\mathcal C\), namely the generic maximal torus of \(G\). Thus
\(L\) is unirational over \(K\) by \SGARefLink{sga3:xiii:lemma:3:4}{XIII~3.4}. Since \SGARefLink{sga3:xiv:6.1}{6.1} makes
\(K\) purely transcendental over \(k\), \(L\) is unirational
over \(k\).

\paragraph{\textbf{Corollary 6.6.}}
\SGARefTarget{sga3:xiv:6.6}
{\itshape Let \(G\) be a smooth connected algebraic group over
\(k\), and let \(H\) be the subgroup generated by the subvariety
\(Z\) of regular semisimple points. Equivalently
\textup{(\SGARefLink{sga3:xii:8.2}{XII~8.2})}, \(H\) is the smallest invariant algebraic
subgroup of \(G\) for which \(G/H\) has reductive rank zero; over
\(\bar k\), \(G/H\) is an extension of an abelian variety by a
smooth connected unipotent group. Then \(H\) is unirational.

In particular, suppose \(G=H\); equivalently
\textup{(\SGARefLink{sga3:xii:8.4}{XII~8.4})}, suppose \(G\) is affine and there is no
nontrivial homomorphism \(G_{\bar k}\to\mathbb G_a\). Then \(G\)
is unirational. If, in addition, \(k\) is infinite, \(G(k)\) is
dense in \(G\).}

This follows immediately from \SGARefLink{sga3:xiv:6.5}{6.5}. Indeed, if \(Z_i\) are smooth
connected unirational \(k\)-preschemes and
\(u_i\colon Z_i\to G\) are morphisms, then the algebraic subgroup
of \(G\) generated by the \(u_i\) is unirational
\textup{(VI\textsubscript{B}~\SGARefLink{prop:sga3-VIB-7.1}{7.1})}.

An interesting special case of either \SGARefLink{sga3:xiv:6.5}{6.5} or \SGARefLink{sga3:xiv:6.6}{6.6} is the
following.

\paragraph{\textbf{Corollary 6.7.}}
\SGARefTarget{sga3:xiv:6.7}
{\itshape Let \(G\) be a smooth connected affine algebraic group
of unipotent rank zero. Then \(G\) is unirational.}

The conclusion of \SGARefLink{sga3:xiv:6.4}{6.4} may be sharpened as follows.

\paragraph{\textbf{Corollary 6.8.}}
\SGARefTarget{sga3:xiv:6.8}
{\itshape Let \(G\) be a smooth connected algebraic group over an
infinite field \(k\), and let \(C\) be a Cartan subgroup of \(G\).
Then the union of the conjugates of \(C\) by regular semisimple
elements of \(G(k)\) is dense in \(G\).}

This follows at once from \SGARefLink{sga3:xiv:6.4}{6.4} and \SGARefLink{sga3:xiii:corollary:3:6}{XIII~3.6}, which says that the
morphism
\[
  \varphi\colon Z\times C\longrightarrow G,\qquad
  \varphi(t,c)=\operatorname{Ad}(t)c,
\]
is dominant. The same result also implies the next corollary,
without any assumption that \(k\) is infinite.

\paragraph{\textbf{Corollary 6.9.}}
\SGARefTarget{sga3:xiv:6.9}
{\itshape Let \(G\) be a smooth connected algebraic group over
\(k\), and let \(H\) be a smooth connected algebraic subgroup
having the same reductive rank and the same nilpotent rank as
\(G\). Equivalently, over \(\bar k\), the group \(H_{\bar k}\)
contains a Cartan subgroup of \(G_{\bar k}\). If \(H\) is
unirational, so is \(G\). If \(H(k)\) is dense in \(H\), then
\(G(k)\) is dense in \(G\).}

The morphism
\[
  \varphi\colon Z\times H\longrightarrow G,\qquad
  \varphi(t,h)=\operatorname{Ad}(t)h,
\]
is dominant. By \SGARefLink{sga3:xiv:6.5}{6.5}, \(Z\) is unirational, and by hypothesis so
is \(H\), hence so are \(Z\times H\) and \(G\). The density
assertion is proved in the same way.

We recover the following well-known result, due to Chevalley in
characteristic zero and to Rosenlicht in characteristic
\(p>0\).

\paragraph{\textbf{Corollary 6.10.}}
\SGARefTarget{sga3:xiv:6.10}
{\itshape Let \(G\) be a smooth connected affine algebraic group
over a perfect field \(k\). Then \(G\) is unirational. If \(k\)
is also infinite, \(G(k)\) is dense in \(G\).}

By \SGARefLink{sga3:xiv:1.1}{1.1}, \(G\) has a Cartan subgroup \(C\). By \SGARefLink{sga3:xiv:6.9}{6.9}, it is enough
to prove that \(C\) is unirational. Since \(k\) is perfect,
Galois descent from the algebraically closed case
\textup{\emph{Bible}~6, \SGARefLink{sga3:xiv:7:lemma2}{Theorem~2}} gives
\[
  C=T\times C_u,
\]
where \(T\) is the maximal torus of \(C\) and \(C_u\) is smooth,
connected, and unipotent. The torus \(T\) is unirational by
XIII~3.4. Moreover, \(C_u\) is \(k\)-unipotent because \(k\) is
perfect, so \SGARefLink{sga3:xiv:6.3}{6.3} applies to \(C_u\).

% END INCLUDED SGA3_English_current_buildable_source_20260823/source/inputs/XIV/tex/components/25_expose_XIV_corollary64_through_corollary610_en.tex

% BEGIN INCLUDED SGA3_English_current_buildable_source_20260823/source/inputs/XIV/tex/components/26_expose_XIV_remarks611_en.tex
% Authority: Expo14.pdf, local pp. 30--31 / combined pp. 864--865.
% This closes Section 6, including the later editorial proof in note 8.
% Exact next cursor: Appendix, local p. 32 / combined p. 866.

\paragraph{\textbf{Remarks 6.11.}}
\SGARefTarget{sga3:xiv:6.11}
\textup{(a)} Rosenlicht gave examples of twisted forms of
\(\mathbb G_a\) over a nonperfect field having only finitely many
rational points; \emph{a fortiori}, they are not unirational.
Chevalley, on the other hand, gave an example of a torus over a
field of characteristic zero that is not rational. By contrast,
Chevalley's theory of semisimple groups shows that, over an
algebraically closed field, every smooth connected affine
algebraic group is rational. In any case, the question of
unirationality arises only for affine algebraic groups: by
Chevalley's structure theorem, a unirational algebraic group is
necessarily affine.

\textup{(b)} With the notation of \SGARefLink{sga3:xiv:6.6}{6.6}, it is tempting to seek a
criterion for the unirationality of \(G\) in terms of \(G/H\),
which is unipotent when \(G\) is affine. Clearly \(G/H\) must be
unirational; is that condition also sufficient? An example of
Rosenlicht shows that a smooth connected unipotent algebraic
group \(U\) may be rational without being \(k\)-unipotent.

\textup{(c)} Over a finite field \(k\), it would be interesting
to study questions of ``density'' such as the following one,
raised by Rosenlicht: if \(G\) is a smooth connected algebraic
group over \(k\), is \(G\) generated by its Cartan
subgroups?\footnote{This question was subsequently answered in
the affirmative by Steinberg.}

\begin{quote}\small
\textbf{Editors' note.}
The editors did not
identify this result in Steinberg's \emph{Collected Papers}; here
is a proof based on the Bruhat decomposition.

Let \(G/k\) be semisimple over the finite field \(k\), and let
\(G^\sharp\) be the subgroup generated by its \(k\)-subtori. The
question is stable under central isogeny and restriction of
scalars, so, as in \SGARefLink{sga3:xiv:7:lemma1}{Lemma~1} of the Appendix below, it is enough to
treat the geometrically simple case. The group \(G\) is
quasi-split by Lang's theorem and admits a Killing pair
\((T,B)\). By definition \(T\subset G^\sharp\), hence \(T\)
normalizes \(G^\sharp\). The big cell
\[
  R_u(B^-)\,T\,R_u(B)
\]
shows that it suffices to prove \(R_u(B)\subset G^\sharp\).
Let \(S\) be the maximal split \(k\)-torus of \(T\). For the
relative root system \(\Phi(G,S)\), choose a basis
\(\Delta_k\), and denote by \(U(\alpha)\) the unipotent subgroup
associated with \(\alpha\)
\textup{(A.~Borel, \emph{Linear Algebraic Groups}, second
edition, 1991, Proposition~21.9)}. Since \(R_u(B)\) is generated
by the \(U(\alpha)\), \(\alpha\in\Delta_k\), it is enough to
show \(U(\alpha)\subset G^\sharp\).

The classification supplies a quasi-split semisimple subgroup
\(G_\alpha\) of type \(A_1\), \({}^{2}A_1\), \({}^{3}A_1\), or
\({}^{2}A_2\), with
\[
  U(\alpha)\subset G_\alpha\subset G.
\]
Thus it is enough to consider \(G=\operatorname{PGL}_2\) or
\(G=\operatorname{SU}_3(K)\), where \(K/k\) is the unique
quadratic extension.

For \(\operatorname{PGL}_2\), since \(G^\sharp\) contains the
standard split torus \(T\), the possibilities up to
\(G(k)\)-conjugacy are
\[
  G^\sharp=T,\qquad G^\sharp=B,\qquad
  G^\sharp=\operatorname{PGL}_2.
\]
The first is excluded because \(G^\sharp\) contains the
\(k\)-torus \(R_{K/k}(\mathbb G_m)/\mathbb G_m\). The preceding
discussion also shows that \(G^\sharp=B\) forces
\(G^\sharp=\operatorname{PGL}_2\). Hence
\(G^\sharp=\operatorname{PGL}_2\).

For \(G=\operatorname{SU}_3(K)\), one has
\(\Delta_k=\{\alpha\}\), and the possibilities up to conjugacy
are
\[
 \begin{split}
  G^\sharp={}&T,\quad B,\quad U(2\alpha)\rtimes T,\\
             &\langle U(2\alpha),U(-2\alpha),T\rangle
               =\operatorname{SL}_2\cdot T,\quad\text{or }G.
 \end{split}
\]
The case \(B\) is excluded as for \(\operatorname{PGL}_2\). The
case \(\operatorname{SL}_2\cdot T\) is excluded because
\(G^\sharp\) contains the \(k\)-torus
\[
  R^1_{K/k}\!\left(R^1_{K_3/K}(\mathbb G_m)\right),
\]
where \(K_3/K\) is the degree-three extension, and the isogeny
class of this torus is irreducible. The same torus excludes
\(T\) and \(U(2\alpha)\rtimes T\). Hence \(G^\sharp=G\).
\end{quote}

For Rosenlicht's question one may reduce to the affine case by
dividing by the center. In the semisimple case, an affirmative
answer would also follow from a strengthening of Chevalley's
existence result for regular points mentioned in \SGARefLink{sga3:xiv:3.20}{3.20}: one would
need a regular element of \(\mathfrak g\) outside
\(\mathfrak h=\operatorname{Lie}(H)\), where \(H\ne G\) is a
prescribed smooth algebraic subgroup of an adjoint semisimple
group \(G\) over the finite field \(k\).

% END INCLUDED SGA3_English_current_buildable_source_20260823/source/inputs/XIV/tex/components/26_expose_XIV_remarks611_en.tex

% BEGIN INCLUDED SGA3_English_current_buildable_source_20260823/source/inputs/XIV/tex/components/27_expose_XIV_appendix_opening_lemma1_en.tex
% Authority: Expo14.pdf, local pp. 31--32 / combined pp. 865--866.
% This opens Serre's Appendix and runs through the setup before Lemma 2.
% Exact next cursor: Lemma 2, local p. 32 / combined p. 866.

\section{Appendix: Existence of Regular Elements over Finite Fields}
\SGARefTarget{sga3:xiv:7}

\begin{center}
\itshape by J.-P. Serre
\end{center}

Throughout the Appendix, \(k\) denotes a finite field,
\(\bar k\) its algebraic closure, and
\(\mathcal G=\operatorname{Gal}(\bar k/k)\). If
\(q=\operatorname{Card}(k)\), then \(\mathcal G\) is
topologically generated by the Frobenius element
\[
  F\colon x\longmapsto x^q.
\]
We shall prove the following theorem.\footnote{The theorem is due
to Chevalley. The editor wishes to express his gratitude to
American Express which, by misplacing a trunk of Chevalley's
manuscripts, obliged him to reconstruct the proof.}

\paragraph{\textbf{Theorem.}}
\SGARefTarget{sga3:xiv:7:theorem}
{\itshape Let \(G\) be an adjoint semisimple group defined over
\(k\), and let \(\mathfrak g\) be its Lie algebra. Then the
\(k\)-Lie algebra \(\mathfrak g(k)\) contains a regular element.}

\paragraph{\textbf{Remarks.}}
\textup{\SGARefLink{sga3:xiv:7:lemma1}{(1)}} Recall that ``adjoint'' means that the center of
\(G\), as a subgroup scheme, is trivial. In view of the Chevalley
seminar, this also means that, for a maximal torus \(T\) of
\(G\), the character group \(X(T)\) of \(T_{\bar k}\) is
generated by the set \(R\) of roots. In particular, the rank of
\(\mathfrak g\) equals \(\dim T\), the rank of \(G\).

\textup{\SGARefLink{sga3:xiv:7:lemma2}{(2)}} The editor does not know whether the adjoint
hypothesis is indispensable.

\paragraph{\textbf{Lemma 1.}}
\SGARefTarget{sga3:xiv:7:lemma1}
{\itshape It is enough to prove the theorem when \(G\) is
geometrically simple.}

Here geometrically simple means that \(G_{\bar k}\) has no smooth
normal subgroup scheme other than \(G_{\bar k}\) and \((e)\);
equivalently, the associated root system is irreducible.

We may first suppose that \(G\) is indecomposable over \(k\).
Then \(G_{\bar k}\) is a product of geometrically simple
components \(S\), permuted transitively by \(\mathcal G\). If
\(\mathcal H\subset\mathcal G\) fixes one component \(S\), and
\(K\) is its fixed field, then \(S\) is defined over \(K\), and
the standard restriction-of-scalars argument gives
\[
  G=R_{K/k}(S).
\]
Similarly,
\[
  \mathfrak g=R_{K/k}(\mathfrak s),
\]
where \(\mathfrak s\) is the Lie algebra of \(S\). If the theorem
holds for \(S\), there is an element
\[
  x\in\mathfrak s(K)=\mathfrak g(k)
\]
regular in \(\mathfrak s\), and one checks immediately that it
is regular in \(\mathfrak g\). This proves the lemma.

Henceforth assume that \(G\) is geometrically simple, and choose
a maximal torus \(T\) of \(G\). Let \(X\) be the character group
of \(T_{\bar k}\), let \(R\) be its root system, \(W\) its Weyl
group, and \(E\) the group of automorphisms of \(X\) preserving
\(R\). Then \(W\) is normal in \(E\).

For another maximal torus \(T'\), write \(X',R',W',E'\) for the
corresponding objects. If
\(y\in G(\bar k)\) satisfies
\[
  yT_{\bar k}y^{-1}=T'_{\bar k},
\]
then \(\operatorname{int}(y)\) identifies
\((X',R',W',E')\) with \((X,R,W,E)\). Replacing \(y\) changes
this identification by an element of \(W\).

Frobenius acts on \(X'\) while preserving \(R'\), and thus
defines an element \(f_{T'}\in E'\). Put \(f=f_T\). After the
above identification \(E'\simeq E\), the element \(f_{T'}\)
becomes an element \(f'\in E\), well defined up to
\(W\)-conjugacy. We now compare \(f'\) with \(f\).

% END INCLUDED SGA3_English_current_buildable_source_20260823/source/inputs/XIV/tex/components/27_expose_XIV_appendix_opening_lemma1_en.tex

% BEGIN INCLUDED SGA3_English_current_buildable_source_20260823/source/inputs/XIV/tex/components/28_expose_XIV_appendix_lemmas2_4_theorem_end_en.tex
% Authority: Expo14.pdf, local pp. 32--33 / combined pp. 866--867.
% Exact next cursor: Proof of Lemma 3, local p. 33 / combined p. 867.

\paragraph{\textbf{Lemma 2.}}
\SGARefTarget{sga3:xiv:7:lemma2}
{\itshape One has \(f'\equiv f\pmod W\). Conversely, every
\(f'\in E\) satisfying this condition arises from a maximal torus
\(T'\) of \(G\).}

Put \(\bar T=T_{\bar k}\) and \(\bar T'=T'_{\bar k}\), and choose
\(y\in G(\bar k)\) such that
\[
  y\bar T y^{-1}=\bar T'.
\]
Since \(y^q\bar T y^{-q}=\bar T'\), the element
\[
  n=y^{-1}y^q
\]
belongs to the normalizer \(N(T)(\bar k)\). The action of \(f'\)
on the points of \(\bar T\) is
\[
 \begin{aligned}
  f'(t)
    &=y^{-1}(yty^{-1})^q y\\
    &=y^{-1}y^q\,t^q\,y^{-q}y
      =nt^qn^{-1}.
 \end{aligned}
\]
If \(w\in W\) is represented by \(n\), then \(f'\) and
\(f\circ w\) have the same action on \(\bar T(\bar k)\), hence
on its characters. Thus \(f'=f\circ w\), and
\(f'\equiv f\pmod W\).

Conversely, represent a given \(w\in W\) by
\(n\in N(T)(\bar k)\). Lang's theorem gives
\[
  n=y^{-1}y^q
\]
for some \(y\in G(\bar k)\). Then
\(y\bar T y^{-1}\) is defined over \(k\), and the preceding
calculation shows that its corresponding element is
\(f\circ w\).

\paragraph{\textbf{Lemma 3.}}
\SGARefTarget{sga3:xiv:7:lemma3}
{\itshape Let \(X,R,W,E\) be as above, with \(R\) irreducible,
and let \(\varphi\in E/W\). There exist a representative
\(f\in E\) of \(\varphi\) and roots
\(\theta_1,\ldots,\theta_n\) such that:
\begin{enumerate}
\item \((\theta_1,\ldots,\theta_n)\) is a basis of \(X\);
\item \(R\) is the union of the orbits of the \(\theta_i\) under
the powers of \(f\): every \(a\in R\) has the form
\(a=f^m\theta_i\) for suitable \(m\) and \(i\).
\end{enumerate}}

The proof will be given below. We first need a lemma of linear
algebra.

\paragraph{\textbf{Lemma 4.}}
\SGARefTarget{sga3:xiv:7:lemma4}
{\itshape Let \(V\) be a vector space over \(k\), and let
\((\theta_1,\ldots,\theta_n)\) be a basis of the dual of
\(V\otimes_k\bar k\). Then there is an \(x\in V\) such that
\(\theta_i(x)\ne0\) for every \(i\).}

Let \(V^*\) be the dual of \(V\), and let \(W_i\) be the
\(\bar k\)-subspace of \(V^*\otimes_k\bar k\) generated by the
Galois conjugates of \(\theta_i\). It descends to a
\(k\)-subspace, again denoted \(W_i\). The natural map
\[
  W_1\otimes\cdots\otimes W_n
     \longrightarrow\bigwedge\nolimits^n V^*
\]
is nonzero: after extension to \(\bar k\), it contains
\(\theta_1\wedge\cdots\wedge\theta_n\ne0\). Hence \(V^*\) has a
basis \(u_1,\ldots,u_n\) with \(u_i\in W_i\). Choose
\(x\in V\) with \(u_i(x)\ne0\) for all \(i\), for example with
\(u_i(x)=1\). If \(\theta_i(x)=0\), every Galois conjugate of
\(\theta_i\) vanishes at \(x\), and so does \(u_i\), a
contradiction.

\paragraph{\textbf{End of the proof of the theorem.}}

By \SGARefLink{sga3:xiv:7:lemma2}{Lemmas 2} and 3, choose a maximal torus \(T\) whose element
\(f\) has the properties of \SGARefLink{sga3:xiv:7:lemma3}{Lemma 3}. Let \(Y\) be the dual of
\(X\). The Lie algebra of \(T\) satisfies
\[
  \mathfrak t(\bar k)\simeq Y\otimes_{\mathbb Z}\bar k,
\]
compatibly with the Galois action. Put
\(V=\mathfrak t(k)\). An element \(x\in V\) is regular precisely
when it is annihilated by no root \(\alpha\in R\), viewed as a
linear form in \(V^*\otimes_k\bar k\).

\SGARefLink{sga3:xiv:7:lemma4}{Lemma 4} supplies \(x\in V\) annihilated by none of the
\(\theta_i\). Every root is a Galois conjugate of one of the
\(\theta_i\), by property \textup{\SGARefLink{sga3:xiv:7:lemma2}{(2)}} of \SGARefLink{sga3:xiv:7:lemma3}{Lemma 3}. Hence no root
annihilates \(x\), and \(x\) is regular. This proves the theorem.

% END INCLUDED SGA3_English_current_buildable_source_20260823/source/inputs/XIV/tex/components/28_expose_XIV_appendix_lemmas2_4_theorem_end_en.tex

% BEGIN INCLUDED SGA3_English_current_buildable_source_20260823/source/inputs/XIV/tex/components/29_expose_XIV_appendix_lemma3_proof_through_case2_en.tex
% Authority: Expo14.pdf, local pp. 33--34 / combined pp. 867--868.
% Exact next cursor: Case (3) in the proof of Lemma 3,
% local p. 34 / combined p. 868.

\paragraph{\textbf{Proof of \SGARefLink{sga3:xiv:7:lemma3}{Lemma 3}.}}

The proof uses Coxeter transformations. Let
\(a_1,\ldots,a_n\) be a simple system of roots of \(R\), and let
\(r_i\) be the reflection corresponding to \(a_i\). Put
\[
  c=r_1\cdots r_n.
\]
Then \(c\in W\). Although \(c\) depends on the simple system and
on its ordering, its conjugacy class does not. It is called the
Coxeter element. We shall use, without proof, the fact that \(1\)
is not an eigenvalue of \(c\).

\paragraph{\textbf{Lemma 5.}}
\SGARefTarget{sga3:xiv:7:lemma5}
Put
\[
  \theta_i=r_n\cdots r_{i+1}(a_i).
\]
Then:
\begin{enumerate}
\item[\textup{\SGAInlineRefTarget{sga3:xiv:7:lemma5:item:a}(a)}] the \(\theta_i\) form a basis of the group
\(X\) generated by the roots;
\item[\textup{\SGAInlineRefTarget{sga3:xiv:7:lemma5:item:b}(b)}] \(\theta_i>0\) and \(c(\theta_i)<0\) for
every \(i\);
\item[\textup{\SGAInlineRefTarget{sga3:xiv:7:lemma5:item:c}(c)}] every root \(a\in R\) such that \(a>0\) and
\(c(a)<0\) is one of the \(\theta_i\);
\item[\textup{\SGAInlineRefTarget{sga3:xiv:7:lemma5:item:d}(d)}] \(R\) is the union of the orbits of the
\(\theta_i\) under the powers of \(c\).
\end{enumerate}

Indeed,
\[
  \theta_i=a_i+\sum_{j>i}m_{ij}a_j,\qquad m_{ij}\in\mathbb Z,
\]
which proves \textup{\SGARefLink{sga3:xiv:7:lemma5:item:a}{(a)}}. Assertions \textup{\SGARefLink{sga3:xiv:7:lemma5:item:b}{(b)}} and
\textup{\SGARefLink{sga3:xiv:7:lemma5:item:c}{(c)}} follow from the observation that \(r_i\) preserves
the sign of every root other than \(\pm a_i\), and reverses the
sign of \(\pm a_i\).

For \textup{\SGARefLink{sga3:xiv:7:lemma5:item:d}{(d)}}, an orbit of \(c\) cannot consist entirely of
positive roots, nor entirely of negative roots. Otherwise the
sum of the roots in that orbit would give a nonzero element of
\(X\) fixed by \(c\), contrary to the fact that \(1\) is not an
eigenvalue. Hence every orbit contains a root \(a>0\) for which
\(c(a)<0\), and \textup{\SGARefLink{sga3:xiv:7:lemma5:item:c}{(c)}} applies.

\paragraph{\textbf{Remark.}}
The proof was recalled only for the reader's convenience; one
could instead refer to the standard texts on Coxeter
transformations, for example Koszul, Séminaire Bourbaki
1959/1960, Exposé~191. They also show that all orbits of \(c\)
have the same number \(h\) of elements and each contains exactly
one \(\theta_i\). In particular,
\[
  h=\frac{\operatorname{Card}(R)}{n}.
\]

We return to \SGARefLink{sga3:xiv:7:lemma3}{Lemma 3} and distinguish three cases.

\medskip
\noindent\textup{\SGARefLink{sga3:xiv:7:lemma1}{(1)}}
Suppose that \(\varphi\in E/W\) is the identity. We must choose
\(f\in W\); take \(f=c\). \SGARefLink{sga3:xiv:7:lemma5}{Lemma 5} shows that this works.

\medskip
\noindent\textup{\SGARefLink{sga3:xiv:7:lemma2}{(2)}}
Suppose that \(R\) has type \(A_n\), where \(n\ge2\) is even, and
that \(\varphi\) is the unique nontrivial element of \(E/W\).
Then \(-1\notin W\), and we take
\[
  f=-c.
\]
The Coxeter element \(c\) has order \(h=n+1\), which is odd. If
\(a\in R\), \SGARefLink{sga3:xiv:7:lemma5}{Lemma 5} gives \(a=c^m\theta_i\) for suitable \(m,i\).
Adding \(h\) to \(m\) if necessary, we may suppose that \(m\) is
even. Then
\[
  a=(-c)^m\theta_i=f^m\theta_i.
\]
Thus the orbits of the \(\theta_i\) under \(f\) fill \(R\).

% END INCLUDED SGA3_English_current_buildable_source_20260823/source/inputs/XIV/tex/components/29_expose_XIV_appendix_lemma3_proof_through_case2_en.tex

% BEGIN INCLUDED SGA3_English_current_buildable_source_20260823/source/inputs/XIV/tex/components/30_expose_XIV_appendix_lemma3_case3_end_en.tex
% Authority: Expo14.pdf, local pp. 34--35 / combined pp. 868--869.
% This closes Lemma 3 and the Appendix.
% Local p. 36 / combined p. 870 is terminal blank.
% Exact next cursor: END OF EXPOSÉ XIV; Exposé XV begins combined p. 871.

\medskip
\noindent\textup{\SGARefLink{sga3:xiv:7:lemma3}{(3)}}
Suppose that \(\varphi\in E/W\) is nontrivial and that \(R\) has
one of the following types:
\[
  A_n\ (n\ge3\text{ odd}),\qquad
  D_4,\qquad D_n\ (n\ge5),\qquad E_6.
\]
These, together with the even \(A_n\) case treated separately
above, are exactly the cases in which \(E/W\) is nontrivial. The
following diagrams record the corresponding Dynkin-diagram
automorphisms.

\begin{center}
\begin{tikzpicture}[
  x=1.05cm,
  y=.92cm,
  >=Latex,
  dynkin node/.style={circle,fill=black,inner sep=2.2pt},
  dynkin edge/.style={line width=.45pt},
  dynkin arrow/.style={line width=.45pt}
]
  % A_n, n odd: reflection of the chain.
  \node[anchor=east] at (-1.4,6.7) {\(A_n,\ n\text{ odd}\geq 3\)};
  \foreach \x/\name in {0/aone,1.3/atwo,2.6/athree,3.9/afour,5.2/afive}
    \node[dynkin node] (\name) at (\x,6.7) {};
  \draw[dynkin edge] (aone)--(atwo)--(athree)--(afour)--(afive);
  \draw[dynkin arrow,<->]
    (aone) .. controls +(1.0,1.35) and +(-1.0,1.35) .. (afive);
  \draw[dynkin arrow,<->]
    (atwo) .. controls +(.65,.72) and +(-.65,.72) .. (afour);

  % D_4: the three outer vertices are cyclically permuted.
  \node[anchor=east] at (-1.4,4.3) {\(D_4\)};
  \node[dynkin node] (dcenter) at (2.6,4.3) {};
  \node[dynkin node] (dtop)    at (1.3,5.15) {};
  \node[dynkin node] (dbottom) at (1.3,3.45) {};
  \node[dynkin node] (dright)  at (4.0,4.3) {};
  \draw[dynkin edge] (dcenter)--(dtop);
  \draw[dynkin edge] (dcenter)--(dbottom);
  \draw[dynkin edge] (dcenter)--(dright);
  \draw[dynkin arrow,->]
    (dright) .. controls +(0,1.05) and +(.75,.48) .. (dtop);
  \draw[dynkin arrow,->]
    (dtop) .. controls +(-1.15,-.18) and +(-1.15,.18) .. (dbottom);
  \draw[dynkin arrow,->]
    (dbottom) .. controls +(.75,-.48) and +(0,-1.05) .. (dright);

  % D_n, n >= 5: interchange of the two terminal vertices.
  \node[anchor=east] at (-1.4,1.85) {\(D_n,\ n\geq 5\)};
  \foreach \x/\name in {0/dnone,1.25/dntwo,2.5/dnthree,3.75/dnbranch}
    \node[dynkin node] (\name) at (\x,1.85) {};
  \node[dynkin node] (dnupper) at (4.85,2.65) {};
  \node[dynkin node] (dnlower) at (4.85,1.05) {};
  \draw[dynkin edge] (dnone)--(dntwo)--(dnthree)--(dnbranch);
  \draw[dynkin edge] (dnbranch)--(dnupper);
  \draw[dynkin edge] (dnbranch)--(dnlower);
  \draw[dynkin arrow,<->]
    (dnupper) .. controls +(.8,-.3) and +(.8,.3) .. (dnlower);

  % E_6: reflection across the central vertical axis.
  \node[anchor=east] at (-1.4,-.75) {\(E_6.\)};
  \foreach \x/\name in {0/eone,1.3/etwo,2.6/ethree,3.9/efour,5.2/efive}
    \node[dynkin node] (\name) at (\x,-.75) {};
  \node[dynkin node] (esix) at (2.6,-1.9) {};
  \draw[dynkin edge] (eone)--(etwo)--(ethree)--(efour)--(efive);
  \draw[dynkin edge] (ethree)--(esix);
  \draw[dynkin arrow,<->]
    (eone) .. controls +(1.0,1.35) and +(-1.0,1.35) .. (efive);
  \draw[dynkin arrow,<->]
    (etwo) .. controls +(.65,.72) and +(-.65,.72) .. (efour);
\end{tikzpicture}
\end{center}

Let \(S\) be a simple system of roots, not yet ordered. The group
\(E\) is the semidirect product of \(W\) with the group
\(\Psi\) of permutations of \(S\) preserving the Cartan matrix,
or equivalently the Dynkin diagram. Thus \(E/W\simeq\Psi\), and
\(\varphi\) corresponds to some \(\psi\in\Psi\).

Inspection of the diagrams shows that each \(\psi\)-orbit in
\(S\) consists of pairwise orthogonal roots. This fails for
\(A_n\) with \(n\) even, which is why that case was separated.
If \(\sigma\) is such an orbit, the reflections \(r_a\),
\(a\in\sigma\), commute; write
\[
  \rho_\sigma=\prod_{a\in\sigma}r_a.
\]
Plainly \(\rho_\sigma\) commutes with \(\psi\).

Choose a total order on \(S\) for which every \(\psi\)-orbit is a
segment, and write \(S=\{a_1,\ldots,a_n\}\). Let
\(\theta_1,\ldots,\theta_n\) and
\(c=r_1\cdots r_n\) be defined as above. The element \(c\) is a
product of the \(\rho_\sigma\), with the orbits in some order, so
it commutes with \(\psi\). Put
\[
  f=c\psi.
\]

The automorphism \(\psi\) permutes the \(\theta_i\). Indeed,
\(\psi(\theta_i)>0\), and, since \(c\) commutes with \(\psi\),
\[
  c(\psi(\theta_i))=\psi(c(\theta_i))<0.
\]
\SGARefLink{sga3:xiv:7:lemma5}{Lemma 5}\textup{\SGARefLink{sga3:xiv:7:lemma5:item:c}{(c)}} therefore gives
\(\psi(\theta_i)=\theta_j\) for some \(j\).

Finally, if \(a\in R\), \SGARefLink{sga3:xiv:7:lemma5}{Lemma 5} gives
\(a=c^m\theta_i\) for some \(m,i\). Hence
\[
  a=f^m(\psi^{-m}\theta_i)=f^m\theta_j
\]
for a suitable \(j\). Thus \(f\) and the \(\theta_i\) satisfy
both conditions of \SGARefLink{sga3:xiv:7:lemma3}{Lemma 3}, completing its proof.

\paragraph{\textbf{Remark.}}
Except in case \textup{\SGARefLink{sga3:xiv:7:lemma2}{(2)}}, one can prove that every orbit of
\(f\) has exactly
\[
  h=\frac{\operatorname{Card}(R)}{n}
\]
elements and contains exactly one \(\theta_i\). In case
\textup{(2)}, some of the \(\theta_i\) are superfluous.

% END INCLUDED SGA3_English_current_buildable_source_20260823/source/inputs/XIV/tex/components/30_expose_XIV_appendix_lemma3_case3_end_en.tex



% END INCLUDED SGA3_English_current_buildable_source_20260823/source/inputs/XIV/body.tex


\SGAStartExpose{XV}
% BEGIN INCLUDED SGA3_English_current_buildable_source_20260823/source/inputs/XV/body.tex
% Mechanically extracted after \mainmatter from the frozen standalone master.

\markboth{Exposé XV. Complements on Sub-tori}{Exposé XV. Complements on Sub-tori}
% BEGIN INCLUDED SGA3_English_current_buildable_source_20260823/source/inputs/XV/tex/components/00_expose_XV_title_and_introduction.tex
% Authority: Expo15.pdf, local p. 1 / combined p. 871 / printed p. 349.

\chapter*{Exposé XV\\Complements on Sub-tori of a Group Prescheme\\
Application to Smooth Groups}
\addcontentsline{toc}{chapter}{Exposé XV. Complements on Sub-tori of a Group Prescheme}
\SGARefTarget{sga3:xv:expose}

\begin{center}
\textit{by M. Raynaud}\footnote{Source-edition note: xy version dated
1 December 2008.}\textsuperscript{,}\footnote{This Exposé and the two following ones
(Exposés XVI and XVII) do not correspond to oral lectures delivered at
the seminar. With some additional material, they develop the substance
of concise manuscript notes written by A. Grothendieck in 1964 for the
present seminar.}
\end{center}

\SGASection{0}{Introduction}
\SGARefTarget{sga3:xv:section:0}

This Exposé supplements and partly recasts Exposés XI and XII; knowledge
of Exposés XIII and XIV is not indispensable. Continuing the effort
begun in XII, we shall work with \(S\)-group preschemes that are not
necessarily affine and not necessarily separated over \(S\).

\SGARefLink{sga3:xv:section:1}{Sections 1}, 2, 3, and 4 are devoted to the study of sub-tori of a group
prescheme. We obtain infinitesimal lifting theorems (§2) and a global
lifting theorem (§4), in which points of finite order (§1) play an
essential role.

\SGARefLink{sga3:xv:section:5}{Sections 5}, 6, and 7 are independent of the preceding sections. The
consideration of infinitesimal neighbourhoods leads to the
representability of the functor of smooth subgroups equal to their
connected normalizer (§5). In §§6 and 7 we turn more particularly to
Cartan subgroups.

Finally, in §8 we give a necessary and sufficient condition for the
functor of sub-tori of a smooth group, or the functor of maximal tori,
to be representable.

% END INCLUDED SGA3_English_current_buildable_source_20260823/source/inputs/XV/tex/components/00_expose_XV_title_and_introduction.tex

% BEGIN INCLUDED SGA3_English_current_buildable_source_20260823/source/inputs/XV/tex/components/01_expose_XV_section1_opening.tex
% Authority: Expo15.pdf, local pp. 1--3 / combined pp. 871--873 /
% printed pp. 350--353.
% Stops immediately before §1.2, Lemma 1.4.

\SGASection{1}{Lifting finite subgroups}
\SGARefTarget{sga3:xv:section:1}

\SGASubsection{1.1}{Finite, smooth, central subgroups of multiplicative type}
\SGARefTarget{sga3:xv:subsection:1:1}

\SGAStatement{Proposition 1.1}{Let \(S\) be an affine scheme, let
\(S_0\) be a closed subscheme of \(S\) defined by a square-zero ideal,
let \(G\) be an \(S\)-group prescheme, and let \(H_0\) be a subgroup
scheme of \(G_0=G\times_S S_0\) that is smooth and of finite type over
\(S_0\), and of multiplicative type. A subgroup scheme of \(G\), of
multiplicative type, lifting \(H_0\) exists if and only if there exists
a subscheme \(H'\) of \(G\), flat over \(S\), lifting \(H_0\).}
\SGARefTarget{sga3:xv:proposition:1:1}

The necessity of the condition is clear; let us prove its sufficiency.
The group \(H_0\) of multiplicative type is quasi-isotrivial (\SGARefLink{sga3:x:corollary:4-5}{X~4.5}).
By Exposé \SGARefLink{sga3:x:proposition:2:1}{X~2.1}, there exist an \(S\)-group \(H\) of multiplicative
type and an \(S_0\)-isomorphism of groups
\[
  u_0:H\times_S S_0\xrightarrow{\sim}H_0.
\]
Since \(H'\) is flat over \(S\) and \(H'\times_S S_0\) is of finite
presentation over \(S_0\) (Exposé \SGARefLink{sga3:ix:proposition:2:1:item:b:01}{IX~2.1(b)}), \(H'\) is of finite
presentation over \(S\). Its fibres are smooth, so \(H'\) is smooth
over \(S\) (EGA IV~17.5.1). Since \(S\) is affine, \(u_0\) therefore
lifts to an \(S\)-morphism of preschemes
\[
  u:H\longrightarrow H'.
\]
It then follows from Exposé \SGARefLink{sga3:iii:theorem:2-1}{III~2.1} and Exposé \SGARefLink{sga3:ix:theorem:3:1}{IX~3.1} that the composite
\[
  v_0:H\times_S S_0\xrightarrow{\sim}H_0\longrightarrow G_0
\]
also lifts to an \(S\)-homomorphism
\[
  v:H\longrightarrow G.
\]
Since \(v_0\) is an immersion, so is \(v\). The image of \(H\) under
\(v\) is therefore a subgroup scheme of \(G\), of multiplicative type,
lifting \(H_0\).

\SGAStatement{Proposition 1.2}{Let \(S\) be a prescheme, let \(S_0\)
be a subprescheme of \(S\) defined by a locally nilpotent sheaf of
ideals, let \(G\) be an \(S\)-group prescheme flat and of finite
presentation over \(S\), and let \(H_0\) be a subgroup scheme of
\(G_0=G\times_S S_0\) that is smooth, finite over \(S_0\), of
multiplicative type, and central. Then there exists a unique subgroup
scheme \(H\) of \(G\), of multiplicative type, lifting \(H_0\).
Moreover, \(H\) is central. (See XVII, Appendix III, 1.)}
\SGARefTarget{sga3:xv:proposition:1:2}

\SGAInlineRefTarget{sga3:xv:proposition:1-2-bis}\SGAStatement{Proposition 1.2 bis}{Let \(S,G,S_0,G_0\) be as above,
let \(H\) be an \(S\)-group scheme of multiplicative type, smooth and
finite over \(S\), and let
\(u_0:H_0=H\times_S S_0\to G_0\) be a central homomorphism. Then
\(u_0\) lifts uniquely to a homomorphism \(u:H\to G\). Moreover,
\(u\) is central.}
\SGARefTarget{sga3:xv:proposition:1:2:bis}

The existence of the lifting \(u\) in \SGARefLink{sga3:xv:proposition:1:2:bis}{Proposition 1.2 bis} follows
readily from \SGARefLink{sga3:xv:proposition:1:2}{Proposition 1.2} by considering the graph of \(u_0\). The
lifting \(u\) is unique and central by Exposé \SGARefLink{sga3:ix:corollary:3:4}{IX~3.4} and Exposé \SGARefLink{sga3:ix:theorem:5:1}{IX~5.1}.

\SGAProof{Proof of \SGARefLink{sga3:xv:proposition:1:2}{Proposition 1.2}}
The uniqueness of \(H\), and the fact that \(H\) is central, follow
from Exposé \SGARefLink{sga3:ix:theorem:5:6:item:b}{IX~5.6(b)} and Exposé \SGARefLink{sga3:ix:corollary:3:4:bis}{IX~3.4 bis}. In view of uniqueness, to
prove existence we may assume that \(S\) is affine and that \(S_0\) is
defined by a square-zero ideal; by \SGARefLink{sga3:xv:proposition:1:1}{Proposition 1.1} it is enough to find
a subscheme of \(G\), flat over \(S\), lifting \(H_0\).

Since \(H_0\) is smooth and finite over \(S_0\), after possibly
restricting \(S\) we may assume that there is an integer \(n>0\),
invertible on \(S\), such that \(H_0={}_nH_0\). Consider the \(n\)-th
power morphism on \(G\):
\[
  u:G\longrightarrow G,\qquad x\longmapsto x^n.
\]
We continue to denote by \({}_nG\) the ``kernel of \(u\)'', that is,
the inverse image under \(u\) of the identity section of \(G\).
Notice that \(u\) is not, in general, a group homomorphism. We first
admit the following lemma.

\SGAStatement{Lemma 1.3}{Let \(k\) be a field, let \(G\) be a group
scheme locally of finite type over \(k\), let \(n\) be an integer prime
to the characteristic of \(k\), and let \(u\) be the \(n\)-th power
morphism on \(G\). Then \(u\) is étale at every point \(x\) of \(G\)
that belongs to the centre of \(G\).}
\SGARefTarget{sga3:xv:lemma:1:3}

Since \(G\) is flat and of finite presentation over \(S\), the
preceding lemma and EGA IV~17.8.2 show that, if \(x\in G\) lies above
\(s\in S\) and belongs to the centre of \(G_s\), then \(u\) is étale
at \(x\). If in addition \(x\in {}_nG\), then \({}_nG\) is étale over
\(S\) at \(x\). By hypothesis, \(H_0\) is central and is contained in
\({}_nG_0\); it is therefore contained in the largest open subset
\(V\) of \({}_nG\) that is étale over \(S\). Since \(H_0\) and
\(V\times_S S_0=V_0\) are étale over \(S_0\), \(H_0\) is an open
subscheme of \(V_0\) (EGA IV~17.8.7 and 17.9.1). The open subprescheme
of \(V\) having the same underlying space as \(H_0\) is consequently
a subscheme of \(G\), flat over \(S\), lifting \(H_0\).

It remains to prove \SGARefLink{sga3:xv:lemma:1:3}{Lemma 1.3}. The largest open subset of \(G\) on
which \(u\) is étale is invariant under extension of the base field
(EGA IV~17.8.2), so we may reduce to the case in which \(x\) is
\(k\)-rational. Let \(t_x\) denote translation by \(x\), a
\(k\)-automorphism of \(G\). Since \(x\) belongs to the centre of
\(G\), the diagram
\[
\begin{tikzcd}[column sep=large,row sep=large]
G \arrow[r,"u"] \arrow[d,"t_x"'] & G \arrow[d,"t_x^n"] \\
G \arrow[r,"u"] & G .
\end{tikzcd}
\]
is commutative. It is therefore enough to prove that \(u\) is étale
at the identity; this was established in VII\textsubscript{A}~\SGARefLink{sga3:xv:definition:8-4}{8.4}.

% END INCLUDED SGA3_English_current_buildable_source_20260823/source/inputs/XV/tex/components/01_expose_XV_section1_opening.tex

% BEGIN INCLUDED SGA3_English_current_buildable_source_20260823/source/inputs/XV/tex/components/02_expose_XV_section1_remainder_and_section2.tex
% Authority: Expo15.pdf, local pp. 3--14 / combined pp. 873--884 /
% printed pp. 353--374.
% Begins at §1.2, Lemma 1.4, and stops immediately before §3.

\SGASubsection{1.2}{Global lifting of finite groups}
\SGARefTarget{sga3:xv:subsection:1:2}

\SGAStatement{\SGAInlineRefTarget{sga3:xv:lemma:1-4}Lemma 1.4}{Let \(A\) be a local ring, separated and
complete for the topology defined by its maximal ideal \(\mathfrak m\),
let
\[
  S=\operatorname{Spec}(A),\qquad
  S_n=\operatorname{Spec}(A/\mathfrak m^n).
\]
Then, for every prescheme \(X\) (respectively, for every \(S\)-prescheme
\(X\)), the canonical map
\[
  \tag{$\ast$}\label{sga3:xv:lemma:1-4:equation:ast}
  \operatorname{Hom}(S,X)
  \longrightarrow
  \varprojlim_n\operatorname{Hom}(S_n,X)
\]
(respectively,
\[
  \tag{$\ast\ast$}\label{sga3:xv:lemma:1-4:equation:ast-ast}
  \Gamma(X/S)
  \longrightarrow
  \varprojlim_n\Gamma(X_n/S_n),
  \qquad X_n=X\times_S S_n)
\]
is bijective.}
\SGARefTarget{sga3:xv:lemma:1:4}

Statement \((\ast\ast)\) is an easy consequence of \((\ast)\). Let us
prove \((\ast)\). Let
\[
  u_n:S_n\longrightarrow X\qquad(n\in\mathbb N)
\]
be a coherent system of morphisms. The image \(y\) of the closed point
of \(S_n\) is independent of \(n\), and \(u_n\) factors through
\(\operatorname{Spec}\mathcal O_y\). On passing to the inverse limit,
the morphisms \(u_n\) define a ring homomorphism
\[
  \widetilde u:\mathcal O_y
  \longrightarrow
  \varprojlim_n(A/\mathfrak m^n)=A.
\]
This proves that \((\ast)\) is surjective; it is injective as soon as
\(A\) is separated for the \(\mathfrak m\)-adic topology.

% Authority page transition: local p. 4 / combined p. 874 / printed p. 354.
\SGAStatement{\SGAInlineRefTarget{sga3:xv:corollary:1-5}Corollary 1.5}{Let \(A\) be a complete noetherian local
ring, let \(\mathfrak m\) be its maximal ideal, let
\[
  S=\operatorname{Spec}(A),\qquad
  S_n=\operatorname{Spec}(A/\mathfrak m^n),
\]
let \(X\) be a finite scheme over \(S\), and let \(Y\) be an
\(S\)-prescheme. Then the canonical map
\[
  \operatorname{Hom}_S(X,Y)
  \longrightarrow
  \varprojlim_n\operatorname{Hom}_{S_n}(X_n,Y_n)
\]
(where \(X_n=X\times_S S_n\), and similarly for \(Y_n\)) is
bijective.}
\SGARefTarget{sga3:xv:corollary:1:5}

Indeed, it follows from EGA II~6.2.5 that \(X\) is a finite disjoint
sum of local \(S\)-schemes finite over \(S\). We are thus reduced to
the case in which \(X\) itself is the spectrum of a complete
noetherian local ring. But
\[
  \operatorname{Hom}_S(X,Y)=\Gamma(Z/X),
  \qquad Z=Y\times_S X,
\]
and the preceding lemma applies.

\SGAStatement{\SGAInlineRefTarget{sga3:xv:proposition:1-6}Proposition 1.6}{Let \(A,S,S_n\) be as above, and let
\(G\) and \(M\) be two \(S\)-group preschemes, with \(M\) finite over
\(S\). Then:
\begin{enumerate}
\item[\normalfont a)]\SGAInlineRefTarget{sga3:xv:proposition:1-6:item:a} The canonical map
\[
  \operatorname{Hom}_{S\text{-}\mathrm{gr}}(M,G)
  \longrightarrow
  \varprojlim_n
  \operatorname{Hom}_{S_n\text{-}\mathrm{gr}}(M_n,G_n)
\]
is bijective.
\item[\normalfont b)] If \(M\) is of multiplicative type and \(G\) is
smooth over \(S\), the canonical map
\[
  \varphi:\operatorname{Hom}_{S\text{-}\mathrm{gr}}(M,G)
  \longrightarrow
  \operatorname{Hom}_{S_0\text{-}\mathrm{gr}}(M_0,G_0)
\]
is surjective. Moreover, if
\(\varphi(u)=\varphi(u')=u_0\), then \(u\) and \(u'\) are conjugate by
an element of \(G(S)\) that reduces to the identity element of
\(G(S_0)\).
\item[\normalfont c)] If \(M\) is of multiplicative type and smooth
over \(S\), if \(G\) is flat and of finite type over \(S\), and if
\(u_0:M_0\to G_0\) is a central homomorphism, then \(u_0\) lifts
uniquely to a homomorphism
\[
  u:M\longrightarrow G.
\]
Moreover, \(u\) is central if \(G\) has connected fibres.
\end{enumerate}}
\SGARefTarget{sga3:xv:proposition:1:6}

\SGAProof{Proof}
\emph{a)} This follows from \SGARefLink{sga3:xv:corollary:1-5}{1.5}, from the fact that
\(M\times_S M\) is finite over \(S\), and from the characterization
of group morphisms as the morphisms for which the familiar diagram
\[
\begin{tikzcd}[column sep=large,row sep=large]
M\times_S M \arrow[r,"u\times u"] \arrow[d]
  & G\times_S G \arrow[d] \\
M \arrow[r,"u"'] & G
\end{tikzcd}
\]
is commutative.

\emph{b)} By Exposé \SGARefLink{sga3:ix:theorem:3:6}{IX~3.6}, one can construct a coherent system of
homomorphisms \(u_n:M_n\to G_n\) lifting a homomorphism
\(u_0:M_0\to G_0\). In view of part a), this proves the first assertion
of b).

If \(u\) and \(u'\) are two liftings of \(u_0\), then \(u_n\) and
\(u'_n\) are conjugate by an element \(g_n\in G(S_n)\) lifting the
identity element of \(G(S_0)\) (Exposé \SGARefLink{sga3:ix:theorem:3:6}{IX~3.6}). The same reference
also implies that the \(g_n\) may be chosen coherently; by \SGARefLink{sga3:xv:lemma:1-4}{1.4} they
therefore arise from a section \(g\in G(S)\). The morphisms \(u\) and
\(\operatorname{Int}(g)u'\) then agree modulo \(\mathfrak m^n\) for
every \(n\), and hence agree by \SGARefLink{sga3:xv:corollary:1-5}{1.5}.

\emph{c)} The existence and uniqueness of \(u\) follow from part a)
and \SGARefLink{sga3:xv:proposition:1:2:bis}{1.2 bis}. If \(G\) has connected fibres, \(u\) is central by
Exposé \SGARefLink{sga3:ix:theorem:5:6:item:a}{IX~5.6(a)}.

% Authority page transition: local p. 5 / combined p. 875 / printed pp. 356--357.
\SGASection{2}{Infinitesimal lifting of sub-tori}
\SGARefTarget{sga3:xv:section:2}

\SGASubsection{2.1}{Statement of the theorem}
\SGARefTarget{sga3:xv:subsection:2:1}

We shall give an infinitesimal lifting theorem for sub-tori of a group
prescheme that does not require smoothness hypotheses (in contrast
with Exposé \SGARefLink{sga3:ix:theorem:3:6:bis}{IX~3.6 bis}) and that also answers a very natural
question:\footnote{Editor's note: The idea of approximating a torus by
its finite subgroups occurs in Grothendieck's proof of the
connectedness of centralizers of tori; see §\SGARefLink{sga3:xv:lemma:4:6}{4.6} of \emph{BIBLE}.}
is it enough to be able to lift ``sufficiently many'' points of finite
order of a sub-torus \(T_0\) in order to ensure that \(T_0\) itself can
be lifted (infinitesimally, of course)?

\SGAStatement{Theorem 2.1}{Let \(S\) be a noetherian affine scheme,
let \(S_0\) be a closed subscheme of \(S\) defined by a square-zero
ideal \(J\), let \(G\) be an \(S\)-group prescheme of finite type, let
\(G_0=G\times_S S_0\), let \(T_0\) be a sub-torus of \(G_0\), and let
\(q>0\) be an integer invertible on \(S\). Suppose that, for every
integer \(n\) equal to a power of \(q\), there exists a subscheme
\(M_n\) of \(G\), flat over \(S\), such that
\[
  M_n\times_S S_0={}_nT_0.
\]
Then there exists a sub-torus \(T\) of \(G\) such that
\(T\times_S S_0=T_0\).}
\SGARefTarget{sga3:xv:theorem:2:1}

\SGARefLink{sga3:xv:theorem:2:1}{Theorem 2.1} will be useful to us through the following two
corollaries.

\SGAStatement{\SGAInlineRefTarget{sga3:xv:corollary:2-2}Corollary 2.2}{Let \(S\) be a locally noetherian
prescheme, let \(S_0\) be a closed subprescheme of \(S\) defined by a
locally nilpotent sheaf of ideals, let \(G\) be an \(S\)-group
prescheme of finite type, let \(T_0\) be a sub-torus of
\(G_0=G\times_S S_0\), and let \(q>0\) be an integer invertible on
\(S\). As \(n\) ranges over the powers of \(q\), let \((M_n)\) be a
coherent system of \(S\)-subgroup schemes of \(G\), of multiplicative
type, lifting the \({}_nT_0\). (The system \((M_n)\) is said to be
coherent if
\[
  M_m={}_m(M_n)
\]
whenever \(m\) divides \(n\).) Then there exists a unique sub-torus
\(T\) of \(G\) such that
\[
  T\times_S S_0=T_0
  \qquad\text{and}\qquad
  {}_nT=M_n
\]
for every \(n\).}
\SGARefTarget{sga3:xv:corollary:2:2}

\SGAStatement{Corollary 2.3}{Let \(G\) be an \(S\)-group prescheme,
flat and of finite presentation over \(S\), let \(S_0\) be a closed
subprescheme of \(S\) defined by a locally nilpotent sheaf of ideals
of finite type, and let \(T_0\) be a central torus of
\(G_0=G\times_S S_0\). Then there exists a unique sub-torus \(T\) of
\(G\) lifting \(T_0\). Moreover, \(T\) is central.}
\SGARefTarget{sga3:xv:corollary:2:3}

\SGAStatement{Remark 2.4}{We leave it to the reader to formulate the
analogues of statements \SGARefLink{sga3:xv:theorem:2:1}{2.1}, \SGARefLink{sga3:xv:corollary:2-2}{2.2}, and \SGARefLink{sga3:xv:corollary:2:3}{2.3} in which, instead of lifting
a sub-torus of \(G_0\), one is given a torus \(T\) over \(S\) and asks
to lift a morphism
\[
  u_0:T_0\longrightarrow G_0.
\]
One reduces to the preceding cases by considering the graph of
\(u_0\).}
\SGARefTarget{sga3:xv:remark:2:4}

Let us show how \SGARefLink{sga3:xv:corollary:2-2}{Corollaries 2.2} and \SGARefLink{sga3:xv:corollary:2:3}{2.3} follow from \SGARefLink{sga3:xv:theorem:2:1}{Theorem 2.1}.

\SGAProof{Proof of \SGARefLink{sga3:xv:corollary:2-2}{Corollary 2.2}}
The uniqueness of \(T\) follows from Exposé \SGARefLink{sga3:ix:corollary:4:8:item:b}{IX~4.8(b)} and Exposé
\SGARefLink{sga3:ix:remark:4:10}{IX~4.10}. To prove existence, we may therefore reduce to the case in
which \(S\) is affine, hence noetherian, and \(S_0\) is defined by a
square-zero ideal.

\SGAStatement{Lemma 2.5}{Let \(G\) be an \(S\)-group prescheme of
finite presentation, and let \(H\) be a subgroup scheme of \(G\), of
multiplicative type and smooth over \(S\). Then
\(\operatorname{Centr}_G(H)\) is representable by an \(S\)-group
subprescheme of \(G\), of finite presentation.}
\SGARefTarget{sga3:xv:lemma:2:5}

When \(G\) is separated over \(S\), the lemma follows from Exposé
\SGARefLink{sga3:viii:examples:6:5:item:e:01}{VIII~6.5(e)}, without any smoothness hypothesis on \(H\). In the present
case, observe that the assertion is local for the fpqc topology. We
may therefore assume that \(H\) is diagonalizable, hence of the form
\(D_S(M)\). We may also assume that \(S\) is affine and then, by EGA
IV~8, noetherian. Since \(H\) is smooth and of finite type over \(S\),
the order \(q\) of the torsion subgroup of \(M\) is invertible on
\(S\) (Exposé \SGARefLink{sga3:viii:proposition:2:1:item:e:01}{VIII~2.1(e)}). It is then immediate (cf. Exposé \SGARefLink{sga3:ix:remark:4:10}{IX~4.10})
that the subgroups \({}_nH\), with \(n\) ranging over the powers of
\(q\), are schematically dense in \(H\) (Exposé \SGARefLink{sga3:ix:definition:4:1}{IX~4.1}). But
\({}_nH\) is a completely split covering of \(S\), that is, it is
isomorphic to a finite disjoint sum of copies of \(S\). Consequently
\[
  \operatorname{Centr}_G({}_nH)=Z_n
\]
is representable as the intersection of the centralizers in \(G\) of
the sections of \({}_nH\) over \(S\). It remains to apply the following
lemma.

\SGAStatement{\SGAInlineRefTarget{sga3:xv:lemma:2-5-bis}Lemma 2.5 bis}{Let \(S\) be a noetherian prescheme, let
\(G\) be an \(S\)-group prescheme of finite type, let \(H\) be a
subgroup of \(G\) of multiplicative type, and let \((H_i)\) be an
increasing filtered family of subgroups of \(G\) of multiplicative
type. Suppose that
\[
  Z_i=\operatorname{Centr}_G(H_i)
\]
is representable by a group subprescheme of \(G\). Then the family
\((Z_i)\) is stationary. If, in addition, the \(H_i\) are schematically
dense in \(H\), then
\[
  Z_i=\operatorname{Centr}_G(H)
\]
for all sufficiently large \(i\).}
\SGARefTarget{sga3:xv:lemma:2:5:bis}

To see that the family \((Z_i)\) is stationary, it is enough to prove
that the family of underlying sets \(\operatorname{ens}(Z_i)\) is
stationary. Indeed, the stationary value will then be a closed subset
of an open subset \(U\) of \(G\). After replacing \(G\) by \(U\), we
are reduced to studying a decreasing filtered family of closed
subpreschemes of a noetherian prescheme. An elementary constructibility
argument reduces us to the case in which \(S\) is integral. We then
have to show that \(\operatorname{ens}(Z_i)\) is stationary over some
nonempty open subset of \(S\). The generic fibre of \(G\) is separated
(Exposé VI\textsubscript{A}~0.3); after restricting \(S\), we may
therefore assume that \(G\) is separated over \(S\) (EGA IV~8). Then
\(Z_i\) is closed in \(G\) (Exposé VIII~6.5(e)).

To prove the last assertion of the lemma, let \(Z\) denote the
stationary value of the family \((Z_i)\). It is clear that
\(\operatorname{Centr}_G(H)\) is a subfunctor of \(Z\); we show that
\(Z\) centralizes \(H\). Let \(E\) be the subprescheme of
\(H\times_S Z\) that is the kernel of the pair of morphisms
\[
\begin{aligned}
  H\times_S Z&\rightrightarrows G,\\
  (h,c)&\longmapsto c,\\
  (h,c)&\longmapsto hch^{-1}.
\end{aligned}
\]
The prescheme \(E\) contains \(H_i\times_S Z\) for every \(i\). The
\(H_i\) are flat over \(S\), so EGA IV~11.10.9 shows that, for every
point \(s\) of \(S\), the \((H_i)_s\) are schematically dense in
\(H_s\), and \(H_i\times_S Z\) is schematically dense in
\(H\times_S Z\). Since \(G_s\) is separated, \(E_s\) is closed in
\((H\times_S Z)_s\), and is therefore equal to it. Thus \(E\) is
closed in \(H\times_S Z\), and hence equals \(H\times_S Z\). This says
that \(Z\) centralizes \(H\).

Returning to the proof of \SGARefLink{sga3:xv:corollary:2-2}{2.2}, \SGARefLink{sga3:xv:lemma:2:5}{Lemma 2.5} shows that
\[
  Z_n=\operatorname{Centr}_G(M_n)
\]
is representable, and \SGARefLink{sga3:xv:lemma:2-5-bis}{Lemma 2.5 bis} shows that the decreasing family
of subpreschemes \((Z_n)\) is stationary; let \(Z\) be its stationary
value. The group \(Z\) contains \(T_0\) and all the \(M_n\). After
replacing \(G\) by \(Z\), we may therefore suppose that the \(M_n\)
are central.

The hypotheses of \SGARefLink{sga3:xv:theorem:2:1}{Theorem 2.1} are now satisfied, so there exists a
sub-torus \(T\) of \(G\) lifting \(T_0\). The groups \({}_nT\) and
\(M_n\) are two liftings of \({}_nT_0\); they are therefore conjugate
(Exposé \SGARefLink{sga3:ix:theorem:3:2:bis}{IX~3.2 bis}), and consequently coincide because \(M_n\) is
central. This torus \(T\) has the required properties.

\SGAProof{Proof of \SGARefLink{sga3:xv:corollary:2:3}{Corollary 2.3}}
The uniqueness of \(T\) follows from Exposé \SGARefLink{sga3:ix:theorem:5:1:bis}{IX~5.1 bis}, and its
centrality follows from Exposé \SGARefLink{sga3:ix:theorem:5:6:item:b}{IX~5.6(b)}. By the usual procedure, this
observation reduces us first to the case in which \(S\) is affine
(and hence \(S_0\) is defined by a nilpotent ideal of finite type),
and then to the noetherian case. After restricting \(S\) if necessary,
we may suppose that there exists an integer \(q\) invertible on \(S\).
\SGARefLink{sga3:xv:corollary:2:3}{Corollary 2.3} then follows from \SGARefLink{sga3:xv:corollary:2-2}{2.2} and \SGARefLink{sga3:xv:proposition:1:2}{1.2}.

\SGAStatement{Remark 2.6}{\SGARefLink{sga3:xv:corollary:2:3}{Corollary 2.3} remains true if the torus
\(T_0\) is replaced by a smooth central subgroup of multiplicative
type of \(G_0\).}
\SGARefTarget{sga3:xv:remark:2:6}

% Authority page transition: local p. 7 / combined p. 877 / printed p. 362.
\SGASubsection{2.2}{Proof of \SGARefLink{sga3:xv:theorem:2:1}{Theorem 2.1}}
\SGARefTarget{sga3:xv:subsection:2:2}

\medskip\noindent\textbf{a) Reduction to the case
\(\boldsymbol{T_0=\mathbb G_{m,S_0}}\).}

By \SGARefLink{sga3:xv:proposition:1:1}{1.1}, we may assume that \(M_n\) is a subgroup of multiplicative
type. Using Exposé \SGARefLink{sga3:ix:theorem:3:2:bis}{IX~3.2 bis}, we may suppose that the family
\((M_n)\) is coherent (\SGARefLink{sga3:xv:corollary:2-2}{2.2}). The centralizer \(Z_n\) of \(M_n\) in
\(G\) is representable (\SGARefLink{sga3:xv:lemma:2:5}{2.5}), and the filtered family \((Z_n)\) is
stationary (\SGARefLink{sga3:xv:lemma:2-5-bis}{2.5 bis}). After replacing \(G\) by \(Z_n\) for sufficiently
large \(n\), we may therefore suppose that \(T_0\) and the \(M_n\) are
central. The uniqueness of a lifting of \(T_0\) is then guaranteed
by Exposé \SGARefLink{sga3:ix:theorem:5:1:bis}{IX~5.1 bis}.

Proceeding as in the proof of \SGARefLink{sga3:xv:proposition:1:1}{1.1}, we may suppose that there exist an
\(S\)-torus \(T\) and an \(S_0\)-isomorphism
\[
  u_0:T\times_S S_0\xrightarrow{\sim}T_0.
\]
Lifting \(u_0\) is equivalent to lifting \(T_0\). By uniqueness, it is
enough to prove the existence of a lifting of \(u_0\) after an affine
faithfully flat extension \(S'\to S\) of finite type (fpqc descent).
This allows us to suppose that
\[
  T=\mathbb G_{m,S}^{\,r}
\]
(Exposé \SGARefLink{sga3:x:corollary:4-5}{X~4.5}). If the restriction of \(u_0\) to each factor
\(\mathbb G_m\) lifts to an \(S\)-morphism---necessarily central---then
these liftings immediately give a lifting of \(u_0\). In short, we may
suppose that \(T_0=\mathbb G_{m,S_0}\).

\medskip\noindent\textbf{b) Definition of the obstruction to the
existence of a lifting of \(\boldsymbol{T_0}\).}

By \SGARefLink{sga3:xv:proposition:1:1}{1.1}, to prove \SGARefLink{sga3:xv:theorem:2:1}{Theorem 2.1} it is enough to find a subscheme of \(G\),
flat over \(S\), lifting \(T_0\). We shall see that the obstruction to
the existence of such a lifting can be defined as an element of a
certain \(\operatorname{Ext}^1\) group of sheaves of modules.

Let \(U\) be an open subset of \(G\) in which \(T_0\) is closed, and
again denote by \(U\) (respectively, \(U_0\)) the open subscheme of
\(G\) (respectively, \(G_0\)) having this open subset as its underlying
space. The sheaf \(\mathcal O_{T_0}\), regarded as a sheaf on \(U\), is
therefore a quotient of \(\mathcal O_{U_0}\). Let
\[
  h:\mathcal O_{U_0}\longrightarrow\mathcal O_{T_0}
\]
be the canonical epimorphism.

\SGAStatement{\SGAInlineRefTarget{sga3:xv:lemma:2-7}Lemma 2.7}{The canonical map
\[
  h'=\operatorname{id}_J\otimes h:
  J\otimes_{S_0}\mathcal O_{U_0}
  \longrightarrow
  J\otimes_{S_0}\mathcal O_{T_0}
\]
factors, necessarily uniquely, as \(i\circ j_U\), where \(j_U\) is
the canonical epimorphism
\[
  J\otimes_{S_0}\mathcal O_{U_0}
  \longrightarrow
  J\mathcal O_U\simeq J\mathcal O_{U_0}.
\]}
\SGARefTarget{sga3:xv:lemma:2:7}

We must show that \(h'(K)=0\), where \(K\) is the kernel of \(j_U\).
For every integer \(n\) equal to a power of \(q\), there is an
epimorphism
\[
  h_n:\mathcal O_{T_0}\longrightarrow\mathcal O_{{}_nT_0}.
\]
Since \({}_nT_0\) lifts to a scheme \(M_n\) flat over \(S\), the
canonical morphism
\[
  j_n:J\otimes_{S_0}\mathcal O_{{}_nT_0}
  \xrightarrow{\sim}J\mathcal O_{M_n}
\]
is an isomorphism. The commutative diagram
\[
\begin{tikzcd}[column sep=large,row sep=large]
K \arrow[r]
  & J\otimes_{S_0}\mathcal O_{U_0}
      \arrow[r,"j_U"] \arrow[d,"h'"]
  & J\mathcal O_U \arrow[r,hook] \arrow[dl,"i"]
  & \mathcal O_U \\
  & J\otimes_{S_0}\mathcal O_{T_0}
      \arrow[d,"h'_n"] & & \\
  & J\otimes_{S_0}\mathcal O_{{}_nT_0}
      \arrow[r,"j_n"',"\sim"] & J\mathcal O_{M_n}
      \arrow[r,hook] & \mathcal O_{M_n}
\end{tikzcd}
\]
shows that \(h'(K)\) is contained in \(\ker h'_n\) for every \(n\).
It is therefore contained in \(\bigcap_n\ker h'_n\), and it remains
to prove that this intersection is zero. Now
\[
  \mathcal O_{T_0}=\mathcal O_{S_0}[\mathbb Z],
  \qquad
  \mathcal O_{{}_nT_0}
     =\mathcal O_{S_0}[\mathbb Z/n\mathbb Z].
\]
Let
\[
  a=\sum_{m\in\mathbb Z}a_m\otimes m
  \in J\otimes_{S_0}\mathcal O_{T_0}.
\]
The \(a_m\) are sections of \(J\), almost all zero. Choose \(n\) large
enough that the indices \(m\) for which \(a_m\ne0\) have distinct
images in \(\mathbb Z/n\mathbb Z\). If \(a\in\ker h'_n\), then
necessarily \(a=0\). Thus \(\bigcap_n\ker h'_n=0\), proving \SGARefLink{sga3:xv:lemma:2-7}{2.7}.

% Authority page transition: local p. 9 / combined p. 879 / printed p. 365.
Let \(K_0\) be the kernel of \(h\). Consider the diagram
\[
\begin{tikzcd}[column sep=large,row sep=large]
  &&& 0 \arrow[d] \\
  & J\mathcal O_U \arrow[d,"\sim"] && K_0 \arrow[d] \\
0 \arrow[r]
  & J\mathcal O_{U_0} \arrow[r] \arrow[d,"i"']
  & \mathcal O_U \arrow[r] & \mathcal O_{U_0}
      \arrow[r] \arrow[d,"h"'] & 0 \\
  & J\otimes_{S_0}\mathcal O_{T_0}
  && \mathcal O_{T_0} \arrow[d] \\
  &&& 0 .
\end{tikzcd}
\]
The sheaf \(\mathcal O_U\) defines an element
\[
  \Phi\in
  \operatorname{Ext}^{1}_{\mathcal O_U}
  (\mathcal O_{U_0},J\mathcal O_{U_0}).
\]
Let \(\Psi\) be the element of
\[
  \operatorname{Ext}^{1}_{\mathcal O_U}
  (K_0,J\otimes_{S_0}\mathcal O_{T_0})
\]
obtained from \(\Phi\), by the bifunctoriality of
\(\operatorname{Ext}^{1}\), through the morphisms
\[
  K_0\longrightarrow\mathcal O_{U_0},
  \qquad
  i:J\mathcal O_{U_0}
    \longrightarrow J\otimes_{S_0}\mathcal O_{T_0}.
\]
It follows from Exposé \SGARefLink{sga3:iii:lemma:4-1}{III~4.1} and the infinitesimal criterion of
flatness (cf. Exposé \SGARefLink{sga3:iii:proposition:4-3}{III~4.3}) that a subscheme of \(U\), flat over
\(S\) and lifting \(T_0\), exists if and only if \(\Psi=0\).

Since \(T_0\) is affine, Exposé \SGARefLink{sga3:iii:proposition:4-5}{III~4.5} and \SGARefLink{sga3:xv:lemma:4:6}{4.6} show that it is enough
to construct, locally on \(U\), a subscheme flat over \(S\) lifting
\(T_0\). It therefore suffices to show that the image of \(\Psi\) in
the sheaf
\[
  \mathscr E=
  \mathcal Ext^{1}_{\mathcal O_U}
  (K_0,J\otimes_{S_0}\mathcal O_{T_0})
\]
is zero. We again denote this element of
\(\Gamma(U,\mathscr E)\) by \(\Psi\).

\medskip\noindent\textbf{c) Reduction to the case in which \(S\) is
local artinian with algebraically closed residue field.}

Since \(K_0\) and \(J\otimes_{S_0}\mathcal O_{T_0}\) are coherent
sheaves, so is \(\mathscr E\); moreover, \(\mathscr E\) is supported
on \(T_0\), because the same is true of
\(J\otimes_{S_0}\mathcal O_{T_0}\). To show that the section \(\Psi\)
of \(\mathscr E\) is zero, it is enough to show that, for every point
\(u\) of \(T_0\), its image in the fibre of \(\mathscr E\) at \(u\)
is zero. Formation of the \(\mathcal Ext^i\) of coherent sheaves
commutes with flat extensions of the base,\footnote{Editor's note:
cf. EGA \(0_{\mathrm{III}}\), 12.3.5.} so we are reduced to proving the
existence of a lifting of
\[
  T_0\cap\operatorname{Spec}\mathcal O_u
\]
flat over \(S\). Let \(s\) be the image of \(u\) in \(S\). We may
replace \(S\) by \(\operatorname{Spec}\mathcal O_s\) and \(G\) by
\(G\times_S\operatorname{Spec}\mathcal O_s\).

After another faithfully flat extension if necessary, we may suppose
that \(\mathcal O_s\) has an algebraically closed residue field
(EGA \(0_{\mathrm{III}}\), 10.3.1). Let \(\mathfrak m\) be the maximal
ideal of \(\mathcal O_s\), and set
\[
  S_n=\operatorname{Spec}(\mathcal O_s/\mathfrak m^n).
\]
Suppose that, for every \(n>0\), the obstruction to lifting
\[
  T_0\times_S S_n=(T_0)_n
\]
to a subscheme of \(G_n=G\times_S S_n\), flat over \(S_n\), has been
shown to vanish, and let \(T_n\) be the unique flat lifting of
\((T_0)_n\) over \(S_n\) that is a sub-torus of \(G_n\). If
\(n>n'\), it is clear that
\[
  T_n\times_{S_n}S_{n'}=T_{n'}.
\]
If \(u\) is a point of \(T_0\) lying over \(s\), the lemma below,
applied to the coherent system of liftings
\[
  T_n\cap\operatorname{Spec}\mathcal O_u
\]
of
\((T_0)_n\cap\operatorname{Spec}\mathcal O_u\), shows that
\(T_0\cap\operatorname{Spec}\mathcal O_u\) indeed has a lifting flat
over \(\mathcal O_s\). We are therefore reduced to proving that
\(\Psi=0\) when \(S=S_n\) is the spectrum of an artinian local ring
with algebraically closed residue field, and to proving the following
lemma.

\SGAStatement{\SGAInlineRefTarget{sga3:xv:lemma:2-8}Lemma 2.8}{Let \(A\to B\) be a local homomorphism of
noetherian local rings, let \(\mathfrak m\) be the maximal ideal of
\(A\), let \(J\) be a square-zero ideal of \(A\), let \(M\) be a
finite-type \(B\)-module, and set
\[
  A_0=A/J,\qquad B_0=B/JB,\qquad M_0=M/JM.
\]
Let \(N_0\) be a quotient \(B_0\)-module of \(M_0\), flat over \(A_0\).
For every \(n>0\), let \(A_n,A_{0,n}\), and so forth denote the objects
obtained by the base extension
\[
  A\longrightarrow A/\mathfrak m^n=A_n,
\]
and let \(\overline{J}_n\) be the image of \(J\) in \(A_n\). For every \(n>0\),
let \(N_n\) be a quotient \(B_n\)-module of \(M_n\), flat over \(A_n\),
lifting \(N_{0,n}\), and suppose that, for \(n\ge n'\), \(N_{n'}\) is
obtained from \(N_n\) by the base extension \(A_n\to A_{n'}\). Then
there exists a quotient \(B\)-module \(N\) of \(M\), flat over \(A\),
lifting \(N_0\).}
\SGARefTarget{sga3:xv:lemma:2:8}

\SGAProof{Proof of \SGARefLink{sga3:xv:lemma:2-8}{2.8}}
Let \(P_0\) be the kernel of the epimorphism \(M_0\to N_0\). For every
\(n>0\), there is a commutative diagram with exact horizontal rows:
\[
\tag{$\ast_n$}\label{sga3:xv:lemma:2-8:equation:ast-n}
\begin{tikzcd}[column sep=2em,row sep=large]
&&&& 0 \arrow[d] & \\
&&&& P_{0,n} \arrow[d] & \\
& 0 \arrow[r]
  & \overline{J}_nM_n \arrow[r] \arrow[dd]
  & M_n \arrow[r] \arrow[dd]
  & M_{0,n} \arrow[r] \arrow[dd] & 0 \\
\overline{J}_n\otimes_{A_{0,n}}M_{0,n}
      \arrow[urr] \arrow[drr] &&&&& \\
& 0 \arrow[r]
  & \overline{J}_n\otimes_{A_{0,n}}N_{0,n} \arrow[r]
  & N_n \arrow[r] \arrow[d]
  & N_{0,n} \arrow[r] \arrow[d] & 0 \\
&&& 0 & 0 &
\end{tikzcd}
\]
Moreover, by hypothesis, the diagram \((\ast)_{n+1}\) reduces modulo
\(\mathfrak m^n\) to \((\ast)_n\).

% Authority page transition: local pp. 10--11 / combined pp. 880--881 /
% printed pp. 367--368.
The Artin--Rees lemma (Bourbaki, \emph{Algèbre commutative}, Chapter
III, §3, Corollary 1) shows that the filtrations on \(JM\),
\(J\otimes_{A_0}M_0\), and \(J\otimes_{A_0}N_0\), defined by the
kernels of the respective morphisms
\[
\begin{gathered}
  JM\longrightarrow \overline{J}_nM_n,\\
  J\otimes_{A_0}M_0
    \longrightarrow \overline{J}_n\otimes_{A_{0,n}}M_{0,n},\\
  J\otimes_{A_0}N_0
    \longrightarrow \overline{J}_n\otimes_{A_{0,n}}N_{0,n},
\end{gathered}
\]
are \(\mathfrak mB\)-good. Passing to the inverse limit in the
diagrams \((\ast)_n\), and writing \(\widehat Q\) for the separated
completion of a \(B\)-module \(Q\) for the \(\mathfrak mB\)-adic
topology, we obtain a commutative diagram in which the two horizontal
rows remain exact:
\[
\tag{$\widehat\ast$}\label{sga3:xv:lemma:2-8:equation:widehat-ast}
\begin{tikzcd}[column sep=large,row sep=large]
&&&& 0 \arrow[d] & \\
&&&& \widehat P_0 \arrow[d] & \\
& 0 \arrow[r]
  & \widehat{JM} \arrow[r] \arrow[dd]
  & \widehat M \arrow[r] \arrow[dd]
  & \widehat M_0 \arrow[r] \arrow[dd] & 0 \\
\widehat{J\otimes_{A_0}M_0}
      \arrow[urr] \arrow[drr] &&&&& \\
& 0 \arrow[r]
  & \widehat{J\otimes_{A_0}N_0} \arrow[r]
  & \varprojlim_n N_n \arrow[r] \arrow[d]
  & \widehat N_0 \arrow[r] \arrow[d] & 0 \\
&&& 0 & 0 &
\end{tikzcd}
\]
Now \((B,\mathfrak mB)\) is a Zariski ring and
\(J\otimes_{A_0}N_0\) is of finite type over \(B\), hence separated
for the \(\mathfrak mB\)-adic topology. The diagram
\((\widehat\ast)\) then shows that the morphism
\[
  J\otimes_{A_0}M_0\longrightarrow J\otimes_{A_0}N_0
\]
induced by the epimorphism \(M_0\to N_0\) factors through \(JM\):
\[
\begin{tikzcd}[column sep=large,row sep=large]
J\otimes_{A_0}M_0 \arrow[rr,"\mathrm{can.}"] \arrow[dr]
  && JM \arrow[dl] \\
& J\otimes_{A_0}N_0 &
\end{tikzcd}
\]

Under these conditions, Exposé \SGARefLink{sga3:iii:lemma:4-1}{III~4.1} and \SGARefLink{sga3:iii:proposition:4-3}{4.3} show that a quotient
\(B\)-module \(N\) of \(M\), flat over \(A\) and lifting \(N_0\),
exists if and only if a certain element
\[
  g\in E=\operatorname{Ext}^{1}_B
  (P_0,J\otimes_{A_0}N_0)
\]
vanishes. More precisely, the exact sequence
\[
  0\longrightarrow JM\longrightarrow M
  \longrightarrow M_0\longrightarrow0
\]
defines an element \(f\in F\), where
\[
  F=\operatorname{Ext}^{1}_B(M_0,JM),
\]
and \(g\) is the image of \(f\) under the natural morphism \(F\to E\)
arising by bifunctoriality from
\[
  P_0\longrightarrow M_0,
  \qquad
  JM\longrightarrow J\otimes_{A_0}N_0.
\]
The diagram \((\widehat\ast)\) and Exposé \SGARefLink{sga3:iii:lemma:4-1}{III~4.1} show that the image
of \(g\) in
\[
  \widehat E\simeq
  \operatorname{Ext}^{1}_{\widehat B}
  (\widehat P_0,\widehat{J\otimes_{A_0}N_0})
\]
is zero. But \(E\) is of finite type over \(B\), and hence separated
for the \(\mathfrak mB\)-adic topology. Consequently \(g=0\), which
completes the proof of \SGARefLink{sga3:xv:lemma:2-8}{2.8}.

\medskip\noindent\textbf{d) Study of \(\boldsymbol{\mathscr E}\).}

We now suppose that \(S\) is the spectrum of an artinian local ring
\(A\) with algebraically closed residue field \(k\). Let \(A_0\) be
the ring of \(S_0\), and let \(\mathfrak m_0\) be the maximal ideal
of \(A_0\). Since \(S_0\) is artinian, \(T_0\) is closed in \(G\)
(Exposé VI\textsubscript{B}~\SGARefLink{cor:sga3-VIB-1.4.2}{1.4.2}). We may therefore take \(U=G\), so
that
\[
  \mathscr E=
  \mathcal Ext^{1}_{\mathcal O_G}
  (K_0,J\otimes_{S_0}\mathcal O_{T_0}).
\]
Let \(B_0\) be the ring of the affine \(S_0\)-scheme
\(T_0=\mathbb G_{m,S_0}\), and let \(\overline B_0\) be the ring of its special
fibre
\[
  T_0\times_{S_0}\operatorname{Spec}(k)=\mathbb G_{m,k}.
\]
The sheaf \(\mathscr E\) is a coherent
\(\mathcal O_{T_0}\)-module, hence is defined by a finite-type
\(B_0\)-module, which we denote by \(E\).

Consider the graded module associated with \(E\) and the
\(\mathfrak m_0B_0\)-adic filtration:
\[
  E_r=\mathfrak m_0^rE/\mathfrak m_0^{r+1}E.
\]
Each \(E_r\) is a finite-type \(\overline B_0\)-module, and \(E_r=0\) for all
sufficiently large \(r\), since \(S_0\) is artinian.

Let \(g\) be a section of \(G\) over \(S\) whose restriction to
\(S_0\) is a section \(g_0\) of \(T_0\). Left translation by \(g\)
defines an ``automorphism of the situation'' from the point of view
of the obstruction problem. In particular, \(g\) induces an
\(S_0\)-automorphism of the sheaf \(\mathscr E\) that fixes the
obstruction \(\Psi\). More precisely, \(g\) defines a semilinear
automorphism of the \(B_0\)-module \(E\), relative to the
\(A_0\)-automorphism of \(B_0\) defined by translation by \(g_0\) in
\(T_0\). Reducing modulo \(\mathfrak m_0^{r+1}\), \(g\) therefore
defines a semilinear automorphism of \(E_r\), relative to the
\(k\)-automorphism of \(\overline B_0\) defined by translation by
\[
  g_0\times_{S_0}\operatorname{Spec}(k)
\]
in \(T_0\times_{S_0}\operatorname{Spec}(k)\).

\SGAStatement{Lemma 2.9}{For every integer \(r\ge0\), \(E_r\) is a
locally free \(\overline B_0\)-module.}
\SGARefTarget{sga3:xv:lemma:2:9}

Let \(x\) be a point of \(T_0\), let \(\kappa(x)\) be its residue
field, and let
\[
  (E_r)_x=E_r\otimes_{\overline B_0}\kappa(x)
\]
be the ``fibre'' of \(E_r\) at \(x\). Denote by \(\ell(x)\) the rank
of \((E_r)_x\) over \(\kappa(x)\), and by \(\ell\) the maximum of
\(\ell(x)\) as \(x\) ranges over the points of \(T_0\). Let \(L\) be
the largest closed subscheme of
\[
  \operatorname{Spec}\overline B_0=\mathbb G_{m,k}
\]
over which \(E_r\) is locally free of rank \(\ell\) (TDTE IV, \SGARefLink{sga3:xv:lemma:3:6}{Lemma
3.6}). Let \(\beta\in L(k)\), which exists because \(L\) is of finite
type over the algebraically closed field \(k\), and let
\(\alpha\in\mathbb G_{m,k}(k)\) be a point whose order is a power
\(n\) of \(q\). The point \(\alpha\) is \(k\)-rational. Since, by
hypothesis, \({}_nT_0\) lifts to a subscheme \(M_n\) étale over \(S\),
there is a section \(a\) of \(G\) over \(S\) lifting \(\alpha\) and
whose restriction to \(S_0\) is a section of \(T_0\).

The preceding observations show that the fibres of \(E_r\) at the
points \(\beta\) and \(\alpha\beta\) of \(\mathbb G_{m,k}(k)\) are
\(k\)-isomorphic. The points whose order is a power of \(q\) are dense
in \(\mathbb G_{m,k}\), as are their translates by \(\beta\). Since
\(L\) is closed in \(\mathbb G_{m,k}\) and \(\mathbb G_{m,k}\) is
reduced, \(L=\mathbb G_{m,k}\), and \(E_r\) is locally free of rank
\(\ell\) over \(\overline B_0\).

\medskip\noindent\textbf{e) Completion of the proof of \SGARefLink{sga3:xv:theorem:2:1}{Theorem 2.1}.}

% Authority page transition: local p. 13 / combined p. 883 / printed p. 372.
Let \(K_0(n)\) be the sheaf of ideals of \(\mathcal O_{G_0}\)
defining the closed subscheme \({}_nT_0\). The sheaf
\(K_0\)\footnote{Editor's note: Introduced after \SGARefLink{sga3:xv:lemma:2-7}{Lemma 2.7}.} is
therefore a subsheaf of \(K_0(n)\). For brevity, set
\[
  R=J\otimes_{S_0}\mathcal O_{T_0},
  \qquad
  R(n)=J\otimes_{S_0}\mathcal O_{{}_nT_0}.
\]
The sheaf \(R(n)\) is a quotient of \(R\), and we have the diagram
\[
\begin{tikzcd}[column sep=large,row sep=large]
&& K_0 \arrow[d] \\
&& K_0(n) \arrow[d] \\
0 \arrow[r]
  & J\mathcal O_G \arrow[r] \arrow[d]
  & \mathcal O_G \arrow[r]
  & \mathcal O_{G_0} \arrow[r] & 0 \\
  & R \arrow[d] \\
  & R(n)
\end{tikzcd}
\]
Using the bifunctoriality of
\(\operatorname{Ext}^{1}_{\mathcal O_G}\), we obtain the commutative
diagram
\[
\begin{tikzcd}[column sep=huge,row sep=large]
\mathcal{E}xt^{1}_{\mathcal O_G}
  (\mathcal O_{G_0},J\mathcal O_G)
  \arrow[dr] \\
&
\mathcal{E}xt^{1}_{\mathcal O_G}(K_0(n),R)
  \arrow[r] \arrow[d]
&
\mathcal{E}xt^{1}_{\mathcal O_G}(K_0,R)=\mathscr E
  \arrow[d] \\
&
\mathcal{E}xt^{1}_{\mathcal O_G}(K_0(n),R(n))
  \arrow[r]
&
\mathcal{E}xt^{1}_{\mathcal O_G}(K_0,R(n)).
\end{tikzcd}
\]
Again let \(\Phi\) denote the element of
\(\operatorname{Ext}^{1}_{\mathcal O_G}
(\mathcal O_{G_0},J\mathcal O_G)\) defined by the exact sequence
\[
  0\longrightarrow J\mathcal O_G
  \longrightarrow\mathcal O_G
  \longrightarrow\mathcal O_{G_0}
  \longrightarrow0,
\]
so that \(\Psi\) is the image of \(\Phi\) in \(\mathscr E\). Since
\({}_nT_0\) lifts to a subscheme \(M_n\) of \(G\), flat over \(S\),
the image of \(\Phi\) in
\[
  \operatorname{Ext}^{1}_{\mathcal O_G}(K_0(n),R(n))
\]
is zero (Exposé \SGARefLink{sga3:iii:lemma:4-1}{III~4.1}). A fortiori, the image of \(\Phi\) in
\(\operatorname{Ext}^{1}_{\mathcal O_G}(K_0,R(n))\), which is also
the image of \(\Psi\), is zero.

\SGAStatement{Lemma 2.10}{For every integer \(n\), the canonical
morphism
\[
  \operatorname{Ext}^{1}_{\mathcal O_G}(K_0,R)
    \otimes_{B_0}\mathcal O_{{}_nT_0}
  \longrightarrow
  \operatorname{Ext}^{1}_{\mathcal O_G}(K_0,R(n))
\]
is injective.}
\SGARefTarget{sga3:xv:lemma:2:10}

Indeed, the affine scheme \(T_0=\mathbb G_{m,S_0}\) has ring
\[
  B_0=A_0[T,T^{-1}].
\]
The subscheme \({}_nT_0\) is defined by the vanishing of
\[
  h(n)=T^n-1.
\]
This section is regular (EGA \(0_{\mathrm{IV}}\), 15.2.2) and remains
regular after every base change \(S'_0\to S_0\). Hence there is an
exact sequence of sheaves
\[
  0\longrightarrow\mathcal O_{T_0}
  \xrightarrow{\,h(n)\,}\mathcal O_{T_0}
  \longrightarrow\mathcal O_{{}_nT_0}
  \longrightarrow0.
\]
Since \({}_nT_0\) is flat over \(S\), tensoring over
\(\mathcal O_{S_0}\) with \(J\) gives the exact sequence
\[
  0\longrightarrow R\xrightarrow{\,h(n)\,}R
  \longrightarrow R(n)\longrightarrow0,
\]
and then the exact sequence of Ext groups
\[
  \cdots\longrightarrow
  \operatorname{Ext}^{1}_{\mathcal O_G}(K_0,R)
  \xrightarrow{\,h(n)\,}
  \operatorname{Ext}^{1}_{\mathcal O_G}(K_0,R)
  \longrightarrow
  \operatorname{Ext}^{1}_{\mathcal O_G}(K_0,R(n)),
\]
which proves the lemma.

The foregoing shows that, for every integer \(n\) equal to a power of
\(q\), the image of \(\Psi\) in
\(\mathscr E/h(n)\mathscr E\) is zero. To prove that \(\Psi=0\), it
is enough to show that
\[
  \Psi\in\mathfrak m_0^r\mathscr E
  \quad\Longrightarrow\quad
  \Psi\in\mathfrak m_0^{r+1}\mathscr E.
\]
Let \(\overline\Psi\) be the image of \(\Psi\) in \(E_r\). There is
an element \(\Psi(n)\in\mathscr E\) such that
\[
  \Psi=\Psi(n)h(n)
  \qquad(n\text{ a power of }q).
\]
We observed that the image \(\overline{h(n)}\) of \(h(n)\) in \(\overline B_0\)
is again \(\overline B_0\)-regular. Since \(E_s\) is locally free over \(\overline B_0\)
for every \(s\) (\SGARefLink{sga3:xv:lemma:2:9}{2.9}), multiplication by \(\overline{h(n)}\) on
\(E_s\) is injective. It follows immediately that
\(\Psi(n)\in\mathfrak m_0^r\mathscr E\). Let
\(\overline{\Psi(n)}\) be its image in \(E_r\); then
\[
  \overline\Psi=\overline{h(n)}\,\overline{\Psi(n)}
  \qquad(n\text{ a power of }q).
\]
Thus the set \(F\) of points of \(\mathbb G_{m,k}\) at which
\(\overline\Psi\) vanishes contains the dense set of points whose
order is a power of \(q\). Moreover, \(F\) is closed because \(E_r\)
is locally free over \(\overline B_0\), and \(\mathbb G_{m,k}\) is reduced.
Therefore \(\overline\Psi=0\).

This completes the proof of \SGARefLink{sga3:xv:theorem:2:1}{Theorem 2.1}.

% Exact stop: immediately before §3, “Characterization of a sub-torus
% by its underlying set”, authority local p. 14 / combined p. 884 /
% printed p. 374.

% END INCLUDED SGA3_English_current_buildable_source_20260823/source/inputs/XV/tex/components/02_expose_XV_section1_remainder_and_section2.tex

% BEGIN INCLUDED SGA3_English_current_buildable_source_20260823/source/inputs/XV/tex/components/03_expose_XV_section3.tex
% Authority: Expo15.pdf, local pp. 14--26 / combined pp. 884--896 /
% printed pp. 374--398.
% Begins at §3, “Characterization of a sub-torus by its underlying
% set”, and stops immediately before §4.

\SGASection{3}{Characterization of a sub-torus by its underlying set}
\SGARefTarget{sga3:xv:section:3}

\SGASubsection{3.1}{Statement of the theorem}
\SGARefTarget{sga3:xv:subsection:3:1}

\noindent\textbf{Notation.}---If \(X\) is a prescheme,
\(\operatorname{ens}(X)\) denotes the underlying set of \(X\). If
\(X\) and \(S'\) are two \(S\)-preschemes,
\[
  X_{S'}=X\times_S S',
  \qquad u:X_{S'}\longrightarrow X
\]
is the canonical morphism, and if \(E\) is a subset of
\(\operatorname{ens}(X)\), we denote by \(E_{S'}\) (or
\(E\times_S S'\)) the subset \(u^{-1}(E)\) of
\(\operatorname{ens}(X_{S'})\).

\SGAStatement{\SGAInlineRefTarget{sga3:xv:theorem:3-1}Theorem 3.1}{Let \(S\) be a locally noetherian
prescheme, let \(G\) be an \(S\)-group prescheme of finite type, let
\(q>0\) be an integer invertible on \(S\), and let \(E\) be a subset
of \(\operatorname{ens}(G)\). Consider the following assertions about
\(E\):
\begin{enumerate}
\item[\normalfont\SGAInlineRefTarget{sga3:xv:theorem:3-1:item:i}(i)] The set \(E\) is the underlying set of a
sub-torus \(T\) of \(G\).
\item[\normalfont\SGAInlineRefTarget{sga3:xv:theorem:3-1:item:ii}(ii)]
  \begin{enumerate}
  \item[\normalfont a)] For every point \(s\) of \(S\), there is a
  sub-torus \(T_s\) of \(G_s\) such that
  \[
    E\cap\operatorname{ens}(G_s)=E_s
  \]
  is the underlying set of \(T_s\).
  \item[\normalfont b)] As \(n\) ranges over the powers of \(q\),
  there is a coherent family (cf.\ \SGARefLink{sga3:xv:corollary:2-2}{2.2}) \(M_n\) of subgroup schemes
  of \(G\), of multiplicative type, such that for every point \(s\)
  of \(S\),
  \[
    (M_n)_s={}_nT_s.
  \]
  \end{enumerate}
\item[\normalfont\SGAInlineRefTarget{sga3:xv:theorem:3-1:item:iii}(iii)]
  \begin{enumerate}
  \item[\normalfont a)] As in \textup{\SGARefLink{sga3:xv:theorem:3-1:item:ii}{(ii)} a)}.
  \item[\normalfont b)] The set \(E\) is locally closed in
  \(\operatorname{ens}(G)\), and the dimension of the fibers of \(E\)
  over \(S\) is locally constant.
  \end{enumerate}
\item[\normalfont\SGAInlineRefTarget{sga3:xv:theorem:3-1:item:iv}(iv)]
  \begin{enumerate}
  \item[\normalfont a)] As in \textup{\SGARefLink{sga3:xv:theorem:3-1:item:ii}{(ii)} a)}.
  \item[\normalfont b)] For every \(S\)-scheme \(S'\) that is the
  spectrum of a complete discrete valuation ring with algebraically
  closed residue field, \(E_{S'}\) is the underlying set of a
  sub-torus of \(G_{S'}\).
  \end{enumerate}
\end{enumerate}
Then the following implications hold:
\begin{enumerate}
\item[\normalfont A)] \(\textup{(i)}\Longleftrightarrow
\textup{(ii)}\Longrightarrow\textup{(iii)}
\Longrightarrow\textup{(iv)}\).
\item[\normalfont B)] If \(G\) is separated over \(S\), then
\(\textup{(iii)}\Longleftrightarrow\textup{(iv)}\), and moreover
\(E\) is closed.
\item[\normalfont C)] Conditions \textup{\SGARefLink{sga3:xv:theorem:3-1:item:i}{(i)}}, \textup{\SGARefLink{sga3:xv:theorem:3-1:item:ii}{(ii)}}, and
\textup{\SGARefLink{sga3:xv:theorem:3-1:item:iii}{(iii)}} (and also \textup{\SGARefLink{sga3:xv:theorem:3-1:item:iv}{(iv)}} if \(G\) is separated over
\(S\)) are equivalent in the following two cases:
  \begin{enumerate}
  \item[\normalfont First case:]
    \begin{enumerate}
    \item[\normalfont a)] The prescheme \(S\) is reduced, or \(G\)
    is flat over \(S\); and
    \item[\normalfont b)] for every point \(s\) of \(S\), \(T_s\)
    is a central torus of \(G_s\).
    \end{enumerate}
  \item[\normalfont Second case:] \(S\) is normal.
  \end{enumerate}
\end{enumerate}
Moreover, in the two cases above, the torus \(T\) with underlying
set \(E\) is unique.}
\SGARefTarget{sga3:xv:theorem:3:1}

% Authority page transition: local p. 15 / combined p. 885 /
% printed pp. 375--376.
\SGAStatement{\SGAInlineRefTarget{sga3:xv:remark:3-2}Remarks 3.2}{\SGAInlineRefTarget{sga3:xv:remark:3-2:item:a}\textup{a)} When \(S\) is reduced, there
is no need in \textup{\SGARefLink{sga3:xv:theorem:3-1:item:ii}{(ii)}} to assume that the family \(M_n\) is
coherent, since this condition follows from the other hypotheses.
Indeed, if \(m\) divides \(n\), the subgroups \({}_mM_n\) and \(M_m\)
are étale over \(S\), hence reduced, and have the same underlying
space; they therefore coincide.

\textup{b)} If \(S\) is no longer assumed normal, it is in general no
longer true that \(\textup{(iii)}\Rightarrow\textup{(i)}\), even when
\(S\) is reduced and geometrically unibranch and \(G\) is a smooth
group scheme over \(S\). Indeed, consider the Borel subgroup of
\(\operatorname{SL}_{2,S}\) consisting of matrices of the form
\[
  \begin{pmatrix}t&u\\0&t^{-1}\end{pmatrix},
\]
where \(S\) is the affine curve over a field \(k\) with ring
\[
  k[x,y]/(y^2-x^3).
\]
Define \(E\) as follows. Above the cusp of \(S\), where \(x=y=0\),
take the diagonal torus \(u=0\). Above the complementary open subset
\(x\ne0\), take the torus obtained from the diagonal torus by
conjugation by
\[
  \begin{pmatrix}1&y/x\\0&1\end{pmatrix}.
\]
The resulting set \(E\) satisfies \textup{\SGARefLink{sga3:xv:theorem:3-1:item:iii}{(iii)} a)}; it is also
closed in \(G\), and the reduced subscheme having \(E\) as its
underlying set has equations
\[
  \begin{cases}
    xu+y(t-t^{-1})=0,\\
    u^2-x(t-t^{-1})^2=0.
  \end{cases}
\]
It is therefore not a sub-torus of \(G\), since its fiber above the
cusp is not reduced.}
\SGARefTarget{sga3:xv:remarks:3:2}

\noindent\textbf{Plan of the proof of \SGARefLink{sga3:xv:theorem:3-1}{3.1}.} In a first part we shall
establish the following ``easy'' implications:
\[
\begin{tikzcd}[column sep=huge,row sep=large]
 & \textup{(ii)} \arrow[dr,Rightarrow] & \\
\textup{(i)} \arrow[ur,Rightarrow] \arrow[dr,Rightarrow] && \textup{(iv)} \\
 & \textup{(iii)} \arrow[ur,Rightarrow] &
\end{tikzcd}
\]
and
\[
  \textup{(iv)}
  \Longrightarrow
  [\,\textup{(iii)}\text{ and }E\text{ is closed}\,]
  \quad\text{if \(G\) is separated over \(S\)}.
\]
The more delicate implications will be proved in three stages:
\begin{enumerate}
\item[\normalfont I)] Reduction of
\(\textup{(iii)}\Rightarrow\textup{(i)}\), under the hypotheses of
C), first case, to the case in which \(S\) is normal.
\item[\normalfont II)] \(\textup{(iii)}\Rightarrow\textup{(ii)}\)
when \(S\) is normal.
\item[\normalfont III)] \(\textup{(ii)}\Rightarrow\textup{(i)}\).
\end{enumerate}

\SGASubsection{3.2}{Proof of the ``easy'' assertions in \SGARefLink{sga3:xv:theorem:3-1}{3.1}}
\SGARefTarget{sga3:xv:subsection:3:2}

The implications \(\textup{(i)}\Rightarrow\textup{(ii)}\) and
\(\textup{(i)}\Rightarrow\textup{(iii)}\) are trivial.

\medskip
\noindent\(\textup{(iii)}\Rightarrow\textup{(iv)}\).
After replacing \(S\) by \(S'\), we may assume that \(S\) is the
spectrum of a discrete valuation ring. Let \(t\) be the generic point
of \(S\), and let \(s\) be its closed point.

Since \(E\) is locally closed, there is a reduced subprescheme of
\(G\) with underlying space \(E\); we also denote it by \(E\). Its
generic fiber \(E_t\) is therefore the sub-torus \(T_t\) of \(G_t\),
of rank \(r\), by \textup{\SGARefLink{sga3:xv:theorem:3-1:item:iii}{(iii)} a)}. Let \(E'\) be the schematic
closure of \(E_t\) in \(E\). The prescheme \(E'\) is irreducible, and
its closed fiber \(E'_s\) is nonempty (it contains the unit section
of \(G_s\)); hence \(E'_s\) has dimension \(r\) (EGA IV~14.3.10).
But \(E'_s\) is a closed subscheme of \(E_s\), while \(E_s\) has
dimension \(r\) by \textup{\SGARefLink{sga3:xv:theorem:3-1:item:iii}{(iii)} b)} and is irreducible because it
has the same underlying space as \(T_s\). Thus \(E'_s\) and \(E_s\)
have the same underlying space, and consequently \(E'=E\). This
proves that \(E\) is flat over \(S\).

% Authority page transition: local pp. 16--17 / combined pp. 886--887 /
% printed pp. 377--379.
Let \(E''\) now be the schematic closure of \(E_t\) in \(G\). Then
\(E''\) is a subgroup prescheme of \(G\), flat over \(S\)
(Exposé \SGARefLink{sga3:viii:lemma:7:1}{VIII~7.1}), which contains \(E'\), and therefore \(E\). Its
closed fiber \(E''_s\) is an algebraic subgroup of \(G_s\) of
dimension \(r\) (\emph{loc.\ cit.}). Since \(T_s\) is a closed
irreducible subset of \(G_s\) of dimension \(r\), \(E_s\) has the
same underlying set as the identity component \((E''_s)^0\) of
\(E''_s\). Let \((E'')^0\) denote the ``identity component'' of
\(E''\), that is, the open subgroup of \(E''\) complementary to the
union of the irreducible components of \(E''_s\) that do not contain
the origin. Then
\[
  E=\operatorname{ens}\bigl((E'')^0\bigr).
\]
Since \(E\) and \(E''\) are reduced, one even has
\(E=(E'')^0\). Finally, \(E\) is a subgroup prescheme of \(G\), flat
and of finite type over \(S\), with connected fibers, hence separated
(Exposé VI\textsubscript{B}~\SGARefLink{sga3:xv:theorem:5-2}{5.2}); its generic fiber is the torus
\(T_t\), and its reduced closed fiber is the torus \(T_s\). It follows
that \(E\) is a torus (Exposé \SGARefLink{sga3:x:theorem:8-8}{X~8.8}).

\medskip
\noindent\(\textup{(ii)}\Rightarrow\textup{(iv)}\).
Again we may assume that \(S\) is the spectrum of a discrete valuation
ring, and we retain the notation introduced above. The schematic
closure \(T\) of \(T_t\) in \(G\) is a subgroup prescheme of \(G\),
flat over \(S\), and it contains \(M_n\) for every \(n\) that is a
power of \(q\). Consequently, its closed fiber is a closed subscheme
of \(G_s\) containing \((M_n)_s\) for every \(n\), hence containing
\(T_s\). For dimension reasons, the ``identity component'' \(T^0\)
of \(T\) has \(E\) as its underlying set, and the preceding argument
shows that \(T^0\) is a sub-torus of \(G\).

\medskip
\noindent\(\textup{(iv)}\Rightarrow
[\,\textup{(iii)}\text{ and }E\text{ is closed}\,]\), if \(G\) is
separated over \(S\). Let us show that \(E\) is
closed.\footnote{Editor's note: The implication
\(\textup{(iv)}\Rightarrow\textup{(iii)}\) will appear in the course
of the proof.} We first prove the following lemma.

\SGAStatement{Lemma 3.3}{If the conditions stated in \SGARefLink{sga3:xv:theorem:3-1}{3.1}
\textup{\SGARefLink{sga3:xv:theorem:3-1:item:iv}{(iv)}} are satisfied, \(E\) is a locally constructible subset
of \(\operatorname{ens}(G)\).}
\SGARefTarget{sga3:xv:lemma:3:3}

By the usual method, we are reduced to the case in which \(S\) is
noetherian and integral, with generic point \(\eta\). After shrinking
\(S\), we may assume (Exposé VI\textsubscript{B}~\SGARefLink{par:sga3-VIB-10.10}{10.10}) that there is
a subgroup scheme \(H\) of \(G\), flat over \(S\) and with connected
fibers, whose generic fiber \(H_\eta\) is \(T_\eta\). To prove \SGARefLink{sga3:xv:lemma:3:3}{Lemma
3.3} it is enough to show that \(\operatorname{ens}(H)=E\).

Let \(s\) be a point of \(S\). By EGA II~7.1.9 there is an
\(S\)-scheme \(S'\), the spectrum of a discrete valuation ring which
we may assume complete with algebraically closed residue field, whose
generic point \(t'\) maps to \(\eta\) and whose closed point \(s'\)
maps to \(s\). By \textup{\SGARefLink{sga3:xv:theorem:3-1:item:iv}{(iv)} b)}, there is a sub-torus \(T_{S'}\)
of \(G_{S'}\) having \(E_{S'}\) as its underlying space. The two
subgroup preschemes \(T_{S'}\) and
\(H_{S'}=H\times_S S'\) of \(G_{S'}\) are flat over \(S'\), have
connected fibers, and have the same generic fiber \(T_{t'}\).
Therefore both coincide with the identity component of the schematic
closure of \(T_{t'}\) in \(G_{S'}\). Hence
\(\operatorname{ens}(H_s)=E_s\), and therefore
\(\operatorname{ens}(H)=E\), proving the lemma.

Now that \(E\) is known to be locally constructible, to see that it
is closed it is enough to prove that every specialization in \(G\)
of a point of \(E\) again belongs to \(E\). By the usual technique,
we are reduced to the case in which \(S\) is the spectrum of a
complete discrete valuation ring with algebraically closed residue
field. In that case the sub-torus of \(G\) with underlying space
\(E\), supplied by \textup{\SGARefLink{sga3:xv:theorem:3-1:item:iv}{(iv)} b)}, is closed in \(G\), because
\(G\) is separated (Exposé \SGARefLink{sga3:viii:proposition:7:12}{VIII~7.12}).

\SGASubsection{3.3}{Continuation of the proof of \SGARefLink{sga3:xv:theorem:3-1}{3.1}}
\SGARefTarget{sga3:xv:subsection:3:3}

\noindent\textbf{I) Reduction of
\(\textup{(iii)}\Rightarrow\textup{(i)}\), case C, first case, to
the case in which \(S\) is normal.}

\smallskip
\noindent\textit{a) Reduction to the case in which \(S\) is affine
and reduced.}\quad
We therefore assume that, for every point \(s\) of \(S\), \(E_s\)
is the underlying space of a central sub-torus of \(G_s\). The
uniqueness of \(T\) then follows from Exposé \SGARefLink{sga3:ix:theorem:5:1:bis}{IX~5.1 bis}, and moreover
\(T\) will necessarily be central (\emph{loc.\ cit.}). In view of
this uniqueness, in order to prove the existence of \(T\) we may
suppose that \(S\) is noetherian and affine, with ring \(A\). If
\(S\) is not reduced, \(G\) is flat over \(S\) by hypothesis. By
\SGARefLink{sga3:xv:corollary:2:3}{2.3} it is then enough to solve the problem for \(S_{\mathrm{red}}\)
and \(G\times_S S_{\mathrm{red}}\). We may therefore suppose in
addition that \(S\) is reduced.

\smallskip
\noindent\textit{b) Reduction to the case in which \(A\) is an
algebra of finite type over \(\mathbb Z\).}\quad
We shall use the following lemma.

\SGAStatement{Lemma 3.4}{Let \(k\) be a field, \(G\) an algebraic
\(k\)-group, and \(E\) a subset of \(G\). Let \(k'/k\) be a field
extension, and suppose that \(E_{k'}\) is the underlying set of a
sub-torus \(T_{k'}\) of \(G_{k'}\). If \(T_{k'}\) is central in
\(G_{k'}\), or if \(k\) is perfect, then \(E\) is the underlying set
of a sub-torus \(T\) of \(G\).}
\SGARefTarget{sga3:xv:lemma:3:4}

\SGAProof{Proof}
By fpqc descent it is enough to prove that the two inverse images of
\(T_{k'}\) in \(G_{k''}\), where
\[
  k''=k'\otimes_k k',
\]
coincide. They have the same underlying set. If \(T_{k'}\) is
central, they therefore coincide by Exposé \SGARefLink{sga3:ix:theorem:5:1:bis}{IX~5.1 bis}. If \(k\) is
perfect, \(k''\) is reduced; the two inverse images, being smooth
over \(k''\), are reduced and consequently coincide.

\SGAStatement{\SGAInlineRefTarget{sga3:xv:remark:3-5}Remarks 3.5}{\textup{(1)} \SGARefLink{sga3:xv:lemma:3:4}{Lemma 3.4} applies in
particular when the residue fields of \(S\) are perfect.

\textup{\SGARefLink{sga3:xv:theorem:5:2}{(2)}} Under the hypotheses of the first case in \SGARefLink{sga3:xv:theorem:3-1}{Theorem 3.1}
C, for every \(S\)-scheme \(S'\) as in \SGARefLink{sga3:xv:theorem:3-1}{3.1} \textup{\SGARefLink{sga3:xv:theorem:3-1:item:iv}{(iv)} b)}, the
sub-torus of \(G_{S'}\) having \(E_{S'}\) as its underlying space is
central.}
\SGARefTarget{sga3:xv:remarks:3:5}

\SGAStatement{Lemma 3.6}{Let \(S\) be the projective limit of a
filtered inverse system of affine schemes \(S_i\), and let \(H\) be
an \(S\)-group of multiplicative type and of finite type. There
exist an index \(i\), a group \(H_i\) of multiplicative type over
\(S_i\), and an \(S\)-isomorphism
\[
  H_i\times_{S_i}S\xrightarrow{\sim}H.
\]
If \(H\) is isotrivial, \(H_i\) may be chosen isotrivial.}
\SGARefTarget{sga3:xv:lemma:3:6}

\SGAProof{Proof}
Since \(H\) is of finite presentation over \(S\), there are an index
\(i\), an \(S_i\)-group prescheme \(H_i\), of finite presentation
over \(S_i\), and an \(S\)-isomorphism
\[
  H_i\times_{S_i}S\xrightarrow{\sim}H
\]
(Exposé VI\textsubscript{B}~\SGARefLink{prop:sga3-VIB-10.16}{10.16}). After increasing \(i\), we may
suppose that \(H_i\) is of multiplicative type
(Exposé \SGARefLink{sga3:ix:proposition:2:11}{IX~2.11}). It remains to prove that \(H_i\) can be chosen
quasi-isotrivial, and isotrivial when \(H\) is isotrivial.

We may thus suppose that \(H\) is trivialized by an étale,
surjective morphism \(S'\to S\) of finite presentation. For \(i\)
sufficiently large, there are an \(S_i\)-prescheme \(S'_i\), an
étale, surjective morphism \(S'_i\to S_i\) of finite presentation,
and an \(S\)-isomorphism
\[
  S'_i\times_{S_i}S\xrightarrow{\sim}S'
\]
(EGA IV~17.16). For \(j\) sufficiently large set
\[
  S'_j=S'_i\times_{S_i}S_j,\qquad
  H'_j=H_j\times_{S_j}S'_j,\qquad
  H'=H\times_S S'.
\]
By the choice of \(S'\), there are a finitely generated abelian
group \(M\) and an \(S'\)-isomorphism of group schemes
\[
  D_{S'}(M)\xrightarrow{\sim}H'.
\]
Since the \(S'_i\) are quasi-compact and
\(S'=\varprojlim S'_i\), Exposé VI\textsubscript{B}~10 shows that
there are an index \(j\) and an \(S'_j\)-isomorphism of group
schemes
\[
  D_{S'_j}(M)\xrightarrow{\sim}H'_j.
\]
This says precisely that \(H_j\) is a quasi-isotrivial group of
multiplicative type. When \(H\) is isotrivial, one argues similarly
using a trivialization of \(H\) by a finite étale morphism
\(S'\to S\).

We can now carry out reduction b). The ring \(A\) of \(S\) is the
filtered direct limit of its subrings \(A_i\) of finite type over
\(\mathbb Z\). Put \(S_i=\operatorname{Spec}A_i\), and let
\[
  u_j:S\longrightarrow S_j,\qquad
  u_{jk}:S_k\longrightarrow S_j\quad(k\geq j)
\]
be the transition morphisms. Since \(S\) is noetherian and \(G\) is
of finite presentation over \(S\), there are an index \(i\) and an
\(S_i\)-group prescheme \(G_i\), of finite type over \(S_i\), such
that \(G\) is \(S\)-isomorphic to \(G_i\times_{S_i}S\). Likewise,
since \(E\) is locally closed in \(G\), we may suppose that there is
a locally closed subset \(E_i\) of \(G_i\) such that
\[
  E=E_i\times_{S_i}S
\]
(EGA IV~8.3.11). For \(j>i\), put
\[
  G_j=G_i\times_{S_i}S_j,\qquad
  E_j=E_i\times_{S_i}S_j,
\]
and let \(Q_j\) be the set of points \(s\) of \(S_j\) for which
\((E_j)_s\) is the underlying set of a central sub-torus of
\((G_j)_s\).

It follows from \SGARefLink{sga3:xv:lemma:3:4}{3.4} that
\[
  Q_k=u_{jk}^{-1}(Q_j)\quad(k\geq j),
\]
and by hypothesis
\[
  \operatorname{ens}(S)=u_j^{-1}(Q_j)\quad(j\geq i).
\]
I further claim that \(Q_j\) is ind-constructible
(EGA IV~1.9.4). Indeed, since \(S_j\) is noetherian, it is enough
(EGA IV~1.9.10) to see the following. Let \(S\) be a noetherian
integral scheme with generic point \(\eta\), and suppose that
\(E_\eta\) is the underlying set of a central sub-torus \(T_\eta\)
of \(G_\eta\). Then there is a neighborhood \(U\) of \(\eta\) such
that \(E_s\) has the same property for every \(s\in U\). After
shrinking \(S\), we may suppose that the center \(Z\) of \(G\) is
representable and that there is a subgroup scheme \(T\) of \(G\)
whose generic fiber is \(T_\eta\)
(Exposé VI\textsubscript{B}, \S10). The two locally closed subsets
\(\operatorname{ens}(T)\) and \(E\) of \(\operatorname{ens}(G)\)
coincide on the generic fiber. We can therefore find a neighborhood
\(U\) of \(\eta\) such that, above \(U\),
\[
  \operatorname{ens}(T)=E
\]
(EGA IV~8.3.11). For the same reasons, we may suppose that \(T\) is
central above \(U\); it is a torus by 3.6.

Now that \(Q_j\) is known to be ind-constructible, EGA IV~8.3.4
implies that \(Q_j=\operatorname{ens}(S_j)\) for \(j\) sufficiently
large. We may then replace \(S,G,E\) by \(S_j,G_j,E_j\), which
reduces us to the case in which \(S\) is an affine reduced scheme of
finite type over \(\mathbb Z\).

\smallskip
\noindent\textit{c) Reduction to the case in which \(S\) is the
spectrum of a complete reduced noetherian local ring.}\quad
Because the torus with underlying space \(E\) is unique, the usual
technique (EGA IV~8) and Lemma 3.6 allow us to replace \(S\) by the
spectrum of the local ring \(A\) at a point of \(S\). Let
\(\widehat S\) be the spectrum of the completion \(\widehat A\) of
\(A\) for the topology defined by its maximal ideal. Since \(A\) is
the localization of an algebra of finite type over \(\mathbb Z\),
\(\widehat S\) is still reduced (EGA IV~7.6.5).

I claim that it is enough to solve the problem after the base
extension \(\widehat S\to S\). Indeed, let \(\widehat T\) be the
sub-torus of \(G_{\widehat S}\) with underlying space
\(E_{\widehat S}\). Its two inverse images in
\(G_{\widehat S\times_S\widehat S}\) are central sub-tori with the
same underlying space; they therefore coincide
(Exposé \SGARefLink{sga3:ix:theorem:5:1:bis}{IX~5.1 bis}). By fpqc descent, \(\widehat T\) comes from a
torus \(T\) of \(G\) that solves the problem (cf.\ \SGARefLink{sga3:xv:lemma:3:4}{3.4}).

\smallskip
\noindent\textit{d) A descent lemma.}\quad
Recall the following properties of finite morphisms, noted in
TDTE I. Let \(S,S'\) be two preschemes and
\(u:S'\to S\) a finite morphism. Then:
\begin{enumerate}
\item[\normalfont\SGAInlineRefTarget{sga3:xv:lemma:3:6:item:1}(1)] The morphism \(u\) is an epimorphism if and
only if the canonical morphism of sheaves
\[
  \mathcal O_S\longrightarrow u_*(\mathcal O_{S'})
\]
is injective.
\item[\normalfont(2)] The morphism \(u\) is an effective
epimorphism (Exposé \SGARefLink{sga3:iv:definition:1:3}{IV~1.3}) if and only if the canonical diagram
\[
  0\longrightarrow\mathcal O_S
  \longrightarrow u_*(\mathcal O_{S'})
  \rightrightarrows
  u_*(\mathcal O_{S'})\otimes_{\mathcal O_S}u_*(\mathcal O_{S'})
\]
is exact.
\item[\normalfont\SGAInlineRefTarget{sga3:xv:lemma:3:6:item:3}(3)] If in addition \(S\) is noetherian and \(u\)
is an epimorphism, \(u\) is the composite of a finite sequence of
finite effective epimorphisms.
\end{enumerate}

We can now prove the following lemma, whose proof uses a non-flat
descent technique.

\SGAStatement{\SGAInlineRefTarget{sga3:xv:lemma:3-7}Lemma 3.7}{Let \(S\) be a locally noetherian
prescheme, \(G\) an \(S\)-group prescheme of finite type, \(S'\) a
prescheme, and \(u:S'\to S\) a finite morphism. For every
\(S\)-prescheme \(T\), let \(\mathrm M(T)\) be the set of subgroups
of \(G_T\) of multiplicative type \textup{(respectively, the set of
central subgroups of \(G_T\) of multiplicative type)}, so
that \(\mathrm M\) is naturally a contravariant functor on
\(\operatorname{Sch}/S\). If \(u\) is an effective epimorphism
\textup{(respectively, an epimorphism)}, then the diagram of
sets
\[
  \tag{*}\label{sga3:xv:lemma:3-7:equation:star}
  \mathrm M(S)\longrightarrow\mathrm M(S')
  \rightrightarrows\mathrm M(S'\times_S S')
\]
is exact.}
\SGARefTarget{sga3:xv:lemma:3:7}

\SGAProof{Proof}
\textit{i) Reduction to the case in which \(u\) is an effective
epimorphism.}\quad
Here we consider the functor of central subgroups of \(G\) of
multiplicative type. The injectivity of
\(\mathrm M(S)\to\mathrm M(S')\) is local on \(S\); once this
injectivity is granted, the exactness of \((*)\) is also local on
\(S\). We may therefore suppose that \(S\) is affine and
noetherian.

Consider the case in which \(u:S'\to S\) is the composite of two
finite epimorphisms
\[
  S'\xrightarrow{v}S''\xrightarrow{w}S.
\]
I claim that if both diagrams
\[
  \tag{*'}\label{sga3:xv:lemma:3-7:equation:star-prime}
  \mathrm M(S'')\longrightarrow\mathrm M(S')
  \rightrightarrows\mathrm M(S'\times_{S''}S')
\]
and
\[
  \tag{*''}\label{sga3:xv:lemma:3-7:equation:star-prime-prime}
  \mathrm M(S)\longrightarrow\mathrm M(S'')
  \rightrightarrows\mathrm M(S''\times_S S'')
\]
are exact, then so is \((*)\).

Indeed, the injectivity of
\(\mathrm M(S)\to\mathrm M(S')\) is clear. If \(T'\) belongs to
\[
  \operatorname{Ker}\bigl(
    \mathrm M(S')\rightrightarrows
    \mathrm M(S'\times_S S')\bigr),
\]
then, a fortiori, \(T'\) belongs to
\[
  \operatorname{Ker}\bigl(
    \mathrm M(S')\rightrightarrows
    \mathrm M(S'\times_{S''}S')\bigr).
\]
By the exactness of \((*')\), it therefore comes from a unique
element \(T''\) of \(\mathrm M(S'')\). It is enough to show that
\(T''\) belongs to
\[
  \operatorname{Ker}\bigl(
    \mathrm M(S'')\rightrightarrows
    \mathrm M(S''\times_S S'')\bigr),
\]
for then the exactness of \((*'')\) descends \(T''\) to \(S\).
Let \(T''_1,T''_2\) be the two inverse images of \(T''\) in
\(\mathrm M(S''\times_S S'')\). These are central subgroups of
\(G_{S''\times_S S''}\) of multiplicative type, so to prove that
\(T''_1=T''_2\) it is enough to see that they have the same fibers
(Exposé \SGARefLink{sga3:ix:theorem:5:1:bis}{IX~5.1 bis}). Consider the commutative diagram
\[
\begin{tikzcd}[column sep=huge,row sep=large]
S' \arrow[d,"v"'] &
  S'\times_S S' \arrow[l,shift left=.55ex]
    \arrow[l,shift right=.55ex] \arrow[d,"v\times v"]\\
S'' &
  S''\times_S S'' . \arrow[l,shift left=.55ex]
    \arrow[l,shift right=.55ex]
\end{tikzcd}
\]
The morphism \(v\) is a finite epimorphism, hence dominant by
property \textup{\SGARefLink{sga3:xv:lemma:3:6:item:1}{(1)}} above, and closed; it is therefore surjective,
and so is \(v\times v\). Let \(x''\) be a point of
\(S''\times_S S''\), and let \(x'\) be a point of
\(S'\times_S S'\) lying above \(x''\). The commutativity of the
diagram shows that the two inverse images of
\((T''_1)_{x''}\) and \((T''_2)_{x''}\) in \(G_{x'}\) coincide.
Thus
\[
  (T''_1)_{x''}=(T''_2)_{x''}
\]
(EGA IV~2.2.15).

The preceding argument and an immediate induction on the number of
factors in a factorization of \(u:S'\to S\) into effective
epimorphisms, supplied by property \textup{\SGARefLink{sga3:xv:lemma:3:6:item:3}{(3)}} above, show that in
order to prove the exactness of \((*)\) for central subgroups of
multiplicative type it is enough to treat the case in which \(u\)
is an effective epimorphism. Finally, another application of
Exposé \SGARefLink{sga3:ix:theorem:5:1:bis}{IX~5.1 bis} shows that it is enough to prove \SGARefLink{sga3:xv:lemma:3-7}{3.7} when
\(\mathrm M\) is the functor of all subgroups of \(G\) of
multiplicative type and \(u\) is an effective epimorphism.

\smallskip
\noindent\textit{ii) Injectivity of
\(\mathrm M(S)\to\mathrm M(S')\).}\quad
Since \(u:S'\to S\) is an epimorphism, the canonical morphism
\[
  \mathcal O_S\longrightarrow u_*(\mathcal O_{S'})
\]
is injective. Moreover, an \(S\)-group of multiplicative type is
flat over \(S\). The injectivity of
\(\mathrm M(S)\to\mathrm M(S')\) is therefore a consequence of the
following more general lemma.

\SGAStatement{\SGAInlineRefTarget{sga3:xv:lemma:3-8}Lemma 3.8}{Let \(f:X\to S\) and \(g:S'\to S\) be
morphisms of preschemes, let \(\mathcal F\) be a quasi-coherent
\(\mathcal O_X\)-module, and put
\[
  X'=X\times_S S',
  \qquad
  \mathcal F'=\mathcal F\otimes_{\mathcal O_S}\mathcal O_{S'}.
\]
Let \(\mathrm Q(\mathcal F)\) be the set of quotient
\(\mathcal O_X\)-modules \(\mathcal G\) of \(\mathcal F\) that are
quasi-coherent and flat over \(S\), and define
\(\mathrm Q(\mathcal F')\) analogously relative to
\(\mathcal F',X'\), and \(S'\). Suppose that \(g\) is quasi-compact
and that
\[
  \mathcal O_S\longrightarrow g_*(\mathcal O_{S'})
\]
is injective. Then the canonical map
\[
  \begin{aligned}
  \mathrm Q(\mathcal F)&\longrightarrow\mathrm Q(\mathcal F'),\\
  \mathcal G&\longmapsto
    \mathcal G\otimes_{\mathcal O_S}\mathcal O_{S'}
  \end{aligned}
\]
is injective.}
\SGARefTarget{sga3:xv:lemma:3:8}

\SGAProof{Proof}
Indeed, we may suppose that \(S\) is affine, and then that \(X\) is
affine. Since \(S'\) is quasi-compact, it has a finite covering by
affine open subsets. Replacing \(S'\) by the disjoint union of these
open subsets, we may therefore suppose that \(S'\) is affine. We are
then reduced to the following lemma, whose proof is immediate.

\SGAStatement{Lemma 3.9}{Let \(A\to A'\) be an injective
homomorphism of rings, let \(M\) be an \(A\)-module, let \(N=M/P\)
be a quotient \(A\)-module flat over \(A\), and put
\[
  M'=M\otimes_A A',
  \qquad
  N'=N\otimes_A A'=M'/P'.
\]
Then \(P\) is the inverse image of \(P'\) under the canonical
homomorphism \(M\to M'\). Thus \(N\) and \(P\) are determined by
\(N'\) and \(P'\).}
\SGARefTarget{sga3:xv:lemma:3:9}

\smallskip
\noindent\textit{iii) Exactness of \SGARefLink{sga3:xv:lemma:3-7}{3.7} \((*)\) at
\(\mathrm M(S')\).}\quad
Let
\[
  T'\in\operatorname{Ker}\bigl(
    \mathrm M(S')\rightrightarrows
    \mathrm M(S'\times_S S')\bigr).
\]
Suppose that iii) has been proved when \(S\) is the spectrum of a
noetherian local ring. For every point \(s\) of \(S\), there then
exists a subgroup \(T_s\) of multiplicative type in
\(G\times_S\operatorname{Spec}(\mathcal O_{S,s})\), obtained by
descent from
\[
  T'\times_{S'}u^{-1}\!\bigl(\operatorname{Spec}(\mathcal O_{S,s})\bigr).
\]
By Exposé VI\textsubscript{B}, \S10, and \SGARefLink{sga3:xv:lemma:3:6}{3.6}, there is an open
neighborhood \(U\) of \(s\) in \(S\) and a subgroup scheme \(T\) of
\(G\) over \(U\), of multiplicative type, that extends \(T_s\).
Let \(U'\) be the inverse image of \(U\) in \(S'\). We then have two
subschemes of \(G'|_{U'}\), namely \(T'|_{U'}\) and
\(T\times_U U'\), which coincide over
\(u^{-1}(\operatorname{Spec}(\mathcal O_{S,s}))\).

Regarding the latter scheme as the projective limit of the
\(u^{-1}(V)\), where \(V\) runs through the open neighborhoods of
\(s\) in \(U\), EGA IV~8 shows that there is an open subset
\(V\subset U\), containing \(s\), over whose inverse image
\(T\times_U U'\) and \(T'\) coincide. Thus, under the stated
hypotheses, \(T'\) descends locally on \(S\); by the uniqueness
proved in ii), it consequently descends globally. In short, it is
enough to prove iii) when \(S\) is the spectrum of a noetherian
local ring.

Let \(\widehat S\) be the spectrum of the completion of the ring of
\(S\), and put
\[
  S''=S'\times_S S',
  \qquad
  \widehat S'=\widehat S\times_S S',
  \qquad
  \widehat S''=\widehat S\times_S S''
      =\widehat S'\times_S S'.
\]
I claim that it is enough to show that
\[
  \tag{\(\widehat *\)}\label{sga3:xv:lemma:3-9:equation:widehat-star}
  \mathrm M(\widehat S)\longrightarrow
  \mathrm M(\widehat S')
  \rightrightarrows
  \mathrm M(\widehat S'')
\]
is exact at \(\mathrm M(\widehat S')\). This follows from the
commutative diagram
\[
\begin{tikzcd}[column sep=large,row sep=large]
\mathrm M(S) \arrow[r] \arrow[d]
  & \mathrm M(S') \arrow[r,shift left=.55ex]
      \arrow[r,shift right=.55ex] \arrow[d]
  & \mathrm M(S'') \arrow[d]\\
\mathrm M(\widehat S) \arrow[r]
      \arrow[d,shift left=.55ex] \arrow[d,shift right=.55ex]
  & \mathrm M(\widehat S') \arrow[r,shift left=.55ex]
      \arrow[r,shift right=.55ex]
      \arrow[d,shift left=.55ex] \arrow[d,shift right=.55ex]
  & \mathrm M(\widehat S'')\\
\mathrm M(\widehat S\times_S\widehat S)
  \arrow[r,"f"]
  & \mathrm M(\widehat S'\times_{S'}\widehat S').
\end{tikzcd}
\]
Here the second row is exact at \(\mathrm M(\widehat S')\) by
hypothesis, the first two columns are exact by fpqc descent, and
\(f\) is injective by ii), applied to the finite epimorphism
\[
  \widehat S'\times_{S'}\widehat S'
  \longrightarrow
  \widehat S\times_S\widehat S
\]
deduced from \(S'\to S\) by the flat base change
\(\widehat S\times_S\widehat S\to S\). The diagram chase is left
to the reader. The characterization of finite effective
epimorphisms recalled in d)\SGARefLink{sga3:xv:theorem:5:2}{(2)} shows that
\(\widehat S'\to\widehat S\), obtained from \(S'\to S\) by the
flat base change \(\widehat S\to S\), is again a finite effective
epimorphism. We are therefore reduced to the case in which \(S\)
is, in addition, the spectrum of a complete noetherian local ring.

Let \(S_0\) be the reduced closed subscheme of \(S\) whose
underlying space is the closed point. Let \(T'\) be as above and
write \(T''\) for its common image in \(\mathrm M(S'')\). Since
\[
  S'_0=S'\times_S S_0
\]
is faithfully flat over \(S_0\), there is a subgroup \(T_0\) of
\(G_{S_0}\), of multiplicative type, whose inverse image in
\(G_{S'_0}\) is \(T'_{S'_0}\), by fpqc descent. Since \(S\) is
complete noetherian local, there exist an \(S\)-group \(T\) of
multiplicative type and an \(S_0\)-morphism
\[
  u_0:T\times_S S_0\longrightarrow T_0
\]
(Exposé \SGARefLink{sga3:x:corollary:3:3}{X~3.3}). The inverse image of \(u_0\) over \(S'_0\) extends
uniquely to an \(S'\)-isomorphism
\[
  u':T_{S'}\xrightarrow{\sim}T',
\]
again by Exposé \SGARefLink{sga3:x:corollary:3:3}{X~3.3}. Notice here that, because \(S'\) is finite
over the complete local scheme \(S\), it is a finite disjoint sum
of complete local schemes.

The two inverse images of \(u'\) over \(S''\) are morphisms from
\(T_{S''}\) to \(T''\) that coincide over \(S_0\times_S S''\);
hence they coincide everywhere by the same result. Since \(T\) is
flat over \(S\) and \(S'\to S\) is a finite effective
epimorphism, TDTE I, page~8, gives an exact diagram
\[
  \operatorname{Hom}_S(T,G)\longrightarrow
  \operatorname{Hom}_{S'}(T_{S'},G_{S'})
  \rightrightarrows
  \operatorname{Hom}_{S''}(T_{S''},G_{S''}).
\]
Thus \(u'\) comes from a morphism of preschemes \(u:T\to G\).
Likewise, \(T\times_S T\) is flat over \(S\), and therefore
\[
  \operatorname{Hom}_S(T\times_S T,G)\longrightarrow
  \operatorname{Hom}_{S'}(T_{S'}\times_{S'}T_{S'},G_{S'})
\]
is injective. It follows at once that \(u\) is a homomorphism of
groups, since \(u'\) is one.

Moreover, \(u\) is a monomorphism. First observe that
\(\operatorname{Ker}(u)\) is flat over \(S\). To prove this we may
suppose that \(S\) is artinian local (EGA
\(0_{\mathrm{III}}\), 10.2.2), hence that \(G\) is separated
(Exposé VI\textsubscript{A}, \SGARefLink{sga3:VIA:0.3}{0.3}); then
\(\operatorname{Ker}(u)\) is of multiplicative type (Exposé
\SGARefLink{sga3:ix:theorem:6:8}{IX~6.8}), and is therefore flat over \(S\). Since
\[
  \operatorname{Ker}(u)\times_S S'
    =\operatorname{Ker}(u')=0,
\]
\SGARefLink{sga3:xv:lemma:3-8}{Lemma~3.8} gives \(\operatorname{Ker}(u)=0\). Finally, a
monomorphism \(u\) is an immersion (\SGARefLink{sga3:viii:remark:7:13}{Exposé VIII, Remarks~7.13}),
and its image \(u(T)\) is an element of \(\mathrm M(S)\) whose
image in \(\mathrm M(S')\) is \(T'\). This proves \SGARefLink{sga3:xv:lemma:3-7}{Lemma~3.7}.

\smallskip
\noindent\textit{e) End of the proof of I).}\quad
By reduction c), we are reduced to the case in which \(S\) is the
spectrum of a complete reduced noetherian local ring \(A\). Let
\(S'\) be the spectrum of the normalization \(A'\) of \(A\).
Nagata's theorem says that \(A'\) is finite over \(A\) (EGA
\(0_{\mathrm{IV}}\), 23.1.5), and \(S'\to S\) is an epimorphism
because \(A\to A'\) is injective. Suppose that \(G_{S'}\) has a
torus \(T'\) whose underlying space is \(E_{S'}\). The two inverse
images of \(T'\) in \(G_{S'\times_S S'}\) are central sub-tori
with the same underlying space, so they coincide (Exposé IX~5
bis). By \SGARefLink{sga3:xv:lemma:3-7}{Lemma~3.7}, \(T'\) therefore comes from a central
sub-torus \(T\) of \(G_S\), which plainly has \(E\) as its
underlying space. It is thus enough to prove the existence of
\(T'\). This reduces us to the normal case and completes the proof
of I).

\medskip
\noindent\textit{II) Proof that \textup{\SGARefLink{sga3:xv:theorem:3-1:item:iii}{(iii)}} implies
\textup{\SGARefLink{sga3:xv:theorem:3-1:item:ii}{(ii)}} when \(S\) is normal.}\quad
We may restrict ourselves to the case in which \(S\) is integral;
let \(t\) be its generic point. For every integer \(n\) that is a
power of \(q\), let \({}_nG\) be the ``kernel'' in \(G\) of
raising to the \(n\)-th power. Since \(E\) is locally closed in
\(G\), by \SGARefLink{sga3:xv:theorem:3-1:item:iii}{(iii)}(b), the intersection of
\(\operatorname{ens}({}_nG)\) with \(E\) is locally closed in
\(\operatorname{ens}(G)\). Let \(E(n)\) denote the reduced
subscheme of \(G\) whose underlying space is
\(\operatorname{ens}({}_nG)\cap E\).

We show that the structural morphism
\[
  E(n)\longrightarrow S
\]
is separated and universally open. Both properties have a
valuative criterion (EGA II~7.2.3 and EGA IV~14.5.8). Let
\(S'\) therefore be an \(S\)-scheme that is the spectrum of a
complete discrete valuation ring with algebraically closed residue
field. We have proved that \SGARefLink{sga3:xv:theorem:3-1:item:iii}{3.1(iii)} implies \SGARefLink{sga3:xv:theorem:3-1:item:iv}{3.1(iv)}, so there is a
sub-torus \(T'\) of \(G_{S'}\) whose underlying space is
\(E_{S'}\). Now \({}_nT'\) is finite étale over \(S'\), hence is
separated and universally open over \(S'\); the same is therefore
true of \(E(n)\times_S S'\), which has the same underlying space.

The fibers of \(E(n)\), moreover, all have the same number of
geometric points, namely \(n^r\) if \(r\) is the rank of the torus
\(T_t\).\footnote{The French edition prints \(r^n\); the intended
cardinality of the \(n\)-torsion of a rank-\(r\) torus is \(n^r\).}
Finally, the generic fiber \(E(n)_t\), being reduced, is equal to
\({}_nT_t\), and hence is étale over \(\operatorname{Spec}(t)\).
Since \(S\) is normal, SGA I~10.11 now shows that \(E(n)\) is an
étale covering of \(S\).

If \(s\) is a point of \(S\), then \(E(n)_s\) is étale over
\(\operatorname{Spec}\kappa(s)\), hence reduced, and consequently
coincides with the group \({}_nT_s\) of multiplicative type. We
show that \(E(n)\) is a group subscheme of \(G\). Let
\[
  m:E(n)\times_S E(n)\longrightarrow G
\]
be the morphism induced by multiplication in \(G\). The underlying
map \(\operatorname{ens}(m)\) factors through \(E(n)\), so \(m\)
itself factors through the scheme \(E(n)\), because its source
\(E(n)\times_S E(n)\) is étale over \(S\), and hence reduced. It
then follows from Exposé \SGARefLink{sga3:x:corollary:4-8:item:a}{X~4.8(a)} that \(E(n)\) is a subgroup of
multiplicative type. As observed already in \SGARefLink{sga3:xv:remark:3-2:item:a}{3.2(a)}, the family of
subgroups \(E(n)\) is necessarily coherent. This proves that
\SGARefLink{sga3:xv:theorem:3-1:item:iii}{(iii)} implies \SGARefLink{sga3:xv:theorem:3-1:item:ii}{(ii)} when \(S\) is normal.

\medskip
\noindent\textit{III) Proof that \textup{\SGARefLink{sga3:xv:theorem:3-1:item:ii}{(ii)}} implies
\textup{\SGARefLink{sga3:xv:theorem:3-1:item:i}{(i)}}.}\quad
In fact, we shall show that there is a unique sub-torus \(T\) of
\(G\), with underlying space \(E\), such that
\[
  {}_nT=M_n
\]
for every \(n\) that is a power of \(q\). Uniqueness follows
immediately from Exposé \SGARefLink{sga3:ix:corollary:4:8:item:b}{IX~4.8(b)}. In proving existence, in view of
this uniqueness, we may successively suppose that:
\begin{enumerate}
\item[\normalfont a)] \(S\) is noetherian.
\item[\normalfont b)] The \(M_n\) are central subgroups. Indeed, it
is enough to replace \(G\) by
\[
  Z_n=\operatorname{Centr}_G(M_n)
\]
for sufficiently large \(n\), by \SGARefLink{sga3:xv:lemma:2:5}{2.5} and \SGARefLink{sga3:xv:lemma:2-5-bis}{2.5 bis}.
\item[\normalfont c)] \(S\) is the spectrum of a local ring.
Suppose that the problem has been solved after every base change
\(\operatorname{Spec}(\mathcal O_{S,s})\to S\), for \(s\in S\),
and let \(T_s\) be the resulting sub-torus of
\(G\times_S\operatorname{Spec}(\mathcal O_{S,s})\). There is then
an open neighborhood \(U\) of \(s\) and a sub-torus \(T\) of
\(G|_U\) extending \(T_s\), by Exposé VI\textsubscript{B}, \S10,
and \SGARefLink{sga3:xv:lemma:3:6}{3.6}. We have proved that \SGARefLink{sga3:xv:theorem:3-1:item:ii}{3.1(ii)} implies \SGARefLink{sga3:xv:theorem:3-1:item:iv}{3.1(iv)}; since \(S\)
is noetherian, \(E\) is therefore constructible by \SGARefLink{sga3:xv:lemma:3:3}{3.3}. After
shrinking \(U\), we may consequently suppose that
\[
  \operatorname{ens}(T)=E\times_S U
\]
(EGA IV~9.5.2). For every \(n\) that is a power of \(q\), the
groups \({}_nT\) and \(M_n\times_S U\) are then two subgroups of
multiplicative type in \(G|_U\), with the same fibers; because
\(M_n\) is central, they coincide by Exposé \SGARefLink{sga3:ix:corollary:5:3:bis}{IX~5.3 bis}. Thus the
problem has a local solution on \(S\), and uniqueness makes these
local solutions into a global one.
\item[\normalfont d)] \(S\) is the spectrum of a complete
noetherian local ring and, if \(s\) is the closed point, \(T_s\)
is trivial. This follows from EGA \(0_{\mathrm{III}}\), 10.3.1,
and fpqc descent.
\item[\normalfont e)] \(S\) is reduced. Apply \SGARefLink{sga3:xv:corollary:2-2}{2.2}.
\item[\normalfont f)] \(S\) is normal. Apply Nagata's finiteness
theorem (EGA \(0_{\mathrm{IV}}\), 23.1.5) and Lemma~3.7 for
central sub-tori.
\end{enumerate}

After these reductions, \(T_s\) is trivial, and hence
\((M_n)_s\) is trivial. It follows from Exposé \SGARefLink{sga3:x:corollary:3:3}{X~3.3} that \(M_n\)
itself is trivial. If \(t\) is any point of \(S\), all the
subgroups \({}_nT_t\) of \(T_t\), for \(n\) a power of \(q\), are
therefore trivial; Exposé \SGARefLink{sga3:x:proposition:1:4}{X~1.4} then shows at once that \(T_t\)
itself is trivial. It remains only to prove the following lemma.

\SGAStatement{Lemma 3.10}{Under the hypotheses of \SGARefLink{sga3:xv:theorem:3-1:item:ii}{3.1(ii)}, suppose
in addition that \(S\) is noetherian and normal and that, for every
point \(s\) of \(S\), \(T_s\) is a trivial torus. Then there exists
a unique sub-torus \(T\) of \(G\), with underlying space \(E\),
such that
\[
  {}_nT=M_n
\]
for every \(n\) that is a power of \(q\). Moreover, \(T\) is
trivial.}
\SGARefTarget{sga3:xv:lemma:3:10}

\SGAProof{Proof}
Uniqueness follows from the fact that \(T\), being smooth over
\(S\), is reduced. To prove existence, we may suppose that \(S\)
is irreducible, with generic point \(\eta\). Let \(r\) be the rank
of \(T_\eta\), put
\[
  T'=\mathbb G_{m,S}^{\,r},
\]
and choose an isomorphism
\[
  u_\eta:T'_\eta\xrightarrow{\sim}T_\eta.
\]
Write \({}_nu_\eta\) for its restriction to \({}_nT'_\eta\).
Since \({}_nT'\) and \(M_n\) are trivial, \({}_nu_\eta\) has a
unique extension to an \(S\)-isomorphism
\[
  {}_nu:{}_nT'\xrightarrow{\sim}M_n.
\]
I claim that, for every point \(s\) of \(S\), there is a
necessarily unique group isomorphism
\[
  u_s:T'_s\xrightarrow{\sim}T_s
\]
that extends \(({}_nu)_s\) for every \(n\) that is a power of
\(q\).

Indeed, let \(S^1\) be an \(S\)-scheme that is the spectrum of a
complete discrete valuation ring with algebraically closed residue
field, whose generic point \(t_1\) maps to \(\eta\) and whose
closed point \(s_1\) maps to \(s\) (EGA II~7.1.9). The proof that
3.1(ii) implies 3.1(iv) gives a sub-torus \(T^1\) of \(G_{S^1}\)
such that
\[
  (M_n)_{S^1}={}_nT^1
\]
for every \(n\) that is a power of \(q\). Since \(S^1\) is normal,
\(T^1\) is isotrivial by Exposé \SGARefLink{sga3:x:theorem:5-16}{X~5.16}. The classification of
isotrivial tori (Exposé \SGARefLink{sga3:x:corollary:1:2}{X~1.2}) and SGA I, Exposé V~8.2, together
with the fact that the generic fiber of \(T^1\) is trivial, show
that \(T^1\) is trivial. We may therefore extend
\[
  u_\eta\times_\eta t_1
\]
to an \(S^1\)-isomorphism of group schemes
\[
  u^1:T'\times_S S^1\xrightarrow{\sim}T^1.
\]
On \({}_nT'\times_S S^1\), the morphisms \(u^1\) and
\({}_nu\times_S S^1\) agree on the generic fiber, and hence agree
everywhere. Restricting \(u^1\) to the closed fiber gives the
required \(u_s\) after the field extension
\(\kappa(s)\to\kappa(s_1)\); it descends to \(\kappa(s)\) by
Exposé \SGARefLink{sga3:ix:corollary:4:8:item:a}{IX~4.8(a)} and fpqc descent.

The scheme \(T'\) is smooth over the normal scheme \(S\), hence is
normal. The extension criterion for rational maps in EGA
IV\textsubscript{4}, 20.4.6, now shows that there is an
\(S\)-morphism
\[
  u:T'\longrightarrow G
\]
whose restriction to each fiber \(T'_s\) is the composite
\[
  T'_s\xrightarrow{u_s}T_s\longrightarrow G_s.
\]
Indeed, for every base change \(S^1\to S\) of the kind used above,
the preceding construction gives precisely the required extension
over \(S^1\).

It remains to show that \(u\) is a homomorphism. Let \(m_{T'}\)
and \(m_G\) denote multiplication in \(T'\) and \(G\),
respectively. The coincidence subscheme of
\[
  m_G\circ(u\times_Su)
  \quad\text{and}\quad
  u\circ m_{T'}
\]
contains every fiber, since every \(u_s\) is a homomorphism. It
therefore has the same underlying space as \(T'\times_S T'\), and
so equals \(T'\times_S T'\), because the latter is smooth over
\(S\) and hence reduced. Finally, \(u\) is a monomorphism because
it is one on every fiber; it is therefore an immersion (Exposé
\SGARefLink{sga3:viii:proposition:7:9}{VIII~7.9}). The image of \(T'\) under \(u\) is a sub-torus of \(G\)
whose underlying set is \(E\).

This completes the proof of \SGARefLink{sga3:xv:theorem:3-1}{Theorem~3.1}.

% END INCLUDED SGA3_English_current_buildable_source_20260823/source/inputs/XV/tex/components/03_expose_XV_section3.tex

% BEGIN INCLUDED SGA3_English_current_buildable_source_20260823/source/inputs/XV/tex/components/04_expose_XV_section4.tex
% Authority: Expo15.pdf, local pp. 26--31 / combined pp. 896--901 /
% printed pp. 398--408.
% Stops immediately before §5.

\SGASection{4}{Characterization of a sub-torus \(T\) by the subgroups
\({}_nT\)}
\SGARefTarget{sga3:xv:section:4}

\SGASubsection{4.1}{Statement of the main theorem}
\SGARefTarget{sga3:xv:subsection:4:1}

\SGAStatement{\SGAInlineRefTarget{sga3:xv:theorem:4-1}Theorem 4.1}{Let \(S\) be a connected locally
noetherian prescheme, let \(G\) be an \(S\)-group prescheme of finite
type over \(S\), let \(q>1\) be an integer invertible on \(S\), and
let \(r\) be a positive integer. For every integer \(n\) that is a
power of \(q\), let \(M(n)\) be a subgroup scheme of \(G\), of
multiplicative type and of type
\((\mathbb Z/n\mathbb Z)^r\). Suppose that:
\begin{enumerate}
\renewcommand{\labelenumi}{\textup{\alph{enumi})}}
\item the family \(M(n)\) is coherent: if \(m\mid n\), then
\[
  {}_mM(n)=M(m);
\]
\item\SGAInlineRefTarget{sga3:xv:theorem:4-1:item:b} there exist a point \(s\in S\) and a sub-torus \(T_s\) of
\(G_s\) such that
\[
  M(n)_s={}_nT_s\qquad\text{for every }n;
\]
\item for every point \(t\in S\), there exists a closed affine
subscheme \(F_t\) of \(G_t\) containing \(M(n)_t\) for every \(n\).
\end{enumerate}
Then there exists a unique sub-torus \(T\) of \(G\) such that
\({}_nT=M(n)\) for every \(n\) that is a power of \(q\).}
\SGARefTarget{sga3:xv:theorem:4:1}

There is an analogous theorem for lifting morphisms.

\SGAStatement{\SGAInlineRefTarget{sga3:xv:theorem:4-1-bis}Theorem 4.1 bis}{Let \(S,G,q\) be as above, let \(T\)
be an \(S\)-torus, and, for every \(n\) that is a power of \(q\), let
\(u(n):{}_nT\to G\) be an \(S\)-homomorphism. Suppose that:
\begin{enumerate}
\renewcommand{\labelenumi}{\textup{\alph{enumi})}}
\item the family \(u(n)\) is coherent: if \(m\mid n\), then
\[
  u(m)=u(n)|_{{}_mT};
\]
\item there exist a point \(s\in S\) and a homomorphism
\(u_s:T_s\to G_s\) such that
\[
  u_s|_{{}_nT_s}=u(n)_s
  \qquad\text{for every }n\text{ that is a power of }q;
\]
\item for every point \(t\in S\), there exists a closed affine
subscheme \(F_t\) of \(G_t\) containing
\(u(n)_t({}_nT_t)\) for every \(n\).
\end{enumerate}
Then there exists a unique homomorphism \(u:T\to G\) whose restriction
to \({}_nT\) equals \(u(n)\) for every \(n\) that is a power of \(q\).}
\SGARefTarget{sga3:xv:theorem:4:1:bis}

\SGAStatement{Remark 4.2}{Using the lower semicontinuity of the
abelian rank of a flat group prescheme of finite type over the
spectrum of a discrete valuation ring (cf. Exposé \SGARefLink{sga3:x:remark:8-7}{X~8.7}), condition
(c) in \SGARefLink{sga3:xv:theorem:4-1}{Theorems 4.1} and \SGARefLink{sga3:xv:theorem:4-1-bis}{4.1 bis} may be weakened by requiring the
closed affine subscheme \(F_t\) only for every maximal point \(t\) of
\(S\).}
\SGARefTarget{sga3:xv:remark:4:2}

Let us deduce \SGARefLink{sga3:xv:theorem:4-1-bis}{Theorem 4.1 bis} from \SGARefLink{sga3:xv:theorem:4-1}{Theorem 4.1}. Put
\(G'=G\times_S T\). For every \(n\) that is a power of \(q\), consider
the homomorphism
\[
  v(n):{}_nT\longrightarrow G'
\]
whose projections to \(G\) and \(T\) are respectively \(u(n)\) and
the canonical immersion \({}_nT\to T\). Thus \(v(n)\) is an immersion;
let \(M(n)\) be its image subgroup. The family \(M(n)\) is coherent
in the sense of \SGARefLink{sga3:xv:theorem:4-1}{Theorem 4.1}. The group \(M(n)_s\) equals
\({}_nT'_s\), where \(T'_s\) is the sub-torus of \(G'_s\) that is the
graph of \(u_s\), and, for every \(t\in S\), the closed affine
subscheme \(F_t\times_tT_t\) of \(G'_t\) contains all the groups
\(M(n)_t\). By \SGARefLink{sga3:xv:theorem:4-1}{Theorem 4.1} there is therefore a sub-torus \(T'\) of
\(G'\) such that \({}_nT'=M(n)\) for every \(n\).

Let \(f:T'\to T\) be the restriction of the projection \(G'\to T\),
and let \(f(n)\) be its restriction to \(M(n)\). Then
\[
  f(n)\circ v(n)=\operatorname{id}_{{}_nT}.
\]
The fibre \(T'_s\) is the graph of \(u_s\) (Exposé \SGARefLink{sga3:ix:corollary:4:8:item:b}{IX~4.8(b)}), so
\(f_s\) is an isomorphism. The groups \(\ker f\) and
\(\operatorname{coker}f\) are of multiplicative type (Exposé \SGARefLink{sga3:ix:proposition:2:7}{IX~2.7})
and of constant type, since \(S\) is connected. They are therefore
trivial, and \(f\) is an isomorphism. If \(v=f^{-1}\), then
\[
  v|_{{}_nT}=v(n).
\]
Composing \(v\) with the projection \(G'\to G\) gives the required
homomorphism \(u:T\to G\). This proves existence; uniqueness follows
from Exposé \SGARefLink{sga3:ix:corollary:4:8:item:a}{IX~4.8(a)}.

\SGASubsection{4.2}{Application}
\SGARefTarget{sga3:xv:subsection:4:2}

We shall generalize Theorem \SGARefLink{sga3:ix:theorem:7:1}{IX~7.1}. Let \(A\) be a complete
noetherian local ring, let \(I\) be its maximal ideal, and put
\[
  S=\operatorname{Spec}A,\qquad
  S_m=\operatorname{Spec}(A/I^m).
\]
For every prescheme \(X\), put \(X_m=X\times_SS_m\).

Let \(G\) be an \(S\)-group prescheme of finite type, let \(T\) be an
\(S\)-torus, let \(q\) be an integer invertible on \(S\), and let
\[
  u_m:T_m\longrightarrow G_m\qquad(m\in\mathbb N)
\]
be a coherent family of homomorphisms. For \(n\) ranging over the
powers of \(q\), denote by \(u_m(n)\) the restriction of \(u_m\) to
\({}_nT_m\), and by
\[
  u(n):{}_nT\longrightarrow G
\]
the unique homomorphism extending the morphisms \(u_m(n)\) for all
\(m\) (\SGARefLink{sga3:xv:proposition:1-6:item:a}{Proposition 1.6(a)}). We call the family \((u_m)_{m\in\mathbb N}\)
\emph{admissible} if, for every point \(t\in S\), there exists a
closed affine subprescheme \(F_t\) of \(G_t\) containing
\(u(n)_t({}_nT_t)\) for every \(n\) that is a power of \(q\). As will
be seen, this property does not depend on \(q\).

\SGAStatement{Proposition 4.3}{With the preceding notation, the
canonical map
\[
 \operatorname{Hom}_{S\text{-}\mathrm{gr}}(T,G)
 \longrightarrow
 \varprojlim_m\operatorname{Hom}_{S_m\text{-}\mathrm{gr}}(T_m,G_m)
\]
identifies its source with the subset of the target consisting of
admissible coherent families.}
\SGARefTarget{sga3:xv:proposition:4:3}

Indeed, apply \SGARefLink{sga3:xv:theorem:4-1-bis}{Theorem 4.1 bis} with \(s\) the closed point of \(S\).

\SGAStatement{Corollary 4.4}{With the preceding notation, suppose in
addition that \(G\) has affine fibres. Then the canonical map
\[
 \operatorname{Hom}_{S\text{-}\mathrm{gr}}(T,G)
 \longrightarrow
 \varprojlim_m\operatorname{Hom}_{S_m\text{-}\mathrm{gr}}(T_m,G_m)
\]
is an isomorphism.}
\SGARefTarget{sga3:xv:corollary:4:4}

Indeed, when \(G\) has affine fibres every coherent family is
admissible.

\SGAStatement{Remark 4.5}{When \(G\) is separated, the ring \(A\) in
Proposition 4.3 may be replaced by a noetherian \(I\)-adic ring: one
uses EGA III~5.4.1 in place of Proposition 1.6(a).}
\SGARefTarget{sga3:xv:remark:4:5}

\SGASubsection{4.3}{Proof of \SGARefLink{sga3:xv:theorem:4-1}{Theorem 4.1}}
\SGARefTarget{sga3:xv:subsection:4:3}

\SGAStatement{Lemma 4.6}{Let \(k\) be a field, let \(G\) be a
\(k\)-algebraic group, let \(q\) be an integer prime to the
characteristic of \(k\), let \(r>0\), and let \(M(n)\), for \(n\)
ranging over the powers of \(q\), be a coherent family of subgroups
of \(G\), of multiplicative type and of type
\((\mathbb Z/n\mathbb Z)^r\). There exists a smallest closed
subscheme \(H\) of \(G\) containing all the \(M(n)\). Moreover, \(H\)
is a smooth, connected, commutative algebraic subgroup of \(G\), and
its formation commutes with extension of the base field.}
\SGARefTarget{sga3:xv:lemma:4:6}

Let \(M\subseteq\operatorname{ens}(G)\) be the union of the underlying
sets of all the \(M(n)\), and let \(H\) be the reduced closed
subscheme of \(G\) whose underlying space is the closure of \(M\).
Each \(M(n)\) is étale over \(k\), hence reduced, and is therefore
contained in \(H\). Thus \(H\) is the smallest closed subscheme
containing all the \(M(n)\).

Let \(\bar k\) be an algebraically closed extension of \(k\). By
construction the \(M(n)\) are schematically dense in \(H\) (Exposé
\SGARefLink{sga3:ix:definition:4:1}{IX~4.1}), hence the \(M(n)_{\bar k}\) are schematically dense in
\(H_{\bar k}\) (Exposé \SGARefLink{sga3:ix:lemma:4:5}{IX~4.5}). Since every \(M(n)_{\bar k}\) is
reduced, \(H_{\bar k}\) is the reduced closed subscheme of
\(G_{\bar k}\) whose underlying space is the closure of \(M_{\bar k}\).
Thus \(H\) is geometrically reduced (EGA IV~4.6.1).

Coherence of the family implies that \(M\) is stable under the group
law. The subset \(M\times_kM\) is dense in
\(\operatorname{ens}(H\times_kH)\), and \(H\times_kH\) is reduced
(EGA IV~4.6.5); it follows that \(H\) is an algebraic subgroup of
\(G\). It is smooth over \(k\), since it is geometrically reduced
(Exposé VI\textsubscript{A}~\SGARefLink{prop:sga3-VIA-1.3.1}{1.3.1}), and it is commutative because all
the \(M(n)\) are commutative.

It remains to prove connectedness. Let \(H^0\) be the identity
component of \(H\), let \(m\) be the number of geometric points of
\(H/H^0\), and let \(q^s\) be the \(q\)-primary part of \(m\). For
every \(n\) that is a power of \(q\), the subgroup
\({}_{q^s}M(n)\) is contained in \(H^0\). Coherence and the prescribed
type of \(M(n)\) give
\[
  {}_{q^s}M(nq^s)=M(n).
\]
Hence \(H^0\) contains every \(M(n)\), and therefore \(H^0=H\).

We can now sharpen condition (c) of \SGARefLink{sga3:xv:theorem:4-1}{Theorem 4.1}. By \SGARefLink{sga3:xv:lemma:4:6}{Lemma 4.6} there
is, for each \(t\in S\), a smallest closed subscheme \(H_t\) of
\(G_t\) containing all the \(M(n)_t\). Its formation commutes with
extension of the residue field. Thus \(H_t\otimes_t\bar t\) is
contained in the affine closed subset \(F_t\), and is affine; hence
\(H_t\) itself is affine. \SGARefLink{sga3:xv:lemma:4:6}{Lemma 4.6} also says that \(H_t\) is a
smooth, connected, commutative algebraic group. The structure theorem
for such affine groups over an algebraically closed field
(\emph{Bible}, 4, \SGARefLink{sga3:xv:section:4}{Theorem 4}) expresses \(H_t\) as a direct product of
a torus \(T_t\) and a unipotent group \(U_t\). Since every \(M(n)_t\)
is contained in \(T_t\), minimality gives \(H_t=T_t\). Consequently
\(H_t\) is a torus.

We begin the proof of \SGARefLink{sga3:xv:theorem:4-1}{Theorem 4.1}.

\medskip\noindent
\textbf{a) Uniqueness.}
This is Exposé \SGARefLink{sga3:ix:corollary:4:8:item:b}{IX~4.8(b)}.

\medskip\noindent
\textbf{b) The case in which \(S\) is the spectrum of a complete
discrete valuation ring \(A\) with algebraically closed residue field
\(k\).}
Let \(K\) be the fraction field of \(A\), let \(s\) and \(t\) be the
closed and generic points of \(S\), let \(J\) be the maximal ideal of
\(A\), and put
\[
  S_m=\operatorname{Spec}(A/J^m),\qquad X_m=X\times_SS_m.
\]

\smallskip\noindent
\textbf{First case: the point in \SGARefLink{sga3:xv:theorem:4-1:item:b}{Theorem 4.1(b)} is the generic point
\(t\).}
Let \(T'\) be the schematic closure of \(T_t\) in \(G\). Its closed
fibre \(T'_s\) is an algebraic subgroup of \(G_s\), of dimension
\(r\), containing every \(M(n)_s\), hence containing \(H_s\). The
torus \(H_s\) contains every \(M(n)_s\), so it has rank at least
\(r\). Consequently \(H_s\) has the same underlying space as the
identity component of \(T'_s\).

The identity component \((T')^0\) of \(T'\) is then a subgroup
prescheme of \(G\), flat and separated (Exposé
VI\textsubscript{B}~\SGARefLink{sga3:xv:theorem:5-2}{5.2}), whose generic fibre \(T_t\) is a torus and
whose reduced closed fibre \(H_s\) is a torus. Hence \((T')^0\) is a
torus (Exposé \SGARefLink{sga3:x:theorem:8-8}{X~8.8}); denote it by \(T\). The groups \({}_nT\) and
\(M(n)\) are smooth over \(S\), hence reduced, and have the same
underlying space. They therefore coincide, and \(T\) is the required
solution.

\smallskip\noindent
\textbf{Second case: the point in \SGARefLink{sga3:xv:theorem:4-1:item:b}{Theorem 4.1(b)} is the closed point
\(s\).}
After replacing \(G\) by the schematic closure in \(G\) of the
smallest algebraic subgroup \(H_t\) containing the family \(M(n)_t\)
(\SGARefLink{sga3:xv:lemma:4:6}{Lemma 4.6}), we may suppose that \(G_t\) is affine.

For every \(m\geq0\), \SGARefLink{sga3:xv:subsection:2:2}{Proposition 2.2} gives a unique sub-torus \(T_m\)
of \(G_m\) lifting \(T_s\) and satisfying
\[
  {}_nT_m=M(n)_m
  \qquad\text{for every }n\text{ that is a power of }q.
\]
Put \(T'=\mathbb G_{m,S}^{\,r}\). Since \(k\) is algebraically
closed, \(T_s\) is split, and there is a \(k\)-isomorphism
\[
  u_s:T'_s\xrightarrow{\sim}T_s.
\]
It lifts uniquely to an \(S_m\)-isomorphism
\(u_m:T'_m\to T_m\) (Exposé \SGARefLink{sga3:ix:corollary:3:3}{IX~3.3}). The coherent family
\((u_m)_{m\in\mathbb N}\) defines a morphism of formal completions
\(\widehat u:\widehat{T'}\to\widehat G\), fitting into
\[
\begin{tikzcd}[column sep=large,row sep=large]
\widehat{T'} \arrow[r,"\widehat u"] \arrow[d,"i"'] &
  \widehat G \arrow[d,"j"'] \\
T' & G ,
\end{tikzcd}
\]
We shall prove that \(\widehat u\) is algebraizable, reducing first
to the case in which \(G\) is affine.

\SGAStatement{Lemma 4.7}{Let \(S=\operatorname{Spec}A\), where \(A\)
is a discrete valuation ring, and let \(X,Y\) be \(S\)-preschemes
that are quasi-compact, quasi-separated, and flat over \(S\). Then the
canonical map
\[
 \Gamma(X,\mathcal O_X)\otimes_A\Gamma(Y,\mathcal O_Y)
 \longrightarrow
 \Gamma(X\times_SY,\mathcal O_{X\times_SY})
\]
is an isomorphism.}
\SGARefTarget{sga3:xv:lemma:4:7}

Let \(f:X\to S\) and \(g:Y\to S\) be the structural morphisms.
Because \(X\) and \(Y\) are quasi-compact and quasi-separated,
EGA I~9.2.2 and EGA IV~1.7.4 show that
\(f_*\mathcal O_X\) and \(g_*\mathcal O_Y\) are quasi-coherent
\(\mathcal O_S\)-algebras. Let \(\widetilde X\) and \(\widetilde Y\)
be the corresponding affine \(S\)-schemes. Since \(Y\) is flat over
\(S\), EGA III~1.4.15 and EGA IV~1.7.21 give an isomorphism
\[
 \Gamma(\widetilde X\times_SY,
        \mathcal O_{\widetilde X\times_SY})
 \xrightarrow{\sim}
 \Gamma(X\times_SY,\mathcal O_{X\times_SY}).
\]
The scheme \(\widetilde X\) is flat over \(S\), because over \(A\)
flatness is equivalent to torsion-freeness. Applying EGA III~1.4.15
again, with \(X\) and \(Y\) interchanged, yields
\[
\begin{aligned}
 \Gamma(\widetilde X,\mathcal O_{\widetilde X})
 \otimes_A\Gamma(\widetilde Y,\mathcal O_{\widetilde Y})
 &\xrightarrow{\sim}
 \Gamma(\widetilde X\times_S\widetilde Y,
        \mathcal O_{\widetilde X\times_S\widetilde Y})\\
 &\xrightarrow{\sim}
 \Gamma(\widetilde X\times_SY,
        \mathcal O_{\widetilde X\times_SY}),
\end{aligned}
\]
which proves the lemma.

The \(S\)-group \(G\) is flat and of finite type over \(S\), hence
quasi-compact and quasi-separated. Applying the lemma to
\(G\times_SG\) gives
\[
 \Gamma(G,\mathcal O_G)\otimes_A\Gamma(G,\mathcal O_G)
 \xrightarrow{\sim}
 \Gamma(G\times_SG,\mathcal O_{G\times_SG}).
\]
Thus the multiplication \(m_G:G\times_SG\to G\) determines a
homomorphism
\[
 \Gamma(G,\mathcal O_G)\longrightarrow
 \Gamma(G,\mathcal O_G)\otimes_A\Gamma(G,\mathcal O_G),
\]
and hence an \(S\)-morphism
\[
  m_{\widetilde G}:\widetilde G\times_S\widetilde G
  \longrightarrow\widetilde G,
\]
where \(\widetilde G\) is the affine \(S\)-scheme corresponding to
\(\Gamma(G,\mathcal O_G)\). The usual identities show that
\(m_{\widetilde G}\) makes \(\widetilde G\) an \(S\)-group scheme and
that the canonical morphism \(v:G\to\widetilde G\) is a homomorphism.

\SGAStatement{Remark 4.8}{The group \(\widetilde G\) plays the role of
a largest affine quotient of \(G\). In fact it is an fpqc quotient of
\(G\), and is therefore of finite type over \(S\) (XVII, Appendix
III, 2).}
\SGARefTarget{sga3:xv:remark:4:8}

In our situation the generic fibre \(G_t\) is affine, so
\(v_t:G_t\to\widetilde G_t\) is an isomorphism (EGA I~9.3.3). Since
\(\widetilde G\) is affine, the coherent family
\[
  w_m=v_m\circ u_m:T'_m\longrightarrow\widetilde G_m
  \qquad(m\in\mathbb N)
\]
comes from a unique \(S\)-homomorphism (Exposé \SGARefLink{sga3:ix:theorem:7:1}{IX~7.1})
\[
  w:T'\longrightarrow\widetilde G.
\]
Let \(T_t\) be the sub-torus
\[
  T_t=v_t^{-1}\bigl(w_t(T'_t)\bigr)
\]
of \(G_t\). Its rank is at most \(r\). We show that it contains every
\(M(n)_t\). Let
\[
 u(n)_m:({}_nT')_m\xrightarrow{\sim}M(n)_m
\]
be the restriction of \(u_m\). This coherent family comes from a
unique \(S\)-isomorphism
\[
  u(n):{}_nT'\xrightarrow{\sim}M(n),
\]
because \(M(n)\) is finite over \(S\). For every \(m\geq0\),
\[
 w_m|_{({}_nT')_m}
  =(v_m\circ u_m)|_{({}_nT')_m}
  =(v\circ u(n))_m.
\]
It follows from \SGARefLink{sga3:xv:proposition:1-6:item:a}{Proposition 1.6(a)} that
\[
  w|_{{}_nT'}=v\circ u(n).
\]
On the generic fibre,
\[
 v_t^{-1}\!\left(w_t(({}_nT')_t)\right)
 =u(n)_t(({}_nT')_t)=M(n)_t.
\]
Thus \(T_t\) contains every \(M(n)_t\), and consequently has rank
\(r\).

We conclude as in the first case. Let \(T''\) be the schematic closure
of \(T_t\) in \(G\). Since \(T_t\) contains every \(M(n)_t\), the
closed fibre \(T''_s\) contains every \(M(n)_s\), and hence contains
\(T_s\) by the density theorem. On the other hand \(T''\) is flat
over \(S\) and \(T_t\) has dimension \(r\), so \(T''_s\) has
dimension \(r\) (Exposé VI\textsubscript{B}~4). Therefore \(T_s\)
has the same underlying space as the identity component of \(T''_s\),
and the identity component of \(T''\) is the required sub-torus of
\(G\).

\medskip\noindent
\textbf{c) End of the proof.}
To prove existence of \(T\), we may assume that \(S\) is reduced
(\SGARefLink{sga3:xv:subsection:2:2}{Proposition 2.2}). By \SGARefLink{sga3:xv:theorem:3-1}{Theorem 3.1}, \SGARefLink{sga3:xv:theorem:3-1:item:ii}{(ii)}\(\Rightarrow\)\SGARefLink{sga3:xv:theorem:3-1:item:i}{(i)}, it is then
enough to show that the set \(U\) of points \(x\in S\) for which there
exists a sub-torus \(T_x\) of \(G_x\) satisfying
\[
  {}_nT_x=M(n)_x
  \qquad\text{for every }n\text{ that is a power of }q
\]
equals \(\operatorname{ens}(S)\). Such a torus is necessarily unique,
and fpqc descent reduces its existence to existence after an extension
of the residue field of \(x\) (Exposé \SGARefLink{sga3:ix:corollary:4:8:item:b}{IX~4.8(b)}).

Since \(S\) is locally noetherian and connected and \(U\) is nonempty
(it contains the point \(s\) of \SGARefLink{sga3:xv:theorem:4-1:item:b}{Theorem 4.1(b)}), it is enough to prove
the following: if \(x,x'\in S\), if \(x\) is a specialization of
\(x'\), and if one of these two points lies in \(U\), then both lie
in \(U\). The usual method (EGA II~7.1.9) reduces this to the case
where \(S\) is the spectrum of a complete discrete valuation ring
with algebraically closed residue field, which was treated in part
(b). This completes the proof of \SGARefLink{sga3:xv:theorem:4-1}{Theorem 4.1}.

% END INCLUDED SGA3_English_current_buildable_source_20260823/source/inputs/XV/tex/components/04_expose_XV_section4.tex

% BEGIN INCLUDED SGA3_English_current_buildable_source_20260823/source/inputs/XV/tex/components/05_expose_XV_section5.tex
% Authority: Expo15.pdf, local pp. 32--38 / combined pp. 902--908 /
% printed pp. 409--421.
% Stops immediately before §6.

\SGASection{5}{Representability of the functor of smooth subgroups
equal to their connected normalizer}
\SGARefTarget{sga3:xv:section:5}

\SGAStatement{\SGAInlineRefTarget{sga3:xv:proposition:5-1}Proposition 5.1}{Let \(S\) be a prescheme, let \(G\)
be an \(S\)-group prescheme of finite presentation, and let \(H\) be
a subgroup prescheme of \(G\), smooth with connected fibres. Then:
\begin{enumerate}
\renewcommand{\labelenumi}{\textup{\alph{enumi})}}
\item the normalizer \(N\) of \(H\) in \(G\) is represented by a
closed subprescheme of \(G\), of finite presentation over \(S\);
\item\SGAInlineRefTarget{sga3:xv:proposition:5-1:item:b} the following conditions are equivalent:
\begin{enumerate}
\renewcommand{\labelenumii}{\textup{\roman{enumii})}}
\item the canonical immersion \(H\to N\) is an open immersion;
\item \(N\) is smooth along the identity section, and its identity
component, which is then representable (Exposé
VI\textsubscript{B}~\SGARefLink{sga3:xv:lemma:3:10}{3.10}), equals \(H\);
\item for every \(s\in S\), one has \(H_s=(N_s)^0\).
\end{enumerate}
\end{enumerate}}
\SGARefTarget{sga3:xv:proposition:5:1}

\SGAProof{Proof}
The group prescheme \(H\) is locally of finite presentation over
\(S\), since it is smooth, and has connected fibres; it is therefore
of finite presentation over \(S\) (Exposé
VI\textsubscript{B}~\SGARefLink{par:sga3-VIA-5.3.3}{5.3.3}). Part (a) follows from Exposé \SGARefLink{sga3:xi:corollary:6:11}{XI~6.11}. The
equivalence in \SGARefLink{sga3:xv:proposition:5-1:item:b}{(b)}, recalled here for reference, was proved in Exposé
VI\textsubscript{B}~\SGARefLink{prop:sga3-VIB-6.5.1}{6.5.1}.

We can now state the main theorem of this section.

\SGAStatement{\SGAInlineRefTarget{sga3:xv:theorem:5-2}Theorem 5.2}{Let \(S\) be a prescheme and \(G\) an
\(S\)-group prescheme of finite presentation over \(S\). Let
\(\mathcal L_G\), or simply \(\mathcal L\), be the \(S\)-functor
defined, for every \(S\)-prescheme \(S'\), by
\[
\mathcal L_G(S')=
\left\{
\begin{array}{c}
\text{subgroup preschemes of \(G_{S'}\), smooth over \(S'\),}\\
\text{with connected fibres, equal to their connected normalizer}
\end{array}
\right\}.
\]
Then \(\mathcal L\) is represented by an \(S\)-prescheme that is the
union of an increasing family \((U_i)_{i\in\mathbb N}\) of open
subpreschemes, each quasi-projective and of finite presentation over
\(S\), and hence in particular separated over \(S\).}
\SGARefTarget{sga3:xv:theorem:5:2}

\medskip\noindent\textbf{Initial reductions.}
For every \(r\geq0\), let \(\mathcal L_r\) be the subfunctor of
\(\mathcal L\) defined by
\[
\mathcal L_r(S')=
\left\{
\begin{array}{c}
H\subseteq G_{S'}:\ H\text{ is smooth with connected fibres,}\\
H\text{ equals its connected normalizer, and }\dim(H/S')=r
\end{array}
\right\}.
\]
The fibre dimension of a smooth group is locally constant. Hence
\(\mathcal L_r\to\mathcal L\) is represented by an immersion that is
both open and closed. It is therefore enough to prove the asserted
representability and finiteness properties for every
\(\mathcal L_r\), for then
\[
  \mathcal L=\coprod_{r\in\mathbb N}\mathcal L_r.
\]

For \(n\geq0\), an \(S\)-prescheme \(S'\), and a subgroup prescheme
\(H\) of \(G_{S'}\), denote by \(H^{(n)}\) the \(n\)-th normal
invariant of the identity section \(S'\to H\) (EGA IV~16.1.2).
Thus \(H^{(n)}\) is a finite-type \(\mathcal O_{S'}\)-module
corresponding to the \(n\)-th infinitesimal neighbourhood of the
identity section of \(H\). If \(H\) is smooth of relative dimension
\(r\), then \(H^{(n)}\) is locally free of a rank
\(\varphi(n,r)\) depending only on \(n,r\). Since \(H\) is a
subprescheme of \(G_{S'}\), there is a canonical epimorphism,
compatible with base change,
\[
 G^{(n)}\otimes_{\mathcal O_S}\mathcal O_{S'}
 \xrightarrow{\sim}G_{S'}^{(n)}
 \longrightarrow H^{(n)}.
\]

Introduce the projective \(S\)-scheme
\[
  P_{\varphi(n,r)}
   =\operatorname{Grass}_{\varphi(n,r)}(G^{(n)})
\]
(EGA I, 2nd ed., 9.7; see also Cartan Seminar 1960/61,
Exposé 14 by A. Grothendieck). The assignment
\[
  H\longmapsto H^{(n)}
\]
therefore defines a canonical morphism
\[
  u_{n,r}:\mathcal L_r\longrightarrow P_{\varphi(n,r)}.
\]
The group \(G\) acts naturally on \(G^{(n)}\), and hence on
\(P_{\varphi(n,r)}\), through
\[
 \operatorname{int}:G\longrightarrow
 \operatorname{Aut}_{S\text{-}\mathrm{gr}}(G),
 \qquad g\longmapsto\operatorname{int}(g).
\]

If \(S'\) is quasi-compact over \(S\) and \(H\in\mathcal L_r(S')\),
then for sufficiently large \(n\) one has (Exposé \SGARefLink{sga3:xi:corollary:6:11}{XI~6.11})
\[
 \operatorname{Norm}_{G_{S'}}(H)
 =\operatorname{Norm}_{G_{S'}}(H^{(n)}).
\]
For each \(n\geq0\), define the subfunctor \(\mathcal L_r^n\) of
\(\mathcal L_r\) by
\[
\mathcal L_r^n(S')=
\left\{
 H\in\mathcal L_r(S')\ \middle|\
 \operatorname{Norm}_{G_{S'}}(H)
 =\operatorname{Norm}_{G_{S'}}(H^{(n)})
\right\}.
\]

\medskip\noindent
\textbf{Representability of \(\mathcal L_r^n\).}
The integer \(r\) remains fixed through the proof, so we write
\(\mathcal L,\mathcal L^n,P_n,u_n\) in place of
\(\mathcal L_r,\mathcal L_r^n,P_{\varphi(n,r)},u_{n,r}\). Let
\[
 v_n:\mathcal L^n\longrightarrow P_n
\]
be the restriction of \(u_n\). The definition of \(\mathcal L^n\)
and \SGARefLink{sga3:xv:proposition:5-1:item:b}{Proposition 5.1(b)}(ii) show that \(v_n\) is a monomorphism.
More precisely:

\SGAStatement{Lemma 5.3}{The morphism \(v_n\) is an immersion of
finite presentation. Consequently \(\mathcal L^n\) is represented by
an \(S\)-prescheme, quasi-projective and of finite presentation over
\(S\).}
\SGARefTarget{sga3:xv:lemma:5:3}

After replacing \(S\) by \(P_n\), the usual method reduces the
assertion to the following. Let \(Q\in P_n(S)\), and define a
subfunctor \(F\) of the final functor \(h_S\) by
\[
F(S')=
\begin{cases}
h_S(S'),&\text{if there exists }H\in\mathcal L^n(S')
                    \text{ with }H^{(n)}=Q_{S'},\\
\varnothing,&\text{otherwise.}
\end{cases}
\]
We must prove that \(F\to S\) is an immersion of finite presentation.

Let \(N\) be the normalizer of \(Q\) in \(G\). It is a subgroup
prescheme of finite presentation over \(S\), namely the inverse image
of \(Q\) under
\[
  G\longrightarrow P_n,\qquad g\longmapsto g\cdot Q.
\]
We claim that \(F\) is equivalently described by
\[
F(S')=
\begin{cases}
h_S(S'),&
\begin{array}{l}
N_{S'}\text{ is smooth along the identity section, of}\\
\text{relative dimension }r,\text{ and }(N_{S'})^{(n)}=Q_{S'},
\end{array}\\[2mm]
\varnothing,&\text{otherwise.}
\end{cases}
\]

Call the first and second descriptions \(F_1\) and \(F_2\).
Suppose \(F_1(S')=h_S(S')\), and let \(H\in\mathcal L^n(S')\) with
\(H^{(n)}=Q_{S'}\). Then
\[
\begin{aligned}
\operatorname{Norm}_{G_{S'}}(H)
 &=\operatorname{Norm}_{G_{S'}}(H^{(n)})\\
 &=\operatorname{Norm}_{G_{S'}}(Q_{S'})
 =N_{S'}.
\end{aligned}
\]
By \SGARefLink{sga3:xv:proposition:5-1:item:b}{Proposition 5.1(b)}(ii), \(N_{S'}\) is smooth along the identity
section and its identity component is \(H\). Thus it has relative
dimension \(r\); since \(H\) is open in \(N_{S'}\),
\[
  (N_{S'})^{(n)}=H^{(n)}=Q_{S'}.
\]
Hence \(F_2(S')=h_S(S')\).

Conversely, suppose \(F_2(S')=h_S(S')\). The identity component of
\(N_{S'}\) is then represented (Exposé VI\textsubscript{B}~\SGARefLink{sga3:xv:lemma:3:10}{3.10}) by
a subgroup prescheme \(H\), smooth over \(S'\), with connected fibres
of dimension \(r\). Since \(H\) is invariant and open in \(N_{S'}\),
\[
\begin{aligned}
N_{S'}&\subseteq\operatorname{Norm}_{G_{S'}}(H)
\subseteq\operatorname{Norm}_{G_{S'}}(H^{(n)})\\
&=\operatorname{Norm}_{G_{S'}}((N_{S'})^{(n)})
=\operatorname{Norm}_{G_{S'}}(Q_{S'})=N_{S'}.
\end{aligned}
\]
All inclusions are equalities. The first says that \(H\) equals its
connected normalizer, and the second says
\(H\in\mathcal L^n(S')\). Hence \(F_1=F_2\).

We retain the second description. Condition (a) is represented by an
immersion of finite presentation, by the following lemma.

\SGAStatement{\SGAInlineRefTarget{sga3:xv:lemma:5-4}Lemma 5.4}{Let \(S\) be a prescheme, let \(X\) be
locally of finite presentation over \(S\), let \(\sigma:S\to X\) be
a section, let \(r\geq0\), and define the subfunctor \(L\) of \(S\)
by
\[
L(S')=
\begin{cases}
h_S(S'),&
\begin{array}{l}
X_{S'}\text{ is smooth along }\sigma_{S'},\text{ of relative}\\
\text{dimension }r\text{ at the points of }\sigma_{S'}(S'),
\end{array}\\[1mm]
\varnothing,&\text{otherwise.}
\end{cases}
\]
Then:
\begin{enumerate}
\renewcommand{\labelenumi}{\textup{\alph{enumi})}}
\item \(L\to S\) is an immersion of finite presentation;
\item let \(\mathcal J\) be the conormal sheaf of \(S\to X\).
If for every \(s\in S\) the vector space
\(\mathcal J\otimes_{\mathcal O_{S,s}}\kappa(s)\) has rank at most
\(r\), then \(L\to S\) is a closed immersion.
\end{enumerate}}
\SGARefTarget{sga3:xv:lemma:5:4}

\SGAProof{Proof}
The functor \(L\) is local on \(S\), so we may suppose \(S\) affine.
After replacing \(X\) by a neighbourhood of \(\sigma(S)\), we may
suppose \(X\) of finite presentation, and then reduce to noetherian
\(S\) by EGA IV~8.9. In part (b), use also compatibility of the
conormal sheaf with base change (EGA IV~16.6.4) and constructibility
of the rank of the fibres of \(\mathcal J\).

For an \(S\)-prescheme \(S'\), let
\(\mathcal J'=\mathcal J_{S'}\), let
\(S^\bullet(\mathcal J')\), canonically isomorphic to
\(S^\bullet(\mathcal J)_{S'}\), be the symmetric algebra of
\(\mathcal J'\) over \(\mathcal O_{S'}\), let
\(\operatorname{Gr}^\bullet(\sigma_{S'})\) be the graded algebra
associated with \(\sigma_{S'}\) (EGA IV~16), and let
\[
 \sigma_{n,S'}:S^n(\mathcal J')
 \longrightarrow\operatorname{Gr}^n(\sigma_{S'})
\]
be the canonical epimorphism. By EGA IV~17.12.3 and
EGA \(0_{\mathrm{IV}}\)~19.5.4, \(X_{S'}\) is smooth along
\(\sigma_{S'}\) of relative dimension \(r\) if and only if:
\begin{enumerate}
\renewcommand{\labelenumi}{\textup{\roman{enumi})}}
\item \(\mathcal J'\) is locally free of rank \(r\);
\item \(\sigma_{n,S'}\) is an isomorphism for every \(n\geq1\).
\end{enumerate}

TDTE IV, \SGARefLink{sga3:xv:lemma:3:6}{Lemma 3.6}, shows that the functor making \(\mathcal J\)
locally free of rank \(r\) is represented by a subprescheme \(S_1\),
closed in \(S\) in case (b). Replacing \(S\) by \(S_1\), assume that
\(\mathcal J\) is locally free of rank \(r\). Proceed by induction on
\(n\). Suppose a closed subprescheme \(S_{n-1}\) represents the
functor making \(\sigma_{q,S'}\) injective for \(q\leq n-1\).
The subfunctor\footnote{The editors replaced ``functor'' by
``subfunctor'' here, and ask whether the intended wording is instead
``which makes \(\sigma_{q,S'}\) injective for \(q\leq n\)''.}
making \(\sigma_{n,S'}\) injective is represented by a closed
subprescheme \(S_n\) of \(S_{n-1}\).

Indeed, after replacing \(S\) by \(S_{n-1}\), all
\(\sigma_{q,S}\) for \(q\leq n-1\) are isomorphisms. The
\((n-1)\)-st normal invariant \(X^{(n-1)}\) of \(\sigma_S\) has a
filtration with successive quotients
 \(\operatorname{Gr}^0(\sigma_S),\ldots,
 \operatorname{Gr}^{n-1}(\sigma_S)\), all locally free; hence
\(X^{(n-1)}\) is flat over \(S\). Compatibility of the normal
invariants with base change gives
\[
\begin{aligned}
\operatorname{Gr}^n(\sigma_{S'})
 &=\ker\!\left(X_{S'}^{(n)}\to X_{S'}^{(n-1)}\right)
  =\operatorname{Gr}^n(\sigma)_{S'},\\
\sigma_{n,S'}&=(\sigma_{n,S})_{S'}.
\end{aligned}
\]
Thus we must represent the functor making \(\sigma_{n,S}\) injective.
Locally we may suppose \(S\) affine and \(S^n(\mathcal J)\) free.
The required functor is represented by the closed subscheme defined
by the ideal generated by the coordinates of
\(\ker\sigma_{n,S}\) in a basis of \(S^n(\mathcal J)\).

Since \(S\) is noetherian, the descending chain \((S_n)\) stabilizes.
Its stable value represents \(L\), proving the lemma.

Return to \(\mathcal L^n\). After replacing \(S\) by the subprescheme
provided by \SGARefLink{sga3:xv:lemma:5-4}{Lemma 5.4}, assume \(N\) smooth along the identity section
of relative dimension \(r\). Then \(\mathcal L^n\) is the coincidence
functor of two sections \(h,g:S\to P_n\), corresponding to \(Q\) and
\(N^{(n)}\). It is therefore represented by the inverse image of the
diagonal of \(P_n\times_SP_n\) under the finite-presentation morphism
\(h\times_Sg\). This proves \SGARefLink{sga3:xv:lemma:5:3}{Lemma 5.3}.

\medskip\noindent
\textbf{The transition morphisms
\(\mathcal L^n\to\mathcal L^m\).}
If \(H\in\mathcal L^n(S')\), then \(H\in\mathcal L^m(S')\) for every
\(m\geq n\). Thus there are natural monomorphisms
\[
  u_n^m:\mathcal L^n\longrightarrow\mathcal L^m
  \qquad(m\geq n).
\]

\SGAStatement{Lemma 5.5}{The morphism \(u_n^m\) is an open immersion.}
\SGARefTarget{sga3:xv:lemma:5:5}

After a base change, it is enough to consider \(H\in\mathcal L^m(S)\)
and
\[
 N=\operatorname{Norm}_G(H)=\operatorname{Norm}_G(H^{(m)}),
 \qquad
 N'=\operatorname{Norm}_G(H^{(n)}).
\]
Let \(D\) be the coincidence functor of \(N,N'\): \(D(S')=h_S(S')\)
if \(N_{S'}=N'_{S'}\), and \(D(S')=\varnothing\) otherwise. We must
show that \(D\to S\) is an open immersion.

The same functor is obtained by requiring that the immersion
\(H\to N'\) become open. One implication follows from
\SGARefLink{sga3:xv:proposition:5-1}{Proposition 5.1}. Conversely, if \(H_{S'}\to N'_{S'}\) is open, then,
because \(H\) has connected fibres, \(H_{S'}\) is the identity
component of \(N'_{S'}\) (Exposé VI\textsubscript{B}~\SGARefLink{sga3:xv:lemma:3:10}{3.10}), hence is
invariant in \(N'_{S'}\). Thus \(N'_{S'}\subseteq N_{S'}\). The
opposite inclusion always holds, so equality follows. Since \(H,N'\)
are of finite presentation and \(H\) is flat over \(S\), Exposé
VI\textsubscript{B}~\SGARefLink{sga3:xv:remark:2:6}{2.6} shows that \(D\to S\) is open.

\medskip\noindent\textbf{End of the proof of \SGARefLink{sga3:xv:theorem:5-2}{Theorem 5.2}.}
The functors \(\mathcal L^n\) are represented, and the transition
morphisms \(u_n^m\) are compatible open immersions. There is therefore
an \(S\)-prescheme \(X\), the union of an increasing sequence of open
subpreschemes \((X_i)_{i\in\mathbb N}\), such that \(X_i\) represents
\(\mathcal L^i\) and \(X_i\to X_j\) corresponds to \(u_i^j\). In the
category of sheaves for the Zariski topology on \(\mathrm{Sch}/S\),
\[
  X=\varinjlim_iX_i,\qquad
  \mathcal L=\varinjlim_i\mathcal L^i.
\]
Thus \(X\) represents \(\mathcal L\), completing the proof.

\SGAStatement{Remark 5.6}{Suppose all residue characteristics of
\(S\) are zero. Then, for every \(r\geq0\),
\(\mathcal L_r=\mathcal L_r^1\), so \(\mathcal L_r\) is represented
by an \(S\)-prescheme quasi-projective and of finite presentation over
\(S\).}
\SGARefTarget{sga3:xv:remark:5:6}

Indeed, for \(H\in\mathcal L_r(S')\), it suffices to prove that
\[
  H\longrightarrow N=\operatorname{Norm}_{G_{S'}}(H^{(1)})
\]
is an open immersion. Because \(H\) is flat and both \(H,N\) are of
finite presentation, this may be checked on the fibres (Exposé
VI\textsubscript{B}~\SGARefLink{sga3:xv:remark:2:6}{2.6}). Cartier's theorem (Exposé
VI\textsubscript{B}~\SGARefLink{sga3:xv:proposition:1-6}{1.6}) gives, for a connected algebraic subgroup
\(H\) of an algebraic group in characteristic zero,
\[
 \operatorname{Norm}_G(H)
 =\operatorname{Norm}_G(\operatorname{Lie}H)
 =\operatorname{Norm}_G(H^{(1)}).
\]
\footnote{For example, one may apply Proposition II.6.1 of
M. Demazure and P. Gabriel, \emph{Groupes algébriques} I,
Masson--North Holland, 1970.}

In nonzero residue characteristic the subfunctors
\(\mathcal L_r^n\) can form a strictly increasing sequence, even when
\(S\) is quasi-compact, and then \(\mathcal L_r\) is not represented
by an \(S\)-prescheme quasi-compact over \(S\). For example, over a
field \(k\) of characteristic \(p>0\), let
\[
 T=\mathbb G_m\times\mathbb G_m,\qquad
 U=\mathbb G_a\times\mathbb G_a,
\]
and let \(G=T\ltimes U\), where
\((t,t')\cdot(u,u')=(tu,t'u')\). For \(n>0\), let \(U_n\subset U\)
be defined by \(u'=u^{p^n}\), and let \(T_n\subset T\) be defined by
\(t'=t^{p^n}\). Then \(T_n\) acts on \(U_n\), and
\(G_n=T_nU_n\) is smooth, connected, and equal to its normalizer in
\(G\). For every \(m\), the distinct groups \(G_n\), \(n\geq m\),
have the same infinitesimal neighbourhood of order \(p^m\).

\SGAStatement{Remark 5.7}{There is a canonical invertible sheaf
\(L\) on \(\mathcal L\) whose restriction to every open subset of
\(\mathcal L\) quasi-compact over \(S\) is \(S\)-ample. One may take
\[
  L=(\det(\operatorname{Lie}H))^{-1},
\]
where \(H\subseteq G_{\mathcal L}\) is the universal subgroup.}
\SGARefTarget{sga3:xv:remark:5:7}

Recall that for a finite-rank locally free sheaf \(\mathcal F\), its
determinant restricts on the open-and-closed locus where
\(\operatorname{rank}\mathcal F=r\) to \(\bigwedge^r\mathcal F\).
It is enough to restrict to \(\mathcal L_r^n\). Let
\(v_n:\mathcal L_r^n\to P_n\) be the canonical immersion, and let
\(Q\) be the universal locally free quotient on the Grassmannian
\(P_n\). By construction,
\[
  v_n^*Q=H^{(n)}.
\]
Since \(\det Q\) is the canonical ample sheaf on \(P_n\) (EGA I,
2nd ed., 9.7), \(\det H^{(n)}\) is ample relative to \(S\)
(EGA II~4.6.13(i) bis).

Let \(\mathcal J=H^{(1)}\) be the conormal sheaf, and let
\(S^q(\mathcal J)\) be the homogeneous degree-\(q\) part of its
symmetric algebra. Smoothness of \(H\) gives
\[
  \det H^{(n)}
  \xrightarrow{\sim}
  \bigotimes_{1\leq q\leq n}\det S^q(\mathcal J).
\]
For every locally free \(\mathcal J\) of rank \(r\) and every \(q>0\)
there is a canonical isomorphism
\[
  \det S^q(\mathcal J)
  \xrightarrow{\sim}
  (\det\mathcal J)^{\otimes s},
\]
where \(s>0\) depends only on \(r,q\). Thus, for some \(s>0\),
\[
  \det H^{(n)}\simeq(\det\mathcal J)^{\otimes s}.
\]
Since \(\det\mathcal J=(\det\operatorname{Lie}H)^{-1}\), EGA
II~4.5.6 proves that \(L\) is \(S\)-ample.

\SGAStatement{Remark 5.8}{Let \(G,H\) be \(S\)-group preschemes of
finite presentation and let \(i:H\to G\) be a monomorphism. Suppose
\(H\) is smooth over \(S\) with connected fibres. The normalizer
\(N=\operatorname{Norm}_G(H)\) is then represented by a closed
subgroup prescheme of \(G\), of finite presentation over \(S\)
(Exposé \SGARefLink{sga3:xi:corollary:6:11}{XI~6.11}). If, in addition, \(N\) is smooth along the identity
section and has the same relative dimension as \(H\), then
\(H=N^0\); in particular \(H\) is a subgroup prescheme of \(G\).}
\SGARefTarget{sga3:xv:remark:5:8}

Indeed, \(N^0\) is represented by an open subgroup prescheme of \(N\),
smooth over \(S\) (Exposé VI\textsubscript{B}~\SGARefLink{sga3:xv:lemma:3:10}{3.10}), and \(i\)
factors through \(N^0\). For every \(s\in S\), the connected smooth
algebraic groups \(H_s,(N^0)_s\) have the same dimension and hence
coincide. Since \(H\) is flat, \(H\to N^0\) is an isomorphism
(EGA IV~17.9.5). Thus the functor \(\mathcal L\) of this section may
equivalently be described as the functor of subgroup preschemes
\(H\) of \(G\) that are smooth over \(S\), have connected fibres,
and equal their connected normalizer.

% END INCLUDED SGA3_English_current_buildable_source_20260823/source/inputs/XV/tex/components/05_expose_XV_section5.tex

% BEGIN INCLUDED SGA3_English_current_buildable_source_20260823/source/inputs/XV/tex/components/06_expose_XV_section6.tex
% Authority: Expo15.pdf, local pp. 38--49 / combined pp. 908--919 /
% printed pp. 421--444.
% Stops immediately before §7.

\SGASection{6}{Functor of Cartan subgroups and functor of parabolic
subgroups}
\SGARefTarget{sga3:xv:section:6}

When \(G\) is a smooth connected algebraic group defined over an
algebraically closed field \(k\), sub-tori of \(G\), maximal
sub-tori, Cartan subgroups (Exposé XII~1), Borel subgroups (Exposé
\SGARefLink{sga3:xiv:4.1}{XIV~4.1}), and parabolic subgroups (Exposé \SGARefLink{sga3:xiv:4.8.1}{XIV~4.8 bis}) have already
been defined. We extend these notions to a group prescheme over an
arbitrary base as follows.

\textbf{\SGAInlineRefTarget{sga3:xv:definition:6-1}Definition 6.1.} \emph{Let \(S\) be a prescheme, let \(G\)
be an \(S\)-group prescheme of finite presentation, let \(H\) be an
\(S\)-group prescheme, and let \(i:H\to G\) be an
\(S\)-monomorphism making \(H\) a subgroup of \(G\). We say that
\(H\) is a sub-torus of \(G\) (respectively, a maximal sub-torus, a
Cartan subgroup, a Borel subgroup, or a parabolic subgroup) if:}

\begin{enumerate}
\renewcommand{\labelenumi}{\textup{\roman{enumi})}}
\item
  \emph{i) \(H\) is smooth over \(S\).}
\item
  \emph{ii) For every geometric point \(\bar s\) above \(S\),
  \(H_{\bar s}\) is a sub-torus of
  \((G_{\bar s})^{0}_{\mathrm{red}}\) (respectively, a maximal
  sub-torus, a Cartan subgroup, a Borel subgroup, or a parabolic
  subgroup).}
\end{enumerate}
\SGARefTarget{sga3:xv:definition:6:1}

\textbf{\SGAInlineRefTarget{sga3:xv:remark:6-1-bis}Remarks 6.1 bis.}

\begin{enumerate}
\renewcommand{\labelenumi}{\textup{\alph{enumi})}}
\item
  If the \(S\)-group \(H\) is a sub-torus of \(G\)
  (respectively, \(\ldots\)), its fibres are connected, and
  consequently \(H\) is of finite presentation over \(S\) (Exposé
  VI\textsubscript{B}~\SGARefLink{par:sga3-VIA-5.3.3}{5.3.3}).
\item
  If \(H\) is a sub-torus of \(G\), then \(H\) is a torus in the
  sense of Exposé IX, as follows immediately from Exposé \SGARefLink{sga3:x:theorem:8-1}{X~8.1}.
  Moreover, the monomorphism \(i:H\to G\) is an immersion (see \SGARefLink{sga3:xv:lemma:8:3}{8.3}
  below).
\item
  If \(G\) is smooth over \(S\), with connected fibres, and if
  \(H\) is a Cartan subgroup of \(G\) (respectively, a Borel
  subgroup or a parabolic subgroup), the monomorphism \(i:H\to G\)
  is an immersion. Thus our definitions coincide with those
  introduced in Exposés XII and XIV. Indeed, \(H\) is then equal to
  its connected normalizer (by \SGARefLink{sga3:xii:theorem:6-6:item:c}{XII~6.6(c)}, \SGARefLink{sga3:xiv:4.8}{XIV~4.8} and \SGARefLink{sga3:xiv:paragraph:4-8-bis}{4.8 bis}),
  and it suffices to apply \SGARefLink{sga3:xv:remark:5:8}{5.8}.
\end{enumerate}

\textbf{Definition 6.1 ter.} \emph{Let \(S\) be a prescheme, let
\(G\) be an \(S\)-group prescheme locally of finite type, and let
\(s\) be a point of \(S\). The \textbf{nilpotent rank} of \(G\) at
\(s\), denoted by \(\rho_n(s)\), is the dimension of the Cartan
subgroups of \((G_s)^0_{\mathrm{red}}\). The reductive rank
\(\rho_r(s)\), unipotent rank \(\rho_u(s)\), and abelian rank
\(\rho_{\mathrm{ab}}(s)\) are defined similarly (see Exposé
\SGARefLink{sga3:x:remark:8-7}{X~8.7}).}
\SGARefTarget{sga3:xv:definition:6:1:ter}

If \(G\) is now a smooth connected algebraic group defined over an
algebraically closed field \(k\), recall that the \textbf{radical}
of \(G\), denoted by \(\operatorname{rad}(G)\), is the largest
invariant smooth connected solvable algebraic subgroup of \(G\);
\(G/\operatorname{rad}(G)\) is then semisimple (use Exposé \SGARefLink{sga3:xii:6.1}{XII~6.1}
to reduce to the affine case). If \(G\) is moreover affine, the
\textbf{unipotent radical} \(\operatorname{rad}^{u}(G)\) is defined
as the largest invariant smooth connected unipotent algebraic
subgroup of \(G\); \(G/\operatorname{rad}^{u}(G)\) is then
reductive.

\textbf{\SGAInlineRefTarget{sga3:xv:proposition:6-2}Proposition 6.2.} \emph{Let \(S\) be a prescheme, let \(G\)
and \(H\) be two \(S\)-group preschemes of finite presentation, and
let \(i:H\to G\) be an \(S\)-monomorphism making \(H\) a subgroup
of \(G\). Let \(P(s)\) be one of the following properties concerning
a point \(s\) of \(S\):}

\begin{enumerate}
\renewcommand{\labelenumi}{\textup{\roman{enumi})}}
\item \emph{\((G_{\bar s})^{0}_{\mathrm{red}}\) is an abelian
variety (respectively, is affine, is a torus, or is unipotent).}
\item \emph{\((H_{\bar s})^{0}_{\mathrm{red}}\) is a maximal torus
of \(G_{\bar s}\).}
\item\SGAInlineRefTarget{sga3:xv:proposition:6-2:item:iii} \emph{\((H_{\bar s})^{0}_{\mathrm{red}}\) is the centralizer
in \((G_{\bar s})^{0}_{\mathrm{red}}\) of a torus of \(G_{\bar s}\)
(respectively, is a Cartan subgroup of \(G_{\bar s}\)).}
\item \emph{\((H_{\bar s})^{0}_{\mathrm{red}}\) is a Borel subgroup
(respectively, a parabolic subgroup) of \(G_{\bar s}\).}
\item \emph{\((H_{\bar s})^{0}_{\mathrm{red}}\) is the radical of
\(G_{\bar s}\) (respectively,
\((G_{\bar s})^{0}_{\mathrm{red}}\) is semisimple).}
\item \emph{\(G_{\bar s}\) is affine and
\((H_{\bar s})^{0}_{\mathrm{red}}\) is the unipotent radical of
\(G_{\bar s}\) (respectively,
\((G_{\bar s})^{0}_{\mathrm{red}}\) is reductive).}
\end{enumerate}

\emph{Then the set \(E\) of points \(s\) of \(S\) for which \(P(s)\)
holds is locally constructible (EGA \(0_{\mathrm{III}}\), 9.1.11).}
\SGARefTarget{sga3:xv:proposition:6:2}

\textbf{\SGAInlineRefTarget{sga3:xv:remark:6-2-1}Remark 6.2.1.} This proposition complements Exposé
VI\textsubscript{B}, \S10. The structure of \(E\) can be described
more precisely by using semicontinuity theorems (see Exposé \SGARefLink{sga3:x:remark:8-7}{X~8.7});
we shall see an example below.

\SGAProof{Proof of \SGARefLink{sga3:xv:proposition:6-2}{6.2}}
Note that if \(S\) is the spectrum of a field, \(E\) is invariant
under extension of that field. A standard reduction (EGA IV~9)
then allows us to reduce to the case where \(S\) is noetherian and
integral, with generic point \(\eta\). One must show that \(E\) or
\(\operatorname{ens}(S)\setminus E\) contains a neighbourhood of
\(\eta\) (EGA IV~9.2.1). One
may assume \(S\) affine of ring \(A\) and field of fractions \(K\). If \(L\) is a finite extension of \(K\), it is immediate that
there exists an \(A\)-subalgebra \(B\) of \(L\), finite over \(A\), having \(L\) as field of fractions. The canonical morphism \(S'
\to S\), where \(S' = \operatorname{Spec} B\), is dominant, of finite presentation, so the image of a non-empty open subset
of \(S'\) contains a non-empty open subset of \(S\) (EGA IV 1.8.4). From the viewpoint that interests us, we may therefore
replace \(S\) by \(S'\), hence replace \(K\) by a finite extension \(L\). Thus we may choose \(L\) so that \((G_{L})_{red}\) and
\((H_{L})_{red}\) are smooth over \(L\) (EGA IV 4.6.6). Possibly restricting \(S'\), we may assume that \(G_{red}\) and
\(H_{\mathrm{red}}\) are group preschemes smooth over \(S\)
(Exposé VI\textsubscript{B}, \S10, and EGA IV~17). In view of the
properties to be proved, we may replace \(G\) and \(H\) by their
reduced identity components (Exposé VI\textsubscript{B}~\SGARefLink{lem:sga3-VIB-10.9}{10.9}), and
hence assume that \(G\) and \(H\) are smooth over \(S\) with
connected fibres. Finally, we may assume that \(H\) is a closed
subgroup prescheme of \(G\) (Exposé VI\textsubscript{B}~10.4).

\emph{Proof of i).} Possibly after a finite extension of \(K\), we
may assume that \(G_{\eta}\) admits a ``Chevalley decomposition'',
i.e. is an extension of an abelian variety \(D_{\eta}\) by a smooth
connected linear algebraic group \(F_{\eta}\) (Séminaire Bourbaki
1956/57, no.~145). By Exposé VI\textsubscript{B}~10.16, there is an
open neighbourhood \(U\) of \(\eta\) over which this generic
extension comes from an extension
\[
  1\longrightarrow F\longrightarrow G|_U\longrightarrow D
  \longrightarrow 1.
\]
One may moreover assume that \(F\) and \(D\) are smooth over \(U\)
with connected fibres, that \(F\) is affine over \(U\), and that
\(D\) is proper over \(U\) (EGA IV~8, 9 and 17). For every point
\(s\) of \(U\), \((D_s,F_s)\) is then the Chevalley decomposition of
\(G_s\). Consequently, \(G_s\) is an abelian variety (respectively,
is affine) if and only if \(F_s\) (respectively, \(D_s\)) is the unit
group; this is a constructible condition (EGA IV~9.2.6.1).

To establish the last two assertions of (i), we may, by the
preceding, assume that \(G\) is affine over \(S\). Let \(q\) be a
prime number invertible on \(S\), and let \({}_qG\) be the
``kernel'' of the \(q\)-th-power morphism in \(G\). It follows
easily from the structure of affine algebraic groups that \(G_s\)
is a torus (respectively, is unipotent) if and only if
\begin{enumerate}
\renewcommand{\labelenumi}{\textup{\alph{enumi})}}
\item \({}_qG_s\) is quasi-finite, which is a constructible
property (EGA IV~9.3.2), and
\item \({}_qG_s\) has \(q^r\) geometric points, where \(r\) is the
relative dimension of \(G\) over \(S\) (respectively,
\({}_qG_s\) has a single point).
\end{enumerate}
The function
\[
 s\longmapsto
 \bigl(\text{number of geometric points of }{}_qG_s\bigr)
\]
is constructible (EGA IV~9.7.9). This completes the proof of (i).

\emph{Proof of iii).} \textbf{a) Case of the centralizer of a
torus.} Suppose
\[
  H_{\eta}=\operatorname{Centr}_{G_{\eta}}(T_{\eta}),
\]
where \(T_{\eta}\) is a torus of \(G_{\eta}\). We shall show that
\(H_s\) is the centralizer of a sub-torus of \(G_s\) for \(s\) in a
neighbourhood of \(\eta\). By (i), after possibly restricting \(S\),
we may assume that \(T_{\eta}\) comes from a sub-torus \(T\) of
\(G\). Then \(Z=\operatorname{Centr}_{G}(T)\) is represented
(Exposé \SGARefLink{sga3:xi:corollary:6:11}{XI~6.11}) by a subgroup prescheme of \(G\). Since \(H\) and
\(Z\) coincide generically, they coincide over a neighbourhood of
\(\eta\). Thus the set \(E\) of points \(s\) of \(S\) for which
\(H_s\) is the centralizer in \(G_s\) of a sub-torus of \(G_s\) is
ind-constructible (EGA IV~1). This will suffice in Lemma~6.6 below
to prove that \(E\) is open in \(S\), and hence, \emph{a fortiori},
locally constructible.

\textbf{b) Case of a Cartan subgroup.} Suppose that \(H_{\eta}\) is
a Cartan subgroup of \(G_{\eta}\). We shall show that \(H_s\) is a
Cartan subgroup of \(G_s\) for every \(s\) in a neighbourhood of
\(\eta\). The group \(H_{\eta}\) is the centralizer in \(G_{\eta}\)
of a torus of \(G_{\eta}\), and is nilpotent (Exposé \SGARefLink{sga3:xii:6.6}{XII~6.6}). By
\SGARefLink{sga3:xii:6.6:item:a}{(a)} and Exposé VI\textsubscript{B}~\SGARefLink{sga3:xv:definition:8-4}{8.4}, \(H_s\) has the same
properties for every \(s\) in some neighbourhood \(U\) of \(\eta\).
For every \(s\in U\), the group \(H_s\) therefore has the same
reductive rank as \(G_s\), and its unique maximal torus is central
(Exposé \SGARefLink{sga3:xii:6.7}{XII~6.7}). Hence \(H_s\) is the centralizer of a maximal torus
of \(G_s\), that is, a Cartan subgroup of \(G_s\).

Suppose now that \(H\) is not a Cartan subgroup of \(G\), and let us show that \(H_{s}\) is not a Cartan subgroup of \(G_{s}\)
for \(s\) in a neighborhood \(U\) of \(\eta\). In view of the assertion provisionally admitted in (a) above, we may restrict
to the case where \(H\) is the centralizer in \(G\) of a sub-torus \(T\). But then \(H_{\eta}\) contains a Cartan subgroup
\(C_{\eta}\) of \(H_{\eta}\). We have just seen that, possibly restricting \(S\), \(C_{\eta}\) extends to a Cartan subgroup \(C\)
of \(G\), which one may assume contained in \(H\). By hypothesis \(H_{\eta}\) strictly contains \(C_{\eta}\), so \(H_{s}\)
strictly contains \(C_{s}\) for \(s\) in a neighborhood \(U\) of \(\eta\) (EGA IV 9.5.2); \emph{a fortiori}, \(H_{s}\) is not

a Cartan subgroup of \(G_{s}\) for \(s\) in \(U\).

\emph{Proof of ii).} Suppose \(H_{\eta}\) is a maximal torus of \(G_{\eta}\) and let \(C_{\eta}\) be its centralizer in \(G_{\eta}\).
By i) and iii), \(H\) is a torus over a neighborhood \(U\) of \(\eta\) and \(C = \operatorname{Centr}_{G|U}(H|U)\) is a Cartan subgroup of
\(G|U\). To prove that \(H\) is a maximal torus of \(G\) over a neighborhood of \(\eta\), we may then replace \(G\) by \(C\), then
by the linear component \(F\) of a Chevalley decomposition of \(C\) (cf.~i)). Let \(q\) be an integer invertible on \(S\),
\({}_{q}F\) the kernel of the \(q\)-th-power morphism in \(F\). Since \(F_{s}\) is affine, nilpotent, smooth and connected,
\(F_{s}\) is the direct product of its maximal torus \(T_{s}\) by a unipotent group (\emph{Bible} 6-04), so
\({}_{q}F_s={}_{q}T_s\). Since \(H_{\eta}\) is a maximal torus,
\({}_{q}H_{\eta}={}_{q}F_{\eta}\), and consequently
\({}_{q}H={}_{q}F\) over a neighborhood \(V\) of \(\eta\).
For every point \(s\) of \(V\), \({}_{q}H_s={}_{q}T_s\), so
\(H_{s} = T_{s}\) is a maximal torus.

Suppose now that \(H_{\eta}\) is not a maximal torus of \(G_{\eta}\). By i), we may restrict to the case where \(H_{\eta}\) is
a torus, then assume that it is contained in a strictly larger torus \(T_{\eta}\). The latter extends to a torus \(T\)
strictly containing \(H\) over a neighborhood \(U\) of \(\eta\). \emph{A fortiori}, \(H_{s}\) is not a maximal torus for \(s \in U\).

\emph{Proof of iv).} Possibly restricting \(S\), we may assume that the center \(Z\) of \(G\) is representable (Exp. VI\textsubscript{B} 10.11)
and flat over \(S\), as well as the quotient \(G/Z\) (\emph{loc. cit.}). The property ``\(H_{s}\) contains \(Z_{s}\)'' is
constructible (EGA IV 9.5.2) and every parabolic subgroup of \(G_{s}\) contains \(Z_{s}\) (Exp. XIV 4.9 a)); this allows us
to replace \(G\) by \(G/Z\), hence assume \(G\) affine over \(S\) (Exp. XII 6.1 and i)). We may further

assume that \(G/H\) is representable, but then \(H_{s}\) is a parabolic subgroup of \(G_{s}\) if and only if \((G/H)_{s}\) is
proper (\emph{Bible} 6 Th. 4 b)), which is an ind-constructible property (EGA IV 9.3.5). So \(E\) is ind-constructible, and
this will suffice for us to prove that \(E\) is open (Lemma 6.6), hence locally constructible.

Let us now examine the case of Borel subgroups. If \(H_{\eta}\) is a Borel subgroup of \(G_{\eta}\), i.e.~a solvable
parabolic subgroup of \(G_{\eta}\), what precedes and Exposé VI\textsubscript{B}, \SGARefLink{prop:sga3-VIB-8.4}{8.4}, entail that these properties remain true at every
point \(s\) of a neighborhood of \(\eta\). If now \(H_{\eta}\) is not a Borel subgroup of \(G_{\eta}\), to prove that the same
holds at points \(s\) of a neighborhood of \(\eta\), we may restrict (in view of what precedes) to the case where \(H_{\eta}\)
is a parabolic subgroup, then assume that \(H_{\eta}\) contains a Borel subgroup \(B_{\eta}\). We have just shown that the
latter extends to a Borel subgroup \(B\) of \(G\) over a neighborhood \(U\) of \(\eta\). Since \(H_{\eta}\) strictly contains
\(B_{\eta}\), \(H_{s}\) strictly contains \(B_{s}\) at every point of an open subset \(V\), and \(H_{s}\) is not a Borel subgroup
of \(G_{s}\) for \(s \in V\).

\emph{Proof of v).} Suppose \(H_{\eta}\) is the radical of \(G_{\eta}\). The group \(H_{\eta}\) is then invariant in \(G_{\eta}\),
solvable (smooth and connected); the same is therefore true for \(H_{s}\) for \(s\) belonging to a neighborhood \(U\) of
\(\eta\) (Exp. VI\textsubscript{B} 8 and 10), so for \(s \in U\), \(H_{s}\) is contained in the radical of \(G_{s}\). Replacing \(G\) by \(G/H\)
(Exp. VI\textsubscript{B} 10), we must

prove that if \(G_{\eta}\) is semisimple, then \(G_{s}\) is semisimple at every point of a neighborhood \(V\) of \(\eta\). Using
i) and ii), one may assume that \(G\) is affine over \(S\) and possesses a maximal torus \(T\). Let \(W\) be the Weyl group of
\(T\) (\SGARefLink{sga3:xii:section:2}{Exp. XII 2}), which is quasi-finite and étale over \(S\), hence finite and étale over an open subset \(V\). It then
follows from the elementary properties of roots (\SGARefLink{sga3:xix:proposition:1-12}{Exp. XIX 1.12}) that \(G\) is semisimple over \(V\).

Suppose now that \(H_{\eta}\) is not the radical of \(G_{\eta}\). Possibly replacing \(K\) by a finite extension \(L\), we may
assume that \(G_{\eta}\) admits a radical \(R_{\eta}\). By what precedes, \(R_{\eta}\) extends to a subgroup prescheme \(R\) of
\(G|_{U}\) such that for every \(s \in U\), \(R_{s}\) is the radical of \(G_{s}\). By hypothesis, \(R_{\eta} \neq H_{\eta}\). So
\(R_{s} \neq H_{s}\) for \(s \in V\). It remains to prove that if \(G_{\eta}\) is not semisimple, neither is \(G_{s}\) at
neighboring points, but this is a particular case of what precedes (take \(H\) = unit group).

\emph{Proof of vi).} The proof is entirely analogous to that of v), in view of i), and is left to the care of the reader.

\textbf{\SGAInlineRefTarget{sga3:xv:corollary:6-3}Corollary 6.3.} \emph{Let \(S_0\) be a quasi-compact
prescheme, let \((S_i)_{i\in L}\) be a projective system of
\(S_0\)-preschemes affine over \(S_0\), and let
\(S=\varprojlim S_i\) (EGA IV~8.2). Let \(G_0\) be a group
prescheme of finite presentation over \(S_0\), set
\[
  G_i=G_0\times_{S_0}S_i,\qquad G=G_0\times_{S_0}S,
\]
and let \(H\) be a subgroup of \(G\). If \(H\) is a sub-torus of
\(G\) (respectively, a maximal sub-torus, a Cartan subgroup, a Borel
subgroup, or a parabolic subgroup), then there exist an index
\(i\in L\) and a subgroup \(H_i\) of \(G_i\) such that
\[
  H=H_i\times_{S_i}S
\]
and \(H_i\) is a sub-torus of \(G_i\) (respectively, \(\ldots\)).}
\SGARefTarget{sga3:xv:corollary:6:3}

Indeed, \(H\) is smooth with connected fibres, hence of finite
presentation over \(S\) (Exposé VI\textsubscript{B}~\SGARefLink{par:sga3-VIA-5.3.3}{5.3.3}). By
Exposé VI\textsubscript{B}, \S10, there exist \(i\in L\) and a
subgroup \(H_i\) of \(G_i\), smooth over \(S_i\), such that
\(H=H_i\times_{S_i}S\). Corollary~6.3 then follows from
Definition~6.1, Proposition~6.2 and EGA IV~9.3.3.

\textbf{Corollary 6.3 bis.} \emph{Let \(S\) be a prescheme and let
\(G\) be an \(S\)-group prescheme of finite presentation. Then the
functions \(\rho_n,\rho_r,\rho_u,\rho_{\mathrm{ab}}\) (cf.~\SGARefLink{sga3:xv:definition:6:1:ter}{6.1 ter})
are locally constructible on \(S\).}
\SGARefTarget{sga3:xv:corollary:6:3:bis}

It suffices to show (EGA IV 9.) that if \(S\) is a noetherian integral scheme with generic point \(\eta\), the functions in
question are constant on a neighborhood of \(\eta\). Possibly replacing \(S\) by a scheme \(S'\) finite over \(S\) and
dominating \(S\), we may assume that \(G_{\eta}\) admits a Cartan subgroup \(C_{\eta}\) with a Chevalley decomposition \(1 \to
L_{\eta} \to C_{\eta} \to A_{\eta} \to 1\). The argument made in \SGARefLink{sga3:xv:proposition:6-2}{6.2} i) proves that this decomposition extends to a
Chevalley decomposition over a neighborhood of \(\eta\):

\[ 1 \longrightarrow L \longrightarrow C \longrightarrow A \longrightarrow 1. \]

Moreover, one may assume that \(C\) is a Cartan subgroup of \(G\) (\SGARefLink{sga3:xv:corollary:6-3}{6.3}) and that the maximal torus \(T_{\eta}\) of \(L_{\eta}\)
extends to a maximal torus \(T\) of \(L\) (\SGARefLink{sga3:xv:corollary:6-3}{6.3}). The corollary follows immediately from this and from the definitions.

\SGARefTarget{sga3:xv:subsection:6:4:0}
\textbf{\SGAInlineRefTarget{sga3:xv:paragraph:6-4-0}6.4.0.} Let \(S\) be a prescheme, \(G\) an \(S\)-prescheme in groups of finite
presentation, and \(\mathcal{L}\) the functor of subgroups of \(G\) that are smooth,
have connected fibres, and are equal to their connected normalizer (cf.~§5).
The functor \(\mathcal{L}\) is representable (\SGARefLink{sga3:xv:theorem:5-2}{5.2} and \SGARefLink{sga3:xv:remark:5:8}{5.8}). The rest of this section
is devoted to certain subfunctors of \(\mathcal{L}\). More precisely, define
\(\mathcal{L}_{C}\) (resp. \(\mathcal{CT}\), resp. \(\mathcal{P}\)) as follows:
for every \(S\)-prescheme \(S'\), the set \(\mathcal{L}_{C}(S')\)
(resp. \(\mathcal{CT}(S')\), resp. \(\mathcal{P}(S')\)) consists of the subgroups
\(H\) of \(G_{S'}\) that are smooth over \(S'\), have connected fibres, are equal
to their connected normalizer, and are such that, for every point \(s'\) of \(S'\),
\(H_{s'}\) contains a Cartan subgroup (\SGARefLink{sga3:xv:definition:6-1}{6.1}) of \(G_{s'}\)
(resp. is the centralizer in \((G_{s'})^{0}_{\mathrm{red}}\) of a subtorus of
\(G_{s'}\), resp. is a parabolic subgroup (\SGARefLink{sga3:xv:definition:6-1}{6.1}) of \(G_{s'}\)).

\SGARefTarget{sga3:xv:theorem:6:4}
\textbf{Theorem 6.4.} \emph{Let \(S\) be a prescheme and \(G\) an \(S\)-prescheme in
groups of finite presentation. The \(S\)-functors \(\mathcal{L}_{C}\),
\(\mathcal{CT}\), and \(\mathcal{P}\) defined in \SGARefLink{sga3:xv:paragraph:6-4-0}{6.4.0} are representable by
\(S\)-preschemes that are quasi-projective and of finite presentation over \(S\).}

\SGARefTarget{sga3:xv:remark:6:5}
\textbf{Remark 6.5.} If \(G\) has smooth fibres---for example, if \(G\) is smooth
over \(S\), or if the residue characteristics of \(S\) are zero
(Exposé VI\textsubscript{B}, \SGARefLink{cor:sga3-VIB-1.6.1}{1.6.1})---then every subgroup \(H\) of \(G\), smooth over
\(S\) with connected fibres, such that \(H_s\) contains a Cartan subgroup of
\((G_s)^0\) for every point \(s\) of \(S\), is necessarily equal to its connected
normalizer (and hence is an element of \(\mathcal{L}_{C}(S)\)). Indeed, by \SGARefLink{sga3:xv:proposition:5-1}{5.1} iii)
we may assume that \(S\) is the spectrum of a field, in which case the property was
noted at the end of the statement of \SGARefLink{sga3:xiii:theorem:2:1}{Exposé XIII, 2.1}.

There are natural monomorphisms
\[
\begin{tikzcd}[column sep=huge,row sep=large]
\mathcal{CT} \arrow[dr,hook] & \\
& \mathcal{L}_{C} \arrow[r,hook] & \mathcal{L}.\\
\mathcal{P} \arrow[ur,hook] &
\end{tikzcd}
\]

Let us show that these monomorphisms are open immersions of finite presentation (which will already prove, in view of
\SGARefLink{sga3:xv:theorem:5-2}{5.2}, that \(\mathcal{L}_{C}\), \(\mathcal{CT}\), \(\mathcal{P}\) are representable by \(S\)-preschemes, union of an increasing
sequence of open subpreschemes, quasi-projective and of finite presentation over \(S\)). Now this will follow, by the
usual technique, from the following lemma:

\SGARefTarget{sga3:xv:lemma:6:6}
\textbf{Lemma 6.6.} \emph{Let \(S\) be a prescheme, \(G\) an \(S\)-prescheme in
groups of finite presentation, and \(H\) a subgroup of \(G\), smooth with connected
fibres. The set of points \(s\) of \(S\) such that \(H_s\) contains a Cartan subgroup
of \(G_s\) (resp. is equal to the centralizer in \((G_s)^0_{\mathrm{red}}\) of a torus
of \(G_s\)) is open in \(S\). If, moreover, \(H\) is equal to its connected
normalizer,\footnote{This hypothesis is in fact superfluous; see Exposé XVII,
Appendix III, 3.} then the set of points \(s\) of \(S\) such that \(H_s\) is a
parabolic subgroup of \(G_s\) is also open in \(S\).}

The assertion to be proved is local on \(S\), which allows us to assume \(S\) affine, then, by EGA IV 8.1 and 6.3, \(S\)
noetherian. Denote by \(E\) the set of points of \(S\) having the property in question. It then follows from the assertions
effectively proved in 6.2 that \(E\) is ind-constructible. But \(S\) is noetherian, so to prove that \(E\) is open, it
suffices to show that \(E\) is stable under generizations (EGA IV 1.10.1). Using EGA II 7.1.9, we are finally led to prove
that if \(S\) is the spectrum of a discrete valuation ring, and if the closed point \(s\) belongs to \(E\), then so does the
generic point \(t\).

We shall need the following lemma:

\SGARefTarget{sga3:xv:lemma:6:7}
\textbf{Lemma 6.7.} \emph{Let \(A\) be a complete noetherian local ring,
\(S=\operatorname{Spec}A\), \(s\) the closed point of \(S\), \(H\) an
\(S\)-prescheme in groups, smooth with connected fibres, and \(T_s\) a subtorus of
\(H_s\). Then:
\begin{enumerate}
\renewcommand{\labelenumi}{\roman{enumi})}
\item there exists a closed subgroup prescheme \(C\) of \(H\), smooth with connected
fibres, such that \(C_s=\operatorname{Centr}_{H_s}(T_s)\);
\item for every point \(t\) of \(S\), \(C_t\) is the centralizer in \(H_t\) of a
subtorus \(T_t\) of \(H_t\).
\end{enumerate}}

\emph{Proof of \SGARefLink{sga3:xv:lemma:6:7}{6.7}.} Let \(T'\) be an \(S\)-torus for which there is an isomorphism
\(u_0:T'_s\xrightarrow{\sim}T_s\) (\SGARefLink{sga3:x:corollary:4-6}{Exposé X, 4.6}). Let \(\mathfrak m\) be the
maximal ideal of \(A\), and put \(A_n=A/\mathfrak m^n\),
\(S_n=\operatorname{Spec}A_n\), \(H_n=H\times_S S_n\), and so on. Since \(H\) is
smooth over \(S\), for every integer \(n>0\) there is an \(S_n\)-group morphism
\[
u_n:T'_n\longrightarrow H_n
\]
lifting \(u_0\) (\SGARefLink{sga3:ix:theorem:3:6}{Exposé IX, 3.6}), and by induction on \(n\) we may assume that
\(u_n\) lifts \(u_{n-1}\). Let \(q\) be a prime invertible on \(S\). For every
integer \(\ell\) that is a power of \(q\), write \({}_{\ell}u_n\) for the restriction
of \(u_n\) to \({}_{\ell}T'_n\). For fixed \(\ell\) and varying \(n\), the morphisms
\({}_{\ell}u_n\) form a projective system and therefore arise from a unique
\(S\)-group morphism
\[
{}_{\ell}u:{}_{\ell}T'\longrightarrow H
\]
(\SGARefLink{sga3:xv:proposition:1-6}{1.6} a)). Since \(H\) is separated (Exposé VI\textsubscript{B}, \SGARefLink{prop:sga3-VIB-5.2}{5.2}) and \(u_0\) is
a monomorphism, \({}_{\ell}u\) is a monomorphism (\SGARefLink{sga3:ix:theorem:6:8}{Exposé IX, 6.8}), and even a closed
immersion because it is finite, \({}_{\ell}T'\) being finite over \(S\). Denote its
image by \(M(\ell)\). The family of subgroups of multiplicative type \(M(\ell)\) is
clearly coherent in the sense of \SGARefLink{sga3:xv:theorem:4-1}{4.1}. Put
\[
C_\ell=\operatorname{Centr}_H(M(\ell)).
\]
This is represented by a closed subgroup prescheme (\SGARefLink{sga3:xv:lemma:2:5}{2.5}), since \(H\) is separated,
and is smooth over \(S\) (\SGARefLink{sga3:xi:corollary:2:4}{Exposé XI, 2.4}). The \(C_\ell\) form a filtered decreasing
family of closed subschemes and are therefore stationary, since \(H\) is noetherian.
Their stationary value is a smooth closed subgroup \(C\) such that
\(C_s=\operatorname{Centr}_{H_s}(T_s)\), by the density theorem (\SGARefLink{sga3:ix:theorem:4:7}{Exposé IX, 4.7}).
It remains to show that for every point \(t\) of \(S\), \(C_t\) is the centralizer
in \(H_t\) of a subtorus \(T_t\); this will also imply that \(C_t\) is connected.
This follows from the more precise lemma below, applied to the family
\(M(\ell)_t\) of subgroups of \(H_t\).

\SGARefTarget{sga3:xv:lemma:6:8}
\textbf{Lemma 6.8.} \emph{Let \(G\) be a connected algebraic group over a field
\(k\), let \(r>0\) be an integer, and let \(q\) be an integer prime to the
characteristic\footnote{Editor's note: the word ``residual'' has been removed.}
of \(k\). Let \(M(\ell)\), where \(\ell\) ranges over the powers of \(q\), be a
coherent family of subgroups of \(G\), of multiplicative type and of type
\((\mathbb Z/\ell\mathbb Z)^r\) (cf.~\SGARefLink{sga3:xv:lemma:4:6}{4.6}). Let \(M\) be the algebraic subgroup of
\(G\) generated by the \(M(\ell)\) (\emph{loc.~cit.}), and let \(T\) be the unique
maximal torus of \(M\) (cf.~\SGARefLink{sga3:xv:lemma:3:4}{3.4}). Then, for all sufficiently large \(\ell\),
\[
\operatorname{Centr}_G(T)=\operatorname{Centr}_G(M)
  =\operatorname{Centr}_G(M(\ell)).
\]}

The last equality is clear. To prove the first, introduce the centre \(Z\) of
\(G\), put \(G'=G/Z\), let \(M'\) (resp. \(T'\)) be the image of \(M\) (resp.
\(T\)) in \(G'\), and let \(K\) be the inverse image of \(T'\) in \(G\), that is,
the algebraic subgroup of \(G\) generated by \(T\) and \(Z\). It plainly suffices
to prove \(K\supset M\), or equivalently \(T'=M'\). Now \(M\) is smooth and
connected (\SGARefLink{sga3:xv:lemma:4:6}{4.6}), while \(G'\) is affine (\SGARefLink{sga3:xii:6.1}{Exposé XII, 6.1}); hence \(M'\) is the
direct product of its maximal torus \(T'\) (\SGARefLink{sga3:xii:6.6}{Exposé XII, 6.6} d)) and a unipotent
group (\emph{Bible} 4, Th.~4), after assuming \(k\) algebraically closed. The image
of \(M(\ell)\) in \(M'\) is therefore contained in \(T'\). Thus the inverse image
of \(T'\) in \(M\) contains \(M(\ell)\) for every \(\ell\), and consequently is
equal to \(M\). Hence \(M'=T'\). This proves \SGARefLink{sga3:xv:lemma:6:8}{6.8}, and therefore \SGARefLink{sga3:xv:lemma:6:7}{6.7}.

We now prove \SGARefLink{sga3:xv:lemma:6:6}{6.6}. We have reduced to the case where \(S\) is the spectrum of a
discrete valuation ring \(A\), which we may moreover assume complete with
algebraically closed residue field. After replacing \(A\), if necessary, by its
normalization in a finite extension of its fraction field, we may assume that
\((G_t)_{\mathrm{red}}\) is smooth.\footnote{Editor's note: details or references
should be supplied here.} To prove \SGARefLink{sga3:xv:lemma:6:6}{6.6} we may clearly replace \(G\) by the identity
component of the schematic closure in \(G\) of \((G_t)^0_{\mathrm{red}}\), and
therefore assume that \(G\) is flat over \(S\), has connected fibres, and has smooth
generic fibre \(G_t\).

\begin{enumerate}
\renewcommand{\labelenumi}{\alph{enumi})}
\item Suppose that \(H_s\) is the centralizer in \((G_s)_{\mathrm{red}}\) of a torus
\(T_s\). We show that \(H_t\) is then the centralizer in \(G_t\) of a subtorus of
\(G_t\). By \SGARefLink{sga3:xv:lemma:6:7}{Lemma 6.7} there is a subgroup scheme \(C\) of \(H\), smooth over \(S\),
whose closed fibre is \(H_s\), such that
\(C_t=\operatorname{Centr}_{H_t}(T_t)\) for a subtorus \(T_t\) of \(H_t\). Since
\(H\) is smooth over \(S\) with connected fibres, dimension considerations give
\(C=H\). With the notation of \SGARefLink{sga3:xv:lemma:6:7}{6.7}, \(H=C=C_\ell\) for all sufficiently large
\(\ell\). Similarly, put
\(C'_\ell=\operatorname{Centr}_G(M(\ell))\) (\SGARefLink{sga3:xv:lemma:2:5}{2.5}), and let \(C'\) be the stationary
value of the \(C'_\ell\) for large \(\ell\) (\SGARefLink{sga3:xv:lemma:2-5-bis}{2.5 bis}). The group scheme \(C'\)
contains \(H\), and
\[
C'_t=\operatorname{Centr}_{G_t}(T_t)\quad\text{(6.8)},\qquad
C'_s=\operatorname{Centr}_{G_s}(T_s).
\]
The hypothesis on \(H_s\) gives \(\dim H_s=\dim C'_s\). Moreover,
\(\dim H_s=\dim H_t\), since \(H\) is smooth over \(S\), and
\(\dim C'_t\leq \dim C'_s\) (Exposé VI\textsubscript{B}, \SGARefLink{prop:sga3-VIB-4.1}{4.1}). Hence
\(\dim H_t=\dim C'_t\). Since \(G_t\) is smooth and connected,
\(C'_t=\operatorname{Centr}_{G_t}(T_t)\) is smooth and connected; finally,
\[
H_t=C'_t=\operatorname{Centr}_{G_t}(T_t).
\]

\item Suppose that \(H_s\) contains a Cartan subgroup \(C_s\) of \(G_s\), that is,
the centralizer in \((G_s)_{\mathrm{red}}\) of a maximal torus \(T_s\) of \(G_s\).
By \SGARefLink{sga3:xv:lemma:6:7}{6.7} there is a subgroup scheme \(C\) of \(H\), smooth over \(S\) with connected
fibres, lifting \(C_s\). It suffices to prove that \(C_t\) contains a Cartan subgroup
of \(G_t\). By part a), applied with \(H=C\), \(C_t\) is the centralizer in \(G_t\)
of a subtorus of \(G_t\), and hence contains a Cartan subgroup of \(G_t\).

\item Suppose that \(H_s\) is a parabolic subgroup of \(G_s\). Let
\(N=\operatorname{Norm}_G(H)\). By hypothesis \(H\) is equal to its connected
normalizer, so \(N_s\) is smooth and consequently
\[
N_s=\operatorname{Norm}_{(G_s)_{\mathrm{red}}}(H_s)=H_s
\qquad\text{(\SGARefLink{sga3:xii:7.10}{Exposé XII, 8 bis}).}
\]
Thus \(N\) is flat over \(S\). We shall see in Exposé XVI that under these
conditions \(G/N\) is representable. Since \(H_s=N_s\) is a parabolic subgroup,
\((G/N)_s\) is proper. Since \(G/N\) is flat over \(S\) with connected fibres,
EGA III, 5.5.1 implies that \(G/N\) is proper over \(S\). Thus
\((G/N)_t=G_t/N_t\) is proper over \(t\), and so is \(G_t/H_t\), because
\(N_t/H_t\) is finite. It remains only to apply the following lemma.
\end{enumerate}

\SGARefTarget{sga3:xv:lemma:6:9}
\textbf{Lemma 6.9.} \emph{Let \(k\) be a field, \(G\) a \(k\)-algebraic group,
\(H\) a smooth connected algebraic subgroup of \(G\), and
\(N=\operatorname{Norm}_G(H)\). If \(\dim H=\dim N\) and \(G/H\) is proper, then
\(H\) is a parabolic subgroup of \(G\).}

Indeed, we may assume that \(k\) is algebraically closed and that \(G\) is smooth
and connected. The centre \(Z\) of \(G\) is contained in \(N\), and
\(\dim H=\dim N\) implies that \(Z'=(Z)^0_{\mathrm{red}}\) is contained in \(H\).
Hence \(G'=G/Z'\) is affine (\SGARefLink{sga3:xii:6.1}{Exposé XII, 6.1}). Replacing \(G\) by \(G'\) and \(H\)
by its image in \(G'\), we reduce to the case in which \(G\) is affine
(\SGARefLink{sga3:xiv:4.9}{Exposé XIV, 4.9}), and \SGARefLink{sga3:xv:lemma:6:9}{Lemma 6.9} then follows from \emph{Bible} 6, Th.~4.
This proves \SGARefLink{sga3:xv:lemma:6:6}{Lemma 6.6}.

To complete the proof of \SGARefLink{sga3:xv:theorem:6:4}{6.4}, we must show that the \(S\)-preschemes representing
\(\mathcal L_C\) (respectively \(\mathcal{CT}\) and \(\mathcal P\)) are of finite
presentation over \(S\). This assertion is local on \(S\), so we may assume first
that \(S\) is affine and then, since \(G\) is of finite presentation, that \(S\) is
noetherian (EGA IV, 8.9). We have just seen that the natural inclusions
\(\mathcal{CT}\to\mathcal L\) and \(\mathcal P\to\mathcal L\) are immersions.
The same is therefore true of
\(\mathcal{CT}\to\mathcal L_C\) and \(\mathcal P\to\mathcal L_C\).
Consequently, it is enough to prove that \(\mathcal L_C\) is represented by an
\(S\)-prescheme of finite presentation.

Resume the notation introduced in \SGARefLink{sga3:xv:theorem:5-2}{5.2}. For every integer \(n\geq 0\), let
\(\mathcal L^n\) be the subfunctor of \(\mathcal L\) defined by
\[
\mathcal L^n(S')
 =\left\{H\in\mathcal L(S')\ \middle|\ 
   \operatorname{Norm}_{G_{S'}}(H)
   =\operatorname{Norm}_{G_{S'}}\bigl(H^{(n)}\bigr)\right\}.
\]
The \(S\)-functor \(\mathcal L^n\) is thus represented by an open subprescheme
of \(\mathcal L\), the disjoint union of the \(\mathcal L_r^n\).
Each \(\mathcal L_r^n\) is of finite presentation over \(S\) (\SGARefLink{sga3:xv:lemma:5:3}{5.3}), and is empty
for
\[
r>\sup_{s\in S}\dim G_s,
\]
which is a finite number because \(S\) is quasi-compact. Hence \(\mathcal L^n\)
is of finite presentation over \(S\). It remains to prove that
\(\mathcal L_C\) is contained in \(\mathcal L^n\) for all sufficiently large
\(n\).

For each point \(s\) of \(S\), let \(d(s)\) be the least integer \(n\), finite or
infinite, such that
\[
\mathcal L_{C,G_s}\subset\mathcal L^n_{G_s}.
\]
It is enough to show that \(d\) is bounded on \(S\). Indeed, if \(M\) is an
upper bound, then \(\mathcal L_C\) is set-theoretically contained in
\(\mathcal L^M\), and hence is contained in \(\mathcal L^M\), since the latter
is open in \(\mathcal L\). A direct constructibility argument reduces us to
showing that, when \(S\) is noetherian and integral with generic point \(t\),
the function \(d\) is bounded on a neighbourhood of \(t\).

\begin{enumerate}
\renewcommand{\labelenumi}{\alph{enumi})}
\item \textbf{Reduction to the case in which \(G\) is smooth over \(S\).}
Proceeding as in \SGARefLink{sga3:xv:proposition:6-2}{6.2}, we see that, after possibly changing \(S\), we may assume
that \(G_{\mathrm{red}}\) is a group prescheme smooth over \(S\); denote it by
\(G'\). Put
\[
X=\mathcal L_{C,G},\qquad X'=\mathcal L_{C,G'},
\]
and let \(H\) (respectively \(H'\)) be the subgroup prescheme of \(G_X\)
(respectively \(G'_{X'}\)) universal for the functor
\(\mathcal L_{C,G}\) (respectively \(\mathcal L_{C,G'}\)).
Since \(H\) is smooth over \(X\), the prescheme
\(H\times_X X_{\mathrm{red}}\) is reduced, and is therefore contained in
\(G'_{X_{\mathrm{red}}}\). It is an element of
\(\mathcal L_{C,G'}(X_{\mathrm{red}})\), and hence defines a canonical
morphism
\[
p:X_{\mathrm{red}}\longrightarrow X'.
\]
It is clear that \(p\) is a monomorphism; we show that it is in fact an
immersion. Let \(N'\) be the normalizer of \(H'\) in \(G_{X'}\).
The set of points \(s\in X'\) for which
\(H'_s\to N'_s\) is an open immersion is an open subset \(U\), and
\[
H'|_U\longrightarrow N'|_U
\]
is an open immersion (Exposé VI\textsubscript{B}, 2.5, and EGA IV, 17.9.5).
It follows from 5.1(iii) that \(U\) is the largest open subset of \(X'\)
over which \(H'\) equals its connected normalizer in \(G_{X'}\).
Thus \(H'|_U\in\mathcal L_{C,G}(U)\), and it follows at once that \(p\)
is an isomorphism of \(X_{\mathrm{red}}\) onto \(U_{\mathrm{red}}\).
If \(\mathcal L_{C,G}\) is known to be of finite type when \(G\) is smooth,
then \(X'\) is of finite type over \(S\); consequently
\(X_{\mathrm{red}}\) is of finite type (since \(S\) is noetherian), and is
therefore contained in \(\mathcal L_G^M\) for all sufficiently large \(M\).
The same then holds for \(X\).

\item \textbf{The case in which \(S\) is the spectrum of an algebraically
closed field \(k\) of characteristic \(p>0\), and \(G\) is a smooth
algebraic group over \(k\).}
Instead of the infinitesimal neighbourhoods \(H^{(n)}\) of a subgroup
prescheme \(H\) of \(G\), we use the radicial subgroups \(F_n(H)\), the
kernels of the iterates of the Frobenius morphism of \(H\)
(Exposé VII\textsubscript{A}, 4). This is legitimate here because \(F_n(H)\)
is contained in the infinitesimal neighbourhood of order \(p^n\) of the
unit section of \(H\). If \(T\) is a subtorus of \(G\), then
\[
F_n(T)={}_{p^n}T,
\]
and duality immediately shows that
\(\operatorname{Centr}_G({}_{p^n}T)=\operatorname{Centr}_G(T)\) for all
sufficiently large \(n\). We have the following more precise proposition.
\end{enumerate}

\SGARefTarget{sga3:xv:proposition:6:10}
\textbf{Proposition 6.10.} \emph{Let \(k\) be a field of characteristic
\(p>0\), let \(G\) be a smooth \(k\)-algebraic group, let \(T\) be a maximal
torus of \(G_{\bar k}\), and let \(m\) be the least integer such that
\[
\operatorname{Centr}_G({}_{p^m}T)^0
 =\operatorname{Centr}_G(T)^0.
\]
Then, for every \(k\)-prescheme \(S\) and every
\(H\in\mathcal L_C(S)\), one has
\[
\operatorname{Norm}_{G_S}(H)
 =\operatorname{Norm}_{G_S}\bigl(F_m(H)\bigr).
\]
In particular, \(\mathcal L_C\) is contained in \(\mathcal L^{p^m}\).}

Since \(H\) is smooth and of finite presentation over \(S\), and has constant
nilpotent rank, namely that of \(G\), we shall see in the next section (\SGARefLink{sga3:xv:theorem:7-3}{7.3})
that the functor \(\mathcal C_H\) of Cartan subgroups of \(H\) is represented
by an \(S\)-prescheme smooth over \(S\). The reader will verify that the proof
given there does not use the finite-type property of \(\mathcal L_{C,H}\).
It follows from \SGARefLink{sga3:xiii:theorem:3:1}{Exposé XIII, 3.1}, that we may consider the open subset \(U\)
of regular points of \(H\).

Let \(S'\) be an \(S\)-prescheme and let \(g\in G(S')\) normalize
\(F_m(H)_{S'}\). To prove that \(g\) normalizes \(H_{S'}\), it is enough to
prove that
\[
\operatorname{int}(g)U_{S'}\subset H_{S'}.
\]
Indeed, \(H_{S'}\cap\operatorname{int}(g)H_{S'}\) would then contain an open
subgroup of \(\operatorname{int}(g)H_{S'}\), and hence would equal it, since
the latter has connected fibres. After replacing \(S'\) by a suitable
covering \(S''\), and then replacing \(S''\) by \(S\), we are reduced to
proving the following: if \(g\in G(S)\) normalizes \(F_m(H)\) and
\(u\in U(S)\), then \(\operatorname{int}(g)u\in H(S)\).

Let \(C\) be the unique Cartan subgroup of \(H\) ``containing'' \(u\)
(\SGARefLink{sga3:xiii:corollary:3:2}{Exposé XIII, 3.2}). It is enough to show that
\[
C'=\operatorname{int}(g)C\subset H.
\]
Since \(H\in\mathcal L_{C,G}(S)\), the group \(C\) is also a Cartan subgroup
of \(G\). The latter admits maximal tori, so \(C\) is the centralizer in
\(G_S^0\) of its unique maximal torus \(T\) (\SGARefLink{sga3:xii:7.1:item:a}{Exposé XII, 7.1(a)}).
By the definition of \(m\), and because any two maximal tori of \(G\) are
locally conjugate for the fpqc topology (\SGARefLink{sga3:xii:7.1}{Exposé XII, 7.1}),
\[
C=\operatorname{Centr}_G(T)^0
 =\operatorname{Centr}_G({}_{p^m}T)^0.
\]
Conjugating by \(g\), we obtain
\[
C'=\operatorname{Centr}_G\bigl(\operatorname{int}(g)({}_{p^m}T)\bigr)^0.
\]
Now \(\operatorname{int}(g)({}_{p^m}T)\) is a subgroup of multiplicative
type of
\[
\operatorname{int}(g)\bigl(F_m(H)\bigr)=F_m(H),
\]
because \(g\) normalizes \(F_m(H)\); it is therefore contained in \(H\).
\SGARefLink{sga3:xiii:theorem:2:1}{Exposé XIII, 2.1} (proved there when the base is a field, but extending at
once to an arbitrary base) then shows that, in order to prove
\(C'\subset H\), it is enough to prove
\[
\operatorname{Lie}C'\subset\operatorname{Lie}H.
\]
But
\[
\operatorname{Lie}C'
 =\operatorname{int}(g)(\operatorname{Lie}C)
 \subset\operatorname{int}(g)(\operatorname{Lie}H).
\]
On the other hand, we may assume \(m\geq1\), and Exposé
VII\textsubscript{A}, \SGARefLink{sga3:VIIA:4.1.2}{4.1.2}, gives
\[
\operatorname{Lie}H
 =\operatorname{Lie}\bigl(F_m(H)\bigr)
 =\operatorname{Lie}\bigl(\operatorname{int}(g)(F_m(H))\bigr)
 =\operatorname{int}(g)(\operatorname{Lie}H).
\]
Thus \(\operatorname{Lie}C'\subset\operatorname{Lie}H\), which completes
the proof of \SGARefLink{sga3:xv:proposition:6:10}{6.10}.

We shall need another description of the integer \(m\) introduced in \SGARefLink{sga3:xv:proposition:6:10}{6.10}.

\SGARefTarget{sga3:xv:lemma:6:11}
\textbf{Lemma 6.11.} \emph{Let \(G\) be a smooth algebraic group over an
algebraically closed field \(k\) of characteristic \(p>0\), let \(T\) be a
maximal torus of \(G\), put \(\mathfrak g=\operatorname{Lie}G\), and let \(R\)
be the set of nonzero characters of \(T\) occurring in the representation of
\(T\) on \(\mathfrak g\) induced by the adjoint representation of \(G\).
For \(r\in R\), let \(e_r\) be the greatest integer \(n\) such that \(p^n\)
divides \(r\) in the character group of \(T\). Then
\[
m=\sup_{r\in R}(e_r+1)
\]
if \(R\neq\varnothing\), and \(m=0\) otherwise.}

Indeed, \(\operatorname{Centr}_G(T)\) is smooth and is contained in
\(\operatorname{Centr}_G({}_{p^m}T)\). Hence
\[
\begin{aligned}
\operatorname{Centr}_G(T)^0
 &=\operatorname{Centr}_G({}_{p^m}T)^0\\
&\Longleftrightarrow\
\operatorname{Lie}\bigl(\operatorname{Centr}_G(T)\bigr)
 =\operatorname{Lie}\bigl(\operatorname{Centr}_G(F_m(T))\bigr)\\
&\Longleftrightarrow\
\mathfrak g^T=\mathfrak g^{{}_{p^m}T}
\qquad\text{(Exposé II, 5.2.3).}
\end{aligned}
\]
With the usual notation,
\[
\mathfrak g=\mathfrak g_0\oplus\bigoplus_{r\in R}\mathfrak g_r.
\]
Thus \(\mathfrak g^T=\mathfrak g_0\), whereas
\[
\mathfrak g^{{}_{p^m}T}
 =\mathfrak g_0\oplus\bigoplus_{r\in R'}\mathfrak g_r,
\]
where \(R'\) is the subset of \(R\) formed by the characters whose restriction
to \({}_{p^m}T\) is zero. A nonzero character \(r\) of \(T\) has trivial
restriction to \({}_{p^m}T\) if and only if \(m\leq e_r\), which proves the
lemma.

\begin{enumerate}
\renewcommand{\labelenumi}{\alph{enumi})}
\setcounter{enumi}{2}
\item \textbf{Return to the proof of 6.4.}
By part a) and the paragraph preceding it, we have reduced to the case in
which \(S\) is a noetherian integral scheme and \(G\) is smooth over \(S\).
We must show that \(d\) is bounded on a neighbourhood of the generic point
\(t\) of \(S\). After possibly changing \(S\), we may assume that \(G\) has
a split maximal torus (\SGARefLink{sga3:xv:proposition:6-2}{6.2})
\[
T=D_S(M).
\]
Let
\[
\mathfrak g=\bigoplus_{\lambda\in M}\mathfrak g^\lambda
\]
be the decomposition of the Lie algebra of \(G\) according to the characters
of \(T\), and let \(R\) be the finite set of nonzero \(\lambda\in M\) for
which \(\mathfrak g^\lambda\neq0\). We distinguish two cases.
\end{enumerate}

\textbf{First case: \(t\), and hence every point of \(S\), has positive
residue characteristic \(p\).}
It is clear from the foregoing that \(d\) is then bounded by \(p^m\), where
\(m\) is defined as in \SGARefLink{sga3:xv:lemma:6:11}{6.11}.

\textbf{Second case: \(t\) has residue characteristic zero.}
For each \(\lambda\in R\), let \(n_\lambda\) be the greatest integer dividing
\(\lambda\) in the group \(M\), and put
\[
n=\prod_{\lambda\in R}n_\lambda.
\]
For every prime number \(q\) dividing \(n\), let \(S_q\) be the closed subset
of \(S\) formed by the points of residue characteristic \(q\), and let \(U\)
be the nonempty open subset, containing \(t\), complementary to the union of
the \(S_q\). If \(s\in U\), then either \(s\) has residue characteristic zero,
in which case \(d(s)=1\) by \SGARefLink{sga3:xv:remark:5:6}{5.6}, or \(s\) has residue characteristic
\(p>0\) not dividing \(n\). In the latter case the integer \(m\) attached to
\(G_s\) in \SGARefLink{sga3:xv:lemma:6:11}{6.11} is at most one. Moreover, Exposé
VII\footnote{Editor's note: the argument should be made explicit.} gives
\[
\operatorname{Norm}_{G_{S'}}\bigl(F_1(H)\bigr)
 =\operatorname{Norm}_{G_{S'}}(\operatorname{Lie}H)
 =\operatorname{Norm}_{G_{S'}}\bigl(H^{(1)}\bigr).
\]
Finally, \SGARefLink{sga3:xv:proposition:6:10}{6.10} shows that, for
\(H\in\mathcal L_{C,G_s}(S')\),
\[
\operatorname{Norm}_{G_{S'}}(H)
 =\operatorname{Norm}_{G_{S'}}\bigl(H^{(1)}\bigr).
\]
Consequently \(d(s)\leq1\), so \(d\) is bounded by \(1\) on \(U\).
This completes the proof of \SGARefLink{sga3:xv:theorem:6:4}{6.4}.

\SGARefTarget{sga3:xv:corollary:6:12}
\textbf{Corollary 6.12.} \emph{Let \(S\) be a prescheme and \(G\) an
\(S\)-prescheme in groups of finite presentation. Suppose that the nilpotent
rank of the fibres of \(G\) (respectively the dimension of their Borel
subgroups) is locally constant on \(S\). Then the functor \(\mathcal C\) of
Cartan subgroups of \(G\) (respectively the functor \(\mathcal B\) of Borel
subgroups of \(G\)) is represented by an \(S\)-prescheme of finite
presentation over \(S\).}

Indeed, after restricting \(S\), we may assume that the nilpotent rank
\(\nu\) of the fibres is constant. Then \(\mathcal C\) is represented by the
open-and-closed subprescheme of the prescheme \(X\) representing
\(\mathcal L_C\) (\SGARefLink{sga3:xv:theorem:6:4}{6.4}) over which the universal subgroup of \(G_X\) for the
functor \(\mathcal L_C\) has relative dimension \(\nu\). The proof is
analogous for \(\mathcal B\), using the representability of \(\mathcal P\).

% END INCLUDED SGA3_English_current_buildable_source_20260823/source/inputs/XV/tex/components/06_expose_XV_section6.tex

% BEGIN INCLUDED SGA3_English_current_buildable_source_20260823/source/inputs/XV/tex/components/07_expose_XV_section7.tex
% Authority: Expo15.pdf, local pp. 50--56 / combined pp. 920--926 /
% printed pp. 445--458.
% Stops immediately before §8.

\SGASection{7}{Cartan subgroups of a smooth group}
\SGARefTarget{sga3:xv:section:7}

\SGAStatement{\SGAInlineRefTarget{sga3:xv:proposition:7-1}Proposition 7.1}{Let \(S\) be a prescheme and let \(G\)
be an \(S\)-group prescheme of finite presentation, smooth over \(S\),
with connected fibres.
\begin{enumerate}
\renewcommand{\labelenumi}{\textup{\roman{enumi})}}
\item\SGAInlineRefTarget{sga3:xv:proposition:7-1:item:i} Let \(\mathcal{CT}\) be the \(S\)-functor defined by
\[
\mathcal{CT}(S')=
\left\{
\begin{array}{c}
\text{subgroup preschemes \(H\) of \(G_{S'}\), smooth over \(S'\),}\\
\text{such that every \(H_{s'}\) is the centralizer in \(G_{s'}\)}\\
\text{of a sub-torus of \(G_{s'}\)}
\end{array}
\right\}.
\]
Then \(\mathcal{CT}\) is represented by an \(S\)-prescheme of finite
presentation, smooth and quasi-projective over \(S\).
\item If \(S\) is artinian and \(H\in\mathcal{CT}(S)\), then \(H\)
is the centralizer in \(G\) of a sub-torus of \(G\). If \(S\) is the
spectrum of a henselian local ring and \(H\in\mathcal{CT}(S)\), then
\(H\) is the centralizer in \(G\) of a subgroup of multiplicative
type that is étale over \(S\).
\item\SGAInlineRefTarget{sga3:xv:proposition:7-1:item:iii} Let \(L\) be an \(S\)-group prescheme smooth and of finite
presentation, let \(i:L\to G\) be an \(S\)-group monomorphism, and
let \(H\in\mathcal{CT}(S)\). Then
\(\operatorname{Transp}_G(H,L)\) and
\(\operatorname{Transpstr}_G(H,L)\)
(Exposé \SGARefLink{sga3:viii:examples:6:5:item:e:01}{VIII~6.5(e)}) are represented by closed subpreschemes of \(G\),
smooth over \(S\).
\item If \(H\in\mathcal{CT}(S)\), then \(H\) is closed in \(G\);
\(N=\operatorname{Norm}_G(H)\) is represented by a closed subgroup
prescheme of \(G\), smooth over \(S\); \(N/H\) is represented by an
\(S\)-group prescheme, separated, étale, and of finite type over
\(S\); and \(G/N\) is represented by an \(S\)-prescheme smooth and
quasi-projective over \(S\).
\item\SGAInlineRefTarget{sga3:xv:proposition:7-1:item:v} Let \(G'\) be an \(S\)-group prescheme of finite presentation
and let \(u:G\to G'\) be faithfully flat. Then \(G'\) satisfies the
same hypotheses as \(G\) (Exposé VI\textsubscript{B}~9). If
\(H\in\mathcal{CT}_G(S)\), its image is represented by a subgroup
prescheme \(H'\in\mathcal{CT}_{G'}(S)\), and \(H\to H'\) is faithfully
flat. If \(H=\operatorname{Centr}_G(T)\) for a torus \(T\), then
\[
 H'=\operatorname{Centr}_{G'}(T'),\qquad T'=u(T).
\]
\item\SGAInlineRefTarget{sga3:xv:proposition:7-1:item:vi} Under the hypotheses of (v), the \(S\)-morphism
\[
 \widetilde u:\mathcal{CT}_G\longrightarrow\mathcal{CT}_{G'},
 \qquad H\longmapsto H'=u(H),
\]
is quasi-compact and faithfully flat. If, moreover, \(\ker u\) is
central, then \(\widetilde u\) is an isomorphism, with inverse
\(H'\mapsto u^{-1}(H')\).
\end{enumerate}}
\SGARefTarget{sga3:xv:proposition:7:1}

\medskip\noindent\textit{Proof of (ii).}
For the first assertion we may suppose that \(S\) is local artinian,
with closed point \(s\). Let \(H\in\mathcal{CT}(S)\), and let \(T_s\)
be the maximal central torus of \(H_s\), which is defined over
\(\kappa(s)\) (cf. \SGARefLink{sga3:xv:lemma:3:4}{Lemma 3.4}). Since
\[
  H_s=\operatorname{Centr}_{G_s}(T_s),
\]
and \(H\) is smooth, \(T_s\) lifts uniquely to a sub-torus \(T\) of
\(H\), central in \(H\) (Exposé \SGARefLink{sga3:ix:theorem:3:6:bis}{IX~3.6 bis} and \SGARefLink{sga3:xv:remark:5:6}{5.6}). Then
\(H'=\operatorname{Centr}_G(T)\) contains \(H\) and has the same
fibre as \(H\). Because \(H\) is flat, \(H=H'\) (Exposé
VI\textsubscript{B}~\SGARefLink{sga3:xv:lemma:2:5}{2.5}).

Suppose now that \(S\) is the spectrum of a henselian local ring,
which the usual reductions allow us to take noetherian. Let \(s\) be
its closed point, let \(T_s\) be the maximal central torus of \(H_s\),
let \(q\) be invertible on \(S\), let \(\ell\) be a power of \(q\),
and let \(T\) be an \(S\)-torus whose closed fibre is isomorphic to
\(T_s\) (Exposé X~4.6). Let \({}_\ell H\) denote the ``kernel'' of
the \(\ell\)-th power morphism in \(H\), and let \(U_\ell\) be the
largest open subset of \({}_\ell H\) étale over \(S\). \SGARefLink{sga3:xv:lemma:1:3}{Lemma 1.3} and
flatness of \(H\) show that \(U_\ell\) contains \({}_\ell T_s\).
Hensel's lemma gives a unique \(S\)-morphism
\[
  u_\ell:{}_\ell T\longrightarrow U_\ell
\]
inducing the canonical immersion
\({}_\ell T_s\to(U_\ell)_s\) on the closed fibre. Uniqueness shows
that \(u_\ell\) is a central homomorphism (Exposé \SGARefLink{sga3:ix:theorem:5:6:item:a}{IX~5.6(a)}).
Proceeding as in \SGARefLink{sga3:xv:lemma:6:6}{6.6} and \SGARefLink{sga3:xv:lemma:6:7}{6.7}, one proves that \(u_\ell\) is an
immersion and, if \(M_\ell=u_\ell({}_\ell T)\), that
\[
  \operatorname{Centr}_G(M_\ell)=H
\]
for all sufficiently large \(\ell\). This is where noetherianity of
\(S\) is used.

\medskip\noindent\textit{Proof of (i).}
If \(H\in\mathcal{CT}(S)\), then \(H\) has connected fibres (Exposé
\SGARefLink{sga3:xii:theorem:6-6:item:b}{XII~6.6(b)}) and equals its connected normalizer (\SGARefLink{sga3:xv:remark:6:5}{6.5}). Hence the
present functor \(\mathcal{CT}\) is the functor of \SGARefLink{sga3:xv:paragraph:6-4-0}{6.4.0}. \SGARefLink{sga3:xv:theorem:6:4}{Theorem 6.4}
therefore represents it by an \(S\)-prescheme of finite presentation,
quasi-projective over \(S\). It remains to prove smoothness.

EGA IV~8 reduces us to affine noetherian \(S\). By Exposé XI~1.5 it
then suffices to prove the following infinitesimal lifting property:
if \(S\) is local artinian, \(S_0\) is defined by a nilpotent ideal,
and \(H_0\in\mathcal{CT}(S_0)\), then \(H_0\) lifts to a smooth
subgroup prescheme \(H\) of \(G\). By part (ii),
\[
 H_0=\operatorname{Centr}_{G_0}(T_0)
\]
for a sub-torus \(T_0\) of \(G_0\). Since \(G\) is smooth, \(T_0\)
lifts to a sub-torus \(T\) of \(G\) (Exposé \SGARefLink{sga3:ix:theorem:3:6:bis}{IX~3.6 bis}), and
\(H=\operatorname{Centr}_G(T)\) works; it is smooth by Exposé \SGARefLink{sga3:xi:corollary:2:4}{XI~2.4}
and \SGARefLink{sga3:xv:lemma:2:5}{Lemma 2.5}.

\medskip\noindent\textit{Proof of (iii).}
Because \(H\) is smooth with connected fibres, Exposé \SGARefLink{sga3:xi:corollary:6:11}{XI~6.11}
represents \(\operatorname{Transp}_G(H,L)\) by a closed subprescheme
of \(G\), of finite presentation over \(S\). Smoothness reduces to
the following lifting problem. Let \(S\) be local artinian, \(S_0\)
a closed subprescheme, and \(g_0\in G(S_0)\) such that
\[
  \operatorname{int}(g_0)H_0\subseteq L_0.
\]
We must lift \(g_0\) to \(g\in G(S)\) with
\(\operatorname{int}(g)H\subseteq L\).

Smoothness of \(G\) gives \(g_1\in G(S)\) lifting \(g_0\). Put
\(H'=\operatorname{int}(g_1)H\). Then
\(H'\in\mathcal{CT}(S)\) and \(H'_0\subseteq L_0\). By (ii),
\(H'=\operatorname{Centr}_G(T')\) for a torus \(T'\subseteq G\).
Since \(L\) is smooth, \(T'_0\subseteq L_0\) lifts to a torus
\(T''\subseteq L\) (Exposé \SGARefLink{sga3:ix:theorem:3:6:bis}{IX~3.6 bis}). The group
\(\operatorname{Centr}_L(T'')\) lies in
\(\operatorname{Centr}_G(T'')\), has the same closed fibre as the
latter, namely \(H'_0\), and is smooth. Hence
\[
 \operatorname{Centr}_L(T'')
 =\operatorname{Centr}_G(T'').
\]
The sub-tori \(T'\) and \(T''\) of \(G\) lift the same \(T'_0\), so
they are conjugate by an element \(h\in G(S)\) reducing to the
identity (Exposé \SGARefLink{sga3:ix:corollary:3:3:bis}{IX~3.3 bis}). Their centralizers are therefore
conjugate. Then \(g=hg_1\) lifts \(g_0\) and satisfies
\(\operatorname{int}(g)H\subseteq L\).

For \(g\in\operatorname{Transp}_G(H,L)(S)\), the equality
\(\operatorname{int}(g)H=L\) holds if and only if
\(\dim H_s=\dim L_s\) for every \(s\in S\). If \(U\) is the
open-and-closed subprescheme over which these dimensions agree, then
\[
 \operatorname{Transpstr}_G(H,L)
 =U\times_S\operatorname{Transp}_G(H,L).
\]

\medskip\noindent\textit{Proof of (iv).}
To prove \(H\) closed in \(G\), reduce first to affine noetherian
\(S\), then, by EGA IV~8, to the spectrum of a complete local ring.
Part (ii) writes \(H\) as the centralizer in \(G\) of a subgroup of
multiplicative type, so \(H\) is closed because \(G\) is separated
(Exposé VI\textsubscript{B}~\SGARefLink{sga3:xv:theorem:5-2}{5.2}).

By (iii),
\[
 N=\operatorname{Norm}_G(H)
  =\operatorname{Transpstr}_G(H,H)
\]
is represented by a subgroup prescheme smooth over \(S\) and closed
in \(G\). The morphism
\[
  G\longrightarrow\mathcal{CT},
  \qquad g\longmapsto\operatorname{int}(g)H
\]
is smooth by (iii), so its image is an open subset \(U\) of
\(\mathcal{CT}\). As in Exposé \SGARefLink{sga3:xi:corollary:5:3}{XI~5.3}, \(G/N\) is represented by
\(U\); it is therefore quasi-projective.

We now treat \(N/H\). By EGA IV~8, assume \(S\) affine noetherian and
then \(S=\operatorname{Spec}A\), with \(A\) local of finite
dimension\footnote{Krull dimension.}. We use induction on
\(\dim S\). When \(\dim S=0\), the assertion is Exposé
VI\textsubscript{A}, §4. If \(N/H\) is representable, it is separated
because \(H\) is closed in \(N\), of finite type and étale because
\(N\) is smooth of finite type and \(H\) is open. It is therefore
quasi-affine (SGA 1, VIII~6.2). Effective descent for quasi-affine
schemes (\emph{loc. cit.}, 7.9) permits replacement of \(A\) by its
completion.

Thus assume \(S\) complete noetherian local, with closed point \(s\),
and put \(U=S\setminus\{s\}\). By induction,
\((N|_U)/(H|_U)\) is represented by a \(U\)-group \(K\). Let \(T_s\)
be the maximal central torus of \(H_s\), choose \(q\) invertible on
\(S\), let \(\ell\) be a power of \(q\), and let \(M_\ell\) be the
unique central subgroup of multiplicative type of \(H\) lifting
\({}_\ell T_s\), as in (ii). Choose \(\ell\) so large that
\[
  \operatorname{Centr}_G(M_\ell)=H,
\]
and put \(N'=\operatorname{Norm}_G(M_\ell)\).

The torus \(T_s\), and hence \((M_\ell)_s\), is characteristic in
\(H_s\). Therefore \(N_s\) normalizes \((M_\ell)_s\), and
\(N_s=N'_s\). The proof of Exposé \SGARefLink{sga3:xi:proposition:5:9}{XI~5.9} represents
\[
  N'/H
  =\operatorname{Norm}_G(M_\ell)/
   \operatorname{Centr}_G(M_\ell)
\]
by an \(S\)-group \(K'\). Since \(N'\) is smooth and
\(N'_s=N_s\), it is an open subgroup of \(N\) containing \(H\)
(Exposé VI\textsubscript{B}~\SGARefLink{sga3:xv:lemma:2:5}{2.5}). Its image over \(U\) is an open
subgroup of \(K\), isomorphic to \(K'|_U\). Glue \(K\) and \(K'\)
along this open subgroup to obtain \(L\), and glue the projections
\(N|_U\to K\) and \(N'\to K'\) to obtain \(p:N\to L\). Then
\((L,p)\) represents \(N/H\).

\medskip\noindent\textit{Proof of (v).}
First suppose \(S=\operatorname{Spec}k\). The image \(H'\) of \(H\)
is a smooth subgroup of \(G'\). The required assertion follows from:

\SGAStatement{Lemma 7.2}{Let \(u:G\to G'\) be an epimorphism of
smooth connected algebraic groups over a field \(k\), let \(T\) be a
torus of \(G\), and let \(T'\) be its image in \(G'\). Then
\[
 u(\operatorname{Centr}_G(T))
 =\operatorname{Centr}_{G'}(T').
\]}
\SGARefTarget{sga3:xv:lemma:7:2}

Put \(H=\operatorname{Centr}_G(T)\),
\(H'=\operatorname{Centr}_{G'}(T')\), and \(H''=u(H)\). We have
\(H''\subseteq H'\). After extending the base field, assume \(k\)
algebraically closed. It suffices to prove that every Cartan subgroup
of \(H'\) lies in \(H''\): then \(H''\) contains the open subset of
regular points of \(H'\), so is open, and equals the connected group
\(H'\).

Let \(C'\) be a Cartan subgroup of \(H'\). It is also a Cartan
subgroup of \(G'\), since \(H'\), being the centralizer of the torus
\(T'\), has the same reductive and nilpotent ranks as \(G'\). Put
\[
  K=(u^{-1}(C'))^0_{\mathrm{red}}.
\]
Because \(T'\) is central in \(H'\), \(C'\) contains \(T'\), and
\(K\) contains \(T\). Let \(C\) be a Cartan subgroup of \(K\)
containing \(T\). The torus \(T\) lies in the unique maximal torus of
\(C\), central in \(C\) (Exposé \SGARefLink{sga3:xii:6.6:item:c}{XII~6.6(c)}); hence
\(C\subseteq H\). Any two Cartan subgroups of \(G\) are conjugate,
and the image of a Cartan subgroup of \(G\) is a Cartan subgroup of
\(G'\) (Exposé \SGARefLink{sga3:xii:6.6}{XII~6.6}). It follows that \(C\) is a Cartan subgroup
of \(G\), and its image is a Cartan subgroup of \(G'\) contained in
\(C'\). Thus \(u(C)=C'\), proving \(C'\subseteq H''\).

This proves (v) over a field. In general, finite presentation and the
usual reductions permit us to assume \(S\) affine noetherian, then
local, and finally the spectrum of a complete noetherian local ring
\(A\), by fpqc descent of subpreschemes of \(G'\).

Use the notation of (ii). Let \(T_s\) be the maximal central torus of
\(H_s\), and let \(M_\ell\subset H\) lift \({}_\ell T_s\), with
\(H=\operatorname{Centr}_G(M_\ell)\). Let \(T'_s=u_s(T_s)\).
Since \(G'\) is separated, \(u(M_\ell)\) is represented by a subgroup
\(M'_\ell\) of multiplicative type (Exposé \SGARefLink{sga3:ix:theorem:6:8}{IX~6.8}). Put
\[
  H'=\operatorname{Centr}_{G'}(M'_\ell),
\]
which is smooth. For every power \(\ell'\) of \(q\), a sufficiently
large \(\ell\) makes \((M'_\ell)_s\) contain
\({}_{\ell'}T'_s\). We may therefore choose \(\ell\) such that
\[
 H'_s=\operatorname{Centr}_{G'_s}(T'_s)=u_s(H_s),
\]
the last equality being \SGARefLink{sga3:xv:lemma:7:2}{Lemma 7.2}.

The restriction \(v:H\to H'\) of \(u\) is flat. Indeed, \(H,H'\) are
flat over \(S\), \(H_s\to H'_s\) is flat, and hence \(v\) is flat
near \(H_s\) (EGA IV~11.3.10 and 11.3.1). It is therefore flat on an
open subgroup of \(H\) (Exposé VI\textsubscript{B}~2.2), and thus
everywhere because the fibres of \(H\) are connected. The
set-theoretic image of \(H\) is consequently an open subset of \(H'\),
necessarily \((H')^0\), and with its induced structure represents
the fpqc image sheaf \(u(H)\). \SGARefLink{sga3:xv:lemma:7:2}{Lemma 7.2} shows that it belongs to
\(\mathcal{CT}_{G'}(S)\).

If \(H=\operatorname{Centr}_G(T)\), then \(T'=u(T)\) is a sub-torus
of \(G'\) (Exposé \SGARefLink{sga3:ix:theorem:6:8}{IX~6.8}). The image \(u(H)\) is contained in
\(\operatorname{Centr}_{G'}(T')\), agrees with it fibrewise by
\SGARefLink{sga3:xv:lemma:7:2}{Lemma 7.2}, and is smooth. Hence
\[
  u(H)=\operatorname{Centr}_{G'}(T').
\]

\medskip\noindent\textit{Proof of (vi).}
Since \(\mathcal{CT}_G,\mathcal{CT}_{G'}\) are smooth over \(S\), it
is enough to prove faithful flatness on geometric fibres. Thus take
\(S=\operatorname{Spec}k\), with \(k\) algebraically closed. Let
\(H'\in\mathcal{CT}_{G'}(k)\), let \(T'\) be its maximal central
torus, choose a sub-torus \(T\subset G\) mapping onto \(T'\), and put
\[
 H=\operatorname{Centr}_G(T),\qquad
 N=\operatorname{Norm}_G(H),\qquad
 N'=\operatorname{Norm}_{G'}(H').
\]
By \SGARefLink{sga3:xv:lemma:7:2}{Lemma 7.2}, \(H'=u(H)\). Part (iv) identifies \(G/N\) and \(G'/N'\)
with open neighbourhoods of \(H\) and \(H'\) in the two
\(\mathcal{CT}\)-spaces. On these neighbourhoods, \(\widetilde u\)
is
\[
  w:G/N\longrightarrow G'/N',
\]
obtained from \(u\) by passage to quotients. This is an epimorphism
of homogeneous spaces under \(G\), hence faithfully flat. Therefore
\(\widetilde u\) is flat and its image is an open subset of
\(\mathcal{CT}_{G'}\) containing all \(k\)-points. Since
\(\mathcal{CT}_{G'}\) is of finite type, \(\widetilde u\) is
surjective and hence faithfully flat. If \(\ker u\) is central, the
remaining assertions are Exposé \SGARefLink{sga3:xii:7.12}{XII~7.12}.

The next theorem generalizes Exposé \SGARefLink{sga3:xii:7.1}{XII~7.1}.

\SGAStatement{\SGAInlineRefTarget{sga3:xv:theorem:7-3}Theorem 7.3}{Let \(S\) be a prescheme and let \(G\)
be an \(S\)-group prescheme of finite presentation, smooth over \(S\),
with connected fibres. Define the functor
\[
 \mathcal C(S')=\{\text{Cartan subgroups of }G_{S'}\}.
\]
\begin{enumerate}
\renewcommand{\labelenumi}{\textup{\roman{enumi})}}
\item The following are equivalent:
\begin{enumerate}
\renewcommand{\labelenumii}{\textup{\alph{enumii})}}
\item \(\mathcal C\) is representable;
\item \(\mathcal C\) is represented by an \(S\)-prescheme smooth,
quasi-projective, and of finite presentation over \(S\), with affine
fibres;
\item \(G\) admits a Cartan subgroup locally for the étale topology;
\item \(G\) admits a Cartan subgroup locally for the faithfully flat
topology;
\item the nilpotent rank of the fibres of \(G\) is locally constant
on \(S\).
\end{enumerate}
\item If these conditions hold, any two Cartan subgroups of \(G\)
are locally conjugate for the étale topology. The regular locus of
the fibres of \(G\) (Exposé \SGARefLink{sga3:xiii:corollary:2:7}{XIII~2.7}) is an open subprescheme
\(G_{\mathrm{reg}}\), of finite presentation over \(S\), and every
section of \(G_{\mathrm{reg}}\) over \(S\) is contained in a unique
Cartan subgroup of \(G\).
\item Let \(G'\) be an \(S\)-group prescheme of finite presentation
and \(u:G\to G'\) faithfully flat. Then, for every Cartan subgroup
\(C\) of \(G\), its image \(C'=u(C)\) is represented by a Cartan
subgroup of \(G'\), and \(C\to C'\) is faithfully flat.
\item Under the conditions of (i) and (iii), the morphism
\[
 \widetilde u:\mathcal C_G\longrightarrow\mathcal C_{G'},
 \qquad C\longmapsto u(C),
\]
is faithfully flat. If \(\ker u\) is central, it is an isomorphism.
\item Further information on transporters and relations with maximal
tori is given in \SGARefLink{sga3:xv:proposition:7-1}{Proposition 7.1} and Exposé \SGARefLink{sga3:xii:7.1}{XII~7.1}.
\end{enumerate}}
\SGARefTarget{sga3:xv:theorem:7:3}

\SGAProof{Proof}
For (i) we prove
\[
 (\mathrm b)\Rightarrow(\mathrm c)\Rightarrow(\mathrm d)
 \Rightarrow(\mathrm e)\Rightarrow(\mathrm b)
 \Rightarrow(\mathrm a)\Rightarrow(\mathrm d).
\]
Assume (b), and let \(s\in S\). Since \(\mathcal C_s\) is smooth
over \(\kappa(s)\), it has points with residue field a finite
separable extension of \(\kappa(s)\). Exposé \SGARefLink{sga3:xi:proposition:1:10}{XI~1.10} gives an open
neighbourhood \(U\) of \(s\) and an étale surjective morphism
\(U'\to U\) such that \(G_{U'}\) admits a Cartan subgroup. Thus
(b)\(\Rightarrow\)(c), while (c)\(\Rightarrow\)(d) is clear.

Assume (d), and let \(s\in S\). There is a flat \(S\)-prescheme
\(S'\), whose image contains \(s\), such that \(G_{S'}\) admits a
Cartan subgroup. For \(s'\in S'\) above \(s\), the nilpotent rank of
the fibres of \(G_{S'}\) is constant on
\(\operatorname{Spec}\mathcal O_{S',s'}\). Hence the nilpotent rank
of the fibres of \(G\) is constant on
\(\operatorname{Spec}\mathcal O_{S,s}\), which is the image of that
local spectrum (EGA IV~2.3.4(ii)). Let the value be \(r\). By 6.3 bis,
the set \(E_r\) where the nilpotent rank equals \(r\) is
ind-constructible; it contains a neighbourhood of \(s\) by
EGA IV~1.10.1. Thus (d)\(\Rightarrow\)(e).

Assume (e). Locally on \(S\), the nilpotent rank is a constant \(r\).
For every \(S'\), the Cartan subgroups of \(G_{S'}\) are exactly the
smooth subgroup preschemes of relative dimension \(r\) whose
geometric fibres are centralizers of tori. By \SGARefLink{sga3:xv:proposition:7-1}{Proposition 7.1},
\(\mathcal C\) is the open-and-closed subprescheme of
\(\mathcal{CT}\) representing the relative-dimension-\(r\)
subfunctor. It has all the properties in (b); affineness of its
fibres is Exposé \SGARefLink{sga3:xii:theorem:7-1:item:d}{XII~7.1(d)}. This proves
(e)\(\Rightarrow\)(b), and (b)\(\Rightarrow\)(a) is immediate.

Finally assume (a). By \SGARefLink{sga3:xv:proposition:6-2:item:iii}{6.2(iii)}, \(\mathcal C\) commutes with filtered
inductive limits of rings, so its representing prescheme is locally
of finite presentation (EGA IV~8.14.2). To prove it smooth, it is
enough to show that, when \(S\) is affine and \(S_0\subseteq S\) is
defined by a nilpotent ideal, every Cartan subgroup \(C_0\) of
\(G_0\) lifts to a Cartan subgroup of \(G\). Existence of \(C_0\)
implies (e) on \(S_0\), hence on \(S\), which has the same underlying
space. Since (e)\(\Rightarrow\)(b), \(\mathcal C\) is smooth.
It is also surjective over \(S\), so (a)\(\Rightarrow\)(d).

For (ii), if \(C,C'\) are Cartan subgroups, then
\(\operatorname{Transp}_G(C,C')\) is represented by a smooth
\(S\)-prescheme (\SGARefLink{sga3:xv:proposition:7-1:item:iii}{Proposition 7.1(iii)}) with nonempty fibres
(Exposé \SGARefLink{sga3:xii:theorem:6-6:item:a}{XII~6.6(a)},(c)). Hensel's lemma (Exposé \SGARefLink{sga3:xi:proposition:1:10}{XI~1.10}) gives local
conjugacy for the étale topology. The assertions about regular points
follow from Exposé \SGARefLink{sga3:xiii:theorem:3:1}{XIII~3.1} and \SGARefLink{sga3:xiii:corollary:3:2}{3.2}.

For (iii), \SGARefLink{sga3:xv:proposition:7-1:item:v}{Proposition 7.1(v)} represents \(u(C)\) by a smooth
subgroup \(C'\) of \(G'\). Its fibres are Cartan subgroups of the
fibres of \(G'\) (Exposé \SGARefLink{sga3:xii:6.6:item:d}{XII~6.6(d)}), so \(C'\) is a Cartan subgroup.

For (iv), faithful flatness is proved as in \SGARefLink{sga3:xv:proposition:7-1:item:vi}{Proposition 7.1(vi)}.
If \(\ker u\) is central and \(C'\) is a Cartan subgroup of \(G'\),
then \(C=u^{-1}(C')\) is smooth with connected fibres
(\SGARefLink{sga3:xv:proposition:7-1:item:vi}{Proposition 7.1(vi)}), and its fibres are Cartan subgroups (Exposé
\SGARefLink{sga3:xii:6.6:item:e}{XII~6.6(f)}). Thus \(C\) is a Cartan subgroup, and \SGARefLink{sga3:xv:proposition:7-1:item:vi}{Proposition
7.1(vi)} completes the proof.

% END INCLUDED SGA3_English_current_buildable_source_20260823/source/inputs/XV/tex/components/07_expose_XV_section7.tex

% BEGIN INCLUDED SGA3_English_current_buildable_source_20260823/source/inputs/XV/tex/components/08_expose_XV_section8_through_theorem8_9.tex
% Authority: Expo15.pdf, local pp. 56--61 / combined pp. 926--931 /
% printed pp. 459--470.
% Stops immediately before the proof of Theorem 8.8(a).

\SGASection{8}{Representability criterion for the functor of sub-tori
of a smooth group}
\SGARefTarget{sga3:xv:section:8}

\medskip\noindent\SGAInlineRefTarget{sga3:xv:paragraph:8-0}\textbf{8.0.}
In this section, if \(S\) is a prescheme and \(G\) an \(S\)-group
prescheme, \(\mathcal T_G\) (or simply \(\mathcal T\), if there is no
ambiguity) denotes the \(S\)-functor
\[
  (\operatorname{Sch}/S)^\circ\longrightarrow\operatorname{Ens}
\]
such that, for every \(S\)-prescheme \(S'\),
\[
  \mathcal T(S')=\{\text{sub-tori of \(G_{S'}\)}\}.
\]
The functor \(\mathcal T_C\) of central sub-tori of \(G\) is defined
similarly.
\SGARefTarget{sga3:xv:paragraph:8:0}

\SGAStatement{\SGAInlineRefTarget{sga3:xv:proposition:8-1}Proposition 8.1}{Let \(S\) be a locally noetherian
prescheme and let \(G\) be an \(S\)-group prescheme of finite type.
Then the following conditions are equivalent:
\begin{enumerate}
\item[\textup{i)}] The functor \(\mathcal T_G\) ``commutes with adic
limits of local artinian rings'': for every \(S\)-prescheme
\(S'=\operatorname{Spec}(A)\), where \(A\) is a complete noetherian
local ring for the topology defined by its maximal ideal
\(\mathfrak m\), the canonical map
\[
  \mathcal T(S')\longrightarrow
  \varprojlim_n\mathcal T(S'_n),
  \qquad S'_n=\operatorname{Spec}(A/\mathfrak m^n),
\]
is bijective.
\item[\textup{ii)}] The assertion in \textup{i)}, restricted to the
case in which \(A\) is a complete discrete valuation ring with
algebraically closed residue field.
\item[\textup{iii)}] The assertion in \textup{ii)}, restricted to the
subfunctor \(\mathcal T^{(1)}\) of \(\mathcal T\) formed by sub-tori
of \(G\) of relative dimension \(1\).
\item[\textup{i bis)}] For every \(S\)-prescheme \(S'\) as in
\textup{i)} and every \(S'\)-torus \(T\), the canonical map
\[
  \operatorname{Hom}_{S'\text{-gr}}(T,G_{S'})
  \longrightarrow
  \varprojlim_n
  \operatorname{Hom}_{S'_n\text{-gr}}(T_{S'_n},G_{S'_n})
\]
is bijective.
\item[\textup{ii bis)}] The assertion in \textup{i bis)}, restricted
to the case in which \(A\) is a complete discrete valuation ring
with algebraically closed residue field.
\item[\textup{iii bis)}] The assertion in \textup{ii bis)}, restricted
to \(T=\mathbb G_{m,S'}\).
\end{enumerate}}
\SGARefTarget{sga3:xv:proposition:8:1}

\medskip\noindent\textbf{Remark 8.2.}
There is an analogous proposition in which one restricts to central
sub-tori of \(G\) and to central homomorphisms from a torus into \(G\).
\SGARefTarget{sga3:xv:remark:8:2}

\medskip\noindent\textit{Proof.}
We shall use the following lemma.

\SGAStatement{Lemma 8.3}{Let \(S\) be a prescheme, \(G\) an
\(S\)-group prescheme, \(T\) an \(S\)-torus, and
\(u:T\to G\) an \(S\)-group morphism. Suppose, moreover, either that
\(G\) is of finite presentation over \(S\), or that \(S\) is locally
noetherian and \(G\) is locally of finite type. Then:
\begin{enumerate}
\renewcommand{\labelenumi}{\textup{\alph{enumi})}}
\item \(\ker u\) is a subgroup of multiplicative type of \(T\);
\item the quotient \(T'=T/\ker u\) is a torus;
\item the canonical monomorphism \(T'\to G\) induced by \(u\) after
passage to the quotient is an immersion.
\end{enumerate}}
\SGARefTarget{sga3:xv:lemma:8:3}

When \(G\) is separated over \(S\), this lemma follows from
Exposé \SGARefLink{sga3:ix:theorem:6:8}{IX~6.8}. In the general case, the usual reduction allows us to
suppose that \(S\) is noetherian. We first show that
\(K=\ker u\) is flat. We may suppose that \(S\) is the spectrum of a
local artinian ring (EGA \(0_{\mathrm{III}}\), 10.2.6); then \(G\) is
separated (Exposé VI\textsubscript{B}~5.2), so \(K\) is of
multiplicative type (Exposé \SGARefLink{sga3:ix:theorem:6:8}{IX~6.8}), and hence is flat over \(S\).
We next show that \(\ker u\) is closed in \(T\). This reduces us to
the case in which \(S\) is the spectrum of a discrete valuation ring
(EGA II~7.2.1). Since \(T\) is flat with connected fibres, \(u\)
factors through the identity component of the schematic closure in
\(G\) of the generic fibre of \(G\). We may thus suppose that \(G\)
is flat with connected fibres; it is then separated
(Exposé VI\textsubscript{B}~\SGARefLink{sga3:xv:theorem:5-2}{5.2}), and \(K\) is closed.
It follows from Exposé \SGARefLink{sga3:x:corollary:4-8:item:b}{X~4.8(b)} that \(K\) is a subgroup of
multiplicative type of \(T\). The quotient \(T'=T/K\) is therefore
representable and is a group of multiplicative type
(Exposé \SGARefLink{sga3:ix:proposition:2:7}{IX~2.7}) whose fibres are tori; hence \(T'\) is a torus.
Finally, the monomorphism \(T'\to G\) is an immersion by
Exposé \SGARefLink{sga3:viii:proposition:7:9}{VIII~7.9}.

\medskip\noindent\textit{Proof of \SGARefLink{sga3:xv:proposition:8-1}{Proposition 8.1}.}
\textup{i)}\(\Rightarrow\)\textup{i bis)}.
Put \(T_n=T_{S'_n}\), \(G_n=G_{S'_n}\), and let
\[
 (u_n)_{n\in\mathbb N}\in
 \varprojlim_n\operatorname{Hom}_{S'_n\text{-gr}}(T_n,G_n).
\]
For every \(n\), \(u_n(T_n)\) is a sub-torus \(T'_n\) of \(G_n\)
(\SGARefLink{sga3:xv:lemma:8:3}{Lemma 8.3}). By hypothesis, there is a unique sub-torus \(T'\) of
\(G\) lifting \(T'_n\) for every \(n\). Since \(T'\) has affine
fibres, the assertion follows from \SGARefLink{sga3:xv:corollary:4:4}{4.4}.

\textup{i bis)}\(\Rightarrow\)\textup{i)}.
Let
\((T_n)_{n\in\mathbb N}\in\varprojlim_n\mathcal T(S'_n)\).
By Exposé \SGARefLink{sga3:x:corollary:4-6}{X~4.6} there are an \(S'\)-torus \(T^0\) and an
\(S'_0\)-isomorphism
\[
  u_0:T^0_0\xrightarrow{\sim}T_0.
\]
Since \(T_n\) is smooth over \(S'_n\), \(u_0\) lifts to an
\(S'_n\)-morphism \(u_n:T^0_n\to T_n\) (Exposé \SGARefLink{sga3:ix:theorem:3:6}{IX~3.6}), and the
family \((u_n)_{n\in\mathbb N}\) may be chosen coherent. It therefore
comes from a morphism \(u:T^0\to G\). By \SGARefLink{sga3:xv:lemma:8:3}{Lemma 8.3}, the image of
\(T^0\) under \(u\) is a sub-torus of \(G\), and it lifts \(T_n\)
for every \(n\).

The implications
\[
 \textup{i)}\Rightarrow\textup{ii)}\Rightarrow\textup{iii)}
 \quad\text{and}\quad
 \textup{i bis)}\Rightarrow\textup{ii bis)}
 \Rightarrow\textup{iii bis)}
\]
are immediate. The implication
\(\textup{iii)}\Rightarrow\textup{iii bis)}\) is proved in the same
way as
\(\textup{i)}\Rightarrow\textup{i bis)}\).
It remains to prove
\[
 \textup{iii bis)}\Rightarrow\textup{ii bis)}
 \Rightarrow\textup{i bis)}.
\]

\textup{ii bis)}\(\Rightarrow\)\textup{i bis)}.
In the terminology of \SGARefLink{sga3:xv:proposition:4:3}{4.3}, assertion \textup{i bis)} holds if and
only if every element
\[
 (u_n)_{n\in\mathbb N}\in
 \varprojlim_n\operatorname{Hom}_{S'_n\text{-gr}}(T_n,G_n)
\]
is admissible. For every point \(t\) of \(S'\) distinct from its
closed point \(s\), there is an \(S\)-scheme \(S''\), the spectrum of
a complete discrete valuation ring with algebraically closed residue
field, whose generic point maps to \(t\) and whose closed point maps
to \(s\) (EGA II~7.1.9). This gives at once a valuative criterion for
a family of morphisms to be admissible, and proves the implication.

\textup{iii bis)}\(\Rightarrow\)\textup{ii bis)}.
Let \(T\) be an \(S\)-torus, where \(S\) is the spectrum of a complete
discrete valuation ring with algebraically closed residue field. The
torus \(T\) is split (Exposé \SGARefLink{sga3:x:corollary:4-6}{X~4.6}), say
\[
  T\simeq\mathbb G_{m,S}^{\,r}.
\]
Let
\[
 (u_n)_{n\in\mathbb N}\in
 \varprojlim_n
 \operatorname{Hom}_{S_n\text{-gr}}
 \bigl((\mathbb G_m)^r_{S_n},G_n\bigr).
\]
By hypothesis, the restriction of \(u_n\) to each factor of
\(\mathbb G_{m,S}^{\,r}\) comes from a group morphism
\(\mathbb G_{m,S}\to G\). Their product, followed by the composition
law \(G^r\to G\), gives a morphism
\[
  v:\mathbb G_{m,S}^{\,r}\longrightarrow G.
\]
Clearly \(u_n=v_n\) for every \(n\). It remains to prove that \(v\)
is a group morphism. Equivalently, two evident morphisms
\[
 f,g:X=\mathbb G_{m,S}^{\,r}\times_S\mathbb G_{m,S}^{\,r}
 \longrightarrow G
\]
must coincide. Let \(Z\) be their coincidence subscheme. Since
\(\mathbb G_{m,S}^{\,r}\) has connected fibres and its image under
\(v\) contains the unit section, the argument of \SGARefLink{sga3:xv:lemma:8:3}{Lemma 8.3} shows
that \(v\) factors through the identity component of \(G\). We may
therefore suppose that \(G\) is separated
(Exposé VI\textsubscript{B}~\SGARefLink{sga3:xv:theorem:5-2}{5.2}), so \(Z\) is closed. On the other
hand, \(v_n\) is a group morphism, hence \(f_n=g_n\) for every \(n\).
Thus \(Z\) contains a neighbourhood of the closed fibre of \(X\)
(EGA I~10.9.4), which is schematically dense in \(X\), since \(X\)
has smooth irreducible fibres (Exposé IX~4.6). Consequently
\(Z=X\), and \(f=g\).

\SGAStatement{\SGAInlineRefTarget{sga3:xv:definition:8-4}Definition 8.4}{Let \(S\) be a locally noetherian
prescheme and \(G\) an \(S\)-group prescheme of finite type. Let
\(\mathcal S\) be the set of \(S\)-schemes \(S'\) that are spectra
of complete discrete valuation rings with algebraically closed
residue field. Write \(s\) and \(t\) for the closed and generic
points of \(S'\), and require that
\((G_t)^0_{\mathrm{red}}\) be smooth and admit a Chevalley
decomposition: an extension of an abelian variety by a smooth
connected linear algebraic subgroup \(L_t\). This decomposition is
then unique. For \(S'\in\mathcal S\), let \(G'\), respectively \(L\),
denote the schematic closure in \(G_{S'}\) of \(G_t\), respectively
of \(L_t\).

We say that ``the abelian part of \(G\) does not degenerate into a
toric part'', or briefly that \(G\) has property \(\mathrm{AT}\), if,
for every \(S'\in\mathcal S\), \(L_s\) has the same reductive rank
as \(G'_s\). Intuitively, suppose that the quotient \(A=G'/L\) is
representable. Then \(A\) is a flat group prescheme and
\((A_t)^0_{\mathrm{red}}\) is an abelian variety; property
\(\mathrm{AT}\) says that \(A_s\) has reductive rank zero, so that
\((A_s)^0_{\mathrm{red}}\) is an extension of an abelian variety by
a unipotent group.

Similarly, if \(G\) has connected fibres, we say that \(G\) has
property \(\mathrm{ATC}\) if, for every \(S'\in\mathcal S\), the
schematic closure \(Z\) of the centre \(Z_t\) of \(G_t\) has
property \(\mathrm{AT}\).}
\SGARefTarget{sga3:xv:definition:8:4}

These technical definitions are justified by the following
proposition.

\SGAStatement{Proposition 8.5}{Let \(S\) be a locally noetherian
prescheme and let \(G\) be an \(S\)-group prescheme of finite type.
\begin{enumerate}
\renewcommand{\labelenumi}{\textup{\alph{enumi})}}
\item\SGAInlineRefTarget{sga3:xv:proposition:8-5:item:a} The functor \(\mathcal T\) of sub-tori of \(G\) commutes with
adic limits of local artinian rings (cf. \SGARefLink{sga3:xv:proposition:8-1}{Proposition 8.1}) if and
only if \(G\) has property \(\mathrm{AT}\).
\item\SGAInlineRefTarget{sga3:xv:proposition:8-5:item:b} If \(G\) has connected fibres, the functor \(\mathcal T_C\) of
central sub-tori of \(G\) commutes with adic limits of local artinian
rings if and only if \(G\) has property \(\mathrm{ATC}\).
\end{enumerate}}
\SGARefTarget{sga3:xv:proposition:8:5}

We need the following technical lemma.

\SGAStatement{Lemma 8.6}{Let \(S\) be the spectrum of a discrete
valuation ring, with closed point \(s\), let \(G\) be an
\(S\)-group prescheme flat and of finite type, and let \(T_s\) be a
sub-torus of \(G_s\). There are an \(S\)-scheme \(S'\), the spectrum
of a discrete valuation ring faithfully flat over \(S\), with closed
point \(s'\), and a subgroup scheme \(C\) of \(G_{S'}\), flat over
\(S'\), commutative, and with connected fibres, such that
\(C_{s'}\) contains \(T_s\times_s s'\).}
\SGARefTarget{sga3:xv:lemma:8:6}

After replacing \(S\) by the spectrum \(S'\) of a discrete valuation
ring faithfully flat over it, we may suppose that
\(T_s=\mathbb G_{m,s}^{\,r}\) and that the transcendence degree of
\(\kappa(s)\) over its prime field is at least \(r\)
(EGA \(0_{\mathrm{III}}\), 10.3.1, and EGA II~7.1.9). There is then
an element \(x\in T_s(s)\) such that every algebraic subgroup of
\(G_s\) containing \(x\) contains \(T_s\)
(cf. Exposé XIII, proof of \SGARefLink{sga3:xiii:theorem:2:1}{2.1},
\(\textup{ii)}\Rightarrow\textup{vii)}\)).
Since \(G\) is flat over \(S\), a quasi-section passes through \(x\)
(EGA IV~14.5.8). After another such replacement of \(S\), we may
therefore suppose that \(x\) lifts to a section
\(\widetilde x\) of \(G\) over \(S\). Let \(C_t\) be the commutative
algebraic subgroup of \(G_t\) generated by \(\widetilde x_t\)
(Exposé VI\textsubscript{B}~7), and let \(C\) be its schematic
closure in \(G\). Clearly \(C_s\) contains \(x\), hence contains
\(T_s\). The identity component of \(C\) is therefore a flat
commutative group scheme with connected fibres having the required
property.

\medskip\noindent\textit{Proof of \SGARefLink{sga3:xv:proposition:8-5:item:a}{Proposition 8.5(a)}.}
Suppose first that \(\mathcal T\) commutes with adic limits of
artinian rings. Let \(S'\in\mathcal S\), and let \(T_s\) be a maximal
torus of \(G'_s\). We must show that \(T_s\subset L_s\). The
formation of \(L\) and \(G'\) plainly commutes with faithfully flat
extensions \(S''\to S'\) of discrete valuation rings. By \SGARefLink{sga3:xv:lemma:8:6}{Lemma 8.6},
we may therefore replace \(S'\) so that there is a flat commutative
subgroup scheme \(C\) of \(G'_{S'}\) for which \(C_s\) contains
\(T_s\). The torus \(T_s\) is central in \(C_s\), and hence lifts
infinitesimally to a central sub-torus (\SGARefLink{sga3:xv:corollary:2:3}{Corollary 2.3}). By the
hypothesis on \(G\), \(T_s\) lifts to a sub-torus \(T\) of \(G\).
Clearly \(T_t\) lies in the linear part \(L_t\) of \(G_t\), and
therefore \(T\subset L\).

Conversely, suppose that \(G\) has property \(\mathrm{AT}\). We prove
condition \SGARefLink{sga3:xv:proposition:8-1}{8.1} \textup{iii bis)}. Let \(S=\operatorname{Spec}(A)\),
where \(A\) is a complete discrete valuation ring with algebraically
closed residue field, let \(\mathfrak m\) be its maximal ideal, put
\(S_n=\operatorname{Spec}(A/\mathfrak m^n)\), and let
\[
 u_n:(\mathbb G_{m,S})_n\longrightarrow
 G_n=G\times_S S_n,\qquad n\in\mathbb N,
\]
be a coherent system of group morphisms. Let \(q\) be a prime
invertible on \(S\). As \(\ell\) ranges over the powers of \(q\),
there is a unique \(S\)-group morphism
\[
 {}_\ell u:{}_\ell\mathbb G_{m,S}\longrightarrow G
\]
lifting the restriction of \(u_n\) to
\({}_\ell(\mathbb G_{m,S})_n\) for every \(n\)
(\SGARefLink{sga3:xv:proposition:1-6:item:a}{Proposition 1.6(a)}). Consequently, if an \(S\)-group morphism
\(u:\mathbb G_{m,S}\to G\) lifting every \(u_n\) exists, its
restriction to \({}_\ell\mathbb G_{m,S}\) is uniquely determined.
The density theorem (Exposé \SGARefLink{sga3:ix:corollary:4:8:item:a}{IX~4.8(a)}) proves both the uniqueness of
\(u\) and that, for its existence, we may make a faithfully flat
extension of the base.

Now \((G_t)^0_{\mathrm{red}}\) admits a Chevalley decomposition,
already defined over a finite extension \(L\) of the fraction field
\(K\) of \(A\). After replacing \(S\) by its normalization in \(L\),
we may suppose that \(S\in\mathcal S\). The morphism
\({}_\ell u_t\) factors through \((G_t)_{\mathrm{red}}\), and
\({}_\ell u\) therefore factors through the schematic closure
\(G''\) of \((G_t)_{\mathrm{red}}\) in \(G'\). The density theorem
then shows that every \(u_n\) factors through \(G''_n\).
Since \(G\) has property \(\mathrm{AT}\), every sub-torus of \(G'_s\),
and hence of \(G''_s\), is contained in \(L_s\); thus \(u_s\)
factors through \(L_s\). In fact every \(u_n\) factors through
\(L_n\). Indeed, \(L_t\) is invariant in \(G''_t\), so \(L\) is
invariant in \(G''\); moreover \(L\) is flat over \(S\), and the
quotient
\[
  H_n=G''_n/L_n
\]
is representable for every \(n\) (Exposé VI\textsubscript{A},
\S4). The image of \((\mathbb G_{m,S})_n\) in \(H_n\) is a
sub-torus (Exposé \SGARefLink{sga3:ix:theorem:6:8}{IX~6.8}) whose closed fibre is zero. It is therefore
zero, and \(u_n\) factors through \(L_n\). Since \(L_t\) is affine,
\SGARefLink{sga3:xv:proposition:4:3}{Proposition 4.3} now shows that the family \((u_n)\) is admissible.

The proof of \SGARefLink{sga3:xv:proposition:8-5:item:b}{Proposition 8.5(b)} is entirely analogous, in view of
\SGARefLink{sga3:xv:remark:8:2}{Remark 8.2} and Exposé \SGARefLink{sga3:ix:theorem:5:6:item:a}{IX~5.6(a)}.

\SGAStatement{Proposition 8.7}{Let \(S\) be a locally noetherian
prescheme and let \(G\) be an \(S\)-group prescheme of finite type.
\begin{enumerate}
\renewcommand{\labelenumi}{\textup{\roman{enumi})}}
\item If \(G\) has connected fibres and has property
\(\mathrm{AT}\), then it has property \(\mathrm{ATC}\).
\item If \(G\) has property \(\mathrm{AT}\), every subgroup
prescheme of \(G\) has property \(\mathrm{AT}\).
\item If \(G\) is flat over \(S\) and the abelian rank of its fibres
is a locally constant function on \(S\), then \(G\) has property
\(\mathrm{AT}\).
\end{enumerate}}
\SGARefTarget{sga3:xv:proposition:8:7}

Assertion \textup{(i)} follows at once from \SGARefLink{sga3:xv:proposition:8:5}{Proposition 8.5} and
Exposé \SGARefLink{sga3:ix:theorem:5:6:item:a}{IX~5.6(a)}.

For \textup{(ii)}, let \(H\) be a subgroup prescheme of \(G\). To
prove that \(H\) has property \(\mathrm{AT}\), we may suppose that
\(S\) is the spectrum of a complete discrete valuation ring. We
must show that every coherent family of \(S_n\)-group morphisms
\[
  u_n:(\mathbb G_{m,S})_n\longrightarrow H_n
\]
is admissible (\SGARefLink{sga3:xv:proposition:8:5}{Proposition 8.5} and \SGARefLink{sga3:xv:proposition:8-1}{8.1} \textup{iii bis)}).
By hypothesis, \(G\) has property \(\mathrm{AT}\), so there is an
\(S\)-group morphism
\[
  u:\mathbb G_{m,S}\longrightarrow G
\]
lifting every \(u_n\). As in the proof of \SGARefLink{sga3:xv:proposition:8:5}{Proposition 8.5}, the
restriction of \(u\) to \({}_\ell\mathbb G_{m,S}\), with \(\ell\)
ranging over powers of a prime \(q\) invertible on \(S\), factors
through \(H\). By density, \(u_t\) factors through \(H_t\) on the
generic fibre. Since \(\mathbb G_{m,S}\) is reduced, \(u\) itself
factors through \(H\).

For \textup{(iii)}, let \(S'\in\mathcal S\), with the notation of
\SGARefLink{sga3:xv:definition:8-4}{Definition 8.4}. The group \(L\) is flat over \(S\), and its generic
fibre \(L_t\) is affine. Under these conditions \(L_s\) is
necessarily affine (Exposé XVII, Appendix III~2). Since \(G\) and
\(L\) are flat, \(G=G'\), and the dimensions of the fibres of \(G\)
and \(L\) are constant on \(S\) (Exposé VI\textsubscript{B}~4).
Consequently
\[
 \operatorname{rk}_{\mathrm{ab}}G_s
 \leq \dim G_s-\dim L_s
 =\dim G_t-\dim L_t
 =\operatorname{rk}_{\mathrm{ab}}G_t.
\]
By hypothesis these inequalities are equalities. Thus
\((G_s/L_s)^0_{\mathrm{red}}\) is an abelian variety and has
reductive rank zero. It follows that \(L_s\) has the same reductive
rank as \(G_s\).

We can now prove the principal theorems of this section.

\SGAStatement{\SGAInlineRefTarget{sga3:xv:theorem:8-8}Theorem 8.8}{Let \(G\) be a group prescheme of finite
type over a locally noetherian prescheme \(S\). Suppose that \(G\)
is flat over \(S\) and has connected fibres.
\begin{enumerate}
\renewcommand{\labelenumi}{\textup{\alph{enumi})}}
\item\SGAInlineRefTarget{sga3:xv:theorem:8-8:item:a} The functor \(\mathcal T_C\) of central sub-tori of \(G\) is
representable if and only if \(G\) has property \(\mathrm{ATC}\).
In that case, \(\mathcal T_C\) is represented by an \(S\)-prescheme
étale and separated over \(S\).
\item\SGAInlineRefTarget{sga3:xv:theorem:8-8:item:b} Under the hypotheses of \textup{(a)}, for every \(S\)-torus
\(T\), the functor
\(\operatorname{Hom}^{\mathrm{cent}}_{S\text{-gr}}(T,G)\) of central
homomorphisms from \(T\) into \(G\) is represented by an
\(S\)-prescheme étale and separated over \(S\).
\end{enumerate}}
\SGARefTarget{sga3:xv:theorem:8:8}

\SGAStatement{Theorem 8.9}{Let \(S\) be a locally noetherian
prescheme and let \(G\) be an \(S\)-group prescheme of finite type,
smooth over \(S\).
\begin{enumerate}
\renewcommand{\labelenumi}{\textup{\alph{enumi})}}
\item\SGAInlineRefTarget{sga3:xv:theorem:8-9:item:a} The functor \(\mathcal T\) of sub-tori of \(G\) is
representable if and only if \(G\) has property \(\mathrm{AT}\).
In that case, \(\mathcal T\) is represented by an \(S\)-prescheme
smooth and separated over \(S\). More precisely, its structural
morphism admits a canonical factorization
\[
  \mathcal T\xrightarrow{u}Y\xrightarrow{v}S,
\]
where \(v\) is smooth and quasi-projective (hence of finite type)
and \(u\) is étale and separated.
\item\SGAInlineRefTarget{sga3:xv:theorem:8-9:item:b} Under the hypotheses of \textup{(a)}, for every \(S\)-torus
\(T\), the functor \(\operatorname{Hom}_{S\text{-gr}}(T,G)\) of
homomorphisms from \(T\) into \(G\) is represented by an
\(S\)-prescheme smooth and separated over \(S\).
\end{enumerate}}
\SGARefTarget{sga3:xv:theorem:8:9}

% END INCLUDED SGA3_English_current_buildable_source_20260823/source/inputs/XV/tex/components/08_expose_XV_section8_through_theorem8_9.tex

% BEGIN INCLUDED SGA3_English_current_buildable_source_20260823/source/inputs/XV/tex/components/09_expose_XV_section8_remainder_and_end.tex
% Authority: Expo15.pdf, local pp. 61--68 / combined pp. 931--938 /
% printed pp. 470--483.
% Completes Exposé XV and stops before Exposé XVI.

\SGAProof{Proof of \SGARefLink{sga3:xv:theorem:8-8:item:a}{Theorem 8.8(a)}}
If the functor \(\mathcal T_C\) is representable, it commutes with
adic limits of artinian rings. Consequently, by \SGARefLink{sga3:xv:remark:8:2}{8.2} and
Proposition 8.5(b), \(G\) has property \(\mathrm{ATC}\). To prove the
converse, we shall use the following result, which will be proved in
EGA VI and may also be found in Murre's exposé,
\emph{Séminaire Bourbaki}, May 1965, no.~294, Theorem~1,
Corollary~2.

\SGAStatement{\SGAInlineRefTarget{sga3:xv:lemma:8-10}Lemma 8.10}{Let \(S\) be a locally noetherian
prescheme and let \(F\) be a contravariant functor from
\(\operatorname{Sch}/S\) to the category of sets. In order that
\(F\) be representable by an \(S\)-prescheme étale and separated
over \(S\), it is necessary and sufficient that \(F\) satisfy the
following five properties:
\begin{enumerate}
\renewcommand{\labelenumi}{\textup{\roman{enumi})}}
\item \(F\) is a sheaf for the fpqc topology (Exposé IV).
\item \(F\) commutes with filtered inductive limits of rings
(Exposé \SGARefLink{sga3:xi:remark:3:2}{XI~3.2}).
\item \(F\) commutes with adic limits of local artinian rings
(\SGARefLink{sga3:xv:proposition:8-1}{8.1}\textup{(i)}).
\item \(F\) satisfies the valuative criterion for separatedness:
for every \(S\)-prescheme \(S'\) that is the spectrum of a discrete
valuation ring, with generic point \(t\), the canonical map
\[
  F(S')\longrightarrow F(t)
\]
is injective.
\item \(F\) is infinitesimally étale (Exposé \SGARefLink{sga3:xi:remark:1:8}{XI~1.8}).
\end{enumerate}}
\SGARefTarget{sga3:xv:lemma:8:10}

We show that the functor \(\mathcal T_C\) of \SGARefLink{sga3:xv:theorem:8-8}{Theorem 8.8} satisfies
the five conditions of \SGARefLink{sga3:xv:lemma:8-10}{Lemma 8.10}.

\begin{enumerate}
\renewcommand{\labelenumi}{\textup{\roman{enumi})}}
\item The functor \(\mathcal T\), and likewise \(\mathcal T_C\), is
a sheaf for the fpqc topology whenever \(G\) is of finite
presentation over \(S\). Indeed, every monomorphism from a torus
into \(G\) is then an immersion (Exposé \SGARefLink{sga3:viii:proposition:7:9}{VIII~7.9}), and the assertion
follows from fpqc descent for subpreschemes.

\item \SGARefLink{sga3:xv:corollary:6-3}{Corollary 6.3} proves that \(\mathcal T\) commutes with filtered
inductive limits of rings when \(G\) is of finite presentation over
\(S\); the same property for \(\mathcal T_C\) follows immediately.

\item By \SGARefLink{sga3:xv:proposition:8:5}{Proposition 8.5}, condition \textup{(iii)} is precisely
equivalent to property \(\mathrm{ATC}\).

\item If \(S\) is the spectrum of a discrete valuation ring, two
sub-tori of \(G\) with the same generic fibre coincide. More
precisely, each is the identity component of the schematic closure
in \(G\) of that common generic fibre.

\item This follows from \SGARefLink{sga3:xv:corollary:2:3}{Proposition 2.3} because \(G\) is flat over
\(S\).
\end{enumerate}

\SGAProof{Proof of \SGARefLink{sga3:xv:theorem:8-8:item:b}{Theorem 8.8(b)}}
Proceeding as in Exposé \SGARefLink{sga3:xi:corollary:4:2}{XI~4.2}, it is enough to prove that the
product group \(T\times_S G\) again has property
\(\mathrm{ATC}\), which is immediate from the definition.

\SGAProof{Proof of \SGARefLink{sga3:xv:theorem:8-9:item:a}{Theorem 8.9(a)}}
After replacing \(G\), if necessary, by its identity component
(Exposé VI\textsubscript{B}~\SGARefLink{sga3:xv:lemma:3:10}{3.10}), we may suppose that \(G\) has
connected fibres. If \(T\) is a sub-torus of \(G\), its centralizer
\(\operatorname{Centr}_G(T)\) is represented (Exposé \SGARefLink{sga3:xi:corollary:6:11}{XI~6.11}) by a
subgroup prescheme of \(G\), smooth over \(S\)
(Exposé \SGARefLink{sga3:xi:corollary:2:4}{XI~2.4}), with connected fibres
(Exposé \SGARefLink{sga3:xii:theorem:6-6:item:b}{XII~6.6(b)}). It is therefore an element of
\(\mathcal{CT}(S)\), where \(\mathcal{CT}\) is the functor defined
in \SGARefLink{sga3:xv:proposition:7-1:item:i}{7.1(i)}. The rule
\[
  T\longmapsto\operatorname{Centr}_G(T)
\]
clearly defines an \(S\)-morphism
\[
  u:\mathcal T\longrightarrow\mathcal{CT}.
\]
Since \(\mathcal{CT}\) is represented by an \(S\)-prescheme smooth
and quasi-projective over \(S\) (\SGARefLink{sga3:xv:proposition:7-1:item:i}{7.1(i)}), it remains only to prove
that \(u\) is representable by a separated étale morphism.

After a suitable base change \(S'\to S\), with
\(S'=\mathcal{CT}\) and hence \(S'\) locally noetherian, we are
reduced to the following problem. Let \(S\) be a locally noetherian
prescheme, let \(G\) be an \(S\)-group prescheme smooth and of
finite type over \(S\), with connected fibres, and let \(H\) be a
subgroup of \(G\), smooth over \(S\), with connected fibres.
Consider the subfunctor \(X\) of \(\mathcal T_{C,H}\) such that,
for every \(S\)-prescheme \(S'\), \(X(S')\) is the set of central
sub-tori \(T\) of \(H_{S'}\) for which
\[
  \operatorname{Centr}_{G_{S'}}(T)=H_{S'}.
\]
We shall show that \(X\) is represented by an \(S\)-prescheme
étale and separated over \(S\).

Indeed, \(G\) has property \(\mathrm{AT}\) by hypothesis, and so
does \(H\) by Proposition 8.7(ii). Since \(H\) has connected
fibres, it also has property \(\mathrm{ATC}\) by
Proposition 8.7(i). Moreover, \(H\) is smooth over \(S\), hence
flat. Theorem 8.8(a) therefore represents \(\mathcal T_{C,H}\) by
an \(S\)-prescheme étale and separated over \(S\). It is enough to
show that the canonical monomorphism
\[
  X\longrightarrow\mathcal T_{C,H}
\]
is represented by an open immersion.

Put \(S'=\mathcal T_{C,H}\), and let \(K\) be the centralizer in
\(G_{S'}\) of the universal central torus of \(H_{S'}\). The group
\(K\) is a smooth group scheme over \(S'\), with connected fibres,
and contains \(H_{S'}\). By definition,
\[
  X=\prod_{K/S'} H_{S'}/K,
\]
which is represented by the open-and-closed subprescheme of \(S'\)
over which \(H_{S'}\) and \(K\) have the same relative dimension.

The proof of \SGARefLink{sga3:xv:theorem:8-9:item:b}{Theorem 8.9(b)} is analogous to that of
\SGARefLink{sga3:xv:theorem:8-8:item:b}{Theorem 8.8(b)}.

\SGAStatement{Corollary 8.11}{Let \(S\) be a prescheme and let
\(G\) be an \(S\)-group prescheme smooth and of finite presentation
over \(S\). If the abelian rank of the fibres of \(G\) is a locally
constant function on \(S\)---in particular, if the fibres of \(G\)
are affine---then the functor of sub-tori of \(G\) is represented
by an \(S\)-prescheme smooth and separated over \(S\). The same
holds for the functor
\(\operatorname{Hom}_{S\text{-gr}}(T,G)\), for every \(S\)-torus
\(T\).}
\SGARefTarget{sga3:xv:corollary:8:11}

The assertion is local on \(S\), so we may suppose that \(S\) is
affine and that the abelian rank of the fibres of \(G\) is constant.
There are then a noetherian affine scheme \(S_0\) and an
\(S_0\)-group prescheme \(G_0\), smooth and of finite type over
\(S_0\), such that
\[
  G\simeq G_0\times_{S_0}S
\]
as \(S\)-group preschemes (Exposé VI\textsubscript{B}, \S10).
The abelian rank of the fibres of an \(S\)-group prescheme of finite
presentation is an ind-constructible function (6.3 bis). A standard
argument (EGA IV~8) therefore allows us, in the present situation,
to choose \(G_0\) so that the abelian rank of its fibres is constant
on \(S_0\). Proposition 8.7(iii) gives property \(\mathrm{AT}\) for
\(G_0\), and \SGARefLink{sga3:xv:theorem:8:8}{Theorem 8.9} then represents its functor of sub-tori by
an \(S_0\)-prescheme smooth and separated over \(S_0\). Base change
gives the assertion for \(G\). The argument for
\(\operatorname{Hom}_{S\text{-gr}}(T,G)\) is analogous.

\medskip\noindent\textbf{Generalization of \SGARefLink{sga3:xv:theorem:8:8}{Theorem 8.9}.}
The functor \(\mathcal T\) of sub-tori of a smooth group \(G\) need
not be representable. We now give sufficient conditions for a
subfunctor of \(\mathcal T\) to be representable.

\SGAStatement{\SGAInlineRefTarget{sga3:xv:proposition:8-12}Proposition 8.12}{Let \(S\) be a locally noetherian
prescheme, let \(G\) be an \(S\)-group prescheme smooth over \(S\)
with connected fibres, and let \(F\) be an \(S\)-subfunctor of the
functor \(\mathcal T\) of sub-tori of \(G\). Suppose that:
\begin{enumerate}
\renewcommand{\labelenumi}{\textup{\roman{enumi})}}
\item\SGAInlineRefTarget{sga3:xv:proposition:8-12:item:i} \(F\) is a sheaf for the fpqc topology (Exposé IV);
\item \(F\) commutes with filtered inductive limits of rings
(Exposé \SGARefLink{sga3:xi:remark:3:2}{XI~3.2});
\item \(F\) commutes with adic limits of local artinian rings
(Exposé \SGARefLink{sga3:xi:remark:3:3}{XI~3.3});
\item\SGAInlineRefTarget{sga3:xv:proposition:8-12:item:iv} \(F\) is infinitesimally smooth over \(S\)
(Exposé \SGARefLink{sga3:xi:remark:1:8}{XI~1.8});
\item\SGAInlineRefTarget{sga3:xv:proposition:8-12:item:v} \(F\) is stable under inner automorphisms of \(G\): for every
\(S\)-prescheme \(S'\),
\[
  T\in F(S'),\quad g\in G(S')
  \quad\Longrightarrow\quad
  \operatorname{int}(g)T\in F(S').
\]
\end{enumerate}
Then \(F\) is represented by an \(S\)-prescheme smooth and
separated over \(S\).}
\SGARefTarget{sga3:xv:proposition:8:12}

\SGAStatement{Proposition 8.12 bis}{Let \(S\) and \(G\) be as in
\SGARefLink{sga3:xv:proposition:8-12}{Proposition 8.12}, let \(T\) be an \(S\)-torus, and let \(F\) be a
subfunctor of \(\operatorname{Hom}_{S\text{-gr}}(T,G)\). Suppose
that \(F\) satisfies conditions \textup{(i)}--\textup{(iv)} of
\SGARefLink{sga3:xv:proposition:8-12}{Proposition 8.12} and is stable under inner automorphisms of \(G\):
for every \(S\)-prescheme \(S'\),
\[
  u\in F(S'),\quad g\in G(S')
  \quad\Longrightarrow\quad
  \operatorname{int}(g)u\in F(S').
\]
Then \(F\) is represented by an \(S\)-prescheme smooth and
separated over \(S\).}
\SGARefTarget{sga3:xv:proposition:8:12:bis}

\SGAProof{Proof of \SGARefLink{sga3:xv:proposition:8-12}{Proposition 8.12}}
The proof of \SGARefLink{sga3:xv:proposition:8:12}{Proposition 8.12 bis} is entirely analogous and is left
to the reader. Let
\[
  u:F\longrightarrow\mathcal{CT}
\]
be the \(S\)-morphism which assigns to every
\(T\in F(S')\) the element
\(\operatorname{Centr}_{G_{S'}}(T)\) of
\(\mathcal{CT}(S')\) (\SGARefLink{sga3:xv:proposition:7-1:item:i}{7.1(i)}). Since \(\mathcal{CT}\) is represented
by an \(S\)-prescheme smooth and quasi-projective over \(S\), we
are reduced to proving that \(u\) is representable.

After the base change \(\mathcal{CT}\to S\), the problem is the
following. Given \(S\) and \(G\) as above and a smooth subgroup
\(H\) of \(G\) with connected fibres, represent the functor \(F_H\)
defined by
\[
  F_H(S')=
  \left\{T\in F(S')\ \middle|\
  H_{S'}=\operatorname{Centr}_{G_{S'}}(T)\right\}.
\]
We show that \(F_H\) is étale and separated over \(S\) by verifying
the five conditions of \SGARefLink{sga3:xv:lemma:8-10}{Lemma 8.10}.

The first three conditions follow directly from
\SGARefLink{sga3:xv:proposition:8-12:item:i}{Proposition 8.12(i)}--(iii). We have already observed that
\(\mathcal T_{C,H}\) is separated and infinitesimally étale.
Thus \(F_H\), being a subfunctor of \(\mathcal T_{C,H}\), is
separated and infinitesimally unramified (Exposé \SGARefLink{sga3:xi:remark:1:8}{XI~1.8}). It remains
to prove that \(F_H\) is infinitesimally smooth.

Let \(S\) be the spectrum of a local artinian ring, let
\(S_0\subseteq S\) be defined by a nilpotent ideal, and let
\(T_0\in F_H(S_0)\). By \SGARefLink{sga3:xv:proposition:8-12:item:iv}{Proposition 8.12(iv)}, \(T_0\) lifts to an
element \(T'\in F(S)\). Since \(H\) is smooth, \(T_0\) also lifts
to a sub-torus \(T''\) of \(H_S\) (Exposé \SGARefLink{sga3:ix:theorem:3:6:bis}{IX~3.6 bis}). This
sub-torus is conjugate to \(T'\) by an element \(g\in G(S)\)
(\emph{loc. cit.}), and hence \(T''\in F(S)\) by
\SGARefLink{sga3:xv:proposition:8-12:item:v}{Proposition 8.12(v)}. The centralizer
\(\operatorname{Centr}_{G_S}(T'')\) is smooth over \(S\) and agrees
with \(H_{S_0}\) over \(S_0\). Therefore \(T''\) lies in the centre
of \(H_S\) (Exposé \SGARefLink{sga3:ix:theorem:5:6:item:a}{IX~5.6(a)}), and its centralizer in \(G\) equals
\(H_S\). Thus \(T''\in F_H(S)\), and we take \(T=T''\).

\SGAStatement{Corollary 8.13}{Let \(S\) be a prescheme and let
\(G\) be an \(S\)-group prescheme smooth and of finite presentation
over \(S\), with connected fibres. Let \(\mathcal T_D\) be the
subfunctor of \(\mathcal T\) for which \(\mathcal T_D(S')\) is the
set of sub-tori \(T\) of \(G_{S'}\) such that, for every
\(s'\in S'\), \(T_{s'}\) is contained in the derived subgroup of
\(G_{s'}\) (Exposé VI\textsubscript{B}~7). Then \(\mathcal T_D\)
is represented by an \(S\)-prescheme smooth and separated over
\(S\).}
\SGARefTarget{sga3:xv:corollary:8:13}

\SGAStatement{Corollary 8.13 bis}{Let \(S\) and \(G\) be as in
\SGARefLink{sga3:xv:corollary:8:13}{Corollary 8.13}, and let \(T\) be an \(S\)-torus. Let
\[
  \operatorname{Hom}^{\mathrm{der}}_{S\text{-gr}}(T,G)
\]
be the subfunctor of \(\operatorname{Hom}_{S\text{-gr}}(T,G)\)
whose \(S'\)-points are the \(S'\)-morphisms
\(u:T_{S'}\to G_{S'}\) such that, for every \(s'\in S'\), the
morphism \(u_{s'}\) factors through the derived group of
\(G_{s'}\). Then this functor is represented by an \(S\)-prescheme
smooth and separated over \(S\).}
\SGARefTarget{sga3:xv:corollary:8:13:bis}

The two corollaries have analogous proofs; we prove
\SGARefLink{sga3:xv:corollary:8:13:bis}{Corollary 8.13 bis}. By the usual procedure, we reduce to the case
in which \(S\) is noetherian. We verify the five conditions of
\SGARefLink{sga3:xv:proposition:8:12}{Proposition 8.12 bis}.

Conditions \textup{(i)} and \textup{(iv)} follow immediately from
the corresponding properties of the functor
\[
  \operatorname{Hom}_{S\text{-gr}}(T,G).
\]
Condition \textup{(v)} holds because the derived group of an
algebraic group is invariant (Exposé VI\textsubscript{B}, \S7).

For condition \textup{(ii)}, a standard reduction (EGA IV~8)
reduces us to the following assertion. Let \(S\) be an integral
noetherian scheme with generic point \(\eta\), and let
\(u:T\to G\) be an \(S\)-group morphism whose generic fibre factors
through the derived group of \(G_\eta\). Then there is a
neighbourhood \(U\) of \(\eta\) such that, for every \(s\in U\),
\(u_s\) factors through the derived group of \(G_s\). This is an
immediate consequence of Exposé VI\textsubscript{B}~\SGARefLink{prop:sga3-VIB-10.12}{10.12}.

For condition \textup{(iii)}, retain the notation of \SGARefLink{sga3:xv:proposition:4:3}{Proposition
4.3}. We must show that an element
\[
  (u_m)_{m\in\mathbb N}\in
  \varprojlim_m
  \operatorname{Hom}^{\mathrm{der}}_{S_m\text{-gr}}(T_m,G_m)
\]
is admissible in the sense of \SGARefLink{sga3:xv:proposition:4:3}{Proposition 4.3} and comes from an
element of
\(\operatorname{Hom}^{\mathrm{der}}_{S\text{-gr}}(T,G)\).
As in \SGARefLink{sga3:xv:proposition:8-1}{8.1}, we reduce immediately to the case in which \(S\) is the
spectrum of a complete discrete valuation ring. Let \(t\) be its
generic point, \(s\) its closed point, and let \(D\) be the
schematic closure in \(G\) of the derived group \(D_t\) of \(G_t\).

\begin{enumerate}
\renewcommand{\labelenumi}{\textup{\alph{enumi})}}
\item The group \(D_s\) contains the derived group of \(G_s\).
Indeed, \(D_t\) is invariant in \(G_t\), and \(G\) is flat over
\(S\), so \(D\) is invariant in \(G\). Moreover, the morphism
\[
  G\times_S G\longrightarrow G,
  \qquad (x,y)\longmapsto xyx^{-1}y^{-1},
\]
factors through \(D_t\) on the generic fibre and therefore factors
through \(D\). Consequently \(G_s/D_s\) is commutative.

\item If
\(u_m\in
\operatorname{Hom}^{\mathrm{der}}_{S_m\text{-gr}}(T_m,G_m)\),
then \(u_m\) factors through \(D_m\). Its closed fibre \(u_s\)
factors through the derived group of \(G_s\), hence through \(D_s\)
by \textup{(a)}. Since \(D_m\) is flat over \(S_m\) and invariant
in \(G_m\), the quotient
\[
  H_m=G_m/D_m
\]
is representable (Exposé VI\textsubscript{A}, \S4). The image of
\(T_m\) in \(H_m\) is a torus (Exposé \SGARefLink{sga3:ix:theorem:6:8}{IX~6.8}) with trivial closed
fibre, hence is trivial. Thus \(u_m\) factors through \(D_m\).

\item The family \((u_m)_{m\in\mathbb N}\) is admissible and lifts
to a morphism \(T\to G\) belonging to
\(\operatorname{Hom}^{\mathrm{der}}_{S\text{-gr}}(T,G)\).
By what precedes and \SGARefLink{sga3:xv:theorem:4-1-bis}{Proposition 4.1 bis}, it is enough to prove that
\(D_t\) is affine. This follows from the next lemma.
\end{enumerate}

\SGAStatement{Lemma 8.14}{Let \(G\) be a smooth connected
algebraic group over a field \(k\). Then the derived group \(D\) of
\(G\) is affine.}
\SGARefTarget{sga3:xv:lemma:8:14}

Formation of \(D\) commutes with extension of the base field
(Exposé VI\textsubscript{B}~7), so we may suppose that \(k\) is
algebraically closed. The group \(G\) is then an extension of an
abelian variety \(A\) by a linear group \(L\). Since \(A\) is
commutative, \(D\) is necessarily contained in \(L\), and is
therefore affine.

\medskip\noindent\textbf{Maximal tori.}

\SGAStatement{\SGAInlineRefTarget{sga3:xv:theorem:8-15}Theorem 8.15}{Let \(S\) be a locally noetherian
prescheme and let \(G\) be an \(S\)-group prescheme smooth and of
finite type over \(S\). The following conditions are equivalent:
\begin{enumerate}
\renewcommand{\labelenumi}{\textup{\roman{enumi})}}
\item The \(S\)-functor \(\mathcal T_M\), whose \(S'\)-points are
the maximal tori of \(G_{S'}\) (Exposé \SGARefLink{sga3:xii:1.3}{XII~1.3}), is representable.
\item The functor \(\mathcal T_M\) is represented by an
\(S\)-prescheme smooth and quasi-projective over \(S\), with affine
fibres.
\item The group \(G\) admits a maximal torus locally for the étale
topology.
\item The group \(G\) admits a maximal torus locally for the
faithfully flat topology.
\item\SGAInlineRefTarget{sga3:xv:theorem:8-15:item:v} The group \(G\) has property \(\mathrm{AT}\), and the
reductive rank of its fibres is a locally constant function on
\(S\).
\end{enumerate}}
\SGARefTarget{sga3:xv:theorem:8:15}

\SGAProof{Proof}
\textup{(ii)}\(\Rightarrow\)\textup{(i)} is clear.

\textup{(i)}\(\Rightarrow\)\textup{(iii)}.
Since \(G\) is of finite presentation over \(S\), \SGARefLink{sga3:xv:corollary:6-3}{Corollary 6.3}
shows that \(\mathcal T_M\) commutes with filtered inductive limits
of rings. If representable, it is therefore locally of finite
presentation (EGA IV~8.14). Moreover, \(\mathcal T_M\) is formally
smooth (Exposé XI~2.1 bis). Hence its representing prescheme is
smooth over \(S\), and \textup{(iii)} follows from Hensel's lemma
(Exposé \SGARefLink{sga3:xi:proposition:1:10}{XI~1.10}).

\textup{(iii)}\(\Rightarrow\)\textup{(iv)} is clear.

\textup{(iv)}\(\Rightarrow\)\textup{(v)}.
Let \(s\in S\). There is an \(S\)-prescheme \(S'\), flat over \(S\)
and with image containing \(s\), such that \(G_{S'}\) has a maximal
torus \(T'\). If \(s'\in S'\) lies over \(s\), the reductive rank of
the fibres of \(G_{S'}\) is constant on
\(\operatorname{Spec}\mathcal O_{S',s'}\). Thus the reductive rank
of the fibres of \(G\) is constant on its image
\(\operatorname{Spec}\mathcal O_{S,s}\)
(EGA IV~2.3.4(ii)). Since this rank is a locally constructible
function on \(S\) (6.3 bis), it is constant on a neighbourhood of
\(s\) (EGA IV~1.10.1).

To prove property \(\mathrm{AT}\), let \(S_1\) be an \(S\)-scheme
which is the spectrum of a discrete valuation ring, whose closed
point \(s_1\) maps to \(s\). The prescheme
\[
  S'_1=S_1\times_S S'
\]
is faithfully flat over \(S_1\), and \(G_{S'_1}\) has a maximal
torus. Let \(A\) and \(A'\) be the rings of \(S_1\) and \(S'_1\),
respectively. Writing \(A'\) as the filtered inductive limit of its
finitely generated \(A\)-subalgebras and applying \SGARefLink{sga3:xv:corollary:6-3}{Corollary 6.3},
we obtain an \(S_1\)-scheme \(S''_1\) such that
\(G_{S''_1}\) has a maximal torus and
\(S''_1\to S_1\) is of finite type and surjective. By
EGA II~7.1.9, we may reduce to the case in which \(S'_1\) is the
spectrum of a discrete valuation ring. Then \(G_{S'_1}\), and
hence \(G_{S_1}\), has property \(\mathrm{AT}\). Since this holds
for every such \(S_1\), \(G\) has property \(\mathrm{AT}\).

\textup{(v)}\(\Rightarrow\)\textup{(i)}.
By \SGARefLink{sga3:xv:theorem:8:8}{Theorem 8.9} the functor \(\mathcal T\) of sub-tori of \(G\) is
representable. If the reductive rank of \(G\) is the constant
function \(r\), then \(\mathcal T_M\) is represented by the
open-and-closed subprescheme of \(\mathcal T\) representing the
subfunctor of tori of rank \(r\).

\textup{(iii)}\(\Rightarrow\)\textup{(ii)}.
Under condition \textup{(iii)}, the results of Exposé \SGARefLink{sga3:xii:7.1}{XII~7.1}
identify \(\mathcal T_M\) canonically with the functor of Cartan
subgroups of \(G\). The assertion follows from Theorem 7.3(i).

\medskip\noindent\textbf{Remark 8.16.}
The prescheme \(\mathcal T_M\) of maximal tori of \(G\) can in fact
be shown to be affine over \(S\),\footnote{Cf. M. Raynaud,
\emph{Faisceaux amples sur les schémas en groupes et les espaces
homogènes} (thesis, forthcoming). Editor's note: see
\emph{Lecture Notes in Mathematics} 119 (1970), Springer,
especially \SGARefLink{sga3:ix:corollary:2:9}{IX~2.9}.} generalizing Exposé \SGARefLink{sga3:xii:5.4}{XII~5.4}.
\SGARefTarget{sga3:xv:remark:8:16}

\SGAStatement{Corollary 8.17}{Let \(S\) be a prescheme and let
\(G\) be an \(S\)-group prescheme smooth and of finite presentation
over \(S\). If the abelian and reductive ranks of the fibres of
\(G\) are locally constant functions on \(S\), then \(G\) satisfies
the equivalent conditions \textup{(i)}--\textup{(iv)} of
\SGARefLink{sga3:xv:theorem:8-15}{Theorem 8.15}.}
\SGARefTarget{sga3:xv:corollary:8:17}

We may suppose that the abelian and reductive ranks of the fibres
of \(G\) are constant. Proceeding as in \SGARefLink{sga3:xv:corollary:8:11}{Corollary 8.11} and using
\SGARefLink{sga3:xv:corollary:6:3:bis}{6.3 bis}, we reduce to the case in which \(S\) is noetherian. Then
\(G\) has property \(\mathrm{AT}\) by \SGARefLink{sga3:xv:proposition:8:7}{Proposition 8.7}, and the
conclusion follows from \SGARefLink{sga3:xv:theorem:8-15:item:v}{Theorem 8.15(v)}.

\SGAStatement{\SGAInlineRefTarget{sga3:xv:corollary:8-18}Corollary 8.18}{Let \(S\) be a prescheme and let
\(G\) be an \(S\)-group prescheme smooth over \(S\). The following
conditions are equivalent:
\begin{enumerate}
\renewcommand{\labelenumi}{\textup{\alph{enumi})}}
\item The unipotent rank \(\rho_u\) and the abelian rank
\(\rho_{\mathrm{ab}}\) of the fibres of \(G\) are locally constant
functions on \(S\).
\item The unipotent rank \(\rho_u\) is locally constant, and \(G\)
admits maximal tori locally for the fpqc topology.
\item The reductive rank \(\rho_r\) and the abelian rank
\(\rho_{\mathrm{ab}}\) of the fibres of \(G\) are locally constant
functions on \(S\).
\end{enumerate}}
\SGARefTarget{sga3:xv:corollary:8:18}

\medskip\noindent\textbf{Remark 8.19.}
Under the hypotheses of \SGARefLink{sga3:xv:corollary:8-18}{Corollary 8.18}, a more precise argument
using lower semicontinuity of the abelian rank, announced in
Exposé \SGARefLink{sga3:x:remark:8-7}{X~8.7}, shows that if two of the three ranks
\(\rho_u,\rho_r,\rho_{\mathrm{ab}}\) are locally constant, then so
is the third.
\SGARefTarget{sga3:xv:remark:8:19}

\SGAProof{Proof of \SGARefLink{sga3:xv:corollary:8-18}{Corollary 8.18}}
After replacing \(G\), if necessary, by its identity component, we
may suppose that \(G\) is of finite presentation over \(S\)
(Exposé VI\textsubscript{B}~\SGARefLink{par:sga3-VIA-5.3.3}{5.3.3}).

\textup{(a)}\(\Rightarrow\)\textup{(c)}.
Let \(s\in S\). Since \(\rho_{\mathrm{ab}}\) is locally constant,
\SGARefLink{sga3:xv:corollary:8:11}{Corollary 8.11} allows us, after an étale extension covering \(S\),
to suppose that \(G\) has a sub-torus \(T\) whose fibre \(T_s\) is
a maximal torus of \(G_s\). Put
\[
  C=\operatorname{Centr}_G(T).
\]
This is a subgroup prescheme of \(G\), smooth over \(S\), with
connected fibres. For every \(t\in S\), the group \(C_t\) contains
a Cartan subgroup \(C'_t\) of \(G_t\). After shrinking \(S\), we
may suppose that \(\rho_u\), \(\rho_{\mathrm{ab}}\), and the
relative dimension of \(C\) over \(S\) are constant. We then have
\[
\begin{aligned}
  \dim C_s=\dim C_t
  &\geq \dim C'_t\\
  &\geq \rho_u(t)+\rho_{\mathrm{ab}}(t)+\dim T_t\\
  &=\rho_u(s)+\rho_{\mathrm{ab}}(s)+\dim T_s\\
  &=\rho_\nu(s)=\dim C_s.
\end{aligned}
\]
It follows that \(C_t=C'_t\). Thus \(T\) is a maximal torus of
\(G\), and in particular \(\rho_r(t)=\rho_r(s)\).

\textup{(c)}\(\Rightarrow\)\textup{(b)} follows from
\SGARefLink{sga3:xv:corollary:8:17}{Corollary 8.17}.

\textup{(b)}\(\Rightarrow\)\textup{(a)}.
Since \(G\) admits maximal tori locally for the fpqc topology, it
also admits Cartan subgroups locally for that topology. Hence, by
Exposé \SGARefLink{sga3:xii:7.3}{XII~7.3}, the nilpotent rank
\[
  \rho_\nu=\rho_u+\rho_r+\rho_{\mathrm{ab}}
\]
is locally constant. The existence of maximal tori makes
\(\rho_r\) locally constant, while \(\rho_u\) is locally constant
by hypothesis; therefore \(\rho_{\mathrm{ab}}\) is locally
constant.\footnote{The final sentence of the printed French
interchanges \(\rho_u\) and \(\rho_{\mathrm{ab}}\). The displayed
identity and condition \textup{(b)} give the implication in the
form written here.}

% END INCLUDED SGA3_English_current_buildable_source_20260823/source/inputs/XV/tex/components/09_expose_XV_section8_remainder_and_end.tex



% END INCLUDED SGA3_English_current_buildable_source_20260823/source/inputs/XV/body.tex


\SGAStartExpose{XVI}
% BEGIN INCLUDED SGA3_English_current_buildable_source_20260823/source/inputs/XVI/body.tex
% Mechanically extracted after \mainmatter from the frozen standalone master.

\markboth{Exposé XVI. Groups of Unipotent Rank Zero}
  {Exposé XVI. Groups of Unipotent Rank Zero}
% BEGIN INCLUDED SGA3_English_current_buildable_source_20260823/source/inputs/XVI/tex/components/01_expose_XVI_section1_immersion_criterion.tex
% SGA 3, Exposé XVI, Loop-1 English reconstruction.
% Current scope: title and all of Section 1.
% Authority: Polo--Gille Expo16.pdf, local pages 1--10
% (combined-reader pages 939--948; printed pages 484--502).
% Local GPU OCR and the Reinhold English Markdown are comparison/locator
% witnesses only. The Polo--Gille page image controls text and formulas.

\setcounter{footnote}{0}
\SGARefTarget{sga3:xvi:chapter:xvi}\chapter*{Exposé XVI\\
Groups of Unipotent Rank Zero}
\addcontentsline{toc}{chapter}{Exposé XVI. Groups of Unipotent Rank Zero}
\markboth{EXPOSÉ XVI. GROUPS OF UNIPOTENT RANK ZERO}
  {EXPOSÉ XVI. GROUPS OF UNIPOTENT RANK ZERO}

\begin{center}
  \textit{by M. Raynaud}\textsuperscript{(*)}
\end{center}
\begingroup
\renewcommand{\thefootnote}{0}
\footnotetext{xy version of December 1, 2008.}
\renewcommand{\thefootnote}{*}
\footnotetext{Cf.\ the note on page 1 of Exposé XV.}
\endgroup

\SGARefTarget{sga3:xvi:section:1}\section*{1. An Immersion Criterion}
\addcontentsline{toc}{section}{1. An Immersion Criterion}

\SGARefTarget{sga3:xvi:subsection:1:1}%
\subsection*{1.1. Examples of Monomorphisms of Group Preschemes
That Are Not Immersions}
\addcontentsline{toc}{subsection}{\SGARefLink{sga3:xvi:subsection:1:1}{1.1}. Examples of Monomorphisms of
Group Preschemes That Are Not Immersions}

We shall construct a prescheme \(S\), two \(S\)-group preschemes
\(G\) and \(H\), and an \(S\)-group monomorphism
\[
  u:G\longrightarrow H
\]
that is not an immersion. The groups \(G\) and \(H\) will be of finite
presentation over \(S\), flat over \(S\), and \(G\) will even be
smooth over \(S\) (cf.\ Exposé VIII, the paragraph and note~\((*)\)
preceding \SGARefLink{sga3:viii:lemma:7:1}{7.1}; see also XVII, Appendix III, 4).

First take \(S\) to be the spectrum of a discrete valuation ring
\(A\) of unequal characteristic and residue characteristic \(2\).
Let \(t\) be the generic point of \(S\) and \(s\) the closed point.

\medskip
\noindent\textit{a)} Take for \(G\) the open subgroup of the constant
group
\[
  G'=(\mathbb Z/2\mathbb Z)_S
\]
obtained by removing the point of the closed fiber
\((\mathbb Z/2\mathbb Z)_s\) distinct from the identity element, and
take for \(H\) the group of multiplicative type
\((\mu_2)_S\) of square roots of unity (\SGARefLink{sga3:i:paragraph:4-4-4}{Exposé I, 4.4.4}). By the
choice of \(S\), \(H_t\) is isomorphic to
\((\mathbb Z/2\mathbb Z)_t\), whereas \(H_s\) is a radicial group.
There is an evident morphism \(G'\to H\), defined by the section
\((-1)\) of \(H\), hence a morphism \(u:G\to H\) which is a
monomorphism but not an immersion (otherwise it would necessarily be
an isomorphism; cf.\ \SGARefLink{sga3:viii:section:7}{Exposé VIII, 7}). Here \SGARefLink{sga3:viii:proposition:7:9}{Exposé VIII, 7.9} does not
apply, since \({}_{2}G=G\) is not finite over \(S\).

\medskip
\noindent\textit{b)} Let \(K\) denote the étale, nonseparated
\(S\)-group prescheme obtained from the identity group by
``doubling'' the unique point of the closed fiber. Take for \(H\) the
product group
\[
  (\mu_2)_S\times_S K.
\]
Let \(a\) (respectively \(b\)) be the unique section of \(\mu_2\)
(respectively \(K\)) over \(S\) distinct from the identity section,
and let \(h=(a,b)\) be the corresponding section of \(H\). The datum
of \(h\) defines an \(S\)-group morphism
\[
  u:G=(\mathbb Z/2\mathbb Z)_S\longrightarrow H.
\]
It is clear that \(u\) is a monomorphism and is not an immersion.
Here \SGARefLink{sga3:viii:proposition:7:9}{Exposé VIII, 7.9} does not apply, since \({}_{2}H=H\) is not
separated over \(S\).

\medskip
\noindent\SGAInlineRefTarget{sga3:xvi:subsection:1:1:item:c}\textit{c)} More interesting is the following example, in
which \(G\) is a smooth group with connected fibers (and even
\(G=(\Ga)_S\)); cf.\ \SGARefLink{sga3:viii:remark:7:10}{Exposé VIII, 7.10}.

Take for \(S\) the spectrum of a discrete valuation ring \(A\) of
equal characteristic \(2\), and let \(\pi\) be a uniformizer of
\(A\) and \(\mathfrak m\) the maximal ideal of \(A\).

Consider the additive group \((\Ga)_S\), the product group
\((\Ga\times\Ga)_S\) with ring \(A[X,Y]\), and let \(G_1\) be the
closed subgroup defined by
\begin{equation}
  X^2+X-\pi Y=0.
  \tag{1}\label{sga3:xvi:eq:1}
\end{equation}
The group \(G_1\) is smooth over \(S\); its generic fiber is
isomorphic to \((\Ga)_t\), and its closed fiber is isomorphic to the
product of \((\Ga)_s\) with \((\mathbb Z/2\mathbb Z)_s\). For the
\((\mathbb Z/2\mathbb Z)_s\)-factor one may take the subgroup \(N_s\)
of \((G_1)_s\) defined by
\begin{equation}
  Y=X^2+X=0.
  \tag{2}\label{sga3:xvi:eq:2}
\end{equation}

Let \(N'\) and \(N''\) be two group subschemes of \(G_1\), isomorphic
to \((\mathbb Z/2\mathbb Z)_S\), whose closed fibers coincide with
\(N_s\) and whose generic fibers are distinct. Thus \(N'\) and
\(N''\) are defined by two sections of \(G_1\) over \(S\), with
coordinates \((a',b')\) and \((a'',b'')\), respectively, such that
\begin{equation}
  \begin{cases}
    a'^2+a'-\pi b'=a''^2+a''-\pi b''=0,\\
    a'\text{ and }a''\equiv 1\pmod{\mathfrak m},
      \qquad a'\ne a'',\\
    b'\text{ and }b''\equiv 0\pmod{\mathfrak m}.
  \end{cases}
  \tag{3}\label{sga3:xvi:eq:3}
\end{equation}
(For example, one may take \((1,0)\) and
\((1+\pi^2,\pi+\pi^3)\).)

The quotient groups
\[
  G'=G_1/N'
  \qquad\text{and}\qquad
  G''=G_1/N''
\]
are representable (\SGARefLink{thm:sga3-V-4.1}{Exposé V, 4.1}). Let \(u\) be the canonical
morphism
\[
  u:G_1\longrightarrow G'\times_S G'',
\]
and let \(H\) be the group subscheme of \(G'\times_S G''\) equal to
the schematic closure in \(G'\times_S G''\) of
\(u_t((G_1)_t)\), so that \(u\) factors through \(H\). The kernel of
\(u\) is
\[
  K=N'\cap N'',
\]
whose generic fiber is the identity group and whose closed fiber is
equal to \(N_s\); in particular, \(K\) is not flat over \(S\). On the
generic fiber, \(u_t\) is therefore an isomorphism from
\((G_1)_t\) onto \(H_t\). On the other hand, \(H_s\) is not smooth.
Indeed, if \(H_s\) were smooth, then, since \(u_s\) is a morphism with
finite kernel and \((G_1)_s\) and \(H_s\) have the same dimension,
namely \(1\), \(u_s\) would be flat. Since \(G_1\) and \(H\) are flat,
\(u\) would then be flat, and consequently \(\Ker u\) would be flat
over \(S\), which it is not. It is clear that the restriction of
\(u\) to the identity component \(G\) of \(G_1\) is a monomorphism:
the fibers of \(K\cap G\) are identity groups, so \(K\cap G\) is the
identity group (Exposé VI\textsubscript{B}, \SGARefLink{lem:sga3-VIB-2.9}{2.9}). But it is not an
immersion, for otherwise it would be an open immersion and \(H_s\)
would be smooth. Here \SGARefLink{sga3:viii:proposition:7:9}{Exposé VIII, 7.9} does not apply, since
\({}_{2}G=G\) is not finite over \(S\).

For lovers of equations, in the preceding example one may take for
\(H\) the closed subgroup of \((\Ga\times_S\Ga)_S\) with ring
\(A[V,W]/(F)\), where
\[
  F=(a''b'-b''a')(a''^2V-a'^2W)-(a''V-a'W)^2
\]
(where \(a',a'',b',b''\) satisfy \eqref{sga3:xvi:eq:3}). For \(G\)
one takes the group \((\Ga)_S\) with ring \(A[T]\), and for the
morphism \(u:G\to H\) the \(S\)-morphism defined by
\[
  V\longmapsto a'(a'T+\pi T^2),
  \qquad
  W\longmapsto a''(a''T+\pi T^2).
\]

\medskip
\noindent\SGARefTarget{sga3:xvi:remark:1:2}%
\textbf{Remark 1.2.} --- The preceding construction is inspired by
the method of Koizumi--Shimura.\footnote{Editor's note:
S.~Koizumi and G.~Shimura, ``Specialization of abelian varieties,''
\emph{Sci.\ Papers Coll.\ Gen.\ Ed.\ Univ.\ Tokyo} 9 (1959),
187--211.} It does not apply when \(S\) is the spectrum of a ring of
unequal characteristic, because the points of finite order of \(G\)
are then ``too rigid.''

\SGARefTarget{sga3:xvi:subsection:1:3}%
\subsection*{1.3. Statement of the Immersion Criterion}
\addcontentsline{toc}{subsection}{1.3. Statement of the Immersion Criterion}

\medskip
\noindent\SGARefTarget{sga3:xvi:theorem:1:3}%
\textbf{Theorem 1.3.} --- \textit{Let \(S\) be a locally noetherian
prescheme, let \(G\) be an \(S\)-group prescheme, smooth and of finite
type over \(S\), possessing Cartan subgroups locally for the fpqc
topology (\SGARefLink{sga3:xv:definition:6-1}{Exposé XV, 6.1} and \SGARefLink{sga3:xv:theorem:7-3}{7.3}~(i)), let \(H\) be an \(S\)-group
prescheme locally of finite type, and let \(u:G\to H\) be an
\(S\)-group monomorphism. Then:}

\begin{enumerate}
\item[\textit{a)}] \textit{If \(G\) has connected fibers, then \(u\)
is an immersion if and only if, for every \(S\)-scheme \(S'\) that is
the spectrum of a complete discrete valuation ring with algebraically
closed residue field, and every Cartan subgroup \(C\) of \(G_{S'}\),
the restriction of \(u_{S'}\) to \(C\) is an immersion.}

\item[\textit{b)}] \textit{The morphism \(u\) is an immersion
(respectively a closed immersion) if and only if, for every \(S'\) as
above and every Cartan subgroup \(C\) of the identity component
\((G_{S'})^0\) of \(G_{S'}\), the restriction of \(u_{S'}\) to the
normalizer \(N\) of \(C\) in \(G_{S'}\) (\SGARefLink{sga3:xi:corollary:6:11}{Exposé XI, 6.11}) is an
immersion (respectively a closed immersion).}
\end{enumerate}

Before proving \SGARefLink{sga3:xvi:theorem:1:3}{1.3}, let us state some applications.

\medskip
\noindent\SGARefTarget{sga3:xvi:corollary:1:4}%
\textbf{Corollary 1.4.} --- \textit{Let \(S\) be a prescheme, let
\(G\) be an \(S\)-group prescheme, smooth and of finite presentation
over \(S\), of unipotent rank zero (\SGARefLink{sga3:xv:definition:6:1:ter}{Exposé XV, 6.1 ter}), possessing
maximal tori locally for the fpqc topology, let \(H\) be an
\(S\)-group prescheme, and let \(u:G\to H\) be an \(S\)-group
monomorphism. Suppose in addition either that \(H\) is of finite
presentation over \(S\), or that \(S\) is locally noetherian and
\(H\) is locally of finite type. Then:}

\begin{enumerate}
\item[\textit{a)}] \textit{If \(G\) has connected fibers, \(u\) is
an immersion.}

\item[\textit{b)}] \textit{If \(H\) is separated over \(S\), and if
for every \(S'\) over \(S\) and every maximal torus \(T\) of
\(G_{S'}\), the Weyl group}
\[
  W=
  \Norm_{G_{S'}}(T)\big/\Centr_{G_{S'}^0}(T)
\]
\textit{is representable (a condition always satisfied when \(G\)
has connected fibers; \SGARefLink{sga3:xv:proposition:7-1}{Exposé XV, 7.1}~(iv)) and finite over \(S\), then
\(u\) is a closed immersion.}
\end{enumerate}

\medskip
\noindent\SGARefTarget{sga3:xvi:corollary:1:5}%
\textbf{Corollary 1.5.} --- \textit{Let \(S\) be a prescheme, let
\(G\) be an \(S\)-group prescheme, smooth and of finite presentation
over \(S\), with connected fibers, let \(H\) be an \(S\)-group
prescheme, and let \(u:G\to H\) be an \(S\)-group monomorphism.
Suppose either that \(H\) is of finite presentation over \(S\), or
that \(S\) is locally noetherian and \(H\) is locally of finite type.
Then:}

\begin{enumerate}
\item[\textit{a)}] \textit{If \(G\) is reductive, that is, affine
over \(S\) with reductive fibers (cf.\ Exposé XIX), \(u\) is an
immersion, and a closed immersion if \(H\) is separated.}

\item[\textit{b)}] \textit{If \(G\) has affine, solvable, connected
fibers of unipotent rank zero, \(u\) is an immersion, and a closed
immersion if \(H\) is separated.}
\end{enumerate}

\medskip
\noindent\textit{Proof of \SGARefLink{sga3:xvi:corollary:1:4}{1.4} from \SGARefLink{sga3:xvi:theorem:1:3}{1.3}.}

\smallskip
\noindent\textit{Reduction to the noetherian case.}
If \(S\) is locally noetherian, the reduction is immediate, since the
properties to be proved are local on \(S\). In the second case \(G\)
and \(H\) are of finite presentation over \(S\). After shrinking
\(S\), we may suppose that \(S\) is affine. By Exposé
VI\textsubscript{B}, \S10, there exist a noetherian scheme \(S_0\), an
\(S_0\)-group prescheme \(G_0\), smooth over \(S_0\) (with connected
fibers in case~a)) and of finite type over \(S_0\), an
\(S_0\)-group prescheme \(H_0\) of finite type (separated in case~b)),
and an \(S_0\)-group morphism
\[
  u_0:G_0\longrightarrow H_0
\]
such that \(G,H,u\) are obtained from \(G_0,H_0,u_0\) by the base
extension \(S\to S_0\). The assertion that \(u\) is a monomorphism
means that \(\Ker u\) is the identity group; we may therefore suppose
that \(u_0\) is a monomorphism.

Since \(G\) has unipotent rank \(\rho_u=0\) and possesses maximal
tori locally for the fpqc topology, the abelian rank \(\rho_{ab}\) of
the fibers of \(G\) is locally constant (\SGARefLink{sga3:xv:corollary:8-18}{Exposé XV, 8.18}). But
\(\rho_u\) and \(\rho_{ab}\) are locally constructible functions
(Exposé XV, 6.3 bis). A standard argument (cf.\ EGA IV, 8.3.4) shows
that one may choose \(S_0\) and \(G_0\) so that the unipotent rank
(respectively the abelian rank) of the fibers of \(G_0\) is zero
(respectively locally constant). Then \(G_0\) possesses maximal tori
locally for the fpqc topology (\SGARefLink{sga3:xv:corollary:8-18}{Exposé XV, 8.18}), and the functor
\(\mathcal{TM}_{G_0}\) of maximal tori of \(G_0\) is representable by
an \(S_0\)-scheme \(X_0\) of finite type over \(S_0\) (\SGARefLink{sga3:xv:theorem:8-15}{Exposé XV,
8.15}). Let \(T_0\) be the ``universal'' maximal torus of
\((G_0)_{X_0}\), let \(X=X_0\times_{S_0}S\), and let
\(T=T_0\times_{S_0}S\), the universal maximal torus for \(G\). By
hypothesis, in case~b), the Weyl group relative to \(T\) is
representable and finite over \(X\). These two properties are
compatible with projective limits of preschemes (Exposé
VI\textsubscript{B}, 10.1~iii, and EGA IV, 8.10.5). We may therefore
choose \(S_0\) so that the Weyl group of \(T_0\) is finite over
\(X_0\). Under these conditions, to prove \SGARefLink{sga3:xvi:corollary:1:4}{1.4} we may replace
\(S,G,u,H\) by \(S_0,G_0,u_0,H_0\), and hence suppose that \(S\) is
noetherian.

Use the valuative criterion furnished by \SGARefLink{sga3:xvi:theorem:1:3}{1.3}. We are reduced to the
case where \(S\) is the spectrum of a discrete valuation ring and
where \(G^0\) possesses a Cartan subgroup \(C\).

\smallskip
\noindent\textit{Proof of \SGARefLink{sga3:xvi:corollary:1:4}{1.4} a).}
We must show that the restriction of \(u\) to \(C\) is an immersion
(\SGARefLink{sga3:xvi:theorem:1:3}{1.3}~a)), which reduces us to the case where \(G=C\) has connected,
nilpotent fibers. Standard reductions (cf.\ the proof of \SGARefLink{sga3:viii:lemma:7:1}{Exposé VIII,
7.1}) reduce us to the case where \(H\) is flat over \(S\), \(u\) is
an isomorphism on the generic fiber, and \(u\) is an isomorphism of
the underlying spaces on the closed fiber. The group \(H\) then has
connected fibers and is therefore separated over \(S\) (Exposé
VI\textsubscript{B}, \SGARefLink{prop:sga3-VIB-5.2}{5.2}). Grant for a moment the following lemma.

\medskip
\noindent\SGARefTarget{sga3:xvi:lemma:1:6}%
\textbf{Lemma 1.6.} --- \textit{Let \(S\) be a prescheme and let
\(G\) be an \(S\)-group prescheme, smooth over \(S\), with connected,
nilpotent fibers and unipotent rank zero. Then:}
\begin{enumerate}
\item[\textit{i)}] \textit{\(G\) is commutative.}
\item[\textit{ii)}] \textit{For every integer \(n>0\),
\({}_{n}G\) is an \(S\)-group prescheme, flat and quasi-finite over
\(S\), and it is finite over \(S\) if and only if the abelian rank,
or the reductive rank, of \(G\) is locally constant on \(S\).}
\item[\textit{iii)}] \textit{For every integer \(q>0\) invertible on
\(S\), the family of subgroups \({}_{q^n}G\) is universally
schematically dense in \(G\) relative to \(S\) (EGA IV, 11.10.8).}
\end{enumerate}

\SGARefLink{sga3:xvi:lemma:1:6}{Lemma 1.6} applies to \(G\), and since \(G\) possesses maximal tori
locally, the reductive rank of the fibers of \(G\) is locally
constant. Thus, by \SGARefLink{sga3:xvi:lemma:1:6}{1.6}~ii), \({}_{n}G\) is finite over \(S\) for
every integer \(n>0\). It follows from \SGARefLink{sga3:viii:proposition:7:9}{Exposé VIII, 7.9} that \(u\) is
an immersion.

\medskip
\noindent\textit{Proof of \SGARefLink{sga3:xvi:lemma:1:6}{1.6}.}
Let us first examine the case where \(S\) is the spectrum of a field.

\medskip
\noindent\SGARefTarget{sga3:xvi:lemma:1:7}%
\textbf{Lemma 1.7.} --- \textit{Let \(k\) be a field of
characteristic \(p\), and let \(G\) be a smooth, connected, nilpotent
algebraic \(k\)-group of unipotent rank zero. Then \(G\) is a
commutative group, an extension of an abelian variety \(A\) by a
torus \(T\). For every integer \(n>0\), \({}_{n}G\) is a finite
group, étale if \((n,p)=1\), defined by a \(k\)-algebra of rank
\(n^{\rho_r+2\rho_{ab}}\). If \((q,p)=1\), the family of subgroups
\({}_{q^n}G\), \(n\in\mathbb N\), is schematically dense in \(G\).}

\smallskip
\noindent\textit{Proof of \SGARefLink{sga3:xvi:lemma:1:7}{1.7}.}
Let \(Z\) be the center of \(G\). The group \(G/Z\) is affine
(\SGARefLink{sga3:xii:6.1}{Exposé XII, 6.1}), smooth, connected, and of unipotent rank zero;
hence it is a torus (BIBLE 4, th.~4). It follows that \(G\) is
commutative (Exposé XII, 6.4). If \(T\) is the unique maximal torus
of \(G\) (Exposé XV, 3.4), Chevalley's theorem (Séminaire Bourbaki
1956/57, no.~145) immediately shows that \(G\) is an extension of an
abelian variety \(A\) by \(T\). For every integer \(n>0\), raising to
the \(n\)-th power is an epimorphism on \(T\), whence the exact
sequence
\[
  0\longrightarrow {}_{n}T\longrightarrow {}_{n}G
  \longrightarrow {}_{n}A\longrightarrow 0.
\]
A classical theorem of Weil (A.~Weil, \emph{Variétés abéliennes et
courbes algébriques}, \S IX, th.~33, cor.~1) says that \({}_{n}A\)
is a finite group defined by a \(k\)-algebra of rank
\(n^{2\rho_{ab}}\). Since \({}_{n}T\) is a finite group of rank
\(n^{\rho_r}\), the asserted structure of \({}_{n}G\) follows.

Now let \(H\) be the smallest closed subscheme of \(G\) containing
\({}_{q^n}G\) for every \(n\). It follows from the preceding
discussion and \SGARefLink{sga3:xv:lemma:4:6}{Exposé XV, 4.6} that \(H\) is a smooth, connected
subgroup, hence of the same type as \(G\). Raising to the \(q\)-th
power on \(H\) is an epimorphism (since \({}_{q}H\) is finite), so
there is an exact sequence
\[
  0\longrightarrow {}_{q}H\longrightarrow {}_{q}G
  \longrightarrow {}_{q}(G/H)\longrightarrow 0.
\]
It follows that \({}_{q}(G/H)=0\), and therefore \(G/H=0\). Thus the
subgroups \({}_{q^n}G\) are schematically dense in \(G\).

\smallskip
\noindent\textit{Continuation of the proof of \SGARefLink{sga3:xvi:lemma:1:6}{1.6}~i).}
To show that \(G\) is commutative, the usual procedure reduces us
first to the case where \(S\) is noetherian, and then to the case
where \(S\) is the spectrum of a local ring with closed point \(s\).
The center of \(G\) is represented by a closed group subscheme \(Z\)
of \(G\) (\SGARefLink{sga3:xi:corollary:6:11}{Exposé XI, 6.11}). To prove that \(Z=G\), it suffices to
prove this after reduction modulo every power of the maximal ideal of
the ring of \(S\): \(Z\) will then be an open subgroup of \(G\), and
hence equal to \(G\), since \(G\) has connected fibers. This reduces
us to the case where \(S\) is local Artinian.

Let \(q\) be an integer invertible on \(S\). For every integer \(n\),
\({}_{q^n}G_s\) is an étale, central subgroup of multiplicative type
of \(G_s\) (\SGARefLink{sga3:xvi:lemma:1:7}{1.7}). Since \(G\) is smooth,
\({}_{q^n}(G_s)\) lifts to an étale, central \(S\)-group subscheme
\(M_n\) of \(G\) (\SGARefLink{sga3:xv:proposition:1:2}{Exposé XV, 1.2}). The family of flat subgroups
\(M_n\), \(n\in\mathbb N\), is then schematically dense in \(G\)
(1.7 and EGA IV, 11.10.9). Since \(Z\) is closed in \(G\) and
contains every \(M_n\), we have \(Z=G\).

\smallskip
\noindent\textit{Proof of \SGARefLink{sga3:xvi:lemma:1:6}{1.6}~ii).}
To see that \({}_{n}G\) is flat and quasi-finite over \(S\), it
suffices to show that the \(n\)-th power map on \(G\) is flat and
quasi-finite. Since \(G\) is flat over \(S\), we are reduced to the
case where \(S\) is the spectrum of a field (EGA IV, 11.3.10), and
we observed in the proof of 1.7 that the \(n\)-th power map is an
epimorphism. Moreover, \({}_{n}G\) is separated over \(S\), since
\(G\) is separated over \(S\) (Exposé VI\textsubscript{B}, \SGARefLink{prop:sga3-VIB-5.2}{5.2}).
For \({}_{n}G\) to be finite over \(S\), it is necessary and
sufficient that the fibers of \({}_{n}G\) be spectra of finite
algebras of locally constant rank. This is seen immediately by
reducing to the case where \(S\) is the spectrum of a henselian local
ring. By \SGARefLink{sga3:xvi:lemma:1:7}{1.7}, this condition is equivalent to saying that the abelian
rank, or the reductive rank, of the fibers of \(G\) is locally
constant.

\smallskip
\noindent\textit{Proof of \SGARefLink{sga3:xvi:lemma:1:6}{1.6}~iii).}
To see that the family of subgroups
\({}_{q^n}G\), \(n\in\mathbb N\), is universally schematically dense
in \(G\), we again reduce to the noetherian case. In view of 1.7 and
1.6~ii), it suffices to apply EGA IV, 11.10.9.

\medskip
\noindent\textit{Proof of \SGARefLink{sga3:xvi:corollary:1:4}{1.4} b).}
We have reduced to the case where \(S\) is the spectrum of a discrete
valuation ring and where \(G^0\) possesses a Cartan subgroup \(C\).
Let
\[
  N=\Norm_G(C)=\Norm_G(T),
\]
where \(T\) is the unique maximal torus of \(C\) (\SGARefLink{sga3:xii:7.1}{Exposé XII,
7.1}~a) and~b)). Since \(H\) is separated over \(S\), it follows from
\SGARefLink{sga3:xvi:lemma:1:6}{1.6}~ii) and \SGARefLink{sga3:viii:proposition:7:12}{Exposé VIII, 7.12} that the restriction of \(u\) to \(C\)
is a closed immersion. On the other hand, to prove that \(u\) is a
closed immersion, it suffices to prove this for \(u|_N\) (\SGARefLink{sga3:xvi:theorem:1:3}{1.3}~b)).
Since, by hypothesis, \(W=N/C\) is represented by a finite group
scheme, the assertion follows from the next lemma, applied to the
exact sequence
\[
  0\longrightarrow C\longrightarrow N\longrightarrow W
  \longrightarrow 1.
\]
(Recall that a proper immersion is a closed immersion.)

\medskip
\noindent\SGARefTarget{sga3:xvi:lemma:1:8}%
\textbf{Lemma 1.8.} --- \textit{Let \(S\) be a prescheme and let
\(G\) be an \(S\)-group prescheme of finite presentation over \(S\),
an extension of an \(S\)-group prescheme \(G''\), proper and of finite
presentation over \(S\), by an \(S\)-group prescheme \(G'\), flat and
of finite presentation over \(S\):}
\[
  1\longrightarrow G'\longrightarrow G\longrightarrow G''
  \longrightarrow 1.
\]
\textit{Let \(H\) be an \(S\)-group prescheme of finite presentation
over \(S\) (or locally of finite presentation if \(S\) is locally
noetherian), and let \(u:G\to H\) be an \(S\)-group morphism. If the
restriction of \(u\) to \(G'\) is proper, then \(u\) is proper.}

\smallskip
\noindent\textit{Proof of \SGARefLink{sga3:xvi:lemma:1:8}{1.8}.}
As usual, we reduce to the case where \(S\) is noetherian (Exposé
VI\textsubscript{B}, \S10, and Exposé XV, 6.2), and we may then apply
the valuative criterion of properness (EGA II, 7.3.8). We therefore
suppose that \(S\) is the spectrum of a discrete valuation ring
\(A\), with closed point \(s\) and generic point \(t\). Let
\(\bar x\in G(t)\) and \(h\in H(S)\) satisfy
\[
  u_t(\bar x)=h(t).
\]
We must show that \(\bar x\) comes from a unique element
\(x\in G(S)\). It is enough to prove the existence and uniqueness of
\(x\) after a faithfully flat extension of the discrete valuation
ring \(A\). Let \(\bar y\) be the image of \(\bar x\) in \(G''(t)\).
Since \(G''\) is proper over \(S\), \(\bar y\) comes from a unique
element \(y\in G''(S)\). Let \(X\) be the inverse image of \(y\) in
\(G\). The prescheme \(X\) is faithfully flat over \(S\), since
\(G'\) is faithfully flat over \(S\), as is the morphism
\(G\to G''\) (Exposé VI\textsubscript{B}, \S9). After a faithfully
flat change of discrete valuation ring, we may suppose that \(X\)
has a section \(\ell\) over \(S\) (EGA IV, 14.5.8). Replacing
\(\bar x\) by \(\bar x\ell_t^{-1}\), we may suppose that
\(\bar x\in G'(t)\). The existence and uniqueness of \(x\) now
follow from the properness of \(u|_{G'}\).

\medskip
\noindent\textit{Proof of \SGARefLink{sga3:xvi:corollary:1:5}{1.5}.}

\smallskip
\noindent\textit{a)}
We shall see in Exposé XIX that, if \(G\) is reductive, then \(G\)
possesses maximal tori locally for the faithfully flat topology, has
a finite Weyl group, and has unipotent rank zero. Assertion~a)
therefore follows from \SGARefLink{sga3:xvi:corollary:1:4}{1.4}.

\smallskip
\noindent\textit{b)}
If \(G\) has affine fibers and unipotent rank zero, then \(G\)
possesses maximal tori locally for the fpqc topology (\SGARefLink{sga3:xv:corollary:8-18}{Exposé XV,
8.18}). It remains only to apply \SGARefLink{sga3:xvi:corollary:1:4}{1.4}, observing that, since \(G\) has
connected, smooth, affine, solvable fibers, its Weyl group is the
identity group (BIBLE 6, th.~1, cor.~3).

\medskip
\noindent\textit{Proof of \SGARefLink{sga3:xvi:theorem:1:3}{1.3} a).}
Since \(u\) is already a monomorphism, to prove that \(u\) is an
immersion it suffices to prove that \(u\) is proper at the points of
\(u(G)\) (EGA IV, 15.7.1); for this we may use the valuative
criterion of local properness (EGA IV, 15.7.5). We are therefore
reduced to the case where \(S\) is the spectrum of a complete
discrete valuation ring with algebraically closed residue field,
closed point \(s\), and generic point \(t\). Since \(G\) is flat over
\(S\), standard reductions (cf.\ the proof of \SGARefLink{sga3:viii:lemma:7:1}{Exposé VIII, 7.1})
reduce us to the case where \(H\) is flat over \(S\), \(u_t\) is an
isomorphism, and \(u_s\) is an isomorphism of underlying spaces. To
say that \(u\) is an immersion is then equivalent to saying that
\(u\) is an isomorphism.

By hypothesis, \(G\) possesses Cartan subgroups locally for the fpqc
topology, so we may speak of the open subset \(U\) of regular points
of \(G\) (\SGARefLink{sga3:xv:theorem:7-3}{Exposé XV, 7.3}~i) and~ii)).

\medskip
\noindent\SGARefTarget{sga3:xvi:lemma:1:9}%
\textbf{Lemma 1.9.} --- \textit{Under the preceding hypotheses, for
\(u\) to be an immersion it suffices that \(u|_U\) be an immersion.}

Indeed, to say that \(u|_U\) is an immersion means that there exist
an open subscheme \(V\) of \(H\) and a closed subscheme \(F\) of
\(V\) such that \(u|_U\) factors through \(F\) and induces an
isomorphism from \(U\) onto \(F\). Since \(H_t\), and therefore
\(V_t\), is irreducible, one necessarily has \(V_t=F_t\). Then \(F\)
contains the schematic closure of \(V_t\) in \(V\), which equals
\(V\) because \(H\) is flat over \(S\). Thus \(V=F\). It follows
that \(u|_U\) is an open immersion, and hence that \(u\) is an open
immersion (Exposé VI\textsubscript{B}, \SGARefLink{cor:sga3-VIB-2.6}{2.6}).

It remains to prove that \(u|_U\) is an immersion. For this we apply
the valuative criterion of local properness. After replacing \(S\)
by the spectrum of a discrete valuation ring faithfully flat over
\(S\), we must show the following: if \(h\) is a section of \(H\)
over \(S\) whose image \(h(S)\) is contained in \(u(U)\), then \(h\)
is the image of a section \(g\) of \(U\) over \(S\). It suffices to
show that \(h\) is contained in the image of a Cartan subgroup \(C\)
of \(G\). Indeed, by hypothesis \(u|_C\) is an immersion, so \(h\)
is the image of a section \(g\) of \(C\), necessarily a section of
\(U\).

Let \(a\in U(s)\) and \(b\in U(t)\) satisfy, respectively,
\[
  u_s(a)=h(s),
  \qquad
  u_t(b)=h(t).
\]
Since \(u_t\) is an isomorphism, \(h(t)\) is a regular point of
\(H_t\), and is therefore contained in a unique Cartan subgroup
\(D_t\) of \(H_t\) (\SGARefLink{rem:sga3-VIB-3.2}{Exposé XIII, 3.2}). Let \(D\) be the schematic
closure of \(D_t\) in \(H\).

\medskip
\noindent\SGARefTarget{sga3:xvi:lemma:1:10}%
\textbf{Lemma 1.10.} ---
\begin{enumerate}
\item[\textit{i)}] \textit{\(D_s\) is a nilpotent algebraic group
(Exposé VI\textsubscript{B}, \S8).}
\item[\textit{ii)}] \textit{\(\dim D_s=\nu\), where \(\nu\) is the
nilpotent rank of \(G\), and also the nilpotent rank of
\((H_s)_{\mathrm{red}}\).}
\item[\textit{iii)}] \textit{\(h(s)\in D_s(\kappa(s))\).}
\end{enumerate}

\noindent\textit{i)}
Since \(H\) is separated (Exposé VI\textsubscript{B}, \SGARefLink{prop:sga3-VIB-5.2}{5.2}), so is
\(D\). Moreover, \(D\) is flat over \(S\) and its generic fiber is
nilpotent. It follows from Exposé VI\textsubscript{B}, \SGARefLink{prop:sga3-VIB-8.4}{8.4} that
\(D_s\) is nilpotent.

\smallskip
\noindent\textit{ii)}
Since \((H_s)_{\mathrm{red}}\simeq G_s\), these two groups have the
same nilpotent rank. Since \(G\) possesses Cartan subgroups locally
for the fpqc topology, \(G_s\) and \(G_t\) have the same nilpotent
rank \(\nu\). Finally, since \(D\) is flat over \(S\),
\[
  \dim D_s=\dim D_t=\nu
\]
(Exposé VI\textsubscript{B}, \S4).

\smallskip
\noindent\textit{iii)}
Since \(h_t\) factors through \(D_t\), the section \(h\) plainly
factors through \(D\), which proves iii).

The next lemma now shows that \((D_s)_{\mathrm{red}}\) is a Cartan
subgroup of
\[
  (H_s)_{\mathrm{red}}\simeq G_s.
\]

\medskip
\noindent\SGARefTarget{sga3:xvi:lemma:1:11}%
\textbf{Lemma 1.11.} --- \textit{Let \(G\) be a smooth, connected
algebraic group over an algebraically closed field \(k\), and let
\(D\) be a smooth, nilpotent algebraic subgroup containing a regular
element \(a\in G(k)\). Then \(D^0\) is contained in a Cartan subgroup
of \(G\). If, in addition, \(\dim D\) equals the nilpotent rank of
\(G\), then \(D\) is a Cartan subgroup of \(G\), and in particular is
connected.}

Let \(Z\) be the center of \(G\), let \(G'=G/Z\), which is affine
(\SGARefLink{sga3:xii:6.1}{Exposé XII, 6.1}), and let \(a'\) and \(D'\) be the images of \(a\)
and \(D\) in \(G'\). The correspondence between Cartan subgroups of
\(G\) and Cartan subgroups of \(G'\) (\SGARefLink{sga3:xii:6.6}{Exposé XII, 6.6}~e)) shows at
once that \(a'\) is a regular element of \(G'\) and that it suffices
to prove the lemma for the pair \((D',G')\). Thus we may suppose that
\(G\) is affine.

Let \(s\) be the semisimple component of \(a\). It belongs to
\(D(k)\) (BIBLE 4, th.~3) and is regular (BIBLE 7, th.~2, cor.~1).
By BIBLE 6, th.~2, \(s\) centralizes \(D^0\); hence \(D^0\) is
contained in the connected centralizer of \(s\), which is a Cartan
subgroup of \(G\) (BIBLE 7, th.~2). If now
\(\dim D\) equals the nilpotent rank of \(G\), then \(D^0\) is a
Cartan subgroup of \(G\), equal to \(\Centr_G(T)\), where \(T\) is
the unique maximal torus of \(D^0\) (\SGARefLink{sga3:xii:6.6}{Exposé XII, 6.6}). But \(D\) is
nilpotent, and therefore centralizes \(T\) (BIBLE 6, th.~2), so
\(D=D^0\).

Return now to the group \(G\). The group
\[
  C_t=u_t^{-1}(D_t)
\]
is the unique Cartan subgroup of \(G_t\) containing \(b\). Let \(C\)
be the schematic closure of \(C_t\) in \(G\).

By functoriality of schematic closure, \(u|_C\) factors through
\(D\), and therefore \(u_s(C_s)\subset D_s\). Since \(u_s\) is an
isomorphism from \(G_s\) onto \((H_s)_{\mathrm{red}}\), since
\[
  \dim C_s=\dim C_t=\nu=\dim D_s
\]
by \SGARefLink{sga3:xvi:lemma:1:10}{1.10}, and since \(D_s\) is connected by \SGARefLink{sga3:xvi:lemma:1:11}{1.11}, \(u_s\) induces an
isomorphism from \((C_s)_{\mathrm{red}}\) onto
\((D_s)_{\mathrm{red}}\). Thus \((C_s)_{\mathrm{red}}\) is a Cartan
subgroup of \(G_s\). The next lemma shows that \(C\) itself is a
Cartan subgroup of \(G\).

\medskip
\noindent\SGARefTarget{sga3:xvi:lemma:1:12}%
\textbf{Lemma 1.12.} --- \textit{Let \(S\) be the spectrum of a
discrete valuation ring, let \(G\) be a smooth \(S\)-group prescheme
of finite type, and let \(C\) be a group subscheme of \(G\), flat
over \(S\), such that the generic fiber \(C_t\) is a Cartan subgroup
of \(G_t\) and the geometric reduced closed fiber
\((C_{\bar s})_{\mathrm{red}}\) is a Cartan subgroup of
\(G_{\bar s}\). Then \(C\) is a Cartan subgroup of \(G\).}

We must prove that \(C\) is smooth over \(S\), and it suffices to
establish this after a faithfully flat extension of the base. The
hypotheses on \(C\) imply that the nilpotent rank of the fibers of
\(G\) is constant; hence we may speak of the open subset \(U\) of
regular points of \(G\) (\SGARefLink{sga3:xv:theorem:7-3}{Exposé XV, 7.3}). Since
\((C_{\bar s})_{\mathrm{red}}\) is a Cartan subgroup of
\(G_{\bar s}\), \(C_{\bar s}\) contains a regular point
\(a_{\bar s}\) of \(G_{\bar s}\). Since \(C\) is flat over \(S\),
after a faithfully flat extension of the base if necessary
(EGA IV, 14.5.8), we may suppose that \(a_{\bar s}\) is an element
of \(G(s)\) which lifts to a section \(a\in C(S)\). Since \(U\) is
open and contains \(a(s)\), it contains \(a\). Let \(C'\) be the
unique Cartan subgroup of \(G\) containing \(a\) (\SGARefLink{sga3:xv:theorem:7-3}{Exposé XV, 7.3}).
Necessarily \(C_t=C'_t\), and consequently \(C'\) and \(C\) coincide
with the identity component of the schematic closure of \(C_t\) in
\(G\) (observe that both \(C'\) and \(C\) are flat over \(S\)).

\smallskip
\noindent\textit{End of the proof of \SGARefLink{sga3:xvi:theorem:1:3}{1.3} a).}
By hypothesis, the restriction of \(u\) to the Cartan subgroup \(C\)
of \(G\) is an immersion. It is clear that \(h(S)\) is contained in
\(u(C)=D\), and hence \(h\) is the image of a section \(g\in C(S)\).
\hfill\(\square\)

\medskip
\noindent\textit{Proof of \SGARefLink{sga3:xvi:theorem:1:3}{1.3} b).}
We shall again use the valuative criterion of local properness
(EGA IV, 15.5) in the case of an immersion (respectively, the
valuative criterion of properness (EGA II, 7.3.6) in the case of a
closed immersion). This reduces us to the case where \(S\) is the
spectrum of a complete discrete valuation ring with algebraically
closed residue field. Let \(s\) be the closed point and \(t\) the
generic point of \(S\). Replacing \(H\) by the schematic closure in
\(H\) of \(u_t(G_t)\), we may suppose that \(H\) is flat over \(S\)
and that \(u_t\) is an isomorphism.

Let \(G^0\) be the connected component of \(G\). By \SGARefLink{sga3:xvi:remark:1:2}{1.2} a), the
restriction \(u|_{G^0}\) is an immersion; denote by \(H^0\) the group
subprescheme of \(H\) that is the image of \(G^0\). Thus
\[
  H_s^0=(H_s)_{\mathrm{red}}^0.
\]

To verify the valuative criterion, we must show that every section
\(h\) of \(H\) over \(S\) such that \(h(S)\) is contained in
\(u(G)\), in the case of an immersion (respectively, every section
\(h\) of \(H\) over \(S\), in the case of a closed immersion), is
the image of a section \(g\) of \(G\) over \(S\).

Let \(C\) be a Cartan subgroup of \(G\), and let
\(D=u(C)\) be the Cartan subgroup of \(H^0\) that is the image of
\(C\). The \(S\)-group
\[
  D'=\operatorname{int}(h)D
\]
has connected fibers, so its underlying space is contained in that
of \(H^0\); on the other hand, since it is smooth over \(S\), it is
reduced, and hence \(D'\) is a smooth group subprescheme of \(H^0\).
The fibers of \(D'\) are Cartan subgroups of the fibers of \(H^0\);
therefore \(D'\) is a Cartan subgroup of \(H^0\), and
\[
  C'=u^{-1}(D')
\]
is a Cartan subgroup of \(G\). But
\(\operatorname{Transpstr}_G(C,C')\) is smooth and surjective over
\(S\) (\SGARefLink{sga3:xii:7.1}{Exposé XII, 7.1} b)). Since \(S\) is henselian with
algebraically closed residue field, there exists \(g\in G(S)\) such
that \(\operatorname{int}(g)C=C'\). Replacing \(h\) by
\(u(g)^{-1}h\), we may suppose that \(h(t)\) normalizes \(D_t\)
(and that \(h(s)\) is the image of an element of the normalizer of
\(C_s\) in \(G_s\), in the case of an immersion). Let
\[
  N=\Norm_G(C).
\]
By hypothesis, \(u|_N\) is an immersion (respectively, a closed
immersion); hence \(h\) is indeed the image of a section of \(N\)
over \(S\), which completes the proof.
\hfill\(\square\)

% END INCLUDED SGA3_English_current_buildable_source_20260823/source/inputs/XVI/tex/components/01_expose_XVI_section1_immersion_criterion.tex

% BEGIN INCLUDED SGA3_English_current_buildable_source_20260823/source/inputs/XVI/tex/components/02_expose_XVI_section2_quotient_representability.tex
% SGA 3, Exposé XVI, Loop-1 English reconstruction.
% Current scope: Section 2 from its opening through the definition of the
% functor F in the proof of Corollary 2.4.
% Authority: Polo--Gille Expo16.pdf, local pages 10--12
% (combined-reader pages 948--950; printed pages 503--507).
% The two diagrams are source-PDF crops used as Loop-1 placeholders; native
% mathematical reconstruction remains a Loop-2 obligation.

\SGARefTarget{sga3:xvi:section:2}%
\section*{2. A Representability Theorem for Quotients}
\addcontentsline{toc}{section}{2. A Representability Theorem for Quotients}

Let us ``recall'' the following result.

\medskip
\noindent\SGARefTarget{sga3:xvi:theorem:2:1}%
\textbf{Theorem 2.1.} --- \textit{Let \(S\) be a prescheme, let
\(X\) and \(Y\) be two \(S\)-preschemes, and let
\(f:X\to Y\) be an \(S\)-morphism. Suppose that one is in one of the
following two cases:}

\begin{enumerate}
\renewcommand{\labelenumi}{\textit{\alph{enumi})}}
\item \textit{The morphism \(f\) is locally of finite presentation.}

\item \textit{The prescheme \(S\) is locally noetherian and \(X\) is
locally of finite type over \(S\).}
\end{enumerate}

\noindent\textit{Then the following conditions are equivalent:}

\begin{enumerate}
\renewcommand{\labelenumi}{\textit{\roman{enumi})}}
\item \textit{There exist an \(S\)-prescheme \(X'\) and a
factorization of \(f\)}
\[
  f:X\xrightarrow{\,f'\,}X'\xrightarrow{\,f''\,}Y,
\]
\textit{where \(f'\) is a faithfully flat \(S\)-morphism, locally
of finite presentation, and \(f''\) is a monomorphism.}

\item \textit{The first projection}
\[
  p_1:X\times_Y X\longrightarrow X
\]
\textit{is a flat morphism.}
\end{enumerate}

\noindent\textit{Moreover, if the preceding conditions hold, then
\((X',f')\) is a quotient of \(X\), for the fpqc topology, by the
equivalence relation defined by \(f\). Consequently the
factorization \(f=f''\circ f'\) in i) is unique up to isomorphism.}

The proof of this delicate theorem will be given in EGA V. One may
also consult the lecture by J.-P. Murre, \textit{Séminaire
Bourbaki}, May 1965, no.~294, theorem~2, corollary~2, which treats
the case where \(Y\) is locally noetherian and \(X\) is of finite
type over \(Y\). We shall see that one can reduce to that case.

Let us first make a few remarks.

\medskip
\noindent\textit{a)} The implication i) \(\Rightarrow\) ii) is
trivial. Indeed, the first projection
\[
  p'_1:X\times_{X'}X\longrightarrow X
\]
factors through \(X\times_Y X\):
\[
  p'_1:X\times_{X'}X\xrightarrow{\,u\,}
  X\times_Y X\xrightarrow{\,p_1\,}X.
\]
The morphism \(u\) is an isomorphism because \(f''\) is a
monomorphism, and \(p'_1\) is flat because \(f'\) is flat; therefore
\(p_1\) is flat.

\medskip
\noindent\textit{b)} The assertions of \SGARefLink{sga3:xvi:theorem:2:1}{2.1} are local on \(Y\)
(hence local on \(S\)); they are also local on \(X\), as follows
readily from the fact that a flat morphism locally of finite
presentation is open (EGA IV, 11.3.1).

\medskip
\noindent\textit{c)} Under the hypotheses of \SGARefLink{sga3:xvi:theorem:2:1}{2.1} a), the foregoing
reduces us to the case where \(X\) and \(Y\) are affine and \(f\) is
of finite presentation. Replacing \(S\) by \(Y\), we may suppose
that \(X\) and \(Y\) are of finite presentation over \(S\). We then
reduce to the noetherian case for \(S\) by EGA IV, 11.2.6.

\medskip
\noindent\textit{d)} Under the hypotheses of \SGARefLink{sga3:xvi:theorem:2:1}{2.1} b), we may suppose
that \(S\), \(X\), and \(Y\) are affine, that \(S\) is noetherian,
and that \(X\) is of finite type over \(S\). Regard \(Y\) as a
filtered projective limit of affine schemes \(Y_i\) of finite type
over \(S\). The schemes \(X\times_{Y_i}X\) form a decreasing
filtered family of closed subschemes of \(X\times_S X\), whose
projective limit is \(X\times_Y X\). Since \(X\times_S X\) is
noetherian, one has
\[
  X\times_{Y_i}X=X\times_Y X
\]
for all sufficiently large \(i\). Thus
\[
  f_i:X\xrightarrow{\,f\,}Y\longrightarrow Y_i
\]
satisfies the hypotheses of \SGARefLink{sga3:xvi:theorem:2:1}{2.1} ii) if \(f\) does. Since the
equivalence relation on \(X\) defined by \(f\) coincides with the
one defined by \(f_i\), it is clear that it suffices to prove
ii) \(\Rightarrow\) i) for \(f_i\). This reduces us to the case
where \(Y\) is of finite type over \(S\).

\subsection*{Application to Group Preschemes}
\addcontentsline{toc}{subsection}{Application to Group Preschemes}

\noindent\SGARefTarget{sga3:xvi:theorem:2:2}%
\textbf{Theorem 2.2.} --- \textit{Let \(S\) be a prescheme and let
\(G\) be an \(S\)-group prescheme locally of finite presentation
over \(S\), acting on an \(S\)-prescheme \(X\). Let
\(\xi\in X(S)\) be a section of \(X\) over \(S\), such that the
stabilizer \(H\) of \(\xi\) in \(G\) is an \(S\)-group
subprescheme of \(G\), flat over \(S\). If \(X\) is locally of
finite type over \(S\), or if \(S\) is locally noetherian, then the
homogeneous quotient space \(G/H\) is representable by an
\(S\)-prescheme locally of finite presentation over \(S\), and the
\(S\)-morphism}
\[
  f:G\longrightarrow X,\qquad g\longmapsto g\cdot\xi
\]
\textit{factors as}
\begin{center}
  % BEGIN INCLUDED SGA3_English_current_buildable_source_20260823/source/inputs/XVI/native_diagrams/exp16/001_theorem_2_2_factorization_triangle.tex
\begin{tikzcd}[column sep=5.6em,row sep=2.6em]
  G \arrow[d,"p"'] \arrow[dr,"f"]
    & {} \\
  G/H \arrow[r,"i"]
    & X.
\end{tikzcd}

% END INCLUDED SGA3_English_current_buildable_source_20260823/source/inputs/XVI/native_diagrams/exp16/001_theorem_2_2_factorization_triangle.tex

\end{center}
\textit{where \(p\) is the canonical projection, a faithfully flat
morphism locally of finite presentation, and \(i\) is a
monomorphism. (For definiteness, we have supposed that \(G\) acts
on the left on \(X\).)}

\medskip
\noindent\textit{Proof.}
The morphism \(f\) makes \(G\) into an \(X\)-prescheme. By the
definition of the stabilizer of \(\xi\), the morphism
\[
  G\times_S H\longrightarrow G\times_X G,\qquad
  (g,h)\longmapsto(g,gh)
\]
is an isomorphism. Since \(H\) is flat over \(S\),
\(G\times_S H\) is flat over \(G\); hence the first projection
\[
  p_1:G\times_X G\longrightarrow G
\]
is flat. Moreover, if \(X\) is locally of finite type over \(S\),
then \(f\) is locally of finite presentation (EGA IV, 1.4.3 v));
otherwise \(S\) is locally noetherian by hypothesis. It therefore
suffices to apply \SGARefLink{sga3:xvi:theorem:2:1}{2.1} to \(f\).

It remains only to show that \(G/H\) is locally of finite
presentation over \(S\), and this follows immediately from
\SGARefLink{sga3:v:proposition:9:1}{Exposé V, 9.1}.
\hfill\(\square\)

\medskip
\noindent\SGARefTarget{sga3:xvi:corollary:2:3}%
\textbf{Corollary 2.3.} --- \textit{Let \(S\) be a prescheme, let
\(G\) and \(H\) be two \(S\)-group preschemes, and let
\(u:G\to H\) be an \(S\)-homomorphism. Suppose that \(G\) is locally
of finite presentation over \(S\), and that either \(H\) is locally
of finite type over \(S\), or \(S\) is locally noetherian. If
\(K=\Ker(u)\) is flat over \(S\), then the quotient group \(G/K\)
is representable by an \(S\)-group prescheme locally of finite
presentation over \(S\), and \(u\) factors as}
\begin{center}
  % BEGIN INCLUDED SGA3_English_current_buildable_source_20260823/source/inputs/XVI/native_diagrams/exp16/002_corollary_2_3_factorization_triangle.tex
\begin{tikzcd}[column sep=5.8em,row sep=4.0em]
  G \arrow[rr,"u"] \arrow[dr,"p"']
    & {}
    & H \\
  {}
    & G/K \arrow[ur,"i"']
    & {}
\end{tikzcd}

% END INCLUDED SGA3_English_current_buildable_source_20260823/source/inputs/XVI/native_diagrams/exp16/002_corollary_2_3_factorization_triangle.tex

\end{center}
\textit{where \(p\) is the canonical projection and \(i\) is a
monomorphism.}

\medskip
\noindent\textit{Proof.}
Apply \SGARefLink{sga3:xvi:theorem:2:2}{2.2} with \(X=H\) and with \(\xi\) the unit section of \(H\).
\hfill\(\square\)

\medskip
\noindent\SGARefTarget{sga3:xvi:corollary:2:4}%
\textbf{Corollary 2.4.} --- \textit{Let \(S\) be a prescheme, let
\(G\) be an \(S\)-group prescheme of finite presentation over
\(S\), let \(H\) be an \(S\)-group prescheme smooth over \(S\) with
connected fibers (and hence of finite presentation over \(S\), by
Exposé VI\textsubscript{B}, \SGARefLink{cor:sga3-VIB-5.5}{5.5}), and let \(i:H\to G\) be a
monomorphism of \(S\)-groups, so that \(H\) is a subgroup of \(G\).
Suppose that}
\[
  N=\Norm_G(H),
\]
\textit{which is representable by a closed group subprescheme of
\(G\), of finite presentation over \(S\) (\SGARefLink{sga3:xi:corollary:6:11}{Exposé XI, 6.11}), is flat
over \(S\). Then \(G/N\) is representable by an \(S\)-prescheme of
finite presentation over \(S\) and quasi-projective over \(S\).}

All the assertions to be proved, except the last one, are local on
\(S\). To establish them, we may therefore suppose that \(S\) is
quasi-compact and that the relative dimension of \(H\) over \(S\)
is constant, equal to \(r\). Proceed as in Exposé XV, \S5. For every
integer \(n>0\), let \(G^{(n)}\) (respectively \(H^{(n)}\)) be the
\(n\)-th normal invariant of the unit section of \(G\)
(respectively of \(H\)) (EGA IV, 16). The sheaf of
\(\mathcal O_S\)-modules \(H^{(n)}\) is a quotient of \(G^{(n)}\).
Since \(H\) is smooth over \(S\) of dimension \(r\), \(H^{(n)}\)
canonically defines an element \(\xi_n\) of
\[
  X_n=\operatorname{Grass}_{\varphi(n,r)}G^{(n)}
\]
(EGA I, second edition, \S9), for a suitable integer
\(\varphi(n,r)\). On the other hand, \(G\) acts naturally on
\(G^{(n)}\), and hence also on \(X_n\), through the representation
\[
  G\longrightarrow\operatorname{Aut}_{S\text{-gr}}(G),
  \qquad g\longmapsto\operatorname{int}(g).
\]
Since \(S\) is quasi-compact, there exists an integer \(m\) such
that, for \(n\geq m\),
\[
  N=\Norm_G(H)=\Norm_G H^{(n)}
  \qquad\text{(Exposé XI, 6.11 b)).}
\]
In other words, \(N\) is the stabilizer of \(\xi_n\) for sufficiently
large \(n\). The representability of \(G/N\) therefore follows from
\SGARefLink{sga3:xvi:theorem:2:2}{2.2}. The fact that \(G/N\) is of finite presentation over \(S\) is
a consequence of \SGARefLink{sga3:v:proposition:9:1}{Exposé V, 9.1}. It remains to prove that \(G/N\) is
quasi-projective over \(S\). For this, we shall exhibit a canonical
invertible sheaf on \(G/N\) that is ample relative to \(S\).
Consider the functor
\[
  \mathscr F:(\mathbf{Sch}/S)^\circ\longrightarrow\mathbf{Ens}
\]
such that, for every \(S\)-prescheme \(T\),
\[
  \mathscr F(T)=
  \left\{
  \begin{array}{c}
  \text{\(T\)-subgroups \(H\) of \(G_T\), representable and smooth
  over \(T\),}\\
  \text{with connected fibers}
  \end{array}
  \right\}.
\]

Such a group \(H\) is of finite presentation over \(S\)
(Exposé VI\textsubscript{B}, 5.5), and
\(H\to G_T\) is a quasi-affine morphism (EGA IV, 8.11.2).
By effective descent for quasi-affine morphisms (SGA 1, VIII, 7.9),
it follows that \(\mathscr F\) is a sheaf for the fpqc topology.
Since \(G/N\) is the sheaf associated with a subfunctor of
\(\mathscr F\), there is a canonical monomorphism
\[
  G/N\longrightarrow\mathscr F.
\]
Consequently there exists a subgroup \(H'\) of
\(G\times_S(G/N)\), representable and smooth over \(G/N\), with
connected fibers, that is ``universal'' for the functor \(G/N\).
I claim that the invertible sheaf
\[
  \mathscr L=\bigl(\det(\operatorname{Lie}H')\bigr)^{-1}
\]
is ample relative to \(S\). In this form the assertion is local on
\(S\), and its proof is analogous to the one given in \SGARefLink{sga3:xv:remark:5:8}{Exposé XV,
5.8}.

The following corollary was announced in \SGARefLink{sga3:xiv:4.8.1}{Exposé XIV, 4.8 bis}.

\medskip
\noindent\SGARefTarget{sga3:xvi:corollary:2:5}%
\textbf{Corollary 2.5.} --- \textit{Let \(S\) be a prescheme, let
\(G\) be an \(S\)-group prescheme smooth and of finite presentation
over \(S\), with connected fibers, and let \(P\) be a parabolic
subgroup of \(G\) (\SGARefLink{sga3:xiv:paragraph:4-8-bis}{Exposé XIV, 4.8 bis}). Then \(G/P\) is
representable by an \(S\)-prescheme smooth and projective over
\(S\).}

It remains only to show (\textit{loc.\ cit.}) that \(G/P\) is
representable by an \(S\)-prescheme that is quasi-projective and of
finite presentation. This follows from \SGARefLink{sga3:xvi:corollary:2:4}{2.4}, since
\[
  P=\Norm_G(P)
\]
(\textit{loc.\ cit.}) and is therefore flat over \(S\).
\hfill\(\square\)

% END INCLUDED SGA3_English_current_buildable_source_20260823/source/inputs/XVI/tex/components/02_expose_XVI_section2_quotient_representability.tex

% BEGIN INCLUDED SGA3_English_current_buildable_source_20260823/source/inputs/XVI/tex/components/03_expose_XVI_section3_flat_center.tex
% SGA 3, Exposé XVI, Loop-1 English reconstruction.
% Current scope: all of Section 3, including the terminal counterexample.
% Authority: Polo--Gille Expo16.pdf, local pages 13--17
% (combined-reader pages 951--955; printed pages 509--518).
% The Proposition 3.1 diagram is a source-PDF crop used as a Loop-1
% placeholder; native mathematical reconstruction remains a Loop-2 task.

\SGARefTarget{sga3:xvi:section:3}%
\section*{3. Groups with Flat Center}
\addcontentsline{toc}{section}{3. Groups with Flat Center}

\noindent\SGARefTarget{sga3:xvi:proposition:3:1}%
\textbf{Proposition 3.1.} --- \textit{Let \(S\) be a prescheme and
let \(G\) be an \(S\)-group prescheme smooth and of finite
presentation over \(S\), with connected fibers. Suppose that the
center \(Z\) of \(G\), which is representable by \SGARefLink{sga3:xi:corollary:6:11}{Exposé XI, 6.11},
is flat over \(S\). For every integer \(n>0\), let \(G^{(n)}\)
denote the locally free \(\mathcal O_S\)-module equal to the
\(n\)-th normal invariant of \(G\) along the unit section, and let
\(\rho_n\) be the natural ``adjoint representation'' of \(G\) in
\(\operatorname{GL}_S(G^{(n)})\). Then:}

\begin{enumerate}
\renewcommand{\labelenumi}{\textit{\alph{enumi})}}
\item \textit{The quotient \(G'=G/Z\) is representable by an
\(S\)-group prescheme that is smooth, of finite presentation over
\(S\), quasi-affine, and has affine connected fibers.}

\item \textit{The representation \(\rho_n\) factors through \(G'\):}
\begin{center}
  % BEGIN INCLUDED SGA3_English_current_buildable_source_20260823/source/inputs/XVI/native_diagrams/exp16/003_proposition_3_1_factorization_triangle.tex
\begin{tikzcd}[column sep=6.0em,row sep=4.0em]
  G \arrow[rr,"\rho_n"] \arrow[dr]
    & {}
    & \operatorname{GL}_S\!\left(G^{(n)}\right) \\
  {}
    & G/Z \arrow[ur,"i_n"']
    & {}
\end{tikzcd}

% END INCLUDED SGA3_English_current_buildable_source_20260823/source/inputs/XVI/native_diagrams/exp16/003_proposition_3_1_factorization_triangle.tex

\end{center}
\textit{If \(S\) is quasi-compact, then \(i_n\) is a monomorphism
for all sufficiently large \(n\).}

\item \textit{The functor \(\mathscr T_{G'}\) of subtori of \(G'\)
is representable by an \(S\)-prescheme smooth over \(S\), and, if
\(S\) is quasi-compact, this prescheme is a sum of a family of
preschemes affine over \(S\).}
\end{enumerate}

\medskip
\noindent\textit{Proof.}
The assertions in a) are local on \(S\), which allows us to suppose
that \(S\) is quasi-compact. By \SGARefLink{sga3:xi:corollary:6:11}{Exposé XI, 6.11} b), one has
\[
  Z=\Ker\rho_n
\]
for all sufficiently large \(n\); this kernel is therefore flat
over \(S\). The representability of \(G/Z\) and the fact that \(i_n\)
is a monomorphism now follow from \SGARefLink{sga3:xvi:corollary:2:3}{2.3}. Since \(G\) is smooth and of
finite presentation over \(S\), the same is true of \(G'\)
(Exposé VI\textsubscript{B}, \S9). The group
\(\operatorname{GL}_S(G^{(n)})\) is affine over \(S\), and a
monomorphism of finite presentation is quasi-affine
(EGA IV, 8.11.2). Hence \(G'\) is quasi-affine over \(S\) and has
affine fibers. Since \(G'\) is smooth over \(S\), the functor
\(\mathscr T_{G'}\) is formally smooth (\SGARefLink{sga3:xi:proposition:2:1:bis}{Exposé XI, 2.1 bis}).
Taking \SGARefLink{sga3:xi:remark:4:6}{Exposé XI, 4.6} and \SGARefLink{sga3:xvi:corollary:4:3}{4.3} into account, the assertions of \SGARefLink{sga3:xvi:proposition:3:1}{3.1} c)
will therefore follow from the next lemma.

\medskip
\noindent\SGARefTarget{sga3:xvi:lemma:3:2}%
\textbf{Lemma 3.2.} --- \textit{Let \(S\) be a prescheme, let \(G\)
and \(H\) be two \(S\)-group preschemes of finite presentation over
\(S\), let \(u:G\to H\) be an \(S\)-group monomorphism, and let
\(\mathscr T_G\) (respectively \(\mathscr T_H\)) be the functor of
subtori of \(G\) (respectively of \(H\)) (cf.\ Exposé XV, \S8).
Then the map \(T\mapsto u(T)\) defines a monomorphism
(\SGARefLink{sga3:xv:lemma:8:3}{Exposé XV, 8.3} c))}
\[
  \widetilde u:\mathscr T_G\longrightarrow\mathscr T_H
\]
\textit{that is representable by a closed immersion of finite
presentation.}

By the usual procedure, we reduce to the following problem. Let
\(T\) be a subtorus of \(H\), let
\[
  T'=G\times_H T
\]
be its inverse image in \(G\), and let \(u_T:T'\to T\) be the
\(S\)-group monomorphism induced by \(u\). One must show that the
\(S\)-functor
\[
  \prod_{T/S}(T'/T)
\]
is representable by a closed subprescheme of \(S\), of finite
presentation. But \(T\), being a torus, is smooth over \(S\) with
connected fibers, and it suffices to apply \SGARefLink{sga3:xi:corollary:6:10}{Exposé XI, 6.10}.
\hfill\(\square\)

\medskip
\noindent\SGARefTarget{sga3:xvi:theorem:3:3}%
\textbf{Theorem 3.3.} --- \textit{Let \(S\) be a prescheme, let
\(G\) be an \(S\)-group prescheme smooth and of finite presentation
over \(S\), with connected fibers, and let \(Z\) be the center of
\(G\). Suppose that \(Z\) is flat over \(S\) and that the unipotent
rank of \(G\) (\SGARefLink{sga3:xv:definition:6:1:ter}{Exposé XV, 6.1 ter}) equals that of \(Z\). Then:}

\begin{enumerate}
\renewcommand{\labelenumi}{\textit{\alph{enumi})}}
\item \textit{The group \(G'=G/Z\) is representable; if \(S\) is
quasi-compact, the canonical morphism}
\[
  i_n:G'\longrightarrow\operatorname{GL}_S(G^{(n)})
\]
\textit{(cf.\ \SGARefLink{sga3:xvi:proposition:3:1}{3.1} b)) is an immersion for all sufficiently large
\(n\).}

\item \textit{The group \(G'\) is quasi-affine over \(S\), with
affine fibers; its center is the unit group, and \(G'\) has
unipotent rank zero.}

\item \textit{Locally for the étale topology, \(G'\) possesses
maximal tori, and these are also Cartan subgroups of \(G'\). The
functor \(\mathscr{TM}_{G'}\) of maximal tori of \(G'\)
(Exposé XV, \S8) is representable by an \(S\)-prescheme smooth and
affine over \(S\).}

\item \textit{Locally for the étale topology, \(G\) possesses
Cartan subgroups, and the functor \(\mathscr C_G\) of Cartan
subgroups of \(G\) is representable by an \(S\)-prescheme smooth
and affine over \(S\).}
\end{enumerate}

\medskip
\noindent\textit{Proof.}
By \SGARefLink{sga3:xvi:proposition:3:1}{3.1}, the group \(G'\) is represented by an \(S\)-prescheme that
is smooth and quasi-affine over \(S\), with affine fibers. In view
of the correspondence between Cartan subgroups of the fibers of
\(G\) and those of the fibers of \(G'\) (\SGARefLink{sga3:xii:6.6}{Exposé XII, 6.6} e)), the
hypothesis on the unipotent rank of \(Z\) implies that \(G'\) has
unipotent rank zero. \SGARefLink{sga3:xv:corollary:8-18}{Exposé XV, 8.18} now shows that, locally for the
étale topology, \(G'\) possesses maximal tori. The fact that \(i_n\)
is an immersion for all sufficiently large \(n\) then follows from
\SGARefLink{sga3:xvi:proposition:3:1}{3.1} b) and \SGARefLink{sga3:xvi:corollary:1:4}{1.4} a). This proves a).

We now prove c). Since \(G'\) has unipotent rank zero, every maximal
torus of \(G'\) is plainly also a Cartan subgroup of \(G'\). The
functor of maximal tori of \(G'\) is represented by a subprescheme
that is both open and closed in the functor of subtori of \(G'\);
moreover it is of finite type over \(S\) (\SGARefLink{sga3:xv:theorem:8-15}{Exposé XV, 8.15}).
It follows from \SGARefLink{sga3:xvi:proposition:3:1}{3.1} c) that this functor is represented by an
\(S\)-prescheme smooth and affine over \(S\).

Since, locally for the étale topology, \(G'\) possesses Cartan
subgroups, the same is true of \(G\); and the functor of Cartan
subgroups of \(G\) is canonically isomorphic to that of \(G'\)
(\SGARefLink{sga3:xv:theorem:7-3}{Exposé XV, 7.3} iv)). Thus c) implies d).

It remains to prove b), and more precisely to prove that the center
\(Z'\) of \(G'\) is the unit group. Since \(Z'\) is representable
(\SGARefLink{sga3:xi:corollary:6:11}{Exposé XI, 6.11}) and of finite presentation over \(S\), it suffices
to show that the fibers of \(Z'\) are reduced to the unit group
(Exposé VI\textsubscript{B}, \SGARefLink{lem:sga3-VIB-2.9}{2.9}). We are thereby reduced to the
case where \(S\) is the spectrum of an algebraically closed field.
The center \(Z'\) is evidently contained in every Cartan subgroup
of \(G'\), hence, by c), in every maximal torus of \(G'\); it is
therefore of multiplicative type. On the other hand, we shall see
in Exposé XVII that, if \(G\) is a connected algebraic group with
center \(Z\), the center \(Z'\) of \(G/Z\) is unipotent. In the
present case, \(Z'\), being at once of multiplicative type and
unipotent, is reduced to the unit group (cf.\ Exposé XVII).
\hfill\(\square\)

\subsection*{Examples of Groups Whose Center Is Flat}
\addcontentsline{toc}{subsection}{Examples of Groups Whose Center Is Flat}

\noindent\SGARefTarget{sga3:xvi:proposition:3:4}%
\textbf{Proposition 3.4.} --- \textit{Let \(S\) be a prescheme, let
\(G\) be an \(S\)-group prescheme of finite presentation over
\(S\), smooth, with connected fibers, and let \(Z\) be the center
of \(G\). Then \(Z\) is flat over \(S\) in either of the following
two cases:}

\begin{enumerate}
\renewcommand{\labelenumi}{\textit{\alph{enumi})}}
\item \textit{The unipotent rank \(\rho_u\) of the fibers of \(G\)
(\SGARefLink{sga3:xv:definition:6:1:ter}{Exposé XV, 6.1 ter}) is zero.}

\item \textit{i) \(S\) is reduced; \quad
ii) the dimension of the fibers of \(Z\) is a locally constant
function on \(S\); \quad
iii) the unipotent rank of \(G\) equals that of \(Z\).}
\end{enumerate}

\medskip
\noindent\textit{Proof of \SGARefLink{sga3:xvi:proposition:3:4}{3.4} a).}
We shall prove at the same time the following proposition.

\medskip
\noindent\SGARefTarget{sga3:xvi:proposition:3:5}%
\textbf{Proposition 3.5.} --- \textit{Under the hypotheses of
\SGARefLink{sga3:xvi:proposition:3:4}{3.4} a), suppose moreover that \(G\) possesses a maximal torus \(T\),
and let}
\[
  C=\Centr_G(T)
\]
\textit{be the Cartan subgroup of \(G\) associated with \(T\)
(\SGARefLink{sga3:xii:7.1}{Exposé XII, 7.1}). Then:}

\begin{enumerate}
\renewcommand{\labelenumi}{\textit{(\roman{enumi})}}
\item \textit{\(Z\cap T\) is a subgroup of multiplicative type of
\(G\).}

\item \textit{\(C\) is commutative and equals \(T\cdot Z\).}

\item \textit{If the quotients \(Z/(Z\cap T)\) and \(C/T\) are
representable, then they are represented by abelian preschemes
(that is, \(S\)-group preschemes smooth over \(S\), whose fibers are
abelian varieties), and the canonical monomorphism}
\[
  Z/(Z\cap T)\longrightarrow C/T
\]
\textit{is an isomorphism.}
\end{enumerate}

Using the general properties of passage to the limit established
in Exposé VI\textsubscript{B}, \S10 and \SGARefLink{sga3:xv:proposition:6-2}{Exposé XV, 6.2}, \SGARefLink{sga3:xvi:proposition:6:3}{6.3}, and
\SGARefLink{sga3:xv:corollary:6:3:bis}{6.3 bis}, we reduce as usual to the case where \(S\) is noetherian
(observe that the assertions of \SGARefLink{sga3:xvi:proposition:3:4}{3.4} and \SGARefLink{sga3:xvi:proposition:3:5}{3.5} are local on \(S\)).
We noted in 3.1 that the hypotheses on \(G\) imply that \(Z\) is
representable. To prove that \(Z\) is flat over \(S\), EGA
0\textsubscript{III}, 10.2.6 reduces us to the case where \(S\) is
local artinian. Since \(G\) is then smooth, it possesses maximal
tori locally for the étale topology (\SGARefLink{sga3:xv:corollary:8:17}{Exposé XV, 8.17}). After a
finite flat extension of \(S\), if necessary (which is permissible
when proving that \(Z\) is flat), we may therefore suppose that
\(G\) possesses a maximal torus \(T\).

Let us show, in the same way, that to establish \SGARefLink{sga3:xvi:proposition:3:5}{3.5} we may reduce
to the case where \(S\) is artinian.

\medskip
\noindent i) Since \(Z\) is closed in \(G\) (\SGARefLink{sga3:xi:corollary:6:11}{Exposé XI, 6.11}),
\(Z\cap T\) is a closed group subscheme of \(T\). It follows from
\SGARefLink{sga3:x:corollary:4-8}{Exposé X, 4.8} b), that \(Z\cap T\) is of multiplicative type if and
only if it is flat over \(S\). As above, it suffices to establish
that \(Z\cap T\) is flat when \(S\) is artinian.

\medskip
\noindent ii) Since \(G\) has unipotent rank zero, every Cartan
subgroup \(C\) of \(G\) satisfies the hypotheses of \SGARefLink{sga3:xvi:lemma:1:6}{Lemma 1.6}, and
is therefore commutative. The equality \(C=T\cdot Z\) will follow
from iii).

\medskip
\noindent iii) If \(C/T\) is representable, then \(C/T\) is smooth
over \(S\) (Exposé VI\textsubscript{B}, \S9), and its fibers are
abelian varieties (\SGARefLink{sga3:xvi:lemma:1:7}{1.7}); thus \(C/T\) is an abelian prescheme. To
show that the canonical monomorphism
\[
  i:Z/(Z\cap T)\longrightarrow C/T
\]
is an isomorphism, it suffices to verify this when \(S\) is local
artinian. Indeed, one then deduces successively that
\(Z/(Z\cap T)\) is flat over \(S\) (EGA
0\textsubscript{III}, 10.2.6), that \(i\) is flat
(EGA IV, 11.3.10), that \(i\) is an open immersion
(EGA IV, 17.9.1), and finally that \(i\) is an isomorphism
(because \(C/T\) has connected fibers).

We henceforth suppose that \(S\) is local artinian, with closed
point \(s\), and that \(G\) possesses a maximal torus \(T\). Put
\[
  C=\Centr_G(T).
\]
The group \(C\) is a Cartan subgroup of \(G\) (\SGARefLink{sga3:xii:7.1}{Exposé XII, 7.1})
and contains \(Z\).

The algebraic group \(M_s=Z_s\cap T_s\) is a central subgroup of
multiplicative type of \(G_s\). Since \(T\) is smooth over \(S\),
\(M_s\) lifts to a group subscheme \(M\) of \(T\), of
multiplicative type (\SGARefLink{sga3:ix:theorem:3:6:bis}{Exposé IX, 3.6 bis}), and contained in the
center of \(G\) (\SGARefLink{sga3:ix:corollary:3:9}{Exposé IX, 3.9}), hence contained in \(Z\cap T\).
Since \(S\) is artinian and \(M\) and \(T\) are flat over \(S\),
the quotient groups
\[
  (Z\cap T)/M,\qquad Z'=Z/M,\qquad A=C/T
\]
are representable (Exposé VI\textsubscript{A}, \S\S4 and 5).
By construction, \(\bigl((Z\cap T)/M\bigr)_s\) is the unit group,
so \(Z\cap T=M\) (Exposé VI\textsubscript{B}, \SGARefLink{lem:sga3-VIB-2.9}{2.9}); in
particular it is flat over \(S\), which proves \SGARefLink{sga3:xvi:proposition:3:5}{3.5} i).

Passing to the quotient gives a canonical monomorphism
\[
  i:Z'\longrightarrow A,
\]
which is therefore a closed immersion (Exposé
VI\textsubscript{B}, \SGARefLink{cor:sga3-VIB-1.4.2}{1.4.2}). We have already observed that \(A\)
is an abelian prescheme. It remains to show that \(i\) is an
isomorphism. This will prove \SGARefLink{sga3:xvi:proposition:3:5}{3.5} iii) and will imply that \(Z'\)
is flat over \(S\), and hence that \(Z\) is flat over \(S\)
(as an extension of \(Z'\) by the flat group \(M\); Exposé
VI\textsubscript{B}, \S9). Since \(Z_s\) has the same abelian rank
as \(G_s\) (\SGARefLink{sga3:xii:6.1}{Exposé XII, 6.1}), \(i_s\) is an epimorphism, hence an
isomorphism. Let \(q>0\) be an integer invertible on \(S\). By the
density theorem of \SGARefLink{sga3:xvi:lemma:1:6}{1.6} iii), to see that \(i(Z')=A\), it suffices
to show that, for every integer \(n\) equal to a power of \(q\),
\(i(Z')\) contains \({}_nA\). Now let \(M_s^0\) be the identity
component of \(M_s\). It follows immediately by duality that
raising to the \(q\)-th power in \(M_s^0\) is an epimorphism.
Consequently, if \(m_0\) is the exponent of \(q\) in the prime
factorization of \(\operatorname{card}(M_s/M_s^0)\), the image of
\({}_{nm_0}Z_s\) in
\[
  A_s\simeq Z_s/M_s
\]
contains \({}_nA_s\). There is therefore a subgroup \(M_s(n)\) of
multiplicative type of \(Z_s\) whose image in \(A_s\) is
\({}_nA_s\). As above, \(M_s(n)\) lifts to a central subgroup
\(M(n)\) of multiplicative type of \(G\). The image of \(M(n)\)
in \(A\) is a subgroup of multiplicative type (\SGARefLink{sga3:ix:theorem:6:8}{Exposé IX, 6.8}),
and is necessarily equal to \({}_nA\), since it is so on the
reduced fiber (\SGARefLink{sga3:ix:theorem:5:1:bis}{Exposé IX, 5.1 bis}). This completes the proof of
\SGARefLink{sga3:xvi:proposition:3:4}{3.4} a) and of \SGARefLink{sga3:xvi:proposition:3:5}{3.5}.

\medskip
\noindent\textit{Proof of \SGARefLink{sga3:xvi:proposition:3:4}{3.4} b).}
The assertion to be proved is local on \(S\), so we may suppose
that \(S\) is affine, with ring \(A\). Regarding \(A\) as the
inductive limit of its finitely generated
\(\mathbb Z\)-subalgebras, we reduce as above to the case where
\(S\) is noetherian and reduced.

To prove that \(Z\) is flat, we may use the valuative criterion
for flatness (EGA IV, 11.8.1), which reduces us to the case where
\(S\) is the spectrum of a complete discrete valuation ring with
algebraically closed residue field, with generic point \(t\) and
closed point \(s\). Let \(Z'\) be the schematic closure of \(Z_t\)
in \(Z\). We must show that \(Z'=Z\), and it even suffices to show
that \(Z_s=Z'_s\) (Exposé VI\textsubscript{B}, \SGARefLink{cor:sga3-VIB-2.6}{2.6}). By Exposé
VI\textsubscript{B}, \S4, one has the inequalities
\[
  \dim Z_t=\dim Z'_t=\dim Z'_s\leq \dim Z_s.
\]
By hypothesis, however, \(\dim Z_t=\dim Z_s\); hence
\(\dim Z'_s=\dim Z_s\), and consequently \(Z'_s\) contains
\((Z_s)^0_{\mathrm{red}}\). It follows that
\[
  G'_s=G_s/Z'_s
\]
is affine (\SGARefLink{sga3:xii:6.1}{Exposé XII, 6.1}). In view of the correspondence between
Cartan subgroups of \(G\) and \(G'\) (\SGARefLink{sga3:xii:6.6}{Exposé XII, 6.6} e)),
hypothesis \SGARefLink{sga3:xvi:proposition:3:4}{3.4} b) iii) implies that the unipotent rank of \(G'_s\)
is zero, and consequently its Cartan subgroups are also its
maximal tori. The image \(Z''_s\) of \(Z_s\) in \(G'_s\) is a
central algebraic subgroup of \(G'_s\), hence is contained in
every Cartan subgroup of \(G'_s\), that is, in every maximal
torus; \(Z''_s\) is therefore a subgroup of multiplicative type
of \(G'_s\). Under these conditions, we shall show in Exposé XVII,
\S7, that there exists a finite subgroup \(M_s\) of multiplicative
type of \(Z_s\) whose image in \(Z_s/Z'_s\) is \(Z''_s\) (we used
this fact in the proof of \SGARefLink{sga3:xvi:proposition:3:4}{3.4} a) when \(Z''_s\) is étale). Since
\(G\) is smooth over \(S\) and \(S\) is the spectrum of a complete
local ring, \(M_s\) lifts to an \(S\)-subgroup \(M\) of
multiplicative type of \(G\) (\SGARefLink{sga3:ix:theorem:3:6:bis}{Exposé IX, 3.6 bis} and \SGARefLink{sga3:xv:proposition:1-6}{Exposé XV,
1.6} b)), which is central (\SGARefLink{sga3:ix:theorem:5:6}{Exposé IX, 5.6} a)). Thus \(M_t\) is
contained in \(Z'_t=Z_t\), and, since \(M\) is flat, \(M\) is
contained in \(Z'\). The group \(M_s\) is therefore contained in
\(Z'_s\), but this implies that \(Z''_s\) is the unit group, that
is, \(Z_s=Z'_s\). This completes the proof of \SGARefLink{sga3:xvi:proposition:3:4}{3.4} b).

\subsection*{Example of a Smooth Group Prescheme with Connected Fibers
Whose Center Is Not Flat}
\addcontentsline{toc}{subsection}{Example of a Smooth Group Prescheme
with Connected Fibers Whose Center Is Not Flat}

Let \(S\) be the spectrum of a discrete valuation ring \(A\), let
\(\pi\) be a uniformizer of \(A\), and let \(t\) and \(s\) be the
generic and closed points of \(S\), respectively. Let \(G\) be the
smooth affine \(S\)-group whose ring is
\[
  B=A[T,T^{-1},U]/F,\qquad F=1-T+\pi U,
\]
the composition law being defined by
\[
  (t,u),(t',u')\longmapsto (tt',\pi uu'+u+u').
\]
The generic fiber \(G_t\) is therefore isomorphic to the
multiplicative group \((\Gm)_t\), while the closed fiber \(G_s\)
is isomorphic to the additive group \((\Ga)_s\). The function
\(T\) is invertible in \(\Gamma(G,\mathcal O_G)\) and defines an
\(S\)-group morphism
\[
  \varphi:G\longrightarrow(\Gm)_S
\]
which is an isomorphism on the generic fiber and the trivial
morphism on the closed fiber. The datum of \(\varphi\) permits one
to construct the semidirect-product group
\[
  H=G\cdot(\Ga)_S
\]
(note that \((\Gm)_S\) acts on \((\Ga)_S\)). The center \(Z\) of
\(H\) is the unit group on the generic fiber and is equal to
\(H_s\) on the closed fiber, and is therefore not flat over \(S\).

% END INCLUDED SGA3_English_current_buildable_source_20260823/source/inputs/XVI/tex/components/03_expose_XVI_section3_flat_center.tex

% BEGIN INCLUDED SGA3_English_current_buildable_source_20260823/source/inputs/XVI/tex/components/04_expose_XVI_section4_affine_fibers.tex
% SGA 3, Exposé XVI, Loop-1 English reconstruction.
% Current scope: all of Section 4.
% Authority: Polo--Gille Expo16.pdf, local pages 18--19
% (combined-reader pages 956--957; printed pages 519--522).

\SGARefTarget{sga3:xvi:section:4}%
\section*{4. Groups with Affine Fibers, of Unipotent Rank Zero}
\addcontentsline{toc}{section}{4. Groups with Affine Fibers,
of Unipotent Rank Zero}

\noindent\SGARefTarget{sga3:xvi:theorem:4:1}%
\textbf{Theorem 4.1.} --- \textit{Let \(S\) be a prescheme and
let \(G\) be an \(S\)-group prescheme smooth over \(S\), of finite
presentation, with affine, connected fibers, of unipotent rank
zero (\SGARefLink{sga3:xv:definition:6:1:ter}{Exposé XV, 6.1 ter}). Then:}

\begin{enumerate}
\renewcommand{\labelenumi}{\textit{\alph{enumi})}}
\item \textit{The center \(Z\) of \(G\) is a group of
multiplicative type (and is therefore a reductive center of \(G\);
\SGARefLink{sga3:xii:4.1}{Exposé XII, 4.1}).}

\item \textit{\(G\) satisfies conditions a) through d) of \SGARefLink{sga3:xvi:theorem:3:3}{3.3}.}

\item \textit{Locally for the étale topology, \(G\) possesses
maximal tori, and these are also Cartan subgroups of \(G\). The
functor \(\mathscr{TM}_G\) of maximal tori of \(G\) is
representable by an \(S\)-prescheme smooth and affine over \(S\).}

\item \textit{\(G\) is quasi-affine over \(S\). Moreover, if
\(T\) is a maximal torus of \(G\), if}
\[
  N=\Norm_G(T)
\]
\textit{is its normalizer in \(G\), and if \(W=N/T\) is the Weyl
group relative to \(T\), the following conditions are equivalent:}
\begin{enumerate}
\renewcommand{\labelenumii}{\textit{(\roman{enumii})}}
\item \textit{\(G\) is affine over \(S\);}
\item \textit{\(G'=G/Z\) is affine over \(S\);}
\item \textit{\(N\) is affine over \(S\);}
\item \textit{\(W\) is affine over \(S\).}
\end{enumerate}
\textit{These conditions are always satisfied if \(S\) is locally
noetherian of dimension at most \(1\).}
\end{enumerate}

\medskip
\noindent\textit{Proof.}
Since \(G\) has affine fibers (and hence abelian rank zero) and
unipotent rank zero, \(G\) possesses maximal tori locally for the
étale topology (\SGARefLink{sga3:xv:corollary:8-18}{Exposé XV, 8.18}), and every maximal torus of \(G\)
is plainly a Cartan subgroup of \(G\). This proves the first part
of c). To see that \(Z\) is of multiplicative type, we may suppose,
by what precedes, that \(G\) possesses a maximal torus \(T\).
Since \(T\) is also a Cartan subgroup, it contains \(Z\) (because
\(T=\Centr_G(T)\) by \SGARefLink{sga3:xii:7.1}{Exposé XII, 7.1}), and
\[
  Z=Z\cap T
\]
is of multiplicative type by \SGARefLink{sga3:xvi:proposition:3:5}{3.5} i). This proves a). Assertion b)
is clear in view of a). On the other hand, the functor
\(\mathscr{TM}_G\) is isomorphic to the functor of Cartan
subgroups of \(G\) (\SGARefLink{sga3:xii:7.1}{Exposé XII, 7.1}), and is therefore smooth and
affine over \(S\) by \SGARefLink{sga3:xvi:theorem:3:3}{3.3} c). This finishes the proof of c). It
remains to prove d).

\medskip
\noindent\textit{Proof of d).}
Since \(G'\) is quasi-affine (\SGARefLink{sga3:xvi:theorem:3:3}{3.3} b)) and \(Z\) is of
multiplicative type, hence affine over \(S\), \(G\) is quasi-affine
(Exposé VI\textsubscript{B}, \S9).

\medskip
\noindent i) \(\Longleftrightarrow\) ii).
If \(G\) is affine, then \(G'\) is affine by \SGARefLink{sga3:ix:proposition:2:3}{Exposé IX, 2.3}. If
\(G'\) is affine, then \(G\) is affine as an extension of an affine
group by an affine group (Exposé VI\textsubscript{B}, \S9).

\medskip
\noindent i) \(\Longleftrightarrow\) iii).
If \(G\) is affine, then \(N\) is affine, since it is closed in
\(G\) (\SGARefLink{sga3:xi:corollary:6:11}{Exposé XI, 6.11} a)). Moreover, \(G/N\) is isomorphic to the
functor \(\mathscr{TM}_G\) (conjugacy of maximal tori; cf.\
\SGARefLink{sga3:xii:7.1}{Exposé XII, 7.1} b)), and is therefore affine over \(S\) by c).
Thus if \(N\) is affine over \(S\), then \(G\) is affine, as an
``extension'' of an affine homogeneous space by an affine group
(Exposé VI\textsubscript{B}, \S9).

\medskip
\noindent iii) \(\Longleftrightarrow\) iv).
If \(N\) is affine, then \(W=N/T\) is affine (\SGARefLink{sga3:ix:proposition:2:3}{Exposé IX, 2.3}), and
the converse follows from Exposé VI\textsubscript{B}, \S9.

Moreover, \(N/T\) is representable by an \(S\)-group prescheme
that is étale, separated, and of finite type over \(S\) (\SGARefLink{sga3:xv:proposition:7-1}{Exposé XV,
7.1} iv)), hence quasi-finite over \(S\). The last assertion of d)
therefore follows from the next lemma, applied with \(X=W\) and
\(Y=S\).

\medskip
\noindent\SGARefTarget{sga3:xvi:lemma:4:2}%
\textbf{Lemma 4.2.} --- \textit{Let \(S\) be a locally noetherian
prescheme of dimension \(1\), let \(X\) and \(Y\) be two
\(S\)-preschemes locally of finite type over \(S\), and let
\(u:X\to Y\) be a quasi-finite and separated \(S\)-morphism.
Suppose that, for every point \(s\) of \(S\), the morphism}
\[
  u_s:X_s\longrightarrow Y_s
\]
\textit{induced from \(u\) by the base change
\(\operatorname{Spec}\kappa(s)\to S\) is finite. Then \(u\) is an
affine morphism.}

\medskip
The assertion to be proved is local on \(Y\), which allows us to
suppose that \(Y\) (and hence also \(X\)) is of finite type over
\(S\). By EGA IV, \S8, it suffices to prove 4.2 when \(S\) is the
spectrum of a local ring \(A\). By fpqc descent we may suppose
that \(A\) is complete, then that \(A\) is reduced (EGA II,
1.6.4), and then that \(A\) is normal (Nagata's theorem, EGA
0\textsubscript{IV}, \S22, and Chevalley's theorem, EGA II,
6.7.1). If \(A\) is a field, then \(u\) is finite by hypothesis,
hence affine. Otherwise \(A\) is a discrete valuation ring; let
\(s\) and \(t\) be the closed and generic points of \(S\), and let
\(\pi\) be a uniformizer of \(A\). Let \(y\) be a point of \(Y_s\).
Applying EGA IV, \S8, once more, we may replace \(Y\) by the
spectrum \(Y'\) of \(\mathcal O_{Y,y}\), and \(X\) by
\[
  X'=X\times_Y Y';
\]
we may even suppose that \(\mathcal O_{Y,y}\) is complete. Since
\(u\) is quasi-finite and separated, EGA II, 6.2.6, shows that
\(X'\) is the disjoint union of two schemes \(X_1\) and \(X_2\),
where \(X_1\) is finite over \(Y'\) (hence affine), while
\(u(X_2)\) does not contain the closed point \(y\) of \(Y'\). I
claim that, under these conditions, \(X_2\) does not meet the
closed fiber \(X'_s\). Indeed, by hypothesis,
\[
  X'_s\longrightarrow Y'_s
\]
is finite, so the restriction of this morphism to the closed
subset \((X_2)_s\) is finite, and its image in \((Y')_s\) is
closed. Since this image does not contain the closed point of the
local scheme \((Y')_s\), \((X_2)_s\) is empty. Since \(u_t\) is
finite, the restriction to \((X_2)_t=X_2\) of the morphism
\[
  X'_t\longrightarrow Y'_t
\]
is finite. Moreover \(Y'_t=Y'_\pi\), so the open immersion
\(Y'_t\to Y'\) is affine. In short, the composite morphism
\[
  X_2\longrightarrow Y'_t\longrightarrow Y'
\]
is affine, and it follows that
\[
  X'=X_1\amalg X_2\longrightarrow Y'
\]
is affine.

\medskip
\noindent\SGARefTarget{sga3:xvi:corollary:4:3}%
\textbf{Corollary 4.3.} --- \textit{Let \(S\) be a prescheme, and
let \(G\) be an \(S\)-group prescheme smooth and of finite
presentation over \(S\), with affine, solvable, connected fibers,
of unipotent rank zero. Then \(G\) is affine over \(S\). If,
moreover, the center of \(G\) is the unit group and \(S\) is
quasi-compact, the canonical morphism}
\[
  i_n:G\longrightarrow\operatorname{GL}_S(G^{(n)})
  \qquad\textit{(cf.\ 3.1 b))}
\]
\textit{is a closed immersion for all sufficiently large \(n\).}

\medskip
To prove that \(G\) is affine over \(S\), we may suppose that
\(G\) possesses a maximal torus \(T\) (\SGARefLink{sga3:xvi:theorem:4:1}{4.1} c)). Since \(G\) has
solvable fibers, the Weyl group relative to \(T\) is the unit
group (BIBLE, 6, Th.\ 1, Cor.\ 3), and condition \SGARefLink{sga3:xvi:theorem:4:1}{4.1} d) iv) is
satisfied. If the center of \(G\) is the unit group, \(i_n\) is a
monomorphism for all sufficiently large \(n\) (\SGARefLink{sga3:xvi:proposition:3:1}{3.1}), and is
therefore a closed immersion (\SGARefLink{sga3:xvi:corollary:1:5}{1.5} b)).

% END INCLUDED SGA3_English_current_buildable_source_20260823/source/inputs/XVI/tex/components/04_expose_XVI_section4_affine_fibers.tex

% BEGIN INCLUDED SGA3_English_current_buildable_source_20260823/source/inputs/XVI/tex/components/05_expose_XVI_section5_reductive_semisimple.tex
% SGA 3, Exposé XVI, Loop-1 English reconstruction.
% Current scope: all of Section 5.
% Authority: Polo--Gille Expo16.pdf, local pages 19--21
% (combined-reader pages 957--959; printed pages 523--525).

\SGARefTarget{sga3:xvi:section:5}%
\section*{5. Application to Reductive and Semisimple Groups}
\addcontentsline{toc}{section}{5. Application to Reductive and
Semisimple Groups}

\noindent\SGARefTarget{sga3:xvi:definition:5:1}%
\textbf{Definition 5.1.} --- \textit{An \(S\)-group prescheme
\(G\) is said to be \emph{reductive} (respectively
\emph{semisimple}) if \(G\) is smooth and affine over \(S\), with
reductive (respectively semisimple) fibers.}

\medskip
\noindent\SGARefTarget{sga3:xvi:5:1:1}%
\textbf{5.1.1.}
Reductive groups will be studied systematically beginning with
Exposé XIX. In this section we shall need the following
properties, which will be proved in Exposé XIX (without using the
developments of the present Exposé).

\begin{enumerate}
\renewcommand{\labelenumi}{\textit{\alph{enumi})}}
\item \textit{Let \(S\) be a prescheme, let \(G\) be an
\(S\)-group prescheme smooth and affine over \(S\), with connected
fibers, and let \(s\) be a point of \(S\) such that \(G_s\) is
reductive (respectively semisimple). Then there exists a
neighborhood \(U\) of \(s\) such that \(G|_U\) is reductive
(respectively semisimple).}

\item \textit{If \(G\) is reductive, then locally for the étale
topology \(G\) possesses maximal tori; and if \(T\) is a maximal
torus of \(G\), the Weyl group}
\[
  W=\Norm_G(T)/T
\]
\textit{is finite over \(S\).}

\item \textit{A reductive group has unipotent rank zero.}
\end{enumerate}

Admitting these facts, we propose to strengthen assertion a)
above.

\medskip
\noindent\SGARefTarget{sga3:xvi:theorem:5:2}%
\textbf{Theorem 5.2.} --- \textit{Let \(S\) be a prescheme, let
\(G\) be an \(S\)-group prescheme smooth and of finite
presentation over \(S\), with connected fibers, and let \(s\) be
a point of \(S\). Then:}

\begin{enumerate}
\renewcommand{\labelenumi}{\textit{(\roman{enumi})}}
\item \textit{If the fibers of \(G\) are affine and \(G_s\) is
reductive, there exists an open neighborhood \(U\) of \(s\) such
that \(G|_U\) is reductive.}

\item \textit{If \(G_s\) is semisimple, there exists an open
neighborhood \(U\) of \(s\) such that \(G|_U\) is semisimple.}
\end{enumerate}

Using \SGARefLink{sga3:xv:proposition:6-2}{Exposé XV, 6.2} i), and Exposé VI\textsubscript{B}, \S10, we
reduce to the case where \(S\) is noetherian.

In case ii), consider the center \(Z\) of \(G\), which is
representable (\SGARefLink{sga3:xi:corollary:6:11}{Exposé XI, 6.11} a)). Since \(G_s\) is semisimple,
it is well known that \(Z_s\) is finite. Consequently (Exposé
VI\textsubscript{B}, \S4), there exists a neighborhood \(U\) of
\(s\) such that \(Z\) is quasi-finite over \(U\). For every point
\(t\) of \(U\), \(G_t\) is then affine (\SGARefLink{sga3:xii:6.1}{Exposé XII, 6.1}).
Shrinking \(S\) if necessary, we may therefore suppose that the
fibers of \(G\) are affine, in case ii) as well as in case i).

In view of a), to prove \SGARefLink{sga3:xvi:theorem:5:2}{5.2} it suffices to prove that \(G\) is
affine over \(S\). Since \(G\) has affine fibers, we know that the
functor of subtori of \(G\) is representable by a smooth
\(S\)-prescheme (\SGARefLink{sga3:xv:corollary:8:11}{Exposé XV, 8.11} and \SGARefLink{sga3:xv:theorem:8:8}{8.9}). After making an étale
extension covering \(s\), which is permissible, we may therefore
suppose that there is a subtorus \(T\) of \(G\) such that \(T_s\)
is a maximal torus of \(G_s\). But then
\[
  C=\Centr_G(T)
\]
is a smooth subgroup of \(G\), with connected fibers, containing
\(T\), such that \(C_s=T_s\) (because \(C_s\) is a Cartan
subgroup of \(G_s\), and \(G_s\) has unipotent rank zero by c)).
It follows that \(C=T\), so \(T\) is a Cartan subgroup and a
maximal torus of \(G\); in particular, the unipotent rank of \(G\)
is zero, and \SGARefLink{sga3:xvi:theorem:4:1}{4.1} applies.

By \SGARefLink{sga3:xvi:theorem:4:1}{4.1} d), it suffices to show that the Weyl group \(W\) relative
to \(T\) is affine over a neighborhood \(U\) of \(s\). In fact,
we shall show that \(W\) is even finite over a neighborhood of
\(s\). Since \(W\) is quasi-finite over \(S\) (\SGARefLink{sga3:xv:proposition:7-1}{Exposé XV,
7.1} iv)), saying that \(W\) is finite over a neighborhood \(U\)
of \(s\) is equivalent to saying that the morphism
\[
  W\longrightarrow S
\]
is proper at \(s\) (EGA IV, 15.7, and EGA III, 4.4.2). To
establish this, the valuative criterion of local properness
(EGA IV, 15.7) reduces us to the case where \(S\) is the spectrum
of a discrete valuation ring, with closed point \(s\). But then,
by the last assertion of \SGARefLink{sga3:xvi:theorem:4:1}{4.1} d), \(G\) is affine over \(S\), and
we may apply property a) recalled above to conclude that \(G\) is
reductive (respectively semisimple). Applying property b), we
conclude that \(W\) is indeed finite over \(S\), which completes
the proof of \SGARefLink{sga3:xvi:theorem:5:2}{5.2}.

% END INCLUDED SGA3_English_current_buildable_source_20260823/source/inputs/XVI/tex/components/05_expose_XVI_section5_reductive_semisimple.tex

% BEGIN INCLUDED SGA3_English_current_buildable_source_20260823/source/inputs/XVI/tex/components/06_expose_XVI_section6_rigidity.tex
% SGA 3, Exposé XVI, Loop-1 English reconstruction.
% Current scope: all of Section 6 and the end of Exposé XVI.
% Authority: Polo--Gille Expo16.pdf, local pages 21--23
% (combined-reader pages 959--961; printed pages 526--530).
% Local authority page 24 (combined-reader page 962) is terminal blank.

\SGARefTarget{sga3:xvi:section:6}%
\section*{6. Applications: Extension of Certain Rigidity Properties
of Tori to Groups of Unipotent Rank Zero}
\addcontentsline{toc}{section}{6. Applications: Extension of Certain
Rigidity Properties of Tori to Groups of Unipotent Rank Zero}

\noindent\SGARefTarget{sga3:xvi:proposition:6:1}%
\textbf{Proposition 6.1.} --- \textit{Let \(S\) be a prescheme and
let \(G\) be an \(S\)-group prescheme smooth over \(S\), of finite
presentation, with connected affine fibers, and whose center is of
multiplicative type (for example, \(G\) of unipotent rank zero;
cf.\ \SGARefLink{sga3:xvi:theorem:4:1}{4.1} a)). Then every closed normal subgroup \(K\) of \(G\),
of finite presentation over \(S\) and quasi-finite over \(S\), is
finite over \(S\).}

\medskip
Indeed, for every geometric point \(\bar s\) above \(S\),
\((K_{\bar s})_{\mathrm{red}}\) is a finite étale subgroup of
\(G_{\bar s}\), normal and therefore central (since \(G_{\bar s}\)
is connected). Consequently
\[
  K'=Z\cap K,
\]
where \(Z\) denotes the center of \(G\), has the same underlying
space as \(K\), and it suffices to show that \(K'\) is finite over
\(S\). Now \(K'\) is a closed quasi-finite subgroup of the group
of multiplicative type \(Z\), and is therefore finite, as is
immediate (locally on \(S\), \(K'\) is contained in \({}_nZ\) for
a suitable integer \(n\), and \({}_nZ\) is finite over \(S\)).

\medskip
\noindent\SGARefTarget{sga3:xvi:corollary:6:2}%
\textbf{Corollary 6.2.} --- \textit{Let \(S\) and \(G\) be as
above, let \(K\) be an \(S\)-group subprescheme of \(G\), of
finite presentation over \(S\), normal and closed in \(G\), and
let \(s\) be a point of \(S\). If \(K_s\) is finite (respectively
is the unit group), there exists an open neighborhood \(U\) of
\(s\) such that \(K|_U\) is finite over \(S\) (respectively is the
unit group).}

\medskip
In view of the upper semicontinuity of the dimension of the fibers
of \(K\) (Exposé VI\textsubscript{B}, \S4), after shrinking \(S\)
we may already suppose that \(K\) is quasi-finite over \(S\), and
hence finite by \SGARefLink{sga3:xvi:proposition:6:1}{6.1}. If, moreover, \(K_s\) is the unit group, then
Nakayama's lemma easily implies that \(K\) is the unit group over
\(\operatorname{Spec}\mathcal O_s\), and hence over a neighborhood
of \(s\).

\medskip
\noindent\SGARefTarget{sga3:xvi:proposition:6:3}%
\textbf{Proposition 6.3.} --- \textit{Let \(S\) be a prescheme,
let \(G\) be an \(S\)-group prescheme smooth and of finite
presentation over \(S\), with affine, connected fibers, of
unipotent rank zero, let \(H\) be an \(S\)-group prescheme, let
\(s\) be a point of \(S\), and let}
\[
  u:G\longrightarrow H
\]
\textit{be an \(S\)-group morphism such that \(\Ker u_s\) is
central. Suppose, moreover, that \(H\) is of finite presentation
over \(S\), or that \(S\) is locally noetherian and \(G\) is
locally of finite type over \(S\). Then there exists an open
neighborhood \(U\) of \(s\) such that, if \(K=\Ker u\), \(K|_U\)
is central and of multiplicative type. Moreover,}
\[
  G'=(G|_U)/(K|_U)
\]
\textit{is representable, and the morphism}
\[
  u':G'\longrightarrow H|_U
\]
\textit{induced from \(u\) by passage to the quotient is an
immersion.}

\medskip
\noindent\textit{Proof.}
Let \(Z\) be the center of \(G\), and put \(K'=K\cap Z\).

\medskip
\noindent a) \(K'\) is a subgroup of multiplicative type over a
neighborhood \(U\) of \(s\). Indeed, after making an étale
extension over a neighborhood \(U\) of \(s\), which is legitimate,
we may suppose that \(G\) possesses a maximal torus \(T\)
(\SGARefLink{sga3:xvi:theorem:4:1}{4.1} c)). But then \(T\cap K\) is a group of multiplicative type
(\SGARefLink{sga3:xv:lemma:8:3}{Exposé XV, 8.3}) whose fiber at \(s\) is central by hypothesis;
therefore \(T\cap K\) is central over a neighborhood of \(s\)
(\SGARefLink{sga3:ix:theorem:5:6}{Exposé IX, 5.6} a)), and consequently coincides with \(K'\).

\medskip
\noindent b) Let us show that \(K'=K\) over a neighborhood \(U\)
of \(s\). By a), we may already suppose that \(K'\) is of
multiplicative type, and hence flat over \(S\). Since \(K'_s=K_s\),
the natural immersion
\[
  K'\longrightarrow K
\]
is then open over a neighborhood \(U\) of \(s\) (Exposé
VI\textsubscript{B}, \SGARefLink{cor:sga3-VIB-2.6}{2.6}). If \(t\in U\), the image of \(K_t\) in
the group
\[
  G'_t=G_t/Z_t
\]
is a finite étale group, normal in \(G'_t\) (since \(K_t\) is
normal in \(G_t\)), hence central, and therefore reduced to the
unit group (\SGARefLink{sga3:xvi:theorem:3:3}{3.3} b)). This says that \(K_t\) is contained in
\(Z_t\), and hence is equal to \(K'_t\), whence \(K=K'\) over
\(U\).

\medskip
\noindent c) The representability of \(G'\) now follows from \SGARefLink{sga3:xvi:corollary:2:3}{2.3};
the fact that \(u'\) is an immersion is contained in \SGARefLink{sga3:xvi:theorem:1:3}{1.3} a), in
view of \SGARefLink{sga3:xvi:theorem:4:1}{4.1} c) applied to \(G'\).

\medskip
\noindent\SGARefTarget{sga3:xvi:proposition:6:4}%
\textbf{Proposition 6.4.} --- \textit{Let \(S\) be a prescheme,
let \(G\) be an \(S\)-group prescheme smooth and of finite
presentation over \(S\), with connected fibers, and let \(s\) be
a point of \(S\) such that \(G_s\) is generated by its subtori
(\SGARefLink{sga3:xii:8.2}{Exposé XII, 8.2}) (for example, \(G_s\) affine of unipotent rank
zero). Let \(H\) be an \(S\)-group prescheme, and let}
\[
  u,v:G\longrightarrow H
\]
\textit{be two \(S\)-homomorphisms such that \(u_s=v_s\) and
\(u_s\) is central. Suppose, moreover, that \(H\) is of finite
presentation over \(S\), or that \(S\) is locally noetherian.
Then there exists a neighborhood \(U\) of \(s\) such that}
\[
  u|_U=v|_U.
\]

\medskip
We reduce as usual to the case where \(S\) is noetherian (to study
the condition ``\(G_s\) is generated by its subtori,'' one uses
Exposé VI\textsubscript{B}, \SGARefLink{prop:sga3-VIB-7.4}{7.4}). Since \(G\) has connected fibers,
to prove that \(u=v\) over a neighborhood \(U\) of \(s\), it
suffices to prove that \(u=v\) after reduction modulo every power
of the maximal ideal of the local ring \(\mathcal O_{S,s}\), which
reduces us to the case where \(S\) is local artinian.

But then the functor of maximal subtori of \(G\) is representable
by an \(S\)-scheme \(X\), smooth and of finite type over \(S\)
(\SGARefLink{sga3:xv:corollary:8:17}{Exposé XV, 8.17}). Let \(T\) be the maximal subtorus of \(G_X\)
universal for the functor \(X\). The hypothesis on \(G_s\) means
that the algebraic subgroup of \(G_s\) generated (Exposé
VI\textsubscript{B}, \S7) by the \(\kappa(s)\)-morphism
\[
  f:T_s\longrightarrow G_s
\]
(the composite of the immersion \(T_s\to G_{X_s}\) and the
canonical projection \(G_{X_s}\to G_s\)) is all of \(G_s\). But
\(X_s\) is geometrically connected (for if \(\overline T\) is a
maximal torus of \(G_{\bar s}\), and \(N\) its normalizer in
\(G_{\bar s}\), then \(X_{\bar s}\) is isomorphic to
\(G_{\bar s}/N\)), and the image of \(T_s\) under \(f\) contains
the unit section of \(G_s\). It follows from Exposé
VI\textsubscript{B}, \SGARefLink{prop:sga3-VIB-7.4}{7.4}, that there exist an integer \(N>0\) and
an \(S\)-morphism
\[
  f^N:T^N\longrightarrow G,
  \qquad
  T^N=T\times_S\cdots\times_S T
  \quad\text{(\(N\) factors),}
\]
where \(f^N\) depends only on \(f\) and the multiplication law of
\(G\), such that \((f^N)_s\) is surjective.

Consider the base change \(X\to S\) and the restrictions of
\(u_X\) and \(v_X\) to the subtorus \(T\) of \(G_X\). By
hypothesis,
\[
  u_{X_s}=v_{X_s},
\]
and
\[
  u_{X_s}:T_{X_s}\longrightarrow H_{X_s}
\]
is a central homomorphism. It then follows from \SGARefLink{sga3:ix:theorem:5:1}{Exposé IX, 5.1},
that
\[
  u_X|_T=v_X|_T.
\]
Making the definition of \(f^N\) explicit, and taking account of
the fact that \(u\) and \(v\) are homomorphisms, one immediately
deduces that
\[
  uf^N=vf^N.
\]
But the \(\kappa(s)\)-morphism \(f_s^N\) is surjective and \(G_s\)
is smooth, hence reduced; consequently \(f_s^N\) is generically
flat. Since \(T\) is smooth over \(S\), hence flat, there exists a
nonempty open subset \(V\) of \(T^N\) such that
\[
  f^N|_V
\]
is flat (EGA IV, 11.3.10).\footnote{The French source prints \(T\)
here; \(T^N\) is the domain of \(f^N\).} The image of \(V\) is an open subset
\(W\) of \(G\) (EGA IV, 11.3.1), and
\[
  f^N:V\longrightarrow W
\]
is faithfully flat and of finite presentation, hence covering for
the fpqc topology. The equality \(uf^N=vf^N\) then implies
\[
  u|_W=v|_W.
\]
Since \(W\) is schematically dense in \(G\) (EGA IV, 11.10.10),
and \(H\) is separated over \(S\) (Exposé VI\textsubscript{A},
\SGARefLink{sga3:VIA:0.3}{0.3}), it follows immediately that \(u=v\).

% END INCLUDED SGA3_English_current_buildable_source_20260823/source/inputs/XVI/tex/components/06_expose_XVI_section6_rigidity.tex



% END INCLUDED SGA3_English_current_buildable_source_20260823/source/inputs/XVI/body.tex


\SGAStartExpose{XVII}
% BEGIN INCLUDED SGA3_English_current_buildable_source_20260823/source/inputs/XVII/body.tex
% Mechanically extracted after \mainmatter from the frozen standalone master.

\markboth{Exposé XVII. Unipotent Algebraic Groups}{Exposé XVII. Unipotent Algebraic Groups}
% BEGIN INCLUDED SGA3_English_current_buildable_source_20260823/source/inputs/XVII/tex/components/00_expose_XVII_title_and_section0.tex
% Authority: Expo17.pdf, local p.1 / combined p.963 / printed pp.531--532.

\chapter*{EXPOSÉ XVII\\[0.4em]
UNIPOTENT ALGEBRAIC GROUPS. EXTENSIONS BETWEEN UNIPOTENT GROUPS AND GROUPS OF MULTIPLICATIVE TYPE}
\addcontentsline{toc}{chapter}{Exposé XVII. Unipotent Algebraic Groups. Extensions}

\begin{center}
\textit{by M. Raynaud}%
\begingroup
\renewcommand{\thefootnote}{*}%
\footnote{Cf. the note on page 1 of Exposé XV.}%
\endgroup
\end{center}

\section*{0.\ Some notation%
\begingroup
\renewcommand{\thefootnote}{0}%
\footnote{Version x.y of 18 November 2008.}%
\endgroup}
\addcontentsline{toc}{section}{0.\ Some notation}
\SGARefTarget{sga3:xvii:section:0}

In this chapter we shall be concerned chiefly with algebraic groups defined over a
field \(k\).  The number \(p\geq 0\) will always denote the characteristic of \(k\);
\(\mathbf F_p\) will denote the prime field with \(p\) elements when \(p>0\);
\(\bar{k}\) will denote an algebraically closed extension of \(k\); and \(q\) will
denote a prime number different from \(p\).

For every \(S\)-prescheme, \((\mathbb G_a)_S\) (respectively
\((\mathbb G_m)_S\)) denotes the additive group (respectively the multiplicative
group) over \(S\) (cf.\ \SGARefLink{sga3:i:subsection:4-3}{Exposé I, 4.3}).  For every integer \(n>0\),
\((\boldsymbol\mu_n)_S\) (respectively \((\mathbb Z/n\mathbb Z)_S\)) denotes the
group of \(n\)-th roots of unity (\SGARefLink{sga3:i:paragraph:4-4-4}{Exposé I, 4.4.4}) (respectively the constant group
over \(S\) associated with the abstract group \(\mathbb Z/n\mathbb Z\)
(\SGARefLink{sga3:i:subsection:4-1}{Exposé I, 4.1})).  The group \((\mathbb Z/n\mathbb Z)_S\) is finite and étale
over \(S\); the group \((\boldsymbol\mu_n)_S\) is flat and finite over \(S\), and
is étale over \(S\) if and only if \(n\) is invertible on \(S\)
(\SGARefLink{sga3:viii:proposition:2:1}{Exposé VIII, 2.1}).

Suppose that \(S\) is a prescheme of characteristic \(p>0\).  For every integer
\(n>0\) and every \(S\)-group prescheme \(G\), we denote by \(F^n(G)\) the
radicial \(S\)-subgroup prescheme of \(G\) equal to the kernel of the \(n\)-th
iterate of the relative Frobenius morphism of \(G\) (Exposé VII\textsubscript{A}).
In particular, if \(G=(\mathbb G_a)_S\), we set
\[
  F^n(G)=(\boldsymbol\alpha_{p^n})_S .
\]
This is a radicial \(S\)-group, flat and finite over \(S\), representing the
following functor: for every \(S\)-prescheme \(S'\),
\((\boldsymbol\alpha_{p^n})_S(S')\) is the set of
\(x'\in\Gamma(S',\mathcal O_{S'})\) such that \(x'^{p^n}=0\).

The already defined group \((\boldsymbol\mu_{p^n})_S\) is canonically isomorphic
to \(F^n((\mathbb G_m)_S)\).

When there is no ambiguity about the base scheme \(S\), we shall simply write
\(\mathbb G_a\), \(\mathbb G_m\), \(\boldsymbol\alpha_{p^n}\), and so on, in
place of \((\mathbb G_a)_S\), \((\mathbb G_m)_S\),
\((\boldsymbol\alpha_{p^n})_S\), and so on.

If \(G\) is a commutative \(S\)-group prescheme, then for every integer \(n>0\),
\({}_nG\) denotes the subgroup prescheme of \(G\) equal to the kernel of
raising to the \(n\)-th power in \(G\).

For the reader's convenience, we have collected in an appendix some properties
of algebraic groups proved in Exposés VI and VII, together with elementary
properties of Hochschild cohomology that will be useful in this chapter.

% END INCLUDED SGA3_English_current_buildable_source_20260823/source/inputs/XVII/tex/components/00_expose_XVII_title_and_section0.tex

% BEGIN INCLUDED SGA3_English_current_buildable_source_20260823/source/inputs/XVII/tex/components/01_expose_XVII_section1_opening.tex
% Authority: Expo17.pdf, local p.2 / combined p.964 / printed pp.533--534.

\SGASection{1}{Definition of unipotent algebraic groups}
\SGARefTarget{sga3:xvii:section:1}

\SGAStatement{Definition 1.1}{%
An algebraic group \(G\) defined over an algebraically closed field \(k\) is
called \emph{unipotent} if \(G\) admits a composition series whose successive
quotients are isomorphic to algebraic subgroups of \((\mathbb G_a)_k\).}
\SGARefTarget{sga3:xvii:definition:1:1}

\SGAStatement{Proposition 1.2}{%
Let \(k\) be an algebraically closed field, let \(K\) be an algebraically closed
extension of \(k\), and let \(G\) be an algebraic group defined over \(k\).
Then \(G\) is unipotent if and only if \(G_K\) is unipotent.}
\SGARefTarget{sga3:xvii:proposition:1:2}

The necessity is clear, since a composition series remains a composition series
after extension of the base.

The proof of sufficiency is standard.  The field \(K\) is the filtered direct
limit of its finitely generated \(k\)-subalgebras.  By Exposé
VI\textsubscript{B}, \S10, one can find a finitely generated \(k\)-subalgebra
\(A\) of \(K\), a composition series \(G_i\) of \(G_S\), where
\(S=\operatorname{Spec}A\), and immersions
\[
  u_i:H_i=G_i/G_{i+1}\longrightarrow(\mathbb G_a)_S .
\]
To prove that \(G\) is unipotent, it therefore suffices to make a \(k\)-valued
base change \(A\to k\).  This is possible because
\(\operatorname{Hom}_{k\text{-alg}}(A,k)\) is nonempty: \(A\) is a nonzero
finitely generated algebra over the algebraically closed field \(k\).

\SGAStatement{Definition 1.3}{%
Let \(G\) be an algebraic group defined over a field \(k\).  We say that \(G\)
is \emph{unipotent} if there exists an algebraically closed extension
\(\bar{k}\) of \(k\) such that \(G_{\bar{k}}\) is unipotent in the sense of
\SGARefLink{sga3:xvii:definition:1:1}{Definition 1.1}.}
\SGARefTarget{sga3:xvii:definition:1:3}

By \SGARefLink{sga3:xvii:proposition:1:2}{Proposition 1.2}, this property is independent of the chosen algebraically
closed extension \(\bar{k}\).

\SGAStatement{Definition 1.4}{%
Let \(G\) be an algebraic group defined over a field \(k\), and let \(H\) be an
algebraic group defined over an extension \(k'\) of \(k\).  We say that \(H\)
is a \emph{form} of \(G\) over \(k'\) if the algebraic groups
\(G_{\bar{k}'}\) and \(H_{\bar{k}'}\) become isomorphic over an algebraic
closure \(\bar{k}'\) of \(k'\).  As above, this property is independent of
the choice of \(\bar{k}'\).  We shall also say that \(H\) is a
\emph{twisted form} of \(G\).}
\SGARefTarget{sga3:xvii:definition:1:4}

We can now describe the algebraic subgroups of \(\mathbb G_a\).

\SGAStatement{Proposition 1.5}{%
Let \(k\) be a field of characteristic \(p\geq0\).  An algebraic subgroup
\(H\) of \((\mathbb G_a)_k\) is of one of the following types:
\begin{enumerate}
\item[\rm(i)] \(H=0\);
\item[\rm(ii)] \(H=\mathbb G_a\);
\item[\rm(iii)]\SGAInlineRefTarget{sga3:xvii:proposition:1:5:item:iii} if \(p>0\), \(H\) is an extension of a twisted constant group
  \((\mathbb Z/p\mathbb Z)^r\) by a radicial group
  \(\boldsymbol\alpha_{p^n}\), where \(r,n\geq0\) and \(r+n>0\).
  If, moreover, \(k\) is perfect, this extension is necessarily trivial.
\end{enumerate}}
\SGARefTarget{sga3:xvii:proposition:1:5}

% END INCLUDED SGA3_English_current_buildable_source_20260823/source/inputs/XVII/tex/components/01_expose_XVII_section1_opening.tex

% BEGIN INCLUDED SGA3_English_current_buildable_source_20260823/source/inputs/XVII/tex/components/02_expose_XVII_section1_remainder_and_section2_opening.tex
% Authority: Expo17.pdf, local pp.3--4 / combined pp.965--966 /
% printed pp.535--539.

\SGAProof{Proof}
If \(\dim H=1\), then clearly \(H=\mathbb G_a\).  Otherwise \(H\) has
dimension zero and is therefore an extension\footnote{Editor's note:
reference to Exposé VII\textsubscript{A}?} of an étale group \(H''\) by its
identity component \(H'\), which is radicial.  To describe \(H''_{\bar{k}}\),
it is enough to know the abstract group \(H''(\bar{k})\), which is isomorphic
to \(H(\bar{k})\).

The latter is a finite subgroup of \(\mathbb G_a(\bar{k})\).  It is therefore
zero when \(p=0\), and otherwise is of the form
\((\mathbb Z/p\mathbb Z)^r\), since it is annihilated by \(p\).  The group
\(H'\) is closed in \(\mathbb G_a\) and is defined by a single equation
(the coordinate ring \(\bar{k}[T]\) of \(\mathbb G_a\) is principal).
Since this equation has \(0\) as its only root over \(\bar{k}\), it is of the
form \(T^n=0\).  Compatibility with the group law implies that
\((T+T')^n\) belongs to the ideal generated by \(T^n\) and \(T'^n\) in
\(\bar{k}[T,T']\).  Thus, if \(p=0\), then \(n=1\) and \(H'\) is the trivial
group; while if \(p>0\), then \(n=p^m\) and
\(H'=\boldsymbol\alpha_{p^m}\).

The last assertion of \SGARefLink{sga3:xvii:proposition:1:5}{Proposition 1.5} follows more generally from the
following lemma.

\SGAStatement{Lemma 1.6}{%
If \(k\) is a perfect field, every extension \(H\) of an étale algebraic group
\(H''\) by a radicial group \(H'\) is trivial.  Moreover, \(H''\) has a unique
lifting in \(H\), namely \(H_{\mathrm{red}}\).}
\SGARefTarget{sga3:xvii:lemma:1:6}

Indeed, since \(k\) is perfect, the reduced \(k\)-scheme
\(H_{\mathrm{red}}\) is an algebraic subgroup of \(H\)
(Exposé VI\textsubscript{A}, \SGARefLink{sga3:VIA:0.2}{0.2}), is geometrically reduced and hence smooth
over \(k\) (Exposé VI\textsubscript{A}, \SGARefLink{prop:sga3-VIA-1.3.1}{1.3.1}), and is therefore étale because
\(H\) has dimension zero.  To see that the canonical projection
\(H_{\mathrm{red}}\to H''\) is an isomorphism, it is enough to check this after
the base-field extension \(k\to\bar{k}\); there it is enough to check that it
induces an isomorphism on \(\bar{k}\)-valued points, which is clear.  The final
assertion follows because every lifting of \(H''\) in \(H\), being étale over
\(k\), is reduced and hence necessarily contained in \(H_{\mathrm{red}}\).

Notice that \(\boldsymbol\alpha_{p^n}\) is an iterated extension of groups
isomorphic to \(\boldsymbol\alpha_p\).  \SGARefLink{sga3:xvii:proposition:1:5}{Proposition 1.5} therefore gives the
following corollary.

\SGAStatement{Corollary 1.7}{%
Let \(G\) be an algebraic group defined over an algebraically closed field
\(k\).  Then \(G\) is unipotent if and only if it has a composition series
whose successive quotients are isomorphic to \(\mathbb G_a\) when \(p=0\),
and to one of the groups
\(\mathbb G_a,\mathbb Z/p\mathbb Z,\boldsymbol\alpha_p\) when \(p>0\).
We shall call these the \emph{elementary unipotent groups}.}
\SGARefTarget{sga3:xvii:corollary:1:7}

\SGASection{2}{First properties of unipotent groups}
\SGARefTarget{sga3:xvii:section:2}

\SGAStatement{Proposition 2.1}{%
A unipotent algebraic group defined over a field \(k\) is affine over \(k\).}
\SGARefTarget{sga3:xvii:proposition:2:1}

By fpqc descent for affine morphisms, it suffices to prove the assertion when
\(k\) is algebraically closed.  In that case, \(G\) is by definition an
iterated extension of affine algebraic groups, and is therefore affine by
Exposé VI\textsubscript{B}, \SGARefLink{concl:sga3-VIB-9.2-viii}{9.2 (viii)}, applied to affine
morphisms.\footnote{Editor's note: verify this reference.}

\SGAStatement{Proposition 2.2}{%
\begin{enumerate}
\item[\rm\SGAInlineRefTarget{sga3:xvii:proposition:2:1:item:i}(i)] The property that an algebraic group is unipotent is invariant
  under extension of the base field.
\item[\rm\SGAInlineRefTarget{sga3:xvii:proposition:2:1:item:ii}(ii)] Every algebraic subgroup of a unipotent group is unipotent.
\item[\rm\SGAInlineRefTarget{sga3:xvii:proposition:2:1:item:iii}(iii)] Every algebraic quotient group of a unipotent group is
  unipotent.
\item[\rm\SGAInlineRefTarget{sga3:xvii:proposition:2:1:item:iv}(iv)] Every extension of a unipotent algebraic group by a unipotent
  algebraic group is likewise unipotent.
\end{enumerate}}
\SGARefTarget{sga3:xvii:proposition:2:2}

\SGAProof{Proof}
Assertion \SGARefLink{sga3:xvii:proposition:2:1:item:i}{(i)} follows immediately from \SGARefLink{sga3:xvii:definition:1:3}{Definitions 1.3} and \SGARefLink{sga3:xvii:proposition:1:2}{Proposition 1.2}.
To prove the remaining assertions, we may suppose that \(k\) is algebraically
closed.  Assertion \SGARefLink{sga3:xvii:proposition:2:1:item:iv}{(iv)} is then immediate from the definition.

Let \(G\) be a unipotent algebraic group, \(G'\) an algebraic subgroup of
\(G\), \(G''\) an algebraic quotient group of \(G\), and let
\(G_i\), \(1\leq i\leq n\), be a composition series of \(G\) such that
\[
  H_i=G_i/G_{i+1}
\]
is an elementary unipotent group (\SGARefLink{sga3:xvii:corollary:1:7}{Corollary 1.7}).

For \SGARefLink{sga3:xvii:proposition:2:1:item:ii}{(ii)}, consider the composition series of \(G'\) induced by that of \(G\):
\(G'_i=G_i\cap G'\).  The group \(G'_i/G'_{i+1}\) identifies with an
algebraic subgroup of \(H_i\), hence is isomorphic to a subgroup of
\(\mathbb G_a\).  It follows that \(G'\) is unipotent.

For \SGARefLink{sga3:xvii:proposition:2:1:item:iii}{(iii)}, consider the composition series of \(G''\) obtained as the image
of that of \(G\): let \(G''_i\) be the image of \(G_i\) in \(G''\).  Then
\(G''_i/G''_{i+1}\) is a quotient of \(H_i\), so it remains only to prove the
following lemma.

\SGAStatement{Lemma 2.3}{%
Let \(H\) be an elementary unipotent group (\SGARefLink{sga3:xvii:corollary:1:7}{Corollary 1.7}) defined over a
field \(k\).  Every algebraic quotient group \(H''\) of \(H\) is either zero
or isomorphic to \(H\) over \(k\).}
\SGARefTarget{sga3:xvii:lemma:2:3}

\SGAProof{Proof}
If \(H=\mathbb G_a\) in characteristic zero, or if
\(H=\boldsymbol\alpha_p\) or \(\mathbb Z/p\mathbb Z\) in characteristic
\(p>0\), it is enough to observe from \SGARefLink{sga3:xvii:proposition:1:5}{Proposition 1.5} that \(H\) has no
algebraic subgroups other than \(0\) and \(H\).  Indeed, a nonzero algebraic
subgroup of \(\boldsymbol\alpha_p\) is defined by a \(k\)-algebra of rank at
least \(p\) over \(k\), by \SGARefLink{sga3:xvii:proposition:1:5:item:iii}{Proposition 1.5(iii)}, and is therefore
\(\boldsymbol\alpha_p\).

Now let \(H=\mathbb G_a\) in characteristic \(p>0\), so that there is an exact
sequence
\[
  0\longrightarrow N\longrightarrow\mathbb G_a
    \longrightarrow H''\longrightarrow0.
\]
If \(N=\mathbb G_a\), then \(H''=0\).  Otherwise, arguing by induction on the
length of a composition series of \(N\), we may suppose that
\(N=\boldsymbol\alpha_p\), or that \(N\) is a form of
\((\mathbb Z/p\mathbb Z)^r\).

\begin{enumerate}
\item[\rm(a)] If \(N\simeq\boldsymbol\alpha_p\), the proof of
  Proposition 1.5(iii) shows that \(N\) is necessarily the kernel of the
  Frobenius morphism \(F\) of \(\mathbb G_a\), and the assertion follows from
  the exact sequence
  \[
    0\longrightarrow\boldsymbol\alpha_p\longrightarrow\mathbb G_a
      \xrightarrow{F}\mathbb G_a\longrightarrow0.
  \]
\item[\rm(b)] If \(N\simeq\mathbb Z/p\mathbb Z\), there is plainly an
  \(a\in k^\times\) such that \(N\) is the closed subscheme of
  \(\mathbb G_a=\operatorname{Spec}k[X]\) defined by
  \(X^p-aX=0\).  The assertion then follows from the exact sequence
  \[
    0\longrightarrow N\longrightarrow\mathbb G_a
      \xrightarrow{P}\mathbb G_a\longrightarrow0,
  \]
  where \(P\) is the Artin--Schreier morphism \(x\mapsto x^p-ax\).
\item[\rm(c)] If \(N\) is a form of \((\mathbb Z/p\mathbb Z)^r\), there is a
  finite Galois extension \(k'/k\) that trivializes \(N\).  By \SGARefLink{sga3:xvii:section:b}{(b)} and an
  evident induction on \(r\), the quotient \(\mathbb G_a/N\) is a form of
  \(\mathbb G_a\) trivialized by \(k'\).  It remains to apply the next lemma.
\end{enumerate}

\SGAStatement{Lemma 2.3 bis}{%
Let \(k\) be a field, and let \(G\) be a \(k\)-algebraic group that is a form
of \(\mathbb G_a\), trivialized by a finite separable extension \(k'/k\).
Then \(G\) is isomorphic to \(\mathbb G_a\).}
\SGARefTarget{sga3:xvii:lemma:2:3:bis}

% END INCLUDED SGA3_English_current_buildable_source_20260823/source/inputs/XVII/tex/components/02_expose_XVII_section1_remainder_and_section2_opening.tex

% BEGIN INCLUDED SGA3_English_current_buildable_source_20260823/source/inputs/XVII/tex/components/03_expose_XVII_section2_continued_page005.tex
% Authority: Expo17.pdf, local p.5 / combined p.967 / printed pp.539--541.

\SGAProof{Proof of \SGARefLink{sga3:xvii:lemma:2:3:bis}{Lemma 2.3 bis}}
Indeed, the group of \(k'\)-automorphisms of the algebraic group
\((\mathbb G_a)_{k'}\) is the group of nonzero homotheties
\((k')^\times\) (\textit{Bible}, Exposé 9, \SGARefLink{sga3:xvii:lemma:3:1}{Lemma 1}), and the Galois
cohomology group \(H^1(k'/k,\mathbb G_m)\) is zero (we may suppose that
\(k'\) is a Galois extension of \(k\)), by Hilbert's Theorem 90.
\SGARefLink{sga3:xvii:lemma:2:3:bis}{Lemma 2.3 bis} then follows from the classification of the \(k\)-forms
of an algebraic group (cf.\ J.-P. Serre, \textit{Cohomologie
galoisienne}, Chapter III, \SGARefLink{sga3:iii:subsection:1-3}{1.3}).

This completes the proof of \SGARefLink{sga3:xvii:proposition:2:2}{2.2}.

\SGAStatement{Proposition 2.4}{Let \(k\) be a field, let \(M\) be a
\(k\)-group of multiplicative type (Exposé VIII) and of finite type,
and let \(U\) be a unipotent \(k\)-algebraic group.  Then:
\begin{enumerate}
\item[\SGAInlineRefTarget{sga3:xvii:proposition:2-4:item:i}(i)]\SGAInlineRefTarget{sga3:xvii:proposition:2:4:item:i} \(\underline{\operatorname{Hom}}_{k\text{-gr}}(M,U)
  =\underline e\)
  (a fortiori \(\operatorname{Hom}_{k\text{-gr}}(M,U)=e\));
\item[(ii)] \(\operatorname{Hom}_{k\text{-gr}}(U,M)=e\).
\end{enumerate}}
\SGARefTarget{sga3:xvii:proposition:2:4}

To prove (i), we must establish that, for every prescheme \(S\) over
\(k\),
\[
  \operatorname{Hom}_{S\text{-gr}}(M_S,U_S)=e.
\]
This follows from the next lemma.

\SGAStatement{Lemma 2.5}{Let \(S\) be a prescheme, let \(M\) be an
\(S\)-group of multiplicative type and of finite type over \(S\), and
let \(U\) be an \(S\)-group prescheme, of finite presentation over
\(S\), whose fibers are unipotent.  Then
\(\operatorname{Hom}_{S\text{-gr}}(M,U)=e\).}
\SGARefTarget{sga3:xvii:lemma:2:5}

\SGAProof{Proof}
Under these hypotheses, for an \(S\)-group morphism \(u:M\to U\) to
be the unit homomorphism, it suffices that the restriction of \(u\) to
the fibers of \(M\) over the points of \(S\) be the zero morphism
(\SGARefLink{sga3:ix:corollary:5:2}{Exposé IX, 5.2}).  We are therefore reduced to the case in which \(S\)
is the spectrum of a field, which we may moreover suppose to be
algebraically closed.  In view of \SGARefLink{sga3:xvii:definition:1:1}{Definition 1.1}, it is enough to take
\(U=\mathbb G_a\), in which case the assertion has already been proved
(\SGARefLink{sga3:xii:4.4.1}{Exposé XII, 4.4.1}).

We now prove Proposition 2.4(ii).  Let \(u:U\to M\) be a \(k\)-group
morphism.  The image \(u(U)\) is represented by an algebraic subgroup
\(U''\) of \(M\) (Exposé VI\textsubscript{B}, \SGARefLink{cor:sga3-VIB-5.4}{5.4}).\footnote{Editor's
note: verify this reference.}  The group \(U''\) is unipotent as a
quotient of a unipotent group (2.2(iii)), and is of multiplicative type
as a subgroup of a group of multiplicative type (cf.\ \textit{Bible},
Exposé 4, Theorem 2, \SGARefLink{sga3:xvii:corollary:C:5:1}{Corollary 1}, or \SGARefLink{sga3:ix:theorem:6:8}{Exposé IX, 6.8}).  Thus \(U''\) is
the unit group by \SGARefLink{sga3:xvii:proposition:2-4:item:i}{Proposition 2.4(i)}.

\SGAStatement{\SGAInlineRefTarget{sga3:xvii:remark:2-6}Remark 2.6}{Keeping the notation of \SGARefLink{sga3:xvii:proposition:2:4}{Proposition 2.4}, it
is no longer true in general that the functor
\(\underline{\operatorname{Hom}}_{k\text{-gr}}(U,M)\) is equal to
\(\underline e\).  Indeed, let \(S\) be a prescheme such that
\(\Gamma(S,\mathcal O_S)\) contains a nonzero element \(\varepsilon\)
with \(\varepsilon^2=0\) (for example, the spectrum of the algebra of
dual numbers over a ring \(A\)).  For every \(S'\) over \(S\), the map
\[
  u\longmapsto 1+\varepsilon_{S'}u
\]
defines, functorially in \(S'\), a homomorphism from the additive group
\(\Gamma(S',\mathcal O_{S'})\) to the multiplicative group
\(\Gamma(S',\mathcal O_{S'}^\times)\).  It therefore defines an
\(S\)-group morphism
\[
  (\mathbb G_a)_S\longrightarrow(\mathbb G_m)_S,
\]
and, since \(\varepsilon\ne0\), this morphism is not zero.}
\SGARefTarget{sga3:xvii:remark:2:6}

Recall (Exposé VII\textsubscript{A}, \(\S3\)) that when \(G\) is a
commutative \(S\)-group prescheme, finite and flat over \(S\), Cartier
duality is available and \(G\) is reflexive in the sense of Exposé
VIII, \(\S1\).  More precisely, the functor
\[
  G\longmapsto
  \underline{\operatorname{Hom}}_{S\text{-gr}}(G,\mathbb G_m)
\]
is represented

% END INCLUDED SGA3_English_current_buildable_source_20260823/source/inputs/XVII/tex/components/03_expose_XVII_section2_continued_page005.tex

% BEGIN INCLUDED SGA3_English_current_buildable_source_20260823/source/inputs/XVII/tex/components/04_expose_XVII_cartier_duality_and_section3_opening_page006.tex
% Authority: Expo17.pdf, local p.6 / combined p.968 / printed pp.541--543.

by a commutative \(S\)-group prescheme \(D(G)\), finite and flat over
\(S\), and the canonical morphism
\[
  G\longrightarrow D(D(G))
\]
is an isomorphism.  In particular:
\begin{enumerate}
\item[(i)] \(D(\boldsymbol\mu_n)\simeq\mathbb Z/n\mathbb Z\), and
  consequently \(D(\mathbb Z/n\mathbb Z)\simeq\boldsymbol\mu_n\);
\item[(ii)] if \(S\) has characteristic \(p>0\), then
  \(D(\boldsymbol\alpha_p)\simeq\boldsymbol\alpha_p\).
\end{enumerate}

\SGASection{3}{Unipotent groups acting on a vector space}
\SGARefTarget{sga3:xvii:section:3}

Recall (\SGARefLink{sga3:i:definition:4-6-1}{Exposé I, 4.6.1}) that if \(S\) is a prescheme and \(M\) is a
sheaf of \(\mathcal O_S\)-modules, \(\mathbf W(M)\) denotes the
\(S\)-functor
\[
  (\operatorname{Sch}/S)^\circ\longrightarrow(\operatorname{Ens})
\]
defined, for every \(S\)-prescheme \(S'\), by
\[
  \mathbf W(M)(S')
  =\Gamma\!\left(S',M\otimes_{\mathcal O_S}\mathcal O_{S'}\right).
\]

Recall also that if an \(S\)-group \(G\) acts on an \(S\)-functor
\(V\), the \(S\)-functor \(V^G\) of invariants of \(V\) under \(G\) is
the subfunctor of \(V\) whose \(S'\)-valued points, for \(S'\) over
\(S\), are those \(x\in V(S')\) such that \(x_{S''}\) is fixed under
\(G(S'')\) for every \(S''\) over \(S'\).

We then have the following lemma.

\SGAStatement{\SGAInlineRefTarget{sga3:xvii:lemma:3-1}Lemma 3.1}{Let \(S\) be a prescheme, and let \(G\) be an
\(S\)-group prescheme, affine over \(S\), defined by the
quasi-coherent \(\mathcal O_S\)-algebra \(A\).  Suppose that \(G\) acts
on a quasi-coherent sheaf \(M\) of \(\mathcal O_S\)-modules.  Let
\[
  \mu:M\longrightarrow A\otimes_{\mathcal O_S}M
\]
be the comorphism defining the action of \(G\) on \(M\) (\SGARefLink{sga3:i:proposition:4-7-2}{Exposé I,
4.7.2}), and let
\[
  \nu:M\longrightarrow A\otimes_{\mathcal O_S}M
\]
be the morphism \(x\mapsto\mu(x)-1\otimes x\).  Then:
\begin{enumerate}
\item[\SGAInlineRefTarget{sga3:xvii:lemma:3-1:item:i}(i)] \(\underline M^{\,G}(S)=\Gamma\operatorname{Ker}(\nu)\);
\item[\SGAInlineRefTarget{sga3:xvii:lemma:3-1:item:ii}(ii)] if \(S\) is the spectrum of a field \(k\), then
  \(\underline M^{\,G}\) is of the form \(\mathbf W(N)\), where \(N\)
  is a vector subspace of \(M\);
\item[\SGAInlineRefTarget{sga3:xvii:lemma:3-1:item:iii}(iii)] if \(S\) is the spectrum of a field \(k\), every element
  \(x\) of \(M\) is contained in a finite-dimensional \(k\)-vector
  subspace of \(M\) stable under the action of \(G\).
\end{enumerate}}
\SGARefTarget{sga3:xvii:lemma:3:1}

\SGAProof{Proof}
Assertion \SGARefLink{sga3:xvii:lemma:3-1:item:iii}{(iii)} is recalled here for reference and was already proved
in Exposé VI\textsubscript{B}, \SGARefLink{cor:sga3-VIB-11.2}{11.2}.

For \SGARefLink{sga3:xvii:lemma:3-1:item:i}{(i)}, it is clear that \(\underline M^{\,G}(S)\) contains
\(\Gamma\operatorname{Ker}(\nu)\).  To prove the reverse inclusion, we
may suppose that \(S\) is affine, with ring \(B\).  Let
\(m\in\underline M^{\,G}(S)\).  For every \(B\)-algebra \(B'\) and
every
\[
  u\in\operatorname{Hom}_{B\text{-alg}}(A,B')
\]
(where \(u\) corresponds to an element of
\(G(\operatorname{Spec}B')\)), we have
\[
  (u\otimes1_M)\nu(m)=0
  \qquad\text{in }B'\otimes_BM.
\]
Taking \(B'=A\) and \(u\) to be the identity of \(A\), we obtain
\(\nu(m)=0\).

For (ii), let \(N\) be the kernel of \(\nu\), equal to
\(\underline M^{\,G}(k)\) by \SGARefLink{sga3:xvii:lemma:3-1:item:i}{(i)}.  Every \(k\)-prescheme \(S\) is flat
over \(k\); hence
\[
  \underline M^{\,G}(S)
  =\Gamma\operatorname{Ker}(\nu\otimes_k S)
  =\Gamma(N\otimes_k S).
\]
Thus \(\underline M^{\,G}\) is isomorphic to the functor
\(\mathbf W(N)\).

\SGAStatement{Proposition 3.2}{Let \(G\) be a unipotent algebraic
group, defined over a field \(k\), acting on a \(k\)-vector space
\(V\).  If \(V\ne0\), then \(V^G\ne0\).}
\SGARefTarget{sga3:xvii:proposition:3:2}

% END INCLUDED SGA3_English_current_buildable_source_20260823/source/inputs/XVII/tex/components/04_expose_XVII_cartier_duality_and_section3_opening_page006.tex

% BEGIN INCLUDED SGA3_English_current_buildable_source_20260823/source/inputs/XVII/tex/components/05_expose_XVII_proposition3_2_proof_and_lemma3_3_page007.tex
% Authority: Expo17.pdf, local p.7 / combined p.969 / printed pp.544--545.

In view of \SGARefLink{sga3:xvii:lemma:3-1:item:ii}{3.1(ii)} and (iii), we may suppose that \(k\) is
algebraically closed and that \(V\) is finite-dimensional over \(k\).

Let
\[
  0\longrightarrow G'\longrightarrow G\longrightarrow G''
    \longrightarrow0
\]
be an exact sequence of algebraic groups, and suppose that \(G\) acts
on a sheaf \(V\) (for the fpqc topology).  Of course
\(V^G\subset V^{G'}\), and the quotient presheaf \(G/G'\) acts
naturally on \(V^{G'}\).  But \(V^{G'}\) is a sheaf, being the kernel
of the familiar pair of morphisms
\[
  V\rightrightarrows\underline{\operatorname{Hom}}(G',V).
\]
Consequently the sheaf associated with \(G/G'\), namely \(G''\),
acts on \(V^{G'}\), and it is immediate that
\[
  \left(V^{G'}\right)^{G''}
  =V^G
  =\left(V^{G'}\right)^{G/G'}.
\]

This remark reduces the proof of \SGARefLink{sga3:xvii:proposition:3:2}{3.2} to the case in which \(G\) is an
elementary unipotent group (\SGARefLink{sga3:xvii:corollary:1:7}{1.7}).

\SGARefTarget{sga3:xvii:proposition:3:2:case:a}
\textit{\SGARefLink{sga3:xvii:appendix:A}{(a)} \(G=\mathbb G_a\), \(p=0\).}
It follows from (BIBLE 4, Proposition 4) that a morphism from
\(\mathbb G_a\) to the linear group \(\operatorname{GL}(V)\) is given
by an exponential map
\[
  T\longmapsto\sum_{q=0}^{\infty}T^q\,\frac{n^q}{q!},
\]
where \(n\) is a nilpotent endomorphism of \(V\).  But then
\(V\ne0\) implies \(\operatorname{Ker}n\ne0\), and it is clear that
every vector of \(V\) annihilated by \(n\) is fixed by \(G\).

Suppose now that \(p>0\).

\SGARefTarget{sga3:xvii:proposition:3:2:case:b}
\textit{\SGARefLink{sga3:xvii:section:b}{(b)} \(G=\boldsymbol\alpha_p\).}
Since \(\boldsymbol\alpha_p\) is a radicial group of height \(1\)
(Exposé VII\textsubscript{A}, \S7), giving a representation of
\(\boldsymbol\alpha_p\) in \(V\) amounts to giving a representation
of the \(p\)-Lie algebra \(\boldsymbol\alpha_p\) in
\(\mathfrak{gl}(V)\) (Appendix II, \SGARefLink{sga3:ii:subsection:2-2}{2.2}), that is, in this case, an
element \(X\in\operatorname{End}(V)\) such that \(X^p=0\)
(Appendix II, \SGARefLink{sga3:xvii:proposition:2:1}{2.1}).  But then
\[
  V\ne0\quad\Longrightarrow\quad
  W=\operatorname{Ker}(X)\ne0,
\]
and, again by Appendix II, \SGARefLink{sga3:ii:subsection:2-2}{2.2}, one has
\(W=V^{\boldsymbol\alpha_p}\).

\SGARefTarget{sga3:xvii:proposition:3:2:case:c}
\textit{\SGARefLink{sga3:xvii:appendix:C}{(c)} \(G=\mathbb Z/p\mathbb Z\).}
A representation of \(G\) in \(V\) is equivalent to the choice of an
element \(x\in\operatorname{Aut}(V)\) such that \(x^p=1\), that is,
\((1-x)^p=0\).  Thus \(x\) is of the form \(1+n\), with \(n\)
nilpotent, and \(W=\operatorname{Ker}n\) is fixed by \(x\).

\SGARefTarget{sga3:xvii:proposition:3:2:case:d}
\textit{(d) \(G=\mathbb G_a\).}
Let \(G_i\), \(i\in I\), be the increasing filtered family of étale
algebraic subgroups of \(\mathbb G_a\), hence groups isomorphic to
\((\mathbb Z/p\mathbb Z)^{r_i}\) (\SGARefLink{sga3:xvii:proposition:1:5}{Proposition 1.5}), and put
\(V_i=V^{G_i}\).  Since \(V\) has finite nonzero dimension and each
\(V_i\) is nonzero by \SGARefLink{sga3:xvii:appendix:C}{(c)}, the decreasing filtered family of the
\(V_i\) is stationary, and
\[
  W=\bigcap_{i\in I}V_i\ne0.
\]
We shall use the following lemma.

\SGAStatement{Lemma 3.3}{The family of étale subgroups of
\((\mathbb G_a)_S\), where \(S\) is a prescheme of characteristic
\(p>0\), is schematically dense in \(G\) (\SGARefLink{sga3:ix:definition:4:1}{Exposé IX, 4.1}).}
\SGARefTarget{sga3:xvii:lemma:3:3}

By \SGARefLink{sga3:ix:lemma:4:4}{Exposé IX, 4.4}, it suffices to prove the lemma when \(S\) is the
spectrum of the prime field \(\mathbb F_p\).  In that case it is
enough to consider the family of étale subgroups \(G_n\), \(n\geq1\),
defined by
\[
  X^{p^n}-X=0.
\]
This family is schematically dense in
\((\mathbb G_a)_{\mathbb F_p}\), since it contains every closed point.

With this established, let us return to the proof of \SGARefLink{sga3:xvii:proposition:3:2:case:d}{3.2(d)}.  If
\(w\in W\), its stabilizer in \(G\) is an algebraic subgroup of
\(\mathbb G_a\) containing \(G_i\) for every \(i\in I\); it is
therefore equal to \(\mathbb G_a\) by \SGARefLink{sga3:xvii:lemma:3:3}{3.3}.  Consequently
\(W=V^G\).

\SGARefLink{sga3:xvii:proposition:3:2}{Proposition 3.2} immediately gives the following.

% END INCLUDED SGA3_English_current_buildable_source_20260823/source/inputs/XVII/tex/components/05_expose_XVII_proposition3_2_proof_and_lemma3_3_page007.tex

% BEGIN INCLUDED SGA3_English_current_buildable_source_20260823/source/inputs/XVII/tex/components/06_expose_XVII_corollary3_4_theorem3_5_definition3_6_page008.tex
% Authority: Expo17.pdf, local p.8 / combined p.970 / printed pp.546--547.

\SGAStatement{Corollary 3.4}{Let \(k\) be a field, and let \(G\) be a
unipotent \(k\)-algebraic group acting on a finite-dimensional
\(k\)-vector space \(V\).  Then \(V\) has a sequence of vector
subspaces \(V_i\), defined over \(k\) and stable under \(G\),
\[
  0=V_0\subset V_1\subset\cdots\subset V_n=V,
\]
such that \(G\) acts trivially on \(V_{i+1}/V_i\).  Moreover, the
\(V_i\) may be chosen so that \(V_{i+1}/V_i\) has dimension \(1\).}
\SGARefTarget{sga3:xvii:corollary:3:4}

We shall now summarize and supplement the properties of unipotent
groups already proved in the following theorem.

\SGAStatement{\SGAInlineRefTarget{sga3:xvii:theorem:3-5}Theorem 3.5}{Let \(G\) be an algebraic group defined
over a field \(k\).  The following properties are equivalent:
\begin{enumerate}
\item[(i)] \SGARefTarget{sga3:xvii:theorem:3:5:item:i}
  \(G\) is unipotent.
\item[(ii)] \SGARefTarget{sga3:xvii:theorem:3:5:item:ii}
  \(G\) has a composition series defined over \(k\), whose successive
  quotients are isomorphic to \(\mathbb G_a\) if \(p=0\), and,
  respectively, to \(\boldsymbol\alpha_p\), \(\mathbb G_a\), or a
  twisted form of \((\mathbb Z/p\mathbb Z)^r\) (\SGARefLink{sga3:xvii:definition:1:4}{1.4}) if \(p>0\).
\item[(iii)] \SGARefTarget{sga3:xvii:theorem:3:5:item:iii}
  As in (ii), with the additional requirement that the composition
  series be central.
\item[(iv)] \SGARefTarget{sga3:xvii:theorem:3:5:item:iv}
  \(G\) has a characteristic composition series
  (Exposé VI\textsubscript{B}) defined over \(k\), whose successive
  quotients are isomorphic to \((\mathbb G_a)^r\) if \(p=0\), and,
  respectively, to \((\boldsymbol\alpha_p)^r\), a twisted form of
  \((\mathbb G_a)^s\), or a twisted form of
  \((\mathbb Z/p\mathbb Z)^t\), taken in this order, if \(p>0\).
\item[(v)] \SGARefTarget{sga3:xvii:theorem:3:5:item:v}
  \(G\) is isomorphic to an algebraic subgroup of the group
  \(\operatorname{Trigstr}(n)_k\) of strictly upper-triangular
  matrices in the linear group \(\operatorname{GL}(n)_k\), for a
  suitable integer \(n\geq0\).
\item[(vi)] \SGARefTarget{sga3:xvii:theorem:3:5:item:vi}
  \(G\) is affine, and for every linear representation of \(G\) on a
  nonzero finite-dimensional \(k\)-vector space \(V\), one has
  \(V^G\ne0\).
\end{enumerate}}
\SGARefTarget{sga3:xvii:theorem:3:5}

\SGAProof{Proof}
The implication \SGARefLink{sga3:xvii:theorem:3:5:item:i}{(i)}\(\Rightarrow\)\SGARefLink{sga3:xvii:theorem:3:5:item:vi}{(vi)} follows from \SGARefLink{sga3:xvii:proposition:2:1}{2.1} and \SGARefLink{sga3:xvii:proposition:3:2}{3.2}.

For \SGARefLink{sga3:xvii:theorem:3:5:item:vi}{(vi)}\(\Rightarrow\)\SGARefLink{sga3:xvii:theorem:3:5:item:v}{(v)}, since the algebraic group \(G\) is
affine, it is an algebraic subgroup of a suitable linear group
\(\operatorname{GL}(V)\) (Exposé VI\textsubscript{B}, \SGARefLink{def:sga3-VIB-11.3}{11.3}).
Applying \SGARefLink{sga3:xvii:corollary:3:4}{3.4} to the representation of \(G\) on \(V\) defined by this
embedding gives \SGARefLink{sga3:xvii:theorem:3:5:item:v}{(v)}.

For \SGARefLink{sga3:xvii:theorem:3:5:item:v}{(v)}\(\Rightarrow\)\SGARefLink{sga3:xvii:theorem:3:5:item:iii}{(iii)}, the algebraic group
\(\operatorname{Trigstr}(n)\) is known to have a central composition
series whose successive quotients are isomorphic to \(\mathbb G_a\).
The composition series induced on \(G\) gives property \SGARefLink{sga3:xvii:theorem:3:5:item:iii}{(iii)}, in view
of \SGARefLink{sga3:xvii:proposition:1:5}{1.5}.

The implications
\[
  \text{(iii)}\Rightarrow\text{(ii)}\Rightarrow\text{(i)}
  \qquad\text{and}\qquad
  \text{(iv)}\Rightarrow\text{(i)}
\]
are clear.

We shall prove \SGARefLink{sga3:xvii:theorem:3:5:item:i}{(i)}\(\Rightarrow\)\SGARefLink{sga3:xvii:theorem:3:5:item:iv}{(iv)} shortly, but first note some
consequences of what has already been proved.

\SGAStatement{Definition 3.6}{A \(p\)-Lie algebra \(\mathfrak g\),
\(p>0\) (cf. Exposé VII\textsubscript{A}, \S5), is called
\emph{unipotent} if the map
\[
  x\longmapsto x^{(p)}
\]
is nilpotent; that is, if for every \(x\in\mathfrak g\) there is an
integer \(n>0\) such that \(x^{(p^n)}=0\).}
\SGARefTarget{sga3:xvii:definition:3:6}

% END INCLUDED SGA3_English_current_buildable_source_20260823/source/inputs/XVII/tex/components/06_expose_XVII_corollary3_4_theorem3_5_definition3_6_page008.tex

% BEGIN INCLUDED SGA3_English_current_buildable_source_20260823/source/inputs/XVII/tex/components/07_expose_XVII_corollaries3_7_3_8_and_lemma3_9_page009.tex
% Authority: Expo17.pdf, local pp.8--9 / combined pp.970--971 /
% printed pp.547--549.

\SGAStatement{Corollary 3.7}{A unipotent algebraic group \(G\) is
nilpotent (Exposé VI\textsubscript{B}, \S8); its Lie algebra
\(\mathfrak g\) is nilpotent (Bourbaki, \emph{Groupes et algèbres de
Lie}, Chap.~I, \S4) and is isomorphic to a Lie algebra of nilpotent
endomorphisms of a finite-dimensional vector space.  In
characteristic \(p>0\), \(\mathfrak g\) is a unipotent \(p\)-Lie
algebra (\SGARefLink{sga3:xvii:definition:3:6}{3.6}).}
\SGARefTarget{sga3:xvii:corollary:3:7}

Since \SGARefLink{sga3:xvii:theorem:3:5:item:i}{(i)} implies (v), it is enough to prove \SGARefLink{sga3:xvii:corollary:3:7}{3.7} when
\(G=\operatorname{Trigstr}(n)\).  We have already used the fact that
\(\operatorname{Trigstr}(n)\) is a nilpotent algebraic group.
Moreover, the Lie algebra \(\mathfrak h\) of
\(\operatorname{Trigstr}(n)\) consists of the upper-triangular
endomorphisms of \(V\) that have zeros on the main diagonal.  These
endomorphisms are therefore nilpotent, and consequently
\(\mathfrak h\) is nilpotent (Bourbaki, \emph{loc.\ cit.},
Chap.~I, \S4, Cor.~3).  If \(p>0\), the \(p\)-th power in the
\(p\)-Lie algebra
\(\mathfrak{gl}(V)=\operatorname{End}(V)\) coincides with the
\(p\)-th power of endomorphisms of \(V\)
(Exposé VII\textsubscript{A}, \SGARefLink{sga3:VIIA:6.4.4}{6.4.4}); hence \(\mathfrak h\) is
unipotent.

\SGAStatement{Corollary 3.8}{Let \(k\) be an algebraically closed
field, and let \(G\) be a smooth affine algebraic group over \(k\).
Then the following properties are equivalent:
\begin{enumerate}
\item[(i)] \SGARefTarget{sga3:xvii:corollary:3:8:item:i}
  \(G\) is unipotent.
\item[(ii)] \SGARefTarget{sga3:xvii:corollary:3:8:item:ii}
  \(G(k)\) consists of unipotent elements
  (BIBLE~4, Prop.~4, Cor.~1); that is, \(G\) is unipotent in the
  sense of BIBLE.
\end{enumerate}}
\SGARefTarget{sga3:xvii:corollary:3:8}

The implication \SGARefLink{sga3:xvii:corollary:3:8:item:i}{(i)}\(\Rightarrow\)\SGARefLink{sga3:xvii:corollary:3:8:item:ii}{(ii)} follows because, by
\SGARefLink{sga3:xvii:proposition:3:2}{3.2}, \SGARefLink{sga3:xvii:corollary:3:8:item:i}{(i)}\(\Rightarrow\)(v), \(G\) is isomorphic to an algebraic
subgroup of a group \(\operatorname{Trigstr}(n)\).

For \SGARefLink{sga3:xvii:corollary:3:8:item:ii}{(ii)}\(\Rightarrow\)\SGARefLink{sga3:xvii:corollary:3:8:item:i}{(i)}, suppose that \(G\) is a unipotent
algebraic group in the sense of BIBLE, and write \(G^0\) for its
identity component.  The maximal tori of \(G^0\), being made up of
unipotent elements, are trivial; thus \(G^0\) is equal to each of its
Cartan subgroups.  It follows that \(G^0\) is solvable
(BIBLE~6, Th.~6), and hence triangularizable
(BIBLE~6, Th.~1).  In short, \(G^0\) is an algebraic subgroup of a
group \(\operatorname{Trigstr}(n)\), and is therefore unipotent in
the sense of this Exposé.

The group \((G/G^0)(k)\) is a finite group consisting of unipotent
elements.  It is therefore trivial if \(p=0\), and is a finite
\(p\)-group if \(p>0\) (BIBLE~4, Prop.~4).  But then \(G/G^0\) is
unipotent in the sense of this Exposé, since it is an iterated
extension of groups isomorphic to \(\mathbb Z/p\mathbb Z\).  This
proves that \(G\) is unipotent.

\SGAProof{Completion of the proof of \SGARefLink{sga3:xvii:theorem:3-5}{3.5}}
It remains to prove \SGARefLink{sga3:xvii:theorem:3:5:item:i}{(i)}\(\Rightarrow\)\SGARefLink{sga3:xvii:theorem:3:5:item:iv}{(iv)}.

\smallskip
\noindent\textit{\SGARefLink{sga3:xvii:appendix:A}{(a)} \(p>0\).}\quad
Consider the increasing sequence of algebraic subgroups of \(G\)
\[
  \{e\}\subset F(G)\subset F^2(G)\subset\cdots\subset
  F^n(G)\subset G^0\subset G.
\]
This gives a characteristic composition series of \(G\)
(Appendix~II, 1), and for \(n\) sufficiently large,
\(G/F^n(G)\) is smooth (Appendix~II, \SGARefLink{sga3:xvii:lemma:3-1}{3.1}).  Thus its successive
quotients are, in order:
\begin{enumerate}
\item[(1)] radicial groups of height \(1\);
\item[(2)] a smooth connected group;
\item[(3)] an étale group.
\end{enumerate}
To prove \SGARefLink{sga3:xvii:lemma:3-1:item:i}{(i)}\(\Rightarrow\)(iv), it is therefore enough to prove the
following lemma.

\SGAStatement{\SGAInlineRefTarget{sga3:xvii:lemma:3-9}Lemma 3.9}{Let \(G\) be a unipotent algebraic group
defined over a field \(k\) of characteristic \(p>0\).  Then \(G\)
has a characteristic composition series, defined over \(k\), whose
successive quotients are isomorphic to:
\begin{enumerate}
\item[(i)] \SGARefTarget{sga3:xvii:lemma:3:9:item:i}
  \((\boldsymbol\alpha_p)^r\), if \(G\) is radicial;
\item[(ii)] \SGARefTarget{sga3:xvii:lemma:3:9:item:ii}
  a twisted form of \((\mathbb G_a)^r\), if \(G\) is smooth and
  connected;
\item[(iii)] \SGARefTarget{sga3:xvii:lemma:3:9:item:iii}
  a twisted form of \((\mathbb Z/p\mathbb Z)^r\), if \(G\) is
  étale.
\end{enumerate}}
\SGARefTarget{sga3:xvii:lemma:3:9}

% END INCLUDED SGA3_English_current_buildable_source_20260823/source/inputs/XVII/tex/components/07_expose_XVII_corollaries3_7_3_8_and_lemma3_9_page009.tex

% BEGIN INCLUDED SGA3_English_current_buildable_source_20260823/source/inputs/XVII/tex/components/08_expose_XVII_lemma3_9_proof_page010.tex
% Authority: Expo17.pdf, local p.10 / combined p.972 /
% printed pp.550--551.

\SGAProof{Proof of \SGARefLink{sga3:xvii:lemma:3:9:item:i}{3.9(i)}}
The group \(G\) is radicial.  Filtering \(G\) by the \(F^n(G)\), we
reduce to the case in which \(G\) is radicial of height \(1\).
Since \(G\) is nilpotent (\SGARefLink{sga3:xvii:theorem:3-5}{3.5}, \SGARefLink{sga3:xvii:theorem:3:5:item:i}{(i)}\(\Rightarrow\)\SGARefLink{sga3:xvii:theorem:3:5:item:iii}{(iii)}), and since
the center of an algebraic group is representable
(\SGARefLink{sga3:viii:examples:6:5:item:e:01}{Exposé VIII, 6.5(e)}), we may consider the ascending central series
of \(G\).  This series is plainly characteristic in \(G\), and
reduces us to the case in which \(G\) is also commutative.

Put \(\mathfrak g=\operatorname{Lie}(G)\).  The \(p\)-th-power
morphism \(\pi\) is then additive on \(\mathfrak g\)
(Exposé VII\textsubscript{A}); we shall reduce to the case in which
it is zero.  For every prescheme \(S\) over \(k\), set
\[
  \mathfrak g_S=\mathfrak g\otimes_k\mathcal O_S,
\]
and let \(\mathfrak h_S\) be the subsheaf of abelian groups of
\(\mathfrak g_S\) that is the image of \(\mathfrak g_S\) under
\(\pi_S\).  Finally, let \(\overline{\mathfrak h}_S\) be the
subsheaf of \(\mathcal O_S\)-modules of \(\mathfrak g_S\) generated
by \(\mathfrak h_S\).  It is clear that
\[
  \overline{\mathfrak h}_S
  =\overline{\mathfrak h}_k\otimes_k\mathcal O_S
\]
and that \(\overline{\mathfrak h}_S\)\footnote[4]{Editor's note:
\(\mathfrak h\) has been changed to \(\overline{\mathfrak h}\).}
is a characteristic \(p\)-Lie subalgebra of \(\mathfrak g_S\);
that is, it is stable under the \(S\)-functor
\(\operatorname{Aut}_{p\text{-}\mathrm{Lie}}(\mathfrak g)\).
It follows from Appendix~II, \SGARefLink{sga3:ii:subsection:2-2}{2.2} that
\(\overline{\mathfrak h}_k\) is the Lie algebra of an algebraic
subgroup \(H\) of \(G\) that is characteristic in \(G\).

Moreover, in view of \SGARefLink{sga3:xvii:corollary:3:7}{3.7} and \SGARefLink{sga3:xvii:lemma:3:9:bis}{Lemma 3.9 bis} below, if
\(G\ne\{e\}\), then \(H\ne G\), since
\(\overline{\mathfrak h}_k\ne\mathfrak g\).  On the other hand, if
\(G/H=G''\), then
\[
  \operatorname{Lie}G''
  =\mathfrak g/\overline{\mathfrak h}_k
  \qquad\text{(Appendix~II, 2.2),}
\]
and consequently the \(p\)-th power is zero in
\(\operatorname{Lie}G''\).  Induction on
\(\dim\operatorname{Lie}G\) therefore reduces us to the case in
which \(\operatorname{Lie}G\) is a \(p\)-Lie algebra whose
\(p\)-th power is zero.  But then \(\operatorname{Lie}G''\) is
isomorphic to \(\operatorname{Lie}((\boldsymbol\alpha_p)^r)\) for
some suitable integer \(r\geq0\) (Appendix~II, \SGARefLink{sga3:xvii:proposition:2:1}{2.1}), and hence
\(G''\) is isomorphic to \((\boldsymbol\alpha_p)^r\)
(Appendix~II, \SGARefLink{sga3:ii:subsection:2-2}{2.2}).  It remains to prove the following lemma.

\SGAStatement{Lemma 3.9 bis}{Let \(k\) be a field of characteristic
\(p>0\), let \(\mathfrak g\) be a commutative unipotent
\(p\)-Lie algebra (\SGARefLink{sga3:xvii:definition:3:6}{3.6}) of finite dimension over \(k\), and let
\(\mathfrak h\) be the \(p\)-Lie subalgebra of \(\mathfrak g\)
generated by the image of the \(p\)-th power in \(\mathfrak g\).
If \(\mathfrak g\ne0\), then \(\mathfrak h\ne\mathfrak g\).}
\SGARefTarget{sga3:xvii:lemma:3:9:bis}

Indeed, since \(\mathfrak g\) is commutative, \(\mathfrak h\) is
simply the \(k\)-vector subspace of \(\mathfrak g\) generated by the
\(X^{(p)}\), \(X\in\mathfrak g\).  If \(\mathfrak g\ne0\) and is
unipotent, there exists an \(X\in\mathfrak g\), \(X\ne0\), such
that \(X^{(p)}=0\).  Let \(X_1,\ldots,X_n\) be a basis of a
complement in \(\mathfrak g\) to the line \(kX\).  The Lie algebra
\(\mathfrak h\) is then the \(k\)-vector subspace of
\(\mathfrak g\) generated by
\[
  X_1^{(p)},\ldots,X_n^{(p)};
\]
it therefore has dimension at most
\(n=\dim\mathfrak g-1\).

\SGAProof{Proof of \SGARefLink{sga3:xvii:lemma:3:9:item:ii}{3.9(ii)}}
The group \(G\) is smooth and connected.  In this case the
descending central series of \(G\) is represented by smooth,
connected, characteristic algebraic subgroups \(G_i\)
(Exposé VI\textsubscript{B}, \SGARefLink{prop:sga3-VIB-8.3}{8.3} and \SGARefLink{prop:sga3-VIB-7.4}{7.4}), and \(G_i=0\) for
sufficiently large \(i\), since \(G\) is nilpotent
(\SGARefLink{sga3:xvii:theorem:3-5}{3.5}, \SGARefLink{sga3:xvii:theorem:3:5:item:i}{(i)}\(\Rightarrow\)\SGARefLink{sga3:xvii:theorem:3:5:item:iii}{(iii)}).  It is enough to prove \SGARefLink{sga3:xvii:lemma:3-9}{3.9} for the
groups \(G_i/G_{i+1}\), which reduces us to the case in which \(G\)
is also commutative.

For every integer \(n>0\), let \(G_n\) be the algebraic subgroup of
\(G\) that is the image of \(G\) under the \(p^n\)-th-power
morphism.  The group \(G_n\) is therefore smooth, connected, and
characteristic; and it follows from \SGARefLink{sga3:xvii:definition:1:1}{Definition 1.1} of unipotent
groups that \(G_n=0\) for sufficiently large \(n\).  Replacing \(G\)
by \(G_n/G_{n+1}\), we may further suppose that \(G\) is annihilated
by raising to the \(p\)-th power.  But then, by J.-P. Serre,
\emph{Groupes algébriques et corps de classes}, Chap.~VII,
Prop.~11, \(G\) is a form of \((\mathbb G_a)^r\) for a suitable
integer \(r\).

% END INCLUDED SGA3_English_current_buildable_source_20260823/source/inputs/XVII/tex/components/08_expose_XVII_lemma3_9_proof_page010.tex

% BEGIN INCLUDED SGA3_English_current_buildable_source_20260823/source/inputs/XVII/tex/components/09_expose_XVII_lemma3_9_iii_and_3_9_ter_page011.tex
% Authority: Expo17.pdf, local p.11 / combined p.973 /
% printed pp.552--553, stopping before Proposition 3.10.

\SGAProof{Proof of \SGARefLink{sga3:xvii:lemma:3:9:item:iii}{3.9(iii)}}
The group \(G\) is étale.  Proceeding as in (ii), we reduce first to
the case in which \(G\) is commutative, and then to the case in which
\(G\) is annihilated by \(p\).  But then
\[
  G_{\overline{k}}\simeq(\mathbb Z/p\mathbb Z)^r.
\]

\SGAProof{Proof of \SGARefLink{sga3:xvii:theorem:3-5}{3.5}, \SGARefLink{sga3:xvii:theorem:3:5:item:i}{(i)}\(\Rightarrow\)\SGARefLink{sga3:xvii:theorem:3:5:item:iv}{(iv)}, in case \SGARefLink{sga3:xvii:section:b}{(b)},
\(p=0\)}
The group \(G\) is then smooth and connected; proceeding as in
\SGARefLink{sga3:xvii:lemma:3:9:item:ii}{3.9(ii)}, we reduce to the case in which \(G\) is also commutative.
We then have the following more precise result.

\SGAStatement{Lemma 3.9 ter}{Let \(k\) be a field of characteristic
\(0\), let \(G\) be a commutative unipotent \(k\)-algebraic group,
and put \(\mathfrak g=\operatorname{Lie}G\).  There is a canonical
isomorphism
\[
  \exp:\mathbf W(\mathfrak g)\xrightarrow{\ \sim\ }G.
\]
The morphism \(\exp\) is the unique homomorphism
\(\mathbf W(\mathfrak g)\to G\) that induces the identity on Lie
algebras.}
\SGARefTarget{sga3:xvii:lemma:3:9:ter}

Since \(G\) is unipotent, it can be realized as an algebraic
subgroup of \(\operatorname{Trigstr}(n)\), for a suitable integer
\(n\) (\SGARefLink{sga3:xvii:theorem:3-5}{3.5}, \SGARefLink{sga3:xvii:theorem:3:5:item:i}{(i)}\(\Rightarrow\)\SGARefLink{sga3:xvii:theorem:3:5:item:v}{(v)}).  Choosing such an embedding
identifies \(\mathfrak g\) with a Lie subalgebra of
\(\mathfrak{gl}(n)\) consisting of nilpotent endomorphisms.  We
thereby obtain a \(k\)-morphism
\[
  \begin{aligned}
  \exp:\mathbf W(\mathfrak g)&\longrightarrow
    \operatorname{GL}(n),\\
  T&\longmapsto\sum_{i\geq0}\frac{T^i}{i!}.
  \end{aligned}
\]

Since \(G\) is commutative, the morphism \(\exp\) is a
homomorphism.  Let \(G'\) be the algebraic group that is the image
of \(\mathbf W(\mathfrak g)\) under \(\exp\).  Under the canonical
identification
\(\operatorname{Lie}\mathbf W(\mathfrak g)=\mathfrak g\), the
tangent linear map of \(\exp\) is simply the inclusion
\(\mathfrak g\hookrightarrow\mathfrak{gl}(n)\).  Hence
\[
  \operatorname{Lie}(G\cap G')
  =\mathfrak g\cap\mathfrak g=\mathfrak g.
\]
Since \(G\cap G'\) is smooth (Exposé VI\textsubscript{B}, \SGARefLink{cor:sga3-VIB-1.6.1}{1.6.1})
and \(G\) is connected---being an iterated extension of groups
\(\mathbb G_a\) (\SGARefLink{sga3:xvii:theorem:3-5}{3.5}, \SGARefLink{sga3:xvii:theorem:3:5:item:i}{(i)}\(\Rightarrow\)\SGARefLink{sga3:xvii:theorem:3:5:item:ii}{(ii)})---one necessarily has
\(G=G'\).  The kernel \(\ker(\exp)\) is an étale unipotent group,
and hence is the trivial group.  Consequently \(\exp\) is an
isomorphism from \(\mathbf W(\mathfrak g)\) onto \(G\).

If \(h:\mathbf W(\mathfrak g)\to G\) is another homomorphism such
that \(\operatorname{Lie}(h)\) is the identity map of
\(\mathfrak g\), then \(h-\exp\) is a homomorphism, since \(G\) is
commutative, and its tangent linear map is zero.  Because \(k\) has
characteristic \(0\) and \(\mathbf W(\mathfrak g)\) is connected,
Cartier's theorem again implies that \(h=\exp\).

% END INCLUDED SGA3_English_current_buildable_source_20260823/source/inputs/XVII/tex/components/09_expose_XVII_lemma3_9_iii_and_3_9_ter_page011.tex

% BEGIN INCLUDED SGA3_English_current_buildable_source_20260823/source/inputs/XVII/tex/components/10_expose_XVII_proposition3_10_through_section4_opening_page012.tex
% Authority: Expo17.pdf, local pp.11--12 / combined pp.973--974 /
% printed pp.553--556, stopping before subsection 4.1.

\SGAStatement{\SGAInlineRefTarget{sga3:xvii:proposition:3-10}Proposition 3.10}{Let \(G\) be an algebraic group
defined over an algebraically closed field \(k\).  The following
properties are equivalent:
\begin{enumerate}
\item[(i)] \SGARefTarget{sga3:xvii:proposition:3:10:item:i}
  Every homomorphism from a group \(M\) of multiplicative type into
  \(G\) is trivial.
\item[(i bis)] \SGARefTarget{sga3:xvii:proposition:3:10:item:i:bis}
  \(G\) has no nontrivial subgroup of multiplicative type.
\item[(ii)] \SGARefTarget{sga3:xvii:proposition:3:10:item:ii}
  \begin{enumerate}
  \item[(a)] \SGARefTarget{sga3:xvii:proposition:3:10:item:ii:a}
    If \(p=0\), then \(G(k)\) has no point of finite order other
    than \(e\).
  \item[(b)] \SGARefTarget{sga3:xvii:proposition:3:10:item:ii:b}
    If \(p\ne0\), then:
    \begin{itemize}
    \item for every prime number \(q\ne p\), if \(x\in G(k)\) and
      \(x^q=e\), then \(x=e\);
    \item \SGARefTarget{sga3:xvii:proposition:3:10:item:ii:b:lie}
      for every \(X\in\mathfrak g=\operatorname{Lie}G\) such that
      \(X^{(p)}=X\), one has \(X=0\).
    \end{itemize}
  \end{enumerate}
\item[(ii bis)] \SGARefTarget{sga3:xvii:proposition:3:10:item:ii:bis}
  \begin{enumerate}
  \item[(a)] \SGARefTarget{sga3:xvii:proposition:3:10:item:ii:bis:a}
    If \(p=0\), the condition is the same as in \SGARefLink{sga3:xvii:proposition:3:10:item:ii}{(ii)}\SGARefLink{sga3:xvii:appendix:A}{(a)}.
  \item[(b)] \SGARefTarget{sga3:xvii:proposition:3:10:item:ii:bis:b}
    If \(p\ne0\), then:
    \begin{itemize}
    \item every finite subgroup of \(G(k)\) is a \(p\)-group;
    \item \SGARefTarget{sga3:xvii:proposition:3:10:item:ii:bis:b:lie}
      for every \(X\in\mathfrak g=\operatorname{Lie}G\) such that
      \(X^{(p)}=X\), one has \(X=0\).
    \end{itemize}
  \end{enumerate}
\end{enumerate}}
\SGARefTarget{sga3:xvii:proposition:3:10}

\SGAProof{Proof}
The equivalence \SGARefLink{sga3:xvii:proposition:3:10:item:i}{(i)}\(\Leftrightarrow\)(i bis) follows because the
image of a group of multiplicative type is of multiplicative type
(\SGARefLink{sga3:ix:proposition:2:7}{Exposé IX, 2.7}).

The equivalence \SGARefLink{sga3:ix:proposition:2:7:item:ii:01}{(ii)}\(\Leftrightarrow\)(ii bis) follows from the
following observation.  If an ordinary finite group \(H\) has order
that is not a power of \(p\), then there exist a prime number
\(q\ne p\) and an element \(x\in H\), distinct from \(e\), such
that \(x^q=e\) (Sylow's theorem; cf. J.-P. Serre,
\emph{Corps locaux}, Chapter IX, \S2).

The implication \SGARefLink{sga3:xvii:proposition:3:10:item:i}{(i)}\(\Rightarrow\)\SGARefLink{sga3:xvii:proposition:3:10:item:ii}{(ii)} follows from the next lemma.

\SGAStatement{\SGAInlineRefTarget{sga3:xvii:lemma:3-11}Lemma 3.11}{Let \(G\) be an algebraic group defined
over a field \(k\).
\begin{enumerate}
\item[(a)] \SGARefTarget{sga3:xvii:lemma:3:11:item:a}
  If \(k\) contains the \(n\)-th roots of unity, where \(n\) is
  prime to \(p\), then
  \[
    \operatorname{Hom}_{k\text{-}\mathrm{gr}}(\mu_n,G)
      \simeq
    \operatorname{Hom}_{k\text{-}\mathrm{gr}}
      (\mathbb Z/n\mathbb Z,G)
      \simeq {}_nG(k),
  \]
  where \({}_nG(k)\) denotes the points of order \(n\) in \(G(k)\).
\item[(b)] \SGARefTarget{sga3:xvii:lemma:3:11:item:b}
  If \(p>0\), then
  \[
    \operatorname{Hom}_{k\text{-}\mathrm{gr}}(\mu_p,G)
      \simeq
    \{X\in\mathfrak g=\operatorname{Lie}G\mid X^{(p)}=X\}.
  \]
\end{enumerate}}
\SGARefTarget{sga3:xvii:lemma:3:11}

Indeed, for \SGARefLink{sga3:xvii:appendix:A}{(a)}, note that \(\mu_n\) is then isomorphic over \(k\)
to \(\mathbb Z/n\mathbb Z\); assertion \SGARefLink{sga3:xvii:section:b}{(b)} follows from Appendix II,
\SGARefLink{sga3:xvii:proposition:2:1}{2.1}.

For \SGARefLink{sga3:xvii:proposition:2:1:item:ii}{(ii)}\(\Rightarrow\)(i bis), \SGARefLink{sga3:xvii:lemma:3-11}{Lemma 3.11} shows that \SGARefLink{sga3:xvii:proposition:2:1:item:ii}{(ii)} is
equivalent to the assertion that, for every prime number \(r\), the
group \(G\) contains no subgroup \(\mu_r\).  The implication now
follows from:

\SGAStatement{Lemma 3.12}{Let \(G\) be a diagonalizable
\(S\)-group, of finite type over \(S\), and distinct from the
trivial group.  There exist a prime number \(r\) and a subgroup of
\(G\) isomorphic to \((\mu_r)_S\).}
\SGARefTarget{sga3:xvii:lemma:3:12}

Write \(G=D_S(M)\), where \(M\) is a finitely generated abelian
group, hence an extension of a free group \(M''\) by a finite group
\(M'\).  If \(M''\ne0\), it is clear that \(M\) has quotients
isomorphic to \(\mathbb Z/r\mathbb Z\) for every integer \(r\).  If
\(M''=0\), then \(M'\) has a quotient isomorphic to
\(\mathbb Z/r\mathbb Z\) for every prime number \(r\) dividing the
order of \(M\).  The lemma follows by duality.

We saw in \SGARefLink{sga3:xvii:proposition:2:4}{Proposition 2.4} that a unipotent group satisfies the
equivalent conditions of \SGARefLink{sga3:xvii:proposition:3-10}{3.10}.  The purpose of the next section is
to prove the converse.

\SGASection{4}{A Characterization of Unipotent Groups}
\SGARefTarget{sga3:xvii:section:4}

As announced, we shall prove that an algebraic group \(G\) defined
over an algebraically closed field \(k\), and containing no
nontrivial subgroup of multiplicative type, is unipotent.  In fact,
it is enough that \(G\) contain none of certain particular
``elementary'' subgroups of multiplicative type, which depend on the
hypotheses imposed on \(G\).  Before stating the general theorem,
let us study some special cases in detail.

% END INCLUDED SGA3_English_current_buildable_source_20260823/source/inputs/XVII/tex/components/10_expose_XVII_proposition3_10_through_section4_opening_page012.tex

% BEGIN INCLUDED SGA3_English_current_buildable_source_20260823/source/inputs/XVII/tex/components/11_expose_XVII_proposition4_1_1_and_lemma4_1_2_page013.tex
% Authority: Expo17.pdf, local p.13 and local p.14 top /
% combined pp.975--976 / printed pp.557--558.

\SGASubsection{4.1}{Smooth, Connected, Affine Algebraic Groups}
\SGARefTarget{sga3:xvii:subsection:4:1}

\SGAStatement{\SGAInlineRefTarget{sga3:xvii:proposition:4-1-1}Proposition 4.1.1}{Let \(k\) be a field, let \(G\)
be a smooth, connected, affine \(k\)-algebraic group, and put
\(\mathfrak g=\operatorname{Lie}G\).  The following properties are
equivalent:
\begin{enumerate}
\item[(i)] \SGARefTarget{sga3:xvii:proposition:4:1:1:item:i}
  \(G\) is unipotent.
\item[(ii)] \SGARefTarget{sga3:xvii:proposition:4:1:1:item:ii}
  \(G\) has a central composition series whose successive quotients
  are forms of \(\mathbb G_a\).
\item[(iii)] \SGARefTarget{sga3:xvii:proposition:4:1:1:item:iii}
  \(G\) has a central, characteristic composition series whose
  successive quotients are forms of \((\mathbb G_a)^r\).
\item[(iv)] \SGARefTarget{sga3:xvii:proposition:4:1:1:item:iv}
  There is an integer \(n>1\) such that \(G_{\overline{k}}\)
  contains no subgroup isomorphic to \(\mu_n\).
\item[(v)] \SGARefTarget{sga3:xvii:proposition:4:1:1:item:v}
  Every maximal torus of \(G\) is the trivial group.
\end{enumerate}
Suppose, moreover, that \(G\) is an algebraic subgroup of a linear
group \(\operatorname{GL}(n)\).  The preceding conditions are then
also equivalent to:
\begin{enumerate}
\item[(vi)] \SGARefTarget{sga3:xvii:proposition:4:1:1:item:vi}
  \(\mathfrak g\subset\mathfrak{gl}(n)\) consists of nilpotent
  endomorphisms.
\item[(vii)] \SGARefTarget{sga3:xvii:proposition:4:1:1:item:vii}
  \(\mathfrak g\) is nilpotent, and its center contains no nonzero
  semisimple endomorphism.
\end{enumerate}}
\SGARefTarget{sga3:xvii:proposition:4:1:1}

\SGAProof{Proof}
The implication \SGARefLink{sga3:xvii:proposition:4:1:1:item:ii}{(ii)}\(\Rightarrow\)\SGARefLink{sga3:xvii:proposition:4:1:1:item:i}{(i)} is clear, and
\SGARefLink{sga3:xvii:proposition:4:1:1:item:i}{(i)}\(\Rightarrow\)\SGARefLink{sga3:xvii:proposition:4:1:1:item:iii}{(iii)} was proved in \SGARefLink{sga3:xvii:lemma:3-9}{3.9}.  The implication
\SGARefLink{sga3:xvii:lemma:3:9:item:iii}{(iii)}\(\Rightarrow\)\SGARefLink{sga3:xvii:lemma:3:9:item:ii}{(ii)} will follow from the next lemma.

\SGAStatement{Lemma 4.1.2}{Let \(k\) be a field, and let \(G\) be
a \(k\)-algebraic group that is a form of \((\mathbb G_a)^r\).
Then:
\begin{enumerate}
\item[(a)] \SGARefTarget{sga3:xvii:lemma:4:1:2:item:a}
  \(G\) can be realized as an algebraic subgroup of
  \((\mathbb G_a)^n\), for a suitable integer \(n\).
\item[(b)] \SGARefTarget{sga3:xvii:lemma:4:1:2:item:b}
  \(G\) has a composition series whose successive quotients are
  forms of \(\mathbb G_a\).
\end{enumerate}}
\SGARefTarget{sga3:xvii:lemma:4:1:2}

Indeed, by hypothesis there is an extension \(k'\) of \(k\) such
that \(G_{k'}\) is isomorphic to \((\mathbb G_{a,k'})^r\).  By the
principle of finite extension (EGA IV, 9.1.1), we may assume that
\(k'/k\) is finite.  For (a), it then suffices to consider the
canonical closed immersion (EGA V\footnote[5]{Editor's note:
provide another reference here\ldots})
\[
  G\longrightarrow
  \prod_{k'/k}(G_{k'})_{/k'}
  \xrightarrow{\ \sim\ }(\mathbb G_{a,k})^n,
  \qquad n=r\deg(k'/k).
\]

To prove \SGARefLink{sga3:xvii:lemma:4:1:2:item:b}{(b)}, in view of \SGARefLink{sga3:xvii:lemma:4:1:2:item:a}{(a)}, we may assume that \(G\) is a closed
subgroup of \(G'=(\mathbb G_a)^n\).  If \(G\ne0\), there is a
hyperplane \(\mathfrak h\) in
\(\mathfrak g'=\operatorname{Lie}G'\) that does not contain
\(\mathfrak g\).  Let \(H\) be the vector subgroup
\(\mathbf W(\mathfrak h)\) of
\(\mathbf W(\mathfrak g')=G'\).  Since \(H\) is defined by one
equation in \(G'\), the intersection \(G\cap H\) is defined by one
equation in \(G\), and
\[
\begin{aligned}
  \dim G-1
  &\leq \dim(G\cap H)\\
  &\leq \dim\operatorname{Lie}(G\cap H)
   =\dim_k(\mathfrak g\cap\mathfrak h)\\
  &=\dim_k\mathfrak g-1
   =\dim G-1.
\end{aligned}
\]
Thus
\(\dim(G\cap H)=\dim\operatorname{Lie}(G\cap H)\), and consequently
\(G\cap H\) is smooth.  The group
\[
  G_1=(G\cap H)^0
\]
is a smooth, connected algebraic subgroup of \(G\) such that
\(G/G_1\) is smooth, connected, and one-dimensional, hence is a
form of \(\mathbb G_a\) (\SGARefLink{sga3:xvii:subsection:4:1}{4.1}, \SGARefLink{sga3:xvii:proposition:4:1:1:item:i}{(i)}\(\Rightarrow\)\SGARefLink{sga3:xvii:proposition:4:1:1:item:iii}{(iii)}).  Induction
on the dimension of \(G\) completes the proof.

% END INCLUDED SGA3_English_current_buildable_source_20260823/source/inputs/XVII/tex/components/11_expose_XVII_proposition4_1_1_and_lemma4_1_2_page013.tex

% BEGIN INCLUDED SGA3_English_current_buildable_source_20260823/source/inputs/XVII/tex/components/12_expose_XVII_corollary4_1_3_through_lemma4_1_5_page014.tex
% Authority: Expo17.pdf, local p.14 / combined p.976 /
% printed pp.558--560, stopping before the proof of Lemma 4.1.5.

Before continuing the proof of \SGARefLink{sga3:xvii:subsection:4:1}{4.1}, note that the equivalence
(i)\(\Leftrightarrow\)(ii), together with \SGARefLink{sga3:xvii:lemma:2:3:bis}{2.3 bis}, gives the
following corollary.

\SGAStatement{Corollary 4.1.3}{If \(k\) is a perfect field, a
smooth, connected \(k\)-algebraic group is unipotent if and only if
it has a composition series whose successive quotients are
isomorphic to \(\mathbb G_a\).}
\SGARefTarget{sga3:xvii:corollary:4:1:3}

\SGAProof{Continuation of the proof of \SGARefLink{sga3:xvii:subsection:4:1}{4.1}}
The implication (i)\(\Rightarrow\)(iv) follows from \SGARefLink{sga3:xiv:2.4:item:i}{2.4(i)}.

For (iv)\(\Rightarrow\)(v), by \SGARefLink{sga3:xiv:4.1}{Exposé XIV, 4.1}, the group \(G\)
has a maximal torus \(T\) defined over \(k\).  If \(r=\dim T\),
then
\[
  ({}_nT)_{\overline{k}}\simeq(\mu_n)^r.
\]
It follows that \(r=0\).

The implication (v)\(\Rightarrow\)(i) follows from the observation
made in the proof of \SGARefLink{sga3:xvii:corollary:3:8}{3.8}.

For (i)\(\Rightarrow\)(vi), by \SGARefLink{sga3:xvii:corollary:3:4}{3.4} the group \(G\) is in fact
contained in an algebraic subgroup of \(\operatorname{GL}(n)\)
isomorphic to \(\operatorname{Trigstr}(n)\).  Hence
\(\mathfrak g\) consists of nilpotent endomorphisms.

For (vi)\(\Rightarrow\)(vii), the Lie algebra \(\mathfrak g\) is
nilpotent by Engel's theorem (Bourbaki, \emph{Groupes et algèbres
de Lie}, Chapter I, \S4, Corollary 3).

For (vii)\(\Rightarrow\)(v), let \(T\) be a maximal subtorus of
\(G\) (\SGARefLink{sga3:xiv:1.1}{Exposé XIV, 1.1}), and let \(\mathfrak t\) be its Lie
algebra.  The embedding of \(G\) in \(\operatorname{GL}(n)\)
defines a representation of \(T\), which is necessarily
semisimple; this can be seen after passing to an algebraic closure
\(\overline{k}\) of \(k\) and applying \SGARefLink{sga3:i:proposition:4-7-3}{Exposé I, 4.7.3}.  Thus, if
\(X\in\mathfrak t\), then \(X\) is a semisimple endomorphism in
\(\mathfrak{gl}(n)\).  It follows at once that
\[
  \operatorname{ad}X:\ Y\longmapsto[X,Y]
\]
is a semisimple endomorphism of \(\mathfrak{gl}(n)\), and hence of
\(\mathfrak g\).  This endomorphism is also nilpotent, since
\(\mathfrak g\) is nilpotent.  Therefore
\(\operatorname{ad}X=0\), so \(X\) is central.  Consequently
\(\mathfrak t\) is central and consists of semisimple
endomorphisms; by hypothesis it is zero.  A fortiori, \(T\) is the
trivial group.

\SGAStatement{Remark 4.1.4}{%
\begin{enumerate}
\item[(a)] \SGARefTarget{sga3:xvii:remark:4:1:4:item:a}
  We shall give below (\SGARefLink{sga3:xvii:proposition:4-3-1}{4.3.1}) an infinitesimal characterization of
  unipotent groups in characteristic \(p>0\) that is independent of
  an embedding in \(\operatorname{GL}(n)\).
\item[(b)] \SGARefTarget{sga3:xvii:remark:4:1:4:item:b}
  When \(k\) is perfect, conditions (ii) and (iii) of \SGARefLink{sga3:xvii:proposition:4-1-1}{4.1.1}
  simplify by virtue of the following lemma.
\end{enumerate}}
\SGARefTarget{sga3:xvii:remark:4:1:4}

\SGAStatement{Lemma 4.1.5}{If \(k\) is a perfect field, every
\(k\)-algebraic group \(G\) that is a form of
\((\mathbb G_a)^r\) is isomorphic to \((\mathbb G_a)^r\).}
\SGARefTarget{sga3:xvii:lemma:4:1:5}

% END INCLUDED SGA3_English_current_buildable_source_20260823/source/inputs/XVII/tex/components/12_expose_XVII_corollary4_1_3_through_lemma4_1_5_page014.tex

% BEGIN INCLUDED SGA3_English_current_buildable_source_20260823/source/inputs/XVII/tex/components/13_expose_XVII_lemma4_1_5_proof_page015.tex
% Authority: Expo17.pdf, local pp.14--15 / combined pp.976--977 /
% printed pp.560--561, stopping before the two-case degree argument.

The lemma follows from \SGARefLink{sga3:xvii:lemma:3:9:ter}{3.9 ter} when the characteristic \(p\) of
\(k\) is zero, and from \SGARefLink{sga3:xvii:lemma:2:3:bis}{2.3 bis} when \(r=1\).  In the general case
\((p>0)\), realize \(G\) as an algebraic subgroup of
\((\mathbb G_a)^n\), for a suitable integer \(n\) (\SGARefLink{sga3:xvii:lemma:4:1:2}{4.1.2}), and
argue by induction on \(n-r\).  If \(r=n\), then indeed
\(G=(\mathbb G_a)^r\).  Otherwise, the quotient
\((\mathbb G_a)^n/G\) is a nonzero smooth, connected, unipotent
group.  By \SGARefLink{sga3:xvii:proposition:4-1-1}{4.1.1}, \SGARefLink{sga3:xvii:proposition:4:1:1:item:i}{(i)}\(\Rightarrow\)\SGARefLink{sga3:xvii:proposition:4:1:1:item:ii}{(ii)}, and \SGARefLink{sga3:xvii:lemma:2:3:bis}{2.3 bis}, it has a
composition series whose quotients are isomorphic to
\(\mathbb G_a\).  It follows that there is a smooth, connected
algebraic subgroup \(G_1\) of \((\mathbb G_a)^n\), containing \(G\),
such that
\[
  (\mathbb G_a)^n/G_1=\mathbb G_a.
\]
By induction, it suffices to show that \(G_1\) is isomorphic to
\((\mathbb G_a)^{n-1}\).

It is immediate that a homomorphism
\[
  \operatorname{Spec}k[X_1,\ldots,X_n]=(\mathbb G_a)^n
  \longrightarrow
  \mathbb G_a=\operatorname{Spec}k[T]
\]
is defined by an additive polynomial of the form
\[
  \sum_{i,j}a_{i,j}X_i^{p^j}.
\]
Since \(G_1\) is smooth, the linear part of this polynomial is
nonzero.  After a linear change of the coordinates \(X_i\), we may
assume that \(G_1\) is the algebraic subgroup of
\((\mathbb G_a)^n\) defined by
\[
  \tag{\(*\)}
  P(X)=-X_1+\sum_{j=1}^{q}a_jX_1^{p^j}
    +Q(X_2,\ldots,X_n)=0,
\]
\SGARefTarget{sga3:xvii:lemma:4:1:5:eq:star}
where
\[
  Q(X_2,\ldots,X_n)
    =\sum_{\substack{i>1\\j>0}}b_{i,j}X_i^{p^j}.
\]

We now argue by induction on the degree of \(P\).  If
\(\deg P=1\), it is clear that
\[
  G_1\xrightarrow{\ \sim\ }(\mathbb G_a)^{n-1}.
\]
Otherwise, since \(k\) is perfect, one has
\(Q(X)=Q_1(X)^p\), and we may define an endomorphism \(v\) of
\((\mathbb G_a)^n\) by
\[
\begin{aligned}
  X_i&\longmapsto X_i &&\text{for \(i>1\)},\\
  X_1&\longmapsto
    \sum_{j=1}^{q}a_j^{1/p}X_1^{p^{j-1}}+Q_1(X).
\end{aligned}
\]
\SGARefTarget{sga3:xvii:lemma:4:1:5:map:v}
It is clear that \(v\) induces an isomorphism on \(G_1\), and that
\(v(G_1)\) is defined in \((\mathbb G_a)^n\) by the equation
\[
  \tag{\(*\)\(_1\)}
  P_1(X)=-X_1+\sum_{j=1}^{q}a_j^{1/p}X_1^{p^j}
    +Q_1(X_2,\ldots,X_n)=0.
\]
\SGARefTarget{sga3:xvii:lemma:4:1:5:eq:star:1}

% END INCLUDED SGA3_English_current_buildable_source_20260823/source/inputs/XVII/tex/components/13_expose_XVII_lemma4_1_5_proof_page015.tex

% BEGIN INCLUDED SGA3_English_current_buildable_source_20260823/source/inputs/XVII/tex/components/14_expose_XVII_lemma4_1_5_completion_pages015_016.tex
% Authority: Expo17.pdf, local pp.15--16 / combined pp.977--978 /
% printed pp.561--562, completing the proof of Lemma 4.1.5.

We distinguish two cases.

\medskip\noindent\textit{First case.}\quad
\(\deg Q=\deg P>p^q\).  Then \(\deg P_1<\deg P\), and the
induction hypothesis applies.

\medskip\noindent\textit{Second case.}\quad
\(\deg P=p^q\), with \(a_q\ne0\).  Then
\[
  \deg P_1=\deg P=p^q
  \qquad\text{and}\qquad
  \deg Q_1<p^q.
\]
We cannot have \(Q=0\), for otherwise \(G_1\) would not be
connected.  If the polynomial \(Q_1\) has no linear part, the
preceding transformation may be repeated.  Continuing in this
fashion, we eventually obtain an equation of the form
\[
  \tag{\(*\)\(_s\)}\label{sga3:xvii:lemma:4-1-5:equation:star-s}
  P_s(X)=-X_1+\sum_{j=1}^{q}a_j^{1/p^s}X_1^{p^j}+Q_s(X),
\]
\SGARefTarget{sga3:xvii:lemma:4:1:5:eq:star:s}
where
\[
  Q_s(X)=Q(X_2,\ldots,X_n)^{1/p^s}
\]
is an additive polynomial with a nonzero linear part, and moreover
\(\deg Q_s<p^q\).  Suppose, for example, that the coefficient of
\(X_2\) in \(Q_s\) is nonzero, and let \(-L\) be the linear part
of \(P_s(X)\).  After a linear change of coordinates, the equation
of \(G_1\) becomes
\[
  P'(X)=-L+\sum_{j=1}^{q}a_j^{1/p^s}X_1^{p^j}
    +Q'(L,X_3,\ldots,X_n),
\]
where \(Q'\) is an additive polynomial without a linear part and
\(\deg Q'<p^q\).  We are thereby reduced to the first
case.\footnote[6]{Editor's note: replace \(X_1\) by \(L\).}


% END INCLUDED SGA3_English_current_buildable_source_20260823/source/inputs/XVII/tex/components/14_expose_XVII_lemma4_1_5_completion_pages015_016.tex

% BEGIN INCLUDED SGA3_English_current_buildable_source_20260823/source/inputs/XVII/tex/components/15_expose_XVII_subsection4_2_through_height1_case_pages016_017.tex
% Authority: Expo17.pdf, local pp.16--17 / combined pp.978--979 /
% printed pp.562--563 and the top of printed p.564.

\SGASubsection{4.2}{Radicial Groups}
\SGARefTarget{sga3:xvii:subsection:4:2}

\SGAStatement{\SGAInlineRefTarget{sga3:xvii:proposition:4-2-1}Proposition 4.2.1}{Let \(G\) be a radicial algebraic
group defined over a field \(k\) of characteristic \(p>0\).  The
following conditions are equivalent:
\begin{enumerate}
\item[(i)] \SGARefTarget{sga3:xvii:proposition:4:2:1:item:i}
  \(G\) is unipotent.
\item[(ii)] \SGARefTarget{sga3:xvii:proposition:4:2:1:item:ii}
  \(G\) has a central composition series whose successive
  quotients are isomorphic to \(\boldsymbol\alpha_p\).
\item[(iii)] \SGARefTarget{sga3:xvii:proposition:4:2:1:item:iii}
  \(G\) has a central, characteristic composition series whose
  successive quotients are isomorphic to
  \((\boldsymbol\alpha_p)^r\).
\item[(iv)] \SGARefTarget{sga3:xvii:proposition:4:2:1:item:iv}
  \(G_{\overline{k}}\) contains no subgroup isomorphic to
  \(\boldsymbol\mu_p\).
\item[(v)] \SGARefTarget{sga3:xvii:proposition:4:2:1:item:v}
  \(\mathfrak g=\operatorname{Lie}G\) is a unipotent
  \(p\)-Lie algebra (\SGARefLink{sga3:xvii:definition:3:6}{3.6}).
\end{enumerate}}
\SGARefTarget{sga3:xvii:proposition:4:2:1}

\SGAProof{Proof}
The implications
\SGARefLink{sga3:xvii:proposition:4:2:1:item:iii}{(iii)}\(\Rightarrow\)\SGARefLink{sga3:xvii:proposition:4:2:1:item:ii}{(ii)}\(\Rightarrow\)\SGARefLink{sga3:xvii:proposition:4:2:1:item:i}{(i)} are clear;
\SGARefLink{sga3:xvii:proposition:4:2:1:item:i}{(i)}\(\Rightarrow\)\SGARefLink{sga3:xvii:proposition:4:2:1:item:iii}{(iii)} is \SGARefLink{sga3:xvii:lemma:3:9:item:i}{3.9(i)}, and
\SGARefLink{sga3:xvii:proposition:4:2:1:item:i}{(i)}\(\Rightarrow\)\SGARefLink{sga3:xvii:proposition:4:2:1:item:iv}{(iv)} follows from \SGARefLink{sga3:xvii:proposition:2:4:item:i}{2.4(i)}.

We shall need the following lemma on abelian \(p\)-Lie algebras.

\SGAStatement{\SGAInlineRefTarget{sga3:xvii:lemma:4-2-2}Lemma 4.2.2}{Let \(\mathfrak g\) be an abelian,
finite-dimensional \(p\)-Lie algebra over a perfect field \(k\).
Then \(\mathfrak g\) can be written uniquely as a direct sum
\[
  \mathfrak g=\mathfrak r\oplus\mathfrak u
\]
of a \(p\)-Lie subalgebra \(\mathfrak r\) on which the \(p\)-th
power is bijective and a unipotent \(p\)-Lie subalgebra
\(\mathfrak u\) (\SGARefLink{sga3:xvii:definition:3:6}{3.6}).  The algebra \(\mathfrak r\) is called the
\emph{reductive part} of \(\mathfrak g\), and \(\mathfrak u\) its
\emph{unipotent part}.  The formation of \(\mathfrak r\) and
\(\mathfrak u\) is compatible with extension of the field \(k\).
If, moreover, \(k\) is algebraically closed, then
\(\mathfrak r\) has a basis \(e_i\) such that
\[
  e_i^{(p)}=e_i.
\]}
\SGARefTarget{sga3:xvii:lemma:4:2:2}

The proof of this lemma is easy and is left to the reader; see
Bourbaki, \emph{Groupes et algèbres de Lie}, Chapter I, \S1,
Exercise 23.  We merely note that \(\mathfrak u\) is the kernel of
a suitable iterate of the map \(X\mapsto X^{(p)}\), while
\(\mathfrak r\) is the image of that same iterate.

We now prove \SGARefLink{sga3:xvii:proposition:4:2:1:item:iv}{(iv)}\(\Rightarrow\)\SGARefLink{sga3:xvii:proposition:4:2:1:item:v}{(v)}.  After extending the base
field if necessary, we may suppose that \(k\) is algebraically
closed.  Let \(X\in\mathfrak g\), and let \(\mathfrak h\) be the
\(p\)-Lie subalgebra of \(\mathfrak g\) generated by \(X\).  The
algebra \(\mathfrak h\) is plainly commutative, and its reductive
part (\SGARefLink{sga3:xvii:lemma:4-2-2}{4.2.2}) is zero.  Otherwise, by the cited result,
\(\mathfrak h\) would contain an element \(Y\ne0\) such that
\(Y^{(p)}=Y\), and consequently, by Appendix II, \SGARefLink{sga3:xvii:proposition:2:1}{2.1} and \SGARefLink{sga3:xvii:proposition:2:2}{2.2},
\(G\) would contain a subgroup isomorphic to
\(\boldsymbol\mu_p\), contrary to the hypothesis.  Thus
\(\mathfrak h\), and hence \(\mathfrak g\), is a unipotent
\(p\)-Lie algebra.

It remains to prove \SGARefLink{sga3:xvii:proposition:4:2:1:item:v}{(v)}\(\Rightarrow\)\SGARefLink{sga3:xvii:proposition:4:2:1:item:i}{(i)}, the least immediate
implication in \SGARefLink{sga3:xvii:proposition:4-2-1}{4.2.1}.

\medskip\noindent\textit{(a) The case in which \(G\) has height
\(1\)} (Exposé VII\textsubscript{A}, \SGARefLink{sga3:VIIA:4.1.3}{4.1.3}).
Since \(G\) is radicial, it is affine and hence isomorphic to an
algebraic subgroup of a linear group \(\operatorname{GL}(V)\)
(Exposé VI\textsubscript{B}, \S11).  This embedding identifies
\(\mathfrak g\) with a \(p\)-Lie subalgebra of
\(\operatorname{End}(V)\), and the \(p\)-th power of \(X\) in
\(\mathfrak g\) agrees with the \(p\)-th power of the endomorphism
\(X\) (Exposé VII\textsubscript{A}, \SGARefLink{sga3:VIIA:6.4.4}{6.4.4}).  Since
\(\mathfrak g\) is unipotent by hypothesis, it is therefore an
algebra of nilpotent endomorphisms of \(V\).  By Engel's theorem
(Bourbaki, \emph{Groupes et algèbres de Lie}, Chapter I, \S4,
Theorem 1), it is a Lie subalgebra of the Lie algebra
\(\mathfrak h\) of the group of strictly upper-triangular matrices
\(\operatorname{Trigstr}(n)\), relative to a suitable basis of
\(V\).  Since \(G\) has height \(1\), Appendix II, \SGARefLink{sga3:ii:subsection:2-2}{2.2} then shows
that \(G\) itself is an algebraic subgroup of
\(\operatorname{Trigstr}(n)\), and hence is unipotent by
\SGARefLink{sga3:xvii:theorem:3:5:item:v}{3.5(v)}\(\Rightarrow\)\SGARefLink{sga3:xvii:proposition:4:2:1:item:i}{(i)}.


% END INCLUDED SGA3_English_current_buildable_source_20260823/source/inputs/XVII/tex/components/15_expose_XVII_subsection4_2_through_height1_case_pages016_017.tex

% BEGIN INCLUDED SGA3_English_current_buildable_source_20260823/source/inputs/XVII/tex/components/16_expose_XVII_proposition4_2_1_general_case_page017.tex
% Authority: Expo17.pdf, local p.17 / combined p.979 /
% printed p.564, completing Proposition 4.2.1 and stopping before 4.3.

\medskip\noindent\textit{\SGARefLink{sga3:xvii:section:b}{(b)} The general case.}\quad
We argue by induction on the height \(h\) of
\(G\).{\renewcommand{\thefootnote}{*}\footnote{That is, the least
integer \(h\) such that \(F^h=\operatorname{id}_G\); see
Appendix II, 1.}}
The case \(h=1\) has just been treated.  Suppose that \(h>1\), and
put
\[
  G'={}_{F}G,
  \qquad
  G''=G/G'.
\]
The group \(G'\) has height \(1\), and
\(\operatorname{Lie}G'=\operatorname{Lie}G\) is unipotent.
Consequently \(G'\) is unipotent by \SGARefLink{sga3:xvii:appendix:A}{(a)}.  To prove that \(G\) is
unipotent, it is therefore enough to show that \(G''\) is unipotent
(\SGARefLink{sga3:xvii:proposition:2:2}{2.2}).  But \(G''\) has height \(h-1\), so the induction
hypothesis reduces us to proving that \(\operatorname{Lie}G''\) is
unipotent.

In view of (iv)\(\Rightarrow\)(v), it suffices to prove that
\(G''_{\overline{k}}\) contains no subgroup isomorphic to
\(\boldsymbol\mu_p\).  Suppose, then, that such a subgroup has been
given, and let \(H\) be its inverse image in
\(G_{\overline{k}}\).  Since \(G'\) is unipotent, we shall prove
in \S5 that the extension
\[
  e\longrightarrow G'_{\overline{k}}
   \longrightarrow H
   \longrightarrow\boldsymbol\mu_p
   \longrightarrow e
\]
is necessarily trivial; the proof of that fact is independent of
the results of the present section.  In short,
\(\boldsymbol\mu_p\) lifts to \(H\).  Being of height \(1\), it is
then necessarily contained in
\[
  G'_{\overline{k}}={}_{F}(G_{\overline{k}}),
\]
which is impossible because \(G'\) is unipotent.

% END INCLUDED SGA3_English_current_buildable_source_20260823/source/inputs/XVII/tex/components/16_expose_XVII_proposition4_2_1_general_case_page017.tex

% BEGIN INCLUDED SGA3_English_current_buildable_source_20260823/source/inputs/XVII/tex/components/17_expose_XVII_proposition4_3_1_and_lemma4_3_2_pages017_018.tex
% Authority: Expo17.pdf, local pp.17--18 / combined pp.979--980 /
% printed pp.564--566, through Lemma 4.3.2.

\SGASubsection{4.3}{Connected Affine Groups in Characteristic \(p>0\)}
\SGARefTarget{sga3:xvii:subsection:4:3}

\SGAStatement{\SGAInlineRefTarget{sga3:xvii:proposition:4-3-1}Proposition 4.3.1}{Let \(G\) be a connected affine
algebraic group over a field \(k\) of characteristic \(p>0\).
The following conditions are equivalent:
\begin{enumerate}
\item[(i)] \SGARefTarget{sga3:xvii:proposition:4:3:1:item:i}
  \(G\) is unipotent.
\item[(ii)] \SGARefTarget{sga3:xvii:proposition:4:3:1:item:ii}
  \(G\) has a composition series whose successive quotients are
  isomorphic to \(\boldsymbol\alpha_p\) and \(\mathbb G_a\),
  taken in that order.
\item[(iii)] \SGARefTarget{sga3:xvii:proposition:4:3:1:item:iii}
  \(G\) has a characteristic composition series whose successive
  quotients are isomorphic to \((\boldsymbol\alpha_p)^r\) and
  \((\mathbb G_a)^s\), taken in that order.
\item[(iv)] \SGARefTarget{sga3:xvii:proposition:4:3:1:item:iv}
  \(G_{\overline{k}}\) contains no subgroup isomorphic to
  \(\boldsymbol\mu_p\).
\item[(v)] \SGARefTarget{sga3:xvii:proposition:4:3:1:item:v}
  \(\mathfrak g=\operatorname{Lie}G\) is unipotent (\SGARefLink{sga3:xvii:definition:3:6}{3.6}).
\item[(vi)] \SGARefTarget{sga3:xvii:proposition:4:3:1:item:vi}
  \(\mathfrak g\) is nilpotent, and the reductive part of the
  center of \(\mathfrak g\) (\SGARefLink{sga3:xvii:lemma:4-2-2}{4.2.2}) is trivial.
\item[(vi bis)]
  \SGARefTarget{sga3:xvii:proposition:4:3:1:item:vi-bis}
  \(\mathfrak g\) is nilpotent, and every subgroup of
  multiplicative type of the identity component of the center of
  \(G\) is trivial.
\item[(vi ter)]
  \SGARefTarget{sga3:xvii:proposition:4:3:1:item:vi-ter}
  \(G\) is nilpotent, and every subgroup of multiplicative type in
  the identity component of the center of \(G\) is trivial.
\end{enumerate}}
\SGARefTarget{sga3:xvii:proposition:4:3:1}

\SGAProof{Proof}
The implications
\SGARefLink{sga3:xvii:proposition:4:3:1:item:iii}{(iii)}\(\Rightarrow\)\SGARefLink{sga3:xvii:proposition:4:3:1:item:ii}{(ii)}\(\Rightarrow\)\SGARefLink{sga3:xvii:proposition:4:3:1:item:i}{(i)} are clear.  To prove
\SGARefLink{sga3:xvii:proposition:4:3:1:item:i}{(i)}\(\Rightarrow\)\SGARefLink{sga3:xvii:proposition:4:3:1:item:iii}{(iii)}, we shall need the following lemma.

\SGAStatement{\SGAInlineRefTarget{sga3:xvii:lemma:4-3-2}Lemma 4.3.2}{Let \(k\) be a field of characteristic
\(p>0\), let \(n\in\mathbb N\), and let \(k'\) be a radicial
extension of \(k\) such that \((k')^{p^n}\) is contained in \(k\).
For every \(k\)-prescheme \(X\) (respectively, every
\(k'\)-prescheme \(X'\)), denote by \(X^{(p^n)}\) (respectively,
\((X')^\phi\)) the \(k\)-prescheme obtained from \(X\)
(respectively, from \(X'\)) by the base change
\[
  F^n:k\longrightarrow k,\qquad x\longmapsto x^{p^n},
  \qquad
  \bigl(\text{respectively, }\;
  \phi:k'\longrightarrow k,\qquad x'\longmapsto(x')^{p^n}\bigr).
\]
Then, for every \(k\)-prescheme \(X\), there is a functorial
isomorphism
\[
  (X_{k'})^\phi\xrightarrow{\sim}X^{(p^n)}.
\]
Consequently, if \(X\) and \(Y\) are two \(k\)-preschemes for
which there exists a \(k'\)-isomorphism
\[
  u':X_{k'}\xrightarrow{\sim}Y_{k'},
\]
then there exists a \(k\)-isomorphism
\[
  v:X^{(p^n)}\xrightarrow{\sim}Y^{(p^n)}.
\]
If, moreover, \(X\) and \(Y\) are endowed with structures of
\(k\)-group preschemes and \(u'\) is a \(k'\)-homomorphism, then
\(v\) is a \(k\)-homomorphism.}
\SGARefTarget{sga3:xvii:lemma:4:3:2}

The lemma follows immediately from the transitivity of base
changes and from the fact that the composite
\[
  k\longrightarrow k'\xrightarrow{\phi}k
\]
is \(F^n\).

% END INCLUDED SGA3_English_current_buildable_source_20260823/source/inputs/XVII/tex/components/17_expose_XVII_proposition4_3_1_and_lemma4_3_2_pages017_018.tex

% BEGIN INCLUDED SGA3_English_current_buildable_source_20260823/source/inputs/XVII/tex/components/18_expose_XVII_proposition4_3_1_proof_page018.tex
% Authority: Expo17.pdf, local p.18 / combined p.980 /
% printed pp.566--567, through (iv) => (i).

\medskip\noindent\textit{Continuation of the proof of \SGARefLink{sga3:xvii:proposition:4-3-1}{4.3.1}.}

\medskip\noindent
\textit{(i)\(\Rightarrow\)\SGARefLink{sga3:xvii:proposition:4:3:1:item:iii}{(iii)}.}\quad
We argue by induction on \(\dim G\).  If \(\dim G=0\), then,
because \(G\) is connected, it is radicial, and \SGARefLink{sga3:xvii:lemma:3:9:item:i}{3.9(i)} applies.
If \(\dim G>0\), there exists an integer \(m\geq0\) such that the
quotient \(G/F_mG\) is a smooth group (Appendix II, \SGARefLink{sga3:xvii:lemma:3-1}{3.1}), plainly
connected and nonzero.  Applying
\SGARefLink{sga3:xvii:proposition:4:1:1:item:i}{4.1.1(i)}\(\Rightarrow\)\SGARefLink{sga3:xvii:lemma:3-1:item:iii}{(iii)} to this quotient, we see that \(G\)
has an algebraic subgroup \(G'\) which is characteristic and
connected, and for which the quotient
\[
  G''=G/G'
\]
is a form of \((\mathbb G_a)^r\), where \(r>0\).  By \SGARefLink{sga3:xvii:lemma:4:1:5}{4.1.5}, if
\(K\) is a perfect closure of \(k\), then
\[
  G''_K\xrightarrow{\sim}(\mathbb G_{a,K})^r.
\]
Since \(G''\) is of finite type over \(k\), there exists a finite
radicial extension \(k'\) of \(k\) such that
\[
  G''_{k'}\xrightarrow{\sim}(\mathbb G_{a,k'})^r
\]
(Exposé VI\textsubscript{B}, \S10).  Choose \(n>0\) such that
\((k')^{p^n}\subset k\).  Retaining the notation of \SGARefLink{sga3:xvii:lemma:4-3-2}{4.3.2}, we
deduce a \(k\)-isomorphism of algebraic groups
\[
  (\mathbb G_{a,k})^r
   =\bigl((\mathbb G_{a,k'})^r\bigr)^\phi
   \xrightarrow{\sim}(G'')^{(p^n)}.
\]
Now consider the Frobenius homomorphism relative to \(G''\)
(Exposé VII\textsubscript{A}, \S4):
\[
  F^n:G''\longrightarrow(G'')^{(p^n)}.
\]
Since \(G''\), and hence also \((G'')^{(p^n)}\), is smooth over
\(k\), while \(F^n\) is radicial, \(F^n\) is an epimorphism for
the fpqc topology.  Thus \((G'')^{(p^n)}\) is identified with
\(G''/F_n(G'')\), and we have shown that
\[
  G''/F_n(G'')\xrightarrow{\sim}(\mathbb G_a)^r
\]
as algebraic groups.  The inverse image \(G'_n\) of \(F_n(G'')\)
in \(G\) is a connected, characteristic subgroup of \(G\), of
dimension strictly less than that of \(G\); the induction
hypothesis therefore applies to \(G'_n\).

\medskip\noindent
\textit{(i)\(\Rightarrow\)(iv).}\quad
This follows from \SGARefLink{sga3:VIIA:2.4:item:i}{2.4(i)}.

\medskip\noindent
\textit{(iv)\(\Rightarrow\)(i).}\quad
Regard \(G\) as an extension of a smooth connected group \(G''\)
by a radicial group \(G'\) (Appendix II, \SGARefLink{sga3:xvii:lemma:3-1}{3.1}).  The group \(G'\)
is unipotent by \SGARefLink{sga3:xvii:proposition:4:2:1:item:iv}{4.2.1(iv)}\(\Leftrightarrow\)\SGARefLink{sga3:xvii:lemma:3-1:item:i}{(i)}.  It is therefore
enough to prove that \(G''\) is unipotent, and for this it is
enough to show that \(G''_{\overline{k}}\) contains no subgroup
isomorphic to \(\boldsymbol\mu_p\), by
\SGARefLink{sga3:xvii:proposition:4:1:1:item:i}{4.1.1(i)}\(\Leftrightarrow\)(iv).  If
\(G''_{\overline{k}}\) contained such a subgroup, it would lift
to \(G_{\overline{k}}\) by the result \SGARefLink{sga3:xvii:subsection:5:1}{5.1}, already used above,
whose proof is given in \S5.  This would contradict (iv).


% END INCLUDED SGA3_English_current_buildable_source_20260823/source/inputs/XVII/tex/components/18_expose_XVII_proposition4_3_1_proof_page018.tex

% BEGIN INCLUDED SGA3_English_current_buildable_source_20260823/source/inputs/XVII/tex/components/19_expose_XVII_proposition4_3_1_completion_through_section4_5_opening_page019.tex
% Authority: Expo17.pdf, local p.19 / combined p.981 /
% printed pp.568--569, through the statement in 4.5.

\medskip\noindent\textit{Continuation of the proof of \SGARefLink{sga3:xvii:proposition:4-3-1}{4.3.1}.}

\medskip\noindent
\textit{(i)\(\Rightarrow\)\SGARefLink{sga3:xvii:proposition:4:3:1:item:v}{(v)}.}\quad
This follows from \SGARefLink{sga3:xvii:corollary:3:7}{3.7}.

\medskip\noindent
\textit{(v)\(\Rightarrow\)\SGARefLink{sga3:xvii:proposition:4:3:1:item:vi}{(vi)}.}\quad
Indeed, since
\[
  (\operatorname{ad}X)^{p^r}
    =\operatorname{ad}\bigl(X^{(p^r)}\bigr)
  \qquad\text{(Exposé VII\textsubscript{A}, 5.2),}
\]
\(\operatorname{ad}X\) is nilpotent when \(\mathfrak g\) is
unipotent; hence \(\mathfrak g\) is nilpotent by Engel's theorem.
Moreover, if \(\mathfrak g\) is unipotent, then its center plainly
is as well, and the reductive part of that center is therefore
trivial (\SGARefLink{sga3:xvii:lemma:4-2-2}{4.2.2}).

\medskip\noindent
\textit{(vi)\(\Rightarrow\)(iv).}\quad
Indeed, if \(G_{\overline{k}}\) contained a subgroup isomorphic to
\(\boldsymbol\mu_p\), its Lie algebra would contain a nonzero
element \(X\) such that
\[
  X^{(p)}=X
  \qquad\text{(Appendix II, 2.1),}
\]
and hence \(X^{(p^r)}=X\) for every \(r>0\).
Since \(\operatorname{ad}X\) is nilpotent because \(\mathfrak g\)
is nilpotent, and
\[
  (\operatorname{ad}X)^{p^r}
    =\operatorname{ad}\bigl(X^{(p^r)}\bigr),
\]
we necessarily have \(\operatorname{ad}X=0\).  Thus \(X\) belongs
to the reductive part of the center of
\(\mathfrak g_{\overline{k}}\), contradicting (vi).

\medskip\noindent
\textit{(i)\(\Rightarrow\)(vi ter).}\quad
This follows from \SGARefLink{sga3:VIIA:2.4:item:i}{2.4(i)} and \SGARefLink{sga3:xvii:theorem:3:5:item:i}{3.5(i)}\(\Rightarrow\)(iii).

\medskip\noindent
\textit{(vi ter)\(\Rightarrow\)(vi bis).}\quad
Indeed, if \(G\) is nilpotent, then so is its subgroup
\({}_{F}G\).  Appendix II, \SGARefLink{sga3:ii:subsection:2-2}{2.2} shows that \({}_{F}G\) is
nilpotent if and only if
\[
  \operatorname{Lie}({}_{F}G)=\operatorname{Lie}G
  \qquad\text{(Exposé VII\textsubscript{A})}
\]
is nilpotent.

\medskip\noindent
\textit{(vi bis)\(\Rightarrow\)(vi).}\quad
Let \(Z\) be the identity component of the center of \(G\), and
let \(\mathfrak r\) be the reductive part of the center of
\(\mathfrak g\).  We must prove that \(\mathfrak r=0\).
It is immediate that \(\mathfrak r\) is a characteristic
\(p\)-Lie subalgebra of
\[
  \mathfrak g=\operatorname{Lie}({}_{F}G);
\]
that is, it is stable under the functor
\(\operatorname{Aut}_{p\text{-Lie}}(\mathfrak g)\).
Consequently \(\mathfrak r\) is the Lie algebra of a
characteristic radicial subgroup \(R\) of \({}_{F}G\)
(Appendix II, \SGARefLink{sga3:ii:subsection:2-2}{2.2}).  The last assertion of \SGARefLink{sga3:xvii:lemma:4-2-2}{4.2.2} and Appendix II,
\SGARefLink{sga3:xvii:proposition:2:1}{2.1} show moreover that \(R\) is a form of
\((\boldsymbol\mu_p)^r\).  Since \(R\) is characteristic in
\({}_{F}G\), which is itself characteristic in \(G\)
(Appendix II, 1), it is invariant in \(G\), and hence central
because \(G\) is connected (\SGARefLink{sga3:ix:corollary:5:5}{Exposé IX, 5.5}).  By (vi bis),
\(R=0\), and therefore \(\mathfrak r=0\).

\SGASubsection{4.4}{Étale Groups}
\SGARefTarget{sga3:xvii:subsection:4:4}

The following proposition is an easy consequence of the Sylow
theorems and of the structure of finite \(q\)-groups
(cf.\ J.-P. Serre, \emph{Corps locaux}, Chapter IX, \S1).

\SGAStatement{Proposition 4.4.1}{Let \(G\) be a finite étale
algebraic group over an algebraically closed field.  Then \(G\)
is unipotent if and only if, for every prime number \(q\) distinct
from the characteristic \(p\) of the field, \(G\) contains no
subgroup isomorphic to \(\boldsymbol\mu_q\).}
\SGARefTarget{sga3:xvii:proposition:4:4:1}

\SGASubsection{4.5}{Abelian Varieties}
\SGARefTarget{sga3:xvii:subsection:4:5}

Let \(G\) be an abelian variety over an algebraically closed field
\(k\).  The following conditions are equivalent:
\begin{enumerate}
\item[(i)] \SGARefTarget{sga3:xvii:subsection:4:5:item:i}
  \(G\) is unipotent.
\item[(ii)] \SGARefTarget{sga3:xvii:subsection:4:5:item:ii}
  \(G=0\).
\item[(iii)] \SGARefTarget{sga3:xvii:subsection:4:5:item:iii}
  There exists an integer \(n\), relatively prime to the
  characteristic \(p\) of \(k\), such that \(G\) contains no
  subgroup isomorphic to \(\boldsymbol\mu_n\).
\end{enumerate}

% END INCLUDED SGA3_English_current_buildable_source_20260823/source/inputs/XVII/tex/components/19_expose_XVII_proposition4_3_1_completion_through_section4_5_opening_page019.tex

% BEGIN INCLUDED SGA3_English_current_buildable_source_20260823/source/inputs/XVII/tex/components/20_expose_XVII_section4_5_proof_through_theorem5_1_1_pages019_021.tex
% Authority: Expo17.pdf, local pp.19--21 / combined pp.981--983 /
% printed pp.569--574, from the proof in 4.5 through the Galois
% cohomology notation following Theorem 5.1.1.

Indeed, if \(G\) is an abelian variety of dimension \(d\), it is
known (cf.\ S. Lang, \emph{Abelian Varieties}, Chapter IV, \S3,
\SGARefLink{sga3:xvii:section:6:site:da2ef89f51}{Theorem 6}) that the group \({}_nG(k)\), for \(n\) prime to \(p\),
is isomorphic to
\[
  (\mathbb Z/n\mathbb Z)^{2d},
\]
and hence to \((\boldsymbol\mu_n)^{2d}\).  Thus
(iii)\(\Rightarrow\)(ii), while
(ii)\(\Rightarrow\)(i)\(\Rightarrow\)(iii) are clear.

\SGASubsection{4.6}{The General Case}
\SGARefTarget{sga3:xvii:subsection:4:6}

If \(G\) and \(H\) are two algebraic groups defined over an
algebraically closed field \(k\), we shall denote by \(P(G,H)\)
the property
\[
  \text{``there is no algebraic subgroup of \(G\) isomorphic to
  \(H\).''}
\]
We then obtain the following characterizations of unipotent
groups.

\SGAStatement{\SGAInlineRefTarget{sga3:xvii:theorem:4-6-1}Theorem 4.6.1}{Let \(G\) be an algebraic group
defined over an algebraically closed field \(k\) of characteristic
\(p\).  Then:
\begin{enumerate}
\item[(i)] \SGARefTarget{sga3:xvii:theorem:4:6:1:item:i}
  If \(G\) is smooth, affine, and connected, then
  \[
    \begin{split}
    G\text{ is unipotent}
      &\Longleftrightarrow
        \text{there exists \(n>1\) such that \(P(G,\boldsymbol\mu_n)\)
        holds}\\
      &\Longleftrightarrow P(G,\mathbb G_m)\text{ holds}.
    \end{split}
  \]
\item[(ii)] \SGARefTarget{sga3:xvii:theorem:4:6:1:item:ii}
  If \(G\) is smooth and connected, then
  \[
    G\text{ is unipotent}
      \Longleftrightarrow
    \text{there exists \(n\) prime to \(p\) such that
      \(P(G,\boldsymbol\mu_n)\) holds}.
  \]
\item[(iii)] \SGARefTarget{sga3:xvii:theorem:4:6:1:item:iii}
  If \(G\) is smooth, then
  \[
    G\text{ is unipotent}
      \Longleftrightarrow
    P(G,\boldsymbol\mu_n)\text{ holds for every prime \(n\ne p\)}.
  \]
\item[(iv)] \SGARefTarget{sga3:xvii:theorem:4:6:1:item:iv}
  If \(G\) is affine and connected and \(p>0\), then
  \[
    G\text{ is unipotent}
      \Longleftrightarrow P(G,\boldsymbol\mu_p)\text{ holds}.
  \]
\item[(v)] \SGARefTarget{sga3:xvii:theorem:4:6:1:item:v}
  If \(G\) is connected and \(p>0\), then
  \[
    \begin{split}
    G\text{ is unipotent}
      \Longleftrightarrow{}&
      \text{there exists \(n\) prime to \(p\) such that}\\[-2pt]
      &P(G,\boldsymbol\mu_n)\text{ and }
        P(G,\boldsymbol\mu_p)\text{ both hold}.
    \end{split}
  \]
\item[(vi)] \SGARefTarget{sga3:xvii:theorem:4:6:1:item:vi}
  For an arbitrary algebraic group \(G\),
  \[
    G\text{ is unipotent}
      \Longleftrightarrow
    P(G,\boldsymbol\mu_n)\text{ holds for every prime \(n\)}.
  \]
\end{enumerate}}
\SGARefTarget{sga3:xvii:theorem:4:6:1}

\SGAProof{Proof}
Assertion (i) follows from \SGARefLink{sga3:xvii:proposition:4-1-1}{4.1.1}, and (iv) from \SGARefLink{sga3:xvii:proposition:4-3-1}{4.3.1}.  We shall
prove (vi); assertions (ii), (iii), and (v) are proved in the same
way and are left to the reader.

Let \(G\) be an algebraic group.  If \(G\) is unipotent, then
\(P(G,\boldsymbol\mu_n)\) holds for every \(n>1\) by 2.4(i).
Conversely, suppose that \(P(G,\boldsymbol\mu_n)\) holds for every
prime \(n\), and let us prove that \(G\) is unipotent.  Write
\(G^0\) for the identity component of \(G\).  If \(G^0\) is
smooth, a classical theorem of
Chevalley\begingroup\renewcommand{\thefootnote}{*}\footnote[0]{%
\emph{Séminaire Bourbaki}, no.\ 145, 1956/57.}\endgroup
shows that \(G^0\) is an extension
of an abelian variety \(A\) by a smooth connected affine group
\(L\).  If \(G\) is not smooth---which entails \(p>0\)
(Exposé VI\textsubscript{B}, \SGARefLink{cor:sga3-VIB-1.6.1}{1.6.1})---there is an integer \(n>0\)
such that
\[
  G''=G^0/{}_{F^n}(G)
\]
is smooth (Appendix II, \SGARefLink{sga3:xvii:lemma:3-1}{3.1}).  Thus \(G''\) is an extension of an
abelian variety \(A\) by a smooth connected linear group \(L''\).
Let \(L\) be the inverse image of \(L''\) in \(G^0\).  It is still
affine and connected, since \({}_{F^n}(G)\) is radicial.  In every
case, therefore, \(G\) has a composition series
\[
  0\subset L\subset G^0\subset G
\]
such that \(L\) is affine and connected, \(G^0/L=A\) is an
abelian variety, and \(G/G^0\) is an étale group.

If \(P(G,\boldsymbol\mu_n)\) holds, then a fortiori
\(P(L,\boldsymbol\mu_n)\) holds, and hence \(L\) is unipotent
(\SGARefLink{sga3:xvii:proposition:4-1-1}{4.1.1} and \SGARefLink{sga3:xvii:proposition:4-3-1}{4.3.1}).  If \(A\ne0\), there are a prime number \(n\)
and a subgroup of \(A\) isomorphic to \(\boldsymbol\mu_n\) (\SGARefLink{sga3:xvii:subsection:4:5}{4.5}).
By \SGARefLink{sga3:xvii:subsection:5:1}{5.1} below, this subgroup lifts to \(G\), contradicting
\(P(G,\boldsymbol\mu_n)\).  Hence \(A=0\).  Finally, if
\(G/G^0\) is not unipotent, there are an integer \(q\) and a
subgroup of \(G/G^0\) isomorphic to \(\boldsymbol\mu_q\) (\SGARefLink{sga3:xvii:proposition:4:4:1}{4.4.1}).
As above, it follows that \(P(G,\boldsymbol\mu_q)\) does not hold.
Thus \(G/G^0\) is unipotent, and consequently so is \(G\).

\SGASection{5}{Extensions of Groups of Multiplicative Type by
Unipotent Groups}
\SGARefTarget{sga3:xvii:section:5}

\SGASubsection{5.1}{Statement of the Theorem}
\SGARefTarget{sga3:xvii:subsection:5:1}

\SGAStatement{\SGAInlineRefTarget{sga3:xvii:definition:5-1-0}Definition 5.1.0}{Let \(k\) be a field and \(G\) an
algebraic \(k\)-group.  Following the terminology introduced by
Rosenlicht (\emph{Questions of rationality for solvable algebraic
groups over non perfect fields}, Annali di Matematica 61 (1963)),
we shall say that \(G\) is \emph{\(k\)-solvable} if it satisfies
the following equivalent conditions:
\begin{enumerate}
\item[(i)] \SGARefTarget{sga3:xvii:definition:5:1:0:item:i}
  \(G\) has a composition series whose successive quotients are
  isomorphic to \(\mathbb G_a\).
\item[(ii)] \SGARefTarget{sga3:xvii:definition:5:1:0:item:ii}
  \(G\) has a characteristic composition series \((G_i)\) such
  that the successive quotients \(G_i/G_{i+1}\) are commutative
  and themselves have composition series whose successive
  quotients are isomorphic to \(\mathbb G_a\).
\end{enumerate}}
\SGARefTarget{sga3:xvii:definition:5:1:0}

The implication \SGARefLink{sga3:xvii:definition:5:1:0:item:i}{(i)}\(\Rightarrow\)\SGARefLink{sga3:xvii:definition:5:1:0:item:ii}{(ii)} is proved in
\emph{loc.\ cit.}  In fact, for the composition series \((G_i)\)
one may take the series introduced in the proof of \SGARefLink{sga3:xvii:lemma:3:9:item:ii}{3.9(ii)}.

\SGAStatement{\SGAInlineRefTarget{sga3:xvii:theorem:5-1-1}Theorem 5.1.1}{Let \(k\) be a field, \(U\) a
unipotent algebraic \(k\)-group, \(H\) a \(k\)-group of
multiplicative type, and \(E\) an algebraic \(k\)-group which is
an extension of \(H\) by \(U\), so that there is an exact sequence
\[
  1\longrightarrow U\longrightarrow E\longrightarrow H
    \longrightarrow1.
\]
Then:
\begin{enumerate}
\item[(i)] \SGARefTarget{sga3:xvii:theorem:5:1:1:item:i}
  The extension \(E\) is trivial in each of the following cases:
  \begin{enumerate}
  \item[(a)] \SGARefTarget{sga3:xvii:theorem:5:1:1:item:i:a}
    \(k\) is algebraically closed;
  \item[(b)] \SGARefTarget{sga3:xvii:theorem:5:1:1:item:i:b}
    \(k\) is perfect and one of the groups \(U\) or \(H\) is
    connected;
  \item[(c)] \SGARefTarget{sga3:xvii:theorem:5:1:1:item:i:c}
    \(U\) is \(k\)-solvable;
  \item[(d)] \SGARefTarget{sga3:xvii:theorem:5:1:1:item:i:d}
    \(U\) is smooth and \(H\) is connected.
  \end{enumerate}
\item[(ii)] \SGARefTarget{sga3:xvii:theorem:5:1:1:item:ii}
  If \(H'\) and \(H''\) are two lifts of \(H\) in \(E\), then
  they are conjugate by an element of \(U(k)\) in each of the
  following cases:
  \begin{enumerate}
  \item[(a)] \SGARefTarget{sga3:xvii:theorem:5:1:1:item:ii:a}
    \(k\) is algebraically closed and \(H\) is smooth;
  \item[(b)] \SGARefTarget{sga3:xvii:theorem:5:1:1:item:ii:b}
    \(k\) is algebraically closed and \(U\) is smooth.
  \end{enumerate}
\end{enumerate}
(During the proof we shall indicate other cases in which the
conclusion of \SGARefLink{sga3:xvii:theorem:5:1:1:item:ii}{(ii)} is true without assuming that \(k\) is
algebraically closed.)}
\SGARefTarget{sga3:xvii:theorem:5:1:1}

If \(U\) is an algebraic group (respectively, a commutative
algebraic group) defined over a field \(k\), we denote by
\(H^1(k,U)\) (respectively, by \(H^i(k,U)\), \(i\geq0\)) the
pointed set (respectively, the \(i\)-th group) of Galois
cohomology of \(k\) with values in \(U\); see J.-P. Serre,
\emph{Cohomologie galoisienne}, Lecture Notes in Mathematics
no.\ 5, Springer.

% END INCLUDED SGA3_English_current_buildable_source_20260823/source/inputs/XVII/tex/components/20_expose_XVII_section4_5_proof_through_theorem5_1_1_pages019_021.tex

% BEGIN INCLUDED SGA3_English_current_buildable_source_20260823/source/inputs/XVII/tex/components/21_expose_XVII_subsection5_2_through_lemma5_2_4_pages022_023.tex
% Authority: Expo17.pdf, local pp.22--23 / combined pp.984--985 /
% printed pp.575--578, through the proof of Lemma 5.2.4(b).

If \(S\) is a prescheme, \(H\) a commutative \(S\)-group
prescheme, and \(G\) an \(S\)-group prescheme acting on \(H\), we
denote by \(H^i(G,H)\)\footnote{Editor's note: \(H_0^i\) has been
replaced by \(H^i\), in agreement with the notation of Appendix I
and Exposé I.} the \(i\)-th Hochschild cohomology group of \(G\)
with values in \(H\) (Appendix I, 1).

To prove \SGARefLink{sga3:xvii:theorem:5-1-1}{5.1.1}, we shall proceed in several steps.

\SGASubsection{5.2}{\texorpdfstring{Proof of \SGARefLink{sga3:xvii:theorem:5:1:1:item:i}{5.1.1(i)} and \SGARefLink{sga3:xvii:theorem:5:1:1:item:ii}{(ii)}
when \(U\) is smooth and \(H\) is étale}{Proof of \SGARefLink{sga3:xvii:theorem:5:1:1:item:i}{5.1.1(i)} and
\SGARefLink{sga3:xvii:theorem:5:1:1:item:ii}{(ii)} when U is smooth and H is etale}}
\SGARefTarget{sga3:xvii:subsection:5:2}

\SGAStatement{Lemma 5.2.1}{With the notation of \SGARefLink{sga3:xvii:theorem:5-1-1}{5.1.1}, suppose
that \(H\) is étale.  The canonical morphism \(E\to H\) has a
section \(s:H\to E\), defined over \(k\), in each of the following
cases:
\begin{enumerate}
\item[(a)] \SGARefTarget{sga3:xvii:lemma:5:2:1:item:a}
  \(k\) is algebraically closed;
\item[(b)] \SGARefTarget{sga3:xvii:lemma:5:2:1:item:b}
  \(k\) is perfect and \(U\) is smooth and connected;
\item[(c)] \SGARefTarget{sga3:xvii:lemma:5:2:1:item:c}
  \(U\) is \(k\)-solvable.
\end{enumerate}}
\SGARefTarget{sga3:xvii:lemma:5:2:1}

Assertion \SGARefLink{sga3:xvii:appendix:A}{(a)} follows simply from the surjectivity of
\(E(k)\to H(k)\).  We have \SGARefLink{sga3:xvii:section:b}{(b)}\(\Rightarrow\)\SGARefLink{sga3:xvii:appendix:C}{(c)} by \SGARefLink{sga3:xvii:lemma:4:1:2:item:b}{4.1.2(b)}
and \SGARefLink{sga3:xvii:lemma:2:3:bis}{2.3 bis}, so it is enough to treat case \SGARefLink{sga3:xvii:appendix:C}{(c)}.  Let \(x\) be a
point of \(H\).  It is an open and closed subset of \(H\), and
the induced subscheme is isomorphic to \(\operatorname{Spec}K\),
where \(K\) is a finite separable extension of \(k\).  Let \(X\)
be the inverse image in \(E\) of the scheme \(x\).  The
\(K\)-scheme \(X\) is a principal homogeneous \(K\)-space under
\(U_K\).  But \(U_K\) has a composition series whose successive
quotients are isomorphic to \((\mathbb G_a)_K\); hence
\[
  H^1(K,U_K)=0
\]
(J.-P. Serre, \emph{Cohomologie galoisienne}, Chapter III,
\SGARefLink{sga3:xvii:section:6}{Proposition 6}).  Consequently \(X\) is trivial and has a
\(K\)-rational point.  Thus, over every point \(x\) of \(H\), we
obtain a \(k\)-section of \(E\to H\), and these give a section
\(H\to E\).

\SGAStatement{\SGAInlineRefTarget{sga3:xvii:lemma:5-2-3}Lemma 5.2.3}{\footnote{Editor's note: there is no
item numbered 5.2.2.}With the notation of \SGARefLink{sga3:xvii:theorem:5-1-1}{5.1.1}, suppose that
\(H\) is étale.  Then:
\begin{enumerate}
\item[(i)] \SGARefTarget{sga3:xvii:lemma:5:2:3:item:i}
  The extension \(E\) is trivial in each of the following cases:
  \begin{enumerate}
  \item[(a)] \SGARefTarget{sga3:xvii:lemma:5:2:3:item:i:a}
    \(U\) is commutative and \(E\to H\) has a section;
  \item[(b)] \SGARefTarget{sga3:xvii:lemma:5:2:3:item:i:b}
    \(k\) is algebraically closed;
  \item[(c)] \SGARefTarget{sga3:xvii:lemma:5:2:3:item:i:c}
    \(k\) is perfect and \(U\) is connected;
  \item[(d)] \SGARefTarget{sga3:xvii:lemma:5:2:3:item:i:d}
    \(U\) is \(k\)-solvable.
  \end{enumerate}
  Moreover, in each of the preceding cases, two lifts \(H'\) and
  \(H''\) of \(H\) in \(E\) are conjugate by an element of
  \(U(k)\).
\item[(ii)] \SGARefTarget{sga3:xvii:lemma:5:2:3:item:ii}
  Let \(R\) be the \(k\)-functor
  \[
    (\operatorname{Sch}/k)^\circ\longrightarrow\operatorname{Ens}
  \]
  for which, for every \(k\)-prescheme \(S\), \(R(S)\) is the set
  of lifts of \(H_S\) in \(E_S\).  If \(U\) is commutative, then
  \(R\) is representable by a nonempty affine \(k\)-scheme.  The
  group \(U\) acts on \(R\) by inner automorphisms, and this action
  makes \(R\) a homogeneous space under \(U\) for the fpqc
  topology (cf.\ Exposé IV).
\end{enumerate}}
\SGARefTarget{sga3:xvii:lemma:5:2:3}

\SGAProof{Proof of \SGARefLink{sga3:xvii:lemma:5:2:3:item:i}{(i)}}
We shall reduce cases \SGARefLink{sga3:xvii:section:b}{(b)}, \SGARefLink{sga3:xvii:appendix:C}{(c)}, and \SGARefLink{sga3:xvii:lemma:5:2:3:item:i:d}{(d)} to case \SGARefLink{sga3:xvii:appendix:A}{(a)}.

\medskip\noindent\textit{Case \SGARefLink{sga3:xvii:section:b}{(b)}.}\quad
Since \(U\) has a characteristic composition series with
commutative successive quotients
(\SGARefLink{sga3:xvii:theorem:3-5}{3.5}, \SGARefLink{sga3:xvii:theorem:3:5:item:i}{(i)}\(\Leftrightarrow\)\SGARefLink{sga3:xvii:theorem:3:5:item:iv}{(iv)}), we immediately reduce to the
case in which \(U\) is commutative.  Moreover, \(E\to H\) has a
section by \SGARefLink{sga3:xvii:lemma:5:2:1:item:a}{5.2.1(a)}, so we are reduced to case \SGARefLink{sga3:xvii:appendix:A}{(a)}.

\medskip\noindent\textit{Case \SGARefLink{sga3:xvii:lemma:5:2:3:item:i:d}{(d)}.}\quad
We argue in the same way, using \SGARefLink{sga3:xvii:definition:5:1:0:item:ii}{5.1.0(ii)} and \SGARefLink{sga3:xvii:lemma:5:2:1:item:c}{5.2.1(c)}.

\medskip\noindent\textit{Case \SGARefLink{sga3:xvii:appendix:C}{(c)}.}\quad
Let \(E_{\mathrm{red}}\) be the reduced subgroup associated with
\(E\).  Since \(H\) is smooth, \(E_{\mathrm{red}}\) is an
extension of \(H\) by
\[
  U_0=U\cap E_{\mathrm{red}},
\]
and we may replace \(E\) by \(E_{\mathrm{red}}\).  The group \(U\),
and therefore \(U_0\), is connected, and \(U_0\) is plainly the
identity component of \(E_{\mathrm{red}}\); in particular it is
smooth.  Thus \(U_0\) is smooth and connected and \(k\) is
perfect, so \(U_0\) is \(k\)-solvable by \SGARefLink{sga3:xvii:remark:4:1:4:item:b}{4.1.4(b)}, reducing us to
case \SGARefLink{sga3:xvii:lemma:5:2:3:item:i:d}{(d)}.

The preceding reductions show that in cases \SGARefLink{sga3:xvii:section:b}{(b)}, \SGARefLink{sga3:xvii:appendix:C}{(c)}, and \SGARefLink{sga3:xvii:lemma:5:2:3:item:i:d}{(d)} we
may suppose that \(U\) has a characteristic composition series
\((U_i)\) such that \(U_i/U_{i+1}\) is commutative and the maps
\[
  U_i(k)\longrightarrow (U_i/U_{i+1})(k)
\]
are surjective.  (In cases \SGARefLink{sga3:xvii:appendix:C}{(c)} and \SGARefLink{sga3:xvii:lemma:5:2:3:item:i:d}{(d)}, the latter assertion
follows from \(H^1(k,\mathbb G_a)=0\).)  An immediate dévissage
then shows that it is enough to prove the conjugacy of two lifts
\(H'\) and \(H''\) in \(E\) when \(U\) is commutative.  In short,
it is enough to prove \SGARefLink{sga3:xvii:lemma:5:2:3:item:i}{(i)}\SGARefLink{sga3:xvii:lemma:5:2:4:item:a}{(a)}.  In that case the extension \(E\)
is trivial provided that \(H^2(H,U)=0\) (Appendix I, \SGARefLink{sga3:i:definition:3-1}{3.1}), and
\(H'\) and \(H''\) are conjugate provided that \(H^1(H,U)=0\).
We now use the following lemma.

\SGAStatement{Lemma 5.2.4}{Let \(S\) be a prescheme, \(U\) a
commutative \(S\)-group prescheme, and \(H\) an étale finite
\(S\)-group prescheme of rank \(n\) acting on \(U\).  Then the
groups \(H^i(H,U)\), for \(i>0\), are annihilated by \(n\) in
each of the following cases:
\begin{enumerate}
\item[(a)] \SGARefTarget{sga3:xvii:lemma:5:2:4:item:a}
  \(H\) is a constant \(S\)-group (\SGARefLink{sga3:i:subsection:4-1}{Exposé I, 4.1});
\item[(b)] \SGARefTarget{sga3:xvii:lemma:5:2:4:item:b}
  \(S\) is the spectrum of a field.
\end{enumerate}}
\SGARefTarget{sga3:xvii:lemma:5:2:4}

\SGAProof{Proof of \SGARefLink{sga3:xvii:appendix:A}{(a)}}
By hypothesis, \(H\) is the constant group associated with an
ordinary group \(\{H\}\) of order \(n\).  It is
clear\footnote{Editor's note: this is Proposition III.6.4.2 of
M. Demazure and P. Gabriel, \emph{Groupes algébriques I},
Masson \& North-Holland (1970).} that \(H^i(H,U)\) is isomorphic
to the \(i\)-th cohomology group
\[
  H^i(\{H\},U(S))
\]
of the group \(\{H\}\) with values in the ordinary group \(U(S)\).
It is classical that these groups are annihilated by \(n\)
(J.-P. Serre, \emph{Corps locaux}, Chapter VIII, \SGARefLink{sga3:viii:section:4}{Proposition 4},
Corollary 1).

\SGAProof{Proof of \SGARefLink{sga3:xvii:section:b}{(b)}}
Let \(x\in H^i(H,U)\), where \(i>0\), and represent it by an
\(i\)-cocycle
\[
  f:H^{(i)}\longrightarrow U,
\]
where \(H^{(i)}\) denotes the product over \(k\) of \(i\) copies
of \(H\).  If \(K\) is a finite Galois extension of \(k\) which
splits \(H\), assertion \SGARefLink{sga3:xvii:lemma:5:2:4:item:a}{(a)} shows that \(nf_K\) is a coboundary.
More precisely, a direct computation shows that, if one defines
the \(K\)-morphism
\[
  F_K:H_K^{(i-1)}\longrightarrow U_K
\]
by
\[
  F_K(h_1,\ldots,h_{i-1})
    =\sum_{h\in H(K)}
      f_K(h_1,\ldots,h_{i-1},h),
\]
then, up to sign,
\[
  d(F_K)=nf_K,
  \qquad d\text{ denoting the coboundary operator}.
\]
An immediate Galois argument shows that \(F_K\) descends to a
\(k\)-morphism \(F:H^{(i-1)}\to U\), and consequently
\[
  d(F)=nf.
\]

% END INCLUDED SGA3_English_current_buildable_source_20260823/source/inputs/XVII/tex/components/21_expose_XVII_subsection5_2_through_lemma5_2_4_pages022_023.tex

% BEGIN INCLUDED SGA3_English_current_buildable_source_20260823/source/inputs/XVII/tex/components/22_expose_XVII_corollary5_2_4_bis_through_proposition5_3_1_opening_page024.tex
% Authority: Expo17.pdf, local p.24 / combined p.986 /
% printed pp.579--580, from Corollary 5.2.4 bis through the
% opening inverse system in the proof of Proposition 5.3.1.

\SGAStatement{Corollary 5.2.4 bis}{With the notation of \SGARefLink{sga3:xvii:lemma:5:2:4}{5.2.4},
suppose in addition that \(U\) is flat and of finite presentation
over \(S\), with unipotent fibres, and that \(n\) is prime to the
residue characteristics of \(S\).  Then, in cases (a) and (b)
above,
\[
  H^i(H,U)=0\qquad(i>0).
\]}
\SGARefTarget{sga3:xvii:corollary:5:2:4:bis}

\SGAProof{Proof}
It is enough to show that the \(n\)-th power map on \(U\) is an
isomorphism, for it will then follow that multiplication by \(n\)
on \(H^i(H,U)\) is both an isomorphism and the zero morphism, and
hence that \(H^i(H,U)=0\).  Under the hypotheses made on \(U\), it
is enough to check that the \(n\)-th power map is an isomorphism
on the fibres of \(U\) (EGA IV, 17.9.5).  We are thus reduced to
the case in which \(S\) is the spectrum of a field \(k\) of
characteristic \(p\).  Since \((n,p)=1\), the \(n\)-th power map
on \(U\) is étale (Exposé VII); it is also a monomorphism (2.4(i)),
and hence an open immersion (EGA IV, 17.9.1).  This already shows
that its restriction to the identity component \(U^0\) is an
isomorphism.  Since \(U/U^0\) is a finite \(p\)-group, the proof
is complete.

This completes the proof of \SGARefLink{sga3:xvii:lemma:5:2:3:item:i}{5.2.3(i)}\SGARefLink{sga3:xvii:appendix:A}{(a)}, since an étale group of
multiplicative type \(H\) has order prime to \(p\).

\SGAProof{Proof of \SGARefLink{sga3:xvii:lemma:5:2:3:item:ii}{5.2.3(ii)}}
It is clear that \(R\) is a sheaf for the fpqc topology.
Moreover, in view of descent for affine schemes, the assertions
of \SGARefLink{sga3:xvii:lemma:5:2:3:item:ii}{5.2.3(ii)} are local for the fpqc topology.  We may therefore
suppose that \(k\) is algebraically closed.  The group \(H\) is
then split and the extension \(E\) is trivial by (i)\SGARefLink{sga3:xvii:section:b}{(b)}; let
\(H'\) be a lift of \(H\) in \(E\).  For every \(k\)-prescheme
\(S\) and every lift \(H''\) of \(H_S\) in \(E_S\), the groups
\(H'_S\) and \(H''\) are conjugate by an element of \(U(S)\),
since
\[
  H^1(H_S,U_S)=0
\]
by \SGARefLink{sga3:xvii:corollary:5:2:4:bis}{5.2.4 bis}.

Let \(U^{H'}\) be the sheaf of invariants of \(U\) under \(H'\).
It is representable by an algebraic subgroup of \(U\)
(\SGARefLink{sga3:viii:examples:6:5:item:d:01}{Exposé VIII, 6.5(d)}).  The preceding remarks show that the
\(k\)-morphism
\[
  U\longrightarrow R,
  \qquad
  u\longmapsto \operatorname{int}(u)H'
  \quad (u\in U(S))
\]
induces a \(k\)-isomorphism
\[
  U/U^{H'}\xrightarrow{\ \sim\ }R.
\]
This proves that \(R\) is representable and affine (\SGARefLink{sga3:xvii:proposition:2:1}{2.1}).

\SGAStatement{Remark 5.2.5}{The assertions of \SGARefLink{sga3:xvii:lemma:5:2:3:item:ii}{5.2.3(ii)} can also
be proved when \(U\) is not commutative, but we shall not need
this in order to prove \SGARefLink{sga3:xvii:theorem:5-1-1}{5.1.1}.}
\SGARefTarget{sga3:xvii:remark:5:2:5}

\SGASubsection{5.3}{\texorpdfstring{The Case Where \(H\) Is
Smooth}{The Case Where H Is Smooth}}
\SGARefTarget{sga3:xvii:subsection:5:3}

\SGAStatement{Proposition 5.3.1}{The assertions of \SGARefLink{sga3:xvii:lemma:5:2:3:item:i}{5.2.3(i)}
remain true when the hypothesis ``\(H\) is étale'' is replaced by
``\(H\) is smooth.''}
\SGARefTarget{sga3:xvii:proposition:5:3:1}

Proceeding as in the proof of \SGARefLink{sga3:xvii:lemma:5:2:3:item:i}{5.2.3(i)}, we reduce to the case in
which \(U\) is also commutative.

Let \(N'\) be the set of positive integers prime to \(p\), ordered
by divisibility.  For every \(n\in N'\), the group \({}_nH\) is
étale, and the family of groups \({}_nH\), \(n\in N'\), is
schematically dense in \(H\), since \(H\) is smooth
(\SGARefLink{sga3:ix:remark:4:10}{Exposé IX, 4.10}).  Let \(E_n\) be the inverse image of
\({}_nH\) in \(E\), so that \(E_n\) is an extension of
\({}_nH\) by \(U\), and let \(R_n\) be the \(k\)-functor of
lifts of \({}_nH\) in \(E_n\) (cf.\ \SGARefLink{sga3:xvii:lemma:5:2:3:item:ii}{5.2.3(ii)}).  If \(n\)
divides \(m\), there is plainly a natural \(k\)-morphism
\[
  R_m\longrightarrow R_n,
\]
so the \(R_n\) form an inverse system of \(k\)-functors.

% END INCLUDED SGA3_English_current_buildable_source_20260823/source/inputs/XVII/tex/components/22_expose_XVII_corollary5_2_4_bis_through_proposition5_3_1_opening_page024.tex

% BEGIN INCLUDED SGA3_English_current_buildable_source_20260823/source/inputs/XVII/tex/components/23_expose_XVII_proposition5_3_1_completion_and_proposition5_4_1_opening_page025.tex
% Authority: Expo17.pdf, local p.25 / combined p.987 /
% printed pp.581--582, completing Proposition 5.3.1 and beginning
% Proposition 5.4.1 through the Frobenius exact sequence.

By \SGARefLink{sga3:xvii:lemma:5:2:3:item:ii}{5.2.3(ii)}, \(R_n\) is representable by a nonempty affine
\(k\)-scheme.  Since a filtered direct limit of nonzero rings is
nonzero, the functor
\[
  R=\varprojlim_n R_n
\]
is representable by a nonempty affine \(k\)-scheme
(EGA IV, 8 and 1.9.1).  There are therefore an extension \(K\) of
\(k\) and a point \(u\in R(K)\).  The image \(u_n\) of \(u\) in
\(R_n(K)\) corresponds to a lift \(H'_n\) of \(({}_nH)_K\) in
\((E_n)_K\).  By construction,
\[
  H'_n={}_n(H'_m)\qquad\text{if \(n\) divides \(m\)}.
\]
Put
\[
  U_n=(U_K)^{H'_n}.
\]
The choice of \(H'_n\) identifies \((R_n)_K\) with \(U_K/U_n\).
The family of groups \(H'_n\) is filtered increasing, and hence
the family \(U_n\) is filtered decreasing.  It is therefore
stationary for sufficiently large \(n\), since \(U_K\) is
noetherian.  It follows that the family \((R_n)_K\) is stationary,
and consequently so is the family \(R_n\).  In short,
\[
  R_n=R
\]
for sufficiently large \(n\).

Under the hypotheses in force, \SGARefLink{sga3:xvii:lemma:5:2:3:item:i}{5.2.3(i)} shows that \(R_n(k)\) is
nonempty.  We may therefore find a coherent system of lifts
\(H'_n\) of \({}_nH\), for \(n\in N'\).  Let \(H'\) be the
smallest closed subscheme of \(E\) containing \(H'_n\) for every
\(n\) (Exposé VI\textsubscript{B}, \S7).  The argument of
\SGARefLink{sga3:xv:lemma:4:6}{Exposé XV, 4.6} shows that \(H'\) is a smooth commutative
algebraic group whose formation commutes with extension of the
base field.  To prove that \(H'\) is a lift of \(H\) in \(E\), we
may therefore suppose that \(k\) is algebraically closed.
According to BIBLE 4, Theorem 4, \(H'\) is then the direct product
of a smooth group \(M\) of multiplicative type and a unipotent
group \(V\).  The groups \(H'_n\) are necessarily contained in
\(M\) by \SGARefLink{sga3:xvii:proposition:2:4}{2.4}; the definition of \(H'\) then gives \(H'=M\).
Thus \(H'\) is of multiplicative type, and hence
\[
  H'\cap U=0.
\]
The morphism \(H'\to H\) is consequently a monomorphism, and the
density theorem of \SGARefLink{sga3:ix:remark:4:10}{Exposé IX, 4.10} shows that it is an
epimorphism.  It is therefore an isomorphism.

Now let \(H'\) and \(H''\) be two lifts of \(H\) in \(E\).  For
each \(n\in N'\), let
\[
  T_n=\operatorname{Transp}({}_nH',{}_nH'')
\]
be the transporter of \({}_nH'\) into \({}_nH''\).  It is
representable by a closed subscheme of \(U\)
(\SGARefLink{sga3:viii:examples:6:5:item:e:01}{Exposé VIII, 6.5(e)}).  The \(T_n\) form a filtered decreasing
family of closed subschemes of \(U\), and they are nonempty by
\SGARefLink{sga3:xvii:lemma:5:2:3:item:i}{5.2.3(i)}\SGARefLink{sga3:xvii:appendix:A}{(a)}.  Let \(T\) be their stationary value.  Under the
hypotheses of \SGARefLink{sga3:xvii:lemma:5:2:3:item:i}{5.2.3(i)}, \(T_n(k)\) is nonempty.  Thus there is an
element \(u\in U(k)\) such that
\[
  {}_nH''=\operatorname{int}(u)({}_nH')
  \qquad\text{for every \(n\in N'\)}.
\]
It follows that
\[
  H''=\operatorname{int}(u)H'
\]
by \SGARefLink{sga3:ix:corollary:4:8:item:b}{Exposé IX, 4.8(b)}.

\SGASubsection{5.4}{\texorpdfstring{The Case Where \(U\) Is
Radicial}{The Case Where U Is Radicial}}
\SGARefTarget{sga3:xvii:subsection:5:4}

\SGAStatement{Proposition 5.4.1}{If \(k\) is a perfect field of
characteristic \(p>0\), and if \(U\) is radicial, then the
extension \(E\) of \SGARefLink{sga3:xvii:theorem:5-1-1}{5.1.1} is trivial.}
\SGARefTarget{sga3:xvii:proposition:5:4:1}

Using a characteristic composition series of \(U\), we may
restrict ourselves to the case
\[
  U=(\boldsymbol\alpha_p)^r
\]
by \SGARefLink{sga3:xvii:lemma:3-9}{3.9}.

It follows from Appendix II, \SGARefLink{sga3:ii:subsection:2-2}{2.2} and \SGARefLink{sga3:xvii:proposition:2:1}{2.1} that there are
isomorphisms of \(k\)-functors
\[
  \operatorname{Aut}_{k\text{-}\mathrm{gr}}
    \bigl((\boldsymbol\alpha_p)^r\bigr)
  \xrightarrow{\ \sim\ }
  \operatorname{Aut}_{p\text{-}\mathrm{Lie}}
    \bigl(\operatorname{Lie}((\boldsymbol\alpha_p)^r)\bigr)
  \xrightarrow{\ \sim\ }
  \operatorname{GL}
    \bigl(\operatorname{Lie}((\boldsymbol\alpha_p)^r)\bigr).
\]
Let \(V\) be a \(k\)-vector space of rank \(r\), let
\(\mathbf W(V)\) be the corresponding vector \(k\)-group
(\SGARefLink{sga3:i:subsection:4-6}{Exposé I, 4.6}), and identify \((\boldsymbol\alpha_p)^r\) with
\({}_F V\).  The preceding remarks show that the action of \(H\)
on \({}_F V\) defined by the extension \(E\) comes from a linear
representation \(v\) of \(H\) on \(V\).  Consider the exact
sequence
\[
  \tag{*}\label{sga3:xvii:proposition:5-4-1:equation:star}
  0\longrightarrow {}_F V\longrightarrow V
    \xrightarrow{F}V^{(p)}\longrightarrow0,
\]
where \(F\) is the Frobenius morphism.

% END INCLUDED SGA3_English_current_buildable_source_20260823/source/inputs/XVII/tex/components/23_expose_XVII_proposition5_3_1_completion_and_proposition5_4_1_opening_page025.tex

% BEGIN INCLUDED SGA3_English_current_buildable_source_20260823/source/inputs/XVII/tex/components/24_expose_XVII_proposition5_4_1_completion_through_proposition5_6_1_statement_pages026_028.tex
% Authority: Expo17.pdf, local pp.26--28 / combined pp.988--990 /
% printed pp.583--587, completing Proposition 5.4.1, proving the
% relevant cases of Theorem 5.1.1(i), and giving the complete
% statement of Proposition 5.6.1.  The next cursor is its proof.

The sequence \((*)\) is an exact sequence of \(H\)-groups,
provided that \(H\) acts on \(V^{(p)}\) through the linear
representation
\[
  H\xrightarrow{v}\operatorname{GL}(V)
    \xrightarrow{F}\operatorname{GL}(V^{(p)}).
\]
Since \(k\) is perfect, the morphism \(F:V\to V^{(p)}\) induces
a surjective map
\[
  V(k)\longrightarrow V^{(p)}(k).
\]
It follows from Appendix I, \SGARefLink{sga3:i:subsection:2-1}{2.1} that the exact sequence \((*)\)
defines an exact sequence
\[
  \tag{**}\label{sga3:xvii:proposition:5-4-1:equation:double-star}
  H^1(H,V^{(p)})\longrightarrow
  \operatorname{Ext}_{\mathrm{alg}}(H,{}_F V)
  \longrightarrow
  \operatorname{Ext}_{\mathrm{alg}}(H,V).
\]

We shall show that
\[
  \operatorname{Ext}_{\mathrm{alg}}(H,V)=0.
\]
Let \(E'\) be an algebraic group which is an extension of \(H\)
by \(V\):
\[
  1\longrightarrow V\longrightarrow E'
    \xrightarrow{h}H\longrightarrow1.
\]
The scheme \(E'\) is a torsor over \(H\) under
\(\mathbb G_a^r\), and hence defines an element of
\[
  H^1(H,\mathcal O_H^r)
\]
in coherent-sheaf cohomology.  Since \(H\) is affine,
\[
  H^1(H,\mathcal O_H^r)=0
\]
by EGA III, \S1.  Thus \(E'\to H\) has a section.  Consequently,
\(\operatorname{Ext}_{\mathrm{alg}}(H,V)\) is isomorphic to
\(H^2(H,V)\) by Appendix I, \SGARefLink{sga3:i:definition:3-1}{3.1}.  But
\[
  H^i(H,V)=H^i(H,\mathbf W(V))=0\qquad(i>0)
\]
by \SGARefLink{sga3:ix:theorem:3:1}{Exposé IX, 3.1}.  The exact sequence \((**)\) now gives
\[
  \operatorname{Ext}_{\mathrm{alg}}(H,{}_F V)=0,
\]
and hence \(E\) is a trivial extension.

\SGASubsection{5.5}{Proof of \SGARefLink{sga3:xvii:theorem:5:1:1:item:i}{5.1.1(i)}}
\SGARefTarget{sga3:xvii:subsection:5:5}

If \(k\) has characteristic \(0\), then \(U\) is \(k\)-solvable
by \SGARefLink{sga3:xvii:corollary:4:1:3}{4.1.3} and \(H\) is smooth.  It follows from \SGARefLink{sga3:xvii:proposition:5:3:1}{5.3.1} and
\SGARefLink{sga3:xvii:lemma:5:2:3:item:i:d}{5.2.3(d)} that \(E\) is a trivial extension.  The same argument
shows that two lifts of \(H\) in \(E\) are conjugate by an
element of \(U(k)\).

Henceforth we suppose that \(k\) is a field of characteristic
\(p>0\).

\SGAProof{Proof of \SGARefLink{sga3:xvii:lemma:5:5:1:item:i}{(i)}\SGARefLink{sga3:xvii:section:b}{(b)}: the case where \(k\) is perfect and
\(U\) is connected}
We shall reduce to the case in which \(U\) is radicial.  Since
\(k\) is perfect, \(H_{\mathrm{red}}\) is smooth; let \(E'\) be
its inverse image in \(E\).  By \SGARefLink{sga3:xvii:proposition:5:3:1}{5.3.1} and \SGARefLink{sga3:xvii:lemma:5:2:3:item:i}{5.2.3(i)}\SGARefLink{sga3:xvii:appendix:C}{(c)}, the
extension
\[
  1\longrightarrow U\longrightarrow E'
    \longrightarrow H_{\mathrm{red}}\longrightarrow1
\]
is trivial.  Let \(H_1\) be a lift of \(H_{\mathrm{red}}\) in
\(E\).  By Appendix II, \SGARefLink{sga3:xvii:lemma:3-1}{3.1}, there is an integer \(n>0\) such
that
\[
  E^{(n)}=E/{}_{F^n}(E)
\]
is smooth.  Let \(E''\) be the algebraic subgroup of \(E\)
generated by \(H_1\) and \({}_{F^n}(E)\); equivalently, \(E''\)
is the inverse image in \(E\) of the image of \(H_1\) in
\(E^{(n)}\).  Let \(H''\) be the image of \(E''\) in \(H\).  We
claim that \(H''=H\).

Indeed, let \(R\) be the image of \({}_{F^n}(E)\) in \(H\), so
that \(H''\) is generated by \(R\) and \(H_{\mathrm{red}}\).
The group \(H/R\) is a quotient of \(E^{(n)}\), and is therefore
smooth.  Consequently the canonical morphism
\[
  H_{\mathrm{red}}\longrightarrow H/R
\]
is an epimorphism.  Thus \(H\) is generated by \(R\) and
\(H_{\mathrm{red}}\), and hence equals \(H''\).  We have obtained
an exact sequence
\[
  \tag{\(\dagger\)}\label{sga3:xvii:subsection:5-5:equation:dagger}
  1\longrightarrow U''=U\cap E''
    \longrightarrow E''\longrightarrow H\longrightarrow1.
\]
But \(E''\) has the same underlying space as \(H_1\), so \(U''\)
is radicial.  It is plainly enough to prove that the extension
\((\dagger)\) is trivial, and this follows from \SGARefLink{sga3:xvii:proposition:5:4:1}{5.4.1}.

\SGAProof{Proof of \SGARefLink{sga3:xvii:lemma:5:5:1:item:i}{(i)}\SGARefLink{sga3:xvii:section:b}{(b)}: the case where \(k\) is perfect and
\(H\) is connected}
The group \(U\) is an extension of an étale group by its identity
component \(U^0\).  The case where \(U\) is connected having just
been treated, it is enough to consider the case where \(U\) is
étale.  We shall use the following more precise lemma.

\SGAStatement{Lemma 5.5.1}{With the notation of \SGARefLink{sga3:xvii:theorem:5-1-1}{5.1.1}, suppose
in addition that \(U\) is étale.  Then:
\begin{enumerate}
\item[(i)] \SGARefTarget{sga3:xvii:lemma:5:5:1:item:i}
  If \(H\) is connected, it has a unique lift in \(E\), namely
  the identity component \(E^0\) of \(E\).
\item[(ii)] \SGARefTarget{sga3:xvii:lemma:5:5:1:item:ii}
  If \(k\) is algebraically closed, then \(E\) is trivial, and
  any two lifts of \(H\) in \(E\) are conjugate by an element of
  \(U(k)\).
\end{enumerate}}
\SGARefTarget{sga3:xvii:lemma:5:5:1}

\SGAProof{Proof of \SGARefLink{sga3:xvii:lemma:5:5:1:item:i}{(i)}}
Since formation of the identity component commutes with extension
of the base field, we may restrict ourselves to the case where
\(k\) is algebraically closed.  If \(H\) is a torus, then \(E\)
is trivial by \SGARefLink{sga3:xvii:proposition:5:3:1}{5.3.1} and \SGARefLink{sga3:xvii:lemma:5:2:3:item:i}{5.2.3(i)}\SGARefLink{sga3:xvii:section:b}{(b)}, and \(E^0\) is plainly the
unique lift of \(H\).  This already proves that in the general
case \(E^0\cap U\) is radicial.  It is also étale, since \(U\) is
étale, and is therefore zero.  The morphism
\[
  E^0\longrightarrow H
\]
is consequently a monomorphism; it is flat because \(E^0\) is
open in \(E\), and it is surjective because \(H\) is connected.
It is therefore an isomorphism.  If \(H'\) is any other lift of
\(H\), then \(H'\) is connected, hence contained in \(E^0\), and
thus equal to \(E^0\).

\SGAProof{Proof of \SGARefLink{sga3:xvii:lemma:5:5:1:item:ii}{(ii)}}
Let \(H^0\) be the identity component of \(H\).  By \SGARefLink{sga3:xvii:lemma:5:5:1:item:i}{(i)}, \(E^0\)
is the unique lift of \(H^0\) in \(E\).  Put
\[
  E'=E/E^0,\qquad H'=H/H^0.
\]
There is then an extension
\[
  1\longrightarrow U\longrightarrow E'
    \longrightarrow H'\longrightarrow1.
\]
Since \(H'\) is étale, this extension is trivial by
\SGARefLink{sga3:xvii:lemma:5:2:3:item:i}{5.2.3(i)}\SGARefLink{sga3:xvii:section:b}{(b)}.  If \(H'_1\) is a lift of \(H'\) in \(E'\), and
\(H_1\) is its inverse image in \(E\), then \(H_1\) plainly
lifts \(H\) in \(E\).  If \(H_2\) is a second lift of \(H\) in
\(E\), then it contains \(E^0\).  By \SGARefLink{sga3:xvii:lemma:5:2:3:item:i}{5.2.3(i)}\SGARefLink{sga3:xvii:section:b}{(b)}, the image of
\(H_2\) in \(E'\) is conjugate to \(H'_1\) by an element
\(u\in U(k)\), and hence
\[
  H_2=\operatorname{int}(u)H_1.
\]

\SGAStatement{Remark 5.5.2}{Under the hypotheses of \SGARefLink{sga3:xvii:lemma:5:5:1:item:i}{5.5.1(i)},
the group \(E^0\) centralizes \(U\).}
\SGARefTarget{sga3:xvii:remark:5:5:2}

\SGAProof{Proof of \SGARefLink{sga3:xvii:theorem:5:1:1:item:i}{5.1.1(i)}\SGARefLink{sga3:xvii:appendix:A}{(a)}}
Using the composition series
\[
  1\longrightarrow U^0\longrightarrow U
    \longrightarrow U/U^0\longrightarrow1,
\]
assertion \SGARefLink{sga3:xvii:lemma:5:5:1:item:i}{(i)}\SGARefLink{sga3:xvii:appendix:A}{(a)} follows by combining \SGARefLink{sga3:xvii:lemma:5:5:1:item:i}{(i)}\SGARefLink{sga3:xvii:section:b}{(b)} with \SGARefLink{sga3:xvii:lemma:5:5:1:item:ii}{5.5.1(ii)}.

Before proving \SGARefLink{sga3:xvii:theorem:5:1:1:item:i}{5.1.1(i)}\SGARefLink{sga3:xvii:appendix:C}{(c)} and (d), we shall first establish
\SGARefLink{sga3:xvii:theorem:5:1:1:item:ii}{5.1.1(ii)}.

\SGAProof{Proof of \SGARefLink{sga3:xvii:theorem:5:1:1:item:ii}{5.1.1(ii)}\SGARefLink{sga3:xvii:appendix:A}{(a)}}
For want of a satisfactory general statement, we shall describe
a number of cases in which, when \(H\) is smooth, two lifts of
\(H\) in \(E\) are conjugate.

\SGAStatement{\SGAInlineRefTarget{sga3:xvii:proposition:5-6-1}Proposition 5.6.1}{With the notation of \SGARefLink{sga3:xvii:theorem:5-1-1}{5.1.1},
suppose in addition that \(H\) is smooth, and let \(H'\) and
\(H''\) be two lifts of \(H\) in \(E\).  Then \(H'\) and \(H''\)
are conjugate by an element of \(U(k)\) in each of the following
cases:
\begin{enumerate}
\item[(a)] \SGARefTarget{sga3:xvii:proposition:5:6:1:item:a}
  \(k\) is algebraically closed.
\item[(b)] \SGARefTarget{sga3:xvii:proposition:5:6:1:item:b}
  \(U\) is commutative.
\item[(c)] \SGARefTarget{sga3:xvii:proposition:5:6:1:item:c}
  \(k\) is perfect and \(U\) is connected.
\item[(d)] \SGARefTarget{sga3:xvii:proposition:5:6:1:item:d}
  \(U\) is \(k\)-solvable.
\item[(e)] \SGARefTarget{sga3:xvii:proposition:5:6:1:item:e}
  The centralizer of \(H'\) in \(U\) is \(k\)-solvable.
\item[(f)] \SGARefTarget{sga3:xvii:proposition:5:6:1:item:f}
  The group of multiplicative type \(H\) is split by a finite
  Galois extension \(K/k\) of degree prime to \(p\).
\end{enumerate}}
\SGARefTarget{sga3:xvii:proposition:5:6:1}

% END INCLUDED SGA3_English_current_buildable_source_20260823/source/inputs/XVII/tex/components/24_expose_XVII_proposition5_4_1_completion_through_proposition5_6_1_statement_pages026_028.tex

% BEGIN INCLUDED SGA3_English_current_buildable_source_20260823/source/inputs/XVII/tex/components/25_expose_XVII_proposition5_6_1_proof_and_section5_7_pages028_031.tex
% Authority: Expo17.pdf, local pp.28--31 / combined pp.990--993 /
% printed pp.587--593, completing Proposition 5.6.1 and all of
% subsection 5.7.  The next cursor is subsection 5.8.

\SGAProof{Proof}
Parts \SGARefLink{sga3:xvii:appendix:A}{(a)}, \SGARefLink{sga3:xvii:section:b}{(b)}, \SGARefLink{sga3:xvii:appendix:C}{(c)}, and (d) follow from \SGARefLink{sga3:xvii:proposition:5:3:1}{5.3.1}.

\smallskip
\noindent
(e) Let \(Z\) be the centralizer of \(H'\) in \(U\), and let \(T\)
be the transporter of \(H'\) into \(H''\).  By \SGARefLink{sga3:xvii:appendix:A}{(a)}, \(T\) is
nonempty; hence \(T\) is a principal homogeneous space under \(Z\).
The hypothesis on \(Z\) therefore implies that \(T\) is trivial.

\smallskip
\noindent
(f) Proceeding as in \SGARefLink{sga3:xvii:proposition:5:3:1}{5.3.1}, we see that it is enough to consider the
case where \(H\) is étale.  Suppose first that \(H\) is
diagonalizable, defined by an abstract group \(M\) of order \(m\)
prime to \(p\).  The two lifts \(H'\) and \(H''\) determine a
one-cocycle of \(M\) with values in \(U(k)\), that is, a map
\[
  h:M\longrightarrow U(k)
\]
such that
\[
  h(mn)=h(m)\,{}^{m}\!h(n)
\]
for every pair \(m,n\) of elements of \(M\).  The groups \(H'\) and
\(H''\) are conjugate by an element of \(U(k)\) if and only if there
exists \(a\in U(k)\) such that
\[
  h(m)=a^{-1}({}^{m}\!a).
\]
Now the abstract group \(U(k)\) has a composition series whose
successive quotients are commutative and annihilated by a power of
\(p\) (in view of \SGARefLink{sga3:xvii:proposition:5:6:1:item:c}{5.6.1(c)}, we may suppose that \(p>0\)).  It follows
at once in this case that \(h\) is a coboundary.

Consider now the general case.  Again let \(Z\) denote the
centralizer of \(H'\) in \(U\), and \(T\) the transporter of \(H'\)
into \(H''\); thus \(T\) is a torsor under \(Z\).  By hypothesis there
is a finite Galois extension \(K/k\), with Galois group \(G\) of
order prime to \(p\), which splits \(H\).  By the preceding case,
\(H'_K\) and \(H''_K\) are conjugate by an element of \(U(K)\), and
therefore \(T_K\) is trivial.  Consequently \(T\) is defined by an
element of
\[
  H^1(G,Z).
\]
For the same reasons as above, the hypothesis on \(G\) implies that
\[
  H^1(G,Z)=e;
\]
hence \(T\) is trivial.

\SGASubsection{5.7}{Proof of \SGARefLink{sga3:xvii:theorem:5:1:1:item:ii}{5.1.1(ii)}\SGARefLink{sga3:xvii:section:b}{(b)}}
\SGARefTarget{sga3:xvii:subsection:5:7}

\SGAStatement{Lemma 5.7.1}{Let \(S\) be a prescheme, let \(G\) be an
\(S\)-group prescheme, separated and smooth over \(S\), and let \(H\)
be an \(S\)-group of multiplicative type acting on \(G\).  Then the
\(S\)-functor
\[
  Z=G^H
\]
of invariants of \(G\) under \(H\) is representable by a closed
smooth \(S\)-group subprescheme of \(G\).}
\SGARefTarget{sga3:xvii:lemma:5:7:1}

The representability of \(G^H\) by a closed group subprescheme of
\(G\) follows from \SGARefLink{sga3:viii:examples:6:5:item:d:01}{Exposé VIII, 6.5(d)}.  Consider the semidirect
product
\[
  K=G\rtimes_S H.
\]
The centralizer of \(H\) in \(K\) is then \(Z\times_S H\).  To prove
that \(Z\) is smooth, we must verify the following.  Suppose that
\(S\) is affine, that \(S_0\) is a closed subscheme of \(S\) defined
by a square-zero ideal, and that \(u_0\in Z(S_0)\).  We must lift
\(u_0\) to an element of \(Z(S)\).

Since \(G\) is smooth over \(S\), there is an element \(u\in G(S)\)
lifting \(u_0\).  Let \(i\) be the canonical immersion of \(H\) into
\(K\), and consider the two \(S\)-group morphisms
\[
  H\rightrightarrows K,\qquad i
  \quad\hbox{and}\quad
  j=\operatorname{int}(u)i.
\]
Since \(u_0\in Z(S_0)\), we have \(i_{S_0}=j_{S_0}\).  By
\SGARefLink{sga3:ix:theorem:3:2}{Exposé IX, 3.2}, there exists \(v\in K(S)\), lifting the identity
section in \(K(S_0)\), such that
\[
  i=\operatorname{int}(v)j
   =\operatorname{int}(vu)i.
\]
Thus
\[
  vu=(u',v')\in (Z\times_S H)(S).
\]
It follows that \(u'\in Z(S)\), and \(u'\) lifts \(u_0\).

\SGAStatement{Lemma 5.7.2}{Let \(k\) be a field and \(H\) a
\(k\)-algebraic group, and let \(P(H)\) denote the following
property:
\begin{quote}
For every smooth unipotent \(k\)-group \(U\), and every extension
\(E\) of \(H\) by \(U\), any two lifts of \(H\) in \(E\) are
conjugate by an element of \(U(k)\).
\end{quote}
Suppose that \(H\) is an algebraic group which is an extension of a
group \(H''\) of multiplicative type by a group \(H'\) of
multiplicative type.  If \(P(H')\) and \(P(H'')\) hold, then
\(P(H)\) holds.}
\SGARefTarget{sga3:xvii:lemma:5:7:2}

Indeed, let \(U\) be a smooth unipotent \(k\)-group, let \(E\) be an
extension of \(H\) by \(U\), and let \(H_1,H_2\) be two lifts of
\(H\) in \(E\).  Let \(H'_1,H'_2\) be the corresponding lifts of
\(H'\).  Since \(P(H')\) holds, there exists \(u\in U(k)\) such that
\[
  H'_2=\operatorname{int}(u)H'_1.
\]
After replacing \(H_1\) by \(\operatorname{int}(u)H_1\), we may
suppose that \(H'_1=H'_2\); denote this common group simply by
\(H'\).  Put
\[
  E'=\operatorname{Centr}_E(H')
     =U^{H'}H_1
     =U^{H'}H_2.
\]
By \SGARefLink{sga3:xvii:lemma:5:7:1}{5.7.1}, \(U^{H'}=U'\) is smooth.  We then have an extension
\[
  1\longrightarrow U'\longrightarrow E'
    \longrightarrow H\longrightarrow1.
\]
By construction \(H'\) is central in \(E'\), and hence invariant.
Passing to the quotient gives an exact sequence
\[
  1\longrightarrow U'\longrightarrow E''
    \longrightarrow H''\longrightarrow1.
\]
Since \(U'\) is smooth and \(P(H'')\) holds, the images of \(H_1\)
and \(H_2\) in \(E''\) are conjugate by an element
\(u\in U'(k)\).  It follows that
\[
  H_1=\operatorname{int}(u)H_2.
\]

To prove \SGARefLink{sga3:xvii:theorem:5:1:1:item:ii}{5.1.1(ii)}\SGARefLink{sga3:xvii:section:b}{(b)}, observe that, since \(k\) is algebraically
closed, \(H\) has a composition series whose successive quotients
are smooth or, when \(p>0\), isomorphic to \(\boldsymbol\mu_p\).
Repeated application of \SGARefLink{sga3:xvii:lemma:5:7:2}{5.7.2} reduces us to the case where \(H\) is
smooth or \(H=\boldsymbol\mu_p\).  In the first case it is enough to
apply \SGARefLink{sga3:xvii:theorem:5:1:1:item:ii}{5.1.1(ii)}(a).  It remains to consider \(H=\boldsymbol\mu_p\).
Since \(U\) is smooth, \(U\) has a characteristic composition series
whose successive quotients are étale or isomorphic to
\(\mathbb G_a^r\), by \SGARefLink{sga3:xvii:lemma:3-9}{3.9}.  If \(U\) is étale, apply \SGARefLink{sga3:xvii:lemma:5:5:1}{5.5.1}.  We are
finally reduced to
\[
  H=\boldsymbol\mu_p,\qquad U=\mathbb G_a^r.
\]

We must prove that
\[
  H^1(\boldsymbol\mu_p,U)=0.
\]
The method used in \SGARefLink{sga3:xvii:proposition:5:4:1}{5.4.1} no longer applies, because
\(\boldsymbol\mu_p\) does not in general act linearly on
\(\mathbb G_a^r\).  Fix the following notation: \(H'\) denotes a
lift of \(H\) in \(E\), and
\[
  \mathfrak e=\operatorname{Lie}E,\qquad
  \mathfrak u=\operatorname{Lie}U,\qquad
  \mathfrak h=\operatorname{Lie}H,\qquad
  \mathfrak h'=\operatorname{Lie}H'.
\]
Let \(X\) be a nonzero element of \(\mathfrak h\) such that
\[
  X^{(p)}=X
\]
(Appendix II, \SGARefLink{sga3:xvii:lemma:3-1}{3.1}), and let \(X'\) be its lift in \(\mathfrak h'\).

Since \(\boldsymbol\mu_p\) is a radicial group of height \(1\), the
set of lifts of \(H\) in \(E\) is in bijective correspondence with
the set \(A\) of elements \(Y\in\mathfrak e\) such that
\[
  Y^{(p)}=Y
\]
and whose image is \(X\) (Appendix II, \SGARefLink{sga3:ii:subsection:2-2}{2.2}).  Likewise, if
\(Y\in A\) corresponds to a lift \(H''\) of \(H\) in \(E\), and if
\(u\in U(k)\), then
\[
  \operatorname{int}(u)H'=H''
  \quad\Longleftrightarrow\quad
  \operatorname{Ad}(u)X'=Y.
\]
Let \(B\) be the subset of \(\mathfrak e\) consisting of the
elements \(\operatorname{Ad}(u)X'\), with \(u\in U(k)\).  Plainly
\(B\subset A\), and it remains only to prove \(A=B\) when \(k\) is
algebraically closed.

\smallskip
\noindent
(a) \emph{Study of \(A\).}
Since \(\mathfrak u\) is commutative and normalized by
\(\mathfrak h'\), Jacobson's formula (Exposé VII\textsubscript{A},
\SGARefLink{sga3:VIIB:5.2:theorem}{5.2}) gives, for \(u\in\mathfrak u\),
\[
  (X'+u)^{(p)}
   =X'{}^{(p)}+u^{(p)}
      +(\operatorname{ad}X')^{p-1}(u)
   =X'+(\operatorname{ad}X')^{p-1}(u).
\]
Consequently,
\[
  X'+u\in A
  \quad\Longleftrightarrow\quad
  u=(\operatorname{ad}X')^{p-1}(u).
\]
Let
\[
  \mathfrak u
   =\bigoplus_{n\in\mathbb Z/p\mathbb Z}\mathfrak u_n
   =\mathfrak u_0\oplus
      \bigoplus_{n\in(\mathbb Z/p\mathbb Z)^\times}\mathfrak u_n
\]
be the canonical decomposition of \(\mathfrak u\) under the action
of \(H'\).  If \(u\in\mathfrak u_0\), then
\(\operatorname{ad}X'(u)=0\).  If
\(u\in\mathfrak u_n\), with
\(n\in(\mathbb Z/p\mathbb Z)^\times\), then
\[
  (\operatorname{ad}X')^{p-1}(u)=u.
\]
Thus \(Y=X'+u\) belongs to \(A\) if and only if
\[
  u\in
  \bigoplus_{n\in(\mathbb Z/p\mathbb Z)^\times}\mathfrak u_n.
\]
In particular, \(A\) is the set of \(k\)-rational points of an
irreducible subscheme of \(\mathbf W(\mathfrak e)\), of dimension
\[
  \operatorname{rk}\mathfrak u-\operatorname{rk}\mathfrak u_0
  =
  \operatorname{rk}\mathfrak u
  -\operatorname{rk}\operatorname{Centr}_{\mathfrak u}(X').
\]

\smallskip
\noindent
(b) \emph{Study of \(B\).}
We shall need the following lemma.

\SGAStatement{Lemma 5.7.3 (Rosenlicht)}{Let \(U\) be a unipotent
algebraic group over a field \(k\), acting on a quasi-affine
\(k\)-scheme \(X\).  Then the orbit of every point \(x\in X(k)\) is
closed in \(X\).  Here, by the orbit of \(x\), we mean the subset of
\(\operatorname{ens}(X)\) which is the image of \(U\) under the
morphism \(g\mapsto g\cdot x\).}
\SGARefTarget{sga3:xvii:lemma:5:7:3}

By functoriality, \(U\) acts on the affine envelope of \(X\), that
is,
\[
  \operatorname{Spec}\Gamma(X,\mathcal O_X),
\]
so we may suppose that \(X\) is affine.  We may then suppose that
\(k\) is algebraically closed, \(X\) is reduced, and \(U\) is smooth
(observe that \(U_{\mathrm{red}}\) acts on \(X_{\mathrm{red}}\) when
\(k\) is perfect).  Let \(Y\) be the schematic image of \(U\)
(EGA I, 9.5.1) under the morphism \(g\mapsto g\cdot x\).  It is a
closed reduced subscheme of \(X\), on which \(U\) acts.  It follows
readily from EGA IV, 1.8.6 that the orbit of \(x\) is an open subset
\(Z\) of \(Y\), dense in \(Y\).  We must prove that
\[
  Z=\operatorname{ens}(Y).
\]
Let \(F\) be the closed reduced subscheme of \(Y\) whose underlying
space is \(Y\setminus Z\).  Thus \(F=V(J)\), where \(J\) is a
nonzero ideal of \(\Gamma(Y,\mathcal O_Y)\).  Since \(U\) is smooth,
\(U\) acts on \(F\), hence on \(J\), and therefore
\[
  J^U\ne0
\]
by \SGARefLink{sga3:xvii:proposition:3:2}{3.2}.  If \(a\) is a nonzero element of \(J^U\), then \(a\) is
constant on the orbit \(Z\), and hence on \(Y\), because \(Z\) is
dense in \(Y\).  The ideal \(J\) consequently contains \(k\), and
\(F=\varnothing\).

Apply the preceding lemma to the group \(U\) acting on the affine
space \(\mathbf W(\mathfrak e)\) through the adjoint representation.
The orbit of \(X'\) is the underlying set of a closed subscheme of
\(\mathbf W(\mathfrak e)\).  Moreover, the stabilizer \(Z\) of
\(X'\) is its centralizer in \(U\), and there is a closed immersion
\[
  U/Z\longrightarrow\mathbf W(\mathfrak e).
\]
Thus the orbit of \(X'\) is the underlying space of a closed
subscheme of \(\mathbf W(\mathfrak e)\) of dimension
\[
  \dim U-\dim Z.
\]

\smallskip
\noindent
\SGARefLink{sga3:xvii:appendix:C}{(c)} \emph{End of the proof of \SGARefLink{sga3:xvii:theorem:5:1:1:item:ii}{5.1.1(ii)}\SGARefLink{sga3:xvii:section:b}{(b)}.}
When \(k\) is algebraically closed, the canonical map
\[
  U(k)\longrightarrow(U/Z)(k)
\]
is surjective.  By \SGARefLink{sga3:xvii:section:b}{(b)}, \(B\) is therefore the set of \(k\)-rational
points of a closed subscheme of \(\mathbf W(\mathfrak e)\) of
dimension \(\dim U-\dim Z\).  In view of \SGARefLink{sga3:xvii:appendix:A}{(a)}, to prove \(A=B\) it is
enough to show that
\[
  \operatorname{rk}\mathfrak u
  -\operatorname{rk}\operatorname{Centr}_{\mathfrak u}(X')
  \leq
  \dim U-\dim\operatorname{Centr}_U(X').
\]
Since \(U\) is smooth,
\[
  \dim U=\operatorname{rk}\mathfrak u.
\]
On the other hand, \SGARefLink{sga3:ii:paragraph:5-3-3}{Exposé II, 5.3.3} gives
\[
  \dim\operatorname{Centr}_U(X')
  \leq
  \operatorname{rk}\operatorname{Centr}_{\mathfrak u}(X'),
\]
which proves the assertion.  Notice in fact that
\[
  \dim\operatorname{Centr}_U(X')
  =
  \operatorname{rk}\operatorname{Centr}_{\mathfrak u}(X');
\]
this also proves once more that \(Z\) is smooth, as in \SGARefLink{sga3:xvii:lemma:5:7:1}{5.7.1}.

% END INCLUDED SGA3_English_current_buildable_source_20260823/source/inputs/XVII/tex/components/25_expose_XVII_proposition5_6_1_proof_and_section5_7_pages028_031.tex

% BEGIN INCLUDED SGA3_English_current_buildable_source_20260823/source/inputs/XVII/tex/components/26_expose_XVII_section5_8_pages031_032.tex
% Authority: Expo17.pdf, local pp.31--32 / combined pp.993--994 /
% printed pp.593--597, all of subsection 5.8.  The next cursor is
% subsection 5.9.

\SGASubsection{5.8}{End of the proof of \SGARefLink{sga3:xvii:theorem:5:1:1:item:i}{5.1.1(i)}}
\SGARefTarget{sga3:xvii:subsection:5:8}

It remains to prove (i)\SGARefLink{sga3:xvii:appendix:C}{(c)} and (i)(d).

\SGAProof{Proof of \SGARefLink{sga3:xvii:theorem:5:1:1:item:i}{5.1.1(i)}(d): \(U\) smooth and \(H\) connected}
We shall need the following lemma.

\SGAStatement{Lemma 5.8.1}{With the notation of \SGARefLink{sga3:xvii:theorem:5-1-1}{5.1.1}, suppose that
\(U\) is smooth and \(H\) is radicial, and let \(H_1\) be an
algebraic subgroup of \(H\) which has a lift \(H'_1\) in \(E\).
Then \(H\) has a lift \(H'\) in \(E\) which contains \(H'_1\).}
\SGARefTarget{sga3:xvii:lemma:5:8:1}

\SGAProof{Proof}
By induction on the height of \(H/H_1\), we may suppose that
\(H/H_1\) has height \(1\).  Put
\[
  C'=\operatorname{Centr}_E(H'_1).
\]
We claim that the canonical morphism \(C'\longrightarrow H\) is an
epimorphism.  To prove this, we may suppose that \(k\) is
algebraically closed.  The extension \(E\) is then trivial, by
\SGARefLink{sga3:xvii:theorem:5:1:1:item:i}{5.1.1(i)}\SGARefLink{sga3:xvii:appendix:A}{(a)}.  Let \(H''\) be a lift of \(H\) in \(E\), and let
\(H''_1\) be the inverse image of \(H_1\) in \(H''\).  The groups
\(H'_1\) and \(H''_1\) are two lifts of \(H_1\) in \(E\); hence they
are conjugate by an element of \(U(k)\), since \(k\) is
algebraically closed and \(U\) is smooth, by \SGARefLink{sga3:xvii:theorem:5:1:1:item:ii}{5.1.1(ii)}\SGARefLink{sga3:xvii:section:b}{(b)}.
Consequently, to prove the assertion about \(C'\), it is enough to
prove it for
\[
  C''=\operatorname{Centr}_E(H''_1).
\]
But \(C''\) contains \(H''\), so the assertion is clear.

It also follows from \SGARefLink{sga3:xvii:lemma:5:7:1}{5.7.1} that \(C'\cap U\) is smooth.  We may
therefore replace \(E\) by \(C'\), and thus suppose that \(H'_1\) is
central.  Passing then to the quotient by \(H'_1\), we may suppose
that \(H_1=0\) and that \(H\) has height \(1\).

Since \(U\) is smooth, we have the following exact sequence of
\(p\)-Lie algebras (Appendix II, \SGARefLink{sga3:xvii:proposition:3:2}{3.2}):
\[
  (*)\qquad
  0\longrightarrow\operatorname{Lie}U
   \longrightarrow\operatorname{Lie}E
   \longrightarrow\operatorname{Lie}H
   \longrightarrow0.
\]
In view of Appendix II, \SGARefLink{sga3:ii:subsection:2-2}{2.2}, to say that \(E\) is trivial is
equivalent to saying that \((*)\) is a trivial extension of
\(p\)-Lie algebras.  Suppose that \(H\ne0\), and hence that
\(\operatorname{Lie}H\ne0\).  Suppose also that we have found a
nonzero \(p\)-Lie subalgebra
\[
  \mathfrak h_1\subset\mathfrak h=\operatorname{Lie}H
\]
which lifts to a \(p\)-Lie subalgebra
\(\mathfrak h'_1\subset\operatorname{Lie}E\).  By the cited result,
there is a subgroup \(H_1\) of \(H\), with
\(\operatorname{Lie}H_1=\mathfrak h_1\), and a lift \(H'_1\) of
\(H_1\) in \(E\), with
\(\operatorname{Lie}H'_1=\mathfrak h'_1\).  Applying once more the
reduction just described, we are brought back to the same problem
with \(H\) replaced by \(H/H_1\).  Since \(H\) has height \(1\),
\[
  \operatorname{Lie}(H/H_1)
   =\operatorname{Lie}H/\operatorname{Lie}H_1
\]
(loc.\ cit.), and consequently
\[
  \operatorname{rk}\operatorname{Lie}(H/H_1)
  \leq \operatorname{rk}\operatorname{Lie}H-1.
\]
Thus, by induction on the rank of \(\operatorname{Lie}H\), it is
enough, when \(\mathfrak h\ne0\), to find a nonzero Lie subalgebra
of \(\operatorname{Lie}H\) which lifts to \(\operatorname{Lie}E\).
This follows from the next lemma.

\SGAStatement{\SGAInlineRefTarget{sga3:xvii:lemma:5-8-2}Lemma 5.8.2}{Let \(k\) be a field of characteristic
\(p>0\), and let
\[
  \varphi:\mathfrak g\longrightarrow\mathfrak h
\]
be a surjective morphism of finite-rank \(p\)-Lie algebras over
\(k\).  Then:
\begin{enumerate}
\item[\SGAInlineRefTarget{sga3:xvii:lemma:5-8-2:item:i}(i)] If \(\mathfrak h_{\overline{k}}\) is reductive
  (\SGARefLink{sga3:xvii:lemma:4-2-2}{4.2.2}) and nonzero, there exists a reductive Lie subalgebra
  \(\mathfrak h_1\) of \(\mathfrak g\) whose image in
  \(\mathfrak h\) is nonzero.
\item[\SGAInlineRefTarget{sga3:xvii:lemma:5-8-2:item:ii}(ii)] If \(k\) is perfect and \(u\) is a unipotent element of
  \(\mathfrak h\), that is, if \(u^{(p^n)}=0\) for some \(n\), then
  \(u\) lifts to a unipotent element of \(\mathfrak g\).
\end{enumerate}}
\SGARefTarget{sga3:xvii:lemma:5:8:2}

\SGAProof{Proof}
Let \(X\) be a nonzero element of \(\mathfrak h\) in case \SGARefLink{sga3:xvii:lemma:5-8-2:item:i}{(i)}, and
let \(X=u\) in case \SGARefLink{sga3:xvii:lemma:5-8-2:item:ii}{(ii)}; choose a lift \(X'\) of \(X\) in
\(\mathfrak g\).  The \(p\)-Lie subalgebra \(\mathfrak h'\) of
\(\mathfrak g\) generated by \(X'\) is abelian (Exposé VII).

\smallskip
\noindent
\emph{Case \SGARefLink{sga3:xvii:lemma:5-8-2:item:i}{(i)}.}
The description given in \SGARefLink{sga3:xvii:lemma:4-2-2}{4.2.2} shows that the reductive part
(loc.\ cit.) of \(\mathfrak h'_{\overline{k}}\) is already defined
over \(k\); denote it by \(\mathfrak r'\).  We claim that the image
of \(\mathfrak r'\) in \(\mathfrak h\) is nonzero.  We may suppose
that \(k\) is algebraically closed, so that
\[
  \mathfrak h'=\mathfrak r'\oplus\mathfrak u',
\]
where \(\mathfrak u'\) is the unipotent part of
\(\mathfrak h'\).  If the image of \(\mathfrak r'\) in
\(\mathfrak h\) were zero, then the image of \(\mathfrak h'\) in
\(\mathfrak h\) would be unipotent, and hence zero because
\(\mathfrak h\) is reductive (\SGARefLink{sga3:VIIA:2.4:item:ii}{2.4(ii)}); but by construction that
image contains \(X\), a contradiction.

\smallskip
\noindent
\emph{Case \SGARefLink{sga3:xvii:lemma:5-8-2:item:ii}{(ii)}.}
The same argument applies, with the roles of \(\mathfrak r'\) and
\(\mathfrak u'\) interchanged.

\smallskip
\noindent
\emph{End of the proof of \SGARefLink{sga3:xvii:lemma:5:8:1}{5.8.1}.}
Suppose that \(H\ne0\).  By \SGARefLink{sga3:xvii:lemma:5-8-2}{5.8.2} there exists a nonzero reductive
\(p\)-Lie subalgebra
\(\mathfrak h'_1\subset\operatorname{Lie}E\).  Since
\(\operatorname{Lie}U\) is unipotent, by \SGARefLink{sga3:xvii:proposition:4:3:1:item:i}{4.3(i)}, we necessarily
have
\[
  (\operatorname{Lie}U)\cap\mathfrak h'_1=0.
\]
Thus \(\mathfrak h'_1\) is a lift of a \(p\)-Lie subalgebra of
\(\operatorname{Lie}H\).

\smallskip
\noindent
\emph{End of the proof of \SGARefLink{sga3:xvii:theorem:5:1:1:item:i}{5.1.1(i)}(d).}
By \SGARefLink{sga3:xvii:lemma:5:8:1}{5.8.1} there is a family of algebraic subgroups
\(H'_n\) (\(n\in\mathbb N\)) of \(E\) such that
\(H'_{n+1}\) contains \(H'_n\), and such that \(H'_n\) lifts
\({}_{F^n}H\).  The descending sequence of subgroups
\[
  \operatorname{Centr}_E(H'_n)
\]
is stationary; let \(C\) be its stationary value.  The center \(Z\)
of \(C\) contains \(H'_n\) for every \(n\).  Hence the image of
\(Z\) in \(H\) contains \({}_{F^n}(H)\) for every \(n\), and is
therefore an open subgroup of \(H\) (Exposé VII\textsubscript{A},
\S 4).  It is consequently equal to \(H\), since \(H\) is
connected.  To prove that \(E\) is trivial, we may thus replace
\(E\) by \(Z\), and hence suppose that \(E\) is commutative.  We
shall see in \SGARefLink{sga3:xvii:theorem:7:2:1}{7.2.1} that \(E\) then contains a maximal subgroup \(M\)
of multiplicative type, whose formation commutes with extension of
the base field.  Since \(E_{\overline{k}}\) is a trivial extension,
by \SGARefLink{sga3:xvii:theorem:5:1:1:item:i}{5.1.1(i)}\SGARefLink{sga3:xvii:appendix:A}{(a)}, and \(U\) is unipotent, it is clear that \(M\) is
the unique lift of \(H\) in \(E\).

\SGAProof{Proof of \SGARefLink{sga3:xvii:theorem:5:1:1:item:i}{5.1.1(i)}\SGARefLink{sga3:xvii:appendix:C}{(c)}: \(U\) \(k\)-solvable}
Since \((H/H^0)(\overline{k})\) has order prime to \(p\), duality
immediately shows that there exists an integer \(n\) such that
\({}_nH\) is an étale subgroup and the canonical morphism
\[
  {}_nH\longrightarrow H/H^0
\]
is an epimorphism.  Thus \(H/{}_nH\) is connected.  By \SGARefLink{sga3:xvii:lemma:5:2:3:item:i:d}{5.2.3(d)},
there exists a lift \(H'\) of \({}_nH\) in \(E\).  As at the
beginning of the proof of \SGARefLink{sga3:xvii:lemma:5:8:1}{5.8.1}, one sees that
\[
  C=\operatorname{Centr}_E(H')
\]
is a subgroup of \(E\) such that \(C\longrightarrow H\) is an
epimorphism and \(C\cap U\) is smooth.  Replacing \(E\) by \(C\)
and passing to the quotient by \(H'\), we are reduced to the case
where \(U\) is smooth and \(H\) is connected, namely case \SGARefLink{sga3:xvii:lemma:5-8-2:item:i}{(i)}(d).

This completes the proof of \SGARefLink{sga3:xvii:theorem:5-1-1}{5.1.1}.

% END INCLUDED SGA3_English_current_buildable_source_20260823/source/inputs/XVII/tex/components/26_expose_XVII_section5_8_pages031_032.tex

% BEGIN INCLUDED SGA3_English_current_buildable_source_20260823/source/inputs/XVII/tex/components/27_expose_XVII_section5_9_counterexamples_pages032_034.tex
% Authority: Expo17.pdf, local pp.32--34 / combined pp.994--996 /
% printed pp.597--601, all of subsection 5.9.  The next cursor is
% section 6, local p.35 / combined p.997 / printed p.602.

\SGASubsection{5.9}{Counterexamples}
\SGARefTarget{sga3:xvii:subsection:5:9}

Let us first indicate a procedure for obtaining nontrivial
extensions of a \(k\)-group \(H\) of multiplicative type by a
unipotent \(k\)-group \(U\).  Suppose that an action of \(H\) on
\(U\) is given, and let
\[
  E=U\rtimes H
\]
be the semidirect product.  Let, moreover, a class in
\(\mathrm H^1(k,U)\) be represented by a \(1\)-cocycle \(a\).
The group \(U\) acts on \(E\) by inner automorphisms, and \(a\)
therefore defines a \(k\)-form of \(E\), denoted by \(E_a\).
Suppose in addition that \(U\) is commutative.  Then \(U\) acts
trivially on \(U\) and on the quotient \(E/U=H\), so \(E_a\) is
again an extension of \(H\) by \(U\).

For simplicity, suppose that \(H\) is an étale diagonalizable
group, so that the Galois group
\(\mathscr G=\operatorname{Gal}(\overline{k}/k)\) acts trivially
on \(H(\overline{k})\).  Its action on \(E_a(\overline{k})\) is
then given by
\[
  {}^g(u,h)
  =\bigl({}^g u+a(g)-{}^h a(g),h\bigr)
  \qquad
  \bigl(g\in\mathscr G,\ u\in U(\overline{k}),\
        h\in H(\overline{k})\bigr).
\]
If \(h\in H(k)\), and \(X\) is its inverse image in \(E_a\), then
\(X\) is a torsor under \(U\) defined by the class of the
\(1\)-cocycle
\[
  g\longmapsto a(g)-{}^h a(g)
\]
of \(\mathscr G\) with values in \(U\).  It follows that if there
is a point \(h\in H(k)\) for which this cocycle is nontrivial,
then the extension \(E_a\) is nontrivial.  We apply this
construction in the following two particular cases.

\smallskip
\noindent
\emph{\SGARefLink{sga3:xvii:appendix:A}{(a)} A nontrivial extension of an étale diagonalizable group
\(H\) by \(\mathbb Z/p\mathbb Z\).}
Take
\[
  H=(\mathbb Z/p\mathbb Z)^\times,
  \qquad U=\mathbb Z/p\mathbb Z,
\]
with \(H\) acting on \(U\) by multiplication.  Let \(k\) be a
field of characteristic \(p\), let \(K/k\) be a Galois extension
whose Galois group \(\mathscr G\) is isomorphic to
\(\mathbb Z/p\mathbb Z\), and take a nonzero element
\[
  a\in\mathrm H^1(\mathscr G,U)
   =\operatorname{Hom}(\mathscr G,U).
\]
Then \(E_a\) is the required example.

\smallskip
\noindent
\emph{\SGARefLink{sga3:xvii:section:b}{(b)} A nontrivial extension of an étale diagonalizable group
\(H\) by a smooth connected unipotent group \(U\).}
Choose a nonperfect field \(k\) admitting a \(k\)-form \(U\) of
\(\mathbb G_a\) such that \(\mathrm H^1(k,U)\ne0\).  For example
(see J.-P.\ Serre, \emph{Cohomologie Galoisienne}), one may take
for \(k\) the fraction field of a discrete valuation ring of
equal characteristic \(p>0\), and for \(U\) the algebraic
subgroup of \(\mathbb G_a\times_k\mathbb G_a\) defined by
\[
  X^p+X+tY^p=0,
\]
where \(t\) is a uniformizer of \(k\).  Indeed, suppose for
simplicity that \(k\) contains the \((p-1)\)-st roots of unity.
There is then an exact sequence of algebraic groups
\[
  0\longrightarrow\mathbb Z/p\mathbb Z
   \longrightarrow U
   \longrightarrow\mathbb G_a
   \longrightarrow0,
  \qquad
  (x,y)\longmapsto y,
\]
and hence an exact cohomology sequence
\[
  \mathbb G_a(k)\xrightarrow{\,d\,}
  \mathrm H^1(k,\mathbb Z/p\mathbb Z)
  \longrightarrow\mathrm H^1(k,U)
  \longrightarrow0,
\]
where \(d\) associates to \(x\) the homogeneous space under
\(\mathbb Z/p\mathbb Z\) defined by
\[
  X^p+X+tx^p=0.
\]
Moreover,
\[
  \mathrm H^1(k,\mathbb Z/p\mathbb Z)\simeq k/P(k),
  \qquad P(x)=x^p+x,
\]
and therefore
\[
  \mathrm H^1(k,U)\simeq k/\bigl(P(k)+tk^p\bigr).
\]
If, in addition, \(p\ne2\), then \(t^2\) does not belong to
\(P(k)+tk^p\); consequently \(\mathrm H^1(k,U)\ne0\).

The group \(\boldsymbol\mu_{p-1}\) acts on \(U\) by
\[
  (h,x,y)\longmapsto(hx,hy).
\]
Let \(\mathscr G\) be the Galois group of the extension \(K/k\)
defined by
\[
  X^p+X+t^2=0,
\]
and let \(a\in\mathrm H^1(\mathscr G,U)\) be the nonzero class
described above.  It follows immediately that \(E_a\) is a
nontrivial extension of \(\boldsymbol\mu_{p-1}\) by \(U\).

\smallskip
\noindent
\emph{\SGAInlineRefTarget{sga3:xvii:subsection:5-9:item:c}\SGARefLink{sga3:xvii:appendix:C}{(c)} A nontrivial extension of \(\mathbb G_m\) by
\(\boldsymbol\alpha_p\).}
By \SGARefLink{sga3:xvii:theorem:5:1:1:item:i}{5.1.1(i)}\SGARefLink{sga3:xvii:section:b}{(b)}, such an extension can exist only over a
nonperfect field.  Let \(k\) therefore be nonperfect, and let
\(G\) be the semidirect product of \(U=\mathbb G_a\) by
\(H=\mathbb G_m\), where \(\mathbb G_m\) acts on
\(\mathbb G_a\) by homotheties.  Since \(U\) is invariant in
\(G\), so is \({}_{F}U\).  Put
\[
  G'=G/{}_{F}U.
\]
The group \(G'\) is isomorphic to
\(\mathbb G_a\rtimes\mathbb G_m\), where \(\mathbb G_m\) acts on
\(\mathbb G_a\) by
\[
  (h,u)\longmapsto h^p u.
\]
The functor \(\mathscr T_G\) (respectively
\(\mathscr T_{G'}\)) of subtori of \(G\) (respectively \(G'\))
(see Exposé XV) is isomorphic to \(\mathbb G_a\), and the morphism
\[
  \mathscr T_G\longrightarrow\mathscr T_{G'}
\]
induced by \(G\longrightarrow G'\) identifies with
\[
  u\longmapsto u^p.
\]
It follows that if \(T'\) is a subtorus of \(G'\) corresponding
to a point \(x\in k\simeq\mathbb G_a(k)\) such that
\(x^{1/p}\notin k\), then its inverse image \(E\) in \(G\) is an
extension of the torus \(T'\) by
\({}_{F}U=\boldsymbol\alpha_p\).  It has no maximal torus defined
over \(k\), and hence is nontrivial.  As a subgroup of
\[
  G=\operatorname{Spec}k[U,T,T^{-1}],
\]
the group \(E\) is defined by
\[
  U^p=x-xT^p.
\]

\SGAStatement{Remark 5.9.1}{This last example shows that a
nonsmooth algebraic group defined over a nonperfect field need
not possess a maximal torus, and thus answers the question posed
in \SGARefLink{sga3:xiv:remark:1-5:item:b}{Exposé XIV, 1.5(b)}.}
\SGARefTarget{sga3:xvii:remark:5:9:1}

\smallskip
\noindent
\emph{(d) A trivial extension with nonconjugate lifts.}
We now give an example of a trivial extension \(E\) of a group
\(H\) of multiplicative type by a unipotent group \(U\), and of
two lifts \(H'\) and \(H''\) of \(H\) which are not conjugate by
an element of \(U(k)\).

Take for \(E\) the semidirect product of \(U=\mathbb G_a\) by
\[
  \boldsymbol\mu_p
  =\operatorname{Spec}k[T]/(T^p-1),
\]
where the action of \(\boldsymbol\mu_p\) on
\(\mathbb G_a=\operatorname{Spec}k[U]\) is defined by the
comorphism
\[
  U\longmapsto(U+U^p)T-U^p.
\]
The centralizer of \(\boldsymbol\mu_p\) in \(U\) is then the
étale group \(Z\) defined by
\[
  U+U^p=0.
\]
It follows that if \(k\) is not algebraically closed, the
canonical map
\[
  U(k)\longrightarrow(U/Z)(k)
\]
is not in general surjective.  In view of \SGARefLink{sga3:xvii:theorem:5:1:1:item:ii}{5.1.1(ii)}\SGARefLink{sga3:xvii:section:b}{(b)}, two
lifts of \(\boldsymbol\mu_p\) in \(E\) therefore need not be
conjugate by an element of \(U(k)\).

Here is another example, with \(k\) algebraically closed of
characteristic \(p>0\).  Let \(G\) be the radicial semidirect
product in which \(\boldsymbol\mu_p\) acts on
\(\boldsymbol\alpha_p\) by homotheties.  There is an exact
sequence of algebraic groups with \(\boldsymbol\mu_p\)-operators
\[
  0\longrightarrow\boldsymbol\alpha_p
   \longrightarrow\mathbb G_a
   \xrightarrow{\,x\longmapsto x^p\,}\mathbb G_a
   \longrightarrow0,
\]
where \(\boldsymbol\mu_p\) acts by homotheties on the first
\(\mathbb G_a\) and trivially on the second.  The exact
cohomology sequence (Appendix I, \SGARefLink{sga3:xvii:proposition:a-1-1}{Proposition 11}) gives
\[
  0\longrightarrow\mathbb G_a(k)
   \longrightarrow\mathrm H^1(\boldsymbol\mu_p,\boldsymbol\alpha_p)
   \longrightarrow\mathrm H^1(\boldsymbol\mu_p,\mathbb G_a).
\]
Since the last term is zero by I, \SGARefLink{sga3:i:theorem:5-3-3}{5.3.3}, the group
\(\mathrm H^1(\boldsymbol\mu_p,\boldsymbol\alpha_p)\) is nonzero.
Thus two lifts of \(\boldsymbol\mu_p\) in \(G\) need not be
conjugate.

% END INCLUDED SGA3_English_current_buildable_source_20260823/source/inputs/XVII/tex/components/27_expose_XVII_section5_9_counterexamples_pages032_034.tex

% BEGIN INCLUDED SGA3_English_current_buildable_source_20260823/source/inputs/XVII/tex/components/28_expose_XVII_section6_1_and_proofs_A_B_pages035_037.tex
% Authority: Expo17.pdf, local pp.35--37 / combined pp.997--999 /
% printed pp.602--608.  This component covers section 6 through the
% proof of Theorem 6.1.1(B), and stops before the proof of 6.1.1(C).

\SGAInlineRefTarget{sga3:xvii:section:6:site:da2ef89f51}\SGASection{6}{Extensions of a Unipotent Group by a Group of
Multiplicative Type}
\SGARefTarget{sga3:xvii:section:6}

\SGASubsection{6.1}{Statement of the Theorem}
\SGARefTarget{sga3:xvii:subsection:6:1}

\SGAStatement{\SGAInlineRefTarget{sga3:xvii:theorem:6-1-1}Theorem 6.1.1}{Let \(k\) be a field, \(U\) a
unipotent algebraic \(k\)-group, \(H\) a \(k\)-group of
multiplicative type, and \(E\) an algebraic \(k\)-group which is
an extension of \(U\) by \(H\), so that there is an exact sequence
\[
  1\longrightarrow H\longrightarrow E\longrightarrow U
    \longrightarrow1.
\]
Then the extension \(E\) is trivial and there exists a unique
lift of \(U\) in \(E\) in each of the following cases:
\begin{enumerate}
\item[(A)] \SGARefTarget{sga3:xvii:theorem:6:1:1:item:A}
  The group \(U\) is smooth and one of the following conditions
  holds:
  \begin{enumerate}
  \item[(i)] \SGARefTarget{sga3:xvii:theorem:6:1:1:item:A:i}
    \(U\) is connected and the canonical morphism
    \(E\longrightarrow U\) admits a section.
  \item[(ii)] \SGARefTarget{sga3:xvii:theorem:6:1:1:item:A:ii}
    \(U\) has a composition series whose successive quotients
    are isomorphic to \(\mathbb G_a\).
  \item[(iii)] \SGARefTarget{sga3:xvii:theorem:6:1:1:item:A:iii}
    \(H\) is étale.
  \item[(iv)] \SGARefTarget{sga3:xvii:theorem:6:1:1:item:A:iv}
    \(k\) is perfect.
  \end{enumerate}
\item[(B)] \SGARefTarget{sga3:xvii:theorem:6:1:1:item:B}
  \(U=\boldsymbol\alpha_p\) and \(k\) is perfect.
\item[(C)] \SGARefTarget{sga3:xvii:theorem:6:1:1:item:C}
  \(E\) is commutative and \(k\) is perfect.
\end{enumerate}}
\SGARefTarget{sga3:xvii:theorem:6:1:1}

\SGASubsection{6.2}{Proof of \SGARefLink{sga3:xvii:theorem:6:1:1:item:A}{6.1.1(A)}}
\SGARefTarget{sga3:xvii:subsection:6:2}

We first establish three lemmas.

\SGAStatement{Lemma 6.2.1}{Let \(S\) be a prescheme, \(E\) an
\(S\)-group prescheme which is an extension of an \(S\)-group
prescheme \(U\), with connected fibres, by an \(S\)-group \(H\)
of multiplicative type and finite type (that is, \(U\) is the
quotient of \(E\) by \(H\) for the fpqc topology).  Then \(E\) is
a central extension.}
\SGARefTarget{sga3:xvii:lemma:6:2:1}

\SGAProof{Proof}
Since \(H\) is commutative, \(U\) acts on \(H\) by inner
automorphisms, through an \(S\)-morphism of groups
\[
  u:U\longrightarrow\operatorname{Aut}_{S\text{-}\mathrm{gr}}(H).
\]
The functor
\(\operatorname{Aut}_{S\text{-}\mathrm{gr}}(H)\) is representable
by an étale \(S\)-scheme (\SGARefLink{sga3:x:corollary:5-10}{Exposé X, 5.10}), and consequently its
identity section is both an open and a closed immersion.  Since
\(U\) has connected fibres, it follows that \(u\) is the unit
morphism.

\SGAStatement{Lemma 6.2.2}{With the notation of \SGARefLink{sga3:xvii:theorem:6-1-1}{6.1.1}, if \(E\)
is trivial, there exists a unique lift of \(U\) in \(E\).}
\SGARefTarget{sga3:xvii:lemma:6:2:2}

\SGAProof{Proof}
Let \(U'\) and \(U''\) be two lifts of \(U\) in \(E\).  To prove
that \(U'=U''\), we may suppose that \(k\) is algebraically closed,
and it is enough to prove that \(\mathrm H^1(U,H)=0\).  If \(U\)
is connected, then \(U\) centralizes \(H\) by \SGARefLink{sga3:xvii:lemma:6:2:1}{6.2.1}, and hence
\[
  \mathrm H^1(U,H)
    =\operatorname{Hom}_{k\text{-}\mathrm{gr}}(U,H)=0
\]
by 2.4(ii).

In the general case, let \(U_1'\) denote the unique lift in \(E\)
of the identity component \(U^0\) of \(U\), and put
\[
  N=\operatorname{Norm}_{E}(U_1').
\]
If \(g\in E(k)\), then
\(\operatorname{int}(g)U_1'\) is a lift of \(U^0\), and is
therefore equal to \(U_1'\); thus \(N(k)\supset E(k)\).
Moreover, \(N\) contains \(H\), by \SGARefLink{sga3:xvii:lemma:6:2:1}{6.2.1}, as well as \(U_1'\),
and it follows at once that \(N=E\).  On passing to the quotient
by \(U_1'\), we are reduced to the case in which \(U\) is étale.
Since \(k\) is algebraically closed, in this case
\(\mathrm H^1(U,H)\) identifies with the ordinary cohomology
group\footnote{Editor's note: see, for example, Proposition
III.6.4.2 of M.\ Demazure and P.\ Gabriel,
\emph{Groupes algébriques I}, Masson \& North-Holland (1970).}
\[
  \mathrm H^1\bigl(U(k),H(k)\bigr),
\]
and is therefore zero, since \(U(k)\) is a finite \(p\)-group and
\(H(k)\) is uniquely \(p\)-divisible.

\SGAStatement{Corollary 6.2.3}{To prove that \(E\) is trivial, it
is enough to prove that it becomes trivial after a finite
separable extension of the base field.  In particular, we may
suppose that \(H\) is diagonalizable.}
\SGARefTarget{sga3:xvii:corollary:6:2:3}

\SGAStatement{Lemma 6.2.4}{In order that every central extension
\(E\) of \(U\) by a diagonalizable \(k\)-group \(H\) be trivial,
it is enough that every central extension of \(U\) by
\(\mathbb G_m\) be trivial.}
\SGARefTarget{sga3:xvii:lemma:6:2:4}

\SGAProof{Proof}
Induction on \(r\) first shows that the hypothesis implies that
\(E\) is trivial when \(H=\mathbb G_m^r\).  In the general case,
\(H\) embeds in \(\mathbb G_m^r\) for a suitable integer \(r\)
(this is immediate by duality); put
\[
  H''=\mathbb G_m^r/H.
\]
We obtain the exact sequence (Appendix I, \SGARefLink{sga3:i:subsection:2-1}{2.1})
\[
  \mathrm Z^1(U,H'')
    \longrightarrow\operatorname{Ext}_{\mathrm{alg}}(U,H)
    \longrightarrow
      \operatorname{Ext}_{\mathrm{alg}}(U,\mathbb G_m^r)=0,
\]
where \(U\) acts trivially on \(H\), \(H''\), and
\(\mathbb G_m^r\).  But
\[
  \mathrm Z^1(U,H'')
    =\operatorname{Hom}_{k\text{-}\mathrm{gr}}(U,H'')=0
\]
by 2.4(ii), and hence
\(\operatorname{Ext}_{\mathrm{alg}}(U,H)=0\).

\SGAProof{Proof of \SGARefLink{sga3:xvii:theorem:6:1:1:item:A}{6.1.1(A)}\SGARefLink{sga3:xvii:theorem:6:1:1:item:A:i}{(i)}}
Since \(E\longrightarrow U\) admits a section, \(E\) defines an
element \(f\) of \(\mathrm H^2(U,H)\) (Appendix I, \SGARefLink{sga3:i:definition:3-1}{3.1}).  We must
prove that \(f\) is zero.  For this it is enough to prove that a
\(2\)-cocycle
\[
  f:U\times_kU\longrightarrow H
\]
is a constant morphism, which follows from the next lemma.

\SGAStatement{Lemma 6.2.5}{Let \(U\) be a smooth connected
unipotent algebraic \(k\)-group, and \(H\) a \(k\)-group of
multiplicative type.  Then every \(k\)-morphism of preschemes
\[
  f:U\longrightarrow H
\]
is constant.}
\SGARefTarget{sga3:xvii:lemma:6:2:5}

\SGAProof{Proof}
We may suppose that \(k\) is algebraically closed.  We argue by
induction on \(\dim U\).  If \(\dim U>0\), then \(U\) has a
composition series (see \SGARefLink{sga3:xvii:lemma:3-9}{3.9})
\[
  1\longrightarrow U'\longrightarrow U
    \xrightarrow{\ \pi\ }U''\longrightarrow1,
  \qquad
  U'\simeq\mathbb G_a,\quad \dim U''<\dim U.
\]
It is enough to show that \(f\) factors through \(U''\).  Since
the graph of the equivalence relation defined by \(\pi\) is
smooth over \(k\), and hence reduced, it is enough to show that
if \(x,y\in U(k)\) have the same image \(z\) in \(U''(k)\), then
\(f(x)=f(y)\).  Now \(\pi^{-1}(z)\) is isomorphic to
\(\mathbb G_a\), so the restriction of \(f\) to
\(\pi^{-1}(z)\) factors through a reduced irreducible component
of \(H\), and hence through a \(k\)-scheme isomorphic to
\(\mathbb G_m^r\).  It remains only to note that every morphism
from \(\mathbb G_a\) to \(\mathbb G_m^r\) is constant, since
every invertible regular function on \(\mathbb G_a\) is
constant.

\SGAProof{Proof of \SGARefLink{sga3:xvii:theorem:6:1:1:item:A}{6.1.1(A)}\SGARefLink{sga3:xvii:theorem:6:1:1:item:A:ii}{(ii)}}
By \SGARefLink{sga3:xvii:lemma:6:2:1}{6.2.1}, \SGARefLink{sga3:xvii:corollary:6:2:3}{6.2.3}, and \SGARefLink{sga3:xvii:lemma:6:2:4}{6.2.4}, we may suppose that
\(H=\mathbb G_m\).  By hypothesis, \(U\) has a composition series
\[
  U\supset U_1\supset\cdots
\]
such that \(U_i/U_{i+1}\) is isomorphic to \(\mathbb G_a\).  Let
\(E_1\) be the inverse image of \(U_1\) in \(E\).  By induction
on \(\dim U\), we may suppose that the extension \(E_1\) is
trivial; let \(U_1'\) be the unique lift of \(U_1\).  Arguing as
in the normalizer argument of \SGARefLink{sga3:xvii:lemma:6:2:2}{6.2.2}, we find that \(U_1'\) is
invariant in \(E\).  After taking the quotient by \(U_1'\), we
are reduced to the case \(U=\mathbb G_a\).  The \(S\)-scheme
\(E\), where \(S=U\), is then a torsor under the \(S\)-group
\(\mathbb G_m\times_kS\), and therefore admits a section, since
\(\operatorname{Pic}(S)=0\) (\SGARefLink{sga3:viii:corollary:4:3}{Exposé VIII, 4.3}).  The extension
\(E\) is then trivial by (A)\SGARefLink{sga3:xvii:theorem:6:1:1:item:A:i}{(i)}.

\SGAProof{Proof of \SGARefLink{sga3:xvii:theorem:6:1:1:item:A}{6.1.1(A)}\SGARefLink{sga3:xvii:theorem:6:1:1:item:A:iii}{(iii)} (\(U\) smooth, \(H\) étale)}
Suppose first that \(U\) is connected.  The group
\(U_{\overline{k}}\) then has a composition series whose
successive quotients are isomorphic to \(\mathbb G_a\), so
\(E_{\overline{k}}\) is trivial by (A)\SGARefLink{sga3:xvii:theorem:6:1:1:item:A:ii}{(ii)}.  Since \(H\) is
étale, the unique lift of \(U_{\overline{k}}\) in
\(E_{\overline{k}}\) is plainly the identity component of
\(E_{\overline{k}}\); this lift is therefore already defined
over \(k\).  In the general case, after quotienting by the
identity component of \(E\), we are reduced to the case where
\(E\) is étale, and then to the case where \(E\) is completely
split (\SGARefLink{sga3:xvii:corollary:6:2:3}{6.2.3}).  Since \(U(k)\) is a \(p\)-group and \(H(k)\) has
order prime to \(p\), we may take the Sylow \(p\)-subgroup of
\(E(k)\) as a lift of \(U(k)\).

\SGAProof{Proof of \SGARefLink{sga3:xvii:theorem:6:1:1:item:A}{6.1.1(A)}\SGARefLink{sga3:xvii:theorem:6:1:1:item:A:iv}{(iv)} (\(U\) smooth, \(k\) perfect)}
If \(U\) is connected, it has a composition series whose
successive quotients are isomorphic to \(\mathbb G_a\)
(\SGARefLink{sga3:xvii:lemma:4:1:2:item:b}{4.1.2(b)}), and we apply (A)\SGARefLink{sga3:xvii:theorem:6:1:1:item:A:ii}{(ii)}.  In the general case, the
preceding argument reduces us to the case in which \(U\) is
étale, and then to the case in which \(U\) is completely split
and \(H\) is diagonalizable (\SGARefLink{sga3:xvii:corollary:6:2:3}{6.2.3}).  Using a characteristic
composition series of \(H\),
\[
  H\supset H^0\supset{}_{F^n}(H)\supset0,
\]
we are reduced to the case in which \(H\) is of one of the
following three types:
\begin{enumerate}
\item[(a)] \(H\) is étale;
\item[(b)] \(H=\mathbb G_m^r\);
\item[(c)] \(H\) is radicial.
\end{enumerate}
In case \SGARefLink{sga3:xvii:appendix:A}{(a)} we apply (A)\SGARefLink{sga3:xvii:theorem:6:1:1:item:A:iii}{(iii)}, and in case \SGARefLink{sga3:xvii:appendix:C}{(c)} we apply \SGARefLink{sga3:xvii:lemma:1:6}{1.6}.
Finally, in case \SGARefLink{sga3:xvii:section:b}{(b)}, Hilbert's \SGARefLink{sga3:viii:corollary:4:5}{Theorem 90} shows that
\(E\longrightarrow U\) admits a section, so it is enough to show
that
\[
  \mathrm H^2(U,\mathbb G_m^r)
    =\mathrm H^2\bigl(U(k),\mathbb G_m^r(k)\bigr)=0.
\]
But \(U(k)\) is a finite \(p\)-group, whereas
\(\mathbb G_m^r(k)\) is uniquely \(p\)-divisible, since \(k\) is
perfect.

\SGASubsection{6.3}{Proof of \SGARefLink{sga3:xvii:theorem:6:1:1:item:B}{6.1.1(B)} and (C)}
\SGARefTarget{sga3:xvii:subsection:6:3}

By \SGARefLink{sga3:xvii:lemma:6:2:1}{6.2.1}, \SGARefLink{sga3:xvii:corollary:6:2:3}{6.2.3}, and \SGARefLink{sga3:xvii:lemma:6:2:4}{6.2.4}, it is enough to prove (B) when
\(H=\mathbb G_m\).  We therefore have an exact sequence
\[
  1\longrightarrow\mathbb G_m\longrightarrow E
    \longrightarrow\boldsymbol\alpha_p\longrightarrow1.
\]
Since \(\mathbb G_m\) is smooth, it gives an exact sequence of
restricted Lie algebras (Appendix II, \SGARefLink{sga3:xvii:proposition:3:2}{3.2})
\begin{equation}
  0\longrightarrow\operatorname{Lie}\mathbb G_m
    \longrightarrow\operatorname{Lie}E
    \longrightarrow\operatorname{Lie}\boldsymbol\alpha_p
    \longrightarrow0.
  \tag{*}\label{sga3:xvii:subsection:6-3:equation:star}
\end{equation}
The group \(\boldsymbol\alpha_p\) has height \(1\), so the
extension \(E\) is trivial if and only if the exact sequence
\((*)\) of restricted Lie algebras splits (Appendix II, \SGARefLink{sga3:ii:subsection:2-2}{2.2}).
The algebra \(\operatorname{Lie}\boldsymbol\alpha_p\) is
generated by an element \(X\ne0\) such that \(X^{(p)}=0\)
(Appendix II, \SGARefLink{sga3:xvii:proposition:2:1}{2.1}).  It is therefore enough to show that \(X\)
lifts to an element \(Z\) of \(\operatorname{Lie}E\) for which
\(Z^{(p)}=0\).  By \SGARefLink{sga3:xvii:lemma:5-8-2:item:ii}{5.8.2(ii)}, there is a unipotent element \(Z\)
of \(\operatorname{Lie}E\) which lifts \(X\).  Since the unipotent
part of \(\operatorname{Lie}E\) plainly has dimension at most
\(1\), it follows necessarily that \(Z^{(p)}=0\).

% END INCLUDED SGA3_English_current_buildable_source_20260823/source/inputs/XVII/tex/components/28_expose_XVII_section6_1_and_proofs_A_B_pages035_037.tex

% BEGIN INCLUDED SGA3_English_current_buildable_source_20260823/source/inputs/XVII/tex/components/29_expose_XVII_proof_C_and_section6_4_examples_page038.tex
\SGAProof{Proof of \SGARefLink{sga3:xvii:theorem:6:1:1:item:C}{6.1.1(C)}}
If \(U_1'\) is a lift in \(E\) of an algebraic subgroup \(U_1\)
of \(U\), then \(U_1'\) is invariant in \(E\), since \(E\) is
assumed commutative, and we may pass to the quotient by \(U_1'\).
Using a composition series of \(U\) (\SGARefLink{sga3:xvii:theorem:3:5:item:ii}{3.5(ii)}), the preceding
remark reduces us to the case in which \(U\) is smooth or equal
to \(\boldsymbol\alpha_p\).  But then \(E\) is trivial by
(A)(iv) and (B).

\SGASubsection{6.4}{Examples of nontrivial extensions of a
unipotent group \(U\) by a group \(H\) of multiplicative type}
\SGARefTarget{sga3:xvii:subsection:6:4}

In view of \SGARefLink{sga3:xvii:theorem:6:1:1:item:A}{6.1.1(A)}(iv), the question arises only in
characteristic \(p>0\).

\begin{enumerate}
\item[(a)] \SGARefTarget{sga3:xvii:subsection:6:4:item:a}
  \(H=\mathbb G_m\), and \(U\) is a nontrivial form of
  \(\mathbb G_a\) (cf.\ Appendix III, \S5).

  Let \(k\) be a nonperfect field of characteristic \(2\), and
  let \(u\in k\) be such that \(u^{1/2}\notin k\).  Consider the
  affine group \(E\) with coordinate ring
  \[
    k\bigl[X,Y,(X^2+uY^2)^{-1}\bigr],
  \]
  whose multiplication is given by the comorphism
  \[
    (X,Y)\longmapsto
      \bigl(XX'+uYY',\,XY'+YX'\bigr).
  \]
  The group \(E\) is smooth, connected, commutative, and of
  dimension \(2\); the subscheme \(Y=0\) defines a subgroup
  \(H\simeq\mathbb G_m\).  The kernel \(K\) of the squaring map
  on \(E\) has equation
  \[
    X^2+uY^2=1,
  \]
  and hence has dimension \(1\).  The group \(K\) contains the
  unipotent radical of \(E\), defined over \(\bar k\), as well as
  the contribution of \(H\), which is isomorphic to
  \(\boldsymbol\mu_2\).  Since \(K\) is reduced over \(\bar k\),
  the unipotent radical of \(E\) is not defined over \(k\), and
  \(U=E/H\) does not lift in \(E\).

  One checks immediately that \(U\) is the form of
  \(\mathbb G_a\) with coordinate ring
  \[
    k[V,W]/(V+uV^2+W^2),
  \]
  the morphism \(E\longrightarrow U\) corresponding to the
  comorphism
  \[
    V\longmapsto\frac{Y^2}{X^2+uY^2},
    \qquad
    W\longmapsto\frac{XY}{X^2+uY^2}.
  \]

\item[(b)] \SGARefTarget{sga3:xvii:subsection:6:4:item:b}
  \(H=\mathbb G_m\), \(U=\mathbb Z/2\mathbb Z\), and \(k\) is
  nonperfect of characteristic \(2\).

  Choose \(k\) and \(u\) as in \SGARefLink{sga3:xvii:appendix:A}{(a)}, and consider the subgroup
  \(E\) of \(\mathrm{GL}_2\) generated by the element
  \[
    X=
    \begin{pmatrix}0&1\\u&0\end{pmatrix}
    =
    \begin{pmatrix}u^{1/2}&0\\0&u^{1/2}\end{pmatrix}
    \begin{pmatrix}0&u^{-1/2}\\u^{1/2}&0\end{pmatrix}.
  \]
  The group \(E\) is an extension of \(\mathbb Z/2\mathbb Z\)
  by \(\mathbb G_m\), but this extension is not trivial because
  the unipotent part of \(X\) is not defined over \(k\).

\item[(c)] \SGARefTarget{sga3:xvii:subsection:6:4:item:c}
  \(H=\boldsymbol\mu_p\), \(U=\boldsymbol\alpha_p\), and \(k\)
  is nonperfect.

  Let \(\mathfrak e\) be the commutative restricted Lie algebra
  generated by two elements \(X\) and \(Y\) such that
  \[
    X^{(p)}=X,\qquad Y^{(p)}=aX.
  \]
  By Appendix II, \SGARefLink{sga3:ii:subsection:2-2}{2.2}, \(\mathfrak e\) is the restricted Lie
  algebra of an algebraic group \(E\) which is an extension of
  \(\boldsymbol\alpha_p\) by \(\boldsymbol\mu_p\).  This
  extension is trivial if and only if there exists \(b\in k\)
  such that \(b^p=a\), for then
  \[
    (Y-bX)^{(p)}=0.
  \]

\item[(d)] \SGARefTarget{sga3:xvii:subsection:6:4:item:d}
  \(H=\boldsymbol\mu_2\), \(U=\boldsymbol\alpha_2\times
  \boldsymbol\alpha_2\), \(E\) noncommutative, and \(k\) a field
  of characteristic \(2\).

  Consider the special linear group \(\mathrm{SL}_{2,k}\), and
  put \(E=F(\mathrm{SL}_{2,k})\).  The group \(E\) is a radicial
  group of height \(1\), whose Lie algebra is generated by three
  elements \(X,Y,Z\) satisfying
  \[
    [X,Y]=Z,\qquad [X,Z]=[Y,Z]=0,
  \]
  \[
    X^{(p)}=Y^{(p)}=0,\qquad Z^{(p)}=Z.
  \]
\end{enumerate}

% END INCLUDED SGA3_English_current_buildable_source_20260823/source/inputs/XVII/tex/components/29_expose_XVII_proof_C_and_section6_4_examples_page038.tex

% BEGIN INCLUDED SGA3_English_current_buildable_source_20260823/source/inputs/XVII/tex/components/30_expose_XVII_section6_4_completion_and_proposition7_1_1_page039.tex
% Authority: Expo17.pdf, local p.39 / combined p.1001 /
% printed pp.611--612.  This component completes 6.4(d), begins
% section 7, and carries the proof of Proposition 7.1.1 through
% part ii(a).  It stops before part ii(b) on authority local p.40.

Consequently, \(E\) is a central extension of
\(U\simeq\boldsymbol\alpha_2\times\boldsymbol\alpha_2\) by
\(\boldsymbol\mu_2\).  Each factor \(\boldsymbol\alpha_2\) of
\(U\) lifts uniquely in \(E\), but \(U\) itself does not lift in
\(E\), since \([X,Y]\ne0\).

\SGASection{7}{Nilpotent Affine Algebraic Groups}
\SGARefTarget{sga3:xvii:section:7}

\SGASubsection{7.1}{Extensions of Groups of Multiplicative Type}
\SGARefTarget{sga3:xvii:subsection:7:1}

\SGAStatement{Proposition 7.1.1}{Let \(S\) be a prescheme, let
\(H\) and \(K\) be two \(S\)-group preschemes of multiplicative
type and finite type, and let \(E\) be an \(S\)-group prescheme
which is an extension of \(K\) by \(H\) (that is, \(K\) is the
quotient of \(E\) by \(H\) for the fpqc topology).  Then \(E\) is
of multiplicative type in either of the following two cases:
\begin{enumerate}
\item[(a)] \SGARefTarget{sga3:xvii:proposition:7:1:1:item:a}
  \(E\) is commutative.
\item[(b)] \SGARefTarget{sga3:xvii:proposition:7:1:1:item:b}
  The fibres of \(K\) are connected.
\end{enumerate}}
\SGARefTarget{sga3:xvii:proposition:7:1:1}

\SGAProof{Proof}
\begin{enumerate}
\item[\SGAInlineRefTarget{sga3:xvii:proposition:7:1:1:item:i}(i)] Suppose first that \(S\) is the spectrum of a field
  \(k\).  The assertion is local for the fpqc topology, so we may
  suppose that \(k\) is algebraically closed and hence that \(H\)
  and \(K\) are diagonalizable.  Notice that \(K\) acts trivially
  on \(H\) by inner automorphisms: this is clear in case \SGARefLink{sga3:xvii:appendix:A}{(a)}, and
  follows from \SGARefLink{sga3:xvii:lemma:6:2:1}{6.2.1} in case \SGARefLink{sga3:xvii:section:b}{(b)}.  Induction on the length of a
  suitable composition series of \(K\) reduces us to the case
  where \(K\) is of one of the following three types:
  \[
    \text{(a) }K=\mathbb G_m,\qquad
    \text{(b) }K=\boldsymbol\mu_p,\qquad
    \text{(c) }K=\boldsymbol\mu_q,\quad(q,p)=1,
  \]
  and in case \SGARefLink{sga3:xvii:proposition:7:1:1:item:c}{(c)}, \(E\) is commutative.  Now embed \(H\) in
  \(\mathbb G_m^r\).  It follows that \(E\) embeds in an extension
  of \(K\) by \(\mathbb G_m^r\), and we may therefore suppose that
  \(H=\mathbb G_m^r\).

  \begin{enumerate}
  \item[(a)] If \(K=\mathbb G_m\), then \(E\) is a smooth,
    connected, affine algebraic group (Exposé VI\textsubscript{B},
    \SGARefLink{prop:sga3-VIB-9.2}{9.2}) of unipotent rank zero, and is therefore a torus.

  \item[(b)] Let \(K=\boldsymbol\mu_p\).  In this case \(E\) is a
    trivial extension.  Indeed, since \(H\) is smooth, the
    canonical morphism
    \[
      \operatorname{Lie}E\longrightarrow
      \operatorname{Lie}\boldsymbol\mu_p
    \]
    is surjective (Appendix II, \SGARefLink{sga3:xvii:proposition:3:2}{3.2}), and the assertion follows
    from \SGARefLink{sga3:xvii:lemma:5-8-2:item:i}{5.8.2(i)}, in view of Appendix II, \SGARefLink{sga3:ii:subsection:2-2}{2.2}.

  \item[\SGAInlineRefTarget{sga3:xvii:proposition:7:1:1:item:c}(c)] Let \(K=\boldsymbol\mu_q\), where \((q,p)=1\), and
    suppose that \(E\) is commutative.  Again the extension \(E\)
    is trivial.  Indeed, let \(x\in E(\bar k)\) lift a generator
    \(\bar x\) of \(\boldsymbol\mu_q(\bar k)\).  The element
    \(x^q\) belongs to \(H(\bar k)\), and hence has the form
    \(y^q\), with \(y\in H(\bar k)\) (notice that
    \(\mathbb G_m^r(\bar k)\) is \(q\)-divisible).  Since \(E\)
    is commutative, \(y^{-1}x\) is a lift of \(\bar x\) having
    order \(q\).
  \end{enumerate}

\item[(ii)] Consider now the general case.  The groups \(H\) and
  \(K\) are flat, affine, and of finite presentation over \(S\)
  (\SGARefLink{sga3:ix:proposition:2:1}{Exposé IX, 2.1}), and consequently so is \(E\) (Exposé
  VI\textsubscript{B}, \SGARefLink{prop:sga3-VIB-9.2}{9.2}).  Using the general technique of
  Exposé VI\textsubscript{B}, \S10, we reduce to the case where
  \(S\) is noetherian.  To prove that \(E\) is of multiplicative
  type, it then suffices to prove that \(E\) is commutative and
  that \({}_nE\) is finite over \(S\) for every \(n\) (Exposé X,
  4.8(b)).

  \begin{enumerate}
  \item[(a)] \(E\) is commutative.  We must verify that the
    morphism
    \[
      E\times_S E\longrightarrow E,\qquad
      (x,y)\longmapsto[x,y]=xyx^{-1}y^{-1}
    \]
    factors through the unit section of \(E\), and it is enough
    to verify this when \(S\) is the spectrum of an artinian
    local ring.  But then \(E\) is of multiplicative type by \SGARefLink{sga3:xvii:proposition:7:1:1:item:i}{(i)}
    and \SGARefLink{sga3:x:corollary:2:3}{Exposé X, 2.3}.

% END INCLUDED SGA3_English_current_buildable_source_20260823/source/inputs/XVII/tex/components/30_expose_XVII_section6_4_completion_and_proposition7_1_1_page039.tex

% BEGIN INCLUDED SGA3_English_current_buildable_source_20260823/source/inputs/XVII/tex/components/31_expose_XVII_proposition7_1_1_completion_and_theorem7_2_1_page040.tex
% Authority: Expo17.pdf, local p.40 / combined p.1002 /
% printed pp.613--614.  This component completes Proposition 7.1.1,
% states Theorem 7.2.1, and carries its proof into the descent
% argument.  The following component completes that argument.

  \item[(b)] \SGARefTarget{sga3:xvii:proposition:7:1:1:proof:ii:item:b}
    The group \({}_nE\) is finite over \(S\).  Indeed, there is an
    exact sequence
    \[
      0\longrightarrow{}_nH\longrightarrow{}_nE
      \longrightarrow{}_nK\xrightarrow{u}H_n,
    \]
    where \(H_n\) is the cokernel of the \(n\)-th power map on
    \(H\).  We know that \(H_n\) is of multiplicative type (\SGARefLink{sga3:ix:proposition:2:7}{Exposé
    IX, 2.7}), and hence separated, and that \({}_nH\) and
    \({}_nK\) are finite over \(S\).  The group \(\ker u\) is a
    closed subgroup of \({}_nK\), and hence is finite over \(S\);
    consequently \({}_nE\) is finite over \(S\), as an extension
    of one finite group by another (Exposé VI\textsubscript{B},
    \SGARefLink{prop:sga3-VIB-9.2}{9.2}).
  \end{enumerate}
\end{enumerate}

\SGASubsection{7.2}{Structure of Commutative Affine Algebraic Groups}
\SGARefTarget{sga3:xvii:subsection:7:2}

\SGAStatement{Theorem 7.2.1}{Let \(k\) be a field and \(G\) a
commutative affine algebraic \(k\)-group.  Then:
\begin{enumerate}
\item[(a)] \SGARefTarget{sga3:xvii:theorem:7:2:1:item:a}
  \(G\) contains a largest subgroup \(M\) of multiplicative type.
  The group \(M\) is characteristic in \(G\), the quotient \(G/M\)
  is unipotent, and the formation of \(M\) commutes with extension
  of the ground field \(k\).
\item[(b)] \SGARefTarget{sga3:xvii:theorem:7:2:1:item:b}
  If \(k\) is perfect, \(G\) is the direct product of \(M\) and a
  unipotent algebraic subgroup \(U\), uniquely so.
\end{enumerate}}
\SGARefTarget{sga3:xvii:theorem:7:2:1}

\SGAProof{Proof}
\begin{enumerate}
\item[(i)] Suppose that \(k\) is algebraically closed.  When \(G\)
  is smooth, \SGARefLink{sga3:xvii:theorem:7:2:1:item:b}{7.2.1(b)} is well known (BIBLE~4, Theorem~4).  If
  \(G\) is radicial of height \(1\), the decomposition of
  \(\operatorname{Lie}G\) described in \SGARefLink{sga3:xvii:lemma:4-2-2}{4.2.2} corresponds, in view
  of \SGARefLink{sga3:xvii:proposition:4:3:1:item:v}{4.3.1(v)} and Appendix II, \SGARefLink{sga3:ii:subsection:2-2}{2.2}, to a decomposition of \(G\) of
  the kind asserted in \SGARefLink{sga3:xvii:theorem:7:2:1:item:b}{7.2.1(b)}.  In general, \(G\) has a
  composition series whose successive quotients are smooth or
  radicial of height \(1\) (Appendix II, \SGARefLink{sga3:xvii:lemma:3-1}{3.1}).  To prove \SGARefLink{sga3:xvii:theorem:7:2:1:item:b}{7.2.1(b)},
  it therefore suffices to observe the following.  Suppose there
  is an exact sequence of commutative algebraic groups
  \[
    0\longrightarrow G'\longrightarrow G\longrightarrow
    G''\longrightarrow0,
  \]
  in which \(G'\), respectively \(G''\), is the product of a group
  of multiplicative type and a unipotent group,
  \[
    G'=M'\cdot U',\qquad G''=M''\cdot U''.
  \]
  Then the same is true of \(G\).  Indeed, consider the exact
  sequence
  \[
    0\longrightarrow G'/U'\longrightarrow G/U'
      \longrightarrow G''\longrightarrow0.
  \]
  By \SGARefLink{sga3:xvii:proposition:7:1:1:item:a}{7.1.1(a)}, the inverse image of \(M''\) in \(G/U'\) is a
  subgroup \(M_1\) of multiplicative type.  The group \(M_1\)
  lifts to a subgroup \(M\) of \(G\) by \SGARefLink{sga3:xvii:theorem:5:1:1:item:i}{5.1.1(i)}\SGARefLink{sga3:xvii:appendix:A}{(a)}.  Likewise,
  now using \SGARefLink{sga3:xvii:theorem:6:1:1:item:C}{6.1.1(C)}, one proves that there is a unipotent
  subgroup \(U\) of \(G\) which is an extension of \(U''\) by
  \(U'\).  Clearly \(G=M\cdot U\).

\item[(ii)] Let \(k\) now be arbitrary.  By \SGARefLink{sga3:xvii:subsection:6:4:item:i}{(i)},
  \(G_{\bar k}\) is the direct product of a group
  \(M_{\bar k}\) of multiplicative type and a unipotent group
  \(U_{\bar k}\).  Put
  \[
    S=\bar k\times_k\bar k,
  \]
  and let \(M_1\) and \(M_2\) be the two inverse images of
  \(M_{\bar k}\) under the two projections
  \(S\rightrightarrows\bar k\).  The group
  \[
    G_S/M_2=(G/M_{\bar k})_S
  \]
  has unipotent fibres; hence the image of \(M_1\) in \(G_S/M_2\)
  is trivial by \SGARefLink{sga3:xvii:proposition:2:4:item:i}{2.4(i)}, and \(M_1\subset M_2\).  The converse
  inclusion follows in the same way, so \(M_1=M_2\).  By fpqc
  descent, \(M_{\bar k}\) therefore descends to an algebraic
  subgroup \(M\) of \(G\).  It is clear that \(M\) is of
  multiplicative type, that \(G/M\) is unipotent, and that the
  formation of \(M\) is compatible with every extension of
  \(k\).  For every \(k\)-prescheme \(S\), every subgroup \(H\) of
  multiplicative type in \(G_S\) is contained in \(M_S\).
  Indeed, its image in the group with unipotent fibres
  \[
    (G/M)_S=G_S/M_S
  \]
  is trivial by \SGARefLink{sga3:xvii:lemma:2:5}{2.5}.

% END INCLUDED SGA3_English_current_buildable_source_20260823/source/inputs/XVII/tex/components/31_expose_XVII_proposition7_1_1_completion_and_theorem7_2_1_page040.tex

% BEGIN INCLUDED SGA3_English_current_buildable_source_20260823/source/inputs/XVII/tex/components/32_expose_XVII_theorem7_2_1_completion_through_lemma7_3_2_page041.tex
% Authority: Expo17.pdf, local p.41 / combined p.1003 /
% printed pp.615--616.  This component completes Theorem 7.2.1,
% gives Remark 7.2.2, states Theorem 7.3.1 and Lemma 7.3.2, and
% carries the lemma's proof through the construction in case (i).
% It stops before the concluding sentence about the generators of N.

  Taking \(H\), in particular, to be the transform of \(M_S\)
  under an automorphism of \(G_S\), we conclude that \(M\) is
  characteristic in \(G\).

  Finally, if \(k\) is perfect, \(G/M\) lifts in \(G\) to a
  unipotent group, uniquely so by \SGARefLink{sga3:xvii:theorem:6:1:1:item:C}{6.1.1(C)}.
\end{enumerate}

\SGAStatement{Remark 7.2.2}{%
\begin{enumerate}
\item[\SGAInlineRefTarget{sga3:xvii:subsection:6:4:item:i}(i)] \SGARefTarget{sga3:xvii:remark:7:2:2:item:i}
  If \(k\) is not perfect, the unipotent component of
  \(G_{\bar k}\) need not be defined over \(k\), as \SGARefLink{sga3:xvii:subsection:6:4:item:a}{Example 6.4(a)}
  shows.

\item[(ii)] \SGARefTarget{sga3:xvii:remark:7:2:2:item:ii}
  In contrast to the multiplicative-type component \(M\), the
  unipotent component \(U\) is not in general characteristic in
  \(G\), whatever the characteristic of \(k\).  Of course, the
  uniqueness of the decomposition in \SGARefLink{sga3:xvii:theorem:7:2:1:item:b}{7.2.1(b)} implies that \(U\)
  is invariant under every \(k\)-automorphism of \(G\).  But if
  \(U\) and \(M\) are such that there exist a \(k\)-prescheme
  \(S\) and a nonzero \(S\)-homomorphism
  \[
    h:U_S\longrightarrow M_S
  \]
  (cf.\ \SGARefLink{sga3:xvii:remark:2-6}{2.6}), then one obtains an \(S\)-automorphism of \(G_S\),
  \[
    (u,m)\longmapsto\bigl(u,h(u)+m\bigr),
  \]
  which does not leave \(U_S\) invariant.

\item[(iii)] \SGARefTarget{sga3:xvii:remark:7:2:2:item:iii}
  If \(G\) is finite over \(k\), then under Cartier duality (\SGARefLink{sga3:xvii:remark:2-6}{2.6})
  the quotient \(G/M\) corresponds to the connected component of
  the identity in the dual \(D(G)\) of \(G\).
\end{enumerate}}
\SGARefTarget{sga3:xvii:remark:7:2:2}

\SGASubsection{7.3}{Structure of Nilpotent Affine Algebraic Groups}
\SGARefTarget{sga3:xvii:subsection:7:3}

\SGAStatement{Theorem 7.3.1}{Let \(k\) be a field and \(G\) a
nilpotent (Exposé VI\textsubscript{B}, \S8), affine, connected
algebraic \(k\)-group.  Then \(G\) has a largest subgroup \(M\) of
multiplicative type.  The group \(M\) is central and
characteristic, and \(G/M\) is a unipotent algebraic group.}
\SGARefTarget{sga3:xvii:theorem:7:3:1}

Let \(Z\) be the centre of \(G\), and let \(M\) be the largest
subgroup of multiplicative type in \(Z\), as supplied by \SGARefLink{sga3:xvii:theorem:7:2:1}{7.2.1}.
Since \(Z\) is characteristic in \(G\), and \(M\) is
characteristic in \(Z\), the group \(M\) is characteristic in
\(G\).  It remains to show that \(G/M\) is unipotent.  By
induction on the length of the ascending central series of \(G\),
this follows from the following more general lemma.

\SGAStatement{Lemma 7.3.2 (Rosenlicht)}{Let \(G\) be a connected
algebraic \(k\)-group and let \(Z\) be its centre.  Then the centre
\(Z'\) of \(G'=G/Z\) is unipotent.}
\SGARefTarget{sga3:xvii:lemma:7:3:2}

\SGAProof{Proof}
We may suppose that \(k\) is algebraically closed.  It is then
enough to show that \(Z'\) contains no subgroup isomorphic to
\(\boldsymbol\mu_\ell\), for any prime number \(\ell\), by
\SGARefLink{sga3:xvii:theorem:4:6:1:item:vi}{4.6.1(vi)}.  Let \(\boldsymbol\mu_\ell\) be such a subgroup of
\(Z'\), and let \(N\) be its inverse image in \(G\).  Since
\(\boldsymbol\mu_\ell\) is central in \(G'\), the group \(N\) is
invariant in \(G\).

\begin{enumerate}
\item[(i)] \SGARefTarget{sga3:xvii:lemma:7:3:2:proof:item:i}
  Suppose that \((\ell,p)=1\).  We can find an element
  \(x\in N(k)\) and an integer \(n\) with the following
  properties:
  \begin{enumerate}
  \item[(a)] \SGARefTarget{sga3:xvii:lemma:7:3:2:proof:item:i:a}
    \(x\) lifts a generator \(\bar x\) of
    \(\boldsymbol\mu_\ell\);
  \item[(b)] \SGARefTarget{sga3:xvii:lemma:7:3:2:proof:item:i:b}
    \(x^{\ell^n}\in Z^0(k)\).
  \end{enumerate}
  Indeed, it suffices to choose a lift of \(\bar x\) whose image
  in \((N/Z^0)(k)\) belongs to the \(\ell\)-Sylow subgroup.
  Raising to the \(\ell\)-th power in the commutative group
  \(Z^0\) is an étale morphism, so \(Z^0(k)\) is
  \(\ell\)-divisible.  Consequently, after multiplying \(x\) by
  an element of \(Z^0(k)\), we may suppose that
  \(x^{\ell^n}=0\).

% END INCLUDED SGA3_English_current_buildable_source_20260823/source/inputs/XVII/tex/components/32_expose_XVII_theorem7_2_1_completion_through_lemma7_3_2_page041.tex

% BEGIN INCLUDED SGA3_English_current_buildable_source_20260823/source/inputs/XVII/tex/components/33_expose_XVII_lemma7_3_2_completion_and_appendixA_opening_page042.tex
  The group \(N\) is then generated by two commuting commutative
  groups, namely \(Z\) and the group generated by \(x\), and is
  therefore commutative.  The group \({}_{\ell^n}N\) is of
  multiplicative type and is characteristic in \(N\); hence it is
  invariant in \(G\), and consequently central because \(G\) is
  connected (\SGARefLink{sga3:ix:corollary:5:5}{Exposé IX, 5.5}).  Thus \({}_{\ell^n}N\) is contained
  in \(Z\), contradicting the fact that its image in \(G'\)
  contains \(\boldsymbol\mu_\ell\).

\item[(ii)] \SGARefTarget{sga3:xvii:lemma:7:3:2:proof:item:ii}
  Suppose that \(\ell=p\).  There is then an integer \(n\) such
  that the image of \(F_nN\) in \(G'\) contains
  \(\boldsymbol\mu_p\), and hence an exact sequence
  \[
    1\longrightarrow K\longrightarrow F_nN
      \longrightarrow\boldsymbol\mu_p\longrightarrow1. \tag{*}\label{sga3:xvii:lemma:7-3-2:equation:star}
  \]
  The group \(K\) is contained in \(Z\), and is therefore
  commutative.  It follows from \SGARefLink{sga3:xvii:theorem:7:2:1}{7.2.1}, \SGARefLink{sga3:xvii:proposition:7:1:1:proof:ii:item:b}{7.1.1(b)}, and \SGARefLink{sga3:xvii:lemma:5:5:1}{5.5.1} that
  there is a subgroup of multiplicative type in \(F_nN\) whose
  image in the quotient by \(K\) is \(\boldsymbol\mu_p\).  As in
  \SGARefLink{sga3:xvii:subsection:6:4:item:i}{(i)}, we deduce that \(F_nN\) is commutative.  Its
  multiplicative-type component (\SGARefLink{sga3:xvii:theorem:7:2:1}{7.2.1}) is characteristic in
  \(F_nN\), hence invariant in \(G\), and therefore central
  (\SGARefLink{sga3:ix:corollary:5:5}{Exposé IX, 5.5}).  Since its image in \(G'\) contains
  \(\boldsymbol\mu_p\), this is a contradiction.
\end{enumerate}

\SGASection{A}{Appendix I. Hochschild Cohomology and Extensions of
Algebraic Groups}
\SGARefTarget{sga3:xvii:appendix:A}

\SGANextSectionTarget{sga3:xvii:subsection:a-1}\SGASubsection{A.1}{Definition of the Cohomology Groups}
\SGARefTarget{sga3:xvii:appendix:A:subsection:1}

Let \(k\) be a field, let \(G\) be an algebraic \(k\)-group, and
let \((\mathrm{Ab})_G\) be the abelian category of commutative
algebraic \(k\)-groups on which \(G\) acts.  If
\(A\in\operatorname{Ob}(\mathrm{Ab})_G\), the functor
\[
  h_A:(\operatorname{Sch}/k)^\circ\longrightarrow
      \operatorname{Ens}
\]
canonically defined by \(A\) is a \(G\)-\(\mathbb Z\)-module in
the sense of \SGARefLink{sga3:i:definition:3-2}{Exposé I, 3.2}.  We may therefore consider the
standard complex \(C^\bullet(G,A)\) of algebraic cochains of \(G\)
with values in \(A\) (\SGARefLink{sga3:i:subsection:5-1}{Exposé I, 5.1}), together with the group
\(Z^i(G,A)\) of \(i\)-cocycles, the group \(B^i(G,A)\) of
\(i\)-coboundaries, and the \(i\)-th cohomology group
\(H^i(G,A)\).  As usual, \(H^0(G,A)\) is identified with the
group \(A^G(k)\), where \(A^G\) is the \(k\)-functor of invariants
of \(A\) under \(G\).  The group \(H^1(G,A)\) classifies the
trivial principal homogeneous spaces under \(A\) on which \(G\)
acts.

In general, the functor
\[
  A\longmapsto H^\bullet(G,A)
\]
from \((\mathrm{Ab})_G\) to \(\mathrm{Ab}\) is not a
cohomological functor.  Nevertheless, we have the following
proposition.

\SGAStatement{\SGAInlineRefTarget{sga3:xvii:proposition:a-1-1}Proposition A.1.1}{Let
\[
  0\longrightarrow A'\xrightarrow{u}A\xrightarrow{v}A''
    \longrightarrow0
\]
be an exact sequence in \((\mathrm{Ab})_G\).  Then:
\begin{enumerate}
\item[(a)] \SGARefTarget{sga3:xvii:proposition:A:1:1:item:a}
  if \(v\) has a section---that is, if there is a \(k\)-morphism
  of preschemes \(s:A''\to A\) such that
  \(v s=1_{A''}\)---then there is the usual exact cohomology
  sequence
  \[
    \cdots\longrightarrow H^i(G,A)\longrightarrow H^i(G,A'')
      \xrightarrow{d}H^{i+1}(G,A')\longrightarrow\cdots;
  \]
\item[(b)] \SGARefTarget{sga3:xvii:proposition:A:1:1:item:b}
  if \(A(k)\to A''(k)\) is surjective, then the sequence
  \[
    0\longrightarrow A'^G(k)\longrightarrow A^G(k)
      \longrightarrow A''^G(k)\xrightarrow{d}H^1(G,A')
      \longrightarrow H^1(G,A)\longrightarrow H^1(G,A'')
      \tag{1}\label{sga3:xvii:proposition:a-1-1:equation:1}
  \]
  is exact.
\end{enumerate}}
\SGARefTarget{sga3:xvii:proposition:A:1:1}

\SGAProof{Proof}
For \SGARefLink{sga3:xvii:appendix:A}{(a)}, observe that the existence of a section implies the
exactness of the sequence of complexes
\[
  0\longrightarrow C^\bullet(G,A')\longrightarrow C^\bullet(G,A)
    \longrightarrow C^\bullet(G,A'')\longrightarrow0.
\]

% END INCLUDED SGA3_English_current_buildable_source_20260823/source/inputs/XVII/tex/components/33_expose_XVII_lemma7_3_2_completion_and_appendixA_opening_page042.tex

% BEGIN INCLUDED SGA3_English_current_buildable_source_20260823/source/inputs/XVII/tex/components/34_expose_XVII_appendixA_1_completion_and_A_2_page043.tex
For (b), if \(x''\in A''{}^G(k)\), its inverse image in \(A\)
is a trivial principal homogeneous space under \(A'\) (because
\(A(k)\to A''(k)\) is assumed to be surjective) on which \(G\)
acts, and therefore defines an element
\(d(x'')\in H^1(G,A')\).  The exactness of the sequence (1) is
then immediate.

\SGANextSectionTarget{sga3:xvii:subsection:a-2}\SGASubsection{A.2}{The Group
\texorpdfstring{\(\operatorname{Ext}_{\mathrm{alg}}(G,A)\)}
{Ext alg (G,A)}}
\SGARefTarget{sga3:xvii:appendix:A:subsection:2}

Let \(A\) and \(G\) be two algebraic \(k\)-groups, and let \(E\)
and \(E'\) be two algebraic \(k\)-groups that are extensions of
\(G\) by \(A\).  These two extensions are said to be
\emph{isomorphic} if there is a \(k\)-morphism of groups
\(u:E\to E'\) that makes the following diagram commutative:
\SGARefTarget{sga3:xvii:appendix:A:diagram:A:2:isomorphic-extensions}
\[
\begin{tikzcd}[column sep=large,row sep=large]
  {} & E \arrow[dr] \arrow[d,"u"] & {} \\
  A \arrow[ur] \arrow[dr] & {} & G \\
  {} & E' \arrow[ur] & {}
\end{tikzcd}
\qquad .
\]

The group \(E\) acts on \(A\) by inner automorphisms.  If \(A\)
is commutative, this action factors through \(G\), so that \(G\)
acts on \(A\).  Conversely, if
\(A\in\operatorname{Ob}(\mathrm{Ab})_G\), we denote by
\(\operatorname{Ext}_{\mathrm{alg}}(G,A)\) the set of classes of
algebraic extensions \(E\) of \(G\) by \(A\) for which the
action of \(G\) on \(A\) defined by \(E\) agrees with the action
coming from the structure of \(A\) as an object of
\((\mathrm{Ab})_G\).

The set \(\operatorname{Ext}_{\mathrm{alg}}(G,A)\) is naturally a
bifunctor, covariant in \(A\) and contravariant in \(G\).  More
precisely:
\begin{enumerate}
\item[(a)] \SGARefTarget{sga3:xvii:appendix:A:subsection:2:item:a}
  if \(f:A\to B\) is a morphism in \((\mathrm{Ab})_G\), and if
  \(E\) represents an element of
  \(\operatorname{Ext}_{\mathrm{alg}}(G,A)\), then
  \(f_*(E)\in\operatorname{Ext}_{\mathrm{alg}}(G,B)\) is defined
  to be the class of the extension of \(G\) by \(B\) equal to
  the quotient of the semidirect product \(B\mathbin{\cdot}E\)
  (where \(E\) acts on \(B\) through \(G\)) by the algebraic
  subgroup that is the image of \(A\) under the morphism
  \((f,i)\).  This quotient is representable by Exposé VI\textsubscript{A},
  \S5, and there is therefore a commutative diagram
  \SGARefTarget{sga3:xvii:appendix:A:diagram:A:2:pushout}
  \[
  \begin{tikzcd}[column sep=large,row sep=large]
    1 \arrow[r] &
    A \arrow[r,"i"] \arrow[d,"f"'] &
    E \arrow[r] \arrow[d] &
    G \arrow[r] \arrow[d,"1_G"] &
    1 \\
    1 \arrow[r] &
    B \arrow[r] &
    f_*(E) \arrow[r] &
    G \arrow[r] &
    1
  \end{tikzcd}
  \qquad .
  \]

\item[(b)] \SGARefTarget{sga3:xvii:appendix:A:subsection:2:item:b}
  if \(g:H\to G\) is a \(k\)-morphism of algebraic
  \(k\)-groups, and if \(E\) is an extension of \(G\) by \(A\),
  the fiber product \(E\times_G H\) is naturally an extension
  of \(H\) by \(A\), denoted by \(g^*(E)\).  Thus there is a
  commutative diagram
  \SGARefTarget{sga3:xvii:appendix:A:diagram:A:2:pullback}
  \[
  \begin{tikzcd}[column sep=large,row sep=large]
    1 \arrow[r] &
    A \arrow[r] \arrow[d,"1_A"'] &
    g^*(E) \arrow[r] \arrow[d] &
    H \arrow[r] \arrow[d,"g"] &
    1 \\
    1 \arrow[r] &
    A \arrow[r] &
    E \arrow[r] &
    G \arrow[r] &
    1
  \end{tikzcd}
  \qquad .
  \]
\end{enumerate}

By adapting the proofs given in J.-P. Serre,
\emph{Groupes algébriques et corps de classe}, Chapter VII, one
endows \(\operatorname{Ext}_{\mathrm{alg}}(G,A)\) with a natural
abelian-group structure, functorial in \(A\) and \(G\).

% END INCLUDED SGA3_English_current_buildable_source_20260823/source/inputs/XVII/tex/components/34_expose_XVII_appendixA_1_completion_and_A_2_page043.tex

% BEGIN INCLUDED SGA3_English_current_buildable_source_20260823/source/inputs/XVII/tex/components/35_expose_XVII_appendixA_sections2_3_and_appendixB_opening_page044.tex
\SGAStatement{\SGAInlineRefTarget{sga3:xvii:proposition:a-2-1}Proposition A.2.1}{Let
\[
  0\longrightarrow A'\longrightarrow A\longrightarrow A''
    \longrightarrow0
\]
be an exact sequence in \((\mathrm{Ab})_G\).  Then:
\begin{enumerate}
\item[(a)] \SGARefTarget{sga3:xvii:proposition:A:2:1:item:a}
  there is a canonical exact sequence of abelian groups
  \[
    Z^1(G,A)\longrightarrow Z^1(G,A'')
      \xrightarrow{d}\operatorname{Ext}_{\mathrm{alg}}(G,A')
      \longrightarrow\operatorname{Ext}_{\mathrm{alg}}(G,A)
      \longrightarrow\operatorname{Ext}_{\mathrm{alg}}(G,A'');
  \]
\item[(b)] \SGARefTarget{sga3:xvii:proposition:A:2:1:item:b}
  if \(A(k)\to A''(k)\) is surjective, one deduces from \SGARefLink{sga3:xvii:appendix:A}{(a)} the
  exact sequence
  \[
  \tag{2}\label{sga3:xvii:proposition:a-2-1:equation:2}
    H^1(G,A)\longrightarrow H^1(G,A'')
      \xrightarrow{d}\operatorname{Ext}_{\mathrm{alg}}(G,A')
      \longrightarrow\operatorname{Ext}_{\mathrm{alg}}(G,A)
      \longrightarrow\operatorname{Ext}_{\mathrm{alg}}(G,A'').
  \]
\end{enumerate}}
\SGARefTarget{sga3:xvii:proposition:A:2:1}

The exact sequence in \SGARefLink{sga3:xvii:appendix:A}{(a)} generalizes the usual exact
\(\operatorname{Hom}(\,\cdot\,,\,\cdot\,)\) sequence valid for
commutative extensions (\emph{loc.\ cit.}), and is proved in the
same way.  We recall only the definition of the coboundary
\[
  d:Z^1(G,A'')\longrightarrow
    \operatorname{Ext}_{\mathrm{alg}}(G,A').
\]
For this purpose, consider the extension
\[
  1\longrightarrow A'\longrightarrow A\mathbin{\cdot}G
    \longrightarrow A''\mathbin{\cdot}G\longrightarrow1,
\]
obtained in the evident way from the exact sequence
\(0\to A'\to A\to A''\to0\).  If \(u\in Z^1(G,A'')\), then
\(u\) defines in the usual way a section homomorphism
\[
  u:G\longrightarrow A''\mathbin{\cdot}G.
\]
One then has
\[
  d(u)=u^*(A\mathbin{\cdot}G).
\]

\SGANextSectionTarget{sga3:xvii:subsection:a-3}\SGASubsection{A.3}{Comparison of
\texorpdfstring{\(H^2(G,A)\) and
\(\operatorname{Ext}_{\mathrm{alg}}(G,A)\)}
{H2(G,A) and Ext alg(G,A)}}
\SGARefTarget{sga3:xvii:appendix:A:subsection:3}

It is well known that, for abstract groups, there is a functorial
isomorphism between the abelian groups \(H^2(G,A)\) and
\(\operatorname{Ext}(G,A)\).  Likewise in the present setting,
if \(A\) is an object of \((\mathrm{Ab})_G\), one can associate
with every \(2\)-cocycle \(u\in Z^2(G,A)\) an algebraic-group
structure on the prescheme \(A\times_kG\), making it an element
of \(\operatorname{Ext}_{\mathrm{alg}}(G,A)\).  Moreover, this
extension is trivial if and only if \(u\in B^2(G,A)\) (see
\SGARefLink{sga3:iii:lemma:1-2-2}{Exposé III, 1.2.2}).  Recall that the composition law on
\(A\times G\) is defined by
\[
  (a,g)(a',g')=a+{}^g a'+u(g,g').
\]

It is clear that the extensions of \(G\) by \(A\) obtained in
this way are not arbitrary, since they possess a section.
Conversely, however, if
\(E\in\operatorname{Ext}_{\mathrm{alg}}(G,A)\) possesses a
section \(s\), then \(E\) is isomorphic to the extension of
\(G\) by \(A\) associated with the \(2\)-cocycle \(u\) such that
\[
  u(g,g')=s(g)s(g')s(gg')^{-1}.
\]
We thus obtain the following proposition.

\SGAStatement{\SGAInlineRefTarget{sga3:xvii:proposition:a-3-1}Proposition A.3.1}{There is a functorial
isomorphism between the bifunctors with values in abelian
groups
\[
  (G,A)\longmapsto H^2(G,A)
  \qquad\text{and}\qquad
  (G,A)\longmapsto\operatorname{Ext}_s(G,A),
\]
where \(\operatorname{Ext}_s(G,A)\) denotes the subgroup of
\(\operatorname{Ext}_{\mathrm{alg}}(G,A)\) formed by the classes
of extensions of \(G\) by \(A\) that possess a section.}
\SGARefTarget{sga3:xvii:proposition:A:3:1}

\SGANextSectionTarget{sga3:xvii:section:b}\SGASection{B}{Appendix II. Review and Supplements on Radicial
Groups}
\SGARefTarget{sga3:xvii:appendix:B}

Let \(p>1\) be a prime number, and let \(S\) be an
\(\mathbb F_p\)-prescheme.

% END INCLUDED SGA3_English_current_buildable_source_20260823/source/inputs/XVII/tex/components/35_expose_XVII_appendixA_sections2_3_and_appendixB_opening_page044.tex

% BEGIN INCLUDED SGA3_English_current_buildable_source_20260823/source/inputs/XVII/tex/components/36_expose_XVII_appendixB_pages045_046_and_appendixC_opening.tex
\SGASubsection{B.1}{The Frobenius morphism}
\SGARefTarget{sga3:xvii:appendix:B:subsection:1}

For every \(S\)-prescheme \(X\) and every integer \(n>0\), let
\[
  P^n:X\longrightarrow X
\]
denote the \(\mathbb F_p\)-endomorphism corresponding to raising
sections of \(\mathcal O_X\) to the \(p^n\)-th power, and let
\(X^{(n)}\) denote the \(S\)-prescheme obtained from \(X\) by base
change along \(P^n:S\to S\).  There is then a unique
\(S\)-morphism
\[
  F^n:X\longrightarrow X^{(n)}
\]
that makes the following diagram commutative:
\[
\begin{tikzcd}[column sep=3.7em,row sep=3.0em]
  {} & {} & X
    \arrow[dll,"P^n"']
    \arrow[dl,"F^n"]
    \arrow[ddl] \\
  X \arrow[d] &
  X^{(n)} \arrow[l] \arrow[d] &
  {} \\
  S &
  S \arrow[l,"P^n"] &
  {}
\end{tikzcd}
\]
% Loop-1 native reconstruction.  This diagram is included in the
% Session-B exhaustive no-overwrite Loop-2 high-zoom successor set.

It is clear that \(F^n\) identifies with the ``\(n\)-th iterate''
of \(F^1=F\), called the Frobenius morphism of \(X/S\).

If \(G\) is an \(S\)-group prescheme, then \(G^{(n)}\) is an
\(S\)-group prescheme and
\[
  F^n:G\longrightarrow G^{(n)}
\]
is an \(S\)-group morphism.  Its kernel
\({}_{F^n}(G)\) is a characteristic subgroup prescheme of \(G\)
(that is, it is stable under the functor
\(\operatorname{Aut}_{S\text{-}\mathrm{gr}}(G)\)), and is radicial
over \(S\).  If \(G\) is a radicial \(S\)-group prescheme, we say
that \(G\) has height at most \(h\) if
\({}_{F^h}(G)=G\).

\SGASubsection{B.2}{Groups and restricted
\texorpdfstring{\(p\)-Lie algebras}{p-Lie algebras}}
\SGARefTarget{sga3:xvii:appendix:B:subsection:2}

If \(G\) is an \(S\)-group prescheme, then
\(\operatorname{Lie}(G)\) (Exposé II) is naturally a restricted
\(p\)-Lie algebra (Exposé VII\textsubscript{A}, §§5 and 6).  In
particular, one has the following result (cf. Exposé
VII\textsubscript{A}).

\SGAStatement{\SGAInlineRefTarget{sga3:xvii:proposition:b-2-1}Proposition B.2.1}{%
\begin{enumerate}
\item[(i)] \(\operatorname{Lie}(\boldsymbol\alpha_p)_S
=\operatorname{Lie}(\mathbb G_a)_S\)\footnote{Editor's note:
  the definition of \(\boldsymbol\alpha_p\), already given at the
  beginning of this Exposé, has been omitted here.}
  is a sheaf of \(\mathcal O_S\)-modules, free of rank \(1\) over
  \(\mathcal O_S\), generated by an element \(X\) such that
  \(X^{(p)}=0\).
\item[(ii)] \(\operatorname{Lie}(\boldsymbol\mu_p)_S
  =\operatorname{Lie}(\mathbb G_m)_S\) is a sheaf of
  \(\mathcal O_S\)-modules, free of rank \(1\) over
  \(\mathcal O_S\), with a canonical basis \(X\) such that
  \(X^{(p)}=X\).
\end{enumerate}}
\SGARefTarget{sga3:xvii:proposition:B:2:1}

We now recall the fundamental result proved in Exposé
VII\textsubscript{A}, §7.

\SGAStatement{\SGAInlineRefTarget{sga3:xvii:theorem:b-2-2}Theorem B.2.2}{Suppose that \(S\) is affine, with
ring \(A\).  Then the functor
\[
  G\longmapsto\operatorname{Lie}G
\]
defines an equivalence between, on the one hand, the category of
\(S\)-group preschemes \(G\) that are of finite presentation and
flat over \(S\), of height \(1\), and whose Lie algebra is locally
free over \(\mathcal O_S\), and, on the other hand, the category
of restricted \(p\)-\(A\)-Lie algebras that are locally free of
finite rank.

Moreover, if \(G\) is as above and \(H\) is any \(S\)-group
prescheme of finite presentation, the canonical morphism
\[
  \operatorname{Hom}_{S\text{-}\mathrm{gr}}(G,H)
  \longrightarrow
  \operatorname{Hom}_{p\text{-}A\text{-}\mathrm{Lie}}
    \bigl(\operatorname{Lie}G(S),\operatorname{Lie}H(S)\bigr)
\]
is an isomorphism.}
\SGARefTarget{sga3:xvii:theorem:B:2:2}

\SGANextSectionTarget{sga3:xvii:subsection:b-3}\SGASubsection{B.3}{Radicial groups and smooth groups}
\SGARefTarget{sga3:xvii:appendix:B:subsection:3}

We now suppose that \(S\) is the spectrum of a field \(k\) of
characteristic \(p\).  Recall that Exposé
VI\textsubscript{A}, §5 proves that, if \(G\) is an algebraic
\(k\)-group and \(H\) an algebraic subgroup of \(G\), then the
sheaf \(G/H\) (for the fpqc topology) is representable.  We also
recall (VII\textsubscript{A}, \SGARefLink{sga3:VIIA:8.3}{8.3}):

\SGAStatement{\SGAInlineRefTarget{sga3:xvii:proposition:b-3-1}Proposition B.3.1}{Let \(G\) be an algebraic
\(k\)-group.  There exists an integer \(m\) such that, for every
\(n\geq m\), the algebraic group
\[
  G/{}_{F^n}(G)
\]
is smooth over \(k\).}
\SGARefTarget{sga3:xvii:proposition:B:3:1}

\SGAStatement{Proposition B.3.2}{Consider an exact sequence of
algebraic \(k\)-groups
\[
  1\longrightarrow G'\longrightarrow G
    \xrightarrow{u}G''\longrightarrow1
\]
and the following assertions:
\begin{enumerate}
\item[(i)] \SGARefTarget{sga3:xvii:proposition:B:3:2:item:i}
  the morphism \(u\) is smooth;
\item[(ii)] \SGARefTarget{sga3:xvii:proposition:B:3:2:item:ii}
  \(G'\) is smooth over \(k\);
\item[(iii)] \SGARefTarget{sga3:xvii:proposition:B:3:2:item:iii}
  for every integer \(n>0\), the sequence
  \[
    1\longrightarrow{}_{F^n}(G')
      \longrightarrow{}_{F^n}(G)
      \longrightarrow{}_{F^n}(G'')
      \longrightarrow1
  \]
  is exact;
\item[(iv)] \SGARefTarget{sga3:xvii:proposition:B:3:2:item:iv}
  the morphism \({}_FG\to{}_FG''\) is an epimorphism;
\item[(v)] \SGARefTarget{sga3:xvii:proposition:B:3:2:item:v}
  the morphism
  \[
    \bigl(\operatorname{Lie}G\bigr)(k)
      \longrightarrow\bigl(\operatorname{Lie}G''\bigr)(k)
  \]
  is surjective.
\end{enumerate}
Then
\[
  \text{\rm(i)}\Longleftrightarrow\text{\rm(ii)}
  \Longrightarrow\text{\rm(iii)}
  \Longrightarrow\text{\rm(iv)}
  \Longleftrightarrow\text{\rm(v)}.
\]
Moreover, if \(G\) is smooth, all five assertions are equivalent.}
\SGARefTarget{sga3:xvii:proposition:B:3:2}

The equivalence \SGARefLink{sga3:xvii:proposition:B:3:2:item:i}{(i)}\(\Longleftrightarrow\)\SGARefLink{sga3:xvii:proposition:B:3:2:item:ii}{(ii)} follows from
Exposé VI\textsubscript{B}, \SGARefLink{concl:sga3-VIB-9.2-vii}{9.2(vii)}.

For \SGARefLink{sga3:xvii:proposition:B:3:2:item:ii}{(ii)}\(\Longrightarrow\)\SGARefLink{sga3:xvii:proposition:B:3:2:item:iii}{(iii)}, if \(G'\) is smooth, then
\[
  F^n:G'\longrightarrow G'^{(n)}
\]
is an epimorphism, and \SGARefLink{sga3:xvii:proposition:B:3:2:item:iii}{(iii)} follows from the snake diagram.
The implication \SGARefLink{sga3:xvii:proposition:B:3:2:item:iii}{(iii)}\(\Longrightarrow\)\SGARefLink{sga3:xvii:proposition:B:3:2:item:iv}{(iv)} is clear, while
\SGARefLink{sga3:xvii:proposition:B:3:2:item:iv}{(iv)}\(\Longleftrightarrow\)\SGARefLink{sga3:xvii:proposition:B:3:2:item:v}{(v)} follows from \SGARefLink{sga3:xvii:theorem:b-2-2}{Theorem B.2.2}.
Finally, suppose that \(G\) is smooth and that \SGARefLink{sga3:xvii:proposition:B:3:2:item:v}{(v)} holds.  Then
\(G''\) is smooth, and \SGARefLink{sga3:xvii:proposition:B:3:2:item:v}{(v)} implies
\[
  \dim G'=\dim_k\bigl(\operatorname{Lie}G'\bigr)(k),
\]
so \(G'\) is smooth.

\SGASection{C}{Appendix III. Remarks and Supplements Concerning
Exposés XV, XVI, and XVII}
\SGARefTarget{sga3:xvii:appendix:C}

\SGANextSectionTarget{sga3:xvii:subsection:c-1}\SGASubsection{C.1}{}
\SGARefTarget{sga3:xvii:appendix:C:subsection:1}

It may be that \SGARefLink{sga3:xvii:proposition:1:2}{Propositions 1.2} and \SGARefLink{sga3:xv:proposition:1-2-bis}{1.2 bis} of Exposé XV remain
valid if the hypothesis that \(H_0\) is smooth over \(S_0\) is
removed.  This is the case, in particular, if \(G\) is finite,
flat, and commutative.

% END INCLUDED SGA3_English_current_buildable_source_20260823/source/inputs/XVII/tex/components/36_expose_XVII_appendixB_pages045_046_and_appendixC_opening.tex

% BEGIN INCLUDED SGA3_English_current_buildable_source_20260823/source/inputs/XVII/tex/components/37_expose_XVII_appendixC_sections2_through5_and_end_pages047_049.tex
\SGASubsection{C.2}{Complement to \SGARefLink{sga3:xv:remark:4:8}{Exposé XV, 4.8}}
\SGARefTarget{sga3:xvii:appendix:C:subsection:2}

The following proposition, as well as \SGARefLink{sga3:xvii:theorem:c-3-1}{Theorems C.3.1} and \SGARefLink{sga3:xvii:theorem:c-4-1}{C.4.1}
below, will appear in a forthcoming article by M. Raynaud on
group schemes over a discrete valuation ring.\footnote{Editor's
note: comments are to be added here, including references to
Exposé VI\textsubscript{B}, \ldots}

\SGAStatement{\SGAInlineRefTarget{sga3:xvii:proposition:c-2-1}Proposition C.2.1}{Let \(S\) be the spectrum of a
discrete valuation ring, let \(t\) be its generic point, let \(G\)
be a flat \(S\)-group prescheme of finite type, and put
\[
  \widetilde G=\operatorname{Spec}\Gamma(G,\mathcal O_G).
\]
Let \(u:G\to\widetilde G\) be the canonical morphism.  Then:
\begin{enumerate}
\item[(1)] \SGARefTarget{sga3:xvii:proposition:C:2:1:item:1}
  \(\widetilde G\) is naturally an \(S\)-group scheme, and \(u\)
  is a homomorphism.
\item[(2)] \SGARefTarget{sga3:xvii:proposition:C:2:1:item:2}
  \(\ker(u)\) is flat over \(S\), and
  \((\widetilde G,u)\) is the fpqc quotient of \(G\) by
  \(\ker(u)\).  Thus \(\widetilde G\) is the largest affine
  quotient of \(G\).
\item[(3)] \SGARefTarget{sga3:xvii:proposition:C:2:1:item:3}
  If \(G_t\) is affine, then \(\ker(u)\) is an étale group over
  \(S\), and it is the trivial group if and only if \(G\) is
  separated over \(S\).  In particular, a flat, separated
  \(S\)-group scheme of finite type whose generic fiber is affine
  is affine.
\end{enumerate}}
\SGARefTarget{sga3:xvii:proposition:C:2:1}

\SGANextSectionTarget{sga3:xvii:subsection:c-3}\SGASubsection{C.3}{}
\SGARefTarget{sga3:xvii:appendix:C:subsection:3}

In the statement of \SGARefLink{sga3:xv:lemma:6:6}{Exposé XV, 6.6}, the hypothesis that \(H\)
be equal to its connected normalizer is unnecessary in order that the
set of points \(s\) of \(S\) for which \(H_s\) is a parabolic
subgroup of \(G_s\) be open.  Indeed, repeat the proof given in
paragraph~c) preceding \SGARefLink{sga3:xv:lemma:6:9}{Lemma 6.9}, and let \(N\) denote the
scheme-theoretic closure of \(N_t\) in \(G\).  Then \(N\) is a
closed subgroup scheme of \(G\), flat over \(S\), which contains
\(H\) and is contained in \(\mathcal N\).  The theorem below
shows that \(G/N\) is representable.  Since \(G_s/N_s\) is proper
and connected, the proof concludes as in \emph{loc.\ cit.}

\SGAStatement{\SGAInlineRefTarget{sga3:xvii:theorem:c-3-1}Theorem C.3.1\footnote{Editor's note: see \SGARefLink{sga3:xvii:appendix:C:subsection:4}{Theorem
4}C in S. Anantharaman, ``Schémas en groupes, espaces homogènes et
espaces algébriques sur une base de dimension 1,'' \emph{Bull.
Soc. Math. France, Mém.} 33 (1976), 5--79.}}{Let \(S\) be the
spectrum of a henselian discrete valuation ring, let \(G\) be an
\(S\)-group prescheme locally of finite type, and let \(H\) be a
closed subgroup prescheme of \(G\), flat over \(S\).  Then \(G/H\)
is representable.}
\SGARefTarget{sga3:xvii:theorem:C:3:1}

\SGANextSectionTarget{sga3:xvii:subsection:c-4}\SGASubsection{C.4}{Complement to \SGARefLink{sga3:xvi:subsection:1:1}{Exposé XVI, 1.1}}
\SGARefTarget{sga3:xvii:appendix:C:subsection:4}

The following theorem refines \SGARefLink{sga3:viii:proposition:7:9}{Exposé VIII, 7.9}.

\SGAStatement{\SGAInlineRefTarget{sga3:xvii:theorem:c-4-1}Theorem C.4.1}{Let \(S\) be the spectrum of a
discrete valuation ring with residue characteristic \(p>0\), and
let \(G\) be a commutative, smooth, separated \(S\)-group scheme
of finite type.  Then the following conditions are equivalent:
\begin{enumerate}
\item[(i)] \SGARefTarget{sga3:xvii:theorem:C:4:1:item:i}
  \({}_pG\) is finite over \(S\).
\item[(ii)] \SGARefTarget{sga3:xvii:theorem:C:4:1:item:ii}
  For every \(S\)-prescheme \(S'\), every separated
  \(S'\)-group prescheme \(H'\) of finite presentation over
  \(S'\), and every \(S'\)-monomorphism
  \[
    u:G_{S'}\longrightarrow H',
  \]
  the morphism \(u\) is an immersion.
\end{enumerate}}
\SGARefTarget{sga3:xvii:theorem:C:4:1}

\SGANextSectionTarget{sga3:xvii:subsection:c-5}\SGASubsection{C.5}{}
\SGARefTarget{sga3:xvii:appendix:C:subsection:5}

Example XVII, \SGARefLink{sga3:xvii:subsection:6:4:item:a}{6.4(a)}, gives a smooth group over a field \(k\)
whose unipotent radical is not defined over \(k\).  The following
proposition gives a general method for obtaining such groups.

\SGAStatement{\SGAInlineRefTarget{sga3:xvii:proposition:c-5-1}Proposition C.5.1}{Let \(k\) be a field, let \(K\)
be a finite radicial extension of \(k\) of degree greater than
\(1\), and let \(H\) be a smooth, connected \(K\)-algebraic group
of dimension \(r\).  Put
\[
  G=\prod_{K/k}H/K,
\]
which is a smooth, connected \(k\)-algebraic group of dimension
\(r[K:k]\).  Let
\[
  u:G_K\longrightarrow H
\]
be the canonical homomorphism, and put \(R=\ker(u)\).  Then:
\begin{enumerate}
\item[(i)] \SGARefTarget{sga3:xvii:proposition:C:5:1:item:i}
  \(u\) is an epimorphism, and \(R\) is a smooth, connected,
  unipotent algebraic group.
\item[(ii)] \SGARefTarget{sga3:xvii:proposition:C:5:1:item:ii}
  If \(U\) is a smooth algebraic subgroup of \(G\) such that
  \(U_K\) contains \(R\), then \(U=G\).
\end{enumerate}}
\SGARefTarget{sga3:xvii:proposition:C:5:1}

\SGAStatement{Corollary}{Retain the notation of the preceding
proposition.
\begin{enumerate}
\item[(a)] \SGARefTarget{sga3:xvii:corollary:C:5:1:item:a}
  If \(H\) is not unipotent, then the unipotent radical of
  \(G_{\bar k}\) is not defined over \(k\).
\item[(b)] \SGARefTarget{sga3:xvii:corollary:C:5:1:item:b}
  If \(H\) is a nonzero abelian variety, then \(G\) is not an
  extension of an abelian variety by a smooth linear group.
\item[(c)] \SGARefTarget{sga3:xvii:corollary:C:5:1:item:c}
  If \(H\) is not solvable, then the solvable radical of
  \(G_{\bar k}\) is not defined over \(k\), and \(G\) has no
  Borel subgroup defined over \(k\).
\end{enumerate}}
\SGARefTarget{sga3:xvii:corollary:C:5:1}

Let us first show how the corollary follows from the proposition.

\smallskip
\noindent\textit{\SGARefLink{sga3:xvii:appendix:A}{(a)}} Let \(U\) be a unipotent radical of \(G\).
Then \(U_K\) is the unipotent radical of \(G_K\), and hence
contains \(R\), since \(R\) is smooth, connected, and unipotent
by \SGARefLink{sga3:xvii:corollary:C:5:1:item:i}{(i)}.  Thus \(U=G\) by \SGARefLink{sga3:xvii:corollary:C:5:1:item:ii}{(ii)}.  But \(G_K\) has \(H\) as a
quotient, and \(H\) is not unipotent by hypothesis.  Therefore
\(G\) is not unipotent, a contradiction.

\smallskip
\noindent\textit{\SGARefLink{sga3:xvii:section:b}{(b)}} Suppose that \(G\) is an extension of an
abelian variety by a linear group \(L\) that is smooth over
\(k\).  Then \(L_K\) necessarily contains \(R\), so \(L=G\).
But \(G\) cannot be a linear group, because \(G_K\) has a nonzero
abelian variety as a quotient.

\smallskip
\noindent\textit{(c)} Let \(S\) be a solvable radical of \(G\),
and let \(B\) be a Borel subgroup of \(G\).  Then \(S_K\) and
\(B_K\) contain \(R\), so \(S=G\) and \(B=G\).  But then \(G\)
is solvable, contradicting the fact that \(G_K\) has a quotient
that is not solvable.

\SGAProof{Proof of \SGARefLink{sga3:xvii:proposition:c-5-1}{Proposition C.5.1}}
\textit{\SGAInlineRefTarget{sga3:xvii:corollary:C:5:1:item:i}(i)} We begin by describing the morphism \(u\).  A
\(K\)-scheme
\[
  f:S\longrightarrow\operatorname{Spec}(K)
\]
gives rise to the diagram
\[
\begin{tikzcd}[column sep=6.0em,row sep=4.0em]
  \operatorname{Spec}(K)
    \arrow[d,"j"'] &
  T
    \arrow[l,"h"']
    \arrow[d,bend right=28,"j_T"'] \\
  \operatorname{Spec}(k) &
  S
    \arrow[l,"g"]
    \arrow[ul,"f"]
    \arrow[u,bend right=28,"s"']
\end{tikzcd}
\]
where
\[
  g=j\circ f,\qquad
  T=\operatorname{Spec}(K)\times_{\operatorname{Spec}(k)}S,
\]
\(h\) and \(j_T\) are the two projections, and \(s\) is the
section of \(T\) over \(S\) such that \(h\circ s=f\).  The map
\[
  u(S):G_K(S)\longrightarrow H(S)
\]
is simply the composite
\[
  G_K(S)\xrightarrow{\sim}H(T)\longrightarrow H(S),
\]
where the last arrow is defined by the section \(s\).

Take, in particular, \(S\) to be the spectrum of an algebraic
closure \(\bar k\) of \(k\), and let \(f\) be the unique
\(k\)-morphism induced by \(K\to\bar k\), so that \(T\) is a
local Artinian scheme.  To prove \SGARefLink{sga3:xvii:corollary:C:5:1:item:i}{(i)}, it is enough to do so after
extending the base field from \(K\) to \(\bar k\).  It is clear
that
\[
  G_S=\prod_{T/S}H_T/T
\]
represents the Greenberg functor of \(H_T\) relative to \(S\)
(M. J. Greenberg, ``Schemata over local rings,'' \emph{Ann. of
Math.} 73 (1961), 624--648).  Under this identification, the
description above shows that \(u_S\) is the canonical transition
morphism
\[
  \operatorname{Green}(H_T)\longrightarrow
  H_S=H_T\times_T S.
\]
Assertion \SGARefLink{sga3:xvii:corollary:C:5:1:item:i}{(i)} now follows from the smoothness of \(H\) over \(K\)
and M. J. Greenberg, ``Schemata over local rings II,''
\emph{Ann. of Math.} 78 (1963), 256--266.

\smallskip
\noindent\textit{\SGAInlineRefTarget{sga3:xvii:corollary:C:5:1:item:ii}(ii)} To prove \SGARefLink{sga3:xvii:corollary:C:5:1:item:ii}{(ii)}, we may suppose that \(k\)
is separably closed.  Let \(U\) be a smooth algebraic subgroup of
\(G\) such that \(U_K\supset R\), and let us prove that \(U=G\).
Since \(G\) is connected, we may suppose that \(U\) is connected.
Let \(V=u(U_K)\), a smooth, connected algebraic subgroup of
\(H\), and put
\[
  V'=\prod_{K/k}V/K,
\]
a smooth, connected algebraic subgroup of \(G\).  The group
\(V'_K\) is an algebraic subgroup of \(G_K\), and the canonical
morphism \(V'_K\to V\) is the restriction of \(u\) to \(V'_K\).
By hypothesis, \(U_K\) contains \(R\), hence \emph{a fortiori}
contains \(\ker(V'_K\to V)\); by construction, moreover, the
image of \(U_K\) in \(H\) is \(V\).  It follows that \(U_K\)
contains \(V'_K\), and hence that \(U\) contains \(V'\).

On the other hand, the canonical isomorphism
\[
  G(k)\lhook\joinrel\longrightarrow G(K)
    \xrightarrow{u(K)}H(K)
\]
plainly sends \(U(k)\) into \(V(K)=V'(k)\).  Thus
\(U(k)\subset V'(k)\).  Since \(U\) is smooth and \(k\) is
separably closed, it follows that \(U\subset V'\), whence
\(U=V'\).  We then have
\[
  ([K:k]-1)\dim V
  =\dim\ker(V'_K\to V)
  =\dim R
  =([K:k]-1)\dim H.
\]
Since \(K\ne k\), we conclude that \(\dim V=\dim H\), hence that
\(V=H\), and finally that \(U=V'=G\).

% END INCLUDED SGA3_English_current_buildable_source_20260823/source/inputs/XVII/tex/components/37_expose_XVII_appendixC_sections2_through5_and_end_pages047_049.tex



% END INCLUDED SGA3_English_current_buildable_source_20260823/source/inputs/XVII/body.tex


\SGAStartExpose{XVIII}
% BEGIN INCLUDED SGA3_English_current_buildable_source_20260823/source/inputs/XVIII/body.tex
% Mechanically extracted after \mainmatter from the frozen standalone master.

\markboth{Exposé XVIII. Weil's Theorem}
  {Exposé XVIII. Weil's Theorem}
% BEGIN INCLUDED SGA3_English_current_buildable_source_20260823/source/inputs/XVIII/tex/components/01_expose_XVIII_introduction_and_section1.tex
% SGA 3, Exposé XVIII, Loop-1 English reconstruction.
% Current scope: title, Introduction, and authority local page 1 opening of §1.
% Authority: Polo--Gille Expo18.pdf, local page 1
% (combined-reader page 1013; printed pages 632--633).
% Local GPU OCR and the Reinhold English Markdown are comparison/locator
% witnesses only. The Polo--Gille page image controls text and formulas.

\setcounter{footnote}{0}
\SGARefTarget{sga3:xviii:chapter:xviii}\chapter*{Exposé XVIII\\
Weil's Theorem on the Construction of a Group from a Rational Law}
\addcontentsline{toc}{chapter}{Exposé XVIII. Weil's Theorem on the
Construction of a Group from a Rational Law}
\markboth{EXPOSÉ XVIII. WEIL'S THEOREM}
  {EXPOSÉ XVIII. WEIL'S THEOREM}

\begin{center}
  \textit{by Michael Artin}
\end{center}
\begingroup
\renewcommand{\thefootnote}{0}
\footnotetext{Version of May 8, 2009.}
\endgroup

\SGARefTarget{sga3:xviii:section:0}\section*{0. Introduction}
\addcontentsline{toc}{section}{0. Introduction}

This Exposé is devoted to the well-known theorem of Weil~[1], which
constructs a group variety from a birational law. It seems that the
generalization of this result to the case where the base is the
spectrum of a discrete valuation ring was already known to several
people; see, for example,~[2]. Here we shall prove the theorem for a
flat group of finite presentation over an arbitrary base prescheme
\(S\). Since the hypotheses must be formulated somewhat more
carefully, the statement is not in the form given by Weil; but when
\(S\) is the spectrum of a field, the form given here is seen to be
essentially equivalent to Weil's. We use the following suggestion of
Grothendieck: given a ``group germ'' \(X\) over a prescheme \(S\), one
should construct the group as the quotient of \(X\times_S X\) by a
suitable equivalence relation.

\SGARefTarget{sga3:xviii:section:1}\section*{1. ``Recollections'' on
Rational Maps}
\addcontentsline{toc}{section}{1. ``Recollections'' on Rational Maps}

Let \(X/S\) be a relative prescheme and let \(U\subseteq X\) be an
open subset. We say that \(U\) is \emph{schematically dense in \(X\)
relative to \(S\)}, or that \(U\hookrightarrow X\) is
\emph{relatively schematically dense}, if, for every base change
\(S'\to S\), the open subset \(U\times_S S'\) is schematically dense
in \(X\times_S S'\). For the definition and properties of this
notion, we refer to Exposé IX, \S4.\begingroup
\renewcommand{\thefootnote}{*}\footnote{See also EGA
IV\textsubscript{3}, 11.9 and 11.10 (especially 11.10.8), where the
term ``universally schematically dense relative to \(S\)'' is used.}
\addtocounter{footnote}{-1}
\endgroup

\SGARefTarget{sga3:xviii:prop:1.1}\paragraph{\SGARefLink{sga3:xviii:prop:1.1}{Proposition 1.1}.}
\textit{(i) A finite intersection, as well as a union of a nonempty
family, of open subsets schematically dense relative to \(S\), is
schematically dense relative to \(S\).}

\textit{(ii) If \(U\subseteq X\) is schematically dense relative to
\(S\), and if \(S\to T\) is a morphism, then
\(U_T\subseteq X_T\) is schematically dense relative to \(T\).}

\textit{(iii) Let \(U\subseteq V\subseteq X\) be open immersions. For
\(U\) to be schematically dense in \(X\) relative to \(S\), it is
necessary and sufficient that it be so in \(V\), and that \(V\) be so
in \(X\).}

\textit{\SGAInlineRefTarget{sga3:xviii:prop:1.1:item:iv}(iv) If \(U\subseteq X\) and \(V\subseteq Y\) are relatively
schematically dense, then \(U\times_S V\) is schematically dense in
\(X\times_S Y\) relative to \(S\).}

In Exposé IX one finds the following criterion:

\SGARefTarget{sga3:xviii:prop:1.2}\paragraph{\SGARefLink{sga3:xviii:prop:1.2}{Proposition 1.2}.}
\textit{Let \(X/S\) be locally of finite presentation and flat, and
let \(U\) be an open subset of \(X\). If, for every \(s\in S\), the
fiber \(U_s\) is schematically dense in \(X_s\), then \(U\) is
schematically dense in \(X\) relative to \(S\).}

In particular:

\SGARefTarget{sga3:xviii:cor:1.3}\paragraph{\SGARefLink{sga3:xviii:cor:1.3}{Corollary 1.3}.}
\textit{Suppose that \(X/S\) is locally of finite presentation and
flat, and that every fiber \(X_s\) is ``without embedded
components.'' Then \(U\) is schematically dense in \(X\) relative to
\(S\) if and only if, for every \(s\in S\), \(U_s\) is dense in
\(X_s\) in the topological sense.}

One easily deduces from the definition the following proposition.

\SGARefTarget{sga3:xviii:prop:1.4}\paragraph{\SGARefLink{sga3:xviii:prop:1.4}{Proposition 1.4}.}
\textit{Let \(U\subseteq X\) be schematically dense relative to \(S\),
and let \(f,g:X\to Y\) be two morphisms of functors over \(S\), where
\(Y\) satisfies one of the following conditions:}

\begin{enumerate}
\item[\SGAInlineRefTarget{sga3:xviii:prop:1.4:item:i}(i)] \textit{\(Y\) is a prescheme over \(S\), and every fiber
  \(Y_s\) is separated.}
\item[\SGAInlineRefTarget{sga3:xviii:prop:1.4:item:ii}(ii)] \textit{\(Y\) is a presheaf\footnote{Editor's note:
  specify for which topology; a priori, the fpqc topology.} over
  \(S\), such that the diagonal morphism
  \(Y\to Y\times_S Y\) is representable by a closed
  immersion.\footnote{Editor's note: Let us recall here the
  definition of a closed immersion. A morphism of \(S\)-functors
  \(F\to G\) is \emph{relatively representable} if, for every
  \(S\)-morphism \(T\to G\), where \(T\) is an \(S\)-prescheme, the
  \(T\)-functor \(F_T=F\times_G T\) is representable by a
  \(T\)-prescheme. The French source prints \(F\times_S T\) here;
  the fiber product over \(G\) is the one dictated by the preceding
  morphism \(T\to G\). A morphism of \(S\)-functors \(F\to G\) is an open
  (respectively, closed) immersion if it is relatively representable
  and if, for every \(S\)-morphism \(T\to G\), where \(T\) is an
  \(S\)-prescheme, the \(T\)-morphism \(F_T\to T\) is an open
  (respectively, closed) immersion.}}
\end{enumerate}

\textit{Then, if \(f=g\) on \(U\), one has \(f=g\).}

\textit{Proof.}
In case \SGARefLink{sga3:xviii:prop:1.4:item:i}{(i)}, the standard argument following IX, \SGARefLink{sga3:ix:definition:4:1}{4.1} shows that
\(f\) and \(g\) are equal on every fiber; hence the subprescheme
\(\operatorname{Ker}(f,g)\) of \(X\) is set-theoretically equal to
\(X\), and is therefore closed. Since it contains \(U\), it is equal
to \(X\), that is, \(f=g\). Case \SGARefLink{sga3:xviii:prop:1.4:item:ii}{(ii)} is proved as in
\textit{loc.\ cit.}\footnote{Editor's note: Give details on this
point\ldots}

\SGARefTarget{sga3:xviii:def:1.5}\paragraph{\SGARefLink{sga3:xviii:def:1.5}{Definition 1.5}.}
\begingroup
\renewcommand{\thefootnote}{*}
\footnote{See EGA IV\textsubscript{4}, 20.5, where the term
``pseudo-morphism from \(X\) to \(Y\) relative to \(S\)'' is used in
order not to conflict with EGA I, 7.12.}
\addtocounter{footnote}{-1}
\endgroup
\textit{Let \(X\) be a prescheme over \(S\), and let \(Y\) be a
presheaf over \(S\). A \emph{rational map}
\(f:X\dashrightarrow Y\) over \(S\) is an equivalence class of
morphisms \(f:U\to Y\) over \(S\), where \(U\) is an open subset of
\(X\) that is schematically dense relative to \(S\). One sets
\[
  (f:U\to Y)\sim(f':U'\to Y)
\]
if and only if there exists an open subset
\(U''\subseteq U'\cap U\), schematically dense relative to \(S\),
such that \(f=f'\) on \(U''\).}

This definition is made in such a way that, for every base change
\(S'\to S\), the rational map \(f\times_S S'\) may be defined in the
obvious way.

If \(U\) is an open subset of \(X\), we say that the rational map
\(f\) is \emph{defined on \(U\)} if there exists a morphism
representing \(f\) whose set of definition contains \(U\).

\SGARefTarget{sga3:xviii:def:1.5.1}\paragraph{\SGARefLink{sga3:xviii:def:1.5.1}{Definition 1.5.1}.}
\footnote{Editor's note: The number \SGARefLink{sga3:xviii:def:1.5.1}{1.5.1} has been added.}
\textit{If \(Y\) satisfies one of conditions \SGARefLink{sga3:xviii:prop:1.4:item:i}{(i)}, \SGARefLink{sga3:xviii:prop:1.4:item:ii}{(ii)} of \SGARefLink{sga3:xviii:prop:1.4}{1.4}, it is
clear that there exists a largest open subset \(U\) of \(X\) on
which \(f\) is defined, and this \(U\) is schematically dense
relative to \(S\). It is called the \emph{domain of definition of
\(f\) over \(S\)}, and is denoted \(\Dom(f)\).}

This notion does not in general commute with base change, but one
has:

\SGARefTarget{sga3:xviii:prop:1.6}\paragraph{\SGARefLink{sga3:xviii:prop:1.6}{Proposition 1.6}.}
\textit{Let \(X\) be an \(S\)-prescheme and \(Y\) an \(S\)-functor
satisfying one of hypotheses \SGARefLink{sga3:xviii:prop:1.4:item:i}{(i)}, \SGARefLink{sga3:xviii:prop:1.4:item:ii}{(ii)} of \SGARefLink{sga3:xviii:prop:1.4}{1.4}. Let
\(f:X\dashrightarrow Y\) be a rational map over \(S\), let
\(S'\to S\) be a flat morphism locally of finite presentation, and
put \(f'=f\times_S S'\). Then}
\[
  \Dom(f')=\Dom(f)\times_S S'.
\]

\textit{Proof.}
Put \(U=\Dom(f)\). It is clear that \(V'=\Dom(f')\) contains
\(U\times_S S'\). Let \(V\) be the image of \(V'\); it is an open
subset of \(X\), because \(X'\to X\) is open. We must prove that
\(V=U\), that is, we must find a morphism \(V\to Y\) representing
\(f\). Put \(S''=S'\times_S S'\), and
\[
  X''=X\times_S S'',\qquad
  U''=U\times_S S'',\qquad
  V''=V'\times_V V'.
\]
Then \(U''\) is schematically dense in \(X''\) relative to \(S''\);
hence \(U''\) is schematically dense in \(V''\) relative to
\(S''\), since \(U''\subseteq V''\subseteq X''\). The restriction
of \(f':V'\to Y'\) to \(U'\) is obtained from \(f:U\to Y\) by base
change. The two morphisms \(V''\to Y\) obtained from \(f'\) by the
two projections agree on \(U''\), and hence agree.\begingroup
\renewcommand{\thefootnote}{\ensuremath{\dagger}}\footnote{The
French source prints \(S\), rather than \(S''\), in the preceding
density assertion and the undefined symbol \(g\), rather than
\(f'\), in this descent assertion. The base-change construction
determines the readings used here.}
\addtocounter{footnote}{-1}
\endgroup Since \(V'\to V\) is flat and
locally of finite presentation, it is an fppf covering (\SGARefLink{sga3:iv:subsection:6:3}{Exposé IV,
6.3}), and the morphism \(V\to Y\) is obtained by descent.

We shall frequently use the following triviality:

\SGARefTarget{sga3:xviii:prop:1.7}\paragraph{\SGARefLink{sga3:xviii:prop:1.7}{Proposition 1.7}.}
\textit{Let \(X/S\) be faithfully flat and locally of finite
presentation, and let \(U\subseteq X\) be a relatively schematically
dense open subset. Then there exists an fppf covering base change
\(S'\to S\), and a section \(x\in X(S')\) contained in \(U(S')\).}

Indeed, \(U/S\) is an fppf covering. Take \(S'=U\), and as the section
take the \(S\)-graph of the inclusion \(U\hookrightarrow X\).

% END INCLUDED SGA3_English_current_buildable_source_20260823/source/inputs/XVIII/tex/components/01_expose_XVIII_introduction_and_section1.tex

% BEGIN INCLUDED SGA3_English_current_buildable_source_20260823/source/inputs/XVIII/tex/components/02_expose_XVIII_section2.tex
\SGARefTarget{sga3:xviii:section:2}\section*{2. Local Determination of a
Group Homomorphism}
\addcontentsline{toc}{section}{2. Local Determination of a Group
Homomorphism}

Let \(G\) and \(H\) be groups, and let \(U\subseteq G\) be a subset
such that \(U\cdot U=G\). If \(f\) and \(g\) are homomorphisms from
\(G\) to \(H\) such that \(f=g\) on \(U\), then \(f=g\). Likewise:

\SGARefTarget{sga3:xviii:prop:2.1}\paragraph{\SGARefLink{sga3:xviii:prop:2.1}{Proposition 2.1}.}
\textit{Let \(S\) be a site (cf.\ Exposé IV\footnote{Editor's note:
that is, a category endowed with a topology; cf.\ IV, \S4.2.}),
let \(G\) be a sheaf of groups on \(S\), and let \(U\subseteq G\) be
a subsheaf of sets such that the morphism \(U\times U\to G\)
induced by multiplication is an epimorphism. If \(f\) and \(g\) are
homomorphisms from \(G\) to a sheaf of groups \(H\) and are equal on
\(U\), then \(f=g\).}

\SGARefTarget{sga3:xviii:cor:2.2}\paragraph{\SGARefLink{sga3:xviii:cor:2.2}{Corollary 2.2}.}
\textit{Let \(G\) be an \(S\)-group prescheme that is locally of
finite presentation and flat, let \(U\subseteq G\) be a relatively
schematically dense open subset, and let \(H\) be a sheaf of groups
on \(S\) for the fppf topology. Then every homomorphism
\(f:G\to H\) is determined by its restriction to \(U\).}

Indeed, since \(G\) is flat over \(S\), the composition law on \(G\)
is a flat morphism (VI\textsubscript{B}, \SGARefLink{concl:sga3-VIB-9.2-xi}{9.2(xi)}); it follows that
\(U\times_S U\to G\) is faithfully flat and locally of finite
presentation, hence an fppf covering, and therefore an epimorphism.

\SGARefTarget{sga3:xviii:prop:2.3}\paragraph{\SGARefLink{sga3:xviii:prop:2.3}{Proposition 2.3}.}
\textit{Let \(G\) be a group prescheme locally of finite
presentation and flat over \(S\), let \(U\subseteq G\) be a
relatively schematically dense open subset, and let \(H\) be a
sheaf of groups for the fppf topology. Write \(m_G\) for the
multiplication on \(G\) and \(m_H\) for that on \(H\).}

\textit{Let \(\bar f:U\to H\) be an \(S\)-morphism, and suppose that
there exists a relatively schematically dense open subset \(V\) of
\(m_G^{-1}(U)\cap(U\times_S U)\) such that the following diagram is
commutative:}

\begin{center}
% BEGIN INCLUDED SGA3_English_current_buildable_source_20260823/source/inputs/XVIII/native_diagrams/exp18/001_proposition_2_3_compatibility_square.tex
\begin{tikzcd}[column sep=large,row sep=large]
  V \arrow[r,"{(\bar f\times\bar f)|_V}"] \arrow[d,"m_G"] &
  H\times_S H \arrow[d,"m_H"]\\
  U \arrow[r,"\bar f"] & H
\end{tikzcd}

% END INCLUDED SGA3_English_current_buildable_source_20260823/source/inputs/XVIII/native_diagrams/exp18/001_proposition_2_3_compatibility_square.tex

\end{center}

\textit{Then there exists a morphism of \(S\)-groups
\(f:G\to H\), necessarily unique, extending \(\bar f\), in each of
the following situations:}

\begin{enumerate}
\item[\SGAInlineRefTarget{sga3:xviii:prop:2.3:item:i}(i)] \textit{\(H\) is representable.}
\item[\SGAInlineRefTarget{sga3:xviii:prop:2.3:item:ii}(ii)] \textit{The diagonal morphism
  \(H\to H\times_S H\) is representable by a closed
  immersion.\footnote{Editor's note: See the note at \SGARefLink{sga3:xviii:prop:1.4}{1.4}.}}
\item[\SGAInlineRefTarget{sga3:xviii:prop:2.3:item:iii}(iii)] \textit{For every section \(a\in U(S)\), the open
  subset
  \(\operatorname{pr}_2((a\times_S U)\cap V)\) is relatively
  schematically dense in \(U\)\begingroup
  \renewcommand{\thefootnote}{*}\footnote{Or in \(G\), which is
  equivalent by \SGARefLink{sga3:xviii:prop:1.1:item:iv}{1.1(iv)}.}
  \addtocounter{footnote}{-1}
  \endgroup, and this assertion remains true after every base
  change \(S'\to S\).}
\end{enumerate}

\begin{proof}
First note that \(V\) is relatively schematically dense in
\(G\times_S G\). Indeed, \(V\) is relatively schematically dense in
\[
  m_G^{-1}(U)\cap(U\times_S G)\cap(G\times_S U),
\]
and each of the three factors is obtained from \(U\subseteq G\) by
an evident base change \(G\times_S G\to G\). The assertion therefore
follows from \SGARefLink{sga3:xviii:prop:1.1}{1.1}.

To construct \(f:G\to H\), since \(U\times_S U\to G\) is an fppf
covering, it is enough to construct a morphism \(U\times_S U\to H\)
whose two pullbacks to
\[
  (U\times_S U)\times_G(U\times_S U)
\]
are equal. Take the morphism
\[
  m_H\circ(\bar f\times\bar f):U\times_S U\longrightarrow H.
\]
For sections \(a,b,c,d\in U(S')\), where \(S'\to S\) is arbitrary,
such that \(ab=cd\), we must prove
\[
  \bar f(a)\bar f(b)=\bar f(c)\bar f(d).
\]
If \((a,b)\) and \((c,d)\) belong to \(V(S')\)\begingroup
\renewcommand{\thefootnote}{\ensuremath{\ddagger}}\footnote{The
French source prints \((a,d)\) and \((c,d)\). The first pair must be
\((a,b)\), as required by the displayed equality and by the
compatibility hypothesis on \(V\).}
\addtocounter{footnote}{-1}
\endgroup, the desired equality follows at once. Thus it holds on
\((V\times_G V)(S')\).

I claim that \(V\times_G V\) is a relatively schematically dense
open subset of
\[
  (U\times_S U)\times_G(U\times_S U).
\]
This will complete the proof in cases \textup{\SGARefLink{sga3:xviii:prop:2.3:item:i}{(i)}} and \textup{\SGARefLink{sga3:xviii:prop:2.3:item:ii}{(ii)}}
by \SGARefLink{sga3:xviii:prop:1.4}{1.4}, because the morphism \(f\) constructed in this way extends
\(\bar f\), and its being a homomorphism is then evident, again by
\SGARefLink{sga3:xviii:prop:1.4}{1.4}. Now
\[
  V\times_G V
  =
  \bigl(V\times_G(G\times_S G)\bigr)
  \cap
  \bigl((G\times_S G)\times_G V\bigr).
\]
By symmetry and \SGARefLink{sga3:xviii:prop:1.1}{1.1}, it is enough to check that
\(V\times_G(G\times_S G)\) is relatively schematically dense in
\((G\times_S G)\times_G(G\times_S G)\). The latter prescheme is
isomorphic over \(S\) to \(G\times_S G\times_S G\), the isomorphism
being given by
\[
  (a,b,c,d)\longmapsto(a,b,c).
\]
What remains is therefore the relative schematic density of
\(V\times_S G\) in \(G\times_S G\times_S G\), which follows from
that of \(V\) in \(G\times_S G\).

It remains to treat case \textup{\SGARefLink{sga3:xviii:prop:2.3:item:iii}{(iii)}}. To prove
\(\bar f(a)\bar f(b)=\bar f(c)\bar f(d)\), we may make an fppf
covering base change. Suppose that, after such a base change
\(S'\to S\), there is a section \(x\in G(S')\) such that
\[
  (b,x),\quad(a,bx),\quad(d,x),\quad(c,dx)
\]
all belong to \(V\). Then
\[
\begin{aligned}
  \bar f(a)\bar f(b)\bar f(x)
    &=\bar f(a)\bar f(bx)
     =\bar f(abx)
     =\bar f(cdx)\\
    &=\bar f(c)\bar f(dx)
     =\bar f(c)\bar f(d)\bar f(x),
\end{aligned}
\]
which gives the desired equality.

To find such an \(x\)\footnote{Editor's note: The beginning of the
following sentence has been added.}, set, for every \(z\in U(S')\),
\[
  V_z=
  \operatorname{pr}_{2/S'}
  \bigl((z\times_{S'}U_{S'})\cap V_{S'}\bigr).
\]
The conditions imposed on \(x\) say exactly that \(x\in W\), where
\[
  W=
  V_b(S')\cap V_d(S')\cap
  b^{-1}V_a(S')\cap d^{-1}V_c(S').
\]
By hypothesis \textup{\SGARefLink{sga3:xviii:prop:2.3:item:iii}{(iii)}}, \(W\) is relatively schematically
dense in \(G\). \SGARefLink{sga3:xviii:prop:1.7}{Proposition 1.7} therefore supplies such a section
\(x\) after an fppf covering base change.

Let us verify that the morphism thus constructed is multiplicative.
Let \(a,b\in G(S')\), where \(S'\to S\) is arbitrary; after this
base change, rename \(S'\) as \(S\). Choose an fppf covering base
extension \(S'\to S\) and a section \(x\in G(S')\) such that
\[
  x\in U(S')
  \qquad\text{and}\qquad
  ax^{-1}\in U(S'),
\]
that is, \(x\in(a^{-1}U)^{-1}(S')\). Choose in addition another
fppf covering extension \(S''\to S'\) and a section
\(y\in G(S'')\) such that
\[
  y,\quad by^{-1},\quad aby^{-1}\in U(S''),
\]
and
\[
  by^{-1}\in V_x,
  \qquad
  xby^{-1}\in V_{ax^{-1}}.
\]
Then, by the definition of \(f\), over \(S''\) we have
\[
  f(a)=\bar f(ax^{-1})\bar f(x),\qquad
  f(b)=\bar f(by^{-1})\bar f(y),\qquad
  f(ab)=\bar f(aby^{-1})\bar f(y).
\]
Furthermore,
\[
\begin{aligned}
  f(a)f(b)
    &=\bar f(ax^{-1})\bar f(x)\bar f(by^{-1})\bar f(y)\\
    &=\bar f(ax^{-1})\bar f(xby^{-1})\bar f(y)\\
    &=\bar f(aby^{-1})\bar f(y)
     =f(ab),
\end{aligned}
\]
which proves multiplicativity.

It is now easy to see that \(f\) extends \(\bar f\). Let
\(a\in U(S')\); after base change, rename \(S'\) as \(S\), and choose
an fppf covering \(S'\to S\) and sections \(x,y\in G(S')\) such that
\[
  (x,y),\qquad (ax,y),\qquad (a,xy)
\]
belong to \(V(S')\). Then
\[
  f(xy)=\bar f(x)\bar f(y)=\bar f(xy),
  \qquad
  f(axy)=\bar f(ax)\bar f(y)=\bar f(axy).
\]
Consequently,
\[
\begin{aligned}
  f(a)
    &=f(axy)f\bigl((xy)^{-1}\bigr)
     =f(axy)f(xy)^{-1}\\
    &=\bar f(axy)\bar f(xy)^{-1}
     =\bar f(a).
\end{aligned}
\]
\end{proof}

\SGARefTarget{sga3:xviii:remark:2.3.1}\paragraph{\SGARefLink{sga3:xviii:remark:2.3.1}{Remark 2.3.1}.}
In many cases hypothesis \textup{\SGARefLink{sga3:xviii:prop:2.3:item:iii}{(iii)}} holds, for it will still
hold if \(U\) is replaced by a smaller open subset \(U'\), still
relatively schematically dense in \(G\). For example:

\SGARefTarget{sga3:xviii:prop:2.4}\paragraph{\SGARefLink{sga3:xviii:prop:2.4}{Proposition 2.4}.}
\textit{In the situation of \SGARefLink{sga3:xviii:prop:2.3}{2.3}, suppose that every geometric fiber
of \(G/S\) is irreducible. Put
\[
  U'=\operatorname{pr}_1V
  \qquad\text{and}\qquad
  V'=V\cap m_G^{-1}(U')\cap(U'\times_S U').
\]
Then \(U'\subseteq U\) is an open subset relatively schematically
dense in \(G\), and \(U'\), \(\bar f|_{U'}\), and \(V'\) satisfy
hypothesis \textup{\SGARefLink{sga3:xviii:prop:2.3:item:iii}{(iii)}}.}

\begin{proof}
The subset \(U'\) is open because
\(G\times_S G\to G\) is flat and locally of finite presentation.
All the other verifications are immediate except for hypothesis
\textup{\SGARefLink{sga3:xviii:prop:2.3:item:iii}{(iii)}}.

Let \(a\in U(S')\); after base change, rename \(S'\) as \(S\). To
verify that
\[
  \operatorname{pr}_2\bigl((a\times_S U')\cap V'\bigr)
\]
is relatively schematically dense in \(U'\), it suffices by
\SGARefLink{sga3:xviii:cor:1.3}{Corollary 1.3} to do so fiber by fiber. Thus it is enough to treat the
case where \(S\) is the spectrum of a field, and in that case it is
enough to prove that \(U'\) is nonempty, since \(G\) is irreducible
and has no embedded components (cf.\ VI\textsubscript{A}, \SGARefLink{prop:sga3-VIA-1.1.1}{1.1.1}).
Now
\[
\begin{aligned}
  \operatorname{pr}_2\bigl((a\times_S U')\cap V'\bigr)
  &=
  \operatorname{pr}_2\bigl((a\times_S U')\cap V\bigr)\\
  &\quad{}\cap
  \operatorname{pr}_2\bigl((a\times_S U')\cap m_G^{-1}(U')\bigr).
\end{aligned}
\]
The second factor on the right is dense in \(G\). It therefore
remains only to prove that
\[
  \operatorname{pr}_2\bigl((a\times_S U')\cap V\bigr)
\]
is dense in \(G\), that is, nonempty; and this is clear because
\(a\in U'=\operatorname{pr}_1V\).
\end{proof}

% END INCLUDED SGA3_English_current_buildable_source_20260823/source/inputs/XVIII/tex/components/02_expose_XVIII_section2.tex

% BEGIN INCLUDED SGA3_English_current_buildable_source_20260823/source/inputs/XVIII/tex/components/03_expose_XVIII_section3.tex
\SGARefTarget{sga3:xviii:section:3}\section*{3. Construction of a
Group from a Rational Law}
\addcontentsline{toc}{section}{3. Construction of a Group from a
Rational Law}

\SGARefTarget{sga3:xviii:paragraph:3.0}\paragraph{\SGARefLink{sga3:xviii:paragraph:3.0}{3.0}.}
Suppose given an \(S\)-prescheme \(X/S\) and a rational map
\[
  X\times_S X\dashrightarrow X
\]
over \(S\). We seek an \(S\)-group \(G/S\) and a birational map
relative to \(S\)\footnote{Editor's note: The notion of a
birational map (relative to \(S\)) should be made explicit, perhaps
in an addition to 1.5.2.}
\[
  X\dashrightarrow G
\]
that commutes with the composition laws. We treat only the case in
which \(X/S\) satisfies the following hypothesis:
\[
\tag{\(\diamond\)}\label{sga3:xviii:section:3:equation:diamond}
\begin{gathered}
  X/S\text{ is faithfully flat and of finite presentation,}\\
  \text{with separated fibers and without embedded components.}
\end{gathered}
\]
(The last two conditions hold for a group
prescheme.\footnote{Editor's note: Cf.\ VI\textsubscript{A}; the
reference is to be specified.})

We shall often omit the symbol \(S\) from fiber products.

Let \(X/S\) be a prescheme having the properties
\(\textup{(\(\diamond\))}\) above, and let \(W\) be a subprescheme
of finite presentation of \(X\times X\times X\) having the
following property:
\[
\tag{*}\label{sga3:xviii:section:3:equation:star}
\begin{gathered}
  \text{the three morphisms \(W\to X\times X\) induced by the
  projections}\\
  \text{from \(X^3\) onto \(X^2\) are relatively schematically
  dense open immersions over \(S\).}
\end{gathered}
\]

\paragraph{Notation.}
We shall use the following terminology. Given sections
\(a,b,c\in X(S')\), where \(S'\to S\) is arbitrary, such that
\((a,b,c)\in W(S')\), we write
\[
  c=ab,\qquad b=a^{-1}c,\qquad a=cb^{-1}.
\]

\SGARefTarget{sga3:xviii:def:3.0.1}\paragraph{\SGARefLink{sga3:xviii:def:3.0.1}{Definition 3.0.1}.}
\footnote{Editor's note: The numbering \SGARefLink{sga3:xviii:def:3.0.1}{3.0.1} and \SGARefLink{sga3:xviii:remark:3.0.2}{3.0.2} has been
introduced.}
Given a section \((a,b)\in X^2(S')\), we say that \(ab\) is
\emph{defined} if, and only if, there exists a section
\(c\in X(S')\) such that \((a,b,c)\in W(S')\); equivalently, if and
only if \((a,b)\in\operatorname{pr}_{12}W(S')\). Similarly, to say
that \(a^{-1}b\) or \(ab^{-1}\) is defined has the analogous
meaning, and we extend this terminology also to products of several
factors.

\SGARefTarget{sga3:xviii:remark:3.0.2}\paragraph{\SGARefLink{sga3:xviii:remark:3.0.2}{Remark 3.0.2}.}
Let us note at once the following fact. By \textup{(*)}, \(W\)
defines a rational map \(X^2\dashrightarrow X\) over \(S\), namely
the one given by \((a,b)\mapsto ab\). This rational map may well
have a domain of definition larger than
\(\operatorname{pr}_{12}W\). Nevertheless, we say that \(ab\) is
defined only if \((a,b)\in\operatorname{pr}_{12}W(S')\).

\SGARefTarget{sga3:xviii:def:3.1}\paragraph{\SGARefLink{sga3:xviii:def:3.1}{Definition 3.1}.}
A \emph{group germ} is an \(S\)-prescheme \(X/S\) having the
properties \(\textup{(\(\diamond\))}\) above, together with a
subprescheme \(W\) of finite presentation of
\(X\times X\times X\), having the following properties:
\begin{enumerate}
\item[\SGAInlineRefTarget{sga3:xviii:def:3.1:item:i}(i)] \(W\) satisfies property \textup{(*)} above.

\item[\SGAInlineRefTarget{sga3:xviii:def:3.1:item:ii}(ii)] For every section \(a\in X(S)\), the sets
\[
\begin{array}{ll}
  \operatorname{pr}_i\bigl((a\times X\times X)\cap W\bigr),
    & i=2,3,\\
  \operatorname{pr}_i\bigl((X\times a\times X)\cap W\bigr),
    & i=3,1,\\
  \operatorname{pr}_i\bigl((X\times X\times a)\cap W\bigr),
    & i=1,2,
\end{array}
\]
are schematically dense in \(X\) relative to \(S\), and this
assertion remains true after every base change \(S'\to S\).
(Condition \SGARefLink{sga3:xviii:def:3.1:item:i}{(i)} implies that these sets are open in \(X\).)
Intuitively, this says that \(ax\) is defined for sufficiently
general \(x\), and similarly for the other products.

\item[\SGAInlineRefTarget{sga3:xviii:def:3.1:item:iii}(iii)] The law is associative: if
\((a,b,c)\in X^3(S')\), for an arbitrary \(S'\to S\), and both
\((ab)c\) and \(a(bc)\) are defined, then
\[
  (ab)c=a(bc).
\]
\end{enumerate}

\paragraph{Remarks.}
\textup{(a)} In \SGARefLink{sga3:xviii:def:3.1:item:i}{(i)}, the condition that \(W\) be schematically
dense in \(X^2\) relative to \(S\) may be replaced by the condition
that, for every \(s\in S\), the fiber \(W_s\) be topologically
dense in \(X_s^2\), by the hypotheses on \(X\) and by \SGARefLink{sga3:xviii:prop:1.2}{1.2}.

\textup{(b)} Condition \SGARefLink{sga3:xviii:def:3.1:item:iii}{(iii)} is equivalent to the following.
Let \(V_1\) (respectively \(V_2\)) be the open subset of \(X^3\)
on which the rational map
\[
  (a,b,c)\longmapsto a(bc)
  \qquad
  \bigl(\text{respectively }(a,b,c)\longmapsto (ab)c\bigr)
\]
is defined\begingroup
\renewcommand{\thefootnote}{*}\footnote{In the sense explained in
\SGARefLink{sga3:xviii:def:3.0.1}{3.0.1}; one could also replace these open subsets by the domains of
definition (cf.\ \SGARefLink{sga3:xviii:def:1.5}{1.5}) of the rational maps under consideration.}
\addtocounter{footnote}{-1}
\endgroup. Then there exists an open subset
\(V\subseteq V_1\cap V_2\), schematically dense in \(X^3\)
relative to \(S\), on which the two maps agree. This follows from
\SGARefLink{sga3:xviii:prop:1.4}{1.4}, since \(V_1\) and \(V_2\) are schematically dense in \(X^3\)
relative to \(S\).

\textup{(c)} Hypothesis \SGARefLink{sga3:xviii:prop:1.4:item:ii}{(ii)} will be used below to ensure that
\(X\) is a subobject of the group prescheme that it defines. In
many cases one can deduce \SGARefLink{sga3:xviii:def:3.1:item:ii}{(ii)} from \SGARefLink{sga3:xviii:def:3.1:item:i}{(i)}, provided that \(X\) is
replaced by a relatively schematically dense open subset \(X'\)
and \(W\) by a relatively schematically dense open subset of
\(W\cap {X'}^3\). More precisely:

\SGARefTarget{sga3:xviii:prop:3.2}\paragraph{\SGARefLink{sga3:xviii:prop:3.2}{Proposition 3.2}.}
\textit{Suppose that every geometric fiber of \(X/S\) is
irreducible, and let \(W\) be a subprescheme of \(X^3\) satisfying
condition \SGARefLink{sga3:xviii:def:3.1:item:i}{(i)} of \SGARefLink{sga3:xviii:def:3.1}{3.1}. Then there exist a relatively schematically
dense open subset \(X'\subseteq X\) and a relatively schematically
dense open subset
\(W'\subseteq W\cap {X'}^3\)\begingroup
\renewcommand{\thefootnote}{\ensuremath{\dagger}}\footnote{The
French source prints \(W\cap(X'\times X')\) in this statement. The
proof immediately defines \(W'=W\cap {X'}^3\), which is the
dimensionally coherent reading used here.}
\addtocounter{footnote}{-1}
\endgroup such that the pair \((X',W')\) satisfies conditions
\SGARefLink{sga3:xviii:def:3.1:item:i}{(i)} and \SGARefLink{sga3:xviii:def:3.1:item:ii}{(ii)}. If \SGARefLink{sga3:xviii:def:3.1:item:iii}{(iii)} holds for \((X,W)\), then it holds for
\((X',W')\).}

\begin{proof}
Set
\[
  X'=\bigcap_{i=1}^{3}\operatorname{pr}_i W.
\]
Each \(\operatorname{pr}_iW\) is open in \(X\), because
\(W\to X^2\) is an open immersion and the projections
\(X^2\to X\) are flat and of finite presentation. Moreover,
\(\operatorname{pr}_iW\) is schematically dense in \(X\) relative
to \(S\), because \(W\) is so in \(X^2\) and the projections
\(X^2\to X\) are surjective; it is enough here to check
topological density. Put
\[
  W'=W\cap {X'}^3.
\]

To verify \SGARefLink{sga3:xviii:def:3.1:item:i}{(i)}, note that
\[
\begin{aligned}
  W'={}&
  \bigl(W\cap(X'\times X'\times X)\bigr)
  \cap\bigl(W\cap(X\times X'\times X')\bigr)\\
  &\cap\bigl(W\cap(X'\times X\times X')\bigr).
\end{aligned}
\]
Now
\[
  W\cap(X'\times X'\times X)
  \simeq \operatorname{pr}_{12}W\cap(X'\times X'),
\]
so this term is schematically dense in \(W\) relative to \(S\).
The other two terms on the right are likewise schematically dense
in \(W\), and hence \(W'\) is schematically dense in \(W\).
Consequently \(\operatorname{pr}_{ij}W'\) is schematically dense
in \(\operatorname{pr}_{ij}W\), hence in \(X\times X\), and
therefore in \(X'\times X'\).

To verify \SGARefLink{sga3:xviii:def:3.1:item:ii}{(ii)}, let \(a\in X'(S')\). After this base change, rename
\(S'\) as \(S\). We must show, for example, that
\[
  \operatorname{pr}_2\bigl((a\times X'\times X')\cap W'\bigr)
\]
is schematically dense in \(X'\) relative to \(S\). By \SGARefLink{sga3:xviii:cor:1.3}{1.3} it is
enough to check this fiber by fiber. Thus it is enough to treat the
case in which \(S\) is the spectrum of a field, and in that case it
is enough to show that the open subset is nonempty, since the
fibers of \(X/S\) are irreducible and have no embedded components.
Since
\[
  \operatorname{pr}_2\bigl((a\times X\times X)\cap W\bigr)
\]
is nonempty, because \(a\) is a section of \(X'\), this open subset
is dense in \(X\). We have
\[
  \operatorname{pr}_2\bigl((a\times X\times X)\cap W\bigr)
  =
  \operatorname{pr}_2\bigl((a\times X)
    \cap\operatorname{pr}_{12}W\bigr).
\]
Therefore
\[
\begin{aligned}
  &\operatorname{pr}_2\bigl((a\times X)
    \cap\operatorname{pr}_{12}W\bigr)\cap X'\\
  &\qquad =
  \operatorname{pr}_2\bigl((a\times X')
    \cap\operatorname{pr}_{12}W\bigr)
  =
  \operatorname{pr}_2\bigl((a\times X'\times X)\cap W\bigr)
\end{aligned}
\]
is dense in \(X\), and hence in \(X'\).

Similarly,
\(\operatorname{pr}_3((a\times X\times X')\cap W)\) is dense in
\(X\), hence dense in
\(\operatorname{pr}_3((a\times X\times X)\cap W)\). Equivalently,
\((a\times X\times X')\cap W\) is dense in
\((a\times X\times X)\cap W\), and therefore
\[
  \operatorname{pr}_2\bigl((a\times X\times X')\cap W\bigr)
\]
is dense in
\(\operatorname{pr}_2((a\times X\times X)\cap W)\), hence in
\(X\), and thus in \(X'\).

Finally,
\[
\begin{aligned}
  \operatorname{pr}_2\bigl((a\times X'\times X')\cap W'\bigr)
  &=
  \operatorname{pr}_2\bigl((a\times X'\times X')\cap W\bigr)\\
  &=
  \operatorname{pr}_2\bigl((a\times X'\times X)\cap W\bigr)
  \cap
  \operatorname{pr}_2\bigl((a\times X\times X')\cap W\bigr),
\end{aligned}
\]
which is indeed dense in \(X'\). This is what was required. The
other assertions in \SGARefLink{sga3:xviii:def:3.1:item:ii}{(ii)} follow by symmetry, and preservation of
condition \SGARefLink{sga3:xviii:def:3.1:item:iii}{(iii)} is immediate.\footnote{Editor's note: The
following sentence has been added.} This proves \SGARefLink{sga3:xviii:prop:3.2}{Proposition 3.2}.
\end{proof}

\SGARefTarget{sga3:xviii:paragraph:3.2.1}\paragraph{\SGARefLink{sga3:xviii:paragraph:3.2.1}{3.2.1}.}
\footnote{Editor's note: The numbering \SGARefLink{sga3:xviii:paragraph:3.2.1}{3.2.1}, as well as \SGARefLink{sga3:xviii:paragraph:3.2.2}{3.2.2},
has been introduced.}
Fix a group germ \((X,W)\) over \(S\). We first make some
preliminary observations about the situation, collected below as
rules. These rules will often be used without explicit mention in
what follows.

Let \(a\in X(S')\), and after base change rename \(S'\) as \(S\).
There is then a birational map \(\varphi\) over \(S\) from \(X\)
to itself, sending a section \(x\) to \(ax\) whenever the latter is
defined. By \SGARefLink{sga3:xviii:def:3.1:item:ii}{3.1(ii)}, the domain of definition of \(\varphi\)
contains the relatively schematically dense open subset
\[
  \operatorname{pr}_2\bigl((a\times X\times X)\cap W\bigr),
\]
and \(\varphi\) defines an isomorphism from this open subset onto
the open subset
\[
  \operatorname{pr}_3\bigl((a\times X\times X)\cap W\bigr),
\]
on which \(a^{-1}x\) is defined. This observation is generalized
in the evident way in Rule 1 below.

\SGARefTarget{sga3:xviii:rule:1}\paragraph{Rule 1.}
Let \(P=P(x,t_1,\ldots,t_n)\) be a ``product'' of the symbols
\(x,t_1,\ldots,t_n\), obtained recursively as follows:
\(P_0=x\), and \(P_{i+1}\) is one of the expressions
\[
  P_i t,\qquad tP_i,\qquad P_i^{-1}t,\qquad
  tP_i^{-1},\qquad P_i t^{-1},\qquad t^{-1}P_i,
\]
where \(t\) is one of the \(t_j\); and \(P_r=P\). Let
\(a_1,\ldots,a_n\in X(S')\). Then there exists a relatively
schematically dense open subset \(U\) of \(X_{S'}\) such that the
product \(P(x,a_1,\ldots,a_n)\) is defined (in the sense of
\SGARefLink{sga3:xviii:remark:3.0.2}{Remark 3.0.2}) for a section \(x\in X(S'')\) if and only if
\(x\in U(S'')\), and the map
\[
  x\longmapsto P(x,(a))
\]
gives an isomorphism from \(U\) onto another relatively
schematically dense open subset of \(X_{S'}\), denoted
\(P(U,(a))\).

\SGARefTarget{sga3:xviii:rule:2}\paragraph{Rule 2.}
Let \(a,b\in X(S')\). Then:
\begin{align*}
  &\text{if \(ab\) is defined, so is \(a^{-1}(ab)\), and }
    a^{-1}(ab)=b;\\
  &\text{if \(a^{-1}b\) is defined, so is \(a(a^{-1}b)\), and }
    a(a^{-1}b)=b;\\
  &\text{if \(ba^{-1}\) is defined, so is \((ba^{-1})a\), and }
    (ba^{-1})a=b.
\end{align*}

\SGARefTarget{sga3:xviii:rule:3}\paragraph{Rule 3.}
Let \(a,b,b'\in X(S')\). If
\[
  ab=ab',\qquad ba=b'a,\qquad a^{-1}b=a^{-1}b',
  \qquad\text{or}\qquad ba^{-1}=b'a^{-1},
\]
then \(b=b'\). It is understood here that an equality entails that
both of its sides are defined.

\SGARefTarget{sga3:xviii:rule:4}\paragraph{Rule 4.}
Let \(a,b,c\in X(S')\). Whenever both sides are defined, one has
\[
  a\bigl((ba)^{-1}c\bigr)=b^{-1}c.
\]
Similarly,
\[
  \bigl(c(ab)^{-1}\bigr)a=cb^{-1},\qquad
  a^{-1}\bigl((ab^{-1})c\bigr)=b^{-1}c,\qquad
  \bigl(c(b^{-1}a)\bigr)a^{-1}=cb^{-1}.
\]

\SGARefTarget{sga3:xviii:rule:5}\paragraph{Rule 5.}
All the following associativity laws hold whenever both sides are
defined:
\[
\begin{aligned}
  (a^{-1}b)c&=a^{-1}(bc),&
  (ab^{-1})c&=a(b^{-1}c),&
  (ab)c^{-1}&=a(bc^{-1}),\\
  (ab)^{-1}c&=b^{-1}(a^{-1}c),&
  (a^{-1}b)c^{-1}&=a^{-1}(bc^{-1}),&
  (ab^{-1})c^{-1}&=a(cb)^{-1}.
\end{aligned}
\]

\SGARefTarget{sga3:xviii:paragraph:3.2.2}\paragraph{\SGARefLink{sga3:xviii:paragraph:3.2.2}{3.2.2}.
Verification of the rules.}
\textup{\SGARefLink{sga3:xviii:section:1}{(1)}} This follows by an evident induction on the length
\(r\) of \(P\), the case \(r=1\) being a direct consequence of
\SGARefLink{sga3:xviii:def:3.1:item:ii}{3.1(ii)}.

\textup{\SGARefLink{sga3:xviii:section:2}{(2)}} This is immediate from the definition.

\textup{\SGARefLink{sga3:xviii:section:3}{(3)}} Indeed, by Rule 2, in the first case for example,
\[
  b=a^{-1}(ab)=a^{-1}(ab')=b'.
\]

\textup{(4)} Let us verify the first relation, for example. On the
right-hand side, left multiplication by \(b\) is defined and gives
\(c\), by Rule 2. Suppose that it is also defined on the
left-hand side. Then
\[
  b\bigl(a((ba)^{-1}c)\bigr)
  =(ba)\bigl((ba)^{-1}c\bigr)
  =c.
\]
Indeed, \(ba\) is defined by hypothesis, since it occurs in the
expression. The middle term is therefore defined and equal to
\(c\) by Rule 2, and it is equal to the left-hand term by
associativity (cf.\ \SGARefLink{sga3:xviii:def:3.1:item:iii}{3.1(iii)}). Rule 3 thus implies that the desired
equality holds whenever this multiplication by \(b\) is defined.

Now fix \(a\) and \(b\). Rule 1 implies that
\(b(a((ba)^{-1}c))\) is defined for \(c\) ``in'' a relatively
schematically dense open subset \(U\) of \(X\). Hence, on this
relatively schematically dense open subset, the two rational maps
\[
  c\longmapsto b^{-1}c
  \qquad\text{and}\qquad
  c\longmapsto a((ba)^{-1}c)
\]
are equal. By \SGARefLink{sga3:xviii:prop:1.4}{1.4}, they are equal on every common domain of
definition, which proves the desired result.

\textup{(5)} The argument is of the same kind as the preceding
one. For example, let us verify
\[
  (ab^{-1})c^{-1}=a(cb)^{-1}.
\]
If right multiplication by \(c\) is defined on the right-hand
side, equality follows from Rule 4. Since it is enough to verify
such a formula on a relatively schematically dense open subset, we
are reduced to the case in which this multiplication is defined.

\SGARefTarget{sga3:xviii:paragraph:3.2.3}\paragraph{\SGARefLink{sga3:xviii:paragraph:3.2.3}{3.2.3}.}
Consider the relation \(R\) on \(X\times X\) defined as follows.
For \(a,b,a',b'\in X(S')\), one has
\[
  (a,b)\sim(a',b')\pmod{R(S')}
\]
if and only if there exist an fppf covering \(S''\to S'\) and a
section \(x\in X(S'')\) such that \((xa)b\) and \((xa')b'\) are
defined and equal. Then \(R\) is an equivalence relation.

Indeed, this relation is evidently symmetric. By Rule 1, the
product \((xa)b\) is defined when \(x\) is ``in'' a suitable
relatively schematically dense open subset. Thus \SGARefLink{sga3:xviii:prop:1.7}{1.7} asserts that
there exist an fppf covering \(S''\to S'\) and an
\(x\in X(S'')\) for which \((xa)b\) is defined. The relation is
therefore reflexive. Its transitivity follows from the next lemma.

\SGARefTarget{sga3:xviii:lemma:3.3}\paragraph{\SGARefLink{sga3:xviii:lemma:3.3}{Lemma 3.3}.}
\textit{Let \(x,y,a,b,a',b'\) be sections of \(X(S')\) such that
\((xa)b\), \((xa')b'\), \((ya)b\), and \((ya')b'\) are defined. If
\[
  (xa)b=(xa')b',
\]
then
\[
  (ya)b=(ya')b'.
\]}

Indeed, the lemma says that one may test
\((a,b)\sim(a',b')\) using any \(x\in X(S')\), with
\(S'\to S\) an fppf covering, for which both products are defined.
Given
\[
  a,b,a',b',a'',b''\in X(S),
\]
Rules 1, \SGARefLink{sga3:xviii:prop:1.1}{1.1}, and \SGARefLink{sga3:xviii:prop:1.7}{1.7} allow us to find an fppf covering
\(S'\to S\) and a section \(x\in X(S')\) such that the three
products in question are defined; transitivity follows.

\begin{proof}[Proof of \SGARefLink{sga3:xviii:lemma:3.3}{Lemma 3.3}]
Let us begin by writing formally
\[
\begin{aligned}
  (za)b
  &=\bigl(((zx^{-1})x)a\bigr)b
   =\bigl((zx^{-1})(xa)\bigr)b
   =(zx^{-1})\bigl((xa)b\bigr),\\
  (za')b'
  &=\bigl(((zx^{-1})x)a'\bigr)b'
   =\bigl((zx^{-1})(xa')\bigr)b'
   =(zx^{-1})\bigl((xa')b'\bigr).
\end{aligned}
\]
The appropriate rules show that these equalities do hold whenever
their terms are defined.\footnote{Editor's note: To be made
precise\ldots} It follows that
\[
  (za)b=(za')b'
\]
whenever all these expressions are defined. Moreover, by Rule 1
and the hypotheses already made, these expressions are defined
when \(z\in X(S'')\) belongs to \(V(S'')\), where \(V\) is a
certain relatively schematically dense open subset of \(X\) (we
have taken \(S'=S\)). Therefore the two rational self-maps of
\(X\) given by
\[
  z\longmapsto(za)b
  \qquad\text{and}\qquad
  z\longmapsto(za')b'
\]
are equal, and consequently
\[
  (ya)b=(ya')b'.
\]
\end{proof}

\SGARefTarget{sga3:xviii:lemma:3.4}\paragraph{\SGARefLink{sga3:xviii:lemma:3.4}{Lemma 3.4}.}
\textit{Consider the rational map
\[
  \varphi\colon X^3\dashrightarrow X
\]
over \(S\), defined by
\[
  (a,b,c)\longmapsto c^{-1}(ab).
\]
Let \(U\) be the domain of definition of \(\varphi\), and let
\(\Gamma\) be the graph of the morphism \(f\colon U\to X\)
induced by \(\varphi\), regarded as a subprescheme of \(X^4\).
Then a section \((a,b,c,d)\in X^4(S)\) belongs to \(\Gamma(S)\)
if and only if
\[
  (a,b)\sim(c,d).
\]}

\begin{proof}
First observe that the rational map \(\varphi\) is also given by
the formula
\[
  (a,b,c)\longmapsto(xc)^{-1}\bigl((xa)b\bigr)
\]
for an arbitrary section \(x\in X(S)\). Equivalently,
\[
  c^{-1}(ab)=(xc)^{-1}\bigl((xa)b\bigr)
\]
whenever both sides are defined. We leave the verification to the
reader.

We now show that \(\varphi\) is defined at a section
\((a,b,c)\in X^3(S)\) if and only if there exists a \(d\) such
that \((a,b)\sim(c,d)\), and that in this case
\[
  d=\varphi(a,b,c).
\]
Suppose first that \((a,b)\sim(c,d)\). To verify that
\(\varphi\) is defined, we may make an fppf-covering base change
(\SGARefLink{sga3:xviii:prop:1.6}{Proposition 1.6}), and may therefore suppose that there exists a
section \(x\in X(S')\) such that \((xa)b\) and \((xc)d\) are
defined and equal. It follows from Rule 2 that
\[
  (xc)^{-1}\bigl((xa)b\bigr)
\]
is defined and equal to \(d\).

Conversely, suppose that \(\varphi\) is defined at
\((a,b,c)\), and put \(d=\varphi(a,b,c)\). Choose an
fppf covering \(S'\to S\) and a section \(x\in X(S')\) such that
\((xa)b\) and \((xc)d\) are defined. We wish to show that they
are equal. It is enough to show that the two rational self-maps of
\(X\) over \(S\), given by
\[
  b\longmapsto(xa)b
  \qquad\text{and}\qquad
  b\longmapsto(xc)\varphi(a,b,c),
\]
are equal; this follows from the observation in the first
paragraph.
\end{proof}

\SGARefTarget{sga3:xviii:prop:3.5}\paragraph{\SGARefLink{sga3:xviii:prop:3.5}{Proposition 3.5}.}
\textit{The equivalence relation
\[
  R\rightrightarrows X^2
\]
is representable, and it is a flat relation of finite
presentation; that is, the projections of \(R\) onto \(X^2\) are
flat morphisms of finite presentation.}

\begin{proof}
We may suppose that \(S\) is affine. Since \((X,W)\) is of finite
presentation over \(S\), we may descend the entire situation to an
\(S\) of finite type over \(\operatorname{Spec}\mathbb Z\), hence
to a noetherian \(S\). We may therefore suppose that \(S\) is
noetherian. It is then immediate that the graph \(\Gamma\) of
\SGARefLink{sga3:xviii:lemma:3.4}{3.4} is of finite presentation over \(X^2\). The projection
\(\left.\operatorname{pr}_{12}\right|_\Gamma\) is flat because
\[
  \Gamma\xrightarrow{\sim}\operatorname{pr}_{123}\Gamma,
\]
the latter being an open subset of \(X^3\), and because the
projection
\[
  \operatorname{pr}_{12}\colon X^3\longrightarrow X^2
\]
is flat, since \(X\) is flat over \(S\).

I claim that \(\Gamma\) represents \(R\). There is something to
prove here, because the domain of definition of a rational map
does not in general commute with base change. Let \(S''\to S\).
What is clear from \SGARefLink{sga3:xviii:lemma:3.4}{3.4} applied over \(S''\) is that
\[
  R(S'')\supseteq\Gamma(S''),
\]
because \(\varphi\times_S S''\) is certainly defined on
\(U\times_S S''\). Now let
\[
  (a,b,c,d)\in R(S'').
\]
We must prove that \((a,b,c,d)\in\Gamma(S'')\). This can be
checked locally, say for the étale topology. We may therefore
suppose that \(S''\) is strictly local, that is, the spectrum of a
henselian ring with separably closed residue field. Moreover, by
applying \SGARefLink{sga3:xviii:prop:1.6}{1.6} and the usual passage-to-the-limit arguments, we are
reduced to the case in which \(S\) is strictly local and
\(S''\to S\) is local.

Suppose that there is a section \(x\in X(S)\) such that over
\(S''\) the products \((xa)b\) and \((xc)d\) are defined. This
implies that
\[
  (xc)^{-1}\bigl((xa)b\bigr)
\]
is defined and equal to \(d\). Now there is a relatively
schematically dense open subset \(V\) of \(X^3\) such that
\((xc)^{-1}((xa)b)\) is defined if and only if
\((a,b,c)\in V(S'')\), and \(V\subseteq U\). Consequently,
\((a,b,c,d)\in\Gamma(S'')\) if such an \(x\) exists. By \SGARefLink{sga3:xviii:prop:1.6}{1.6}, we
may make an fppf-covering base change \(S'\to S\) in order to find
such an \(x\). Since \(S\) is strictly local, we can
find\begingroup
\renewcommand{\thefootnote}{*}\footnote{Combining Exposé
VI\textsubscript{A}, 1.1.1 and EGA IV\textsubscript{4}, 17.16.2.}
\addtocounter{footnote}{-1}
\endgroup{} a faithfully flat, local, finite morphism
\(S'\to S\) and a section \(x\in X(S')\) that passes through an
arbitrary closed point of the closed fiber of \(X/S\). To say
that \((xa)b\) and \((xc)d\) are defined means that, over
\(S''\), \(x\) lies in a certain relatively schematically dense
open subset; this can be checked on the closed fiber of
\(X_{S''}\). Thus the required section exists.
\end{proof}

Let \(G\) now be the quotient of \(X^2\) by \(R\), as a sheaf for
the fppf topology. We define a composition law on \(G\) as
follows. Let
\[
  (g,g')\in G(S')
\]
be represented by a section
\[
  ((a,b),(c,d))\in (X^2\times_S X^2)(S''),
\]
where \(S''\to S'\) is an fppf covering. Suppose moreover that
there is a section \(x\in X(S'')\) such that \(a(b(cx))\) and
\(x^{-1}d\) are defined; this is permitted by Rule 1 and \SGARefLink{sga3:xviii:prop:1.7}{1.7}. We
define \(gg'\) to be the class in \(G(S')\) represented by the
section
\[
  \bigl(a(b(cx)),x^{-1}d\bigr)\in X^2(S'').
\]

Let us verify that \(gg'\) is independent of the choice of the
section \(x\) and of the representative
\(((a,b),(c,d))\). Let
\[
  (a',b')\sim(a,b),\qquad (c',d')\sim(c,d),
\]
and let \(x'\) be such that \(a'(b'(c'x'))\) and
\({x'}^{-1}d'\) are defined. We may suppose that all these are
sections over \(S''\). We must prove that
\[
  \bigl(a(b(cx)),x^{-1}d\bigr)
  \sim
  \bigl(a'(b'(c'x')),{x'}^{-1}d'\bigr);
\]
that is, for a suitable fppf covering \(S'''\to S''\) and a
suitable section \(z\in X(S''')\), we must have
\[
 \bigl(z(a(b(cx)))\bigr)(x^{-1}d)
 =
 \bigl(z(a'(b'(c'x')))\bigr)({x'}^{-1}d').
\]
Whenever all the products involved are defined,
\[
\begin{aligned}
 \bigl(z(a(b(cx)))\bigr)(x^{-1}d)
 &=\bigl((za)(b(cx))\bigr)(x^{-1}d)\\
 &=\bigl(((za)b)(cx)\bigr)(x^{-1}d)\\
 &=\bigl((((za)b)c)x\bigr)(x^{-1}d)\\
 &=\bigl(((za)b)c\bigr)\bigl(x(x^{-1}d)\bigr)\\
 &=\bigl(((za)b)c\bigr)d,
\end{aligned}
\]
and the same identities hold with primes. By Rule 1 and \SGARefLink{sga3:xviii:prop:1.7}{1.7},
such a \(z\) exists. It therefore remains to prove that
\[
  \bigl(((za)b)c\bigr)d
  =
  \bigl(((za')b')c'\bigr)d'.
\]
But \((za)b=(za')b'\), since
\((a,b)\sim(a',b')\) by \SGARefLink{sga3:xviii:lemma:3.3}{3.3}; the desired equality now follows
from \((c,d)\sim(c',d')\).

Consider the natural morphism
\[
  i\colon X\longrightarrow G
\]
defined as follows. For \(a\in X(S')\), choose an fppf covering
\(S''\to S'\) and a section \(b\in X(S'')\) such that
\(ab^{-1}\) is defined, and put
\[
  i(a)=\text{the class in \(G\) of }\bigl(ab^{-1},b\bigr).
\]
It is readily verified that this class, which a priori lies in
\(G(S'')\), is independent of the choice of \(b\), and therefore
defines a well-determined element of \(G(S')\).

The reader may enjoy verifying the following proposition.

\SGARefTarget{sga3:xviii:prop:3.6}\paragraph{\SGARefLink{sga3:xviii:prop:3.6}{Proposition 3.6}.}
\textit{The morphism \(i\) commutes with the composition laws on
\(X\) and \(G\). In other words, if \(a,b\in X(S')\) are
sections for which \(ab\) is defined, then
\[
  i(a)i(b)=i(ab).
\]}

The goal of this section is the following theorem.

\SGARefTarget{sga3:xviii:thm:3.7}\paragraph{\SGARefLink{sga3:xviii:thm:3.7}{Theorem 3.7}.}
\footnote{Editor's note: If \(X\) is smooth and separated over
\(S\), faithfully flat and of finite presentation over \(S\),
then \(G/S\) is representable by a smooth \(S\)-group scheme of
finite type over \(S\). This is Theorem 6.6.1 of
\emph{Néron Models}, by Bosch--Lütkebohmert--Raynaud, Springer
(1990).}
\textit{Let \((X,W)\) be a group germ over \(S\), with \(X/S\)
faithfully flat and of finite presentation, and with separated
fibers having no embedded component. With the notation above:}
\begin{enumerate}
\item[\SGAInlineRefTarget{sga3:xviii:thm:3.7:item:i}(i)] \textit{\(G\) is a sheaf of groups.}
\item[(ii)] \textit{\(i\colon X\to G\) is representable by an
open immersion.}
\item[\SGAInlineRefTarget{sga3:xviii:theorem:3-7:item:iii}\SGAInlineRefTarget{sga3:xviii:thm:3.7:item:iii}(iii)] \textit{\(G\) is representable fppf-locally over
\(S\).}
\item[\SGAInlineRefTarget{sga3:xviii:thm:3.7:item:iv}(iv)] \textit{If \(G/S\) is representable, then it is a
flat group of finite presentation, and
\(i\colon X\to G\) is schematically dense relative to \(S\).}
\end{enumerate}

Observe that \(G\) is plainly characterized by properties
\textup{\SGARefLink{sga3:xviii:thm:3.7:item:i}{(i)}}--\textup{\SGARefLink{sga3:xviii:thm:3.7:item:iv}{(iv)}}; one may therefore forget its
explicit construction.

The proof proceeds in several stages.

\SGARefTarget{sga3:xviii:lemma:3.8}\paragraph{\SGARefLink{sga3:xviii:lemma:3.8}{Lemma 3.8}.}
\begin{enumerate}
\item[\SGAInlineRefTarget{sga3:xviii:lemma:3.8:item:i}(i)] \textit{Let \(c\) be a section of \(X(S)\). Then the
morphism from the prescheme \(c\times X\), which is
\(S\)-isomorphic to \(X\), into \(G\), given by
\[
  c\times X\lhook\joinrel\longrightarrow X\times X
  \longrightarrow G,
\]
is a monomorphism.}

\item[\SGAInlineRefTarget{sga3:xviii:lemma:3.8:item:ii}(ii)] \textit{Let \(\{c_i\}_{i\in I}\) be sections of
\(X(S)\), put
\[
  Z=\coprod_i c_i\times X,
\]
and let \(R'\rightrightarrows Z\) be the equivalence relation
induced on \(Z\) by the evident morphism \(Z\to X^2\). Then
\(R'\) is a ``gluing'' of the \(c_i\times X\simeq X\).}
\end{enumerate}

\begin{proof}
\textup{\SGARefLink{sga3:xviii:lemma:3.8:item:i}{(i)}} For two sections \((c,a)\) and \((c,a')\) of
\((c\times X)(S')\) to have the same image in \(G\), one must
have \((c,a)\sim(c,a')\); that is, for a suitable section
\(x\in X(S'')\), where \(S''\to S'\) is an fppf covering,
\[
  (xc)a=(xc)a'
\]
\begingroup
\renewcommand{\thefootnote}{\ensuremath{\dagger}}\footnote{The
French source prints \(\sim\) between these two sections of
\(X\). Equality is the relation-defining condition and is
required for the following application of Rule 3.}
\addtocounter{footnote}{-1}
\endgroup.
It follows from Rule 3 that \(a=a'\).

\textup{\SGARefLink{sga3:xviii:lemma:3.8:item:ii}{(ii)}} Let \(c_i,c_j\) be two sections, and consider the
birational self-map \(\psi_{ji}\) of \(X\) over \(S\), given by
\[
  x\longmapsto c_j^{-1}(c_i x).
\]
\begingroup
\renewcommand{\thefootnote}{\ensuremath{\ddagger}}\footnote{The
French source introduces this map as \(\psi_{ji}\), but calls it
\(\varphi_{ji}\) thereafter. We use \(\psi_{ji}\) consistently.}
\addtocounter{footnote}{-1}
\endgroup
For \(y\in X(S)\), this is the same map as
\[
  x\longmapsto (yc_j)^{-1}\bigl((yc_i)x\bigr),
\]
as is readily seen. Moreover, \(\psi_{ji}\) is defined at a
section \(b\in X(S')\) if and only if there exists a section
\(b'\) such that
\[
  (c_i,b)\sim(c_j,b'),
\]
and then \(b'=\psi_{ji}(b)\). Let \(U_{ji}\) be the domain of
definition of \(\psi_{ji}\) over \(S\). It remains to prove that
this domain of definition is universal: if \(b\in X(S'')\), for
an arbitrary \(S''\to S\), and \(\psi_{ji}\) is defined at
\(b\), then \(b\in U_{ji}(S'')\). Equivalently, one must prove
that if \(b,b'\in X(S'')\) satisfy
\[
  (c_i,b)\sim(c_j,b'),
\]
then \(b\in U_{ji}(S'')\). We leave to the reader the
verification of this fact, which is analogous to that in \SGARefLink{sga3:xviii:prop:3.5}{3.5}.
\end{proof}

\SGARefTarget{sga3:xviii:lemma:3.9}\paragraph{\SGARefLink{sga3:xviii:lemma:3.9}{Lemma 3.9}.}
\textit{Suppose that \(\{c_i\}_{i\in I}\) are sections of
\(X(S)\) such that
\[
  \coprod_i c_i\times X\longrightarrow G
\]
is surjective as a morphism of sheaves. Then \(G\) is
representable, flat, and of finite presentation over \(S\), and
the structural morphism
\[
  X^2\longrightarrow G
\]
is flat and of finite presentation.}

\begin{proof}
The representability of \(G\) is an immediate consequence of
\SGARefLink{sga3:xviii:lemma:3.8}{Lemma 3.8}, and it follows that
\[
  \coprod_i c_i\times X\longrightarrow G
\]
is an open covering. To prove that \(X^2\to G\) is flat, it is
enough to do so locally, hence to prove that the rational map
\[
  X^2\dashrightarrow c_i\times X
\]
induces a flat morphism on its domain of definition. This rational
map is given by
\[
  (a,b)\longmapsto
  \bigl(c_i,\bigl((xc_i)^{-1}(xa)\bigr)b\bigr),
\]
where \(x\) is an arbitrary section. If it is defined at
\((a,b)\in X^2(S')\), one can find an fppf covering
\(S''\to S'\) and a section \(x\in X(S'')\) such that
\(\bigl((xc_i)^{-1}(xa)\bigr)b\) is defined. It is then readily
seen that this morphism is flat. The same argument shows that it
is locally of finite presentation, and hence is fppf-covering.

By construction the relation \(R\) is effective. \SGARefLink{sga3:xviii:prop:3.5}{Proposition 3.5}
therefore shows that \(X^2\to G\) becomes a morphism of finite
presentation after the fppf-covering base change
\(G\leftarrow X^2\). Thus \(X^2\to G\) is of finite presentation.

Let us prove that \(G\to S\) is flat and of finite presentation.
It is flat and locally of finite presentation because \(G\) is
covered by the \(c_i\times X\simeq X\). Moreover, \(X^2/S\) is
quasi-compact and \(X^2\to G\) is surjective, so \(G/S\) is
quasi-compact. To prove that \(G\to S\) is quasi-separated, note
the following cartesian diagram:
\[
% BEGIN INCLUDED SGA3_English_current_buildable_source_20260823/source/inputs/XVIII/native_diagrams/exp18/002_cartesian_relation_square.tex
\begin{tikzcd}[column sep=large,row sep=large]
  R \arrow[r,"\alpha"] \arrow[d,"\gamma"'] &
  X^2\times_S X^2 \arrow[d,"\beta"]\\
  G \arrow[r,"\Delta"] & G\times_S G.
\end{tikzcd}

% END INCLUDED SGA3_English_current_buildable_source_20260823/source/inputs/XVIII/native_diagrams/exp18/002_cartesian_relation_square.tex

\]
Here \(\gamma\) is surjective and \(\alpha,\beta\) are
quasi-compact. Hence \(\beta\alpha\) is quasi-compact, and
therefore so is \(\Delta\).\footnote{Editor's note: Another way
to conclude here is by faithfully flat descent
\textup{(EGA IV\textsubscript{2}, 2.7.1)}, since \(\beta\) is an
fppf covering.}
\end{proof}

\SGARefTarget{sga3:xviii:lemma:3.10}\paragraph{\SGARefLink{sga3:xviii:lemma:3.10}{Lemma 3.10}.}
\textit{Let \(\{c_i\}_{i\in I}\) be sections of \(X(S)\). For
\[
  \coprod_i c_i\times X\longrightarrow G
\]
to be surjective as a morphism of fppf sheaves, it is enough that
the following condition hold: for every \(S'\to S\) and every
\((a,b)\in X^2(S')\), there exist an open covering
\(\{S'_\nu\}_{\nu\in N}\) of \(S'\) and a function
\[
  N\longrightarrow I,\qquad \nu\longmapsto i(\nu),
\]
such that \(\bigl(c_{i(\nu)}^{-1}a\bigr)b\) is defined over
\(S'_\nu\).}

\begin{proof}
Let \(S''\to S\) be arbitrary, and let \(g\in G(S'')\). Choose
an fppf covering \(S'\to S''\) and a section
\((a,b)\in X^2(S')\) representing \(g\). Take the open covering
\(\{S'_\nu\}\) supplied by the hypothesis. Over each \(S'_\nu\)
one has
\[
  (a,b)\sim
  \bigl(c_{i(\nu)},\bigl(c_{i(\nu)}^{-1}a\bigr)b\bigr).
\]
Thus, locally over \(S'_\nu\), the element \(g\) is represented
by a section of \(c_{i(\nu)}\times X\). Since the family
\(\{S'_\nu\to S''\}\) is fppf-covering, the desired surjectivity
follows.
\end{proof}

\SGARefTarget{sga3:xviii:lemma:3.11}\paragraph{\SGARefLink{sga3:xviii:lemma:3.11}{Lemma 3.11}.}
\textit{Let \(Y/S\) be a prescheme of finite presentation, and
let \(\{a_i\}_{i\in I}\) be sections of \(Y(S)\). Let
\(s_0,s_1\) be points of \(S\), with \(s_0\) a specialization of
\(s_1\), and let \(Y_j\) be the fiber of \(Y/S\) at \(s_j\).
Let \(C_j\) be the closure in \(Y_j\) of the set of points
\(\{a_i(s_j)\}_{i\in I}\). Then
\[
  \dim C_1\geq \dim C_0.
\]}

\begin{proof}
It is enough to verify the assertion after a base change
\(S'\to S\) together with chosen points \(s'_0,s'_1\) such that
\(s'_j\mapsto s_j\) and \(s'_0\) is a specialization of
\(s'_1\). By EGA II, 7.1.4, we are therefore reduced to the case
in which \(S\) is the spectrum of a valuation ring \(A\),
\(s_0\) is its closed point, and \(s_1\) is its generic point.

Let \(V\) be the closure of \(C_1\) in \(Y\). Clearly
\(C_0\subseteq V\). The lemma is consequently a special case of
the ``well-known'' fact that an irreducible closed subprescheme
\(V\) of a prescheme \(Y/S\) of finite presentation satisfies
\[
  \dim(V\times_S s_1)\geq \dim(V\times_S s_0)
\]
when \(S\) is the spectrum of a valuation ring and
\(V\times_S s_1\neq\varnothing\) (EGA IV, 13.1.6).
\end{proof}

\SGARefTarget{sga3:xviii:lemma:3.12}\paragraph{\SGARefLink{sga3:xviii:lemma:3.12}{Lemma 3.12}.}
\textit{Suppose that \(S\) is the spectrum of a local ring, with
closed point \(s_0\), and let \(\{c_i\}_{i\in I}\) be sections
such that the closure \(C_0\), in the closed fiber \(X_0\), of
the set \(\{c_i(s_0)\}_{i\in I}\) has dimension
\[
  \dim C_0=\dim X_0=n.
\]
Then the condition of \SGARefLink{sga3:xviii:lemma:3.9}{Lemma 3.9} is satisfied.}

\begin{proof}
First note that all the fibers of \(X/S\) have the same
dimension. This follows from EGA IV, 12.1.1(i), and from the fact
that \(X\) has a rational composition law. Lemma 3.11 therefore
implies that, for every morphism \(S_1\to S\), with \(S_1\) the
spectrum of a field, the closure of
\(\{c_i\times_S S_1\}_{i\in I}\) in \(X_{S_1}\) has dimension
\(n\).

Let us verify the condition of \SGARefLink{sga3:xviii:lemma:3.10}{Lemma 3.10}. Let
\((a,b)\in X^2(S')\). For \((c_i^{-1}a)b\) to be defined, it is
necessary and sufficient that \(c_i\) lie in a certain open
subset \(U\subseteq X_{S'}\) that is schematically dense relative
to \(S'\) (Rule 1). We must show that this holds for a suitable
\(i\), locally on \(S'\). It is therefore enough to treat the
case in which \(S'\) is the spectrum of a local ring; membership
of \(c_i\) in \(U(S')\) can then be checked on the closed fiber.
We are thus reduced to \(S'=\operatorname{Spec}k\), where \(k\)
is a field.

With the notation above, take \(S'=S_1\). Then
\[
  \dim C_1=\dim X_{S_1},
\]
and \(U\) is relatively schematically dense in \(X_{S_1}\).
Hence \(U\cap C_1\neq\varnothing\), and therefore
\[
  U\cap\{c_i\times_S S_1\}_{i\in I}\neq\varnothing.
\]
This proves the assertion.
\end{proof}

\noindent\textit{Proof of \SGARefLink{sga3:xviii:thm:3.7}{Theorem 3.7}.}
The proof is now easy. First note the following consequence of
the finiteness assertion in \SGARefLink{sga3:xviii:lemma:3.9}{Lemma 3.9}. Let \(\{A_i\}\) be a
directed system of rings over \(S\), put
\[
  \widetilde A=\varinjlim_i A_i,
\]
and suppose that the hypotheses of \SGARefLink{sga3:xviii:lemma:3.9}{Lemma 3.9} hold over
\(S=\operatorname{Spec}\widetilde A\). Then the object
representing the quotient \(G\) of
\[
  R\rightrightarrows X^2
\]
descends to some \(S_i=\operatorname{Spec}A_i\), together with
the finiteness and flatness properties stated in Lemma 3.9. This
is the usual passage to the limit (EGA IV, 8 and 11).

It follows that, in proving \SGARefLink{sga3:xviii:theorem:3-7:item:iii}{Theorem 3.7(iii)} and \SGARefLink{sga3:xviii:thm:3.7:item:iv}{(iv)}, we may
restrict ourselves to the case
\[
  S=\operatorname{Spec}A,
\]
where \(A\) is a strictly local ring. Let \(x_0\) be a closed
point of the closed fiber \(X_0\) of \(X/S\). There
exists\begingroup
\renewcommand{\thefootnote}{\ensuremath{\ast}}\footnote{Cf. the
note at the foot of page 15 in Exposé VII.}
\addtocounter{footnote}{-1}
\endgroup{} a local, finite, free extension \(A'\) of \(A\), and
a section of
\[
  X'=X\times_S S'
\]
passing through the unique point \(x'_0\) of \(X'\) lying over
\(x_0\). Notice that \(X'_0\to X_0\) is radicial because \(S\)
is strictly local, and hence the residue field of \(A\) is
separably closed.

It follows that there exists a directed system \(\{A_i\}\) of
local rings, flat and finite over \(A\), such that, on putting
\[
  \widetilde A=\varinjlim_i A_i,\qquad
  \widetilde X=X\times_S\operatorname{Spec}\widetilde A,
\]
the prescheme \(\widetilde X\) has a set of sections inducing a
dense set on its closed fiber \(\widetilde X_0\). Indeed, for
each closed point \(x_0\) of \(X_0\), choose an extension
\(A(x_0)\) such that the corresponding base change of \(X\)
admits a section ``passing through \(x_0\)'', and take for the
directed system the finite tensor products of the \(A(x_0)\).
Notice that \(\widetilde A\) is local, being a direct limit of
local rings. Thus \SGARefLink{sga3:xviii:lemma:3.12}{Lemmas 3.12} and \SGARefLink{sga3:xviii:lemma:3.10}{3.10} give the quotient \(G\)
over \(\operatorname{Spec}\widetilde A\), hence, by the preceding
remarks, over one of the \(\operatorname{Spec}A_i\), and therefore
fppf-locally.

In fact, it follows from the constructions that fppf-locally one
can find a finite set of sections
\(\{c_i\}\), \(i=1,\ldots,r\), such that \(G\) is covered by the
\(c_i\times X\). Assertions \SGARefLink{sga3:xviii:lemma:3.8:item:ii}{(ii)} and \SGARefLink{sga3:xviii:thm:3.7:item:iv}{(iv)} follow easily from this
fact, and we leave the verification of \SGARefLink{sga3:xviii:lemma:3.8:item:i}{(i)} to the reader.

\SGARefTarget{sga3:xviii:corollary:3.13}\paragraph{\SGARefLink{sga3:xviii:corollary:3.13}{Corollary 3.13}.}
\textit{The group sheaf \(G\) determined by a group germ
\((X,W)\) over \(S\) is representable in each of the following
situations:}
\begingroup
\renewcommand{\labelenumi}{\textup{(\roman{enumi})}}
\begin{enumerate}
\item \textit{\(S\) is Artinian.}
\item \textit{For every local scheme \(S'\to S\) at a closed
point \(s\) of \(S\), the scheme \(X_{S'}\) has a set of sections
whose restrictions to the closed fiber form a set with closure
of dimension \(\dim X_s\).}
\item\SGAInlineRefTarget{sga3:xviii:corollary:3-13:item:iii} \textit{\(S\) is strictly local and \(X/S\) is smooth.}
\item \textit{There exists an fppf covering \(S'\to S\) such
that \(G_{S'}\) is representable and affine over \(S\).}
\end{enumerate}
\endgroup

Indeed, \SGARefLink{sga3:xviii:lemma:3.8:item:ii}{(ii)} follows from \SGARefLink{sga3:xviii:lemma:3.10}{Lemmas 3.10} and \SGARefLink{sga3:xviii:lemma:3.12}{3.12}; \SGARefLink{sga3:xviii:thm:3.7:item:iii}{(iii)} follows
directly from \SGARefLink{sga3:xviii:lemma:3.8:item:ii}{(ii)} and Hensel's lemma; \SGARefLink{sga3:xviii:thm:3.7:item:iv}{(iv)} follows from descent
for affine schemes; and \SGARefLink{sga3:xviii:lemma:3.8:item:i}{(i)} follows from descent for group
schemes, which is possible here because every finite subset of a
group over a field is known to be contained in an affine open
subset (Exposé VI).

\SGARefTarget{sga3:xviii:bibliography}
\begin{thebibliography}{9}
\bibitem{Weil1948}
A.~Weil,
\emph{Variétés abéliennes et courbes algébriques},
Hermann, Paris, 1948.

\bibitem{Yanagihara1963}
H.~Yanagihara,
\emph{Reduction of group varieties and transformation spaces},
Journ. Sci. Hiroshima Univ., Ser. A-I, vol.~27, no.~1,
June 1963.
\end{thebibliography}

% END INCLUDED SGA3_English_current_buildable_source_20260823/source/inputs/XVIII/tex/components/03_expose_XVIII_section3.tex



% END INCLUDED SGA3_English_current_buildable_source_20260823/source/inputs/XVIII/body.tex


\clearpage
\SGAInputRoot{GUIDE}
% BEGIN INCLUDED SGA3_English_current_buildable_source_20260823/source/inputs/GUIDE/body.tex
% Exact R15 special component.
% BEGIN INCLUDED SGA3_English_current_buildable_source_20260823/source/inputs/GUIDE/guide/tex/SGA3_Tome_III_Mathematical_Guide_body.tex
\chapter*{Tome III Mathematical Guide and Contents}
\addcontentsline{toc}{chapter}{Tome III Mathematical Guide and Contents}
\markboth{Tome III Mathematical Guide}{Tome III Mathematical Guide}
% BEGIN INCLUDED SGA3_English_current_buildable_source_20260823/source/inputs/GUIDE/guide/tex/components/01_tomeIII_contents_english.tex
% SGA 3, Tome III, English Loop-1 reconstruction of the four-page
% contents/guide leaf. The fourth authority page is blank.

\thispagestyle{empty}
\begin{center}
  \vspace*{14mm}
  {\Large\MakeUppercase{Contents}}
\end{center}
\vspace{16mm}

\begingroup
\small
\GuideChapter{\textbf{XIX. Reductive Groups---Generalities}, by
  M.~Demazure}{1}
\GuideSection{1. Review of groups over an algebraically closed field}{1}
\GuideSection{2. Reductive group schemes. Definitions and first
  properties}{8}
\GuideSection{3. Roots and root systems of reductive group schemes}{12}
\GuideSection{4. Roots and vector group schemes}{15}
\GuideSection{5. An instructive example}{20}
\GuideSection{6. Local existence of maximal tori. The Weyl group}{23}
\GuideSection{Bibliography}{25}

\GuideChapter{\textbf{XX. Reductive Groups of Semisimple Rank \(1\)}, by
  M.~Demazure}{27}
\GuideSection{1. Elementary systems. The groups \(U_{\alpha}\) and
  \(U_{-\alpha}\)}{27}
\GuideSection{2. Structure of elementary systems}{35}
\GuideSection{3. The Weyl group}{44}
\GuideSection{4. The isomorphism theorem}{51}
\GuideSection{5. Examples of elementary systems; applications}{53}
\GuideSection{6. Generators and relations for an elementary system}{58}

\GuideChapter{\textbf{XXI. Root Data}, by M.~Demazure}{63}
\GuideSection{1. Generalities}{63}
\GuideSection{2. Relations between two roots}{67}
\GuideSection{3. Simple roots, positive roots}{71}
\GuideSection{4. Reduced root data of semisimple rank \(2\)}{84}
\GuideSection{5. The Weyl group: generators and relations}{86}
\GuideSection{6. Morphisms of root data}{90}
\GuideSection{7. Structure}{101}
\GuideSection{Bibliography}{108}

\GuideChapterSplit{XXII. Reductive Groups: Splittings, Subgroups,
  Quotient Groups}{by M.~Demazure}{109}
\GuideSection{1. Roots and coroots. Split groups and root data}{109}
\endgroup

\clearpage
\begingroup
\small
\GuideSection{2. Existence of a splitting. Type of a reductive group}{114}
\GuideSection{3. The Weyl group}{116}
\GuideSection{4. Homomorphisms of split groups}{118}
\GuideSection{5. Subgroups of type \textup{(R)}}{128}
\GuideSection{6. The derived group}{168}
\GuideSection{Bibliography}{175}

\GuideChapter{\textbf{XXIII. Reductive Groups: Uniqueness of Pinned
  Groups}, by M.~Demazure}{177}
\GuideSection{1. Pinnings}{177}
\GuideSection{2. Generators and relations for a pinned group}{182}
\GuideSection{3. Groups of semisimple rank \(2\)}{190}
\GuideSection{4. Uniqueness of pinned groups: fundamental theorem}{202}
\GuideSection{5. Corollaries of the fundamental theorem}{206}
\GuideSection{6. Chevalley systems}{209}
\GuideSection{Bibliography}{213}

\GuideChapter{\textbf{XXIV. Automorphisms of Reductive Groups}, by
  M.~Demazure}{215}
\GuideSection{1. Scheme of automorphisms of a reductive group}{216}
\GuideSection{2. Automorphisms and subgroups}{224}
\GuideSection{3. Dynkin scheme of a reductive group. Quasi-split
  groups}{228}
\GuideSection{4. Isotriviality of reductive groups and of principal
  bundles under reductive groups}{237}
\GuideSection{5. Canonical decomposition of an adjoint or simply
  connected group}{243}
\GuideSection{6. Automorphisms of Borel subgroups of reductive
  groups}{248}
\GuideSection{7. Representability of the functors
  \(\operatorname{Hom}_{S\text{-grp.}}(G,H)\), for \(G\) reductive}{250}
\GuideSection{8. Appendix: Cohomology of a smooth group over a henselian
  ring. Cohomology and the functor \(\prod\)}{261}
\GuideSection{Bibliography}{266}

\GuideChapter{\textbf{XXV. The Existence Theorem}, by M.~Demazure}{267}
\GuideSection{1. Statement of the theorem}{267}
\GuideSection{2. Existence theorem: construction of a group chunk}{269}
\GuideSection{3. Existence theorem: end of the proof}{274}
\GuideSection{4. Appendix}{276}
\GuideSection{Bibliography}{277}

\GuideChapter{\textbf{XXVI. Parabolic Subgroups of Reductive Groups}, by
  M.~Demazure}{279}
\GuideSection{1. Review. Levi subgroups}{279}
\GuideSection{2. Structure of the unipotent radical of a parabolic
  subgroup}{285}
\GuideSection{3. Scheme of parabolic subgroups of a reductive group}{289}
\GuideSection{4. Relative position of two parabolic subgroups}{296}
\GuideSection{5. Conjugacy theorem}{311}
\GuideSection{6. Parabolic subgroups and split tori}{320}
\GuideSection{7. Relative root datum}{322}
\endgroup

\clearpage
\begingroup
\small
\GuideSection{Bibliography}{330}
\endgroup

\clearpage
\thispagestyle{empty}
\null

% END INCLUDED SGA3_English_current_buildable_source_20260823/source/inputs/GUIDE/guide/tex/components/01_tomeIII_contents_english.tex


% END INCLUDED SGA3_English_current_buildable_source_20260823/source/inputs/GUIDE/guide/tex/SGA3_Tome_III_Mathematical_Guide_body.tex


% END INCLUDED SGA3_English_current_buildable_source_20260823/source/inputs/GUIDE/body.tex


\SGAStartExpose{XIX}
% BEGIN INCLUDED SGA3_English_current_buildable_source_20260823/source/inputs/XIX/body.tex
% Mechanically extracted after \mainmatter from the frozen standalone master.

\markboth{Exposé XIX. Reductive Groups---Generalities}{Exposé XIX. Reductive Groups---Generalities}
% BEGIN INCLUDED SGA3_English_current_buildable_source_20260823/source/inputs/XIX/tex/components/00_expose_XIX_title_intro_and_section1_opening_page001.tex
% Authority: Exp19-13oct24.pdf, local p.1 / combined p.1035 / printed p.637.

\chapter*{EXPOSÉ XIX\\[0.4em]
REDUCTIVE GROUPS---GENERALITIES}
\addcontentsline{toc}{chapter}{Exposé XIX. Reductive Groups---Generalities}

\begin{center}
\textit{by M. Demazure}
\end{center}

\begingroup
\renewcommand{\thefootnote}{0}%
\footnotetext{Editor's note: version of 13 October 2024.}%
\endgroup

The remainder of this Seminar is devoted to the study of reductive groups.  Its
principal aim is to generalize Chevalley's classical results (the
\emph{Bible} and \emph{Tôhoku})%
\footnote{Editor's note: see the bibliographical references at the end of this
Exposé.  In particular, the 2005 re-edition of the Chevalley Seminar
1956--1958, cited as \emph{Bible}, was revised by P.\ Cartier; see [Ch05].}
to arbitrary base schemes.  The two central results are the uniqueness theorem
(Exposé XXIII, Theorem 4.1 and \SGARefLink{sga3:xix:paragraph:5-1}{Corollaries 5.1}--5.10) and the existence theorem
(\SGARefLink{sga3:xxv:theorem:1:1}{Exposé XXV, Theorem 1.1}) for pinned reductive group schemes corresponding to
prescribed root data.  The proofs used here are inspired by Chevalley's; the
language of schemes makes it possible to give them greater effectiveness.

The results of the first volume of the \emph{Bible} (Exposés 1--13) will be
used systematically.  By contrast, we shall prove directly over an arbitrary
scheme the results of the second volume (in particular the isomorphism
theorem); knowledge of the proofs over an algebraically closed field is
therefore not absolutely indispensable.

In proving these two fundamental results, we shall use only the most elementary
results from the theory of groups of multiplicative type, found essentially in
Exposés VIII and IX; on the other hand, we shall make essential use of the
results of Exposé XVIII.%
\footnote{Editor's note: more precisely, Proposition \SGARefLink{sga3:xviii:prop:2.3}{XVIII 2.3} (extension of a
``generic homomorphism'' between groups) is used in \SGARefLink{sga3:xxii:prop:4.1.11}{XXII 4.1.11} and then in the
proof of uniqueness in XXIII, as well as in XXIV.  Theorem \SGARefLink{sga3:xviii:thm:3.7}{XVIII 3.7}
(construction of a group from a group germ) is used only in XXV.}
The reader interested specifically in the existence and uniqueness theorems
may, on a first reading, skip Exposés X--XVII.

\SGASection{1}{Review of groups over an algebraically closed field}
\SGARefTarget{sga3:xix:section:1}

\noindent\SGAInlineRefTarget{sga3:xix:paragraph:1-1}\textbf{1.1.}\quad
Throughout this subsection, \(k\) will denote an algebraically closed field.
As announced above, the only results from the \emph{Bible} used below are
contained in its first volume (Exposés 1--13).

% END INCLUDED SGA3_English_current_buildable_source_20260823/source/inputs/XIX/tex/components/00_expose_XIX_title_intro_and_section1_opening_page001.tex

% BEGIN INCLUDED SGA3_English_current_buildable_source_20260823/source/inputs/XIX/tex/components/01_expose_XIX_section1_1_through_1_3_page002.tex
% Authority: Exp19-13oct24.pdf, local p.2 / combined p.1036 / printed pp.638--639.

All the results of the \emph{Bible} (\emph{loc.\ cit.}) concerning semisimple
groups hold more generally for reductive groups (defined below), and their
proofs are identical after the following harmless modifications:
\begin{itemize}
\item Exposé 9, §4, Definition 3: see 1.6.1 below;
\item Exposé 12, §4, Theorem 2(e): omit ``finite'';
\item Exposé 13, §3, Theorem 2: replace ``rank'' by ``semisimple rank'';
\item Exposé 13, §4, Corollary 2 to Theorem 3: replace ``rank'' by
  ``reductive rank''.
\end{itemize}

\noindent\SGAInlineRefTarget{sga3:xix:paragraph:1-2}\textbf{1.2.}\quad
Let \(G\) be a smooth affine connected \(k\)-group.  The \emph{radical} of
\(G\) (\emph{Bible}, §9.4, Proposition 2)%
\footnote{Editor's note: the reference 9-06 (that is, Exposé 9, p.\ 6) has
been replaced by §9.4 (that is, Exposé 9, §4), which applies both to the
\emph{Bible} and to [Ch05].  When the numbering in [Ch05] differs from that
in the \emph{Bible}, as happens in Exposé 6, both references will be given
explicitly.}
is the reduced subgroup associated with the identity component of the
intersection of the Borel subgroups of \(G\); it is also the largest smooth
connected solvable normal subgroup of \(G\).  We denote it by
\(\operatorname{rad}(G)\).

The \emph{unipotent radical} of \(G\) is the unipotent part of the radical of
\(G\); it is also the largest smooth connected unipotent normal subgroup of
\(G\).  We denote it by \(\operatorname{rad}^{u}(G)\).

\noindent\SGAInlineRefTarget{sga3:xix:paragraph:1-3}\textbf{1.3.}\quad
Let \(G\) be a smooth affine connected \(k\)-group, and let \(Q\) be a torus
in \(G\).  Then the centralizer
\(\underline{\operatorname{Centr}}_{G}(Q)\) of \(Q\) in \(G\) is a closed
subgroup of \(G\) (\SGARefLink{sga3:viii:corollary:6:7}{Exposé VIII, 6.7}), smooth (\SGARefLink{sga3:xi:corollary:2:4}{Exposé XI, 2.4}), and connected
(cf.\ \emph{Bible}, §6.6, Theorem 6(a), or [Ch05], §6.7, Theorem 7(a)).

We have the fundamental relation
\begin{equation}
  \operatorname{rad}^{u}\!\bigl(\underline{\operatorname{Centr}}_{G}(Q)\bigr)
  =
  \operatorname{rad}^{u}(G)\cap
  \underline{\operatorname{Centr}}_{G}(Q).
  \tag{1.3.1}\label{sga3:xix:eq:1:3:1}
\end{equation}

First,%
\footnote{Editor's note: the argument that follows has been expanded.}
\[
  U=\operatorname{rad}^{u}(G)\cap
  \underline{\operatorname{Centr}}_{G}(Q)
\]
is a normal unipotent subgroup of
\(\underline{\operatorname{Centr}}_{G}(Q)\); let us show that it is smooth
and connected.  Let \(Q\) act on \(\operatorname{rad}^{u}(G)\) by inner
automorphisms.  Then \(U\) is precisely the scheme of fixed points of this
action.  This scheme is smooth and connected by \SGARefLink{sga3:xix:lemma:1:4}{Lemma 1.4} below.

Consequently, \(U\) is a closed subgroup of
\(\operatorname{rad}^{u}(\underline{\operatorname{Centr}}_{G}(Q))\).
On the other hand, by the \emph{Bible}, Exposé 12, §3, Corollary to
Theorem 1,
\[
  \operatorname{rad}^{u}\!\bigl(
    \underline{\operatorname{Centr}}_{G}(Q)\bigr)(k)
  \subset
  \operatorname{rad}^{u}(G)(k).
\]
Equality \eqref{sga3:xix:eq:1:3:1} follows.

Let us record one special case of the preceding result.  If \(G\) is a smooth
affine connected \(k\)-group and \(T\) is a maximal torus of \(G\), then
\begin{equation}
  \underline{\operatorname{Centr}}_{G}(T)
  =
  T\cdot\bigl(
    \underline{\operatorname{Centr}}_{G}(T)\cap
    \operatorname{rad}^{u}(G)\bigr).
  \tag{1.3.2}\label{sga3:xix:eq:1:3:2}
\end{equation}
Indeed (cf.\ \emph{Bible}, §6.6, Theorem 6(c), or [Ch05], §6.7,
Theorem 7(c)), \(\underline{\operatorname{Centr}}_{G}(T)\) is solvable and
hence is the semidirect product of its maximal torus \(T\) and its unipotent
radical.

By the density theorem (cf.\ \emph{Bible}, §6.5, Theorem 5(a), or [Ch05],
§6.6, Theorem 6(a)), the union of
\(\underline{\operatorname{Centr}}_{G}(T)\), as \(T\) ranges over the
maximal tori of \(G\), contains a dense open subset of \(G\).  It therefore
follows from \eqref{sga3:xix:eq:1:3:2} that:

% END INCLUDED SGA3_English_current_buildable_source_20260823/source/inputs/XIX/tex/components/01_expose_XIX_section1_1_through_1_3_page002.tex

% BEGIN INCLUDED SGA3_English_current_buildable_source_20260823/source/inputs/XIX/tex/components/02_expose_XIX_corollary1_3_3_through_section1_5_page003.tex
% Authority: Exp19-13oct24.pdf, local p.3 / combined p.1037 / printed p.639.

\SGAStatement{Corollary 1.3.3}{%
If \(G\) is a smooth affine connected \(k\)-group, then the union of
\(T\cdot\operatorname{rad}^{u}(G)\), where \(T\) ranges over the maximal tori
of \(G\), contains a dense open subset of \(G\).}
\SGARefTarget{sga3:xix:corollary:1:3:3}

\par\medskip\noindent
\textbf{\textit{\SGARefLink{sga3:xix:notation:1:4:0}{Notation 1.4.0}.}}---%
\footnote{Editor's note: this notation, which appeared in \SGARefLink{sga3:xix:lemma:2-3}{2.3} of the original,
has been inserted here; moreover, the statement and proof of \SGARefLink{sga3:xix:lemma:1-4}{1.4} have been
expanded.}
\textit{We shall systematically use the following notation.  If \(S\) is a
scheme and \(s\) is a point of \(S\), then \(\bar{s}\) denotes the spectrum of
an algebraic closure \(\overline{\kappa(s)}\) of \(\kappa(s)\).}
\SGARefTarget{sga3:xix:notation:1:4:0}

\par\medskip\noindent
\SGAInlineRefTarget{sga3:xix:lemma:1-4}\textbf{\textit{\SGARefLink{sga3:xix:lemma:1-4}{Lemma 1.4}.}}---\footnotemark[5]
\textit{Let \(S\) be a scheme and let \(Q\) be an \(S\)-torus acting on an
\(S\)-group scheme \(H\), separated and smooth over \(S\).
\begin{enumerate}
\item[\SGAInlineRefTarget{sga3:xix:notation:1:4:0:item:i}(i)] The fixed-point functor \(H^{Q}\) is representable by a closed
  subscheme of \(H\), smooth over \(S\).
\item[(ii)] If, moreover, \(H\) is affine over \(S\) and has connected fibers,
  then the same holds for \(H^{Q}\).
\end{enumerate}}
\SGARefTarget{sga3:xix:lemma:1:4}

\SGAProof{Proof}
First, by \SGARefLink{sga3:viii:examples:6:5:item:d:01}{VIII 6.5(d)},%
\footnote{Editor's note: see also the additions made in
VI\textsubscript{B}, \SGARefLink{ex:sga3-VIB-6.2.4-d}{6.2.4(d)} and \SGARefLink{concl:sga3-VIB-6.5.5-a}{6.5.5(a)}.}
since \(H\) is separated over \(S\) and \(Q\) is essentially free over \(S\),
\(H^{Q}\) is representable by a closed subscheme of \(H\).  In particular, if
\(H\) is affine over \(S\), then so is \(H^{Q}\).

Now consider the semidirect product \(G=H\cdot Q\); it is smooth and separated
over \(S\).  The centralizer \(\underline{\operatorname{Centr}}_{G}(Q)\) is
then representable by a closed subscheme of \(G\), of finite presentation over
\(G\), by \SGARefLink{sga3:xi:corollary:6:11:item:a}{Exposé XI 6.11(a)},\footnotemark[6] and it is smooth over \(S\) by
\SGARefLink{sga3:xi:corollary:2:4}{Exposé XI 2.4}.

Since
\[
  \underline{\operatorname{Centr}}_{G}(Q)=H^{Q}\cdot Q,
\]
it follows that \(H^{Q}\) is smooth over \(S\).  Indeed, using the section
\[
  \sigma:H^{Q}\longrightarrow
  \underline{\operatorname{Centr}}_{G}(Q)
\]
deduced from the identity section \(\epsilon_{Q}\) of \(Q\), one sees at once
that \(H^{Q}\) is formally smooth over \(S\); and since \(\epsilon_{Q}\) is of
finite presentation, \(H^{Q}\) is locally of finite presentation over \(S\).
This proves \SGARefLink{sga3:xix:notation:1:4:0:item:i}{(i)}.

Finally, suppose that \(H\) is affine over \(S\) and has connected fibers.
Then every geometric fiber \(G_{\bar{s}}\) of \(G\) is a smooth affine
connected \(\kappa(\bar{s})\)-group.  Hence, by the first assertion of \SGARefLink{sga3:xix:paragraph:1-3}{1.3},
the group
\[
  \underline{\operatorname{Centr}}_{G_{\bar{s}}}(Q_{\bar{s}})
  =
  \bigl(\underline{\operatorname{Centr}}_{G}(Q)\bigr)_{\bar{s}}
  =
  (H^{Q})_{\bar{s}}\cdot Q_{\bar{s}}
\]
is also connected, and this implies that \((H^{Q})_{\bar{s}}\) is connected.
\hfill\(\square\)

\par\medskip\noindent\SGAInlineRefTarget{sga3:xix:paragraph:1-5}\textbf{1.5.}\quad
Recall that the \emph{reductive rank} of a smooth affine \(k\)-group \(G\) is
the common dimension of its maximal tori.  We denote it by
\(\operatorname{rgred}(G/k)\) or \(\operatorname{rgred}(G)\).  The equality
\(\operatorname{rgred}(G/k)=0\) holds if and only if \(G\) is unipotent
(that is, \(G=\operatorname{rad}^{u}(G)\)), by the \emph{Bible}, §6.4,
Corollary 1 to \SGARefLink{sga3:xix:lemma:1:4}{Theorem 4}, or [Ch05], §6.5, Corollary 1 to \SGARefLink{sga3:xix:theorem:2:5}{Theorem 5}.

If \(H\) is a normal subgroup of the smooth affine \(k\)-group \(G\), then the
quotient \(G/H\) is affine and smooth (Exposé VI\textsubscript{B}, \SGARefLink{thm:sga3-VIB-11.17}{11.17} and
VI\textsubscript{A}, \SGARefLink{thm:sga3-VIA-3.3.2:item:iii}{3.3.2(iii)}).  Moreover (\emph{Bible}, §7.3,
Theorem 3(a),(c)),
\[
  \operatorname{rgred}(G)
  =
  \operatorname{rgred}(G/H)+\operatorname{rgred}(H).
\]

% END INCLUDED SGA3_English_current_buildable_source_20260823/source/inputs/XIX/tex/components/02_expose_XIX_corollary1_3_3_through_section1_5_page003.tex

% BEGIN INCLUDED SGA3_English_current_buildable_source_20260823/source/inputs/XIX/tex/components/03_expose_XIX_definition1_6_1_through_section1_7_page004.tex
\par\medskip\noindent
\SGAInlineRefTarget{sga3:xix:definition:1-6-1}\textbf{\textit{\SGARefLink{sga3:xix:definition:1:6:1}{Definition 1.6.1}.}}---\footnote{Editor's note: the former
number 1.6 has been replaced by numbers \SGARefLink{sga3:xix:definition:1-6-1}{1.6.1} through \SGARefLink{sga3:xix:remark:1-6-3}{1.6.3}.  Note that in
this Seminar every reductive \(k\)-group (respectively, every semisimple
\(k\)-group; cf.\ \SGARefLink{sga3:xix:paragraph:1-8}{1.8}) is connected by definition.}
\textit{A \(k\)-group \(G\) is said to be \emph{reductive} if it is smooth,
affine, and connected, and if \(\operatorname{rad}(G)\) is a torus; that is,
if
\[
  \operatorname{rad}^{u}(G)=\{e\}.
\]}
\SGARefTarget{sga3:xix:definition:1:6:1}

\par\medskip\noindent
\textbf{\textit{\SGARefLink{sga3:xix:lemma:1:6:2}{Lemma 1.6.2}.}}---
\textit{Let \(G\) be a reductive \(k\)-group.
\begin{enumerate}
\item[\SGAInlineRefTarget{sga3:xix:definition:1:6:1:item:i}(i)]\SGAInlineRefTarget{sga3:xix:lemma:1:6:2:item:i} If \(Q\) is a torus of \(G\), then
  \(\underline{\operatorname{Centr}}_{G}(Q)\) is reductive.
\item[\SGAInlineRefTarget{sga3:xix:definition:1:6:1:item:ii}(ii)]\SGAInlineRefTarget{sga3:xix:lemma:1:6:2:item:ii} In particular, if \(T\) is a maximal torus of \(G\), then
  \(\underline{\operatorname{Centr}}_{G}(T)=T\).
\item[\SGAInlineRefTarget{sga3:xix:definition:1:6:1:item:iii}(iii)] The center of \(G\) is contained in every maximal torus and is
  therefore diagonalizable.
\item[\SGAInlineRefTarget{sga3:xix:definition:1:6:1:item:iv}(iv)]\SGAInlineRefTarget{sga3:xix:lemma:1:6:2:item:iv} The radical of \(G\) is the unique maximal torus of
  \(\underline{\operatorname{Centr}}(G)\).
\end{enumerate}}
\SGARefTarget{sga3:xix:lemma:1:6:2}

\SGAProof{Proof}
Indeed, \SGARefLink{sga3:xix:definition:1:6:1:item:i}{(i)} follows from (\SGARefLink{sga3:xix:eq:1:3:1}{1.3.1}), \SGARefLink{sga3:xix:definition:1:6:1:item:ii}{(ii)} from (\SGARefLink{sga3:xix:eq:1:3:2}{1.3.2}), and \SGARefLink{sga3:xix:definition:1:6:1:item:iii}{(iii)} follows from
\SGARefLink{sga3:xix:definition:1:6:1:item:ii}{(ii)}.  Finally, the maximal torus of
\(\underline{\operatorname{Centr}}(G)\), that is, the neutral component
\(\underline{\operatorname{Centr}}(G)^{0}\), is a smooth connected solvable
subgroup normal in \(G\), and hence is contained in
\(\operatorname{rad}(G)\).  Conversely, \(\operatorname{rad}(G)\) is a torus
normal in \(G\), and hence central (\emph{Bible}, §\SGARefLink{sga3:xix:paragraph:4-3}{4.3}, Corollary to
Proposition 2); this proves \SGARefLink{sga3:xix:definition:1:6:1:item:iv}{(iv)}.
\hfill\(\square\)

\par\medskip\noindent
\SGAInlineRefTarget{sga3:xix:remark:1-6-3}\textbf{\textit{\SGARefLink{sga3:xix:remark:1:6:3}{Remark 1.6.3}.}}---
\textit{If \(G\) is reductive and
\(\dim(G)\ne\operatorname{rgred}(G)\), then
\[
  \dim(G)-\operatorname{rgred}(G)\geq 2.
\]
Indeed, this difference is always even (cf.\ \SGARefLink{sga3:xix:paragraph:1-10}{1.10} below).}
\SGARefTarget{sga3:xix:remark:1:6:3}

\par\medskip\noindent\SGAInlineRefTarget{sga3:xix:paragraph:1-7}\textbf{1.7.}\quad
Let \(G\) be a smooth affine connected \(k\)-group, and let \(H\) be a
smooth connected normal subgroup.  Then
\begin{equation}
  \operatorname{rad}(H)
  =
  \bigl(\operatorname{rad}(G)\cap H\bigr)_{\mathrm{red}}^{0},
  \qquad
  \operatorname{rad}^{u}(H)
  =
  \bigl(\operatorname{rad}^{u}(G)\cap H\bigr)_{\mathrm{red}}^{0},
  \tag{1.7.1}
\end{equation}
as is seen immediately.  In particular, if \(G\) is reductive, then so is
\(H\).
\SGARefTarget{sga3:xix:equation:1:7:1}

Let \(f:G\longrightarrow G'\) be a surjective morphism of smooth affine
connected \(k\)-groups.  Then
\begin{equation}
  f\bigl(\operatorname{rad}^{u}(G)\bigr)
  =
  \operatorname{rad}^{u}(G').
  \tag{1.7.2}
\end{equation}
In particular, if \(G\) is reductive, then so is \(G'\).
\SGARefTarget{sga3:xix:equation:1:7:2}

Let us prove (\SGARefLink{sga3:xix:equation:1:7:2}{1.7.2}).  First,
\(f\bigl(\operatorname{rad}^{u}(G)\bigr)\) is contained in
\(\operatorname{rad}^{u}(G')\).  On introducing
\[
  H=
  \bigl(f^{-1}(\operatorname{rad}^{u}(G'))\bigr)_{\mathrm{red}}^{0},
\]
which contains \(\operatorname{rad}^{u}(G)\), we have
\(\operatorname{rad}^{u}(H)=\operatorname{rad}^{u}(G)\); thus we are reduced
to the case \(G=H\), that is, to the case in which \(G'\) is unipotent.
As \(T\) runs over the maximal tori of \(G\), the union of the groups
\(T\operatorname{rad}^{u}(G)\) is dense in \(G\).  Hence the union of the
groups \(f(T)\,f\bigl(\operatorname{rad}^{u}(G)\bigr)\) is dense in \(G'\).
But \(f(T)\) consists of semisimple elements, so \(f(T)=\{e\}\), since
\(G'\) is unipotent.  It follows that
\(f\bigl(\operatorname{rad}^{u}(G)\bigr)\) is dense in \(G'\).  Therefore,
by \emph{Bible}, §\SGARefLink{sga3:xix:paragraph:5:4}{5.4}, \SGARefLink{sga3:xix:lemma:1:4}{Lemma 4}, or [Ch05], §6.1, \SGARefLink{sga3:xix:lemma:3:1}{Lemma 1},\footnote{Editor's
note: the following argument has been supplied in detail.}
\(f\bigl(\operatorname{rad}^{u}(G)\bigr)\) is an open subgroup of \(G'\).
Since \(G'\) is connected, it follows that
\[
  f\bigl(\operatorname{rad}^{u}(G)\bigr)=G'.
\]

\noindent\textit{N.B.}\quad Under the same hypotheses one can prove that
\[
  f\bigl(\operatorname{rad}(G)\bigr)=\operatorname{rad}(G').
\]

% END INCLUDED SGA3_English_current_buildable_source_20260823/source/inputs/XIX/tex/components/03_expose_XIX_definition1_6_1_through_section1_7_page004.tex

% BEGIN INCLUDED SGA3_English_current_buildable_source_20260823/source/inputs/XIX/tex/components/04_expose_XIX_sections1_8_through_1_10_pages005_006.tex
\par\medskip\noindent\SGAInlineRefTarget{sga3:xix:paragraph:1-8}\textbf{1.8.}\quad
A \(k\)-group \(G\) is said to be \emph{semisimple} if it is smooth,
affine, and connected, and if
\[
  \operatorname{rad}(G)=\{e\},
\]
that is, if \(G\) is reductive and
\(\underline{\operatorname{Centr}}(G)\) is finite.  If \(G\) is an
arbitrary smooth affine connected \(k\)-group, then
\(G/\operatorname{rad}(G)\) is semisimple (\emph{Bible}, §9.4,
Proposition 2), and \(G/\operatorname{rad}^{u}(G)\) is reductive.  The
\emph{semisimple rank of \(G\)}, denoted by
\(\operatorname{rgss}(G/k)\), or simply by \(\operatorname{rgss}(G)\), is
the reductive rank of \(G/\operatorname{rad}(G)\).
\SGARefTarget{sga3:xix:section:1:8}

Thus, if \(G\) is reductive,
\[
  \operatorname{rgred}(G)
  =
  \operatorname{rgss}(G)+\dim\bigl(\operatorname{rad}(G)\bigr).
\]

If \(G\) is a smooth affine connected \(k\)-group and \(Q\) is a central
subtorus of \(G\), then \(G/Q\) is semisimple if and only if \(G\) is
reductive and \(Q=\operatorname{rad}(G)\).  Indeed, if \(G/Q\) is
semisimple, then \(Q\supset\operatorname{rad}(G)\); hence
\(\operatorname{rad}(G)\) is a torus, so \(G\) is reductive, and
\(\operatorname{rad}(G)\) is the maximal torus of
\(\underline{\operatorname{Centr}}(G)\), whence
\(\operatorname{rad}(G)=Q\).  If \(G\) is reductive and \(Q\) is a central
subgroup, then \(G/Q\) is reductive and
\[
  \operatorname{rgss}(G)=\operatorname{rgss}(G/Q).
\]

\par\medskip\noindent\SGAInlineRefTarget{sga3:xix:paragraph:1-9}\textbf{1.9.}\quad
Let \(K\) be an algebraically closed extension of \(k\), and let \(G\) be
a smooth affine connected \(k\)-group.  Then \(G\) is reductive
(respectively, semisimple) if and only if \(G_K\) is, and
\[
  \operatorname{rgred}(G/k)=\operatorname{rgred}(G_K/K),
  \qquad
  \operatorname{rgss}(G/k)=\operatorname{rgss}(G_K/K).
\]
\SGARefTarget{sga3:xix:section:1:9}

\par\medskip\noindent\SGAInlineRefTarget{sga3:xix:paragraph:1-10}\textbf{1.10.}\quad
Let \(G\) be a smooth connected \(k\)-group and let \(T\) be a torus of
\(G\).\footnote{Editor's note: the following sentence has been added.}
Write \(\mathfrak g\) for the \(k\)-Lie algebra of \(G\), that is,
\(\mathfrak g=\operatorname{Lie}(G)=\operatorname{Lie}(G/k)(k)\), and
similarly write \(\mathfrak t=\operatorname{Lie}(T)\).  Under the action
of \(T\), through \(T\longrightarrow G\) and the adjoint representation
of \(G\), the Lie algebra \(\mathfrak g\) decomposes as
\begin{equation}
  \mathfrak g
  =
  \mathfrak g^{0}
  \oplus
  \bigoplus_{\alpha\in R}\mathfrak g^{\alpha},
  \tag{1.10.1}\label{sga3:xix:paragraph:1-10:equation:1-10-1}
\end{equation}
where the elements \(\alpha\in R\) are nontrivial characters of \(T\)
and the spaces \(\mathfrak g^{\alpha}\) are nonzero.  The characters
\(\alpha\in R\) are called the \emph{roots of \(G\) relative to \(T\)}.
By \SGARefLink{sga3:ii:theorem:5-2-3:item:i}{Exposé II, 5.2.3(i)},
\begin{equation}
  \mathfrak g^{0}
  =
  \operatorname{Lie}\bigl(\underline{\operatorname{Centr}}_{G}(T)\bigr).
  \tag{1.10.2}
\end{equation}
\SGARefTarget{sga3:xix:equation:1:10:1}
\SGARefTarget{sga3:xix:equation:1:10:2}

In particular,\footnote{Editor's note: the following argument has been
given in detail.} since
\(\underline{\operatorname{Centr}}_{G}(T)\) is connected (by \SGARefLink{sga3:xix:paragraph:1-3}{1.3} when
\(G\) is affine, and by \SGARefLink{sga3:xii:6.6:item:b}{Exposé XII, 6.6(b)}, in general), \(T\) is its own
centralizer if and only if
\(\mathfrak g^{0}=\operatorname{Lie}(T)\).  This condition implies that
\(T\) is maximal and that
\(\underline{\operatorname{Centr}}(G)\subset T\); hence, by \SGARefLink{sga3:xii:8.8:item:d:02}{Exposé XII,
8.8(d)},
\begin{equation}
  \underline{\operatorname{Centr}}(G)
  =
  \operatorname{Ker}\bigl(T\longrightarrow
    \operatorname{GL}(\mathfrak g)\bigr)
  =
  \bigcap_{\alpha\in R}\operatorname{Ker}(\alpha).
  \tag{1.10.3}
\end{equation}
\SGARefTarget{sga3:xix:equation:1:10:3}
Moreover, \(\underline{\operatorname{Centr}}(G)\) is then affine, and
therefore the morphism
\(G\longrightarrow G/\underline{\operatorname{Centr}}(G)\) is affine.
Since the latter group is affine (\SGARefLink{sga3:xii:6.1}{Exposé XII, 6.1}),\footnote{Editor's
note: see also Addition VI\textsubscript{B}, \SGARefLink{cor:sga3-VIB-6.2.6}{6.2.6}.} \(G\) is affine.

When \(G\) is reductive and \(T\) is maximal, the roots in the preceding
sense coincide with the roots in the sense of \emph{Bible}, §12.2,
Definition 1.  The latter are indeed roots in the sense of this number
(\emph{Bible}, §13.2, Theorem 1(c)), and their number is
\(\dim(G)-\dim(T)\) (\emph{Bible}, §13.4, Corollary 2 to Theorem 3).
Moreover, if \(\alpha\) is a root, then so is \(-\alpha\)
(\emph{Bible}, §12.2, Corollary to Proposition 1).  As usual, we write
the group law of
\(\operatorname{Hom}_{k\text{-}\mathrm{gr}}
  (T,\mathbb G_{m,k})\)
indifferently additively or multiplicatively.  It follows that, for
reductive \(G\),
\begin{equation}
  \dim(G)-\operatorname{rgred}(G)=\operatorname{Card}(R)
  \tag{1.10.4}
\end{equation}
is always even.
\SGARefTarget{sga3:xix:equation:1:10:4}

The semisimple rank of the reductive group \(G\) is the rank of the
subset \(R\) of the \(\mathbb Q\)-vector space
\[
  \operatorname{Hom}_{k\text{-}\mathrm{gr}}
    (T,\mathbb G_{m,k})\otimes_{\mathbb Z}\mathbb Q
\]
(\emph{Bible}, §13.3, Theorem 2).
\SGARefTarget{sga3:xix:section:1:10}

% END INCLUDED SGA3_English_current_buildable_source_20260823/source/inputs/XIX/tex/components/04_expose_XIX_sections1_8_through_1_10_pages005_006.tex

% BEGIN INCLUDED SGA3_English_current_buildable_source_20260823/source/inputs/XIX/tex/components/05_expose_XIX_lemma1_11_opening_page006.tex
\par\medskip\noindent
\SGAInlineRefTarget{sga3:xix:lemma:1-11}\textbf{\textit{\SGARefLink{sga3:xix:lemma:1:11}{Lemma 1.11}.}}---
\textit{Let \(k\) be an algebraically closed field, \(G\) a smooth affine
connected \(k\)-group, and \(T\) a torus of \(G\).  Put
\[
  W(T)
  =
  \underline{\operatorname{Norm}}_{G}(T)/
  \underline{\operatorname{Centr}}_{G}(T),
\]
the Weyl group of \(G\) relative to \(T\).  The following conditions are
equivalent:
\begin{enumerate}
\item[(i)] \(G\) is reductive, \(T\) is maximal, and
  \(\operatorname{rgss}(G)=1\).
\item[(ii)] \(G\) is reductive, \(T\) is maximal, and \(G\ne T\); there
  exists a subtorus \(Q\) of \(T\), of codimension \(1\) in \(T\), that
  is central in \(G\).
\item[(iii)] \(G\) is not solvable and
  \(\dim(G)-\dim(T)\leq 2\).
\item[(iv)] \(W(T)\ne\{e\}\) and
  \(\dim(G)-\dim(T)\leq 2\).
\item[(v)] \(W(T)=(\mathbb Z/2\mathbb Z)_k\) and
  \(\dim(G)-\dim(T)=2\).
\end{enumerate}
Moreover, under these conditions there are exactly two roots of \(G\)
relative to \(T\), and they are opposite.  Under the conditions of
\textup{(ii)}, \(Q=\operatorname{rad}(G)\).}
\SGARefTarget{sga3:xix:lemma:1:11}

\SGAProof{Proof}
The implication (v)\(\Rightarrow\)(iv) is clear, and
(iv)\(\Rightarrow\)\SGARefLink{sga3:xix:lemma:1:11:item:iii}{(iii)} follows from \emph{Bible}, §6.1, Corollary 3 to
Theorem 1, or [Ch05], §6.2, Corollary 3 to Theorem 2.

Let us prove \SGARefLink{sga3:xix:lemma:1:11:item:iii}{(iii)}\(\Rightarrow\)\SGARefLink{sga3:xix:lemma:1:11:item:ii}{(ii)}.  Let \(U\) be the unipotent radical
of \(G\).  The group \(G/U\) is reductive and is not a torus (otherwise
\(G\) would be solvable); hence, by (\SGARefLink{sga3:xix:equation:1:10:4}{1.10.4}),
\[
  \dim(G/U)-\operatorname{rgred}(G/U)
  =
  \operatorname{Card}(R)
  \geq 2.
\]
On the other hand,
\[
  \operatorname{rgred}(G)
  =
  \operatorname{rgred}(G/U)+\operatorname{rgred}(U)
  =
  \operatorname{rgred}(G/U),
\]
and therefore
\[
\begin{aligned}
  \dim(G)-\dim(U)-\operatorname{rgred}(G)
    &=\dim(G/U)-\operatorname{rgred}(G/U)\\
    &\geq 2
     \geq \dim(G)-\dim(T)
     \geq \dim(G)-\operatorname{rgred}(G).
\end{aligned}
\]
It follows that \(\dim(U)=0\), so \(G\) is reductive, and that
\(\operatorname{rgred}(G)=\dim(T)\); thus \(T\) is maximal and
\(\dim(G)-\dim(T)=2\).  By \emph{Bible}, §10.4, Proposition 8, there
exists a singular torus \(Q\) of codimension \(1\) in \(T\).  Then
\(\underline{\operatorname{Centr}}_{G}(Q)\) is reductive (1.6.2(i)) and
not solvable (by the definition of a singular torus), and hence
\[
  \dim\bigl(\underline{\operatorname{Centr}}_{G}(Q)\bigr)
  -\operatorname{rgred}(G)\geq 2.
\]
This proves that
\(G=\underline{\operatorname{Centr}}_{G}(Q)\), and that \(Q\) is central
in \(G\).

Let us prove \SGARefLink{sga3:xix:lemma:1:11:item:ii}{(ii)}\(\Rightarrow\)\SGARefLink{sga3:xix:lemma:1:11:item:i}{(i)}.  We know that \(G/Q\) is reductive
by \SGARefLink{sga3:xix:paragraph:1-7}{1.7} and that \(\operatorname{rgred}(G/Q)=1\); hence
\(\operatorname{rgss}(G/Q)\) is either \(0\) or \(1\).  The first case is
impossible, since it would imply \(\operatorname{rgss}(G)=0\), and
therefore \(G=T\).  Thus
\[
  \operatorname{rgred}(G/Q)=\operatorname{rgss}(G/Q)=1,
\]
which proves that \(G/Q\) is semisimple.  Consequently \(Q\) is the
radical of \(G\) and \(\operatorname{rgss}(G)=1\).

% END INCLUDED SGA3_English_current_buildable_source_20260823/source/inputs/XIX/tex/components/05_expose_XIX_lemma1_11_opening_page006.tex

% BEGIN INCLUDED SGA3_English_current_buildable_source_20260823/source/inputs/XIX/tex/components/06_expose_XIX_lemma1_11_completion_and_proposition1_12_pages007_008.tex
Let us finally prove (i)\(\Rightarrow\)(v).  If \(Q\) is the radical of
\(G\), then
\(\dim(T)-\dim(Q)=1\), and \(Q\) is central in \(G\).  Hence
\(G=\underline{\operatorname{Centr}}_{G}(Q)\), which proves that \(Q\)
is a singular torus.  By \emph{Bible}, §11.3, Theorem 2,
\[
  W_G(T)=(\mathbb Z/2\mathbb Z)_k,
\]
and, by \emph{Bible}, §12.1, Lemma 1,
\(\dim(G)-\dim(T)=2\).  Thus \(G\) has at most two roots relative to
\(T\), while it has at least two, opposite to one another, by \SGARefLink{sga3:xix:paragraph:1-10}{1.10}.

\par\medskip\noindent
\SGAInlineRefTarget{sga3:xix:proposition:1-12}\textbf{\textit{\SGARefLink{sga3:xix:proposition:1:12}{Proposition 1.12}.}}---
\textit{Let \(k\) be an algebraically closed field, \(G\) a smooth
connected \(k\)-group, \(T\) a torus of \(G\), and \(R\) the set of
roots of \(G\) relative to \(T\).  Write
\[
  \mathfrak g
  =
  \mathfrak g^{0}
  \oplus
  \bigoplus_{\alpha\in R}\mathfrak g^{\alpha},
  \qquad
  \mathfrak g^{\alpha}\ne0,
\]
for the decomposition of \(\mathfrak g\) under
\(\operatorname{Ad}(T)\).  For every \(\alpha\in R\), put
\[
  T_{\alpha}=\operatorname{Ker}(\alpha)^{0},
  \qquad
  Z_{\alpha}=\underline{\operatorname{Centr}}_{G}(T_{\alpha}).
\]
\footnote{Editor's note: thus \(T_{\alpha}\) is the identity component
of \(\operatorname{Ker}(\alpha)\).}
The following conditions are equivalent:
\begin{enumerate}
\item[\SGAInlineRefTarget{sga3:xix:lemma:1:11:item:i}(i)] \(G\) is affine and reductive, and \(T\) is maximal.
\item[\SGAInlineRefTarget{sga3:xix:lemma:1:11:item:ii}(ii)] \(\mathfrak g^{0}=\mathfrak t\), and every \(Z_{\alpha}\),
  for \(\alpha\in R\), is reductive.
\item[\SGAInlineRefTarget{sga3:xix:lemma:1:11:item:iii}(iii)] \(\mathfrak g^{0}=\mathfrak t\); every
  \(\mathfrak g^{\alpha}\), for \(\alpha\in R\), has dimension \(1\);
  if \(\alpha,\beta\in R\) and \(\beta=q\alpha\) for some
  \(q\in\mathbb Q\), then \(q=\pm1\); and, for every \(\alpha\in R\),
  there exists \(w_{\alpha}\in G(k)\) which normalizes \(T\),
  centralizes \(T_{\alpha}\), but does not centralize \(T\).
\end{enumerate}
Moreover, under these conditions every \(Z_{\alpha}\) has semisimple
rank \(1\), and
\[
  \operatorname{Lie}(Z_{\alpha})
  =
  \mathfrak t\oplus
  \mathfrak g^{\alpha}\oplus
  \mathfrak g^{-\alpha}.
\]}
\SGARefTarget{sga3:xix:proposition:1:12}

\SGAProof{Proof}
Suppose first that \(G\) is affine and reductive and that \(T\) is
maximal.  Every \(Z_{\alpha}\) is then affine and reductive by \SGARefLink{sga3:xix:lemma:1:6:2:item:i}{1.6.2(i)},
with maximal torus \(T\); furthermore, \(T\) is its own centralizer by
\SGARefLink{sga3:xix:lemma:1:6:2:item:ii}{1.6.2(ii)}.  Thus \(\mathfrak g^{0}=\mathfrak t\), proving
\SGARefLink{sga3:xix:proposition:1:12:item:i}{(i)}\(\Rightarrow\)\SGARefLink{sga3:xix:proposition:1:12:item:ii}{(ii)}.

Conversely, the equality \(\mathfrak g^{0}=\mathfrak t\) implies that
\(T\) is maximal and that \(G\) is affine, by \SGARefLink{sga3:xix:paragraph:1-10}{1.10}.  Hence every
\(Z_{\alpha}\) is affine, smooth, and connected, by \SGARefLink{sga3:xix:paragraph:1-3}{1.3}.  In all cases,
\SGARefLink{sga3:ii:lemma:5-2-2}{Exposé II, 5.2.2} gives
\[
  \operatorname{Lie}(Z_{\alpha})
  =
  \mathfrak g^{T_{\alpha}}
  =
  \mathfrak g^{0}
  \oplus
  \bigoplus_{\beta\in R\cap\mathbb Q\alpha}
    \mathfrak g^{\beta}.
\]
Consequently,
\[
  \operatorname{Lie}(Z_{\alpha})
  \supset
  \mathfrak t\oplus
  \mathfrak g^{\alpha}\oplus
  \mathfrak g^{-\alpha},
\]
so \(Z_{\alpha}\ne T\).  Since \(T_{\alpha}\) is a central subtorus of
\(Z_{\alpha}\), of codimension \(1\) in \(T\), \SGARefLink{sga3:xix:lemma:1-11}{Lemma 1.11} applied to
\(Z_{\alpha}\) proves the equivalence of \SGARefLink{sga3:xix:proposition:1:12:item:ii}{(ii)} and (iii), together with
the additional assertions.

It remains to prove \SGARefLink{sga3:xix:proposition:1:12:item:ii}{(ii)}\(\Rightarrow\)\SGARefLink{sga3:xix:proposition:1:12:item:i}{(i)}.  We already know that \SGARefLink{sga3:xix:proposition:1:12:item:ii}{(ii)}
implies that \(T\) is maximal and that \(G\) is affine.  Let \(U\) be
the unipotent radical of \(G\).  It is invariant in \(G\), and its Lie
algebra is invariant under \(\operatorname{Ad}(T)\); hence
\[
  \operatorname{Lie}(U)
  =
  \bigl(\operatorname{Lie}(U)\cap\mathfrak g^{0}\bigr)
  \oplus
  \bigoplus_{\alpha\in R}
    \bigl(\operatorname{Lie}(U)\cap\mathfrak g^{\alpha}\bigr).
\]
By (\SGARefLink{sga3:xix:eq:1:3:1}{1.3.1}),
\[
\begin{aligned}
  U\cap T
    &=
    U\cap\underline{\operatorname{Centr}}_{G}(T)
     =
    \operatorname{rad}^{u}
      \bigl(\underline{\operatorname{Centr}}_{G}(T)\bigr)
     =
    \operatorname{rad}^{u}(T)
     =
    \{e\},\\
  U\cap Z_{\alpha}
    &=
    U\cap\underline{\operatorname{Centr}}_{G}(T_{\alpha})
     =
    \operatorname{rad}^{u}
      \bigl(\underline{\operatorname{Centr}}_{G}(T_{\alpha})\bigr)
     =
    \operatorname{rad}^{u}(Z_{\alpha})
     =
    \{e\}.
\end{aligned}
\]
It follows\footnote{Editor's note: if \(U\) and \(L\) are two subgroup
schemes of \(G\), then
\(\operatorname{Lie}(U)\cap\operatorname{Lie}(L)
=\operatorname{Lie}(U\cap L)\).} that
\[
\begin{aligned}
  \operatorname{Lie}(U)\cap\mathfrak g^{0}
    &=
    \operatorname{Lie}(U)\cap\mathfrak t
     =
    \operatorname{Lie}(U\cap T)
     =
    0,\\
  \operatorname{Lie}(U)\cap\mathfrak g^{\alpha}
    &\subset
    \operatorname{Lie}(U)\cap\operatorname{Lie}(Z_{\alpha})
     =
    \operatorname{Lie}(U\cap Z_{\alpha})
     =
    0.
\end{aligned}
\]
Thus \(\operatorname{Lie}(U)=0\), and therefore \(U=\{e\}\); in other
words, \(G\) is reductive.

\par\medskip\noindent
\SGAInlineRefTarget{sga3:xix:corollary:1-13}\textbf{\textit{\SGARefLink{sga3:xix:corollary:1:13}{Corollary 1.13}.}}---
\textit{Let \(k\), \(G\), \(T\), \(R\), and
\[
  \mathfrak g
  =
  \mathfrak g^{0}
  \oplus
  \bigoplus_{\alpha\in R}\mathfrak g^{\alpha}
\]
be as in \SGARefLink{sga3:xix:proposition:1-12}{Proposition 1.12}.  For every \(\alpha\in R\), let
\(T_{\alpha}\) be the maximal torus of \(\operatorname{Ker}(\alpha)\)
and put
\(Z_{\alpha}=\underline{\operatorname{Centr}}_{G}(T_{\alpha})\).
The following conditions are equivalent:
\begin{enumerate}
\item[\SGAInlineRefTarget{sga3:xix:proposition:1:12:item:i}(i)] \(G\) is affine and semisimple, and \(T\) is maximal.
\item[\SGAInlineRefTarget{sga3:xix:proposition:1:12:item:ii}(ii)] \(\mathfrak g^{0}=\mathfrak t\), every \(Z_{\alpha}\) is
  reductive, and
  \(\bigcap_{\alpha\in R}\operatorname{Ker}(\alpha)\) is finite.
\end{enumerate}
\footnote{Editor's note: this follows from the equalities
\(\operatorname{rad}(G)
=\underline{\operatorname{Centr}}(G)^{0}\) and
\(\underline{\operatorname{Centr}}(G)
=\bigcap_{\alpha\in R}\operatorname{Ker}(\alpha)\), established in
\SGARefLink{sga3:xix:lemma:1:6:2:item:iv}{1.6.2(iv)} and (\SGARefLink{sga3:xix:equation:1:10:3}{1.10.3}).}}
\SGARefTarget{sga3:xix:corollary:1:13}

% END INCLUDED SGA3_English_current_buildable_source_20260823/source/inputs/XIX/tex/components/06_expose_XIX_lemma1_11_completion_and_proposition1_12_pages007_008.tex

% BEGIN INCLUDED SGA3_English_current_buildable_source_20260823/source/inputs/XIX/tex/components/07_expose_XIX_section2_scholium2_1_and_2_2_pages008_009.tex
\SGASection{2}{Reductive group schemes: definitions and first properties}
\SGARefTarget{sga3:xix:section:2}

\par\medskip\noindent
\textbf{\textit{\SGARefLink{sga3:xix:scholium:2:1}{Scholium 2.1}.}}---
\textit{For an \(S\)-group scheme \(G\), the following properties are
equivalent:
\begin{enumerate}
\item[\SGAInlineRefTarget{sga3:xix:section:2:item:i}(i)] \(G\) is affine and smooth over \(S\), with connected
  fibers.
\item[\SGAInlineRefTarget{sga3:xix:section:2:item:ii}(ii)] \(G\) is affine and flat over \(S\), of finite presentation
  over \(S\), with geometrically integral fibers.
\end{enumerate}
These properties are stable under base change and are local for the
fpqc topology.}
\SGARefTarget{sga3:xix:scholium:2:1}

\footnote{Editor's note: the following details have been added.}
Indeed, suppose that \SGARefLink{sga3:xix:section:2:item:i}{(i)} holds.  Since \(G\) is affine and smooth over
\(S\), it is of finite presentation over \(S\); and since its fibers
are smooth and connected, they are geometrically integral, by
Exposé VI\textsubscript{A}, \SGARefLink{prop:sga3-VIA-2.4}{2.4}.

Conversely, if \SGARefLink{sga3:xix:section:2:item:ii}{(ii)} holds, the fibers of \(G\) are connected and
geometrically reduced, and hence smooth
(Exposé VI\textsubscript{A}, 1.3.1).  It follows that \(G\) is smooth
over \(S\), by EGA IV\textsubscript{4}, 17.5.2.

These properties are plainly stable under base change: see EGA II,
1.5.1 for ``affine'', EGA IV\textsubscript{1}, 1.6.2(iii) for ``of
finite presentation'', EGA IV\textsubscript{2}, 2.1.4 for ``flat'',
and EGA IV\textsubscript{2}, 4.6.5(i) for ``with geometrically integral
fibers''.

They are also clearly local for the Zariski topology.  It is therefore
enough to check that, if \(S'\longrightarrow S\) is faithfully flat and
quasi-compact and \(G_{S'}\longrightarrow S'\) has the indicated
properties, then so does \(G\longrightarrow S\).  This follows from
EGA IV\textsubscript{2}, 2.5.1 for ``flat'', 2.7.1(vi) and (xiii) for
``of finite presentation'' and ``affine'', and 4.6.5(i) for ``with
geometrically integral fibers''; for ``smooth'', see also
EGA IV\textsubscript{4}, 17.7.3(ii).

\par\medskip\noindent\SGAInlineRefTarget{sga3:xix:paragraph:2-2}\textbf{2.2.}\quad
Let \(G\) be an \(S\)-group scheme satisfying the preceding conditions,
and let \(Q\) be a torus of \(G\) in the sense of \SGARefLink{sga3:ix:definition:1:3}{Exposé IX,
Definition 1.3}.\footnote{Editor's note: the following argument has been
given in detail.}
By \SGARefLink{sga3:xi:corollary:6:11:item:a}{Exposé XI, 6.11(a)} and \SGARefLink{sga3:xix:remark:2-4}{2.4},
\(\underline{\operatorname{Centr}}_{G}(Q)\) is represented by a closed
subgroup scheme of \(G\); it is therefore affine over \(S\), and it is
of finite presentation and smooth over \(S\).  Moreover, every
geometric fiber \(G_{\bar s}\) is a smooth affine connected
\(\kappa(\bar s)\)-group.  The first assertion of \SGARefLink{sga3:xix:paragraph:1-3}{1.3} therefore shows
that the same is true of
\[
  \underline{\operatorname{Centr}}_{G_{\bar s}}(Q_{\bar s})
  =
  \bigl(\underline{\operatorname{Centr}}_{G}(Q)\bigr)_{\bar s}.
\]
\SGARefTarget{sga3:xix:section:2:2}

% END INCLUDED SGA3_English_current_buildable_source_20260823/source/inputs/XIX/tex/components/07_expose_XIX_section2_scholium2_1_and_2_2_pages008_009.tex

% BEGIN INCLUDED SGA3_English_current_buildable_source_20260823/source/inputs/XIX/tex/components/08_expose_XIX_lemma2_3_and_remark2_4_page009.tex
\par\medskip\noindent
\SGAInlineRefTarget{sga3:xix:lemma:2-3}\textbf{\textit{\SGARefLink{sga3:xix:lemma:2:3}{Lemma 2.3}.}}---
\textit{Let \(S\) be a scheme, let \(G\) be an \(S\)-group scheme that
is smooth and affine over \(S\), with connected fibers, and let \(T\)
be a torus of \(G\).  The set of points \(s\in S\) for which
\(G_{\bar s}\) is reductive, of semisimple rank \(1\), with maximal
torus \(T_{\bar s}\), is an open subset \(U\) of \(S\).}
\SGARefTarget{sga3:xix:lemma:2:3}

\SGAProof{Proof}
Since \(G\) and \(T\) are flat over \(S\), the function
\[
  s
  \longmapsto
  \dim(G_{\bar s}/\bar s)-\dim(T_{\bar s}/\bar s)
\]
is locally constant.  Let \(U_1\) be the open subset of \(S\) on which
this function has the value \(2\).

\footnote{Editor's note: the following details have been added.  Note
also that §6 is independent of the rest of this Exposé.}
By \SGARefLink{sga3:xix:proposition:6-3}{6.3} the Weyl group
\[
  W_G(T)
  =
  \underline{\operatorname{Norm}}_{G}(T)/
  \underline{\operatorname{Centr}}_{G}(T)
\]
is represented by an \(S\)-group scheme that is étale and separated
over \(S\), and the function
\[
  s\longmapsto
  \text{the number of points of \(W_G(T)_{\bar s}\)}
\]
is lower semicontinuous.  Let \(U_2\) be the locus on which this number
is greater than \(1\); it is open.  By Lemma 1.11, the locus on which
\(G_{\bar s}\) is reductive, of semisimple rank \(1\), with maximal
torus \(T_{\bar s}\), is
\[
  U=U_1\cap U_2.
\]
Furthermore, \(W_G(T)_{\bar s}\) has exactly two points for every
\(s\in U\).  Consequently, by SGA 1, I, 10.9 and
EGA IV\textsubscript{3}, 15.5.1 and IV\textsubscript{4}, 18.12.4,
\(W_G(T)_U\) is finite and étale over \(U\).

\par\medskip\noindent
\SGAInlineRefTarget{sga3:xix:remark:2-4}\textbf{\textit{\SGARefLink{sga3:xix:remark:2:4}{Remark 2.4}.}}---
\textit{The group \(W_G(T)_U\) is isomorphic to
\((\mathbb Z/2\mathbb Z)_U\).}
\SGARefTarget{sga3:xix:remark:2:4}

Indeed, the preceding argument shows that it is finite and étale over
\(U\).  Since the automorphism functor of
\((\mathbb Z/2\mathbb Z)_U\) is the trivial group, it is enough to
verify the assertion locally.  It is immediate whenever the algebra
\(\mathscr A\) defining \(W_G(T)_U\) is a free
\(\mathcal O_U\)-module.\footnote{Editor's note: the hypothesis implies
that \(\mathscr A\) is locally free of rank \(2\).  Since the
augmentation ideal \(\mathscr I\) is a direct summand of \(\mathscr A\),
after replacing \(U\) by a sufficiently small affine open
\(\operatorname{Spec}(R)\), one may suppose that
\(I=\Gamma(U,\mathscr I)\) is a free \(R\)-module of rank \(1\).
If \(x\) is a generator of \(I\), then \(x^2=ax\) for some \(a\in R\);
the assumption that \(A=\Gamma(U,\mathscr A)\) is étale implies that
\(a\) is invertible, which gives the asserted local identification.}

% END INCLUDED SGA3_English_current_buildable_source_20260823/source/inputs/XIX/tex/components/08_expose_XIX_lemma2_3_and_remark2_4_page009.tex

% BEGIN INCLUDED SGA3_English_current_buildable_source_20260823/source/inputs/XIX/tex/components/09_expose_XIX_notations2_5_0_and_theorem2_5_page010.tex
\par\medskip\noindent
\textbf{\textit{Notations and reminders 2.5.0.}}---%
\footnote{Editor's note: this paragraph of notations and reminders has
been added.}
Let \(S\) be a scheme, let \(G\) be an \(S\)-group scheme, and let
\(\epsilon\colon S\to G\) be the unit section of \(G\).  We saw in
II, §\SGARefLink{sga3:xix:paragraph:4-11}{4.11} that the functor
\(\underline{\operatorname{Lie}}(G/S)\) is represented by the vector
bundle
\(\mathbb V(\omega^1_{G/S})\), where
\(\omega^1_{G/S}=\epsilon^{*}(\Omega^1_{G/S})\), and we write
\[
  \mathfrak g
  =
  \operatorname{Lie}(G/S)
  =
  \bigl(\omega^1_{G/S}\bigr)^{\vee}
\]
for the sheaf of sections of this vector bundle.  Suppose in addition
that \(G\) is smooth over \(S\).  Then \(\omega^1_{G/S}\), and hence
\(\operatorname{Lie}(G/S)\), are finite locally free
\(\mathcal O_S\)-modules, and (cf. I, \SGARefLink{sga3:i:corollary:4-6-5-1}{4.6.5.1})
\[
  \underline{\operatorname{Lie}}(G/S)
  =
  \mathbb W\bigl(\operatorname{Lie}(G/S)\bigr);
\]
that is, for every \(S\)-scheme \(S'\),
\[
  \underline{\operatorname{Lie}}(G/S)(S')
  =
  \Gamma\!\left(
    S',
    \operatorname{Lie}(G/S)\otimes_{\mathcal O_S}\mathcal O_{S'}
  \right).
\]
By II, \SGARefLink{sga3:ii:corollary:4-1-1-1}{4.1.1.1}, the adjoint action of \(G\) gives
\(\underline{\operatorname{Lie}}(G/S)
=\mathbb W(\operatorname{Lie}(G/S))\) the structure of a
\(G\)-\(\mathcal O_S\)-module; thus
\(\operatorname{Lie}(G/S)\) is a \(G\)-\(\mathcal O_S\)-module
(cf. I, \SGARefLink{sga3:i:definition:4-7-1}{4.7.1}).  If \(G\) is moreover affine over \(S\), this says,
by I, \SGARefLink{sga3:i:proposition:4-7-2}{4.7.2}, that \(\operatorname{Lie}(G/S)\) is an
\(A(G)\)-comodule.

If \(T\) is a torus over \(S\) (IX, \SGARefLink{sga3:ix:definition:1:3}{Definition 1.3}), we say that \(T\)
is \emph{split} (``trivial'' in the original) if it is isomorphic to
\(\mathbb G_{m,S}^{\,r}\) for some integer \(r\geq 0\), and that \(T\)
is \emph{trivialized} if such an isomorphism has been fixed, or more
generally if an isomorphism
\[
  T\simeq D_S(M)
\]
has been fixed, where \(M\) is a free abelian group of rank \(r\).

Finally, recall (XII, \SGARefLink{sga3:xii:1.3}{Definition 1.3}) that a torus \(T\) of \(G\) is
called \emph{maximal} if, for every \(s\in S\), \(T_{\bar s}\) is a
maximal torus of \(G_{\bar s}\).
\SGARefTarget{sga3:xix:notations:2:5:0}

\par\medskip\noindent
\SGAInlineRefTarget{sga3:xix:theorem:2-5}\textbf{\textit{\SGARefLink{sga3:xix:theorem:2:5}{Theorem 2.5}.}}---
\textit{Let \(S\) be a scheme, let \(G\) be an \(S\)-group scheme that
is affine and of finite presentation over \(S\), with connected
fibers, and let \(s_0\) be a point of \(S\).  Suppose that \(G\) is
flat over \(S\) at \(\epsilon(s_0)\), and that the geometric fiber
\(G_{\bar s_0}\) is (reduced and) reductive
\textup{(resp. semisimple)}.  Then there exist an open subset
\(U\subset S\), containing \(s_0\), and an étale surjective morphism
\(S'\to U\) such that:}
\SGARefTarget{sga3:xix:theorem:2:5}

\begin{enumerate}
\item[\textit{\SGAInlineRefTarget{sga3:xix:theorem:2:5:item:i}(i)}]
\textit{\(G_U\) is smooth, with reductive
\textup{(resp. semisimple)} geometric fibers, and its reductive rank
and semisimple rank are constant.}

\item[\textit{\SGAInlineRefTarget{sga3:xix:theorem:2:5:item:ii}(ii)}]
\textit{\(G_{S'}\) has a split\footnote{Editor's note: the word
``split'' has been added.} maximal torus \(T\), and the Weyl group
\textup{(cf. \SGARefLink{sga3:xix:proposition:6-3}{6.3})}}
\[
  W_{G_{S'}}(T)
  =
  \underline{\operatorname{Norm}}_{G_{S'}}(T)/
  \underline{\operatorname{Centr}}_{G_{S'}}(T)
  =
  \underline{\operatorname{Norm}}_{G_{S'}}(T)/T
\]
\textit{is finite over \(S'\).}
\end{enumerate}


% END INCLUDED SGA3_English_current_buildable_source_20260823/source/inputs/XIX/tex/components/09_expose_XIX_notations2_5_0_and_theorem2_5_page010.tex

% BEGIN INCLUDED SGA3_English_current_buildable_source_20260823/source/inputs/XIX/tex/components/10_expose_XIX_theorem2_5_proof_pages010_012.tex
\SGAProof{Proof}\footnote{Editor's note: the proof has been expanded.}
Write \(\pi\colon G\to S\) for the structure morphism and
\(\epsilon\colon S\to G\) for the unit section.  Since \(G\) is flat
over \(S\) at \(\epsilon(s_0)\), and since \(G_{\bar s_0}\) is reduced
(hence smooth over \(\bar s_0\), by VI\textsubscript{A}, \SGARefLink{prop:sga3-VIA-1.3.1}{1.3.1}),
\(G\) is smooth over \(S\) at \(\epsilon(s_0)\)
(EGA IV\textsubscript{4}, 17.5.1).  Thus there is an open
neighborhood \(V\) of \(\epsilon(s_0)\) on which \(\pi\) is smooth.
Then \(S'=\epsilon^{-1}(V)\) is open in \(S\), and \(G_{S'}\) is
smooth over \(S'\) at the points of \(\epsilon(S')\).  Since \(G\)
has connected fibers, VI\textsubscript{B}, \SGARefLink{thm:sga3-VIB-3.10}{3.10} shows that \(G_{S'}\)
is smooth over \(S'\).  Replacing \(S\) by \(S'\), we may therefore
suppose that \(G\) is smooth over \(S\).

By \SGARefLink{sga3:xi:theorem:4:1}{Exposé XI, Theorem 4.1}, the functor of subgroups of multiplicative
type of \(G\) is represented by an \(S\)-scheme \(\mathscr M\), smooth
and separated over \(S\).  Let \(r_0\) be the rank of
\(G_{\bar s_0}\), and consider the open subscheme
\(\mathscr M_{r_0}\subset\mathscr M\) representing the functor of
subtori of rank \(r_0\) of \(G\) (IX, \SGARefLink{sga3:ix:definition:1:4}{1.4}).  Smoothness implies that
\(\mathscr M_{r_0}\) has a rational point over a finite separable
extension of \(\kappa(s_0)\) (EGA IV\textsubscript{4}, 17.15.10(iii)).
Consequently there is an étale morphism \(S'\to S\), together with a
point \(s'_0\) above \(s_0\), such that \(G_{s'_0}\) has a torus of
rank \(r_0\).  Replacing \(S\) by \(S'\), we may therefore suppose
that \(G_{s_0}\) has a torus of rank \(r_0\).

By ``Hensel's lemma'' (XI, \SGARefLink{sga3:xi:proposition:1:10}{1.10}), the smoothness of
\(\mathscr M_{r_0}\) allows this torus to be lifted to an \(S'\)-torus
\(T\) of \(G\), where \(S'\to S\) is étale and is equipped with a
point \(s'_0\) above \(s_0\).  By \SGARefLink{sga3:x:corollary:4-5}{Exposé X, 4.5}
(see also \SGARefLink{sga3:xix:proposition:6-1}{6.1}\footnote{Editor's note: §\SGARefLink{sga3:xix:proposition:6-1}{6.1} is independent of the rest
of this Exposé.}), there is an étale morphism \(f\colon S''\to S'\)
that splits \(T\) and for which \(f^{-1}(s'_0)\ne\varnothing\).
Because an étale morphism is open, and the assertions in \textup{\SGARefLink{sga3:xix:theorem:2:5:item:i}{(i)}}
are local for the étale topology, we may thus suppose that \(G\) has
a split torus \(T\),\footnote{Editor's note: throughout this volume,
``trivial torus'' has been replaced by ``split torus''.} maximal at
\(s_0\).  Write \(T=D_S(M)\), and let
\[
  \mathfrak g=\prod_{m\in M}\mathfrak g^{m}
\]
be the decomposition of
\(\mathfrak g=\operatorname{Lie}(G/S)\) under
\(\operatorname{Ad}(T)\) (\SGARefLink{sga3:i:proposition:4-7-3}{Exposé I, 4.7.3}).

Put \(\mathfrak t=\operatorname{Lie}(T)\), and, for every \(m\in M\),
write
\[
  \mathfrak g^{m}(s_0)
  =
  \mathfrak g^{m}\otimes_{\mathcal O_S}\kappa(s_0).
\]
\footnote{Editor's note: the preceding sentence has been added.}
Let \(R\) be the set of \(m\in M\), \(m\ne0\), such that
\(\mathfrak g^{m}(s_0)\ne0\).\footnote{Editor's note: since
\(\mathfrak g\) is a finite locally free \(\mathcal O_S\)-module,
\(R\) is finite.}  Since \(G_{\bar s_0}\) is reductive,
\(\mathfrak g^{0}(s_0)=\mathfrak t(s_0)\), and hence
\[
  \mathfrak g(s_0)
  =
  \mathfrak t(s_0)
  \oplus
  \prod_{\alpha\in R}\mathfrak g^{\alpha}(s_0).
\]
Since the modules in question are locally free, after shrinking \(S\)
we may suppose that the \(\mathfrak g^{\alpha}\) are free and that
\[
  \mathfrak g
  =
  \mathfrak t
  \oplus
  \prod_{\alpha\in R}\mathfrak g^{\alpha}.
\]

Recall from \SGARefLink{sga3:xii:1.12}{Exposé XII, 1.12} that a group of multiplicative type has a
unique maximal torus (this also follows almost immediately by descent,
the diagonalizable case being evident).  Let \(T_\alpha\) be the
maximal torus of \(\ker(\alpha)\), and put
\[
  Z_\alpha=\underline{\operatorname{Centr}}_G(T_\alpha).
\]
\footnote{Editor's note: the following argument has been expanded.}
By \SGARefLink{sga3:xix:paragraph:2-2}{2.2}, \(Z_\alpha\) is affine and smooth over \(S\), with connected
fibers; and by \SGARefLink{sga3:xix:proposition:1-12}{1.12}, its fiber
\((Z_\alpha)_{\bar s_0}\) is reductive, of semisimple rank \(1\), with
maximal torus \(T_{\bar s_0}\).  By \SGARefLink{sga3:xix:lemma:2-3}{2.3} there is therefore an open
subset \(U_\alpha\subset S\), containing \(s_0\), on which
\(Z_\alpha\) has reductive fibers.

Set
\[
  U=\bigcap_{\alpha\in R}U_\alpha.
\]
By \SGARefLink{sga3:xix:proposition:1-12}{1.12}, for every \(s\in U\), \(G_{\bar s}\) is reductive, with
maximal torus \(T_{\bar s}\), and the set of roots of \(G_{\bar s}\)
relative to \(T_{\bar s}\) is identified with \(R\).  Consequently
\[
  \operatorname{rgred}(G_{\bar s})
  =
  \dim(T)
  =
  \operatorname{rg}(M),
  \qquad
  \operatorname{rgss}(G_{\bar s})
  =
  \operatorname{rg}(R)
  \quad\text{(cf. 1.10).}
\]
This proves \textup{\SGARefLink{sga3:xix:theorem:2:5:item:i}{(i)}} and the first assertion of \textup{\SGARefLink{sga3:xix:theorem:2:5:item:ii}{(ii)}}.
It remains to prove that \(W_{G_U}(T_U)\) is finite over \(U\), that
is, that it has the same number of points in every geometric fiber
(SGA 1, I, 10.9; EGA IV\textsubscript{3}, 15.5.1; and
EGA IV\textsubscript{4}, 8.12.4).

For this it is enough to observe that, for \(s\in U\), the geometric
fiber of this group is determined by the configuration \(R\subset M\):
it is the constant group associated with the abstract Weyl group of
this root system.  In particular it is independent of \(s\); see the
\emph{Bible}, §11.3, Theorem 2 (and also Exposé XXII, no. 3).
\hfill\(\square\)


% END INCLUDED SGA3_English_current_buildable_source_20260823/source/inputs/XIX/tex/components/10_expose_XIX_theorem2_5_proof_pages010_012.tex

% BEGIN INCLUDED SGA3_English_current_buildable_source_20260823/source/inputs/XIX/tex/components/11_expose_XIX_corollary2_6_through_remark2_9_page012.tex
\par\medskip\noindent
\textbf{\textit{\SGARefLink{sga3:xix:corollary:2:6}{Corollary 2.6}.}}---
\textit{Let \(G\) be an \(S\)-group scheme that is affine and smooth
over \(S\), with connected fibers.  The set of points \(s\in S\) for
which \(G_{\bar s}\) is reductive \textup{(resp. semisimple)} is an
open subset \(U\) of \(S\), and the functions}
\[
  s\longmapsto\operatorname{rgred}(G_{\bar s}/\bar s),
  \qquad
  s\longmapsto\operatorname{rgss}(G_{\bar s}/\bar s)
\]
\textit{are locally constant on \(U\).}
\SGARefTarget{sga3:xix:corollary:2:6}

\par\medskip\noindent
\textbf{\textit{\SGARefLink{sga3:xix:definition:2:7}{Definition 2.7}.}}---
\textit{An \(S\)-group \textup{(that is, an \(S\)-group scheme)}
\(G\) is called \emph{reductive} \textup{(resp. \emph{semisimple})}
if it is affine and smooth over \(S\), with connected reductive
\textup{(resp. semisimple)} geometric fibers.}
\SGARefTarget{sga3:xix:definition:2:7}

For an \(S\)-group scheme \(G\), the property of being reductive
(resp. semisimple) is stable under base change and local for the fpqc
topology.

\par\medskip\noindent
\textbf{2.8.}---
Let \(G\) be a reductive \(S\)-group.  For every torus
(resp. maximal torus) \(Q\) of \(G\),
\[
  Z(Q)=\underline{\operatorname{Centr}}_G(Q)
\]
is reductive (resp. is equal to \(Q\)).  This follows at once from \SGARefLink{sga3:xix:paragraph:2-2}{2.2}
and 1.6.  Applying \SGARefLink{sga3:xix:theorem:2-5}{2.5} to
\(\underline{\operatorname{Centr}}_G(Q)\), we deduce that \(Q\) is
contained, locally for the étale topology, in a maximal torus.
\SGARefTarget{sga3:xix:paragraph:2:8}

\par\medskip\noindent
\textbf{\textit{Remark 2.9.}}---
\textit{Using, as usual, the technique of EGA IV\textsubscript{3},
§8, the finite-presentation hypotheses, and Theorem 2.5, one sees
that if \(G\) is a reductive group over \(S\), there is an open cover
\(\{U_i\}\) of \(S\) such that each \(G_{U_i}\) is obtained by base
change from a reductive group over a noetherian ring \textup{(in fact,
over a finitely generated \(\mathbb Z\)-algebra)}.  Likewise, if \(G\)
has a split maximal torus \(T\), one may further suppose that
\(T_{U_i}\) comes from a split maximal torus of the preceding
group, \ldots}
\SGARefTarget{sga3:xix:remark:2:9}

\SGASection{3}{Roots and root systems of reductive group schemes}
\SGARefTarget{sga3:xix:section:3}


% END INCLUDED SGA3_English_current_buildable_source_20260823/source/inputs/XIX/tex/components/11_expose_XIX_corollary2_6_through_remark2_9_page012.tex

% BEGIN INCLUDED SGA3_English_current_buildable_source_20260823/source/inputs/XIX/tex/components/12_expose_XIX_section3_1_through_3_4_page013.tex
\par\medskip\noindent
\textbf{\SGAInlineRefTarget{sga3:xix:paragraph:3-1}3.1.}---
Let \(S\) be a scheme, and let \(T\) be an \(S\)-torus acting linearly
on a finite locally free \(\mathcal O_S\)-module \(F\) (cf. I, §\SGARefLink{sga3:xix:lemma:4-7}{4.7}).
For every character \(\alpha\) of \(T\), that is,
\[
  \alpha\in
  \operatorname{Hom}_{S\text{-gr.}}(T,\mathbb G_{m,S}),
\]
define a subfunctor of \(\mathbb W(F)\) by
\[
  \mathbb W(F)^\alpha(S')
  =
  \left\{
    x\in\mathbb W(F)(S')
    \ \middle|\
    t\cdot x=\alpha(t)x
    \text{ for every }t\in T(S''),\ S''\to S'
  \right\}.
\]
\SGARefTarget{sga3:xix:paragraph:3:1}

\par\medskip\noindent
\textbf{\textit{Lemma.}}---%
\footnote{Editor's note: this statement has been made into a lemma in
order to bring it out.}
\textit{Then
\(\mathbb W(F)^\alpha=\mathbb W(F^\alpha)\), where \(F^\alpha\) is a
submodule of \(F\), locally a direct summand of \(F\), and hence also
locally free.}
\SGARefTarget{sga3:xix:lemma:3:1}

Indeed, the assertion is local for the fpqc topology, and we may
suppose that \(T=D_S(M)\), where \(M\) is a finite-type (free) abelian
group.  Then \(\alpha\) is identified with a locally constant function
from \(S\) to \(M\) (\SGARefLink{sga3:viii:corollary:1:3}{Exposé VIII, 1.3}), and, after restricting \(S\) if
necessary, we may suppose that this function is constant.  We are then
reduced to \SGARefLink{sga3:i:proposition:4-7-3}{Exposé I, 4.7.3}.

\par\medskip\noindent
\SGAInlineRefTarget{sga3:xix:definition:3-2}\textbf{\textit{\SGARefLink{sga3:xix:definition:3:2}{Definition 3.2}.}}---
\textit{Let \(S\) be a scheme, let \(G\) be an \(S\)-group scheme
smooth over \(S\) and with connected fibers,\footnote{Editor's note:
the original also assumed that \(G\) was of finite presentation over
\(S\), an assumption that does not appear to be used below.  In any
case, since \(G\) is smooth over \(S\) and has connected fibers, it is
quasi-compact and separated over \(S\) (VI\textsubscript{B}, \SGARefLink{cor:sga3-VIB-5.5}{5.5}), and
hence of finite presentation over \(S\).} and let \(T\) be a torus of
\(G\).  Put \(\mathfrak g=\operatorname{Lie}(G/S)\), and let \(T\)
act on \(\mathfrak g\) through the adjoint representation of \(G\).}

\textit{A character \(\alpha\) of \(T\) is called a \emph{root} of
\(G\) with respect to \(T\) if the following equivalent conditions
are satisfied:}
\SGARefTarget{sga3:xix:definition:3:2}

\begin{enumerate}
\item[\textit{(i)}]
\textit{For every \(s\in S\), \(\alpha_s\) is a root of \(G_s\) with
respect to \(T_s\) (\SGARefLink{sga3:xix:paragraph:1-10}{1.10}).}

\item[\textit{(ii)}]
\textit{\(\alpha\) is nontrivial on every fiber, and
\(\mathfrak g^\alpha(s)\ne0\) for every \(s\in S\).}
\end{enumerate}

\par\medskip\noindent
\textbf{\textit{\SGARefLink{sga3:xix:lemma:3:3}{Lemma 3.3}.}}---
\textit{Let \(S\) be a scheme, let \(T\) be an \(S\)-torus, and let
\(\alpha\) be a character of \(T\).  The following conditions are
equivalent:}
\SGARefTarget{sga3:xix:lemma:3:3}

\begin{enumerate}
\item[\textit{\SGAInlineRefTarget{sga3:xix:lemma:3:3:item:i}(i)}]
\textit{\(\alpha\) is nontrivial on every fiber; that is, for every
\(s\in S\), \(\alpha_s\) is distinct from the unit character of
\(T_s\).}

\item[\textit{\SGAInlineRefTarget{sga3:xix:lemma:3:3:item:ii}(ii)}]
\textit{For every \(S'\to S\) with \(S'\ne\varnothing\),
\(\alpha_{S'}\) is distinct from the unit character of \(T_{S'}\).}

\item[\textit{\SGAInlineRefTarget{sga3:xix:lemma:3:3:item:iii}(iii)}]
\textit{The morphism \(\alpha\) is faithfully flat.}
\end{enumerate}

\footnote{Editor's note: the original has been simplified; it invoked
\SGARefLink{sga3:ix:corollary:5:2}{Exposé IX, 5.2}.  Note, however, that loc.\ cit., \SGARefLink{sga3:xix:lemma:5-3}{5.3} shows that, if
\(S\) is connected, the conditions of the lemma hold as soon as
\textup{\SGARefLink{sga3:xix:lemma:3:3:item:i}{(i)}} holds at one point \(s\) of \(S\).}
It is clear that \SGARefLink{sga3:xix:lemma:3:3:item:ii}{(ii)} implies \SGARefLink{sga3:xix:lemma:3:3:item:i}{(i)}, and it is easy to see that \SGARefLink{sga3:xix:lemma:3:3:item:iii}{(iii)}
implies \SGARefLink{sga3:xix:lemma:3:3:item:i}{(i)}.  Condition \SGARefLink{sga3:xix:lemma:3:3:item:i}{(i)} implies \SGARefLink{sga3:xix:lemma:3:3:item:ii}{(ii)}: if \(s'\in S'\) lies over
\(s\), and \(\alpha_{s'}\) is the unit character, then so is
\(\alpha_s\).  Finally, to prove that \SGARefLink{sga3:xix:lemma:3:3:item:i}{(i)} implies \SGARefLink{sga3:xix:lemma:3:3:item:iii}{(iii)}, we reduce to
the case in which \(T\) is diagonalizable and conclude by \SGARefLink{sga3:viii:corollary:3:2:item:a:01}{Exposé VIII,
3.2(a)}.

\par\medskip\noindent
\textbf{3.4.}---%
\footnote{Editor's note: what follows has been modified so as to recall
the hypotheses of \SGARefLink{sga3:xix:definition:3-2}{3.2} and to add that \(T\) is a maximal torus.}
Let \(G\) be a reductive \(S\)-group scheme, let \(T\) be a maximal
torus of \(G\), and let \(\alpha\) be a root of \(G\) with respect to
\(T\).  Then, by \SGARefLink{sga3:xix:notations:2:5:0}{2.5.0} and \SGARefLink{sga3:xix:proposition:1-12}{1.12}, \(\mathfrak g^\alpha\) is a locally
free \(\mathcal O_S\)-module of rank one.  Moreover, by \SGARefLink{sga3:xix:paragraph:1-10}{1.10},
\(-\alpha\) is also a root of \(G\) with respect to \(T\).  In
particular, if \(G\) has semisimple rank one, \SGARefLink{sga3:xix:lemma:1-11}{1.11} and \SGARefLink{sga3:xix:proposition:1-12}{1.12} give the
following statement.
\SGARefTarget{sga3:xix:paragraph:3:4}

% END INCLUDED SGA3_English_current_buildable_source_20260823/source/inputs/XIX/tex/components/12_expose_XIX_section3_1_through_3_4_page013.tex

% BEGIN INCLUDED SGA3_English_current_buildable_source_20260823/source/inputs/XIX/tex/components/13_expose_XIX_lemma3_5_through_root_functor_page014.tex
\par\medskip\noindent
\textbf{\textit{\SGARefLink{sga3:xix:lemma:3:5}{Lemma 3.5}.}}---
\textit{Let \(S\) be a scheme, let \(G\) be a reductive \(S\)-group
of semisimple rank one, and let \(T\) be a maximal torus of \(G\).  If
\(\alpha\) is a root of \(G\) with respect to \(T\), then so is
\(-\alpha\), and}
\[
  \mathfrak g
  =
  \mathfrak t
  \oplus \mathfrak g^\alpha
  \oplus \mathfrak g^{-\alpha},
\]
\textit{where \(\mathfrak g^\alpha\) and
\(\mathfrak g^{-\alpha}\) are locally free of rank one.}
\SGARefTarget{sga3:xix:lemma:3:5}

\par\medskip\noindent
\SGAInlineRefTarget{sga3:xix:definition:3-6}\textbf{\textit{\SGARefLink{sga3:xix:definition:3:6}{Definition 3.6}.}}---
\textit{Let \(S\) be a scheme, let \(G\) be a reductive \(S\)-group,
let \(T\) be a maximal torus of \(G\), and let \(R\) be a set of roots
of \(G\) with respect to \(T\).  We say that \(R\) is a \emph{root
system} of \(G\) with respect to \(T\) if the following equivalent
conditions are satisfied:}
\SGARefTarget{sga3:xix:definition:3:6}

\begin{enumerate}
\item[\textit{(i)}]
\textit{For every \(s\in S\), the map
\(\alpha\mapsto\alpha_s\) is a bijection from \(R\) onto the set of
roots of \(G_s\) with respect to \(T_s\).}

\item[\textit{\SGAInlineRefTarget{sga3:xix:definition:3:6:item:ii}(ii)}]
\textit{The elements of \(R\) are distinct on every fiber
\textup{(that is, if \(\alpha,\alpha'\in R\) and
\(\alpha\ne\alpha'\), then
\(\alpha-\alpha'=\alpha{\alpha'}^{-1}\) is nontrivial on every
fiber)}, and, for every \(s\in S\),}
\[
  \dim(G_s/\kappa(s))-\dim(T_s/\kappa(s))=\operatorname{Card}(R).
\]

\item[\textit{\SGAInlineRefTarget{sga3:xix:definition:3:6:item:iii}(iii)}]
\textit{One has}
\[
  \mathfrak g
  =
  \mathfrak t
  \oplus
  \bigoplus_{\alpha\in R}\mathfrak g^\alpha.
\]
\end{enumerate}

The equivalence of these conditions is immediate.

\par\medskip\noindent
\SGAInlineRefTarget{sga3:xix:lemma:3-7}\textbf{\textit{\SGARefLink{sga3:xix:lemma:3:7}{Lemma 3.7}.}}---
\textit{Let \(S\) be a scheme, let \(G\) be a reductive \(S\)-group,
let \(T\) be a maximal torus of \(G\), and let \(R\) be a root system
of \(G\) with respect to \(T\).  Every root of \(G\) with respect to
\(T\) is locally on \(S\) equal to an element of \(R\).}
\SGARefTarget{sga3:xix:lemma:3:7}

\SGAProof{Proof}
This follows from condition \SGARefLink{sga3:xix:definition:3:6:item:iii}{(iii)}.

Put
\[
  \mathcal M
  =
  \underline{\operatorname{Hom}}_{S\text{-gr.}}
  (T,\mathbb G_{m,S});
\]
this is a twisted constant \(S\)-group scheme (\SGARefLink{sga3:x:theorem:5-6}{Exposé X, 5.6}).  If
\(G\) has a root system \(R\) with respect to \(T\), then the inclusion
\(R\hookrightarrow\mathcal M(S)\) defines a morphism
\[
  R_S\longrightarrow\mathcal M,
\]
where \(R_S\) is the constant \(S\)-scheme defined by \(R\).  By
condition \SGARefLink{sga3:xix:definition:3:6:item:ii}{(ii)}, this morphism is an open and closed immersion whose
image is precisely
\[
  \bigcup_{\alpha\in R}\alpha(S),
\]
where each \(\alpha\in R\) is regarded as a section
\(S\to\mathcal M\).

Let \(\mathcal R\) be the functor of roots of \(G\) with respect to
\(T\).  By definition, \(\mathcal R(S')\) is the set of roots of
\(G_{S'}\) with respect to \(T_{S'}\), for every \(S'\to S\).  If
\(S'=\varnothing\), set \(\mathcal R(\varnothing)=\{e\}\); if
\(S'\ne\varnothing\), the inclusion
\(R\hookrightarrow\mathcal M(S')\) identifies \(R\) with a root
system of \(G_{S'}\) with respect to \(T_{S'}\).  Thus, by \SGARefLink{sga3:xix:lemma:3-7}{3.7},
\[
  \mathcal R(S')
  =
  \operatorname{Hom}_{\mathrm{loc.\ const.}}(S',R),
\]
which shows that \(\mathcal R\) is represented by \(R_S\).

Now drop the assumption that \(G\) necessarily has a root system with
respect to \(T\).  The functor \(\mathcal R\) is in any case an fpqc
subsheaf of \(\mathcal M\).  Locally for this topology, \(G\) does
have a root system with respect to \(T\) (for example, make \(T\)
split and use the proof of \SGARefLink{sga3:xix:theorem:2:5}{Theorem 2.5}).  By \SGARefLink{sga3:iv:proposition:4:6:8}{Exposé IV, 4.6.8} and the
descent theory of open (respectively closed)
subschemes,\footnote{Editor's note: see SGA 1, VIII, 4.4 or
EGA IV\textsubscript{2}, 2.7.1.} we obtain the following proposition.

% END INCLUDED SGA3_English_current_buildable_source_20260823/source/inputs/XIX/tex/components/13_expose_XIX_lemma3_5_through_root_functor_page014.tex

% BEGIN INCLUDED SGA3_English_current_buildable_source_20260823/source/inputs/XIX/tex/components/14_expose_XIX_proposition3_8_through_3_10_page015.tex
\par\medskip\noindent
\SGAInlineRefTarget{sga3:xix:proposition:3-8}\textbf{\textit{\SGARefLink{sga3:xix:proposition:3:8}{Proposition 3.8}.}}---
\textit{Let \(S\) be a scheme, let \(G\) be an \(S\)-group, and let
\(T\) be a maximal torus of \(G\).  The functor \(\mathcal R\) of
roots of \(G\) with respect to \(T\) is represented by a finite
twisted constant \(S\)-scheme \textup{(\SGARefLink{sga3:x:definition:5-1}{Exposé X, 5.1})} which is an
open and closed subscheme of}
\[
  \underline{\operatorname{Hom}}_{S\text{-gr.}}
  (T,\mathbb G_{m,S}).
\]

\textit{A subset
\(R\subset
\operatorname{Hom}_{S\text{-gr.}}(T,\mathbb G_{m,S})\)
is a root system of \(G\) with respect to \(T\) if and only if the
morphism}
\[
  R_S
  \longrightarrow
  \underline{\operatorname{Hom}}_{S\text{-gr.}}
  (T,\mathbb G_{m,S}),
\]
\textit{defined by the preceding inclusion, induces an isomorphism}
\[
  R_S\xrightarrow{\ \sim\ }\mathcal R.
\]
\SGARefTarget{sga3:xix:proposition:3:8}

\par\medskip\noindent
\textbf{\SGAInlineRefTarget{sga3:xix:paragraph:3-9}3.9.}---
Let \(S\) be a scheme, let \(G\) be a reductive \(S\)-group, let \(T\)
be a maximal torus of \(G\), and let \(\alpha\) be a root of \(G\)
with respect to \(T\), that is, a section of \(\mathcal R\).  Consider
the kernel \(\operatorname{Ker}(\alpha)\) of \(\alpha\), its unique
maximal torus \(T_\alpha\), and the centralizer of the latter,
\[
  Z_\alpha=\underline{\operatorname{Centr}}_G(T_\alpha).
\]
This is a closed \(S\)-subgroup of \(G\), reductive (\SGARefLink{sga3:xix:paragraph:2:8}{2.8}), and of
semisimple rank one (\SGARefLink{sga3:xix:proposition:1-12}{1.12}).  Moreover,
\[
  \operatorname{Lie}(Z_\alpha/S)
  =
  \mathfrak t
  \oplus\mathfrak g^\alpha
  \oplus\mathfrak g^{-\alpha};
\]
hence \(\{\alpha,-\alpha\}\) is a root system of \(Z_\alpha\) with
respect to \(T\).
\SGARefTarget{sga3:xix:paragraph:3:9}

\par\medskip\noindent
\textbf{3.10.}---
Let \(S\) be a scheme, let \(G\) be a reductive \(S\)-group, let \(T\)
be a maximal torus of \(G\), and let \(\alpha\) be a root of \(G\)
with respect to \(T\).  If \(q\in\mathbb Q\) and \(q\alpha\) is a root
of \(G\) with respect to \(T\), then \(q=1\) or \(q=-1\).  This follows
at once from \SGARefLink{sga3:xix:proposition:1-12}{1.12}.
\SGARefTarget{sga3:xix:paragraph:3:10}

% END INCLUDED SGA3_English_current_buildable_source_20260823/source/inputs/XIX/tex/components/14_expose_XIX_proposition3_8_through_3_10_page015.tex

% BEGIN INCLUDED SGA3_English_current_buildable_source_20260823/source/inputs/XIX/tex/components/15_expose_XIX_section4_1_and_lemma4_2_pages015_016.tex
\SGASection{4}{Roots and vector group schemes}
\SGARefTarget{sga3:xix:section:4}

\par\medskip\noindent
\textbf{\SGAInlineRefTarget{sga3:xix:paragraph:4-1}4.1.}---
Let \(S\) be a scheme and let \(F\) be a finite locally free
\(\mathcal O_S\)-module.  The \(S\)-scheme \(\mathbb W(F)\) is smooth
over \(S\).  Its Lie algebra is canonically isomorphic to \(F\).
Indeed, there is a canonical isomorphism
\[
  \mathbb W(F)
  \xrightarrow{\ \sim\ }
  \operatorname{Lie}(\mathbb W(F)/S)
  =
  \mathbb W\bigl(\operatorname{Lie}(\mathbb W(F)/S)\bigr)
\]
\textup{(\SGARefLink{sga3:ii:corollary:4-4-1}{Exposé II, 4.4.1} and \SGARefLink{sga3:xix:scholium:4-4-2}{4.4.2})}.  We shall always identify
\(F\) with \(\operatorname{Lie}(\mathbb W(F)/S)\).
\SGARefTarget{sga3:xix:paragraph:4:1}

\par\medskip\noindent
\textbf{\textit{\SGARefLink{sga3:xix:lemma:4:2}{Lemma 4.2}.}}---
\textit{Let \(S\) be a scheme and let \(V\) be a vector bundle over
\(S\), smooth over \(S\).  There is then a unique isomorphism of
\(\mathcal O_S\)-modules}
\[
  \exp\colon
  \mathbb W\bigl(\operatorname{Lie}(V/S)\bigr)
  \xrightarrow{\ \sim\ } V
\]
\textit{which induces the identity on Lie algebras.}
\footnote{Editor's note: the proof in the original edition has been
retained; it may also be expanded as follows.

Let \(\mathscr F\) be a quasi-coherent \(\mathcal O_S\)-module and put
\(V=\mathbb V(\mathscr F)\).  Write \(\pi\colon V\to S\) for the
projection and \(\epsilon\colon S\to V\) for the zero section.  Then
\(\Omega^1_{V/S}=\pi^*\mathscr F\), whence
\[
  \tag{1}\label{sga3:xix:paragraph:4-1:equation:1}
  \omega^1_{V/S}
  =
  \epsilon^*\Omega^1_{V/S}
  \simeq
  \epsilon^*\pi^*\mathscr F
  \simeq
  \mathscr F,
\]
and therefore
\(\operatorname{Lie}(V/S)=(\omega^1_{V/S})^\vee
\simeq\mathscr F^\vee\).
If \(V\) is smooth over \(S\), then by \textup{(1)}
\(\mathscr F\) is finite locally free, and hence
\[
  \tag{2}\label{sga3:xix:paragraph:4-1:equation:2}
  V=\mathbb V(\mathscr F)
  \simeq
  \mathbb W(\mathscr F^\vee)
  \simeq
  \mathbb W\bigl(\operatorname{Lie}(V/S)\bigr).
\]
Now let \(U\) be an \(S\)-scheme equipped with a section
\(\tau\colon S\to U\), and suppose that an isomorphism
\(\phi\colon V\xrightarrow{\sim}U\) of pointed \(S\)-schemes has been
given, that is, \(\phi\circ\epsilon=\tau\) \textup{(for example, when
\(U\) is an \(S\)-group)}.  Then \(\phi\) induces an isomorphism
\[
  d\phi\colon
  \mathscr F^\vee
  \xrightarrow{\ \sim\ }
  \operatorname{Lie}(U/S),
\]
and therefore an isomorphism \(\psi\) displayed by
\[
\begin{tikzcd}[
  ampersand replacement=\&,
  column sep=huge,
  row sep=large
]
  \mathbb W(\mathscr F^\vee)
    \arrow[r,equal]
  \&
  \mathbb V(\mathscr F)
    \arrow[r,"\phi","\sim"']
  \&
  U
  \\
  \&
  \mathbb W\bigl(\operatorname{Lie}(U/S)\bigr)
    \arrow[ul,"\sim","\mathbb W(d\phi^{-1})"']
    \arrow[ur,"\sim","\psi"']
  \&
\end{tikzcd}
\]
which identifies \(U\) with
\(\mathbb W(\operatorname{Lie}(U/S))\).  Since the functor
\(\mathbb W\) is fully faithful, there is a unique isomorphism
\[
  \exp\colon
  \mathbb W\bigl(\operatorname{Lie}(U/S)\bigr)
  \xrightarrow{\ \sim\ }U
\]
such that
\(\operatorname{Lie}(\psi^{-1}\circ\exp)=\operatorname{id}\).
In fact, one sees directly that
\(\operatorname{Lie}(\psi)=\operatorname{id}\), so that
\(\exp=\psi\).}
\SGARefTarget{sga3:xix:lemma:4:2}

\SGAProof{Proof}
We have \(V=\mathbb V(\mathscr F)\) for some quasi-coherent
\(\mathcal O_S\)-module \(\mathscr F\).  Since \(V\) is smooth over
\(S\), the module
\(\mathscr F\simeq\omega^1_{V/S}\) is finite locally free, and hence
\[
  \operatorname{Lie}(V/S)
  =
  \mathbb V(\omega^1_{V/S})
  \simeq
  \mathbb W\bigl(\operatorname{Lie}(V/S)\bigr).
\]
Moreover, by Exposé II, loc.\ cit., there is a canonical isomorphism
\[
  V
  \xrightarrow{\ \sim\ }
  \operatorname{Lie}(V/S)
  \simeq
  \mathbb W\bigl(\operatorname{Lie}(V/S)\bigr),
\]
and the uniqueness of \(\exp\) follows at once because
\(\mathbb W\) is fully faithful.
\hfill\(\square\)

% END INCLUDED SGA3_English_current_buildable_source_20260823/source/inputs/XIX/tex/components/15_expose_XIX_section4_1_and_lemma4_2_pages015_016.tex

% BEGIN INCLUDED SGA3_English_current_buildable_source_20260823/source/inputs/XIX/tex/components/16_expose_XIX_notations4_3_through_scholium4_4_2_page016.tex
\par\medskip\noindent
\textbf{\SGAInlineRefTarget{sga3:xix:paragraph:4-3}4.3. Notation.}---
If \(V\) is a vector bundle over \(S\), let \(V^\times\) denote the
open subset obtained from \(V\) by removing the zero section.  Write
the group law of \(V\) multiplicatively.  The action of
\(\mathcal O_S\) on \(V\) which defines its module structure will
then be written exponentially:
\[
  \mathcal O_S\times_S V\longrightarrow V,
  \qquad
  (x,v)\longmapsto v^x.
\]
Thus
\[
  (vv')^x=v^x(v')^x,
  \qquad
  v^{x+x'}=v^xv^{x'},
  \qquad
  v^{xx'}=(v^x)^{x'}.
\]
In particular, upon restricting the operator action to
\(\mathbb G_{m,S}\), the open subset \(V^\times\) is stable and is
therefore an object with operator group \(\mathbb G_{m,S}\).  We
shall continue to write this action as
\((z,v)\longmapsto v^z\).
\SGARefTarget{sga3:xix:paragraph:4:3}

\par\medskip\noindent
\SGAInlineRefTarget{sga3:xix:definition:4-4-1}\textbf{\textit{\SGARefLink{sga3:xix:definition:4:4:1}{Definition 4.4.1}.}}\footnote{Editor's note: for
purposes of later references, paragraph 4.4 has been divided into
\SGARefLink{sga3:xix:definition:4-4-1}{4.4.1} and \SGARefLink{sga3:xix:scholium:4-4-2}{4.4.2}.}---
\textit{Let \(L\) be an invertible module on \(S\), and let
\(\mathbb W(L)\) be the associated vector bundle.  Then
\(\mathbb W(L)^\times\) is a \textup{(locally trivial)} principal
homogeneous bundle under \(\mathbb G_{m,S}\).  We write}
\[
  \Gamma(S,L)^\times=\mathbb W(L)^\times(S).
\]
\SGARefTarget{sga3:xix:definition:4:4:1}

\par\medskip\noindent
\SGAInlineRefTarget{sga3:xix:scholium:4-4-2}\textbf{\textit{\SGARefLink{sga3:xix:scholium:4:4:2}{Scholium 4.4.2}.}}---
\textit{There are natural bijections among the isomorphisms of
\(\mathcal O_S\)-modules
\(\mathcal O_S\simeq L\), the isomorphisms of
\(\mathcal O_S\)-modules
\(\mathcal O_S\simeq\mathbb W(L)\),\footnote{Editor's note: here the
vector bundle \(\mathbb W(L)\) is identified with the
\(\mathcal O_S\)-module functor which it represents.} and the
sections \(S\to\mathbb W(L)^\times\).}

The correspondence is
\[
  f\longmapsto\mathbb W(f)\longmapsto\mathbb W(f)(1).
\]
It is compatible with base change.  Thus
\(\mathbb W(L)^\times\) may be regarded as the ``scheme of
trivializations of \(\mathbb W(L)\)''.
\SGARefTarget{sga3:xix:scholium:4:4:2}


% END INCLUDED SGA3_English_current_buildable_source_20260823/source/inputs/XIX/tex/components/16_expose_XIX_notations4_3_through_scholium4_4_2_page016.tex

% BEGIN INCLUDED SGA3_English_current_buildable_source_20260823/source/inputs/XIX/tex/components/17_expose_XIX_paragraph4_5_through_corollary4_9_page017.tex
\par\medskip\noindent
\textbf{\SGAInlineRefTarget{sga3:xix:paragraph:4-5}4.5.}---
Let \(S\) be a scheme, let \(T\) be an \(S\)-torus, let \(P\) be an
\(S\)-group with operator group \(\mathbb G_{m,S}\)
\textup{(for example, a vector bundle)}, and let \(\alpha\) be a
character of \(T\).  We write \(T\mathbin{\cdot_\alpha}P\) for the
semidirect product of \(P\) by \(T\), where \(T\) acts on \(P\)
through the composite morphism
\[
  T
  \xrightarrow{\ \alpha\ }
  \mathbb G_{m,S}
  \longrightarrow
  \underline{\operatorname{Aut}}_{S\text{-gr.}}(P).
\]
\SGARefTarget{sga3:xix:paragraph:4:5}

\par\medskip\noindent
\SGAInlineRefTarget{sga3:xix:definition:4-6}\textbf{\textit{\SGARefLink{sga3:xix:definition:4:6}{Definition 4.6}.}}---
\textit{Let \(S\) be a scheme, let \(G\) be an \(S\)-group scheme,
let \(T\) be a subgroup of \(G\), let \(\alpha\) be a character of
\(T\), and let \(L\) be an \(\mathcal O_S\)-module.  Let}
\[
  p\colon\mathbb W(L)\longrightarrow G
\]
\textit{be a morphism of \(S\)-group functors.}\footnote{Editor's
note: the edition uses a bold \(\mathbf W(L)\) here because, for an
arbitrary \(\mathcal O_S\)-module \(L\), the functor need not be
representable.  Below, however, \(L\) is assumed finite locally free
\textup{(indeed, invertible)}, in which case the functor is
represented by the vector bundle \(\mathbb V(L^\vee)\), and the
notation \(\mathbb W(L)\) is used without this distinction.}
\textit{We say that \(p\) is \emph{normalized by \(T\) with
multiplier \(\alpha\)} if it satisfies the following equivalent
conditions:}

\begin{enumerate}
\item[\textup{(i)}]
\(p\) is a morphism of objects with operator group \(T\), where \(T\)
acts on \(\mathbb W(L)\) through \(\alpha\), and on \(G\) by inner
automorphisms.  In other words, for every \(S'\to S\), every
\(t\in T(S')\), and every
\[
  x\in\mathbb W(L)(S')
  =
  \Gamma(S',L\otimes_{\mathcal O_S}\mathcal O_{S'}),
\]
one has
\[
  \operatorname{int}(t)p(x)=p\bigl(\alpha(t)x\bigr).
\]

\item[\textup{(ii)}]
The morphism
\[
  T\mathbin{\cdot_\alpha}\mathbb W(L)\longrightarrow G,
  \qquad
  (t,x)\longmapsto t\,p(x),
\]
defined by multiplication in \(G\), is a group morphism.
\end{enumerate}
\SGARefTarget{sga3:xix:definition:4:6}

\par\medskip\noindent
\SGAInlineRefTarget{sga3:xix:lemma:4-7}\textbf{\textit{\SGARefLink{sga3:xix:lemma:4:7}{Lemma 4.7}.}}---
\textit{Under the hypotheses of \SGARefLink{sga3:xix:definition:4-6}{4.6}, suppose in addition that \(G\)
is smooth with connected fibers,\footnote{Editor's note: compare
editor's note \textup{(28)}.} that \(T\) is a maximal torus of \(G\),
and that \(L\) is invertible.  If \(p\) is a monomorphism and
\(\alpha\) is nonzero on every fiber, then \(\alpha\) is a root of
\(G\) with respect to \(T\).}
\SGARefTarget{sga3:xix:lemma:4:7}

Indeed,
\[
  \operatorname{Lie}(p)\colon L\longrightarrow\mathfrak g
\]
is a monomorphism which factors through \(\mathfrak g^\alpha\).

\par\medskip\noindent
\SGAInlineRefTarget{sga3:xix:proposition:4-8}\textbf{\textit{\SGARefLink{sga3:xix:proposition:4:8}{Proposition 4.8}.}}---
\textit{Under the hypotheses of \SGARefLink{sga3:xix:lemma:4-7}{4.7}, suppose that \(G\) is reductive
and that \(p\) is a monomorphism.  Then \(\alpha\) is a root of \(G\)
with respect to \(T\), and \(\operatorname{Lie}(p)\) induces an
isomorphism}
\[
  \operatorname{Lie}(p)\colon
  L\xrightarrow{\ \sim\ }\mathfrak g^\alpha.
\]
\SGARefTarget{sga3:xix:proposition:4:8}

Indeed, by \SGARefLink{sga3:xix:lemma:4-7}{4.7} and the fact that \(\mathfrak g^\alpha\) is
invertible, it is enough to prove that \(\alpha\) is nonzero on every
fiber.  Let \(s\in S\) be such that
\(\alpha_{\bar s}=0\) \textup{(\(=1\) in multiplicative notation)}.
If \(X\) is a nonzero section of
\(L\otimes_S\kappa(\bar s)\), then \(p(X)\) is a unipotent element
different from the identity in \(G(\bar s)\) which centralizes
\(T_{\bar s}\).  This is impossible, since \(T_{\bar s}\) is its own
centralizer.

\par\medskip\noindent
\SGAInlineRefTarget{sga3:xix:corollary:4-9}\textbf{\textit{\SGARefLink{sga3:xix:corollary:4:9}{Corollary 4.9}.}}---
\textit{Under the hypotheses of \SGARefLink{sga3:xix:proposition:4-8}{4.8}, there exists a monomorphism of
groups with operator group \(T\) \textup{(that is, normalized by
\(T\) with multiplier \(\alpha\))}
\[
  \mathbb W(\mathfrak g^\alpha)\longrightarrow G
\]
which induces on Lie algebras the canonical morphism
\(\mathfrak g^\alpha\to\mathfrak g\).}
\SGARefTarget{sga3:xix:corollary:4:9}

We shall soon see that \SGARefLink{sga3:xix:corollary:4-9}{4.9} in fact holds whenever \(G\) is a
reductive group and \(\alpha\) is a root of \(G\) with respect to
\(T\), and that such a morphism is unique.


% END INCLUDED SGA3_English_current_buildable_source_20260823/source/inputs/XIX/tex/components/17_expose_XIX_paragraph4_5_through_corollary4_9_page017.tex

% BEGIN INCLUDED SGA3_English_current_buildable_source_20260823/source/inputs/XIX/tex/components/18_expose_XIX_reminder4_10_through_proposition4_12_statement_page018.tex
\par\medskip\noindent
\textbf{\textit{Reminder 4.10.}}---
\textit{Let \(k\) be an algebraically closed field, let \(G\) be a
reductive \(k\)-group, let \(T\) be a maximal torus of \(G\), and let
\(\alpha\) be a root of \(G\) with respect to \(T\).  There exists a
monomorphism}
\[
  p\colon\mathbb G_{a,k}\longrightarrow G
\]
\textit{normalized by \(T\) with multiplier \(\alpha\).}

See the \emph{Bible}, §13.1, Theorem 1.
\SGARefTarget{sga3:xix:reminder:4:10}

\par\medskip\noindent
\textbf{\SGAInlineRefTarget{sga3:xix:paragraph:4-11}4.11.}---
We finish this section with a technical result which will be useful
to us.  Let \(S\) be a scheme, and let \(L\) be an invertible
\(\mathcal O_S\)-module.  Let \(q>0\) be an integer such that
\(x\longmapsto x^q\) is an endomorphism of the \(S\)-group
\(\mathbb G_{a,S}\).  If \(S\ne\varnothing\), then \(q=1\), or
\(q=p^n\), where \(p\) is a prime number which vanishes on \(S\).
This follows at once from the elementary fact that the greatest
common divisor of the binomial coefficients
\(\binom qi\), for \(i\ne0,q\), is \(p\) if \(q=p^n\) with \(p\)
prime, and is \(1\) otherwise.

The morphism defined by taking the \(q\)-th power,
\[
  L\longrightarrow L^{\otimes q},
\]
is a morphism of sheaves of abelian groups.  By base change it
defines a morphism of \(S\)-group schemes
\[
  \mathbb W(L)\longrightarrow\mathbb W(L^{\otimes q}).
\]
In particular, if \(L'\) is another invertible module and an
\(\mathcal O_S\)-linear morphism
\[
  h\colon L^{\otimes q}\longrightarrow L'
\]
is given, one obtains a morphism of \(S\)-group schemes
\[
  \mathbb W(L)\longrightarrow\mathbb W(L'),
  \qquad
  x\longmapsto h(x^q).
\]
\SGARefTarget{sga3:xix:paragraph:4:11}

With this notation, we have the following result.

\par\medskip\noindent
\textbf{\textit{\SGARefLink{sga3:xix:proposition:4:12}{Proposition 4.12}.}}---
\textit{Let \(S\) be a scheme, let \(T\) \textup{(resp. \(T'\))} be
an \(S\)-torus, let \(L\) \textup{(resp. \(L'\))} be an invertible
\(\mathcal O_S\)-module, and let \(\alpha\)
\textup{(resp. \(\alpha'\))} be a character of \(T\)
\textup{(resp. \(T'\))}.  Let \(f\colon T\to T'\) be a group
morphism and let
\[
  g\colon\mathbb W(L)\longrightarrow\mathbb W(L')
\]
be a morphism of \(S\)-schemes \textup{(not necessarily a morphism
of \(S\)-groups)} satisfying}
\[
  \tag{\(\star\)}\label{sga3:xix:paragraph:4-11:equation:star}
  g\bigl(\alpha(t)x\bigr)
  =
  \alpha'\bigl(f(t)\bigr)g(x)
\]
\textit{for every \(S'\to S\), every \(x\in\mathbb W(L)(S')\), and
every \(t\in T(S')\).  Let \(s_0\in S\) be such that
\(\alpha_{\bar s_0}\ne0\).}
\footnote{Editor's note: the phrases ``consider the semidirect
product \(T\mathbin{\cdot_\alpha}\mathbb W(L)\), respectively
\(T'\mathbin{\cdot_{\alpha'}}\mathbb W(L')\)'' have been deleted.
On the other hand, in part \textup{\SGARefLink{sga3:xix:paragraph:4:11:item:b}{(b)}} the hypothesis that \(g\) be
a group morphism, together with \textup{\((\star)\)}, is equivalent
to requiring
\[
  (t,x)\longmapsto\bigl(f(t),g(x)\bigr)
\]
to be a group morphism from
\(T\mathbin{\cdot_\alpha}\mathbb W(L)\) to
\(T'\mathbin{\cdot_{\alpha'}}\mathbb W(L')\).}

\begin{enumerate}
\item[\textup{(a)}]\SGAInlineRefTarget{sga3:xix:proposition:4:12:item:a}
\textit{Suppose that \(g\) sends the zero section to the zero section
and that, for every integer \(n>0\),}
\[
  (\alpha'\circ f)_{\bar s_0}\ne n\alpha_{\bar s_0}.
\]
\textit{Then \(g=0\) in a neighborhood of \(s_0\).}

\item[\textup{\SGAInlineRefTarget{sga3:xix:proposition:4-12:item:b}\SGAInlineRefTarget{sga3:xix:paragraph:4:11:item:b}(b)}]
\textit{Suppose that \(g\) is a group morphism and that
\(g_{\bar s_0}\ne0\).  There then exist an open subset \(U\subset S\)
containing \(s_0\) and an integer \(q>0\) such that
\(x\longmapsto x^q\) is an endomorphism of
\(\mathbb G_{a,U}\) and}
\[
  (\alpha'\circ f)_U=q\alpha_U.
\]
\end{enumerate}
\SGARefTarget{sga3:xix:proposition:4:12}


% END INCLUDED SGA3_English_current_buildable_source_20260823/source/inputs/XIX/tex/components/18_expose_XIX_reminder4_10_through_proposition4_12_statement_page018.tex

% BEGIN INCLUDED SGA3_English_current_buildable_source_20260823/source/inputs/XIX/tex/components/19_expose_XIX_proposition4_12_part_c_and_proof_page019.tex
\begin{enumerate}
\item[\textup{\SGAInlineRefTarget{sga3:xix:proposition:4:12:item:c}(c)}]
\textit{Suppose that}
\[
  (\alpha'\circ f)_{\bar s_0}=q\alpha_{\bar s_0},
\]
\textit{where \(q>0\) is an integer such that
\(x\longmapsto x^q\) is an endomorphism of
\(\mathbb G_{a,S}\).  There then exist an open subset \(U\subset S\)
containing \(s_0\) and a unique morphism of
\(\mathcal O_S\)-modules}
\[
  h\colon L^{\otimes q}|_U\longrightarrow L'|_U
\]
\textit{such that \(g_U\) is the composite}
\[
  \mathbb W(L)_U
  \xrightarrow{\ x\mapsto x^q\ }
  \mathbb W(L^{\otimes q})_U
  \xrightarrow{\ \mathbb W(h)\ }
  \mathbb W(L')_U.
\]
\end{enumerate}

\SGAProof{Proof of \SGARefLink{sga3:xix:proposition:4:12}{Proposition 4.12}.}
Let us prove \textup{(a)}.  Since the conclusion is local on \(S\),
we may suppose that
\(\mathbb W(L)=\mathbb W(L')=\mathbb G_{a,S}\), and hence that \(g\)
is given by a polynomial
\[
  g(X)=\sum_{n\geq0}a_nX^n,
  \qquad
  a_n\in\Gamma(S,\mathcal O_S).
\]
Condition \textup{\((\star)\)}, which relates \(f\) and \(g\), is the
identity in \(\Gamma(S',\mathcal O_{S'})[X]\)
\[
  \sum_{n\geq0}a_n\alpha'\bigl(f(t)\bigr)X^n
  =
  \sum_{n\geq0}a_n\alpha(t)^nX^n.
\]
Thus, for every \(n\geq0\), every \(S'\to S\), and every
\(t\in T(S')\),
\[
  a_n\bigl(\alpha'(f(t))-\alpha(t)^n\bigr)=0.
\]

For every \(n\geq0\), let \(S_n\) be the set of points \(s\in S\)
such that
\[
  (\alpha'\circ f)_{\bar s}=n\alpha_{\bar s}.
\]
It is known \textup{(\SGARefLink{sga3:ix:corollary:5:3}{Exposé IX, 5.3})} that the \(S_n\) are both open
and closed and that
\[
  (\alpha'\circ f)_{S_n}=n\alpha_{S_n}.
\]
Moreover, since \(\alpha_{\bar s_0}\ne0\), after shrinking \(S\) we
may suppose that \(\alpha\) is nonzero on every fiber \textup{(same
reference)}; it follows that the \(S_n\) are pairwise disjoint.
After shrinking \(S\) once more, we may therefore suppose that one
of the following two cases holds: \(S=S_n\) for some \(n\), or every
\(S_n\) is empty.

Let \(m\geq0\) be such that \(S_m=\varnothing\).  We claim that
\(a_m=0\).  Indeed, \(\alpha'\circ f\) and \(m\alpha\) are distinct
on every fiber of \(S\), and we have:

\par\medskip\noindent
\textbf{\textit{\SGARefLink{sga3:xix:lemma:4:13}{Lemma 4.13}.}}---
\textit{Let \(S\) be a scheme, let \(T\) be an \(S\)-torus, and let
\(\alpha,\alpha'\) be two characters of \(T\) which are distinct on
every fiber.  There exist an fpqc covering family
\(\{S_i\to S\}\) and, for every \(i\), an element \(t_i\in T(S_i)\)
such that}
\[
  \alpha(t_i)-\alpha'(t_i)=1.
\]
\SGARefTarget{sga3:xix:lemma:4:13}

The lemma is immediate by reduction first to the diagonalizable
case, and then to \(T=\mathbb G_{m,S}\).

We return to the proof of \SGARefLink{sga3:xix:proposition:4:12}{Proposition 4.12}.  We have just proved
\textup{(a)}.  In cases \textup{(b)} and \textup{\SGARefLink{sga3:xix:proposition:4:12:item:c}{(c)}}, there is an
integer \(n\) such that \(S=S_n\) \textup{(\(n=q\) in
\textup{\SGARefLink{sga3:xix:proposition:4:12:item:c}{(c)}})}.  The preceding argument gives \(a_m=0\) for
\(m\ne n\), and therefore
\[
  g(X)=a_nX^n,
  \qquad
  a_n\in\Gamma(S,\mathcal O_S).
\]
This proves \textup{(c)} at once.  In case \textup{(b)}, we know
that \(a_n(s_0)\ne0\); we may therefore suppose that \(a_n\) is
invertible on \(S\).  It follows that \(x\longmapsto x^n\) is an
endomorphism of \(\mathbb G_{a,S}\), by the hypothesis in
\textup{(b)}, and this completes the proof.


% END INCLUDED SGA3_English_current_buildable_source_20260823/source/inputs/XIX/tex/components/19_expose_XIX_proposition4_12_part_c_and_proof_page019.tex

% BEGIN INCLUDED SGA3_English_current_buildable_source_20260823/source/inputs/XIX/tex/components/20_expose_XIX_section5_1_through_lemma5_3_page020.tex
\SGANextSectionTarget{sga3:xix:section:5}\SGASection{5}{An instructive example}

\par\medskip\noindent
\textbf{\SGAInlineRefTarget{sga3:xix:paragraph:5-1}5.1.}---
Let \(k\) be an algebraically closed field of characteristic \(0\).
Put \(A=k[t]\), the polynomial ring in one variable over \(k\), and
let \(S=\operatorname{Spec}(A)\).  Consider the following Lie algebra
\(\mathfrak g\) over \(\mathcal O_S\).  As a module it is free of
rank \(3\), with basis \(\{X,Y,H\}\), and its multiplication table is
\[
  [X,Y]=2tH,
  \qquad
  [H,X]=X,
  \qquad
  [H,Y]=-Y.
\]
For \(s\in S\), \(s\ne s_0\), where \(s_0\) is the point defined by
\(t=0\), the fiber
\[
  \mathfrak g(s)=\mathfrak g\otimes_A\kappa(s)
\]
is isomorphic to the Lie algebra of
\(\mathrm{PGL}_{2,\kappa(s)}\).  For \(s=s_0\), it is a solvable Lie
algebra.
\SGARefTarget{sga3:xix:paragraph:5:1}

\par\medskip\noindent
\textbf{\SGAInlineRefTarget{sga3:xix:paragraph:5-2}5.2.}---
Let \(G_1\) be the group scheme of automorphisms of \(\mathfrak g\):
for every \(S'\to S\), \(G_1(S')\) is the group of automorphisms of
the \(\mathcal O_{S'}\)-Lie algebra
\(\mathfrak g\otimes_{\mathcal O_S}\mathcal O_{S'}\).  It is a closed
subscheme of the group \(\mathrm{GL}(\mathfrak g)\) of automorphisms
of the \(\mathcal O_S\)-module \(\mathfrak g\).

Let \(S'\to S\), and let
\[
  u\in M_3\bigl(\Gamma(S',\mathcal O_{S'})\bigr)
\]
be viewed as an endomorphism of the \(\mathcal O_{S'}\)-module
\(\mathfrak g\otimes_{\mathcal O_S}\mathcal O_{S'}\):
\[
  \begin{aligned}
  u(X)&=aX+bY+eH,\\
  u(Y)&=b'X+a'Y+e'H,\\
  u(H)&=cX+c'Y+dH.
  \end{aligned}
\]
It is immediate that \(u\) is a section of \(G_1\) if and only if
\(\det(u)\) is invertible and the following relations hold:
\[
\begin{alignedat}{4}
  \textup{(1)}\quad&a(d-1)=ec,
  &\qquad&
  \textup{(1')}\quad&a'(d-1)=e'c',\\
  \textup{(2)}\quad&b(d+1)=ec',
  &&
  \textup{(2')}\quad&b'(d+1)=e'c,\\
  \textup{(3)}\quad&e=2t(bc-ac'),
  &&
  \textup{(3')}\quad&e'=2t(b'c'-a'c),\\
  \textup{(4)}\quad&2tc=eb'-ae',
  &&
  \textup{(4')}\quad&2tc'=be'-ea',\\
  \textup{(5)}\quad&2t(aa'-bb')=2td.&&&
\end{alignedat}
\]
\footnote{Editor's note: the equality
\([u(X),u(Y)]=2t\,u(H)\) gives relations \textup{(4)},
\textup{(4')}, and \textup{(5)}; the equality
\([u(H),u(X)]=u(X)\) gives \textup{(1)--(3)}; and the equality
\([u(H),u(Y)]=-u(Y)\) gives \textup{(1')--(3')}.}
\SGARefTarget{sga3:xix:paragraph:5:2}

\par\medskip\noindent
\SGAInlineRefTarget{sga3:xix:lemma:5-3}\textbf{\textit{\SGARefLink{sga3:xix:lemma:5:3}{Lemma 5.3}.}}---
\textit{Relations \textup{(1)}, \textup{(1')}, \textup{(2)}, and
\textup{(2')} imply}
\[
  \det(u)=aa'(2-d)+bb'(d+2),
  \qquad
  aa'-bb'=d\det(u).
\]
\SGARefTarget{sga3:xix:lemma:5:3}

Indeed, the first assertion follows at once by substituting
relations \textup{(1)}, \textup{(1')}, \textup{(2)}, and
\textup{(2')} in the expansion of \(\det(u)\):
\[
\begin{aligned}
  \det(u)
  &=aa'd+be'c+b'c'e-a'ec-ae'c'-bb'd\\
  &=aa'd+bb'(d+1)+bb'(d+1)
    -aa'(d-1)-bb'd-aa'(d-1)\\
  &=aa'(d-d+1-d+1)+bb'(d+1+d+1-d)\\
  &=aa'(2-d)+bb'(d+2).
\end{aligned}
\]
Multiplying this relation by \(d\), we obtain
\[
  d\det(u)=aa'(2d-d^2)+bb'(d^2+2d).
\]


% END INCLUDED SGA3_English_current_buildable_source_20260823/source/inputs/XIX/tex/components/20_expose_XIX_section5_1_through_lemma5_3_page020.tex

% BEGIN INCLUDED SGA3_English_current_buildable_source_20260823/source/inputs/XIX/tex/components/21_expose_XIX_lemma5_3_completion_through_lemma5_7_page021.tex
On the other hand, the relation
\(\textup{(1)}\times\textup{(1')}=
\textup{(2)}\times\textup{(2')}\) gives at once
\[
  aa'(d-1)^2=bb'(d+1)^2.
\]
Combining the last two relations immediately yields the second
formula of \SGARefLink{sga3:xix:lemma:5:3}{Lemma 5.3}.

\par\medskip\noindent
\textbf{5.4.}---
Now consider
\[
  G=G_1\cap\mathrm{SL}(\mathfrak g).
\]
This is the closed subgroup of \(G_1\) defined by the equation
\(\det(u)=1\); hence it is an affine group over \(S\).
\SGARefTarget{sga3:xix:paragraph:5:4}

\par\medskip\noindent
\SGAInlineRefTarget{sga3:xix:proposition:5-5}\textbf{\textit{\SGARefLink{sga3:xix:proposition:5:5}{Proposition 5.5}.}}---
\textit{The group \(G\) is smooth over \(S\).}
\SGARefTarget{sga3:xix:proposition:5:5}

To prove the proposition, we shall need the following lemmas.

\par\medskip\noindent
\textbf{\textit{\SGARefLink{sga3:xix:lemma:5:6}{Lemma 5.6}.}}---
\textit{Let \(U\) be the open subset of}
\[
  \operatorname{End}_A(\mathfrak g)
  \simeq
  \mathbb W\bigl(M_3(\mathcal O_S)\bigr)
\]
\textit{defined by the condition that \(ad\) be invertible; that is,
the principal open subset
\(\operatorname{End}_A(\mathfrak g)_f\), where
\(f(u)=ad\).  Then \(U\cap G\) is the closed subscheme of \(U\)
defined by the six equations}
\[
  \textup{(1)},\quad
  \textup{(2)},\quad
  \textup{(2')},\quad
  \textup{(3)},\quad
  \textup{(3')},\quad
  \textup{(D)}\colon aa'-bb'=d.
\]
\SGARefTarget{sga3:xix:lemma:5:6}

It is clear that these six relations are satisfied by every point
of \(G\), by \SGARefLink{sga3:xix:lemma:5-3}{Lemma 5.3}, and therefore by every point of \(U\cap G\).
Conversely, suppose that \(u\in U(S')\), for some \(S'\to S\),
satisfies the six conditions in the statement.  We must show that
\(\det(u)=1\), and that \(u\) also satisfies \textup{(1')},
\textup{(4)}, \textup{(4')}, and \textup{(5)}.

First, \textup{(D)} implies \textup{(5)}.  Relations \textup{(2)}
and \textup{(2')} give
\[
  bb'(d+1)=bce'=b'c'e.
\]
On the other hand, from \textup{(3)} and \textup{(3')}, writing
\(2t(bc-ac')(b'c'-a'c)\) in two ways, we obtain
\[
  (bc-ac')e'=(b'c'-a'c)e.
\]
Together with the preceding relation, this gives
\[
  ac'e'=a'ce.
\]
By \textup{(1)}, however,
\[
  a'ce=a'a(d-1).
\]
Thus
\[
  a\bigl(a'(d-1)-e'c'\bigr)=0,
\]
and this implies \textup{(1')}, since \(a\) is invertible.

Consequently, \textup{(1)}, \textup{(2)}, \textup{(2')}, and
\textup{(1')} hold.  \SGARefLink{sga3:xix:lemma:5:3}{Lemma 5.3} and \textup{(D)} now give
\[
  d\bigl(\det(u)-1\bigr)=0.
\]
Since \(d\) is invertible, it follows that \(\det(u)=1\).

It remains to prove \textup{(4)} and \textup{(4')}.  We prove
\textup{(4')}; the other calculation follows by symmetry.  From
\textup{(3)}, \textup{(3')}, and \textup{(D)}, we have
\[
  a'e+be'=-2t(aa'-bb')c'=-2tdc'.
\]
Combining \textup{(1')} and \textup{(2)}, we also have
\[
  a'e+be'=-d(be'-ea'),
\]
\footnote{Editor's note: the sign has been corrected.}
which proves \textup{(4')}, since \(d\) is invertible.

\par\medskip\noindent
\SGAInlineRefTarget{sga3:xix:lemma:5-7}\textbf{\textit{\SGARefLink{sga3:xix:lemma:5:7}{Lemma 5.7}.}}---
\textit{The group \(G\) is smooth over \(S\) along the identity
section.}
\SGARefTarget{sga3:xix:lemma:5:7}


% END INCLUDED SGA3_English_current_buildable_source_20260823/source/inputs/XIX/tex/components/21_expose_XIX_lemma5_3_completion_through_lemma5_7_page021.tex

% BEGIN INCLUDED SGA3_English_current_buildable_source_20260823/source/inputs/XIX/tex/components/22_expose_XIX_lemma5_7_proof_through_paragraph5_11_page022.tex
\SGAProof{Proof}
By 5.6 and SGA 1, II 4.10, it is enough to prove that the
differentials of the functions
\[
\begin{gathered}
  a(d-1)-ec,\qquad
  b(d+1)-ec',\qquad
  b'(d+1)-e'c,\\
  e-2t(bc-ac'),\qquad
  e'-2t(b'c'-a'c),\qquad
  aa'-bb'-d
\end{gathered}
\]
at the points of the identity section of \(G\) are linearly
independent.  If the differential of a lower-case letter is denoted
by the corresponding capital letter, these differentials are
\[
  D,\qquad
  2B,\qquad
  2B',\qquad
  E+2tC',\qquad
  E'+2tC,\qquad
  A+A'-D,
\]
\footnote{Editor's note: \(A+A'+B\) has been corrected to
\(A+A'-D\).}
and they are indeed linearly independent modulo every
\((t-\lambda)\), \(\lambda\in k\).
\footnote{Editor's note: ``\((t-a)\), \(a\in A\)'' has been
corrected to ``\((t-\lambda)\), \(\lambda\in k\)''.}

\par\medskip\noindent
\textbf{\textit{\SGARefLink{sga3:xix:lemma:5:8}{Lemma 5.8}.}}---
\textit{For \(s\in S\), \(s\ne s_0\), the fiber \(G_s\) is connected
and semisimple.}
\SGARefTarget{sga3:xix:lemma:5:8}

\SGAProof{Proof}
Indeed, since \(s\ne s_0\), \(\mathfrak g(s)\) is isomorphic to the
Lie algebra of \(\mathrm{PGL}_{2,\kappa(s)}\), and, on the other hand,
\(G_s=(G_1)_s\).  The automorphism group of the Lie algebra of
\(\mathrm{PGL}_2\) over a field of characteristic \(0\) is known to
be \(\mathrm{PGL}_2\) itself, which is connected and semisimple.
\footnote{Editor's note: the preceding sentence has been slightly
modified.}

\par\medskip\noindent
\textbf{\textit{\SGARefLink{sga3:xix:lemma:5:9}{Lemma 5.9}.}}---
\textit{The fiber \(G_{s_0}\) is solvable and has two connected
components, of the following form:}
\[
\begin{aligned}
  G^0_{s_0}
  &=
  \left\{
  \begin{pmatrix}
    a&0&0\\
    0&a^{-1}&0\\
    c&c'&1
  \end{pmatrix}
  \right\},
  &
  G^-_{s_0}
  &=
  \left\{
  \begin{pmatrix}
    0&b&0\\
    b^{-1}&0&0\\
    c&c'&-1
  \end{pmatrix}
  \right\}.
\end{aligned}
\]
\SGARefTarget{sga3:xix:lemma:5:9}

\SGAProof{Proof}
We have \(e=e'=0\), since \(t=0\) at \(s_0\).  The equations
\textup{(1)}, \textup{(1')}, \(\ldots\), \textup{(5)}, and
\textup{(D)} can then be solved immediately.

\par\medskip\noindent
\SGAInlineRefTarget{sga3:xix:lemma:5-10}\textbf{\textit{\SGARefLink{sga3:xix:lemma:5:10}{Lemma 5.10}.}}---
\textit{The matrix}
\[
  w=
  \begin{pmatrix}
    0&1&0\\
    1&0&0\\
    0&0&-1
  \end{pmatrix}
\]
\textit{is a section of \(G\) over \(S\), such that
\(w(s_0)\in G^-_{s_0}\).}
\SGARefTarget{sga3:xix:lemma:5:10}

Let us now prove \SGARefLink{sga3:xix:proposition:5-5}{Proposition 5.5}.  Denote by \(G^0\) the union of the
identity components of the fibers of \(G\), that is, the complement
of \(G^-_{s_0}\).  Since \(G\) is smooth over \(S\) along the identity
section by \SGARefLink{sga3:xix:lemma:5-7}{Lemma 5.7}, \(G^0\) is an open subgroup of \(G\), smooth
over \(S\), by Exposé VI\textsubscript{B}, \SGARefLink{thm:sga3-VIB-3.10}{3.10}.  By translation,
\(G\) is also plainly smooth at the points of \(w(S)\); hence \(G\)
is smooth over \(S\).
\footnote{Editor's note: the preceding sentence has been modified,
and the reference to Exposé VI\textsubscript{B}, \SGARefLink{thm:sga3-VIB-3.10}{3.10} has been
added.}

\par\medskip\noindent
\textbf{\SGAInlineRefTarget{sga3:xix:paragraph:5-11}5.11.}---
Consider the morphism
\[
  \mathbb G_{m,S}\longrightarrow G^0,
  \qquad
  z\longmapsto
  \begin{pmatrix}
    z&0&0\\
    0&z^{-1}&0\\
    0&0&1
  \end{pmatrix}.
\]
It is a monomorphism defining a torus \(T\) of \(G^0\).  We claim
that
\[
  T=\operatorname{Centr}_G(T)
   =\operatorname{Centr}_{G^0}(T).
\]
\SGARefTarget{sga3:xix:paragraph:5:11}

% END INCLUDED SGA3_English_current_buildable_source_20260823/source/inputs/XIX/tex/components/22_expose_XIX_lemma5_7_proof_through_paragraph5_11_page022.tex

% BEGIN INCLUDED SGA3_English_current_buildable_source_20260823/source/inputs/XIX/tex/components/23_expose_XIX_paragraph5_11_completion_through_proposition6_1_opening_page023.tex
It is enough to verify the first equality.  Since the two sides are
smooth \(S\)-subgroups of \(G\), it is enough to check that they have
the same geometric points.  In the fibers over points \(s\ne s_0\),
this follows because \(\mathrm{PGL}_{2,\kappa(s)}\) is reductive and
\(T_s\) is a maximal torus in it for dimensional reasons
\textup{(cf.\ \SGARefLink{sga3:xix:lemma:1-11}{1.11})}.  In the fiber over \(s_0\), the verification is
immediate.  In particular, \(T\) is a maximal torus of both \(G\)
and \(G^0\).

\par\medskip\noindent
\textbf{5.12.}---
The section \(w\) of \(G\) defined in \SGARefLink{sga3:xix:lemma:5-10}{Lemma 5.10} normalizes \(T\).
It follows immediately \textup{(cf.\ \SGARefLink{sga3:xix:remark:2-4}{2.4})} that the Weyl group of
\(G\) is isomorphic to \((\mathbb Z/2\mathbb Z)_S\), and is therefore
finite over \(S\):
\[
  W_G(T)=\operatorname{Norm}_G(T)/T
        =(\mathbb Z/2\mathbb Z)_S.
\]
By contrast, \(W_{G^0}(T)\) is not finite over \(S\): it ``lacks one
point'' above \(s_0\).
\SGARefTarget{sga3:xix:paragraph:5:12}

\par\medskip\noindent
\textbf{\SGAInlineRefTarget{sga3:xix:paragraph:5-13}5.13.}---
The open immersion \(G^0\to G\) is not a closed immersion, since
\(G^0\) is dense in \(G\); it is nevertheless an affine morphism
\textup{(and hence \(G^0\) is affine over \(S\))}.  Indeed,
\(G^0_{s_0}\) is closed in \(G_{s_0}\), which is closed in \(G\), so
its complement \(U\) in \(G\) is open.  The open subsets \(G^0\) and
\(U\) cover \(G\), and it is enough to verify that the immersions
\[
  G^0\longrightarrow G^0
  \qquad\text{and}\qquad
  G^0\cap U\longrightarrow U
\]
are affine.  This is trivial for the first; for the second, observe
that \(U\cap G^0\) is defined in \(U\) by the condition \(b\ne0\).

We have thus constructed a smooth affine \(S\)-group \(G^0\), with
connected fibers, possessing a self-centralizing maximal torus \(T\)
whose Weyl group \(W_{G^0}(T)\) is not finite \textup{(compare
\SGARefLink{sga3:xix:theorem:2:5}{Theorem 2.5})}.
\SGARefTarget{sga3:xix:paragraph:5:13}

\SGANextSectionTarget{sga3:xix:section:6}\SGASection{6}{Local existence of maximal tori.  The Weyl group}

In the proof of \SGARefLink{sga3:xix:theorem:2:5}{Theorem 2.5} we used a result from Exposé XI on the
local existence, for the étale topology, of maximal tori.  The proof
in Exposé XI uses a deep representability result \textup{(\SGARefLink{sga3:xi:theorem:4:1}{XI 4.1})}.
In the particular case under consideration, one can give another,
far more elementary proof, based on the ideas of Exposé XII, no.~7.

\par\medskip\noindent
\SGAInlineRefTarget{sga3:xix:proposition:6-1}\textbf{\textit{\SGARefLink{sga3:xix:proposition:6:1}{Proposition 6.1}.}}---
\textit{Let \(S\) be a scheme, let \(G\) be a smooth affine
\(S\)-group scheme with connected fibers over \(S\), and let
\(s_0\in S\) be a point such that the maximal tori of the geometric
fiber \(G_{\bar s_0}\) are self-centralizing.  There exist an étale
morphism \(S'\to S\) whose image contains \(s_0\) and a split maximal
torus \(T\) of \(G_{S'}\).}
\SGARefTarget{sga3:xix:proposition:6:1}

\SGAProof{Proof}
First, we may suppose that \(S\) is affine.
\footnote{Editor's note: the preceding sentence and the reference
to EGA IV below have been added.}
Since \(G\) is of finite presentation over \(S\), we may successively
suppose that \(S\) is noetherian, local, and then henselian with
separably closed residue field \textup{(cf.\ EGA IV, 8.12,
\S\S 8.8 and 18.8)}.  Thus let
\[
  S=\operatorname{Spec}(A),
  \qquad
  k=\kappa(s_0),
\]
where \(A\) is henselian and \(k\) is separably closed.  Choose a
maximal torus \(T_0\) of \(G_0=G_k\).  Such a torus exists, for
example because the scheme of maximal tori of \(G_0\) is smooth over
\(k\) \textup{(\SGARefLink{sga3:xii:7.1:item:c}{Exposé XII, 7.1(c)})}.  Since \(k\) is separably
closed, \(T_0\) is split \textup{(cf.\ \SGARefLink{sga3:x:proposition:1:4}{Exposé X, 1.4})}, and is
therefore given by a group monomorphism
\[
  f_0\colon \mathbb G_{m,k}^{\,r}\longrightarrow G_0.
\]

% END INCLUDED SGA3_English_current_buildable_source_20260823/source/inputs/XIX/tex/components/23_expose_XIX_paragraph5_11_completion_through_proposition6_1_opening_page023.tex

% BEGIN INCLUDED SGA3_English_current_buildable_source_20260823/source/inputs/XIX/tex/components/24_expose_XIX_proposition6_1_proof_through_proposition6_3_opening_page024.tex
Let \(m>1\) be an integer prime to the characteristic of \(k\).
By \SGARefLink{sga3:viii:corollary:6:7}{Exposé VIII, 6.7}, for every \(h>0\),
\(\operatorname{Centr}_{G_0}({}_{m^h}T_0)\) is represented by a
closed subscheme of \(G_0\).  The subgroup schemes
\({}_{m^h}T_0\) are schematically dense in \(T_0\)
\textup{(cf.\ \SGARefLink{sga3:ix:remark:4:10}{Exposé IX, 4.10})}, and \(G_0\) is noetherian.
Consequently there is an \(h\) such that
\[
  \operatorname{Centr}_{G_0}({}_{m^h}T_0)
  =
  \operatorname{Centr}_{G_0}(T_0)
  =
  T_0.
\]
Put \(n=m^h\).  Since \(n\) is invertible on \(S\),
\({}_n\mathbb G_{m,S}\) is isomorphic to
\((\mathbb Z/n\mathbb Z)_S\), so \(f_0\) defines a group
monomorphism
\[
  u_0\colon(\mathbb Z/n\mathbb Z)_k^{\,r}\longrightarrow G_0
\]
such that \(\operatorname{Centr}_{G_0}(u_0)=T_0\).

Now the \(S\)-functor
\[
  P=
  \operatorname{Hom}_{S\text{-}\mathrm{gr}}
  \bigl((\mathbb Z/n\mathbb Z)_S^{\,r},G\bigr)
\]
is represented by an \(S\)-scheme of finite type: it is a closed
subscheme of \(r\) copies of
\[
  {}_nG
  =
  \operatorname{Hom}_{S\text{-}\mathrm{gr}}
  \bigl((\mathbb Z/n\mathbb Z)_S,G\bigr)
  =
  \operatorname{Ker}\bigl(G\xrightarrow{\,n\,}G\bigr).
\]
But \(P\) is smooth over \(S\) \textup{(\SGARefLink{sga3:ix:theorem:3:6}{Exposé IX, 3.6})}.
Therefore \(u_0\in P(k)\) lifts to a section \(u\in P(S)\), by
Hensel's lemma \textup{(\SGARefLink{sga3:xi:corollary:1:11}{Exposé XI, 1.11})}:
\[
  u\colon(\mathbb Z/n\mathbb Z)_S^{\,r}\longrightarrow G.
\]

Consider \(H=\operatorname{Centr}_G(u)\).  It is a closed subgroup
scheme of \(G\), by \SGARefLink{sga3:viii:examples:6:5:item:e:01}{Exposé VIII, 6.5(e)}, and its special fiber is
\(H_0=T_0\).  Moreover, \(H\) is smooth over \(S\).  Indeed, let
\(S'=\operatorname{Spec}(A')\) be an affine \(S\)-scheme, let
\[
  u'\colon({}_n\mathbb G_{m,S'})^r\longrightarrow G_{S'}
\]
be obtained from \(u\) by base change, put
\(S'_J=\operatorname{Spec}(A'/J)\), where \(J\) is a square-zero
ideal, and let \(x\in H(S'_J)\).  Since \(G\) is smooth, \(x\) lifts
to an element \(g\in G(S')\).  Then
\[
  v=\operatorname{int}(g)(u')
\]
satisfies \(v_J=u'_J\).  By \SGARefLink{sga3:ix:theorem:3:2}{Exposé IX, 3.2}, there is therefore an
element \(g'\in G(S')\) such that \(g'_J=e\) and
\(\operatorname{int}(g')(v)=u'\).  It follows that \(h=g'g\) belongs
to \(H(S')\) and satisfies \(h_J=x\).
\footnote{Editor's note: the reference to Exposé VIII, 6.5(e) was
added above, and the preceding smoothness argument has been
expanded.}

Let \(H^0\) be the identity component of \(H\).  It is a subgroup
scheme of \(G\), smooth and with connected fibers, whose special
fiber is a torus.  By \SGARefLink{sga3:x:theorem:8-1}{Exposé X, 8.1}, it is a torus, necessarily split
by \SGARefLink{sga3:x:corollary:4-6}{Exposé X, 4.6}.  Put \(T=H^0\), and let
\[
  C=\operatorname{Centr}_G(T).
\]
This is a closed subgroup of \(G\) \textup{(\SGARefLink{sga3:viii:examples:6:5:item:e:01}{Exposé VIII, 6.5(e)})}
and is smooth \textup{(\SGARefLink{sga3:xi:corollary:2:4}{Exposé XI, 2.4})}.  Consider \(C^0\)
\textup{(in fact \(C^0=C\), but we do not need this)}.  Then
\(C^0\supset T\), and both are smooth groups with connected fibers.
They coincide at \(s_0\), hence in a neighborhood of \(s_0\).
After shrinking \(S\), we may therefore suppose that \(C^0=T\);
then, a fortiori, \(T\) is maximal.

\par\medskip\noindent
\SGAInlineRefTarget{sga3:xix:remark:6-2}\textbf{\textit{\SGARefLink{sga3:xix:remark:6:2}{Remark 6.2}.}}---
\textit{The proof shows in particular that the reductive rank of
\(G_{\bar s}\) is constant in a neighborhood of \(s=s_0\).}
\SGARefTarget{sga3:xix:remark:6:2}

\par\medskip\noindent
\SGAInlineRefTarget{sga3:xix:proposition:6-3}\textbf{\textit{\SGARefLink{sga3:xix:proposition:6:3}{Proposition 6.3}.}}---
\textit{Let \(S\) be a scheme, let \(G\) be an \(S\)-group scheme
that is smooth and of finite presentation over \(S\), and let \(Q\)
be a subtorus of \(G\).}
\SGARefTarget{sga3:xix:proposition:6:3}

% END INCLUDED SGA3_English_current_buildable_source_20260823/source/inputs/XIX/tex/components/24_expose_XIX_proposition6_1_proof_through_proposition6_3_opening_page024.tex

% BEGIN INCLUDED SGA3_English_current_buildable_source_20260823/source/inputs/XIX/tex/components/25_expose_XIX_proposition6_3_completion_and_bibliography_page025.tex
\begin{enumerate}
\item[\textup{(i)}]
\textit{\(\operatorname{Centr}_G(Q)\) and
\(\operatorname{Norm}_G(Q)\) are represented by closed subgroup
schemes, smooth \textup{(and hence of finite presentation)} over
\(S\).}

\item[\textup{(ii)}]
\textit{\(\operatorname{Centr}_G(Q)\) is an open and closed
subscheme of \(\operatorname{Norm}_G(Q)\).  The quotient}
\[
  W_G(Q)
  =
  \operatorname{Norm}_G(Q)/\operatorname{Centr}_G(Q)
\]
\textit{is represented by an open subgroup scheme of
\(\operatorname{Aut}_{S\text{-}\mathrm{gr}}(Q)\); it is therefore a
quasi-finite, étale, separated \(S\)-group scheme.}

\item[\textup{(iii)}]
\textit{For every \(s\in S\), put}
\[
  w(s)
  =
  \left|
  \operatorname{Norm}_{G(\bar s)}(Q(\bar s))/
  \operatorname{Centr}_{G(\bar s)}(Q(\bar s))
  \right|.
\]
\textit{Then \(s\longmapsto w(s)\) is lower semicontinuous, and is
constant in a neighborhood of \(s\) if and only if \(W_G(Q)\) is
finite over \(S\) in a neighborhood of \(s\).}
\end{enumerate}

\SGAProof{Proof}
By \SGARefLink{sga3:xi:corollary:6:11}{Exposé XI, 6.11},
\(\operatorname{Centr}_G(Q)\) and \(\operatorname{Norm}_G(Q)\) are
represented by closed subschemes of \(G\) of finite presentation.
They are smooth by \SGARefLink{sga3:xi:corollary:2:4}{Exposé XI, 2.4} and 2.4 bis, which proves
\textup{(i)}.  Assertions \textup{(ii)} and \textup{(iii)} are then
proved as in \SGARefLink{sga3:xi:proposition:5:9}{Exposé XI, 5.9} and 5.10; those proofs in fact use only
\textup{(i)}, and not the deeper theorems of \SGARefLink{sga3:xi:theorem:4:1}{Exposé XI, 4.1} and 4.2.

\section*{Bibliography}
\addcontentsline{toc}{section}{Bibliography}

\begin{sloppypar}
\begin{description}
\item[{[Bible]}]
C.~Chevalley, with the collaboration of P.~Cartier,
A.~Grothendieck, and M.~Lazard,
\textit{Classification des groupes de Lie algébriques}, 1956--58.

\item[{[Ch05]}]
C.~Chevalley,
\textit{Classification des groupes algébriques semi-simples}, with
the collaboration of P.~Cartier, A.~Grothendieck, and M.~Lazard,
\textit{Collected Works}, vol.~3, Springer, 2005.

\item[{[Tô55]}]
C.~Chevalley, \textit{Sur certains groupes simples},
\textit{Tôhoku Math.\ J.}\ (2) \textbf{7} (1955), 14--66.

\item[{[TO70]}]
J.~Tate and F.~Oort, \textit{Group schemes of prime order},
\textit{Ann.\ scient.\ Éc.\ Norm.\ Sup.}\ (4), t.~\textbf{3}
(1970), 1--21.
\end{description}
\end{sloppypar}

% END INCLUDED SGA3_English_current_buildable_source_20260823/source/inputs/XIX/tex/components/25_expose_XIX_proposition6_3_completion_and_bibliography_page025.tex



% END INCLUDED SGA3_English_current_buildable_source_20260823/source/inputs/XIX/body.tex


\SGAStartExpose{XX}
\chapter*{Expos\'e XX\\Reductive Groups of Semisimple Rank 1}
\addcontentsline{toc}{chapter}{Expos\'e XX. Reductive Groups of Semisimple Rank 1}
\begin{center}
  \textit{by M. Demazure}
\end{center}
% BEGIN INCLUDED SGA3_English_current_buildable_source_20260823/source/inputs/XX/body.tex
% Mechanically extracted after \mainmatter from the frozen standalone master.

\markboth{Exposé XX. Reductive Groups of Semisimple Rank 1}
  {Exposé XX. Reductive Groups of Semisimple Rank 1}
% BEGIN INCLUDED SGA3_English_current_buildable_source_20260823/source/inputs/XX/tex/components/00_expose_XX_recollection11_through_14_en.tex
% Authority: Exp20-13oct24.pdf, local pp. 1--3 / combined pp. 1060--1062.
% Continuous English stops immediately before Theorem 1.5.

\SGARefTarget{sga3:xx:expose}

\setcounter{footnote}{0}
\footnotetext{Editors' note: version of 13 October 2024.}

\section{Elementary systems. The groups \(U_\alpha\) and \(U_{-\alpha}\)}
\SGARefTarget{sga3:xx:section:1}

\paragraph{\textbf{Recollection \SGARefLink{sga3:xx:section:1}{1.1}.}}
\SGARefTarget{sga3:xx:1.1}
\emph{Let \(S=\Spec(k)\), where \(k\) is an algebraically closed
field, and let \(G\) be a reductive \(S\)-group of semisimple rank
one and \(T\) a maximal torus of \(G\), necessarily split. Then}
\[
  \mathfrak g=\mathfrak t\oplus\mathfrak g^\alpha
  \oplus\mathfrak g^{-\alpha},
\]
\emph{where \(\alpha\) and \(-\alpha\) are the roots of \(G\) with
respect to \(T\). Moreover, there exist two group monomorphisms}
\[
  p_\alpha\colon\mathbb G_{a,S}\longrightarrow G,
  \qquad
  p_{-\alpha}\colon\mathbb G_{a,S}\longrightarrow G
\]
\emph{such that}
\[
  t\,p_\alpha(x)\,t^{-1}=p_\alpha\bigl(\alpha(t)x\bigr),
  \qquad
  t\,p_{-\alpha}(x)\,t^{-1}
    =p_{-\alpha}\bigl(\alpha(t)^{-1}x\bigr)
\]
\emph{for every \(S'\to S\), \(t\in T(S')\), and
\(x\in\mathbb G_a(S')\), and such that the morphism}
\[
  \mathbb G_{a,S}\times_S T\times_S\mathbb G_{a,S}
    \longrightarrow G,
  \qquad
  (y,t,x)\longmapsto p_{-\alpha}(y)\,t\,p_\alpha(x),
\]
\emph{is radicial and dominant \textup{(\emph{Bible}, Section 13.4,
Corollary 2 to Theorem 3)}.}

Since its tangent map at the identity is bijective, this morphism is
also separable, hence birational. Zariski's Main Theorem
\textup{(EGA III\(_1\), 4.4.9)} therefore shows that it is an open
immersion.

\paragraph{\textbf{Lemma 1.2.}}
\SGARefTarget{sga3:xx:1.2}
\emph{Let \(S\) be a scheme, \(G\) an \(S\)-group scheme, \(T\) a
torus of \(G\), \(Q\) a subtorus of \(T\), and \(\alpha\) a
character of \(T\) whose restriction to \(Q\) is nontrivial on every
fiber. Let
\[
  p_\alpha\colon\mathbb G_{a,S}\longrightarrow G
  \quad\textup{and}\quad
  p_{-\alpha}\colon\mathbb G_{a,S}\longrightarrow G
\]
be group morphisms normalized by \(T\), with multipliers \(\alpha\)
and \(-\alpha\), respectively. Suppose that the morphism
\[
  u\colon\mathbb G_{a,S}\times_S T\times_S\mathbb G_{a,S}
    \longrightarrow G,
  \qquad
  u(y,t,x)=p_{-\alpha}(y)\,t\,p_\alpha(x),
\]
defined on points by the displayed formula, is an open immersion.
Finally, let \(q\geq 0\) be an integer and let
\[
  p\colon\mathbb G_{a,S}\longrightarrow G
\]
be a group morphism such that, for every \(S'\to S\),
\(t\in Q(S')\), and \(x\in\mathbb G_a(S')\),
\[
  \intaut(t)^q\bigl(p(x)\bigr)=p\bigl(\alpha(t)x\bigr).
\]
Then there is a unique \(\nu\in\mathbb G_a(S)\) such that
\[
  p(x)=p_\alpha(\nu x^q).
\]}

Indeed, let \(\Omega\) be the image of \(u\), and put
\(U=p^{-1}(\Omega)\). This is an open subset of
\(\mathbb G_{a,S}\) containing the zero section. For every section
\(t\) of \(Q\), multiplication by \(\alpha(t)\) leaves \(U\)
globally invariant. We have \(U=\mathbb G_{a,S}\). It is enough to
check this when \(S\) is the spectrum of an algebraically closed
field \(k\); then
\[
  \alpha\colon Q(k)\longrightarrow k^\times
\]
is surjective, so \(U(k)\supset k^\times\), and hence
\(U=\mathbb G_{a,k}\).

There are consequently three morphisms
\[
  a\colon\mathbb G_{a,S}\longrightarrow\mathbb G_{a,S},
  \qquad
  b\colon\mathbb G_{a,S}\longrightarrow T,
  \qquad
  c\colon\mathbb G_{a,S}\longrightarrow\mathbb G_{a,S}
\]
such that
\[
  p(x)=p_{-\alpha}\bigl(a(x)\bigr)\,b(x)\,
       p_\alpha\bigl(c(x)\bigr).
\]
The condition on \(p\) becomes
\[
  a\bigl(\alpha(t)x\bigr)=\alpha(t)^{-q}a(x),\qquad
  b\bigl(\alpha(t)x\bigr)=b(x),\qquad
  c\bigl(\alpha(t)x\bigr)=\alpha(t)^q c(x).
\]
For the same reason, for every \(S'\to S\) and
\(z\in\mathbb G_m(S')\),
\[
  a(zx)=z^{-q}a(x),\qquad b(zx)=b(x),\qquad c(zx)=z^q c(x),
\]
and therefore
\[
  z^q a(z)=a(1),\qquad b(z)=b(1),\qquad c(z)=z^q c(1).
\]
Since \(\mathbb G_{m,S}\) is schematically dense in
\(\mathbb G_{a,S}\), for every \(S'\to S\) and
\(x\in\mathbb G_a(S')\) we obtain
\[
\begin{aligned}
  x^q a(x)&=a(1)=a(0)=0, &&\text{hence }a=0,\\
  c(x)&=x^q c(1)=\nu x^q,
    &&\text{for some }\nu\in\mathbb G_a(S),\\
  b(x)&=b(1)=b(0)=e, &&\text{hence }b=e.
\end{aligned}
\]
This completes the proof.

\paragraph{\textbf{Definition 1.3.}}
\SGARefTarget{sga3:xx:1.3}
\emph{An \(S\)-elementary system is a triple \((G,T,\alpha)\) in
which}
\begin{enumerate}
\item[\textup{(i)}] \emph{\(G\) is a reductive \(S\)-group of
  semisimple rank one \textup{(\SGARefLink{sga3:xix:definition:2:7}{Exposé XIX, 2.7})};}
\item[\textup{(ii)}] \emph{\(T\) is a maximal torus of \(G\);}
\item[\textup{(iii)}] \emph{\(\alpha\) is a root of \(G\) with
  respect to \(T\) \textup{(\SGARefLink{sga3:xix:definition:3-2}{Exposé XIX, 3.2})}.}
\end{enumerate}
Thus there is a direct-sum decomposition
\textup{(\SGARefLink{sga3:xix:lemma:3:5}{Exposé XIX, 3.5})}
\[
  \mathfrak g=\mathfrak t\oplus\mathfrak g^\alpha
  \oplus\mathfrak g^{-\alpha},
\]
where \(\mathfrak g^\alpha\) and \(\mathfrak g^{-\alpha}\) are
locally free of rank one.\footnote{As
\(\mathcal O_S\)-modules.}

\paragraph{\textbf{1.4.}}
\SGARefTarget{sga3:xx:1.4}
If \((G,T,\alpha)\) is an \(S\)-elementary system, then
\((G_{S'},T_{S'},\alpha_{S'})\) is an \(S'\)-elementary system for
every \(S'\to S\). If \((G,T,\alpha)\) is an \(S\)-elementary
system, then so is \((G,T,-\alpha)\).

If \(S\) is a scheme, \(G\) a reductive \(S\)-group, \(T\) a
maximal torus of \(G\), and \(\alpha\) a root of \(G\) with respect
to \(T\), then \((Z_\alpha,T,\alpha)\) is an \(S\)-elementary
system \textup{(\SGARefLink{sga3:xix:paragraph:3-9}{Exposé XIX, 3.9})}.

Let \((G,T,\alpha)\) be an \(S\)-elementary system. The invertible
module \(\mathfrak g^\alpha\) is canonically endowed with a
\(T\)-module structure. Hence the vector bundle
\(W(\mathfrak g^\alpha)\) is also a \(T\)-module. On the other hand,
the inner automorphisms by elements of \(T\) make \(G\) a group
equipped with a group of operators \(T\).

% END INCLUDED SGA3_English_current_buildable_source_20260823/source/inputs/XX/tex/components/00_expose_XX_recollection11_through_14_en.tex

% BEGIN INCLUDED SGA3_English_current_buildable_source_20260823/source/inputs/XX/tex/components/01_expose_XX_theorem15_through_lemma18_en.tex
% Authority: Exp20-13oct24.pdf, local pp. 3--5 / combined pp. 1062--1064.
% Continuous English stops immediately before the proof of Lemma 1.7.

\paragraph{\textbf{Theorem 1.5.}}
\SGARefTarget{sga3:xx:1.5}
\emph{Let \((G,T,\alpha)\) be an \(S\)-elementary system.}
\begin{enumerate}
\item[\textup{\SGAInlineRefTarget{sga3:xx:1.5:item:i}(i)}] \emph{There exists a unique morphism of groups
  with operator group \(T\)
  \[
    \exp\colon W(\mathfrak g^\alpha)\longrightarrow G
  \]
  whose induced morphism on Lie algebras is the canonical map
  \(\mathfrak g^\alpha\to\mathfrak g\).\footnote{It will be seen
  below, in \SGARefLink{sga3:xx:5.9}{Corollary 5.9}, that \(\exp\) is an isomorphism from
  \(W(\mathfrak g^\alpha)\) onto a closed subgroup of \(G\).}
  Equivalently, \(\exp\) is the unique morphism such that, for every
  \(S'\to S\), \(t\in T(S')\), and
  \(X,X'\in W(\mathfrak g^\alpha)(S')\),
  \[
  \begin{aligned}
    \exp(X+X')&=\exp(X)\exp(X'),\\
    \intaut(t)\bigl(\exp(X)\bigr)&=\exp\bigl(\alpha(t)X\bigr),\\
    \Lie(\exp)(X)&=X.
  \end{aligned}
  \]}
\item[\textup{\SGAInlineRefTarget{sga3:xx:1.5:item:ii}(ii)}] \emph{Define analogously, in the
  \(S\)-elementary system \((G,T,-\alpha)\),
  \[
    \exp\colon W(\mathfrak g^{-\alpha})\longrightarrow G.
  \]
  Then the morphism
  \[
    W(\mathfrak g^{-\alpha})\times_S T
      \times_S W(\mathfrak g^\alpha)\longrightarrow G,
    \qquad
    (Y,t,X)\longmapsto \exp(Y)\,t\,\exp(X),
  \]
  defined on points by the displayed formula, is an open
  immersion.}
\end{enumerate}

Suppose for the moment that the required exponential morphisms
exist, and prove the other assertions. We first prove \textup{\SGARefLink{sga3:xx:1.5:item:ii}{(ii)}}.
Both sides are of finite presentation and flat over \(S\), so it is
enough to work when \(S\) is the spectrum of an algebraically closed
field \textup{(SGA 1, I 5.7 and VIII 5.5)}. Put \(S=\Spec(k)\), and
choose
\[
  X\in\Gamma(S,\mathfrak g^\alpha)^\times,
  \qquad
  Y\in\Gamma(S,\mathfrak g^{-\alpha})^\times.
\]
It is enough to prove that
\[
  \mathbb G_{a,k}\times_k T\times_k\mathbb G_{a,k}
    \longrightarrow G,
  \qquad
  (y,t,x)\longmapsto \exp(yY)\,t\,\exp(xX),
\]
is an open immersion. By \SGARefLink{sga3:xx:section:1}{1.1} and \SGARefLink{sga3:xx:1.2}{1.2}, there exist \(a,b\in k\)
such that
\[
  \exp(yY)=p_{-\alpha}(ay),
  \qquad
  \exp(xX)=p_\alpha(bx).
\]
Since
\(\exp\colon W(\mathfrak g^{-\alpha})\to G\) induces a monomorphism
on Lie algebras, \(a\neq0\); likewise \(b\neq0\). The assertion is
therefore reduced to \SGARefLink{sga3:xx:section:1}{1.1}.

The uniqueness of \(\exp\) may be proved locally on \(S\). We reduce
to the case in which \(\mathfrak g^\alpha\) and
\(\mathfrak g^{-\alpha}\) are free, and then apply \SGARefLink{sga3:xx:1.2}{1.2} with
\(Q=T\) and \(q=1\).

It remains to prove existence. Faithfully flat descent, together
with the uniqueness just proved, shows that it is enough to prove
existence locally on \(S\) for the fpqc topology. The usual
finite-presentation arguments reduce first to \(S\) noetherian and
then to \(S\) noetherian local. We may therefore assume that
\[
  S=\Spec(A),
\]
where \(A\) is a complete noetherian local ring with algebraically
closed residue field \(k\). Let
\[
  p_0\colon\mathbb G_{a,k}\longrightarrow G_k
\]
be a monomorphism of \(k\)-groups normalized by \(T_k\), with
multiplier
\(\alpha_0=\alpha\otimes_A k\); such a morphism exists by \SGARefLink{sga3:xx:section:1}{1.1}.
By \SGARefLink{sga3:xx:section:1}{1.1} and \SGARefLink{sga3:xx:1.2}{1.2}, the associated morphism
\[
  T_k\mathbin{\cdot_{\alpha_0}}\mathbb G_{a,k}
    \longrightarrow G_k
\]
is an immersion and hence, in particular, a monomorphism. We
provisionally admit the following two lemmas.

\paragraph{\textbf{Lemma 1.6.}}
\SGARefTarget{sga3:xx:1.6}
\emph{Let \(S\) be a scheme, \(G\) an \(S\)-group of finite
presentation, \(T\) an \(S\)-torus, \(\alpha\) a character of \(T\)
that is nontrivial on every fiber, and \(s_0\in S\). Let
\[
  f\colon T\mathbin{\cdot_\alpha}\mathbb G_{a,S}\longrightarrow G
\]
be an \(S\)-group morphism such that \(f_{s_0}\) is a monomorphism
and the restriction of \(f\) to \(T\) is a monomorphism. There
exists an open neighborhood \(U\) of \(s_0\) such that
\(f|_U\) is a monomorphism.}

\paragraph{\textbf{Lemma 1.7.}}
\SGARefTarget{sga3:xx:1.7}
\emph{Let \(A\) be a complete noetherian local ring with
algebraically closed residue field \(k\), let
\((G,T,\alpha)\) be an \(A\)-elementary system, and let
\[
  p_0\colon\mathbb G_{a,k}\longrightarrow G_k
\]
be a morphism of \(k\)-groups normalized by \(T_k\), with
multiplier \(\alpha\otimes_A k\). There exists a group morphism
\[
  p\colon\mathbb G_{a,A}\longrightarrow G
\]
normalized by \(T\), with multiplier \(\alpha\).}

Let \(p\) be the morphism asserted by \SGARefLink{sga3:xx:1.7}{1.7}, and let
\[
  f\colon T\mathbin{\cdot_\alpha}\mathbb G_{a,S}\longrightarrow G
\]
be the corresponding morphism. It satisfies the hypotheses of \SGARefLink{sga3:xx:1.6}{1.6},
so it is a monomorphism; in particular, \(p\) is a monomorphism.
The theorem now follows from \SGARefLink{sga3:xix:corollary:4:9}{Exposé XIX, 4.9}.

\medskip
\noindent\emph{Proof of \SGARefLink{sga3:xx:1.6}{1.6}.}
Let
\[
  \epsilon\colon S\longrightarrow
    T\mathbin{\cdot_\alpha}\mathbb G_{a,S}
\]
be the identity section. Since \(f\) is unramified at
\(\epsilon(s_0)\), it is unramified at \(\epsilon(s)\) for all
\(s\) in an open neighborhood \(U\) of \(s_0\). Hence \(f|_U\) is
unramified \textup{(\SGARefLink{sga3:x:lemma:3:5}{Exposé X, 3.5})},\footnote{See also Exposé
VI\(_B\), \SGARefLink{sga3:xx:2.5}{2.5}.} and its kernel \(\Ker(f)_U\) is unramified over
\(U\). To prove that \(f|_U\) is a monomorphism, it is therefore
enough\footnote{By EGA IV\(_4\), 17.9.1.} to prove that
\(\Ker(f)_U\) is radicial over \(U\), which is a set-theoretic
question. We are reduced to the following assertion.

\paragraph{\textbf{Lemma 1.8.}}
\SGARefTarget{sga3:xx:1.8}
\emph{Let \(k\) be an algebraically closed field, and let \(N\) be
an invariant subgroup of
\(T\mathbin{\cdot_\alpha}\mathbb G_{a,k}\), where \(\alpha\) is a
nontrivial character of the torus \(T\). Suppose that \(N\) is
étale over \(k\) and that \(N\cap T=\{e\}\). Then \(N=\{e\}\).}

We have
\[
  \intaut(t')(t,x)=\bigl(t,\alpha(t')x\bigr).
\]
If \((t,x)\) is a point of \(N\) with \(x\neq0\), then
\((t,zx)\) is also a point of \(N\) for every \(z\in k^\times\).
Thus \((t,x)\) is not isolated, contradicting the quasi-finiteness
of \(N\). Set-theoretically, therefore, \(N\subset T\), and the
hypothesis \(N\cap T=\{e\}\) completes the proof.

% END INCLUDED SGA3_English_current_buildable_source_20260823/source/inputs/XX/tex/components/01_expose_XX_theorem15_through_lemma18_en.tex

% BEGIN INCLUDED SGA3_English_current_buildable_source_20260823/source/inputs/XX/tex/components/02_expose_XX_proof17_through_proof110_en.tex
% Authority: Exp20-13oct24.pdf, local pp. 5--6 / combined pp. 1064--1065.
% Continuous English stops immediately before 1.12, the proof of Lemma 1.9.

\medskip
\noindent\emph{Proof of \SGARefLink{sga3:xx:1.7}{1.7}.}
Let \(\mathfrak m\) be the radical of \(A\), and put
\[
  S_n=\Spec(A/\mathfrak m^{n+1}),\qquad n\geq 0.
\]
We first show, by induction on \(n\), that \(p_0\) can be extended,
for every \(n\), to a morphism of \(S_n\)-groups
\[
  p_n\colon\mathbb G_{a,S_n}\longrightarrow G_{S_n},
\]
normalized by \(T_{S_n}\) with multiplier \(\alpha_{S_n}\), in such
a way that
\[
  p_{n+1}\times_{S_{n+1}}S_n=p_n.
\]

Let \(H=T\mathbin{\cdot_\alpha}\mathbb G_{a,S}\). Denote by \(f_n\)
the morphism \(H_{S_n}\to G_{S_n}\) defined by \(p_n\). We
provisionally admit the following lemma.

\paragraph{\textbf{Lemma 1.9.}}
\SGARefTarget{sga3:xx:1.9}
\emph{Let \((G,T,\alpha)\) be a \(k\)-elementary system, where \(k\)
is algebraically closed, and let
\[
  p\colon\mathbb G_{a,k}\longrightarrow G
\]
be a monomorphism normalized by \(T\), with multiplier \(\alpha\).
Then}
\[
  \mathrm H^2\bigl(T\mathbin{\cdot_\alpha}\mathbb G_{a,k},
    \mathfrak g\bigr)=0.
\]
\emph{Here \(T\mathbin{\cdot_\alpha}\mathbb G_{a,k}\) acts on
\(\mathfrak g\) through the morphism
\(T\mathbin{\cdot_\alpha}\mathbb G_{a,k}\to G\) defined by \(p\)
and through the adjoint representation of \(G\).}

By \SGARefLink{sga3:iii:corollary:2-8}{Exposé III, 2.8}, \(f_n\) therefore extends to a morphism of
\(S_{n+1}\)-groups
\[
  f'_{n+1}\colon H_{S_{n+1}}\longrightarrow G_{S_{n+1}}.
\]
The morphism \(f'_{n+1}\) and the canonical immersion of
\(T_{S_{n+1}}\) into \(G_{S_{n+1}}\) have the same restriction to
\(T_{S_n}\). By \SGARefLink{sga3:iii:corollary:2-5}{Exposé III, 2.5}, there is an element
\(g\in G(S_{n+1})\) such that
\[
  g\times_{S_{n+1}}S_n=e
\]
and such that
\[
  f_{n+1}=\intaut(g)\circ f'_{n+1}
\]
restricts to \(T_{S_{n+1}}\) by the canonical immersion. Let
\(p_{n+1}\) be the restriction of \(f_{n+1}\) to
\(\mathbb G_{a,S_{n+1}}\). It is normalized by \(T_{S_{n+1}}\),
has multiplier \(\alpha_{S_{n+1}}\), and extends \(p_n\).

We have thus constructed a coherent system \((f_n)\); it remains to
algebraize it. For this we use:

\paragraph{\textbf{Lemma 1.10.}}
\SGARefTarget{sga3:xx:1.10}
\emph{Let \(A\) be a complete noetherian local ring,
\(\mathfrak m\) its maximal ideal,
\[
  S=\Spec(A),\qquad S_n=\Spec(A/\mathfrak m^{n+1}),
\]
let \(T\) be an \(S\)-torus, let \(\alpha\) be a nonzero character
of \(T\), and let \(X\) be an affine \(S\)-scheme on which \(T\)
acts. Let \(T\) act on \(\mathbb G_{a,S}\) through \(\alpha\).
Let \(q\geq0\), and let \((f_n)_{n\geq0}\) be a coherent system of
morphisms}
\[
  f_n\colon\mathbb G_{a,S_n}^{\,q}\longrightarrow X_{S_n}
\]
\emph{of objects with operators \(T_{S_n}\). There exists a unique
morphism of objects with operators \(T\)}
\[
  f\colon\mathbb G_{a,S}^{\,q}\longrightarrow X
\]
\emph{that induces the \(f_n\) \textup{(compare \SGARefLink{sga3:ix:theorem:7:1}{Exposé IX, 7.1})}.}

\paragraph{\textbf{Corollary 1.11.}}
\SGARefTarget{sga3:xx:1.11}
\emph{If \(X\) is a group with operator group \(T\), and if the
\(f_n\) are group morphisms, then \(f\) is also a group morphism.}

Indeed, apply the uniqueness assertion of the lemma to the two
morphisms
\(\mathbb G_{a,S}^{\,2q}\to X\) deduced from \(f\) in the usual
way.

\medskip
\noindent\emph{Proof of \SGARefLink{sga3:xx:1.10}{1.10}.}
Suppose that \(T\) is split, as it is in the application of \SGARefLink{sga3:xx:1.10}{1.10} to
the proof of \SGARefLink{sga3:xx:1.5}{1.5}. By \SGARefLink{sga3:i:corollary:4-7-3-1}{Exposé I, 4.7.3.1}, the functor
\[
  X\longmapsto\mathscr A(X)
\]
defines an equivalence between the category of affine \(S\)-schemes
equipped with an action of \(T\) and the category opposite to that
of \(M\)-graded \(S\)-algebras, where
\[
  M=\operatorname{Hom}_{S\text{-gr.}}(T,\mathbb G_{m,S}).
\]
Thus there are gradings
\[
  B=\mathscr A(X)=\bigoplus_{m\in M}B_m,
  \qquad
  C=\mathscr A(\mathbb G_{a,S}^{\,q})
    =\bigoplus_{m\in M}C_m.
\]
Each \(C_m\) is visibly a finite free \(A\)-module. Indeed,
\(C_m=0\) if \(m\) is not a multiple of \(\alpha\), whereas, if
\(m=d\alpha\), then \(C_m\) is isomorphic to the \(A\)-module of
homogeneous polynomials of degree \(d\) in \(q\) variables. Put
\[
\begin{aligned}
  \widehat B_m
    &=\varprojlim_n B_m\otimes_A(A/\mathfrak m^{n+1}),\\
  \widehat C_m
    &=\varprojlim_n C_m\otimes_A(A/\mathfrak m^{n+1}),\\
  \widehat B&=\bigoplus_{m\in M}\widehat B_m,
  &\widehat C&=\bigoplus_{m\in M}\widehat C_m.
\end{aligned}
\]
There are canonical morphisms of \(M\)-graded algebras
\[
  g_B\colon B\longrightarrow\widehat B,
  \qquad
  g_C\colon C\longrightarrow\widehat C.
\]
The preceding observation shows that \(g_C\) is an isomorphism.
Giving a coherent system \((f_n)\) as in the statement is
equivalent to giving a morphism of graded \(A\)-algebras
\[
  \widehat F\colon\widehat B\longrightarrow\widehat C.
\]
Finding a morphism \(f\) as in the statement is equivalent to
finding a morphism of graded \(A\)-algebras \(F\colon B\to C\)
making the diagram
\[
\begin{tikzcd}[column sep=large,row sep=large]
B \arrow[r,"F"] \arrow[d,"g_B"']
  & C \arrow[d,"g_C"']\\
\widehat B \arrow[r,"\widehat F"]
  & \widehat C
\end{tikzcd}
\]
commute. Since \(g_C\) is an isomorphism, the existence and
uniqueness of \(F\) are immediate. This proves \SGARefLink{sga3:xx:1.10}{1.10}.

To complete the proof of \SGARefLink{sga3:xx:1.5}{1.5}, it remains only to prove \SGARefLink{sga3:xx:1.9}{1.9}.

% END INCLUDED SGA3_English_current_buildable_source_20260823/source/inputs/XX/tex/components/02_expose_XX_proof17_through_proof110_en.tex

% BEGIN INCLUDED SGA3_English_current_buildable_source_20260823/source/inputs/XX/tex/components/03_expose_XX_par112_through_cor121_en.tex
% Authority: Exp20-13oct24.pdf, local pp. 6--9 / combined pp. 1065--1068.
% Continuous English stops immediately before Section 2 and Theorem 2.1.

\paragraph{\textbf{1.12. Proof of \SGARefLink{sga3:xx:1.9}{1.9}.}}\SGAInlineRefTarget{sga3:xx:paragraph:1-12}
We have
\[
  \mathfrak g
    =\mathfrak t\oplus\mathfrak g^\alpha\oplus\mathfrak g^{-\alpha}.
\]
As explained in \SGARefLink{sga3:xx:1.9}{1.9}, regard \(\mathfrak g\) as a
\(T\mathbin{\cdot_\alpha}\mathbb G_{a,k}\)-module. Clearly
\(\mathfrak t\oplus\mathfrak g^\alpha\) is a submodule of
\(\mathfrak g\), and the quotient is isomorphic to
\(\mathfrak g^{-\alpha}\) as a \(k\)-vector space and even as a
\(T\)-module. The group \(\mathbb G_{a,k}\) acts trivially on this
one-dimensional quotient, since every group morphism
\(\mathbb G_{a,k}\to\mathbb G_{m,k}\) is trivial. Likewise,
\(\mathfrak g^\alpha\) is a submodule of
\(\mathfrak t\oplus\mathfrak g^\alpha\); the quotient is isomorphic
to \(\mathfrak t\) as a \(T\)-module, and
\(\mathbb G_{a,k}\) acts trivially on it. In summary:

\paragraph{\textbf{Lemma 1.13.}}
\SGARefTarget{sga3:xx:1.13}
\emph{Under the hypotheses of \SGARefLink{sga3:xx:1.9}{1.9}, \(\mathfrak g\), as a
\(T\mathbin{\cdot_\alpha}\mathbb G_{a,k}\)-module, has a
composition series whose successive quotients are}
\[
  \mathfrak g^{-\alpha},\qquad \mathfrak t,\qquad
  \mathfrak g^\alpha,
\]
\emph{regarded as
\(T\mathbin{\cdot_\alpha}\mathbb G_{a,k}\)-modules through the
projection
\(T\mathbin{\cdot_\alpha}\mathbb G_{a,k}\to T\).}

We are thus reduced to computing the cohomology of
\(T\mathbin{\cdot_\alpha}\mathbb G_{a,k}\) acting on \(W(k)\)
through the projection
\(T\mathbin{\cdot_\alpha}\mathbb G_{a,k}\to T\) and a character
\(\beta\) of \(T\) (here \(\beta=0,\alpha\), or \(-\alpha\)).
\footnote{The editors added the sentence that follows.}
Write \(k[x_1,\ldots,x_n]\) for the polynomial algebra over \(k\)
in \(n\) variables, and \(k_q[x_1,\ldots,x_n]\) for its subspace
of homogeneous polynomials of degree \(q\).

\paragraph{\textbf{Lemma 1.14.}}
\SGARefTarget{sga3:xx:1.14}
\emph{With the preceding notation,}
\[
  \mathrm H^n\bigl(T\mathbin{\cdot_\alpha}\mathbb G_{a,k},k\bigr)
    =\mathrm H^n(C^\ast_{\alpha,\beta}),
\]
\emph{where the complex \(C^\ast_{\alpha,\beta}\) is defined by}
\[
  C^n_{\alpha,\beta}=
  \begin{cases}
    k_q[x_1,\ldots,x_n],
      &\text{if \(\beta=q\alpha\), with \(q\in\mathbb N^\ast\),}\\
    0,&\text{otherwise,}
  \end{cases}
\]
\emph{and}
\[
\begin{aligned}
(\delta f)(x_1,x_2,\ldots,x_{n+1})
  ={}&f(x_2,\ldots,x_{n+1})\\
   &+\sum_{i=1}^{n}(-1)^i
       f(x_1,\ldots,x_i+x_{i+1},\ldots,x_{n+1})\\
   &+(-1)^{n+1}f(x_1,\ldots,x_n).
\end{aligned}
\]

Indeed, by \SGARefLink{sga3:i:proposition:5-3-2}{Exposé I, 5.3.2}, the functor
\(M\mapsto\mathrm H^0(T,M)\) is exact on the category of
\(T\)-modules, and the \(\mathrm H^q(T,-)\) vanish. It follows, as
in the usual theory of group cohomology, that
\(\mathrm H^n(T\mathbin{\cdot_\alpha}\mathbb G_{a,k},k)\) may be
computed as the \(n\)-th cohomology group of the complex of
cochains of \(\mathbb G_{a,k}\) in \(k\) that are invariant under
\(T\), that is, those satisfying
\[
  f\bigl(\alpha(t)x_1,\ldots,\alpha(t)x_n\bigr)
    =\beta(t)f(x_1,\ldots,x_n).
\]
This is precisely the complex just displayed.

To prove \SGARefLink{sga3:xx:1.9}{1.9}, it is therefore enough to show that
\(\mathrm H^2(C^\ast_{\alpha,\beta})=0\) for
\(\beta=0,\alpha,-\alpha\), which is immediate.

\paragraph{\textbf{Remark 1.15.}}
\SGARefTarget{sga3:xx:1.15}
The groups \(\mathrm H^n(C^\ast_{\alpha,\beta})\) can be computed
explicitly when \(\beta=q\alpha\); see M.~Lazard,
\emph{Lois de groupes et analyseurs}, Annales E.N.S., 1955. In
particular,
\[
  \mathrm H^n(C^\ast_{\alpha,q\alpha})=0
  \qquad\text{for \(n>q\).}
\]

\paragraph{\textbf{Notations 1.16.}}
\SGARefTarget{sga3:xx:1.16}
The image of the canonical immersion
\[
  W(\mathfrak g^{-\alpha})\times_S T
    \times_S W(\mathfrak g^\alpha)\longrightarrow G
\]
will be denoted by \(\Omega\). It is an open subset of \(G\)
containing the identity section. The images of
\[
\begin{gathered}
  W(\mathfrak g^{-\alpha}),\qquad W(\mathfrak g^\alpha),\qquad
  W(\mathfrak g^{-\alpha})\times_S T,\qquad
  T\times_S W(\mathfrak g^\alpha)
\end{gathered}
\]
will be denoted, respectively,\footnote{The editors replaced
\(P_{-\alpha}\) and \(P_\alpha\) by \(U_{-\alpha}\) and
\(U_\alpha\).}
\[
  U_{-\alpha},\qquad U_\alpha,\qquad
  U_{-\alpha}\mathbin{\cdot}T,\qquad T\mathbin{\cdot}U_\alpha.
\]
Then \(U_\alpha\) (respectively \(U_{-\alpha}\)) is a subgroup of
\(G\), canonically endowed with a vector-bundle structure, and
\[
  \intaut(t)(x)=x^{\alpha(t)}
  \qquad
  \bigl(\text{respectively }x^{-\alpha(t)}\bigr)
\]
for every \(S'\to S\), \(t\in T(S')\), and
\(x\in U_\alpha(S')\) (respectively \(x\in U_{-\alpha}(S')\)).

There are canonical isomorphisms
\[
  T\mathbin{\cdot}U_\alpha\simeq T\mathbin{\cdot_\alpha}U_\alpha,
  \qquad
  T\mathbin{\cdot}U_{-\alpha}
    \simeq T\mathbin{\cdot_{-\alpha}}U_{-\alpha}.
\]
The open subset \(\Omega\) is stable under \(\intaut(T)\), and
\[
  \intaut(t')\bigl(y\mathbin{\cdot}t\mathbin{\cdot}x\bigr)
    =y^{-\alpha(t')}\mathbin{\cdot}t
       \mathbin{\cdot}x^{\alpha(t')}.
\]

\paragraph{\textbf{Corollary 1.17.}}
\SGARefTarget{sga3:xx:1.17}
\emph{We have}
\[
  \Lie(U_\alpha/S)=\mathfrak g^\alpha,\qquad
  \Lie(U_{-\alpha}/S)=\mathfrak g^{-\alpha}.
\]
\emph{The isomorphisms}
\[
  W(\mathfrak g^\alpha)\xrightarrow[\sim]{\exp}U_\alpha,
  \qquad
  W(\mathfrak g^{-\alpha})\xrightarrow[\sim]{\exp}U_{-\alpha}
\]
\emph{are those of \SGARefLink{sga3:xix:lemma:4:2}{Exposé XIX, 4.2}.}

\paragraph{\textbf{Corollary 1.18.}}
\SGARefTarget{sga3:xx:1.18}
\emph{The open subset \(\Omega\) is relatively schematically dense
in \(G\) \textup{(cf. Exposé XVIII, \S1)}.}

This follows immediately from Exposé XVIII, 1.3.\footnote{See also
EGA IV\(_3\), 11.10.10.}

\paragraph{\textbf{Corollary 1.19.}}
\SGARefTarget{sga3:xx:1.19}
\emph{The center of \(G\) is}
\[
  \operatorname{Centr}(G)=\Ker(\alpha).
\]
\emph{It is therefore a closed subgroup of \(G\), of
multiplicative type and of finite type.}

The second assertion follows from the first by \SGARefLink{sga3:ix:proposition:2:7}{Exposé IX, 2.7}. Let
us prove the first. The inner automorphism defined by a section of
\(\Ker(\alpha)\) acts trivially on \(\Omega\), by the last formula
of \SGARefLink{sga3:xx:1.16}{1.16}, and hence on \(G\), by \SGARefLink{sga3:xx:1.18}{1.18}. Conversely, if
\(g\in G(S)\) centralizes \(G\), then it centralizes \(T\) and
\(U_\alpha\). It is therefore a section of \(T\), by \SGARefLink{sga3:xix:paragraph:2:8}{Exposé XIX,
2.8}, and annihilates \(\alpha\). Since the same argument applies
after every base change, we indeed have
\(\operatorname{Centr}(G)=\Ker(\alpha)\).

\paragraph{\textbf{Corollary 1.20.}}
\SGARefTarget{sga3:xx:1.20}
\emph{A monomorphism
\[
  p_\alpha\colon\mathbb G_{a,S}\longrightarrow G
\]
normalized by \(T\), with multiplier \(\alpha\), exists if and
only if the \(\mathscr O_S\)-module \(\mathfrak g^\alpha\) is
free. More precisely, the assignments}
\[
  X_\alpha\longmapsto
    \bigl(x\longmapsto\exp(xX_\alpha)\bigr),
  \qquad
  p_\alpha\longmapsto\Lie(p_\alpha)
\]
\emph{give a bijection between
\(\Gamma(S,\mathfrak g^\alpha)^\times\) and the set of such
monomorphisms \(p_\alpha\), which is also the set of isomorphisms
of vector group schemes}
\[
  \mathbb G_{a,S}\xrightarrow{\sim}U_\alpha.
\]
\footnote{Indeed,
\(\Lie(\mathbb G_{a,S})=\mathscr O_S\), and
\(\Lie(p_\alpha)\) is an element of
\(\operatorname{Hom}_{\mathscr O_S}
  (\mathscr O_S,\mathfrak g^\alpha)
=\Gamma(S,\mathfrak g^\alpha)\).}

\paragraph{\textbf{Corollary 1.21.}}
\SGARefTarget{sga3:xx:1.21}
\emph{The subgroups \(U_\alpha\) and \(U_{-\alpha}\) of \(G\)
commute on no fiber.}

Indeed, if \((U_\alpha)_s\) and \((U_{-\alpha})_s\) commute, then
\(\Omega_s\) is a subgroup of \(G_s\), and hence
\(\Omega_s=G_s\).\footnote{By \SGARefLink{sga3:xx:1.18}{1.18}.} It follows that \(G_s\) is
solvable, contrary to the hypothesis that \(G_s\) is reductive of
semisimple rank \(1\).

% END INCLUDED SGA3_English_current_buildable_source_20260823/source/inputs/XX/tex/components/03_expose_XX_par112_through_cor121_en.tex

% BEGIN INCLUDED SGA3_English_current_buildable_source_20260823/source/inputs/XX/tex/components/04_expose_XX_theorem21_prop22_partA_en.tex
% Authority: Exp20-13oct24.pdf, local pp. 9--11 / combined pp. 1068--1070.
% Continuous English stops immediately before part B of the proof of 2.2.

\section{Structure of elementary systems}
\SGARefTarget{sga3:xx:section:2}

\paragraph{\textbf{Theorem 2.1.}}
\SGARefTarget{sga3:xx:2.1}
\emph{Let \(S\) be a scheme and \((G,T,\alpha)\) an
\(S\)-elementary system. There exist a morphism of
\(\mathscr O_S\)-modules}
\[
  \mathfrak g^\alpha\otimes_{\mathscr O_S}\mathfrak g^{-\alpha}
    \longrightarrow\mathscr O_S,
  \qquad
  (X,Y)\longmapsto\langle X,Y\rangle,
\]
\emph{and a morphism of \(S\)-groups}
\[
  \alpha^\ast\colon\mathbb G_{m,S}\longrightarrow T
\]
\emph{such that, for every \(S'\to S\) and all}
\[
  X\in\Gamma\bigl(S',
    \mathfrak g^\alpha\otimes_{\mathscr O_S}\mathscr O_{S'}\bigr),
  \qquad
  Y\in\Gamma\bigl(S',
    \mathfrak g^{-\alpha}\otimes_{\mathscr O_S}\mathscr O_{S'}\bigr),
\]
\emph{one has}
\[
  \exp(X)\mathbin{\cdot}\exp(Y)\in\Omega(S')
  \quad\Longleftrightarrow\quad
  1+\langle X,Y\rangle\in\mathbb G_m(S').
\]
\emph{Under these conditions,}
\SGARefTarget{sga3:xx:2.1:F}
\begin{equation}\tag{F}
\begin{aligned}
  \exp(X)\mathbin{\cdot}\exp(Y)
   ={}&\exp\left(\frac{Y}{1+\langle X,Y\rangle}\right)
       \alpha^\ast\bigl(1+\langle X,Y\rangle\bigr)\\
     &{}\exp\left(\frac{X}{1+\langle X,Y\rangle}\right).
\end{aligned}
\end{equation}
\emph{Moreover, the morphisms
\((X,Y)\mapsto\langle X,Y\rangle\) and \(\alpha^\ast\) are uniquely
determined. The first is an isomorphism, and therefore places
\(\mathfrak g^\alpha\) and \(\mathfrak g^{-\alpha}\) in duality;
furthermore,}
\[
  \alpha\circ\alpha^\ast=2,
\]
\emph{where \(2\) denotes squaring on \(\mathbb G_{m,S}\).}

In view of the uniqueness assertions, it is enough to prove the
theorem locally on \(S\). We may therefore assume that
\(\mathfrak g^\alpha\) and \(\mathfrak g^{-\alpha}\) are free on
\(S\). Choose
\[
  X\in\Gamma(S,\mathfrak g^\alpha)^\times,\qquad
  Y\in\Gamma(S,\mathfrak g^{-\alpha})^\times,
\]
and put
\[
  p_\alpha(x)=\exp(xX),\qquad
  p_{-\alpha}(y)=\exp(yY)
\]
for \(x,y\in\mathbb G_a(S')\), \(S'\to S\). By \SGARefLink{sga3:xx:1.5}{1.5} and \SGARefLink{sga3:xx:1.21}{1.21}, it
is enough to prove the following proposition.

\paragraph{\textbf{Proposition 2.2.}}
\SGARefTarget{sga3:xx:2.2}
\emph{Let \(S\) be a scheme, \(G\) an \(S\)-group, \(T\) a torus
of \(G\), and \(\alpha\) a character of \(T\) that is nontrivial
on every fiber. Let}
\[
  p_\alpha\colon\mathbb G_{a,S}\longrightarrow G,\qquad
  p_{-\alpha}\colon\mathbb G_{a,S}\longrightarrow G
\]
\emph{be group monomorphisms normalized by \(T\), with
multipliers \(\alpha\) and \(-\alpha\), respectively. Suppose
that:}
\begin{enumerate}
\item[\textup{(i)}] \emph{the morphism}
  \[
    \mathbb G_{a,S}\times_S T\times_S\mathbb G_{a,S}
      \longrightarrow G,
    \qquad
    (y,t,x)\longmapsto p_{-\alpha}(y)\,t\,p_\alpha(x),
  \]
  \emph{is an open immersion; denote its image by \(\Omega\);}
\item[\textup{(ii)}] \emph{for every \(s\in S\), the groups}
  \[
    (p_\alpha)_s\bigl(\mathbb G_{a,\kappa(s)}\bigr)
    \quad\text{and}\quad
    (p_{-\alpha})_s\bigl(\mathbb G_{a,\kappa(s)}\bigr)
  \]
  \emph{do not commute.}
\end{enumerate}
\emph{Then there exist unique elements}
\[
  a\in\mathbb G_a(S),\qquad
  \alpha^\ast\in
    \operatorname{Hom}_{S\text{-gr.}}(\mathbb G_{m,S},T)
\]
\emph{with the following properties: for every \(S'\to S\) and
all \(x,y\in\mathbb G_a(S')\),}
\[
  p_\alpha(x)p_{-\alpha}(y)\in\Omega(S')
  \quad\Longleftrightarrow\quad
  1+axy\in\mathbb G_m(S'),
\]
\emph{and, under this condition,}
\[
  p_\alpha(x)p_{-\alpha}(y)
   =p_{-\alpha}\left(\frac{y}{1+axy}\right)
      \alpha^\ast(1+axy)
      p_\alpha\left(\frac{x}{1+axy}\right).
\]
\emph{Moreover, \(a\) is invertible, that is,
\(a\in\mathbb G_m(S)\), and \(\alpha\circ\alpha^\ast=2\).}

\medskip
\noindent\emph{Proof.}
\textup{A)} Consider the morphism
\[
  \mathbb G_{a,S}^{\,2}\longrightarrow G,\qquad
  (x,y)\longmapsto p_\alpha(x)p_{-\alpha}(y).
\]
Let \(U\) be the inverse image of \(\Omega\). It is an open subset
of \(\mathbb G_{a,S}^{\,2}\) containing
\(0\times_S\mathbb G_{a,S}\) and
\(\mathbb G_{a,S}\times_S0\). There are therefore unique
morphisms of \(S\)-schemes
\[
  A\colon U\longrightarrow\mathbb G_{a,S},\qquad
  C\colon U\longrightarrow\mathbb G_{a,S},\qquad
  B\colon U\longrightarrow T
\]
satisfying the set-theoretic identity
\[
  p_\alpha(u)p_{-\alpha}(v)
    =p_{-\alpha}\bigl(A(u,v)\bigr)\,
      B(u,v)\,p_\alpha\bigl(C(u,v)\bigr).
\]
Immediately,
\[
\begin{gathered}
  A(0,v)=v,\qquad A(u,0)=0,\qquad
  C(u,0)=u,\qquad C(0,v)=0,\\
  B(u,0)=B(0,v)=e.
\end{gathered}
\]

Let \(S'\) be a separated \(S\)-scheme and let \(t\in T(S')\).
Since \(\Omega_{S'}\) is stable under \(\intaut(t)\), the last
formula of \SGARefLink{sga3:xx:1.16}{1.16} shows that \(U_{S'}\) is stable under the
automorphism
\[
  (x,y)\longmapsto
    \bigl(\alpha(t)x,\alpha(t)^{-1}y\bigr)
\]
of \(\mathbb G_{a,S'}^{\,2}\), and
\[
\begin{aligned}
 A\bigl(\alpha(t)u,\alpha(t)^{-1}v\bigr)
   &=\alpha(t)^{-1}A(u,v),\\
 C\bigl(\alpha(t)u,\alpha(t)^{-1}v\bigr)
   &=\alpha(t)C(u,v),\\
 B\bigl(\alpha(t)u,\alpha(t)^{-1}v\bigr)
   &=B(u,v).
\end{aligned}
\]
Because \(\alpha\) is faithfully flat, it follows that for every
\(S'\to S\) and every \(z\in\mathbb G_m(S')\), the open subset
\(U_{S'}\) is stable under
\((x,y)\mapsto(zx,z^{-1}y)\), and
\[
\begin{aligned}
  A(zu,z^{-1}v)&=z^{-1}A(u,v),\\
  C(zu,z^{-1}v)&=zC(u,v),\\
  B(zu,z^{-1}v)&=B(u,v).
\end{aligned}
\]

Suppose first that \(v\) is invertible. Taking \(z=v\), we see
that if \((u,v)\) is a section of \(U\), then \((uv,1)\) is also
a section of \(U\), and
\[
  A(uv,1)=v^{-1}A(u,v),\qquad
  B(uv,1)=B(u,v).
\]
Let \(V\) be the open subset of \(\mathbb G_{a,S}^{\,2}\) defined
by\footnote{The editors added the condition
``\((1,uv)\in U(S')\)''.}
\[
  (u,v)\in V(S')
  \quad\Longleftrightarrow\quad
  (u,v),\ (uv,1),\ \text{and }(1,uv)\text{ belong to }U(S').
\]

Since \(U\) contains
\(0\times_S\mathbb G_{a,S}\) and
\(\mathbb G_{a,S}\times_S0\), the open subset \(V\) is a
neighborhood of the zero section of \(\mathbb G_{a,S}^{\,2}\).
The preceding calculation shows that, on
\(V\cap(\mathbb G_{a,S}\times_S\mathbb G_{m,S})\), the two
morphisms
\[
  (u,v)\longmapsto A(u,v),\qquad
  (u,v)\longmapsto vA(uv,1)
\]
coincide, as do
\[
  (u,v)\longmapsto B(u,v),\qquad
  (u,v)\longmapsto B(uv,1).
\]
Since
\(\mathbb G_{a,S}\times_S\mathbb G_{m,S}\) is schematically dense
in \(\mathbb G_{a,S}^{\,2}\), these morphisms coincide on \(V\).

We have \(A(0,1)=1\). Hence there is an open subset \(W_1\) of
\(\mathbb G_{a,S}\), containing the zero section, such that
\(A(x,1)\) is invertible for every section \(x\) of \(W_1\). Put
\[
  A(x,1)^{-1}=F(x).
\]
If \((u,v)\in V(S')\)\footnote{Here and below, the editors replaced
\(U(S')\) by \(V(S')\).} and \(uv\in W_1(S')\), then
\[
  A(u,v)=vA(uv,1)=vF(uv)^{-1}.
\]
The same argument applied to \(C\) gives an open subset \(W_2\)
of \(\mathbb G_{a,S}\), containing the zero section, and an
element \(E\in\mathscr O(W_2)^\times\)\footnote{The editors use
\(E\) for the element denoted \(G\) in the original, since \(G\)
already denotes the \(S\)-group under consideration.} such that
\[
  C(u,v)=uC(1,uv)=uE(uv)^{-1}
\]
whenever \((u,v)\in V(S')\) and \(uv\in W_2(S')\).

Consequently, with \(W=W_1\cap W_2\), there are \(S\)-morphisms
\[
\begin{array}{lll}
  F\colon W\to\mathbb G_{m,S},&\qquad&F(0)=1,\\
  H\colon W\to T,&&H(0)=e,\\
  E\colon W\to\mathbb G_{m,S},&&E(0)=1,
\end{array}
\]
such that, if \((u,v)\in V(S')\) and \(uv\in W(S')\), then
\SGARefTarget{sga3:xx:2.2:plus}
\begin{equation}\tag{+}
  p_\alpha(u)p_{-\alpha}(v)
    =p_{-\alpha}\bigl(vF(uv)^{-1}\bigr)\,
       H(uv)\,p_\alpha\bigl(uE(uv)^{-1}\bigr).
\end{equation}

% END INCLUDED SGA3_English_current_buildable_source_20260823/source/inputs/XX/tex/components/04_expose_XX_theorem21_prop22_partA_en.tex

% BEGIN INCLUDED SGA3_English_current_buildable_source_20260823/source/inputs/XX/tex/components/05_expose_XX_prop22_partsBC_lemma23_en.tex
% Authority: Exp20-13oct24.pdf, local pp. 11--13 / combined pp. 1070--1072.
% Continuous English stops after the proof of Lemma 2.3, before part D
% of the proof of Proposition 2.2.

\textup{B)} We now use the associativity of \(G\) to write
\[
  p_\alpha(u)p_{-\alpha}(v)p_{-\alpha}(w)
    =p_\alpha(u)p_{-\alpha}(v+w).
\]
There is an open subset \(L\) of
\(\mathbb G_{a,S}^{\,3}\), containing the zero section, such that
\((u,v,w)\in L(S')\) is equivalent to
\[
\begin{gathered}
  (u,v)\in V(S'),\qquad
  \bigl(uE(uv)^{-1},w\bigr)\in V(S'),\qquad
  (u,v+w)\in V(S'),\\
  uv\in W(S'),\qquad
  uwE(uv)^{-1}\in W(S'),\qquad
  u(v+w)\in W(S').
\end{gathered}
\]
Using formula \((+)\), we immediately obtain, for
\((u,v,w)\in L(S')\),
\SGARefTarget{sga3:xx:2.2:1}
\begin{equation}\tag{1}
  E(uv+uw)=E\bigl(uwE(uv)^{-1}\bigr)E(uv),
\end{equation}
\SGARefTarget{sga3:xx:2.2:2}
\begin{equation}\tag{2}
  H(uv+uw)=H\bigl(uwE(uv)^{-1}\bigr)H(uv),
\end{equation}
and
\SGARefTarget{sga3:xx:2.2:3}
\begin{equation}\tag{3}
\begin{split}
  (v+w)F(uv+uw)^{-1}
   ={}&\alpha\bigl(H(uv)^{-1}\bigr)\,
        wF\bigl(uwE(uv)^{-1}\bigr)^{-1}\\
     &{}+vF(uv)^{-1}.
\end{split}
\end{equation}

It follows immediately from the definition of \(L\) that
\((1,0,0)\in L(S)\). Consider therefore
\[
  L\cap_T
    \bigl(1\times_S\mathbb G_{a,S}\times_S\mathbb G_{a,S}\bigr)
      =1\times_S M.
\]
The scheme \(M\) is an open subset of
\(\mathbb G_{a,S}^{\,2}\) containing the section \((0,0)\).
For \((v,w)\in M(S')\), the elements
\[
  v,\qquad wE(v)^{-1},\qquad v+w
\]
belong to \(W(S')\), and
\SGARefTarget{sga3:xx:2.2:1-prime}
\begin{equation}\tag{1'}
  E(v+w)=E\bigl(wE(v)^{-1}\bigr)E(v),
\end{equation}
\SGARefTarget{sga3:xx:2.2:2-prime}
\begin{equation}\tag{2'}
  H(v+w)=H\bigl(wE(v)^{-1}\bigr)H(v),
\end{equation}
\SGARefTarget{sga3:xx:2.2:3-prime}
\begin{equation}\tag{3'}
\begin{split}
  (v+w)F(v+w)^{-1}
   ={}&\alpha\bigl(H(v)\bigr)^{-1}
        wF\bigl(wE(v)^{-1}\bigr)^{-1}\\
     &{}+vF(v)^{-1}.
\end{split}
\end{equation}

Finally, consider the morphism from \(M\) to
\(\mathbb G_{a,S}^{\,2}\) defined set-theoretically by
\[
  (v,w)\longmapsto\bigl(v,wE(v)^{-1}\bigr).
\]
\footnote{The editors corrected what follows.}
It preserves the section \((0,0)\) and induces an isomorphism of
\(M\) onto an open subset \(N\) of
\(\mathbb G_{a,S}^{\,2}\) containing the zero section; the inverse
isomorphism is
\[
  (x,y)\longmapsto\bigl(x,yE(x)\bigr).
\]
\footnote{Thus the editors made the change of variables
\(x=v,\ y=wE(v)^{-1}\), or equivalently
\(v=x,\ w=yE(x)\).}
We have therefore proved the following assertion:

There is an open subset \(N\) of
\(\mathbb G_{a,S}^{\,2}\), containing the zero section, such that,
if \((x,y)\in N(S')\), then \(x\), \(y\), and
\(x+yE(x)\)\footnote{The editors corrected \(yE(x)\) to
\(x+yE(x)\).} belong to \(W(S')\), and
\SGARefTarget{sga3:xx:2.2:1-double-prime}
\begin{equation}\tag{1''}
  E\bigl(x+yE(x)\bigr)=E(x)E(y),
\end{equation}
\SGARefTarget{sga3:xx:2.2:2-double-prime}
\begin{equation}\tag{2''}
  H\bigl(x+yE(x)\bigr)=H(x)H(y),
\end{equation}
\SGARefTarget{sga3:xx:2.2:3-double-prime}
\begin{equation}\tag{3''}
\begin{split}
  \bigl(x+yE(x)\bigr)F\bigl(x+yE(x)\bigr)^{-1}
   ={}&xF(x)^{-1}\\
     &{}+\alpha\bigl(H(x)\bigr)^{-1}yE(x)F(y)^{-1}.
\end{split}
\end{equation}

\textup{C)} Arguing in the same way with left associativity, we
obtain the following assertion.\footnote{That is, one writes the
equalities resulting from
\(p_\alpha(t)p_\alpha(u)p_{-\alpha}(v)
  =p_\alpha(t+u)p_{-\alpha}(v)\), and sets
\(v=1,\ x=u,\ t=yF(u)\), or equivalently
\(y=tF(u)^{-1}\).}

There is an open subset \(N'\) of
\(\mathbb G_{a,S}^{\,2}\), containing the zero section, such that,
if \((x,y)\in N'(S')\), then \(x\), \(y\), and
\(x+yF(x)\)\footnote{The editors corrected \(yF(x)\) to
\(x+yF(x)\).} belong to \(W(S')\), and
\SGARefTarget{sga3:xx:2.2:4-double-prime}
\begin{equation}\tag{4''}
  F\bigl(x+yF(x)\bigr)=F(x)F(y),
\end{equation}
\SGARefTarget{sga3:xx:2.2:5-double-prime}
\begin{equation}\tag{5''}
  H\bigl(x+yF(x)\bigr)=H(x)H(y),
\end{equation}
\SGARefTarget{sga3:xx:2.2:6-double-prime}
\begin{equation}\tag{6''}
\begin{split}
  \bigl(x+yF(x)\bigr)E\bigl(x+yF(x)\bigr)^{-1}
   ={}&xE(x)^{-1}\\
     &{}+\alpha\bigl(H(x)\bigr)^{-1}yF(x)E(y)^{-1}.
\end{split}
\end{equation}

We are thus led to solve the ``functional equation'' \((1'')\).

\paragraph{\textbf{Lemma 2.3.}}
\SGARefTarget{sga3:xx:2.3}
\emph{Let \(S\) be a scheme, let \(W\) be an open subset of
\(\mathbb G_{a,S}\) containing the zero section, and let}
\[
  F\colon W\longrightarrow\mathbb G_{m,S}
\]
\emph{be an \(S\)-morphism. Suppose that \(F(0)=1\) and that there
is an open subset \(N\) of \(\mathbb G_{a,S}^{\,2}\), containing
the zero section, such that, for \((x,y)\in N(S')\), the elements
\(x\), \(y\), and \(x+yF(x)\) belong to \(W(S')\), and}
\SGARefTarget{sga3:xx:2.3:dagger}
\begin{equation}\tag{\(\dagger\)}
  F\bigl(x+yF(x)\bigr)=F(x)F(y).
\end{equation}
\begin{enumerate}
\item[\textup{\SGAInlineRefTarget{sga3:xx:2.3:item:i}(i)}] \emph{If \(S=\operatorname{Spec}(k)\) for a
  field \(k\), there is an \(a\in k\) such that
  \(F(x)=1+ax\).}
\item[\textup{\SGAInlineRefTarget{sga3:xx:2.3:item:ii}(ii)}] \emph{If
  \(a=F'(0)\in\Gamma(S,\mathscr O_S)\) is invertible, then
  \(F(x)=1+ax\).}
\end{enumerate}

\medskip
\noindent\emph{Proof.}
Under the hypotheses, we may differentiate the given equation at
\(x=0\), respectively at \(y=0\), and obtain
\SGARefTarget{sga3:xx:2.3:star}
\begin{equation}\tag{\(\ast\)}
  F'(y)\bigl(1+yF'(0)\bigr)=F'(0)F(y)
  \qquad\text{for }(0,y)\in N(S'),
\end{equation}
respectively
\[
  F'(x)F(x)=F(x)F'(0)
  \qquad\text{for }(x,0)\in N(S').
\]
Since \(F\) takes its values in \(\mathbb G_m\), the second
relation gives
\SGARefTarget{sga3:xx:2.3:star-prime}
\begin{equation}\tag{\(\ast'\)}
  F'(x)=F'(0)\qquad\text{for }(x,0)\in N(S').
\end{equation}
The first relation therefore gives
\[
  F'(0)\bigl(1+yF'(0)\bigr)=F'(0)F(y)
  \qquad\text{for }(y,0),(0,y)\in N(S').
\]
If \(a=F'(0)\) is invertible, it follows that
\[
  F(y)=1+ay
\]
for \(y\) a section of an open subset of \(W\) containing the zero
section, hence schematically dense in \(W\). This proves
\textup{\SGARefLink{sga3:xx:2.3:item:ii}{(ii)}}, and also proves \textup{\SGARefLink{sga3:xx:2.3:item:i}{(i)}} when \(F'(0)\ne0\).

Suppose \(F'(0)=0\). Then \((\ast')\) shows that \(F'(x)=0\)
for \(x\) near \(0\), and hence \(F'=0\) by schematic density.
If \(k\) has characteristic zero, \(F\) is a rational function
with zero derivative, and is therefore constant and equal to
\(F(0)=1\).

Suppose now that \(k\) has characteristic \(p\), and that \(F\)
is not constant.\footnote{The editors corrected what follows.}
There are an integer \(n>0\) and a rational function
\(F_1\in k(X)\) such that \(F_1'(X)\ne0\) and
\[
  F(X)=F_1\bigl(X^{p^n}\bigr)=F_1(X)^{p^n}.
\]
Substitution in the functional equation gives
\SGARefTarget{sga3:xx:2.3:dagger-one}
\begin{equation}\tag{\(\dagger_1\)}
  F_1\bigl(x+yF_1(x)^{p^n}\bigr)=F_1(x)F_1(y).
\end{equation}
Differentiating at \(x=0\) gives
\SGARefTarget{sga3:xx:2.3:star-one}
\begin{equation}\tag{\(\ast_1\)}
  F_1'(y)=F_1'(0)F_1(y),
\end{equation}
whereas differentiation of \((\dagger_1)\) at \(y=0\) gives
\SGARefTarget{sga3:xx:2.3:star-prime-one}
\begin{equation}\tag{\(\ast'_1\)}
  F_1'(x)F_1(x)^{p^n}=F_1(x)F_1'(0).
\end{equation}
By hypothesis, \(F_1'(X)\) is invertible in \(k(X)\); these two
equalities therefore imply
\[
  F_1(X)^{p^n}=1.
\]
Thus \(F_1\) is constant, contrary to the choice of \(F_1\).
Consequently \(F\) is constant and equal to \(1=F(0)\).
\qed

% END INCLUDED SGA3_English_current_buildable_source_20260823/source/inputs/XX/tex/components/05_expose_XX_prop22_partsBC_lemma23_en.tex

% BEGIN INCLUDED SGA3_English_current_buildable_source_20260823/source/inputs/XX/tex/components/06_expose_XX_prop22_partsD_through_G_en.tex
% Authority: Exp20-13oct24.pdf, local pp. 13--15 / combined pp. 1072--1074.
% Continuous English stops after the proof of Proposition 2.2.

\textup{D)} Suppose first that \(S\) is the spectrum of a field.
If \(F'(0)=0\), then \(F=1\). Formula \((5'')\) then gives
\[
  H(x+y)=H(x)H(y),
\]
which shows that \(H\) extends to a group morphism
\[
  \mathbb G_{a,S}\longrightarrow T
\]
\textup{(\SGARefLink{sga3:xviii:prop:2.3}{Exposé XVIII, 2.3})}, necessarily the constant morphism
with value \(e\). On the other hand, \SGARefLink{sga3:xx:2.3}{Lemma 2.3} also gives
\[
  E(x)=1+bx
\]
for some \(b\in k\). Formula \((6'')\) therefore gives, for
\((x,y)\in N(S')\),
\[
  (x+y)E(x+y)^{-1}=xE(x)^{-1}+yE(y)^{-1}.
\]
Thus,\footnote{The editors added the sentence that follows. One
may also see by direct calculation that the preceding equality
implies
\(0=xyb\bigl(2+(x+y)b\bigr)\), hence
\(0=b\bigl(2+(x+y)b\bigr)\), and finally \(b=0\).}
again by \SGARefLink{sga3:xviii:prop:2.3}{Exposé XVIII, 2.3}, the map
\[
  x\longmapsto xE(x)^{-1}
\]
extends to a morphism of \(k\)-groups
\(\mathbb G_{a,k}\to\mathbb G_{a,k}\). Consequently
\[
  \frac{x}{1+bx}=cx
\]
for some \(c\in k\), whence \(b=0\) and \(c=1\).

Thus \(F,H,E\) are constant with values \((1,e,1)\) in a
neighborhood of the zero section, and hence everywhere. Formula
\((+)\) then shows that \(U_\alpha\) and \(U_{-\alpha}\) commute,
contrary to hypothesis \textup{\SGARefLink{sga3:xx:2.3:item:ii}{(ii)}}.

For arbitrary \(S\), we have therefore proved that \(F'(0)\) is
nonzero on every fiber, and hence invertible. The same is plainly
true of \(E'(0)\). \SGARefLink{sga3:xx:2.3}{Lemma 2.3} now shows that there are
\(a,b\in\mathbb G_m(S)\) such that
\SGARefTarget{sga3:xx:2.2:diamond-one}
\begin{equation}\tag{\(\diamondsuit_1\)}
  F(x)=1+ax,\qquad E(x)=1+bx
  \qquad\text{for }x\in W(S').
\end{equation}

\textup{E)} The rest is now easy. Substitution of the preceding
results in \((3'')\) gives
\[
  y\,\alpha\bigl(H(x)\bigr)(1+ay)
   =y\bigl(1+ax+ay(1+bx)\bigr)(1+ax).
\]
This formula holds for every section \((x,y)\) of \(N\). Since
\(\mathbb G_{a,S}\times_S\mathbb G_{m,S}\) is schematically dense
in \(\mathbb G_{a,S}^{\,2}\), it follows that
\[
  (1+ay)\alpha\bigl(H(x)\bigr)
   =\bigl(1+ax+ay(1+bx)\bigr)(1+ax).
\]
Setting \(y=0\) gives
\[
  \alpha\bigl(H(x)\bigr)=(1+ax)^2.
\]
Substitution in the preceding equality gives\footnote{In the three
preceding equalities, the editors corrected the original by
replacing the factor \(1+bx\) on the right by \(1+ax\).
Substitution of the third equality in the second, and use of the
invertibility of \(1+ax\), gives precisely
\(a^2xy=abxy\).}
\[
  a^2xy=abxy.
\]
Since \(\mathbb G_{m,S}\) is schematically dense in
\(\mathbb G_{a,S}\), we obtain \(a^2=ab\). As \(a\) is invertible,
\SGARefTarget{sga3:xx:2.2:diamond-two}
\begin{equation}\tag{\(\diamondsuit_2\)}
  a=b,\qquad \alpha\bigl(H(x)\bigr)=(1+ax)^2.
\end{equation}

Because \(a\) is invertible, the map \(x\mapsto1+ax\) is an
automorphism of \(\mathbb G_{a,S}\). There are therefore an open
subset \(W'\) of \(\mathbb G_{a,S}\), containing the section \(1\),
and a morphism
\[
  P\colon W'\longrightarrow T
\]
such that
\[
  P(1+ax)=H(x).
\]
\footnote{This is the change of variables
\(x'=1+ax\), or \(x=(x'-1)/a\).}
Substitution in \((2')\) gives immediately, for
\((x,y)\in N(S')\),
\[
  P(1+ax+ay)
   =P\left(\frac{1+ax+ay}{1+ax}\right)P(1+ax).
\]
Thus there is an open neighborhood of \(1\) in
\(\mathbb G_{m,S}\) on which
\[
  P(x)P(y)=P(xy).
\]
By \SGARefLink{sga3:xviii:prop:2.3}{Exposé XVIII, 2.3}, there is a group morphism
\SGARefTarget{sga3:xx:2.2:diamond-three}
\begin{equation}\tag{\(\diamondsuit_3\)}
  \alpha^\ast\colon\mathbb G_{m,S}\longrightarrow T
\end{equation}
extending \(P\). Since
\(\alpha(H(x))=(1+ax)^2\) near the zero section, we have
\(\alpha(\alpha^\ast(z))=z^2\) near the section \(1\), and hence
\SGARefTarget{sga3:xx:2.2:diamond-four}
\begin{equation}\tag{\(\diamondsuit_4\)}
  \alpha\circ\alpha^\ast=2.
\end{equation}

\textup{F)}\footnote{The editors slightly modified what follows,
since an open subset \(V\) had already been introduced in part A.}
Combining \((+)\) and
\((\diamondsuit_1)\)--\((\diamondsuit_4)\), we see that there are
\[
  a\in\mathbb G_m(S),\qquad
  \alpha^\ast\in
    \operatorname{Hom}_{S\text{-gr.}}(\mathbb G_{m,S},T)
\]
such that \(\alpha\circ\alpha^\ast=2\), and such that, whenever
\((u,v)\in V(S')\) and \(uv\in W(S')\), the element \(1+auv\)
is invertible and
\[
  p_\alpha(u)p_{-\alpha}(v)
   =p_{-\alpha}\left(\frac{v}{1+auv}\right)
      \alpha^\ast(1+auv)
      p_\alpha\left(\frac{u}{1+auv}\right).
\]

Let \(V'\) be the open subset of \(\mathbb G_{a,S}^{\,2}\) defined
by the invertibility of \(1+auv\); thus
\[
  V'=\bigl(\mathbb G_{a,S}^{\,2}\bigr)_f,
  \qquad f(u,v)=1+auv.
\]
The two sides of the preceding formula define morphisms
\(V'\to G\) which agree in a neighborhood of the zero section,
and hence agree on \(V'\). Thus the formula holds for every
section \((u,v)\) of \(V'\), and in particular \(V'\subset U\),
where \(U\) is the open subset introduced at the beginning of
part A.

We prove that \(U=V'\). Returning to the notation of part A, there
is a morphism
\[
  A\colon U\longrightarrow\mathbb G_{a,S}
\]
whose restriction to \(V'\) is
\[
  A(u,v)=v(1+auv)^{-1}.
\]
The equality \(U=V'\) is a set-theoretic question and may therefore
be checked when \(S=\operatorname{Spec}(k)\) for a field \(k\).
It then reduces to the evident assertion that the domain of
definition of the rational map
\[
  \mathbb G_{a,k}^{\,2}\dashrightarrow\mathbb G_{a,k},
  \qquad
  (X,Y)\longmapsto\frac{Y}{1+aXY},
\]
is the open subset defined by \(1+aXY\).

\textup{G)} We have proved the existence of \(a\) and
\(\alpha^\ast\), as well as the two asserted supplementary
properties. It remains to prove uniqueness. Let
\(a'\) and \(\alpha^{\ast\prime}\) also satisfy the required
conditions. For \(u,v\in\mathbb G_a(S')\), we immediately have
\[
\begin{gathered}
  1+auv\text{ invertible}
    \ \Longrightarrow\
  1+a'uv\text{ invertible},\\
  \frac{v}{1+auv}=\frac{v}{1+a'uv}.
\end{gathered}
\]
Hence, for every \(u\in\mathbb G_a(S')\),
\[
  1+au\text{ invertible}
    \ \Longrightarrow\
  1+au=1+a'u,
\]
which proves \(a=a'\). With the same notation,
\[
  1+au\text{ invertible}
    \ \Longrightarrow\
  \alpha^\ast(1+au)=\alpha^{\ast\prime}(1+au),
\]
and therefore \(\alpha^\ast=\alpha^{\ast\prime}\).
\qed

% END INCLUDED SGA3_English_current_buildable_source_20260823/source/inputs/XX/tex/components/06_expose_XX_prop22_partsD_through_G_en.tex

% BEGIN INCLUDED SGA3_English_current_buildable_source_20260823/source/inputs/XX/tex/components/07_expose_XX_cor24_through_lemma213_en.tex
% Authority: Exp20-13oct24.pdf, local pp. 15--17 / combined pp. 1074--1076.
% Continuous English stops after Lemma 2.13, before Lemma 2.14.

\paragraph{\textbf{Corollary 2.4.}}
\SGARefTarget{sga3:xx:2.4}
\emph{Let}
\[
  \exp(Y)t\exp(X)
  \quad\text{and}\quad
  \exp(Y')t'\exp(X')
\]
\emph{be two elements of \(\Omega(S')\). Their product belongs to
\(\Omega(S')\) if and only if}
\[
  u=1+\langle X,Y'\rangle
\]
\emph{is invertible; in that case,}
\SGARefTarget{sga3:xx:2.4:F-prime}
\begin{equation}\tag{F'}
\begin{split}
  &\exp(Y)t\exp(X)\cdot\exp(Y')t'\exp(X')\\
  &\qquad
   =\exp\bigl(Y+u^{-1}\alpha(t)^{-1}Y'\bigr)\,
      tt'\alpha^\ast(u)\,
      \exp\bigl(u^{-1}\alpha(t')^{-1}X+X'\bigr).
\end{split}
\end{equation}

\paragraph{\textbf{Remark 2.5.}}
\SGARefTarget{sga3:xx:2.5}
\emph{Formula \textup{(F)} of \SGARefLink{sga3:xx:2.1}{Theorem 2.1} may also be written
without using the morphisms \(\exp\). Transporting the duality
\(\mathfrak g_\alpha\otimes\mathfrak g_{-\alpha}\to\mathscr O_S\)
through those morphisms gives a canonical pairing of vector
bundles}
\[
  U_\alpha\times_S U_{-\alpha}\longrightarrow\mathbb G_{a,S},
\]
\emph{which we continue to denote by
\((x,y)\mapsto\langle x,y\rangle\). Thus}
\[
  \langle\exp X,\exp Y\rangle=\langle X,Y\rangle.
\]
\emph{If \(x\in U_\alpha(S')\), \(y\in U_{-\alpha}(S')\), and
\(1+\langle x,y\rangle\in\mathbb G_m(S')\), then}
\[
  x\cdot y
   =y^{(1+\langle x,y\rangle)^{-1}}\,
      \alpha^\ast\bigl(1+\langle x,y\rangle\bigr)\,
      x^{(1+\langle x,y\rangle)^{-1}}.
\]

\paragraph{\textbf{Corollary 2.6.}}
\SGARefTarget{sga3:xx:2.6}
\emph{The pairing}
\[
  W(\mathfrak g_\alpha)\times_S W(\mathfrak g_{-\alpha})
    \longrightarrow\mathbb G_{a,S}
\]
\emph{induces a pairing of principal
\(\mathbb G_{m,S}\)-bundles}
\[
  W(\mathfrak g_\alpha)^\times
    \times_S W(\mathfrak g_{-\alpha})^\times
      \longrightarrow\mathbb G_{m,S}.
\]
\emph{We denote this pairing by
\((X,Y)\mapsto\langle X,Y\rangle\), or simply
\((X,Y)\mapsto XY\).}

For every \(X\in\Gamma(S,\mathfrak g_\alpha)^\times\), there is
therefore a unique
\[
  X^{-1}\in\Gamma(S,\mathfrak g_{-\alpha})^\times
\]
such that \(XX^{-1}=1\). Moreover,
\[
  (zX)^{-1}=z^{-1}X^{-1}.
\]
The morphism
\[
  s\colon W(\mathfrak g_\alpha)^\times
    \longrightarrow W(\mathfrak g_{-\alpha})^\times
\]
so defined is an isomorphism of schemes, compatible with the
isomorphism \(s\colon z\mapsto z^{-1}\) of their operator groups.

\paragraph{\textbf{Definition 2.6.1.}}
\SGARefTarget{sga3:xx:2.6.1}
\emph{The elements \(X\) and \(s(X)=X^{-1}\) are said to be
paired.}

Apply \SGARefLink{sga3:xx:2.4}{Corollary 2.4} with \(Y=0=X'\) and
\[
  Y'=aX^{-1},\qquad a\in\mathscr O_S(S).
\]
Then \(u=1+a\), and
\[
  u^{-1}Y'
   =u^{-1}(u-1)X^{-1}
   =(1-u^{-1})X^{-1}.
\]
Hence:

\paragraph{\textbf{Corollary 2.7.}}
\SGARefTarget{sga3:xx:2.7}
\emph{Let \(X\in\Gamma(S,\mathfrak g_\alpha)^\times\) and
\(u\in\Gamma(S,\mathscr O_S)^\times\). Then}
\[
  \alpha^\ast(u)
   =\exp\bigl((u^{-1}-1)X^{-1}\bigr)\,
      \exp(X)\,
      \exp\bigl((u-1)X^{-1}\bigr)\,
      \exp(-u^{-1}X).
\]

\paragraph{\textbf{Definition 2.8.}}
\SGARefTarget{sga3:xx:2.8}
\emph{The morphism \(\alpha^\ast\) is called the coroot associated
with the root \(\alpha\).}

\paragraph{\textbf{Remark 2.9.}}
\SGARefTarget{sga3:xx:2.9}
\emph{If \((G,T,\alpha)\) is an \(S\)-elementary system, then so
is \((G,T,-\alpha)\). \SGARefLink{sga3:xx:2.1}{Theorem 2.1} therefore gives a duality
between \(\mathfrak g_{-\alpha}\) and \(\mathfrak g_\alpha\), and
a coroot \((-\alpha)^\ast\). Taking the inverse of formula
\textup{(F)} gives immediately}
\[
  \langle X,Y\rangle=\langle Y,X\rangle,
  \qquad
  (-\alpha)^\ast=-\alpha^\ast.
\]

We now pass to the Lie algebra of \(G\). The root \(\alpha\) and
the coroot \(\alpha^\ast\) define linear forms
\[
  \mathscr O_S\xrightarrow{\ \overline{\alpha^\ast}\ }
    \mathfrak t\xrightarrow{\ \bar\alpha\ }\mathscr O_S.
\]
Put
\[
  H_\alpha=\overline{\alpha^\ast}(1).
\]
We call \(\bar\alpha\) the infinitesimal root associated with
\(\alpha\), and \(H_\alpha\) the corresponding infinitesimal
coroot.

\paragraph{\textbf{Lemma 2.10.}}
\SGARefTarget{sga3:xx:2.10}
\emph{Let \(S'\to S\), and let}
\[
  X,X'\in W(\mathfrak g_\alpha)(S'),\qquad
  H\in W(\mathfrak t)(S'),\qquad
  Y,Y'\in W(\mathfrak g_{-\alpha})(S'),\qquad
  t\in T(S').
\]
\emph{Then}
\[
\tag{1}\label{sga3:xx:lemma:2-10:equation:1}
  \operatorname{Ad}(t)H=H,\qquad
  \operatorname{Ad}(t)X=\alpha(t)X,\qquad
  \operatorname{Ad}(t)Y=\alpha(t)^{-1}Y;
\]
\[
  \tag{2}\label{sga3:xx:lemma:2-10:equation:2}
\begin{aligned}
  \operatorname{Ad}(\exp X)H
    &=H-\bar\alpha(H)X,&
  \operatorname{Ad}(\exp X)X'&=X',\\
  \operatorname{Ad}(\exp X)Y
    &=Y+\langle X,Y\rangle H_\alpha
        -\langle X,Y\rangle X;
\end{aligned}
\]
\[
  \tag{2'}\label{sga3:xx:lemma:2-10:equation:2-prime}
\begin{aligned}
  \operatorname{Ad}(\exp Y)H
    &=H+\bar\alpha(H)Y,&
  \operatorname{Ad}(\exp Y)Y'&=Y',\\
  \operatorname{Ad}(\exp Y)X
    &=X+\langle X,Y\rangle H_{-\alpha}
        -\langle X,Y\rangle Y;
\end{aligned}
\]
\[
\tag{3}\label{sga3:xx:lemma:2-10:equation:3}
\begin{gathered}
  [H,X]=\bar\alpha(H)X,\qquad
  [H,Y]=-\bar\alpha(H)Y,\\
  [X,Y]=\langle X,Y\rangle H_\alpha;
\end{gathered}
\]
\begin{align}
  H_{-\alpha}&=-H_\alpha, \tag{4}\label{sga3:xx:lemma:2-10:equation:4}\\
  \bar\alpha(H_\alpha)&=2. \tag{5}\label{sga3:xx:lemma:2-10:equation:5}
\end{align}

The proof of these formulas is either immediate or an immediate
consequence of formula \textup{(F)} of \SGARefLink{sga3:xx:2.1}{2.1}.

\paragraph{\textbf{Corollary 2.11.}}
\SGARefTarget{sga3:xx:2.11}
\emph{Suppose that \(H_\alpha\) is nonzero on every fiber; by
\textup{\SGARefLink{sga3:xx:lemma:2-10:equation:5}{(5)}}, this holds in particular if \(2\) is invertible on
\(S\). Then}
\[
  X_\alpha\in\Gamma(S,\mathfrak g_\alpha)^\times,\qquad
  X_{-\alpha}\in\Gamma(S,\mathfrak g_{-\alpha})^\times
\]
\emph{are paired if and only if}
\[
  [X_\alpha,X_{-\alpha}]=H_\alpha.
\]

\paragraph{\textbf{2.12.}}
\SGARefTarget{sga3:xx:2.12}
Let \((G,T,\alpha)\) be an \(S\)-elementary system. By \SGARefLink{sga3:xx:1.19}{1.19},
\[
  \operatorname{Centr}(G)=\ker(\alpha)
\]
is a group of multiplicative type and finite type. If \(Q\) is a
subgroup of multiplicative type of \(\operatorname{Centr}(G)\),
then \(G/Q\) is affine over \(S\) \textup{(\SGARefLink{sga3:ix:corollary:2:5}{Exposé IX, 2.5})}, smooth
over \(S\) \textup{(Exposé VI\textsubscript{B}, \SGARefLink{prop:sga3-VIB-9.2}{9.2})}, and has
connected reductive fibers of semisimple rank \(1\)
\textup{(\SGARefLink{sga3:xix:paragraph:1-8}{Exposé XIX, 1.8})}.

Put \(G'=G/Q\). It is a reductive \(S\)-group of semisimple rank
\(1\), and \(T'=T/Q\) is a maximal torus of \(G'\). The open
subset
\[
  U_{-\alpha}\cdot T\cdot U_\alpha
\]
of \(G\) is stable under \(Q\), and its quotient is isomorphic to
\[
  U_{-\alpha}\times_S(T/Q)\times_S U_\alpha.
\]
If \(\alpha'\) denotes the character of \(T'\) induced by
\(\alpha\), the canonical morphism \(G\to G'\) induces
isomorphisms
\[
  \mathfrak g_\alpha\xrightarrow{\sim}\mathfrak g'_{\alpha'},
  \qquad
  \mathfrak g_{-\alpha}\xrightarrow{\sim}\mathfrak g'_{-\alpha'}.
\]
In particular, \(\alpha'\) is a root of \(G'\) relative to \(T'\).
Writing \(\alpha/Q\) for the character
\[
  T/Q\longrightarrow\mathbb G_{m,S}
\]
induced by \(\alpha\), we obtain:

\paragraph{\textbf{Lemma 2.13.}}
\SGARefTarget{sga3:xx:2.13}
\emph{If \(Q\) is a subgroup of multiplicative type of
\(\ker(\alpha)\), then}
\[
  (G/Q,T/Q,\alpha/Q)
\]
\emph{is an elementary system.}

% END INCLUDED SGA3_English_current_buildable_source_20260823/source/inputs/XX/tex/components/07_expose_XX_cor24_through_lemma213_en.tex

% BEGIN INCLUDED SGA3_English_current_buildable_source_20260823/source/inputs/XX/tex/components/08_expose_XX_lemma214_through_theorem31_statement_en.tex
% Authority: Exp20-13oct24.pdf, local pp. 18--19 / combined pp. 1077--1078.
% Native Loop-2 redraws; the source-derived raster crops remain private witnesses.
% Continuous English stops immediately before the proof of Theorem 3.1.

\paragraph{\textbf{Lemma 2.14.}}
\SGARefTarget{sga3:xx:2.14}
\emph{Under the preceding conditions, the following diagrams are
commutative:}

\begin{center}
  \begin{tikzcd}[column sep=huge,row sep=large]
    \mathrm W(\mathfrak g^\alpha)
      \arrow[r,"\exp"]
      \arrow[d,"\mathrm{can}"',"\sim"]
      & G
        \arrow[d,"\mathrm{can}"']
      & \mathrm W(\mathfrak g^{-\alpha})
        \arrow[l,"\exp"]
        \arrow[d,"\mathrm{can}"',"\sim"] \\
    \mathrm W(\mathfrak g'^{\alpha'})
      \arrow[r,"\exp"]
      & G'
      & \mathrm W(\mathfrak g'^{-\alpha'})
        \arrow[l,"\exp"]
  \end{tikzcd}

  \medskip
  \begin{tikzcd}[column sep=huge,row sep=large]
    \mathfrak g^\alpha\otimes\mathfrak g^{-\alpha}
      \arrow[r,"\sim"]
      \arrow[d,"\mathrm{can}"',"\sim"]
      & \mathcal O_S
        \arrow[d,"\mathrm{id}"] \\
    \mathfrak g'^{\alpha'}\otimes\mathfrak g'^{-\alpha'}
      \arrow[r,"\sim"]
      & \mathcal O_S
  \end{tikzcd}

  \medskip
  \begin{tikzcd}[column sep=huge,row sep=large]
      & T
        \arrow[dr,"\alpha"]
        \arrow[dd,"\mathrm{can}"] & \\
    \mathbb G_{m,S}
      \arrow[ur,"\alpha^\ast"]
      \arrow[dr,"\alpha'^\ast"']
      & & \mathbb G_{m,S} \\
      & T'
        \arrow[ur,"\alpha'"'] &
  \end{tikzcd}
\end{center}

\section{The Weyl group}
\SGARefTarget{sga3:xx:section:3}

\paragraph{\textbf{Notations 3.0.}}\footnote{The editors added the
number \SGARefLink{sga3:xx:3.0}{3.0} for later references.}
\SGARefTarget{sga3:xx:3.0}
If \((G,T,\alpha)\) is an \(S\)-elementary system, write
\[
  N=\operatorname{Norm}_G(T),
  \qquad
  W=\operatorname{Norm}_G(T)/T
\]
\textup{(cf.\ Exposé XIX, 6.3)}.  The group \(N\) is a closed
subgroup of \(G\), smooth over \(S\).  Write
\[
  N^\times=N-T
\]
for the open subscheme of \(N\) induced on the complement of
\(T\).\footnote{The editors replaced \(Q\) by the more suggestive
notation \(N^\times\).}

Let \(R\) be the unique maximal torus of \(\Ker(\alpha)\), and let
\(T'\) be the image of
\[
  \alpha^\ast\colon\mathbb G_{m,S}\longrightarrow T;
\]
it is a one-dimensional subtorus of \(T\).  The morphism
\[
  T'\times_S R\longrightarrow T
\]
induced by multiplication in \(T\) is surjective, hence faithfully
flat.  Indeed, it is enough to check this on the geometric fibers,
where it follows at once from the formula
\[
  \alpha\circ\alpha^\ast=2.
\]

\paragraph{\textbf{Theorem 3.1.}}
\SGARefTarget{sga3:xx:3.1}
\emph{With the preceding notation:}

\smallskip
\noindent\emph{\textup{(i)} \(W\) is isomorphic to the constant group
\((\mathbb Z/2\mathbb Z)_S\).}

\smallskip
\noindent\emph{\textup{(ii)} \(N^\times\) is a locally trivial
principal homogeneous bundle under \(T\), on the left by the law
\((t,q)\mapsto tq\), and on the right by the law
\((q,t)\mapsto qt\).}

\smallskip
\noindent\emph{\textup{\SGAInlineRefTarget{sga3:xx:3.1:item:iii}(iii)} One has}
\[
  \intaut(w)t=t\cdot\alpha^\ast\bigl(\alpha(t)^{-1}\bigr)
\]
\emph{for \(w\in N^\times(S')\), \(t\in T(S')\), and \(S'\to S\).
In the decomposition \(T_{S'}=T'_{S'}\cdot R_{S'}\),
\(\intaut(w)\) induces the identity on \(R_{S'}\) and inversion on
\(T'_{S'}\).  One has the relations}
\[
  \alpha\circ\intaut(w)=\alpha^{-1},
  \qquad
  \intaut(w)\circ\alpha^\ast=(\alpha^\ast)^{-1}.
\]

\smallskip
\noindent\emph{\textup{\SGAInlineRefTarget{sga3:xx:3.1:item:iv}(iv)} For
\(X\in W(\mathfrak g^\alpha)^\times(S')\), \(S'\to S\), set}
\[
  w_\alpha(X)
    =\exp(X)\exp(-X^{-1})\exp(X).
\]
\emph{Then \(w_\alpha(X)\in N^\times(S')\), and the morphism
\(w_\alpha\colon W(\mathfrak g^\alpha)^\times\to N^\times\) so
defined satisfies}
\[
  w_\alpha(zX)
    =\alpha^\ast(z)w_\alpha(X)
    =w_\alpha(X)\alpha^\ast(z)^{-1}
\]
\emph{for \(z\in\mathbb G_m(S')\),
\(X\in W(\mathfrak g^\alpha)^\times(S')\), and \(S'\to S\).}

\smallskip
\noindent\emph{\textup{\SGAInlineRefTarget{sga3:xx:3.1:item:v}(v)} For
\(X,Y\in W(\mathfrak g^\alpha)^\times(S')\), \(S'\to S\), one has,
with the notation of \SGARefLink{sga3:xx:2.6}{2.6},}
\[
  w_\alpha(X)w_\alpha(Y)=\alpha^\ast(-XY^{-1}).
\]
\footnote{The editors corrected the first formula of the original.}
\emph{In particular,}
\[
  w_\alpha(X)^2
    =\alpha^\ast(-1)
    \in {}_2T(S)\cap\operatorname{Centr}(G)(S),
\]
\[
  w_\alpha(X)^{-1}
    =w_\alpha(-X)
    =\alpha^\ast(-1)w_\alpha(X).
\]

\smallskip
\noindent\emph{\textup{\SGAInlineRefTarget{sga3:xx:3.1:item:vi}(vi)} If, similarly, one defines for
\(Y\in W(\mathfrak g^{-\alpha})^\times(S')\)}
\[
  w_{-\alpha}(Y)
    =\exp(Y)\exp(-Y^{-1})\exp(Y),
\]
\emph{then, in addition to the formulas analogous to the preceding
ones,}
\[
  w_{-\alpha}(X^{-1})
    =w_\alpha(X)^{-1}
    =w_\alpha(-X),
\qquad
  w_\alpha(X)w_{-\alpha}(Y)=\alpha^\ast(XY).
\]

% END INCLUDED SGA3_English_current_buildable_source_20260823/source/inputs/XX/tex/components/08_expose_XX_lemma214_through_theorem31_statement_en.tex

% BEGIN INCLUDED SGA3_English_current_buildable_source_20260823/source/inputs/XX/tex/components/09_expose_XX_theorem31_proof_through_remark34_en.tex
% Authority: Exp20-13oct24.pdf, local pp. 19--21 / combined pp. 1078--1080.
% Continuous English stops immediately before paragraph 3.5.

\noindent\emph{Proof.}
Assertion \textup{(i)} was already proved in \SGARefLink{sga3:xix:remark:2:4}{Exposé XIX, 2.4}. It
follows at once that \(N^\times\) is indeed a principal homogeneous
bundle under \(T\) for the laws defined in \textup{(ii)}; its local
triviality for the Zariski topology follows in particular from
\textup{\SGARefLink{sga3:xx:3.1:item:iv}{(iv)}}.

Let us prove \textup{(iii)}. If \(w\in N^\times(S)\), it is clear
that \(\alpha\circ\intaut(w)\) is a root of \(G\) relative to \(T\),
and hence is locally equal to \(\alpha\) or \(-\alpha\). Since it is
\(-\alpha\) on every fiber \textup{(Bible, 12-05, proof of the
corollary to Proposition 1)}, one has
\[
  \alpha\circ\intaut(w)=-\alpha.
\]
By transport of structure, it follows that
\[
  -\alpha^\ast
    =\intaut(w)^{-1}\circ\alpha^\ast
    =\intaut(w)\circ\alpha^\ast,
\]
since \(\intaut(w)^2=\intaut(w^2)\) and \(w^2\) is a section of
\(T\). Thus \(\intaut(w)\) induces inversion on \(T'\); since \(R\)
is central, it induces the identity on \(R\). The formula in
\textup{(iii)} defines a morphism \(T\to T\) having the same
properties, and therefore coincides with \(\intaut(w)\).

Let us prove \textup{\SGARefLink{sga3:xx:3.1:item:iv}{(iv)}}. Successively,
\[
\begin{aligned}
  w_\alpha(X)t w_\alpha(X)^{-1}
    &={}\exp(X)\exp(-X^{-1})\exp(X)t
       \exp(-X)\exp(X^{-1})\exp(-X)\\
    &={}\exp(X)\exp(-X^{-1})
       \exp\bigl(X-\alpha(t)X\bigr)
       \exp\bigl(\alpha(t)^{-1}X^{-1}\bigr)
       \exp\bigl(-\alpha(t)X\bigr)t.
\end{aligned}
\]
Applying formula \textup{(F)} gives
\[
\begin{aligned}
  \exp(-X^{-1})\exp\bigl((1-\alpha(t))X\bigr)
    ={}&\exp\bigl((\alpha(t)^{-1}-1)X\bigr)
       \alpha^\ast\bigl(\alpha(t)^{-1}\bigr)\\
     &{}\cdot\exp\bigl(-\alpha(t)^{-1}X^{-1}\bigr).
\end{aligned}
\]
Substitution in the preceding relation yields
\[
\begin{aligned}
  \intaut(w_\alpha(X))t
    &={}\exp\bigl(\alpha(t)^{-1}X\bigr)
       \alpha^\ast\bigl(\alpha(t)^{-1}\bigr)
       \exp\bigl(-\alpha(t)X\bigr)t\\
    &={}\exp(aX)\alpha^\ast\bigl(\alpha(t)^{-1}\bigr)t,
\end{aligned}
\]
where
\[
  a=\alpha(t)^{-1}
     -(\alpha\circ\alpha^\ast)\bigl(\alpha(t)^{-1}\bigr)\alpha(t).
\]
But \(\alpha\circ\alpha^\ast=2\), which gives \(a=0\) at once and
shows that \(w_\alpha(X)\in N^\times(S')\).

We now prove the second assertion of \textup{(iv)}. One has%
\footnote{The first equality follows from \SGARefLink{sga3:xx:1.5}{1.5} \textup{\SGARefLink{sga3:xx:1.5:item:i}{(i)}}, which,
together with \(\alpha\circ\alpha^\ast=2\), gives
\(\alpha^\ast(z)\exp(X)\alpha^\ast(z)^{-1}=\exp(z^2X)\) and
\(\alpha^\ast(z)\exp(X^{-1})\alpha^\ast(z)^{-1}
=\exp(z^{-2}X)\). The third equality follows from formula
\textup{(F)}, and the fourth again from these formulas. Finally, an
analogous calculation gives
\(w_\alpha(X)\alpha^\ast(z^{-1})=w_\alpha(zX)\).}
\[
\begin{aligned}
  \alpha^\ast(z)w_\alpha(X)
    &={}\exp(z^2X)\exp(-z^{-2}X^{-1})
       \exp(z^2X)\alpha^\ast(z)\\
    &={}\exp(zX)\exp\bigl((z^2-z)X\bigr)
       \exp(-z^{-2}X^{-1})\exp(z^2X)\alpha^\ast(z)\\
    &={}\exp(zX)\exp(-z^{-1}X^{-1})\alpha^\ast(z)^{-1}
       \exp\bigl((z^3-z^2)X\bigr)\exp(z^2X)\alpha^\ast(z)\\
    &={}\exp(zX)\exp(-z^{-1}X^{-1})\exp(zX)
       =w_\alpha(zX).
\end{aligned}
\]

Let us prove \textup{\SGARefLink{sga3:xx:3.1:item:v}{(v)}}. By the preceding result, the first
formula in \textup{\SGARefLink{sga3:xx:3.1:item:v}{(v)}} follows at once from the second; we prove
the latter:
\[
\begin{aligned}
  w_\alpha(X)^2
    &={}\exp(X)\exp(-X^{-1})\exp(2X)\exp(-X^{-1})\exp(X)\\
    &={}\exp(X)\exp(-X^{-1})\exp(X^{-1})
       \alpha^\ast(-1)\exp(-2X)\exp(X)\\
    &={}\exp(X)\alpha^\ast(-1)\exp(-X)
       =\alpha^\ast(-1),
\end{aligned}
\]
since
\[
  \alpha\bigl(\alpha^\ast(-1)\bigr)=(-1)^2=1.
\]
This also proves that
\(\alpha^\ast(-1)\in\operatorname{Centr}(G)(S)\).

Finally, let us prove \textup{\SGARefLink{sga3:xx:3.1:item:vi}{(vi)}}. The first assertion is a
special case of the second, so we prove the latter. Its two sides
define morphisms
\[
  W(\mathfrak g^\alpha)^\times
    \times_S W(\mathfrak g^{-\alpha})^\times
  \longrightarrow G.
\]
To prove that they coincide, it suffices to do so on a nonempty
open subset of every fiber \textup{(\SGARefLink{sga3:xviii:prop:1.4}{Exposé XVIII, 1.4})}. We may
therefore suppose that \(1+XY\) is invertible. Then, successively,
\[
\begin{aligned}
  w_\alpha(X)w_{-\alpha}(Y)
    ={}&\exp(X)\exp(-X^{-1})\exp(X)
        \exp(Y)\\
       &\qquad{}\exp(-Y^{-1})\exp(Y)\\
    ={}&\exp(X)\exp(-X^{-1})
        \exp\left(\frac{Y}{1+XY}\right)
        \alpha^\ast(1+XY)\\
       &\qquad{}
        \exp\left(\frac{X}{1+XY}\right)
        \exp(-Y^{-1})\exp(Y)\\
    ={}&\exp(X)
        \exp\left(\frac{-X^{-1}}{1+XY}\right)
        \alpha^\ast(1+XY)
        \exp\left(\frac{-Y^{-1}}{1+XY}\right)\exp(Y)\\
    ={}&\exp(-X^{-2}Y^{-1})
        \alpha^\ast\left(\frac{XY}{1+XY}\right)
        \exp(X+Y^{-1})\alpha^\ast(1+XY)\\
       &{}\qquad\cdot
        \exp\left(\frac{-Y^{-1}}{1+XY}\right)\exp(Y)\\
    ={}&\exp(-X^{-2}Y^{-1})\alpha^\ast(XY)
        \exp\left(\frac{Y^{-1}+X}{(1+XY)^2}\right)\\
       &\qquad{}
        \exp\left(\frac{-Y^{-1}}{1+XY}\right)\exp(Y)\\
    ={}&\alpha^\ast(XY)\exp(-Y)\exp(Y)
        =\alpha^\ast(XY).
\end{aligned}
\]

\paragraph{\textbf{Corollary 3.2.}}
\SGARefTarget{sga3:xx:3.2}
\emph{Let \(n\in\mathbb Z\), \(n\ne0\). For every \(w\in G(S)\),
the following conditions are equivalent:}

\smallskip
\noindent\emph{\textup{(i)} \(w\in N^\times(S)\);}

\smallskip
\noindent\emph{\textup{\SGAInlineRefTarget{sga3:xx:3.2:item:ii}(ii)}
\(\intaut(w)\circ n\alpha^\ast=-n\alpha^\ast\), where
\((n\alpha^\ast)(z)=\alpha^\ast(z)^n\).}

\smallskip
One has \textup{(i)}\(\Rightarrow\)\textup{\SGARefLink{sga3:xx:3.2:item:ii}{(ii)}} by \SGARefLink{sga3:xx:3.1}{Theorem 3.1}
\textup{\SGARefLink{sga3:xx:3.1:item:iii}{(iii)}}. Conversely, one may suppose that \(N^\times\) has a
section, and one is reduced to proving:

\paragraph{\textbf{Lemma 3.3.}}
\SGARefTarget{sga3:xx:3.3}
\emph{For \(n\ne0\), one has}
\[
  \operatorname{Centr}_G(n\alpha^\ast)=T.
\]

Indeed, the image \(T'\) of \(n\alpha^\ast\) is a subtorus of
\(G\). It follows from \SGARefLink{sga3:xix:paragraph:2:8}{Exposé XIX, 2.8} that
\(\operatorname{Centr}_G(n\alpha^\ast)\) is a reductive subgroup
of \(G\) containing \(T\). On every fiber,
\[
  \operatorname{Centr}_{G_s}(n\alpha^\ast_s)\ne G_s,
\]
and hence
\[
  \operatorname{Centr}_{G_s}(n\alpha^\ast_s)=T_s
\]
by \SGARefLink{sga3:xix:remark:1:6:3}{Exposé XIX, 1.6.3}.\footnote{The hypothesis
\(\operatorname{Centr}_{G_s}(n\alpha^\ast_s)\ne G_s\) implies
\[
  \dim\operatorname{Centr}_{G_s}(n\alpha^\ast_s)-\dim T_s<2;
\]
this difference is even by the cited result.}
Therefore
\[
  \operatorname{Centr}_G(n\alpha^\ast)=T,
\]
since these are smooth subgroups of \(G\).

\paragraph{\textbf{Remark 3.4.}}
\SGARefTarget{sga3:xx:3.4}
\emph{The construction of \(w_\alpha\), and the fact that
\(w_\alpha(X)\) normalizes \(T\), depend only on formula
\textup{(F)}. In particular, if \(G\) is an \(S\)-group satisfying
the conditions of \SGARefLink{sga3:xx:2.2}{2.2}, then \(\operatorname{Norm}_G(T)\) differs
from \(T\) on every fiber. Consequently, if \(G\) is an affine
\(S\)-group with connected fibers satisfying the conditions of \SGARefLink{sga3:xx:2.2}{2.2},
then \(G\) is reductive of semisimple rank \(1\). Indeed, it is
smooth in a neighborhood of the identity section, hence smooth, and
the criterion of \SGARefLink{sga3:xix:lemma:1:11}{Exposé XIX, 1.11} applies.}

% END INCLUDED SGA3_English_current_buildable_source_20260823/source/inputs/XX/tex/components/09_expose_XX_theorem31_proof_through_remark34_en.tex

% BEGIN INCLUDED SGA3_English_current_buildable_source_20260823/source/inputs/XX/tex/components/10_expose_XX_par35_through_cor38_en.tex
% Authority: Exp20-13oct24.pdf, local pp. 21--23 / combined pp. 1080--1082.
% Continuous English stops immediately before Remark 3.9.

\paragraph{\textbf{3.5.}}
\SGARefTarget{sga3:xx:3.5}
Before stating the following theorem, let us make a few remarks.  As
usual, we identify \(\mathfrak g^{-\alpha}\) with
\((\mathfrak g^\alpha)^{\otimes-1}\).  Similarly, we identify
\[
  \operatorname{Hom}_{\mathscr O_S}
    (\mathfrak g^{-\alpha},\mathfrak g^\alpha)
  \quad\text{with}\quad
  (\mathfrak g^\alpha)^{\otimes2},
\]
and hence
\[
  \operatorname{Isom}_{\mathscr O_S\text{-mod.}}
    \bigl(W(\mathfrak g^{-\alpha}),W(\mathfrak g^\alpha)\bigr)
  \simeq
  W\bigl((\mathfrak g^\alpha)^{\otimes2}\bigr)^\times.
\]
If \(w\in N^\times(S)\), then
\(\operatorname{Ad}(w)\) permutes
\(\mathfrak g^\alpha\) and \(\mathfrak g^{-\alpha}\)
\textup{(\SGARefLink{sga3:xx:3.1}{3.1} \textup{\SGARefLink{sga3:xx:3.1:item:iii}{(iii)}})}, and therefore defines an isomorphism
\[
  a_\alpha(w)\colon
  \mathfrak g^{-\alpha}\xrightarrow{\ \sim\ }\mathfrak g^\alpha.
\]
We thus identify it with a section
\[
  a_\alpha(w)\in
  \Gamma\bigl(S,(\mathfrak g^\alpha)^{\otimes2}\bigr)^\times.
\]
This construction is compatible with base change and therefore
defines a morphism
\[
  a_\alpha\colon
  N^\times\longrightarrow
  W\bigl((\mathfrak g^\alpha)^{\otimes2}\bigr)^\times
\]
such that
\[
  a_\alpha(w)Y=\operatorname{Ad}(w)Y
\]
for every \(w\in N^\times(S')\),
\(Y\in\Gamma(S',\mathfrak g^{-\alpha})^\times\), and \(S'\to S\).

\paragraph{\textbf{Theorem 3.6.}}
\SGARefTarget{sga3:xx:3.6}
\emph{\textup{(i)} For every \(S'\to S\),
\(w\in N^\times(S')\), and
\(Y\in W(\mathfrak g^{-\alpha})(S')\), one has}
\[
  \intaut(w)\exp(Y)=\exp\bigl(a_\alpha(w)Y\bigr).
\]

\smallskip
\noindent\emph{\textup{(ii)} One has}
\[
  a_\alpha(tw)=\alpha(t)a_\alpha(w),
  \qquad
  a_\alpha(wt)=\alpha(t)^{-1}a_\alpha(w).
\]

\smallskip
\noindent\emph{\textup{\SGAInlineRefTarget{sga3:xx:3.6:item:iii}(iii)} If one similarly defines}
\[
  a_{-\alpha}\colon
  N^\times\longrightarrow
  W\bigl((\mathfrak g^{-\alpha})^{\otimes2}\bigr)^\times,
\]
\emph{then}
\[
  a_{-\alpha}(w)=a_\alpha(w)^{-1}.
\]
\footnote{That is, \(a_{-\alpha}(w)\) and \(a_\alpha(w)\) are paired;
cf.\ \SGARefLink{sga3:xx:2.6.1}{2.6.1}.}

\smallskip
\noindent\emph{\textup{\SGAInlineRefTarget{sga3:xx:3.6:item:iv}(iv)} For every
\(X\in W(\mathfrak g^\alpha)^\times(S')\), \(S'\to S\), one has}
\[
  a_\alpha\bigl(w_\alpha(X)\bigr)=-X^2.
\]

\smallskip
\noindent\emph{Proof.}
Assertion \textup{\SGARefLink{sga3:xx:1.5:item:i}{(i)}} is immediate from the characterization of the
morphisms \(\exp\) given in \SGARefLink{sga3:xx:1.5}{1.5}.  Assertion \textup{\SGARefLink{sga3:xx:1.5:item:ii}{(ii)}} is
immediate, as is \textup{\SGARefLink{sga3:xx:3.6:item:iii}{(iii)}}.  Let us prove \textup{\SGARefLink{sga3:xx:3.6:item:iv}{(iv)}}.  Let
\[
  X\in\Gamma(S',\mathfrak g^\alpha)^\times,
  \qquad
  Z\in\Gamma(S',\mathfrak g^\alpha).
\]
By definition,\footnote{The editors supplied the details that follow.}
\[
\begin{aligned}
  a_\alpha\bigl(w_\alpha(X)\bigr)^{-1}(Z)
    &={}\operatorname{Ad}\bigl(w_\alpha(X)\bigr)(Z)\\
    &={}\operatorname{Ad}(\exp(X))
       \operatorname{Ad}(\exp(-X^{-1}))
       \operatorname{Ad}(\exp(X))(Z).
\end{aligned}
\]
Applying formulas \textup{(2\(^{\prime}\))} and \textup{(2)} of
\SGARefLink{sga3:xx:2.10}{Lemma 2.10}, together with
\[
  H_{-\alpha}=-H_\alpha,
  \qquad
  \bar\alpha(H_\alpha)=2
\]
\textup{(loc.\ cit., \textup{\SGARefLink{sga3:xx:lemma:2-10:equation:4}{(4)}} and \textup{\SGARefLink{sga3:xx:lemma:2-10:equation:5}{(5)}})}, and
\(\langle X,X^{-1}\rangle=1\) \textup{(\SGARefLink{sga3:xx:2.6}{2.6})}, the expression on the
right is, successively,
\[
\begin{aligned}
  &\operatorname{Ad}(\exp(X))
    \operatorname{Ad}(\exp(-X^{-1}))(Z)\\
  &\quad={}\operatorname{Ad}(\exp(X))
       \bigl(Z+\langle X^{-1},Z\rangle(H_\alpha-X^{-1})\bigr)\\
  &\quad={}Z+\langle X^{-1},Z\rangle
       (H_\alpha-2X-X^{-1}-H_\alpha+X)\\
  &\quad={}Z-\langle X^{-1},Z\rangle X
       -\langle X^{-1},Z\rangle X^{-1}.
\end{aligned}
\]
But
\[
  Z=\langle X^{-1},Z\rangle X,
  \qquad
  \langle X^{-1},Z\rangle X^{-1}=X^{-2}Z.
\]
Thus
\[
  a_\alpha\bigl(w_\alpha(X)\bigr)^{-1}=-X^{-2},
\]
and consequently
\[
  a_\alpha\bigl(w_\alpha(X)\bigr)=-X^2.
\]

\paragraph{\textbf{Corollary 3.7.}}
\SGARefTarget{sga3:xx:3.7}
\emph{In particular, one has}
\[
  \intaut\bigl(w_\alpha(X)\bigr)\exp(X)=\exp(-X^{-1}).
\]
\emph{Hence, by the definition of \(w_\alpha(X)\),}
\[
  w_\alpha(X)\exp(X)w_\alpha(X)^{-1}
    =\exp(-X)w_\alpha(X)\exp(-X),
\]
\emph{and therefore, by an immediate calculation,}
\[
  \bigl(w_\alpha(X)\exp(X)\bigr)^3=e.
\]

\paragraph{\textbf{Corollary 3.8.}}
\SGARefTarget{sga3:xx:3.8}
\emph{Let \(X\in\Gamma(S,\mathfrak g^\alpha)^\times\), and let
\(n\in\mathbb Z\), \(n\ne0\).  Then \(w_\alpha(X)\) is the unique
section \(w\in G(S)\) satisfying}

\smallskip
\noindent\emph{\textup{\SGAInlineRefTarget{sga3:xx:3.8:item:i}(i)}
\(\intaut(w)\circ n\alpha^\ast=-n\alpha^\ast\);}

\smallskip
\noindent\emph{\textup{\SGAInlineRefTarget{sga3:xx:3.8:item:ii}(ii)}
\(\bigl(w\exp(X)\bigr)^3=e\).}

\smallskip
We know that \(w_\alpha(X)\) satisfies these conditions.  Conversely,
let \(w\in G(S)\) satisfy \textup{\SGARefLink{sga3:xx:3.8:item:i}{(i)}} and \textup{\SGARefLink{sga3:xx:3.8:item:ii}{(ii)}}.  By \SGARefLink{sga3:xx:3.2}{3.2} and
\SGARefLink{sga3:xx:3.1}{3.1} \textup{\SGARefLink{sga3:xx:3.2:item:ii}{(ii)}}, there exists \(t\in T(S)\) such that
\[
  w=w_\alpha(X)t.
\]
Set \(u=\exp(X)\).  Then
\[
\begin{aligned}
  wuw^{-1}
    &=w_\alpha(X)t\exp(X)t^{-1}w_\alpha(X)^{-1}\\
    &=\exp\bigl(-\alpha(t)X^{-1}\bigr),
\end{aligned}
\]
whereas
\[
\begin{aligned}
  u^{-1}wu^{-1}
    &={}\exp(-X)w_\alpha(X)t\exp(-X)\\
    &={}\exp(-X)w_\alpha(X)\exp(-X)
       \exp\bigl(X-\alpha(t)X\bigr)t\\
    &={}\exp(-X^{-1})\exp\bigl(X-\alpha(t)X\bigr)t\\
    &={}\exp(-X^{-1})t
       \exp\bigl((\alpha(t)^{-1}-1)X\bigr).
\end{aligned}
\]
Now
\[
  (wu)^3=e
  \quad\Longleftrightarrow\quad
  wuw^{-1}=u^{-1}wu^{-1}.
\]
Comparing the two decompositions of this element in
\[
  U_{-\alpha}\cdot T\cdot U_\alpha
\]
gives \(t=e\).

% END INCLUDED SGA3_English_current_buildable_source_20260823/source/inputs/XX/tex/components/10_expose_XX_par35_through_cor38_en.tex

% BEGIN INCLUDED SGA3_English_current_buildable_source_20260823/source/inputs/XX/tex/components/11_expose_XX_remark39_through_prop311_en.tex
\paragraph{\textbf{Remark 3.9.}}
\SGARefTarget{sga3:xx:3.9}
We may summarize a number of the results of this section by the
following diagram of principal homogeneous bundles (on the left):

\begin{center}
  \begin{minipage}{\textwidth}
  \centering
  \begin{tikzcd}[column sep=huge]
    \mathrm W(\mathfrak g^\alpha)^\times
      \arrow[r,"w_\alpha"]
      & N^\times
        \arrow[r,"a_\alpha"]
      & \mathrm W\bigl((\mathfrak g^\alpha)^{\otimes 2}\bigr)^\times
  \end{tikzcd}

  \smallskip
  \begin{tikzcd}[column sep=huge]
    \mathbb G_{m,S}
      \arrow[r,"\alpha^\ast"]
      & T
        \arrow[r,"\alpha"]
      & \mathbb G_{m,S}
  \end{tikzcd}
  \end{minipage}
\end{center}

Notice that \(a_\alpha\) is faithfully flat (since \(\alpha\) is),
and that \(w_\alpha\) is a monomorphism if and only if
\(\alpha^\ast\) is a monomorphism.  We leave to the reader the task
of writing down the corresponding diagrams for right
principal-bundle structures, analogous diagrams for the root
\(-\alpha\), and of studying the relations among these various
diagrams.

\paragraph{\textbf{Lemma 3.10.}}
\SGARefTarget{sga3:xx:3.10}
\emph{Let \(S\) be a scheme, and let \(q>0\) be an integer such that
\(x\mapsto x^q\) defines an endomorphism of \(\mathbb G_{a,S}\).
Let \((G,T,\alpha)\) and \((G',T',\alpha')\) be two elementary
\(S\)-systems, and let \(f\colon G\to G'\) be a morphism of
\(S\)-groups.  Let}
\[
  h\colon(\mathfrak g^\alpha)^{\otimes q}
    \xrightarrow{\ \sim\ }\mathfrak g'^{\,\alpha'}
\]
\emph{be an isomorphism of \(\mathcal O_S\)-modules, and let}
\[
  h^\vee\colon(\mathfrak g^{-\alpha})^{\otimes q}
    \xrightarrow{\ \sim\ }\mathfrak g'^{-\alpha'}
\]
\emph{be the contragredient isomorphism.  Suppose that, for every
\(S'\to S\) and every \(X\in W(\mathfrak g^\alpha)(S')\),}
\[
  f\bigl(\exp(X)\bigr)=\exp\bigl(h(X^q)\bigr).
\]
\emph{Then the following conditions are equivalent:}

\smallskip
\noindent\emph{\textup{\SGAInlineRefTarget{sga3:xx:3.10:item:i}(i)}
\(f\bigl(\alpha^\ast(z)\bigr)=\alpha'^\ast(z)^q\).}

\smallskip
\noindent\emph{\textup{\SGAInlineRefTarget{sga3:xx:3.10:item:ii}(ii)}
\(f\bigl(w_\alpha(Z)\bigr)=w_{\alpha'}\bigl(h(Z^q)\bigr)\).}

\smallskip
\noindent\emph{\textup{\SGAInlineRefTarget{sga3:xx:3.10:item:iii}(iii)}
\(f\bigl(\exp(Y)\bigr)=\exp\bigl(h^\vee(Y^q)\bigr)\).}

\smallskip
Each condition is to be read as holding for every \(S'\to S\) and
every
\[
  z\in\mathbb G_m(S'),\qquad
  Z\in W(\mathfrak g^\alpha)^\times(S'),\qquad
  Y\in W(\mathfrak g^{-\alpha})(S'),
\]
respectively.  Indeed, \textup{\SGARefLink{sga3:xx:3.8:item:i}{(i)}} implies \textup{\SGARefLink{sga3:xx:3.8:item:ii}{(ii)}} by \SGARefLink{sga3:xx:3.8}{3.8},
\textup{\SGARefLink{sga3:xx:3.8:item:ii}{(ii)}} implies \textup{\SGARefLink{sga3:xx:3.10:item:iii}{(iii)}} by \SGARefLink{sga3:xx:3.7}{3.7}, and \textup{\SGARefLink{sga3:xx:3.10:item:iii}{(iii)}}
implies \textup{\SGARefLink{sga3:xx:3.8:item:i}{(i)}} by \SGARefLink{sga3:xx:2.7}{2.7}.

\paragraph{\textbf{Proposition 3.11.}}
\SGARefTarget{sga3:xx:3.11}
\emph{Let \(S\) be a scheme, and let \(q\in\mathbb Z\), \(q>0\), be
such that \(x\mapsto x^q\) defines an endomorphism of
\(\mathbb G_{a,S}\).  Let \((G,T,\alpha)\) and
\((G',T',\alpha')\) be two elementary \(S\)-systems, and let
\(f\colon G\to G'\) be a morphism of \(S\)-groups.  The following
conditions on \(f\) are equivalent:}

\smallskip
\noindent\emph{\textup{\SGARefLink{sga3:xx:3.10:item:i}{(i)}} The restriction of \(f\) to \(T\)
factors through a morphism \(f_T\colon T\to T'\) making the diagram}

\begin{center}
  \begin{tikzcd}[column sep=huge,row sep=large]
    \mathbb G_{m,S}
      \arrow[r,"\alpha^\ast"]
      \arrow[d,"q"']
      & T
        \arrow[r,"\alpha"]
        \arrow[d,"f_T"']
      & \mathbb G_{m,S}
        \arrow[d,"q"'] \\
    \mathbb G_{m,S}
      \arrow[r,"\alpha'^\ast"]
      & T'
        \arrow[r,"\alpha'"]
      & \mathbb G_{m,S}
  \end{tikzcd}
\end{center}

\noindent\emph{commutative.}

\smallskip
\noindent\emph{\textup{\SGAInlineRefTarget{sga3:xx:3.11:item:ii}(ii)} There exists a unique isomorphism of
\(\mathcal O_S\)-modules}
\[
  h\colon(\mathfrak g^\alpha)^{\otimes q}
    \xrightarrow{\ \sim\ }\mathfrak g'^{\,\alpha'}
\]
\emph{such that}
\[
  f\bigl(\exp(X)\bigr)=\exp\bigl(h(X^q)\bigr),
  \qquad
  f\bigl(\exp(Y)\bigr)=\exp\bigl(h^\vee(Y^q)\bigr)
\]
\emph{for every \(S'\to S\),
\(X\in W(\mathfrak g^\alpha)(S')\), and
\(Y\in W(\mathfrak g^{-\alpha})(S')\).  Consequently, \(f\) also
satisfies the equivalent conditions of \SGARefLink{sga3:xx:3.10}{3.10}.}

\smallskip
We first prove that \textup{\SGARefLink{sga3:xx:3.10:item:ii}{(ii)}} implies \textup{\SGARefLink{sga3:xx:3.10:item:i}{(i)}}.  By \SGARefLink{sga3:xx:3.10}{3.10},
condition \textup{\SGARefLink{sga3:xx:3.10:item:ii}{(ii)}} gives
\[
  f\circ\alpha^\ast=q\alpha'^\ast;
\]
hence, by \SGARefLink{sga3:xx:3.3}{3.3}, \(f|_T\) factors through \(T'\).  It remains to prove
\[
  \alpha'\bigl(f(t)\bigr)=\alpha(t)^q,
\]
which follows at once from the fact that \(f\) induces a morphism of
groups
\[
  T\cdot U_\alpha\longrightarrow T'\cdot U_{\alpha'}.
\]

We now prove that \textup{\SGARefLink{sga3:xx:3.10:item:i}{(i)}} implies \textup{\SGARefLink{sga3:xx:3.11:item:ii}{(ii)}}.  Let
\[
  X\in\Gamma(S,\mathfrak g^\alpha),
  \qquad
  Y\in\Gamma(S,\mathfrak g^{-\alpha}),
\]
and set
\[
  p_+(x)=f\bigl(\exp(xX)\bigr),
  \qquad
  p_-(y)=f\bigl(\exp(yY)\bigr).
\]
These are morphisms of groups
\[
  p_+,p_-\colon\mathbb G_{a,S}\longrightarrow G'.
\]
Now
\[
\begin{aligned}
  \intaut\bigl(\alpha'^\ast(z)\bigr)^q\bigl(p_+(x)\bigr)
    &={}\intaut\bigl(f_T(\alpha^\ast(z))\bigr)
       \bigl(f(\exp(xX))\bigr)\\
    &={}f\Bigl(\intaut\bigl(\alpha^\ast(z)\bigr)
       \bigl(\exp(xX)\bigr)\Bigr)\\
    &={}f\bigl(\exp(z^2xX)\bigr)
     =p_+(z^2x).
\end{aligned}
\]
Applying \SGARefLink{sga3:xx:1.2}{Lemma 1.2} with
\(Q=\alpha'^\ast(\mathbb G_{m,S})\), we obtain a section
\[
  X'\in\Gamma(S,\mathfrak g'^{\,\alpha'})
\]
such that
\[
  f\bigl(\exp(xX)\bigr)=p_+(x)=\exp(x^qX').
\]
Similarly, there exists a section
\[
  Y'\in\Gamma(S,\mathfrak g'^{-\alpha'})
\]
such that
\[
  f\bigl(\exp(yY)\bigr)=\exp(y^qY').
\]
Since \(f\) is a group morphism, it respects formula \textup{(F)};
hence
\[
  X^qY^q=(XY)^q=X'Y'.
\]
It follows readily that \(X^q\mapsto X'\) and \(Y^q\mapsto Y'\)
define the announced isomorphisms \(h\) and \(h^\vee\).

% END INCLUDED SGA3_English_current_buildable_source_20260823/source/inputs/XX/tex/components/11_expose_XX_remark39_through_prop311_en.tex

% BEGIN INCLUDED SGA3_English_current_buildable_source_20260823/source/inputs/XX/tex/components/12_expose_XX_prop312_through_cor44_en.tex
\paragraph{\textbf{Proposition 3.12.}}
\SGARefTarget{sga3:xx:3.12}
\emph{Let \((G,T,\alpha)\) be an elementary \(S\)-system and let
\(w\in N^\times(S)\).  Put}
\[
  \Omega_0=\Omega\cap\intaut(w^{-1})(\Omega).
\]
\emph{Let \(d\) be the function on \(\Omega\) defined by}
\[
  d\bigl(\exp(Y)\cdot t\cdot\exp(X)\bigr)
    =\alpha(t)^{-1}+XY.
\]
\emph{Then \(\Omega_0=\Omega_d\), and, for
\(\exp(Y)\cdot t\cdot\exp(X)\in\Omega_0(S')\), one has the following
formula, where
\(z=d(\exp(Y)\cdot t\cdot\exp(X))\):}
\begin{equation}
\tag{\(\star\)}\label{sga3:xx:proposition:3-12:equation:star}
  \intaut(w)\bigl(\exp(Y)\cdot t\cdot\exp(X)\bigr)
  =
  \exp\bigl(z^{-1}a_\alpha(w)^{-1}X\bigr)\cdot
  t\alpha^\ast(z)\cdot
  \exp\bigl(z^{-1}a_\alpha(w)Y\bigr).
\end{equation}
\emph{Moreover,}
\[
  d\circ\intaut(w)=d^{-1}.
\]

Indeed, one immediately has\footnote{The editors corrected the
original by interchanging \(a_\alpha(w)\) and
\(a_\alpha(w)^{-1}\), and supplied the details of the proof of
\(d\circ\intaut(w)=d^{-1}\).}
\[
\begin{aligned}
  \intaut(w)\bigl(\exp(Y)\cdot t\cdot\exp(X)\bigr)
   &=
   \exp\bigl(a_\alpha(w)Y\bigr)\cdot
   t\alpha^\ast\bigl(\alpha(t)^{-1}\bigr)\cdot
   \exp\bigl(a_\alpha(w)^{-1}X\bigr)\\
   &=
   \exp\bigl(a_\alpha(w)Y\bigr)\cdot
   \exp\bigl(\alpha(t)a_\alpha(w)^{-1}X\bigr)\cdot
   t\alpha^\ast\bigl(\alpha(t)^{-1}\bigr).
\end{aligned}
\]
By \SGARefLink{sga3:xx:2.1}{2.1}, this is a section of \(\Omega\) if and only if
\(1+\alpha(t)XY\) is invertible, which proves
\(\Omega_0=\Omega_d\).  Applying formula \textup{(F)} of the same
paragraph then gives \((\star)\) by an immediate computation.
Finally, \((\star)\) gives
\[
\begin{aligned}
  (d\circ\intaut(w))
    \bigl(\exp(Y)\cdot t\cdot\exp(X)\bigr)
   &=
   \alpha\bigl(t\alpha^\ast(z)\bigr)^{-1}
     +z^{-2}XY\\
   &=z^{-2}\bigl(\alpha(t)^{-1}+XY\bigr)
    =z^{-1},
\end{aligned}
\]
which proves the last assertion.

\smallskip
\noindent\textbf{N.B.}\ The function \(d\) is independent of the
choice of \(w\).

\section{The Isomorphism Theorem}

\paragraph{\textbf{Theorem 4.1.}}
\SGARefTarget{sga3:xx:4.1}
\emph{Let \(S\) be a scheme and let \(q\in\mathbb Z\), \(q>0\), be
such that \(x\mapsto x^q\) is an endomorphism of
\(\mathbb G_{a,S}\).  Let \((G,T,\alpha)\) and
\((G',T',\alpha')\) be two elementary \(S\)-systems.  Let}
\[
  h\colon(\mathfrak g^\alpha)^{\otimes q}
      \xrightarrow{\ \sim\ }\mathfrak g'^{\,\alpha'},
  \qquad
  h^\vee\colon(\mathfrak g^{-\alpha})^{\otimes q}
      \xrightarrow{\ \sim\ }\mathfrak g'^{-\alpha'}
\]
\emph{be mutually contragredient isomorphisms.  Let
\(f_T\colon T\to T'\) be a morphism of \(S\)-groups making the
following diagram commute:}

\begin{center}
  \begin{tikzcd}[column sep=huge,row sep=large]
    \mathbb G_{m,S}
      \arrow[r,"q"]
      \arrow[d,"\alpha^\ast"']
      & \mathbb G_{m,S}
        \arrow[d,"\alpha'^\ast"] \\
    T
      \arrow[r,"f_T"]
      \arrow[d,"\alpha"']
      & T'
        \arrow[d,"\alpha'"] \\
    \mathbb G_{m,S}
      \arrow[r,"q"]
      & \mathbb G_{m,S}
  \end{tikzcd}
\end{center}

\emph{There exists a unique morphism of \(S\)-groups
\(f\colon G\to G'\) extending \(f_T\) and satisfying}
\[
  f\bigl(\exp(X)\bigr)=\exp\bigl(h(X^q)\bigr)
\]
\emph{for every \(S'\to S\) and
\(X\in W(\mathfrak g^\alpha)(S')\).  Moreover, this morphism also
satisfies}
\[
  f\bigl(\exp(Y)\bigr)=\exp\bigl(h^\vee(Y^q)\bigr),
  \qquad
  f\bigl(w_\alpha(Z)\bigr)
    =w_{\alpha'}\bigl(h(Z^q)\bigr)
\]
\emph{for every \(S'\to S\),
\(Y\in\Gamma(S',\mathfrak g^{-\alpha})\), and
\(Z\in\Gamma(S',\mathfrak g^\alpha)^\times\).}

If \(f\colon G\to G'\) extends \(f_T\), then
\[
  f\circ\alpha^\ast=(\alpha'^\ast)^q.
\]
If, in addition, \(f\) satisfies the second displayed condition in
the statement, then it satisfies the other two conditions as well by
\SGARefLink{sga3:xx:3.10}{3.10}.  It follows that \(f\) is determined on \(\Omega\) by
\[
  f\bigl(\exp(Y)t\exp(X)\bigr)
  =
  \exp\bigl(h^\vee(Y^q)\bigr)\,
  f_T(t)\,
  \exp\bigl(h(X^q)\bigr).
\]
Since \(\Omega\) is schematically dense in \(G\), this already proves
the uniqueness of \(f\).  To prove existence, it is enough, by
\SGARefLink{sga3:xviii:prop:2.3}{Exposé XVIII, 2.3}, to show that the preceding formula defines a
``generically multiplicative'' morphism from \(\Omega\) to \(G'\).
By \SGARefLink{sga3:xx:2.4}{2.4}, this amounts to checking
\[
  \alpha'\circ f=\alpha^q,
\]
which follows from the fact that \(f\) extends \(f_T\).

\paragraph{\textbf{Scholium 4.2.}}
\SGARefTarget{sga3:xx:4.2}
\SGARefLink{sga3:xx:4.1}{Theorem 4.1} may also be interpreted as follows.  Let
\(\mathcal E\) be the category of elementary \(S\)-systems, and let
\(\mathcal D\) be the category of tuples
\[
  \bigl(
    \mathbb G_{m,S}\xrightarrow{\alpha^\ast}T
    \xrightarrow{\alpha}\mathbb G_{m,S},\,L
  \bigr),
\]
where \(T\) is a torus, \(\alpha\) and \(\alpha^\ast\) are group
morphisms such that \(\alpha\circ\alpha^\ast=2\), and \(L\) is an
invertible \(\mathcal O_S\)-module.  The reader may specify the
morphisms in the two categories under consideration.  Define a
functor
\[
\begin{aligned}
  \mathcal E&\longrightarrow\mathcal D,\\
  (G,T,\alpha)&\longmapsto
  \bigl(
    \mathbb G_{m,S}\xrightarrow{\alpha^\ast}T
    \xrightarrow{\alpha}\mathbb G_{m,S},\,
    \mathfrak g^\alpha
  \bigr).
\end{aligned}
\]
The preceding theorem says that this functor is fully faithful.  In
fact, it is an equivalence of categories, as will be seen in the next
section.  We already have:

\paragraph{\textbf{Corollary 4.3.}}
\SGARefTarget{sga3:xx:4.3}
\emph{If \(q=1\) and \(f_T\) is an isomorphism, then \(f\) is an
isomorphism.}

\paragraph{\textbf{Corollary 4.4.}}
\SGARefTarget{sga3:xx:4.4}
\emph{If \(q=1\) and \(f_T\) is faithfully flat with kernel \(Q\)
\textup{(cf.\ \SGARefLink{sga3:ix:proposition:2:7}{Exposé IX, 2.7})}, then \(f\) is faithfully flat and
quasi-compact with kernel \(Q\), and hence identifies \(G'\) with
\(G/Q\).}

Indeed, if \(f_T\) is faithfully flat with kernel \(Q\), then
\[
  Q=\Ker(f_T)
  \subset\Ker(\alpha'\circ f_T)
  =\Ker(\alpha).
\]
Introduce the elementary \(S\)-system
\((G/Q,T/Q,\alpha/Q)\) of \SGARefLink{sga3:xx:2.13}{2.13}.  By \SGARefLink{sga3:xx:2.14}{2.14}, we are reduced to proving
that \(f/Q\) induces an isomorphism from \(G/Q\) onto \(G'\), which
follows immediately from \SGARefLink{sga3:xx:4.3}{4.3}.

% Authority: Exp20-13oct24.pdf, proposition 3.12 through corollary 4.4.
% Continuous English stops immediately before section 5.1.

% END INCLUDED SGA3_English_current_buildable_source_20260823/source/inputs/XX/tex/components/12_expose_XX_prop312_through_cor44_en.tex

% BEGIN INCLUDED SGA3_English_current_buildable_source_20260823/source/inputs/XX/tex/components/13_expose_XX_section5_par51_through_par56_en.tex
\section{Examples of Elementary Systems; Applications}

\paragraph{\textbf{5.1.}}
\SGARefTarget{sga3:xx:5.1}
Let \(S\) be a scheme and \(L\) an invertible
\(\mathcal O_S\)-module.  Consider the group \(\mathrm{GL}_L\) over
\(S\) defined by
\[
 \mathrm{GL}_L(S')
 =
 \left\{
 \begin{pmatrix}a&b\\c&d\end{pmatrix}
 \ \middle|\
 \begin{array}{l}
 a,d\in\mathbb G_a(S'),\quad b\in W(L)(S'),\\
 c\in W(L^{-1})(S'),\quad ad-bc\in\mathbb G_m(S')
 \end{array}
 \right\},
\]
equipped with the usual matrix-multiplication law.  It is locally
isomorphic to \(\mathrm{GL}_{2,S}\).  It is therefore an affine smooth
\(S\)-group scheme with connected fibers.

\smallskip
\noindent\textit{Remark.}
Let \(L'\) and \(L''\) be invertible sheaves on \(S\) such that
\[
 L=L'\otimes L''^{-1}.
\]
Then there is an isomorphism of \(S\)-groups
\[
 \mathrm{GL}_L\xrightarrow{\ \sim\ }\mathrm{GL}(L'\oplus L'')
\]
defined as follows: if \(x\) (respectively \(y\)) is a section of
\(L'\) (respectively \(L''\)) on an open subset \(V\) of \(S\), then
\[
 \begin{pmatrix}a&b\\c&d\end{pmatrix}
 \begin{pmatrix}x\\y\end{pmatrix}
 =
 \begin{pmatrix}ax+by\\cx+dy\end{pmatrix}.
\]
\footnote{The editors corrected
\(L''\otimes L'^{-1}\) to \(L'\otimes L''^{-1}\) and expanded the
sentence that follows.}

\paragraph{\textbf{5.2.}}
\SGARefTarget{sga3:xx:5.2}
Denote by \(\mathrm{SL}_L\) the closed subgroup of
\(\mathrm{GL}_L\) defined by \(ad-bc=1\).  It too is an affine smooth
\(S\)-group scheme with connected fibers; under the preceding
isomorphism it is isomorphic to \(\mathrm{SL}(L'\oplus L'')\).

Likewise, consider the morphism
\[
 \mathbb G_{m,S}\longrightarrow\mathrm{GL}_L,
 \qquad
 z\longmapsto\begin{pmatrix}z&0\\0&z\end{pmatrix}.
\]
It is a central monomorphism.  Passage to the quotient gives a smooth
affine \(S\)-group \(\mathrm P_L\) with connected fibers
\textup{(cf.\ \SGARefLink{sga3:viii:corollary:5:7}{Exposé VIII, 5.7})}.  By passing to the quotient in the
isomorphism of the preceding remark, one sees that \(\mathrm P_L\)
identifies with the automorphism group of the projective bundle
\(\mathbb P(L'\oplus L'')\) \textup{(cf.\ EGA II, 4.2.7)}.
Write \(i\) and \(p\) for the canonical morphisms
\[
 \mathrm{SL}_L\xrightarrow{\ i\ }\mathrm{GL}_L
   \xrightarrow{\ p\ }\mathrm P_L.
\]
Here \(i\) is a closed immersion, while \(p\) is faithfully flat and
affine.

\paragraph{\textbf{5.3.}}
\SGARefTarget{sga3:xx:5.3}
Consider the group morphisms
\[
\begin{aligned}
 t_G\colon\mathbb G_{m,S}^{\,2}&\longrightarrow\mathrm{GL}_L,
 &t_G(z,z')&=\begin{pmatrix}z&0\\0&z'\end{pmatrix},\\
 t_S\colon\mathbb G_{m,S}&\longrightarrow\mathrm{SL}_L,
 &t_S(z)&=\begin{pmatrix}z&0\\0&z^{-1}\end{pmatrix},\\
 t_P\colon\mathbb G_{m,S}&\longrightarrow\mathrm P_L,
 &t_P(z)&=p\bigl(t_G(z,1)\bigr).
\end{aligned}
\]
They are group monomorphisms and define in each group a split torus of
relative codimension \(2\).  For \(s\in S\), let
\[
 X\in\Gamma(s,L\otimes s)^\times.
\]
The section
\(\begin{psmallmatrix}0&X\\-X^{-1}&0\end{psmallmatrix}\) of
\(\mathrm{GL}_{L,s}\) normalizes \(t_G(\mathbb G_{m,s}^{\,2})\) and
does not centralize it.  \SGARefLink{sga3:xix:definition:1:6:1}{Exposé XIX, 1.6} therefore shows that
\(\mathrm{GL}_L\) is reductive of semisimple rank \(1\), with maximal
torus \(t_G(\mathbb G_{m,S}^{\,2})\).  The same argument shows that
\(\mathrm{SL}_L\) and \(\mathrm P_L\) are reductive of semisimple rank
\(1\), with respective maximal tori
\(t_S(\mathbb G_{m,S})\) and \(t_P(\mathbb G_{m,S})\).

\paragraph{\textbf{5.4.}}
\SGARefTarget{sga3:xx:5.4}
The Lie algebras of these groups, and the adjoint actions of the
chosen maximal tori, are determined in the usual way.  For
\(\mathrm{GL}_L\), \SGARefLink{sga3:ii:proposition:4-8}{Exposé II, 4.8} immediately gives
\[
 \Lie(\mathrm{GL}_L/S)
 =
 \left\{
 \begin{pmatrix}a&b\\c&d\end{pmatrix}
 \ \middle|\
 a,d\in\mathcal O_S,\ b\in L,\ c\in L^{-1}
 \right\},
\]
with the usual bracket, and
\[
 \operatorname{Ad}\bigl(t_G(z,z')\bigr)
 \begin{pmatrix}a&b\\c&d\end{pmatrix}
 =
 \begin{pmatrix}
 a&zz'^{-1}b\\
 z'z^{-1}c&d
 \end{pmatrix}.
\]
Put \(\mathfrak g=\Lie(\mathrm{GL}_L/S)\), and define
\[
 \alpha_G\colon t_G(\mathbb G_{m,S}^{\,2})\longrightarrow
 \mathbb G_{m,S},
 \qquad
 \alpha_G\bigl(t_G(z,z')\bigr)=zz'^{-1}.
\]
The preceding identity shows that \(\alpha_G\) is a root of
\(\mathrm{GL}_L\) relative to \(t_G(\mathbb G_{m,S}^{\,2})\), and
that
\[
 u\colon L\longrightarrow\mathfrak g,
 \quad X\longmapsto\begin{pmatrix}0&X\\0&0\end{pmatrix},
 \qquad
 u^-\colon L^{-1}\longrightarrow\mathfrak g,
 \quad Y\longmapsto\begin{pmatrix}0&0\\Y&0\end{pmatrix}
\]
identify \(L\) with \(\mathfrak g^{\alpha_G}\) and \(L^{-1}\) with
\(\mathfrak g^{-\alpha_G}\), respectively.  Thus
\[
 \bigl(\mathrm{GL}_L,t_G(\mathbb G_{m,S}^{\,2}),\alpha_G\bigr)
\]
is an elementary \(S\)-system.

Similarly, setting
\[
 \alpha_S(t_S(z))=z^2,
 \qquad
 \alpha_P(t_P(z))=z,
\]
one finds elementary \(S\)-systems
\[
 \bigl(\mathrm{SL}_L,t_S(\mathbb G_{m,S}),\alpha_S\bigr),
 \qquad
 \bigl(\mathrm P_L,t_P(\mathbb G_{m,S}),\alpha_P\bigr),
\]
and corresponding identifications of \(L\) and \(L^{-1}\) with the
root summands of their Lie algebras.

\paragraph{\textbf{5.5.}}
\SGARefTarget{sga3:xx:5.5}
Set
\[
 \exp\begin{pmatrix}0&X\\0&0\end{pmatrix}
   =\begin{pmatrix}1&X\\0&1\end{pmatrix}.
\]
This defines a morphism
\[
 W(\mathfrak g^{\alpha_G})\longrightarrow\mathrm{GL}_L
\]
whose tangent map is the canonical inclusion; it is therefore the
unique morphism of this kind by \SGARefLink{sga3:xx:1.5}{1.5}.  Likewise, set
\[
 \exp\begin{pmatrix}0&0\\Y&0\end{pmatrix}
   =\begin{pmatrix}1&0\\Y&1\end{pmatrix}.
\]
Direct calculation of formula \textup{(F)} gives
\[
 \left\langle
 \begin{pmatrix}0&X\\0&0\end{pmatrix},
 \begin{pmatrix}0&0\\Y&0\end{pmatrix}
 \right\rangle
 =XY,
 \qquad
 \alpha_G^\ast(z)
 =\begin{pmatrix}z&0\\0&z^{-1}\end{pmatrix}
 =t_G(z,z^{-1}).
\]
\footnote{The editors expanded the calculations that follow.}

The open subset \(N^\times=N_G^\times\) defined before \SGARefLink{sga3:xx:3.1}{3.1} is
\[
 N_G^\times(S')
 =
 \left\{
 \begin{pmatrix}0&P\\Q&0\end{pmatrix}
 \ \middle|\
 P\in W(\mathfrak g^\alpha)^\times(S'),\
 Q\in W(\mathfrak g^{-\alpha})^\times(S')
 \right\}.
\]
For \(X\in W(\mathfrak g^\alpha)^\times(S')\),
\[
 w_{\alpha_G}(X)
 =\begin{pmatrix}0&X\\-X^{-1}&0\end{pmatrix}.
\]
Moreover, if
\[
 w=\begin{pmatrix}0&P\\Q&0\end{pmatrix}\in N_G^\times(S'),
\]
then
\[
 a_{\alpha_G}(w)=PQ^{-1}
   \in W\bigl((\mathfrak g^\alpha)^{\otimes2}\bigr)^\times(S');
\]
that is, for \(Y\in W(\mathfrak g^{-\alpha})^\times(S')\),
\[
 a_{\alpha_G}(w)(Y)=PQ^{-1}Y
   \in W(\mathfrak g^\alpha)^\times(S').
\]

\paragraph{\textbf{5.6.}}
\SGARefTarget{sga3:xx:5.6}
We leave the analogous calculations for \(\mathrm{SL}_L\) and
\(\mathrm P_L\) to the reader.  One obtains the same duality formula
and the coroots
\[
 \alpha_S^\ast(z)=t_S(z),
 \qquad
 \alpha_P^\ast(z)=t_P(z^2).
\]
Let \(p_T\) be the morphism induced by
\(p\colon\mathrm{GL}_L\to\mathrm P_L\) on
\(t_S(\mathbb G_{m,S})\), namely
\[
 p_T(t_S(z))=t_P(z^2).
\]
Thus the following diagram is commutative:
\footnote{Here \(t_S\) and \(t_P\) are isomorphisms.}
\begin{center}
\begin{tikzpicture}[>=stealth,
  every node/.style={inner sep=2pt},
  every path/.style={line width=.45pt}]
  \node (top) at (0,3.0) {$\mathbb G_{m,S}$};
  \node (left) at (-5.5,0) {$\mathbb G_{m,S}$};
  \node (ts) at (-1.9,0) {$t_S(\mathbb G_{m,S})$};
  \node (tp) at (1.9,0) {$t_P(\mathbb G_{m,S})$};
  \node (right) at (5.5,0) {$\mathbb G_{m,S}$};
  \node (bottom) at (0,-3.0) {$\mathbb G_{m,S}$};

  \draw[->] (top) -- node[pos=.48,above left] {$\mathrm{id}$} (left);
  \draw[->] (top) -- node[pos=.55,above left] {$\alpha_S^\ast$} (ts);
  \draw[->] (top) -- node[pos=.55,above right] {$\alpha_P^\ast$} (tp);
  \draw[->] (top) -- node[pos=.48,above right] {$2$} (right);

  \draw[->] (left) -- node[above] {$t_S$} (ts);
  \draw[->] (ts) -- node[above] {$p_T$} (tp);
  \draw[->] (right) -- node[above] {$t_P$} (tp);

  \draw[->] (left) -- node[pos=.48,below left] {$2$} (bottom);
  \draw[->] (ts) -- node[pos=.55,below left] {$\alpha_S$} (bottom);
  \draw[->] (tp) -- node[pos=.55,below right] {$\alpha_P$} (bottom);
  \draw[->] (right) -- node[pos=.48,below right] {$\mathrm{id}$} (bottom);
\end{tikzpicture}
\end{center}
Its central portion is the commutative diagram of \SGARefLink{sga3:xx:4.1}{4.1}
\footnote{With \(q=1\).}
associated with the canonical morphism
\(p\circ i\colon\mathrm{SL}_L\to\mathrm P_L\), which induces a
morphism between the preceding elementary \(S\)-systems.

% Authority: Exp20-13oct24.pdf, section 5.1 through 5.6.

% END INCLUDED SGA3_English_current_buildable_source_20260823/source/inputs/XX/tex/components/13_expose_XX_section5_par51_through_par56_en.tex

% BEGIN INCLUDED SGA3_English_current_buildable_source_20260823/source/inputs/XX/tex/components/14_expose_XX_par57_through_theorem511_en.tex
\paragraph{\textbf{5.7.}}
\SGARefTarget{sga3:xx:5.7}
Let \((G,T,\alpha)\) now be an arbitrary elementary \(S\)-system.
Consider the commutative diagram
\begin{center}
\begin{tikzpicture}[>=stealth,
  every node/.style={inner sep=2pt},
  every path/.style={line width=.45pt}]
  \node (top) at (0,2.9) {$\mathbb G_{m,S}$};
  \node (left) at (-3.7,0) {$\mathbb G_{m,S}$};
  \node (mid) at (0,0) {$T$};
  \node (right) at (3.7,0) {$\mathbb G_{m,S}$};
  \node (bottom) at (0,-2.9) {$\mathbb G_{m,S}$};

  \draw[->] (top) -- node[pos=.48,above left] {$\mathrm{id}$} (left);
  \draw[->] (top) -- node[left] {$\alpha^\ast$} (mid);
  \draw[->] (top) -- node[pos=.48,above right] {$2$} (right);
  \draw[->] (left) -- node[above] {$\alpha^\ast$} (mid);
  \draw[->] (mid) -- node[above] {$\alpha$} (right);
  \draw[->] (left) -- node[pos=.48,below left] {$2$} (bottom);
  \draw[->] (mid) -- node[left] {$\alpha$} (bottom);
  \draw[->] (right) -- node[pos=.48,below right] {$\mathrm{id}$} (bottom);
\end{tikzpicture}
\end{center}
Combining this diagram with that of \SGARefLink{sga3:xx:5.6}{5.6} gives a commutative diagram
\begin{center}
\begin{tikzpicture}[>=stealth,
  every node/.style={inner sep=2pt},
  every path/.style={line width=.45pt}]
  \node (top) at (0,2.9) {$\mathbb G_{m,S}$};
  \node (left) at (-4.8,0) {$t_S(\mathbb G_{m,S})$};
  \node (mid) at (0,0) {$T$};
  \node (right) at (4.8,0) {$t_P(\mathbb G_{m,S})$};
  \node (bottom) at (0,-2.9) {$\mathbb G_{m,S}$};

  \draw[->] (top) -- node[pos=.5,above left] {$\alpha_S^\ast$} (left);
  \draw[->] (top) -- node[left] {$\alpha^\ast$} (mid);
  \draw[->] (top) -- node[pos=.5,above right] {$\alpha_P^\ast$} (right);
  \draw[->] (left) -- node[above] {$\alpha^\ast\circ t_S^{-1}$} (mid);
  \draw[->] (mid) -- node[above] {$t_P\circ\alpha$} (right);
  \draw[->] (left) -- node[pos=.5,below left] {$\alpha_S$} (bottom);
  \draw[->] (mid) -- node[left] {$\alpha$} (bottom);
  \draw[->] (right) -- node[pos=.5,below right] {$\alpha_P$} (bottom);
\end{tikzpicture}
\end{center}
Using \SGARefLink{sga3:xx:4.1}{4.1}, one obtains:

\paragraph{\textbf{Proposition 5.8.}}
\SGARefTarget{sga3:xx:5.8}
\emph{Let \(S\) be a scheme and let \((G,T,\alpha)\) be an
elementary \(S\)-system.  Put \(L=\mathfrak g^\alpha\), so that
\(L^{-1}=\mathfrak g^{-\alpha}\).}

\emph{\textup{(i)} There exists a unique group morphism
\(f\colon\mathrm{SL}_L\to G\) satisfying the following equivalent
conditions:}
\[
\begin{array}{ll}
\textup{(a)}
&
f\begin{pmatrix}z&0\\0&z^{-1}\end{pmatrix}=\alpha^\ast(z),
\qquad
f\begin{pmatrix}1&X\\0&1\end{pmatrix}=\exp(X);
\\[1.1em]
\textup{(b)}
&
f\begin{pmatrix}1&X\\0&1\end{pmatrix}=\exp(X),
\qquad
f\begin{pmatrix}1&0\\Y&1\end{pmatrix}=\exp(Y);
\\[1.1em]
\textup{(c)}
&
f\begin{pmatrix}1&X\\0&1\end{pmatrix}=\exp(X),
\qquad
f\begin{pmatrix}0&X\\-X^{-1}&0\end{pmatrix}=w_\alpha(X).
\end{array}
\]

\emph{\textup{(ii)} There exists a unique group morphism
\(g\colon G\to\mathrm P_L\) satisfying}
\[
 g(t)=
 \begin{pmatrix}\alpha(t)&0\\0&1\end{pmatrix},
 \qquad
 g(\exp(X))
 =p\begin{pmatrix}1&X\\0&1\end{pmatrix}.
\]
\emph{Moreover,}
\[
 g(\exp(Y))
 =p\begin{pmatrix}1&0\\Y&1\end{pmatrix},
 \qquad
 g(w_\alpha(X))
 =p\begin{pmatrix}0&X\\-X^{-1}&0\end{pmatrix}.
\]
\emph{The morphism \(g\) is faithfully flat and quasi-compact, with}
\[
 \Ker(g)=\Ker(\alpha)=\operatorname{Centr}(G),
\]
\emph{and \(g\circ f\) is the canonical morphism
\(\mathrm{SL}_L\to\mathrm P_L\).}

Notice that condition \textup{(b)} in \textup{(i)} gives an explicit
description of the duality between
\(\mathfrak g^\alpha\) and \(\mathfrak g^{-\alpha}\).

\paragraph{\textbf{Corollary 5.9.}}
\SGARefTarget{sga3:xx:5.9}
\emph{Let \((G,T,\alpha)\) be an elementary \(S\)-system.  The
subgroups}
\[
 T\cdot U_\alpha,\qquad T\cdot U_{-\alpha},
 \qquad U_\alpha,\qquad U_{-\alpha}
\]
\emph{are closed.}

Since \(U_\alpha\) is a closed subgroup scheme of
\(T\cdot U_\alpha\), it suffices to prove the assertion for the
latter.  By Noether's theorem \textup{(\SGARefLink{sga3:iv:proposition:5:3:1}{Exposé IV, 5.3.1} and \SGARefLink{sga3:iv:lemma:6:4:1}{6.4.1})},
it is enough to show that
\[
 (T\cdot U_\alpha)/\Ker(\alpha)
\]
is a closed subgroup of \(G/\Ker(\alpha)\).  By \SGARefLink{sga3:xx:5.8}{5.8}, this reduces to
the evident fact that the subgroup of \(\mathrm P_L\) defined by
\(c=0\) is closed; one may equivalently work in \(\mathrm{GL}_L\),
by another application of Noether's theorem.  Consequently, the
exponential morphisms of \SGARefLink{sga3:xx:1.5}{Theorem 1.5}\textup{(i)} are closed
immersions.

\smallskip
\noindent\textbf{N.B.}\ The corollary also follows from the fact that
\(T\cdot U_\alpha\) and \(T\cdot U_{-\alpha}\) are ``Borel
subgroups'' of \(G\) \textup{(cf.\ \SGARefLink{sga3:xii:7.10}{Exposé XII, 7.10})}.

\paragraph{\textbf{5.10.}}
\SGARefTarget{sga3:xx:5.10}
Let \(L\) be an invertible \(\mathcal O_S\)-module, and let
\[
 \mathbb G_{m,S}\xrightarrow{\ \alpha^\ast\ }T
 \xrightarrow{\ \alpha\ }\mathbb G_{m,S}
\]
be a diagram of groups, where \(T\) is a torus,\footnote{This
qualification is supplied by the editors.} such that
\(\alpha\circ\alpha^\ast=2\).  Let \(R\) be the maximal torus of
\(\Ker(\alpha)\), and put \(K=\alpha^{\ast-1}(R)\).  Then \(K\) is a
subgroup of multiplicative type of \(\mathbb G_{m,S}\); since
\(\alpha\circ\alpha^\ast=2\), it is even a subgroup of
\(\mu_{2,S}\).  In particular,
\[
 K\longrightarrow\mathrm{SL}_L,
 \qquad
 z\longmapsto\begin{pmatrix}z&0\\0&z^{-1}\end{pmatrix}
\]
is central.  Hence there is a central group monomorphism
\[
 K\longrightarrow R\times\mathrm{SL}_L,
 \qquad
 z\longmapsto
 \left(
 \alpha^\ast(z),
 \begin{pmatrix}z&0\\0&z^{-1}\end{pmatrix}
 \right).
\]

Let
\[
 G=(R\times\mathrm{SL}_L)/K.
\]
This is an affine smooth group over \(S\), with connected fibers.
The sequence
\[
 1\longrightarrow K\longrightarrow
 R\times t_S(\mathbb G_{m,S})
 \xrightarrow{\ u\ }T\longrightarrow1,
 \qquad
 u(x,t_S(z))=x\alpha^\ast(z),
\]
is exact.  The image of \(R\times t_S(\mathbb G_{m,S})\) in \(G\)
is therefore a torus \(T'\) isomorphic to \(T\).  If \(\alpha'\) is
the character of \(T'\) transported from \(\alpha\), one readily
checks that \((G,T',\alpha')\) is an elementary \(S\)-system, that
\(\mathfrak g^{\alpha'}\) is isomorphic to \(L\), and that
\(\alpha'^\ast\) is obtained from \(\alpha^\ast\) under the
isomorphism \(T\xrightarrow{\sim}T'\).  Thus the object
\[
 \left(
 \mathbb G_{m,S}\xrightarrow{\ \alpha'^\ast\ }T'
 \xrightarrow{\ \alpha'\ }\mathbb G_{m,S},\,
 \mathfrak g^{\alpha'}
 \right)
\]
of the category \(\mathcal D\) defined in \SGARefLink{sga3:xx:4.2}{4.2} is isomorphic to
\[
 \left(
 \mathbb G_{m,S}\xrightarrow{\ \alpha^\ast\ }T
 \xrightarrow{\ \alpha\ }\mathbb G_{m,S},\,L
 \right).
 \]
We have therefore proved:

\paragraph{\textbf{Theorem 5.11.}}
\SGARefTarget{sga3:xx:5.11}
\emph{With the notation of \SGARefLink{sga3:xx:4.2}{4.2}, the functor}
\[
 (G,T,\alpha)\longmapsto
 \left(
 \mathbb G_{m,S}\xrightarrow{\ \alpha^\ast\ }T
 \xrightarrow{\ \alpha\ }\mathbb G_{m,S},\,
 \mathfrak g^\alpha
 \right)
\]
\emph{is an equivalence of categories between
\(\mathcal E\) and \(\mathcal D\).}

% Authority: Exp20-13oct24.pdf, section 5.7 through theorem 5.11.

% END INCLUDED SGA3_English_current_buildable_source_20260823/source/inputs/XX/tex/components/14_expose_XX_par57_through_theorem511_en.tex

% BEGIN INCLUDED SGA3_English_current_buildable_source_20260823/source/inputs/XX/tex/components/15_expose_XX_section6_through_lemma621_en.tex
\section{Generators and Relations for an Elementary System}

\paragraph{\textbf{6.1.}}
\SGARefTarget{sga3:xx:6.1}
Let \(S\) be a scheme and let \((G,T,\alpha)\) be an elementary
\(S\)-system.  Let
\[
 X\in W(\mathfrak g^\alpha)^\times(S),
 \qquad u=\exp(X).
\]
As observed in \SGARefLink{sga3:xx:3.8}{3.8}, the element \(w=w_\alpha(X)\) satisfies, in
particular,
\[
 (wu)^3=e.
\]
\footnote{The editors added the sentence that follows.}
Let \(s_\alpha\) denote the automorphism of \(T\) induced by
\(\intaut(w)\).  By \SGARefLink{sga3:xx:3.1}{Theorem 3.1}\textup{\SGARefLink{sga3:xx:3.1:item:iii}{(iii)}}, for every \(S'\to S\)
and every \(t\in T(S')\),
\[
 s_\alpha(t)=\intaut(w)(t)
 =t\,\alpha^\ast\bigl(\alpha(t)^{-1}\bigr).
\]

\paragraph{\textbf{Theorem 6.2.}}
\SGARefTarget{sga3:xx:6.2}
\emph{Let \(H\) be a sheaf of groups over \(S\) for the fppf
topology.  Let}
\[
 f_T\colon T\longrightarrow H,
 \qquad
 f_\alpha\colon U_\alpha\longrightarrow H
\]
\emph{be group morphisms, and let \(h\in H(S)\).  There exists a
\textup{(necessarily unique)} group morphism}
\[
 f\colon G\longrightarrow H
\]
\emph{extending \(f_T\) and \(f_\alpha\) and satisfying \(f(w)=h\)
if and only if the following conditions hold:}

\emph{\textup{(i)} For every \(S'\to S\),
\(t\in T(S')\), and \(x\in U_\alpha(S')\),}
\begin{equation}
\tag{1}\label{sga3:xx:theorem:6-2:equation:1}
 f_T(t)f_\alpha(x)f_T(t)^{-1}
 =f_\alpha(txt^{-1})
 =f_\alpha\bigl(x^{\alpha(t)}\bigr).
\end{equation}
\emph{Equivalently, \(f_T\) and \(f_\alpha\) extend to a group
morphism \(T\cdot U_\alpha\to H\).}

\emph{\textup{(ii)} For every \(S'\to S\) and \(t\in T(S')\),}
\begin{equation}
\tag{2}\label{sga3:xx:theorem:6-2:equation:2}
 h f_T(t)h^{-1}
 =f_T(s_\alpha(t))
 =f_T\!\left(t\,\alpha^\ast(\alpha(t)^{-1})\right).
\end{equation}

\emph{\textup{(iii)} The following two relations hold in \(H(S)\):}
\begin{equation}
\tag{3}\label{sga3:xx:theorem:6-2:equation:3}
 h^2=f_T\bigl(\alpha^\ast(-1)\bigr),
\end{equation}
\begin{equation}
\tag{4}\label{sga3:xx:theorem:6-2:equation:4}
 \bigl(hf_\alpha(u)\bigr)^3=e.
\end{equation}

\smallskip
\noindent\emph{Proof.}
Write the laws on \(U_\alpha\) and \(U_{-\alpha}\) additively, while
writing their vector-space structures multiplicatively.  If \(f\)
has the asserted properties, then for
\(y\in U_{-\alpha}(S')\) one necessarily has
\[
 f(y)
 =f(w^{-1}wyw^{-1}w)
 =h f_\alpha(w^{-1}yw)h^{-1}.
\]
Define a morphism \(f_{-\alpha}\colon U_{-\alpha}\to H\) by
\begin{equation}
\tag{\(\ast_1\)}\label{sga3:xx:theorem:6-2:equation:ast-1}
 f_{-\alpha}(y)=h f_\alpha(w^{-1}yw)h^{-1}.
\end{equation}
It is a group morphism.  On the big cell \(\Omega\), the morphism
\(f\) is forced to satisfy
\[
 f(ytx)=f_{-\alpha}(y)f_T(t)f_\alpha(x).
\]
This proves uniqueness.  The conditions in the statement are
manifestly necessary; let us prove that they are sufficient.

Relation \textup{\SGARefLink{sga3:xx:theorem:6-2:equation:4}{(4)}} gives
\[
 h f_\alpha(u)h^{-1}h^2
 =
 f_\alpha(-u)h^{-1}f_\alpha(-u).
\]
By \textup{\SGARefLink{sga3:xx:theorem:6-2:equation:3}{(3)}} and \textup{\SGARefLink{sga3:xx:theorem:6-2:equation:1}{(1)}},
\(h^2=f_T(\alpha^\ast(-1))\) commutes with \(f_\alpha(-u)\), and
hence
\[
 h f_\alpha(u)h^{-1}
 =
 f_\alpha(-u)h f_\alpha(-u).
\]
By definition,
\[
 h f_\alpha(u)h^{-1}=f_{-\alpha}(wuw^{-1}),
\]
whereas \SGARefLink{sga3:xx:3.7}{3.7} gives
\begin{equation}
\tag{\(\ast_2\)}\label{sga3:xx:theorem:6-2:equation:ast-2}
 wuw^{-1}=-\widetilde u,
\end{equation}
with \(\widetilde u\) the element paired with \(u\).  Therefore
\begin{equation}
\tag{\(\ast_3\)}\label{sga3:xx:theorem:6-2:equation:ast-3}
 f_{-\alpha}(-\widetilde u)
 =
 f_\alpha(-u)h f_\alpha(-u).
\end{equation}

Let \(t\) be a section of \(T\) over a variable \(S'\to S\).
Apply \(\intaut(f_T(t))\) to this identity.  On the left one
obtains\footnote{The editors corrected the calculation that follows.}
\[
\begin{aligned}
 f_T(t)f_{-\alpha}(-\widetilde u)f_T(t)^{-1}
 &=
 f_T(t)h f_\alpha(u)h^{-1}f_T(t)^{-1}\\
 &=
 h\bigl(h^{-1}f_T(t)h\bigr)f_\alpha(u)
   \bigl(h^{-1}f_T(t)^{-1}h\bigr)h^{-1}\\
 &=
 h f_T(s_\alpha(t))f_\alpha(u)
   f_T(s_\alpha(t))^{-1}h^{-1}\\
 &=
 h f_\alpha\bigl(\alpha(s_\alpha(t))u\bigr)h^{-1}\\
 &=
 h f_\alpha\bigl(\alpha(t)^{-1}u\bigr)h^{-1}\\
 &=
 f_{-\alpha}\bigl(-\alpha(t)^{-1}\widetilde u\bigr).
\end{aligned}
\]
Here we used \textup{\SGARefLink{sga3:xx:theorem:6-2:equation:2}{(2)}}, \textup{\SGARefLink{sga3:xx:theorem:6-2:equation:1}{(1)}},
\(s_\alpha(t)=t\alpha^\ast(\alpha(t)^{-1})\), and
\(\alpha\circ\alpha^\ast=2\).

The right-hand side of \((\ast_3)\), after conjugation, becomes
\[
 f_\alpha(-\alpha(t)u)\,
 f_T(t)h f_T(t)^{-1}h^{-1}\,
 h f_\alpha(-\alpha(t)u).
\]
Since
\[
 h f_T(t)^{-1}h^{-1}
 =f_T(s_\alpha(t^{-1}))
  =f_T\bigl(t^{-1}\,\alpha^\ast(\alpha(t))\bigr),
\]
this equals
\[
 f_\alpha(-\alpha(t)u)\,
 f_T\bigl(\alpha^\ast(\alpha(t))\bigr)\,
 h f_\alpha(-\alpha(t)u).
\]
Comparison gives
\[
 f_{-\alpha}\bigl(-\alpha(t)^{-1}\widetilde u\bigr)
 =
 f_\alpha(-\alpha(t)u)\,
 f_T\bigl(\alpha^\ast(\alpha(t))\bigr)\,
 h f_\alpha(-\alpha(t)u).
\]
Since \(\alpha\colon T\to\mathbb G_{m,S}\) is faithfully flat and
\(H\) is a separated presheaf, \(-\alpha(t)^{-1}\) may be replaced
by an arbitrary section of \(\mathbb G_{m,S}\).  Thus:

\paragraph{\textbf{Lemma 6.2.1.}}
\SGARefTarget{sga3:xx:6.2.1}
\emph{For every \(S'\to S\) and every
\(z\in\mathbb G_m(S')\),}
\[
 f_{-\alpha}(z\widetilde u)
 =
 f_\alpha(z^{-1}u)\,
 f_T\bigl(\alpha^\ast(-z^{-1})\bigr)\,
 h f_\alpha(z^{-1}u).
\]

% Authority: Exp20-13oct24.pdf, section 6 through lemma 6.2.1.

% END INCLUDED SGA3_English_current_buildable_source_20260823/source/inputs/XX/tex/components/15_expose_XX_section6_through_lemma621_en.tex

% BEGIN INCLUDED SGA3_English_current_buildable_source_20260823/source/inputs/XX/tex/components/16_expose_XX_lemma622_through_end_en.tex
Let \(x,y\in\mathbb G_a(S')\), where \(S'\to S\), and suppose that
\(y\) and \(1+xy\) are invertible.  Applying \SGARefLink{sga3:xx:6.2.1}{Lemma 6.2.1} first with
\(z=y\), one obtains\footnote{The editors expanded the calculations
that follow.}
\[
 f_\alpha(xu)f_{-\alpha}(y\widetilde u)
 =
 f_\alpha\bigl((x+y^{-1})u\bigr)\,
 f_T\bigl(\alpha^\ast(-y^{-1})\bigr)\,
 h f_\alpha(y^{-1}u).
\]
Now \(x+y^{-1}=y^{-1}(1+xy)\).  Applying the lemma with
\(z=y/(1+xy)\) gives
\[
\begin{aligned}
 f_\alpha\bigl((x+y^{-1})u\bigr)
 &=
 f_{-\alpha}\!\left(\frac{y}{1+xy}\widetilde u\right)
 f_\alpha\bigl(-(x+y^{-1})u\bigr)\\
 &\qquad{}\cdot h\,
 f_T\!\left(\alpha^\ast\!\left(-\frac{y}{1+xy}\right)\right).
\end{aligned}
\]
Substitution in the preceding identity yields
\[
\begin{aligned}
 f_\alpha(xu)f_{-\alpha}(y\widetilde u)
 &=
 f_{-\alpha}\!\left(\frac{y}{1+xy}\widetilde u\right)
 f_\alpha\bigl(-(x+y^{-1})u\bigr)h^{-1}\\
 &\qquad{}\cdot
 f_T\bigl(\alpha^\ast(1+xy)^{-1}\bigr)
 h f_\alpha(y^{-1}u).
\end{aligned}
\]
Since \(h^{-1}f_T(t)h=f_T(s_\alpha(t))\) by \textup{(2)}
\textup{(recall that \(s_\alpha^2=\operatorname{id}\))}, and
\(s_\alpha\circ\alpha^\ast=-\alpha^\ast\), this is
\[
\begin{aligned}
 f_{-\alpha}\!\left(\frac{y}{1+xy}\widetilde u\right)
 f_\alpha\bigl(-(x+y^{-1})u\bigr)
 f_T\bigl(\alpha^\ast(1+xy)\bigr)
 f_\alpha(y^{-1}u).
\end{aligned}
\]
For \(x'\in U_\alpha(S')\) and
\(z\in\mathbb G_m(S')\), one has
\[
 f_\alpha(x')f_T(\alpha^\ast(z))
 =
 f_T(\alpha^\ast(z))f_\alpha(z^{-2}x').
\]
Consequently,
\[
\begin{aligned}
 f_\alpha(xu)f_{-\alpha}(y\widetilde u)
 &=
 f_{-\alpha}\!\left(\frac{y}{1+xy}\widetilde u\right)
 f_T\bigl(\alpha^\ast(1+xy)\bigr)\\
 &\qquad{}\cdot
 f_\alpha\!\left(
 \left[
 -\frac{y^{-1}(1+xy)}{(1+xy)^2}+y^{-1}
 \right]u
 \right)\\
 &=
 f_{-\alpha}\!\left(\frac{y}{1+xy}\widetilde u\right)
 f_T\bigl(\alpha^\ast(1+xy)\bigr)
 f_\alpha\!\left(\frac{x}{1+xy}u\right).
\end{aligned}
\]
We have proved:

\paragraph{\textbf{Lemma 6.2.2.}}
\SGARefTarget{sga3:xx:6.2.2}
\emph{Let \(S'\to S\).  If
\(a\in U_\alpha(S')\), \(b\in U_{-\alpha}^\times(S')\), and
\(1+ab\in\mathbb G_m(S')\), then}
\[
 f_\alpha(a)f_{-\alpha}(b)
 =
 f_{-\alpha}\!\left(\frac{b}{1+ab}\right)
 f_T\bigl(\alpha^\ast(1+ab)\bigr)
 f_\alpha\!\left(\frac{a}{1+ab}\right).
\]

By schematic density, this formula remains valid for
\(b\in U_{-\alpha}(S')\), with \(1+ab\) invertible.  Define
\[
 f\colon\Omega\longrightarrow H,
 \qquad
 f(ytx)=f_{-\alpha}(y)f_T(t)f_\alpha(x).
\]
It follows at once from \SGARefLink{sga3:xx:6.2.2}{Lemma 6.2.2}, condition \SGARefLink{sga3:xx:6.2}{6.2}\textup{(i)}, and
formula \textup{(F')} of \SGARefLink{sga3:xx:2.4}{2.4} that whenever
\(g,g'\in\Omega(S')\) and \(gg'\in\Omega(S')\), one has
\[
 f(gg')=f(g)f(g').
\]
By \SGARefLink{sga3:xviii:prop:2.3}{Exposé XVIII, 2.3}\textup{\SGARefLink{sga3:xviii:prop:2.3:item:iii}{(iii)}} and \SGARefLink{sga3:xx:2.4}{2.4},\footnote{Each geometric
fiber of \(G\) is connected, for example by \SGARefLink{sga3:xx:section:1}{1.1}.} there is therefore
a group morphism \(G\to H\) extending \(f\).  Denote it again by
\(f\).  It remains only to prove that \(f(w_\alpha)=h\).  Since
\[
 w_\alpha=u\cdot(-\widetilde u)\cdot u,
\]
formula \((\ast_3)\) gives\footnote{The editors simplified the
original argument by invoking \((\ast_3)\).}
\[
 f(w_\alpha)
 =f_\alpha(u)f_{-\alpha}(-\widetilde u)f_\alpha(u)
 =h.
\]
This completes the proof of \SGARefLink{sga3:xx:6.2}{Theorem 6.2}.

\paragraph{\textbf{Remark 6.3.}}
\SGARefTarget{sga3:xx:6.3}
\emph{These results will be completed in \SGARefLink{sga3:xxiii:subsection:3-5}{Exposé XXIII, 3.5}.}

% Authority: Exp20-13oct24.pdf, lemma 6.2.2 through the end of Exposé XX.

% END INCLUDED SGA3_English_current_buildable_source_20260823/source/inputs/XX/tex/components/16_expose_XX_lemma622_through_end_en.tex



% END INCLUDED SGA3_English_current_buildable_source_20260823/source/inputs/XX/body.tex


\SGAStartExpose{XXI}
\counterwithin{theorem}{subsection}
\renewcommand{\thesubsection}{\thesection.\arabic{subsection}}
% BEGIN INCLUDED SGA3_English_current_buildable_source_20260823/source/inputs/XXI/body.tex
% Mechanically extracted after \mainmatter from the frozen standalone master.

\markboth{Exposé XXI. Root Data}{Exposé XXI. Root Data}
% BEGIN INCLUDED SGA3_English_current_buildable_source_20260823/source/inputs/XXI/tex/components/00_expose_XXI_opening_through_scholium1113_en.tex
\chapter*{Exposé XXI. Root Data}
\addcontentsline{toc}{chapter}{Exposé XXI. Root Data}

\noindent
This Exposé\footnote{Editors' note: Version of 13 October 2024.}
collects, in the absence of a suitable reference,\footnote{%
Editors' note: For results on root systems (Sections 1--5), one may
consult \emph{Groupes et algèbres de Lie}, Chapter VI.}
known results on root data (that is, ``abstract'' root systems), most of
which will be used below.

\medskip
\noindent\textbf{Notation.}
We denote by \(\Qp\) the set of positive (or zero) rational numbers;
\(\mathbb Z\cap\Qp=\mathbb N\). Let \(V\) be a \(\mathbb Q\)-vector
space. If \(A\) (respectively \(B\)) is a subset of \(\mathbb Q\)
(respectively of \(V\)), we denote by \(A\cdot B\) the image of
\(A\otimes B\) under the morphism
\(\mathbb Q\otimes V\to V\), in other words the set of linear
combinations of elements of \(B\) with coefficients in \(A\). We write
\(-B=\{-1\}B\). We denote by \(E-F\) the set of elements of \(E\)
that do not belong to \(F\).

\section{Generalities}

\SGAInlineRefTarget{sga3:xxi:subsection:1-1}
\subsection{Definitions and first properties}

\begin{definition}\label{sga3:XXI:1.1.1}
Let \(M\) and \(M^\ast\) be two free finite-type \(\mathbb Z\)-modules
in duality. Write
\[
  V=M\otimes\mathbb Q,\qquad V^\ast=M^\ast\otimes\mathbb Q;
\]
these are two \(\mathbb Q\)-vector spaces in duality. We identify \(M\)
(respectively \(M^\ast\)) with a subset of \(V\) (respectively
\(V^\ast\)). The canonical bilinear form on
\(M^\ast\times M\) (respectively \(V^\ast\times V\)) is denoted
\((\, ,\,)\).

Let \(R\) be a finite subset of \(M\). Suppose given a map
\(\alpha\mapsto\alpha^\ast\) from \(R\) to \(M^\ast\); the set of
\(\alpha^\ast\), for \(\alpha\in R\), is denoted \(R^\ast\).
To every \(\alpha\in R\) associate the endomorphism \(s_\alpha\)
(respectively \(s_\alpha^\ast\)) of \(M\) and \(V\) (respectively of
\(M^\ast\) and \(V^\ast\)) given by
\begin{align}
 s_\alpha(x)&=x-(\alpha^\ast,x)\alpha,
 &s_\alpha&=\operatorname{id}-\alpha^\ast\otimes\alpha,
 \tag{1}\label{sga3:XXI:eq:1}\\
 s_\alpha^\ast(u)&=u-(u,\alpha)\alpha^\ast,
 &s_\alpha^\ast&=\operatorname{id}-\alpha\otimes\alpha^\ast.
 \tag{1${}^\ast$}\label{sga3:XXI:eq:1star}
\end{align}
The pair \((R,R^\ast)\) (more precisely, the pair
\((R,R\to M^\ast)\)) is called a \emph{root datum} in
\((M,M^\ast)\), and \((M,M^\ast,R,R^\ast)\) is called a root datum,
if the following axioms hold:
\begin{align*}
 (\mathrm{DR\,I})\quad&
    (\alpha^\ast,\alpha)=2
    &&\text{for every }\alpha\in R,\\
 (\mathrm{DR\,II})\quad&
    s_\alpha(R)\subset R,\qquad
    s_\alpha^\ast(R^\ast)\subset R^\ast
    &&\text{for every }\alpha\in R.
\end{align*}
The set \(R\) is called the \emph{root system} of the root datum
\(\cR=(M,M^\ast,R,R^\ast)\). The elements of \(R\) (respectively
\(R^\ast\)) are called the \emph{roots} (respectively
\emph{coroots}) of the root datum.
\end{definition}

\begin{remark}\label{sga3:XXI:1.1.2}
Axiom \((\mathrm{DR\,I})\) is equivalent to any one of the following
properties:
\begin{align}
 s_\alpha s_\alpha&=\operatorname{id},
 &s_\alpha^\ast s_\alpha^\ast&=\operatorname{id},
 \tag{2,2${}^\ast$}\label{sga3:XXI:eq:2}\\
 s_\alpha(\alpha)&=-\alpha,
 &s_\alpha^\ast(\alpha^\ast)&=-\alpha^\ast.
 \tag{3,3${}^\ast$}\label{sga3:XXI:eq:3}
\end{align}
\end{remark}

\begin{remark}\label{sga3:XXI:1.1.3}
Axioms \((\mathrm{DR\,I})\) and \((\mathrm{DR\,II})\) imply
\[
 R=-R,\qquad R^\ast=-R^\ast,\qquad
 0\notin R,\qquad 0\notin R^\ast.
\]
\end{remark}

\begin{lemma}\label{sga3:XXI:1.1.4}
The map \(R\to R^\ast\) is a bijection. More generally, if
\(\alpha,\beta\in R\) and
\((\alpha^\ast,x)=(\beta^\ast,x)\) for every \(x\in R\), then
\(\alpha=\beta\).
\end{lemma}

\begin{proof}
Indeed, one then has
\[
 s_\beta(\alpha)=\alpha-2\beta,\qquad
 s_\alpha(\beta)=\beta-2\alpha.
\]
It follows at once that
\[
 s_\beta s_\alpha(\alpha)
   =2\beta-\alpha=\alpha+2(\beta-\alpha),\qquad
 s_\beta s_\alpha(\beta-\alpha)
   =s_\beta(\beta-\alpha)=\beta-\alpha.
\]
Hence, by \((\mathrm{DR\,II})\),
\[
 (s_\beta s_\alpha)^n(\alpha)
   =\alpha+2n(\beta-\alpha)\in R.
\]
Since \(R\) is finite, \(\beta-\alpha=0\).
\end{proof}

\begin{corollary}\label{sga3:XXI:1.1.5}
The inverse map \(R^\ast\to R\) defines a root datum
\[
 \cR^\ast=(M^\ast,M,R^\ast,R),
\]
called the \emph{dual} of \(\cR\).\footnote{Editors' note:
\((\alpha^\ast)^\ast=\alpha\).}
\end{corollary}

\begin{definition}\label{sga3:XXI:1.1.6}
We denote by \(\Gamma_0(R)\) the subgroup of \(M\) generated by \(R\).
We denote by \(\VR(R)\) the vector subspace of \(V\) generated by \(R\),
that is,
\[
 \VR(R)=\Gamma_0(R)\otimes\mathbb Q.
\]
Applying these definitions to \(\cR^\ast\), one likewise obtains
\(\Gamma_0(R^\ast)\) and \(\VR(R^\ast)\).

The \emph{reductive rank} of \(\cR\) is the number
\[
 \rgred(\cR)
 =\rank(M)=\dim(V)=\dim(V^\ast)=\rank(M^\ast)=\rgred(\cR^\ast).
\]
The \emph{semisimple rank} of \(\cR\) is the number
\[
 \rgss(\cR)
 =\rank(R)=\rank(\Gamma_0(R))=\dim(\VR(R)).
\]
Thus \(\rgss(\cR)\leq\rgred(\cR)\).

We shall see below that
\(\rgss(\cR)=\rgss(\cR^\ast)\), that is, that \(\VR(R)\) and
\(\VR(R^\ast)\) have the same dimension.
\end{definition}

\begin{definition}\label{sga3:XXI:1.1.7}\leavevmode\par
The root datum \(\cR\) is called \emph{semisimple} (respectively
\emph{trivial}) if
\(\rgss(\cR)=\rgred(\cR)\) (respectively \(\rgss(\cR)=0\)).
Thus \(\cR\) is trivial if and only if \(R\) is empty. The trivial root
datum of reductive rank zero is denoted
\[
 0=(\{0\},\{0\},\varnothing,\varnothing).
\]
\end{definition}

\begin{definition}\label{sga3:XXI:1.1.8}
We denote by \(W(\cR)\) the group of transformations of \(M\) generated
by the \(s_\alpha\), \(\alpha\in R\). It is called the \emph{Weyl
group} of \(\cR\). We write
\[
 W^\ast(\cR)=W(\cR^\ast).
\]
Then \(W(\cR)\) acts on \(R\), \(\Gamma_0(R)\), \(\VR(R)\), \(M\),
and \(V\). If \(w\in W(\cR)\) and \(x\in M\) (respectively
\(x\in V\)), then
\[
 wx-x\in\Gamma_0(R)
 \quad\text{(respectively }wx-x\in\VR(R)\text{)};
\]
this follows immediately from formula \eqref{sga3:XXI:eq:1}.
Likewise for \(W^\ast(\cR)\).
\end{definition}

\begin{lemma}\label{sga3:XXI:1.1.9}
For every \(\alpha\in R\), \(x\in V\), and \(u\in V^\ast\), one has
\begin{equation}
 (s_\alpha^\ast(u),s_\alpha(x))=(u,x).
 \tag{4}\label{sga3:XXI:eq:4}
\end{equation}
\end{lemma}

\begin{proof}
Taking the scalar product of \eqref{sga3:XXI:eq:1} and
\eqref{sga3:XXI:eq:1star}, the left-hand side is
\[
 (u,x)+(u,\alpha)(\alpha^\ast,x)
       \bigl((\alpha^\ast,\alpha)-2\bigr)=(u,x).
\]
\end{proof}

\begin{remark}\label{sga3:XXI:1.1.10}
If one assumes \(0\notin R\) and \(0\notin R^\ast\), then
Lemma~\ref{sga3:XXI:1.1.9} is equivalent to
\((\mathrm{DR\,I})\).
\end{remark}

\begin{corollary}\label{sga3:XXI:1.1.11}
Let \(h\mapsto h^\vee\) denote the isomorphism from
\(\GL(M)\) onto \(\GL(M^\ast)\) that associates to \(h\) its
contragredient. Formula \eqref{sga3:XXI:eq:4} may then also be written
\begin{equation}
 (s_\alpha)^\vee=s_\alpha^\ast.
 \tag{5}\label{sga3:XXI:eq:5}
\end{equation}
\end{corollary}

\begin{corollary}\label{sga3:XXI:1.1.12}
The preceding isomorphism induces an isomorphism from \(W(\cR)\) onto
\(W^\ast(\cR)\).
\end{corollary}

\begin{scholium}\label{sga3:XXI:1.1.13}
By the preceding result, we shall identify \(W\) and \(W^\ast\), and
regard \(W\) as a group of transformations of
\[
 R,\ R^\ast,\ M,\ M^\ast,\ \Gamma_0(R),\ \Gamma_0(R^\ast),\
 V,\ V^\ast,\ \VR(R),\ \VR(R^\ast).
\]
We shall write \(s_\alpha\) for \(s_\alpha^\ast\).
\end{scholium}

% END INCLUDED SGA3_English_current_buildable_source_20260823/source/inputs/XXI/tex/components/00_expose_XXI_opening_through_scholium1113_en.tex

% BEGIN INCLUDED SGA3_English_current_buildable_source_20260823/source/inputs/XXI/tex/components/01_expose_XXI_section12_through_application212_en.tex
\subsection{The map \texorpdfstring{\(p\)}{p}}

\begin{lemma}\label{sga3:XXI:1.2.1}
Let \(p:M\to M^\ast\) (respectively \(V\to V^\ast\)) be the linear map
defined by
\begin{equation}
 p(x)=\sum_{u\in R^\ast}(u,x)u.
 \tag{6}\label{sga3:XXI:eq:6}
\end{equation}
Write \(\ell(x)=(p(x),x)\). Then:
\begin{align}
 \ell(x)&\geq0,\qquad \ell(\alpha)>0
     &&\text{for }\alpha\in R,
 \tag{7}\label{sga3:XXI:eq:7}\\
 \ell(wx)&=\ell(x)
     &&\text{for }w\in W,
 \tag{8}\label{sga3:XXI:eq:8}\\
 (p(x),y)&=(p(y),x)
     &&\text{for }x,y\in V,
 \tag{9}\label{sga3:XXI:eq:9}\\
 \ell(\alpha)\alpha^\ast&=2p(\alpha)
     &&\text{for }\alpha\in R.
 \tag{10}\label{sga3:XXI:eq:10}
\end{align}
\end{lemma}

\begin{proof}
The first three relations are immediate.\footnote{Editors' note:
Formula \eqref{sga3:XXI:eq:8} follows from
\eqref{sga3:XXI:eq:4}, \eqref{sga3:XXI:eq:2}, and
\((\mathrm{DR\,II})\).}
We prove the last one. By \eqref{sga3:XXI:eq:1}, for \(u\in V^\ast\),
\begin{align*}
 (u,\alpha)^2\alpha^\ast
  &=(u,\alpha)u-(u,\alpha)s_\alpha(u)\\
  &=(u,\alpha)u+(u,s_\alpha(\alpha))s_\alpha(u)\\
  &=(u,\alpha)u+(s_\alpha(u),\alpha)s_\alpha(u)
       \qquad\bigl(s_\alpha^2=\operatorname{id}\bigr).
\end{align*}
Since \(s_\alpha\) permutes \(R^\ast\) by
\((\mathrm{DR\,II})\), summing over \(u\in R^\ast\) gives the result.
\end{proof}

\begin{scholium}\label{sga3:XXI:1.2.2}
Relation \eqref{sga3:XXI:eq:10} says that the map
\(\alpha\mapsto\ell(\alpha)\alpha^\ast\) is the restriction to \(R\)
of a linear map from \(M\) to \(M^\ast\). In particular,
\[
 (-\alpha)^\ast=-\alpha^\ast.
\]
\end{scholium}

\begin{corollary}\label{sga3:XXI:1.2.3}
The map \(p\) induces an isomorphism from \(\VR(R)\) onto
\(\VR(R^\ast)\).
\end{corollary}

\begin{proof}
The map \(p\) sends \(\VR(R)\) onto \(\VR(R^\ast)\). Hence
\[
 \dim\VR(R)\geq\dim\VR(R^\ast).
\]
Applying this inequality to the dual root datum gives
\[
 \dim\VR(R)=\dim\VR(R^\ast),
\]
so the surjective map \(p\) is also bijective.
\end{proof}

\begin{corollary}\label{sga3:XXI:1.2.4}
One has \(\dim\VR(R)=\dim\VR(R^\ast)\), and therefore
\[
 \rgss(\cR)=\rgss(\cR^\ast).
\]
\end{corollary}

\begin{corollary}\label{sga3:XXI:1.2.5}
The bilinear form \((u,x)\) is nondegenerate on
\(\VR(R^\ast)\times\VR(R)\), and thus puts these
\(\mathbb Q\)-vector spaces in duality.
\end{corollary}

\begin{proof}
If \((u,x)=0\) for every \(u\in R^\ast\), then \(p(x)=0\).
\end{proof}

\begin{corollary}\label{sga3:XXI:1.2.6}
The symmetric bilinear form \((p(x),y)\) is positive and nondegenerate
on \(\VR(R)\).
\end{corollary}

\begin{corollary}\label{sga3:XXI:1.2.7}
The group \(W\) acts faithfully on \(R\) (and hence on the other sets
listed in Scholium~\ref{sga3:XXI:1.1.13}).
\end{corollary}

\begin{proof}
Let \(u\in V^\ast\), and suppose \(w\in W\) satisfies
\(w(\alpha)=\alpha\) for every \(\alpha\in R\). We prove \(w(u)=u\).
For every \(\alpha\in R\),
\[
 (w(u)-u,\alpha)
 =(w(u),\alpha)-(u,\alpha)
 =(u,w^{-1}(\alpha))-(u,\alpha)=0.
\]
But \(w(u)-u\in\VR(R^\ast)\). Since it is orthogonal to every root, it
vanishes by Corollary~\ref{sga3:XXI:1.2.5}.
\end{proof}

\begin{corollary}\label{sga3:XXI:1.2.8}
The group \(W\) is finite.
\end{corollary}

\begin{proposition}\label{sga3:XXI:1.2.9}
The operations of \(W\) respect the correspondence between roots and
coroots. In other words, for \(\alpha\in R\) and \(w\in W\),
\[
 w(\alpha^\ast)=w(\alpha)^\ast.
\]
\end{proposition}

\begin{proof}
It suffices to verify this for \(w=s_\beta\), \(\beta\in R\), that is,
to verify
\[
 s_\beta(\alpha^\ast)=s_\beta(\alpha)^\ast.
\]
Now \(\ell(s_\beta(\alpha))s_\beta(\alpha)^\ast/2\) equals
\begin{align*}
 p(s_\beta(\alpha))
 &=\sum_{u\in R^\ast}(u,s_\beta(\alpha))u\\
 &=\sum_{u\in R^\ast}(s_\beta(u),\alpha)u\\
 &=\sum_{u\in R^\ast}(u,\alpha)s_\beta(u)
  =s_\beta(p(\alpha)).
\end{align*}
Since \(\ell(s_\beta(\alpha))=\ell(\alpha)\), we obtain
\[
 s_\beta(\alpha)^\ast
 =s_\beta\!\left(\frac{2p(\alpha)}{\ell(\alpha)}\right)
 =s_\beta(\alpha^\ast).
\]
\end{proof}

\begin{corollary}\label{sga3:XXI:1.2.10}
If \(\alpha\in R\) and \(w\in W\), then
\[
 ws_\alpha w^{-1}=s_{w(\alpha)}.
\]
\end{corollary}

\begin{proof}
For \(x\in V\),
\[
 ws_\alpha w^{-1}(x)
 =x-(\alpha^\ast,w^{-1}(x))w(\alpha)
 =x-(w(\alpha^\ast),x)w(\alpha).
\]
By Proposition~\ref{sga3:XXI:1.2.9}, the last expression is
\[
 \begin{aligned}
 x-(w(\alpha)^\ast,x)w(\alpha)
   &=s_{w(\alpha)}(x).
 \end{aligned}
\]
\end{proof}

\begin{corollary}\label{sga3:XXI:1.2.11}\leavevmode\par
Let \(R'\subset R\). The quadruple
\((M,M^\ast,R',R'^\ast)\) is a root datum if and only if
\(\alpha,\beta\in R'\) implies \(s_\alpha(\beta)\in R'\).
\end{corollary}

\section{Relations between two roots}

\subsection{Proportional roots}

\begin{proposition}\label{sga3:XXI:2.1.1}
Let \(\alpha\) and \(\beta\) be two roots. The following conditions are
equivalent:
\begin{enumerate}
\item[(i)] there exists \(k\in\mathbb Q\) such that
  \(\alpha=k\beta\);
\item[(ii)] \(s_\alpha=s_\beta\).
\end{enumerate}
Under these conditions,
\(\alpha^\ast=k^{-1}\beta^\ast\), and \(k\) is one of
\[
 1,\ -1,\ 2,\ -2,\ \tfrac12,\ -\tfrac12.
\]
\end{proposition}

\begin{proof}
Assume first \SGARefLink{sga3:XXI:2.1.1:item:i:02}{(i)}. Then
\[
 \alpha^\ast
 =\ell(\alpha)^{-1}2p(\alpha)
 =k^{-2}\ell(\beta)^{-1}2k\,p(\beta)
 =k^{-1}\beta^\ast.
\]
This immediately gives \(s_\alpha=s_\beta\). Conversely, if
\(s_\alpha=s_\beta\), then
\[
 \alpha=s_\alpha(-\alpha)
       =s_\beta(-\alpha)
       =(\beta^\ast,\alpha)\beta-\alpha,
\]
so
\[
 2\alpha=(\beta^\ast,\alpha)\beta,
\qquad(\beta^\ast,\alpha)\in\mathbb Z;
\]
thus \SGARefLink{sga3:XXI:2.1.1:item:ii:02}{(ii)} implies \SGARefLink{sga3:XXI:2.1.1:item:i:02}{(i)}. Finally, if \(\alpha=k\beta\), then
\(\alpha^\ast=k^{-1}\beta^\ast\), whence
\[
 (\alpha^\ast,\beta)=2k^{-1},
 \qquad(\beta^\ast,\alpha)=2k.
\]
Thus \(2k\) and \(2k^{-1}\) are integers, which gives the stated list.
\end{proof}

\begin{application}\label{sga3:XXI:2.1.2}
Root data \(\cR\) with \(\rgss(\cR)=1\) are of one of the following two
types:
\begin{enumerate}
\item[\SGAInlineRefTarget{sga3:XXI:2.1.1:item:i:02}(i)] Type \(A_1\): there exists \(\alpha\in M\) such that the
  roots are \(\alpha\) and \(-\alpha\). The coroots are then
  \(\alpha^\ast\) and \(-\alpha^\ast\).
\item[\SGAInlineRefTarget{sga3:XXI:2.1.1:item:ii:02}(ii)] Type \(A'_1\): there exists \(\alpha\in M\) such that the
  roots are
  \(\alpha,-\alpha,2\alpha,-2\alpha\). The coroots are then
  \(\alpha^\ast,-\alpha^\ast,\alpha^\ast/2,-\alpha^\ast/2\).
\end{enumerate}
\end{application}

% END INCLUDED SGA3_English_current_buildable_source_20260823/source/inputs/XXI/tex/components/01_expose_XXI_section12_through_application212_en.tex

% BEGIN INCLUDED SGA3_English_current_buildable_source_20260823/source/inputs/XXI/tex/components/02_expose_XXI_definition213_through_definition223_en.tex
\begin{definition}\label{sga3:XXI:2.1.3}
A root \(\alpha\in R\) is called \emph{indivisible} if
\(\alpha/2\notin R\). The root datum \(\cR\) is called
\emph{reduced} if every root is indivisible.
\end{definition}

The root datum \(\cR\) is reduced if and only if \(\cR^\ast\) is
reduced. If \(\alpha\) is indivisible and \(w\in W\), then
\(w(\alpha)\) is indivisible.

\begin{definition}\label{sga3:XXI:2.1.4}
Let \(\alpha\in R\). If \(\alpha\) is indivisible, set
\(\operatorname{ind}(\alpha)=\alpha\); otherwise set
\(\operatorname{ind}(\alpha)=\alpha/2\).
\end{definition}

\begin{corollary}\label{sga3:XXI:2.1.5}
If \(\alpha\in R\) is indivisible and \(k\alpha\in R\), where
\(k\in\mathbb Q\), then \(k\in\mathbb Z\).
\end{corollary}

\begin{proposition}\label{sga3:XXI:2.1.6}
Let \(\cR\) be a root datum. Then
\[
 \operatorname{ind}(\cR)
 =\bigl(M,M^\ast,\operatorname{ind}(R),
         \operatorname{ind}(R)^\ast\bigr)
\]
is a reduced root datum, and
\[
 W\bigl(\operatorname{ind}(\cR)\bigr)=W(\cR).
\]
\end{proposition}

\begin{proof}
The quadruple \(\operatorname{ind}(\cR)\) is a root datum by
Corollary~\ref{sga3:XXI:1.2.11}, since the Weyl group permutes the
indivisible roots. The second assertion follows from
Proposition~\ref{sga3:XXI:2.1.1}.
\end{proof}

\begin{remark}\label{sga3:XXI:2.1.7}
If \(\cR\) is not reduced, then
\[
 \operatorname{ind}(R)^\ast\neq\operatorname{ind}(R^\ast),
 \qquad
 \operatorname{ind}(\cR^\ast)\neq\operatorname{ind}(\cR)^\ast.
\]
\end{remark}

\subsection{Orthogonal roots}

\begin{lemma}\label{sga3:XXI:2.2.1}
Let \(\alpha\) and \(\beta\) be two roots. Then
\begin{equation}
 \ell(\alpha)(\alpha^\ast,\beta)
 =\ell(\beta)(\beta^\ast,\alpha).
 \tag{11}\label{sga3:XXI:eq:11}
\end{equation}
\end{lemma}

\begin{proof}
This follows immediately from Lemma~\ref{sga3:XXI:1.2.1},
formulas \eqref{sga3:XXI:eq:9} and \eqref{sga3:XXI:eq:10}.
\end{proof}

\begin{corollary}\label{sga3:XXI:2.2.2}
Let \(\alpha,\beta\in R\). The following conditions are equivalent:
\begin{enumerate}
\item[\SGAInlineRefTarget{sga3:XXI:2.2.2:item:i}(i)] \((\alpha^\ast,\beta)=0\);
\item[(i bis)] \((\beta^\ast,\alpha)=0\);
\item[(ii)] \((p(\alpha),\beta)=0\);
\item[(iii)] \(s_\alpha(\beta)=\beta\);
\item[(iii bis)] \(s_\beta(\alpha)=\alpha\);
\item[(iv)] \(s_\alpha(\beta^\ast)=\beta^\ast\);
\item[(iv bis)] \(s_\beta(\alpha^\ast)=\alpha^\ast\);
\item[\SGAInlineRefTarget{sga3:XXI:2.2.2:item:v}(v)] \(s_\alpha\neq s_\beta\), and \(s_\alpha\) and \(s_\beta\)
  commute.
\end{enumerate}
\end{corollary}

\begin{proof}
All the equivalences are immediate except those involving \SGARefLink{sga3:XXI:2.2.2:item:v}{(v)}. We
first show that \SGARefLink{sga3:XXI:2.2.2:item:i}{(i)}, together with its equivalent condition (i bis),
implies \SGARefLink{sga3:XXI:2.2.2:item:v}{(v)}. One has
\[
 s_\alpha s_\beta(x)
 =x-(\beta^\ast,x)\beta-(\alpha^\ast,x)\alpha
   +(\beta^\ast,x)(\alpha^\ast,\beta)\alpha.
\]
If \((\alpha^\ast,\beta)=0\), then
\((\beta^\ast,\alpha)=0\), and consequently
\[
 s_\alpha s_\beta(x)=s_\beta s_\alpha(x).
\]

Conversely, suppose \SGARefLink{sga3:XXI:2.2.2:item:v}{(v)}. Then
\[
 s_\alpha=s_\beta s_\alpha s_\beta
          =s_{s_\beta(\alpha)}
 \qquad\text{by Corollary~\ref{sga3:XXI:1.2.10}.}
\]
By Proposition~\ref{sga3:XXI:2.1.1}, there is a
\(k\in\mathbb Q^\ast\) such that
\[
 \alpha=k\,s_\beta(\alpha)
       =k\alpha-k(\beta^\ast,\alpha)\beta.
\]
Since \(s_\alpha\neq s_\beta\), the roots \(\alpha\) and \(\beta\)
are not proportional, again by Proposition~\ref{sga3:XXI:2.1.1}.
It follows that \((\beta^\ast,\alpha)=0\).
\end{proof}

\begin{definition}\label{sga3:XXI:2.2.3}
Two roots satisfying the equivalent conditions of
Corollary~\ref{sga3:XXI:2.2.2} are called \emph{orthogonal}.
\end{definition}

% END INCLUDED SGA3_English_current_buildable_source_20260823/source/inputs/XXI/tex/components/02_expose_XXI_definition213_through_definition223_en.tex

% BEGIN INCLUDED SGA3_English_current_buildable_source_20260823/source/inputs/XXI/tex/components/03_expose_XXI_remark224_through_lemma311_en.tex
\begin{remark}\label{sga3:XXI:2.2.4}
The roots \(\alpha\) and \(\beta\) are orthogonal if and only if the
coroots \(\alpha^\ast\) and \(\beta^\ast\) are orthogonal.
\end{remark}

\begin{lemma}\label{sga3:XXI:2.2.5}
If \(\alpha\) and \(\beta\) are two orthogonal roots, then
\(\alpha+\beta\in R\) if and only if \(\alpha-\beta\in R\).
\end{lemma}

\begin{proof}
Indeed,
\[
 s_\beta(\alpha+\beta)
 =s_\beta(\alpha)+s_\beta(\beta)=\alpha-\beta.
\]
\end{proof}

\begin{lemma}\label{sga3:XXI:2.2.6}
Let \(\alpha\) and \(\beta\) be two nonorthogonal roots. For a coroot
\(\gamma^\ast\), define \(\ell(\gamma^\ast)\) to be the value of
\(\ell\) on the root \(\gamma^\ast\) of \(\cR^\ast\). Then
\begin{equation}
 \ell(\alpha)\ell(\alpha^\ast)
 =\ell(\beta)\ell(\beta^\ast).
 \tag{12}\label{sga3:XXI:eq:12}
\end{equation}
\end{lemma}

\begin{proof}
Multiply formula \eqref{sga3:XXI:eq:11} by the corresponding formula
for \(\cR^\ast\), and use
\((\gamma^\ast)^\ast=\gamma\) for every \(\gamma\in R\). This gives
\[
 (\beta^\ast,\alpha)(\alpha^\ast,\beta)
 \ell(\alpha)\ell(\alpha^\ast)
 =
 (\beta^\ast,\alpha)(\alpha^\ast,\beta)
 \ell(\beta)\ell(\beta^\ast).
\]
The common integer factor is nonzero because the roots are not
orthogonal, so it may be cancelled.
\end{proof}

\subsection{General case}

\begin{proposition}\label{sga3:XXI:2.3.1}
If \(\alpha\) and \(\beta\) are any two roots, then
\begin{equation}
 0\leq(\alpha^\ast,\beta)(\beta^\ast,\alpha)\leq4.
 \tag{13}\label{sga3:XXI:eq:13}
\end{equation}
If \(\alpha\) and \(\beta\) are neither proportional nor orthogonal,
then
\[
 1\leq(\alpha^\ast,\beta)(\beta^\ast,\alpha)\leq3.
\]
\end{proposition}

\begin{proof}
One has
\[
 4(p(\alpha),\beta)^2
 =\ell(\alpha)\ell(\beta)
   (\alpha^\ast,\beta)(\beta^\ast,\alpha).
\]
On the other hand, by Corollary~\ref{sga3:XXI:1.2.6}, the symmetric
bilinear form \((p(x),y)\) is positive and nondegenerate on
\(\VR(R)\). Hence
\[
 (p(x),y)^2\leq\ell(x)\ell(y).
\]
\footnote{Editors' note: equality holds if and only if \(x\) and
\(y\) are proportional.}
The first inequality follows. In the nonproportional, nonorthogonal
case, the product is a nonzero integer and cannot equal \(4\), which
gives the second inequality.
\end{proof}

\begin{corollary}\label{sga3:XXI:2.3.2}
Let \(\alpha\) and \(\beta\) be two nonorthogonal roots. If
\(\ell(\alpha)=\ell(\beta)\), then there is a \(w\in W\) such that
\(w(\alpha)=\beta\).
\end{corollary}

\begin{proof}
If \(\alpha\) and \(\beta\) are proportional, equality of their
\(\ell\)-values gives \(\alpha=\beta\) or \(\alpha=-\beta\); take
\(w=1\) or \(w=s_\alpha\), respectively. If they are neither
proportional nor orthogonal, formula \eqref{sga3:XXI:eq:11} and
Proposition~\ref{sga3:XXI:2.3.1} give
\[
 (\beta^\ast,\alpha)=(\alpha^\ast,\beta)=\pm1.
\]
If both integers are \(1\), take \(w=s_\beta s_\alpha s_\beta\). If
both are \(-1\), take \(w=s_\alpha s_\beta\).
\end{proof}

\begin{corollary}\label{sga3:XXI:2.3.3}
Let \(\alpha\) and \(\beta\) be roots. If
\(\alpha\neq\beta\) and \((\beta^\ast,\alpha)>0\), then
\(\alpha-\beta\) is a root. Respectively, if
\(\alpha\neq-\beta\) and \((\beta^\ast,\alpha)<0\), then
\(\alpha+\beta\) is a root.
\end{corollary}

\begin{proof}
The second case follows from the first by replacing \(\beta\) by
\(-\beta\). Suppose first that \(\alpha\) and \(\beta\) are
proportional and that \((\beta^\ast,\alpha)>0\). Then
\[
 \beta=\alpha,\qquad 2\beta=\alpha,\qquad\text{or}\qquad
 \beta=2\alpha.
\]
The first case is excluded. In the other two cases,
\(\alpha-\beta=\beta\in R\) and
\(\alpha-\beta=-\alpha\in R\), respectively.

If \(\alpha\) and \(\beta\) are not proportional, then
\((\alpha^\ast,\beta)\) and \((\beta^\ast,\alpha)\) are strictly
positive integers whose product is at most \(3\). Thus one of them
equals \(1\). If \((\beta^\ast,\alpha)=1\), then
\[
 \alpha-\beta=s_\beta(\alpha)\in R;
\]
if \((\alpha^\ast,\beta)=1\), then
\[
 \alpha-\beta=-s_\alpha(\beta)\in R.
\]
\end{proof}

\begin{lemma}\label{sga3:XXI:2.3.4}
Let \(\alpha\) and \(\beta\) be two nonproportional roots. If
\(\beta-\alpha\) is not a root, then
\(\beta+k\alpha\in R\) for \(k=-(\alpha^\ast,\beta)\), but not for
\(k=-(\alpha^\ast,\beta)+1\).
\end{lemma}

\begin{proof}
Indeed,
\[
 \beta-(\alpha^\ast,\beta)\alpha=s_\alpha(\beta)\in R,
\]
whereas
\[
 \beta+\bigl(-(\alpha^\ast,\beta)+1\bigr)\alpha
 =s_\alpha(\beta-\alpha)\notin R.
\]
\end{proof}

\begin{proposition}\label{sga3:XXI:2.3.5}
Let \(\alpha\) and \(\beta\) be two nonproportional roots. The set of
integers \(k\) for which \(\beta+k\alpha\in R\) is an interval
\([p,q]\), with \(p\leq0\) and \(q\geq0\), and
\[
 p+q=-(\alpha^\ast,\beta).
\]
\end{proposition}

\begin{proof}
For the first assertion it suffices, for example, to show that if
\(\beta+k\alpha\in R\) and \(k>0\) is an integer, then
\(\beta+(k-1)\alpha\in R\). This is trivial if \(k=1\). If
\(k\geq2\), then
\[
 (\alpha^\ast,\beta+k\alpha)
 =(\alpha^\ast,\beta)+2k\geq-3+4>0,
\]
and the assertion follows from Corollary~\ref{sga3:XXI:2.3.3}.

Let \([p,q]\) be the interval in question. Applying
Lemma~\ref{sga3:XXI:2.3.4} to \(\beta+p\alpha\), one obtains
\[
 q-p
 =-(\alpha^\ast,\beta+p\alpha)
 =-(\alpha^\ast,\beta)-2p.
\]
\footnote{Editors' note: in this equality, \(+2p\) was corrected to
\(-2p\).}
The asserted formula follows.
\end{proof}

\begin{remark}\label{sga3:XXI:2.3.6}
The preceding formula contains the qualitative statements of
Lemma~\ref{sga3:XXI:2.2.5} and
Corollary~\ref{sga3:XXI:2.3.3}.
\end{remark}

\begin{complements}\label{sga3:XXI:2.3.7}
\footnote{Editors' note: these complements, useful for the proof of
Proposition~\ref{sga3:XXI:3.1.5}, were added by the editors.}
By formulas \eqref{sga3:XXI:eq:9} and
\eqref{sga3:XXI:eq:10}, and by
Corollary~\ref{sga3:XXI:1.2.6}, the bilinear form on \(\VR(R)\)
defined by
\[
 \langle x,y\rangle=(p(x),y)
\]
is symmetric and positive definite. For every \(\alpha\in R\) and
\(y\in\VR(R)\),
\[
 \langle\alpha,y\rangle
 =\frac{\ell(\alpha)}2(\alpha^\ast,y),
 \qquad
 \ell(\alpha)=\langle\alpha,\alpha\rangle.
\]
Consequently, Corollary~\ref{sga3:XXI:2.3.3} gives the following
corollary.

\medskip
\noindent\textbf{Corollary.}\quad
\SGARefTarget{sga3:XXI:2.3.7:corollary}
Let \(\alpha\neq\beta\) in \(R\). If \(\alpha-\beta\notin R\), then
\[
 \langle\alpha,\beta\rangle\leq0.
\]
\end{complements}

\section{Simple roots, positive roots}

\subsection{Systems of simple roots}

\begin{lemma}\label{sga3:XXI:3.1.1}
Let \(\alpha\) and \(\alpha_i\), \(i=1,\ldots,n\), be roots, and
suppose that \(\alpha\) is distinct from the \(\alpha_i\). If
\[
 \alpha=\sum_{i=1}^n q_i\alpha_i,
 \qquad q_i\in\Qp,
\]
then there is an index \(i\) such that
\[
 q_i\neq0,\qquad(\alpha^\ast,\alpha_i)>0,\qquad
 \alpha-\alpha_i\in R.
\]
\end{lemma}

\begin{proof}
Write
\[
 2=(\alpha^\ast,\alpha)
   =\sum_{i=1}^n q_i(\alpha^\ast,\alpha_i).
\]
This proves the first two assertions. The third then follows from
Corollary~\ref{sga3:XXI:2.3.3}.
\end{proof}

% END INCLUDED SGA3_English_current_buildable_source_20260823/source/inputs/XXI/tex/components/03_expose_XXI_remark224_through_lemma311_en.tex

% BEGIN INCLUDED SGA3_English_current_buildable_source_20260823/source/inputs/XXI/tex/components/04_expose_XXI_proposition312_through_remark319_en.tex
\begin{proposition}\label{sga3:XXI:3.1.2}
Let \(\alpha\) and \(\alpha_i\), \(i=1,\ldots,n\), be roots. If
\[
 \alpha=\sum_{i=1}^n m_i\alpha_i,
 \qquad m_i\in\mathbb N,
\]
then there is a sequence \(\beta_1,\ldots,\beta_m\) of roots chosen
from among the \(\alpha_i\), not necessarily pairwise distinct, such
that, on setting
\[
 \gamma_p=\sum_{i=1}^p\beta_i,\qquad p=1,\ldots,m,
\]
one has \(\gamma_p\in R\) and \(\gamma_m=\alpha\).
\end{proposition}

\begin{proof}
We argue by induction on the integer
\(m'=\sum_i m_i\). If \(\alpha\) equals one of the \(\alpha_i\), say
\(\alpha_{i_0}\), which necessarily happens when \(m'=1\), take
\[
 m=1,\qquad\beta_1=\alpha_{i_0}.
\]
Otherwise, Lemma~\ref{sga3:XXI:3.1.1} gives an index \(i\) for which
\(m_i\neq0\) and \(\alpha-\alpha_i\) is a root. Then
\(m_i-1\in\mathbb N\) and
\[
 \alpha-\alpha_i
 =(m_i-1)\alpha_i+\sum_{j\neq i}m_j\alpha_j.
\]
Apply the induction hypothesis to \(\alpha-\alpha_i\), and append
\(\alpha_i\) to the resulting sequence.
\end{proof}

\begin{corollary}\label{sga3:XXI:3.1.3}
Let \(R'\subset R\). The following conditions are equivalent:
\begin{enumerate}
\item[\SGAInlineRefTarget{sga3:XXI:3.1.3:item:i}(i)] if \(\alpha,\beta\in R'\) and \(\alpha+\beta\in R\), then
  \(\alpha+\beta\in R'\);
\item[\SGAInlineRefTarget{sga3:XXI:3.1.3:item:ii}(ii)] \((\mathbb N\cdot R')\cap R\subset R'\).
\end{enumerate}
\end{corollary}

\begin{proof}
Clearly \SGARefLink{sga3:XXI:3.1.3:item:ii}{(ii)} implies \SGARefLink{sga3:XXI:3.1.3:item:i}{(i)}. The converse follows immediately from
Proposition~\ref{sga3:XXI:3.1.2}.
\end{proof}

\begin{definition}\label{sga3:XXI:3.1.4}
A set of roots satisfying the equivalent conditions of
Corollary~\ref{sga3:XXI:3.1.3} is called \emph{closed}.
\end{definition}

\begin{proposition}\label{sga3:XXI:3.1.5}
Let \(\Delta\subset R\) be a set of roots. The following conditions
are equivalent:
\begin{enumerate}
\item[\SGAInlineRefTarget{sga3:XXI:3.1.5:item:i}(i)] the elements of \(\Delta\) are indivisible and linearly
  independent, and
  \[
   R\subset(\Qp\cdot\Delta)\cup(-\Qp\cdot\Delta);
  \]
\item[\SGAInlineRefTarget{sga3:XXI:3.1.5:item:ii}(ii)] the elements of \(\Delta\) are linearly independent, and
  \[
   R\subset(\mathbb N\cdot\Delta)\cup(-\mathbb N\cdot\Delta);
  \]
\item[\SGAInlineRefTarget{sga3:XXI:2.3.7:corollary:item:iii}(iii)] every root can be written uniquely as a linear
  combination of the elements of \(\Delta\), with integer
  coefficients all of the same sign.
\end{enumerate}
\end{proposition}

\begin{proof}
\footnote{Editors' note: in the argument that follows, the order was
changed and the proof of the implication \SGARefLink{sga3:XXI:2.3.7:corollary:item:iii}{(iii)}\(\Rightarrow\)\SGARefLink{sga3:XXI:3.1.5:item:i}{(i)} was
spelled out.}
It is clear that \SGARefLink{sga3:XXI:3.1.5:item:ii}{(ii)} implies \SGARefLink{sga3:XXI:2.3.7:corollary:item:iii}{(iii)}. Write
\(\alpha_1,\ldots,\alpha_n\) for the distinct elements of \(\Delta\).

Assume \SGARefLink{sga3:XXI:3.1.5:item:i}{(i)}, and let \(\alpha\in R\). There is a unique expression
\[
 \varepsilon\alpha=\sum_i q_i\alpha_i,
 \qquad q_i\in\Qp,\qquad\varepsilon=\pm1.
\]
We show that the \(q_i\) are integers. This is immediate when all but
one of them vanish, by Corollary~\ref{sga3:XXI:2.1.5}. Otherwise,
\(\varepsilon\alpha\) is distinct from the \(\alpha_i\), and
Lemma~\ref{sga3:XXI:3.1.1} gives an \(i_0\) such that
\[
 \alpha'=\varepsilon\alpha-\alpha_{i_0}\in R,
 \qquad q_{i_0}\neq0.
\]
Thus
\[
 \alpha'
 =(q_{i_0}-1)\alpha_{i_0}
  +\sum_{i\neq i_0}q_i\alpha_i.
\]
At least one \(q_i\), with \(i\neq i_0\), is nonzero, so \SGARefLink{sga3:XXI:3.1.5:item:i}{(i)} forces
\(q_{i_0}-1\geq0\). Repeating this operation with \(\alpha'\), after
finitely many steps we find that all \(q_i\) are integers. This proves
\SGARefLink{sga3:XXI:3.1.5:item:ii}{(ii)}.

Assume now \SGARefLink{sga3:XXI:2.3.7:corollary:item:iii}{(iii)}. We first prove that
\(\alpha_1,\ldots,\alpha_n\) are linearly independent over
\(\mathbb Q\). Otherwise there would be an equality
\[
 x=\sum_{i\in I}a_i\alpha_i
  =\sum_{j\in J}b_j\alpha_j,
\]
where \(I\) and \(J\) are disjoint subsets of
\(\{1,\ldots,n\}\), at least one of them, say \(I\), is nonempty, and
\(a_i,b_j\in\mathbb N^\ast\). By the corollary in
Complements~\ref{sga3:XXI:2.3.7}, one has
\(\langle\alpha_i,\alpha_j\rangle\leq0\) for \(i\in I\) and \(j\in J\).
Therefore \(\langle x,x\rangle\leq0\), and hence \(x=0\).

Choose \(i_0\in I\). The two expressions
\[
 \alpha_{i_0}
 =\alpha_{i_0}+\sum_{i\in I}a_i\alpha_i
 =\alpha_{i_0}+\sum_{j\in J}b_j\alpha_j
\]
then imply, by \SGARefLink{sga3:XXI:2.3.7:corollary:item:iii}{(iii)}, that \(a_i=0=b_j\) for all \(i,j\), a
contradiction. Thus the elements of \(\Delta\) are linearly
independent over \(\mathbb Q\).

It remains to show that they are indivisible. Suppose that
\(\beta\in\Delta\) is divisible. Then \(\beta/2\in R\), so
\[
 \frac{\beta}{2}
 =\varepsilon\sum_{\alpha\in\Delta}m_\alpha\alpha,
 \qquad m_\alpha\in\mathbb N,\qquad\varepsilon=\pm1.
\]
Consequently,
\[
 \beta=2\varepsilon\sum_{\alpha\in\Delta}m_\alpha\alpha.
\]
Uniqueness gives \(m_\alpha=0\) for \(\alpha\neq\beta\) and
\(2\varepsilon m_\beta=1\), a contradiction.
\end{proof}

\begin{definition}\label{sga3:XXI:3.1.6}
A set \(\Delta\) of roots satisfying the conditions of
Proposition~\ref{sga3:XXI:3.1.5} is called a
\emph{system of simple roots}, or a \emph{base} of \(\cR\).
\end{definition}

If \(w\in W\) and \(\Delta\) is a system of simple roots, then
\(w(\Delta)\) is a system of simple roots.

\begin{remark}\label{sga3:XXI:3.1.7}
This definition involves only \(R\), not \(\cR\); indeed, it involves
only \(\operatorname{ind}(R)\).
\end{remark}

\begin{remark}\label{sga3:XXI:3.1.8}
If \(\Delta\) is a system of simple roots, then \(\Delta\) is a basis
of the free abelian group \(\Gamma_0(R)\). Hence
\[
 \operatorname{Card}(\Delta)=\rgss(\cR).
\]
\end{remark}

\begin{remark}\label{sga3:XXI:3.1.9}
The conditions of Proposition~\ref{sga3:XXI:3.1.5} are therefore
equivalent to the following:
\begin{enumerate}
\item[(i\('\))] the elements of \(\Delta\) are indivisible, their
  number is at most \(\rgss(\cR)\), and
  \[
   R\subset(\Qp\cdot\Delta)\cup(-\Qp\cdot\Delta);
  \]
\item[(ii\('\))] \(\operatorname{Card}(\Delta)\leq\rgss(\cR)\), and
  \[
   R\subset(\mathbb N\cdot\Delta)\cup(-\mathbb N\cdot\Delta).
  \]
\end{enumerate}
\end{remark}

% END INCLUDED SGA3_English_current_buildable_source_20260823/source/inputs/XXI/tex/components/04_expose_XXI_proposition312_through_remark319_en.tex

% BEGIN INCLUDED SGA3_English_current_buildable_source_20260823/source/inputs/XXI/tex/components/05_expose_XXI_corollary3110_through_definition325_en.tex
\begin{corollary}\label{sga3:XXI:3.1.10}
If \(\Delta\) is a system of simple roots, then
\(\operatorname{ind}(\Delta^\ast)\) is a system of simple coroots
\emph{(i.e.}\ a system of simple roots of \(\cR^\ast\)\emph{)}.
\end{corollary}

\begin{proof}
If
\[
 \beta=\sum_{\alpha\in\Delta}a_\alpha\alpha,
 \qquad a_\alpha\in\mathbb N,
\]
then
\[
 p(\beta)=\sum_{\alpha\in\Delta}a_\alpha p(\alpha).
\]
Consequently, by \eqref{sga3:XXI:eq:10},
\[
 \beta^\ast
 =\sum_{\alpha\in\Delta}
   a_\alpha\frac{\ell(\alpha)}{\ell(\beta)}\alpha^\ast,
\]
which shows that \(\operatorname{ind}(\Delta^\ast)\) satisfies
condition (i\('\)) of Remark~\ref{sga3:XXI:3.1.9}.
\end{proof}

By Proposition~\ref{sga3:XXI:2.1.1}, if \(\alpha^\ast\) is not
indivisible, then
\[
 \alpha^\ast=2u^\ast,
 \qquad
 u^\ast\in\operatorname{ind}(\Delta^\ast)
 \quad\text{and}\quad
 u=2\alpha.
\]
It follows that:\footnote{Editors' note: The preceding sentence was
added, and in Corollary 3.1.11 the expression
\(4a_\alpha\ell(\alpha)\) was replaced by
\(2a_\alpha\ell(\alpha)\).}

\begin{corollary}\label{sga3:XXI:3.1.11}
Let \(\Delta\) be a system of simple roots, and let
\[
 \beta=\sum_{\alpha\in\Delta}a_\alpha\alpha,
 \qquad a_\alpha\in\mathbb Z,
\]
be the expression of \(\beta\) relative to \(\Delta\). Then
\(\ell(\beta)\) divides \(2a_\alpha\ell(\alpha)\); it even divides
\(a_\alpha\ell(\alpha)\) if \(\alpha^\ast\) is indivisible.
\end{corollary}

Before continuing to state the properties of systems of simple roots,
we show that such systems exist.

\subsection{Systems of positive roots}

\begin{definition}\label{sga3:XXI:3.2.1}
A subset \(R_+\subset R\) is called a \emph{system of positive roots}
of \(\cR\) \emph{(or of \(R\); see
Remark~\ref{sga3:XXI:3.2.2}\emph{)}} if it satisfies:
\begin{enumerate}
\item[(P 1)] \(R_+\cap(-R_+)=\varnothing\);
\item[(P 2)] \(R_+\cup(-R_+)=R\);
\item[(P 3)] \((\Qp\cdot R_+)\cap R\subset R_+\).
\end{enumerate}
In particular, such a set is closed. We shall see later, in
Corollary~\ref{sga3:XXI:3.3.8}, that a closed set satisfying
(P 1) and (P 2) also satisfies (P 3), and hence is a system of
positive roots. If \(w\in W\) and \(R_+\) is a system of positive
roots, then \(w(R_+)\) is a system of positive roots.
\end{definition}

\begin{remark}\label{sga3:XXI:3.2.2}
This definition involves only \(R\). We shall also say that \(R_+\)
is a system of positive roots of \(R\).
\end{remark}

\begin{remark}\label{sga3:XXI:3.2.3}
Conditions (P 1) and (P 2) immediately give
\[
 \operatorname{Card}(R_+)=\frac{\operatorname{Card}(R)}{2}.
\]
It follows that if \(R_+\) and \(R'_+\) are two systems of positive
roots and \(R_+\subset R'_+\), then \(R_+=R'_+\).
\end{remark}

\begin{remark}\label{sga3:XXI:3.2.4}
If \(R_+\) is a system of positive roots, then \(R^\ast_+\) is a
system of positive coroots \emph{(i.e.}\ a system of positive roots
of \(R^\ast\)\emph{)}.
\end{remark}

This follows immediately from Lemma~\ref{sga3:XXI:1.1.4} and
Scholium~\ref{sga3:XXI:1.2.2}.

\begin{definition}\label{sga3:XXI:3.2.5}
Let \(\Delta\) be a system of simple roots. Set
\[
 \mathscr{P}(\Delta)
 =(\Qp\cdot\Delta)\cap R
 =(\mathbb N_+\cdot\Delta)\cap R.
\]
\end{definition}

% END INCLUDED SGA3_English_current_buildable_source_20260823/source/inputs/XXI/tex/components/05_expose_XXI_corollary3110_through_definition325_en.tex

% BEGIN INCLUDED SGA3_English_current_buildable_source_20260823/source/inputs/XXI/tex/components/06_expose_XXI_proposition326_through_lemma3210_en.tex
\begin{proposition}\label{sga3:XXI:3.2.6}
If \(\Delta\) is a system of simple roots, then
\(\mathscr{P}(\Delta)\) is a system of positive roots. If \(\Delta\)
is a system of simple roots and \(R_+\) is a system of positive roots,
then
\[
 \Delta\subset R_+
 \quad\Longleftrightarrow\quad
 R_+=\mathscr{P}(\Delta).
\]
\end{proposition}

\begin{proof}
The first assertion is immediate. If \(\Delta\subset R_+\), then
\(\mathscr{P}(\Delta)\subset R_+\) by (P 3), and therefore
\(\mathscr{P}(\Delta)=R_+\) by
Remark~\ref{sga3:XXI:3.2.3}. The rest is trivial.
\end{proof}

\begin{remark}\label{sga3:XXI:3.2.7}
Systems of positive roots exist: endow \(\VR(R)\) with the structure
of a totally ordered vector space. The set of roots that are
\(\geq0\) for this order is a system of positive roots.
\end{remark}

\begin{theorem}\label{sga3:XXI:3.2.8}
Let \(R_+\) be a system of positive roots. There exists a unique
system of simple roots \(\mathscr{S}(R_+)\) such that
\[
 \mathscr{S}(R_+)\subset R_+,
 \qquad\text{i.e.}\qquad
 R_+=\mathscr{P}\bigl(\mathscr{S}(R_+)\bigr).
\]
\end{theorem}

Uniqueness follows immediately from:

\begin{lemma}\label{sga3:XXI:3.2.9}
Let \(\Delta\) be a system of simple roots. An element
\(\alpha\in\mathscr{P}(\Delta)\) belongs to \(\Delta\) if and only if
\(\alpha\) is not the sum of two elements of
\(\mathscr{P}(\Delta)\).
\end{lemma}

\begin{proof}
This follows immediately from the definitions and
Proposition~\ref{sga3:XXI:3.1.2}.
\end{proof}

We now prove the existence of \(\mathscr{S}(R_+)\). Consider the set
of subsets \(T\) of \(R_+\) such that
\[
 T=\operatorname{ind}(T)
 \qquad\text{and}\qquad
 \Qp\cdot T\supset R_+.
\]
This set is nonempty, since it contains
\(\operatorname{ind}(R_+)\). Let \(\Delta\) be a minimal element of
this set under inclusion. We show that \(\Delta\) is a system of
simple roots; by Proposition~\ref{sga3:XXI:3.1.5}\SGARefLink{sga3:XXI:3.1.5:item:i}{(i)}, it is enough to
show that \(\Delta\) is a linearly independent subset of
\(M\otimes\mathbb Q\).

\begin{lemma}\label{sga3:XXI:3.2.10}
If \(\alpha,\beta\in\Delta\) and
\[
 q\alpha+q'\beta\in R,
 \qquad q,q'\in\mathbb Q,
\]
then \(qq'\geq0\).
\end{lemma}

\begin{proof}
If \(qq'<0\), then, after interchanging \(\alpha\) and \(\beta\) if
necessary, we may write
\[
 a\alpha-b\beta\in R_+,
 \qquad a,b\in\mathbb Q^\ast_+.
\]
By the defining property of \(\Delta\), there is a relation
\[
 a\alpha-b\beta
   =\sum_{\gamma\in\Delta}c(\gamma)\gamma,
 \qquad c(\gamma)\in\Qp.
\]
If \(a\leq c(\alpha)\), this becomes
\[
 -\beta
 =\frac{c(\alpha)-a}{b}\alpha
  +\sum_{\gamma\neq\alpha}\frac{c(\gamma)}{b}\gamma.
\]
Thus \(-\beta\in(\Qp\cdot\Delta)\cap R\), which is contained in
\(R_+\) by (P 3). This gives
\(\beta\in R_+\cap(-R_+)\), contradicting (P 1).

If instead \(a>c(\alpha)\), then
\[
 \bigl(a-c(\alpha)\bigr)\alpha
 =b\beta+\sum_{\gamma\neq\alpha}c(\gamma)\gamma,
\]
which shows that
\[
 \Delta\subset\Qp\cdot\bigl(\Delta\setminus\{\alpha\}\bigr),
\]
contrary to the minimality of \(\Delta\).
\end{proof}

% END INCLUDED SGA3_English_current_buildable_source_20260823/source/inputs/XXI/tex/components/06_expose_XXI_proposition326_through_lemma3210_en.tex

% BEGIN INCLUDED SGA3_English_current_buildable_source_20260823/source/inputs/XXI/tex/components/07_expose_XXI_lemma3211_through_lemma3216_en.tex
Recall from Complements~\ref{sga3:XXI:2.3.7} that the bilinear form on
\(\VR(R)\) defined by
\[
 \langle x,y\rangle=(p(x),y)
\]
is a Euclidean inner product. Moreover, for \(\alpha\in R\),
\[
 \langle\alpha,y\rangle
 =\frac{\ell(\alpha)}{2}(\alpha^\ast,y),
\]
so \(\langle\alpha,y\rangle\) and \((\alpha^\ast,y)\) have the same
sign.

\begin{lemma}\label{sga3:XXI:3.2.11}
If \(\alpha,\beta\in\Delta\), then
\[
 (\beta^\ast,\alpha)\leq0
 \qquad\text{and hence}\qquad
 \langle\beta,\alpha\rangle\leq0.
\]
\end{lemma}

\begin{proof}
Indeed,
\[
 s_\beta(\alpha)
 =\alpha-(\beta^\ast,\alpha)\beta\in R,
\]
and therefore \((\beta^\ast,\alpha)\leq0\) by
Lemma~\ref{sga3:XXI:3.2.10}.
\end{proof}

We now show that \(\Delta\) is linearly independent. Otherwise,
Lemma~\ref{sga3:XXI:3.2.11} would imply, exactly as in the proof of
Proposition~\ref{sga3:XXI:3.1.5}, the existence of a nontrivial
relation
\[
 \sum_{\alpha\in\Delta}m(\alpha)\alpha=0,
 \qquad m(\alpha)\in\mathbb N.
\]
If \(m(\alpha_0)\geq1\), then
\[
 -\alpha_0
 =\bigl(m(\alpha_0)-1\bigr)\alpha_0
  +\sum_{\alpha\neq\alpha_0}m(\alpha)\alpha.
\]
By (P 3), this gives
\(\alpha_0\in R_+\cap(-R_+)\), contradicting (P 1).
Thus \(\Delta\) is a base of \(\cR\), completing the proof of
Theorem~\ref{sga3:XXI:3.2.8}.

\begin{corollary}\label{sga3:XXI:3.2.12}
Let \(R_+\) be a system of positive roots, let \(\Delta\) be a base of
\(\cR\), and let \(w\in W\). Then
\[
 \Delta\subset R_+
 \quad\Longleftrightarrow\quad
 R_+=\mathscr{P}(\Delta)
 \quad\Longleftrightarrow\quad
 \Delta=\mathscr{S}(R_+),
\]
\[
 \mathscr{P}\bigl(\operatorname{ind}(\Delta^\ast)\bigr)
   =\mathscr{P}(\Delta)^\ast,
 \qquad
 \mathscr{S}(R^\ast_+)
   =\operatorname{ind}\bigl(\mathscr{S}(R_+)^\ast\bigr),
\]
and
\[
 \mathscr{S}\bigl(w(R_+)\bigr)=w\bigl(\mathscr{S}(R_+)\bigr),
 \qquad
 \mathscr{P}\bigl(w(\Delta)\bigr)=w\bigl(\mathscr{P}(\Delta)\bigr).
\]
\end{corollary}

\begin{definition}\label{sga3:XXI:3.2.13}
Once a system of simple roots \(\Delta\) has been chosen, the elements
of \(\mathscr{P}(\Delta)\) are called \emph{positive}. Once a system
of positive roots \(R_+\) has been chosen, the elements of
\(\mathscr{S}(R_+)\) are called \emph{simple}.
\end{definition}

\begin{corollary}\label{sga3:XXI:3.2.14}
Let \(R_+\) be a system of positive roots, and let \(\alpha\in R_+\).
The following conditions are equivalent:
\begin{enumerate}
\item[\SGAInlineRefTarget{sga3:XXI:3.2.14:item:i}(i)] \(\alpha\) is simple, that is,
  \(\alpha\in\mathscr{S}(R_+)\);
\item[\SGAInlineRefTarget{sga3:XXI:3.2.14:item:ii}(ii)] \(\alpha\) is not the sum of two elements of \(R_+\);
\item[(iii)] \(R_+\setminus\{\alpha\}\) is closed.
\end{enumerate}
\end{corollary}

\begin{proof}
The equivalence of \SGARefLink{sga3:XXI:3.2.14:item:i}{(i)} and \SGARefLink{sga3:XXI:3.2.14:item:ii}{(ii)} follows immediately from
Lemma~\ref{sga3:XXI:3.2.9}; the equivalence of \SGARefLink{sga3:XXI:3.2.14:item:ii}{(ii)} and \SGARefLink{sga3:XXI:2.3.7:corollary:item:iii}{(iii)} is
immediate from Definition~\ref{sga3:XXI:3.1.4}.
\end{proof}

\begin{definition}\label{sga3:XXI:3.2.15}
Let \(\Delta\) be a system of simple roots. The sum of the
coefficients in the expression of a root \(\alpha\) relative to
\(\Delta\) is called the \emph{order of \(\alpha\) relative to
\(\Delta\)} and is denoted \(\operatorname{ord}_\Delta(\alpha)\).
\end{definition}

One has
\[
 \alpha\in\mathscr{P}(\Delta)
 \quad\Longleftrightarrow\quad
 \operatorname{ord}_\Delta(\alpha)>0,
 \qquad
 \alpha\in\Delta
 \quad\Longleftrightarrow\quad
 \operatorname{ord}_\Delta(\alpha)=1.
\]

\begin{lemma}\label{sga3:XXI:3.2.16}
Let \(\Delta\) be a system of simple roots and
\(\alpha\in\mathscr{P}(\Delta)\). There is a sequence
\(\alpha_i\in\Delta\), \(i=1,\ldots,m\), such that
\[
 \alpha_1+\alpha_2+\cdots+\alpha_p\in\mathscr{P}(\Delta)
 \quad\text{for }p=1,\ldots,m,
\]
and
\[
 \alpha_1+\alpha_2+\cdots+\alpha_m=\alpha.
\]
Moreover, for every sequence satisfying these conditions,
\[
 m=\operatorname{ord}_\Delta(\alpha).
\]
\end{lemma}

\begin{proof}
This is immediate from Proposition~\ref{sga3:XXI:3.1.2}.
\end{proof}

% END INCLUDED SGA3_English_current_buildable_source_20260823/source/inputs/XXI/tex/components/07_expose_XXI_lemma3211_through_lemma3216_en.tex

% BEGIN INCLUDED SGA3_English_current_buildable_source_20260823/source/inputs/XXI/tex/components/08_expose_XXI_lemma331_through_corollary339_en.tex
\subsection{Characterization and conjugacy of systems of positive roots}

\begin{lemma}\label{sga3:XXI:3.3.1}
If \(\alpha\in\mathscr{P}(\Delta)\), \(\beta\in\Delta\), and
\(\alpha\) and \(\beta\) are not proportional, then
\[
 s_\beta(\alpha)\in\mathscr{P}(\Delta).
\]
\end{lemma}

\begin{proof}
One has
\[
 s_\beta(\alpha)=\alpha-(\beta^\ast,\alpha)\beta.
\]
At least one simple root other than \(\beta\) occurs with a nonzero,
and hence strictly positive, coefficient in the expression of
\(\alpha\). It occurs with the same coefficient in the expression of
\(s_\beta(\alpha)\); therefore \(s_\beta(\alpha)\) is also positive.
\end{proof}

\begin{corollary}\label{sga3:XXI:3.3.2}
If \(\beta\in\Delta\), then the reflection \(s_\beta\) permutes the
elements of \(\mathscr{P}(\Delta)\) that are not proportional to
\(\beta\).
\end{corollary}

\begin{lemma}\label{sga3:XXI:3.3.3}
Let
\[
 \alpha\in\mathscr{P}(\Delta)\setminus\Delta
\]
be indivisible. There exists \(\beta\in\Delta\) such that
\[
 s_\beta(\alpha)\in\mathscr{P}(\Delta)
 \qquad\text{and}\qquad
 \operatorname{ord}_\Delta\bigl(s_\beta(\alpha)\bigr)
   <\operatorname{ord}_\Delta(\alpha).
\]
\end{lemma}

\begin{proof}
By Lemma~\ref{sga3:XXI:3.1.1}, there is a \(\beta\in\Delta\) such
that \((\beta^\ast,\alpha)>0\). Since \(\alpha\notin\Delta\) and
\(\alpha\) is indivisible, it cannot be proportional to \(\beta\).
Thus \(s_\beta(\alpha)\in\mathscr{P}(\Delta)\) by
Lemma~\ref{sga3:XXI:3.3.1}, and
\[
 \operatorname{ord}_\Delta\bigl(s_\beta(\alpha)\bigr)
 =\operatorname{ord}_\Delta(\alpha)-(\beta^\ast,\alpha)
 <\operatorname{ord}_\Delta(\alpha).
\]
\end{proof}

\begin{corollary}\label{sga3:XXI:3.3.4}
If \(\alpha\in\mathscr{P}(\Delta)\) is indivisible, then there is a
sequence \(\beta_p\in\Delta\), \(p=1,\ldots,q\), such that
\[
 \alpha_p
 =s_{\beta_p}\cdots s_{\beta_1}(\alpha)
 \in\mathscr{P}(\Delta)
 \quad\text{for }p=1,\ldots,q,
\]
and \(\alpha_q\in\Delta\).
\end{corollary}

\begin{proof}
Apply Lemma~\ref{sga3:XXI:3.3.3} repeatedly, inducting on
\(\operatorname{ord}_\Delta(\alpha)\).
\end{proof}

\begin{proposition}\label{sga3:XXI:3.3.5}
The Weyl group is generated by the \(s_\alpha\), for
\(\alpha\in\Delta\). Every indivisible root is conjugate under an
element of the Weyl group to a simple root.
\end{proposition}

\begin{proof}
The second assertion follows immediately from
Corollary~\ref{sga3:XXI:3.3.4}. The first then follows from
Corollary~\ref{sga3:XXI:1.2.10} and
Proposition~\ref{sga3:XXI:2.1.1}.
\end{proof}

\begin{proposition}\label{sga3:XXI:3.3.6}
Let \(R_+\) be a system of positive roots, and let \(P'\subset R\) be
closed and satisfy (P 2). Then there exists \(w\in W\) such that
\[
 w(R_+)\subset P'.
\]
\end{proposition}

We first state the corollaries.

\begin{corollary}\label{sga3:XXI:3.3.7}
The Weyl group acts transitively on the set of systems of positive
roots, and on the set of systems of simple roots.
\end{corollary}

\begin{corollary}\label{sga3:XXI:3.3.8}
A subset of \(R\) is a system of positive roots if and only if it is
closed and satisfies (P 1) and (P 2).
\end{corollary}

\begin{corollary}\label{sga3:XXI:3.3.9}
Suppose that \(\Gamma_0(R)\) is endowed with an ordered-group
structure for which every root is either \(>0\) or \(<0\). Then the
set of positive roots for this order is a system of positive roots.
\end{corollary}

% END INCLUDED SGA3_English_current_buildable_source_20260823/source/inputs/XXI/tex/components/08_expose_XXI_lemma331_through_corollary339_en.tex

% BEGIN INCLUDED SGA3_English_current_buildable_source_20260823/source/inputs/XXI/tex/components/09_expose_XXI_proof336_through_definition33101_en.tex
\begin{proof}[Proof of Proposition~\ref{sga3:XXI:3.3.6}]
Choose \(w\in W\) such that
\(\operatorname{Card}(w(R_+)\cap P')\) is maximal. Replacing \(R_+\)
by \(w(R_+)\), we may assume that
\begin{equation}
 \operatorname{Card}(R_+\cap P')
 \geq
 \operatorname{Card}\bigl(R_+\cap s_\alpha(P')\bigr)
 \quad\text{for every }\alpha\in\mathscr{S}(R_+)=\Delta.
 \tag{\(\ast\)}\label{sga3:XXI:3.3.6:star}
\end{equation}
We show that \(R_+\subset P'\). Otherwise, since \(P'\) is closed,
there is an \(\alpha\in\Delta\) such that \(\alpha\notin P'\). As
\(P'\) satisfies (P 2), one then has \(-\alpha\in P'\), and also
\(-2\alpha\in P'\) if \(2\alpha\) is a root.

For every \(\beta\in R_+\cap P'\), one has \(\beta\neq\alpha\). If
\(2\alpha\) is not a root, then
\(s_\alpha(\beta)\in R_+\) by Lemma~\ref{sga3:XXI:3.3.1}; hence
\[
 s_\alpha(R_+\cap P')
 \subset R_+\cap s_\alpha(P').
\]
But \(R_+\cap s_\alpha(P')\) also contains
\(\alpha=s_\alpha(-\alpha)\), contradicting
\eqref{sga3:XXI:3.3.6:star}. If \(2\alpha\) is a root, the same
argument applies.
\end{proof}

Thus, to study closed sets of roots satisfying (P 2), it is enough to
consider those that contain a system of positive roots.

\begin{proposition}\label{sga3:XXI:3.3.10}
Let \(R_+\) be a system of positive roots, and let \(P'\) be a closed
subset of \(R\) containing \(R_+\). Put
\[
 \Delta=\mathscr{S}(R_+),
 \qquad
 \Delta'=\Delta\cap(-P').
\]
Then \(P'\) is the union of \(R_+\) and the set of roots that are
linear combinations of the elements of \(\Delta'\) with negative
coefficients.
\end{proposition}

\begin{proof}
We argue by induction on the order of a root
\(\gamma\in P'\setminus R_+\). If
\(\operatorname{ord}_\Delta(\gamma)=-1\), then
\(-\gamma\in\Delta\), and hence \(\gamma\in-\Delta'\).

Suppose that \(\operatorname{ord}_\Delta(\gamma)<-1\). By
Proposition~\ref{sga3:XXI:3.1.2}, there exists \(\beta\in\Delta\)
such that \(-\gamma-\beta\in R\). Then
\[
 0<\operatorname{ord}_\Delta(-\gamma-\beta)
  =\operatorname{ord}_\Delta(-\gamma)-1
  <\operatorname{ord}_\Delta(-\gamma).
\]
The first inequality shows that \(-\gamma-\beta\in R_+\).
Consequently, since \(R_+\cap(-R_+)=\varnothing\), and since
\(\gamma+\beta\) is the sum of two roots of \(P'\), the root
\(\gamma+\beta\) belongs to \(P'\setminus R_+\) and satisfies
\[
 \operatorname{ord}_\Delta(\gamma+\beta)
 >\operatorname{ord}_\Delta(\gamma).
\]
By induction, \(\gamma+\beta\) is a linear combination of the
elements of \(\Delta'\) with negative coefficients. Since
\(\gamma=(\gamma+\beta)-\beta\), it remains only to show that
\(\beta\in-P'\). But \(\gamma\in P'\) and
\(-(\gamma+\beta)\in R_+\subset P'\), so
\[
 -\beta=\gamma-(\gamma+\beta)\in P'.
\]
\end{proof}

\begingroup
\renewcommand{\thetheorem}{3.3.10.1}
\begin{definition}\label{sga3:XXI:3.3.10.1}
\footnote{Editors' note: This definition was inserted here.}
A subset \(R'\) of \(R\) is called \emph{symmetric} if
\[
 -R'=R'.
\]
\end{definition}
\endgroup
\addtocounter{theorem}{-1}

% END INCLUDED SGA3_English_current_buildable_source_20260823/source/inputs/XXI/tex/components/09_expose_XXI_proof336_through_definition33101_en.tex

% BEGIN INCLUDED SGA3_English_current_buildable_source_20260823/source/inputs/XXI/tex/components/10_expose_XXI_scholium3311_through_remark345_en.tex
\begin{scholium}\label{sga3:XXI:3.3.11}
Let \(P'\) be a set of roots satisfying (P 2) and closed. There exist a
system of simple roots \(\Delta\) and a subset \(\Delta'\) of
\(\Delta\) such that
\[
 P'=R\cap\bigl(\mathbb{N}\cdot\Delta
          \,\cup\, -\mathbb{N}\cdot\Delta'\bigr).
\]
If
\[
 R_{\Delta'}=(\mathbb{Z}\cdot\Delta')\cap R,
\]
which is a closed and symmetric subset of \(R\), then \(P'\) is the
disjoint union of \(R_{\Delta'}\) and the closed subset of
\(\mathscr{P}(\Delta)\) formed by the
\[
 \alpha\in\mathscr{P}(\Delta)\setminus
           \mathbb{N}\cdot\Delta',
\]
that is, the positive roots whose decomposition with respect to
\(\Delta\) ``contains at least one root of
\(\Delta\setminus\Delta'\)''.
\end{scholium}

\subsection{Closed and symmetric sets of roots}
\label{sga3:XXI:3.4}

\begin{proposition}\label{sga3:XXI:3.4.1}
Let
\[
 \cR=(M,M^\ast,R,R^\ast)
\]
be a root datum, and let \(R'\) be a closed and symmetric subset of
\(R\). Then:
\begin{enumerate}
\item[(i)] \(\cR'=(M,M^\ast,R',R'^\ast)\) is a root datum;
\item[(ii)] for every system of positive roots \(R_+\) of \(R\),
  \(R'_+=R_+\cap R'\) is a system of positive roots of \(R'\);
\item[(iii)] the Weyl group \(W(\cR')\) of \(\cR'\) is the subgroup
  of \(W(\cR)\) generated by the \(s_\alpha\), for
  \(\alpha\in R'\).
\end{enumerate}
\end{proposition}

\begin{proof}
The first assertion is immediate from
Corollary~\ref{sga3:XXI:1.2.11}; the second follows from
Corollary~\ref{sga3:XXI:3.3.8}; and the third is evident.
\end{proof}

\begin{corollary}\label{sga3:XXI:3.4.2}
Let \(w\in W(\cR')\). The order of \(w\) is the smallest integer
\(n>0\) such that
\[
 w^n(\alpha')=\alpha'
 \qquad\text{for every }\alpha'\in R'.
\]
\end{corollary}

\begin{proof}
Apply Corollary~\ref{sga3:XXI:1.2.7} to the root datum \(\cR'\).
\end{proof}

\begin{corollary}\label{sga3:XXI:3.4.3}
Let \(\alpha\) and \(\beta\) be two nonproportional roots. Let \(n\)
be the smallest integer \(>0\) such that
\[
 (s_\alpha s_\beta)^n(\alpha)=\alpha,
 \qquad
 (s_\alpha s_\beta)^n(\beta)=\beta.
\]
Then the subgroup \(W_{\alpha,\beta}\) of \(W\) generated by
\(s_\alpha\) and \(s_\beta\) is defined by the relations
\[
 s_\alpha^2=1,\qquad
 s_\beta^2=1,\qquad
 (s_\alpha s_\beta)^n=1.
\]
\end{corollary}

\begin{proof}
In view of Corollary~\ref{sga3:XXI:3.4.2}, it remains only to verify
the following lemma.
\end{proof}

\begin{lemma}\label{sga3:XXI:3.4.4}
Let \(G\) be the group generated by two elements \(x\) and \(y\)
subject to the relations \(x^2=y^2=1\). Every normal subgroup of
\(G\) containing neither \(x\) nor \(y\) is generated, as a normal
subgroup, by an element of the form \((xy)^n\), where \(n>0\).
\end{lemma}

\begin{proof}
\footnote{Editors' note: The original has been corrected in what
follows.}
Every element of \(G\) can be written as \((xy)^n\), as
\((yx)^n=(xy)^{-n}\), or in one of the forms
\[
\begin{array}{c@{\qquad}c@{\qquad}c}
 (a) & x(yx)^{2n} & \text{or } y(xy)^{2n+1},\\[2mm]
 (b) & y(xy)^{2n} & \text{or } x(yx)^{2n+1},
\end{array}
\]
where \(n\in\mathbb{N}\). The elements of type (a), respectively of
type (b), are conjugate to \(x\), respectively to \(y\).
\end{proof}

\begin{remark}\label{sga3:XXI:3.4.5}
The integer \(n\) is computed immediately. Put
\[
 (\alpha^\ast,\beta)=p,\qquad
 (\beta^\ast,\alpha)=q.
\]
Then\footnote{Editors' note: The original has been corrected in what
follows.}
\[
 (s_\alpha s_\beta)(\alpha)=(pq-1)\alpha-q\beta,
 \qquad
 (s_\alpha s_\beta)(\beta)=p\alpha-\beta.
\]
Thus \(n\) is the order of the matrix
\[
 \begin{pmatrix}
  pq-1 & p\\
  -q   & -1
 \end{pmatrix}.
\]
If \(pq=0\), then \(p=q=0\) by
Corollary~\ref{sga3:XXI:2.2.2}, whence \(n=2\). Otherwise,
Proposition~\ref{sga3:XXI:2.3.1} shows that \(pq\) equals \(1\),
\(2\), or \(3\), and one finds respectively \(n=3\), \(4\), or
\(6\).

\smallskip
\noindent
\emph{N.B.} Expressing the finiteness of the order of the preceding
matrix recovers inequality \SGARefLink{sga3:XXI:eq:13}{(13)} of
Proposition~\ref{sga3:XXI:2.3.1}.
\end{remark}

% END INCLUDED SGA3_English_current_buildable_source_20260823/source/inputs/XXI/tex/components/10_expose_XXI_scholium3311_through_remark345_en.tex

% BEGIN INCLUDED SGA3_English_current_buildable_source_20260823/source/inputs/XXI/tex/components/11_expose_XXI_definition346_through_lemma349_en.tex
\begin{definition}\label{sga3:XXI:3.4.6}
Let \(\Delta\) be a system of simple roots and let
\(\Delta'\subset\Delta\). We write
\[
 R_{\Delta'}
   =R\cap(\mathbb{Q}\cdot\Delta')
   =R\cap(\mathbb{Z}\cdot\Delta').
\]
\end{definition}

\begin{lemma}\label{sga3:XXI:3.4.7}
The set \(R_{\Delta'}\) is closed and symmetric, and \(\Delta'\) is
a system of simple roots of the root datum
\[
 (M,M^\ast,R_{\Delta'},R_{\Delta'}^\ast).
\]
Its Weyl group is the subgroup \(W_{\Delta'}\) of \(W\) generated by
the \(s_\alpha\), for \(\alpha\in\Delta'\). Moreover,
\[
 \Delta\cap R_{\Delta'}=\Delta'.
\]
\end{lemma}

\begin{proof}
Trivial.
\end{proof}

\begin{proposition}\label{sga3:XXI:3.4.8}
Let \(R'\subset R\). The following conditions are equivalent:
\begin{enumerate}
\item[\SGAInlineRefTarget{sga3:XXI:3.4.8:item:i}(i)] There exists a vector subspace \(V'\) of \(V\), or of
  \(\VR(R)\), such that \(R'=R\cap V'\).
\item[\SGAInlineRefTarget{sga3:XXI:3.4.8:item:ii}(ii)] There exist a system of simple roots \(\Delta\) of \(R\)
  and a subset \(\Delta'\) of \(\Delta\) such that
  \(R'=R_{\Delta'}\).
\end{enumerate}
More precisely, under these conditions, every system of simple roots
\(\Delta'\) of \((M,M^\ast,R',R'^\ast)\) is contained in a system of
simple roots \(\Delta\) of \(R\), and one has
\(R'=R_{\Delta'}\).
\end{proposition}

\begin{proof}
The implication \SGARefLink{sga3:XXI:3.4.8:item:ii}{(ii)} \(\Rightarrow\) \SGARefLink{sga3:XXI:3.4.8:item:i}{(i)} is clear. Suppose \SGARefLink{sga3:XXI:3.4.8:item:i}{(i)}.
Then \(R'\) is closed and symmetric, so
\((M,M^\ast,R',R'^\ast)\) is a root datum. Let \(\Delta'\) be a
system of simple roots of this datum, and put
\[
 R'_+=\mathscr{P}(\Delta').
\]
If \(V'=V\), then \(\Delta'\) is a system of simple roots of
\(\cR\), and there is nothing to prove. Otherwise, there exists
\(x\in V^\ast\) such that
\[
 (x,V')=\{0\},
 \qquad
 (x,\alpha)\ne0
 \quad\text{for every }\alpha\in R\setminus R'.
\]
Put
\[
 R_+(x)=\{\alpha\in R\mid(x,\alpha)>0\},
 \qquad
 R_+=R_+(x)\cup R'_+.
\]
For every \(\alpha\in R\), one has
\[
\begin{aligned}
 (x,\alpha)>0 &\Longleftrightarrow \alpha\in R_+(x),\\
 (x,\alpha)<0 &\Longleftrightarrow \alpha\in-R_+(x),\\
 (x,\alpha)=0 &\Longleftrightarrow \alpha\in R'.
\end{aligned}
\]
It follows at once from Corollary~\ref{sga3:XXI:3.3.8} that \(R_+\)
is a system of positive roots of \(R\). Put
\(\Delta=\mathscr{S}(R_+)\). It is enough to prove that
\(\Delta'\subset\Delta\). Otherwise, let
\(\alpha\in\Delta'\setminus\Delta\). By
Corollary~\ref{sga3:XXI:3.2.14}, there exist
\(\beta,\gamma\in R_+\) such that
\[
 \alpha=\beta+\gamma.
\]
If \(\beta,\gamma\in R_+(x)\), then \(\alpha\in R_+(x)\), which is
impossible because \((x,\Delta')=0\). If one of \(\beta,\gamma\),
say \(\beta\), belongs to \(R'_+\), then, since \(R'\) is symmetric
and closed,
\[
 \gamma=\alpha-\beta\in R_+\cap R'=R'_+.
\]
But then \(\alpha\) is not simple in \(R'_+\), again a
contradiction.
\end{proof}

\begin{lemma}\label{sga3:XXI:3.4.9}
Under the preceding conditions, let
\[
 \alpha\in\mathscr{P}(\Delta)\setminus R'.
\]
For every \(w\in W_{\Delta'}\), one has
\[
 w(\alpha)\in\mathscr{P}(\Delta)\setminus R'.
\]
\end{lemma}

\begin{proof}
It suffices to verify the assertion for
\(w=s_\beta\), \(\beta\in\Delta'\). It then follows from
Lemma~\ref{sga3:XXI:3.3.1} and from
\(s_\beta(R')=R'\).
\end{proof}

% END INCLUDED SGA3_English_current_buildable_source_20260823/source/inputs/XXI/tex/components/11_expose_XXI_definition346_through_lemma349_en.tex

% BEGIN INCLUDED SGA3_English_current_buildable_source_20260823/source/inputs/XXI/tex/components/12_expose_XXI_lemma3410_through_proposition354_en.tex
\begin{lemma}\label{sga3:XXI:3.4.10}
Let \(w\in W\). Under the conditions of
Proposition~\ref{sga3:XXI:3.4.8}, the following conditions are
equivalent:
\begin{enumerate}
\item[\SGAInlineRefTarget{sga3:XXI:3.4.10:item:i}(i)] \(w\in W_{\Delta'}\).
\item[\SGAInlineRefTarget{sga3:XXI:3.4.10:item:ii}(ii)] For every \(m\in M\), one has \(w(m)-m\in V'\).
\item[(iii)] For every \(\alpha\in R\), one has
  \(w(\alpha)-\alpha\in V'\).
\end{enumerate}
\end{lemma}

\begin{proof}
Clearly \SGARefLink{sga3:XXI:3.4.10:item:i}{(i)} implies \SGARefLink{sga3:XXI:3.4.10:item:ii}{(ii)}, and \SGARefLink{sga3:XXI:3.4.10:item:ii}{(ii)} implies \SGARefLink{sga3:XXI:2.3.7:corollary:item:iii}{(iii)}. Suppose \SGARefLink{sga3:XXI:2.3.7:corollary:item:iii}{(iii)}.
Write
\[
 w=s_{\alpha_n}\cdots s_{\alpha_1},
 \qquad \alpha_i\in\Delta,
\]
and argue by induction on \(n\) that every \(\alpha_i\) belongs to
\(\Delta'\). We may assume that
\[
 w'=s_{\alpha_{n-1}}\cdots s_{\alpha_1}\in W_{\Delta'}.
\]
Then, for every \(\alpha\in R\),
\[
 w(\alpha)-\alpha
   =w'(\alpha)-\alpha-(\alpha_n^\ast,w'(\alpha))\alpha_n.
\]
Taking \(\alpha=w'^{-1}(\alpha_n)\), one obtains
\(2\alpha_n\in V'\), hence \(\alpha_n\in\Delta'\). Therefore
\[
 w=s_{\alpha_n}w'\in W_{\Delta'}.
\]
\end{proof}

\subsection{Miscellaneous remarks}
\label{sga3:XXI:3.5}

\begin{proposition}\label{sga3:XXI:3.5.1}
Let \(R_+\) be a system of positive roots. Put
\[
 \rho_{R_+}
   =\frac12\sum_{\alpha\in\operatorname{ind}(R_+)}\alpha.
\]
Then
\[
 (\beta^\ast,\rho_{R_+})=1
 \qquad\text{for every }\beta\in\mathscr{S}(R_+).
\]
\end{proposition}

\begin{proof}
One can write
\[
 2\rho_{R_+}
   =\beta+
     \sum_{\substack{\alpha\in\operatorname{ind}(R_+)\\
                     \alpha\ne\beta}}\alpha.
\]
Hence
\[
 s_\beta(2\rho_{R_+})=2\rho_{R_+}-2\beta
\]
by Lemma~\ref{sga3:XXI:3.3.2}.
\end{proof}

\begin{corollary}\label{sga3:XXI:3.5.2}
Put
\[
 \rho_{R_+}^\ast
   =\frac12
     \sum_{\alpha^\ast\in\operatorname{ind}(R_+^\ast)}
     \alpha^\ast.
\]
Then:
\begin{enumerate}
\item[(i)] \((\rho_{R_+}^\ast,\alpha)>0\) for every
  \(\alpha\in R_+\); equivalently,
  \(\rho_{R_+}^\ast\in\mathcal{C}(R_+)\), in the notation of
  Definition~\ref{sga3:XXI:3.6.8}.
\item[(ii)] For every \(\alpha\in R\),
  \[
   \operatorname{ord}_{\mathscr{S}(R_+)}(\alpha)
      =(\rho_{R_+}^\ast,\alpha).
  \]
\end{enumerate}
\footnote{Editors' note: By Lemma~\ref{sga3:XXI:3.1.5}, it is enough
to verify the formula for \(\alpha\in\Delta\). Proposition
\ref{sga3:XXI:3.5.1}, applied to \(R^\ast\), then gives
\((\rho_{R_+}^\ast,\alpha)=1
=\operatorname{ord}_{\mathscr{S}(R_+)}(\alpha)\).}
\end{corollary}

\begin{remark}\label{sga3:XXI:3.5.3}
If \(w\in W\), then
\[
 \rho_{w(R_+)}=w(\rho_{R_+}),
 \qquad
 \rho_{w(R_+)}^\ast=w(\rho_{R_+}^\ast).
\]
\end{remark}

\begin{proposition}\label{sga3:XXI:3.5.4}
Let \(\alpha\) and \(\gamma\) be two nonproportional roots, with
\(\alpha\) indivisible. There exist a system of simple roots
containing \(\alpha\) and a root \(\beta\) such that
\[
 \gamma=a\alpha+b\beta,
 \qquad a,b\in\mathbb{N}.
\]
\end{proposition}

\begin{proof}
Choose a basis of the vector space \(\VR(R)\) containing
\(\alpha_1=\alpha\) and \(\alpha_2=\gamma\), and consider the
lexicographic order with respect to this basis. Let \(R_+\) be the set
of roots \(>0\). It is clear that \(R_+\) is a system of positive
roots and\footnote{Editors' note: The original was modified here to
take account of the case in which \(2\alpha\) is a root.} that the
least element of \(R_+\) not proportional to \(\alpha\) is simple.
This element has the form
\[
 \beta=p\alpha+q\gamma,
 \qquad 0<q\le1.
\]
Consequently
\[
 \gamma=q^{-1}\beta-q^{-1}p\alpha.
\]
Since \(q^{-1}>0\), one has
\[
 q^{-1}\in\mathbb{N}^{\ast},
 \qquad
 -q^{-1}p\in\mathbb{N}.
\]
\end{proof}

% END INCLUDED SGA3_English_current_buildable_source_20260823/source/inputs/XXI/tex/components/12_expose_XXI_lemma3410_through_proposition354_en.tex

% BEGIN INCLUDED SGA3_English_current_buildable_source_20260823/source/inputs/XXI/tex/components/13_expose_XXI_proposition355_through_lemma361_en.tex
We conclude with two remarks on the group \(\Gamma_0(R)\).

\begin{proposition}\label{sga3:XXI:3.5.5}
Let \(G\) be an abelian group, and let \(f:R\to G\) be a map
satisfying:
\begin{enumerate}
\item[(i)] if \(\alpha\in R\), then
  \(f(-\alpha)=-f(\alpha)\);
\item[(ii)] if \(\alpha,\beta,\alpha+\beta\in R\), then
  \[
   f(\alpha+\beta)=f(\alpha)+f(\beta).
  \]
\end{enumerate}
There exists a unique group homomorphism
\[
 \overline f:\Gamma_0(R)\longrightarrow G
\]
such that \(\overline f(\alpha)=f(\alpha)\) for every
\(\alpha\in R\).
\end{proposition}

\begin{proof}
Let \(\Delta\) be a system of simple roots of \(R\), and write
\[
 \beta=\sum_{\alpha\in\Delta}a(\alpha)\alpha.
\]
Lemma~\ref{sga3:XXI:3.2.16} gives at once
\[
 f(\beta)=\sum_{\alpha\in\Delta}a(\alpha)f(\alpha).
\]
Now \(\Delta\) is a basis of \(\Gamma_0(R)\).
\end{proof}

\begin{proposition}\label{sga3:XXI:3.5.6}
Let \(\Delta\) be a system of simple roots. There exists on
\(\Gamma_0(R)\) a totally ordered group structure such that the roots
\(>0\) are the elements of \(\mathscr{P}(\Delta)\), and such that
\[
 \alpha\longmapsto\operatorname{ord}_\Delta(\alpha)
\]
is increasing.
\end{proposition}

\begin{proof}
Let
\[
 \alpha_1,\alpha_2,\ldots,\alpha_n,
 \qquad n=\rgss(\cR),
\]
be the elements of \(\Delta\). Every \(x\in\Gamma_0(R)\) has a
decomposition
\[
 x=\sum_{i=1}^{n}m_i(x)\alpha_i.
\]
Take the lexicographic order with respect to the functions
\[
 \sum_i m_i,\quad m_n,\quad m_{n-1},\quad\ldots,\quad m_2.
\]
\end{proof}

\begin{remark}\label{sga3:XXI:3.5.7}
The first roots, in order, are
\[
 \alpha_1,\alpha_2,\ldots,\alpha_n;
\]
then come, when they are roots,
\[
 2\alpha_1,\quad
 \alpha_1+\alpha_2,\quad
 \alpha_1+\alpha_3,\quad\ldots.
\]
\end{remark}

\subsection{Weyl chambers}
\label{sga3:XXI:3.6}

\begin{lemma}\label{sga3:XXI:3.6.1}
Let \(V\) be a finite-dimensional real vector
space.\footnote{Editors' note: The original \(\mathbb{Q}\) has been
replaced by \(\mathbb{R}\).} Let \(f_i\) be independent linear forms.
Put
\[
 C=\{x\in V\mid f_i(x)>0\}.
\]
Then \(C\) is a maximal convex subset of
\[
 X=V\setminus\bigcup_i f_i^{-1}(0).
\]
\end{lemma}

\begin{proof}
Trivial.
\end{proof}

% END INCLUDED SGA3_English_current_buildable_source_20260823/source/inputs/XXI/tex/components/13_expose_XXI_proposition355_through_lemma361_en.tex

% BEGIN INCLUDED SGA3_English_current_buildable_source_20260823/source/inputs/XXI/tex/components/14_expose_XXI_definition362_through_definition368_en.tex
\begin{definition}\label{sga3:XXI:3.6.2}
A subset \(C\subset V\) described by the construction of
Lemma~\ref{sga3:XXI:3.6.1} is called, here, a \emph{chamber} of \(V\).
\end{definition}

\begin{definition}\label{sga3:XXI:3.6.3}
A hyperplane \(H\) of \(V\) is called a \emph{wall} of \(C\) if
\[
 H\cap(\overline C\setminus C)
\]
contains a nonempty open subset of \(H\).
\end{definition}

\begin{remark}\label{sga3:XXI:3.6.4}
For a convex subset, the closure can be described without reference
to the topology of \(V\): it is the set of endpoints of all open
segments contained in the given subset.
\end{remark}

\begin{lemma}\label{sga3:XXI:3.6.5}
Under the conditions of Lemma~\ref{sga3:XXI:3.6.1},
\[
 \overline C=\{x\in V\mid f_i(x)\ge0\}.
\]
The walls of \(C\) are the hyperplanes \(f_i^{-1}(0)\).
\end{lemma}

\begin{proof}
The first assertion is clear. The second follows from
\[
 \overline C\setminus C\subset\bigcup_i f_i^{-1}(0)
\]
and from the fact that the \(f_i^{-1}(0)\) are visibly walls of
\(C\).
\end{proof}

\begin{proposition}\label{sga3:XXI:3.6.6}
Let \(C\) be a chamber of \(V\). If
\[
 H_i,\qquad i=1,2,\ldots,n,
\]
are the distinct walls of \(C\), then, for every system of linear
forms \(\{u_i\}\) such that \(H_i=u_i^{-1}(0)\), there exist
\(\varepsilon_i\in\{-1,+1\}\) for which
\[
 C=\{x\in V\mid\varepsilon_i u_i(x)>0\}.
\]
For every wall \(H\) of \(C\),
\[
 H\cap C=\varnothing,
 \qquad
 H\cap\overline C=\overline{C_H},
\]
where \(C_H\) is a chamber in \(H\). The walls of \(C_{H_i}\) are
the \(H_j\cap H_i\), for \(j\ne i\).
\end{proposition}

\begin{proof}
This follows immediately from the lemma.
\end{proof}

\begin{definition}\label{sga3:XXI:3.6.7}
The \(C_{H_i}\) are called the \emph{faces} of \(C\).
\end{definition}

Let now
\[
 \cR=(M,M^\ast,R,R^\ast)
\]
be a root datum, and put
\[
 V_{\mathbb{R}}^\ast=M^\ast\otimes\mathbb{R}.
\]

\begin{definition}\label{sga3:XXI:3.6.8}
\footnote{Editors' note: The definition of the hyperplanes
\(\mathcal{H}_\alpha\) was added here.}
For every \(\alpha\in R\), put
\[
 \mathcal{H}_\alpha
   =\{x\in V_{\mathbb{R}}^\ast\mid(x,\alpha)=0\}.
\]
Write
\[
 X=V_{\mathbb{R}}^\ast
   \setminus\bigcup_{\alpha\in R}\mathcal{H}_\alpha.
\]
For every \(x\in X\), put
\[
 R_+(x)=\{\alpha\in R\mid(x,\alpha)>0\}.
\]
For every system of positive roots \(R_+\), put
\[
 \mathcal{C}(R_+)
   =\{x\in V_{\mathbb{R}}^\ast
       \mid (x,\alpha)>0\text{ for every }\alpha\in R_+\}.
\]
\end{definition}

% END INCLUDED SGA3_English_current_buildable_source_20260823/source/inputs/XXI/tex/components/14_expose_XXI_definition362_through_definition368_en.tex

% BEGIN INCLUDED SGA3_English_current_buildable_source_20260823/source/inputs/XXI/tex/components/15_expose_XXI_proposition369_through_remark3614_en.tex
\begin{proposition}\label{sga3:XXI:3.6.9}
\begin{enumerate}
\item[\SGAInlineRefTarget{sga3:XXI:3.6.9:item:i}(i)] For every \(x\in X\), \(R_+(x)\) is a system of positive
roots. For every system of positive roots \(R_+\),
\(\mathcal{C}(R_+)\) is a chamber in \(V_{\mathbb R}^\ast\). The
\(\mathcal{C}(R_+)\) are the maximal convex subsets of \(X\).
\item[\SGAInlineRefTarget{sga3:XXI:3.6.9:item:ii}(ii)] Let \(\Delta\) be a system of simple roots. One has
\[
 \mathcal{C}(\mathscr{P}(\Delta))
 =\{x\in V_{\mathbb R}^\ast
       \mid (x,\alpha)>0\text{ for every }\alpha\in\Delta\}.
\]
The walls of \(\mathcal{C}(\mathscr{P}(\Delta))\) are the
hyperplanes
\[
 \mathcal{H}_\alpha=\alpha^{-1}(0),
 \qquad \alpha\in\Delta;
\]
its faces are
\[
 C_\alpha
 =\{x\in V_{\mathbb R}^\ast
      \mid (x,\alpha)=0,\ 
      (x,\beta)>0\text{ for }\beta\in\Delta,\ \beta\ne\alpha\}.
\]
\item[\SGAInlineRefTarget{sga3:XXI:3.6:item:iii}(iii)] One has the equivalence
\[
 R_+(x)=R_+
 \quad\Longleftrightarrow\quad
 x\in\mathcal{C}(R_+).
\]
\end{enumerate}
\end{proposition}

\begin{proof}
It is first clear that \(R_+(x)\) is a system of positive roots.
Since \(x\in\mathcal{C}(R_+(x))\), the union of the
\(\mathcal{C}(R_+(x))\) is \(X\). Property \SGARefLink{sga3:XXI:3.6:item:iii}{(iii)} is immediate; it
follows that the \(\mathcal{C}(R_+)\) form a partition of \(X\).
The first assertion of \SGARefLink{sga3:XXI:3.6.9:item:ii}{(ii)} is evident. It follows at once that
\(\mathcal{C}(R_+)\) is a chamber in \(V_{\mathbb R}^\ast\), which
proves the remainder of \SGARefLink{sga3:XXI:3.6.9:item:i}{(i)}. Put
\[
 C=\mathcal{C}(\mathscr{P}(\Delta)).
\]
It remains only to note that
\[
 \overline C\cap
 \left(V_{\mathbb R}^\ast\setminus
       \bigcup_{\alpha\in R}\alpha^{-1}(0)\right)
 =
 \overline C\cap
 \left(V_{\mathbb R}^\ast\setminus
       \bigcup_{\alpha\in\Delta}\alpha^{-1}(0)\right)
\]
to complete the proof of \SGARefLink{sga3:XXI:3.6.9:item:ii}{(ii)} by
Lemma~\ref{sga3:XXI:3.6.1}.
\end{proof}

\begin{definition}\label{sga3:XXI:3.6.10}
The \(\mathcal{C}(R_+)\) are called the \emph{Weyl chambers} of the
root datum. Conversely, for every Weyl chamber \(C\), put
\[
 R_+(C)=R_+(x)
\]
for any \(x\in C\).
\end{definition}

\begin{corollary}\label{sga3:XXI:3.6.11}
The maps
\[
 R_+\longmapsto\mathcal{C}(R_+),
 \qquad
 C\longmapsto R_+(C)
\]
give a bijective correspondence between systems of positive roots
and Weyl chambers.
\end{corollary}

This correspondence is invariant under the Weyl group:

\begin{lemma}\label{sga3:XXI:3.6.12}
If \(w\in W(\mathcal{R})\), then
\[
 \mathcal{C}(w(R_+))=w(\mathcal{C}(R_+)).
\]
\end{lemma}

\begin{corollary}\label{sga3:XXI:3.6.13}
The correspondences
\[
 \Delta\longleftrightarrow R_+\longleftrightarrow C
\]
are isomorphisms of homogeneous spaces under \(W(\mathcal{R})\).
\end{corollary}

We shall see later that these homogeneous spaces are principal
(\SGARefLink{sga3:XXI:5.5}{5.5}).

\begin{remark}\label{sga3:XXI:3.6.14}
If \(C\) is a Weyl chamber, then \(-C\) is also one; it is called the
\emph{opposite} of \(C\). Consequently there exists
\(w_0\in W\)---and in fact exactly one, cf.
\SGARefLink{sga3:XXI:5.5}{5.5}---such that
\[
 w_0(C)=-C.
\]
It is called the \emph{symmetry} of the root datum relative to the
Weyl chamber \(C\) (or relative to \(R_+(C)\), or to
\(\mathscr{S}(R_+(C))\), and so on).
\end{remark}

% END INCLUDED SGA3_English_current_buildable_source_20260823/source/inputs/XXI/tex/components/15_expose_XXI_proposition369_through_remark3614_en.tex

% BEGIN INCLUDED SGA3_English_current_buildable_source_20260823/source/inputs/XXI/tex/components/16_expose_XXI_section4_rank_two_tables_en.tex
\section{Reduced root data of semisimple rank \(2\)}
\label{sga3:XXI:4}

\renewcommand{\thetheorem}{\thesection.\arabic{theorem}}
\setcounter{theorem}{0}

\SGARefTarget{sga3:XXI:4.0}
\noindent\textbf{4.0.}\footnote{Editors' note: The numbering \(4.0\)
was added for later references.}
Let \(\mathcal{R}\) be a root datum of semisimple rank \(2\). Let
\(\{\alpha,\beta\}\) be a system of simple roots. Suppose that
\(\ell(\alpha)\le\ell(\beta)\). By
\ref{sga3:XXI:2.3.1} and \ref{sga3:XXI:3.2.1}, there are then four
possibilities:
\[
\begin{array}{c|c|c|c|c}
\text{Type}
 &\ell(\beta)/\ell(\alpha)
 &\ell(\beta^\ast)/\ell(\alpha^\ast)
 &(\beta^\ast,\alpha)
 &(\alpha^\ast,\beta)\\ \hline
A_1\times A_1&-&-&0&0\\
A_2&1&1&-1&-1\\
B_2&2&1/2&-1&-2\\
G_2&3&1/3&-1&-3
\end{array}
\]
It follows from Remark~\ref{sga3:XXI:3.4.5} that the order of
\(s_\alpha s_\beta\) is respectively \(2,3,4,6\).

We study each of these systems separately and give the list of
indivisible roots.

\medskip
\noindent\emph{Type \(A_1\times A_1\).}
The indivisible roots are
\[
 \alpha,\quad\beta,\quad-\alpha,\quad-\beta.
\]
The corresponding coroots are
\[
 \alpha^\ast,\quad\beta^\ast,\quad-\alpha^\ast,\quad-\beta^\ast.
\]

\medskip
\noindent\emph{Type \(A_2\).}
The indivisible positive roots are as follows:
\[
\begin{array}{c|ccc}
\text{root }\gamma
 &\alpha&\beta&\alpha+\beta\\ \hline
\ell(\gamma)/\ell(\alpha)
 &1&1&1\\ \hline
\text{coroot }\gamma^\ast
 &\alpha^\ast&\beta^\ast&\alpha^\ast+\beta^\ast\\ \hline
\ell(\gamma^\ast)/\ell(\beta^\ast)
 &1&1&1
\end{array}
\]
The half-sum of the indivisible positive roots is
\[
 \rho=\alpha+\beta.
\]

\medskip
\noindent\emph{Type \(B_2\).}
The indivisible positive roots are as follows:
\[
\begin{array}{c|cccc}
\text{root }\gamma
 &\alpha&\beta&\alpha+\beta&2\alpha+\beta\\ \hline
\ell(\gamma)/\ell(\alpha)
 &1&2&1&2\\ \hline
\text{coroot }\gamma^\ast
 &\alpha^\ast&\beta^\ast
 &2\alpha^\ast+\beta^\ast&\alpha^\ast+\beta^\ast\\ \hline
\ell(\gamma^\ast)/\ell(\beta^\ast)
 &2&1&2&1
\end{array}
\]
The half-sum of the indivisible positive roots is
\[
 \rho=\frac{4\alpha+3\beta}{2}.
\]

\medskip
\noindent\emph{Type \(G_2\).}
The indivisible positive roots are as follows:
\begin{center}
\resizebox{\textwidth}{!}{%
$\displaystyle
\begin{array}{c|cccccc}
\text{root }\gamma
 &\alpha&\beta&\alpha+\beta&2\alpha+\beta
 &3\alpha+\beta&3\alpha+2\beta\\ \hline
\ell(\gamma)/\ell(\alpha)
 &1&3&1&1&3&3\\ \hline
\text{coroot }\gamma^\ast
 &\alpha^\ast&\beta^\ast&\alpha^\ast+3\beta^\ast
 &2\alpha^\ast+3\beta^\ast&\alpha^\ast+\beta^\ast
 &\alpha^\ast+2\beta^\ast\\ \hline
\ell(\gamma^\ast)/\ell(\beta^\ast)
 &3&1&3&3&1&1
\end{array}
$%
}
\end{center}
The half-sum of the indivisible positive roots is
\[
 \rho=5\alpha+3\beta.
\]

% END INCLUDED SGA3_English_current_buildable_source_20260823/source/inputs/XXI/tex/components/16_expose_XXI_section4_rank_two_tables_en.tex

% BEGIN INCLUDED SGA3_English_current_buildable_source_20260823/source/inputs/XXI/tex/components/17_expose_XXI_proposition41_through_lemma42_en.tex
\begin{proposition}\label{sga3:XXI:4.1}
\footnote{The following numbers \(4.1\) to \(4.4\) are used in the
proof of Theorem~\ref{sga3:XXI:5.1}. There are now simpler proofs of
that theorem; see \emph{Groupes et algèbres de Lie}, Chapter~V,
\S3, Theorem~1. Editors' note: The reference was made explicit and
this note, which occurred at \(4.2\) in the original, was moved here.}
Let \(n\) be the order of \(s_\alpha s_\beta\). Put \(u_0=0\), and,
for \(p\in\mathbb N\), define
\[
\left\{
\begin{aligned}
u_{2p+1}&=u_{2p}+(s_\alpha s_\beta)^p(\alpha),\\
u_{2p+2}&=u_{2p+1}+(s_\alpha s_\beta)^p s_\alpha(\beta).
\end{aligned}
\right.
\]
Then:
\footnote{Editors' note: Let \(\rho\) be the half-sum of the positive
roots, cf. Proposition~\ref{sga3:XXI:3.5.1}. If, as below, one puts
\(w_0=1\), \(w_1=s_\alpha\), \(w_2=s_\beta s_\alpha\),
\(w_3=s_\alpha s_\beta s_\alpha\), and so on, then
\(u_k=\rho-w_k^{-1}(\rho)\). This proves \SGARefLink{sga3:XXI:4.1:item:i}{(i)}, since \(w_{2n}=1\).}
\begin{enumerate}
\item[\SGAInlineRefTarget{sga3:XXI:4.1:item:i}(i)] \(u_{k+2n}=u_k\), for every \(k\in\mathbb N\).
\item[\SGAInlineRefTarget{sga3:XXI:4.1:item:ii}(ii)] \(u_0=0\), \(u_1=\alpha\), \(u_{2n-1}=\beta\),
and \(u_{2n}=0\).
\item[(iii)] If \(1<k<2n-1\), then
\[
 u_k=a_k\alpha+b_k\beta,
 \qquad a_k,b_k\in\mathbb N^\ast.
\]
\end{enumerate}
\end{proposition}

\begin{proof}
Assertion \SGARefLink{sga3:XXI:4.1:item:i}{(i)} follows from \((s_\alpha s_\beta)^n=1\) and
\(u_{2n}=0\).

We prove \SGARefLink{sga3:XXI:4.1:item:ii}{(ii)} and \SGARefLink{sga3:XXI:3.6:item:iii}{(iii)}. Direct calculation gives, in the four cases,
the following sequences of values:
\footnote{Editors' note: The original was corrected here to agree
with the convention \(\ell(\alpha)\le\ell(\beta)\) of \(4.0\).}
\[
\begin{array}{c|l}
A_1\times A_1&0,\alpha,\beta+\alpha,\beta,0,\\[2pt]
A_2&0,\alpha,\beta+2\alpha,2\beta+2\alpha,
       2\beta+\alpha,\beta,0,\\[2pt]
B_2&0,\alpha,\beta+3\alpha,2\beta+4\alpha,
       3\beta+4\alpha,3\beta+3\alpha,2\beta+\alpha,\beta,0,\\[2pt]
G_2&0,\alpha,\beta+4\alpha,2\beta+6\alpha,
       4\beta+9\alpha,5\beta+10\alpha,6\beta+10\alpha,\\
   &\quad 6\beta+9\alpha,5\beta+6\alpha,4\beta+4\alpha,
       2\beta+\alpha,\beta,0.
\end{array}
\]
\end{proof}

\begin{lemma}\label{sga3:XXI:4.2}
Put
\[
 w_{2p}=(s_\beta s_\alpha)^p,
 \qquad
 w_{2p+1}=s_\alpha(s_\beta s_\alpha)^p,
\]
so that \(w_0,\ldots,w_{2n-1}\) are the distinct elements of \(W\).
Let \(u_0,\ldots,u_{2n-1}\) be the elements of \(V\) defined in
Proposition~\ref{sga3:XXI:4.1}. For every \(x\in V^\ast\), one has
\[
 x-w_k(x)=n_k\alpha^\ast+m_k\beta^\ast,
\]
where \(n_k,m_k\in\mathbb Q\) and
\[
 n_k+m_k=(x,u_k).
\]
\footnote{Editors' note: In view of the preceding editor's note, this
follows immediately from Propositions~\ref{sga3:XXI:3.5.1} and
\ref{sga3:XXI:1.1.9}: one has
\[
 n_k+m_k=(x-w_k(x),\rho)
 =(x,\rho-w_k^{-1}(\rho))=(x,u_k).
\]}
\end{lemma}

\begin{proof}
The proof is by induction on \(k\). If \(k=0\), the formula is
trivial. We carry out, for example, the passage from \(w_{2p}\) to
\(w_{2p+1}\). Since \(w_{2p+1}=s_\alpha w_{2p}\),
\[
\begin{aligned}
x-w_{2p+1}(x)
 &=x-w_{2p}(x)+w_{2p}(x)-s_\alpha w_{2p}(x)\\
 &=n_{2p}\alpha^\ast+m_{2p}\beta^\ast
   +(w_{2p}(x),\alpha)\alpha^\ast.
\end{aligned}
\]
Consequently
\[
\begin{aligned}
n_{2p+1}+m_{2p+1}
 &=n_{2p}+m_{2p}
   +\bigl((s_\beta s_\alpha)^p(x),\alpha\bigr)\\
 &=(x,u_{2p})+\bigl(x,(s_\alpha s_\beta)^p(\alpha)\bigr)\\
 &=(x,u_{2p+1}).
\end{aligned}
\]
\end{proof}

% END INCLUDED SGA3_English_current_buildable_source_20260823/source/inputs/XXI/tex/components/17_expose_XXI_proposition41_through_lemma42_en.tex

% BEGIN INCLUDED SGA3_English_current_buildable_source_20260823/source/inputs/XXI/tex/components/18_expose_XXI_corollary43_through_lemma52_statement_en.tex
\begin{corollary}\label{sga3:XXI:4.3}
\footnote{Editors' note: The statement was corrected.}
Let \(x\in V^\ast\). For every \(w\in W\), put
\[
 x-w(x)=a_w\alpha^\ast+b_w\beta^\ast.
\]
If \((x,\alpha)\ge0\) and \((x,\beta)\ge0\), then
\[
 a_w+b_w\ge0.
\]
If moreover \((x,\alpha)>0\) (respectively \((x,\beta)>0\)), then
\[
 a_w+b_w>0
\]
for \(w\ne1,s_\beta\) (respectively for \(w\ne1,s_\alpha\)).
\end{corollary}

\begin{proof}
This follows at once from Proposition~\ref{sga3:XXI:4.1} and
Lemma~\ref{sga3:XXI:4.2}.
\end{proof}

\begin{corollary}\label{sga3:XXI:4.4}
Let \(\mathcal{R}\) be any root datum, and let \(\Delta\) be a system
of simple roots. Let \(\gamma\) be a positive root, let
\(\alpha,\beta\) be two simple roots, and let \(W_{\alpha,\beta}\) be
the subgroup of \(W\) generated by \(s_\alpha\) and \(s_\beta\). If
\[
 \operatorname{ord}_\Delta(s_\alpha(\gamma))
 <\operatorname{ord}_\Delta(\gamma),
 \qquad
 \operatorname{ord}_\Delta(s_\beta(\gamma))
 \le\operatorname{ord}_\Delta(\gamma),
\]
then, for every \(w\in W_{\alpha,\beta}\),
\[
 \operatorname{ord}_\Delta(w(\gamma))
 \le\operatorname{ord}_\Delta(\gamma).
\]
Moreover,
\[
 \operatorname{ord}_\Delta(w(\gamma))
 <\operatorname{ord}_\Delta(\gamma)
\]
if \(w\ne1,s_\beta\).
\end{corollary}

\begin{proof}
Consider the dual root datum \(\mathcal{R}^\ast\), and then the datum
\[
 (M^\ast,M,R'^\ast,R'),
\]
where \(R'\) is the set of roots that are rational linear
combinations of \(\alpha\) and \(\beta\). Applying
Corollary~\ref{sga3:XXI:4.3} to this datum gives the stated result.
\footnote{Editors' note: The hypotheses say that
\((\gamma,\alpha^\ast)>0\) and \((\gamma,\beta^\ast)\ge0\).
Consequently
\(\gamma-w(\gamma)\in\mathbb N\alpha+\mathbb N\beta\) for every
\(w\in W_{\alpha,\beta}\), and it is nonzero if
\(w\notin\{1,s_\beta\}\).}
\end{proof}

\section{The Weyl group: generators and relations}
\label{sga3:XXI:5}
\setcounter{theorem}{0}

Let \(\mathcal{R}\) be a root datum. Since the Weyl group is the same
for this datum and for the corresponding reduced datum, we may assume
\(\mathcal{R}\) reduced when studying the Weyl group.

Let
\[
 \Delta=\{\alpha_1,\ldots,\alpha_n\}
\]
be a system of simple roots, where \(n=\rgss(\mathcal{R})\). Let
\(n_{ij}\) be the order of \(s_{\alpha_i}s_{\alpha_j}\) in \(W\).
In particular \(n_{ii}=1\), and Remarks~\ref{sga3:XXI:3.4.2} and
\ref{sga3:XXI:3.4.3} show that the subgroup \(W_{ij}\) of \(W\)
generated by \(s_{\alpha_i}\) and \(s_{\alpha_j}\) is defined by
\[
 s_{\alpha_i}^{\,2}=s_{\alpha_j}^{\,2}
 =(s_{\alpha_i}s_{\alpha_j})^{n_{ij}}=1.
\]

\begin{theorem}\label{sga3:XXI:5.1}
The group \(W\) is the group generated by the elements
\[
 s_{\alpha_i},\qquad i=1,2,\ldots,n,
\]
subject to the relations
\[
 (s_{\alpha_i}s_{\alpha_j})^{n_{ij}}=1.
\]
\end{theorem}

\begin{proof}
We have already seen that the theorem is true when \(n=2\); we shall
use this fact in the proof. Introduce the group \(\widetilde W\)
generated by elements
\[
 T_i,\qquad i=1,2,\ldots,n,
\]
subject to the relations
\[
 (T_iT_j)^{n_{ij}}=1.
\]
In particular \(n_{ii}=1\), whence \(T_i^2=1\). Let
\[
 p:\widetilde W\longrightarrow W
\]
be the group homomorphism sending \(T_i\) to \(s_{\alpha_i}\). We know
that \(p\) is surjective; we shall prove that it is injective.
\end{proof}

\begin{lemma}\label{sga3:XXI:5.2}
For every \(\alpha\in\mathscr{P}(\Delta)\), one can define uniquely an
element \(T_\alpha\in\widetilde W\) with the following properties:
\begin{enumerate}
\item[(i)] \(p(T_\alpha)=s_\alpha\);
\item[(ii)] \(T_{\alpha_i}=T_i\);
\item[(iii)] if \(\beta\) and \(\alpha\) are positive roots such that
\(s_{\alpha_i}(\alpha)=\beta\), then
\[
 T_iT_\alpha T_i=T_\beta.
\]
\end{enumerate}
\end{lemma}

% END INCLUDED SGA3_English_current_buildable_source_20260823/source/inputs/XXI/tex/components/18_expose_XXI_corollary43_through_lemma52_statement_en.tex

% BEGIN INCLUDED SGA3_English_current_buildable_source_20260823/source/inputs/XXI/tex/components/19_expose_XXI_lemma52_proof_partA_en.tex
\begin{proof}
First observe that it follows from
\ref{sga3:XXI:1.2.10} and \ref{sga3:XXI:3.3.6} that (i) is a
consequence of (ii) and \SGARefLink{sga3:XXI:3.6:item:iii}{(iii)}, and that (ii) and \SGARefLink{sga3:XXI:3.6:item:iii}{(iii)} determine the
\(T_\alpha\) completely. We construct them by induction on
\(\operatorname{ord}_\Delta(\alpha)\). If
\(\operatorname{ord}_\Delta(\alpha)=1\), then \(\alpha\in\Delta\),
and we put \(T_\alpha=T_i\) if \(\alpha=\alpha_i\).

Consider the hypothesis
\[
\begin{split}
(\mathrm{H}_p)\quad&
\text{there exist \(T_\alpha\), for
\(\alpha\in\mathscr{P}(\Delta)\) with
\(\operatorname{ord}_\Delta(\alpha)\le p\), satisfying (ii),}\\
&\text{and satisfying (iii) whenever
\(\operatorname{ord}_\Delta(\alpha)\le p\) and
\(\operatorname{ord}_\Delta(\beta)\le p\).}
\end{split}
\]
It holds for \(p=1\). Indeed, if \(\alpha\) and
\(s_{\alpha_i}(\alpha)=\beta\) are simple, then \(\alpha_i\) and
\(\alpha\) are orthogonal. Thus, if
\(\alpha_j=\alpha=\beta\), one has \(n_{ij}=2\), whence
\[
 T_iT_jT_i=T_j.
\]

Suppose \(p>1\) and \((\mathrm{H}_{p-1})\) holds.

\medskip
\noindent\textbf{A) Construction of the \(T_\alpha\) for
\(\operatorname{ord}_\Delta(\alpha)\le p\).}
It is enough to do this when
\(\operatorname{ord}_\Delta(\alpha)=p\). There exists
\(\alpha_i\in\Delta\) such that
\[
 s_{\alpha_i}(\alpha)\in\mathscr{P}(\Delta),
 \qquad
 \operatorname{ord}_\Delta(s_{\alpha_i}(\alpha))<p
\]
by Lemma~\ref{sga3:XXI:3.3.3}. Put
\begin{equation}
\tag{\(\star\)}
\label{sga3:XXI:5.2:construction}
 T_\alpha=T_iT_{s_{\alpha_i}(\alpha)}T_i.
\end{equation}
We verify that \(T_\alpha\) depends only on \(\alpha\). Let
\(\alpha_j\in\Delta\) also satisfy
\[
 s_{\alpha_j}(\alpha)\in\mathscr{P}(\Delta),
 \qquad
 \operatorname{ord}_\Delta(s_{\alpha_j}(\alpha))<p.
\]
We must prove
\begin{equation}
\tag{\(+\)}
\label{sga3:XXI:5.2:independence}
 T_iT_{s_{\alpha_i}(\alpha)}T_i
 =T_jT_{s_{\alpha_j}(\alpha)}T_j.
\end{equation}
We distinguish two cases.

\begin{enumerate}
\item[(1)] Suppose that \(\alpha\) is a linear combination of
\(\alpha_i\) and \(\alpha_j\). The same is then true of
\(s_{\alpha_i}(\alpha)\) and \(s_{\alpha_j}(\alpha)\); by
\((\mathrm{H}_{p-1})\), the elements
\(T_{s_{\alpha_i}(\alpha)}\) and
\(T_{s_{\alpha_j}(\alpha)}\) are words in \(T_i,T_j\). Since the
projection of \eqref{sga3:XXI:5.2:independence} to \(W\) holds and
the theorem is true for \(n=2\), the restriction of \(p\) to the
subgroup of \(\widetilde W\) generated by \(T_i,T_j\) is injective.
Thus \eqref{sga3:XXI:5.2:independence} holds.

\item[(2)] Suppose that \(\alpha\) is not a linear combination of
\(\alpha_i\) and \(\alpha_j\). Then, for
\(w\in W_{\alpha_i,\alpha_j}\), all the \(w(\alpha)\) are positive
(cf. Lemma~\ref{sga3:XXI:3.4.9}). The relation to be verified can
also be written
\begin{equation}
\tag{\({+}{+}\)}
\label{sga3:XXI:5.2:rank-two-relation}
 (T_iT_j)^{n_{ij}-1}T_{s_{\alpha_i}(\alpha)}
 (T_jT_i)^{n_{ij}-1}
 =T_{s_{\alpha_j}(\alpha)}.
\end{equation}
Corollary~\ref{sga3:XXI:4.4} shows that all the \(w(\alpha)\) have
order \(<p\) for
\(w\in W_{\alpha_i,\alpha_j}\), \(w\ne1\). We may therefore apply
\((\mathrm{H}_{p-1})\) \(2(n_{ij}-1)\) times, which proves the
relation.
\end{enumerate}

\medskip
\noindent\textbf{B) Verification of \((\mathrm{H}_p)\).}
\footnote{Editors' note: The argument in the original has been
expanded in what follows.}
We must verify that, if \(\alpha_j\in\Delta\) and
\(\beta=s_{\alpha_j}(\alpha)\) satisfies
\(\operatorname{ord}_\Delta(\beta)\le p\), then
\[
 T_jT_\alpha T_j=T_\beta.
\]
If \(\operatorname{ord}_\Delta(\beta)<p\), this follows from the
preceding construction, since \(T_j^2=1\). We may therefore suppose
\[
 \operatorname{ord}_\Delta(\beta)
 =p=\operatorname{ord}_\Delta(\alpha).
\]
In this case \(\alpha\) and \(\alpha_j\) are orthogonal, so
\(\beta=\alpha\), and it remains to prove
\begin{equation}
\tag{\(\dagger\)}
\label{sga3:XXI:5.2:orthogonal}
 T_jT_\alpha T_j=T_\alpha.
\end{equation}
By \eqref{sga3:XXI:5.2:construction},
\[
 T_\alpha=T_iT_{s_{\alpha_i}(\alpha)}T_i,
\]
so it remains to verify
\begin{equation}
\tag{\({+}{+}{+}\)}
\label{sga3:XXI:5.2:orthogonal-expanded}
 T_jT_iT_{s_{\alpha_i}(\alpha)}T_iT_j
 =T_\alpha
 =T_iT_{s_{\alpha_i}(\alpha)}T_i.
\end{equation}
\end{proof}

% END INCLUDED SGA3_English_current_buildable_source_20260823/source/inputs/XXI/tex/components/19_expose_XXI_lemma52_proof_partA_en.tex

% BEGIN INCLUDED SGA3_English_current_buildable_source_20260823/source/inputs/XXI/tex/components/20_expose_XXI_lemma52_partB_through_theorem56_statement_en.tex
\begin{proof}[Proof of Lemma~\ref{sga3:XXI:5.2}, continued]
Put \(m=n_{ij}\), \(s_i=s_{\alpha_i}\), and
\(s_j=s_{\alpha_j}\). We have
\[
 T_jT_i=(T_iT_j)^{m-1}.
\]
By Corollary~\ref{sga3:XXI:4.4},
\(\operatorname{ord}_\Delta(w(\alpha))<p\) for every
\(w\in W_{ij}\) distinct from
\[
 1=(s_is_j)^m
 \quad\text{and}\quad
 s_j=s_i(s_is_j)^{m-1}.
\]
The induction hypothesis therefore gives
\[
\begin{split}
 &T_j(T_iT_j)^{m-2}T_{s_i(\alpha)}
       (T_jT_i)^{m-2}T_j\\
 &\qquad
 =T_{s_j(s_is_j)^{m-2}s_i(\alpha)}
 =T_{s_is_j(\alpha)}
 =T_{s_i(\alpha)}.
\end{split}
\]
The last equality holds because \(s_j(\alpha)=\alpha\). Hence
\[
 (T_iT_j)^{m-1}T_{s_i(\alpha)}(T_jT_i)^{m-1}
 =T_iT_{s_i(\alpha)}T_i,
\]
which proves \eqref{sga3:XXI:5.2:orthogonal-expanded}.
\end{proof}

\begin{lemma}\label{sga3:XXI:5.3}
Let \(h\in\widetilde W\). Write
\[
 h=T_{\alpha_1}\cdots T_{\alpha_m},
\]
where the \(\alpha_i\in\Delta\), not necessarily distinct, are chosen
so that \(m\) is minimal. Then
\[
 p(h)(\alpha_m)\in-\mathscr{P}(\Delta).
\]
\end{lemma}

\begin{proof}
Since
\[
 p(T_{\alpha_m})(\alpha_m)
 =s_{\alpha_m}(\alpha_m)=-\alpha_m,
\]
if \(p(h)(\alpha_m)\) were positive, there would be an index
\(k\), \(1\le k\le m-1\), such that
\[
 u=s_{\alpha_{k+1}}\cdots s_{\alpha_m}(\alpha_m)
 =-s_{\alpha_{k+1}}\cdots s_{\alpha_{m-1}}(\alpha_m)
 \in-\mathscr{P}(\Delta)
\]
and \(s_{\alpha_k}(u)\in\mathscr{P}(\Delta)\). Necessarily
\(u=-\alpha_k\), by Lemma~\ref{sga3:XXI:3.3.1}, whence
\[
 s_{\alpha_{k+1}}\cdots s_{\alpha_{m-1}}(\alpha_m)=\alpha_k.
\]
Property \SGARefLink{sga3:XXI:3.6:item:iii}{(iii)} of Lemma~\ref{sga3:XXI:5.2} then gives
\[
 T_{\alpha_k}T_{\alpha_{k+1}}\cdots
 T_{\alpha_{m-1}}T_{\alpha_m}
 =
 T_{\alpha_{k+1}}\cdots T_{\alpha_{m-1}},
\]
contradicting the minimality of \(m\).
\end{proof}

Let now \(h\in\widetilde W\) satisfy
\[
 p(h)\bigl(\mathscr{P}(\Delta)\bigr)
 \subset\mathscr{P}(\Delta).
\]
Lemma~\ref{sga3:XXI:5.3} gives \(p(h)=1\). This proves
Theorem~\ref{sga3:XXI:5.1} and also the following two corollaries.

\begin{corollary}\label{sga3:XXI:5.4}
If \(R_+\) is a system of positive roots and \(w\in W\) satisfies
\(w(R_+)=R_+\), then \(w=1\).
\end{corollary}

\begin{corollary}\label{sga3:XXI:5.5}
The Weyl group acts freely and transitively on the set of systems of
positive roots, respectively on the set of systems of simple roots,
respectively on the set of Weyl chambers.
\end{corollary}

Choose a system of simple roots \(\Delta\), and put
\[
 R^+=\mathscr{P}(\Delta).
\]
\footnote{Editors' note: The superscript notation \(R^+\), rather
than \(R_+\), is used here so that
\(R^+_{\alpha,\beta}=R_{\alpha,\beta}\cap R^+\).}
For every pair of simple roots
\((\alpha,\beta)\in\Delta\times\Delta\), let
\(R_{\alpha,\beta}\) be the set of roots that are linear combinations
of \(\alpha\) and \(\beta\). Put
\[
 R^+_{\alpha,\beta}=R^+\cap R_{\alpha,\beta},
\]
and let \(W_{\alpha,\beta}\) be the Weyl group of
\(R_{\alpha,\beta}\), that is, the subgroup of \(W\) generated by
\(s_\alpha\) and \(s_\beta\).

\begin{theorem}[Tits]\label{sga3:XXI:5.6}
Let \(\alpha\) and \(\beta\) be two simple roots, and let \(w\in W\)
satisfy \(w(\alpha)=\beta\). There exist a sequence of simple roots
\(\alpha_0,\ldots,\alpha_m\) and a sequence
\(w_0,\ldots,w_{m-1}\) of elements of \(W\) satisfying:
\begin{enumerate}
\item[(i)] \(\alpha_0=\alpha\) and \(\alpha_m=\beta\);
\item[(ii)] \(w=w_{m-1}w_{m-2}\cdots w_0\);
\item[(iii)] \(w_i(\alpha_i)=\alpha_{i+1}\) for
  \(0\le i\le m-1\);
\item[(iv)] for every \(i\), \(0\le i\le m-1\), such that
  \(\alpha_i\ne\alpha_{i+1}\), one has
  \(w_i\in W_{\alpha_i,\alpha_{i+1}}\);
\item[(v)] for every \(i\), \(0\le i\le m-1\), such that
  \(\alpha_i=\alpha_{i+1}\), there exists a simple root \(\beta_i\)
  such that \(w_i\in W_{\alpha_i,\beta_i}\).
\end{enumerate}
\end{theorem}

% END INCLUDED SGA3_English_current_buildable_source_20260823/source/inputs/XXI/tex/components/20_expose_XXI_lemma52_partB_through_theorem56_statement_en.tex

% BEGIN INCLUDED SGA3_English_current_buildable_source_20260823/source/inputs/XXI/tex/components/21_expose_XXI_theorem56_proof_en.tex
\begin{proof}
Put\footnote{Editors' note: In what follows, the original argument has
been expanded, and equality \((1)\) has been corrected.}
\[
 M(w)
 =\operatorname{Card}\bigl(R^+\cap w^{-1}(-R^+)\bigr)
 =\operatorname{Card}\{\gamma\in R^+\mid w(\gamma)\in-R^+\}.
\]
If \(M(w)=0\), then \(w(R^+)=R^+\), so
\(w=1\) by Corollary~\ref{sga3:XXI:5.4}; the theorem is then
trivial: \(m=0\), and assertions (iii)--(v) are empty.

We argue by induction on \(M(w)\). Suppose \(M(w)>0\). There exists
\(\beta_0\in\Delta\) such that \(w(\beta_0)\in-R^+\). Put
\(\alpha_0=\alpha\), and consider
\[
 A=w^{-1}(R^+)\cap R_{\alpha_0,\beta_0}.
\]
This is a system of positive roots of
\(R_{\alpha_0,\beta_0}\). Hence there exists
\(w_0\in W_{\alpha_0,\beta_0}\) such that
\[
 w_0^{-1}\bigl(R^+_{\alpha_0,\beta_0}\bigr)=A.
\]
Put \(w'=ww_0^{-1}\). Lemma~\ref{sga3:XXI:3.4.9} gives
\[
 R^+-R^+_{\alpha_0,\beta_0}
 =
 w_0\bigl(R^+-R^+_{\alpha_0,\beta_0}\bigr),
\]
and therefore
\begin{equation}
\tag{1}
\label{sga3:XXI:5.6:outside-rank-two}
\begin{split}
 &(R^+-R^+_{\alpha_0,\beta_0})\cap w^{-1}(-R^+)\\
 &\qquad=
 w_0\Bigl((R^+-R^+_{\alpha_0,\beta_0})
             \cap w'^{-1}(-R^+)\Bigr).
\end{split}
\end{equation}
On the other hand,
\begin{equation}
\tag{2}
\label{sga3:XXI:5.6:beta-negative}
 \beta_0\in
 R^+_{\alpha_0,\beta_0}\cap w^{-1}(-R^+).
\end{equation}
Since
\(w_0(R_{\alpha_0,\beta_0})=R_{\alpha_0,\beta_0}\), one has
\[
 w_0(-A)
 =R_{\alpha_0,\beta_0}\cap w'^{-1}(-R^+),
\]
and hence
\begin{equation}
\tag{\(2'\)}
\label{sga3:XXI:5.6:rank-two-empty}
\begin{split}
 R^+_{\alpha_0,\beta_0}\cap w'^{-1}(-R^+)
 &=
 R^+_{\alpha_0,\beta_0}\cap w_0(-A)\\
 &=
 R^+_{\alpha_0,\beta_0}\cap(-R^+_{\alpha_0,\beta_0})
 =\varnothing.
\end{split}
\end{equation}
It follows from \eqref{sga3:XXI:5.6:outside-rank-two},
\eqref{sga3:XXI:5.6:beta-negative}, and
\eqref{sga3:XXI:5.6:rank-two-empty} that
\[
 M(w')<M(w).
\]

Put \(\alpha_1=w_0(\alpha_0)\). We show that
\(\alpha_1\in\Delta\), or equivalently that
\(\alpha_0\in w_0^{-1}(\Delta)\). Since
\(w(\alpha_0)\in\Delta\), we have
\[
 \alpha_0\in
 w^{-1}(\Delta)\cap R_{\alpha_0,\beta_0}.
\]
Thus \(\alpha_0\) is a simple root of
\[
 A
 =w^{-1}(R^+)\cap R_{\alpha_0,\beta_0}
 =w_0^{-1}(R^+_{\alpha_0,\beta_0}),
\]
and consequently belongs to
\[
 w_0^{-1}\bigl(\Delta\cap R^+_{\alpha_0,\beta_0}\bigr)
 =
 w_0^{-1}\{\alpha_0,\beta_0\}
\]
by Lemma~\ref{sga3:XXI:3.4.8}. Hence
\(\alpha_1=w_0(\alpha_0)\) equals \(\alpha_0\) or \(\beta_0\).
If \(\alpha_1\ne\alpha_0\), then
\(\alpha_1=\beta_0\) and
\(w_0\in W_{\alpha_0,\alpha_1}\).

Finally,
\[
 \beta=w'(\alpha_1)
 \qquad\text{and}\qquad
 M(w')<M(w),
\]
so the induction hypothesis completes the proof.
\end{proof}

% END INCLUDED SGA3_English_current_buildable_source_20260823/source/inputs/XXI/tex/components/21_expose_XXI_theorem56_proof_en.tex

% BEGIN INCLUDED SGA3_English_current_buildable_source_20260823/source/inputs/XXI/tex/components/22_expose_XXI_section61_morphisms_en.tex
\section{Morphisms of root data}\label{sga3:XXI:6}
\renewcommand{\thetheorem}{\thesubsection.\arabic{theorem}}

\subsection{Definition}\label{sga3:XXI:6.1}

Let
\[
 \cR=(M,M^\ast,R,R^\ast),
 \qquad
 \cR'=(M',M'^\ast,R',R'^\ast)
\]
be two root data. Let \(f:M'\to M\) be a linear map and
\({}^tf:M^\ast\to M'^\ast\) its transpose.

\begin{definition}\label{sga3:XXI:6.1.1}
The map \(f\) is called a \emph{morphism} from \(\cR'\) to \(\cR\),
and is denoted
\[
 f:\cR'\longrightarrow\cR,
\]
if \(f\) induces a bijection from \(R'\) onto \(R\), and
\({}^tf\) induces a bijection from \(R^\ast\) onto \(R'^\ast\).
\end{definition}

Then \({}^tf\) is a morphism of the dual root data:
\[
 {}^tf:\cR^\ast\longrightarrow\cR'^\ast.
\]
If \(f:\cR'\to\cR\) is a morphism and
\(\alpha=f(\alpha')\) for \(\alpha'\in R'\), then
\[
 \alpha'^\ast={}^tf(\alpha^\ast).
\]
Indeed, if \(p\) and \(p'\) denote the maps of
\ref{sga3:XXI:1.2.1} for \(\cR\) and \(\cR'\), respectively, then
\[
 p'={}^tf\circ p\circ f,
\]
which gives the assertion at once. The reader may verify the
statements below, almost all of which are immediate.

\begin{proposition}\label{sga3:XXI:6.1.2}
Let \(f:\cR'\to\cR\) be a morphism of root data. If
\(\alpha'\in R'\) and \(\alpha=f(\alpha')\), then
\[
 \alpha'^\ast={}^tf(\alpha^\ast).
\]
Moreover, \(f\) induces isomorphisms
\[
 R'\xrightarrow{\sim}R,\qquad
 \Gamma_0(R')\xrightarrow{\sim}\Gamma_0(R),\qquad
 \VR(R')\xrightarrow{\sim}\VR(R),
\]
and \({}^tf\) induces isomorphisms
\[
 R^\ast\xrightarrow{\sim}R'^\ast,\qquad
 \Gamma_0(R^\ast)\xrightarrow{\sim}\Gamma_0(R'^\ast),\qquad
 \VR(R^\ast)\xrightarrow{\sim}\VR(R'^\ast).
\]
The last is the transpose of the corresponding morphism induced by
\(f\). The map
\[
 s_{\alpha'}\longmapsto s_{f(\alpha')}
\]
extends to an isomorphism
\[
 W(\cR')\xrightarrow{\sim}W(\cR)
\]
compatible with the actions of these two groups on the sets of
\ref{sga3:XXI:1.1.13}.
\end{proposition}

\begin{proposition}\label{sga3:XXI:6.1.3}
The maps
\[
 \Delta'\longmapsto f(\Delta'),\qquad
 R'_+\longmapsto f(R'_+),\qquad
 C'\longmapsto({}^tf\otimes\mathbb R)^{-1}(C')
\]
define bijective correspondences between the systems of simple roots,
the systems of positive roots, and the Weyl chambers for \(\cR'\) and
\(\cR\). These correspondences are compatible with the actions of the
Weyl groups and with the correspondences
\[
 \mathscr{S}(R_+)\longleftrightarrow R_+
 \longleftrightarrow\mathscr{C}(R_+).
\]
\end{proposition}

\begin{lemma}\label{sga3:XXI:6.1.4}
Morphisms compose. A morphism \(f:\cR'\to\cR\) is an isomorphism if
and only if the linear map \(f:M'\to M\) is bijective.
\end{lemma}

% END INCLUDED SGA3_English_current_buildable_source_20260823/source/inputs/XXI/tex/components/22_expose_XXI_section61_morphisms_en.tex

% BEGIN INCLUDED SGA3_English_current_buildable_source_20260823/source/inputs/XXI/tex/components/23_expose_XXI_section62_isogenies_en.tex
\subsection{Isogenies}\label{sga3:XXI:6.2}

\begin{definition}\label{sga3:XXI:6.2.1}
A morphism \(f:\cR'\to\cR\) of root data is called an
\emph{isogeny} if \(f:M'\to M\) is injective with finite cokernel.
\end{definition}

If \(f\) is an isogeny, then \({}^tf\) is an isogeny.

\begin{definition}\label{sga3:XXI:6.2.2}
Let \(f:\cR'\to\cR\) be an isogeny. Put
\[
 K(f)=\operatorname{Coker}\bigl(M'\xrightarrow{f}M\bigr).
\]
\end{definition}

\begin{lemma}\label{sga3:XXI:6.2.3}
There is a natural pairing
\[
 K(f)\times K({}^tf)\longrightarrow\mathbb Q/\mathbb Z
\]
that puts these two finite groups in duality.
\end{lemma}

\begin{proof}
This is classical.
\end{proof}

\begin{lemma}\label{sga3:XXI:6.2.4}
If \(f:\cR'\to\cR\) is a morphism, then
\[
 \rgss(\cR')=\rgss(\cR).
\]
If \(f\) is moreover an isogeny, then
\[
 \rgred(\cR')=\rgred(\cR).
\]
\end{lemma}

\begin{proof}
This is immediate.
\end{proof}

\begin{lemma}\label{sga3:XXI:6.2.5}
Every morphism of semisimple root data is an isogeny.
\end{lemma}

\begin{proof}
The map \(f\) must induce an isomorphism from
\(V'=\VR(R')\) onto \(V=\VR(R)\).
\end{proof}

If \(\cR'\) and \(\cR\) are semisimple, every isogeny
\(f:\cR'\to\cR\) defines a commutative diagram:
\begin{center}
\begin{tikzcd}[column sep=huge,row sep=large]
\Gamma_0(R') \arrow[r,"\sim"] \arrow[d]
  & \Gamma_0(R) \arrow[d] \\
M' \arrow[r,"f"'] & M
\end{tikzcd}
\end{center}
If \(M=\Gamma_0(R)\), then \(f\) is necessarily an isomorphism.

\begin{definition}\label{sga3:XXI:6.2.6}
A root datum is called \emph{adjoint}, respectively
\emph{simply connected}, if
\[
 M=\Gamma_0(R),
 \qquad\text{respectively}\qquad
 M^\ast=\Gamma_0(R^\ast).
\]
\end{definition}

An adjoint or simply connected root datum is therefore semisimple.
Moreover, \(\cR\) is adjoint, respectively simply connected, if and
only if \(\cR^\ast\) is simply connected, respectively adjoint. The
preceding result gives the following proposition.

\begin{proposition}\label{sga3:XXI:6.2.7}
Let \(\cR\) be a semisimple root datum. The following conditions are
equivalent:
\begin{enumerate}
\item[(i)] \(\cR\) is adjoint, respectively simply connected;
\item[(ii)] every isogeny \(\cR'\to\cR\), respectively
  \(\cR\to\cR'\), is an isomorphism.
\end{enumerate}
\end{proposition}

\begin{proposition}\label{sga3:XXI:6.2.8}
Let \(\cR\) be an adjoint, respectively simply connected, root datum.
Every indivisible root, respectively coroot, is an indivisible element
of \(M\), respectively of \(M^\ast\).
\end{proposition}

\begin{proof}
Every indivisible root belongs to a basis of \(\Gamma_0(R)\), by
\ref{sga3:XXI:3.3.5}.
\end{proof}

% END INCLUDED SGA3_English_current_buildable_source_20260823/source/inputs/XXI/tex/components/23_expose_XXI_section62_isogenies_en.tex

% BEGIN INCLUDED SGA3_English_current_buildable_source_20260823/source/inputs/XXI/tex/components/24_expose_XXI_section63_radical_coradical_en.tex
\subsection{Radical and coradical}\label{sga3:XXI:6.3}

Let \(\cR\) be a root datum. Put
\[
 N=\{x\in M\mid(\alpha^\ast,x)=0
       \text{ for every }\alpha^\ast\in R^\ast\},
\qquad
 N^\ast=M^\ast/\bigl(\VR(R^\ast)\cap M^\ast\bigr).
\]

\begin{lemma}\label{sga3:XXI:6.3.1}
Consider the canonical morphisms
\[
 N\longrightarrow M,
 \qquad
 M^\ast\longrightarrow N^\ast.
\]
They are transposes of one another, and \(N^\ast\) is identified with
the dual of \(N\).
\end{lemma}

\begin{proof}
This follows immediately from \ref{sga3:XXI:1.2.5}.
\end{proof}

\begin{definition}\label{sga3:XXI:6.3.2}
The \emph{coradical} of \(\cR\), denoted
\(\operatorname{corad}(\cR)\), is the trivial root datum
\[
 \operatorname{corad}(\cR)=(N,N^\ast,\varnothing,\varnothing).
\]
\end{definition}

If
\[
 \cR^0=(M,M^\ast,\varnothing,\varnothing),
\]
which is a trivial root datum, there is therefore a morphism
\[
 \operatorname{corad}(\cR)\longrightarrow\cR^0.
\]

\begin{definition}\label{sga3:XXI:6.3.3}
The \emph{radical} of \(\cR\), denoted
\(\operatorname{rad}(\cR)\), is the trivial root datum
\footnote{Editors' note: Equivalently, if
\[
 P=\{y\in M^\ast\mid(y,\alpha)=0\text{ for every }\alpha\in R\},
 \qquad
 P^\ast=M/\bigl(\VR(R)\cap M\bigr),
\]
then
\(\operatorname{rad}(\cR)=(P^\ast,P,\varnothing,\varnothing)\).}
\[
 \operatorname{rad}(\cR)=\operatorname{corad}(\cR^\ast)^\ast.
\]
\end{definition}

Thus there is a diagram
\begin{center}
\begin{tikzcd}[column sep=huge,row sep=large]
\operatorname{corad}(\cR) \arrow[rr] \arrow[dr]
  && \operatorname{rad}(\cR) \\
& \cR^0 \arrow[ur] &
\end{tikzcd}
\end{center}
whose transpose is the corresponding diagram for \(\cR^\ast\).

\begin{lemma}\label{sga3:XXI:6.3.4}
The canonical morphism
\[
 u:\operatorname{corad}(\cR)\longrightarrow\operatorname{rad}(\cR)
\]
is an isogeny.
\end{lemma}

\begin{definition}\label{sga3:XXI:6.3.5}
Put
\[
 N(\cR)=K(u)
 =
 M\Big/
 \left(
   \bigl(\VR(R)\cap M\bigr)
   +\bigcap_{\alpha\in R}\ker(\alpha^\ast)
 \right).
\]
There is then a canonical pairing
\[
 N(\cR)\times N(\cR^\ast)\longrightarrow\mathbb Q/\mathbb Z.
\]
\end{definition}

\begin{lemma}\label{sga3:XXI:6.3.6}
One has
\[
 \rgred\bigl(\operatorname{rad}(\cR)\bigr)
 =
 \rgred\bigl(\operatorname{corad}(\cR)\bigr)
 =
 \rgred(\cR)-\rgss(\cR),
\]
and the following conditions are equivalent:
\begin{enumerate}
\item[(i)] \(\cR\) is semisimple;
\item[(ii)] \(\operatorname{rad}(\cR)=0\);
\item[(iii)] \(\operatorname{corad}(\cR)=0\).
\end{enumerate}
\end{lemma}

% END INCLUDED SGA3_English_current_buildable_source_20260823/source/inputs/XXI/tex/components/24_expose_XXI_section63_radical_coradical_en.tex

% BEGIN INCLUDED SGA3_English_current_buildable_source_20260823/source/inputs/XXI/tex/components/25_expose_XXI_section64_products_en.tex
\subsection{Products of root data}\label{sga3:XXI:6.4}

\begin{definition}\label{sga3:XXI:6.4.1}
Let
\[
 \cR=(M,M^\ast,R,R^\ast),
 \qquad
 \cR'=(M',M'^\ast,R',R'^\ast)
\]
be two root data. The \emph{product root datum} of \(\cR\) and
\(\cR'\), denoted
\[
 \cR''=\cR\times\cR',
\]
is the root datum \((M'',M''^\ast,R'',R''^\ast)\), where
\[
\begin{aligned}
 M''&=M\times M',
 &M''^\ast&=M^\ast\times M'^\ast,\\
 R''&=(R\times0)\cup(0\times R'),
 &R''^\ast&=(R^\ast\times0)\cup(0\times R'^\ast),
\end{aligned}
\]
and the map \(\alpha\mapsto\alpha^\ast\) is the evident one.
\end{definition}

\begin{proposition}\label{sga3:XXI:6.4.2}
Under the preceding hypotheses there are canonical isomorphisms
\[
 \Gamma_0(R'')\simeq\Gamma_0(R)\times\Gamma_0(R'),
 \qquad
 \VR(R'')\simeq\VR(R)\times\VR(R'),
\]
\[
 W(\cR'')\simeq W(\cR)\times W(\cR'),
\]
and similarly for the other associated constructions. Moreover,
\[
 \rgred(\cR'')=\rgred(\cR)+\rgred(\cR'),
 \qquad
 \rgss(\cR'')=\rgss(\cR)+\rgss(\cR').
\]
There is also a canonical isomorphism of root data
\[
 (\cR\times\cR')^\ast\simeq\cR^\ast\times\cR'^\ast.
\]
\end{proposition}

The preceding definitions extend immediately to products of any finite
number of factors.

\begin{proposition}\label{sga3:XXI:6.4.3}
Let
\[
 \cR=\cR_1\times\cR_2\times\cdots\times\cR_n
\]
be a product of root data. The following conditions are equivalent:
\begin{enumerate}
\item[(i)] \(\cR\) is semisimple, respectively simply connected,
  respectively adjoint, respectively reduced;
\item[(ii)] every \(\cR_i\) is semisimple, respectively simply
  connected, respectively adjoint, respectively reduced.
\end{enumerate}
\end{proposition}

Consider the following special case. Let \(\cR_0\) be a trivial root
datum and \(\cR_1\) a semisimple root datum. There is then a commutative
diagram
\begin{center}
\begin{tikzcd}[column sep=large,row sep=large]
& \cR_0\times\cR_1 \arrow[dr] & \\
\cR_1 \arrow[ur] \arrow[rr,"\operatorname{id}"]
  && \cR_1
\end{tikzcd}
\end{center}

\begin{lemma}\label{sga3:XXI:6.4.4}
There are canonical isomorphisms
\begin{center}
\begin{tikzcd}[column sep=huge,row sep=large]
\operatorname{corad}(\cR_0\times\cR_1)
  \arrow[rr,"\sim"] \arrow[dr,"\sim"']
  && \operatorname{rad}(\cR_0\times\cR_1) \\
& \cR_0 \arrow[ur,"\sim"'] &
\end{tikzcd}
\end{center}
In particular,
\[
 N(\cR_0\times\cR_1)=0.
\]
\end{lemma}

We shall see later that, conversely, if \(N(\cR)=0\), then \(\cR\) is
the product of a semisimple root datum and a trivial root datum.

% END INCLUDED SGA3_English_current_buildable_source_20260823/source/inputs/XXI/tex/components/25_expose_XXI_section64_products_en.tex

% BEGIN INCLUDED SGA3_English_current_buildable_source_20260823/source/inputs/XXI/tex/components/26_expose_XXI_section65_induced_through_lemma654_en.tex
\subsection{Induced and coinduced root data}\label{sga3:XXI:6.5}

Let \(\cR=(M,M^\ast,R,R^\ast)\) be a root datum, and let
\(N\subset M\) be a subgroup containing the roots, that is,
\[
 \Gamma_0(R)\subset N\subset M.
\]
The canonical linear map \(i_N:N\to M\) gives, by transposition, a
linear map
\[
 {}^ti_N:M^\ast\longrightarrow N^\ast.
\]
Put
\[
 R_N=R,
 \qquad
 R_N^\ast={}^ti_N(R^\ast).
\]

\begin{lemma}\label{sga3:XXI:6.5.1}
The quadruple
\[
 \cR_N=(N,N^\ast,R_N,R_N^\ast)
\]
is a root datum, and \(i_N\) is a morphism.
\end{lemma}

\begin{proof}
First, \({}^ti_N\) induces an isomorphism from \(R^\ast\) onto
\(R_N^\ast\). Indeed, if \(\alpha,\beta\in R\) and
\[
 {}^ti_N(\alpha^\ast)={}^ti_N(\beta^\ast),
\]
then
\[
 (\alpha^\ast,x)=(\beta^\ast,x)
\]
for every \(x\in N\), in particular for \(x\in R\). Hence
\(\alpha=\beta\) by \ref{sga3:XXI:1.1.4}. The rest follows without
difficulty.
\end{proof}

\begin{definition}\label{sga3:XXI:6.5.2}
The root datum \(\cR_N\) is called the root datum
\emph{induced by \(\cR\) on \(N\)}.
\end{definition}

\begin{lemma}\label{sga3:XXI:6.5.3}
Let \(f:\cR'\to\cR\) be a morphism, and put
\[
 N=f(M')\subset M.
\]
Then \(f\) factors uniquely through \(i_N\).
\end{lemma}

In particular, up to isomorphism, the isogenies
\(\cR'\to\cR\) correspond bijectively to the finite-index subgroups of
\[
 M/\Gamma_0(R).
\]
This sharpens \ref{sga3:XXI:6.2.7}.

Now let \(N^\ast\subset M^\ast\) be a subgroup containing \(R^\ast\).
Define the root datum \emph{coinduced by \(\cR\) on \(N^\ast\)} by
\[
 \cR^{N^\ast}=(\cR^\ast_{N^\ast})^\ast.
\]
There is a canonical morphism
\[
 p^{N^\ast}:\cR\longrightarrow\cR^{N^\ast}.
\]

\begin{lemma}\label{sga3:XXI:6.5.4}
Let \(f:\cR'\to\cR\) be a morphism. There exist subgroups
\(N\subset M\) and \(N'^\ast\subset M'^\ast\) such that \(f\) factors
as
\begin{center}
\begin{tikzcd}[column sep=huge,row sep=large]
\cR' \arrow[r,"f"] \arrow[d,"p^{N'^\ast}"']
  & \cR \\
\cR'^{\,N'^\ast} \arrow[r,"f_0","\sim"']
  & \cR_N \arrow[u,"i_N"']
\end{tikzcd}
\end{center}
where \(f_0\) is an isomorphism.
\end{lemma}

\begin{proof}
Take \(N=f(M')\), as in \ref{sga3:XXI:6.5.3}. The resulting morphism
\(M'\to N\) is surjective, so its transpose is injective. Take its
image for \(N'^\ast\).
\end{proof}

% END INCLUDED SGA3_English_current_buildable_source_20260823/source/inputs/XXI/tex/components/26_expose_XXI_section65_induced_through_lemma654_en.tex

% BEGIN INCLUDED SGA3_English_current_buildable_source_20260823/source/inputs/XXI/tex/components/27_expose_XXI_section65_special_cases_en.tex
We now treat some special cases. If \(N=\Gamma_0(R)\), write
\[
 \cR_N=\operatorname{ad}(\cR).
\]
If \(N=\VR(R)\cap M\), write
\[
 \cR_N=\operatorname{ss}(\cR).
\]
Thus there is a diagram
\[
 \operatorname{ad}(\cR)\longrightarrow
 \operatorname{ss}(\cR)\longrightarrow\cR.
\]
Put
\[
 \operatorname{der}(\cR)=\operatorname{ss}(\cR^\ast)^\ast,
 \qquad
 \operatorname{sc}(\cR)=\operatorname{ad}(\cR^\ast)^\ast.
\]
By duality there is a diagram
\[
 \cR\longrightarrow
 \operatorname{der}(\cR)\longrightarrow
 \operatorname{sc}(\cR).
\]

\begin{proposition}\label{sga3:XXI:6.5.5}
\begin{enumerate}
\item[(i)] In the first row of the diagram
\begin{center}
\begin{tikzcd}[column sep=large,row sep=large]
\operatorname{ad}(\cR) \arrow[r]
  & \operatorname{ss}(\cR) \arrow[rr] \arrow[dr]
  && \operatorname{der}(\cR) \arrow[r]
  & \operatorname{sc}(\cR) \\
&& \cR \arrow[ur] &&
\end{tikzcd}
\end{center}
the four root data are semisimple and the three morphisms are
isogenies.
\item[(ii)] The datum \(\operatorname{ad}(\cR)\) is adjoint, and
\(\cR\) is adjoint if and only if
\(\operatorname{ad}(\cR)\to\cR\) is an isomorphism.
\item[(iii)] The datum \(\operatorname{sc}(\cR)\) is simply connected,
and \(\cR\) is simply connected if and only if
\(\cR\to\operatorname{sc}(\cR)\) is an isomorphism.
\item[(iv)] The following conditions are equivalent:
\begin{enumerate}
\item[(a)] \(\cR\) is semisimple;
\item[(b)] \(\operatorname{ss}(\cR)\to\cR\) is an isomorphism;
\item[(c)] \(\cR\to\operatorname{der}(\cR)\) is an isomorphism.
\end{enumerate}
\end{enumerate}
\end{proposition}

Consider the morphism
\[
 \operatorname{ss}(\cR)\longrightarrow\operatorname{der}(\cR).
\]
The definitions immediately give the following result.

\begin{lemma}\label{sga3:XXI:6.5.6}
Let
\[
 h:\operatorname{ss}(\cR)\longrightarrow\operatorname{der}(\cR)
\]
be the canonical isogeny. Then
\[
 K(h)\simeq N(\cR).
\]
\end{lemma}

\SGARefTarget{sga3:XXI:6.5.7}
Consider other special cases of induced root data. Put
\[
 N=
 \{x\in M\mid(\alpha^\ast,x)=0\text{ for every }\alpha\in R\}
 \times\Gamma_0(R).
\]
The sum is direct by \ref{sga3:XXI:1.2.5}. It follows that
\(\cR_N\) is identified with
\[
 \operatorname{ad}(\cR)\times\operatorname{corad}(\cR).
\]
One may replace \(\Gamma_0(R)\) by \(\VR(R)\cap M\), and then dualize
both constructions. This gives the following diagram of root data:
\begin{center}
\resizebox{\textwidth}{!}{%
\begin{tikzpicture}[
  >=stealth,
  every node/.style={font=\small},
  x=1cm,y=1cm
]
\node (at) at (0,2.1) {$\operatorname{ad}(\cR)$};
\node (st) at (4,2.1) {$\operatorname{ss}(\cR)$};
\node (dt) at (10,2.1) {$\operatorname{der}(\cR)$};
\node (ct) at (14,2.1) {$\operatorname{sc}(\cR)$};

\node (am) at (0,0) {$\operatorname{ad}(\cR)\times\operatorname{corad}(\cR)$};
\node (sm) at (4,0) {$\operatorname{ss}(\cR)\times\operatorname{corad}(\cR)$};
\node (rm) at (7,0) {$\cR$};
\node (dm) at (10,0) {$\operatorname{der}(\cR)\times\operatorname{rad}(\cR)$};
\node (cm) at (14,0) {$\operatorname{sc}(\cR)\times\operatorname{rad}(\cR)$};

\node (ab) at (0,-2.1) {$\operatorname{ad}(\cR)$};
\node (sb) at (4,-2.1) {$\operatorname{ss}(\cR)$};
\node (db) at (10,-2.1) {$\operatorname{der}(\cR)$};
\node (cb) at (14,-2.1) {$\operatorname{sc}(\cR)$};

\draw[->] (at) -- (st);
\draw[->] (st) -- (dt);
\draw[->] (dt) -- (ct);
\draw[->] (am) -- (sm);
\draw[->] (sm) -- (rm);
\draw[->] (rm) -- (dm);
\draw[->] (dm) -- (cm);
\draw[->] (ab) -- (sb);
\draw[->] (sb) -- (db);
\draw[->] (db) -- (cb);

\draw[->] (at) -- (am);
\draw[->] (am) -- (ab);
\draw[->] (st) -- (sm);
\draw[->] (sm) -- (sb);
\draw[->] (dt) -- (dm);
\draw[->] (dm) -- (db);
\draw[->] (ct) -- (cm);
\draw[->] (cm) -- (cb);
\end{tikzpicture}
}
\end{center}
It is commutative, as is checked at once, and is self-dual in the
evident sense. The horizontal morphisms are isogenies, and the
composites of the vertical arrows are the identity.

\begin{lemma}\label{sga3:XXI:6.5.8}
Let \(h_1\) and \(h_2\) be the canonical isogenies
\[
 \operatorname{ss}(\cR)\times\operatorname{corad}(\cR)
 \xrightarrow{h_1}\cR
 \xrightarrow{h_2}
 \operatorname{der}(\cR)\times\operatorname{rad}(\cR).
\]
Then
\[
 K(h_1)\simeq K(h_2)\simeq N(\cR).
\]
\end{lemma}

\begin{proof}
This is immediate from the definitions.
\end{proof}

\begin{corollary}\label{sga3:XXI:6.5.9}
For a root datum \(\cR\), the following conditions are equivalent:
\begin{enumerate}
\item[(i)] \(N(\cR)=0\), that is,
  \(\operatorname{corad}(\cR)\to\operatorname{rad}(\cR)\) is an
  isomorphism;
\item[(ii)] the canonical map
  \(\operatorname{ss}(\cR)\to\operatorname{der}(\cR)\) is an
  isomorphism;
\item[(iii)] \(h_1\) is an isomorphism;
\item[(iv)] \(h_2\) is an isomorphism;
\item[(v)] \(\cR\) is the product of a semisimple root datum and a
  trivial root datum.
\end{enumerate}
\end{corollary}

\begin{corollary}\label{sga3:XXI:6.5.10}
For every root datum \(\cR\), there are isogenies
\[
 \operatorname{ad}(\cR)\times\cR_0
 \longrightarrow\cR\longrightarrow
 \operatorname{sc}(\cR)\times\cR_0,
\]
where \(\cR_0\) is the trivial root datum of rank
\(\rgred(\cR)-\rgss(\cR)\).
\end{corollary}

\begin{lemma}\label{sga3:XXI:6.5.11}
Let \(\cR\) be a root datum, let \(\Delta\) be a system of simple
roots, and let \(\Delta'\subset\Delta\). Consider the root datum
\[
 \cR_{\Delta'}=(M,M^\ast,R_{\Delta'},R^\ast_{\Delta'})
\]
of \ref{sga3:XXI:3.4.7}.
\begin{enumerate}
\item[(i)] If \(\cR\) is simply connected, then
  \(\operatorname{der}(\cR_{\Delta'})\) is simply connected.
\item[(ii)] If \(\cR\) is adjoint, then
  \(\operatorname{ss}(\cR_{\Delta'})\) is adjoint.
\end{enumerate}
\end{lemma}

\begin{proof}
The two assertions are equivalent by duality. The second reduces to
\[
 M\cap\VR(R_{\Delta'})=\Gamma_0(R_{\Delta'}).
\]
When \(M=\Gamma_0(R)\), both sides are the subgroup of \(M\) generated
by \(\Delta'\).
\end{proof}

% END INCLUDED SGA3_English_current_buildable_source_20260823/source/inputs/XXI/tex/components/27_expose_XXI_section65_special_cases_en.tex

% BEGIN INCLUDED SGA3_English_current_buildable_source_20260823/source/inputs/XXI/tex/components/28_expose_XXI_section66_weights_partA_en.tex
\subsection{Weights}\label{sga3:XXI:6.6}

\begin{definition}\label{sga3:XXI:6.6.1}
For a root datum \(\cR\), put
\footnote{Editors' note: The notation \(\Gamma(\cR)\) has been replaced
by \(\Lambda(\cR)\) to avoid any possible confusion with
\(\Gamma_0(R)\).}
\[
 \Lambda(\cR)=
 \{x\in\VR(R)\mid(\alpha^\ast,x)\in\mathbb Z
   \text{ for every }\alpha^\ast\in R^\ast\}.
\]
The elements of \(\Lambda(\cR)\) are called the \emph{weights} of
\(\cR\). The weights of \(\cR^\ast\) are called the \emph{coweights}
of \(\cR\).
\end{definition}

One has
\[
 \Gamma_0(R)\subset\Lambda(\cR),
\]
and \(\Lambda(\cR)\) is stable under \(W(\cR)\).

\begin{lemma}\label{sga3:XXI:6.6.2}
The bilinear map
\[
 V^\ast\times V\longrightarrow\mathbb Q
\]
induces a duality
\[
 \Gamma_0(R^\ast)\times\Lambda(\cR)\longrightarrow\mathbb Z.
\]
\end{lemma}

\begin{proof}
Immediate.
\end{proof}

\begin{corollary}\label{sga3:XXI:6.6.3}
Let
\[
 \Delta^\ast=(\alpha_1^\ast,\ldots,\alpha_n^\ast)
\]
be a system of simple coroots, and let \(p_i\in\VR(R)\),
\(i=1,\ldots,n\), be defined by
\[
 (\alpha_i^\ast,p_j)=\delta_{ij}.
\]
Thus
\[
 s_{\alpha_i}(p_i)=p_i-\alpha_i,
 \qquad
 s_{\alpha_i}(p_j)=p_j\quad(i\ne j).
\footnote{Caution: when the datum is not reduced, the \(\alpha_i\)
need not form a system of simple roots.}
\]
Then \(\Lambda(\cR)\) is the free abelian group generated by the
\(p_i\).
\end{corollary}

The \(p_i\) are called the \emph{fundamental weights} corresponding to
the system of simple coroots \(\Delta^\ast\).

\begin{corollary}\label{sga3:XXI:6.6.4}
For every \(\alpha^\ast\in\Delta^\ast\),
\[
 \left(\alpha^\ast,\sum_i p_i\right)=1.
\]
Consequently,
\[
 \sum_i p_i=\rho_{R_+},
 \qquad
 R_+=\mathscr P\bigl(\operatorname{ind}(\Delta)\bigr),
\]
in the notation of \ref{sga3:XXI:3.5.1}.
\end{corollary}

\begin{corollary}\label{sga3:XXI:6.6.5}
For every \(x\in\VR(R)\),
\[
 x=\sum_i(\alpha_i^\ast,x)p_i.
\]
\end{corollary}

Since
\[
 R^\ast\subset\Gamma_0(R^\ast),
 \qquad
 R\subset\Lambda(\cR),
\]
the quadruple
\[
 \bigl(\Lambda(\cR),\Gamma_0(R^\ast),R,R^\ast\bigr)
\]
is a root datum.

\begin{corollary}\label{sga3:XXI:6.6.6}
The canonical morphism
\[
 \Gamma_0(R^\ast)\longrightarrow M^\ast
\]
is the transpose of
\[
 x\longmapsto\sum_i(\alpha_i^\ast,x)p_i.
\]
This map defines a morphism of root data, and there is a commutative
diagram
\begin{center}
\begin{tikzcd}[column sep=huge,row sep=large]
\cR \arrow[rr] \arrow[dr]
  && \operatorname{sc}(\cR) \arrow[dl,"\wr"'] \\
& \bigl(\Lambda(\cR),\Gamma_0(R^\ast),R,R^\ast\bigr) &
\end{tikzcd}
\end{center}
\end{corollary}

% END INCLUDED SGA3_English_current_buildable_source_20260823/source/inputs/XXI/tex/components/28_expose_XXI_section66_weights_partA_en.tex

% BEGIN INCLUDED SGA3_English_current_buildable_source_20260823/source/inputs/XXI/tex/components/29_expose_XXI_section66_partB_through_lemma671_en.tex
Thus we have an explicit description of \(\operatorname{sc}(\cR)\)
in terms of the weights of \(\cR\). Likewise, one obtains a
commutative diagram
\begin{center}
\begin{tikzcd}[column sep=huge,row sep=large]
\operatorname{ad}(\cR) \arrow[rr] \arrow[d,"\wr"]
  && \cR \\
\bigl(\Gamma_0(R),\Lambda(\cR^\ast),R,R^\ast\bigr)
  \arrow[urr] &&
\end{tikzcd}
\end{center}

\begin{corollary}\label{sga3:XXI:6.6.7}
The datum \(\cR\) is simply connected if and only if
\[
 M=\Lambda(\cR).
\]
\end{corollary}

\begin{remark}\label{sga3:XXI:6.6.8}
One has
\[
 \Lambda(\cR)\cap M=\VR(R)\cap M.
\]
Consequently, \(\cR\) is semisimple if and only if
\[
 M\subset\Lambda(\cR).
\]
\end{remark}

The results of \S\ref{sga3:XXI:6.5} also give:

\begin{corollary}\label{sga3:XXI:6.6.9}
The datum \(\cR\) is the product of a simply connected datum and a
trivial datum if and only if
\[
 M\supset\Lambda(\cR).
\]
\end{corollary}

Consider now the canonical isogeny
\[
 f:\operatorname{ad}(\cR)\longrightarrow\operatorname{sc}(\cR),
\]
and put \(Z(\cR)=K(f)\). One has
\[
 Z(\cR)\simeq Z(\operatorname{sc}\cR)
 \simeq Z(\operatorname{ad}\cR).
\]

\begin{corollary}\label{sga3:XXI:6.6.10}
There is a canonical isomorphism
\[
 Z(\cR)=\Lambda(\cR)/\Gamma_0(R).
\]
More precisely, there is an exact sequence of
\(W(\cR)\)-modules
\[
 0\longrightarrow\Gamma_0(R)\longrightarrow\Lambda(\cR)
 \longrightarrow Z(\cR)\longrightarrow0.
\]
\end{corollary}

\begin{corollary}\label{sga3:XXI:6.6.11}
There is a canonical pairing
\[
 Z(\cR^\ast)\times Z(\cR)\longrightarrow\mathbb Q/\mathbb Z
\]
which puts these groups in duality.
\end{corollary}

\begin{remark}\label{sga3:XXI:6.6.12}
One has
\[
 Z(\cR\times\cR')\simeq Z(\cR)\times Z(\cR').
\]
In particular, consider simply connected data
\(\cR_i\), \(i=1,2,\ldots,n\), and a trivial datum \(\cR_0\).
Put
\[
 \cR=\cR_0\times\cR_1\times\cdots\times\cR_n.
\]
If
\[
 \cR=(M,M^\ast,R,R^\ast),
 \qquad
 \cR_0=(M_0,M_0^\ast,\varnothing,\varnothing),
\]
then
\[
 M/\Gamma_0(R)\simeq
 M_0\times Z(\cR_1)\times\cdots\times Z(\cR_n).
\]
\end{remark}

\subsection{Automorphisms}\label{sga3:XXI:6.7}

By \ref{sga3:XXI:6.1.4}, an automorphism of \(\cR\) is an
automorphism \(u\) of \(M\) such that
\[
 u(R)=R,
 \qquad
 {}^tu(R^\ast)=R^\ast.
\]
In particular, every \(w\in W(\cR)\) defines an automorphism of
\(\cR\).

\begin{lemma}\label{sga3:XXI:6.7.1}
The group \(W(\cR)\) is a normal subgroup of
\(\operatorname{Aut}(\cR)\). More precisely, if
\(u\in\operatorname{Aut}(\cR)\) and \(\alpha\in R\), then
\[
 us_\alpha u^{-1}=s_{u(\alpha)}.
\]
\end{lemma}

\begin{proof}
The proof is the same as that of \ref{sga3:XXI:1.2.10}.
\end{proof}

% END INCLUDED SGA3_English_current_buildable_source_20260823/source/inputs/XXI/tex/components/29_expose_XXI_section66_partB_through_lemma671_en.tex

% BEGIN INCLUDED SGA3_English_current_buildable_source_20260823/source/inputs/XXI/tex/components/30_expose_XXI_proposition672_through_lemma677_en.tex
\begin{proposition}\label{sga3:XXI:6.7.2}
Let \(\Delta\) be a system of simple roots, and put
\[
 E_\Delta(\cR)=
 \{u\in\operatorname{Aut}(\cR)\mid u(\Delta)=\Delta\}.
\]
Then \(\operatorname{Aut}(\cR)\) is the semidirect product of
\(W(\cR)\) by \(E_\Delta(\cR)\).
\end{proposition}

\begin{proof}
This follows at once from the fact that \(W(\cR)\) acts simply
transitively on the systems of simple roots and from the fact that,
if \(\Delta\) is a system of simple roots of \(\cR\), then
\(u(\Delta)\) is a system of simple roots for every automorphism
\(u\) of \(\cR\).
\end{proof}

We shall later give a simpler description of \(E_\Delta(\cR)\) in
the case of reduced irreducible root data.\footnote{Editors' note:
When \(\cR\) is simply connected or adjoint, see \SGARefLink{sga3:XXI:7.4.5}{7.4.5}.}

\begin{definition}\label{sga3:XXI:6.7.3}
We denote by \(\operatorname{Aut}^{s}(\cR)\)\footnote{Editors' note:
The superscript \(s\) is intended to suggest ``semisimple.''}
the set of \(u\in\operatorname{Aut}(\cR)\) for which the following
diagram is commutative:
\begin{center}
\begin{tikzcd}[column sep=huge,row sep=large]
\cR \arrow[rr,"u"] \arrow[dr]
  && \cR \arrow[dl] \\
& \operatorname{rad}(\cR) &
\end{tikzcd}
\end{center}
We write
\[
 E_\Delta^{s}(\cR)
 =E_\Delta(\cR)\cap\operatorname{Aut}^{s}(\cR).
\]
\end{definition}

\begin{remark}\label{sga3:XXI:6.7.4}
If \(u\in\operatorname{Aut}(\cR)\), then
\(u\in\operatorname{Aut}^{s}(\cR)\) if and only if
\[
 (u-\operatorname{id})(M)\subset\VR(R).
\]
In particular,
\[
 W(\cR)\subset\operatorname{Aut}^{s}(\cR).
\]
\end{remark}

\begin{proposition}\label{sga3:XXI:6.7.5}
For every system of simple roots \(\Delta\), the group
\(\operatorname{Aut}^{s}(\cR)\) is the semidirect product of
\(W(\cR)\) by \(E_\Delta^{s}(\cR)\).
\end{proposition}

By functoriality, every automorphism of \(\cR\) induces an
automorphism of \(\operatorname{ad}(\cR)\). Thus there is a
canonical homomorphism
\[
 \operatorname{Aut}(\cR)\longrightarrow
 \operatorname{Aut}(\operatorname{ad}(\cR)).
\]

\begin{lemma}\label{sga3:XXI:6.7.6}
The homomorphism
\[
 \operatorname{Aut}^{s}(\cR)\longrightarrow
 \operatorname{Aut}(\operatorname{ad}(\cR))
\]
is injective.
\end{lemma}

\begin{proof}
Let \(u\) be an automorphism of \(M\) such that
\[
 (u-\operatorname{id})(M)\subset\VR(R)
\]
and
\[
 {}^tu(\alpha^\ast)=\alpha^\ast
 \qquad(\alpha^\ast\in R^\ast).
\]
For every \(x\in M\), one has
\[
 (\alpha^\ast,u(x)-x)
 =({}^tu(\alpha^\ast)-\alpha^\ast,x)=0.
\]
Hence \(u(x)-x=0\), by \ref{sga3:XXI:1.2.5}.
\end{proof}

\begin{lemma}\label{sga3:XXI:6.7.7}
The group \(\operatorname{Aut}^{s}(\cR)\) is finite.
\end{lemma}

\begin{proof}
It suffices to prove that \(\operatorname{Aut}(\cR)\) is finite when
\(\cR\) is adjoint. In that case \(M\) is generated by \(R\), so
every automorphism of \(\cR\) is determined by the permutation it
induces on \(R\).
\end{proof}

% END INCLUDED SGA3_English_current_buildable_source_20260823/source/inputs/XXI/tex/components/30_expose_XXI_proposition672_through_lemma677_en.tex

% BEGIN INCLUDED SGA3_English_current_buildable_source_20260823/source/inputs/XXI/tex/components/31_expose_XXI_remark678_through_example686_en.tex
\begin{remark}\label{sga3:XXI:6.7.8}
It is immediate that \(\operatorname{Aut}(\cR)\) (respectively
\(E_\Delta(\cR)\)) is finite if and only if
\[
 \rgred(\cR)-\rgss(\cR)\leq1.
\]
\end{remark}

\subsection{\(p\)-morphisms of reduced root data}
\label{sga3:XXI:6.8}

Throughout this subsection, \(p\) is a fixed positive integer.

\begin{definition}\label{sga3:XXI:6.8.1}
Let
\[
 \cR=(M,M^\ast,R,R^\ast),
 \qquad
 \cR'=(M',M'^\ast,R',R'^\ast)
\]
be two reduced root data. A group homomorphism
\[
 f:M'\longrightarrow M
\]
is called a \emph{\(p\)-morphism from \(\cR'\) to \(\cR\)} if
there exist a bijection
\[
 u:R\xrightarrow{\sim}R'
\]
and a map
\[
 q:R\longrightarrow\{p^n\mid n\in\mathbb N\}
\]
such that:
\begin{enumerate}
\item[(i)] \(f(u(\alpha))=q(\alpha)\alpha\) for every
  \(\alpha\in R\);
\item[(ii)] \({}^tf(\alpha^\ast)=q(\alpha)u(\alpha)^\ast\) for
  every \(\alpha\in R\).\footnote{Editors' note: These two
  conditions imply
  \[
   q(\alpha)\bigl(u(\alpha)^\ast,u(\beta)\bigr)
   =q(\beta)(\alpha^\ast,\beta)
  \]
  for all \(\alpha,\beta\in R\).}
\end{enumerate}
\end{definition}

\begin{corollary}\label{sga3:XXI:6.8.2}
A \(1\)-morphism is precisely a morphism.
\end{corollary}

\begin{corollary}\label{sga3:XXI:6.8.3}
The transpose of a \(p\)-morphism is a \(p\)-morphism.
\end{corollary}

\begin{lemma}\label{sga3:XXI:6.8.4}
If \(w\in W(\cR)\) and \(\alpha\in R\), then
\[
 q(w(\alpha))=q(\alpha).
\]
The map
\[
 s_\alpha\longmapsto s_{u(\alpha)}
\]
extends to an isomorphism
\[
 u:W(\cR)\xrightarrow{\sim}W(\cR')
\]
such that
\[
 u(w(\alpha))=u(w)\bigl(u(\alpha)\bigr).
\]
\end{lemma}

\begin{proof}
It suffices to prove, for \(\alpha,\beta\in R\), that
\[
 u(s_\alpha(\beta))=s_{u(\alpha)}u(\beta),
 \qquad
 q(s_\alpha(\beta))=q(\beta).
\]
Successively, one has
\begin{align*}
 f\bigl(s_{u(\alpha)}u(\beta)\bigr)
  &=f(u(\beta))
    -\bigl(u(\alpha)^\ast,u(\beta)\bigr)f(u(\alpha))\\
  &=q(\beta)\beta
    -q(\beta)q(\alpha)^{-1}
       (\alpha^\ast,\beta)q(\alpha)\alpha\\
  &=q(\beta)\bigl(\beta-(\alpha^\ast,\beta)\alpha\bigr)\\
  &=q(\beta)s_\alpha(\beta).
\end{align*}
If
\[
 \gamma=u^{-1}\bigl(s_{u(\alpha)}u(\beta)\bigr),
\]
then
\[
 q(\gamma)\gamma=f(u(\gamma))=q(\beta)s_\alpha(\beta).
\]
The roots \(\gamma\) and \(s_\alpha(\beta)\) are proportional over
\(\mathbb Q\), hence equal or opposite. Since \(q(\gamma)\) and
\(q(\beta)\) are positive, it follows that
\[
 q(\gamma)=q(\beta),
 \qquad
 \gamma=s_\alpha(\beta).
\]
\end{proof}

\begin{definition}\label{sga3:XXI:6.8.5}
The numbers \(q(\alpha)\) are called the \emph{root exponents} of
\(f\).
\end{definition}

\begin{example}\label{sga3:XXI:6.8.6}
Let \(\cR\) be a reduced root datum and let
\(q=p^n\), \(n\in\mathbb N\). Multiplication by \(q\),
\[
 M\longrightarrow M,\qquad x\longmapsto qx,
\]
is a \(p\)-morphism whose root exponents are all equal to \(q\)
(and for which \(u=\operatorname{id}\)). It is denoted
\[
 q:\cR\longrightarrow\cR.
\]
\end{example}

% END INCLUDED SGA3_English_current_buildable_source_20260823/source/inputs/XXI/tex/components/31_expose_XXI_remark678_through_example686_en.tex

% BEGIN INCLUDED SGA3_English_current_buildable_source_20260823/source/inputs/XXI/tex/components/32_expose_XXI_proposition687_through_proposition711_en.tex
\begin{proposition}\label{sga3:XXI:6.8.7}
In the notation of \ref{sga3:XXI:6.8.1}, \(u\) induces an
isomorphism from the set of systems of simple roots (respectively
positive roots) of \(R\) onto the corresponding set for \(R'\).
\end{proposition}

\begin{proof}
This follows immediately from \ref{sga3:XXI:3.1.5}
(respectively \ref{sga3:XXI:3.2.1}).
\end{proof}

\section{Structure}\label{sga3:XXI:7}

\subsection{Decomposition of a root datum}\label{sga3:XXI:7.1}

\begin{proposition}\label{sga3:XXI:7.1.1}
Let \(\cR\) be a root datum and let \(\Delta\) be a system of
simple roots.
\begin{enumerate}
\item[\SGAInlineRefTarget{sga3:XXI:7.1.1:item:i}(i)] Let \(R'\) and \(R''\) be two closed symmetric sets of
  roots forming a partition of \(R\). If
  \[
   \Delta'=\Delta\cap R',
   \qquad
   \Delta''=\Delta\cap R'',
  \]
  then
  \[
   R'=R_{\Delta'},
   \qquad
   R''=R_{\Delta''},
  \]
  and every root of \(\Delta'\) is orthogonal to every root of
  \(\Delta''\).
\item[\SGAInlineRefTarget{sga3:XXI:7.1.1:item:ii}(ii)] Let \(\Delta'\) and \(\Delta''\) be two orthogonal
  subsets forming a partition of \(\Delta\). Then
  \(R_{\Delta'}\) and \(R_{\Delta''}\) form a partition of \(R\).
\end{enumerate}
\end{proposition}

\begin{proof}
We first prove \SGARefLink{sga3:XXI:7.1.1:item:i}{(i)}.

\begin{lemma}\label{sga3:XXI:7.1.2}
Under the hypotheses of \textup{\SGARefLink{sga3:XXI:7.1.1:item:i}{(i)}}, if
\(\alpha\), \(\beta\), and \(\alpha+\beta\) are roots, then all
three belong to \(R'\), or all three belong to \(R''\).
\end{lemma}

Suppose, for example, that \(\alpha+\beta\in R'\). One cannot have
\(\alpha,\beta\in R''\), since \(R''\) is closed. Suppose therefore
that \(\alpha\in R'\). Then \(-\alpha\in R'\), and
\[
 \beta=(\beta+\alpha)-\alpha\in R'.
\]

We now prove \(R'=R_{\Delta'}\) by induction on the order of a
positive root \(\alpha\in R'\cap\mathscr P(\Delta)\). If
\(\operatorname{ord}_\Delta(\alpha)=1\), then
\[
 \alpha\in R'\cap\Delta=\Delta'.
\]
If \(\operatorname{ord}_\Delta(\alpha)>1\), there is
\(\beta\in\Delta\) such that \(\alpha-\beta\in R\). By the lemma,
\[
 \beta\in\Delta',
 \qquad
 \alpha-\beta\in R'.
\]
Hence \(\alpha-\beta\in R_{\Delta'}\) by induction, and finally
\[
 \alpha=(\alpha-\beta)+\beta\in R_{\Delta'}.
\]

We next prove that \(\Delta'\) and \(\Delta''\) are orthogonal. If
\(\alpha\in\Delta'\) and \(\beta\in\Delta''\), then
\[
 (\beta^\ast,\alpha)\leq0.
\]
If \((\beta^\ast,\alpha)\neq0\), then \(\beta+\alpha\) is a root,
contrary to the lemma.

We now prove \SGARefLink{sga3:XXI:7.1.1:item:ii}{(ii)}. If \(\Delta'\) or \(\Delta''\) is empty, the
assertion is immediate. Otherwise, if \(R_{\Delta'}\) and
\(R_{\Delta''}\) do not form a partition of \(R\), there is a root
\(\alpha\) of the form
\[
 \alpha=\sum_i m_i'\alpha_i'
       +\sum_j m_j''\alpha_j'',
 \qquad
 m_i',m_j''\in\mathbb Z_+,
\]
where the \(\alpha_i'\) (respectively \(\alpha_j''\)) belong to
\(\Delta'\) (respectively \(\Delta''\)). Applying
\ref{sga3:XXI:3.1.2}, and interchanging \(\Delta'\) and
\(\Delta''\) if necessary, yields a relation
\[
 \delta=\gamma+\beta,
 \qquad
 \gamma\in R_{\Delta'},
 \quad
 \beta\in\Delta'',
 \quad
 \delta\in R.
\]
But \((\beta^\ast,\gamma)=0\), so
\(\gamma-\beta\) is also a root by \ref{sga3:XXI:2.2.5}, which is
impossible.
\end{proof}

% END INCLUDED SGA3_English_current_buildable_source_20260823/source/inputs/XXI/tex/components/32_expose_XXI_proposition687_through_proposition711_en.tex

% BEGIN INCLUDED SGA3_English_current_buildable_source_20260823/source/inputs/XXI/tex/components/33_expose_XXI_proposition713_through_lemma723_en.tex
\begin{proposition}\label{sga3:XXI:7.1.3}
Let \(\cR\) be a root datum. The following conditions are
equivalent:
\begin{enumerate}
\item[\SGAInlineRefTarget{sga3:XXI:7.1.3:item:i}(i)] There is no nontrivial partition of \(R\) into two
  closed symmetric subsets.
\item[\SGAInlineRefTarget{sga3:XXI:7.1:item:ii}(ii)] For some (respectively every) system of simple roots
  \(\Delta\) of \(R\), there is no partition of \(\Delta\) into
  two nonempty orthogonal subsets.
\item[\SGAInlineRefTarget{sga3:XXI:7.1.3:item:iii}(iii)] The natural representation of \(W(\cR)\) on
  \(\VR(R)\) is irreducible.
\item[\SGAInlineRefTarget{sga3:XXI:7.1.3:item:iv}(iv)] For every pair \((\alpha,\beta)\) of roots, there is a
  sequence of roots
  \[
   \alpha_0,\alpha_1,\ldots,\alpha_n,
   \qquad
   \alpha_0=\alpha,\quad\alpha_n=\beta,
  \]
  such that \(\alpha_i\) and \(\alpha_{i+1}\) are nonorthogonal
  for \(i=0,\ldots,n-1\).
\end{enumerate}
\end{proposition}

\begin{proof}
The equivalence \SGARefLink{sga3:XXI:7.1.3:item:i}{(i)}\(\Leftrightarrow\)\SGARefLink{sga3:XXI:7.1:item:ii}{(ii)} follows from
\ref{sga3:XXI:7.1.1}, and \SGARefLink{sga3:XXI:7.1.3:item:iv}{(iv)}\(\Rightarrow\)\SGARefLink{sga3:XXI:7.1.1:item:ii}{(ii)} is immediate.
Conversely, if \SGARefLink{sga3:XXI:7.1.1:item:ii}{(ii)} holds for \(\Delta\), then \SGARefLink{sga3:XXI:7.1.3:item:iv}{(iv)} holds when
\(\alpha,\beta\in\Delta\). Every root is nonorthogonal to some
simple root, by \ref{sga3:XXI:3.1.1}, for example. Moreover,
\SGARefLink{sga3:XXI:7.1.3:item:iii}{(iii)}\(\Rightarrow\)\SGARefLink{sga3:XXI:7.1.3:item:i}{(i)}: under the hypotheses of
\ref{sga3:XXI:7.1.1}, the space \(\VR(R')\) is stable under
\(W(\cR)\). It remains to prove \SGARefLink{sga3:XXI:7.1.1:item:i}{(i)}\(\Rightarrow\)\SGARefLink{sga3:XXI:7.1.3:item:iii}{(iii)}.

Let \(H\) be a vector subspace of \(\VR(R)\), stable under
\(W(\cR)\). For every \(\alpha\in R\), the equality
\(s_\alpha(H)=H\) immediately gives
\[
 \alpha\in H
 \qquad\text{or}\qquad
 \alpha^\ast\in H^\perp,
\]
where \(H^\perp\) denotes the orthogonal complement of \(H\) in
\(\VR(R^\ast)\), which is in duality with \(\VR(R)\). Put
\[
 R'=\{\alpha\in R\mid\alpha\in H\},
 \qquad
 R''=\{\alpha\in R\mid\alpha^\ast\in H^\perp\}.
\]
This gives a partition of \(R\) into two closed symmetric subsets.
\end{proof}

\begin{definition}\label{sga3:XXI:7.1.4}
A root datum (respectively root system) satisfying the equivalent
conditions of \ref{sga3:XXI:7.1.3} and having nonzero semisimple
rank is called \emph{irreducible}.
\end{definition}

\begin{corollary}\label{sga3:XXI:7.1.5}
For every root datum \(\cR\), there is a unique partition, up to
order, of \(R\) into closed symmetric irreducible subsets.
\end{corollary}

\begin{corollary}\label{sga3:XXI:7.1.6}
Every adjoint (respectively simply connected) root datum is a
product of irreducible adjoint (respectively simply connected) root
data.
\end{corollary}

\begin{proof}
It suffices to consider the adjoint case. The assertion then follows
from the fact that, under the hypotheses of
\ref{sga3:XXI:7.1.1},
\[
 \Gamma_0(R)=\Gamma_0(R')\times\Gamma_0(R'').
\]
\end{proof}

\begin{corollary}\label{sga3:XXI:7.1.7}
For every root datum (respectively reduced root datum) \(\cR\),
there is an isogeny
\[
 \cR\longrightarrow\cR',
\]
where \(\cR'\) is a product of a trivial root datum and irreducible
simply connected root data (respectively irreducible reduced simply
connected root data).
\end{corollary}

\subsection{Properties of irreducible root data}
\label{sga3:XXI:7.2}

\begin{definition}\label{sga3:XXI:7.2.1}
Let \(\cR\) be an irreducible root datum. For every \(\alpha\in R\),
put
\[
 \operatorname{long}(\alpha)
 =\ell(\alpha)/\ell(\alpha_0),
\]
where \(\alpha_0\in R\) is such that \(\ell(\alpha_0)\) is minimal.
The number \(\operatorname{long}(\alpha)\) is called the
\emph{length} of \(\alpha\).
\end{definition}

\begin{lemma}\label{sga3:XXI:7.2.2}
Let \(\cR\) be an irreducible root datum. The Weyl group acts
transitively on the set of roots of a given length.
\end{lemma}

\begin{proof}
Let \(\alpha,\beta\in R\). Since the representation of \(W(\cR)\)
on \(\VR(R)\) is irreducible, \(\alpha\) cannot be orthogonal to
all \(w(\beta)\), \(w\in W(\cR)\). Hence some \(w(\beta)\) is
nonorthogonal to \(\alpha\). Since
\[
 \ell(w(\beta))=\ell(\beta),
\]
the assertion follows from \ref{sga3:XXI:2.3.2}.
\end{proof}

\begin{lemma}\label{sga3:XXI:7.2.3}
If \(\cR\) is irreducible and reduced, then
\[
 \operatorname{long}(R)
 \in\bigl\{\{1\},\{1,2\},\{1,3\}\bigr\}.
\]
\end{lemma}

\begin{proof}
By the observation used above, for every
\(\alpha,\beta\in R\) there is \(w\in W(\cR)\) such that
\(w(\beta)\) is nonorthogonal to \(\alpha\). Therefore, by
\ref{sga3:XXI:2.3.1},
\[
 \frac{\ell(\alpha)}{\ell(\beta)}
 \in\left\{1,2,3,\frac12,\frac13\right\}.
\]
Thus \(\operatorname{long}(\alpha)\) is \(1\), \(2\), or \(3\).
But lengths \(2\) and \(3\) cannot both occur, since their ratio
\(2/3\) is not in the displayed set.
\end{proof}

% END INCLUDED SGA3_English_current_buildable_source_20260823/source/inputs/XXI/tex/components/33_expose_XXI_proposition713_through_lemma723_en.tex

% BEGIN INCLUDED SGA3_English_current_buildable_source_20260823/source/inputs/XXI/tex/components/34_expose_XXI_remark724_through_remark729_en.tex
\begin{remark}\label{sga3:XXI:7.2.4}
Reasoning similarly gives the following result.  If \(\cR\) is
irreducible and nonreduced, with \(\rgss(\cR)>1\), then
\[
 \operatorname{long}(R)=\{1,2,4\}.
\]
If \(R_i=\operatorname{long}^{-1}(i)\), then
\[
 \operatorname{ind}(R)=R_1\cup R_2,\qquad R_4=2R_1,
\]
and two nonproportional roots of \(R_1\) are orthogonal.
Conversely, suppose that \(R\) is an irreducible reduced root
system such that \(\operatorname{long}(R)=\{1,2\}\).  Put
\(R_i=\operatorname{long}^{-1}(i)\), and suppose that two
nonproportional roots of \(R_1\) are orthogonal.  Then
\(R\cup2R_1\) is irreducible and nonreduced, and
\[
 \operatorname{ind}(R\cup2R_1)=R.
\]
\end{remark}

\begin{lemma}\label{sga3:XXI:7.2.5}
If \(\cR\) is an irreducible root datum, then \(\cR^\ast\) is
irreducible as well, and the product
\[
 \operatorname{long}(\alpha)\operatorname{long}(\alpha^\ast)
\]
is constant as \(\alpha\) runs through \(R\).
\end{lemma}

\begin{proof}
This follows immediately from \ref{sga3:XXI:7.1.3}\,\SGARefLink{sga3:XXI:7.2:item:iv}{(iv)} and
\ref{sga3:XXI:2.2.6}.
\end{proof}

\begin{definition}\label{sga3:XXI:7.2.6}
Let \(\cR\) be any root datum.  The \emph{length} of
\(\alpha\in R\), denoted \(\operatorname{long}(\alpha)\), is the
length of \(\alpha\) in its irreducible component.
\end{definition}

\begin{lemma}\label{sga3:XXI:7.2.7}
There is a unique group homomorphism
\[
 u:\Gamma_0(R)\longrightarrow\Gamma_0(R^\ast)
\]
such that
\[
 u(\alpha)=\operatorname{long}(\alpha)\alpha^\ast
 \qquad(\alpha\in R).
\]
\end{lemma}

\begin{proof}
By \ref{sga3:XXI:3.5.5}, it is enough to verify that, whenever
\(\alpha,\beta,\alpha+\beta\in R\),
\[
 \operatorname{long}(\alpha)\alpha^\ast+
 \operatorname{long}(\beta)\beta^\ast
 =
 \operatorname{long}(\alpha+\beta)(\alpha+\beta)^\ast.
\]
But \(\alpha,\beta,\alpha+\beta\) belong to the same irreducible
component of \(R\), by \ref{sga3:XXI:7.1.2}; the assertion is
therefore reduced to \ref{sga3:XXI:1.2.2}.
\end{proof}

\begin{remark}\label{sga3:XXI:7.2.8}
Let \(u\) be as in \ref{sga3:XXI:7.2.7}.  For
\(\alpha,\beta\in R\),
\[
 (u(\alpha),\beta)=(u(\beta),\alpha).
\]
Indeed, this amounts to
\[
 \operatorname{long}(\alpha)(\alpha^\ast,\beta)
 =
 \operatorname{long}(\beta)(\beta^\ast,\alpha),
\]
which is clear when \(\alpha\) and \(\beta\) are orthogonal.  If
they are not orthogonal, they lie in the same irreducible component
of \(R\), and the assertion reduces to \ref{sga3:XXI:1.2.1},
formula~\SGARefLink{sga3:XXI:eq:9}{(9)}.
\end{remark}

\begin{remark}\label{sga3:XXI:7.2.9}
The symmetric bilinear form
\[
 (x,y)\longmapsto(u(x),y)
\]
is positive and nondegenerate on \(\Gamma_0(R)\).
\end{remark}

\begin{proof}
Let \(R_i\) be the irreducible components of \(R\).  Then
\[
 \Gamma_0(R)=\prod_i\Gamma_0(R_i),
\]
and the bilinear form \((u(x),y)\) is the product, on the groups
\(\Gamma_0(R_i)\), of the forms
\[
 2\ell(\alpha_i)^{-1}(p(x),y),
\]
where \(\ell(\alpha_i)\) is the minimum of \(\ell(\alpha)\) for
\(\alpha\in R_i\).  These symmetric bilinear forms are positive
and nondegenerate by \ref{sga3:XXI:1.2.6}.
\end{proof}

% END INCLUDED SGA3_English_current_buildable_source_20260823/source/inputs/XXI/tex/components/34_expose_XXI_remark724_through_remark729_en.tex

% BEGIN INCLUDED SGA3_English_current_buildable_source_20260823/source/inputs/XXI/tex/components/35_expose_XXI_section73_cartan_matrix_en.tex
\subsection{Cartan matrix}\label{sga3:XXI:7.3}

Let \(\cR\) be a root datum.  If \(\Delta\) is a system of simple
roots, the \emph{Cartan matrix} of \(\cR\) relative to \(\Delta\)
is the square matrix indexed by \(\Delta\) and defined by
\[
 a_{\alpha,\beta}=(\alpha^\ast,\beta),
 \qquad \alpha,\beta\in\Delta.
\]
If \(\Delta'\) is another system of simple roots and
\(w\in W(\cR)\) satisfies \(w(\Delta)=\Delta'\), then
\[
 (w(\alpha)^\ast,w(\beta))=(\alpha^\ast,\beta).
\]
Thus the Cartan matrix relative to \(\Delta'\) is obtained from
that relative to \(\Delta\) by the isomorphism
\(\Delta\to\Delta'\) of index sets defined by \(w\).  Consequently,
up to isomorphism of its index set, the Cartan matrix depends only
on \(\cR\).

\begin{proposition}\label{sga3:XXI:7.3.1}
The Cartan matrix has the following properties.
\begin{enumerate}
\item[\SGAInlineRefTarget{sga3:XXI:7.3.1:item:i}(i)] \(a_{\alpha,\alpha}=2\), and
  \(a_{\alpha,\beta}\leq0\) for \(\alpha\neq\beta\).
\item[\SGAInlineRefTarget{sga3:XXI:7.3.1:item:ii}(ii)] \(a_{\alpha,\beta}=0\) implies
  \(a_{\beta,\alpha}=0\).
\item[\SGAInlineRefTarget{sga3:XXI:7.2:item:iii}(iii)] There are strictly positive integers
  \(m_\alpha=\operatorname{long}(\alpha)\) such that the matrix
  \[
   (m_\alpha a_{\alpha,\beta})
  \]
  is symmetric, positive, and nondegenerate.
\item[\SGAInlineRefTarget{sga3:XXI:7.2:item:iv}(iv)] The principal minors of
  \((a_{\alpha,\beta})_{\alpha,\beta\in\Delta}\), namely
  \[
   \det(a_{\alpha,\beta})_{\alpha,\beta\in\Delta'}
   \qquad(\Delta'\subset\Delta),
  \]
  are strictly positive.
\item[\SGAInlineRefTarget{sga3:XXI:7.2:item:v}(v)] One has
  \[
   s_\alpha(\beta)=\beta-a_{\alpha,\beta}\alpha,
   \qquad
   s_\alpha(\beta^\ast)
   =\beta^\ast-a_{\beta,\alpha}\alpha^\ast.
  \]
\end{enumerate}
\end{proposition}

\begin{proof}
Part \SGARefLink{sga3:XXI:7.2:item:v}{(v)} is a definition; \SGARefLink{sga3:XXI:7.3.1:item:i}{(i)} follows from
\ref{sga3:XXI:3.2.11}, \SGARefLink{sga3:XXI:7.3.1:item:ii}{(ii)} from \ref{sga3:XXI:2.2.2}, and
\SGARefLink{sga3:XXI:7.2:item:iii}{(iii)} from \ref{sga3:XXI:7.2.9}.  Part \SGARefLink{sga3:XXI:7.2:item:iv}{(iv)} follows immediately
from \SGARefLink{sga3:XXI:7.2:item:iii}{(iii)} and
\[
 \det(m_\alpha a_{\alpha,\beta})_{\alpha,\beta\in\Delta'}
 =
 \left(\prod_{\alpha\in\Delta'}m_\alpha\right)
 \det(a_{\alpha,\beta})_{\alpha,\beta\in\Delta'}.
\]
\end{proof}

\begin{proposition}\label{sga3:XXI:7.3.2}
Let \(\cR\) and \(\cR'\) be two reduced simply connected
(respectively adjoint) root data, let \(\Delta\) and \(\Delta'\)
be systems of simple roots of \(\cR\) and \(\cR'\), and let
\(u:\Delta\to\Delta'\) be an isomorphism.  Denote the corresponding
Cartan matrices by \((a_{\alpha,\beta})\) and
\((a'_{\alpha',\beta'})\), and suppose that
\[
 a'_{u(\alpha),u(\beta)}=a_{\alpha,\beta}.
\]
Then there is a unique isomorphism from \(\cR\) onto \(\cR'\)
that induces \(u\) on \(\Delta\).
\end{proposition}

\begin{proof}
It is enough to give the proof in the adjoint case.  Then
\(M=\Gamma_0(R)\) and \(M'=\Gamma_0(R')\) are the free abelian
groups generated by \(\Delta\) and \(\Delta'\).  Hence there is a
unique group isomorphism \(M\to M'\) inducing \(u\) on \(\Delta\);
denote it again by \(u\).

We show that \(u(R)\subset R'\).  Every root \(\alpha\) of \(R\)
can be written
\[
 \alpha=s_{\alpha_1}\cdots s_{\alpha_n}(\alpha_{n+1}),
 \qquad \alpha_i\in\Delta.
\]
By the hypothesis on \(u\) and the relations
\ref{sga3:XXI:7.3.1}\,\SGARefLink{sga3:XXI:7.2:item:v}{(v)},
\[
 u(\alpha)=
 s_{u(\alpha_1)}\cdots s_{u(\alpha_n)}
 \bigl(u(\alpha_{n+1})\bigr).
\]
It remains to prove that \({}^tu(R'^\ast)\subset R^\ast\).  This
follows from the fact that the elements of \(M^\ast\)
(respectively \(M'^\ast\)) are determined by their duality with
\(R\) or \(\Delta\) (respectively \(R'\) or \(\Delta'\)), by
\ref{sga3:XXI:1.2.5}.
\end{proof}

\begin{corollary}\label{sga3:XXI:7.3.3}
A reduced simply connected or adjoint root datum is determined up
to isomorphism by its Cartan matrix.
\end{corollary}

\begin{corollary}\label{sga3:XXI:7.3.4}
Let \(\cR\) be a reduced simply connected (respectively adjoint)
root datum and let \(\Delta\) be a system of simple roots.  The
group \(E_\Delta(\cR)\) is identified with the group of
automorphisms of the set \(\Delta\) that preserve the Cartan
matrix.
\end{corollary}

\begin{remark}\label{sga3:XXI:7.3.5}
The existence of a root datum corresponding to a given Cartan
matrix satisfying, for example, \SGARefLink{sga3:XXI:7.3.1:item:i}{(i)}, \SGARefLink{sga3:XXI:7.3.1:item:ii}{(ii)}, and \SGARefLink{sga3:XXI:7.2:item:iv}{(iv)} is not easily
settled directly without using the classification.
\end{remark}

% END INCLUDED SGA3_English_current_buildable_source_20260823/source/inputs/XXI/tex/components/35_expose_XXI_section73_cartan_matrix_en.tex

% BEGIN INCLUDED SGA3_English_current_buildable_source_20260823/source/inputs/XXI/tex/components/36_expose_XXI_section74_dynkin_diagrams_en.tex
\subsection{Dynkin diagrams}\label{sga3:XXI:7.4}

\begin{definition}\label{sga3:XXI:7.4.1}
A \emph{Dynkin-diagram structure} on a finite set \(\Delta\)
(the word ``scheme'' being banished for obvious reasons) consists
of a set of pairs of distinct elements of \(\Delta\), called
\emph{linked pairs}, together with a map
\[
 \Delta\longrightarrow\{1,2,3\}.
\]
The notion of isomorphism between such structures is evident.
\end{definition}

\begin{definition}\label{sga3:XXI:7.4.2}
Let \(\cR\) be a root datum and let \(\Delta\) be a system of
simple roots.  The \emph{Dynkin diagram} of \(\cR\) relative to
\(\Delta\) is the set \(\Delta\), in which two simple roots are
linked if and only if they are nonorthogonal, and each root is
assigned its length.
\end{definition}

\begin{proposition}\label{sga3:XXI:7.4.3}
The Dynkin diagram and the Cartan matrix determine one another
uniquely.
\end{proposition}

\begin{proof}
The equivalence
\[
 \alpha\text{ and }\beta\text{ are not linked}
 \quad\Longleftrightarrow\quad
 a_{\alpha,\beta}=0
\]
and the relation
\[
 \operatorname{long}(\alpha)a_{\alpha,\beta}
 =
 \operatorname{long}(\beta)a_{\beta,\alpha},
\]
with \(\inf\operatorname{long}(\alpha)=1\) in each connected
component of the diagram, determine the \(a_{\alpha,\beta}\) from
the links and lengths, and conversely.  The details are left to the
reader.
\end{proof}

\begin{corollary}\label{sga3:XXI:7.4.4}
A reduced simply connected or adjoint root datum is determined by
its Dynkin diagram.
\end{corollary}

\begin{corollary}\label{sga3:XXI:7.4.5}
Let \(\cR\) be a reduced simply connected or adjoint root datum and
let \(\Delta\) be a system of simple roots.  The group
\(E_\Delta(\cR)\) is identified with the group of automorphisms of
the Dynkin diagram of \(\cR\) relative to \(\Delta\); that is, with
the group of permutations of \(\Delta\) preserving lengths and
links.
\end{corollary}

\begin{remark}\label{sga3:XXI:7.4.6}
The usual method\footnote{See Bourbaki, \emph{Groupes et algèbres
de Lie}, Chapter~VI, no.~4.2, or the Sophus Lie Seminar.}
\footnote{Editors' note: for a slightly different proof, see also
[De80].}
classifies the connected Dynkin diagrams and shows that each
actually corresponds to an irreducible reduced simply connected
root datum.  One obtains the familiar types:
\begin{center}
% BEGIN INCLUDED SGA3_English_current_buildable_source_20260823/source/inputs/XXI/tex/figures/xxi_p44_dynkin_classification_native.tex
\begin{tikzpicture}[
  x=.86cm,y=.78cm,
  dyn/.style={circle,fill=black,inner sep=1.35pt},
  wt/.style={font=\small,above=2pt},
  typ/.style={font=\normalsize,anchor=west}
]
% A_n
\foreach \x in {0,1,2,4,5,6} {
  \node[dyn] (A\x) at (\x,0) {};
  \node[wt] at (A\x) {$1$};
}
\draw (A0)--(A1)--(A2);
\draw (A4)--(A5)--(A6);
\node at (3,0) {$\cdots$};
\node[typ] at (7.3,0) {$A_n,\qquad n\geq1.$};

% B_n
\foreach \x in {0,1,2,4,5} {
  \node[dyn] (B\x) at (\x,-1.55) {};
  \node[wt] at (B\x) {$2$};
}
\node[dyn] (B6) at (6,-1.55) {};
\node[wt] at (B6) {$1$};
\draw (B0)--(B1)--(B2);
\draw (B4)--(B5)--(B6);
\node at (3,-1.55) {$\cdots$};
\node[typ] at (7.3,-1.55) {$B_n,\qquad n\geq2.$};

% C_n
\foreach \x in {0,1,2,4,5} {
  \node[dyn] (C\x) at (\x,-3.10) {};
  \node[wt] at (C\x) {$1$};
}
\node[dyn] (C6) at (6,-3.10) {};
\node[wt] at (C6) {$2$};
\draw (C0)--(C1)--(C2);
\draw (C4)--(C5)--(C6);
\node at (3,-3.10) {$\cdots$};
\node[typ] at (7.3,-3.10) {$C_n,\qquad n\geq3.$};

% D_n
\foreach \x in {0,1,2,4} {
  \node[dyn] (D\x) at (\x,-4.65) {};
  \node[wt] at (D\x) {$1$};
}
\node[dyn] (Du) at (5,-4.25) {};
\node[wt] at (Du) {$1$};
\node[dyn] (Dl) at (5,-5.05) {};
\node[wt] at (Dl) {$1$};
\draw (D0)--(D1)--(D2);
\draw (D4)--(Du);
\draw (D4)--(Dl);
\node at (3,-4.65) {$\cdots$};
\node[typ] at (7.3,-4.65) {$D_n,\qquad n\geq4.$};

% E_n
\foreach \x in {0,1,3,4,5} {
  \node[dyn] (E\x) at (\x,-6.35) {};
  \node[wt] at (E\x) {$1$};
}
\node[dyn] (Ev) at (3,-7.15) {};
\node[right=2pt,font=\small] at (Ev) {$1$};
\draw (E0)--(E1);
\draw (E3)--(E4)--(E5);
\draw (E3)--(Ev);
\node at (2,-6.35) {$\cdots$};
\node[typ] at (7.3,-6.65) {$E_n,\qquad n=6,7,8.$};

% F_4
\foreach \x/\w in {0/1,1/1,2/2,3/2} {
  \node[dyn] (F\x) at (\x,-8.15) {};
  \node[wt] at (F\x) {$\w$};
}
\draw (F0)--(F1)--(F2)--(F3);
\node[typ] at (7.3,-8.15) {$F_4.$};

% G_2
\node[dyn] (G0) at (2,-9.70) {};
\node[wt] at (G0) {$1$};
\node[dyn] (G1) at (3,-9.70) {};
\node[wt] at (G1) {$3$};
\draw (G0)--(G1);
\node[typ] at (7.3,-9.70) {$G_2.$};
\end{tikzpicture}

% END INCLUDED SGA3_English_current_buildable_source_20260823/source/inputs/XXI/tex/figures/xxi_p44_dynkin_classification_native.tex

\end{center}
By \ref{sga3:XXI:7.4.5}, the corresponding group
\(E_\Delta(\cR)\) is obtained at once:
\[
\begin{array}{ll}
E_\Delta(\cR)=\{e\},
 & A_1,\ B_n,\ C_n,\ E_7,\ E_8,\ F_4,\ G_2,\\[2pt]
E_\Delta(\cR)=\mathbb Z/2\mathbb Z,
 & A_n\ (n\geq2),\ D_n\ (n\geq5),\ E_6,\\[2pt]
E_\Delta(\cR)=\mathfrak S_3,
 & D_4.
\end{array}
\]
\end{remark}

% END INCLUDED SGA3_English_current_buildable_source_20260823/source/inputs/XXI/tex/components/36_expose_XXI_section74_dynkin_diagrams_en.tex

% BEGIN INCLUDED SGA3_English_current_buildable_source_20260823/source/inputs/XXI/tex/components/37_expose_XXI_section75_p_morphisms_and_bibliography_en.tex
\subsection{Complements on \(p\)-morphisms}\label{sga3:XXI:7.5}

Let \(f:\cR\to\cR'\) be a \(p\)-morphism
(cf.\ \ref{sga3:XXI:6.8}).  It is clear from the definitions that
the bijection \(u:R\xrightarrow{\sim}R'\) associated with \(f\)
makes the systems of simple roots, the systems of positive roots,
the irreducible components, and so forth of \(R\) and \(R'\)
correspond.  We may therefore suppose, for simplicity, that
\(R\) and \(R'\) are irreducible.

\begin{lemma}\label{sga3:XXI:7.5.1}
If \(R\) and \(R'\) are irreducible, there is a \(k\in\mathbb Q\)
such that, for every \(\alpha\in R\),
\[
 k\,\operatorname{long}(u(\alpha))
 =
 q(\alpha)^2\operatorname{long}(\alpha).
\]
\end{lemma}

\begin{proof}
One has
\[
 \operatorname{long}(\alpha)(\alpha^\ast,\beta)
 =
 \operatorname{long}(\beta)(\beta^\ast,\alpha)
\]
and, similarly,
\[
 \operatorname{long}(u(\alpha))
   (u(\alpha)^\ast,u(\beta))
 =
 \operatorname{long}(u(\beta))
   (u(\beta)^\ast,u(\alpha)).
\]
It follows immediately that, for nonorthogonal \(\alpha,\beta\),
\[
 \frac{q(\alpha)^2\operatorname{long}(\alpha)}
      {\operatorname{long}(u(\alpha))}
 =
 \frac{q(\beta)^2\operatorname{long}(\beta)}
      {\operatorname{long}(u(\beta))}.
\]
The conclusion then follows from \ref{sga3:XXI:7.1.3}\,\SGARefLink{sga3:XXI:7.1.3:item:iv}{(iv)}.
\end{proof}

\begin{remark}\label{sga3:XXI:7.5.2}
It follows from \ref{sga3:XXI:7.2.2} and
\ref{sga3:XXI:6.8.4} that \(q(\alpha)\) depends only on
\(\operatorname{long}(\alpha)\).  One then sees readily that, if
\(q(\alpha)\) is not constant, the product
\(q(\alpha)\operatorname{long}(\alpha)\) is constant; hence
\(p=2\) or \(3\).  Inspection of the preceding diagrams gives the
four possible cases (the same letter denotes a Dynkin diagram and
the corresponding reduced simply connected root datum):
\[
\begin{array}{llll}
p=2, & B_n\xrightarrow{\,f_1\,}C_n,
     & C_n\xrightarrow{\,f_2\,}B_n,
     & (C_2=B_2),\\
p=2, & F_4\xrightarrow{\,g\,}F_4,\\
p=3, & G_2\xrightarrow{\,h\,}G_2.
\end{array}
\]
The reader will observe that
\(f_1\circ f_2\), \(f_2\circ f_1\), \(g\circ g\), and
\(h\circ h\) are \(p\)-morphisms of the form described in
\ref{sga3:XXI:6.8.6}.
\end{remark}

\medskip
\noindent\textbf{7.5.3.}\label{sga3:XXI:7.5.3}
The preceding description shows immediately that, if \(\cR\) and
\(\cR'\) are two reduced root data of semisimple rank at most \(2\)
and there is a \(p\)-morphism from \(\cR'\) to \(\cR\), then
\(\cR\) and \(\cR'\) have the same type.  More precisely, one has
the following table.

\medskip
\noindent\emph{Notation.}
Let \(f:\cR'\to\cR\) be a \(p\)-morphism.  The symbols \(q\) and
\(q_1\) denote arbitrary positive powers of \(p\).  For rank-two
systems we use the notation of Section~\ref{sga3:XXI:4}, with
\(\alpha,\beta\) the simple roots and
\(\ell(\alpha)\leq\ell(\beta)\).

\begin{center}
\begingroup
\renewcommand{\arraystretch}{1.35}
\resizebox{\linewidth}{!}{%
\begin{tabular}{|c|c|c|c|}
\hline
\textbf{Type} & \(p\) & \textbf{values of \(f\)}
 & \textbf{values of \({}^tf\)}\\
\hline
Trivial & arbitrary & --- & ---\\
\hline
\(A_1\) & arbitrary
 & \(f(\alpha')=q\alpha\)
 & \({}^tf(\alpha^\ast)=q\alpha'^\ast\)\\
\hline
\(A_1\times A_1\) & arbitrary
 & \(\begin{array}{c}f(\alpha')=q\alpha\\
                      f(\beta')=q_1\beta\end{array}\)
 & \(\begin{array}{c}{}^tf(\alpha^\ast)=q\alpha'^\ast\\
                      {}^tf(\beta^\ast)=q_1\beta'^\ast\end{array}\)\\
\hline
\(A_2,B_2,G_2\) & arbitrary
 & \(\begin{array}{c}f(\alpha')=q\alpha\\
                      f(\beta')=q\beta\end{array}\)
 & \(\begin{array}{c}{}^tf(\alpha^\ast)=q\alpha'^\ast\\
                      {}^tf(\beta^\ast)=q\beta'^\ast\end{array}\)\\
\hline
\(B_2\) & \(p=2\)
 & \(\begin{array}{c}f(\alpha')=q\beta\\
                      f(\beta')=2q\alpha\end{array}\)
 & \(\begin{array}{c}{}^tf(\alpha^\ast)=2q\beta'^\ast\\
                      {}^tf(\beta^\ast)=q\alpha'^\ast\end{array}\)\\
\hline
\(G_2\) & \(p=3\)
 & \(\begin{array}{c}f(\alpha')=q\beta\\
                      f(\beta')=3q\alpha\end{array}\)
 & \(\begin{array}{c}{}^tf(\alpha^\ast)=3q\beta'^\ast\\
                      {}^tf(\beta^\ast)=q\alpha'^\ast\end{array}\)\\
\hline
\end{tabular}}
\endgroup
\end{center}

\section*{Bibliography}
\addcontentsline{toc}{section}{Bibliography}
\label{sga3:XXI:bibliography}

\begin{description}
\item[{[BLie]}] N.~Bourbaki,
  \emph{Groupes et algèbres de Lie}, Chapters IV--VI,
  Hermann, 1968.
\item[{[De80]}] M.~Demazure,
  ``A, B, C, D, E, F, etc.,'' pp.~221--227 in
  \emph{Séminaire sur les singularités des surfaces}
  (Palaiseau, 1976--1977), eds.\ M.~Demazure,
  H.~C.~Pinkham, and B.~Teissier,
  Lecture Notes in Mathematics \textbf{777},
  Springer-Verlag, 1980.
\end{description}

% END INCLUDED SGA3_English_current_buildable_source_20260823/source/inputs/XXI/tex/components/37_expose_XXI_section75_p_morphisms_and_bibliography_en.tex



% END INCLUDED SGA3_English_current_buildable_source_20260823/source/inputs/XXI/body.tex

\counterwithin{theorem}{section}

\SGAStartExpose{XXII}
% BEGIN INCLUDED SGA3_English_current_buildable_source_20260823/source/inputs/XXII/body.tex
% Mechanically extracted after \mainmatter from the frozen standalone master.

\markboth{EXPOSÉ XXII. REDUCTIVE GROUPS}
  {EXPOSÉ XXII. REDUCTIVE GROUPS}
% BEGIN INCLUDED SGA3_English_current_buildable_source_20260823/source/inputs/XXII/tex/components/01_expose_XXII_title_section1.tex
% SGA 3, Exposé XXII, Loop-1 English reconstruction.
% Current scope: authority local pages 1--63 / combined-reader pages 1141--1203.
% Local GPU OCR and Reinhold's English Markdown are locator/comparison
% witnesses only. The Polo--Gille page image controls text and formulas.

\setcounter{footnote}{0}
\SGARefTarget{sga3:xxii:chapter:xxii}\chapter*{Exposé XXII\\
Reductive Groups: Splittings, Subgroups, Quotient Groups}
\addcontentsline{toc}{chapter}{Exposé XXII. Reductive Groups:
Splittings, Subgroups, Quotient Groups}
\markboth{EXPOSÉ XXII. REDUCTIVE GROUPS}
  {EXPOSÉ XXII. REDUCTIVE GROUPS}

\begin{center}
  \textit{by Michel Demazure}
\end{center}

\begingroup
\renewcommand{\thefootnote}{0}
\footnote{Editor's note: Version dated October 13, 2024.}
\addtocounter{footnote}{-1}
\endgroup
This Exposé consists of two parts. The first (\S\S1--\SGARefLink{sga3:xxii:subsection:5-5}{5.5}) gathers the
technical results needed for the proofs of the uniqueness and
existence theorems. The second (\S5.6 to the end) will not be used in
that proof; the end of \S5 will be used in particular in Exposé XXVI,
which is devoted to parabolic subgroups; \S6 establishes, in the
scheme-theoretic setting, the classical results concerning the
derived group of a reductive group.

\SGARefTarget{sga3:xxii:section:1}\section*{1. Roots and Coroots.
Split Groups and Root Data}
\addcontentsline{toc}{section}{1. Roots and Coroots. Split Groups and
Root Data}

\SGARefTarget{sga3:xxii:thm:1.1}\paragraph{\SGARefLink{sga3:xxii:thm:1.1}{Theorem 1.1}.}
\textit{Let \(S\) be a scheme, \(G\) a reductive \(S\)-group,
\(T\) a maximal torus of \(G\), and \(\alpha\) a root of \(G\)
relative to \(T\).}

\begin{enumerate}
\item[\SGAInlineRefTarget{sga3:xxii:thm:1.1:item:i}(i)] \textit{There exists a unique \(T\)-equivariant
homomorphism of group schemes}
\[
  \exp_{\alpha}:W(\mathfrak g^{\alpha})\longrightarrow G
\]
\textit{which induces on Lie algebras the canonical morphism
\(\mathfrak g^{\alpha}\to\mathfrak g\). This homomorphism is a closed
immersion. The corresponding homomorphism}
\[
  T\mathbin{\cdot_{\alpha}}W(\mathfrak g^{\alpha})
  \longrightarrow G
\]
\textit{is likewise a closed immersion.}

\textit{If
\(p_{\alpha}:\mathbb G_{a,S}\to G\) is a monomorphism normalized by
\(T\), with multiplier \(\alpha\), there exists a unique
\(X_{\alpha}\in\Gamma(S,\mathfrak g^{\alpha})^{\times}\)\footnote{Editor's
note: The set
\(\Gamma(S,\mathfrak g^{\alpha})^{\times}\) is defined in \SGARefLink{sga3:xix:definition:4-4-1}{Exposé XIX,
4.4.1}.} such that
\(p_{\alpha}(x)=\exp_{\alpha}(xX_{\alpha})\); one has
\(\Lie(p_{\alpha})(1)=X_{\alpha}\), and the preceding two formulas
establish a bijective correspondence between
\(\Gamma(S,\mathfrak g^{\alpha})^{\times}\) and the set of
monomorphisms \(\mathbb G_{a,S}\to G\) normalized by \(T\), with
multiplier \(\alpha\).}

\item[(ii)] \textit{There exists a unique duality, denoted
\((X,Y)\longmapsto XY\),}
\[
  \mathfrak g^{\alpha}\otimes_{\mathcal O_S}\mathfrak g^{-\alpha}
  \xrightarrow{\;\sim\;}\mathcal O_S,
\]
\textit{and a unique homomorphism of group schemes}
\[
  \alpha^*:\mathbb G_{m,S}\longrightarrow T,
\]
\textit{such that formula \textup{(F)} of \SGARefLink{sga3:xx:2.1}{Exposé XX, 2.1} holds. One
has}
\[
  \alpha\circ\alpha^*=2,
  \qquad
  (-\alpha)^*=-\alpha^*,
\]
\textit{and \(\alpha^*\) is given by the formula of \SGARefLink{sga3:xx:2.7}{Exposé XX, 2.7}.}
\end{enumerate}

Indeed, a homomorphism normalized by \(T\), with multiplier \(\alpha\),
necessarily factors through the closed subgroup
\[
  Z_\alpha=\underline{\operatorname{Centr}}_G(T_\alpha)
\]
of \(G\) (cf. \SGARefLink{sga3:xix:paragraph:3-9}{Exposé XIX, 3.9}). Now \((Z_\alpha,T,\alpha)\) is an
elementary \(S\)-system (\SGARefLink{sga3:xx:1.4}{Exposé XX, 1.4}), and one is reduced to the
results of Exposé XX (\SGARefLink{sga3:xxii:def:1.5}{1.5}, \SGARefLink{sga3:xxii:prop:2.1}{2.1}, and \SGARefLink{sga3:xxii:section:5.9}{5.9}).

\SGARefTarget{sga3:xxii:rem:1.2}\paragraph{\SGARefLink{sga3:xxii:rem:1.2}{Remark 1.2}.}
Part \SGARefLink{sga3:xxii:thm:1.1:item:i}{(i)} of \SGARefLink{sga3:xxii:thm:1.1}{Theorem 1.1} remains valid if one assumes only that
\(\alpha\) is a character of \(T\), nontrivial on every fiber. Indeed,
there is then a decomposition
\[
  S=S'\amalg S'',
\]
such that \(\alpha|_{S'}\) is a root of \(G_{S'}\) relative to
\(T_{S'}\), and
\(\mathfrak g^\alpha|_{S''}=0\). If \(S=S'\), one is reduced to \SGARefLink{sga3:xxii:thm:1.1}{1.1};
if \(S=S''\), the result is trivial; the general case follows at once.

\SGARefTarget{sga3:xxii:not:1.3}\paragraph{\SGARefLink{sga3:xxii:not:1.3}{Notations 1.3}.}
As in Exposé XX, let \(U_\alpha\) denote the image of
\(W(\mathfrak g^\alpha)\); it is a closed subgroup of \(G\),
canonically endowed with a vector-group structure. We call it the
\emph{vector group associated with the root \(\alpha\)}. We say that
\(\alpha^*\) is the \emph{coroot associated with \(\alpha\)}. Sections
\[
  X_\alpha\in\Gamma(S,\mathfrak g^\alpha)
  \quad\text{and}\quad
  X_{-\alpha}\in\Gamma(S,\mathfrak g^{-\alpha})
\]
are said to be \emph{paired} if \(X_\alpha X_{-\alpha}=1\). Then
\(X_\alpha\in\Gamma(S,\mathfrak g^\alpha)^\times\), and likewise for
\(X_{-\alpha}\). The corresponding homomorphisms \(p_\alpha\) and
\(p_{-\alpha}\) are contragredient to one another in the sense of
\SGARefLink{sga3:xx:1.5}{Exposé XX, 1.5}\footnote{Editor's note: That is, if
\(\langle\, ,\,\rangle\) denotes the duality between
\(\mathfrak g^\alpha\) and \(\mathfrak g^{-\alpha}\), and if
\(\mathfrak g^\alpha\) is identified with \(U_\alpha\) via
\(p_\alpha(x)=\exp(xX_\alpha)\), and likewise for \(-\alpha\), then
\(\langle p_\alpha(x),p_{-\alpha}(y)\rangle=xy\).}, and one has
\[
  p_\alpha(x)p_{-\alpha}(y)
  =
  p_{-\alpha}\left(\frac{y}{1+xy}\right)
  \alpha^*(1+xy)
  p_\alpha\left(\frac{x}{1+xy}\right).
\]

\SGARefTarget{sga3:xxii:prop:1.4}\paragraph{\SGARefLink{sga3:xxii:prop:1.4}{Proposition 1.4}.}
\textit{Under the hypotheses of \SGARefLink{sga3:xxii:thm:1.1}{1.1}, let
\(w\in\underline{\operatorname{Norm}}_G(T)(S)\). Then}
\[
  \beta=\alpha\circ\operatorname{int}(w)^{-1}
  :T\longrightarrow\mathbb G_{m,S}
\]
\textit{is a root of \(G\) relative to \(T\),
\(\beta^*=\operatorname{int}(w)\circ\alpha^*\) is the corresponding
coroot, and the following diagram is commutative:}

\begin{center}
  % BEGIN INCLUDED SGA3_English_current_buildable_source_20260823/source/inputs/XXII/native_diagrams/exp22/diag_001_root_transport_square.tex
\begin{tikzcd}[column sep=5.2em,row sep=3.2em]
  W(\mathfrak g^\alpha)
    \arrow[r,"\exp_\alpha"]
    \arrow[d,"\operatorname{Ad}(w)"']
  &
  G
    \arrow[d,"\operatorname{int}(w)"]
  \\
  W(\mathfrak g^\beta)
    \arrow[r,"\exp_\beta"]
  &
  G.
\end{tikzcd}

% END INCLUDED SGA3_English_current_buildable_source_20260823/source/inputs/XXII/native_diagrams/exp22/diag_001_root_transport_square.tex

\end{center}

Trivial: transport of structure.

\SGARefTarget{sga3:xxii:def:1.5}\paragraph{\SGARefLink{sga3:xxii:def:1.5}{Definitions 1.5}.}
\textup{(a)} Under the hypotheses of \SGARefLink{sga3:xxii:thm:1.1}{1.1}, let \(s_\alpha\) denote the
automorphism of \(T\) defined by
\[
  s_\alpha(t)=t\mathbin{\cdot}\alpha^*(\alpha(t))^{-1}.
\]
Let \((\, ,\,)\) denote the canonical pairing
\[
  \underline{\operatorname{Hom}}_{S\text{-gr}}
    (\mathbb G_{m,S},T)
  \times
  \underline{\operatorname{Hom}}_{S\text{-gr}}
    (T,\mathbb G_{m,S})
  \longrightarrow
  \underline{\operatorname{Hom}}_{S\text{-gr}}
    (\mathbb G_{m,S},\mathbb G_{m,S})
  =\mathbb Z_S.
\]
Then \(s_\alpha\) acts on
\(\underline{\operatorname{Hom}}_{S\text{-gr}}
(T,\mathbb G_{m,S})\), respectively on
\(\underline{\operatorname{Hom}}_{S\text{-gr}}
(\mathbb G_{m,S},T)\), by the following formulas, where \(\chi\),
respectively \(u\), denotes an arbitrary section of the corresponding
scheme:
\[
  s_\alpha(\chi)=\chi-(\alpha^*,\chi)\alpha,
  \qquad
  s_\alpha(u)=u-(u,\alpha)\alpha^*.
\]
One has
\[
  s_\alpha\circ s_\alpha=\operatorname{id},
  \qquad
  s_{-\alpha}=s_\alpha.
\]

\textup{\SGAInlineRefTarget{sga3:xxii:def:1.5:item:b}(b)} If
\(X\in\Gamma(S,\mathfrak g^\alpha)^\times\), then the inner
automorphism \(w_\alpha(X)\) of \(T\), defined by
\[
  w_\alpha(X)
  =
  \exp_\alpha(X)\exp_{-\alpha}(-X^{-1})\exp_\alpha(X)
\]
(cf. \SGARefLink{sga3:xx:3.1}{Exposé XX, 3.1}), coincides with \(s_\alpha\) (loc. cit.).
\SGARefLink{sga3:xxii:prop:1.4}{Proposition 1.4} therefore gives the following.

\SGARefTarget{sga3:xxii:cor:1.6}\paragraph{\SGARefLink{sga3:xxii:cor:1.6}{Corollary 1.6}.}
\textit{Let \(S\) be a scheme, \(G\) a reductive \(S\)-group, \(T\)
a maximal torus of \(G\), and \(\alpha,\beta\) two roots of \(G\)
relative to \(T\). Then}
\[
  s_\alpha(\beta)=\beta-(\alpha^*,\beta)\alpha
\]
\textit{is a root of \(G\) relative to \(T\), and its corresponding
coroot is}
\[
  s_\alpha(\beta)^*
  =
  s_\alpha(\beta^*)
  =
  \beta^*-(\beta^*,\alpha)\alpha^*.
\]

\SGARefTarget{sga3:xxii:cor:1.7}\paragraph{\SGARefLink{sga3:xxii:cor:1.7}{Corollary 1.7}.}
\textit{Under the preceding conditions,
\(\alpha^*=\beta^*\) implies \(\alpha=\beta\).}

Indeed, if \(\alpha^*=\beta^*\), one has (cf. \SGARefLink{sga3:xx:1.4}{Exposé XXI, 1.4})
\[
  s_\beta(\alpha)=\alpha-2\beta,
  \qquad
  s_\alpha(\beta)=\beta-2\alpha,
\]
and hence immediately
\[
  (s_\beta s_\alpha)^n(\alpha)
  =
  \alpha+2n(\beta-\alpha).
\]
If \(\beta\ne\alpha\), there exists \(s\in S\) such that
\(\alpha_s\ne\beta_s\). The preceding formula would then give
infinitely many distinct roots of \(G_{\bar s}\) relative to
\(T_{\bar s}\), which is impossible.

\SGARefTarget{sga3:xxii:def:1.8.0}\paragraph{\SGARefLink{sga3:xxii:def:1.8.0}{Definitions 1.8.0}.}
\footnote{Editor's note: We added the numbering \SGARefLink{sga3:xxii:def:1.8.0}{1.8.0} to highlight
these definitions.}
If \(u:\mathbb G_{m,S}\to T\) is a homomorphism of groups, we say
that \(u\) is a \emph{coroot} of \(G\) relative to \(T\) if there
exists a root \(\alpha\) of \(G\) relative to \(T\) such that
\(\alpha^*=u\). Consider the functor \(\mathcal R^*\) of coroots of
\(G\) relative to \(T\), defined by
\[
  \mathcal R^*(S')
  =
  \{\text{coroots of \(G_{S'}\) relative to \(T_{S'}\)}\}.
\]
If \(\mathcal R\) is the functor of roots of \(G\) relative to \(T\)
(\SGARefLink{sga3:xix:proposition:3:8}{Exposé XIX, 3.8}), there is a canonical morphism
\(\mathcal R\to\mathcal R^*\). By \SGARefLink{sga3:xxii:cor:1.7}{1.7} and \SGARefLink{sga3:xix:proposition:3-8}{Exposé XIX, 3.8}, one has:

\SGARefTarget{sga3:xxii:cor:1.8}\paragraph{\SGARefLink{sga3:xxii:cor:1.8}{Corollary 1.8}.}
\textit{The canonical morphism
\(\mathcal R\to\mathcal R^*\) is an isomorphism. In particular,
\(\mathcal R^*\) is represented by a twisted finite constant
\(S\)-scheme that is an open and closed subscheme of}
\[
  \underline{\operatorname{Hom}}_{S\text{-gr}}
    (\mathbb G_{m,S},T).
\]

This leads to the following definition.

\SGARefTarget{sga3:xxii:def:1.9}\paragraph{\SGARefLink{sga3:xxii:def:1.9}{Definition 1.9}.}
\textit{Let \(S\) be a scheme and \(T\) an \(S\)-torus. A
\emph{twisted root datum} in \(T\) consists of}
\begin{enumerate}
\item[(i)] \textit{a finite subscheme \(\mathcal R\) of}
  \[
    \underline{\operatorname{Hom}}_{S\text{-gr}}
      (T,\mathbb G_{m,S});
  \]
\item[(ii)] \textit{a finite subscheme \(\mathcal R^*\) of}
  \[
    \underline{\operatorname{Hom}}_{S\text{-gr}}
      (\mathbb G_{m,S},T);
  \]
\item[\SGAInlineRefTarget{sga3:xxii:def:1.9:item:iii}(iii)] \textit{an isomorphism
  \(\mathcal R\xrightarrow{\sim}\mathcal R^*\), denoted
  \(\alpha\mapsto\alpha^*\),}
\end{enumerate}
\textit{satisfying the following conditions:}
\begin{description}
\item[\textup{(DR 1)}]
  \textit{For every \(S'\to S\) and every
  \(\alpha\in\mathcal R(S')\), one has
  \(\alpha\circ\alpha^*=2\).}
\item[\textup{(DR 2)}]
  \textit{For every \(S'\to S\) and all
  \(\alpha,\beta\in\mathcal R(S')\), one has}
  \[
    \alpha-(\beta^*,\alpha)\beta\in\mathcal R(S'),
    \qquad
    \alpha^*-(\alpha^*,\beta)\beta^*\in\mathcal R^*(S').
  \]
\end{description}
\textit{Moreover, the twisted root datum is said to be
\emph{reduced} if, for every nonempty \(S'\to S\) and every
\(\alpha\in\mathcal R(S')\), one has
\(2\alpha\notin\mathcal R(S')\).}

\SGARefTarget{sga3:xxii:prop:1.10}\paragraph{\SGARefLink{sga3:xxii:prop:1.10}{Proposition 1.10}.}
\textit{Let \(S\) be a scheme, \(G\) a reductive \(S\)-group, \(T\)
a maximal torus of \(G\), and let \(\mathcal R\), respectively
\(\mathcal R^*\), be the scheme of roots, respectively coroots, of
\(G\) relative to \(T\). Then \((\mathcal R,\mathcal R^*)\) is a
reduced twisted root datum in \(T\).}

It remains only to verify that this twisted root datum is reduced.
This was done in \SGARefLink{sga3:xix:paragraph:3:10}{Exposé XIX, 3.10}.

\SGARefTarget{sga3:xxii:par:1.11}\paragraph{\SGARefLink{sga3:xxii:par:1.11}{1.11}.}
Let \(T=\mathrm D_S(M)\) be a \emph{trivialized} torus. If \(M^*\)
denotes the abelian group dual to \(M\), there are canonical
isomorphisms (cf. \SGARefLink{sga3:viii:corollary:1:5}{Exposé VIII, 1.5})
\[
  \underline{\operatorname{Hom}}_{S\text{-gr}}
    (T,\mathbb G_{m,S})
  \xrightarrow{\sim} M_S,
  \qquad
  \underline{\operatorname{Hom}}_{S\text{-gr}}
    (\mathbb G_{m,S},T)
  \xrightarrow{\sim} M^*_S,
\]
and hence group isomorphisms
\[
  \operatorname{Hom}_{S\text{-gr}}(T,\mathbb G_{m,S})
  \xrightarrow{\sim}
  \operatorname{Hom}_{\mathrm{loc.const.}}(S,M),
\]
\[
  \operatorname{Hom}_{S\text{-gr}}(\mathbb G_{m,S},T)
  \xrightarrow{\sim}
  \operatorname{Hom}_{\mathrm{loc.const.}}(S,M^*).
\]
A character of \(T\), respectively a group homomorphism
\(\mathbb G_{m,S}\to T\), is said to be \emph{constant} relative to
the given trivialization if the preceding isomorphism carries it to
a constant map from \(S\) to \(M\), respectively to \(M^*\).

\SGARefTarget{sga3:xxii:par:1.12}\paragraph{\SGARefLink{sga3:xxii:par:1.12}{1.12}.}
With the same notation, let \((M,M^*,R,R^*)\) be a root datum
(Exposé XXI). Then \((R_S,R^*_S)\) is a twisted root datum in \(T\).
Conversely, if \((\mathcal R,\mathcal R^*)\) is a twisted root datum
in a torus \(T\), a \emph{splitting} of this root datum means the
data of an ordinary root datum \((M,M^*,R,R^*)\) and an isomorphism
\[
  T\simeq\mathrm D_S(M)
\]
that carries \((\mathcal R,\mathcal R^*)\) to \((R_S,R^*_S)\).

\SGARefTarget{sga3:xxii:def:1.13}\paragraph{\SGARefLink{sga3:xxii:def:1.13}{Definition 1.13}.}
\textit{Let \(S\) be a scheme, \(G\) a reductive \(S\)-group, and
\(T\) a maximal torus of \(G\). A \emph{splitting of \(G\) relative
to \(T\)} consists of}
\begin{enumerate}
\item[\SGAInlineRefTarget{sga3:xxii:def:1.13:item:i}(i)] \textit{an abelian group \(M\) and an isomorphism
  \(T\simeq\mathrm D_S(M)\);}
\item[\SGAInlineRefTarget{sga3:xxii:def:1.13:item:ii}(ii)] \textit{a system of roots \(R\) of \(G\) relative to
  \(T\) \textup{(\SGARefLink{sga3:xix:definition:3-6}{Exposé XIX, 3.6})},}
\end{enumerate}
\textit{satisfying the following two conditions:}
\begin{description}
\item[\textup{(D 1)}]
  \textit{\(S\) is nonempty, and the roots \(\alpha\in R\)
  \textup{(respectively, the corresponding coroots)} identify with
  constant functions from \(S\) to \(M\) \textup{(respectively, to
  \(M^*\))}.}
\item[\textup{(D 2)}]
  \textit{The \(\mathfrak g^\alpha\), \(\alpha\in R\), are free
  \(\mathcal O_S\)-modules.}
\end{description}

We say that \(G\) is \emph{splittable relative to \(T\)} if a
splitting of \(G\) relative to \(T\) exists. A \emph{splitting of
\(G\)} consists of a maximal torus \(T\) of \(G\) together with a
splitting of \(G\) relative to \(T\). We say that \(G\) is
\emph{splittable} if it admits a splitting. A \emph{split
\(S\)-group} is a reductive \(S\)-group equipped with a splitting;
we denote it by a symbol of the form \((G,T,M,R)\), or simply by
\(G\) when there is no risk of confusion.

Condition \textup{(D 1)} canonically identifies \(R\), respectively
\(R^*\), with a subset of \(M\), respectively of \(M^*\).

\SGARefTarget{sga3:xxii:prop:1.14}\paragraph{\SGARefLink{sga3:xxii:prop:1.14}{Proposition 1.14}.}
\textit{Let \(S\) be a nonempty scheme and let \((G,T,M,R)\) be a
split \(S\)-group. Then}
\[
  \mathcal R(G,T,M,R)=(M,M^*,R,R^*)
\]
\textit{is a reduced root datum \textup{(\SGARefLink{sga3:xxi:subsection:1-1}{Exposé XXI, 1.1} and
\SGARefLink{sga3:XXI:2.1.3}{2.1.3})}; it is a splitting of the twisted root datum of \SGARefLink{sga3:xxii:prop:1.10}{1.10}.}

This follows trivially from \SGARefLink{sga3:xxii:prop:1.10}{1.10} and \SGARefLink{sga3:xix:lemma:3:7}{Exposé XIX, 3.7}.

For brevity we shall sometimes write
\(\mathcal R(G,T,M,R)=\mathcal R(G)\). We shall systematically use
the notation \(V\), \(V(\mathcal R)\), \(W,\ldots\) of Exposé XXI.

\SGARefTarget{sga3:xxii:rem:1.15}\paragraph{\SGARefLink{sga3:xxii:rem:1.15}{Remark 1.15}.}
\textup{(a)} If \(S\) is nonempty and connected, respectively if
\(\operatorname{Pic}(S)=0\), condition \textup{(D 1)}, respectively
condition \textup{(D 2)}, is automatically satisfied.

\SGAInlineRefTarget{sga3:xxii:rem:1.15:item:b}\textup{\SGARefLink{sga3:xxii:def:1.5:item:b}{(b)}} If \((G,T,M,R)\) is a split \(S\)-group, then for every
nonempty \(S'\to S\), \((G_{S'},T_{S'},M,R)\) is a split
\(S'\)-group and
\[
  \mathcal R(G,T,M,R)
  =
  \mathcal R(G_{S'},T_{S'},M,R).
\]

\SGARefTarget{sga3:xxii:par:1.16}\paragraph{\SGARefLink{sga3:xxii:par:1.16}{1.16}.}
Let \(T=\mathrm D_S(M)\) be a trivialized torus. The Lie algebra
\(\mathfrak t\) of \(T\) identifies canonically \textup{(\SGARefLink{sga3:ii:corollary:5-1-1}{Exposé II,
5.1.1})} with
\[
  \mathfrak t\simeq M^*\otimes\mathcal O_S.
\]
For every homomorphism of groups \(u:T\to\mathbb G_{m,S}\),
\(\Lie(u)\) is a linear form
\[
  \Lie(u):
  \mathfrak t\longrightarrow
  \mathcal O_S=\Lie(\mathbb G_{m,S}/S).
\]
In particular, if \(u\) is defined by an element \(\alpha\in M\),
then \(\Lie(u)\) is the linear form \(\bar\alpha\) on
\(M^*\otimes\mathcal O_S\) defined by
\[
  \bar\alpha(m\otimes x)=(m,\alpha)x.
\]

Symmetrically, for every homomorphism of groups
\(h:\mathbb G_{m,S}\to T\), \(\Lie(h)\) is an
\(\mathcal O_S\)-linear morphism
\[
  \mathcal O_S=\Lie(\mathbb G_{m,S}/S)
  \longrightarrow\mathfrak t,
\]
canonically determined by the section
\[
  H=\Lie(h)(1)\in\Gamma(S,\mathfrak t).
\]
In particular, if \(h\) is defined by an element \(m\in M^*\), then
\[
  H=\Lie(h)(1)=m\otimes 1.
\]
Comparing the two definitions, we find in particular
\[
  \bar\alpha(H)=(h,\alpha)\cdot 1
  \in\Gamma(S,\mathcal O_S).
\]

\SGARefTarget{sga3:xxii:par:1.17}\paragraph{\SGARefLink{sga3:xxii:par:1.17}{1.17}.}
These definitions apply in particular when \(T\) is the maximal
torus of a split group. Every root \(\alpha\in R\) defines an
\emph{infinitesimal root}
\[
  \bar\alpha\in
  \operatorname{Hom}_{\mathcal O_S}(\mathfrak t,\mathcal O_S),
  \qquad
  \bar\alpha(m\otimes x)=(m,\alpha)x.
\]
Every coroot \(\alpha^*\), \(\alpha\in R\), defines an
\emph{infinitesimal coroot}
\[
  H_\alpha\in\Gamma(S,\mathfrak t),
  \qquad
  H_\alpha=\alpha^*\otimes 1.
\]
For \(\alpha,\beta\in R\), we have
\[
  \bar\alpha(H_\beta)=(\beta^*,\alpha)\cdot 1,
\]
and in particular
\[
  \bar\alpha(H_\alpha)=2.
\]
Consequently, if \(2\) is invertible on \(S\), then
\(\bar\alpha\) and \(H_\alpha\) are nonzero on every fiber.

\SGARefTarget{sga3:xxii:section:2}\section*{2. Existence of a
Splitting. Type of a Reductive Group}
\addcontentsline{toc}{section}{2. Existence of a Splitting. Type of
a Reductive Group}

\SGARefTarget{sga3:xxii:prop:2.1}\paragraph{\SGARefLink{sga3:xxii:prop:2.1}{Proposition 2.1}.}
\textit{Let \(S\) be a scheme, \(G\) a reductive \(S\)-group, and
\(T\) a maximal torus of \(G\). Suppose that \(T\) is split. Then
\(G\) is locally splittable relative to \(T\): for every
\(s_0\in S\), there exists an open neighborhood \(U\) of \(s_0\)
such that the \(U\)-group \(G_U\) is splittable relative to
\(T_U\).}

Indeed, write \(T\simeq\mathrm D_S(M)\) and
\[
  \mathfrak g=\coprod_{m\in M}\mathfrak g^m.
\]
Let
\[
  R=\{m\in M\mid m\ne0,\ \mathfrak g^m(s_0)\ne0\}.
\]
After restricting \(S\) and replacing it by an open neighborhood
of \(s_0\), we may suppose that the \(\mathfrak g^\alpha\),
\(\alpha\in R\), are free and that the \(\mathfrak g^m\), for
\(m\ne0\) and \(m\notin R\), vanish. We then have
\[
  \mathfrak g
  =
  \mathfrak t\oplus
  \coprod_{\alpha\in R}\mathfrak g^\alpha,
\]
where the \(\mathfrak g^\alpha\) are free of rank \(1\). It follows
that \(R\) is a system of roots of \(G\) relative to \(T\)
\textup{(\SGARefLink{sga3:xix:definition:3-6}{Exposé XIX, 3.6})}. The coroots \(\alpha^*\) corresponding
to the \(\alpha\in R\) then identify with locally constant
functions on \(S\) with values in \(M^*\). Shrinking \(S\) once
more, we may suppose that they are constant, and the proof is
complete.

Observe that the proof immediately gives:

\SGARefTarget{sga3:xxii:prop:2.2}\paragraph{\SGARefLink{sga3:xxii:prop:2.2}{Proposition 2.2}.}
\textit{Let \(S\) be a nonempty connected scheme such that
\(\operatorname{Pic}(S)=0\), for example
\(\operatorname{Spec}(\mathbb Z)\) or a local scheme
\textup{(in particular, the spectrum of a field)}. If \(G\) is a
reductive \(S\)-group possessing a split maximal torus \(T\), then
\(G\) is splittable relative to \(T\).}

We deduce immediately from \SGARefLink{sga3:xxii:prop:2.1}{2.1} and from the fact that a reductive
group locally possesses maximal tori for the étale topology
\textup{(\SGARefLink{sga3:xix:theorem:2:5}{Exposé XIX, 2.5})}:

\SGARefTarget{sga3:xxii:cor:2.3}\paragraph{\SGARefLink{sga3:xxii:cor:2.3}{Corollary 2.3}.}
\textit{Let \(S\) be a scheme and \(G\) a reductive \(S\)-group
\textup{(respectively, let \(T\) also be a maximal torus of \(G\))}.
Then \(G\) is locally splittable
\textup{(respectively, locally splittable relative to \(T\))}
for the étale topology on \(S\).}

\SGARefTarget{sga3:xxii:cor:2.4}\paragraph{\SGARefLink{sga3:xxii:cor:2.4}{Corollary 2.4}.}
\textit{Let \(k\) be a field and \(G\) a reductive \(k\)-group.
There exists a finite separable extension \(K/k\) such that \(G_K\)
is splittable.}

\SGARefTarget{sga3:xxii:rem:2.5}\paragraph{\SGARefLink{sga3:xxii:rem:2.5}{Remark 2.5}.}
Using \SGARefLink{sga3:xxii:prop:2.1}{2.1} and \SGARefLink{sga3:xix:remark:2:9}{Exposé XIX, Remark 2.9}, one immediately proves the
following result: let \(G=(G,T,M,R)\) be a split \(S\)-group. There
exists an open covering \((U_i)\) of \(S\) such that each split group
\(G_{U_i}\) is obtained by base change from a split group over a
noetherian ring \textup{(and in fact over a finitely generated
\(\mathbb Z\)-algebra)}. We shall moreover prove that every split
group over \(S\) is already obtained from a split
\(\mathbb Z\)-group \textup{(Exposé XXV)}.

\SGARefTarget{sga3:xxii:par:2.6}\paragraph{\SGARefLink{sga3:xxii:par:2.6}{2.6}.}
Let \(k\) be an algebraically closed field and \(G\) a reductive
\(k\)-group. We know \textup{(for example by \SGARefLink{sga3:xxii:cor:2.4}{2.4})} that splittings of
\(G\) exist. Let \((G,T,M,R)\) and \((G,T',M',R')\) be two splittings
of \(G\). Then the root data
\[
  \mathcal R(G,T,M,R)
  \qquad\text{and}\qquad
  \mathcal R(G,T',M',R')
\]
are isomorphic.

Indeed, we first see that it is enough to reduce to the case
\(T=T'\): there exists \(g\in G(k)\) such that
\(T'=\operatorname{int}(g)T\), and one readily verifies that
transporting a splitting by an automorphism of \(G\) produces a root
datum isomorphic to the original one. But since
\(S=\operatorname{Spec}(k)\) is connected, the isomorphism
\[
  \mathrm D_k(M)\xrightarrow{\sim}T
  \xrightarrow{\sim}\mathrm D_k(M')
\]
comes from a unique isomorphism \(M\simeq M'\). For the same reason,
there exists at most one root system of \(G\) relative to \(T\).

\SGARefTarget{sga3:xxii:def:2.6.1}\paragraph{\SGARefLink{sga3:xxii:def:2.6.1}{Definition 2.6.1}.}
\footnote{Editor's note: The number \SGARefLink{sga3:xxii:def:2.6.1}{2.6.1} was added for later
references.}
If \(G\) is a reductive \(k\)-group, where \(k\) is algebraically
closed, the \emph{type of \(G\)} is the isomorphism class of the root
datum defined by any splitting of \(G\). If \(G\) is a torus of type
\(M\) in the sense of \SGARefLink{sga3:ix:definition:1:4}{Exposé IX, 1.4}, then the type of \(G\) as a
reductive group is given by the trivial root datum
\[
  (M,M^*,\varnothing,\varnothing).
\]

By \SGARefLink{sga3:xxii:rem:1.15:item:b}{1.15(b)}\footnote{Editor's note: We have corrected the original,
which referred to \SGARefLink{sga3:xxii:par:1.17}{1.17}.}, the type is invariant under an
\textup{(algebraically closed)} extension of the base field.

\SGARefTarget{sga3:xxii:def:2.7}\paragraph{\SGARefLink{sga3:xxii:def:2.7}{Definition 2.7}.}
If \(G\) is a reductive \(S\)-group and \(s\in S\), the \emph{type of
\(G\) at \(s\)} is the type of the reductive
\(\bar s\)-group \(G_{\bar s}\).

For every \(S'\to S\) and every \(s'\in S'\) mapping to \(s\in S\),
the type of \(G_{S'}\) at \(s'\) is equal to the type of \(G\) at
\(s\).

If \(G\) is splittable and \((G,T,M,R)\) is a splitting of \(G\),
then the type of \(G\) at \(s\) is the isomorphism class of
\(\mathcal R(G,T,M,R)\), by \SGARefLink{sga3:xxii:rem:1.15:item:b}{1.15(b)}\footnotemark[5]. The following
then follows immediately from \SGARefLink{sga3:xxii:cor:2.3}{2.3}:

\SGARefTarget{sga3:xxii:prop:2.8}\paragraph{\SGARefLink{sga3:xxii:prop:2.8}{Proposition 2.8}.}
\textit{Let \(G\) be a reductive \(S\)-group, where
\(S\ne\varnothing\). The function}
\[
  s\longmapsto\text{\textit{the type of \(G\) at \(s\)}}
\]
\textit{is locally constant on \(S\). In particular, there exists a
partition of \(S\) into nonempty open subschemes on each of which
\(G\) has constant type. More precisely, let \(E\) be the set of
types of the fibers of \(G\), and, for every \(t\in E\), let \(S_t\)
be the set of points \(s\in S\) where \(G\) has type \(t\). Then
\((S_t)_{t\in E}\) is a partition of \(S\), and each \(S_t\) is open
and closed \textup{(and nonempty)}.}

\SGARefTarget{sga3:xxii:section:3}\section*{3. The Weyl Group}
\addcontentsline{toc}{section}{3. The Weyl Group}

\SGARefTarget{sga3:xxii:par:3.1}\paragraph{\SGARefLink{sga3:xxii:par:3.1}{3.1}.}
Let \(S\) be a scheme, \(G\) a reductive \(S\)-group, and \(T\) a
maximal torus of \(G\). Then
\[
  W_G(T)=\underline{\operatorname{Norm}}_G(T)/T
\]
is a finite étale \(S\)-group \textup{(\SGARefLink{sga3:xix:theorem:2-5}{Exposé XIX, 2.5})}. The
morphism \(n\mapsto\operatorname{int}(n)\) induces, upon passage to
the quotient, a canonical monomorphism \textup{(which is moreover an
open immersion)}
\[
  W_G(T)\longrightarrow
  \underline{\operatorname{Aut}}_{S\text{-gr.}}(T).
\]

\SGARefTarget{sga3:xxii:par:3.2}\paragraph{\SGARefLink{sga3:xxii:par:3.2}{3.2}.}
Suppose now that \(G\) is splittable relative to \(T\). Choose a
splitting, say \((G,T,M,R)\). We then have a canonical isomorphism
\textup{(\SGARefLink{sga3:viii:corollary:1:5}{Exposé VIII, 1.5})}
\[
  \underline{\operatorname{Aut}}_{S\text{-gr.}}(T)
  \simeq
  \bigl(\operatorname{Aut}_{\mathrm{gr.}}(M)\bigr)_S.
\]
In particular, if \(W\) is the Weyl group of the root datum
\(\mathcal R(G)\) \textup{(\SGARefLink{sga3:XXI:1.1.8}{Exposé XXI, 1.1.8})}, we have a
monomorphism
\[
  W_S\longrightarrow
  \underline{\operatorname{Aut}}_{S\text{-gr.}}(T).
\]

\SGARefTarget{sga3:xxii:par:3.3}\paragraph{\SGARefLink{sga3:xxii:par:3.3}{3.3}.}
For every root \(\alpha\in R\), the reflection \(s_\alpha\in W\)
acts on \(M\) by
\[
  s_\alpha(x)=x-(\alpha^*,x)\alpha,
\]
and hence acts on \(T\), by the preceding morphism, according to
\[
  s_\alpha(t)=t\cdot\alpha^*\bigl(\alpha(t)\bigr)^{-1}.
\]
On the other hand, since \(\mathfrak g^\alpha\) is assumed to be
free, there exists
\[
  X\in\Gamma(S,\mathfrak g^\alpha)^\times.
\]
Consider
\[
  w_\alpha(X)\in
  \underline{\operatorname{Norm}}_G(T)(S)
  \qquad\textup{(Exposé XX, 3.1)}.
\]
We have \textup{(loc.\ cit.)}
\[
  \operatorname{int}\bigl(w_\alpha(X)\bigr)(t)=s_\alpha(t).
\]
Since \(W\) is generated by the \(s_\alpha\), \(\alpha\in R\), it
follows from the preceding remarks that, if \(W\) and
\(\underline{\operatorname{Norm}}_G(T)(S)/T(S)\) are regarded as
groups of automorphisms of \(T\), then
\[
  W
  \subset
  \underline{\operatorname{Norm}}_G(T)(S)/T(S)
  \subset
  W_G(T)(S).
\]

By definition of the constant group \(W_S\) associated with \(W\)
\textup{(cf.\ \SGARefLink{sga3:i:subsection:1:8}{Exposé I, 1.8})}, we therefore have a commutative
diagram:

\begin{center}
  % BEGIN INCLUDED SGA3_English_current_buildable_source_20260823/source/inputs/XXII/native_diagrams/exp22/diag_002_weyl_commutative_triangle.tex
\begin{tikzcd}[column sep=5.0em,row sep=3.0em]
  W_S
    \arrow[rr]
    \arrow[dr]
  &&
  W_G(T)
    \arrow[dl]
  \\
  &
  \underline{\operatorname{Aut}}_{S\text{-gr.}}(T).
  &
\end{tikzcd}

% END INCLUDED SGA3_English_current_buildable_source_20260823/source/inputs/XXII/native_diagrams/exp22/diag_002_weyl_commutative_triangle.tex

\end{center}

\SGARefTarget{sga3:xxii:prop:3.4}\paragraph{\SGARefLink{sga3:xxii:prop:3.4}{Proposition 3.4}.}
\textit{Let \(S\) be a scheme, let \((G,T,M,R)\) be a split
\(S\)-group, and let \(W\) be the Weyl group of the root datum
\(\mathcal R(G)\). Then the canonical monomorphism}
\[
  W_S\longrightarrow
  W_G(T)=\underline{\operatorname{Norm}}_G(T)/T
\]
\textit{is an isomorphism.}

Indeed, these are étale groups over \(S\); it therefore suffices to
verify that, for every \(s\in S\), the morphism
\[
  W_S(\bar s)\longrightarrow W_G(T)(\bar s)
\]
is an isomorphism.\footnote{Editor's note: Indeed, since \(W_S\) and
\(W_G(T)=\underline{\operatorname{Norm}}_G(T)/T\) are étale over
\(S\), the morphism
\(f:W_S\to W_G(T)=\underline{\operatorname{Norm}}_G(T)/T\) is étale
\textup{(EGA IV\textsubscript{4}, 17.3.4)}. If, moreover, each
\(f_s\) is an isomorphism, then, by \textup{loc.\ cit.}, 17.9.1,
\(f\) is a surjective open immersion, hence an isomorphism.}
The latter assertion follows, for example, from \emph{Bible},
\S11.3, Theorem 2.

\SGARefTarget{sga3:xxii:rem:3.5}\paragraph{\SGARefLink{sga3:xxii:rem:3.5}{Remark 3.5}.}
Using \SGARefLink{sga3:xxii:cor:2.3}{2.3}, the preceding proposition gives a new proof of the fact
that the Weyl group of a maximal torus of a reductive \(S\)-group
\(G\) is finite over \(S\) \textup{(Exposé XIX, 2.5(ii))}.
\footnote{Editor's note: Indeed, let \(T\) be a maximal torus of
\(G\). The fact that \(W_G(T)\) is finite over \(S\) is local for
the fpqc topology \textup{(EGA IV\textsubscript{2}, 2.7.1)}, hence
\emph{a fortiori} for the étale topology. By 2.3, we may therefore
suppose that \(G\) is split relative to \(T\), in which case the
assertion follows from \SGARefLink{sga3:xxii:prop:3.4}{3.4}.}

\SGARefTarget{sga3:xxii:par:3.6}\paragraph{\SGARefLink{sga3:xxii:par:3.6}{3.6}.}
Under the hypotheses of \SGARefLink{sga3:xxii:par:3.1}{3.1}, for every \(w\in W_G(T)(S)\), let
\(N_w\)\footnote{Editor's note: We have replaced \(Q_w\) by \(N_w\),
just as in \SGARefLink{sga3:xx:3.0}{Exposé XX, 3.0}, \(Q\) was replaced by \(N^\times\).}
denote the fiber product in the following diagram:

\begin{center}
  % BEGIN INCLUDED SGA3_English_current_buildable_source_20260823/source/inputs/XXII/native_diagrams/exp22/diag_003_Nw_fiber_square.tex
\begin{tikzcd}[column sep=5.2em,row sep=3.2em]
  N_w
    \arrow[r]
    \arrow[d]
  &
  \underline{\operatorname{Norm}}_G(T)
    \arrow[d]
  \\
  S
    \arrow[r,"w"]
  &
  W_G(T).
\end{tikzcd}

% END INCLUDED SGA3_English_current_buildable_source_20260823/source/inputs/XXII/native_diagrams/exp22/diag_003_Nw_fiber_square.tex

\end{center}

It is an open and closed subscheme of
\(\underline{\operatorname{Norm}}_G(T)\), and is a principal
homogeneous bundle under \(T\) on the left \textup{(respectively,
on the right)} for the law
\[
  (t,q)\longmapsto tq
  \qquad
  \textup{(respectively, \((q,t)\longmapsto qt\))}.
\]
If \(n\in N_w(S)\), then
\[
  N_{ww'}=n\mathbin{\cdot}N_{w'},
  \qquad
  N_{w'w}=N_{w'}\mathbin{\cdot}n.
\]

\SGARefTarget{sga3:xxii:par:3.7}\paragraph{\SGARefLink{sga3:xxii:par:3.7}{3.7}.}
In particular, if \(\alpha\) is a root of \(G\) with respect to
\(T\), then \(N_{s_\alpha}\) is none other than what was denoted by
\(N^\times\) in \SGARefLink{sga3:xx:3.0}{Exposé XX, 3.0}. If \(\mathfrak g^\alpha\) is free
over \(S\), we therefore have
\[
  N_{s_\alpha}(S)\ne\varnothing.
\]
By \SGARefLink{sga3:xxii:prop:3.4}{3.4} and condition \textup{(D\,2)} of the splitting, we deduce:

\SGARefTarget{sga3:xxii:cor:3.8}\paragraph{\SGARefLink{sga3:xxii:cor:3.8}{Corollary 3.8}.}
\textit{Under the hypotheses of \SGARefLink{sga3:xxii:prop:3.4}{3.4}, the morphism}
\[
  \underline{\operatorname{Norm}}_G(T)(S)
  \longrightarrow
  W_G(T)(S)
  =
  \operatorname{Hom}_{\mathrm{loc.const.}}(S,W)
\]
\textit{is surjective. In particular, for every \(w\in W\), there
exists \(n_w\in\underline{\operatorname{Norm}}_G(T)(S)\) such that}
\[
  \operatorname{int}(n_w)|_T=w.
\]

\SGARefTarget{sga3:xxii:section:4}
\section*{4. Homomorphisms of Split Groups}
\addcontentsline{toc}{section}{4. Homomorphisms of Split Groups}

\SGARefTarget{sga3:xxii:subsection:4.1}
\subsection*{4.1. The ``Big Cell''}
\addcontentsline{toc}{subsection}{\SGARefLink{sga3:xxii:subsection:4.1}{4.1}. The ``Big Cell''}

\SGARefTarget{sga3:xxii:par:4.1.1}\paragraph{\SGARefLink{sga3:xxii:par:4.1.1}{4.1.1}.}
Let \((G,T,M,R)\) be a split reductive \(S\)-group. Choose a system
of positive roots \(R_+\) of the root datum \(\mathcal R(G)\)
\textup{(\SGARefLink{sga3:XXI:3.2.1}{Exposé XXI, 3.2.1})}. Put \(R_-=-R_+\).

Choose a total ordering on \(R_+\) \textup{(respectively, \(R_-\))}
and consider the morphism induced by multiplication in \(G\):
\[
  u:
  \left(\prod_{\alpha\in R_-}U_\alpha\right)
  \times_S T
  \times_S
  \left(\prod_{\alpha\in R_+}U_\alpha\right)
  \longrightarrow G.
\]
This is an open immersion. Indeed, since both sides are flat and of
finite presentation over \(S\), it suffices to verify this on every
geometric fiber \textup{(SGA 1, I, 5.7 and VIII, 5.4)}. We are thus
reduced to the case in which \(S\) is the spectrum of an
algebraically closed field. But by \emph{Bible}, \S13.4,
Corollary 2 to Theorem 3, \(u\) is radicial and dominant. Since the
tangent map of \(u\) at the origin is an isomorphism
\textup{(definition of a root system)}, \(u\) is birational. As
\(G\) is normal, Zariski's ``Main Theorem''
\textup{(EGA III\textsubscript{1}, 4.4.9)} applies, and \(u\) is an
open immersion.

Let us show that the image \(\Omega\) of this open immersion is
independent of the ordering chosen on \(R_+\)
\textup{(respectively, \(R_-\))}. Since this amounts to comparing
open subsets of \(G\), we are reduced to proving that they have the
same geometric points, and may therefore again suppose that \(S\)
is the spectrum of an algebraically closed field. The assertion is
then precisely \emph{Bible}, \S13, Proposition 1(c) and
Theorem 1(a).

We have therefore proved:

\SGARefTarget{sga3:xxii:prop:4.1.2}\paragraph{\SGARefLink{sga3:xxii:prop:4.1.2}{Proposition 4.1.2}.}
\textit{Let \((G,T,M,R)\) be a split \(S\)-group, and let \(R_+\)
be a system of positive roots of \(R\). There exists an open subset
\(\Omega_{R_+}\) of \(G\) such that, for every total ordering on
\(R_+\) \textup{(respectively, \(R_-\))}, the morphism induced by
multiplication in \(G\)}
\[
  \left(\prod_{\alpha\in R_-}U_\alpha\right)
  \times_S T
  \times_S
  \left(\prod_{\alpha\in R_+}U_\alpha\right)
  \longrightarrow G
\]
\textit{is an open immersion with image \(\Omega_{R_+}\).}

\SGARefTarget{sga3:xxii:rem:4.1.3}\paragraph{\SGARefLink{sga3:xxii:rem:4.1.3}{Remark 4.1.3}.}
\SGARefLink{sga3:xxii:prop:4.1.2}{Proposition 4.1.2} may be restated as follows: for every
\(\alpha\in R\), choose an isomorphism of vector groups
\[
  p_\alpha:\mathbb G_{a,S}\xrightarrow{\;\sim\;}U_\alpha
  \qquad\textup{(cf.\ 1.19)}.
\]
Then the morphism \textup{(where
\(N=\operatorname{Card}(R_+)=\operatorname{Card}(R_-)\))}
\[
  \mathbb G_{a,S}^{N}
  \times_S T
  \times_S
  \mathbb G_{a,S}^{N}
  \longrightarrow G
\]
defined on points by
\[
  \bigl((x_\alpha)_{\alpha\in R_-},t,
        (x_\alpha)_{\alpha\in R_+}\bigr)
  \longmapsto
  \prod_{\alpha\in R_-}p_\alpha(x_\alpha)\mathbin{\cdot}t
  \mathbin{\cdot}\prod_{\alpha\in R_+}p_\alpha(x_\alpha)
\]
is an open immersion whose image depends only on \(R_+\), and not
on the choice of the \(p_\alpha\) or of the orderings on \(R_+\)
and \(R_-\).

\SGARefTarget{sga3:xxii:notation:4.1.4}\paragraph{\SGARefLink{sga3:xxii:notation:4.1.4}{Notation 4.1.4}.}
We write\footnote{Editor's note: It is called the ``big cell''
corresponding to \(R_+\).}
\[
  \Omega_{R_+}
  =
  \prod_{\alpha\in R_-}U_\alpha
  \mathbin{\cdot}T\mathbin{\cdot}
  \prod_{\alpha\in R_+}U_\alpha.
\]

\SGARefTarget{sga3:xxii:prop:4.1.5}\paragraph{\SGARefLink{sga3:xxii:prop:4.1.5}{Proposition 4.1.5}.}
\textit{The scheme \(\Omega_{R_+}\) is of finite presentation over
\(S\) \textup{(hence quasi-compact in \(G\))} and is universally
schematically dense in \(G\) relative to \(S\)
\textup{(cf.\ \SGARefLink{sga3:xviii:section:1}{Exposé XVIII, 1})}.}

The first assertion is trivial. Thus,\footnote{Editor's note: We
have expanded the original in what follows.}
\(\Omega_{R_+}\) is flat and of finite presentation over \(S\), and
contains the identity section; it therefore meets every fiber of
\(G\) in a nonempty, hence dense, open subset. The second assertion
thus follows from Exposé XVIII, 1.3.\footnote{Editor's note: See also
EGA IV\textsubscript{3}, 11.10.10.}

\SGARefTarget{sga3:xxii:cor:4.1.6}\paragraph{\SGARefLink{sga3:xxii:cor:4.1.6}{Corollary 4.1.6}.}
\textit{Let \((G,T,M,R)\) be a split reductive \(S\)-group. Then}
\[
  \underline{\operatorname{Centr}}(G)
  =
  \bigcap_{\alpha\in R}\operatorname{Ker}(\alpha).
\]
\textit{Consequently,
\(\underline{\operatorname{Centr}}(G)\) is representable by a
closed diagonalizable subgroup of \(G\).}
\footnote{Editor's note: We have added the preceding sentence.}

The second assertion follows immediately from the first. To prove
the first, one may invoke \SGARefLink{sga3:xii:4.8}{Exposé XII, 4.8} and \SGARefLink{sga3:xii:4.11}{4.11}; alternatively,
one may argue directly as follows.\footnote{Editor's note: We have
expanded the original in what follows.}
Let \(S'\to S\). If \(t\in T(S')\) and \(\alpha(t)=1\) for every
\(\alpha\in R\), then \(\operatorname{int}(t)\) induces the identity
on \(T_{S'}\) and on every \((U_\alpha)_{S'}\), \(\alpha\in R\),
and therefore also on \((\Omega_{R_+})_{S'}\). By schematic density
it consequently induces the identity on \(G_{S'}\), whence
\[
  t\in\underline{\operatorname{Centr}}(G)(S').
\]
Conversely, since
\[
  \underline{\operatorname{Centr}}_{G_{S'}}(T_{S'})
  =
  T_{S'}
  \qquad\textup{(cf.\ Exposé XIX, 2.8)},
\]
if \(g\in G(S')\) centralizes \(T_{S'}\) and the
\((U_\alpha)_{S'}\), then \(g\) is a section of \(T_{S'}\) on
which every \(\alpha\in R\) vanishes.

\SGARefTarget{sga3:xxii:cor:4.1.7}\paragraph{\SGARefLink{sga3:xxii:cor:4.1.7}{Corollary 4.1.7}.}
\textit{Let \(S\) be a scheme and \(G\) a reductive \(S\)-group.
Then the center of \(G\) is representable by a closed subgroup of
\(G\) of multiplicative type. It is also ``the intersection of the
maximal tori of \(G\)'' in the following sense: for every
\(S'\to S\), \(\underline{\operatorname{Centr}}(G)(S')\) is the set
of \(g\in G(S')\) such that, for every \(S''\to S'\), the inverse
image of \(g\) in \(G(S'')\) belongs to \(T(S'')\) for every maximal
torus \(T\) of \(G_{S''}\).}

In view of \SGARefLink{sga3:xxii:cor:2.3}{2.3}, the first assertion follows from \SGARefLink{sga3:xxii:cor:4.1.6}{4.1.6} by
descent.\footnote{Editor's note: Indeed, representability of the
center by a closed subscheme of \(G\) is local for the fpqc topology
\textup{(SGA 1, VIII, 5.2 and 5.4)}, hence \emph{a fortiori} for the
étale topology; the same is true of the property of being of
multiplicative type.}
Let us prove the second assertion. Let \(H\) denote ``the
intersection of the maximal tori of \(G\)'' in the preceding sense.
Clearly,\footnote{Editor's note: We have expanded the original in
what follows.}
\[
  \underline{\operatorname{Centr}}(G)\subset H.
\]
By descent, it is therefore enough to prove
\(\underline{\operatorname{Centr}}(G)=H\) when \(G\) is split.
Since \(H\) is contained in the intersection of the maximal tori of
\(G\) in the usual sense, this follows from the following remark:
if \((G,T,M,R)\) is a splitting, \(\alpha\in R\), and
\(X\in\Gamma(S,\mathfrak g^\alpha)^\times\), then
\[
  \operatorname{int}\bigl(\exp_\alpha(X)\bigr)(T)\cap T
  =
  \operatorname{Ker}(\alpha),
\]
as a trivial calculation shows. See also \SGARefLink{sga3:xii:8.6}{Exposé XII, 8.6} and \SGARefLink{sga3:xii:8.8}{8.8}
for a more general statement.

\SGARefTarget{sga3:xxii:rem:4.1.8}\paragraph{\SGARefLink{sga3:xxii:rem:4.1.8}{Remark 4.1.8}.}
In what follows, in the split case we shall systematically identify
\(T\) with \(D_S(M)\). Then
\[
  \underline{\operatorname{Centr}}(G)
  =
  D_S\bigl(M/\Gamma_0(R)\bigr),
\]
where \(\Gamma_0(R)\) is the subgroup of \(M\) generated by \(R\)
\textup{(cf.\ \SGARefLink{sga3:XXI:1.1.6}{Exposé XXI, 1.1.6})}.
\par\noindent
If
\(\{\alpha_1,\alpha_2,\ldots,\alpha_n\}\) is a system of simple
roots of \(R\), then one immediately has
\textup{(cf.\ \SGARefLink{sga3:xx:1.19}{Exposé XX, 1.19})}
\[
  \underline{\operatorname{Centr}}(G)
  =
  \bigcap_i\operatorname{Ker}(\alpha_i)
  =
  \bigcap_i\underline{\operatorname{Centr}}(Z_{\alpha_i}).
\]

\SGARefTarget{sga3:xxii:prop:4.1.9}\paragraph{\SGARefLink{sga3:xxii:prop:4.1.9}{Proposition 4.1.9}.}
\textit{Let \(S\) be a scheme, let \((G,T,M,R)\) be a split
\(S\)-group, let \(Q\) be an \(S\)-torus, let \(\alpha_0\) be a
character of \(Q\), let \(\mathcal L\) be an invertible
\(\mathcal O_S\)-module, and let}
\[
  f:Q\longrightarrow T,
  \qquad
  p:W(\mathcal L)\longrightarrow G
\]
\textit{be group homomorphisms satisfying the following
set-theoretic relation:}
\[
  p\bigl(\alpha_0(q)x\bigr)
  =
  \operatorname{int}\bigl(f(q)\bigr)\mathbin{\cdot}p(x)
\]
\textit{for every \(q\in Q(S')\), \(x\in W(\mathcal L)(S')\), and
\(S'\to S\). Suppose that \(f\) separates the elements of \(R\) in
the following sense: if \(\alpha,\alpha'\in R\) and
\(m,m'\in\mathbb Z\), then}\footnote{Editor's note: This is
equivalent to the following hypothesis: if
\(\alpha,\alpha'\in R\), \(m,m'\in\mathbb Z\), and
\[
  (m\alpha\circ f)_{\bar s}
  =
  (m'\alpha'\circ f)_{\bar s}
\]
for every geometric point \(\bar s\) of \(S\), then
\(m\alpha=m'\alpha'\). In particular, this separation hypothesis is
stable under base change.}
\[
  m\alpha\circ f=m'\alpha'\circ f
  \quad\Longrightarrow\quad
  m\alpha=m'\alpha'.
\]
\textit{Finally, let \(s\in S\) be such that
\((\alpha_0)_{\bar s}\ne e\) and \(p_{\bar s}\ne e\).}

\textit{There then exist an open subset \(U\) of \(S\) containing
\(s\), an integer \(q>0\) such that \(x\mapsto x^q\) is an
endomorphism of \(\mathbb G_{a,U}\), a root \(\alpha\in R\), and an
isomorphism of \(\mathcal O_U\)-modules}
\[
  h:(\mathcal L|_U)^{\otimes q}
  \xrightarrow{\;\sim\;}
  \mathfrak g^\alpha|_U
\]
\textit{such that}
\[
  \textup{(i)}\quad
  (\alpha\circ f)|_U=(q\alpha_0)|_U,
\]
\[
  \textup{(ii)}\quad
  p(X)=\exp_\alpha\bigl(h(X^q)\bigr)
  \quad\textup{for every \(X\in W(\mathcal L)(S')\), \(S'\to U\).}
\]
\textit{Moreover, once \(U\) has been chosen, \(q\), \(\alpha\),
and \(h\) are uniquely determined.}

After shrinking \(S\), we may suppose that \(\alpha_0\) is nonzero
on every fiber of \(S\). Choose a system of positive roots \(R_+\)
of \(R\), and put
\[
  V=p^{-1}(\Omega_{R_+}).
\]
This is an open subset of \(W(\mathcal L)\) containing the zero
section and stable under multiplication by every
\(\alpha_0(q)\), where \(q\in Q(S')\) and \(S'\to S\). Since
\(\alpha_0\) is nontrivial on every fiber, it follows immediately
that \(V=W(\mathcal L)\), and hence that \(p\) factors through
\(\Omega_{R_+}\). Choose arbitrary orderings of \(R_+\) and
\(R_-\); every product below is understood to be taken in the
chosen ordering. There are therefore unique morphisms
\[
  a_\alpha:W(\mathcal L)\longrightarrow U_\alpha,
  \qquad \alpha\in R,
  \qquad
  b:W(\mathcal L)\longrightarrow T
\]
such that
\[
  p(x)
  =
  \prod_{\alpha\in R_-}a_\alpha(x)\mathbin{\cdot}b(x)
  \mathbin{\cdot}\prod_{\alpha\in R_+}a_\alpha(x).
\]
Writing out the covariance condition under \(Q\), one immediately
obtains
\[
  a_\alpha\bigl(\alpha_0(q)x\bigr)
  =
  \alpha\bigl(f(q)\bigr)a_\alpha(x),
  \qquad \alpha\in R,
\]
\[
  b\bigl(\alpha_0(q)x\bigr)=b(x)
\]
for every \(x\in W(\mathcal L)(S')\), \(q\in Q(S')\), and
\(S'\to S\). The second condition immediately gives \(b=e\).

Now let \(\alpha\in R\) be such that
\((a_\alpha)_{\bar s}\ne e\); such an \(\alpha\) exists because
\(p_{\bar s}\ne e\) by hypothesis. Applying \SGARefLink{sga3:xix:proposition:4:12:item:a}{Exposé XIX, 4.12(a)},
we find an integer \(n>0\) such that
\[
  (\alpha\circ f)_{\bar s}=(n\alpha_0)_{\bar s}.
\]
After shrinking \(S\), we may suppose that
\(\alpha\circ f=n\alpha_0\) (\SGARefLink{sga3:ix:corollary:5:3}{Exposé IX, 5.3}). But then, for every
\(\alpha'\in R\) with \(\alpha'\ne\alpha\), one has
\[
  (\alpha'\circ f)_{\bar s}\ne m\alpha_0
  \qquad\textup{for every integer \(m>0\),}
\]
by the hypothesis on \(f\) and the fact that the only roots
proportional to \(\alpha\) are \(\alpha\) and \(-\alpha\).
Applying \SGARefLink{sga3:xix:proposition:4:12:item:a}{Exposé XIX, 4.12(a)}, now to \(a_{\alpha'}\), we find that
\(a_{\alpha'}\) vanishes in a neighborhood of \(s\).
% The French authority prints uppercase S here; the local context
% and the ensuing shrinking argument require the point s.
Since \(R\) is finite, after shrinking \(S\) once more we may
suppose that \(a_{\alpha'}=0\) for every
\(\alpha'\in R\), \(\alpha'\ne\alpha\). Thus \(p=a_\alpha\), and
we may apply \SGARefLink{sga3:xix:proposition:4-12:item:b}{Exposé XIX, 4.12(b)}, followed by \SGARefLink{sga3:xix:proposition:4:12:item:c}{4.12(c)}, which gives
the asserted result. The uniqueness assertions are clear.

\SGARefTarget{sga3:xxii:rem:4.1.10}\paragraph{\SGARefLink{sga3:xxii:rem:4.1.10}{Remark 4.1.10}.}
The condition imposed on \(f\) in \SGARefLink{sga3:xxii:prop:4.1.9}{4.1.9} is satisfied in particular
when \(f\) is surjective \textup{(that is, faithfully flat)}.

\SGARefTarget{sga3:xxii:prop:4.1.11}\paragraph{\SGARefLink{sga3:xxii:prop:4.1.11}{Proposition 4.1.11}.}
\textit{Let \((G,T,M,R)\) be a split \(S\)-group, let \(R_+\) be a
system of positive roots of \(R\), and let \(\Omega_{R_+}\) be the
corresponding ``big cell.''}

\textit{\textup{\SGARefLink{sga3:xxii:def:1.13:item:i}{(i)}} Let \(H\) be a group-valued \(S\)-functor
that is separated\footnote{Editor's note: Recall
\textup{(cf.\ \SGARefLink{sga3:iv:corollary:4:3:5}{Exposé IV, 4.3.5})} that an \(S\)-presheaf \(H\) is
separated for a topology \(\mathcal T\) if, for every
\(S'\to S\) and every family of \(S\)-morphisms
\((S_i'\to S')_{i\in I}\) covering for \(\mathcal T\), the map
\[
  H(S')\longrightarrow\prod_i H(S_i')
\]
is injective.}
for the fppf topology. If
\(f,g:G\rightrightarrows H\) are two group homomorphisms that
coincide on \(\Omega_{R_+}\), then \(f=g\).}

\SGAInlineRefTarget{sga3:xxii:prop:4.1.11:item:ii}\textit{\textup{\SGARefLink{sga3:xxii:def:1.13:item:ii}{(ii)}} Let \(H\) be an \(S\)-sheaf of groups for the
fppf topology, and let \(f:\Omega_{R_+}\to H\) be an \(S\)-morphism
with the following property: for every \(S'\to S\) and every
\(x,y\in\Omega_{R_+}(S')\) such that
\(xy\in\Omega_{R_+}(S')\), one has
\(f(xy)=f(x)f(y)\). Then there exists a group homomorphism
\(\bar f:G\to H\), unique by \textup{\SGARefLink{sga3:xxii:def:1.13:item:i}{(i)}}, extending \(f\).}

Indeed, by \SGARefLink{sga3:xxii:prop:4.1.5}{4.1.5}, \textup{\SGARefLink{sga3:xxii:def:1.13:item:i}{(i)}} \textup{(respectively
\textup{\SGARefLink{sga3:xxii:def:1.13:item:ii}{(ii)}})} follows immediately from \SGARefLink{sga3:xviii:cor:2.2}{Exposé XVIII, 2.2}
\textup{(respectively \SGARefLink{sga3:xxii:cor:2.3}{2.3} and \SGARefLink{sga3:xxii:cor:2.4}{2.4})}.

\SGARefTarget{sga3:xxii:rem:4.1.12}\paragraph{\SGARefLink{sga3:xxii:rem:4.1.12}{Remark 4.1.12}.}
If \(\alpha\in R_+\), then
\[
  \tag{\(\dagger\)}\label{sga3:xxii:subsection:4-1:equation:dagger}
  \Omega_{R_+}\cap Z_\alpha
  =
  U_{-\alpha}\mathbin{\cdot}T\mathbin{\cdot}U_\alpha.
\]
Indeed,\footnote{Editor's note: We have expanded the original in
what follows.} for every \(S'\to S\), if
\[
  g=
  \prod_{\beta\in R_-}p_\beta(x_\beta)
  \mathbin{\cdot}t\mathbin{\cdot}
  \prod_{\beta\in R_+}p_\beta(x_\beta)
\]
is an element of \(\Omega_{R_+}(S')\), and if
\(t'\in T_\alpha(S')\), then
\[
  t'gt'^{-1}
  =
  \prod_{\beta\in R_-}p_\beta\bigl(\beta(t')x_\beta\bigr)
  \mathbin{\cdot}t\mathbin{\cdot}
  \prod_{\beta\in R_+}p_\beta\bigl(\beta(t')x_\beta\bigr).
\]
% The French authority prints S'' for the base of t', but g and
% the displayed conjugation both live over S'; English corrects it.
Since \(\alpha\) and \(-\alpha\) are the only elements of \(R\)
whose restriction to \(T_\alpha\) is \(1\), it follows that \(g\)
belongs to \(Z_\alpha=\operatorname{Centr}(T_\alpha)\) if and only
if \(x_\beta=0\) for \(\beta\ne\pm\alpha\).

From \((\dagger)\) and \SGARefLink{sga3:xx:2.1}{Exposé XX, 2.1}, one deduces that if
\(X\in\Gamma(S,\mathfrak g^\alpha)\) and
\(Y\in\Gamma(S,\mathfrak g^{-\alpha})\), then
\[
  \exp_\alpha(X)\exp_{-\alpha}(Y)\in\Omega_{R_+}(S)
  \quad\Longleftrightarrow\quad
  1+XY\ \text{is invertible}.
\]
% The French authority prints Gamma(X,g^alpha); the section base is S.

\SGARefTarget{sga3:xxii:section:4.2}
\subsection*{\SGAInlineRefTarget{sga3:xxii:subsection:4-2}4.2. Morphisms of Split Groups}
\addcontentsline{toc}{subsection}{\SGARefLink{sga3:xxii:subsection:4-2}{4.2}. Morphisms of Split Groups}

\SGARefTarget{sga3:xxii:def:4.2.1}\paragraph{\SGARefLink{sga3:xxii:def:4.2.1}{Definition 4.2.1}.}
Let \(S\) be a nonempty scheme, and let
\((G,T,M,R)\) and \((G',T',M',R')\) be two split \(S\)-groups.
An \(S\)-group homomorphism \(f:G\to G'\) is said to be
\emph{compatible with the splittings}, or to define a
\emph{morphism of split groups}, if the restriction of \(f\) to
\(T\) factors as a morphism \(f_T:T\to T'\) of the form
\(f_T=D_S(h)\), where \(h:M'\to M\) is a group homomorphism
satisfying the following condition: there exist a bijection
\(d:R\xrightarrow{\sim}R'\)\footnote{Editor's note: We have
replaced the notation \(u:R\to R'\) by \(d:R\to R'\).} and, for
every \(\alpha\in R\), an integer \(q(\alpha)>0\) such that
\(x\mapsto x^{q(\alpha)}\) is an endomorphism of
\(\mathbb G_{a,S}\) and
\[
  h\bigl(d(\alpha)\bigr)=q(\alpha)\alpha,
  \qquad
  {}^t h(\alpha^*)=q(\alpha)d(\alpha)^*.
\]

\SGARefTarget{sga3:xxii:rem:4.2.2}\paragraph{\SGARefLink{sga3:xxii:rem:4.2.2}{Remark 4.2.2}.}
It is immediate that \(h\), \(d\), and the \(q(\alpha)\), for
\(\alpha\in R\), are uniquely determined by \(f\). We write
\(h=\mathscr R(f)\). The \(q(\alpha)\) are the radicial exponents
of \(f\) \textup{(or of \(h\))}.

Let \(p\) be the prime number, if one exists, that vanishes on
\(S\); put \(p=1\) if no prime number vanishes on \(S\). Then
\(\mathscr R(f)\) is a \(p\)-morphism of reduced root data in the
sense of \SGARefLink{sga3:XXI:6.8}{Exposé XXI, 6.8}. We have thus defined a functor
\(\mathscr R\) from the category of split \(S\)-groups to the
category of reduced root data \textup{(equipped with
\(p\)-morphisms)}.

\SGARefTarget{sga3:xxii:prop:4.2.3}\paragraph{\SGARefLink{sga3:xxii:prop:4.2.3}{Proposition 4.2.3}.}
\textit{Under the hypotheses of \SGARefLink{sga3:xxii:def:4.2.1}{4.2.1}, the following properties
hold:}

\textit{\textup{\SGARefLink{sga3:xxii:def:1.13:item:i}{(i)}} For every \(\alpha\in R\), there exists a
unique isomorphism of \(\mathcal O_S\)-modules}
\[
  f_\alpha:
  (\mathfrak g^\alpha)^{\otimes q(\alpha)}
  \xrightarrow{\;\sim\;}
  \mathfrak g'^{\,d(\alpha)}
\]
\textit{such that}
\[
  f\bigl(\exp_\alpha(X)\bigr)
  =
  \exp_{d(\alpha)}
  \bigl(f_\alpha(X^{q(\alpha)})\bigr)
\]
\textit{for every \(X\in W(\mathfrak g^\alpha)(S')\) and every
\(S'\to S\).}

\textit{\textup{\SGARefLink{sga3:xxii:def:1.13:item:ii}{(ii)}} For every \(\alpha\in R\), one has
\(q(-\alpha)=q(\alpha)\), and \(f_\alpha\) and \(f_{-\alpha}\) are
contragredient to one another.}

\textit{\textup{\SGARefLink{sga3:xxii:def:1.9:item:iii}{(iii)}} For every \(\alpha\in R\), every
\(Z\in W(\mathfrak g^\alpha)^*(S')\), and every \(S'\to S\), one
has}
\[
  f\bigl(w_\alpha(Z)\bigr)
  =
  w_{d(\alpha)}\bigl(Z^{q(\alpha)}\bigr).
\]

By hypothesis, the following diagram is commutative:

\SGARefTarget{sga3:xxii:diagram:4.2.3}
\begin{center}
  % BEGIN INCLUDED SGA3_English_current_buildable_source_20260823/source/inputs/XXII/native_diagrams/exp22/diag_004_split_morphism_grid.tex
\begin{tikzcd}[column sep=4.6em,row sep=3.2em]
  \mathbb G_{m,S}
    \arrow[r,"\alpha^*"]
    \arrow[d,"q(\alpha)"']
  &
  T
    \arrow[r,"\alpha"]
    \arrow[d,"f_T"']
  &
  \mathbb G_{m,S}
    \arrow[d,"q(\alpha)"]
  \\
  \mathbb G_{m,S}
    \arrow[r,"d(\alpha)^*"]
  &
  T'
    \arrow[r,"d(\alpha)"]
  &
  \mathbb G_{m,S}
\end{tikzcd}

% END INCLUDED SGA3_English_current_buildable_source_20260823/source/inputs/XXII/native_diagrams/exp22/diag_004_split_morphism_grid.tex

\end{center}

It follows that \(f\) maps \(\operatorname{Ker}(\alpha)\) into
\(\operatorname{Ker}(d(\alpha))\), hence \(T_\alpha\) into
\(T'_{d(\alpha)}\), and hence \(Z_\alpha\) into
\(Z'_{d(\alpha)}\). It remains only to apply \SGARefLink{sga3:xx:3.10}{Exposé XX, 3.10} and
\SGARefLink{sga3:xx:3.11}{3.11}, to the groups \(Z_\alpha\) and \(Z'_{d(\alpha)}\).

\begingroup
\emergencystretch=2em
\SGARefTarget{sga3:xxii:prop:4.2.4}\paragraph{\SGARefLink{sga3:xxii:prop:4.2.4}{Proposition 4.2.4}.}
\textit{The morphism \(f\) induces a morphism \(f_N\) from
\(\underline{\operatorname{Norm}}_G(T)\) to
\(\underline{\operatorname{Norm}}_{G'}(T')\), hence a morphism
\(f_W\) from \(W_G(T)\) to \(W_{G'}(T')\); the latter is an
isomorphism. More precisely, let}
\[
  \bar d:W\bigl(\mathscr R(G)\bigr)=W
  \longrightarrow
  W'=W\bigl(\mathscr R(G')\bigr)
\]
\textit{be the isomorphism extending
\(s_\alpha\mapsto s_{d(\alpha)}\) \textup{(\SGARefLink{sga3:XXI:6.8.4}{Exposé XXI, 6.8.4})}.
Then there is a commutative diagram of isomorphisms:}
\endgroup

\SGARefTarget{sga3:xxii:diagram:4.2.4}
\begin{center}
  % The French source prints W_G(T') in the upper-right node; the
  % proposition and construction require W_{G'}(T').
  % BEGIN INCLUDED SGA3_English_current_buildable_source_20260823/source/inputs/XXII/native_diagrams/exp22/diag_005_weyl_isomorphism_square.tex
\begin{tikzcd}[column sep=5.2em,row sep=3.2em]
  W_G(T)
    \arrow[r,"f_W","\sim"']
  &
  W_{G'}(T')
  \\
  W_S
    \arrow[u,"\sim"']
    \arrow[r,"\bar d_S","\sim"']
  &
  W'_S.
    \arrow[u,"\sim"']
\end{tikzcd}

% END INCLUDED SGA3_English_current_buildable_source_20260823/source/inputs/XXII/native_diagrams/exp22/diag_005_weyl_isomorphism_square.tex

\end{center}

This follows immediately from \SGARefLink{sga3:xxii:prop:3.4}{3.4}, \SGARefLink{sga3:XXI:6.8.4}{Exposé XXI, 6.8.4}, and
\textup{\SGARefLink{sga3:xxii:def:1.9:item:iii}{(iii)}} above.

\SGARefTarget{sga3:xxii:rem:4.2.5}\paragraph{\SGARefLink{sga3:xxii:rem:4.2.5}{Remark 4.2.5}.}
With the notation of \SGARefLink{sga3:xxii:prop:4.2.3}{4.2.3}, the restriction of \(f\) to
\(\Omega_{R_+}\), for a system of positive roots \(R_+\), can be
written explicitly: it maps \(\Omega_{R_+}\) into
\(\Omega'_{d(R_+)}\) \textup{(\(d(R_+)\) is a system of positive
roots of \(R'\) by \SGARefLink{sga3:XXI:6.8.7}{Exposé XXI, 6.8.7})}, and is given by the
set-theoretic formula
\[
\begin{aligned}
  f\Biggl(
    \prod_{\alpha\in R_-}\exp_\alpha(X_\alpha)
    \mathbin{\cdot}t\mathbin{\cdot}
    \prod_{\alpha\in R_+}\exp_\alpha(X_\alpha)
  \Biggr)
  &=
  \prod_{\alpha\in R_-}
    \exp_{d(\alpha)}
    \Bigl(f_\alpha(X_\alpha^{q(\alpha)})\Bigr)
  \mathbin{\cdot}f_T(t)\mathbin{\cdot}\\
  &\qquad
  \prod_{\alpha\in R_+}
    \exp_{d(\alpha)}
    \Bigl(f_\alpha(X_\alpha^{q(\alpha)})\Bigr).
\end{aligned}
\]

\SGARefTarget{sga3:xxii:prop:4.2.6}\paragraph{\SGARefLink{sga3:xxii:prop:4.2.6}{Proposition 4.2.6}.}
\textit{\textup{\SGARefLink{sga3:xxii:def:1.13:item:i}{(i)}} The morphism \(f\) is surjective
\textup{(\(=\) faithfully flat in the present case; cf.\
Exposé \(\mathrm{VI_B}\), 3.11)} if and only if \(f_T\) is
surjective.}

\textit{\textup{\SGAInlineRefTarget{sga3:xxii:prop:4.2.6:item:ii}(ii)} One has
\(\operatorname{Ker}(f)\subset\Omega_{R_+}\).}

We prove \textup{\SGARefLink{sga3:xxii:def:1.13:item:i}{(i)}}. If \(f\) is surjective, then
\(f_T(T)=f(T)\) is a maximal torus of \(G'\). Indeed, \(f(T)\) is
a subtorus of a maximal torus \(T'\) \textup{(\SGARefLink{sga3:ix:theorem:6:8}{Exposé IX, 6.8})}; to
verify that \(f(T)=T'\), one reduces to the case of an algebraically
closed field, where this is \emph{Bible}, \S7.3, Theorem 3\textup{(a)}.

If \(f_T\) is surjective, the preceding formula shows that \(f\)
induces a surjection from \(\Omega=\Omega_{R_+}\) onto
\(\Omega'=\Omega'_{d(R_+)}\).\footnote{Editor's note: We have given
the reference to Exposé VI in greater detail below.} Since the fibers
of \(G'\) are connected, it follows \textup{(cf.\ Exposé
\(\mathrm{VI_A}\), 0.5)} that \(f\) is surjective.

We prove \textup{\SGARefLink{sga3:xxii:prop:4.2.6:item:ii}{(ii)}}, and for this purpose admit a result that will
be proved below \textup{(\SGARefLink{sga3:xxii:thm:5.7.4}{5.7.4})}. For each \(w\in W\), choose an
\(n_w\in\underline{\operatorname{Norm}}_G(T)(S)\) representing it.
Then the open subschemes \(n_w\Omega\), for \(w\in W\), cover \(G\).
It is therefore enough to prove that
\[
  \operatorname{Ker}(f)\cap n_w\Omega\ne\varnothing
  \quad\Longrightarrow\quad w=1.
\]
If \(x\in\Omega(S')\), where \(S'\to S\), and \(f(n_wx)=1\), then
\(f(x)=f(n_w)^{-1}\). But \(f(x)\in\Omega'(S')\), whereas
\[
  f(n_w)^{-1}\in
  \underline{\operatorname{Norm}}_{G'}(T')(S').
\]
By \SGARefLink{sga3:xxii:prop:4.2.4}{4.2.4}, we are reduced to proving the following.

\SGARefTarget{sga3:xxii:lem:4.2.7}\paragraph{\SGARefLink{sga3:xxii:lem:4.2.7}{Lemma 4.2.7}.}
\textit{Under the hypotheses of \SGARefLink{sga3:xxii:prop:4.1.2}{4.1.2}, one has}
\[
  \Omega\cap\underline{\operatorname{Norm}}_G(T)=T.
\]

Let
\[
  x=
  \prod_{\alpha\in R_-}p_\alpha(x_\alpha)
  \mathbin{\cdot}t\mathbin{\cdot}
  \prod_{\alpha\in R_+}p_\alpha(x_\alpha)
  =
  vtu
  \in\Omega(S').
\]
If \(x\) normalizes \(T_{S'}\), then, for every \(t'\in T(S')\),
\[
  xt'x^{-1}=t''\in T(S'),
\]
that is, \(xt'=t''x\), which may be written
\[
  v(tt')(t'^{-1}ut')
  =
  (t''vt''^{-1})(t''t)u.
\]
This gives \(t'^{-1}ut'=u\), hence
\[
  u\in\operatorname{Centr}_G(T)(S')=T(S'),
\]
and therefore \(u=1\). Similarly, \(v=1\).

\SGARefTarget{sga3:xxii:cor:4.2.8}\paragraph{\SGARefLink{sga3:xxii:cor:4.2.8}{Corollary 4.2.8}.}
\textit{One has}
\[
  \operatorname{Ker}(f)
  =
  \prod_{\alpha\in R_-}K_\alpha
  \mathbin{\cdot}\operatorname{Ker}(f_T)\mathbin{\cdot}
  \prod_{\alpha\in R_+}K_\alpha,
\]
\textit{where, for each \(\alpha\in R\), \(K_\alpha\) denotes the
finite \(S\)-group}
\[
  K_\alpha
  =
  \operatorname{Ker}\!\left(
    U_\alpha\longrightarrow U_\alpha^{\otimes q(\alpha)}
  \right)
  \simeq
  \boldsymbol{\alpha}_{q(\alpha),S}.
\]

To apply this corollary, we make the following definition.

\SGARefTarget{sga3:xxii:def:4.2.9}\paragraph{\SGARefLink{sga3:xxii:def:4.2.9}{Definition 4.2.9}.}
Let \(S\) be a scheme and let \(G\) and \(G'\) be two reductive
\(S\)-groups. A faithfully flat and finite morphism of \(S\)-groups
\(f:G\to G'\) \textup{(i.e.\ one that is surjective and has finite
kernel over \(S\))} is called an \emph{isogeny}. If, moreover,
\(\operatorname{Ker}(f)\) is a central subgroup of \(G\), then \(f\)
is called a \emph{central isogeny}.

\SGARefTarget{sga3:xxii:prop:4.2.10}\paragraph{\SGARefLink{sga3:xxii:prop:4.2.10}{Proposition 4.2.10}.}
\textit{Let \(f:G\to G'\) be a morphism of split groups. For \(f\)
to be an isogeny \textup{(respectively, a central isogeny)}, it is
necessary and sufficient that \(f_T\) be an isogeny, i.e.\ that
\(\mathscr R(f)\) be injective with finite cokernel
\textup{(respectively, and that \(q(\alpha)=1\) for every
\(\alpha\in R\))}.}

Indeed, by \SGARefLink{sga3:xxii:cor:4.2.8}{4.2.8}, \(\operatorname{Ker}(f)\) is finite over \(S\)
if and only if \(\operatorname{Ker}(f_T)\) is finite over \(S\),
and \(\operatorname{Ker}(f)\subset T\) if and only if every
\(q(\alpha)\) is equal to \(1\). In that case
\(\operatorname{Ker}(f)\) is central, since it is of multiplicative
type and invariant; cf.\ \SGARefLink{sga3:ix:corollary:5:5}{Exposé IX, 5.5}.

\SGARefTarget{sga3:xxii:rem:4.2.11}\paragraph{\SGARefLink{sga3:xxii:rem:4.2.11}{Remark 4.2.11}.}
\textup{(a)} Thus \(f:G\to G'\) is a central isogeny if and only if
\[
  \mathscr R(f):\mathscr R(G')\longrightarrow\mathscr R(G)
\]
is an isogeny in the sense of \SGARefLink{sga3:XXI:6.2}{Exposé XXI, 6.2}. Moreover, in this
case, with the notation of \emph{loc.\ cit.}, one has
\[
  \operatorname{Ker}(f)
  =
  D_S\!\left(K\bigl(\mathscr R(f)\bigr)\right),
  \qquad
  K\bigl(\mathscr R(f)\bigr)
  =
  \operatorname{Coker}\bigl(\mathscr R(f)\bigr).
\]

\textup{\SGARefLink{sga3:xxii:def:1.5:item:b}{(b)}} If \(G\) and \(G'\) are semisimple, every morphism of
split groups \(f:G\to G'\) is an isogeny.\footnote{Editor's note:
Indeed, since \(G\) and \(G'\) are semisimple, \(f_T\) is
surjective with finite kernel; hence \(f\) is faithfully flat with
finite kernel by \SGARefLink{sga3:xxii:prop:4.2.6}{4.2.6}\textup{(i)} and \SGARefLink{sga3:xxii:cor:4.2.8}{4.2.8}.}

\textup{(c)} If \(f:G\to G'\) is faithfully flat and finite and
\(G\) is reductive \textup{(respectively, semisimple)}, then \(G'\)
is so as well. Indeed, \(G'\) is of finite presentation over \(S\)
\textup{(Exposé V, 9.1)}, affine over \(S\)
\textup{(EGA II, 6.7.1)}, smooth over \(S\)
\textup{(Exposé VI, 9.2)}, and has connected reductive
\textup{(respectively, semisimple)} geometric fibers by \SGARefLink{sga3:xix:paragraph:1-7}{Exposé XIX,
1.7}.

\SGARefLink{sga3:xxii:def:4.2.1}{Definition 4.2.1} may seem arbitrary. It is justified by the
following proposition, which, for simplicity, we shall state for
semisimple groups.

A morphism \(f:G\to G'\) of reductive \(S\)-groups is said to be
\emph{splittable} if there exist splittings of \(G\) and \(G'\)
with which \(f\) is compatible. We then have the following.

\SGARefTarget{sga3:xxii:prop:4.2.12}\paragraph{\SGARefLink{sga3:xxii:prop:4.2.12}{Proposition 4.2.12}.}
\textit{Let \(S\) be a scheme, let \(G\) and \(G'\) be two
semisimple \(S\)-groups, and let \(f:G\to G'\) be a morphism of
groups. Let \(s\in S\). The following conditions are equivalent:}

\begin{enumerate}
\item[\textup{\SGAInlineRefTarget{sga3:xxii:prop:4.2.12:item:i}(i)}]
\(\operatorname{Ker}(f_{\bar s})\) is finite
\textup{(\(\Longleftrightarrow e\) is isolated in
\(\operatorname{Ker}(f)(\bar s)\))}, and \(f_{\bar s}\) is
surjective, i.e.\ \(f_{\bar s}\) is an isogeny.

\item[\textup{\SGAInlineRefTarget{sga3:xxii:prop:4.2.12:item:ii}(ii)}]
\(f_{\bar s}\) is splittable.

\item[\textup{\SGAInlineRefTarget{sga3:xxii:prop:4.2.12:item:iii}(iii)}]
There exists an étale morphism \(S'\to S\) covering \(s\) such that
\(f_{S'}:G_{S'}\to G'_{S'}\) is splittable.
\end{enumerate}

Clearly, \textup{\SGARefLink{sga3:xxii:prop:4.2.12:item:iii}{(iii)}} is equivalent to \textup{\SGARefLink{sga3:xxii:prop:4.2.12:item:ii}{(ii)}}, and
\textup{\SGARefLink{sga3:xxii:prop:4.2.12:item:ii}{(ii)}} implies \textup{\SGARefLink{sga3:xxii:prop:4.2.12:item:i}{(i)}} by \SGARefLink{sga3:xxii:rem:4.2.11}{4.2.11}\textup{\SGARefLink{sga3:xxii:def:1.5:item:b}{(b)}}. This is
where the hypothesis that \(G\) and \(G'\) are semisimple enters;
the other implications are valid for reductive groups.

We now prove that \textup{\SGARefLink{sga3:xxii:prop:4.2.12:item:i}{(i)}} implies \textup{\SGARefLink{sga3:xxii:prop:4.2.12:item:iii}{(iii)}}. We may
suppose that \(G\) and \(G'\) are split in such a way that \(f\)
induces a morphism
\[
  f_T:T\longrightarrow T'
\]
\textup{(\SGARefLink{sga3:xxii:cor:2.3}{2.3} and \SGARefLink{sga3:xix:paragraph:2:8}{Exposé XIX, 2.8})}. After restricting \(S\), we may
suppose that \(f_T=D_S(h)\), where \(h:M'\to M\) is a morphism of
groups. Let \(\alpha\in R\), and consider the composite morphism
\[
  p:
  W(\mathfrak g^\alpha)
  \xrightarrow{\ \exp_\alpha\ }
  G
  \xrightarrow{\ f\ }
  G'.
\]
Since \(\operatorname{Ker}(p_{\bar s})\) is finite,
\(p_{\bar s}\ne e\). On the other hand, \(f_{T_{\bar s}}\) is
surjective. We may therefore apply \SGARefLink{sga3:xxii:prop:4.1.9}{4.1.9}: there exist an open
subset \(V_\alpha\) of \(S\) containing \(s\), a root
\(\alpha'\in R'\), and an integer \(q(\alpha)\) such that
\(x\mapsto x^{q(\alpha)}\) is an endomorphism of
\(\mathbb G_{a,V_\alpha}\), and an isomorphism of
\(\mathcal O_{V_\alpha}\)-modules
\[
  f_\alpha:
  \left.(\mathfrak g^\alpha)^{\otimes q(\alpha)}
  \right|_{V_\alpha}
  \xrightarrow{\ \sim\ }
  \left.\mathfrak g'^{\alpha'}\right|_{V_\alpha},
\]
such that
\[
  f\bigl(\exp_\alpha(X_\alpha)\bigr)
  =
  \exp_{d(\alpha)}
  \bigl(f_\alpha(X_\alpha^{q(\alpha)})\bigr)
\]
and
\[
  \alpha'\circ f_T=h(\alpha')=q(\alpha)\alpha.
\]
We may replace \(S\) by the intersection of the \(V_\alpha\), for
\(\alpha\in R\). Put \(\alpha'=d(\alpha)\). It is clear that
\[
  d:R\longrightarrow R'
\]
is a bijection, since the kernel of \(h\) is finite
\textup{(\(f_{T_{\bar s}}\) being surjective)}. It remains only to
prove that
\[
  f_T\circ\alpha^*=q(\alpha)\alpha'^*,
\]
which follows by a trivial modification of the argument used in
\SGARefLink{sga3:xx:3.11}{Exposé XX, 3.11}.

In any event, as we have seen in the course of the proof,
\textup{\SGARefLink{sga3:xxii:prop:4.2.12:item:i}{(i)}} implies \textup{\SGARefLink{sga3:xxii:prop:4.2.12:item:iii}{(iii)}}. We have therefore proved the
following.

\SGARefTarget{sga3:xxii:cor:4.2.13}\paragraph{\SGARefLink{sga3:xxii:cor:4.2.13}{Corollary 4.2.13}.}
\textit{Let \(S\) be a scheme and let \(f:G\to G'\) be an isogeny
of reductive groups. Then \(f\) is locally splittable for the
étale topology.}

\SGARefTarget{sga3:xxii:section:4.3}\section*{4.3. Central
Quotients of Reductive Groups}
\addcontentsline{toc}{section}{\SGARefLink{sga3:xxii:section:4.3}{4.3}. Central Quotients of Reductive
Groups}

We first consider a special case.

\SGARefTarget{sga3:xxii:prop:4.3.1}\paragraph{\SGARefLink{sga3:xxii:prop:4.3.1}{Proposition 4.3.1}.}
\textit{Let \(S\) be a scheme, let \((G,T,M,R)\) be a split
\(S\)-group, let \(N\) be a subgroup of \(M\) containing \(R\),
and put}
\[
  Q=D_S(M/N)\subset\operatorname{Centr}(G).
\]
\textit{Then:}

\begin{enumerate}
\item[\textup{(i)}]
\textit{\(G'=G/Q\) is a reductive \(S\)-group, and \(T'=T/Q\) is a
maximal torus of \(G'\);}

\item[\textup{\SGAInlineRefTarget{sga3:xxii:prop:4.3.1:item:ii}(ii)}]
\textit{if \(T'\) is identified with \(D_S(N)\), then \(R\subset N\)
is a root system of \(G'\) relative to \(T'\),
\((G',T',N,R)\) is a splitting of \(G'\), and
\(\mathscr R(G')\) is canonically identified with the induced root
datum \textup{(\SGARefLink{sga3:XXI:6.5}{Exposé XXI, 6.5})} \(\mathscr R(G)_N\);}
% The Polo--Gille authority prints “un déploiement de G” here.
% Since the tuple begins with G' and describes the quotient, G' is
% required; the English corrects this transparent source typo.

\item[\textup{(iii)}]
\textit{the canonical morphism \(G\to G'\) is compatible with the
splittings, has root exponents equal to \(1\), and, by
functoriality, induces the canonical morphism \textup{(loc. cit.)}}
\[
  \mathscr R(G)_N\longrightarrow\mathscr R(G).
\]
\end{enumerate}

We know that \(G'=G/Q\) is representable by an affine group scheme
over \(S\) \textup{(\SGARefLink{sga3:viii:corollary:5:7}{Exposé VIII, 5.7})}, smooth over \(S\)
\textup{(Exposé VI\textsubscript{B}, \SGARefLink{prop:sga3-VIB-9.2}{9.2})}, and with connected
reductive geometric fibers \textup{(as quotients of reductive
groups; cf.\ \SGARefLink{sga3:xix:paragraph:1-7}{Exposé XIX, 1.7})}. Thus \(G'\) is a reductive
\(S\)-group.

It is clear that
\[
  T'=T/Q\simeq D_S(N)
\]
is a maximal torus of \(G'\). Next, choose a system \(R_+\) of
positive roots of \(R\). The open subset \(\Omega_{R_+}\) of \SGARefLink{sga3:xxii:subsection:4-2}{4.2} is
stable under \(Q\), and there is a canonical isomorphism
\[
  \Omega_{R_+}/Q
  \simeq
  \left(\prod_{\alpha\in R_-}U_\alpha\right)
  \mathop{\times}_S(T/Q)
  \mathop{\times}_S
  \left(\prod_{\alpha\in R_+}U_\alpha\right).
\]
Moreover, \(\Omega_{R_+}/Q\) is an open subset of \(G'\) containing
the identity section \textup{(cf.\ \SGARefLink{sga3:iv:proposition:4:7:2}{Exposé IV, 4.7.2} and \SGARefLink{sga3:iv:lemma:6:4:1}{6.4.1})}.

It follows that, if \(\mathfrak g'\) denotes the Lie algebra of
\(G'\) and \(\bar\alpha\) the character of \(T/Q\) induced by
\(\alpha\) \textup{(or, equivalently, defined by \(\alpha\in N\)
under the identification \(T/Q=D_S(N)\))}, then the canonical
morphism \(\mathfrak g\to\mathfrak g'\) induces, for every
\(\alpha\in R\), an isomorphism
\[
  \mathfrak g^\alpha
  \xrightarrow{\ \sim\ }
  \mathfrak g'^{\bar\alpha}.
\]

We have thus proved that \(R\) is a root system of \(G'\) relative
to \(T'\), and the proof is completed without difficulty.

\SGARefTarget{sga3:xxii:cor:4.3.2}\paragraph{\SGARefLink{sga3:xxii:cor:4.3.2}{Corollary 4.3.2}.}
\textit{Let \(S\) be a scheme, let \(G\) be a reductive
\(S\)-group, and let \(Q\) be a normal subgroup of \(G\) of
multiplicative type. Then \(Q\) is central in \(G\), the quotient
\(G/Q\) is representable by a reductive \(S\)-group, and the
canonical morphism \(G\to G/Q\) is locally splittable for the
étale topology \textup{(with root exponents equal to \(1\))}.}

The first assertion follows from \SGARefLink{sga3:ix:corollary:5:5}{Exposé IX, 5.5}; the others are
local for the étale topology and reduce to \SGARefLink{sga3:xxii:prop:4.3.1}{4.3.1}.

\SGARefTarget{sga3:xxii:def:4.3.3}\paragraph{\SGARefLink{sga3:xxii:def:4.3.3}{Definition 4.3.3}.}
Let \(G\) be a reductive \(S\)-group. We say that \(G\) is
\emph{adjoint} \textup{(respectively, \emph{simply connected})} if,
for every \(s\in S\), the type of \(G\) at \(s\) is given by an
adjoint \textup{(respectively, simply connected)} root datum, i.e.\
\textup{(\SGARefLink{sga3:XXI:6.2.6}{Exposé XXI, 6.2.6})} one for which \(M\) is generated by
\(R\) \textup{(respectively, \(M^*\) is generated by \(R^*\))}.

\SGARefTarget{sga3:xxii:prop:4.3.4}\paragraph{\SGARefLink{sga3:xxii:prop:4.3.4}{Proposition 4.3.4}.}
\textit{\textup{\SGAInlineRefTarget{sga3:xxii:prop:4.3.4:item:i}(i)} An adjoint \textup{(respectively, simply
connected)} reductive group is semisimple.}

\textit{\textup{\SGARefLink{sga3:xxii:prop:4.3.1:item:ii}{(ii)}} If \(T\) is a maximal torus of the adjoint
\textup{(respectively, simply connected)} reductive group \(G\),
and if \(\alpha\) is a root of \(G\) relative to \(T\), then the
infinitesimal root \(\bar\alpha\) is nonzero on every fiber
\textup{(respectively, \(\alpha^*\) is a monomorphism and the
infinitesimal coroot \(H_\alpha\) is nonzero on every fiber)}.}

Indeed, \textup{\SGARefLink{sga3:xxii:prop:4.3.4:item:i}{(i)}} is trivial; \textup{\SGARefLink{sga3:xxii:prop:4.3.1:item:ii}{(ii)}} may be checked on
the geometric fibers and follows at once from \SGARefLink{sga3:XXI:6.2.8}{Exposé XXI, 6.2.8}.

\SGARefTarget{sga3:xxii:prop:4.3.5}\paragraph{\SGARefLink{sga3:xxii:prop:4.3.5}{Proposition 4.3.5}.}
\textit{\textup{\SGARefLink{sga3:xxii:prop:4.3.4:item:i}{(i)}} The reductive group \(G\) is adjoint if and
only if}
\[
  \operatorname{Centr}(G)=\{e\}_S.
\]

\textit{\textup{\SGARefLink{sga3:xxii:prop:4.3.1:item:ii}{(ii)}} For every reductive group \(G\), the quotient
group \(G/\operatorname{Centr}(G)\) is an adjoint reductive group.}

Indeed, we may suppose that \(G\) is split. Then \textup{\SGARefLink{sga3:xxii:prop:4.3.4:item:i}{(i)}} is
immediate, since
\[
  \operatorname{Centr}(G)=D_S\bigl(M/\Gamma_0(R)\bigr),
\]
and \textup{\SGARefLink{sga3:xxii:prop:4.3.1:item:ii}{(ii)}} follows from \SGARefLink{sga3:xxii:prop:4.3.1}{4.3.1}.

\SGARefTarget{sga3:xxii:def:4.3.6}\paragraph{\SGARefLink{sga3:xxii:def:4.3.6}{Definition 4.3.6}.}
Let \(G\) be a reductive \(S\)-group. The \emph{adjoint group} of
\(G\), denoted by \(\operatorname{ad}(G)\), is the group
\(G/\operatorname{Centr}(G)\). The \emph{radical} of \(G\), denoted
by \(\operatorname{rad}(G)\), is the maximal torus
\textup{(unique by \SGARefLink{sga3:xii:1.12}{Exposé XII, 1.12})} of
\(\operatorname{Centr}(G)\). The \emph{semisimple group associated
with \(G\)} is the quotient \(G/\operatorname{rad}(G)\).

The preceding definitions are compatible with base change. If
\(s\in S\), then
\(\operatorname{rad}(G)_{\bar s}\) is indeed the radical of
\(G_{\bar s}\) in the usual sense \textup{(\SGARefLink{sga3:xii:1.6}{Exposé XIX, 1.6})}.

\SGARefTarget{sga3:xxii:par:4.3.7}\paragraph{\SGARefLink{sga3:xxii:par:4.3.7}{4.3.7}.}
If \((G,T,M,R)\) is a split group, then
\[
  \operatorname{rad}(G)=D_S(M/N),
  \qquad
  N=M\cap V(R).
\]
Thus the semisimple group associated with \(G\)
\textup{(as also the adjoint group of \(G\))} is equipped with a
canonical splitting \textup{(\SGARefLink{sga3:xxii:prop:4.3.1}{4.3.1})}, and we have a diagram of
split groups
\[
  G\longrightarrow \operatorname{ss}(G)
   \longrightarrow \operatorname{ad}(G)
\]
corresponding to the canonical diagram of root data
\textup{(\SGARefLink{sga3:XXI:6.5.5}{Exposé XXI, 6.5.5})}
\[
  \operatorname{ad}\bigl(\mathscr R(G)\bigr)
  \longrightarrow
  \operatorname{ss}\bigl(\mathscr R(G)\bigr)
  \longrightarrow
  \mathscr R(G).
\]

\SGARefTarget{sga3:xxii:rem:4.3.8}\paragraph{\SGARefLink{sga3:xxii:rem:4.3.8}{Remark 4.3.8}.}
\textit{Let \((G,T,M,R)\) be a split adjoint
\textup{(respectively, simply connected)} \(S\)-group, and let
\(\Delta\) be a system of simple roots of \(R\). Then the family
\(\{\alpha\}_{\alpha\in\Delta}\), respectively
\(\{\alpha^*\}_{\alpha\in\Delta}\), induces an isomorphism}
\[
  T\xrightarrow{\ \sim\ }(\mathbb G_{m,S})^\Delta,
  \qquad\textit{respectively}\qquad
  (\mathbb G_{m,S})^\Delta\xrightarrow{\ \sim\ }T.
\]

Indeed,
\[
  M=\Gamma_0(R)
  \qquad
  \textup{(respectively, \(M^*=\Gamma_0(R^*)\))},
\]
and \(\Delta\) \textup{(respectively,
\(\{\alpha^*\}_{\alpha\in\Delta}\))} is a basis of the free
abelian group \(\Gamma_0(R)\)
\textup{(respectively, \(\Gamma_0(R^*)\))}; see \SGARefLink{sga3:XXI:3.1.8}{Exposé XXI, 3.1.8}.
% The French abbreviates the simply-connected counterpart by
% “(resp. ...)”; the English expands the formally dual statement.

\SGARefTarget{sga3:xxii:rem:4.3.9}\paragraph{\SGARefLink{sga3:xxii:rem:4.3.9}{Remark 4.3.9}.}
The radical of a reductive group is a \emph{characteristic}
subgroup \textup{(i.e.\ stable under
\(\underline{\operatorname{Aut}}_{S\text{-gr.}}(G)\))}, by its
definition.

\SGARefTarget{sga3:xxii:section:5}\section*{5. Subgroups of Type
\textup{(R)}}
\addcontentsline{toc}{section}{5. Subgroups of Type (R)}

We are especially interested in reductive groups, but some of the
results that we shall establish hold more generally for a wider
class of groups: groups of type \textup{(RR)}.

\SGARefTarget{sga3:xxii:section:5.1}\subsection*{\SGAInlineRefTarget{sga3:xxii:subsection:5-1}5.1. Groups of
Type \textup{(RR)}}
\addcontentsline{toc}{subsection}{\SGARefLink{sga3:xxii:subsection:5-1}{5.1}. Groups of Type (RR)}

\SGARefTarget{sga3:xxii:def:5.1.1}\paragraph{\SGARefLink{sga3:xxii:def:5.1.1}{Definition 5.1.1}.}
Let \(S\) be a scheme and \(G\) an \(S\)-group scheme. We say that
\(G\) is \emph{of type \textup{(RR)}} if it satisfies the following
conditions:

\begin{enumerate}
\item[\textup{\SGAInlineRefTarget{sga3:xxii:def:5.1.1:item:i}(i)}]
\(G\) is smooth and of finite presentation over \(S\), with
connected fibers.

\item[\textup{\SGAInlineRefTarget{sga3:xxii:def:5.1.1:item:ii}(ii)}]
\(G\) locally possesses maximal tori for the fpqc topology.

\item[\textup{\SGAInlineRefTarget{sga3:xxii:def:5.1.1:item:iii}(iii)}]
For every \(s\in S\), every maximal torus \(T\) of
\(G_{\bar s}\), and every root \(\alpha\) of \(G_{\bar s}\)
relative to \(T_{\bar s}\) \textup{(\SGARefLink{sga3:xix:paragraph:1-10}{Exposé XIX, 1.10})},
\[
  \Lie(G_{\bar s})^\alpha
\]
has rank \(1\) \textup{(as a vector space over
\(\overline{\kappa(s)}\))}.

\item[\textup{\SGAInlineRefTarget{sga3:xxii:def:5.1.1:item:iv}(iv)}]
For every \(s\in S\) and every maximal torus \(T\) of
\(G_{\bar s}\), let \(R\) denote the set of roots of
\(G_{\bar s}\) relative to \(T\), and let \(\Gamma_0(R)\) be the
subgroup of
\[
  \operatorname{Hom}_{\bar s\text{-gr.}}
    (T,\mathbb G_{m,\bar s})
\]
generated by \(R\). Then the
content\footnote{Editor's note: The content of the root \(\alpha\)
is the positive generator of the ideal
\(\{f(\alpha)\mid f\in\Gamma_0(R)^*\}\) of \(\mathbb Z\); it is the
greatest integer \(c>0\) such that
\(\alpha/c\in\Gamma_0(R)\).}
of every root \(\alpha\in R\) in the free abelian group
\(\Gamma_0(R)\) \textup{(and hence a positive integer)} is
invertible on \(S\).
\end{enumerate}

\SGARefTarget{sga3:xxii:reminder:5.1.1.1}
\paragraph{Reminder \SGARefLink{sga3:xxii:reminder:5.1.1.1}{5.1.1.1}.}
\footnote{Editor's note: This reminder was added.}
Let us recall that, if \(G\) is a smooth connected algebraic group
over an algebraically closed field \(k\), a \emph{Cartan subgroup}
of \(G\) is the centralizer of a maximal torus of \(G\)
\textup{(\SGARefLink{sga3:xii:1.0}{Exposé XII, 1.0})}, and such a subgroup is smooth and
connected. For this, as well as for other characterizations of
Cartan subgroups, see \emph{Bible}, \S7.1, Theorem 1 in the case
where \(G\) is affine, and \SGARefLink{sga3:xii:6.6}{Exposé XII, Theorem 6.6} in the general
case.

If \(S\) is an arbitrary scheme and \(G\) is a smooth \(S\)-group
of finite type, a \emph{Cartan subgroup} of \(G\) is a smooth
\(S\)-subgroup \(C\) of \(G\) such that, for every \(s\in S\),
\(C_{\bar s}\) is a Cartan subgroup of \(G_{\bar s}\)
\textup{(\SGARefLink{sga3:xii:3.1}{Exposé XII, Definition 3.1})}.

\SGARefTarget{sga3:xxii:rem:5.1.2}\paragraph{\SGARefLink{sga3:xxii:rem:5.1.2}{Remark 5.1.2}.}
\textup{\SGAInlineRefTarget{sga3:xxii:rem:5.1.2:item:a}(a)} By \SGARefLink{sga3:xii:7.1}{Exposé XII, 7.1} \textup{(where the hypothesis that
\(G\) is separated is satisfied here because \(G\) has connected
fibers; cf.\ Exposé VI\textsubscript{B}, \SGARefLink{cor:sga3-VIB-5.5}{5.5})}, conditions
\textup{\SGARefLink{sga3:xxii:def:5.1.1:item:i}{(i)}} and \textup{\SGARefLink{sga3:xxii:def:5.1.1:item:ii}{(ii)}} imply that \(G\) locally possesses,
for the étale topology, maximal tori \textup{(respectively, Cartan
subgroups)}, which are locally conjugate for the étale topology.

\textup{\SGARefLink{sga3:xxii:def:1.5:item:b}{(b)}} The Cartan subgroups of \(G\) have connected fibers
\textup{(\SGARefLink{sga3:xii:6.6}{Exposé XII, 6.6})}.

\textup{\SGARefLink{sga3:xii:6.6:item:c}{(c)}} If \(G\) is affine over \(S\), conditions
\textup{(i)} and \textup{\SGARefLink{sga3:xxii:def:5.1.1:item:ii}{(ii)}} are respectively equivalent to:
\begin{enumerate}
\item[\textup{(i$'$)}]
\(G\) is smooth over \(S\) and has connected fibers;

\item[\textup{(ii$'$)}]
the reductive rank of the fibers of \(G\) is locally constant
\textup{(\SGARefLink{sga3:xii:1.7}{Exposé XII, 1.7})}.
\end{enumerate}

\textup{\SGAInlineRefTarget{sga3:xxii:rem:5.1.2:item:d}(d)} By \SGARefLink{sga3:xii:8.8}{Exposé XII, 8.8} \textup{\SGARefLink{sga3:xii:8.8:item:c:02}{(c)}} and \textup{\SGARefLink{sga3:xxii:rem:5.1.2:item:d}{(d)}},
\(G\) has a reductive center \(Z\), and for every \(s\in S\), in
the notation of \textup{\SGARefLink{sga3:xxii:def:5.1.1:item:iv}{(iv)}},\footnote{Editor's note: The original
said ``under the conditions of \textup{\SGARefLink{sga3:xxii:def:5.1.1:item:iv}{(iv)}},'' but the last
condition of \textup{\SGARefLink{sga3:xxii:def:5.1.1:item:iv}{(iv)}} does not appear to be used here.}
\[
  Z_{\bar s}
  =
  \bigcap_{\alpha\in R}\operatorname{Ker}(\alpha),
\]
whence
\[
  \operatorname{Hom}_{\bar s\text{-gr.}}
    \bigl((T/Z)_{\bar s},\mathbb G_{m,\bar s}\bigr)
  \simeq
  \Gamma_0(R).
\]

\textup{\SGAInlineRefTarget{sga3:xxii:rem:5.1.2:item:e}(e)} Condition \textup{\SGARefLink{sga3:xxii:def:5.1.1:item:iv}{(iv)}} holds in particular in the
following two cases:
\begin{enumerate}
\item[\textup{(1)}]
\(S\) has characteristic \(0\);

\item[\textup{(2)}]
every root is an indivisible element of the group generated by the
roots.
\end{enumerate}

\textup{(f)} The property of being of type \textup{(RR)} is stable
under base change and is local for the fpqc topology.

Remarks \textup{(c)} and \textup{\SGARefLink{sga3:xxii:rem:5.1.2:item:e}{(e)}} immediately give the
following proposition.

\SGARefTarget{sga3:xxii:prop:5.1.3}
\paragraph{\SGARefLink{sga3:xxii:prop:5.1.3}{Proposition 5.1.3}.}
\textit{A reductive group is of type \textup{(RR)}.}

\SGARefTarget{sga3:xxii:prop:5.1.4}
\paragraph{\SGARefLink{sga3:xxii:prop:5.1.4}{Proposition 5.1.4}.}
\textit{Let \(S\) be a scheme, let \(G\) be an \(S\)-group of type
\textup{(RR)}, and let \(Q\) be a central subgroup of \(G\), of
finite presentation over \(S\), such that the quotient \(G/Q\) is
representable \textup{(for example, \(G\) affine over \(S\) and
\(Q\) of multiplicative type \textup{(\SGARefLink{sga3:ix:proposition:2:3}{Exposé IX, 2.3})}, or \(S\)
artinian \textup{(Exposé VI\textsubscript{A}, \SGARefLink{thm:sga3-VIA-3.3.2}{3.3.2})})}. Then
\(G/Q\) is an \(S\)-group of type \textup{(RR)}.}

Indeed, \(G/Q\) is smooth over \(S\)
\textup{(Exposé VI\textsubscript{B}, \SGARefLink{prop:sga3-VIB-9.2}{9.2})}, of finite presentation
over \(S\) \textup{(\SGARefLink{sga3:v:proposition:9:1}{Exposé V, 9.1})}, and has connected fibers, so
condition \textup{\SGARefLink{concl:sga3-VIB-9.2-i}{(i)}} holds. On the other hand, condition
\textup{\SGARefLink{concl:sga3-VIB-9.2-ii}{(ii)}} follows from \SGARefLink{sga3:xii:7.6}{Exposé XII, 7.6}; it remains to verify
conditions \textup{\SGARefLink{sga3:xxii:def:5.1.1:item:iii}{(iii)}} and \textup{\SGARefLink{sga3:xxii:def:5.1.1:item:iv}{(iv)}}.

Put \(G'=G/Q\), let
\[
  u:G\longrightarrow G'
\]
be the canonical morphism, and let \(T'=u(T)\) be the maximal torus
of \(G'\) that is the image of \(T\)
\textup{(cf.\ \SGARefLink{sga3:xii:7.1}{Exposé XII, 7.1} \textup{\SGARefLink{sga3:xii:7.1:item:e}{(e)}})}. For each
\(\alpha\in R\), again denote by \(\alpha\) the character of
\(T'\) induced by \(\alpha\); indeed,
\[
  Q\cap T
  \subset
  \bigcap_{\alpha\in R}\operatorname{Ker}(\alpha)
\]
by \SGARefLink{sga3:xxii:rem:5.1.2}{5.1.2} \textup{\SGARefLink{sga3:xxii:rem:5.1.2:item:d}{(d)}}. We first prove:

\SGARefTarget{sga3:xxii:lemma:5.1.5}
\paragraph{\SGARefLink{sga3:xxii:lemma:5.1.5}{Lemma 5.1.5}.}
\textit{Under the hypotheses of \SGARefLink{sga3:xxii:prop:5.1.4}{5.1.4}, let \(T=D_S(M)\) be a
trivialized maximal torus of \(G\), and suppose that the
decomposition of \(\mathfrak g=\Lie(G)\) under
\(\operatorname{Ad}(T)\) has the form}
\[
  \mathfrak g
  =
  \mathfrak g^0
  \oplus
  \coprod_{\alpha\in R}\mathfrak g^\alpha,
  \qquad
  R\subset M-\{0\},
\]
\textit{where, for every \(s\in S\),
\(\mathfrak g^\alpha(s)\ne0\) for \(\alpha\in R\)
\textup{(thus \(\mathfrak g^\alpha\) is an invertible
\(\mathcal O_S\)-module for every \(\alpha\in R\), and \(R\) is
the set of roots of \(G_{\bar s}\) relative to \(T_{\bar s}\), for
every \(s\in S\))}. Then the Lie algebra \(\mathfrak g'\) of \(G'\)
decomposes under \(\operatorname{Ad}(T')\) as follows:}
\[
  \mathfrak g'
  =
  \mathfrak g'^{\,0}
  \oplus
  \coprod_{\alpha\in R}\mathfrak g'^{\,\alpha},
\]
\textit{and \(\Lie(u)\) induces an isomorphism from
\(\mathfrak g^\alpha\) onto \(\mathfrak g'^{\,\alpha}\).}

Indeed, let
\[
  p=\Lie(u):\mathfrak g\longrightarrow\mathfrak g'.
\]
One immediately has
\(p(\mathfrak g^\alpha)\subset\mathfrak g'^{\,\alpha}\) for every
\(\alpha\in R\), and
\(p(\mathfrak g^0)\subset\mathfrak g'^{\,0}\). Since
\[
  \operatorname{Ker}(p)
  =
  \Lie(Q)
  \subset
  \Lie\bigl(\operatorname{Centr}_G(T)\bigr)
  =
  \mathfrak g^0,
\]
\(p\) induces a monomorphism from \(\mathfrak g^\alpha\) into
\(\mathfrak g'^{\,\alpha}\), for every \(\alpha\in R\).

To prove the lemma, it is enough to do so when \(S\) is the
spectrum of an algebraically closed field; in view of the preceding
remarks, it is then enough to prove that
\[
  \operatorname{rk}(\mathfrak g')
  =
  \operatorname{rk}(\mathfrak g'^{\,0})
  +
  \operatorname{Card}(R).
\]
Set
\[
  C=\operatorname{Centr}_G(T),
  \qquad
  C'=\operatorname{Centr}_{G'}(T').
\]
By \SGARefLink{sga3:xii:7.1}{Exposé XII, 7.1} \textup{\SGARefLink{sga3:xii:7.1:item:e}{(e)}}, \(u\) induces a faithfully flat
morphism \(C\to C'\) with kernel \(Q\). Hence
\[
  \dim C'+\dim Q=\dim C.
\]
But \(G\), \(G'\), \(C\), and \(C'\) are smooth, and therefore
\[
\begin{aligned}
  \dim G
  &=
  \operatorname{rk}(\mathfrak g)
   =
  \operatorname{rk}(\mathfrak g^0)+\operatorname{Card}(R)
   =
  \dim C+\operatorname{Card}(R)
  \\
  &=
  \dim Q+\dim C'+\operatorname{Card}(R)
   =
  \dim Q+\operatorname{rk}(\mathfrak g'^{\,0})
  +\operatorname{Card}(R),
  \\
  \operatorname{rk}(\mathfrak g')
  &=
  \dim G'
   =
  \dim G-\dim Q.
\end{aligned}
\]
It follows that
\[
  \operatorname{rk}(\mathfrak g')
  =
  \operatorname{rk}(\mathfrak g'^{\,0})
  +
  \operatorname{Card}(R),
\]
which is the desired relation.

We return to the proof of \SGARefLink{sga3:xxii:prop:5.1.4}{5.1.4}. We may suppose that \(S\) is the
spectrum of an algebraically closed field. Let \(T\) be a maximal
torus of \(G\). Applying \SGARefLink{sga3:xxii:lemma:5.1.5}{5.1.5}, we immediately obtain
\textup{\SGARefLink{sga3:xxii:def:5.1.1:item:iii}{(iii)}} and \textup{\SGARefLink{sga3:xxii:def:5.1.1:item:iv}{(iv)}} for \(G/Q\).

To use the preceding proposition, we introduce a definition.

\SGARefTarget{sga3:xxii:def:5.1.6}\paragraph{\SGARefLink{sga3:xxii:def:5.1.6}{Definition 5.1.6}.}
We say that the \(S\)-group scheme \(G\) is \emph{of type
\textup{(RA)}} if it is of type \textup{(RR)} and also satisfies
condition \textup{(iv$'$)} below, which is stronger than
\textup{(iv)}:
\begin{enumerate}
\item[\textup{(iv$'$)}]
For every \(s\in S\) and every maximal torus \(T\) of
\(G_{\bar s}\), the content of every root of \(G_{\bar s}\)
relative to \(T\), in
\[
  \operatorname{Hom}_{\bar s\text{-gr.}}
    (T,\mathbb G_{m,\bar s}),
\]
is invertible on \(S\).
\end{enumerate}

\SGARefTarget{sga3:xxii:rem:5.1.7}\paragraph{\SGARefLink{sga3:xxii:rem:5.1.7}{Remarks 5.1.7}.}
\textup{\SGARefLink{sga3:xxii:rem:5.1.2:item:a}{(a)}} An adjoint reductive \(S\)-group is of type
\textup{(RA)}.

\textup{\SGARefLink{sga3:xxii:def:1.5:item:b}{(b)}} If \(S\) has characteristic \(0\), every group of
type \textup{(RR)} is of type \textup{(RA)}.

\textup{(c)} The property of being of type \textup{(RA)} is stable
under base change and is local for the fpqc topology.

Remark \textup{\SGARefLink{sga3:xxii:rem:5.1.2:item:a}{(a)}} above admits the following generalization.

\SGARefTarget{sga3:xxii:prop:5.1.8}\paragraph{\SGARefLink{sga3:xxii:prop:5.1.8}{Proposition 5.1.8}.}
\textit{Let \(S\) be a scheme, let \(G\) be an \(S\)-group of type
\textup{(RR)}, and let \(Z\) be its reductive center. Suppose that
the quotient \(G/Z\) is representable \textup{(for example, if
\(G\) is affine over \(S\), or if \(S\) is artinian)}. Then \(G/Z\)
is of type \textup{(RA)}.}

Indeed, this follows immediately from \SGARefLink{sga3:xxii:prop:5.1.4}{5.1.4}, \SGARefLink{sga3:xxii:lemma:5.1.5}{5.1.5}, and \SGARefLink{sga3:xxii:rem:5.1.2}{5.1.2}
\textup{\SGARefLink{sga3:xxii:rem:5.1.2:item:d}{(d)}}.

\SGARefTarget{sga3:xxii:rem:5.1.9}\paragraph{\SGARefLink{sga3:xxii:rem:5.1.9}{Remark 5.1.9}.}
If \(G\) is of type \textup{(RR)} and \(T\) is a maximal torus of
\(G\), \SGARefLink{sga3:xix:proposition:6-3}{Exposé XIX, 6.3} applies. In particular,
\(W_G(T)\) is étale, quasi-finite, and separated over \(S\).

\SGARefTarget{sga3:xxii:section:5.2}
\subsection*{\SGAInlineRefTarget{sga3:xxii:subsection:5-2}5.2. Subgroups of Type \textup{(R)}}
\addcontentsline{toc}{subsection}{\SGARefLink{sga3:xxii:subsection:5-2}{5.2}. Subgroups of Type (R)}

\SGARefTarget{sga3:xxii:def:5.2.1}\paragraph{\SGARefLink{sga3:xxii:def:5.2.1}{Definition 5.2.1}.}
Let \(S\) be a scheme, let \(G\) be a smooth \(S\)-group scheme of
finite presentation with connected fibers,\footnote{Editor's note:
The hypothesis that \(G\) (respectively \(H\)) is of finite
presentation over \(S\) is automatically satisfied: since \(G\)
(respectively \(H\)) is smooth over \(S\) with connected fibers, it
is quasi-compact and separated over \(S\) (Exposé
VI\textsubscript{B}, \SGARefLink{cor:sga3-VIB-5.5}{5.5}), and hence of finite presentation over
\(S\).} and let \(H\) be a group subscheme of \(G\). We say that
\(H\) is \emph{of type \textup{(R)}} if:
\begin{enumerate}
\item[\textup{\SGAInlineRefTarget{sga3:xxii:def:5.2.1:item:i}(i)}]
\(H\) is smooth and of finite presentation over \(S\), and has
connected fibers;\footnotemark[\value{footnote}]

\item[\textup{\SGAInlineRefTarget{sga3:xxii:def:5.2.1:item:ii}(ii)}]
for every \(s\in S\), \(H_{\bar s}\) contains a Cartan subgroup of
\(G_{\bar s}\).
\end{enumerate}
This notion is stable under base change and is local for the fpqc
topology.

\SGARefTarget{sga3:xxii:recall:5.2.2}\paragraph{Recall \SGARefLink{sga3:xxii:recall:5.2.2}{5.2.2}.}
\textup{(cf.\ \SGARefLink{sga3:xii:7.9}{Exposé XII, 7.9})} Under the preceding conditions:
\begin{enumerate}
\item[\textup{\SGAInlineRefTarget{sga3:xxii:def:5.2.1:item:a}(a)}]
\[
  H=\underline{\operatorname{Norm}}_G(H)^0.
\]

\item[\textup{\SGAInlineRefTarget{sga3:xxii:def:5.2.1:item:b}(b)}]
If \(G\) locally possesses Cartan subgroups (respectively maximal
tori) for the étale topology, then so does \(H\), and the Cartan
subgroups (respectively maximal tori) of \(H\) are Cartan subgroups
(respectively maximal tori) of \(G\).
\end{enumerate}

\SGARefTarget{sga3:xxii:examples:5.2.3}\paragraph{\SGARefLink{sga3:xxii:examples:5.2.3}{Examples 5.2.3}.}
\textup{\SGARefLink{sga3:xxii:def:5.2.1:item:a}{(a)}} \emph{Borel subgroups.} A Borel subgroup of \(G\) is a
subgroup \(H\) of type \textup{(R)} whose geometric fibers are
Borel subgroups of the corresponding fibers of
\(G\).\footnote{Editor's note: Equivalently, \(H\) is a smooth
subgroup of \(G\) such that every geometric fiber \(H_{\bar s}\) is
a Borel subgroup of \(G_{\bar s}\) (since every Borel subgroup of
\(G_{\bar s}\) is connected and contains a Cartan subgroup of
\(G_{\bar s}\)).}

\textup{(b)} \emph{Parabolic subgroups.} A parabolic subgroup of
\(G\) is a subgroup of type \textup{(R)} whose geometric fibers
contain Borel subgroups.

Other examples are given by the following propositions.

\SGARefTarget{sga3:xxii:prop:5.2.4}\paragraph{\SGARefLink{sga3:xxii:prop:5.2.4}{Proposition 5.2.4}.}
\textit{Let \(G\) be as in \SGARefLink{sga3:xxii:def:5.2.1}{5.2.1}, let \(K\subset H\) be two group
subschemes of \(G\), and suppose that \(H\) is of type
\textup{(R)}. Then \(K\) is a subgroup of type \textup{(R)} of
\(H\) if and only if it is a subgroup of type \textup{(R)} of
\(G\).}

Indeed,\footnote{Editor's note: We have expanded the argument that
follows.} let \(s\in S\). Since \(H\) is of type \textup{(R)},
every maximal torus of \(H_{\bar s}\) is a maximal torus of
\(G_{\bar s}\), and the same therefore holds for Cartan subgroups.

\SGARefTarget{sga3:xxii:prop:5.2.5}\paragraph{\SGARefLink{sga3:xxii:prop:5.2.5}{Proposition 5.2.5}.}
\textit{Let \(G\) be as in \SGARefLink{sga3:xxii:def:5.2.1}{5.2.1}, let \(T\) be a maximal torus of
\(G\), let \(Q\) be a subtorus of \(T\), and put
\(Z=\underline{\operatorname{Centr}}_G(Q)\). If \(H\) is a
subgroup of type \textup{(R)} of \(G\) containing \(T\), then
\(H\cap Z\) is a subgroup of type \textup{(R)} of \(Z\).}

First recall that \(Z\) is a closed group subscheme of \(G\), of
finite presentation (\SGARefLink{sga3:xi:corollary:6:11}{Exposé XI, 6.11}), with connected fibers
(\SGARefLink{sga3:xii:6.6}{Exposé XII, 6.6}), and smooth over \(S\) (\SGARefLink{sga3:xi:corollary:2:4}{Exposé XI, 2.4}), so it
satisfies the conditions imposed in \SGARefLink{sga3:xxii:def:5.2.1}{Definition 5.2.1}. Likewise,
\(H\cap Z\) is of finite presentation, smooth, and has connected
fibers (since
\(H\cap Z=\underline{\operatorname{Centr}}_H(Q)\)); moreover,
\[
  H\cap Z\supset\underline{\operatorname{Centr}}_G(T).
\]

\SGARefTarget{sga3:xxii:prop:5.2.6}\paragraph{\SGARefLink{sga3:xxii:prop:5.2.6}{Proposition 5.2.6}.}
\textit{Let \(S\) be a scheme, let \(G\) be an \(S\)-group of type
\textup{(RR)} \textup{(respectively \textup{(RA)})}, and let \(H\)
be a subgroup of type \textup{(R)} of \(G\). Then \(H\) is an
\(S\)-group of type \textup{(RR)} \textup{(respectively
\textup{(RA)})}.}

Indeed, \textup{\SGARefLink{sga3:xxii:def:5.2.1:item:i}{(i)}} is clear, \textup{\SGARefLink{sga3:xxii:def:5.2.1:item:ii}{(ii)}} follows from \SGARefLink{sga3:xxii:recall:5.2.2}{5.2.2}
\textup{\SGARefLink{sga3:xxii:def:5.2.1:item:b}{(b)}}, and \textup{\SGARefLink{sga3:xxii:def:5.1.1:item:iii}{(iii)}} and \textup{\SGARefLink{sga3:xxii:def:5.1.1:item:iv}{(iv)}}
\textup{(respectively \textup{(iv$'$)})} are to be verified when
\(S\) is the spectrum of an algebraically closed field. Then \(H\)
contains a maximal torus \(T\) of \(G\), and therefore also
\[
  C=\underline{\operatorname{Centr}}_G(T),\footnote{Editor's note:
By hypothesis, \(H\) contains
\(\underline{\operatorname{Centr}}_G(T')\) for some maximal torus
\(T'\) of \(G\). The tori \(T\) and \(T'\) are then conjugate in
\(H\), and hence \(H\) also contains
\(C=\underline{\operatorname{Centr}}_G(T)\).}
\]
and the assertions to be proved follow immediately from:

\SGARefTarget{sga3:xxii:lemma:5.2.7}\paragraph{\SGARefLink{sga3:xxii:lemma:5.2.7}{Lemma 5.2.7}.}
\textit{Let \(S\) be a scheme, let \(G\) be an \(S\)-group of type
\textup{(RR)}, and let \(T\) be a maximal torus of \(G\) equipped
with a trivialization \(T\simeq D_S(M)\). Suppose that}
\[
  \mathfrak g
  =
  \mathfrak g^0
  \oplus
  \coprod_{\alpha\in R}\mathfrak g^\alpha
\]
\textit{\textup{(the \(\mathfrak g^\alpha\) are then invertible
\(\mathcal O_S\)-modules)}. Let \(H\) be a subgroup of type
\textup{(R)} containing
\(C=\underline{\operatorname{Centr}}_G(T)\)
\textup{(that is, containing \(T\))}. Then
\(\mathfrak h=\Lie(H/S)\) is locally on \(S\) of the form}
\[
  \mathfrak g^0
  +
  \coprod_{\alpha\in R'}\mathfrak g^\alpha
  =
  \mathfrak g_{R'}.
\]
\textit{More precisely, for each \(s\in S\), put}
\[
  R'(s)
  =
  \bigl\{
    \alpha\in R
    \mathrel{\big|}
    \mathfrak g^\alpha(s)\subset\mathfrak h(s)
  \bigr\}.
\]
\textit{Then \(R'(s)\) is a locally constant function of \(s\). If
\(U\) is an open subset of \(S\) on which \(R'(s)=R'\), then}
\[
  \mathfrak h_U
  =
  \mathfrak g_U^0
  \oplus
  \coprod_{\alpha\in R'}\mathfrak g_U^\alpha.
\]

Indeed, \(\mathfrak h\) is a submodule of \(\mathfrak g\), locally
a direct factor, containing \(\mathfrak g^0\), and stable under
\(T\).

\SGARefTarget{sga3:xxii:section:5.3}
\subsection*{\SGAInlineRefTarget{sga3:xxii:subsection:5-3}5.3. Strict Transporter of Two Subgroups of Type
\textup{(R)}. Applications}
\addcontentsline{toc}{subsection}{\SGARefLink{sga3:xxii:subsection:5-3}{5.3}. Strict Transporter of Two
Subgroups of Type (R). Applications}

\SGARefTarget{sga3:xxii:recall:5.3.0}\paragraph{Recall \SGARefLink{sga3:xxii:recall:5.3.0}{5.3.0}.}
\footnote{Editor's note: This reminder was added; it is used in the
proofs of \SGARefLink{sga3:xxii:prop:5.3.1}{5.3.1} and \SGARefLink{sga3:xxii:cor:5.3.4}{5.3.4}.} Let \(S\) be a scheme, let \(G\) be a
smooth \(S\)-group, put \(\mathfrak g=\Lie(G/S)\), and let
\(\mathfrak h\) be an \(\mathcal O_S\)-submodule of \(\mathfrak g\)
that is locally a direct factor. The \(\mathcal O_S\)-algebra
\[
  \mathscr A=\operatorname{Sym}(\omega^1_{G/S})
\]
is locally free; hence the \(S\)-scheme
\[
  \underline{\operatorname{Lie}}(G/S)
  =
  W(\mathfrak g)
  =
  \operatorname{Spec}(\mathscr A)
\]
is essentially free in the sense of \SGARefLink{sga3:viii:definition:6:1}{Exposé VIII, 6.1}. Since
\(W(\mathfrak h)\) is a closed subscheme of
\(\underline{\operatorname{Lie}}(G/S)\), of finite presentation
over \(\underline{\operatorname{Lie}}(G/S)\), it follows from
\SGARefLink{sga3:viii:examples:6:5}{Exposé VIII, 6.5} \textup{\SGARefLink{sga3:viii:examples:6:5:item:a:01}{(a)}} that
\[
  N=\underline{\operatorname{Norm}}_G(\mathfrak h)
\]
is represented by a closed group subscheme of \(G\), of finite
presentation over \(G\). (See also the additions \SGARefLink{sga3:xxii:par:6.2.3}{6.2.3} and \SGARefLink{sga3:xxii:prop:6.2.4}{6.2.4}
\textup{\SGARefLink{sga3:viii:examples:6:5:item:a:01}{(a)}} in Exposé VI\textsubscript{B}.) On the other hand, by
\SGARefLink{sga3:ii:theorem:5-3-1}{Exposé II, 5.3.1},
\[
  \underline{\operatorname{Lie}}(N/S)
  =
  \underline{\operatorname{Norm}}_
    {\underline{\operatorname{Lie}}(G/S)}(\mathfrak h).
\]

Finally, by Exposé VI\textsubscript{B}, \SGARefLink{thm:sga3-VIB-3.10}{3.10}, if \(N\) is smooth
over \(S\) at the points of the unit section, then the group
subfunctor \(N^0\) (defined in Exposé VI\textsubscript{B}, \SGARefLink{def:sga3-VIB-3.1}{3.1}) is
represented by an open group subscheme of \(N\), smooth over \(S\).

\SGARefTarget{sga3:xxii:prop:5.3.1}\paragraph{\SGARefLink{sga3:xxii:prop:5.3.1}{Proposition 5.3.1}.}
\textit{Let \(S\) be a scheme, let \(G\) be an \(S\)-group of type
\textup{(RA)} \textup{(\SGARefLink{sga3:xxii:def:5.1.6}{5.1.6})}, let \(H\) be a subgroup of type
\textup{(R)} of \(G\), and let
\(\mathfrak g\supset\mathfrak h\) be their Lie algebras. Then
\(\underline{\operatorname{Norm}}_G(\mathfrak h)\)
\textup{(which, by \SGARefLink{sga3:xxii:recall:5.3.0}{5.3.0}, is representable by a closed subscheme of
\(G\), of finite presentation over \(S\))} is smooth over \(S\) at
every point of the unit section, and}
\[
  \underline{\operatorname{Norm}}_G(\mathfrak h)^0=H.
\]

\footnote{Editor's note: We have expanded the original in the
argument that follows.} \emph{Proof.} Put
\[
  N=\underline{\operatorname{Norm}}_G(\mathfrak h),
  \qquad
  \mathfrak n=\Lie(N/S).
\]
We have \(H\subset N\), and by \SGARefLink{sga3:ii:theorem:5-3-1}{Exposé II, 5.3.1}, for every
\(s\in S\),
\[
  \mathfrak h(s)
  \subset
  \mathfrak n(s)
  =
  \operatorname{Norm}_{\mathfrak g(s)}(\mathfrak h(s)).
\]
Now, by \SGARefLink{sga3:xxii:lemma:5.3.2}{5.3.2} below,
\[
  \mathfrak h(s)
  =
  \operatorname{Norm}_{\mathfrak g(s)}(\mathfrak h(s)),
\]
and since \(H\) is smooth over \(S\),
\[
  \dim_{\kappa(s)}\mathfrak h(s)=\dim H_s
\]
\textup{(cf.\ \textup{[DG70]}, \S II.5, Theorem 2.1)}. Hence
\[
  \dim_{\kappa(s)}\mathfrak n(s)
  =
  \dim_{\kappa(s)}\mathfrak h(s)
  =
  \dim H_s
  \leq
  \dim N_s.
\]
It follows that
\[
  N_s^0=H_s^0=H_s,
\]
since \(H\) has connected fibers. Consequently the group subfunctor
\(N^0\) \textup{(defined in Exposé VI\textsubscript{B}, \SGARefLink{def:sga3-VIB-3.1}{3.1})} is
represented by the smooth \(S\)-group \(H\). This proves \SGARefLink{sga3:xxii:prop:5.3.1}{5.3.1},
modulo the following lemma.

\SGARefTarget{sga3:xxii:lemma:5.3.2}\paragraph{\SGARefLink{sga3:xxii:lemma:5.3.2}{Lemma 5.3.2}.}
\textit{Under the conditions of \SGARefLink{sga3:xxii:lemma:5.2.7}{5.2.7}, if \(G\) is of type
\textup{(RA)}, then, for every \(s\in S\),}
\[
  \operatorname{Norm}_{\mathfrak g(s)}(\mathfrak h(s))
  =
  \mathfrak h(s).
\]

Indeed, we are reduced to the case where \(S\) is the spectrum of a
field, and hence where
\[
  \mathfrak h=\mathfrak g_{R'}
\]
for some \(R'\subset R\). But we already have
\[
  \operatorname{Transp}_{\mathfrak g}
    (\mathfrak t,\mathfrak h)
  =
  \mathfrak h.
\]
Indeed, if \(H\in\mathfrak t\) and
\(X\in\mathfrak g^\alpha\), then
\[
  [H,X]=\bar\alpha(H)X,
\]
where \(\bar\alpha:\mathfrak t\to\mathcal O_S\) is the morphism
derived from \(\alpha\). Condition \textup{(iv$'$)} says precisely
that \(\bar\alpha\ne0\) for every \(\alpha\in R\).

\SGARefTarget{sga3:xxii:cor:5.3.3}\paragraph{\SGARefLink{sga3:xxii:cor:5.3.3}{Corollary 5.3.3}.}
\textit{Let \(S\) be a scheme, let \(G\) be an \(S\)-group of type
\textup{(RA)}, let \(H\) and \(H'\) be two subgroups of type
\textup{(R)} of \(G\), and let \(\mathfrak h\) and
\(\mathfrak h'\) be their Lie algebras. Then}
\[
  H=H'
  \quad\Longleftrightarrow\quad
  \mathfrak h=\mathfrak h'.
\]

\SGARefTarget{sga3:xxii:cor:5.3.4}\paragraph{\SGARefLink{sga3:xxii:cor:5.3.4}{Corollary 5.3.4}.}
\textit{Under the conditions of \SGARefLink{sga3:xxii:lemma:5.2.7}{5.2.7}, with \(G\) of type
\textup{(RA)}, the maps}
\[
  H\longmapsto\Lie(H/S),
  \qquad
  \mathfrak h\longmapsto
    \underline{\operatorname{Norm}}_G(\mathfrak h)^0
\]
\textit{give a bijective correspondence between the set of
subgroups of type \textup{(R)} of \(G\) containing \(T\), and the
set of Lie subalgebras of \(\mathfrak g\) containing
\(\mathfrak g^0\), stable under \(T\), and whose normalizer in
\(G\) is smooth over \(S\) at every point of the unit section.}

Indeed,\footnote{Editor's note: The proof that follows was added.}
let \(\mathfrak h\) be a Lie subalgebra of \(\mathfrak g\)
satisfying the properties above. By \SGARefLink{sga3:xxii:recall:5.3.0}{5.3.0},
\[
  H=\underline{\operatorname{Norm}}_G(\mathfrak h)^0
\]
is a smooth \(S\)-group scheme. Moreover, since
\[
  C=\underline{\operatorname{Centr}}_G(T)
\]
stabilizes every \(\mathfrak g^\alpha\) and has connected fibers
\textup{(\SGARefLink{sga3:xii:6.6}{Exposé XII, 6.6})}, we have \(C\subset H\). Therefore
\(H\) is a subgroup of \(G\) of type \textup{(R)}. By \SGARefLink{sga3:ii:theorem:5-3-1}{Exposé II,
5.3.1},
\[
  \underline{\operatorname{Lie}}(H)
  =
  \underline{\operatorname{Norm}}_{\mathfrak g}(\mathfrak h).
\]
Finally, by the proof of \SGARefLink{sga3:xxii:lemma:5.3.2}{5.3.2},
\[
  \underline{\operatorname{Norm}}_{\mathfrak g}(\mathfrak h)
  =
  \mathfrak h.
\]

\SGARefTarget{sga3:xxii:cor:5.3.5}\paragraph{\SGARefLink{sga3:xxii:cor:5.3.5}{Corollary 5.3.5}.}
\textit{Let \(S\) be a scheme, let \(G\) be an \(S\)-group of type
\textup{(RR)} \textup{(\SGARefLink{sga3:xxii:def:5.1.1}{5.1.1})}, let \(T\) be a maximal torus of
\(G\), and let \(H\) and \(H'\) be two subgroups of type
\textup{(R)} of \(G\), both containing \(T\). Then}
\[
  H=H'
  \quad\Longleftrightarrow\quad
  \mathfrak h=\mathfrak h'.
\]

By virtue of the finite-presentation hypotheses, we reduce as usual
\textup{(cf.\ EGA IV\textsubscript{3}, \S8, and Exposé
VI\textsubscript{B}, \S10)} to the case where \(S\) is noetherian.
It is then enough to verify that
\(\mathfrak h=\mathfrak h'\) implies
\[
  H_{S'}=H'_{S'}
\]
for every \(S'\) that is the spectrum of an artinian quotient of a
local ring of \(S\);\footnote{Editor's note: Indeed, let \(g\in G\),
let \(s\) be its image in \(S\), let \(\mathfrak m\) be the maximal
ideal of \(\mathcal O_{G,g}\), let \(\mathfrak n\) be that of
\(\mathcal O_{S,s}\), and let \(I\) (respectively \(I'\)) be the
kernel of the morphism from \(\mathcal O_{G,g}\) to
\(\mathcal O_{H,g}\) (respectively \(\mathcal O_{H',g}\)), the
latter being the zero ring if \(g\notin H\) (respectively
\(g\notin H'\)). Since \(\mathcal O_{G,g}\) is noetherian, \(I\)
and \(I'\) are closed for the \(\mathfrak m\)-adic topology, hence
\emph{a fortiori} for the \(\mathfrak n\)-adic topology. It is
therefore enough to prove
\[
  I+\mathfrak n^n\mathcal O_{G,g}
  =
  I'+\mathfrak n^n\mathcal O_{G,g}
\]
for every \(n\in\mathbb N\).} we are therefore reduced to the case
where \(S\) is artinian, in which case we may apply \SGARefLink{sga3:xxii:prop:5.1.8}{5.1.8}.

Let
\[
  u:G\longrightarrow G'=G/Z
\]
be the canonical morphism, and let \(T'=T/Z\) be the maximal torus
of \(G'\) corresponding to \(T\). By \SGARefLink{sga3:xii:7.12}{Exposé XII, 7.12}, there exist
subgroups \(H_1\) and \(H'_1\) of type \textup{(R)} of \(G'\),
containing \(T'\), such that
\[
  H=u^{-1}(H_1),
  \qquad
  H'=u^{-1}(H'_1).
\]
It is enough to prove that \(H_1=H'_1\). But by \SGARefLink{sga3:xxii:lemma:5.2.7}{5.2.7} and \SGARefLink{sga3:xxii:lemma:5.1.5}{5.1.5},
\[
  \Lie(H_1)=\Lie(H'_1),
\]
and we are reduced to \SGARefLink{sga3:xxii:cor:5.3.3}{5.3.3}.

\SGARefTarget{sga3:xxii:rem:5.3.6}\paragraph{\SGARefLink{sga3:xxii:rem:5.3.6}{Remark 5.3.6}.}
The fact that \(H\) and \(H'\) contain the same maximal torus is
essential for the validity of \SGARefLink{sga3:xxii:cor:5.3.5}{5.3.5} when \(G\) is not of type
\textup{(RA)}. For example, consider the maximal tori of
\(\operatorname{SL}_{2,k}\), where \(k\) has characteristic
\(2\).\footnote{Editor's note: Indeed, if \(k\) is algebraically
closed, all maximal tori of \(G=\operatorname{SL}_{2,k}\) are
conjugate under \(G(k)\), and they all have as Lie algebra the line
\(k\,\mathrm{id}\subset M_2(k)\), which is invariant under the
adjoint action.}

\SGARefTarget{sga3:xxii:cor:5.3.7}\paragraph{\SGARefLink{sga3:xxii:cor:5.3.7}{Corollary 5.3.7}.}
\textit{Let \(S\) be a scheme, let \(G\) be an \(S\)-group of type
\textup{(RR)}, let \(T\) be a maximal torus of \(G\), and let \(H\)
and \(H'\) be two subgroups of type \textup{(R)} of \(G\), both
containing \(T\). The set \(U\) of \(s\in S\) such that
\(H_s=H'_s\) is open and closed in \(S\), and}
\[
  H_U=H'_U.
\]

Indeed, this follows immediately from \SGARefLink{sga3:xxii:cor:5.3.5}{5.3.5} and \SGARefLink{sga3:xxii:lemma:5.2.7}{5.2.7}.

\SGARefTarget{sga3:xxii:cor:5.3.8}\paragraph{\SGARefLink{sga3:xxii:cor:5.3.8}{Corollary 5.3.8}.}
\textit{The ``functor of subgroups of type \textup{(R)} containing
\(T\)'', where \(T\) is a given maximal torus in a group \(G\) of
type \textup{(RR)}, is formally unramified \textup{(\SGARefLink{sga3:xi:definition:1:1}{Exposé XI,
1.1})}.}

\SGARefTarget{sga3:xxii:thm:5.3.9}\paragraph{\SGARefLink{sga3:xxii:thm:5.3.9}{Theorem 5.3.9}.}
\textit{Let \(G\) be an \(S\)-group of type \textup{(RR)}
\textup{(\SGARefLink{sga3:xxii:def:5.1.1}{5.1.1})}, and let \(H\) and \(H'\) be two subgroups of type
\textup{(R)} \textup{(\SGARefLink{sga3:xxii:def:5.2.1}{5.2.1})}. Let
\(\underline{\operatorname{Transt}}_G(H,H')\) be the strict
transporter of \(H\) into \(H'\), defined by}
\[
  \underline{\operatorname{Transt}}_G(H,H')(S')
  =
  \bigl\{
    g\in G(S')
    \mathrel{\big|}
    \operatorname{int}(g)H_{S'}=H'_{S'}
  \bigr\}.
\]
\textit{Then \(\underline{\operatorname{Transt}}_G(H,H')\) is
representable by a closed subscheme of \(G\), smooth and of finite
presentation over \(S\).}

The fact that
\(\underline{\operatorname{Transt}}_G(H,H')\) is representable by
a closed subscheme of \(G\), of finite presentation over \(S\),
follows from \SGARefLink{sga3:xi:corollary:6:11}{Exposé XI, 6.11} \textup{\SGARefLink{sga3:xi:corollary:6:11:item:a}{(a)}}. To prove that it is
smooth over \(S\), it is necessary to prove that if \(S\) is affine,
if \(S_0\) is the closed subscheme defined by a nilpotent ideal
\(J\), and if
\[
  g_0\in G(S_0),
  \qquad
  \operatorname{int}(g_0)H_0=H'_0,
\]
then there exists \(g\in G(S)\), lifting \(g_0\), such that
\[
  \operatorname{int}(g)H=H'.
\]
Since smoothness is local for the étale topology, we may suppose
that \(H\) contains a maximal torus \(T\) of \(G\).

Then \(T_0\) is a maximal torus of \(H_0\), and hence
\(\operatorname{int}(g_0)T_0\) is a maximal torus of \(H'_0\).
By \SGARefLink{sga3:ix:theorem:3:6:bis}{Exposé IX, 3.6 bis}, there exists a torus \(T'\) of \(H'\) such
that
\[
  T'_0=\operatorname{int}(g_0)T_0;
\]
by \SGARefLink{sga3:ix:corollary:3:3:bis}{Exposé IX, 3.3 bis}, there therefore exists \(g\in G(S)\), lifting
\(g_0\), such that
\[
  \operatorname{int}(g)T=T'.
\]
After replacing \(H\) by \(\operatorname{int}(g)H\), we may
therefore suppose that \(H\) and \(H'\) contain the same maximal
torus \(T\) and that \(H_0=H'_0\). But then \(H=H'\) by \SGARefLink{sga3:xxii:cor:5.3.7}{5.3.7}.
\qed

\SGARefTarget{sga3:xxii:cor:5.3.10}\paragraph{\SGARefLink{sga3:xxii:cor:5.3.10}{Corollary 5.3.10}.}
\textit{Let \(G\) be an \(S\)-group of type \textup{(RR)}, and let
\(H\) be a subgroup of type \textup{(R)} of \(G\). Then
\(\underline{\operatorname{Norm}}_G(H)\) is representable by a
closed group subscheme of \(G\), of finite presentation and smooth
over \(S\).}

Using now the reasoning employed in Exposé XI to deduce \SGARefLink{sga3:xi:corollary:5:4:bis}{5.4 bis}
from \SGARefLink{sga3:xi:corollary:5:2:bis}{5.2 bis}, we obtain:

\SGARefTarget{sga3:xxii:cor:5.3.11}\paragraph{\SGARefLink{sga3:xxii:cor:5.3.11}{Corollary 5.3.11}.}
\textit{Under the hypotheses of \SGARefLink{sga3:xxii:thm:5.3.9}{5.3.9}, the following conditions are
equivalent:}
\begin{enumerate}
  \item[\textup{\SGAInlineRefTarget{sga3:xxii:cor:5.3.11:item:i}(i)}] \textit{\(H\) and \(H'\) are locally conjugate
  in \(G\) for the étale topology;}

  \item[\textup{(i bis)}] \textit{the same holds for the fpqc
  topology;}

  \item[\textup{(ii)}] \textit{for every \(s\in S\),
  \(H_{\bar s}\) and \(H'_{\bar s}\) are conjugate by an element of
  \(G(\bar s)\);}

  \item[\textup{(ii bis)}] \textit{the structural morphism}
  \[
    \underline{\operatorname{Transt}}_G(H,H')
      \longrightarrow S
  \]
  \textit{is surjective;}

  \item[\textup{\SGAInlineRefTarget{sga3:xxii:cor:5.3.11:item:iii}(iii)}]
  \textit{\(\underline{\operatorname{Transt}}_G(H,H')\) is a
  principal homogeneous bundle under the action of the smooth
  \(S\)-group scheme of finite presentation
  \(\underline{\operatorname{Norm}}_G(H)\).}
\end{enumerate}

We merely note that the nontrivial implication
\textup{\SGARefLink{sga3:xxii:cor:5.3.11:item:iii}{(iii)}}\(\Rightarrow\)\textup{\SGARefLink{sga3:xxii:cor:5.3.11:item:i}{(i)}} is Hensel's lemma.

Using now \textit{Bible}, \S6.4, Theorem 4
\textup{(= [Ch05], \S6.5, Theorem 5)} and \S9.3,
Theorem 1, we obtain from \SGARefLink{sga3:xxii:cor:5.3.10}{5.3.10} and \SGARefLink{sga3:xxii:cor:5.3.11}{5.3.11}:

\SGARefTarget{sga3:xxii:cor:5.3.12}\paragraph{\SGARefLink{sga3:xxii:cor:5.3.12}{Corollary 5.3.12}.}
\textit{Let \(G\) be an \(S\)-group of type \textup{(RR)}. The
Borel subgroups of \(G\) are closed in \(G\), equal to their
normalizers, and locally conjugate for the étale topology.}

\SGARefTarget{sga3:xxii:def:5.3.13}\paragraph{\SGARefLink{sga3:xxii:def:5.3.13}{Definition 5.3.13}.}
\textit{Let \(S\) be a scheme, and let \(G\) be a smooth
\(S\)-group scheme of finite presentation with connected
fibers.\footnote{Editor's note: See note 25.} A \emph{Killing pair}
of \(G\) is a pair \(T\subset B\), where \(T\) is a maximal torus
of \(G\) and \(B\) is a Borel subgroup of \(G\) containing \(T\).}

Using now the conjugacy of maximal tori in \(B\)
\textup{(cf.\ \SGARefLink{sga3:xxii:rem:5.1.2}{5.1.2} \textup{\SGARefLink{sga3:xxii:rem:5.1.2:item:a}{(a)}} and \SGARefLink{sga3:xxii:prop:5.2.6}{5.2.6}, for example)}, we have:

\SGARefTarget{sga3:xxii:cor:5.3.14}\paragraph{\SGARefLink{sga3:xxii:cor:5.3.14}{Corollary 5.3.14}.}
\textit{Let \(G\) be an \(S\)-group of type \textup{(RR)}. The
Killing pairs of \(G\) are locally conjugate for the étale topology.}

\SGARefTarget{sga3:xxii:cor:5.3.15}\paragraph{\SGARefLink{sga3:xxii:cor:5.3.15}{Corollary 5.3.15}.}
\textit{Let \(G\) be an \(S\)-group of type \textup{(RR)}. Let
\(T\) be a maximal torus of \(G\), and let}
\[
  W_G(T)
  =
  \underline{\operatorname{Norm}}_G(T)/
  \underline{\operatorname{Centr}}_G(T)
\]
\textit{be the corresponding Weyl group \textup{(\SGARefLink{sga3:xix:proposition:6:3}{Exposé XIX,
6.3})}. The ``functor of Borel subgroups of \(G\) containing \(T\)''
is formally principal homogeneous under \(W_G(T)\).}

This follows immediately from \SGARefLink{sga3:xxii:cor:5.3.14}{5.3.14} and from the fact that, if
\(B\) is a Borel subgroup of \(G\) containing \(T\), then
\[
  \underline{\operatorname{Norm}}_G(T)\cap B
  =
  \underline{\operatorname{Centr}}_G(T),
\]
cf.\ \SGARefLink{sga3:xiv:4.4}{Exposé XIV, 4.4}.

\SGARefTarget{sga3:xxii:prop:5.3.16}\paragraph{\SGARefLink{sga3:xxii:prop:5.3.16}{Proposition 5.3.16}.}
\textit{Let \(G\) be an \(S\)-group of type \textup{(RR)}, let
\(H\) be a subgroup of type \textup{(R)}, and let
\(N=\underline{\operatorname{Norm}}_G(H)\) be its normalizer
\textup{(\SGARefLink{sga3:xxii:cor:5.3.10}{5.3.10})}. Let \(T\) be a maximal torus of \(H\), and let
\(W_H(T)\) and \(W_N(T)\) be the corresponding Weyl groups
\textup{(étale, quasi-finite, and separated by \SGARefLink{sga3:xix:proposition:6:3}{Exposé XIX, 6.3})}.
There is the following exact sequence of sheaves \textup{(for the
étale topology)}, whose morphisms are induced by}
\[
  \underline{\operatorname{Norm}}_H(T)
  \longrightarrow
  \underline{\operatorname{Norm}}_N(T)
  \longrightarrow
  N/H:
\]
\[
  1
  \longrightarrow W_H(T)
  \longrightarrow W_N(T)
  \longrightarrow N/H
  \longrightarrow 1.
\]

The only nontrivial point is that the last arrow is an epimorphism.
Let \(n\in N(S')\), where \(S'\to S\). The two maximal tori \(T\)
and \(\operatorname{int}(n)T\) of \(H\) are locally conjugate in
\(H\) for the étale topology. Thus there is a covering family
\(\{S'_i\to S'\}\) and, for every \(i\), an
\(h_i\in H(S'_i)\) such that
\[
  \operatorname{int}(h_i)T=\operatorname{int}(n)T.
\]
Consequently
\[
  nh_i^{-1}\in \underline{\operatorname{Norm}}_N(T),
\]
which gives the desired result.

\SGARefTarget{sga3:xxii:rem:5.3.17}\paragraph{\SGARefLink{sga3:xxii:rem:5.3.17}{Remark 5.3.17}.}
The group \(W_N(T)\) may be described as follows. Suppose that we
have reduced to the situation of \SGARefLink{sga3:xxii:lemma:5.2.7}{5.2.7}, with
\(\mathfrak h=\mathfrak g_{R'}\). Then \(W_N(T)\) is equal to
\(\underline{\operatorname{Norm}}_W(R')\), the sheaf of sections
of \(W=W_G(T)\) which, acting on \(R\), normalize \(R'\). Indeed,
by \SGARefLink{sga3:xxii:cor:5.3.5}{5.3.5},
\[
  \underline{\operatorname{Norm}}_N(T)
  =
  \underline{\operatorname{Norm}}_G(H)
    \cap \underline{\operatorname{Norm}}_G(T)
  =
  \underline{\operatorname{Norm}}_G(\mathfrak h)
    \cap \underline{\operatorname{Norm}}_G(T).
\]

\SGARefTarget{sga3:xxii:cor:5.3.18}\paragraph{\SGARefLink{sga3:xxii:cor:5.3.18}{Corollary 5.3.18}.}
\textit{Let \(G\) be an \(S\)-group of type \textup{(RR)}, and let
\(H\) be a subgroup of type \textup{(R)}. Suppose that ``the Weyl
groups of \(G\) are finite'', that is, that for every \(S'\to S\)
and every maximal torus \(T\) of \(G_{S'}\), the étale
\(S'\)-scheme}
\[
  \underline{\operatorname{Norm}}_{G_{S'}}(T)/
  \underline{\operatorname{Centr}}_{G_{S'}}(T)
\]
\textit{is finite \textup{(cf.\ \SGARefLink{sga3:xix:proposition:6:3}{Exposé XIX, 6.3} \textup{\SGARefLink{sga3:xxii:cor:5.3.18:item:iii}{(iii)}})}.
Then the following conditions are equivalent:}
\begin{enumerate}
  \item[\textup{\SGAInlineRefTarget{sga3:xxii:cor:5.3.18:item:i}(i)}] \textit{\(H\) is closed in \(G\);}

  \item[\textup{\SGAInlineRefTarget{sga3:xxii:cor:5.3.18:item:ii}(ii)}]
  \textit{\(\underline{\operatorname{Norm}}_G(H)/H\) is
  representable by a finite étale \(S\)-scheme;}

  \item[\textup{\SGAInlineRefTarget{sga3:xxii:cor:5.3.18:item:iii}(iii)}]
  \textit{``the Weyl groups of \(H\) are finite''.}
\end{enumerate}

Indeed, we may suppose that \(H\) has a maximal torus \(T\). By
\SGARefLink{sga3:xxii:cor:5.3.10}{5.3.10}, \(N=\underline{\operatorname{Norm}}_G(H)\) is closed in
\(G\); therefore \(W_N(T)\) is closed in \(W_G(T)\), and hence is
finite over \(S\). Clearly \textup{\SGARefLink{sga3:xxii:cor:5.3.18:item:i}{(i)}} implies \textup{\SGARefLink{sga3:xxii:cor:5.3.18:item:iii}{(iii)}}, and
\textup{\SGARefLink{sga3:xxii:cor:5.3.18:item:iii}{(iii)}} implies \textup{\SGARefLink{sga3:xxii:cor:5.3.18:item:ii}{(ii)}} by the exact sequence of
\SGARefLink{sga3:xxii:prop:5.3.16}{5.3.16}. Finally, \textup{\SGARefLink{sga3:xxii:cor:5.3.18:item:ii}{(ii)}} implies \textup{\SGARefLink{sga3:xxii:cor:5.3.18:item:i}{(i)}}, since if
\(N/H\) is finite, it is separated, so \(H\) is closed in \(N\),
and hence in \(G\).

\SGARefTarget{sga3:xxii:rem:5.3.19}\paragraph{\SGARefLink{sga3:xxii:rem:5.3.19}{Remark 5.3.19}.}
When \(G\) is reductive, the preceding conditions on \(H\) seem
always to be satisfied. We prove them below in most cases.

\SGARefTarget{sga3:xxii:section:5.4}
\subsection*{\SGAInlineRefTarget{sga3:xxii:subsection:5-4}5.4. Subgroups of Type \textup{(R)} of a Split
Reductive Group \textup{(Generalities)}}
\addcontentsline{toc}{subsection}{\SGARefLink{sga3:xxii:subsection:5-4}{5.4}. Subgroups of Type (R) of a
Split Reductive Group (Generalities)}

\SGARefTarget{sga3:xxii:5.4.1}\paragraph{\SGARefLink{sga3:xxii:5.4.1}{5.4.1}.}
If \(H\) is a subgroup of type \textup{(R)} of the reductive group
\(G\), then \(H\) locally for the étale topology contains a maximal
torus of \(G\) \textup{(\SGARefLink{sga3:xxii:recall:5.2.2}{5.2.2})}. By \SGARefLink{sga3:xxii:cor:2.3}{2.3}, locally for the étale
topology we may suppose that \(G\) is split relative to this torus.
Thus let \((G,T,M,R)\) be a split \(S\)-group, and let \(H\) be a
subgroup of type \textup{(R)} of \(G\) containing \(T\). By \SGARefLink{sga3:xxii:cor:5.3.5}{5.3.5},
such a subgroup is characterized by its Lie algebra, which
\textup{(\SGARefLink{sga3:xxii:lemma:5.2.7}{5.2.7})} is locally on \(S\) of the form
\(\mathfrak g_{R'}\):
\[
  \mathfrak g_{R'}
  =
  \mathfrak t
  \oplus
  \prod_{\alpha\in R'}\mathfrak g^\alpha.
\]

\SGARefTarget{sga3:xxii:def:5.4.2}\paragraph{\SGARefLink{sga3:xxii:def:5.4.2}{Definition 5.4.2}.}
\textit{Let \((G,T,M,R)\) be a split \(S\)-group. A subset \(R'\)
of \(R\) is said to be of type \textup{(R)} if
\(\mathfrak g_{R'}\) is the Lie algebra of a subgroup of type
\textup{(R)} of \(G\) containing \(T\). This subgroup, uniquely
determined by \(R'\), is denoted \(H_{R'}\).}

\SGARefTarget{sga3:xxii:lem:5.4.3}\paragraph{\SGARefLink{sga3:xxii:lem:5.4.3}{Lemma 5.4.3}.}
\textit{Under the preceding conditions, the following equivalences
hold:}
\[
\begin{aligned}
  H\cap Z_\alpha=T
    &\quad\Longleftrightarrow\quad
      \alpha\notin R',\ -\alpha\notin R',\\
  H\supset U_\alpha
    &\quad\Longleftrightarrow\quad
      \alpha\in R',\\
  H\cap U_\alpha=e
    &\quad\Longleftrightarrow\quad
      \alpha\notin R',\\
  H\supset Z_\alpha
    &\quad\Longleftrightarrow\quad
      \alpha\in R',\ -\alpha\in R'.
\end{aligned}
\]

Indeed, \(H\cap Z_\alpha\) is a subgroup of type \textup{(R)} of
\(Z_\alpha\), by \SGARefLink{sga3:xxii:prop:5.2.5}{5.2.5}; but a subgroup of type \textup{(R)} of
\(Z_\alpha\) containing \(T\) is locally equal to one of the
following subgroups:
\[
  T,\qquad T\cdot U_\alpha,\qquad
  T\cdot U_{-\alpha},\qquad Z_\alpha,
\]
by \SGARefLink{sga3:xxii:cor:5.3.5}{5.3.5}.

\SGARefTarget{sga3:xxii:lem:5.4.4}\paragraph{\SGARefLink{sga3:xxii:lem:5.4.4}{Lemma 5.4.4}.}
\textit{Under the conditions of \SGARefLink{sga3:xxii:def:5.4.2}{5.4.2}, let \(R_+\) be a system of
positive roots; choose orderings on \(R'\cap R_+\) and
\(R'\cap -R_+\). The morphism}
\[
  \Omega_{R_+,R'}
  =
  \prod_{\alpha\in R'\cap -R_+}U_\alpha
  \mathop{\times}\limits_S T
  \mathop{\times}\limits_S
  \prod_{\alpha\in R'\cap R_+}U_\alpha
  \longrightarrow G,
\]
\textit{induced by multiplication in \(G\), induces an open
immersion}
\[
  \Omega_{R_+,R'}\longrightarrow H_{R'}.
\]

Indeed, by \SGARefLink{sga3:xxii:lem:5.4.3}{5.4.3} this morphism factors through \(H_{R'}\), and
therefore induces an immersion
\(\Omega_{R_+,R'}\to H_{R'}\). We then argue as in \SGARefLink{sga3:xxii:par:4.1.1}{4.1.1}.

\SGARefTarget{sga3:xxii:prop:5.4.5}\paragraph{\SGARefLink{sga3:xxii:prop:5.4.5}{Proposition 5.4.5}.}
\textit{Let \((G,T,M,R)\) be a split \(S\)-group. Let \(R'\) and
\(R''\) be two subsets of \(R\) of type \textup{(R)}.}
\begin{enumerate}
  \item[\textup{\SGAInlineRefTarget{sga3:xxii:prop:5.4.5:item:i}(i)}]
  \textit{\(H_{R'}\cap H_{R''}\) is smooth at every point of the
  unit section, \(R'\cap R''\) is of type \textup{(R)}, and}
  \[
    \bigl(H_{R'}\cap H_{R''}\bigr)^0
    =
    H_{R'\cap R''}.
  \]

  \item[\textup{\SGAInlineRefTarget{sga3:xxii:prop:5.4.5:item:ii}(ii)}]
  \textit{The following equivalence holds:}
  \[
    H_{R'}\subset H_{R''}
    \quad\Longleftrightarrow\quad
    R'\subset R''.
  \]
\end{enumerate}

Indeed, \textup{\SGARefLink{sga3:xxii:prop:5.4.5:item:ii}{(ii)}} follows immediately from \textup{\SGARefLink{sga3:xxii:prop:5.4.5:item:i}{(i)}}. To
prove \textup{\SGARefLink{sga3:xxii:prop:5.4.5:item:i}{(i)}}, it suffices to prove that
\(H_{R'}\cap H_{R''}\) is smooth at every point of the unit section:
its neutral component \textup{(cf.\ Exposé
VI\textsubscript{B}, \SGARefLink{thm:sga3-VIB-3.10}{3.10})} will then be a subgroup of type
\textup{(R)} containing \(T\), and hence equal to
\(H_{R'\cap R''}\). But
\[
  \Omega_{R_+,R'}\cap\Omega_{R_+,R''}
  =
  \Omega_{R_+,R'\cap R''}
\]
is an open subset of \(H_{R'}\cap H_{R''}\), contains the unit
section, and is smooth over \(S\).

\SGARefTarget{sga3:xxii:cor:5.4.6}\paragraph{\SGARefLink{sga3:xxii:cor:5.4.6}{Corollary 5.4.6}.}
\textit{Let \(S\) be a scheme, \(G\) a reductive \(S\)-group, \(T\)
a maximal torus of \(G\), and \(s\) a point of \(S\). If \(H\) and
\(H'\) are two subgroups of type \textup{(R)} containing \(T\), and
if \(H_s\subset H'_s\), then there is an open subset \(U\) of \(S\)
containing \(s\) such that \(H_U\subset H'_U\).}

Indeed, we may suppose that \(G\) is split relative to \(T\), and
the assertion then follows immediately from \SGARefLink{sga3:xxii:prop:5.4.5}{5.4.5} \textup{\SGARefLink{sga3:xxii:prop:5.4.5:item:ii}{(ii)}}.

We are thus led to ask which subsets \(R'\) of \(R\) are of type
\textup{(R)}. We may suppose that the group is adjoint; one must
then verify that \(\mathfrak g_{R'}\) is a Lie algebra and that its
normalizer is smooth at every point of the unit section. The most
important case is given by the following theorem.

\SGARefTarget{sga3:xxii:thm:5.4.7}\paragraph{\SGARefLink{sga3:xxii:thm:5.4.7}{Theorem 5.4.7}.}
\textit{Every closed subset \(R'\) of \(R\) is of type
\textup{(R)}.}

Recall \textup{(\SGARefLink{sga3:XXI:3.1.4}{Exposé XXI, 3.1.4})} that a subset \(R'\) of \(R\)
is said to be \emph{closed} if
\[
  \alpha,\beta\in R',\quad \alpha+\beta\in R
  \quad\Longrightarrow\quad
  \alpha+\beta\in R'.
\]

\SGARefTarget{sga3:xxii:rem:5.4.8}\paragraph{\SGARefLink{sga3:xxii:rem:5.4.8}{Remark 5.4.8}.}
We shall see later \textup{(\SGARefLink{sga3:XXIII:6.6}{Exposé XXIII, 6.6})} that, if
\(6\cdot1_S\ne0\)\footnote{Editor's note: In fact, it suffices
\textup{(cf.\ loc.\ cit.)} that \(2\) and \(3\) be nonzero on
\(S\).}---for example, if \(S\) has residue characteristics
different from \(2\) and \(3\)---the fact that
\(\mathfrak g_{R'}\) is a Lie algebra already implies that \(R'\)
is closed. In that case a subset \(R'\) is of type \textup{(R)} if
and only if it is closed. Thus \SGARefLink{sga3:xxii:thm:5.4.7}{Theorem 5.4.7} gives all the subsets
of type \textup{(R)} that are ``independent of the
characteristic.''

Let us first prove the following lemma.

\SGARefTarget{sga3:xxii:lem:5.4.9}\paragraph{\SGARefLink{sga3:xxii:lem:5.4.9}{Lemma 5.4.9}.}
\textit{For each \(\alpha\in R\), choose an element}
\[
  X_\alpha\in\Gamma(S,\mathfrak g^\alpha)^\times.
\]
\textit{Let \(\alpha,\beta\in R\), with
\(\alpha+\beta\ne0\), and let \(q\) be the largest integer \(i\)
such that \(\beta+i\alpha\in R\).\begingroup
\renewcommand{\thefootnote}{*}\footnote{Source note: the French
text prints \(\alpha+i\beta\) at this point; the displayed
expansion and the proof require \(\beta+i\alpha\).}
\addtocounter{footnote}{-1}\endgroup There exist
uniquely determined sections}
\[
  M_{\alpha,\beta,i}\in\Gamma(S,\mathcal O_S)
\]
\textit{such that}
\[
  \operatorname{Ad}\bigl(\exp_\alpha(xX_\alpha)\bigr)(X_\beta)
  =
  X_\beta+
  \sum_{i=1}^{q}
    M_{\alpha,\beta,i}x^iX_{\beta+i\alpha}
\]
\textit{for every \(x\in\mathbb G_a(S')\), where \(S'\to S\).}

The map
\[
  x\longmapsto
  \operatorname{Ad}\bigl(\exp_\alpha(xX_\alpha)\bigr)(X_\beta)
\]
defines a morphism
\[
  \mathbb G_{a,S}\longrightarrow
  \mathbf W(\mathfrak g)\simeq\mathbb G_{a,S}^{m}.
\]
There are therefore uniquely determined sections
\(Y_n\in\Gamma(S,\mathfrak g)\) such that
\[
  \operatorname{Ad}\bigl(\exp_\alpha(xX_\alpha)\bigr)(X_\beta)
  =
  \sum_{n\geq0}x^nY_n.
\]
Applying the inner automorphism defined by a section \(t\) of
\(T\), we immediately find
\[
  \operatorname{Ad}(t)(Y_n)
  =
  \beta(t)\alpha(t)^nY_n,
\]
and hence \(Y_n\in\Gamma(S,\mathfrak g^{\beta+n\alpha})\). Since
\(\alpha\) and \(\beta\) are not proportional, none of the
\(\beta+n\alpha\) is zero. Thus \(Y_n=0\) for \(n>q\), and
\[
  Y_n=M_{\alpha,\beta,n}X_{\beta+n\alpha}
\]
for \(0\leq n\leq q\), where
\(M_{\alpha,\beta,n}\in\mathbb G_a(S)\) is uniquely determined.
Setting \(x=0\) in the formula obtained gives \(Y_0=X_\beta\),
which completes the proof.

\SGARefTarget{sga3:xxii:rem:5.4.10}\paragraph{\SGARefLink{sga3:xxii:rem:5.4.10}{Remark 5.4.10}.}
Differentiating the preceding formula at \(x=0\), we find
\[
  [X_\alpha,X_\beta]
  =
  \begin{cases}
    N_{\alpha,\beta}X_{\alpha+\beta},
      &\text{where \(N_{\alpha,\beta}=M_{\alpha,\beta,1}\), if
        \(\alpha+\beta\in R\),}\\
    0,
      &\text{if \(\alpha+\beta\notin R\) and
        \(\alpha+\beta\ne0\).}
  \end{cases}
\]

Let us now prove \SGARefLink{sga3:xxii:thm:5.4.7}{5.4.7}. If \(R'\) is a closed subset of \(R\), then
\(\mathfrak g_{R'}\) is a Lie subalgebra of \(\mathfrak g\), by
\SGARefLink{sga3:xxii:rem:5.4.10}{5.4.10} and \SGARefLink{sga3:xx:2.10}{Exposé XX, 2.10}, formula \textup{(3)}. By \SGARefLink{sga3:xxii:lem:5.4.9}{5.4.9} and
\SGARefLink{sga3:xx:2.10}{Exposé XX, 2.10}, formula \textup{\SGARefLink{sga3:xxii:section:2}{(2)}}, \(U_\alpha\) normalizes
\(\mathfrak g_{R'}\) for every \(\alpha\in R'\). Choose a system of
positive roots \(R_+\), and consider the open subset
\(\Omega_{R_+}\) of \SGARefLink{sga3:xxii:prop:4.1.2}{4.1.2}. Let \(\Omega_{R_+,R'}\) be the closed
subscheme of \(\Omega_{R_+}\) defined by
\[
  \Omega_{R_+,R'}
  =
  \left(
    \prod_{\alpha\in R'\cap -R_+}U_\alpha
  \right)
  \cdot T\cdot
  \left(
    \prod_{\alpha\in R'\cap R_+}U_\alpha
  \right).
\]
The canonical immersion \(\Omega_{R_+,R'}\to G\) factors through
\[
  i:\Omega_{R_+,R'}
  \longrightarrow
  \underline{\operatorname{Norm}}_G(\mathfrak g_{R'}).
\]
Suppose that \(G\) is adjoint. The tangent map of \(i\) at the
points of the unit section is bijective by \SGARefLink{sga3:xxii:lemma:5.3.2}{5.3.2}. In particular,
\(i\) is étale at every point of the unit section, and hence is a
local immersion\footnote{Editor's note: By EGA
\(\mathrm{IV}_4\), 17.11.2, \(i\) is étale at every point of the
unit section \textup{(and
\(\underline{\operatorname{Norm}}_G(\mathfrak g_{R'})\) is smooth
at every point of the unit section)}. Moreover, let \(V\) be the
largest open subset of \(\Omega_{R_+,R'}\) on which \(i\) is
étale. Since \(i\) is a monomorphism, \(i_V\) is an open immersion
\textup{(ibid., 17.9.1)}.} at every point of the unit section.
Therefore
\(\underline{\operatorname{Norm}}_G(\mathfrak g_{R'})\) is smooth
at every point of the unit section, as was to be proved.

\SGARefTarget{sga3:xxii:section:5.5}
\subsection*{\SGAInlineRefTarget{sga3:xxii:subsection:5-5}5.5. Borel Subgroups of a Split Reductive Group}
\addcontentsline{toc}{subsection}{\SGARefLink{sga3:xxii:subsection:5-5}{5.5}. Borel Subgroups of a Split
Reductive Group}

\SGARefTarget{sga3:xxii:prop:5.5.1}\paragraph{\SGARefLink{sga3:xxii:prop:5.5.1}{Proposition 5.5.1}.}
\textit{Let \((G,T,M,R)\) be a split \(S\)-group. For every system
of positive roots \(R_+\) of \(R\), the subgroup \(H_{R_+}\)
\textup{(which exists by \SGARefLink{sga3:xxii:thm:5.4.7}{5.4.7})} is a Borel subgroup of \(G\).
Moreover, for every ordering on \(R_+\), the morphism induced by
multiplication in \(G\),}
\[
  T\mathop{\times}\limits_S
  \prod_{\alpha\in R_+}U_\alpha
  \longrightarrow G,
\]
\textit{is a closed immersion with image \(H_{R_+}\). We write
\(B_{R_+}=H_{R_+}\).}

By the definition of Borel subgroups, the first assertion may be
verified after replacing \(S\) by the spectrum of an algebraically
closed field. Let \(B\) be the Borel subgroup of \(G\) containing
\(T\) and corresponding to the system of positive roots \(R_+\)
\textup{(\emph{Bible}, \S10.4, Proposition 9)}. The Lie algebra of
\(B\) is \(\mathfrak g_{R_+}\); therefore \(B=H_{R_+}\) by \SGARefLink{sga3:xxii:cor:5.3.5}{5.3.5}.

For the second assertion, the morphism in the statement induces an
open immersion
\[
  i:
  T\mathop{\times}\limits_S
  \prod_{\alpha\in R_+}U_\alpha
  \longrightarrow H_{R_+}
\]
by \SGARefLink{sga3:xxii:lem:5.4.4}{5.4.4}. But \(i\) is surjective
\textup{(\emph{Bible}, \S15.1, Corollary 1 to Proposition 1)}.

\SGARefTarget{sga3:xxii:cor:5.5.2}\paragraph{\SGARefLink{sga3:xxii:cor:5.5.2}{Corollary 5.5.2}.}
\textit{Choose an arbitrary ordering on \(R_+\), and for each
\(\alpha\in R_+\) choose an}
\[
  X_\alpha\in\Gamma(S,\mathfrak g^\alpha)^\times.
\]
\textit{Let \(\alpha,\beta\in R_+\). For every pair
\((i,j)\in\mathbb N^*\times\mathbb N^*\) such that
\(i\alpha+j\beta\in R\), there is a unique section}
\[
  C_{i,j,\alpha,\beta}\in\Gamma(S,\mathcal O_S)
\]
\textit{such that, for all \(x,y\in\mathbb G_a(S')\), where
\(S'\to S\), one has}
\[
  \exp_\alpha(xX_\alpha)
  \exp_\beta(yX_\beta)
  \exp_\alpha(xX_\alpha)^{-1}
  =
  \exp_\beta(yX_\beta)
  \prod_{\substack{i,j\in\mathbb N^*\\ i\alpha+j\beta\in R}}
    \exp_{i\alpha+j\beta}
      \bigl(C_{i,j,\alpha,\beta}x^iy^jX_{i\alpha+j\beta}\bigr).
\]

If \(\alpha=\beta\), the assertion is trivial. Suppose therefore
that \(\alpha\ne\beta\). By the proposition there are unique
morphisms
\[
  F_0:\mathbb G_{a,S}^{\,2}\longrightarrow T,
  \qquad
  F_\gamma:\mathbb G_{a,S}^{\,2}\longrightarrow\mathbb G_{a,S}
  \quad(\gamma\in R_+)
\]
such that
\[
  \exp(xX_\alpha)\exp(yX_\beta)\exp(xX_\alpha)^{-1}
  =
  F_0(x,y)
  \prod_{\gamma\in R_+}
    \exp\bigl(F_\gamma(x,y)X_\gamma\bigr).
\]
Let \(t\in T(S')\), where \(S'\to S\). Applying
\(\operatorname{int}(t)\) to this formula, we immediately obtain
\[
  \tag{1}\label{sga3:xxii:subsection:5-5:equation:1}
  F_0\bigl(\alpha(t)x,\beta(t)y\bigr)=F_0(x,y),
\]
\[
  \tag{2}\label{sga3:xxii:subsection:5-5:equation:2}
  F_\gamma\bigl(\alpha(t)x,\beta(t)y\bigr)
  =
  \gamma(t)F_\gamma(x,y).
\]
Since \(\alpha\) and \(\beta\) are two linearly independent
characters of \(T\) over \(\mathbb Q\), the first relation implies,
as usual, that \(F_0\) is constant, and hence
\(F_0(x,y)=e\). Write next
\[
  F_\gamma(x,y)=\sum a_{ij}x^iy^j,
  \qquad a_{ij}\in\Gamma(S,\mathcal O_S).
\]
Substituting in relation \textup{\SGARefLink{sga3:xxii:subsection:5-5:equation:2}{(2)}} and comparing the
polynomials on both sides gives
\[
  a_{ij}
  \bigl(\alpha(t)^i\beta(t)^j-\gamma(t)\bigr)=0.
\]
If \(\gamma\ne i\alpha+j\beta\), we know
\textup{(\SGARefLink{sga3:xix:lemma:4:13}{Exposé XIX, 4.13})} that there exist a faithfully flat
quasi-compact morphism \(S'\to S\) and an element \(t\in T(S')\)
such that
\[
  \alpha(t)^i\beta(t)^j-\gamma(t)=1.
\]
Thus \(a_{ij}=0\) over \(S'\), and hence over \(S\). If
\(\gamma=i\alpha+j\beta\), set
\(a_{ij}=C_{i,j,\alpha,\beta}\). Setting \(x=0\)
\textup{(respectively, \(y=0\))}, one obtains
\[
  C_{0,1,\alpha,\beta}=1
  \qquad
  \textup{(respectively, \(C_{1,0,\alpha,\beta}=0\)).}
\]

\SGARefTarget{sga3:xxii:rem:5.5.3}\paragraph{\SGARefLink{sga3:xxii:rem:5.5.3}{Remark 5.5.3}.}
Differentiating with respect to \(y\) at \(y=0\) and comparing
with \SGARefLink{sga3:xxii:lem:5.4.9}{5.4.9}, one finds
\[
  C_{i,1,\alpha,\beta}=M_{\alpha,\beta,i}.
\]

\SGARefTarget{sga3:xxii:cor:5.5.4}\paragraph{\SGARefLink{sga3:xxii:cor:5.5.4}{Corollary 5.5.4}.}
\textit{Let \(S\) be a scheme, let \(G\) be a reductive
\(S\)-group, let \(T\) be a maximal torus of \(G\), and let
\(\alpha\ne\beta\) be two roots of \(G\) relative to \(T\) such
that \(\alpha+\beta\) is nontrivial on every fiber. Order
arbitrarily the set of roots \(i\alpha+j\beta\), where
\((i,j)\in\mathbb N^*\times\mathbb N^*\). For all
\(i,j\in\mathbb N^*\) such that \(i\alpha+j\beta\in R\), there is
a unique morphism of \(\mathcal O_S\)-modules}
\[
  f_{\alpha,\beta,i,j}:
  (\mathfrak g^\alpha)^{\otimes i}
  \otimes
  (\mathfrak g^\beta)^{\otimes j}
  \longrightarrow
  \mathfrak g^{i\alpha+j\beta}.
\]
\textit{such that, for every \(S'\to S\) and all
\(X\in\mathbf W(\mathfrak g^\alpha)(S')\),
\(Y\in\mathbf W(\mathfrak g^\beta)(S')\), one has
\textup{(the exponentials in the right-hand side being understood
in the sense of \SGARefLink{sga3:xxii:rem:1.2}{1.2})}\footnote{Editor's note: That is, if
\(\mathfrak g^{i\alpha+j\beta}=0\) on a connected component of
\(S\), the corresponding exponential is \(1\) on that component.}}
\[
  \exp_\alpha(X)\exp_\beta(Y)\exp_\alpha(-X)
  =
  \exp_\beta(Y)
  \prod_{(i,j)}
    \exp_{i\alpha+j\beta}
      \bigl(f_{\alpha,\beta,i,j}(X^{\otimes i}\otimes
      Y^{\otimes j})\bigr).
\]

The assertion is local for the fpqc topology. By \S2 we may
therefore suppose that \(G\) is split relative to \(T\), with
\(\alpha\) and \(\beta\) constant in the splitting. Since
\(\alpha+\beta\ne0\), there is a system of positive roots \(R_+\)
containing \(\alpha\) and \(\beta\) \textup{(\SGARefLink{sga3:XXI:3.5.4}{Exposé XXI, 3.5.4})},
and we are reduced to \SGARefLink{sga3:xxii:cor:5.5.2}{5.5.2}.

\SGARefTarget{sga3:xxii:cor:5.5.5}\paragraph{\SGARefLink{sga3:xxii:cor:5.5.5}{Corollary 5.5.5}.}
\textit{Let \(S\) be a scheme and \(G\) a reductive
\(S\)-group.}
\begin{enumerate}
\item[\SGAInlineRefTarget{sga3:xxii:cor:5.5.5:item:i}(i)] \textit{\(G\) has Borel subgroups locally for the
étale topology. If \(T\) is a maximal torus of \(G\), then,
locally for the étale topology, \(G\) also has Borel subgroups
containing \(T\).}
\item[\SGAInlineRefTarget{sga3:xxii:corollary:5-5-5:item:ii}\SGAInlineRefTarget{sga3:xxii:cor:5.5.5:item:ii}(ii)] \textit{If \(T\) is a maximal torus of \(G\), the
``functor of Borel subgroups of \(G\) containing \(T\)'' is
represented by a principal homogeneous bundle under \(W_G(T)\).}
\item[\SGAInlineRefTarget{sga3:xxii:cor:5.5.5:item:iii}(iii)] \textit{If \((G,T,M,R)\) is split, every Borel
subgroup \(B\) of \(G\) containing \(T\) is locally on \(S\) of
the form \(B_{R_+}\), where \(R_+\) is a system of positive roots
of \(R\).}
\item[\SGAInlineRefTarget{sga3:xxii:cor:5.5.5:item:iv}(iv)] \textit{If \(T\subset B\) is a Killing pair of \(G\),
there is an étale covering family \(\{S_i\to S\}\), and for every
\(i\) a splitting \((G_{S_i},T_{S_i},M_i,R_i)\) and a system of
positive roots \(R_{i+}\) of \(R_i\), such that
\(B_{S_i}=B_{R_{i+}}\).}
\end{enumerate}

Indeed, \textup{\SGARefLink{sga3:xxii:cor:5.5.5:item:i}{(i)}} follows from \SGARefLink{sga3:xxii:cor:2.3}{2.3} and \SGARefLink{sga3:xxii:prop:5.5.1}{5.5.1},
\textup{\SGARefLink{sga3:xxii:cor:5.5.5:item:ii}{(ii)}} from \textup{\SGARefLink{sga3:xxii:cor:5.5.5:item:i}{(i)}} and \SGARefLink{sga3:xxii:cor:5.3.15}{5.3.15},
\textup{\SGARefLink{sga3:xxii:cor:5.5.5:item:iii}{(iii)}} from \textup{\SGARefLink{sga3:xxii:cor:5.5.5:item:ii}{(ii)}} and \SGARefLink{sga3:xxii:prop:5.5.1}{5.5.1}, and
\textup{\SGARefLink{sga3:xxii:cor:5.5.5:item:iv}{(iv)}} from \textup{\SGARefLink{sga3:xxii:cor:5.5.5:item:iii}{(iii)}} and \SGARefLink{sga3:xxii:cor:2.3}{2.3}.

\SGARefTarget{sga3:xxii:lem:5.5.6}\paragraph{\SGARefLink{sga3:xxii:lem:5.5.6}{Lemma 5.5.6}.}
\textit{Choose on the group \(\Gamma_0(R)\) generated by the
roots a totally ordered group structure such that the roots
\(>0\) are the elements of \(R_+\)
\textup{(cf. \SGARefLink{sga3:XXI:3.5.6}{Exposé XXI, 3.5.6})}\footnote{Editor's note: Such an
ordering is necessarily compatible with
\(\operatorname{ord}_\Delta\), where
\(\Delta=\mathscr S(R_+)\) \textup{(cf. \SGARefLink{sga3:XXI:3.2.15}{Exposé XXI, 3.2.15})}.}
Let \(\alpha_1<\cdots<\alpha_N\) be the elements of \(R_+\).
Consider the isomorphism}
\[
  f:
  T\mathop{\times}\limits_S U_{\alpha_1}
  \mathop{\times}\limits_S\cdots
  \mathop{\times}\limits_S U_{\alpha_N}
  \longrightarrow B_{R_+}
\]
\textit{induced by multiplication in \(G\). For
\(i=1,\ldots,N\), put}
\[
  U_{\geq i}
  =
  f\left(
    U_{\alpha_i}\mathop{\times}\limits_S\cdots
    \mathop{\times}\limits_S U_{\alpha_N}
  \right).
\]
\begin{enumerate}
\item[\SGAInlineRefTarget{sga3:xxii:lem:5.5.6:item:i}(i)] \textit{Each \(U_{\geq i}\) is an invariant subgroup
of \(B_{R_+}\).}
\item[\SGAInlineRefTarget{sga3:xxii:lem:5.5.6:item:ii}(ii)] \textit{For \(1\leq i\leq N-1\),
\(U_{\geq i}\) identifies with the semidirect product}
\[
  U_{\geq i}=U_{\alpha_i}\cdot U_{\geq i+1}.
\]
\item[\SGAInlineRefTarget{sga3:xxii:lem:5.5.6:item:iii}(iii)] \textit{\(B_{R_+}\) identifies with the semidirect
product}
\[
  B_{R_+}=T\cdot U_{\geq1}.
\]
\item[\SGAInlineRefTarget{sga3:xxii:lem:5.5.6:item:iv}(iv)] \textit{For \(1\leq i\leq N-1\), the inner
automorphisms of \(U_{\geq1}\) act trivially on
\(U_{\geq i}/U_{\geq i+1}\), identified with \(U_{\alpha_i}\)
by \textup{\SGARefLink{sga3:xxii:lem:5.5.6:item:ii}{(ii)}}.}
\end{enumerate}

We first prove, by descending induction on \(i\), the following
assertion:
\[
  U_{\geq i}\ \text{is an invariant subgroup of }B_{R_+},
  \quad\text{the semidirect product of }
  U_{\alpha_i}\text{ and }U_{\geq i+1}.
\]
This is true for \(i=N\); suppose it has been proved for \(i+1\),
and let us prove it for \(i\). As schemes, we have
\[
  U_{\geq i}=U_{\alpha_i}\cdot U_{\geq i+1}.
\]
We must show that \(U_{\geq i}\) is invariant under the inner
automorphisms of \(B_{R_+}\). This is clear for
\(\operatorname{int}(t)\), where \(t\in T(S)\); it is therefore
enough to consider \(\operatorname{int}(x)\), where
\(x\in U_\alpha(S)\) and \(\alpha\in R_+\). Since
\(U_{\geq i+1}\) is invariant, it is enough to prove that
\[
  \operatorname{int}(x)(U_{\alpha_i})\subset U_{\geq i}.
\]
By \SGARefLink{sga3:xxii:cor:5.5.2}{5.5.2}, if \(y\in U_{\alpha_i}(S')\), then
\[
  y^{-1}xyx^{-1}\in U_{\geq i+1}(S'),
\]
which implies that
\[
  \operatorname{int}(x)(y)\in U_{\geq i}(S').
\]

It remains to prove that \(U_{\geq i}\) is a subgroup. Let
\(x,y\in U_{\geq i}(S)\), and write
\[
  x=px',\qquad y=qy',
\]
where \(p,q\in U_{\alpha_i}(S)\) and
\(x',y'\in U_{\geq i+1}(S)\). We then have
\[
  xy=px'qy'=pq(q^{-1}x'q)y'
  \in U_{\alpha_i}(S')U_{\geq i+1}(S'),
\]
and likewise
\[
  x^{-1}
  =
  p^{-1}\bigl(p{x'}^{-1}p^{-1}\bigr)
  \in U_{\alpha_i}(S')U_{\geq i+1}(S').
\]
We have thus proved \textup{\SGARefLink{sga3:xxii:lem:5.5.6:item:i}{(i)}} and \textup{\SGARefLink{sga3:xxii:lem:5.5.6:item:ii}{(ii)}}, as well as
\textup{\SGARefLink{sga3:xxii:lem:5.5.6:item:iv}{(iv)}} along the way. Assertion \textup{\SGARefLink{sga3:xxii:lem:5.5.6:item:iii}{(iii)}} is a trivial
consequence of \SGARefLink{sga3:xxii:prop:5.5.1}{5.5.1}.

\SGARefTarget{sga3:xxii:lem:5.5.7}\paragraph{\SGARefLink{sga3:xxii:lem:5.5.7}{Lemma 5.5.7}.}
\textit{With the preceding notation, choose, for each
\(1\leq i\leq N\), an element}
\[
  X_i\in\Gamma(S,\mathfrak g^{\alpha_i})^\times,
\]
\textit{and consider the isomorphism}
\[
  a:\mathbb G_{a,S}^{\,N}\longrightarrow U_{\geq1}
\]
\textit{defined on underlying sets by}
\[
  a(x_1,\ldots,x_N)
  =
  \exp_{\alpha_1}(x_1X_1)\cdots
  \exp_{\alpha_N}(x_NX_N).
\]
\textit{There is a unique family \((Q_i)_{1\leq i\leq N}\) of
polynomials}
\[
  Q_i
  =
  Q_i(x_1,\ldots,x_N,y_1,\ldots,y_N)
\]
\textit{with coefficients in \(\Gamma(S,\mathcal O_S)\), such
that}
\[
\begin{split}
  &a(x_1,\ldots,x_N)a(y_1,\ldots,y_N)\\
  &\qquad =
  a\bigl(
    Q_1(x_1,\ldots,x_N,y_1,\ldots,y_N),
    \ldots,
    Q_N(x_1,\ldots,x_N,y_1,\ldots,y_N)
  \bigr).
\end{split}
\]
\textit{Moreover, the coefficients of these polynomials belong
to the subring of \(\Gamma(S,\mathcal O_S)\) generated by the
\(C_{i,j,\alpha,\beta}\) of \SGARefLink{sga3:xxii:cor:5.5.2}{5.5.2}
\(\bigl(\alpha,\beta\in R_+,\ i,j\in\mathbb N^*\bigr)\), and
each \(Q_i\) is of the form}
\[
\begin{split}
  Q_i(x_1,\ldots,x_N,y_1,\ldots,y_N)
  ={}&x_i+y_i\\
  &+
  Q_i'(x_1,\ldots,x_{i-1},y_1,\ldots,y_{i-1}).
\end{split}
\]

The existence and uniqueness of the \(Q_i\) are immediate from
the fact that \(a\) is an isomorphism of schemes. Let
\(z,z',z''\) denote sections of \(U_{\geq i+1}\). We have
\[
\begin{split}
  &a(x_1,\ldots,x_N)a(y_1,\ldots,y_N)\\
  &\quad =
  a(x_1,\ldots,x_{i-1},0,\ldots,0)
  \exp_{\alpha_i}(x_iX_i)z\\
  &\qquad\quad{}\cdot
  a(y_1,\ldots,y_{i-1},0,\ldots,0)
  \exp_{\alpha_i}(y_iX_i)z'.
\end{split}
\]
Using \SGARefLink{sga3:xxii:lem:5.5.6}{5.5.6} \textup{\SGARefLink{sga3:xxii:lem:5.5.6:item:i}{(i)}} and \textup{\SGARefLink{sga3:xxii:lem:5.5.6:item:iv}{(iv)}}, the right-hand side
can be written
\[
\begin{split}
  &a(x_1,\ldots,x_{i-1},0,\ldots,0)
  a(y_1,\ldots,y_{i-1},0,\ldots,0)\\
  &\qquad{}\cdot
  \exp_{\alpha_i}\bigl((x_i+y_i)X_i\bigr)z'',
\end{split}
\]
and, applying \SGARefLink{sga3:xxii:lem:5.5.6}{5.5.6} once again, this gives
\[
  Q_i(x_1,\ldots,x_N,y_1,\ldots,y_N)
  =
  x_i+y_i+
  Q_i'(x_1,\ldots,x_{i-1},y_1,\ldots,y_{i-1}),
\]
where
\[
  Q_i'(x_1,\ldots,x_{i-1},y_1,\ldots,y_{i-1})
  =
  Q_i(x_1,\ldots,x_{i-1},0,\ldots,0,
      y_1,\ldots,y_{i-1},0,\ldots,0);
\]
this is precisely the required form.\footnote{Editor's note: The
French source prints \(x_i\) as the first argument on the right-hand
side of this formula. We have written \(x_1\), as required by the
defining specialization and by the preceding calculation.}

It remains to prove the assertion concerning the coefficients of the
polynomials \(Q_i\). Let \(A\) be the subring of
\(\Gamma(S,\mathcal O_S)\) generated by the
\(C_{i,j,\alpha,\beta}\), where
\(\alpha,\beta\in R_+\) and \(i,j\in\mathbb N^*\).
We prove by descending induction on \(i\) that, if
\[
  x_1=\cdots=x_{i-1}=0
  \qquad\text{and}\qquad
  y_1=\cdots=y_{i-1}=0,
\]
that is, if \(a(x_1,\ldots,x_N)\) and \(a(y_1,\ldots,y_N)\)
are sections of \(U_{\geq i}\), then the polynomials
\[
  R_j(x_i,\ldots,x_N,y_i,\ldots,y_N)
  =
  Q_j(x_1,\ldots,x_N,y_1,\ldots,y_N)
\]
have coefficients in \(A\). This is trivial for \(i=N\), and
also for \(j<i\), since \(R_j=0\) when \(j<i\). Let \(i<N\);
suppose the assertion proved for \(i+1\), and let us prove it
for \(i\) and \(j\geq i\). We have
\[
\begin{split}
  a(0,\ldots,0,x_i,\ldots,x_N)
  &=
  \exp(x_iX_i)
  a(0,\ldots,0,x_{i+1},\ldots,x_N)\\
  &=
  \exp(x_iX_i)Z_i.
\end{split}
\]
Similarly, write
\[
  a(0,\ldots,y_i,\ldots,y_N)=\exp(y_iX_i)T_i.
\]
Then
\[
\begin{split}
  &a(0,\ldots,x_i,\ldots,x_N)
   a(0,\ldots,y_i,\ldots,y_N)\\
  &\qquad =
  \exp\bigl((x_i+y_i)X_i\bigr)
  \operatorname{int}\bigl(\exp(-y_iX_i)\bigr)(Z_i)\cdot T_i.
\end{split}
\]
Now\footnote{Editor's note: In the second conjugating exponential
the French source prints the undefined symbol \(Y_i\). We have
written \(X_i\), as required by the preceding identity and by the
definition of \(Z_i\).}
\[
\begin{split}
  \operatorname{int}\bigl(\exp(-y_iX_i)\bigr)(Z_i)
  &=
  \operatorname{int}\bigl(\exp(-y_iX_i)\bigr)
  \bigl(
    \exp(x_{i+1}X_{i+1})\cdots\exp(x_NX_N)
  \bigr)
\end{split}
\]
is a product of \(N-i-1\) sections of \(U_{\geq i+1}\), whose
coefficients in the decomposition
\[
  U_{\geq i+1}
  =
  U_{\alpha_{i+1}}\cdots U_{\alpha_N}
\]
are polynomials in \(y_i,x_{i+1},\ldots,x_N\) with coefficients
in \(A\), by \SGARefLink{sga3:xxii:cor:5.5.2}{5.5.2}. Applying the induction hypothesis, we deduce
that the coefficients of
\[
  \operatorname{int}\bigl(\exp(-y_iX_i)\bigr)(Z_i)\cdot T_i
\]
are likewise polynomials with coefficients in \(A\), which
completes the proof.

The preceding induction immediately gives a proof of the following.

\SGARefTarget{sga3:xxii:lem:5.5.8}\paragraph{\SGARefLink{sga3:xxii:lem:5.5.8}{Lemma 5.5.8}.}
\textit{With the notation of \SGARefLink{sga3:xxii:lem:5.5.6}{5.5.6}, let, for each
\(i=1,\ldots,N\),}
\[
  f_i:U_{\alpha_i}\longrightarrow H
\]
\textit{be a group morphism, where \(H\) is an \(S\)-group
functor. For the morphism}
\[
  f:U_{\geq1}\longrightarrow H
\]
\textit{defined by}
\[
  f\bigl(
    \exp(x_1X_1)\cdots\exp(x_NX_N)
  \bigr)
  =
  f_1\bigl(\exp(x_1X_1)\bigr)\cdots
  f_N\bigl(\exp(x_NX_N)\bigr)
\]
\textit{to be a group morphism, it is necessary and sufficient
that, for every pair \(i<j\),}
\[
\begin{split}
  &f_j\bigl(\exp(x_jX_j)\bigr)
   f_i\bigl(\exp(x_iX_i)\bigr)
   f_j\bigl(\exp(-x_jX_j)\bigr)\\
  &\qquad =
  f\bigl(
    \exp(x_jX_j)\exp(x_iX_i)\exp(-x_jX_j)
  \bigr).
\end{split}
\]

\SGARefTarget{sga3:xxii:sec:5.6}
\subsection*{\SGAInlineRefTarget{sga3:xxii:subsection:5-6}5.6. Subgroups of Type \textup{(R)} with Solvable
Fibers}

\SGARefTarget{sga3:xxii:prop:5.6.1}
\paragraph{\SGARefLink{sga3:xxii:prop:5.6.1}{Proposition 5.6.1}.}
\textit{Let \((G,T,M,R)\) be a split \(S\)-group, let \(R'\) be
a subset of \(R\) of type \textup{(R)} \textup{(\SGARefLink{sga3:xxii:def:5.4.2}{5.4.2})}, and let
\(H_{R'}\) be the corresponding subgroup of \(G\). The following
conditions are equivalent:}
\begin{enumerate}
\item[\textup{\SGAInlineRefTarget{sga3:xxii:prop:5.6.1:item:i}(i)}]
\textit{\(H_{R'}\) has solvable geometric fibers.}

\item[\textup{\SGAInlineRefTarget{sga3:xxii:prop:5.6.1:item:ii}(ii)}]
\textit{There exists a system of positive roots \(R_+\) such that
\(R'\subset R_+\), and hence \(H_{R'}\subset B_{R_+}\)
\textup{(cf.\ \SGARefLink{sga3:xxii:prop:5.4.5}{5.4.5})}.}

\item[\textup{\SGAInlineRefTarget{sga3:xxii:prop:5.6.1:item:iii}(iii)}]
\textit{\(R'\cap(-R')=\varnothing\).}

\item[\textup{\SGAInlineRefTarget{sga3:xxii:prop:5.6.1:item:iv}(iv)}]
\textit{For every ordering of \(R'\), the morphism induced by
multiplication in \(G\)}
\[
  T\times_S\prod_{\alpha\in R'}U_\alpha
  \longrightarrow H_{R'}
\]
\textit{is an isomorphism.}

\item[\textup{\SGAInlineRefTarget{sga3:xxii:prop:5.6.1:item:v}(v)}]
\textit{\(
  H_{R'}\cap\underline{\operatorname{Norm}}_G(T)=T.
\)}

\item[\textup{\SGAInlineRefTarget{sga3:xxii:prop:5.6.1:item:vi}(vi)}]
\textit{For every subset \(R''\) of \(R\) of type \textup{(R)},
one has \textup{(cf.\ \SGARefLink{sga3:xxii:prop:5.4.5}{5.4.5})}}
\[
  H_{R'}\cap
  \underline{\operatorname{Norm}}_G(H_{R''})
  =
  H_{R'\cap R''}.
\]
\end{enumerate}

We shall prove these equivalences according to the following logical
scheme.

\SGARefTarget{sga3:xxii:diagram:5.6.1:implication-scheme}
\begin{center}
  % BEGIN INCLUDED SGA3_English_current_buildable_source_20260823/source/inputs/XXII/native_diagrams/exp22/diag_006_prop561_implication_scheme.tex
\begin{tikzcd}[column sep=5.2em,row sep=2.9em]
  \textup{(iii)}
    \arrow[dd,Rightarrow]
  &
  \textup{(ii)}
    \arrow[l,Rightarrow]
    \arrow[r,Rightarrow]
    \arrow[d,Rightarrow]
  &
  \textup{(vi)}
    \arrow[dd,Rightarrow]
  \\
  &
  \textup{(iv)}
    \arrow[dr,Rightarrow]
  &
  \\
  \textup{(i)}
    \arrow[ur,Rightarrow]
  &&
  \textup{(v)}.
    \arrow[ll,Rightarrow]
\end{tikzcd}

% END INCLUDED SGA3_English_current_buildable_source_20260823/source/inputs/XXII/native_diagrams/exp22/diag_006_prop561_implication_scheme.tex

\end{center}

Clearly \textup{\SGARefLink{sga3:xxii:prop:5.6.1:item:ii}{(ii)}} implies \textup{\SGARefLink{sga3:xxii:prop:5.6.1:item:iii}{(iii)}}, and \textup{\SGARefLink{sga3:xxii:prop:5.6.1:item:vi}{(vi)}}
implies \textup{\SGARefLink{sga3:xxii:prop:5.6.1:item:v}{(v)}} \textup{(take \(R''=\varnothing\))}. By \SGARefLink{sga3:xxii:cor:5.4.6}{5.4.6},
it is enough to verify that \textup{\SGARefLink{sga3:xxii:prop:5.6.1:item:i}{(i)}} implies \textup{\SGARefLink{sga3:xxii:prop:5.6.1:item:ii}{(ii)}} on
the geometric fibers. If \(S\) is the spectrum of an algebraically
closed field, then \(H_{R'}\) is contained in a Borel subgroup
containing \(T\), and this subgroup is therefore of the form
\(H_{R_+}\) \textup{(\SGARefLink{sga3:xxii:cor:5.5.5}{5.5.5}\textup{\SGARefLink{sga3:xxii:cor:5.5.5:item:iii}{(iii)}})}.

Likewise, the implication from \textup{\SGARefLink{sga3:xxii:prop:5.6.1:item:iii}{(iii)}} to \textup{\SGARefLink{sga3:xxii:prop:5.6.1:item:i}{(i)}} may be
checked on the geometric fibers. Suppose that \textup{\SGARefLink{sga3:xxii:prop:5.6.1:item:iii}{(iii)}} holds.
If \(H_{R'}\) were not solvable, there would exist a subtorus \(Q\)
of \(T\), of codimension \(1\) in \(T\), such that
\[
  \underline{\operatorname{Centr}}_{H_{R'}}(Q)
\]
is not solvable \textup{(\emph{Bible}, \S10.4, Proposition 8)}.
But \(\underline{\operatorname{Centr}}_{H_{R'}}(Q)\) has Lie algebra
\(\mathfrak g_{R''}\), where \(R''\) is the set of roots in \(R'\)
which vanish on \(Q\). Thus, by \textup{\SGARefLink{sga3:xxii:prop:5.6.1:item:iii}{(iii)}},
\[
  R''=\varnothing\qquad\text{or}\qquad R''=\{\alpha\}.
\]
Consequently,
\(\underline{\operatorname{Centr}}_{H_{R'}}(Q)\), which is a subgroup
of \(G\) of type \textup{(R)}, is \(T\) or \(T\cdot U_\alpha\), and
is therefore solvable, contrary to the hypothesis.

Likewise, the implication from \textup{\SGARefLink{sga3:xxii:prop:5.6.1:item:ii}{(ii)}} to \textup{\SGARefLink{sga3:xxii:prop:5.6.1:item:iv}{(iv)}} may be
checked on the geometric fibers \textup{(since the schemes in
question are flat \(S\)-schemes of finite presentation)}. By
\emph{Bible}, \S13.2, Theorem 1\textup{\SGARefLink{sga3:xxii:rem:5.1.2:item:d}{(d)}}, the morphism in question
is bijective. It induces an isomorphism on the tangent spaces at the
origin, and the conclusion follows as usual \textup{(cf.\ \SGARefLink{sga3:xxii:par:4.1.1}{4.1.1})}.

The implication from \textup{\SGARefLink{sga3:xxii:prop:5.6.1:item:iv}{(iv)}} to \textup{\SGARefLink{sga3:xxii:prop:5.6.1:item:v}{(v)}} follows from
\SGARefLink{sga3:xxii:lem:4.2.7}{4.2.7}. To prove that \textup{\SGARefLink{sga3:xxii:prop:5.6.1:item:v}{(v)}} implies \textup{\SGARefLink{sga3:xxii:prop:5.6.1:item:i}{(i)}}, we may again
reduce to the case in which \(S\) is the spectrum of an algebraically
closed field; the conclusion then follows from \emph{Bible}, \S10.3,
the corollary to \SGARefLink{sga3:xxii:section:6}{Proposition 6}, and \S9.3, Corollary 3 to Theorem 1.

It remains only to prove the implication from \textup{\SGARefLink{sga3:xxii:prop:5.6.1:item:ii}{(ii)}} to
\textup{\SGARefLink{sga3:xxii:prop:5.6.1:item:vi}{(vi)}}. We may reduce to the case in which \(G\) is adjoint.
Then, by \SGARefLink{sga3:xxii:cor:5.3.3}{5.3.3},
\[
  \underline{\operatorname{Norm}}_G(H_{R''})
  =
  \underline{\operatorname{Norm}}_G(\mathfrak g_{R''})
  \subset
  \underline{\operatorname{Transp}}_G
    (\mathfrak t,\mathfrak g_{R''}).
  \tag{39}\label{sga3:xxii:subsection:5-6:equation:39}
\]
\footnote{Editor's note: The inclusion follows from
\(\mathfrak t\subset\mathfrak g_{R''}\).}
By \SGARefLink{sga3:xxii:prop:5.4.5}{5.4.5}, it is enough to prove
\[
  H_{R'}(S)\cap
  \underline{\operatorname{Transp}}_{G(S)}
    (\mathfrak t,\mathfrak g_{R''})
  \subset
  H_{R'\cap R''}(S).
  \tag{x}\label{sga3:xxii:subsection:5-6:equation:x}
\]
We first prove a lemma.

\SGARefTarget{sga3:xxii:lem:5.6.2}
\paragraph{\SGARefLink{sga3:xxii:lem:5.6.2}{Lemma 5.6.2}.}
\textit{With the notation of \SGARefLink{sga3:xxii:lem:5.5.7}{5.5.7}, let}
\[
  u=\exp(x_1X_1)\cdots\exp(x_NX_N),
\]
\textit{where \(x_i\in\mathbb G_a(S)\). Let \(m\), with
\(1\leq m\leq N\), be such that \(x_i=0\) for \(i<m\).}
\begin{enumerate}
\item[\textup{\SGAInlineRefTarget{sga3:xxii:lem:5.6.2:item:a}(a)}]
\textit{If \(H\in\Gamma(S,\mathfrak t)\), the component of
\(\operatorname{Ad}(u)H\) in \(\mathfrak g^{\alpha_m}\) is}
\[
  -\overline{\alpha_m}(H)x_mX_m.
\]

\item[\textup{(b)}]
\textit{If \(Y\in\Gamma(S,\mathfrak g^{-\alpha_m})\), the component
of \(\operatorname{Ad}(u)Y\) in \(\mathfrak t\) is, with the notation
of \SGARefLink{sga3:xx:2.6}{Exposé XX, 2.6},}
\[
  x_m\langle X_m,Y\rangle H_{\alpha_m}.
\]
\end{enumerate}

\textit{Proof.}
Set
\[
  \mathfrak u
  =
  \mathfrak g^{\alpha_{m+1}}+\cdots+\mathfrak g^{\alpha_N}.
\]
By \SGARefLink{sga3:xxii:lem:5.4.9}{5.4.9}, one has
\[
  \operatorname{Ad}\bigl(\exp(x_iX_i)\bigr)\mathfrak u
  \subset\mathfrak u
  \qquad\text{for }i>m.
\]
By \SGARefLink{sga3:xx:2.10}{Exposé XX, 2.10},
\[
  \operatorname{Ad}\bigl(\exp(x_iX_i)\bigr)H
  =
  H-\overline{\alpha_i}(H)x_iX_i.
\]
It follows immediately that, modulo \(\mathfrak u\),
\[
  \operatorname{Ad}(u)H
  =
  \operatorname{Ad}\bigl(\exp(x_mX_m)\bigr)H
  =
  H-\overline{\alpha_m}(H)x_mX_m,
\]
which proves the first assertion.

Likewise, set\footnote{Editor's note: The original has been corrected
in what follows.}
\[
  \mathfrak n
  =
  \mathfrak g^{\alpha_1}+\cdots+\mathfrak g^{\alpha_N},
  \qquad
  u_1
  =
  \exp(x_{m+1}X_{m+1})\cdots\exp(x_NX_N).
\]
For \(i>m\), one has \(\alpha_i>\alpha_m\); hence, by \SGARefLink{sga3:xxii:lem:5.4.9}{5.4.9}, modulo
\(\mathfrak n\),
\[
  \operatorname{Ad}(u_1)Y\equiv Y,
  \qquad\text{and therefore}\qquad
  \operatorname{Ad}(u)Y
  \equiv
  \operatorname{Ad}\bigl(\exp(x_mX_m)\bigr)Y.
\]
Applying \SGARefLink{sga3:xx:2.10}{Exposé XX, 2.10}, one therefore obtains, modulo
\(\mathfrak n\),
\[
  \operatorname{Ad}(u)Y-Y
  \equiv
  x_m\langle X_m,Y\rangle H_{\alpha_m},
\]
which proves the second assertion.

Let us return to the proof of inclusion \textup{\SGARefLink{sga3:xxii:subsection:5-6:equation:x}{(x)}}. Suppose that
there exists
\[
  h\in H_{R'}(S),
  \qquad
  h\notin H_{R'\cap R''}(S),
\]
such that
\[
  \operatorname{Ad}(h)\mathfrak t\subset\mathfrak g_{R''}.
\]
We may write
\[
  h=t\exp(x_1X_1)\cdots\exp(x_NX_N).
\]
Since \(h\notin H_{R'\cap R''}(S)\), there is a least \(m\) such
that
\[
  t\exp(x_1X_1)\cdots\exp(x_{m-1}X_{m-1})
  \in H_{R'\cap R''}(S),
  \qquad
  \alpha_m\notin R'',
  \qquad
  x_m\neq0.
\]
Then
\[
  h'=\exp(x_mX_m)\cdots\exp(x_NX_N)
\]
still satisfies the conditions imposed above on \(h\). By \SGARefLink{sga3:xxii:lem:5.6.2}{Lemma
5.6.2}, for every \(H\in\Gamma(S,\mathfrak t)\), the component of
\(\operatorname{Ad}(h')H\) in \(\mathfrak g^{\alpha_m}\) is
\[
  -\overline{\alpha_m}(H)x_mX_m.
\]
\footnote{Editor's note: In this sentence the French source prints
\(\overline{\alpha_m}(X)\); \SGARefLink{sga3:xxii:lem:5.6.2}{Lemma 5.6.2}\textup{\SGARefLink{sga3:xxii:lem:5.6.2:item:a}{(a)}} and the following
sentence require \(H\).}
It therefore follows from the hypothesis on \(h\) and from the
definition of \(m\) that
\[
  \overline{\alpha_m}(H)=0
  \qquad\text{for every }H\in\Gamma(S,\mathfrak t).
\]
This is impossible, since \(G\) is adjoint and
\(\overline{\alpha_m}\) is nonzero on every fiber.

\SGARefTarget{sga3:xxii:rem:5.6.2bis}
\SGAInlineRefTarget{sga3:xxii:remark:5-6-2-bis}\paragraph{\SGARefLink{sga3:xxii:remark:5-6-2-bis}{Remark 5.6.2 bis}.}
With the notation of \SGARefLink{sga3:xxii:lem:5.6.2}{Lemma 5.6.2}, suppose that
\(\operatorname{Ad}(u)\) induces the identity on \(\mathfrak t\) and
on \(\mathfrak g^{-\alpha_m}\). Then \(x_m=0\). Indeed,
\[
  x_m\overline{\alpha_m}=0,
  \qquad
  x_mH_{\alpha_m}=0.
\]
If \(\alpha_m\notin2M\), then \(\overline{\alpha_m}\) is nonzero on
every fiber; if \(\alpha_m\in2M\), then
\(\alpha_m^\ast\notin2M^\ast\), and \(H_{\alpha_m}\) is nonzero on
every fiber. Thus \(x_m=0\). It follows that \(u=e\) if
\(\operatorname{Ad}(u)\) acts trivially on \(\mathfrak g\).

\SGARefTarget{sga3:xxii:rem:5.6.3}
\paragraph{\SGARefLink{sga3:xxii:rem:5.6.3}{Remark 5.6.3}.}
Let \(H\) be a subgroup of type \textup{(R)} of a reductive
\(S\)-group \(G\), and suppose that its geometric fibers are
solvable. Then \(H\) is closed in \(G\), and
\[
  \underline{\operatorname{Norm}}_G(H)/H
\]
is represented by a separated finite étale \(S\)-scheme. This
follows from \SGARefLink{sga3:xxii:cor:5.3.18}{5.3.18} and either \SGARefLink{sga3:xxii:rem:3.5}{3.5} or \SGARefLink{sga3:xix:theorem:2:5}{Exposé XIX, 2.5}\textup{\SGARefLink{sga3:xix:theorem:2:5:item:ii}{(ii)}}.

\SGARefTarget{sga3:xxii:cor:5.6.4}
\paragraph{\SGARefLink{sga3:xxii:cor:5.6.4}{Corollary 5.6.4}.}
\textit{Let \((G,T,M,R)\) be a split reductive \(S\)-group. If
\(R'\subset R\) is closed and \(R'\cap(-R')=\varnothing\), then
\(R'\) is contained in a system of positive roots.}
\footnote{Editor's note: This is in fact a statement about root
systems, supplementing Exposé XXI \textup{(cf.\ \emph{Bible},
VI, \S1.7, Proposition 22)}, and it is proved here indirectly.}

\textit{Proof.}
By \SGARefLink{sga3:xxii:thm:5.4.7}{5.4.7}, \(R'\) is of type \textup{(R)}, so the result follows from
\SGARefLink{sga3:xxii:prop:5.6.1}{Proposition 5.6.1}.

\SGARefTarget{sga3:xxii:cor:5.6.5}
\paragraph{\SGARefLink{sga3:xxii:cor:5.6.5}{Corollary 5.6.5}.}
\textit{Under the hypotheses of \SGARefLink{sga3:xxii:prop:5.6.1}{Proposition 5.6.1}, multiplication
in \(G\) induces an isomorphism}
\[
  \prod_{\alpha\in R'}U_\alpha
  \xrightarrow{\ \sim\ } U_{R'},
\]
\textit{where \(U_{R'}\) is a closed subgroup scheme of \(G\),
smooth over \(S\), with connected unipotent geometric fibers, and is
independent of the chosen ordering of \(R'\). Moreover, \(H_{R'}\)
is the semidirect product \(T\cdot U_{R'}\), with \(U_{R'}\)
invariant in \(H_{R'}\).}

\textit{Proof.}
If \(R'\subset R_+\), then
\[
  U_{R'}=H_{R'}\cap U_{\geq1},
\]
with the notation of \SGARefLink{sga3:xxii:lem:5.5.6}{5.5.6}, is a closed subgroup of \(G\), of finite
presentation, and invariant in \(H_{R'}\). By \SGARefLink{sga3:xxii:prop:5.6.1}{Proposition
5.6.1}\textup{\SGARefLink{sga3:xxii:prop:5.6.1:item:iv}{(iv)}}, one has \(H_{R'}=T\cdot U_{R'}\), from which all
the assertions follow.

\SGARefTarget{sga3:xxii:rem:5.6.6}
\paragraph{\SGARefLink{sga3:xxii:rem:5.6.6}{Remark 5.6.6}.}
In particular, \(U_{R_+}\) is the group \(U_{\geq1}\) of \SGARefLink{sga3:xxii:lem:5.5.6}{5.5.6}.

Let us draw some corollaries from the preceding results concerning
groups of type \(U_{R'}\).

\SGARefTarget{sga3:xxii:cor:5.6.7}
\paragraph{\SGARefLink{sga3:xxii:cor:5.6.7}{Corollary 5.6.7}.}
\textit{Let \((G,T,M,R)\) be a split reductive \(S\)-group, and let
\(R'\) and \(R''\) be two subsets of \(R\) of type \textup{(R)},
with \(R'\cap(-R')=\varnothing\).}
\begin{enumerate}
\item[\textup{\SGAInlineRefTarget{sga3:xxii:cor:5.6.7:item:i}(i)}]
\textit{One has}
\[
  U_{R'}\cap
  \underline{\operatorname{Norm}}_G(H_{R''})
  =
  U_{R'\cap R''}.
\]

\item[\textup{\SGAInlineRefTarget{sga3:xxii:cor:5.6.7:item:ii}(ii)}]
\textit{Suppose that \(R'\) is closed. If the conditions
\(\alpha\in R'\), \(\beta\in R''\), and
\(\alpha+\beta\in R\) imply that \(\alpha+\beta\in R'\), then
\(H_{R''}\) normalizes \(U_{R'}\).}
\end{enumerate}

\textit{Proof.}
Assertion \textup{\SGARefLink{sga3:xxii:cor:5.6.7:item:i}{(i)}} follows at once from \SGARefLink{sga3:xxii:cor:5.6.5}{5.6.5} and
\SGARefLink{sga3:xxii:prop:5.6.1}{5.6.1}\textup{\SGARefLink{sga3:xxii:prop:5.6.1:item:vi}{(vi)}}. To prove \textup{\SGARefLink{sga3:xxii:cor:5.6.7:item:ii}{(ii)}}, it is enough, by \SGARefLink{sga3:xxii:lem:5.4.4}{5.4.4},
to show that \(T\) and each \(U_\beta\), for \(\beta\in R''\),
normalize \(U_{R'}\).\footnote{Editor's note: Indeed, since \(G\)
has connected fibers, it is separated over \(S\) by Exposé
VI\textsubscript{B}, \SGARefLink{cor:sga3-VIB-5.5}{5.5}; hence, by \SGARefLink{sga3:xi:corollary:6:11}{Exposé XI, 6.11} \textup{(see
also Addendum \SGARefLink{cor:sga3-VIB-6.5.5}{6.5.5} in Exposé VI\textsubscript{B})},
\(\underline{\operatorname{Norm}}_G(U_{R'})\) is represented by a
closed subgroup scheme \(N\) of \(G\). If \(N\) contains \(T\) and
the \(U_\beta\), for \(\beta\in R''\), it contains the big cell
\(\Omega_{R_+,R''}\). This cell is schematically dense in
\(H_{R''}\) by 5.4.4 and EGA IV\textsubscript{3}, 11.10.10
\textup{(the fibers of \(H_{R''}\) are integral, and
\(\Omega_{R_+,R''}\) contains a nonempty open subset of every
fiber)}. It follows that \(H_{R''}\subset N\), and therefore that
\(H_{R''}\) normalizes \(U_{R'}\).}
For \(T\) this is trivial; for \(U_\beta\) it follows from \SGARefLink{sga3:xxii:cor:5.5.2}{5.5.2} and
\SGARefLink{sga3:XXI:3.1.2}{Exposé XXI, 3.1.2}.\footnote{Editor's note: One must show that, under
the hypotheses of \textup{\SGARefLink{sga3:xxii:cor:5.6.7:item:ii}{(ii)}}, if \(\alpha\in R'\) and
\(\beta\in R''\), then every root of the form \(i\alpha+j\beta\),
with \(i,j\in\mathbb N^\ast\), belongs to \(R'\). For this purpose
the reference Exposé XXI, 2.3.5, has been replaced by Exposé XXI,
3.1.2. This may also be seen directly by inspecting the root systems
of rank \(2\).}

\SGARefTarget{sga3:xxii:cor:5.6.8}
\paragraph{\SGARefLink{sga3:xxii:cor:5.6.8}{Corollary 5.6.8}.}
\textit{Let \((G,T,M,R)\) be a split \(S\)-group, let \(R_+\) be a
system of positive roots, and let \(\alpha\) be a simple root of
\(R_+\) \textup{(that is, an element of \(R_+\) such that
\(R_+\setminus\{\alpha\}\) is closed)}. Set}
\[
  U_{\widehat\alpha}=U_{R_+\setminus\{\alpha\}}.
\]
\textit{Then:}
\begin{enumerate}
\item[\textup{\SGAInlineRefTarget{sga3:xxii:cor:5.6.8:item:i}(i)}]
\textit{\(U_{\widehat\alpha}\) is a normal subgroup of \(B_{R_+}\);}

\item[\textup{\SGAInlineRefTarget{sga3:xxii:cor:5.6.8:item:ii}(ii)}]
\textit{\(U_{R_+}\) is the semidirect product of
\(U_{\widehat\alpha}\) by \(U_\alpha\);}

\item[\textup{\SGAInlineRefTarget{sga3:xxii:cor:5.6.8:item:iii}(iii)}]
\textit{\(U_{-\alpha}\) normalizes \(U_{\widehat\alpha}\);}

\item[\textup{\SGAInlineRefTarget{sga3:xxii:cor:5.6.8:item:iv}(iv)}]
\textit{\(Z_\alpha\) normalizes \(U_{\widehat\alpha}\).}
\end{enumerate}
\textit{If one similarly defines}
\[
  U_{-\widehat\alpha}
  =
  U_{R_-\setminus\{-\alpha\}},
  \qquad
  R_-=-R_+,
\]
\textit{then}
\[
  \Omega_{R_+}
  =
  U_{-\widehat\alpha}\cdot U_{-\alpha}\cdot T\cdot
  U_\alpha\cdot U_{\widehat\alpha}.
\]

\textit{Proof.}
Assertion \textup{\SGARefLink{sga3:xxii:cor:5.6.8:item:ii}{(ii)}} follows from \SGARefLink{sga3:xxii:cor:5.6.5}{5.6.5}, and \textup{\SGARefLink{sga3:xxii:cor:5.6.8:item:i}{(i)}} from
\SGARefLink{sga3:xxii:cor:5.6.7}{5.6.7}\textup{\SGARefLink{sga3:xxii:cor:5.6.8:item:ii}{(ii)}}. Likewise, \textup{\SGARefLink{sga3:xxii:cor:5.6.8:item:iii}{(iii)}} follows from \SGARefLink{sga3:xxii:cor:5.5.2}{5.5.2}:
indeed, if \(\beta\in R_+\), with \(\beta\neq\alpha\), no
combination
\[
  i(-\alpha)+j\beta,
  \qquad i,j>0,
\]
can be negative, since \(\beta\) contains at least one simple root
different from \(\alpha\). Assertion \textup{\SGARefLink{sga3:xxii:cor:5.6.8:item:iv}{(iv)}} then follows from
\textup{\SGARefLink{sga3:xxii:cor:5.6.8:item:i}{(i)}} and \textup{\SGARefLink{sga3:xxii:cor:5.6.8:item:iii}{(iii)}}, because
\[
  U_{-\alpha}\cdot T\cdot U_\alpha
\]
is schematically dense in \(Z_\alpha\). Finally, the last assertion
follows from \textup{\SGARefLink{sga3:xxii:cor:5.6.8:item:ii}{(ii)}} and its analogue for \(U_{R_-}\).

Let us return to the general situation.

\SGARefTarget{sga3:xxii:prop:5.6.9}
\paragraph{\SGARefLink{sga3:xxii:prop:5.6.9}{Proposition 5.6.9}.}
\textit{Let \(S\) be a scheme, let \(G\) be a reductive \(S\)-group,
and let \(H\) be a subgroup of type \textup{(R)} with solvable
geometric fibers.}
\begin{enumerate}
\item[\textup{\SGAInlineRefTarget{sga3:xxii:prop:5.6.9:item:i}(i)}]
\textit{The group}
\[
  \mathrm D_S(H)
  =
  \underline{\operatorname{Hom}}_{S\text{-gr.}}
    (H,\mathbb G_{m,S})
\]
\textit{is represented by a twisted constant \(S\)-group whose type
at \(s\in S\) is}
\[
  \mathbb Z^{\operatorname{rgred}(G_s)}.
\]
\textit{The biduality morphism \textup{(Exposé VIII, \S1)}}
\[
  f:H\longrightarrow\mathrm D_S\bigl(\mathrm D_S(H)\bigr)
\]
\textit{is smooth and surjective.}

\item[\textup{\SGAInlineRefTarget{sga3:xxii:prop:5.6.9:item:ii}(ii)}]
\textit{The kernel \(H^u\) of \(f\) is the largest closed normal
subgroup scheme of \(H\), smooth over \(S\), with connected
unipotent geometric fibers. It is called the unipotent part of
\(H\), and is also denoted}
\[
  H^u=\operatorname{rad}^u(H).
\]
\textit{Moreover, \(H^u\) is the sheaf of commutators of \(H\):
every morphism of groups from \(H\) to an \(S\)-presheaf of
commutative groups that is separated for the fppf topology vanishes
on \(H^u\), and therefore factors through}
\[
  H/H^u=\mathrm D_S\bigl(\mathrm D_S(H)\bigr).
\]

\item[\textup{\SGAInlineRefTarget{sga3:xxii:prop:5.6.9:item:iii}(iii)}]
\textit{If \(T\) is a maximal torus of \(H\), the morphism
\(T\longrightarrow H\) induces isomorphisms}
\[
  \mathrm D_S(H)\xrightarrow{\ \sim\ }\mathrm D_S(T),
  \qquad
  T\xrightarrow{\ \sim\ }\mathrm D_S\bigl(\mathrm D_S(H)\bigr).
\]
\textit{Moreover, \(H\) is identified with the semidirect product
of \(H^u\) by \(T\).}

\item[\textup{\SGAInlineRefTarget{sga3:xxii:prop:5.6.9:item:iv}(iv)}]
\textit{In the situation of \SGARefLink{sga3:xxii:prop:5.6.1}{5.6.1}, if \(H=H_{R'}\), then}
\[
  H^u=U_{R'}.
\]
\end{enumerate}

\textit{Proof.}
The assertions of the proposition are local for the étale topology
\textup{(\SGARefLink{sga3:x:corollary:5-5}{Exposé X, 5.5})}. We may therefore reduce to the situation
of \SGARefLink{sga3:xxii:prop:5.6.1}{5.6.1}. We then know from \SGARefLink{sga3:xxii:cor:5.6.5}{5.6.5} that \(H_{R'}\) is the semidirect
product of \(U_{R'}\) by \(T\). Let us show that \(U_{R'}\) is the
sheaf of commutators of \(H_{R'}\). Since
\(H_{R'}/U_{R'}=T\) is commutative, it suffices to prove that every
morphism of groups
\[
  \phi:H_{R'}\longrightarrow V
\]
as in \textup{\SGARefLink{sga3:xxii:prop:5.6.9:item:ii}{(ii)}} vanishes on \(U_{R'}\). It is enough to prove
that \(\phi\) vanishes on each \(U_\alpha\), for \(\alpha\in R'\).
If \(t\in T(S')\) and \(X\in W(\mathfrak g^\alpha)(S')\), then
\[
  1
  =
  \phi\bigl(t\exp_\alpha(X)t^{-1}\exp_\alpha(-X)\bigr)
  =
  \phi\bigl(\exp_\alpha((\alpha(t)-1)X)\bigr).
\]
Since \(\alpha:T\longrightarrow\mathbb G_{m,S}\) is faithfully
flat, it follows at once that \(\phi\) vanishes on
\(U_\alpha^\times\). But every section of \(U_\alpha\) is locally
the sum of two sections of \(U_\alpha^\times\). Thus
\[
  \underline{\operatorname{Hom}}_{S\text{-gr.}}(H,V)
  =
  \underline{\operatorname{Hom}}_{S\text{-gr.}}
    (H/U_{R'},V)
\]
for every \(V\) as above. Applying this result to
\(V=\mathbb G_{m,S}\), we immediately obtain \textup{\SGARefLink{sga3:xxii:prop:5.6.9:item:i}{(i)}} and
\textup{\SGARefLink{sga3:xxii:prop:5.6.9:item:iii}{(iii)}}, then \textup{\SGARefLink{sga3:xxii:prop:5.6.9:item:iv}{(iv)}} and the second assertion of
\textup{\SGARefLink{sga3:xxii:prop:5.6.9:item:ii}{(ii)}}.

It remains to prove the first assertion of \textup{\SGARefLink{sga3:xxii:prop:5.6.9:item:ii}{(ii)}}. The only
nontrivial point is that every closed normal subgroup \(U\) of
\(H\), smooth over \(S\), with connected unipotent geometric
fibers, is a subgroup of \(H^u\). We first have:

\SGARefTarget{sga3:xxii:lem:5.6.10}
\paragraph{\SGARefLink{sga3:xxii:lem:5.6.10}{Lemma 5.6.10}.}
\textit{Let \(G\) be a reductive \(S\)-group, \(T\) a maximal
torus, and \(U\) a subgroup scheme of \(G\), smooth over \(S\),
with unipotent geometric fibers, and normalized by \(T\). Then}
\[
  U\cap T=e.
\]

\textit{Proof.}
Indeed, since
\[
  T=\underline{\operatorname{Centr}}_G(T),
\]
one has \(U\cap T=U^T\), where the right-hand side denotes the
invariants under \(\operatorname{int}(T)\). By \SGARefLink{sga3:xix:lemma:1-4}{Exposé XIX, 1.4},
\(U\cap T\) is therefore smooth over \(S\). It is also radicial
over \(S\): for every \(s\in S\), the group
\[
  U(\overline s)\cap T(\overline s)
\]
consists of elements that are both unipotent and semisimple. This
proves the lemma.

We return to the proof of \SGARefLink{sga3:xxii:prop:5.6.9}{5.6.9}\textup{\SGARefLink{sga3:xxii:prop:5.6.9:item:ii}{(ii)}}. If \(U\) is a normal
subgroup of \(H\) as above, then the semidirect product \(T\cdot U\)
is a subgroup of \(G\) of type \textup{(R)}, with solvable
geometric fibers. We may therefore suppose it to be of the form
\(H_{R''}\), with \(R''\subset R'\). It is enough to prove that
\(U=U_{R''}\), and we are thus reduced to the case
\(H=T\cdot U\). Since \(H/U\) is commutative, the French text says
that \(U\) is a subsheaf of the sheaf of commutators of \(H\),
which is \(H^u\).\footnote{Editor's note: The implication in the
French is stated in the wrong direction: commutativity of \(H/U\)
gives \(H^u\subset U\). The desired inclusion follows because
\(H/H^u\) is a torus whereas \(U\) has connected unipotent
geometric fibers.}
\hfill\(\square\)

Let us note that we have in fact just proved:

\SGARefTarget{sga3:xxii:prop:5.6.11}
\paragraph{\SGARefLink{sga3:xxii:prop:5.6.11}{Proposition 5.6.11}.}
\textit{Let \(S\) be a scheme, \(G\) a reductive \(S\)-group, and
\(T\) a maximal torus of \(G\). The maps}
\[
  H\longmapsto H^u,
  \qquad
  U\longmapsto T\cdot U
\]
\textit{are mutually inverse bijections between the set of subgroups
\(H\) of \(G\) of type \textup{(R)}, containing \(T\), and having
solvable geometric fibers, and the set of subgroups \(U\) of \(G\),
smooth over \(S\), normalized by \(T\), and having connected
unipotent geometric fibers.}\footnote{Editor's note: The hypothesis
that \(U\) be closed has been omitted, since it is automatic. Indeed,
for such a \(U\), one has \(U\cap T=e\) by \SGARefLink{sga3:xxii:lem:5.6.10}{5.6.10}, so the semidirect
product \(H=T\cdot U\) is a subgroup of \(G\) of type \textup{(R)}
\textup{(cf. \SGARefLink{sga3:xxii:def:5.2.1}{5.2.1})}, with solvable geometric fibers. Thus, by \SGARefLink{sga3:xxii:rem:5.6.3}{5.6.3},
\(H\) is closed in \(G\); and since \(U\) is closed in \(H\), it is
closed in \(G\).}

In particular, when \((G,T,M,R)\) is split, the groups \(H_{R'}\)
and \(U_{R'}\) correspond.

\SGARefTarget{sga3:xxii:cor:5.6.12}
\paragraph{\SGARefLink{sga3:xxii:cor:5.6.12}{Corollary 5.6.12}.}
\textit{Let \(S\) be a scheme, let \((G,T,M,R)\) be a split
\(S\)-group \textup{(respectively, and let \(R_+\) be a system of
positive roots of \(R\) defining the Borel subgroup \(B\))}.
Every smooth group subscheme of \(G\) with connected unipotent
geometric fibers \textup{(respectively, every smooth group subscheme
of \(B^u\))}, normalized by \(T\), is locally on \(S\) of the form
\(U_{R'}\), where \(R'\) is a subset of \(R\) contained in a system
of positive roots \textup{(respectively, a subset of \(R_+\))} and
is of type \textup{(R)}.}

\textit{Proof.}
For the ``respectively'' case, it is enough to observe that the
geometric fibers of the given group are unipotent and connected by
the \emph{Bible}, \S13.2, Theorem 1\textup{\SGARefLink{sga3:xxii:rem:5.1.2:item:d}{(d)}}.
\hfill\(\square\)

\SGARefLink{sga3:xxii:prop:5.6.9}{Proposition 5.6.9} also has the following corollary:

\SGARefTarget{sga3:xxii:cor:5.6.13}
\paragraph{\SGARefLink{sga3:xxii:cor:5.6.13}{Corollary 5.6.13}.}
\textit{Let \(S\) be a scheme, \(G\) a reductive \(S\)-group, and
\(H\) a subgroup of type \textup{(R)} with solvable geometric
fibers. Let \(\underline{\operatorname{Tor}}(H)\) be the functor of
maximal tori of \(H\):}
\[
  \underline{\operatorname{Tor}}(H)(S')
  =
  \{\,\text{maximal tori of }H_{S'}\,\}.
\]
\textit{Then \(\underline{\operatorname{Tor}}(H)\) is represented by
an affine smooth \(S\)-scheme, which is a principal homogeneous
bundle under \(H^u\) for the action}
\[
  (h,T)\longmapsto\operatorname{int}(h)T.
\]
\footnote{Editor's note: This reproves and sharpens \SGARefLink{sga3:xii:1.10}{Exposé XII, 1.10}
\textup{(for \(G\) reductive)}.}

\textit{Proof.}
Indeed, if \(T\) and \(T'\) are two maximal tori of \(H_{S'}\), there
exists a unique section \(h\in H^u(S')\) such that
\(\operatorname{int}(h)T=T'\). The uniqueness of \(h\) follows at
once from the equality
\[
  \underline{\operatorname{Norm}}_G(T)\cap H^u=e
\]
\textup{(cf., for example, \SGARefLink{sga3:xxii:prop:5.6.1}{5.6.1})}; it is therefore enough to prove
the existence of \(h\) locally for the étale topology. By \SGARefLink{sga3:xxii:prop:5.2.6}{5.2.6} and
\SGARefLink{sga3:xxii:rem:5.1.2}{5.1.2}\textup{\SGARefLink{sga3:xxii:rem:5.1.2:item:a}{(a)}}, we may suppose that \(T\) and \(T'\) are conjugate
by a section of \(H\), whence the desired conclusion because
\[
  H=H^u\cdot T
\]
by \SGARefLink{sga3:xxii:prop:5.6.9}{5.6.9}\textup{\SGARefLink{sga3:xxii:prop:5.6.9:item:iii}{(iii)}}. It follows that
\(\underline{\operatorname{Tor}}(H)\) is a principal homogeneous
sheaf under \(H^u\), which is affine and smooth over \(S\); the
assertion follows at once.\footnote{Editor's note: By SGA 1, VIII,
2.1, and EGA \(\mathrm{IV}_4\), 17.7.1.}
\hfill\(\square\)

\SGARefTarget{sga3:xxii:sec:5.7}
\subsection*{\SGAInlineRefTarget{sga3:xxii:subsection:5-7}5.7. Bruhat's Theorem}
\addcontentsline{toc}{subsection}{\SGARefLink{sga3:xxii:subsection:5-7}{5.7}. Bruhat's Theorem}

\SGARefTarget{sga3:xxii:recall:5.7.1}
\paragraph{Recall \SGARefLink{sga3:xxii:recall:5.7.1}{5.7.1}.}
\textit{Let \(k\) be an algebraically closed field, \(G\) a
reductive \(k\)-group, \(B\) a Borel subgroup of \(G\), \(T\) a
maximal torus of \(B\), and
\(N=\underline{\operatorname{Norm}}_G(T)\). Then}
\[
  G(k)=B(k)N(k)B(k);
\]
\textit{this is Bruhat's theorem \textup{(\emph{Bible}, \S13.4,
Corollary 1 to Theorem 3)}. More precisely, with the notation of
\SGARefLink{sga3:xxii:par:3.6}{3.6}, the sets}
\[
  B(k)N_w(k)B(k)=B^u(k)N_w(k)B^u(k)
\]
\textit{form a partition of \(G(k)\) as \(w\) ranges over
\((N/T)(k)\). If \(B'\) is another Borel subgroup of \(G\)
containing \(T\), the sets \(B'(k)N_w(k)B(k)\) likewise form a
partition of \(G(k)\). Indeed, if \(y\in N(k)\) is such that
\(\operatorname{int}(y)B=B'\), then}
\[
  yB(k)N_w(k)B(k)=B'(k)N_{yw}(k)B(k).
\]

\SGARefTarget{sga3:xxii:def:5.7.2}
\paragraph{\SGARefLink{sga3:xxii:def:5.7.2}{Definition 5.7.2}.}
\textit{Let \((G,T,M,R)\) be a split \(S\)-group, let \(R_-\) be a
system of positive roots of \(R\), and let \(B'=B_{R_-}\) be the
Borel subgroup it defines. For \(w\in W\), put
\textup{(cf. \SGARefLink{sga3:xxii:cor:5.6.5}{5.6.5})}}
\[
  R_-^w=R_-\cap w(R_-),
  \qquad
  B'{}^u_w=U_{R_-^w}
  =\prod_{\alpha\in R_-^w}U_\alpha.
\]
\textit{If
\(n_w\in\underline{\operatorname{Norm}}_G(T)(S)\) is a
representative of \(w\) \textup{(\SGARefLink{sga3:xxii:cor:3.8}{3.8})}, one may also write}
\[
  B'{}^u_w=B'{}^u\cap\operatorname{int}(n_w)B'{}^u.
\]

\SGARefTarget{sga3:xxii:lem:5.7.3}
\paragraph{\SGARefLink{sga3:xxii:lem:5.7.3}{Lemma 5.7.3}.}
\textit{Let \((G,T,M,R)\) be a split \(S\)-group, let \(R_+\) be a
system of positive roots of \(R\), let \(R_-=-R_+\), and let \(B\)
\textup{(respectively \(B'\))} be the Borel subgroup of \(G\)
defined by \(R_+\) \textup{(respectively \(R_-\))}. Let \(w\in W\),
and let \(N_w\) and \(B'{}^u_w\) be the corresponding subschemes of
\(G\) \textup{(\SGARefLink{sga3:xxii:cor:3.8}{3.8} and \SGARefLink{sga3:xxii:def:5.7.2}{5.7.2})}.}
\begin{enumerate}
\item[\textup{\SGAInlineRefTarget{sga3:xxii:lem:5.7.3:item:i}(i)}]
\textit{The sheaf \(B'\cdot N_w\cdot B\), the image of the
morphism}
\[
  B'\mathop{\times}\limits_S N_w
  \mathop{\times}\limits_S B
  \longrightarrow G
\]
\textit{induced by multiplication in \(G\), is represented by a
subscheme of \(G\) \textup{(in fact, by a closed subscheme of the
open subscheme \(n_w\Omega_{R_+}\))}.}

\item[\textup{\SGAInlineRefTarget{sga3:xxii:lem:5.7.3:item:ii}(ii)}]
\textit{The morphism}
\[
  B'{}^u_w\mathop{\times}\limits_S N_w
  \mathop{\times}\limits_S B^u
  \longrightarrow G,
\]
\textit{induced by multiplication in \(G\), is an immersion whose
image is the preceding subscheme.}
\end{enumerate}

\textit{Proof.}
Let us first show that the morphism in \textup{\SGARefLink{sga3:xxii:lem:5.7.3:item:ii}{(ii)}} is an
immersion. By definition, \(\operatorname{int}(n_w)^{-1}\) induces
a closed immersion of \(B'{}^u_w\) into \(B'{}^u\); hence the
morphism
\[
  (u,b)\longmapsto n_w^{-1}u\,n_wb
\]
induces a closed immersion
\[
  B'{}^u_w\mathop{\times}\limits_S B
  \longrightarrow \Omega_{R_+}.
\]
It follows immediately that the morphism in \textup{\SGARefLink{sga3:xxii:lem:5.7.3:item:ii}{(ii)}} induces a
closed immersion of its source into the open subscheme
\(n_w\Omega_{R_+}\). To prove \textup{\SGARefLink{sga3:xxii:lem:5.7.3:item:i}{(i)}}, it is enough to see that
\[
  B'(S)N_w(S)B(S)
  =
  B'{}^u_w(S)N_w(S)B^u(S).
\]
Now, if \(\alpha\in R\), then
\(\operatorname{int}(n_w)U_\alpha(S)=U_{w(\alpha)}(S)\). Thus, if
\(w^{-1}(\alpha)\in R_+\),
\[
\begin{aligned}
  U_\alpha(S)N_w(S)B(S)
  &=U_\alpha(S)n_wT(S)B^u(S)\\
  &=n_wU_{w^{-1}(\alpha)}(S)T(S)B^u(S)\\
  &=n_wB(S)=N_w(S)B^u(S).
\end{aligned}
\]
In view of the definition of \(R_-^w\), this immediately gives the
desired assertion.
\hfill\(\square\)

\SGARefTarget{sga3:xxii:thm:5.7.4}
\paragraph{\SGARefLink{sga3:xxii:thm:5.7.4}{Theorem 5.7.4}.}
\textit{Let \(S\) be a scheme, let \((G,T,M,R)\) be a split
\(S\)-group, let \(B\) be the Borel subgroup defined by the system
of positive roots \(R_+\), and let \(B'\) be the Borel subgroup
defined by \(R_-=-R_+\).}
\begin{enumerate}
\item[\textup{(i)}]
\textit{\textup{(Bruhat's theorem)} The schemes
\(B'{}^u_w\cdot N_w\cdot B\), as \(w\) ranges over \(W\), form a
partition of the underlying set of \(G\).}

\item[\textup{\SGAInlineRefTarget{sga3:xxii:thm:5.7.4:item:ii}(ii)}]
\textit{For each \(w\in W\), let \(n_w\) be a representative of
\(w\) in \(\underline{\operatorname{Norm}}_G(T)(S)\)
\textup{(\SGARefLink{sga3:xxii:cor:3.8}{3.8})}. Then the open subschemes}
\[
  n_w\Omega=n_wB'{}^u\cdot T\cdot B^u
\]
\textit{as \(w\) ranges over \(W\), form an open covering of \(G\).}
\end{enumerate}

\textit{Proof.}
Both assertions may be checked on the geometric fibers, where they
follow from \SGARefLink{sga3:xxii:recall:5.7.1}{5.7.1} and \SGARefLink{sga3:xxii:lem:5.7.3}{5.7.3}.
\hfill\(\square\)

\SGARefTarget{sga3:xxii:rem:5.7.5}
\paragraph{\SGARefLink{sga3:xxii:rem:5.7.5}{Remark 5.7.5}.}
Assertion \textup{\SGARefLink{sga3:xxii:lem:5.7.3:item:i}{(i)}} implies that if \(S\) is the spectrum of a
field, then \(G(S)\) is the disjoint union of the sets
\[
  B'{}^u_w(S)\cdot T(S)\cdot B^u(S).
\]
The corresponding assertion for an arbitrary \(S\), even a local
or Artinian one, is plainly false. Note, however, that
\textup{\SGARefLink{sga3:xxii:thm:5.7.4:item:ii}{(ii)}} implies that if \(S\) is local, then \(G(S)\) is
indeed the union of the \(n_w\Omega(S)\). In fact:

\SGARefTarget{sga3:xxii:cor:5.7.6}
\paragraph{\SGARefLink{sga3:xxii:cor:5.7.6}{Corollary 5.7.6}.}
\textit{Let \(\Delta\) be a system of simple roots of the split
group \(G\) over the local scheme \(S\).}
\begin{enumerate}
\item[\textup{\SGAInlineRefTarget{sga3:xxii:cor:5.7.6:item:i}(i)}]
\textit{Then \(G(S)\) is generated by \(T(S)\) and the groups
\(U_\alpha(S)\), for \(\alpha\in\Delta\cup-\Delta\).}

\item[\textup{\SGAInlineRefTarget{sga3:xxii:cor:5.7.6:item:ii}(ii)}]
\textit{If \(G\) is simply connected \textup{(\SGARefLink{sga3:xxii:def:4.3.3}{4.3.3})}, then
\(G(S)\) is already generated by the groups \(U_\alpha(S)\), for
\(\alpha\in\Delta\cup-\Delta\).}
\end{enumerate}

\textit{Proof.}
Let \(H\) be the subgroup of \(G(S)\) generated by the
\(U_\alpha(S)\), for \(\alpha\in\Delta\cup-\Delta\). First observe
that \(H\) contains a representative of every \(s_\alpha\), for
\(\alpha\in\Delta\), in
\(\underline{\operatorname{Norm}}_G(T)(S)\) \textup{(\SGARefLink{sga3:xx:3.1}{Exposé XX,
3.1})}, and hence a representative \(n_w\) of every \(w\in W\).
Since every \(\alpha\in R\) can be written as \(w(\alpha_0)\), with
\(w\in W\) and \(\alpha_0\in\Delta\), we have
\[
  U_\alpha(S)
  =
  \operatorname{int}(n_w)U_{\alpha_0}(S)
  \subset H.
\]
The subgroup generated by \(H\) and \(T(S)\) therefore contains
\(\Omega(S)\), and hence, by the preceding observation, is all of
\(G(S)\).

Now suppose that \(G\) is simply connected. Let us prove that
\(H\supset T(S)\). By \SGARefLink{sga3:xx:2.7}{Exposé XX, 2.7}, \(H\) contains
\(\alpha^*(\mathbb G_m(S))\) for every \(\alpha\in\Delta\), and it
remains only to apply \SGARefLink{sga3:xxii:rem:4.3.8}{4.3.8}.
\hfill\(\square\)

\SGARefTarget{sga3:xxii:rem:5.7.6.1}
\paragraph{\SGARefLink{sga3:xxii:rem:5.7.6.1}{Remark 5.7.6.1}.}
Instead of taking, for each \(\alpha\in\Delta\), both \(U_\alpha(S)\)
and \(U_{-\alpha}(S)\), it is enough to take \(U_\alpha(S)\) together
with a representative \(w_\alpha\) of the reflection \(s_\alpha\), or
else \(U_\alpha(S)\) together with a section of
\(U_{-\alpha}^{\times}\), \ldots.

\SGARefTarget{sga3:xxii:cor:5.7.7}
\paragraph{\SGARefLink{sga3:xxii:cor:5.7.7}{Corollary 5.7.7}.}
\textit{Suppose that \(G\) has semisimple rank \(1\), and choose
\(u_\alpha\in U_\alpha^\times(S)\). Then \(\Omega\) and
\(u_\alpha\Omega\) form an open covering of \(G\).}

\textit{Proof.}
If \(u_{-\alpha}\) is the section of \(U_{-\alpha}\) paired with
\(u_\alpha\) \textup{(cf.\ \SGARefLink{sga3:xxii:not:1.3}{1.3})}, then \SGARefLink{sga3:xxii:thm:5.7.4}{5.7.4} \textup{\SGARefLink{sga3:xxii:thm:5.7.4:item:ii}{(ii)}} gives
\[
  G
  =
  \Omega\cup
  u_{-\alpha}^{-1}u_\alpha u_{-\alpha}^{-1}\Omega,
\]
and hence
\[
  G
  =
  u_{-\alpha}G
  =
  u_{-\alpha}\Omega\cup
  u_\alpha u_{-\alpha}^{-1}\Omega
  =
  \Omega\cup u_\alpha\Omega.
\]
\hfill\(\square\)

\SGARefTarget{sga3:xxii:cor:5.7.8}
\paragraph{\SGARefLink{sga3:xxii:cor:5.7.8}{Corollary 5.7.8}.}
\textit{Let \(S\) be a scheme and let \(G\) be a reductive
\(S\)-group. Then \(G\) is essentially free over \(S\) in the sense
of \SGARefLink{sga3:viii:definition:6:1}{Exposé VIII, 6.1}.}\footnote{Editor's note: See also the additions
made in Exposé VI\textsubscript{B}, \SGARefLink{def:sga3-VIB-6.2.1}{6.2.1}--\SGARefLink{sga3:xxii:par:6.2.6}{6.2.6} and
\SGARefLink{thm:sga3-VIB-6.5.2}{6.5.2}--\SGARefLink{cor:sga3-VIB-6.5.5}{6.5.5}.}

\textit{Proof.}
The assertion is local for the fpqc topology, so we may assume that
\(G\) is split. It then admits an open covering by open subschemes
isomorphic to
\[
  \mathbb G_{a,S}^{N}\mathop{\times}\limits_S\mathbb G_{m,S}^{n},
\]
and these are essentially free.
\hfill\(\square\)

\SGARefTarget{sga3:xxii:lem:5.7.9}
\paragraph{\SGARefLink{sga3:xxii:lem:5.7.9}{Lemma 5.7.9}.}
\textit{Under the hypotheses of \SGARefLink{sga3:xxii:thm:5.7.4}{5.7.4}, let \(\alpha\) be a simple
root of \(R_+\), and let \(u_\alpha\in U_\alpha^\times(S)\). For
every \(v\in U_{-\alpha}(S)\), one has}
\[
  \Omega\cdot v\subset\Omega\cup u_\alpha\cdot\Omega.
\]

\textit{Proof.}
We are comparing two open subschemes of \(G\), so it is enough to do
so when \(S\) is the spectrum of a field \(k\). We must therefore
prove
\[
  \Omega(k)v\subset\Omega(k)\cup u_\alpha\Omega(k).
\]
Now
\[
\begin{aligned}
  \Omega(k)v
  &=B'{}^u(k)T(k)B^u(k)v\\
  &=U_{-\widehat\alpha}(k)U_{-\alpha}(k)T(k)
    U_\alpha(k)U_{\widehat\alpha}(k)v\\
  &\subset
    U_{-\widehat\alpha}(k)Z_\alpha(k)
    U_{\widehat\alpha}(k)v.
\end{aligned}
\]
Here we use the decomposition of \SGARefLink{sga3:xxii:cor:5.6.8}{5.6.8}. Applying \SGARefLink{sga3:xxii:cor:5.6.8}{5.6.8}
\textup{\SGARefLink{sga3:xxii:cor:5.6.8:item:iii}{(iii)}} and then \SGARefLink{sga3:xxii:cor:5.7.7}{5.7.7} to the group \(Z_\alpha\), we obtain
\[
\begin{aligned}
  \Omega(k)v
  &\subset
    U_{-\widehat\alpha}(k)Z_\alpha(k)v
    U_{\widehat\alpha}(k)\\
  &\subset
    U_{-\widehat\alpha}(k)Z_\alpha(k)
    U_{\widehat\alpha}(k)\\
  &\subset
    U_{-\widehat\alpha}(k)U_{-\alpha}(k)T(k)
    U_\alpha(k)U_{\widehat\alpha}(k)\\
  &\qquad\cup
    U_{-\widehat\alpha}(k)u_\alpha U_{-\alpha}(k)T(k)
    U_\alpha(k)U_{\widehat\alpha}(k).
\end{aligned}
\]
Applying \SGARefLink{sga3:xxii:cor:5.6.8}{5.6.8} \textup{\SGARefLink{sga3:xxii:cor:5.6.8:item:iii}{(iii)}} once more, now with \(R_-\) in place
of \(R_+\), gives the result.
\hfill\(\square\)

\SGARefTarget{sga3:xxii:prop:5.7.10}
\paragraph{\SGARefLink{sga3:xxii:prop:5.7.10}{Proposition 5.7.10}.}
\textit{Under the hypotheses of \SGARefLink{sga3:xxii:thm:5.7.4}{5.7.4}, choose, for each simple
root \(\alpha\), an element
\(u_\alpha\in U_\alpha^\times(S)\). Let \(U_1\) be the submonoid of
\(B^u(S)\) generated by the \(u_\alpha\). Then the open subschemes
\(u\Omega\), for \(u\in U_1\), form an open covering of \(G\).}

\textit{Proof.}
Once again, we may assume that \(S\) is the spectrum of a field
\(k\). In view of \SGARefLink{sga3:xxii:cor:5.7.6}{5.7.6}, it is enough to prove that
\[
  \bigcup_{u\in U_1}u\Omega(k)
\]
is stable under right multiplication by \(T(k)\), \(U_\alpha(k)\),
and \(U_{-\alpha}(k)\), for every simple root \(\alpha\). This is
trivial in the first two cases, and in the last it follows from the
lemma.
\hfill\(\square\)

\SGARefTarget{sga3:xxii:rem:5.7.11}
\paragraph{\SGARefLink{sga3:xxii:rem:5.7.11}{Remark 5.7.11}.}
Let us point out a special case of \SGARefLink{sga3:xxii:def:5.7.2}{5.7.2}. If \(w=s_\alpha\) is the
reflection with respect to the simple root \(\alpha\), then
\[
  R_-\cap s_\alpha(R_-)=R_-\setminus\{-\alpha\}
\]
\textup{(\SGARefLink{sga3:XXI:3.3.1}{Exposé XXI, 3.3.1})}, and hence, in the notation of \SGARefLink{sga3:xxii:cor:5.6.8}{5.6.8},
\[
  B'{}^u_{s_\alpha}=U_{-\widehat\alpha}.
\]

\SGARefTarget{sga3:xxii:rem:5.7.12}
\paragraph{\SGARefLink{sga3:xxii:rem:5.7.12}{Remark 5.7.12}.}
In fact, the proof of \SGARefLink{sga3:xxii:prop:5.7.10}{5.7.10} immediately gives the following
statement. Under the hypotheses of \SGARefLink{sga3:xxii:prop:5.7.10}{5.7.10}, let \(\Gamma\) be a
submonoid of \(G(S)\). The open subschemes \(g\Omega\), for
\(g\in\Gamma\), form an open covering of \(G\) if and only if, for
every \(s\in S\) and every simple root \(\alpha\), one has
\[
  (u_\alpha)_{\bar{s}}B'{}^u(\bar{s})
  \subset
  \Gamma\cdot B'{}^u(\bar{s})\cdot T(\bar{s})\cdot B^u(\bar{s}).
\]

\SGARefTarget{sga3:xxii:rem:5.7.13}
\paragraph{\SGARefLink{sga3:xxii:rem:5.7.13}{Remark 5.7.13}.}
By \SGARefLink{sga3:xxii:cor:5.5.5}{5.5.5} \textup{\SGARefLink{sga3:xxii:cor:5.5.5:item:iii}{(iii)}}, arguing as in \SGARefLink{sga3:xxii:recall:5.7.1}{5.7.1}, one immediately
obtains the following variant of \SGARefLink{sga3:xxii:thm:5.7.4}{5.7.4}. Let \((G,T,M,R)\) be a
split \(S\)-group, and let \(B\) and \(B'\) be two Borel subgroups
of \(G\) containing \(T\). For every \(w\in W\), the sheaf
\(B'\cdot N_w\cdot B\) is representable by a subscheme of \(G\);
these subschemes, as \(w\) ranges over \(W\), form a partition of
the underlying set of \(G\). One can also give the analogue of
\SGARefLink{sga3:xxii:lem:5.7.3}{5.7.3} \textup{\SGARefLink{sga3:xxii:lem:5.7.3:item:ii}{(ii)}}: one must put
\[
  B'{}^u_w
  =
  B'{}^u\cap\operatorname{int}(n_w)\widetilde B{}^u,
\]
where \(\widetilde B\) is the Borel subgroup opposite to \(B\) with
respect to \(T\) \textup{(cf.\ \SGARefLink{sga3:xxii:prop:5.9.2}{5.9.2})}.

\SGARefTarget{sga3:xxii:prop:5.7.14}
\paragraph{\SGARefLink{sga3:xxii:prop:5.7.14}{Proposition 5.7.14}.}
\textit{Let \(S\) be a scheme, let \(G\) be a reductive
\(S\)-group, and let}
\[
  \operatorname{Ad}\colon
  G\longrightarrow GL_{\mathcal O_S}(\mathfrak g)
\]
\textit{be its adjoint representation. Then}
\[
  \operatorname{Ker}(\operatorname{Ad})
  =
  \underline{\operatorname{Centr}}(G).
\]
\textit{Equivalently, the canonical homomorphism induced by
\(\operatorname{Ad}\) upon passage to the quotient,}
\[
  \overline{\operatorname{Ad}}\colon
  G/\underline{\operatorname{Centr}}(G)
  =
  \operatorname{ad}(G)
  \longrightarrow
  GL_{\mathcal O_S}(\mathfrak g),
\]
\textit{is a monomorphism.}\footnote{Editor's note: It is even a
closed immersion, by \SGARefLink{sga3:xvi:corollary:1:5}{Exposé XVI, 1.5}\textup{\SGARefLink{sga3:xxii:lem:5.6.2:item:a}{(a)}}.}

\textit{Proof.}
We may assume that \(G\) is split. Choose a total order on
\(\Gamma_0(R)\) compatible with its group structure, and let
\(R_+\) be the set of positive roots. By \SGARefLink{sga3:xxii:thm:5.7.4}{5.7.4} \textup{\SGARefLink{sga3:xxii:thm:5.7.4:item:ii}{(ii)}} and
\SGARefLink{sga3:xxii:cor:4.1.6}{4.1.6}, it is enough to prove that, if \(n_w\) represents
\(w\in W\), if
\[
  u\in U(S),\qquad t\in T(S),\qquad v\in U_-(S),
\]
and if \(\operatorname{Ad}(n_wvtu)=\operatorname{id}\), then
\(w=e\), \(v=e\), and \(u=e\). For each \(m\in R\cup\{0\}\), put
\[
  \mathfrak g^{>m}
  =
  \bigoplus_{n>m}\mathfrak g^n,
  \qquad
  \mathfrak g^{<m}
  =
  \bigoplus_{n<m}\mathfrak g^n.
\]
Let \(X\in\Gamma(S,\mathfrak g^m)\); write
\[
  \operatorname{Ad}(tu)X
  =
  \operatorname{Ad}(v^{-1}n_w^{-1})X.
\]
Now
\[
\begin{aligned}
  \operatorname{Ad}(t)\operatorname{Ad}(u)X-m(t)X
  &\in\Gamma(S,\mathfrak g^{>m}),\\
  \operatorname{Ad}(v^{-1}n_w^{-1})X
  -\operatorname{Ad}(n_w^{-1})X
  &\in\Gamma\bigl(S,\mathfrak g^{<w^{-1}(m)}\bigr).
\end{aligned}
\]
If \(w\ne e\), there is an \(\alpha\in R\) such that
\(w^{-1}(\alpha)<\alpha\). On taking \(m=\alpha\), this gives a
contradiction, since
\[
\begin{aligned}
  \operatorname{Ad}(tu)X
  &\in
  \Gamma\bigl(S,\mathfrak g^\alpha+\mathfrak g^{>\alpha}\bigr)
  \cap
  \Gamma\bigl(
    S,
    \mathfrak g^{w^{-1}(\alpha)}
    +\mathfrak g^{<w^{-1}(\alpha)}
  \bigr)\\
  &=0.
\end{aligned}
\]
Thus \(w=e\), and we may choose \(n_w=e\). We then have
\[
  \operatorname{Ad}(v^{-1})X-X
  \in
  \Gamma\bigl(
    S,
    \mathfrak g^{<m}\cap
    (\mathfrak g^m+\mathfrak g^{>m})
  \bigr)
  =0,
\]
whence \(\operatorname{Ad}(v)X=X\) for every
\(X\in\Gamma(S,\mathfrak g^m)\), and therefore
\(\operatorname{Ad}(v)=\operatorname{id}\). Similarly,
\(\operatorname{Ad}(u)=\operatorname{id}\). The result now follows
from \SGARefLink{sga3:xxii:remark:5-6-2-bis}{5.6.2 bis}.
\hfill\(\square\)

\SGARefTarget{sga3:xxii:section:5.8}
\section*{5.8. Schemes Associated with a Reductive Group}
\addcontentsline{toc}{section}{\SGARefLink{sga3:xxii:section:5.8}{5.8}. Schemes Associated with a
Reductive Group}

\SGARefTarget{sga3:xxii:thm:5.8.1}
\paragraph{\SGARefLink{sga3:xxii:thm:5.8.1}{Theorem 5.8.1}.}
\textit{Let \(S\) be a scheme and let \(G\) be a reductive
\(S\)-group. Let \(\mathcal H\) be the functor of subgroups of type
\textup{(R)} of \(G\): for every \(S'\to S\),
\(\mathcal H(S')\) is the set of subgroups of type \textup{(R)} of
\(G_{S'}\) \textup{(cf.\ \SGARefLink{sga3:xxii:def:5.2.1}{5.2.1})}. Then \(\mathcal H\) is
representable by a quasi-projective \(S\)-scheme of finite
presentation over \(S\).}

\textit{Proof.}\footnote{Editor's note: The original is expanded in
what follows.}
Let
\[
  G'=G/\underline{\operatorname{Centr}}(G)
\]
be the adjoint group of \(G\) \textup{(\SGARefLink{sga3:xxii:def:4.3.6}{4.3.6})}, and let
\(u\colon G\to G'\) be the quotient morphism. By \SGARefLink{sga3:xii:7.12}{Exposé XII, 7.12},
the map
\[
  H'\longmapsto u^{-1}(H')
\]
is a bijection from the set of subgroups of type \textup{(R)} of
\(G'\) onto the set of subgroups of type \textup{(R)} of \(G\);
this remains true after every base change. Thus, replacing \(G\)
by \(G'\), we may assume that \(G\) is adjoint. Consider the
morphism
\[
  u\colon
  \mathcal H
  \longrightarrow
  \underline{\operatorname{Grass}}(\mathfrak g)
\]
that assigns to every subgroup of type \textup{(R)} its Lie algebra
\textup{(a submodule of \(\mathfrak g\) that is locally a direct
summand\footnote{Editor's note: Cf.\ \SGARefLink{sga3:ii:corollary:4-11-8}{Exposé II, 4.11.8}.})}.
Then \(u\) is a monomorphism by \SGARefLink{sga3:xxii:cor:5.3.3}{5.3.3}. It is enough to prove that
it is representable by an immersion of finite presentation; in
other words, it is enough to prove the following assertion. For
every \(S_1\to S\), given a locally direct-summand submodule
\(\mathfrak h\) of \(\mathfrak g_{S_1}\), the morphisms
\(S'\to S_1\) for which \(\mathfrak h_{S'}\) is the Lie algebra of
a subgroup of type \textup{(R)} of \(G_{S'}\) are exactly those
that factor through a certain subscheme \(\Sigma\) of finite
presentation over \(S_1\). Replacing \(S_1\) by \(S\), we reduce
to \(S_1=S\), and we may moreover assume that \(S\) is affine;
then there exists an affine Noetherian scheme \(S_0\) such that
\(G\) \textup{(respectively, \(\mathfrak h\))} is obtained by base
change from an adjoint reductive \(S_0\)-group \(G_0\)
\textup{(respectively, from a locally direct-summand submodule
\(\mathfrak h_0\) of
\(\mathfrak g_0=\operatorname{Lie}(G_0/S_0)\))}. It is enough to
show that there exists a subscheme \(\Sigma_0\) of \(S_0\) having
the required properties, for then
\(\Sigma=\Sigma_0\times_{S_0}S\). Replacing \(S\) by \(S_0\), we
may therefore suppose that \(S\) is affine and Noetherian
\textup{(note that then every subscheme of \(S\) is of finite
presentation over \(S\))}. Finally, after replacing \(S\) by a
sufficiently small open subscheme, we may suppose that
\(\mathfrak g\) is free of rank \(n\) and that \(\mathfrak h\) is
a direct summand, free of rank \(r\).

We must first express the condition that \(\mathfrak h_{S'}\) be
a Lie subalgebra of \(\mathfrak g_{S'}\), that is, that the
morphism induced by the Lie bracket
\[
  \varphi\colon
  \mathfrak h\otimes\mathfrak h
  \xrightarrow{[\, ,\,]}
  \mathfrak g
\]
factor through \(\mathfrak h\). If
\((e_1,\ldots,e_n)\) is a basis of \(\mathfrak g\) such that
\((e_1,\ldots,e_r)\) is a basis of \(\mathfrak h\), then
\(\varphi\) is given by sections \(a_k^{ij}\) of
\(\mathcal O_S\), where \(i,j=1,\ldots,r\) and
\(k=1,\ldots,n\). The preceding condition is equivalent to
requiring that \(S'\to S\) factor through the closed subscheme of
\(S\) defined by the equations
\[
  a_k^{ij}=0
  \qquad
  \text{for \(k=r+1,\ldots,n\) and \(i,j=1,\ldots,r\).}
\]
We replace \(S\) by this closed subscheme.

Then, by \SGARefLink{sga3:xxii:recall:5.3.0}{5.3.0},
\[
  N=\underline{\operatorname{Norm}}_G(\mathfrak h)
\]
is a closed subgroup scheme of \(G\), of finite presentation over
\(G\). We must now express \textup{(cf.\ \SGARefLink{sga3:xxii:prop:5.3.1}{5.3.1})} the condition
that \(N_{S'}\) be smooth at every point of the unit section, of
relative dimension \(r=\operatorname{rank}(\mathfrak h)\), and
that the inclusion of \(\mathfrak h_{S'}\) in
\(\operatorname{Lie}(N_{S'}/S')\) be an equality.

Since \(N\) is affine over \(S\)\footnote{Editor's note: The
original then indicated that, denoting by
\(\mathfrak n_{N/G}\) the inverse image of the conormal sheaf
\(\mathcal N_{N/G}\) under the unit section \(S\to N\), the
condition that
\(\mathfrak n_{N_{S'}/G_{S'}}\to
\omega^1_{G_{S'}/S'}\)
be universally injective is equivalent to the condition that
\(S'\to S\) factor through a certain open subscheme of \(S\). We
were unable to justify this point, since formation of
\(\mathcal N_{N/G}\) does not commute with base change, and we
have replaced that argument by the one that follows, suggested by
O.\ Gabber.}
\textup{(being closed in \(G\))}, the unit section
\(\varepsilon\colon S\to N\) is a closed immersion. Thus
\(\varepsilon(S)\) is defined by a quasi-coherent ideal
\(\mathcal J\) of \(\mathcal A(N)\). Note that
\[
  \mathfrak n_{S/N}=\mathcal J/\mathcal J^2
  \simeq \varepsilon^*(\Omega^1_{N/S}),
\]
so its formation commutes with every base change \(S'\to S\).
By the equivalence \textup{\((c')\Longleftrightarrow(a)\)} in
\(\mathrm{EGA}\ \mathrm{IV}_4\), 17.12.1 \textup{(applied to
\(f\colon N\to S\) and \(j=\varepsilon\))}, the morphism
\(N_{S'}\to S'\) is smooth, of relative dimension \(r\), at every
point of \(\varepsilon(S')\) if and only if
\[
  \mathfrak n_{S'/N'}
  =
  \mathfrak n_{S/N}\otimes_{\mathcal O_S}\mathcal O_{S'}
\]
is locally free of rank \(r\), and the morphism
\[
  \varphi_n(S')\colon
  \operatorname{Sym}^n(\mathfrak n_{S'/N'})
  \longrightarrow
  \mathcal J_{S'}^n/\mathcal J_{S'}^{n+1}
\]
is an isomorphism for every \(n\geq1\). Put
\[
  K_n(S')=\ker\varphi_n(S').
\]
By \(\mathrm{TDTE}\) I, Lemma 3.6,
\(\mathfrak n_{S/N}\otimes\mathcal O_{S'}\) is locally free of
rank \(r\) if and only if \(S'\to S\) factors through a certain
subscheme \(Z\) of \(S\). Replacing \(S\) by \(Z\), we may
therefore suppose that
\(\mathfrak n_{S/N}=\mathcal J/\mathcal J^2\) is locally free of
rank \(r\). Then, for every \(S'\to S\),
\[
  (\mathcal J^2/\mathcal J^3)\otimes\mathcal O_{S'}
  =
  \mathcal J_{S'}^2/\mathcal J_{S'}^3,
\]
and consequently
\[
  K_2(S')=K_2(S)\otimes\mathcal O_{S'}.
\]
It follows that \(\varphi_2(S')\) is an isomorphism if and only if
\(S'\to S\) factors through the closed subscheme \(S_2\) of \(S\)
defined by the ideal generated by the image of
\[
  \operatorname{Sym}^2(\mathfrak n_{S/N})^*
  \otimes K_2(S)
  \longrightarrow
  \mathcal O_S.
\]
Over \(S_2\), the module \(\mathcal J^2/\mathcal J^3\) is then
isomorphic to \(\operatorname{Sym}^2(\mathfrak n_{S/N})\), hence
is locally free; the same argument shows that
\(\varphi_3(S')\) is an isomorphism if and only if \(S'\to S\)
factors through a certain closed subscheme \(S_3\) of \(S\), and
so on. We thus obtain that \(N_{S'}\to S'\) is smooth, of relative
dimension \(r\), at every point of \(\varepsilon(S')\) if and only
if \(S'\to S\) factors through the closed subscheme
\[
  Z=\bigcap_n S_n.
\]
But then, for every \(S'\to Z\),
\(\operatorname{Lie}(N_{S'}/S')\) is locally a direct summand of
rank \(r\) of \(\mathfrak g_{S'}\), and hence the inclusion
\[
  \mathfrak h_{S'}
  \subset
  \operatorname{Lie}(N_{S'}/S')
\]
is an equality. We then put \(H=N^0\).

Replacing \(S\) by \(Z\), it remains only to express the condition
that \(H_{s'}\) have the same reductive rank as \(G_{s'}\) at every
point \(s'\in S'\), or, equivalently, that \(H_s\) have the same
reductive rank as \(G_s\) at every point \(s\) in the
set-theoretic image of \(S'\) in \(S\). This condition defines an
open subset of \(S\) \textup{(\SGARefLink{sga3:xix:remark:6-2}{Exposé XIX, 6.2})}.
\hfill\(\square\)

\SGARefTarget{sga3:xxii:rem:after-5.8.1}
\paragraph{Remark.}
In general, the scheme \(\mathcal H\) is not smooth over \(S\). It
is nevertheless smooth if \(6\) is invertible on \(S\), or if
there exists a prime number \(p\) such that \(p\cdot1_S=0\)
\textup{(that is, if \(S\) has characteristic \(p>0\))}.

\SGARefTarget{sga3:xxii:cor:5.8.2}
\paragraph{\SGARefLink{sga3:xxii:cor:5.8.2}{Corollary 5.8.2}.}
\textit{Let \(S\) be a scheme, let \(G\) be a reductive
\(S\)-group, and let \(H\) be a subgroup of type \textup{(R)} of
\(G\). Recall from \SGARefLink{sga3:xxii:cor:5.3.10}{5.3.10} that
\(\underline{\operatorname{Norm}}_G(H)\) is representable by a
closed subgroup scheme of \(G\), smooth over \(S\). Then the
quotient sheaf
\[
  G/\underline{\operatorname{Norm}}_G(H)
\]
is representable by a quasi-projective \(S\)-scheme, smooth and of
finite presentation over \(S\) \textup{(in fact, by an open
subscheme of \(\mathcal H\))}.}

Indeed, consider the morphism
\[
  f\colon G\longrightarrow\mathcal H
\]
defined set-theoretically by
\[
  f(g)=\operatorname{int}(g)H.
\]
By \SGARefLink{sga3:xxii:thm:5.3.9}{5.3.9}, this morphism is smooth and of finite presentation, hence
open. Let \(V=f(G)\), equipped with its open-subscheme structure in
\(\mathcal H\). The morphism \(f\colon G\to V\) is a covering
morphism with kernel \(\underline{\operatorname{Norm}}_G(H)\).
Consequently
\(G/\underline{\operatorname{Norm}}_G(H)\) is representable by
\(V\) \textup{(cf.\ \SGARefLink{sga3:iv:lemma:4:6:5}{Exposé IV, 4.6.5})}.
\hfill\(\square\)

\SGARefTarget{sga3:xxii:cor:5.8.3}
\paragraph{\SGARefLink{sga3:xxii:cor:5.8.3}{Corollary 5.8.3}.}
{\itshape Let \(S\) be a scheme and let \(G\) be a reductive
\(S\)-group. Consider the functors
\(\underline{\operatorname{Tor}}(G)\),
\(\underline{\operatorname{Bor}}(G)\), and
\(\underline{\operatorname{Kil}}(G)\), defined by
\[
\begin{aligned}
  \underline{\operatorname{Tor}}(G)(S')
  &=
  \{\text{maximal tori of \(G_{S'}\)}\},\\
  \underline{\operatorname{Bor}}(G)(S')
  &=
  \{\text{Borel subgroups of \(G_{S'}\)}\},\\
  \underline{\operatorname{Kil}}(G)(S')
  &=
  \{\text{Killing pairs of \(G_{S'}\)
  \textup{ (cf.\ 5.3.13)}}\}.
\end{aligned}
\]
\textup{\SGARefLink{sga3:xxii:cor:5.7.6:item:i}{(i)}}
\(\underline{\operatorname{Tor}}(G)\),
\(\underline{\operatorname{Bor}}(G)\), and
\(\underline{\operatorname{Kil}}(G)\) are representable by smooth
\(S\)-schemes of finite presentation, with geometrically integral
fibres, and are respectively affine, projective, and affine over
\(S\).

\textup{\SGARefLink{sga3:xxii:cor:5.7.6:item:ii}{(ii)}}
The canonical morphism
\[
  \underline{\operatorname{Kil}}(G)
  \longrightarrow
  \underline{\operatorname{Tor}}(G)
\]
\textup{(respectively,
\(\underline{\operatorname{Kil}}(G)\to
\underline{\operatorname{Bor}}(G)\))}
is finite étale and surjective
\textup{(respectively, affine, smooth, and surjective)}.

\SGAInlineRefTarget{sga3:xxii:cor:5.8.3:item:iii}\textup{\SGARefLink{sga3:xxii:prop:5.6.9:item:iii}{(iii)}}
Let \(T\) be a maximal torus of \(G\)
\textup{(respectively, let \(B\) be a Borel subgroup of \(G\), or
let \(B\supset T\) be a Killing pair of \(G\))}. The morphism
\[
\begin{aligned}
  G&\longrightarrow\underline{\operatorname{Tor}}(G),
  &\text{resp.}\quad
  G&\longrightarrow\underline{\operatorname{Bor}}(G),
  &\text{resp.}\quad
  G&\longrightarrow\underline{\operatorname{Kil}}(G),
\end{aligned}
\]
defined by
\[
\begin{aligned}
  g&\longmapsto\operatorname{int}(g)T,
  &\text{respectively}\quad
  g&\longmapsto\operatorname{int}(g)B,
  &\text{respectively}\quad
  g&\longmapsto
  \bigl(\operatorname{int}(g)B,\operatorname{int}(g)T\bigr)
\end{aligned}
\]}

{\itshape induces an isomorphism}
\[
\begin{aligned}
  G/\underline{\operatorname{Norm}}_G(T)
  &\xrightarrow{\;\sim\;}
  \underline{\operatorname{Tor}}(G),
  &\text{resp.}\quad
  G/B
  &\xrightarrow{\;\sim\;}
  \underline{\operatorname{Bor}}(G),
  &\text{resp.}\quad
  G/T
  &\xrightarrow{\;\sim\;}
  \underline{\operatorname{Kil}}(G).
\end{aligned}
\]

One first sees that \textup{\SGARefLink{sga3:xxii:prop:5.6.9:item:iii}{(iii)}} follows from the conjugacy theorem
for maximal tori \textup{(respectively, Borel subgroups, respectively,
Killing pairs)} and from the facts that
\[
  \underline{\operatorname{Norm}}_G(B)=B,
  \qquad
  \underline{\operatorname{Norm}}_G(B)
  \cap
  \underline{\operatorname{Norm}}_G(T)=T,
\]
all of which were established above \textup{(\SGARefLink{sga3:xxii:rem:5.1.2}{5.1.2}, \SGARefLink{sga3:xxii:cor:5.3.12}{5.3.12}, \SGARefLink{sga3:xxii:cor:5.3.14}{5.3.14},
\SGARefLink{sga3:xxii:prop:5.6.1}{5.6.1})}.

It follows first that the canonical morphisms
\[
  \underline{\operatorname{Tor}}(G)\longrightarrow\mathcal H,
  \qquad
  \underline{\operatorname{Bor}}(G)\longrightarrow\mathcal H
\]
are representable, locally for the étale topology, by open immersions
\textup{(\SGARefLink{sga3:xxii:cor:5.8.2}{5.8.2} and \SGARefLink{sga3:xxii:rem:5.1.2}{5.1.2}, respectively \SGARefLink{sga3:xxii:cor:5.5.5}{5.5.5})}. Hence, by descent,
\(\underline{\operatorname{Tor}}(G)\) and
\(\underline{\operatorname{Bor}}(G)\) are representable by open
subschemes of \(\mathcal H\). Likewise,
\(\underline{\operatorname{Kil}}(G)\) is locally, for the étale
topology, representable by a scheme affine over the base
\textup{(\SGARefLink{sga3:ix:proposition:2:3}{Exposé IX, 2.3})}, and hence is representable by an affine
\(S\)-scheme, by descent for affine schemes.\footnote{Editor's note:
Cf.\ SGA 1, VIII, 2.1.}

The assertions of \textup{\SGARefLink{sga3:xxii:cor:5.7.6:item:ii}{(ii)}} follow at once from \SGARefLink{sga3:xxii:cor:5.5.5}{5.5.5}\textup{\SGARefLink{sga3:xxii:cor:5.5.5:item:ii}{(ii)}}
and \SGARefLink{sga3:xxii:cor:5.6.13}{5.6.13}. It already follows that
\(\underline{\operatorname{Tor}}(G)\) is affine over \(S\)
\textup{(EGA II, 6.7.1)}. It remains only to prove that
\(\underline{\operatorname{Bor}}(G)\) is projective over \(S\). We
already know that it is quasi-projective, so it remains to prove that
it is proper;\footnote{Editor's note: We may suppose that \(S\) is
affine and, since \(\underline{\operatorname{Bor}}(G)\) is of finite
presentation over \(S\) by \textup{(i)}, reduce to the case in which
\(S\) is Noetherian; we are then under the hypotheses of EGA
IV\(_3\), 15.7.10.} but it has connected fibres, and therefore, by EGA
IV\(_3\), 15.7.10, one is reduced to proving this on geometric fibres.
If \(S\) is the spectrum of an algebraically closed field, then
\(\underline{\operatorname{Bor}}(G)=G/B\) by \textup{\SGARefLink{sga3:xxii:prop:5.6.9:item:iii}{(iii)}}, and the
conclusion follows from \emph{Bible}, \S6.4, Theorem 4
\textup{(or [Ch05], \S6.5, Theorem 5)}.
\hfill\(\square\)

\SGARefTarget{sga3:xxii:rem:5.8.4}
\paragraph{\SGARefLink{sga3:xxii:rem:5.8.4}{Remark 5.8.4}.}
Under the hypotheses of \SGARefLink{sga3:xxii:cor:5.8.3}{5.8.3}, let \(Q\) be a central subgroup of
multiplicative type of \(G\). The evident morphisms define
isomorphisms
\[
  \underline{\operatorname{Tor}}(G)
  \simeq
  \underline{\operatorname{Tor}}(G/Q),
  \qquad
  \underline{\operatorname{Bor}}(G)
  \simeq
  \underline{\operatorname{Bor}}(G/Q),
  \qquad
  \underline{\operatorname{Kil}}(G)
  \simeq
  \underline{\operatorname{Kil}}(G/Q).
\]

\SGARefTarget{sga3:xxii:cor:5.8.5}
\paragraph{\SGARefLink{sga3:xxii:cor:5.8.5}{Corollary 5.8.5}.}
{\itshape Let \(S\) be a scheme, \(G\) a reductive \(S\)-group, and
\(P\) a subgroup scheme of \(G\), smooth and of finite presentation
over \(S\). The following conditions are equivalent:

\textup{\SGARefLink{sga3:xxii:cor:5.7.6:item:i}{(i)}} For every \(s\in S\), \(P_{\bar s}\) is a parabolic
subgroup of \(G_{\bar s}\) \textup{(that is, the quotient scheme
\(G_{\bar s}/P_{\bar s}\) is proper over \(\bar s\), or equivalently
\(P_{\bar s}\) contains a Borel subgroup of \(G_{\bar s}\); cf.\
\emph{Bible}, \S6.4, Theorem 4, or [Ch05], \S6.5, Theorem 5)}.

\textup{\SGARefLink{sga3:xxii:cor:5.7.6:item:ii}{(ii)}} The quotient sheaf \(G/P\) is representable by an
\(S\)-scheme smooth and projective over \(S\).

Moreover, under these conditions, \(P\) is closed in \(G\), has
connected fibres, and
\[
  P=\underline{\operatorname{Norm}}_G(P).
\]}

Clearly \textup{\SGARefLink{sga3:xxii:cor:5.7.6:item:ii}{(ii)}} implies \textup{\SGARefLink{sga3:xxii:cor:5.7.6:item:i}{(i)}}. If \textup{\SGARefLink{sga3:xxii:cor:5.7.6:item:i}{(i)}} holds,
then
\[
  P(\bar s)
  =
  \underline{\operatorname{Norm}}_{G(\bar s)}(P_{\bar s})
\]
and \(P_{\bar s}\) is connected \textup{(for the first assertion,
cf.\ \emph{Bible}, \S12.3, \SGARefLink{sga3:xxii:section:4}{Lemma 4};\footnote{Editor's note: We have
supplied the details that follow.} the second assertion}
\textup{follows from this, since
\(P'=P_{\bar s}^{0}\) is a parabolic subgroup of \(G_{\bar s}\)
normalized by \(P_{\bar s}\), whence
\(P'(\bar s)=P(\bar s)\) and therefore \(P'=P_{\bar s}\))}. It
follows that \(P\) is of type \textup{(R)} and that
\[
  P=\underline{\operatorname{Norm}}_G(P),
\]
so that \(P\) is closed in \(G\). By \SGARefLink{sga3:xxii:cor:5.8.2}{5.8.2},
\[
  G/P=G/\underline{\operatorname{Norm}}_G(P)
\]
is representable by a quasi-projective \(S\)-scheme. Its fibres are
connected and proper, and it is therefore projective by the argument
of \SGARefLink{sga3:xxii:cor:5.8.3}{5.8.3}.
\hfill\(\square\)

\SGARefTarget{sga3:xxii:rem:5.8.6}
\paragraph{\SGARefLink{sga3:xxii:rem:5.8.6}{Remark 5.8.6}.}
The statements \SGARefLink{sga3:xxii:thm:5.8.1}{5.8.1}, \SGARefLink{sga3:xxii:cor:5.8.2}{5.8.2}, and \SGARefLink{sga3:xxii:cor:5.8.5}{5.8.5} are valid for an \(S\)-group
of type \textup{(RA)}, or for an \(S\)-group of type \textup{(RR)}
satisfying \SGARefLink{sga3:xxii:prop:5.1.8}{5.1.8}.\footnote{Editor's note: Indeed, the proof of \SGARefLink{sga3:xxii:thm:5.8.1}{5.8.1}
uses only \SGARefLink{sga3:xxii:cor:5.3.3}{5.3.3} \textup{(which is valid for a group of type
\textup{(RA)})} and \SGARefLink{sga3:xix:remark:6:2}{Exposé XIX, 6.2}, which, by \SGARefLink{sga3:xii:1.7:item:b}{Exposé XII, 1.7(b)}, is
also valid for groups of type \textup{(RR)}.}

\SGARefTarget{sga3:xxii:rem:5.8.7}
\paragraph{\SGARefLink{sga3:xxii:rem:5.8.7}{Remark 5.8.7}.}
Through the inner automorphisms of \(G\), there are canonical actions
\[
\begin{aligned}
  G&\longrightarrow
  \underline{\operatorname{Aut}}_S
  \bigl(\underline{\operatorname{Tor}}(G)\bigr),
  &
  G&\longrightarrow
  \underline{\operatorname{Aut}}_S
  \bigl(\underline{\operatorname{Bor}}(G)\bigr),
  &
  G&\longrightarrow
  \underline{\operatorname{Aut}}_S
  \bigl(\underline{\operatorname{Kil}}(G)\bigr),
\end{aligned}
\]
which, in the situation of \SGARefLink{sga3:xxii:cor:5.8.3}{5.8.3}\textup{\SGARefLink{sga3:xxii:prop:5.6.9:item:iii}{(iii)}}, are identified with
the canonical actions
\[
\begin{aligned}
  G&\longrightarrow
  \underline{\operatorname{Aut}}_S
  \bigl(G/\underline{\operatorname{Norm}}_G(T)\bigr),
  &
  G&\longrightarrow
  \underline{\operatorname{Aut}}_S(G/B),
  &
  G&\longrightarrow
  \underline{\operatorname{Aut}}_S(G/T).
\end{aligned}
\]
In particular, one obtains
\[
\begin{aligned}
  \ker\!\left(
    G\longrightarrow
    \underline{\operatorname{Aut}}_S
    \bigl(\underline{\operatorname{Tor}}(G)\bigr)
  \right)
  &=
  \ker\!\left(
    G\longrightarrow
    \underline{\operatorname{Aut}}_S
    \bigl(\underline{\operatorname{Bor}}(G)\bigr)
  \right)\\
  &=
  \ker\!\left(
    G\longrightarrow
    \underline{\operatorname{Aut}}_S
    \bigl(\underline{\operatorname{Kil}}(G)\bigr)
  \right)\\
  &=
  \underline{\operatorname{Centr}}(G).
\end{aligned}
\]
Indeed, it is clear that \(\underline{\operatorname{Centr}}(G)\)
acts trivially on each of the three schemes. Conversely, the kernel
of
\[
  G\longrightarrow
  \underline{\operatorname{Aut}}_S
  \bigl(\underline{\operatorname{Kil}}(G)\bigr)
\]
is ``the intersection of the maximal tori of \(G\)'' in the sense of
\SGARefLink{sga3:xxii:cor:4.1.7}{4.1.7}, and hence equals
\(\underline{\operatorname{Centr}}(G)\) \textup{(loc.\ cit.)}. For
\(\underline{\operatorname{Bor}}(G)\), one observes that ``the
intersection of the Borel subgroups of \(G\)'' is also ``the
intersection of its maximal tori'' \textup{(see the following
section)}. For \(\underline{\operatorname{Tor}}(G)\), one uses
\SGARefLink{sga3:xii:4.11}{Exposé XII, 4.11}.

\SGARefTarget{sga3:xxii:section:5.9}
\section*{5.9. Properties Peculiar to Borel Subgroups}
\addcontentsline{toc}{section}{\SGARefLink{sga3:xxii:section:5.9}{5.9}. Properties Peculiar to Borel
Subgroups}

Most of these properties will be generalized in Exposé XXVI to
parabolic subgroups.

\SGARefTarget{sga3:xxii:def:5.9.1}
\paragraph{\SGARefLink{sga3:xxii:def:5.9.1}{Definition 5.9.1}.}
Let \(S\) be a scheme, \(G\) a reductive \(S\)-group, and \(B,B'\)
two Borel subgroups of \(G\). We say that \(B\) and \(B'\) are
\emph{in general position} \textup{(or that \(B'\) is in general
position relative to \(B\))} if \(B\cap B'\) is a torus
\textup{(necessarily maximal)} of \(G\).

If \(T\) is a maximal torus of \(G\) contained in \(B\) and \(B'\),
we say that \(B\) and \(B'\) are \emph{opposite} \textup{(relative
to \(T\))} if
\[
  B\cap B'=T.
\]

\SGARefTarget{sga3:xxii:prop:5.9.2}
\paragraph{\SGARefLink{sga3:xxii:prop:5.9.2}{Proposition 5.9.2}.}
{\itshape Let \(S\) be a scheme, \(G\) a reductive \(S\)-group, \(B\)
a Borel subgroup of \(G\), and \(T\) a maximal torus of \(B\). There
exists a unique Borel subgroup \(B'\) of \(G\), opposite to \(B\)
relative to \(T\).

If \((G,T,M,R)\) is a splitting of \(G\) relative to \(T\), and if
\(B=B_{R_+}\) \textup{(\SGARefLink{sga3:xxii:prop:5.5.1}{5.5.1})}, then
\[
  B'=B_{-R_+}.
\]}

By faithfully flat descent, it is enough to prove the proposition in
the split case, with \(B=B_{R_+}\) \textup{(\SGARefLink{sga3:xxii:cor:5.5.5:item:iv}{5.5.5(iv)})}. Then
\(B_{-R_+}\) is indeed opposite to \(B\) \textup{(\SGARefLink{sga3:xxii:prop:4.1.2}{4.1.2})}. Let us
show that it is the only Borel subgroup of \(G\) containing \(T\)
that is opposite to \(B\). If \(B'\) is a Borel subgroup of \(G\)
containing \(T\), then locally on \(S\) it has the form
\(B_{R'_+}\), where \(R'_+\) is another system of positive roots of
\(R\) \textup{(\SGARefLink{sga3:xxii:cor:5.5.5:item:iii}{5.5.5(iii)})}. If \(R'_+\ne -R_+\), there exists
\(\alpha\in R'_+\cap R_+\), and hence
\[
  U_\alpha\subset B_{R_+}\cap B_{R'_+}.
\]
\hfill\(\square\)

\SGARefTarget{sga3:xxii:prop:5.9.3}
\paragraph{\SGARefLink{sga3:xxii:prop:5.9.3}{Proposition 5.9.3}.}
{\itshape Let \(S\) be a scheme, \(G\) a reductive \(S\)-group, and
\(B\) a Borel subgroup of \(G\).

\textup{\SGARefLink{sga3:xxii:cor:5.7.6:item:i}{(i)}} If \(B'\) is a Borel subgroup of \(G\), the following
conditions are equivalent:

\begin{enumerate}
\item[\textup{(a)}] \(B'\) is in general position relative to \(B\)
\textup{(\SGARefLink{sga3:xxii:def:5.9.1}{5.9.1})}.

\item[\textup{\SGAInlineRefTarget{sga3:xxii:prop:5.9.3:item:b}(b)}] \(B'^{u}\cap B^{u}=e\).

\item[\textup{(b\(^{\prime}\))}] \(B'^{u}\cap B=e\).

\item[\textup{\SGAInlineRefTarget{sga3:xxii:prop:5.9.3:item:c}(c)}] Multiplication in \(G\) induces an open immersion
\[
  B'^{u}\times_S B\longrightarrow G.
\]

\item[\textup{(c\(^{\prime}\))}] The canonical morphism
\[
  B'^{u}\longrightarrow G/B
\]
is an open immersion.
\end{enumerate}

\textup{\SGARefLink{sga3:xxii:cor:5.7.6:item:ii}{(ii)}} The functor \(\underline{\operatorname{Opp}}(B)\),
\[
  S'\longmapsto
  \left\{
    \begin{array}{c}
      \text{Borel subgroups of \(G_{S'}\) in general position}\\
      \text{relative to \(B_{S'}\)}
    \end{array}
  \right\},
\]
is representable by an open subscheme of
\(\underline{\operatorname{Bor}}(G)\) \textup{(\SGARefLink{sga3:xxii:cor:5.8.3}{5.8.3})}. The morphism
\[
  \underline{\operatorname{Opp}}(B)
  \longrightarrow
  \underline{\operatorname{Tor}}(B)
\]
defined by \(B'\mapsto B\cap B'\) is an isomorphism. In particular
\textup{(\SGARefLink{sga3:xxii:cor:5.6.13}{5.6.13})}, the inner automorphisms of \(B^u\) equip
\(\underline{\operatorname{Opp}}(B)\) with the structure of a
principal homogeneous bundle under \(B^u\).}

Let us first examine \textup{\SGARefLink{sga3:xxii:cor:5.7.6:item:i}{(i)}}. We have
\(\textup{(a)}\Rightarrow\textup{(c)}\). Indeed,
\textup{\SGARefLink{sga3:xxii:prop:5.9.3:item:c}{(c)}} is local for the étale topology; by \SGARefLink{sga3:xxii:cor:5.5.5:item:iv}{5.5.5(iv)}, one
reduces to the case in which \(G\) is split relative to
\(B\cap B'\) and \(B\) has the form \(B_{R_+}\). By \SGARefLink{sga3:xxii:prop:5.9.2}{5.9.2},
\[
  B'^{u}=U_{-R_+},
\]
and one is reduced to \SGARefLink{sga3:xxii:prop:4.1.2}{4.1.2}.

Trivially,
\[
  \textup{(c\(^{\prime}\))}
  \Longleftrightarrow
  \textup{(c)}
  \Longrightarrow
  \textup{(b\(^{\prime}\))}
  \Longrightarrow
  \textup{(b)}.
\]
It remains to prove
\(\textup{(b)}\Rightarrow\textup{(a)}\). Let us first prove
\textup{\SGARefLink{sga3:xxii:cor:5.7.6:item:ii}{(ii)}}. The second assertion is a formal consequence of
\SGARefLink{sga3:xxii:prop:5.9.2}{5.9.2}, and the third follows at once from it by \SGARefLink{sga3:xxii:cor:5.6.13}{5.6.13}. Let us prove
the first. It is local for the étale topology, and we may therefore
suppose that \(B\) contains a maximal torus \(T\). Let \(B'_0\) be
the Borel subgroup opposite to \(B\) relative to \(T\)
\textup{(\SGARefLink{sga3:xxii:prop:5.9.2}{5.9.2})}.

By what precedes, the morphism
\[
  B^u\longrightarrow\underline{\operatorname{Bor}}(G)
\]
induced by the canonical morphism
\[
  G\longrightarrow G/B'_0
  \longrightarrow\underline{\operatorname{Bor}}(G)
\]
\textup{(\SGARefLink{sga3:xxii:cor:5.8.3}{5.8.3})} induces an isomorphism
\[
  B^u\xrightarrow{\;\sim\;}
  \underline{\operatorname{Opp}}(B).
\]
Thus there is a commutative diagram
\SGARefTarget{sga3:xxii:diag:5.9.3-opp-square}
\[
  % BEGIN INCLUDED SGA3_English_current_buildable_source_20260823/source/inputs/XXII/native_diagrams/exp22/diag_007_opposite_borel_square.tex
\begin{tikzcd}[column sep=5.0em,row sep=3.3em]
  B^u
    \arrow[r]
    \arrow[d,"\sim"]
  &
  G/B'_0
    \arrow[d,"\sim"]
  \\
  \underline{\operatorname{Opp}}(B)
    \arrow[r]
  &
  \underline{\operatorname{Bor}}(G).
\end{tikzcd}

% END INCLUDED SGA3_English_current_buildable_source_20260823/source/inputs/XXII/native_diagrams/exp22/diag_007_opposite_borel_square.tex

\]
Now the morphism \(B^u\to G/B'_0\) is an open immersion
\textup{(by \textup{\SGARefLink{sga3:xxii:cor:5.7.6:item:i}{(i)}}, \(\textup{(a)}\Rightarrow
\textup{(c\(^{\prime}\))}\))}, which completes the proof of
\textup{\SGARefLink{sga3:xxii:cor:5.7.6:item:ii}{(ii)}}. We record at once the following corollary.

\SGARefTarget{sga3:xxii:cor:5.9.4}
\paragraph{\SGARefLink{sga3:xxii:cor:5.9.4}{Corollary 5.9.4}.}
{\itshape Let \(G\) be a reductive \(S\)-group and let \(B,B'\) be
two Borel subgroups of \(G\). If \(s\in S\) is such that
\(B_{\bar s}\) and \(B'_{\bar s}\) are in general position, there
exists an open subset \(V\) of \(S\), containing \(s\), such that
\(B_V\) and \(B'_V\) are in general position.}

It therefore remains only to prove
\(\textup{(b)}\Rightarrow\textup{(a)}\). By the preceding corollary,
it is enough to do so when \(S\) is the spectrum of an algebraically
closed field \(k\). We may suppose that \(G\) is split relative to a
maximal torus \(T\) of \(B\). Let \(B'_0\) be the Borel subgroup
opposite to \(B\). Since the Borel subgroups of \(G\) are conjugate
under \(G(k)\), there exists \(g\in G(k)\) such that
\(\operatorname{int}(g)B'_0=B'\). By the Bruhat theorem
\textup{(\SGARefLink{sga3:xxii:thm:5.7.4}{5.7.4})}, we may write
\[
  g=bnb',
  \qquad
  b\in B(k),\quad
  b'\in B'_0(k),\quad
  n\in\underline{\operatorname{Norm}}_G(T)(k).
\]
Thus
\[
  B'=\operatorname{int}(b)\operatorname{int}(n)B'_0
\]
and
\[
  B'\cap B
  =
  \operatorname{int}(b)
  \bigl(\operatorname{int}(n)B'_0\cap B\bigr).
\]
If \(n\notin T(k)\), then
\[
  \operatorname{int}(n)(B'_0)^u\cap B^u\ne e
\]
\textup{(cf. the proof of \SGARefLink{sga3:xxii:prop:5.9.2}{5.9.2})}; it follows that
\textup{\SGARefLink{sga3:xxii:prop:5.9.3:item:b}{(b)}} implies
\[
  B'\cap B
  =
  \operatorname{int}(b)(B'_0\cap B)
  =
  \operatorname{int}(b)T.
\]
\hfill\(\square\)

\SGARefTarget{sga3:xxii:prop:5.9.5}
\paragraph{\SGARefLink{sga3:xxii:prop:5.9.5}{Proposition 5.9.5}.}
{\itshape Let \(S\) be a scheme, \(G\) a reductive \(S\)-group,
\(B\) a Borel subgroup of \(G\), and \(B^u\) its unipotent part.
There exists a sequence of subgroups of \(B\),
\[
  U_0=B^u\supset U_1\supset\cdots\supset U_n\supset\cdots,
\]
having the following properties:
\begin{enumerate}
\item[\textup{(i)}] Each \(U_i\) is smooth, has connected fibers,
and is characteristic in \(B\); the inner automorphisms of \(B^u\)
act trivially on the \textup{(sheaf)} quotients \(U_i/U_{i+1}\).

\item[\textup{(ii)}] For each \(i\geq 0\), there exists a locally
free \(\mathcal O_S\)-module \(\mathcal E_i\) and an isomorphism of
sheaves of groups on \(S\),
\[
  U_i/U_{i+1}\xrightarrow{\;\sim\;}W(\mathcal E_i).
\]

\item[\textup{\SGAInlineRefTarget{sga3:xxii:prop:5.9.5:item:iii}(iii)}] For every \(s\in S\),
\[
  (U_n)_s=e
  \qquad\text{for}\qquad
  n\geq\dim(B^u_s).
\]
\end{enumerate}}

Suppose first that there exists a splitting \((G,T,M,R)\) of \(G\)
and a system \(R_+\) of positive roots of \(R\) such that
\(B=B_{R_+}\). Denote by \(\Delta\) the set of simple roots of
\(R_+\). For each \(\alpha\in R_+\), denote by
\(\operatorname{ord}(\alpha)\) the sum of the coefficients of
\(\alpha\) in the basis \(\Delta\) of \(\Gamma_0(R)\); this is the
\emph{order} of \(\alpha\) relative to \(R_+\). We have
\[
  \operatorname{ord}(\alpha)\leq\operatorname{Card}(R_+).
\]
For every \(i>0\), let \(R^{(i)}\) be the set of roots of order
\(>i\). This is a closed set of positive roots, and we may therefore
construct \textup{(\SGARefLink{sga3:xxii:cor:5.6.5}{5.6.5})}
\[
  U_i=U_{R^{(i)}}.
\]
If \(\alpha\in R_+\) and \(\beta\in R^{(i)}\), then
\(\alpha+\beta\in R^{(i+1)}\). It follows from \SGARefLink{sga3:xxii:cor:5.5.2}{5.5.2} that each
\(U_i\) is an invariant subgroup of \(B\), and that the inner
automorphisms of \(B^u\) act trivially on \(U_i/U_{i+1}\). Moreover,
this group identifies with
\[
  \prod_{\operatorname{ord}(\alpha)=i+1}U_\alpha
\]
and is therefore endowed with a vector-group structure.

If \(B\) is of the form \(B_{R'_+}\) for another splitting
\((G,T',M',R')\) of \(G\), let us show that the groups \(U'_i\)
constructed as above using the new splitting coincide with the
\(U_i\), and that the vector-group structures on the successive
quotients likewise coincide. By \SGARefLink{sga3:xxii:cor:5.6.13}{5.6.13}, there exists
\(b\in B^u(S)\) such that \(T'=\operatorname{int}(b)T\); the
assertion to be proved is local on \(S\), and we may therefore
suppose that the isomorphism
\[
  T\xrightarrow{\;\sim\;}T'
\]
induced by \(\operatorname{int}(b)\) arises by duality from an
isomorphism of root data
\[
  h:(M',M'{}^*,R',R'{}^*)
  \xrightarrow{\;\sim\;}
  (M,M^*,R,R^*).
\]
The roots in \(R'_+\) are then the roots
\[
  \alpha\circ\operatorname{int}(b)=h(\alpha),
  \qquad \alpha\in R_+,
\]
and its simple roots are the \(h(\alpha)\), with
\(\alpha\in\Delta\). Consequently
\[
  \operatorname{ord}\bigl(h(\alpha)\bigr)
  =
  \operatorname{ord}(\alpha).
\]
By transport of structure, the vector groups
\(U'_{h(\alpha)}\) are the groups
\(\operatorname{int}(b)U_\alpha\). It follows that
\[
  \operatorname{int}(b)U_i=U'_i.
\]
But \(U_i\) is invariant in \(B\), so that \(U_i=U'_i\).
Likewise, the isomorphism
\[
  \operatorname{int}(b):
  U_i/U_{i+1}\xrightarrow{\;\sim\;}U'_i/U'_{i+1}
\]
is the identity.

We now pass to the general case. Let \(\{S_i\to S\}\) be an étale
covering family such that each
\((G_i,T_i,M_i,R_i)\) is a splitting, and choose a system
\(R_{i+}\) of positive roots such that
\[
  B\times_S S_i=B_{R_{i+}}
\]
\textup{(\SGARefLink{sga3:xxii:cor:5.5.5}{5.5.5}\textup{\SGARefLink{sga3:xxii:cor:5.5.5:item:iii}{(iii)}})}. For every \(i\), the preceding
construction gives a sequence
\[
  B_{S_i}^u=U_{i,0}
  \supset U_{i,1}\supset\cdots
  \supset U_{i,j}\supset\cdots
\]
and vector-group structures on its successive quotients.
By descent, it is enough to verify, for every \(i,i'\) and every
\(j\), that
\[
  U_{i,j}\times_{S_i}S_{ii'}
  =
  U_{i',j}\times_{S_{i'}}S_{ii'},
  \qquad
  S_{ii'}=S_i\times_S S_{i'},
\]
and that the vector-group structures on
\[
  (U_{i,j}/U_{i,j+1})\times_{S_i}S_{ii'}
  \quad\text{and}\quad
  (U_{i',j}/U_{i',j+1})\times_{S_{i'}}S_{ii'}
\]
coincide.\begingroup
\renewcommand{\thefootnote}{*}\footnote{Source note: in the second
base change the French source prints \(S_i\); the preceding equality
and \(S_{ii'}=S_i\times_S S_{i'}\) require \(S_{i'}\).}
\addtocounter{footnote}{-1}\endgroup
If \(S_{ii'}=\varnothing\), there is nothing to prove. Otherwise we
are in the situation considered above: \(B\times_S S_{ii'}\) is
defined by the system of positive roots \(R_{i+}\), respectively
\(R_{i'+}\), in the splittings
\[
  (G_{S_{ii'}},
   T_i\times_{S_i}S_{ii'},M_i,R_i),
  \qquad
  (G_{S_{ii'}},
   T_{i'}\times_{S_{i'}}S_{ii'},M_{i'},R_{i'}).
\]
This proves the proposition.
\hfill\(\square\)

\SGARefTarget{sga3:xxii:cor:5.9.6}
\paragraph{\SGARefLink{sga3:xxii:cor:5.9.6}{Corollary 5.9.6}.}
{\itshape If \(S\) is affine, then
\[
  H^1(S,B^u)=e;
\]
that is, every principal homogeneous bundle under \(B^u\) has a
section.}

Indeed, \(S\) is a disjoint union of subschemes on each of which
\(B^u\) has constant relative dimension. In view of
\SGARefLink{sga3:xxii:prop:5.9.5}{5.9.5}\textup{\SGARefLink{sga3:xxii:prop:5.9.5:item:iii}{(iii)}}, we may therefore suppose that \(U_n=e\).
By TDTE I, B.1.1,\setcounter{footnote}{55}\footnote{Editor's note:
This is the corollary on
page 18 of TDTE I.} one has
\[
  H^1(S,U_i/U_{i+1})
  =
  H^1\bigl(S,W(\mathcal E_i)\bigr)
  =
  0.
\]
It follows by induction that \(H^1(S,B^u)=e\).

\SGARefTarget{sga3:xxii:cor:5.9.7}
\paragraph{\SGARefLink{sga3:xxii:cor:5.9.7}{Corollary 5.9.7}.}
{\itshape If \(S\) is affine, then \(B\) has maximal tori. If \(T\)
is a maximal torus of \(B\), then
\[
  H^1(S,T)=H^1(S,B).
\]}

The first assertion follows at once from \SGARefLink{sga3:xxii:cor:5.9.6}{5.9.6} and \SGARefLink{sga3:xxii:cor:5.6.13}{5.6.13}; the
second follows from it in the standard way.\footnote{Editor's note:
Indeed, there is an exact sequence
\[
  H^1(S,B^u)\longrightarrow H^1(S,B)
  \xrightarrow{\;\pi\;}H^1(S,B/B^u)=H^1(S,T);
\]
see \textup{[Se64]}, I \S5.5, Proposition 38, or
\textup{[Gi71]}, III Proposition 3.3.1. Now
\(H^1(S,T)\to H^1(S,B)\) is a section of \(\pi\), so that
\(\pi\) is surjective; on the other hand,
\(H^1(S,B^u)=0\) by \SGARefLink{sga3:xxii:cor:5.9.6}{5.9.6}.}

\SGARefTarget{sga3:xxii:cor:5.9.8}
\paragraph{\SGARefLink{sga3:xxii:cor:5.9.8}{Corollary 5.9.8}.}
{\itshape If \(G\) is a reductive \(S\)-group, the canonical
morphism \textup{(cf.\ \SGARefLink{sga3:xxii:cor:5.8.3}{5.8.3})}
\[
  \underline{\operatorname{Kil}}(G)
  \longrightarrow
  \underline{\operatorname{Bor}}(G)
\]
has sections over every affine open subset.}

\SGARefTarget{sga3:xxii:cor:5.9.9}
\paragraph{\SGARefLink{sga3:xxii:cor:5.9.9}{Corollary 5.9.9}.}
{\itshape Under the hypotheses of \SGARefLink{sga3:xxii:prop:5.9.5}{5.9.5}, suppose that \(S\) is
affine. Then there exists a locally free \(\mathcal O_S\)-module
\(\mathcal E\) such that \(B^u\), as a scheme, is \(S\)-isomorphic
to \(W(\mathcal E)\).}

We prove by induction on \(i\) that \(B^u/U_i\) is \(S\)-isomorphic
to
\[
  W(\mathcal E_0\oplus\cdots\oplus\mathcal E_{i-1}).
\]
This is clear for \(i=0\); suppose that \(i\geq1\). Then
\(B^u/U_i\) is a principal homogeneous bundle with base
\[
  X=B^u/U_{i-1}
\]
under the group \((U_{i-1}/U_i)_X\). Since \(B^u/U_{i-1}\) is
affine by the induction hypothesis, and since
\[
  U_{i-1}/U_i=W(\mathcal E_{i-1}),
\]
this bundle is trivial. We therefore have, at least, an isomorphism
of \(S\)-schemes
\[
  B^u/U_i
  \xrightarrow{\;\sim\;}
  (B^u/U_{i-1})\times_S W(\mathcal E_{i-1})
  =
  W(\mathcal E_0\oplus\cdots\oplus\mathcal E_{i-1}).
\]
The conclusion now follows at once from condition
\SGARefLink{sga3:xxii:prop:5.9.5}{5.9.5}\textup{\SGARefLink{sga3:xxii:prop:5.9.5:item:iii}{(iii)}}.

\SGARefTarget{sga3:xxii:cor:5.9.10}
\paragraph{\SGARefLink{sga3:xxii:cor:5.9.10}{Corollary 5.9.10}.}
{\itshape Let \(S\) be a semilocal scheme, let \(\{s_i\}\) be its
closed points, and let \(B\) be a Borel subgroup of the reductive
\(S\)-group \(G\). The canonical map
\[
  B^u(S)
  \longrightarrow
  \prod_i B^u\bigl(\operatorname{Spec}\kappa(s_i)\bigr)
\]
is surjective.}

Indeed, if \(S=\operatorname{Spec}(A)\),
\(\kappa(s_i)=A/\mathfrak p_i\), and \(\mathcal E\) is defined by
the \(A\)-module \(E\), then
\[
  B^u(S)=E\otimes A,
  \qquad
  B^u\bigl(\operatorname{Spec}\kappa(s_i)\bigr)
  =
  E\otimes_A(A/\mathfrak p_i).
\]
The assertion follows at once from the surjectivity of
\[
  A\longrightarrow\prod_i A/\mathfrak p_i.
\]

\SGARefTarget{sga3:xxii:section:5.10}
\section*{5.10. Subgroups of Type \textup{(R)} with Reductive
Fibers}
\addcontentsline{toc}{section}{\SGARefLink{sga3:xxii:section:5.10}{5.10}. Subgroups of Type (R) with
Reductive Fibers}

\SGARefTarget{sga3:xxii:prop:5.10.1}
\paragraph{\SGARefLink{sga3:xxii:prop:5.10.1}{Proposition 5.10.1}.}
{\itshape Let \((G,T,M,R)\) be a split \(S\)-group, let \(R'\) be a
subset of \(R\) of type \textup{(R)} \textup{(\SGARefLink{sga3:xxii:def:5.4.2}{5.4.2})}, and let
\(H_{R'}\) be the corresponding subgroup of \(G\). The following
conditions are equivalent:
\begin{enumerate}
\item[\textup{\SGAInlineRefTarget{sga3:xxii:prop:5.10.1:item:i}(i)}] \(H_{R'}\) is reductive \textup{(that is, its
geometric fibers are reductive)}.

\item[\textup{\SGAInlineRefTarget{sga3:xxii:prop:5.10.1:item:ii}(ii)}] \(R'=-R'\), that is, \(R'\) is symmetric.
\end{enumerate}
Moreover, under these conditions, \((H_{R'},T,M,R')\) is a
splitting of \(H_{R'}\); for every system \(R_+\) of positive roots
of \(R\),
\[
  R'_+=R'\cap R_+
\]
is a system of positive roots of \(R'\), and
\[
  B_{R_+}\cap H_{R'}=H_{R'_+}.
\]
This is a Borel subgroup of \(H_{R'}\), whose unipotent part is
\[
  U_{R_+}\cap H_{R'}=U_{R'_+}.
\]}

Clearly \textup{\SGARefLink{sga3:xxii:prop:5.10.1:item:i}{(i)}} implies \textup{\SGARefLink{sga3:xxii:prop:5.10.1:item:ii}{(ii)}}: it is enough to check
this fiber by fiber, and \(R'\) is the root system of \(H_{R'}\)
relative to \(T\). To prove that \textup{\SGARefLink{sga3:xxii:prop:5.10.1:item:ii}{(ii)}} implies
\textup{\SGARefLink{sga3:xxii:prop:5.10.1:item:i}{(i)}}, observe, by \SGARefLink{sga3:xxii:lem:5.4.3}{5.4.3}, that
\[
  H_{R'}\cap Z_\alpha
  =
  \underline{\operatorname{Centr}}_{H_{R'}}(T_\alpha)
  =
  Z_\alpha
\]
for every \(\alpha\in R'\), and apply the criterion of \SGARefLink{sga3:xix:proposition:1-12}{Exposé XIX,
1.12}.

If \(R_+\) is a system of positive roots of \(R\), then
\(R'_+=R_+\cap R'\) is plainly a closed subset of \(R'\) such that
\[
  R'_+\cup(-R'_+)=R',
  \qquad
  R'_+\cap(-R'_+)=\varnothing;
\]
it is therefore a system of positive roots of \(R'\). The other two
assertions follow respectively from \SGARefLink{sga3:xxii:prop:5.6.1}{5.6.1}\textup{\SGARefLink{sga3:xxii:prop:5.6.1:item:vi}{(vi)}} and
\SGARefLink{sga3:xxii:cor:5.6.7}{5.6.7}\textup{\SGARefLink{sga3:xxii:cor:5.6.7:item:i}{(i)}}.
\hfill\(\square\)

\SGARefTarget{sga3:xxii:cor:5.10.2}
\paragraph{\SGARefLink{sga3:xxii:cor:5.10.2}{Corollary 5.10.2}.}
{\itshape Let \(S\) be a scheme, \(G\) a reductive \(S\)-group,
and \(H\) a reductive \(S\)-subgroup such that, for every
\(s\in S\), \(G_{\bar s}\) and \(H_{\bar s}\) have the same
reductive rank. Then \(H\) is closed in \(G\),
\(\underline{\operatorname{Norm}}_G(H)\) is smooth over \(S\), and
\[
  \underline{\operatorname{Norm}}_G(H)/H
\]
is representable by a finite étale \(S\)-scheme.

If \(T\) is a maximal torus of \(H\), and \(B\) a Borel subgroup of
\(G\) containing \(T\), then \(B\cap H\) is a Borel subgroup of
\(H\), whose unipotent part is
\[
  (B\cap H)^u=B^u\cap H.
\]}

The first assertions follow at once from \SGARefLink{sga3:xxii:cor:5.3.10}{5.3.10} and \SGARefLink{sga3:xxii:cor:5.3.18}{5.3.18}, using
the fact that the Weyl groups of \(G\) and \(H\) are finite
\textup{(\SGARefLink{sga3:xix:theorem:2:5}{Exposé XIX, 2.5})}. The other assertions are local for the
étale topology and reduce to the case studied in \SGARefLink{sga3:xxii:prop:5.10.1}{5.10.1}.
\hfill\(\square\)

\SGARefTarget{sga3:xxii:prop:5.10.3}
\paragraph{\SGARefLink{sga3:xxii:prop:5.10.3}{Proposition 5.10.3}.}
{\itshape Let \(S\) be a scheme and \(G\) a reductive \(S\)-group.
\begin{enumerate}
\item[\textup{\SGAInlineRefTarget{sga3:xxii:prop:5.10.3:item:a}(a)}] If \(Q\) is a torus of \(G\), then
\(\underline{\operatorname{Centr}}_G(Q)\) is a subgroup of \(G\) of
type \textup{(R)} with reductive fibers. If \(Q\subset Q'\) are two
tori of \(G\), then
\[
  \underline{\operatorname{Centr}}_G(Q)
  \supset
  \underline{\operatorname{Centr}}_G(Q').
\]

\item[\textup{\SGAInlineRefTarget{sga3:xxii:prop:5.10.3:item:b}(b)}] If \(H\) is a subgroup of \(G\) of type
\textup{(R)} with reductive fibers, then
\(\operatorname{rad}(H)\) \textup{(\SGARefLink{sga3:xxii:def:4.3.6}{4.3.6})} is a torus of \(G\). If
\(H\subset H'\) are two subgroups of \(G\) of type \textup{(R)}
with reductive fibers, then
\[
  \operatorname{rad}(H)\supset\operatorname{rad}(H').
\]

\item[\textup{\SGAInlineRefTarget{sga3:xxii:prop:5.10.3:item:c}(c)}] If \(Q\) is a torus of \(G\), then
\[
  \operatorname{rad}\bigl(
    \underline{\operatorname{Centr}}_G(Q)\bigr)
  \supset Q
  \quad\text{and}\quad
  \underline{\operatorname{Centr}}_G\!\left(
    \operatorname{rad}\bigl(
      \underline{\operatorname{Centr}}_G(Q)\bigr)\right)
  =
  \underline{\operatorname{Centr}}_G(Q).
\]

\item[\textup{\SGAInlineRefTarget{sga3:xxii:prop:5.10.3:item:d}(d)}] If \(H\) is a subgroup of \(G\) of type
\textup{(R)} with reductive fibers, then
\[
  \underline{\operatorname{Centr}}_G\bigl(\operatorname{rad}(H)\bigr)
  \supset H
  \quad\text{and}\quad
  \operatorname{rad}\!\left(
    \underline{\operatorname{Centr}}_G\bigl(
      \operatorname{rad}(H)\bigr)\right)
  =
  \operatorname{rad}(H).
\]
\end{enumerate}}

Indeed, \textup{\SGARefLink{sga3:xxii:prop:5.10.3:item:a}{(a)}} follows at once from \SGARefLink{sga3:xix:paragraph:2:8}{Exposé XIX, 2.8}. To prove
\textup{\SGARefLink{sga3:xxii:prop:5.10.3:item:b}{(b)}}, it is enough to observe that
\(\operatorname{rad}(H')\subset H\), since \(H\) contains,
fpqc-locally, a maximal torus of \(G\), hence of \(H'\). The first
assertion of \textup{\SGARefLink{sga3:xxii:prop:5.10.3:item:c}{(c)}} \textup{(respectively, of \textup{\SGARefLink{sga3:xxii:prop:5.10.3:item:d}{(d)}})}
is trivial, and the second follows by the usual argument.
\hfill\(\square\)

This proposition leads to the following definition:

\SGARefTarget{sga3:xxii:def:5.10.4}
\paragraph{\SGARefLink{sga3:xxii:def:5.10.4}{Definition 5.10.4}.}
Let \(S\) be a scheme, \(G\) a reductive \(S\)-group, \(H\) a
reductive subgroup of \(G\) of type \textup{(R)}, and \(Q\) a
subtorus of \(G\).\footnote{Editor's note: We have modified the
original by introducing the term ``\(C\)-critical torus'' in place
of ``critical torus'', in order to avoid confusion in later
references \textup{(cf.\ \SGARefLink{sga3:xxvi:item:3:9}{Exposé XXVI, 3.9})}. We have also expanded
the statement of \SGARefLink{sga3:xxii:cor:5.10.5}{5.10.5} and added \SGARefLink{sga3:xxii:rem:5.10.5.1}{Remark 5.10.5.1}.}
\begin{enumerate}
\item[\textup{1)}] \(H\) is called a \emph{critical subgroup} if it
is the centralizer of its radical.
\item[\textup{2)}] \(Q\) is called a \emph{\(C\)-critical torus} if
it is the radical of its centralizer.
\end{enumerate}

\SGARefLink{sga3:xxii:prop:5.10.3}{Proposition 5.10.3} then gives:

\SGARefTarget{sga3:xxii:cor:5.10.5}
\paragraph{\SGARefLink{sga3:xxii:cor:5.10.5}{Corollary 5.10.5}.}
\begin{enumerate}
\item[\textup{(i)}] For every subtorus \(Q\) of \(G\),
\(\underline{\operatorname{Centr}}_G(Q)\) is critical.

\item[\textup{(ii)}] For every subgroup \(H\) of \(G\) of type
\textup{(R)} with reductive fibers, \(\operatorname{rad}(H)\) is a
\(C\)-critical torus of \(G\).

\item[\textup{(iii)}] The maps
\[
  Q\longmapsto\underline{\operatorname{Centr}}_G(Q),
  \qquad
  H\longmapsto\operatorname{rad}(H)
\]
are mutually inverse bijections between the set of \(C\)-critical
tori of \(G\) and the set of its critical reductive subgroups of
type \textup{(R)}.

\item[\textup{(iv)}] If \(Q\) is a torus of \(G\), then
\[
  \operatorname{rad}\bigl(
    \underline{\operatorname{Centr}}_G(Q)\bigr)
\]
is the smallest \(C\)-critical torus of \(G\) containing \(Q\).

\item[\textup{(v)}] If \(H\) is a reductive subgroup of \(G\) of
type \textup{(R)}, then
\[
  \underline{\operatorname{Centr}}_G\bigl(\operatorname{rad}(H)\bigr)
\]
is the smallest critical reductive subgroup of \(G\) of type
\textup{(R)} containing \(H\).
\end{enumerate}

\SGARefTarget{sga3:xxii:rem:5.10.5.1}
\paragraph{\SGARefLink{sga3:xxii:rem:5.10.5.1}{Remark 5.10.5.1}.}
\textup{(58)}
\begin{enumerate}
\item[\textup{1)}] A torus \(T\) of \(G\) is a critical subgroup of
\(G\) if and only if it is a maximal torus.
\item[\textup{2)}] Henceforth, ``critical torus'' means
``\(C\)-critical torus''.
\end{enumerate}

\SGARefTarget{sga3:xxii:prop:5.10.6}
\paragraph{\SGARefLink{sga3:xxii:prop:5.10.6}{Proposition 5.10.6}.}
{\itshape Let \((G,T,M,R)\) be a split \(S\)-group, and let
\(R'\) be a subset of \(R\). The following conditions are
equivalent:
\begin{enumerate}
\item[\textup{\SGAInlineRefTarget{sga3:xxii:prop:5.10.6:item:i}(i)}] \(R'\) is of type \textup{(R)}, and \(H_{R'}\)
is reductive and critical.

\item[\textup{(ii)}] There exist a system \(\Delta\) of simple
roots of \(R\) and a subset \(\Delta'\) of \(\Delta\) such that
\(R'\) is the set of those elements of \(R\) that are linear
combinations of elements of \(\Delta'\).

\item[\textup{\SGAInlineRefTarget{sga3:xxii:prop:5.10.6:item:iii}(iii)}] \(R'\) is closed and symmetric, and every
system of simple roots of \(R'\) is the intersection with \(R'\) of
a system of simple roots of \(R\).
\end{enumerate}}

Indeed, by \SGARefLink{sga3:XXI:3.4.8}{Exposé XXI, 3.4.8}, conditions \textup{\SGARefLink{sga3:XXI:3.4.8:item:ii}{(ii)}} and
\textup{\SGARefLink{sga3:xxii:prop:5.10.6:item:iii}{(iii)}} are equivalent; they are also equivalent to the
condition that \(R'\) be the intersection of \(R\) with a
\(\mathbf Q\)-vector subspace of \(M\otimes\mathbf Q\). This last
condition follows from \textup{\SGARefLink{sga3:xxii:prop:5.10.6:item:i}{(i)}}: if
\[
  H_{R'}=\underline{\operatorname{Centr}}_G(Q),
\]
then \(R'\) is the set of elements of \(R\) that vanish on \(Q\)
\textup{(\SGARefLink{sga3:ii:theorem:5-2-3:item:ii}{Exposé II, 5.2.3(ii)})}. Conversely, this condition implies
\textup{\SGARefLink{sga3:xxii:prop:5.10.6:item:i}{(i)}}, because \(\operatorname{rad}(H_{R'})\) is the maximal
torus of
\[
  \bigcap_{\alpha\in R'}\operatorname{Ker}(\alpha).
\]
Consequently,
\[
  \underline{\operatorname{Centr}}_G\bigl(
    \operatorname{rad}(H_{R'})\bigr)
  =
  H_{R''},
\]
where \(R''\) is the intersection of \(R\) with the vector subspace
spanned by \(R'\); hence \(R''=R'\).
\hfill\(\square\)

\SGARefTarget{sga3:xxii:para:5.10.7}
\paragraph{\SGARefLink{sga3:xxii:para:5.10.7}{5.10.7}.}
Let us summarize some of the preceding results. Let
\((G,T,M,R)\) be a split \(S\)-group, let \(\Delta\) be a system
of simple roots of \(R\), and let \(R_+\) be the corresponding
system of positive roots. Choose a subset \(\Delta'\) of
\(\Delta\), denote by \(R'\) the set of those elements of \(R\)
which are linear combinations of elements of \(\Delta'\), and put
\[
  R'_+=R'\cap R_+.
\]
Let \(T_{\Delta'}\) be the maximal torus of
\[
  \bigcap_{\alpha\in\Delta'}\operatorname{Ker}(\alpha)
\]
and put
\[
  Z_{\Delta'}
  =
  \underline{\operatorname{Centr}}_G(T_{\Delta'}).
\]
Then \(Z_{\Delta'}\) is a reductive subgroup of \(G\), with radical
\(T_{\Delta'}\); \((Z_{\Delta'},T,M,R')\) is a split
\(S\)-group; \(B_{R_+}\cap Z_{\Delta'}\) is the Borel subgroup of
\(Z_{\Delta'}\) defined by the system \(R'_+\) of positive roots
\textup{(or, equivalently, by the system \(\Delta'\) of simple
roots)}, and its unipotent part is
\[
  U_{R_+}\cap Z_{\Delta'}=U_{R'_+}.
\]

\SGARefTarget{sga3:xxii:rem:5.10.8}
\paragraph{\SGARefLink{sga3:xxii:rem:5.10.8}{Remark 5.10.8}.}
Under the conditions of \SGARefLink{sga3:xxii:def:5.10.4}{5.10.4}, let \(Q\) be a critical torus of
\(G\), and let
\[
  L=\underline{\operatorname{Centr}}_G(Q)
\]
be its centralizer. Since \(Q=\operatorname{rad}(L)\), \(Q\) is a
characteristic subgroup of \(L\); it follows at once that
\[
  \underline{\operatorname{Norm}}_G(L)
  =
  \underline{\operatorname{Norm}}_G(Q),
\]
and hence also that
\[
  \underline{\operatorname{Norm}}_G(L)/L
  =
  \underline{\operatorname{Norm}}_G(Q)/
  \underline{\operatorname{Centr}}_G(Q)
  =
  W_G(Q).
\]
By \SGARefLink{sga3:xxii:cor:5.10.2}{5.10.2}, we deduce:

\SGARefTarget{sga3:xxii:prop:5.10.9}
\paragraph{\SGARefLink{sga3:xxii:prop:5.10.9}{Proposition 5.10.9}.}
{\itshape Let \(S\) be a scheme, let \(G\) be a reductive
\(S\)-group, and let \(Q\) be a critical torus of \(G\). The Weyl
group \(W_G(Q)\) is finite \textup{(étale)} over \(S\).}

\SGARefTarget{sga3:xxii:rem:5.10.10}
\paragraph{\SGARefLink{sga3:xxii:rem:5.10.10}{Remark 5.10.10}.}
Under the conditions of \SGARefLink{sga3:xxii:para:5.10.7}{5.10.7}, one may describe explicitly
\[
  W_G(T_{\Delta'})
  =
  \underline{\operatorname{Norm}}_G(Z_{\Delta'})/Z_{\Delta'}.
\]
This is the constant group associated with the quotient \(W_1/W_2\),
where \(W_1\) is the subgroup of \(W\) consisting of those elements
which normalize the subgroup of \(M\) generated by \(\Delta'\), and
\(W_2\) is the subgroup of \(W\) generated by the \(s_\alpha\), for
\(\alpha\in\Delta'\).

\SGARefTarget{sga3:xxii:section:5.11}
\section*{5.11. Subgroups of Type \textup{(RC)}}
\addcontentsline{toc}{section}{\SGARefLink{sga3:xxii:section:5.11}{5.11}. Subgroups of Type (RC)}

\SGARefTarget{sga3:xxii:def:5.11.1}
\paragraph{\SGARefLink{sga3:xxii:def:5.11.1}{Definition 5.11.1}.}
Let \(S\) be a scheme and let \(G\) be a reductive \(S\)-group. A
group subscheme \(H\) of \(G\) is said to be \emph{of type
\textup{(RC)}} if it is of type \textup{(R)}, that is, if it
satisfies the following two conditions from \SGARefLink{sga3:xxii:def:5.2.1}{5.2.1}:
\begin{enumerate}
\item[\textup{(i)}] \(H\) is smooth over \(S\), with connected
fibers;

\item[\textup{(ii)}] for every \(s\in S\), \(H_{\bar s}\) contains
a maximal torus of \(G_{\bar s}\);
\end{enumerate}
and if it also satisfies the following condition:
\begin{enumerate}
\item[\textup{\SGAInlineRefTarget{sga3:xxii:def:5.11.1:item:iii}(iii)}] for every \(s\in S\) and every maximal torus
\(T\) of \(H_{\bar s}\), the set of roots of \(H_{\bar s}\)
relative to \(T\) is a closed subset of the set of all roots of
\(G_{\bar s}\) relative to \(T\).
\end{enumerate}

\SGARefTarget{sga3:xxii:rem:5.11.2}
\paragraph{\SGARefLink{sga3:xxii:rem:5.11.2}{Remark 5.11.2}.}
As we already observed in \SGARefLink{sga3:xxii:rem:5.4.8}{5.4.8}, condition \textup{\SGARefLink{sga3:xxii:def:5.11.1:item:iii}{(iii)}} follows
from the other two when \(6\) is invertible on
\(S\).\footnote{Editor's note: That is, when every residue
characteristic of \(S\) is greater than \(3\).}

\SGARefTarget{sga3:xxii:lem:5.11.3}
\paragraph{\SGARefLink{sga3:xxii:lem:5.11.3}{Lemma 5.11.3}.}
{\itshape Let \((G,T,M,R)\) be a split \(S\)-group, and let \(R'\)
be a closed subset of \(R\). Put
\[
  R_1=\{\alpha\in R'\mid-\alpha\in R'\},
  \qquad
  R_2=\{\alpha\in R'\mid-\alpha\notin R'\}.
\]
Then \(R_1\) and \(R_2\) are closed. Consider the groups
\(H_{R'}\), \(H_{R_1}\), and \(U_{R_2}\) \textup{(\SGARefLink{sga3:xxii:thm:5.4.7}{5.4.7} and
\SGARefLink{sga3:xxii:cor:5.6.5}{5.6.5})}, which are smooth and have connected fibers.
\begin{enumerate}
\item[\textup{\SGAInlineRefTarget{sga3:xxii:lem:5.11.3:item:i}(i)}] \(U_{R_2}\) is normal in \(H_{R'}\), and
\(H_{R'}\) is the semidirect product of \(U_{R_2}\) by
\(H_{R_1}\).

\item[\textup{\SGAInlineRefTarget{sga3:xxii:lem:5.11.3:item:ii}(ii)}] The group \(H_{R_1}\) is reductive, and
\(U_{R_2}\) has connected unipotent geometric fibers. Every normal
subgroup scheme of \(H_{R'}\) that is smooth over \(S\) and has
connected unipotent geometric fibers is contained in \(U_{R_2}\).
Every reductive subgroup of \(H_{R'}\) that contains \(T\) is
contained in \(H_{R_1}\).

\item[\textup{\SGAInlineRefTarget{sga3:xxii:lem:5.11.3:item:iii}(iii)}]
\[
  U_{R_2}\cap\underline{\operatorname{Norm}}_G(H_{R_1})=e.
\]
\end{enumerate}}

Let us first prove \textup{\SGARefLink{sga3:xxii:lem:5.11.3:item:iii}{(iii)}}, which follows from \SGARefLink{sga3:xxii:cor:5.6.7:item:i}{5.6.7(i)}. The
first assertion of \textup{\SGARefLink{sga3:xxii:lem:5.11.3:item:i}{(i)}} follows from \SGARefLink{sga3:xxii:cor:5.6.7:item:ii}{5.6.7(ii)}. Since
\[
  U_{R_2}\cap H_{R_1}=e
\]
by \textup{\SGARefLink{sga3:xxii:lem:5.11.3:item:iii}{(iii)}}, the semidirect product \(H_{R_1}\cdot U_{R_2}\)
is a subgroup of \(H_{R'}\). But these are two subgroups of \(G\)
of type \textup{(R)}, containing \(T\), and having the same Lie
algebra \(\mathfrak g_{R'}\); they coincide by \SGARefLink{sga3:xxii:cor:5.3.5}{5.3.5}, which proves
\textup{\SGARefLink{sga3:xxii:lem:5.11.3:item:i}{(i)}}.

For \textup{\SGARefLink{sga3:xxii:lem:5.11.3:item:ii}{(ii)}}, the first two assertions are merely \SGARefLink{sga3:xxii:prop:5.10.1}{5.10.1} and
\SGARefLink{sga3:xxii:cor:5.6.5}{5.6.5}. Let \(U\) be a subgroup scheme of \(H_{R'}\), smooth and of
finite presentation over \(S\), normal in \(H_{R'}\) \textup{(and
hence normalized by \(T\))}, and with connected unipotent geometric
fibers. By \SGARefLink{sga3:xxii:cor:5.6.12}{5.6.12}, locally one has \(U=U_{R''}\), where
\(R''\subset R'\) and
\[
  R''\cap(-R'')=\varnothing.
\]
If \(U\) were not contained in \(U_{R_2}\), then \(R''\) would not
be contained in \(R_2\); there would therefore exist
\(\alpha\in R''\) such that \(-\alpha\in R'\). Then
\(Z_\alpha\subset H_{R'}\) by \SGARefLink{sga3:xxii:lem:5.4.3}{5.4.3}, so \(Z_\alpha\) normalizes
\(U\). But \(U\) contains \(U_\alpha\), and \(Z_\alpha\) has a
section \(w\) such that
\[
  \operatorname{int}(w)U_\alpha=U_{-\alpha};
\]
thus \(-\alpha\in R''\), a contradiction.

Finally, if \(L\) is a reductive subgroup of \(H_{R'}\) containing
\(T\), locally one has \(L=H_{R'''}\), where \(R'''\) is symmetric
and contained in \(R'\); hence \(R'''\subset R_1\).
\hfill\(\square\)

\SGARefTarget{sga3:xxii:prop:5.11.4}
\paragraph{\SGARefLink{sga3:xxii:prop:5.11.4}{Proposition 5.11.4}.}
{\itshape Let \(S\) be a scheme, let \(G\) be a reductive
\(S\)-group, and let \(H\) be a subgroup scheme of \(G\) of type
\textup{(RC)}.
\begin{enumerate}
\item[\textup{\SGAInlineRefTarget{sga3:xxii:prop:5.11.4:item:i}(i)}] \(H\) is closed in \(G\), and
\(\underline{\operatorname{Norm}}_G(H)/H\) is representable by a
finite étale \(S\)-group scheme.

\item[\textup{(ii)}] \(H\) has a largest normal subgroup scheme
that is smooth and of finite presentation over \(S\), with
connected unipotent geometric fibers. It is called the
\emph{unipotent radical} of \(H\) and is denoted
\(\operatorname{rad}^{u}(H)\). The quotient sheaf
\[
  H/\operatorname{rad}^{u}(H)
\]
is representable by a reductive \(S\)-group.

\item[\textup{(iii)}] If \(T\) is a maximal torus of \(H\), then
\(H\) has a reductive subgroup \(L\), containing \(T\), of type
\textup{(RC)}, with the following properties:
\begin{enumerate}
\item[\textup{\SGAInlineRefTarget{sga3:xxii:prop:5.11.4:item:a}(a)}] every reductive subgroup of \(H\) containing
\(T\) is contained in \(L\);

\item[\textup{\SGAInlineRefTarget{sga3:xxii:prop:5.11.4:item:b}(b)}] \(H\) is the semidirect product
\[
  H=L\cdot\operatorname{rad}^{u}(H);
\]
in other words, the canonical morphism
\[
  L\longrightarrow H/\operatorname{rad}^{u}(H)
\]
is an isomorphism.
\end{enumerate}
\end{enumerate}

Moreover, \(L\) is the unique reductive subgroup of \(H\) that
contains \(T\) and satisfies either one of the two preceding
conditions. Finally,
\[
  \underline{\operatorname{Norm}}_H(L)=L,
  \qquad
  \underline{\operatorname{Norm}}_H(T)
    =\underline{\operatorname{Norm}}_L(T),
  \qquad
  W_H(T)=W_L(T);
\]
in particular, \(W_H(T)\) is finite over \(S\).}

\textit{Proof.}
First observe that \textup{\SGARefLink{sga3:xxii:prop:5.11.4:item:i}{(i)}} is local for the étale topology.
Thus, by \SGARefLink{sga3:xxii:cor:5.3.18}{Corollary 5.3.18}, \textup{\SGARefLink{sga3:xxii:cor:5.3.18:item:i}{(i)}} follows from the last
assertion of \textup{\SGARefLink{sga3:xxii:cor:5.3.18:item:iii}{(iii)}}.

The assertions of \textup{\SGARefLink{sga3:xxii:cor:5.3.18:item:ii}{(ii)}} are local for the étale topology.
We may therefore suppose ourselves in the situation of \SGARefLink{sga3:xxii:lem:5.11.3}{5.11.3},
where they follow immediately from \textup{\SGARefLink{sga3:xxii:lem:5.11.3:item:i}{(i)}} and \textup{\SGARefLink{sga3:xxii:lem:5.11.3:item:ii}{(ii)}}.

By virtue of the uniqueness assertions contained there,
\textup{\SGARefLink{sga3:xxii:lem:5.11.3:item:iii}{(iii)}} is likewise local for the étale topology, and we may
again reduce to the situation of \SGARefLink{sga3:xxii:lem:5.11.3}{5.11.3}, where properties
\textup{\SGARefLink{sga3:xxii:prop:5.11.4:item:a}{(a)}} and \textup{\SGARefLink{sga3:xxii:prop:5.11.4:item:b}{(b)}} have been verified. The uniqueness of
an \(L\) satisfying \textup{\SGARefLink{sga3:xxii:prop:5.11.4:item:a}{(a)}} is immediate; in view of
\textup{\SGARefLink{sga3:xxii:prop:5.11.4:item:a}{(a)}}, the uniqueness of an \(L\) satisfying \textup{\SGARefLink{sga3:xxii:prop:5.11.4:item:b}{(b)}}
is clear. The equality
\[
  \underline{\operatorname{Norm}}_H(L)=L
\]
is precisely \SGARefLink{sga3:xxii:lem:5.11.3:item:iii}{5.11.3(iii)}. If a section of \(H\) normalizes \(T\),
then it normalizes \(L\), by the uniqueness of \(L\), and hence is
a section of \(L\) by what we have just proved. This proves the
second equality; the third is then immediate.

\SGARefTarget{sga3:xxii:prop:5.11.5}
\paragraph{\SGARefLink{sga3:xxii:prop:5.11.5}{Proposition 5.11.5}.}
{\itshape Let \(S\) be a scheme, let \(G\) be a reductive
\(S\)-group, and let \(\mathcal H_c\) be the functor of subgroups
of \(G\) of type \textup{(RC)}, which is a subfunctor of the
functor \(\mathcal H\) of \SGARefLink{sga3:xxii:thm:5.8.1}{5.8.1}.
\begin{enumerate}
\item[\textup{\SGAInlineRefTarget{sga3:xxii:prop:5.11.5:item:i}(i)}] \(\mathcal H_c\) is representable by an open
subscheme of \(\mathcal H\), smooth, quasi-projective, and of finite
presentation over \(S\).

\item[\textup{\SGAInlineRefTarget{sga3:xxii:prop:5.11.5:item:ii}(ii)}] There exist a finite étale \(S\)-scheme
\(\mathcal C\ell_c\) and a morphism
\[
  \operatorname{cl}\colon\mathcal H_c\longrightarrow\mathcal C\ell_c
\]
that is smooth, quasi-projective, of finite presentation,
surjective, and has connected geometric fibers, with the following
property:

For every \(S'\to S\) and all
\(H,H'\in\mathcal H_c(S')\), one has
\(\operatorname{cl}(H)=\operatorname{cl}(H')\) if and only if
\(H\) and \(H'\) are conjugate in \(G\) locally for the étale
topology \textup{(or, equivalently by \SGARefLink{sga3:xxii:cor:5.3.11}{5.3.11}, if for every
\(s\in S\), \(H_{\bar s}\) and \(H'_{\bar s}\) are conjugate by an
element of \(G(\bar s)\))}.

\item[\textup{\SGAInlineRefTarget{sga3:xxii:prop:5.11.5:item:iii}(iii)}] \(\mathcal C\ell_c\) and
\(\operatorname{cl}\) are determined \textup{(up to a unique
isomorphism)} by the preceding conditions.

\item[\textup{\SGAInlineRefTarget{sga3:xxii:prop:5.11.5:item:iv}(iv)}] If \((G,T,M,R)\) is a splitting of \(G\), let
\(E\) be the set of conjugacy classes modulo \(W\) of closed
subsets of \(R\). Then there is an isomorphism
\[
  \mathcal C\ell_c\xrightarrow{\;\sim\;}E_S
\]
such that, for every closed subset \(R'\) of \(R\),
\(\operatorname{cl}(H_{R'})\) corresponds to the canonical image
of \(R'\) in
\[
  E_S(S)=\operatorname{Hom}_{\mathrm{loc.const.}}(S,E).
\]
\end{enumerate}}

It is first of all clear that \(\mathcal H_c\) is a sheaf for the
étale topology, and that \textup{\SGARefLink{sga3:xxii:prop:5.11.5:item:ii}{(ii)}} implies that
\(\mathcal C\ell_c\) is simply the quotient sheaf of
\(\mathcal H_c\) by the equivalence relation defined by
conjugation.

This first implies \textup{\SGARefLink{sga3:xxii:prop:5.11.5:item:iii}{(iii)}}, as well as the fact that it
suffices to verify \textup{\SGARefLink{sga3:xxii:prop:5.11.5:item:i}{(i)}} and \textup{\SGARefLink{sga3:xxii:prop:5.11.5:item:ii}{(ii)}} locally for the
étale topology. We are therefore reduced to the situation of
\textup{\SGARefLink{sga3:xxii:prop:5.11.5:item:iv}{(iv)}}; let us first construct a morphism
\[
  f\colon\mathcal H_c\longrightarrow E_S.
\]
It suffices to construct a map
\(\mathcal H_c(S)\to E_S(S)\), functorial in \(S\); let \(H\) be a
subgroup of \(G\) of type \textup{(RC)}. Since \(H\) locally for
the étale topology has maximal tori, and since the maximal tori of
\(G\) are locally conjugate for the étale topology, there exist a
covering family \((S_i\to S)\), elements \(g_i\in G(S_i)\), and
closed subsets \(R_i\) of \(R\) such that
\[
  \operatorname{int}(g_i)(H\times_S S_i)
  =H_{R_i}\times_S S_i.
\]
Each \(R_i\) defines a section \(\eta_i\) of \(E_{S_i}\), that is,
an element of \(E_S(S_i)\). It remains to prove that the family
\((\eta_i)\) comes from a section \(\eta=f(H)\) of \(E_S\) over
\(S\), and that it depends only on \(H\).

For this it is enough to prove that \(H_{R'}\) and \(H_{R''}\) are
locally conjugate for the étale topology if and only if \(R'\) and
\(R''\) are conjugate by an element of the Weyl group \(W\), which
is clear.

For every \(\eta\in E\), there exists
\(H_0\in\mathcal H_c(S)\) such that \(f(H_0)=\eta\): take
\(H_0=H_{R'}\), where \(R'\) is a closed subset whose image is
\(\eta\). If \(H\in\mathcal H_c(S')\), for \(S'\to S\), then \(H\)
is locally conjugate to \(H_0\) if and only if \(f(H)=\eta\).
Consequently,
\[
  f^{-1}(\eta)\simeq
  G\big/\underline{\operatorname{Norm}}_G(H_0),
\]
which, by \SGARefLink{sga3:xxii:cor:5.8.2}{5.8.2}, is an open subscheme of \(\mathcal H\), smooth,
quasi-projective, and of finite presentation over \(S\), with
nonempty connected fibers. Since \(E_S\) is the disjoint union of
the open subschemes that are the images of the sections
corresponding to the elements \(\eta\in E\), \(\mathcal H_c\)
identifies with the disjoint union of the \(f^{-1}(\eta)\). This
proves \textup{\SGARefLink{sga3:xxii:prop:5.11.5:item:i}{(i)}} and \textup{\SGARefLink{sga3:xxii:prop:5.11.5:item:ii}{(ii)}}; assertion \textup{\SGARefLink{sga3:xxii:prop:5.11.5:item:iv}{(iv)}}
follows from the construction.

\SGARefTarget{sga3:xxii:cor:5.11.6}
\paragraph{\SGARefLink{sga3:xxii:cor:5.11.6}{Corollary 5.11.6}.}
{\itshape Let \(u\in\mathcal C\ell_c(S')\), where \(S'\to S\).
Then \(\operatorname{cl}^{-1}(u)\) is a smooth,
quasi-projective \(S'\)-scheme of finite presentation, with
nonempty connected fibers. It is an open subscheme of
\(\mathcal H_c\) and a ``homogeneous'' scheme under \(G_{S'}\)
\textup{(acting by inner automorphisms)}.

In particular, if
\(H\in\operatorname{cl}^{-1}(u)(S')\), the morphism
\[
  G_{S'}\longrightarrow(\mathcal H_c)_{S'},\qquad
  g\longmapsto\operatorname{int}(g)H
\]
identifies
\[
  G_{S'}\big/
  \underline{\operatorname{Norm}}_{G_{S'}}(H)
  \quad\text{with}\quad
  \operatorname{cl}^{-1}(u).
\]}

\SGARefTarget{sga3:xxii:examples:5.11.7}
\paragraph{\SGARefLink{sga3:xxii:examples:5.11.7}{Examples 5.11.7}.}
There are two canonical sections \(u_t\) and \(u_b\) of
\(\mathcal C\ell_c\). The section \(u_t\) corresponds to maximal
tori, for which \(R'=\varnothing\). The section \(u_b\) corresponds
to Borel subgroups, for which \(R'\) is a system of positive roots.
The schemes
\(\operatorname{cl}^{-1}(u_t)\) and
\(\operatorname{cl}^{-1}(u_b)\) are respectively
\(\underline{\operatorname{Tor}}(G)\) and
\(\underline{\operatorname{Bor}}(G)\), introduced in \SGARefLink{sga3:xxii:cor:5.8.3}{5.8.3}.
Other examples will be given in Exposé XXVI.

\SGARefTarget{sga3:xxii:rem:5.11.8}
\paragraph{\SGARefLink{sga3:xxii:rem:5.11.8}{Remark 5.11.8}.}
One can construct an \(S\)-scheme \(\mathcal C\ell\), of finite
presentation and unramified over \(S\), and a smooth surjective
morphism
\[
  \mathcal H\longrightarrow\mathcal C\ell
\]
with connected geometric fibers, having properties analogous to
those of \SGARefLink{sga3:xxii:prop:5.11.5}{5.11.5}\textup{\SGARefLink{sga3:xxii:prop:5.11.5:item:ii}{(ii)}} and \textup{\SGARefLink{sga3:xxii:prop:5.11.5:item:iii}{(iii)}}.

\SGARefTarget{sga3:xxii:section:6}
\section*{6. The Derived Group}
\addcontentsline{toc}{section}{6. The Derived Group}

\SGARefTarget{sga3:xxii:section:6.1}
\subsection*{\SGAInlineRefTarget{sga3:xxii:subsection:6-1}6.1. Preliminaries}
\addcontentsline{toc}{subsection}{\SGARefLink{sga3:xxii:subsection:6-1}{6.1}. Preliminaries}

In this section, we fix a scheme \(S\), a split \(S\)-group
\((G,T,M,R)\), and a system \(R_+\) of positive roots of \(R\).
We write
\[
  B=B_{R_+},\qquad B^-=B_{R_-},\qquad
  U=B^u,\qquad U^-=(B^-)^u,
\]
and
\[
  \Omega=\Omega_{R_+}=U^-\cdot T\cdot U.
\]

\SGARefTarget{sga3:xxii:par:6.1.1}
\paragraph{\SGARefLink{sga3:xxii:par:6.1.1}{6.1.1}.}
Let \(T'\) denote the subtorus of \(T\) that is the ``image of the
family \(\alpha^*\), \(\alpha\in R\)''; in other words, \(T'\) is
the image of the homomorphism
\[
  (\mathbb G_{m,S})^R\longrightarrow T
\]
defined by
\[
  (z_\alpha)_{\alpha\in R}
  \longmapsto
  \prod_{\alpha\in R}\alpha^*(z_\alpha).
\]
It follows at once that, if \(\Delta\) denotes the set of simple
roots of \(R_+\), the homomorphism
\[
  (\mathbb G_{m,S})^\Delta\longrightarrow T'
\]
defined in the same way is surjective and has finite kernel. If
we identify \(T\) with \(D_S(M)\), then \(T'\) identifies with
\(D_S(M/N)\), where
\[
  N=M\cap V(R^*)^\perp
\]
\textup{(we write \(V(R^*)^\perp\) for the orthogonal complement
of \(V(R^*)\) under the duality between \(V\) and \(V^*\))}.

\SGARefTarget{sga3:xxii:lemma:6.1.2}
\paragraph{\SGARefLink{sga3:xxii:lemma:6.1.2}{Lemma 6.1.2}.}
{\itshape The homomorphism defined by multiplication in \(T\),
\[
  \operatorname{rad}(G)\times_S T'\longrightarrow T,
\]
is an isogeny \textup{(cf.\ \SGARefLink{sga3:xxii:def:4.2.9}{4.2.9})}.}

Indeed, the canonical homomorphism
\[
  \operatorname{rad}(G)\longrightarrow T/T'
\]
comes by duality from the homomorphism of commutative groups
\[
  M\cap V(R^*)^\perp
  \longrightarrow
  M\big/\bigl(M\cap V(R)\bigr),
\]
which is immediately seen to be injective with finite cokernel
\textup{(cf.\ \SGARefLink{sga3:XXI:6.3}{Exposé XXI, 6.3})}.

\SGARefTarget{sga3:xxii:def:6.1.3}
\paragraph{\SGARefLink{sga3:xxii:def:6.1.3}{Definition 6.1.3}.}
We put
\[
  \Omega'=U^-\cdot T'\cdot U;
\]
this is a closed subscheme of
\(\Omega=U^-\cdot T\cdot U\).

\SGARefTarget{sga3:xxii:lemma:6.1.4}
\paragraph{\SGARefLink{sga3:xxii:lemma:6.1.4}{Lemma 6.1.4}.}
{\itshape Let \(\alpha\) be a simple root and let
\(w_\alpha\in\underline{\operatorname{Norm}}_G(T)(S)\) lift
\(s_\alpha\). Then
\[
  \operatorname{int}(w_\alpha)\Omega'\cap\Omega\subset\Omega'.
\]}

It is enough to prove that, if \(g\in\Omega'(S)\) and
\(\operatorname{int}(w_\alpha)g\in\Omega(S)\), then
\(\operatorname{int}(w_\alpha)g\in\Omega'(S)\). By \SGARefLink{sga3:xxii:cor:5.6.8}{5.6.8}, write
\[
  g=a\exp_{-\alpha}(Y)t\exp_\alpha(X)b,
\]
where
\[
  a\in U_{-\widehat\alpha}(S),\qquad
  Y\in\Gamma(S,\mathfrak g^{-\alpha}),\qquad
  t\in T'(S),\qquad
  X\in\Gamma(S,\mathfrak g^\alpha),\qquad
  b\in U_{\widehat\alpha}(S).
\]
We then have
\[
  \operatorname{int}(w_\alpha)g
  =\operatorname{int}(w_\alpha)a\cdot
   \operatorname{int}(w_\alpha)
   \bigl(\exp_{-\alpha}(Y)t\exp_\alpha(X)\bigr)\cdot
   \operatorname{int}(w_\alpha)b.
\]
By \SGARefLink{sga3:xxii:cor:5.6.8}{5.6.8}\textup{\SGARefLink{sga3:xxii:cor:5.6.8:item:iv}{(iv)}},
\[
  \operatorname{int}(w_\alpha)a\in U_{-\widehat\alpha}(S),
  \qquad
  \operatorname{int}(w_\alpha)b\in U_{\widehat\alpha}(S).
\]
It follows that, on putting
\[
  h=\exp_{-\alpha}(Y)t\exp_\alpha(X),
\]
we have the equivalences
\[
  \operatorname{int}(w_\alpha)g\in\Omega(S)
  \quad\Longleftrightarrow\quad
  \operatorname{int}(w_\alpha)h\in\Omega(S),
\]
and
\[
  \operatorname{int}(w_\alpha)g\in\Omega'(S)
  \quad\Longleftrightarrow\quad
  \operatorname{int}(w_\alpha)h\in\Omega'(S).
\]
We are therefore reduced to the case \(g=h\). Since, by
\textup{(\SGARefLink{sga3:xxii:rem:4.1.12}{4.1.12})},
\[
  Z_\alpha\cap\Omega
  =U_{-\alpha}\cdot T\cdot U_\alpha,
  \qquad
  Z_\alpha\cap\Omega'
  =U_{-\alpha}\cdot T'\cdot U_\alpha,
\]
we are reduced to proving the following assertion:
\[
  \operatorname{int}(w_\alpha)h
  \in\bigl(U_{-\alpha}\cdot T\cdot U_\alpha\bigr)(S)
  \quad\Longrightarrow\quad
  \operatorname{int}(w_\alpha)h
  \in\bigl(U_{-\alpha}\cdot T'\cdot U_\alpha\bigr)(S).
\]
But this follows at once from \SGARefLink{sga3:xx:3.12}{Exposé XX, 3.12}, which shows that the
\(T\)-component of \(\operatorname{int}(w_\alpha)h\) is of the form
\[
  t\cdot\alpha^*(z)\in T'(S).
\]

\SGARefTarget{sga3:xxii:lemma:6.1.5}
\paragraph{\SGARefLink{sga3:xxii:lemma:6.1.5}{Lemma 6.1.5}.}
{\itshape For every
\(w\in\underline{\operatorname{Norm}}_G(T)(S)\), there exists an
open subscheme \(V_w\) of \(G\), containing the unit section, such
that
\[
  \operatorname{int}(w)\Omega'\cap V_w\subset\Omega'.
\]}

For each simple root \(\alpha\), choose an element
\(n_\alpha\in\underline{\operatorname{Norm}}_G(T)(S)\) lifting
\(s_\alpha\). For every point \(s\in S\), there exist an open
subscheme \(V\) of \(S\) containing \(s\), an element \(t\in T(V)\),
and, over \(V\), a relation
\[
  w=n_{\alpha_1}\cdots n_{\alpha_p}t,
\]
where the \(\alpha_i\) are simple. We may plainly restrict ourselves
to proving the result when \(V=S\); we argue by induction on \(p\).
If \(p=0\), then \(w\in T(S)\), and we take \(V_w=G\). Suppose,
therefore, that \(w=n_\alpha w'\), where \(w'\) satisfies the
conclusion of the lemma. There is then an open subscheme \(V_{w'}\)
of \(G\), containing the unit section, such that
\[
  \operatorname{int}(w')\Omega'\cap V_{w'}\subset\Omega'.
\]
We may then write
\[
\begin{aligned}
  \operatorname{int}(w)\Omega'
    \cap\operatorname{int}(n_\alpha)V_{w'}\cap\Omega
  &=
  \operatorname{int}(n_\alpha)
  \bigl(\operatorname{int}(w')\Omega'\cap V_{w'}\bigr)\cap\Omega\\
  &\subset
  \operatorname{int}(n_\alpha)\Omega'\cap\Omega
  \subset\Omega',
\end{aligned}
\]
by \SGARefLink{sga3:xxii:lemma:6.1.4}{6.1.4}. We now take
\[
  V_w=\operatorname{int}(n_\alpha)V_{w'}\cap\Omega,
\]
and the proof is complete.

\SGARefTarget{sga3:xxii:lemma:6.1.6}
\paragraph{\SGARefLink{sga3:xxii:lemma:6.1.6}{Lemma 6.1.6}.}
{\itshape There exists an open subscheme \(V_0\) of \(G\),
containing the unit section, such that for every \(S'\to S\),
\[
  U(S')U^-(S')\cap V_0(S')\subset\Omega'(S').
\]}

Let \(n_0\) be an element of
\(\underline{\operatorname{Norm}}_G(T)(S)\) lifting the symmetry
\(w_0\) of the Weyl group,\footnote{Editor's note: The symmetry
\(w_0\) is defined in \SGARefLink{sga3:XXI:3.6.14}{Exposé XXI, 3.6.14}.} that is, such that
\[
  \operatorname{int}(n_0)U=U^-
  \qquad\text{\textup{(cf.\ \SGARefLink{sga3:XXI:3.6.14}{Exposé XXI, 3.6.14})}}.
\]
Then \(n_0^2\in T(S)\). We show that the open subscheme
\(V_0=V_{n_0}\) of \SGARefLink{sga3:xxii:lemma:6.1.5}{6.1.5} has the required property. Indeed,
\[
\begin{aligned}
  U(S')U^-(S')
  &=
  \operatorname{int}(n_0)
  \bigl(
    \operatorname{int}(n_0)^{-1}U(S')
    \cdot\operatorname{int}(n_0)^{-1}U^-(S')
  \bigr)\\
  &=
  \operatorname{int}(n_0)\bigl(U^-(S')\cdot U(S')\bigr)
  \subset\operatorname{int}(n_0)\Omega'(S').
\end{aligned}
\]
Consequently,
\[
  U(S')U^-(S')\cap V_0(S')
  \subset
  \operatorname{int}(n_0)\Omega'(S')\cap V_0(S')
  \subset\Omega'(S').
\]

\SGARefTarget{sga3:xxii:lemma:6.1.7}
\paragraph{\SGARefLink{sga3:xxii:lemma:6.1.7}{Lemma 6.1.7}.}
{\itshape Consider the morphism
\[
  f:\Omega=U^-\cdot T\cdot U\longrightarrow T/T'
\]
obtained by composing the second projection with the canonical
morphism from \(T\) to \(T/T'\). Then \(f\) is
``generically multiplicative'': there exists an open subscheme
\(V\) of \(\Omega\times_S\Omega\), containing the unit section
\textup{(and hence relatively schematically dense,
\SGARefLink{sga3:xviii:cor:1.3}{Exposé XVIII, 1.3})}, such that for every \(S'\to S\) and every
\((x,y)\in V(S')\), one has \(xy\in\Omega(S')\) and
\[
  f(xy)=f(x)f(y).
\]}

Let \(x,y\) be sections of \(\Omega\) over \(S'\). Write
\[
  x=utv,\qquad y=u't'v',
\]
where
\[
  u,u'\in U^-(S'),\qquad t,t'\in T(S'),\qquad v,v'\in U(S').
\]
Let \(V_0\) be the open subscheme of \SGARefLink{sga3:xxii:lemma:6.1.6}{6.1.6}, and let \(V\) be the
open subscheme of \(\Omega\times_S\Omega\) defined by the condition
\[
  vu'\in V_0(S').
\]
This open subscheme has the required property. Indeed, if
\((x,y)\in V(S')\), then
\[
  xy=(utv)(u't'v')=(ut)(vu')(t'v').
\]
Since \(vu'\in\Omega'(S')\), we obtain
\[
  xy
  \in U^-(S')\,t\Omega'(S')t'\,U(S')
  \subset U^-(S')\,tt'T'(S')\,U(S'),
\]
and therefore \(xy\in\Omega(S')\) and
\[
  f(xy)=f(tt')=f(t)f(t')=f(x)f(y).
\]

\SGARefTarget{sga3:xxii:prop:6.1.8}
\paragraph{\SGARefLink{sga3:xxii:prop:6.1.8}{Proposition 6.1.8}.}
\textit{There exists a group morphism}
\[
  f:G\longrightarrow T/T'
\]
\textit{inducing the canonical projection on \(T\). The kernel
\(\operatorname{Ker}(f)\) of \(f\) is a closed subgroup scheme of
\(G\), smooth over \(S\), with connected fibers. Every group
morphism from \(G\) to an \(S\)-presheaf of commutative groups
that is separated for the fppf topology vanishes on
\(\operatorname{Ker}(f)\).}

The first assertion follows at once from \SGARefLink{sga3:xxii:prop:4.1.11}{4.1.11}. We immediately
have
\[
  \operatorname{Ker}(f)\cap\Omega=\Omega',
\]
which proves that \(\operatorname{Ker}(f)\) is smooth over \(S\)
at every point of the unit section.\footnote{Editor's note: We have
added ``at every point of the unit section,'' as well as the
reference to Exposé \(\mathrm{VI_B}\), \SGARefLink{thm:sga3-VIB-3.10}{3.10}, at the end of the
proof. Moreover, in the following sentence ``prescheme'' has been
replaced by ``presheaf.''}
By \SGARefLink{sga3:xxii:prop:5.6.9}{5.6.9}\textup{\SGARefLink{sga3:xxii:prop:5.6.9:item:ii}{(ii)}}, every morphism
\(\phi\) from \(G\) to an \(S\)-presheaf of commutative groups
that is separated for the fppf topology vanishes on \(U\) and on
\(U^-\). By \SGARefLink{sga3:xx:2.7}{Exposé XX, 2.7}, \(\phi\) also vanishes on \(T'\), and
hence on \(\Omega'\). Using the notation of \SGARefLink{sga3:xxii:prop:5.7.10}{5.7.10}, we see that the
monoid \(U_1\) is contained in \(\operatorname{Ker}(f)(S)\), which
shows that
\[
  \operatorname{Ker}(f)
  =
  \bigcup_{u\in U_1}u\Omega'.
\]
It follows, on the one hand, that every \(\phi\) as above vanishes
on \(\operatorname{Ker}(f)\), and, on the other hand, that
\(\operatorname{Ker}(f)\) has connected fibers and is therefore
smooth over \(S\), by Exposé \(\mathrm{VI_B}\), 3.10.

\SGARefTarget{sga3:xxii:section:6.2}
\subsection*{\SGAInlineRefTarget{sga3:xxii:subsection:6-2}6.2. The Derived Group of a Reductive Group}
\addcontentsline{toc}{subsection}{\SGARefLink{sga3:xxii:subsection:6-2}{6.2}. The Derived Group of a
Reductive Group}

\SGARefTarget{sga3:xxii:thm:6.2.1}
\paragraph{\SGARefLink{sga3:xxii:thm:6.2.1}{Theorem 6.2.1}.}
\textit{Let \(S\) be a scheme and \(G\) a reductive \(S\)-group.}
\begin{enumerate}
\item[\textup{\SGAInlineRefTarget{sga3:xxii:thm:6.2.1:item:i}(i)}]
\textit{The group}
\[
  \mathrm D_S(G)
  =
  \underline{\operatorname{Hom}}_{S\text{-gr.}}
    (G,\mathbb G_{m,S})
\]
\textit{is represented by a twisted constant \(S\)-group whose
type at \(s\in S\) is}
\[
  \mathbb Z^{
    \operatorname{rgred}(G_s)-\operatorname{rgss}(G_s)}.
\]

\item[\textup{\SGAInlineRefTarget{sga3:xxii:thm:6.2.1:item:ii}(ii)}]
\textit{Write}
\[
  \operatorname{corad}(G)
  =
  \mathrm D_S\bigl(\mathrm D_S(G)\bigr),
\]
\textit{which is therefore an \(S\)-torus. The biduality morphism
\textup{(cf.\ Exposé VIII, \S1)}}
\[
  f_0:G\longrightarrow\operatorname{corad}(G)
\]
\textit{is smooth and surjective.}

\item[\textup{\SGAInlineRefTarget{sga3:xxii:thm:6.2.1:item:iii}(iii)}]
\textit{The composite morphism}
\[
  \operatorname{rad}(G)\longrightarrow G
  \longrightarrow\operatorname{corad}(G)
\]
\textit{is an isogeny \textup{(cf.\ \SGARefLink{sga3:xxii:def:4.2.9}{4.2.9})}.}

\item[\textup{(iv)}]
\textit{The kernel of \(f_0\), denoted}
\[
  \operatorname{der}(G)=\operatorname{Ker}(f_0),
\]
\textit{is a closed subgroup scheme of \(G\), semisimple over
\(S\), called the derived group of \(G\). If \(G\) is semisimple,
then \(\operatorname{der}(G)=G\).}

\item[\textup{\SGAInlineRefTarget{sga3:xxii:thm:6.2.1:item:v}(v)}]
\textit{Every group morphism from \(G\) to an \(S\)-presheaf of
commutative groups that is separated for the fppf topology
vanishes on \(\operatorname{der}(G)\), and therefore factors
through \(f_0\).}
\end{enumerate}

\textit{Proof.}
All the assertions of the theorem are local for the étale
topology; we may therefore reduce to the case in which \(G\) is
split over \(S\). Consider the morphism \(f\) of \SGARefLink{sga3:xxii:prop:6.1.8}{6.1.8}. By the
last assertion of \SGARefLink{sga3:xxii:prop:6.1.8}{6.1.8}, we immediately have an isomorphism
\[
  \underline{\operatorname{Hom}}_{S\text{-gr.}}
    (G,\mathbb G_{m,S})
  \xrightarrow{\ \sim\ }
  \underline{\operatorname{Hom}}_{S\text{-gr.}}
    (T/T',\mathbb G_{m,S}),
\]
which proves \textup{\SGARefLink{sga3:xxii:thm:6.2.1:item:i}{(i)}}, then \textup{\SGARefLink{sga3:xxii:thm:6.2.1:item:ii}{(ii)}}, and yields a
commutative diagram
\[
  % BEGIN INCLUDED SGA3_English_current_buildable_source_20260823/source/inputs/XXII/native_diagrams/exp22/diag_008_coradical_triangle.tex
\begin{tikzcd}[column sep=5.0em,row sep=3.2em]
  G
    \arrow[r,"f_0"]
    \arrow[dr,"f"']
  &
  \operatorname{corad}(G)
    \arrow[d,"\sim"]
  \\
  &
  T/T'.
\end{tikzcd}

% END INCLUDED SGA3_English_current_buildable_source_20260823/source/inputs/XXII/native_diagrams/exp22/diag_008_coradical_triangle.tex

\]
Assertion \textup{\SGARefLink{sga3:xxii:thm:6.2.1:item:v}{(v)}} now follows from \SGARefLink{sga3:xxii:prop:6.1.8}{6.1.8}, and
\textup{\SGARefLink{sga3:xxii:thm:6.2.1:item:iii}{(iii)}} from \SGARefLink{sga3:xxii:lemma:6.1.2}{6.1.2}. We also have
\(\operatorname{Ker}(f)=\operatorname{Ker}(f_0)\), which, by
\SGARefLink{sga3:xxii:prop:6.1.8}{6.1.8}, implies that \(\operatorname{der}(G)\) is smooth over
\(S\) with connected fibers. It remains to verify that its fibers
are semisimple. They are reductive by \SGARefLink{sga3:xix:paragraph:1-7}{Exposé XIX, 1.7}, as invariant
subgroups of reductive groups. By \textup{\SGARefLink{sga3:xxii:thm:6.2.1:item:iii}{(iii)}},
\(\operatorname{rad}(G)\cap\operatorname{der}(G)\) is finite,
which indeed implies that the fibers of \(\operatorname{der}(G)\)
are semisimple.

\SGARefTarget{sga3:xxii:rem:6.2.2}
\paragraph{\SGARefLink{sga3:xxii:rem:6.2.2}{Remark 6.2.2}.}
\SGAInlineRefTarget{sga3:xxii:rem:6.2.2:item:a}\textup{a)} By construction, when \(G\) is split,
\(\operatorname{der}(G)\) is the fppf subsheaf of \(G\) generated
by the \(U_\alpha\), \(\alpha\in R\). It even suffices to take the
\(U_\alpha\), \(\alpha\in\Delta\cup-\Delta\), where \(\Delta\) is
a base of \(R\).

\textup{b)}\footnote{Editor's note: In what follows, we have
expanded the original and removed the assertion that
``\(\operatorname{der}(G)\) is the fppf-separated presheaf of
commutators of \(G\).''}
Let \(C\) be the presheaf of commutators of \(G\), that is, the
\(S\)-functor in groups that assigns to every \(S'\to S\) the
commutator subgroup of \(G(S')\) (that is, the subgroup of
\(G(S')\) generated by the elements
\(xyx^{-1}y^{-1}\), for \(x,y\in G(S')\)), and let
\(\widetilde C\) be the associated fppf sheaf. Since the quotient
\(G/\operatorname{der}(G)=T/T'\) is commutative,
\(\operatorname{der}(G)\) contains \(C\), and hence
\(\widetilde C\) \textup{(cf.\ \SGARefLink{sga3:iv:corollary:4:3:12}{Exposé IV, 4.3.12})}.

On the other hand, the quotient presheaf \(G/\widetilde C\) is
separated \textup{(\SGARefLink{sga3:iv:lemma:4:4:8:1}{Exposé IV, 4.4.8.1})}, and therefore, by
\textup{(v)}, we have
\(\operatorname{der}(G)\subset\widetilde C\). Thus
\(\operatorname{der}(G)=\widetilde C\); in other words,
\(\operatorname{der}(G)\) is the fppf sheaf of commutators of
\(G\).

Let us finally note that, since \(C\) is a subpresheaf of \(G\),
it is separated, but it is not equal to
\(\operatorname{der}(G)\) in general: for example,
\(\operatorname{der}(\mathrm{SL}_2)=\mathrm{SL}_2\), whereas
\(\mathrm{SL}_2(\mathbb F_2)\simeq\mathfrak S_3\) is not equal
to its commutator subgroup.

\textup{c)} When \(S\) is the spectrum of an algebraically closed
field \(k\), \(\operatorname{der}(G)(k)\) is the commutator
subgroup of \(G(k)\) \textup{(Exposé \(\mathrm{VI_B}\), 7.10)}.

\SGARefTarget{sga3:xxii:par:6.2.3}
\paragraph{\SGARefLink{sga3:xxii:par:6.2.3}{6.2.3}.}
Consider now the two exact sequences
\[
\begin{gathered}
  1\longrightarrow\operatorname{rad}(G)
  \longrightarrow G
  \longrightarrow\operatorname{ss}(G)
  \longrightarrow 1,\\
  1\longrightarrow\operatorname{der}(G)
  \longrightarrow G
  \longrightarrow\operatorname{corad}(G)
  \longrightarrow 1.
\end{gathered}
\]
Since \(\operatorname{rad}(G)\) is central in \(G\),
multiplication in \(G\) defines a group morphism
\[
  u:\operatorname{rad}(G)\times_S\operatorname{der}(G)
  \longrightarrow G,
\]
which is covering by \SGARefLink{sga3:xxii:thm:6.2.1}{6.2.1}\textup{\SGARefLink{sga3:xxii:thm:6.2.1:item:iii}{(iii)}}, and hence surjective and
flat \textup{(Exposé \(\mathrm{VI_B}\), 9.2\textup{(xi)})}.
\footnote{Editor's note: Indeed, \(G\) is the fppf quotient of
\(\operatorname{rad}(G)\times_S\operatorname{der}(G)\) by
\(\operatorname{Ker}(u)\), which is a group of multiplicative
type and hence is flat over \(S\). Therefore, by Exposé
\(\mathrm{VI_B}\), 9.2\textup{(xi)}, \(u\) is flat.}
Its kernel is isomorphic to
\(\operatorname{rad}(G)\cap\operatorname{der}(G)\), which is also
the kernel of
\[
  \operatorname{rad}(G)\longrightarrow\operatorname{corad}(G);
\]
it is a finite subgroup of multiplicative type of
\(\operatorname{rad}(G)\). Similarly, the kernel of the morphism
\[
  G\longrightarrow
  \operatorname{corad}(G)\times_S\operatorname{ss}(G)
\]
is isomorphic to
\(\operatorname{der}(G)\cap\operatorname{rad}(G)\).

\SGARefTarget{sga3:xxii:prop:6.2.4}
\paragraph{\SGARefLink{sga3:xxii:prop:6.2.4}{Proposition 6.2.4}.}
\textit{Let \(G\) be a reductive \(S\)-group. The three morphisms}
\[
\begin{gathered}
  \operatorname{rad}(G)\times_S\operatorname{der}(G)
    \longrightarrow G,\\
  G\longrightarrow
    \operatorname{corad}(G)\times_S\operatorname{ss}(G),\\
  \operatorname{rad}(G)\longrightarrow\operatorname{corad}(G)
\end{gathered}
\]
\textit{are central isogenies, and their kernels are isomorphic.}

\SGARefTarget{sga3:xxii:cor:6.2.5}
\paragraph{\SGARefLink{sga3:xxii:cor:6.2.5}{Corollary 6.2.5}.}
\textit{The following conditions are equivalent:}
\begin{enumerate}
\item[\textup{(i)}]
\textit{\(G\) is the product of a semisimple group and a torus.}

\item[\textup{(ii)}]
\textit{The morphism}
\[
  \operatorname{rad}(G)\times_S\operatorname{der}(G)
  \xrightarrow{\ \sim\ }G
\]
\textit{is an isomorphism.}

\item[\textup{(iii)}]
\textit{The morphism}
\[
  G\xrightarrow{\ \sim\ }
  \operatorname{corad}(G)\times_S\operatorname{ss}(G)
\]
\textit{is an isomorphism.}

\item[\textup{(iv)}]
\textit{\(\operatorname{rad}(G)\cap\operatorname{der}(G)=e\).}
\end{enumerate}

\SGARefTarget{sga3:xxii:par:6.2.6}
\paragraph{\SGARefLink{sga3:xxii:par:6.2.6}{6.2.6}.}
Let us return temporarily to the case of a split group. We retain
the notation of \SGARefLink{sga3:xxii:subsection:6-1}{6.1}. Put
\[
  N=M\cap V(R^*)^\perp,
  \qquad\text{so that}\qquad
  T'=\mathrm D_S(M/N).
\]
Then \(U^-\cdot T'\cdot U\) is an open neighborhood of the unit
section in \(\operatorname{der}(G)\). Consequently,
\[
  \operatorname{Lie}\bigl(\operatorname{der}(G)/S\bigr)
  =
  \mathfrak t'\oplus
  \coprod_{\alpha\in R}\mathfrak g^\alpha.
\]
Since the characters induced on \(T'\) by the
\(\alpha\in R\) are nonzero and distinct
\textup{(\SGARefLink{sga3:XXI:1.2.5}{Exposé XXI, 1.2.5}---which we already used in \SGARefLink{sga3:xxii:lemma:6.1.2}{6.1.2})},
\(R\) is a root system of \(\operatorname{der}(G)\) relative to
\(T'\). It is then immediate \textup{(since
\(U_\alpha\subset\operatorname{der}(G)\))} that the
\(\exp\)-morphisms of \(\operatorname{der}(G)\) ``are'' those of
\(G\), and likewise for the coroots.

It follows that:

\SGARefTarget{sga3:xxii:prop:6.2.7}
\paragraph{\SGARefLink{sga3:xxii:prop:6.2.7}{Proposition 6.2.7}.}
\textit{With the preceding notation,}
\[
  \bigl(\operatorname{der}(G),T',M/N,R\bigr)
\]
\textit{is a split group with root datum
\(\operatorname{der}\bigl(\mathcal R(G)\bigr)\). The canonical
morphism \(\operatorname{der}(G)\to G\) induces, by functoriality,
the canonical morphism of root data}
\[
  \mathcal R(G)\longrightarrow
  \operatorname{der}\bigl(\mathcal R(G)\bigr)
\]
\textit{of \SGARefLink{sga3:XXI:6.5}{Exposé XXI, 6.5}.}

\textit{N.B.} As an exercise, the reader may construct the diagram
of split groups corresponding to the three left-hand columns of
the diagram of root data in \SGARefLink{sga3:XXI:6.5.7}{Exposé XXI, 6.5.7}.

\SGARefTarget{sga3:xxii:prop:6.2.8}
\paragraph{\SGARefLink{sga3:xxii:prop:6.2.8}{Proposition 6.2.8}.}
\textit{Let \(S\) be a scheme, \(G\) a reductive \(S\)-group, and
\(\operatorname{der}(G)\) its derived group.}
\begin{enumerate}
\item[\textup{\SGAInlineRefTarget{sga3:xxii:prop:6.2.8:item:i}(i)}]
\textit{For every maximal torus \(T\) of \(G\),
\(T\cap\operatorname{der}(G)\) is a maximal torus of
\(\operatorname{der}(G)\). For every maximal torus \(T'\) of
\(\operatorname{der}(G)\),}
\[
  \underline{\operatorname{Centr}}_G(T')
  =
  \operatorname{rad}(G)\cdot T'
\]
\textit{is a maximal torus of \(G\). These two constructions are
inverse to one another and establish a bijective correspondence
between the maximal tori of \(G\) and those of
\(\operatorname{der}(G)\).}

\item[\textup{\SGAInlineRefTarget{sga3:xxii:prop:6.2.8:item:ii}(ii)}]
\textit{For every Borel subgroup \(B\) of \(G\),
\(B\cap\operatorname{der}(G)\) is a Borel subgroup \(B'\) of
\(\operatorname{der}(G)\). We have \(B'{}^u=B^u\). For every
Borel subgroup \(B'\) of \(\operatorname{der}(G)\),}
\[
  \underline{\operatorname{Norm}}_G(B')
  =
  \operatorname{rad}(G)\cdot B'
\]
\textit{is a Borel subgroup of \(G\). These constructions are
inverse to one another and establish a bijective correspondence
between the Borel subgroups of \(G\) and those of
\(\operatorname{der}(G)\).}
\end{enumerate}

By the local conjugacy theorem for maximal tori and the
construction of the derived group, the only assertion that
remains to be proved in \textup{\SGARefLink{sga3:xxii:prop:6.2.8:item:i}{(i)}} is the following: if \(T\)
is a maximal torus of \(G\), then
\[
  T
  =
  \bigl(T\cap\operatorname{der}(G)\bigr)\cdot
    \operatorname{rad}(G)
  =
  \underline{\operatorname{Centr}}_G
    \bigl(T\cap\operatorname{der}(G)\bigr).
\]
The first equality is immediate \textup{(since we may reduce to
the split case)}; the second follows at once, because
\(\operatorname{rad}(G)\) is central in \(G\), and hence
\[
  T
  =
  \underline{\operatorname{Centr}}_G(T)
  =
  \underline{\operatorname{Centr}}_G
    \bigl(T\cap\operatorname{der}(G)\bigr).
\]
The argument for \textup{\SGARefLink{sga3:xxii:prop:6.2.8:item:ii}{(ii)}} is the same.

\SGARefTarget{sga3:xxii:section:6.3}
\subsection*{\SGAInlineRefTarget{sga3:xxii:subsection:6-3}6.3. Subgroups with Commutative Quotients}
\addcontentsline{toc}{subsection}{\SGARefLink{sga3:xxii:subsection:6-3}{6.3}. Subgroups with Commutative
Quotients}

\SGARefTarget{sga3:xxii:par:6.3.1}
\paragraph{\SGARefLink{sga3:xxii:par:6.3.1}{6.3.1}.}
Let \(G\) be a reductive \(S\)-group. If \(H\) is a subsheaf in
groups of \(G\), the following conditions are equivalent:
\begin{itemize}
\item \(H\) contains \(\operatorname{der}(G)\);
\item \(H\) is normal and \(G/H\) is commutative.
\end{itemize}
In this case, the canonical morphism
\[
  f_0:G\longrightarrow\operatorname{corad}(G)
\]
maps \(H\) onto a subsheaf \(f_0(H)\) of
\(\operatorname{corad}(G)\), and we have
\[
\begin{gathered}
  G/H\simeq\operatorname{corad}(G)/f_0(H),
  \qquad
  H/\operatorname{der}(H)\simeq f_0(H),\\
  \operatorname{der}(G)=\operatorname{der}(H),
  \qquad
  H=f_0^{-1}\bigl(f_0(H)\bigr).
\end{gathered}
\]

Since \(\operatorname{der}(G)\) is smooth over \(S\) with
connected fibers,\setcounter{footnote}{63}\footnote{Editor's note:
We have expanded the original in what follows. In particular, we
have added the conclusion, implicit in the original, that every
subgroup of \(G\) that is smooth with connected fibers and contains
\(\operatorname{der}(G)\) is closed in \(G\).} it follows from
\SGARefLink{sga3:iv:proposition:5:3:1}{Exposé IV, 5.3.1} and \SGARefLink{sga3:xxii:par:6.3.1}{6.3.1}, and Exposé \(\mathrm{IV_B}\), 9.2,
that the assignment
\[
  H\longmapsto f_0(H)
\]
establishes a bijective correspondence between the group
subschemes \textup{(respectively, the closed group subschemes)} of
\(G\) that contain \(\operatorname{der}(G)\), are smooth over
\(S\), and have connected fibers, and the group subschemes
\textup{(respectively, the closed group subschemes)} of
\(\operatorname{corad}(G)\) that are smooth over \(S\) and have
connected fibers.

Now, if \(H'\) is a group subscheme of
\(\operatorname{corad}(G)\), smooth over \(S\) with connected
fibers, then \(H'\) is of finite presentation over \(S\)
\textup{(Exposé \(\mathrm{VI_B}\), \SGARefLink{sga3:xxii:subsection:5-5}{5.5})}, and its fibers are tori
\textup{(because those of \(\operatorname{corad}(G)\) are)}.
Hence, by \SGARefLink{sga3:x:corollary:8-2}{Exposé X, 8.2}, \(H'\) is a subtorus of
\(\operatorname{corad}(G)\), and is therefore closed in
\(\operatorname{corad}(G)\) \textup{(\SGARefLink{sga3:ix:remark:2:6}{Exposé IX, 2.6})}.

Consequently, every subgroup of \(G\) that is smooth with
connected fibers and contains \(\operatorname{der}(G)\) is closed
in \(G\).

\SGARefTarget{sga3:xxii:par:6.3.2}
\paragraph{\SGARefLink{sga3:xxii:par:6.3.2}{6.3.2}.}
If \(H\) is a closed group subscheme of \(G\), smooth over \(S\),
with connected fibers, and normal in \(G\), then \(H\) is
reductive. If, moreover, \(H\supset\operatorname{der}(G)\), then
\[
  \operatorname{der}(H)=\operatorname{der}(G)
  \qquad\text{and}\qquad
  f_0(H)\simeq\operatorname{corad}(H).
\]
We have therefore proved the following proposition.

\SGARefTarget{sga3:xxii:prop:6.3.3}
\paragraph{\SGARefLink{sga3:xxii:prop:6.3.3}{Proposition 6.3.3}.}
\textit{Let \(G\) be a reductive \(S\)-group. Every group
subscheme \(H\) of \(G\) that is normal in \(G\), has commutative
quotient \textup{(that is, contains
\(\operatorname{der}(G)\))}, is smooth over \(S\), and has
connected fibers}\footnote{Editor's note: We have deleted the
hypothesis that \(H\) be retrocompact in \(G\), which is automatic:
by Exposé \(\mathrm{VI_B}\), \SGARefLink{sga3:xxii:subsection:5-5}{5.5}, both \(G\) and \(H\) are
separated and quasi-compact over \(S\); hence
\(H\hookrightarrow G\) is quasi-compact by EGA
\(\mathrm{IV}_1\), 1.1.2\textup{(v)}.}
\textit{is closed and reductive. We have}
\[
  \operatorname{der}(H)=\operatorname{der}(G),
  \qquad
  f_0(H)\simeq\operatorname{corad}(H),
\]
\textit{and}
\[
  G/H\simeq
  \operatorname{corad}(G)/\operatorname{corad}(H),
  \qquad
  H=
  \bigl(H\cap\operatorname{rad}(G)\bigr)
  \cdot\operatorname{der}(G).
\]
\textit{Moreover, the assignment \(H\mapsto f_0(H)\) defines a
bijection between the set of subgroups \(H\) of \(G\) having the
preceding properties and the set of subtori of
\(\operatorname{corad}(G)\).}\footnote{Editor's note: The French
source prints \(\operatorname{corad}(H)\) here; the codomain of
\(f_0\) and the preceding correspondence require
\(\operatorname{corad}(G)\).}

By another application of Noether's theorem \textup{(\SGARefLink{sga3:iv:proposition:5:3:1}{Exposé IV,
5.3.1} and \SGARefLink{sga3:xxii:par:6.3.1}{6.3.1})}, we deduce the following.

\SGARefTarget{sga3:xxii:prop:6.3.4}
\paragraph{\SGARefLink{sga3:xxii:prop:6.3.4}{Proposition 6.3.4}.}
\textit{Let \(S\) be a scheme, \(G\) a reductive \(S\)-group, and
\(T\) a maximal torus of \(G\). For every subgroup \(H\) of \(G\)
as above, \(T\cap H\) is a maximal torus of
\(H\),}\footnote{Editor's note: The French source prints \(G\)
here; since \(T\cap H\subset H\), and by the formulas that follow,
the intended group is \(H\).}
\textit{and}
\[
  G/H\simeq T/(T\cap H),
  \qquad
  H=(T\cap H)\cdot\operatorname{der}(G).
\]
\textit{Moreover, the assignment \(H\mapsto T\cap H\) is a
bijection between the set of subgroups \(H\) of \(G\) as above
and the set of subtori of \(T\) containing
\(T\cap\operatorname{der}(G)\).}

\clearpage
\SGARefTarget{sga3:xxii:bibliography}
\section*{Bibliography}
\addcontentsline{toc}{section}{Bibliography}

\begin{description}
\SGARefTarget{sga3:xxii:bib:BLie}
\item[{[BLie]}]
N.~Bourbaki, \emph{Groupes et algèbres de Lie}, Chapters
IV--VI, Hermann, 1968.

\SGARefTarget{sga3:xxii:bib:Ch05}
\item[{[Ch05]}]
C.~Chevalley, \emph{Classification des groupes algébriques
semi-simples} \textup{(with the collaboration of P.~Cartier,
A.~Grothendieck, and M.~Lazard)}, \emph{Collected Works},
vol.~3, Springer, 2005.

\SGARefTarget{sga3:xxii:bib:DG70}
\item[{[DG70]}]
M.~Demazure and P.~Gabriel, \emph{Groupes algébriques},
Masson \& North-Holland, 1970.

\SGARefTarget{sga3:xxii:bib:Gi71}
\item[{[Gi71]}]
J.~Giraud, \emph{Cohomologie non abélienne},
Springer-Verlag, 1971.

\SGARefTarget{sga3:xxii:bib:Se64}
\item[{[Se64]}]
J.-P.~Serre, \emph{Cohomologie galoisienne},
Springer-Verlag, 1964; 5th ed., 1994.
\end{description}

% END INCLUDED SGA3_English_current_buildable_source_20260823/source/inputs/XXII/tex/components/01_expose_XXII_title_section1.tex



% END INCLUDED SGA3_English_current_buildable_source_20260823/source/inputs/XXII/body.tex


\SGAStartExpose{XXIII}
\renewcommand{\SGAProof}{%
  \par\smallskip\noindent\textit{Proof.}\quad
}
% BEGIN INCLUDED SGA3_English_current_buildable_source_20260823/source/inputs/XXIII/body.tex
% Mechanically extracted after \mainmatter from the frozen standalone master.

\markboth{Exposé XXIII. Uniqueness of Pinned Groups}
         {Exposé XXIII. Uniqueness of Pinned Groups}
% BEGIN INCLUDED SGA3_English_current_buildable_source_20260823/source/inputs/XXIII/tex/components/00_expose_XXIII_title_and_definition11.tex
\chapter*{Exposé XXIII\\Reductive Groups: Uniqueness of Pinned Groups}
\addcontentsline{toc}{chapter}{Exposé XXIII. Uniqueness of Pinned Groups}
\begin{center}
\textit{by M. Demazure}
\end{center}

\NDE{Version of 13 October 2024.}
The aim of this Exposé is to prove the uniqueness theorem
(Theorem~4.1). It was proved by Chevalley in the case of an
algebraically closed field; the reduction-to-rank-two method used here
is likewise due to Chevalley (see \textit{Bible}, Exposés~23 and~24).
Along the way, we obtain an explicit description of reductive groups
by generators and relations (3.5).

\SGANextSectionTarget{sga3:xxiii:section:1}\SGASection{1}{Pinnings}

\SGARefTarget{sga3:XXIII:1.1}
\SGAStatement{Definition 1.1}{Let \(S\) be a scheme and let
\((G,T,M,R)\) be an \(S\)-split group (XXII, \SGARefLink{sga3:xxii:def:1.13}{1.13}). A \emph{pinning}
\NDE{Demazure tells us that behind this terminology lies the image of
the butterfly (supplied to him by Grothendieck): the body is a maximal
torus \(T\), the wings are two Borel subgroups opposite with respect to
\(T\); one spreads the butterfly by laying out the wings and then fixes
elements in the additive groups (pins) to rigidify the situation,
that is, to eliminate the automorphisms.}
of this split group is the datum of a system \(\Delta\) of simple
roots of \(R\), together with, for each \(\alpha\in\Delta\), a section
\[
  X_\alpha\in\Gamma(S,\mathfrak g^\alpha)^\times .
\]}

In other words, a pinning of the reductive group \(G\) over the
nonempty scheme \(S\) is the datum of:
\begin{enumerate}
  \item[(i)] a maximal torus \(T\);
  \item[(ii)] an abelian group \(M\) and an isomorphism
    \(T\simeq D_S(M)\);
  \item[(iii)] a system of roots \(R\) of \(G\) with respect to \(T\);
  \item[(iv)] a system \(\Delta\) of simple roots of \(R\);
  \item[\SGAInlineRefTarget{sga3:XXIII:1.1:item:v}(v)] an element
    \(X_\alpha\in\Gamma(S,\mathfrak g^\alpha)^\times\) for every
    \(\alpha\in\Delta\), that is, an element
    \[
      u_\alpha=\exp_\alpha(X_\alpha)\in U_\alpha^\times(S)
      \qquad\text{for every }\alpha\in\Delta,
    \]
\end{enumerate}
satisfying condition (D1) of \SGARefLink{sga3:xxii:def:1.13}{Exposé~XXII, 1.13} (indeed, condition
(D2) of \textit{loc.\ cit.} is automatic\NDE{It is implied by
condition \SGARefLink{sga3:XXIII:1.1:item:v}{(v)}, namely the existence of a section
\(X_\alpha\in\Gamma(S,\mathfrak g^\alpha)^\times\).}).

Every split group admits a pinning; in particular, every reductive
group is locally pinnable for the étale topology.

% END INCLUDED SGA3_English_current_buildable_source_20260823/source/inputs/XXIII/tex/components/00_expose_XXIII_title_and_definition11.tex

% BEGIN INCLUDED SGA3_English_current_buildable_source_20260823/source/inputs/XXIII/tex/components/01_expose_XXIII_par12_and_definition13.tex
\SGARefTarget{sga3:XXIII:1.2}
\SGAParagraph{1.2}
If \(G\) is a pinned \(S\)-group, that is, an \(S\)-split group
equipped with a pinning, it is canonically endowed with the system of
positive roots \(R^+\) defined by \(\Delta\), the corresponding Borel
subgroup \(B=B_{R^+}\), the opposite Borel subgroup \(B^-=B_{R^-}\),
the unipotent groups \(U=B_u\) and \(U^-=(B^-)_{u}\), the open
subscheme \(U^-\cdot T\cdot U\), and so forth. Likewise, for each
\(\alpha\in\Delta\), there is a canonical isomorphism of vector groups
\[
  p_\alpha:\mathbb G_{a,S}\xrightarrow{\ \sim\ }U_\alpha,
  \qquad x\longmapsto \exp_\alpha(xX_\alpha)=u_\alpha^x,
\]
normalized by \(T\) with multiplier \(\alpha\), and whose datum is
equivalent to that of \(X_\alpha\) (\SGARefLink{sga3:xxii:thm:1.1}{Exposé~XXII, 1.1}).

By duality one obtains an
\(X_{-\alpha}\in\Gamma(S,\mathfrak g^{-\alpha})^\times\) and an
isomorphism
\[
  p_{-\alpha}:\mathbb G_{a,S}\xrightarrow{\ \sim\ }U_{-\alpha}
\]
which is contragredient to the preceding one (\SGARefLink{sga3:xxii:not:1.3}{Exposé~XXII, 1.3}). We
set (\SGARefLink{sga3:xx:3.1}{Exposé~XX, 3.1})
\[
  w_\alpha=w_\alpha(X_\alpha)
  =p_\alpha(1)p_{-\alpha}(-1)p_\alpha(1)
  =p_{-\alpha}(-1)p_\alpha(1)p_{-\alpha}(-1).
\]
Then (\textit{loc.\ cit.}, \SGARefLink{sga3:xxiii:subsection:3-1}{3.1} and \SGARefLink{sga3:xx:3.7}{3.7})
\[
  w_\alpha^2=\alpha^\ast(-1),\qquad
  \operatorname{int}(w_\alpha)t
    =s_\alpha(t)=t\cdot\alpha^\ast\!\bigl(\alpha(t)^{-1}\bigr),
\]
\[
\begin{aligned}
  \operatorname{int}(w_\alpha)p_\alpha(x)
    &=p_{-\alpha}(-x)=p_{-\alpha}(x)^{-1},&
  \operatorname{Ad}(w_\alpha)X_\alpha&=-X_{-\alpha},\\
  \operatorname{int}(w_\alpha)p_{-\alpha}(x)
    &=p_\alpha(-x)=p_\alpha(x)^{-1},&
  \operatorname{Ad}(w_\alpha)X_{-\alpha}&=-X_\alpha .
\end{aligned}
\]
We shall systematically use the preceding notation in what follows.

\SGARefTarget{sga3:XXIII:1.3}
\SGAStatement{Definition 1.3}{Let \(S\) be a scheme, and suppose that
\[
  (G,T,M,R,\Delta,(X_\alpha))
  \quad\text{and}\quad
  (G',T',M',R',\Delta',(X'_{\alpha'}))
\]
are two pinned \(S\)-groups.
The morphism of split groups (\SGARefLink{sga3:xxii:def:4.2.1}{Exposé~XXII, 4.2.1})
\[
  f:G\longrightarrow G'
\]
is said to be \emph{compatible with the pinnings}, or to define a
\emph{morphism of pinned groups}, if the associated bijection
\(d:R\to R'\) (cf.\ \textit{loc.\ cit.}) satisfies
\(d(\Delta)=\Delta'\) and if, for every \(\alpha\in\Delta\),
\[
  f\bigl(\exp_\alpha(X_\alpha)\bigr)
    =\exp_{d(\alpha)}\bigl(X'_{d(\alpha)}\bigr),
  \qquad\text{i.e.}\qquad
  f(u_\alpha)=u'_{d(\alpha)} .
\]}
\NDE{The bijection \(R\xrightarrow{\sim}R'\) is denoted by \(d\)
(rather than by \(u\)) in order to avoid the notation
\(u'_{u(\alpha)}\).}

% END INCLUDED SGA3_English_current_buildable_source_20260823/source/inputs/XXIII/tex/components/01_expose_XXIII_par12_and_definition13.tex

% BEGIN INCLUDED SGA3_English_current_buildable_source_20260823/source/inputs/XXIII/tex/components/02_expose_XXIII_par14_through_scholium16.tex
\SGARefTarget{sga3:XXIII:1.4}
\SGAParagraph{1.4}
If \(q(\alpha)\) denotes the integer of \textit{loc.\ cit.}, then
\[
  f\bigl(p_\alpha(x)\bigr)
    =p'_{d(\alpha)}\bigl(x^{q(\alpha)}\bigr)
    \qquad (\alpha\in\Delta),
\]
and hence
\[
  f\bigl(p_{-\alpha}(x)\bigr)
    =p'_{-d(\alpha)}\bigl(x^{q(\alpha)}\bigr),
  \qquad
  f(w_\alpha)=w'_{d(\alpha)} .
\]
Recall (\SGARefLink{sga3:xxii:subsection:4-2}{Exposé~XXII, 4.2}) that, for every \(\alpha\in R\), every
\(z\in\mathbb G_m(S')\), and every \(t\in T(S')\),
\[
  f\bigl(\alpha^\ast(z)\bigr)
    =d(\alpha)^\ast(z)^{q(\alpha)}
    =d(\alpha)^\ast\bigl(z^{q(\alpha)}\bigr),
  \qquad
  d(\alpha)\bigl(f(t)\bigr)=\alpha(t)^{q(\alpha)} .
\]

\SGARefTarget{sga3:XXIII:1.5}
\SGAParagraph{1.5}
We call a root datum equipped with a system of simple roots a
\emph{pinned root datum}, and a \(p\)-morphism of root data
(\SGARefLink{sga3:XXI:6.8}{Exposé~XXI, 6.8}) that sends simple roots to simple roots a
\emph{\(p\)-morphism of pinned root data}.

If \(G\) is a pinned \(S\)-group, we denote by \(\mathscr R(G)\) the
corresponding pinned root datum (the root datum of \SGARefLink{sga3:xxii:prop:1.14}{Exposé~XXII, 1.14}
equipped with \(\Delta\)). Let \(p\) be the integer defined in
\SGARefLink{sga3:xxii:rem:4.2.2}{Exposé~XXII, 4.2.2}. Then:

\SGARefTarget{sga3:XXIII:1.6}
\SGAStatement{Scholium 1.6}{The correspondence
\(G\longmapsto\mathscr R(G)\) defines a functor from the category of
pinned reductive \(S\)-groups to the category of pinned root data
(whose morphisms are the \(p\)-morphisms).}

% END INCLUDED SGA3_English_current_buildable_source_20260823/source/inputs/XXIII/tex/components/02_expose_XXIII_par14_through_scholium16.tex

% BEGIN INCLUDED SGA3_English_current_buildable_source_20260823/source/inputs/XXIII/tex/components/03_expose_XXIII_par17_through_remark193.tex
\SGASubsection{1.7}{The pinned groups \(Z_{\Delta_1}\)}
\SGARefTarget{sga3:XXIII:1.7}

Let \(\Delta_1\) be a subset of the system \(\Delta\) of simple roots
of the pinned group \(G\). Let \(T_{\Delta_1}\) be the maximal torus
of \(\bigcap_{\alpha\in\Delta_1}\ker(\alpha)\), and set
\[
  Z_{\Delta_1}=\operatorname{Centr}_G(T_{\Delta_1}).
\]
Put \(R'=\mathbb Z\Delta_1\cap R\). It is known
(\SGARefLink{sga3:xxii:para:5.10.7}{Exposé~XXII, 5.10.7}) that \(Z_{\Delta_1}\) is a reductive \(S\)-group
with radical \(T_{\Delta_1}\), that \((T,M,R')\) is a splitting of
this group, and that \(\Delta_1\) is a system of simple roots. Thus
\[
  (Z_{\Delta_1},T,M,R',\Delta_1,(X_\alpha)_{\alpha\in\Delta_1})
\]
is a pinned \(S\)-group. We shall always give \(Z_{\Delta_1}\) this
pinning. In particular, we shall consider
\[
  Z_\alpha=Z_{\{\alpha\}},\qquad
  Z_{\alpha\beta}=Z_{\{\alpha,\beta\}}.
\]
Write \(B_{\Delta_1}=B\cap Z_{\Delta_1}\). It is known
(\textit{loc.\ cit.}) that this is the canonical Borel subgroup
\NDE{That is, the Borel subgroup of \(Z_{\Delta_1}\) corresponding to
\(R'^+=R'\cap R^+\).}
of \(Z_{\Delta_1}\), and that its unipotent part is
\(U_{\Delta_1}=U\cap Z_{\Delta_1}\). In particular,
\[
  B_\alpha=T\cdot U_\alpha.
\]
We write
\[
  U_{\alpha\beta}=U_{\{\alpha,\beta\}}
  =U\cap Z_{\alpha\beta}
  =\prod_{\gamma\in R_{\alpha\beta}^{+}}U_\gamma,
\]
where \(R_{\alpha\beta}^{+}\) is the set of positive roots that are
linear combinations of \(\alpha\) and \(\beta\).

Now let \(f:G\to G'\) be a morphism of pinned \(S\)-groups. If
\(d:R\to R'\) is the corresponding bijection and
\(\Delta_1\subset\Delta\), then
\(d(\Delta_1)=\Delta'_1\subset\Delta'\). Clearly \(f\) maps
\(T_{\Delta_1}\) into \(T'_{\Delta'_1}\), and hence
\(Z_{\Delta_1}\) into \(Z'_{\Delta'_1}\). The corresponding
\(S\)-morphism
\[
  f_{\Delta_1}:Z_{\Delta_1}\longrightarrow Z'_{\Delta'_1}
\]
is compatible with the canonical pinnings. It defines a morphism of
pinned root data
\[
  \mathscr R(f_{\Delta_1}):
  \mathscr R(Z_{\Delta_1},T,M,\ldots)
  \longrightarrow
  \mathscr R(Z'_{\Delta'_1},T',M',\ldots),
\]
and the underlying morphisms \(M'\to M\) of \(\mathscr R(f)\) and
\(\mathscr R(f_{\Delta_1})\) coincide.

\SGASubsection{1.8}{Study of the group
\(\operatorname{Norm}_G(T)\)}
\SGARefTarget{sga3:XXIII:1.8}

For every pair \((\alpha,\beta)\) of simple roots, let
\(n_{\alpha\beta}\) be the order of \(s_\alpha s_\beta\) in the Weyl
group \(W\). In particular \(n_{\alpha\alpha}=1\), and therefore
\((w_\alpha w_\beta)^{n_{\alpha\beta}}\in T(S)\).

\SGARefTarget{sga3:XXIII:1.8.1}
\SGAStatement{Definition 1.8.1}{For every
\((\alpha,\beta)\in\Delta\times\Delta\), set
\[
  t_{\alpha\beta}=(w_\alpha w_\beta)^{n_{\alpha\beta}}\in T(S).
\]
Moreover, set (\SGARefLink{sga3:xx:3.1}{Exposé~XX, 3.1})
\[
  t_\alpha=t_{\alpha\alpha}=w_\alpha^2=\alpha^\ast(-1)\in T(S).
\]}

\SGARefTarget{sga3:XXIII:1.8.2}
\SGAStatement{Proposition 1.8.2}{Let \(H\) be an \(S\)-functor in
groups that carries direct sums of schemes to products (for example,
a sheaf for the Zariski topology). Let
\[
  f_T:T\longrightarrow H
\]
be a morphism of groups, and let \(h_\alpha\in H(S)\)
(\(\alpha\in\Delta\)). There exists a necessarily unique morphism of
groups
\[
  f_N:\operatorname{Norm}_G(T)\longrightarrow H
\]
inducing \(f_T\) on \(T\) and satisfying \(f_N(w_\alpha)=h_\alpha\)
for \(\alpha\in\Delta\) if and only if the following two conditions
hold:
\begin{enumerate}
  \item[\SGAInlineRefTarget{sga3:XXIII:1.8.2:item:i}(i)] for every \(\alpha\in\Delta\),
    \[
      f_T(s_\alpha(t))=h_\alpha f_T(t)h_\alpha^{-1}
    \]
    for every \(t\in T(S')\), \(S'\to S\); equivalently,
    \(f_T\circ s_\alpha=\operatorname{int}(h_\alpha)\circ f_T\);
  \item[\SGAInlineRefTarget{sga3:XXIII:1.8.2:item:ii}(ii)] for every
    \((\alpha,\beta)\in\Delta\times\Delta\),
    \[
      f_T(t_{\alpha\beta})
        =(h_\alpha h_\beta)^{n_{\alpha\beta}}.
    \]
\end{enumerate}}

\SGAProof
Equip \((\mathrm{Sch})\) with the topology \(\mathcal T\) generated
by the pretopology whose covering families are direct sums. The
hypothesis says that \(H\) is a \(\mathcal T\)-sheaf. Let \(L\) be
the free group on the generators \((m_\alpha)_{\alpha\in R}\), and
let \(L_1\) be the invariant subgroup generated by
\[
  (m_\alpha m_\beta)^{n_{\alpha\beta}},
  \qquad(\alpha,\beta)\in\Delta\times\Delta.
\]
Let \(g:L\to W\) be defined by \(g(m_\alpha)=s_\alpha\). It is known
(\SGARefLink{sga3:XXI:5.1}{Exposé~XXI, 5.1}) that \(g\) induces an isomorphism
\(\overline g:L/L_1\xrightarrow{\sim}W\). Let \(L\) act on \(T\)
through \(g\), equivalently by \(m_\alpha\cdot t=s_\alpha(t)\).
Writing \(L_S\) for the constant group defined by \(L\), form the
semidirect product \(N=T\cdot L_S\). There is a unique morphism of
\(S\)-groups
\[
  h:N=T\cdot L_S\longrightarrow\operatorname{Norm}_G(T)
\]
such that \(h(m_\alpha)=w_\alpha\) and \(h(t)=t\) for every
\(t\in T(S')\), \(S'\to S\).

Let \(N_1\) be the invariant subgroup sheaf of \(N\) generated by
\[
  t_{\alpha\beta}^{-1}(m_\alpha m_\beta)^{n_{\alpha\beta}},
  \qquad(\alpha,\beta)\in\Delta\times\Delta.
\]
Clearly \(N_1\subset\ker h\), so \(h\) induces
\[
  h_1:N/N_1\longrightarrow\operatorname{Norm}_G(T).
\]
We prove that \(h_1\) is an isomorphism. Since \(h|_T\) is the
canonical immersion, hence a monomorphism, the canonical morphism
\(T\to N/N_1\) is also a monomorphism and induces an isomorphism
\(T\xrightarrow{\sim}TN_1/N_1\). Thus \(h_1\) induces an
isomorphism \(TN_1/N_1\xrightarrow{\sim}T\). It remains only to show
that
\[
  h_2:N/TN_1\longrightarrow\operatorname{Norm}_G(T)/T
\]
is an isomorphism. But \(TN_1\) is the invariant
\(\mathcal T\)-subgroup sheaf of \(N\) generated by \(T\) and the
\((m_\alpha m_\beta)^{n_{\alpha\beta}}\); hence it is the subgroup
sheaf generated by \(T\) and \(L_1\), namely \(T\cdot(L_1)_S\).
Therefore \(h_2\) identifies with
\[
  \overline g:L_S/(L_1)_S\longrightarrow W_S,
\]
which is an isomorphism by construction.

The proof of the proposition is now immediate. Necessity is clear.
Conversely, condition \SGARefLink{sga3:XXIII:1.8.2:item:i}{(i)} gives a morphism
\[
  u:N\longrightarrow H
\]
such that \(u(m_\alpha)=h_\alpha\) for \(\alpha\in\Delta\) and
\(u|_T=f_T\). Condition \SGARefLink{sga3:XXIII:1.8.2:item:ii}{(ii)} says that \(u\) vanishes on \(N_1\),
and the result follows.

\SGASubsection{1.9}{Faithfulness of the functor \(\mathscr R\)}
\SGARefTarget{sga3:XXIII:1.9}

\SGARefTarget{sga3:XXIII:1.9.1}
\SGAStatement{Proposition 1.9.1}{The functor \(\mathscr R\) of \SGARefLink{sga3:XXIII:1.6}{1.6} is
faithful: if
\[
  f,g:G\rightrightarrows G'
\]
are two morphisms of pinned groups such that
\(\mathscr R(f)=\mathscr R(g)\), then \(f=g\).}

Indeed, \(f\) and \(g\) agree on \(T\), on every \(U_\alpha\)
(\(\alpha\in\Delta\)), and on every \(U_{-\alpha}\)
(\(\alpha\in\Delta\)); it therefore suffices to apply:

\SGARefTarget{sga3:XXIII:1.9.2}
\SGAStatement{Lemma 1.9.2}{Let \(S\) be a scheme, let \(G\) be a
pinned reductive \(S\)-group, and let \(H\) be an \(S\)-presheaf in
groups separated for the fppf topology. Let
\[
  f,g:G\rightrightarrows H
\]
be two morphisms of \(S\)-groups. The following conditions are
equivalent:
\begin{enumerate}
  \item[\SGAInlineRefTarget{sga3:XXIII:1.9.2:item:i}(i)] \(f=g\);
  \item[\SGAInlineRefTarget{sga3:XXIII:1.9.2:item:ii}(ii)] \(f\) and \(g\) agree on \(T\), on every \(U_\alpha\)
    (\(\alpha\in\Delta\)), and on every \(U_{-\alpha}\)
    (\(\alpha\in\Delta\));
  \item[\SGAInlineRefTarget{sga3:XXIII:1.9.2:item:iii}(iii)] \(f\) and \(g\) agree on \(T\) and on every
    \(U_\alpha\) (\(\alpha\in\Delta\)), and
    \(f(w_\alpha)=g(w_\alpha)\) for every \(\alpha\in\Delta\).
\end{enumerate}}

\SGAProof
The implication \SGARefLink{sga3:XXIII:1.9.2:item:i}{(i)}\(\Rightarrow\)\SGARefLink{sga3:XXIII:1.9.2:item:ii}{(ii)} is trivial, while
\SGARefLink{sga3:XXIII:1.9.2:item:ii}{(ii)}\(\Rightarrow\)\SGARefLink{sga3:XXIII:1.9.2:item:iii}{(iii)} follows immediately from the definition of
\(w_\alpha\) in \SGARefLink{sga3:XXIII:1.2}{1.2}. It remains to prove
\SGARefLink{sga3:XXIII:1.9.2:item:iii}{(iii)}\(\Rightarrow\)\SGARefLink{sga3:XXIII:1.9.2:item:i}{(i)}. If \(\alpha\in R\), there is a sequence
\(\{\alpha_i\}\subset\Delta\) such that
\[
  \alpha=s_{\alpha_1}s_{\alpha_2}\cdots
    s_{\alpha_n}(\alpha_{n+1}),
\]
and hence
\[
  U_\alpha
    =\operatorname{int}(w_{\alpha_1}\cdots w_{\alpha_n})
       U_{\alpha_{n+1}} .
\]
Thus \(f\) and \(g\) agree on every \(U_\alpha\). They consequently
agree on \(\Omega\), and therefore agree everywhere
(\SGARefLink{sga3:xxii:prop:4.1.11}{Exposé~XXII, 4.1.11}).

\SGARefTarget{sga3:XXIII:1.9.3}
\SGAStatement{Remark 1.9.3}{If \(G\) is semisimple, the hypothesis
that \(f\) and \(g\) agree on \(T\) may be omitted in \SGARefLink{sga3:XXIII:1.9.2:item:ii}{(ii)} and \SGARefLink{sga3:XXIII:1.9.2:item:iii}{(iii)}.
Indeed, \(G\), as an fppf sheaf, is generated by the \(U_\alpha\),
\(\alpha\in R\) (\SGARefLink{sga3:xxii:rem:6.2.2:item:a}{Exposé~XXII, 6.2.2(a)}).}

% END INCLUDED SGA3_English_current_buildable_source_20260823/source/inputs/XXIII/tex/components/03_expose_XXIII_par17_through_remark193.tex

% BEGIN INCLUDED SGA3_English_current_buildable_source_20260823/source/inputs/XXIII/tex/components/04_expose_XXIII_theorem21_through_corollary24.tex
\SGASection{2}{Generators and relations for a pinned group}
\SGARefTarget{sga3:XXIII:2}

Throughout this section \(G\) is a fixed pinned \(S\)-group. For
\(\alpha,\beta\in\Delta\), we use the notation
\(Z_\alpha,Z_{\alpha\beta},U_{\alpha\beta},R^+_{\alpha\beta}\)
of \SGARefLink{sga3:XXIII:1.7}{1.7}.

\SGARefTarget{sga3:XXIII:2.1}
\SGAStatement{Theorem 2.1}{Let \(S\) be a scheme, let \(G\) be a
pinned \(S\)-group, and let \(H\) be an fppf \(S\)-sheaf in groups.
Suppose given morphisms of groups
\[
  f_N:\operatorname{Norm}_G(T)\longrightarrow H,\qquad
  f_\alpha:U_\alpha\longrightarrow H\quad(\alpha\in R).
\]
There exists a necessarily unique morphism of groups
\[
  f:G\longrightarrow H
\]
extending \(f_N\) and all the \(f_\alpha\), \(\alpha\in R\), if and
only if the following conditions hold:
\begin{enumerate}
  \item[\SGAInlineRefTarget{sga3:XXIII:2.1:item:i}(i)] for every \(\alpha\in\Delta\) and \(\beta\in R\),
    \[
      \operatorname{int}\bigl(f_N(w_\alpha)\bigr)\circ f_\beta
        =f_{s_\alpha(\beta)}\circ\operatorname{int}(w_\alpha);
    \]
  \item[\SGAInlineRefTarget{sga3:XXIII:2.1:item:ii}(ii)] for every \(\alpha\in\Delta\), there exists a morphism
    of groups
    \[
      F_\alpha:Z_\alpha\longrightarrow H
    \]
    extending \(f_\alpha\), \(f_{-\alpha}\), and
    \(f_N|_{\operatorname{Norm}_{Z_\alpha}(T)}\);
  \item[\SGAInlineRefTarget{sga3:XXIII:2.1:item:iii}(iii)] for every
    \((\alpha,\beta)\in\Delta\times\Delta\), there exists a morphism
    of groups \(U_{\alpha\beta}\to H\) inducing \(f_\gamma\) on
    \(U_\gamma\) for every \(\gamma\in R^+_{\alpha\beta}\)
    (so \(U_\gamma\subset U_{\alpha\beta}\)).
\end{enumerate}}

\SGASubsection{2.1.1}{Proof}
\SGARefTarget{sga3:XXIII:2.1.1}

The conditions are plainly necessary. Choose a totally ordered group
structure on \(\Gamma_0(R)\) for which the positive elements of \(R\)
are precisely \(R^+\) (\SGARefLink{sga3:XXI:3.5.6}{Exposé~XXI, 3.5.6}); every product indexed by a
subset of \(R\) will be taken in this order. Let \(f_T\) denote the
restriction of \(f_N\) to \(T\), and define a morphism of sets
\[
  f:\Omega\longrightarrow H
\]
by
\[
  f\left(
    \prod_{\alpha\in R^-}y_\alpha\cdot t\cdot
    \prod_{\alpha\in R^+}x_\alpha
  \right)
  =
  \prod_{\alpha\in R^-}f_\alpha(y_\alpha)\cdot f_T(t)\cdot
  \prod_{\alpha\in R^+}f_\alpha(x_\alpha).
\]
Every morphism satisfying the theorem must extend this map. Conversely,
any morphism of groups \(f:G\to H\) extending it also extends \(f_N\):
since it extends \(f_\alpha\) and \(f_{-\alpha}\), one has
\[
  f(w_\alpha)=F_\alpha(w_\alpha)=f_N(w_\alpha),
\]
and it extends \(f_T\) by hypothesis. By \SGARefLink{sga3:xxii:prop:4.1.11:item:ii}{Exposé~XXII, 4.1.11(ii)}, it
therefore remains to prove:

\SGARefTarget{sga3:XXIII:2.1.2}
\SGAStatement{Proposition 2.1.2}{The morphism
\(f:\Omega\to H\) defined above is \emph{generically multiplicative}:
for every \(S'\to S\) and every \(x,y\in\Omega(S')\) such that
\(xy\in\Omega(S')\), one has \(f(xy)=f(x)f(y)\).}

\SGARefTarget{sga3:XXIII:2.1.3}
\SGAStatement{Lemma 2.1.3}{Let
\(n\in\operatorname{Norm}_G(T)(S')\), \(S'\to S\), and let
\(\alpha,\beta\in R\) satisfy
\(\operatorname{int}(n)U_\alpha=U_\beta\), equivalently
\(n(\alpha)=\beta\). Then
\[
  \operatorname{int}\bigl(f_N(n)\bigr)\circ f_\alpha
    =f_\beta\circ\operatorname{int}(n).
\]}

It suffices to verify this formula for a system of generators of
\(\operatorname{Norm}_G(T)\). It holds for every \(w_\alpha\),
\(\alpha\in\Delta\), by \SGARefLink{sga3:XXIII:2.1:item:i}{2.1(i)}, so it remains to treat
\(n\in T(S')\). This is immediate from \SGARefLink{sga3:XXIII:2.1:item:ii}{2.1(ii)} when \(\alpha\) is
simple. Otherwise choose \(w\in W\) with
\(w^{-1}(\alpha)\in\Delta\). Writing \(w\) as a product of simple
reflections\NDE{``Fundamental symmetries'' was replaced by ``simple
reflections.''}, one reduces to showing that if the formula holds for
\(\alpha\) and every \(n\), it also holds for
\(w_{\alpha_0}(\alpha)\) and \(t\in T(S')\), where
\(\alpha_0\in\Delta\). By \SGARefLink{sga3:XXIII:2.1:item:i}{2.1(i)},
\[
\begin{aligned}
  \operatorname{int}\bigl(f_N(t)\bigr)\circ f_{w_{\alpha_0}(\alpha)}
   &=
  \operatorname{int}\bigl(f_N(tw_{\alpha_0})\bigr)\circ
  f_\alpha\circ\operatorname{int}(w_{\alpha_0}^{-1})\\
   &=
  f_{w_{\alpha_0}(\alpha)}\circ\operatorname{int}(tw_{\alpha_0})
    \circ\operatorname{int}(w_{\alpha_0}^{-1}).
\end{aligned}
\]

\SGARefTarget{sga3:XXIII:2.1.4}
\SGAStatement{Lemma 2.1.4}{The restriction of \(f\) to \(U\)
(respectively \(U^-\)) is a morphism of groups. In particular, this
restriction is independent of the chosen ordering of the roots.}

It suffices to prove the assertion for \(U\). By \SGARefLink{sga3:xxii:lem:5.5.8}{Exposé~XXII, 5.5.8},
it is enough to show that, for every pair \(\alpha<\beta\) of positive
roots and every \(x_\alpha\in U_\alpha(S')\),
\(x_\beta\in U_\beta(S')\), \(S'\to S\),
\[
  f_\beta(x_\beta)f_\alpha(x_\alpha)f_\beta(x_\beta^{-1})
   =f(x_\beta x_\alpha x_\beta^{-1}).
\]
By \SGARefLink{sga3:xxii:cor:5.5.2}{Exposé~XXII, 5.5.2}, there are elements
\(x_\gamma\in U_\gamma(S')\), where
\(\gamma=i\alpha+j\beta\in R\), \(i>0\), \(j\geq0\)
\NDE{The printed \(i,j>0\) was corrected to \(i>0,j\geq0\), since
\(x_\alpha\) occurs in the product on the right.}, such that
\[
  x_\beta x_\alpha x_\beta^{-1}=\prod_\gamma x_\gamma.
\]
It remains to verify
\[
  f_\beta(x_\beta)f_\alpha(x_\alpha)f_\beta(x_\beta^{-1})
   =
  \prod_{\substack{\gamma=i\alpha+j\beta\\i>0,\ j\geq0}}
  f_\gamma(x_\gamma).
\]
By \SGARefLink{sga3:XXI:3.5.4}{Exposé~XXI, 3.5.4}, there is a \(w\in W\) such that
\(w(\alpha)=\alpha_0\in\Delta\) and
\(w(\beta)\in A_{\alpha_0\beta_0}\), in the notation of \SGARefLink{sga3:XXIII:1.7}{1.7}, for
some \(\beta_0\in\Delta\). Lift \(w\) to
\(n\in\operatorname{Norm}_G(T)(S)\) by \SGARefLink{sga3:xxii:cor:3.8}{Exposé~XXII, 3.8}. It suffices
to verify the relation after conjugation by \(f_N(n)\); by \SGARefLink{sga3:XXIII:2.1.3}{2.1.3},
this reduces to \(\alpha,\beta\in R^+_{\alpha_0\beta_0}\), where
condition \SGARefLink{sga3:XXIII:2.1:item:iii}{2.1(iii)} applies.

\SGARefTarget{sga3:XXIII:2.1.5}
\SGAStatement{Lemma 2.1.5}{Let
\(n\in\operatorname{Norm}_G(T)(S)\) satisfy
\(\operatorname{int}(n)U=U^-\), so that \(n\) represents the
symmetry of the Weyl group\NDE{Relative to \(\Delta\).}
(\SGARefLink{sga3:XXI:3.6.14}{Exposé~XXI, 3.6.14}). For every \(u\in U(S')\), \(S'\to S\)
(respectively \(u\in U^-(S')\)), one has
\[
  f(nun^{-1})=f_N(n)f(u)f_N(n^{-1}).
\]}
This follows immediately from \SGARefLink{sga3:XXIII:2.1.3}{2.1.3} and \SGARefLink{sga3:XXIII:2.1.4}{2.1.4}.

\SGARefTarget{sga3:XXIII:2.1.6}
\SGAStatement{Lemma 2.1.6}{Let \(u\in B(S')\),
\(v\in B^-(S')\), and \(g\in\Omega(S')\), where \(S'\to S\). Then
\[
  f(vgu)=f(v)f(g)f(u).
\]}

Write
\[
  v=v_1t_1,\qquad g=v_2t_2u_2,\qquad u=t_3u_3,
\]
with \(v_i\in U^-(S')\), \(t_i\in T(S')\), and \(u_i\in U(S')\).
On the one hand,
\[
  f(v)f(g)f(u)
  =f(v_1)f_T(t_1)f(v_2)f_T(t_2)f(u_2)f_T(t_3)f(u_3).
\]
On the other hand,
\[
\begin{aligned}
 f(vgu)
 &=f(v_1t_1v_2t_1^{-1}\,t_1t_2t_3\,
       t_3^{-1}u_2t_3u_3)\\
 &=f(v_1\cdot t_1v_2t_1^{-1})f_T(t_1t_2t_3)
       f(t_3^{-1}u_2t_3\cdot u_3).
\end{aligned}
\]
Using \SGARefLink{sga3:XXIII:2.1.4}{2.1.4} on the two outer factors gives
\[
 f(vgu)=f(v_1)f(t_1v_2t_1^{-1})f_T(t_1t_2t_3)
   f(t_3^{-1}u_2t_3)f(u_3).
\]
It remains only to use
\[
\begin{aligned}
 f(t_1v_2t_1^{-1})&=f_T(t_1)f(v_2)f_T(t_1)^{-1},\\
 f(t_3^{-1}u_2t_3)&=f_T(t_3)^{-1}f(u_2)f_T(t_3),
\end{aligned}
\]
which follow from the definition of \(f\) and \SGARefLink{sga3:XXIII:2.1.3}{2.1.3}.

\SGARefTarget{sga3:XXIII:2.1.7}
\SGAStatement{Lemma 2.1.7}{Let \(\alpha\in\Delta\) and
\[
  g\in\Omega(S')\cap\operatorname{int}(w_\alpha)^{-1}\Omega(S'),
  \qquad S'\to S.
\]
Then
\[
  f(w_\alpha gw_\alpha^{-1})
   =f_N(w_\alpha)f(g)f_N(w_\alpha)^{-1}.
\]}

By \SGARefLink{sga3:xxii:cor:5.6.8}{Exposé~XXII, 5.6.8}, write
\[
  g=ax_{-\alpha}tx_\alpha b,
\]
where \(a\in U^-_{-\widehat\alpha}(S')\),
\(x_{-\alpha}\in U_{-\alpha}(S')\), \(t\in T(S')\),
\(x_\alpha\in U_\alpha(S')\), and
\(b\in U_{\widehat\alpha}(S')\). By \SGARefLink{sga3:XXIII:2.1.3}{2.1.3} and \SGARefLink{sga3:XXIII:2.1.4}{2.1.4}, and because
\(s_\alpha\) permutes the positive roots other than \(\alpha\)
(\SGARefLink{sga3:XXI:3.3.2}{Exposé~XXI, 3.3.2}),
\[
  \operatorname{int}(w_\alpha)a\in U^-_{-\widehat\alpha}(S'),
  \qquad
  \operatorname{int}(w_\alpha)b\in U_{\widehat\alpha}(S').
\]
The desired formula holds for \(g=a\) and \(g=b\). By \SGARefLink{sga3:XXIII:2.1.6}{2.1.6} it
therefore suffices to treat
\(g=x_{-\alpha}tx_\alpha\in Z_\alpha(S')\); in that case everything
takes place in \(Z_\alpha\), and condition \SGARefLink{sga3:XXIII:2.1:item:ii}{2.1(ii)} applies.

\SGARefTarget{sga3:XXIII:2.1.8}
\SGAStatement{Lemma 2.1.8}{Let
\(n\in\operatorname{Norm}_G(T)(S)\). For every \(S'\to S\) and every
\(g\in\Omega(S')\) such that
\(\operatorname{int}(n)g\in\Omega(S')\),
\[
  f(ngn^{-1})=f_N(n)f(g)f_N(n)^{-1}.
\]}

This is immediate for \(n\in T(S)\), by \SGARefLink{sga3:XXIII:2.1.3}{2.1.3}. The two sides define
morphisms from
\(\Omega\cap\operatorname{int}(n)^{-1}\Omega\) to \(H\). It is
enough to find an open \(V_n\subset\Omega\), containing the unit
section, such that \(\operatorname{int}(n)V_n\subset\Omega\), and on
which \(f\circ\operatorname{int}(n)\) and
\(\operatorname{int}(f_N(n))\circ f\) agree. By the structure of
\(\operatorname{Norm}_G(T)\), it suffices to show that if this is true
for \(n'\), then it is true for \(n=n'w_\alpha\),
\(\alpha\in\Delta\). Put
\[
  V_n=\Omega\cap\operatorname{int}(w_\alpha)^{-1}V_{n'}.
\]
Then
\(\operatorname{int}(n)V_n\subset\operatorname{int}(n')V_{n'}
\subset\Omega\). If \(g\in V_n(S')\), then
\[
  \operatorname{int}(n)g
   =\operatorname{int}(n')\operatorname{int}(w_\alpha)g.
\]
Now \(\operatorname{int}(w_\alpha)g\in V_{n'}(S')\), so the
induction hypothesis gives
\[
  f\bigl(\operatorname{int}(n')\operatorname{int}(w_\alpha)g\bigr)
   =\operatorname{int}(f_N(n'))f(\operatorname{int}(w_\alpha)g).
\]
\SGARefLink{sga3:XXIII:2.1.7}{Lemma 2.1.7} gives
\[
  f(\operatorname{int}(w_\alpha)g)
   =\operatorname{int}(f_N(w_\alpha))f(g),
\]
and the conclusion follows.

We now prove \SGARefLink{sga3:XXIII:2.1.2}{Proposition 2.1.2}. Let \(x,x'\in\Omega(S')\), and write
\[
  x=vtu,\qquad x'=v't'u',
\]
so \(xx'=vt(uv')t'u'\). By \SGARefLink{sga3:XXIII:2.1.6}{2.1.6} and
\(U^-(S')\Omega(S')U(S')=\Omega(S')\), it suffices to show that
\(f(uv')=f(u)f(v')\) whenever \(uv'\in\Omega(S')\). Choose
\(n\in\operatorname{Norm}_G(T)(S')\) as in \SGARefLink{sga3:XXIII:2.1.5}{2.1.5}. Then
\[
\begin{aligned}
 f(u)&=f_N(n)^{-1}f(nun^{-1})f_N(n),\\
 f(v')&=f_N(n)^{-1}f(nv'n^{-1})f_N(n),
\end{aligned}
\]
and hence
\[
 f(u)f(v')=f_N(n)^{-1}f(nun^{-1})f(nv'n^{-1})f_N(n).
\]
Since \(nun^{-1}\in U^-(S')\) and \(nv'n^{-1}\in U(S')\),
\[
 f(nun^{-1})f(nv'n^{-1})
   =f\bigl((nun^{-1})(nv'n^{-1})\bigr)
   =f(nuv'n^{-1}).
\]
Thus
\[
  f(u)f(v')=f_N(n)^{-1}f(nuv'n^{-1})f_N(n).
\]
If \(uv'\in\Omega(S')\), \SGARefLink{sga3:XXIII:2.1.8}{Lemma 2.1.8} finishes the proof.

\SGARefTarget{sga3:XXIII:2.2}
\SGAStatement{Remark 2.2}{Instead of prescribing \(f_N\), one may
prescribe a morphism \(f_T:T\to H\) and sections
\(h_\alpha\in H(S)\), \(\alpha\in\Delta\), satisfying the conditions
of \SGARefLink{sga3:XXIII:1.8.2}{1.8.2}. Condition \SGARefLink{sga3:XXIII:1.8.2:item:ii}{(ii)} of \SGARefLink{sga3:XXIII:2.1}{Theorem 2.1} must then be replaced by:
\begin{quote}
\emph{(ii\('\)) There exists a morphism of groups
\(F_\alpha:Z_\alpha\to H\) extending \(f_\alpha,f_{-\alpha},f_T\)
and satisfying \(F_\alpha(w_\alpha)=h_\alpha\).}
\end{quote}}

We now give the preceding theorem a more explicit form, retaining its
notation.

\SGARefTarget{sga3:XXIII:2.3}
\SGAStatement{Theorem 2.3}{Let \(S\) be a scheme, let \(G\) be a
pinned \(S\)-group, and let \(H\) be an fppf \(S\)-sheaf in groups.
Suppose given morphisms of groups
\[
  f_T:T\longrightarrow H,\qquad
  f_\alpha:U_\alpha\longrightarrow H\quad(\alpha\in\Delta),
\]
and sections \(h_\alpha\in H(S)\), \(\alpha\in\Delta\). There exists
a necessarily unique morphism of groups \(f:G\to H\) inducing
\(f_T\) and the \(f_\alpha\) and satisfying
\(f(w_\alpha)=h_\alpha\) for all \(\alpha\in\Delta\) if and only if:
\begin{enumerate}
  \item[\SGAInlineRefTarget{sga3:XXIII:2.3:item:i}(i)] for every \(S'\to S\), \(t\in T(S')\),
    \(\alpha\in\Delta\), and \(x\in U_\alpha(S')\),
    \begin{equation}
    \tag{1}\label{sga3:XXIII:2.3:eq1}
      f_T(t)f_\alpha(x)f_T(t)^{-1}
       =f_\alpha(\operatorname{int}(t)x);
    \end{equation}
  \item[\SGAInlineRefTarget{sga3:XXIII:2.3:item:ii}(ii)] for every \(\alpha\in\Delta\), \(S'\to S\), and
    \(t\in T(S')\),
    \begin{equation}
    \tag{2}\label{sga3:XXIII:2.3:eq2}
      h_\alpha f_T(t)h_\alpha^{-1}
       =f_T(s_\alpha(t))
       =f_T\!\bigl(t\alpha^\ast(\alpha(t)^{-1})\bigr);
    \end{equation}
  \item[(iii)] for every
    \((\alpha,\beta)\in\Delta\times\Delta\),
    \begin{equation}
    \tag{3}\label{sga3:XXIII:2.3:eq3}
      (h_\alpha h_\beta)^{n_{\alpha\beta}}
       =f_T(t_{\alpha\beta});
    \end{equation}
  \item[(iv)] for every \(\alpha\in\Delta\), recalling that
    \(u_\alpha=\exp_\alpha(X_\alpha)\),
    \begin{equation}
    \tag{4}\label{sga3:XXIII:2.3:eq4}
      (h_\alpha f_\alpha(u_\alpha))^3=e;
    \end{equation}
  \item[\SGAInlineRefTarget{sga3:XXIII:2.3:item:v}(v)] for every distinct
    \(\alpha,\beta\in\Delta\), there exists a morphism
    \(f_{\alpha\beta}:U_{\alpha\beta}\to H\) such that
    \begin{equation}
    \tag{5}\label{sga3:XXIII:2.3:eq5}
      f_{\alpha\beta}|_{U_\alpha}=f_\alpha,\qquad
      f_{\alpha\beta}|_{U_\beta}=f_\beta,
    \end{equation}
    and, for every \(\gamma\in R^+_{\alpha\beta}\) with
    \(\gamma\ne\alpha\) (respectively \(\gamma\ne\beta\)) and every
    \(x\in U_\gamma(S')\), \(S'\to S\),
    \begin{equation}
    \tag{6}\label{sga3:XXIII:2.3:eq6}
      \operatorname{int}(h_\alpha)f_{\alpha\beta}(x)
       =f_{\alpha\beta}(\operatorname{int}(w_\alpha)x)
    \end{equation}
    (respectively
    \(\operatorname{int}(h_\beta)f_{\alpha\beta}(x)
     =f_{\alpha\beta}(\operatorname{int}(w_\beta)x)\)).
\end{enumerate}}

\SGAProof
Uniqueness follows from \SGARefLink{sga3:XXIII:1.9.2}{1.9.2}. We prove existence.

\SGARefTarget{sga3:XXIII:2.3.1}
\SGAStatement{Lemma 2.3.1}{There exists a morphism
\[
  f_N:\operatorname{Norm}_G(T)\longrightarrow H
\]
extending \(f_T\) and satisfying \(f_N(w_\alpha)=h_\alpha\).}
This is precisely \SGARefLink{sga3:XXIII:1.8.2}{1.8.2}, in view of conditions
\eqref{sga3:XXIII:2.3:eq2} and \eqref{sga3:XXIII:2.3:eq3}.

\SGARefTarget{sga3:XXIII:2.3.2}
\SGAStatement{Lemma 2.3.2}{For every \(\alpha\in\Delta\), there is a
morphism of groups \(F_\alpha:Z_\alpha\to H\) extending \(f_T\) and
\(f_\alpha\), satisfying \(F_\alpha(w_\alpha)=h_\alpha\), and hence
extending \(f_N|_{\operatorname{Norm}_{Z_\alpha}(T)}\).}
This follows from \SGARefLink{sga3:xx:6.2}{Exposé~XX, 6.2} and conditions
\eqref{sga3:XXIII:2.3:eq1}, \eqref{sga3:XXIII:2.3:eq3}, and
\eqref{sga3:XXIII:2.3:eq4}.

\SGARefTarget{sga3:XXIII:2.3.3}
\SGAStatement{Lemma 2.3.3}{Let
\((\alpha,\beta)\in\Delta\times\Delta\), \(\alpha\ne\beta\). Let
\(n\in\operatorname{Norm}_{Z_{\alpha\beta}}(T)(S)\) satisfy
\(n(\alpha)=\alpha\) (respectively \(n(\alpha)=\beta\)), equivalently
\(\operatorname{int}(n)U_\alpha=U_\alpha\) (respectively
\(U_\beta\)). Then, for every \(S'\to S\) and \(x\in U_\alpha(S')\),
\[
  \operatorname{int}(f_N(n))f_\alpha(x)
   =f_\alpha(\operatorname{int}(n)x)
\]
(respectively the right-hand side is
\(f_\beta(\operatorname{int}(n)x)\)).}

There exist \(t\in T(S)\) and a sequence
\(\{\alpha_i\}\subset\{\alpha,\beta\}\) such that
\[
  n=tw_{\alpha_k}\cdots w_{\alpha_1}.
\]
Condition \eqref{sga3:XXIII:2.3:eq1} treats \(n=t\), so we may omit
\(t\). Proceed by induction on \(k\), and set
\[
  \gamma_i=s_{\alpha_i}\cdots s_{\alpha_1}(\alpha).
\]
If every \(\gamma_i\) is positive, condition \SGARefLink{sga3:XXIII:2.3:item:v}{(v)} yields
\[
  \operatorname{int}\bigl(f_N(w_{\alpha_i}\cdots w_{\alpha_1})\bigr)
     f_\alpha(x)
   =
  f_{\alpha\beta}\bigl(
     \operatorname{int}(w_{\alpha_i}\cdots w_{\alpha_1})x\bigr),
\]
and \(i=k\) gives the result. Otherwise there is a \(j<k\) such that
\(\gamma_j\in R^+\) but \(\gamma_{j+1}\notin R^+\). Since
\(\gamma_{j+1}=s_{\alpha_j}(\gamma_j)\), \SGARefLink{sga3:XXI:3.3.2}{Exposé~XXI, 3.3.2} gives
\(\gamma_j=\alpha_j\), hence \(\gamma_j\) is \(\alpha\) or \(\beta\).
Decompose
\[
 n=n'n'',\qquad
 n'=w_{\alpha_k}\cdots w_{\alpha_{j+1}},\qquad
 n''=w_{\alpha_j}\cdots w_{\alpha_1}.
\]
Both \(n'\) and \(n''\) satisfy the lemma's hypotheses and involve
fewer than \(k\) reflections, so induction applies.

\SGARefTarget{sga3:XXIII:2.3.4}
\SGAStatement{Lemma 2.3.4}{Let \(\alpha\in R\). Suppose
\(n,n'\in\operatorname{Norm}_G(T)(S)\) and
\(\beta,\beta'\in\Delta\) satisfy
\(n(\alpha)=\beta\) and \(n'(\alpha)=\beta'\). Then, for every
\(x\in U_\alpha(S')\), \(S'\to S\),
\[
 \operatorname{int}(f_N(n)^{-1})f_\beta(nx n^{-1})
 =
 \operatorname{int}(f_N(n')^{-1})f_{\beta'}(n'x n'^{-1}).
\]}

It suffices to prove that if \(n(\alpha)=\beta\), with
\(\alpha,\beta\in\Delta\), then
\[
  \operatorname{int}(f_N(n))\circ f_\alpha
    =f_\beta\circ\operatorname{int}(n).
\]
By Tits's lemma (\SGARefLink{sga3:XXI:5.6}{Exposé~XXI, 5.6}), there are simple roots
\[
  \alpha_0=\alpha,\alpha_1,\ldots,\alpha_m=\beta
\]
and elements \(w_i\in W\), \(0\leq i<m\), such that
\[
  n=w_{m-1}\cdots w_0,\qquad
  w_i(\alpha_i)=\alpha_{i+1},
\]
and for each \(i\) there is a simple root \(\beta_i\) such that
\[
  w_i\in W_{\alpha_i\beta_i},\qquad
  \alpha_{i+1}=\alpha_i\ \text{or}\ \beta_i.
\]
Induction reduces the assertion to \SGARefLink{sga3:XXIII:2.3.3}{Lemma 2.3.3}.

\SGARefTarget{sga3:XXIII:2.3.5}
\SGAStatement{Lemma 2.3.5}{There is a family of morphisms of groups
\(f'_\gamma:U_\gamma\to H\), \(\gamma\in R\), such that:
\begin{enumerate}
  \item[\SGAInlineRefTarget{sga3:XXIII:2.3.5:item:i}(i)] for \(\alpha\in\Delta\),
    \(f'_\alpha=f_\alpha\) and
    \(f'_{-\alpha}=F_\alpha|_{U_{-\alpha}}\);
  \item[\SGAInlineRefTarget{sga3:XXIII:2.3.5:item:ii}(ii)] if \(\alpha,\beta\in\Delta\) and
    \(\gamma\in R^+_{\alpha\beta}\), then
    \(f'_\gamma=f_{\alpha\beta}|_{U_\gamma}\);
  \item[\SGAInlineRefTarget{sga3:XXIII:2.3.5:item:iii}(iii)] if \(n\in\operatorname{Norm}_G(T)(S)\) and
    \(\alpha,\beta\in R\) satisfy \(n(\alpha)=\beta\), then
    \[
      \operatorname{int}(f_N(n))f'_\alpha(x)
       =f'_\beta(\operatorname{int}(n)x)
    \]
    for every \(x\in U_\alpha(S')\), \(S'\to S\).
\end{enumerate}}

For every root \(\alpha\), choose
\(n\in\operatorname{Norm}_G(T)(S)\) with \(n(\alpha)\in\Delta\), and
define \(f'_\alpha(x)\) by the expression in \SGARefLink{sga3:XXIII:2.3.4}{Lemma 2.3.4}. This is
independent of \(n\), and \(f'_\alpha\) is a morphism of groups.
Property \SGARefLink{sga3:XXIII:2.3.5:item:iii}{(iii)} holds by construction. The first part of \SGARefLink{sga3:XXIII:2.3.5:item:i}{(i)} follows
by taking \(n=e\), and the second by taking \(n=w_\alpha\) and using
\SGARefLink{sga3:XXIII:2.3.2}{2.3.2}. If \(\gamma\in R^+_{\alpha\beta}\), choose
\(n\in\operatorname{Norm}_{Z_{\alpha\beta}}(T)(S)\) with
\(n(\gamma)=\alpha\) or \(\beta\); then \SGARefLink{sga3:XXIII:2.3.5:item:iii}{(iii)}, together with
conditions \eqref{sga3:XXIII:2.3:eq5} and
\eqref{sga3:XXIII:2.3:eq6}, proves \SGARefLink{sga3:XXIII:2.3.5:item:ii}{(ii)}.

\SGARefTarget{sga3:XXIII:2.3.6}
\SGAParagraph{2.3.6}
We finish the proof by verifying the conditions of \SGARefLink{sga3:XXIII:2.1}{Theorem 2.1} for
\(f_N\) and \(f'_\alpha\), \(\alpha\in R\):
\begin{itemize}
  \item \SGARefLink{sga3:XXIII:2.1:item:i}{2.1(i)} follows from \SGARefLink{sga3:XXIII:2.3.5:item:iii}{2.3.5(iii)};
  \item \SGARefLink{sga3:XXIII:2.1:item:ii}{2.1(ii)} follows from \SGARefLink{sga3:XXIII:2.3.5:item:i}{2.3.5(i)} and \SGARefLink{sga3:XXIII:2.3.2}{2.3.2};
  \item \SGARefLink{sga3:XXIII:2.1:item:iii}{2.1(iii)} follows from \SGARefLink{sga3:XXIII:2.3.5:item:ii}{2.3.5(ii)} and the fact that
    \(f_{\alpha\beta}\) is a morphism of groups.
\end{itemize}

An immediate corollary is:

\SGARefTarget{sga3:XXIII:2.4}
\SGAStatement{Corollary 2.4}{Let \(S\) be a scheme, let \(G\) be a
pinned \(S\)-group of semisimple rank at least \(1\), and let \(H\)
be an fppf \(S\)-sheaf in groups. For each
\((\alpha,\beta)\in\Delta\times\Delta\), let
\[
  F_{\alpha\beta}:Z_{\alpha\beta}\longrightarrow H
\]
be a morphism of groups. There exists a morphism \(f:G\to H\)
inducing the \(F_{\alpha\beta}\) if and only if, for every
\((\alpha,\beta)\in\Delta\times\Delta\),
\[
  F_{\alpha\beta}|_{Z_\alpha}=F_{\alpha\alpha}.
\]}

Necessity is clear. Conversely, put
\[
  f_T=F_{\alpha\alpha}|_T;
\]
this is independent of \(\alpha\), since
\[
  F_{\alpha\alpha}|_T
   =F_{\alpha\beta}|_T
   =F_{\beta\beta}|_T.
\]
Set
\[
  f_\alpha=F_{\alpha\alpha}|_{U_\alpha},\qquad
  h_\alpha=F_{\alpha\alpha}(w_\alpha),\qquad
  f_{\alpha\beta}=F_{\alpha\beta}|_{U_{\alpha\beta}}.
\]
The conditions of \SGARefLink{sga3:XXIII:2.3}{Theorem 2.3} hold: (1), (2), and (4) hold in
\(Z_\alpha\), while (3), (5), and (6) hold in \(Z_{\alpha\beta}\).
Thus there is a morphism \(f:G\to H\) extending \(f_T\) and all
\(f_\alpha\), and satisfying \(f(w_\alpha)=h_\alpha\). It agrees
with \(F_{\alpha\beta}\) on the generators of \(Z_{\alpha\beta}\);
hence \(f|_{Z_{\alpha\beta}}=F_{\alpha\beta}\).

% END INCLUDED SGA3_English_current_buildable_source_20260823/source/inputs/XXIII/tex/components/04_expose_XXIII_theorem21_through_corollary24.tex

% BEGIN INCLUDED SGA3_English_current_buildable_source_20260823/source/inputs/XXIII/tex/components/05_expose_XXIII_cor25_through_par313.tex
\SGARefTarget{sga3:XXIII:2.5}
\SGAStatement{Corollary 2.5}{Let \(S\) be a scheme; let \(G\) and
\(G'\) be two pinned \(S\)-groups of semisimple rank \(2\); let
\(q>0\) be an integer such that \(x\mapsto x^q\) is an endomorphism
of \(\mathbf G_{a,S}\); and let
\[
  h:R(G')\longrightarrow R(G)
\]
be a \(q\)-morphism of pinned root data. For each
\(\alpha\in R^+\), choose
\[
  X_\alpha\in\Gamma(S,\mathfrak g^\alpha)^\times,
  \qquad
  X'_{d(\alpha)}
    \in\Gamma(S,\mathfrak g'^{\,d(\alpha)})^\times,
\]
extending the canonical choices for \(\alpha\in\Delta\). Suppose:
\begin{enumerate}
  \item[\SGAInlineRefTarget{sga3:XXIII:2.5:item:i}(i)] if \(\Delta=\{\alpha,\beta\}\), then
    \[
      D_S(h)(t_{\alpha\beta})=t'_{d(\alpha)d(\beta)};
    \]
  \item[\SGAInlineRefTarget{sga3:XXIII:2.5:item:ii}(ii)] for every \(\alpha\in\Delta\) and every
    \(\beta\in R^+\), \(\beta\ne\alpha\), if
    \(z\in\mathbf G_m(S)\) is defined by
    \[
      \operatorname{Ad}(w_\alpha)X_\beta
        =zX_{s_\alpha(\beta)},
    \]
    then
    \[
      \operatorname{Ad}(w'_{d(\alpha)})X'_{d(\beta)}
        =z^{q(\beta)}X'_{d(s_\alpha(\beta))};
    \]
  \item[(iii)] there is a morphism of groups \(f:U\to U'\) such
    that, for every \(\alpha\in R^+\), every \(S'\to S\), and every
    \(x\in\mathbf G_a(S')\),
    \[
      f(\exp(xX_\alpha))
        =\exp\!\bigl(x^{q(\alpha)}X'_{d(\alpha)}\bigr).
    \]
\end{enumerate}
Then there is a morphism of pinned groups \(g:G\to G'\) such that
\(R(g)=h\).}

Indeed, define \(f_\alpha:U_\alpha\to G'\) by
\[
  f_\alpha(\exp(xX_\alpha))
    =\exp\!\bigl(x^{q(\alpha)}X'_{d(\alpha)}\bigr),
\]
and put
\[
  f_T=D_S(h),\qquad h_\alpha=w'_{d(\alpha)}.
\]
The conditions of \SGARefLink{sga3:XXIII:2.3}{Theorem 2.3} hold: note that
\[
  q(s_\alpha(\beta))=q(\beta)
\]
by \SGARefLink{sga3:XXI:6.8.4}{Exposé XXI, 6.8.4}, and that
\[
  D_S(h)(t_\alpha)=t'_{d(\alpha)}
\]
always holds. The assertion follows at once.

\SGARefTarget{sga3:XXIII:2.6}
\SGAStatement{Remark 2.6}{Condition (v) of \SGARefLink{sga3:XXIII:2.3}{Theorem 2.3} may be made
more explicit as follows. One must first verify:
\begin{enumerate}
  \item[\SGAInlineRefTarget{sga3:XXIII:2.6:item:a}(a)] for every word in \(w_\alpha,w_\beta\) whose
    corresponding word carries \(\alpha\) or \(\beta\) to
    \(\alpha\) or \(\beta\), the corresponding relation of the type
    occurring in \SGARefLink{sga3:XXIII:2.3.3}{2.3.3} holds.
\end{enumerate}
The proof of \SGARefLink{sga3:XXIII:2.3.3}{2.3.3} shows that it suffices to check the words that are
minimal, in the sense that no nontrivial initial subword satisfies
the imposed conditions.

If \SGARefLink{sga3:XXIII:2.6:item:a}{(a)} holds, define \(f_\gamma:U_\gamma\to H\) for every
\(\gamma\in R^+_{\alpha\beta}\) as in \SGARefLink{sga3:XXIII:2.3.5}{2.3.5}. One must then require:
\begin{enumerate}
  \item[\SGAInlineRefTarget{sga3:XXIII:2.6:item:b}(b)] the map defined by the \(f_\gamma\),
    \[
      U_{\alpha\beta}
        =\prod_{\gamma\in R^+_{\alpha\beta}}U_\gamma
        \longrightarrow H,
    \]
    is a morphism of groups.
\end{enumerate}
By \SGARefLink{sga3:xxii:lem:5.5.8}{Exposé XXII, 5.5.8}, condition \SGARefLink{sga3:XXIII:2.6:item:b}{(b)} follows from:
\begin{enumerate}
  \item[(b\('\))] the preceding map respects the commutation
    relations between \(U_\gamma\) and \(U_\delta\), for
    \(\gamma,\delta\in R^+_{\alpha\beta}\), \(\gamma>\delta\);
    namely, the relations involving
    \(C_{i,j,\gamma,\delta}\) in \SGARefLink{sga3:xxii:cor:5.5.2}{Exposé XXII, 5.5.2}.
\end{enumerate}}

Conversely, the conjunction of \SGARefLink{sga3:XXIII:2.6:item:a}{(a)} and (b\('\)) is equivalent to
condition (v).

These conditions may even be reduced to conditions involving only
\[
  h_\alpha,\quad h_\beta,\quad
  f_\alpha(u_\alpha),\quad f_\beta(u_\beta)
\]
in \(H\). For example, if
\(s_\alpha s_\beta s_\alpha(\beta)=\alpha\), a condition of type
(a) reads
\begin{equation}
  \operatorname{int}(h_\alpha h_\beta h_\alpha)f_\beta(x)
    =f_\alpha\!\left(
       \operatorname{int}(w_\alpha w_\beta w_\alpha)x
     \right).
  \tag{1}
  \label{sga3:XXIII:2.6:eq1}
\end{equation}
Here \(x\in U_\beta(S')\), \(S'\to S\). In particular, for
\(x=u_\beta\),
\[
  \operatorname{int}(w_\alpha w_\beta w_\alpha)u_\beta
    =u_\alpha^z
\]
for a certain \(z\in\mathbf G_m(S)\), and
\begin{equation}
  \operatorname{int}(h_\alpha h_\beta h_\alpha)f_\beta(u_\beta)
    =f_\alpha(u_\alpha)^z.
  \tag{1\('\)}
  \label{sga3:XXIII:2.6:eq1prime}
\end{equation}

Conversely, taking conditions \SGARefLink{sga3:XXIII:2.5:item:i}{(i)}--(iv) of \SGARefLink{sga3:XXIII:2.3}{2.3} into account,
\eqref{sga3:XXIII:2.6:eq1prime} implies
\eqref{sga3:XXIII:2.6:eq1}. If \(t\in T(S')\), conjugating
\eqref{sga3:XXIII:2.6:eq1prime} by \(t\), and using \SGARefLink{sga3:XXIII:2.3:item:i}{(i)} and \SGARefLink{sga3:XXIII:2.3:item:ii}{(ii)},
gives \eqref{sga3:XXIII:2.6:eq1} for
\[
  x=\operatorname{int}(t)u_\beta=u_\beta^{\beta(t)}.
\]
The morphism \(\beta:T\to\mathbf G_{m,S}\) is faithfully flat, so
\eqref{sga3:XXIII:2.6:eq1} holds for
\(x\in U_\beta(S')^\times\). It is additive in \(x\), and every
section of \(U_\beta\) is locally the sum of two invertible sections.
Thus
\[
  \text{(1')}+\text{(i)}+\text{(ii)}
    \quad\Longrightarrow\quad\text{(1)}.
\]

The same argument applies to conditions of type \SGARefLink{sga3:XXIII:2.6:item:b}{(b)}. One uses the
fact that, for two distinct positive roots \(\gamma,\gamma'\), hence
linearly independent over \(\mathbf Q\), the morphism
\[
  T\longrightarrow\mathbf G_{m,S}^2
\]
with components \(\gamma,\gamma'\) is faithfully flat. We leave the
details of this transposition to the reader.

\SGANextSectionTarget{sga3:xxiii:section:3}\SGASection{3}{Groups of Semisimple Rank \(2\)}
\SGANextSectionTarget{sga3:xxiii:subsection:3-1}\SGASubsection{3.1}{Generalities}

\SGARefTarget{sga3:XXIII:3.1.1}
\SGAStatement{Lemma 3.1.1}{Let \(S\) be a scheme, let \(G\) be a
pinned \(S\)-group, and let \(\alpha,\beta\) be two roots of \(G\)
such that \(\alpha+\beta\ne0\).
\begin{enumerate}
  \item[\SGAInlineRefTarget{sga3:xxiii:subsection:3-1:item:i}\SGAInlineRefTarget{sga3:xxiii:lemma:3-1-1:item:i}(i)] If \(\alpha+\beta\notin R\), then
    \[
      \exp(X_\alpha)\exp(X_\beta)
        =\exp(X_\beta)\exp(X_\alpha)
    \]
    for every \(S'\to S\),
    \(X_\alpha\in W(\mathfrak g^\alpha)(S')\), and
    \(X_\beta\in W(\mathfrak g^\beta)(S')\).
  \item[\SGAInlineRefTarget{sga3:xxiii:lemma:3-1-1:item:ii}\SGAInlineRefTarget{sga3:xxiii:subsection:3-1:item:ii}(ii)] If neither \(\alpha+\beta\) nor \(\alpha-\beta\) is a
    root, then
    \[
      w_\alpha(z_\alpha)\exp(X_\beta)w_\alpha(z_\alpha)^{-1}
        =\exp(X_\beta)
    \]
    for every \(X_\beta\in W(\mathfrak g^\beta)(S')\) and
    \(z_\alpha\in W(\mathfrak g^\alpha)^\times(S')\), and
    \[
      w_\alpha(z_\alpha)w_\beta(z_\beta)
        =w_\beta(z_\beta)w_\alpha(z_\alpha)
    \]
    for every
    \(z_\alpha\in W(\mathfrak g^\alpha)^\times(S')\) and
    \(z_\beta\in W(\mathfrak g^\beta)^\times(S')\).
  \item[\SGAInlineRefTarget{sga3:XXIII:3.1.1:item:iii}\SGAInlineRefTarget{sga3:xxiii:lemma:3-1-1:item:iii}(iii)] Let
    \(X_\alpha\in W(\mathfrak g^\alpha)^\times(S')\),
    \(X_\beta\in W(\mathfrak g^\beta)^\times(S')\), and
    \(w\in\operatorname{Norm}_G(T)(S')\) satisfy
    \(w(\alpha)=\beta\). Define \(z\in\mathbf G_m(S')\) by
    \[
      \operatorname{Ad}(w)X_\alpha=zX_\beta.
    \]
    Then
    \[
    \begin{aligned}
      \operatorname{int}(w)\exp(xX_\alpha)
        &=\exp(xzX_\beta),\\
      \operatorname{int}(w)\exp(yX_\alpha^{-1})
        &=\exp(yz^{-1}X_\beta^{-1}),\\
      \operatorname{int}(w)w_\alpha(X_\alpha)
        &=\beta^*(z)w_\beta(X_\beta).
    \end{aligned}
    \]
    In particular, if \(z=\pm1\), then
    \[
      \operatorname{int}(w)w_\alpha(X_\alpha)
        =w_\beta(X_\beta)^z.
    \]
  \item[\SGAInlineRefTarget{sga3:XXIII:3.1.1:item:iv}(iv)] If
    \(t_\alpha=\alpha^*(-1)\) and \(t_\beta=\beta^*(-1)\), then
    \[
      s_\alpha(t_\beta)
        =t_\beta t_\alpha^{(\beta^*,\alpha)},
      \qquad
      \beta(t_\alpha)=(-1)^{(\alpha^*,\beta)}.
    \]
\end{enumerate}}

\SGAProof
Part \SGARefLink{sga3:xxiii:subsection:3-1:item:i}{(i)} follows immediately from \SGARefLink{sga3:xxii:cor:5.5.2}{Exposé XXII, 5.5.2}. Part \SGARefLink{sga3:xxiii:subsection:3-1:item:ii}{(ii)}
follows from \SGARefLink{sga3:xx:3.1}{Exposé XX, 3.1} and from \SGARefLink{sga3:xxiii:subsection:3-1:item:i}{(i)} applied to
\[
  (\beta,\alpha),\quad(\beta,-\alpha),\quad
  (-\beta,\alpha),\quad(-\beta,-\alpha).
\]
Part \SGARefLink{sga3:XXIII:3.1.1:item:iii}{(iii)} is immediate from the definitions; for its last assertion,
observe that
\[
  \beta^*(-1)=w_\beta(X_\beta)^{-2}.
\]
Finally, \SGARefLink{sga3:XXIII:3.1.1:item:iv}{(iv)} is tautological.

\SGARefTarget{sga3:XXIII:3.1.2}
\SGAStatement{Proposition 3.1.2 (Groups of type
\(A_1\times A_1\))}{Let \(S\) be a scheme and \(G\) a pinned
\(S\)-group of type \(A_1\times A_1\). Write
\(\Delta=R^+=\{\alpha,\beta\}\).
\begin{enumerate}
  \item[\SGAInlineRefTarget{sga3:XXIII:3.1.2:item:i}(i)]
    \[
      t_{\alpha\beta}=(w_\alpha w_\beta)^2
        =t_\alpha t_\beta
        =(w_\beta w_\alpha)^2=t_{\beta\alpha}.
    \]
  \item[\SGAInlineRefTarget{sga3:XXIII:3.1.2:item:ii}(ii)]\SGAInlineRefTarget{sga3:xxiii:proposition:3-1-2:item:ii}
    \[
      \operatorname{Ad}(w_\alpha)X_\beta=X_\beta,
      \qquad
      \operatorname{Ad}(w_\beta)X_\alpha=X_\alpha.
    \]
  \item[\SGAInlineRefTarget{sga3:XXIII:3.1.2:item:iii}(iii)] The groups \(U_\alpha\) and \(U_\beta\) commute; in
    other words, \(U\) is commutative.
\end{enumerate}}

\SGAProof
\SGARefLink{sga3:xxiii:lemma:3-1-1:item:ii}{Lemma 3.1.1(ii)} gives \(w_\alpha w_\beta=w_\beta w_\alpha\), whence
\[
  (w_\alpha w_\beta)^2=w_\alpha^2w_\beta^2=t_\alpha t_\beta,
\]
which proves \SGARefLink{sga3:XXIII:3.1.2:item:i}{(i)}. The same lemma gives \SGARefLink{sga3:XXIII:3.1.2:item:ii}{(ii)}, and \SGARefLink{sga3:XXIII:3.1.2:item:iii}{(iii)} is
\SGARefLink{sga3:xxiii:lemma:3-1-1:item:i}{Lemma 3.1.1(i)}.

\SGARefTarget{sga3:XXIII:3.1.3}
\SGAParagraph{3.1.3}
Condition (v) of \SGARefLink{sga3:XXIII:2.3}{Theorem 2.3} takes the following explicit form. Put
\[
  v_\alpha=f_\alpha(u_\alpha)
  \qquad(\alpha\in\Delta).
\]
Then the two groups of conditions are
\[
\text{(A)}\quad
\begin{cases}
  h_\alpha v_\beta h_\alpha^{-1}=v_\beta,\\
  h_\beta v_\alpha h_\beta^{-1}=v_\alpha,
\end{cases}
\qquad
\text{(B)}\quad
v_\alpha v_\beta=v_\beta v_\alpha.
\]

% END INCLUDED SGA3_English_current_buildable_source_20260823/source/inputs/XXIII/tex/components/05_expose_XXIII_cor25_through_par313.tex

% BEGIN INCLUDED SGA3_English_current_buildable_source_20260823/source/inputs/XXIII/tex/components/06_expose_XXIII_type_A2_and_B2.tex
\SGANextSectionTarget{sga3:xxiii:subsection:3-2}\SGASubsection{3.2}{Groups of Type \(A_2\)}

\SGARefTarget{sga3:XXIII:3.2.1}
\SGAStatement{Proposition 3.2.1}{Let \(S\) be a scheme and \(G\) a
pinned \(S\)-group of type \(A_2\). Write
\[
  \Delta=\{\alpha,\beta\},
  \qquad
  R^+=\{\alpha,\beta,\alpha+\beta\}.
\]
\begin{enumerate}
  \item[\SGAInlineRefTarget{sga3:XXIII:3.2.1:item:i}\SGAInlineRefTarget{sga3:xxiii:proposition:3-2-1:item:i}(i)]
    \[
      t_{\alpha\beta}=(w_\alpha w_\beta)^3
        =e=(w_\beta w_\alpha)^3=t_{\beta\alpha}.
    \]
  \item[\SGAInlineRefTarget{sga3:xxiii:proposition:3-2-1:item:ii}\SGAInlineRefTarget{sga3:XXIII:3.2.1:item:ii}(ii)] Put
    \(\operatorname{Ad}(w_\beta)X_\alpha=X_{\alpha+\beta}\). Then
    \[
    \begin{aligned}
      \operatorname{Ad}(w_\alpha)X_\beta&=-X_{\alpha+\beta},\\
      \operatorname{Ad}(w_\alpha)X_{\alpha+\beta}&=X_\beta,\\
      \operatorname{Ad}(w_\beta)X_{\alpha+\beta}&=-X_\alpha.
    \end{aligned}
    \]
  \item[\SGAInlineRefTarget{sga3:XXIII:3.2.1:item:iii}(iii)] Put
    \[
      p_{\alpha+\beta}(x)
        =\exp(xX_{\alpha+\beta})
        =\operatorname{int}(w_\beta)p_\alpha(x).
    \]
    Then
    \[
      p_\beta(y)p_\alpha(x)
        =p_\alpha(x)p_\beta(y)p_{\alpha+\beta}(xy).
    \]
\end{enumerate}}

\SGARefTarget{sga3:XXIII:3.2.2}
\SGAParagraph{3.2.2}
The proof occupies \SGARefLink{sga3:XXIII:3.2.2}{3.2.2}--\SGARefLink{sga3:XXIII:3.2.7}{3.2.7}. First,
\[
  (\beta^*,\alpha)=(\alpha^*,\beta)=-1,
\]
so \(\alpha(t_\beta)=\beta(t_\alpha)=-1\). Put
\[
  \operatorname{Ad}(w_\beta)X_\alpha=X_{\alpha+\beta}.
\]
Then
\[
  \operatorname{Ad}(w_\beta)X_{\alpha+\beta}
    =\alpha(t_\beta)X_\alpha=-X_\alpha.
\]
Write
\[
  \operatorname{Ad}(w_\alpha)X_\beta=zX_{\alpha+\beta},
  \qquad z\in\mathbf G_m(S).
\]
It follows that
\[
  \operatorname{Ad}(w_\alpha)X_{\alpha+\beta}
    =\beta(t_\alpha)z^{-1}X_\beta=-z^{-1}X_\beta.
\]
By \SGARefLink{sga3:xxii:cor:5.5.2}{Exposé XXII, 5.5.2}, there is a unique
\(A\in\mathbf G_a(S)\) such that
\begin{equation}
  p_\beta(y)p_\alpha(x)
    =p_\alpha(x)p_\beta(y)p_{\alpha+\beta}(Axy).
  \tag{+}
  \label{sga3:XXIII:3.2:eqplus}
\end{equation}
Thus \SGARefLink{sga3:XXIII:3.2.1:item:ii}{(ii)} and \SGARefLink{sga3:XXIII:3.2.1:item:iii}{(iii)} reduce to proving \(A=-z=1\).

\SGARefTarget{sga3:XXIII:3.2.3}
\SGAParagraph{3.2.3}
Conjugating \eqref{sga3:XXIII:3.2:eqplus} by \(w_\beta\) gives
\begin{equation}
  p_{-\beta}(-y)p_{\alpha+\beta}(x)
    =p_{\alpha+\beta}(x)p_{-\beta}(-y)p_\alpha(-Axy).
  \tag{++}
  \label{sga3:XXIII:3.2:eqplusplus}
\end{equation}

\SGARefTarget{sga3:XXIII:3.2.4}
\SGAParagraph{3.2.4}
By definition,
\[
  w_\beta p_\alpha(x)w_\beta^{-1}=p_{\alpha+\beta}(x).
\]
Expanding \(w_\beta\), this is
\[
\begin{aligned}
  &p_\beta(1)p_{-\beta}(-1)p_\beta(1)p_\alpha(x)
    p_\beta(-1)p_{-\beta}(1)p_\beta(-1)\\
  &\hspace{55mm}=p_{\alpha+\beta}(x).
\end{aligned}
\]
Since \(p_\beta\) commutes with \(p_{\alpha+\beta}\), this may be
written
\[
  p_\beta(1)p_\alpha(x)p_\beta(-1)
    =p_{-\beta}(1)p_{\alpha+\beta}(x)p_{-\beta}(-1).
\]
Using \eqref{sga3:XXIII:3.2:eqplus} on the left and
\eqref{sga3:XXIII:3.2:eqplusplus} on the right gives
\[
\begin{aligned}
  p_\alpha(x)p_\beta(1)p_{\alpha+\beta}(Ax)p_\beta(-1)
   ={}&p_{\alpha+\beta}(x)p_{-\beta}(1)\\
     &{}\cdot p_\alpha(Ax)p_{-\beta}(-1).
\end{aligned}
\]
Because \(\alpha+2\beta\), \(\alpha-\beta\), and
\(2\alpha+\beta\) are not roots, this reduces to
\[
  p_\alpha(x)p_{\alpha+\beta}(Ax)
    =p_\alpha(Ax)p_{\alpha+\beta}(x).
\]
\SGARefLink{sga3:xxii:rem:4.1.3}{Exposé XXII, 4.1.3} now gives \(A=1\).

\SGARefTarget{sga3:XXIII:3.2.5}
\SGAParagraph{3.2.5}
Conjugating \eqref{sga3:XXIII:3.2:eqplus} by \(w_\alpha\), and using
\(A=1\), gives
\begin{equation}
  p_{\alpha+\beta}(zy)p_{-\alpha}(-x)
   =p_{-\alpha}(-x)p_{\alpha+\beta}(zy)
      p_\beta(-z^{-1}xy).
  \tag{+++}
  \label{sga3:XXIII:3.2:eqplusplusplus}
\end{equation}

\SGARefTarget{sga3:XXIII:3.2.6}
\SGAParagraph{3.2.6}
As before,
\[
  w_\alpha p_\beta(y)w_\alpha^{-1}=p_{\alpha+\beta}(zy).
\]
Since \(p_\alpha\) and \(p_{\alpha+\beta}\) commute,
\[
  p_\alpha(1)p_\beta(y)p_\alpha(-1)
    =p_{-\alpha}(1)p_{\alpha+\beta}(zy)p_{-\alpha}(-1).
\]
Using \eqref{sga3:XXIII:3.2:eqplus} and
\eqref{sga3:XXIII:3.2:eqplusplusplus}, this becomes
\[
  p_\beta(y)p_{\alpha+\beta}(-y)
    =p_{\alpha+\beta}(zy)p_\beta(-z^{-1}y).
\]
The two factors commute, and hence \(z=-1\).

\SGARefTarget{sga3:XXIII:3.2.7}
\SGAParagraph{3.2.7}
It remains to prove \SGARefLink{sga3:XXIII:3.2.1:item:i}{(i)}. One has
\[
  w_\alpha w_\beta w_\alpha
    =w_\alpha w_\beta w_\alpha^{-1}w_\alpha^2
    =w_{\alpha+\beta}^{-1}t_\alpha,
\]
and therefore
\[
\begin{aligned}
  w_\beta w_\alpha w_\beta w_\alpha w_\beta
    &=w_\beta w_{\alpha+\beta}^{-1}t_\alpha w_\beta\\
    &=w_\beta w_{\alpha+\beta}^{-1}w_\beta^{-1}
      \,s_\beta(t_\alpha)t_\beta\\
    &=w_\alpha t_\alpha t_\beta t_\beta
      =w_\alpha t_\alpha=w_\alpha^{-1}.
\end{aligned}
\]
Thus
\[
  (w_\alpha w_\beta)^3=(w_\beta w_\alpha)^3=e.
\]

\SGARefTarget{sga3:XXIII:3.2.8}
\SGAStatement{Remark 3.2.8}{Put
\(v_\alpha=f_\alpha(u_\alpha)\). Condition (v) of \SGARefLink{sga3:XXIII:2.3}{2.3} becomes
\[
\text{(A)}\quad
\begin{cases}
  \operatorname{int}(h_\alpha h_\beta h_\alpha)v_\beta=v_\beta,\\
  \operatorname{int}(h_\beta h_\alpha h_\beta)v_\alpha=v_\alpha,
\end{cases}
\]
and
\[
\text{(B)}\quad
\begin{cases}
  v_\beta v_\alpha
    =v_\alpha v_\beta\,(h_\beta v_\alpha h_\beta^{-1}),\\
  v_\beta(h_\beta v_\alpha h_\beta^{-1})
    =(h_\beta v_\alpha h_\beta^{-1})v_\beta,\\
  v_\alpha(h_\beta v_\alpha h_\beta^{-1})
    =(h_\beta v_\alpha h_\beta^{-1})v_\alpha.
\end{cases}
\]
Writing \(v_{\alpha+\beta}=\operatorname{int}(h_\beta)v_\alpha\),
the latter conditions are equivalently
\[
\text{(B)}\quad
\begin{cases}
  v_\beta v_\alpha=v_\alpha v_\beta v_{\alpha+\beta},\\
  v_\alpha v_{\alpha+\beta}=v_{\alpha+\beta}v_\alpha,\\
  v_\beta v_{\alpha+\beta}=v_{\alpha+\beta}v_\beta.
\end{cases}
\]}

\SGANextSectionTarget{sga3:xxiii:subsection:3-3}\SGASubsection{3.3}{Groups of Type \(B_2\)}

\SGARefTarget{sga3:XXIII:3.3.1}
\SGAInlineRefTarget{sga3:xxiii:proposition:3-3-1}\SGAStatement{Proposition 3.3.1}{Let \(S\) be a scheme and \(G\) a
pinned \(S\)-group of type \(B_2\). Write
\[
  \Delta=\{\alpha,\beta\},
  \qquad
  R^+=\{\alpha,\beta,\alpha+\beta,2\alpha+\beta\}.
\]
\begin{enumerate}
  \item[\SGAInlineRefTarget{sga3:xxiii:proposition:3-3-1:item:i}(i)]
    \[
      t_{\alpha\beta}=(w_\alpha w_\beta)^4
        =t_\alpha=(w_\beta w_\alpha)^4=t_{\beta\alpha}.
    \]
  \item[\SGAInlineRefTarget{sga3:XXIII:3.3.1:item:ii}\SGAInlineRefTarget{sga3:xxiii:proposition:3-3-1:item:ii}(ii)] If
    \[
      \operatorname{Ad}(w_\beta)X_\alpha=X_{\alpha+\beta},
      \qquad
      \operatorname{Ad}(w_\alpha)X_\beta=X_{2\alpha+\beta},
    \]
    then
    \[
    \begin{aligned}
      \operatorname{Ad}(w_\alpha)X_{\alpha+\beta}
        &=-X_{\alpha+\beta},&
      \operatorname{Ad}(w_\alpha)X_{2\alpha+\beta}&=X_\beta,\\
      \operatorname{Ad}(w_\beta)X_{\alpha+\beta}
        &=-X_\alpha,&
      \operatorname{Ad}(w_\beta)X_{2\alpha+\beta}&=X_{2\alpha+\beta}.
    \end{aligned}
    \]
  \item[\SGAInlineRefTarget{sga3:XXIII:3.3.1:item:iii}(iii)] Put
    \[
    \begin{aligned}
      p_{\alpha+\beta}(x)
        &=\exp(xX_{\alpha+\beta})
          =\operatorname{int}(w_\beta)p_\alpha(x),\\
      p_{2\alpha+\beta}(x)
        &=\exp(xX_{2\alpha+\beta})
          =\operatorname{int}(w_\alpha)p_\beta(x).
    \end{aligned}
    \]
    Then
    \[
    \begin{aligned}
      p_\beta(y)p_\alpha(x)
        &=p_\alpha(x)p_\beta(y)p_{\alpha+\beta}(xy)
          p_{2\alpha+\beta}(x^2y),\\
      p_{\alpha+\beta}(y)p_\alpha(x)
        &=p_\alpha(x)p_{\alpha+\beta}(y)
          p_{2\alpha+\beta}(2xy).
    \end{aligned}
    \]
\end{enumerate}}

\SGARefTarget{sga3:XXIII:3.3.2}
\SGAParagraph{3.3.2}
The proof occupies \SGARefLink{sga3:XXIII:3.3.2}{3.3.2}--\SGARefLink{sga3:XXIII:3.3.6}{3.3.6}. We have
\[
  (\beta^*,\alpha)=-1,\qquad(\alpha^*,\beta)=-2,
\]
so \(\alpha(t_\beta)=-1\) and \(\beta(t_\alpha)=1\). In particular,
\(\beta(t_\alpha)=\alpha(t_\alpha)=1\), so \(t_\alpha\) is a section
of \(\operatorname{Centr}(G)\). Put
\[
  \operatorname{Ad}(w_\beta)X_\alpha=X_{\alpha+\beta},
  \qquad
  \operatorname{Ad}(w_\alpha)X_\beta=X_{2\alpha+\beta}.
\]
Then
\[
\begin{aligned}
  \operatorname{Ad}(w_\beta)X_{\alpha+\beta}
    &=\alpha(t_\beta)X_\alpha=-X_\alpha,\\
  \operatorname{Ad}(w_\alpha)X_{2\alpha+\beta}
    &=\beta(t_\alpha)X_\beta=X_\beta,\\
  \operatorname{Ad}(w_\beta)X_{2\alpha+\beta}
    &=X_{2\alpha+\beta}.
\end{aligned}
\]
For some \(k\in\mathbf G_m(S)\),
\[
  \operatorname{Ad}(w_\alpha)X_{\alpha+\beta}
    =kX_{\alpha+\beta}.
\]
By \SGARefLink{sga3:xxii:cor:5.5.2}{Exposé XXII, 5.5.2}, there are
\(A,B,C\in\mathbf G_a(S)\) such that
\begin{align}
  p_\beta(y)p_\alpha(x)
    &=p_\alpha(x)p_\beta(y)p_{\alpha+\beta}(Axy)
      p_{2\alpha+\beta}(Bx^2y), \tag{1}\label{sga3:XXIII:3.3:eq1}\\
  p_{\alpha+\beta}(y)p_\alpha(x)
    &=p_\alpha(x)p_{\alpha+\beta}(y)
      p_{2\alpha+\beta}(Cxy). \tag{2}\label{sga3:XXIII:3.3:eq2}
\end{align}
It remains to prove
\[
  A=B=1,\qquad C=2,\qquad k=-1.
\]

\SGARefTarget{sga3:XXIII:3.3.3}
\SGAParagraph{3.3.3}
Conjugating \eqref{sga3:XXIII:3.3:eq1} by \(w_\alpha\) gives
\begin{equation}
\begin{aligned}
  p_{2\alpha+\beta}(y)p_{-\alpha}(-x)
    ={}&p_{-\alpha}(-x)p_{2\alpha+\beta}(y)\\
      &{}\cdot p_{\alpha+\beta}(Akxy)p_\beta(Bx^2y).
\end{aligned}
\tag{3}\label{sga3:XXIII:3.3:eq3}
\end{equation}
Transforming \eqref{sga3:XXIII:3.3:eq2} in the same way gives
\begin{equation}
  p_{\alpha+\beta}(ky)p_{-\alpha}(-x)
   =p_{-\alpha}(-x)p_{\alpha+\beta}(ky)p_\beta(Cxy).
\tag{4}\label{sga3:XXIII:3.3:eq4}
\end{equation}
Conjugating \eqref{sga3:XXIII:3.3:eq1} by \(w_\beta\) gives
\begin{equation}
\begin{aligned}
  p_{-\beta}(-y)p_{\alpha+\beta}(x)
    ={}&p_{\alpha+\beta}(x)p_{-\beta}(-y)\\
      &{}\cdot p_\alpha(-Axy)p_{2\alpha+\beta}(Bx^2y).
\end{aligned}
\tag{5}\label{sga3:XXIII:3.3:eq5}
\end{equation}

\SGARefTarget{sga3:XXIII:3.3.4}
\SGAParagraph{3.3.4}
Write
\[
  w_\beta p_\alpha(x)w_\beta^{-1}=p_{\alpha+\beta}(x).
\]
Because \(\alpha+2\beta\) is not a root,
\[
  p_\beta(1)p_\alpha(x)p_\beta(-1)
    =p_{-\beta}(1)p_{\alpha+\beta}(x)p_{-\beta}(-1).
\]
Using \eqref{sga3:XXIII:3.3:eq1} on the left and
\eqref{sga3:XXIII:3.3:eq5} on the right yields
\[
\begin{aligned}
  &p_\alpha(x)p_\beta(1)p_{\alpha+\beta}(Ax)
    p_{2\alpha+\beta}(Bx^2)p_\beta(-1)\\
  &\quad=p_{\alpha+\beta}(x)p_{-\beta}(1)p_\alpha(Ax)
    p_{2\alpha+\beta}(-Bx^2)p_{-\beta}(-1).
\end{aligned}
\]
After commuting the indicated factors and applying
\eqref{sga3:XXIII:3.3:eq2}, this becomes
\[
\begin{aligned}
  p_\alpha(x)p_{\alpha+\beta}(Ax)p_{2\alpha+\beta}(Bx^2)
  ={}&p_\alpha(Ax)p_{\alpha+\beta}(x)\\
     &{}\cdot p_{2\alpha+\beta}((AC-B)x^2).
\end{aligned}
\]
Hence
\[
  A=1,\qquad C=2B.
\]

\SGARefTarget{sga3:XXIII:3.3.5}
\SGAParagraph{3.3.5}
Now write
\[
  w_\alpha p_\beta(y)w_\alpha^{-1}=p_{2\alpha+\beta}(y).
\]
Since \(3\alpha+\beta\) is not a root,
\[
  p_\alpha(1)p_\beta(y)p_\alpha(-1)
    =p_{-\alpha}(1)p_{2\alpha+\beta}(y)p_{-\alpha}(-1).
\]
Using \eqref{sga3:XXIII:3.3:eq1} on the left and
\eqref{sga3:XXIII:3.3:eq3} on the right gives
\[
\begin{aligned}
  p_\beta(y)p_{\alpha+\beta}(-Ay)p_{2\alpha+\beta}(By)
  ={}&p_{2\alpha+\beta}(y)p_{\alpha+\beta}(Aky)p_\beta(By).
\end{aligned}
\]
The three groups commute, so
\[
  B=1,\qquad -A=Ak.
\]
Consequently
\[
  A=B=1,\qquad C=2,\qquad k=-1.
\]

\SGARefTarget{sga3:XXIII:3.3.6}
\SGAParagraph{3.3.6}
We have proved \SGARefLink{sga3:XXIII:3.3.1:item:ii}{(ii)} and \SGARefLink{sga3:XXIII:3.3.1:item:iii}{(iii)}. Since
\[
  s_\beta(t_\alpha)
    =t_\alpha t_\beta^{(\alpha^*,\beta)}
    =t_\alpha,
\]
one obtains successively
\[
\begin{aligned}
  w_\alpha w_\beta w_\alpha
    &=w_{2\alpha+\beta}t_\alpha,\\
  w_\beta w_\alpha w_\beta w_\alpha w_\beta
    &=w_{2\alpha+\beta}t_\alpha t_\beta,\\
  w_\alpha w_\beta w_\alpha w_\beta w_\alpha w_\beta w_\alpha
    &=w_\beta^{-1}t_\alpha.
\end{aligned}
\]
\NDE{The French prints
\((\alpha^*,\beta)=2\); the convention used immediately above gives
\(-2\). The conclusion is unchanged because \(t_\beta^2=e\).}
It follows that
\[
  (w_\beta w_\alpha)^4=t_\alpha,
  \qquad
  (w_\alpha w_\beta)^4=s_\beta(t_\alpha)=t_\alpha.
\]

\SGARefTarget{sga3:XXIII:3.3.7}
\SGAStatement{Remark 3.3.7}{Put
\[
  v_{\alpha+\beta}=\operatorname{int}(h_\beta)v_\alpha,
  \qquad
  v_{2\alpha+\beta}=\operatorname{int}(h_\alpha)v_\beta.
\]
Condition \SGARefLink{sga3:XXIII:2.3:item:v}{(v)} of \SGARefLink{sga3:XXIII:2.3}{2.3} becomes
\[
\text{(A)}\quad
\begin{cases}
  \operatorname{int}(h_\alpha h_\beta h_\alpha)v_\beta=v_\beta,\\
  \operatorname{int}(h_\beta h_\alpha h_\beta)v_\alpha=v_\alpha,
\end{cases}
\]
and
\[
\text{(B)}\quad
\begin{cases}
  v_\beta v_\alpha
    =v_\alpha v_\beta v_{\alpha+\beta}v_{2\alpha+\beta},\\
  v_{\alpha+\beta}v_\alpha
    =v_\alpha v_{\alpha+\beta}v_{2\alpha+\beta}^{\,2},\\
  v_{\alpha+\beta}v_\beta=v_\beta v_{\alpha+\beta},\\
  v_{2\alpha+\beta}v_\alpha=v_\alpha v_{2\alpha+\beta},\\
  v_{2\alpha+\beta}v_\beta=v_\beta v_{2\alpha+\beta},\\
  v_{2\alpha+\beta}v_{\alpha+\beta}
    =v_{\alpha+\beta}v_{2\alpha+\beta}.
\end{cases}
\]}

% END INCLUDED SGA3_English_current_buildable_source_20260823/source/inputs/XXIII/tex/components/06_expose_XXIII_type_A2_and_B2.tex

% BEGIN INCLUDED SGA3_English_current_buildable_source_20260823/source/inputs/XXIII/tex/components/07_expose_XXIII_type_G2.tex
\SGANextSectionTarget{sga3:xxiii:subsection:3-4}\SGASubsection{3.4}{Groups of Type \(G_2\)}

\SGARefTarget{sga3:XXIII:3.4.1}
\SGAInlineRefTarget{sga3:xxiii:proposition:3-4-1}\SGAStatement{Proposition 3.4.1}{Let \(S\) be a scheme and \(G\) a
pinned \(S\)-group of type \(G_2\). Write
\[
  \Delta=\{\alpha,\beta\},\qquad
  R^+=\{\alpha,\beta,\alpha+\beta,2\alpha+\beta,
                 3\alpha+\beta,3\alpha+2\beta\}.
\]
\begin{enumerate}
  \item[\SGAInlineRefTarget{sga3:xxiii:proposition:3-4-1:item:i}\SGAInlineRefTarget{sga3:XXIII:3.4.1:item:i}(i)]
    \[
      t_{\alpha\beta}=(w_\alpha w_\beta)^6
        =e=(w_\beta w_\alpha)^6=t_{\beta\alpha}.
    \]
  \item[\SGAInlineRefTarget{sga3:XXIII:3.4.1:item:ii}\SGAInlineRefTarget{sga3:xxiii:proposition:3-4-1:item:ii}(ii)] If
    \[
    \begin{aligned}
      \operatorname{Ad}(w_\beta)X_\alpha&=X_{\alpha+\beta},&
      \operatorname{Ad}(w_\alpha)X_{\alpha+\beta}&=X_{2\alpha+\beta},\\
      \operatorname{Ad}(w_\alpha)X_\beta&=-X_{3\alpha+\beta},&
      \operatorname{Ad}(w_\beta)X_{3\alpha+\beta}&=X_{3\alpha+2\beta},
    \end{aligned}
    \]
    then
    \[
    \begin{aligned}
      \operatorname{Ad}(w_\alpha)X_{2\alpha+\beta}
        &=-X_{\alpha+\beta},&
      \operatorname{Ad}(w_\alpha)X_{3\alpha+\beta}&=X_\beta,\\
      \operatorname{Ad}(w_\alpha)X_{3\alpha+2\beta}
        &=X_{3\alpha+2\beta},&
      \operatorname{Ad}(w_\beta)X_{\alpha+\beta}&=-X_\alpha,\\
      \operatorname{Ad}(w_\beta)X_{2\alpha+\beta}
        &=X_{2\alpha+\beta},&
      \operatorname{Ad}(w_\beta)X_{3\alpha+2\beta}&=-X_{3\alpha+\beta}.
    \end{aligned}
    \]
  \item[\SGAInlineRefTarget{sga3:XXIII:3.4.1:item:iii}(iii)] Put
    \[
    \begin{aligned}
      p_{\alpha+\beta}(x)
        &=\exp(xX_{\alpha+\beta})
          =\operatorname{int}(w_\beta)p_\alpha(x),\\
      p_{2\alpha+\beta}(x)
        &=\exp(xX_{2\alpha+\beta})
          =\operatorname{int}(w_\alpha w_\beta)p_\alpha(x),\\
      p_{3\alpha+\beta}(x)
        &=\exp(xX_{3\alpha+\beta})
          =\operatorname{int}(w_\alpha)p_\beta(-x),\\
      p_{3\alpha+2\beta}(x)
        &=\exp(xX_{3\alpha+2\beta})
          =\operatorname{int}(w_\beta w_\alpha)p_\beta(-x).
    \end{aligned}
    \]
    Then
    \[
    \begin{aligned}
      p_\beta(y)p_\alpha(x)
        ={}&p_\alpha(x)p_\beta(y)p_{\alpha+\beta}(xy)
          p_{2\alpha+\beta}(x^2y)p_{3\alpha+\beta}(x^3y)
          p_{3\alpha+2\beta}(x^3y^2),\\
      p_{\alpha+\beta}(y)p_\alpha(x)
        ={}&p_\alpha(x)p_{\alpha+\beta}(y)
          p_{2\alpha+\beta}(2xy)p_{3\alpha+\beta}(3x^2y)
          p_{3\alpha+2\beta}(3xy^2),\\
      p_{2\alpha+\beta}(y)p_\alpha(x)
        ={}&p_\alpha(x)p_{2\alpha+\beta}(y)
          p_{3\alpha+\beta}(3xy),\\
      p_{3\alpha+\beta}(y)p_\beta(x)
        ={}&p_\beta(x)p_{3\alpha+\beta}(y)
          p_{3\alpha+2\beta}(-xy),\\
      p_{2\alpha+\beta}(y)p_{\alpha+\beta}(x)
        ={}&p_{\alpha+\beta}(x)p_{2\alpha+\beta}(y)
          p_{3\alpha+2\beta}(3xy).
    \end{aligned}
    \]
\end{enumerate}}

\SGARefTarget{sga3:XXIII:3.4.2}
\SGAParagraph{3.4.2}
The proof occupies \SGARefLink{sga3:XXIII:3.4.2}{3.4.2}--\SGARefLink{sga3:XXIII:3.4.9}{3.4.9}. Since
\[
  (\beta^*,\alpha)=-1,\qquad(\alpha^*,\beta)=-3,
\]
we have \(\beta(t_\alpha)=\alpha(t_\beta)=-1\). Define the four
nonsimple root vectors as in \SGARefLink{sga3:XXIII:3.4.1:item:ii}{(ii)}. Then
\[
\begin{aligned}
  \operatorname{Ad}(w_\beta)X_{\alpha+\beta}
    &=\alpha(t_\beta)X_\alpha=-X_\alpha,\\
  \operatorname{Ad}(w_\alpha)X_{2\alpha+\beta}
    &=\alpha(t_\alpha)\beta(t_\alpha)X_{\alpha+\beta}
      =-X_{\alpha+\beta},\\
  \operatorname{Ad}(w_\alpha)X_{3\alpha+\beta}
    &=-\beta(t_\alpha)X_\beta=X_\beta,\\
  \operatorname{Ad}(w_\beta)X_{3\alpha+2\beta}
    &=\alpha(t_\beta)^3\beta(t_\beta)X_{3\alpha+\beta}
      =-X_{3\alpha+\beta}.
\end{aligned}
\]
Since \((3\alpha+2\beta)\pm\alpha\) and
\((2\alpha+\beta)\pm\beta\) are not roots,
\[
  \operatorname{Ad}(w_\alpha)X_{3\alpha+2\beta}=X_{3\alpha+2\beta},
  \qquad
  \operatorname{Ad}(w_\beta)X_{2\alpha+\beta}=X_{2\alpha+\beta}.
\]

\SGARefTarget{sga3:XXIII:3.4.3}
\SGAParagraph{3.4.3}
By \SGARefLink{sga3:xxii:cor:5.5.2}{Exposé XXII, 5.5.2}, there are sections
\[
  A,B,C,D,E,F,G,H,J\in\mathbf G_a(S)
\]
such that
\begin{align}
 p_\beta(y)p_\alpha(x)
  ={}&p_\alpha(x)p_\beta(y)p_{\alpha+\beta}(Axy)
     p_{2\alpha+\beta}(Bx^2y)p_{3\alpha+\beta}(Cx^3y)
     p_{3\alpha+2\beta}(Dx^3y^2), \tag{1}\label{sga3:XXIII:3.4:eq1}\\
 p_{\alpha+\beta}(y)p_\alpha(x)
  ={}&p_\alpha(x)p_{\alpha+\beta}(y)p_{2\alpha+\beta}(Exy)
     p_{3\alpha+\beta}(Fx^2y)p_{3\alpha+2\beta}(Gxy^2),
     \tag{2}\label{sga3:XXIII:3.4:eq2}\\
 p_{2\alpha+\beta}(y)p_\alpha(x)
  ={}&p_\alpha(x)p_{2\alpha+\beta}(y)p_{3\alpha+\beta}(Hxy),
     \tag{3}\label{sga3:XXIII:3.4:eq3}\\
 p_{3\alpha+\beta}(y)p_\beta(x)
  ={}&p_\beta(x)p_{3\alpha+\beta}(y)p_{3\alpha+2\beta}(Jxy).
     \tag{4}\label{sga3:XXIII:3.4:eq4}
\end{align}

\SGARefTarget{sga3:XXIII:3.4.4}
\SGAParagraph{3.4.4}
Conjugating (1), (3), and (4) by \(w_\beta\) gives
\begin{align}
 p_{-\beta}(-y)p_{\alpha+\beta}(x)
  ={}&p_{\alpha+\beta}(x)p_{-\beta}(-y)p_\alpha(-Axy)
     p_{2\alpha+\beta}(Bx^2y)\notag\\
    &{}\cdot p_{3\alpha+2\beta}(Cx^3y)
     p_{3\alpha+\beta}(-Dx^3y^2), \tag{5}\label{sga3:XXIII:3.4:eq5}\\
 p_{2\alpha+\beta}(y)p_{\alpha+\beta}(x)
  ={}&p_{\alpha+\beta}(x)p_{2\alpha+\beta}(y)
     p_{3\alpha+2\beta}(Hxy), \tag{6}\label{sga3:XXIII:3.4:eq6}\\
 p_{3\alpha+2\beta}(y)p_{-\beta}(-x)
  ={}&p_{-\beta}(-x)p_{3\alpha+2\beta}(y)
     p_{3\alpha+\beta}(-Jxy). \tag{7}\label{sga3:XXIII:3.4:eq7}
\end{align}
Conjugating (1) by \(w_\alpha\) gives
\begin{align}
 p_{3\alpha+\beta}(-y)p_{-\alpha}(-x)
  ={}&p_{-\alpha}(-x)p_{3\alpha+\beta}(-y)
     p_{2\alpha+\beta}(Axy)p_{\alpha+\beta}(-Bx^2y)\notag\\
    &{}\cdot p_\beta(Cx^3y)p_{3\alpha+\beta}(Dx^3y^2).
     \tag{8}\label{sga3:XXIII:3.4:eq8}
\end{align}

\SGARefTarget{sga3:XXIII:3.4.5}
\SGAParagraph{3.4.5}
The identity
\[
  w_\beta p_\alpha(x)w_\beta^{-1}=p_{\alpha+\beta}(x)
\]
may be written, since \(\alpha+2\beta\) is not a root,
\begin{equation}
  p_\beta(1)p_\alpha(x)p_\beta(-1)
   =p_{-\beta}(1)p_{\alpha+\beta}(x)p_{-\beta}(-1).
  \tag{9}\label{sga3:XXIII:3.4:eq9}
\end{equation}
Using (1) and then (4), its left side becomes
\begin{align}
 &p_\alpha(x)p_{\alpha+\beta}(Ax)p_{2\alpha+\beta}(Bx^2)
   p_{3\alpha+\beta}(Cx^3)
   p_{3\alpha+2\beta}((D-CJ)x^3). \tag{10}\label{sga3:XXIII:3.4:eq10}
\end{align}
Using (5) and then \SGARefLink{sga3:XXIII:3.4:eq7}{(7)}, its right side becomes
\[
 p_{\alpha+\beta}(x)p_\alpha(Ax)p_{2\alpha+\beta}(-Bx^2)
 p_{3\alpha+2\beta}(-Cx^3)p_{3\alpha+\beta}((CJ-D)x^3).
\]
Applying (2) transforms this into
\begin{align}
 &p_\alpha(Ax)p_{\alpha+\beta}(x)
 p_{2\alpha+\beta}((AE-B)x^2)
 p_{3\alpha+\beta}((A^2F+CJ-D)x^3)\notag\\
 &\hspace{35mm}\cdot
 p_{3\alpha+2\beta}((AG-C)x^3). \tag{12}\label{sga3:XXIII:3.4:eq12}
\end{align}
Comparing \SGARefLink{sga3:XXIII:3.4:eq10}{(10)} and \SGARefLink{sga3:XXIII:3.4:eq12}{(12)} gives
\[
  A=1,\qquad E=2B,\qquad C+D=F+CJ,\qquad F=G.
\]

\SGARefTarget{sga3:XXIII:3.4.6}
\SGAParagraph{3.4.6}
Write
\[
  w_\alpha p_\beta(y)w_\alpha^{-1}=p_{3\alpha+\beta}(-y).
\]
Since \(4\alpha+\beta\) is not a root,
\begin{equation}
 p_\alpha(1)p_\beta(y)p_\alpha(-1)
  =p_{-\alpha}(1)p_{3\alpha+\beta}(-y)p_{-\alpha}(-1).
 \tag{13}\label{sga3:XXIII:3.4:eq13}
\end{equation}
Using (1), the left side is
\[
 p_\beta(y)p_{\alpha+\beta}(-Ay)p_{2\alpha+\beta}(By)
 p_{3\alpha+\beta}(-Cy)p_{3\alpha+2\beta}(-Dy^2).
\]
Using \SGARefLink{sga3:XXIII:3.4:eq8}{(8)}, (6), and (4), the right side is
\[
\begin{aligned}
 p_\beta(Cy)p_{\alpha+\beta}(-By)p_{2\alpha+\beta}(Ay)
 p_{3\alpha+\beta}(-y)
 p_{3\alpha+2\beta}((D-CJ-ABH)y^2).
\end{aligned}
\]
Thus
\[
  C=1,\qquad A=B,\qquad D-CJ-ABH=-D.
\]
Together with the preceding identities, this yields
\[
 A=B=C=1,\quad E=2,\quad F=G,\quad
 D+1=F+J,\quad 2D=H+J.
\]

\SGARefTarget{sga3:XXIII:3.4.7}
\SGAParagraph{3.4.7}
From
\[
  w_\beta p_{3\alpha+\beta}(x)w_\beta^{-1}
    =p_{3\alpha+2\beta}(x)
\]
one obtains
\[
  p_{3\alpha+\beta}(x)p_{3\alpha+2\beta}(-Jx)
   =p_{3\alpha+2\beta}(x)p_{3\alpha+\beta}(-Jx),
\]
and therefore \(J=-1\).

\SGARefTarget{sga3:XXIII:3.4.8}
\SGAParagraph{3.4.8}
Finally,
\[
  w_\alpha p_{\alpha+\beta}(y)w_\alpha^{-1}=p_{2\alpha+\beta}(y),
\]
or equivalently
\[
 p_\alpha(1)p_{\alpha+\beta}(y)p_\alpha(-1)
  =p_{-\alpha}(1)p_\alpha(-1)p_{2\alpha+\beta}(y)
    p_\alpha(1)p_{-\alpha}(-1).
\]
Using (2) on the left and (3) on the right gives
\[
\begin{aligned}
 &p_{\alpha+\beta}(y)p_{2\alpha+\beta}(-Ey)
   p_{3\alpha+\beta}(Fy)p_{3\alpha+2\beta}(-Gy^2)\\
 &\qquad=p_{-\alpha}(1)p_{2\alpha+\beta}(y)
   p_{3\alpha+\beta}(Hy)p_{-\alpha}(-1).
\end{aligned}
\]
Commuting \(p_{-\alpha}(-1)\) first past
\(p_{3\alpha+\beta}(Hy)\), and then past \(p_{2\alpha+\beta}(y)\),
introduces no new \(p_{3\alpha+\beta}\)-factor. Comparing those
factors gives \(F=H\). Hence
\[
  2D=D+1,\qquad D=1,\qquad F=G=H=2D-J=3.
\]
This determines all the coefficients and proves \SGARefLink{sga3:XXIII:3.4.1:item:iii}{(iii)}.

\SGARefTarget{sga3:XXIII:3.4.9}
\SGAParagraph{3.4.9}
For \SGARefLink{sga3:XXIII:3.4.1:item:i}{(i)}, one obtains successively
\[
\begin{aligned}
 w_\alpha w_\beta w_\alpha
  &=w_{3\alpha+\beta}^{-1}t_\alpha,\\
 w_\beta(w_\alpha w_\beta)^2
  &=w_{3\alpha+2\beta}^{-1}t_\alpha,\\
 w_\alpha(w_\beta w_\alpha)^3
  &=w_{3\alpha+2\beta}^{-1},\\
 w_\beta(w_\alpha w_\beta)^4
  &=w_{3\alpha+\beta}t_\beta,\\
 w_\alpha(w_\beta w_\alpha)^5
  &=w_\beta^{-1}.
\end{aligned}
\]
Consequently
\[
  (w_\alpha w_\beta)^6=(w_\beta w_\alpha)^6=e.
\]

\SGARefTarget{sga3:XXIII:3.4.10}
\SGAStatement{Remark 3.4.10}{Condition \SGARefLink{sga3:XXIII:2.3:item:v}{(v)} of \SGARefLink{sga3:XXIII:2.3}{2.3} consists first of
\[
\text{(A)}\quad
\begin{cases}
 \operatorname{int}(h_\alpha h_\beta h_\alpha h_\beta h_\alpha)v_\beta
   =v_\beta,\\
 \operatorname{int}(h_\beta h_\alpha h_\beta h_\alpha h_\beta)v_\alpha
   =v_\alpha,
\end{cases}
\]
and, after setting
\[
\begin{aligned}
 v_{\alpha+\beta}&=\operatorname{int}(h_\beta)v_\alpha,&
 v_{2\alpha+\beta}&=\operatorname{int}(h_\alpha h_\beta)v_\alpha,\\
 v_{3\alpha+\beta}&=\operatorname{int}(h_\alpha)v_\beta^{-1},&
 v_{3\alpha+2\beta}&=\operatorname{int}(h_\beta h_\alpha)v_\beta^{-1},
\end{aligned}
\]
of the commutation relations
\[
\text{(B)}\quad
\begin{aligned}
 v_\beta v_\alpha
  &=v_\alpha v_\beta v_{\alpha+\beta}v_{2\alpha+\beta}
    v_{3\alpha+\beta}v_{3\alpha+2\beta},\\
 v_{\alpha+\beta}v_\alpha
  &=v_\alpha v_{\alpha+\beta}v_{2\alpha+\beta}^{\,2}
    v_{3\alpha+\beta}^{\,3}v_{3\alpha+2\beta}^{\,3},\\
 v_{2\alpha+\beta}v_\alpha
  &=v_\alpha v_{2\alpha+\beta}v_{3\alpha+\beta}^{\,3},\\
 v_{3\alpha+\beta}v_\alpha&=v_\alpha v_{3\alpha+\beta},&
 v_{3\alpha+2\beta}v_\alpha&=v_\alpha v_{3\alpha+2\beta},\\
 v_{\alpha+\beta}v_\beta&=v_\beta v_{\alpha+\beta},&
 v_{2\alpha+\beta}v_\beta&=v_\beta v_{2\alpha+\beta},\\
 v_{3\alpha+\beta}v_\beta
  &=v_\beta v_{3\alpha+\beta}v_{3\alpha+2\beta}^{-1},\\
 v_{3\alpha+2\beta}v_\beta&=v_\beta v_{3\alpha+2\beta},\\
 v_{2\alpha+\beta}v_{\alpha+\beta}
  &=v_{\alpha+\beta}v_{2\alpha+\beta}v_{3\alpha+2\beta}^{\,3},\\
 v_{3\alpha+\beta}v_{\alpha+\beta}
  &=v_{\alpha+\beta}v_{3\alpha+\beta},\\
 v_{3\alpha+2\beta}v_{\alpha+\beta}
  &=v_{\alpha+\beta}v_{3\alpha+2\beta},\\
 v_{3\alpha+\beta}v_{2\alpha+\beta}
  &=v_{2\alpha+\beta}v_{3\alpha+\beta},\\
 v_{3\alpha+2\beta}v_{2\alpha+\beta}
  &=v_{2\alpha+\beta}v_{3\alpha+2\beta},\\
 v_{3\alpha+2\beta}v_{3\alpha+\beta}
  &=v_{3\alpha+\beta}v_{3\alpha+2\beta}.
\end{aligned}
\]}

% END INCLUDED SGA3_English_current_buildable_source_20260823/source/inputs/XXIII/tex/components/07_expose_XXIII_type_G2.tex

% BEGIN INCLUDED SGA3_English_current_buildable_source_20260823/source/inputs/XXIII/tex/components/08_expose_XXIII_generators_relations_explicit.tex
\SGANextSectionTarget{sga3:xxiii:subsection:3-5}\SGASubsection{3.5}{Explicit Form of the Generators-and-Relations
Theorem}

\SGARefTarget{sga3:XXIII:3.5.1}
\SGAStatement{Theorem 3.5.1}{Let \(S\) be a scheme and \(G\) a
pinned \(S\)-group, with maximal torus \(T\), system of simple roots
\(\Delta\), and elements
\[
  u_\alpha\in U_\alpha(S)^\times,\qquad
  w_\alpha\in\operatorname{Norm}_G(T)(S)
  \qquad(\alpha\in\Delta)
\]
defined by the pinning. Let \(H\) be an fppf \(S\)-sheaf in groups,
let
\[
  f_T:T\to H,\qquad f_\alpha:U_\alpha\to H
  \quad(\alpha\in\Delta)
\]
be morphisms of groups, and let \(h_\alpha\in H(S)\). Put
\[
  v_\alpha=f_\alpha(u_\alpha).
\]
There exists a morphism of groups \(f:G\to H\) extending \(f_T\) and
all the \(f_\alpha\), and satisfying
\(f(w_\alpha)=h_\alpha\), if and only if:
\begin{enumerate}
  \item[(i)] for every \(S'\to S\), \(\alpha\in\Delta\),
    \(t\in T(S')\), and \(x\in U_\alpha(S')\),
    \begin{equation}
      \operatorname{int}(f_T(t))f_\alpha(x)
       =f_\alpha(\operatorname{int}(t)x)
       =f_\alpha(x^{\alpha(t)}); \tag{1}\label{sga3:xxiii:theorem:3-5-1:equation:1}
    \end{equation}
  \item[(ii)] for every \(\alpha\in\Delta\), every \(S'\to S\), and
    every \(t\in T(S')\),
    \begin{equation}
      \operatorname{int}(h_\alpha)f_T(t)
       =f_T(s_\alpha(t))
       =f_T\!\left(t\alpha^*(\alpha(t)^{-1})\right); \tag{2}\label{sga3:xxiii:theorem:3-5-1:equation:2}
    \end{equation}
  \item[(iii)] for every \(\alpha\in\Delta\),
    \begin{equation}
      h_\alpha^2=f_T(\alpha^*(-1)), \tag{\(3_1\)}\label{sga3:xxiii:theorem:3-5-1:equation:3-1}
    \end{equation}
    and
    \begin{equation}
      (h_\alpha v_\alpha)^3=e; \tag{4}\label{sga3:xxiii:theorem:3-5-1:equation:4}
    \end{equation}
  \item[\SGAInlineRefTarget{sga3:XXIII:3.5.1:item:iv}(iv)] for distinct \(\alpha,\beta\in\Delta\):
    \begin{enumerate}
      \item[(a)] the appropriate braid relation is
        \begin{equation}
        \tag{\(3_2\)}\label{sga3:xxiii:theorem:3-5-1:equation:3-2}
        \begin{cases}
        (h_\alpha h_\beta)^2
          =f_T(\alpha^*(-1)\beta^*(-1)),
          &(\alpha^*,\beta)=0,\\
        (h_\alpha h_\beta)^3=e,
          &(\alpha^*,\beta)=-1,\\
        (h_\alpha h_\beta)^4=f_T(\alpha^*(-1)),
          &(\alpha^*,\beta)=-2,\\
        (h_\alpha h_\beta)^6=e,
          &(\alpha^*,\beta)=-3;
        \end{cases}
        \end{equation}
      \item[(b)] the corresponding relations (A) and (B) of
        \SGARefLink{sga3:XXIII:3.1.3}{3.1.3}, \SGARefLink{sga3:XXIII:3.2.8}{3.2.8}, \SGARefLink{sga3:XXIII:3.3.7}{3.3.7}, or \SGARefLink{sga3:XXIII:3.4.10}{3.4.10} hold.
    \end{enumerate}
\end{enumerate}
When it exists, \(f\) is necessarily unique.}

This follows immediately from \SGARefLink{sga3:XXIII:2.3}{Theorem 2.3}, \SGARefLink{sga3:XXIII:2.6}{Remark 2.6}, and the
calculations in each rank-two case.

\SGARefTarget{sga3:XXIII:3.5.2}
\SGAStatement{Remark 3.5.2}{The result may equivalently be presented
by giving morphisms
\[
  a_\alpha:T\!\cdot U_\alpha\longrightarrow H,
  \qquad
  b_\alpha:\operatorname{Norm}_{Z_\alpha}(T)\longrightarrow H
  \qquad(\alpha\in\Delta),
\]
and putting
\[
  h_\alpha=b_\alpha(w_\alpha),
  \qquad
  v_\alpha=a_\alpha(u_\alpha).
\]
The required conditions are:
\begin{enumerate}
  \item[(1)] all the \(a_\alpha\) and \(b_\alpha\) have the same
    restriction to \(T\);
  \item[(2)] condition (4) and all the conditions in \SGARefLink{sga3:XXIII:3.5.1:item:iv}{(iv)} of
    \SGARefLink{sga3:XXIII:3.5.1}{Theorem 3.5.1} hold.
\end{enumerate}}

\SGARefTarget{sga3:XXIII:3.5.3}
\SGAParagraph{3.5.3}
Further applications will be given in the next Exposé. Here is one.
\SGARefLink{sga3:XXIII:3.5.1}{Theorem 3.5.1} describes \(G\) by generators and relations in the
category of fppf \(S\)-sheaves. More explicitly, for every
\(S'\to S\), let \(H(S')\) be the group generated by
\[
  T(S'),\qquad U_\alpha(S')\ (\alpha\in R),\qquad
  w_\alpha\ (\alpha\in R),
\]
subject to relations analogous to
\[
  (1),\ (2),\ (3_1),\ (4),\ (3_2),\ (A),\ (B).
\]
Then \(G\) is the sheaf associated with the presheaf
\(S'\mapsto H(S')\).

In particular, if \(S'\) is the spectrum of an algebraically closed
field \(k\), then \(G(S')=H(S')\). This follows immediately from the
Nullstellensatz in the form: an fppf-covering sieve over an
algebraically closed field is trivial. Thus \SGARefLink{sga3:XXIII:3.5.1}{Theorem 3.5.1} gives an
explicit presentation by generators and relations of the abstract
group \(G(k)\).

% END INCLUDED SGA3_English_current_buildable_source_20260823/source/inputs/XXIII/tex/components/08_expose_XXIII_generators_relations_explicit.tex

% BEGIN INCLUDED SGA3_English_current_buildable_source_20260823/source/inputs/XXIII/tex/components/09_expose_XXIII_fundamental_theorem.tex
\SGANextSectionTarget{sga3:xxiii:section:4}\SGASection{4}{Uniqueness of Pinned Groups: Fundamental Theorem}

\SGARefTarget{sga3:XXIII:4.1}
\SGAInlineRefTarget{sga3:xxiii:theorem:4-1}\SGAStatement{Theorem 4.1}{Let \(S\) be a nonempty scheme. The
functor \(R\) of \SGARefLink{sga3:XXIII:1.6}{1.6} is fully faithful. More precisely, let \(G,G'\)
be pinned \(S\)-groups, let \(p>0\) be an integer such that
\(x\mapsto x^p\) is an endomorphism of \(\mathbf G_{a,S}\), and let
\[
  h:R(G')\longrightarrow R(G)
\]
be a \(p\)-morphism of pinned root data. There is a unique morphism
of pinned groups \(f:G\to G'\) such that \(R(f)=h\).}

Uniqueness was proved in \SGARefLink{sga3:XXIII:1.9}{1.9}, so only existence remains. By
hypothesis there are a bijection \(d:R\xrightarrow{\sim}R'\) and a
map
\[
  q:R\longrightarrow\{p^n\mid n\in\mathbf N\}
\]
such that, for every \(\alpha\in R\),
\[
  h(d(\alpha))=q(\alpha)\alpha,
  \qquad
  {}^th(\alpha^*)=q(\alpha)d(\alpha)^*.
\]
In particular, \(G\) and \(G'\) have the same semisimple rank.

\SGARefTarget{sga3:XXIII:4.1.1}
\SGAParagraph{4.1.1}
Suppose that both semisimple ranks are \(0\). Then \(G,G'\) are
tori:
\[
  G=T=D_S(M),\qquad G'=T'=D_S(M'),
\]
and \(h:M'\to M\) is an ordinary homomorphism. Take \(f=D_S(h)\).

\SGARefTarget{sga3:XXIII:4.1.2}
\SGAParagraph{4.1.2}
Suppose that both semisimple ranks are \(1\), and put
\[
  f_T=D_S(h):T\longrightarrow T'.
\]
Writing \(\alpha'=d(\alpha)\), the diagram
\[
\begin{tikzcd}[column sep=large,row sep=large]
\mathbf G_{m,S}\ar[r,"\alpha^*"]\ar[d,"q(\alpha)"]
  &T\ar[r,"\alpha"]\ar[d,"f_T"']
  &\mathbf G_{m,S}\ar[d,"q(\alpha)"]\\
\mathbf G_{m,S}\ar[r,"\alpha'^*"]
  &T'\ar[r,"\alpha'"]
  &\mathbf G_{m,S}
\end{tikzcd}
\]
commutes. Apply \SGARefLink{sga3:xx:4.1}{Exposé XX, 4.1}.

\SGARefTarget{sga3:XXIII:4.1.3}
\SGAParagraph{4.1.3}
Suppose that both semisimple ranks are \(2\). \SGARefLink{sga3:XXI:7.5.3}{Exposé XXI, 7.5.3}
lists every possible \(h\); we verify the conditions of \SGARefLink{sga3:XXIII:2.5}{Corollary 2.5}
in each case. Write
\[
  \Delta=\{\alpha,\beta\},\qquad
  \Delta'=\{\alpha',\beta'\},\qquad
  d(\alpha)=\alpha',\quad d(\beta)=\beta'.
\]

\SGARefTarget{sga3:XXIII:4.1.4}
\SGAParagraph{4.1.4. \(G\) and \(G'\) of type
\(A_1\times A_1\)}
Here
\[
  h(\alpha')=q\alpha,\qquad h(\beta')=q_1\beta.
\]
Conditions \SGARefLink{sga3:XXIII:2.5:item:ii}{2.5(ii)} and (iii) follow from \SGARefLink{sga3:xxiii:proposition:3-1-2:item:ii}{Proposition 3.1.2(ii)},(iii).
For (i),
\[
\begin{aligned}
 D_S(h)t_{\alpha\beta}
  &=D_S(h)(t_\alpha t_\beta)\\
  &={}^th(\alpha^*)(-1)\,{}^th(\beta^*)(-1)\\
  &=\alpha'^*((-1)^q)\,\beta'^*((-1)^{q_1}).
\end{aligned}
\]
The assumption on \(x\mapsto x^p\), with \(S\ne\varnothing\), implies
that \(p=1\) or \(S\) has characteristic \(p\). Thus
\((-1)^q=-1\): if \(q\) is even, then \(p=2\) and \(1=-1\).
Consequently
\[
  D_S(h)t_{\alpha\beta}
   =\alpha'^*(-1)\beta'^*(-1)=t'_{\alpha'\beta'}.
\]

\SGARefTarget{sga3:XXIII:4.1.5}
\SGAParagraph{4.1.5. \(G\) and \(G'\) of type \(A_2\)}
Here
\[
  h(\alpha')=q\alpha,\qquad h(\beta')=q\beta.
\]
Put
\[
  X_{\alpha+\beta}=\operatorname{Ad}(w_\beta)X_\alpha,
  \qquad
  X'_{\alpha'+\beta'}=\operatorname{Ad}(w'_{\beta'})X'_{\alpha'}.
\]
Condition (i) follows as above from \SGARefLink{sga3:xxiii:proposition:3-2-1:item:i}{Proposition 3.2.1(i)}, and (ii)
follows immediately from \SGARefLink{sga3:xxiii:proposition:3-2-1:item:ii}{Proposition 3.2.1(ii)}. For (iii), the map
\[
\begin{aligned}
 p_\alpha(x)p_\beta(y)p_{\alpha+\beta}(z)
 \longmapsto
 p'_{\alpha'}(x^q)p'_{\beta'}(y^q)p'_{\alpha'+\beta'}(z^q)
\end{aligned}
\]
must be a homomorphism. The only nontrivial relation is
\[
 p_\beta(y)p_\alpha(x)
  =p_\alpha(x)p_\beta(y)p_{\alpha+\beta}(xy),
\]
which maps to
\[
 p'_{\beta'}(y^q)p'_{\alpha'}(x^q)
  =p'_{\alpha'}(x^q)p'_{\beta'}(y^q)
    p'_{\alpha'+\beta'}(x^qy^q).
\]

\SGARefTarget{sga3:XXIII:4.1.6}
\SGAParagraph{4.1.6}
The same argument treats types \(B_2\) and \(G_2\) when the
radicial exponents are equal, using \SGARefLink{sga3:xxiii:proposition:3-3-1}{Propositions 3.3.1} and \SGARefLink{sga3:xxiii:proposition:3-4-1}{3.4.1}.
It remains to treat the two exceptional cases in \SGARefLink{sga3:XXI:7.5.3}{Exposé XXI, 7.5.3}.

\SGARefTarget{sga3:XXIII:4.1.7}
\SGAParagraph{4.1.7. Type \(B_2\), characteristic \(2\),
\(q(\alpha)=2q\), \(q(\beta)=q\)}
The positive roots are
\[
 \{\alpha,\beta,\alpha+\beta,2\alpha+\beta\}
\]
and
\[
 \{\alpha',\beta',\alpha'+\beta',\alpha'+2\beta'\},
\]
where \(\alpha\) and \(\beta'\) are the short simple roots. We have
\[
\begin{aligned}
 h(\alpha')&=2q\alpha,&h(\beta')&=q\beta,\\
 h(\alpha'+\beta')&=q(2\alpha+\beta),&
 h(\alpha'+2\beta')&=2q(\alpha+\beta),
\end{aligned}
\]
so
\[
\begin{aligned}
 d(\alpha+\beta)&=\alpha'+2\beta',&q(\alpha+\beta)&=2q,\\
 d(2\alpha+\beta)&=\alpha'+\beta',&q(2\alpha+\beta)&=q.
\end{aligned}
\]
Put
\[
\begin{aligned}
 X_{\alpha+\beta}&=\operatorname{Ad}(w_\beta)X_\alpha,&
 X_{2\alpha+\beta}&=\operatorname{Ad}(w_\alpha)X_\beta,\\
 X'_{\alpha'+\beta'}&=\operatorname{Ad}(w'_{\alpha'})X'_{\beta'},&
 X'_{\alpha'+2\beta'}&=\operatorname{Ad}(w'_{\beta'})X'_{\alpha'}.
\end{aligned}
\]

For (i), characteristic \(2\) gives \(-1=1\), hence
\[
 t_{\alpha\beta}=t_\alpha=\alpha^*(-1)=e
   =\beta'^*(-1)=t'_{\beta'}=t'_{\alpha'\beta'}.
\]
For (ii), the two defining identities are
\[
\begin{aligned}
 \operatorname{Ad}(w_\alpha)X_\beta
  &=X_{2\alpha+\beta},&
 \operatorname{Ad}(w'_{\alpha'})X'_{\beta'}
  &=X'_{\alpha'+\beta'}=X'_{d(2\alpha+\beta)},\\
 \operatorname{Ad}(w_\beta)X_\alpha
  &=X_{\alpha+\beta},&
 \operatorname{Ad}(w'_{\beta'})X'_{\alpha'}
  &=X'_{\alpha'+2\beta'}=X'_{d(\alpha+\beta)}.
\end{aligned}
\]
\SGARefLink{sga3:xxiii:proposition:3-3-1:item:ii}{Proposition 3.3.1(ii)}, together with \(-1=1\), gives the four
remaining identities:
\[
\begin{aligned}
 \operatorname{Ad}(w_\alpha)X_{\alpha+\beta}&=X_{\alpha+\beta},&
 \operatorname{Ad}(w'_{\alpha'})X'_{d(\alpha+\beta)}
   &=X'_{d(\alpha+\beta)},\\
 \operatorname{Ad}(w_\alpha)X_{2\alpha+\beta}&=X_\beta,&
 \operatorname{Ad}(w'_{\alpha'})X'_{d(2\alpha+\beta)}
   &=X'_{d(\beta)},\\
 \operatorname{Ad}(w_\beta)X_{\alpha+\beta}&=X_\alpha,&
 \operatorname{Ad}(w'_{\beta'})X'_{d(\alpha+\beta)}
   &=X'_{d(\alpha)},\\
 \operatorname{Ad}(w_\beta)X_{2\alpha+\beta}&=X_{2\alpha+\beta},&
 \operatorname{Ad}(w'_{\beta'})X'_{d(2\alpha+\beta)}
   &=X'_{d(2\alpha+\beta)}.
\end{aligned}
\]

For (iii), the sole nontrivial relations are
\[
\begin{aligned}
 p_\beta(y)p_\alpha(x)
  &=p_\alpha(x)p_\beta(y)p_{\alpha+\beta}(xy)
    p_{2\alpha+\beta}(x^2y),\\
 p'_{\alpha'}(y')p'_{\beta'}(x')
  &=p'_{\beta'}(x')p'_{\alpha'}(y')
    p'_{\alpha'+\beta'}(x'y')
    p'_{\alpha'+2\beta'}(x'y'^2).
\end{aligned}
\]
The proposed map is
\[
\begin{aligned}
 &p_\alpha(x)p_\beta(y)p_{\alpha+\beta}(z)p_{2\alpha+\beta}(t)\\
 &\quad\longmapsto
 p'_{\alpha'}(x^{2q})p'_{\beta'}(y^q)
 p'_{\alpha'+2\beta'}(z^{2q})p'_{\alpha'+\beta'}(t^q).
\end{aligned}
\]
It respects the first relation because the second, with
\(y'=x^{2q}\) and \(x'=y^q\), reads
\[
\begin{aligned}
 p'_{\beta'}(y^q)p'_{\alpha'}(x^{2q})
  ={}&p'_{\alpha'}(x^{2q})p'_{\beta'}(y^q)
   p'_{\alpha'+2\beta'}((xy)^{2q})\\
   &{}\cdot p'_{\alpha'+\beta'}((x^2y)^q).
\end{aligned}
\]

\SGARefTarget{sga3:XXIII:4.1.8}
\SGAParagraph{4.1.8. Type \(G_2\), characteristic \(3\),
\(q(\alpha)=3q\), \(q(\beta)=q\)}
The positive roots are
\[
\{\alpha,\beta,\alpha+\beta,2\alpha+\beta,3\alpha+\beta,
  3\alpha+2\beta\}
\]
and
\[
\{\alpha',\beta',\alpha'+\beta',\alpha'+2\beta',
  \alpha'+3\beta',2\alpha'+3\beta'\},
\]
where again \(\alpha\) and \(\beta'\) are short. We have
\[
\begin{aligned}
 h(\alpha')&=3q\alpha,&h(\beta')&=q\beta,\\
 h(\alpha'+\beta')&=q(3\alpha+\beta),&
 h(\alpha'+2\beta')&=q(3\alpha+2\beta),\\
 h(\alpha'+3\beta')&=3q(\alpha+\beta),&
 h(2\alpha'+3\beta')&=3q(2\alpha+\beta),
\end{aligned}
\]
and hence
\[
\begin{aligned}
 d(\alpha+\beta)&=\alpha'+3\beta',&q(\alpha+\beta)&=3q,\\
 d(2\alpha+\beta)&=2\alpha'+3\beta',&q(2\alpha+\beta)&=3q,\\
 d(3\alpha+\beta)&=\alpha'+\beta',&q(3\alpha+\beta)&=q,\\
 d(3\alpha+2\beta)&=\alpha'+2\beta',&q(3\alpha+2\beta)&=q.
\end{aligned}
\]
Define
\[
\begin{aligned}
 X_{\alpha+\beta}&=\operatorname{Ad}(w_\beta)X_\alpha,&
 X_{2\alpha+\beta}&=\operatorname{Ad}(w_\alpha)X_{\alpha+\beta},\\
 X_{3\alpha+\beta}&=-\operatorname{Ad}(w_\alpha)X_\beta,&
 X_{3\alpha+2\beta}&=\operatorname{Ad}(w_\beta)X_{3\alpha+\beta},\\
 X'_{\alpha'+\beta'}&=-\operatorname{Ad}(w'_{\alpha'})X'_{\beta'},&
 X'_{\alpha'+2\beta'}&=\operatorname{Ad}(w'_{\beta'})X'_{\alpha'+\beta'},\\
 X'_{\alpha'+3\beta'}&=\operatorname{Ad}(w'_{\beta'})X'_{\alpha'},&
 X'_{2\alpha'+3\beta'}&=\operatorname{Ad}(w'_{\alpha'})X'_{\alpha'+3\beta'}.
\end{aligned}
\]
Condition (i) follows from \SGARefLink{sga3:xxiii:proposition:3-4-1:item:i}{Proposition 3.4.1(i)}. Four instances of
(ii) hold by definition:
\[
\begin{aligned}
 \operatorname{Ad}(w_\beta)X_\alpha&=X_{\alpha+\beta},&
 \operatorname{Ad}(w'_{\beta'})X'_{\alpha'}&=X'_{d(\alpha+\beta)},\\
 \operatorname{Ad}(w_\alpha)X_\beta&=-X_{3\alpha+\beta},&
 \operatorname{Ad}(w'_{\alpha'})X'_{\beta'}&=-X'_{d(3\alpha+\beta)},\\
 \operatorname{Ad}(w_\beta)X_{3\alpha+\beta}&=X_{3\alpha+2\beta},&
 \operatorname{Ad}(w'_{\beta'})X'_{d(3\alpha+\beta)}
   &=X'_{d(3\alpha+2\beta)},\\
 \operatorname{Ad}(w_\alpha)X_{\alpha+\beta}&=X_{2\alpha+\beta},&
 \operatorname{Ad}(w'_{\alpha'})X'_{d(\alpha+\beta)}
   &=X'_{d(2\alpha+\beta)}.
\end{aligned}
\]
The remaining instances follow from \SGARefLink{sga3:xxiii:proposition:3-4-1:item:ii}{Proposition 3.4.1(ii)}; for
example,
\[
\begin{aligned}
 \operatorname{Ad}(w_\alpha)X_{2\alpha+\beta}&=-X_{\alpha+\beta},&
 \operatorname{Ad}(w'_{\alpha'})X'_{d(2\alpha+\beta)}
   &=-X'_{d(\alpha+\beta)},\\
 \operatorname{Ad}(w_\beta)X_{3\alpha+2\beta}&=-X_{3\alpha+\beta},&
 \operatorname{Ad}(w'_{\beta'})X'_{d(3\alpha+2\beta)}
   &=-X'_{d(3\alpha+\beta)}.
\end{aligned}
\]

For (iii), in characteristic \(3\), the nontrivial relations in
\(U\) are
\[
\begin{aligned}
 p_\beta(y)p_\alpha(x)
  ={}&p_\alpha(x)p_\beta(y)p_{\alpha+\beta}(xy)
  p_{2\alpha+\beta}(x^2y)p_{3\alpha+\beta}(x^3y)
  p_{3\alpha+2\beta}(x^3y^2),\\
 p_{\alpha+\beta}(y)p_\alpha(x)
  ={}&p_\alpha(x)p_{\alpha+\beta}(y)p_{2\alpha+\beta}(-xy),\\
 p_{3\alpha+\beta}(y)p_\beta(x)
  ={}&p_\beta(x)p_{3\alpha+\beta}(y)p_{3\alpha+2\beta}(-xy).
\end{aligned}
\]
The corresponding relations in \(U'\) are
\[
\begin{aligned}
 p'_{\alpha'}(y')p'_{\beta'}(x')
  ={}&p'_{\beta'}(x')p'_{\alpha'}(y')
  p'_{\alpha'+\beta'}(-x'y')p'_{\alpha'+2\beta'}(-x'^2y')\\
  &{}\cdot p'_{\alpha'+3\beta'}(-x'^3y')
  p'_{2\alpha'+3\beta'}(-x'^3y'^2),\\
 p'_{\alpha'+\beta'}(y')p'_{\beta'}(x')
  ={}&p'_{\beta'}(x')p'_{\alpha'+\beta'}(y')
  p'_{\alpha'+2\beta'}(x'y'),\\
 p'_{\alpha'+3\beta'}(y')p'_{\alpha'}(x')
  ={}&p'_{\alpha'}(x')p'_{\alpha'+3\beta'}(y')
  p'_{2\alpha'+3\beta'}(x'y').
\end{aligned}
\]
Define \(\phi:U\to U'\) by
\[
\begin{aligned}
 &\phi\!\left(
 p_\alpha(x)p_\beta(y)p_{\alpha+\beta}(t)p_{2\alpha+\beta}(u)
 p_{3\alpha+\beta}(v)p_{3\alpha+2\beta}(w)\right)\\
 &\quad=
 p'_{\alpha'}(x^{3q})p'_{\beta'}(y^q)
 p'_{\alpha'+3\beta'}(t^{3q})p'_{2\alpha'+3\beta'}(u^{3q})
 p'_{\alpha'+\beta'}(v^q)p'_{\alpha'+2\beta'}(w^q).
\end{aligned}
\]
Substituting \(x^{3q},y^q\) in the three displayed relations for
\(U'\) yields exactly the images of the three relations for \(U\);
thus \(\phi\) is a homomorphism.

\SGARefTarget{sga3:XXIII:4.1.9}
\SGAParagraph{4.1.9. Semisimple rank greater than \(2\)}
For \(\alpha\in\Delta\), put \(\alpha'=d(\alpha)\). For every
\((\alpha,\beta)\), consider the pinned rank-at-most-two groups
\[
  Z_{\alpha\beta},\qquad Z'_{\alpha'\beta'}.
\]
The homomorphism \(M'\to M\) underlying \(h\) induces a
\(p\)-morphism
\[
  h_{\alpha\beta}:R(Z'_{\alpha'\beta'})
    \longrightarrow R(Z_{\alpha\beta}).
\]
By the preceding cases, there is a morphism of pinned groups
\[
  f_{\alpha\beta}:Z_{\alpha\beta}
    \longrightarrow Z'_{\alpha'\beta'}
\]
with \(R(f_{\alpha\beta})=h_{\alpha\beta}\). The restrictions
\[
  f_{\alpha\beta}|_{Z_\alpha},
  \qquad f_{\alpha\alpha}:Z_\alpha\to Z'_{\alpha'}
\]
correspond to the same morphism of pinned root data, and therefore
coincide by uniqueness. \SGARefLink{sga3:XXIII:2.4}{Corollary 2.4} now glues the
\(f_{\alpha\beta}\) to a morphism \(f:G\to G'\). It is a morphism of
pinned groups and satisfies \(R(f)=h\), completing the proof.

% END INCLUDED SGA3_English_current_buildable_source_20260823/source/inputs/XXIII/tex/components/09_expose_XXIII_fundamental_theorem.tex

% BEGIN INCLUDED SGA3_English_current_buildable_source_20260823/source/inputs/XXIII/tex/components/10_expose_XXIII_corollaries.tex
\SGANextSectionTarget{sga3:xxiii:section:5}\SGASection{5}{Corollaries of the Fundamental Theorem}

The most important is the following.

\SGARefTarget{sga3:XXIII:5.1}
\SGAStatement{Corollary 5.1}{Let \(S\) be a nonempty scheme, let
\(G,G'\) be pinned \(S\)-groups, and let
\[
  h:R(G')\xrightarrow{\sim}R(G)
\]
be an isomorphism of pinned root data. There is a unique isomorphism
of pinned \(S\)-groups \(f:G\xrightarrow{\sim}G'\) such that
\(R(f)=h\).}

This also follows from \SGARefLink{sga3:XXIII:3.5.1}{Theorem 3.5.1}, because its relations can be
written using only \(R(G)\). It also follows from the elementary part
of the proof of \SGARefLink{sga3:xxiii:theorem:4-1}{Theorem 4.1}, without the exceptional isogenies of
\SGARefLink{sga3:XXIII:4.1.7}{4.1.7} and \SGARefLink{sga3:XXIII:4.1.8}{4.1.8}.

\SGARefTarget{sga3:XXIII:5.2}
\SGAStatement{Corollary 5.2 (Uniqueness theorem)}{Let \(S\) be a
scheme and let \(G,G'\) be two splittable \(S\)-groups. If they have
the same type, then they are isomorphic.}

\SGARefTarget{sga3:XXIII:5.3}
\SGAStatement{Corollary 5.3}{Let \(S\) be a scheme and \(G,G'\) two
splittable \(S\)-groups. The following are equivalent:
\begin{enumerate}
  \item[\SGAInlineRefTarget{sga3:XXIII:5.3:item:i}(i)] \(G\) and \(G'\) are isomorphic;
  \item[\SGAInlineRefTarget{sga3:XXIII:5.3:item:ii}(ii)] they are locally isomorphic for the fpqc topology;
  \item[\SGAInlineRefTarget{sga3:XXIII:5.3:item:iii}(iii)] there is \(s\in S\) such that \(G_s\) and \(G'_s\)
    have the same type.
\end{enumerate}}

Clearly \SGARefLink{sga3:XXIII:5.3:item:i}{(i)} implies \SGARefLink{sga3:XXIII:5.3:item:ii}{(ii)}, which implies \SGARefLink{sga3:XXIII:5.3:item:iii}{(iii)}. Conversely, \SGARefLink{sga3:XXIII:5.3:item:iii}{(iii)}
gives
\[
  R(G)=R(G_s)=R(G'_s)=R(G'),
\]
so \SGARefLink{sga3:XXIII:5.2}{Corollary 5.2} applies.

\SGARefTarget{sga3:XXIII:5.4}
\SGAStatement{Corollary 5.4 (Uniqueness of Chevalley schemes)}{Let
\(G,G'\) be reductive groups over \(\mathbf Z\), each possessing a
split maximal torus. The following are equivalent:
\begin{enumerate}
  \item[(i)] \(G\) and \(G'\) are isomorphic;
  \item[(ii)] there is \(s\in\operatorname{Spec}\mathbf Z\) such
    that \(G_s\) and \(G'_s\) have the same type;
  \item[(iii)] \(G_{\mathbf C}\) and \(G'_{\mathbf C}\) have the same
    type.
\end{enumerate}}

Indeed, \(G\) and \(G'\) are splittable by \SGARefLink{sga3:xxii:prop:2.2}{Exposé XXII, 2.2}.
\NDE{In fact every \(\mathbf Z\)-torus is split: every
\(\mathbf Z\)-torus is isotrivial, and every étale cover of
\(\operatorname{Spec}\mathbf Z\) is trivial.}

\SGARefTarget{sga3:XXIII:5.5}
\SGAStatement{Corollary 5.5 (Existence of outer automorphisms)}{Let
\(S\) be a scheme, \(G\) a pinned \(S\)-group, and \(h\) an
automorphism of the pinned root datum \(R(G)\). There is a unique
automorphism \(u\) of \(G\) preserving the pinning and satisfying
\(R(u)=h\).}

\SGARefTarget{sga3:XXIII:5.5bis}
\SGAStatement{Corollary 5.5 bis}{Let \(S\) be a scheme, \(G\) a
split reductive \(S\)-group, and \(R^+\) a system of positive roots.
For every simple root \(\alpha\), choose an isomorphism of vector
groups
\[
  p_\alpha:\mathbf G_{a,S}\xrightarrow{\sim}U_\alpha.
\]
Let \(h\) be an automorphism of \(M\) permuting the positive roots
and the corresponding coroots: if \(\alpha\in R^+\), then
\[
  h(\alpha)\in R^+,\qquad h^\vee(\alpha^*)=h(\alpha)^*.
\]
There is a unique automorphism \(u\) of \(G\) inducing \(D_S(h)\) on
\(T\) and satisfying \(u\circ p_\alpha=p_{h(\alpha)}\) for every
simple root \(\alpha\).}

\SGARefTarget{sga3:XXIII:5.6}
\SGAStatement{Corollary 5.6}{Let \(G,G'\) be reductive
\(S\)-groups. The following are equivalent:
\begin{enumerate}
  \item[\SGAInlineRefTarget{sga3:XXIII:5.6:item:i}(i)] \(G,G'\) are locally isomorphic for the fpqc topology;
  \item[(i bis)] they are locally isomorphic for the étale topology;
  \item[(ii)] \(G_s,G'_s\) are isomorphic for every \(s\in S\);
  \item[\SGAInlineRefTarget{sga3:XXIII:5.6:item:iii}(iii)] the functions \(s\mapsto\operatorname{type}(G_s)\)
    and \(s\mapsto\operatorname{type}(G'_s)\) are equal.
\end{enumerate}}

The implication (i bis)\(\Rightarrow\)\SGARefLink{sga3:XXIII:5.6:item:i}{(i)} is immediate;
(i)\(\Rightarrow\)(ii) follows from the principle of finite
extension (EGA IV\(_3\), 9.1.4); and
(ii)\(\Rightarrow\)(iii) is immediate. For
\SGARefLink{sga3:XXIII:5.6:item:iii}{(iii)}\(\Rightarrow\)(i bis), one may suppose \(G,G'\) splittable
(\SGARefLink{sga3:xxii:cor:2.3}{Exposé XXII, 2.3}), and then apply \SGARefLink{sga3:XXIII:5.3}{Corollary 5.3}.

\SGARefTarget{sga3:XXIII:5.7}
\SGAStatement{Corollary 5.7}{Let \(S\) be a scheme, \(G\) a
reductive \(S\)-group, and \(G'\) an affine smooth \(S\)-group with
connected fibers. If \(s\in S\) and \(G_s\simeq G'_s\), then there
is an étale morphism \(S'\to S\), covering \(s\), such that
\(G_{S'}\simeq G'_{S'}\).}

By Exposés XIX, \SGARefLink{sga3:XXIII:2.5}{2.5} and XXII, \SGARefLink{sga3:xxii:prop:2.1}{2.1}, one may suppose \(G,G'\)
reductive and splittable, and \SGARefLink{sga3:XXIII:5.3}{Corollary 5.3} applies.

\SGARefTarget{sga3:XXIII:5.8}
\SGAStatement{Corollary 5.8}{Let \(k\) be a field and \(G,G'\)
reductive \(k\)-groups. The following are equivalent:
\begin{enumerate}
  \item[(i)] \(G,G'\) have the same type;
  \item[(ii)] \(G\otimes_k\bar k\) and
    \(G'\otimes_k\bar k\) are isomorphic;
  \item[(iii)] for some finite separable extension \(K/k\),
    \(G\otimes_kK\) and \(G'\otimes_kK\) are isomorphic.
\end{enumerate}}

\SGARefTarget{sga3:XXIII:5.9}
\SGAStatement{Corollary 5.9}{Let \(S\) be a nonempty scheme and
\(R\) a root datum. The following are equivalent:
\begin{enumerate}
  \item[\SGAInlineRefTarget{sga3:XXIII:5.9:item:i}(i)] there is a pinned \(S\)-group of type \(R\);
  \item[(ii)] there is an \(S\)-group of type \(R\);
  \item[\SGAInlineRefTarget{sga3:XXIII:5.9:item:iii}(iii)] fpqc-locally, there is a reductive \(S\)-group of
    type \(R\).
\end{enumerate}}

Only \SGARefLink{sga3:XXIII:5.9:item:iii}{(iii)}\(\Rightarrow\)\SGARefLink{sga3:XXIII:5.9:item:i}{(i)} needs proof. Suppose, for simplicity,
that there are a faithfully flat quasi-compact morphism \(S'\to S\)
and a reductive \(S'\)-group \(G'\) of type \(R\). We may suppose
\(G'\) splittable and choose a pinning \(E'\), writing
\(R=R(G',E')\). The two pullbacks of \((G',E')\) to
\[
  S''=S'\times_SS'
\]
are pinned groups \((G''_1,E''_1)\), \((G''_2,E''_2)\), with
canonical isomorphisms
\[
  p_i:R(G''_i,E''_i)\xrightarrow{\sim}R.
\]
Thus
\[
  p=p_2^{-1}\circ p_1:
  R(G''_1,E''_1)\xrightarrow{\sim}R(G''_2,E''_2).
\]
The uniqueness theorem gives a unique isomorphism
\[
  f:(G''_1,E''_1)\xrightarrow{\sim}(G''_2,E''_2)
\]
with \(R(f)=p\). This is a descent datum on \(G'\). Its cocycle
condition may be checked after applying \(R\), since \(R\) is fully
faithful, and then follows from the construction of \(p\). Because
\(G'\) is affine, the datum is effective. The pinning is stable
under descent, so it descends to a pinned \(S\)-group \((G,E)\) of
type \(R\).

\medskip
\noindent\textit{N.B.}\quad
In the language of fibered categories, the preceding proof becomes
both shorter and more transparent.

\SGARefTarget{sga3:XXIII:5.10}
\SGAStatement{Corollary 5.10}{Let \(S\) be a nonempty scheme and
\(R\) a pinned root datum for which there exists a reductive
\(S\)-group of type \(R\). Then there is a pinned \(S\)-group of
type \(R\), unique up to unique isomorphism.}

\SGARefTarget{sga3:XXIII:5.11}
\SGAStatement{Definition 5.11}{Under these hypotheses, denote the
unique pinned \(S\)-group of type \(R\) by
\(\operatorname{\acute Ep}_S(R)\), its canonical maximal torus by
\(T_S(R)\), its canonical Borel subgroup by \(B_S(R)\), and so on.}

For every morphism \(S'\to S\), with \(S'\ne\varnothing\), one may
identify
\[
  \operatorname{\acute Ep}_{S'}(R)
    =\operatorname{\acute Ep}_S(R)\times_SS'.
\]
If \(\operatorname{\acute Ep}_{\operatorname{Spec}\mathbf Z}(R)\)
exists, write it \(\operatorname{\acute Ep}(R)\); then
\[
  \operatorname{\acute Ep}_S(R)
    =\operatorname{\acute Ep}(R)
      \times_{\operatorname{Spec}\mathbf Z}S.
\]
The group \(\operatorname{\acute Ep}(R)\) is called the
\emph{Chevalley group scheme of type \(R\)}.

\SGARefTarget{sga3:XXIII:5.12}
\SGAParagraph{5.12}
Thus an \(S\)-sheaf in groups \(G\) is a reductive \(S\)-group of
type \(R\) if and only if it is locally isomorphic, for the étale or
fpqc topology, to \(\operatorname{\acute Ep}_S(R)\). Similarly, by
the conjugacy theorems, giving a reductive \(S\)-group of type \(R\)
with a maximal torus is locally equivalent to
\[
  (\operatorname{\acute Ep}_S(R),T_S(R)).
\]
The analogous statements hold for Borel subgroups and Killing pairs.

% END INCLUDED SGA3_English_current_buildable_source_20260823/source/inputs/XXIII/tex/components/10_expose_XXIII_corollaries.tex

% BEGIN INCLUDED SGA3_English_current_buildable_source_20260823/source/inputs/XXIII/tex/components/11_expose_XXIII_chevalley_systems_and_bibliography.tex
\SGANextSectionTarget{sga3:xxiii:section:6}\SGASection{6}{Chevalley Systems}

The explicit calculations of \SGARefLink{sga3:xxiii:section:3}{Section 3} have important numerical
consequences.

\SGARefTarget{sga3:XXIII:6.1}
\SGAStatement{Definition 6.1}{Let \(S\) be a scheme and
\((G,T,M,R)\) a split \(S\)-group. A \emph{Chevalley system} of
\(G\) is a family
\[
  (X_\alpha)_{\alpha\in R},
  \qquad
  X_\alpha\in\Gamma(S,\mathfrak g^\alpha)^\times,
\]
satisfying:
\begin{description}
  \item[(SC)] for every pair \(\alpha,\beta\in R\),
  \[
    \operatorname{Ad}(w_\alpha(X_\alpha))X_\beta
      =\pm X_{s_\alpha(\beta)}.
  \]
\end{description}}

Recall that
\[
  w_\alpha(X_\alpha)
    =\exp(X_\alpha)\exp(-X_\alpha^{-1})\exp(X_\alpha).
\]
Condition (SC) implies, in particular,
\[
  X_{-\alpha}=\pm X_\alpha^{-1}
\]
by the identity
\[
  \operatorname{Ad}(w_\alpha(X_\alpha))X_\alpha=-X_\alpha^{-1}.
\]

\SGARefTarget{sga3:XXIII:6.2}
\SGAStatement{Proposition 6.2}{Every split group admits a Chevalley
system. More precisely, if
\[
  \bigl(\Delta,(X'_\alpha)_{\alpha\in\Delta}\bigr)
\]
is a pinning of \((G,T,M,R)\), there is a Chevalley system
\((X_\alpha)_{\alpha\in R}\) with \(X_\alpha=X'_\alpha\) for
\(\alpha\in\Delta\).}

First, it suffices to verify (SC) for \(\alpha\in\Delta\). Indeed,
for every \(\alpha\in R\), choose a sequence
\(\alpha_i\in\Delta\) such that
\[
  \alpha=s_{\alpha_1}\cdots s_{\alpha_n}(\alpha_{n+1}).
\]
Repeated use of (SC) gives
\[
 X_\alpha
  =\pm\operatorname{Ad}\!\left(
    w_{\alpha_1}(X_{\alpha_1})\cdots
    w_{\alpha_n}(X_{\alpha_n})\right)X_{\alpha_{n+1}}.
\]
By \SGARefLink{sga3:xxiii:lemma:3-1-1:item:iii}{Lemma 3.1.1(iii)},
\[
\begin{aligned}
 &w_{\alpha_1}(X_{\alpha_1})\cdots w_{\alpha_n}(X_{\alpha_n})
  w_{\alpha_{n+1}}(X_{\alpha_{n+1}})
  w_{\alpha_n}(X_{\alpha_n})^{-1}\cdots
  w_{\alpha_1}(X_{\alpha_1})^{-1}\\
 &\hspace{42mm}=\alpha^*(\pm1)w_\alpha(X_\alpha).
\end{aligned}
\]
Now
\[
 w_{\alpha_i}(X_{\alpha_i})^{-1}
  =\alpha_i^*(-1)w_{\alpha_i}(X_{\alpha_i}),
\]
and
\[
 \beta(\gamma^*(-1))=(-1)^{(\gamma^*,\beta)}=\pm1.
\]
It follows that (SC) for every pair \((\alpha_i,\gamma)\) implies
(SC) for every \((\alpha,\beta)\).

For each \(\alpha\in R\), choose such a sequence and define
\[
 X_\alpha=
 \operatorname{Ad}\!\left(
  w_{\alpha_1}(X_{\alpha_1})\cdots
  w_{\alpha_n}(X_{\alpha_n})\right)X'_{\alpha_{n+1}}.
\]
It remains to prove the following lemma.

\SGARefTarget{sga3:XXIII:6.3}
\SGAStatement{Lemma 6.3}{Let
\[
 (G,T,M,R,\Delta,(X_\alpha)_{\alpha\in\Delta})
\]
be a pinned \(S\)-group. If \(\alpha_i\), \(0\le i\le n+1\), are
simple roots such that
\[
  s_{\alpha_1}\cdots s_{\alpha_n}(\alpha_{n+1})=\alpha_0,
\]
then
\[
  \operatorname{Ad}(w_{\alpha_1}\cdots w_{\alpha_n})
    X_{\alpha_{n+1}}=\pm X_{\alpha_0}.
\]}

Arguing as in \SGARefLink{sga3:XXIII:2.3.4}{2.3.4}, using Tits's lemma, reduces the proof to:
\begin{enumerate}
  \item[\SGAInlineRefTarget{sga3:XXIII:6.3:item:a}(a)] \(G\) has semisimple rank at most \(2\); or
  \item[\SGAInlineRefTarget{sga3:XXIII:6.3:item:b}(b)] \(w_{\alpha_1}\cdots w_{\alpha_n}\) is a section of
    \(T\).
\end{enumerate}
Case \SGARefLink{sga3:XXIII:6.3:item:a}{(a)} is exactly the calculation in \SGARefLink{sga3:XXIII:3.1.2:item:i}{3.1.2(i)}, \SGARefLink{sga3:xxiii:proposition:3-2-1:item:i}{3.2.1(i)},
\SGARefLink{sga3:xxiii:proposition:3-3-1:item:i}{3.3.1(i)}, or \SGARefLink{sga3:xxiii:proposition:3-4-1:item:i}{3.4.1(i)}.

In case \SGARefLink{sga3:XXIII:6.3:item:b}{(b)}, it suffices to show that if
\[
  s_{\alpha_1}\cdots s_{\alpha_n}=1,
\]
then \(t=w_{\alpha_1}\cdots w_{\alpha_n}\) satisfies
\(\alpha(t)=\pm1\) for every root \(\alpha\). By the structure of the
Weyl group, the word in the free group generated by the
\(s_\alpha\) lies in the normal subgroup generated by
\[
  (s_\alpha s_\beta)^{n_{\alpha\beta}}.
\]
Thus one reduces to a word
\[
 s_{\alpha_1}\cdots s_{\alpha_i}
 (s_{\alpha_{i+1}}s_{\alpha_{i+2}})
 ^{n_{\alpha_{i+1}\alpha_{i+2}}}
 s_{\alpha_i}\cdots s_{\alpha_1},
\]
for which
\[
  t=s_{\alpha_1}\cdots s_{\alpha_i}
      (t_{\alpha_{i+1}\alpha_{i+2}}).
\]
The assertion is now immediate from the values of
\(t_{\alpha_1\alpha_2}\) computed in \SGARefLink{sga3:XXIII:3.1.2:item:i}{3.1.2(i)}, \SGARefLink{sga3:xxiii:proposition:3-2-1:item:i}{3.2.1(i)},
\SGARefLink{sga3:xxiii:proposition:3-3-1:item:i}{3.3.1(i)}, and \SGARefLink{sga3:xxiii:proposition:3-4-1:item:i}{3.4.1(i)}.

\SGARefTarget{sga3:XXIII:6.4}
\SGAStatement{Proposition 6.4}{Let \((G,T,M,R)\) be a split
\(S\)-group with Chevalley system \((X_\alpha)_{\alpha\in R}\).
Let \(\alpha,\beta\) be nonproportional roots, with
\[
  \operatorname{length}(\alpha)\le\operatorname{length}(\beta),
  \qquad
  |(\beta^*,\alpha)|\le|(\alpha^*,\beta)|.
\]
Let \(q\) and \(p-1\) be the nonnegative integers for which the roots
of the form \(\beta+k\alpha\) are
\[
  \{\beta-(p-1)\alpha,\ldots,\beta,\ldots,\beta+q\alpha\}.
\]
Thus \(-(\alpha^*,\beta)=q-p+1\). Writing
\(p_\gamma(x)=\exp(xX_\gamma)\), the commutation relation between
\(U_\alpha\) and \(U_\beta\) is given by:
\[
\begin{array}{c|l}
(p,q)&p_\beta(y)p_\alpha(x)p_\beta(-y)p_\alpha(-x)\\ \hline
(-,0)&e\\
(1,1)&p_{\alpha+\beta}(\pm xy)\\
(1,2)&p_{\alpha+\beta}(\pm xy)p_{2\alpha+\beta}(\pm x^2y)\\
(1,3)&p_{\alpha+\beta}(\pm xy)p_{2\alpha+\beta}(\pm x^2y)
       p_{3\alpha+\beta}(\pm x^3y)p_{3\alpha+2\beta}(\pm x^3y^2)\\
(2,1)&p_{\alpha+\beta}(\pm2xy)\\
(2,2)&p_{\alpha+\beta}(\pm2xy)p_{2\alpha+\beta}(\pm3x^2y)
       p_{\alpha+2\beta}(\pm3xy^2)\\
(3,1)&p_{\alpha+\beta}(\pm3xy)
\end{array}
\]
These are all possibilities, since the root string has length at
most \(3\).}

\SGAProof
By \SGARefLink{sga3:XXI:3.5.4}{Exposé XXI, 3.5.4}, there is a system of simple roots \(\Delta\)
such that \(\alpha\in\Delta\), and there are
\(\alpha'\in\Delta\), \(a,b\in\mathbf Q_+\), with
\[
  \beta=a\alpha+b\alpha'.
\]
The desired relation lies in \(U_{\alpha\alpha'}\), so it follows
from the rank-two calculations of Section 3 and condition (SC).

\SGARefTarget{sga3:XXIII:6.5}
\SGAStatement{Corollary 6.5 (Chevalley's rule)}{Let
\((X_\alpha)_{\alpha\in R}\) be a Chevalley system of the split
\(S\)-group \(G\). If \(\alpha,\beta,\alpha+\beta\in R\), then
\[
  [X_\alpha,X_\beta]=\pm pX_{\alpha+\beta},
\]
where \(p\) is the least positive integer for which
\(\beta-p\alpha\) is not a root.}

The assertion is symmetric in \(\alpha,\beta\); after arranging
\(\operatorname{length}(\alpha)\le\operatorname{length}(\beta)\),
it follows from \SGARefLink{sga3:XXIII:6.4}{Proposition 6.4}.

\SGARefTarget{sga3:XXIII:6.6}
\SGAStatement{Corollary 6.6}{Let \(S\) be a scheme such that
\(6\cdot1_S\ne0\), and let \((G,T,M,R)\) be a split \(S\)-group.
If \(R'\subset R\) is such that
\begin{equation}
  \mathfrak g_{R'}
   =\mathfrak t\oplus\bigoplus_{\alpha\in R'}\mathfrak g^\alpha
  \tag{*}\label{sga3:XXIII:6.6:star}
\end{equation}
is a Lie subalgebra of \(\mathfrak g\), then \(R'\) is a closed
subset of \(R\).}

Choose a Chevalley system and let \(\alpha,\beta\in R'\) with
\(\alpha+\beta\in R\). By \SGARefLink{sga3:XXIII:6.5}{Corollary 6.5} and \SGARefLink{sga3:XXIII:6.4}{Proposition 6.4},
\[
  [X_\alpha,X_\beta]=\pm pX_{\alpha+\beta},
  \qquad p\in\{1,2,3\}.
\]
Since neither \(2\) nor \(3\) is zero on \(S\),
\eqref{sga3:XXIII:6.6:star} implies
\(\mathfrak g^{\alpha+\beta}\subset\mathfrak g_{R'}\), and hence
\(\alpha+\beta\in R'\).
\NDE{It suffices that \(2\) and \(3\) are nonzero on \(S\); for
example the result applies to \(S=\operatorname{Spec}(\mathbf Z/6\mathbf Z)\).
In the excluded characteristics, short-root subsystems furnish
standard counterexamples.}

\SGARefTarget{sga3:XXIII:6.7}
\SGAParagraph{6.7}
The exact signs in this section can be determined by studying
\(\operatorname{Norm}_G(T)\), and more precisely the extended Weyl
group
\[
  W^*=N(\mathbf Z),\qquad
  N=\operatorname{Norm}_{\operatorname{\acute{E}p}_{\mathbf Z}(R)}
         (T_{\mathbf Z}(R)).
\]
It is an extension of \(W(R)\) by an abelian group of type
\((2,2,\ldots,2)\), which accounts for the signs.

\SGARefTarget{sga3:XXIII:6.7.1}
\SGAStatement{Remark 6.7.1}{By 3.1.2(i) and 3.\(n\).1,
\(n=2,3,4\), the elements \(w_\alpha,w_\beta\in N(\mathbf Z)\)
satisfy the braid relations
\[
  w_\alpha w_\beta\cdots=w_\beta w_\alpha\cdots
\]
with \(n_{\alpha\beta}\) factors on each side.}

See also Tits \cite{Ti66}, Bourbaki \cite{BLie}, IX.4, Ex.~12, and
Springer \cite[9.3.2]{Sp98}.

\begin{thebibliography}{99}
\bibitem[BLie]{BLie}
N.~Bourbaki,
\emph{Groupes et algèbres de Lie}, Chap.~IX,
Masson, 1982.

\bibitem[Ja87]{Ja87}
J.~C. Jantzen,
\emph{Representations of Algebraic Groups},
Academic Press, 1987; 2nd ed., American Mathematical Society, 2003.

\bibitem[Sp98]{Sp98}
T.~A. Springer,
\emph{Linear Algebraic Groups}, 2nd ed.,
Birkhäuser, 1998.

\bibitem[St62]{St62}
R.~Steinberg,
\emph{Générateurs, relations et revêtements de groupes algébriques},
Colloque sur la théorie des groupes algébriques (Bruxelles, 1962),
Université de Louvain and Gauthier--Villars, 1962;
pp.~133--147 in \emph{R. Steinberg Collected Papers},
American Mathematical Society, 1997.

\bibitem[St67]{St67}
R.~Steinberg,
\emph{Lectures on Chevalley Groups},
Yale University, 1967.

\bibitem[Ta83]{Ta83}
M.~Takeuchi,
\emph{A hyperalgebraic proof of the isomorphism and isogeny
theorems for reductive groups},
J. Algebra \textbf{85} (1983), 179--196.

\bibitem[Ti66]{Ti66}
J.~Tits,
\emph{Normalisateurs de tores I. Groupes de Coxeter étendus},
J. Algebra \textbf{4} (1966), 96--116.
\end{thebibliography}

% END INCLUDED SGA3_English_current_buildable_source_20260823/source/inputs/XXIII/tex/components/11_expose_XXIII_chevalley_systems_and_bibliography.tex



% END INCLUDED SGA3_English_current_buildable_source_20260823/source/inputs/XXIII/body.tex


\SGAStartExpose{XXIV}
% BEGIN INCLUDED SGA3_English_current_buildable_source_20260823/source/inputs/XXIV/body.tex
% Mechanically extracted after \mainmatter from the frozen standalone master.

\markboth{EXPOSÉ XXIV. AUTOMORPHISMS OF REDUCTIVE GROUPS}
  {EXPOSÉ XXIV. AUTOMORPHISMS OF REDUCTIVE GROUPS}
% BEGIN INCLUDED SGA3_English_current_buildable_source_20260823/source/inputs/XXIV/tex/components/01_expose_XXIV_title_intro.tex
% SGA 3, Exposé XXIV, Loop-1 English reconstruction.
% Current scope: authority local pages 1--3 / combined-reader pages 1246--1248.
% Local GPU OCR and Reinhold's English Markdown are locator/comparison
% witnesses only. The Polo--Gille page image controls text and formulas.

\setcounter{footnote}{0}
\SGARefTarget{sga3:xxiv:chapter:xxiv}\chapter*{Exposé XXIV\\
Automorphisms of Reductive Groups}
\addcontentsline{toc}{chapter}{Exposé XXIV. Automorphisms of Reductive
Groups}
\markboth{EXPOSÉ XXIV. AUTOMORPHISMS OF REDUCTIVE GROUPS}
  {EXPOSÉ XXIV. AUTOMORPHISMS OF REDUCTIVE GROUPS}

\begin{center}
  \textit{by Michel Demazure}
\end{center}

\begingroup
\renewcommand{\thefootnote}{0}
\footnote{Editor's note: Version dated October 13, 2024.}
\addtocounter{footnote}{-1}
\endgroup
The first part of this Exposé (nos.~1 to 5) is a direct consequence of
the existence, for a reductive group, of ``sufficiently many outer
automorphisms'', a result which is itself a consequence of the weakest
form of the isomorphism theorem for pinned groups. The second part
(nos.~6 and 7) sets out two applications of the more precise results of
the preceding Exposé; in particular, no.~7 uses the generators-and-
relations theorem in its explicit form. Finally, we have included in
an appendix (no.~8) results on ``Galois'' cohomology used in the text.

Let us specify our cohomological notation. If \(S\) is a scheme and
\(G\) an \(S\)-group scheme, we write \(H^1(S,G)\) for the first
cohomology set of \(S\) with coefficients in \(G\), computed for the
\textup{(fpqc)} topology; this is also the set of isomorphism classes
of principal homogeneous \textup{(fpqc)} sheaves under \(G\). We write
\(H^1_{\mathrm{\acute et}}(S,G)\) for the corresponding set for the
étale topology; it is therefore the part of \(H^1(S,G)\) formed by
classes of homogeneous sheaves under \(G\) which are quasi-isotrivial
\textup{(that is, locally trivial for the étale topology)}. We write
\(\Fib(S,G)\) for the part of \(H^1(S,G)\) formed by classes of
representable sheaves \textup{(principal homogeneous bundles)}. Thus
we have the inclusions
\[
  H^1_{\mathrm{\acute et}}(S,G)\subset H^1(S,G),
  \qquad
  \Fib(S,G)\subset H^1(S,G).
\]

If every principal homogeneous sheaf under \(G\) is representable
\textup{(for example, if \(G\) is quasi-affine over \(S\); cf.\ SGA~1,
VIII, \SGARefLink{sga3:viii:proposition:7:9}{7.9})}, then
\[
  \Fib(S,G)=H^1(S,G).
\]

If \(S'\to S\) is a covering morphism for the \textup{(fpqc)}
topology, we write \(H^1(S'/S,G)\) for the kernel of the canonical map
\[
  H^1(S,G)\longrightarrow H^1(S',G_{S'}).
\]
It is known that \(H^1(S'/S,G)\) can be computed simplicially
\textup{(TDTE~I, \S A.4)}, which implies that, when \(S'\to S\) is a
covering for the étale topology, \(H^1(S'/S,G)\) is also the kernel of
\[
  H^1_{\mathrm{\acute et}}(S,G)
  \longrightarrow
  H^1_{\mathrm{\acute et}}(S',G_{S'}).
\]

Finally, following \SGARefLink{sga3:viii:corollary:4:5}{Exp.~VIII, 4.5}, the name ``\SGARefLink{sga3:viii:corollary:4:5}{Theorem~90}'' is given to
the following assertion: ``every principal homogeneous sheaf under
\(G_{m,S}\) is representable and locally trivial.'' This assertion is
equivalent to
\[
  H^1(S,G_{m,S})=\operatorname{Pic}(S),
\]
or again to
\[
  H^1(S,G_{m,S})=0
\]
when \(S\) is local \textup{(or, more generally, semilocal)}.

\SGARefTarget{sga3:xxiv:section:1}
\section*{1. Scheme of Automorphisms of a Reductive Group}
\addcontentsline{toc}{section}{1. Scheme of Automorphisms of a Reductive Group}

\SGARefTarget{sga3:xxiv:1.0}
\subsection*{1.0.}
We must first make precise certain definitions from the preceding
Exposé. Let \(S\) be a nonempty scheme, let \(G\) be a reductive
\(S\)-group, and let
\[
  \mathscr R=(M,M^*,R,R^*,\Delta)
\]
be a pinned reduced root datum \textup{(\SGARefLink{sga3:XXIII:1.5}{Exp.~XXIII, 1.5})}. A
\emph{pinning of \(G\) of type \(\mathscr R\)}, or an
\(\mathscr R\)-\emph{pinning of \(G\)}, consists of:
\begin{enumerate}
\item[(i)] an isomorphism of \(D_S(M)\) onto a maximal torus \(T\) of
  \(G\) \textup{(or, equivalently, a monomorphism
  \(D_S(M)\to G\) whose image is a maximal torus \(T\) of \(G\))},
  identifying \(R\) with a root system of \(G\) relative to \(T\)
  \textup{(\SGARefLink{sga3:xix:definition:3-6}{Exp.~XIX, 3.6})} and \(R^*\) with the corresponding set of
  coroots;
\item[(ii)] for each \(\alpha\in\Delta\), an element
  \[
    X_\alpha\in\Gamma(S,\mathfrak g^\alpha)^\times.
  \]
\end{enumerate}
The group \(G\) possesses an \(\mathscr R\)-pinning if and only if it
is split and of type \(\mathscr R\) \textup{(\SGARefLink{sga3:xxii:def:2.7}{Exp.~XXII, 2.7})}.

If \(u:G\to G'\) is an isomorphism of reductive \(S\)-groups, then
every \(\mathscr R\)-pinning \(\mathcal E\) of \(G\) determines, by
``transport of structure'', an \(\mathscr R\)-pinning
\(u(\mathcal E)\) of \(G'\). If
\(v:\mathscr R'\to\mathscr R\) is an isomorphism of pinned root data,
then every \(\mathscr R'\)-pinning \(\mathcal E'\) of \(G'\)
determines, by transport of structure, an \(\mathscr R\)-pinning
\(v(\mathcal E')\) of \(G\).

Call a \emph{pinned group} a triple \((G,\mathscr R,\mathcal E)\),
where \(G\) is a reductive \(S\)-group, \(\mathscr R\) is a pinned
reduced root datum, and \(\mathcal E\) is a pinning of \(G\) of type
\(\mathscr R\). An isomorphism from the pinned group
\((G,\mathscr R,\mathcal E)\) onto the pinned group
\((G',\mathscr R',\mathcal E')\) is a pair \((u,v)\), where
\(u:G\to G'\) is an isomorphism and
\(v:\mathscr R'\to\mathscr R\) is an isomorphism of pinned root data,
such that
\[
  u(\mathcal E)=v(\mathcal E').
\]
\footnote{Editor's note: Thus
\(\Lie(u)(X_{v(\alpha)})=X'_\alpha\) for every
\(\alpha\in R'\).}

\textit{N.B.} If \(S\) is nonempty, \(v\) is uniquely determined by
\(u\), and by an abuse of language \(u\) will also be called an
isomorphism of the pinned groups. In particular, if
\((G,\mathscr R,\mathcal E)\) is a pinned group, an automorphism of
\((G,\mathscr R,\mathcal E)\) is therefore an automorphism \(u\) of
\(G\) for which there exists an automorphism \(v\) of \(\mathscr R\)
such that \(u(\mathcal E)=v(\mathcal E)\). It is thus an automorphism
of \(G\) which normalizes \(T\), induces on \(T\) an automorphism of
the form \(D_S(h)\), where \(h\) is an automorphism of \(M\), and
permutes the elements \(X_\alpha\), \(\alpha\in\Delta\).
\textup{(As is readily seen, the preceding conditions in fact
characterize the automorphisms of
\((G,\mathscr R,\mathcal E)\).)}

There is an evident contravariant functor
\[
  \Phi:\quad
  (G,\mathscr R,\mathcal E)\longmapsto\mathscr R,
  \qquad
  (u,v)\longmapsto v,
\]
and the principal result of the preceding Exposé
\textup{(\SGARefLink{sga3:xxiii:theorem:4-1}{Exp.~XXIII, 4.1})} shows that this functor is fully faithful
\textup{(we shall moreover see in the following Exposé that it is an
equivalence of categories)}. In particular, it follows that the
automorphism group of \((G,\mathscr R,\mathcal E)\) is canonically
isomorphic to the automorphism group of the pinned root datum
\(\mathscr R\) \textup{(cf.\ \SGARefLink{sga3:XXIII:5.5}{Exp.~XXIII, 5.5})}.

\SGARefTarget{sga3:xxiv:1.1}
\subsection*{1.1.}
Let \(S\) be a scheme; endow \((\mathbf{Sch})/S\) with the
\textup{(fpqc)} topology, and consider the \(S\)-sheaf of groups
\(\underline{\Aut}_{S\text{-}\mathrm{gr.}}(G)\), where \(G\) is an
\(S\)-group scheme. There is an exact sequence of \(S\)-sheaves of
groups
\[
  1\longrightarrow \underline{\Centr}(G)
  \longrightarrow G
  \xrightarrow{\mathrm{int}}
  \underline{\Aut}_{S\text{-}\mathrm{gr.}}(G),
\]
which defines a monomorphism
\[
  j:G/\underline{\Centr}(G)
  \longrightarrow \underline{\Aut}_{S\text{-}\mathrm{gr.}}(G).
\]
The image sheaf of \(j\) is the sheaf of inner automorphisms of \(G\).
An automorphism \(u\) of \(G\) is inner if and only if there is a
covering family \(\{S_i\to S\}\) and, for every \(i\), an element
\(g_i\in G(S_i)\) such that
\[
  \operatorname{int}(g_i)=u_{S_i}.
\]
In this case, if \(v\) is another automorphism of \(G\), one sees at
once that
\[
  \operatorname{int}(v)u=vuv^{-1}
\]
is the inner automorphism defined by the family
\(g'_i=v(g_i)\). It follows that the image of \(j\) is normal in
\(\underline{\Aut}_{S\text{-}\mathrm{gr.}}(G)\). The quotient sheaf of
groups, denoted \(\underline{\Autext}(G)\), is the sheaf of outer
automorphisms of \(G\). Thus there is an exact sequence
\[
  1\longrightarrow G/\underline{\Centr}(G)
  \longrightarrow \underline{\Aut}_{S\text{-}\mathrm{gr.}}(G)
  \longrightarrow \underline{\Autext}(G)
  \longrightarrow 1.
\]
All the preceding definitions are compatible with base change. They
are naturally valid in any site.

\SGARefTarget{sga3:xxiv:1.2}
\subsection*{1.2.}
Let \(S\) be a scheme and let \((G,\mathscr R,\mathcal E)\) be a
pinned reductive \(S\)-group. Let \(E\) be the \textup{(abstract)}
automorphism group of the pinned root datum \(\mathscr R\), that is,
the group of automorphisms of \(\mathscr R\) which normalize
\(\Delta\). By \SGARefLink{sga3:XXIII:5.5}{Exp.~XXIII, 5.5}, there is a canonical monomorphism
\[
  E\longrightarrow \underline{\Aut}_{S\text{-}\mathrm{gr.}}(G)
\]
which associates with \(h\in E\) the unique automorphism \(u\) of the
pinned group \(G\) such that \(\Phi(u)=h\). This monomorphism
canonically defines a monomorphism of sheaves
\[
  \tag{*}\label{sga3:xxiv:subsection:1-2:equation:star}
  a:E_S\longrightarrow \underline{\Aut}_{S\text{-}\mathrm{gr.}}(G).
\]
An automorphism \(u\) of \(G\) is a section of the image sheaf of \(a\)
if and only if the following conditions hold:
\begin{enumerate}
\item[(i)] \(u\) normalizes \(T\). It is then known that \(u\)
  permutes the roots of \(G\) relative to \(T\). If \(\alpha\in R\),
  then
  \[
    u(\alpha):t\longmapsto\alpha\bigl(u^{-1}(t)\bigr)
  \]
  is therefore locally on \(S\) of the form
  \(t\longmapsto\beta(t)\), with \(\beta\in R\). The second condition
  may then be written as follows:
\item[(ii)] If \(\alpha\in\Delta\) and \(U\) is an open subset of
  \(S\) such that \(u(\alpha)_U\in R\), then
  \(u(\alpha)_U\in\Delta\) and
  \[
    \Lie(u_U)(X_\alpha)_U=(X_{u(\alpha)})_U.
  \]
\end{enumerate}
It follows at once from the definitions that the sections of
\(a(E_S)\) normalize the subgroups of \(G\) defined by the pinning:
\[
  T,\quad B,\quad B^{-},\quad U,\quad U^{-}.
\]

With these definitions in place, we have:

% END INCLUDED SGA3_English_current_buildable_source_20260823/source/inputs/XXIV/tex/components/01_expose_XXIV_title_intro.tex

% BEGIN INCLUDED SGA3_English_current_buildable_source_20260823/source/inputs/XXIV/tex/components/02_expose_XXIV_theorem13_lemma15_uniqueness.tex
% SGA 3, Exposé XXIV, authority local page 4 / combined-reader page 1249.
% Personally checked against the Polo--Gille page image, band by band.

\SGARefTarget{sga3:xxiv:1.3}
\paragraph{\SGARefLink{sga3:xxiv:1.3}{Theorem 1.3}.}
\emph{Let \(S\) be a scheme and \(G\) a reductive \(S\)-group. Consider
the canonical exact sequence of \(S\)-sheaves of groups}
\footnote{Editor's note: Recall that
\(\operatorname{ad}(G)=G/\underline{\Centr}(G)\) denotes the adjoint
group of \(G\); cf.\ \SGARefLink{sga3:xxii:def:4.3.6}{Exp.~XXII, 4.3.6}.}
\[
  1\longrightarrow \operatorname{ad}(G)
  \longrightarrow \underline{\Aut}_{S\text{-}\mathrm{gr.}}(G)
  \xrightarrow{p}\underline{\Autext}(G)
  \longrightarrow 1.
\]
\emph{Then:}
\begin{enumerate}
\item[\SGAInlineRefTarget{sga3:xxiv:1.3:item:i}(i)] \(\underline{\Aut}_{S\text{-}\mathrm{gr.}}(G)\) is
  representable by a smooth separated \(S\)-scheme;
\item[\SGAInlineRefTarget{sga3:xxiv:1.3:item:ii}(ii)] \(\underline{\Autext}(G)\) is representable by a twisted
  constant \(S\)-scheme of finite generation \textup{(\SGARefLink{sga3:x:definition:5-1}{Exp.~X, 5.1})};
\item[\SGAInlineRefTarget{sga3:xxiv:1.3:item:iii}(iii)] if \(G\) is split, the preceding exact sequence is split.
  More precisely, for every pinning of \(G\), the morphism
  \textup{(cf.\ \SGARefLink{sga3:xxiv:1.2}{1.2}~\((*)\))}
  \[
    p\circ a:E_S\longrightarrow\underline{\Autext}(G)
  \]
  is an isomorphism.
\end{enumerate}

Let us first show how the theorem follows from the following lemma.

\SGARefTarget{sga3:xxiv:1.4}
\paragraph{\SGARefLink{sga3:xxiv:1.4}{Lemma 1.4}.}
\emph{Under the hypotheses of \textup{\SGARefLink{sga3:xxiv:1.3:item:iii}{(iii)}},
\(\underline{\Aut}_{S\text{-}\mathrm{gr.}}(G)\) is the semidirect
product}
\[
  a(E_S)\cdot\operatorname{ad}(G).
\]
\footnote{Editor's note: Here and below, the notation
\(\operatorname{int}(G)\) has been replaced by
\(\operatorname{ad}(G)\).}

The lemma immediately implies the theorem when \(G\) is split. Since
\(G\) is locally split for the étale topology \textup{(\SGARefLink{sga3:xxii:cor:2.3}{Exp.~XXII,
2.3})}, hence also for the \textup{(fppf)} topology, and since the
latter is ``of effective descent'' for the fibred category of twisted
constant morphisms \textup{(\SGARefLink{sga3:x:corollary:5-5}{Exp.~X, 5.5})}, assertion \textup{\SGARefLink{sga3:xxiv:1.3:item:ii}{(ii)}}
follows in the general case \textup{(cf.\ \SGARefLink{sga3:iv:proposition:4:6:8}{Exp.~IV, 4.6.8})}. To deduce
\textup{\SGARefLink{sga3:xxiv:1.3:item:i}{(i)}}, observe that \(\operatorname{ad}(G)\) is affine over
\(S\); thus \(p\) is affine whenever
\(\underline{\Aut}_{S\text{-}\mathrm{gr.}}(G)\) is representable, and
the result follows by descent for affine schemes.
\footnote{Editor's note: Cf.\ SGA~1, VIII, \SGARefLink{sga3:viii:proposition:2:1}{2.1}.}

It remains only to prove \SGARefLink{sga3:xxiv:1.4}{1.4}. For this, it suffices to prove:

\SGARefTarget{sga3:xxiv:1.5}
\paragraph{\SGARefLink{sga3:xxiv:1.5}{Lemma 1.5}.}
\emph{If \((\mathscr R,\mathcal E)\) and
\((\mathscr R',\mathcal E')\) are two pinnings of the reductive
\(S\)-group \(G\), there exists a unique inner automorphism \(u\) of
\(G\) over \(S\) carrying one pinning to the other; that is, such that
there exists an isomorphism
\[
  v:\mathscr R'\xrightarrow{\sim}\mathscr R
\]
for which \(u(\mathcal E)=v(\mathcal E')\); cf.\ \SGARefLink{sga3:xxiv:1.0}{1.0}.}

\SGARefTarget{sga3:xxiv:1.5.1}
\subsubsection*{1.5.1. Uniqueness.}
It suffices to prove that, if \(G\) is a pinned \(S\)-group and
\(\operatorname{int}(g)\) is an automorphism of the pinned group
\textup{(\(g\in G(S)\))}, then \(\operatorname{int}(g)=\operatorname{id}\).
First,
\[
  \operatorname{int}(g)T=T,\qquad
  \operatorname{int}(g)B=B,
\]
and hence
\[
  g\in
  \underline{\operatorname{Norm}}_G(T)(S)
  \cap\underline{\operatorname{Norm}}_G(B)(S)
  =T(S)
\]
\textup{(cf., for example, \SGARefLink{sga3:xxii:prop:5.6.1}{Exp.~XXII, 5.6.1})}. It follows that
\(\operatorname{int}(g)\) normalizes every \(U_\alpha\) and that
\[
  \Lie(\operatorname{int}(g))X_\alpha
  =\operatorname{Ad}(g)X_\alpha
  =\alpha(g)X_\alpha
\]
for every \(\alpha\in\Delta\). Thus \(\alpha(g)=1\) for
\(\alpha\in\Delta\), and therefore
\[
  g\in
  \bigcap_{\alpha\in\Delta}(\ker\alpha)(S)
  =\underline{\Centr}(G)(S)
\]
\textup{(\SGARefLink{sga3:xxii:rem:4.1.8}{Exp.~XXII, 4.1.8})}. \qed

% END INCLUDED SGA3_English_current_buildable_source_20260823/source/inputs/XXIV/tex/components/02_expose_XXIV_theorem13_lemma15_uniqueness.tex

% BEGIN INCLUDED SGA3_English_current_buildable_source_20260823/source/inputs/XXIV/tex/components/03_expose_XXIV_lemma15_existence_corollaries16_18.tex
% SGA 3, Exposé XXIV, authority local page 5 / combined-reader page 1250.
% Personally checked against the Polo--Gille page image, band by band.

\SGARefTarget{sga3:xxiv:1.5.2}
\subsubsection*{1.5.2. Existence.}
It suffices to prove existence locally for the \textup{(fpqc)}
topology. Let
\[
  (T,M,R,\Delta,(X_\alpha)_{\alpha\in\Delta})
  \quad\text{and}\quad
  (T',M',R',\Delta',(X'_{\alpha'})_{\alpha'\in\Delta'})
\]
be the two pinnings. By conjugacy of maximal tori, we may assume
\(T=T'\). After restricting \(S\), we may assume that the isomorphism
\[
  D_S(M)\simeq D_S(M')
\]
comes from an isomorphism \(M\simeq M'\) carrying \(R\) onto \(R'\);
we are thus reduced to the situation
\[
  T=T',\qquad M=M',\qquad R=R'.
\]
Since the systems of simple roots are conjugate under the Weyl group
\textup{(\SGARefLink{sga3:XXI:3.3.7}{Exp.~XXI, 3.3.7})}, we may also assume
\(\Delta=\Delta'\). For every \(\alpha\in\Delta\), there is then a
scalar \(z_\alpha\in\mathbb G_m(S)\) such that
\[
  X'_\alpha=z_\alpha X_\alpha,
\]
and it is enough to construct, locally for the \textup{(fpqc)}
topology, a section \(t\) of \(T\) such that
\(\alpha(t)=z_\alpha\) for every \(\alpha\in\Delta\). But the
morphism
\[
  T\longrightarrow(\mathbb G_{m,S})^\Delta
\]
whose components are the \(\alpha\), \(\alpha\in\Delta\), is dual to
an injection \(\mathbb Z^\Delta\to M\), and is therefore faithfully
flat. This completes the proof of \SGARefLink{sga3:xxiv:1.5.2}{1.5.2}, and hence of \SGARefLink{sga3:xxiv:1.3}{1.3}.

\SGARefTarget{sga3:xxiv:1.6}
\paragraph{\SGARefLink{sga3:xxiv:1.6}{Corollary 1.6}.}
\emph{Let \(S\) be a scheme and \(G\) a reductive \(S\)-group. The
following conditions are equivalent:}
\begin{enumerate}
\item[\SGAInlineRefTarget{sga3:xxiv:1.6:item:i}(i)] \(\underline{\Aut}_{S\text{-}\mathrm{gr.}}(G)\) is affine
  \emph{\textup{(respectively of finite type, respectively of finite
  presentation, respectively quasi-compact)}} over \(S\);
\item[\SGAInlineRefTarget{sga3:xxiv:1.6:item:ii}(ii)] \(\underline{\Autext}(G)\) is finite over \(S\);
\item[\SGAInlineRefTarget{sga3:xxiv:1.6:item:iii}(iii)] for every \(s\in S\),
  \[
    \operatorname{rg}_{\mathrm{red}}(G_s)
    -\operatorname{rg}_{\mathrm{ss}}(G_s)\leq 1.
  \]
\end{enumerate}

Indeed, since \(\operatorname{ad}(G)\) is affine, flat, and of finite
presentation over \(S\), the morphism
\[
  p:\underline{\Aut}_{S\text{-}\mathrm{gr.}}(G)
  \longrightarrow\underline{\Autext}(G)
\]
is affine, faithfully flat, and of finite presentation.

If \(\underline{\Autext}(G)\) is finite over \(S\), it is affine over
\(S\), and hence so is
\(\underline{\Aut}_{S\text{-}\mathrm{gr.}}(G)\); this proves
\textup{\SGARefLink{sga3:xxiv:1.6:item:ii}{(ii)}}\(\Rightarrow\)\textup{\SGARefLink{sga3:xxiv:1.6:item:i}{(i)}}. If
\(\underline{\Aut}_{S\text{-}\mathrm{gr.}}(G)\) is quasi-compact over
\(S\), it is of finite presentation over \(S\), since it is in any
case locally of finite presentation and separated over \(S\). By
\SGARefLink{sga3:v:proposition:9:1}{Exp.~V, 9.1}, \(\underline{\Autext}(G)\) is then of finite presentation
over \(S\), and hence finite; this proves
\textup{\SGARefLink{sga3:xxiv:1.6:item:i}{(i)}}\(\Rightarrow\)\textup{\SGARefLink{sga3:xxiv:1.6:item:ii}{(ii)}}. Finally, to prove the
equivalence of \textup{\SGARefLink{sga3:xxiv:1.6:item:ii}{(ii)}} and \textup{\SGARefLink{sga3:xxiv:1.6:item:iii}{(iii)}}, one may suppose that
\(G\) is split, and the assertion reduces to \SGARefLink{sga3:XXI:6.7.8}{Exp.~XXI, 6.7.8}.

\SGARefTarget{sga3:xxiv:1.7}
\paragraph{\SGARefLink{sga3:xxiv:1.7}{Corollary 1.7}.}
\emph{Let \(S\) be a scheme and \(G\) a reductive \(S\)-group. Then}
\[
  \underline{\Aut}_{S\text{-}\mathrm{gr.}}(G)^0
  \simeq\operatorname{ad}(G).
\]

\SGARefTarget{sga3:xxiv:1.8}
\paragraph{\SGARefLink{sga3:xxiv:1.8}{Corollary 1.8}.}
\emph{Let \(S\) be a scheme, \(G\) a reductive \(S\)-group, and
\(H\) an \(S\)-group scheme that is smooth, affine, and has connected
fibres over \(S\). Then the \(S\)-functor}
\[
  \underline{\operatorname{Isom}}_{S\text{-}\mathrm{gr.}}(G,H)
\]
\emph{is representable by a smooth separated \(S\)-scheme, which is
affine over \(S\) if \(G\) is semisimple.}

Indeed, let \(U\) be the set of points \(s\in S\) for which \(H_s\)
is reductive; this is an open subset \textup{(\SGARefLink{sga3:xix:corollary:2:6}{Exp.~XIX, 2.6})}. If
\(S'\) is an \(S\)-scheme, \(H_{S'}\) is reductive if and only if
\(S'\to S\) factors through \(U\). It follows that the canonical
morphism
\[
  \underline{\operatorname{Isom}}_{S\text{-}\mathrm{gr.}}(G,H)
  \longrightarrow S
\]
factors through \(U\). We may therefore assume \(S=U\), and we are
reduced to the following corollary.

% END INCLUDED SGA3_English_current_buildable_source_20260823/source/inputs/XXIV/tex/components/03_expose_XXIV_lemma15_existence_corollaries16_18.tex

% BEGIN INCLUDED SGA3_English_current_buildable_source_20260823/source/inputs/XXIV/tex/components/04_expose_XXIV_corollaries19_110_remark111.tex
% SGA 3, Exposé XXIV, authority local page 6 / combined-reader page 1251.
% Personally checked against the Polo--Gille page image, band by band.

\SGARefTarget{sga3:xxiv:1.9}
\paragraph{\SGARefLink{sga3:xxiv:1.9}{Corollary 1.9}.}
\emph{Let \(S\) be a scheme, and let \(G\) and \(G'\) be two
reductive \(S\)-groups. Then}
\[
  F=\underline{\operatorname{Isom}}_{S\text{-}\mathrm{gr.}}(G,G')
\]
\emph{is representable by a smooth separated \(S\)-scheme, affine
if \(G\) or \(G'\) is semisimple. Moreover, \(S\) is the disjoint
union of two open subschemes \(S_1\) and \(S_2\) such that
\(F_{S_1}=\varnothing\), while \(F_{S_2}\) is a principal homogeneous
bundle on the left \textup{(respectively, on the right)} under
\(\underline{\Aut}_{S_2\text{-}\mathrm{gr.}}(G'_{S_2})\)
\textup{(respectively,
\(\underline{\Aut}_{S_2\text{-}\mathrm{gr.}}(G_{S_2})\)).}}

Indeed, let \(S_2\) be the set of points of \(S\) at which \(G\) and
\(G'\) have the same type, and let \(S_1\) be its complement. Since
the type of a reductive group is a locally constant function, \(S_1\)
and \(S_2\) are open. Clearly \(F_{S_1}=\varnothing\), and we may
assume \(S=S_2\). By \SGARefLink{sga3:XXIII:5.6}{Exp.~XXIII, 5.6}, \(F\) is a principal
homogeneous sheaf under
\(\underline{\Aut}_{S\text{-}\mathrm{gr.}}(G)\), locally trivial for
the étale topology. It follows that
\[
  F_0=F/\operatorname{ad}(G)
\]
is a principal homogeneous sheaf under \(\underline{\Autext}(G)\),
locally trivial for the étale topology, and hence representable
\textup{(\SGARefLink{sga3:x:corollary:5-5}{Exp.~X, 5.5})}. Since \(\operatorname{ad}(G)\) is affine,
\(F\) is therefore representable as well.

We note that the proof also gives the following result.

\SGARefTarget{sga3:xxiv:1.10}
\paragraph{\SGARefLink{sga3:xxiv:1.10}{Corollary 1.10}.}
\emph{Let \(S\) be a scheme, and let \(G\) and \(G'\) be two
reductive \(S\)-groups having the same type at every \(s\in S\).
Then \(\operatorname{ad}(G)\) acts freely on the right on
\(\underline{\operatorname{Isom}}_{S\text{-}\mathrm{gr.}}(G,G')\).
The quotient sheaf}
\[
  \underline{\operatorname{Isomext}}(G,G')
  =
  \underline{\operatorname{Isom}}_{S\text{-}\mathrm{gr.}}(G,G')
  /\operatorname{ad}(G)
\]
\emph{is representable by a twisted constant \(S\)-scheme; it is a
principal homogeneous bundle under \(\underline{\Autext}(G)\), and
is therefore finite over \(S\) if \(G\) is semisimple. Moreover, the
isomorphism}
\[
  \underline{\operatorname{Isom}}_{S\text{-}\mathrm{gr.}}(G,G')
  \simeq
  \underline{\operatorname{Isom}}_{S\text{-}\mathrm{gr.}}(G',G),
  \qquad u\longmapsto u^{-1},
\]
\emph{induces an isomorphism}
\[
  \underline{\operatorname{Isomext}}(G,G')
  \simeq
  \underline{\operatorname{Isomext}}(G',G).
\]

\SGARefTarget{sga3:xxiv:1.11}
\paragraph{\SGARefLink{sga3:xxiv:1.11}{Remark 1.11}.}
If
\(\underline{\operatorname{Isomext}}(G,G')(S)\neq\varnothing\), one
says that \(G\) is an \emph{inner twisted form} of \(G'\). Then \(G'\)
is an inner twisted form of \(G\), and the structure group of
\(\underline{\operatorname{Isom}}_{S\text{-}\mathrm{gr.}}(G,G')\)
can be reduced to \(\operatorname{ad}(G)\). More precisely, let
\[
  u\in\underline{\operatorname{Isomext}}(G,G')(S)
\]
and regard \(u\) as a section
\(u:S\to\underline{\operatorname{Isomext}}(G,G')\). Denote by
\[
  \underline{\operatorname{Isomint}}_u(G,G')
\]
the inverse image, under the canonical morphism
\[
  \underline{\operatorname{Isom}}_{S\text{-}\mathrm{gr.}}(G,G')
  \longrightarrow
  \underline{\operatorname{Isomext}}(G,G'),
\]
of the closed subscheme of
\(\underline{\operatorname{Isomext}}(G,G')\) that is the image of
\(u\). The natural action of \(\operatorname{ad}(G)\) on
\(\underline{\operatorname{Isomint}}_u(G,G')\) makes this scheme a
principal homogeneous bundle. By extension of the structure
group\footnote{The source's editorial note corrects
\(\operatorname{ad}(G)\to\underline{\Autext}(G)\) to
\(\operatorname{ad}(G)\to\underline{\Aut}(G)\).}
\[
  \operatorname{ad}(G)\longrightarrow\underline{\Aut}(G),
\]
one recovers
\(\underline{\operatorname{Isom}}_{S\text{-}\mathrm{gr.}}(G,G')\)
from \(\underline{\operatorname{Isomint}}_u(G,G')\).

By Hensel's lemma \textup{(\SGARefLink{sga3:xi:corollary:1:11}{Exp.~XI, 1.11})}, \SGARefLink{sga3:xxiv:1.8}{Corollary~1.8} immediately
gives:

% END INCLUDED SGA3_English_current_buildable_source_20260823/source/inputs/XXIV/tex/components/04_expose_XXIV_corollaries19_110_remark111.tex

% BEGIN INCLUDED SGA3_English_current_buildable_source_20260823/source/inputs/XXIV/tex/components/05_expose_XXIV_corollaries112_115.tex
% SGA 3, Exposé XXIV, authority local page 7 / combined-reader page 1252.
% Personally checked against the Polo--Gille page image, band by band.

\SGARefTarget{sga3:xxiv:1.12}
\paragraph{\SGARefLink{sga3:xxiv:1.12}{Corollary 1.12}.}
\emph{Let \(S\) be a henselian local scheme, let \(G\) be a
reductive \(S\)-group, let \(G'\) be a smooth affine \(S\)-group with
connected fibres, and let \(s\) be the closed point of \(S\). If
\(G_s\) and \(G'_s\) are isomorphic algebraic groups over
\(\kappa(s)\), then \(G\) and \(G'\) are isomorphic. More precisely,
every \(\kappa(s)\)-isomorphism}
\[
  G_s\simeq G'_s
\]
\emph{comes from an \(S\)-isomorphism \(G\simeq G'\).}

Applying \SGARefLink{sga3:xxiv:1.7}{Corollary~1.7} \textup{(respectively, \SGARefLink{sga3:xxiv:1.12}{Corollary~1.12})} to the
scheme of dual numbers over a field, \SGARefLink{sga3:iii:remark:2-10}{Exp.~III, 2.10}
\textup{(respectively, \SGARefLink{sga3:xxiv:3.10}{3.10})} yields part \textup{(i)}
\textup{(respectively, part \textup{(ii)})} of the next corollary.

\SGARefTarget{sga3:xxiv:1.13}
\paragraph{\SGARefLink{sga3:xxiv:1.13}{Corollary 1.13}.}
\emph{Let \(k\) be a field and \(G\) a reductive \(k\)-group.}
\begin{enumerate}
\item[\SGAInlineRefTarget{sga3:xxiv:1.13:item:i}(i)] \emph{If \(G\) is adjoint, then}
\[
  H^1(G,\Lie(G/k))=0.
\]
\footnote{The source's editorial note corrects the original, which
stated \textup{\SGARefLink{sga3:xxiv:1.13:item:i}{(i)}} without assuming \(G\) adjoint. By the
description of \(H^i(G,-)\) as the derived functors of
\(\operatorname{Hom}_G(k,-)\) \textup{(\SGARefLink{sga3:i:theorem:5-3-1}{Exp.~I, 5.3.1}; see also
[Ja87], I, 4.16)}, one has
\(H^i(G,V)=\operatorname{Ext}^i_G(k,V)\) for every \(G\)-module
\(V\). If \(\operatorname{char}(k)=p>0\) and
\(G=\operatorname{SL}_{p,k}\), then
\(\Lie(\operatorname{GL}_p/k)\) is a nontrivial extension of \(k\)
by \(\mathfrak g=\Lie(\operatorname{SL}_p/k)\), so
\(H^1(G,\mathfrak g)\neq0\). See also the addition \SGARefLink{sga3:xxiv:1.15.1}{1.15.1} below.}
\item[\SGAInlineRefTarget{sga3:xxiv:1.13:item:ii}(ii)] \emph{One has}
\[
  H^2(G,\Lie(G/k))=0.
\]
\end{enumerate}

\SGARefTarget{sga3:xxiv:1.14}
\paragraph{\SGARefLink{sga3:xxiv:1.14}{Remark 1.14}.}
\begin{enumerate}
\item[(i)] The assertion concerning \(H^1\) was known to Chevalley;
the assertion concerning \(H^2\) had been proved in most cases of
the classification by Chevalley.
\item[(ii)] In fact, the conjunction of \SGARefLink{sga3:xxiv:1.13}{Corollary~1.13} and the
uniqueness theorem over an algebraically closed field is essentially
equivalent to the uniqueness theorem. A direct proof of
\SGARefLink{sga3:xxiv:1.13}{Corollary~1.13} would therefore give a way to deduce the general
uniqueness theorem from Chevalley's uniqueness theorem over a field.
\end{enumerate}

The existence of reductive groups of every type over every scheme
\textup{(Exp.~XXV)} shows that the obstructions to lifting a
reductive \(k\)-group \(G\) over artinian rings with residue field
\(k\), which by \SGARefLink{sga3:iii:corollary:3-8}{Exp.~III, 3.8} are elements of
\[
  H^3(G,\Lie(G/k)\otimes V)
  \simeq H^3(G,\Lie(G/k))\otimes V,
\]
where \(V\) is a certain \(k\)-vector space endowed with the trivial
action of \(G\), vanish. This suggests that
\[
  H^3(G,\Lie(G/k))=0.
\]
A direct proof, if the assertion is true, would likewise doubtless
give a way to deduce the general existence theorem from the
existence theorem over a field \textup{(Chevalley's Tôhoku)}.

\SGARefTarget{sga3:xxiv:1.15}
\paragraph{\SGARefLink{sga3:xxiv:1.15}{Corollary 1.15}.}
\emph{Let \(k\) be a field,\footnote{The source's editorial note
replaces ``scheme'' with ``field.''} and let \(G\) be a reductive
\(k\)-group. Regard \(k\) as a trivial \(G\)-module. Then}
\[
  H^1(G,k)=H^2(G,k)=0.
\]
\footnote{Because of the correction made in \SGARefLink{sga3:xxiv:1.13}{1.13}, the source's
editors supply another proof for \(H^1\). Moreover, a theorem of
G.~Kempf gives \(H^i(G,k)=0\) for every \(i>0\); cf.\ [Ja87],
II, \SGARefLink{sga3:ii:proposition:4-5}{4.5} and 4.11.}
Since
\[
  H^i(G_{\bar k},\bar k)=H^i(G,k)\otimes_k\bar k,
\]
we may assume \(k\) algebraically closed. An element of \(H^1(G,k)\)
is simply a morphism of \(k\)-groups
\[
  \varphi:G\longrightarrow\mathbb G_{a,k}.
\]
Then
\[
  \varphi(G)=G/\ker(\varphi)
\]
is a smooth connected reductive subgroup of
\(\mathbb G_{a,k}\) \textup{(cf.\ XIX, \SGARefLink{sga3:xix:paragraph:1-7}{1.7})}, and is therefore
trivial. Thus \(H^1(G,k)=0\).

% END INCLUDED SGA3_English_current_buildable_source_20260823/source/inputs/XXIV/tex/components/05_expose_XXIV_corollaries112_115.tex

% BEGIN INCLUDED SGA3_English_current_buildable_source_20260823/source/inputs/XXIV/tex/components/06_expose_XXIV_corollary115_through_proposition28_en.tex
% Exposé XXIV, authority local pages 8--12 / combined-reader pages 1253--1257.
% Authority local p. 8 / combined-reader p. 1253.
Now consider the reductive \(k\)-group
\[
  H=G\times_k\mathbb G_{m,k}.
\]
There is an \(H\)-stable decomposition
\[
  \Lie(H/k)=\Lie(G/k)\oplus k.
\]
For every \(H\)-module \(V\),
\[
  H^i(H,V)=H^i\bigl(G,H^0(\mathbb G_{m,k},V)\bigr).
\]
Indeed, the \(H^i(H,-)\) are the derived functors of
\[
  H^0(H,-)
  =H^0\bigl(G,H^0(\mathbb G_{m,k},-)\bigr),
\]
and the functor \(H^0(\mathbb G_{m,k},-)\) is exact
(cf.\ \SGARefLink{sga3:i:theorem:5-3-1}{Exp.~I, 5.3.1} and \SGARefLink{sga3:i:theorem:5-3-3}{5.3.3}). In particular, for every \(i\),
\[
  H^i(H,\Lie(H/k))
  =
  H^i(G,\Lie(G/k))\oplus H^i(G,k).
\]
Applying \SGARefLink{sga3:xxiv:1.13:item:ii}{1.13(ii)} to the reductive group \(H\) gives
\(H^2(G,k)=0\).

\SGARefTarget{sga3:xxiv:1.15.1}
\medskip\noindent\textbf{Corollary 1.15.1.} ---
\footnote{Editor's note: This corollary, drawn from remarks of
O.~Gabber, has been added to make \SGARefLink{sga3:xxiv:1.13:item:i}{1.13(i)} more precise.}
{\itshape Let \(k\) be a field, \(G\) a reductive \(k\)-group, \(Z\)
its center, and \(\mathscr R=(M,M^*,R,R^*)\) the type of \(G\). Then
\[
  H^1(G,\Lie(G/k))
  \simeq
  \operatorname{Ext}^1_{\mathbb Z}(M/\Gamma_0(R),k).
\]
In particular, \(H^1(G,\Lie(G/k))=0\) if and only if \(Z\) is smooth
over \(k\).}

Let \(\mathfrak g\), \(\mathfrak z\), and
\(\mathfrak g_{\mathrm{ad}}\) be the Lie algebras of \(G\), \(Z\), and
\(G_{\mathrm{ad}}=G/Z\), respectively. It follows from \SGARefLink{sga3:xxiv:1.15}{1.15} and its
proof that
\[
  H^1(G,\mathfrak g)
  =
  H^1(G_{\mathrm{ad}},\mathfrak g)
  \simeq
  H^1(G_{\mathrm{ad}},\mathfrak g/\mathfrak z).
\]
Set \(C=\operatorname{Coker}(\mathfrak g\to
\mathfrak g_{\mathrm{ad}})\). There is an exact sequence
\[
  0\longrightarrow\mathfrak g/\mathfrak z
  \longrightarrow\mathfrak g_{\mathrm{ad}}
  \longrightarrow C\longrightarrow0.
\]
Since \(H^1(G_{\mathrm{ad}},\mathfrak g_{\mathrm{ad}})=0\) by \SGARefLink{sga3:xxiv:1.13}{1.13}
and
\[
  H^0(G_{\mathrm{ad}},\mathfrak g_{\mathrm{ad}})=0
\]
(cf.\ \SGARefLink{sga3:ii:theorem:5-2-3}{Exp.~II, 5.2.3}), we obtain
\[
  H^1(G_{\mathrm{ad}},\mathfrak g/\mathfrak z)
  \simeq H^0(G_{\mathrm{ad}},C).
\]
To calculate the right-hand side, we may assume \(k\) algebraically
closed. Let \((G,T,M,R)\) be a splitting of \(G\). Then
\[
  Z=D_k(M/\Gamma_0(R))
\]
and \(C\) is identified with
\[
  \operatorname{Coker}\!\left(
    \operatorname{Hom}_{\mathbb Z}(M,k)
    \longrightarrow
    \operatorname{Hom}_{\mathbb Z}(\Gamma_0(R),k)
  \right)
  \simeq
  \operatorname{Ext}^1_{\mathbb Z}(M/\Gamma_0(R),k),
\]
with trivial \(G_{\mathrm{ad}}\)-action. The corollary follows, since
\(Z\) is not smooth over \(k\) exactly when
\(\operatorname{char}(k)=p>0\) and \(M/\Gamma_0(R)\) has
\(p\)-torsion (\SGARefLink{sga3:ix:proposition:2:1}{Exp.~IX, 2.1}).

\SGARefTarget{sga3:xxiv:1.16}
\medskip\noindent\textbf{Definition 1.16.} ---
{\itshape Let \(S\) be a scheme and \(G\) a reductive \(S\)-group. A
\emph{form of \(G\) over \(S\)} is an \(S\)-group scheme \(G'\)
locally isomorphic to \(G\) for the \textup{(fpqc)} topology. It is
equivalent to require that \(G'\) be locally isomorphic to \(G\) for
the étale topology (cf.\ \SGARefLink{sga3:XXIII:5.6}{Exp.~XXIII, 5.6}), or that \(G'\) be a
reductive \(S\)-group of the same type as \(G\) at every point of
\(S\).}

\SGARefTarget{sga3:xxiv:1.17}
\medskip\noindent\textbf{Corollary 1.17.} ---
{\itshape Let \(S\) be a scheme and \(G\) a reductive \(S\)-group.
\begin{enumerate}
\item[(i)] The functor
\[
  G'\longmapsto
  \underline{\operatorname{Isom}}_{S\text{-}\mathrm{gr}}(G,G')
\]
is an equivalence between the category of forms of \(G\) over \(S\)
and the category of principal homogeneous bundles under
\(\underline{\Aut}_{S\text{-}\mathrm{gr}}(G)\).
\item[(ii)] If \(S'\to S\) is a covering morphism, forms of \(G\)
trivialized by \(S'\) correspond to bundles trivialized by \(S'\).
\item[\SGAInlineRefTarget{sga3:xxiv:1.17:item:iii}(iii)] Every principal homogeneous sheaf under
\(\underline{\Aut}_{S\text{-}\mathrm{gr}}(G)\) is representable and
quasi-isotrivial, that is, locally trivial for the étale topology.
\end{enumerate}}

% Authority local p. 9 / combined-reader p. 1254.
The first assertion is formal in the category of sheaves, for example
for the \textup{(fpqc)} topology. Every sheaf locally isomorphic to
\(G\) for the \textup{(fpqc)} topology is representable, since \(G\)
is affine over \(S\), and is locally isomorphic to \(G\) for the
étale topology. Finally, for every form \(G'\) of \(G\), the
\(S\)-sheaf
\(\underline{\operatorname{Isom}}_{S\text{-}\mathrm{gr}}(G,G')\) is
representable by \SGARefLink{sga3:xxiv:1.8}{1.8}. The corollary follows.

\SGARefTarget{sga3:xxiv:1.18}
\medskip\noindent\textbf{Corollary 1.18.} ---
{\itshape The set of isomorphism classes of forms of the reductive
group \(G\) over \(S\) is isomorphic to
\[
  H^1\!\left(S,\underline{\Aut}_{S\text{-}\mathrm{gr}}(G)\right)
  =
  H^1_{\mathrm{\acute et}}\!
    \left(S,\underline{\Aut}_{S\text{-}\mathrm{gr}}(G)\right)
  =
  \Fib\!\left(S,\underline{\Aut}_{S\text{-}\mathrm{gr}}(G)\right).
\]
If \(S'\to S\) is a covering morphism, the subset consisting of forms
trivialized by \(S'\) is isomorphic to
\[
  H^1\!\left(S'/S,\underline{\Aut}_{S\text{-}\mathrm{gr}}(G)\right).
\]}

\SGARefTarget{sga3:xxiv:1.19}
\medskip\noindent\textbf{Corollary 1.19.} ---
{\itshape Let \(S\) be a scheme and \(\mathscr R\) a pinned reduced root
datum such that \(\operatorname{\acute{E}p}_S(\mathscr R)\)
exists,\footnote{Editor's note: Cf.\ XXIII, \SGARefLink{sga3:XXIII:5.11}{Definition 5.11}.} a
condition which is automatically satisfied by Exp.~XXV. Put
\[
  A_S(\mathscr R)
  =
  \underline{\Aut}_{S\text{-}\mathrm{gr}}
    (\operatorname{\acute{E}p}_S(\mathscr R))
  =
  \operatorname{ad}(\operatorname{\acute{E}p}_S(\mathscr R))
  \cdot E(\mathscr R)_S.
\]
\begin{enumerate}
\item[\SGAInlineRefTarget{sga3:xxiv:1.19:item:i}(i)] The set of isomorphism classes of reductive \(S\)-groups
of type \(\mathscr R\) (\SGARefLink{sga3:xxii:def:2.7}{Exp.~XXII, 2.7}) is isomorphic, by
\SGARefLink{sga3:XXIII:5.12}{Exp.~XXIII, 5.12}, to
\[
  H^1(S,A_S(\mathscr R))
  =
  H^1_{\mathrm{\acute et}}(S,A_S(\mathscr R))
  =
  \Fib(S,A_S(\mathscr R)).
\]
\item[\SGAInlineRefTarget{sga3:xxiv:1.19:item:ii}(ii)] If \(S'\to S\) is a covering morphism, the subset
consisting of classes of groups splittable over \(S'\) is isomorphic
to \(H^1(S'/S,A_S(\mathscr R))\).
\end{enumerate}}

\SGARefTarget{sga3:xxiv:1.20}
\medskip\noindent\textbf{Remark 1.20.} ---
With the preceding notation, every reductive \(S\)-group \(G\) of
type \(\mathscr R\) has a canonically associated right principal
homogeneous bundle under \(A_S(\mathscr R)\):
\[
  \underline{\operatorname{Isom}}_{S\text{-}\mathrm{gr}}
    (\operatorname{\acute{E}p}_S(\mathscr R),G)=P.
\]
The scheme \(P\) may be interpreted as the ``scheme of pinnings of
\(G\) of type \(\mathscr R\)'' (cf.\ \SGARefLink{sga3:xxiv:1.0}{1.0}). It is also a principal
homogeneous bundle on the left under
\(\underline{\Aut}_{S\text{-}\mathrm{gr}}(G)\), a structure visible
immediately from the description above.

\SGARefTarget{sga3:xxiv:1.21}
\medskip\noindent\textbf{Proposition 1.21.} ---
{\itshape Let \(S\) be a henselian local scheme and \(s\) its closed
point. The functor
\[
  G\longmapsto G_s
\]
induces a bijection from the set of isomorphism classes of reductive
\(S\)-groups onto the set of isomorphism classes of reductive
\(\kappa(s)\)-groups. In particular, for every reductive \(S\)-group
\(G\), there exists a finite étale surjective morphism \(S'\to S\)
such that \(G_{S'}\) is splittable.}

% Authority local p. 10 / combined-reader p. 1255.
Using the existence of \(\operatorname{\acute{E}p}_S(\mathscr R)\)
(Exp.~XXV), it is enough to prove that, with
\(H=A_S(\mathscr R)\), the canonical map
\[
  \Fib(S,H)\longrightarrow\Fib(\kappa(s),H_s)
\]
is bijective, and that every element of \(\Fib(S,H)\) has the property
just stated. Every finite subset of \(H\) is contained in an affine
open subscheme: this is clear for a constant group, and \(H\) is
affine over a constant group. The result proved in the Appendix,
\SGARefLink{sga3:xxiv:proposition:8-1}{8.1}, therefore applies.

\SGARefTarget{sga3:xxiv:section:2}
\section*{2. Automorphisms and Subgroups}

Introduce the following notation. If
\[
  \mathcal H=\underline{\Aut}_{S\text{-}\mathrm{gr}}(G)
\]
and \(X\) is a subfunctor of \(G\), set
\[
  \underline{\Aut}_{S\text{-}\mathrm{gr}}(G,X)
  =\operatorname{Norm}_{\mathcal H}(X),
  \qquad
  \underline{\Aut}_{S\text{-}\mathrm{gr}}(G,\operatorname{id}_X)
  =\underline{\Centr}_{\mathcal H}(X).
\]
If \(Y\) is a second subfunctor of \(G\), define
\[
  \underline{\Aut}_{S\text{-}\mathrm{gr}}(G,X,Y)
  =
  \underline{\Aut}_{S\text{-}\mathrm{gr}}(G,X)
  \cap
  \underline{\Aut}_{S\text{-}\mathrm{gr}}(G,Y).
\]
If \(G'\) is a second \(S\)-group and \(X'\) a subfunctor of \(G'\),\footnote{Translator's note: The French text prints \(G\) here; \(G'\) is required by the typing of the definition.}
write
\[
  \underline{\operatorname{Isom}}_{S\text{-}\mathrm{gr}}
    (G,X;G',X')
\]
for the subfunctor of
\(\underline{\operatorname{Isom}}_{S\text{-}\mathrm{gr}}(G,G')\)
defined, for every \(S'\to S\), by
\[
  \left\{
    u\in
    \underline{\operatorname{Isom}}_{S\text{-}\mathrm{gr}}(G,G')(S')
    \,\middle|\,
    u(X_{S'})=X'_{S'}
  \right\}.
\]
The notation
\(\underline{\operatorname{Isom}}_{S\text{-}\mathrm{gr}}
(G,X,Y;G',X',Y')\), and so on, is defined similarly.\footnote{Editor's
note: The preceding definitions have been made explicit. The original
only said that
\(\underline{\Aut}_{S\text{-}\mathrm{gr}}(G,X,Y),\ldots\) and
\(\underline{\operatorname{Isom}}_{S\text{-}\mathrm{gr}}
(G,X;G',X'),\ldots\) were defined similarly.}

\SGARefTarget{sga3:xxiv:2.1}
\medskip\noindent\textbf{Proposition 2.1.} ---
{\itshape Let \(S\) be a scheme, \(G\) a reductive \(S\)-group, and
\(T\) a maximal torus of \(G\) (respectively, \(B\) a Borel subgroup
of \(G\), or \(B\supset T\) a Killing couple of \(G\)). Let
\(T_{\mathrm{ad}}\) (respectively, \(B_{\mathrm{ad}}\)) denote the
corresponding maximal torus (respectively, Borel subgroup) of
\(\operatorname{ad}(G)\):
\[
  B_{\mathrm{ad}}
  \simeq B/\underline{\Centr}(G)
  =B/\underline{\Centr}(B),
  \qquad
  T_{\mathrm{ad}}\simeq T/\underline{\Centr}(G).
\]
Then
\(\underline{\Aut}_{S\text{-}\mathrm{gr}}(G,T)\)
(respectively,
\(\underline{\Aut}_{S\text{-}\mathrm{gr}}(G,B)\), or
\(\underline{\Aut}_{S\text{-}\mathrm{gr}}(G,B,T)\))
is represented by a closed subscheme of
\(\underline{\Aut}_{S\text{-}\mathrm{gr}}(G)\), smooth over \(S\),
and the exact sequence of \SGARefLink{sga3:xxiv:1.3}{Theorem 1.3} induces the exact sequences
\[
  1\longrightarrow
  \operatorname{Norm}_{\operatorname{ad}(G)}(T_{\mathrm{ad}})
  \longrightarrow
  \underline{\Aut}_{S\text{-}\mathrm{gr}}(G,T)
  \longrightarrow
  \underline{\Autext}(G)
  \longrightarrow1,
\]
\[
  1\longrightarrow B_{\mathrm{ad}}
  \longrightarrow
  \underline{\Aut}_{S\text{-}\mathrm{gr}}(G,B)
  \longrightarrow
  \underline{\Autext}(G)
  \longrightarrow1,
\]
\[
  1\longrightarrow T_{\mathrm{ad}}
  \longrightarrow
  \underline{\Aut}_{S\text{-}\mathrm{gr}}(G,B,T)
  \longrightarrow
  \underline{\Autext}(G)
  \longrightarrow1.
\]}

By descent of closed subschemes,\footnote{Editor's note: Cf.\ SGA~1,
VIII, 1.9.} we reduce immediately to the case in which \(G\) is
pinned and \(B\supset T\) is its canonical Killing couple
(cf.\ \SGARefLink{sga3:xxii:cor:5.5.5:item:iv}{Exp.~XXII, 5.5.5(iv)}). Since the group \(E\) of \SGARefLink{sga3:xxiv:1.2}{1.2} normalizes
\(B\) and \(T\), the result follows at once from the normalization
theorems in \(\operatorname{ad}(G)\) (\SGARefLink{sga3:xxii:cor:5.3.12}{Exp.~XXII, 5.3.12} and 5.6.1).

% Authority local p. 11 / combined-reader p. 1256.

Using now the conjugacy theorems (cf.\ \SGARefLink{sga3:XXIII:5.12}{Exp.~XXIII, 5.12}), and arguing
as in \S1, one obtains:

\SGARefTarget{sga3:xxiv:2.2}
\medskip\noindent\textbf{Corollary 2.2.} ---
{\itshape Let \(S\) be a scheme, let \(G\) and \(G'\) be two reductive
\(S\)-groups having the same type at every point, and let
\(B\supset T\) (respectively, \(B'\supset T'\)) be a Killing couple
of \(G\) (respectively, \(G'\)).

\textup{\SGARefLink{sga3:xxiv:1.19:item:i}{(i)}} The \(S\)-functor
\[
  \underline{\operatorname{Isom}}_{S\text{-}\mathrm{gr}}
    (G,T;G',T')
\]
is represented by a closed smooth subscheme of
\(\underline{\operatorname{Isom}}_{S\text{-}\mathrm{gr}}(G,G')\)
which is a principal homogeneous space under
\(\underline{\Aut}_{S\text{-}\mathrm{gr}}(G,T)\).
Moreover,
\(\operatorname{Norm}_{\operatorname{ad}(G)}(T_{\mathrm{ad}})\)
acts freely on this scheme, and there is a canonical isomorphism
\[
  \underline{\operatorname{Isom}}_{S\text{-}\mathrm{gr}}
    (G,T;G',T')
  \big/
  \operatorname{Norm}_{\operatorname{ad}(G)}(T_{\mathrm{ad}})
  \simeq
  \underline{\operatorname{Isomext}}(G,G').
\]

\textup{\SGARefLink{sga3:xxiv:1.19:item:ii}{(ii)}} The \(S\)-functor
\[
  \underline{\operatorname{Isom}}_{S\text{-}\mathrm{gr}}
    (G,B;G',B')
\]
is represented by a closed smooth subscheme of
\(\underline{\operatorname{Isom}}_{S\text{-}\mathrm{gr}}(G,G')\)
which is a principal homogeneous space under
\(\underline{\Aut}_{S\text{-}\mathrm{gr}}(G,B)\).
Moreover, \(B_{\mathrm{ad}}\) acts freely on this scheme, and there
is a canonical isomorphism
\[
  \underline{\operatorname{Isom}}_{S\text{-}\mathrm{gr}}
    (G,B;G',B')
  \big/B_{\mathrm{ad}}
  \simeq
  \underline{\operatorname{Isomext}}(G,G').
\]

\textup{\SGARefLink{sga3:xxiv:1.17:item:iii}{(iii)}} The \(S\)-functor
\[
  \underline{\operatorname{Isom}}_{S\text{-}\mathrm{gr}}
    (G,B,T;G',B',T')
\]
is represented by a closed smooth subscheme of
\(\underline{\operatorname{Isom}}_{S\text{-}\mathrm{gr}}(G,G')\),
which is a principal homogeneous space under
\(\underline{\Aut}_{S\text{-}\mathrm{gr}}(G,B,T)\).
Moreover, \(T_{\mathrm{ad}}\) acts freely on this scheme, and there
is a canonical isomorphism
\[
  \underline{\operatorname{Isom}}_{S\text{-}\mathrm{gr}}
    (G,B,T;G',B',T')
  \big/T_{\mathrm{ad}}
  \simeq
  \underline{\operatorname{Isomext}}(G,G').
\]}

Arguing once again as in \S1, one obtains:

\SGARefTarget{sga3:xxiv:2.3}
\medskip\noindent\textbf{Corollary 2.3.} ---
{\itshape Let \(S\) be a scheme, let \(G\) be a reductive \(S\)-group,
and let \(B\supset T\) be a Killing couple of \(G\). The functor}
\[
  (G',T')
  \longmapsto
  \underline{\operatorname{Isom}}_{S\text{-}\mathrm{gr}}
    (G,T;G',T'),
\]
{\itshape respectively}
\[
  (G',B')
  \longmapsto
  \underline{\operatorname{Isom}}_{S\text{-}\mathrm{gr}}
    (G,B;G',B'),
\]
{\itshape respectively}
\[
  (G',B',T')
  \longmapsto
  \underline{\operatorname{Isom}}_{S\text{-}\mathrm{gr}}
    (G,B,T;G',B',T'),
\]
{\itshape is an equivalence between the category of pairs
\((G',T')\) (respectively, pairs \((G',B')\), respectively, triples
\((G',B',T')\)), where \(G'\) is a form of \(G\) and \(T'\) is a
maximal torus of \(G'\) (respectively, \(B'\) is a Borel subgroup of
\(G'\), respectively, \(B'\supset T'\) is a Killing couple of
\(G'\)), and the category of principal homogeneous bundles under
the \(S\)-group \(H\), where
\(H=\underline{\Aut}_{S\text{-}\mathrm{gr}}(G,T)\)
(respectively,
\(H=\underline{\Aut}_{S\text{-}\mathrm{gr}}(G,B)\),
respectively,
\(H=\underline{\Aut}_{S\text{-}\mathrm{gr}}(G,B,T)\)).
Moreover, every principal homogeneous sheaf under \(H\) is
representable and quasi-isotrivial, so that}
\[
  H^1(S,H)=H^1_{\mathrm{\acute et}}(S,H)=\Fib(S,H).
\]

% Authority local p. 12 / combined-reader p. 1257.

\SGARefTarget{sga3:xxiv:2.4}
\medskip\noindent\textbf{Remark 2.4.} ---
Under the hypotheses of \SGARefLink{sga3:xxiv:2.2}{2.2}, the morphism denoted set-theoretically
by \(u\mapsto u(T)\) (respectively, \(u\mapsto u(B)\), respectively,
\(u\mapsto(u(B),u(T))\)) induces an isomorphism
\[
  \underline{\operatorname{Isom}}_{S\text{-}\mathrm{gr}}(G,G')
  \big/
  \underline{\Aut}_{S\text{-}\mathrm{gr}}(G,T)
  \simeq \operatorname{Tor}(G'),
\]
respectively,
\[
  \underline{\operatorname{Isom}}_{S\text{-}\mathrm{gr}}(G,G')
  \big/
  \underline{\Aut}_{S\text{-}\mathrm{gr}}(G,B)
  \simeq \operatorname{Bor}(G'),
\]
respectively,
\[
  \underline{\operatorname{Isom}}_{S\text{-}\mathrm{gr}}(G,G')
  \big/
  \underline{\Aut}_{S\text{-}\mathrm{gr}}(G,B,T)
  \simeq \operatorname{Kil}(G').
\]
The proof is immediate: it suffices to prove the assertion locally
for \textup{(fpqc)}, so one may assume \(G\simeq G'\), and the result
then reduces to \SGARefLink{sga3:xxii:cor:5.8.3:item:iii}{Exp.~XXII, 5.8.3(iii)}.

\SGARefTarget{sga3:xxiv:2.5}
\medskip\noindent\textbf{Remark 2.5.} ---
The preceding results are immediately interpreted in terms of
reduction of the structure group. If \(G'\) is a form of \(G\),
corresponding to the principal bundle
\(\underline{\operatorname{Isom}}_{S\text{-}\mathrm{gr}}(G,G')\),
to give a reduction of the structure group of this bundle to
\(\underline{\Aut}_{S\text{-}\mathrm{gr}}(G,T)\) amounts to giving
a maximal torus \(T'\) of \(G'\). The indicated bijections are,
on the one hand, that of \SGARefLink{sga3:xxiv:2.4}{2.4} and, on the other hand, the map
\[
  T'\longmapsto
  \underline{\operatorname{Isom}}_{S\text{-}\mathrm{gr}}
    (G,T;G',T').
\]
The same applies to Borel subgroups and Killing couples.

\SGARefTarget{sga3:xxiv:2.6}
\medskip\noindent\textbf{Proposition 2.6.} ---
{\itshape Let \(S\) be a scheme, let \(G\) and \(G'\) be two reductive
\(S\)-groups having the same type at every point, and let \(T\)
(respectively, \(T'\)) be a maximal torus of \(G\) (respectively,
\(G'\)). Then \(T_{\mathrm{ad}}\) acts freely on
\(\underline{\operatorname{Isom}}_{S\text{-}\mathrm{gr}}
(G,T;G',T')\), and the quotient}
\[
  P=
  \underline{\operatorname{Isom}}_{S\text{-}\mathrm{gr}}
    (G,T;G',T')
  \big/T_{\mathrm{ad}}
\]
{\itshape is representable; it is a principal homogeneous bundle under}
\[
  A=
  \underline{\Aut}_{S\text{-}\mathrm{gr}}(G,T)
  \big/T_{\mathrm{ad}},
\]
{\itshape where \(A\) is represented by a twisted constant \(S\)-scheme
and is an extension of \(\underline{\Autext}(G)\) by}
\[
  W_{\operatorname{ad}(G)}(T_{\mathrm{ad}})
  =
  \operatorname{Norm}_{\operatorname{ad}(G)}(T_{\mathrm{ad}})
  \big/T_{\mathrm{ad}}.
\]
{\itshape Moreover, if \(A\) is made to act on \(T\) in the evident
way, the bundle associated with \(P\) is precisely \(T'\).}\footnote{Editor's
note: Consequently, if \(P\) has a section, then \(T\) and \(T'\)
are isomorphic; in particular, \(T_P\) and \(T'_P\) are isomorphic.
See the addition \SGARefLink{sga3:xxiv:proposition:4-5-1}{4.5.1}, where this observation is developed in the
case \(G=\operatorname{\acute{E}p}_S(\mathscr R)\).}

The first part of the proposition follows immediately from the
preceding results. To prove the second, observe that there is an
evident morphism
\[
  P\times_S T\longrightarrow T',
  \qquad (u,t)\longmapsto u(t).
\]
To show that, after passage to the quotient by \(A\), it induces an
isomorphism, one may once again assume
\((G,T)\simeq(G',T')\), in which case the assertion is immediate.

In an entirely analogous way, one has:

\SGARefTarget{sga3:xxiv:2.7}
\medskip\noindent\textbf{Proposition 2.7.} ---
{\itshape Let \(S\) be a scheme, let \(G\) and \(G'\) be two reductive
\(S\)-groups having the same type at every point, and let
\(B\supset T\) (respectively, \(B'\supset T'\)) be a Killing couple
of \(G\) (respectively, \(G'\)). If
\[
  \underline{\Aut}_{S\text{-}\mathrm{gr}}(G,B,T)
  \big/T_{\mathrm{ad}}
  \simeq \underline{\Autext}(G)
\]
is made to act on \(T\) in the evident way, the bundle associated
with \(\underline{\operatorname{Isomext}}(G,G')\) is precisely
\(T'\).}

\SGARefTarget{sga3:xxiv:2.8}
\medskip\noindent\textbf{Corollary 2.8.} ---
{\itshape Let \(G\) and \(G'\) be two reductive \(S\)-groups which are
inner twisted forms of one another, and let \(B\supset T\)
(respectively, \(B'\supset T'\)) be a Killing couple of \(G\)
(respectively, \(G'\)). Then \(T\) and \(T'\) are
isomorphic.}\footnote{Editor's note: This follows from \SGARefLink{sga3:xxiv:2.7}{2.7}, since
\(\underline{\operatorname{Isomext}}(G,G')\) has a section
(cf.\ the preceding editor's note).}


% END INCLUDED SGA3_English_current_buildable_source_20260823/source/inputs/XXIV/tex/components/06_expose_XXIV_corollary115_through_proposition28_en.tex

% BEGIN INCLUDED SGA3_English_current_buildable_source_20260823/source/inputs/XXIV/tex/components/07_expose_XXIV_section2_conclusion_and_section3_through_39_en.tex
% Exposé XXIV, authority local pages 13--16 / combined-reader pages 1258--1261.
% Authority local p. 13 / combined-reader p. 1258.

\SGARefTarget{sga3:xxiv:2.9}
\medskip\noindent\textbf{Remark 2.9.} ---
It is not true in general that \(B\) and \(B'\) are isomorphic;
they are nevertheless inner twisted forms of one another
(cf.\ \S5).

One may develop ``\(\operatorname{Isomint}\)'' variants\footnote{Editor's
note: Recall that
\(\underline{\operatorname{Isomint}}_u(G,G')\) is defined in \SGARefLink{sga3:xxiv:1.11}{1.11}.}
of the preceding results. We mention one:

\SGARefTarget{sga3:xxiv:2.10}
\medskip\noindent\textbf{Proposition 2.10.} ---
{\itshape Let \(S\) be a scheme, let \(G\) and \(G'\) be two reductive
\(S\)-groups having the same type at every point, let \(B\supset T\)
(respectively, \(B'\supset T'\)) be a Killing couple of \(G\)
(respectively, \(G'\)), and let
\(u\in\underline{\operatorname{Isomext}}(G,G')(S)\). Consider}
\[
\begin{aligned}
  &\underline{\operatorname{Isomint}}_u
     (G,B,T;G',B',T')\\
  &\qquad =
  \underline{\operatorname{Isom}}_{S\text{-}\mathrm{gr}}
     (G,B,T;G',B',T')
  \cap
  \underline{\operatorname{Isomint}}_u(G,G').
\end{aligned}
\]
{\itshape It is a smooth affine \(S\)-scheme which is a principal
homogeneous bundle under \(T_{\mathrm{ad}}\). In particular,}
\[
  \underline{\operatorname{Isomint}}_u
    (G,B,T;G',B',T')(S)\ne\varnothing
\]
{\itshape if and only if the corresponding element of
\(H^1(S,T_{\mathrm{ad}})\) is zero.}

To conclude this section, we prove two results that will be useful
later.

\SGARefTarget{sga3:xxiv:2.11}
\medskip\noindent\textbf{Proposition 2.11.} ---
{\itshape Let \(S\) be a scheme, let \(G\) be a reductive \(S\)-group,
and let \(T\) be a maximal torus of \(G\). The evident morphism}
\[
  T_{\mathrm{ad}}
  \xrightarrow{\ \sim\ }
  \underline{\Aut}_{S\text{-}\mathrm{gr}}(G,\operatorname{id}_T)
\]
{\itshape is an isomorphism.}

This follows from the preceding statements. A direct proof was also
given in the course of the proof of \SGARefLink{sga3:xxiv:1.5.2}{1.5.2}.

\SGARefTarget{sga3:xxiv:2.12}
\medskip\noindent\textbf{Corollary 2.12.} ---
{\itshape Under the preceding hypotheses, there is an equivalence
between the category of pairs \((G',f)\), where \(G'\) is a form of
\(G\) and \(f\) is an isomorphism from \(T\) onto a maximal torus of
\(G'\), and the category of principal homogeneous bundles under
\(T_{\mathrm{ad}}\).}

\SGARefTarget{sga3:xxiv:2.13}
\medskip\noindent\textbf{Corollary 2.13.} ---
{\itshape If \(H^1(S,T_{\mathrm{ad}})=0\), and if \(G'\) is a form of
\(G\) having a maximal torus isomorphic to \(T\), then \(G'\) is
isomorphic to \(G\).}

\SGARefTarget{sga3:xxiv:2.14}
\medskip\noindent\textbf{Corollary 2.14.} ---
{\itshape Let \(S\) be a scheme such that \(\operatorname{Pic}(S)=0\),
and let \(G\) be a reductive \(S\)-group of constant type. Then \(G\)
is splittable if and only if \(G\) has a split maximal torus.}

Let \(G\) be a reductive group and let
\(\operatorname{rad}(G)\) be its radical (\SGARefLink{sga3:xxii:rem:4.3.9}{Exp.~XXII, 4.3.9}).
Since \(\operatorname{rad}(G)\) is central and characteristic in
\(G\), there is a canonical morphism
\[
  q:\underline{\Autext}(G)
  \longrightarrow
  \underline{\Aut}_{S\text{-}\mathrm{gr}}
    \bigl(\operatorname{rad}(G)\bigr).
\]

% Authority local p. 14 / combined-reader p. 1259.

\SGARefTarget{sga3:xxiv:2.15}
\medskip\noindent\textbf{Proposition 2.15.} ---
{\itshape Let \(G\) be a reductive \(S\)-group. The following sequence
is exact:}
\[
  1\longrightarrow
  \operatorname{ad}(G)
  \longrightarrow
  \underline{\Aut}_{S\text{-}\mathrm{gr}}
    \bigl(G,\operatorname{id}_{\operatorname{rad}(G)}\bigr)
  \xrightarrow{p}
  \underline{\Autext}(G)
  \xrightarrow{q}
  \underline{\Aut}_{S\text{-}\mathrm{gr}}
    \bigl(\operatorname{rad}(G)\bigr),
\]
{\itshape and \(\ker(q)=\operatorname{im}(p)\) is an open and closed
subscheme of \(\underline{\Autext}(G)\), finite over \(S\).}

One may assume that \(G\) is split. The first assertion is immediate,
and the second follows from \SGARefLink{sga3:XXI:6.7.5}{Exp.~XXI, 6.7.5} and 6.7.7.

Writing
\[
  H=
  \underline{\Aut}_{S\text{-}\mathrm{gr}}
    \bigl(G,\operatorname{id}_{\operatorname{rad}(G)}\bigr),
\]
one obtains:

\SGARefTarget{sga3:xxiv:2.16}
\medskip\noindent\textbf{Corollary 2.16.} ---
{\itshape There is an equivalence between the category of pairs
\((G',f)\), where \(G'\) is a form of \(G\) and \(f\) is an
isomorphism from \(\operatorname{rad}(G)\) onto the radical of \(G'\),
and the category of principal bundles under a certain \(S\)-group
scheme \(H\), where \(H\) is such that there is an exact sequence}
\[
  1\longrightarrow \operatorname{ad}(G)
  \longrightarrow H
  \longrightarrow F
  \longrightarrow 1,
\]
{\itshape in which the \(S\)-group \(F\) is étale and finite over
\(S\).}

\SGARefTarget{sga3:xxiv:section:3}
\section*{3. Dynkin Scheme of a Reductive Group. Quasi-Split Groups}

\SGARefTarget{sga3:xxiv:3.1}
\medskip\noindent\textbf{3.1.} ---
Recall (\SGARefLink{sga3:XXI:7.4.1}{Exp.~XXI, 7.4.1}) that a Dynkin diagram is a finite set endowed
with the structure defined by a set of pairs of distinct elements
(bonds) and a map to \(\{1,2,3\}\) (lengths). To every pinned reduced
root datum \(\mathscr R\) there is associated a Dynkin diagram
\(\Delta(\mathscr R)\), whose underlying set is the set of simple
roots.

\SGARefTarget{sga3:xxiv:3.2}
\medskip\noindent\textbf{3.2.} ---
Let \(S\) be a scheme. A Dynkin \(S\)-scheme is a finite twisted
constant \(S\)-scheme \(\Delta\), endowed with the structure defined
by a subscheme \(L\) of \(\Delta\times_S\Delta\) having empty
intersection with the diagonal, and by a morphism
\[
  \Delta\longrightarrow\{1,2,3\}_S.
\]
If \(S'\) is a connected \(S\)-scheme,\footnote{Editor's note: The
hypothesis that \(S'\) be connected has been added.}
\(\Delta(S')\) is naturally endowed with the structure of a Dynkin
diagram.

The following notions are defined immediately: an isomorphism of two
Dynkin schemes, the base extension of a Dynkin scheme, and the
constant Dynkin scheme associated with a Dynkin diagram.

Every étale descent datum on a Dynkin scheme is effective.

\SGARefTarget{sga3:xxiv:3.3}
\medskip\noindent\textbf{3.3.} ---
We shall associate a Dynkin \(S\)-scheme with every reductive
\(S\)-group \(G\). First suppose that \(G\) is splittable over \(S\).
For every pinning \(E\) of \(G\), let \(\Delta(E)\) be the constant
Dynkin scheme associated with the pinned root datum defined by \(E\).
If \(E\) and \(E'\) are two pinnings of \(G\), then by \SGARefLink{sga3:xxiv:1.5}{1.5} there is
a unique inner automorphism of \(G\) over \(S\) carrying \(E\) to
\(E'\). This automorphism of \(G\) defines an isomorphism
\[
  a_{EE'}:\Delta(E)\xrightarrow{\ \sim\ }\Delta(E').
\]
The \(a_{EE'}\) plainly form a transitive system, so the
\(\Delta(E)\) may be identified (that is, one takes their inductive
limit); the result is a constant Dynkin scheme denoted by
\(\operatorname{Dyn}(G)\). If now \(G\) is an arbitrary reductive
\(S\)-group, there is an étale covering family \(\{S_i\to S\}\)
such that

% Authority local p. 15 / combined-reader p. 1260.

\(G_{S_i}\) is splittable. Arguing as above, one therefore obtains a
canonical descent datum on the \(\operatorname{Dyn}(G_{S_i})\), and
by descent this constructs a Dynkin \(S\)-scheme
\(\operatorname{Dyn}(G)\).

\SGARefTarget{sga3:xxiv:3.4}
\medskip\noindent\textbf{3.4.} ---
This construction has the following properties (which, moreover,
essentially characterize it):
\begin{enumerate}
\item[\SGAInlineRefTarget{sga3:xxiv:3.4:item:i}(i)] With every reductive \(S\)-group \(G\) is associated a
Dynkin scheme \(\operatorname{Dyn}(G)\); with every isomorphism
\[
  u:G\xrightarrow{\ \sim\ }G'
\]
is functorially associated an isomorphism
\[
  \operatorname{Dyn}(u):
  \operatorname{Dyn}(G)\xrightarrow{\ \sim\ }\operatorname{Dyn}(G').
\]

\item[\SGAInlineRefTarget{sga3:xxiv:3.4:item:ii}(ii)] If \(S'\) is an \(S\)-scheme and \(G\) is a reductive
\(S\)-group, then
\[
  \operatorname{Dyn}(G\times_S S')
  \simeq
  \operatorname{Dyn}(G)\times_S S'.
\]

\item[(iii)] If \(E\) is a pinning of \(G\) defining the pinned root
datum with Dynkin diagram \(\Delta\), then
\[
  \operatorname{Dyn}(G)\simeq\Delta_S.
\]

\item[\SGAInlineRefTarget{sga3:xxiv:3.4:item:iv}(iv)] If \(u\) is an inner automorphism of \(G\), then
\(\operatorname{Dyn}(u)\) is the identity automorphism of
\(\operatorname{Dyn}(G)\).
\end{enumerate}

\SGARefTarget{sga3:xxiv:3.5}
\medskip\noindent\textbf{3.5.} ---
Let \(S\) be a scheme and \(G\) a reductive \(S\)-group. It is clear
that the functor
\[
  \underline{\Aut}_{\mathrm{Dyn}}
    \bigl(\operatorname{Dyn}(G)\bigr)
\]
of automorphisms of \(\operatorname{Dyn}(G)\) as a Dynkin scheme is
represented by a finite twisted constant \(S\)-scheme. By \SGARefLink{sga3:xxiv:3.4:item:i}{3.4(i)}
and \SGARefLink{sga3:xxiv:3.4:item:ii}{(ii)}, there is a canonical morphism
\[
  \underline{\Aut}_{S\text{-}\mathrm{gr}}(G)
  \longrightarrow
  \underline{\Aut}_{\mathrm{Dyn}}
    \bigl(\operatorname{Dyn}(G)\bigr),
\]
which, by \SGARefLink{sga3:xxiv:3.4:item:iv}{(iv)}, factors through a morphism
\[
  \underline{\Autext}(G)
  \longrightarrow
  \underline{\Aut}_{\mathrm{Dyn}}
    \bigl(\operatorname{Dyn}(G)\bigr).
\]
More generally, if \(G\) and \(G'\) are two reductive \(S\)-groups,
there is a canonical morphism
\[
  \underline{\operatorname{Isomext}}(G,G')
  \longrightarrow
  \underline{\operatorname{Isom}}_{\mathrm{Dyn}}
    \bigl(\operatorname{Dyn}(G),\operatorname{Dyn}(G')\bigr).
\]
In particular, if \(G'\) is an inner twisted form of \(G\) (\SGARefLink{sga3:xxiv:1.11}{1.11}),
the Dynkin schemes \(\operatorname{Dyn}(G)\) and
\(\operatorname{Dyn}(G')\) are isomorphic.

\SGARefTarget{sga3:xxiv:3.6}
\medskip\noindent\textbf{3.6.} ---
If \(G\) is semisimple (respectively, adjoint or simply connected),
the morphism
\[
  \underline{\Autext}(G)
  \longrightarrow
  \underline{\Aut}_{\mathrm{Dyn}}
    \bigl(\operatorname{Dyn}(G)\bigr)
\]
is a monomorphism (respectively, an isomorphism). Indeed, one may
assume that \(G\) is pinned, and the assertion reduces to the
corresponding result for pinned reduced root data
(cf.\ \SGARefLink{sga3:XXI:7.4.5}{Exp.~XXI, 7.4.5}).

There is an analogous result for the \(\operatorname{Isom}\) functors.
It follows in particular that two semisimple adjoint (respectively,
simply connected) \(S\)-groups are inner twisted forms of one another
if and only if their Dynkin schemes are isomorphic.

% Authority local p. 16 / combined-reader p. 1261.

\SGARefTarget{sga3:xxiv:3.7}
\medskip\noindent\textbf{3.7.} ---
One can give a different construction of the Dynkin scheme associated
with a reductive group. Let \(\mathscr R\) be a pinned reduced root
datum and \(G\) a reductive \(S\)-group of type \(\mathscr R\), and
write \(\Delta(\mathscr R)\) for the Dynkin diagram defined by
\(\mathscr R\). There is, by \SGARefLink{sga3:xxiv:3.5}{3.5}, a canonical morphism
\[
  A_S(\mathscr R)
  =
  \underline{\Aut}_{S\text{-}\mathrm{gr}}
    \bigl(\operatorname{\acute{E}p}_S(\mathscr R)\bigr)
  \longrightarrow
  \underline{\Aut}_{\mathrm{Dyn}}
    \bigl(\Delta(\mathscr R)_S\bigr).
\]
The reductive \(S\)-group \(G\) corresponds (\SGARefLink{sga3:xxiv:1.17}{1.17}) to a bundle
\[
  \underline{\operatorname{Isom}}_{S\text{-}\mathrm{gr}}
    \bigl(\operatorname{\acute{E}p}_S(\mathscr R),G\bigr),
\]
principal homogeneous under \(A_S(\mathscr R)\). The associated
bundle under
\(\underline{\Aut}_{\mathrm{Dyn}}(\Delta(\mathscr R)_S)\)
corresponds to a form over \(S\) of \(\Delta(\mathscr R)_S\): it is
\(\operatorname{Dyn}(G)\). In other words, this associated bundle is
precisely
\[
  \underline{\operatorname{Isom}}_{\mathrm{Dyn}}
    \bigl(\Delta(\mathscr R)_S,\operatorname{Dyn}(G)\bigr).
\]
In this form the proof is immediate.

\SGARefTarget{sga3:xxiv:3.8}
\medskip\noindent\textbf{3.8. Dynkin Scheme and Killing Couples.} ---
Let \(S\) be a scheme, \(G\) a reductive \(S\)-group, and
\(B\supset T\) a Killing couple of \(G\). There is a canonical
morphism
\[
  i:\operatorname{Dyn}(G)
  \longrightarrow
  \underline{\operatorname{Hom}}_{S\text{-}\mathrm{gr}}
    (T,\mathbb G_{m,S})
\]
which identifies \(\operatorname{Dyn}(G)\) with the ``scheme of
simple roots of \(B\) relative to \(T\).'' This morphism is defined
immediately by descent from the pinned case. Observe also that the
data of \(T\) and \(i\) allow one to reconstruct \(B\) (the
``one-to-one correspondence between systems of simple roots and
systems of positive roots'').

It follows from the preceding description of
\(D=\operatorname{Dyn}(G)\) that there is a canonical root of \(B_D\)
relative to \(T_D\): this root \(\alpha_D\) is the image, under
\(i(D)\), of the identity morphism of \(D\). One obtains from it a
canonical invertible \(\mathcal O_D\)-module \(\mathfrak g_D\):
\[
  \mathfrak g_D
  =
  \bigl(\mathfrak g\otimes_{\mathcal O_S}\mathcal O_D\bigr)^{\alpha_D}.
\]
In the pinned case,
\[
  D=\coprod_{\alpha\in\Delta}S_\alpha,
\]
where every \(S_\alpha\) is a copy of \(S\), and
\(\mathfrak g_D\) is then the \(\mathcal O_D\)-module inducing
\(\mathfrak g_\alpha\) on \(S_\alpha\), for every
\(\alpha\in\Delta\).

\SGARefTarget{sga3:xxiv:3.9}
\medskip\noindent\textbf{3.9. Quasi-Pinnings. Quasi-Pinned Groups.} ---
If \(G\) is a reductive \(S\)-group, a \emph{quasi-pinning of \(G\)}
means the data
\begin{enumerate}
\item[(i)] of a Killing couple \(B\supset T\) of \(G\);
\item[(ii)] of a section
\[
  X\in\Gamma\bigl(\operatorname{Dyn}(G),\mathfrak g_D\bigr)^\times.
\]
\end{enumerate}
A reductive \(S\)-group is called \emph{quasi-splittable} if it has a
quasi-pinning. A \emph{quasi-pinned group} is a reductive group
equipped with a quasi-pinning.

Let \(B\supset T\) be a Killing couple of the reductive \(S\)-group
\(G\). Then \(G\) is quasi-pinnable relative to this Killing couple
if and only if \(\mathfrak g_D\) has a section nonzero at every point,
that is, if the element of \(\operatorname{Pic}(\operatorname{Dyn}(G))\)
defined by \(\mathfrak g_D\) is zero. Suppose in particular that
\(S\) is semilocal. Then \(\operatorname{Dyn}(G)\) is also semilocal,
and hence \(\operatorname{Pic}(\operatorname{Dyn}(G))=0\). Thus:

\SGARefTarget{sga3:xxiv:3.9.1}
\medskip\noindent\textbf{Proposition 3.9.1.} ---
{\itshape Let \(S\) be a semilocal scheme and \(G\) a reductive
\(S\)-group. Then \(G\) is quasi-splittable if and only if it has a
Borel subgroup.}


% END INCLUDED SGA3_English_current_buildable_source_20260823/source/inputs/XXIV/tex/components/07_expose_XXIV_section2_conclusion_and_section3_through_39_en.tex

% BEGIN INCLUDED SGA3_English_current_buildable_source_20260823/source/inputs/XXIV/tex/components/08_expose_XXIV_39_proof_through_corollary312_en.tex
% Exposé XXIV, authority local pages 17--20 / combined-reader pages 1262--1265.
% Authority local p. 17 / combined-reader p. 1262.

\footnote{Editor's note: The reference to Exp.~XXII, 5.9.7 has been
given in detail.}
Indeed, \(S\) is affine, so by the first assertion of
\SGARefLink{sga3:xxii:cor:5.9.7}{Exp.~XXII, 5.9.7}, if \(G\) has a Borel subgroup \(B\), it also has a
Killing couple \(B\supset T\). Then, because
\(\operatorname{Pic}(D)=0\), the module \(\mathfrak g_D\) has a
nowhere-zero section \(X\).

Continue to let \(A\) be a semilocal ring and
\(S=\operatorname{Spec}(A)\). For every reductive \(S\)-group \(G\),
the morphism
\[
  \operatorname{Bor}(G)\longrightarrow S
\]
is surjective (because \(G_s\) has Borel subgroups for every
\(s\in S\)), smooth, and projective (\SGARefLink{sga3:xxii:cor:5.8.3}{Exp.~XXII, 5.8.3}). It therefore
has sections after a finite, surjective, étale extension of the
base.\footnote{Editor's note: The original referred to EGA~IV,
\S24, which did not appear. The point is to use ``Bertini's
theorem.'' We give the details of the argument, which were indicated
to us by O.\ Gabber. Let \(X\to S\) be a surjective smooth projective
morphism. Replacing \(S\) by a connected component, one may assume
that \(X/S\) has relative dimension \(d\) at every point
(cf.\ EGA~IV\textsubscript{4}, 17.10.2). One may also assume that
\(X\) is a closed subscheme of a projective space
\(\mathbb P^n_S=\mathbb P(E)\), where \(E\) is a free \(A\)-module
of rank \(n+1\) (cf.\ EGA~II, 5.3.3). Let
\(s_1,\ldots,s_r\) be the closed points of \(S\). By Bertini's
theorem (see, for example, \textup{[Jou83]}, I, 6.10), there is an
open subset \(U\) of
\[
  P=\mathbb P(E^\ast)^d
\]
with nonempty fibers such that, for every point
\(u=(f_1,\ldots,f_d)\) of \(U_{s_i}\), the intersection of
\(X_{\kappa(u)}\) with the \(d\) hyperplanes of
\(\mathbb P^n_{\kappa(u)}\) defined by \(u\) is étale over
\(\kappa(u)\). After shrinking \(U\), one may assume that \(U\) is
the complement in affine space \(\mathbb A^{nd}_S\) of the zero
locus \(V(Q)\) of a polynomial \(Q\) of degree \(m>0\). An easy
induction on the number of variables then shows that \(V(Q)\) cannot
contain every rational point of
\(\mathbb A^{nd}_{\kappa(s_i)}\) if
\(\lvert\kappa(s_i)\rvert>m\). Since the morphism
\[
  A\longrightarrow\prod_i\kappa(s_i)
\]
is surjective, one can find a monic polynomial \(R\in A[X]\) of
degree \(m\) whose image in \(\kappa(s_i)[X]\) is irreducible when
\(\kappa(s_i)\) is finite and has \(m\) distinct roots when
\(\kappa(s_i)\) is infinite. Put
\[
  A'=A[X]/(R),
  \qquad
  S'=\operatorname{Spec}(A').
\]
Then \(S'\to S\) is étale, finite, and surjective, and the residue
field at each of its closed points \(s'_1,\ldots,s'_t\) has
cardinality at least \(2m\).
Thus \(U\) has a rational point \(u_i\) above every closed point of
\(S'\), and, since
\[
  A'\longrightarrow\prod_i\kappa(s'_i)
\]
is surjective \emph{(Translator's note: the French text prints
\(\kappa(s_i)\); the intended closed points of \(S'\) are
\(s'_i\))}, these points lift to a section \(u\) of \(P_{S'}\).
Let \(Z\) be the intersection of \(X_{S'}\) with the \(d\)
hyperplanes of \(\mathbb P^n_{S'}\) defined by \(u\), and let \(V\)
be the open subset of \(Z\) consisting of the points at which \(Z\)
is étale over \(S'\). By EGA~IV\textsubscript{3}, 11.3.8, \(V\)
contains the fibers \(Z_{s'_i}\) for every \(i\). Since
\(\pi:Z\to S'\) is proper, it follows that the closed set
\(\pi(Z-V)\) is empty, whence \(V=Z\). Thus \(Z\to S'\) is
surjective, étale, and proper, hence finite, as is the composite
\(Z\to S\); this gives the desired section of \(X\to S\).}
It follows that:

\SGARefTarget{sga3:xxiv:3.9.2}
\medskip\noindent\textbf{Corollary 3.9.2.} ---
{\itshape Let \(S\) be a semilocal scheme and \(G\) a reductive
\(S\)-group. There is an étale, finite, and surjective morphism
\(S'\to S\) such that \(G_{S'}\) is quasi-splittable.}

\SGARefTarget{sga3:xxiv:3.9.3}
\medskip\noindent\textbf{Remark 3.9.3.} ---
Under the preceding hypotheses, let \(T\) be a maximal torus of
\(G\) (cf.\ \SGARefLink{sga3:xiv:3.20}{Exp.~XIV, 3.20}). One may moreover assume that \(G_{S'}\)
is quasi-splittable relative to \(T_{S'}\): it suffices to apply the
preceding argument to the ``scheme of Borel subgroups containing
\(T\),'' which is finite and étale over \(S\)
(\SGARefLink{sga3:xxii:corollary:5-5-5:item:ii}{Exp.~XXII, 5.5.5(ii)}).

\SGARefTarget{sga3:xxiv:3.10}
\medskip\noindent\textbf{3.10.} ---
Let \(E\) and \(E'\) be two quasi-pinnings of the reductive
\(S\)-group \(G\). There is a unique inner automorphism of \(G\)
carrying \(E\) to \(E'\). Indeed, one reduces immediately to the
split case, where the assertion has already been proved (\SGARefLink{sga3:xxiv:1.5}{1.5}; it
suffices to observe that for an inner automorphism of \(G\), preserving
a pinning is equivalent to preserving the underlying quasi-pinning).
One concludes, as in \S1, that a

% Authority local p. 18 / combined-reader p. 1263.

quasi-pinning of the reductive \(S\)-group \(G\) defines a splitting
\(h\) of the exact sequence:

\[
  % BEGIN INCLUDED SGA3_English_current_buildable_source_20260823/source/inputs/XXIV/native_diagrams/exp24/001_split_autext_exact_sequence.tex
\begin{tikzcd}[column sep=small]
1 \arrow[r]
&
\underline{\operatorname{ad}}(G) \arrow[r]
&
\underline{\Aut}_{S\text{-}\mathrm{gr.}}(G)
  \arrow[r,shift right=.7ex,"p"']
&
\underline{\Autext}(G)
  \arrow[l,shift right=.7ex,dotted,"h"']
  \arrow[r]
&
1
\end{tikzcd}

% END INCLUDED SGA3_English_current_buildable_source_20260823/source/inputs/XXIV/native_diagrams/exp24/001_split_autext_exact_sequence.tex

\]

the image of \(h\) being the subgroup of
\(\underline{\Aut}_{S\text{-}\mathrm{gr}}(G)\) that preserves the
quasi-pinning.

Likewise, if \(G\) and \(G'\) are two quasi-pinned \(S\)-groups, one
defines naturally the subfunctor
\[
  \underline{\operatorname{Isom}}_{S\text{-}\mathrm{gr},\,q\text{-}\acute{e}p}
    (G,G')
\]
of
\(\underline{\operatorname{Isom}}_{S\text{-}\mathrm{gr}}(G,G')\).
The projection from
\(\underline{\operatorname{Isom}}_{S\text{-}\mathrm{gr}}(G,G')\)
to \(\underline{\operatorname{Isomext}}(G,G')\) induces an
isomorphism
\[
  (*)\qquad
  \underline{\operatorname{Isom}}_{S\text{-}\mathrm{gr},\,q\text{-}\acute{e}p}
    (G,G')
  \xrightarrow{\ \sim\ }
  \underline{\operatorname{Isomext}}(G,G').
\]

\SGARefTarget{sga3:xxiv:3.11}
\medskip\noindent\textbf{Theorem 3.11.} ---
{\itshape Let \(S\) be a scheme, let \(\mathscr R\) be a pinned reduced
root datum such that
\(\operatorname{\acute{E}p}_S(\mathscr R)\) exists
(cf.\ Exp.~XXV), and let \(E\) be its group of automorphisms.
Consider the following three categories:
\begin{enumerate}
\item[\SGAInlineRefTarget{sga3:xxiv:3.11:item:i}(i)] the category \(\mathsf{Rev}\) of principal Galois coverings
of \(S\) with group \(E\) (the morphisms are isomorphisms of
\(E\)-bundles);
\item[\SGAInlineRefTarget{sga3:xxiv:3.11:item:ii}(ii)] the category \(\mathsf{Redext}\), whose objects are the
reductive \(S\)-groups of type \(\mathscr R\) (\SGARefLink{sga3:xxii:def:2.7}{Exp.~XXII, 2.7}), the
morphisms from \(G\) to \(G'\) being the elements of
\(\underline{\operatorname{Isomext}}(G,G')(S)\);
\item[(iii)] the category \(\mathsf{Q\acute{e}p}\) of quasi-pinned
reductive \(S\)-groups of type \(\mathscr R\) (the morphisms are the
isomorphisms preserving the quasi-pinnings).
\end{enumerate}
These three categories are equivalent. More precisely, there is a
diagram of functors, commutative up to functorial
isomorphisms:}

\[
  % BEGIN INCLUDED SGA3_English_current_buildable_source_20260823/source/inputs/XXIV/native_diagrams/exp24/002_rev_qep_redext_triangle.tex
\begin{tikzcd}[column sep=large,row sep=large]
\mathsf{Rev}
  \arrow[rr,"q\acute{e}p"]
&&
\mathsf{Q\acute{e}p}
  \arrow[dl,"i"]
\\
&
\mathsf{Redext}
  \arrow[ul,"\operatorname{rev}"]
&
\end{tikzcd}

% END INCLUDED SGA3_English_current_buildable_source_20260823/source/inputs/XXIV/native_diagrams/exp24/002_rev_qep_redext_triangle.tex

\]

We describe these three functors below, leaving to the reader the
verification that the diagram commutes.\footnote{Editor's note: In
\SGARefLink{sga3:xxiv:3.11.4}{3.11.4} we have sketched the verification that the functors \(i\),
\(\operatorname{rev}\), and \(q\acute{e}p\) are fully faithful and
that the composite
\(\operatorname{rev}\circ i\circ q\acute{e}p\) is isomorphic to the
identity functor of \(\mathsf{Rev}\).}

\SGARefTarget{sga3:xxiv:3.11.1}
\medskip\noindent\textbf{3.11.1. The functor \(i\).} ---
This is the evident functor:
\[
  i(G)=G,
  \qquad
  i(f)=\text{the image of \(f\) under the isomorphism 3.10 \((*)\)}.
\]

\SGARefTarget{sga3:xxiv:3.11.2}
\medskip\noindent\textbf{3.11.2. The functor
\(q\acute{e}p\).} ---
Let \(S'\) be a principal Galois covering of \(S\) with group \(E\).
Let
\[
  \bigl(G,T,M,\mathscr R,\Delta,
    (X_\alpha)_{\alpha\in\Delta}\bigr)
\]
be a pinning of
\(G=\operatorname{\acute{E}p}_S(\mathscr R)\), and let
\((X'_\alpha)_{\alpha\in\Delta}\) be the corresponding pinning of
\[
  G'=G_{S'}=\operatorname{\acute{E}p}_{S'}(\mathscr R).
\]
By \SGARefLink{sga3:XXIII:5.5bis}{Exp.~XXIII, 5.5 bis}, there is a unique group homomorphism
\[
  \theta:E
  \longrightarrow
  \underline{\Aut}_{S\text{-}\mathrm{gr}}
    \bigl(\operatorname{\acute{E}p}_S(\mathscr R)\bigr)
  =A(\mathscr R)(S)
\]

% Authority local p. 19 / combined-reader p. 1264.

such that, for every \(h\in E\), the automorphism \(\theta(h)\)
induces \(D_S(h)\) on \(T\), and
\[
  \Lie(\theta(h))(X_\alpha)=X_{h(\alpha)}
  \qquad(\alpha\in\Delta).
\]
Taking into account the action of \(E\) on \(S'\), this yields an
action of \(E\) on \(G'\), compatible with the action of \(E\) on
\(S'\) and preserving the canonical quasi-pinning of
\(G'\).\footnote{Editor's note: The preceding discussion expands the
original in order to show that the action of \(E\) on \(G'\) is
obtained by combining the natural action of \(E\) on
\(\operatorname{\acute{E}p}_{S'}(\mathscr R)\) with the given action
on \(S'\). The canonical Killing couple of \(G'\) is preserved by
this action, as is the quasi-pinning given by the element
\(\widetilde X'\) of
\[
  \Gamma\!\left(
    \Delta_{S'},
    \Lie(G'/S')_{\Delta_{S'}}
  \right),
\]
equal to \(X'_\alpha\) on the copy of \(S'\) indexed by \(\alpha\).
Indeed, since \(E\) acts on \(\Delta_{S'}\) by permuting the copies
of \(S'\), one has \(h(\widetilde X')=\widetilde X'\) for every
\(h\in E\).}
Since \(S'\to S\) is a covering for the \textup{(fpqc)} topology, it
is an effective descent morphism for the fibered category of affine
morphisms. Write
\[
  q\acute{e}p(S')
  =
  Q\text{-}\operatorname{\acute{E}p}_{S'/S}(\mathscr R)
\]
for the quasi-pinned \(S\)-group obtained by Galois
descent.\footnote{Editor's note: See, for example, TDTE~I, p.~22,
\SGARefLink{sga3:xxiv:section:1}{Example~1}. One may also describe \(q\acute{e}p(S')\) as the twist of
\(G=\operatorname{\acute{E}p}_S(\mathscr R)\) by the \(E\)-torsor
\(S'/S\), that is, the \textup{(fpqc)} sheaf quotient of
\(S'\times_SG\) by the right action of \(E\) defined by
\[
  (s',g)\cdot h=(s'h,h^{-1}(g)).
\]
Since \(G\) is affine over \(S\), and since \(E\) acts on \(G\) by
group automorphisms, this sheaf is represented by an \(S\)-group
\(G^\sharp\), which is a ``twisted'' form of \(G\), and
\(D^\sharp=\operatorname{Dyn}(G^\sharp)\) is the twist of the
constant Dynkin scheme \(\Delta_S\) by the torsor \(S'/S\). Since
\(E\) normalizes \(B\) and \(T\), one obtains likewise a Killing
couple \((B^\sharp,T^\sharp)\) of \(G^\sharp\). On the other hand,
put
\(\mathfrak g=\Lie(G/S)\) and
\(\mathfrak g^\sharp=\Lie(G^\sharp/S)\). For every \(U\to S\), the
sections of \(\mathfrak g^\sharp\) over \(U\) are the
\(E\)-equivariant \(S\)-morphisms
\[
  U\times_SS'\longrightarrow W(\mathfrak g).
\]
Since \(D^\sharp\times_SS'\) is \(E\)-isomorphic to
\(\Delta\times S'\), endowed with the action
\[
  (\alpha,s')\cdot h=(h^{-1}(\alpha),s'h),
\]
the morphism \((\alpha,s')\mapsto X_\alpha\) gives a section of
\(\mathfrak g^\sharp\) over \(D^\sharp\) which is a quasi-pinning,
that is, a nowhere-zero section of
\[
  \bigl(
    \mathfrak g^\sharp\otimes_{\mathcal O_S}\mathcal O_{D^\sharp}
  \bigr)^{D^\sharp}
\]
(cf.\ \SGARefLink{sga3:xxiv:3.8}{3.8}).
}

\SGARefTarget{sga3:xxiv:3.11.3}
\medskip\noindent\textbf{3.11.3. The functor
\(\operatorname{rev}\).} ---
Let \(G\) be a reductive \(S\)-group of type \(\mathscr R\). Write
\[
  \operatorname{rev}(G)
  =
  \underline{\operatorname{Isomext}}
    \bigl(\operatorname{\acute{E}p}_S(\mathscr R),G\bigr).
\]
This is a principal homogeneous bundle under \(E_S\)
(cf.\ \SGARefLink{sga3:xxiv:1.10}{1.10} and \SGARefLink{sga3:xxiv:1.19}{1.19}), that is, an object of \(\mathsf{Rev}\).

\SGARefTarget{sga3:xxiv:3.11.4}
\medskip\noindent\textbf{3.11.4.} ---
\footnote{Editor's note: This subsection number has been added.}
It is clear from the isomorphism \SGARefLink{sga3:xxiv:3.10}{3.10} \((*)\) that the functor \(i\)
is fully faithful. On the other hand, if \(G_0,G,G'\) are reductive
\(S\)-groups, the natural morphism
\[
  \underline{\operatorname{Isom}}(G_0,G)
  \times_S
  \underline{\operatorname{Isom}}(G,G')
  \longrightarrow
  \underline{\operatorname{Isom}}(G_0,G')
\]
induces a morphism
\[
  (\star)\qquad
  \underline{\operatorname{Isomext}}(G_0,G)
  \times_S
  \underline{\operatorname{Isomext}}(G,G')
  \longrightarrow
  \underline{\operatorname{Isomext}}(G_0,G'),
\]
since, for every \(S'\to S\), if
\[
  g_0\in G_0(S'),\quad g\in G(S'),\quad
  u\in\underline{\operatorname{Isom}}(G_0,G)(S'),\quad
  v\in\underline{\operatorname{Isom}}(G,G')(S'),
\]
one has
\[
  v\circ\operatorname{int}(g)\circ u\circ\operatorname{int}(g_0)
  =
  v\circ u\circ
  \operatorname{int}\bigl(u^{-1}(g)g_0\bigr).
\]
Taking \(G_0=\operatorname{\acute{E}p}_S(\mathscr R)\), one obtains
a map
\[
\begin{aligned}
  \underline{\operatorname{Isomext}}(G,G')(S)
  \longrightarrow
  \operatorname{Hom}_{\mathsf{Rev}}\bigl(
  &\underline{\operatorname{Isomext}}
    (\operatorname{\acute{E}p}_S(\mathscr R),G),\\
  &\underline{\operatorname{Isomext}}
    (\operatorname{\acute{E}p}_S(\mathscr R),G')
  \bigr),
\end{aligned}
\]
and, to show that this map is a bijection, one may assume by
\textup{(fpqc)} descent that \(G\) and \(G'\) are pinned, in which
case it is evident.

% Authority local p. 20 / combined-reader p. 1265.

Similarly, if \(S_1,S_2\) are two objects of \(\mathsf{Rev}\), and
if one writes
\[
  S_i\times^E\operatorname{\acute{E}p}_S(\mathscr R)
\]
for the \(S\)-group \(q\acute{e}p(S_i)\) obtained by twisting
\(\operatorname{\acute{E}p}_S(\mathscr R)\) by the \(E\)-torsor
\(S_i\) (cf.\ editor's note~(21)), there is a natural map
\[
\begin{aligned}
  \operatorname{Hom}_{\mathsf{Rev}}(S_1,S_2)
  \longrightarrow
  \operatorname{Hom}_{\mathsf{Q\acute{e}p}}\bigl(
    &S_1\times^E\operatorname{\acute{E}p}_S(\mathscr R),\\
    &S_2\times^E\operatorname{\acute{E}p}_S(\mathscr R)
  \bigr).
\end{aligned}
\]
To show that this map is a bijection, one may assume by
\textup{(fpqc)} descent that \(S_1\) and \(S_2\) are trivial
\(E\)-bundles, in which case the assertion is evident.

Thus the three functors in the diagram are fully faithful. Moreover,
for every object \(S'\) of \(\mathsf{Rev}\), there is a natural
morphism
\[
  S'
  \longrightarrow
  \underline{\operatorname{Isomext}}\!\left(
    \operatorname{\acute{E}p}_S(\mathscr R),
    S'\times^E\operatorname{\acute{E}p}_S(\mathscr R)
  \right),
\]
and to see that it is an isomorphism, one may assume by
\textup{(fpqc)} descent that \(S'\) is a trivial \(E\)-bundle, in
which case the assertion is evident.

There is therefore an isomorphism of functors
\[
  \operatorname{Id}_{\mathsf{Rev}}
  \longrightarrow
  \operatorname{rev}\circ i\circ q\acute{e}p.
\]
Since \(\operatorname{rev}\) is fully faithful, for every reductive
\(S\)-group \(H\) (not necessarily quasi-split) one obtains a
bijection
\[
  \underline{\operatorname{Isomext}}\!\left(
    Q\text{-}\operatorname{\acute{E}p}_{\operatorname{rev}(H)/S}
      (\mathscr R),
    H
  \right)(S)
  \simeq
  \operatorname{Aut}_{\mathsf{Rev}}\bigl(\operatorname{rev}(H)\bigr),
\]
functorial in \(H\). In particular, the identity morphism of
\(\operatorname{rev}(H)\) corresponds to a ``canonical'' element
\[
  u_0\in
  \underline{\operatorname{Isomext}}\!\left(
    Q\text{-}\operatorname{\acute{E}p}_{\operatorname{rev}(H)/S}
      (\mathscr R),
    H
  \right)(S).
\]
Moreover, every element \(u\) of this same set corresponds to an
automorphism \(\varphi_u\) of \(\operatorname{rev}(H)\), and
\(q\acute{e}p(\varphi_u)\) is the unique automorphism of the pinned
\(S\)-group
\[
  Q\text{-}\operatorname{\acute{E}p}_{\operatorname{rev}(H)/S}
    (\mathscr R)
\]
such that
\[
  u_0\circ q\acute{e}p(\varphi_u)=u.
\]
This proves \SGARefLink{sga3:xxiv:3.11}{Theorem 3.11}.

We develop one of the corollaries of \SGARefLink{sga3:xxiv:3.11}{3.11}.

\SGARefTarget{sga3:xxiv:3.12}
\medskip\noindent\textbf{Corollary 3.12.} ---
{\itshape For every reductive \(S\)-group \(G\), there is a quasi-pinned
\(S\)-group \(G_{q\text{-}\acute{e}p}\) and an ``outer isomorphism''}
\[
  u\in
  \underline{\operatorname{Isomext}}
    (G_{q\text{-}\acute{e}p},G)(S).
\]
{\itshape The pair \((G_{q\text{-}\acute{e}p},u)\) is unique up to a
unique isomorphism.}

Indeed, one may assume that \(G\) has constant type \(\mathscr R\),
and one takes
\[
  G_{q\text{-}\acute{e}p}
  =
  Q\text{-}\operatorname{\acute{E}p}_{\operatorname{rev}(G)/S}
    (\mathscr R).
\]

\SGARefTarget{sga3:xxiv:3.12.1}
\medskip\noindent\textbf{3.12.1.} ---
Observe that the data of the pair
\((G_{q\text{-}\acute{e}p},u)\) canonically define the \(S\)-scheme
(cf.\ \SGARefLink{sga3:xxiv:1.11}{1.11})
\[
  Q=
  \underline{\operatorname{Isomint}}
    (G_{q\text{-}\acute{e}p},G),\footnote{Editor's note: We write
    \(\underline{\operatorname{Isomint}}
    (G_{q\text{-}\acute{e}p},G)\) instead of
    \(\underline{\operatorname{Isomint}}_u
    (G_{q\text{-}\acute{e}p},G)\) (cf.\ 1.11), since the pair
    \((G_{q\text{-}\acute{e}p},u)\) is unique up to a unique
    isomorphism.}
\]
which is a principal homogeneous bundle under
\(\operatorname{ad}(G_{q\text{-}\acute{e}p})\), and whose data are
``equivalent'' to those of the inner twisted form \(G\) of
\(G_{q\text{-}\acute{e}p}\). Moreover, \(Q\) is precisely the
``scheme of quasi-pinnings of \(G\),'' a definition that explains its
structure as a principal homogeneous bundle (on the left) under
\(\operatorname{ad}(G)\) (\SGARefLink{sga3:xxiv:3.10}{3.10}); compare \SGARefLink{sga3:xxiv:1.20}{1.20}.

% END INCLUDED SGA3_English_current_buildable_source_20260823/source/inputs/XXIV/tex/components/08_expose_XXIV_39_proof_through_corollary312_en.tex

% BEGIN INCLUDED SGA3_English_current_buildable_source_20260823/source/inputs/XXIV/tex/components/09_expose_XXIV_proposition313_through_remark3168_en.tex
% Authority coverage: Polo--Gille Exposé XXIV, local pages 21--24.
% Lead-reviewed against the authority PDF; the Reinhold Markdown lineage was comparison-only.

\textbf{\SGAInlineRefTarget{sga3:xxiv:proposition:3-13}Proposition 3.13.} \emph{Let \(S\) be a scheme, \(G\) a semisimple adjoint (resp. simply connected) \(S\)-group, \(B \supset T\) a
Killing couple of \(G\), \(\operatorname{Dyn}(G)\) the Dynkin \(S\)-scheme}

\emph{of \(G\). There exists a canonical isomorphism of \(S\)-group schemes}

\[
  T\xrightarrow{\sim}
  \prod_{\operatorname{Dyn}(G)/S}
  \mathbb G_{m,\operatorname{Dyn}(G)}.
\]

\emph{(Recall (cf.~Exp. II, § 1) that the right-hand side is by definition the \(S\)-functor which to \(S' \to S\) associates
\(G_{m}(\operatorname{Dyn}(G) \times_{S} S')\), or, what amounts to the same, \(G_{m}(Dyn(G_{S'}))\).)}

\emph{First proof.} Let us do it for simplicity in the adjoint case. Consider the composite morphism

\[
  T\longrightarrow\prod_{D/S}T_D
  \longrightarrow\prod_{D/S}\mathbb G_{m,D}.
\]

where the first morphism is the canonical morphism, the second is \(\prod_{D/S} \alpha_{D}\) (we have written \(D =
\operatorname{Dyn}(G)\)). To verify that this morphism is an isomorphism, one may assume \(G\) split; now, in this case, this is none
other than the morphism \(T \to (G_{m, S})^{\Delta}\) with components \(\alpha\), for \(\alpha \in \Delta\), and the latter is
an isomorphism (\SGARefLink{sga3:xxii:rem:4.3.8}{Exp. XXII, 4.3.8}).

\emph{Second proof.} By \SGARefLink{sga3:xxiv:2.8}{2.8}, \SGARefLink{sga3:xxiv:3.5}{3.5} and \SGARefLink{sga3:xxiv:3.11}{3.11}, one may assume that \(G = Q-\acute{E}p_{S'/S}(R)\), \(T\) being the canonical maximal
torus. On \(S'\), one has by \SGARefLink{sga3:xxii:rem:4.3.8}{Exp. XXII 4.3.8}, an isomorphism \(T_{S'} \xrightarrow{\sim} (G_{m, S'})^{\Delta}\), defined by
the simple roots (resp. the simple coroots). The group \(E\) operates on the right-hand side by permutation of \(\Delta\).
Now the bundle associated to \(S'/S\) by \(E \to \operatorname{Aut}(\Delta)\) is \(\operatorname{Dyn}(G)\) (\SGARefLink{sga3:xxiv:3.7}{3.7}), and one concludes at once.

Using Prop. \SGARefLink{sga3:xxiv:proposition:8-4}{8.4} of the appendix, one deduces:

\textbf{\SGAInlineRefTarget{sga3:xxiv:corollary:3-14}Corollary 3.14.} \emph{Under the preceding conditions, one has}

\[
  H^1(S,T)\xrightarrow{\sim}
  H^1\bigl(\operatorname{Dyn}(G),\mathbb G_m\bigr)
  \xrightarrow{\sim}\operatorname{Pic}\bigl(\operatorname{Dyn}(G)\bigr).
\]

\emph{In particular, \(H^{1}(S, T) = 0\) when \(S\) is semi-local.}

\textbf{\SGAInlineRefTarget{sga3:xxiv:remark:3-15}Remark 3.15.} Let \(S\) be a scheme, \(G\) (resp. \(G'\)) an \(S\)-reductive group, \(B \supset T\) (resp. \(B' \supset T'\)) a
Killing couple of \(G\) (resp. \(G'\)), \(u \in Isomext(G, G')(S)\). Set (cf.~\SGARefLink{sga3:xxiv:2.10}{2.10})

\[
  P=\underline{\operatorname{Isomint}}_u(G,B,T;G',B',T').
\]

this is a principal homogeneous bundle under \(T_{ad}\) (by \((f, t) \mapsto f \circ int(t)\)).

Let on the other hand \(D = \operatorname{Dyn}(G) = \operatorname{Dyn}(G')\) (identified via \(u\) (\SGARefLink{sga3:xxiv:3.5}{3.5})), and let \(L = Isom_{D}(g_{D}, g'_{D})\) be the
principal homogeneous bundle under \(G_{m,D}\) defined by the invertible \(\mathcal O_D\)-module

\[
  \operatorname{Hom}_{\mathcal O_D}(\mathfrak g_D,\mathfrak g'_D)
  =\mathfrak g'_D\otimes(\mathfrak g_D)^\vee .
\]

Each \(f \in P(S')\) defines an isomorphism of \(g_{D} \otimes_{O_{S}} O_{S'}\) onto \(g'_{D} \otimes_{O_{S}} O_{S'}\), whence
a canonical morphism

\[ P \to \prod_{D/S} L. \]

This morphism is an isomorphism, compatible with the isomorphism of operator groups

\[
  T_{\mathrm{ad}}\xrightarrow{\sim}
  \prod_{D/S}\mathbb G_{m,D}.
\]

defined above. Indeed, it suffices to verify it in the case where the two groups are pinned, where it is easy.

It follows in particular that in the isomorphism

\[ H^{1}(S, T_{ad}) \xrightarrow{\sim} \operatorname{Pic}(\operatorname{Dyn}(G)) \]

of \SGARefLink{sga3:xxiv:corollary:3-14}{3.14}, the class of the bundle \(P\) is transformed into \(c\ell(g'^{D}) - c\ell(g_{D})\). The bundle \(P\) is therefore
trivial if and only if \(g'^{D}\) and \(g_{D}\) are isomorphic.

If \((G, B, T)\) is quasi-splittable, for example if one takes for \(G\) the group \(G'_{q-\acute{e}p.}\), with its canonical
Killing couple, it follows that the image of the class of \(P\) is none other than the obstruction to the quasi-splitting
of \(G'\) defined in \SGARefLink{sga3:xxiv:3.9}{3.9}.

\subsection*{\SGAInlineRefTarget{sga3:xxiv:subsection:3-16}3.16. Symmetries}

\SGAInlineRefTarget{sga3:xxiv:subsubsection:3-16-1}\subsubsection*{3.16.1.}

Let \(S\) be a scheme, \(G\) an \(S\)-reductive group, \(B \supset T\) a Killing couple of \(G\). Write \(D = \operatorname{Dyn}(G)\). Recall (\SGARefLink{sga3:xxii:prop:5.9.2}{Exp.
XXII, 5.9.2}) that there exists a unique Borel subgroup \(B^{-}\) of \(G\) such that \(B \cap B^{-} = T\). If \(X \in \Gamma(D,
g_{D})^{\times}\) defines a quasi-pinning of \(G\) relative to \((B, T)\) (\SGARefLink{sga3:xxiv:3.9}{3.9}), then \(Y = -X^{-1} \in \Gamma(D,
(g_{D})^{\vee})^{\times}\) defines a quasi-pinning of \(G\) relative to \((B^{-}, T)\); one says that this is the \emph{opposite
quasi-pinning}.

If \(R\) is a pinned reduced root datum and if \(E\) is an \(R\)-pinning of the \(S\)-reductive group \(G\), one defines an
\(R\)-pinning \(E^{-}\) said to be \emph{opposite to \(E\)} in the following way: one keeps the same maximal torus \(T\), one takes
the opposite of the isomorphism \(D_{S}(M) \xrightarrow{\sim} T\), and one ``pins'' by \(Y_{\alpha} = -X^{-1}_{\alpha} \in
\Gamma(S, g^{-\alpha})^{\times}\), for \(\alpha \in \Delta\). The quasi-pinning underlying \(E^{-}\) is the quasi-pinning
opposite to the quasi-pinning underlying \(E\).

\textbf{Remark.} In the notation of \SGARefLink{sga3:xx:3.1}{Exp. XX, 3.1}, if one sets

\[
  w_\alpha(X_\alpha)
  =\exp(X_\alpha)\exp(-X_\alpha^{-1})\exp(X_\alpha).
\]

one has \(w_{\alpha}(X_{\alpha}) = w_{-\alpha}(Y_{\alpha})\) (loc. cit., 3.1 (vi)).

\textbf{\SGAInlineRefTarget{sga3:xxiv:proposition:3-16-2}Proposition 3.16.2.} \emph{Let \(S\) be a scheme, \(G\) an \(S\)-reductive group.}

\emph{\SGAInlineRefTarget{sga3:xxiv:proposition:3-16-2:item:i}(i) Let \(T\) be a maximal torus of \(G\); there exists a unique}

\[
  i_T\in
  \left(
    \underline{\operatorname{Aut}}_{S\text{-}\mathrm{gr.}}(G,T)/T_{\mathrm{ad}}
  \right)(S)
  \subset
  \underline{\operatorname{Aut}}_{S\text{-}\mathrm{gr.}}(T).
\]

\emph{such that \(i_{T}(t) = t^{-1}\) for every section \(t\) of \(T\).}

\emph{\SGAInlineRefTarget{sga3:xxiv:proposition:3-16-2:item:ii}(ii) Let \((B, T)\) be a Killing couple of \(G\); there exists a unique section}

\[
  w_{B,T}\in
  \bigl(\operatorname{Norm}_G(T)/T\bigr)(S)=W_G(T)(S).
\]

\emph{such that \(int(w_{B, T})(B) = B^{-}\) (with the obvious abuse of language).}

\emph{\SGAInlineRefTarget{sga3:xxiv:proposition:3-16-2:item:iii}(iii) Let \(Q = (B, T, X)\) be a quasi-pinning of \(G\), \(Q^{-} = (B^{-}, T, Y)\) the opposite quasi-pinning; there exists a
unique inner automorphism of \(G\)}

\[ n_{Q} \in Norm_{ad(G)}(T_{ad})(S) \subset \underline{\operatorname{Aut}}_{S\text{-}\mathrm{gr.}}(G) \]

\emph{such that \(n_{Q}(Q) = Q^{-}\), that is, \(n_{Q}(T) = T\), \(n_{Q}(B) = B^{-}\), \(n_{Q}(X) = Y\).}

\emph{\SGAInlineRefTarget{sga3:xxiv:proposition:3-16-2:item:iv}(iv) Let \((R, E)\) be a pinning of \(G\), \((R, E^{-})\) the opposite pinning; there exists a unique automorphism of \(G\)}

\[
  u_E\in
  \underline{\operatorname{Aut}}_{S\text{-}\mathrm{gr.}}(G,T)
  \subset
  \underline{\operatorname{Aut}}_{S\text{-}\mathrm{gr.}}(G).
\]

\emph{such that \(u_{E}(E) = E^{-}\), i.e.~\(u_{E}(t) = t^{-1}\) for every section \(t\) of \(T\), and \(Ad(u_{E}) X_{\alpha} =
Y_{\alpha}\) for every \(\alpha \in \Delta\).}

\emph{Proof.} \SGARefLink{sga3:xxiv:proposition:3-16-2:item:ii}{(ii)} follows from \SGARefLink{sga3:xxii:corollary:5-5-5:item:ii}{Exp. XXII, 5.5.5 (ii)}, then \SGARefLink{sga3:xxiv:proposition:3-16-2:item:iii}{(iii)} follows from \SGARefLink{sga3:xxiv:3.10}{3.10}, and \SGARefLink{sga3:xxiv:proposition:3-16-2:item:iv}{(iv)} follows from \SGARefLink{sga3:xxiii:theorem:4-1}{Exp. XXIII, 4.1}.
Finally to prove \SGARefLink{sga3:xxiv:proposition:3-16-2:item:i}{(i)}, one may assume \(G\) pinned. Existence follows from \SGARefLink{sga3:xxiv:proposition:3-16-2:item:iv}{(iv)} for example, uniqueness from the fact that
an automorphism of \(G\) inducing the identity on \(T\) is given by a section of \(T_{ad}\) (\SGARefLink{sga3:xxiv:2.11}{2.11}).

\textbf{\SGAInlineRefTarget{sga3:xxiv:corollary:3-16-3}Corollary 3.16.3.} \emph{One has}

\[
  i_T^2=w_{B,T}^2=n_Q^2=u_E^2=e.
\]

\emph{Moreover, \(i_{T}\) (resp. \(u_{E}\)) is \(\neq e\) if \(G \neq e\), and \(w_{B, T}\) (resp. \(n_{Q}\)) is \(\neq e\) if \(G\) is not a
torus.}

\textbf{\SGAInlineRefTarget{sga3:xxiv:corollary:3-16-4}Corollary 3.16.4.} \emph{In the situation of \SGARefLink{sga3:xxiv:proposition:3-16-2:item:iii}{(iii)} (resp. \SGARefLink{sga3:xxiv:proposition:3-16-2:item:iv}{(iv)}), \(n_{Q}\) projects onto \(w_{B, T}\) (resp. \(u_{E}\) projects
onto \(i_{T}\)) by the canonical morphism}

\[ Norm_{ad(G)}(T_{ad}) \to W_{ad(G)}(T_{ad}) \simeq W_{G}(T), \]

\emph{resp.}

\[
  \underline{\operatorname{Aut}}_{S\text{-}\mathrm{gr.}}(G,T)
  \longrightarrow
  \underline{\operatorname{Aut}}_{S\text{-}\mathrm{gr.}}(G,T)/T_{\mathrm{ad}}.
\]

\textbf{\SGAInlineRefTarget{sga3:xxiv:corollary:3-16-5}Corollary 3.16.5.} \emph{The preceding definitions are compatible with base change, and are functorial under isomorphism
(in an obvious sense).}

\textbf{\SGAInlineRefTarget{sga3:xxiv:proposition:3-16-6}Proposition 3.16.6.} \emph{\SGAInlineRefTarget{sga3:xxiv:proposition:3-16-6:item:i}(i) One can define uniquely for each reductive group \(G\) over a scheme \(S\) an element}

\[ s_{G} \in Autext(G)(S) \]

\emph{in such a way that this construction is functorial in \(G\) under isomorphism, is compatible with base change, and is
such that whenever \(T\) is a maximal torus of the \(S\)-reductive group \(G\), \(s_{G}\) is the image of \(i_{T}\) under the
canonical morphism}

\[
  \underline{\operatorname{Aut}}_{S\text{-}\mathrm{gr.}}(G,T)/T_{\mathrm{ad}}
  \longrightarrow\underline{\underline{\operatorname{Autext}}}(G).
\]

\emph{\SGAInlineRefTarget{sga3:xxiv:proposition:3-16-6:item:ii}(ii) One has \(s^{2}_{G} = e\), and \(s_{G}\) is a central element of \(Autext(G)\).}

\emph{\SGAInlineRefTarget{sga3:xxiv:proposition:3-16-6:item:iii}(iii) Under the conditions of \SGARefLink{sga3:xxiv:proposition:3-16-2:item:ii}{3.16.2 (ii)}, if one identifies \(\underline{\operatorname{Aut}}_{S\text{-}\mathrm{gr.}}(G, B, T) / T_{ad}\) with
\(Autext(G)\) (\SGARefLink{sga3:xxiv:2.2}{2.2}), one has}

\[
  w_{B,T}i_T=i_Tw_{B,T}=s_G.
\]

\emph{\SGAInlineRefTarget{sga3:xxiv:proposition:3-16-6:item:iv}(iv) Under the conditions of \SGARefLink{sga3:xxiv:proposition:3-16-2:item:iv}{3.16.2 (iv)}, if one identifies \(\operatorname{Aut}_{S-gr.-\acute{e}p}(G, R, E)\) with
\(Autext(G)\) (\SGARefLink{sga3:xxiv:1.3:item:iii}{1.3 (iii)}), one has}

\[
  n_Eu_E=u_En_E=s_G.
\]

\emph{Proof.} \SGARefLink{sga3:xxiv:proposition:3-16-6:item:i}{(i)} is proved without difficulty by descent. On the other hand, since \(i_{T}\) is evidently a central section of
square \(e\) in \(\underline{\operatorname{Aut}}_{S\text{-}\mathrm{gr.}}(T)\), \SGARefLink{sga3:xxiv:proposition:3-16-6:item:ii}{(ii)} follows immediately; \SGARefLink{sga3:xxiv:proposition:3-16-6:item:iii}{(iii)} is a consequence of \SGARefLink{sga3:xxiv:proposition:3-16-6:item:iv}{(iv)} by descent.
Finally, under the conditions of \SGARefLink{sga3:xxiv:proposition:3-16-6:item:iv}{(iv)}, it is clear that \(n_{E} u_{E} = u_{E} n_{E}\) and that this automorphism of \(G\)
respects the pinning; modulo the identification made, it is therefore equal to its image in \(Autext(G)\); but \(n_{E}\) is
inner and \(u_{E}\) projects onto \(s_{G}\).

\textbf{\SGAInlineRefTarget{sga3:xxiv:remark:3-16-7}Remark 3.16.7.} \SGARefLink{sga3:xxiv:proposition:3-16-6:item:i}{(i)} One determines \(s_{G}\) explicitly in each case of the classification thanks to \SGARefLink{sga3:xxiv:proposition:3-16-6:item:iii}{(iii)}: for each
irreducible pinned root datum it suffices to compose the symmetry through the origin with the symmetry in the Weyl group
(i.e.~the element \(w_{0}\) of the Weyl group such that \(w_{0}(\Delta) = -\Delta\)). One finds the following results: one
has \(s_{G} = e\) except for \(A_{n}\) (\(n \geqslant 2\)), \(D_{n}\) (\(n\) odd) and \(E_6\), in which case \(s_{G}\) is the unique
non-trivial ``outer automorphism''.

\begin{enumerate}
\def\labelenumi{(\roman{enumi})}
\setcounter{enumi}{1}
\item
  The automorphism \(u_{E}\) is the one used to construct ``the compact real forms'' in the theory of semisimple Lie
  algebras.
\end{enumerate}

\textbf{\SGAInlineRefTarget{sga3:xxiv:remark:3-16-8}Remark 3.16.8.} We defined in \SGARefLink{sga3:xxiv:subsubsection:3-16-1}{3.16.1} an involution in the \(S\)-scheme
\[
Q = Isomint(G_{q-\acute{e}p.}, G)
\]
of quasi-pinnings of \(G\) (cf.~\SGARefLink{sga3:xxiv:3.12.1}{3.12.1}); by transport of structure from \(G_{q-\acute{e}p.}\) to \(G\), one sees at once that
this involution is given by the action of an element of \(ad(G_{q-\acute{e}p.})(S)\): the element \(n_{0}\) defined (\SGARefLink{sga3:xxiv:proposition:3-16-2:item:iii}{3.16.2
(iii)}) by the canonical quasi-pinning of \(G_{q-\acute{e}p.}\).

In the same way, one has defined an involution in the \(S\)-scheme

\[ P = Isom_{S-gr.}(\acute{E}p_{S}(R), G) \]

of \(R\)-pinnings of \(G\) (cf.~\SGARefLink{sga3:xxiv:1.20}{1.20}). Arguing as before, one sees that this involution is given by the action of the
automorphism \(u_{0}\) of \(\acute{E}p_{S}(R)\) defined (\SGARefLink{sga3:xxiv:proposition:3-16-2:item:iv}{3.16.2 (iv)}) by the canonical pinning of \(\acute{E}p_{S}(R)\).

% END INCLUDED SGA3_English_current_buildable_source_20260823/source/inputs/XXIV/tex/components/09_expose_XXIV_proposition313_through_remark3168_en.tex

% BEGIN INCLUDED SGA3_English_current_buildable_source_20260823/source/inputs/XXIV/tex/components/10_expose_XXIV_section4_through_theorem441_setup_en.tex
% Authority coverage: Polo--Gille Exposé XXIV, local pages 25--29.
% Lead-reviewed against the authority PDF; authority images remain explicit native-diagram debt.

\section*{\SGAInlineRefTarget{sga3:xxiv:section:4}4. Isotriviality of reductive groups and of principal bundles under reductive groups}

\subsection*{\SGAInlineRefTarget{sga3:xxiv:subsection:4-1}4.1. Definitions. Isotriviality theorem}

\textbf{\SGAInlineRefTarget{sga3:xxiv:definition:4-1-1}Definition 4.1.1.} \emph{Let \(S\) be a scheme, \(G\) an \(S\)-group scheme, \(P\) a principal homogeneous bundle under \(G\). One
says that \(P\) is} locally isotrivial \emph{(resp.} semi-locally isotrivial) if for every point \(s \in S\) (resp. every
finite set \(F\) of points of \(S\) contained in an affine open set) there exist an open set \(U\) of \(S\) containing \(s\)
(resp. \(F\)) and a finite surjective étale morphism \(S' \to U\) such that \(P_{S'}\) is trivial.

\textbf{\SGAInlineRefTarget{sga3:xxiv:definition:4-1-2}Definition 4.1.2.} \emph{Let \(S\) be a scheme, \(G\) an \(S\)-reductive group. One says that \(G\) is} locally isotrivial
\emph{(resp.} semi-locally isotrivial) if for every point \(s \in S\) (resp. every finite set \(F\) of points of \(S\) contained
in an affine open set) there exist an open set \(U\) of \(S\) containing \(s\) (resp. \(F\)) and a finite surjective étale
morphism \(S' \to U\) such that \(G_{S'}\) is splittable.*

\textbf{\SGAInlineRefTarget{sga3:xxiv:remark:4-1-3}Remark 4.1.3.} (i) The equivalence of categories \SGARefLink{sga3:xxiv:definition:4-1-1}{4.1.1} of \SGARefLink{sga3:xxiv:1.17}{1.17} respects by definition local (resp. semi-local)
isotriviality.

\begin{enumerate}
\def\labelenumi{(\roman{enumi})}
\setcounter{enumi}{1}
\item
  Add to the conditions of \SGARefLink{sga3:xxiv:definition:4-1-1}{4.1.1}: \(G\) locally of finite presentation over \(S\). Then the principal bundle \(P\) (or the
  reductive group \(G\)) is locally isotrivial (resp. semi-locally isotrivial) if and only if for every \(S' \to S\), \(S'\)
  local (resp. semi-local), \(P_{S'}\) is isotrivial (or \(G_{S'}\) isotrivial), that is to say, if there exists \(S'' \to S'\)
  finite surjective étale such that \(P_{S''}\) is trivial (or \(G_{S''}\) split).
\item
  In the case of tori, \SGARefLink{sga3:xxiv:definition:4-1-2}{definition 4.1.2} coincides with that of \SGARefLink{sga3:ix:definition:1:1}{Exp. IX, 1.1}.
\end{enumerate}

\subsubsection*{4.1.4.}

Recall (\SGARefLink{sga3:xxii:section:4.3}{Exp. XXII, 4.3} and \SGARefLink{sga3:xxiv:lemma:6-2}{6.2}) that for every reductive group \(G\), we have introduced the groups \(rad(G)\), \(corad(G)\)
and \(d\acute{e}r(G)\). The groups \(rad(G)\) and \(corad(G)\) are tori, which are split when \(G\) is split; moreover, there
exists an isogeny \(rad(G) \to corad(G)\). The \(S\)-group \(d\acute{e}r(G)\) is semisimple, one has \(G / d\acute{e}r(G) =
corad(G)\); it follows that for every principal homogeneous bundle \(P\) under \(G\), \(P / d\acute{e}r(G)\) is a principal
homogeneous bundle under \(corad(G)\). This said, one has:

\textbf{\SGAInlineRefTarget{sga3:xxiv:theorem:4-1-5}Theorem 4.1.5.} \emph{Let \(S\) be a scheme, \(G\) an \(S\)-reductive group of constant type.}

\emph{\SGAInlineRefTarget{sga3:xxiv:theorem:4-1-5:item:i}(i) The following conditions are equivalent:}

\emph{\SGAInlineRefTarget{sga3:xxiv:theorem:4-1-5:item:a}(a) \(G\) is locally (resp. semi-locally) isotrivial.}

\emph{\SGAInlineRefTarget{sga3:xxiv:theorem:4-1-5:item:b}(b) The torus \(rad(G)\) is.}

\emph{(b') The torus \(corad(G)\) is.}

\emph{\SGAInlineRefTarget{sga3:xxiv:theorem:4-1-5:item:c}(c) The Galois covering of \(S\) associated to \(G\) (\SGARefLink{sga3:xxiv:3.11}{3.11}) is.}

\emph{If \(T\) is a maximal torus of \(G\), these conditions are also equivalent to}

\emph{\SGAInlineRefTarget{sga3:xxiv:theorem:4-1-5:item:d}(d) The torus \(T\) is locally (resp. semi-locally) isotrivial.}

\emph{\SGAInlineRefTarget{sga3:xxiv:theorem:4-1-5:item:ii}(ii) Let \(P\) be a principal homogeneous bundle under \(G\); for \(P\) to be locally (resp. semi-locally) isotrivial, it is
necessary and sufficient that the \(corad(G)\)-principal bundle \(P / d\acute{e}r(G)\) be so.}

\textbf{\SGAInlineRefTarget{sga3:xxiv:corollary:4-1-6}Corollary 4.1.6.} \emph{The conditions of \SGARefLink{sga3:xxiv:theorem:4-1-5:item:i}{(i)} are satisfied when \(G\) is semisimple or when \(S\) is locally noetherian and
normal (or more generally geometrically unibranch). The conditions of \SGARefLink{sga3:xxiv:theorem:4-1-5:item:ii}{(ii)} are satisfied when \(G\) is locally (resp.
semi-locally) isotrivial.}

For \SGARefLink{sga3:xxiv:theorem:4-1-5:item:i}{(i)}, the first assertion is trivial from \SGARefLink{sga3:xxiv:theorem:4-1-5:item:b}{(b)}, the second follows from \SGARefLink{sga3:xxiv:theorem:4-1-5:item:c}{(c)} and \SGARefLink{sga3:x:corollary:5-14}{Exp. X, 5.14} and 5.15. For \SGARefLink{sga3:xxiv:theorem:4-1-5:item:ii}{(ii)}, it
suffices to note that by virtue of \SGARefLink{sga3:viii:corollary:4:5}{theorem 90}, a principal bundle under a split torus is semi-locally isotrivial.

\subsection*{\SGAInlineRefTarget{sga3:xxiv:subsection:4-2}4.2. Proof: the semisimple case}

Let us first prove, for later reference:

\textbf{\SGAInlineRefTarget{sga3:xxiv:proposition:4-2-1}Proposition 4.2.1.} \emph{Let \(S\) be a scheme, \(G\) an \(S\)-reductive group, \(T\) a maximal torus of \(G\) (resp. \(B\) a Borel
subgroup, resp. \(B \supset T\) a Killing couple of \(G\)), \(P\) a principal homogeneous bundle under \(G\), \(G'\) the twisted
form of \(G\) associated to \(P\) (via the morphism \(int : G \to \operatorname{Aut}_{S-gr.}(G)\)). One has canonical
isomorphisms}

\[
  P/\operatorname{Norm}_G(T)\xrightarrow{\sim}\operatorname{Tor}(G'),
  \qquad
  P/B\xrightarrow{\sim}\operatorname{Bor}(G'),
  \qquad
  P/T\xrightarrow{\sim}\operatorname{Kil}(G').
\]

By construction, \(G'\) is the quotient of \(P \times_{S} G\) by a certain action of \(G\) (\(((p, g')) g = (pg, g^{-1} g'
g)\)); one therefore has a morphism \(P \times_{S} G \to G'\), that is, a morphism

\[ P \to \operatorname{Hom}_{S}(G, G'), \]

which, as one sees at once, factors through a morphism

\[
  f:P\longrightarrow\operatorname{Isom}_{S\text{-}\mathrm{gr.}}(G,G').
\]

(to verify this assertion, one may assume \(G = P\), in which case one has \(G' = G\), \(f = int\)). The set-theoretic map \(p
\mapsto f(p)(T)\) defines a morphism

\[ P \to Tor(G'). \]

To verify that this morphism induces an isomorphism \(P / Norm_{G}(T) \xrightarrow{\sim} Tor(G')\) as announced, one may
again assume \(P = G\) in which case one is reduced to \SGARefLink{sga3:xxii:cor:5.8.3:item:iii}{Exp. XXII, 5.8.3 (iii)}. (In fact, \emph{loc. cit.} should be replaced by
the statement above.) One argues similarly for \(\operatorname{Bor}\) and \(\operatorname{Kil}\).

\textbf{\SGAInlineRefTarget{sga3:xxiv:proposition:4-2-2}Proposition 4.2.2.} \emph{Let \(S\) be a semi-local scheme, \(G\) an \(S\)-semisimple group of constant type.}

\emph{\SGAInlineRefTarget{sga3:xxiv:proposition:4-2-2:item:i}(i) \(G\) is isotrivial.}

\emph{(ii) Every principal bundle under \(G\) is isotrivial.}

Let us prove \SGARefLink{sga3:xxiv:proposition:4-2-2:item:i}{(i)}. Up to making a finite surjective étale extension of the base, one may, by \SGARefLink{sga3:xxiv:3.9.2}{3.9.2}, assume \(G\)
quasi-split. But then \(G\) is isotrivial by construction (\SGARefLink{sga3:xxiv:3.11}{3.11}, the group \(E\) being finite). To prove \SGARefLink{sga3:xxiv:3.11:item:ii}{(ii)}, one may, by
\SGARefLink{sga3:xxiv:3.11:item:i}{(i)}, assume \(G\) split; one is then reduced to:

\textbf{\SGAInlineRefTarget{sga3:xxiv:lemma:4-2-3}Lemma 4.2.3.} \emph{Let \(S\) be a semi-local scheme. Every principal bundle under a split reductive group is isotrivial.}

Indeed, with the notation of \SGARefLink{sga3:xxiv:proposition:4-2-1}{4.2.1}, where \(B \supset T\) denotes the canonical Killing couple of the split \(G\), the
\(S\)-scheme \(Kil(G')\) possesses a section, after finite surjective étale extension of the base, by \SGARefLink{sga3:xxiv:3.9.2}{3.9.2}. One may
therefore reduce the structure group of \(G\) to \(T\); but \(T\) is split, so \(H^{1}(S, T) = 0\) (theorem 90).

\textbf{\SGAInlineRefTarget{sga3:xxiv:corollary:4-2-4}Corollary 4.2.4.} \emph{Let \(S\) be a scheme and}

\[ 1 \to G \to G' \to G'' \to 1 \]

\emph{an exact sequence of \(S\)-group schemes (for (fpqc)), \(G\) being semisimple of constant type. Let \(P\) be a principal
homogeneous bundle under \(G'\), suppose the sheaf\footnote{N.D.E.: We have replaced ``bundle'' by ``sheaf'', and then \(G'\) by \(G''\).} \(P / G\) associated representable (for example \(G''\)
affine over \(S\)). For \(P\) to be locally isotrivial (resp. semi-locally isotrivial), it is necessary and sufficient that
\(P / G\) be so (as bundle under \(G''\)).}

If \(P\) is trivial, \(P / G\) is too, which shows that the condition is necessary. Conversely, suppose \(S\) local (resp.
semi-local) and \(P / G\) isotrivial, so trivialized by a finite surjective étale extension \(S'\) of \(S\). Extending the
base to \(S'\), one may reduce the structure group of \(P_{S'}\) to \(G_{S'}\). But \(S'\) is still semi-local and \(G_{S'}\)
semisimple of constant type, so the obtained bundle is isotrivial (\SGARefLink{sga3:xxiv:proposition:4-2-2}{4.2.2}).

\subsection*{\SGAInlineRefTarget{sga3:xxiv:subsection:4-3}4.3. Proof: general case}

Let us first note that \SGARefLink{sga3:xxiv:theorem:4-1-5:item:ii}{4.1.5 (ii)} follows at once from the application of \SGARefLink{sga3:xxiv:corollary:4-2-4}{4.2.4} to the exact sequence

\[ 1 \to d\acute{e}r(G) \to G \to corad(G) \to 1. \]

Let us therefore prove \SGARefLink{sga3:xxiv:proposition:4-2-2:item:i}{(i)}. If \(G\) is split, \(rad(G)\) and \(corad(G)\) are split, as is \(rev(G)\); so \SGARefLink{sga3:xxiv:theorem:4-1-5:item:a}{(a)} implies \SGARefLink{sga3:xxiv:theorem:4-1-5:item:b}{(b)},
(b'), and \SGARefLink{sga3:xxiv:theorem:4-1-5:item:c}{(c)}.

\subsubsection*{4.3.1.}

One has \SGARefLink{sga3:xxiv:theorem:4-1-5:item:c}{(c)} \(\Rightarrow\) \SGARefLink{sga3:xxiv:theorem:4-1-5:item:a}{(a)}. Let \(R\) be the type of \(G\); one has an exact sequence

\[ 1 \to ad(\acute{E}p_{S}(R)) \to A_{S}(R) \to E_{S} \to 1. \]

Applying \SGARefLink{sga3:xxiv:corollary:4-2-4}{4.2.4} to the bundle \(P = Isom_{S-gr.}(\acute{E}p_{S}(R), G)\) and to the associated bundle \(rev(G) = P /
ad(\acute{E}p_{S}(R))\), one has \SGARefLink{sga3:xxiv:theorem:4-1-5:item:c}{(c)} \(\Rightarrow\) \SGARefLink{sga3:xxiv:theorem:4-1-5:item:a}{(a)}.

\subsubsection*{4.3.2.}

One has \SGARefLink{sga3:xxiv:theorem:4-1-5:item:b}{(b)} \(\Rightarrow\) \SGARefLink{sga3:xxiv:theorem:4-1-5:item:a}{(a)}. It suffices to prove that if \(rad(G)\) is split, \(G\) is semi-locally isotrivial. Now let \(G_{0} =
\acute{E}p_{S}(R)\); consider the category of pairs \((G',f)\), where \(G'\) is a form of \(G_0\) and \(f\) is an isomorphism of
\(rad(G_{0})\) onto \(rad(G)\).

One knows (\SGARefLink{sga3:xxiv:2.16}{2.16}) that this category is equivalent to the category of principal homogeneous bundles under a certain
\(S\)-group \(H\) extension of a finite twisted constant group by a semisimple group. It suffices to prove that every
principal bundle under \(H\) is semi-locally isotrivial; this follows at once from \SGARefLink{sga3:xxiv:corollary:4-2-4}{4.2.4}.

\subsubsection*{4.3.3.}

One has (b') \(\Rightarrow\) \SGARefLink{sga3:xxiv:theorem:4-1-5:item:a}{(a)}. One can argue as previously (it will moreover be the same group \(H\) that arises). One can also see
that \SGARefLink{sga3:xxiv:theorem:4-1-5:item:b}{(b)} and (b') are equivalent: a torus isogenous to a locally split torus is also locally split (cf.~\SGARefLink{sga3:ix:proposition:2:11:item:iii:01}{Exp. IX, 2.11
(iii)}).

\subsubsection*{4.3.4.}

One has \SGARefLink{sga3:xxiv:theorem:4-1-5:item:d}{(d)} \(\Rightarrow\) \SGARefLink{sga3:xxiv:theorem:4-1-5:item:a}{(a)}. It suffices to prove that a group of constant type possessing a split maximal torus is semi-locally
isotrivial; this follows at once from \SGARefLink{sga3:xxiv:2.14}{2.14}.

\subsubsection*{4.3.5.}

One has \SGARefLink{sga3:xxiv:theorem:4-1-5:item:a}{(a)} \(\Rightarrow\) \SGARefLink{sga3:xxiv:theorem:4-1-5:item:d}{(d)}. It suffices to prove that a maximal torus of a split group is semi-locally isotrivial. More
precisely:

\textbf{\SGAInlineRefTarget{sga3:xxiv:lemma:4-3-6}Lemma 4.3.6.} \emph{Let \(S\) be a semi-local scheme, \(G\) an \(S\)-reductive group, \(T_0\) a split maximal torus of \(G\), \(W_{0}
= Norm_{G}(T_{0}) / T_{0}\) (this is a locally constant \(S\)-group, constant if \(G\) is of constant type, by \SGARefLink{sga3:xxiv:2.14}{2.14}), \(T\) a
maximal torus of \(G\).}

\emph{There exists a morphism \(S' \to S\) which is principal homogeneous under \(W_0\) (hence étale finite and surjective, and
even principal Galois if \(G\) is of constant type) such that \(T_{S'}\) is conjugate to \((T_{0})_{S'}\) by an element of
\(G(S')\) (and so in particular split).}

Indeed, one knows that \(Transp_{G}(T_{0}, T)\) is a principal homogeneous bundle under \(Norm_{G}(T_{0})\) (cf.~for example
\SGARefLink{sga3:xi:corollary:5:4:bis}{Exp. XI, 5.4 bis}). Set \(S' = Transp_{G}(T_{0}, T) / T_{0}\); this is a principal homogeneous bundle under \(W_0\).
Extending the base from \(S\) to \(S'\), one can reduce the structure group of \(Transp_{G}(T_{0}, T)\) to \(T_0\). Now \(S'\) is
semi-local and \(T_0\) split, so \(Transp_{G}(T_{0}, T)\) possesses a section over \(S'\).\footnote{N.D.E.: by ``theorem 90''.} QED

\subsection*{\SGAInlineRefTarget{sga3:xxiv:subsection:4-4}4.4. Use of the existence of maximal tori}

Using Grothendieck's existence theorem for maximal tori (\SGARefLink{sga3:xiv:3.20}{Exp. XIV, 3.20}), one can considerably sharpen the preceding
results. Let us state at once:

\textbf{\SGAInlineRefTarget{sga3:xxiv:theorem:4-4-1}Theorem 4.4.1.} \emph{Let \(S\) be a semi-local scheme, \(R\) a pinned reduced root datum, \(W\) its Weyl group, \(E\) the group
of its automorphisms. (Recall that \(E\) operates naturally on \(W\) and that the semi-direct product \(A = W \cdot E\) is
identified with the group of automorphisms of \(R\) non-pinned, cf.~\SGARefLink{sga3:XXI:6.7.2}{Exp. XXI, 6.7.2}).}

\emph{(i) Every principal homogeneous bundle under \(\acute{E}p_{S}(R)\) is trivialized by a principal Galois covering \(S' \to
S\) of group \(W\).}

\emph{(ii) Let \(G\) be an \(S\)-reductive group of type \(R\); let \(rev(G) = Isomext(\acute{E}p_{S}(R), G)\) be the associated
Galois covering of \(S\) with group \(E\). Let \(W_0\) be the form of \(W_S\) associated to \(rev(G)\). There exists a morphism
\(S' \to S\), which is a principal homogeneous bundle under \(W_0\), such that \(G_{S'}\) is quasi-splittable (i.e.~possesses
a Borel subgroup).}

\emph{\SGAInlineRefTarget{sga3:xxiv:theorem:4-4-1:item:iii}(iii) Every \(S\)-reductive group of type \(R\) is split by a principal Galois covering \(\tilde{S} \to S\) with group \(A\).}

Let us first state:

\textbf{\SGAInlineRefTarget{sga3:xxiv:proposition:4-4-2}Proposition 4.4.2.} \emph{Let \(S\) be a scheme, \(G\) an \(S\)-reductive group, \(T\) a maximal torus of \(G\), \(P\) a principal
homogeneous bundle under \(G\), \(G'\) the twisted form of \(G\) associated to \(P\) (one then has, cf.~\SGARefLink{sga3:xxiv:proposition:4-2-1}{4.2.1}, a canonical
isomorphism \(P / Norm_{G}(T) \xrightarrow{\sim} Tor(G')\)). Let \(T'\) be a maximal torus of \(G'\), defining a composite
morphism}

\[
  S\longrightarrow\operatorname{Tor}(G')
  \xrightarrow{\sim}P/\operatorname{Norm}_G(T).
\]

\emph{Consider the canonical morphisms}

\[
  P\longrightarrow P/T\longrightarrow P/\operatorname{Norm}_G(T).
\]

\emph{and take their inverse images by the preceding morphism:}

\[
  % BEGIN INCLUDED SGA3_English_current_buildable_source_20260823/source/inputs/XXIV/native_diagrams/exp24/003_inverse_image_grid.tex
\begin{tikzcd}[column sep=large,row sep=large]
P \arrow[r]
&
P/T \arrow[r]
&
P/\operatorname{Norm}_G(T)
\\
H \arrow[u] \arrow[r]
&
S' \arrow[u] \arrow[r]
&
S \arrow[u]
\end{tikzcd}

% END INCLUDED SGA3_English_current_buildable_source_20260823/source/inputs/XXIV/native_diagrams/exp24/003_inverse_image_grid.tex

\]

\emph{Then \(S'\) (resp. \(H\)) is a principal homogeneous bundle over \(S\) (resp. \(S'\)) under \(W_{G}(T)\) (resp. \(T_{S'}\)).
Moreover, if one makes \(W_{G}(T)\) operate on \(T\) in the obvious way, the bundle associated to \(S'\) is isomorphic to
\(T'\).}

The first two assertions are trivial, the last is proved like the corresponding assertion of \SGARefLink{sga3:xxiv:2.6}{2.6}.

\textbf{\SGAInlineRefTarget{sga3:xxiv:corollary:4-4-3}Corollary 4.4.3.} \emph{Let \(S\) be a semi-local scheme, \(G\) an \(S\)-reductive group, \(T\) a maximal torus of \(G\). Suppose
one of the two following conditions is satisfied:}

\emph{\SGAInlineRefTarget{sga3:xxiv:corollary:4-4-3:item:i}(i) \(T\) is split.}

\emph{\SGAInlineRefTarget{sga3:xxiv:corollary:4-4-3:item:ii}(ii) \(T\) is contained in a Borel subgroup of \(G\), and \(G\) is either adjoint, or simply connected.}

\emph{Let moreover \(P\) be a principal homogeneous bundle under \(G\). There exists an \(S\)-scheme \(S'\), which is a principal
homogeneous bundle under \(W_{G}(T)\), such that \(P_{S'}\) is trivial.}

Indeed, if \(G'\) is the form of \(G\) associated to \(P\), then \(G'\) possesses a maximal torus \(T'\) (\SGARefLink{sga3:xiv:3.20}{Exp. XIV, 3.20}).
Resuming the notation of the preceding proposition, one sees that \(H^{1}(S', T_{S'}) = 0\) (by theorem 90 for \SGARefLink{sga3:xxiv:corollary:4-4-3:item:i}{(i)}, by
\SGARefLink{sga3:xxiv:corollary:3-14}{3.14} for \SGARefLink{sga3:xxiv:corollary:4-4-3:item:ii}{(ii)}). The morphism \(H \to S'\) possesses a section, so \(P_{S'}\) also possesses a section over \(S'\). QED

Let us now prove the theorem. Assertion \SGARefLink{sga3:xxiv:corollary:4-4-3:item:i}{(i)} is a particular case of the preceding corollary (take \(G =
\acute{E}p_{S}(R)\), endowed with its canonical split torus). Let us prove \SGARefLink{sga3:xxiv:corollary:4-4-3:item:ii}{(ii)}. One knows (\SGARefLink{sga3:xxiv:3.12}{3.12}) that \(G\) is an inner
twisted form of

\[ G_{0} = Q-\acute{E}p_{rev(G)/S}(R). \]

If \(T_0\) is the canonical maximal torus of \(G_0\), \(W_{G_{0}}(T_{0})\) is the group \(W_0\) described in the statement. The
form \(G\) of \(G_0\) corresponds to a principal homogeneous bundle \(P\) under \(ad(G_{0})\) (\(P = Isomint(G_{0}, G)\)). The
group \(W_{ad(G_{0})}(T^{ad}_{0})\) is canonically isomorphic to \(W_0\), and one obtains the wanted result by applying
\SGARefLink{sga3:xxiv:corollary:4-4-3}{4.4.3} to the situation \((ad(G_{0}), T^{ad}_{0}, P)\), hypothesis \SGARefLink{sga3:xxiv:corollary:4-4-3:item:ii}{(ii)} of \SGARefLink{sga3:xxiv:corollary:4-4-3}{4.4.3} being well verified. Let us finally prove
\SGARefLink{sga3:xxiv:theorem:4-4-1:item:iii}{(iii)}. Resume the notation of \SGARefLink{sga3:xxiv:corollary:4-4-3:item:ii}{(ii)}; one has a diagram

\[
  % BEGIN INCLUDED SGA3_English_current_buildable_source_20260823/source/inputs/XXIV/native_diagrams/exp24/004_rev_base_triangle.tex
\begin{tikzcd}[column sep=large,row sep=large]
&
\operatorname{rev}(G) \arrow[d,"E_S"']
\\
S'
  \arrow[r,"W_0"']
&
S
\end{tikzcd}

% END INCLUDED SGA3_English_current_buildable_source_20260823/source/inputs/XXIV/native_diagrams/exp24/004_rev_base_triangle.tex

\]

One knows that \(G_{S'}\) is isomorphic to \((G_{0})_{S'}\) and that \((G_{0})_{rev(G)}\) is splittable. If one sets
\(\tilde{S} = S' \times_{S} rev(G)\), \(G_{\tilde{S}}\) is splittable, and it only remains to verify that \(\tilde{S}\) is
indeed a principal Galois covering of \(S\) with group \(A\), which follows from the following more general lemma (naturally
valid in any site):

% END INCLUDED SGA3_English_current_buildable_source_20260823/source/inputs/XXIV/tex/components/10_expose_XXIV_section4_through_theorem441_setup_en.tex

% BEGIN INCLUDED SGA3_English_current_buildable_source_20260823/source/inputs/XXIV/tex/components/11_expose_XXIV_lemma444_through_section54_en.tex
% Authority coverage: Polo--Gille Exposé XXIV, local pages 29--32.
% Lead-reviewed against the authority PDF; the next source cursor is Proposition 5.5 on local page 32.


\textbf{\SGAInlineRefTarget{sga3:xxiv:lemma:4-4-4}Lemma 4.4.4.} \emph{Let \(S\) be a scheme, \(G\) and \(H\) two \(S\)-group schemes, \(G \to \operatorname{Aut}_{S-gr.}(H)\) an
action of \(G\) on \(H\), \(E\) a \(G\)-principal homogeneous bundle, \(F\) an \(H_0\)-principal homogeneous bundle, where \(H_0\) is
the form of \(H\) associated to \(E\). Then \(E \times_{S} F\) is endowed naturally with a structure of principal homogeneous
bundle under the semi-direct product \(H \cdot G\).}

Write \((e, g) \mapsto eg\) (resp. \((f, u) \mapsto f u\)) for the action (on the right) of \(G\) on \(E\) (resp. of \(H_0\) on
\(F\)). Write

\[
  p:E\times_S H\longrightarrow H_0.
\]

for the canonical projection (\(H_0\) is by definition the quotient of \(E \times_{S} H\) by \(G\) acting on it by the formula
\((e, h) g = (eg, g^{-1} h g)\)). Consider the morphism

\[
  r:E\times_S F\times_S H\times_S G\longrightarrow E\times_S F.
\]

defined set-theoretically by

\[
  r(e,f,h,g)=\bigl(eg,f\cdot p(e,h)\bigr).
\]

The morphism \(r\) indeed defines an action of the semi-direct product \(H \cdot G\) on \(E \times_{S} F\). Indeed, one has
set-theoretically

\[
  r\bigl(r(e,f,h,g),h',g'\bigr)
  =\bigl(egg',f\,p(e,h)\,p(eg,h')\bigr).
\]

but

\[
  p(e,h)p(eg,h')
  =p(e,h)p(e,gh'g^{-1})
  =p(e,hgh'g^{-1}).
\]

whence

\[
  r\bigl(r(e,f,h,g),h',g'\bigr)
  =r(e,f,hgh'g^{-1},gg').
\]

which had to be proved.

To prove now that this law is indeed a law of principal homogeneous bundle, one may assume that \(E\) and \(F\) are trivial,
in which case one sees at once that \(E \times_{S} F\) is also a trivial bundle.

\subsection*{\SGAInlineRefTarget{sga3:xxiv:subsection:4-5}4.5. Isotriviality of maximal tori of semisimple groups}
\footnote{\textit{Editors' note.} This no.~\SGARefLink{sga3:xxiv:subsection:4-5}{4.5} has been added.}

Let us bring out here the following result, contained implicitly in \SGARefLink{sga3:xxiv:2.6}{2.6} (cf.~the N.D.E. (13)) and which was mentioned at
the end of \SGARefLink{sga3:ix:comments:1:2}{Exp. IX 1.2}. Recall (\SGARefLink{sga3:XXIII:5.11}{XXIII 5.11}) that, for every scheme \(S\) and every reduced root datum \(R\), one writes
\(\acute{E}p_{S}(R)\) for the unique pinned \(S\)-group of type \(R\) (which exists by Exp. XXV) and \(T_{S}(R)\) its canonical
maximal torus.

\textbf{\SGAInlineRefTarget{sga3:xxiv:proposition:4-5-1}Proposition 4.5.1.} \emph{Let \(G\) be an \(S\)-semisimple group of constant type \(R\); assume that \(G\) possesses a maximal
torus \(T\) (which is the case if \(S\) is semi-local, by \SGARefLink{sga3:xiv:3.20}{XIV 3.20}).}

\emph{\SGAInlineRefTarget{sga3:xxiv:proposition:4-5-1:item:i}(i) Then \(P = Isom_{S-gr.}(\acute{E}p_{S}(R), T_{S}(R); G, T) / T_{S}(R)_{ad}\) is a principal \(S\)-bundle under the
finite group \(\operatorname{Aut}(R)\), and \(T_P\) is a split maximal torus of \(G_P\). Consequently, \(T\) is isotrivial.}

\emph{\SGAInlineRefTarget{sga3:xxiv:proposition:4-5-1:item:ii}(ii) Moreover, for every finite set \(F\) of points of \(P\) contained in an affine open set, there exists an open set \(U\)
of \(P\) containing \(F\) such that \(G_U\) is split.}

\emph{Proof.} The exact sequence

\[
  1\longrightarrow T_S(\mathscr R)_{\mathrm{ad}}
  \longrightarrow
  \underline{\operatorname{Aut}}_{S\text{-}\mathrm{gr.}}
  \bigl(\operatorname{\acute{E}p}_S(\mathscr R),T_S(\mathscr R)\bigr)
  \longrightarrow
  \underline{\operatorname{Aut}}_{S\text{-}\mathrm{gr.}}\bigl(T_S(\mathscr R)\bigr).
\]

induces for every \(S' \to S\) a morphism of bundles \(P_{S'} \to Isom_{S'-gr.}(T_{S'}(R), T_{S'})\), so if \(P_{S'}\)
possesses a section, then \(T_{S'}(R)\) and \(T_{S'}\) are isomorphic. This is in particular the case for \(S' = P\).

Note moreover that \(\operatorname{Aut}_{S-gr.}(\acute{E}p_{S}(R), T_{S}(R)) / T_{S}(R)_{ad}\) is none other than the
constant \(S\)-group \(\operatorname{Aut}(R)_{S}\), and that \(P\) is identified with the \(\operatorname{Aut}(R)_{S}\)-torsor
\(Isom(R, \tilde{R}(G, T))\), where \(\tilde{R}(G, T)\) denotes the twisted root datum associated to \((G, T)\) (cf.~\SGARefLink{sga3:xxii:prop:1.10}{XXII
1.10}).

Then \SGARefLink{sga3:xxiv:proposition:4-5-1:item:ii}{(ii)} follows from \SGARefLink{sga3:xxiv:proposition:4-5-1:item:i}{(i)} and from \SGARefLink{sga3:xxiv:2.14}{2.14}, taking into account the fact that the \(g_{\alpha}\) (\(\alpha\) ranging over the
roots of \(R\)) are free \(\mathcal O_U\)-modules of rank 1 on a sufficiently small affine open set \(U\) containing \(F\).

\section*{\SGAInlineRefTarget{sga3:xxiv:section:5}5. Canonical decomposition of an adjoint or simply connected group}

In this section, we shall use the results of no 1 to generalize to the case of schemes a classical decomposition of
adjoint (resp. simply connected) groups. So as not to overburden the exposition indefinitely, the proofs are sketched
and the details are left to the reader; in fact it is always a matter of absolutely standard proofs in the theory of
principal bundles: reduction of the structure group, twisting, etc.

\subsection*{\SGAInlineRefTarget{sga3:xxiv:subsection:5-1}5.1.}

Recall (\SGARefLink{sga3:XXI:7.4}{Exp. XXI, 7.4}) that a Dynkin diagram is a disjoint union of its connected components, which are Dynkin diagrams.
Moreover, every non-empty connected Dynkin diagram corresponding to a root datum is isomorphic to one of the standard
diagrams (\(A_{n}\), \(B_{n}\), \ldots, \(G_2\)) which were exhibited in \SGARefLink{sga3:XXI:7.4.6}{Exp. XXI, 7.4.6}. In the sequel, we shall be interested
only in Dynkin diagrams whose connected components are of one of the preceding types. Let \(T\) be the set of these
standard diagrams. For every Dynkin diagram \(D\), let \(n(t) = n_{D}(t)\) be the number of connected components of \(D\)
isomorphic to \(t\), where \(t \in T\). The \emph{type} of \(D\) is by definition \(\sum_{t \in T} n_{D}(t) \cdot t\).

A Dynkin diagram of type \(t\) is said to be \emph{simple of type \(t\)}, a Dynkin diagram of type \(n \cdot t\) is said to be
\emph{isotypic of type \(t\)}. Let \(D_0\) be the set of connected components of \(D\) and let

\[ a : D_{0} \to T \]

be the obvious map. The type of \(D\) is none other than \(\sum_{x \in D_{0}} a(x)\).

\subsection*{\SGAInlineRefTarget{sga3:xxiv:subsection:5-2}5.2.}

Let \(S\) be a scheme, \(D\) an \(S\)-Dynkin scheme (verifying the restrictive condition stated above). The cokernel of the
pair of morphisms \(L \Rightarrow D\) (\(L =\)\footnote{N.D.E.: We have replaced \(R\) by \(L\), as in \SGARefLink{sga3:xxiv:3.2}{3.2}.} scheme of bonds of \(D\)) is written \(D_0\). This is the ``scheme
of connected components'' of \(D\) (it

exists trivially when \(D\) is constant; the general case is deduced from this by descent); it is a finite twisted
constant \(S\)-scheme. One has a canonical morphism

\[ a : D_{0} \to T_{S}. \]

For \(t \in T\), set \(a^{-1}(t) = D_{0, t}\); this is a subscheme of \(D_0\), whose inverse image in \(D\), written \(D_{t}\), is
the \emph{isotypic component of type \(t\)} of the Dynkin scheme \(D\). Each \(D_{t}\) is a subscheme of \(D\), and one has

\[
  D=\coprod_{t\in T}D_t.
\]

Note that the degree of \(D_{0, t}\) at \(s \in S\) is \(n(s, t)\), if the type of \(D\) at \(s\) is \(\sum_{t} n(s, t) \cdot t\).

\subsection*{\SGAInlineRefTarget{sga3:xxiv:subsection:5-3}5.3.}

In what follows, we shall consider only semisimple adjoint (resp. simply connected) groups. To simplify the language, we
shall state the results for adjoint groups; all statements will remain valid if one substitutes everywhere ``simply
connected'' for ``adjoint''.

Recall that an adjoint reduced root datum is determined up to isomorphism by the type of its Dynkin diagrams. We shall
therefore say that an adjoint root datum \(R\) (resp. a semisimple adjoint group \(G\)) is \emph{of type \(\sum n(t) \cdot t\)} if
its Dynkin diagrams are (resp. if its type is given by an adjoint root datum of type \(\sum n(t) \cdot t\)). We shall say
that \(R\) or \(G\) is \emph{simple of type \(t\)} (resp. \emph{isotypic of type \(t\)}) if its type is \(t\) (resp. \(n \cdot t\), \(n > 0\)).

If \(G\) is a semisimple adjoint group, we shall use the symbols \(Dyn_{0}(G)\) and \(Dyn_{0, t}(G)\) in the sense defined in
\SGARefLink{sga3:xxiv:subsection:5-2}{5.2}.

\subsection*{\SGAInlineRefTarget{sga3:xxiv:subsection:5-4}5.4.}

Let \(t_i\), \(i=1,2,\ldots,n\), be distinct elements of \(T\), and let \(R_i\) be a pinned adjoint root datum isotypic of
type \(t_{i}\). Consider the product \(R = R_{1} \times \cdot\cdot\cdot \times R_{n}\) (\SGARefLink{sga3:XXI:6.4.1}{Exp. XXI, 6.4.1}). Let \(S\) be a
scheme such that the different \(\acute{E}p_{S}(R_{i})\) exist (cf.~Exp. XXV). One verifies at once that there exists a
canonical isomorphism

\[
  \tag{*}\label{sga3:xxiv:subsection:5-4:equation:star}
  \operatorname{\acute{E}p}_S(\mathscr R)
  =
  \operatorname{\acute{E}p}_S(\mathscr R_1)
  \times_S\cdots\times_S
  \operatorname{\acute{E}p}_S(\mathscr R_n).
\]



% END INCLUDED SGA3_English_current_buildable_source_20260823/source/inputs/XXIV/tex/components/11_expose_XXIV_lemma444_through_section54_en.tex

% BEGIN INCLUDED SGA3_English_current_buildable_source_20260823/source/inputs/XXIV/tex/components/12_expose_XXIV_proposition55_through_remark631_en.tex
% Authority coverage: Polo--Gille Exposé XXIV, local pages 32--37.
% Reconstructed manually from the authority PDF.  The Reinhold Markdown
% lineage was used only as a declared comparison translation.

\par\medskip\noindent\SGAInlineRefTarget{sga3:xxiv:proposition:5-5}\textbf{Proposition 5.5.---}\enspace
\emph{Let \(S\) be a scheme and \(G\) an adjoint semisimple
\(S\)-group.  There is a unique decomposition}
\[
  G\simeq\prod_{t\in\mathrm T}G_t,
\]
\emph{where \(G_t\) is an adjoint semisimple \(S\)-group, isotypic of
type \(t\).  Moreover, the preceding decomposition induces
isomorphisms}
\[
  \operatorname{Dyn}_t(G)\simeq\operatorname{Dyn}(G_t),
  \qquad
  \underline{\operatorname{Aut}}_{S\text{-}\mathrm{gr.}}(G)
  \simeq
  \prod_{t\in\mathrm T}
  \underline{\operatorname{Aut}}_{S\text{-}\mathrm{gr.}}(G_t).
\]

This was proved above when \(G\) is split.  In the general case one may
suppose that \(G\) has constant type \(\mathscr R\).  Using the
preceding decomposition of \(A_S(\mathscr R)\) and \SGARefLink{sga3:xxiv:1.17}{1.17}, one obtains
the desired decomposition of \(G\).  The other assertions then follow
by descent.

\par\medskip\noindent\SGAInlineRefTarget{sga3:xxiv:remark:5-6}\textbf{Remark 5.6.---}\enspace
More generally, if \(G\) and \(H\) are two adjoint semisimple groups,
there are canonical isomorphisms in the following commutative diagram:
\[
  % BEGIN INCLUDED SGA3_English_current_buildable_source_20260823/source/inputs/XXIV/native_diagrams/exp24/012_isomorphism_dynkin_comparison.tex
\begin{tikzcd}[column sep=huge,row sep=large]
\underline{\operatorname{Isom}}_{S\text{-}\mathrm{gr.}}(G,H)
  \arrow[r,"\sim"] \arrow[d]
&
\displaystyle\prod_{t\in\mathrm T}
\underline{\operatorname{Isom}}_{S\text{-}\mathrm{gr.}}(G_t,H_t)
  \arrow[d]
\\
\underline{\operatorname{Isomext}}(G,H)
  \arrow[r,"\sim"] \arrow[d,"\wr"']
&
\displaystyle\prod_{t\in\mathrm T}
\underline{\operatorname{Isomext}}(G_t,H_t)
  \arrow[d,"\wr"]
\\
\underline{\operatorname{Isom}}_{\mathrm{Dyn}}
  \bigl(\operatorname{Dyn}(G),\operatorname{Dyn}(H)\bigr)
  \arrow[r,"\sim"]
&
\displaystyle\prod_{t\in\mathrm T}
\underline{\operatorname{Isom}}_{\mathrm{Dyn}}
  \bigl(\operatorname{Dyn}(G_t),\operatorname{Dyn}(H_t)\bigr).
\end{tikzcd}

% END INCLUDED SGA3_English_current_buildable_source_20260823/source/inputs/XXIV/native_diagrams/exp24/012_isomorphism_dynkin_comparison.tex

\]

\par\medskip\noindent\SGAInlineRefTarget{sga3:xxiv:remark:5-7}\textbf{Remark 5.7.---}\enspace
One can give the following intrinsic characterization of \(G_t\),
which we state without proof: it is the largest reductive subgroup of
\(G\) that is invariant (and indeed characteristic) and isotypic of
type \(t\).

The preceding proposition reduces the study of adjoint semisimple
groups to that of isotypic adjoint semisimple groups.  We now undertake
this study.

\subsection*{\SGAInlineRefTarget{sga3:xxiv:subsection:5-8}5.8.}

Let \(\mathscr R\) be a pinned reduced adjoint simple root datum of type
\(t\), let \(I\) be a nonempty finite set, and let
\(\mathscr R^I\) be the product root datum of copies
\(\mathscr R_i\) of \(\mathscr R\), for \(i\in I\).  Denote by \(E\)
the group of automorphisms of \(\mathscr R\), identified with the group
of automorphisms of its Dynkin diagram \(D\).  The Dynkin diagram of
\(\mathscr R^I\) is identified with \(D^I\), and hence there is an
exact sequence
\[
  1\longrightarrow\operatorname{Aut}(D)^I
  \longrightarrow\operatorname{Aut}(D^I)
  \longrightarrow\operatorname{Aut}(I)
  \longrightarrow1,
\]
where \(\operatorname{Aut}(I)\) is the group of permutations of \(I\).
It follows that the canonical isomorphism
\[
  \operatorname{\acute Ep}_S(\mathscr R)^I
  \simeq
  \operatorname{\acute Ep}_S(\mathscr R^I)
\]
induces an exact sequence
\[
  1\longrightarrow A_S(\mathscr R)^I
  \longrightarrow A_S(\mathscr R^I)
  \longrightarrow\underline{\operatorname{Aut}}(I_S)
  \longrightarrow1.
\]
The last group is the constant \(S\)-group associated with
\(\operatorname{Aut}(I)\).  Notice also that \(I\) is canonically
identified with the set of connected components of \(D^I\).

Let \(G\) be an \(S\)-semisimple group of type \(\mathscr R^I\), thus
defining (cf.\ \SGARefLink{sga3:xxiv:1.17}{1.17}) a principal homogeneous bundle \(P\) under
\(A_S(\mathscr R^I)\).  The definitions by descent of
\(\operatorname{Dyn}(G)\) (\SGARefLink{sga3:xxiv:3.7}{3.7}) and of
\(\operatorname{Dyn}_0(G)\) (\SGARefLink{sga3:xxiv:subsection:5-2}{5.2}) show that the bundle associated with
\(P\) by
\[
  A_S(\mathscr R^I)\longrightarrow
  \underline{\operatorname{Aut}}(I_S)
\]
is precisely
\(\underline{\operatorname{Isom}}_S(I_S,\operatorname{Dyn}_0(G))\),
corresponding to the form \(\operatorname{Dyn}_0(G)\) of \(I_S\).
Using again the equivalence of categories \SGARefLink{sga3:xxiv:1.17}{1.17} and the preceding exact
sequence, one deduces formally that there is a reductive
\(\operatorname{Dyn}_0(G)\)-group \(G_0\) of type \(\mathscr R\), and
an \(S\)-isomorphism
\[
  G\simeq\prod_{\operatorname{Dyn}_0(G)/S}G_0.
\]

\par\medskip\noindent\SGAInlineRefTarget{sga3:xxiv:proposition:5-9}\textbf{Proposition 5.9.---}\enspace
\emph{Let \(S\) be a scheme and \(G\) an adjoint semisimple
\(S\)-group isotypic of type \(t\).  There exist an adjoint semisimple
\(\operatorname{Dyn}_0(G)\)-group \(G_0\), simple of type \(t\), and a
unique \(S\)-isomorphism}
\[
  G\simeq\prod_{\operatorname{Dyn}_0(G)/S}G_0.
\]
\emph{Moreover, this isomorphism induces an exact sequence}
\[
  1\longrightarrow
  \prod_{\operatorname{Dyn}_0(G)/S}
    \underline{\operatorname{Aut}}_{\mathrm{gr.}}(G_0)
  \longrightarrow
  \underline{\operatorname{Aut}}_{S\text{-}\mathrm{gr.}}(G)
  \longrightarrow
  \underline{\operatorname{Aut}}_S\bigl(\operatorname{Dyn}_0(G)\bigr)
  \longrightarrow1.
\]

One may plainly suppose that \(G\) has constant type \(n\,t\).  The
first assertion has already been proved (and uniqueness is clear).
The second then follows from the split case by descent.

Combining \SGARefLink{sga3:xxiv:proposition:5-5}{5.5} and \SGARefLink{sga3:xxiv:proposition:5-9}{5.9} gives:

\par\medskip\noindent\SGAInlineRefTarget{sga3:xxiv:proposition:5-10}\textbf{Proposition 5.10.---}\enspace
\emph{Let \(S\) be a scheme, let \(G\) be an adjoint semisimple
\(S\)-group, and put \(\mathfrak D=\operatorname{Dyn}(G)\).}

\emph{(i) There is a canonical decomposition}
\[
  G\simeq
  \prod_{t\in\mathrm T}
  \prod_{\mathfrak D_{0,t}/S}G_{0,t}
  \simeq
  \prod_{\operatorname{Dyn}_0(G)/S}G_0,
\]
\emph{where each \(G_{0,t}\) is a simple adjoint
\(\operatorname{Dyn}_{0,t}(G)\)-group of type \(t\); equivalently,
\(G_0\) is an adjoint semisimple
\(\operatorname{Dyn}_0(G)\)-group whose type at
\(x\in\operatorname{Dyn}_0(G)\) is \(a(x)\in\mathrm T\), where
\(a:\operatorname{Dyn}_0(G)\to\mathrm T_S\) was defined in \SGARefLink{sga3:xxiv:subsection:5-2}{5.2}.}

\emph{(ii) The preceding decompositions induce isomorphic exact
sequences, written vertically below.  Here
\(\underline{\operatorname{Aut}}_{S,a}(\mathfrak D_0)\) denotes the
scheme of automorphisms of \(\operatorname{Dyn}_0(G)\) commuting with
\(a:\operatorname{Dyn}_0(G)\to\mathrm T_S\):}
\[
  % BEGIN INCLUDED SGA3_English_current_buildable_source_20260823/source/inputs/XXIV/native_diagrams/exp24/013_automorphism_exact_columns.tex
\begin{tikzcd}[column sep=large,row sep=large]
1 \arrow[d] & 1 \arrow[d] \\
\displaystyle
\prod_{t\in\mathrm T}\prod_{\mathfrak D_{0,t}/S}
  \underline{\operatorname{Aut}}_{\mathrm{gr.}}(G_{0,t})
  \arrow[r,"\sim"] \arrow[d]
&
\displaystyle
\prod_{\operatorname{Dyn}_0(G)/S}
  \underline{\operatorname{Aut}}_{\mathrm{gr.}}(G_0)
  \arrow[d]
\\
\underline{\operatorname{Aut}}_{S\text{-}\mathrm{gr.}}(G)
  \arrow[r,equal] \arrow[d]
&
\underline{\operatorname{Aut}}_{S\text{-}\mathrm{gr.}}(G)
  \arrow[d]
\\
\displaystyle\prod_{t\in\mathrm T}
  \underline{\operatorname{Aut}}_S
    \bigl(\operatorname{Dyn}_{0,t}(G)\bigr)
  \arrow[r,"\sim"] \arrow[d]
&
\underline{\operatorname{Aut}}_{S,a}
  \bigl(\operatorname{Dyn}_0(G)\bigr)
  \arrow[d]
\\
1 & 1
\end{tikzcd}

% END INCLUDED SGA3_English_current_buildable_source_20260823/source/inputs/XXIV/native_diagrams/exp24/013_automorphism_exact_columns.tex

\]

\par\medskip\noindent\SGAInlineRefTarget{sga3:xxiv:corollary:5-11}\textbf{Corollary 5.11.---}\enspace
\emph{Under the preceding hypotheses, the following three categories
are equivalent:}
\begin{enumerate}
\item[(i)] the category of principal homogeneous bundles under \(G\);
\item[(ii)] the category of principal homogeneous bundles under \(G_0\);
\item[(iii)] the product, for \(t\in\mathrm T\), of the categories of
principal homogeneous bundles under the \(G_{0,t}\).
\end{enumerate}
This follows formally from the preceding decompositions and \SGARefLink{sga3:xxiv:proposition:8-4}{8.4}.

\par\medskip\noindent\SGAInlineRefTarget{sga3:xxiv:corollary:5-12}\textbf{Corollary 5.12.---}\enspace
\emph{There are canonical isomorphisms}
\[
  \underline{\operatorname{Tor}}(G)
  \simeq
  \prod_{t\in\mathrm T}
  \prod_{\mathfrak D_{0,t}/S}
    \underline{\operatorname{Tor}}(G_{0,t})
  \simeq
  \prod_{\operatorname{Dyn}_0(G)/S}
    \underline{\operatorname{Tor}}(G_0),
\]
\emph{and likewise with \(\underline{\operatorname{Tor}}\) replaced
by \(\underline{\operatorname{Bor}}\), respectively by
\(\underline{\operatorname{Kil}}\).}

This is immediate from the split case.

\par\medskip\noindent\SGAInlineRefTarget{sga3:xxiv:remark:5-13}\textbf{Remark 5.13.---}\enspace
The canonical morphism
\[
  \operatorname{Dyn}(G)\longrightarrow\operatorname{Dyn}_0(G)
\]
allows one to regard \(\operatorname{Dyn}(G)\) as a Dynkin scheme over
\(\operatorname{Dyn}_0(G)\); in fact, it is the Dynkin scheme
\(\operatorname{Dyn}(G_0)\) of \(G_0\).

Likewise, if \(T\subset B\) is a Killing couple of \(G\), corresponding
canonically to a Killing couple \(T_0\subset B_0\) of \(G_0\), then the
obstructions to the quasi-splitting of \(G\) and \(G_0\), which lie by
\SGARefLink{sga3:xxiv:3.9}{3.9} in
\[
  \operatorname{Pic}\bigl(\operatorname{Dyn}(G)\bigr)
  =
  \operatorname{Pic}\bigl(\operatorname{Dyn}(G_0)\bigr),
\]
coincide.  Hence:

\par\medskip\noindent\SGAInlineRefTarget{sga3:xxiv:corollary:5-14}\textbf{Corollary 5.14.---}\enspace
\emph{The following conditions are equivalent:}
\begin{enumerate}
\item[(i)] \(G\) is quasi-splittable;
\item[(ii)] \(G_0\) is quasi-splittable;
\item[(iii)] each \(G_{0,t}\), \(t\in\mathrm T\), is quasi-splittable.
\end{enumerate}

\section*{\SGAInlineRefTarget{sga3:xxiv:section:6}6. Automorphisms of Borel Subgroups of Reductive Groups}

\par\medskip\noindent\SGAInlineRefTarget{sga3:xxiv:lemma:6-1}\textbf{Lemma 6.1.---}\enspace
\emph{Let \(S\) be a scheme, let \((G,T,M,R)\) be a split \(S\)-group,
let \(\Delta\) be a system of simple roots of \(R\), and let \(B\) be
the corresponding Borel subgroup.  Then \(B^u\) is generated, as an
\textup{(fppf)} sheaf of groups, by the \(U_\alpha\),
\(\alpha\in\Delta\).}

Let \(H\) be the subgroup sheaf of \(B^u\) generated by the
\(U_\alpha\), \(\alpha\in\Delta\).  We prove
\(H\supset U_\beta\), for \(\beta\in R_+\), by induction on
\(\operatorname{ord}(\beta)=\operatorname{ord}_\Delta(\beta)>0\)
(cf.\ \SGARefLink{sga3:XXI:3.2.15}{Exp.~XXI, 3.2.15}).  The assertion holds by hypothesis when
\(\operatorname{ord}(\beta)=1\).  Suppose
\(\operatorname{ord}(\beta)>1\), and suppose
\(U_\gamma\subset H\) whenever
\(\operatorname{ord}(\gamma)<\operatorname{ord}(\beta)\).  There is an
\(\alpha\in\Delta\) such that \(\beta-\alpha\in R_+\)
(\SGARefLink{sga3:XXI:3.1.1}{Exp.~XXI, 3.1.1}).  Let \(p\) be the largest integer such that
\(\beta-p\alpha=\beta'\in R_+\).  Then
\(U_\alpha,U_{\beta'}\subset H\) and \(\beta'-\alpha\notin R\).
Thus it remains to prove:

\par\medskip\noindent\SGAInlineRefTarget{sga3:xxiv:lemma:6-2}\textbf{Lemma 6.2.---}\enspace
\emph{Keep the notation of \SGARefLink{sga3:XXIII:6.4}{Exp.~XXIII, 6.4}.  Suppose \(p=1\), that is,
\(\beta-\alpha\) is not a root.  If \(H\) is a subgroup sheaf of
\(B^u\) such that \(U_\alpha,U_\beta\subset H\), then
\(U_{i\alpha+j\beta}\subset H\) whenever
\(i\alpha+j\beta\in R\), \(i>0\), \(j>0\).}

Distinguish four cases according as \(q=0,1,2,3\).  In what follows,
\(x\) and \(y\) denote arbitrary sections of
\(\mathbb G_{a,S'}\), where \(S'\to S\).

For \(q=0\), there is nothing to prove.

For \(q=1\), one has
\[
 p_{\alpha+\beta}(\pm x)
 =p_\beta(-y)p_\alpha(-x)p_\beta(y)p_\alpha(x)\in H(S'),
\]
and hence \(U_{\alpha+\beta}\subset H\).

For \(q=2\), one similarly has
\[
 p_{\alpha+\beta}(\pm xy)
 p_{2\alpha+\beta}(\pm x^2y)
 =p_\beta(-y)p_\alpha(-x)p_\beta(y)p_\alpha(x)\in H(S').
\]
Thus, after changing some signs,
\[
 F(x,y)=p_{\alpha+\beta}(xy)p_{2\alpha+\beta}(x^2y)\in H(S').
\]
Since \(p_{\alpha+\beta}\) and \(p_{2\alpha+\beta}\) commute,
\[
 F(x,1)F(1,-x)=p_{2\alpha+\beta}(x^2-x).
\]
\setcounter{footnote}{27}
\footnote{\textit{Editors' note.} We corrected \(F(-1,x)\) to
\(F(1,-x)\).}
If \(a\in\mathbb G_a(S)\), the equation \(X^2-X=a\) defines a free
surjective extension of \(S\), namely
\(\operatorname{Spec}\mathcal O_S[X]/(X^2-X-a)\).  Consequently
\(p_{2\alpha+\beta}(a)\in H(S')\), so
\(U_{2\alpha+\beta}\subset H\), and then
\(U_{\alpha+\beta}\subset H\), since \(F(1,a)\in H(S')\).

For \(q=3\), one similarly has
\[
 F(x,y)=
 p_{\alpha+\beta}(xy)
 p_{2\alpha+\beta}(x^2y)
 p_{3\alpha+\beta}(x^3y)
 p_{3\alpha+2\beta}(x^3y^2)\in H(S'),
\]
and, since
\[
 p_{3\alpha+2\beta}(\pm x)
 =F(1,1)^{-1}p_\beta(-x)F(1,1)p_\beta(x)\in H(S'),
\]
one obtains \(U_{3\alpha+2\beta}\subset H\), and therefore
\[
 K(x,y)=
 p_{\alpha+\beta}(xy)
 p_{2\alpha+\beta}(x^2y)
 p_{3\alpha+\beta}(x^3y)\in H(S').
\]
Computing
\[
 K(x+y,1)K(-x,1)K(1,y)^{-1}
\]
modulo \(U_{3\alpha+2\beta}\), one finds
\[
 p_{2\alpha+\beta}(2x^2+y^2+2xy-y)
 p_{3\alpha+\beta}(y^3+3xy^2+3x^2y-y)\in H(S').
\]
If \(a\in\mathbb G_a(S)\), the ``equations''
\[
  x^2=-xy-y+1-a,\qquad
  y^2=3y-2+3a
\]
define a free surjective extension of \(S\).  One then has
\(y^3+3xy^2+3x^2y-y=0\), and the preceding expression gives
\(p_{2\alpha+\beta}(a)\in H(S')\).
\footnote{\textit{Editors' note.} We corrected
\(p_{3\alpha+\beta}(a)\) to \(p_{2\alpha+\beta}(a)\) and supplied the
end of the argument.}
Thus \(H\) contains \(U_{2\alpha+\beta}\) and
\(U_{3\alpha+2\beta}\), and
\[
 E(x,y)=p_{\alpha+\beta}(xy)p_{3\alpha+\beta}(x^3y)\in H(S').
\]
Since \(p_{\alpha+\beta}(xy)\) and \(p_{3\alpha+\beta}\) commute, the
preceding computation applies:
\[
 E(1,x)E(1,-x)=p_{3\alpha+\beta}(x^3-x).
\]
Hence \(U_{3\alpha+\beta}\subset H\), and then
\(U_{\alpha+\beta}\subset H\).

\par\medskip\noindent\SGAInlineRefTarget{sga3:xxiv:remark:6-2-1}\textbf{Remark 6.2.1.---}\enspace
The preceding proof shows that the hypothesis
``\(H\) contains \(U_\alpha\) and \(U_\beta\)'' could have been
replaced by: ``\(H\) contains \(U_\alpha\) or \(U_\beta\), and \(H\)
is invariant under \(\operatorname{int}(U_\alpha)\) and
\(\operatorname{int}(U_\beta)\).''

\par\medskip\noindent\SGAInlineRefTarget{sga3:xxiv:theorem:6-3}\textbf{Theorem 6.3.---}\enspace
\emph{Let \(S\) be a scheme, let \(G\) and \(G'\) be two semisimple
\(S\)-groups, and let \(B\) (respectively \(B'\)) be a Borel subgroup
of \(G\) (respectively \(G'\)).  Every isomorphism
\(B\xrightarrow{\sim}B'\) is induced by a unique isomorphism
\(G\xrightarrow{\sim}G'\).}

The assertion is local for the étale topology, so we may suppose that
\(G\) is split: \(G=(G,T,M,R)\), with \(B\) defined by the system
\(R_+\) of positive roots of \(R\).  The given isomorphism
\(u:B\xrightarrow{\sim}B'\) induces an isomorphism from \(T\) onto a
maximal torus \(T'\) of \(B'\), hence of \(G'\).  The given
isomorphism \(T\simeq D_S(M)\) gives an isomorphism
\(T'\simeq D_S(M)\) in which the elements of \(R_+\) become the
constant roots of \(B'\) relative to \(T'\), and hence the elements of
\(R\) become the constant roots of \(G'\) relative to \(T'\).  Since
\(G\) and \(G'\) are semisimple, the coroots are determined by duality
(\SGARefLink{sga3:XXI:1.2.5}{Exp.~XXI, 1.2.5}); this proves that \((T',M,R)\) is a splitting of
\(G'\) such that \(\mathscr R(G)=\mathscr R(G')\).

Applying \SGARefLink{sga3:XXIII:5.1}{Exp.~XXIII, 5.1} (the uniqueness theorem), one obtains a unique
isomorphism \(G\xrightarrow{\sim}G'\) agreeing with \(u\) on \(T\)
and on the \(U_\alpha\), \(\alpha\in\Delta\).  By \SGARefLink{sga3:xxiv:lemma:6-1}{6.1}, its restriction
to \(B\) is \(u\).

\par\medskip\noindent\SGAInlineRefTarget{sga3:xxiv:remark:6-3-1}\textbf{Remark 6.3.1.---}\enspace
(i) Using \SGARefLink{sga3:xxii:prop:4.1.9}{Exp.~XXII, 4.1.9} and arguing as in \emph{loc.\ cit.}, 4.2.12,
one may replace ``isomorphism'' in the theorem by ``isogeny''
(respectively, ``central isogeny'').

(ii) The theorem is false for reductive groups.  For example, take
\(G=G'=\mathrm{SL}_2\times\mathbb G_m\), identified with the group
\[
 \left\{
 \begin{pmatrix}
   a&b&0\\
   c&d&0\\
   0&0&h
 \end{pmatrix}
 \ \middle|\ ad-bc=1,\ h\text{ invertible}
 \right\}.
\]

% END INCLUDED SGA3_English_current_buildable_source_20260823/source/inputs/XXIV/tex/components/12_expose_XXIV_proposition55_through_remark631_en.tex

% BEGIN INCLUDED SGA3_English_current_buildable_source_20260823/source/inputs/XXIV/tex/components/13_expose_XXIV_remark631_continuation_through_bibliography_en.tex
% Exposé XXIV, authority local pp. 38--54 / combined-reader pp. 1283--1299.
% Lead-audited against the Polo--Gille authority at 2400 dpi.
% Reinhold English Markdown was comparison-only and is credited in STATUS.
% Provisional authority images preserve diagram content pending native reconstruction.

\setcounter{footnote}{29}

take \(B=B'\) to be the Borel subgroup defined by \(c=0\), and let \(u\)
be the following automorphism of \(B\):
\[
  u(a,b,d,h)=(a,b,d,ah).
\]

\par\medskip\noindent\SGAInlineRefTarget{sga3:xxiv:corollary:6-4}\textbf{Corollary 6.4.---}\enspace
\emph{The functor
\[
  (G,B)\longmapsto B
\]
from the category of pairs \((G,B)\), where \(G\) is an
\(S\)-semisimple group and \(B\) is a Borel subgroup of \(G\), to the
category of \(S\)-group schemes, with isomorphisms as morphisms, is
fully faithful.}

\par\medskip\noindent\SGAInlineRefTarget{sga3:xxiv:corollary:6-5}\textbf{Corollary 6.5.---}\enspace
\emph{Let \(S\) be a scheme, \(G\) an \(S\)-semisimple group, \(B\) a
Borel subgroup of \(G\), and \(B'\) an \(S\)-group locally isomorphic to
\(B\) for the \textup{(fpqc)} topology. Then \(B'\) is locally
isomorphic to \(B\) for the local finite étale
topology,\footnote{\textit{Editors' note.} This topology is denoted
\textup{(étf)} in \SGARefLink{sga3:iv:subsection:6:3}{Exposé IV, 6.3}. In other words, for every \(s\in S\)
there exist an open neighborhood \(U\) of \(s\) and a finite surjective
étale morphism \(V\to U\) such that
\(B'\times_S V\simeq B\times_S V\).} and there exists an
\(S\)-semisimple group \(G'\) of which \(B'\) is a Borel subgroup.
Moreover, \(G'\) is unique up to a unique isomorphism inducing the
identity on \(B'\).}

\par\medskip\noindent\SGAInlineRefTarget{sga3:xxiv:corollary:6-6}\textbf{Corollary 6.6.---}\enspace
\emph{Let \(S\) be a scheme, \(G\) an \(S\)-semisimple group, and \(B\)
a Borel subgroup of \(G\). Then
\(\underline{\operatorname{Aut}}_{S\text{-}\mathrm{gr.}}(B)\) is
representable by an affine smooth \(S\)-scheme,
\(\underline{\operatorname{Autext}}(B)\) is representable by an étale
finite \(S\)-scheme, and the evident morphisms induce an isomorphism of
exact sequences.}

\[
  % BEGIN INCLUDED SGA3_English_current_buildable_source_20260823/source/inputs/XXIV/native_diagrams/exp24/005_borel_automorphism_exact_sequences.tex
\begin{tikzcd}[column sep=small,row sep=large]
1 \arrow[r]
&
\underline{B}_{\mathrm{ad}}(G)
  \arrow[r]
  \arrow[d,"\wr" description]
&
\underline{\Aut}_{S\text{-}\mathrm{gr.}}(G,B)
  \arrow[r]
  \arrow[d,"\wr" description]
&
\underline{\Autext}(G)
  \arrow[r]
  \arrow[d,"\wr" description]
&
1
\\
1 \arrow[r]
&
\underline{B}_{\mathrm{ad}}
  \arrow[r]
&
\underline{\Aut}_{S\text{-}\mathrm{gr.}}(G)
  \arrow[r]
&
\underline{\Autext}(B)
  \arrow[r]
&
1 .
\end{tikzcd}

% END INCLUDED SGA3_English_current_buildable_source_20260823/source/inputs/XXIV/native_diagrams/exp24/005_borel_automorphism_exact_sequences.tex

\]

This follows immediately from \SGARefLink{sga3:xxiv:2.1}{2.1} and the preceding results. We leave
to the reader the task of developing the analogues of the results of
nos.~1, 2, 3, and 4 in the foregoing setting.

\par\medskip\noindent\SGAInlineRefTarget{sga3:xxiv:remark:6-7}\textbf{Remark 6.7.---}\enspace
If \(S\) is a scheme and \(B\) an \(S\)-group, then \(B\) can be a
Borel subgroup of only one semisimple group, uniquely determined by
\(B\). It remains to characterize those \(S\)-groups \(B\) which are
Borel subgroups of semisimple
groups.\footnote{\textit{Editors' note.} One may ask whether every
smooth affine \(S\)-group whose geometric fibers are Borel subgroups of
semisimple groups is itself a Borel subgroup of a semisimple
\(S\)-group. We have suppressed the counterexample in the original,
with \(S\) the scheme of dual numbers over a field \(k\), because it
was erroneous, as M.~Demazure pointed out.}

\section*{\SGAInlineRefTarget{sga3:xxiv:section:7}7. Representability of the functors
\(\underline{\operatorname{Hom}}_{S\text{-}\mathrm{gr.}}(G,H)\), for
\(G\) reductive}

\subsection*{\SGAInlineRefTarget{sga3:xxiv:subsection:7-1}7.1. The split case}

\subsubsection*{7.1.1.}

Let \(S\) be a scheme, \(G\) a pinned \(S\)-group, \(T\) its maximal
torus, and \(\Delta\) the set of simple roots. For
\(\alpha\in\Delta\), let
\[
  u_\alpha\in U_\alpha^\times(S),
  \qquad
  w_\alpha\in\operatorname{Norm}_G(T)(S)
\]
be the elements defined by the pinning.

Let \(H\) be an \(S\)-group scheme. We are interested in the functor
\(\underline{\operatorname{Hom}}_{S\text{-}\mathrm{gr.}}(G,H)\), and
more precisely in the morphism
\[
  q:\underline{\operatorname{Hom}}_{S\text{-}\mathrm{gr.}}(G,H)
  \longrightarrow
  \underline{\operatorname{Hom}}_{S\text{-}\mathrm{gr.}}(T,H).
\]
Let
\[
  f_T:T\longrightarrow H
\]
be an \(S\)-group homomorphism and put \(q^{-1}(f_T)=\mathscr F\).
Thus
\[
  \mathscr F(S')
  =
  \left\{
    f\in\operatorname{Hom}_{S'\text{-}\mathrm{gr.}}
      (G_{S'},H_{S'})
    \,\middle|\,
    f|_{T_{S'}}=(f_T)_{S'}
  \right\}.
\]
There is a morphism of \(S\)-functors
\[
  i:\mathscr F\longrightarrow H^{2n},
\]
where \(n=\operatorname{Card}(\Delta)\), defined on points by
\[
  i(f)=\bigl(f(u_\alpha),f(w_\alpha)\bigr)_{\alpha\in\Delta}.
\]
By \SGARefLink{sga3:XXIII:1.9}{Exposé XXIII, 1.9}, \(i\) is a monomorphism.

\par\medskip\noindent\SGAInlineRefTarget{sga3:xxiv:proposition:7-1-2}\textbf{Proposition 7.1.2.---}\enspace
\emph{If \(H\) is separated over \(S\), then \(\mathscr F\) is
representable and \(i\) is a closed immersion.}

The usual technique of relative
representability\footnote{\textit{Editors' note.} See, for example, the
proof of \SGARefLink{sga3:xxii:thm:5.8.1}{Exposé XXII, 5.8.1}.} shows that it is enough to prove the
following. Given sections
\[
  v_\alpha,h_\alpha\in H(S),
  \qquad \alpha\in\Delta,
\]
the \(S\)-schemes \(S'\) for which there exists an \(S'\)-homomorphism
\[
  f:G_{S'}\longrightarrow H_{S'}
\]
extending \((f_T)_{S'}\) and satisfying
\[
  f(u_\alpha)=v_\alpha,
  \qquad
  f(w_\alpha)=h_\alpha
\]
are exactly those which factor through a certain closed subscheme of
\(S\). We may plainly suppose \(S\) affine.

For convenience, say that a morphism \(Y\to
X\)\footnote{\textit{Editors' note.} We have corrected \(X\to Y\) to
\(Y\to X\).} of affine schemes satisfies condition \textup{(L)} if
\(Y\) is the spectrum of an \(\mathcal O(X)\)-algebra which is a free
\(\mathcal O(X)\)-module. Within the category of affine schemes,
products and composites of morphisms satisfying \textup{(L)} again
satisfy \textup{(L)}, and \textup{(L)} is stable under base change.

\par\medskip\noindent\SGAInlineRefTarget{sga3:xxiv:lemma:7-1-3}\textbf{Lemma 7.1.3.---}\enspace
\emph{Suppose that \(S\) is affine, and let \(\alpha\in\Delta\).
Consider the morphism
\[
  a:T\times_S T\longrightarrow G_{a,S},
  \qquad
  a(t,t')=\alpha(t)+\alpha(t').
\]
Then:
\begin{enumerate}
\item[(i)] \(a\) is faithfully flat and of finite presentation;
\item[(ii)] if \(R\) is the fiber square of \(a\), the structure
  morphism \(R\to S\) satisfies \textup{(L)}.
\end{enumerate}}

The morphism \(a\) is surjective, since
\(\alpha:T\to G_{m,S}\) is surjective, and it is plainly of finite
presentation. Since \(\alpha\) satisfies \textup{(L)}, it is enough,
for both assertions, to show that
\[
  u:G_{m,S}^2\longrightarrow G_{a,S},
  \qquad
  u(x,y)=x+y
\]
satisfies \textup{(L)}. Equivalently, for every ring \(A\), the ring
\[
  A[X,Y,X^{-1},Y^{-1}]
\]
must be free over its subring \(A[X+Y]\).\footnote{\textit{Editors'
note.} We have simplified the proof in the original.} Now \(A[X,Y]\)
is free with basis \(\{1,X\}\) over
\[
  B=A[X+Y,XY],
\]
because \(X^2-(X+Y)X+XY=0\). Hence
\[
  A[X^{\pm1},Y^{\pm1}]
  =A[X,Y]_{XY}
\]
is free over
\[
  B_{XY}=A[X+Y,(XY)^{\pm1}]
\]
with basis \(\{1,X\}\), and free over \(A[X+Y]\) with basis
\[
  (XY)^p,\quad X(XY)^p,
  \qquad p\in\mathbb Z.
\]

\par\medskip\noindent\SGAInlineRefTarget{sga3:xxiv:lemma:7-1-4}\textbf{Lemma 7.1.4.---}\enspace
\emph{Let \(\alpha\in\Delta\), and let
\[
  b:T\times_S T\longrightarrow H
\]
be the morphism defined on points by
\[
  b(t,t')
  =
  \operatorname{int}(f_T(t))(v_\alpha)\,
  \operatorname{int}(f_T(t'))(v_\alpha).
\]
Let \(f_\alpha:U_\alpha\to H\) be an \(S\)-morphism. The following
conditions are equivalent:\footnote{\textit{Editors' note.} In what
follows the group law of \(U_\alpha\) is written additively. If
\(p_\alpha:G_{a,S}\xrightarrow{\sim}U_\alpha\) is the isomorphism with
\(p_\alpha(1)=u_\alpha\), and \(x=p_\alpha(z)\), then
\(\alpha(t)x\) denotes \(p_\alpha(\alpha(t)z)\), while
\(a(t,t')u_\alpha\) denotes
\(p_\alpha(a(t,t'))=\alpha(t)u_\alpha+\alpha(t')u_\alpha\).}
\begin{enumerate}
\item[\SGAInlineRefTarget{sga3:xxiv:lemma:7-1-4:item:i}(i)] \(f_\alpha\) is a group homomorphism,
  \(f_\alpha(u_\alpha)=v_\alpha\), and
  \[
    f_T(t)f_\alpha(x)f_T(t)^{-1}
    =
    f_\alpha(\alpha(t)x)
  \]
  for all \(t\in T(S')\), \(x\in U_\alpha(S')\), and \(S'\to S\);
\item[\SGAInlineRefTarget{sga3:xxiv:lemma:7-1-4:item:ii}(ii)] \(f_\alpha(u_\alpha)=v_\alpha\), and
  \[
    f_\alpha(a(t,t')u_\alpha)=b(t,t')
  \]
  for all \(t,t'\in T(S')\) and \(S'\to S\).
\end{enumerate}}

If \(f_\alpha\) satisfies \SGARefLink{sga3:xxiv:lemma:7-1-4:item:i}{(i)}, then
\[
  f_\alpha(\alpha(t)u_\alpha)
  =
  \operatorname{int}(f_T(t))(v_\alpha),
\]
which immediately gives \SGARefLink{sga3:xxiv:lemma:7-1-4:item:ii}{(ii)}. Conversely, suppose \SGARefLink{sga3:xxiv:lemma:7-1-4:item:ii}{(ii)}. Since \(a\)
is an \textup{(fpqc)} covering, to prove the last condition in \SGARefLink{sga3:xxiv:lemma:7-1-4:item:i}{(i)} it
suffices to show that, for \(t,t',t''\in T(S')\),
\[
  f_T(t)\,
  f_\alpha(a(t',t'')u_\alpha)\,
  f_T(t)^{-1}
  =
  f_\alpha(\alpha(t)a(t',t'')u_\alpha).
\]
This is the same as
\[
  f_T(t)b(t',t'')f_T(t)^{-1}=b(tt',tt''),
\]
which follows directly from the definition of \(b\).

It remains to prove that \(f_\alpha\) is a group homomorphism. The
preceding property gives
\[
\begin{aligned}
 f_\alpha(\alpha(t)u_\alpha)
 f_\alpha(\alpha(t')u_\alpha)
 &=
 \bigl(f_T(t)f_\alpha(u_\alpha)f_T(t)^{-1}\bigr)
 \bigl(f_T(t')f_\alpha(u_\alpha)f_T(t')^{-1}\bigr)\\
 &=b(t,t')\\
  &=f_\alpha\bigl(\alpha(t)u_\alpha+
                   \alpha(t')u_\alpha\bigr).
\end{aligned}
\]
Thus \(f_\alpha(x+x')=f_\alpha(x)f_\alpha(x')\) whenever \(x,x'\) are
sections of the schematically dense open subset
\(U_\alpha^\times\subset U_\alpha\). \SGARefLink{sga3:xviii:prop:1.4}{Exposé XVIII, 1.4} then gives the
result.

\subsubsection*{7.1.5.}

Fix provisionally an \(\alpha\in\Delta\). Since \(a\) is faithfully
flat and quasi-compact, \(G_{a,S}\) is identified with the quotient of
\(T\times_S T\) by the equivalence relation \(R\) defined by \(a\).
Write \(i_1,i_2:R\to T\times_S T\) for its two projections.

\[
  % BEGIN INCLUDED SGA3_English_current_buildable_source_20260823/source/inputs/XXIV/native_diagrams/exp24/006_equivalence_relation_pair.tex
\begin{tikzcd}[column sep=huge]
R
  \arrow[r,shift left=.7ex,"i_1"]
  \arrow[r,shift right=.7ex,"i_2"']
&
T\times_S T
\end{tikzcd}

% END INCLUDED SGA3_English_current_buildable_source_20260823/source/inputs/XXIV/native_diagrams/exp24/006_equivalence_relation_pair.tex

\]

For \(b\) to factor through \(a\),\footnote{\textit{Editors' note.}
That is, for condition \SGARefLink{sga3:xxiv:lemma:7-1-4:item:ii}{7.1.4(ii)} to hold.} it is necessary and
sufficient that
\[
  b\circ i_1=b\circ i_2.
\]
Let \(K\) be the equalizer of \(b\circ i_1\) and \(b\circ i_2\).

\[
  % BEGIN INCLUDED SGA3_English_current_buildable_source_20260823/source/inputs/XXIV/native_diagrams/exp24/007_kernel_pair.tex
\begin{tikzcd}[column sep=huge]
R
  \arrow[r,shift left=.7ex,"b\circ i_1"]
  \arrow[r,shift right=.7ex,"b\circ i_2"']
&
H
\end{tikzcd}

% END INCLUDED SGA3_English_current_buildable_source_20260823/source/inputs/XXIV/native_diagrams/exp24/007_kernel_pair.tex

\]

Then the required condition is \(K=R\). Since \(H\) is separated over
\(S\), \(K\) is a closed subscheme of \(R\). On the other hand,
\(R\) is essentially free over \(S\), by \SGARefLink{sga3:xxiv:lemma:7-1-3}{7.1.3}.

\subsubsection*{7.1.6.}

It follows that, for an \(S\)-scheme \(S'\), an \(f_\alpha\) satisfying
\SGARefLink{sga3:xxiv:lemma:7-1-4:item:i}{7.1.4(i)} exists over \(S'\), and is then unique, if and only if
\[
  K_{S'}=R_{S'}
\]
and the resulting morphism satisfies
\[
  f_\alpha(u_\alpha)=v_\alpha.
\]
The first condition is equivalent to requiring \(S'\to S\) to factor
through a certain closed subscheme of \(S\) (\SGARefLink{sga3:viii:theorem:6:4}{Exposé VIII,
6.4}\footnote{\textit{Editors' note.} See also the addition in Exposé
VI\textsubscript{B}, \SGARefLink{thm:sga3-VIB-6.2.3}{6.2.3}.}). After replacing \(S\) by this closed
subscheme, the second condition again defines a closed subscheme of
\(S\), since it is the equality of two sections of the separated
\(S\)-scheme \(H\).

After replacing \(S\) by the resulting closed subscheme, we may thus
suppose that a morphism \(f_\alpha:U_\alpha\to H\) satisfying \SGARefLink{sga3:xxiv:lemma:7-1-4:item:i}{7.1.4(i)}
exists. Intersecting the closed subschemes obtained for all
\(\alpha\in\Delta\), we may suppose this condition holds for every
\(\alpha\in\Delta\).

\subsubsection*{7.1.7.}

Similarly, for each \(\alpha\in\Delta\), consider the two
\(S\)-group homomorphisms
\[
  f_T\circ\operatorname{int}(w_\alpha),
  \qquad
  \operatorname{int}(h_\alpha)\circ f_T
  :T\longrightarrow H.
\]
Since \(H\) is separated over \(S\) and \(T\) is essentially free over
\(S\), the same reasoning shows that, after replacing \(S\) by a
closed subscheme, one may suppose that
\[
  f_T\circ\operatorname{int}(w_\alpha)
  =
  \operatorname{int}(h_\alpha)\circ f_T
\]
for every \(\alpha\in\Delta\).

\subsubsection*{7.1.8.}

The generators-and-relations theorem (\SGARefLink{sga3:xxiii:subsection:3-5}{Exposé XXIII, 3.5}) now shows that
a group homomorphism \(f:G\to H\) with the required properties exists
if and only if a certain finite set of algebraic relations among the
sections
\[
  h_\alpha,\qquad
  v_\alpha,\qquad
  f_T\bigl(\alpha^*(-1)\bigr),
  \qquad \alpha\in\Delta,
\]
is satisfied:
\[
  R_i\!\left(
    (h_\alpha)_\alpha,
    (v_\alpha)_\alpha,
    \bigl(f_T(\alpha^*(-1))\bigr)_\alpha
  \right)=e,
  \qquad i=1,\ldots,n.
\]
Because \(H\) is separated over \(S\), these relations again define a
closed subscheme of \(S\), proving \SGARefLink{sga3:xxiv:proposition:7-1-2}{Proposition 7.1.2}.

\par\medskip\noindent\SGAInlineRefTarget{sga3:xxiv:corollary:7-1-9}\textbf{Corollary 7.1.9.---}\enspace
\emph{Let \(S\) be a scheme, \(G\) a split reductive \(S\)-group, \(T\)
its maximal torus, and \(H\) a separated \(S\)-group scheme. Let
\(\mathcal P\) be a property of morphisms satisfying:
\begin{enumerate}
\item[(i)] every closed immersion has \(\mathcal P\);
\item[(ii)] a composite of morphisms having \(\mathcal P\) has
  \(\mathcal P\);
\item[(iii)] \(\mathcal P\) is stable under base change;
\item[(iv)] the structure morphism \(H\to S\) has \(\mathcal P\).
\end{enumerate}
Then the canonical morphism
\[
  \underline{\operatorname{Hom}}_{S\text{-}\mathrm{gr.}}(G,H)
  \longrightarrow
  \underline{\operatorname{Hom}}_{S\text{-}\mathrm{gr.}}(T,H)
\]
is relatively representable by a separated morphism having
\(\mathcal P\).\footnote{\textit{Editors' note.} More explicitly, for
every representable
\(S'\to\underline{\operatorname{Hom}}_{S\text{-}\mathrm{gr.}}(T,H)\),
the fiber product
\(\underline{\operatorname{Hom}}_{S\text{-}\mathrm{gr.}}(G,H)
\times S'\) is representable, and its morphism to \(S'\) is separated
and has \(\mathcal P\).}}

Indeed, we may suppose \(G\) pinned. The structure morphism
\(H^{2n}\to S\) has \(\mathcal P\), and \SGARefLink{sga3:xxiv:proposition:7-1-2}{7.1.2} gives the assertion.

\par\medskip\noindent\SGAInlineRefTarget{sga3:xxiv:corollary:7-1-10}\textbf{Corollary 7.1.10.---}\enspace
\emph{Let \(S\) be a scheme, \(G\) a split reductive \(S\)-group, and
\(H\) a smooth \(S\)-group scheme with affine
fibers.\footnote{\textit{Editors' note.} We have suppressed the
superfluous hypothesis that \(H\) be quasi-separated.} Then
\(\underline{\operatorname{Hom}}_{S\text{-}\mathrm{gr.}}(G,H)\) is
representable by an \(S\)-scheme locally of finite presentation and
separated over \(S\).}

Since \(H\) is smooth over \(S\), its neutral component \(H^0\) is
defined; it has affine fibers and is smooth, separated, and of finite
presentation over \(S\) (Exposé VI\textsubscript{B}, \SGARefLink{thm:sga3-VIB-3.10}{3.10} and \SGARefLink{sga3:xxiv:proposition:5-5}{5.5}).
Since \(G\) has connected fibers, \(H\) may be replaced by \(H^0\).
As \(G\) and \(H\) are then of finite presentation, we may suppose
\(S\) noetherian and apply \SGARefLink{sga3:xxiv:corollary:7-1-9}{7.1.9} with \(\mathcal P\) equal to the
property ``of finite type''. By \SGARefLink{sga3:xv:theorem:8:8}{Exposé XV, 8.9},
\(\underline{\operatorname{Hom}}_{S\text{-}\mathrm{gr.}}(T,H)\) is
representable by a separated \(S\)-scheme locally of finite type.

\par\medskip\noindent\SGAInlineRefTarget{sga3:xxiv:remark:7-1-11}\textbf{Remark 7.1.11.---}\enspace
If \(H\) is affine over \(S\), the reference to Exposé XV may be
replaced by \SGARefLink{sga3:xi:corollary:4:2}{Exposé XI, 4.2}.

\subsection*{\SGAInlineRefTarget{sga3:xxiv:subsection:7-2}7.2. General case}

\par\medskip\noindent\SGAInlineRefTarget{sga3:xxiv:proposition:7-2-1}\textbf{Proposition 7.2.1.---}\enspace
\emph{Let \(S\) be a scheme, \(G\) a reductive \(S\)-group, and \(T\) a
maximal torus of \(G\). Let \(H\) be an \(S\)-group scheme whose
structure morphism \(H\to S\) satisfies:
\begin{quote}
\textup{(+)} Every point \(s\in S\) has an open neighborhood \(U\)
such that \(H_U\to U\) is quasi-projective.
\end{quote}
Then the canonical morphism
\[
  \underline{\operatorname{Hom}}_{S\text{-}\mathrm{gr.}}(G,H)
  \longrightarrow
  \underline{\operatorname{Hom}}_{S\text{-}\mathrm{gr.}}(T,H)
\]
is relatively representable by a morphism satisfying \textup{(+)}.}

When \(G\) is split relative to \(T\), apply 7.1.9 with
\(\mathcal P=\textup{(+)}\). When \(G\) is locally isotrivial, for
example when \(G\) is semisimple (4.2.2), note that the assertion is
local for the local finite étale topology. Indeed,
quasi-projectivity is local for the global finite étale topology and
ensures effectiveness of descent for that topology (SGA~1, VIII, 7.7).
The general case follows from the next lemma.

\par\medskip\noindent\SGAInlineRefTarget{sga3:xxiv:lemma:7-2-2}\textbf{Lemma 7.2.2.---}\enspace
\emph{Let \(S\) be a scheme, \(G\) a reductive \(S\)-group, \(T\) a
maximal torus of \(G\), \(G'\) the derived group of \(G\) (Exposé XXII,
6.2), and \(T'=T\cap G'\) the maximal torus of \(G'\) corresponding to
\(T\) (Exposé XXII, 6.2.8). The natural diagram involving
\(G,G',T,T'\) is a pushout in the category of \(S\)-group sheaves.
Equivalently, for every \(S\)-group sheaf \(H\), the corresponding
square of Hom-functors is cartesian.}

\[
  % BEGIN INCLUDED SGA3_English_current_buildable_source_20260823/source/inputs/XXIV/native_diagrams/exp24/008_pushout_square.tex
\begin{tikzcd}[column sep=large,row sep=large]
G & G' \arrow[l] \\
T \arrow[u] & T' \arrow[l] \arrow[u]
\end{tikzcd}

% END INCLUDED SGA3_English_current_buildable_source_20260823/source/inputs/XXIV/native_diagrams/exp24/008_pushout_square.tex

\]

\[
  % BEGIN INCLUDED SGA3_English_current_buildable_source_20260823/source/inputs/XXIV/native_diagrams/exp24/009_hom_cartesian_square.tex
\begin{tikzcd}[column sep=large,row sep=large]
\underline{\operatorname{Hom}}_{S\text{-}\mathrm{gr.}}(G,H)
  \arrow[r]
  \arrow[d]
&
\underline{\operatorname{Hom}}_{S\text{-}\mathrm{gr.}}(G',H)
  \arrow[d]
\\
\underline{\operatorname{Hom}}_{S\text{-}\mathrm{gr.}}(T,H)
  \arrow[r]
&
\underline{\operatorname{Hom}}_{S\text{-}\mathrm{gr.}}(T',H)
\end{tikzcd}

% END INCLUDED SGA3_English_current_buildable_source_20260823/source/inputs/XXIV/native_diagrams/exp24/009_hom_cartesian_square.tex

\]

If \(\operatorname{rad}(G)\) is the radical of \(G\), then
\(\operatorname{rad}(G)\subset T\), so
\[
  \operatorname{rad}(G)\cap G'
  =
  \operatorname{rad}(G)\cap T'
  =
  K.
\]
Multiplication in \(G\) induces isogenies (\SGARefLink{sga3:xxii:subsection:6-2}{Exposé XXII, 6.2})
\[
  i:G'\times_S\operatorname{rad}(G)\longrightarrow G,
  \qquad
  j:T'\times_S\operatorname{rad}(G)\longrightarrow T.
\]
Let
\[
  f_{G'}:G'\to H,\qquad
  f_T:T\to H,\qquad
  f_{T'}:T'\to H
\]
be \(S\)-group homomorphisms such that
\[
  f_{G'}|_{T'}=f_T|_{T'}=f_{T'}.
\]
We show that there is a unique \(S\)-group homomorphism \(f:G\to H\)
inducing \(f_{G'}\) and \(f_T\). Put
\[
  f_{\mathrm{rad}}
  =
  f_T|_{\operatorname{rad}(G)}
\]
\footnote{\textit{Editors' note.} We have replaced \(f_r\) by
\(f_{\mathrm{rad}}\).}
and define
\[
  f_1=f_{G'}\cdot f_{\mathrm{rad}}:
  G'\times_S\operatorname{rad}(G)\longrightarrow H.
\]
For \(f\) to exist, and it is plainly then unique, it is necessary and
sufficient that \(f_1\) be trivial on the kernel of \(i\), that is,
that \(f_{G'}\) and \(f_{\mathrm{rad}}\) coincide on \(K\). The
existence of \(f_T\) similarly shows that \(f_{T'}\) and
\(f_{\mathrm{rad}}\) coincide on \(K\). Finally, \(f_{G'}\) and
\(f_{T'}\) coincide on \(K\subset T'\), proving the lemma.

Arguing as in \SGARefLink{sga3:xxiv:corollary:7-1-10}{7.1.10}, \SGARefLink{sga3:xxiv:proposition:7-2-1}{Proposition 7.2.1} gives:

\par\medskip\noindent\SGAInlineRefTarget{sga3:xxiv:corollary:7-2-3}\textbf{Corollary 7.2.3.---}\enspace
\emph{Let \(S\) be a scheme, \(G\) a reductive \(S\)-group, and \(H\)
a smooth, quasi-projective \(S\)-group scheme with affine fibers. Then
\(\underline{\operatorname{Hom}}_{S\text{-}\mathrm{gr.}}(G,H)\) is
representable by an \(S\)-scheme locally of finite presentation and
separated over \(S\).}

\subsection*{\SGAInlineRefTarget{sga3:xxiv:subsection:7-3}7.3. Phenomena specific to characteristic zero}

If \(G\) and \(H\) are smooth \(S\)-groups, with Lie algebras
\(\mathfrak g\) and \(\mathfrak h\), there is a canonical morphism
\[
  \Lie:
  \underline{\operatorname{Hom}}_{S\text{-}\mathrm{gr.}}(G,H)
  \longrightarrow
  \underline{\operatorname{Hom}}_{\mathcal O_S\text{-}\mathrm{Lie}}
    (\mathfrak g,\mathfrak h),
\]
where the functor on the right has the evident meaning.

\par\medskip\noindent\SGAInlineRefTarget{sga3:xxiv:proposition:7-3-1}\textbf{Proposition 7.3.1.---}\enspace
\emph{Let \(S\) be a scheme of characteristic zero, that is, a
\(\mathbb Q\)-scheme; let \(G\) be a reductive \(S\)-group, and \(H\) a
smooth, quasi-projective \(S\)-group scheme with affine fibers.
\begin{enumerate}
\item[\SGAInlineRefTarget{sga3:xxiv:proposition:7-3-1:item:i}(i)]
\(\underline{\operatorname{Hom}}_{S\text{-}\mathrm{gr.}}(G,H)\) is
representable by a smooth separated \(S\)-scheme.
\item[(ii)] If \(G\) is semisimple, this scheme is affine and of finite
presentation over \(S\).
\item[\SGAInlineRefTarget{sga3:xxiv:proposition:7-3-1:item:iii}(iii)] If \(G\) is simply connected (\SGARefLink{sga3:xxii:def:4.3.3}{Exposé XXII, 4.3.3}), the
canonical morphism
\[
  \Lie:
  \underline{\operatorname{Hom}}_{S\text{-}\mathrm{gr.}}(G,H)
  \longrightarrow
  \underline{\operatorname{Hom}}_{\mathcal O_S\text{-}\mathrm{Lie}}
    (\mathfrak g,\mathfrak h)
\]
is an isomorphism.
\end{enumerate}}

\par\medskip\noindent\SGAInlineRefTarget{sga3:xxiv:lemma:7-3-2}\textbf{Lemma 7.3.2.---}\enspace
\emph{Let \(k\) be a field of characteristic zero, \(G\) a reductive
\(k\)-group, \(V\) a finite-dimensional \(k\)-vector space, and
\(G\to\operatorname{GL}(V)\) a linear representation. Then
\[
  H^0(G,V)=H^0(\mathfrak g,V),
  \qquad
  H^i(G,V)=0\quad(i>0).
\]}

The first equality is valid in general for a smooth connected
group.\footnote{\textit{Editors' note.} The inclusion
\(V^G\subset V^{\mathfrak g}\) is clear. Conversely,
\(C=\operatorname{Centr}_G(V^{\mathfrak g})\) is a closed subgroup of
\(G\), smooth because \(\operatorname{char}(k)=0\). By \SGARefLink{sga3:ii:theorem:5-3-1}{Exposé II,
5.3.1},
\(\Lie(C)=\operatorname{Centr}_{\mathfrak g}(V^{\mathfrak g})
=\mathfrak g\); because \(C\) is smooth and \(G\) connected, \(C=G\).
Thus \(V^{\mathfrak g}\subset V^G\). See also \([{\rm DG70}]\),
\S II.6, Proposition 2.1(c).} For a reductive group it can also be
proved as follows. We may suppose \(k\) algebraically closed, hence
\(G\) split and generated by subgroups isomorphic to
\(G_{m,k}\),\footnote{\textit{Editors' note.} The union of the maximal
tori of \(G\) is dense in \(G\); cf.\ \emph{Bible}, \S6.5,
Theorem~5, or \([{\rm Ch05}]\), \S6.6, Theorem~6.} and the assertion
for this group is immediate.

The first equality shows that the functor \(H^0(G,V)\) is exact in
\(V\) when \(G\) is semisimple: \(\mathfrak g\) is then a semisimple
Lie algebra, and one applies \([{\rm BLie}]\), \S I.6,
Exercise~1(b). The assertion remains true when \(G\) is reductive.
Indeed, introducing the radical \(R\) of
\(G\)\footnote{\textit{Editors' note.} Observe that \(R\) is a torus.}
and the semisimple quotient \(G/R\), one has
\[
  H^0(G,V)=H^0\bigl(G/R,H^0(R,V)\bigr),
\]
so \(H^0(G,-)\) is a composite of two exact functors. \SGARefLink{sga3:i:theorem:5-3-1}{Exposé I, 5.3.1}
then gives \(H^i(G,V)=0\) for \(i>0\).

\par\medskip\noindent\SGAInlineRefTarget{sga3:xxiv:remark:7-3-3}\textbf{Remark 7.3.3.---}\enspace
Under the preceding hypotheses, if \(G\) is semisimple, then
\[
  H^1(\mathfrak g,V)=H^2(\mathfrak g,V)=0;
\]
see \([{\rm BLie}]\), \emph{loc.\ cit.}, (b) and (d).

\subsubsection*{7.3.4.}

We prove the proposition. By \SGARefLink{sga3:xxiv:corollary:7-2-3}{7.2.3},
\(\underline{\operatorname{Hom}}_{S\text{-}\mathrm{gr.}}(G,H)\) is
already representable by an \(S\)-scheme locally of finite
presentation and separated over \(S\). To show it smooth, it suffices
to show it infinitesimally smooth (\SGARefLink{sga3:xi:remark:1:8}{Exposé XI,
1.8}),\footnote{\textit{Editors' note.} This means that for every local
artinian \(S\)-scheme \(S'\), with closed point \(s'\), every
homomorphism of \(\kappa(s')\)-groups
\(G_{s'}\to H_{s'}\) lifts to an \(S'\)-group homomorphism
\(G_{S'}\to H_{S'}\). By \SGARefLink{sga3:iii:corollary:2-8}{Exposé III, 2.8}, this follows from the
vanishing of \(H^2(G_{s'},V)\), where
\(V=\Lie(H_{s'}/\kappa(s'))\).} and this follows from \SGARefLink{sga3:iii:corollary:2-8:item:i}{Exposé III,
2.8(i)} and \SGARefLink{sga3:xxiv:lemma:7-3-2}{7.3.2}. This proves \SGARefLink{sga3:xxiv:proposition:7-3-1:item:i}{(i)}.

We next show that \SGARefLink{sga3:iii:corollary:2-8:item:ii}{(ii)} follows from \SGARefLink{sga3:xxiv:proposition:7-3-1:item:iii}{(iii)}. First it is enough to prove
that
\(\underline{\operatorname{Hom}}_{S\text{-}\mathrm{gr.}}(G,H)\) is
affine over \(S\), since it is then of finite presentation, being
smooth by \SGARefLink{sga3:xxiv:proposition:7-3-1:item:i}{(i)}. In any case,
\[
  \underline{\operatorname{Hom}}_{\mathcal O_S\text{-}\mathrm{Lie}}
    (\mathfrak g,\mathfrak h)
\]
is represented by a closed subscheme of
\[
  \underline{\operatorname{Hom}}_{\mathcal O_S\text{-}\mathrm{mod}}
    (\mathfrak g,\mathfrak h)
  \simeq
  \mathrm W\!\left(
    \mathfrak g^\vee\otimes_{\mathcal O_S}\mathfrak h
  \right),
\]
which is an affine \(S\)-scheme. Thus the desired conclusion follows
when \(G\) is simply connected.

In general, we may suppose \(G\) split, so that
\(G\simeq\operatorname{\acute Ep}_S(\mathscr R)\). Introduce the
simply connected root datum \(\operatorname{sc}(\mathscr R)\)
(\SGARefLink{sga3:XXI:6.5.5}{Exposé XXI, 6.5.5}). The existence theorem (\SGARefLink{sga3:xxv:theorem:1:1}{Exposé XXV, 1.1}) gives a
simply connected \(S\)-group \(\widetilde G\) and a central isogeny
\[
  \widetilde G\longrightarrow G.
\]
Its kernel \(K\) is a finite diagonalizable \(S\)-group, hence a
twisted constant \(S\)-group because \(S\) has characteristic zero.
The functor
\(\underline{\operatorname{Hom}}_{S\text{-}\mathrm{gr.}}(K,H)\) is
plainly represented by a separated \(S\)-scheme. For example, if
\(K\simeq(\mathbb Z/n\mathbb Z)_S\), then it is isomorphic to
\[
  \ker\!\left(H\xrightarrow{\,n\,}H\right).
\]
There is an exact sequence of pointed \(S\)-schemes
\[
  1\longrightarrow
  \underline{\operatorname{Hom}}_{S\text{-}\mathrm{gr.}}(G,H)
  \xrightarrow{u}
  \underline{\operatorname{Hom}}_{S\text{-}\mathrm{gr.}}
    (\widetilde G,H)
  \longrightarrow
  \underline{\operatorname{Hom}}_{S\text{-}\mathrm{gr.}}(K,H).
\]
Thus \(u\) is a closed immersion, and
\(\underline{\operatorname{Hom}}_{S\text{-}\mathrm{gr.}}(G,H)\) is
affine over \(S\).

\SGAInlineRefTarget{sga3:xxiv:subsubsection:7-3-5}\subsubsection*{7.3.5.}

It remains to prove \SGARefLink{sga3:xxiv:proposition:7-3-1:item:iii}{(iii)}. As in the proof of \SGARefLink{sga3:xxiv:proposition:7-3-1:item:i}{(i)}, and using \SGARefLink{sga3:xxiv:remark:7-3-3}{7.3.3},
one shows that
\[
  \underline{\operatorname{Hom}}_{\mathcal O_S\text{-}\mathrm{Lie}}
    (\mathfrak g,\mathfrak h)
\]
is smooth over \(S\). We may therefore suppose
\(S=\operatorname{Spec}(k)\), where \(k\) is algebraically closed of
characteristic zero. It is enough to show that \(\Lie\) is bijective
on \(k\)-points and induces a bijection on tangent spaces at
corresponding points. First we show that
\[
  \operatorname{Hom}_{k\text{-}\mathrm{gr.}}(G,H)
  \longrightarrow
  \operatorname{Hom}_{k\text{-}\mathrm{Lie}}
    (\mathfrak g,\mathfrak h)
\]
is bijective.

If \(u:G\to H\) is a \(k\)-group homomorphism, its graph is a connected
subgroup of \(G\times_k H\), determined by its Lie algebra, which is
the graph of \(\Lie(u)\). Hence the map is injective. Conversely, if
\(v:\mathfrak g\to\mathfrak h\) is a Lie algebra homomorphism, its
graph \(\mathfrak k\) is a Lie subalgebra of
\(\mathfrak g\times\mathfrak h\), isomorphic to \(\mathfrak g\).
Because \(\mathfrak g\) is its own derived algebra, so is
\(\mathfrak k\). A theorem of Chevalley
\([{\rm Ch51}]\), \S14, Theorem~15, therefore shows that
\(\mathfrak k\) is the Lie algebra of a connected subgroup \(K\) of
\(G\times_k H\).\footnote{\textit{Editors' note.} See also
\([{\rm Bo91}]\), II 7.9. We have also added the following sentence.}

Since \(\mathfrak k\simeq\mathfrak g\) is semisimple, \(K\) is a
semisimple \(k\)-group (the French prints \(G\) here; the argument
and the following sentence require \(K\)).
Moreover,
\[
  \dim K=\dim\mathfrak k=\dim\mathfrak g=\dim G,
\]
and \(K\cap H\) is finite because its Lie algebra is zero. Thus the
canonical morphism
\[
  \operatorname{pr}_1:K\longrightarrow G
\]
is finite and dominant; because \(G\) is connected, it is surjective,
and hence an isogeny. Since \(G\) is simply connected and \(K\) is
semisimple, \SGARefLink{sga3:XXI:6.2.7}{Exposé XXI, 6.2.7} shows that \(\operatorname{pr}_1\) is an
isomorphism. Consequently \(K\) is the graph of a \(k\)-group
homomorphism \(u:G\to H\) with \(\Lie(u)=v\).

Finally, let \(u:G\to H\) be a \(k\)-group homomorphism. The tangent
space at \(u\) to
\[
  \underline{\operatorname{Hom}}_{k\text{-}\mathrm{gr.}}(G,H)
\]
identifies with \(Z^1(G,\mathfrak h)\) (\SGARefLink{sga3:ii:subsection:4-2}{Exposé II, 4.2}), where
\(G\) acts on \(\mathfrak h\) through
\(\operatorname{Ad}_H\circ u\). Similarly, the tangent space to
\[
  \underline{\operatorname{Hom}}_{k\text{-}\mathrm{Lie}}
    (\mathfrak g,\mathfrak h)
\]
at \(\Lie(u)\) identifies with \(Z^1(\mathfrak g,\mathfrak h)\). We
must show that
\[
  Z^1(G,\mathfrak h)\longrightarrow
  Z^1(\mathfrak g,\mathfrak h)
\]
is bijective.

\[
  % BEGIN INCLUDED SGA3_English_current_buildable_source_20260823/source/inputs/XXIV/native_diagrams/exp24/010_group_lie_cohomology_rows.tex
\begin{tikzcd}[column sep=small,row sep=large]
0 \arrow[r]
&
\mathfrak h^G
  \arrow[r]
  \arrow[d]
&
\mathfrak h
  \arrow[r]
  \arrow[d,equal]
&
Z^1(G,\mathfrak h)
  \arrow[r]
  \arrow[d]
&
H^1(G,\mathfrak h)
  \arrow[d]
\\
0 \arrow[r]
&
\mathfrak h^{\mathfrak g}
  \arrow[r]
&
\mathfrak h
  \arrow[r]
&
Z^1(\mathfrak g,\mathfrak h)
  \arrow[r]
&
H^1(\mathfrak g,\mathfrak h)
\end{tikzcd}

% END INCLUDED SGA3_English_current_buildable_source_20260823/source/inputs/XXIV/native_diagrams/exp24/010_group_lie_cohomology_rows.tex

\]

By \SGARefLink{sga3:xxiv:lemma:7-3-2}{7.3.2} and \SGARefLink{sga3:xxiv:remark:7-3-3}{7.3.3},
\[
  \mathfrak h^G=\mathfrak h^{\mathfrak g},
  \qquad
  H^1(G,\mathfrak h)=0=H^1(\mathfrak g,\mathfrak h),
\]
which gives the desired conclusion.

\par\medskip\noindent\SGAInlineRefTarget{sga3:xxiv:remark:7-3-6}\textbf{Remark 7.3.6.---}\enspace
It is presumably true that, if \(S\) is locally noetherian, \(G\) a
simply connected semisimple \(S\)-group, \(H\) a smooth \(S\)-group
scheme, and \(\widehat G,\widehat H\) their formal completions along
the identity section, then every homomorphism
\(\widehat G\to\widehat H\) comes from a unique homomorphism \(G\to H\).
This would generalize \SGARefLink{sga3:xxiv:proposition:7-3-1:item:iii}{7.3.1(iii)}. When \(S=\operatorname{Spec}(k)\)
and \(H\) is affine and of finite type, this follows from a theorem of
Dieudonné \([{\rm Di57}]\), \S15, Theorem~4.\footnote{\textit{Editors'
note.} See Exposé VII\textsubscript{B} for the definition of the formal
groups \(\widehat G\) and \(\widehat H\). Suppose that
\(S=\operatorname{Spec}(k)\), where \(k\) is a field. By loc. cit.,
2.2.1, giving a morphism of formal \(k\)-groups
\(v:\widehat G\to\widehat H\) is equivalent to giving a morphism of
Hopf \(k\)-algebras
\(\phi:\operatorname{Dist}(G)\to\operatorname{Dist}(H)\), where
\(\operatorname{Dist}(G)\) denotes the algebra of distributions (at
the origin) of \(G\); see Exposé VII\textsubscript{A}, \SGARefLink{sga3:VIIA:2.1}{2.1},
\([{\rm DG70}]\), \S II.4, \SGARefLink{sga3:xxiv:lemma:6-1}{6.1}, or \([{\rm Ja87}]\), I 7.7. It is called
the hyperalgebra of \(G\) in \([{\rm Di57}]\) and \([{\rm Ta75}]\).
Theorem~4 of \([{\rm Di57}]\) (see also \([{\rm Ta75}]\), 0.3.4 (f)
and (g)) generalizes in this context the theorem of Chevalley used in
\SGARefLink{sga3:xxiv:subsubsection:7-3-5}{7.3.5}: one obtains a closed connected \(k\)-subgroup
\(K\subset G\times H\) such that \(\widehat K\) equals the graph of
\(v\). Since \(\operatorname{Dist}(K)\to\operatorname{Dist}(G)\) is an
isomorphism, \(K\to G\) is finite étale. It follows that \(K\) is
semisimple and then, by \SGARefLink{sga3:XXI:6.2.7}{Exposé XXI, 6.2.7}, that \(K\to G\) is an
isomorphism (because \(G\) is simply connected; see also
\([{\rm Ta75}]\), \SGARefLink{sga3:xxiv:1.8}{1.8} and \SGARefLink{sga3:xxiv:2.2}{2.2}). More generally, loc. cit. studies the
\(k\)-groups \(G\) having property \textup{(SC)}: every finite étale
morphism of \(k\)-groups \(G'\to G\), with \(G'\) connected, is an
isomorphism. Finally, the preceding argument shows that every
finite-dimensional \(\operatorname{Dist}(G)\)-module is, uniquely, a
\(G\)-module; for an extension to a split reductive \(k\)-group \(G\)
(or even to a Borel subgroup of \(G\)), see \([{\rm Ja87}]\), II 1.20
(and the references given there).}

\subsection*{\SGAInlineRefTarget{sga3:xxiv:subsection:7-4}7.4. An example}

As an example, we shall determine
\[
  \underline{\operatorname{Hom}}_{S\text{-}\mathrm{gr.}}
    (\mathrm{SL}_{2,S},\mathrm{SL}_{2,S}).
\]

\subsubsection*{7.4.1.}

Recall (Exposé XX, no.~5) that \(\mathrm{SL}_{2,S}\) is the
\(S\)-group scheme of matrices
\[
  \begin{pmatrix}a&b\\c&d\end{pmatrix}
\]
over \(S\) satisfying \(ad-bc=1\). A maximal torus \(T\) is defined
by the monomorphism
\[
  \alpha^*:\mathbb G_{m,S}\longrightarrow\mathrm{SL}_{2,S},
  \qquad
  \alpha^*(z)=
  \begin{pmatrix}z&0\\0&z^{-1}\end{pmatrix}.
\]
A root of \(G\) with respect to \(T\) is defined by
\[
  \alpha\bigl(\alpha^*(z)\bigr)=z^2,
\]
and a corresponding monomorphism is
\[
  p:\mathbb G_{a,S}\longrightarrow\mathrm{SL}_{2,S},
  \qquad
  p(x)=
  \begin{pmatrix}1&x\\0&1\end{pmatrix}.
\]
Finally, putting \(u=p(1)\), the corresponding representative of the
Weyl group is
\[
  w=
  \begin{pmatrix}0&1\\-1&0\end{pmatrix}.
\]

Recall also (\SGARefLink{sga3:xx:6.2}{Exposé XX, 6.2}) that \(\mathrm{SL}_{2,S}\) is generated by
\(T\), \(p(\mathbb G_{a,S})\), and \(w\), subject to the relations
\[
  \alpha^*(z)p(x)\alpha^*(z^{-1})=p(z^2x),
\]
\[
  w\alpha^*(z)w^{-1}=\alpha^*(z^{-1}),
  \qquad
  w^2=\alpha^*(-1),
  \qquad
  (wu)^3=1.
\]

\subsubsection*{7.4.2.}

Let \(f\) be an endomorphism of \(\mathrm{SL}_{2,S}\). It first
defines a homomorphism
\[
  f\circ\alpha^*:\mathbb G_{m,S}\longrightarrow\mathrm{SL}_{2,S}.
\]
The closed subgroup
\(\ker(f\circ\alpha^*)\subset\mathbb G_{m,S}\) is, locally on \(S\),
either \(\mu_{n,S}\), with \(n\geq1\), or
\(\mathbb G_{m,S}\). After restricting \(S\), we may therefore assume
that there are an integer \(n\in\mathbb N\) and a monomorphism
\[
  f':\mathbb G_{m,S}\longrightarrow\mathrm{SL}_{2,S}
\]
such that
\[
  f\circ\alpha^*(z)=f'(z^n).
\]
By conjugacy of monomorphisms
\(\mathbb G_{m,S}\to\mathrm{SL}_{2,S}\), after an étale covering
extension of the base there is a section \(g\) of \(G\) such that
\[
  f\circ\alpha^*(z)
  =
  \operatorname{int}(g)\circ\alpha^*(z^n).
\]
Replacing \(f\) by an inner conjugate, we are thus reduced to the case
\[
  f\circ\alpha^*(z)=\alpha^*(z^n)
\]
for some \(n\in\mathbb N\).

\subsubsection*{7.4.3.}

Consider now
\[
  f\circ p:\mathbb G_{a,S}\longrightarrow\mathrm{SL}_{2,S}.
\]
It satisfies
\[
  \alpha^*(z^n)(f\circ p)(x)\alpha^*(z^n)^{-1}
  =
  (f\circ p)(z^2x).
\]
If \(n=0\), it follows at once that \(f\circ p\) is invariant under
the homotheties of \(\mathbb G_{a,S}\), and hence is constant. If
\(n\neq0\), \SGARefLink{sga3:xxii:prop:4.1.9}{Exposé XXII, 4.1.9} applies: the map
\(x\mapsto x^n\) is an endomorphism of \(\mathbb G_{a,S}\), and there
is a \(\lambda\in\mathbb G_{m,S}\) such that
\[
  f\circ p(x)=p(\lambda x^n).
\]
This relation is also valid for \(n=0\), on taking \(\lambda=1\).
After extending \(S\) once more, choose
\(z\in\mathbb G_{m,S}\) with \(z^2=\lambda\). Replacing \(f\) by
\(\operatorname{int}(\alpha^*(z))\circ f\), we are reduced to
\[
  f\circ\alpha^*(z)=\alpha^*(z^n),
  \qquad
  f\circ p(x)=p(x^n).
\]

\subsubsection*{7.4.4.}

Finally, one must have
\[
  f(w)\alpha^*(z^n)f(w)^{-1}=\alpha^*(z^n)^{-1},
  \qquad
  (f(w)u)^3=1.
\]
By \SGARefLink{sga3:xx:3.8}{Exposé XX, 3.8}, this implies \(f(w)=w^n\). Since \(G\) is generated
by \(T\), \(p(\mathbb G_{a,S})\), and \(w\), it follows that, for every
\(S'\to S\) and every
\(\begin{pmatrix}a&b\\c&d\end{pmatrix}\in G(S')\),
\[
  f\!\begin{pmatrix}a&b\\c&d\end{pmatrix}
  =
  \begin{pmatrix}a^n&b^n\\c^n&d^n\end{pmatrix}.
\]

\subsubsection*{7.4.5.}

Summarizing, locally on \(S\) for the étale topology, every
endomorphism \(f\) of \(\mathrm{SL}_{2,S}\) can be written
\[
  f=\phi\circ F_n,
\]
where \(\phi\) is an inner automorphism, \(n\geq0\), and
\[
  F_n\!\begin{pmatrix}a&b\\c&d\end{pmatrix}
  =
  \begin{pmatrix}a^n&b^n\\c^n&d^n\end{pmatrix}.
\]
If \(f=\phi\circ F_n\), the integer \(n\) is determined by any fiber
\(f_s\), for example by \(\ker(f_s)\). Consequently \(n\) is a
locally constant function on \(S\).

\subsubsection*{7.4.6.}

It follows that, for every endomorphism \(f\) of
\(\mathrm{SL}_{2,S}\), the scheme \(S\) has a canonical decomposition
into open and closed subschemes
\[
  S_0,\qquad S_1,\qquad S_{p^n},
\]
where \(p^n\) runs through the positive powers of prime numbers, such
that:
\begin{enumerate}
\item[(i)] \(f_{S_0}\) is the zero morphism;
\item[(ii)] \(f_{S_1}\) is an isomorphism, necessarily an inner
  automorphism;
\item[(iii)] \(S_{p^n}\) has characteristic \(p\), and
  \(f_{S_{p^n}}\) has a unique decomposition
  \[
    f_{S_{p^n}}=\phi\circ F_p^n,
  \]
  where \(\phi\) is an inner automorphism and \(F_p\) is the Frobenius
  endomorphism of \(\mathrm{SL}_{2,\mathbb F_p}\).
\end{enumerate}

\subsubsection*{7.4.7.}

Equivalently,
\[
  \underline{\operatorname{Hom}}_{\mathbb Z\text{-}\mathrm{gr.}}
    (\mathrm{SL}_{2,\mathbb Z},\mathrm{SL}_{2,\mathbb Z})
\]
is the disjoint-union scheme consisting of:
\begin{enumerate}
\item[(i)] one scheme isomorphic to \(\operatorname{Spec}(\mathbb Z)\);
\item[(ii)] one scheme isomorphic to
  \[
    \underline{\operatorname{Aut}}_{\mathbb Z\text{-}\mathrm{gr.}}
      (\mathrm{SL}_{2,\mathbb Z})
    \simeq
    \operatorname{ad}(\mathrm{SL}_{2,\mathbb Z});
  \]
\item[(iii)] for each prime \(p\) and each integer \(n>0\), one scheme
  isomorphic to
  \[
    \underline{\operatorname{Aut}}_{\mathbb F_p\text{-}\mathrm{gr.}}
      (\mathrm{SL}_{2,\mathbb F_p})
    \simeq
    \operatorname{ad}(\mathrm{SL}_{2,\mathbb F_p}).
  \]
\end{enumerate}

\subsubsection*{7.4.8.}

In particular:
\begin{enumerate}
\item[(i)]
\(\underline{\operatorname{Hom}}_{\mathbb F_p\text{-}\mathrm{gr.}}
(\mathrm{SL}_{2,\mathbb F_p},\mathrm{SL}_{2,\mathbb F_p})\)
has infinitely many connected components and is therefore not
quasi-compact.
\item[(ii)] If \(S\) is a scheme of unequal characteristics, then
\(\underline{\operatorname{Hom}}_{S\text{-}\mathrm{gr.}}
(\mathrm{SL}_{2,S},\mathrm{SL}_{2,S})\) is not flat over \(S\).
\end{enumerate}

\section*{\SGAInlineRefTarget{sga3:xxiv:section:8}8. Appendix: Cohomology of a smooth group\\
over a henselian ring. Cohomology and the functor \(\prod\)}

\par\medskip\noindent\SGAInlineRefTarget{sga3:xxiv:proposition:8-1}\textbf{Proposition 8.1.---}\enspace
\emph{Let \(S\) be a henselian local scheme, \(s\) its closed point,
and \(G\) a smooth \(S\)-group scheme such that every finite subset of
\(G\) is contained in an affine open subset.%
\begingroup
\renewcommand{\thefootnote}{*}%
\footnote{This last condition is in fact unnecessary; see
A.~Grothendieck, \textit{Le groupe de Brauer III}, in
\textit{Dix Exposés sur la cohomologie des schémas}, North-Holland,
1968, Theorem~11.7 and Remarks~11.8.3.}%
\endgroup
\addtocounter{footnote}{-1}
\begin{enumerate}
\item[(i)] If \(P\) is a principal homogeneous bundle under \(G\),
  there is a finite surjective étale morphism \(S'\to S\) which
  trivializes \(P\). Moreover,
  \[
    \operatorname{Fib}(S,G)\simeq H^1_{\acute et}(S,G).
  \]
\item[\SGAInlineRefTarget{sga3:xxiv:proposition:8-1:item:ii}(ii)] For every finite surjective étale morphism \(S'\to S\),
  the canonical map
  \[
    H^1(S'/S,G)
    \longrightarrow
    H^1(S'\otimes_S\kappa(s)/\kappa(s),G_s)
  \]
  is bijective.
\item[(iii)] The canonical map
  \[
    \operatorname{Fib}(S,G)
    \longrightarrow
    \operatorname{Fib}(\kappa(s),G_s)
  \]
  is bijective.
\end{enumerate}}

\subsubsection*{8.1.2.}

If \(K\) is a finite separable extension of \(\kappa(s)\), there is a
finite surjective étale morphism \(S'\to S\) such that
\[
  K\simeq S'\otimes_S\kappa(s).
\]
\footnote{\textit{Editors' note.} The scheme \(S'\) is still local and
henselian. We have also retained the numbering of the original: there
is no no.~8.1.1.}
If \(P\) is a principal homogeneous bundle under \(G\), then \(P\) is
smooth over \(S\), hence \(P_s\) is smooth over \(\kappa(s)\).
Consequently there is a finite separable extension \(K/\kappa(s)\)
such that \(P_K\) has a section (EGA IV\textsubscript{4}, 17.15.10).
Representing \(K\) as above, Hensel's lemma
(EGA IV\textsubscript{4}, 18.5.17) shows that \(P_{S'}\) has a section.
This proves the first assertion of (i).

Conversely, if \(P\) is a principal homogeneous sheaf under \(G\) for
the étale topology, then a finite surjective étale morphism
\(S'\to S\) trivializes \(P\): every étale covering family of a
henselian local scheme is dominated by a family consisting of one
morphism of that form. By the descent hypothesis imposed on \(G\),
\(P\) is representable (SGA~1, VIII, 7.6). This completes the proof of
(i), and shows that (ii) implies (iii). It remains to prove (ii).

\subsubsection*{8.1.3.}

The map in (ii) is injective. Let \(P\) and \(Q\) be two principal
homogeneous bundles under \(G\), trivialized by \(S'\), and consider
the \(S\)-group sheaf
\[
  \mathcal H=
  \underline{\operatorname{Isom}}_{G\text{-}\mathrm{bundles}}(P,Q).
\]
Since \(\mathcal H_{S'}\) is isomorphic to \(G_{S'}\),
\(\mathcal H\) is representable by the second hypothesis on \(G\), as
above. If \(\mathcal H(\kappa(s))\neq\varnothing\), Hensel's lemma
gives \(\mathcal H(S)\neq\varnothing\), and hence \(P\simeq Q\).

\subsubsection*{8.1.4.}

To prove surjectivity in \SGARefLink{sga3:xxiv:proposition:8-1:item:ii}{(ii)}, it is equivalent to prove that
\[
  Z^1(S'/S,G)
  \longrightarrow
  Z^1(S'\otimes_S\kappa(s)/\kappa(s),G_s)
\]
is surjective. Let \(Z^1\) be the \(S\)-functor defined by
\[
  Z^1(T)=Z^1(S'\times_ST/T,G_T).
\]
The preceding map is then
\[
  Z^1(S)\longrightarrow Z^1(\kappa(s)).
\]
Another application of Hensel's lemma reduces the assertion to
showing that \(Z^1\) is representable by a smooth \(S\)-scheme.

\subsubsection*{8.1.5.}

We first show that \(Z^1\) is representable by an \(S\)-scheme locally
of finite presentation. For \(i=0,1,\ldots\), define the \(S\)-functor
\(C_i\) by
\[
  C_i(T)=C_i(S'\times_ST/T,G_T).
\]
Equivalently,
\[
  C_i(T)
  =
  G\bigl((S'/S)^{i+1}\times_ST\bigr),
  \qquad
  C_i
  =
  \underline{\operatorname{Hom}}_S((S'/S)^{i+1},G).
\]
Since \(Z^1\) is obtained from \(C_1\) and \(C_2\) by fiber products,
it is enough to prove that \(C_i\), for \(i=1,2\), is representable by
an \(S\)-scheme locally of finite presentation.

\subsubsection*{8.1.6.}

Let \(T\to S\) be a finite surjective étale morphism which splits
\(S'\). Then
\[
  C_i\times_ST
  =
  \underline{\operatorname{Hom}}_T
    ((S'\times_ST/T)^{i+1},G_T)
\]
is represented by a product of \(n\) copies of \(G_T\), where \(n\)
is the degree of \((S'/S)^{i+1}\). The hypothesis on \(G\), applied
once again, shows that \(C_i\) is represented by an \(S\)-scheme
locally of finite presentation (SGA~1, VIII, \textit{loc. cit.}).

\subsubsection*{8.1.7.}

It remains to prove that \(Z^1\) is smooth. By definition, if \(T\) is
an affine \(S\)-scheme and \(T_0\) is the closed subscheme defined by
a square-zero ideal \(J\), we must show that
\[
  Z^1(T)\longrightarrow Z^1(T_0)
\]
is surjective. Since \(G\) is smooth, the canonical map
\(G(T)\to G(T_0)\) is surjective, and it is enough to prove that
\[
  H^1(S'\times_ST/T,G_T)
  \longrightarrow
  H^1(S'\times_ST_0/T_0,G_{T_0})
\]
is bijective. After a slight change of notation and a harmless
generalization of the hypotheses, this reduces to the following
lemma.

\par\medskip\noindent\SGAInlineRefTarget{sga3:xxiv:lemma:8-1-8}\textbf{Lemma 8.1.8.---}\enspace
\emph{Let \(S\) and \(S'\) be affine schemes, let \(S'\to S\) be
faithfully flat, let \(J\) be a square-zero ideal on \(S\), and let
\(S_0\) be the closed subscheme which it defines. Let \(G\) be a
smooth \(S\)-group. Put
\[
  S'_0=S'\times_SS_0,
  \qquad
  G_0=G\times_SS_0.
\]
Then the canonical map
\[
  H^1(S'/S,G)\longrightarrow H^1(S'_0/S_0,G_0)
\]
is bijective.}

\par\medskip\noindent\SGAInlineRefTarget{sga3:xxiv:remark:8-1-9}\textbf{Remark 8.1.9.---}\enspace
If \(G\) is commutative, the same assertion holds for every
\(H^i\), \(i>0\), with an analogous proof.

\smallskip\noindent\emph{Proof.}
Let \(M_0\) be the quasi-coherent \(\mathcal O_{S_0}\)-module
\[
  M_0
  =
  \Lie(G_0/S_0)\otimes_{\mathcal O_{S_0}}J.
\]
For every \(S_0\)-scheme \(T_0\), put
\[
  M_0(T_0)
  =
  H^0\!\left(
    T_0,
    M_0\otimes_{\mathcal O_{S_0}}\mathcal O_{T_0}
  \right),
\]
and define the \(S\)-group functors
\[
  M=\prod_{S_0/S}M_0,
  \qquad
  \overline G=\prod_{S_0/S}G_0
\]
by
\[
  M(T)=M_0(T_0),
  \qquad
  \overline G(T)=G_0(T_0),
  \qquad
  T_0=T\times_SS_0.
\]
By \SGARefLink{sga3:iii:corollary:0-9}{Exposé III, 0.9} and (\SGARefLink{sga3:iii:remark:0-6:equation:0-6-2}{0.6.2}), for every affine \(S\)-scheme \(T\)
there is a functorial exact sequence
\[
  1\longrightarrow M(T)\longrightarrow G(T)
  \longrightarrow\overline G(T)\longrightarrow1.
\]
We must study the map
\[
  H^1(S'/S,G)
  \longrightarrow
  H^1(S'_0/S_0,G_0)
  =
  H^1(S'/S,\overline G).
\]

Suppose first that \(G\) is commutative. There is then an exact
cohomology sequence
\[
  H^1(S'/S,M)
  \longrightarrow H^1(S'/S,G)
  \longrightarrow H^1(S'/S,\overline G)
  \longrightarrow H^2(S'/S,M).
\]
Moreover,
\[
  H^i(S'/S,M)
  =
  H^i(S'_0/S_0,M_0)
  =
  H^i(S'_0/S_0,\mathcal M_0)
  =
  0
  \qquad(i\neq0),
\]
the last equality following from TDTE I, B, Lemma~1.1.

Now suppose that \(G\) is not commutative. We use the exact sequence
of nonabelian cohomology. If
\(u\in Z^1(S'/S,\overline G)\), the elements of \(H^1(S'/S,G)\) whose
image in \(H^1(S'/S,\overline G)\) equals the class of \(u\) lie in the
image of the corresponding coboundary map
\[
  H^1(S'/S,M_u)\longrightarrow H^1(S'/S,G),
\]
where \(M_u\) is \(M\) twisted by \(u\). Similarly, for
\(v\in Z^1(S'/S,\overline G)\), there is an obstruction
\[
  \Delta(v)\in H^2(S'/S,M_v)
\]
which vanishes exactly when the class of \(v\) is in the image of
\(H^1(S'/S,G)\). It therefore suffices to prove
\[
  H^1(S'/S,M_u)=0,
  \qquad
  H^2(S'/S,M_v)=0.
\]

Recall (\SGARefLink{sga3:iii:proposition:0-8}{Exposé III, 0.8}) that the action of \(G\) on \(M\) defined by
the exact sequence above is induced functorially by the adjoint
representation
\[
  \operatorname{ad}:G_0
  \longrightarrow
  \underline{\operatorname{Aut}}_{\mathcal O_{S_0}}
    (\Lie(G_0/S_0)).
\]
The cocycle \(u\), respectively \(v\), therefore defines a descent
datum on \(\Lie(G_0/S_0)\); write \(L_u\), respectively \(L_v\), for
the resulting quasi-coherent \(\mathcal O_{S_0}\)-module. For
\(T\to S\),
\[
  M_u(T)
  =
  H^0\!\left(
    T_0,
    L_u\otimes_{\mathcal O_{S_0}}J
      \otimes_{\mathcal O_{S_0}}\mathcal O_{T_0}
  \right),
\]
and similarly with \(v\) in place of \(u\). Hence
\[
  H^1(S'/S,M_u)=H^1(S'_0/S_0,L_u\otimes J),
\]
\[
  H^2(S'/S,M_v)=H^2(S'_0/S_0,L_v\otimes J),
\]
and both groups vanish by the result already used.
\hfill\(\square\)

\par\medskip\noindent\SGAInlineRefTarget{sga3:xxiv:proposition:8-2}\textbf{Proposition 8.2.---}\enspace
\emph{Let \(\mathcal C\) be a category with fiber products, endowed
with a topology coarser than the canonical topology. Let \(S'\to S\)
be a morphism of \(\mathcal C\), let \(G'\) be an \(S'\)-group sheaf,
and let
\[
  G=\prod_{S'/S}G'
\]
be the corresponding \(S\)-group sheaf. Let
\[
  H^1_S(S',G')\subset H^1(S',G')
\]
be the subset of classes of principal homogeneous sheaves under
\(G'\) which are trivialized by a sieve on \(S'\) obtained by base
change from a suitable covering sieve on \(S\). The canonical map
\[
  H^1(S,G)\longrightarrow H^1(S',G'),
  \qquad
  P\longmapsto P\times_SS',
\]
induces a bijection
\[
  H^1(S,G)\xrightarrow{\sim}H^1_S(S',G').
\]
Its inverse is induced by
\[
  P'\longmapsto\prod_{S'/S}P'.
\]}

For every object \(X\) of \(\mathcal C/S\), the definition gives an
isomorphism, functorial in \(X\),
\[
  \operatorname{Hom}_S(X,G)
  \xrightarrow{\sim}
  \operatorname{Hom}_{S'}(X\times_SS',G').
\]
Thus, for every \(S\)-object \(T\), there is a bijection functorial in
\(T\),
\[
  H^1(T/S,G)
  \xrightarrow{\sim}
  H^1(T'/S',G'),
  \qquad
  T'=T\times_SS'.
\]
Replacing the single morphism \(T\to S\) by an arbitrary covering
family of \(S\), and then passing to the inductive limit, gives the
first part of the proposition. The second follows immediately.

\par\medskip\noindent\SGAInlineRefTarget{sga3:xxiv:lemma:8-3}\textbf{Lemma 8.3.---}\enspace
\emph{Under the hypotheses of \SGARefLink{sga3:xxiv:proposition:8-2}{8.2}, the equality
\[
  H^1_S(S',G')=H^1(S',G')
\]
is local on \(S\). More precisely, suppose that there is a covering
family \(\{S_i\to S\}\) such that, for every \(i\),
\[
  H^1_{S_i}(S'\times_SS_i,G')
  =
  H^1(S'\times_SS_i,G').
\]
Then \(H^1_S(S',G')=H^1(S',G')\).}

Indeed, let \(P'\) be a principal homogeneous sheaf under \(G'\), and
put
\[
  P'_i=P'\times_{S'}(S'\times_SS_i).
\]
By hypothesis, there is a covering family
\(\{S_{ij}\to S_i\}\) such that, for each \(j\),
\[
  P'\times_{S'}(S'\times_SS_{ij})
\]
has a section. But \(\{S_{ij}\to S\}\) is a covering family of \(S\),
and \(P'\) is therefore trivialized by the covering family of \(S'\)
obtained from it by base change.

\par\medskip\noindent\SGAInlineRefTarget{sga3:xxiv:proposition:8-4}\textbf{Proposition 8.4.---}\enspace
\emph{Let \(S'\to S\) be a finite étale morphism of schemes. Let
\(G'\) be an \(S'\)-group sheaf and let
\[
  G=\prod_{S'/S}G'
\]
be the corresponding \(S\)-group sheaf. For the étale topology,
respectively the local finite étale topology, respectively the
\textup{(fpqc)} topology, the following functors induce mutually
inverse bijections.}

\[
  % BEGIN INCLUDED SGA3_English_current_buildable_source_20260823/source/inputs/XXIV/native_diagrams/exp24/011_weil_restriction_torsor_functors.tex
\begin{tikzcd}[column sep=huge,row sep=large]
\mathcal P
  \arrow[r,mapsto]
&
\mathcal P\times_S S'
\\
\displaystyle\prod_{S'/S}\mathcal P'
&
\mathcal P'
  \arrow[l,mapsto]
\end{tikzcd}

% END INCLUDED SGA3_English_current_buildable_source_20260823/source/inputs/XXIV/native_diagrams/exp24/011_weil_restriction_torsor_functors.tex

\]

\emph{In particular,
\[
  H^1(S,G)\simeq H^1(S',G').
\]}

By \SGARefLink{sga3:xxiv:proposition:8-2}{8.2}, it is enough to show
\[
  H^1_S(S',G')=H^1(S',G').
\]
By \SGARefLink{sga3:xxiv:lemma:8-3}{8.3}, this may be checked locally for the local finite étale
topology. We may therefore suppose that \(S'\) is a finite disjoint
union of copies of \(S\), say \(I_S\), where \(I\) is a finite set.
Then \(G'\) is given by a family \((G_i)_{i\in I}\) of sheaves on
\(S\), and
\[
  H^1(S',G')
  \simeq
  \prod_{i\in I}H^1(S,G_i).
\]
On the other hand,
\[
  H^1(S,G)
  \simeq
  \prod_{i\in I}H^1(S,G_i).
\]
It follows from \SGARefLink{sga3:xxiv:proposition:8-2}{8.2} that
\(H^1_S(S',G')=H^1(S',G')\).
\hfill\(\square\)

\par\medskip\noindent\SGAInlineRefTarget{sga3:xxiv:remark:8-5}\textbf{Remarks 8.5.---}\enspace
\SGARefLink{sga3:xxiv:proposition:8-2}{Propositions 8.2} and \SGARefLink{sga3:xxiv:proposition:8-4}{8.4} may be interpreted through the following
exact sequence, where \(f:S'\to S\) is the given morphism:
\[
  1\longrightarrow H^1(S,f_*G')
  \longrightarrow H^1(S',G')
  \longrightarrow H^0(S,R^1f_*G').
\]
In the commutative case, this exact sequence follows from the Leray
spectral sequence. It remains valid in the noncommutative
case.\footnote{\textit{Editors' note.} See \S V.3.1.4 of
\([{\rm Gi71}]\).}

In this form one sees that the result remains valid when \(f\) is
assumed only finite, or even merely integral, provided that the
topology is the étale topology. Indeed, for every \(G'\),
\[
  R^1f_*(G')=\text{the final sheaf}
\]
by SGA~4, VIII, 5.3.\footnote{\textit{Editors' note.} This reference
also points to \S V.3.1.4 of \([{\rm Gi71}]\).}

The result becomes false for a topology such as \textup{(fpqc)} or
\textup{(fppf)}, even if
\[
  S=\operatorname{Spec}(k),
  \qquad
  S'=\operatorname{Spec}(k[t]/t^2),
\]
where \(k\) is an algebraically closed field of characteristic
\(p\neq0\), and \(G'=\mu_p\) or \(\alpha_p\).

Likewise, 8.2 becomes false, even for the étale topology, if the
finiteness hypothesis on \(f\) is omitted. For example, take \(f\)
to be an open immersion with
\[
  S=\operatorname{Spec}(V),
\]
where \(V\) is a complete discrete valuation ring with algebraically
closed residue field, take \(S'\) to be the open subscheme consisting
of the generic point, and take \(G'\) to be the constant group
\((\mathbb Z/n\mathbb Z)_{S'}\), where \(n\) is prime to the residue
characteristic of \(V\). Then
\[
  H^1(S,G)=0,
  \qquad
  H^1(S',G')\neq0.
\]
Replacing \(S'\) by \(S\amalg S'\) gives an analogous example in which
\(S'\to S\) is étale and surjective, hence a cover for the topology
under consideration.

\section*{Bibliography}

\noindent[BLie]\enspace
N.~Bourbaki, \textit{Groupes et algèbres de Lie}, Chap.~I, Hermann,
1960.\par

\smallskip\noindent[Ch51]\enspace
C.~Chevalley, \textit{Théorie des groupes de Lie}, t.~II,
\textit{Groupes algébriques}, Hermann, 1951.\par

\smallskip\noindent[Di57]\enspace
J.~Dieudonné, Lie groups and Lie hyperalgebras over a field of
characteristic \(p>0\). VI, \textit{Amer. J. Math.} \textbf{79}
(1957), no.~2, 331--388.\footnote{\textit{Editors' note.} The editors
have added, after these three references which appeared in the
original, the references that follow.}\par

\smallskip\noindent[Bo91]\enspace
A.~Borel, \textit{Linear Algebraic Groups}, 2nd ed.,
Springer-Verlag, 1991.\par

\smallskip\noindent[Ch05]\enspace
C.~Chevalley, \textit{Classification des groupes algébriques
semi-simples}, with the collaboration of P.~Cartier,
A.~Grothendieck, and M.~Lazard, \textit{Collected Works}, vol.~3,
Springer, 2005.\par

\smallskip\noindent[DG70]\enspace
M.~Demazure and P.~Gabriel, \textit{Groupes algébriques}, Masson and
North-Holland, 1970.\par

\smallskip\noindent[Gi71]\enspace
J.~Giraud, \textit{Cohomologie non abélienne}, Springer-Verlag,
1971.\par

\smallskip\noindent[Ja87]\enspace
J.~C.~Jantzen, \textit{Representations of Algebraic Groups}, Academic
Press, 1987; 2nd ed., American Mathematical Society, 2003.\par

\smallskip\noindent[Jou83]\enspace
J.-P.~Jouanolou, \textit{Théorèmes de Bertini et applications},
Birkhäuser, 1983.\par

\smallskip\noindent[Ta75]\enspace
M.~Takeuchi, On coverings and hyperalgebras of affine algebraic
groups, \textit{Trans. Amer. Math. Soc.} \textbf{211} (1975),
249--275.\par

% HARD STOP: end of Exposé XXIV authority local page 54.

% END INCLUDED SGA3_English_current_buildable_source_20260823/source/inputs/XXIV/tex/components/13_expose_XXIV_remark631_continuation_through_bibliography_en.tex



% END INCLUDED SGA3_English_current_buildable_source_20260823/source/inputs/XXIV/body.tex


\SGAStartExpose{XXV}
% BEGIN INCLUDED SGA3_English_current_buildable_source_20260823/source/inputs/XXV/body.tex
% Mechanically extracted after \mainmatter from the frozen standalone master.

\markboth{Exposé XXV. The Existence Theorem}{Exposé XXV. The Existence Theorem}
% BEGIN INCLUDED SGA3_English_current_buildable_source_20260823/source/inputs/XXV/tex/components/00_title_introduction_and_section1.tex
\chapter*{Exposé XXV\\The Existence Theorem}
\addcontentsline{toc}{chapter}{Exposé XXV. The Existence Theorem}
\begin{center}
\textit{by M. Demazure}
\end{center}

\footnote{Editors' note: Version of 13 October 2024.}
For completeness, in this Exposé we give a proof of the existence theorem for split
groups. Like Chevalley's original proof (``Sur certains schémas de groupes
semi-simples'', \textit{Séminaire Bourbaki}, May 1961, no.~219), it rests on the
existence of complex semisimple algebraic groups of every possible type. Cartier
gave the principle of a more satisfactory proof, establishing directly the
existence of a simply connected split semisimple \(\mathbb Z\)-group associated
with a given Cartan matrix (unpublished).\footnote{Editors' note: Such a principle
was also sketched by B.~Kostant \emph{[Ko66]}; a complete proof, using quantum
groups, was given more recently by G.~Lusztig \emph{[Lu09]}.}
Let us point out, however, that the difficulty is not to give an explicit
construction of a group scheme, but to verify that the group so constructed
really satisfies the required conditions---essentially, that its fibers are
indeed smooth and reductive.

\SGASection{1}{Statement of the theorem}
\SGARefTarget{sga3:xxv:section:1}

\begin{sgastatement}{Theorem 1.1}
Let \(S\) be a nonempty scheme. The functor
\[
  G\longmapsto R(G)
\]
is an equivalence from the category of pinned reductive \(S\)-groups to the
category of pinned reduced root data.
\end{sgastatement}
\SGARefTarget{sga3:xxv:theorem:1:1}

By the uniqueness theorem (Exposé XXIII, 4.1), the preceding statement is
equivalent to the following.

\begin{sgastatement}{\SGAInlineRefTarget{sga3:xxv:corollary:1-2}Corollary 1.2 (Existence of split groups)}
For every reduced root datum \(R\), there exists a split reductive
\(\mathbb Z\)-group \(G\) such that \(R(G)\simeq R\).\footnote{Editors' note:
The work \emph{[BT84]} of F.~Bruhat and J.~Tits also contains a variant of
Chevalley's construction (\emph{loc.\ cit.}, 2.2.3--2.2.5 and \S3.2), which
in particular gives a smooth affine \(\mathbb Z\)-group \(G\) with connected
fibers, possessing a split maximal torus, and whose generic fiber is a reductive
\(\mathbb Q\)-group of type \(R\). The reductivity of the geometric fibers of
\(G\) follows from the study of the unipotent radical of a special fiber in
\emph{loc.\ cit.}, 4.6.12 and 4.6.15 (valid for more general \(G\), associated
with a valued root datum); but it is simpler to deduce it from the description
of \(\mathfrak g=\operatorname{Lie}(G)\), from Exposé XIX, 1.12(iii), and from
the existence of the elements \(w_\alpha(X)\) of \SGARefLink{sga3:xx:3.1:item:iv}{Exposé XX, 3.1(iv)}; see
\emph{[BT84]}, 3.2.1.}
\end{sgastatement}
\SGARefTarget{sga3:xxv:corollary:1:2}

In particular:

\begin{sgastatement}{Corollary 1.3}
Let \(k\) be a field. For every splittable reductive \(k\)-group \(G_k\), there
exists a reductive \(\mathbb Z\)-group \(G\) such that
\[
  G\otimes_{\mathbb Z}k\simeq G_k.
\]
\end{sgastatement}
\SGARefTarget{sga3:xxv:corollary:1:3}

Conversely, first observe that to prove \SGARefLink{sga3:xxv:corollary:1-2}{1.2} it is enough, by \SGARefLink{sga3:xxii:prop:4.3.1}{Exposé XXII, 4.3.1},
and \SGARefLink{sga3:XXI:6.5.10}{Exposé XXI, 6.5.10}, to consider the case in which the root datum \(R\) is
simply connected (and even irreducible, if desired, by \SGARefLink{sga3:XXI:7.1.6}{Exposé XXI, 7.1.6}).

On the other hand, under the hypotheses of \SGARefLink{sga3:xxv:corollary:1:3}{1.3} the group scheme \(G\) has
constant type (since \(\operatorname{Spec}(\mathbb Z)\) is connected), hence type
\(R(G_k)\). By \SGARefLink{sga3:XXIII:5.9}{Exposé XXIII, 5.9}, the validity of \SGARefLink{sga3:xxv:corollary:1:3}{1.3} for a given group \(G_k\)
therefore entails the existence of a split \(\mathbb Z\)-group of type
\(R(G_k)\).\footnote{Editors' note: In fact, \SGARefLink{sga3:XXIII:5.9}{Exposé XXIII, 5.9} is unnecessary:
the present Exposé constructs, for every semisimple \(\mathbb Q\)-group
\(G_{\mathbb Q}\), a semisimple \(\mathbb Z\)-group \(G\) of the same type as
\(G_{\mathbb Q}\), equipped with a split maximal torus \(T\); hence \(G\) is
split by \SGARefLink{sga3:xxii:prop:2.2}{Exposé XXII, 2.2}.}

Thus, to prove \SGARefLink{sga3:xxv:corollary:1-2}{1.2} and hence \SGARefLink{sga3:xxv:theorem:1:1}{1.1}, it is enough to prove \SGARefLink{sga3:xxv:corollary:1:3}{1.3} when \(k\) has
characteristic zero (for example, \(k=\mathbb C\)) and \(G_k\) is simply
connected (and in particular semisimple), together with the following
proposition.

\begin{sgastatement}{Proposition 1.4}
For every simply connected reduced root datum \(R\), there exists a semisimple
algebraic \(\mathbb C\)-group of type \(R\).
\end{sgastatement}
\SGARefTarget{sga3:xxv:proposition:1:4}

One may prove \SGARefLink{sga3:xxv:proposition:1:4}{1.4} as follows. First, there exists a complex semisimple Lie
algebra \(\mathfrak g\) of type \(R\); see, for example, N.~Jacobson,
\textit{Lie Algebras}, Chapter VII, Theorem 5.\footnote{Editors' note:
Let \(\Delta\) be a base of \(R\), and let \(\widetilde{\mathfrak g}\) be the
complex Lie algebra generated by
\((e_\alpha,f_\alpha,h_\alpha)_{\alpha\in\Delta}\), subject to the relations
\([h_\alpha,h_\beta]=0\), \([e_\alpha,f_\beta]=h_\alpha\) if
\(\beta=\alpha\) and \(=0\) otherwise,
\([h_\alpha,e_\beta]=(\alpha^\ast,\beta)e_\beta\), and
\([h_\alpha,f_\beta]=-(\alpha^\ast,\beta)f_\beta\).
In \emph{loc.\ cit.}, \(\mathfrak g\) is defined as the quotient of
\(\widetilde{\mathfrak g}\) by the intersection of the kernels of its
finite-dimensional irreducible representations. In \emph{[Se66]},
\S VI.5, Theorem 9 (see also \emph{[BLie]}, VIII, \S4.3, \SGARefLink{sga3:xxv:section:1}{Theorem 1}), it is
shown that \(\mathfrak g\) is the quotient of \(\widetilde{\mathfrak g}\) by
the relations
\[
 \operatorname{ad}(e_\alpha)^{1-(\alpha^\ast,\beta)}(e_\beta)=0,\qquad
 \operatorname{ad}(f_\alpha)^{1-(\alpha^\ast,\beta)}(f_\beta)=0.
\]
For an explicit description of the structure constants, and in particular
the choice of signs, see \SGARefLink{sga3:XXIII:6.5}{Exposé XXIII, 6.5} and 6.7, as well as
\emph{[Ti66]}, \S4, \SGARefLink{sga3:xxv:section:1}{Theorem 1}, and, for types \(A,D,E\),
\emph{[Sp98]}, 10.2.5.}
Then \(G=\operatorname{Aut}(\mathfrak g)^0\) is a semisimple algebraic
\(\mathbb C\)-group of type \(\operatorname{ad}(R)\).\footnote{Editors' note:
One may suppose \(R\) irreducible, hence \(\mathfrak g\) simple. By a general
argument, \(\operatorname{Lie}(G)\) is the complex Lie algebra of derivations
of \(\mathfrak g\); see \emph{[DG70]}, \S II.4, Proposition 2.3. These
derivations are all inner; see \emph{[BLie]}, I, \S6.1, Corollary 3 to
\SGARefLink{sga3:xxv:section:1}{Proposition 1}. Thus \(\operatorname{Lie}(G)=\mathfrak g\), so \(G\) has no
invariant subgroup of positive dimension and is semisimple. Its root system
is consequently the same as that of \(\mathfrak g\). Moreover, the center of
\(G\) acts both trivially and faithfully on \(\mathfrak g\), and is therefore
trivial; hence \(G\) is adjoint.}
By \textit{Bible}, \S23.1, \SGARefLink{sga3:xxv:section:1}{Proposition 1}, it follows that there exists a
semisimple \(\mathbb C\)-group of type \(R\).

The remainder of this Exposé is devoted to the proof of \SGARefLink{sga3:xxv:corollary:1:3}{1.3} when \(k\) has
characteristic zero and \(G_k\) is semisimple. The proof proceeds in two
stages: the construction of a ``piece of \(\mathbb Z\)-group scheme''
(\S2), and the study of the group obtained by applying ``Weil's theorem''
(\S3).

To avoid confusion, in \S2 we shall not use the abbreviated notation \(X_k\)
for the \(k\)-scheme \(X\otimes_{\mathbb Z}k\), where \(X\) is a
\(\mathbb Z\)-scheme.


% END INCLUDED SGA3_English_current_buildable_source_20260823/source/inputs/XXV/tex/components/00_title_introduction_and_section1.tex

% BEGIN INCLUDED SGA3_English_current_buildable_source_20260823/source/inputs/XXV/tex/components/01_section2_1_through_2_5.tex
\SGASection{2}{Existence theorem: construction of a piece of group}
\SGARefTarget{sga3:xxv:section:2}

\SGANextSectionTarget{sga3:xxv:subsection:2-1}\SGASubsection{2.1}{}
\SGARefTarget{sga3:xxv:paragraph:2:1}

Choose once and for all a splitting of \(G_k\), denoted
\[
  (G_k,T_k,M,R)
\]
(see \SGARefLink{sga3:xxii:def:1.13}{Exposé XXII, 1.13}), a system \(\Delta\) of simple roots of \(R\)
(defining the system \(R^+\) of positive roots), and a Chevalley system
\((X_{\alpha,k})_{\alpha\in R}\) of \(G_k\) (\SGARefLink{sga3:XXIII:6.1}{Exposé XXIII, 6.1} and \SGARefLink{sga3:XXIII:6.2}{6.2})
satisfying the additional condition (see \SGARefLink{sga3:xx:2.6}{Exposé XX, 2.6})
\[
  X_{\alpha,k}X_{-\alpha,k}=1\qquad(\alpha\in R).
\]
Finally, choose on the subgroup of \(M\) generated by \(R\) a total order
compatible with the group structure and such that the roots \(>0\) are the
elements of \(R^+\). We then write the roots as
\[
 -\alpha_n<-\alpha_{n-1}<\cdots<-\alpha_1
 <\alpha_1<\alpha_2<\cdots<\alpha_n.
\]
For \(\alpha\in R\), let \(U_{\alpha,k}\) denote the vector group
corresponding to the root \(\alpha\), and let
\[
 p_{\alpha,k}\colon\mathbb G_{a,k}\xrightarrow{\ \sim\ }U_{\alpha,k}
\]
be the isomorphism of vector groups defined by \(X_{\alpha,k}\).

\SGANextSectionTarget{sga3:xxv:subsection:2-2}\SGASubsection{2.2}{}
\SGARefTarget{sga3:xxv:paragraph:2:2}

The splitting of \(G_k\) includes, in particular, an isomorphism of
\(k\)-groups
\[
 T_k\simeq D_k(M).
\]
Put \(T=D(M)\). This is a \(\mathbb Z\)-torus, and the preceding isomorphism
may be regarded as an isomorphism
\[
 T_k\simeq T\otimes_{\mathbb Z}k.
\]
We have
\[
 \operatorname{Hom}_{\mathbb Z\text{-gr.}}(T,\mathbb G_m)=M,
\]
and shall regard the elements of \(R\subset M\) as characters of \(T\).
Likewise, the elements of \(R^\ast\) will be regarded as morphisms of
\(\mathbb Z\)-groups \(\mathbb G_m\to T\).

\SGANextSectionTarget{sga3:xxv:subsection:2-3}\SGASubsection{2.3}{}
\SGARefTarget{sga3:xxv:paragraph:2:3}

For each \(\alpha\in R^+\), let \(\mathbb G_a(\alpha)\) be a copy of
\(\mathbb G_a\), and consider the \(\mathbb Z\)-scheme
\[
 U=\mathbb G_a(\alpha_1)\times\cdots\times\mathbb G_a(\alpha_n).
\]
If \(U_k\) denotes the unipotent part of the Borel subgroup \(B_k\) of
\(G_k\) defined by \(R^+\), let
\[
 a\colon U\otimes_{\mathbb Z}k\xrightarrow{\ \sim\ }U_k
\]
be the isomorphism of \(k\)-schemes defined by
\[
 a(x_1,\ldots,x_n)
 =p_{\alpha_1,k}(x_1)\cdots p_{\alpha_n,k}(x_n).
\]

\SGANextSectionTarget{sga3:xxv:subsection:2-4}\SGASubsection{2.4}{}
\SGARefTarget{sga3:xxv:paragraph:2:4}

The group law of \(U_k\) is expressed by relations of the form
\[
 a(x_1,\ldots,x_n)a(y_1,\ldots,y_n)=a(z_1,\ldots,z_n),
\]
where, for \(h=1,\ldots,n\), each \(z_h\) is a polynomial
\[
 z_h=x_h+y_h+
 Q_h(x_1,\ldots,x_{h-1},y_1,\ldots,y_{h-1}),
\]
whose coefficients are integers (\SGARefLink{sga3:xxii:lem:5.5.8}{Exposé XXII, 5.5.8}, and \SGARefLink{sga3:XXIII:6.4}{Exposé XXIII,
6.4}). Moreover,
\[
 Q_h(x_1,\ldots,x_{h-1},0,\ldots,0)=0.
\]

Equip \(U\) with the composition law defined by these formulas, which are
indeed ``defined over \(\mathbb Z\)''. Since this law induces the group law
on \(U_k\), it is associative, and \((0)\) is a unit element. Indeed, these
two assertions are expressed by relations among the polynomials \(Q_h\),
and \(\mathbb Z\to k\) is injective. This is a group law: if \((x_i)\) is
a section of \(U\) over an arbitrary scheme \(S\), its inverse \((y_i)\)
is computed recursively by
\[
 y_i=-x_i-Q_i(x_1,\ldots,x_{i-1},y_1,\ldots,y_{i-1}),
\]
formulas which are again ``defined over \(\mathbb Z\)''.

In summary, we have constructed a group law on \(U\) for which the
isomorphism \(a\) above is an isomorphism of groups.

For each \(\alpha\in R^+\), consider the morphism
\[
 p_\alpha\colon\mathbb G_a\longrightarrow U
\]
defined by \(p_\alpha(x)=(x_i)\), where
\[
 x_i=
 \begin{cases}
 x,&\alpha_i=\alpha,\\
 0,&\alpha_i\ne\alpha.
 \end{cases}
\]
It is a closed immersion and a group homomorphism; denote its image by
\(U_\alpha\). Since
\[
 (x_i)=p_{\alpha_1}(x_1)\cdots p_{\alpha_n}(x_n),
\]
\(U\) is identified with the product
\[
 U=U_{\alpha_1}U_{\alpha_2}\cdots U_{\alpha_n}.
\]

\SGANextSectionTarget{sga3:xxv:subsection:2-5}\SGASubsection{2.5}{}
\SGARefTarget{sga3:xxv:paragraph:2:5}

Let \(T=D(M)\) act on each \(U_\alpha\) through the character \(\alpha\).
One verifies immediately that this defines an action of \(T\) on the
group \(U\), and one may form the semidirect product \(B=T\ltimes U\).
There is a canonical isomorphism of \(k\)-groups
\[
 B\otimes_{\mathbb Z}k\xrightarrow{\ \sim\ }B_k.
\]

If an arbitrary order is now placed on \(R^+\), the morphism
\[
 \prod_{\alpha\in R^+}U_\alpha\longrightarrow U
\]
defined by multiplication in \(U\) is still an isomorphism. Indeed, both
sides are flat \(\mathbb Z\)-schemes of finite presentation, so it is
enough to verify the assertion on geometric fibers. We are then reduced
to Lazard's theory (\textit{Bible}, \S13.1, \SGARefLink{sga3:xxv:paragraph:2:1}{Proposition 1}): regard \(U\)
as a group with operators \(T\), and use the fact that the \(U_\alpha\)
are pairwise nonisomorphic as groups with operators, since the characters
\(\alpha\in R^+\) of \(T\) are pairwise distinct on every fiber.


% END INCLUDED SGA3_English_current_buildable_source_20260823/source/inputs/XXV/tex/components/01_section2_1_through_2_5.tex

% BEGIN INCLUDED SGA3_English_current_buildable_source_20260823/source/inputs/XXV/tex/components/02_section2_6_through_2_9.tex
\SGANextSectionTarget{sga3:xxv:subsection:2-6}\SGASubsection{2.6}{}
\SGARefTarget{sga3:xxv:paragraph:2:6}

Replacing \(R^+\) by \(R^-=-R^+\), one similarly constructs the groups
\(U^-\), \(B^-\), and \(U_\alpha\) for \(\alpha\in R^-\), together with
isomorphisms
\[
 p_\alpha\colon\mathbb G_a\xrightarrow{\ \sim\ }U_\alpha,
 \qquad \alpha\in R^-.
\]
Finally introduce the product scheme
\[
 \Omega=U^-\times T\times U.
\]
There is a canonical isomorphism of \(k\)-schemes
\[
 \Omega\otimes_{\mathbb Z}k
 \xrightarrow{\ \sim\ }
 U^-_k\times_kT_k\times_kU_k
 \simeq\Omega_k,
\]
where \(\Omega_k\) is the ``big cell'' of \(G_k\) (\SGARefLink{sga3:xxii:prop:4.1.2}{Exposé XXII, 4.1.2}).

From now on we identify \(\Omega\otimes_{\mathbb Z}k\) with \(\Omega_k\)
by this isomorphism, and regard \(U^-\), \(T\), and \(U\) as subschemes
of \(\Omega\) through their unit sections. Write
\[
 e=((0),e,(0))
\]
for the ``unit section'' of \(\Omega\). Our aim is now to put a
piece-of-group law on \(\Omega\).

\begin{sgastatement}{\SGAInlineRefTarget{sga3:xxv:lemma:2-7}Lemma 2.7}
Let \(\alpha\in\Delta\), and let \(w_{\alpha,k}\) be the element of
\(\operatorname{Norm}_{G_k}(T_k)(k)\) defined by \(X_{\alpha,k}\); recall
that, by definition,
\[
 w_{\alpha,k}
 =p_{-\alpha,k}(-1)p_{\alpha,k}(1)p_{-\alpha,k}(-1).
\]
There exist an open subset \(V_\alpha\) of \(\Omega\), containing the
section \(e\), and a morphism
\[
 h_\alpha\colon V_\alpha\longrightarrow\Omega
\]
such that:
\begin{enumerate}[label=\textup{(\roman*)},leftmargin=2.2em]
\item \(h_\alpha(e)=e\);
\item \((h_\alpha)\otimes_{\mathbb Z}k\) coincides with the restriction
of \(\operatorname{int}(w_{\alpha,k})\) to
\(V_\alpha\otimes_{\mathbb Z}k\subset G_k\);
\item \(T\subset V_\alpha\), and \(h_\alpha\) sends \(T\) into \(T\).
For every \(\beta\in R\), one has \(U_\beta\subset V_\alpha\), and
\(h_\alpha\) sends \(U_\beta\) into \(U_{s_\alpha(\beta)}\).
\end{enumerate}
\end{sgastatement}
\SGARefTarget{sga3:xxv:lemma:2:7}

By the definition of a Chevalley system (\SGARefLink{sga3:XXIII:6.1}{Exposé XXIII, 6.1}), for every
\(\beta\in R\) there exists an integer \(e_\beta=\pm1\) such that
\[
 \operatorname{int}(w_{\alpha,k})p_{\beta,k}(x)
 =p_{s_\alpha(\beta),k}(e_\beta x)
\]
for every \(x\in\mathbb G_a(S)\) and every \(S\to\operatorname{Spec}(k)\).

Let \(S\) be any scheme. Write an arbitrary element of \(\Omega(S)\) in
the form
\[
\left(
 \left(\prod_{\substack{\beta\in R^-\\\beta\ne-\alpha}}
 p_\beta(x_\beta)\right)p_{-\alpha}(x_{-\alpha}),
 \ t,\
 p_\alpha(x_\alpha)
 \left(\prod_{\substack{\beta\in R^+\\\beta\ne\alpha}}
 p_\beta(x_\beta)\right)
\right),
\]
where arbitrary orders have been chosen on \(R^-\setminus\{-\alpha\}\)
and \(R^+\setminus\{\alpha\}\), as permitted by \SGARefLink{sga3:xxv:subsection:2-5}{2.5}. Define a morphism
\(d\colon\Omega\to\mathbb A^1_{\mathbb Z}\) by
\[
 d(u)=\alpha(t)+x_\alpha x_{-\alpha}.
\]

Let \(V_\alpha=\Omega_d\), the open subset of \(\Omega\) defined by the
invertibility of \(d(u)\). It contains \(e\), \(T\), and every \(U_\beta\),
\(\beta\in R\). Define
\[
 h_\alpha\colon V_\alpha\longrightarrow\Omega,\qquad
 h_\alpha(u)=(a(u),b(u),c(u)),
\]
where
\begin{align*}
a(u)
 &=
\left(\prod_{\substack{\beta\in R^-\\\beta\ne-\alpha}}
 p_{s_\alpha(\beta)}(e_\beta x_\beta)\right)
p_{-\alpha}\!\left(-x_\alpha d(u)^{-1}\right),\\
b(u)&=t\cdot\alpha^\ast(d(u)),\\
c(u)
 &=p_\alpha\!\left(-x_{-\alpha}d(u)^{-1}\right)
\left(\prod_{\substack{\beta\in R^+\\\beta\ne\alpha}}
 p_{s_\alpha(\beta)}(e_\beta x_\beta)\right).
\end{align*}
Since \(s_\alpha\) permutes the positive roots distinct from \(\alpha\)
(respectively the negative roots distinct from \(-\alpha\)), \(c(u)\)
(respectively \(a(u)\)) is a section of \(U\) (respectively \(U^-\)),
so the preceding morphism is well defined. It plainly satisfies (i)
and (iii). Assertion (ii) follows immediately from the definition of
the \(e_\beta\), \(\beta\in R\), and \SGARefLink{sga3:xx:3.12}{Exposé XX, 3.12}.

\begin{sgastatement}{\SGAInlineRefTarget{sga3:xxv:lemma:2-8}Lemma 2.8}
There exist open subsets \(V,V'\) of \(\Omega\) and morphisms
\[
 h\colon V\longrightarrow\Omega,\qquad
 h'\colon V'\longrightarrow\Omega
\]
with the following properties:
\begin{enumerate}[label=\textup{(\roman*)},leftmargin=2.2em]
\item \(V\) and \(V'\) contain \(e\), and \(h(e)=h'(e)=e\);
\item\SGAInlineRefTarget{sga3:xxv:lemma:2-8:item:ii} the morphism induced by
\[
 h'\circ h\colon h^{-1}(V')\longrightarrow\Omega
\]
is the restriction of the identity morphism;
\item\SGAInlineRefTarget{sga3:xxv:lemma:2-8:item:iii} the \(k\)-morphisms \(h\otimes_{\mathbb Z}k\) and
\(h'\otimes_{\mathbb Z}k\) are the restrictions, to
\(V\otimes_{\mathbb Z}k\) and \(V'\otimes_{\mathbb Z}k\), of
automorphisms of the group \(G_k\);
\item\SGAInlineRefTarget{sga3:xxv:lemma:2-8:item:iv} \(V\) and \(V'\) contain \(U,T,U^-\); the morphisms \(h,h'\)
send \(U\) into \(U^-\), \(U^-\) into \(U\), and \(T\) into \(T\).
\end{enumerate}
\end{sgastatement}
\SGARefTarget{sga3:xxv:lemma:2:8}

Let \(w^0\) be the element of the Weyl group of \(G_k\) that sends \(R^+\)
to \(R^-\). Write
\[
 w^0=s_{\alpha_n}\cdots s_{\alpha_1},\qquad \alpha_i\in\Delta
\]
(with no relation to the numbering in \SGARefLink{sga3:xxv:subsection:2-1}{2.1}), and put
\[
 w_0=w_{\alpha_n,k}\cdots w_{\alpha_1,k}
 \in\operatorname{Norm}_{G_k}(T_k)(k).
\]
Define recursively, for \(i\leq n\), an open subset \(V_i\) of \(\Omega\)
and a morphism \(h_i\colon V_i\to\Omega\) by
\[
 V_0=\Omega,\qquad h_0=\operatorname{id},
\]
and, for \(i=0,\ldots,n-1\),
\[
 V_{i+1}=h_i^{-1}(V_{\alpha_{i+1}}),\qquad
 h_{i+1}=h_{\alpha_{i+1}}\circ h_i,
\]
where \(V_{\alpha_j}\) and \(h_{\alpha_j}\) are those of \SGARefLink{sga3:xxv:lemma:2-7}{2.7}.

Take \(V=V_n\) and \(h=h_n\). The conditions (i), (iii), and (iv) concerning
\(V\) and \(h\) hold: for (i) and (iii) this follows immediately from
\SGARefLink{sga3:xxv:lemma:2-8}{2.8}, and for (iv) from the fact that \(h\otimes_{\mathbb Z}k\) is the
restriction of \(\operatorname{int}(w_0)\) to
\(V\otimes_{\mathbb Z}k\).

Since \((w^0)^2=1\), we also have
\[
 w^0=s_{\alpha_1}\cdots s_{\alpha_n}
 =(s_{\alpha_1}^{\,3})\cdots(s_{\alpha_n}^{\,3}).
\]
Put
\[
 w'_0=(w_{\alpha_1,k})^3\cdots(w_{\alpha_n,k})^3.
\]
Repeating the construction above gives an open subset \(V'\) and a
morphism \(h'\) that also satisfy (i), (iii), and (iv). Moreover,
\(h'\otimes_{\mathbb Z}k\) is the restriction of
\(\operatorname{int}(w'_0)\) to \(V'\otimes_{\mathbb Z}k\). For every
simple root \(\alpha\in\Delta\), however,
\((w_{\alpha,k})^4=e\) (\SGARefLink{sga3:xx:3.1}{Exposé XX, 3.1}), hence \(w'_0w_0=e\).
It follows that \(h'\circ h\) induces the identity morphism on a nonempty
open subset of \(\Omega\otimes_{\mathbb Z}k\). Since \(\Omega\) is smooth
and of finite presentation over \(\mathbb Z\),
\(\Omega\otimes_{\mathbb Z}k\) is schematically dense in \(\Omega\);
this proves (ii).

\begin{sgastatement}{\SGAInlineRefTarget{sga3:xxv:proposition:2-9}Proposition 2.9}
There exist an open subset \(V_1\) of \(\Omega\times\Omega\), an open
subset \(V_2\) of \(\Omega\), and morphisms
\[
 \pi\colon V_1\longrightarrow\Omega,\qquad
 \sigma\colon V_2\longrightarrow\Omega
\]
having the following properties:
\begin{enumerate}[label=\textup{(\roman*)},leftmargin=2.2em]
\item if \(x\in\Omega(S)\), then \((e,x)\) and \((x,e)\) are sections
of \(V_1\), and
\[
 \pi(e,x)=\pi(x,e)=x;
\]
\item \(V_2\) contains \(e\), and \(\sigma(e)=e\);
\item\SGAInlineRefTarget{sga3:xxv:proposition:2-9:item:iii} \(\pi_k\) and \(\sigma_k\) are the restrictions of the morphisms
\(G_k\times_kG_k\to G_k\) and \(G_k\to G_k\) defined respectively by
multiplication and inversion.
\end{enumerate}
\end{sgastatement}
\SGARefTarget{sga3:xxv:proposition:2:9}

\SGAProof
Let \((v,t,u)\) and \((v',t',u')\) be two sections of \(\Omega\).
By \SGARefLink{sga3:xxv:lemma:2-8:item:iv}{2.8(iv)}, \(h(u)\) is a section of \(U^-\) and \(h(v')\) a section
of \(U\); thus \((h(u),e,h(v'))\) may be regarded as a section of
\(\Omega\). Define \(V_1\) as the open subset of \(\Omega\times\Omega\)
specified by
\[
\bigl((v,t,u),(v',t',u')\bigr)\in V_1(S)
\quad\Longleftrightarrow\quad
(h(u),e,h(v'))\in V'(S),
\]
using the notation of \SGARefLink{sga3:xxv:lemma:2-8}{2.8}. If the pair on the left is a section of
\(V_1\), then \(h'(h(u),e,h(v'))\) is defined and is a section of
\(\Omega\); write
\[
 h'(h(u),e,h(v'))=(v'',t'',u'').
\]
Set
\[
 \pi\bigl((v,t,u),(v',t',u')\bigr)
 =
 \bigl(v\cdot t v''t^{-1},\ tt''t',\
 t'^{-1}u''t'\cdot u'\bigr).
\]
Assertion (i) follows immediately from \SGARefLink{sga3:xxv:lemma:2-8:item:ii}{2.8(ii)}. To verify the assertion
about \(\pi\) in (iii), observe that
\[
 h'(h(u),e,h(v'))=uv'=v''t''u''
\]
whenever \(u\in U(S)\), \(v'\in U^-(S)\), and \(S\to k\), by
\SGARefLink{sga3:xxv:lemma:2-8:item:iii}{2.8(iii)} and (ii).

The morphism \(\sigma\) is constructed similarly. If \((v,t,u)\) is a
section of \(\Omega\), then \(h(u^{-1})\) is a section of \(U^-\),
\(h(v^{-1})\) a section of \(U\), and \(h(t^{-1})\) a section of \(T\).
Thus
\[
 (h(u^{-1}),h(t^{-1}),h(v^{-1}))
\]
is a section of \(\Omega\).
Define an open subset \(V_2\) of \(\Omega\) by
\[
 (v,t,u)\in V_2(S)
 \quad\Longleftrightarrow\quad
 (h(u^{-1}),h(t^{-1}),h(v^{-1}))\in V'(S),
\]
and define
\[
 \sigma(v,t,u)=h'(h(u^{-1}),h(t^{-1}),h(v^{-1})).
\]
The required properties of \(\sigma\) are verified in the same way.

% END INCLUDED SGA3_English_current_buildable_source_20260823/source/inputs/XXV/tex/components/02_section2_6_through_2_9.tex

% BEGIN INCLUDED SGA3_English_current_buildable_source_20260823/source/inputs/XXV/tex/components/03_section2_10_and_section3.tex
\begin{sgastatement}{Corollary 2.10}
\(\pi\) is ``generically associative'' and \(\sigma\) is a ``generic
inverse'': if \(x,y,z\in\Omega(S)\), and if the expressions below are
defined (as they always are over an open subset of \(\Omega\) containing
the unit section), then
\[
 \pi(x,\pi(y,z))=\pi(\pi(x,y),z),\qquad
 \pi(x,\sigma(x))=e=\pi(\sigma(x),x).
\]
\end{sgastatement}
\SGARefTarget{sga3:xxv:corollary:2:10}

Indeed, the two sides of each formula define morphisms between smooth
\(\mathbb Z\)-schemes of finite presentation, and the morphisms agree on
the generic fibers by \SGARefLink{sga3:xxv:proposition:2-9:item:iii}{2.9(iii)}.\footnote{The French source reads
``by \SGARefLink{sga3:xxv:proposition:2-9:item:iii}{2.10(iii)}''; the intended reference is Proposition 2.9(iii).}

\begin{sgastatement}{\SGAInlineRefTarget{sga3:xxv:corollary:2-11}Corollary 2.11}
Let \(\alpha\in R\), and let \(S\) be a scheme. Suppose that
\(x,y\in\mathbb G_a(S)\) satisfy
\[
 (p_\alpha(x),p_{-\alpha}(y))\in V_1(S),
 \qquad 1+xy\in\mathbb G_m(S).
\]
These conditions define an open subset of \(\mathbb G_{a,S}^2\)
containing the section \((0,0)\). Then
\[
 \pi(p_\alpha(x),p_{-\alpha}(y))
 =
 \left(
 p_{-\alpha}\!\left(\frac{y}{1+xy}\right),
 \alpha^\ast(1+xy),
 p_\alpha\!\left(\frac{x}{1+xy}\right)
 \right).
\]
\end{sgastatement}
\SGARefTarget{sga3:xxv:corollary:2:11}

The proof is the same as above, by \SGARefLink{sga3:xx:2.1}{Exposé XX, 2.1}.

\SGASection{3}{Existence theorem: conclusion of the proof}
\SGARefTarget{sga3:xxv:section:3}

For convenience, make the following definition.

\begin{sgastatement}{\SGAInlineRefTarget{sga3:xxv:definition:3-1}Definition 3.1}
Let \(S\) be a scheme and \(G\) an \(S\)-group scheme. The group \(G\)
is called \emph{admissible} if there is an open immersion of
\(S\)-schemes
\[
 i\colon\Omega_S=\Omega\times S\longrightarrow G
\]
such that:
\begin{enumerate}[label=\textup{(\roman*)},leftmargin=2.2em]
\item the diagram
\[
\begin{tikzcd}[column sep=4.5em,row sep=3.2em]
(V_1)_S
  \arrow[r]
  \arrow[d,"\pi"]
&
\Omega_S\times_S\Omega_S
  \arrow[r,"i\times_S i"]
&
G\times_SG
  \arrow[d,"\pi_G"]
\\
\Omega_S
  \arrow[rr,"i_S"]
&&
G
\end{tikzcd}
\]
is commutative, where \(\pi_G\) denotes the multiplication morphism
of \(G\);
\item there is a finite set of sections \(a_j\in\Omega(S)\) such that
the open subsets \(i(a_j)\,i(\Omega_S)\) cover \(G\).
\end{enumerate}
\end{sgastatement}
\SGARefTarget{sga3:xxv:definition:3:1}

By ``Weil's theorem'' (\SGARefLink{sga3:xviii:corollary:3-13:item:iii}{Exposé XVIII, 3.13(iii)} and (iv)), we have:

\begin{sgastatement}{Lemma 3.2}
If every admissible \(S\)-group is affine whenever \(S\) is étale and
of finite type over \(\mathbb Z\), then there exists an admissible
affine \(\mathbb Z\)-group.
\end{sgastatement}
\SGARefTarget{sga3:xxv:lemma:3:2}

We now have:

\begin{sgastatement}{Lemma 3.3}
Let \(S\) be a scheme and \(G\) an admissible \(S\)-group. Then \(G\)
is smooth and of finite presentation over \(S\), with affine,
connected, semisimple fibers.
\end{sgastatement}
\SGARefTarget{sga3:xxv:lemma:3:3}

Since \(\Omega_S\) is smooth and of finite presentation over \(S\),
with connected fibers, the same is true of \(G\) by condition (ii).
To prove \SGARefLink{sga3:xxv:lemma:3:3}{3.3}, we may therefore suppose that \(S\) is the spectrum of a
field \(K\). Identify \(\Omega_K\) with its image in \(G\). Plainly
\(\Omega_K\) is the product
\[
 \left(\prod_{\alpha\in R^-}U_{\alpha,K}\right)
 T_K
 \left(\prod_{\alpha\in R^+}U_{\alpha,K}\right)
\]
of the subgroups \(T_K\) and \(U_{\alpha,K}\), \(\alpha\in R\), of
\(G\). Thus the Lie algebra of \(G\) is identified with the direct sum
\[
 \operatorname{Lie}(T_K)
 \oplus
 \bigoplus_{\alpha\in R}\operatorname{Lie}(U_{\alpha,K}).
\]
The inner automorphism defined by a section of \(T_K\) acts on
\(U_{\alpha,K}\), and hence on \(\operatorname{Lie}(U_{\alpha,K})\),
through the character
\[
 \alpha\in R\subset M
 \simeq\operatorname{Hom}_{K\text{-gr.}}(T_K,\mathbb G_{m,K}).
\]
The preceding decomposition of \(\operatorname{Lie}(G)\) is therefore
exactly its decomposition under the adjoint action of \(T\). Thus the
roots of \(G_K\) with respect to \(T_K\) are the elements
\(\alpha\in R\).

Apply \SGARefLink{sga3:xix:corollary:1-13}{Exposé XIX, 1.13}. Let \(T_\alpha\) be the maximal torus of
\(\ker(\alpha)\subset T_K\), and put
\[
 Z_\alpha=\operatorname{Centr}_G(T_\alpha).
\]
It is enough to prove that each \(Z_\alpha\) is reductive. But
\[
 Z_\alpha\cap\Omega_K
 =
 \left(
 \prod_{\substack{\beta\in R^-\\\beta|_{T_\alpha}=e}}
 U_{\beta,K}
 \right)
 T_K
 \left(
 \prod_{\substack{\beta\in R^+\\\beta|_{T_\alpha}=e}}
 U_{\beta,K}
 \right).
\]
The roots vanishing on \(T_\alpha\) are the rational multiples of
\(\alpha\), hence precisely \(\alpha\) and \(-\alpha\). Consequently,
\[
 Z_\alpha\cap\Omega_K
 =U_{-\alpha,K}\,T_K\,U_{\alpha,K}.
\]
By \SGARefLink{sga3:xx:3.4}{Exposé XX, 3.4}, to prove \(Z_\alpha\) reductive it is enough to show
that \(U_{\alpha,K}\) and \(U_{-\alpha,K}\) do not commute, and this
follows immediately from \SGARefLink{sga3:xxv:corollary:2-11}{2.11}.

It follows from \SGARefLink{sga3:xxv:lemma:3:2}{3.2} and \SGARefLink{sga3:xxv:lemma:3:3}{3.3} that the proof will be complete once the
following lemma is proved.

\begin{sgastatement}{Lemma 3.4}
Let \(S\) be a locally noetherian scheme of dimension at most \(1\),
and let \(G\) be a smooth \(S\)-group of finite type, with affine,
connected, semisimple fibers. Then \(G\) is affine (and hence
semisimple).
\end{sgastatement}
\SGARefTarget{sga3:xxv:lemma:3:4}

\noindent\textit{Nota.}\quad
In Exposé XVI it was shown that \SGARefLink{sga3:xxv:lemma:3:4}{3.4} holds without any hypothesis on
\(S\), but the proof is comparatively delicate. Since only the special
case \SGARefLink{sga3:xxv:lemma:3:4}{3.4} is needed here, we give a direct proof.

Consider the Lie algebra \(\mathfrak g\) of \(G\), a locally free
\(\mathcal O_S\)-module, and the adjoint representation
\[
 \operatorname{Ad}\colon
 G\longrightarrow\operatorname{GL}_{\mathcal O_S}(\mathfrak g).
\]
To prove that \(G\) is affine over \(S\), it is enough to prove that
\(\operatorname{Ad}\) is affine. Since \(G\) is smooth with connected
fibers, it is separated over \(S\) (Exposé VI\textsubscript{B}, \SGARefLink{cor:sga3-VIB-5.5}{5.5}),
so \(\operatorname{Ad}\) is separated. By the result proved in the
Appendix (see \SGARefLink{sga3:xxv:proposition:4:1}{4.1}), it suffices to show that \(\operatorname{Ad}\) is
quasi-finite. We are thus reduced to the case in which \(S\) is the
spectrum of a field. In that case \(G\) is affine and therefore
semisimple, and the assertion is \SGARefLink{sga3:xxii:prop:5.7.14}{Exposé XXII, 5.7.14}.

% END INCLUDED SGA3_English_current_buildable_source_20260823/source/inputs/XXV/tex/components/03_section2_10_and_section3.tex

% BEGIN INCLUDED SGA3_English_current_buildable_source_20260823/source/inputs/XXV/tex/components/04_section4_and_bibliography.tex
\SGASection{4}{Appendix}
\SGARefTarget{sga3:xxv:section:4}

We used the following proposition in the course of the proof.

\begin{sgastatement}{Proposition 4.1}
Let \(S\) be a locally noetherian scheme of dimension at most \(1\);
let \(G\) and \(H\) be two \(S\)-group schemes of finite type; and let
\(f\colon G\to H\) be a quasi-finite separated group homomorphism.
If \(G\) is flat over \(S\),\footnote{Editors' note: We have added the
flatness hypothesis, which had been omitted.} then \(f\) is affine.
\end{sgastatement}
\SGARefTarget{sga3:xxv:proposition:4:1}

We shall give the proof only when \(G\) is smooth over \(S\), which is
indeed the hypothesis satisfied in the application above.

\SGANextSectionTarget{sga3:xxv:subsection:4-2}\SGASubsection{4.2}{}
\SGARefTarget{sga3:xxv:paragraph:4:2}

By EGA II, 1.6.4, we may suppose that \(S\) is reduced. By the usual
techniques of passage to the limit,\footnote{Editors' note: See
EGA IV\textsubscript{3}, 8.10.5(viii).} we may suppose \(S\) local. If
\(\dim(S)=0\), the assertion is trivial.\footnote{Editors' note:
Indeed, if \(S=\operatorname{Spec}(k)\), with \(k\) a field, then \(f\)
is the composite of the projection \(p\colon G\to G/\ker(f)\) and a
closed immersion \(i\). Since \(\ker(f)\) is finite over \(k\), \(p\)
is finite (Exposé VI\textsubscript{B}, 9.2).}
Suppose therefore that \(\dim(S)=1\). By faithfully flat descent, we may
suppose that \(S\) is complete with algebraically closed residue field.
After replacing \(S\) by its normalization \(\widetilde S\), if
necessary, we may suppose \(S\) normal (EGA II, 6.7.1, and
EGA \(0_{\mathrm{IV}}\), 23.1.5).\footnote{Editors' note:
Indeed, the irreducible components \(S_1,\ldots,S_r\) of \(S\) are
complete integral noetherian local schemes. By a theorem of Nagata
(EGA \(0_{\mathrm{IV}}\), 23.1.5), the normalization
\(\widetilde S_i\) is finite over \(S_i\), and hence
\(\widetilde S\) is finite over \(S\). A theorem of Chevalley
(EGA II, 6.7.1) then shows that if
\(f\times_S\widetilde S\) is affine, so is \(f\).}
We are thus reduced to the case in which \(S\) is the spectrum of a
complete discrete valuation ring \(A\) with algebraically closed
residue field.

\SGANextSectionTarget{sga3:xxv:subsection:4-3}\SGASubsection{4.3}{}
\SGARefTarget{sga3:xxv:paragraph:4:3}

Let \(\eta\) (respectively \(s\)) be the generic (respectively closed)
point of \(S\). Consider the image \(f_\eta(G_\eta)\) of \(G_\eta\)
in \(H_\eta\). It is a closed subgroup scheme of \(H_\eta\). Let
\(H'\) be the schematic closure of \(f_\eta(G_\eta)\) in \(H\).
Since \(H'\to H\) is affine (it is a closed immersion), we may replace
\(H\) by \(H'\), and thus suppose that \(H\) is flat over \(S\) and
that \(f_\eta\) is surjective. Since \(f_\eta\) is finite and \(G,H\)
are flat over \(S\), we have
\[
 \dim(G_s)=\dim(G_\eta)=\dim(H_\eta)=\dim(H_s).
\]

\SGANextSectionTarget{sga3:xxv:subsection:4-4}\SGASubsection{4.4}{}
\SGARefTarget{sga3:xxv:paragraph:4:4}

Let \(H_s^0,\ldots,H_s^n\) be the irreducible components of \(H_s\),
where \(H_s^0\) denotes the identity component, and let
\(z_0,\ldots,z_n\) be their generic points. Since each local ring
\(\mathcal O_{H,z_i}\) has dimension at most \(1\), the morphism
\[
 G\times_H\mathcal O_{H,z_i}\longrightarrow\mathcal O_{H,z_i}
\]
is affine, being quasi-finite and separated (see \SGARefLink{sga3:xvi:lemma:4:2}{Exposé XVI, Lemma 4.2}).
Thus \(G\) is affine over \(H\) in a neighborhood of \(z_i\). Let \(V\)
be the largest open subset of \(H\) over which \(G|_V\) is affine.
Then \(V\) contains every \(z_i\),\footnote{Editors' note: The following
sentence has been added.} and therefore contains at least one closed
point \(y_i\in H_s^i\). Moreover,
\(\kappa(y_i)=\kappa(s)\), since \(\kappa(s)\) is algebraically closed.

On the other hand, \(V\) is plainly stable under translation by any
element \(g\in G(S)\). We have \(\dim(G_s)=\dim(H_s)\), and \(f_s\) is
finite, so
\[
 f(s)\colon G_s^0(s)\longrightarrow H_s^0(s)
\]
is surjective. Since \(A\) is complete and \(G\) is smooth over \(S\),
the canonical map \(G^0(S)\to G_s^0(s)\) is surjective. Since
\(H_s^0(s)\) acts transitively on each \(H_s^i(s)\), it follows that
\(V\supset H_s(s)\), and therefore, since \(\kappa(s)\) is algebraically
closed, \(V\supset H_s\).\footnote{Editors' note: When \(G\) is not
assumed smooth over \(S\), one may argue as follows. Let \(h_0\) be a
rational point of \(H_s^0\), and let \(g_0\) be a rational point of
\(G_s^0\) such that \(f(g_0)=h_0\). By Exposé VI\textsubscript{B},
5.6.1, there is a commutative diagram
\[
\begin{tikzcd}[ampersand replacement=\&,column sep=4.0em,row sep=3.0em]
S'' \arrow[r,"g"] \arrow[d,"\pi"'] \arrow[dr,"\phi"]
  \& G \arrow[d]\\
S' \arrow[r,"w"] \& S
\end{tikzcd}
\]
in which \(w\) is étale and surjective, \(\pi\) is finite and
surjective, and \(\phi^{-1}(s)\) consists of a single point \(s''\)
such that \(g(s'')=g_0\). Then \(G_{S''}\to H_{S''}\) is affine above
a neighborhood of \(h_0y_i\), and the same is true of
\(G_{S'}\to H_{S'}\) (EGA II, 6.7.1), and hence of \(G\to H\) by
faithfully flat descent (EGA IV\textsubscript{2}, 2.7.1(xiii)).
Thus \(h_0y_i\in V\), from which it follows that \(V\) contains
\(H_s(s)\), and therefore \(H_s\).}

Clearly \(V\supset H_\eta\), since \(f_\eta\) is finite. Consequently,
\(V=H\). \SGAQED

\section*{Bibliography\footnote{Editors' note: Additional references
cited in this Exposé.}}
\addcontentsline{toc}{section}{Bibliography}
\SGARefTarget{sga3:xxv:bibliography}

\begin{description}[style=multiline,leftmargin=4.8em,labelwidth=4.0em]
\item[\textup{[BLie]}]
N.~Bourbaki, \textit{Groupes et algèbres de Lie}, Chapters I and
VII--VIII, Hermann, 1960 and 1975.

\item[\textup{[BT84]}]
F.~Bruhat and J.~Tits, \textit{Groupes réductifs sur un corps local II.
Schémas en groupes. Existence d'une donnée radicielle valuée},
Publ.\ Math.\ I.H.É.S. \textbf{60} (1984), 5--184.

\item[\textup{[DG70]}]
M.~Demazure and P.~Gabriel, \textit{Groupes algébriques},
Masson \& North-Holland, 1970.

\item[\textup{[Ko66]}]
B.~Kostant, \textit{Groups over \(\mathbb Z\)}, pp.~90--98 in
\textit{Algebraic groups and their discontinuous subgroups}
(A.~Borel and G.~D.~Mostow, eds.), Proc.\ Symp.\ Pure Math.\
\textbf{IX}, American Mathematical Society, 1966.

\item[\textup{[Lu09]}]
G.~Lusztig, \textit{Study of a \(\mathbb Z\)-form of the coordinate
ring of a reductive group}, J.\ Amer.\ Math.\ Soc.\ \textbf{22}
(2009), no.~3, 739--769.

\item[\textup{[Se66]}]
J.-P.~Serre, \textit{Algèbres de Lie semi-simples complexes},
Benjamin, 1966.

\item[\textup{[Sp98]}]
T.~A.~Springer, \textit{Linear algebraic groups}, 2nd ed.,
Birkhäuser, 1998.

\item[\textup{[Ti66]}]
J.~Tits, \textit{Sur les constantes de structure et le théorème
d'existence des algèbres de Lie semi-simples},
Publ.\ Math.\ I.H.É.S. \textbf{31} (1966), 21--58.
\end{description}

% END INCLUDED SGA3_English_current_buildable_source_20260823/source/inputs/XXV/tex/components/04_section4_and_bibliography.tex



% END INCLUDED SGA3_English_current_buildable_source_20260823/source/inputs/XXV/body.tex


\SGAStartExpose{XXVI}
\providecommand{\rad}{\operatorname{rad}}
\providecommand{\Transt}{\underline{\operatorname{Transt}}}
\providecommand{\Tor}{\underline{\operatorname{Tor}}}
\providecommand{\Lev}{\underline{\operatorname{Lev}}}
\providecommand{\Par}{\underline{\operatorname{Par}}}
\providecommand{\Crit}{\underline{\operatorname{Crit}}}
\providecommand{\Bor}{\underline{\operatorname{Bor}}}
\providecommand{\Gen}{\underline{\operatorname{Gen}}}
\providecommand{\Opp}{\underline{\operatorname{Opp}}}
\providecommand{\Stand}{\underline{\operatorname{Stand}}}
\providecommand{\TypeStand}{\underline{\operatorname{TypeStand}}}
\providecommand{\Of}{\operatorname{Of}}
\providecommand{\Int}{\operatorname{int}}
\renewcommand{\Centr}{\underline{\operatorname{Centr}}}
\renewcommand{\Norm}{\underline{\operatorname{Norm}}}
% BEGIN INCLUDED SGA3_English_current_buildable_source_20260823/source/inputs/XXVI/body.tex
% Mechanically extracted after \mainmatter from the frozen standalone master.

\markboth{Exposé XXVI. Parabolic Subgroups}{Exposé XXVI. Parabolic Subgroups}
% BEGIN INCLUDED SGA3_English_current_buildable_source_20260823/source/inputs/XXVI/tex/components/00_title_introduction_and_section1.tex
\chapter*{Exposé XXVI\\Parabolic Subgroups of Reductive Groups}
\addcontentsline{toc}{chapter}{Exposé XXVI. Parabolic Subgroups of Reductive Groups}
\begin{center}
\large M. Demazure
\end{center}

\medskip

This Exposé studies the parabolic subgroups of an \(S\)-reductive group \(G\).
Its essential result is the conjugacy theorem (\SGARefLink{sga3:xxvi:cor:5:4}{5.4}).  Its essential tool is
the notion of the transversal position of two parabolic subgroups, a notion
studied systematically in §4.  Another fact plays an important role: the
decomposition of the unipotent radical \(\rad^{u}(P)\) of a parabolic subgroup
\(P\) into successive extensions of vector groups
(\SGARefLink{sga3:xxvi:proposition:2-1}{2.1}).\footnote{\emph{Editor's note.} Since the Levi subgroups of \(P\) form a
torsor under \(\rad^{u}(P)\) (\SGARefLink{sga3:xxvi:cor:1:9}{1.9}), this entails, when \(S\) is semi-local,
that \(P\) has a Levi subgroup and hence a maximal torus (\SGARefLink{sga3:xxvi:cor:2:4}{2.4}).}

Various schemes associated with \(G\) are studied in §3; §6 deals with the
split subtori\footnote{\emph{Editor's note.} The terminology ``trivial
torus'' has been replaced by ``split torus''.} of \(G\) and their relations
with parabolic subgroups.

Finally, in §7, we explain briefly how, over a semi-local base, one formulates
the ``relative theory'' of reductive groups as presented over fields in the
article of A.~Borel and J.~Tits, \emph{Groupes réductifs}, Publications
Mathématiques de l'IHÉS, no.~27.  In that article, cited as [BT65] below, the
reader will also find, for a base field, other results not touched upon here.
\footnote{\emph{Editor's note.} Version of 13 October 2024.}

\SGASection{1}{Recollections. Levi subgroups}
\SGARefTarget{sga3:xxvi:section:1}

\begin{sgastatement}{Definition 1.1}
Let \(S\) be a scheme, \(G\) an \(S\)-reductive group, and \(P\) a
sub-\(S\)-group-scheme of \(G\).  One says that \(P\) is a
\emph{parabolic subgroup} of \(G\) if:
\begin{enumerate}[label=\textup{(\roman*)}]
\item \(P\) is smooth over \(S\);
\item for every \(s\in S\), the quotient \(\bar s\)-scheme
\(G_{\bar s}/P_{\bar s}\) is proper (that is, by the Bible, §\SGARefLink{sga3:xxvi:corollary:6-4}{6.4},
\SGARefLink{sga3:xxvi:prop:1:4}{Theorem~4} \(=\) [Ch05], §\SGARefLink{sga3:xxvi:subsec:6:5}{6.5}, Theorem~5, \(P_{\bar s}\) contains a
Borel subgroup of \(G_{\bar s}\)).
\end{enumerate}
\end{sgastatement}
\SGARefTarget{sga3:xxvi:def:1:1}

\begin{sgastatement}{Proposition 1.2 \textup{(Exp.\ XXII, \SGARefLink{sga3:xxii:cor:5.8.5}{5.8.5})}}
Let \(S\) be a scheme, \(G\) an \(S\)-reductive group, and \(P\) a parabolic
subgroup of \(G\).  Then \(P\) is closed in \(G\), has connected fibers, and
\[
  P=\Norm_G(P).
\]
Moreover, the quotient sheaf \(G/P\) is representable by an \(S\)-scheme
that is smooth and projective over \(S\).
\end{sgastatement}
\SGARefTarget{sga3:xxvi:prop:1:2}

\begin{sgastatement}{\SGAInlineRefTarget{sga3:xxvi:proposition:1-3}Proposition 1.3 \textup{(Exp.\ XXII, \SGARefLink{sga3:xxii:thm:5.3.9}{5.3.9} and 5.3.11)}}
Let \(S\) be a scheme, \(G\) an \(S\)-reductive group, and \(P,P'\) two
parabolic subgroups of \(G\).  The following conditions are equivalent:
\begin{enumerate}[label=\textup{(\roman*)}]
\item \(P\) and \(P'\) are conjugate in \(G\), locally for the étale
\textup{(respectively, fpqc)} topology.
\item For every \(s\in S\), \(P_{\bar s}\) and \(P'_{\bar s}\) are conjugate
by an element of \(G(\bar s)\).
\item The strict transporter \(\Transt_G(P,P')\) of \(P\) into \(P'\), defined
for every \(S'\to S\) by
\[
 \Transt_G(P,P')(S')
 =\{g\in G(S')\mid \Int(g)P_{S'}=P'_{S'}\},
\]
is a closed subscheme of \(G\), smooth and of finite presentation over
\(S\), and is a principal homogeneous bundle on the right under \(P\) and
on the left under \(P'\).
\end{enumerate}
\end{sgastatement}
\SGARefTarget{sga3:xxvi:prop:1:3}

\begin{sgastatement}{Proposition 1.4}
Let \(S\) be a nonempty scheme, let \((G,T,M,R)\) be a split
\(S\)-reductive group, and let \(R'\) be a subset of \(R\).  The following
conditions on \(R'\) are equivalent:
\begin{enumerate}[label=\textup{(\roman*)}]
\item
\[
 \mathfrak g_{R'}=\mathfrak t\oplus
 \bigoplus_{\alpha\in R'}\mathfrak g^\alpha
\]
is the Lie algebra of a parabolic subgroup of \(G\) containing \(T\)
\textup{(necessarily unique, Exp.\ XXII, \SGARefLink{sga3:xxii:cor:5.3.5}{5.3.5})}.
\item \(R'\) is of type \textup{(R)} \textup{(Exp.\ XXII, \SGARefLink{sga3:xxii:def:5.4.2}{5.4.2})} and
contains a system of positive roots.
\item \(R'\) is a closed subset of \(R\) and satisfies: if
\(\alpha\in R-R'\), then \(-\alpha\in R'\); equivalently,
\(R=R'\cup(-R')\).
\item There exist a system \(\Delta\) of simple roots and a subset
\(A\subset\Delta\) such that \(R'\) is the union of the positive roots and
the negative roots that are linear combinations of elements of \(A\).
\item\SGAInlineRefTarget{sga3:xxvi:prop:1:4:item:v} \(R'\) contains a system of simple roots of \(R\); moreover, if
\(\Delta\subset R'\) is a system of simple roots of \(R\) and
\[
 A=(-R')\cap\Delta,
\]
then \(R'\) is the union of the positive roots and the negative roots that
are linear combinations of elements of \(A\).
\end{enumerate}
\end{sgastatement}
\SGARefTarget{sga3:xxvi:prop:1:4}

One has (i)\(\Leftrightarrow\)(ii) by Exp.\ XXII, \SGARefLink{sga3:xxii:prop:5.4.5:item:ii}{5.4.5(ii)} and \SGARefLink{sga3:xxii:prop:5.5.1}{5.5.1},
and (iii)\(\Rightarrow\)(ii) by Exp.\ XXI, \SGARefLink{sga3:XXI:3.3.6}{3.3.6} and Exp.\ XXII, \SGARefLink{sga3:xxii:thm:5.4.7}{5.4.7}.
Clearly (v)\(\Rightarrow\)(iv)\(\Rightarrow\)(iii), while
(iii)\(\Rightarrow\)(v) follows from Exp.\ XXI, \SGARefLink{sga3:XXI:3.3.6}{3.3.6} and 3.3.10.

It remains to prove that (i) implies that \(R'\) is closed in \(R\).  This
last assertion may be checked on any geometric fiber, so we may suppose that
\(S\) is the spectrum of an algebraically closed field.

Let \(P\) be the parabolic subgroup of \(G\) containing \(T\) whose Lie
algebra is \(\mathfrak g_{R'}\).  Since the Borel subgroups of \(P\) are the
Borel subgroups of \(G\) contained in \(P\), the Bible, §12.3, \SGARefLink{sga3:xxvi:theorem:4:2:1}{Theorem~1},
and Exp.\ XXII, \SGARefLink{sga3:xxii:prop:5.4.5:item:i}{5.4.5(i)}, show that, if \(U\) denotes the unipotent radical
of \(P\), then \(T\cdot U\) is the subgroup of \(G\) containing \(T\) whose
Lie algebra is
\[
 \mathfrak g_{R''}=\mathfrak t\oplus
 \bigoplus_{\alpha\in R''}\mathfrak g^\alpha,
\]
where \(R''\) is the intersection of the systems of positive roots of \(R\)
contained in \(R'\).  In particular, \(R''\) is closed and
\(R''\cap(-R'')=\varnothing\).

On the other hand, \(H=P/U\) is reductive, the canonical image \(\bar T\)
of \(T\) is a maximal torus of \(H\) (and \(T\to\bar T\) is an
isomorphism), and there is an isomorphism of \(\bar T\)-modules, or
equivalently of \(M\)-graded vector spaces,
\[
 \operatorname{Lie}(H)\simeq\mathfrak t\oplus
 \bigoplus_{\alpha\in R_s}\mathfrak g^\alpha,
\]
where \(R_s\) is the complement of \(R''\) in \(R'\).  Thus \(R_s\)
identifies naturally with the roots of \(H\) relative to \(\bar T\), and in
particular \(R_s=-R_s\).  Consequently,
\[
 R''=\{\alpha\in R'\mid-\alpha\notin R'\},
 \qquad
 R_s=\{\alpha\in R'\mid-\alpha\in R'\}.
\]

We now prove that \(R'\) is closed.  Let \(\alpha,\beta\in R'\) with
\(\alpha+\beta\in R\).  If \(\alpha,\beta\in R''\), then
\(\alpha+\beta\in R''\) because \(R''\) is closed.  If
\(\alpha\in R_s\), \(\beta\in R''\), and \(\alpha+\beta\notin R'\), then
\(\alpha+\beta\in-R''\), whence
\[
 -\alpha=-(\alpha+\beta)+\beta\in R'',
\]
because \(R''\) is closed.  This gives
\(-\alpha\in R''\cap R_s\), a contradiction.  It remains to consider
\(\alpha,\beta\in R_s\).  If \(\alpha+\beta\notin R'\), then
\(\alpha+\beta\in-R''\).  Since \(\alpha+\beta\ne0\), some system of
positive roots of the root system \(R_s\) contains both \(\alpha\) and
\(\beta\), and hence some Borel subgroup of \(H=P/U\) contains the canonical
images of \(U_\alpha\) and \(U_\beta\).  Its inverse image in \(P\) is a
Borel subgroup containing \(U_\alpha\), \(U_\beta\), \(U\), and therefore
\(U_{-(\alpha+\beta)}\), which is impossible.

\begin{sgastatement}{Corollary 1.5}
A parabolic subgroup of a reductive group is of type \textup{(RC)}
\textup{(Exp.\ XXII, \SGARefLink{sga3:xxii:def:5.11.1}{5.11.1})}.
\end{sgastatement}
\SGARefTarget{sga3:xxvi:cor:1:5}

\begin{sgastatement}{\SGAInlineRefTarget{sga3:xxvi:proposition:1-6}Proposition 1.6 \textup{(Exp.\ XXII, \SGARefLink{sga3:xxii:prop:5.11.4}{5.11.4})}}
Let \(S\) be a scheme, \(G\) an \(S\)-reductive group, and \(P\) a subgroup
of \(G\) of type \textup{(RC)}.\footnote{\emph{Editor's note.}
``Parabolic subgroup'' has been replaced by ``subgroup of type (RC)'', as
in Exp.\ XXII, \SGARefLink{sga3:xxii:prop:5.11.4}{5.11.4}, because the statement will later be applied
(\SGARefLink{sga3:xxvi:proposition:4-5-1}{4.5.1}, \SGARefLink{sga3:xxvi:cor:6:17}{6.17}) to \(P\cap P'\) when \(P,P'\) are parabolic and
\(P\cap P'\) is of type (RC).}
\begin{enumerate}[label=\textup{(\roman*)}]
\item \(P\) has a largest invariant sub-group-scheme that is smooth and of
finite presentation over \(S\), with connected unipotent geometric fibers.
It is a characteristic subgroup of \(P\), called the \emph{unipotent
radical} of \(P\) and denoted \(\rad^u(P)\).  The quotient sheaf
\(P/\rad^u(P)\) is representable by an \(S\)-reductive group.
\item\SGAInlineRefTarget{sga3:xxvi:proposition:1-6:item:ii} If \(T\) is a maximal torus of \(P\), then \(P\) has a reductive
subgroup \(L\) containing \(T\) such that:
\begin{enumerate}[label=\textup{(\alph*)}]
\item every reductive subgroup of \(P\) containing \(T\) is contained in
\(L\);
\item \(P\) is the semidirect product \(L\cdot\rad^u(P)\), that is, the
canonical morphism \(L\to P/\rad^u(P)\) is an isomorphism.
\end{enumerate}
Moreover, \(L\) is the unique subgroup \textup{(respectively, reductive
subgroup)} of \(P\) containing \(T\) and satisfying (b) \textup{(respectively,
(a))}.  Finally,
\[
 \Norm_P(L)=L,\qquad \Norm_P(T)=\Norm_L(T).
\]
\end{enumerate}
\end{sgastatement}
\SGARefTarget{sga3:xxvi:prop:1:6}

\paragraph{\SGARefLink{sga3:xxvi:prop:1:6}{1.7}.}
A subgroup \(L\) of \(P\) satisfying condition \SGARefLink{sga3:xxvi:proposition:1-6:item:ii}{1.6(ii)}(b) is called a
\emph{Levi subgroup} of \(P\).  It is a maximal reductive subgroup of
\(P\).  Indeed, it is reductive because it is isomorphic to
\(P/\rad^u(P)\).  If \(L'\) is a reductive subgroup of \(P\) containing
\(L\), then, locally for the fpqc topology, \(L\) has a maximal torus
\(T\), and \SGARefLink{sga3:xxvi:proposition:1-6:item:ii}{1.6(ii)} gives \(L'=L\).
\SGARefTarget{sga3:xxvi:item:1:7}

If \(L\) and \(L'\) are two Levi subgroups of \(P\), they are conjugate in
\(P\), locally for the fpqc topology.  Indeed, locally one may choose
maximal tori \(T\subset L\) and \(T'\subset L'\); after a further fpqc
localization they are conjugate in \(P\), so one may assume \(T=T'\), and
then \SGARefLink{sga3:xxvi:proposition:1-6:item:ii}{1.6(ii)} gives \(L=L'\).  Since
\(P=L\cdot\rad^u(P)\) and \(\Norm_P(L)=L\), one obtains:

\begin{sgastatement}{Corollary 1.8}
Let \(P\) be a subgroup of type \textup{(RC)} of the \(S\)-reductive group
\(G\).  If \(L,L'\) are two Levi subgroups of \(P\), there is a unique
\(u\in\rad^u(P)(S)\) such that \(\Int(u)L=L'\).
\end{sgastatement}
\SGARefTarget{sga3:xxvi:cor:1:8}

Let \(\Lev(P)\) denote the functor of Levi subgroups of \(P\): for
\(S'\to S\), \(\Lev(P)(S')\) is the set of Levi subgroups of \(P_{S'}\).
\SGARefLink{sga3:xxvi:cor:1:8}{Corollary~1.8} gives:

\begin{sgastatement}{Corollary 1.9}
Let \(P\) be a subgroup of type \textup{(RC)} of the \(S\)-reductive group
\(G\).  Then \(\Lev(P)\) is a principal homogeneous bundle under the
\(S\)-group \(\rad^u(P)\); in particular, it is representable by an
\(S\)-scheme that is smooth and affine over \(S\), with integral geometric
fibers.
\end{sgastatement}
\SGARefTarget{sga3:xxvi:cor:1:9}

Likewise, \SGARefLink{sga3:xxvi:proposition:1-6}{1.6} immediately yields:

\begin{sgastatement}{Corollary 1.10}
Let \(P\) be a subgroup of type \textup{(RC)} of the \(S\)-reductive group
\(G\).  The functor \(\Tor(P)\) of maximal tori of \(P\) is representable
by a smooth affine \(S\)-scheme.  The relation \(L\supset T\) defines a
morphism
\[
 \Tor(P)\longrightarrow\Lev(P),
\]
whose fiber above \(L\in\Lev(P)(S)\) identifies with \(\Tor(L)\)
\textup{(Exp.\ XXII, \SGARefLink{sga3:xxii:cor:5.8.3}{5.8.3})}.
\end{sgastatement}
\SGARefTarget{sga3:xxvi:cor:1:10}

The first assertion of \SGARefLink{sga3:xxvi:cor:1:10}{1.10} follows from the other two and Exp.\ XXII,
\SGARefLink{sga3:xxii:cor:5.8.3}{5.8.3}.

\begin{sgastatement}{\SGAInlineRefTarget{sga3:xxvi:definition:1-11}Definition 1.11}
Let \(S\) be a nonempty scheme, \(G\) an \(S\)-reductive group, \(P\) a
parabolic subgroup of \(G\), and
\[
 \mathcal E=(T,M,R,\Delta,(X_\alpha)_{\alpha\in\Delta})
\]
a pinning of \(G\).  One says that \(\mathcal E\) is \emph{adapted} to
\(P\), or is a \emph{pinning of the pair \((G,P)\)}, if \(P\supset T\) and
the Lie algebra of \(P\) is of the form
\[
 \mathfrak t\oplus\bigoplus_{\alpha\in R'}\mathfrak g^\alpha,
\]
where \(R'\subset R\) contains \(R^+\).
\end{sgastatement}
\SGARefTarget{sga3:xxvi:def:1:11}

In particular, if \(T\subset B\) is the Killing pair defined by the
pinning, then \(T\subset B\subset P\).  Under these conditions set
\(\Delta(P)=\Delta\cap(-R')\).  Exp.\ XXII, \SGARefLink{sga3:xxii:lem:5.4.3}{5.4.3} gives
\[
 \alpha\in\Delta(P)
 \ \Longleftrightarrow\
 \alpha\in\Delta\text{ and }U_{-\alpha}\subset P
 \ \Longleftrightarrow\
 \alpha\in\Delta\text{ and }U_{-\alpha}\cap P\ne e.
\]

\SGARefLink{sga3:xxvi:prop:1:4:item:v}{Proposition~1.4(v)} and Exp.\ XXII, \SGARefLink{sga3:xxii:lem:5.11.3}{5.11.3} and 5.10.6 now immediately give:

\begin{sgastatement}{\SGAInlineRefTarget{sga3:xxvi:proposition:1-12}Proposition 1.12}
Let \(S\) be a scheme, \(G\) an \(S\)-reductive group, \(P\) a parabolic
subgroup of \(G\), and
\((T,M,R,\Delta,(X_\alpha)_{\alpha\in\Delta})\) a pinning of \(G\) adapted
to \(P\).  Let \(\Delta(P)\subset\Delta\) be as above.
\begin{enumerate}[label=\textup{(\roman*)}]
\item The unipotent radical of \(P\) is
\[
 \rad^u(P)=U_{R''}=\prod_{\alpha\in R''}U_\alpha,
\]
where \(R''\) is the set of positive roots whose expansion in \(\Delta\)
contains at least one element of
\(\Delta-\Delta(P)\)\footnote{\emph{Editor's note.}
\(\Delta(P)\) has been corrected to \(\Delta-\Delta(P)\).} with a
nonzero coefficient.
\item The unique Levi subgroup \(L\) of \(P\) containing \(T\) is
\[
 Z_{\Delta(P)}=\Centr_G(T_{\Delta(P)}),
\]
where \(T_{\Delta(P)}\) is the maximal torus of
\(\bigcap_{\alpha\in\Delta(P)}\ker(\alpha)\).  Moreover,
\(T_{\Delta(P)}=\rad(L)\).
\end{enumerate}
\end{sgastatement}
\SGARefTarget{sga3:xxvi:prop:1:12}

\begin{sgastatement}{Corollary 1.13}
Every Levi subgroup \(L\) of a parabolic subgroup \(P\) of a reductive group
\(G\) is a \emph{critical subgroup} of \(G\); that is
\textup{(Exp.\ XXII, \SGARefLink{sga3:xxii:def:5.10.4}{5.10.4})},
\[
 L=\Centr_G(\rad(L)).
\]
\end{sgastatement}
\SGARefTarget{sga3:xxvi:cor:1:13}

This follows immediately from \SGARefLink{sga3:xxvi:proposition:1-12}{1.12} and the following lemma, which is
contained in \SGARefLink{sga3:xxvi:prop:1:3}{1.4} and Exp.\ XXII, \SGARefLink{sga3:xxii:5.4.1}{5.4.1}:

\begin{sgastatement}{Lemma 1.14}
Locally for the étale topology, every pair \((G,P)\), where \(P\) is a
parabolic subgroup of the reductive group \(G\), admits a pinning in the
sense of \SGARefLink{sga3:xxvi:definition:1-11}{1.11}.
\end{sgastatement}
\SGARefTarget{sga3:xxvi:lemma:1:14}

We also record:

\begin{sgastatement}{Proposition 1.15}
Let \(S\) be a scheme, \(G\) an \(S\)-reductive group, \(P\) a parabolic
subgroup of \(G\), and \(\mathcal E,\mathcal E'\) two pinnings of \(G\)
adapted to \(P\).  The unique inner automorphism of \(G\) over \(S\) that
transforms \(\mathcal E\) into \(\mathcal E'\) \textup{(Exp.\ XXIV, \SGARefLink{sga3:xxiv:1.5}{1.5})}
comes from \(P\), via
\[
 P\longrightarrow P/\Centr(P)
 =P/\Centr(G)\longrightarrow G/\Centr(G).
\]
\end{sgastatement}
\SGARefTarget{sga3:xxvi:prop:1:15}

Indeed, it suffices to argue as in Exp.\ XXIV, \SGARefLink{sga3:xxiv:1.5}{1.5}, using:

\begin{sgastatement}{Lemma 1.16 \textup{(Exp.\ XXII, \SGARefLink{sga3:xxii:cor:5.3.14}{5.3.14} and \SGARefLink{sga3:xxii:prop:5.2.6}{5.2.6})}}
The maximal tori \textup{(respectively, Borel subgroups, respectively,
Killing pairs)} of a parabolic subgroup \(P\) of an \(S\)-reductive group
\(G\) are conjugate in \(P\), locally for the étale topology.
\end{sgastatement}
\SGARefTarget{sga3:xxvi:lemma:1:16}

\begin{sgastatement}{Proposition 1.17}
Let \(S\) be a scheme, \(G\) an \(S\)-reductive group, \(P,P'\) two
parabolic subgroups of \(G\), and \(B\) a Borel subgroup contained in both
\(P\) and \(P'\).  If \(P\) and \(P'\) are conjugate in \(G\), locally for
the étale topology, then \(P=P'\).
\end{sgastatement}
\SGARefTarget{sga3:xxvi:prop:1:17}

Indeed, we may assume that \(g\in G(S)\) satisfies \(\Int(g)P=P'\).  The
groups \(B\) and \(\Int(g)^{-1}B\) are then two Borel subgroups of \(P\).
After extending \(S\), \SGARefLink{sga3:xxvi:lemma:1:16}{Lemma~1.16} lets us choose \(p\in P(S)\) with
\(\Int(p)\Int(g^{-1})B=B\).  Hence
\[
 pg^{-1}\in\Norm_G(B)(S)=B(S),
\]
so \(g\in B(S)\cdot p\subset P(S)\), and therefore
\(P'=\Int(g)P=P\).

\begin{sgastatement}{Remark 1.18}
If \(P,P'\) are two parabolic subgroups of \(G\) containing a common Borel
subgroup, then \(P\cap P'\) is again a parabolic subgroup of \(G\).  Indeed,
it is smooth along the unit section \textup{(Exp.\ XXII, \SGARefLink{sga3:xxii:prop:5.4.5}{5.4.5})} and
contains a Borel subgroup.
\end{sgastatement}
\SGARefTarget{sga3:xxvi:rem:1:18}

\begin{sgastatement}{\SGAInlineRefTarget{sga3:xxvi:proposition:1-19}Proposition 1.19}
Let \(S\) be a scheme, \(G\) an \(S\)-reductive group, and \(G'\) its
derived group \textup{(Exp.\ XXII, \SGARefLink{sga3:xxii:thm:6.2.1}{6.2.1})}.
\begin{enumerate}[label=\textup{(\roman*)}]
\item The assignments
\[
 P\longmapsto P'=P\cap G',
 \qquad
 P'\longmapsto P'\cdot\rad(G)=\Norm_G(P')
\]
are mutually inverse bijections between the parabolic subgroups of \(G\)
and those of \(G'\).  Moreover,
\(\rad^u(P)=\rad^u(P')\).
\item Let \(P\) be a parabolic subgroup of \(G\) and \(P'=P\cap G'\).
The assignments
\[
 L\longmapsto L'=L\cap G'=L\cap P',
 \qquad
 L'\longmapsto L'\cdot\rad(G)=\Centr_G(\rad(L'))
\]
are mutually inverse bijections between the Levi subgroups of \(P\) and
those of \(P'\).  Moreover,
\[
 \rad(L')=(\rad(L)\cap G')^0.
\]
\end{enumerate}
\end{sgastatement}
\SGARefTarget{sga3:xxvi:prop:1:19}

The proof, for example by reduction to the split case, is straightforward
and is left to the reader, together with the immediate proof of:

\begin{sgastatement}{\SGAInlineRefTarget{sga3:xxvi:proposition:1-20}Proposition 1.20}
Let \(S\) be a scheme, \(G\) an \(S\)-reductive group, \(P\) a parabolic
subgroup of \(G\), and \(L\) a Levi subgroup of \(P\).  The assignments
\[
 Q\longmapsto Q\cap L=Q',
 \qquad
 Q'\longmapsto Q'\cdot\rad^u(P)
\]
are mutually inverse bijections between the parabolic subgroups of \(G\)
contained in \(P\) and the parabolic subgroups of \(L\).  Moreover, the
Levi subgroups of \(Q'\) are precisely the Levi subgroups of \(Q\)
contained in \(L\).
\end{sgastatement}
\SGARefTarget{sga3:xxvi:prop:1:20}

\SGARefLink{sga3:xxvi:proposition:1-6}{Proposition~1.6} may be completed as follows.

\begin{sgastatement}{\SGAInlineRefTarget{sga3:xxvi:proposition:1-21}Proposition 1.21}
Let \(S\) be a scheme, \(G\) an \(S\)-reductive group, and \(P\) a
parabolic subgroup of \(G\).
\begin{enumerate}[label=\textup{(\roman*)}]
\item \(P\) has a largest invariant subgroup that is smooth and of finite
presentation over \(S\), with connected solvable geometric fibers.  It is
a characteristic subgroup of \(P\), called the \emph{radical} of \(P\)
and denoted \(\rad(P)\).  The quotient sheaf \(P/\rad(P)\) is representable
by an \(S\)-semisimple group.
\item If \(L\) is a Levi subgroup of \(P\), then \(\rad(P)\) is the
semidirect product of \(\rad^u(P)\) and \(\rad(L)\); one has
\[
 \rad(L)=L\cap\rad(P),
 \qquad
 L=\Centr_G(L\cap\rad(P)),
 \qquad
 P/\rad(P)\simeq L/\rad(L).
\]
\end{enumerate}
\end{sgastatement}
\SGARefTarget{sga3:xxvi:prop:1:21}

Indeed, assertion (i) being local, one may suppose that \(P\) has a Levi
subgroup \(L\).  It is then immediate that
\[
 R=\rad^u(P)\cdot\rad(L)
\]
has the properties asserted in (i).  For (ii), it only remains to prove
that \(\rad(L)=L\cap\rad(P)\); this follows at once because
\(L\cap\rad(P)\) is smooth with connected fibers, \(L\) being the
centralizer of a torus.\footnote{\emph{Editor's note.} See \SGARefLink{sga3:xix:lemma:1:4}{Exposé XIX,
1.4}.}

% END INCLUDED SGA3_English_current_buildable_source_20260823/source/inputs/XXVI/tex/components/00_title_introduction_and_section1.tex

% BEGIN INCLUDED SGA3_English_current_buildable_source_20260823/source/inputs/XXVI/tex/components/01_section2.tex
\SGASection{2}{Structure of the unipotent radical of a parabolic subgroup}
\SGARefTarget{sga3:xxvi:section:2}

\begin{sgastatement}{\SGAInlineRefTarget{sga3:xxvi:proposition:2-1}Proposition 2.1}
Let \(S\) be a scheme, \(G\) an \(S\)-reductive group, \(P\) a parabolic
subgroup of \(G\), and \(\rad^u(P)\) its unipotent radical.  There is a
sequence of sub-group-schemes of \(\rad^u(P)\)
\[
 \rad^u(P)=U_0\supset U_1\supset U_2\supset\cdots
 \supset U_n\supset\cdots
\]
having the following properties:
\begin{enumerate}[label=\textup{(\roman*)}]
\item Every \(U_i\) is smooth, with connected fibers, characteristic and
closed in \(P\).  The commutator of a section of \(U_i\) and a section of
\(U_j\) is a section of \(U_{i+j+1}\), over a varying \(S'\to S\).
\item For every \(i\geq0\), there is a locally free
\(\mathcal O_S\)-module \(\mathcal E_i\) and an isomorphism of sheaves of
\(S\)-groups
\[
 U_i/U_{i+1}\xrightarrow{\sim}W(\mathcal E_i).
\]
Moreover, the automorphisms of \(P\), over a varying \(S'\to S\), act
linearly on \(W(\mathcal E_i)\).
\item For every \(s\in S\), one has \(U_{n,s}=e\) whenever
\[
 n>\dim\bigl(\rad^u(P)_s\bigr).
\]
\end{enumerate}
\end{sgastatement}
\SGARefTarget{sga3:xxvi:prop:2:1}

\paragraph{\SGARefLink{sga3:xxvi:prop:2:1}{2.1.1}.}
Suppose first that the pair \((G,P)\) admits a pinning.  Let
\((T,M,R,\Delta,\ldots)\) be a pinning of \(G\) adapted to \(P\), and let
\(\Delta(P)\subset\Delta\) be determined by \(P\).  Write
\(\alpha_1,\ldots,\alpha_p\) for the elements of \(\Delta(P)\), and
\(\beta_1,\ldots,\beta_q\) for the elements of
\(\Delta-\Delta(P)\).  Every root \(\gamma\in R\) has a unique expression
\[
 \gamma
 =a_1\alpha_1+\cdots+a_p\alpha_p
  +b_1\beta_1+\cdots+b_q\beta_q.
\]
Set
\[
 a(\gamma)=b_1+\cdots+b_q.
\tag{7}\label{sga3:xxvi:proposition:2-1:equation:7}
\]
\footnote{\emph{Editor's note.} The original
\(a_1+\cdots+a_p\) has been corrected to \(b_1+\cdots+b_q\).}
The definitions immediately give (cf.\ \SGARefLink{sga3:xxvi:proposition:1-12}{1.12}):
\[
\begin{aligned}
 U_\gamma\subset P
   &\ \Longleftrightarrow\ a(\gamma)\geq0,\\
 U_\gamma\subset\rad^u(P)
   &\ \Longleftrightarrow\ a(\gamma)>0,\\
 a(n\gamma+m\gamma')
   &=n\,a(\gamma)+m\,a(\gamma')
   \qquad(n,m\in\mathbb Z).
\end{aligned}
\]

For \(i\geq0\), let \(R_i\) be the set of roots \(\gamma\in R\) such that
\(a(\gamma)>i\).  Each \(R_i\) is a closed set of roots satisfying
\(R_i\cap(-R_i)=\varnothing\).  Consider (Exp.\ XXII, \SGARefLink{sga3:xxii:cor:5.6.5}{5.6.5}) the
\(S\)-group
\[
 U_i=U_{R_i}=\prod_{\gamma\in R_i}U_\gamma.
\]
It is a closed sub-group-scheme of \(G\), smooth over \(S\), with connected
fibers.

Let \(\alpha,\beta\in R\).  Consider the commutation relation of
Exp.\ XXII, \SGARefLink{sga3:xxii:cor:5.5.2}{5.5.2}:
\[
 p_\alpha(x)p_\beta(y)p_\alpha(-x)
 =
 p_\beta(y)
 \prod_{n,m\in\mathbb N^*}
 p_{n\alpha+m\beta}
 \bigl(C_{n,m,\alpha,\beta}x^ny^m\bigr),
\]
where every
\[
 p_\gamma:\mathbb G_{a,S}\xrightarrow{\sim}U_\gamma
\]
is an isomorphism of vector groups.  If \(a(\alpha)>i\) and
\(a(\beta)>j\), then, for \(n,m>0\),
\[
\begin{aligned}
 a(n\alpha+m\beta)
 &=n\,a(\alpha)+m\,a(\beta)\\
 &\geq n(i+1)+m(j+1)>i+j+1.
\end{aligned}
\]
Hence the commutator of a section of \(U_\alpha\) and a section of
\(U_\beta\) belongs to \(U_{i+j+1}\); therefore
\[
 (U_i,U_j)\subset U_{i+j+1}.
\]
In particular \(U_i/U_{i+1}\) is commutative, and identifies naturally
with
\[
 U_i/U_{i+1}
 \simeq
 \prod_{a(\gamma)=i+1}U_\gamma
 \simeq W(\mathcal E_i),
\]
where \(\mathcal E_i\) is the direct sum of the
\(\mathfrak g^\gamma\) for which \(a(\gamma)=i+1\).

Return to the commutation formula and suppose
\(a(\alpha)\geq0\), \(a(\beta)>i\).  If \(n,m\in\mathbb N^*\), then
either
\[
 a(n\alpha+m\beta)>i+1,
\]
or \(a(n\alpha+m\beta)=i+1\), in which case necessarily \(m=1\).
This first proves that \(\Int(p_\alpha(x))\) preserves \(U_i\), and hence
also \(U_{i+1}\).  It then shows that, modulo \(U_{i+1}\), the expression
for \(\Int(p_\alpha(x))p_\beta(y)\) contains only terms
\[
 p_{n\alpha+\beta}
 \bigl(C_{n,1,\alpha,\beta}x^ny\bigr),
\]
which are linear in \(y\).  Thus the inner automorphisms defined by
sections of \(U_\alpha\) act linearly on
\(U_i/U_{i+1}\simeq W(\mathcal E_i)\).  The same is clear for inner
automorphisms defined by sections of \(T\).  Since \(P\) is generated by
\(T\) and by the \(U_\alpha\) with \(a(\alpha)\geq0\), it follows that:
\begin{enumerate}[label=\textup{(\roman*)}]
\item every \(U_i\) is invariant in \(P\);
\item the inner automorphisms defined by sections of \(P\) act linearly on
\(U_i/U_{i+1}\simeq W(\mathcal E_i)\).
\end{enumerate}
\SGARefTarget{sga3:xxvi:item:2:1:1}

\paragraph{\SGARefLink{sga3:xxvi:item:2:1:1}{2.1.2}.}
Let \((T',M',R',\Delta',\ldots)\) now be another pinning of \(G\) adapted
to \(P\).  By \SGARefLink{sga3:xxvi:prop:1:15}{1.15} there is an inner automorphism of \(G\), coming from
\(P\), that carries the first pinning to the second.  After extending
\(S\), it may be written \(\Int(p)\) for some \(p\in P(S)\).
Repeating the construction with the second pinning produces groups
\(U'_i\) and vector-group isomorphisms
\[
 U'_i/U'_{i+1}\simeq W(\mathcal E'_i)
\]
obtained from \(U_i\) and
\(U_i/U_{i+1}\simeq W(\mathcal E_i)\) by transport of structure under
\(\Int(p)\).  The preceding invariance and linearity statements therefore
give \(U'_i=U_i\), and the two vector-space structures on
\[
 U_i/U_{i+1}=U'_i/U'_{i+1}
\]
coincide.

Thus the \(U_i\) and the vector structures on their successive quotients
are independent of the chosen pinning and, in particular, invariant under
every automorphism of \(P\).  This proves the proposition when \((G,P)\)
admits a pinning.  Assertion (iii) is immediate: by Exp.\ XXI, \SGARefLink{sga3:XXI:3.1.2}{3.1.2}, the
set \(\{a(\gamma)\mid\gamma\in R\}\) is an interval in \(\mathbb Z\), so
\(\dim U_{i,s}=\dim U_{i+1,s}\) cannot occur unless \(U_{i,s}=e\).
\SGARefTarget{sga3:xxvi:item:2:1:2}

\paragraph{\SGARefLink{sga3:xxvi:item:2:1:2}{2.1.3}.}
In general there is an fpqc covering family \(\{S_j\to S\}\) such that
every pair \((G_{S_j},P_{S_j})\) admits a pinning (\SGARefLink{sga3:xxvi:lemma:1:14}{1.14}).  What precedes
gives descent data on the \(\rad^u(P_{S_j})_i\), compatible with the
vector structures on the quotients.  The result follows by descent for
closed subschemes and for locally free modules.\footnote{\emph{Editor's
note.} See SGA~1, \SGARefLink{sga3:viii:corollary:1:3}{Exposé VIII, 1.3}, \SGARefLink{sga3:xxvi:cor:1:9}{1.9}, and \SGARefLink{sga3:xxvi:cor:1:10}{1.10}.}
\SGARefTarget{sga3:xxvi:item:2:1:3}

\begin{sgastatement}{Corollary 2.2}
Let \(S\) be affine, \(G\) an \(S\)-reductive group, and \(P\) a parabolic
subgroup of \(G\).  Then
\[
 H^1(S,\rad^u(P))=0;
\]
that is, every principal homogeneous bundle under \(\rad^u(P)\) is
trivial.
\end{sgastatement}
\SGARefTarget{sga3:xxvi:cor:2:2}

Indeed, \(S\) is the disjoint union of subschemes on each of which
\(\rad^u(P)\) has constant relative dimension.  By 2.1(iii) one may thus
assume that \(U_n=e\) for some \(n\).  Since
\[
 H^1(S,U_i/U_{i+1})
 =H^1(S,W(\mathcal E_i))=0
\]
by TDTE I, B, \SGARefLink{sga3:xxvi:def:1:1}{1.1} (or SGA~1, \SGARefLink{sga3:xi:corollary:5:1}{Exposé XI, 5.1}), the conclusion follows.

\begin{sgastatement}{Corollary 2.3}
Under the preceding hypotheses, \(P\) has a Levi subgroup \(L\).  If
\(L\) is any Levi subgroup of \(P\), the canonical map
\[
 H^1(S,L)\longrightarrow H^1(S,P)
\]
is bijective \textup{(see the introduction to Exp.\ XXIV for the
definition of \(H^1(S,-)\))}.
\end{sgastatement}
\SGARefTarget{sga3:xxvi:cor:2:3}

The first assertion follows from \SGARefLink{sga3:xxvi:cor:2:2}{2.2} and \SGARefLink{sga3:xxvi:cor:1:9}{1.9}.  The map
\(H^1(S,L)\to H^1(S,P)\) is surjective because
\(P=L\cdot\rad^u(P)\).  For injectivity it is enough to show that
\[
 H^1\bigl(S,\rad^u(P)_Q\bigr)=0
\]
for every principal homogeneous \(L\)-bundle \(Q\), where the subscript
denotes twisting by \(Q\).  One may either repeat the proof of \SGARefLink{sga3:xxvi:cor:2:2}{2.2}, using
the \(L\)-invariance of the vector structures on \(U_i/U_{i+1}\), or
observe that \(\rad^u(P)_Q\) is the unipotent radical of the parabolic
subgroup \(P_Q\) of \(G_Q\) and apply \SGARefLink{sga3:xxvi:cor:2:2}{2.2} to \(P_Q\).

\begin{sgastatement}{Corollary 2.4}
Let \(S\) be a semi-local scheme, \(G\) an \(S\)-reductive group, and
\(P\) a parabolic subgroup of \(G\).  There is a maximal torus \(T\) of
\(G\) contained in \(P\).
\end{sgastatement}
\SGARefTarget{sga3:xxvi:cor:2:4}

Indeed, by \SGARefLink{sga3:xxvi:cor:2:3}{2.3} the group \(P\) has a Levi subgroup \(L\), and it suffices
to know that \(L\) has a maximal torus, which follows from Exp.\ XIV,
\SGARefLink{sga3:xiv:3.20}{3.20}.

\begin{sgastatement}{Corollary 2.5}
Let \(S\) be affine, \(G\) an \(S\)-reductive group, and \(P\) a parabolic
subgroup of \(G\).  There is a locally free \(\mathcal O_S\)-module
\(\mathcal E\) such that \(\rad^u(P)\), as an \(S\)-scheme, is isomorphic
to \(W(\mathcal E)\).
\end{sgastatement}
\SGARefTarget{sga3:xxvi:cor:2:5}

We prove by induction on \(i\) that there is an isomorphism of \(S\)-schemes
\[
 \rad^u(P)/U_i
 \simeq
 W(\mathcal E_0\oplus\mathcal E_1\oplus\cdots\oplus\mathcal E_{i-1}).
\]
This is clear for \(i=0\).  If \(i>0\), then \(\rad^u(P)/U_i\) is a
principal homogeneous bundle over \(\rad^u(P)/U_{i-1}\) under
\[
 (U_{i-1}/U_i)_{\rad^u(P)/U_{i-1}}
 \simeq
 W\!\left(
 \mathcal E_{i-1}\otimes
 \mathcal O_{\rad^u(P)/U_{i-1}}
 \right).
\]
The base is affine, for example by induction, so the bundle is trivial
(TDTE~I or SGA~1, Exposé XI, \emph{loc.\ cit.}).  Hence
\[
 \rad^u(P)/U_i
 \simeq
 \bigl(\rad^u(P)/U_{i-1}\bigr)\times_S(U_{i-1}/U_i),
\]
which completes the induction.

\begin{sgastatement}{Corollary 2.6}
Let \(S\) be semi-local, let \(\{s_i\}\) be its closed points, let \(G\)
be an \(S\)-reductive group, and let \(P\) be a parabolic subgroup of
\(G\).  The canonical map
\[
 \rad^u(P)(S)\longrightarrow
 \prod_i\rad^u(P)(\operatorname{Spec}\kappa(s_i))
\]
is surjective.
\end{sgastatement}
\SGARefTarget{sga3:xxvi:cor:2:6}

Indeed, write \(S=\operatorname{Spec}A\) and
\(\kappa(s_i)=A/\mathfrak p_i\).  If \(\mathcal E\) corresponds to the
projective, hence flat, \(A\)-module \(E\), it remains to prove that
\[
 E\longrightarrow\prod_i E\otimes_A A/\mathfrak p_i
\]
is surjective.  It suffices to do so for \(E=A\), when it is well known
(Bourbaki, \emph{Algèbre commutative}, Chapter~II, §1, no.~2,
Proposition~5).

\begin{sgastatement}{Corollary 2.7}
Let \(k\) be an infinite field, \(G\) a \(k\)-reductive group, and \(P\)
a parabolic subgroup of \(G\).  Then \(\rad^u(P)(k)\) is dense in
\(\rad^u(P)\).
\end{sgastatement}
\SGARefTarget{sga3:xxvi:cor:2:7}

\begin{sgastatement}{Corollary 2.8}
Let \(S\) be semi-local, let \(\{s_i\}\) be its closed points, let \(G\)
be an \(S\)-reductive group, let \(P\) be a parabolic subgroup of \(G\),
and, for every \(i\), let \(L_i\) be a Levi subgroup of \(P_{s_i}\).
There is a Levi subgroup \(L\) of \(P\) whose fiber at every \(s_i\) is
\(L_i\).
\end{sgastatement}
\SGARefTarget{sga3:xxvi:cor:2:8}

Indeed, let \(L_0\) be a Levi subgroup of \(P\) (\SGARefLink{sga3:xxvi:cor:2:3}{2.3}).  For every \(i\),
choose
\[
 u_i\in\rad^u(P)(\operatorname{Spec}\kappa(s_i))
\quad\text{with}\quad
 \Int(u_i)L_{0,s_i}=L_i
\]
by \SGARefLink{sga3:xxvi:cor:1:8}{1.8}.  If \(u\in\rad^u(P)(S)\) induces the \(u_i\), as \SGARefLink{sga3:xxvi:cor:2:6}{2.6} permits,
then \(L=\Int(u)L_0\) has the required fibers.

\begin{sgastatement}{\SGAInlineRefTarget{sga3:xxvi:corollary:2-9}Corollary 2.9}
In the setting of \SGARefLink{sga3:xxvi:proposition:2-1}{2.1}, let \(H\) also be a sub-group-scheme of \(G\),
smooth and of finite presentation over \(S\), with connected fibers, such
that \(P\cap H\) contains, locally for the fpqc topology, a maximal torus
of \(G\).  For every \(i\geq0\) there is a locally direct-factor
submodule \(\mathcal F_i\subset\mathcal E_i\) for which the isomorphism
\[
 U_i/U_{i+1}\xrightarrow{\sim}W(\mathcal E_i)
\]
induces an isomorphism
\[
 (U_i\cap H)/(U_{i+1}\cap H)\xrightarrow{\sim}W(\mathcal F_i).
\]
\end{sgastatement}
\SGARefTarget{sga3:xxvi:cor:2:9}

Indeed, \(H\) is a subgroup of \(G\) of type (R) (Exp.\ XXII, \SGARefLink{sga3:xxii:def:5.2.1}{5.2.1}).
The assertion is fpqc-local, so we may suppose that \(G\) is split
relative to a maximal torus of \(P\cap H\), and even reduce to \SGARefLink{sga3:xxvi:prop:2:1}{2.1.1} with
\(H\) defined by a subset \(R'\subset R\).  In the notation used there,
Exp.\ XXII, \SGARefLink{sga3:xxii:cor:5.6.7:item:ii}{5.6.7(ii)} gives
\[
 U_i\cap H=\prod_{\alpha\in R_i\cap R'}U_\alpha.
\]
Thus \((U_i\cap H)/(U_{i+1}\cap H)\) identifies with
\[
 \prod_{\substack{\alpha\in R'\\a(\alpha)=i+1}}U_\alpha,
\]
which proves the assertion.

\begin{sgastatement}{Corollary 2.10}
In the setting of \SGARefLink{sga3:xxvi:corollary:2-9}{2.9}, the conclusions of \SGARefLink{sga3:xxvi:cor:2:2}{2.2}, \SGARefLink{sga3:xxvi:cor:2:5}{2.5}, \SGARefLink{sga3:xxvi:cor:2:6}{2.6}, and \SGARefLink{sga3:xxvi:cor:2:7}{2.7} remain
valid with \(\rad^u(P)\) replaced by \(\rad^u(P)\cap H\).
\end{sgastatement}
\SGARefTarget{sga3:xxvi:cor:2:10}

\begin{sgastatement}{Corollary 2.11}
Let\footnote{\emph{Editor's note.} \SGARefLink{sga3:xxvi:cor:2:11}{Corollary 2.11} has been added; it is
used in \SGARefLink{sga3:xxvi:proposition:4-5-1}{4.5.1} and \SGARefLink{sga3:xxvi:cor:6:17}{6.17}.} \(G\) be an \(S\)-reductive group, \(P\) a parabolic subgroup, and
\(H\) a subgroup of \(P\) of type \textup{(RC)} such that
\[
 \rad^u(H)=\rad^u(P)\cap H.
\]
Then statements \SGARefLink{sga3:xxvi:cor:2:2}{2.2} through \SGARefLink{sga3:xxvi:cor:2:8}{2.8} remain valid with \(P\) replaced by
\(H\).
\end{sgastatement}
\SGARefTarget{sga3:xxvi:cor:2:11}

% END INCLUDED SGA3_English_current_buildable_source_20260823/source/inputs/XXVI/tex/components/01_section2.tex

% BEGIN INCLUDED SGA3_English_current_buildable_source_20260823/source/inputs/XXVI/tex/components/02_section3.tex
\SGASection{3}{Scheme of parabolic subgroups of a reductive group}
\SGARefTarget{sga3:xxvi:section:3}

\SGAInlineRefTarget{sga3:xxvi:paragraph:3-1}\paragraph{\SGARefLink{sga3:xxvi:paragraph:3-1}{3.1}.}
Let \(E\) be a finite twisted constant \(S\)-scheme (Exp.\ X, \SGARefLink{sga3:x:definition:5-1}{5.1}).
Consider the \(S\)-functor \(\Of(E)\), where \(\Of(E)(S')\) is the set of
open-and-closed subschemes of \(E_{S'}\), or equivalently the set of
open-and-closed subsets of \(E_{S'}\).  Then \(\Of(E)\) is representable
by a finite twisted constant \(S\)-scheme.  Indeed, if \(E=A_S\), for a
finite set \(A\), then
\[
 \Of(E)\simeq\mathcal P(A)_S,
\]
where \(\mathcal P(A)\) denotes the power set of \(A\); the general case
follows by descent of open-and-closed subschemes.  Plainly,
\[
 \Of(E_{S'})=\Of(E)_{S'},
 \qquad
 \Of(E\times_SE')=\Of(E)\times_S\Of(E').
\]
\SGARefTarget{sga3:xxvi:item:3:1}

\paragraph{\SGARefLink{sga3:xxvi:item:3:1}{3.2}.}
Let \(S\) be a scheme and \(G\) an \(S\)-reductive group.  Define the
functor \(\Par(G)\) of parabolic subgroups of \(G\) by
\[
 \Par(G)(S')
 =\{\text{parabolic subgroups of }G_{S'}\}.
\]
In particular \(G\in\Par(G)(S)\), and
\(\underline{\operatorname{Bor}}(G)\subset\Par(G)\).  We shall define a
morphism
\[
 t:\Par(G)\longrightarrow\Of(\underline{\operatorname{Dyn}}(G))
\]
with the following properties:
\begin{enumerate}[label=\textup{(\roman*)}]
\item \(t\) is functorial in \(G\), with respect to isomorphisms, and
commutes with base change.
\item If \((T,M,R,\Delta,\ldots)\) is a pinning of \(G\) adapted to a
parabolic subgroup \(P\), then the canonical isomorphism
\[
 \underline{\operatorname{Dyn}}(G)\simeq\Delta_S
\qquad\textup{(Exp.\ XXIV, 3.4(iii))}
\]
carries \(t(P)\) onto \(\Delta(P)_S\), in the notation of \SGARefLink{sga3:xxvi:definition:1-11}{1.11} and \SGARefLink{sga3:xxvi:proposition:1-12}{1.12}.
\end{enumerate}

First suppose that \(P\) is a parabolic subgroup of \(G\) and that
\((T,M,R,\Delta,\ldots)\) is a pinning of \(G\) adapted to \(P\).  Define
\(t(P)\) by (ii).  The resulting subscheme of
\(\underline{\operatorname{Dyn}}(G)\) is independent of the pinning.
Indeed, if \((T',M',R',\Delta',\ldots)\) is another pinning adapted to
\(P\), then the unique inner automorphism of \(G\) carrying the first
pinning to the second comes from \(P\), by \SGARefLink{sga3:xxvi:prop:1:15}{1.15}.  The canonical
isomorphism \(\Delta\xrightarrow{\sim}\Delta'\) therefore carries
\(\Delta(P)\) onto \(\Delta'(P)\).

Without assuming that \((G,P)\) admits a pinning, \SGARefLink{sga3:xxvi:lemma:1:14}{1.14} and the definition
of \(\underline{\operatorname{Dyn}}(G)\) (Exp.\ XXIV, \SGARefLink{sga3:xxiv:3.3}{3.3}) allow one to
descend a unique open-and-closed subscheme \(t(P)\) such that
\[
 t(P)_{S'}=t(P_{S'})
\]
whenever \((G,P)_{S'}\) admits a pinning.
\SGARefTarget{sga3:xxvi:item:3:2}

\begin{sgastatement}{Theorem 3.3}
Let \(S\) be a scheme, \(G\) an \(S\)-reductive group, and
\[
 t:\Par(G)\longrightarrow\Of(\underline{\operatorname{Dyn}}(G))
\]
the morphism just defined.
\begin{enumerate}[label=\textup{(\roman*)}]
\item Two parabolic subgroups \(P,P'\) of \(G\) are conjugate locally for
the fpqc topology \textup{(cf.\ \SGARefLink{sga3:xxvi:proposition:1-3}{1.3})} if and only if \(t(P)=t(P')\).
\item\SGAInlineRefTarget{sga3:xxvi:thm:3:3:item:ii} \(\Par(G)\) is representable, and \(t\) is smooth and projective,
with integral geometric fibers.
\end{enumerate}
\end{sgastatement}
\SGARefTarget{sga3:xxvi:thm:3:3}

By 3.2(i), together with the fact that inner automorphisms of \(G\) act
trivially on \(\underline{\operatorname{Dyn}}(G)\)
(Exp.\ XXIV, \SGARefLink{sga3:xxiv:3.4:item:iv}{3.4(iv)}), conjugate parabolics have the same type.
Conversely, suppose \(t(P)=t(P')\).  Locally one may first assume that
both \((G,P)\) and \((G,P')\) admit pinnings (\SGARefLink{sga3:xxvi:lemma:1:14}{1.14}).  By the conjugacy of
pinnings in \(G\) (Exp.\ XXIV, \SGARefLink{sga3:xxiv:1.5}{1.5}), one may then choose a single pinning
\((T,M,R,\Delta)\) adapted to both \(P\) and \(P'\).  Equality of types
gives \(\Delta(P)=\Delta(P')\), hence \(P=P'\) by \SGARefLink{sga3:xxvi:prop:1:4:item:v}{1.4(v)}.  This proves
(i).

For (ii), use the notation of Exp.\ XXII, \SGARefLink{sga3:xxii:prop:5.11.5}{5.11.5}.\footnote{\emph{Editor's
note.} Recall that \(\mathscr H_c=\mathscr H_c(G)\) denotes the functor of
subgroups of \(G\) of type (RC), that
\(\mathcal C\ell_c=\mathcal C\ell_c(G)\) denotes the functor of their
``conjugacy classes'', and that
\(c\ell:\mathscr H_c\to\mathcal C\ell_c\) is the canonical projection.}
There is a canonical morphism \(\Par(G)\to\mathscr H_c\), which fits into
the cartesian square
\[
\begin{tikzcd}[column sep=large,row sep=large]
 \Par(G) \arrow[r,"t"] \arrow[d,hook]
   & \Of(\underline{\operatorname{Dyn}}(G)) \arrow[d,hook]\\
 \mathscr H_c \arrow[r,"c\ell"]
   & \mathcal C\ell_c .
\end{tikzcd}
\]
The vertical arrows are monomorphisms.  By \emph{loc.\ cit.},
\(\mathscr H_c\) is representable and \(c\ell\) is smooth,
quasi-projective, and of finite presentation, with integral geometric
fibers; hence the same is true of \(t\).

It remains to prove that \(t\) is proper.  This is fpqc-local, so one may
reduce to a pinned group \(G=(G,T,M,R,\Delta,\ldots)\).  Then
\(\underline{\operatorname{Dyn}}(G)\simeq\Delta_S\), and it is enough to
show that
\[
 t^{-1}((\Delta_1)_S)
\]
is proper over \(S\) for every \(\Delta_1\subset\Delta\).  Let \(P_1\)
be the parabolic containing \(T\) with \(\Delta(P_1)=\Delta_1\).  Part
(i) shows that \(g\mapsto\Int(g)P_1\) induces an isomorphism
\[
 G/\Norm_G(P_1)\xrightarrow{\sim}t^{-1}((\Delta_1)_S).
\]
By \SGARefLink{sga3:xxvi:prop:1:2}{1.2}, \(G/\Norm_G(P_1)=G/P_1\), which is projective over \(S\).

\begin{sgastatement}{Definition 3.4}
\(\Of(\underline{\operatorname{Dyn}}(G))\) is called the
\emph{scheme of types of parabolics} of \(G\), and \(t(P)\) is called
the \emph{type} of \(P\).
\end{sgastatement}
\SGARefTarget{sga3:xxvi:def:3:4}

\begin{sgastatement}{Corollary 3.5}
The \(S\)-functor \(\Par(G)\) is representable by an \(S\)-scheme that is
smooth and projective over \(S\).  The factorization
\[
 \Par(G)\longrightarrow
 \Of(\underline{\operatorname{Dyn}}(G))
 \longrightarrow S
\]
is the Stein factorization \textup{(EGA III, 4.3.3)} of the structural
morphism \(\Par(G)\to S\).
\end{sgastatement}
\SGARefTarget{sga3:xxvi:cor:3:5}

\begin{sgastatement}{\SGAInlineRefTarget{sga3:xxvi:corollary:3-6}Corollary 3.6}
For every
\(t\in\Of(\underline{\operatorname{Dyn}}(G))(S)\), the \(S\)-scheme
\[
 \Par_t(G)=t^{-1}(t)
\]
of parabolic subgroups of type \(t\) is smooth and projective over \(S\)
and homogeneous under \(G\).  If \(P\) is a parabolic subgroup of \(G\),
there is a canonical isomorphism
\[
 G/P\xrightarrow{\sim}\Par_{t(P)}(G).
\]
Moreover,
\[
 \Par_{\varnothing}(G)=\underline{\operatorname{Bor}}(G),
 \qquad
 \Par_{\underline{\operatorname{Dyn}}(G)}(G)\xrightarrow{\sim}S.
\]
\end{sgastatement}
\SGARefTarget{sga3:xxvi:cor:3:6}

\begin{sgastatement}{\SGAInlineRefTarget{sga3:xxvi:remark:3-7}Remark 3.7}
The \(S\)-scheme \(\Of(\underline{\operatorname{Dyn}}(G))\) has a natural
order structure, the domination relation being set-theoretic inclusion of
subschemes.  It is a lattice; in particular, the infimum of two
open-and-closed subschemes of
\(\underline{\operatorname{Dyn}}(G)_{S'}\) is their intersection.

If \(B\) is a Borel subgroup of \(G\), let \(X\) be the functor of
parabolic subgroups containing \(B\).  The morphism
\[
 X\longrightarrow\Of(\underline{\operatorname{Dyn}}(G))
\]
induced by \(t\) is an isomorphism of ordered schemes.  This follows from
\(P\subset Q\Rightarrow t(P)\subset t(Q)\) and the next lemma.
\end{sgastatement}
\SGARefTarget{sga3:xxvi:rem:3:7}

\begin{sgastatement}{Lemma 3.8}
Let \(S\) be a scheme, \(G\) an \(S\)-reductive group, \(P\) a parabolic
subgroup of \(G\), and
\[
 t'\in\Of(\underline{\operatorname{Dyn}}(G))(S)
\quad\text{with}\quad
 t(P)\subset t'.
\]
There is a unique parabolic subgroup \(P'\) of \(G\) containing \(P\)
and satisfying \(t(P')=t'\).
\end{sgastatement}
\SGARefTarget{sga3:xxvi:lemma:3:8}

After extending the base, one may suppose that \(P\) contains a Borel
subgroup \(B\) of \(G\).  Uniqueness then follows from \SGARefLink{sga3:xxvi:prop:1:17}{1.17}.  For
existence one may work in the split case, where the assertion is immediate
from §1.

\begin{sgastatement}{Remarks 3.8.1}
\begin{enumerate}[label=\textup{(\roman*)}]
\item The analogue of \SGARefLink{sga3:xxvi:lemma:3:8}{3.8} obtained by reversing the inclusions is false.
It would imply, for example, that every group of type \(A_1\) has a Borel
subgroup, which is not so \textup{(Exp.\ XX, §5)}.
\item The preceding discussion shows that \(t(P)\subset t(Q)\) means that
locally for the fpqc, or étale, topology, \(P\) is conjugate to a subgroup
of \(Q\).  It is enough to check this on geometric fibers.  \SGARefLink{sga3:xxvi:section:5}{Section~5}
will show that the étale topology may be replaced by the Zariski topology.
\end{enumerate}
\end{sgastatement}
\SGARefTarget{sga3:xxvi:rem:3:8:1}

\SGAInlineRefTarget{sga3:xxvi:paragraph:3-9}\paragraph{\SGARefLink{sga3:xxvi:paragraph:3-9}{3.9}.}
The preceding discussion can be repeated for critical subgroups.  Recall
(Exp.\ XXII, \SGARefLink{sga3:xxii:def:5.10.4}{5.10.4} and \SGARefLink{sga3:xxii:cor:5.10.5}{5.10.5}) that a reductive subgroup \(H\) of a
reductive group \(G\) is \emph{critical} if
\[
 H=\Centr_G(\rad(H)),
\]
that a subtorus \(Q\subset G\) is a \emph{\(C\)-critical torus} if
\[
 Q=\rad(\Centr_G(Q)),
\]
and that critical subgroups and \(C\)-critical tori are in bijection via
\[
 H\longmapsto\rad(H),
 \qquad
 Q\longmapsto\Centr_G(Q).
\]
\footnote{\emph{Editor's note.} ``Critical torus'' has here been replaced
by ``\(C\)-critical torus''.  Below, ``critical torus'' is again used for
``\(C\)-critical torus''.}

If \((G,T,M,R)\) is split, the subgroup of \(G\) containing \(T\) and
corresponding to \(R'\subset R\) (Exp.\ XXII, \SGARefLink{sga3:xxii:def:5.4.2}{5.4.2}) is critical if and
only if \(R'\) is \emph{vectorial}, meaning that it is the intersection
of \(R\) with a vector subspace of \(M\otimes\mathbb Q\)
(Exp.\ XXII, \SGARefLink{sga3:xxii:prop:5.10.6}{5.10.6}).

For arbitrary \(G\), Exp.\ XXII, \SGARefLink{sga3:xxii:prop:5.11.5}{5.11.5} defines a finite étale \(S\)-scheme
\(\mathcal C\ell_{\mathrm{crit}}\); in the split case it is the constant
scheme associated with the vectorial subsets of \(R\), modulo the action
of the Weyl group.  If \(\Crit(G)\) denotes the functor of critical
subgroups, there is a canonical morphism
\[
 c\ell:\Crit(G)\longrightarrow\mathcal C\ell_{\mathrm{crit}}
\]
in the cartesian square
\[
\begin{tikzcd}[column sep=large,row sep=large]
 \Crit(G) \arrow[r,"c\ell"] \arrow[d,hook]
   & \mathcal C\ell_{\mathrm{crit}} \arrow[d,hook]\\
  \mathscr H_c \arrow[r,"c\ell"]
   & \mathcal C\ell_c .
\end{tikzcd}
\]
\footnote{\emph{Editor's note.} See \SGARefLink{sga3:xxvi:item:3:2}{3.3}, editor's note, for the notation
\(\mathscr H_c,\mathcal C\ell_c\).}
\SGARefTarget{sga3:xxvi:item:3:9}

\begin{sgastatement}{Proposition 3.10}
Let \(S\) be a scheme, \(G\) an \(S\)-reductive group,
\(\Crit(G)\) the functor of its critical subgroups, and let
\(\mathcal C\ell_{\mathrm{crit}}\) and
\[
 c\ell:\Crit(G)\longrightarrow\mathcal C\ell_{\mathrm{crit}}
\]
be as above.
\begin{enumerate}[label=\textup{(\roman*)}]
\item Critical subgroups \(H,H'\) of \(G\) are conjugate, locally for the
fpqc topology, if and only if \(c\ell(H)=c\ell(H')\).
\item \(\Crit(G)\) is representable, and \(c\ell\) is smooth and affine,
with integral geometric fibers.
\end{enumerate}
\end{sgastatement}
\SGARefTarget{sga3:xxvi:prop:3:10}

The proof is the same as for \SGARefLink{sga3:xxvi:item:3:2}{3.3} except for affineness of \(c\ell\).  It
is enough to prove that \(\Crit(G)\) is affine over \(S\).
The functor \(\Crit(G)\) identifies naturally with the functor of critical
tori of \(G\), giving a canonical monomorphism
\[
 \Crit(G)\longrightarrow\mathcal M,
\]
where \(\mathcal M\) is the scheme of subgroups of multiplicative type of
\(G\) (Exp.\ XI, \SGARefLink{sga3:xi:theorem:4:1}{4.1}).  By Exp.\ XII, \SGARefLink{sga3:xii:5.3}{5.3}, it is enough to show that this
is an open-and-closed immersion.  After the base change
\(\mathcal M\to S\), this reduces to the following: if \(Q\) is a subgroup
of multiplicative type in \(G\), the \(S'\to S\) for which \(Q_{S'}\) is
a critical torus of \(G_{S'}\) are precisely those factoring through an
open-and-closed subscheme of \(S\).  The condition is the conjunction of:
\begin{enumerate}[label=\textup{(\arabic*)}]
\item \(Q\) is a torus;
\item once \(Q\) is a torus, the torus \(\rad(\Centr_G(Q))\) has the same
relative dimension as \(Q\).
\end{enumerate}
Both conditions are of the required kind.

\begin{sgastatement}{Corollary 3.11}
The \(S\)-functor \(\Crit(G)\) is representable by an \(S\)-scheme that
is smooth and affine over \(S\).
\end{sgastatement}
\SGARefTarget{sga3:xxvi:cor:3:11}

\begin{sgastatement}{Corollary 3.12}
Let \(H\) be a critical subgroup of the \(S\)-reductive group \(G\).
Then \(G/\Norm_G(H)\) and \(G/H\) are representable by \(S\)-schemes that
are affine and smooth over \(S\).
\end{sgastatement}
\SGARefTarget{sga3:xxvi:cor:3:12}

The first assertion follows from \SGARefLink{sga3:xxvi:cor:3:11}{3.11}, and the second from the first and
Exp.\ XXII, \SGARefLink{sga3:xxii:cor:5.10.2}{5.10.2}.

\begin{sgastatement}{Corollary 3.13}
Let \(Q\) be a subtorus of the \(S\)-reductive group \(G\).  Then
\(G/\Centr_G(Q)\) is representable by an \(S\)-scheme that is smooth and
affine over \(S\).  The same holds for \(G/\Norm_G(Q)\) when \(Q\) is a
critical subtorus of \(G\).
\end{sgastatement}
\SGARefTarget{sga3:xxvi:cor:3:13}

Indeed, \(H=\Centr_G(Q)\) is critical (Exp.\ XXII, \SGARefLink{sga3:xxii:cor:5.10.5}{5.10.5}), and when
\(Q\) is critical one has
\(\Norm_G(H)=\Norm_G(Q)\) (\emph{loc.\ cit.}, 5.10.8).

\SGAInlineRefTarget{sga3:xxvi:paragraph:3-14}\paragraph{\SGARefLink{sga3:xxvi:paragraph:3-14}{3.14}.}
The conjugacy of Levi subgroups of parabolic subgroups of \(G\) gives a
unique morphism
\[
 u:\Of(\underline{\operatorname{Dyn}}(G))
 \longrightarrow\mathcal C\ell_{\mathrm{crit}}
\]
such that, for every parabolic subgroup \(P\) and every Levi subgroup
\(L\subset P\),
\[
 c\ell(L)=u(t(P)),
\]
compatibly with every base change.
\SGARefTarget{sga3:xxvi:item:3:14}

\paragraph{\SGARefLink{sga3:xxvi:item:3:14}{3.15}.}
Let \(S\) be a scheme and \(G\) an \(S\)-reductive group.  Consider the
\(S\)-functors\footnote{\emph{Editor's note.} The original \(LT\) has
been changed to \(CT\).}
\[
\begin{aligned}
 PL(S')={}&\{(P,L)\mid P\supset L,\ P\text{ parabolic in }G_{S'},\
                    L\text{ a Levi subgroup of }P\},\\
 PT(S')={}&\{(P,T)\mid P\supset T,\ P\text{ parabolic in }G_{S'},\
                    T\text{ a maximal torus of }P\},\\
 CT(S')={}&\{(C,T)\mid C\supset T,\ C\text{ critical in }G_{S'},\
                    T\text{ a maximal torus of }C\},\\
 PLT(S')={}&\{(P,L,T)\mid P\supset L\supset T,\ (P,T)\in PT(S'),\
                    L\text{ a Levi subgroup of }P\}.
\end{aligned}
\]
There are evident morphisms between these functors and
\(\Par(G),\Crit(G)\), and \(\Tor(G)\).  They form the following
commutative truncated cube.

\begin{figure}[htbp]
\centering
\begin{tikzpicture}[
  >=to,
  every node/.style={font=\small},
  morph/.style={->,line width=.45pt}
]
  \node (CT)     at (5.2,7.2) {\(\underline{CT}\)};
  \node (PLT)    at (9.0,7.2) {\(\underline{PLT}\)};
  \node (Crit)   at (0.8,5.2) {\(\Crit(G)\)};
  \node (PL)     at (7.1,5.2) {\(\underline{PL}\)};
  \node (Cl)     at (1.0,3.1) {\(\mathcal C\ell_{\mathrm{crit}}\)};
  \node (Tor)    at (5.2,3.1) {\(\Tor(G)\)};
  \node (PT)     at (9.0,3.1) {\(\underline{PT}\)};
  \node (S)      at (3.3,1.4) {\(S\)};
  \node (Of)     at (2.0,-.2) {\(\Of(\underline{\operatorname{Dyn}}(G))\)};
  \node (Par)    at (7.1,-.2) {\(\Par(G)\)};

  \draw[morph] (PLT) -- node[above]{\(g\) (ét.)} (CT);
  \draw[morph] (CT) -- node[above left]{\(a\)} node[below right]{(aff.)} (Crit);
  \draw[morph] (PLT) -- node[above left]{\(b\)} node[below right]{(aff.)} (PL);
  \draw[morph] (PL) -- node[above]{\(f\) (ét.)} (Crit);
  \draw[morph] (CT) -- node[right]{\(k\) (ét.)} (Tor);
  \draw[morph] (PLT) -- node[left]{\(\wr\)}
    node[right]{\(e\) (isom.)} (PT);
  \draw[morph] (Crit) -- node[left]{\(c\ell\) (aff.)} (Cl);
  \draw[morph] (PL) -- node[right]{\(d\) (aff.)} (Par);
  \draw[morph] (PT) -- node[above]{\(h\) (ét.)} (Tor);
  \draw[morph] (PT) -- node[right]{\(c\) (aff.)} (Par);
  \draw[morph] (S) -- node[above left]{\(q\) (ét.)} (Cl);
  \draw[morph] (Tor) -- node[above right]{\(r\) (aff.)} (S);
  \draw[morph] (Of) -- node[above left]{\(u\) (ét.)} (Cl);
  \draw[morph] (Of) -- node[right]{\(p\) (ét.)} (S);
  \draw[morph] (Par) -- node[above]{\(t\) (proj.)} (Of);
\end{tikzpicture}
\refstepcounter{figure}
\label{sga3:xxvi:fig:3:15:1}
\par\smallskip
{\small Figure 3.15.1}
\end{figure}
\SGARefTarget{sga3:xxvi:item:3:15}

\begin{sgastatement}{Theorem 3.16 \textup{(cf.\ Figure 3.15.1)}}
\begin{enumerate}[label=\textup{(\roman*)}]
\item Every morphism in the diagram is smooth, surjective, and of finite
presentation.
\item Every morphism except \(t\) is affine; \(t\) is projective.
\item Every morphism is either finite étale or has integral geometric
fibers: \(f,g,h,k,p,q,u\) are finite étale;
\(a,b,c,d,r,c\ell\) have integral geometric fibers; and \(e\) is an
isomorphism.
\item The square \((a,b,f,g)\) is cartesian.
\end{enumerate}
\end{sgastatement}
\SGARefTarget{sga3:xxvi:thm:3:16}

\SGAProof
It is clear from \SGARefLink{sga3:xxvi:proposition:1-6:item:ii}{1.6(ii)} that \(e\) is an isomorphism, and (iv) is
immediate.  The morphism \(a\) is smooth and affine with integral
geometric fibers.  Indeed, after the base change \(\Crit(G)\to S\), this
reduces to the same statement for
\[
 \Tor(L_0)\longrightarrow S,
\]
where \(L_0\) is the universal critical subgroup; \(L_0\) is reductive by
definition, so Exp.\ XXII, \SGARefLink{sga3:xxii:cor:5.8.3}{5.8.3} applies.  By (iv), \(b\) has the same
properties.

The morphism \(d\) is likewise smooth and affine with integral geometric
fibers, by \SGARefLink{sga3:xxvi:cor:1:9}{1.9}; the same is therefore true of
\[
 c=dbe^{-1}.
\]
The morphism \(r\) has these properties by Exp.\ XXII, \SGARefLink{sga3:xxii:cor:5.8.3}{5.8.3}, and
\(c\ell\) by \SGARefLink{sga3:xxvi:item:3:9}{3.10}.

We already know that \(p\) and \(q\) are finite étale and surjective
(\SGARefLink{sga3:xxvi:paragraph:3-1}{3.1} and \SGARefLink{sga3:xxvi:paragraph:3-9}{3.9}).  If \(f\) and \(k\) are finite étale and surjective, then
so is \(g\), by (iv), and so is \(h=kge^{-1}\).  Since the assertions
about \(t\) were proved in \SGARefLink{sga3:xxvi:thm:3:3:item:ii}{3.3(ii)}, it remains only to prove the claim for
\(f\) and \(k\).  We do so for \(k\), the proof for \(f\) being the same.

If \(T\) is a maximal torus of \(G\), let \(\mathcal C\) be the functor of
critical subgroups containing \(T\).  We must show that \(\mathcal C\) is
represented by a finite étale \(S\)-scheme with nonempty fibers.  We may
suppose that \(G\) is split relative to \(T\).  If \(E\) is the set of
vectorial subsets of the root system \(R\), then \(\mathcal C\) is
represented by \(E_S\), by 3.9.
\SGAQED

\begin{sgastatement}{Corollary 3.17}
All the functors in the preceding diagram are representable by
\(S\)-schemes that are smooth over \(S\), and all are affine over \(S\)
except \(\Par(G)\).
\end{sgastatement}
\SGARefTarget{sga3:xxvi:cor:3:17}

\begin{sgastatement}{\SGAInlineRefTarget{sga3:xxvi:remark:3-18}Remark 3.18}
\begin{enumerate}[label=\textup{(\roman*)}]
\item The fact that \(f:PL\to\Crit(G)\) is étale and surjective says that
a subgroup of \(G\) is critical if and only if, locally for the étale
topology, it is a Levi subgroup of a parabolic subgroup.

In general the map \(PL(S)\to\Crit(G)(S)\) need not be surjective: a
critical subgroup \(C\) may fail to arise over \(S\) from a parabolic
subgroup.  For example, a maximal torus need not be contained in a Borel
subgroup, as for a nonsplit form of \(\operatorname{SL}_2\)
\textup{(Exp.\ XX, §5)}.
\item Likewise,
\[
 u:\Of(\underline{\operatorname{Dyn}}(G))
 \longrightarrow\mathcal C\ell_{\mathrm{crit}}
\]
need not be an isomorphism.  Two parabolic subgroups of different types
may have Levi subgroups of the same type.  In a group of type \(A_2\),
for instance, the two types of parabolics whose Levi subgroups have
semisimple rank \(1\) correspond to the two vertices of the diagram,
whereas there is only one type of critical subgroup of rank \(1\).

The analogous example for type \(A_3\) shows that, even over an
algebraically closed field, nonisomorphic parabolic subgroups may have
Levi subgroups of the same type.\footnote{\emph{Editor's note.}
Concretely, in \(\operatorname{GL}_4\), the parabolics
\[
 P=\left\{\begin{pmatrix}
 *&*&*&*\\0&*&*&*\\0&0&*&*\\0&0&0&*
 \end{pmatrix}\right\},
 \qquad
 P'=\left\{\begin{pmatrix}
 *&*&*&*\\0&*&*&*\\0&*&*&*\\0&0&0&*
 \end{pmatrix}\right\}
\]
are not isomorphic (their unipotent radicals are not isomorphic), while
their Levi subgroups
\[
 L=\left\{\begin{pmatrix}
 *&*&0&0\\ *&*&0&0\\0&0&*&0\\0&0&0&*
 \end{pmatrix}\right\},
 \quad
 L'=\left\{\begin{pmatrix}
 *&0&0&0\\0&*&*&0\\0&*&*&0\\0&0&0&*
 \end{pmatrix}\right\}
\]
are conjugate by
\(
\left(\begin{smallmatrix}
0&0&1&0\\0&1&0&0\\1&0&0&0\\0&0&0&1
\end{smallmatrix}\right).
\)}
\end{enumerate}
\end{sgastatement}
\SGARefTarget{sga3:xxvi:rem:3:18}

We end this section with an application to principal bundles.

\begin{sgastatement}{\SGAInlineRefTarget{sga3:xxvi:lemma:3-20}Lemma 3.20}
Let\footnote{\emph{Editor's note.} The numbering of the original is
retained: there is no §3.19.} \(S\) be a scheme, \(G\) an \(S\)-reductive group, \(F\) a principal
homogeneous bundle under \(G\), and \(G_F\) the corresponding twisted form.
Identify \(\underline{\operatorname{Dyn}}(G)\) with
\(\underline{\operatorname{Dyn}}(G_F)\) \textup{(Exp.\ XXIV, \SGARefLink{sga3:xxiv:3.5}{3.5})}.
For every parabolic subgroup \(P\) of \(G\), there is a canonical
isomorphism
\[
 F/P\xrightarrow{\sim}\Par_{t(P)}(G_F).
\]
In particular, the structural group of \(F\) reduces to \(P\) if and only
if \(G_F\) has a parabolic subgroup of type \(t(P)\).
\end{sgastatement}
\SGARefTarget{sga3:xxvi:lemma:3:20}

The proof is exactly the same as that of Exp.\ XXIV, \SGARefLink{sga3:xxiv:proposition:4-2-1}{4.2.1}.

\SGAInlineRefTarget{sga3:xxvi:paragraph:3-21}\paragraph{\SGARefLink{sga3:xxvi:paragraph:3-21}{3.21}.}
If \(S\) is a scheme, \(G\) an \(S\)-reductive group, and
\[
 t\in\Of(\underline{\operatorname{Dyn}}(G))(S),
\]
write \(H^1_t(S,G)\) for the subset of \(H^1(S,G)\) consisting of the
classes of principal \(G\)-bundles \(F\) for which the associated group
\(G_F\) has a parabolic subgroup of type \(t\).  If \(G\) itself has a
parabolic subgroup \(P\) of type \(t\), then \(H^1_t(S,G)\) is precisely
the image of \(H^1(S,P)\) in \(H^1(S,G)\), and hence does not depend on
the chosen \(P\).
\SGARefTarget{sga3:xxvi:item:3:21}

% END INCLUDED SGA3_English_current_buildable_source_20260823/source/inputs/XXVI/tex/components/02_section3.tex

% BEGIN INCLUDED SGA3_English_current_buildable_source_20260823/source/inputs/XXVI/tex/components/03_section4_1_through_4_2_4.tex
\SGASection{4}{Relative position of two parabolic subgroups}
\SGARefTarget{sga3:xxvi:section:4}

\SGASubsection{4.1}{A preliminary result}
\SGARefTarget{sga3:xxvi:subsection:4:1}

\begin{sgastatement}{Lemma 4.1.1}
Let \(k\) be a field, \(G\) a reductive \(k\)-group, and \(P,P'\) two
parabolic subgroups of \(G\).
\begin{enumerate}[label=\textup{(\roman*)}]
\item The intersection \(P\cap P'\) is smooth, has the same reductive
rank as \(G\), and contains a maximal torus \(T\) of \(G\).
\item The set \(R'\) of roots of \(P\cap P'\) relative to \(T\) is a
closed subset of the set \(R\) of roots of \(G\) relative to \(T\).
\end{enumerate}
\end{sgastatement}
\SGARefTarget{sga3:xxvi:lemma:4:1:1}
\footnote{\emph{Editor's note.} Part \textup{(ii)}, which will be useful
in \SGARefLink{sga3:xxvi:proposition:4-5-1}{4.5.1}, was added by the editors.}

\SGAProof
Suppose first that \(k\) is algebraically closed.  Let \(B\) (respectively
\(B'\)) be a Borel subgroup of \(P\) (respectively \(P'\)).  There is an
element \(g\in G(k)\) such that \(\Int(g)B=B'\).  On the other hand, if
\(T_0\) is a maximal torus of \(B\) and \(N=\Norm_G(T_0)\), Bruhat's
theorem (\emph{Bible}, §13.4, Corollary 1 to Theorem 3) gives
\[
 G(k)=B(k)N(k)B(k).
\]
Thus there are \(b,b_1\in B(k)\) and \(n\in N(k)\) such that
\(g=bnb_1\), and hence
\(\Int(b)\Int(n)B=B'\).  Consequently
\[
 P\cap P'\supset B\cap B'
  =\Int(b)\bigl(B\cap\Int(n)B\bigr)
  \supset\Int(b)T_0.
\]

Now let \(k\) be arbitrary.  The preceding result shows that
\(P_{\bar k}\cap P'_{\bar k}\) contains a maximal torus of
\(G_{\bar k}\).  By Exp.\ XXII, \SGARefLink{sga3:xxii:prop:5.4.5}{5.4.5},
\((P\cap P')_{\bar k}\) is smooth along the identity section; since the
base is a field, it is therefore smooth (Exp.\ VI\textsubscript{A},
\SGARefLink{prop:sga3-VIA-1.3.1}{1.3.1}), and so \(P\cap P'\) is smooth.  By Exp.\ VI\textsubscript{A},
\SGARefLink{prop:sga3-VIA-2.3.1}{2.3.1}, the identity component \((P\cap P')^0\) is an open subgroup of
\(P\cap P'\), smooth over \(k\).  Applying Exp.\ XIV, \SGARefLink{sga3:xiv:1.1}{1.1} to it proves
\textup{(i)}.

Finally, the set \(R_P\) (respectively \(R_{P'}\)) of roots of \(P\)
(respectively \(P'\)) relative to \(T\) is closed, and
\[
 R'=R_P\cap R_{P'}.
\]
This proves \textup{(ii)}.
\SGAQED

\begin{sgastatement}{Remark 4.1.2}
One can prove \textup{([BT65], \SGARefLink{sga3:xxvi:subsection:4:5}{4.5})} that \(P\cap P'\) is connected; this
fact will be used in \SGARefLink{sga3:xxvi:proposition:4-5-1}{4.5.1}.
\end{sgastatement}
\SGARefTarget{sga3:xxvi:remark:4:1:2}
\footnote{\emph{Editor's note.} In view of the details added in \SGARefLink{sga3:xxvi:proposition:4-5-1}{4.5.1},
the original wording ``we shall not use this fact'' was changed.}

\begin{sgastatement}{\SGAInlineRefTarget{sga3:xxvi:remark:4-1-3}Remark 4.1.3}
The preceding lemma is false over an arbitrary scheme.  For example, let
\(G\) be a reductive group over an algebraically closed field \(k\), and
let \(B\) be a Borel subgroup of \(G\).  Take \(S=G\) as base and
consider the Borel subgroups \(B_1,B_2\) of \(G_G\), where
\[
 B_1=B_G,\qquad B_2=\Int(g_0)B_1,
\]
and \(g_0\) is the canonical (diagonal) section of \(G_G\).  For every
\(g\in G(k)\), the fiber of \(B_1\cap B_2\) at \(g\) is
\(B\cap\Int(g)B\).  If \(B\ne G\), the dimension of this fiber varies
with \(g\), so \(B_1\cap B_2\) cannot be smooth over \(G\).
\end{sgastatement}
\SGARefTarget{sga3:xxvi:remark:4:1:3}

\SGASubsection{4.2}{Transversal position}
\SGARefTarget{sga3:xxvi:subsection:4:2}

\begin{sgastatement}{\SGAInlineRefTarget{sga3:xxvi:theorem:4-2-1}Theorem 4.2.1}
Let \(S\) be a scheme, \(G\) an \(S\)-reductive group, and \(P,Q\) two
parabolic subgroups of \(G\).  The following conditions on the pair
\((P,Q)\) are equivalent:
\begin{enumerate}[leftmargin=3.2em,labelsep=.7em]
\item[\textup{\SGAInlineRefTarget{sga3:xxvi:theorem:4-2-1:item:i}(i)}]
\(\operatorname{Lie}(P/S)+\operatorname{Lie}(Q/S)
  =\operatorname{Lie}(G/S)\).
\item[\textup{\SGAInlineRefTarget{sga3:xxvi:theorem:4-2-1:item:ii}(ii)}]
The canonical morphism \(P\times_SQ\to G\) is smooth.
\item[\textup{(ii\(^{\prime}\))}]
The canonical morphism \(P\to G/Q\) is smooth.
\item[\textup{\SGAInlineRefTarget{sga3:xxvi:theorem:4-2-1:item:iii}(iii)}]
The canonical morphism \(P\times_SQ\to G\) is open.
\item[\textup{(iii\(^{\prime}\))}]
The canonical morphism \(P\to G/Q\) is open.
\item[\textup{\SGAInlineRefTarget{sga3:xxvi:theorem:4-2-1:item:iv}(iv)}]
The canonical morphism \(P\times_SQ\to G\) is fiberwise dominant.
\item[\textup{(iv\(^{\prime}\))}]
The canonical morphism \(P\to G/Q\) is fiberwise dominant.
\item[\textup{\SGAInlineRefTarget{sga3:xxvi:theorem:4-2-1:item:v}(v)}]
For every \(s\in S\), the group \(P_s\cap Q_s\) has minimum dimension,
that is,
\[
 \dim(P_s\cap Q_s)=\dim P_s+\dim Q_s-\dim G_s.
\]
\item[\textup{\SGAInlineRefTarget{sga3:xxvi:theorem:4-2-1:item:vi}(vi)}]
There is an étale covering family \(\{S_i\to S\}\) and, for every \(i\),
a Borel subgroup \(B_i\) of \(P_{S_i}\) and a Borel subgroup \(B'_i\)
of \(Q_{S_i}\), such that \(B_i\cap B'_i\) is a maximal torus of
\(G_{S_i}\).
\item[\textup{\SGAInlineRefTarget{sga3:xxvi:theorem:4-2-1:item:vii}(vii)}]
There is an étale covering family \(\{S_i\to S\}\) and, for every \(i\),
a splitting \((T_i,\ldots,R_i)\) of \(G_{S_i}\) and a system
\(R_i^+\) of positive roots of \(R_i\), such that \(P_{S_i}\)
(respectively \(Q_{S_i}\)) is the subgroup of type \textup{(R)} of
\(G_{S_i}\) containing \(T_i\) and defined by a subset \(R_i^{(1)}\)
(respectively \(R_i^{(2)}\)) of \(R_i\) containing \(R_i^+\)
(respectively \(-R_i^+\)); see Exp.\ XXII, \SGARefLink{sga3:xxii:def:5.4.2}{5.4.2} and 5.2.1.
\end{enumerate}
\end{sgastatement}
\SGARefTarget{sga3:xxvi:theorem:4:2:1}

\SGAProof
We prove the theorem according to the following logical diagram:
\[
\begin{tikzcd}[column sep=large,row sep=large]
 \textup{(i)} \arrow[d,Rightarrow] &
 \textup{(vii)} \arrow[l,Rightarrow] &
 \textup{(vi)} \arrow[l,Rightarrow] \\
 \textup{(ii\(^{\prime}\))} \arrow[d,Rightarrow] \arrow[r,Rightarrow] &
 \textup{(iii\(^{\prime}\))} \arrow[r,Rightarrow] &
 \textup{(iv\(^{\prime}\))} \arrow[r,Rightarrow] &
 \textup{(v)} \arrow[ul,Rightarrow] \\
 \textup{(ii)} \arrow[r,Rightarrow] &
 \textup{(iii)} \arrow[r,Rightarrow] &
 \textup{(iv)} \arrow[ur,Rightarrow] &
 {} .
\end{tikzcd}
\]

The implications \textup{\SGARefLink{sga3:xxvi:theorem:4-2-1:item:ii}{(ii)}} \(\Rightarrow\) \textup{\SGARefLink{sga3:xxvi:theorem:4-2-1:item:iii}{(iii)}} and
\textup{(ii\(^{\prime}\))} \(\Rightarrow\)
\textup{(iii\(^{\prime}\))} are immediate.  If \textup{\SGARefLink{sga3:xxvi:theorem:4-2-1:item:iii}{(iii)}} holds,
the set-theoretic image of \(P\times_SQ\to G\) is an open subset of
\(G\) containing the identity section.  Since the fibers of \(G\) are
connected, this image is dense in every fiber; thus \textup{\SGARefLink{sga3:xxvi:theorem:4-2-1:item:iv}{(iv)}}
holds.  The implication
\textup{(iii\(^{\prime}\))} \(\Rightarrow\)
\textup{(iv\(^{\prime}\))} is identical.

Condition \textup{(ii\(^{\prime}\))} implies \textup{\SGARefLink{sga3:xxvi:theorem:4-2-1:item:ii}{(ii)}}, by the
cartesian diagram
\[
\begin{tikzcd}[column sep=large,row sep=large]
 P\times_SQ \arrow[r] \arrow[d,"\operatorname{pr}_1"'] & G \arrow[d]\\
 P \arrow[r] & G/Q .
\end{tikzcd}
\]
On the other hand, \textup{(iv)} or \textup{(iv\(^{\prime}\))}
implies \textup{(v)} by dimension theory (EGA
IV\textsubscript{2}, 5.6.6).  Indeed, one may assume
\(S=\operatorname{Spec}(k)\), with \(k\) algebraically closed; every
nonempty fiber of the morphism in \textup{\SGARefLink{sga3:xxvi:theorem:4-2-1:item:iv}{(iv)}} at a point of \(G(k)\),
or of the morphism in \textup{(iv\(^{\prime}\))} at a point of
\((G/Q)(k)\), is isomorphic to \(P\cap Q\).

The implication \textup{\SGARefLink{sga3:xxvi:theorem:4-2-1:item:vi}{(vi)}} \(\Rightarrow\) \textup{\SGARefLink{sga3:xxvi:theorem:4-2-1:item:vii}{(vii)}} follows
from Exp.\ XXII, \SGARefLink{sga3:xxii:prop:5.5.1}{5.5.1} \textup{\SGARefLink{sga3:xxvi:theorem:4-2-1:item:iv}{(iv)}} and \SGARefLink{sga3:xxii:prop:5.9.2}{5.9.2}.  Condition
\textup{\SGARefLink{sga3:xxvi:theorem:4-2-1:item:vii}{(vii)}} implies \textup{\SGARefLink{sga3:xxvi:theorem:4-2-1:item:i}{(i)}}, because the latter may be checked
fpqc-locally.  Under \textup{\SGARefLink{sga3:xxvi:theorem:4-2-1:item:vii}{(vii)}} we may therefore suppose that
\(G\) is split, \(P\supset B_{R^+}\), and
\(Q\supset B_{-R^+}\); then already
\[
 \operatorname{Lie}(B_{R^+})+\operatorname{Lie}(B_{-R^+})
 =\operatorname{Lie}(G).
\]

Let us prove that \textup{\SGARefLink{sga3:xxvi:theorem:4-2-1:item:i}{(i)}} implies
\textup{(ii\(^{\prime}\))}.  Let \(u:P\to G/Q\) be the canonical
morphism.  Smoothness can be checked on geometric fibers, since both
\(P\) and \(G/Q\) are smooth over \(S\); hence we may suppose that
\(S\) is the spectrum of an algebraically closed field.  The morphism
\(u\) is compatible with the evident action of \(P\):
\(u(pp')=p\,u(p')\).  A translation argument (compare the proof of
Exp.\ VI\textsubscript{B}, 1.3) therefore reduces smoothness to the
point \(e\in P(k)\).  By SGA 1, II, 4.7, it is enough that the tangent
map at \(e\) be surjective.  This map is naturally identified with
\[
 \operatorname{Lie}(P)\longrightarrow
 \operatorname{Lie}(G)/\operatorname{Lie}(Q),
\]
which is surjective under \textup{\SGARefLink{sga3:xxvi:theorem:4-2-1:item:i}{(i)}}.

It remains to prove \textup{\SGARefLink{sga3:xxvi:theorem:4-2-1:item:v}{(v)}} \(\Rightarrow\) \textup{\SGARefLink{sga3:xxvi:theorem:4-2-1:item:vi}{(vi)}}.
Suppose first that \(S\) is the spectrum of an algebraically closed
field.  By \SGARefLink{sga3:xxvi:lemma:4:1:1}{4.1.1} there is a maximal torus \(T\) contained in \(P\) and
\(Q\).  Let \(R\), \(R_1\), and \(R_2\) be the respective sets of roots
of \(G\), \(P\), and \(Q\) relative to \(T\).  By Exp.\ XXII, \SGARefLink{sga3:xxii:lem:5.4.4}{5.4.4} and
\SGARefLink{sga3:xxii:prop:5.4.5}{5.4.5},
\[
\begin{aligned}
 \dim G&=\dim T+\operatorname{Card}(R),&
 \dim P&=\dim T+\operatorname{Card}(R_1),\\
 \dim Q&=\dim T+\operatorname{Card}(R_2),&
 \dim(P\cap Q)&=\dim T+\operatorname{Card}(R_1\cap R_2).
\end{aligned}
\]
Condition \textup{\SGARefLink{sga3:xxvi:theorem:4-2-1:item:v}{(v)}} is therefore equivalent to
\[
 \operatorname{Card}(R_1\cap R_2)
 =\operatorname{Card}(R_1)+\operatorname{Card}(R_2)
  -\operatorname{Card}(R),
\]
that is, to \(R_1\cup R_2=R\).  By Exp.\ XXII, \SGARefLink{sga3:xxii:prop:5.9.2}{5.9.2} and \SGARefLink{sga3:xxii:prop:5.4.5}{5.4.5}, it is
enough to prove that \(R_1\cap(-R_2)\) contains a system of positive
roots of \(R\).  We are thus reduced to the following lemma.

\begin{sgastatement}{Lemma 4.2.2}
Let \(R\) be a ``root system''---for example, the set of roots of a root
datum in the sense of Exposé XXI.  Let \(R_1,R_2\) be two closed subsets
of \(R\), each containing a system of positive roots.  If
\(R_1\cup R_2=R\), then \(R_1\cap(-R_2)\) contains a system of positive
roots.
\end{sgastatement}
\SGARefTarget{sga3:xxvi:lemma:4:2:2}

Indeed, \(R_3=R_1\cap(-R_2)\) is closed, so by Exp.\ XXI, \SGARefLink{sga3:XXI:3.3.6}{3.3.6} it is
enough to show that \(R_3\cup(-R_3)=R\).  Now
\[
 R_1\cup(-R_1)=R=R_2\cup(-R_2),
\]
and the conclusion follows from the following elementary fact.  If
\(A,A',B,B'\) are four subsets of a set \(E\), and
\[
 A\cup A'=B\cup B'=A\cup B=A'\cup B'=E,
\]
then
\[
 (A\cap B')\cup(A'\cap B)=E.
\]

This proves \textup{\SGARefLink{sga3:xxvi:theorem:4-2-1:item:v}{(v)}} \(\Rightarrow\) \textup{\SGARefLink{sga3:xxvi:theorem:4-2-1:item:vi}{(vi)}} over the
spectrum of an algebraically closed field.  Return now to the general
case and suppose \textup{\SGARefLink{sga3:xxvi:theorem:4-2-1:item:v}{(v)}} holds.  Let \(s\in S\).  By the preceding
argument, there is a Borel subgroup \(\bar B\) of \(P_{\bar s}\) and a
Borel subgroup \(\bar B'\) of \(Q_{\bar s}\) such that
\(\bar B\cap\bar B'\) is a maximal torus of \(G_{\bar s}\).

The \(S\)-scheme
\[
 \Bor(P)\simeq\Bor\bigl(P/\rad^u(P)\bigr)
\]
of Borel subgroups of \(P\) is smooth.  Applying ``Hensel's lemma''
(Exp.\ XI, \SGARefLink{sga3:xi:proposition:1:10}{1.10}) and working étale-locally around \(s\), we may suppose
that there is a Borel subgroup \(B\) of \(P\) specializing to
\(\bar B\).  Similarly, there is a Borel subgroup \(B'\) of \(Q\)
specializing to \(\bar B'\).  Since
\(\bar B\cap\bar B'\) is a maximal torus of \(G_{\bar s}\), some open
neighborhood \(U\) of \(s\) has the property that
\(B_U\cap B'_U\) is a maximal torus of \(G_U\)
(Exp.\ XXII, \SGARefLink{sga3:xxii:cor:5.9.4}{5.9.4}).  This proves \textup{\SGARefLink{sga3:xxvi:theorem:4-2-1:item:vi}{(vi)}}.
\SGAQED

\begin{sgastatement}{Definition 4.2.3}
A pair \((P,Q)\) satisfying the equivalent conditions
\textup{\SGARefLink{sga3:xxvi:theorem:4-2-1:item:i}{(i)}}--\textup{\SGARefLink{sga3:xxvi:theorem:4-2-1:item:vii}{(vii)}} of \SGARefLink{sga3:xxvi:theorem:4-2-1}{Theorem 4.2.1} is said to be in
\emph{transversal position}.  One also says that \(P\) is in
transversal position relative to \(Q\), or, by an abuse of language,
that \(P\) and \(Q\) are in (mutual) transversal position.
\end{sgastatement}
\SGARefTarget{sga3:xxvi:def:4:2:3}

In view of \textup{\SGARefLink{sga3:xxvi:theorem:4-2-1:item:vi}{(vi)}}, this definition agrees, for Borel subgroups,
with that of Exp.\ XXII, \SGARefLink{sga3:xxii:def:5.9.1}{5.9.1}.

\begin{sgastatement}{\SGAInlineRefTarget{sga3:xxvi:corollary:4-2-4}Corollary 4.2.4}
Let \(S\) be a scheme, \(G\) an \(S\)-reductive group, and \(P,Q\) two
parabolic subgroups of \(G\).
\begin{enumerate}[label=\textup{(\roman*)}]
\item The pair \((P,Q)\) is in transversal position if and only if, for
every \(s\in S\), the pair \((P_{\bar s},Q_{\bar s})\) is in transversal
position.  If \(S'\to S\) is surjective and
\((P_{S'},Q_{S'})\) is in transversal position, then so is \((P,Q)\).
\item There is an open subscheme \(U\subset S\) with the following
property: a morphism \(S'\to S\) factors through \(U\) if and only if
\((P_{S'},Q_{S'})\) is in transversal position.
\item Consider the subfunctors
\[
 \Gen(G)\subset\Par(G)\times_S\Par(G),\qquad
 \Gen(/Q)\subset\Par(G),\qquad
 \Gen(P/Q)\subset G,
\]
defined as follows.  For \(S'\to S\), the set \(\Gen(G)(S')\) consists
of pairs of parabolic subgroups of \(G_{S'}\) in transversal position;
\(\Gen(/Q)(S')\) consists of parabolic subgroups of \(G_{S'}\) in
transversal position relative to \(Q_{S'}\); and \(\Gen(P/Q)(S')\)
consists of the \(g\in G(S')\) such that \(\Int(g)P_{S'}\) is in
transversal position relative to \(Q_{S'}\).

Each of these functors is represented by an open subscheme,
universally schematically dense over \(S\), of the corresponding
\(S\)-scheme
\[
 \Par(G)\times_S\Par(G),\qquad\Par(G),\qquad G,
\]
respectively.
\end{enumerate}
\end{sgastatement}
\SGARefTarget{sga3:xxvi:cor:4:2:4}
\footnote{\emph{Editor's note.} The original expression ``relatively
dense'' was replaced by ``universally schematically dense over \(S\)'';
compare EGA IV\textsubscript{3}, Definition 11.10.8.}

Assertion \textup{\SGARefLink{sga3:xxvi:theorem:4-2-1:item:i}{(i)}} follows immediately from condition
\SGARefLink{sga3:xxvi:theorem:4-2-1}{4.2.1}\textup{\SGARefLink{sga3:xxvi:theorem:4-2-1:item:i}{(i)}}.  For \textup{\SGARefLink{sga3:xxvi:theorem:4-2-1:item:ii}{(ii)}}, take
\[
 U=S-\operatorname{Supp}(\operatorname{Coker}u),
\]
where
\[
 u:\operatorname{Lie}(P)\oplus\operatorname{Lie}(Q)
   \longrightarrow\operatorname{Lie}(G)
\]
is the canonical morphism.\footnote{\emph{Editor's note.} The second
summand, printed as \(\operatorname{Lie}(G)\), has been corrected to
\(\operatorname{Lie}(Q)\); this is the map used in \SGARefLink{sga3:xxvi:theorem:4-2-1}{Theorem
4.2.1}\textup{\SGARefLink{sga3:xxvi:theorem:4-2-1:item:i}{(i)}}.}

There are cartesian diagrams
\[
\begin{tikzcd}[column sep=large,row sep=large]
 G \arrow[r,"f"] & \Par(G)\\
 \Gen(P/Q) \arrow[u,hook] \arrow[r] & \Gen(/Q)\arrow[u,hook]
\end{tikzcd}
\qquad
\begin{tikzcd}[column sep=large,row sep=large]
 \Par(G) \arrow[r,"f'"] & \Par(G)\times_S\Par(G)\\
 \Gen(/Q) \arrow[u,hook] \arrow[r] & \Gen(G)\arrow[u,hook],
\end{tikzcd}
\]
where \(f(g)=\Int(g)P\) and \(f'(R)=(R,Q)\).  It therefore suffices to
prove \textup{\SGARefLink{sga3:xxvi:theorem:4-2-1:item:iii}{(iii)}} for \(\Gen(G)\).

Let \(P_0\) be the canonical parabolic subgroup of
\(G_{\Par(G)}\), and set
\[
 X=\Par(G)\times_S\Par(G),\qquad
 P=\operatorname{pr}_1^*(P_0),\qquad
 Q=\operatorname{pr}_2^*(P_0).
\]
Applying \textup{\SGARefLink{sga3:xxvi:theorem:4-2-1:item:ii}{(ii)}} to the parabolic subgroups \(P,Q\) of \(G_X\)
constructs an open subscheme \(U\subset X\), which is immediately
identified with \(\Gen(G)\).  The density assertion may be checked on
geometric fibers,\footnote{\emph{Editor's note.} Compare EGA
IV\textsubscript{3}, 11.10.10.} so assume
\(S=\operatorname{Spec}(k)\), where \(k\) is algebraically closed.
Since \(\Par(G)\) is smooth, it is enough to show that \(\Gen(G)\)
meets every irreducible component of
\(\Par(G)\times_S\Par(G)\).  By \SGARefLink{sga3:xxvi:item:3:2}{3.3}, this amounts to the following:
given \(t,t'\in\Of(\underline{\operatorname{Dyn}}(G))(S)\), find a pair
\((P,P')\) in transversal position such that \(t(P)=t\) and
\(t(P')=t'\).

Choose Borel subgroups \(B,B'\) of \(G\) such that \(B\cap B'\) is a
maximal torus (split \(G\) and apply Exp.\ XXII, \SGARefLink{sga3:xxii:prop:5.9.2}{5.9.2}).  By \SGARefLink{sga3:xxvi:lemma:3:8}{3.8} there
are parabolic subgroups \(P\supset B\) and \(P'\supset B'\) with
\(t(P)=t\) and \(t(P')=t'\).  They have the required types and are in
transversal position by \SGARefLink{sga3:xxvi:theorem:4-2-1}{4.2.1}\textup{\SGARefLink{sga3:xxvi:theorem:4-2-1:item:vi}{(vi)}}.
\SGAQED

% END INCLUDED SGA3_English_current_buildable_source_20260823/source/inputs/XXVI/tex/components/03_section4_1_through_4_2_4.tex

% BEGIN INCLUDED SGA3_English_current_buildable_source_20260823/source/inputs/XXVI/tex/components/04_section4_2_5_through_4_2_7.tex
\begin{sgastatement}{\SGAInlineRefTarget{sga3:xxvi:corollary:4-2-5}Corollary 4.2.5}
Let \(S\) be a scheme, \(G\) an \(S\)-reductive group, and \(P,Q\) two
parabolic subgroups of \(G\), with \((P,Q)\) in transversal position.
\begin{enumerate}[label=\textup{(\roman*)}]
\item Let \(P',Q'\) be parabolic subgroups of \(G\) of the same types as
\(P,Q\), respectively.  The pair \((P',Q')\) is in transversal position
if and only if, étale-locally, it is conjugate to \((P,Q)\).
\textup{(In §5 we shall see that the étale topology can be replaced by
the Zariski topology.)}
\item The canonical morphism \(P\times_SQ\to G\) induces a smooth
surjective morphism
\[
 P\times_SQ\longrightarrow\Gen(Q/P)
\]
and an isomorphism
\[
 (P\times_SQ)/(P\cap Q)\xrightarrow{\sim}\Gen(Q/P),
\]
where \(R=P\cap Q\) acts on \(P\times_SQ\) by
\[
 (p,q)r=(pr,r^{-1}q).
\]
\item The canonical morphism
\[
 P\longrightarrow\Par_{t(Q)}(G),\qquad p\longmapsto\Int(p)Q,
\]
induces a smooth surjective morphism
\[
 P\longrightarrow\Gen(/P)\cap\Par_{t(Q)}(G)
\]
and an isomorphism
\[
 P/(P\cap Q)\xrightarrow{\sim}
 \Gen(/P)\cap\Par_{t(Q)}(G).
\]
\item The canonical morphism
\[
 G\longrightarrow\Par_{t(P)}(G)\times_S\Par_{t(Q)}(G),\qquad
 g\longmapsto\bigl(\Int(g)P,\Int(g)Q\bigr),
\]
induces a smooth surjective morphism
\[
 G\longrightarrow
 \Gen(G)\cap
 \bigl(\Par_{t(P)}(G)\times_S\Par_{t(Q)}(G)\bigr)
\]
and an isomorphism
\[
 G/(P\cap Q)\xrightarrow{\sim}
 \Gen(G)\cap
 \bigl(\Par_{t(P)}(G)\times_S\Par_{t(Q)}(G)\bigr).
\]
\end{enumerate}
\end{sgastatement}
\SGARefTarget{sga3:xxvi:cor:4:2:5}

\SGAProof
Let us prove \textup{(i)}.  The stated condition is plainly sufficient.
Conversely, suppose \((P',Q')\) is in transversal position.  Since
\(P\) and \(P'\) are conjugate étale-locally, we may assume \(P=P'\).
It is then enough to show that two parabolic subgroups \(Q,Q'\), of the
same type and both in transversal position relative to \(P\), are
étale-locally conjugate by a section of \(P\).

By \SGARefLink{sga3:xxvi:theorem:4-2-1}{4.2.1}\textup{\SGARefLink{sga3:xxvi:theorem:4-2-1:item:vi}{(vi)}}, after another étale localization there are Borel
subgroups \(B,B'\subset P\), \(B_1\subset Q\), and \(B'_1\subset Q'\)
such that
\[
 B\cap B_1=T,\qquad B'\cap B'_1=T'
\]
are maximal tori of \(G\).  The Killing pairs \((B,T)\) and
\((B',T')\) of \(P\) are étale-locally conjugate in \(P\) (\SGARefLink{sga3:xxvi:lemma:1:16}{1.16}).
We may therefore assume \(B=B'\) and \(T=T'\).  Then
\(B_1=B'_1\) by Exp.\ XXII, \SGARefLink{sga3:xxii:prop:5.9.2}{5.9.2}, and hence \(Q=Q'\) by \SGARefLink{sga3:xxvi:lemma:3:8}{3.8}.

Parts \textup{(ii)}, \textup{(iii)}, and \textup{(iv)} are proved in
parallel.  We prove \textup{(ii)}.  Let \(g\in\Gen(Q/P)(S)\); thus
\(g\in G(S)\) and \(\Int(g)Q\) is in transversal position relative to
\(P\).  The preceding argument shows that \(Q\) and \(\Int(g)Q\) are,
étale-locally, conjugate by a section of \(P\).  Working locally for
this topology, choose \(p\in P(S)\) such that
\(\Int(g)Q=\Int(p)Q\).  Then
\[
 p^{-1}g\in\Norm_G(Q)(S)=Q(S),
\]
so \(g=pq\) for some \(q\in Q(S)\).  Thus the morphism in
\textup{\SGARefLink{sga3:xxvi:theorem:4-2-1:item:ii}{(ii)}} is an étale covering.  Comparing with
\SGARefLink{sga3:xxvi:theorem:4-2-1}{4.2.1}\textup{\SGARefLink{sga3:xxvi:theorem:4-2-1:item:ii}{(ii)}}, it is smooth and surjective.

Finally, the equivalence relation on \(P\times_SQ\) defined by
\(P\times_SQ\to G\) is exactly the relation associated with the action
of \(R=P\cap Q\),
\[
 (p,q)r=(pr,r^{-1}q).
\]
This proves the quotient assertion, since a smooth surjective morphism
is, in particular, an effective epimorphism.
\SGAQED

\begin{sgastatement}{Remark 4.2.6}
If \(P,Q\) are in transversal position, the open subset \(\Gen(Q/P)\)
of \(G\) will often be denoted \(P\cdot Q\); this notation is justified
by \SGARefLink{sga3:xxvi:corollary:4-2-5}{4.2.5}\textup{(ii)}.
\end{sgastatement}
\SGARefTarget{sga3:xxvi:remark:4:2:6}

\begin{sgastatement}{\SGAInlineRefTarget{sga3:xxvi:proposition:4-2-7}Proposition 4.2.7}
Let \(S\) be a scheme, \(G\) an \(S\)-reductive group, and \(P,Q\) two
parabolic subgroups of \(G\), with \((P,Q)\) in transversal position.
\begin{enumerate}[label=\textup{(\roman*)}]
\item The group \(P\cap Q\) is smooth over \(S\) (and in fact has
connected fibers by \SGARefLink{sga3:xxvi:remark:4:1:2}{4.1.2}).  Its identity component
\((P\cap Q)^0\) \textup{(Exp.\ VI\textsubscript{B}, \SGARefLink{thm:sga3-VIB-3.10}{3.10})} is a
subgroup of type \textup{(RC)} of \(G\)
\textup{(Exp.\ XXII, \SGARefLink{sga3:xxii:def:5.11.1}{5.11.1})}, and its unipotent radical
\textup{(loc.\ cit., 5.11.4)} decomposes as the direct product
\[
 \rad^u\bigl((P\cap Q)^0\bigr)
 =
 \bigl(\rad^u(P)\cap Q\bigr)
 \times_S
 \bigl(P\cap\rad^u(Q)\bigr).
\]
\item If \(S\) is affine, then
\[
 H^1\bigl(S,\rad^u((P\cap Q)^0)\bigr)=0.
\]
If \(S\) is semilocal, \(P\cap Q\) contains a maximal torus of \(G\).
\end{enumerate}
\end{sgastatement}
\SGARefTarget{sga3:xxvi:prop:4:2:7}

\SGAProof
The group \(P\cap Q\) is smooth by \SGARefLink{sga3:xxvi:theorem:4-2-1}{4.2.1}\textup{\SGARefLink{sga3:xxvi:theorem:4-2-1:item:ii}{(ii)}} and the
cartesian diagram
\[
\begin{tikzcd}[column sep=large,row sep=large]
 P\times_SQ \arrow[r] & G\\
 P\cap Q \arrow[u] \arrow[r] & S\arrow[u]
\end{tikzcd}.
\]

The remaining assertions about \((P\cap Q)^0\) may be checked
étale-locally.  By \SGARefLink{sga3:xxvi:theorem:4-2-1}{4.2.1}\textup{\SGARefLink{sga3:xxvi:theorem:4-2-1:item:vii}{(vii)}}, choose a splitting
\((G,T,M,R)\) of \(G\) such that \(P,Q\) contain \(T\) and are defined
by subsets \(R_1,R_2\) of \(R\), where \(R_1\) contains a system
\(R^+\) of positive roots and \(R_2\) contains the opposite system
\(R^-\).  Let \(\Delta\) be the set of simple roots of \(R^+\), and put
\[
 A_1=\Delta\cap(-R_1),\qquad
 A_2=\Delta\cap R_2,\qquad
 A=A_1\cap A_2.
\]
By \SGARefLink{sga3:xxvi:prop:1:3}{1.4}\textup{(v)} and \SGARefLink{sga3:xxvi:proposition:1-12}{1.12},
\[
\begin{aligned}
 R_1&=R^+\cup\bigl(R^-\cap(-\mathbb N A_1)\bigr),&
 R_2&=\bigl(R^+\cap\mathbb N A_2\bigr)\cup R^-,\\
 \rad^u(P)&=
   \prod_{\substack{\alpha\in R_1\\\alpha\notin-R_1}}U_\alpha,&
 \rad^u(Q)&=
   \prod_{\substack{\alpha\in R_2\\\alpha\notin-R_2}}U_\alpha.
\end{aligned}
\]
Exp.\ XXII, \SGARefLink{sga3:xxii:cor:5.6.7}{5.6.7} therefore gives
\[
\begin{aligned}
 \rad^u(P)\cap Q
  &=
  \prod_{\substack{\alpha\in R_1\cap R_2\\\alpha\notin-R_1}}U_\alpha
  =\prod_{\alpha\in K_2}U_\alpha,\\
 \rad^u(Q)\cap P
  &=
  \prod_{\substack{\alpha\in R_1\cap R_2\\\alpha\notin-R_2}}U_\alpha
  =\prod_{\alpha\in K_1}U_\alpha,
\end{aligned}
\]
where \(K_2\) consists of the positive roots that are linear
combinations of elements of \(A_2\), but not of elements of \(A\), and
\(K_1\) consists of the negative roots that are linear combinations of
elements of \(A_1\), but not of elements of \(A\).  If
\(\alpha\in K_2\) and \(\beta\in K_1\), then \(\alpha+\beta\) is neither
a root nor zero.  Hence the two groups displayed above commute.

By Exp.\ XXII, \SGARefLink{sga3:xxii:prop:5.4.5}{5.4.5}, the group
\[
 H=(P\cap Q)^0
\]
is defined by the set of roots
\[
 R_1\cap R_2
 =
 \bigl(R^+\cap\mathbb N A_2\bigr)
 \cup
 \bigl(R^-\cap(-\mathbb N A_1)\bigr).
\]
This set is closed, so \(H\) is of type \textup{(RC)} by definition
(Exp.\ XXII, \SGARefLink{sga3:xxii:def:5.11.1}{5.11.1}).  By loc.\ cit., \SGARefLink{sga3:xxii:lem:5.11.3}{5.11.3} and \SGARefLink{sga3:xxii:prop:5.11.4}{5.11.4},
\[
 \rad^u(H)=\prod_{\alpha\in K}U_\alpha,
\]
where \(K\) consists of the \(\alpha\in R_1\cap R_2\) such that
\(\alpha\notin-(R_1\cap R_2)\).  The symmetric part of
\(R_1\cap R_2\) is \(R\cap\mathbb ZA\), and therefore
\(K=K_1\cup K_2\).  This proves \textup{(i)}.

The first assertion of \textup{(ii)} follows from \textup{(i)} and
\SGARefLink{sga3:xxvi:cor:2:10}{2.10}.  For the second, the reductive group
\[
 (P\cap Q)^0/\rad^u((P\cap Q)^0)
\]
has a maximal torus \(\bar T\) when the base is semilocal.  The inverse
image \(N\) of \(\bar T\) in \((P\cap Q)^0\) is a subgroup of type
\textup{(R)} of \(G\), with solvable fibers, and
\[
 N^u=\rad^u((P\cap Q)^0)
\]
(Exp.\ XXII, \SGARefLink{sga3:xxii:prop:5.6.9}{5.6.9}).  The scheme of maximal tori of \(N\) is a
principal homogeneous bundle under \(N^u\) (loc.\ cit., \SGARefLink{sga3:xxii:cor:5.6.13}{5.6.13}), and
therefore has a section because \(H^1(S,N^u)=0\).
\SGAQED

% END INCLUDED SGA3_English_current_buildable_source_20260823/source/inputs/XXVI/tex/components/04_section4_2_5_through_4_2_7.tex

% BEGIN INCLUDED SGA3_English_current_buildable_source_20260823/source/inputs/XXVI/tex/components/05_section4_3.tex
\SGASubsection{4.3}{Opposite parabolic subgroups}
\SGARefTarget{sga3:xxvi:subsection:4:3}

\SGAInlineRefTarget{sga3:xxvi:paragraph:4-3-1}\paragraph{\SGARefLink{sga3:xxvi:paragraph:4-3-1}{4.3.1}.}
Let \(G\) be an \(S\)-reductive group.  Exp.\ XXIV, \SGARefLink{sga3:xxiv:proposition:3-16-6}{3.16.6} defines a
canonical ``outer automorphism'' \(s_G\) of \(G\), of order at most
\(2\).\footnote{\emph{Editor's note.} ``Of order \(2\)'' was replaced
by ``of order at most \(2\)'', since \(s_G\) may be trivial, for example
when \(G\) has type \(A_1,B_n,C_n,\ldots\).}  It induces a canonical
automorphism, again of order at most \(2\), of
\(\underline{\operatorname{Dyn}}(G)\), and hence of
\(\Of(\underline{\operatorname{Dyn}}(G))\).  We denote all of these
automorphisms by \(s_G\), or simply by \(s\).  Two types of parabolic
subgroups
\[
 t,t'\in\Of(\underline{\operatorname{Dyn}}(G))(S)
\]
are said to be \emph{opposite} if \(t=s_G(t')\).
\SGARefTarget{sga3:xxvi:item:4:3:1}

\begin{sgastatement}{Theorem 4.3.2}
Let \(S\) be a scheme, \(G\) an \(S\)-reductive group, and \(P\) a
parabolic subgroup of \(G\).

\textup{(a)} If \(L\) is a Levi subgroup of \(P\), there is a unique
parabolic subgroup \(P'\) of \(G\) such that \(P\cap P'=L\).

\textup{(b)} For every parabolic subgroup \(Q\) of \(G\), the following
conditions are equivalent:
\begin{enumerate}[leftmargin=3.2em,labelsep=.7em]
\item[\textup{\SGAInlineRefTarget{sga3:xxvi:item:4:3:1:item:i}(i)}]
For every \(s\in S\), the identity component
\(\bigl((P\cap Q)_{\bar s}\bigr)^0\), which is smooth by \SGARefLink{sga3:xxvi:lemma:4:1:1}{4.1.1}, is
reductive.
\item[\textup{\SGAInlineRefTarget{sga3:xxvi:item:4:3:1:item:ii}(ii)}]
The group \(P\cap Q\) is a Levi subgroup of both \(P\) and \(Q\).
\item[\textup{\SGAInlineRefTarget{sga3:xxvi:item:4:3:1:item:iii}(iii)}]
The groups \(P,Q\) have opposite types and the pair \((P,Q)\) is in
transversal position.
\item[\textup{\SGAInlineRefTarget{sga3:xxvi:item:4:3:1:item:iv}(iv)}]
The groups \(P,Q\) have opposite types and
\(\rad^u(P)\cap Q=e\).
\item[\textup{\SGAInlineRefTarget{sga3:xxvi:item:4:3:1:item:v}(v)}]
\[
 \rad^u(P)\cap Q=\rad^u(Q)\cap P=e.
\]
\item[\textup{\SGAInlineRefTarget{sga3:xxvi:item:4:3:1:item:vi}(vi)}]
The canonical morphism \(\rad^u(P)\times_SQ\to G\) is an open immersion.
\item[\textup{(vi\(^{\prime}\))}]
The canonical morphism \(\rad^u(P)\to G/Q\) is an open immersion.
\item[\textup{\SGAInlineRefTarget{sga3:xxvi:item:4:3:1:item:vii}(vii)}]
There is an étale covering family \(\{S_i\to S\}\) and, for every \(i\),
a splitting \((T_i,M_i,R_i)\) of \(G_{S_i}\) and a subset
\(R_i^{(1)}\subset R_i\), such that \(P_{S_i}\) is the subgroup of type
\textup{(R)} containing \(T_i\) and defined by \(R_i^{(1)}\), while
\(Q_{S_i}\) is the subgroup defined by
\[
 R_i^{(2)}=-R_i^{(1)}.
\]
\end{enumerate}
\end{sgastatement}
\SGARefTarget{sga3:xxvi:theorem:4:3:2}

\SGAProof
We first prove part \textup{(b)}.  Conditions \textup{\SGARefLink{sga3:xxvi:theorem:4-2-1:item:iii}{(iii)}} and
\textup{\SGARefLink{sga3:xxvi:theorem:4-2-1:item:vii}{(vii)}} are equivalent by \SGARefLink{sga3:xxvi:theorem:4-2-1}{4.2.1}\textup{\SGARefLink{sga3:xxvi:theorem:4-2-1:item:vii}{(vii)}} and the
definition of \(s_G\) in the split case (Exp.\ XXIV,
\SGARefLink{sga3:xxiv:proposition:3-16-2}{3.16.2}\textup{\SGARefLink{sga3:xxiv:proposition:3-16-2:item:iv}{(iv)}}).  Plainly \textup{\SGARefLink{sga3:xxiv:proposition:3-16-2:item:ii}{(ii)}} implies \textup{\SGARefLink{sga3:xxiv:proposition:3-16-2:item:i}{(i)}}, and
\textup{(vi\(^{\prime}\))} implies \textup{\SGARefLink{sga3:xxvi:theorem:4-2-1:item:vi}{(vi)}} by base change along
\(G\to G/Q\).

Suppose now that \textup{\SGARefLink{sga3:xxvi:item:4:3:1:item:vii}{(vii)}} holds.  Since every assertion to be
proved is étale-local, assume that \(G=(G,T,M,R)\) is split, that \(P\)
is defined by \(R'\subset R\), and that \(Q\) is defined by \(-R'\).
Let \(L\) be the subgroup of type \textup{(R)} containing \(T\) and
defined by \(R'\cap(-R')\).  By Exp.\ XXII, \SGARefLink{sga3:xxii:lem:5.11.3}{5.11.3}, \(L\) is a Levi
subgroup of both \(P\) and \(Q\).  Moreover,
\[
 P=L\cdot\rad^u(P),\qquad Q=L\cdot\rad^u(Q),
\]
and Exp.\ XXII, \SGARefLink{sga3:xxii:cor:5.6.7}{5.6.7} gives
\[
 \rad^u(P)\cap Q=Q\cap\rad^u(P)=e.
\]
Thus \(P\cap Q=L\), proving \textup{\SGARefLink{sga3:xxii:cor:5.6.7:item:ii}{(ii)}} and \textup{\SGARefLink{sga3:xxvi:item:4:3:1:item:v}{(v)}}.

Since \(P,Q\) are in transversal position, the canonical morphism
\(P\to G/Q\) induces an open immersion
\[
 P/(P\cap Q)\longrightarrow G/Q
\]
by \SGARefLink{sga3:xxvi:theorem:4-2-1}{4.2.1}.  The canonical morphism
\[
 \rad^u(P)\longrightarrow P/(P\cap Q)=P/L
\]
is an isomorphism, so \textup{(vi\(^{\prime}\))} follows.  Hence all
the conditions follow from \textup{\SGARefLink{sga3:xxvi:theorem:4-2-1:item:vii}{(vii)}}.

It remains to show that each of \textup{\SGARefLink{sga3:xxvi:item:4:3:1:item:i}{(i)}}, \textup{\SGARefLink{sga3:xxvi:item:4:3:1:item:iv}{(iv)}},
\textup{\SGARefLink{sga3:xxvi:item:4:3:1:item:v}{(v)}}, and \textup{\SGARefLink{sga3:xxvi:item:4:3:1:item:vi}{(vi)}} implies \textup{\SGARefLink{sga3:xxvi:item:4:3:1:item:vii}{(vii)}}.  Since the
equivalence of \textup{\SGARefLink{sga3:xxvi:item:4:3:1:item:ii}{(ii)}} and \textup{\SGARefLink{sga3:xxvi:item:4:3:1:item:iii}{(iii)}} is already known, this
may be checked on geometric fibers.  We may suppose that \(S\) is the
spectrum of an algebraically closed field.  By \SGARefLink{sga3:xxvi:lemma:4:1:1}{4.1.1}, a maximal torus
\(T\) of \(G\) is contained in \(P\cap Q\).  Let \(R,R_1,R_2\) be the
sets of roots of \(G,P,Q\), respectively, relative to \(T\).  Write
\[
 R_1^a=\{\alpha\in R_1\mid-\alpha\notin R_1\},
\]
and define \(R_2^a\) similarly.  We must prove \(R_1=-R_2\).

\begin{itemize}[leftmargin=2em]
\item Condition \textup{\SGARefLink{sga3:xxvi:item:4:3:1:item:i}{(i)}} says that \(R_1\cap R_2\) is symmetric.
If \(\alpha\in R_1\) and \(\alpha\notin R_2\), then
\(\alpha\in-R_2\); if \(\alpha\in R_2\), then
\[
 \alpha\in R_1\cap R_2=-(R_1\cap R_2)\subset-R_2.
\]
Thus \(R_1\subset-R_2\), and symmetry gives \(R_1=-R_2\).

\item Condition \textup{\SGARefLink{sga3:xxvi:item:4:3:1:item:iv}{(iv)}} gives
\(\operatorname{Card}(R_1)=\operatorname{Card}(R_2)\) and
\(R_1^a\cap R_2=\varnothing\).  The latter condition is equivalent to
\(R_2\subset-R_1\), and equality of cardinalities then gives
\(R_2=-R_1\).

\item Condition \textup{\SGARefLink{sga3:xxvi:item:4:3:1:item:v}{(v)}} gives
\[
 R_1^a\cap R_2=R_2^a\cap R_1=\varnothing.
\]
Hence \(R_2\subset-R_1\) and \(R_1\subset-R_2\), so again
\(R_2=-R_1\).

\item Condition \textup{\SGARefLink{sga3:xxvi:item:4:3:1:item:vi}{(vi)}} gives
\[
 \operatorname{Lie}(\rad^u(P))\oplus\operatorname{Lie}(Q)
 =\operatorname{Lie}(G).
\]
Thus \(R\) is the disjoint union of \(R_2\) and \(R_1^a\), which implies
\(R_2=-R_1\).
\end{itemize}

This proves part \textup{(b)}.  For part \textup{(a)}, the implication
\textup{\SGARefLink{sga3:xxvi:item:4:3:1:item:vii}{(vii)}} \(\Rightarrow\) \textup{\SGARefLink{sga3:xxvi:item:4:3:1:item:ii}{(ii)}} already gives the desired
parabolic subgroup \(P'\) étale-locally.  It remains only to prove
uniqueness, and this too is étale-local.  We may suppose that \(G\) is
split relative to a maximal torus \(T\) of \(L\), with \(P\) and \(P'\)
defined by subsets \(R_1\) and \(R'_1\) of \(R\).  By hypothesis
\(R_1\cap R'_1\) is symmetric.  The preceding argument gives
\(R'_1=-R_1\); hence \(P'\) is determined by \(P\) and \(L\).
\SGAQED

\begin{sgastatement}{\SGAInlineRefTarget{sga3:xxvi:definition:4-3-3}Definition 4.3.3}
Two parabolic subgroups of \(G\) satisfying the equivalent conditions
\textup{(i)}--\textup{\SGARefLink{sga3:xxvi:item:4:3:1:item:vii}{(vii)}} of \SGARefLink{sga3:xxvi:item:4:3:1}{Theorem 4.3.2} are said to be
\emph{opposite}.  If \(P\) is a parabolic subgroup of \(G\) and \(L\)
is a Levi subgroup of \(P\), the \emph{parabolic subgroup opposite to
\(P\) relative to \(L\)} is the unique parabolic subgroup \(Q\) of
\(G\) such that \(P\cap Q=L\).

If \(T\) is a maximal torus of \(P\), the \emph{parabolic subgroup
opposite to \(P\) relative to \(T\)} is the unique \(Q\) such that
\(P\cap Q\) is the unique Levi subgroup of \(P\) containing \(T\)
\textup{(\SGARefLink{sga3:xxvi:proposition:1-6}{1.6})}; equivalently, it is the unique \(Q\) such that
\(P\cap Q\) contains \(T\) and \(P,Q\) are opposite.
\end{sgastatement}
\SGARefTarget{sga3:xxvi:def:4:3:3}

By \SGARefLink{sga3:xxvi:item:4:3:1}{4.3.2}\textup{\SGARefLink{sga3:xxvi:item:4:3:1:item:iii}{(iii)}}, \SGARefLink{sga3:xxvi:corollary:4-2-4}{Corollaries 4.2.4} and \SGARefLink{sga3:xxvi:corollary:4-2-5}{4.2.5} immediately yield
parallel results.  We record one sample.

\begin{sgastatement}{\SGAInlineRefTarget{sga3:xxvi:corollary:4-3-4}Corollary 4.3.4}
Let \(S\) be a scheme, \(G\) an \(S\)-reductive group, and \(P,Q\) two
parabolic subgroups of \(G\).
\begin{enumerate}[label=\textup{(\roman*)}]
\item The groups \(P,Q\) are opposite if and only if
\(P_{\bar s},Q_{\bar s}\) are opposite for every \(s\in S\).  If
\(S'\to S\) is surjective and \(P_{S'},Q_{S'}\) are opposite, then
\(P,Q\) are opposite.
\item Let \(\Opp(G)(S')\) be the set of pairs of opposite parabolic
subgroups of \(G_{S'}\), for \(S'\to S\).  This functor is represented
by an open subscheme of \(\Par(G)^2\).  Let \(\Opp(/P)(S')\) be the set
of parabolic subgroups of \(G_{S'}\) opposite to \(P_{S'}\).  This
functor is represented by an open subscheme, universally schematically
dense over \(S\), of
\[
 \Par_{s(t(P))}(G).
\]
\item Suppose \(P,Q\) are opposite, and let \(P',Q'\) be parabolic
subgroups of \(G\), with \(P'\) of the same type as \(P\).  The groups
\(P',Q'\) are opposite if and only if, étale-locally, the pair
\((P',Q')\) is conjugate to \((P,Q)\).
\textup{(In §5 the étale topology will be replaced by the Zariski
topology.)}
\end{enumerate}
\end{sgastatement}
\SGARefTarget{sga3:xxvi:cor:4:3:4}
\footnote{\emph{Editor's note.} Here again, ``relatively dense'' was
replaced by ``universally schematically dense over \(S\)''; compare
EGA IV\textsubscript{3}, Definition 11.10.8.}

\begin{sgastatement}{\SGAInlineRefTarget{sga3:xxvi:corollary:4-3-5}Corollary 4.3.5}
Let \(S\) be a scheme, \(G\) an \(S\)-reductive group, and \(P\) a
parabolic subgroup of \(G\).
\begin{enumerate}[label=\textup{(\roman*)}]
\item The morphism
\[
 \Opp(/P)\longrightarrow\Lev(P),\qquad Q\longmapsto P\cap Q,
\]
is an isomorphism.  The scheme \(\Opp(/P)\) is a principal homogeneous
bundle under \(\rad^u(P)\), acting by inner automorphisms.  If \(S\) is
affine, \(G\) has a parabolic subgroup opposite to \(P\).
\item Suppose \(S\) is semilocal, with closed points \(\{s_i\}\).  For
each \(i\), let \(Q_i\) be a parabolic subgroup of \(G_{s_i}\) opposite
to \(P_{s_i}\).  There is a parabolic subgroup \(Q\) of \(G\), opposite
to \(P\), such that \(Q_{s_i}=Q_i\) for every \(i\).
\item The morphism
\[
 \Opp(G)\longrightarrow\mathrm{PL},\qquad
 (P,Q)\longmapsto(P,P\cap Q),
\]
is an isomorphism.
\end{enumerate}
\end{sgastatement}
\SGARefTarget{sga3:xxvi:cor:4:3:5}

These assertions follow from \SGARefLink{sga3:xxvi:item:4:3:1}{Theorem 4.3.2}\textup{(a)} and from \SGARefLink{sga3:xxvi:cor:1:9}{1.9},
\SGARefLink{sga3:xxvi:cor:2:3}{2.3}, and \SGARefLink{sga3:xxvi:cor:2:8}{2.8}.

\begin{sgastatement}{Remark 4.3.6}
Let \(P,Q\) be opposite parabolic subgroups of \(G\), and let
\(P\cdot Q\) be the open subscheme of \(G\), introduced in \SGARefLink{sga3:xxvi:remark:4:2:6}{4.2.6}, which
is the sheaf-theoretic image of \(P\times_SQ\).  The product morphism
\(G\times_SG\to G\) induces isomorphisms
\[
 \rad^u(P)\times_SQ
 \xrightarrow{\sim}P\cdot Q
 \xleftarrow{\sim}P\times_S\rad^u(Q).
\]
\end{sgastatement}
\SGARefTarget{sga3:xxvi:remark:4:3:6}

Indeed, this follows from \SGARefLink{sga3:xxvi:item:4:3:1}{4.3.2} (or \SGARefLink{sga3:xxvi:corollary:4-2-5}{4.2.5}\textup{\SGARefLink{sga3:xxvi:item:4:3:1:item:ii}{(ii)}}) and from the fact
that \(P\cap Q\) is a Levi subgroup of both \(P\) and \(Q\):
\[
 P=\rad^u(P)\cdot(P\cap Q),\qquad
 Q=\rad^u(Q)\cdot(P\cap Q).
\]
Likewise, there is a commutative diagram
\[
\begin{tikzcd}[column sep=large,row sep=large]
 \rad^u(P) \arrow[r,hook] \arrow[d,"\sim"] & G/Q\arrow[d,"\sim"]\\
 \Opp(/P) \arrow[r,hook] & \Par_{t(Q)}(G),
\end{tikzcd}
\]
in which the vertical arrows are induced by
\(g\mapsto\Int(g)Q\).

% END INCLUDED SGA3_English_current_buildable_source_20260823/source/inputs/XXVI/tex/components/05_section4_3.tex

% BEGIN INCLUDED SGA3_English_current_buildable_source_20260823/source/inputs/XXVI/tex/components/06_section4_4.tex
\SGASubsection{4.4}{Osculating position}
\SGARefTarget{sga3:xxvi:subsection:4:4}

\begin{sgastatement}{\SGAInlineRefTarget{sga3:xxvi:proposition:4-4-1}Proposition 4.4.1}
Let \(S\) be a scheme, \(G\) an \(S\)-reductive group, and \(P,Q\)
two parabolic subgroups of \(G\).  The following conditions are
equivalent:
\begin{enumerate}[label=\textup{(\roman*)}]
\item \(P\cap Q\) is a parabolic subgroup of \(G\).
\item Étale-locally, \(P\cap Q\) contains a Borel subgroup of \(G\).
\item Étale-locally, \(P\cap Q\) contains a maximal torus of \(G\);
moreover, for every \(S'\to S\) and every maximal torus \(T\) of
\(G_{S'}\) contained in both \(P_{S'}\) and \(Q_{S'}\), the parabolic
subgroup opposite to \(P_{S'}\) relative to \(T\) is in transverse
position relative to \(Q_{S'}\).
\item There is an étale covering family \(\{S_i\to S\}\) and, for
each \(i\), a maximal torus \(T_i\) of \(G_{S_i}\), contained in both
\(P_{S_i}\) and \(Q_{S_i}\), such that the parabolic subgroup opposite
to \(P_{S_i}\) relative to \(T_i\) is in transverse position relative
to \(Q_{S_i}\).
\item There is an étale covering family \(\{S_i\to S\}\) and, for
each \(i\), a splitting \((T_i,M_i,R_i)\) of \(G_{S_i}\) such that
\(P_{S_i}\), respectively \(Q_{S_i}\), is the subgroup of type
\textup{(R)} of \(G_{S_i}\) containing \(T_i\) and defined by a set
of roots \(R_i^{(1)}\), respectively \(R_i^{(2)}\), with
\(R_i^{(1)}\cap R_i^{(2)}\) containing a system of positive roots
of \(R_i\).
\end{enumerate}
Moreover, when these conditions hold,
\[
 t(P\cap Q)=t(P)\cap t(Q)
\]
in the notation of \SGARefLink{sga3:xxvi:item:3:1}{3.2}.
\end{sgastatement}
\SGARefTarget{sga3:xxvi:prop:4:4:1}

The implications \textup{(v)}\(\Rightarrow\)\textup{(ii)} and
\textup{(iii)}\(\Rightarrow\)\textup{(iv)} are immediate.  Moreover,
\textup{(ii)}\(\Rightarrow\)\textup{(i)} follows from \SGARefLink{sga3:xxvi:rem:1:18}{1.18}.  To prove
\textup{(iv)}\(\Rightarrow\)\textup{(v)}, one may suppose that \(G\)
is split and that \(P\), respectively \(Q\), is defined by a set of
roots \(R_1\), respectively \(R_2\).  The parabolic subgroup opposite
to \(P\) is then defined by \(-R_1\), and the assertion reduces to
\SGARefLink{sga3:xxvi:lemma:4:2:2}{Lemma 4.2.2}.  The implication
\textup{(i)}\(\Rightarrow\)\textup{(iii)} is proved in the same way
after splitting.  Finally, the last assertion may be proved
étale-locally: one may suppose that \(P\cap Q\) contains a Borel
subgroup \(B\) of \(G\), and the assertion then reduces to \SGARefLink{sga3:xxvi:remark:3-7}{3.7}.

\begin{sgastatement}{Definition 4.4.2}
Two parabolic subgroups of \(G\) satisfying the equivalent conditions
\textup{(i)}--\textup{(v)} of \SGARefLink{sga3:xxvi:proposition:4-4-1}{4.4.1} are said to be in
\emph{osculating position}.
\end{sgastatement}
\SGARefTarget{sga3:xxvi:def:4:4:2}

\begin{sgastatement}{Corollary 4.4.3}
Let \(P,Q\) be parabolic subgroups in osculating position, and let
\(P',Q'\) be parabolic subgroups of \(G\) of the same types as \(P,Q\),
respectively.  The subgroups \(P',Q'\) are in osculating position if
and only if, étale-locally, the pair \((P',Q')\) is conjugate to the
pair \((P,Q)\).
\end{sgastatement}
\SGARefTarget{sga3:xxvi:cor:4:4:3}

It is enough to show that, if both \(P\) and \(P'\) are in osculating
position relative to \(Q\), then they are conjugate, étale-locally,
by a section of \(Q\).  Now \(P\cap Q\) and \(P'\cap Q\) are parabolic
subgroups of the same type contained in \(Q\).  By part \textup{(ii)}
of the lemma below, they are therefore conjugate étale-locally by a
section of \(Q\).  We may consequently suppose
\(P\cap Q=P'\cap Q\); part \textup{(i)} of that lemma then gives
\(P=P'\).

\begin{sgastatement}{Lemma 4.4.4}
Let \(P,P',Q\) be three parabolic subgroups of the \(S\)-reductive
group \(G\).
\begin{enumerate}[label=\textup{(\roman*)}]
\item The equality \(P=P'\) holds if and only if \(P\) and \(P'\) are
in osculating position and have the same type.
\item Suppose \(P\subset Q\) and \(P'\subset Q\).  If \(g\in G(S)\)
is such that \(\Int(g)P\) and \(P'\) are in osculating position, then
\(g\in Q(S)\).
\end{enumerate}
\end{sgastatement}
\SGARefTarget{sga3:xxvi:lemma:4:4:4}

Part \textup{(i)} follows immediately from the last assertion of
\SGARefLink{sga3:xxvi:proposition:4-4-1}{4.4.1}.  To prove \textup{(ii)}, observe that \(Q\) and \(\Int(g)Q\)
both contain \(P'\cap\Int(g)P\), and hence are in osculating position.
They coincide by \textup{(i)}, so that
\[
 g\in\Norm_G(Q)(S)=Q(S).
\]

Assertions \textup{(iii)} and \textup{(iv)} of \SGARefLink{sga3:xxvi:proposition:4-4-1}{Proposition 4.4.1} give
at once:

\begin{sgastatement}{Corollary 4.4.5}
Let \(P,P',Q\) be three parabolic subgroups of the \(S\)-reductive
group \(G\), containing the same maximal torus \(T\) of \(G\).
Suppose that \(P\) and \(P'\) are opposite relative to \(T\).  Then
\(Q\) is in osculating position relative to \(P\) if and only if it is
in transverse position relative to \(P'\).  Under these conditions,
\(P\cap Q\) is also in transverse position relative to \(P'\).
\end{sgastatement}
\SGARefTarget{sga3:xxvi:cor:4:4:5}

\begin{sgastatement}{Corollary 4.4.6}
Let \(P,Q\) be two parabolic subgroups of \(G\) containing the same
maximal torus \(T\).  The subgroups \(P,Q\) are in transverse position
if and only if there are parabolic subgroups \(P'\subset P\) and
\(Q'\subset Q\) of \(G\), opposite relative to \(T\).  One may even
choose
\[
 t(P')=t(P)\cap s(t(Q)).
\]
\end{sgastatement}
\SGARefTarget{sga3:xxvi:cor:4:4:6}
\footnote{\emph{Editor's note.} The involution \(s\) was defined in
\SGARefLink{sga3:xxvi:paragraph:4-3-1}{4.3.1}.}

The condition is plainly sufficient, by \SGARefLink{sga3:xxvi:theorem:4-2-1}{4.2.1}\textup{\SGARefLink{sga3:xxvi:theorem:4-2-1:item:i}{(i)}} and
\SGARefLink{sga3:xxvi:item:4:3:1}{4.3.2}\textup{\SGARefLink{sga3:xxvi:item:4:3:1:item:iii}{(iii)}}.  Conversely, let \(P^-\), respectively \(Q^-\),
be the parabolic subgroup opposite to \(P\), respectively \(Q\),
relative to \(T\).  By \SGARefLink{sga3:xxvi:cor:4:4:5}{4.4.5}, \(P^-\cap Q\) is in transverse position
relative to both \(P\) and \(Q^-\), and therefore also relative to
\(P\cap Q^-\), by another application of \SGARefLink{sga3:xxvi:cor:4:4:5}{4.4.5}.  Moreover,
\[
\begin{aligned}
 t(P^-\cap Q)
 &=t(P^-)\cap t(Q)\\
 &=s(t(P))\cap s(t(Q^-))\\
 &=s\bigl(t(P)\cap t(Q^-)\bigr)\\
 &=s\bigl(t(P\cap Q^-)\bigr).
\end{aligned}
\]
It follows that
\[
 P^-\cap Q=P',\qquad P\cap Q^-=Q'
\]
are opposite, by \SGARefLink{sga3:xxvi:item:4:3:1}{4.3.2}\textup{\SGARefLink{sga3:xxvi:item:4:3:1:item:iii}{(iii)}}.  Since \(P'\cap Q'\supset T\),
they are indeed opposite relative to \(T\).

% END INCLUDED SGA3_English_current_buildable_source_20260823/source/inputs/XXVI/tex/components/06_section4_4.tex

% BEGIN INCLUDED SGA3_English_current_buildable_source_20260823/source/inputs/XXVI/tex/components/07_section4_5.tex
\SGASubsection{4.5}{Standard position}
\SGARefTarget{sga3:xxvi:subsection:4:5}

In this subsection we briefly indicate how some of the preceding
results generalize.

\begin{sgastatement}{\SGAInlineRefTarget{sga3:xxvi:proposition:4-5-1}Proposition 4.5.1}
Let \(P_1,P_2\) be two parabolic subgroups of the \(S\)-reductive
group \(G\).  The following conditions are equivalent:
\begin{enumerate}[label=\textup{(\roman*)}]
\item \(P_1\cap P_2\) is smooth.
\item \(P_1\cap P_2\) is a subgroup of type \textup{(R)}, or of type
\textup{(RC)}, of \(G\).
\item Fpqc-locally, \(P_1\cap P_2\) contains a maximal torus of \(G\).
\item Zariski-locally, \(P_1\cap P_2\) contains a maximal torus of
\(G\).
\end{enumerate}
When \(S\) is semilocal, these conditions are also equivalent to:
\begin{enumerate}[label=\textup{(\roman*)},start=5]
\item \(P_1\cap P_2\) contains a maximal torus of \(G\).
\end{enumerate}
\end{sgastatement}
\SGARefTarget{sga3:xxvi:prop:4:5:1}

\SGAProof
Clearly, \textup{(iv)} implies \textup{(iii)}, and \textup{(v)}
implies \textup{(iv)} when \(S\) is semilocal; condition \textup{(ii)}
implies \textup{(i)} by Exp.\ XXII, \SGARefLink{sga3:xxii:def:5.2.1}{Definition 5.2.1}.  We shall prove
\[
 \textup{(i)}\Rightarrow\textup{(ii)}
 \Rightarrow\textup{(iii)},
\]
and then show that \textup{(iii)} implies \textup{\SGARefLink{sga3:xxii:def:5.2.1:item:i}{(i)}} and
\textup{(iv)}, as well as \textup{(v)} when \(S\) is semilocal.
Put \(K=P_1\cap P_2\).  By \SGARefLink{sga3:xxvi:lemma:4:1:1}{4.1.1} and \textup{[BT65]}, \SGARefLink{sga3:xxvi:subsection:4:5}{4.5}, every
geometric fiber \(K_{\bar s}\) contains a maximal torus
\(T_{\bar s}\) of \(G_{\bar s}\), is connected, and has, relative to
\(T_{\bar s}\), a closed subset of the roots of \(G_{\bar s}\).
Consequently, if \(K\) is smooth over \(S\), then it is a subgroup of
type \textup{(RC)}, hence a fortiori of type \textup{(R)}; see Exp.\
XXII, \SGARefLink{sga3:xxii:def:5.2.1}{5.2.1} and 5.11.1.  Thus \textup{\SGARefLink{sga3:xxii:def:5.2.1:item:i}{(i)}} and \textup{\SGARefLink{sga3:xxii:def:5.2.1:item:ii}{(ii)}} are
equivalent.

If \(K\) is of type \textup{(R)}, it contains a maximal torus of \(G\)
étale-locally, by Exp.\ XXII, \SGARefLink{sga3:xxii:prop:2.2}{2.2} and Exp.\ XIX, \SGARefLink{sga3:xxvi:proposition:6-1}{6.1}.  Hence
\textup{\SGARefLink{sga3:xxii:def:5.2.1:item:ii}{(ii)}} implies \textup{(iii)}.

Assume \textup{(iii)}, and let us prove that \(K\) is smooth.  By
fpqc descent we may suppose that
\[
 G=(G,T,M,R)
\]
is split, where \(T\) is a maximal torus contained in \(P_1\cap P_2\),
and that there are closed subsets \(R_1,R_2\) of \(R\) such that
\[
 P_1=H_{R_1},\qquad P_2=H_{R_2}.
\]
Since \(K\) has connected fibers, Exp.\ XXII, \SGARefLink{sga3:xxii:prop:5.4.5}{5.4.5} gives
\[
 K=H_{R_1\cap R_2};
\]
therefore \(K\) is smooth over \(S\) and of type \textup{(RC)}.

Let \(R_1^s\) be the symmetric part of \(R_1\), and put
\(R_1^a=R_1-R_1^s\); then
\[
 \rad^u(P_1)=H_{R_1^a}.
\]
As observed in the proof of \textup{[BT65]}, \SGARefLink{sga3:xxvi:subsection:4:4}{4.4},
\[
 R'=(R_1^s\cap R_2)\cup R_1^a
\]
is a closed subset of \(R\) such that \(R'\cup(-R')=R\).  Hence, by
Exp.\ XXI, \SGARefLink{sga3:XXI:3.3.6}{3.3.6}, \(R'\) contains a system of positive roots of \(R\).
Moreover, the symmetric part of \(R'\) is \(R_1^s\cap R_2^s\), which
is contained in \(R_1\cap R_2\).

It follows that
\[
 P'=K\cdot\rad^u(P_1)
\]
is a parabolic subgroup of \(G\), and that \(K\) is a subgroup of
type \textup{(RC)} of \(P'\) satisfying
\[
 \rad^u(K)=K\cap\rad^u(P').
\]
By \SGARefLink{sga3:xxvi:cor:2:11}{2.11}, \(K\) therefore has a Levi subgroup \(L\).  Exp.\ XIV, \SGARefLink{sga3:xiv:3.20}{3.20}
and \SGARefLink{sga3:xxvi:paragraph:3-21}{3.21} then show that \(K\) satisfies \textup{(iv)}, and also
\textup{(v)} when \(S\) is semilocal.  This proves \SGARefLink{sga3:xxvi:proposition:4-5-1}{Proposition 4.5.1}.
\SGAQED
\footnote{\emph{Editor's note.} The proof of the equivalence, together
with condition \textup{(v)}, which was used implicitly in \SGARefLink{sga3:xxvi:cor:6:17}{6.17} of the
original edition, was added in the Polo--Gille edition.  The original
\SGARefLink{sga3:xxvi:proposition:4-5-1}{4.5.1} was consequently divided into \SGARefLink{sga3:xxvi:proposition:4-5-1}{Proposition 4.5.1} and \SGARefLink{sga3:xxvi:def:4:5:1:1}{Definition
4.5.1.1}.}

\begin{sgastatement}{Definition 4.5.1.1}
When the preceding conditions hold, \(P_1\) and \(P_2\) are said to be
in \emph{mutual standard position}.  This is the case, for example,
when they are in transverse position, when they are in osculating
position, or when the base is the spectrum of a field.  The notion is
stable under base extension and is fpqc-local.
\end{sgastatement}
\SGARefTarget{sga3:xxvi:def:4:5:1:1}
\footnote{\emph{Editor's note.} See the preceding editorial note.  For
an example of parabolic subgroups that are not in standard position,
see \SGARefLink{sga3:xxvi:counterexamples:7:11}{7.11} below.}

\SGAInlineRefTarget{sga3:xxvi:paragraph:4-5-2}\paragraph{\SGARefLink{sga3:xxvi:paragraph:4-5-2}{4.5.2}.}
\SGARefTarget{sga3:xxvi:item:4:5:2}
Let \((P,Q)\) and \((P',Q')\) be two pairs of parabolic subgroups of
\(G\) in standard position.  Define the subfunctor \(H\) of \(G\) by
letting \(H(S')\) be the set of \(g\in G(S')\) such that
\[
 \Int(g)P=P',\qquad \Int(g)Q=Q'.
\]
It is a closed subscheme of \(G\), smooth over \(S\) and formally a
principal homogeneous space under \(P\cap Q\).  It follows that the
following conditions are equivalent:
\begin{enumerate}[label=\textup{(\roman*)}]
\item \((P,Q)\) and \((P',Q')\) are conjugate fpqc-locally;
\item \((P,Q)\) and \((P',Q')\) are conjugate étale-locally;
\item \((P,Q)\) and \((P',Q')\) are conjugate on every geometric
fiber.
\end{enumerate}
The pairs \((P,Q)\) and \((P',Q')\) are then said to have the same
\emph{type of mutual position}.  This notion is stable under base
change and is fpqc-local.

\paragraph{4.5.3.}
\SGARefTarget{sga3:xxvi:item:4:5:3}
Let \(\Stand(G)\) be the subfunctor of
\(\Par(G)\times_S\Par(G)\) formed by the pairs in mutual standard
position.  The functor \(\Stand(G)\) is representable.  There is a
finite étale \(S\)-scheme \(\TypeStand\), the \emph{scheme of types of
mutual standard position}, and a smooth morphism of finite
presentation with irreducible geometric fibers, hence in particular
faithfully flat,
\[
 t_2:\Stand(G)\longrightarrow\TypeStand.
\]
This morphism is a quotient of \(\Stand(G)\) by the action of \(G\):
two sections of \(\Stand(G)\), over any \(S'\to S\), have the same
type of mutual position if and only if they have the same image under
\(t_2\).  There is a commutative diagram
\[
\begin{tikzcd}[column sep=huge,row sep=large]
 \Stand(G) \arrow[r,"t_2"] \arrow[d]
   & \TypeStand \arrow[d,"q"]\\
 \Par(G)\times_S\Par(G) \arrow[r,"t\times t"]
   & \Of(\underline{\operatorname{Dyn}}(G))
       \times_S\Of(\underline{\operatorname{Dyn}}(G)).
\end{tikzcd}
\]
The morphism \(q\) may be described by descent as follows.  Let
\((P,Q)\) be a pair of parabolic subgroups of \(G\) in mutual standard
position, and let \(T\) be a maximal torus of \(P\cap Q\).  The
morphism
\[
 \Norm_G(T)\longrightarrow\Stand(G),\qquad
 n\longmapsto(P,\Int(n)Q),
\]
defined set-theoretically as displayed, induces an isomorphism
\[
 W_P(T)\backslash W_G(T)/W_Q(T)
 \xrightarrow{\sim}q^{-1}\bigl(t(P),t(Q)\bigr).
\]
The left-hand side denotes the sheaf of double cosets.  These
assertions are immediate; note in particular that
\[
 t_2^{-1}\bigl(t_2(P,Q)\bigr)\simeq G/(P\cap Q).
\]

\SGAInlineRefTarget{sga3:xxvi:paragraph:4-5-4}\paragraph{\SGARefLink{sga3:xxvi:paragraph:4-5-4}{4.5.4}.}
\SGARefTarget{sga3:xxvi:item:4:5:4}
Fix a parabolic subgroup \(P\) of \(G\), and let \(\Par(G;P)\) be the
functor of parabolic subgroups of \(G\) in standard position relative
to \(P\).  For every
\[
 t\in\Of(\underline{\operatorname{Dyn}}(G))(S),
\]
put
\[
 \Par_t(G;P)=\Par(G;P)\cap\Par_t(G).
\]
Both functors are obtained from \(\Stand(G)\) by fiber products, and
are therefore represented by \(S\)-schemes, smooth and of finite
presentation over \(S\), with nonempty fibers.  There is a canonical
morphism induced by \(t_2\),
\[
 t_P:\Par_t(G;P)\longrightarrow q^{-1}\bigl(t(P),t\bigr),
 \qquad t_P(Q)=t_2(P,Q),
\]
which is smooth and of finite presentation, with irreducible
geometric fibers.  The canonical morphism
\[
 \Par_t(G;P)\longrightarrow\Par_t(G)
\]
is a surjective monomorphism, and may therefore be viewed as a
\emph{cell decomposition} of \(\Par_t(G)\), indexed by the set of
connected components of \(q^{-1}(t(P),t)\).

\SGAInlineRefTarget{sga3:xxvi:paragraph:4-5-5}\paragraph{4.5.5.}
\SGARefTarget{sga3:xxvi:item:4:5:5}
Suppose that the type \(t\) is of the form \(t(Q)\), where \(Q\) is a
parabolic subgroup of \(G\) in standard position relative to \(P\),
and suppose that \(P\cap Q\) contains a maximal torus \(T\).  Then
\[
 \Par_t(G)\simeq G/Q,\qquad
 q^{-1}\bigl(t(P),t\bigr)
 \simeq W_P(T)\backslash W_G(T)/W_Q(T),
\]
and there is a diagram
\[
\begin{tikzcd}[column sep=huge,row sep=large]
 \Par_{t(Q)}(G;P)
   \arrow[r,"f"] \arrow[d,"i"']
   & W_P(T)\backslash W_G(T)/W_Q(T)\\
 \Par_{t(Q)}(G) \arrow[r,"\sim"]
   & G/Q .
\end{tikzcd}
\]
Here \(i\) is a surjective monomorphism, and \(f\) is smooth and of
finite presentation, with irreducible geometric fibers.  Moreover, if
\(Q_1,Q_2\) are two sections of \(\Par_{t(Q)}(G;P)\), over some
\(S'\to S\), that is, two parabolic subgroups of \(G_{S'}\) which are
fpqc-locally conjugate to \(Q\) and are in standard position relative
to \(P_{S'}\), then \(Q_1,Q_2\) are fpqc-locally conjugate by a
section of \(P\) if and only if
\[
 f(Q_1)=f(Q_2).
\]
If \(S\) is the spectrum of an algebraically closed field \(k\), this
gives
\[
 P(k)\backslash G(k)/Q(k)
 \simeq
 W_P(T)(k)\backslash W_G(T)(k)/W_Q(T)(k).
\]

More generally, suppose that the scheme
\(W_P(T)\backslash W_G(T)/W_Q(T)\) is constant, of the form \(E_S\),
as happens for example when \(G\) is split relative to \(T\) and
\(S\) is connected.  The subschemes \(f^{-1}(e)\), for \(e\in E\),
then decompose \(\Par_{t(Q)}(G;P)\) into open and closed subschemes.
They are homogeneous spaces under \(P\), smooth and of finite
presentation over \(S\), with irreducible geometric fibers.

\SGAInlineRefTarget{sga3:xxvi:paragraph:4-5-6}\paragraph{\SGARefLink{sga3:xxvi:paragraph:4-5-6}{4.5.6}.}
\SGARefTarget{sga3:xxvi:item:4:5:6}
Return to the general situation of \SGARefLink{sga3:xxvi:paragraph:4-5-4}{4.5.4}.  The scheme
\(q^{-1}(t(P),t)\) always has two distinguished sections,
corresponding respectively to the types ``transverse position'' and
``osculating position''.  The inverse image of the first section is
a relatively dense open subscheme of \(\Par_t(G;P)\), as seen above;
it is the cell of greatest relative dimension in the decomposition.

The inverse image of the second section is presumably a closed
subscheme of \(\Par_t(G;P)\); it is the cell of least relative
dimension in the decomposition.
\footnote{\emph{Editor's note.} In this paragraph the Polo--Gille
edition corrects \(t^{-1}(t(P),t)\) to \(q^{-1}(t(P),t)\), and twice
corrects \(\Par_t(G)\) to \(\Par_t(G;P)\).}

% END INCLUDED SGA3_English_current_buildable_source_20260823/source/inputs/XXVI/tex/components/07_section4_5.tex

% BEGIN INCLUDED SGA3_English_current_buildable_source_20260823/source/inputs/XXVI/tex/components/08_section5_1.tex
\SGASection{5}{Conjugacy theorem}
\SGARefTarget{sga3:xxvi:section:5}

\begin{sgastatement}{Theorem 5.1}
Let \(S\) be a semilocal scheme, \(G\) an \(S\)-reductive group, and
\(P,P'\) two opposite parabolic subgroups of \(G\) \textup{(\SGARefLink{sga3:xxvi:definition:4-3-3}{4.3.3})}.
Then
\[
 \rad^u(P)(S)\cdot P'\cdot P=G.
\]
In other words, as \(u\) runs through \(\rad^u(P)(S)\), the open
subsets \(uP'\cdot P\) of \(G\) \textup{(\SGARefLink{sga3:xxvi:remark:4:3:6}{4.3.6})} cover all of \(G\).
\end{sgastatement}
\SGARefTarget{sga3:xxvi:theorem:5:1}

The proof is divided into several steps.

\SGAInlineRefTarget{sga3:xxvi:paragraph:5-1-1}\paragraph{\SGARefLink{sga3:xxvi:paragraph:5-1-1}{5.1.1}.}
\SGARefTarget{sga3:xxvi:item:5:1:1}
It is enough to prove the theorem when \(S\) is the spectrum of a
field \(k\); this follows at once from \SGARefLink{sga3:xxvi:cor:2:6}{2.6}.

\SGAInlineRefTarget{sga3:xxvi:paragraph:5-1-2}\paragraph{\SGARefLink{sga3:xxvi:paragraph:5-1-2}{5.1.2}.}
\SGARefTarget{sga3:xxvi:item:5:1:2}
Put \(L=P\cap P'\).  Suppose that \(L\) has a Borel subgroup \(B_L\),
and let \(T\) be a maximal torus of \(B_L\), as in Exp.\ XXII, \SGARefLink{sga3:xxii:cor:5.9.7}{5.9.7}.
One checks immediately that
\[
 B=B_L\cdot\rad^u(P)
\]
is a Borel subgroup of \(P\).  Let \(B'\) be the Borel subgroup of
\(G\) opposite to \(B\) relative to \(T\), so that \(B\cap B'=T\).
After splitting \(G\) relative to \(T\), one immediately sees that
\(B'\subset P'\).  We claim that
\[
 B^u(S)\cdot B'\cdot B
 \subset \rad^u(P)(S)\cdot P'\cdot P.                 \tag{x}\label{sga3:xxvi:theorem:5-1:equation:x}
\]
Indeed,
\[
 B^u(S)\subset P(S)
 =\rad^u(P)(S)\cdot L(S)
 \subset\rad^u(P)(S)\cdot P'(S),
\]
so it is enough to observe that \(B'\cdot B\subset P'\cdot P\).
Thus \textup{(x)} reduces \SGARefLink{sga3:xxvi:theorem:5:1}{Theorem 5.1} to the pair \((B,B')\).

\SGAInlineRefTarget{sga3:xxvi:paragraph:5-1-3}\paragraph{\SGARefLink{sga3:xxvi:paragraph:5-1-3}{5.1.3}.}
\SGARefTarget{sga3:xxvi:item:5:1:3}
The theorem holds when \(k\) is algebraically closed: the hypothesis
of \SGARefLink{sga3:xxvi:paragraph:5-1-2}{5.1.2} is then satisfied, and Exp.\ XXII, \SGARefLink{sga3:xxii:prop:5.7.10}{5.7.10} applies.

\SGAInlineRefTarget{sga3:xxvi:paragraph:5-1-4}\paragraph{\SGARefLink{sga3:xxvi:paragraph:5-1-4}{5.1.4}.}
\SGARefTarget{sga3:xxvi:item:5:1:4}
The theorem also holds when \(k\) is infinite.  Indeed, by \SGARefLink{sga3:xxvi:cor:2:7}{2.7},
\[
 \rad^u(P)(k)\ \text{is dense in}\ \rad^u(P)(\bar k),
\]
and the theorem holds over \(\bar k\).

\paragraph{5.1.5.}
\SGARefTarget{sga3:xxvi:item:5:1:5}
We are therefore reduced to the case of a finite field \(k\).
The scheme \(\Bor(L)\) is a smooth homogeneous space under \(L\);
Lang's theorem \textup{(Amer.\ J. Math.\ \textbf{78} \textup{(1956)})}
therefore shows that \(L\) has a Borel subgroup \(B_L\).
\footnote{\emph{Editor's note.} See also \textup{[DG70]},
\(\S\)III.5, \SGARefLink{sga3:xxvi:theorem:7-4}{7.4}.}
By 5.1.2, we may consequently suppose that \(P=B\) and \(P'=B'\)
are Borel subgroups.  We write \(T=B\cap B'\).

\SGAInlineRefTarget{sga3:xxvi:paragraph:5-1-6}\paragraph{5.1.6.}
\SGARefTarget{sga3:xxvi:item:5:1:6}
Let \(K\) be an algebraic closure of \(k\).  Choose a pinning of
\((G_K,B_K,T_K)\), and let \(R^+\), respectively \(\Delta\), be the
set of positive, respectively simple, roots.  By Exp.\ XXII, \SGARefLink{sga3:xxii:def:5.7.2}{5.7.2},
it is enough to prove, for every \(\alpha\in\Delta\), that
\[
 u_\alpha B'^{\,u}(K)
 \subset B^u(k)B'^{\,u}(K)B(K).                       \tag{1}\label{sga3:xxvi:theorem:5-1:equation:1}
\]
Let \(\alpha_i\) be the distinct conjugates of \(\alpha\) over \(k\).
They belong to \(\Delta\), since \(B\) is defined over \(k\).  Let
\(R'\) be the set of roots which are linear combinations of the
\(\alpha_i\), and put \(R'_- = R_-\cap R'\).  Since \(R'\) is defined
over \(k\), there is a subtorus \(Q\) of \(T\) such that \(Q_K\) is
the maximal torus of the common kernel of the \(\alpha_i\).

Put
\[
 Z=\Centr_G(Q),\qquad B_Z=B\cap Z,\qquad B'_Z=B'\cap Z
\]
\textup{(Exp.\ XXII, \SGARefLink{sga3:xxii:cor:5.10.2}{5.10.2})}.  It is enough to verify the desired
assertion in \(Z\), namely
\[
 u_\alpha B_Z'^{\,u}(K)
 \subset B_Z^u(k)B_Z'^{\,u}(K)B_Z(K).                 \tag{2}\label{sga3:xxvi:theorem:5-1:equation:2}
\]
Indeed,
\[
 (B_Z'^{\,u})_K=\prod_{\beta\in R'_-}U_\beta.
\]
Let \(R''=R-R'\), and set
\[
 V=\prod_{\beta\in R''\cap R_-}U_\beta.
\]
Then
\[
 (B'^{\,u})_K=(B_Z'^{\,u})_K\cdot V,
\]
and \((B_Z)_K\) normalizes \(V\), by Exp.\ XXII, \SGARefLink{sga3:xxii:cor:5.6.7}{5.6.7}.  From
\textup{\SGARefLink{sga3:xxvi:theorem:5-1:equation:2}{(2)}} we therefore obtain successively
\[
\begin{aligned}
 u_\alpha B'^{\,u}(K)
 &=u_\alpha B_Z'^{\,u}(K)V(K)\\
 &\subset B_Z^u(k)B_Z'^{\,u}(K)B_Z(K)V(K)\\
 &\subset B_Z^u(k)B_Z'^{\,u}(K)V(K)B_Z(K),
\end{aligned}
\]
which immediately gives \textup{\SGARefLink{sga3:xxvi:theorem:5-1:equation:1}{(1)}}.  We are thus reduced to the case
\(G=Z\), that is, to the case in which
\(\operatorname{Gal}(K/k)\) acts transitively on the simple roots.

\SGAInlineRefTarget{sga3:xxvi:paragraph:5-1-7}\paragraph{\SGARefLink{sga3:xxvi:paragraph:5-1-7}{5.1.7}.}
\SGARefTarget{sga3:xxvi:item:5:1:7}
The assertion to be proved says equivalently that \(G/B\) is the
union of the translates, under \(B^u(k)\), of the open image of
\(B'^{\,u}\).  This assertion is unchanged when \(G\) is replaced by
its adjoint group, or indeed by any group having the same adjoint
group.  We may therefore suppose that \(G\) is adjoint.

\SGAInlineRefTarget{sga3:xxvi:paragraph:5-1-8}\paragraph{\SGARefLink{sga3:xxvi:paragraph:5-1-8}{5.1.8}.}
\SGARefTarget{sga3:xxvi:item:5:1:8}
Consider the Dynkin diagram of \(G_K\).  The Galois group acts
transitively on it.  But this Galois group has only cyclic quotients,
whereas the Dynkin diagram has no cycles.  The diagram is consequently
of type \(nA_1\), with \(n\geq0\), or \(mA_2\), with \(m\geq0\).
Using the canonical decomposition of Exp.\ XXIV, \SGARefLink{sga3:xxiv:proposition:5-9}{5.9}, we may write
\[
 G=\prod_{D/k}G_0,
\]
where \(D\) is a finite \(k\)-scheme and \(G_0\) is either a torus or
is of type \(A_1\) or \(A_2\).

By Exp.\ XXIV, \SGARefLink{sga3:xxiv:corollary:5-12}{5.12}, the group \(B\) comes from a Borel subgroup
\(B_0\) of \(G_0\), and \(T\) comes from a maximal torus \(T_0\) of
\(G_0\); the group \(B'\) comes from the Borel subgroup \(B'_0\) of
\(G_0\) opposite to \(B_0\) relative to \(T_0\).  We have
\[
 B^{u}(k)=B_0^u(D),\qquad
 B'^{\,u}\cdot T\cdot B^u
 =\prod_{D/k}B_0'^{\,u}\cdot T_0\cdot B_0^u.
\]
It is therefore enough to prove the assertion for the triple
\((G_0,B_0,T_0)\).

\SGAInlineRefTarget{sga3:xxvi:paragraph:5-1-9}\paragraph{\SGARefLink{sga3:xxvi:paragraph:5-1-9}{5.1.9}.}
\SGARefTarget{sga3:xxvi:item:5:1:9}
We may consequently suppose that \(G\) has type
\(\varnothing\), \(A_1\), or \(A_2\).  Since \(G\) has a Borel
subgroup \(B\), it is quasi-split relative to \(B\), by Exp.\ XXIV,
\SGARefLink{sga3:xxiv:3.9.1}{3.9.1}, and hence split if it has type \(\varnothing\) or \(A_1\).
The theorem is already known in the split case, by Exp.\ XXII,
\SGARefLink{sga3:xxii:prop:5.7.10}{5.7.10}, so only type \(A_2\) remains.

By Exp.\ XXIV, \SGARefLink{sga3:xxiv:3.11}{3.11}, there is a morphism
\(E\to\operatorname{Spec}(k)\), a principal Galois bundle under the
group \(\mathbb Z/2\mathbb Z\) of automorphisms of the Dynkin diagram
of type \(A_2\), such that
\[
 G=\operatorname{Q\text{-}\acute{E}p}_{E/\operatorname{Spec}(k)}
 (A_2).
\]
If \(E\) has a section, \(G\) is split and the theorem follows.
Otherwise \(E=\operatorname{Spec}(k')\), where \(k'/k\) is a quadratic
extension.  Finally, as observed in \SGARefLink{sga3:xxvi:paragraph:5-1-7}{5.1.7}, we may suppose that \(G\)
is simply connected, that is, that \(G\) is a form of
\(\operatorname{SL}_{3,k}\).

\SGAInlineRefTarget{sga3:xxvi:paragraph:5-1-10}\paragraph{\SGARefLink{sga3:xxvi:paragraph:5-1-10}{5.1.10}.}
\SGARefTarget{sga3:xxvi:item:5:1:10}
We are now in the following situation: \(k\) is a finite field and
\(k'/k\) is a quadratic extension.  Pin
\(\operatorname{SL}_{3,k'}\), the group of \(3\times3\) matrices of
determinant \(1\), as follows.  Its maximal torus is the group of
diagonal matrices, its Borel subgroup is the group of upper triangular
matrices, and its root vectors are
\[
 u_\alpha=
 \begin{pmatrix}1&1&0\\0&1&0\\0&0&1\end{pmatrix},
 \qquad
 u_\beta=
 \begin{pmatrix}1&0&0\\0&1&1\\0&0&1\end{pmatrix}.
\]
The big cell \(\Omega\) is given by
\[
 \begin{pmatrix}a&b&c\\d&e&f\\g&h&i\end{pmatrix}
 \in\Omega(S)
 \quad\Longleftrightarrow\quad
 a\ \text{and}\ ae-bd\ \text{are invertible},
\]
and
\[
 B^u(k)=
 \left\{
 \left.
 \begin{pmatrix}1&\bar x&z\\0&1&x\\0&0&1\end{pmatrix}
 \ \right|\
 x,z\in k',\ z+\bar z=x\bar x
 \right\}.
\]

We must prove inclusion \textup{\SGARefLink{sga3:xxvi:theorem:5-1:equation:1}{(1)}} of \SGARefLink{sga3:xxvi:paragraph:5-1-6}{5.1.6}.  Thus, for every
\(a,b,c\in K\), we must find \(x,z\in k'\) such that
\(z+\bar z=x\bar x\) and
\[
 \begin{pmatrix}1&\bar x&z\\0&1&x\\0&0&1\end{pmatrix}
 \begin{pmatrix}1&1&0\\0&1&0\\0&0&1\end{pmatrix}
 \begin{pmatrix}1&0&0\\a&1&0\\b&c&1\end{pmatrix}
 \in\Omega(K).                                         \tag{3}\label{sga3:xxvi:theorem:5-1:equation:3}
\]
If \(a\ne-1\), take \(x=z=0\).  If \(a=-1\), the required conditions
are
\[
 \begin{cases}
 z+\bar z=x\bar x,\\
 bz-\bar x\ne0,\\
 (b+c)\bar z+bx-1\ne0.
 \end{cases}                                           \tag{4}\label{sga3:xxvi:theorem:5-1:equation:4}
\]
Let \(q\geq2\) be the number of elements of \(k\).  For every
\(m\in k\), the equation \(z+\bar z=m\), with \(z\in k'\), has
\(q\) solutions.

If \(b=0\), take \(x=1\).  It remains to solve
\[
 z+\bar z=1,\qquad c\bar z\ne1,
\]
which is always possible by the preceding observation.

If \(b\ne0\), take \(x=0\).  It remains to solve
\[
 z+\bar z=0,\qquad z\ne0,\qquad (b+c)z\ne1.
\]
This is always possible when \(q\geq3\).  When \(k=\mathbb F_2\), it
is possible if \(b+c\ne1\), by taking \(z=1\).

It remains to consider \(k=\mathbb F_2\), \(b+c=1\), and \(b\ne0\).
The system becomes
\[
 z+\bar z=x\bar x,\qquad bz\ne\bar x,\qquad
 \bar z+bx\ne1.
\]
If \(b=1\), or respectively \(b\notin k'\), take \(x=1\); the last two
conditions become
\[
 z\ne b^{-1},\qquad z\ne1-b,
\]
and follow from any solution of \(z+\bar z=1\).  Finally, if
\(b\in k'-k\), take \(x=b\) and \(z=\bar b\).  This completes the
proof.
\SGAQED

% END INCLUDED SGA3_English_current_buildable_source_20260823/source/inputs/XXVI/tex/components/08_section5_1.tex

% BEGIN INCLUDED SGA3_English_current_buildable_source_20260823/source/inputs/XXVI/tex/components/09_section5_2_through_5_5.tex
\begin{sgastatement}{\SGAInlineRefTarget{sga3:xxvi:corollary:5-2}Corollary 5.2}
Let \(S\) be a semilocal scheme, \(G\) an \(S\)-reductive group, and
\(P,P'\) two opposite parabolic subgroups of \(G\).  The canonical map
\[
 \rad^u(P)(S)\cdot\rad^u(P')(S)\longrightarrow(G/P)(S)
\]
is surjective.  In particular,
\[
 (G/P)(S)=G(S)/P(S).
\]
Every parabolic subgroup \(Q\) of \(G\) having the same type as \(P\)
is of the form
\[
 Q=\Int(uu')P
\]
with \(u\in\rad^u(P)(S)\) and \(u'\in\rad^u(P')(S)\).
\end{sgastatement}
\SGARefTarget{sga3:xxvi:cor:5:2}

The second assertion is plainly equivalent to the first, so let us
prove the latter.  Let \(s_i\) be the closed points of \(S\), let
\(V\) be the open subscheme of \(G/P\) which is the image of \(P'\)
\textup{(and is isomorphic to \(\rad^u(P')\), by \SGARefLink{sga3:xxvi:remark:4:3:6}{4.3.6})}, and take
\(x\in(G/P)(S)\).  By \SGARefLink{sga3:xxvi:theorem:5:1}{5.1}, for every \(i\) there is
\[
 u_i\in\rad^u(P)\bigl(\kappa(s_i)\bigr)
\]
such that \(u_ix_{s_i}\) is a section of \(V_{s_i}\).  Choose
\(u\in\rad^u(P)(S)\) lifting all the \(u_i\), as permitted by \SGARefLink{sga3:xxvi:cor:2:6}{2.6}.
Then \(ux\) is a section of \(V\), since this may be checked on the
closed fibers.  Finally,
\[
 \rad^u(P')(S)\xrightarrow{\sim}V(S)
\]
is bijective, and the result follows.

\begin{sgastatement}{Corollary 5.3}
Let \(S\) be a semilocal scheme, \(G\) an \(S\)-reductive group, and
\(P,P',Q\) three parabolic subgroups of \(G\).  There is a
\(g\in G(S)\) such that \(\Int(g)Q\) is in transverse position
relative to both \(P\) and \(P'\).
\end{sgastatement}
\SGARefTarget{sga3:xxvi:cor:5:3}

In the notation of \SGARefLink{sga3:xxvi:corollary:4-2-4}{4.2.4}\textup{(ii)}, we must show that the
universally schematically dense open subscheme
\[
 \Gen(Q/P)\cap\Gen(Q/P')
\]
of \(G\) has a section over \(S\).
\footnote{\emph{Editor's note.} The Polo--Gille edition replaces
``relatively dense'' by ``universally schematically dense over
\(S\)''; compare EGA IV\textsubscript{3}, Definition 11.10.8.}
Choose a parabolic subgroup \(Q_1\) opposite to \(Q\), by
\SGARefLink{sga3:xxvi:corollary:4-3-5}{4.3.5}\textup{(i)}, and put
\[
 U=\rad^u(Q),\qquad U'=\rad^u(Q_1).
\]
We shall find the desired \(g\) in \(U(S)U'(S)\).  In this form,
\SGARefLink{sga3:xxvi:corollary:4-2-4}{4.2.4}\textup{(i)} and \SGARefLink{sga3:xxvi:cor:2:6}{2.6} show that it suffices to work on the fibers
over the closed points of \(S\); we may therefore suppose
\(S=\operatorname{Spec}(k)\).

If \(k\) is algebraically closed, there is a \(g\in G(k)\) with the
required property.  By \SGARefLink{sga3:xxvi:corollary:5-2}{5.2} it may be written
\[
 g=uu'q,\qquad u\in U(k),\quad u'\in U'(k),\quad q\in Q(k),
\]
and \(\Int(uu')Q=\Int(g)Q\).

Suppose \(k\) is infinite.  Let \(V\) be the open subset of
\(U\times_kU'\) defined by the cartesian diagram
\[
\begin{tikzcd}[column sep=huge,row sep=large]
 U\times_kU' \arrow[r] & G\\
 V \arrow[u,hook] \arrow[r]
   & \Gen(Q/P)\cap\Gen(Q/P') \arrow[u,hook].
\end{tikzcd}
\]
Since \(V(\bar k)\ne\varnothing\), as just shown, \(V\) is dense in
\(U\times_kU'\); it therefore has a \(k\)-point by \SGARefLink{sga3:xxvi:cor:2:7}{2.7}.

Suppose finally that \(k\) is finite.  Lang's theorem, as used in
5.1.5, shows that \(P\), respectively \(P'\), has a Borel subgroup
\(B\), respectively \(B'\), since
\[
 \Bor(P)\simeq\Bor(P/\rad^u(P)),\qquad
 \Bor(P')\simeq\Bor(P'/\rad^u(P'))
\]
are smooth.
\footnote{\emph{Editor's note.} The displayed French source prints
\(\Bor(Q)\) in the second isomorphism; the parallel argument requires
\(\Bor(P')\), which is used here.}
Let \(B_1\) be a Borel subgroup opposite to \(B\).  By \SGARefLink{sga3:xxvi:corollary:5-2}{5.2} there are
\[
 a\in B^u(k),\qquad a_1\in B_1^u(k)
\]
such that \(\Int(aa_1)B=B'\).  Then
\[
 B_0=\Int(aa_1)B_1=\Int(a)B_1
\]
is opposite to both \(B'\) and \(B\).  Let \(Q_0\) be the unique
parabolic subgroup of \(G\) containing \(B_0\) and having the same
type as \(Q\), as in \SGARefLink{sga3:xxvi:lemma:3:8}{3.8}.  It is in transverse position relative to
\(P\) and \(P'\), by \SGARefLink{sga3:xxvi:theorem:4-2-1}{4.2.1}\textup{\SGARefLink{sga3:xxvi:theorem:4-2-1:item:vi}{(vi)}}.  On the other hand, \SGARefLink{sga3:xxvi:corollary:5-2}{5.2}
gives
\[
 Q_0=\Int(uu')Q,\qquad u\in U(k),\quad u'\in U'(k),
\]
which proves the assertion.

\begin{sgastatement}{Corollary 5.4}
Let \(S\) be a semilocal scheme, \(G\) an \(S\)-reductive group, and
\(P,Q\) two parabolic subgroups of \(G\).  There is a \(g\in G(S)\)
such that \(\Int(g)P\) is in osculating position relative to \(Q\);
equivalently, \(\Int(g)P\cap Q\) is a parabolic subgroup of \(G\).
\end{sgastatement}
\SGARefTarget{sga3:xxvi:cor:5:4}

Choose a parabolic subgroup \(P'\) opposite to \(P\), by
\SGARefLink{sga3:xxvi:corollary:4-3-5}{4.3.5}\textup{(i)}.  By \SGARefLink{sga3:xxvi:cor:5:3}{5.3}, there is a parabolic subgroup \(P'_1\) of
the same type as \(P'\), in transverse position relative to both
\(P\) and \(Q\).  Let \(T\) be a maximal torus of \(P'_1\cap Q\),
as in \SGARefLink{sga3:xxvi:proposition:4-2-7}{4.2.7}\textup{(ii)}, and let \(P_1\) be the parabolic subgroup
opposite to \(P'_1\) relative to \(T\).  Then \(P_1\) and \(Q\) are
in osculating position, by \SGARefLink{sga3:xxvi:cor:4:4:5}{4.4.5}.  Since \(P\) and \(P_1\) are both
opposite to \(P'_1\), there is a
\[
 g\in\rad^u(P'_1)(S)
\]
such that \(\Int(g)P=P_1\), by \SGARefLink{sga3:xxvi:corollary:4-3-5}{4.3.5}\textup{(i)}.
\SGAQED

For the same reason there is a \(u\in\rad^u(P)(S)\) such that
\(\Int(u)P'=P'_1\).  Hence \(g=\Int(u)u'\) for some
\(u'\in\rad^u(P')(S)\), and
\[
 P_1=\Int(uu'u^{-1})P=\Int(uu')P.
\]
This also gives another proof of \SGARefLink{sga3:xxvi:corollary:5-2}{5.2}.

\SGARefLink{sga3:xxvi:cor:5:3}{Corollaries 5.3} and \SGARefLink{sga3:xxvi:cor:5:4}{5.4} are the essential results of this section.
We first record some consequences of \SGARefLink{sga3:xxvi:cor:5:4}{5.4}.

\begin{sgastatement}{\SGAInlineRefTarget{sga3:xxvi:corollary:5-5}Corollary 5.5}
Let \(S\) be a semilocal scheme and \(G\) an \(S\)-reductive group.
\begin{enumerate}[label=\textup{(\roman*)}]
\item If \(P,Q\) are parabolic subgroups of \(G\) and
\(t(P)\subset t(Q)\), then there is a \(g\in G(S)\) such that
\(\Int(g)P\subset Q\).
\item Let
\[
 P_1\supset P_2\supset\cdots\supset P_n,\qquad
 P'_1\supset P'_2\supset\cdots\supset P'_n
\]
be two chains of parabolic subgroups of \(G\), with
\(t(P_i)=t(P'_i)\) for every \(i\).  There is a \(g\in G(S)\) such
that \(\Int(g)P_i=P'_i\) for every \(i\).
\item Let \(P,Q,P',Q'\) be four parabolic subgroups of \(G\), with
\(t(P)=t(P')\) and \(t(Q)=t(Q')\).  If the pairs \((P,Q)\) and
\((P',Q')\) are both in transverse position, respectively both in
osculating position, then there is a \(g\in G(S)\) such that
\[
 \Int(g)P=P',\qquad\Int(g)Q=Q'.
\]
\item Let \(P,P'\) be parabolic subgroups of the same type, and let
\(L\), respectively \(L'\), be a Levi subgroup of \(P\), respectively
\(P'\).  There is a \(g\in G(S)\) such that
\[
 \Int(g)P=P',\qquad\Int(g)L=L'.
\]
\end{enumerate}
\end{sgastatement}
\SGARefTarget{sga3:xxvi:cor:5:5}

\SGAProof
Assertion \textup{(i)} follows immediately from \SGARefLink{sga3:xxvi:cor:5:4}{5.4}.  Assertion
\textup{(ii)} is proved by induction on \(n\), the case \(n=0\)
being trivial.  We may suppose that \(P_i=P'_i\) for
\(i=1,\ldots,n-1\).  By \SGARefLink{sga3:xxvi:corollary:5-2}{5.2} there is a \(g\in G(S)\) such that
\(\Int(g)P_n=P'_n\).  But \(P_n\) and \(\Int(g)P_n\) are both
contained in \(P_{n-1}=P'_{n-1}\), and hence
\[
 g\in P_{n-1}(S)
\]
by \SGARefLink{sga3:xxvi:lemma:4:4:4}{4.4.4}\textup{(ii)}.  It follows that
\(\Int(g)P_i=P'_i\) for every \(i\).

Assertion \textup{(iv)} follows at once from \SGARefLink{sga3:xxvi:corollary:5-2}{5.2} and \SGARefLink{sga3:xxvi:cor:1:8}{1.8}.  Let us
prove \textup{(iii)} first in the transverse case.  When the types of
\(P\) and \(Q\) are opposite, the assertion follows from
\textup{(iv)} and \SGARefLink{sga3:xxvi:definition:4-3-3}{4.3.3}\textup{(iii)}.  In general, by
\SGARefLink{sga3:xxvi:proposition:4-2-7}{4.2.7}\textup{(iii)} and \SGARefLink{sga3:xxvi:cor:4:4:6}{4.4.6} one can find parabolic subgroups
\[
 P_1\subset P,\quad P'_1\subset P',\quad
 Q_1\subset Q,\quad Q'_1\subset Q'
\]
such that \(P_1,P'_1\) are opposite, \(Q_1,Q'_1\) are opposite, and
\(t(P_1)=t(P'_1)\).  There is consequently a \(g\in G(S)\) such that
\[
 \Int(g)P_1=P'_1,\qquad\Int(g)Q_1=Q'_1.
\]
After conjugating, we may suppose \(P_1=P'_1\) and \(Q_1=Q'_1\).
Then \(P,P'\) are in osculating position and have the same type, so
\(P=P'\) by \SGARefLink{sga3:xxvi:lemma:4:4:4}{4.4.4}\textup{(i)}.  Likewise \(Q=Q'\).

It remains to prove \textup{(iii)} in the osculating case.  By \SGARefLink{sga3:xxvi:corollary:5-2}{5.2}
we may suppose \(P=P'\).  Again by \SGARefLink{sga3:xxvi:corollary:5-2}{5.2} there is a \(g\in G(S)\) such
that \(\Int(g)Q=Q'\).  But then \(g\in P(S)\), by
\SGARefLink{sga3:xxvi:lemma:4:4:4}{4.4.4}\textup{(ii)}, and therefore
\[
 \Int(g)P=P=P'.
\]
\SGAQED

% END INCLUDED SGA3_English_current_buildable_source_20260823/source/inputs/XXVI/tex/components/09_section5_2_through_5_5.tex

% BEGIN INCLUDED SGA3_English_current_buildable_source_20260823/source/inputs/XXVI/tex/components/10_section5_6_through_5_11.tex
\begin{sgastatement}{Definition 5.6}
Let \(S\) be a scheme, \(G\) an \(S\)-reductive group, and \(P\) a
parabolic subgroup of \(G\).  The subgroup \(P\) is said to be
\emph{minimal} if every parabolic subgroup \(Q\) of \(G\) contained
in \(P\) is equal to \(P\).
\end{sgastatement}
\SGARefTarget{sga3:xxvi:def:5:6}

Notice that this notion is not, in general, stable under passage to
the fibers.

\begin{sgastatement}{\SGAInlineRefTarget{sga3:xxvi:corollary:5-7}Corollary 5.7}
Let \(S\) be a semilocal scheme and \(G\) an \(S\)-reductive group.
\begin{enumerate}[label=\textup{(\roman*)}]
\item Let
\[
 t,t'\in\Of\bigl(\operatorname{Dyn}(G)\bigr)(S).
\]
If \(G\) has a parabolic subgroup of type \(t\) and one of type
\(t'\), then it has a parabolic subgroup of type \(t\cap t'\).  In
particular, the set of \(t(P)\), as \(P\) ranges over the parabolic
subgroups of \(G\), has a least element \(t_{\min}\).
\item Every parabolic subgroup of \(G\) contains a minimal parabolic
subgroup.  A parabolic subgroup of \(G\) is minimal if and only if it
has type \(t_{\min}\).  Any two minimal parabolic subgroups of \(G\)
are conjugate by an element of \(G(S)\).
\end{enumerate}
\end{sgastatement}
\SGARefTarget{sga3:xxvi:cor:5:7}

This follows immediately from \SGARefLink{sga3:xxvi:cor:5:4}{5.4} and \SGARefLink{sga3:xxvi:corollary:5-5}{5.5}\textup{(i)}.

\begin{sgastatement}{Remark 5.8}
A parabolic subgroup opposite to a minimal parabolic subgroup is
also minimal.  Consequently,
\[
 s(t_{\min})=t_{\min}.
\]
\end{sgastatement}
\SGARefTarget{sga3:xxvi:rem:5:8}

\begin{sgastatement}{\SGAInlineRefTarget{sga3:xxvi:corollary:5-9}Corollary 5.9}
Let \(S\) be a scheme, \(G\) an \(S\)-reductive group, and \(P\) a
parabolic subgroup of \(G\).  The canonical morphism
\[
 G\longrightarrow G/P=X
\]
makes \(G\) a Zariski-locally trivial \(P_X\)-bundle over \(X\).
If \(L\) is a Levi subgroup of \(P\), the canonical morphism
\textup{(cf.\ \SGARefLink{sga3:xxvi:cor:3:12}{3.12})}
\[
 G\longrightarrow G/L=Y
\]
makes \(G\) a Zariski-locally trivial \(L_Y\)-bundle over \(Y\).
\end{sgastatement}
\SGARefTarget{sga3:xxvi:cor:5:9}

It is enough to prove the following lifting assertion.  Given a
morphism \(S'\to S\), with \(S'\) local, every morphism \(S'\to X\)
\textup{(respectively \(S'\to Y\))} lifts to a morphism \(S'\to G\).
In other words, we may assume that \(S\) is local, and must show that
\[
 G(S)\longrightarrow X(S),
 \qquad\text{respectively}\qquad
 G(S)\longrightarrow Y(S),
\]
is surjective.

The first assertion was proved in \SGARefLink{sga3:xxvi:corollary:5-2}{5.2}.  For the second, let
\(y\in Y(S)\).  Its canonical image in \(X(S)\) is the image of some
\(g\in G(S)\).  The image \(y'\) of \(g\) in \(Y(S)\) therefore has
the same image in \(X(S)\) as \(y\).  There is consequently a unique
\[
 u\in\rad^u(P)(S)
\]
such that \(y'u=y\), and the image of \(gu\) in \(Y(S)\) is indeed
\(y\).

\begin{sgastatement}{\SGAInlineRefTarget{sga3:xxvi:corollary:5-10}Corollary 5.10}
Let \(S\) be a semilocal scheme and \(G\) an \(S\)-reductive group.
\begin{enumerate}[label=\textup{(\roman*)}]
\item Let \(P\) be a parabolic subgroup of \(G\), and \(L\) a Levi
subgroup of \(P\).  The canonical maps \textup{(cf.\ 3.21)}
induce bijections
\[
 H^1(S,L)\xrightarrow{\ \sim\ }H^1(S,P)
       \xrightarrow{\ \sim\ }H^1_{t(P)}(S,G).
\]
\item If
\[
 t,t'\in\Of\bigl(\operatorname{Dyn}(G)\bigr)(S),
\]
then \textup{(cf.\ 3.21)}
\[
 H^1_t(S,G)\cap H^1_{t'}(S,G)=H^1_{t\cap t'}(S,G).
\]
\item If \(P\) and \(Q\) are two parabolic subgroups of \(G\) in
osculating position, the following canonical diagram is cartesian
and all its arrows are injective:
\[
\begin{tikzcd}[column sep=huge,row sep=large]
 H^1(S,P) \arrow[r] & H^1(S,G)\\
 H^1(S,P\cap Q) \arrow[u] \arrow[r] & H^1(S,Q). \arrow[u]
\end{tikzcd}
\]
\end{enumerate}
\end{sgastatement}
\SGARefTarget{sga3:xxvi:cor:5:10}

Let us prove \textup{(i)}.  The map
\[
 H^1(S,L)\longrightarrow H^1(S,P)
\]
is bijective by \SGARefLink{sga3:xxvi:cor:2:3}{2.3}, while
\[
 H^1(S,P)\longrightarrow H^1_{t(P)}(S,G)
\]
is surjective by 3.21.  We must prove that it is injective, or
equivalently that the canonical map
\[
 H^1(S,P)\longrightarrow H^1(S,G)
\]
is injective.  Let \(Q\) be a principal homogeneous bundle under
\(P\), let \(Q_1\) be the associated principal homogeneous bundle
under \(G\), and let \(P'\) and \(G'\) be the corresponding twisted
forms of \(P\) and \(G\).  Clearly \(G'\) is an \(S\)-reductive group
and \(P'\) is a parabolic subgroup of \(G'\).

The elements of \(H^1(S,P)\) having the same image in \(H^1(S,G)\)
as the class of \(Q\) are naturally identified with the kernel of
\[
 H^1(S,P')\longrightarrow H^1(S,G').
\]
By the exact sequence in cohomology, this kernel is in turn
identified with the set of \(G'(S)\)-orbits in
\((G'/P')(S)\).\footnote{For these arguments in nonabelian
cohomology, see Giraud's thesis, \textup{[Gi71]}, \S~III.3.}
But \(G'(S)\) acts transitively on \((G'/P')(S)\), by \SGARefLink{sga3:xxvi:corollary:5-2}{5.2}.

For \textup{(ii)}, let \(Q\) be a principal homogeneous bundle under
\(G\), and let \(G_Q\) be the corresponding twisted form of \(G\).
By definition \textup{(3.21)}, we must prove that \(G_Q\) has a
parabolic subgroup of type \(t\cap t'\) if and only if it has
parabolic subgroups of types \(t\) and \(t'\).  This is precisely the
combination of \SGARefLink{sga3:xxvi:lemma:3:8}{3.8} and \SGARefLink{sga3:xxvi:corollary:5-7}{5.7}\textup{(i)}.  Finally, \textup{(iii)}
follows immediately from \textup{(i)} and \textup{(ii)}.

We now state a consequence of \SGARefLink{sga3:xxvi:cor:5:3}{5.3}.

\begin{sgastatement}{Corollary 5.11}
Let \(S\) be a semilocal scheme, \(G\) an \(S\)-reductive group, and
\(P\) a parabolic subgroup of \(G\).  If \(P\ne G\), there exist at
least three distinct parabolic subgroups of \(G\) having the same
type as \(P\).  In other words,
\[
 P\ne G\qquad\Longrightarrow\qquad
 \bigl(G(S):P(S)\bigr)\geq 3.
\]
\end{sgastatement}
\SGARefTarget{sga3:xxvi:cor:5:11}

Indeed, let \(P'\) be a parabolic subgroup opposite to \(P\), as in
\SGARefLink{sga3:xxvi:corollary:4-3-5}{4.3.5}\textup{(i)}.  Since \(P\ne G\), one has
\(\rad^u(P')\ne e\), for example by \SGARefLink{sga3:xxvi:item:4:3:1}{4.3.2}.  By \SGARefLink{sga3:xxvi:proposition:2-1}{2.1},
\(\rad^u(P')(S)\ne e\); choose
\[
 u\in\rad^u(P')(S),\qquad u\ne e.
\]
Then \(\Int(u)P\ne P\).  By \SGARefLink{sga3:xxvi:cor:5:3}{5.3} there is a parabolic subgroup \(P_1\)
of the same type as \(P\), opposite to both \(P\) and \(\Int(u)P\).
Thus \(P_1\), \(P\), and \(\Int(u)P\) are three distinct parabolic
subgroups of \(G\), all of the same type.
\SGAQED

% END INCLUDED SGA3_English_current_buildable_source_20260823/source/inputs/XXVI/tex/components/10_section5_6_through_5_11.tex

% BEGIN INCLUDED SGA3_English_current_buildable_source_20260823/source/inputs/XXVI/tex/components/11_section6_1_through_6_8.tex
\SGASection{6}{Parabolic subgroups and split tori}
\SGARefTarget{sga3:xxvi:sec:6}

\begin{sgastatement}{\SGAInlineRefTarget{sga3:xxvi:proposition:6-1}Proposition 6.1}
Let \(S\) be a scheme, \(G\) an \(S\)-reductive group,
\(\mathfrak g=\operatorname{Lie}(G/S)\), and \(Q\) a split subtorus
of \(G\).  Write \(Q=D_S(M)\), and let
\[
 \mathfrak g=\bigoplus_{\alpha\in M}\mathfrak g^\alpha
\]
be the decomposition of \(\mathfrak g\) under the action of \(Q\).
Let \(M_1\subset M\) satisfy
\[
 0\in M_1,\qquad
 \alpha,\beta\in M_1\Longrightarrow\alpha+\beta\in M_1.
\]
\begin{enumerate}[label=\textup{(\roman*)}]
\item There is a unique smooth subgroup \(H_{M_1}\) of \(G\), with
connected fibers, which contains \(\Centr_G(Q)\) and whose Lie
algebra is
\[
 \bigoplus_{\alpha\in M_1}\mathfrak g^\alpha.
\]
\item The following implications hold:
\[
\begin{aligned}
 M_1=\{0\}
   &\Longrightarrow H_{M_1}=\Centr_G(Q),\\
 M_1=-M_1
   &\Longrightarrow H_{M_1}\text{ is reductive},\\
 M_1\cup(-M_1)=M
   &\Longrightarrow
   \begin{gathered}
    H_{M_1}\text{ and }H_{-M_1}\text{ are opposite parabolic}\\
    \text{subgroups of \(G\), with common Levi subgroup }
    H_{M_1\cap(-M_1)}.
   \end{gathered}
\end{aligned}
\]
\end{enumerate}
\end{sgastatement}
\SGARefTarget{sga3:xxvi:prop:6:1}

Assertions \textup{(i)} and \textup{(ii)} are local for the fpqc
topology.  We may therefore suppose that \(Q\) is contained in a
maximal torus \(T\) of \(G\), and may moreover split \(G\) relative
to \(T\).  Assertion \textup{(i)} then follows immediately from
\SGARefLink{sga3:xxii:cor:5.3.5}{Exposé XXII, 5.3.5}, \SGARefLink{sga3:xxii:prop:5.4.5}{5.4.5}, and \SGARefLink{sga3:xxii:thm:5.4.7}{5.4.7}.  The assertions in
\textup{(ii)} follow from \SGARefLink{sga3:xxii:cor:5.3.5}{Exposé XXII, 5.3.5}, \SGARefLink{sga3:xxii:prop:5.10.1}{5.10.1}, and \SGARefLink{sga3:xxii:lem:5.11.3}{5.11.3},
together with \SGARefLink{sga3:xxvi:prop:1:3}{1.4} and \SGARefLink{sga3:xxvi:item:4:3:1}{4.3.2} of the present Exposé.

\begin{sgastatement}{Corollary 6.2}
Let \(S\) be a scheme, \(G\) an \(S\)-reductive group, and \(Q\) a
split subtorus of \(G\).  There is a parabolic subgroup of \(G\)
having \(\Centr_G(Q)\) as a Levi subgroup.
\end{sgastatement}
\SGARefTarget{sga3:xxvi:cor:6:2}

Indeed, write \(Q=D_S(M)\), choose a total order on the group \(M\),
and let \(M_1\) be the set of its positive elements.  The subgroup
\(H_{M_1}\) has the required property.

\begin{sgastatement}{Corollary 6.3}
If the \(S\)-reductive group \(G\) has a noncentral split subtorus,
then it has a proper parabolic subgroup \textup{(that is, one
different from \(G\))}.
\end{sgastatement}
\SGARefTarget{sga3:xxvi:cor:6:3}

Combining \SGARefLink{sga3:xxvi:cor:6:2}{6.2} with \SGARefLink{sga3:xxvi:corollary:5-9}{5.9} and \SGARefLink{sga3:xxvi:corollary:5-10}{5.10} gives the following result.

\begin{sgastatement}{\SGAInlineRefTarget{sga3:xxvi:corollary:6-4}Corollary 6.4}
Let \(S\) be a scheme, \(G\) an \(S\)-reductive group, and \(Q\) a
split subtorus of \(G\).  The canonical morphism
\[
 G\longrightarrow G/\Centr_G(Q)
\]
is a locally trivial fibration.

If \(S\) is semilocal, the map
\[
 G(S)\longrightarrow\bigl(G/\Centr_G(Q)\bigr)(S)
\]
is surjective, and the map
\[
 H^1\bigl(S,\Centr_G(Q)\bigr)\longrightarrow H^1(S,G)
\]
is injective.
\end{sgastatement}
\SGARefTarget{sga3:xxvi:cor:6:4}

\SGASubsection{6.5}{}
\SGARefTarget{sga3:xxvi:subsec:6:5}

Suppose that \(S\) is connected.  If \(T\) is an \(S\)-torus and
\(T'\), \(T''\) are two split subtori of \(T\), their product
\[
 T'\cdot T''
\]
is again a split subtorus of \(T\).\footnote{\emph{Editor's note.}
The Polo--Gille edition adopts multiplicative notation here,
replacing the source's ``sum \(T'+T''\)'' by ``product
\(T'\cdot T''\)''.}  Indeed, it is identified with the quotient of
\(T'\times_S T''\) by \(T'\cap T''\), and this quotient is split by
\SGARefLink{sga3:ix:proposition:2:11}{Exposé IX, 2.11}.  It follows that \(T\) has a largest split subtorus,
which we denote by \(T_{\mathrm{spl}}\).

\begin{sgastatement}{\SGAInlineRefTarget{sga3:xxvi:lemma:6-6}Lemma 6.6}
Let \(S\) be a connected scheme, \(T\) an isotrivial \(S\)-torus,
and \(T_{\mathrm{spl}}\) its largest split subtorus.  The following
conditions are equivalent:
\begin{enumerate}[label=\textup{(\roman*)}]
\item There is a homomorphism
\[
 T\longrightarrow\mathbf G_{m,S}
\]
different from the trivial homomorphism.
\item \(T_{\mathrm{spl}}\ne e\).
\end{enumerate}
\end{sgastatement}
\SGARefTarget{sga3:xxvi:lem:6:6}

Since \(T\) is isotrivial, there are a finite group \(\Gamma\), a
connected principal Galois covering \(S'\to S\) with group
\(\Gamma\), and an isomorphism
\[
 T_{S'}\simeq D_{S'}(M).
\]
The group \(M\) is then a \(\Gamma\)-module, and there is a natural
isomorphism
\[
 \operatorname{Hom}_S(T,\mathbf G_{m,S})
   =H^0(\Gamma,M).
\]

Let \(V\) be the vector subspace of \(M\otimes\mathbf Q\) generated
by the elements
\[
 \gamma(m)-m,\qquad \gamma\in\Gamma,\quad m\in M.
\]
One checks immediately that
\[
 (T_{\mathrm{spl}})_{S'}
   \simeq D_{S'}\bigl(M/(M\cap V)\bigr).
\]
Thus \textup{(i)} is equivalent to
\[
 H^0(\Gamma,M)\ne0,
\]
or equivalently to
\[
 H^0(\Gamma,M\otimes\mathbf Q)\ne0,
\]
whereas \textup{(ii)} is equivalent to \(M\ne M\cap V\), or to
\(M\otimes\mathbf Q\ne V\).  But
\[
 M\otimes\mathbf Q
 =H^0(\Gamma,M\otimes\mathbf Q)\oplus V.
\]
This follows at once by considering the projector from
\(M\otimes\mathbf Q\) onto its invariants which sends \(x\) to the
average of its transforms under \(\Gamma\).

\begin{sgastatement}{Lemma 6.7}
Let \(S\) be a scheme, \(G\) an \(S\)-reductive group, \(P\ne G\) a
parabolic subgroup of \(G\), \(L\) a Levi subgroup of \(P\), and
\(Q=\rad(L)\).  There is a homomorphism
\[
 Q\longrightarrow\mathbf G_{m,S}
\]
different from the trivial homomorphism.
\end{sgastatement}
\SGARefTarget{sga3:xxvi:lem:6:7}

Let \(U=\rad^u(P)\).  It is invariant under \(\Int(P)\), hence under
\(\Int(Q)\).  Consider the invertible \(\mathcal O_S\)-module
\[
 \det\bigl(\operatorname{Lie}(U)\bigr),
\]
the highest exterior power of the locally free
\(\mathcal O_S\)-module \(\operatorname{Lie}(U)\).  The adjoint
representation defines a homomorphism
\[
 f:Q\longrightarrow
 \operatorname{Aut}\!\left(\det\bigl(\operatorname{Lie}(U)\bigr)\right)
 =\mathbf G_{m,S}.
\]
Since \(P\ne G\), one has \(U\ne e\).  Choose \(s\in S\) such that
\(U_s\ne e\).  After splitting \(G_s\) relative to a maximal torus
containing \(Q_s\), it is immediate that \(f_s\ne e\).

\begin{sgastatement}{Proposition 6.8}
Let \(S\) be a connected semilocal scheme, \(G\) an \(S\)-reductive
group, \(P\) a parabolic subgroup of \(G\), \(L\) a Levi subgroup of
\(P\), and \(Q=\rad(L)\).  Let \(Q_{\mathrm{spl}}\) be the largest
split subtorus of \(Q\), equivalently the largest central split
subtorus of \(L\).  Then
\[
 L=\Centr_G(Q_{\mathrm{spl}}).
\]
\end{sgastatement}
\SGARefTarget{sga3:xxvi:prop:6:8}

Put \(L'=\Centr_G(Q_{\mathrm{spl}})\).  This is a reductive subgroup
of \(G\) containing \(L\), and
\[
 P'=P\cap L'
\]
is a parabolic subgroup of \(L'\) with Levi subgroup \(L\), by
\SGARefLink{sga3:xxvi:proposition:1-20}{1.20}.  If \(L'\ne L\), then \(L'\not\subset P\), since \(L\) is a
maximal reductive subgroup of \(P\) by \SGARefLink{sga3:xxvi:prop:1:6}{1.7}.  Hence \(P'\ne L'\).

Let \(G_1\) be the derived group of \(L'\), and put
\(P_1=P'\cap G_1\).  By \SGARefLink{sga3:xxvi:proposition:1-19}{1.19}, \(P_1\) is a parabolic subgroup of the
semisimple group \(G_1\), the group \(L_1=L\cap G_1\) is a Levi
subgroup of \(P_1\), and
\[
 Q_1=\rad(L_1)
    =\bigl(\rad(L)\cap G_1\bigr)^0
    =\bigl(Q\cap G_1\bigr)^0.
\]
The group \(L_1\) has a maximal torus \(T_1\), by \SGARefLink{sga3:xiv:3.20}{Exposé XIV, 3.20},
and \(T_1\) is isotrivial, by \SGARefLink{sga3:xxiv:theorem:4-1-5}{Exposé XXIV, 4.1.5}.  Therefore its
subtorus \(Q_1\) is also isotrivial, by \SGARefLink{sga3:ix:proposition:2:11}{Exposé IX, 2.11}.  Since
\(P_1\ne G_1\), \SGARefLink{sga3:xxvi:lem:6:7}{Lemmas 6.7} and \SGARefLink{sga3:xxvi:lemma:6-6}{6.6} give
\[
 (Q_1)_{\mathrm{spl}}\ne e.
\]
It follows that
\[
 \bigl(Q_{\mathrm{spl}}\cap G_1\bigr)^0\ne e,
\]
and hence that \(Q_{\mathrm{spl}}\not\subset\rad(L')\), because
\(\rad(L')\cap G_1\) is finite.  This contradicts the definition of
\(L'\).
\SGAQED

% END INCLUDED SGA3_English_current_buildable_source_20260823/source/inputs/XXVI/tex/components/11_section6_1_through_6_8.tex

% BEGIN INCLUDED SGA3_English_current_buildable_source_20260823/source/inputs/XXVI/tex/components/12_section6_9_through_6_18.tex
\begin{sgastatement}{\SGAInlineRefTarget{sga3:xxvi:corollary:6-9}Corollary 6.9}
Let \(S\) be a connected semilocal scheme, \(G\) an
\(S\)-reductive group, and \(Q\) a critical subtorus of \(G\), that
is, a subtorus such that
\[
 \rad\bigl(\Centr_G(Q)\bigr)=Q.
\]
The group \(\Centr_G(Q)\) is a Levi subgroup of a parabolic subgroup
of \(G\) if and only if
\[
 \Centr_G(Q)=\Centr_G(Q_{\mathrm{spl}}),
\]
or equivalently,
\[
 \operatorname{Lie}(G)^Q
   =\operatorname{Lie}(G)^{Q_{\mathrm{spl}}}.
\]
\end{sgastatement}
\SGARefTarget{sga3:xxvi:cor:6:9}

This follows from \SGARefLink{sga3:xxvi:cor:6:2}{6.2} and \SGARefLink{sga3:xxvi:prop:6:8}{6.8}.

\begin{sgastatement}{\SGAInlineRefTarget{sga3:xxvi:corollary:6-10}Corollary 6.10}
Let \(S\) be a connected semilocal scheme, \(G\) an
\(S\)-reductive group, and \(L\) a subgroup of \(G\).  The following
conditions are equivalent:
\begin{enumerate}[label=\textup{(\roman*)}]
\item There is a parabolic subgroup of \(G\) having \(L\) as a Levi
subgroup.
\item There is a split subtorus of \(G\) whose centralizer is \(L\).
\item There is a homomorphism
\[
 \mathbf G_{m,S}\longrightarrow G
\]
whose centralizer is \(L\).
\end{enumerate}
\end{sgastatement}
\SGARefTarget{sga3:xxvi:cor:6:10}

The implication \textup{(i)}\(\Rightarrow\)\textup{(ii)} follows
from \SGARefLink{sga3:xxvi:prop:6:8}{6.8}, and \textup{(iii)}\(\Rightarrow\)\textup{(i)} from \SGARefLink{sga3:xxvi:proposition:6-1}{6.1}.
It remains to prove
\textup{(ii)}\(\Rightarrow\)\textup{(iii)}.  Suppose therefore that
\[
 L=\Centr_G(Q),\qquad Q=D_S(M).
\]
Write \(\mathfrak g=\operatorname{Lie}(G/S)\) and
\[
 \mathfrak g=\bigoplus_{\alpha\in M}\mathfrak g^\alpha,
\]
and let \(R\) be the set of
\(\alpha\in M\setminus\{0\}\) for which
\(\mathfrak g^\alpha\ne0\).

Since \(R\) is finite and does not contain \(0\), there is a
homomorphism
\[
 u:M\longrightarrow\mathbf Z
\]
such that \(u(\alpha)\ne0\) for every \(\alpha\in R\).  By duality,
\(u\) yields a homomorphism \(\mathbf G_{m,S}\to Q\), and hence a
homomorphism
\[
 f:\mathbf G_{m,S}\longrightarrow G.
\]
One has
\[
 \Centr_G(f)\supset\Centr_G(Q).
\]
These are two smooth subgroups of \(G\), with connected fibers, and
their Lie algebras coincide, since both are equal to
\(\mathfrak g^0\).  They are therefore equal, by the usual
argument.

\begin{sgastatement}{\SGAInlineRefTarget{sga3:xxvi:corollary:6-11}Corollary 6.11}
Let \(S\) be a connected semilocal scheme and \(G\) an
\(S\)-reductive group.  The maps
\[
 L\longmapsto\rad(L)_{\mathrm{spl}},
 \qquad
 Q\longmapsto\Centr_G(Q)
\]
are mutually inverse, order-reversing bijections between:
\begin{itemize}
\item the set of subgroups \(L\) of \(G\) which are Levi subgroups
of parabolic subgroups of \(G\); and
\item the set of split subtori \(Q\) of \(G\) such that
\[
 \rad\bigl(\Centr_G(Q)\bigr)_{\mathrm{spl}}=Q.
\]
\end{itemize}
\end{sgastatement}
\SGARefTarget{sga3:xxvi:cor:6:11}

\begin{sgastatement}{Corollary 6.12}
Let \(S\) be a connected semilocal scheme and \(G\) an
\(S\)-reductive group.  Consider the following assertions:
\begin{enumerate}[label=\textup{(\roman*)}]
\item \(G\) has a parabolic subgroup different from \(G\).
\item \(G\) has a noncentral split subtorus.
\item[\textup{(ii bis)}] \(G\) has a noncentral split subtorus of
relative dimension \(1\).
\item There is a group homomorphism
\[
 \mathbf G_{a,S}\longrightarrow G
\]
which is a closed immersion.
\end{enumerate}
Then
\[
 \textup{(i)}
 \Longleftrightarrow\textup{(ii)}
 \Longleftrightarrow\textup{(ii bis)}
 \Longrightarrow\textup{(iii)}.
\]
\end{sgastatement}
\SGARefTarget{sga3:xxvi:cor:6:12}

The only new assertion is
\textup{(i)}\(\Rightarrow\)\textup{(iii)}.  Let \(P\ne G\) be a
parabolic subgroup of \(G\).  Then \(U=\rad^u(P)\ne e\).  Consider
the last nontrivial subgroup \(U_n\) in the composition series of
\(U\), as in \SGARefLink{sga3:xxvi:proposition:2-1}{2.1}.  There is an isomorphism
\[
 U_n\simeq W(E_n),
\]
where \(E_n\) is a locally free \(\mathcal O_S\)-module, and hence
is free because \(S\) is semilocal and connected.  Since
\(E_n\ne0\), there is a monomorphism
\[
 \mathcal O_S\longrightarrow E_n
\]
whose image is locally a direct summand.  It induces a closed
immersion
\[
 \mathbf G_{a,S}=W(\mathcal O_S)
   \lhook\joinrel\longrightarrow W(E_n)\simeq U_n,
\]
which proves \textup{(iii)}.

\begin{sgastatement}{Remark 6.12.1}
When \(S\) is the spectrum of a field of characteristic \(0\), the
Jacobson--Morozov theorem gives
\textup{(iii)}\(\Rightarrow\)\textup{(ii bis)}.  The four preceding
conditions are then equivalent \textup{(``Godement's criterion'';
cf.\ \textup{[BT65]}, 8.5)}.\footnote{More generally, this remains
true when \(S\) is the spectrum of a perfect field \textup{(Tits)};
see \textup{[BT71]}, Corollary 3.7.}
\end{sgastatement}
\SGARefTarget{sga3:xxvi:rem:6:12:1}

\begin{sgastatement}{Definition 6.13}
Let \(S\) be a connected semilocal scheme and \(G\) an
\(S\)-reductive group.  The group \(G\) is said to be
\emph{anisotropic} if it contains no nontrivial split subtorus.
\end{sgastatement}
\SGARefTarget{sga3:xxvi:def:6:13}

\begin{sgastatement}{Corollary 6.14}
Let \(S\) be a connected semilocal scheme.  An \(S\)-reductive group
\(G\) is anisotropic if and only if it has no parabolic subgroup
\(P\ne G\) and its radical is anisotropic.
\end{sgastatement}
\SGARefTarget{sga3:xxvi:cor:6:14}

Combining \SGARefLink{sga3:xxvi:lemma:6-6}{6.6} with \SGARefLink{sga3:xxiv:theorem:4-1-5}{Exposé XXIV, 4.1.5}, and \SGARefLink{sga3:xxii:subsection:6-2}{Exposé XXII, 6.2} gives:

\begin{sgastatement}{Corollary 6.15}
Let \(S\) be a connected semilocal scheme and \(G\) an isotrivial
\(S\)-reductive group \textup{(for example, \(G\) semisimple, or
\(S\) normal; cf.\ \SGARefLink{sga3:x:theorem:5-16}{Exposé X, 5.16})}.  The group \(G\) is anisotropic
if and only if it has no parabolic subgroup \(P\ne G\) and
\[
 \operatorname{Hom}_{S\text{-}\mathrm{gr}}
   (G,\mathbf G_{m,S})=e.
\]
\end{sgastatement}
\SGARefTarget{sga3:xxvi:cor:6:15}

\begin{sgastatement}{Proposition 6.16}
Let \(S\) be a connected semilocal scheme and \(G\) an
\(S\)-reductive group.  The maximal split subtori of \(G\) are
precisely the largest central split subtori of the Levi subgroups of
the minimal parabolic subgroups of \(G\).  Any two such tori are
conjugate by an element of \(G(S)\).
\end{sgastatement}
\SGARefTarget{sga3:xxvi:prop:6:16}

Let \(Q\) be a maximal split subtorus of \(G\).  By \SGARefLink{sga3:xxvi:cor:6:2}{6.2},
\[
 L=\Centr_G(Q)
\]
is a Levi subgroup of a parabolic subgroup \(P\) of \(G\).  Since
\[
 Q\subset\rad\bigl(\Centr_G(Q)\bigr)_{\mathrm{spl}},
\]
the maximality of \(Q\) gives
\[
 Q=\rad\bigl(\Centr_G(Q)\bigr)_{\mathrm{spl}}.
\]
By \SGARefLink{sga3:xxvi:corollary:6-11}{6.11}, \(L\) is minimal among the Levi subgroups of parabolic
subgroups of \(G\).  Hence \(P\) is a minimal parabolic subgroup of
\(G\), by \SGARefLink{sga3:xxvi:proposition:1-20}{1.20}.  It follows from \SGARefLink{sga3:xxvi:corollary:5-7}{5.7} and \SGARefLink{sga3:xxvi:corollary:5-5}{5.5}\textup{(iv)} that any
two such subtori \(Q\) are conjugate by an element of \(G(S)\).
The conjugacy of the tori \(Q\) and of the pairs \((P,L)\) also
proves the first assertion.

\begin{sgastatement}{Corollary 6.17}
Let \(S\) be a connected semilocal scheme, and let \(P,P'\) be two
minimal parabolic subgroups in standard position
\textup{(\SGARefLink{sga3:xxvi:def:4:5:1:1}{4.5.1.1})}.  Then \(P\cap P'\) contains a common Levi
subgroup of \(P\) and \(P'\).
\end{sgastatement}
\SGARefTarget{sga3:xxvi:cor:6:17}

Indeed, \(P\cap P'\) contains a maximal torus \(T\) of \(G\), by
\SGARefLink{sga3:xxvi:proposition:4-5-1}{4.5.1}\textup{(v)}.  Let \(L\) be the unique Levi subgroup of \(P\)
which contains \(T\).  By \SGARefLink{sga3:xxvi:proposition:1-21}{1.21},
\[
 \rad(P)\cap T=\rad(P)\cap L=\rad(L).
\]
Thus \(\rad(P)\cap T\) contains
\(\rad(L)_{\mathrm{spl}}\), which is a maximal split subtorus of
\(G\), and is therefore necessarily equal to
\(T_{\mathrm{spl}}\).  Hence
\[
 L=\Centr_G(T_{\mathrm{spl}}),
\]
and, by symmetry, \(L\) is also a Levi subgroup of \(P'\).

\begin{sgastatement}{Remark 6.18}
It follows from \SGARefLink{sga3:xxvi:proposition:1-21}{1.21} that a parabolic subgroup \(P\) of \(G\) is
minimal if and only if \(\rad(P)\) contains a maximal split subtorus
of \(G\).  In that case, if \(T\) is a maximal torus of \(P\), the
torus \(T_{\mathrm{spl}}\) is a maximal split subtorus both of \(G\)
and of \(\rad(P)\), and \(\Centr_G(T_{\mathrm{spl}})\) is a Levi
subgroup of \(P\).  Moreover, every Levi subgroup of \(P\) is
obtained in this way.
\end{sgastatement}
\SGARefTarget{sga3:xxvi:rem:6:18}

% END INCLUDED SGA3_English_current_buildable_source_20260823/source/inputs/XXVI/tex/components/12_section6_9_through_6_18.tex

% BEGIN INCLUDED SGA3_English_current_buildable_source_20260823/source/inputs/XXVI/tex/components/13_section7_opening_through_7_4.tex
\SGASection{7}{Relative root datum}
\SGARefTarget{sga3:xxvi:sec:7}

Throughout this section, \(S\) denotes a nonempty connected
semilocal scheme, \(G\) an \(S\)-reductive group, \(Q\) a maximal
split subtorus of \(G\), and \(L\) the centralizer of \(Q\) in \(G\):
\[
 L=\Centr_G(Q).
\]

\SGASubsection{7.1}{}
\SGARefTarget{sga3:xxvi:subsec:7:1}

Since \(Q\) is the largest central split subtorus of \(L\), every
element of \(G(S)\) which normalizes \(L\) also normalizes \(Q\).
Thus \textup{(cf.\ \SGARefLink{sga3:xxvi:subsec:7:1:1}{7.1.1})}
\[
 \Norm_G(L)(S)=\Norm_G(Q)(S).
\]
On the other hand, \SGARefLink{sga3:xxvi:corollary:6-4}{6.4} shows that
\[
 G(S)\longrightarrow(G/L)(S)
\]
is surjective.  Consequently there is a canonical identification
\[
 W_G(Q)(S)
 =\bigl(\Norm_G(Q)/\Centr_G(Q)\bigr)(S)
 \simeq\Norm_G(L)(S)/L(S).
\]

Let
\[
 M=\operatorname{Hom}_{S\text{-}\mathrm{gr}}
      (Q,\mathbf G_{m,S}),
\]
so that \(Q\simeq D_S(M)\) canonically.  Denote by \(W\) the group
of automorphisms of \(M\) defined by \(W_G(Q)(S)\).  Thus
\[
 W\simeq W_G(Q)(S)\simeq\Norm_G(L)(S)/L(S).
\]

\SGASubsection{7.1.1}{}
\SGARefTarget{sga3:xxvi:subsec:7:1:1}

In general one does not have
\(\Norm_G(L)=\Norm_G(Q)\).  For example, let \(S\) be the spectrum of
a field \(k\) having a quadratic extension \(k'\), and take for
\(G\) the unitary group
\[
 {}^2\!A_{2,k'/k}
\]
\textup{(cf.\ \SGARefLink{sga3:xxiv:3.11.2}{Exposé XXIV, 3.11.2})}.  The minimal parabolic
subgroups of \(G\) are its Borel subgroups, and their Levi subgroups
are maximal tori; hence
\[
 \bigl(\Norm_G(L)/L\bigr)(k)=\mathfrak S_3.
\]
Since \(G\) is not split, its maximal split tori have dimension at
most \(1\), and are therefore isomorphic to \(\mathbf G_{m,k}\).
The group \(\Norm_G(Q)/L\) acts faithfully on \(Q\), and
\[
 \bigl(\Norm_G(Q)/L\bigr)(k)=\mathbf Z/2\mathbf Z.
\]
\footnote{\emph{Editor's note.} The French edition prints
\(\Norm_G(Q)/Q\) in this final equality, although the immediately
preceding action and the definition of the relative Weyl group use
\(\Norm_G(Q)/L\).}

\SGASubsection{7.2}{}
\SGARefTarget{sga3:xxvi:subsec:7:2}

Let \(P\) be a parabolic subgroup of \(G\) with Levi subgroup \(L\);
such a subgroup exists by \SGARefLink{sga3:xxvi:cor:6:2}{6.2}.  It is necessarily minimal, by \SGARefLink{sga3:xxvi:rem:6:18}{6.18}.
Using the conjugacy of the minimal parabolic subgroups of \(G\)
\textup{(\SGARefLink{sga3:xxvi:corollary:5-7}{5.7})}, the conjugacy of the Levi subgroups of a parabolic
subgroup \textup{(\SGARefLink{sga3:xxvi:cor:1:8}{1.8})}, and the equalities
\[
 P=\Norm_G(P),\qquad \Norm_G(L)\cap P=L
\]
from \SGARefLink{sga3:xxvi:proposition:1-6}{1.6}, we obtain the following conclusion: the set of
\textup{(minimal)} parabolic subgroups of \(G\) having Levi subgroup
\(L\) is a principal homogeneous space under \(W\).

\SGASubsection{7.3}{}
\SGARefTarget{sga3:xxvi:subsec:7:3}

The Lie algebra of \(G\) decomposes under the action of \(Q\) as
\[
 \operatorname{Lie}(G)
 =\operatorname{Lie}(L)
  \oplus\bigoplus_{\alpha\in R}\operatorname{Lie}(G)^\alpha,
\]
where \(R\) is the set of nonzero characters \(\alpha\) of \(Q\)
such that \(\operatorname{Lie}(G)^\alpha\ne0\).  These are the roots
of \(G\) relative to \(Q\).

Let
\[
 M^*=\operatorname{Hom}_{S\text{-}\mathrm{gr}}
       (\mathbf G_{m,S},Q).
\]
It is in duality with \(M\), and \(W\) acts on it naturally by
transport of structure.

\begin{sgastatement}{\SGAInlineRefTarget{sga3:xxvi:theorem:7-4}Theorem 7.4}
With the notation of \SGARefLink{sga3:xxvi:subsec:7:3}{7.3}, there is a unique map
\[
 R\longrightarrow M^*,\qquad \alpha\longmapsto\alpha^*,
\]
which makes \((M,M^*)\) into a root datum \textup{(\SGARefLink{sga3:xxi:subsection:1-1}{Exposé XXI,
1.1})} with Weyl group \(W\).

Moreover, parabolic subgroups \(P\) of \(G\) having Levi subgroup
\(L\) correspond bijectively to systems of positive roots \(R^+\)
in \(R\) by the relation
\[
 \operatorname{Lie}(P)
 =\operatorname{Lie}(L)
  \oplus\bigoplus_{\alpha\in R^+}\operatorname{Lie}(G)^\alpha.
\]
\end{sgastatement}
\SGARefTarget{sga3:xxvi:thm:7:4}

\SGASubsection{7.4.1}{}
\SGARefTarget{sga3:xxvi:subsec:7:4:1}

Suppose first that the existence of the desired map
\(\alpha\mapsto\alpha^*\) has been proved.  By \SGARefLink{sga3:XXI:3.4.10}{Exposé XXI, 3.4.10},
\(s_\alpha\) is the unique element of \(W\) such that
\[
 s_\alpha(m)-m
\]
is a rational multiple of \(\alpha\) for every \(m\in M\).
Consequently \(s_\alpha\) is determined by \(\alpha\).  Since
\[
 \alpha^*(m)\alpha=m-s_\alpha(m),
\]
\footnote{\emph{Editor's note.} Here and again in \SGARefLink{sga3:xxvi:subsec:7:4:4}{7.4.4}, the
Polo--Gille edition restores the factor \(\alpha\) omitted after
\(\alpha^*(m)\) in the original text.}
the coroot \(\alpha^*\) is also determined by \(\alpha\), proving
the uniqueness of the map.

\SGASubsection{7.4.2}{}
\SGARefTarget{sga3:xxvi:subsec:7:4:2}

Let \(\alpha\in R\).  Denote by \(L_\alpha\), \(H_\alpha\), and
\(H_{-\alpha}\), respectively, the unique smooth subgroups of \(G\)
with connected fibers which contain \(L\) and satisfy
\textup{(cf.\ \SGARefLink{sga3:xxvi:proposition:6-1}{6.1}\textup{(i)})}
\[
\begin{aligned}
 \operatorname{Lie}(L_\alpha)
  &=\operatorname{Lie}(L)
    \oplus\bigoplus_{\gamma\in\mathbf Z\alpha\cap R}
       \operatorname{Lie}(G)^\gamma,\\
 \operatorname{Lie}(H_\alpha)
  &=\operatorname{Lie}(L)
    \oplus\bigoplus_{\gamma\in\mathbf N\alpha\cap R}
       \operatorname{Lie}(G)^\gamma,\\
 \operatorname{Lie}(H_{-\alpha})
  &=\operatorname{Lie}(L)
    \oplus\bigoplus_{\gamma\in-\mathbf N\alpha\cap R}
       \operatorname{Lie}(G)^\gamma.
\end{aligned}
\]
The group \(L_\alpha\) is a reductive subgroup of \(G\), and
\(H_\alpha,H_{-\alpha}\) are opposite parabolic subgroups of
\(L_\alpha\), with common Levi subgroup \(L\), by
\SGARefLink{sga3:xxvi:proposition:6-1}{6.1}\textup{(ii)}.

By \SGARefLink{sga3:xxvi:subsec:7:2}{7.2} there is therefore an element
\[
 s_\alpha\in\Norm_{L_\alpha}(Q)(S)/L(S)\subset W
\]
such that \(s_\alpha(H_\alpha)=H_{-\alpha}\).  One has
\[
 s_\alpha(\alpha)=-\alpha,
\]
because \(\alpha\), respectively \(-\alpha\), is the common divisor
of the elements of \(R\) occurring in \(H_\alpha\), respectively
\(H_{-\alpha}\).  Moreover \(s_\alpha^2=\mathrm{id}\), because
\(s_\alpha^2(H_\alpha)\) and \(H_\alpha\) are both opposite to
\(H_{-\alpha}\) relative to \(L\).  We have thus constructed an
element \(s_\alpha\in W\) with the properties
\[
\tag{x}\label{sga3:xxvi:subsection:7-4-2:equation:x}
 s_\alpha(\alpha)=-\alpha,\qquad s_\alpha^2=\mathrm{id},
\]
and
\[
\tag{xx}\label{sga3:xxvi:subsection:7-4-2:equation:xx}
 s_\alpha\text{ can be represented by an element of }L_\alpha(S).
\]
The construction is canonical in \(\alpha\); in particular,
\[
\tag{xxx}\label{sga3:xxvi:subsection:7-4-2:equation:xxx}
 w s_\alpha w^{-1}=s_{w(\alpha)}
 \qquad\text{for every }w\in W.
\]

\SGASubsection{7.4.3}{}
\SGARefTarget{sga3:xxvi:subsec:7:4:3}

We now prove
\[
\tag{xxxx}\label{sga3:xxvi:subsection:7-4-3:equation:xxxx}
 s_\alpha(m)-m\in\mathbf Z\alpha
 \qquad\text{for every }m\in M.
\]
Since \(S\) is connected, this assertion is local for the fpqc
topology.  We may therefore suppose that
\[
 L_\alpha=G_1
\]
is split relative to a maximal torus \(T_1\) of \(L\).  Let
\((G_1,T_1,M_1,R_1)\) be such a splitting.  The monomorphism
\(Q\to T_1\) identifies \(M\) with a quotient of \(M_1\); denote the
canonical map by
\[
 p:M_1\longrightarrow M.
\]

The image of \(R_1\) under \(p\), possibly including zero, consists
of the roots of \(G_1\) relative to \(Q\), hence of the elements of
\(R\) which are integral multiples of \(\alpha\).  Thus
\[
 p(R_1)\subset\mathbf Z\alpha.
\]
By \textup{\SGARefLink{sga3:xxvi:subsection:7-4-2:equation:xx}{(xx)}}, an element of \(\Norm_{G_1}(L)(S)\) induces
\(s_\alpha\) on \(Q\).  \SGARefLink{sga3:xxii:rem:5.10.10}{Exposé XXII, 5.10.10}, then gives a section
\[
 w\in(W_1)_S(S)
\]
which induces \(s_\alpha\) on \(Q\), where \(W_1\) is the Weyl group
of the root datum \((M_1,R_1,\ldots)\).  After restricting \(S\) if
necessary, we may suppose that \(w\in W_1\) and that
\[
 p\bigl(w(m_1)\bigr)=s_\alpha\bigl(p(m_1)\bigr)
 \qquad(m_1\in M_1).
\]
By definition, \(w\) is a product of reflections associated with
elements of \(R_1\).  Hence \(w(m_1)-m_1\) is an integral linear
combination of elements of \(R_1\).  It follows that
\[
 s_\alpha\bigl(p(m_1)\bigr)-p(m_1)
\]
is an integral linear combination of elements of
\(p(R_1)\subset\mathbf Z\alpha\), and is therefore an integral
multiple of \(\alpha\), as required.

\SGASubsection{7.4.4}{}
\SGARefTarget{sga3:xxvi:subsec:7:4:4}

We may now define \(\alpha^*\in M^*\) by
\[
 \alpha^*(m)\alpha=m-s_\alpha(m).
\]
By \textup{(x)},
\[
 \langle\alpha^*,\alpha\rangle=2.
\]
Moreover, \textup{\SGARefLink{sga3:xxvi:subsection:7-4-2:equation:xxx}{(xxx)}} implies that for every
\((\alpha,\alpha')\in R\times R\),
\[
 s_\alpha(\alpha')\in R,\qquad
 s_\alpha(\alpha'^*)=s_\alpha(\alpha')^*.
\]
By \SGARefLink{sga3:xxi:subsection:1-1}{Exposé XXI, 1.1}, the map
\(\alpha\mapsto\alpha^*\) therefore defines a root datum on
\((M,M^*)\).

\SGASubsection{7.4.5}{}
\SGARefTarget{sga3:xxvi:subsec:7:4:5}

Let \(W'\) be the Weyl group of this root datum, that is, the group
of transformations of \(M\) generated by the \(s_\alpha\).  Then
\(W'\subset W\).

Choose a total order on the free abelian group \(M\), and put
\[
 R^+=\{\alpha\in R\mid\alpha\geq0\}.
\]
This is a system of positive roots in \(R\).  Let \(w\in W\), and
represent it by an element
\[
 n\in\Norm_G(L)(S)=\Norm_G(Q)(S).
\]
Put \(P=H_{R^+}\), in the notation of \SGARefLink{sga3:xxvi:proposition:6-1}{6.1}.  Then \(P\) is a
parabolic subgroup of \(G\) with Levi subgroup \(L\), and
\[
 \Int(n)P=H_{w(R^+)}.
\]
It follows from \SGARefLink{sga3:xxvi:subsec:7:3}{7.3} that \(w(R^+)=R^+\) implies \(w=e\).  The group
\(W'\) acts transitively on the systems of positive roots in \(R\),
by \SGARefLink{sga3:XXI:3.3.7}{Exposé XXI, 3.3.7}, while the stabilizer of \(R^+\) in \(W\) is
trivial.  Therefore \(W=W'\).

It follows as well that \(W=W'\) acts simply transitively both on
the systems of positive roots in \(R\) and on the parabolic
subgroups of \(G\) having Levi subgroup \(L\).  This proves the last
assertion of \SGARefLink{sga3:xxvi:theorem:7-4}{7.4}.
\SGAQED

% END INCLUDED SGA3_English_current_buildable_source_20260823/source/inputs/XXVI/tex/components/13_section7_opening_through_7_4.tex

% BEGIN INCLUDED SGA3_English_current_buildable_source_20260823/source/inputs/XXVI/tex/components/14_section7_5_through_7_9.tex
\SGASubsection{7.5}{}
\SGARefTarget{sga3:xxvi:subsec:7:5}

Let \(P,P_1\) be two minimal parabolic subgroups of \(G\), with
Levi subgroups \(L,L_1\), and let \(Q,Q_1\) be the largest central
split subtori of \(L,L_1\), respectively.  Then the pairs
\((P,Q)\) and \((P_1,Q_1)\) are conjugate: there is a \(g\in G(S)\)
such that
\[
 \Int(g)P=P_1,\qquad \Int(g)Q=Q_1.
\]
Indeed, \(P\) and \(P_1\) are conjugate by \SGARefLink{sga3:xxvi:corollary:5-7}{5.7}, so we may suppose
\(P=P_1\).  Then \(L\) and \(L_1\) are conjugate by an element of
\(P(S)\), by \SGARefLink{sga3:xxvi:cor:1:8}{1.8}.  Moreover, if \(g,g'\) are two elements of \(G(S)\)
which conjugate \((P,Q)\) to \((P_1,Q_1)\), then \(g'^{-1}g\)
normalizes \(P\) and \(Q\), hence also \(P\) and \(L\).  But
\[
 \Norm_G(P)\cap\Norm_G(L)
 =P\cap\Norm_G(L)
 =L
 =\Centr_G(Q).
\]
Consequently the isomorphism
\[
 Q\xrightarrow{\ \sim\ }Q_1
\]
induced by \(\Int(g)\) is independent of \(g\).

Let \(\mathcal R,\mathcal R_1\) be the root data defined by \SGARefLink{sga3:xxvi:theorem:7-4}{7.4} in
\[
 \operatorname{Hom}_{S\text{-}\mathrm{gr}}(Q,\mathbf G_{m,S})
 \quad\text{and}\quad
 \operatorname{Hom}_{S\text{-}\mathrm{gr}}(Q_1,\mathbf G_{m,S}),
\]
and let \(R^+,R_1^+\) be the systems of positive roots corresponding
to \(P,P_1\).  The canonical isomorphism \(Q\simeq Q_1\) just
constructed carries \((\mathcal R,R^+)\) to
\((\mathcal R_1,R_1^+)\).  We may therefore define the
\emph{based relative root datum} of \(G\) over \(S\), by identifying
the various pairs \((\mathcal R,R^+)\) through the transitive system
of isomorphisms above.\footnote{A based root datum is a root datum
together with the choice of a system of positive roots, equivalently
of simple roots; cf.\ \SGARefLink{sga3:XXIII:1.5}{Exposé XXIII, 1.5}.}

From now on, denote this based root datum by
\[
 \mathcal R(G/S)
 =(\underline M,\underline M^*,\underline R,
   \underline R^*,\underline R^+).
\]
For every pair \(P\supset Q\) as above there is therefore a
canonical isomorphism
\[
 \underline M\xrightarrow{\ \sim\ }
 \operatorname{Hom}_{S\text{-}\mathrm{gr}}(Q,\mathbf G_{m,S})
\]
which carries \(\underline R\), respectively \(\underline R^+\), to
the roots of \(G\), respectively of \(P\), relative to \(Q\), and
carries \(W(\mathcal R)\) to \(W_G(Q)(S)\).

\SGASubsection{7.6}{}
\SGARefTarget{sga3:xxvi:subsec:7:6}

Let again \(Q\) be a maximal split torus of \(G\), let \(P\) be a
minimal parabolic subgroup of \(G\) with Levi subgroup
\(\Centr_G(Q)\), let
\[
 (M,M^*,R,R^*,R^+)
\]
be the corresponding based root datum of \SGARefLink{sga3:xxvi:theorem:7-4}{7.4}, and let \(\Delta\)
be the set of simple roots in \(R^+\).  For every
\(A\subset\Delta\), let
\[
 R_A=R^+\cup(\mathbf Z A\cap R^-).
\]
Thus \(R_A\) consists of the positive roots together with the
negative roots which are integral linear combinations of elements
of \(A\).  It is a closed set of roots, by \SGARefLink{sga3:XXI:3.1.4}{Exposé XXI, 3.1.4}, and
every closed set containing \(R^+\) is uniquely of this form, by
\SGARefLink{sga3:XXI:3.3.10}{Exposé XXI, 3.3.10}.

By \SGARefLink{sga3:xxvi:proposition:6-1}{6.1} there is a unique smooth subgroup \(P_A\) of \(G\), with
connected fibers, which contains \(\Centr_G(Q)\) and satisfies
\[
 \operatorname{Lie}(P_A)
 =\operatorname{Lie}(G)^0
  \oplus\bigoplus_{\alpha\in R_A}\operatorname{Lie}(G)^\alpha.
\]
The following result now follows immediately from \SGARefLink{sga3:xxvi:proposition:6-1}{6.1}, the
conjugacy of minimal parabolic subgroups, and the fact that the set
of roots of a parabolic subgroup of \(G\) containing \(Q\) is
closed \textup{(which follows from \SGARefLink{sga3:xxvi:prop:1:3}{1.4} after splitting)}.

\begin{sgastatement}{\SGAInlineRefTarget{sga3:xxvi:proposition:7-7}Proposition 7.7}
\begin{enumerate}[label=\textup{(\roman*)}]
\item The map
\[
 A\longmapsto P_A
\]
is an inclusion-preserving bijection from the set of subsets of
\(\Delta\) onto the set of parabolic subgroups of \(G\) containing
\(P\).
\item Every parabolic subgroup of \(G\) is conjugate by an element
of \(G(S)\) to a unique \(P_A\).
\end{enumerate}
\end{sgastatement}
\SGARefTarget{sga3:xxvi:prop:7:7}

\SGASubsection{7.8}{}
\SGARefTarget{sga3:xxvi:subsec:7:8}

Let \(P\supset Q\) be as above.  Consider the relative root datum of
\(G\) over \(S\) defined in \SGARefLink{sga3:xxvi:subsec:7:5}{7.5} and the canonical isomorphism
\[
 f:(M,M^*,R,R^*,R^+)
   \xrightarrow{\ \sim\ }
   (\underline M,\underline M^*,\underline R,
    \underline R^*,\underline R^+).
\]
The set \(\Delta\) of simple roots in \(R^+\) is carried to the set
\(\underline\Delta\) of simple roots in \(\underline R^+\).
Thus every subset \(A\subset\Delta\) is carried to a subset
\[
 f(A)\subset\underline\Delta.
\]

\begin{sgastatement}{Definition 7.9.0}
Let \(H\) be any parabolic subgroup of \(G\).  By
\SGARefLink{sga3:xxvi:proposition:7-7}{7.7}\textup{(ii)}, it is conjugate to a unique \(P_A\).  Put
\[
 t_r(H)=f(A)\subset\underline\Delta.
\]
The conjugacy theorems show immediately that \(t_r(H)\) is
independent of the choice of the pair \(P\supset Q\).  It is called
the \emph{relative type} of \(H\).
\end{sgastatement}
\SGARefTarget{sga3:xxvi:def:7:9:0}

\begin{sgastatement}{\SGAInlineRefTarget{sga3:xxvi:proposition:7-9}Proposition 7.9}
\begin{enumerate}[label=\textup{(\roman*)}]
\item The map \(H\mapsto t_r(H)\) induces a bijection from the set of
\(G(S)\)-conjugacy classes of parabolic subgroups of \(G\) onto the
set of subsets of \(\underline\Delta\).
\item Let \(H\) be a parabolic subgroup of \(G\), \(P\) a minimal
parabolic subgroup contained in \(H\), \(Q\) the largest central
split torus of a Levi subgroup of \(P\), \(\Delta\) the set of
simple roots of \(P\) relative to \(Q\), and
\[
 f:\Delta\xrightarrow{\ \sim\ }\underline\Delta
\]
the canonical isomorphism.  For every \(\alpha\in\Delta\),
\[
 f(\alpha)\in t_r(H)
 \quad\Longleftrightarrow\quad
 \operatorname{Lie}(H)^{-\alpha}\ne0,
\]
and
\[
 H=P_A,\qquad A=f^{-1}\bigl(t_r(H)\bigr).
\]
\item If \(H,H'\) are two parabolic subgroups of \(G\) containing
\(P\), then \textup{(see \SGARefLink{sga3:xxvi:rem:3:8:1}{3.8.1}\textup{(ii)} and
\SGARefLink{sga3:xxvi:corollary:5-5}{5.5}\textup{(i)} for other equivalent conditions)}
\[
 t_r(H)\subset t_r(H')
 \quad\Longleftrightarrow\quad
 t(H)\subset t(H').
\]
\end{enumerate}
\end{sgastatement}
\SGARefTarget{sga3:xxvi:prop:7:9}

% END INCLUDED SGA3_English_current_buildable_source_20260823/source/inputs/XXVI/tex/components/14_section7_5_through_7_9.tex

% BEGIN INCLUDED SGA3_English_current_buildable_source_20260823/source/inputs/XXVI/tex/components/15_section7_10_and_7_11.tex
\SGASubsection{7.10}{}
\SGARefTarget{sga3:xxvi:subsec:7:10}

The relative positions of two minimal parabolic subgroups can be
described as follows.  We use the notation \(t_2(P,P_1)\) of 4.5.2.

\begin{enumerate}[label=\textup{(\arabic*)}]
\item Let \(P,P_1,P',P'_1\) be four minimal parabolic subgroups of
\(G\).  Then
\[
 t_2(P,P_1)=t_2(P',P'_1),
\]
that is, the pairs \((P,P_1)\) and \((P',P'_1)\) are fpqc-locally
conjugate, if and only if there is a \(g\in G(S)\) such that
\[
 \Int(g)P=P',\qquad \Int(g)P_1=P'_1.
\]

\item Fix a minimal parabolic subgroup \(P\) of \(G\) with Levi
subgroup \(L\), and let \(T\) be a maximal torus of \(L\).  Consider
the scheme
\[
 \Par_{t_{\min}}(G;P)
\]
of minimal parabolic subgroups of \(G\) in standard position
relative to \(P\).  There is a morphism \textup{(cf.\ \SGARefLink{sga3:xxvi:paragraph:4-5-5}{4.5.5})}
\[
 f:\Par_{t_{\min}}(G;P)
   \longrightarrow W_P(T)\backslash W_G(T)/W_P(T)
\]
whose fibers are the orbits of \(P\) in
\(\Par_{t_{\min}}(G;P)\).  By \textup{(1)}, \(f\) therefore induces
a monomorphism
\[
 P(S)\backslash\Par_{t_{\min}}(G;P)(S)
 \lhook\joinrel\longrightarrow
 \bigl(W_P(T)\backslash W_G(T)/W_P(T)\bigr)(S).
\]
Its image is identified with \(W\).  This is the Bruhat theorem:
each orbit of \(P(S)\) in \(\Par_{t_{\min}}(G;P)(S)\) contains
exactly one parabolic subgroup of \(G\) with Levi subgroup \(L\),
that is, one of the form
\[
 \Int(n)P,\qquad n\in\Norm_G(L)(S).
\]

\item Equivalently, let \(E(S)\) be the set of elements \(g\in G(S)\)
such that \(\Int(g)P\) and \(P\) are in mutual standard position.
Then \(E(S)\) is partitioned into double cosets modulo \(P(S)\),
indexed by \(W\):
\[
 E(S)=P(S)\,W\,P(S).
\]
Putting \(U=\rad^u(P)\), one may also write
\[
 E(S)=U(S)\cdot\Norm_G(L)(S)\cdot U(S).
\]
This exhibits a partition of \(E(S)\) into double cosets modulo
\(U(S)\), indexed by \(\Norm_G(L)(S)\).

\item If \(S\) is the spectrum of a field, then \(E(S)=G(S)\), and
one recovers \textup{[BT65]}, 5.15.
\end{enumerate}

\begin{sgastatement}{Counter\SGARefLink{sga3:xxvi:counterexamples:7:11}{examples 7.11}}
Let
\[
 S=\operatorname{Spec}(\mathbf Z/4\mathbf Z),
 \qquad G=\operatorname{SL}_{2,S}.
\]
Let \(B\) be the usual Borel subgroup consisting of the matrices
\[
 \begin{pmatrix}a&b\\c&d\end{pmatrix}
 \qquad\text{with }c=0,
\]
and put
\[
 g=\begin{pmatrix}-1&1\\2&1\end{pmatrix}\in G(S),
 \qquad B'=\Int(g)B.
\]
Then
\[
 B(S)=B'(S),
\]
whereas \(B\cap B'\) contains no maximal torus.
\footnote{Indeed, \(B'\) is the subgroup of \(G\) defined by
\(c=2(a+b)\).  Thus \(B\cap B'\) is not flat over \(S\), and hence
contains no maximal torus, by \SGARefLink{sga3:xxvi:proposition:4-5-1}{4.5.1}.}

This shows both that two distinct minimal parabolic subgroups may
have the same group of sections and that, in general, one cannot
decide from the groups \(P(S)\) and \(P'(S)\) alone whether two
minimal parabolic subgroups \(P,P'\) are in standard position.  In
particular, the subset \(E(S)\) of \(G(S)\) does not appear to be
definable solely from the configuration
\[
 \{G(S),P(S),\Norm_G(L)(S)\}.
\]
In the example above, this subset is defined by \(c\ne2\).
\footnote{More generally, for every \(S\)-scheme \(S'\), the set
\(E(S')\) is defined by the condition that \(c\) is either zero or
invertible.}
\end{sgastatement}
\SGARefTarget{sga3:xxvi:counterexamples:7:11}

% END INCLUDED SGA3_English_current_buildable_source_20260823/source/inputs/XXVI/tex/components/15_section7_10_and_7_11.tex

% BEGIN INCLUDED SGA3_English_current_buildable_source_20260823/source/inputs/XXVI/tex/components/16_section7_12_through_bibliography.tex
\SGASubsection{7.12}{}
\SGARefTarget{sga3:xxvi:subsec:7:12}

We now study how \(\mathcal R(G/S)\) varies with \(S\).  Let \(S'\)
be an \(S\)-scheme which is likewise nonempty, connected, and
semilocal.  If \(Q\) is a maximal split torus of \(G\), then
\(Q_{S'}\) is a split torus of \(G_{S'}\).  Choose a maximal split
torus \(Q'\) of \(G_{S'}\) containing \(Q_{S'}\), and put
\[
\begin{aligned}
 M&=\operatorname{Hom}_{S\text{-}\mathrm{gr}}
       (Q,\mathbf G_{m,S})
   \simeq
   \operatorname{Hom}_{S'\text{-}\mathrm{gr}}
       (Q_{S'},\mathbf G_{m,S'}),\\
 M'&=\operatorname{Hom}_{S'\text{-}\mathrm{gr}}
       (Q',\mathbf G_{m,S'}).
\end{aligned}
\]
The monomorphism \(Q_{S'}\to Q'\) induces an epimorphism
\[
 u:M'\longrightarrow M.
\]
Write
\[
 L=\Centr_G(Q),\qquad L'=\Centr_{G_{S'}}(Q').
\]
Then \(L'\subset L_{S'}\).

If \(H\) is a subgroup of \(G\) containing \(L\), then \(H_{S'}\)
contains \(L'\), and
\[
\begin{aligned}
 \operatorname{Lie}(H)
  &=\operatorname{Lie}(L)
    \oplus\bigoplus_{\alpha\in R_H}
       \operatorname{Lie}(H)^\alpha,\\
 \operatorname{Lie}(H_{S'})
  &=\operatorname{Lie}(L')
    \oplus\bigoplus_{\alpha'\in R'_{H_{S'}}}
       \operatorname{Lie}(H_{S'})^{\alpha'}.
\end{aligned}
\]
Here \(R_H\), respectively \(R'_{H_{S'}}\), denotes the set of
roots of \(H\), respectively \(H_{S'}\), relative to \(Q\),
respectively \(Q'\).  It follows immediately that
\[
 R_H\subset u\bigl(R'_{H_{S'}}\bigr)
       \subset R_H\cup\{0\}.
\]

Taking \(H=G\), we first obtain
\[
 R\subset u(R')\subset R\cup\{0\}.
\]
Next take for \(H\) a minimal parabolic subgroup \(P\) with Levi
subgroup \(L\).  Then \(R'_{H_{S'}}\) contains a system of positive
roots in \(R'\).  By \SGARefLink{sga3:xxvi:theorem:7-4}{7.4}, there is therefore a minimal parabolic
subgroup \(P'\) of \(G_{S'}\), with Levi subgroup \(L'\), contained
in \(P_{S'}\).  We have constructed the diagram of inclusions
\[
\begin{array}{ccccc}
 Q_{S'} & \subset & L_{S'} & \subset & P_{S'}\\[1.2ex]
 \cap   &         & \cup   &         & \cup\\[1.2ex]
 Q'     & \subset & L'     & \subset & P'.
\end{array}
\]

If \(R^+\), respectively \(\Delta\), is the system of positive,
respectively simple, roots in \(R\) defined by \(P\), and if
\(R'^+,\Delta'\) are defined similarly, then
\[
 R^+\subset u(R'^+)\subset R^+\cup\{0\},
 \qquad
 \Delta\subset u(\Delta')\subset\Delta\cup\{0\}.
\]

Now let
\[
 w\in W\simeq\Norm_G(Q)(S)/\Centr_G(Q)(S),
\]
and choose a representative \(n\in\Norm_G(Q)(S)\).  Since
\(\Int(n)Q=Q\), one has \(\Int(n)L=L\), and hence
\(\Int(n)L_{S'}=L_{S'}\).  The tori \(Q'\) and \(\Int(n)Q'\) are
therefore two maximal split tori of \(L_{S'}\), so they are
conjugate by an element \(x\in L(S')\).  We may choose \(x\) so that
\[
 \Int(nx)Q'=Q'.
\]
Thus \(n'=nx\) belongs to \(\Norm_{G_{S'}}(Q')(S')\).  Let \(w'\)
be the image of \(n'\) in
\[
 W'\simeq
 \Norm_{G_{S'}}(Q')(S')/\Centr_{G_{S'}}(Q')(S').
\]
The action of \(w'\) on \(M'\) is compatible with
\(u:M'\to M\), and the induced action on \(M\) is the action of
\(w\).

Using the definition of relative root data together with the
conjugacy theorems gives the following result.

\begin{sgastatement}{\SGAInlineRefTarget{sga3:xxvi:theorem:7-13}Theorem 7.13}
Let \(S,S'\) be nonempty connected semilocal schemes, let
\(S'\to S\) be a morphism of schemes, let \(G\) be an
\(S\)-reductive group, and write
\[
\begin{aligned}
 \mathcal R(G/S)
  &=(\underline M,\underline M^*,\underline R,
     \underline R^*,\underline R^+),\\
 \mathcal R(G_{S'}/S')
  &=(\underline M',\underline M'^*,\underline R',
     \underline R'^*,\underline R'^+)
\end{aligned}
\]
for the based relative root data.  There is a canonical
homomorphism
\[
 u:\underline M'\longrightarrow\underline M
\]
with the following properties:
\begin{enumerate}[label=\textup{(\roman*)}]
\item The homomorphism \(u\) is surjective.
\item For every \(w\in W\), there is an element \(w'\in W'\)
compatible with \(u\) and inducing \(w\) on \(\underline M\).
\item For every subset \(X\subset\underline M\), put
\[
 X^\wedge=X\cap(\underline M\setminus\{0\}).
\]
Then
\[
 u(\underline R'^+)^\wedge=\underline R^+,
 \qquad
 u(\underline\Delta')^\wedge=\underline\Delta.
\]
\item For every parabolic subgroup \(H\), consider
\[
 t_r(H)\subset\underline\Delta,
 \qquad
 t_r(H_{S'})\subset\underline\Delta'.
\]
Then
\[
\begin{aligned}
 t_r(H_{S'})
 &=u^{-1}\bigl(t_r(H)\cup\{0\}\bigr)
       \cap\underline\Delta'\\
 &=\left\{\alpha'\in\underline\Delta'
       \ \middle|\ 
       u(\alpha')\in t_r(H)\text{ or }u(\alpha')=0
   \right\}.
\end{aligned}
\]
\end{enumerate}
\end{sgastatement}
\SGARefTarget{sga3:xxvi:thm:7:13}

\begin{sgastatement}{Remark 7.14}
If \(G\) is splittable over \(S\), its maximal split tori are
maximal tori, and the relative notions introduced here coincide
with the previously defined absolute notions.  The preceding
theorem therefore describes the relative root datum
\(\mathcal R(G/S)\) and the relative type \(t_r\) in terms of the
absolute root datum and absolute type of \(G_{S'}\), where \(S'\) is
chosen so that \(G_{S'}\) is splittable
\textup{(cf.\ \SGARefLink{sga3:xxiv:theorem:4-4-1}{Exposé XXIV, 4.4.1})}.  See \textup{[BT65]}, \SGARefLink{sga3:xxvi:cor:6:11}{6.12}
and following, for this description.
\end{sgastatement}
\SGARefTarget{sga3:xxvi:rem:7:14}

\SGASubsection{7.15}{}
\SGARefTarget{sga3:xxvi:subsec:7:15}

Let \(S\) be a Henselian local scheme, \(s_0\) its closed point, and
\(S_0\) the spectrum of the residue field at \(s_0\), identified
with a closed subscheme of \(S\).  For every object \(X\) over
\(S\), write \(X_0\) for its base change to \(S_0\).  Let \(G\) be
an \(S\)-reductive group.

For every parabolic subgroup \(\overline P\) of \(G\), the group
\(\overline P_0\) is a parabolic subgroup of \(G_0\).  Conversely,
for every parabolic subgroup \(P\) of \(G_0\), there is a parabolic
subgroup \(\overline P\) of \(G\) such that
\[
 \overline P_0=P.
\]
This follows from Hensel's lemma and from the fact that \(\Par(G)\)
is a smooth \(S\)-scheme.  In particular, by \SGARefLink{sga3:xxvi:corollary:5-7}{5.7},
\(\overline P\) is minimal if and only if \(\overline P_0\) is
minimal.

Once such a subgroup \(\overline P\) has been chosen, the same
argument shows that every maximal split subtorus of
\(\overline P_0\) is of the form \(T_0\), where \(T\) is a maximal
split subtorus of \(\overline P\).  It follows readily that the
relative root data of \(G\) over \(S\) and of \(G_0\) over \(S_0\)
are canonically isomorphic.  Thus the theory of parabolic subgroups
of \(G\) reduces to the corresponding theory for \(G_0\).

Conversely, every \(S_0\)-reductive group is of the form \(G_0\), by
\SGARefLink{sga3:xxiv:1.21}{Exposé XXIV, Proposition 1.21}.  This makes it possible to reduce the
study of the parabolic subgroups of an \(S_0\)-reductive group to
the corresponding study over \(S\).

\chapter*{Bibliography}
\addcontentsline{toc}{chapter}{Bibliography}
\SGARefTarget{sga3:xxvi:bibliography}

\begin{description}[leftmargin=4.5em,style=nextline]
\item[\textup{[BT65]}]
A.~Borel and J.~Tits, \emph{Groupes réductifs},
Publ.\ Math.\ I.H.É.S.\ \textbf{27} \textup{(1965)}, 55--150.
\footnote{\emph{Editor's note.} This reference appeared in the
original; the references which follow were added in the
Polo--Gille edition.}

\item[\textup{[BT71]}]
A.~Borel and J.~Tits,\\
\emph{Éléments unipotents et sous-groupes
paraboliques de groupes réductifs. I},
Invent.\ Math.\ \textbf{12} \textup{(1971)}, 95--104.

\item[\textup{[Ch05]}]
C.~Chevalley, \emph{Classification des groupes algébriques
semi-simples}, with the collaboration of P.~Cartier,
A.~Grothendieck, and M.~Lazard, Collected Works, vol.~3,
Springer, 2005.

\item[\textup{[DG70]}]
M.~Demazure and P.~Gabriel, \emph{Groupes algébriques},
Masson \& North-Holland, 1970.

\item[\textup{[Gi71]}]
J.~Giraud, \emph{Cohomologie non abélienne},
Springer-Verlag, 1971.
\end{description}

% END INCLUDED SGA3_English_current_buildable_source_20260823/source/inputs/XXVI/tex/components/16_section7_12_through_bibliography.tex



% END INCLUDED SGA3_English_current_buildable_source_20260823/source/inputs/XXVI/body.tex


\clearpage
\SGAInputRoot{TERMINAL_INDEX}
% BEGIN INCLUDED SGA3_English_current_buildable_source_20260823/source/inputs/TERMINAL_INDEX/body.tex
% Exact R15 special component.
% BEGIN INCLUDED SGA3_English_current_buildable_source_20260823/source/inputs/TERMINAL_INDEX/index/tex/SGA3_TomeIII_Index_body.tex
\chapter*{Index --- Tome III}
\addcontentsline{toc}{chapter}{Index --- Tome III}
\noindent\emph{Page numbers refer to the Polo--Gille re-edition.}
\medskip
\begin{multicols}{2}
\small
\setlength{\parskip}{1.2pt}
\interlinepenalty=10000
% BEGIN INCLUDED SGA3_English_current_buildable_source_20260823/source/inputs/TERMINAL_INDEX/index/tex/components/00_tomeIII_index_entries.tex
\Idx{\(\operatorname{ad}(G)\)}{127}
\Idx{Adjoint (canonical decomposition)}{243}
\Idx{Adjoint \(S\)-reductive group}{127}
\Idx{Adjoint root datum}{91}
\Idx{\(\operatorname{ad}(\mathscr R)\)}{95, 98}
\Idx{\(\alpha^*\)}{110}
\Idx{Paired sections of \(\mathfrak g^\alpha\) and
  \(\mathfrak g^{-\alpha}\)}{110}
\Idx{\(A_S(\mathscr R)\)}{223}
\Idx{\(\underline{\operatorname{Aut}}_{S\text{-}\mathrm{gr}}(G)\)}{217}
\Idx{Automorphisms of a reductive group}{216}
\Idx{Automorphisms of a root datum}{98}
\Idx{Automorphisms of the Borel subgroups of reductive groups}{248}
\Idx{\(\operatorname{Aut}(\mathscr R)\)}{98}
\Idx{\(\operatorname{Aut}^s(\mathscr R)\)}{99}
\Idx{\(B^{\mathrm{ad}}\)}{224}
\Idx{Borel subgroups}{131}
\Idx{\(\underline{\operatorname{Bor}}(G)\)}{156}
\Idx{Bruhat decomposition}{149, 151}
\Idx{Cartan subgroups}{128}
\Idx{Center of an \(S\)-reductive group}{119}
\Idx{Weyl chambers}{83}
\Idx{Chevalley group scheme}{209}
\Idx{\(\operatorname{Cl}_{\mathrm{crit}}\)}{292}
\Idx{Closed set of roots}{71}
\Idx{Coroot of a reductive group scheme}{110, 111}
\Idx{Coroot of a root datum}{64}
\Idx{Infinitesimal coroot \(H_\alpha\)}{114}
\Idx{\(\operatorname{corad}(G)\)}{171}
\Idx{Coradical torus}{171}
\Idx{Coradical of a root datum}{92}
\Idx{\(\operatorname{corad}(\mathscr R)\)}{92}
\Idx{Killing pair}{135}
\Idx{\(\underline{\operatorname{Crit}}(G)\)}{292}
\Idx{Critical subgroup}{292}
\Idx{\(C\)-critical torus}{292}
\Idx{\(C^i(T)\)}{262}
\Idx{\(\underline{\mathrm{CT}}\)}{293}
\Idx{Splitting of a reductive group}{113}
\Idx{Split reductive group}{113}
\Idx{Split torus}{10}
\Idx{Splittable reductive group}{113}
\Idx{\(\operatorname{der}(\mathscr R)\)}{95}
\Idx{Dynkin diagram of a root datum}{105}
\Idx{Root datum}{63}
\Idx{Adjoint root datum}{91}
\Idx{Dual root datum}{64}
\Idx{Induced or coinduced root datum}{94}
\Idx{Irreducible root datum}{102}
\Idx{Reduced root datum}{68}
\Idx{Relative root datum}{325}
\Idx{Semisimple root datum}{64}
\Idx{Simply connected root datum}{91}
\Idx{Twisted root datum}{112}
\Idx{Trivial root datum}{64}
\Idx{Based root data}{179}
\Idx{\(\underline{\operatorname{Dyn}}(G)\)}{228}
\Idx{\(\underline{\operatorname{Dyn}}_0(G)\)}{244}
\Idx{\(\underline{\operatorname{Dyn}}_{0,t}(G)\)}{244}
\Idx{\(E_\Delta(\mathscr R)\)}{98}
\Idx{\(E^s_\Delta(\mathscr R)\)}{99}
\Idx{Generation of \(G(S)\) by the \(U_\alpha(S)\), \(S\) local}{151}
\Idx{Pinnings}{177}
\Idx{Pinned \(S\)-reductive group}{178}
\Idx{Essentially free \(S\)-scheme}{132, 152}
\Idx{\(\exp\)}{29}
\Idx{\(\exp_\alpha\)}{109}
\Idx{Exponential map}{29}
\Idx{\((\mathrm F)\), formula}{35}
\Idx{\(\mathscr F^\alpha\)}{12}
\Idx{\(\underline{\operatorname{Fib}}(S,G)\)}{215}
\Idx{Form of \(G\) over \(S\)}{222}
\Idx{\(\mathfrak g^\alpha\)}{5}
\Idx{\(\Gamma_0(\mathscr R)\)}{64}
\Idx{Generators and relations for a pinned group}{182}
\Idx{\(\underline{\operatorname{Gen}}(G)\)}{299}
\Idx{\(\underline{\operatorname{Gen}}(P/Q)\)}{299}
\Idx{\(\underline{\operatorname{Gen}}(/Q)\)}{299}
\Idx{\(G^{q\text{-}\text{ép}.}\)}{233}
\Idx{Big cell}{118, 271}
\Idx{Pinned group}{216}
\Idx{Group of type \textup{(RA)}}{130}
\Idx{Group of type \textup{(RR)}}{128}
\Idx{Derived group of an \(S\)-reductive group}{171}
\Idx{Quasi-pinned group}{230}
\Idx{Extended Weyl group}{212}
\Idx{Weyl group of a torus}{6}
\Idx{Weyl group of a root datum}{65}
\Idx{Generators and relations for the Weyl group of a root datum}{86}
\Idx{Weyl group of \((G,T)/S\)}{116}
\Idx{Weyl group in semisimple rank one}{44}
\Idx{Reductive group over an algebraically closed field}{4}
\Idx{Reductive group over an arbitrary base}{12}
\Idx{Semisimple group associated with an \(S\)-reductive group}{127}
\Idx{Semisimple group over an algebraically closed field}{5}
\Idx{Semisimple group over an arbitrary base}{12}
\Idx{\(G_S^{\text{Ép}}(\mathscr R)\)}{209}
\Idx{Indivisible root}{68}
\Idx{Central isogeny of \(S\)-reductive groups}{124}
\Idx{Isogeny of \(S\)-reductive groups}{124}
\Idx{Isogeny of root data}{90}
\Idx{\(\underline{\operatorname{Isom}}^{\mathrm{ext}}(G,G')\)}{220}
\Idx{\(\underline{\operatorname{Isom}}^{\mathrm{int},u}(G,G')\)}{226}
\Idx{Locally isotrivial}{237}
\Idx{Semilocally isotrivial}{237}
\Idx{Isotypic component}{243}
\Idx{\(\underline{\operatorname{Kil}}(G)\)}{156}
\Idx{\(\Lambda(\mathscr R)\)}{97}
\Idx{\(\underline{\operatorname{Lev}}(P)\)}{282}
\Idx{Cartan matrix of a root datum}{104}
\Idx{Morphisms of root data}{90}
\Idx{Morphisms of pinned groups}{178}
\Idx{Morphisms of split groups}{122}
\Idx{\(\underline{\operatorname{Opp}}(B)\)}{158}
\Idx{\(\underline{\operatorname{Opp}}(G)\)}{305}
\Idx{\(\underline{\operatorname{Opp}}(/P)\)}{305}
\Idx{Opposite parabolic subgroups}{303}
\Idx{\(\operatorname{ord}_\Delta(\alpha)\)}{76}
\Idx{\(\operatorname{Of}(E)\)}{289}
\Idx{\(p_\alpha\)}{27, 109}
\Idx{\(\underline{\operatorname{Par}}(G)\)}{289}
\Idx{\(\underline{\mathrm{PL}}\)}{293}
\Idx{\(\underline{\mathrm{PLT}}\)}{293}
\Idx{\(p\)-morphisms of based root data}{179}
\Idx{\(p\)-morphisms of reduced root data}{100}
\Idx{Weights of a root datum}{97}
\Idx{Fundamental weights}{97}
\Idx{Borel subgroups in general position}{158}
\Idx{Relative position of two parabolic subgroups}{296}
\Idx{Standard position of two parabolic subgroups}{308}
\Idx{Transverse position of two parabolic subgroups}{297}
\Idx{\(\underline{\mathrm{PT}}\)}{293}
\Idx{Quasi-splittable \(S\)-reductive group}{230}
\Idx{\(\underline{\mathrm{Q\acute ep}}\)}{232}
\Idx{Quasi-pinnings}{230}
\Idx{Central quotients of reductive groups}{126}
\Idx{Chevalley's rule}{211}
\Idx{Root of a reductive group scheme}{13, 109}
\Idx{Root of a root datum}{64}
\Idx{Infinitesimal root \(\overline\alpha\)}{114}
\Idx{\(\rad^u(G)\)}{2}
\Idx{\(\rad(G)\)}{127}
\Idx{Radical of an \(S\)-reductive group}{127}
\Idx{Radical of an algebraic group}{2}
\Idx{Radical of a parabolic subgroup}{284}
\Idx{Radical of a root datum}{92}
\Idx{Unipotent radical}{2}
\Idx{Unipotent radical of a type \textup{(RC)} subgroup}{166, 281}
\Idx{Unipotent radical of a parabolic subgroup}{281, 283}
\Idx{\(\operatorname{rad}(\mathscr R)\)}{92}
\Idx{\(\operatorname{rad}(G)\)}{2}
\Idx{\(\operatorname{rad}^u(P)\)}{281, 283}
\Idx{Reductive rank of a smooth affine \(k\)-group}{4}
\Idx{Reductive rank of a root datum}{64}
\Idx{Semisimple rank of a smooth affine \(k\)-group}{5}
\Idx{Semisimple rank of a root datum}{64}
\Idx{\(\operatorname{rgred}(G)\)}{4}
\Idx{\(\operatorname{rgred}(\mathscr R)\)}{64}
\Idx{\(\operatorname{rgss}(G)\)}{5}
\Idx{\(\operatorname{rgss}(\mathscr R)\)}{64}
\Idx{\(\mathscr R,\mathscr R^*\), root data}{64}
\Idx{\(\underline{\operatorname{Redext}}\)}{232}
\Idx{\(\underline{\operatorname{Rev}}\)}{232}
\Idx{\(\mathscr R(G)\)}{179}
\Idx{\(\mathcal R\), scheme of roots}{15}
\Idx{\(s_\alpha(t)\)}{110}
\Idx{Dynkin scheme of a reductive group}{228}
\Idx{Scheme of coroots of an \(S\)-reductive group}{111}
\Idx{Scheme of roots of an \(S\)-reductive group}{15}
\Idx{Scheme of parabolic subgroups of a reductive group}{289}
\Idx{Scheme of types of parabolic subgroups of a reductive group}{291}
\Idx{Reductive group scheme}{8}
\Idx{Henselian local scheme}{220, 223, 261, 330}
\Idx{Semilocal scheme}{162, 230, 239, 241, 287, 288, 311, 318, 322}
\Idx{\(\operatorname{sc}(\mathscr R)\)}{95}
\Idx{\(\operatorname{ss}(\mathscr R)\)}{95}
\Idx{\(s_G\)}{236}
\Idx{Simply connected \(S\)-reductive group}{127}
\Idx{Simply connected canonical decomposition}{243}
\Idx{Simply connected root datum}{91}
\Idx{Subgroups with commutative quotients}{174}
\Idx{Cartan subgroups}{128}
\Idx{Minimal parabolic subgroups}{316}
\Idx{Borel subgroups}{131}
\Idx{Borel subgroups of a split reductive group}{139}
\Idx{Levi subgroups}{282}
\Idx{Subgroups of type \textup{(R)}}{131}
\Idx{Type \textup{(R)} subgroups with reductive fibers}{162}
\Idx{Type \textup{(R)} subgroups of a split reductive group}{137}
\Idx{Subgroups of type \textup{(RC)}}{165}
\Idx{Parabolic subgroups}{131, 279}
\Idx{Critical reductive subgroups}{163}
\Idx{\(\underline{\operatorname{Stand}}(G)\)}{309}
\Idx{Root system}{64}
\Idx{Root system of a reductive group scheme}{14}
\Idx{Positive root system}{73}
\Idx{Simple root system}{71}
\Idx{Chevalley systems}{269}
\Idx{Elementary system}{28}
\Idx{Generators and relations for an elementary system}{58}
\Idx{Chevalley systems}{209}
\Idx{\(t_{\alpha\beta}=(w_\alpha w_\beta)^{n_{\alpha\beta}}\)}{180}
\Idx{\(T^{\mathrm{ad}}\)}{224}
\Idx{\(T_\alpha\)}{7}
\Idx{Uniqueness theorem}{207}
\Idx{Theorem 90}{215}
\Idx{Existence theorem}{267}
\Idx{Bruhat theorem}{149, 151}
\Idx{Conjugacy theorem}{311}
\Idx{Fundamental theorem}{202}
\Idx{Coradical torus of an \(S\)-reductive group}{171}
\Idx{Critical torus}{163}
\Idx{Split torus}{10}
\Idx{Maximal torus of an \(S\)-group \(G\)}{10, 23}
\Idx{Trivialized torus}{10}
\Idx{\(\underline{\operatorname{Tor}}(G)\)}{156}
\Idx{Strict transporter of two type \textup{(R)} subgroups}{134}
\Idx{Strict transporter of parabolic subgroups}{280}
\Idx{Braid relations}{212}
\Idx{Type of a reductive group}{115}
\Idx{Type at a point \(s\)}{115}
\Idx{\(\mathbf t\)}{289}
\Idx{\(\mathbf t(P)\)}{289}
\Idx{\(U_\alpha\)}{33, 110}
\Idx{\(V^\times\)}{16}
\Idx{\(w_\alpha\)}{7}
\Idx{\(w_\alpha(X)\)}{111}
\Idx{\(W(T)\)}{6}
\Idx{\(W(\mathscr F)^\alpha\)}{12}
\Idx{\(W(\mathscr R)\)}{65, 98}
\Idx{\(W^*(\mathscr R)=W(\mathscr R^*)\)}{65}
\Idx{\(W^*\)}{212}
\Idx{\(X_\alpha\)}{109}
\Idx{\(Z_\alpha\)}{7}
\Idx{\(Z(\mathscr R)\)}{98}
\Idx{\(Z^1(S'/S,G)\)}{262}

% END INCLUDED SGA3_English_current_buildable_source_20260823/source/inputs/TERMINAL_INDEX/index/tex/components/00_tomeIII_index_entries.tex

\end{multicols}

% END INCLUDED SGA3_English_current_buildable_source_20260823/source/inputs/TERMINAL_INDEX/index/tex/SGA3_TomeIII_Index_body.tex


% END INCLUDED SGA3_English_current_buildable_source_20260823/source/inputs/TERMINAL_INDEX/body.tex


\clearpage
\chapter*{Cumulative boundary}
\addcontentsline{toc}{chapter}{Cumulative boundary}
This reader ends with the terminal Tome-III index.  It contains no
material beyond SGA~3 and makes no archive, publication, public-license,
or public-readback claim.

\end{document}

% SGA-CUMULATIVE-END 03_SGA3_EN_COMPLETE_READER.tex

% SGA-CUMULATIVE-BEGIN 04_SGA4_EN_COMPLETE_READER.tex
% Release engine: XeLaTeX (the LuaTeX compatibility layer breaks Xy-pic output).
\documentclass[12pt]{../fichiers-configuration-latex/amsbook}
\usepackage[a4paper,left=3cm,right=3cm,marginpar=0.5cm,marginparsep=0.8cm,vmargin=2.8cm]{geometry}
% BEGIN INCLUDED SGA4_English_current_buildable_source_20260906/source/fichiers-configuration-latex/paquets.tex
\usepackage{polyglossia}
%\usepackage[francais]{babel}
\usepackage{fontspec}
\usepackage[math-style=ISO,bold-style=upright,sans-style=upright]{unicode-math}
%\usepackage[mathletters]{ucs}
%\usepackage[utf8]{inputenc}
%\usepackage[T1]{fontenc}
\usepackage[unicode,hypertexnames=false,colorlinks=true,linkcolor=blue,citecolor=Sepia,urlcolor=magenta]{hyperref}
\setmainlanguage{french}
\setmainfont[Ligatures=TeX]{Linux Libertine O}
% variante possible : TeXGyrePagella ?
\setmathfont{XITS Math}

% maths
\usepackage{amsmath}
%\usepackage{amssymb}
%\usepackage{mathrsfs}
%\usepackage{stmaryrd}
%\usepackage{euscript}

% couleurs
\usepackage[usenames,dvipsnames]{xcolor}

% graphiques
\usepackage{graphicx}
\usepackage[all]{xy}
%\usepackage{tikz}
%\usetikzlibrary{matrix,arrows,calc}
%\tikzset{column sep=3em,row sep=2em,text height=1.5ex, text depth=0.25ex,>=angle 90}
%\tikzset{description/.style={fill=white,inner sep=1pt}}

% hyperliens, 'etiquettes
\usepackage{xr-hyper}
\usepackage{../fichiers-configuration-latex/smfthm}
%\ifx\cacheretiquettes\undefined
%\usepackage[inline]{showlabels} \else \fi % pour rendre visibles les 'etiquettes
%\usepackage{../fichiers-configuration-latex/makeidx} % pour faire un index
%\usepackage[pagebackref,unicode]{hyperref}  % pour voir o`u sont cit'ees les r'ef'erences
% hyperref is loaded once above with the cumulative-reader destination policy.
%\hypersetup{breaklinks=true} %permet le retour `a la ligne dans les liens trop longs. Incompatible avec pdflatex
%\renewcommand*{\backref}[1]{{\small [cit\'e page(s)~#1.]}}

\usepackage{xspace}
\usepackage{pifont} % pour expose XVII

% END INCLUDED SGA4_English_current_buildable_source_20260906/source/fichiers-configuration-latex/paquets.tex

% BEGIN INCLUDED SGA4_English_current_buildable_source_20260906/source/fichiers-configuration-latex/numerotation.tex
\renewcommand{\thefootnote}{\roman{footnote}} % pour mettre les notes en bas de page sous la forme (ii) et non (2), qui ressemble trop (?) a un exposant.

\newcounter{pageoriginal}
\newcommand{\pageoriginale}{\refstepcounter{pageoriginal}\marginpar{\footnotesize\textbf{\thepageoriginal}}\xspace}
\let\pageoriginaled\pageoriginale

\usepackage{ifthen}
\newboolean{original}
\def\sisi{\ifthenelse{\boolean{original}}}
\setboolean{original}{false}

\newcounter{cnote}
\def\nde#1{\ifthenelse{\boolean{original}}{}{\stepcounter{cnote}\def\thefootnote{\arabic{cnote}}\protect\footnote[\value{cnote}]{N.D.E.: #1}}}
\def\ndemark{\ifthenelse{\boolean{original}}{}{\stepcounter{cnote}\def\thefootnote{\arabic{cnote}}\protect\footnotemark[\value{cnote}]}}
\def\ndetext#1{\ifthenelse{\boolean{original}}{}{\def\thefootnote{\arabic{cnote}}\protect\footnotetext[\value{cnote}]{N.D.E.: #1}}}

% END INCLUDED SGA4_English_current_buildable_source_20260906/source/fichiers-configuration-latex/numerotation.tex

% BEGIN INCLUDED SGA4_English_current_buildable_source_20260906/source/fichiers-configuration-latex/macros.tex
\setmainlanguage{french}
\setotherlanguage{english}


%% Fixed, engine-independent release identifier.  The inherited dynamic Git
%% probe used Lua io.popen, exposed build-state-dependent output, and printed
%% a visible question mark outside a Git checkout.
\def\showgitstatus{English release r2, 2026-07-23}


% macros pour LaTeX
\renewcommand{\mathsf}{}

% Police différente pour U+2113 SCRIPT SMALL L
\setmathfont[range="2113]{latinmodern-math.otf}

% Give every inherited semantic label its own stable PDF destination instead
% of inheriting whichever automatic equation/page anchor happened to be live.
% The legacy label remains as a compatibility alias used by existing edges.
\makeatletter
\def\etiquette[#1]#2{%
  \Hy@raisedlink{\hyper@anchorstart{sga4.semantic.#2}\hyper@anchorend}%
  \def\@currentHref{sga4.semantic.#2}%
  \def\@currentlabel{#1}%
  \label{#2}%
}
\makeatother
\def\eqhyperref[#1]#2{(\hyperref[#1]{#2})}
\def\pagehyperref[#1]#2{(page \hyperref[#1]{#2})}
\def\hyperrefext[#1][#2]#3{\hyperref[#2]{#1 #3}}
\def\eqhyperrefext[#1][#2]#3{\hyperref[#2]{#1 #3}}     %cf. exposes 17
\def\neweqhyperrefext[#1][#2]#3{\hyperref[#2]{#1 #3}}  %cf. exposes 17

\def\shat#1{\protect\widehat{#1}}
\def\stilde#1{\protect\widetilde{#1}}
%%%%%%%%%%%%%%%%%%


\def\etcn{etc.\xspace}
\def\etco{etc...\xspace}
\def\etc{\sisi{\etco}{\etcn}}
\def\ldotsn{\sisi{...}{\etcn}}

\def\cf{{cf.}\kern.3em}
\def\Cf{{Cf.}\kern.3em}
\def\ie{{i.e.}\kern.3em}
\def\iev{{i.e.},\xspace}
\def\loccito{loc.\kern2pt cit.\xspace}
\def\loccitn{\textit{loc.\kern2pt cit.}\xspace}
\def\loccit{\sisi{\loccito}{\loccitn}}
\def\EGA{\textup{EGA}\kern.2em}
\def\SGA{\textup{SGA}\kern.3em}

\def\quoi{}

\def\quav{\quad,\quad}
\def\quaet{\quad \text{et}\quad}
\def\fibreproduct#1#2#3{#1 ×_{#3} #2}
\let\fprod\fibreproduct
\def\ccheck#1#2{\mbox{${\check{\check{\rm #1}}}^{\lower2pt\hbox{$\scriptstyle #2$}}$}}
\def\sheaf#1{\protect\mathcal{#1}}
\def\vieuxsheaf#1{{#1}} % cf. problème avec sheaf{a\Phi_i} par exemple...
\DeclareMathOperator{\Hom}{Hom}
\DeclareMathOperator{\can}{can}
\DeclareMathOperator{\Ind}{Ind}
\DeclareMathOperator{\ind}{ind}
%\DeclareMathOperator{\Hhomtop}{\mathcal{H}\!\mathit{omtop}}
\DeclareMathOperator{\Hhomtop}{\mathbf{Homtop}}
\DeclareMathOperator{\Point}{Point}
\DeclareMathOperator{\Pro}{Pro}
\DeclareMathOperator{\Top}{Top}
\DeclareMathOperator{\Mod}{Mod}
\DeclareMathOperator{\Cart}{Cart}
%\DeclareMathOperator{\Mmorsite}{\mathcal{M}\!\mathit{orsite}}
\DeclareMathOperator{\Mmorsite}{\mathbf{Morsite}}
\DeclareMathOperator{\loc}{loc}
\DeclareMathOperator{\Loc}{Loc}
\DeclareMathOperator{\Pic}{Pic}
\DeclareMathOperator{\Et}{Et}
\DeclareMathOperator{\card}{card}
\DeclareMathOperator{\Ens}{Ens}
\DeclareMathOperator{\h}{\underline{H}}
\DeclareMathOperator{\Fonct}{Fonct}
\DeclareMathOperator{\Ker}{Ker}
\DeclareMathOperator{\but}{but}
\DeclareMathOperator{\Ob}{Ob}
\DeclareMathOperator{\ob}{ob}
\DeclareMathOperator{\Lib}{Lib}
\DeclareMathOperator{\esj}{esj}
\DeclareMathOperator{\id}{id}
\DeclareMathOperator{\Id}{Id}
\DeclareMathOperator{\ad}{ad}
\DeclareMathOperator{\Iim}{Im}
\DeclareMathOperator{\pr}{pr}
\DeclareMathOperator{\Cat}{Cat}
\DeclareMathOperator{\car}{car}
\DeclareMathOperator{\SysInd}{SysInd}
\DeclareMathOperator{\Sup}{Sup}
\DeclareMathOperator{\Coker}{Coker}
\DeclareMathOperator{\Filt}{Filt}
\DeclareMathOperator{\acc}{acc}
\DeclareMathOperator{\sur}{sur}
\DeclareMathOperator{\cat}{cat}
\DeclareMathOperator{\Ord}{Ord}
\DeclareMathOperator{\Gr}{Gr}
\DeclareMathOperator{\Cov}{Cov}
\DeclareMathOperator{\Caf}{Caf}
\DeclareMathOperator{\coker}{coker}
\DeclareMathOperator{\coim}{coim}
\DeclareMathOperator{\im}{im}
\DeclareMathOperator{\Ouv}{Ouv}
\DeclareMathOperator{\diag}{diag}
\DeclareMathOperator{\Esp}{Esp}
\DeclareMathOperator{\TOP}{TOP}
\DeclareMathOperator{\Unif}{Unif}
\DeclareMathOperator{\Iind}{\mathbf{Ind}}
\DeclareMathOperator{\Hhomcart}{\mathbf{Homcart}}
\DeclareMathOperator{\Hhomtopess}{\mathbf{Homtopess}}
\DeclareMathOperator{\Ppoints}{\mathbf{Points}}
\DeclareMathOperator{\Ppoint}{\mathbf{Point}}
\DeclareMathOperator{\Ttors}{\mathbf{Tors}}
\DeclareMathOperator{\Ssex}{\mathbf{Sex}}
\DeclareMathOperator{\Ffib}{\mathbf{Fib}}
\DeclareMathOperator{\Vvois}{\mathbf{Vois}}
\DeclareMathOperator{\Ppointess}{\mathbf{Pointess}}
\DeclareMathOperator{\scrHom}{\mathbf{Hom}}
\DeclareMathOperator{\sheafPic}{\mathbf{Pic}}
\DeclareMathOperator{\sheafAut}{\mathbf{Aut}}
\DeclareMathOperator{\Hhomtoplocan}{\mathbf{Homtoplocan}}
\DeclareMathOperator{\Hhomtopan}{\mathbf{Homtopan}}
\DeclareMathOperator{\Pptgeom}{\mathbf{Ptgéom}}
\DeclareMathOperator{\Morsite}{Morsite}
\DeclareMathOperator{\Fib}{Fib}
\DeclareMathOperator{\Fais}{Fais}
\DeclareMathOperator{\kar}{kar}
\DeclareMathOperator{\Groupes}{Groupes}
\DeclareMathOperator{\Restr}{Restr}
\DeclareMathOperator{\Isom}{Isom}
\DeclareMathOperator{\Prol}{Prol}
\DeclareMathOperator{\Res}{Res}
\DeclareMathOperator{\Tors}{Tors}
\DeclareMathOperator{\esp}{esp}
\DeclareMathOperator{\es}{es}
\DeclareMathOperator{\od}{od}
\DeclareMathOperator{\gr}{gr}
\DeclareMathOperator{\sob}{sob}
\DeclareMathOperator{\dfn}{dfn}
\DeclareMathOperator{\Sch}{Sch}
\DeclareMathOperator{\ouv}{ouv}
\DeclareMathOperator{\cons}{cons}
\DeclareMathOperator{\Nerf}{Nerf}
\DeclareMathOperator{\Points}{Points}
\DeclareMathOperator{\Espsob}{Espsob}
\DeclareMathOperator{\Homtop}{Homtop}
\DeclareMathOperator{\Pt}{Pt}
\DeclareMathOperator{\Quot}{Quot}
\DeclareMathOperator{\glob}{glob}
\DeclareMathOperator{\Homfib}{Homfib}
\DeclareMathOperator{\Qouv}{Qouv}
\DeclareMathOperator{\Quotouv}{Quotouv}
\DeclareMathOperator{\ens}{ens}
\DeclareMathOperator{\ab}{ab}
\DeclareMathOperator{\an}{an}
\DeclareMathOperator{\fs}{fs}
\DeclareMathOperator{\red}{red}
\DeclareMathOperator{\Ab}{Ab}
\DeclareMathOperator{\cosq}{cosq}
\DeclareMathOperator{\Ann}{Ann}
\DeclareMathOperator{\Zar}{Zar}


%%%%%%%%%%%%
%%%%% SGA4-2
\DeclareMathOperator{\Sex}{Sex}
\DeclareMathOperator{\vois}{vois}
\DeclareMathOperator{\Vois}{Vois}
\DeclareMathOperator{\Ext}{Ext}
\DeclareMathOperator{\ext}{ext}
\DeclareMathOperator{\Fer}{Fer}
\DeclareMathOperator{\Add}{Add}
\DeclareMathOperator{\Cocart}{Cocart}
\DeclareMathOperator{\cosk}{cosk}
\DeclareMathOperator{\plat}{plat}

\DeclareMathOperator{\faisc}{faisc}
\DeclareMathOperator{\sk}{sk}
\DeclareMathOperator{\ch}{ch}

\def\HH{\mathbf{H}}
\def\qc{\mathrm{qc}}
\def\et{\mathrm{et}}
\def\coh{\mathrm{coh}}

\def\Hvr#1{%
\begin{array}[b]{l}
{\kern2pt\scriptstyle r}\\[-8pt]
{\kern1pt\scriptstyle \vee}\\[-5pt]
H^{#1}
\end{array}}

\def\Hvv#1{%
\begin{array}[b]{l}
\kern2pt\rotatebox{90}{${\scriptscriptstyle\ll}$}\\[-6pt]
H^{#1}
\end{array}}

\def\Hvi#1{%
\begin{array}[b]{l}
{\kern0.2pt\scriptscriptstyle \infty}\\[-8pt]
{\kern1pt\scriptstyle \vee}\\[-5pt]
H^{#1}
\end{array}}

\def\Hv#1{%
\begin{array}[b]{l}
{\kern1pt\scriptstyle \vee}\\[-5pt]
H^{#1}
\end{array}}

\DeclareMathOperator{\cart}{cart}
\DeclareMathOperator{\Ffaisc}{\mathbf{Faisc}}


\def\Ev#1{%
\begin{array}[b]{l}
{\kern1pt\scriptstyle \vee}\\[-5pt]
E'_{#1}
\end{array}}

\DeclareMathOperator{\Cont}{Cont}
\DeclareMathOperator{\Ccont}{\mathscr{C}\!\mathit{ont}}
\DeclareMathOperator{\Site}{Site}
\newcommand\MM{\mathbf{M}}
\def\limtop#1{{\displaystyle{\varprojlim_{#1}\!\mbox{\rm top }}}}
\def\Limtop{\mathop{\varprojlim\!\mbox{\rm top}}}
\DeclareMathOperator{\spec}{spec}
\def\ss{\scriptstyle{\mathscr{S}}}
\DeclareMathOperator{\Gal}{Gal}
\def\smallm{\scriptstyle{\mathscr{M}}}
\DeclareMathOperator{\End}{End}
\DeclareMathOperator{\ass}{ass}
\DeclareMathOperator{\zar}{zar}
\DeclareMathOperator{\Ogg}{Ogg}
\DeclareMathOperator{\res}{res}
\DeclareMathOperator{\tr}{tr}
\DeclareMathOperator{\Supp}{Supp}
\DeclareMathOperator{\codim}{codim}
\DeclareMathOperator{\rk}{rk}
\DeclareMathOperator{\cur}{cur}
\DeclareMathOperator{\cdqc}{cdqc}
\newcommand\DD{\mathbf{D}}
\DeclareMathOperator{\cl}{cl}
\def\JJ{\mathbf{J}}
\def\FF{\mathbf{F}}
\def\BB{\mathbf{B}}
\DeclareMathOperator{\Alget}{Alget}
\DeclareMathOperator{\Ass}{Ass}
\DeclareMathOperator{\supp}{supp}
\DeclareMathOperator{\rad}{rad}
\DeclareMathOperator{\prof}{prof}
\DeclareMathOperator{\cint}{int}
\DeclareMathOperator{\sic}{sic}
\DeclareMathOperator{\Tor}{Tor}
\DeclareMathOperator{\tors}{tors}
\DeclareMathOperator{\Sym}{Sym}
\DeclareMathOperator{\parf}{parf}
\DeclareMathOperator{\tor}{tor}
\DeclareMathOperator{\defn}{defn}
\DeclareMathOperator{\Div}{Div}
\DeclareMathOperator{\Pol}{Pol}
\DeclareMathOperator{\Fex}{Fex}
\DeclareMathOperator{\Lie}{Lie}
\DeclareMathOperator{\tot}{tot}
\DeclareMathOperator{\Aut}{Aut}
\DeclareMathOperator{\HOM}{HOM}
\DeclareMathOperator{\pch}{pch}
\DeclareMathOperator{\Ch}{Ch}
\DeclareMathOperator{\EXT}{EXT}
\DeclareMathOperator{\PIC}{PIC}
\DeclareMathOperator{\Rf}{Rf}
\DeclareMathOperator{\inj}{inj}
\newcommand\VV{\mathbf{V}}

\def\Xv{%
\begin{array}[b]{l}
{\kern1pt\scriptstyle \vee}\\[-5pt]
X
\end{array}}

%%%%%%%
\newcommand\TT{\mathbf{T}}
\newcommand\XX{\mathbf{X}}
\newcommand\EE{\mathbf{E}}
%%%%%%

\def\qedsymbol{\textup{C.Q.F.D.}}
\let\CQFD\qed
\let\cqfd\qed

%%%%%\bigdot
\font\bigsymb=cmsy10 at 4pt
\def\bigdot{{\kern1.2pt\raise 1.5pt\hbox{\bigsymb\char15}}}
%%%%%

\def\Lright{\mathop{\vcenter{\hbox{\hbox{$\longrightarrow$}\lower8pt \hbox{\kern-16pt\hbox{$\longrightarrow$}}}}}}

\def\Aa{{\mathcal A}}
\def\Bb{{\mathcal B}}
\def\Cc{{\mathcal C}}
\def\Dd{{\mathcal D}}
\def\Ee{{\mathcal E}}
\def\Ff{{\mathcal F}}
\def\Gg{{\mathcal G}}
\def\Hh{{\mathcal H}}
\def\Ii{{\mathcal I}}
\def\Jj{{\mathcal J}}
\def\Kk{{\mathcal K}}
\def\Ll{{\mathcal L}}
\def\Mm{{\mathcal M}}
\def\Nn{{\mathcal N}}
\def\Oo{{\mathcal O}}
\def\Pp{{\mathcal P}}
\def\Qq{{\mathcal Q}}
\def\Rr{{\mathcal R}}
\def\Ss{{\mathcal S}}
\def\Tt{{\mathcal T}}
\def\Uu{{\mathcal U}}
\def\Vv{{\mathcal V}}
\def\Ww{{\mathcal W}}
\def\Xx{{\mathcal X}}
\def\Yy{{\mathcal Y}}
\def\Zz{{\mathcal Z}}

\newcommand{\CC}{\mathbf{C}}
\newcommand{\GG}{\mathbf{G}}
\newcommand{\KK}{\mathbf{K}}
\newcommand{\LL}{\mathbf{L}}
\newcommand{\NN}{\mathbf{N}}
\newcommand{\PP}{\mathbf{P}}
\newcommand\QQ{\mathbf{Q}}
\newcommand\RR{\mathbf{R}}
\newcommand\ZZ{\mathbf{Z}}

\def\to{\mathchoice{\longrightarrow}{\rightarrow}{\rightarrow}{\rightarrow}}
\def\from{\mathchoice{\longleftarrow}{\leftarrow}{\leftarrow}{\leftarrow}}
\def\mto{\mathchoice{\longmapsto}{\mapsto}{\mapsto}{\mapsto}}
\def\mfrom{\mathrel{\from\joinrel\mapstochar\,}}
\def\hto{\mathrel{\lhook\joinrel\to}}
\def\hfrom{\mathrel{\from\joinrel\rhook}}

\def\implies{\mathchoice{\Longrightarrow}{\Rightarrow}{\Rightarrow}{\Rightarrow}}
\def\implied{\mathchoice{\Longleftarrow}{\Leftarrow}{\Leftarrow}{\Leftarrow}}
\def\iff{\mathchoice{\Longleftrightarrow}{\Leftrightarrow}{\Leftrightarrow}{\Leftrightarrow}}

\def\To#1{\mathchoice{\xrightarrow{\labelstyle~~#1~}}{\stackrel{#1}{\longrightarrow}}{}{}}
\def\Mto#1{\mathrel{\mapstochar\To{#1}}}
\def\From#1{\mathchoice{\xleftarrow{\labelstyle~~#1~}}{\stackrel{#1}{\longleftarrow}}{}{}}

\newcommand{\isomto}{\overset{\sim}{\longrightarrow}}
\newcommand{\isomfrom}{\overset{\sim}{\longleftarrow}}

\newcommand{\IFF}{\unskip~$\iff$~\ignorespaces}
\newcommand{\THEN}{\unskip~$\implies$~\ignorespaces}

\let\leq\leqslant
\let\geq\geqslant
\let\le\leqslant
\let\ge\geqslant

\def\cf{{cf.}\kern.3em}
\def\Cf{{Cf.}\kern.3em}
\def\ie{{i.e.}\kern.3em}
\def\iev{{i.e.},\xspace}
%\def\loccito{loc.\kern2pt cit.\xspace}
%\def\loccitn{\textit{loc.\kern2pt cit.}\xspace}
%\def\loccit{\sisi{\loccito}{\loccitn}}
\def\EGA{\textup{EGA}\kern.2em}
\def\SGA{\textup{SGA}\kern.3em}

\let\oldS\S
\def\S{\oldS\kern .15em}
\def\ptbl{.\kern .2em}
\def\resp{resp.\kern.3em}
\def\Exp{Exp.\kern.3em}

\let\bar\overline
\let\hat\widehat
\let\tilde\widetilde
\let\dpl\displaystyle
\let\ldots\dots
\let\emptyset\varnothing
\let\skpt\skippointrait
\let\epsilon\varepsilon

\def\sfrac#1#2{{#1}/{#2}}

\DeclareMathOperator{\cd}{cd}
%\DeclareMathOperator{\Hhom}{\mathcal{H}\!\mathit{om}}
\DeclareMathOperator{\Hhom}{\mathbf{Hom}}
%\DeclareMathOperator{\Eext}{\mathcal{E}\!\mathit{xt}}
\DeclareMathOperator{\Eext}{\mathbf{Ext}}
\DeclareMathOperator{\Spec}{Spec}
\DeclareMathOperator{\Tr}{Tr}

\let\emptyset=\varnothing

%%%%%%%%%%%%%%%
%%% expose XVII
%%%%%%%%%%%%%%%

\makeatletter
\let\old@optexpose\optexpose
\newcommand{\empty@optexpose}[1]{}
\newcommand{\paren@optexpose}[1]{\old@optexpose{#1}(}
\newcommand{\refnoexpose}[1]{\let\optexpose\empty@optexpose\ref{#1}\let\optexpose\old@optexpose}%\refnoexpose{X.cor5.2} engendre "5.2" au lieu de "X 5.2"
\newcommand{\neweqref}[1]{\let\optexpose\paren@optexpose\textup{\ref{#1})}\let\optexpose\old@optexpose}%\neweqref{XII.eq4.1} engendre "XII (4.1)" au lieu de "(XII 4.1)"
\makeatother
\newcommand{\refvb}[2][VI${}_{\text{B}}$ ]{\sisi{#1\refnoexpose{#2}}{\ref{#2}}}

\newcommand{\leftexp}[2]{{\vphantom{#2}}^{#1}{#2}}
\newcommand{\AAA}{\mathbf{A}}
\newcommand{\Sf}{\mathfrak{S}}
\DeclareMathOperator*{\quotprojlim}{\text{``$\varprojlim$''}}
\DeclareMathOperator*{\quotinjlim}{\text{``$\varinjlim$''}}
\DeclareMathOperator{\Fl}{F\sisi{\ell}{l}}
\DeclareMathOperator{\Pl}{P\sisi{\ell}{l}}
%\DeclareMathOperator{\cl}{c\sisi{\ell}{l}}
\DeclareMathOperator{\tdf}{tdf}
\DeclareMathOperator{\rg}{rg}
\newcommand{\iInd}{\sisi{I}{i}nd}
\newcommand{\iind}{ind}
\newcommand{\ppro}{pro}
\newcommand{\tTor}{\sisi{T}{t}or}
\newcommand{\ttor}{tor}
\newcommand{\chap}{Chap.}
\newcommand{\Ffex}{\Ff\mathit{ex}}
\newcommand{\sheafk}{\underline{k}}
\newcommand{\sheafn}{\underline{n}}
\newcommand{\sheafGamma}{\underline{\Gamma}}

% Repeated references to an already typeset footnote must not use
% \footnotemark without a matching \footnotetext: Hyperref would then create
% a live annotation to a destination that cannot exist.  This helper keeps the
% current printed footnote mark and links it to an explicit stable note target.
\makeatletter
\newcommand{\sgarepeatedfootnotemark}[1]{%
  \begingroup
  \edef\@thefnmark{\thefootnote}%
  \hyperlink{#1}{\@makefnmark}%
  \endgroup}
\makeatother

\sisi{
\DeclareMathOperator{\simarrow}{\rule{.15mm}{1.8mm}\kern -3.4pt\rightsquigarrow}
\newcommand{\sheafell}{\underline{\ell}}
\newcommand{\sheafZZ}{\underline{\ZZ}}
}{
\newcommand{\nderef}[1]{la N.D.E.~(\ref{#1}) page~\pageref{#1}}
}

%%%%%%%%%%%%%%%%%
%%%%% expose XVII
%%%%%%%%%%%%%%%%%

% END INCLUDED SGA4_English_current_buildable_source_20260906/source/fichiers-configuration-latex/macros.tex

\providecommand{\SGAArchiveOnly}[1]{}

\hypersetup{
  pdftitle={SGA 4: Theory of Topoi and Étale Cohomology of Schemes -- English Translation, Exposés I--XIX including V bis},
  pdfauthor={M. Artin; A. Grothendieck; J.-L. Verdier},
  pdfsubject={English translation of SGA 4 proper, Exposés I--XIX including V bis},
  pdfkeywords={SGA 4, topoi, étale cohomology, English translation, exposés I--XIX},
  pdfcreator={XeLaTeX}
}

\begin{document}
\frontmatter
\begin{center}
{\Large Theory of Topoi and Étale Cohomology of Schemes.\\
Volumes 1 to 3}\\
\end{center}

\vskip1cm
\begin{center}
M.~Artin, A.~Grothendieck, J.-L. Verdier
\end{center}

\vskip1cm
\begin{center}
With the participation of N.~Bourbaki, P.~Deligne, B.~Saint-Donat
\end{center}

\selectlanguage{english}
\SGAArchiveOnly{%
\begin{center}
\small Edition: \showgitstatus
\end{center}
}

\clearpage
\SGAArchiveOnly{%
\thispagestyle{empty}
\vspace*{\fill}
\begin{center}
\begin{minipage}{0.88\textwidth}
\small
\textbf{Source and rights notice.} This reader is based on the frozen
Orgogozo SGA~4 TeX snapshot at commit \texttt{71766d9}; the frozen source
archive has SHA-256
\texttt{DA2D939D3BD66B03E7BEFF4353550F7A27982847AE4F798D0872F1DB5D64C7DC}.
No license grant is asserted here for the underlying French text, this
English translation, or third-party TeX support files. Rights remain with
their respective holders; inclusion and provenance metadata do not by
themselves confer reuse permission.
\end{minipage}
\end{center}
\vspace*{\fill}
\clearpage
}
\tableofcontents

\mainmatter

\def\danslelivre{1}

\def\thechapter{\Roman{chapter}}
\def\chaptername{Exposé}

\makeatletter
\def\input@path{{../01/}}
\makeatother
% BEGIN INCLUDED SGA4_English_current_buildable_source_20260906/source/01/01_English.tex
\ifx\danslelivre\undefined
% Release engine: XeLaTeX (the LuaTeX compatibility layer breaks Xy-pic output).
\documentclass[10pt]{../fichiers-configuration-latex/smfart}
% BEGIN INCLUDED SGA4_English_current_buildable_source_20260906/source/fichiers-configuration-latex/paquets.tex
\usepackage{polyglossia}
%\usepackage[francais]{babel}
\usepackage{fontspec}
\usepackage[math-style=ISO,bold-style=upright,sans-style=upright]{unicode-math}
%\usepackage[mathletters]{ucs}
%\usepackage[utf8]{inputenc}
%\usepackage[T1]{fontenc}
\usepackage[unicode,hypertexnames=false,colorlinks=true,linkcolor=blue,citecolor=Sepia,urlcolor=magenta]{hyperref}
\setmainlanguage{french}
\setmainfont[Ligatures=TeX]{Linux Libertine O}
% variante possible : TeXGyrePagella ?
\setmathfont{XITS Math}

% maths
\usepackage{amsmath}
%\usepackage{amssymb}
%\usepackage{mathrsfs}
%\usepackage{stmaryrd}
%\usepackage{euscript}

% couleurs
\usepackage[usenames,dvipsnames]{xcolor}

% graphiques
\usepackage{graphicx}
\usepackage[all]{xy}
%\usepackage{tikz}
%\usetikzlibrary{matrix,arrows,calc}
%\tikzset{column sep=3em,row sep=2em,text height=1.5ex, text depth=0.25ex,>=angle 90}
%\tikzset{description/.style={fill=white,inner sep=1pt}}

% hyperliens, 'etiquettes
\usepackage{xr-hyper}
\usepackage{../fichiers-configuration-latex/smfthm}
%\ifx\cacheretiquettes\undefined
%\usepackage[inline]{showlabels} \else \fi % pour rendre visibles les 'etiquettes
%\usepackage{../fichiers-configuration-latex/makeidx} % pour faire un index
%\usepackage[pagebackref,unicode]{hyperref}  % pour voir o`u sont cit'ees les r'ef'erences
% hyperref is loaded once above with the cumulative-reader destination policy.
%\hypersetup{breaklinks=true} %permet le retour `a la ligne dans les liens trop longs. Incompatible avec pdflatex
%\renewcommand*{\backref}[1]{{\small [cit\'e page(s)~#1.]}}

\usepackage{xspace}
\usepackage{pifont} % pour expose XVII

% END INCLUDED SGA4_English_current_buildable_source_20260906/source/fichiers-configuration-latex/paquets.tex

% BEGIN INCLUDED SGA4_English_current_buildable_source_20260906/source/fichiers-configuration-latex/numerotation.tex
\renewcommand{\thefootnote}{\roman{footnote}} % pour mettre les notes en bas de page sous la forme (ii) et non (2), qui ressemble trop (?) a un exposant.

\newcounter{pageoriginal}
\newcommand{\pageoriginale}{\refstepcounter{pageoriginal}\marginpar{\footnotesize\textbf{\thepageoriginal}}\xspace}
\let\pageoriginaled\pageoriginale

\usepackage{ifthen}
\newboolean{original}
\def\sisi{\ifthenelse{\boolean{original}}}
\setboolean{original}{false}

\newcounter{cnote}
\def\nde#1{\ifthenelse{\boolean{original}}{}{\stepcounter{cnote}\def\thefootnote{\arabic{cnote}}\protect\footnote[\value{cnote}]{N.D.E.: #1}}}
\def\ndemark{\ifthenelse{\boolean{original}}{}{\stepcounter{cnote}\def\thefootnote{\arabic{cnote}}\protect\footnotemark[\value{cnote}]}}
\def\ndetext#1{\ifthenelse{\boolean{original}}{}{\def\thefootnote{\arabic{cnote}}\protect\footnotetext[\value{cnote}]{N.D.E.: #1}}}

% END INCLUDED SGA4_English_current_buildable_source_20260906/source/fichiers-configuration-latex/numerotation.tex

% BEGIN INCLUDED SGA4_English_current_buildable_source_20260906/source/fichiers-configuration-latex/macros.tex
\setmainlanguage{french}
\setotherlanguage{english}


%% Fixed, engine-independent release identifier.  The inherited dynamic Git
%% probe used Lua io.popen, exposed build-state-dependent output, and printed
%% a visible question mark outside a Git checkout.
\def\showgitstatus{English release r2, 2026-07-23}


% macros pour LaTeX
\renewcommand{\mathsf}{}

% Police différente pour U+2113 SCRIPT SMALL L
\setmathfont[range="2113]{latinmodern-math.otf}

% Give every inherited semantic label its own stable PDF destination instead
% of inheriting whichever automatic equation/page anchor happened to be live.
% The legacy label remains as a compatibility alias used by existing edges.
\makeatletter
\def\etiquette[#1]#2{%
  \Hy@raisedlink{\hyper@anchorstart{sga4.semantic.#2}\hyper@anchorend}%
  \def\@currentHref{sga4.semantic.#2}%
  \def\@currentlabel{#1}%
  \label{#2}%
}
\makeatother
\def\eqhyperref[#1]#2{(\hyperref[#1]{#2})}
\def\pagehyperref[#1]#2{(page \hyperref[#1]{#2})}
\def\hyperrefext[#1][#2]#3{\hyperref[#2]{#1 #3}}
\def\eqhyperrefext[#1][#2]#3{\hyperref[#2]{#1 #3}}     %cf. exposes 17
\def\neweqhyperrefext[#1][#2]#3{\hyperref[#2]{#1 #3}}  %cf. exposes 17

\def\shat#1{\protect\widehat{#1}}
\def\stilde#1{\protect\widetilde{#1}}
%%%%%%%%%%%%%%%%%%


\def\etcn{etc.\xspace}
\def\etco{etc...\xspace}
\def\etc{\sisi{\etco}{\etcn}}
\def\ldotsn{\sisi{...}{\etcn}}

\def\cf{{cf.}\kern.3em}
\def\Cf{{Cf.}\kern.3em}
\def\ie{{i.e.}\kern.3em}
\def\iev{{i.e.},\xspace}
\def\loccito{loc.\kern2pt cit.\xspace}
\def\loccitn{\textit{loc.\kern2pt cit.}\xspace}
\def\loccit{\sisi{\loccito}{\loccitn}}
\def\EGA{\textup{EGA}\kern.2em}
\def\SGA{\textup{SGA}\kern.3em}

\def\quoi{}

\def\quav{\quad,\quad}
\def\quaet{\quad \text{et}\quad}
\def\fibreproduct#1#2#3{#1 ×_{#3} #2}
\let\fprod\fibreproduct
\def\ccheck#1#2{\mbox{${\check{\check{\rm #1}}}^{\lower2pt\hbox{$\scriptstyle #2$}}$}}
\def\sheaf#1{\protect\mathcal{#1}}
\def\vieuxsheaf#1{{#1}} % cf. problème avec sheaf{a\Phi_i} par exemple...
\DeclareMathOperator{\Hom}{Hom}
\DeclareMathOperator{\can}{can}
\DeclareMathOperator{\Ind}{Ind}
\DeclareMathOperator{\ind}{ind}
%\DeclareMathOperator{\Hhomtop}{\mathcal{H}\!\mathit{omtop}}
\DeclareMathOperator{\Hhomtop}{\mathbf{Homtop}}
\DeclareMathOperator{\Point}{Point}
\DeclareMathOperator{\Pro}{Pro}
\DeclareMathOperator{\Top}{Top}
\DeclareMathOperator{\Mod}{Mod}
\DeclareMathOperator{\Cart}{Cart}
%\DeclareMathOperator{\Mmorsite}{\mathcal{M}\!\mathit{orsite}}
\DeclareMathOperator{\Mmorsite}{\mathbf{Morsite}}
\DeclareMathOperator{\loc}{loc}
\DeclareMathOperator{\Loc}{Loc}
\DeclareMathOperator{\Pic}{Pic}
\DeclareMathOperator{\Et}{Et}
\DeclareMathOperator{\card}{card}
\DeclareMathOperator{\Ens}{Ens}
\DeclareMathOperator{\h}{\underline{H}}
\DeclareMathOperator{\Fonct}{Fonct}
\DeclareMathOperator{\Ker}{Ker}
\DeclareMathOperator{\but}{but}
\DeclareMathOperator{\Ob}{Ob}
\DeclareMathOperator{\ob}{ob}
\DeclareMathOperator{\Lib}{Lib}
\DeclareMathOperator{\esj}{esj}
\DeclareMathOperator{\id}{id}
\DeclareMathOperator{\Id}{Id}
\DeclareMathOperator{\ad}{ad}
\DeclareMathOperator{\Iim}{Im}
\DeclareMathOperator{\pr}{pr}
\DeclareMathOperator{\Cat}{Cat}
\DeclareMathOperator{\car}{car}
\DeclareMathOperator{\SysInd}{SysInd}
\DeclareMathOperator{\Sup}{Sup}
\DeclareMathOperator{\Coker}{Coker}
\DeclareMathOperator{\Filt}{Filt}
\DeclareMathOperator{\acc}{acc}
\DeclareMathOperator{\sur}{sur}
\DeclareMathOperator{\cat}{cat}
\DeclareMathOperator{\Ord}{Ord}
\DeclareMathOperator{\Gr}{Gr}
\DeclareMathOperator{\Cov}{Cov}
\DeclareMathOperator{\Caf}{Caf}
\DeclareMathOperator{\coker}{coker}
\DeclareMathOperator{\coim}{coim}
\DeclareMathOperator{\im}{im}
\DeclareMathOperator{\Ouv}{Ouv}
\DeclareMathOperator{\diag}{diag}
\DeclareMathOperator{\Esp}{Esp}
\DeclareMathOperator{\TOP}{TOP}
\DeclareMathOperator{\Unif}{Unif}
\DeclareMathOperator{\Iind}{\mathbf{Ind}}
\DeclareMathOperator{\Hhomcart}{\mathbf{Homcart}}
\DeclareMathOperator{\Hhomtopess}{\mathbf{Homtopess}}
\DeclareMathOperator{\Ppoints}{\mathbf{Points}}
\DeclareMathOperator{\Ppoint}{\mathbf{Point}}
\DeclareMathOperator{\Ttors}{\mathbf{Tors}}
\DeclareMathOperator{\Ssex}{\mathbf{Sex}}
\DeclareMathOperator{\Ffib}{\mathbf{Fib}}
\DeclareMathOperator{\Vvois}{\mathbf{Vois}}
\DeclareMathOperator{\Ppointess}{\mathbf{Pointess}}
\DeclareMathOperator{\scrHom}{\mathbf{Hom}}
\DeclareMathOperator{\sheafPic}{\mathbf{Pic}}
\DeclareMathOperator{\sheafAut}{\mathbf{Aut}}
\DeclareMathOperator{\Hhomtoplocan}{\mathbf{Homtoplocan}}
\DeclareMathOperator{\Hhomtopan}{\mathbf{Homtopan}}
\DeclareMathOperator{\Pptgeom}{\mathbf{Ptgéom}}
\DeclareMathOperator{\Morsite}{Morsite}
\DeclareMathOperator{\Fib}{Fib}
\DeclareMathOperator{\Fais}{Fais}
\DeclareMathOperator{\kar}{kar}
\DeclareMathOperator{\Groupes}{Groupes}
\DeclareMathOperator{\Restr}{Restr}
\DeclareMathOperator{\Isom}{Isom}
\DeclareMathOperator{\Prol}{Prol}
\DeclareMathOperator{\Res}{Res}
\DeclareMathOperator{\Tors}{Tors}
\DeclareMathOperator{\esp}{esp}
\DeclareMathOperator{\es}{es}
\DeclareMathOperator{\od}{od}
\DeclareMathOperator{\gr}{gr}
\DeclareMathOperator{\sob}{sob}
\DeclareMathOperator{\dfn}{dfn}
\DeclareMathOperator{\Sch}{Sch}
\DeclareMathOperator{\ouv}{ouv}
\DeclareMathOperator{\cons}{cons}
\DeclareMathOperator{\Nerf}{Nerf}
\DeclareMathOperator{\Points}{Points}
\DeclareMathOperator{\Espsob}{Espsob}
\DeclareMathOperator{\Homtop}{Homtop}
\DeclareMathOperator{\Pt}{Pt}
\DeclareMathOperator{\Quot}{Quot}
\DeclareMathOperator{\glob}{glob}
\DeclareMathOperator{\Homfib}{Homfib}
\DeclareMathOperator{\Qouv}{Qouv}
\DeclareMathOperator{\Quotouv}{Quotouv}
\DeclareMathOperator{\ens}{ens}
\DeclareMathOperator{\ab}{ab}
\DeclareMathOperator{\an}{an}
\DeclareMathOperator{\fs}{fs}
\DeclareMathOperator{\red}{red}
\DeclareMathOperator{\Ab}{Ab}
\DeclareMathOperator{\cosq}{cosq}
\DeclareMathOperator{\Ann}{Ann}
\DeclareMathOperator{\Zar}{Zar}


%%%%%%%%%%%%
%%%%% SGA4-2
\DeclareMathOperator{\Sex}{Sex}
\DeclareMathOperator{\vois}{vois}
\DeclareMathOperator{\Vois}{Vois}
\DeclareMathOperator{\Ext}{Ext}
\DeclareMathOperator{\ext}{ext}
\DeclareMathOperator{\Fer}{Fer}
\DeclareMathOperator{\Add}{Add}
\DeclareMathOperator{\Cocart}{Cocart}
\DeclareMathOperator{\cosk}{cosk}
\DeclareMathOperator{\plat}{plat}

\DeclareMathOperator{\faisc}{faisc}
\DeclareMathOperator{\sk}{sk}
\DeclareMathOperator{\ch}{ch}

\def\HH{\mathbf{H}}
\def\qc{\mathrm{qc}}
\def\et{\mathrm{et}}
\def\coh{\mathrm{coh}}

\def\Hvr#1{%
\begin{array}[b]{l}
{\kern2pt\scriptstyle r}\\[-8pt]
{\kern1pt\scriptstyle \vee}\\[-5pt]
H^{#1}
\end{array}}

\def\Hvv#1{%
\begin{array}[b]{l}
\kern2pt\rotatebox{90}{${\scriptscriptstyle\ll}$}\\[-6pt]
H^{#1}
\end{array}}

\def\Hvi#1{%
\begin{array}[b]{l}
{\kern0.2pt\scriptscriptstyle \infty}\\[-8pt]
{\kern1pt\scriptstyle \vee}\\[-5pt]
H^{#1}
\end{array}}

\def\Hv#1{%
\begin{array}[b]{l}
{\kern1pt\scriptstyle \vee}\\[-5pt]
H^{#1}
\end{array}}

\DeclareMathOperator{\cart}{cart}
\DeclareMathOperator{\Ffaisc}{\mathbf{Faisc}}


\def\Ev#1{%
\begin{array}[b]{l}
{\kern1pt\scriptstyle \vee}\\[-5pt]
E'_{#1}
\end{array}}

\DeclareMathOperator{\Cont}{Cont}
\DeclareMathOperator{\Ccont}{\mathscr{C}\!\mathit{ont}}
\DeclareMathOperator{\Site}{Site}
\newcommand\MM{\mathbf{M}}
\def\limtop#1{{\displaystyle{\varprojlim_{#1}\!\mbox{\rm top }}}}
\def\Limtop{\mathop{\varprojlim\!\mbox{\rm top}}}
\DeclareMathOperator{\spec}{spec}
\def\ss{\scriptstyle{\mathscr{S}}}
\DeclareMathOperator{\Gal}{Gal}
\def\smallm{\scriptstyle{\mathscr{M}}}
\DeclareMathOperator{\End}{End}
\DeclareMathOperator{\ass}{ass}
\DeclareMathOperator{\zar}{zar}
\DeclareMathOperator{\Ogg}{Ogg}
\DeclareMathOperator{\res}{res}
\DeclareMathOperator{\tr}{tr}
\DeclareMathOperator{\Supp}{Supp}
\DeclareMathOperator{\codim}{codim}
\DeclareMathOperator{\rk}{rk}
\DeclareMathOperator{\cur}{cur}
\DeclareMathOperator{\cdqc}{cdqc}
\newcommand\DD{\mathbf{D}}
\DeclareMathOperator{\cl}{cl}
\def\JJ{\mathbf{J}}
\def\FF{\mathbf{F}}
\def\BB{\mathbf{B}}
\DeclareMathOperator{\Alget}{Alget}
\DeclareMathOperator{\Ass}{Ass}
\DeclareMathOperator{\supp}{supp}
\DeclareMathOperator{\rad}{rad}
\DeclareMathOperator{\prof}{prof}
\DeclareMathOperator{\cint}{int}
\DeclareMathOperator{\sic}{sic}
\DeclareMathOperator{\Tor}{Tor}
\DeclareMathOperator{\tors}{tors}
\DeclareMathOperator{\Sym}{Sym}
\DeclareMathOperator{\parf}{parf}
\DeclareMathOperator{\tor}{tor}
\DeclareMathOperator{\defn}{defn}
\DeclareMathOperator{\Div}{Div}
\DeclareMathOperator{\Pol}{Pol}
\DeclareMathOperator{\Fex}{Fex}
\DeclareMathOperator{\Lie}{Lie}
\DeclareMathOperator{\tot}{tot}
\DeclareMathOperator{\Aut}{Aut}
\DeclareMathOperator{\HOM}{HOM}
\DeclareMathOperator{\pch}{pch}
\DeclareMathOperator{\Ch}{Ch}
\DeclareMathOperator{\EXT}{EXT}
\DeclareMathOperator{\PIC}{PIC}
\DeclareMathOperator{\Rf}{Rf}
\DeclareMathOperator{\inj}{inj}
\newcommand\VV{\mathbf{V}}

\def\Xv{%
\begin{array}[b]{l}
{\kern1pt\scriptstyle \vee}\\[-5pt]
X
\end{array}}

%%%%%%%
\newcommand\TT{\mathbf{T}}
\newcommand\XX{\mathbf{X}}
\newcommand\EE{\mathbf{E}}
%%%%%%

\def\qedsymbol{\textup{C.Q.F.D.}}
\let\CQFD\qed
\let\cqfd\qed

%%%%%\bigdot
\font\bigsymb=cmsy10 at 4pt
\def\bigdot{{\kern1.2pt\raise 1.5pt\hbox{\bigsymb\char15}}}
%%%%%

\def\Lright{\mathop{\vcenter{\hbox{\hbox{$\longrightarrow$}\lower8pt \hbox{\kern-16pt\hbox{$\longrightarrow$}}}}}}

\def\Aa{{\mathcal A}}
\def\Bb{{\mathcal B}}
\def\Cc{{\mathcal C}}
\def\Dd{{\mathcal D}}
\def\Ee{{\mathcal E}}
\def\Ff{{\mathcal F}}
\def\Gg{{\mathcal G}}
\def\Hh{{\mathcal H}}
\def\Ii{{\mathcal I}}
\def\Jj{{\mathcal J}}
\def\Kk{{\mathcal K}}
\def\Ll{{\mathcal L}}
\def\Mm{{\mathcal M}}
\def\Nn{{\mathcal N}}
\def\Oo{{\mathcal O}}
\def\Pp{{\mathcal P}}
\def\Qq{{\mathcal Q}}
\def\Rr{{\mathcal R}}
\def\Ss{{\mathcal S}}
\def\Tt{{\mathcal T}}
\def\Uu{{\mathcal U}}
\def\Vv{{\mathcal V}}
\def\Ww{{\mathcal W}}
\def\Xx{{\mathcal X}}
\def\Yy{{\mathcal Y}}
\def\Zz{{\mathcal Z}}

\newcommand{\CC}{\mathbf{C}}
\newcommand{\GG}{\mathbf{G}}
\newcommand{\KK}{\mathbf{K}}
\newcommand{\LL}{\mathbf{L}}
\newcommand{\NN}{\mathbf{N}}
\newcommand{\PP}{\mathbf{P}}
\newcommand\QQ{\mathbf{Q}}
\newcommand\RR{\mathbf{R}}
\newcommand\ZZ{\mathbf{Z}}

\def\to{\mathchoice{\longrightarrow}{\rightarrow}{\rightarrow}{\rightarrow}}
\def\from{\mathchoice{\longleftarrow}{\leftarrow}{\leftarrow}{\leftarrow}}
\def\mto{\mathchoice{\longmapsto}{\mapsto}{\mapsto}{\mapsto}}
\def\mfrom{\mathrel{\from\joinrel\mapstochar\,}}
\def\hto{\mathrel{\lhook\joinrel\to}}
\def\hfrom{\mathrel{\from\joinrel\rhook}}

\def\implies{\mathchoice{\Longrightarrow}{\Rightarrow}{\Rightarrow}{\Rightarrow}}
\def\implied{\mathchoice{\Longleftarrow}{\Leftarrow}{\Leftarrow}{\Leftarrow}}
\def\iff{\mathchoice{\Longleftrightarrow}{\Leftrightarrow}{\Leftrightarrow}{\Leftrightarrow}}

\def\To#1{\mathchoice{\xrightarrow{\labelstyle~~#1~}}{\stackrel{#1}{\longrightarrow}}{}{}}
\def\Mto#1{\mathrel{\mapstochar\To{#1}}}
\def\From#1{\mathchoice{\xleftarrow{\labelstyle~~#1~}}{\stackrel{#1}{\longleftarrow}}{}{}}

\newcommand{\isomto}{\overset{\sim}{\longrightarrow}}
\newcommand{\isomfrom}{\overset{\sim}{\longleftarrow}}

\newcommand{\IFF}{\unskip~$\iff$~\ignorespaces}
\newcommand{\THEN}{\unskip~$\implies$~\ignorespaces}

\let\leq\leqslant
\let\geq\geqslant
\let\le\leqslant
\let\ge\geqslant

\def\cf{{cf.}\kern.3em}
\def\Cf{{Cf.}\kern.3em}
\def\ie{{i.e.}\kern.3em}
\def\iev{{i.e.},\xspace}
%\def\loccito{loc.\kern2pt cit.\xspace}
%\def\loccitn{\textit{loc.\kern2pt cit.}\xspace}
%\def\loccit{\sisi{\loccito}{\loccitn}}
\def\EGA{\textup{EGA}\kern.2em}
\def\SGA{\textup{SGA}\kern.3em}

\let\oldS\S
\def\S{\oldS\kern .15em}
\def\ptbl{.\kern .2em}
\def\resp{resp.\kern.3em}
\def\Exp{Exp.\kern.3em}

\let\bar\overline
\let\hat\widehat
\let\tilde\widetilde
\let\dpl\displaystyle
\let\ldots\dots
\let\emptyset\varnothing
\let\skpt\skippointrait
\let\epsilon\varepsilon

\def\sfrac#1#2{{#1}/{#2}}

\DeclareMathOperator{\cd}{cd}
%\DeclareMathOperator{\Hhom}{\mathcal{H}\!\mathit{om}}
\DeclareMathOperator{\Hhom}{\mathbf{Hom}}
%\DeclareMathOperator{\Eext}{\mathcal{E}\!\mathit{xt}}
\DeclareMathOperator{\Eext}{\mathbf{Ext}}
\DeclareMathOperator{\Spec}{Spec}
\DeclareMathOperator{\Tr}{Tr}

\let\emptyset=\varnothing

%%%%%%%%%%%%%%%
%%% expose XVII
%%%%%%%%%%%%%%%

\makeatletter
\let\old@optexpose\optexpose
\newcommand{\empty@optexpose}[1]{}
\newcommand{\paren@optexpose}[1]{\old@optexpose{#1}(}
\newcommand{\refnoexpose}[1]{\let\optexpose\empty@optexpose\ref{#1}\let\optexpose\old@optexpose}%\refnoexpose{X.cor5.2} engendre "5.2" au lieu de "X 5.2"
\newcommand{\neweqref}[1]{\let\optexpose\paren@optexpose\textup{\ref{#1})}\let\optexpose\old@optexpose}%\neweqref{XII.eq4.1} engendre "XII (4.1)" au lieu de "(XII 4.1)"
\makeatother
\newcommand{\refvb}[2][VI${}_{\text{B}}$ ]{\sisi{#1\refnoexpose{#2}}{\ref{#2}}}

\newcommand{\leftexp}[2]{{\vphantom{#2}}^{#1}{#2}}
\newcommand{\AAA}{\mathbf{A}}
\newcommand{\Sf}{\mathfrak{S}}
\DeclareMathOperator*{\quotprojlim}{\text{``$\varprojlim$''}}
\DeclareMathOperator*{\quotinjlim}{\text{``$\varinjlim$''}}
\DeclareMathOperator{\Fl}{F\sisi{\ell}{l}}
\DeclareMathOperator{\Pl}{P\sisi{\ell}{l}}
%\DeclareMathOperator{\cl}{c\sisi{\ell}{l}}
\DeclareMathOperator{\tdf}{tdf}
\DeclareMathOperator{\rg}{rg}
\newcommand{\iInd}{\sisi{I}{i}nd}
\newcommand{\iind}{ind}
\newcommand{\ppro}{pro}
\newcommand{\tTor}{\sisi{T}{t}or}
\newcommand{\ttor}{tor}
\newcommand{\chap}{Chap.}
\newcommand{\Ffex}{\Ff\mathit{ex}}
\newcommand{\sheafk}{\underline{k}}
\newcommand{\sheafn}{\underline{n}}
\newcommand{\sheafGamma}{\underline{\Gamma}}

% Repeated references to an already typeset footnote must not use
% \footnotemark without a matching \footnotetext: Hyperref would then create
% a live annotation to a destination that cannot exist.  This helper keeps the
% current printed footnote mark and links it to an explicit stable note target.
\makeatletter
\newcommand{\sgarepeatedfootnotemark}[1]{%
  \begingroup
  \edef\@thefnmark{\thefootnote}%
  \hyperlink{#1}{\@makefnmark}%
  \endgroup}
\makeatother

\sisi{
\DeclareMathOperator{\simarrow}{\rule{.15mm}{1.8mm}\kern -3.4pt\rightsquigarrow}
\newcommand{\sheafell}{\underline{\ell}}
\newcommand{\sheafZZ}{\underline{\ZZ}}
}{
\newcommand{\nderef}[1]{la N.D.E.~(\ref{#1}) page~\pageref{#1}}
}

%%%%%%%%%%%%%%%%%
%%%%% expose XVII
%%%%%%%%%%%%%%%%%

% END INCLUDED SGA4_English_current_buildable_source_20260906/source/fichiers-configuration-latex/macros.tex


%\externaldocument{../01/01.tex}
\externaldocument{../02/01.tex}
\externaldocument{../03/03.tex}
\externaldocument{../04/04.tex}

\setboolean{original}{false}

\def\labelenumi{(\theenumi)}
\def\labelenumii{(\theenumii)}
\def\labelenumiii{(\theenumiii)}
\def\labelenumiv{(\theenumiv)}

\begin{document}

\begin{center}
I. Presheaves\\
A. Grothendieck and J.-L. Verdier
\end{center}

\selectlanguage{english}version: \showgitstatus

\tableofcontents

\else
\chapter{Presheaves}
\begin{center}
A. Grothendieck and J.-L. Verdier
\end{center}
\fi
\phantomsection\MakeLinkTarget*{sga4.exp.I}\label{sga4.exp.I}

% BEGIN INCLUDED SGA4_English_current_buildable_source_20260906/source/01/drafts/SGA4_I_LINES_46_100_English.texfrag
In\pageoriginale numbers 0 through 5 of this exposé, we present the
elementary, and most often well-known, properties of categories of
presheaves\footnote[1]{The reader may also consult SGA 3 I \S\S 1 to 3.}.
These properties are used constantly throughout the rest of the seminar,
and knowledge of them is essential for understanding the subsequent
exposés. The proofs are immediate; most often, they are omitted. In
numbers 6 through 9, we take up a few themes used on various occasions
later. The reader in a hurry %%
may omit them on a first reading. Number 10 fixes the %%
terminology used. Appendix 11 is %%
due to N. Bourbaki.


\setcounter{section}{-1}
\section{Universes}\etiquette[0]{I.sec0}

A \emph{universe} is a nonempty set\nde{This definition contradicts
Bourbaki's appendix \textit{below}. However, the usefulness of empty
universes seems debatable...} $𝒰$ having the following properties:
\begin{itemize}
\item[($𝒰$ 1)] If $x\in 𝒰$ and $y\in x$, then $y\in 𝒰$.

\item[($𝒰$ 2)] If $x$, $y\in 𝒰$, then $\{x,y\}\in 𝒰$.

\item[($𝒰$ 3)] If $x\in 𝒰$, then $\Pp(x)\in 𝒰$.

\item[($𝒰$ 4)] If $(x_i, i\in I\sisi{\in 𝒰}{})$ is a
  family of elements of $𝒰$\sisi{}{ and $I\in𝒰$}, then $\cup_{i\in I}x_i \in
  𝒰$.
\end{itemize}

From\pageoriginale the preceding axioms one readily deduces the following
properties:
\begin{itemize}
\item If $x\in 𝒰$, the set $\{x\}$ belongs to $𝒰$.

\item If $x$ is a subset of $y\in 𝒰$, then $x\in 𝒰$.

\item If $x$, $y\in 𝒰$, the ordered pair $(x,y)=\{\{x\},\{x,y\}\}$ (Kuratowski's
  definition) is an element of $𝒰$.

\item If $x$, $y\in 𝒰$, the union $x\cup y$ and %%
the product $x\times
   y$ are elements of $𝒰$.

\item If $(x_i, i\in I\in 𝒰)$ is a family
  of elements of $𝒰$,
   the product $\prod_{i\in I}x_i$ is an element of $𝒰$.

\item If $x\in 𝒰$, then $\card(x)< \card(𝒰)$ (strictly). In
   particular, the relation $𝒰\in 𝒰$ does not hold.
\end{itemize}

% END INCLUDED SGA4_English_current_buildable_source_20260906/source/01/drafts/SGA4_I_LINES_46_100_English.texfrag


% BEGIN INCLUDED SGA4_English_current_buildable_source_20260906/source/01/drafts/SGA4_I_LINES_102_196_English.texfrag
We can therefore carry out all the usual operations of set theory
starting from the elements of a universe without thereby causing the
final result to cease to be an element of the universe.

The first advantage of the notion of a universe is that it provides a
definition of the usual categories: the category of sets belonging to
the universe $𝒰$ ($𝒰$-$\Ens$), the category of
topological spaces belonging to the universe $𝒰$, the category of
commutative groups belonging to the universe $𝒰$
($𝒰$-$\Ab$), the category of categories belonging to
the universe $𝒰$\ldots.

However, the only known universe is the set of symbols of the form
$\{\{∅\}, \{\{∅\}, ∅\}\}$ \etc. (all
the elements of this universe are finite sets, and this universe
is countable). In particular, no universe is known that
contains an element of infinite cardinality. We are therefore led
to add
to the axioms of set theory the axiom: $(𝒰 A)$
\emph{For every set $x$ there exists a universe $𝒰$ such
that $x\in
  𝒰$.}

Since the intersection\pageoriginale of a family of universes is a
universe, it follows immediately that every set is
an element
of a smallest universe. One can show that axiom $(𝒰
A)$ is independent of the axioms of set theory.

We shall also add the axiom:


\noindent \emph{$(𝒰 B)$ Let $R\{x\}$ be a relation and let
$𝒰$ be a universe. If there
  exists an element $y\in 𝒰$ such that $R\{y\}$, then
  $\tau_xR\{x\}\in 𝒰$.}

The consistency of axioms $(𝒰 A)$ and $(𝒰 B)$
relative to the other axioms of set theory has not been
proved and, it seems, cannot be proved.

Let $𝒰$ be a universe and let $c(𝒰)$ be the supremum
of the cardinalities of the elements of $𝒰(c(𝒰)\leq
\card(𝒰))$. The cardinal $c(𝒰)$ has the following
properties:

$(FI)$ If $a<c(𝒰)$, then $2^a<c(𝒰)$.

$(FII)$ If $(a_i, i\in I)$ is a family of cardinals strictly
less than $c(𝒰)$ and if $\card(I)$ is %%
strictly less than $c(𝒰)$, then $\Sigma_{i\in I} a_i <
c(𝒰)$.

The cardinals that have properties $(FI)$ and $(FII)$ are
called \emph{strongly inaccessible cardinals.}

The cardinal 0 and the countably infinite cardinal are strongly
inaccessible.

Axiom $(𝒰 A)$ implies:

\emph{$(𝒰 A')$ Every cardinal is strictly less than a
  strongly inaccessible cardinal.}

Conversely, one can show (11) that the consistency of
$(UA')$ implies the consistency of $(UA)$, and that the
consistency of these axioms entails that of axiom
$(𝒰 B)$.

Call\pageoriginale an \emph{Artinian set} any set $E$
for which there is no infinite family $(x_n, n\in N)$ such
that $x_0 \in E$, $x_{n+1}\in x_n$. One can then show
[\loccit] that there is a one-to-one correspondence between the
strongly inaccessible cardinals and the Artinian universes, defined
as follows: to every strongly inaccessible cardinal $c$ we
associate the unique Artinian universe $U_c$ such that
$$
\card(𝒰_c)=c{\quoi}.
$$
Let $c<c'$ be two strongly inaccessible cardinals. We have:
$$
𝒰_c\in 𝒰_{c'}{\quoi}.
$$
In particular, the Artinian universes of cardinality less than a
given cardinal form a well-ordered set (with respect to the
membership relation). Axiom $(𝒰 A')$ is equivalent to
the axiom:

$(𝒰 A'')$. \emph{Every Artinian set is an element of an
Artinian
  universe.}

Observe that all the usual sets $(\ZZ, \QQ,\ldots)$ are
Artinian sets.

% END INCLUDED SGA4_English_current_buildable_source_20260906/source/01/drafts/SGA4_I_LINES_102_196_English.texfrag


% BEGIN INCLUDED SGA4_English_current_buildable_source_20260906/source/01/drafts/SGA4_I_LINES_197_410_English.texfrag
\section{$𝒰$-categories. Presheaves of sets}\etiquette[1]{I.sec1}

\setcounter{subsection}{-1}
\subsection{}\etiquette[1.0]{I.sec1.0}
\emph{Throughout the remainder of the seminar and unless expressly stated
  otherwise, the universes under consideration will contain an element
  of infinite cardinality.} Let $𝒰$ be a universe. A set is said to be
  $𝒰$-\emph{small} (or, when no confusion can result,
  \emph{small}) if it is \emph{isomorphic} to an element of
  $𝒰$. We shall also use the terminology: small group, small ring,
  small category\nde{A category $C$ is regarded as a set of arrows
  (equipped with the subset of objects, regarded as the set of identity
  arrows, and with the ``source, target'' maps). Thus, the expressions ``$C$ is an element of $𝒰$'' or
  ``$C$ is $𝒰$-small'' make sense.}
  \ldots We shall often assume, without
  explicit mention, that the schemes, topological spaces, indexing
  sets\ldots with which we work are $\in 𝒰$,
or\pageoriginale at least have cardinality $\in 𝒰$;
  nevertheless, many of the categories with which we shall work
  will not be $\in 𝒰$.

\begin{enonce*}{Definition 1.1}\etiquette[1.1]{I.def1.1}
Let $𝒰$ be a universe and $C$ a category. We say that $C$
is a $𝒰$-category if, for every pair $(x,y)$ of objects
of $C$, the set $\Hom_C(x,y)$ is $𝒰$-small.
\end{enonce*}

\setcounter{subsection}{1}
\subsubsection{}\etiquette[1.1.1]{I.sec1.1.1}
Let $C$ and $D$ be two categories and let $\text{Fonct}(C,D)$ be the
category of functors from $C$ to $D$. The following assertions
are immediate:
\begin{enumerate}
\item[a)] If $C$ and $D$ are elements of a universe $𝒰$
  (\resp  are $𝒰$-small), the category $\text{Fonct}(C,D)$ is an element
  of $𝒰$ (\resp  is $𝒰$-small).

\item[b)] If $C$ is $𝒰$-small and $D$ is a $𝒰$-category,
  $\text{Fonct}(C,D)$ is a $𝒰$-category.
\end{enumerate}

\begin{enonce*}[remark]{Remark 1.1.2}\etiquette[1.1.2]{I.rem1.1.2}
Let $D$ be a category having the following properties:
\begin{itemize}
\item[(C1)] The set $\ob(D)$ is contained in the universe $𝒰$.

\item[(C2)] For every pair $(x,y)$ of objects of $D$, the set
  $\Hom_D(x,y)$ is an element of $𝒰$.
\end{itemize}
(The usual categories constructed from a universe
$𝒰$ have these two properties: $𝒰$-$\Ens$,
$𝒰$-$\Ab$,\ldots). Let $C$ be a category belonging to
$𝒰$. Then the category $\text{Fonct}(C,D)$ does not
in general have properties (C1) and (C2). For example, the category
$\text{Fonct}(C,𝒰$-$\Ens$) has neither property
(C1) nor property (C2). This justifies the adopted definition of
$𝒰$-category, in preference to the more restrictive notion given
by conditions (C1) and (C2) above.
\end{enonce*}

\begin{enonce*}{Definition 1.2}\etiquette[1.2]{I.def1.2}
Let\pageoriginale $C$ be a category. The category of %%
presheaves of sets on $C$ relative to the universe $𝒰$
(or, when no confusion can result, the category of
presheaves on $C$) is the category of contravariant functors on
$C$ with values in the category of $𝒰$-sets.
\end{enonce*}

We denote by $\widehat{C}_{𝒰}$ (or more
simply, when no confusion can result, by
$\widehat{C}$\,) the category of presheaves
of sets on $C$ relative to the universe $𝒰$. The objects of
$\widehat{C}_{𝒰}$ are called
$𝒰$-presheaves (or simply presheaves) on $C$.
When $C$ is $𝒰$-small, the category
$\widehat{C}_{𝒰}$ is a
$𝒰$-category. When $C$ is a $𝒰$-category,
$\widehat{C}_{𝒰}$ is not necessarily a
$𝒰$-category.


\begin{enonce*}[remark]{Construction-Definition
    1.3}\etiquette[1.3]{I.condef1.3}
Let $x$ be an object of a $𝒰$-category $C$. The
$𝒰$-\emph{functor represented by} $x$ is the functor
$h_{𝒰}(x):C^\circ\rightarrow 𝒰$-$\Ens$ whose
construction is as follows \footnote[1]{$C^\circ$ will always denote the
category opposite
  to $C$.}. Let $y$ be an object of $C$.
\begin{enumerate}
\item[a)] If $\Hom_C(y,x)$ is an element of $𝒰$, then:
      $$
      h_{𝒰}(x)(y)=\Hom_C(y,x)
      $$

\item[b)] Suppose that $\Hom_C(y,x)$ is not an element of $𝒰$
  and let $R(Z,x,y)$ be the relation: ``The set $Z$ is the codomain of an
  isomorphism $\Hom_C(y,x)\xrightarrow{\sim}Z$''. We then set:
  $$
  h_{𝒰}(x)(y)=\tau_ZR(Z){\quoi}.
  $$
\end{enumerate}
(By axiom $(𝒰 B)$, $h_{𝒰}(x)(y)$ is an
element of $𝒰$). Let $R'(u,x,y)$ be the relation: « $u$ is
a bijection
$$
u:\Hom_C(y,x)\xrightarrow{\sim} h_{𝒰}(x)(y)\mbox{
»}{\quoi}.
$$
We\pageoriginale then set:
$$
\varphi(y,x)=\tau_uR'(u){\quoi}.
$$

It should be noted that, in cases (a) and (b), we have $a$ canonical
isomorphism:
$$
\varphi(y,x):\Hom_C(y,x)\xrightarrow{\sim}
h_{𝒰}(x)(y){\quoi}.
$$
(In case a) $\varphi(y,x)$ is the identity). Now let
$u:y\rightarrow y'$ be an arrow of $C$. By composition of morphisms, the
morphism $u$ defines a map:
$$
\Hom_C(u,x):\Hom_C(y',x)\longrightarrow \Hom_C(y,x){\quoi}.
$$
We then set
$$
h_{𝒰}(x)(u)=\varphi(y,x)\Hom_C(u,x)\varphi(y',x)^{{}-1}{\quoi}.
$$

It is immediate that $h_{𝒰}(x)$, so defined, is
a functor $C^\circ\rightarrow 𝒰$-$\Ens$.
\end{enonce*}

\setcounter{subsection}{3} \setcounter{subsubsection}{0}
\subsubsection{}\etiquette[1.3.1]{I.sec1.3.1}
Similarly, using the isomorphisms $\varphi$, for every
morphism $y:x\rightarrow x'$ we define a morphism of functors:
$$
h_{𝒰}(y):h_{𝒰}(x)\longrightarrow h_{𝒰}(x')
$$
and it is immediate that this defines a functor
$$
h_{𝒰}:C\longrightarrow
\widehat{C}_{𝒰}{\quoi}.
$$

\subsubsection{}\etiquette[1.3.2]{I.sec1.3.2}
Now let $𝒱$ be a universe such that $𝒰\subset
𝒱$. There is then a canonical fully faithful functor:
$$
𝒰\text{-}\Ens\hookrightarrow 𝒱\text{-}\Ens
$$
whence a fully faithful functor:
$$
\widehat{C}_{𝒰}\hookrightarrow
\widehat{C}_{𝒱}
$$
and\pageoriginale the diagram:

%\def\objectstyle{\mathrm}--here
\[
\xymatrix@R=1.75cm@C=3.5cm@M=.2cm{
C\ar[r]^{h_{\underline{U}}} \ar[dr]_{h_{\underline{V}}}& C{\widehat{\phantom{i}}}_{\underline{U}}\ar@{_{(}->}[d]\\
& C\widehat{\phantom{i}}_{\underline{V}} }
\]
commutes up to canonical isomorphism. When $C$ is an
element of $𝒰$, this canonical isomorphism is the identity.

\subsubsection{}\etiquette[1.3.3]{I.sec1.3.3}
\emph{In practice, the universe $𝒰$ is fixed
  once and for all and is not mentioned.} We then use the
  notations $\widehat{C}$ (for the category of
  $𝒰$-presheaves of sets) and
$$
h:C\longrightarrow \widehat{C}{\quoi}.
$$
For every object $x$ of $C$, the presheaf $h(x)$ is called the
\emph{presheaf represented by} $x$, and we shall always identify
the value at $y\in \ob(C)$ of the presheaf $h(x)$ with $\Hom_C(y,x)$.

\begin{enonce*}{Proposition 1.4}\etiquette[1.4]{I.prop1.4}
Let $C$ be a $𝒰$-category, $F$ a presheaf on $C$, and
$X$ an object of $C$. There is an isomorphism functorial in $X$ and
in $F$:
$$
i:\Hom_{\widehat{C}}(h(X),F)\xleftarrow{\sim}
F(X){\quoi},
$$
where the map $i$, when $F$ is of the form $h(Y)$, is none other
than the map:
$$
h:\Hom(h(X), h(Y))\longleftarrow \Hom(X,Y){\quoi}.
$$
In particular, the functor $h$ is fully faithful.
\end{enonce*}

\setcounter{subsection}{4} \setcounter{subsubsection}{0}
\subsubsection{}\etiquette[1.4.1]{I.sec1.4.1}
This proposition justifies the customary abuse of language identifying
an object of $C$ with the corresponding contravariant functor. A
presheaf\pageoriginale isomorphic to an object in the image of $h$ (or,
using the abuse of language noted above, isomorphic to an object
of $C$) is called a \emph{representable presheaf.}

\subsubsection{}\etiquette[1.4.2]{I.sec1.4.2}
Let $𝒱$ be a universe containing the universe $𝒰$. The
category of $𝒰$-presheaves is a \emph{full} subcategory
of the category of $𝒱$-presheaves, and consequently
a $𝒰$-presheaf is representable if and only
if its image in the category of $𝒱$-presheaves is a
representable $𝒱$-presheaf.

% END INCLUDED SGA4_English_current_buildable_source_20260906/source/01/drafts/SGA4_I_LINES_197_410_English.texfrag


% BEGIN INCLUDED SGA4_English_current_buildable_source_20260906/source/01/drafts/SGA4_I_LINES_411_824_English.texfrag
\section{Projective and inductive limits}\etiquette[2]{I.sec2}

Let $C$ be a $𝒰$-category and $I$ a small category.
Denote by $\Fonct(I,C)$ the category of functors from $I$ to $C$. The
category $\Fonct(I,C)$ is a $𝒰$-category. To every object
$X$ of $C$, associate the subcategory $𝒳$ of $C$ whose
only object is $X$ and whose only arrow is the identity of $X$.
Let $i_X:𝒳\rightarrow C$ denote the inclusion functor.
There is a unique functor $e_X:I\rightarrow 𝒳$, and
we denote by $k_X:I\rightarrow C$ the functor $i_X\circ
e_X$. (We say that $k_X$ is the constant functor with value $X$.) The
correspondence $X\mapsto k_X$ is clearly functorial in $X$,
which allows us to define, for every functor $G:I\rightarrow
C$, the presheaf $X\mapsto \Hom_{\Fonct(I,C)}(k_X,G)$.

\begin{enonce*}{Definition 2.1}\etiquette[2.1]{I.def2.1}
The projective limit of $G$, denoted by $\varprojlim_I G$, is the
presheaf:
$$
X\longmapsto \Hom_{\Fonct(I,C)}(k_X,G){\quoi}.
$$
When the presheaf $\varprojlim_I G$ is representable, we
also denote by $\varprojlim_I G$ an\pageoriginale object of $C$
representing it. Thus the object $\varprojlim_I G$ is defined only up to
isomorphism. When no confusion can result, we use
the abbreviated notation $\varprojlim G$.
\end{enonce*}

As $G$ varies, $\varprojlim G$ is a functor from the category
$\Fonct(I,C)$ to $\widehat{C}$.

\setcounter{subsection}{1} \setcounter{subsubsection}{0}
\subsubsection{}\etiquette[2.1.1]{I.sec2.1.1}
By symmetry (reversing the direction of the arrows
in $C$), one similarly defines the inductive limit of a functor: it is a
covariant functor on $C$ with values in the category of
$𝒰$-sets. We shall use the notations $\varinjlim_I
G$ or $\varinjlim G$.

Observe that products, fiber products, and kernels are projective
limits. Similarly, sums, amalgamated sums, and cokernels are
%%
inductive limits.

\begin{enonce*}{Definition 2.2}\etiquette[2.2]{I.def2.2}
Let $I$ and $C$ be two categories and $G:I\rightarrow C$ a functor.
The projective limit of $G$ is said to be representable if there exists
a universe $𝒰$ such that:
\begin{enumerate}
\renewcommand{\labelenumi}{\theenumi)}
\item The category $I$ is $𝒰$-small.

\item The category $C$ is a $𝒰$-category.

\item The presheaf $\varprojlim G$, with values in the
  category of $𝒰$-sets, is representable.
\end{enumerate}
\end{enonce*}

It follows from ${\rm n^\circ}$ 1 that the object $\varprojlim G$
representing the presheaf $\varprojlim G$ is independent, up to
isomorphism, of the universe $𝒰$. Note also that there
exists a smallest universe $𝒰'$ having properties
(1) and (2), and that the presheaf $\varprojlim G$ necessarily takes
values in the category of $𝒰'$-sets.

\setcounter{subsection}{2} \setcounter{subsubsection}{0}
\subsubsection{}\etiquette[2.2.1]{I.sec2.2.1}
Let\pageoriginale $I$ and $C$ be two categories. We say that
\emph{the $I$-projective limits in $C$ are
  representable} if, for every functor $G:I\rightarrow C$, the
  projective limit of $G$ is representable. Finally, let $C$ be a
  category and $𝒰$ a universe. We say that the
  $𝒰$-\emph{projective
  limits in $C$ are representable} if, for every
  category $I$ that is $𝒰$-small and every functor $G:I\rightarrow
  C$, the projective limit of $G$ is representable.

\begin{enonce*}{Proposition 2.3}\etiquette[2.3]{I.prop2.3}
Let $C$ be a category and $𝒰$ a universe\sisi{ that is
nonempty}{}. The following assertions are equivalent:
\begin{enumerate}
\renewcommand{\theenumi}{\roman{enumi}}
\renewcommand{\labelenumi}{\theenumi)}
\item The $𝒰$-projective limits in $C$ are representable.

\item Products indexed by a small set are
  representable, and kernels of pairs of arrows are
  representable.
\item Products indexed by a small set are
  representable, and fiber products are representable.
\end{enumerate}
\end{enonce*}

\begin{enonce*}[remark]{Proof}
It suffices to observe that there is an isomorphism, functorial in
$G$,
$$
\varprojlim G=\Ker (\prod_{i\in \ob(I)} %%
G(i)\rightrightarrows \prod_{u\in \rm{F\ell}(I)} G(\but(u))){\quoi},
$$
where the pair of arrows is defined by the morphisms
\begin{align*}
\prod_{i\in \Ob(I)}  G(i) %%
 & \xrightarrow{\pr_{\but(u)}}{} G(\but(u))\\
\prod_{i\in \Ob(I)} G(i) %%
 &\xrightarrow{\pr_{\text{source}(u)}}{} G(\text{source
}(u))\xrightarrow{G(u)} G(\but(u)){\quoi}.
\end{align*}
Moreover, it is clear that $\Ker (X  \displaystyle \Lright^{u}_{v}Y)=
X\prod_{X\times Y} X$, where the two morphisms from $X$ to $X\times Y$
are $\id_X \times u$ and $\id_X\times v$.
\end{enonce*}

\setcounter{subsection}{3} \setcounter{subsubsection}{0}
\subsubsection{}\etiquette[2.3.1]{I.sec2.3.1}
There are of course analogous definitions and assertions for
inductive limits, which we shall not spell out. The same applies to
finite projective\pageoriginale and inductive limits (\ie
those relative to a %%
finite category $I$).


\begin{enonce*}{Corollary 2.3.2}\etiquette[2.3.2]{I.cor2.3.2}
Denote by $𝒰$-$\Ens$ the category of
$𝒰$-sets and by $𝒰$-$\Ab$ the category of
abelian group objects in $𝒰$-$\Ens$. The
$𝒰$-projective and inductive limits in
$𝒰$-$\Ens$ and in $𝒰$-$\Ab$ are representable.
\end{enonce*}

\begin{enonce*}{Proposition 2.3.3}\etiquette[2.3.3]{I.prop2.3.3}
Let $I$ be a small category, $F:I\rightarrow 𝒰$-$\Ens$
a functor, and \linebreak %%
$G\subset\ob(I)$ a small set of objects of $I$ such that, for every
object $X$ of $I$, there exists an object $Y\in G$ and a morphism
$Y\rightarrow X$ (\resp $X\rightarrow Y$). Then $\varprojlim_I F$
(\resp $\varinjlim_I F$) is representable, and
$\card(\varprojlim_I F)\leq \Pi_{Y\in G}\card(F(Y))$ (\resp
$\card(\varinjlim_I F)\leq \Sigma_{Y\in G}\card(F(Y))$).
\end{enonce*}

\begin{enonce*}[remark]{Proof}
One immediately verifies that $\varprojlim_I F$ (\resp $\varinjlim_I
F$) is isomorphic to a subobject (\resp a quotient) of $\Pi_{Y\in
G}F(Y)$ (\resp $\amalg_{Y\in G}F(Y)$); hence the proposition.
\end{enonce*}

\begin{enonce*}{Definition 2.4.1}\etiquette[2.4.1]{I.def2.4.1}
Let $C$ be a category in which finite projective (\resp
inductive) limits are representable, and let $u:C\rightarrow C'$ be a
functor. We say that $u$ is left exact (\resp right exact) if it
``commutes'' with finite projective (\resp inductive) limits.
A functor that is both left and right exact is called an exact functor.
\end{enonce*}

\setcounter{subsection}{4} \setcounter{subsubsection}{1}
\subsubsection{}\etiquette[2.4.2]{I.sec2.4.2}
It follows from \hyperref[I.prop2.3]{2.3} that, for a functor to be exact\sisi{}{ on the left}, it is necessary and %% étiquette !
sufficient that it take the terminal object (= empty product) to the
terminal object, the product of two objects to the product of their
images, and the kernel of a pair of arrows to the kernel of the
image pair; equivalently, that it take the terminal object to the
terminal object and fiber products to fiber products\sisi{. (W}{ (w}e
assume that in $C$ finite $\varprojlim$ are
representable\sisi{.)}{).}

\setcounter{subsection}{5} \setcounter{subsubsection}{-1}
\subsubsection{}\etiquette[2.5.0]{I.sec2.5.0}
Let\pageoriginale $I$, $J$, and $C$ be three categories, and let $G:I\times
J\rightarrow C$ be a functor (\ie a functor from $I$ with values in the
category of functors from $J$ to $C$). Suppose that the projective
limits of the functors
\begin{align*}
G_i:&J\longrightarrow C\qquad i\in \ob(I)\\
    & j\longmapsto G(i\times j)
\end{align*}
are representable, and that the functor
$$
i\longmapsto \varprojlim G_i
$$
admits a representable projective limit. It is then clear that
the functor $G$ admits a representable projective limit and that there is
a canonical isomorphism
$$
\varprojlim_{I\times J} G\\
\xrightarrow{\sim}\varprojlim_I \varprojlim_J G_i{\quoi},
$$
and consequently a canonical isomorphism
$$
\varprojlim_I \varprojlim_J G_i\xrightarrow{\sim}\varprojlim_J
\varprojlim_I G_j{\quoi}.
$$
Hereafter we shall say more briefly that projective limits
commute with projective limits. Similarly, inductive limits
commute with inductive limits. But in general it is not true
that inductive limits commute with projective limits.

\setcounter{subsection}{5}
\begin{enonce*}{Definition 2.5}\etiquette[2.5]{I.def2.5}
Let $C$ be a category with representable fiber products, $I$ a
category, $G:I\rightarrow C$ a functor, $g:G\rightarrow X$ a
morphism from $G$ to an object of $C$ (\ie a morphism from $G$ to
the constant functor associated with $X$), and $m:Y\rightarrow X$ a morphism
in $C$. Let\pageoriginale $\fprod{G}{Y}{X}:I\rightarrow C$ be the
functor
$$
i\longmapsto \fprod{G(i)}{Y}{X} {\quoi}.
$$
We say that the inductive limit of $G$ is universal if, for every
object $X$, every morphism $g:G\rightarrow X$, and every morphism
$m:Y\rightarrow X$,
\begin{enumerate}
\item[a)] the inductive limit of the functor $\fprod{G}{Y}{X}$ is
  representable,

\item[b)] the canonical morphism $\varinjlim (\fprod{G}{Y}{X}
  )\rightarrow \fprod{(\varinjlim G)}{Y}{X}$ is an isomorphism.
\end{enumerate}
\end{enonce*}

\begin{enonce*}{Proposition 2.6}\etiquette[2.6]{I.prop2.6}
Let $𝒰$ be a universe. The inductive limits in
$𝒰$-$\Ens$ that are representable (in particular, the
$𝒰$-inductive limits $\hyperref[I.sec2.2.1]{(2.2.1)}$ in $𝒰$-$\Ens$) %% étiquette !
are universal.
\end{enonce*}
We shall also use another commutation result for
projective and inductive limits, which we now present.

\begin{enonce*}{Definition 2.7}\etiquette[2.7]{I.def2.7}
A category $I$ is pseudo-filtered if it has the following
properties:
\begin{itemize}
\item[PS 1)] Every diagram of the form
\[
\xymatrix@R=.5cm@C=1cm{
       & j\\
i \ar[ru]\ar[dr]     &  \\
       & j'
}
\]
can be embedded in a commutative diagram
\[
\xymatrix@R=.5cm@C=1cm{
  & j \ar[rd]&\\
i\ar[ur]\ar[rd] &   & k\\
  & j'\ar[ru]&
}
\]

\item[PS 2)] Every diagram of the form
\[
\xymatrix@R=.75cm@C=1.5cm@M=.25cm{
i\ar@<2.5pt>[r]^{u}\ar@<-2.5pt>[r]_{v} & j }
\]
%%
  can \pageoriginale be embedded in a diagram
\[
\xymatrix@C=1.75cm@M=.25cm{ i\ar@<2.5pt>[r]^{u}\ar@<-2.5pt>[r]_{v} &
j\ar[r]%%
^{w} & k\\
}
\]
  such that
  $$
  w\circ u=w\circ v{\quoi}.
  $$
\end{itemize}
A category $I$ is said to be filtered if it is pseudo-filtered,
nonempty, and connected, \ie if any two objects of $I$ can
be joined by a sequence of arrows (with no condition imposed
on the directions of the arrows); in the presence of {\rm
PS 2}, this also means that $I\neq ∅$ and that, for any two objects $a$, $b$ of
$I$, there is always an object $c$ of $I$ and arrows $a\sisi{$
$}{\rightarrow}c$ and $b\sisi{$ $}{\rightarrow}c$. We also say
that a category $I$ is cofiltered if $I^\circ$ is filtered.
\end{enonce*}

\begin{enonce*}{Example 2.7.1}\etiquette[2.7.1]{I.exem2.7.1}
If, in $I$, amalgamated sums (\resp sums of two objects)
and cokernels of pairs of parallel arrows are representable, then $I$
is pseudo-filtered (\resp filtered).
\end{enonce*}

\begin{enonce*}{Proposition 2.8}\etiquette[2.8]{I.prop2.8}
Let $𝒰$ be a universe. Filtered $𝒰$-inductive limits
in $𝒰$-$\Ens$ commute with finite projective
limits.
\end{enonce*}

We are immediately reduced to proving that filtered $𝒰$-inductive
limits commute with fiber products. The proof is
left to the reader. One may use the description of the limit
given by the following

\begin{enonce*}{Lemma 2.8.1}\etiquette[2.8.1]{I.lem2.8.1}
Let $I$ be a small filtered category and $i\mapsto X_i$ a
functor from $I$ to $𝒰$-$\Ens$. On the disjoint union
$\amalg_{i\in
 \ob I} X_i$, let $R$ be the relation:
\begin{itemize}
\item[(R)] Two elements $x_i\in X_i$ and $x_j\in X_j$ are
  related if there exists an object $k\in \ob I$ and two morphisms
  $u:i\rightarrow k$ and $v:j\rightarrow k$ such that the images in
  $X_k$ of $x_i$ and $x_j$ under the transition maps $u$ and
  $v$, respectively, are equal.
\end{itemize}
Then:\pageoriginale
\begin{enumerate}
\renewcommand{\labelenumi}{\theenumi)}
\item $R$ is an equivalence relation.

\item The quotient $\amalg_{i\in \ob I}X_i/R$ is canonically isomorphic
  to $\varinjlim_I X_i$.

\item Two elements $x_i\in X_i$ and $x_j\in X_j$ are
  respectively equivalent under $R$ to two elements of the
  same $X_k$.

\item For every $i\in \ob I$, two elements $\alpha$ and $\beta$ of $X_i$ are
  equivalent under $R$ if and only if there is a morphism
  $u:i\rightarrow j$ such that $u(\alpha)=u(\beta)$.
\end{enumerate}
\end{enonce*}

Assertion 1) follows from (PS 1) \hyperref[I.def2.7]{(2.7)}. To prove 2), one %% étiquette !
checks that $\amalg_{i\in \ob I}X_i/R$ has the universal
property of the inductive limit (only (PS 1) is used).
Assertion 3) follows from the fact that $I$ is connected. Assertion 4)
follows from (PS 2).

\begin{enonce*}{Corollary 2.9}\etiquette[2.9]{I.cor2.9}
Let $%%
\gamma$ be a kind of algebraic structure ``defined by finite
projective limits.'' (The reader is asked to give a mathematical meaning
to the preceding phrase. We merely note that the structures of groups,
abelian groups, rings, modules, etc. %%
are structures of this kind.) Denote by $𝒰$-$\gamma$ the
category of $\gamma$-objects in $𝒰$-$\Ens$. The functor that
associates with each object of $𝒰$-$\gamma$ its underlying set
commutes with filtered $𝒰$-inductive limits. Consequently,
filtered $𝒰$-inductive limits in $𝒰$-$\gamma$
commute with finite projective limits.
\end{enonce*}

\begin{enonce*}{Corollary 2.10}\etiquette[2.10]{I.cor2.10}
Pseudo-filtered $𝒰$-inductive limits in $𝒰$-$\Ab$
commute with finite projective limits.
\end{enonce*}

(One reduces to filtered limits by decomposing the indexing category
into connected components.)

We now record\pageoriginale a result that we shall use
constantly: let $C$ and $C'$ be two $𝒰$-categories, and let $u$ and
$v$ be two functors $\xymatrix@C=1.25cm{ C  \ar@<-.6ex>[r]_u & C'
\ar@<-.6ex>[l]_v }$, with $u$ \emph{left adjoint} to $v$. Recall
that this means that there is an isomorphism between the
bifunctors with values in $𝒰$-$\Ens$:
$$
\Hom_{C'}(u(X),X')\xrightarrow{\sim}\Hom_{C}(X,v(X')){\quoi}.
$$

\begin{enonce*}{Proposition 2.11}\etiquette[2.11]{I.prop2.11}
The functor $u$ commutes with representable inductive limits; the
functor $v$ commutes with representable projective limits.
\end{enonce*}

This assertion means that, for every category $I$ and every
functor $G:I\rightarrow C'$ such that the projective limit of $G$ is
representable, the functor $v\circ G$ has a representable
projective limit and there is a canonical isomorphism
$$
v(\varprojlim G)\longrightarrow \varprojlim(v\circ G){\quoi}.
$$
The reader may spell out the assertion concerning the
functor $u$.

Finally, we record a computation of projective limits in the case where the
indexing category has products.

\begin{enonce*}{Proposition 2.12}\etiquette[2.12]{I.prop2.12}
Let $I$ and $C$ be two categories, and $𝒱$ a universe. Assume
that the $𝒱$-projective limits in $C$ are
representable and that there is in $I$ a family of objects
$(i_\alpha)_{\alpha\in A}$, where $A$ is an element of $𝒱$,
such that:
\begin{enumerate}
\renewcommand{\labelenumi}{\theenumi)}
\item the products $i_\alpha\times i_\beta$ are representable
  for every pair $(\alpha,\beta)\in A\times A$,

\item every object of $I$ admits a morphism to at least one of the $i_\alpha$.
\end{enumerate}


Then, for every contravariant functor $F$ from $I$ to $C$,
$$
F:I^\circ\longrightarrow C{\quoi},
$$
the projective limit of $F$ exists, and there is an isomorphism
functorial in $F$:
$$
\varprojlim F\xrightarrow{\sim} \Ker \left(\prod_{\alpha\in
  A} F(i_\alpha)\rightrightarrows \prod_{(\alpha,\beta)\in A\times A}
F(i_\alpha \times i_\beta)\right){\quoi},
$$
where the\pageoriginale two arrows are defined by the projections
of the products $i_\alpha\times i_\beta$ onto their factors.
\end{enonce*}

% END INCLUDED SGA4_English_current_buildable_source_20260906/source/01/drafts/SGA4_I_LINES_411_824_English.texfrag


% BEGIN INCLUDED SGA4_English_current_buildable_source_20260906/source/01/drafts/SGA4_I_LINES_825_1024_English.texfrag
\section{Exactness properties of the category of presheaves}\etiquette[3]{I.sec3}

Let $C$ be a $𝒰$-category and let $\widehat{C}$ be the
category of presheaves on $C$. All the exactness properties of
$\widehat{C}$ follow from the following proposition:

\begin{enonce*}{Proposition 3.1}\etiquette[3.1]{I.prop3.1}
The $𝒰$-projective and inductive limits in
$\widehat{C}$ are representable. For every object $X$
of $C$, the functor on $\widehat{C}$
$$
F\longmapsto F(X)\qquad F\in \widehat{C}
$$
commutes with inductive and projective limits.

In other words, ``inductive and projective limits in
$\widehat{C}$ are computed argument by argument.''
We mention a few corollaries.
\end{enonce*}

\begin{enonce*}{Corollary 3.2}\etiquette[3.2]{I.cor3.2}
Let $\gamma$ be an algebraic structure defined by finite projective
limits. The category of contravariant functors with values in
$𝒰$-$\gamma$ is equivalent to the category of $\gamma$-objects of
$\widehat{C}$.
\end{enonce*}

\begin{enonce*}{Corollary 3.3}\etiquette[3.3]{I.cor3.3}
A morphism of $\widehat{C}$ that is both a monomorphism
and an epimorphism is an isomorphism. A morphism of
$\widehat{C}$ factors uniquely as an epimorphism followed
by a monomorphism. The representable inductive limits in
$\widehat{C}$ are universal. The filtered
$𝒰$-inductive limits in $\widehat{C}$ commute with
finite projective limits. The canonical functor $C\rightarrow
\widehat{C}$ commutes with representable projective
limits.
\end{enonce*}

More\pageoriginale generally, one may say that the category
$\widehat{C}$ inherits all the properties of the category of
$𝒰$-sets that involve inductive and projective limits.

\setcounter{subsection}{4} \setcounter{subsubsection}{-1}
\subsubsection{}\etiquette[3.4.0]{I.sec3.4.0}
Let $F$ be an object of $\widehat{C}$. We denote by
$C/F$ the following category. The objects of $C/F$ are pairs consisting
of an object $X$ of $C$ and a morphism $u$ from $X$ to $F$.
Let $(X,u)$ and $(Y,v)$ be two objects. A morphism from $(X,u)$ to
$(Y,v)$ is a morphism $g$ from $X$ to $Y$ such that the following
diagram is commutative:
\[
\xymatrix@=1.25cm{
X \ar[rr]^{g}\ar[rd]_{u}&   & Y\ar[ld]^{v}\\
            & F &.
}
\]

\begin{enonce*}{Proposition 3.4}\etiquette[3.4]{I.prop3.4}
With the notation of \hyperref[I.sec3.4.0]{(3.4.0)}, the source functor $C/F\rightarrow
\widehat{C}$ admits a representable inductive limit in
$\widehat{C}$. The canonical morphism
$$
\mathop{\underrightarrow{\text{\rm lim source}}}_{C/F}\;
(.)\longrightarrow F
$$
is an isomorphism.
\end{enonce*}


\begin{enonce*}{Corollary 3.5}\etiquette[3.5]{I.cor3.5}
Let $F$ and $H$ be two objects of $\widehat{C}$. There is
a canonical isomorphism
$$
\Hom(F,H)\xrightarrow{\sim}\varprojlim_{(X,u)\in \ob(C/F)}
H(X){\quoi}.
$$
\end{enonce*}

\begin{enonce*}[remark]{Proof}
The corollary follows immediately from \hyperref[I.prop3.4]{3.4}
and \hyperref[I.def1.2]{1.2}.
\end{enonce*}

\begin{enonce*}{Proposition 3.6}\etiquette[3.6]{I.prop3.6}
Let\pageoriginale $C$ be a $𝒰$-category, let $𝒱$ be a
universe containing $𝒰$, and let
$\nobreak{i:\widehat{C}_{𝒰}\rightarrow
\widehat{C}_{𝒱}}$ be the natural inclusion functor of the
corresponding categories of presheaves \hyperref[I.condef1.3]{(1.3)}. The functor $i$
commutes with inductive and projective limits. Moreover, for every
object $F$ of $\widehat{C}_{𝒰}$, every subobject of $i(F)$
is isomorphic to the image under $i$ of a unique subobject of $F$
(thus giving a bijection from the set of subobjects of $F$ to the
set of subobjects of $i(F)$).
\end{enonce*}

\section{Sieves}\etiquette[4]{I.sec4}

\begin{enonce*}{Definition 4.1}\etiquette[4.1]{I.def4.1}
Let $C$ be a category. A \emph{sieve} in the category $C$ is a full
subcategory $D$ of $C$ with the following property: every object of
$C$ from which there exists a morphism to an object of $D$ belongs to
$D$. If $X$ is an object of $C$, the sieves of $X$ are (by an abuse of
language) the sieves in the category $C/X$.
\end{enonce*}

Let $𝒰$ be a universe such that $C$ is a $𝒰$-category.
Let $\widehat{C}$ be the corresponding category of presheaves.
To every sieve of $X$ one associates a subobject of $X$ in
$\widehat{C}$ as follows: to every object $Y$ of $C$, one
assigns the set of morphisms $f:Y\rightarrow X$ such that the object
$(Y,f)$ belongs to the sieve.

\begin{enonce*}{Proposition 4.2}\etiquette[4.2]{I.prop4.2}
The map defined above establishes a %%
bijection between the set of sieves of $X$ and the set of subobjects
of $X$ in $\widehat{C}$.
\end{enonce*}

\begin{enonce*}[remark]{Proof}
We shall only \sisi{show that it}{describe what} is the inverse map. To every
subfunctor $R$ of $X$ we associate the category $C/R$ of objects of
$C$ over $R$ \sisi{(3.4)}{(\hyperref[I.prop3.4]{3.4})}. One immediately
checks that $C/R$ is a sieve of $X$.
\end{enonce*}

\begin{enonce*}[remark]{Remark 4.2.1}\etiquette[4.2.1]{I.rem4.2.1}
One\pageoriginale sees likewise that the sieves of $C$ are in canonical
bijective correspondence with the set of subfunctors of the ``terminal
functor'' on $C$ (the ``terminal'' object of $\widehat{C}$).
\end{enonce*}

\setcounter{subsection}{2}
\subsection{}\etiquette[4.3]{I.sec4.3}
Let $C$ be a $𝒰$-category. By an abuse of language, we shall
also call the subobjects of $X$ in the category $\widehat{C}$
the sieves of $X$. This abuse of language allows us, for every presheaf
$F$ and every sieve $R$ of $X$, to define
$\Hom_{\widehat{C}}(R,F)$ as the set of %%
morphisms from the functor $R$ to $F$. There is moreover a canonical
isomorphism, functorial in $F$ \hyperref[I.cor3.5]{(3.5)}:
$$
\Hom_{\widehat{C}}(R,F)\xrightarrow{\sim}\varprojlim_{C/R}
F(.){\quoi},
$$
which also gives a direct definition\nde{More simply, $C/R$ is identified
with $R$, so that one has the formula $
\Hom_{\widehat{C}}(R,F)\xrightarrow{\sim}\varprojlim_{R}
F{\quoi}, $ where, by an abuse of notation, $F$ denotes the composite of
$F$ and the tautological ``source'' functor $R\rightarrow C/X\rightarrow
C$.}. Likewise, \hyperref[I.prop4.2]{Proposition 4.2} allows us to transfer the usual operations
on functors to sieves. We mention the following:

\subsubsection{}\etiquette[4.3.1]{I.4.3.1}
\emph{Base change.} Let $R$ be a sieve of $X$ and let
$f:Y\rightarrow X$ be a morphism of objects of $C$. The fiber product
$\fprod{R}{Y}{X}$ is a sieve of $Y$, called the \emph{sieve
deduced
  from $R$ by base change.} The corresponding subcategory of $C/Y$
is the inverse image of the subcategory of $C/X$ defined by $R$ under
the canonical functor $C/Y\rightarrow C/X$ defined by $f$.

\subsubsection{}\etiquette[4.3.2]{I.sec4.3.2}
\emph{Order relation, intersection, union.} The inclusion relation on
the subfunctors of $X$ is an order relation. The union and intersection
of a family of sieves indexed by an arbitrary set may be defined as the
least upper bound and the greatest lower bound of the family of
corresponding subpresheaves.

\subsubsection{}\etiquette[4.3.3]{I.sec4.3.3}
\emph{Image, generated sieve.}\pageoriginale Let
$(F_\alpha)_{\alpha\in A}$ be a family of presheaves and, for every
$\alpha\in A$, let $f_\alpha:F_\alpha\rightarrow X$ be a morphism,
where $X$ is an object of $C$. The \emph{image} of this family of
morphisms is the union of the images of the $f_\alpha$. The image of
this family is therefore a sieve of $X$. In particular, if all the
$F_\alpha$ are objects of $C$, the image sieve will be called the
sieve generated by the morphisms $f_\alpha$. The category $C/R$ is
the full subcategory of $C/X$ formed by the objects $X'\rightarrow X$
over $X$ for which there exists an $X$-morphism from $X'$ to one of the
$F_{\alpha}$.

The reader may, as an exercise, translate the relations and operations
defined here on subfunctors into the language of the categories $C/R$.
One will then observe that these relations and operations do not depend
on the universe $𝒰$ for which $C$ is a $𝒰$-category, and that
they are therefore defined for every category, without having to specify
the universe to which the morphism sets belong; this could moreover have
been foreseen \textit{a priori} from \hyperref[I.prop3.6]{3.6}.

% END INCLUDED SGA4_English_current_buildable_source_20260906/source/01/drafts/SGA4_I_LINES_825_1024_English.texfrag


% BEGIN INCLUDED SGA4_English_current_buildable_source_20260906/source/01/drafts/SGA4_I_LINES_1025_1161_English.texfrag
\section{Functoriality of categories of presheaves}\etiquette[5]{I.sec5}

\setcounter{subsection}{-1}
\subsection{}\etiquette[5.0]{I.sec5.0}
Let $C$, $C'$, and $D$ be three categories, and let $u:C\rightarrow C'$ be a
functor. We denote by $u^\ast$ the functor \sisi{$$
u^\ast:\Hhom({C'}^\circ,D)\longrightarrow \Hhom(C^\circ, D),
G\longmapsto G\circ u{\quoi},
$$}{$$u^\ast:\left\{\begin{array}{ccc}
                      \Hhom({C'}^\circ,D) & \longrightarrow & \Hhom(C^\circ, D) \\
                      G & \longmapsto & G\circ u{\quoi}
                    \end{array}
\right.$$} obtained by composition with the functor $u$. The functor
$u^\ast$ commutes with inductive and projective limits.

\begin{enonce*}{Proposition 5.1}\etiquette[5.1]{I.prop5.1}
Suppose\pageoriginale that $C$ is small and that the
$𝒰$-inductive (
\resp projective) limits in $D$ are representable. The functor $u^\ast$
has a left adjoint $u_{!}$ (\resp a right adjoint $u_\ast$). Thus there
is an isomorphism
\begin{gather*}
\Hom_{\Hhom(C^\circ,D)}(F,u^\ast G)\xrightarrow{\sim}
\Hom_{\Hhom({C'}^\circ,D)}(u_{!}F,G)\\
(\mbox{\resp  } \Hom_{\Hhom(C^\circ,D)}(u^\ast G,F)\simeq
\Hom_{\Hhom({C'}^\circ,D)}(G,u_\ast F){\quoi}).
\end{gather*}
\end{enonce*}

\begin{enonce*}[remark]{Proof}
We shall give only the proof of the existence of the left adjoint functor.
The \resp part of the proposition will then follow formally from the
isomorphisms
$$
\Hhom(C^\circ,
D)\xrightarrow{\sim}\Hhom(C,D^\circ)^\circ\xrightarrow{\sim}\Hhom((C^\circ)^\circ,D^\circ)^\circ{\quoi}.
$$
\end{enonce*}

Let $Y$ be an object of $C'$. Denote by $I^Y_u$ the following category.
The objects of $I^Y_u$ are pairs $(X,m)$, where $X$ is an object of $C$
and $m$ is a morphism $Y\rightarrow u(X)$. Let $(X,m)$ and $(X',m')$
be two objects of $I^Y_u$. A morphism from $(X,m)$ to $(X',m')$ is a
morphism $\xi:X\rightarrow X'$ such that $m'=u(\xi)m$. Composition of
morphisms is defined in the evident way.

Let $f:Y\rightarrow Y'$ be a morphism of $C'$. By composition, the
morphism $f$ defines a functor $I^f_u:I^{Y'}_u\rightarrow I^Y_u$.

There is moreover a functor $\pr_Y:I^Y_u\rightarrow C$ that assigns the
object $X$ to the object $(X,m)$. Notice that the diagram
\[\vcenter{
\xymatrix@=1.5cm{
       &I^{Y'}_u \ar[rr]^{I^f_u} \ar[rd]_{\pr_{Y'}}& & I^Y_u \ar[ld]^{\pr_Y}\\
       & &   C   & \\
}}\leqno{(*)}
\]
is then commutative.

Now\pageoriginale let $F$ be a presheaf on $C$, and set
\hypertarget{sga4.exp.I.equation.5.1.1}{}$$
\displaylines{\rlap{(5.1.1)}\hfill u_{!}
  F(Y)=\varinjlim_{I^Y_u} F\sisi{}{\circ}\pr_Y(.){\quoi}.\hfill}
$$
The commutativity of diagram ($^\ast$) and the functoriality of the
inductive limit make $u_{!}F$ a presheaf on $C'$. Let us show that the %%
functor $u_{!}$ is a left adjoint of the functor $u^\ast$. For this, we
show that for every presheaf $G$ on $C'$, there is a functorial
isomorphism
$$
\Hom(u_{!}F,G)\xrightarrow{\sim}\Hom(F,u^\ast G){\quoi}.
$$
Let $\xi\in \Hom(u_{!}F,G)$. Thus, for every object $X$ of $C$, there
is a morphism
$$
\xi_X:u_{!}F(u(X))\longrightarrow G(u(X)){\quoi}.
$$
But $(u(X),\id_{u(X)})$ is an object of $I^{u(X)}_u$, and by the
definition of the inductive limit there is a canonical morphism
$$
F(X)\longrightarrow u_{!}F(u(X)){\quoi}.
$$
For every object $X$ of $C$, we therefore obtain a morphism
$$
\eta_X:F(X)\longrightarrow G(u(X))
$$
which is visibly functorial in $X$. Hence a morphism
$$
\eta:F\longrightarrow u^\ast G{\quoi}.
$$

Conversely, let $\eta\in \Hom(F,u^\ast G)$. For every object $Y$ of
$C'$, this gives a morphism of functors
$$
\eta_Y:F\sisi{}{\circ}\pr_Y\longrightarrow (u^\ast G)\circ\pr_Y{\quoi},
$$
and hence, by composition with the evident morphism from the functor
$(u^\ast G)\circ\pr_Y$ to the constant functor $G(Y)$, a morphism
$$
\xi_Y:u_{!}F(Y)\longrightarrow G(Y)
$$
which is functorial in $Y$. Hence a morphism
$$
\xi:u_{!}F\longrightarrow G{\quoi}.
$$

The\pageoriginale reader will check that the two maps thus defined are
inverse to one another, thereby completing the proof.

\begin{enonce*}{Proposition 5.2}\etiquette[5.2]{I.prop5.2}
Suppose that in $D$ the $𝒰$-inductive limits are representable,
finite projective limits are representable, and filtered
$𝒰$-inductive limits commute with finite projective limits. Suppose
also that finite projective limits are representable in $C$ and that
the functor $u$ is left exact $(\hyperref[I.cor2.3.2]{2.3.2})$. Then
finite projective limits are representable in $\Hhom(C^\circ,D)$ and
in $\Hhom({C'}^\circ,D)$, and the functor $u_{!}$ is left exact.
\end{enonce*}

\begin{enonce*}[remark]{Proof}
The first assertion is trivial. Let us prove the second. By the proof of
\hyperref[I.prop5.1]{5.1}, for every presheaf $F$ on $C$ and every object $Y$ of $C'$, one has
$$
u_{!}F(Y)\xrightarrow{\sim} \varinjlim_{I^Y_u}
F\sisi{}{\circ}\pr_Y{\quoi}.
$$
 It therefore suffices to show that the category $(I^Y_u)^\circ$
 satisfies axioms (PS 1) and (PS 2) \hyperref[I.def2.7]{(2.7)}, and that this category is
 connected. The verification is left to the reader.
\end{enonce*}

% END INCLUDED SGA4_English_current_buildable_source_20260906/source/01/drafts/SGA4_I_LINES_1025_1161_English.texfrag


% BEGIN INCLUDED SGA4_English_current_buildable_source_20260906/source/01/drafts/SGA4_I_LINES_1162_1346_English.texfrag
\setcounter{subsection}{2}
\subsection{}\etiquette[5.3]{I.sec5.3}
Specializing these results to the case where $D$ is the category of
$𝒰$-sets, one obtains a sequence of three functors:
$$
u_{!}, u^\ast, u_\ast {\quoi},
$$
which is a ``sequence of adjoint functors'' in the sense that
for any two consecutive functors in the sequence, the one on the right is
right adjoint to the other. Their\pageoriginale essential
properties are summarized in the following proposition:

\begin{enonce*}{Proposition 5.4}\etiquette[5.4]{I.prop5.4}
Let $C$ be a small category, $C'$ a $𝒰$-category, and
$u:C\rightarrow C'$ a functor.
\begin{enumerate}
\renewcommand{\labelenumi}{\theenumi)}
\item The functor $u^\ast:\widehat{C'}\rightarrow
  \widehat{C}$ commutes with inductive and
  projective limits.

\item The functor $u_\ast:\widehat{C}\rightarrow
  \widehat{C'}$ commutes with projective limits. For
  every presheaf $F$ on $C$ and every object $Y$ of $C'$, one has:
  $$
  u_\ast F(Y)\xrightarrow{\sim}
  \Hom_{\widehat{C}}(u^\ast(Y),F){\quoi}.
  $$

\hypertarget{sga4.exp.I.proposition.5.4.item.3}{}\item The functor $u_{!}:\widehat{C}\rightarrow
  \widehat{C'}$ commutes with inductive limits. The
  functor $u_{!}$ is defined only up to isomorphism, but it
  can always be chosen so that the diagram
\[
\xymatrix@C=2.75cm@R=1.5cm{
  C \ar[r]^{u}\ar[d]_{h}& C' \ar[d]^{h^\prime}\\
C{\widehat{\phantom{\;i}}}\ar[r]^{u_{!}}&C'{\widehat{\phantom{\;i}}}}
\]
($h$ and $h'$ are the canonical inclusion functors) is
commutative.

For every presheaf $F$ on $C$, one has:
$$
u_{!}F\xrightarrow{\sim} \varinjlim_{C/F} h'\sisi{}{\circ}u{\quoi}.
$$

\item If finite projective limits are representable in
  $C$ and if $u$ is left exact $\hyperref[I.cor2.3.2]{(2.3.2)}$, the functor $u_{!}$
  is left exact.
\end{enumerate}
\end{enonce*}

\begin{enonce*}[remark]{Proof}
Assertion (1) is trivial. Assertion (2) follows from the fact that
$u_\ast$ is a right adjoint functor, by (\hyperref[I.prop2.11]{2.11}) and
(\hyperref[I.prop1.4]{1.4}). The same is true of assertion (3), but one
also applies (\hyperref[I.prop3.4]{3.4}). Finally, assertion (4) is
simply \sisi{5.2}({\hyperref[I.prop5.2]{5.2}}), which can be applied thanks to
(\hyperref[I.def2.7]{2.7}).
\end{enonce*}

\begin{enonce*}{Proposition 5.5}\etiquette[5.5]{I.prop5.5}
Let\pageoriginale $C$ and $C'$ be two small categories and let
$\xymatrix@C=1cm{
 C \ar@<-2.5pt>[r]_{u} & C'\ar@<-2.5pt>[l]_{v}
}$ be a pair of functors, where $v$ is left adjoint to $u$. There
are then isomorphisms, compatible with the adjunction
isomorphisms:
\[
\xymatrix@R=.75cm{
v^{\ast} \ar[rr]^{\sim}&    & u_!      &\\
v_{\ast} \ar[rr]^{\sim}&    & u^{\ast} & . }
\]
\end{enonce*}

\begin{enonce*}[remark]{Proof}
It suffices to exhibit an isomorphism $v^\ast\xrightarrow{\sim}
u_{!}$; the other isomorphism follows from it by adjunction. Let $F$
be a presheaf on $C$ and let $Y$ be an object of $C'$. Then:
$$
v^\ast F(Y)\xrightarrow{\sim} \Hom(v(Y),F){\quoi}.
$$
Next, using (\hyperref[I.prop3.4]{3.4}):
$$
\Hom(v(Y),F)\xrightarrow{\sim} \varinjlim_{C/F} \Hom(v(Y),.){\quoi}.
$$
But $v$ is left adjoint to $u$, and hence:
$$
\varinjlim_{C/F}\Hom(v(Y),.)\xrightarrow{\sim}\varinjlim_{C/F}\Hom(Y,u(.)){\quoi}.
$$
Now using (\hyperref[I.prop5.4]{5.4}.3)), one obtains:
$$
\varinjlim_{C/F}\Hom(Y,u(.))\xrightarrow{\sim}
\Hom(Y,u_{!}F)\xrightarrow{\sim} u_{!}F(Y){\quoi}.
$$
Thus, for every object $Y$ of $C'$, we have determined an isomorphism
$v^\ast F(Y)\xrightarrow{\sim} u_{!}F(Y)$ that is clearly
functorial in $Y$ and in $F$, \cqfd
\end{enonce*}

\begin{enonce*}{Corollary 5.5.1}\etiquette[5.5.1]{I.cor5.5.1}
Let $u:C\rightarrow C'$ be a functor that admits a left adjoint.
The functor $u_{!}:\widehat{C}\rightarrow
\widehat{C'}$ commutes with projective limits
(recall that it commutes with inductive limits by $\hyperlink{sga4.exp.I.proposition.5.4.item.3}{(1.4.3)}$).
\end{enonce*}

\begin{enonce*}[remark]{Remark 5.5.2}\etiquette[5.5.2]{I.rem5.5.2}
One\pageoriginale thus obtains a ``sequence of \emph{four}
adjoint functors'' (\cf \hyperref[I.sec5.3]{5.3}):
$$
v_{!}, v^\ast=u_{!}, v_\ast = u^\ast, u_\ast{\quoi},
$$
of which the first three (\resp  the last three) therefore commute with
$\varinjlim$ (\resp  $\varprojlim$).
\end{enonce*}

\begin{enonce*}{Proposition 5.6}\etiquette[5.6]{I.prop5.6}
The hypotheses are those of $\hyperref[I.prop5.4]{5.4}$. The following
conditions are equivalent:
\begin{enumerate}
\renewcommand{\theenumi}{\roman{enumi}}
\renewcommand{\labelenumi}{\theenumi)}
\item The functor $u$ is fully faithful.

\item The functor $u_{!}$ is fully faithful.

\item The adjunction morphism
  $\id_{\widehat{C}}\rightarrow u^\ast u_{!}$ is an
  isomorphism.

\item The functor $u_\ast$ is fully faithful.

\item The adjunction morphism $u^\ast u_\ast\rightarrow
  \id_{\widehat{C}}$ is an isomorphism.
\end{enumerate}
\end{enonce*}

\begin{enonce*}[remark]{Proof}
It is clear that ii) $\Leftrightarrow$ iii) and iv) $\Leftrightarrow$
v) (general properties of adjoint functors), and that ii)
$\Rightarrow$ i) (\hyperref[I.prop5.4]{5.4}.3)). Let us show that i) $\Rightarrow$
iii). The functors $\id_{\widehat{C}}$, $u^\ast$, and
$u_{!}$ commute with inductive limits. By \hyperref[I.prop3.4]{(3.4)}, it therefore suffices
to prove that $H\rightarrow u^\ast u_{!} H$ is an
isomorphism when $H$ is representable, which is evident.
Let us show that %%
iii) is equivalent to %%
v). For every object $H$ (\resp  $K$) of $\widehat{C}$,
denote by $\Phi(H):H\rightarrow u^\ast u_{!} H$ (\resp  by
$\Psi(K):u^\ast u_\ast\rightarrow K$) the adjunction morphism. We
then have a commutative diagram:
\[
\xymatrix@C=3.5cm@L=6pt{
\Hom_{\widehat{C}}(H,u^{\ast}u_{\ast}
K)\ar@{-}[d]_{\wr}\ar[r]^{\quad \Hom(H,\Psi(K))}   & \Hom_{\widehat{C}}(H,K)\\
\Hom_{\widehat{C}}(u_{!}H,u_{\ast}K)\ar@{-}[d]_{\wr} &  \\
\Hom_{\widehat{C}}(u^{\ast}u_{!}H,K)\ar[uru]_{\quad
   \Hom(\Phi(H),K)} & .
}
\]

It\pageoriginale follows that $\Phi(H)$ is an isomorphism for every $H$
if and only if $\Psi(𝒦)$ is an isomorphism for every
$K$, \cqfd
\end{enonce*}

\begin{enonce*}[remark]{Remark 5.7}\etiquette[5.7]{I.rem5.7}
\begin{enumerate}
\renewcommand{\theenumi}{\alph{enumi}}
\renewcommand{\labelenumi}{\theenumi)}
\item The equivalences ii) $\Leftrightarrow$ iii)
  $\Leftrightarrow$ iv) $\Leftrightarrow$ v) are general results
  on adjoint functors.

\item The explicit form of $u_{!}$ \resp  $u_\ast$ given in the
  proof of \hyperref[I.prop1.4]{1.4} immediately shows that, under the general
  hypotheses of \hyperref[I.def1.1]{1.1}, if $u$ is fully faithful, then the
  adjunction morphism $\id\rightarrow u^\ast u_{!}$ (\resp  $u^\ast
  u_\ast\rightarrow \id$) is an isomorphism, \ie that $u_{!}$
  (\resp  $u_\ast$) is fully faithful.
\end{enumerate}
\end{enonce*}

\setcounter{subsection}{8} \setcounter{subsubsection}{-1}

% END INCLUDED SGA4_English_current_buildable_source_20260906/source/01/drafts/SGA4_I_LINES_1162_1346_English.texfrag


% BEGIN INCLUDED SGA4_English_current_buildable_source_20260906/source/01/drafts/SGA4_I_LINES_1347_1496_English.texfrag
\subsubsection{}\etiquette[5.8.0]{I.sec5.8.0}
Let $\gamma$ be a type of algebraic structure defined by finite
projective limits, let $𝒰$-$\gamma$-$\Ens$ be the category of
$\gamma$-objects of $𝒰$-$\Ens$, and let esj:
$𝒰$-$\gamma$-$\Ens\rightarrow 𝒰$-$\Ens$ be the underlying-set functor
(for simplicity, we suppose that the type of structure under consideration
has only one underlying set). Let $C$ be a category.
Composition with $\esj$ gives a functor denoted by
$$
\mbox{esj}^{\widehat{\phantom{,}}}:\Hhom(C^\circ,
𝒰\text{-}\gamma\text{-}\Ens)\longrightarrow
\widehat{C}{\quoi}.
$$
Since projective limits in $\widehat{C}$ are computed argument by
argument, the functor $\mbox{esj}^{\widehat{\phantom{,}}}$ factors as
an equivalence
\hypertarget{sga4.exp.I.equation.5.8.1}{}$$
\displaylines{\rlap{(5.8.1)}\hfill \Hhom(C^\circ,
  𝒰\text{-}\gamma\text{-}\Ens)\xrightarrow{\approx}
  \widehat{C}_\gamma{\quoi},\hfill}
$$
where $\widehat{C}_\gamma$ is the category of $\gamma$-objects of
$\widehat{C}$, followed by a functor, still denoted by
$$
\mbox{esj}^{\widehat{\phantom{,}}}:\widehat{C}_\gamma\longrightarrow
\widehat{C}{\quoi},
$$
and called the ``\emph{underlying presheaf of sets}'' functor.

\setcounter{subsubsection}{1}
\subsubsection{}\etiquette[5.8.2]{I.sec5.8.2}
Suppose\pageoriginale that the functor
$\esj:𝒰$-$\gamma$-$\Ens\rightarrow 𝒰$-$\Ens$ has a left adjoint
$\Lib:𝒰$-$\Ens\rightarrow 𝒰$-$\gamma$-$\Ens$
\footnote[1]{The functor $\Lib$ is the ``free generated
$\gamma$-object'' functor. Examples: free group, free abelian group,
free $A$-module, \etc} (in fact, one can show that this condition is
always satisfied). Composition with $\Lib$ gives a functor
$$
\Lib^{\widehat{\phantom{,}}}:\widehat{C}\longrightarrow
\Hhom(C^\circ, 𝒰\text{-}\gamma\text{-}\Ens)
$$
and, after composition with the equivalence \hyperlink{sga4.exp.I.equation.5.8.1}{(5.8.1)}, a functor still
denoted by
$$
\Lib^{\widehat{\phantom{,}}}:\widehat{C}\longrightarrow
\widehat{C}_\gamma
$$
and called the ``freely generated presheaf of $\gamma$-objects'' functor.
The functor $\Lib^{\widehat{\phantom{,}}}$ is left adjoint to the functor
$\esj^{\widehat{\phantom{,}}}$.

\begin{enonce*}{Proposition 5.8.3}\etiquette[5.8.3]{I.prop5.8.3}
Let $\gamma$ be a type of algebraic structure defined by finite
projective limits such that, in the category of $\gamma$-objects of
$𝒰$-$\Ens$, the $𝒰$-inductive limits are representable
\footnote[2]{One can show that this condition is always satisfied.};
let $C$ be a category belonging to $𝒰$, let $C'$ be a
$𝒰$-category, and let $u:C\rightarrow C'$ be a functor. Denote by
$\widehat{C}_\gamma$ (\resp $\widehat{C'}_\gamma$) the category of
$\gamma$-objects of $\widehat{C}$ (\resp $\widehat{C'}$), and by
${u^\ast}^\gamma$ the functor on $\gamma$-objects induced by the functor
$u^\ast$. It follows from \hyperref[I.prop5.1]{5.1} and the equivalence
\hyperlink{sga4.exp.I.equation.5.8.1}{5.8.1} that the functor ${u^\ast}^\gamma$ has a left (\resp right)
adjoint. This functor is denoted by $u_{!_\gamma}$ (\resp
$u_{\ast_\gamma}$).
\begin{enumerate}
\item The functor ${u^\ast}^\gamma$ commutes with inductive and
  projective limits. The diagram
\[\vcenter{
\xymatrix@C=3cm@R=1.5cm{
 \widehat{C'}_\gamma \ar[r]^{u^{\ast^{\gamma}}} \ar[d]_{\esj'\widehat{\phantom{a}}}&
\widehat{C}_\gamma  \ar[d]_{\esj^{\widehat{\phantom{a}}}}\\
 \widehat{C'} \ar[r]^{u_{\ast}}& \widehat{C}
}}\leqno{(*)}
\]
is\pageoriginale commutative ($\esj'^{\widehat{\phantom{,}}}$ and
$\esj^{\widehat{\phantom{,}}}$ denote the ``underlying-set'' functors).

\item The functor $u_{\ast_\gamma}$ commutes with projective limits. The
  diagram
\[\vcenter{
\xymatrix@C=3cm@R=1.5cm{
 \widehat{C}_\gamma \ar[r]^{u_{\ast_\gamma}} \ar[d]_{\esj\widehat{\phantom{a}}}&
  \widehat{C'}_\gamma\ar[d]_{\esj'\widehat{\phantom{a}}}\\
\widehat{C} \ar[r]^{u_{\ast}}&
\widehat{C'} }}\leqno{({*}{*})}
\]
commutes up to isomorphism.

\item The functor $u_{!_\gamma}$ commutes with inductive limits. Suppose
  that $\esj^{\widehat{\phantom{,}}}$
  (\resp ${\esj'}^{\widehat{\phantom{,}}}$) has a left adjoint
  $\Lib^{\widehat{\phantom{,}}}$
  (\resp ${\Lib'}^{\widehat{\phantom{,}}}$) $\hyperref[I.sec5.8.2]{(5.8.2)}$. The diagram
\[\vcenter{
\xymatrix@C=3.5cm@R=1.75cm{
 \widehat{C} \ar[r]^{u_{!}} \ar[d]_{\Lib\widehat{\phantom{a}}}&
  \widehat{C'} \ar[d]_{\Lib'\widehat{\phantom{a}}}\\
\widehat{C}_{\gamma} \ar[r]^{u_{!\gamma}} &
\widehat{C'}_{\gamma} }}\leqno{({*}{*}{*})}
\]
commutes up to isomorphism.
\end{enumerate}

Suppose that $u_{!}$ commutes with finite projective limits
(\hyperref[I.prop5.4]{5.4}) and (\hyperref[I.prop5.6]{5.6}). Then the diagram
\[\vcenter{
\xymatrix@C=3.5cm@R=1.5cm{
 \widehat{C}_\gamma \ar[r]^{u_{!\gamma}} \ar[d]_{\esj\widehat{\phantom{a}}}&
  \widehat{C'}_\gamma\ar[d]_{\esj'\widehat{\phantom{a}}}\\
\widehat{C} \ar[r]^{u_{!}}& \widehat{C'}
}}\leqno{({*}{*}{*}{*})}
\]
commutes up to isomorphism, and $u_{!\gamma}$ commutes with finite
projective limits.
\end{enonce*}

\begin{enonce*}[remark]{Proof}
Assertion (1) is evident. So is assertion (2), since $u_\ast$ is a right
adjoint and therefore (\hyperref[I.prop2.11]{2.11}) commutes with projective
limits and, in particular, with finite projective limits; this gives the
commutativity of diagram $({\ast}{\ast})$. The commutativity of diagram
$({\ast}{\ast}{\ast})$ follows from the uniqueness, up\pageoriginale to
isomorphism, of the left adjoint functor, and finally the commutativity of
diagram $({\ast}{\ast}{\ast}{\ast})$ follows immediately from the fact that
$u_{!}$ commutes with finite projective limits.
\end{enonce*}

\begin{enonce*}[remark]{Notation 5.9}\etiquette[5.9]{I.not5.9}
By an abuse of notation, the functors ${u^\ast}^\gamma$ and
$u_{\ast\gamma}$ will often be denoted by $u^\ast$ and $u_\ast$. In view
of the commutativity of diagrams $(\ast)$ and $({\ast}{\ast})$, this
cannot cause confusion. On the other hand, when $u_{!}$ does not commute
with finite projective limits, diagram $({\ast}{\ast}{\ast}{\ast})$ does
not commute up to isomorphism, and the notations $u_{!}$ and
$u_{!\gamma}$ must be used to avoid confusion.
\end{enonce*}

\setcounter{subsection}{9}

% END INCLUDED SGA4_English_current_buildable_source_20260906/source/01/drafts/SGA4_I_LINES_1347_1496_English.texfrag


% BEGIN INCLUDED SGA4_English_current_buildable_source_20260906/source/01/drafts/SGA4_I_LINES_1497_1738_English.texfrag
\subsection{}\etiquette[5.10]{I.sec5.10}
Let $C$ be a small category. For every object $X$ of
$\widehat{C}$, let $C/X$ denote the category of arrows with target
$X$ and source an object of $C$. The source functor defines a functor
$j_X:C/X\rightarrow C$. Let $m:Y\rightarrow X$ be a morphism of
$\widehat{C}$. Composition of morphisms defines a functor
$j_m:C/Y\rightarrow C/X$. The diagram
\[
\xymatrix@=1.5cm{
C/Y \ar[rr]^{j_m} \ar[rd]_{j_{Y}}&   & C/X\ar[ld]^{j_{X}}\\
    & C &
}
\]
is commutative. It follows from \hyperref[I.prop5.1]{5.1} that, for every
object $F$ of $(C/X)^{\widehat{\phantom{,}}}$ and every object $Y$ of
$C$, one has
\hypertarget{sga4.exp.I.equation.5.10.1}{}$$
\displaylines{\rlap{(5.10.1)}\hfill j_{X!}F(Y)=\coprod_{u\in
    \Hom_{\widehat{C}}(Y,X)} F(u){\quoi}.\hfill}
$$
Formula \hyperlink{sga4.exp.I.equation.5.10.1}{(5.10.1)} defines $j_{X!}$ when $C$ is a $𝒰$-category, and
one checks that the functor $j_{X!}$ thus defined is always right
adjoint\pageoriginale to the functor
$j^\ast_X:\widehat{C}\rightarrow
(C/X)^{\widehat{\phantom{,}}}$ (\hyperref[I.sec5.0]{5.0}).

\begin{enonce*}{Proposition 5.11}\etiquette[5.11]{I.prop5.11}
Let $C$ be a $𝒰$-category and let $X$ be a presheaf on $C$.
\begin{enumerate}
\renewcommand{\labelenumi}{\theenumi)}
\item The functor
$$
j_{X!}:(C/X)^{\widehat{\phantom{,}}}\longrightarrow
\widehat{C}
$$
factors through the category $\widehat{C}/X$:
$$
(C/X)^{\widehat{\phantom{,}}}\xrightarrow{e_X}
\widehat{C}/X\longrightarrow
\widehat{C}{\quoi}.
$$
The functor $e_X$ is an equivalence of categories.

\item The functor $e_{X}\circ j^\ast_X:\widehat{C}
\to  \widehat{C}/X$ is canonically isomorphic to the functor
  $H\mapsto (H\times X \xrightarrow{\pr_2} X)$.
\end{enumerate}
\end{enonce*}

\begin{enonce*}[remark]{Proof}
\begin{enumerate}
\renewcommand{\labelenumi}{\theenumi)}
\item Let $f$ be the terminal object of $(C/X)^{\widehat{\phantom{,}}}$.
There is a canonical isomorphism $f\xrightarrow{\sim}
\varinjlim_{Y\in \ob C/X}Y$, and hence
$j_{X!}(f)=\varinjlim_{Y\in \ob C/X}j_X(Y)$ \hyperref[I.prop5.4]{(5.4)}. But
$j_{X!}(f)\simeq X$, which gives the factorization. To show that $e_X$
is an equivalence, it is enough to exhibit a quasi-inverse functor. To
every object $H\rightarrow X$ of $\widehat{C}/X$, associate the
presheaf on $C/X$
$$
(Y\rightarrow X)\longmapsto
  \Hom_{\widehat{C}/X}((Y\rightarrow X), (H\rightarrow
  X)){\quoi}.
$$

\item The functor $e_X\circ j^\ast_X$ is right adjoint to the forgetful
functor, and therefore the functor $e_X\circ j^\ast$ is canonically
isomorphic to the functor $H\mapsto (H\times
  H\xrightarrow{\pr_2}X)$.
\end{enumerate}
\end{enonce*}

\setcounter{subsection}{11}
\subsection{}\etiquette[5.12]{I.sec5.12}
Let $m:Y\rightarrow X$ be a morphism of $\widehat{C}$.
By \hyperref[I.prop5.11]{(5.11)}, the morphism $m$ is canonically isomorphic to the image under
$e_X$ of an object of $(C/X)^{\widehat{\phantom{,}}}$, which we denote
by $[m]$. By restriction to the subcategories, the functor $e_X$ defines
an equivalence
$$
(C/X)/[m]\xrightarrow{e_m}C/Y{\quoi}.
$$
The\pageoriginale diagram
\[
\xymatrix{
  (C/X)/[m]\ar[rr]^{e_m} \ar[rd]_{j_{[m]}} &       & C/Y \ar[ld]^{j_m}\\
          & C/X   &
}
\]
commutes up to canonical isomorphism.

\subsection{}\etiquette[5.13]{I.sec5.13}

We record a result that will be useful in Exp.~VI. Let
$u:C\rightarrow C'$ be a functor between small categories. For every
object $H$ of $\widehat{C}$, denote by $u/H:C/H\rightarrow C'/u_{!}H$
the functor that assigns to every morphism $m:X\rightarrow H$ the
morphism $u_{!}X\xrightarrow{u_{!}m}u_{!}H$ (we know
(\hyperref[I.prop5.4]{5.4}) that we may always put $u_{!}X=uX$). The
following diagram is commutative:
\hypertarget{sga4.exp.I.equation.5.13.1}{}\[\vcenter{
\xymatrix@C=3cm@R=1.75cm{
C/H \ar[r]^{u/H}\ar[d]_{J_H} & C'/u_{!}H \ar[d]_{j_{u_{!}H}} \\
C\ar[r]^{u}   & C'\quad. }}\leqno{(5.13.1)}
\]
We therefore have a diagram commutative up to isomorphism:
%\[\vcenter{
%\xymatrix@C=2.5cm@R=1.5cm{ (C/H){\widehat{\phantom{a}}}
%\ar[r]^{(u/H)_{!}} \ar[d]_{(j_H)_!}&
%  (C'/u^·H){\widehat{\phantom{a}}} \ar[d]_{(j_{u!H})_{!}}\\
%  C{\widehat{\phantom{\;a}}} \ar[r]^{u_!}& C'{\widehat{\phantom{\;a}}}\quad,
%}}\leqno{(5.13.2)}
%\]
and since $(u/H)_{!}$ takes the terminal object of
$(C/H)^{\widehat{\phantom{,}}}$ to the terminal object of
$(C'/u_{!}H)^{\widehat{\phantom{,}}}$, the diagram
\hypertarget{sga4.exp.I.equation.5.13.3}{}\[\vcenter{
\xymatrix@C=3.5cm@R=1.5cm{ (C/H)\widehat{\phantom{a}}
\ar[r]^{(u/H)_!} \ar[d]_{e_H}&
  (C'/u_!H)\widehat{\phantom{a}} \ar[d]_{e_{u_{!}H}}\\
C\widehat{\phantom{a}}/H  \ar[r]^{u_{!}/H} &
C'\widehat{\phantom{a}}/u_{!}H }}\leqno{(5.13.3)}
\]
commutes\pageoriginale up to isomorphism \hyperref[I.def6.1]{(6.1)}.

\begin{enonce*}{Proposition 5.14}\etiquette[5.14]{I.prop5.14}
Let $u:C\rightarrow C'$ be a functor between small categories. Suppose
that $u$ has the following property:

{\rm(PPF)} For every object $X$ of $C$, the functor
$u/X:C/X\rightarrow C'/uX$ is fully faithful.

Then:
\begin{enumerate}
\renewcommand{\labelenumi}{\theenumi)}
\item Let $f$ be the terminal object of $\widehat{C}$. The functor $u$
factors as
$$
C=C/f\xrightarrow{u/f}C'/u_{!}f\xrightarrow{j_{u_{!}f}} C'{\quoi}.
$$
The functor $u/f$ is fully faithful.

\item The functor $u_{!}:\widehat{C}\rightarrow
  \widehat{C'}$ factors as
$$
\widehat{C}=(C/f)^{\widehat{\phantom{,}}}\xrightarrow{(u/f)_{!}}
(C'/u_{!}f)^{\widehat{\phantom{,}}}\xrightarrow{e_{u_{!}f}}\widehat{C'}/u_{!}f\longrightarrow
\widehat{C'}{\quoi},
$$
where the functor $(u/f)_{!}$ is fully faithful, the functor
$e_{u_{!}f}$ is an equivalence, and the functor
$\widehat{C'}/u_{!}f\rightarrow\widehat{C'}$ is the forgetful functor.

\item In particular, the functor $u_{!}$ is faithful, and therefore the
unit morphism $\id\xrightarrow{\Phi}u^\ast u_{!}$ is a monomorphism.
Moreover, for every morphism $\alpha:H\rightarrow K$ of
$\widehat{C}$, the diagram
\[
\xymatrix@C=3.5cm@R=1.5cm@M=.2cm{
H \ar@{^{(}->}[r]^{\Phi(H)}\ar[d]_{\alpha} & u^{*}u_{!}H \ar[d]^{u.u^{\bigdot}(\alpha)}\\
K \ar@{^{(}->}[r]^{\Phi(K)} &  u^{*}u_!K }
\]
is cartesian.
\end{enumerate}
\end{enonce*}

\begin{enonce*}[remark]{Proof}
\begin{enumerate}
\renewcommand{\labelenumi}{\theenumi)}
\item The\pageoriginale factorization comes from diagram \hyperlink{sga4.exp.I.equation.5.13.1}{(5.13.1)}. The
functor $u$ is faithful. Hence $u/f$ is faithful. Let us show that it is
fully faithful. Let $X$ and $Y$ be two objects of $C$, let
$\can_X:uX\rightarrow u_{!}f$ (\resp $\can_Y:uY\rightarrow u_{!}f$)
be the canonical morphisms, and let
\[
\xymatrix@=1.25cm{
uX \ar[rr]^{m} \ar[rd]_{\can_X}&               & uY \ar[ld]^{\can_Y}\\
   & u^{\bigdot}f &
}
\]
be a morphism of $C'/u_{!}f$. We have
$u_{!}f=\varinjlim_{Z\in \ob C}uZ$, and therefore \hyperref[I.prop3.1]{(3.1)}
$$
\Hom_{\widehat{C'}}(uX,u_{!}f)=\varinjlim_{Z\in
  \ob C}\Hom_{C'}(uX, uZ){\quoi}.
$$
By the definition of the inductive limit, to say that
$\can_Ym=\can_X$ is equivalent to saying that there exist
\begin{enumerate}
\renewcommand{\theenumii}{\alph{enumii}}
\renewcommand{\labelenumii}{\theenumii)}
\item a finite sequence of objects $X_i$ of $C$, $i\in [0,n]$,
  with $X_\circ=X$ and $X_n=Y$,

\item for every $i$, a morphism $m_i:uX\rightarrow
  uX_i$ $(m_\circ=\id_{X}, m_n=m)$,

\item for every pair $(i,i+1)$, a morphism $f_i:X_i\rightarrow
  X_{i+1}$ or else $f_i:X_{i+1}\rightarrow X_i$,
\end{enumerate}
such that the diagrams
\begin{displaymath}
\xymatrix@R=1.25cm@C=1.25cm{
   &  uX \ar[rd]^{\;\;m_{i+1} \qquad\quad\text{ or else}} \ar[ld]_{m_i}&\\
  uX_i\ar[rr]^{u(f_i)} &          &  uX_{i+1}
}\qquad \xymatrix@R=1.25cm@C=1.25cm{
      &  uX \ar[rd]^{\;\;m_{i+1}} \ar[ld]_{m_i}&\\
 uX_i &                                        &
   uX_{i+1}\ar[ll]_{u(f_i)}\quad ,
}
\end{displaymath}
are commutative.

It then follows immediately, by induction on $i$ and using property
(PPF), that $m_i$ is of the form $u(p_i)$. In particular $m=u(p)$, and
hence $u/f$ is fully faithful.

\item The\pageoriginale factorization is immediate. The functor
$(u/f)_{!}$ is fully faithful by \hyperref[I.prop5.6]{5.6}. The other
assertions follow from \hyperref[I.prop5.11]{5.11}.

\item The functor $u_{!}$ is the composite of the faithful forgetful
functor and fully faithful functors. It is therefore faithful. It follows
from the general properties of adjoint functors that the unit morphism
$\id\rightarrow u^\ast u_{!}$ is a monomorphism. By 2), the functor
$u_{!}$ appears as the composite of a fully faithful functor
$v:\widehat{C}\to\widehat{C'}/u_{!}f$ and the forgetful functor. The
functor $u^\ast$, right adjoint to $u_{!}$, is therefore the composite of
the ``product with $u_{!}f$'' functor, right adjoint to the forgetful
functor, and a functor $w$ right adjoint to $v$. Moreover, since $v$ is
fully faithful, the unit morphism $\id\rightarrow wv$ is an isomorphism.
The last assertion follows readily.
\end{enumerate}
\end{enonce*}


% END INCLUDED SGA4_English_current_buildable_source_20260906/source/01/drafts/SGA4_I_LINES_1497_1738_English.texfrag


% BEGIN INCLUDED SGA4_English_current_buildable_source_20260906/source/01/drafts/SGA4_I_LINES_1739_1946_English.texfrag
\section{Faithful functors and conservative functors}\etiquette[6]{I.sec6}\pageoriginale

\begin{enonce*}{Definition 6.1}\etiquette[6.1]{I.def6.1}
Let $E$ be a category and let $(\varphi_i:E\rightarrow F_i)_{i\in I}$
be a family of functors
$$
\varphi_i:E\longrightarrow F_i{\quoi}.
$$
The family of functors $(\varphi_i)$ is said to be faithful if, for every
pair of objects $X$, $Y$ of $E$ and every pair of arrows
$u$, $v:X\rightrightarrows Y$, the relation
$\varphi_i(u)=\varphi_i(v)$ for every $i\in I$ implies $u=v$ (in other
words, if the map $\Hom(X,Y)\rightarrow
\Pi_i\Hom(\varphi_i(X),\varphi_i(Y))$ defined by $(\varphi_i)$ is
injective). The family of functors $(\varphi_i)$ is said to be
conservative if every arrow $u$ of $E$ for which $\varphi_i(u)$ is an
isomorphism for every $i\in I$ is an isomorphism. We say that
$(\varphi_i)$ is conservative for monomorphisms (\resp for epimorphisms,
\resp \ldots) if the preceding condition holds whenever $u$ is a
monomorphism (\resp an epimorphism, \resp \ldots).
\end{enonce*}

\setcounter{subsection}{1} \setcounter{subsubsection}{0}
\subsubsection{}\etiquette[6.1.1]{I.sec6.1.1}
If we introduce the unique functor
$$
\varphi:E\longrightarrow F=\prod_{i\in I} F_i
$$
defined by the family of functors $(\varphi_i)$, it is clear that this
family is faithful (\resp conservative, \resp conservative for
monomorphisms, \resp\ldots) if and only if the functor $\varphi$ is
faithful (\resp conservative, \resp\ldots) (by which we mean that the
family consisting only of $\varphi$ is faithful, \resp conservative,
\resp\ldots). Thus, without major inconvenience, we could restrict
ourselves below to a family consisting of a single functor. Nevertheless,
for ease of future reference, we shall state the following results for
families.

The\pageoriginale notions of \hyperref[I.def6.1]{6.1} are especially useful when the
$\varphi_i$ satisfy suitable exactness properties, and in that case they
tend to coincide:

\begin{enonce*}{Proposition 6.2}\etiquette[6.2]{I.prop6.2}
The notation is that of \hyperref[I.prop5.1]{5.1}.
\begin{enumerate}
\renewcommand{\theenumi}{\roman{enumi}}
\renewcommand{\labelenumi}{(\theenumi)}
\item If the kernels of double arrows, or the cokernels of double arrows,
  are representable in $E$, and if the $\varphi_i$ commute with them,
  then one has the implication
$$
(\varphi_i)\mbox{ conservative } \Longrightarrow (\varphi_i)\mbox{
    faithful}{\quoi}.
$$

\item Suppose that fiber products (\resp amalgamated sums) are
  representable in $E$, and that the $\varphi_i$ commute with them.
  Suppose that $(\varphi_i)$ is faithful or conservative; then for every
  arrow $u$ of $E$, $u$ is a monomorphism (\resp an epimorphism) if and
  only if the same holds for $\varphi_i(u)$ for every $i\in I$.

\item Suppose that fiber products and amalgamated sums are representable
  in $E$ and that the $\varphi_i$ commute with them, and that every arrow
  in $E$ which is a bimorphism (\ie both a monomorphism and an
  epimorphism) is an isomorphism. Then one has the implication
$$
(\varphi_i)\mbox{ faithful } \Longrightarrow (\varphi_i)\mbox{
    conservative}{\quoi}.
$$

\item Suppose that fiber products (\resp amalgamated sums) are
  representable in $E$, and that the $\varphi_i$ commute with them. Then,
  if $(\varphi_i)$ is conservative for monomorphisms (\resp for
  epimorphisms), it is in fact conservative.

\item Let $D$ be a type of diagram, let $F:d\mapsto F(d)$ be a diagram of
  type $D$ in $E$, let $X$ be an object of $E$, and let
  $u=(u_d)_{d\in D}$ be a family of arrows
$$
X\longrightarrow F(d)\quad (\mbox{\resp  } F(d)\longrightarrow
X){\quoi}.
$$
Suppose\pageoriginale that $(\varphi_i)$ is conservative, that projective
(\resp inductive) limits of type $D$ are representable in $E$, and that
the $\varphi_i$ commute with them. Then $u$ makes $X$ a projective
(\resp inductive) limit of $F$ in $E$ if and only if, for every $i\in I$,
$\varphi_i(u)$ makes $\varphi_i(X)$ a projective (\resp inductive) limit
of $\varphi_i(D)$ in $E_i$.
\end{enumerate}
\end{enonce*}

\noindent \emph{Proof.}
\begin{enumerate}
\renewcommand{\theenumi}{\roman{enumi}}
\renewcommand{\labelenumi}{(\theenumi)}
\item For the non-\resp statement, given a double arrow
$u$, $v:X\rightrightarrows Y$, it suffices to express the equality $u=v$
by the condition that the inclusion $\Ker(u,v)\rightarrow X$ is an
isomorphism. Here and below, we omit repetition of the dual argument for
the dual statement.

\item If $(\varphi_i)$ is faithful, the condition that
$u:X\rightarrow Y$ be a monomorphism is expressed by the equality
$\pr_1=\pr_2$ for the fiber product $\fprod{X}{X}{Y}$. If $(\varphi_i)$
is conservative, it is expressed by the condition that the diagonal
morphism $\delta:X\rightarrow \fprod{X}{X}{Y}$ be an isomorphism.

\setcounter{enumi}{3}
\item Since in the latter argument the morphism $\delta$ is a
monomorphism, we see that it would in fact have been enough to suppose
that $(\varphi_i)$ is conservative for monomorphisms. But this then
implies that $(\varphi_i)$ is conservative without qualification.
Indeed, if $u\in %%
F\ell\,E$ is such that the $\varphi_i(u)$ are isomorphisms, the preceding
argument shows that these are monomorphisms, and hence they are
isomorphisms by the hypothesis on $(\varphi_i)$.

\setcounter{enumi}{2}
\item This is a trivial consequence of (ii).

\setcounter{enumi}{4}
\item This is a trivial consequence of the definitions.
\end{enumerate}

We record the following consequence of (i), (ii), and (iv):

\begin{enonce*}{Corollary 6.3}\etiquette[6.3]{I.cor6.3}
Suppose\pageoriginale that fiber products and amalgamated sums are
representable in $E$ and that the $\varphi_i$ commute with them, and that
kernels of double arrows or cokernels of double arrows are representable
and that the $\varphi_i$ commute with them. (It suffices, for example,
that finite projective limits and finite inductive limits be representable
in $E$, and that the $\varphi_i$ be exact functors.) Then the following
conditions are equivalent:
\begin{itemize}
\item[a)] $(\varphi_i)$ is faithful.

\item[b)] $(\varphi_i)$ is conservative.

\item[c)] $(\varphi_i)$ is conservative for monomorphisms.

\item[${\rm c')}$] $(\varphi_i)$ is conservative for epimorphisms.
\end{itemize}
\end{enonce*}
We also record, for reference:

\begin{enonce*}{Proposition 6.4}\etiquette[6.4]{I.prop6.4}
Let $\varphi:E\rightarrow F$ be a functor with a right adjoint $\psi$
(thus $\Hom(\varphi(X),Y)\simeq \Hom(X,\psi(Y))$). The functor $\varphi$
(\resp $\psi$) is faithful if and only if, for every element $X$ of $E$
(\resp every element $Y$ of $F$), the unit morphism
$X\rightarrow \psi\varphi(X)$ is a monomorphism. The functor $\varphi$
(\resp $\psi$) is fully faithful if and only if the preceding unit
morphism is an isomorphism.
\end{enonce*}
Indeed, if $X'$, $X$ are two objects of $E$, the map
\begin{equation*}
\Hom(X',X)\longrightarrow
  \Hom(\varphi(X'),\varphi(X))\tag{$\ast$}
\end{equation*}
is identified with the map obtained from
\begin{equation*}
X\longrightarrow \psi
  \varphi (X)\tag{${\ast}{\ast}$}
\end{equation*}
by applying the functor $\Hom(X',-)$. Thus, for $(\ast)$ to be a
monomorphism\pageoriginale (\resp an isomorphism) for every $X'$, with
$X$ fixed, it is necessary and sufficient that the unit morphism
$({\ast}{\ast})$ be a monomorphism (\resp an isomorphism).

\begin{enonce*}{Proposition 6.5}\etiquette[6.5]{I.prop6.5}
Let $p:E'\rightarrow E$ be a functor. The following conditions are %%
equivalent:
\begin{enumerate}
\renewcommand{\theenumi}{\roman{enumi}}
\renewcommand{\labelenumi}{(\theenumi)}
\item $p$ is faithful, conservative, and fibered {\rm(SGA 1 VI 6.1).}

\item $p$ is a fibered functor whose fibers are discrete categories.

\item For every $X'\in \ob E'$, the functor $E'_{/X'}\rightarrow
  E_{/p(X)}$ induced by $p$ is an equivalence of categories that is
  surjective on objects.

\item (When $E$ is a $𝒰$-category.) There exists an element
  $F\in \ob\widehat{C}$ and an equivalence of categories over $E$
  {\rm(SGA 1 VI 4.3)} $E'\xrightarrow{\approx}E_{/F}$ (where $E_{/F}$
  is the full subcategory of $\widehat{E}_{/F}$ formed by the arrows
  $X\rightarrow F$ whose source belongs to $E$).
\end{enumerate}
\end{enonce*}

\setcounter{subsection}{5} \setcounter{subsubsection}{0}

% END INCLUDED SGA4_English_current_buildable_source_20260906/source/01/drafts/SGA4_I_LINES_1739_1946_English.texfrag


% BEGIN INCLUDED SGA4_English_current_buildable_source_20260906/source/01/drafts/SGA4_I_LINES_1947_2034_English.texfrag
\subsubsection{}\etiquette[6.5.1]{I.sec6.5.1}
Recall that a category $C$ is said to be \emph{discrete} if it is a
\emph{groupoid} (\ie every arrow in it is invertible) and if it
is \emph{rigid} (\ie the automorphism group of every object
is reduced to the trivial group); equivalently, the
category is equivalent to the category $C'$ defined by a
set $I$ (with $\text{ob}C'=I$, and with the identity
arrows as its only arrows). When $C$ is already assumed to be a
groupoid, saying that $C$ is discrete amounts to saying that for
two objects $X$, $Y$ of $C$, there exists at most one arrow from $X$
to $Y$, \ie that $C$ is isomorphic to the category defined by a
preordered set.

The equivalence\pageoriginale of conditions (i) and (ii) of
\hyperref[I.prop6.5]{6.5} is an immediate consequence of the preceding observations
and of the following

\begin{enonce*}{Lemma 6.5.2}\etiquette[6.5.2]{I.lem6.5.2}
Let $p:E'\rightarrow E$ be a fibered functor. Then:
\begin{enumerate}
\renewcommand{\theenumi}{\roman{enumi}}
\renewcommand{\labelenumi}{(\theenumi)}
\item For $p$ to be conservative, it is necessary and sufficient that its
  fiber categories be groupoids.

\item For $p$ to be faithful, it is necessary and sufficient that its
  fiber categories be ordered categories.
\end{enumerate}
\end{enonce*}

\noindent \emph{Proof of} \hyperref[I.lem6.5.2]{6.5.2}.
\begin{enumerate}
\renewcommand{\theenumi}{\roman{enumi}}
\renewcommand{\labelenumi}{(\theenumi)}
\item Suppose $p$ is conservative. For every arrow $u$ in a fiber
  $E'_X$, $p(u)=\id_X$ is an isomorphism, so $u$ is an
  isomorphism in $E'$, hence also in $E'_X$ (since an inverse of
  $u$ in $E'$ will clearly be an inverse in $E'_X$). Thus $E'_X$
  is a groupoid. Conversely, suppose that the $E'_X$ are
  groupoids, and let $u'$ be an arrow of $E'$ such that $p(u')$
  is an isomorphism; let us prove that $u$ is an isomorphism. To
  this end, note that, since $p$ is fibered, we can factor
  $u':X'\rightarrow Y'$ as a composite $X'\rightarrow
  u^\ast(Y')\rightarrow Y'$, where the first arrow is an
  $X$-morphism (N.B. $X=p(X')$, $u=p(u')$), and the second is a
  cartesian morphism over $u$. The first arrow is
  an isomorphism since $E'_X$ is a groupoid, and the second
  is as well, since a \emph{cartesian} %%
morphism of a fibered category
  is clearly an isomorphism whenever its projection is an isomorphism.

\item Suppose $p$ is faithful, and let $X'$, $Y'$ be two objects of a
  fiber category $E_X$. Then two arrows from $X'$ to $Y'$ lie
  over the same arrow $\id_X$ of $E$, and hence are identical,
  so $E_X$ is ordered. Conversely,\pageoriginale suppose that the
  fiber categories are ordered, and let us prove that $p$ is
  faithful. Thus let $u'$, $v':X'\rightrightarrows Y'$ be
  arrows of $E'$ lying over the same arrow $u:X\rightarrow Y$
  of $E$. They then factor as $X'\rightrightarrows
  u^\ast(Y)\rightarrow Y'$, where the two arrows
  $X'\rightrightarrows u^\ast(Y')$ are arrows of $E'_X$ with
  the same source and target; they are therefore equal, and hence
  $u'=v'$, \cqfd
\end{enumerate}

Let us return to the proof of \hyperref[I.prop6.5]{6.5}. We have proved (i)
$\Leftrightarrow$ (ii). On the other hand, (ii) %%
$\Leftrightarrow$ (iv) is fairly clear: indeed, on the one hand, the
category $E_{/F}$ is fibered over $E$, with fiber categories the
discrete categories defined by the sets $F(X)$, as
follows immediately from the definitions; on the other hand, if $p$ is as
in (ii), then by the sorites of SGA 1 VI 8 the fibered category
$E'$ over $E$ is $E$-equivalent to the split category over $E$
defined by the functor $E\rightarrow (\Cat)$ determined by the
functor $F:E\rightarrow (\Ens)$, which assigns to every $X\in \ob E$
the set of isomorphism classes of objects of $E'_X$. Now this
split category is $E$-isomorphic to the category $E_{/F}$. Since
(iv) $\Rightarrow$ (iii) is clear, it remains to prove (iii)
$\Rightarrow$ (i). Now it is clear that for $p$ to be faithful
(\resp conservative), it is necessary and sufficient that the induced functors
$E'_{/X'}\rightarrow E_{/p(X)}$ be so; \textit{a fortiori}, it
suffices that these be fully faithful; therefore it remains only to
prove that (iii) implies that $p$ is fibered. But one
again sees that a functor $p$ is fibered if and only if the
induced functors $E'_{/X'}\rightarrow E_{/p(X')}$ are
fibered. This is in particular the case if they are equivalences of categories
that are surjective on objects.


% END INCLUDED SGA4_English_current_buildable_source_20260906/source/01/drafts/SGA4_I_LINES_1947_2034_English.texfrag


% BEGIN INCLUDED SGA4_English_current_buildable_source_20260906/source/01/drafts/SGA4_I_LINES_2035_2353_English.texfrag
\section{Generating and cogenerating
  subcategories}\etiquette[7]{I.sec7}\pageoriginale

\begin{enonce*}{Definition 7.1}\etiquette[7.1]{I.def7.1}
Let $E$ be a category and $C$ a full subcategory of $E$. We
say that $C$ is a subcategory of $E$ \emph{generating by strict
epimorphisms (\resp  by epimorphisms)} if, for every object
$X$ of $E$, the family of arrows of $E$ with target $X$ and source
$X'\in \ob C$ is epimorphic (\resp  strictly epimorphic)
(\hyperref[I.condef1.3]{1.3}). We say that $C$ is a subcategory of $E$
\emph{generating (\resp  generating for monomorphisms, \resp
generating for strict monomorphisms)} if, for every arrow
$u:Y\rightarrow X$ (\resp  every monomorphism of $E$, \resp  every
strict monomorphism of $E$ $\hyperref[I.sec10.5]{(10.5)}$) such that, for every $X'\in
\ob C$, the corresponding map $\Hom(X',Y)\rightarrow
\Hom(X',X)$ is bijective, $u$ is an isomorphism. Finally, we say
that a family $(X_i)$ of objects of $E$ is generating by
strict epimorphisms (\resp \ldots) if the full subcategory $C$
of $E$ generated by this family is generating by strict epimorphisms
(\resp \ldots).
\end{enonce*}

\setcounter{subsection}{1} \setcounter{subsubsection}{0}
\subsubsection{}\etiquette[7.1.1]{I.sec7.1.1}
Observe that, in terms of the family $(h_{X'})_{X'\in \ob C}$ of
functors
$$
h_{X'}:E\longrightarrow (\Ens)\; (X'\in \ob C)
$$
represented by the $X'\in \ob C$, the condition that
$C$ be generating (\resp  generating for
monomorphisms,
\resp  generating for strict monomorphisms) can be expressed by
saying that the family $(h_{X'})$ is \emph{conservative} (\resp  conservative
for monomorphisms, \resp  conservative for strict monomorphisms)
(\hyperref[I.def6.1]{6.1}). It also follows immediately from the
definitions that $C$ is generating by epimorphisms
if\pageoriginale and only if the family $(h_{X'})_{X'\in \ob C}$
is \emph{faithful} (\hyperref[I.def6.1]{6.1}). We shall also give below
(\hyperref[I.prop7.2]{7.2} (i)) an analogous interpretation of the condition
that $C$ be generating by strict epimorphisms.

\subsubsection{}\etiquette[7.1.2]{I.sec7.1.2}
As with the notions introduced in \hyperref[I.def6.1]{6.1}, the notions
of \hyperref[I.def7.1]{7.1} are especially useful when $E$ has suitable
exactness properties, in which case the various notions
introduced have a marked tendency to be all equivalent
(\hyperref[I.cor7.3]{7.3}). Thus the question of which of
these notions in \hyperref[I.def7.1]{7.1} should be regarded as the most
important scarcely arises; in the most important cases, these
notions coincide, and the term ``generating
subcategory''
may therefore be interpreted indifferently as
referring to
any one of the properties considered in \hyperref[I.def7.1]{7.1}
(for example the first, which is the strongest of all, as we
shall see (\hyperref[I.prop7.2]{7.2} (ii))).

\subsubsection{}\etiquette[7.1.3]{I.sec7.1.3}
Suppose that $E$ is a $𝒰$-category, and consider the
canonical functor
\begin{equation*}
\varphi:E\longrightarrow
  \widehat{C}=\Hhom(C, 𝒰\text{-}\Ens)\tag{7.1.3.1}
\end{equation*}\etiquette[7.1.3.1]{I.sec7.1.3.1}
which is the composite of the functors $E\rightarrow \widehat{E}\rightarrow
\widehat{C}$, where the first functor is
the canonical functor (\hyperref[I.sec1.3.3]{1.3.3}), and
the second is
the restriction functor to $C$. Observe that it is plainly
equivalent to say that the preceding functor $\varphi$ is
conservative (\resp  faithful), or to say that the family of functors
$h_{X'}:X\rightarrow \Hom(X',X)=\varphi(X)(X')$, for variable $X'\in\ob C$,
is a conservative (\resp  faithful) family, that is,
also (\hyperref[I.sec7.1.2]{7.1.2}), that $C$ is generating (\resp
generating by
epimorphisms). Similarly, $\varphi$ is\pageoriginale conservative for
monomorphisms (\resp  for strict monomorphisms) if and
only if the family of the $h_{X'}(X'\in \ob C)$ is conservative for
monomorphisms (\resp  for strict monomorphisms), \ie if
and only if the subcategory $C$ of $E$ is generating for
monomorphisms (\resp  for strict monomorphisms).

\begin{enonce*}{Proposition 7.2}\etiquette[7.2]{I.prop7.2}
Let $E$ be a $𝒰$-category and $C$ a full subcategory.
\begin{enumerate}
\renewcommand{\theenumi}{\roman{enumi}}
\renewcommand{\labelenumi}{(\theenumi)}
\item The following conditions are equivalent:
\begin{enumerate}
\renewcommand{\theenumii}{\alph{enumii}}
\renewcommand{\labelenumii}{\theenumii)}
\item $C$ is a subcategory generating by strict
  epimorphisms.

\item For every $X\in \ob E$, let $C_{/X}$ denote the
  full subcategory of $E_{/X}$ consisting of arrows
  $X'\rightarrow X$ with source $X'\in \ob C$. The natural arrow from the
  inclusion functor $j:C_{/X}\rightarrow E$ to the constant functor
  on $C_{/X}$ with value $X$ makes $X$ an inductive limit
  of $j$:
$$
X\xleftarrow{\sim} \lim_{C_{/X}}X'{\quoi}.
$$

\item The canonical functor $\varphi$ of (\hyperref[I.sec7.1.3.1]{7.1.3.1}) is fully
  faithful.
\end{enumerate}

\item Among the notions of \hyperref[I.def7.1]{7.1} there are the following
implications~:
\[
\xymatrix@C=-.7cm@R=.45cm{
  &  \txt{{\rm 1)} $C$ generating by\\ strict epimorphisms\\
  \emph{($\varphi$ fully faithful)}} \ar@{=>}[rd]&\\
 \txt{{\rm 2)} $C$ generating by\\ epimorphisms\\
 \emph{($\varphi$ faithful)}} &    &
 \txt{{\rm 3)} $C$ generating \phantom{aaa}\\
 \emph{($\varphi$ conservative)}} \ar@{=>}[d]\\
    &      &  \txt{{\rm 4)} $C$ generating for\\ monomorphisms\\
  \emph{($\varphi$ conservative for mon.)}}\ar@{=>}[d]\\
    &      &  \txt{{\rm 5)} $C$ generating for\\ strict monomorphisms\\
  \emph{($\varphi$ conservative for strict mon.).}}\\
}
\]

\item There\pageoriginale are the following conditional implications:
\begin{enumerate}
\renewcommand{\theenumii}{\alph{enumii}}
\renewcommand{\labelenumii}{\theenumii)}
\item If in $E$ epimorphic families of arrows are
  strictly epimorphic, then $2)\Rightarrow 1)$. If in $E$
  monomorphisms are strict, then $5)\Rightarrow 4)$.

\item If in $E$ the kernels of pairs of arrows (\resp  fiber
  products) are representable, then $3)\Rightarrow
  2)$ (\resp  $4)\Rightarrow 3)$).

\item If in $E$ every family of morphisms $X_i\rightarrow X$ with
  common target $X$ factors as a strictly epimorphic
  (\resp  epimorphic) family $X_i\rightarrow Y$ followed by a monomorphism
  (\resp  by a strict monomorphism) $Y\rightarrow X$, then
  $4)\Rightarrow 1)$ (\resp  $5)\Rightarrow 2)$).
\end{enumerate}
\end{enumerate}
\end{enonce*}
Let us immediately record the following.

\begin{enonce*}{Corollary 7.3}\etiquette[7.3]{I.cor7.3}
All the notions considered in $\hyperref[I.def6.1]{6.1}$ (and repeated in the
implication diagram in {\rm(ii)} above) are equivalent
in each of the following two cases:
\begin{enumerate}
\renewcommand{\theenumi}{\roman{enumi}}
\renewcommand{\labelenumi}{(\theenumi)}
\item In $E$, the kernels of pairs of arrows and the fiber
  products are representable, the monomorphisms are strict, and
  the epimorphic families of arrows are strictly
  epimorphic.

\item In $E$, every family $(X_i\rightarrow X)_{i\in I}$ of
  arrows with common target $X$ factors as an
  epimorphic family $(X_i\rightarrow Y)$ followed by a monomorphism
  $Y\rightarrow X$, every monomorphism of $E$ is strict, and every
  epimorphic family of arrows of $E$ is strict.
\end{enumerate}
\end{enonce*}

Indeed,\pageoriginale in case (i) we have 4) $\Rightarrow$ 3) and
3) $\Rightarrow$ 2) by b), and 5) $\Rightarrow$ 4) and 2)
$\Rightarrow$ 1) by a). In case (ii), by c), we have the
implications 4) $\Rightarrow$ 1) and 5) $\Rightarrow$ 2). We therefore
conclude by the implication diagram \hyperref[I.prop6.2]{6.2} (ii).

\medskip
\noindent \emph{Proof of} \hyperref[I.prop7.2]{7.2}.
\begin{enumerate}
\renewcommand{\theenumi}{\roman{enumi}}
\renewcommand{\labelenumi}{(\theenumi)}
\item The implication b) $\Rightarrow$ a) follows at once from the
  definitions. Let us prove a) $\Rightarrow$ b). Thus, under
  hypothesis a), we must prove that, for all $X$, $Y\in \ob E$,
  every system of arrows
$$
u_{X'}:X'\longrightarrow Y
$$
indexed by the $X'\in \ob C_{/X}$, such that
$u_{X'}f=u_{X''}$ for every arrow $f:X''\rightarrow X'$ in
$C_{/X}$, factors through an arrow (necessarily unique by
hypothesis a)) $X\rightarrow Y$. By hypothesis a), it suffices
to check that, for every object $Z$ of $E_{/X}$ and every pair of
morphisms $v':Z\rightarrow X'$, $v'':Z\rightarrow X''$ in
$E_{/X}$, with $X'$ and $X''$ in $C_{/X}$, we have
$u_{X'}v'=u_{X''}v''$. But, by hypothesis a), the family of
arrows $w:X'''\rightarrow Z$, with $X'''\in \ob C$, is
epimorphic, and it therefore suffices to check that, for every such
$w$, we have $(u_{X'}v')w=(u_{X''}v'')w$, which can also be written
$u_{X'}(v'w)=u_{X''}(v''w)$ and follows at once from the hypothesis made
on the family of the $u$.

Let us now prove the equivalence of conditions b) and c. For this,
observe that, for every $X\in \ob E$, the object $\varphi(X)$ of
$\widehat{C}$ is the inductive limit in $\widehat{C}$ of the canonical
functor $\widehat{C}_{/\varphi(X)}\rightarrow \widehat{C}$ (\hyperref[I.prop3.4]{3.4});
but $\widehat{C}_{/\varphi(X)}$ is canonically isomorphic to
$C_{/X}$, so that in $\widehat{C}$ we have
$$
\varphi(X)=\varinjlim_{C_{/X}}X'{\quoi},
$$
where\pageoriginale the object $X'$ of $C$ is identified with the functor
$\in \widehat{C}$ that it represents. Consequently, for a second
object $Y$ of $E$, we have a canonical isomorphism
\begin{align*}
\Hom(\varphi(X),\varphi(Y))\simeq\varprojlim_{C_{/X}}
\Hom(X',\varphi(Y))&\simeq
\varprojlim_{C_{/X}}\varphi(Y)(X')\hyperref[I.prop1.4]{(1.4)}\\
&\simeq \varprojlim_{C_{/C}} \Hom(X'%%
,Y){\quoi}.\tag{$\ast$}
\end{align*}
This being so, the map
$$
\Hom(X,Y)\longrightarrow \Hom(\varphi(X),\varphi(Y))
$$
defined by $\varphi$ is none other, via the isomorphism between the
extreme terms of $(\ast)$, than the map
$$
\Hom(X,Y)\longrightarrow \varprojlim_{C_{/X}} \Hom(X',Y)
$$
deduced from the inductive system of arrows $X'\rightarrow X$ indexed by
$C_{/X}$ considered in \hyperref[I.prop7.2]{7.2} (i) b). Thus the first
map is bijective for every $Y$ (with $X$ fixed) if and
only if $X$ is an inductive limit of the inclusion functor
$j:C_{/X}\rightarrow E$, which proves the equivalence of b) and c).

\item The implications 1) $\Rightarrow$ 2) and 3) $\Rightarrow$ 4)
  $\Rightarrow$ 5) are trivial by virtue of the
  definitions. The implication 1) $\Rightarrow$ 3) is obtained by
  interpreting property 1) as the full faithfulness of
  $\varphi$ by (i), and by observing that a fully
  faithful functor is conservative. But we have already observed
  (\hyperref[I.sec7.1.1]{7.1.1}) that 3)
  means that $\varphi$ is conservative.

\item Assertion\pageoriginale a) is a tautology. Assertion
  b) follows from \hyperref[I.prop6.2]{6.2} (i) (\resp  \hyperref[I.prop6.2]{6.2} (iv)) applied to the functor
  $\varphi$, taking into account that the latter is left exact. Finally, let us prove
  c). For $X\in \ob E$, consider the family of all
  the morphisms $X_i\rightarrow X$, with $X_i\in \ob C$; by
  hypothesis on $E$ it factors as a family $X_i\rightarrow
  Y$ that is effective epimorphic (\resp  epimorphic), followed by a
  monomorphism (\resp  by a strict monomorphism) $Y\rightarrow
  X$. Then hypothesis 4) (\resp  5)) implies that $Y\rightarrow
  X$ is an isomorphism, so the family considered is
  strictly epimorphic (\resp  epimorphic), \cqfd
\end{enumerate}

\begin{enonce*}{Proposition 7.4}\etiquette[7.4]{I.prop7.4}
Let $E$ be a category, $C$ a full subcategory of $E$, and $X$
an object of $E$. Suppose that $C$ is generating in $E$ (\resp  $C$
is generating for monomorphisms, and that the fiber product of two
subobjects of $X$ over $X$ is representable in $E$). Then a
strict subobject $\hyperref[I.sec10.11]{(10.11)}$ (\resp  a subobject) $X'$ of $X$ is
determined %%
when, for every $T\in \ob C$, one knows the subset of $\Hom(T,X)$
which is the image of $\Hom(T,X')$. Consequently, the cardinality of the set of
strict subobjects (\resp  of the set of subobjects) of $X$
is bounded above by $\prod_{T\in \ob C} 2^{\card
  \Hom(T,X)}$, and if $X\in \ob C$, it is bounded above by $2^{\card%%
F\ell\,C}$.
\end{enonce*}

Let us first prove the parenthetical assertion. Let $X'$ and $X''$ be two
subobjects of $X$ such that, for every $T\in \ob C$, the images of
$\Hom(T,X')$ and $\Hom(T,X'')$ in $\Hom(T,X)$ are equal. They
are therefore also equal to the image of $\Hom(T,X''')$, where $X'''$ is the
fiber product of $X'$ and $X''$ over $X$. Since $C$ is generating
for monomorphisms, it follows that the
monomorphisms\pageoriginale $X'''\rightarrow X'$ and
$X'''\rightarrow X''$ are isomorphisms, so $X'$ and $X''$
are equal, each being equal to the subobject $X'''$ of $X$.

Let us prove the nonparenthetical assertion. By definition of the notion of
strict subobject, it suffices to check that knowledge of the
subset $\Hom(T,X')$ of $\Hom(T,X)$ for every $T\in\ob C$ implies
knowledge of those pairs of arrows $\xymatrix{
X\ar@<2pt>[r]^{u,v}\ar@<-2pt>[r]& T }$ such that $ui=vi$, where
$i:X'\rightarrow X$ is the canonical injection. But since $C$ is
generating, the relation $ui=vi$ is equivalent to the relation
$(ui)f=(vi)f$ for every $f\in \Hom(T,X')$, \ie to $ug=vg$ for every
$g\in \Hom(T,X)$ coming from $\Hom(T,X')$ (\ie of the form if, with
$f\in \Hom(T,X')$), which proves our assertion.

\begin{enonce*}{Corollary 7.5}\etiquette[7.5]{I.cor7.5}
Let $E$ be a category, $C$ a full generating subcategory,
and
$X$ an object of $C$. Then a strict quotient (\hyperref[I.sec10.8]{10.8}) $X'$ of $X$
is determined when, for every $T\in \ob C$, one knows the subset of
$\Hom(T,X)^2$ consisting of pairs $(u,v)$ such that $qu=qv$, where
$q:X\rightarrow X'$ is the canonical morphism. Thus the cardinality of
the set of strict quotients of $X$ is bounded above by $\prod_{T\in
\ob C} 2^{\card\Hom(T,X)^2}$.
\end{enonce*}

Indeed, by %%
definition, a strict quotient $X'$ of $X$ is determined when one
knows, for every object $Y$ of $E$, the subset of $\Hom(Y,X)\times
\Hom(Y,X)$ consisting of pairs $(u,v)$ such that $qu=qv$, where
$q:X\rightarrow X'$ is the canonical morphism. But the relation
$qu=qv$ is equivalent to the relation $(qu)f=(qv)f$ for every morphism
$f:T\rightarrow Y$ with source $T$ in $C$, since $C$ is
generating. This relation may also be written $q(uf)=q(vf)$, which
proves the first assertion of \hyperref[I.cor7.5]{7.5}. The second follows
at once.

\setcounter{subsection}{5} \setcounter{subsubsection}{0}

% END INCLUDED SGA4_English_current_buildable_source_20260906/source/01/drafts/SGA4_I_LINES_2035_2353_English.texfrag


% BEGIN INCLUDED SGA4_English_current_buildable_source_20260906/source/01/drafts/SGA4_I_LINES_2354_2668_English.texfrag
\subsubsection{}\etiquette[7.5.1]{I.sec7.5.1}
We\pageoriginale can generalize \hyperref[I.cor7.5]{7.5} by
introducing, for a
family of objects $(X_i)_{i\in I}$ of $E$, the notion of a
\emph{strict quotient
  in $E$ of the family}, by which we mean a strictly epimorphic family
$(p_i:X_i\rightarrow X')_{i\in I}$ of
morphisms of $E$, --- with the understanding, as for ordinary
quotients, that we identify two such families
$(p_i:X_i\rightarrow X')$ and $(q_i:X_i\rightarrow X'')$ if one can
find an isomorphism $v:X'\rightarrow X''$ (necessarily
unique) such that $fp_i=q_i$ for every $i\in I$. With this
terminology, the proof of \hyperref[I.cor7.5]{7.5} applies {\it ne
varietur} to give the

\begin{enonce*}[remark]{Variant 7.5.2}\etiquette[7.5.2]{I.var7.5.2}
Let $E$, $C$ be as in \hyperref[I.cor7.5]{7.5}, and let $(X_i)_{i\in I}$ be a family
of objects of $E$. Then a strict quotient $X'$ of $(X_i)_{i\in I}$
in $E$ (\hyperref[I.sec7.5.1]{7.5.1}) is determined by the following data: for every
$T\in \ob C$,
and every pair $(i,j)\in I\times I$, the subset of $\Hom(T,X_i)\times
\Hom(T,X_j)$ consisting of the pairs $(u,v)$ such that $p_iu=p_jv$, where
for every $i\in I$, $p_i:X_i\rightarrow X'$ denotes the canonical
morphism. Consequently, the cardinality of the set of strict
quotients of $(X_i)_{i\in I}$ in $E$ is bounded above by $ \prod_{T\in
\ob C} \prod_{i, j\in I}2^{\card \Hom(T,X_i)\times
  \card\Hom(T,X_j)}{\quoi}.
$
\end{enonce*}

\setcounter{subsubsection}{2}
\subsubsection{}\etiquette[7.5.3]{I.sec7.5.3}
It is immediately clear that the analogous conclusion holds if one
assumes only that $C$ is generating for strict monomorphisms,
provided that the products $X_i\times X_j$
are representable and that one restricts to the \emph{effective} quotients
of the family $(X_{i})_{i\in I}$, \ie to the strict quotients
$X'$ such that the fiber products $\fprod{X_i}{X_{j}}{X'}$
are representable in $E$ (which is no restriction if
$E$ is stable under fiber products).

\begin{enonce*}{Proposition 7.6}\etiquette[7.6]{I.prop7.6}
Let\pageoriginale $E$ be a category, $C$ a full subcategory
of $E$ generating by epimorphisms
(\hyperref[I.def7.1]{7.1}), $\Pi_\circ=$ an
infinite cardinal $\geq \card  F\ell  C$, $\Pi\geq \Pi_\circ$ a
cardinal, and $(u_i:X_i\rightarrow X)_{i\in I}$ a universally epimorphic
family (\hyperref[I.sec10.3]{10.3}) in $E$ consisting of pullbackable morphisms
(\hyperref[I.sec10.7]{10.7}) with sources $X_i\in\ob C$, and such that $\card I\leq \Pi$.
Then $\card F\ell  C_{/X}\leq \Pi^{\Pi_\circ}$.
\end{enonce*}

Let $T\in \ob C$; it will suffice to prove that
\begin{equation*}
\card\Hom(T,X)\leq
  \Pi^{\Pi_\circ}{\quoi},\tag{7.6.1}
\end{equation*}\MakeLinkTarget*{sga4.exp.I.equation.7.6.1}\label{I.eq7.6.1}
for one then obtains successively
$$
\card\ob C_{/X}\leq(\Pi^{\Pi_\circ})\times\card \ob C_{/X}\leq
(\Pi^{\Pi_\circ})\times \Pi_\circ=\Pi^{\Pi_\circ}{\quoi},
$$
and finally $\card F\ell C_{/X}\leq ((\Pi^{\Pi_\circ})^2)\times
\Pi_\circ=\Pi^{\Pi_\circ}$, since for two objects $S$, $T$ of
$C_{/X}$, one has $\card \Hom_{C_{/X}}(S,T)\leq
\car\Hom_C(S,T)\leq\Pi_\circ$. To prove \ref{I.eq7.6.1}, first note the

\begin{enonce*}{Lemma 7.6.2}\etiquette[7.6.2]{I.lem7.6.2}
Let $J$ be a set such that $\card J=\Pi_\circ$. For every object
$T$ of $C$, and every morphism $f:T\rightarrow X$, there exists an
epimorphic family $(v_j:S_j\rightarrow T)_{j\in J}$, with sources
$S_j\in\ob C$, and, for every $j\in J$, an $i(j)\in I$ and a
$g_j:S_j\rightarrow X_{i(j)}$ such that $u_{i(j)}g_j=fv_j$.
\end{enonce*}

Indeed, the family of $\fprod{X_i}{T}{X}\rightarrow$ is
epimorphic by hypothesis; on the other hand, since $C$ is
generating
by strict epimorphisms, for every $i$ the family of arrows
$S\rightarrow \fprod{X_i}{T}{X}$ with source $S\in \ob C$ is
epimorphic, and hence, by transitivity, the family of arrows
$v:S\rightarrow T$ with source in $\ob C$ which factor through one
of the $\fprod{X_i}{T}{X}$ is epimorphic. These are precisely the arrows
$v:S\rightarrow T$ with source in $C$ for which there exist an
$i\in I$ and a $g:S\rightarrow X_i$ such that $u_ig=fv$. Since
the set of\pageoriginale these arrows $v:S\rightarrow T$ is
contained in $F\ell$ $C$, and therefore has cardinality $\leq \Pi_\circ$, it can
be indexed by the index set $J$, whence the lemma.

Now note that a morphism $f:T\rightarrow X$ is determined once
the $fv_j$ are known, and these are determined once the $g_j$ are known.
Thus $\card\Hom(T,X)$ is bounded above by the cardinality of the set of
families $(i(j),S_j,v_j,g_j)_{j\in J}$; since the cardinality of
the set of maps $j\rightarrow I$ is $\leq
\Pi^{\Pi_\circ}$, and since, for a fixed such map $j\mapsto
i(j)$, the cardinality of the set of corresponding families
$S_j$, $f_j$, $v_j$ is bounded above by that of the set of
maps from $J$ to $F\ell  C\times F\ell  C$, which is $\leq
\Pi^{\Pi_\circ}_\circ$ because $\card(F\ell C\times F\ell
C)=\Pi_\circ^2=\Pi_\circ$, we find that the left-hand side of
\hyperref[I.eq7.6.1]{(7.6.1)} is bounded above by $\Pi^{\Pi_\circ}\times
\Pi_\circ^{\Pi_\circ}=\Pi^{\Pi_\circ}$, which completes the
proof of \hyperref[I.prop7.6]{7.6}.

\begin{enonce*}{Proposition 7.7}\etiquette[7.7]{I.prop7.7}
Let $E$ be a category in which fiber products are representable,
$(X_\alpha)_{\alpha\in A}$ a generating family of objects of $E$,
and $(\varphi_i)_{i\in I}$ a family of functors
$\varphi_i:E\rightarrow E_i$ commuting with fiber products. The family
$(\varphi_i)$ is conservative (\hyperref[I.prop5.1]{5.1}) if and only if,
for every $\alpha \in A$ and every subobject $X'$ of
$X_\alpha$ distinct from $X_\alpha$, there exists an $i\in I$ such that
$\varphi_i(X')\rightarrow \varphi_i(X_\alpha)$ is not an
isomorphism. In this case, if $E$ is a $𝒰$-category and
if $A$ is $𝒰$-small, there exists a
$𝒰$-small subset $J$ of $I$ such that $(\varphi_j)_{j\in J}$ is
already a conservative family of functors.
\end{enonce*}

The necessity of the stated
conservativity condition being
obvious, let us prove its sufficiency. By \hyperref[I.def6.1]{6.1} (iv),
it suffices to
prove that $(\varphi_i)$ is conservative for monomorphisms.
Let $u:Y'\rightarrow Y$\pageoriginale be a monomorphism in $E$ which
is not an isomorphism; we must prove that there exists $i\in I$ such
that $\varphi_i(u)$ is not an isomorphism. By the hypothesis on
$(X_\alpha)$, there exist an $\alpha$ and a morphism
$X_\alpha\rightarrow Y$ which does not factor through $Y'$, in
other words, such that the inverse image $X'$ of $Y'$ is a
subobject of $X_\alpha$ distinct from $X_\alpha$. By hypothesis, there
exists an $i\in I$ such that $\varphi_i(X')\rightarrow \varphi(X)$ is not
an isomorphism. Since $\varphi_i(X')\simeq
\fprod{\varphi(X_\alpha)}{\varphi_i(Y')}{\varphi_{i(Y)}}$, it
indeed follows that $\varphi_i(Y')\rightarrow \varphi_i(Y)$ is not
an isomorphism.

The second assertion of \hyperref[I.prop7.7]{7.7} follows immediately from the
first, in view of \hyperref[I.prop7.4]{7.4}.

\begin{enonce*}{Corollary 7.7.1}\etiquette[7.7.1]{I.cor7.7.1}
Let $E$ be a $𝒰$-category in which fiber products are
representable, and which admits a $𝒰$-small generating family
of objects. Then every generating family
$(Y_i)_{i\in I}$ of $E$ contains a $𝒰$-small generating subfamily
$(Y_j)_{j\in J}(\card J\in 𝒰)$.
\end{enonce*}

Indeed, it suffices to apply \hyperref[I.prop7.7]{7.7} to a
$𝒰$-small generating family $(X_\alpha)_{\alpha\in A}$ of
$E$ and to the family of functors $\varphi_i$ $(i\in I)$ represented
by the $Y_i$.

\begin{enonce*}{Proposition 7.8}\etiquette[7.8]{I.prop7.8}
Let $E$ be a category, $C$ a full subcategory generating
by strict epimorphisms, $D$ a category, and $\Hhom'(E,D)$ the
full subcategory of $\Hhom(E,D)$ consisting of the functors which
commute with inductive limits of type $C_{/X}$, where $X$ is an
arbitrary object of $E$. Then the functor $F\rightarrow F| C$
induces a fully faithful functor
$$
\Hhom'(E,D)\longrightarrow \Hhom(C,D){\quoi}.
$$
Consequently,\pageoriginale if $\Hhom(C,D)$ is a
$𝒰$-category $(\hyperref[I.def1.1]{1.1}.)$, for example {\rm(\hyperref[I.sec1.1.1]{1.1.1}. b))} if $C$
is $𝒰$-small and $D$ is a $𝒰$-category, then
$\Hhom'(E,D)$ is also a $𝒰$-category.
\end{enonce*}

Let $F$, $G:E\rightarrow D$ be two functors in the first
member, and let $u:F\rightarrow G$ be a natural transformation. If $X=\varinjlim
X_i$ in $E$ and if $F$ and $G$ commute with the inductive limit
under consideration, then $u(X):F(X)\rightarrow G(X)$ is identified with the inductive limit
of the morphisms $u(X_i):F(X_i)\rightarrow G(X_i)$, and is therefore determined
once the $u(X_i)$ are known. This shows that the functor considered
in \hyperref[I.prop7.8]{7.8} is faithful, in view of the implication a)
$\Rightarrow$ b) in \hyperref[I.prop7.2]{7.2} (i). Conversely, let $v:F|
\sisi{D}{C}\rightarrow G|\sisi{D}{C}$ be a natural transformation; let us prove that it comes from a
natural transformation $u:F\rightarrow G$. For every $X\in \ob
E$, define
$$
u(X):F(X)=\varinjlim_{C_{/X}}F(X_i)\longrightarrow
G(X)=\varinjlim_{C_{/X}}G(X_i)
$$
as the inductive limit of the $v(X_i)$. It is immediate that one
thus obtains maps functorial in $X$, hence a
natural transformation $u:F\rightarrow G$, \sisi{in}{and} finally that the morphism induced by $u$ from
$F|C\sisi{'}{}$ to $G|C\sisi{'}{}$ is $v$, which completes the proof.


\setcounter{subsection}{8}
\subsection{Cogenerating families and subcategories}\etiquette[7.9]{I.sec7.9}

Let $E$ be a category and $C$ a full subcategory of $E$. We
say that $C$ is cogenerating by strict monomorphisms (\resp
cogenerating by monomorphisms, \resp cogenerating, \resp
cogenerating for epimorphisms, \resp cogenerating for
strict epimorphisms) if the full subcategory $C^\circ$ of
$E^\circ$ is generating by strict epimorphisms (\resp\ldotsn).
The terminology for families is analogous. All the results of this
section concerning\pageoriginale the notion of a generating subcategory
and its variants (\hyperref[I.def7.1]{7.1}) therefore yield
corresponding results
for the dual notions, which we leave to the reader
to formulate for their personal satisfaction.

\begin{enonce*}{Proposition 7.10}\etiquette[7.10]{I.prop7.10}
Let $E$ be a category, $C$ a full
generating subcategory
$\hyperref[I.def7.1]{(7.1)}$, and $D$ a full subcategory of $E$. The subcategory $D$ is
cogenerating $\hyperref[I.sec7.9]{(7.9)}$ if and only if, for every pair of parallel
arrows $\xymatrix{ T\ar@<2pt>[r]^{u,v}\ar@<-2pt>[r]& X }$ in $E$
with source $T\in \ob C$ and $u\neq v$, there exists an arrow
$w:X\rightarrow I$ with target $I\in \ob D$ such that $wu\neq wv$.
\end{enonce*}

By definition, to say that $D$ is cogenerating means that for
every pair of parallel arrows $\xymatrix{ Y\ar@<2pt>[r]^{u,v}\ar@<-2pt>[r]& X
}$ in $E$ such that $u\neq v$, there exists an arrow
$w:X\rightarrow I$, with $I\in\ob D$, such that $wu\neq wv$. Thus
\hyperref[I.prop7.10]{7.10} simply means that it suffices to test this
property when $Y\in \ob C$. Now, since $C$ is generating,
the hypothesis $u\neq v$ implies that there exists $f:T\rightarrow Y$ such
that $uf\neq vf$, whence, by hypothesis, there exists a
$w:X\rightarrow I$ with target $I\in\ob D$ such that $w(uf)\neq w(vf)$,
and hence $wu=wv$, \cqfd

\begin{enonce*}{Corollary 7.11}\etiquette[7.11]{I.cor7.11}
With the notation of \hyperref[I.prop7.10]{7.10}, suppose that the
objects $I$ of $D$ are injective objects of $E$, \ie such that
for every monomorphism $X\rightarrow Y$ in $E$, the corresponding map
$\Hom(Y,I)\rightarrow \Hom(X,I)$ is surjective.
Suppose moreover that every pair of parallel arrows $\xymatrix{
T\ar@<2pt>[r]^{u,v}\ar@<-2pt>[r]& X }$ in $E$ factors as an
epimorphic (\resp effective epimorphic) pair of parallel arrows
$\xymatrix{ Y\ar@<2pt>[r]^{u',v'}\ar@<-2pt>[r]& X }$ followed by a
monomorphism $i:X'\rightarrow X$. Then, in the criterion of
\hyperref[I.prop7.10]{7.10} for $D$ to be cogenerating, one may restrict
to the pairs $(u,v)$ which are epimorphic (\resp
effective epimorphic).
\end{enonce*}

We\pageoriginale conclude:

\begin{enonce*}{Corollary 7.12}\etiquette[7.12]{I.cor7.12}
Let $E$ be a $𝒰$-category admitting a small
generating subcategory $C$, and such that every object of $E$ is the
source of a monomorphism into an injective object of $E$. Suppose
moreover that, for every $T\in\ob C$, the coproduct $T\amalg T$ in $E$
is representable, and that every morphism with source $T\amalg T$
factors as an effective epimorphism followed by a monomorphism.
Then $E$ admits a small full cogenerating subcategory $D$.
More precisely, one may take $D$ such that $\card \ob D\leq
\prod_{T,T'\in \ob C}2^{\card(\Hom(T',T\coprod T))^2}$.
\end{enonce*}

Indeed, by \hyperref[I.cor7.11]{7.11}, it suffices, for every $T\in \ob C$ and every
effective quotient $X$ of $T\amalg T$, to choose an embedding
of $X$ into an injective object $I$ of $E$, and to take for $D$ the
full subcategory of $E$ spanned by these $I$. The conclusion
then follows from \hyperref[I.cor7.5]{7.5}.

\begin{enonce*}[remark]{Examples 7.13}\etiquette[7.13]{I.exem7.13}
To construct small cogenerating subcategories in terms
of small generating subcategories, we are thus led to
seek conditions under which a $𝒰$-category $E$
has ``enough injectives'', \ie every object
embeds into an injective object (by a monomorphism). It is well
known [Tohoku] that this condition holds in a
abelian $𝒰$-category with exact (small) filtered inductive
limits (axiom $AB$ 5 of \loccit) which admits a small
generating family. Moreover, the construction of \loccit is not
essentially tied to abelian categories, and also works
in the category of sheaves of $𝒰$-sets
on a topological space $X\in 𝒰$. We shall not state
here the exactness properties which make the construction in
question work, and shall merely point out that, in\pageoriginale the
special case of the category of sheaves of sets on $X$, one
immediately reduces to the case of \loccit as follows. Note
that if $\mathcal{O}_X$ is a ring on $X$, then every injective
object $I$ of the category of $\mathcal{O}_X$-modules is also
injective as an object of the category of sheaves of sets.
Indeed, if $F$ is a %%
sheaf of sets, and $\mathcal{O}_X[F]$ the ``free $\mathcal{O}_X$-module generated by $F$'', there is by
definition a homomorphism of sheaves of sets
\begin{equation*}
F\longrightarrow \mathcal{O}_X[F]\tag{$\ast$}
\end{equation*}
giving rise to an isomorphism, functorial in $F$,
$$
\Hom_{\mathcal{O}_X}(\mathcal{O}_X[F], I)\xrightarrow{\sim} \Hom
(F,I){\quoi}.
$$
Since the functor $F\mapsto \mathcal{O}_X[F]$ manifestly takes
monomorphisms to monomorphisms, it follows immediately
that if $I$ is an injective $\mathcal{O}_X$-module, it is also an
injective sheaf of sets. It follows that if the morphism
$(\ast)$ is a monomorphism (as is the case if one takes for
$\mathcal{O}_X$ a constant ring with value a ring $A\neq 0$),
then an embedding of $\mathcal{O}_X[F^\ast]$ into an
injective $\mathcal{O}_X$-module $I$ gives an embedding of $F$ into
the injective sheaf of sets underlying $I$.
\end{enonce*}

The same argument applies, without change, to the category
of sheaves of $𝒰$-sets on a $𝒰$-site, which
will be introduced in the following exposé.

The chief interest of the existence of a small
cogenerating subcategory
lies in the representability criterion
\hyperref[I.cor8.12.7]{8.12.7}
below.

\setcounter{section}{7}

% END INCLUDED SGA4_English_current_buildable_source_20260906/source/01/drafts/SGA4_I_LINES_2354_2668_English.texfrag


% BEGIN INCLUDED SGA4_English_current_buildable_source_20260906/source/01/drafts/SGA4_I_LINES_2669_2936_English.texfrag
\section{Ind-objects and pro-objects}\etiquette[8]{I.sec8}\pageoriginale

\subsection{Cofinal functors and cofinal subcategories}\etiquette[8.1]{I.sec8.1}

\begin{enonce*}{Definition 8.1.1}\etiquette[8.1.1]{I.def8.1.1}
Let
\hypertarget{sga4.exp.I.equation.8.1.1.1}{}$$
\displaylines{\rlap{$(8.1.1.1)$}\hfill \varphi:I\longrightarrow
I'\hfill}
$$
be a functor. The functor $\varphi$ is said to be cofinal if, for every
category $C$ and every functor $u:I'\rightarrow C$, regarding
$\varinjlim u$ and $\varinjlim u\circ\varphi$ as objects of
$\Hhom(G(𝒱\text{-}\Ens))^\circ$ (\hyperref[I.def2.1]{2.1},
\hyperref[I.sec2.3.1]{2.3.1}) (where $𝒱$ is a universe such that
$I$, $I'\in 𝒱$ and $C$ is a $𝒱$-category), the canonical morphism
\hypertarget{sga4.exp.I.equation.8.1.1.2}{}$$
\displaylines{\rlap{$(8.1.1.2)$}\hfill \varinjlim
  u\circ\varphi\longrightarrow \varinjlim u\hfill}
$$
is an isomorphism. When $\varphi$ is the inclusion functor of a
subcategory of $I'$, the subcategory $I$ is said to be cofinal if
$\varphi$ is cofinal.
\end{enonce*}

\setcounter{subsubsection}{1}
\subsubsection{}\etiquette[8.1.2]{I.sec8.1.2}
Notice that for given $\varphi$ and $u$, the bijectivity of \hyperlink{sga4.exp.I.equation.8.1.1.2}{(8.1.1.2)}
does not depend on the choice of the universe $𝒱$.

Returning to the definition of the terms occurring in \hyperlink{sga4.exp.I.equation.8.1.1.2}{(8.1.1.2)}
(\hyperref[I.def2.1]{2.1}, \hyperref[I.sec2.3.1]{2.3.1}), we see that to
say that $\varphi$ is cofinal also means that, for every universe $𝒱$
such that $I$ and $I'$ are $𝒱$-small, and every functor
$$
v:{I'}^\circ\longrightarrow (𝒱\text{-}\Ens){\quoi},
$$
the canonical homomorphism
\hypertarget{sga4.exp.I.equation.8.1.2.1}{}$$
\displaylines{\rlap{$(8.1.2.1)$}\hfill \varprojlim v\longrightarrow
  \varprojlim v\circ\varphi\hfill}
$$
is an isomorphism, \ie a bijection. To see that this condition is indeed
necessary, observe, on putting $v=v^\circ$, that it also says that the
condition of \hyperref[I.def8.1.1]{Definition 8.1.1} is satisfied when
$C=(𝒱\text{-}\Ens)^\circ$ (which\pageoriginale implies that the
inductive limits considered in \hyperlink{sga4.exp.I.equation.8.1.1.2}{(8.1.1.2)} are representable in $C$).

\setcounter{paragraph}{1}
\paragraph{}\etiquette[8.1.2.2]{I.sec8.1.2.2}
It follows immediately from the definition that \emph{the composite of
two cofinal functors is cofinal.}

\begin{enonce*}{Proposition 8.1.3}\etiquette[8.1.3]{I.prop8.1.3}
Let $\varphi:I\rightarrow I'$ be a functor.
\begin{enumerate}
\renewcommand{\theenumi}{\alph{enumi}}
\renewcommand{\labelenumi}{\theenumi)}
\item For $\varphi$ to be cofinal, it is necessary that $\varphi$
  satisfy the condition:

\begin{itemize}
\item[F 1)] For every object $i'$ of $I'$, there exists an object $i$ of
  $I$ such that $\varphi(i)$ is an upper bound of $i'$ (\ie such that
  $\Hom(i',\varphi(i))\neq \phi$).
\end{itemize}

\item Suppose that $I$ is filtered. For $\varphi$ to be cofinal, it is
  necessary and sufficient that $\varphi$ satisfy condition {\rm (F 1)}
  and the following condition:
\begin{itemize}
\item[F 2)] For every object $i$ of $I$ and every double arrow
$\xymatrix{ i'\ar@<2pt>[r]^{f',g'}\ar@<-2pt>[r]& \varphi(i) }$ in $I'$
with target $\varphi(i)$, there exists an arrow $h:i\rightarrow j$ in
$I$ such that $\varphi(h)f'=\varphi(h)g'$.
\end{itemize}
Moreover, if $\varphi$ is cofinal, $I'$ is filtered.

\item Suppose that $I'$ is filtered and $\varphi$ is fully faithful. For
  $\varphi$ to be cofinal, it is necessary and sufficient that it satisfy
  condition {\rm F 1)} of {\rm a)}. This implies that $I$ is filtered.
\end{enumerate}
\end{enonce*}

\noindent \emph{Proof.} For necessity in a) and b), use Definition
\hyperref[I.def8.1.1]{8.1.1} only in the case where $u$ is the canonical
inclusion functor (\hyperref[I.sec1.3.3]{1.3.3})
$$
u:I'\hookrightarrow \widehat{I}'
$$
(having chosen a universe $𝒰$ such that $I$ and $I'$ are
$𝒰$-small). We may then interpret \hyperlink{sga4.exp.I.equation.8.1.1.2}{(8.1.1.2)} as an arrow of
$\widehat{I}'$ (\hyperref[I.prop3.1]{3.1}) whose target is the
\emph{terminal} functor\pageoriginale on $I'$ (\hyperref[I.prop3.4]{3.4}).
It is immediately clear that condition F 1) says that this arrow is an
epimorphism of $\widehat{I}'$, \ie that it is surjective at each
argument (taking into account that inductive limits in $\widehat{I}'$
are computed ``argument by argument''). This proves a). Now suppose that
$I$ is filtered and $\varphi$ is cofinal. Then $I'$ is filtered: indeed,
the condition that two objects of $I'$ be bounded above by a third follows
at once from the same condition on $I$ and from F 1), and it remains only
to prove condition PS 2) of \hyperref[I.def2.7]{2.7}, \ie that for every
double arrow
$$
f',g':i'\rightrightarrows j'
$$
in $I'$, there exists an arrow $h':j'\rightarrow k'$ of $I'$ such that
$h'f'=h'g'$. By F~1), we may suppose that $k'$ is of the form
$\varphi(i)$. But since $\varphi$ is cofinal, for every object $i'$ of
$I'$ the set $\varinjlim_i\Hom(i',\varphi(i))$ consists of a single
element. The standard description of \emph{filtered} inductive limits in
(Ens) (\hyperref[I.lem2.8.1]{2.8.1}) therefore implies that there exists
an arrow $h:i\rightarrow j$ in $I$ such that
$\varphi(h)f'=\varphi(h)g'$. This completes the proof that $I'$ is
filtered and at the same time proves condition F~2).

To prove that the conditions in b) are sufficient for $\varphi$ to be
cofinal, use the form \hyperref[I.sec8.1.2]{8.1.2} of the definition. One
sees at once that condition F 1) implies that \hyperlink{sga4.exp.I.equation.8.1.2.1}{(8.1.2.1)} is a
monomorphism, \ie is injective at each argument, while condition F 2)
(together with F 1)) ensures that it is bijective (taking into account
the argument-by-argument computation of projective limits in
$\widehat{{I'}^\circ}$). This proves b).

Finally,\pageoriginale if $I'$ is filtered and $\varphi$ is fully
faithful, it is immediate that condition F 1) implies that $I$ is
filtered and implies condition F 2). Thus c) follows from a) and b).

\setcounter{subsubsection}{3}
\subsubsection{}\etiquette[8.1.4]{I.sec8.1.4}
When $I'$ is a filtered category and $I$ is a full subcategory of $I'$,
we see from \hyperref[I.prop8.1.3]{8.1.3} c) that the condition that $I$ be cofinal in $I'$
depends only on the subset $\Ob I$ of the preordered set $OB\,I'$ (for
the preorder relation ``$x\leq y$'' $\Leftrightarrow$
``$\Hom(x,y)\neq ∅$''). If $J'$ is a preordered set and $J$ a subset of
$J'$, we shall sometimes say that $J$ is a \emph{cofinal subset of} $J'$
when every element of $J'$ is bounded above by an element of $J$. When
$J'$ is filtered, this therefore means that the inclusion functor
$𝒥\hookrightarrow 𝒥'$ for the associated categories is cofinal.

\subsubsection{}\etiquette[8.1.5]{I.sec8.1.5}
In what follows, we shall use the notion of cofinal functor only when the
categories $I$ and $I'$ are filtered. Classically, one even restricted
oneself to categories associated with preordered sets (\ie categories in
which there is at most one arrow with given source and target). It turns
out, however, that this restriction is inconvenient in applications,
since the ``natural'' filtered categories arising in many applications
are \emph{not} ordered categories. The following result, %%
due to P. \textsc{Deligne}, shows nevertheless that there is no essential
difference between the two points of view:

\begin{enonce*}{Proposition 8.1.6}\etiquette[8.1.6]{I.prop8.1.6}
Let\pageoriginale $I$ be a small filtered category. Then there exist a
small ordered set $E$ and a cofinal functor $\varphi:ℰ\rightarrow I$,
where $ℰ$ denotes the category associated with $E$.
\end{enonce*}

First suppose that the preordered set $\Ob I$ has no greatest element.
A subdiagram of $I$ is a pair $D=(\mathcal{O},F)$ consisting of a subset
$F$ of $F\ell(I)$ and a subset $\mathcal{O}$ of $\Ob I$, such that for
every $f\in F$, the source and target of $f$ belong to $\mathcal{O}$. An
element $e$ of $\mathcal{O}$ is called a \emph{terminal object} of the
subdiagram $D$ if, for every $x\in D'$, the set
$\Hom(x,e)\cap F$ of arrows of $D$ with source $x$ and target $y$ has
exactly one element $f_x$, if $f_x=f_y\circ g$ for every arrow
$g:x\rightarrow y$ of $D$, and if $f_e=\id_e$. Let $E$ be the set of
\emph{finite} subdiagrams $D$ of $I$ having a \emph{unique} terminal
element $\varphi(D)$, and order $E$ by inclusion. If $D$, $D'$ are two
elements of $E$ and $D\subset D'$, there is a unique arrow
$\varphi(D',D)$ of $D'$ with source $\varphi(D)$ and target
$\varphi(D')$. If $D\subset D'\subset D''$ are inclusions in $E$, one
clearly has $\varphi(D'',D)=\varphi(D'',D')\varphi(D',D)$; one also has
$\varphi(D,D)=\id_{\varphi(D)}$. Thus we obtain a functor
$\varphi:ℰ\rightarrow I$, and it remains to prove that $E$ is filtered
and $\varphi$ is cofinal. In view of \hyperref[I.prop8.1.3]{8.1.3} b),
three verifications are required:
\begin{enumerate}
\renewcommand{\labelenumi}{\theenumi)}
\item Condition F 1) of \hyperref[I.prop8.1.3]{8.1.3}: for every $i\in I$, take for $D$ the
  subdiagram of $I$ consisting only of the object $i$ and its identity
  arrow; then $i\leq \varphi(D)=i$. This condition F 1) also shows that
  $E\neq ∅$.

\item Condition F 2): For every $i\in \Ob I$, every $D\in E$, and every
  double arrow $i\rightrightarrows \varphi(D)$, find a diagram $D'\in E$
  containing $D$ such that
  $\varphi(D',D):\varphi(D)\rightarrow \varphi(D')$ \emph{equalizes}
  the double arrow. Since $I$ is filtered, there\pageoriginale is an
object $j$ of $I$ and an arrow $f:\varphi(D)\rightarrow j$ that equalizes
the double arrow. Since $I$ has no greatest element, we may suppose that
$j$ is a \emph{strict} upper bound of $\varphi(D)$, which implies that it
is distinct from every other object of $D$. Let $D'$ be the subdiagram
of $I$ whose set of objects is $\mathcal{O}\cup\{j\}$ and whose set of
arrows consists of the arrows of $D$, together with the composites
$x\xrightarrow{f_x}\varphi(D)\xrightarrow{f}j$, where $x$ is an object
of $D$ and $f_x$ is the unique arrow of $D$ with source $x$ and target
$\varphi(D)$, and $\id_j$. It is then clear that $D'$ is a finite
subdiagram of $I$, that it has $j$ as its unique terminal object, hence
that $D'\in E$, and that $D'$ satisfies the required condition.

\renewcommand{\labelenumi}{$3^\circ$}
\item Two subdiagrams $D$, $D'\subset E$ are contained in the same
  $D''\in E$. Indeed, one may find a strict upper bound $j$ of
  $\varphi(D)$ and $\varphi(D')$,
\[
\xymatrix@C=1.5cm@R=1cm{
\varphi(D)\ar[rd]^f && \\
           &j &\\
\varphi(D')\ar[ru]_{f'}&\quad , }
\]
and take for $D''$ the subdiagram whose set of objects is the union of
the sets of objects of $D$ and $D'$ with $\{j\}$, and whose set of arrows
is the union of the sets of arrows of $D$ and $D'$, the set of composites
$f\circ f_x$ ($x$ an object of $D$) and $f'\circ f'_{x'}$ ($x'$ an
object of $D'$), and $\{\id_j\}$. This indeed gives a finite subdiagram
of $E$. Let us show that, after replacing $j$, $f$, $f'$ by $j'$, $gf$,
$gf'$ for a suitable $g:j\rightarrow j'$, we may arrange that $j$ is a
terminal object of $D'$. Equivalently, for every $x$ that is an object
of both $D$ and $D'$, one must have $ff_x=f'f'_{x'}$. This amounts to a
\emph{finite} set of double arrows\pageoriginale to be equalized by a
$g:j\rightarrow j'$, which is possible since $I$ is filtered.
\end{enumerate}

This completes the proof in the case considered. The general case is
reduced to it by introducing the filtered category $\NN$ associated with
the ordered set $\NN$ of natural integers, and observing that the category
$\NN\times I$ is filtered, has no greatest element, and that the
projection $\NN\times I\rightarrow I$ is a cofinal functor.

\begin{enonce*}{Corollary 8.1.7}\etiquette[8.1.7]{I.cor8.1.7}
Let $I$ be a $𝒰$-category. There exist a small filtered ordered set $E$
and a cofinal functor $ℰ\rightarrow I$ if and only if $I$ is filtered
and $\Ob I$ has a small cofinal subset $I'$ (\hyperref[I.sec8.1.4]{8.1.4}).
\end{enonce*}

Necessity follows from \hyperref[I.prop8.1.3]{8.1.3} a), and sufficiency follows from \hyperref[I.prop8.1.6]{8.1.6}
applied to the full subcategory of $I$ defined by $I'$.

\begin{enonce*}{Definition 8.1.8}\etiquette[8.1.8]{I.def8.1.8}
Let $I$ be a filtered category. We say that $I$ is essentially small if
$I$ is a $𝒰$-category and satisfies the equivalent conditions of
\hyperref[I.cor8.1.7]{8.1.7}.
\end{enonce*}

\setcounter{subsection}{1}

% END INCLUDED SGA4_English_current_buildable_source_20260906/source/01/drafts/SGA4_I_LINES_2669_2936_English.texfrag


% BEGIN INCLUDED SGA4_English_current_buildable_source_20260906/source/01/drafts/SGA4_I_LINES_2937_3118_English.texfrag
\subsection{Ind-objects and ind-representable functors}\etiquette[8.2]{I.sec8.2}

\subsubsection{}\etiquette[8.2.1]{I.sec8.2.1}
In what follows, we fix a $𝒰$-category $C$, which we always regard as
embedded in the category $\widehat{C}$ of $𝒰$-presheaves by means of
the canonical functor \hyperref[I.sec1.3.1]{(1.3.1)}
\hypertarget{sga4.exp.I.equation.8.2.1.1}{}$$
\displaylines{\rlap{$(8.2.1.1)$}\hfill h:C\hookrightarrow
  \widehat{C}{\quoi}.\hfill}
$$
An \emph{ind-object of} $C$ (or an \emph{inductive system} of $C$) is
any functor
\hypertarget{sga4.exp.I.equation.8.2.1.2}{}$$
\displaylines{\rlap{$(8.2.1.2)$}\hfill \varphi:I\longrightarrow
  C{\quoi},\hfill}
$$
where\pageoriginale $I$ is a \emph{filtered} category
(\hyperref[I.def2.7]{2.7}), called the \emph{index category} of the
ind-object. (Apart from this, $I$ is arbitrary; in particular, we do not
require $I$ to be a $𝒰$-category.) It will often be convenient to denote
an ind-object $\varphi$ by the indexed notation $(X_i)_{i\in \ob I}$, or
even (by another abuse of notation) $(X_i)_{i\in I}$, or $(X_i)_i$, or
$(X_i)$, where $X_i=\varphi(i)$ for $i\in \ob I$. This notation is no
more and no less abusive than the notation classically used for inductive
systems indexed by filtered ordered sets (where the transition morphisms
are likewise not mentioned). The $X_i$ are called the \emph{component
objects} of the ind-object $𝒳$. Note that, in the general case considered
here, the ``\emph{transition morphisms}'' $\varphi(f)$ of the inductive
system are indexed by the set $F\ell I$ of arrows of $I$, which in
general is not identified with a subset of $\Ob I\times\Ob I$; in other
words, for given $i$ and $j$, with $j$ an upper bound of $i$, there may
exist more than one transition morphism from $X_i$ to $X_j$.

\subsubsection{}\etiquette[8.2.2]{I.sec8.2.2}
The most useful ind-objects $𝒳=(X_i)_{i\in I}$ of $C$ are those for
which the index category $I$ is essentially small \hyperref[I.def8.1.8]{(8.1.8)}; such an
ind-object is called \emph{essentially small}. If $(X_i)_{i\in I}$ is
such an object, then, using the fact that small inductive limits are
representable in $\widehat{C}$ \hyperref[I.prop3.1]{(3.1)}, we may consider
\hypertarget{sga4.exp.I.equation.8.2.2.1}{}$$
\displaylines{\rlap{$(8.2.2.1)$}\hfill
L(𝒳)\mathrel{\mathop=^{\rm def}}
  \varinjlim h\cdot 𝒳 = \varinjlim_i X_i \in \ob
  \widehat{C}{\quoi},\hfill}
$$
where the last limit is taken \emph{in} $\widehat{C}$. Thus $L(𝒳)$ is
the presheaf
\hypertarget{sga4.exp.I.equation.8.2.2.2}{}$$
\displaylines{\rlap{$(8.2.2.2)$}\hfill L(𝒳):Y\longmapsto
  \varinjlim_i \Hom(Y,X_i)\hfill}
$$
on\pageoriginale $C$. This presheaf is called the \emph{presheaf
ind-represented by the ind-object $𝒳=(X_i)_{i\in I}$ of $C$.} A presheaf
$F$ on $C$ is called an \emph{ind-representable presheaf} if it is
isomorphic to a presheaf $L(𝒳)$ ind-represented by an essentially small
ind-object. It follows moreover from \hyperref[I.prop8.1.6]{8.1.6} that,
in this definition, we may suppose that $I$ is a small category, and even
that $I$ is associated with an ordered set $\in 𝒰$.

\subsubsection{}\etiquette[8.2.3]{I.sec8.2.3}
Consider an ind-object $𝒳=(X_i)_{i\in I}$ given by a functor
$\varphi:I\rightarrow C$, and let
$$
u:I'\longrightarrow I
$$
be a functor, where $I'$ is a second filtered category. Then the functor
$\varphi\circ u$ is an ind-object $𝒳'$ of $C$ with index category $I'$;
in indexed notation it is $(X_{u(i')})_{i'\in I'}$. It is called the
ind-object of $C$ \emph{deduced from $(X_i)_{i\in I}$ by a change of
index categories} using the functor $u$. If $I$ and $I'$ (\ie $𝒳$ and
$𝒳'$) are essentially small, one obtains a canonical morphism
\hypertarget{sga4.exp.I.equation.8.2.3.1}{}$$
\displaylines{\rlap{$(8.2.3.1)$}\hfill
  L(𝒳')=\varinjlim_{i'}X_{u(i')}\longrightarrow L(𝒳)=\varinjlim_i
  X_i\hfill}
$$
between the presheaves ind-represented by $𝒳'$ and $𝒳$, respectively.
When $u$ is a cofinal functor (\hyperref[I.def8.1.1]{8.1.1}), $𝒳$ is
essentially small if and only if $𝒳'$ is \hyperref[I.cor8.1.7]{(8.1.7)}, and the preceding
homomorphism \hyperlink{sga4.exp.I.equation.8.2.3.1}{(8.2.3.1)} is an isomorphism
\hypertarget{sga4.exp.I.equation.8.2.3.2}{}$$
\displaylines{\rlap{$(8.2.3.2)$}\hfill L(𝒳')\simeq
  L(𝒳){\quoi}.\hfill}
$$

\subsubsection{}\etiquette[8.2.4]{I.sec8.2.4}
Let\pageoriginale
\hypertarget{sga4.exp.I.equation.8.2.4.1}{}$$
\displaylines{\rlap{$(8.2.4.1)$}\hfill 𝒳=(X_i)_{i\in
    I},𝒴=(Y_j)_{j\in I}\hfill}
$$
be two essentially small ind-objects of $C$, indexed by the (filtered and
essentially small) categories $I$ and $J$. A \emph{morphism} from the
ind-object $𝒳$ to the ind-object $𝒴$ is any morphism
\hypertarget{sga4.exp.I.equation.8.2.4.2}{}$$
\displaylines{\rlap{$(8.2.4.2)$}\hfill L(𝒳)\longrightarrow
  L(𝒴)\hfill}
$$
between the presheaves that they ind-represent. We put
\hypertarget{sga4.exp.I.equation.8.2.4.3}{}$$
\displaylines{\rlap{$(8.2.4.3)$}\hfill
\Hom_{\text{ind}-\ob}(𝒳,𝒴)=\Hom_{\widehat{C}}(L(𝒳),L(𝒴)){\quoi}.
   \hfill}
$$
(The subscript ind-ob on $\Hom$ is omitted when there is no risk of
confusion.) Composition of morphisms of ind-objects of $𝒞$ is defined as
composition of the morphisms of the presheaves that they ind-represent.
If $𝒱\supset 𝒰$ is a universe containing $𝒰$, and if we restrict to
the essentially small ind-objects that belong to $𝒱$, these form the set
of objects of a category whose set of arrows consists of triples
$(𝒳,𝒴,u)$, where $𝒳$ and $𝒴$ are essentially small ind-objects
$\in 𝒱$ and $u:𝒳\rightarrow 𝒴$ is a morphism of ind-objects, \ie a
morphism $L(𝒳)\rightarrow L(𝒴)$. The category thus obtained is denoted by
\hypertarget{sga4.exp.I.equation.8.2.4.4}{}$$
\displaylines{\rlap{$(8.2.4.4)$}\hfill
  \Ind_{𝒱}(C,𝒰){\quoi}.\hfill}
$$
When $𝒱=𝒰$, it is denoted simply by
\hypertarget{sga4.exp.I.equation.8.2.4.5}{}$$
\displaylines{\rlap{$(8.2.4.5)$}\hfill \Ind_{𝒰}(C)\text{ or
}
  \Ind(C){\quoi},\hfill}
$$
and is called the \emph{category of ind-objects of} $C$ relative to
$𝒰$, with the mention of $𝒰$ omitted when no confusion is likely.

Clearly,\pageoriginale if $𝒱\subset 𝒱'$ are two universes containing
$𝒰$, then $\Ind_{𝒱}(C,𝒰)$ is a full subcategory of
$\Ind_{𝒱'}(,𝒰)$; it follows from \hyperref[I.sec8.2.3]{8.2.3} that the inclusion functor is an
\emph{equivalence of categories}:
\hypertarget{sga4.exp.I.equation.8.2.4.6}{}$$
\displaylines{\rlap{$(8.2.4.6)$}\hfill
\Ind_{𝒱}(C,𝒰)\xrightarrow{\approx}\Ind_{𝒱'}(C,𝒰){\quoi}.\hfill}
$$
This shows in particular that \emph{every essentially small ind-object of
$C$ defines an element of $\Ind(C)$, up to unique isomorphism.} This
observation justifies the fairly common abuse of language of identifying
every essentially small ind-object of $C$ with an object of $\Ind(C)$.

It follows at once from the definitions that, for every universe $𝒱$,
there is a canonical functor
\hypertarget{sga4.exp.I.equation.8.2.4.7}{}$$
\displaylines{\rlap{$(8.2.4.7)$}\hfill
  L:\Ind_{𝒱}(C,𝒰)\longrightarrow \widehat{C}\hfill}
$$
which is fully faithful. By \hyperlink{sga4.exp.I.equation.8.2.4.6}{(8.2.4.6)}, these functors are clearly
determined up to unique isomorphism once one of them is known, and in
particular once the canonical functor
\hypertarget{sga4.exp.I.equation.8.2.4.8}{}$$
\displaylines{\rlap{$(8.2.4.8)$}\hfill L:\Ind(C)\longrightarrow
  \widehat{C}{\quoi}.\hfill}
$$
is known. Since this functor is fully faithful, it is often used to
identify an object of the first member with the functor that it
ind-represents, or even to identify $\Ind(C)$ with its essential image in
$\widehat{C}$, consisting of the ind-representable presheaves. Note,
however, that this identification has distinctly more disadvantages than
the analogous identification of $C$ with a full subcategory of
$\widehat{C}$ by the functor $h$ \hyperlink{sga4.exp.I.equation.8.2.1.1}{(8.2.1.1)}, since, unlike the latter, the
functor \hyperlink{sga4.exp.I.equation.8.2.4.8}{(8.2.4.8)} is not in general injective on objects.


% END INCLUDED SGA4_English_current_buildable_source_20260906/source/01/drafts/SGA4_I_LINES_2937_3118_English.texfrag


% BEGIN INCLUDED SGA4_English_current_buildable_source_20260906/source/01/drafts/SGA4_I_LINES_3119_3227_English.texfrag
\subsubsection{}\etiquette[8.2.5]{I.sec8.2.5}
Let us express\pageoriginale Definition \hyperlink{sga4.exp.I.equation.8.2.4.3}{(8.2.4.3)} explicitly in terms
of the indexed expressions \hyperlink{sga4.exp.I.equation.8.2.4.1}{(8.2.4.1)} of the ind-objects under
consideration. Taking Definition \hyperlink{sga4.exp.I.equation.8.2.2.1}{(8.2.2.1)} into account, one obtains
$$
\Hom(𝒳,𝒴)\simeq \varprojlim_i
\Hom(X_i,𝒴)\mathrel{\mathop=^{\rm dfn}} \varprojlim_i
\Hom(X_i, L(𝒴))=\varprojlim_i\varinjlim_j \Hom_C(X_i,
Y_j){\quoi},
$$
where the last equality follows from \hyperlink{sga4.exp.I.equation.8.2.2.2}{(8.2.2.2)} applied to $𝒴$.
Thus one obtains a canonical bijection
\hypertarget{sga4.exp.I.equation.8.2.5.1}{}$$
\displaylines{\rlap{$(8.2.5.1)$}\hfill \Hom((X_i)_{i\in
I},(Y_j)_{j\in
    J})\simeq \varprojlim_{i\in I} \varinjlim_{j\in
    J}\Hom_C(X_i,Y_j){\quoi}.\hfill}
$$
We leave to the reader the task of expressing composition of morphisms of
ind-objects explicitly in terms of this formula. Notice that it follows
at once from this formula that the set $\Hom(𝒳,𝒴)$ of homomorphisms of
ind-objects is $𝒰$-\emph{small}, \ie that the categories \hyperlink{sga4.exp.I.equation.8.2.4.4}{(8.2.4.4)} are
$𝒰$-categories. Equivalently (by \hyperlink{sga4.exp.I.equation.8.2.4.6}{(8.2.4.6)}), the \emph{category
$\Ind(C)$ is a $𝒰$-category}, \ie the second member of \hyperlink{sga4.exp.I.equation.8.2.5.1}{(8.2.5.1)} is
small when $I$, $J\in 𝒰$. This follows at once from the fact that $C$
is a $𝒰$-category, so the sets $\Hom_C(X_i,Y_j)$ are small.

\subsubsection{}\etiquette[8.2.6]{I.sec8.2.6}
Let $I$ be an essentially small filtered category. It follows from
\hyperlink{sga4.exp.I.equation.8.2.5.1}{(8.2.5.1)}, or from the functoriality of the inductive limit of a functor
$I\rightarrow \widehat{C}$ with respect to that functor, that there is a
canonical functor
\hypertarget{sga4.exp.I.equation.8.2.6.1}{}$$
\displaylines{\rlap{$(8.2.6.1)$}\hfill \Hhom(I,C)\longrightarrow
  \Ind(C){\quoi}.\hfill}
$$
Note that this functor is in general \emph{not} fully faithful, nor even
faithful.

\begin{enonce*}[remark]{Remark 8.2.7}\etiquette[8.2.7]{I.sec8.2.7}
The\pageoriginale definitions of \hyperref[I.sec8.2.4]{8.2.4} make clear
the notion of \emph{isomorphism} of two (essentially small) ind-objects,
which also means the isomorphism of the presheaves that they
ind-represent. Note that this is a very weak notion of isomorphism when
compared with the notion of isomorphism in categories of the form
$\Hhom(I,C)$. Thus, many fairly natural conditions that may be imposed on
an ind-object (such as having its components in a given strictly full
subcategory, or being strict (\hyperref[I.sec8.12]{8.12} below))
\emph{are not} invariant under isomorphism. Hence, if one wishes to work
with notions invariant under isomorphism of ind-objects, the notions in
question must be ``saturated'' by passing to the strictly full
subcategory of $\Ind(C)$ generated by the subcategory of $\Ind(C)$ formed
by the objects satisfying the condition under consideration. See, for
example, the notion of essentially constant ind-object introduced below
(\hyperref[I.sec8.4]{8.4}).
\end{enonce*}

\begin{enonce*}[remark]{Exercise 8.2.8}\etiquette[8.2.8]{I.exer8.2.8}
Let $𝒰\subset 𝒱$ be two universes and let $C$ be a
$𝒰$-category belonging to $𝒱$. Denote by $\SysInd(C)$ the following
category:
\begin{enumerate}
\renewcommand{\theenumi}{\alph{enumi}}
\renewcommand{\labelenumi}{\theenumi)}
\item The objects of $\SysInd(C)$ are the functors $\varphi:I\rightarrow C$
where $I$ is an essentially $𝒰$-small filtered category belonging to
$𝒱$.

\item Let $(I,\varphi)$ and $(J,\psi)$ be two objects of $\SysInd(C)$.
A morphism in $\SysInd(C)$ from $(I,\varphi)$ to $(J,\psi)$ is a pair
$(m,u)$, where $m:I\rightarrow J$ is a functor and
$u:\varphi\rightarrow \psi\circ m$ is a morphism of functors. Composition
of morphisms in $\SysInd(C)$ is defined in the evident way.
\end{enumerate}
\end{enonce*}

Let\pageoriginale $S$ be the set of morphisms
$(m,u):(I,\varphi)\rightarrow (J,\Psi)$ of $\SysInd(C)$ such that $m$ is
cofinal and $u$ is an isomorphism. Let $(I,\varphi)$ be an object of
$\SysInd(C)$. The category $F\ell(I)$ of arrows of $I$ maps to $I$ by two
natural functors: source and target. These two functors are moreover
linked by the canonical morphism of functors $v:$ source $\rightarrow$
target. Hence there are two morphisms in $\SysInd(C)$ from
$(F\ell(I),\varphi\circ\text{ source})$ to $(I,\varphi)$:
$p_1(I,\varphi)=(\text{source, } \id)$ and
$p_2(I,\varphi)=(\mbox{target, } \varphi^\ast v:\varphi\circ
\text{source }\rightarrow \varphi\circ{\text{ target}})$. Let
$p:\SysInd(C)\rightarrow\Ind(C)$ be the evident functor. Let $B$ be a
category. Show that the functor $F\mapsto F\circ p$
$$
\Hhom(\Ind(C),B)\longrightarrow \Hhom(\SysInd(C),B)
$$
is fully faithful, and that a functor $G:\SysInd(C)\rightarrow B$ belongs
to its essential image if and only if it has the following two
properties:
\begin{enumerate}
\renewcommand{\labelenumi}{\theenumi)}
\item For every $s\in S$, $F(s)$ is an isomorphism of $B$.

\item For every object $(I,\varphi)$ of $\SysInd(C)$,
$F(p_1(I,\varphi))=F(p_2(I,\varphi))$.
\end{enumerate}


% END INCLUDED SGA4_English_current_buildable_source_20260906/source/01/drafts/SGA4_I_LINES_3119_3227_English.texfrag


% BEGIN INCLUDED SGA4_English_current_buildable_source_20260906/source/01/drafts/SGA4_I_LINES_3228_3408_English.texfrag
\subsection{Characterization of ind-representable functors}\etiquette[8.3]{I.sec8.3}

\begin{enonce*}{Proposition 8.3.1}\etiquette[8.3.1]{I.prop8.3.1}
Let $F$ be an ind-representable presheaf on $C$. Then $F$ is
left exact, \ie (\hyperref[I.cor2.3.2]{2.3.2}) for every finite category
$J$ and every
functor $\varphi:J\rightarrow C$ such that $\varinjlim \varphi$ is
representable (\ie $\varprojlim \varphi^\circ$ is representable),
the canonical map
$$
F(\varinjlim \varphi)\longrightarrow \varprojlim F\circ\varphi
$$
is bijective.
\end{enonce*}

Indeed,\pageoriginale since representable presheaves are
clearly left exact, it follows at once
from \hyperref[I.prop2.8]{2.8} that the same is true
of every filtered inductive limit of such functors, \cqfd

\setcounter{subsubsection}{1}
\subsubsection{}\etiquette[8.3.2]{I.sec8.3.2}
Given a presheaf $F$ on $C$, we shall often have occasion to work
with the category
\hypertarget{sga4.exp.I.equation.8.3.2.1}{}$$
\displaylines{\rlap{$(8.3.2.1)$}\hfill C_{/F}\hookrightarrow
\widehat{C}_{/F}{\quoi},\hfill}
$$
which denotes the full subcategory of the latter category formed by
arrows $X\rightarrow F$ \emph{whose source lies in} $C$. It follows from
\hyperref[I.prop1.4]{1.4} that the objects of this category may be identified with pairs
$(X,u)$, where $X$ is an object of $C$ and $u\in F(X)$. A morphism from
$(X,u)$ to $(X',u')$ is then interpreted as an arrow
$f:X\rightarrow X'$ in $C$ such that $F(f)(u')=u$. The
functor from $\widehat{C}_{/F}$ to $\widehat{C}$ that ``forgets the
target'' induces a functor (called \emph{canonical}).
\hypertarget{sga4.exp.I.equation.8.3.2.2}{}$$
\displaylines{\rlap{$(8.3.2.2)$}\hfill C_{/F}\longrightarrow
C{\quoi},\hfill}
$$
already considered in \hyperref[I.prop3.4]{3.4}, where it is proved that
the inductive limit of this
functor exists in $\widehat{C}$ (without any smallness condition on
$C_{/F}$!) and is canonically isomorphic to $F$.


\setcounter{paragraph}{2}
\paragraph{}\etiquette[8.3.2.3]{I.sec8.3.2.3}
Note that if finite inductive limits are
representable in $C$, and if $F$ is left exact (\ie takes them to
finite projective limits in $(\Ens)$), then the same is true
in $C_{/F}$, and a fortiori (\hyperref[I.exem2.7.1]{2.7.1}) $C_{/F}$ is filtered. If,
more generally, sums of two objects and
cokernels of pairs of parallel arrows are representable in $C$, and if $F$
takes them to products \resp kernels, then\pageoriginale $C_{/F}$
is closed under the same types of finite inductive limits, hence
it is filtered if and only if it is nonempty, \ie if and
only if the functor $F$ is not the constant functor with
value $∅$.

\begin{enonce*}{Theorem 8.3.3}\etiquette[8.3.3]{I.them8.3.3}
Let $C$ be a $𝒰$-category, and let $F\in \ob \widehat{C}$,
$F:C^\circ\rightarrow (𝒰\text{-}\Ens)$ be a
$𝒰$-presheaf on $C$. The following conditions are
equivalent:
\begin{itemize}
\item[(i)] $F$ is ind-representable (\hyperref[I.sec8.2.2]{8.2.2}).

\item[(ii)] The category $C_{/F}$ (\hyperref[I.sec8.3.2]{8.3.2}) is filtered and
essentially small.

\item[(iii)] (if finite inductive limits are
representable in $C$.) The functor $F$ is left exact, and $\ob C_{/F}$
has a small cofinal subset (\hyperref[I.sec8.1.4]{8.1.4}).

\item[(iii bis)] (If sums of two objects and
cokernels of pairs of parallel arrows are representable in $C$.) The functor $F$
takes sums of two objects of $C$ to products, and cokernels of
pairs of parallel arrows of $C$ to kernels; the category $C_{/F}$ is nonempty,
\ie the functor $F$ is not the constant functor with value
$∅$; finally, there exists a small family of objects of
$C_{/F}$ such that every object $X$ of $C_{/F}$ is dominated by an
object $X_i$, \ie admits an $F$-morphism $X_i\rightarrow X$.

\item[(iv)] (if the category $C$ is equivalent to a small
category.) The category $C_{/F}$ is filtered.

\item[(v)] (If the category $C$ is equivalent to a small
category, and if filtered inductive limits are representable in
it.) The functor $F$ is left exact.
\end{itemize}
\end{enonce*}

\noindent \emph{Proof.} (i)\pageoriginale $\Rightarrow$
(ii). Suppose that $F$ is ind-represented by $(X_i)_{i\in I}$, with $I$
small. Let us prove that $C_{/F}$ is filtered. Let $X$, $X'$ be two
objects of $C_{/F}$, \ie objects of $C$ equipped with morphisms
$X\rightarrow F$, $X'\rightarrow F$; let us prove that they have
a common upper bound in $C_{/F}$. The morphisms $X\rightarrow
F$, $X'\rightarrow F$ come from morphisms $X\rightarrow X_i$,
$X'\rightarrow X'_i$, and, after replacing $i$, $i'$ by a
common upper bound in $\Ob I$, we may assume that $i=i'$, and take
as a common upper bound of $X$, $X'$ the object $X_i$ equipped with the
canonical morphism $X_i\rightarrow F$. Now let $\xymatrix{
X\ar@<2pt>[r]^{f,g}\ar@<-2pt>[r]& X' \ar@<2pt>[r]^{h} & F }$
 be a pair of parallel arrows in $C_{/F}$; let us prove that it is
equalized by an arrow $X'\rightarrow X''$ in $C_{/F}$. Now
$h:X'\rightarrow F$ is given by a morphism $h_i:X'\rightarrow
X_i$, and the condition $hf=hg$ means that there exists
$\alpha:i\rightarrow j$ in $I$ such that
$\varphi(\alpha)(h_if)=\varphi(\alpha)(h_ig)$, \ie after
replacing $h_i$ by $\varphi(\alpha)h_i$, this morphism equalizes $f$ and $g$.
This proves that $C_{/F}$ is filtered. It is then immediate
that it is essentially small, since the objects
$(X_i\rightarrow F)_{\in \Ob I}$ form a small family in
$\Ob C_{/F}$ that is cofinal.

(ii) $\Rightarrow$ (i). Put $I=C_{/F}$, and consider the
canonical functor \hyperlink{sga4.exp.I.equation.8.3.2.1}{(8.3.2.1)}
$$
\varphi:I=C_{/F}\longrightarrow F{\quoi}.
$$
We know that the presheaf ind-represented by this ind-object of $C$ is
$F$ (\hyperref[I.prop3.4]{3.4}), so $F$ is ind-representable by
definition (\hyperref[I.sec8.2.2]{8.2.2}).

(ii) $\Leftrightarrow$ (iii). Indeed, (ii) $\Rightarrow$ (iii) because
an ind-representable functor is left exact \hyperref[I.prop8.3.1]{(8.3.1)}, and (iii)
$\Rightarrow$ (ii) because we noted (\hyperref[I.prop8.3.1]{8.3.1}) that if $F$
is left exact, then $C_{/F}$ is filtered, and\pageoriginale
we apply Definition \hyperref[I.def8.1.8]{8.1.8} to conclude that this
category is essentially small.

(ii) $\Leftrightarrow$ (iii bis) is proved in the same way as (ii)
$\Leftrightarrow$ (iii).

The equivalences (ii) $\Leftrightarrow$ (iv) and (iii)
$\Leftrightarrow$ (v) are immediate, since if $C$ is equivalent to
a small category, then $C_{/F}$ is clearly likewise equivalent
to a small category and \textit{a fortiori} it is
automatically essentially small as soon as it is filtered.
This completes the proof of \hyperref[I.them8.3.3]{8.3.3}.

\begin{enonce*}[remark]{Remark 8.3.4}\etiquette[8.3.4]{I.rem8.3.4}
Let $F:C'\rightarrow C''$ be a functor between $𝒰$-categories.
It appears that the notion of left exactness
(\hyperref[I.cor2.3.2]{2.3.2}) is of
little practical interest unless finite projective
limits are representable in $C'$. In the case where
$C''=(𝒰\text{-}\Ens)$, and where $C'$ is $𝒰$-small,
writing $C'$ in the form $C^\circ$ (thus $C={C'}^\circ)$, one
should regard the ``right notion'' that replaces,
in the general case, that of left exactness as the
ind-representability of $F$ (regarded as a
presheaf on
$C^\circ$; \ie the pro-representability of $F$,
regarded as a
functor on $C'$, \cf \hyperref[I.sec8.10]{8.10}). This notion does indeed coincide
with that
of left exactness when finite projective
limits are representable in $C'$ %%
 \textit{i.e.} finite inductive limits in $C$ are
                   %%
representable (\hyperref[I.them8.3.3]{8.3.3}). When $C'$ is no longer assumed
$𝒰$-small, or at least equivalent to a
$𝒰$-small category, the two notions for a functor $F$,
left exactness and pro-representability, no longer necessarily coincide,
even if finite projective limits in $C'$ are
representable (\cf \hyperref[I.exer8.12.9]{8.12.9}). In fact, it seems that
in\pageoriginale this case, left exact functors that are not
ind-representable should be regarded as pathological, while
the ``good'' objects remain the ind-representable functors;
compare \hyperref[I.cor8.13.3]{8.13.3}. For the case of functors
$F:C'\rightarrow C''$, with $C'$ and $C''$ once again arbitrary
$𝒰$-categories, we shall develop below
\hyperref[I.sec8.11.5]{(8.11.5)}, more generally, a notion that ``improves'' upon
left exactness (namely, that of a functor admitting a
pro-adjoint).
\end{enonce*}

% END INCLUDED SGA4_English_current_buildable_source_20260906/source/01/drafts/SGA4_I_LINES_3228_3408_English.texfrag


% BEGIN INCLUDED SGA4_English_current_buildable_source_20260906/source/01/drafts/SGA4_I_LINES_3409_3598_English.texfrag
\subsection{Constant and essentially constant ind-objects}\etiquette[8.4]{I.sec8.4}

Choose a terminal category $\sheaf{e}$, \ie one such that $\Ob
% 2013-5-24 compilation problem with \sheaf \XXX
\sheaf{e}$ and $F\ell$ $\sheaf{e}$ each consist of a single element (for
example, one may take for $\sheaf{e}$ the full subcategory
of $(\Ens)$ on the empty set, if a canonical choice is
desired). This is obviously a small filtered category. If
$X$ is an object of $C$, we associate with it the constant ind-object indexed
by $\sheaf{e}$, with value $X$, denoted by $c(X)$. It is clear that, as
$X$ varies, this gives a fully faithful functor
\hypertarget{sga4.exp.I.equation.8.4.1}{}$$
\displaylines{\rlap{$(8.4.1)$}\hfill c:C\longrightarrow
\Ind(C){\quoi},\hfill}
$$
which is moreover injective on objects, and by which we shall identify
$C$ with a full subcategory of $\Ind(C)$. Note that the
composite functor
\hypertarget{sga4.exp.I.equation.8.4.2}{}$$
\displaylines{\rlap{$(8.4.2)$}\hfill
C\xrightarrow{c}\Ind(C)\xrightarrow{L}\widehat{C}\hfill}
$$
is the canonical functor $h$ \hyperlink{sga4.exp.I.equation.8.2.1.1}{(8.2.1.1)}.

\setcounter{subsubsection}{2}
\subsubsection{}\etiquette[8.4.3]{I.sec8.4.3}
More\pageoriginale generally, let $I$ be a filtered category.
The constant functor on $I$ whose value is an object $X$ of $C$ is
also called \emph{the constant ind-object with value $X$ indexed by
$I$.} It is clear that, if $I$ is essentially small, this
ind-object ind-represents the functor $h(X)$ represented by $X$, hence
this ind-object is isomorphic to $c(X)$.

\subsubsection{}\etiquette[8.4.4]{I.sec8.4.4}
An ind-object $𝒳$ of $C$ is said to be \emph{essentially
constant} if it is isomorphic (in a category
$\Ind_{𝒱}(C,𝒱')$, for suitable $𝒱$, $𝒱'$
) to an ind-object of the form $c(X)$, $X\in \ob C$.
The object $X$, which is then determined up to canonical isomorphism,
is called the \emph{value} of the essentially constant ind-object
under consideration. Thus, by definition, the functor $c$ of \hyperlink{sga4.exp.I.equation.8.4.1}{(8.4.1)} induces an
equivalence of $C$ with the full subcategory of $\Ind(C)$ consisting
of the ind-objects that are essentially constant, this
subcategory being the essential image of the functor $c$.

Clearly a constant ind-object is essentially constant,
whereas the converse holds only in the trivial case in which $C$ is
empty or a one-point category.

\subsection{Filtered inductive limits in $\Ind(C)$}\etiquette[8.5]{I.sec8.5}

\begin{enonce*}{Proposition 8.5.1}\etiquette[8.5.1]{I.prop8.5.1}
Let $C$ be a $𝒰$-category. In $\Ind(C)$, small
filtered inductive limits are representable, and the
canonical functor $\hyperlink{sga4.exp.I.equation.8.2.4.8}{(8.2.4.8)}$
$$
L:\Ind(C)\longrightarrow \widehat{C}
$$
commutes with them.
\end{enonce*}

Since\pageoriginale $L$ is fully faithful, the
assertions of the proposition are equivalent to the following:

\begin{enonce*}{Corollary 8.5.2}\etiquette[8.5.2]{I.cor8.5.2}
Every small filtered inductive limit (in $\widehat{C}$) of
ind-representable presheaves is ind-representable.
\end{enonce*}

This follows easily from criterion \hyperref[I.them8.3.3]{8.3.3} (ii). The
details of the verification are left to the reader.

\setcounter{subsubsection}{2}
\subsubsection{}\etiquette[8.5.3]{I.sec8.5.3}
The ``tautological'' case of filtered inductive limits in
$\Ind(C)$ is that in which one starts with an essentially small
ind-object $𝒳=(X_i)_{i\in I}$ of $C$. Then, in $\widehat{C}$
and hence also in $\Ind(C)$ (or in $\widehat{C}$), one has:
\hypertarget{sga4.exp.I.equation.8.5.3.1}{}$$
\displaylines{\rlap{$(8.5.3.1)$}\hfill 𝒳=\varinjlim_I
X_i{\quoi}.\hfill}
$$
When writing this formula, care must be taken that it is not
an inductive limit in $C$, and that even when the latter
exists, it is not isomorphic in $\Ind(C)$ to $𝒳$: indeed,
the canonical functor \hyperlink{sga4.exp.I.equation.8.4.1}{(8.4.1)} $c:C\rightarrow \Ind(C)$
\emph{does not in general commute with filtered inductive limits.}
(\Cf \hyperref[I.exer8.5.5]{8.5.5} below.)

To avoid this possible confusion in writing
\hyperlink{sga4.exp.I.equation.8.5.3.1}{(8.5.3.1)}, ``some authors'' (= {P. Deligne}) prefer
to write it as
\hypertarget{sga4.exp.I.equation.8.5.3.2}{}$$
\displaylines{\rlap{$(8.5.3.2)$}\hfill
%𝒳=\lq\lq{\varinjlim_I}\textrm{ »\,}
𝒳=\textrm{« }{\varinjlim_I}\textrm{ »\,}X_i{\quoi},\hfill}
$$
the quotation marks indicating that the inductive limit is
taken in a category of ind-objects $\Ind(C)$. By extension, one
should then denote by $\textrm{« }{\lim}\textrm{ »}$ every
inductive-limit operation in $\Ind(C)$ (without the
components of the inductive system under consideration in $\Ind(C)$
necessarily lying in the image, or the essential image, of
$c:C\rightarrow \Ind(C)$).

\subsubsection{}\etiquette[8.5.4]{I.sec8.5.4}
To\pageoriginale compute the inductive or projective limit of a
functor
\hypertarget{sga4.exp.I.equation.8.5.4.1}{}$$
\displaylines{\rlap{$(8.5.4.1)$}\hfill \varphi:J\longrightarrow
\Ind(C){\quoi},\hfill}
$$
where $J$ is now an arbitrary category, it is often
convenient to make $\varphi$ explicit by means of a functor
\hypertarget{sga4.exp.I.equation.8.5.4.2}{}$$
\displaylines{\rlap{$(8.5.4.2)$}\hfill \Phi:J\times I\longrightarrow
C{\quoi},\hfill}
$$
where $I$ is an essentially small filtered category, or
preferably a small one. Such a $\Phi$ indeed defines a functor
$$
J\longmapsto (i\longmapsto \Phi (j,i)):J\longmapsto
\Hhom(I,C){\quoi},
$$
whence, by composing with the canonical arrow \hyperlink{sga4.exp.I.equation.8.2.6.1}{(8.2.6.1)}
$\Hhom(I,C)\rightarrow \Ind(C)$, a functor
\hypertarget{sga4.exp.I.equation.8.5.4.3}{}$$
\displaylines{\rlap{$(8.5.4.3)$}\hfill J\longmapsto (i\longmapsto
\Phi(j,i)):J\longrightarrow \Ind(C){\quoi}.\hfill}
$$
We shall say that $\Phi$ is an \emph{indexed expression of
$\varphi$, indexed by the (filtered) category $I$}, if the corresponding
functor \hyperlink{sga4.exp.I.equation.8.5.4.3}{(8.5.4.3)} is isomorphic to $\varphi$. We shall study
below (\hyperref[I.sec8.8]{8.8}) general conditions for the existence
of an indexed expression
for a given functor $\varphi$, and here we shall restrict ourselves to
starting from a functor $\varphi$ given in indexed form
\hyperlink{sga4.exp.I.equation.8.5.4.3}{(8.5.4.3)}, in the case where $J$ is a small
filtered category,
in order to indicate in this case how to compute
$\textrm{«}{\varinjlim}\textrm{»\,}\varphi=\textrm{«}{\varinjlim_j}\textrm{»\,}\varphi(j)$.
Formula \hyperlink{sga4.exp.I.equation.8.5.3.2}{(8.5.3.2)}, together with the associativity formula for
inductive limits (\hyperref[I.sec2.5.0]{2.5.0}), then immediately gives the ``\emph{tautological
computation}'' of
$\textrm{«}{\varinjlim}\textrm{»\,}\varphi$:
\hypertarget{sga4.exp.I.equation.8.5.4.4}{}$$
\displaylines{\rlap{$(8.5.4.4)$}\hfill
\textrm{«}{\varinjlim_J}\textrm{»\,}\varphi=\textrm{«}{\varinjlim_{J\times
I}}\textrm{»\,}\Phi{\quoi},\hfill}
$$
\ie\pageoriginale
\hypertarget{sga4.exp.I.equation.8.5.4.5}{}$$
\displaylines{\rlap{$(8.5.4.5)$}\hfill
\textrm{«}{\varinjlim_j}\textrm{»\,}(i\longmapsto \Phi
(j,i))=\textrm{«}{\varinjlim_{j,i}}\textrm{»\,}\Phi(j,i){\quoi},\hfill}
$$
Thus, the sought-for inductive system is none other than $\Phi$ itself
(the indexing category being $J\times I$, which is indeed filtered
since $J$ and $I$ are).

\begin{enonce*}[remark]{Exercise 8.5.5}\etiquette[8.5.5]{I.exer8.5.5}
Let $C$ be a $𝒰$-category
\begin{enumerate}
\renewcommand{\theenumi}{\alph{enumi}}
\renewcommand{\labelenumi}{\theenumi)}
\item Prove that the following conditions are equivalent:

\begin{itemize}
\item[(i)] In $C$, small inductive limits are representable,
and the functor $c:C\rightarrow \Ind(C)$ commutes with them.

\item[(iii)] The functor $c$ is an equivalence of categories.

\item[(ii)] In $C$, small inductive limits are
representable, and for every $X\in \Ob C$, the functor $Y\mapsto
\Hom(X,Y)$ commutes with these limits.

\item[(iv)] Small filtered inductive limits are
representable in $C$, and the functor $\lim c:\Ind C\rightarrow C$
(\cf \hyperlink{sga4.exp.I.equation.8.7.1.5}{8.7.1.5}) is fully faithful (or equivalently, an
equivalence).
\end{itemize}

\item If $C$ is a finite category, prove that the preceding
conditions are equivalent to the following: for every
projector in $C$
(\hyperref[I.sec10.6]{10.6}), its image is representable in $C$.
\end{enumerate}
\end{enonce*}

% END INCLUDED SGA4_English_current_buildable_source_20260906/source/01/drafts/SGA4_I_LINES_3409_3598_English.texfrag


% BEGIN INCLUDED SGA4_English_current_buildable_source_20260906/source/01/drafts/SGA4_I_LINES_3599_3796_English.texfrag
\subsection{Extension of a functor to ind-objects}\etiquette[8.6]{I.sec8.6}

\subsubsection{}\etiquette[8.6.1]{I.sec8.6.1}
Let
\hypertarget{sga4.exp.I.equation.8.6.1.1}{}$$
\displaylines{\rlap{$(8.6.1.1)$}\hfill f:C\longrightarrow C'\hfill}
$$
be a functor between $𝒰$-categories. For every ind-object
$$
\varphi:I\longrightarrow C
$$
of\pageoriginale $C$ ($I$ a filtered category), the composite functor
$f\circ\varphi$ is an ind-object of $C'$, also denoted by
$\Ind(f)(\varphi)$. In indexed notation, if $\varphi$ is written as
$𝒳=(X_i)_{i\in I}$, one obtains
\hypertarget{sga4.exp.I.equation.8.6.1.2}{}$$
\displaylines{\rlap{$(8.6.1.2)$}\hfill \Ind(f)(X_i)_{i\in
I}=(f(X_i))_{i\in I}{\quoi}.\hfill}
$$
Let $𝒴=(Y_j)_{j\in J}$ be a second inductive system of $C$. One then
obtains an evident map, also denoted by $u\rightarrow\Ind(f)(u)$:
\hypertarget{sga4.exp.I.equation.8.6.1.3}{}$$
\displaylines{\rlap{$(8.6.1.3)$}\hfill
\Hom(𝒳,𝒴)=\varprojlim_i \varinjlim_j \Hom
(X_i,Y_j)\rightarrow\hfill\cr \hfill
\Hom(\Ind(f)(𝒳),\Ind(f)(𝒴))\simeq \varprojlim_i
\varinjlim_j \Hom(f(X_i),f(Y_j)){\quoi}.\hfill}
$$
Referring to the unwritten formula for composition of morphisms of
ind-objects (\hyperref[I.sec8.2.5]{8.2.5}), one immediately sees that
these maps are compatible with composition of morphisms of ind-objects.
Thus, for every universe $𝒱$, we have obtained a functor
\hypertarget{sga4.exp.I.equation.8.6.1.4}{}$$
\displaylines{\rlap{$(8.6.1.4)$}\hfill
\Ind(f):\Ind_{𝒱}(C,𝒰)\longrightarrow
\Ind_{𝒱}(C',𝒰)\hfill}
$$
(notation of \hyperlink{sga4.exp.I.equation.8.2.4.5}{(8.2.4.5)}), and in particular a functor
\hypertarget{sga4.exp.I.equation.8.6.1.5}{}$$
\displaylines{\rlap{$(8.6.1.5)$}\hfill
\Ind(f):\Ind(C)\longrightarrow \Ind(C'){\quoi}.\hfill}
$$
If there is a second functor
$$
g:C'\longrightarrow C''{\quoi},
$$
then one clearly has an identity of functors between $\Ind$-categories
(or $\Ind_{𝒱}$-categories, as desired):
\hypertarget{sga4.exp.I.equation.8.6.1.6}{}$$
\displaylines{\rlap{$(8.6.1.6)$}\hfill \Ind(gf)=\Ind(g)\circ
\Ind(f){\quoi},\hfill}
$$
and,\pageoriginale when $C'=C$, the perennial formula
\hypertarget{sga4.exp.I.equation.8.6.1.7}{}$$
\displaylines{\rlap{$(8.6.1.7)$}\hfill
\Ind(\id_C)=\id_{\Ind(C)}{\quoi}.\hfill}
$$
One may therefore say, heuristically, that the category $\Ind(C)$
\emph{depends functorially on $C$} (covariantly)
\footnote[1]{For a contravariant dependence, under certain conditions,
\cf \hyperref[I.sec8.11]{8.11}.}. We leave to the reader the task of making explicit even a
2-\emph{functorial} dependence by defining, for every pair of
$𝒰$-categories $C$, $C'$, a canonical functor
\hypertarget{sga4.exp.I.equation.8.6.1.8}{}$$
\displaylines{\rlap{$(8.6.1.8)$}\hfill
\Iind(f):\Hhom(C,C')\longrightarrow \Hhom(\Ind(C),
\Ind(C')){\quoi},\hfill}
$$
and specifying the functorial nature of the identical isomorphisms
\hyperlink{sga4.exp.I.equation.8.6.1.6}{(8.6.1.6)} and \hyperlink{sga4.exp.I.equation.8.6.1.7}{(8.6.1.7)}.

\subsubsection{}\etiquette[8.6.2]{I.sec8.6.2}
Return to the functor \hyperlink{sga4.exp.I.equation.8.6.1.1}{(8.6.1.1)}, and consider the corresponding functor
\hyperref[I.prop5.1]{(5.1)}
\hypertarget{sga4.exp.I.equation.8.6.2.1}{}$$
\displaylines{\rlap{$(8.6.2.1)$}\hfill
f_{!}:\widehat{C}\longrightarrow \widehat{C'}{\quoi},\hfill}
$$
and the diagram of functors
\hypertarget{sga4.exp.I.equation.8.6.2.2}{}\[\vcenter{
\xymatrix@C=3.25cm@R=1.25cm{
\Ind(C) \ar[r]^{\ind(f)} \ar[d]_{L_C}& \Ind(C')\ar[d]^{L_{C'}}\\
\hat{C} \ar[r]^{f_!}& \hat{C}' }}\leqno{(8.6.2.2)}
\]
where the vertical arrows denote the canonical functors $L$ of
\hyperlink{sga4.exp.I.equation.8.2.4.8}{(8.2.4.8)}. Since the functor $f_{!}$ ``extends $f$'' and commutes with
inductive limits \hyperlink{sga4.exp.I.proposition.5.4.item.3}{(5.4.3)}, and the functors $L_C$ and $L_{C'}$ commute
with small filtered inductive limits (\hyperref[I.prop8.5.1]{8.5.1}), it
follows from \hyperlink{sga4.exp.I.equation.8.5.3.2}{(8.5.3.2)} that, for every ind-object $(X_i)_{i\in I}$ of
$C$, there is a canonical isomorphism in $\widehat{C'}$
$$
f_{!}(L_C((X_i)_{i\in I}))\simeq \varinjlim_i f_{!}L_C(X_i)\simeq
\varinjlim_i L_{C'}f(X_i)=L_{C'}(f(X_i)_{i\in I})
$$
\ie\pageoriginale a canonical isomorphism
$$
f_{!}L_C(𝒳)\simeq L_{C'}(\Ind(f)(𝒳)){\quoi}.
$$
We leave to the reader the verification that the latter is functorial,
\ie that it corresponds to a canonical isomorphism
\hypertarget{sga4.exp.I.equation.8.6.2.3}{}$$
\displaylines{\rlap{$(8.6.2.3)$}\hfill f_{!}L_C\simeq L_{C'}\circ
\Ind(f){\quoi}.\hfill}
$$
In other words, the \emph{square \hyperlink{sga4.exp.I.equation.8.6.2.2}{(8.6.2.2)} commutes up to canonical
isomorphism.} Taking into account the fact that $f_{!}$ commutes with
small inductive limits, and using \hyperref[I.prop8.5.1]{8.5.1}, we conclude:

\begin{enonce*}{Proposition 8.6.3}\etiquette[8.6.3]{I.prop8.6.3}
Let $f:C\rightarrow C'$ be a functor between $𝒰$-categories. Then the
functor
$$
\Ind(f):\Ind(C)\longrightarrow \Ind(C')
$$
commutes with filtered inductive limits. Moreover, it makes the following
diagram commutative:
\[
\xymatrix@C=2.5cm@R=1.5cm{
C       \ar[r]^{f}       \ar[d]_{c_{C}}& C' \ar[d]^{c_{C'}}\\
\Ind(C) \ar[r]^{\Ind(f)} & \Ind(C')\quad , }
\]
where the vertical arrows $c_C$, $c_{C'}$ are the canonical functors
$\hyperlink{sga4.exp.I.equation.8.4.1}{(8.4.1)}$.
\end{enonce*}
The last assertion is trivial from the definitions and has been included
for ease of reference. Notice moreover that the properties stated in
\hyperref[I.prop8.6.3]{8.6.3} \emph{characterize the functor} $\Ind(f)$ up to unique
isomorphism,\pageoriginale as the functor induced by $f_{!}$ \hyperlink{sga4.exp.I.equation.8.6.2.1}{(8.6.2.1)},
as follows from the proof of \hyperlink{sga4.exp.I.equation.8.6.2.3}{(8.6.2.3)} just given.

\begin{enonce*}{Proposition 8.6.4}\etiquette[8.6.4]{I.prop8.6.4}
The notation is that of $\hyperref[I.prop8.6.3]{8.6.3}$.
\begin{enumerate}
\renewcommand{\theenumi}{\alph{enumi}}
\renewcommand{\labelenumi}{\theenumi)}
\item The functor $\Ind(f)$ is faithful (\resp fully faithful) if and
only if $f$ is.

\item The functor $\Ind(f)$ is an equivalence of categories if and only
if $f$ is fully faithful and every object of $C'$ is isomorphic to a
direct factor \hyperref[I.sec10.6]{(10.6)} of an object in the image of $f$.
\end{enumerate}
\end{enonce*}

\noindent \emph{Proof.}
\begin{enumerate}
\renewcommand{\theenumi}{\alph{enumi}}
\renewcommand{\labelenumi}{\theenumi)}
\item Necessity clearly follows from the fact that $f$ is induced by
$\Ind(f)$. For sufficiency, it is enough to use the form \hyperlink{sga4.exp.I.equation.8.6.1.3}{(8.6.1.3)} of
$\Ind(f)$ on the sets $\Hom(𝒳,𝒴)$, recalling that filtered inductive
limits of sets and arbitrary projective limits take monomorphisms to
monomorphisms and isomorphisms to isomorphisms.

\item We may already suppose that $f$, hence $\Ind(f)$, is fully faithful.
Since every object of $\Ind(C')$ is a small filtered inductive limit of
objects of $C'$, the full faithfulness of $\Ind(f)$ implies that, for this
functor to be essentially surjective, it is equivalent that every object
$X'$ of $C'$ belong to the essential image. But if there is an
isomorphism
$$
X'\xrightarrow{\sim} \textrm{« }{\varinjlim}\textrm{ »\,}f(X_i){\quoi},
$$
this isomorphism factors through one of the $f(X_i)$, showing that $X'$
is a direct factor of $f(X_i)$ and proving necessity in b). For
sufficiency, using the full faithfulness of $\Ind(f)$, the assertion
follows immediately.
\end{enumerate}

\begin{enonce*}{Corollary 8.6.5}\etiquette[8.6.5]{I.cor8.6.5}
In\pageoriginale $\Ind(E)$ the images of projectors
(\hyperref[I.sec10.6]{10.6}) are representable, and the functor
$\Ind(f):\Ind(C)\rightarrow \Ind(C')$ commutes with the formation of
these images.
\end{enonce*}

Indeed, this follows from the fact that small filtered inductive limits
are representable in $\Ind(C)$ (\hyperref[I.prop8.5.1]{8.5.1}) and that
$\Ind(f)$ commutes with them (\hyperref[I.prop8.6.3]{8.6.3}), taking into
account that the image of a projector $p:X\rightarrow X$ may be
interpreted as the inductive limit of a filtered inductive system indexed
by $\NN$
$$
X\xrightarrow{P}X\xrightarrow{P}X\longrightarrow\ldots
$$
or, equivalently, as the inductive limit of the evident functor on the
filtered category $P$ having one object and one nonidentity arrow $\Pi$
such that $\Pi^2=\Pi$.


% END INCLUDED SGA4_English_current_buildable_source_20260906/source/01/drafts/SGA4_I_LINES_3599_3796_English.texfrag


% BEGIN INCLUDED SGA4_English_current_buildable_source_20260906/source/01/drafts/SGA4_I_LINES_3797_3955_English.texfrag
\subsection{}\etiquette[8.7]{I.sec8.7}
\emph{The functor $\varinjlim_C:\Ind(C)\rightarrow C$.
Universal characterizations of the category $\Ind(C)$.}

\subsubsection{}\etiquette[8.7.1]{I.sec8.7.1}
Let us return to a $𝒰$-category $C$ and the canonical functor
\hypertarget{sga4.exp.I.equation.8.7.1.1}{}$$
\displaylines{\rlap{$(8.7.1.1)$}\hfill c:C\longrightarrow
\Ind(C){\quoi}.\hfill}
$$
Let $𝒳=(X_i)_{i\in I}$ be an object of $\Ind(C)$. It follows
immediately from the definition of representable inductive limits
(\hyperref[I.def2.1]{2.1}, \hyperref[I.sec2.1.1]{2.1.1}) that
$\varinjlim_iX_i$ \emph{is representable in $C$ if and only if the left
adjoint of $c$ is defined at $𝒳$}, and that in this case
$\varinjlim_iX_i$ (computed in $C$) is precisely the value at $𝒳$ of
that left adjoint:
\hypertarget{sga4.exp.I.equation.8.7.1.2}{}$$
\displaylines{\rlap{$(8.7.1.2)$}\hfill \Hom_C(\varinjlim_i X_i,
Y)\simeq \Hom_{\Ind(C)}((X_i)_{i\in I}, c(Y)){\quoi}.\hfill}
$$
If\pageoriginale we denote by
\hypertarget{sga4.exp.I.equation.8.7.1.3}{}$$
\displaylines{\rlap{$(8.7.1.3)$}\hfill \Ind(C)'\subset
\Ind(C)\hfill}
$$
the full subcategory of $\Ind(C)$ formed by the ind-objects of $C$ that
admit an inductive limit \emph{in} $C$, the preceding observation shows
that this subcategory is \emph{strictly full}, and that there is a
canonical functor
\hypertarget{sga4.exp.I.equation.8.7.1.4}{}$$
\displaylines{\rlap{$(8.7.1.4)$}\hfill
\textstyle{\varinjlim}_C:\Ind(C)'\longrightarrow C{\quoi},\hfill}
$$
whose value at each object $(X_i)_{i\in I}$ of $\Ind(C)'$ is its
inductive limit in $C$. In the special case where small filtered
inductive limits are representable in $C$, one therefore obtains a
natural functor
\hypertarget{sga4.exp.I.equation.8.7.1.5}{}$$
\displaylines{\rlap{$(8.7.1.5)$}\hfill
{\varinjlim}_C:\Ind(C)\longrightarrow C{\quoi}.\hfill}
$$

\setcounter{paragraph}{5}
\paragraph{}\etiquette[8.7.1.6]{I.sec8.7.1.6}
Of course, the preceding functors may also be defined on categories of
the form $\Ind_{𝒱}(C,𝒰)'$ and $\Ind_{𝒱}(C,𝒰)$, but in view of
$\hyperlink{sga4.exp.I.equation.8.2.4.6}{(8.2.4.6)}$ they are already determined (up to unique isomorphism) by
the preceding functors corresponding to the case $𝒱=𝒰$.

\paragraph{}\etiquette[8.7.1.7]{I.sec8.7.1.7}
The preceding construction of $\varinjlim_C$ as a left adjoint functor
immediately implies that this functor \emph{commutes with arbitrary
inductive limits}, and in particular with small filtered inductive limits
(the latter being representable in $\Ind(C)$
(\hyperref[I.prop8.5.1]{8.5.1})). Notice also that $\Ind(C)'$ always
contains the essential image of $c$ (formed by the essentially constant
ind-objects), and that there is a canonical isomorphism, functorial in
$X\in\ob\Ind(C)'$:
\hypertarget{sga4.exp.I.equation.8.7.1.8}{}$$
\displaylines{\rlap{$(8.7.1.8)$}\hfill {\varinjlim}_C c(X)\simeq
X{\quoi},\hfill}
$$
\ie,\pageoriginale in the favorable case where $C$ is stable under small
inductive limits, there is a canonical isomorphism
\hypertarget{sga4.exp.I.equation.8.7.1.9}{}$$
\displaylines{\rlap{$(8.7.1.9)$}\hfill {\varinjlim}_C \circ c\simeq
\id_C{\quoi}.\hfill}
$$

\subsubsection{}\etiquette[8.7.2]{I.sec8.7.2}

Now consider a functor
\hypertarget{sga4.exp.I.equation.8.7.2.1}{}$$
\displaylines{\rlap{$(8.7.2.1)$}\hfill f:C\longrightarrow
E{\quoi},\hfill}
$$
where $E$ is a $𝒰$-category in which small inductive limits are
representable, so that there is a functor
$$
{\varinjlim}_E:\Ind(E)\longrightarrow E{\quoi}.
$$
Since we have also defined \hyperref[I.sec8.6]{(8.6)}
$$
\Ind(f):\Ind(C)\longrightarrow \Ind(E){\quoi},
$$
we may consider the composite
\hypertarget{sga4.exp.I.equation.8.7.2.2}{}$$
\displaylines{\rlap{$(8.7.2.2)$}\hfill \bar{f}={\varinjlim}_E\circ
\Ind(f):\Ind(C)\longrightarrow E{\quoi},\hfill}
$$
sometimes called the \emph{canonical extension of $f$ to ind-objects}
(but not to be confused with $\Ind(f)$!). As a composite of two functors
that commute with small filtered inductive limits
(\hyperref[I.prop8.6.3]{8.6.3}, \hyperref[I.sec8.7.1.7]{8.7.1.7}), this
functor itself commutes with small filtered inductive limits. Moreover,
the isomorphism \hyperlink{sga4.exp.I.equation.8.7.1.9}{(8.7.1.9)} applied to $E$, together with \hyperlink{sga4.exp.I.equation.8.6.2.3}{(8.6.2.3)}, shows
that $\bar f$ \emph{extends} $f$ up to isomorphism, \ie there is a
canonical isomorphism
\hypertarget{sga4.exp.I.equation.8.7.2.3}{}$$
\displaylines{\rlap{$(8.7.2.3)$}\hfill \bar{f}\circ c_C\simeq
f{\quoi}.\hfill}
$$
It\pageoriginale is moreover fairly clear, in view of the fact that every
object of $\Ind(C)$ is a small filtered inductive limit of objects of $C$,
that the two preceding properties still \emph{characterize} $\bar f$ up
to unique isomorphism. This can be sharpened to obtain at the same time a
universal characterization, up to equivalence, of $\Ind(C)$ among the
$𝒰$-categories $E$ in which small filtered inductive limits are
representable. If $E$ and $F$ are two such $𝒰$-categories, denote by
\hypertarget{sga4.exp.I.equation.8.7.2.4}{}$$
\displaylines{\rlap{$(8.7.2.4)$}\hfill \Hhom(E,F)'\subset
\Hhom(E,F)\hfill}
$$
the full subcategory of $\Hhom(E,F)$ formed by the functors that commute
with small filtered inductive limits. Then:

\begin{enonce*}{Proposition 8.7.3}\etiquette[8.7.3]{I.prop8.7.3}
Let $C$ be a $𝒰$-category, and use the notation \hyperlink{sga4.exp.I.equation.8.7.2.4}{(8.7.2.4)} above. Then
the canonical functor
$$
c_C:C\longrightarrow \Ind(C)
$$
is $2$-universal among the functors with source $C$ and target a
$𝒰$-category having representable small filtered inductive limits. In
other words, $\Ind(C)$ is such a category
(\hyperref[I.sec8.2.5]{8.2.5}, \hyperref[I.prop8.5.1]{8.5.1}), and if $E$
is such a category, the functor
\hypertarget{sga4.exp.I.equation.8.7.3.1}{}$$
\displaylines{\rlap{$(8.7.3.1)$}\hfill g\longmapsto g\circ
c_C:\Hhom(\Ind(C),E)'\longrightarrow \Hhom(C,E)\hfill}
$$
is an equivalence of categories.
\end{enonce*}

To prove the latter point, it is enough to verify that the map
$f\mapsto\bar f$ defined by \hyperlink{sga4.exp.I.equation.8.7.2.2}{(8.7.2.2)} can be refined to a \emph{functor}
\hypertarget{sga4.exp.I.equation.8.7.3.2}{}$$
\displaylines{\rlap{$(8.7.3.2)$}\hfill f\longmapsto
\bar{f}={\varinjlim}_E\cdot \Ind(f):\Hhom(C,E)\longrightarrow
\Hhom(\Ind(C),E)'{\quoi},\hfill}
$$
and\pageoriginale that the latter is a quasi-inverse of \hyperlink{sga4.exp.I.equation.8.7.3.1}{(8.7.3.1)}. The
details of the essentially trivial verification are left to the reader.
Alternatively, one may invoke \hyperref[I.prop7.8]{7.8} to conclude first that \hyperlink{sga4.exp.I.equation.8.7.3.1}{(8.7.3.1)} is
fully faithful, and then \hyperlink{sga4.exp.I.equation.8.7.2.3}{(8.7.2.3)} to conclude that it is essentially
surjective.

\setcounter{subsubsection}{3}

% END INCLUDED SGA4_English_current_buildable_source_20260906/source/01/drafts/SGA4_I_LINES_3797_3955_English.texfrag


% BEGIN INCLUDED SGA4_English_current_buildable_source_20260906/source/01/drafts/SGA4_I_LINES_3956_4068_English.texfrag
\subsubsection{}\etiquette[8.7.4]{I.sec8.7.4}
Return to a functor between $𝒰$-categories
\hypertarget{sga4.exp.I.equation.8.7.4.1}{}$$
\displaylines{\rlap{$(8.7.4.1)$}\hfill f:C\longrightarrow
E{\quoi},\hfill}
$$
where $E$ is a $𝒰$-category in which small filtered inductive limits
are representable, giving a functor
\hypertarget{sga4.exp.I.equation.8.7.4.2}{}$$
\displaylines{\rlap{$(8.7.4.2)$}\hfill \bar{f}:(X_i)_{i\in I}
\longmapsto \mathop{\lim}_i f(X_i):\Ind(C)\longrightarrow
E{\quoi}.\hfill}
$$
We propose to study conditions on $f$ ensuring that $\bar f$ is fully
faithful, \resp an equivalence of categories. Since $f$ is isomorphic to
the composite $\bar f\circ c_C$, and $c_C:C\rightarrow\Ind(C)$ is fully
faithful, we see that for $\bar f$ to be fully faithful, it is
\emph{necessary} that $f$ be so. Replacing $C$ by its essential image in
$E$, we therefore lose essentially no generality by supposing that $f$ is
the inclusion functor of a subcategory $C$ of $E$, as we shall do below
to simplify the notation.

\begin{enonce*}{Proposition 8.7.5}\etiquette[8.7.5]{I.prop8.7.5}
The notation is that of \hyperref[I.sec8.7.4]{8.7.4}.
\begin{enumerate}
\renewcommand{\theenumi}{\alph{enumi}}
\renewcommand{\labelenumi}{\theenumi)}
\item The functor $\bar f$ is fully faithful if and only if:
\begin{itemize}
\item[(i)] Every object $X$ of $C$ satisfies the following condition:

{\rm(PF)} For every small filtered inductive system $(Y_i)_{i\in I}$ in
$C$ whose inductive limit $\varinjlim_iY_i$ exists in\pageoriginale $E$,
the canonical map
\hypertarget{sga4.exp.I.equation.8.7.5.1}{}$$
\displaylines{\rlap{$(8.7.5.1)$}\hfill \varinjlim_i \Hom
(X,Y_i)\longrightarrow \Hom(X,\varinjlim_i Y_i)\hfill}
$$
is bijective.
\end{itemize}

\item The functor $\bar f$ is an equivalence of categories if and only if
$C$ satisfies condition {\rm(i)} and the following two conditions:
\begin{itemize}
\item[(ii)] $C$ is a generating subcategory of $E$ by strict
epimorphisms (\hyperref[I.def7.1]{7.1}).

\item[\rm(iii)] For every object $X$ of $E$, the full subcategory
$C_{/X}$ of $E_{/X}$ formed by the objects $X'$ over $X$ whose source is
in $\Ob C$ is filtered and essentially small.
\end{itemize}
\end{enumerate}

Condition {\rm(iii)} is satisfied in particular if $C$ is equivalent to
a small category, the finite inductive limits in $C$ are representable,
and the inclusion functor $f:C\rightarrow E$ commutes with them (\ie for
every finite category $J$ and every functor $\varphi:J\rightarrow C$, the
inductive limit of $f\varphi$ is representable and is isomorphic in $E$
to an object of $C$).
\end{enonce*}


\noindent \emph{Proof}.
\begin{enumerate}
\renewcommand{\theenumi}{\alph{enumi}}
\renewcommand{\labelenumi}{\theenumi)}
\item With the notation of condition (i), let $𝒴$ denote the ind-object
$(Y_i)$ and let $𝒳$ denote the constant ind-object defined by $X$. Then
\hyperlink{sga4.exp.I.equation.8.7.5.1}{(8.7.5.1)} is precisely the canonical map
$$
\Hom(𝒳, 𝒴)\longrightarrow \Hom
(\bar{f}(𝒳),\bar{f}(𝒴)){\quoi},
$$
so if $\bar f$ is fully faithful, this map is indeed bijective.
Conversely, if $𝒳=(X_j)_{j\in J}$ and $𝒴=(Y_i)_{i\in I}$ are arbitrary
ind-objects of $E$, then the map
$u:\Hom(𝒳,𝒴)\rightarrow\Hom(\bar f(𝒳),\bar f(𝒴))$ is\pageoriginale the
projective limit over $j$ of the maps
$u_j:\Hom(𝒳_j,𝒴)\rightarrow\Hom(\bar f(𝒳_j),𝒴)$, where $𝒳_j$ is the
constant ind-object with value $X_j$. Thus $u$ is bijective if the $u_j$
are, and (i) therefore implies that $\bar f$ is fully faithful.

\item Suppose that (i), (ii), and (iii) hold, and let us prove that
$\bar f$ is an equivalence. It remains to prove that it is essentially
surjective, hence that every $X\in\ob E$ belongs to the essential image.
But hypothesis (ii) means that $\varinjlim_{C_{/X}}Z\rightarrow X$ is an
isomorphism, while (iii) says that the category $C_{/X}$ is filtered and
essentially small. Thus $X$ is the image of the ind-object of $E$ defined
by the natural functor $C_{/X}\rightarrow E$. Conversely, suppose that
$\bar f$ is an equivalence. In view of a), it remains to verify (ii) and
(iii). We may suppose that $E=\Ind(C)$, with $C$ identified with the
subcategory $c(C)$ of $E$. But we know (\hyperref[I.them8.3.3]{8.3.3} (ii)) that for every
$X\in\Ind(C)$, identified if desired with the functor $F$ on $C^\circ$
that it ind-represents, $C_{/X}=C_{/F}$ is an essentially small filtered
category (which proves (iii)), and that the inductive limit in
$\widehat{C}$ of the functor $C_{/F}\rightarrow\widehat{C}$ is $F$. It
is therefore also so in the full subcategory $E$ of $\widehat C$, since
$F=X$ belongs to $E$, which proves (ii). It remains to prove the last
assertion of b). But the hypothesis on $C$ clearly implies that finite
inductive limits are representable in the category $C_{/X}$; \textit{a
fortiori}, the category $C_{/X}$ is filtered.
\end{enumerate}


% END INCLUDED SGA4_English_current_buildable_source_20260906/source/01/drafts/SGA4_I_LINES_3956_4068_English.texfrag


% BEGIN INCLUDED SGA4_English_current_buildable_source_20260906/source/01/drafts/SGA4_I_LINES_4069_4270_English.texfrag
\begin{enonce*}[remark]{Remark 8.7.5.2}\etiquette[8.7.5.2]{I.rem8.7.5.2}
The\pageoriginale preceding argument shows more generally that, when
condition (i) is satisfied, the essential image of the fully faithful
functor $\bar f$ \hyperlink{sga4.exp.I.equation.8.7.4.2}{(8.7.4.2)} consists of the $X\in\ob E$ such that the
category $C_{/X}$ is filtered and essentially small (a condition
automatically satisfied if $C$ satisfies the conditions stated at the end
of b)), \emph{and} such that the inductive limit of the canonical functor
$C_{/X}\rightarrow E$ is $X$.
\end{enonce*}

\begin{enonce*}{Corollary 8.7.6}\etiquette[8.7.6]{I.cor8.7.6}
Let $E_{PF}$ be the full subcategory of $E$ formed by the objects
satisfying condition $(PF)$ of \hyperref[I.prop8.7.5]{8.7.5} {\rm a)}.
Then, for every functor $J\rightarrow E_{PF}$, where $J$ is a finite
category, that admits an inductive limit $X$ in $E$, one has
$X\in\ob E_{PF}$. The canonical functor
$(X_i)_{i\in I}\rightarrow\displaystyle{\mathop{\varinjlim}_i}{}_EX_i$
deduced from the inclusion $g:E_{PF}\hookrightarrow E$
\hypertarget{sga4.exp.I.equation.8.7.6.1}{}$$
\displaylines{\rlap{$(8.7.6.1)$}\hfill
\bar{g}:\Ind(E_{PF})\longrightarrow E\hfill}
$$
is fully faithful. If finite inductive limits are representable in $E$
and $E_{PF}$ is equivalent to a small category, then the essential image
of the functor $\bar f$ consists of the objects $X$ of $E$ such that the
canonical morphism $\varinjlim_{(E_{PF})_{/X}}Y\rightarrow X$ is an
isomorphism. If, moreover, the subcategory $E_{PF}$ of $E$ is generating
by strict epimorphisms (\hyperref[I.def7.1]{7.1}), then the functor
\hyperlink{sga4.exp.I.equation.8.7.6.1}{(8.7.6.1)} is an equivalence of categories.
\end{enonce*}

All the assertions are evident, in the order in which they are stated,
in view of \hyperref[I.prop8.7.5]{8.7.5}, using
\hyperref[I.prop2.8]{2.8} for the first assertion.

\begin{enonce*}{Corollary 8.7.7}\etiquette[8.7.7]{I.cor8.7.7}
Suppose\pageoriginale that finite inductive limits are representable in
$E$. Let $C$ be a full subcategory of $E$, equivalent to a small
category, and consider the functor
$(X_i)_{i\in I}\rightarrow
{\displaystyle{\mathop{\varinjlim}_i}{}_E}X_i$:
\hypertarget{sga4.exp.I.equation.8.7.7.1}{}$$
\displaylines{\rlap{$(8.7.7.1)$}\hfill \bar{f}:\Ind(C)\rightarrow
E{\quoi}.\hfill}
$$
Let $E_{PF}$ be as in $\hyperref[I.cor8.7.6]{8.7.6}$. The following conditions are equivalent:
\begin{itemize}
\item[(i)] The functor $\bar f:\Ind(C)\rightarrow E$ is an equivalence.

\item[(ii)] The category $C$ is generating in $E$ by strict
epimorphisms, is contained in $E_{PF}$, and every object of $E_{PF}$ is
isomorphic in $E$ (or in $E_{PF}$, which is equivalent
(\hyperref[I.cor8.7.6]{8.7.6})) to a direct factor of an object of $C$.

When the subcategory $C$ of $E$ is stable under direct factors
(\hyperref[I.sec10.6]{10.6}), the preceding conditions are also
equivalent to the following:

\item[(ii bis)] $C=E_{PF}$, and $C$ is generating in $E$ by strict
epimorphisms.

\item[(iii)] The subcategory $C$ of $E$ is generating by strict
epimorphisms, is contained in $E_{PF}$, and finite inductive limits are
representable in it.

When, moreover, finite inductive limits are representable in $E$, these
conditions are also equivalent to

\item[(iii bis)] The subcategory $C$ of $E$ is generating by strict
epimorphisms, is contained in $E_{PF}$, and is stable in $E$ under finite
inductive limits (or, equivalently, under finite sums and cokernels of
double arrows).
\end{itemize}
\end{enonce*}

We\pageoriginale already know (\hyperref[I.prop8.7.5]{8.7.5}) that
condition (i) implies $C\subset E_{PF}$ and that $C$ is generating by
strict epimorphisms. Let us also prove that every object $X$ of $E_{PF}$
is then isomorphic to a direct factor of an object of $C$. Indeed, we know
that $X$ is a filtered inductive limit $\varinjlim_iX_i$ in $E$ of
objects of $C$, and by the definition of $E_{PF}$, for a suitable $i$ the
canonical homomorphism $X_i\rightarrow X$ has a left inverse, so $X$ is
indeed a direct factor of the object $X_i$ of $C$. This proves
(i) $\Rightarrow$ (ii) (without any hypothesis on $C$ or $E$, moreover).
Let us prove (ii) $\Rightarrow$ (i). Since $C$ is equivalent to a small
category and every object of $E_{PF}$ is a direct factor of an object of
$C$, we see that $E_{PF}$ is also equivalent to a small category. Thus,
by \hyperref[I.cor8.7.6]{8.7.6}, the functor \hyperlink{sga4.exp.I.equation.8.7.6.1}{(8.7.6.1)} is an equivalence
of categories, and we conclude by \hyperref[I.prop8.6.4]{8.6.4} b) applied
to the inclusion $E\hookrightarrow E_{PF}$.

When every direct factor in $E$ of an object of $C$ belongs to $C$, it is
clear that (ii) $\Leftrightarrow$ (ii bis). On the other hand, (ii bis)
implies (iii) \resp (iii bis) by \hyperref[I.cor8.7.6]{8.7.6}. Finally, (iii) (and
\textit{a fortiori} (iii bis)) implies (i) by
\hyperref[I.prop8.7.5]{8.7.5}. This completes the proof.

\begin{enonce*}[remark]{Exercise 8.7.8}\etiquette[8.7.8]{I.exer8.7.8}
(\emph{Karoubi envelopes}). Let $C$ be a $𝒰$-category, let
$E=\Ind(C)$, and let $E_{PF}\supset C$ be the full subcategory of $E$
defined in \hyperref[I.cor8.7.6]{8.7.6}.
\begin{enumerate}
\renewcommand{\theenumi}{\alph{enumi}}
\renewcommand{\labelenumi}{\theenumi)}
\item Prove that images of projectors are representable in $E_{PF}$ and
that every object of $E_{PF}$ is a direct factor of an object of $C$.

\item Call a category $F$ \emph{Karoubian} if images of projectors are
representable in it. Let
$$
\varphi:C\longrightarrow K
$$
be\pageoriginale a fully faithful functor such that $K$ is Karoubian and
every object of $K$ is a direct factor of an object in the image of $C$.
Prove that $\varphi$ is 2-\emph{universal} among functors from $C$ with
values in Karoubian categories, more precisely: for every Karoubian
category $F$, the functor
$$
g\longmapsto g\circ \varphi :\Hhom(K,F)\longrightarrow \Hhom(K,F)
$$
is an equivalence of categories. (Use the fact that every functor commutes
with images of projectors.) The category $K$ equipped with $\varphi$,
determined up to equivalence (itself determined up to unique isomorphism)
by the preceding properties, is called the \emph{Karoubi envelope}
$\widetilde C$ of $C$. (Compare also IV \hyperref[I.cor7.5]{7.5}.)

\item Deduce from a) and b) that $\Ind(C)_{PF}$, equipped with the functor
$C\rightarrow\Ind(C)_{PF}$ induced by $c:C\rightarrow\Ind(C)$, is a
Karoubi envelope of $C$.

\item Show that every functor $f:C\rightarrow C'$ extends in an
essentially unique way to a functor
$\widetilde f:\widetilde C\rightarrow\widetilde{C'}$ of Karoubi
envelopes, and that if these Karoubi envelopes are taken as in c),
$\widetilde f$ is the functor induced by
$\Ind(f):\Ind(C)\rightarrow\Ind(C')$. Show that the conditions of
\hyperref[I.sec8.5.4]{8.5.4} %%
b) %%
are also equivalent to the following: $\widetilde f$ is an equivalence of
categories.
\end{enumerate}
\end{enonce*}

\begin{enonce*}[remark]{Exercise 8.7.9}\etiquette[8.7.9]{I.exer8.7.9}
Let $E$ be a $𝒰$-category.
\begin{enumerate}
\renewcommand{\theenumi}{\alph{enumi}}
\renewcommand{\labelenumi}{\theenumi)}
\item Show that the following conditions are equivalent:

\begin{itemize}
\item[(i)] $E$ is equivalent to a category of the form $\Ind(C)$, where
$C$ is a $𝒰$-category.

\item[(i bis)] As\pageoriginale in (i), with $C$ moreover Karoubian
(\hyperref[I.exer8.7.8]{8.7.8} b)).

\item[(ii)] The subcategory $E_{PF}$ of $E$
(\hyperref[I.cor8.7.6]{8.7.6}) is generating by strict epimorphisms, and
for every object $X$ of $E$, $(E_{PF})_{/X}$ is an essentially small
filtered category.
\end{itemize}

Show that if $C$ is a \emph{Karoubian} $𝒰$-category such that $E$ is
equivalent to $\Ind(C)$, then $C$ is equivalent to $E_{PF}$ (by an
equivalence determined up to unique isomorphism). Establish a
2-equivalence between the 2-category formed by the categories $E$
satisfying the preceding conditions and the functors between them that
\emph{commute with small filtered inductive limits}, and the 2-category
formed by \emph{Karoubian} $𝒰$-categories and arbitrary functors between
them.

\item Show that the following conditions are equivalent:

\begin{itemize}
\item[(i)] $E$ is equivalent to a category of the form $\Ind(C)$, where
$C$ is a $𝒰$-category in which finite inductive limits are
representable.

\item[(i bis)] As in (i), but with $C$ moreover Karoubian.

\item[(ii)] Finite inductive limits in $E$ are representable, the
subcategory $E_{PF}$ of $E$ is generating by strict epimorphisms, and for
every object $X$ of $E$, there exists a small subset $I$ of
$\Ob E_{PF/X}$ such that every object of $E_{PF/X}$ is bounded above by
an object of $I$. (Use \hyperref[I.prop8.9.5]{8.9.5} b).)
\end{itemize}
\end{enumerate}
\end{enonce*}


% END INCLUDED SGA4_English_current_buildable_source_20260906/source/01/drafts/SGA4_I_LINES_4069_4270_English.texfrag


% BEGIN INCLUDED SGA4_English_current_buildable_source_20260906/source/01/drafts/SGA4_I_LINES_4271_4625_English.texfrag
\subsection{Indexed representation of a functor $J\rightarrow
\Ind(C)$}\etiquette[8.8]{I.sec8.8}\pageoriginale

\subsubsection{}\etiquette[8.8.1]{I.sec8.8.1}
Let
$$
\varphi:J\longrightarrow \Ind(C)\text{ \ie } \varphi\in \Ob
\Hhom(J,\Ind(C))
$$
be a functor, where $C$ is a $𝒰$-category. Recall \hyperref[I.sec8.5.4]{(8.5.4)}
that an \emph{indexed representation of} $\varphi$ is, by
definition, an ind-object of $\Hhom(J,C)$
$$
\psi:I\longrightarrow \Hhom(J,C){\quoi},
$$
with I an essentially small filtered category, together with a specified
isomorphism from its inductive limit in $\Hhom(J,\Ind(C))$ to $\varphi$.
Thus $\varphi$ admits an indexed
representation if and only if it is isomorphic to an essentially
small filtered inductive limit in $\Hhom(J,\Ind(C))$
of objects of the full subcategory $\Hhom(J,C)$. We shall say that
the category $J$ is \emph{admissible for} $C$ if every functor
$\varphi:J\rightarrow \Ind(C)$ admits an indexed representation.

Since in $\Ind(C)$ small filtered inductive limits are representable
(\hyperref[I.prop8.5.1]{8.5.1}), the same is true in
$\Hhom(J,\Ind(C))$, and they are computed “argument by argument”.
Consequently, the inclusion functor
\hypertarget{sga4.exp.I.equation.8.8.1.1}{}$$
\displaylines{\rlap{$(8.8.1.1)$}\hfill \Hhom(J,C)\hookrightarrow
\Hhom(J,\Ind(C))\hfill}
$$
extends canonically to a functor
\hypertarget{sga4.exp.I.equation.8.8.1.2}{}$$
\displaylines{\rlap{$(8.8.1.2)$}\hfill
\Ind(\Hhom(J,C))\longrightarrow \Hhom(J,\Ind(C))\hfill}
$$
\hyperref[I.sec8.7.2]{(8.7.2)}. The objects of the essential image of this functor are
       then precisely the $\varphi$ that admit an
       indexed representation; thus, to say that $J$ is admissible
       for $C$ also means that the functor \hyperlink{sga4.exp.I.equation.8.8.1.2}{(8.8.1.2)} is
       essentially surjective. In this connection, we note:

\begin{enonce*}{Proposition 8.8.2}\etiquette[8.8.2]{I.prop8.8.2}
If\pageoriginale the category $J$ is equivalent to a
finite category,
the canonical functor $\hyperlink{sga4.exp.I.equation.8.8.1.2}{(8.8.1.2)}$ is fully faithful; thus
$J$ is admissible for $C$ if and only if this functor is an equivalence
of categories.
\end{enonce*}

By \hyperref[I.prop8.7.5]{8.7.5} a), it suffices to prove that for
every small filtered inductive system $(\varphi_i)_{i\in I}$ in
$\Hhom(J,C)$ and every object $\varphi$ of $\Hhom(C,J)$, the
canonical map
\hypertarget{sga4.exp.I.equation.8.8.2.1}{}$$
\displaylines{\rlap{$(8.8.2.1)$}\hfill \varinjlim_i \Hom(\varphi,
\varphi_i)\longrightarrow \Hom(\varphi, \varinjlim_i
\varphi_i)\hfill}
$$
is bijective, where $\varinjlim_i \varphi_i$ denotes the inductive
limit taken in $\Hom(J,\Ind(C))$. Now, if $\varphi$, $\psi$
are two functors $J\rightarrow C$, there is an exact diagram
of sets, functorial in $\varphi$ and $\psi$:
$$
\Hom(\varphi,\psi)\longrightarrow \prod_{X\in \Ob J}
\Hom(\varphi(X),\psi(X))\rightrightarrows \prod_{\substack{f\in
F\ell J\\ f:X\rightarrow Y}} \Hom(\varphi(X), \psi (Y)){\quoi}.
$$
Since filtered inductive limits commute with kernels of
double arrows and with finite products, the bijectivity of
\hyperlink{sga4.exp.I.equation.8.8.2.1}{(8.8.2.1)}
follows when $J$ is finite. The case in which $J$ is equivalent to a
finite category reduces immediately to the preceding case.

\begin{enonce*}{Proposition 8.8.3}\etiquette[8.8.3]{I.prop8.8.3}
Let $\varphi:J\rightarrow \Ind(C)$ be a functor, with $J$
equivalent to a small category. For $\varphi$ to admit an
indexed representation, the following two conditions are sufficient,
and they are also necessary when $J$
is equivalent to a finite category:
\begin{enumerate}
\renewcommand{\theenumi}{\alph{enumi}}
\renewcommand{\labelenumi}{\theenumi)}
\item The category $\Hhom(J,C)_{/\varphi}$ (formed by the
arrows in $\Hhom(J,\Ind(C))$ with target $\varphi$ and source in
$\Hom(J,C)$) is filtered.

\item For\pageoriginale every object $j$ of $J$, the functor
\hypertarget{sga4.exp.I.equation.8.8.3.1}{}$$
\displaylines{\rlap{$(8.8.3.1)$}\hfill \psi\longmapsto
\psi(j):\Hhom(J,C)_{/\varphi}\longrightarrow C_{/\varphi(j)}\hfill}
$$
\end{enumerate}
is cofinal, \ie satisfies conditions {\rm %%
F 1), F 2)} of $\hyperref[I.prop8.1.3]{8.1.3}$.
\end{enonce*}

\noindent \emph{Proof}. Suppose a) and b) hold;
let us prove that $\varphi$ admits an indexed representation. We may
clearly suppose that $J$ is small. Let
\hypertarget{sga4.exp.I.equation.8.8.3.2}{}$$
\displaylines{\rlap{$(8.8.3.2)$}\hfill
I=\Hhom(J,C)_{/\varphi}{\quoi},\hfill}
$$
which is a filtered category by hypothesis, and first suppose that
$I$ is essentially small, so that the “inclusion” of $I$
in $\Hhom(J,C)$ defines an ind-object of $\Hhom(J,C)$, $i\mapsto
\psi_i$. Moreover, there is a canonical homomorphism
\hypertarget{sga4.exp.I.equation.8.8.3.3}{}$$
\displaylines{\rlap{$(8.8.3.3)$}\hfill \varinjlim_i
\psi_i\longrightarrow \varphi{\quoi},\hfill}
$$
given argument by argument by the homomorphism
$$
\varinjlim_I \psi_i (j)\longrightarrow
\varphi(j)=\varinjlim_{C_{/\varphi(j)}}X
$$
induced by the functor \hyperlink{sga4.exp.I.equation.8.8.3.1}{(8.8.3.1)}. Since %%
the latter is cofinal, we conclude that \hyperlink{sga4.exp.I.equation.8.8.3.3}{(8.8.3.3)} is an
isomorphism, which proves the assertion. In the case where we do not
suppose $I$ essentially small, it is enough to construct an
essentially small full subcategory $I'$ such that
\hyperlink{sga4.exp.I.equation.8.8.3.3}{(8.8.3.3)} remains an isomorphism when $\varinjlim_{I'}$ is used in
place of $\varinjlim_I$. For this purpose, it is enough that the composite functors
$$
I'\longrightarrow I\longrightarrow C_{/\varphi(j)}
$$
induced\pageoriginale by the functors \hyperlink{sga4.exp.I.equation.8.8.3.1}{(8.8.3.1)} all be
cofinal. We then conclude by the following lemma, whose
proof is immediate from criterion \hyperref[I.prop8.1.3]{(8.1.3 b)}) and is
left to the reader:

\begin{enonce*}{Lemma 8.8.3.4}\etiquette[8.8.3.4]{I.lem8.8.3.4}
Let $I$ be a filtered $𝒰$-category, and let
$$
f_j:I\longrightarrow I_j\quad (j\in J)
$$
be a small family of cofinal functors from $I$ to essentially small
filtered categories. Then there exists a
full subcategory $I'$ of $I$ which is filtered and
essentially small, and such that the functors induced by the
$f_j$ are cofinal.
\end{enonce*}

Finally, let us prove necessity in \hyperref[I.prop8.8.3]{8.8.3} when
$J$ is assumed equivalent to a finite category. An indexed representation
of $\varphi$ by means of an essentially small filtered indexing
category $I$ defines a functor $\psi$
\hypertarget{sga4.exp.I.equation.8.8.3.5}{}\[
\xymatrix@C=1.25cm@R=2cm{
 I \ar[rr]^(.4){\psi} \ar@{-->}[rrd]^{\psi_j}&&  \Hhom(J,C)_{/\varphi} \ar@{-->}[d] &\\
                        &     &   C_{/\varphi(j)}\quad ,\\
}\tag{8.8.3.5}
\]
and it is enough to prove that $\psi$ and each of the functors $\psi_j$
induced by it are cofinal: by \hyperref[I.prop8.1.3]{8.1.3} b), it
will then indeed follow that $\Hhom(J,C)_{/\varphi}$ is filtered, and by
definition \hyperref[I.def8.1.1]{8.1.1} that the vertical arrow in \hyperlink{sga4.exp.I.equation.8.8.3.5}{(8.8.3.5)} is
also cofinal. By \hyperref[I.prop8.8.2]{8.8.2}, we may identify $\varphi$ with
an object of $\Ind(E)$, where $E=\Hhom(J,C)$, and the fact that the functor
$$
\psi:I\longrightarrow E_{/\varphi}{\quoi},
$$
arising\pageoriginale from the ind-object $\varphi$ of $E$ indexed by $I$
is cofinal is a general fact, which follows immediately
from criteria (F 1) and (F 2) of \hyperref[I.prop8.1.3]{8.1.3}. The same
argument (with $E$, $\varphi$ replaced by $C$, $\varphi(j)$)
shows that $\psi_j$ is cofinal, which completes the proof.

\begin{enonce*}[remark]{Remark 8.8.4}\etiquette[8.8.4]{I.rem8.8.4}
\begin{enumerate}
\renewcommand{\theenumi}{\alph{enumi}}
\renewcommand{\labelenumi}{\theenumi)}
\item In the case where $J$ is equivalent to a finite
category, if $\varphi$ admits an indexed representation, we have seen that
\hyperlink{sga4.exp.I.equation.8.8.3.5}{(8.8.3.5)} is a cofinal functor, which implies that the category
$\Hhom(J,C)_{/\varphi}$ is itself essentially small.

\item Suppose that in $C$ finite inductive limits are representable.
Then the same is true in $E=\Hhom(J,C)$, and
these limits are computed argument by argument and are also
inductive limits in $\Hhom(J,\widehat{C})$ and \textit{a
fortiori} in $\Hhom(J,\Ind(C))$. It follows immediately that
condition a) of \hyperref[I.prop8.8.3]{8.8.3} is then automatically
satisfied, and it remains only to consider condition b). I do not know whether
it is automatically satisfied when $J$ is in addition assumed
small.
\end{enumerate}
\end{enonce*}
We now come to the main result of this section:

\begin{enonce*}{Proposition 8.8.5}\etiquette[8.8.5]{I.prop8.8.5}
Let $J$ be a category equivalent to a finite category, and
suppose in addition that $J$ is rigid, \ie that for every $j\in \Ob J$, every
endomorphism of $j$ is the identity. Then $J$ is admissible for
$C$ (for any $𝒰$-category $C$), \ie every
functor $J\rightarrow \Ind(C)$ admits an indexed representation.
\end{enonce*}

Upon \pageoriginale %%
replacing $J$ by an equivalent category if necessary, we may suppose
that $J$ is \emph{skeletal}, \ie that two isomorphic objects of $J$
are equal. Then $J$ is in fact finite. We proceed by
induction on $\text{card}\Ob J$; the case where this number is $\leq 0$
is trivial, so we shall suppose it is $\geq 1$. Let $j_\circ$ be a
\emph{maximal} object of $J$, \ie such that for every arrow
$j_\circ\rightarrow j$, there exists an arrow $j\rightarrow j_\circ$;
in view of the fact that $J$ is rigid, this implies that
$j_\circ\rightarrow j$ is an isomorphism, and since $J$ is
skeletal, it implies that $j_\circ=j$. Let $J'$ therefore be the full
subcategory obtained from $J$ by removing the object $j_\circ$, and let $J''$ be the
full subcategory consisting of the object $j_\circ$. The proposition
then follows from the

\begin{enonce*}{Lemma 8.8.5.1}\etiquette[8.8.5.1]{I.lem8.8.5.1}
Let $J$ be a finite category, and let $J'$ and $J''$ be two full
subcategories such that $\Ob J=\Ob J'\cup \Ob J''$, and such that for every
$j'\in \Ob J'$ and $J''\in \Ob J''$, one has
$\Hom(j'',j')=∅$. If $J'$ and $J''$ are admissible for
$C$, then so is $J$.
\end{enonce*}

It is enough to prove criteria a) and b) of \hyperref[I.prop8.8.3]{8.8.3},
which amounts to carrying out six elementary verifications, namely the
last two filteredness conditions for
$\Hhom(J,C)_{/\varphi}$, and conditions
F 1) and
F 2) for the functors \hyperlink{sga4.exp.I.equation.8.8.3.1}{(8.8.3.1)}, successively in the cases $j\in \Ob J'$ and
$j\in \Ob J''$; the latter conditions moreover imply
that $\Hhom(J,C)_{/\varphi}$ is nonempty. The rather
tedious verification presents no difficulty and is left to the
reader. (The editor sought in vain for an elegant proof
that would circumvent these unpleasant
verifications.)

\begin{enonce*}[remark]{Exercise 8.8.6}\etiquette[8.8.6]{I.exer8.8.6}
\begin{enumerate}
\renewcommand{\theenumi}{\alph{enumi}}
\renewcommand{\labelenumi}{\theenumi)}
\item Let\pageoriginale $J$, $J'$ be two categories admissible
for $C$, one of them being finite. Prove that $J\times J'$ is
admissible for $C$. (Use \hyperref[I.prop8.8.2]{8.8.2}.)

\item Prove that a small discrete category is admissible for
every category $C$.

\item Let $J$ be a category satisfying the following conditions:
(i) $J$ is rigid, (ii) $\Ob J$ is countable, (iii) For any two
objects $j$, $j'$ of $J$, $\Hom(j,j')$ is finite, (iv)
For every object $j$ of $J$, the set of objects of $J$ that are
bounded above by $j$, \ie that are sources of an arrow with target $j$, is
finite. Show that $J$ is admissible for every
$𝒰$-category $C$. (First show that $J$ can be written
as a filtered union of finite full subcategories $J_n$,
such that every object of $J$ bounded above by an object of $J_n$ belongs to
$J_n$. Then verify criteria a), b) of \hyperref[I.prop8.8.3]{8.8.3}.)
\end{enumerate}
\end{enonce*}

\begin{enonce*}[remark]{Exercise 8.8.7}\etiquette[8.8.7]{I.exer8.8.7}
In our notation we identify an ordered set with the
category it defines.
\begin{enumerate}
\renewcommand{\theenumi}{\alph{enumi}}
\renewcommand{\labelenumi}{\theenumi)}
\item Let $C$ be an ordered set. Let $C^{\widetilde{\phantom{,}}}$
be the set of subsets $A$ of $C$ that are filtered and such that
$x\in A$ and $y\leq x$ imply $y\in A$. Let $L_\circ:C\rightarrow
C^{\widetilde{\phantom{,}}}$ be the map that associates to every $x\in
C$ the set $L_\circ(x)$ consisting of those $y\in A$ such that $y\leq x$. Show
that $L$ is injective and that the order on $C$ is
induced by that on $C^{\widetilde{\phantom{,}}}$. For every $A\in
C^{\widetilde{\phantom{,}}}$, consider the inclusion
$f(A):A\rightarrow C$, %%
which is an ind-object of $C$; for every ind-object $\varphi:I\rightarrow
C$ of $C$, let $g(\varphi)\in C^{\widetilde{\phantom{,}}}$ be the
subset of $C$ consisting of the elements of $C$ bounded above by an element of
the form $\varphi(i)$. Show\pageoriginale that one thus obtains
two mutually quasi-inverse equivalences
\[
\xymatrix@C=1.5cm{ \Ind(C)
\ar@{<->}[r]^(.5){\displaystyle\mathop{f,g}_{\approx}}&
C^{\widetilde{\phantom{a}}}, }
\]
transforming the functor $L_\circ$ into the canonical functor
$L:C\rightarrow \Ind(C)$.

\item Suppose that for every $x\in C$, the set of elements
bounding $x$ from above (\resp from below) is finite. Show that every ind-object
(\resp pro-object) of $C$ is essentially constant, and indeed that
for every such $\varphi:I\rightarrow C$ (\resp
$\varphi:I^\circ\rightarrow C$) there exists an $i_\circ\in \Ob I$ such
that the functor induced on $I_{/i_\circ}$ is constant (\ie
takes every arrow to an isomorphism).

\item Let $C\subset\NN^\circ\times
\NN$ be formed by the pairs of natural numbers $(j,i)$ such that $i\geq
j$. Show that $\NN^{\widetilde{\phantom{,}}}$ is isomorphic to
the ordered set obtained from $\NN$ (identified with an ordered subset
of $\NN^{\widetilde{\phantom{,}}}$ as in a)) by
adjoining a greatest element $\infty$. Show that
$C^{\widetilde{\phantom{,}}}$ is isomorphic to $\NN^\circ\times
\NN^{\widetilde{\phantom{,}}}$. Consider the functor
$$
\varphi:\NN^\circ\longrightarrow C^{\widetilde{\phantom{,}}}\quav
\varphi(j)=(j,\infty){\quoi}.
$$
\end{enumerate}
\end{enonce*}

\noindent Show that:
\begin{enumerate}
\renewcommand{\labelenumi}{\theenumi)}
\item the projective limit of $\varphi$ is not representable in
$C^{\widetilde{\phantom{,}}}$, although filtered projective limits
in $C$ are representable (and are essentially
constant, by b)).

\item For every functor $\psi:\NN^\circ\rightarrow
C$, one has $\Hom(\psi,\varphi)=∅$, and \textit{a fortiori} the
functor $\varphi$ admits no indexed
representation.
\end{enumerate}

\begin{enonce*}[remark]{Exercise 8.8.8}\etiquette[8.8.8]{I.exer8.8.8}
Let\pageoriginale $X=\NN\amalg \NN$ be the disjoint union of two
copies of $\NN$, let $s$ be the symmetry of $X$, and for every $i\in \NN$,
let $X_i$ be the subset $\Delta_{i+1}\amalg\Delta_i$ of $\NN$ (where
$\Delta_i$ denotes the subset of $\NN$ formed by the $j$ such that $j\leq
i$). Let $C$ be the subcategory of the category of subsets
of $X$ whose objects are the $X_i$ and whose arrows are the
maps between the $X_i$ induced by the identity of $X$
or by its symmetry $s$, and let $C^{\widetilde{\phantom{,}}}$ be the
category defined analogously, but with the additional
object $X$.

\begin{enumerate}
\renewcommand{\theenumi}{\alph{enumi}}
\renewcommand{\labelenumi}{\theenumi)}
\item Show that $\Ind(C)$ is equivalent to
$C^{\widetilde{\phantom{,}}}$, with the canonical functor $C\rightarrow
\Ind(C)$ corresponding to the inclusion $C\rightarrow
C^{\widetilde{\phantom{,}}}$. Show that there exists no pair
$(Y,f)$ consisting of an object $Y$ of $C$ and an endomorphism $f$ of $Y$, together
with a morphism of objects with endomorphism in $\Ind(C)$ from $(Y,f)$ to
$(X,s)$.

\item Deduce that if $J$ is the category having a single object
and, in addition to the identity arrow, a single arrow of order 2, then
the functor $J\rightarrow \Ind(C)$ defined by $(X,s)$ admits no
indexed representation.
\end{enumerate}
\end{enonce*}
 

% END INCLUDED SGA4_English_current_buildable_source_20260906/source/01/drafts/SGA4_I_LINES_4271_4625_English.texfrag


% BEGIN INCLUDED SGA4_English_current_buildable_source_20260906/source/01/drafts/SGA4_I_LINES_4626_4853_English.texfrag
\subsection{Exactness properties of $\Ind(C)$}\etiquette[8.9]{I.sec8.9}

\begin{enonce*}{Proposition 8.9.1}\etiquette[8.9.1]{I.prop8.9.1}
Let $C$ be a $𝒰$-category.
\begin{enumerate}
\renewcommand{\theenumi}{\alph{enumi}}
\renewcommand{\labelenumi}{\theenumi)}
\item The canonical functors $L$ \hyperlink{sga4.exp.I.equation.8.2.4.8}{(8.2.4.8)} and $c$ \hyperlink{sga4.exp.I.equation.8.4.1}{(8.4.1)}
$$
C\xrightarrow{c}\Ind(C)\xrightarrow{L}\widehat{C}
$$
commute with projective limits.

\item Suppose that finite inductive limits are representable in $C$
and that $C$ is equivalent to a small
category.
Then small projective limits are representable in
$\Ind(C)$.

\item If\pageoriginale small projective limits
(\resp finite projective limits) are representable in $C$, then the
same is true in $\Ind(C)$. Let $J$ be a category that is finite
and rigid, or discrete; if projective limits of type $J$ are
representable in $C$, projective limits of the same type are
representable in $\Ind(C)$.

\item Small filtered inductive limits in $\Ind(C)$ (which
are representable by virtue of \hyperref[I.prop8.5.1]{8.5.1}) are left
exact, \ie commute with finite projective limits.
\end{enumerate}
\end{enonce*}

\noindent \emph{Proof}.
\begin{enumerate}
\renewcommand{\theenumi}{\alph{enumi}}
\renewcommand{\labelenumi}{\theenumi)}
\item The fact that $L$ commutes with projective limits follows from
the fact that it is fully faithful and from the computation of projective
limits in $\widehat{C}$ ``argument by argument,'' which
implies that, to verify that a projective system of morphisms
$F\rightarrow F_i$ $(i\in I)$ in $\widehat{C}$ exhibits $F$ as a
projective limit of the $F_i$, it suffices to verify that the analogous
assertion holds for the %%
projective systems of maps of sets
$\Hom(X,F)\rightarrow \Hom(X,F_i)$, for every $X\in \Ob C$. But
$C\subset \Ind(C)$.

The fact that $c$ commutes with projective limits follows formally
from the fact that $L$ does so, as does the composite $L\circ
c:C\rightarrow \widehat{C}$.

\item In view of a), the assertion amounts to saying that every
projective limit of ind-representable presheaves is
ind-representable, which follows at once from the criterion
\hyperref[I.them8.3.3]{8.3.3} (v),
since a projective limit of left exact functors
is left exact (which follows from the fact that ``projective
limits commute with projective limits'' (\hyperref[I.sec2.5.0]{2.5.0})).

\item This\pageoriginale follows formally from the analogous
  property of
$\widehat{C}$ (\hyperref[I.cor3.3]{3.3}), and from the fact that $L$ is conservative (being
fully faithful) and commutes with limits of the type under consideration (a) and
\hyperref[I.prop8.5.1]{8.5.1}).

\item It is well known (\cf \hyperref[I.prop2.3]{2.3}) that the
  representability of
finite projective limits is equivalent to that of projective limits
of the types

\begin{center}
empty $, \ldots ,\quad \vcenter{ \xymatrix@C=.25cm@R=.5cm{
                     \bigdot \ar[rd]&       &\bigdot\ar[ld]\\
                             &\bigdot&
}}$
\end{center}

\noindent (corresponding to certain finite ordered sets),
and the representability of small projective limits amounts to
that of products and finite projective limits. Thus the
second assertion made in (c) implies the first.
\end{enumerate}

\emph{Suppose first that $I$ is finite.} Using the result
\hyperref[I.prop8.8.5]{8.8.5} on the representability of %%
functors $\varphi:J\rightarrow \Ind(C)$ in indexed form
(\hyperref[I.sec8.5.4]{8.5.4})
\hypertarget{sga4.exp.I.equation.8.9.1.1}{}$$
\displaylines{\rlap{$(8.9.1.1)$}\hfill \phi:J\times I\longrightarrow
C{\quoi},\hfill}
$$
the existence of $\varprojlim \varphi$ is therefore a special case of the
following more precise and more general result:

\begin{enonce*}{Corollary \sisi{8.8.2}{8.9.2}}\etiquette[8.9.2]{I.cor8.9.2}
Consider a functor $\varphi:J\rightarrow \Ind(C)$ given in
indexed form $\Phi$ \hyperlink{sga4.exp.I.equation.8.9.1.1}{(8.9.1.1)}, where $J$ is a finite
category.
If for every $i\in \Ob I$, the partial functor $j\rightarrow
\Phi(j,i)$ has a projective (\resp inductive) limit representable
in $C$, then $\varphi$ has a projective (\resp inductive) limit
representable in $\Ind(C)$, and there is a canonical isomorphism
\hypertarget{sga4.exp.I.equation.8.9.2.1}{}$$
\displaylines{\rlap{$(8.9.2.1)$}\hfill
\varprojlim\varphi\simeq\textrm{« }{\varinjlim_i}\textrm{ »} \varprojlim_j
\Phi(j,i)\hfill}
$$
(\resp\pageoriginale
\hypertarget{sga4.exp.I.equation.8.9.2.2}{}$$
\displaylines{\rlap{$(8.9.2.2)$} \hfill \varinjlim \varphi \simeq
\textrm{« }{\varinjlim_i}\textrm{ »} \varprojlim_j\Phi(j,i)){\quoi},\hfill}
$$
where $\varprojlim_j$ (\resp. $\varinjlim_j$) is computed in $C$.
\end{enonce*}

For the first formula, one simply uses the fact that in $\Ind(C)$
filtered inductive limits commute with $\varprojlim_j$, and
that $c$ commutes with $\varprojlim_J$ \hyperref[I.prop8.9.1]{(8.9.1 d)} and a)):
$$
\varprojlim \varphi=\varprojlim_j \textrm{« }{\varinjlim_i}\textrm{ »\,} \Phi
(j,i) \xleftarrow{\sim} \textrm{« }{\varinjlim_i}\textrm{ »\,}
{\varprojlim_j}^{\Ind(C)} \Phi(j,i)\simeq \textrm{« }{\varinjlim_i}\textrm{ »}%%
\varprojlim_j{^C \ \Phi(j,i){\quoi}}.
$$
The second is proved in the same way, using the commutation of
inductive limits of type $I$ with inductive limits of type $J$
(\hyperref[I.sec2.5.0]{2.5.0}) and the fact that $c$ commutes with inductive limits of type
$J$, proved in \hyperref[I.rem8.9.4]{8.9.4} a) below.

\emph{The case where $J$ is discrete}. This case is contained in the following more
precise assertion:

\begin{enonce*}{Corollary 8.9.3}\etiquette[8.9.3]{I.cor8.9.3}
Let the $I_\alpha$ be essentially small filtered categories
indexed by a small set $A$, and for every $\alpha \in A$,
let $X(\alpha)=(X(\alpha)_{i_\alpha})_{i_\alpha \in I_\alpha}$ be an
ind-object of $C$ indexed by $I_\alpha$. Let $I$ be the product category
of the $I_\alpha$, and suppose that for every
$\sheaf{i}=(i_\alpha)_{\alpha\in A}\in I$, the product
$\prod_{\alpha\in A}%%
X (\alpha)_{i_\alpha}$ is representable in $C$. Then the product
$\prod_{\alpha \in A} %%
X (\alpha)$ is representable in $\Ind(C)$, and it is
canonically isomorphic to the ind-object indexed by $I$ given by the
formula
\hypertarget{sga4.exp.I.equation.8.9.3.1}{}$$
\displaylines{\rlap{$(8.9.3.1)$}\hfill \prod_{\alpha \in A}\;
\mathop{\textrm{« }\varinjlim\textrm{ »}}_{i_\alpha\in I_\alpha}\;
X(\alpha)_{i_\alpha}\simeq
\mathop{\textrm{« }\varinjlim\textrm{ »}}_{\sheaf{i}=(i_\alpha)_{\alpha\in
A}\in I}\; \left(\prod_{\alpha\in
A}A(\alpha)_{i_\alpha}\right){\quoi},\hfill}
$$
where\pageoriginale the product on the right-hand side is the product
computed in $C$. (Note that $I$ is filtered and essentially
small, since the $I_\alpha$ are.)
\end{enonce*}

Indeed, it is well known (and immediate by reduction to the case in which one
works in the category of sets) that the formula in question
holds when the limits are computed in $\widehat{C}$. The
conclusion then follows from the fact that $L$ commutes with the limits
in question (\hyperref[I.prop8.9.1]{8.9.1} a) and \hyperref[I.prop8.5.1]{8.5.1}).

\begin{enonce*}[remark]{Remarks 8.9.4}\etiquette[8.9.4]{I.rem8.9.4}
The proof given for \hyperref[I.cor8.9.2]{8.9.2} and \hyperref[I.cor8.9.3]{8.9.3} shows, more
generally,
that if for every $i\in \Ob I$, $\varprojlim_j \Phi(j,i)$ (\resp
$\varinjlim_j \Phi(j,i)$) \emph{computed in} $\Ind(C)$ is
representable, then the same is true of $\varprojlim_i\varphi$
(\resp of $\varinjlim_j\varphi$). This, together with the argument of c), shows
that for a given category %%
$J$ arising from a finite or discrete ordered set (or more
generally, one that is $C$-admissible (\hyperref[I.prop8.3.1]{8.3.1})), projective
limits (\resp inductive limits) of type $J$ are
representable in $\Ind(C)$ if (and only if) for every
functor $\varphi:J\rightarrow C$, the projective (\resp
inductive) limit of $\varphi$ \emph{computed in} $\Ind(C)$ is
representable. In the nonparenthetical case, this also means, by virtue
of (a), that every projective limit of type $J$ of representable
presheaves on $C$ is ind-representable.
\end{enonce*}

\begin{enonce*}{Proposition 8.9.5}\etiquette[8.9.5]{I.prop8.9.5}
Let $C$ be a $𝒰$-category.
\begin{enumerate}
\renewcommand{\theenumi}{\alph{enumi}}
\renewcommand{\labelenumi}{\theenumi)}
\item The canonical functor
$$
c:C\longrightarrow \Ind(C)
$$
is right exact (hence exact, in view of $\hyperref[I.prop8.9.1]{8.9.1}$ {\rm a)}).

\item If finite inductive limits (\resp finite sums) are
representable in $C$, then small inductive limits (\resp
small sums) are representable in $\Ind(C)$. Let $J$ be a
preordered\pageoriginale set that is finite or discrete; if inductive limits
of type $J$ are representable in $C$, the same is
true in $\Ind(C)$.
\end{enumerate}
\end{enonce*}

\noindent \emph{Proof.}
\begin{enumerate}
\renewcommand{\theenumi}{\alph{enumi}}
\renewcommand{\labelenumi}{\theenumi)}
\item Suppose that in $C$ one has $X\simeq \varinjlim_J X_j$, where
$J$ is a finite category and $X$ and the $X_j$ belong to $C$. Then,
for every ind-object $𝒴=(Y_i)_{i\in I}$ of $C$, one has:
\begin{gather*}
\Hom(X,𝒴)\simeq \varinjlim_i \Hom(X,Y_i)\simeq \varinjlim_i
\varprojlim_j \Hom(X_j,Y_i)\\
\mathrel{\mathop\simeq^{\hyperref[I.prop2.8]{2.8}}}\varprojlim_j \varinjlim_i \Hom(X_j,
Y_i) \simeq \varprojlim_j \Hom (X_j,𝒴){\quoi},
\end{gather*}
and the desired conclusion follows by comparing the outer
terms.

\item We already noted in \hyperref[I.rem8.8.4]{8.8.4} that the
  assertions made
follow from a) and \hyperref[I.cor8.9.2]{8.9.2}, at least for finite inductive
limits. To prove the stated conclusions in the infinite
cases, we are reduced to the case of sums, which can be interpreted
as a filtered inductive limit of finite sums, and we therefore conclude
by \hyperref[I.prop8.5.1]{8.5.1}.
\end{enumerate}

% END INCLUDED SGA4_English_current_buildable_source_20260906/source/01/drafts/SGA4_I_LINES_4626_4853_English.texfrag


% BEGIN INCLUDED SGA4_English_current_buildable_source_20260906/source/01/drafts/SGA4_I_LINES_4854_5078_English.texfrag
\begin{enonce*}[remark]{Remark 8.9.6}\etiquette[8.9.6]{I.rem8.9.6}
We have already observed that the functor $c$ does not in general commute with
filtered inductive limits, and hence does not commute with infinite sums
either, in contrast to what happens for $L:\Ind(C)\rightarrow
\widehat{C}$. On the other hand, apart from its commutation with
\emph{filtered} inductive limits (\hyperref[I.prop8.5.1]{8.5.1}), $L$ has
practically no commutation properties with any other type of inductive
limit (initial object, sum of two objects, cokernels of double arrows).
Thus, if $C$ admits an initial object $∅_C$, which is therefore an initial
object of $\Ind(C)$ by a), $∅_C$ is never an initial object of
$\widehat{C}$ (\ie identical to the constant presheaf with value $∅$),
since $\Hom(∅_C,∅_C)\neq ∅$!
\end{enonce*}

\begin{enonce*}{Proposition 8.9.7}\etiquette[8.9.7]{I.prop8.9.7}
Let\pageoriginale
$$
f:C\longrightarrow C'
$$
be a functor between $𝒰$-categories, and let $J$ be a finite rigid
category $\hyperref[I.prop8.8.5]{(8.8.5)}$. If inductive (\resp projective) limits of type $J$
are representable in $C$ and $f$ commutes with them, then
$\Ind(f):\Ind(C)\rightarrow \Ind(C')$ also commutes with limits of this
type.
\end{enonce*}

This follows at once from \hyperref[I.prop8.8.5]{8.8.5} and from the
calculation \hyperref[I.cor8.9.2]{8.9.2} of finite limits in a category of
ind-objects, for a functor represented in indexed form.

\begin{enonce*}{Corollary 8.9.8}\etiquette[8.9.8]{I.cor8.9.8}
If finite inductive (\resp projective) limits are representable in $C$,
and if $f$ is right exact (\resp left exact), then the same is true of
$\Ind(f)$; in the unparenthesized case, $\Ind(f)$ even commutes with all small
inductive limits.
\end{enonce*}

The first assertion follows from \hyperref[I.prop8.9.7]{8.9.7}. The second
follows from the first, taking into account that $\Ind(f)$ commutes with
filtered inductive limits (\hyperref[I.prop8.6.3]{8.6.3}).

\begin{enonce*}[remark]{Exercise 8.9.9}\etiquette[8.9.9]{I.exer8.9.9}
Let $C$ be a $𝒰$-category.
\begin{enumerate}
\renewcommand{\theenumi}{\alph{enumi}}
\renewcommand{\labelenumi}{\theenumi)}
\item Suppose that finite sums are representable in $C$. Show that if
\emph{finite} sums in $C$ are disjoint (\resp universal) (\cf II \hyperref[II.def4.5]{4.5}),
then in $\Ind(C)$ \emph{small} sums are disjoint (\resp universal).

\item Suppose that every morphism in $C$ factors as an epimorphism
followed by a monomorphism (\resp as an effective epimorphism followed\pageoriginale by
a monomorphism, \resp as an epimorphism followed by an effective
monomorphism, \resp as an effective epimorphism followed by an effective
monomorphism). Show that $\Ind(C)$ satisfies the same property. In the
last case, it follows that every epimorphism of $\Ind(C)$ is effective,
every monomorphism of $\Ind(C)$ is effective, and every bimorphism of
$\Ind(C)$ is an isomorphism; show that in this case every epimorphism
(\resp monomorphism) of $\Ind(C)$ can be represented by an inductive
system of epimorphisms (\resp monomorphisms) of $C$. If, in addition,
every epimorphism (\resp every monomorphism) in $C$ is assumed to be
\emph{universal}, the same property holds in $\Ind(C)$.

\item Suppose that $C$ is additive (\resp abelian); then so is $\Ind(C)$.

\item Let $X=\textrm{« }\varinjlim_I\textrm{ »\,}X_i$ be an ind-object of $C$,
and let $R\rightrightarrows X$ be an equivalence relation on $X$. For
every $i\in \Ob I$, let $R_i\rightrightarrows X_i$ be the equivalence
relation on $X_i$ (regarded as an object of $\Ind(C)$) induced by $R$
via $X_i\rightarrow X$. (We assume that the fiber product
$R_i=R\times_{X_XX}X_i \times X_i$ in $(\Ind(C))$ is representable by an
object of $\Ind(C)$, as is the case if finite projective limits are
representable in $C$.) Show that, for $R$ to be effective, it suffices
that the $R_i$ be effective.

\item Let $X$ be an object of $C$, and let $R$ be an equivalence relation
on $c(X)$, \ie on $X$ regarded as an object of $\Ind(C)$. Suppose that
fiber products are representable in $C$. For $R$ to be effective, it is
necessary that $R$ be of the form
$\textrm{« }\varinjlim_i\textrm{ »\,} R_i$, where the $R_i$ are equivalence
relations on $X$ (regarded as an object of $C$), and this condition
is\pageoriginale sufficient if equivalence relations in $C$ are assumed
to be effective.

\item Suppose that every morphism in $C$ factors as an effective
epimorphism followed by an effective monomorphism, and that finite
projective limits are representable. Let $X$ be an object of $C$, and
let $R$ be a subobject of $X\times X$, regarded as an object of
$\Ind(C)$. Write $R$ in the form $\textrm{« }\varinjlim\mbox{''
}Z_i$, where $(Z_i)_{i\in I}$ is a filtered increasing family of
subobjects of $X\times X$ \emph{in} $C$ (this is possible by b)). Show
that, for $R$ to be an equivalence relation on $X$, it is necessary and
sufficient that, for every $i\in \Ob I$, there exists an index
$j\in \Ob I$ whose corresponding subobject contains $s(Z_i)$ and
$Z_i\circ Z_i$, where $s$ is the symmetry of $X\times X$.

\item Suppose that in $C$ finite projective limits and finite sums are
representable, that every equivalence relation there is effective, and
that every morphism there factors as an effective
epimorphism followed by an effective monomorphism. Show that equivalence
relations in $\Ind(C)$ are universal if and only if $C$ satisfies the
following condition:
\begin{itemize}
\item[%%
ST)] For every object $X$ of $C$ and every subobject $Z$ of $X\times
X$
in $C$, if the sequence of subobjects $Z_n$ $(n\geq 0)$ of $X\times X$
in $C$ is defined recursively by $Z_\circ=Z$,
$Z_n=\Sup(Z_n,s(Z_n),Z_n\circ Z_n)$ (where $\Sup$ is taken in the set of
subobjects of $X\times X$, which exists by the hypotheses made on C),
then the sequence $(Z_n)_{n\geq 0}$ is \emph{stationary.}
\end{itemize}

\setcounter{enumi}{10}
\item Show that equivalence relations in $\Ind((\Ens))$ are not
necessarily effective.

$^\ast$ NB. For criteria ensuring that $\Ind(C)$ is a topos, \cf VI
8.9.9.${}_\ast$
\end{enumerate}
\end{enonce*}

\begin{enonce*}[remark]{Exercise 8.9.10}\etiquette[8.9.10]{I.exer8.9.10}
Let\pageoriginale $C$ be a $𝒰$-category. Put
$\Pro(C)=(\Ind(C^\circ))^\circ$ (\cf \hyperref[I.sec8.11]{8.11}).
\begin{enumerate}
\renewcommand{\theenumi}{\alph{enumi}}
\renewcommand{\labelenumi}{\theenumi)}
\item Suppose that finite sums (\resp small sums) are representable in
$C$. Show that if they are disjoint (\cf II \hyperref[II.def4.5]{4.5}), then $\Pro(C)$
satisfies the same %%
condition.

\item Let $J$ be a small set such that sums indexed by $J$ are
representable in $C$, and hence also in $\Pro(C)$. Let
$(X(\alpha))_{\alpha\in J}$ be a family of objects of $\Pro(C)$, with
$X(\alpha)=\textrm{« }\varprojlim_{I_\alpha}\textrm{ »\,}X_{i_\alpha}$.
Show that, for the sum of the $X(\alpha)$ in $\Ind(C)$ to be universal,
it suffices that the same hold for each of the families
$\varprojlim_J X_{i_\alpha}$, where $(i_\alpha)_{a\in I}\in
\prod_{\alpha\in J}I_\alpha$. Deduce that, if in $C$ sums of type $J$
are universal, then for the same to hold in $\Pro(C)$ it is necessary
and sufficient that, for every family $(X(\alpha))_{\alpha\in J}$ as
above, with $I_\alpha=I$ for every $\alpha\in J$, the canonical
homomorphism in $\Pro(C)$
$$
\textrm{« }\varprojlim_I\textrm{ »\,}\amalg_\alpha X(\alpha)_i\longleftarrow
\textrm{« }\varprojlim_{I^J}\textrm{ »\,} \amalg_{\alpha\in
J}X(\alpha)_{i_\alpha}
$$
be an isomorphism. Deduce that in $\Pro((\Ens))$, sums of type $J$ are
universal if and only if $J$ is finite.
\end{enumerate}
\end{enonce*}

\begin{enonce*}[remark]{Exercise 8.9.11}\etiquette[8.9.11]{I.exer8.9.11}
Let $C$ be a $𝒰$-category.
\begin{enumerate}
\renewcommand{\theenumi}{\alph{enumi}}
\renewcommand{\labelenumi}{\theenumi)}
\item Let $(X_i\rightarrow X)_{i\in I}$ be a family of morphisms in
$C$. Show that, for it to be epimorphic in $\Ind(C)$, it suffices that
there exists a \emph{finite} subfamily that is epimorphic in $C$. Prove
that this condition is also necessary when it is assumed that in $C$
every finite family of morphisms with target $X_i$ factors as an epimorphic
family,\pageoriginale followed by an effective monomorphism \hyperref[I.sec10.5]{(10.5)}. (For
this last assertion, let $P$ be the set of finite subsets of $I$, ordered
by inclusion, and let $X'=(X'_J)_{J\in P}$ be the ind-object formed by the
images of the finite subfamilies of the given family; finally, let
$X''=X\amalg_{X'}X=\textrm{« }\lim\textrm{ »\,}X\amalg_{X'_J}X$. Show
that the two canonical morphisms $X\rightrightarrows X''$ are distinct,
but coincide on the $X'_J$.)

\item Suppose that every finite family of morphisms with the same target
in $C$ factors as an epimorphic family followed by an effective
monomorphism. Prove that $\Ind(C)$ admits a small subcategory generating
by epimorphisms (\hyperref[I.def7.1]{7.1}) if and only if $C$ admits a
small subcategory $C'$ such that every object of $C$ is the target of a
\emph{finite} epimorphic family with source in $C'$. If $C$ is assumed
to admit a small generating subcategory (\hyperref[I.def7.1]{7.1}), then
$\Ind(C)$ admits a small subcategory generating by epimorphisms if and
only if $C$ is equivalent to a small category. (For this last statement,
use \hyperref[I.var7.5.2]{7.5.2}).

\item $\Ind((\Ens))$ does not admit a small generating subcategory.
\end{enumerate}
\end{enonce*}

\begin{enonce*}[remark]{Exercise 8.9.12}\etiquette[8.9.12]{I.8.9.12}
Let $C$ be a $𝒰$-category, and consider
$\Pro(C)=\Ind(C^\circ)^\circ$.
\begin{enumerate}
\renewcommand{\theenumi}{\alph{enumi}}
\renewcommand{\labelenumi}{\theenumi)}
\item Show that if $\Pro(C)$ admits a small subcategory $P'$ generating
by epimorphisms \hyperref[I.def7.1]{(7.1)}, then there exists a small subcategory $C'$ of $C$
that is generating by epimorphisms in $\Pro(C)$. (Take the category of
components of the objects of $P'$.)

\item Show\pageoriginale that the category $\Pro((\Ens))$ has no small
subcategory generating by epimorphisms. (If $C'$ is as in a), choose a
set $X$ whose cardinality strictly exceeds the cardinalities of the
objects of $C'$, and consider the pro-set $𝒳$ formed by the complements
in $X$ of subsets of cardinality $<\card(X)$. Show that, for every
nonempty set $Y$ of cardinality $<\card(X)$, one has $\Hom(Y,𝒳)=∅$, but
that $𝒳$ is not isomorphic to the constant pro-set $∅$.)
\end{enumerate}
\end{enonce*}

% END INCLUDED SGA4_English_current_buildable_source_20260906/source/01/drafts/SGA4_I_LINES_4854_5078_English.texfrag


% BEGIN INCLUDED SGA4_English_current_buildable_source_20260906/source/01/drafts/SGA4_I_LINES_5079_5229_English.texfrag
\subsection{Dual notions: pro-objects, pro-representable functors}\etiquette[8.10]{I.sec8.10}

Let $C$ be a $𝒰$-category. A \emph{pro-object} of $C$ is a functor
\hypertarget{sga4.exp.I.equation.8.10.1}{}$$
\displaylines{\rlap{$(8.10.1)$}\hfill \varphi:I^\circ\longrightarrow
C{\quoi},\hfill}
$$
where $I$ is an essentially small filtered category (called the index
category), and, as usual, $I^\circ$ denotes the opposite category. Let
$𝒱$ be a universe such that $𝒱\supset 𝒰$. One should note that the
pro-objects of $C$ indexed by $I\in 𝒱$ are in one-to-one
correspondence with the ind-objects of $C^\circ$ indexed by $I\in
𝒱$: to every such ind-object $\psi:I\rightarrow C^\circ$ one
associates the pro-object $\varphi=\psi^\circ:I^\circ\rightarrow C$.
Following this correspondence, one defines ``by transport of structure
and reversal of arrows'' the notion of a morphism between pro-objects
of $C$ from the analogous notion \hyperlink{sga4.exp.I.equation.8.2.4.2}{(8.2.4.2)}, \hyperlink{sga4.exp.I.equation.8.2.5.1}{(8.2.5.1)} for ind-objects,
and consequently defines the category of pro-objects of $C$ indexed by
$I\in 𝒱$, together with a canonical isomorphism
\hypertarget{sga4.exp.I.equation.8.10.2}{}$$
\displaylines{\rlap{$(8.10.2)$}\hfill
\Pro_{𝒱}(C,𝒰)\simeq \Ind_{𝒱}(C^\circ,
𝒰)^\circ{\quoi};\hfill}
$$
for\pageoriginale $𝒱=𝒰$ one simply speaks of the \emph{category of
pro-objects of} $C$, denoted $\Pro_{𝒰}(C)$ or simply $\Pro(C)$,
which is thus defined by
\hypertarget{sga4.exp.I.equation.8.10.3}{}$$
\displaylines{\rlap{$(8.10.3)$}\hfill
\Pro(C)=\Pro_{𝒰}(C,𝒰)\simeq
\Ind(C^\circ)^\circ{\quoi}.\hfill}
$$
If two pro-objects are given in \emph{indexed form}
\hypertarget{sga4.exp.I.equation.8.10.4}{}\begin{equation}
𝒳=(X_i)_{i\in I}\quav 𝒴=(Y_j)_{j\in
J}{\quoi},\tag{8.10.4}
\end{equation}
formula \hyperlink{sga4.exp.I.equation.8.2.5.1}{(8.2.5.1)} takes here the form
\hypertarget{sga4.exp.I.equation.8.10.5}{}\begin{equation}
\Hom(𝒳,𝒴)\simeq \varprojlim_j \varinjlim_i
\Hom(X_i,Y_j){\quoi}.\tag{8.10.5}
\end{equation}
The functor \hyperlink{sga4.exp.I.equation.8.4.1}{(8.4.1)}, applied to $C^\circ$, gives, upon passing to
opposite categories, a canonical functor (denoted by the same letter
$c$ when there is no risk of confusion) that identifies $C$ with a
full subcategory of $\Pro(C)$:
\hypertarget{sga4.exp.I.equation.8.10.6}{}\begin{equation}
c:C\longrightarrow\Pro(C){\quoi};\tag{8.10.6}
\end{equation}
it is in order to have such a functor, rather than
$C\rightarrow\Pro(C)^\circ$, that we ``reversed the arrows'' in the
definition of morphisms of pro-objects starting from the analogous
definition for ind-objects. Pro-objects in the essential image of
\hyperlink{sga4.exp.I.equation.8.10.6}{(8.10.6)} are again called \emph{essentially constant pro-objects.}

Put
\hypertarget{sga4.exp.I.equation.8.10.7}{}\begin{equation}
\check{C}=(C^\circ)^{\widehat{\phantom{,}}}=\Hhom(C,(\Ens)){\quoi},\tag{8.10.7}
\end{equation}
then the functor \hyperlink{sga4.exp.I.equation.8.2.4.7}{(8.2.4.7)} may be regarded as a canonical fully faithful
functor
\hypertarget{sga4.exp.I.equation.8.10.8}{}\begin{equation}
L:\Pro_{𝒱}(C,𝒰)\longrightarrow
\check{C}^\circ{\quoi},\tag{8.10.8}
\end{equation}\pageoriginale
and in particular one obtains
\hypertarget{sga4.exp.I.equation.8.10.9}{}\begin{equation}
L:\Pro(C)\hookrightarrow \check{C}^\circ{\quoi}.\tag{8.10.9}
\end{equation}

We leave it to the reader to translate, as the need arises, the results
of the preceding sections and of those that follow from the language of
ind-objects into that of pro-objects. Depending on the mathematical
context, one language or the other is more useful, not to mention the
cases in which both notions arise simultaneously, for example when one
must consider complex categories such as $\Pro(\Ind(C))$ or
$\Ind(\Pro(C))$. We confine ourselves to giving a few additional
indications, in order to fix the terminology and notation and to make
the ``yoga'' precise.


\setcounter{subsubsection}{9}
\subsubsection{}\etiquette[8.10.10]{I.sec8.10.10}
A \emph{covariant functor $F:C\rightarrow (\Ens)$} is called a
\emph{pro-representable functor} if it belongs to the essential image
of \hyperlink{sga4.exp.I.equation.8.10.9}{(8.10.9)} (or of \hyperlink{sga4.exp.I.equation.8.10.8}{8.10.8}, which amounts to the same thing); the full
subcategory of $\check{C}$ formed by these functors is therefore
equivalent to $\Pro(C)^\circ$. In general one will prefer to work in
the opposite category, the essential image \emph{as a subcategory of}
\hyperlink{sga4.exp.I.equation.8.10.9}{(8.10.9)}, which is therefore equivalent to $\Pro(C)$ itself. In
accordance with this usage, it is often preferable to regard covariant
functors
$$
F:C\longrightarrow (\Ens)
$$
as objects of $\Hhom(C,(\Ens))^\circ=\check{C}^\circ$, and consequently
to write the category of morphisms
$$
X\longrightarrow F\mbox{ in } \check{C}^\circ{\quoi},
$$
\ie\pageoriginale of pairs $(X,u)$ consisting of an object $X$ of $C$
and an element $u\in F(X)$, as
\hypertarget{sga4.exp.I.equation.8.10.11}{}\begin{equation}
{}_{\displaystyle F}\backslash C{\quoi}  %%
\ .\tag{8.10.11}
\end{equation}
With this convention, one thus obtains a \emph{covariant} functor
(source functor)
\hypertarget{sga4.exp.I.equation.8.10.12}{}\begin{equation}
{}_{\displaystyle F}\backslash C\rightarrow C{\quoi},\tag{8.10.12}
\end{equation}
and \hyperref[I.prop3.4]{3.4} is rewritten in the form
\hypertarget{sga4.exp.I.equation.8.10.13}{}\begin{equation}
F\xrightarrow{\sim}%%
\textrm{``}\varprojlim_{({}_{\displaystyle F}\backslash C)^\circ}\textrm{''\,}
X\tag{8.10.13}
\end{equation}

\setcounter{subsubsection}{13}
\subsubsection{}\etiquette[8.10.14]{I.sec8.10.14}
Criterion \hyperref[I.them8.3.3]{8.3.3}, applied to $C^\circ$, becomes
a pro-representability criterion for $F$: it is necessary and
sufficient that $({}_F\backslash C)^\circ$ be filtered and essentially
small, the latter condition being superfluous if $C$ is equivalent to a
small category; if, moreover, finite projective limits are representable
in $C$, this is equivalent to saying that $F$ commutes with them, \ie is
\emph{left exact.}

\subsubsection{}\etiquette[8.10.15]{I.sec8.10.15}
In $\Pro(C)$ small filtered projective limits are representable and the
functor \hyperlink{sga4.exp.I.equation.8.10.9}{(8.10.9)} commutes with them; in other words, this functor carries
projective limits in $\Pro(C)$ to \emph{inductive} limits in
$\check{C}=\Hhom(C,(\Ens))$, just like its composite
$C\rightarrow\check{C}^\circ$ with \hyperlink{sga4.exp.I.equation.8.10.6}{(8.10.6)}, which shares with it the
unfortunate property of being contravariant when regarded as taking
values in $\check{C}$. One should note that pro-representable\pageoriginale
functors from $C$ to $(\Ens)$ are the functors that are small filtered
\emph{inductive} limits of representable functors (and not projective
limits, as the terminology might perhaps suggest).


% END INCLUDED SGA4_English_current_buildable_source_20260906/source/01/drafts/SGA4_I_LINES_5079_5229_English.texfrag


% BEGIN INCLUDED SGA4_English_current_buildable_source_20260906/source/01/drafts/SGA4_I_LINES_5230_5389_English.texfrag
\subsection{Ind-adjoints and pro-adjoints}\etiquette[8.11]{I.sec8.11}

\subsubsection{}\etiquette[8.11.1]{I.sec8.11.1}
Let
\hypertarget{sga4.exp.I.equation.8.11.1.1}{}\begin{equation}
f:C\longrightarrow C'\tag{8.11.1.1}
\end{equation}
be a functor between $𝒰$-categories, whence a functor
$F \sisi{→}{↦} F\circ f$
\hypertarget{sga4.exp.I.equation.8.11.1.2}{}\begin{equation}
f^\ast:\widehat{C}'\longrightarrow\widehat{C}{\quoi}.\tag{8.11.1.2}
\end{equation}
We say that $f$ \emph{admits an ind-adjoint} if the preceding functor
carries ind-representable functors to ind-representable functors. Since
$f^\ast$ commutes with inductive limits, and since the full subcategory
of $\widehat{C}'$ formed by the ind-representable functors is stable
under small filtered inductive limits, it is equivalent to say that
$f$ admits an ind-adjoint, or that $f^\ast$ carries objects of $C'$
(identified with a full subcategory of $\widehat{C}'$) to
ind-representable functors on $C$, \ie
\hypertarget{sga4.exp.I.equation.8.11.1.3}{}\begin{equation}
X\longmapsto \Hom(f(X),Y')\mbox{ is ind-representable for every }
Y'\in \Ob C'{\quoi}.\tag{8.11.1.3}
\end{equation}
Another way to express the condition that $f$ admit an ind-adjoint is
to say that there exists a functor
\hypertarget{sga4.exp.I.equation.8.11.1.4}{}\begin{equation}
g:\Ind(C')\longrightarrow \Ind(C)\tag{8.11.1.4}
\end{equation}
essentially induced by $f^\ast$, \ie such that there is an isomorphism
of bifunctors
\hypertarget{sga4.exp.I.equation.8.11.1.5}{}\begin{equation}
\Hom_{\Ind(C)}(X,g(𝒴'))\simeq
\Hom_{\Ind(C')}(f(X),𝒴'){\quoi},\tag{8.11.1.5}
\end{equation}
where\pageoriginale $X\in \Ob C$ and $𝒴'\in \Ob \Ind(C')$. In fact,
it even suffices to have a functor
\hypertarget{sga4.exp.I.equation.8.11.1.6}{}\begin{equation}
g_\circ:C'\longrightarrow \Ind(C)\tag{8.11.1.6}
\end{equation}
and an isomorphism of bifunctors
\hypertarget{sga4.exp.I.equation.8.11.1.7}{}\begin{equation}
\Hom_{\Ind(C)}(X,g_\circ(Y'))\simeq \Hom_{C'}(f(X),Y')\tag{8.11.1.7}
\end{equation}
in $X\in \Ob C$, $Y'\in \Ob C'$. The functor $g$ (respectively
$g_\circ$) is evidently determined up to unique isomorphism by $f$,
and conversely $f$ can be reconstructed up to canonical isomorphism
from the knowledge of this $g$ (respectively $g_\circ$). It is clear
from \hyperlink{sga4.exp.I.equation.8.11.1.5}{(8.11.1.5)} that the functor $g$ commutes with inductive limits and
that it ``extends'' the functor $g_\circ$; this therefore determines it
up to unique isomorphism in terms of $g_\circ$
(\hyperref[I.sec8.7.2]{8.7.2}). The functor $g$, and sometimes also the
functor $g_\circ$ that it extends, is called the \emph{ind-adjoint
functor of} $f$ (there is no need here to specify: on the right, since
the other one, if it exists, will be called the pro-adjoint; cf.
\hyperref[I.sec8.11.5]{8.11.5} below).

Of course, when $f$ admits a right adjoint
$$
f^{\ad}:C'\longrightarrow C{\quoi},
$$
it admits an ind-adjoint $g$, and the latter is canonically isomorphic
to the canonical extension of $f^{\ad}$ to ind-objects:
\hypertarget{sga4.exp.I.equation.8.11.1.8}{}\begin{equation}
g\simeq\Ind(f^{\ad}){\quoi}.\tag{8.11.1.8}
\end{equation}
The notion of an ind-adjoint is thus a natural generalization of the
notion of a right adjoint.

Now consider\pageoriginale the canonical extension
\hypertarget{sga4.exp.I.equation.8.11.1.9}{}\begin{equation}
\Ind(f):\Ind(C)\longrightarrow\Ind(C'){\quoi}.\tag{8.11.1.9}
\end{equation}
We then have:

\begin{enonce*}{Proposition 8.11.2}\etiquette[8.11.2]{I.prop8.11.2}
For the functor $f:C\rightarrow C'$ to admit an ind-adjoint, it is
necessary and sufficient that the functor $\Ind(f)$ $\hyperlink{sga4.exp.I.equation.8.11.1.9}{(8.11.1.9)}$ admit
a right adjoint. The latter is canonically isomorphic to the
ind-adjoint $\hyperlink{sga4.exp.I.equation.8.11.1.4}{(8.11.1.4)}$.
\end{enonce*}

Sufficiency and the last assertion are immediate from the
ind-adjunction formula \hyperlink{sga4.exp.I.equation.8.11.1.5}{(8.11.1.5)}. For necessity, note that by passing
to the projective limit of this formula over the $X_i$, one obtains an
adjunction isomorphism for the ind-objects $𝒳=(X_i)_{i\in I}$ and
$𝒴'$:
\hypertarget{sga4.exp.I.equation.8.11.2.1}{}\begin{equation}
\Hom_{\Ind(C)}(𝒳, g(𝒴'))\simeq
\Hom_{\Ind(C')}(\Ind(f)(𝒳),𝒴'){\quoi},\tag{8.11.2.1}
\end{equation}
\cqfd

\begin{enonce*}{Corollary 8.11.3}\etiquette[8.11.3]{I.cor8.11.3}
If $f$ admits an ind-adjoint, then $\Ind(f)$ commutes with inductive
limits, and the ind-adjoint $g$ commutes with projective limits.
\end{enonce*}

\begin{enonce*}{Proposition 8.11.4}\etiquette[8.11.4]{I.prop8.11.4}
For the functor $f:C\rightarrow C'$ to admit an ind-adjoint, it is
necessary that $f$ be \sisi{left}{right} exact, and this condition is
sufficient when $C$ is equivalent to a small category.
\end{enonce*}

This follows from criterion \hyperlink{sga4.exp.I.equation.8.11.1.3}{(8.11.1.3)}, and from
\hyperref[I.prop8.3.1]{8.3.1} and \hyperref[I.them8.3.3]{8.3.3} (iv).

For another criterion in terms of the notion of an accessible functor,
cf. \hyperref[I.cor8.13.3]{8.13.3}.

\setcounter{subsubsection}{4}\etiquette[8.11.5]{I.sec8.11.5}
\subsubsection{}
Now\pageoriginale consider the functor $F\mapsto F\circ f$
\hypertarget{sga4.exp.I.equation.8.11.5.1}{}\begin{equation}
f^{\circ\ast}:\check{C}'\longrightarrow \check{C}\tag{8.11.5.1}
\end{equation}
induced by \sisi{$F$}{$f$}. We say that $f$ \emph{admits a pro-adjoint}
if the preceding functor carries pro-representable functors to
pro-representable functors, \ie if there exists a functor (called the
\emph{pro-adjoint functor of $f$})
\hypertarget{sga4.exp.I.equation.8.11.5.2}{}\begin{equation}
g:\Pro(\sisi{C}{C'})\longrightarrow \Pro(\sisi{C'}{C}){\quoi},\tag{8.11.5.2}
\end{equation}
and an isomorphism of bifunctors
\hypertarget{sga4.exp.I.equation.8.11.5.3}{}\begin{equation}
\Hom_{\Pro(C)}(g(𝒴'),𝒳)\simeq
\Hom_{\Pro(C')}(𝒴'\sisi{)}{},f(𝒳)){\quoi}.\tag{8.11.5.3}
\end{equation}
Of course, to say that $f$ admits a pro-adjoint means that $f^\circ$
admits an ind-adjoint, so that the notions and results for ind-adjoints
translate immediately into terms of pro-adjoints. We merely point out
that $f$ admits a pro-adjoint if and only if
$\Pro(f):\Pro(C)\rightarrow\Pro(C')$ admits a \emph{\sisi{right}{left}}
adjoint, and that in this case the preceding functor $g$ is such a
\sisi{right}{left} adjoint of $\Pro(\sisi{C}{f})$; and that for this it
is necessary that $f$ be \emph{left} exact, this condition also being
sufficient when $C$ is equivalent to a small category. In this case,
$f$ is therefore exact if and only if it admits both an ind-adjoint and
a pro-adjoint.

\begin{enonce*}[remark]{Example 8.11.6}\etiquette[8.11.6]{I.exem8.11.6}
Consider the case of a functor
$$
f:C\longrightarrow (\Ens){\quoi}.
$$
If this functor admits a pro-adjoint, it is pro-representable, and the
converse is true if and only if the full subcategory of $\check{C}$
formed by the pro-representable functors is stable under small products;
this\pageoriginale is the case in particular if $C$ is equivalent to a
small category (\hyperref[I.sec8.10.14]{8.10.14}) or if small products
are representable in $C$ (\hyperref[I.prop8.9.5]{8.9.5} b), applied to
$C^\circ$). Up to minor qualifications, one may thus say that for a
functor $f:C\rightarrow C'$, the existence of a pro-adjoint is the
\emph{natural generalization of the notion of pro-representability of
$f$}, which is defined when $C'=(\Ens)$.
\end{enonce*}


% END INCLUDED SGA4_English_current_buildable_source_20260906/source/01/drafts/SGA4_I_LINES_5230_5389_English.texfrag


% BEGIN INCLUDED SGA4_English_current_buildable_source_20260906/source/01/drafts/SGA4_I_LINES_5390_5467_English.texfrag
\subsection{Strict ind-objects and pro-objects. Application to a
representability criterion}\etiquette[8.12]{I.sec8.12}

\subsubsection{}\etiquette[8.12.1]{I.sec8.12.1}
Let
$$
𝒳=(X_i)_{i\in I}
$$
be an ind-object of the $𝒰$-category $C$, and let $F$ be the presheaf
that it ind-represents. It is then immediately clear that it is
equivalent to say that the canonical morphisms
$$
X_i\longrightarrow F=\textrm{``}\varinjlim_i\textrm{''\,}X_i
$$
are monomorphisms in $\Ind(C)$ (or, equivalently, in $\widehat{C}$,
\ie monomorphisms of functors ``argument by argument''), or to say that
for every arrow $i\rightarrow j$ of $I$, the corresponding transition
arrow
$$
X_i\longrightarrow X_j
$$
is a monomorphism. When these conditions are satisfied, and if moreover
$I$ is an ordered category, we say that $𝒳$ is a \emph{strict
ind-object}. One\pageoriginale should note that this condition is not
invariant under isomorphism of ind-objects; an ind-object will be called
\emph{essentially strict} if it is isomorphic to a strict ind-object. A
presheaf will be called \emph{strictly ind-representable} if it is
ind-representable by a strict ind-object; thus
$F=\textrm{``}\varinjlim_i\textrm{''\,}X_i$ is strictly
ind-representable if and only if $𝒳=(X_i)_{i\in I}$ is essentially
strict.

\paragraph{}\etiquette[8.12.1.1]{I.8.12.1.1}
Let $F$ be a presheaf on $C$, and consider the full subcategory of
$C_{/F}$ formed by the arrows $X\rightarrow F$ with source in $C$ that
are monomorphisms. We shall call it the \emph{category of representable
subfunctors of $F$}; it is the category associated with the ordered set
of representable subfunctors of $F$, ordered by the order induced from
that of the set of subobjects of $F$. The following then follows
immediately from the definitions:

\begin{enonce*}{Proposition 8.12.2}\etiquette[8.12.2]{I.prop8.12.2}
For the presheaf $F$ on $C$ to be strictly ind-representable, it is
necessary and sufficient that the (ordered) category $I$ of
representable subfunctors of $F$ be filtered and essentially small, and
that one have
\hypertarget{sga4.exp.I.equation.8.12.2.1}{}\begin{equation*}
\varinjlim_I X_i \xrightarrow{\sim}F{\quoi}.\tag{8.12.2.1}
\end{equation*}
When, for every object $X$ of $C$, the set of subobjects of $X$ in $C$
is small (for example, if $C$ admits a small generating subcategory
$\hyperref[I.prop7.4]{(7.4)}$), this implies that the category of representable subfunctors is
itself small.
\end{enonce*}

\setcounter{subsubsection}{2} \setcounter{paragraph}{1}
\paragraph{}\etiquette[8.12.2.2]{I.sec8.12.2.2}
If\pageoriginale $F$ is strictly ind-representable, there is thus a
preferred way of ind-representing it by an ind-object, and the latter is
even a strict ind-object: one takes the representation \hyperlink{sga4.exp.I.equation.8.12.2.1}{(8.12.2.1)}. A
strict ind-object $𝒴$ is called \emph{saturated} if it is isomorphic to
an ind-object of the form $(X_i)$ occurring in \hyperlink{sga4.exp.I.equation.8.12.2.1}{(8.12.2.1)}, for a
suitable $F$; thus, up to unique isomorphism, there exists a single
saturated strict ind-object isomorphic to the given strict ind-object,
namely the one considered in \hyperref[I.prop8.12.2]{8.12.2}, taking
$F=\varinjlim 𝒴$ (the limit in $\widehat{C}$).

\paragraph{}\etiquette[8.12.2.3]{I.sec8.12.2.3}
Suppose that it is already known that one can find a small cofinal subset
in the set $\Ob C_{/F}$, as is the case in particular if $F$ is
ind-representable; it then follows that the same condition is satisfied
in the full subcategory $I$ considered in
\hyperref[I.prop8.12.2]{8.12.2}, so in criterion
\hyperref[I.prop8.12.2]{8.12.2} one may omit the condition that $I$ be
essentially small.


% END INCLUDED SGA4_English_current_buildable_source_20260906/source/01/drafts/SGA4_I_LINES_5390_5467_English.texfrag


% BEGIN INCLUDED SGA4_English_current_buildable_source_20260906/source/01/drafts/SGA4_I_LINES_5468_5594_English.texfrag
\subsubsection{}\etiquette[8.12.3]{I.sec8.12.3}
Let $F$ be a presheaf on $C$, and consider an object
$$
u:X\longrightarrow F
$$
of $C_{/F}$. We say that $u$ (or the pair $(X,u)$) is \emph{minimal}
if, for every factorization of $u$ as
\hypertarget{sga4.exp.I.equation.8.12.3.1}{}\begin{equation}
X\xrightarrow{p}X'\xrightarrow{u'}F{\quoi},\tag{8.12.3.1}
\end{equation}
with $p$ a strict epimorphism (\hyperref[I.sec10.2]{10.2}), $p$ is an
isomorphism. Regarding $u$ as an element of $F(X)$ \hyperref[I.prop1.4]{(1.4)}, to say that
$u$ is minimal thus means that\pageoriginale every strict epimorphism
$p:X\rightarrow X'$ such that $u\in
\Iim(F(p):F(X')\rightarrow F(X))$ is an isomorphism. This notion is
clarified by part a) of the following lemma:

\begin{enonce*}{Lemma 8.12.4}\etiquette[8.12.4]{I.lem8.12.4}
Let $F$ be a presheaf on $C$.
\begin{enumerate}
\renewcommand{\theenumi}{\alph{enumi}}
\renewcommand{\labelenumi}{\theenumi)}
\item Suppose that $F$ carries cokernels to kernels, and consider a
morphism
$$
u:X\longrightarrow F{\quoi},
$$
with $X\in \Ob C$. For $u$ to be a monomorphism, it is sufficient,
when cokernels of pairs of arrows are representable in $C$, that $u$ be
minimal; this condition is also necessary if fiber products are
representable in $C$.

\item Suppose that finite $\varinjlim$ are representable in $C$. For
the subcategory of representable subfunctors of $F$
(\hyperref[I.8.12.1.1]{8.12.1.1}) to be filtered (not necessarily
small) and to have $F$ as its inductive limit, it is sufficient that
$F$ be left exact and that every morphism $u:X\rightarrow F$, with
$X\in \Ob C$, factor as
$$
X\rightarrow X'\xrightarrow{u'}F{\quoi},
$$
with $u'$ minimal; this condition is also necessary if fiber products
are representable in $C$.

\item Suppose that $C$ admits a small generating subcategory
(\hyperref[I.def7.1]{7.1}), and that $I$, the category of representable
subfunctors of $F$, is filtered; then $I$ is small.
\end{enumerate}
\end{enonce*}

\noindent \emph{Proof}.\pageoriginale
\begin{enumerate}
\renewcommand{\theenumi}{\alph{enumi}}
\renewcommand{\labelenumi}{\theenumi)}
\item Suppose $u$ is minimal; let us prove that $u$ is a monomorphism,
\ie that for every pair of arrows $v,v':Y\rightrightarrows X$ such that
$uv=uv'$, one has $v=v'$. Indeed, if $X'=\Coker(v,v')$, then $u$
factors as $X\xrightarrow{p}X'\xrightarrow{u'}F$ (since $F$ carries
cokernels to kernels), and since $p$ is a strict epimorphism by
construction, it follows that $p$ is an isomorphism, \ie $v=v'$.
Suppose that $u$ is a monomorphism; let us prove that it is minimal.
Consider a factorization \hyperlink{sga4.exp.I.equation.8.12.3.1}{(8.12.3.1)}; since $p$ is a strict epimorphism,
it is the cokernel of the canonical pair of arrows
$v,v':\fprod{X}{X}{X'}\rightrightarrows X$, and since
$(u'p)v=(u'p)v'$ and $u'p=u$ is a monomorphism, one has $v=v'$, hence
$p$ is an isomorphism.

\item Sufficiency: Since $F$ is left exact, $C_{/F}$ is filtered. Since
we know that $F=\varinjlim_{C_{/F}}X$, by
\hyperref[I.prop8.1.3]{8.1.3} c) it remains to prove that the full
subcategory $I$ of $C_{/F}$ consisting of the subobjects of $F$ is
cofinal in $C_{/F}$; by the ``sufficient'' part of a), this is precisely
what is ensured by the hypothesis that every object of $C_{/F}$ is
dominated by a ``minimal'' object. Necessity: Since every filtered
inductive limit of right exact functors is again right exact, the first
condition is trivially necessary. The second then follows from the
``necessary'' part of a).

\item A representable subfunctor $X\hookrightarrow F$ of $F$ is known
once one knows the full subcategory $C'_{/X}$ of $C'_{/F}$, where $C'$
is a fixed small generating subcategory of $C$. (Use the hypothesis
that $I$ is filtered.) Since $C'_{/F}$ is small, the set of its full
subcategories is small, whence the conclusion.
\end{enumerate}

\begin{enonce*}{Proposition 8.12.5}\etiquette[8.12.5]{I.prop8.12.5}
Let\pageoriginale $C$ be a $𝒰$-category in which finite inductive
limits are representable, and which admits a small generating
subcategory (\hyperref[I.def7.1]{7.1}). Let $F$ be a presheaf on $C$.
For $F$ to be strictly ind-representable, it is sufficient that it
satisfy the following two conditions, and these conditions are also
necessary if fiber products are representable in $C$:
\begin{enumerate}
\renewcommand{\theenumi}{\alph{enumi}}
\renewcommand{\labelenumi}{\theenumi)}
\item $F$ is left exact.

\item Every pair $(X,u)$, with $X\in \Ob C$ and $u\in F(X)$, is
dominated in $C_{/F}$ by a minimal pair $\hyperref[I.sec8.12.3]{(8.12.3)}$, \ie there exists a
minimal pair $(X',u')$ $(X'\in \Ob C, u'\in F(X'))$ and a morphism
$f:X\rightarrow X'$ such that $u=F(f)(u')$.
\end{enumerate}
\end{enonce*}

Sufficiency follows from \hyperref[I.prop8.12.2]{8.12.2} and
\hyperref[I.lem8.12.4]{8.12.4} b), c), and necessity from
\hyperref[I.prop8.3.1]{8.3.1} and \hyperref[I.lem8.12.4]{8.12.4} a).

\begin{enonce*}[remark]{Remark 8.12.6}\etiquette[8.12.6]{I.rem8.12.6}
When $F$ carries amalgamated sums in $C$ to fiber products, it is
immediately clear that for a given pair $(X,u)$ as in b), the set of
strict quotients $X'$ of $X$ such that
$u\in\Iim(F(X')\rightarrow F(X))$ is filtered decreasingly; this
implies that if the set of strict quotients of $X$ is Artinian, then
condition b) of \hyperref[I.prop8.12.5]{8.12.5} is automatically
satisfied. Thus, if the preceding condition on $X$ is satisfied for
every object $X$ of $C$ (one then also says, sometimes, that the objects
of $C^\circ$ are \emph{Artinian}), it follows from \hyperref[I.prop8.12.5]{8.12.5} that $F$ is
strictly ind-representable if and only if $F$ is left exact; in this
case, $F$ is therefore strictly ind-representable as soon as it is
ind-representable (\hyperref[I.prop8.3.1]{8.3.1}).
\end{enonce*}


% END INCLUDED SGA4_English_current_buildable_source_20260906/source/01/drafts/SGA4_I_LINES_5468_5594_English.texfrag


% BEGIN INCLUDED SGA4_English_current_buildable_source_20260906/source/01/drafts/SGA4_I_LINES_5595_5699_English.texfrag
\begin{enonce*}{Corollary 8.12.7}\etiquette[8.12.7]{I.cor8.12.7}
Let\pageoriginale $C$ be a $𝒰$-category in which small projective
limits are representable, and which admits a small cogenerating
subcategory (\hyperref[I.sec7.9]{7.9},
\hyperref[I.exem7.13]{7.13}). Then a functor
$F:C\rightarrow (\Ens)$ is representable if and only if $F$ commutes
with small projective limits.
\end{enonce*}
Necessity is clear. For sufficiency, apply
\hyperref[I.prop8.12.5]{8.12.5} to the opposite category $C^\circ$. To
prove first that $F$ is (strictly) pro-representable, it remains to show
that every pair $(X,u)$, $X\in \Ob C$ and $u\in F(X)$, is dominated by
a ``minimal'' pair. Now the family $(X_i)_{i\in I}$ of \emph{strict}
subobjects of $X$ such that
$u\in\Iim(F(X_i)\rightarrow F(X))$ is small
(\hyperref[I.cor7.5]{7.5} in dual form). Since $F$ is left exact, this
family is cofiltered (\hyperref[I.rem8.12.6]{8.12.6}), and since $F$
commutes with small projective limits, one sees that
$X'=\varprojlim_i X_i$ is a least object of this family. (Use the fact
that, being left exact, $F$ carries monomorphisms to monomorphisms.) If
$u'\in F(X')$ is the unique element whose image in $F(X)$ is $u$, one
then sees that $(X',u')$ is a minimal pair dominating $(X,u)$. This
proves that $F$ is pro-representable, and the conclusion then follows
from the

\begin{enonce*}{Lemma 8.12.8.1}\etiquette[8.12.8.1]{I.lem8.12.8.1}
Let $F:C\rightarrow(\Ens)$ be a functor, where $C$ is a $𝒰$-category
in which small projective limits are representable. For $F$ to be
representable, it is necessary and sufficient that it be
pro-representable and commute with small projective limits.
\end{enonce*}

Necessity\pageoriginale is clear; let us prove sufficiency. To say that
$F$ is representable evidently means that ${}_F\backslash C$ admits an
initial object (in fact, the initial objects of ${}_F\backslash C$ are
precisely the isomorphisms $F\rightarrow X$, \ie the representation
data for $F$). Now, since $F$ is pro-representable,
$({}_F\backslash C)^\circ$ is filtered and equivalent to a small
category, so $X=\varprojlim_{({}_F\backslash C)^\circ} X$ is
representable in $C$, and since $F$ commutes with the limit under
consideration, it follows that $X$ is an initial object of
${}_F\backslash C$, \cqfd

\begin{enonce*}{Corollary 8.12.8}\etiquette[8.12.8]{I.cor8.12.8}
Under the hypotheses on $C$ of $\hyperref[I.cor8.12.7]{8.12.7}$, let $f:C\rightarrow C'$ be a
functor from $C$ to a $𝒰$-category $C'$. For $f$ to admit a left
adjoint, it is necessary and sufficient that $f$ commute with small
projective limits.
\end{enonce*}

Indeed, this reduces immediately to \hyperref[I.cor8.12.7]{8.12.7}, by
applying that statement to composite functors of the form
$X\mapsto \Hom(Y',f(X))$.

\begin{enonce*}[remark]{Examples 8.12.9}\etiquette[8.12.9]{I.exer8.12.9}
As was pointed out in \hyperref[I.exem7.13]{7.13}, the hypotheses on
$C$ in \hyperref[I.cor8.12.7]{8.12.7} and
\hyperref[I.cor8.12.8]{8.12.8} are satisfied if $C$ is the category of
$𝒰$-sheaves of sets on a topological space $X\in 𝒰$ (or, more
generally, on a $𝒰$-site (II \hyperref[II.def3.0.2]{(3.0.2)}). We give an instructive
example\footnote[1]{(due to H. BASS).} showing that the hypothesis of
the existence of a small cogenerating subcategory $D$ of $C$ is not
superfluous in \hyperref[I.cor8.12.7]{8.12.7}. Take for $C$ the
category of groups that are elements of $𝒰$. Let $J$ be the set of
isomorphism classes of simple groups in $C$; for every $j\in J$ choose
a simple group $G_j$ in the class $j$, and let $I$ be the filtered
ordered set of $𝒰$-small subsets of $J$, and, for $i\in I$, let
$X_i=\amalg_{j\in i}G_j$. The $X_i$ then form a\pageoriginale
projective system $(X_i)_{i\in I}$ in $C$, whose transition morphisms
are epimorphisms, and the corresponding functor
\begin{equation}
F(X)=\varinjlim_i \Hom(X_i,X)\tag{$\ast$}
\end{equation}
takes its values in $(𝒰\text{-}\Ens)$ (although the index set $I$
clearly does not have cardinality in $𝒰$); more precisely, let us show
that for every small full subcategory $C_\circ$ of $C$, there exists an
$i_\circ\in I$ such that the restriction of $F$ to $C_\circ$ is
represented by $X_{i_\circ}$, which will prove both that $F$ takes
values in $𝒰\text{-}\Ens$ and that it commutes with small projective
limits. To prove our assertion, it suffices to note that for every
$X\in \Ob C_\circ$, the cardinality of the set $J(X)$ of $j\in J$ for
which there exists a nontrivial homomorphism from $G_j$ to $X$ is
necessarily small, since such a morphism is necessarily a monomorphism
($G_j$ being simple); consequently, if $i_\circ$ is the subset of $J$
given by the union of the $J(X)$ for $X\in\ob C_\circ$, then $i_\circ$
is small, \ie $i_\circ\in I$, and it has the required property. On the
other hand, it is clear that $F$ is not representable, since
$\card I\not\in 𝒰$. From this and \hyperref[I.cor8.12.7]{8.12.7} one therefore concludes that
\emph{the category $C$ of groups in $𝒰$ admits no small full
cogenerating subcategory.} Since, on the other hand, the object $\ZZ$
of $C$ is a generator, it then follows from the proof of
\hyperref[I.cor7.12]{7.12} that there exists a two-generator group $G$
that cannot be embedded in an injective object of the category $C$ of
groups in $𝒰$. It moreover seems plausible that $C$ has no injective
objects other than the trivial groups.
\end{enonce*}


% END INCLUDED SGA4_English_current_buildable_source_20260906/source/01/drafts/SGA4_I_LINES_5595_5699_English.texfrag


% BEGIN INCLUDED SGA4_English_current_buildable_source_20260906/source/01/drafts/SGA4_I_LINES_5700_5804_English.texfrag
\subsection{Pro-representable functors and accessible
  functors}\etiquette[8.13]{I.sec8.13}\pageoriginale

\subsubsection{}\etiquette[8.13.1]{I.sec8.13.1}
In this section, we use some notions and results from the following
section, in particular \hyperref[I.prop9.11]{9.11} and
\hyperref[I.prop9.13]{9.13}, to obtain a pro-representability criterion that
we shall use (incidentally) in IV \hyperref[IV.theo9.5.4]{9.16}. In what follows, $C$ denotes a
$𝒰$-category satisfying condition $L$ of
\hyperref[I.def9.1]{9.1} b), this condition being satisfied, for example, if
small filtered inductive limits are representable in $C$.

\begin{enonce*}{Proposition 8.13.2}\etiquette[8.13.2]{I.prop8.13.2}
Let $C$ be as above, and let
$$
f:C\longrightarrow(\Ens)
$$
be a functor.
\begin{enumerate}
\renewcommand{\theenumi}{\alph{enumi}}
\renewcommand{\labelenumi}{\theenumi)}
\item Suppose that every object of $C$ is accessible (\hyperref[I.def9.3]{9.3}). If $f$
is pro-representable, then $f$ is accessible (\hyperref[I.def9.2]{9.2}) and
left exact.

\item Suppose that finite projective limits are representable in
$C$, and that $C$ admits a cardinal filtration
(\hyperref[I.def9.12]{9.12}). If $f$ is accessible and left exact, then $f$ is
pro-representable.
\end{enumerate}
\end{enonce*}

\noindent \emph{Proof.}
\begin{enumerate}
\renewcommand{\theenumi}{\alph{enumi}}
\renewcommand{\labelenumi}{\theenumi)}
\item The hypothesis on $C$ means that the covariant representable
functors from $C$ to $(\Ens)$ are accessible. The same is therefore
true of every small inductive limit of such functors
(\hyperref[I.prop9.6]{9.6} (i)), and hence also of every
pro-representable functor.

\item By \hyperref[I.them8.3.3]{8.3.3} (iii), it remains to prove that
$\Ob({}_F\backslash C)^\circ$ contains a small cofinal subcategory.
By hypothesis, there exists a cardinal $\Pi$ such that $f$ is
$\Pi$-accessible. Let $(X,u)$,\pageoriginale $u\in
F(X)$, be an object of ${}_F\backslash C$. With the notation of \hyperref[I.def9.12]{9.12} c),
we then have $X=\varinjlim_i X_i$, with $I$ filtered and large relative to
$\Pi$ and the $X_i$ in $C'=\Filt^{\Pi}(C)$, whence
$F(X)\xleftarrow{\sim}\varinjlim_i F(X_i)$. This shows that the
small subcategory $({}_F\backslash C')^\circ$ is cofinal in
$({}_F\backslash C)^\circ$, and completes the proof.
\end{enumerate}

\begin{enonce*}{Corollary 8.13.3}\etiquette[8.13.3]{I.cor8.13.3}
Let $C$ be a $𝒰$-category satisfying the following
conditions:
\begin{enumerate}
\renewcommand{\theenumi}{\alph{enumi}}
\renewcommand{\labelenumi}{\theenumi)}
\item Finite projective limits are representable in $C$.

\item Small filtered inductive limits are representable in
$C$.

\item The functor $\Ker$ on the category of double arrows of
$C$ is accessible (for example, it commutes with small filtered
$\varinjlim$).

\item Every strict epimorphism of $C$ is a universal strict epimorphism $\hyperref[I.sec10.2]{(10.2)}$.

\item There exists a small subcategory of $C$ generating by
strict epimorphisms $\hyperref[I.def7.1]{(7.1)}$.
\end{enumerate}
Under these conditions, a functor $f:C\rightarrow (\Ens)$ is
pro-representable if and only if it is left exact and
accessible (\hyperref[I.def9.2]{9.2}).
\end{enonce*}

Indeed, the conditions of \hyperref[I.prop8.13.2]{8.13.2} a) and b) on $C$ are
satisfied, by \hyperref[I.prop9.11]{9.11} and \hyperref[I.prop9.13]{9.13}
respectively.


\begin{enonce*}{Corollary 8.13.4}\etiquette[8.13.4]{I.cor8.13.4}
Let $f:C\rightarrow C'$ be a functor between $𝒰$-categories
satisfying conditions {\rm a)} through {\rm e)} of \hyperref[I.cor8.13.3]{8.13.3}. For
$f$ to admit a pro-adjoint, it is necessary and sufficient that $f$ be
left exact and accessible.
\end{enonce*}

Since\pageoriginale the representable functors
$h_{Y'}:X'\rightarrow \Hom(Y',X'):C'\rightarrow(\Ens)$ are left exact
and accessible (\hyperref[I.prop9.11]{9.11}), if $f$ has the same
properties, so do its composites with the preceding functors;
these are therefore pro-representable by \hyperref[I.cor8.13.3]{8.13.3},
\ie $f$ admits a pro-adjoint. Conversely, suppose that $f$ admits a
pro-adjoint, and let $C'_1$ be a small generating subcategory
of $C$, and $\Pi$ a cardinal such that the $Y'\in \Ob C'$ are
$\Pi$-accessible; to prove that $f$ is $\Pi$-accessible, it
therefore suffices to prove the same of its composites with the
$h_{Y'}$, $Y'\in \Ob C'$; by hypothesis, these composites are
pro-representable, hence accessible (\hyperref[I.cor8.13.3]{8.13.3}),
and hence $\Pi$-accessible provided $\Pi$ is taken sufficiently large, \cqfd

% END INCLUDED SGA4_English_current_buildable_source_20260906/source/01/drafts/SGA4_I_LINES_5700_5804_English.texfrag


% BEGIN INCLUDED SGA4_English_current_buildable_source_20260906/source/01/drafts/SGA4_I_LINES_5805_5940_English.texfrag
\section[Accessible functors and small generating subcategories]{Accessible functors, cardinal filtrations, and the construction
of small generating subcategories}\etiquette[9]{I.sec9}

The present section, which is more technical in nature than the other
sections of this exposé, will be used in this seminar only in IV 9 and
VI 4, which are not used elsewhere in the Seminar. The reader should
therefore omit the present section, at least on a first reading!

\setcounter{subsection}{-1}
\subsection{}\etiquette[9.0]{I.sec9.0}
All categories considered in this section are assumed to be
$𝒰$-categories. Except for the small index categories $I,J\ldots$ that
we shall have occasion to use, the developments that follow will apply
primarily to ``large'' categories $E,F\ldots$ that are stable under
small filtered inductive limits. Most often, however, it will suffice
that a slightly weaker condition be satisfied\pageoriginale (condition
$L$ in \hyperref[I.def9.1]{9.1} below). All cardinals considered in the
present section are assumed to belong to $𝒰$.

Following a suggestion of P. \textsc{Deligne}, we shall study, for a
functor $f:E\rightarrow F$ between large categories, a condition that
$f$ commute with certain types of filtered inductive limits. This
condition is remarkably stable, and is satisfied by the most important
functors occurring in nature. The applications we have in mind for our
seminar are 9.13.3, 9.13.4 (used in VI 4), and especially
\hyperref[I.cor9.25]{9.25}, which gives, in a nontrivial case, the
existence of a small generating family in a category of sections of a
fibered category; this result will be used in IV \hyperref[IV.theo9.5.4]{9.16}.

\begin{enonce*}{Definition 9.1}\etiquette[9.1]{I.def9.1}
\begin{enumerate}
\renewcommand{\theenumi}{\alph{enumi}}
\renewcommand{\labelenumi}{\theenumi)}
\item Let $I$ be a preordered set and $pi$ a cardinal. We say that $I$
is large relative to $pi$ if $I$ is filtered and every subset of $I$
of cardinality $\leq \pi$ has an upper bound in $I$.

\item Let $E$ be a category. If $pi$ is a cardinal, we say that $E$
satisfies condition $L_{\pi}$ if, for every small ordered set $I$ large
relative to $\pi$, $E$ is stable under inductive limits of type $I$. We
say that $E$ satisfies condition $L$ if there exists a cardinal
$\pi\in 𝒰$ such that $E$ satisfies condition $L_{\pi}$.
\end{enumerate}
\end{enonce*}

\setcounter{subsection}{1}
\subsubsection{}\etiquette[9.1.1]{I.sec9.1.1}
When, in \hyperref[I.def9.1]{9.1} a), one has $\pi\geq 2$, the second
stated condition already implies that $I$ is filtered, and if $\pi$ is
finite, for $I$ to be large relative to $\pi$ simply means that $I$ is
filtered. In what follows we shall be concerned almost exclusively with
the case in which $\pi$ is infinite. Note that if $\pi,\pi'$ are two
cardinals such that $\pi'\geq\pi$, then $I$ being large relative to
$\pi'$ evidently implies that it is large relative to $\pi$.

\subsubsection{}\etiquette[9.1.2]{I.sec9.1.2}
As\pageoriginale announced in \hyperref[I.sec9.0]{9.0}, conditions
$L_{\pi}$ and $L$ should be regarded as technical variants of the
stronger condition of stability under small filtered inductive limits.
It is clear that if $\pi,\pi'$ are cardinals such that
$\pi'\geq\pi$, then condition $L_{\pi}$ implies condition $L_{\pi'}$.

\begin{enonce*}{Definition 9.2}\etiquette[9.2]{I.def9.2}
Let $f:E\rightarrow F$ be a functor. If $\Pi$ is a cardinal, we say
that $f$ is $\Pi$-preaccessible (respectively $\Pi$-accessible) if $E$
satisfies $L_{\Pi}$ (\hyperref[I.def9.1]{9.1}) and if, for every ordered
set $I\in 𝒰$ large relative to $\Pi$, and every inductive system
$(X_i)_{i\in I}$ in $E$ of type $I$, the canonical morphism
$$
\varinjlim f(X_i)\longrightarrow f(\varinjlim X_i)
$$
is a monomorphism (respectively an isomorphism). We say that $f$ is
preaccessible (respectively accessible) (relative to the universe
$𝒰$) if there exists a cardinal $\Pi\in 𝒰$ such that $f$ is
$\Pi$-preaccessible (respectively $\Pi$-accessible).
\end{enonce*}

The category of $\Pi$-accessible (respectively accessible) functors
from $E$ to $F$ will be denoted $\Hhom(E,F)_{\Pi}$, respectively
$\Hhom(E,F)_{\acc}$.

\setcounter{subsection}{2} \setcounter{subsubsection}{0}
\subsubsection{}\etiquette[9.2.1]{I.sec9.2.1}
Evidently, a functor commuting with small filtered $\varinjlim$ (for
example, a functor admitting a right adjoint) is $\Pi$-accessible for
every cardinal $\Pi\geq 2$.

\begin{enonce*}{Definition 9.3}\etiquette[9.3]{I.def9.3}
Let $E$ be a category, $X$ an object of $E$,
$$
h^\circ_X:E\longrightarrow (𝒰\text{-}\Ens)
$$
the covariant functor it represents, and $\Pi\in 𝒰$ a cardinal. We say
that $X$ is a \emph{$\Pi$-preaccessible (respectively
$\Pi$-accessible) object of} $E$ if the functor
$h^\circ_X$\pageoriginale is $\Pi$-preaccessible (respectively
$\Pi$-accessible); we say that $X$ is \emph{preaccessible (respectively
accessible)} if there exists a cardinal $\Pi\in 𝒰$ such that $X$ is a
$\Pi$-preaccessible (respectively $\Pi$-accessible) object of $E$.
\end{enonce*}

\setcounter{subsection}{3} \setcounter{subsubsection}{0}
\subsubsection{}\etiquette[9.3.1]{I.sec9.3.1}
We denote by $E_{\Pi}$ the full subcategory of $E$ formed by the
$\Pi$-accessible objects of $E$.

When \hyperref[I.def9.3]{Definition 9.3} is applied to a category of the form
$\Hhom(C,F)$, the terminology just introduced presents \textit{a
priori} an ambiguity with the analogous terminology introduced in \hyperref[I.def9.2]{9.2},
when the objects of $\Hhom(C,F)$ are interpreted as functors; there
does not, however, seem to be any serious risk of confusion.

\begin{enonce*}{Definition 9.4}\etiquette[9.4]{I.def9.4}
Let $E$ be a category and let $\pi\in 𝒰$ be a cardinal. We say that
$E$ is a $\pi$-preaccessible (respectively $\pi$-accessible) category
if there exists in $E$ a small full subcategory $C$ that is generating
(\hyperref[I.def7.1]{7.1}) and whose objects are $\pi$-preaccessible
(respectively $\pi$-accessible) (\hyperref[I.def9.3]{9.3}). We say that
$E$ is a preaccessible (respectively accessible) category if there
exists a cardinal $\pi\in 𝒰$ such that $E$ is $\pi$-preaccessible
(respectively $\pi$-accessible).
\end{enonce*}

For important examples, cf. \hyperref[I.rem9.11.3]{9.11.3} below.

% END INCLUDED SGA4_English_current_buildable_source_20260906/source/01/drafts/SGA4_I_LINES_5805_5940_English.texfrag


% BEGIN INCLUDED SGA4_English_current_buildable_source_20260906/source/01/drafts/SGA4_I_LINES_5941_6010_English.texfrag
\begin{enonce*}{Proposition 9.5}\etiquette[9.5]{I.prop9.5}
Let $f:E\rightarrow F$ be a functor between categories such that $F$
is accessible and every object of $E$ is accessible. Then, if
$f$ admits a left adjoint, $f$ is accessible.
\end{enonce*}

Indeed,\pageoriginale by hypothesis, $F$ admits a small conservative
family of representable functors $F\rightarrow
(𝒰\text{-}\Ens)$ that are accessible, so we are immediately reduced
to showing that the composite of $f$ with each of the preceding functors
is accessible. But since $f$ admits a left adjoint,
these composites are functors $E\rightarrow
(𝒰\text{-}\Ens)$ that are \emph{representable,} and hence accessible
by hypothesis on $E$.

\begin{enonce*}{Proposition 9.6}\etiquette[9.6]{I.prop9.6}
Let $E$ and $F$ be two categories, let $\pi\in 𝒰$ be a cardinal, and
consider the full subcategory $\Hhom(E,F)_{\Pi}$ of
$\Hhom(E,F)$ formed by the $\Pi$-accessible functors $\hyperref[I.def9.2]{(9.2)}$ from $E$
to $F$.
\begin{enumerate}
\renewcommand{\theenumi}{\roman{enumi}}
\renewcommand{\labelenumi}{(\theenumi)}
\item This subcategory is stable under every type of $\varinjlim$
that is representable in~$F$.

\item Suppose that $F$ is $\pi$-accessible (\hyperref[I.def9.4]{9.4}), and let $J$ be a
category such that $\card  F\ell \  J\leq \pi$ and that the
$\varprojlim$ of type $J$ are representable in $F$, so the
$\varprojlim$ of type $J$ are representable in $\Hhom(E,F)$.
Then the subcategory $\Hhom(E,F)_{\Pi}$ is stable under the
$\varprojlim$ of type $J$.
\end{enumerate}
\end{enonce*}

\begin{enonce*}{Corollary 9.7}\etiquette[9.7]{I.cor9.7}
Let $E$ and $F$ be two categories, with $F$ accessible (\hyperref[I.def9.4]{9.4}).
Then the full subcategory $\Hhom(E,F)_{\acc}$ of $\Hhom(E,F)$
formed by the accessible functors is stable under every type of inductive
or projective limit, indexed by a small category
$J$, that is representable in $F$ (and hence in $\Hhom(E,F)$).
\end{enonce*}

Proof of \hyperref[I.prop9.6]{9.6}. Assertion (i) follows trivially from
the commutation of the functor $\varinjlim$ with arbitrary inductive
limits. For (ii), let $(f_j)_{j\in J}$ be a projective system
of $\Pi$-accessible functors $E\rightarrow F$, and let $f=\varprojlim
f_j$\pageoriginale be its projective limit, computed ``argument
by argument.'' We prove that $f$ is $\Pi$-accessible, \ie that
for every ordered set $I\in 𝒰$ large relative to $\Pi$, and
every inductive system $(X_i)_{i\in I}$ in $E$, the
canonical morphism
$$
\varinjlim_i ((\varprojlim_j f_j)(X_i))\longrightarrow
(\varprojlim_j f_j) (\varinjlim_i X_i)
$$
is an isomorphism. Now the ``argument-by-argument'' computation of the
functor $\varprojlim_j f_j$ allows us to identify the preceding canonical
morphism with the canonical morphism
$$
\varinjlim_i \varprojlim_j f_j (X_i)\longrightarrow \varprojlim_j
\varinjlim_i f_j (X_i)
$$
associated with the bifunctor
$$
(i,j)\longmapsto f_j(X_i):I\times J\longrightarrow F{\quoi}.
$$
Thus \hyperref[I.prop9.6]{9.6} is a consequence of the following more
general assertion:


% END INCLUDED SGA4_English_current_buildable_source_20260906/source/01/drafts/SGA4_I_LINES_5941_6010_English.texfrag


% BEGIN INCLUDED SGA4_English_current_buildable_source_20260906/source/01/drafts/SGA4_I_LINES_6011_6089_English.texfrag
\begin{enonce*}{Corollary 9.8}\etiquette[9.8]{I.cor9.8}
Let $F$ be a $\Pi$-accessible category, $I$ an ordered set large
relative to $\Pi$, $J$ a small category such that $\card F\ell\leq
\Pi$ and such that projective limits of type $J$ are representable in
$F$, and $h:I\times J\rightarrow F$ a functor; then the canonical
morphism
\hypertarget{sga4.exp.I.equation.9.8.1}{}\begin{equation}
\varinjlim_i \varprojlim_j h(i,j)\longrightarrow \varprojlim_j
\varinjlim_i h(i,j)\tag{9.8.1}
\end{equation}
is an isomorphism. In other words, the functor
\hypertarget{sga4.exp.I.equation.9.8.2}{}\begin{equation}
\varinjlim_i:\Hhom(I,F)\longrightarrow F\tag{9.8.2}
\end{equation}
commutes with projective limits of type $J$ (for every small category
$J$\pageoriginale such that $\card F\ell\,(J)\leq\Pi$ and such that
projective limits of type $J$ are representable in $F$). Equivalently,
for every $J$ as above, the functor
\hypertarget{sga4.exp.I.equation.9.8.3}{}\begin{equation}
\varinjlim_j:\Hhom(J,F)\longrightarrow F\tag{9.8.3}
\end{equation}
is $\Pi$-accessible.
\end{enonce*}

By hypothesis, $F$ admits a conservative family of covariant
representable $\Pi$-accessible functors
$g:F\rightarrow(𝒰\text{-}\Ens)$, so it is immediately clear that it
suffices to prove \hyperref[I.cor9.8]{9.8} in the case $F=𝒰\text{-}\Ens$. We shall then
prove the assertion in the form that \hyperlink{sga4.exp.I.equation.9.8.2}{(9.8.2)} commutes with projective
limits of type $J$, where $\card F\ell\,J\leq\Pi$. Since $I$ is
filtered, we know that \hyperlink{sga4.exp.I.equation.9.8.2}{(9.8.2)} commutes with finite projective limits
\hyperref[I.prop2.8]{(2.8)}. By a standard argument (cf. \hyperref[I.prop2.3]{2.3}), it therefore remains to prove
that it commutes with products indexed by a set $J$ such that
$\card J\leq\Pi$. This reduces us to proving the bijectivity of
\hyperlink{sga4.exp.I.equation.9.8.1}{(9.8.1)} when $J$ is discrete. Let us prove injectivity: consider two
elements $a,b$ of the first member; they therefore come from
$\prod_jh(i_\circ,j)$ %%
for a suitable $i_\circ\in I$, say $\prod_j a(i_\circ,j)$ and
$\prod_j b(i_\circ,j)$; suppose that the elements
$\prod_jh(\infty,j)$ of the second member that they define are equal,
\ie that for every $j$, there exists $i(j)\geq i_\circ$ such that
$a(i_\circ,j)$ and $b(i_\circ,j)$ have the same image in $h(i,j)$.
Since $I$ is large relative to $\card J$, it follows that there exists
a common upper bound $i_1\in I$ of all the $i(j)$, which implies that
the two elements under consideration of $\prod_jh(i_\circ,j)$ have the
same image in $\prod_jh(i_1,j)$ and therefore define %%
the same element of the first member of \hyperlink{sga4.exp.I.equation.9.8.1}{(9.8.1)}, establishing
injectivity.\pageoriginale For surjectivity, let
$\prod_j a(\infty,j)$ be an element of the second member; thus, for
every $j\in J$, $a(\infty,j)$ comes from an element of $h(i(j),j)$ for
a suitable $i(j)\in I$. As before, one can find a common upper bound
$i\in I$ of the $i(j)$, so $\prod_j a(\infty,j)$ comes from an element
$\prod_j a(i,j)$ of $\prod_jh(i,j)$, and hence lies in the image of
\hyperlink{sga4.exp.I.equation.9.8.1}{(9.8.1)}. This completes the proof.

\begin{enonce*}{Corollary 9.9}\etiquette[9.9]{I.cor9.9}
Let $F$ be a $\Pi$-accessible (respectively accessible) category. Then,
for every category $J$ such that $\card F\ell\,J\leq\Pi$ (respectively
every small category $J$), and for every functor
$(X_j)_{j\in J}:J\rightarrow F$ whose projective limit in $F$ is
representable, if the $X_j$ are $\Pi$-accessible (respectively
accessible), then so is $\varprojlim X_j$.
\end{enonce*}

Indeed, the functor $F\rightarrow(𝒰\text{-}\Ens)$ represented by
$\varprojlim X_j$ is the projective limit of the functors represented
by the $X_j$, and one applies \hyperref[I.prop9.6]{9.6} (ii) to the
projective system formed by these functors.

\begin{enonce*}[remark]{Remark 9.10}\etiquette[9.10]{I.rem9.10}
In the statements \hyperref[I.prop9.6]{9.6},
\hyperref[I.cor9.7]{9.7}, \hyperref[I.cor9.8]{9.8}, and
\hyperref[I.cor9.9]{9.9}, one may everywhere replace the words
``$\Pi$-accessible'', ``accessible'' by ``$\Pi$-preaccessible'',
``preaccessible''. Indeed, the proof given also proves this variant of
the preceding statements.
\end{enonce*}


% END INCLUDED SGA4_English_current_buildable_source_20260906/source/01/drafts/SGA4_I_LINES_6011_6089_English.texfrag


% BEGIN INCLUDED SGA4_English_current_buildable_source_20260906/source/01/drafts/SGA4_I_LINES_6090_6222_English.texfrag
\begin{enonce*}{Proposition 9.11}\etiquette[9.11]{I.prop9.11}
Let $E$ be a category satisfying condition $L$
(\hyperref[I.def9.1]{9.1}), and let $C$ be a small full generating
subcategory (\hyperref[I.def7.1]{7.1}). Suppose that the following
condition is satisfied:
\begin{enumerate}
\renewcommand{\theenumi}{\alph{enumi}}
\renewcommand{\labelenumi}{\theenumi)}
\item $E$ is stable under kernels of pairs of arrows, and the functor
$\Ker$ on the category of pairs of arrows in $E$ is accessible (for
example, it\pageoriginale commutes with small filtered inductive
limits).

Then every object of $E$ is preaccessible (\hyperref[I.def9.3]{9.3}),
\textit{a fortiori} $E$ is preaccessible. Suppose that $C$ is even
generating by strict epimorphisms (\hyperref[I.def7.1]{7.1}), and
suppose that $E$ satisfies the following conditions:

\item $E$ is stable under fiber products.

\item Every small strictly epimorphic family in $E$ is universally
strictly epimorphic (\hyperref[I.sec10.3]{10.3}), or small direct sums
are representable in $E$ and every strict epimorphism of $E$ is a
universal strict epimorphism.
\end{enumerate}

Then every object of $E$ is accessible, \textit{a fortiori} $E$ is
accessible.
\end{enonce*}

The fact that, under (a), every object of $E$ is preaccessible follows
from the fact that for every object $X$ of $E$, the set of strict
subobjects of $X$ is small (\hyperref[I.prop7.4]{7.4}), and from the

\begin{enonce*}{Lemma 9.11.1}\etiquette[9.11.1]{I.lem9.11.1}
Under the conditions of (\hyperref[I.prop9.11]{9.11}a), let
$X\in\ob E$ and let $\Pi$ be a cardinal such that $E$ satisfies
$L_{\Pi}$ (\hyperref[I.def9.1]{9.1}), the functor $\Ker$ in
(\hyperref[I.prop9.11]{9.11}a) is $\Pi$-accessible, and the set of
strict subobjects of $X$ has cardinality bounded by $\Pi$. Then $X$ is
$\Pi$-preaccessible.
\end{enonce*}

Indeed, let $I$ be an ordered set large relative to $\Pi$, let
$(Y_i)_{i\in I}$ be an inductive system in $E$ with inductive limit
$Y$, and let $(u_i,v_i:X\rightrightarrows Y_i)$ be a pair of arrows
such that the composite pair
$ \xymatrix{ X \ar@<2pt>[r]^{u,v} \ar@<-2pt>[r]& Y }$ satisfies
$u=v$; let us prove that there exists $j\geq i$ in $I$ such that
$u_j=v_j$. To this end,\pageoriginale consider the inductive system of
pairs of arrows $(u_j,v_j:X\rightrightarrows Y_j)_{j\geq i}$, whose
inductive limit is $(u,v)$. By hypothesis one has
$X=\Ker(u,v)=\varinjlim_j\Ker(u_j,v_j)=\varinjlim_jX_j$, where
$X_j=\Ker(u_j,v_j)$. Now the $X_j$ are strict subobjects of $X$, so
there exists a subset $I'$ of $I$, formed by indices $j\geq i$, such
that $\card I'\leq\Pi$ and every $X_j$ is equal to one of the $X_{i'}$
$(i'\in I')$. Since $I$ is large relative to $\Pi$, there exists an
upper bound $j$ of $I'$ in $I$. Then $X_j$ contains all the $X_{j'}$
for $j'\geq i$, so $X_j\rightarrow X$ is an epimorphism (since the
family of the $X_{j'}\rightarrow X$ is epimorphic), and hence an
isomorphism because it is a strict monomorphism. Thus $u_j=v_j$, which
proves \hyperref[I.lem9.11.1]{9.11.1}.

The second assertion of \hyperref[I.prop9.11]{9.11} follows from the
first and from the

\begin{enonce*}{Lemma 9.11.2}\etiquette[9.11.2]{I.lem9.11.2}
Under the conditions of (\hyperref[I.prop9.11]{9.11}a), b), c), let
$X$ be an object of $E$, and let $\Pi$ be an infinite cardinal such
that $E$ satisfies $L_{\Pi}$ $\hyperref[I.def9.1]{(9.1)}$, that
$\card\ob C_{/X}\leq\Pi$, that for two objects $X'\rightarrow X$ and
$X''\rightarrow X$ of $C_{/X}$, $\fprod{X'}{X''}{X}$ is
$\Pi$-preaccessible, and finally that $X$ is $\Pi$-preaccessible.
Then $X$ is $\Pi$-accessible.
\end{enonce*}

With the notation of the proof of \hyperref[I.lem9.11.1]{9.11.1}, it
suffices to prove that every morphism $u:X\rightarrow Y$ comes from a
morphism $u_i:X\rightarrow Y_i$. Consider the family of morphisms
$Y_i\rightarrow Y$, which is strictly epimorphic; by b) and c), the
family of morphisms $\fprod{Y_i}{X}{Y}\rightarrow X$ is also strictly
epimorphic; on the other hand, for every $i$, since $C$ is generating
by strict epimorphisms, the family of arrows
$$
T\longrightarrow \fprod{Y_i}{X}{Y}
$$
with\pageoriginale source in $C$ is strictly epimorphic. In the second
alternative considered in c), one can find such a strictly epimorphic
arrow whose source is a (small) sum of objects of $C$. By the
transitivity of the notion of a universally strictly epimorphic family
(II \hyperref[II.def2.5]{2.5}), it then follows that the family of arrows
$T_\alpha\xrightarrow{f_\alpha}X$ with source in $C$ that factor
through one of the $\fprod{Y_i}{X}{Y}$ is strictly epimorphic. Let $J$
be the index set of this family; its cardinality is bounded by
$\card\ob C_{/X}\leq\Pi$. For every $\alpha\in J$, choose an
$i=i(\alpha)\in I$ and an $X$-morphism
$T_\alpha\rightarrow\fprod{Y_i}{X}{Y}$, or, equivalently, a
$v_\alpha:T_\alpha\rightarrow Y_i$ such that
$p_iv_\alpha=uf_\alpha$, where $p_i:Y_i\rightarrow Y$ is the
canonical morphism. Since $I$ is large relative to $\Pi$, one can
choose $i(\alpha)$ independently of $\alpha$, say equal to $i$. For
every pair of indices $\alpha,\beta\in J$, consider the composites
$$
\pr_1v_\alpha,\pr_2v_\beta:
\fprod{T_\alpha}{T_\beta}{X}\rightrightarrows Y_i{\quoi}.
$$
Their composites with $Y_i\rightarrow Y$ are equal, so, since
$\fprod{T_\alpha}{T_\beta}{X}$ is $\Pi$-preaccessible by hypothesis,
there exists an index $i'=i(\alpha,\beta)\geq i$ such that the
composites of the arrows under consideration with
$Y_i\rightarrow Y_{i'}$ are equal. Since the set of pairs
$\alpha,\beta$ has cardinality $\leq\Pi^2=\Pi$ ($\Pi$ being
infinite), it follows once again that $i'$ can be chosen independently
of $\alpha,\beta$. We may evidently suppose that $i'=i$. But then,
since the family $(f_\alpha:T_\alpha\rightarrow X)$ is strictly
epimorphic, one can find a morphism $u_i:X\rightarrow Y_i$ such that
$u_if_\alpha=v_\alpha$. One then has $p_iu_i=u$, since for every
$\alpha$ one has
$(p_iu_i)f_\alpha=p_i(u_if_\alpha)=p_iv_\alpha=uf_\alpha$, and the
family of the $f_\alpha$ is epimorphic. This completes the proof of
\hyperref[I.lem9.11.2]{9.11.2}.

\begin{enonce*}[remark]{Remark 9.11.3}\etiquette[9.11.3]{I.rem9.11.3}
As\pageoriginale a special case of \hyperref[I.prop9.11]{9.11}, one
sees that every object in the category $E$ is accessible in each of the
following two cases: 1) $E$ is a $𝒰$-abelian category with exact
filtered inductive limits and admitting a small generating subcategory
(here it is the \emph{second} alternative of c) that applies). 2) $E$
is the category of sheaves of sets on a topological space $X\in 𝒰$.
More generally, it suffices that $E$ be the category of sheaves of sets
on a $𝒰$-site (II \hyperref[II.def2.1]{2.1}), or that $E$ be a $𝒰$-topos (IV \hyperref[IV.def1.1]{1.1}).
\end{enonce*}


% END INCLUDED SGA4_English_current_buildable_source_20260906/source/01/drafts/SGA4_I_LINES_6090_6222_English.texfrag


% BEGIN INCLUDED SGA4_English_current_buildable_source_20260906/source/01/drafts/SGA4_I_LINES_6223_6376_English.texfrag
\begin{enonce*}{Definition 9.12}\etiquette[9.12]{I.def9.12}
Let $E$ be a category. A cardinal filtration of $E$ is an
increasing filtration $(\Filt^{\Pi}(E))_{\Pi\geq \Pi_\circ}$ of $E$
by strictly full subcategories $\Filt^{\Pi}(E)$,
indexed by the cardinals $\Pi\in 𝒰$ such that $\Pi\geq
\Pi_\circ$ (where $\Pi_\circ$ is a fixed infinite cardinal, depending on
the cardinal filtration under consideration), and satisfying the following
conditions:
\begin{enumerate}
\renewcommand{\theenumi}{\alph{enumi}}
\renewcommand{\labelenumi}{\theenumi)}
\item For every $\Pi\geq\Pi_\circ$, $\Filt^{\Pi}(E)$ is
equivalent to a small category.

\item $E$ satisfies $L_{\Pi_\circ}$ (\hyperref[I.def9.3]{9.3}), and for every
$\Pi\geq \Pi_\circ$, $\Filt^{\Pi}(E)$ is stable in $E$ under
filtered inductive limits indexed by ordered sets
$I$ large relative to $\Pi_\circ$ such that $\card I\leq \Pi$.

\item For every $\Pi\geq \Pi_\circ$, and every $X\in \ob E$, one can
find an isomorphism
$$
X\simeq \varinjlim_I X_i{\quoi},
$$
where $(X_i)_{i\in I}$ is an inductive system in $E$ indexed by an
ordered set $I$\pageoriginale large relative to $\Pi$
(\hyperref[I.def9.1]{9.1}), and where
the $X_i$ belong to $\Filt^{\Pi}(E)$. Moreover, if $X\in \ob
\Filt^{\Pi}(E)$, $\Pi'\geq \Pi$, one may choose $I$ such that $\card
I\leq {\Pi'}^{\Pi}$.
\end{enumerate}
\end{enonce*}

\setcounter{subsection}{12} \setcounter{subsubsection}{0}
\subsubsection{}\etiquette[9.12.1]{I.9.12.1}
Note that it follows from b) and c) that every $X\in \ob E$ belongs
to some $\Filt^{\Pi}(E)$, for $\Pi$ sufficiently large.

\begin{enonce*}[remark]{Example 9.12.2}\etiquette[9.12.2]{I.exem9.12.2}
Take $E=(𝒰\text{-}\Ens)$, $\Phi_\circ=\aleph_\circ$,
$\Filt^{\Pi}(E)=$ the full subcategory of $E$ consisting of the sets
such that $\card X\leq \Pi$. More generally:
\end{enonce*}

\begin{enonce*}{Proposition 9.13}\etiquette[9.13]{I.prop9.13}
Let $E$ be a $𝒰$-category. Suppose that $E$ is stable
under small filtered $\varinjlim$, under sums of two objects, and under
cokernels of pairs of arrows, that $E$ is stable under fiber products,
and that the epimorphisms in $E$ are universal strict
epimorphisms. Let $C$ be a small full subcategory of $E$
that is generating by strict epimorphisms. Let $\Pi_\circ$ be an
infinite cardinal $\geq \card F\ell\,  C$. For every cardinal $\Pi\geq
\Pi_\circ$, let $\Filt^{\Pi}(E)$ be the strictly full
subcategory of $E$ consisting of the objects $X$ of $E$ for which there exists a
strictly epimorphic family $(X_i\rightarrow X)_{i \in I}$ with target
$X$, such that $\card I\leq \Pi$ and $X_i\in \ob C$ for every
$i\in I$. Then $(\Filt^{\Pi}(E))_{\Pi\geq \Pi_\circ}$ is a
cardinal filtration of $E$. Moreover, for every $X\in \ob
\Filt^{\Pi}(E)$, the cardinality of the set of arrows of $C_{/X}$
(and \textit{a fortiori}, of the set of objects of $C_{/X}$) is
bounded by $\Pi^{\Pi_\circ}$.
\end{enonce*}

This last assertion is precisely \hyperref[I.prop7.6]{7.6}.

To show that this is a cardinal filtration, one must prove
conditions a), b), and c) of \hyperref[I.def9.12]{9.12}. Condition a)
follows at once from
\hyperref[I.var7.5.2]{7.5.2}. For b), suppose that $X=\varinjlim_I X_i$, with
$\card I\leq \Pi$ and the\pageoriginale $X_i\in \ob \Filt^{\Pi}(E)$.
Thus the family of canonical morphisms $X_i\rightarrow X$ is
strictly epimorphic, and by the hypothesis on the $X_i$, for
every $i\in I$ there exists a strictly epimorphic family $(X_{ij}\rightarrow
X_i)_{j\in J_i}$, with $\card J_i\leq \Pi$. Taking the sums
$\amalg_i X_i$ and $\amalg_{i,j}X_{ij}$, one sees by the transitivity of
universal strict epimorphisms (II \hyperref[II.def2.5]{2.5}) that the family of
composites $X_{ij}\rightarrow X_i\rightarrow X$ (indexed by
the disjoint union of the $J_i$, for $i\in I$) is strictly epimorphic.
But $\card J\leq \Pi^2=\Pi$, whence the conclusion. (NB. we have
used only the fact that $\varinjlim_I X_i$ exists, and not the
fact that $I$ is large relative to $\Pi_\circ$, or even that it is filtered.)
Finally, let us prove c). Note that for every $X\in \ob E$, the family of
arrows $X_i\rightarrow X$ with source in $\ob C$ is strictly
epimorphic, since $C$ is generating by strict
epimorphisms,
and is plainly small, so there exists a cardinal $\Pi\geq \Pi_\circ$
such that $X\in \ob \Filt^{\Pi}(E)$. It remains to prove that if
$\Pi'\geq \Pi\geq \Pi_\circ$ are cardinals, and if $X\in \ob
\Filt^{\Pi'}(E)$, then there is an isomorphism
$$
X\simeq \varinjlim_I X_i{\quoi},
$$
with $I$ large relative to $\Pi$, $\card I\leq {\Pi'}^{\Pi}$, and the $X_i$
in $\ob \Filt^{\Pi}(E)$. Now, since $C$ is generating by
strict epimorphisms, one has
\begin{equation*}
X\simeq \varinjlim_{C_{/X}} T\tag{$\ast$}{\quoi}.
\end{equation*}
Let us also note %%
that by successively adjoining to $C$ sums and cokernels of
pairs of arrows in $C$, and doing the same for the category thus obtained,
and so on, one reduces to the case where $C$ is stable under sums of two objects
in $E$ and under cokernels of pairs of arrows in $E$, without
losing the hypothesis $\Pi_\circ \geq \card F\ell\, C$. We may
suppose moreover that, if $E$ contains a\pageoriginale chosen initial
object $∅_E$, then $∅_E\in \ob C$. This implies
that the category $C_{/X}$ is stable under sums of two objects and
cokernels of pairs of arrows. It is then filtered, since it is
nonempty: if it were empty, the relation $(\ast)$
would show that $X$ is an initial object, and hence $C_{/X}$ would contain the
arrow $∅_E\rightarrow X$, a contradiction. Let
$I$ be the set, ordered by inclusion, of full subcategories
$i$ of $C_{/X}$ that are filtered and such that $\card \ob i\leq
\Pi$. For every $i\in I$, let $X_i\in \ob E$ be the inductive limit of the
composite functor $i\rightarrow C_{/X}\rightarrow E$, which exists by
hypothesis. Then plainly
$$
\varinjlim_I X_i \simeq \varinjlim_{C/X} T\simeq X{\quoi}.
$$
On the other hand, we have already noted that $\card \ob C_{/X}\leq
{\Pi'}^{\Pi_\circ}$. It follows that $I$ is large relative to $\Pi$, and
that $\card I\leq
({\Pi'}^{\Pi_\circ})^{\Pi}={\Pi'}^{\Pi\Pi_\circ}={\Pi'}^{\Pi}$ by
the following lemma, whose proof is left to the reader
(where we take $I=C_{/X}$, $c={\Pi'}^{\Pi_\circ}$):

\begin{enonce*}{Lemma 9.13.1}\etiquette[9.13.1]{I.lem9.13.1}
Let $J$ be a filtered category, and let $c$ and $\Pi$ be two cardinals such
that $c\geq \Sup(\card  F\ell\,  J,\Pi)$, and suppose that $\card
\Hom(j,J')\leq \Pi_\circ$ for every pair of objects $j$, $j'$ of
$J$. Let $I$ be the set of full filtered subcategories $i$
of $J$ such that $\card \ob i\leq \Pi$. Then, ordered by
inclusion, $I$ is large relative to $\Pi$, and one has $\card I\leq
c^{\Pi}$.
\end{enonce*}

This completes the proof of \hyperref[I.prop9.13]{9.13}.

\begin{enonce*}[remark]{Remark 9.13.2}\etiquette[9.13.2]{I.rem9.13.2}
A\pageoriginale slight additional effort should make it possible to
replace in \hyperref[I.prop9.13]{9.13} the hypothesis that $E$ is stable under
small filtered inductive limits by the hypothesis that $E$
satisfies $L$ (\hyperref[I.def9.1]{9.1}), provided $E$ is stable under
small sums, or that in $E$ every epimorphic \emph{family}
is universally epimorphic. In the proof
of c), one must then choose $\Pi_\circ$ such that $E$ satisfies $L_{\Pi_\circ}$,
and restrict oneself to those full subcategories $i$ of $C_{/X}$ that are
not only filtered, but such that the preordered set $\ob
i$ is large relative to $\Pi_\circ$. Using \hyperref[I.sec8]{Section 8} and the hypothesis that
$E$ satisfies $L_{\Pi_\circ}$, one then finds that the $X_i$
exist, and one should be able to conclude by a suitable variant of
\hyperref[I.lem9.13.1]{9.13.1}, which the editor has not verified.
\end{enonce*}


% END INCLUDED SGA4_English_current_buildable_source_20260906/source/01/drafts/SGA4_I_LINES_6223_6376_English.texfrag


% BEGIN INCLUDED SGA4_English_current_buildable_source_20260906/source/01/drafts/SGA4_I_LINES_6377_6439_English.texfrag
\begin{enonce*}{Proposition 9.14}\etiquette[9.14]{I.prop9.14}
Let $f:E\rightarrow F$ be an accessible functor
(\hyperref[I.def9.2]{9.2}) between two categories equipped with cardinal
filtrations $(\Filt^{\Pi}(E))_{\Pi\geq\Pi_\circ}$ and
$(\Filt^{\Pi}(F))_{\Pi\geq\Pi'_\circ}$. Then there exists a cardinal
$\Pi_1\geq\Sup(\Pi_\circ,\Pi'_\circ)$ such that for every cardinal
$\Pi\geq\Pi_1$, one has
\hypertarget{sga4.exp.I.equation.9.14.1}{}\begin{equation}
f(\Filt^{\Pi}(E))\subset \Filt^{\Pi^{\Pi_1}}(F){\quoi}.\tag{9.14.1}
\end{equation}
In particular, if $\Pi=2^c$ with $c\geq\Pi_1$, whence
$\Pi^{\Pi_1}=2^{c\Pi_1}=2^c=\Pi$, one has
\hypertarget{sga4.exp.I.equation.9.14.2}{}\begin{equation}
f(\Filt^{\Pi}(E))\subset \Filt^{\Pi}(F){\quoi}.\tag{9.14.2}
\end{equation}
\end{enonce*}

We immediately point out the

\begin{enonce*}{Corollary 9.15}\etiquette[9.15]{I.cor9.15}
Let\pageoriginale $E$ be a category equipped with two cardinal
filtrations\linebreak $(\Filt^{\Pi}(E))_{\Pi\geq\Pi_\circ}$ and
$({\Filt'}^{\Pi}(E))_{\Pi\geq %%
\Pi'_\circ}$. Then there exists a cardinal
$\Pi_1\geq\Sup(\Pi_\circ,\Pi'_\circ)$ such that, for every cardinal
$c\geq\Pi_1$, putting $\Pi=2^c$, one has
$\Filt^{\Pi}(E)={\Filt'}^{\Pi}(E)$.
\end{enonce*}

\noindent \emph{Proof of} \hyperref[I.prop9.14]{9.14}. Let $c$ be a cardinal such that
$c\geq\Sup(\Pi_\circ,\Pi'_\circ)$ and such that $f$ is
$c$-accessible. Let, moreover, $\Pi_1\geq c$ be such that
\begin{equation}
f(\Filt^c(E))\subset \Filt^{\Pi_1}(F){\quoi}.\tag{$\ast$}
\end{equation}
Such a $\Pi_1$ exists because $\Filt^c(E)$ is equivalent to a small
category \hyperref[I.def9.12]{(9.12 a)}), and hence so is $f(\Filt^c(E))$, so that one can
apply \hyperref[I.9.12.1]{9.12.1} to the objects of the latter to find a
$\Filt^{\Pi}(F)$ containing them all %%
(taking into account that the $\Filt^{\Pi}(F)$ are full
subcategories). Let then $\Pi\geq\Pi_1$ and
$X\in\ob\Filt^{\Pi}(E)$; let us prove that
$f(X)\in\Filt^{\Pi^{\Pi_1}}(F)$. Indeed, write
$$
X=\varinjlim_I X_i{\quoi},
$$
where the $X_i\in\ob\Filt^c(E)$, $I$ is large relative to $c$, and
$\card I\leq\Pi^c\leq\Pi^{\Pi_1}$
(\hyperref[I.def9.12]{9.12}c)). Since $f$ is $c$-accessible and $I$ is
large relative to $c$, one has
$$
f(X)\simeq \varinjlim_I f(X_i){\quoi},
$$
and since $\card I\leq\Pi^{\Pi_1}$ and
$f(X_i)\in\ob\Filt^{\Pi_1}(F)\subset
\ob\Filt^{\Pi^{\Pi_1}}(F)$ by $(\ast)$, one has
$f(X)\in\Filt^{\Pi^{\Pi_1}}(F)$ by
\hyperref[I.def9.12]{9.12} b), \cqfd

\setcounter{subsection}{15} \setcounter{subsubsection}{0}
\subsubsection{}\etiquette[9.15.1]{I.sec9.15.1}
The notion of cardinal filtration \hyperref[I.def9.12]{9.12} is of
interest mainly when the objects of $E$ are accessible. We point out
that it follows from \hyperref[I.prop9.11]{9.11} that this condition is
satisfied when, in addition to the hypotheses\pageoriginale of
\hyperref[I.prop9.13]{9.13}, one supposes that the functor $\Ker$ on
the category of pairs of arrows in $E$ is accessible. We also point
out:


% END INCLUDED SGA4_English_current_buildable_source_20260906/source/01/drafts/SGA4_I_LINES_6377_6439_English.texfrag


% BEGIN INCLUDED SGA4_English_current_buildable_source_20260906/source/01/drafts/SGA4_I_LINES_6440_6522_English.texfrag
\begin{enonce*}{Proposition 9.16}\etiquette[9.16]{I.prop9.16}
Let $E$ be a category equipped with a cardinal filtration
$(\Filt^{\Pi}(E))_{\Pi\geq\Pi_\circ}$. Suppose that the objects of
$\Filt^{\Pi_\circ}(E)$ are accessible objects of $E$; then
there exists a cardinal $\Pi_1\geq\Pi_\circ$ in $𝒰$ such that
the objects of $\Filt^{\Pi_\circ}(E)$ are $\Pi_1$-accessible; if
$\Pi_1$ is chosen in this way, then for every cardinal $\Pi\geq \Pi_1$,
we have (with the notation of \hyperref[I.sec9.3.1]{9.3.1}):
\hypertarget{sga4.exp.I.equation.9.16.1}{}\begin{equation}
E_{\Pi}\subset \Filt^{\Pi}(E)\subset
E_{(\Pi^{\Pi_\circ})}{\quoi}.\tag{9.16.1}
\end{equation}
\end{enonce*}

The existence of $\Pi_1$ follows immediately from the fact that
$\Filt^{\Pi_\circ}(E)$ is equivalent to a small
category (\hyperref[I.def9.12]{9.12}a)). Now let $\Pi\geq \Pi_1$. If $X$
  belongs to $\Filt^{\Pi}(E)$,
write $X\simeq \varinjlim_I X_i$ with $\card I\leq
\Pi^{\Pi_\circ}$, $X_i\in \ob \Filt^{\Pi_\circ}(E)$ for every $i$
(\hyperref[I.def9.12]{9.12}c)). It then follows from \hyperref[I.cor9.9]{9.9} that $X$ is
$\Pi^{\Pi_\circ}$-accessible, whence the second inclusion \hyperlink{sga4.exp.I.equation.9.16.1}{(9.16.1)}.
Suppose that $X$ is $\Pi$-accessible, and write $X\simeq
\varinjlim_I X_i$, with $I$ large relative to $\Pi$ and the $X_i\in \ob
\Filt^{\Pi}(E)$ (\hyperref[I.def9.12]{9.12}c)). By the hypothesis on $X$, the given
isomorphism $X\xrightarrow{\sim} \varinjlim_I X_i$ factors through one of the
$X_i$, so $X$ is isomorphic to a direct factor of that $X_i$. We
conclude that $X\in \ob \Filt^{\Pi}(E)$, and hence obtain the first inclusion
\hyperlink{sga4.exp.I.equation.9.16.1}{(9.16.1)}, by the following lemma.

\begin{enonce*}{Lemma 9.16.2}\etiquette[9.16.2]{I.lem9.16.2}
Every object of $E$ that is a direct factor of an object of
$\Filt^{\Pi}(E)$ belongs to $\Filt^{\Pi}(E)$.
\end{enonce*}

Indeed, if $X$ is the image of a projector $p$ on an object $Y$ of
$E$ (\ie an endomorphism $p$ such that $p^2=p$), and if $I$ is an
ordered\pageoriginale  filtered set, $X$ is the inductive limit
of the filtered inductive system $(Y_i)_{i\in I}$ defined by $Y_i=Y$
for every $i\in I$, $p:Y_i\rightarrow Y_j$ if $i<j$. Taking $I$
large relative to $\Pi_\circ$ and $\card I\leq \Pi$, we thus see that if
$Y\in \Filt^{\Pi}(E)$, then the same holds for $X$, by
\hyperref[I.def9.12]{9.12} b).

\begin{enonce*}{Corollary 9.17}\etiquette[9.17]{I.cor9.17}
Under the conditions of \hyperref[I.prop9.16]{9.16}, for every cardinal $c\geq
\Pi_1$, setting $\Pi=2^c$, $\Filt^{\Pi}(E)$ coincides with the
strictly full subcategory $E_{\Pi}$ of $E$ consisting of the
$\Pi$-accessible objects.
\end{enonce*}

\begin{enonce*}{Corollary 9.18}\etiquette[9.18]{I.cor9.18}
Let $E$ be a category satisfying condition $L_{\Pi}$
(\hyperref[I.def9.1]{9.1}), where $\Pi$ is a cardinal, and let $C$ be a
full subcategory
of $E$. Denote by $\Ind(C)_{\Pi}$ the full subcategory of
the category $\Ind(C)$ of ind-objects of $C$ (\hyperref[I.sec8.2]{8.2})
consisting of
ind-objects of the form $(X_i)_{i\in I}$, where $I$ is an ordered set
large relative to $\Pi$. Consider the canonical functor
\hypertarget{sga4.exp.I.equation.9.18.1}{}\begin{equation}
(X_i)_{i\in I}\longmapsto \varinjlim_I
X_i:\Ind(C)_{\Pi}\longrightarrow E{\quoi}.\tag{9.18.1}
\end{equation}
\begin{enumerate}
\renewcommand{\theenumi}{\alph{enumi}}
\renewcommand{\labelenumi}{\theenumi)}
\item This functor is fully faithful if and only if every object $X$
of $C$ is a $\Pi$-accessible object
(\hyperref[I.def9.3]{9.3}) of $E$.

\item Assume the conditions of \hyperref[I.cor9.17]{9.17}, in particular
$\Pi=2^c$, and take $C=\Filt^{\Pi}(E)$. Then the functor \hyperlink{sga4.exp.I.equation.9.18.1}{(9.18.1)}
is an equivalence of categories.
\end{enumerate}
\end{enonce*}

Assertion a) is an immediate generalization of \hyperref[I.prop8.7.5]{8.7.5}
a), and is proved in the same way. Part b) then follows from
\hyperref[I.cor9.17]{9.17} and condition \hyperref[I.def9.12]{9.12} c) on
cardinal filtrations, which implies that the functor under consideration is
essentially surjective.


% END INCLUDED SGA4_English_current_buildable_source_20260906/source/01/drafts/SGA4_I_LINES_6440_6522_English.texfrag


% BEGIN INCLUDED SGA4_English_current_buildable_source_20260906/source/01/drafts/SGA4_I_LINES_6523_6659_English.texfrag
\begin{enonce*}{Corollary 9.19}\etiquette[9.19]{I.cor9.19}
Under\pageoriginale the conditions of \hyperref[I.cor9.17]{9.17}, let
$C=\Filt^{\Pi}(E)$, let $F$ be a $𝒰$-category, and consider the
functor
\hypertarget{sga4.exp.I.equation.9.19.1}{}\begin{equation}
\Hhom(E,F)_{\Pi}\longrightarrow \Hhom(C,F)\tag{9.19.1}
\end{equation}
induced by the “restriction to $C$” functor $f\mapsto f| C$,
where the source of \hyperlink{sga4.exp.I.equation.9.19.1}{(9.19.1)} is the category of
$\Pi$-accessible functors from $E$ to $F$ (\hyperref[I.def9.2]{9.2}). The
preceding functor is
fully faithful. If $F$ satisfies condition $L_{\Pi}$
(\hyperref[I.def9.1]{9.1}), then the functor \hyperlink{sga4.exp.I.equation.9.19.1}{(9.19.1)} is an equivalence of
categories.
\end{enonce*}

The first assertion is proved as in \hyperref[I.prop7.8]{7.8}, using
\hyperref[I.cor9.18]{9.18} b). The second is obtained by constructing a
quasi-inverse functor to \hyperlink{sga4.exp.I.equation.9.19.1}{(9.19.1)}, associating to every
functor $f:C\rightarrow F$ the functor $(X_i)_{i\in I}\rightarrow
\varinjlim_i f(X_i)$ from $\Ind(E_{\Pi})$ to $F$, and using
\hyperref[I.cor9.18]{9.18} b) to deduce from it a functor $\bar{f}:E\rightarrow
F$. It remains only to %%
show that the%%
latter is $\Pi$-accessible. But this is proved in the same way as the
analogous assertion \hyperref[I.prop8.7.3]{8.7.3}.

\begin{enonce*}{Corollary 9.20}\etiquette[9.20]{I.cor9.20}
Let $E$ be a category admitting a cardinal filtration and such
that every object of $E$ is accessible (\cf \hyperref[I.prop9.16]{9.16}), and let $F$
be a category. Then:
\begin{enumerate}
\renewcommand{\theenumi}{\alph{enumi}}
\renewcommand{\labelenumi}{\theenumi)}
\item The category $\Hhom(E,F)_{acc}$ of accessible functors
from $E$ to $F$ (\hyperref[I.def9.2]{9.2}) is a
$𝒰$-category. ({\rm NB} we
recall (\hyperref[I.sec9.0]{9.0}) that the given categories $E$, $F$
are assumed
to be $𝒰$-categories.)

\item Suppose that $F$ is stable under small filtered inductive limits.
For every full subcategory $C$ of $E$ equivalent
to a small category, with inclusion functor $i:C\rightarrow E$,
consider the corresponding functor
$$
i_{!}:\Hhom(C,F)\longrightarrow \Hhom(E,F)
$$
$\hyperref[I.prop5.1]{(5.1)}$.\pageoriginale A functor $f:E\rightarrow F$ is
accessible if and only if there exists a small
subcategory $C$ of $E$ such that $f$ belongs to the essential image
of the preceding functor $i_{!}$.
\end{enumerate}
\end{enonce*}

\noindent \emph{Proof}.
\begin{enumerate}
\renewcommand{\theenumi}{\alph{enumi}}
\renewcommand{\labelenumi}{\theenumi)}
\item It suffices to prove that for every cardinal $\Pi$ such that $E$
satisfies $L_{\Pi}$, the full subcategory $\Hhom(E,F)_{\Pi}$ of
$\Hhom(E,F)_{acc}$ is a $𝒰$-category. It is
clearly enough to verify this for cardinals of the form $2^c$, with
$c$ sufficiently large. But then this follows from \hyperref[I.cor9.19]{9.19}, since
($C=\Filt^{\Pi}(E)$ being essentially small) $\Hhom(C,F)$ is
obviously a $𝒰$-category.

\item By transitivity of the formation of the functors $i_{!}$, one
may restrict in the statement to subcategories $C$ of the form
$\Filt^{\Pi}(E)$, where $\Pi$ is as in \hyperref[I.cor9.17]{9.17}. One then sees
readily that the composite functor $\Hhom(\Filt^{\Pi}(E),F)\rightarrow
\Hhom(E,F)_{\Pi}\rightarrow \Hhom(E,F)$, where the first arrow is
quasi-inverse to \hyperlink{sga4.exp.I.equation.9.19.1}{(9.19.1)}, and the second is the inclusion, is
none other than the functor $i_{!}$, up to isomorphism. Thus assertion
b) follows from \hyperref[I.cor9.19]{9.19}.
\end{enumerate}

\begin{enonce*}[remark]{*Exercise 9.20.1}\etiquette[9.20.1]{I.exer9.20.1}
(The present exercise uses the notions of site and topos,
developed in Exposés II and IV.) Let $E$ be a
$𝒰$-topos, let $𝒱$ be a universe such that $𝒰\in
𝒱$, let $C$ be a small full generating subcategory of the
$𝒰$-topos $E$, and let $\Pi_\circ\in 𝒰$ be an infinite cardinal
such that $\Pi_\circ\in \card F\ell\,  C$. For every cardinal $\Pi\geq
\Pi_\circ$, $\Pi\in 𝒰$, let $\Filt^{\Pi}(E)$ be the
strictly full subcategory of $E$ consisting of the objects $X$ such
that there exists a strictly epimorphic family\pageoriginale
$(X_i\rightarrow X)_{i\in I}$ with target $X$, such that $\card I\leq
\Pi$ and $X_i\in \Ob C$ for every $i\in I$.
\begin{enumerate}
\renewcommand{\theenumi}{\alph{enumi}}
\renewcommand{\labelenumi}{\theenumi)}
\item Show that $(\Filt^{\Pi}(E))_{\Pi\geq \Pi_\circ}$ is a
cardinal filtration of the $𝒰$-category $E$, and that one can
choose $\Pi_\circ$ such that for every $\Pi\geq \Pi_\circ$, one has
the inclusions
$$
E_{\Pi}\subset \Filt^{\Pi}(E)\subset E_{\Pi^{\Pi_\circ}}{\quoi},
$$
where, for every cardinal $c$, $E_c$ denotes the strictly full
subcategory of $E$ consisting of the $c$-accessible objects.

\item Choosing a cardinal $\Pi\geq \Pi_\circ$ such that
$\Pi=\Pi^{\Pi_\circ}$ (for example, $\Pi$ of the form $2^c$, with
$c>\Pi_\circ$), and putting $C=\Filt^{\Pi}(E)$, show that there is an
equivalence of categories
$$
\Ind(C)_{\Pi}\xrightarrow{\approx} E{\quoi},
$$
(notation $\Ind(C)_{\Pi}$ of \hyperref[I.cor9.18]{9.18}).

\item Retain the notation of b), and let $𝒱$ be a universe such
that $𝒰\subset 𝒱$, let $%%
E^{\sim{\phantom{,}}}_{\underline{V}}$ be the $𝒱$-topos of
$𝒱$-sheaves on the site $E$ (endowed with its canonical
topology). Denote by $\Ind(C,𝒱)_{\Pi}$ the category of
$𝒱\text{-}\Ind$-objects of $C$ indexed by%%
 preordered index sets \emph{large relative to} $\Pi$. Show that there is
an equivalence of categories
$$
\Ind(C,𝒱)_{\Pi}\xrightarrow{\approx}%%
E^{\tilde{\phantom{,}}}_{𝒱}{\quoi}.
$$

\item Let $c\in 𝒱$ be a cardinal, and let $\Ind(E,𝒱)'_c$ be the
full subcategory of $\Ind(E,𝒱)$ consisting of the
$𝒱$-ind-objects of $E$ indexed by a preordered set that
is large relative to every cardinal $<c$. Taking $c=\card 𝒰$,
show\pageoriginale that there is an equivalence of categories
$$
\displaylines{\hfill \Ind(E,𝒱)'_{c}\xrightarrow{\approx}%%
E^{\tilde{\phantom{,}}}_{𝒱}{\quoi}.\hfill{\llap{\textbf{*}}}}
$$
\end{enumerate}
\end{enonce*}

\setcounter{subsection}{20}

% END INCLUDED SGA4_English_current_buildable_source_20260906/source/01/drafts/SGA4_I_LINES_6523_6659_English.texfrag


% BEGIN INCLUDED SGA4_English_current_buildable_source_20260906/source/01/drafts/SGA4_I_LINES_6660_6750_English.texfrag
\subsection{}\etiquette[9.21]{I.sec9.21}
The present section \hyperref[I.sec9.21]{9.21} develops technical
preliminaries for the proof of Theorem \hyperref[I.them9.22]{9.22}
below, which is the main result of the present section
\hyperref[I.sec9]{9}. Let
\hypertarget{sga4.exp.I.equation.9.21.1}{}\begin{equation}
p:E\rightarrow B\tag{9.21.1}
\end{equation}
be a fibered functor, where $B$ is a small category, and where the
inverse-image functors
$$
f^\ast:E_\beta\longrightarrow E_\alpha
$$
associated with the arrows $f:\alpha\rightarrow\beta$ of $B$ are
accessible (\hyperref[I.def9.2]{9.2}). In particular, the fiber
categories $E_\alpha$ $(\alpha\in\ob B)$ satisfy condition $L$
(\hyperref[I.def9.1]{9.1}), and therefore, since $B$ is small, there
exists a cardinal $\Pi'_\circ\in 𝒰$ such that all the categories
$E_\alpha$ satisfy condition $L_{\Pi'_\circ}$. Let
$\Pi_\circ\in 𝒰$ be an infinite cardinal $>\Pi'_\circ$, so that one
has
\hypertarget{sga4.exp.I.equation.9.21.2}{}\begin{equation}
\Pi_{\circ}>\Pi'_{\circ}, E_\alpha \mbox{ satisfies }
L_{\Pi'_{\circ}} \mbox{ for every } \alpha\in \ob B.\tag{9.21.2}
\end{equation}
Moreover, the hypotheses made immediately imply the existence of a
cardinal $c\in 𝒰$ satisfying the following conditions:
\hypertarget{sga4.exp.I.equation.9.21.3}{}\begin{equation}
\left\{
\begin{tabular}{rp{6cm}}
a) & $c$ is infinite,\\
b) & $c\geq \card F\ell\,B$,\\
c) & for every arrow $f:\alpha\rightarrow\beta$ in $B$, the functor
$f^\ast:E_\beta\rightarrow E_\alpha$ is $c$-accessible.
\end{tabular}\right.
\tag{9.21.3}
\end{equation}
Suppose\pageoriginale moreover that, for every $\alpha\in\ob B$, one
can find a full subcategory $C_\alpha$ of $E_\alpha$ satisfying the
following conditions:
\hypertarget{sga4.exp.I.equation.9.21.4}{}\begin{equation}
\left\{
\begin{tabular}{rp{6cm}}
a) & For every $X\in\ob C_\alpha$, $X$ is $c$-accessible in
$E_\alpha$ \hyperref[I.def9.3]{(9.3)}.\\
b) & For every $X\in\ob E_\alpha$, one can find an isomorphism
$X\simeq\varinjlim_I X_i$ in $E_\alpha$, where the $X_i$ belong to
$C_\alpha$ and $I$ is an ordered set large relative to $c$.
\end{tabular}\right.
\tag{9.21.4}
\end{equation}

Let $\Pi$ be a cardinal such that
\hypertarget{sga4.exp.I.equation.9.21.5}{}\begin{equation}
\Pi\geq \Pi_\circ\quav\mbox{ hence }
\Pi>\Pi'_\circ{\quoi},\tag{9.21.5}
\end{equation}
and, for every $\alpha\in\ob B$, let
\hypertarget{sga4.exp.I.equation.9.21.6}{}\begin{equation}
C^{\Pi}_{\alpha}\subset E_{\alpha}\tag{9.21.6}
\end{equation}
be formed by the objects that can be represented in the form
$\varinjlim_I X_i$, where $I$ is an ordered set large relative to
$\Pi'_\circ$, such that $\card I\leq\Pi$, and where the $X_i$ belong
to $C_\alpha$. It is then clear, by \hyperref[I.var7.5.2]{7.5.2}, that
$C^{\Pi}_{\alpha}$ is essentially small, \ie is equivalent to a small
category.

Let
\hypertarget{sga4.exp.I.equation.9.21.7}{}\begin{equation}
F=\Hhom_B(B,E)\tag{9.21.7}
\end{equation}
be the category of sections of $E$ over $B$, and let
\hypertarget{sga4.exp.I.equation.9.21.8}{}\begin{equation}
F^{\Pi}\subset F\tag{9.21.8}
\end{equation}
be the strictly full subcategory of $F$ formed by the sections
$X:\alpha\mapsto X(\alpha)$ such that for every $\alpha\in\ob B$ one
has
$$
X(\alpha)\in C^{\Pi}_{\alpha}{\quoi}.
$$
Since\pageoriginale the $C^{\Pi}_{\alpha}$ are essentially small, it
is clear that the same is true of the category $F^{\Pi}$. We shall
show that this category is generating, and more precisely:


% END INCLUDED SGA4_English_current_buildable_source_20260906/source/01/drafts/SGA4_I_LINES_6660_6750_English.texfrag


% BEGIN INCLUDED SGA4_English_current_buildable_source_20260906/source/01/drafts/SGA4_I_LINES_6751_6838_English.texfrag
\begin{enonce*}{Lemma 9.21.9}\etiquette[9.21.9]{I.lem9.21.9}
Under the preceding conditions and with the preceding notation, the
following statements hold:
\begin{enumerate}
\renewcommand{\theenumi}{\roman{enumi}}
\renewcommand{\labelenumi}{(\theenumi)}
\item Every object of $F$ is accessible (\hyperref[I.def9.3]{9.3}). If
$d$ is a cardinal $\geq\Pi'_\circ$, and if
$\alpha\mapsto X(\alpha)$ is an element of $F$ such that, for every
$\alpha$, $X(\alpha)$ is $d$-accessible in $E_\alpha$, then $X$ is
$d$-accessible in $F$.

\item Suppose that $\Pi\leq c$, or that $\Pi'_\circ$ is finite (\ie
the $E_\alpha$ are stable under small filtered $\varinjlim$). Then
every object $X$ of $F$ is isomorphic to an object of the form
$\varinjlim_I X_i$, where the $X_i$ belong to $F^{\Pi}$ and $I$ is
large relative to $\Pi$.
\end{enumerate}
\end{enonce*}

To prove (i), note that by \hyperlink{sga4.exp.I.equation.9.21.4}{(9.21.4)} and
\hyperref[I.cor9.9]{9.9}, every object of $E_\alpha$ is accessible. On
the other hand, $F$ satisfies condition $L_{\Pi'_\circ}$ by the


\begin{enonce*}{Lemma 9.21.10}\etiquette[9.21.10]{I.lem9.21.10}
Let $p:E\rightarrow B$ be a fibered functor, let
$F=\Hhom_B(B,E)$, let $I$ be a category, and let $i\mapsto X_i$ be a
functor from $I$ to $F$. For $\varinjlim_I X_i$ to be representable in
$F$, it suffices that, for every $\alpha\in\ob B$,
$\varinjlim_I X_i(\alpha)$ be representable in the fiber category
$E_\alpha$; when this is so, $\varinjlim_I X_i$ is ``computed
argument by argument''. In particular, if the fiber categories satisfy
condition $L_{\Pi'_\circ}$ (where $\Pi'_\circ$ is a given cardinal),
then so does $F$.
\end{enonce*}

We\pageoriginale leave the details of the (easy) proof of
\hyperref[I.lem9.21.10]{9.21.10} to the reader, merely noting that it
is convenient to use the following result, whose proof is immediate:

\begin{enonce*}{Corollary 9.21.10.1}\etiquette[9.21.10.1]{I.cor9.21.10.1}
With the hypotheses and notation of \hyperref[I.lem9.21.10]{9.21.10}
for $p:E\rightarrow B$ and for $I$, for every $\alpha\in\ob B$, the
inclusion functor $E_\alpha\rightarrow E$ commutes with inductive
limits of type $I$.
\end{enonce*}

Returning now to the general conditions of
\hyperref[I.lem9.21.9]{9.21.9}, let $X$ be an object of $F$. Since, for
every $\alpha\in\ob B$, $X(\alpha)$ is accessible in $E_\alpha$,
there exists a cardinal $d\in 𝒰$ such that, for every $\alpha$,
$X(\alpha)$ is $d$-accessible in $E_\alpha$. We may choose
$d\geq\Pi'_\circ$, so that $F$ satisfies $L_d$ by
\hyperref[I.lem9.21.10]{9.21.10}. Using
\hyperref[I.lem9.21.10]{9.21.10} once more to compute the inductive
limits $\varinjlim_IY_i$ in $F$, with $I$ large relative to $d$, one
immediately sees that $X$ is $d$-accessible. This proves (i).

We shall now prove \hyperref[I.lem9.21.9]{9.21.9} (ii) in several
steps (\hyperlink{sga4.exp.I.equation.9.21.11}{(9.21.11)} to (\hyperref[I.lem9.21.16]{9.21.16})).

Let then
$$
X:\alpha\longmapsto X(\alpha)
$$
be an object of $F$. By \hyperlink{sga4.exp.I.equation.9.21.4}{(9.21.4 b)}), for every $\alpha\in\ob B$ one
can find an ordered set $I_\alpha$ large relative to $c$, and an
isomorphism
$X(\alpha)=\varinjlim_{i\in I_\alpha}X(\alpha)_i$, where
$i\mapsto X(\alpha)_i$ is an inductive system of type $I_\alpha$ in
$E_\alpha$, with the $X(\alpha)_i$ in $C_\alpha$. Consider the product
ordered set $I=\Pi_\alpha I_\alpha$; it is clear that it is large
relative to $c$, and that the preceding inductive
systems\pageoriginale give rise to inductive systems in the $E_\alpha$,
\hypertarget{sga4.exp.I.equation.9.21.11}{}\begin{equation}
i\longmapsto X(\alpha)_i\in \ob C_{\alpha}, i\in \ob
I{\quoi},\tag{9.21.11}
\end{equation}
indexed by the same ordered set $I$, and to isomorphisms
\hypertarget{sga4.exp.I.equation.9.21.12}{}\begin{equation}
X(\alpha)\xrightarrow{\sim}\varinjlim_I X(\alpha)_i\mbox{ in }
E_{\alpha}{\quoi},\tag{9.21.12}
\end{equation}
for every $\alpha\in\ob B$. In what follows we suppose that data
\hyperlink{sga4.exp.I.equation.9.21.11}{(9.21.11)} and \hyperlink{sga4.exp.I.equation.9.21.12}{(9.21.12)} have been fixed.


% END INCLUDED SGA4_English_current_buildable_source_20260906/source/01/drafts/SGA4_I_LINES_6751_6838_English.texfrag


% BEGIN INCLUDED SGA4_English_current_buildable_source_20260906/source/01/drafts/SGA4_I_LINES_6839_6963_English.texfrag
\begin{enonce*}{Lemma 9.21.13}\etiquette[9.21.13]{I.lem9.21.13}
Under the preceding conditions, one can find a map
$\varphi:I\rightarrow I$ such that $\varphi(i)\geq i$ for every
$i\in I$, and a map $(i,f)\mapsto\lambda(i,f)$ from $I\times %%
\mathrm{F}{}l\,B$ to $%%
\mathrm{F}{}l\,E$, satisfying the following conditions:
\begin{enumerate}
\renewcommand{\theenumi}{\alph{enumi}}
\renewcommand{\labelenumi}{\theenumi)}
\item For $i\in I$ and $(f:\alpha\rightarrow\beta)\in F\ell B$,
$\lambda(i,f)$ is an $f$-morphism from $X(\alpha)_i$ to
$X(\beta)_{\varphi(i)}$, making the diagram
\[
\xymatrix@C=3cm@R=.8cm{
 X(\alpha)\ar[r]^{X(f)} & X(\beta)\\
           &  X(\beta)_{\varphi(i)}\ar[u]\\
X(\alpha)_{i} \ar[uu]\ar[ru]^{\lambda(i,f)}&
                  \\
\alpha \ar[r]^{f}& \beta\quad , }
\]
commutative, where the vertical arrows are the canonical morphisms
deduced from \hyperlink{sga4.exp.I.equation.9.21.12}{(9.21.12)}.

\item For\pageoriginale every $i\in I$, and every
$(f:\alpha\rightarrow\beta)\in %%
\mathrm{F}{}l\,B$, the diagram
\[
\xymatrix@C=3cm@R=.9cm@L=6pt{
                      & X(\beta)_{\varphi^{2}(i)}\\
  X(\alpha)_\varphi(i) \ar[ru]^{\lambda(\varphi(i),f)}  & X(\beta)_{\varphi(i)} \ar[u]\\
X(\alpha)_i  \ar[ru]_{\lambda(i,f)} \ar[u]         \\
\alpha \ar[r]^{f} & \beta\quad . }
\]
is commutative.

\item For every pair of consecutive arrows
$\alpha\xrightarrow{f}\beta\xrightarrow{g}\gamma$ of $B$, the
following diagram is commutative:
\[
\xymatrix@C=2cm@L=6pt@R=.5cm{
X(\alpha) \ar[r]^{X(f)}& X(\beta) \ar[r]^{X(g)}            & X(\gamma)\\
          &                      & X(\gamma)_{\varphi^{2}(i)}\ar[u]\\
          & X(\beta)_{\varphi(i)} \ar[uu]\ar[ru]^{\lambda(\varphi(i),g)} & X(\gamma)_{\varphi(i)}\ar[u]\\
X(\alpha)_i \ar[uuu]\ar[ru]^{\lambda(i,f)} \ar[rru]_{\lambda(i,gf)}&
 \\
\alpha \ar[r]^{f}& \beta \ar[r]^{g}& \gamma\quad , }
\]
where the vertical arrows are those deduced from \hyperlink{sga4.exp.I.equation.9.21.12}{(9.21.12)}.
\end{enumerate}
\end{enonce*}

One\pageoriginale should moreover note that the commutativity of the
two upper trapezoids in c) is already contained in a), so c) in fact
asserts the commutativity of the lower triangle of the diagram under
consideration.

Let us prove \hyperref[I.lem9.21.13]{9.21.13}. Let $i\in I$, and let
$f:\alpha\rightarrow\beta$ be an arrow in $B$. Consider the composite
$X(\alpha)_i\rightarrow X(\alpha)\xrightarrow{X(f)}X(\beta)\simeq
\varinjlim_jX(\beta)_j$; it defines a morphism in $E$:
$$
X(\alpha)_i\rightarrow f^\ast(X(\beta))\simeq
f^\ast(\varinjlim_jX(\beta)_j)\simeq
\varinjlim_jf^\ast(X(\beta)_j){\quoi},
$$
where the last equality follows from the fact that $f^\ast$ is
$c$-accessible \hyperlink{sga4.exp.I.equation.9.21.3}{(9.21.3 c)}%%
 and that $I$ is large relative to $c$. By \hyperlink{sga4.exp.I.equation.9.21.4}{(9.21.4 a)}), since
$X(\alpha)_i\in\ob C_\alpha$, the morphism under consideration
factors through some $f^\ast(X(\beta)_j)$, where $j$ \textit{a priori}
depends on $i$ and on $f\in %%
\mathrm{F}{}l\,(B)$. But by \hyperlink{sga4.exp.I.equation.9.21.3}{(9.21.3 b)}), and since $I$ is large
relative to $c$, one may choose $j$ independently of $f$, say
$j=\varphi(i)$. Since $I$ is filtered, one may suppose
$\varphi(i)\geq i$. One thus obtains, for $i\in I$ and
$f\in %%
\mathrm{F}{}l\,B$, a morphism
$$
X(\alpha)_i\longrightarrow f^\ast(X(\beta)_{\varphi(i)}){\quoi},
$$
or, equivalently, an $f$-morphism
$$
\lambda(i,f):X(\alpha)_i\longrightarrow
X(\beta)_{\varphi(i)}{\quoi},
$$
which by construction makes the diagram
\[\vcenter{
\xymatrix@C=3cm{
X(\alpha)\ar[r]^{X(f)} & X(\beta)\\
          & X(\beta)_{\varphi(i)} \ar[u]\\
X(\alpha)_i \ar[uu] \ar[ru]^{\lambda(i,f)}&\\
\alpha \ar[r]^{f} & \beta \quad . }}\leqno{D(f)}
\]
commutative. Now\pageoriginale, for given $i,f,g$, consider the lower
triangle of the diagram in \hyperref[I.lem9.21.13]{9.21.13} c); it is not clear that it is
commutative, but the two composites
$X(\alpha)_i\rightrightarrows X(\gamma)_{\varphi^2(i)}$ become equal
after composition with
$X(\gamma)_{\varphi^2(i)}\rightarrow X(\gamma)$, as follows from the
commutativity of the two upper trapezoids and of the outer trapezoid,
which are respectively the diagrams $D(f)$, $D(g)$, and $D(gf)$.
Thus, since $X(\alpha)_i$ is $c$-accessible \hyperlink{sga4.exp.I.equation.9.21.4}{(9.21.4 a)}) and
$X(\gamma)\simeq\varinjlim_{j\in I}X(\gamma)_j$, with $I$ large
relative to $c$, it follows that one can find an element
$j\geq\varphi^2(i)$ of $I$ such that the two arrows under consideration
become equal after composition with
$X(\gamma)_{\varphi^2(i)}\rightarrow X(\gamma)_j$. A priori, $j$
depends on $i,f,g$. But, for fixed $i$, the set of possible pairs
$f,g$ has cardinality $\leq c^2=c$, by \hyperlink{sga4.exp.I.equation.9.21.3}{(9.21.3 a)} and b)), and hence,
since $I$ is large relative to $c$, one may choose $j$ independently of
$f$ and $g$, say
$j=\varphi'(i)\geq\varphi^2(i)\geq\varphi(i)$. For every $i\in I$
and $f:\alpha\rightarrow\beta$, let then
$$
\lambda'(i,f):X(\alpha)_i\longrightarrow X(\beta)_{\varphi'(i)}
$$
be the composite $f$-morphism
$X(\alpha)_i\rightarrow X(\beta)_{\varphi(i)}\rightarrow
X(\beta)_{\varphi'(i)}$, where the second arrow is the transition
morphism. It is then immediate from the construction that
$(\varphi',\lambda')$ satisfies conditions \hyperref[I.lem9.21.13]{9.21.13} a) and c) for
$(\varphi,\lambda)$. Proceeding in the same way for condition b), one
sees that $\varphi'$ may be chosen so that this condition is also
satisfied for $(\varphi',\lambda')$. This completes the proof of
\hyperref[I.lem9.21.13]{(9.21.13)}.

\setcounter{subsection}{21} \setcounter{subsubsection}{13}

% END INCLUDED SGA4_English_current_buildable_source_20260906/source/01/drafts/SGA4_I_LINES_6839_6963_English.texfrag


% BEGIN INCLUDED SGA4_English_current_buildable_source_20260906/source/01/drafts/SGA4_I_LINES_6964_7051_English.texfrag
\subsubsection{}\etiquette[9.21.14]{I.sec9.21.14}
Now\pageoriginale let
$$
J\subset I
$$
be a subset of $I$ satisfying the following conditions:
\begin{enumerate}
\renewcommand{\theenumi}{\alph{enumi}}
\renewcommand{\labelenumi}{\theenumi)}
\item $J$ is filtered,

\item for every $j\in J$, one has $\varphi(j)\in J$,

\item For every $\alpha\in\ob B$,
$X_J(\alpha)=\varinjlim_{i\in J}X(\alpha)_i$ is representable in
$E_\alpha$.
\end{enumerate}
Conditions a) and c) are satisfied in particular if $J$ is large
relative to $\Pi'_\circ$ \hyperlink{sga4.exp.I.equation.9.21.2}{(9.21.2)}. It then follows from
\hyperref[I.cor9.21.10.1]{9.21.10.1} that, for every
$\alpha\in\ob B$, $X_J(\alpha)$ is the inductive limit
$\varinjlim_{i\in J}X(\alpha)_i$ in $E$. Now, for an arrow
$f:\alpha\rightarrow\beta$ of $B$, the arrows $\lambda(i,f)$, with
$i$ varying in $J$, define, by (\hyperref[I.lem9.21.13]{9.21.13} b)),
a morphism of ind-objects from $(X(\alpha)_i)_{i\in J}$ to
$(X(\beta)_i)_{i\in J}$, whence a morphism
$$
X_J(f):X_J(\alpha)\rightarrow X_J(\beta)
$$
on the inductive limits, which is manifestly an $f$-morphism. Using
(\hyperref[I.lem9.21.13]{9.21.13} (c)), one obtains the transitivity
relations
$$
X_J(gf)=X_J(g)X_J(f){\quoi},
$$
so that a section $X_J\in\ob F$ of $E$ over $B$ has been defined.
Finally, (\hyperref[I.lem9.21.13]{9.21.13} a)) shows us that the
canonical morphisms
$$
X_J(\alpha)\longrightarrow X(\alpha)\quav \alpha\in \ob B{\quoi},
$$
are functorial in $\alpha$, so there is a canonical morphism
$$
u_J:X_J\longrightarrow X{\quoi}.
$$
On the\pageoriginale other hand, if
$$
J'\supset J
$$
is another subset of $I$ satisfying conditions a), b), c) above, one
obtains a canonical morphism
$$
u_{J',J}:X_J\longrightarrow X_{J'}{\quoi},
$$
so that the $X_J$ form an inductive system in $F$, parametrized by the
set (ordered by inclusion) of subsets $J$ of $I$ satisfying the
conditions under consideration. Finally, the morphisms $u_J$ above
define a morphism from this inductive system to (the constant inductive
system defined by) $X$.


\setcounter{subsection}{21} \setcounter{subsubsection}{14}
\subsubsection{}\etiquette[9.21.15]{I.sec9.21.15}
Now let $K$ be a set of subsets $J$ of $I$ satisfying conditions a),
b), c) considered in \hyperref[I.sec9.21.14]{9.21.14}, and suppose
that $K$ is filtered and has union $I$. It is then clear that one has
$$
\varinjlim_{J\in K}X_J\xrightarrow{\sim}X{\quoi},
$$
using \hyperref[I.cor9.21.10.1]{9.21.10.1}, which reduces us to
checking that there is an isomorphism argument by argument.

Take, for example, for $K$ the set of all subsets $J$ of $I$ that are
large relative to $\Pi'_\circ$, stable under $\varphi$, and such that
$\card J\leq\Pi$. Then, by the definition \hyperlink{sga4.exp.I.equation.9.21.6}{(9.21.6)} of
$C^{\Pi}_\alpha$, for every $\alpha\in\ob B$ one has
$X_J(\alpha)\in\ob C^{\Pi}_\alpha$, hence $X_J\in\ob F^{\Pi}$.
Consequently, \hyperref[I.lem9.21.9]{9.21.9} (ii) will be proved if we
establish that $K$ is large relative to $\Pi$ (and hence filtered) and
has union $I$. For this, it will evidently suffice to prove that every
subset $S$ of $I$ such that $\card S\leq\Pi$ is contained in a
$J\in K$. This will indeed follow\pageoriginale from the preliminary
hypothesis stated in \hyperref[I.lem9.21.9]{9.21.9} (ii), and from the


% END INCLUDED SGA4_English_current_buildable_source_20260906/source/01/drafts/SGA4_I_LINES_6964_7051_English.texfrag


% BEGIN INCLUDED SGA4_English_current_buildable_source_20260906/source/01/drafts/SGA4_I_LINES_7052_7130_English.texfrag
\begin{enonce*}{Lemma 9.21.16}\etiquette[9.21.16]{I.lem9.21.16}
Let $I$ be an ordered set, let $\Pi$ be an infinite cardinal, and let
$\varphi:I\rightarrow I$ be a map from $I$ to itself.
\begin{enumerate}
\renewcommand{\theenumi}{\alph{enumi}}
\renewcommand{\labelenumi}{\theenumi)}
\item Suppose $I$ is filtered. Then, for every subset $S$ of $I$ such
that $\card S\leq\Pi$, there exists a filtered subset $J\supset S$ of
$I$ such that $\varphi(J)\subset J$ and $\card(J)\leq\Pi$.

\item Let $\Pi'_\circ$ be a cardinal $<\Pi$, and suppose that $I$ is
large relative to $\Pi$. Then, for every subset $S$ of $I$ such that
$\card S\leq\Pi$, there exists a subset $J\supset S$ of $I$ large
relative to $\Pi'_\circ$, such that $\varphi(J)\subset J$ and
$\card J\leq\Pi$.
\end{enumerate}
\end{enonce*}

Let us prove b), for example (the proof of a) being similar and
simpler). Let $P$ be the set of subsets of $I$ of cardinality
$\leq\Pi$, and let $T\mapsto i_T:P\rightarrow I$ be a map such that,
for $T\in P$, $i_T$ is an upper bound of $T$ in $I$; the existence of
this map merely expresses the hypothesis that $I$ is large relative to
$\Pi$. Let $A$ be a well-ordered set such that $\card A=\Pi$, and such
that, for every $a\in A$, the set of $a'\leq a$ has cardinality
$<\Pi$. It follows in particular, since $\Pi'_\circ<\Pi$, that $A$
\emph{is large relative to} $\Pi'_\circ$. Define, by transfinite
recursion, maps
$$
S\longmapsto S_a:P\longrightarrow P\quad(a\in A){\quoi},
$$
by the following formulas:

$S_\circ=S$, $S_{a+1}=S_a\cup\varphi(S_a)\cup\{i_{S_a}\}$,
$S_a=\bigcup_{a'<a}S_{a'}$ if $a$ is a limit ordinal.

Then put,\pageoriginale for $S\in P$,
$$
J(S)=\bigcup_{a\in A}S_a{\quoi}.
$$
It is then clear that $J(S)\supset S$, that $J(S)$ is stable under
$\varphi$, that $\card J(S)\leq\Pi$ (since $\card A=\Pi$), and finally
that $J(S)$ is large relative to $\Pi'_\circ$ (using the fact that $A$
is large relative to $\Pi'_\circ$).

This proves \hyperref[I.lem9.21.16]{9.21.16} and completes the proof of
\hyperref[I.lem9.21.9]{9.21.9}.

We also point out a variant of \hyperref[I.lem9.21.9]{9.21.9}:

\begin{enonce*}{Lemma 9.21.17}\etiquette[9.21.17]{I.lem9.21.17}
The notation is that of \hyperref[I.lem9.21.9]{9.21.9}. Let $F'$ denote
the full subcategory $\Hhomcart_B(B,E)$ of $F=\Hhom_B(B,E)$ formed by
the cartesian sections of $E$ over $B$. Then:
\begin{enumerate}
\renewcommand{\theenumi}{\roman{enumi}}
\renewcommand{\labelenumi}{(\theenumi)}
\item Every object of $F'$ is accessible. More precisely, let
$\alpha\mapsto X(\alpha)$ be an element of $F'$, and let $d$ be a
cardinal such that $d\geq\Pi'_\circ$, $d\geq c$, and such that, for
every $\alpha\in\ob B$, $X(\alpha)$ is a $d$-accessible object of
$E_\alpha$. Then $X$ is a $d$-accessible object of $F'$.

\item Suppose $\Pi\leq c$ and that the functors
$f^\ast:E_\beta\rightarrow E_\alpha$ are
$\Pi'_\circ$-accessible, or that the categories $E_\alpha$ are stable
under small filtered inductive limits and the functors
$f^\ast:E_\beta\rightarrow E_\alpha$ commute with these inductive
limits. Then every object $X$ of $F'$ is isomorphic to an object of the
form $\varinjlim_I X_i$, where, for every $i\in I$, $X_i$ is an object
of the full subcategory ${F'}^{\Pi}=F'\cap F^{\Pi}$ of $F'$, and where
$I$ is large relative to $\Pi$.
\end{enumerate}
\end{enonce*}

The\pageoriginale proof being entirely analogous to that of
\hyperref[I.lem9.21.9]{9.21.9}, we shall merely indicate the points at
which the latter must be modified. First note:


% END INCLUDED SGA4_English_current_buildable_source_20260906/source/01/drafts/SGA4_I_LINES_7052_7130_English.texfrag


% BEGIN INCLUDED SGA4_English_current_buildable_source_20260906/source/01/drafts/SGA4_I_LINES_7131_7193_English.texfrag
\begin{enonce*}{Lemma 9.21.18}\etiquette[9.21.18]{I.lem9.21.18}
The statement of \hyperref[I.lem9.21.10]{9.21.10} remains valid when
$F=\Hhom_B(B,E)$ is replaced there by $F'=\Hhomcart_B(B,E)$, provided
that the inverse image functors $f^\ast:E_{\beta}\rightarrow
E_{\alpha}$ commute with inductive limits of type $I$.
\end{enonce*}

Thus, under the conditions of \hyperref[I.lem9.21.17]{9.21.17}, if $d$ is a cardinal such that
$d\geq c$ and $d\geq \Pi'_{\circ}$, it follows from the above and
from hypothesis \hyperlink{sga4.exp.I.equation.9.21.3}{(9.21.3 c)}) that, for every
ordered set $I$ large
relative to $d$, inductive limits of type $I$ are representable
in $F'$ and are computed argument by argument, whence the conclusion
\hyperref[I.lem9.21.17]{9.21.17} (i) follows by an immediate argument.

To prove (ii), one must supplement \hyperref[I.lem9.21.13]{9.21.13}:

\begin{enonce*}{Lemma 9.21.19}\etiquette[9.21.19]{I.lem9.21.19}
Under the conditions of \hyperref[I.lem9.21.13]{9.21.13}, suppose that $X\in \ob F'$ \ie
that $X$ is a cartesian section of $E$ over $B$. Then the conclusion can be
strengthened by asserting that there exists a function
$μ$ which assigns to every $i\in I$ and every arrow $f:\alpha\rightarrow
\beta$ of $B$ a morphism $μ(i,f):f^\ast
X(\beta)_i\rightarrow X(\alpha)_{\varphi(i)}$ in $E_{\alpha}$, such
that:

\begin{enumerate}
\renewcommand{\theenumi}{\alph{enumi}}
\renewcommand{\labelenumi}{\theenumi)}
\setcounter{enumi}{3}
\item For every $i\in I$ and every arrow $f:\alpha\rightarrow \beta$
of $B$, the following two diagrams are commutative:

\begin{displaymath}
\xymatrix@L=5pt@R=1.25cm@C=2.2cm{
                       & f^{\ast}X(\beta)_{\varphi^{2}(i)}\\
X(\alpha)_{\varphi(i)}\ar[ru]^{\lambda(\varphi (i),f)} &\\
                       & f^{\ast}X(\beta)_i \ar[lu]^{\mu(i,f)}\ar[uu]
} \qquad\qquad \xymatrix@L=5pt@R=1.35cm@C=2cm{
X(\alpha)_{\varphi^2(i)} & \\
                         & f^{\ast}X(\alpha)_{\varphi(i)} \ar[lu]_{\mu(\varphi(i),f)}\\
x(\alpha)_i \ar[uu]\ar[ru]_{\lambda(i,j)}             &. }
\end{displaymath}
\end{enumerate}
\end{enonce*}

To\pageoriginale prove this, one proceeds as in \hyperref[I.lem9.21.13]{9.21.13}, by
considering the composite morphism $f^\ast(X(\beta)_i)\rightarrow
f^\ast(X(\beta))\xrightarrow{X(f)^{-1}}X(\alpha)$, where (as already
in the formulation of condition d)) we identify notationally
an $f$-morphism $R\rightarrow S$ of $E$ with the corresponding morphism
$R\rightarrow f^\ast(S)$ of $E_{\alpha}$. Since
$X(\alpha)=\varinjlim_j X(\alpha)_j$, and $I$ is large relative to $c$, one
can factor the preceding morphism through one of the $X(\alpha)_j$, where
$j$ \textit{a priori} depends on $i$ \emph{and} on $f$, but can be
chosen independently of $f$, say $\psi(i)$. After enlarging the
function $\varphi$ of \hyperref[I.lem9.21.13]{9.21.13}, one may suppose that $\psi=\varphi$.
Proceeding as for conditions b) and c) of
\hyperref[I.lem9.21.13]{9.21.13}, one sees
that, after further enlarging the map $\varphi$, one may
suppose that the two diagrams of \hyperref[I.lem9.21.19]{9.21.19} d) are commutative.

\setcounter{subsection}{21} \setcounter{subsubsection}{19}

% END INCLUDED SGA4_English_current_buildable_source_20260906/source/01/drafts/SGA4_I_LINES_7131_7193_English.texfrag


% BEGIN INCLUDED SGA4_English_current_buildable_source_20260906/source/01/drafts/SGA4_I_LINES_7194_7240_English.texfrag
\subsubsection{}\etiquette[9.21.20]{I.sec9.21.20}
With the map $\varphi$ chosen as in
\hyperref[I.lem9.21.19]{9.21.19}, let us repeat
the argument of \hyperref[I.sec9.21.14]{9.21.14}, except that we must assume that $J$
satisfies, in addition to the stated conditions a), b), c), the following condition:
\begin{enumerate}
\renewcommand{\theenumi}{\alph{enumi}}
\renewcommand{\labelenumi}{\theenumi)}
\setcounter{enumi}{3}
\item For every arrow $f:\alpha\rightarrow \beta$ in $B$, the
functor $f^\ast:E_{\beta}\rightarrow E_{\alpha}$ commutes with
inductive limits of type $J$.
\end{enumerate}

We\pageoriginale claim that, under this additional condition,
the section $X_J$ of $E$ over $B$ is cartesian, \ie for every
arrow $f:\alpha\rightarrow \beta$ of $B$, the morphism
\begin{equation}
X_J(\alpha)\longrightarrow f^\ast X_J(\beta)\tag{$\ast$}
\end{equation}
is an isomorphism. Indeed, by condition d)
above, the target of the morphism under consideration is identified with
$\varinjlim_{i\in J}f^\ast(X(\beta)_i)$, and the morphism is obtained
by taking $\varinjlim_J$ of the morphism of ind-objects
in $E_{\alpha}$
$$
(X(\alpha)_i)_{i\in J}\longrightarrow (f^\ast)(X(\beta)_i)_{i\in
J}{\quoi},
$$
induced by the morphisms $\lambda(i,f):X(\alpha)_i\rightarrow
f^\ast(X(\beta)_{\varphi(i)})$ for $i\in J$. Now the conditions
set out in \hyperref[I.lem9.21.19]{9.21.19} ensure that
the preceding homomorphism
of ind-objects is in fact an isomorphism, an inverse being obtained
from the homomorphism induced by the system of $\mu(i,f)$, for $i\in J$.
This proves our assertion that $(\ast)$ is an %%
isomorphism, and hence that $X\in \ob F'$.

\subsubsection{}\etiquette[9.21.21]{I.sec9.21.21}
We shall now assume that the functors
$f^\ast:E_{\beta}\rightarrow E_{\alpha}$ are
$\Pi'_{\circ}$-\emph{accessible.} Then the conditions on $J$
considered in \hyperref[I.sec9.21.20]{9.21.20} are satisfied if $J$ is large relative to
$\Pi'_{\circ}$, stable under $\varphi$, and such that $\card J\leq \Pi$,
and we complete the proof of \hyperref[I.lem9.21.17]{9.21.17} just as in
\hyperref[I.sec9.21.15]{9.21.15}.

% END INCLUDED SGA4_English_current_buildable_source_20260906/source/01/drafts/SGA4_I_LINES_7194_7240_English.texfrag


% BEGIN INCLUDED SGA4_English_current_buildable_source_20260906/source/01/drafts/SGA4_I_LINES_7241_7301_English.texfrag
\begin{enonce*}{Theorem 9.22}\etiquette[9.22]{I.them9.22}
Let\pageoriginale $p:E\rightarrow B$ be a fibered functor, where $B$
is a category essentially equivalent to a small
category, and where for every arrow $f:\alpha\rightarrow \beta$ in
$B$ the inverse image functor $f^\ast:E_{\beta}\rightarrow
E_{\alpha}$ is accessible (\hyperref[I.def9.2]{9.2}). Suppose moreover that for every
$\alpha\in \ob E$, the fiber category $E_{\alpha}$ admits a
cardinal filtration (\hyperref[I.def9.12]{9.12}) and that every object of $E_{\alpha}$ is
accessible (\hyperref[I.def9.3]{9.3}). Then each of the categories
$F=\Hhom_B(B,E)$ and
$F'=\Hhomcart_B(B,E)$ has a small full subcategory
generating by strict epimorphisms (\hyperref[I.def7.1]{7.1}), and
every object of either category is accessible.
\end{enonce*}

It is immediate that the statement is not essentially changed when
$B$ is replaced by a full subcategory $B_{\circ}$ such that the
inclusion functor $B_{\circ}\rightarrow B$ is an equivalence,
and $E$ by $E_{\circ}=E\times_B B_{\circ}$. This allows us to
assume that $B$ is a small category.

For every $\alpha\in \ob B$, let $(\Filt^{\Pi}(E))_{\Pi\geq
\Pi_{\alpha}}$ be a cardinal filtration of $E_{\alpha}$. Let
$c_{\circ}\in 𝒰$ be a cardinal satisfying the following
conditions:
\hypertarget{sga4.exp.I.equation.9.22.1}{}\begin{equation}
\begin{cases}
& c_{\circ}\geq {\displaystyle{\mathop{\Sup}_{\alpha\in ob
B}}}\Pi_{\alpha}{\quav} c_{\circ}\geq \card %%
\mathrm{F}{}l\, B;\\
& \mbox{for every arrow $f:\alpha\rightarrow \beta$ of $B$,
$f^\ast:E_{\beta}\rightarrow E_{\alpha}$ is
$c_{\circ}$-accessible;}\\
& \mbox{for every $\alpha\in \ob B$, the objects of
$\Filt^{\Pi_{\alpha}}(E_{\alpha})$ are $c_\circ$-accessible.}
\end{cases}\tag{9.22.1}
\end{equation}
Set
$$
c=2^{c_{\circ}}{\quoi},
$$
so that $c>c_{\circ}$, and, for every $\alpha\in \ob B$, let
$$
C_{\alpha}=\Filt^c(E_{\alpha}){\quoi}.
$$
We\pageoriginale claim that the preliminary conditions of
\hyperref[I.lem9.21.9]{9.21.9} and
\hyperref[I.lem9.21.17]{9.21.17} are satisfied on taking
$\Pi=\Pi_{\circ}=c$. This is clear
for \hyperlink{sga4.exp.I.equation.9.21.2}{(9.21.2)} and \hyperlink{sga4.exp.I.equation.9.21.3}{(9.21.3)}. For \hyperlink{sga4.exp.I.equation.9.21.4}{(9.21.4 a)}),
this follows from \hyperref[I.cor9.17]{9.17},
and \hyperlink{sga4.exp.I.equation.9.21.4}{(9.21.4 b)}) follows from condition \hyperref[I.def9.12]{9.12} c)
for a cardinal
filtration. Let us also note that condition \hyperref[I.def9.12]{9.12} b) on
cardinal filtrations implies that for every $\alpha\in \ob B$, we
have $C^{\Pi}_{\alpha}=C_{\alpha}=\Filt^C(E_{\alpha})$, where
$C^{\Pi}_{\alpha}$ is defined in \hyperlink{sga4.exp.I.equation.9.21.6}{(9.21.6)}. Consequently,
\hyperref[I.them9.22]{9.22} follows
from \hyperref[I.lem9.21.9]{9.21.9} and \hyperref[I.lem9.21.17]{9.21.17}, and more
precisely, we have proved the


% END INCLUDED SGA4_English_current_buildable_source_20260906/source/01/drafts/SGA4_I_LINES_7241_7301_English.texfrag


% BEGIN INCLUDED SGA4_English_current_buildable_source_20260906/source/01/drafts/SGA4_I_LINES_7302_7379_English.texfrag
\begin{enonce*}{Corollary 9.23}\etiquette[9.23]{I.cor9.23}
Under the conditions of \hyperref[I.them9.22]{9.22}, if
$c_\circ\in 𝒰$ is a cardinal satisfying conditions \hyperlink{sga4.exp.I.equation.9.22.1}{(9.22.1)}, and if
$c=2^{c_\circ}$, then the subcategory $F^c$ of $F$ (respectively
${F'}^c$ of $F'$) formed by the $X$ such that
$X(\alpha)\in\ob\Filt^c(E_\alpha)$ for every $\alpha\in\ob B$ is
essentially small and generating by strict epimorphisms; more
precisely, every object $X$ of $F$ (respectively $F'$) is isomorphic to
an object of the form $\varinjlim_I X_i$, where the $X_i$ belong to
$F^c$ (respectively ${F'}^c$), and where $I$ is an ordered set large
relative to $c$.
\end{enonce*}

Under slightly stronger hypotheses in \hyperref[I.them9.22]{9.22}, one
can moreover make \hyperref[I.them9.22]{9.22} and
\hyperref[I.cor9.23]{9.23} considerably more precise:

\begin{enonce*}{Corollary 9.24}\etiquette[9.24]{I.cor9.24}
Under the conditions of \hyperref[I.them9.22]{9.22}, suppose that each
of the categories $E_\alpha$ $(\alpha\in\ob B)$ is stable under small
filtered inductive limits, and that, for every $\Pi\geq\Pi_\alpha$,
$\Filt^{\Pi}(E_\alpha)$ is stable under filtered inductive limits
indexed by filtered ordered sets $I$ such that $\card I\leq\Pi$ (which
slightly strengthens condition \hyperref[I.def9.12]{9.12} {\rm b)} for
cardinal\pageoriginale filtrations). In the case where $F'$ is being
considered, suppose moreover that, for every arrow
$f:\alpha\rightarrow\beta$ of $B$, the functor
$f^\ast:E_\beta\rightarrow E_\alpha$ commutes with small filtered
inductive limits. Finally, let $c_\circ\in 𝒰$ be a cardinal satisfying
conditions \hyperlink{sga4.exp.I.equation.9.22.1}{(9.22.1)}, and, for every cardinal
$\Pi\geq c=2^{c_\circ}$, consider the strictly full subcategory
$F^{\Pi}$ (respectively ${F'}^{\Pi}$) of $F$ (respectively $F'$)
formed by the $X$ such that
$X(\alpha)\in\Filt^{\Pi}(E_\alpha)$ for every $\alpha\in\ob B$. Then
the subcategories under consideration define a cardinal filtration
(\hyperref[I.def9.12]{9.12}) of $F$ (respectively $F'$).
\end{enonce*}

Conditions a), b), c) of \hyperref[I.def9.12]{9.12} must be checked.
Conditions a) and b) follow from the analogous conditions for the given
cardinal filtrations of the $E_\alpha$, and from
\hyperref[I.lem9.21.10]{9.21.10} (respectively
\hyperref[I.lem9.21.18]{9.21.18}, taking account of the second of
conditions \hyperlink{sga4.exp.I.equation.9.22.1}{(9.22.1)}). It remains to prove c), whose first part follows
immediately from \hyperref[I.lem9.21.9]{9.21.9} (respectively
\hyperref[I.lem9.21.17]{9.21.17}), by taking $\Pi'_\circ=0$ and
$\Pi_\circ=c$. It remains to prove that if
$\Pi'\geq\Pi\geq c$ are two cardinals, then for every $X$ in
$F^{\Pi'}$ (respectively ${F'}^{\Pi'}$) one has
$X\simeq\varinjlim_{i\in K}X_i$, with the $X_i$ in $F^{\Pi}$
(respectively ${F'}^{\Pi}$), $K$ large relative to $\Pi$, and
$\card K\leq{\Pi'}^{\Pi}$. For this, return to the proofs of
\hyperref[I.lem9.21.9]{9.21.9} (ii) and
\hyperref[I.lem9.21.17]{9.21.17} (ii). By condition
\hyperref[I.def9.12]{9.12} c) on the cardinal filtration of $E_\alpha$,
one may suppose that each $I_\alpha$ has cardinality $\leq\Pi'$, and
therefore one will have
$$
\card I\leq ({\Pi'}^{\Pi})^{\card \ob B}\leq
({\Pi'}^{\Pi})^c={\Pi'}^{\Pi c}={\Pi'}^{\Pi}{\quoi}.
$$
The index set $K$ used in \hyperref[I.sec9.21.15]{9.21.15} (respectively \hyperref[I.sec9.21.21]{9.21.21}) is the set of
subsets $J$ of $I$ that are filtered (\ie large relative to
$\Pi'_\circ=0$), stable under $\varphi$, and such that
$\card J\leq\Pi$. It is then immediate that
$$
\card K\leq (\card I)^{\Pi}\leq
({\Pi'}^{\Pi})^{\Pi}={\Pi'}^{\Pi^2}={\Pi'}^{\Pi}{\quoi},
$$
which\pageoriginale completes the proof of \hyperref[I.cor9.24]{9.24}.

To conclude, it is appropriate to give a statement freed from the
somewhat technical hypotheses of \hyperref[I.them9.22]{9.22}, replacing
them by the conjunction, for the $E_\alpha$, of the hypotheses that
occur in \hyperref[I.prop9.11]{9.11} (which ensure the accessibility of
the objects of the $E_\alpha$) and in \hyperref[I.prop9.13]{9.13}
(which ensure the existence of a cardinal filtration in $E_\alpha$):


% END INCLUDED SGA4_English_current_buildable_source_20260906/source/01/drafts/SGA4_I_LINES_7302_7379_English.texfrag


% BEGIN INCLUDED SGA4_English_current_buildable_source_20260906/source/01/drafts/SGA4_I_LINES_7380_7423_English.texfrag
\begin{enonce*}{Corollary 9.25}\etiquette[9.25]{I.cor9.25}
Let $p:E\rightarrow B$ be a fibered functor, where $B$ is a category
equivalent to a small category, and where, for every arrow
$f:\alpha\rightarrow\beta$ of $B$, the functor
$f^\ast:E_\beta\rightarrow E_\alpha$ is accessible
(\hyperref[I.def9.2]{9.2}). Suppose moreover that, for every
$\alpha\in\ob B$, the fiber category $E_\alpha$ satisfies the
following conditions:
\begin{enumerate}
\renewcommand{\theenumi}{\alph{enumi}}
\renewcommand{\labelenumi}{\theenumi)}
\item $E_\alpha$ admits a small full subcategory generating by strict
epimorphisms (\hyperref[I.def7.1]{7.1}).

\item $E_\alpha$ is stable under fiber products, kernels of pairs of
arrows, sums of two objects, and cokernels of pairs of arrows, and
finally under small filtered inductive limits
\footnote[1]{This last condition is probably unnecessary; cf. \hyperref[I.rem9.13.2]{9.13.2}.}.

\item The functor $\Ker(u,v)$ on the category of pairs of arrows of
$E_\alpha$ is accessible (for example, it commutes with small filtered
inductive limits).

\item Every strict epimorphism of $E_\alpha$
(\hyperref[I.sec10.2]{10.2}) is a universal strict epimorphism.
\end{enumerate}

Under these conditions, each of the categories $F=\Hhom_B(B,E)$ and
$F'=\Hhomcart_B(B,E)$ admits a small subcategory generating by strict
epimorphisms, and all its objects are accessible
(\hyperref[I.def9.3]{9.3}).
\end{enonce*}
Moreover,\pageoriginale \hyperref[I.cor9.24]{9.24} gives us:

\begin{enonce*}{Corollary 9.26}\etiquette[9.26]{I.cor9.26}
Under the conditions of $\hyperref[I.cor9.25]{9.25}$, and in the case where one considers
$F'=\Hhomcart_B(B,E)$, suppose that the functors
$f^\ast:E_\beta\rightarrow E_\alpha$ commute with small filtered
inductive limits. Then $F$ (respectively $F'$) admits a cardinal
filtration, which can be made explicit by the procedure of
\hyperref[I.cor9.24]{9.24}.
\end{enonce*}


% END INCLUDED SGA4_English_current_buildable_source_20260906/source/01/drafts/SGA4_I_LINES_7380_7423_English.texfrag


% BEGIN INCLUDED SGA4_English_current_buildable_source_20260906/source/01/drafts/SGA4_I_LINES_7424_7512_English.texfrag
\section{Glossary}\etiquette[10]{I.sec10}

For the reader's convenience, we collect here the definitions of some
terms used in the preceding sections. We denote a category by $C$.

\subsection{Cartesian, cocartesian}\etiquette[10.1]{I.sec10.1}

A diagram in $C$
\[
\xymatrix@C=3.25cm@R=1.5cm{
A\ar[r] \ar[d]& B \ar[d]\\
C\ar[r] & D }
\]
is called \emph{cartesian} if it is commutative and the canonical
morphism from $A$ to the fiber product $C\times_D B$ is an isomorphism.
It is called \emph{cocartesian} if the corresponding diagram in the
category opposite to \sisi{$𝒜$}{$C$} is cartesian.

\subsection{Epimorphism, strict epimorphism}\etiquette[10.2]{I.sec10.2}

\etc\pageoriginale \Cf \hyperref[I.sec10.3]{10.3}.

\subsection{Epimorphic family, strictly epimorphic
family}\etiquette[10.3]{I.sec10.3}

\etc: A family
$$
f_i:A_i\longrightarrow B, i\in I{\quoi},
$$
of arrows with the same target in $C$ is called an \emph{epimorphic
family} if, for every object $C$ of $C$, the map induced by the $f_i$,
\hypertarget{sga4.exp.I.equation.10.3.1}{}\begin{equation}
\Hom_{C}(B,C)\longrightarrow \prod_i
\Hom_{C}(A_i,C),\tag{10.3.1}
\end{equation}
is injective.

The family under consideration is called \emph{strictly epimorphic} if
the image of the map \hyperlink{sga4.exp.I.equation.10.3.1}{(10.3.1)} consists of the families $(g_i)$ such
that, for every object $R$ of $C$, every pair of indices $i,j\in I$,
and every pair of arrows $u:R\rightarrow A_i$, $v:R\rightarrow A_j$
with $f_iu=f_jv$, one also has $g_iu=g_jv$. When the fiber products
$A_i\times_BA_j$ are representable for $i,j\in I$, this is equivalent
to saying that the image of \hyperlink{sga4.exp.I.equation.10.3.1}{(10.3.1)} consists of the $(g_i)$ such that,
for every pair of indices $i,j\in I$, one has
$g_i\pr_1=g_j\pr_2$, where $\pr_1,\pr_2$ are the two projections of
$A_i\times_BA_j$. The family $(f_i)_{i\in I}$ is called an
\emph{effective epimorphic family} if the preceding condition is
satisfied, \ie if it is strictly epimorphic and the fiber products
$A_i\times_BA_j$ are representable.

The family $(f_i)_{i\in I}$ is called a \emph{universally epimorphic}
(respectively \emph{universally effective epimorphic}) family if the
morphisms $f_i$ are pullbackable (\hyperref[I.sec10.3]{10.3}), and if,
for every arrow $B'\rightarrow B$, the family of arrows
$f'_i:A'_i=A_i\times_BB'\rightarrow B'$, $i\in I$, deduced from the
family $(f_i)_{i\in I}$ by the base change $B'\rightarrow B$, is
epimorphic (respectively effective epimorphic).

The\pageoriginale notion of a \emph{universally strictly epimorphic
family} will be defined in (II \hyperref[II.def2.5]{2.5}, \hyperref[II.prop2.6]{2.6}). Note that when the morphisms
$f_i$ are pullbackable (for example, if fiber products are representable
in $C$), this notion coincides with that of a universally effective
epimorphic family (use II \hyperref[II.cor2.4]{2.4}). In the general case, the six variants of
the notion of epimorphicity under consideration are related by the
logical implications:
\[
\xymatrix@C=-1cm@R=.5cm{
  &   & \txt<3cm>{universally effective epimorphic} \ar@{=>}[lld]\ar@{=>}[rrd] &     & &\\
\txt<3cm>{effective epimorphic} \ar@{=>}[rdd]       &   &    &     &  \txt<3cm>{universally strictly epimorphic} \ar@{=>}[rd] \ar@{=>}[ld]&\\
 &   &    &   \txt{strictly\\ epimorphic}
         \ar@{=>}[lld]&     & \txt{universally \\epimorphic} \ar@{=>}[llld]\\
         & \txt{epimorphic}&    &     &      &
}
\]
A morphism $f:A\rightarrow B$ of $C$ is called an \emph{epimorphism
(respectively a strict epimorphism, respectively an effective
epimorphism, respectively a universal epimorphism, respectively a
universal effective epimorphism, respectively a universal strict
epimorphism)} if the family of morphisms consisting of the single
element $f$ is epimorphic (respectively~\ldots).


% END INCLUDED SGA4_English_current_buildable_source_20260906/source/01/drafts/SGA4_I_LINES_7424_7512_English.texfrag


% BEGIN INCLUDED SGA4_English_current_buildable_source_20260906/source/01/drafts/SGA4_I_LINES_7513_7636_English.texfrag
\subsection{Monomorphic family, strictly monomorphic
  family}\etiquette[10.4]{I.sec10.4}

\etc A family of arrows $f_i:A_i\rightarrow B$ of $C$ is called
\emph{monomorphic (respectively strictly monomorphic,\ldots)} if,
regarded as a family of arrows of the opposite category $C^{\circ}$,
it is epimorphic (respectively strictly epimorphic,\ldots). An arrow
$f:A\rightarrow B$ of $C$ is called a \emph{monomorphism (respectively
a strict monomorphism,\ldots)} if the family consisting only of $f$
is monomorphic (respectively\pageoriginale strictly
monomorphic,\ldots), \ie if $f$, regarded as an arrow of the opposite
category $C^{\circ}$, is an epimorphism (respectively a strict
epimorphism,\ldots).

\subsection{Monomorphism, strict
  monomorphism}\etiquette[10.5]{I.sec10.5}

\etc \Cf \hyperref[I.sec10.4]{10.4}.

\subsection{Projector, image of a projector, direct factor of an
object}\etiquette[10.6]{I.sec10.6}

Let $f:X\rightarrow X$ be an endomorphism of an object of $C$. The
endomorphism $f$ is called a \emph{projector} if $f^2=f$. The pair
$(f,\id_X)$ then has a representable cokernel $X'$ if and only if it
has a representable kernel $X''$; when this is the case, there exists a
unique isomorphism $u:X'\rightarrow X''$ in $C$ such that $iup=f$,
where $p:X\rightarrow X'$ and $i:X''\rightarrow X$ are the canonical
morphisms.\nde{Representability of the kernel, or of the cokernel, is
also equivalent to the existence of morphisms $i$, $p$, and an
isomorphism~$u$ such that $iup=f$ and $pi=u^{-1}$. In particular, this
property is preserved by every functor.}
One then generally identifies $X'$ with
$X''$ and calls it the \emph{image of the projector} $f$; the
projector $f$ is said to \emph{admit an image} if
$\Ker(f,\id_X)$ is representable, \ie if $\Coker(f,\id_X)$ is
representable. A subobject (respectively a quotient) $Y$ of $X$ is
called a \emph{direct factor} of $X$ if there exists a projector $f$
on $X$ that factors through $Y$ and has $Y$ as its image as a
subobject (respectively as a quotient) of $X$. Sometimes, by abuse of
language, an object of $C$ is called a \emph{direct factor of} $X$ if
it is isomorphic to the image of a projector on $X$.

\subsection{Pullbackable}\etiquette[10.7]{I.sec10.7}

An arrow $f:A\rightarrow B$ of $C$ is called \emph{pullbackable} if,
for every arrow $B'\rightarrow B$ of $C$, the fiber product
$A\times_B B'$ is representable in $C$.

\subsection{Quotient, strict quotient}\etiquette[10.8]{I.sec10.8}

\etc\pageoriginale

Let $A$ be an object of $C$. Two epimorphisms $f:A\rightarrow B$ and
$f':A\rightarrow B'$ with source $A$ are called \emph{equivalent} if
there exists an isomorphism $u:B\xrightarrow{\sim}B'$ such that
$uf=f'$. This defines an equivalence relation on the set of
epimorphisms with source $A$, whose classes are called the
\emph{quotients, or quotient objects,} of $A$. For each quotient of
$A$, one generally chooses an element of this class, say
$f:A\rightarrow B$, and one often speaks (by abuse of language) of
the quotient $B$ of $A$ (or of the quotient $f:A\rightarrow B$ of
$A$). The quotient $B$ is called a \emph{strict quotient}
(respectively an \emph{effective quotient, respectively a universal
quotient,}\ldots) if the morphism $f:A\rightarrow B$ is a strict
epimorphism (respectively an effective epimorphism, respectively a
universal epimorphism,\ldots) (\hyperref[I.sec10.3]{10.3}).

\subsection{Equivalence relations}\etiquette[10.9]{I.sec10.9}

Let $F$ and $G$ be two presheaves of sets on $C$. A diagram
$p_1$, $p_2:F\rightrightarrows G$ is called an \emph{equivalence
relation on} $G$ if, for every object $A$ of $C$, the map
$$
(p_1(A),p_2(A)):F(A)\longrightarrow G(A)\times G(A)
$$
induces a bijection from $F(A)$ onto the graph of an equivalence
relation on the set $G(A)$.

A diagram $p_1$, $p_2:B\rightrightarrows C$ in $C$ is called an
equivalence relation on $C$ if the corresponding diagram of
presheaves is an equivalence relation. When finite products and fiber
products are representable in $C$, a functor $\Phi:C\rightarrow C'$
commuting with finite products and fiber products carries equivalence
relations on $C$ to equivalence relations on $\Phi(C)$.

Let\pageoriginale $\pi:C\rightarrow D$ be a morphism of $C$ such that
the fiber product $C\times_D C$ is representable. Then the diagram
$\pr_1$, $\pr_2:C\times_D C\rightrightarrows C$ is an equivalence
relation on~$C$.

\subsection{Effective equivalence relation, universally effective
equivalence relation}\etiquette[10.10]{I.sec10.10}

An equivalence relation $p_1$, $p_2:R\rightrightarrows A$ on an object
$A$ of $C$ is called an \emph{effective equivalence relation} if
there exists a morphism $\pi:A\rightarrow B$ such that the square
\[
\xymatrix@C=3.75cm@R=1.5cm{
R\ar[r]^{p_2} \ar[d]_{p_1}& A \ar[d]^{\pi}\\
A\ar[r] & B }
\]
is cartesian and cocartesian. (The morphism $\pi$ is the cokernel of
the pair $(p_1,p_2)$ and is therefore determined up to unique
isomorphism.) The morphism $\pi$ is then an effective epimorphism; if
it is a universal effective epimorphism (\hyperref[I.sec10.3]{10.3}),
the \emph{equivalence relation} is said to be \emph{universally
effective.}

\subsection{Subobject, strict subobject}\etiquette[10.11]{I.sec10.11}

\etc These are the notions dual to those of quotient and strict
quotient \etc (\hyperref[I.sec10.8]{10.8}).

\pdfbookmark[1]{Bibliography}{sga4.exp.I.bibliography.main}
\begingroup
\makeatletter
\def\Hy@writebookmark#1#2#3#4#5{}
\makeatother
\begin{thebibliography}{1}
\bibitem[1]{sga4.exp.I.bib.I.1} \sisi{B. Mitchell, Theory of Categories, Academic Press
(1965).}{\textsc{B. Mitchell}, ``Theory of Categories,'' Academic
Press (1965).}
\end{thebibliography}
\endgroup


\newpage

\footnotetext[1]{We reproduce here, with his permission, secret papers
of N. \textsc{Bourbaki}. The references in this text are to his
learned work.}

% END INCLUDED SGA4_English_current_buildable_source_20260906/source/01/drafts/SGA4_I_LINES_7513_7636_English.texfrag


% BEGIN INCLUDED SGA4_English_current_buildable_source_20260906/source/01/drafts/SGA4_I_LINES_7637_7820_English.texfrag
\pdfbookmark[1]{II. Appendix: Universes (by N. Bourbaki)}{sga4.exp.I.appendix.II}
\begingroup
\makeatletter
\def\Hy@writebookmark#1#2#3#4#5{}
\makeatother
\section*{II. Appendix: Universes (by
  N. \textsc{Bourbaki}($^*$))}\pageoriginale
\endgroup
\markboth{II. APPENDIX: UNIVERSES}{II. APPENDIX: UNIVERSES}

\hypersetup{bookmarksdepth=2}
\makeatletter
\let\sgaFourSavedToclevelSection\toclevel@section
\def\toclevel@section{2}
\makeatother

\setcounter{section}{0}
\section{Definition and first properties of universes}%
\hypertarget{sga4.exp.I.appendix.II.sec.1}{}

\hypertarget{sga4.exp.I.appendix.II.sec.1.definition.1}{}
\begin{enonce*}{Definition 1}
A set $U$ is called a universe if it satisfies the following conditions:
\begin{itemize}
\hypertarget{sga4.exp.I.appendix.II.sec.1.definition.1.condition.U.I}{}
\item[(U.I)] if $x\in U$ and $y\in x$, then $y\in U$;

\hypertarget{sga4.exp.I.appendix.II.sec.1.definition.1.condition.U.II}{}
\item[(U.II)] if $x$, $y\in U$, then $\{x,y\}\in U$;

\hypertarget{sga4.exp.I.appendix.II.sec.1.definition.1.condition.U.III}{}
\item[(U.III)] if $x\in U$, then $\mathcal{P}(x)\in U$;

\hypertarget{sga4.exp.I.appendix.II.sec.1.definition.1.condition.U.IV}{}
\item[(U.IV)] if $(x_{\alpha})_{\alpha\in I}$ is a family
of elements of $U$, and if $I\in U$, then the union
$\bigcup_{\alpha\in I}x_{\alpha}$ belongs to $U$;

\hypertarget{sga4.exp.I.appendix.II.sec.1.definition.1.condition.U.question}{}
\item[(U.?)] if $x$, $y\in U$, then the ordered pair $(x,y)\in U$.
\end{itemize}
\end{enonce*}
\noindent N.B. Since it has, I believe, been decided that in future
editions the ordered pair will be defined à la \sisi{Kuratowaki}{Kuratowski}
by $(x,y)=\{\{x,y\}, \{x\}\}$, condition (U.?) is unnecessary, since
it follows from \hyperlink{sga4.exp.I.appendix.II.sec.1.definition.1.condition.U.II}{(U.II)}.

\hypertarget{sga4.exp.I.appendix.II.sec.1.examples}{}
\begin{enonce*}[remark]{Examples}
\begin{enumerate}
\renewcommand{\labelenumi}{\theenumi)}
\hypertarget{sga4.exp.I.appendix.II.sec.1.example.1}{}
\item The \emph{empty} set is a universe, denoted by $U_\circ$.

\hypertarget{sga4.exp.I.appendix.II.sec.1.example.2}{}
\item Consider the nonempty finite words formed with the four
symbols $''\{'',''\}'','',''$ and $''∅''$ (\cf Alg. I).
Define, by induction\pageoriginale on the length $n$ of such a word,
the notion of a \emph{meaningful} word:
\begin{enumerate}
\renewcommand{\theenumii}{\alph{enumii}}
\renewcommand{\labelenumii}{(\theenumii)}
\item For $n=1$, only the word $∅$ is meaningful;

\item for a word $A$ of length $n$ to be meaningful, it is necessary and
sufficient that there exist $p$ distinct meaningful words
$A_1,\ldots,A_p$ $(p\geq 1)$ of lengths $<n$ such that
$$
A=\{A_1,A_2,\ldots,A_p\}{\quoi}.
$$
\end{enumerate}
\end{enumerate}
\end{enonce*}
\noindent For example, $∅$, $\{∅\}$,
$\{\{∅\},\{\{∅\}\},\{∅,\{∅\}\}\}$
are meaningful words. It is clear, by induction on the
length, that every meaningful word denotes a term of Set
Theory; for example, if $A=\{A_1,A_2,\ldots,A_p\}$, A
denotes the set whose elements are $A_1$, $A_2,\ldots,$ and
$A_p$; these terms are plainly \emph{finite sets.} Let
$U_1$ be the set of sets thus obtained. It is readily checked
that $U_1$ satisfies conditions \hyperlink{sga4.exp.I.appendix.II.sec.1.definition.1.condition.U.I}{(U.I)}, \hyperlink{sga4.exp.I.appendix.II.sec.1.definition.1.condition.U.II}{(U.II)}, \hyperlink{sga4.exp.I.appendix.II.sec.1.definition.1.condition.U.III}{(U.III)}, and \hyperlink{sga4.exp.I.appendix.II.sec.1.definition.1.condition.U.IV}{(U.IV)}
of \hyperlink{sga4.exp.I.appendix.II.sec.1.definition.1}{Definition 1} (but not that idiotic condition (U.?), which is no
problem if one is willing to decanulate the ordered pair). Thus $U_1$ is a
\emph{universe}. Note that the elements of $U_1$ are finite, and
that $U_1$ is countable.

In the statements that follow, $U$ \emph{denotes a universe.}

\hypertarget{sga4.exp.I.appendix.II.sec.1.proposition.1}{}
\begin{enonce*}{Proposition 1}
If $X\in U$ and if $y\subset x$, then $y\in U$.
\end{enonce*}
Indeed, $\mathcal{P}(x)\in U$ by \hyperlink{sga4.exp.I.appendix.II.sec.1.definition.1.condition.U.III}{(U.III)}, whence $y\in
\mathcal{P}(x)$ and $y\in U$ by \hyperlink{sga4.exp.I.appendix.II.sec.1.definition.1.condition.U.I}{(U.I)}.

\hypertarget{sga4.exp.I.appendix.II.sec.1.corollary.to.proposition.1}{}
\begin{enonce*}{Corollary}
If $x\in U$, every quotient set $y$ of $x$ is an element of $U$.
\end{enonce*}
Indeed, $y$ is a subset of $\mathcal{P}(x)$. Hence $y\in U$ by
\hyperlink{sga4.exp.I.appendix.II.sec.1.definition.1.condition.U.III}{(U.III)} and \hyperlink{sga4.exp.I.appendix.II.sec.1.proposition.1}{Proposition 1}.

\hypertarget{sga4.exp.I.appendix.II.sec.1.proposition.2}{}
\begin{enonce*}{Proposition 2}
If\pageoriginale $x\in U$, then $\{x\}\in U$.
\end{enonce*}
This follows from \hyperlink{sga4.exp.I.appendix.II.sec.1.definition.1.condition.U.II}{(U.II)}, applied with $y=x$.

\hypertarget{sga4.exp.I.appendix.II.sec.1.proposition.3}{}
\begin{enonce*}{Proposition 3}
Every ordered pair, ordered triple, and ordered quadruple of elements of $U$ is an
element of $U$.
\end{enonce*}
This is true for ordered pairs by the N.B. or \hyperlink{sga4.exp.I.appendix.II.sec.1.definition.1.condition.U.question}{(U.?)}. The case of ordered
triples follows, since $(x,y,z)=((x,y),z)$, and then that of ordered
quadruples, since $(x,y,z,t)=((x,y,z),t)$.

\hypertarget{sga4.exp.I.appendix.II.sec.1.proposition.4}{}
\begin{enonce*}{Proposition 4}
If $X$, $Y\in U$, then $X\times Y\in U$.
\end{enonce*}
Indeed, for $x\in X$, $\{x\}\times Y$ is the union of the
family $\{(x,y)\}_{y\in Y}$; since $\{(x,y)\}\in U$ by
\hyperlink{sga4.exp.I.appendix.II.sec.1.definition.1.condition.U.I}{(U.I)}, \hyperlink{sga4.exp.I.appendix.II.sec.1.definition.1.condition.U.II}{(U.II)}, and \hyperlink{sga4.exp.I.appendix.II.sec.1.proposition.3}{Proposition 3}, we have $\{x\}\times Y\in U$ by
(U.IV). Finally, $X\times Y$ is the union of the family $(\{x\}\times
Y)_{x\in X}$; it is therefore an element of $U$, again by \hyperlink{sga4.exp.I.appendix.II.sec.1.definition.1.condition.U.IV}{(U.IV)}.

\hypertarget{sga4.exp.I.appendix.II.sec.1.corollary.to.proposition.4.1}{}
\begin{enonce*}{Corollary 1}
If $X$, $Y$, $Z,\ldots$ are elements of $U$, all the sets
of the scale constructed on $X$, $Y$, $Z,\ldots$ are elements
of $U$ {\rm(\cf Chapter IV, \S)}.
\end{enonce*}

This indeed follows by successive applications of \hyperlink{sga4.exp.I.appendix.II.sec.1.proposition.4}{Proposition 4} and
of (U.III).

\hypertarget{sga4.exp.I.appendix.II.sec.1.corollary.to.proposition.4.2}{}
\begin{enonce*}{Corollary 2}
If $(x_{\alpha})_{\alpha \in I}$ is a family of elements of $U$
and if $I\in U$, the sum set $\Sigma_{\alpha \in I} x_{\alpha}$
is an element of $U$.
\end{enonce*}
Indeed, this sum set is a subset of the product
$(\bigcup_{\alpha\in I} x_{\alpha})\times I$ (Chapter II), a product
both of whose factors are elements of $U$ (by (U.IV) and
the hypothesis). We then apply \hyperlink{sga4.exp.I.appendix.II.sec.1.proposition.4}{Propositions 4} and \hyperlink{sga4.exp.I.appendix.II.sec.1.proposition.1}{1}.

\hypertarget{sga4.exp.I.appendix.II.sec.1.proposition.5}{}
\begin{enonce*}{Proposition 5}
If\pageoriginale $X$ and $Y$ are elements of $U$, every
correspondence between $X$ and $Y$ (in particular, every map from
$X$ to $Y$) is an element of $U$.
\end{enonce*}
Indeed, such a correspondence $C$ is an ordered triple $(X,Y,\Gamma)$,
where $\Gamma$ is a subset of $X\times Y$ (the graph of $C$) (Chapter
II, \S). We have $\Gamma\in U$ by \hyperlink{sga4.exp.I.appendix.II.sec.1.proposition.4}{Propositions 4} and 1. Hence $C\in U$
by \hyperlink{sga4.exp.I.appendix.II.sec.1.proposition.3}{Proposition 3}.

\hypertarget{sga4.exp.I.appendix.II.sec.1.proposition.6}{}
\begin{enonce*}{Proposition 6}
If $X$, $Y\in U$, every set $Z$ of correspondences between $X$ and
$Y$ (in particular, of maps from $X$ to $Y$) is an element of
$U$.
\end{enonce*}
Indeed, $Z$ is a subset of $\{X\}\times \{Y\}\times
\mathcal{P}(X\times Y)$. Now this product is an element of $U$
by
\hyperlink{sga4.exp.I.appendix.II.sec.1.proposition.2}{Proposition 2} and the \hyperlink{sga4.exp.I.appendix.II.sec.1.corollary.to.proposition.4.1}{corollary to Proposition 4}. Therefore $Z\in
U$ by \hyperlink{sga4.exp.I.appendix.II.sec.1.proposition.1}{Proposition 1}.

\hypertarget{sga4.exp.I.appendix.II.sec.1.corollary.to.proposition.6}{}
\begin{enonce*}{Corollary}
If $(x_{\alpha})_{\alpha \in I}$ is a family of elements of $U$,
and if $I\in U$, then $\Pi_{\alpha\in I}x_{\alpha} \in U$.
\end{enonce*}
Indeed, this product is a set of maps from $I$ to
$\bigcup_{\alpha\in I}x_{\alpha}$, and this union is an element of
$U$ by \hyperlink{sga4.exp.I.appendix.II.sec.1.definition.1.condition.U.IV}{(U.IV)}.

\hypertarget{sga4.exp.I.appendix.II.sec.1.proposition.7}{}
\begin{enonce*}{Proposition 7}
If $X$ is a subset of $U$ whose cardinality is at most that of an
element of $U$, then $X$ is an element of $U$.
\end{enonce*}
Let $I$ be an element of $U$ such that $\card(X)\leq \card(I)$. There is a
surjection $i\mapsto x_i$ from a subset $I'$ of $I$ onto $X$. Then
$X$ is the union of the family $(\{x_i\})_{i\in I'}$; now this
union is an element of $U$ by \hyperlink{sga4.exp.I.appendix.II.sec.1.proposition.1}{Proposition 1}, \hyperlink{sga4.exp.I.appendix.II.sec.1.proposition.2}{Proposition 2}, and \hyperlink{sga4.exp.I.appendix.II.sec.1.definition.1.condition.U.IV}{(U.IV)}.

\hypertarget{sga4.exp.I.appendix.II.sec.1.corollary.to.proposition.7}{}
\begin{enonce*}{Corollary}
If\pageoriginale $U$ is nonempty, every finite subset of $U$ is an
element of $U$, and $U$ has elements of arbitrarily large finite cardinality.
\end{enonce*}
On the other hand, if $U=∅$, then $∅$ is a finite subset
of $∅$ but not an element of $∅$.

Indeed, if $U$ is nonempty, \hyperlink{sga4.exp.I.appendix.II.sec.1.proposition.2}{Proposition 2} shows that a one-element
set $x_{\circ}$ belongs to $U$. By induction on $n$,
put $x_{n+1}=\mathcal{P}(x_n)$. We %%
have $x_n\in U$ by \hyperlink{sga4.exp.I.appendix.II.sec.1.definition.1.condition.U.III}{(U.III)}, and $\card(x_{n+1})=2^{\card(x_n)}$, so
that $U$ has elements of arbitrarily large finite cardinality.

\begin{enonce*}[remark]{Remark}
It follows from \hyperlink{sga4.exp.I.appendix.II.sec.1.proposition.2}{Proposition 2} and the \hyperlink{sga4.exp.I.appendix.II.sec.1.corollary.to.proposition.7}{corollary to Proposition 7} that every
nonempty universe $U$ contains the universe $U_1$ of \hyperlink{sga4.exp.I.appendix.II.sec.1.example.2}{Example 2}; indeed,
$∅\in U$ by \hyperlink{sga4.exp.I.appendix.II.sec.1.proposition.1}{Proposition 1}. Thus $U_1$ is
the intersection of all nonempty universes. More generally:
\end{enonce*}

\hypertarget{sga4.exp.I.appendix.II.sec.1.proposition.8}{}
\begin{enonce*}{Proposition 8}
If $(U_{\lambda})_{\lambda\in L}$ is a nonempty family
of universes, then $U=\bigcap_{\lambda\in L}U_{\lambda}$ is a
universe.
\end{enonce*}
This follows at once from \hyperlink{sga4.exp.I.appendix.II.sec.1.definition.1}{Definition 1}.

% END INCLUDED SGA4_English_current_buildable_source_20260906/source/01/drafts/SGA4_I_LINES_7637_7820_English.texfrag


% BEGIN INCLUDED SGA4_English_current_buildable_source_20260906/source/01/drafts/SGA4_I_LINES_7821_7923_English.texfrag
\section{Universes and species of structures}%
\hypertarget{sga4.exp.I.appendix.II.sec.2}{}

Let $\mathcal{E}$ be a species of structure; suppose, to fix ideas
and simplify the exposition, that every structure of species
$\mathcal{E}$ is defined on an underlying set. Let $(X,S)$ be a
structure of species $\mathcal{E}$ (where $X$ is the underlying set
and $S$ the structure), and let $U$ be a \emph{universe.} If $X\in U$,
then all the constituent objects of the structure $S$ on $X$ are
elements of $U$ (by \hyperlink{sga4.exp.I.appendix.II.sec.1.corollary.to.proposition.4.1}{Cor. 1 of Prop. 4}, \hyperlink{sga4.exp.I.appendix.II.sec.1}{${\rm n^{\circ}}$ 1}, and by
\hyperlink{sga4.exp.I.appendix.II.sec.1.definition.1.condition.U.I}{(U.I)}), so the structure $(X,S)$ is an element of $U$.

Suppose\pageoriginale now that $\mathcal{E}$ is a species of structure
with \emph{morphisms}. If $X$ and $X'$ are elements of $U$ endowed
with structures of species $\mathcal{E}$, then the set of morphisms
from $X$ to $X'$ is again an element of $U$ (\hyperlink{sga4.exp.I.appendix.II.sec.1}{${\rm n^{\circ}}$ 1},
\hyperlink{sga4.exp.I.appendix.II.sec.1.proposition.6}{Prop. 6}).

Now consider the \emph{category} $(U\text{-}\mathcal{E})$ defined in
\hyperlink{sga4.exp.I.appendix.II.sec.1}{\S 1}, ${\rm n^{\circ}}$ 2, Ex. c). Since the objects and arrows of
$(U\text{-}\mathcal{E})$ are elements of $U$,
$(U\text{-}\mathcal{E})$ is a \emph{pair of subsets} of $U$, endowed
with a quadruple of maps; in general, it is not an element of $U$.
The notation $X\in (U\text{-}\mathcal{E})$ will mean that $X$ is an
element of $U$ endowed with a structure of species $\mathcal{E}$.

The category $(U\text{-}\mathcal{E})$ is \emph{stable} under many
operations:

\begin{enumerate}
\renewcommand{\theenumi}{\alph{enumi}}
\renewcommand{\labelenumi}{\theenumi)}
\item If $X\in(U\text{-}\mathcal{E})$, and if $X'$ is a subset of $X$ that
admits an \emph{induced structure}, then $X'\in
(U\text{-}\mathcal{E})$. This follows from \hyperlink{sga4.exp.I.appendix.II.sec.1.proposition.1}{Prop. 1}, \hyperlink{sga4.exp.I.appendix.II.sec.1}{${\rm
n^{\circ}}$ 1}.

\item If $X\in (U\text{-}\mathcal{E})$, and if $X''$ is a quotient set of
$X$ that admits a \emph{quotient structure}, then $X''\in
(U\text{-}\mathcal{E})$ (\hyperlink{sga4.exp.I.appendix.II.sec.1.corollary.to.proposition.1}{Cor. of Prop. 1}, \hyperlink{sga4.exp.I.appendix.II.sec.1}{${\rm n^{\circ}}$ 1}).

\item If $X$, $Y\in (U\text{-}\mathcal{E})$, and if $X\times Y$ admits a
\emph{product structure}, then $X\times Y\in
(U\text{-}\mathcal{E})$ (\hyperlink{sga4.exp.I.appendix.II.sec.1.proposition.4}{Prop. 4}, \hyperlink{sga4.exp.I.appendix.II.sec.1}{${\rm n^{\circ}}$ 1}). More
generally, if $(X_i)_{i\in I}$ is a family of elements of
$(U\text{-}\mathcal{E})$, \emph{if} $I\in U$, and if $\Pi_{i\in
I}X_i$ admits a \emph{product structure}, then $\Pi_{i\in I}X_i
\in (U\text{-}\mathcal{E})$ (\hyperlink{sga4.exp.I.appendix.II.sec.1.corollary.to.proposition.6}{Cor. of Prop. 6}). The analogous
assertions hold for \emph{sum structures} (\hyperlink{sga4.exp.I.appendix.II.sec.1.corollary.to.proposition.4.2}{Cor. 2 of Prop. 4}).

\item Let $I$ be a preordered set, and let $(X_i, f_{ij})_{i,j\in
I}$ be a \emph{projective system} (\resp \emph{inductive system}) of
sets endowed with structures of species $\mathcal{E}$ and morphisms;
let $L$ be the limit of this system, in the sense of Chap. III. If
$X_i\in U$ for every $i$, if $I\in U$, and if $L$ admits a
projective-limit structure (\resp an inductive-limit structure), then
$L\in (U\text{-}\mathcal{E})$: indeed, $L$ is a\pageoriginale subset
of $\Pi_{i\in I}X_i$ (\resp a quotient set of $\Sigma_{i\in I}X_i$),
and is therefore an element of $U$ by \hyperlink{sga4.exp.I.appendix.II.sec.1}{${\rm n^{\circ}}$ 1}.
\end{enumerate}

Here are a few more particular examples:
\begin{itemize}
\hypertarget{sga4.exp.I.appendix.II.sec.2.example.star.1}{}
\item[* 1)] Let $X$, $Y\in (U\text{-}\Top)$ be two \emph{locally
compact} spaces. Then the set $\mathscr{C}(X,Y)$ of continuous maps
from $X$ to $Y$, endowed with the topology of compact convergence
(General Topology $X$), is an element of $(U\text{-}\Top)$. Indeed,
$\mathscr{C}(X,Y)\in U$ by \hyperlink{sga4.exp.I.appendix.II.sec.1.proposition.6}{Prop. 6} of \hyperlink{sga4.exp.I.appendix.II.sec.1}{${\rm
n^{\circ}}$ 1}.${}_{\mbox{*}}$

\hypertarget{sga4.exp.I.appendix.II.sec.2.example.star.2}{}
\item[* 2)] Let $X$, $Y\in (U\text{-}\Ab)$. Then the group
$\Hom_{\ZZ}(X,Y)$ is an element of $(U\text{-}\Ab)$ by \hyperlink{sga4.exp.I.appendix.II.sec.1.proposition.6}{Prop. 6} of
\hyperlink{sga4.exp.I.appendix.II.sec.1}{${\rm n^{\circ}}$ 1}. If, in addition, $\ZZ\in U$, the tensor product
$X\otimes_{\ZZ}Y$ is an element of $(U\text{-}\Ab)$; indeed, this
tensor product is a quotient of $\ZZ^{X\times Y}$, which is a subset
of $\ZZ^{X\times Y}$, which in turn is an element of $U$ by \hyperlink{sga4.exp.I.appendix.II.sec.1.proposition.4}{Props. 4}
and \hyperlink{sga4.exp.I.appendix.II.sec.1.proposition.6}{6} of \hyperlink{sga4.exp.I.appendix.II.sec.1}{${\rm n^{\circ}}$ 1}.${}_{\mbox{*}}$

\hypertarget{sga4.exp.I.appendix.II.sec.2.example.star.3}{}
\item[* 3)] Let $X\in (U\text{-}\Unif)$ be a uniform space. Then its
\emph{completion} $\widehat{X}$ is an element of $(U\text{-}\Unif)$.
Indeed, $\widehat{X}$ is a set of equivalence classes of Cauchy
filters on $X$; now a filter on $X$ is an element of
$\mathcal{P}(\mathcal{P}(X))$, so an equivalence class of filters is
an element of $\mathcal{P}(\mathcal{P}(\mathcal{P}(X)))$, whence
$\widehat{X}$ is an element of
$\mathcal{P}(\mathcal{P}(\mathcal{P}(\mathcal{P}(X))))$, and hence an
element of $U$ by \hyperlink{sga4.exp.I.appendix.II.sec.1.definition.1.condition.U.III}{(U.III)} and \hyperlink{sga4.exp.I.appendix.II.sec.1.definition.1.condition.U.I}{(U.I)}.${}_{\mbox{*}}$
\end{itemize}

N.B. Moral: Bourbaki will have to take care to ``canonicalize'' its
constructions properly. For example, in constructions in which an
element $\infty$ is adjoined (Alexandroff compactification, projective
field), it will be advantageous to take $\infty =∅$\pageoriginale
(since $∅$ is an element of every nonempty universe) and to form the
disjoint union of $\{∅\}$ and the given set. Of course, one should
also ``canonicalize'' the two-element set used to construct disjoint
unions; thus $\{∅, \{∅\}\}$ seems to me a good candidate, since it
belongs to every nonempty universe.

% END INCLUDED SGA4_English_current_buildable_source_20260906/source/01/drafts/SGA4_I_LINES_7821_7923_English.texfrag


% BEGIN INCLUDED SGA4_English_current_buildable_source_20260906/source/01/drafts/SGA4_I_LINES_7924_7993_English.texfrag
\section{Universes and categories}%
\hypertarget{sga4.exp.I.appendix.II.sec.3}{}

Let $U$ be a universe and $C$ a category. We write $C\in U$ if
$\Ob(C)\in U$ and $%%
\mathrm{F}{}l\,(C)\in U$. This notation is justified by the fact that
$C$ is the sextuple consisting of $\Ob(C)$, $\mathrm{F}{}l\,(C)$, and
the four structure maps; thus, if $\Ob(C)\in U$ and
$\mathrm{F}{}l\,(C)\in U$, these four maps are elements of $U$
(\hyperlink{sga4.exp.I.appendix.II.sec.1}{${\rm n^{\circ}}$ 1}, \hyperlink{sga4.exp.I.appendix.II.sec.1.corollary.to.proposition.4.1}{Corollary 1} of \hyperlink{sga4.exp.I.appendix.II.sec.1.proposition.4}{Proposition 4}), and hence so is
the sextuple $C$ (\hyperlink{sga4.exp.I.appendix.II.sec.1}{${\rm n^{\circ}}$ 1}, \hyperlink{sga4.exp.I.appendix.II.sec.1.proposition.3}{Proposition 3}, \textit{cum
grano salis}). This is moreover a special case of \hyperlink{sga4.exp.I.appendix.II.sec.2}{${\rm n^{\circ}}$ 2},
if one considers the species of structure ``$\cat$'' of a category.
With the notation of \hyperlink{sga4.exp.I.appendix.II.sec.2}{${\rm n^{\circ}}$ 2}, the relation $C\in U$ is
also written $C\in(U\text{-}\cat)$.

One should note that a category such as $(U\text{-}\Ens)$,
$(U\text{-}\Ord)$, $(U\text{-}\Top)$, or $(U\text{-}\Gr)$ is
\emph{not}, in general, an element of $U$.

Let $C,D$ be two categories and $U$ a universe such that $C,D\in U$.
If $C'$ is a \emph{subcategory} of $C$, then $C'\in U$; the
\emph{product category} $C\times D$, the sum category of $C$ and $D$,
and the opposite categories $C^\circ$ and $D^\circ$ are also elements
of $U$ (cf. \hyperlink{sga4.exp.I.appendix.II.sec.2}{${\rm n^{\circ}}$ 2}). The \emph{functor category}
$E=\Hhom(C,D)$ is likewise an element of $U$: indeed, $\Ob(E)$ is a set
of pairs of maps $\Ob(C)\rightarrow\Ob(D)$,
$F\ell\,(C)\rightarrow F\ell\,(D)$, and hence $\Ob(E)\in U$ by
\hyperlink{sga4.exp.I.appendix.II.sec.1}{${\rm n^{\circ}}$ 1};\pageoriginale as for $F\ell\,(E)$, it is a set
of natural transformations, that is, of maps
$\Ob(C)\rightarrow F\ell\,(D)$, whence $F\ell\,(E)\in U$ by
\hyperlink{sga4.exp.I.appendix.II.sec.1.proposition.6}{Proposition 6} of \hyperlink{sga4.exp.I.appendix.II.sec.1}{${\rm n^{\circ}}$ 1}.

\hypertarget{sga4.exp.I.appendix.II.sec.3.proposition.9}{}
\begin{enonce*}{PROPOSITION 9}
Let $\mathcal{E}$ be a species of structure with morphisms, and let
$U,U'$ be two universes such that $U'\subset U$. Then
$(U'\text{-}\mathcal{E})$ is a full subcategory of
$(U\text{-}\mathcal{E})$.
\end{enonce*}

Indeed, if $x$ and $y$ are objects of $(U'\text{-}\mathcal{E})$, they
are by definition objects of $(U\text{-}\mathcal{E})$. As for
$\Hom(x,y)$, it is the same set in both categories (cf. \hyperlink{sga4.exp.I.appendix.II.sec.1}{\S 1},
${\rm n^{\circ}}$ 2, Example c)).

N.B. One should note that the hypothesis that $U$ and $U'$ are
universes %%
is unnecessary, so that \hyperlink{sga4.exp.I.appendix.II.sec.3.proposition.9}{Proposition 9} would be more
appropriately placed in \hyperlink{sga4.exp.I.appendix.II.sec.1}{\S 1}, ${\rm n^{\circ}}$ 4. But the Tribe asked
that it be put here.

\begin{enonce*}[remark]{Remark}
It may happen that the axioms of the species of structure $\mathcal{E}$
imply that the \emph{cardinals} of sets endowed with structures of
species $\mathcal{E}$ are \emph{bounded} by a fixed cardinal, say
$\sheaf{\sheaf{c}}$ (* for example, the species of structure of finite
group, finitely generated group, finitely generated module over a fixed
ring $A$, or finitely generated algebra over $A$ ${}_{\mbox{*}}$).
Suppose then that there exists an element $z$ of $U'$ such that
$\sheaf{\sheaf{c}}\leq\card(z)$. Then, for every set $x$ endowed with
a structure of species $\mathcal{E}$, there exists an element $x'$ of
$U'$ equipotent to $x$, for example a subset of $z$; endow $x'$ with
the structure deduced from that of $x$ by transport of structure; this
gives an element of $(U'\text{-}\mathcal{E})$. It follows that
\emph{the inclusion functor from $(U'\text{-}\mathcal{E})$ to
$(U\text{-}\mathcal{E})$ is then essentially surjective {\rm
(\hyperlink{sga4.exp.I.appendix.II.sec.4}{\S 4}, ${\rm n^{\circ}}$ 2, Definition 2)} and is therefore an
equivalence of categories} (\hyperlink{sga4.exp.I.appendix.II.sec.4}{\S 4}, ${\rm n^{\circ}}$ 3, Theorem 1).
\end{enonce*}

% END INCLUDED SGA4_English_current_buildable_source_20260906/source/01/drafts/SGA4_I_LINES_7924_7993_English.texfrag


% BEGIN INCLUDED SGA4_English_current_buildable_source_20260906/source/01/drafts/SGA4_I_LINES_7994_8025_English.texfrag
\section{The universe axiom}%
\hypertarget{sga4.exp.I.appendix.II.sec.4}{}\pageoriginale

The very pleasant stability results of \hyperlink{sga4.exp.I.appendix.II.sec.2}{${\rm n^{\circ}}$ 2} are of
interest only if they can be applied to something other than the two
small universes $U_\circ$ and $U_1$ described in \hyperlink{sga4.exp.I.appendix.II.sec.1}{${\rm n^{\circ}}$ 1}.
We shall therefore add the following axiom to the axioms of Set Theory:

\hypertarget{sga4.exp.I.appendix.II.sec.4.axiom.A.6}{}
(A.6) \emph{For every set $x$, there exists a universe $U$ such that
$x\in U$.}

This axiom implies that, if $(x_\alpha)_{\alpha\in I}$ is a
\emph{family} of sets, there exists a universe $U$ such that
$x_\alpha\in U$ for every $\alpha\in I$: indeed, it suffices to apply
\hyperlink{sga4.exp.I.appendix.II.sec.4.axiom.A.6}{(A.6)} to $x=\bigcup_{\alpha\in I}x_\alpha$ and then apply \hyperlink{sga4.exp.I.appendix.II.sec.1.proposition.1}{Proposition
1} of \hyperlink{sga4.exp.I.appendix.II.sec.1}{${\rm n^{\circ}}$ 1}.

In particular, given a \emph{category} $C$, there exists a universe
$U$ such that $C\in U$ in the sense of \hyperlink{sga4.exp.I.appendix.II.sec.3}{${\rm n^{\circ}}$ 3}: apply the
preceding result to the family $(\Ob(C),F\ell\,(C))$. This applies to
categories of the form $(V\text{-}\mathcal{E})$, where $\mathcal{E}$
is a species of structure with morphisms and $V$ is a set, for example
a universe; in general one has $V\neq U$.

For example, if $V$ is a universe, then
$\Ob(V\text{-}\Ens)=V$ and $F\ell\,(V\text{-}\Ens)\subset V$
(\hyperlink{sga4.exp.I.appendix.II.sec.1}{${\rm n^{\circ}}$ 1}, \hyperlink{sga4.exp.I.appendix.II.sec.1.proposition.5}{Proposition 5}). The universes $U$ such that
$(V\text{-}\Ens)\in U$ are therefore those such that $V\in U$. But
the relation $V\in V$ is impossible for a universe: indeed, for every
subset $A$ of $V$, one would have $A\in V$ (\hyperlink{sga4.exp.I.appendix.II.sec.1}{${\rm n^{\circ}}$ 1},
\hyperlink{sga4.exp.I.appendix.II.sec.1.proposition.1}{Proposition 1}), whence $\card(\mathcal{P}(V))\leq\card(V)$, which is
impossible.

% END INCLUDED SGA4_English_current_buildable_source_20260906/source/01/drafts/SGA4_I_LINES_7994_8025_English.texfrag


% BEGIN INCLUDED SGA4_English_current_buildable_source_20260906/source/01/drafts/SGA4_I_LINES_8026_8095_English.texfrag
\section{Universes and strongly inaccessible cardinals}%
\hypertarget{sga4.exp.I.appendix.II.sec.5}{}\pageoriginale

Let $U$ be a universe. By condition (U.I) of \hyperlink{sga4.exp.I.appendix.II.sec.1.definition.1}{Definition 1}, every
element $x$ of $U$ is a subset of $U$; hence
$\card(x)\leq\card(U)$. Since the cardinals of the subsets of $U$ form
a well-ordered set, the cardinal
\hypertarget{sga4.exp.I.appendix.II.sec.5.equation.1.cardinal.c.U}{}
\begin{equation}
\sheaf{\sheaf{c}}(U)={\sup}_{x\in U}\card(x)\tag{1}
\end{equation}
exists. Note that, for every cardinal
$\sheaf{\sheaf{c}}<\sheaf{\sheaf{c}}(U)$, there exists an element $x$
of $U$ such that $\card(x)=\sheaf{\sheaf{c}}$; indeed, by definition
there exists $y\in U$ such that
$\sheaf{\sheaf{c}}\leq\card(y)\leq\sheaf{\sheaf{c}}(U)$, and one takes
for $x$ a suitable subset of $y$. Conversely, if $x\in U$, then
$\card(x)<\sheaf{\sheaf{c}}(U)$; indeed,
$\mathcal{P}(x)\in U$ by \hyperlink{sga4.exp.I.appendix.II.sec.1.definition.1.condition.U.III}{(U.III)}, whence
$2^{\card(x)}\leq\sheaf{\sheaf{c}}(U)$. Thus the cardinal
$\sheaf{\sheaf{c}}(U)$ has the following properties:

\begin{enumerate}
\renewcommand{\labelenumi}{\theenumi)}
\item If $\sheaf{\sheaf{c}}$ is a cardinal
$<\sheaf{\sheaf{c}}(U)$, then
$2^{\sheaf{\sheaf{c}}}<\sheaf{\sheaf{c}}(U)$; indeed, if $x$ is an
element of $U$ of cardinality $\sheaf{\sheaf{c}}$, then
$\mathcal{P}(x)\in U$ by \hyperlink{sga4.exp.I.appendix.II.sec.1.definition.1.condition.U.III}{(U.III)}, whence
$2^{\sheaf{\sheaf{c}}}<\sheaf{\sheaf{c}}(U)$.

\item If $(\sheaf{\sheaf{c}}_\lambda)_{\lambda\in I}$ is a family of
cardinals such that
$\sheaf{\sheaf{c}}_\lambda<\sheaf{\sheaf{c}}(U)$ for every
$\lambda\in I$ and $\card(I)<\sheaf{\sheaf{c}}(U)$, then the sum
cardinal $\Sigma_{\lambda\in I}\sheaf{\sheaf{c}}_\lambda$ is
$<\sheaf{\sheaf{c}}(U)$. Indeed, let $x_\lambda\in U$ be such that
$\card(x_\lambda)=\sheaf{\sheaf{c}}_\lambda$; after replacing $I$ by
an equipotent index set if necessary, one may suppose that $I\in U$;
then the sum set of the $x_\lambda$ is an element of $U$ (\hyperlink{sga4.exp.I.appendix.II.sec.1.corollary.to.proposition.4.2}{Corollary 2}
of \hyperlink{sga4.exp.I.appendix.II.sec.1.proposition.4}{Proposition 4} of \hyperlink{sga4.exp.I.appendix.II.sec.1}{${\rm n^{\circ}}$ 1}), which proves our assertion.
\end{enumerate}

We make the following definition:

\hypertarget{sga4.exp.I.appendix.II.sec.5.definition.2}{}
\begin{enonce*}{DEFINITION 2}
A\pageoriginale cardinal $\sheaf{\sheaf{d}}$ is called strongly
inaccessible if:
\begin{itemize}
\hypertarget{sga4.exp.I.appendix.II.sec.5.definition.2.condition.FI.1}{}
\item[(FI.1)] If $\sheaf{\sheaf{c}}$ is a cardinal such that
$\sheaf{\sheaf{c}}<\sheaf{\sheaf{d}}$, then
$2^{\sheaf{\sheaf{c}}}<\sheaf{\sheaf{d}}$.

\hypertarget{sga4.exp.I.appendix.II.sec.5.definition.2.condition.FI.2}{}
\item[(FI.2)] If $(\sheaf{\sheaf{c}}_\lambda)_{\lambda\in I}$ is a
family of cardinals such that
$\sheaf{\sheaf{c}}_\lambda<\sheaf{\sheaf{d}}$ for every
$\lambda\in I$, and if $\card(I)<\sheaf{\sheaf{d}}$, then
$\Sigma_{\lambda\in I}\sheaf{\sheaf{c}}_\lambda<
\sheaf{\sheaf{d}}$.
\end{itemize}
\end{enonce*}

\begin{enonce*}[remark]{Examples}
The cardinal $0$ and the countably infinite cardinal are strongly
inaccessible. No nonzero finite cardinal is strongly inaccessible.
Thus we have just proved that the universe axiom \hyperlink{sga4.exp.I.appendix.II.sec.4.axiom.A.6}{(A.6)} implies the
statement:
\begin{itemize}
\hypertarget{sga4.exp.I.appendix.II.sec.5.axiom.A-prime.6}{}
\item[${\rm (A'.6)}$] \emph{Every cardinal is strictly bounded above
by a strongly inaccessible cardinal.}
\end{itemize}
Conversely:
\end{enonce*}

% END INCLUDED SGA4_English_current_buildable_source_20260906/source/01/drafts/SGA4_I_LINES_8026_8095_English.texfrag


% BEGIN INCLUDED SGA4_English_current_buildable_source_20260906/source/01/drafts/SGA4_I_LINES_8096_8244_English.texfrag
\hypertarget{sga4.exp.I.appendix.II.sec.5.theorem.1}{}
\begin{enonce*}{THEOREM %%
1} Statement \hyperlink{sga4.exp.I.appendix.II.sec.5.axiom.A-prime.6}{$(A'.6)$} implies the universe axiom (A.6).
\end{enonce*}
Indeed, let $A$ be a set. We must construct a universe $U$ of which
$A$ is an element. Define recursively a sequence $(A_n)_{n\geq0}$ of
sets by:
\begin{itemize}
\hypertarget{sga4.exp.I.appendix.II.sec.5.theorem.1.proof.equation.2}{}
\item[(2)] $A_0=A$, $A_{n+1}=$ the union of the elements of $A_n=$
the set of elements of the elements of $A_n$.
\end{itemize}
\noindent Put $B=\bigcup^\infty_{n=0}A_n$. By \hyperlink{sga4.exp.I.appendix.II.sec.5.axiom.A-prime.6}{$(A'.6)$}, let
$\sheaf{\sheaf{c}}$ be a strongly inaccessible cardinal such that
$\card(B)<\sheaf{\sheaf{c}}$.

There exists a \emph{well-ordered} set $I$ such that
$\card(I)=\sheaf{\sheaf{c}}$. After replacing $I$ by its smallest
segment of cardinality $\sheaf{\sheaf{c}}$ if necessary, we may suppose
that \emph{every segment of $I$ distinct from $I$ has cardinality
$<\sheaf{\sheaf{c}}$.} We\pageoriginale denote by $\mathcal{E}$ the
\emph{least element of $I$}. It follows from the hypothesis on segments
that $I$ has no greatest element (otherwise we would remove it), and
hence that every element of $I$ has a successor; for $\alpha\in I$, we
denote the \emph{successor} of $\alpha$ by $s(\alpha)$.

Now define, by \emph{transfinite recursion}, a family
$(B_\alpha)_{\alpha\in I}$ of sets by:
\hypertarget{sga4.exp.I.appendix.II.sec.5.theorem.1.proof.equation.3}{}
\begin{equation}
\begin{cases}
& B_{\mathcal{E}}=B\\
& B_{s(\alpha)}=B_{\alpha}\cup \mathcal{P}(B_{\alpha})\\
& B_{\alpha}={\displaystyle{\bigcup_{\beta<\alpha}}}B_{\beta}\mbox{
if $\alpha$ has no predecessor.}
\end{cases}\tag{3}
\end{equation}
Put $U=\bigcup_{\alpha\in I}B_\alpha$. We shall show that $U$ is the
desired universe.

First, $A\in U$, since $A=A_\circ$ is a subset of
$B=B_{\mathcal{E}}$, and therefore an element of
$\mathcal{P}(B_{\mathcal{E}})\subset B_{s(\mathcal{E})}$.

Next note that, for every $\alpha\in I$, one has:
\hypertarget{sga4.exp.I.appendix.II.sec.5.theorem.1.proof.equation.4}{}
\begin{equation}
\card(B_{\alpha})<\sheaf{\sheaf{c}}{\quoi}.\tag{4}
\end{equation}
Indeed, proceed by transfinite recursion on $\alpha$. This is true for
$\alpha=\mathcal{E}$ by hypothesis. From \hyperlink{sga4.exp.I.appendix.II.sec.5.theorem.1.proof.equation.4}{(4)} one obtains
$\card(B_{s(\alpha)})\leq
\card(B_\alpha)+2^{\card(B_\alpha)}<
\sheaf{\sheaf{c}}+\sheaf{\sheaf{c}}$ (by \hyperlink{sga4.exp.I.appendix.II.sec.5.definition.2.condition.FI.1}{(FI.1)})
$=\sheaf{\sheaf{c}}$ (since $\sheaf{\sheaf{c}}$ is infinite).
Finally, if $\alpha$ has no predecessor, the set $I'$ of
$\beta<\alpha$ has cardinality $<\sheaf{\sheaf{c}}$ by construction;
hence (if $\card(B_\beta)<\sheaf{\sheaf{c}}$ for every
$\beta<\alpha$),
$\card(B_\alpha)=\card(\bigcup_{\beta\in I'}B_\beta)
\leq\Sigma_{\beta\in I'}\card(B_\beta)<\sheaf{\sheaf{c}}$ (by
(FI.2)). This proves \hyperlink{sga4.exp.I.appendix.II.sec.5.theorem.1.proof.equation.4}{(4)}.

Let us now\pageoriginale show that
\hypertarget{sga4.exp.I.appendix.II.sec.5.theorem.1.proof.equation.5}{}
\begin{equation}
\card(x)<\sheaf{\sheaf{c}}\mbox{ for every }x\in U{\quoi}.\tag{5}
\end{equation}
Indeed, it suffices to show that, for every $\alpha\in I$, one has
``$\card(x)<\sheaf{\sheaf{c}}$ for every $x\in B_\alpha$''. Proceed
again by transfinite recursion on $\alpha$. This is true for
$\alpha=\mathcal{E}$, since, if $x\in B_{\mathcal{E}}=B$, there exists
$n\geq0$ such that $x\in A_n$; hence $x\subset A_{n+1}$ by \hyperlink{sga4.exp.I.appendix.II.sec.5.theorem.1.proof.equation.2}{(2)},
$x\subset B$, and
$\card(x)\leq\card(B)<\sheaf{\sheaf{c}}$. The case where $\alpha$ has
no predecessor is evident. Finally, if $x\in B_\alpha$ and
$\alpha=s(\beta)$, either $x\in B_\beta$, in which case the assertion
``$\card(x)<\sheaf{\sheaf{c}}$'' is true by recursion, or
$x\in\mathcal{P}(B_\beta)$, whence $x\subset B_\beta$ and the
assertion ``$\card(x)<\sheaf{\sheaf{c}}$'' is true by \hyperlink{sga4.exp.I.appendix.II.sec.5.theorem.1.proof.equation.4}{(4)}.

We are now in a position to prove that $U$ is indeed a \emph{universe}:

\begin{itemize}
\hypertarget{sga4.exp.I.appendix.II.sec.5.theorem.1.proof.clause.U.I}{}
\item[(U.I)] Let $x\in U$ and $y\in x$. We must show %%
that $y\in U$. In other words, we must show that, for every
$\alpha\in I$, one has:
$$
\textrm{``}x\in B_{\alpha}\mbox{ and }y\in x\Longrightarrow y\in
B_{\alpha}\textrm{''\,}{\quoi}.
$$
Proceed again by transfinite recursion on $\alpha$. This is true for
$\alpha=\mathcal{E}$, since $x\in B_{\mathcal{E}}=B$ implies
$x\in A_n$ for some $n$; hence, if $y\in x$, then $y\in A_{n+1}$,
whence $y\in B=B_{\mathcal{E}}$. Passage to an element $\alpha$ with no
predecessor is evident. Finally, pass from $\alpha$ to $s(\alpha)$: if
$x\in B_{s(\alpha)}$ and $y\in x$, either $x\in B_\alpha$, whence
$y\in B_\alpha\subset B_{s(\alpha)}$ by recursion, or
$x\in\mathcal{P}(B_\alpha)$%%
, whence again $y\in B_\alpha\subset B_{s(\alpha)}$.

\hypertarget{sga4.exp.I.appendix.II.sec.5.theorem.1.proof.clause.U.II}{}
\item[(U.II)] Let $x,y\in U$. Since the family
$(B_\alpha)_{\alpha\in I}$ is increasing by \hyperlink{sga4.exp.I.appendix.II.sec.5.theorem.1.proof.equation.3}{(3)}, there exists
$\alpha\in I$ such that $x,y\in B_\alpha$. Then
$\{x,y\}\in\mathcal{P}(B_\alpha)\subset B_{s(\alpha)}\subset U$.

\hypertarget{sga4.exp.I.appendix.II.sec.5.theorem.1.proof.clause.U.III}{}
\item[(U.III)] Let\pageoriginale $x\in U$. Then there exists
$\alpha\in I$ such that $x\in B_\alpha$. It therefore suffices to
show that, for every $\alpha\in I$, one has:
$$
\textrm{``}x\in B_{\alpha}\Longrightarrow \mathcal{P}(x)\in
B_{s(s(\alpha))}\textrm{''\,}{\quoi}.
$$
Proceed again by transfinite recursion on $\alpha$. If
$\alpha=\mathcal{E}$, then $x\in A_n$ for some $n$, whence
$y\subset A_{n+1}\subset B=B_{\mathcal{E}}$ for every subset $y$ of
$x$; thus $y\in\mathcal{P}(B_{\mathcal{E}})\subset
B_{s(\mathcal{E})}$ for every $y\in\mathcal{P}(x)$, whence
$\mathcal{P}(x)\subset B_{s(\mathcal{E})}$ and
$\mathcal{P}(x)\in\mathcal{P}(B_{s(\mathcal{E})})\subset
B_{s(s(\mathcal{E}))}$, so our assertion is true for
$\alpha=\mathcal{E}$. %%
Passage to an element $\alpha$ with no predecessor is immediate,
since $\beta<\alpha$ then implies $s(s(\beta))<\alpha$. Finally, pass
from $\alpha$ to $s(\alpha)$; let $x\in B_{s(\alpha)}$; if
$x\in B_\alpha$, then
$\mathcal{P}(x)\in B_{s(s(\alpha))}\subset B_{s(s(s(\alpha)))}$ by
recursion; if $x\in\mathcal{P}(B_\alpha)$, then $x\subset B_\alpha$,
whence $\mathcal{P}(x)\subset\mathcal{P}(B_\alpha)$ and
$\mathcal{P}(x)\in\mathcal{P}(\mathcal{P}(B_\alpha))\in
B_{s(s(s(\alpha)))}$.

\hypertarget{sga4.exp.I.appendix.II.sec.5.theorem.1.proof.clause.U.IV}{}
\item[(U.IV)] Let $(x_\lambda)_{\lambda\in K}$ be a family of
elements of $U$ such that $K\in U$. We must show that the union
$x=\bigcup_{\lambda\in K}x_\lambda$ is an element of $U$. For every
$\lambda\in K$, choose an $\alpha(\lambda)\in I$ such that
$x_\lambda\in B_{\alpha(\lambda)}$. Let us show that the set
$\alpha(K)$ of the $\alpha(\lambda)$ is \emph{bounded above} in $I$:
indeed, if it were not, one would have
$$
I=\bigcup_{\lambda%%
\in K}[\mathcal{E},\alpha(\lambda)]{\quoi},
$$
which would contradict \hyperlink{sga4.exp.I.appendix.II.sec.5.definition.2.condition.FI.2}{(FI.2)}, since
$\card(K)<\sheaf{\sheaf{c}}$ and
$\card([\mathcal{E},\alpha(\lambda)])<\sheaf{\sheaf{c}}$ for every
$\lambda\in K$. Let then $\beta$ be an upper bound of $\alpha(K)$;
one has $x_\lambda\in B_\beta$ for every $\lambda\in K$, since the
family $(B_\alpha)_{\alpha\in I}$ is increasing. Every element of
$x_\lambda$ then belongs to $B_\beta$ by what was proved in \hyperlink{sga4.exp.I.appendix.II.sec.5.theorem.1.proof.clause.U.I}{(U.I)},
whence $x=\bigcup_{\lambda\in K}x_\lambda\subset B_\beta$. By \hyperlink{sga4.exp.I.appendix.II.sec.5.theorem.1.proof.equation.3}{(3)}, it
follows that $x\in B_{s(\beta)}$, whence $x\in U$. \cqfd
\end{itemize}

% END INCLUDED SGA4_English_current_buildable_source_20260906/source/01/drafts/SGA4_I_LINES_8096_8244_English.texfrag


% BEGIN INCLUDED SGA4_English_current_buildable_source_20260906/source/01/drafts/SGA4_I_LINES_8245_8358_English.texfrag
\hypertarget{sga4.exp.I.appendix.II.sec.5.definition.3}{}
\begin{enonce*}{DEFINITION 3}
Let\pageoriginale $x$, $y$ be two sets and let $n$ be an integer
$\geq 0$. We say that $y$ is a component of order $n$ of $x$ if there
exists a sequence $(x_j)_{j=0,\ldots,n}$ such that $x_0=x$, $x_n=y$,
and $x_{j+1}\in x_j$ for $j=0,\ldots,n-1$.
\end{enonce*}

Thus $x$ is the only component of order 0 of $x$. The components of
order 1 (\resp 2) of $x$ are the elements of $x$ (\resp the elements
of the elements of $x$). We say that $y$ is a \emph{component} of $x$
if there exists $n\geq 0$ such that $y$ is a component of order $n$ of
$x$. The relation ``$y$ is a component of $x$'' is a \emph{preorder
relation}. By the selection--union scheme (Chapter II, \ldots), the
relation ``$y$ is a component of $x$'' is collectivizing with respect
to $y$, so that the components of $x$ form a \emph{set}.

\hypertarget{sga4.exp.I.appendix.II.sec.5.definition.4}{}
\begin{enonce*}{DEFINITION 4}
Let $\sheaf{\sheaf{c}}$ be a cardinal. A set $x$ is said to be of
type $\sheaf{\sheaf{c}}$ (\resp of strict type
$\sheaf{\sheaf{c}}$, of finite type) if all the components of $x$ have
cardinalities $\leq \sheaf{\sheaf{c}}$ (\resp
$<\sheaf{\sheaf{c}}$, finite cardinalities).
\end{enonce*}

\begin{enonce*}[remark]{Examples}
The elements of the universe $U_1$ (\hyperlink{sga4.exp.I.appendix.II.sec.1}{${\rm n^{\circ}}$ 1}, \hyperlink{sga4.exp.I.appendix.II.sec.1.example.2}{Example 2})
are all of finite type. If $\sheaf{\sheaf{c}}$ is a cardinal and $x$
is of type $\sheaf{\sheaf{c}}$ (\resp of strict type
$\sheaf{\sheaf{c}}$, of finite type), then every component of $x$ and
every subset of $x$ is of type $\sheaf{\sheaf{c}}$ (\resp of strict
type $\sheaf{\sheaf{c}}$, of finite type); likewise
$\mathcal{P}(x)$ is of type $2^{\sheaf{\sheaf{c}}}$ (\resp of strict
type $2^{\sheaf{\sheaf{c}}}$, of finite type). If $x$ is of type
$\sheaf{\sheaf{c}}$ and if $\sheaf{\sheaf{c}}\leq
\sheaf{\sheaf{c}}'$, then $x$ is of type $\sheaf{\sheaf{c}}'$.
\end{enonce*}

\hypertarget{sga4.exp.I.appendix.II.sec.5.lemma.of.section.5}{}
\begin{enonce*}{LEMMA}
Let\pageoriginale $\sheaf{\sheaf{c}}$ be an uncountable strongly
inaccessible cardinal, and let $X$ be a set of strict type
$\sheaf{\sheaf{c}}$. Then there exists a cardinal
$\sheaf{\sheaf{d}}<\sheaf{\sheaf{c}}$ such that $X$ is of type
$\sheaf{\sheaf{d}}$.
\end{enonce*}

Indeed, for every integer $n$, let $\sheaf{\sheaf{d}}_n$ denote the
cardinality of the set $X_n$ of components of order $n$ of $X$. We
have $\sheaf{\sheaf{d}}_0=1$,
$\sheaf{\sheaf{d}}_1=\card(X)<\sheaf{\sheaf{c}}$, and, since
$X_n=\bigcup_{y\in X_{n-1}} y$, induction on $n$ using \hyperlink{sga4.exp.I.appendix.II.sec.5.definition.2.condition.FI.2}{(FI.2)} gives
$\sheaf{\sheaf{d}}_n<\sheaf{\sheaf{c}}$ for every $n$. The
$\sheaf{\sheaf{d}}_n$ are therefore bounded above by the cardinal
$\sheaf{\sheaf{d}}=\Sigma_{j\geq 0}\sheaf{\sheaf{d}}_j$, which is
$<\sheaf{\sheaf{c}}$ by (FI.2) and the hypothesis that
$\sheaf{\sheaf{c}}$ is uncountable. If $Y$ is a component of order
$n$ of $X$, then $\card(Y)\leq \card(X_{n+1})\leq
\sheaf{\sheaf{d}}$. \cqfd

\hypertarget{sga4.exp.I.appendix.II.sec.5.proposition.10}{}
\begin{enonce*}{PROPOSITION 10}
Let $U$ be a universe and let $\sheaf{\sheaf{c}}$ be a strongly
inaccessible cardinal. Then the set $U'$ of those $x\in U$ that are
of strict type $\sheaf{\sheaf{c}}$ is a universe.
\end{enonce*}
Indeed, if $x$, $y\in U'$ and if $z\in x$, then plainly $z\in U'$
and $\{x,y\}\in U'$. If $x\in U'$, then
$\card(x)<\sheaf{\sheaf{c}}$, whence
$\card(\mathcal{P}(x))<\sheaf{\sheaf{c}}$ by \hyperlink{sga4.exp.I.appendix.II.sec.5.definition.2.condition.FI.1}{(FI.1)}; since a
component of $\mathcal{P}(x)$ is either $\mathcal{P}(x)$ itself, a
subset of $x$, or a component of $x$, we have
$\mathcal{P}(x)\in U'$. Finally, if
$(x_{\alpha})_{\alpha\in I}$ is a family of elements of $U'$ such that
$I\in U'$, then
$\card(\bigcup_{\alpha}x_{\alpha})<\sheaf{\sheaf{c}}$ by \hyperlink{sga4.exp.I.appendix.II.sec.5.definition.2.condition.FI.2}{(FI.2)};
since every component of order $>0$ of
$\bigcup_{\alpha}x_{\alpha}$ is a component of some $x_{\alpha}$, we
indeed have $\bigcup_{\alpha}x_{\alpha}\in U'$. \cqfd

\medskip
\noindent \emph{Remark on the cardinal} $\sheaf{\sheaf{c}}(U)$.
Let $U$ be a universe and let $\sheaf{\sheaf{c}}(U)$ be the cardinal
defined by \hyperlink{sga4.exp.I.appendix.II.sec.5.equation.1.cardinal.c.U}{(1)}, \ie
$$
\sheaf{\sheaf{c}}(U)={\sup}_{x\in U}\card(x){\quoi}.
$$
We\pageoriginale plainly have $\sheaf{\sheaf{c}}(U)\leq \card(U)$,
since, as we recall, $x\in U$ implies $x\subset U$. But the equality
$\sheaf{\sheaf{c}}(U)=\card(U)$ does not always hold. For example, let
$(x_{\alpha})_{\alpha\in I}$ be an \emph{uncountable} family of
symbols subject to the axioms:
$$
\text{«} x_{\alpha}=\{x_{\alpha}\}\mbox{ for every } \alpha\in I
\text{»} \quav \text{«} x_{\alpha}\neq x_{\beta}\mbox{ if }
\alpha\neq \beta \text{»}
$$
(in other words, $x_{\alpha}$ is a one-element set, namely itself).
As in \hyperlink{sga4.exp.I.appendix.II.sec.1.example.2}{Example 2} of \hyperlink{sga4.exp.I.appendix.II.sec.1}{${\rm n^{\circ}}$ 1}, form the ``significant words''
constructed from the symbols $∅$, $x_{\alpha}$, $\{,\}$ %%
and so forth. Each denotes a \emph{finite} set. Let $U$ be the set of
these words. One readily verifies that the conditions of \hyperlink{sga4.exp.I.appendix.II.sec.1.definition.1}{Definition 1}
are satisfied, so that $U$ is a \emph{universe}. Since the significant
words are finite sequences of elements of a set of cardinality
$\card(I)+4=\card(I)$, we have $\card(U)=\card(I)$ (Chapter III; in
fact, $\card(U)\geq \card(I)$ suffices, and this is obvious). On the
other hand, $\sheaf{\sheaf{c}}(U)$ is the countably infinite cardinal,
and hence $\sheaf{\sheaf{c}}(U)<\card(U)$.

The strict inequality $\sheaf{\sheaf{c}}(U)<\card(U)$ is due to our
having introduced sets $x_{\alpha}$ such that $x_{\alpha}\in
x_{\alpha}$. If one rules out horrors of this kind, one obtains
$\card(U)=\sheaf{\sheaf{c}}(U)$ for every universe $U$, as well as
other very pretty results ``which cannot possibly be useful.''
This is what we shall do in the next section.

% END INCLUDED SGA4_English_current_buildable_source_20260906/source/01/drafts/SGA4_I_LINES_8245_8358_English.texfrag


% BEGIN INCLUDED SGA4_English_current_buildable_source_20260906/source/01/drafts/SGA4_I_LINES_8359_8513_English.texfrag
\section{Artinian sets and universes}%
\hypertarget{sga4.exp.I.appendix.II.sec.6}{}

\hypertarget{sga4.exp.I.appendix.II.sec.6.definition.5}{}
\begin{enonce*}{DEFINITION 5}
A set $x$ is said to be Artinian if there is no infinite sequence
$(x_n)_{n\geq 0}$ such that $x_{\circ}=x$ and $x_{n+1}\in
x_n$ for every $n\geq 0$.
\end{enonce*}

\begin{enonce*}[remark]{Examples}
The sets $∅$, $\{∅\}$,
$\{∅,\{∅\}\}$, and, more generally, the elements of
the universe $U_1$ of \hyperlink{sga4.exp.I.appendix.II.sec.1}{${\rm n^{\circ}}$ 1} are Artinian.
\end{enonce*}

In\pageoriginale pictorial terms, if $x$ is Artinian, and if one takes
an element $x_1$ of $x$, then an element $x_2$ of $x_1$, \etc, the
process must stop, and one arrives at a component $x_n$ of $x$
that is \emph{empty}. In other words, Artinian sets are
``built from $∅$''; this will be made more precise
later (\cf \hyperlink{sga4.exp.I.appendix.II.sec.6.theorem.2.proof.claim.star}{(*)} in the proof of \hyperlink{sga4.exp.I.appendix.II.sec.6.theorem.2}{Theorem 2}).

Artinian sets plainly have the following properties:
\begin{itemize}
\hypertarget{sga4.exp.I.appendix.II.sec.6.property.AR.I}{}
\item[(AR.I)] If $x$ is Artinian, every subset of $x$ and every
component of $x$ is Artinian. For $x$ to be Artinian, it is necessary
and sufficient that every element of $x$ be Artinian.

\hypertarget{sga4.exp.I.appendix.II.sec.6.property.AR.II}{}
\item[(AR.II)] If $x$ and $y$ are Artinian, then $\{x,y\}$ is
Artinian.

\hypertarget{sga4.exp.I.appendix.II.sec.6.property.AR.III}{}
\item[(AR.III)] If $x$ is Artinian, $\mathcal{P}(x)$ is Artinian.

\hypertarget{sga4.exp.I.appendix.II.sec.6.property.AR.IV}{}
\item[(AR.IV)] Every union of Artinian sets is an Artinian set.
\end{itemize}

These properties immediately yield the following result:

\hypertarget{sga4.exp.I.appendix.II.sec.6.proposition.11}{}
\begin{enonce*}{PROPOSITION 11}
If $U$ is a universe, the set of Artinian elements of $U$ is
a universe, necessarily Artinian.
\end{enonce*}

\hypertarget{sga4.exp.I.appendix.II.sec.6.corollary.to.proposition.11}{}
\begin{enonce*}{COROLLARY}
If $x$ is an Artinian set, there exists an Artinian universe $U$
such that $x\in U$.
\end{enonce*}

Indeed, by axiom \hyperlink{sga4.exp.I.appendix.II.sec.4.axiom.A.6}{(A.6)}, $x$ is an element of a universe $V$; take
for $U$ the set of Artinian elements of $V$.

The following proposition is even less useful than the rest of ${\rm
n^{\circ}}$ 6:

\hypertarget{sga4.exp.I.appendix.II.sec.6.proposition.12}{}
\begin{enonce*}{PROPOSITION 12}
Let\pageoriginale $A$ be an Artinian set. Then:
\begin{enumerate}
\renewcommand{\theenumi}{\alph{enumi}}
\renewcommand{\labelenumi}{\theenumi)}
\hypertarget{sga4.exp.I.appendix.II.sec.6.proposition.12.part.a}{}
\item For every $x\in A$, one has $x\not\in x$;

\hypertarget{sga4.exp.I.appendix.II.sec.6.proposition.12.part.b}{}
\item If $x$, $y\in A$, one cannot have both $x\in y$ and $y\in
x$;

\hypertarget{sga4.exp.I.appendix.II.sec.6.proposition.12.part.c}{}
\item For every $x\in A$, the relation ``$y$ is a component of $z$'' on the components $y$, $z$ of $x$ is an order relation;

\hypertarget{sga4.exp.I.appendix.II.sec.6.proposition.12.part.d}{}
\item For every nonempty element $x$ of $A$, there exists $y\in x$
such that $x\cap y=∅$.
\end{enumerate}
\end{enonce*}
Indeed, the negation of a) (\resp of b)) entails the existence of an
infinite sequence $(x_n)_{n\geq 0}$ contradicting \hyperlink{sga4.exp.I.appendix.II.sec.6.definition.5}{Definition 5}, namely
$(x,x,x,\ldots)$ (\resp $(x,y,x,y,x,y,\ldots)$). The negation of c)
means that there exist $x\in A$, distinct components $y$, $z$ of $x$,
and membership chains:
$$
y\in y_1\in\ldots \in y_q\in z\quav z\in z_1\in\ldots\in z_r\in
y{\quoi};
$$
whence, as in \hyperlink{sga4.exp.I.appendix.II.sec.6.proposition.12.part.b}{b)}, an infinite sequence $(x_n)_{n\geq 0}$
contradicting \hyperlink{sga4.exp.I.appendix.II.sec.6.definition.5}{Definition 5}. Finally, if d) is false, there exists a
nonempty element $x$ of $A$ such that $x\cap y\neq ∅$ for
every $y\in x$; set $x_\circ=x$, and take for $x_1$ an
element of $x$; since $x\cap x_1\neq ∅$, take for
$x_2$ an element of $x\cap x_1$, and so on; more formally, define
by induction an infinite sequence $(x_n)_{n\geq 1}$
of elements of $x$ by means of $x_1\in x$, $x_{n+1}\in x_n\cap x$
for $n\geq 1$; then the sequence $(x_n)_{n\geq 0}$ contradicts
\hyperlink{sga4.exp.I.appendix.II.sec.6.definition.5}{Definition 5}.

\begin{enonce*}[remark]{Remark}
Let $B$ be a set. Then $B$ is Artinian if and only if every set $A$
whose elements are sets of components of $B$ satisfies \hyperlink{sga4.exp.I.appendix.II.sec.6.proposition.12.part.d}{condition d)}
of \hyperlink{sga4.exp.I.appendix.II.sec.6.proposition.12}{Proposition 12}. Indeed, necessity
follows from \hyperlink{sga4.exp.I.appendix.II.sec.6.property.AR.I}{(AR.I)} and\pageoriginale \hyperlink{sga4.exp.I.appendix.II.sec.6.proposition.12}{Proposition 12}. Conversely,
if $B$ is not Artinian, there exists an infinite sequence $(x_n)_{n\geq
0}$ with $x_{\circ}=B$ and $x_{n+1}\in x_n$ for every $n\geq 0$; then
take for $A$ the set whose sole element is the set $X$ of the $x_n$;
thus $A$ contains a nonempty element $X$ such that, for every
element $y$ of $X$, one has $y\cap X\neq ∅$ (indeed, $y$
is of the form $x_n$, and one has $x_{n+1}\in y\cap X$).
\end{enonce*}

\hypertarget{sga4.exp.I.appendix.II.sec.6.theorem.2}{}
\begin{enonce*}{THEOREM 2}
Let $\sheaf{\sheaf{c}}$ be an infinite cardinal. Then:
\begin{itemize}
\hypertarget{sga4.exp.I.appendix.II.sec.6.theorem.2.part.a}{}
\item[a)] The relation ``$x$ is an Artinian set of type $\sheaf{\sheaf{c}}$
(\resp of strict type $\sheaf{\sheaf{c}}$)'' %%
is collectivizing with respect to $x$; the set
$A_{\sheaf{\sheaf{c}}}$ of Artinian sets of type
$\sheaf{\sheaf{c}}$ has cardinality $2^{\sheaf{\sheaf{c}}}$.

\hypertarget{sga4.exp.I.appendix.II.sec.6.theorem.2.part.b}{}
\item[b)] If $\sheaf{\sheaf{c}}$ is strongly inaccessible, the set
$U_{\sheaf{\sheaf{c}}}$ of Artinian sets of strict type
$\sheaf{\sheaf{c}}$ is a universe of cardinality $\sheaf{\sheaf{c}}$;
the cardinal $\sheaf{\sheaf{c}}(U_{\sheaf{\sheaf{c}}})=\sup_{x\in
U_{\sheaf{\sheaf{c}}}}\card(x)$ is $\sheaf{\sheaf{c}}$.

\hypertarget{sga4.exp.I.appendix.II.sec.6.theorem.2.part.c}{}
\item[c)] If a universe $U$ contains an element of cardinality $\sheaf{\sheaf{c}}$,
every Artinian set of type $\sheaf{\sheaf{c}}$ belongs to $U$
(in other words, $A_{\sheaf{\sheaf{c}}}\subset U$).

\hypertarget{sga4.exp.I.appendix.II.sec.6.theorem.2.part.c-prime}{}
\item[${\rm c')}$] If a universe $U$ is nonempty, every Artinian
set of finite type is an element of $U$.
\end{itemize}
\end{enonce*}
Before proving \hyperlink{sga4.exp.I.appendix.II.sec.6.theorem.2}{Theorem 2}, let us deduce from it a few illuminating
corollaries:

\hypertarget{sga4.exp.I.appendix.II.sec.6.corollary.to.theorem.2.1}{}
\begin{enonce*}{COROLLARY 1}
If a universe $U$ is Artinian, then $\card(U)$ is strongly
inaccessible, and $U$ is the set of Artinian sets of strict type
$\card(U)$.
\end{enonce*}

Indeed,\pageoriginale set $\sheaf{\sheaf{c}}(U)=\sup_{x\in
U}\card(x)$. This is a strongly inaccessible cardinal (beginning of
\hyperlink{sga4.exp.I.appendix.II.sec.5}{${\rm n^{\circ}}$ 5}). First suppose that it is uncountable; then,
for every infinite cardinal $\sheaf{\sheaf{c}}<\sheaf{\sheaf{c}}(U)$,
every Artinian set of type $\sheaf{\sheaf{c}}$ is an element of
$U$ by \hyperlink{sga4.exp.I.appendix.II.sec.6.theorem.2.part.c}{c)}; hence every Artinian set of strict type
$\sheaf{\sheaf{c}}(U)$ is an element of $U$ by the \hyperlink{sga4.exp.I.appendix.II.sec.5.lemma.of.section.5}{lemma of ${\rm
n^{\circ}}$ 5}. This last assertion remains valid if
$\sheaf{\sheaf{c}}(U)$ is countable, by \hyperlink{sga4.exp.I.appendix.II.sec.6.theorem.2.part.c-prime}{${\rm c'}$)}. Conversely,
by (U.I), every set in $U$ is of strict type
$\sheaf{\sheaf{c}}(U)$. Thus $U$ is the universe
$U_{\sheaf{\sheaf{c}}(U)}$ of \hyperlink{sga4.exp.I.appendix.II.sec.6.theorem.2.part.b}{b)}, whence
$\sheaf{\sheaf{c}}(U)=\card(U)$ by b).

It follows from \hyperlink{sga4.exp.I.appendix.II.sec.6.corollary.to.theorem.2.1}{Corollary 1} that an Artinian universe is uniquely
\emph{determined} by its cardinality (which is moreover strongly
inaccessible). Thus there is a ``one-to-one correspondence'' between
Artinian universes and strongly inaccessible cardinals. In particular:

% END INCLUDED SGA4_English_current_buildable_source_20260906/source/01/drafts/SGA4_I_LINES_8359_8513_English.texfrag


% BEGIN INCLUDED SGA4_English_current_buildable_source_20260906/source/01/drafts/SGA4_I_LINES_8514_8642_English.texfrag
\hypertarget{sga4.exp.I.appendix.II.sec.6.corollary.to.theorem.2.2}{}
\begin{enonce*}{COROLLARY 2}
The inclusion relation $U\subset U'$ on Artinian universes is a
well-ordering relation.
\end{enonce*}

Indeed, the relation $\sheaf{\sheaf{c}}\leq \sheaf{\sheaf{c}}'$ on
cardinals is a well-ordering relation (Chapter III), and, with the
notation of \hyperlink{sga4.exp.I.appendix.II.sec.6.theorem.2.part.b}{b)} of \hyperlink{sga4.exp.I.appendix.II.sec.6.theorem.2}{Theorem 2}, the relations $\sheaf{\sheaf{c}}\leq
\sheaf{\sheaf{c}}'$ and $U_{\sheaf{\sheaf{c}}}\subset
U_{\sheaf{\sheaf{c}}'}$ are equivalent.

Observe that \hyperlink{sga4.exp.I.appendix.II.sec.6.theorem.2.part.b}{b)} of \hyperlink{sga4.exp.I.appendix.II.sec.6.theorem.2}{Theorem 2} gives a second proof of \hyperlink{sga4.exp.I.appendix.II.sec.5.theorem.1}{Theorem 1}
(\hyperlink{sga4.exp.I.appendix.II.sec.5}{${\rm n^{\circ}}$ 5}).

We now turn to the proof of \hyperlink{sga4.exp.I.appendix.II.sec.6.theorem.2}{Theorem 2}. Given a set $A$, we shall call
a \emph{chain} of $A$ any finite sequence
$(x_j)_{j=0,\ldots,n}$ such that $x_0=A$ and $x_{j+1}\in x_j$
for $j=0,\ldots,n-1$. The chains of $A$ form a\pageoriginale
\emph{set}, by the selection--union scheme; denote it by
$G(A)$. Given a chain $X=(x_j)_{j=0,\ldots,n}$, the chains of the
form $(x_i)_{i=0,\ldots,q}$ with $q\leq n$ will be called less than
$X$; this endows $G(A)$ with the structure of an \emph{ordered set}.
For $A$ to be Artinian, it is necessary and sufficient that $G(A)$
be a ``\emph{Noetherian}'' ordered set (i.e. one satisfying the
equivalent conditions of Chapter III, \S 6, ${\rm n^{\circ}}$ 5).

We shall prove that:

\hypertarget{sga4.exp.I.appendix.II.sec.6.theorem.2.proof.claim.star}{}
\noindent (*) \emph{If $A$ is Artinian, it is uniquely determined by
the isomorphism class of the ordered set $G(A)$.}

Indeed, given an ordered set $G$ and an element $g\in G$, let $S(g)$
denote the set of all $g'\in G$ such that $g<g'$ and such that
$g\leq h\leq g'$ implies $h=g$ or $h=g'$ (in other words, the set
of ``immediate successors'' of $g$). Consider the map $\theta$ that
assigns to every chain $X=(x_0,\ldots,x_n)$ in $G(A)$ the set $x_n$;
then, for every $X\in G(A)$, we have:
$$
\theta(X)=\{\theta(X')| X'\in S(X)\}{\quoi}.
$$
Since $A$ is the image under $\theta$ of the least element $X_0=(A)$
of $G(A)$, it will suffice to show that $\theta$ is uniquely
determined by the isomorphism class of the ordered set $G(A)$. This
follows from the following lemma:

\hypertarget{sga4.exp.I.appendix.II.sec.6.lemma.1}{}
\begin{enonce*}{LEMMA 1}
Let\pageoriginale $G$ be a Noetherian ordered set and let $\varphi$
be a function with domain $G$ such that, for every $g\in G$, one has:
\hypertarget{sga4.exp.I.appendix.II.sec.6.lemma.1.equation.1}{}
\begin{equation}
\varphi(g)=\{\varphi(g')| g'\in S(g)\}{\quoi}.\tag{1}
\end{equation}
Then $\varphi$ is uniquely determined. Moreover, if $U$ is a universe
containing an element equipotent to $G$, then $\varphi$ takes its
values in $U$.
\end{enonce*}

The hypothesis implies that $\varphi(g)=∅$ if $g$ is a maximal
element of $G$.

Indeed, let $\varphi$ and $\varphi'$ be two functions satisfying
\hyperlink{sga4.exp.I.appendix.II.sec.6.lemma.1.equation.1}{(1)}. If $\varphi\neq \varphi'$, the set of all $g\in G$ such that
$\varphi(g)\neq \varphi'(g)$ is nonempty and hence, since $G$ is
Noetherian, has a \emph{maximal} element $h$. We then have
$\varphi(g)=\varphi'(g)$ for every $g>h$, in particular for every
immediate successor $g\in S(h)$; hence $\varphi(h)=\varphi'(h)$ by
\hyperlink{sga4.exp.I.appendix.II.sec.6.lemma.1.equation.1}{(1)}, which is a contradiction. Thus indeed
$\varphi=\varphi'$. Now let $U$ be a universe containing an element
$x$ equipotent to $G$; let us show that $\varphi$ takes its values in
$U$. Otherwise, let $h$ be a maximal element among those $g\in G$
such that $\varphi(g)\not\in U$. Then $\varphi(g')\in U$ for every
$g'\in S(h)$, so that
$$
\varphi(h)=\{\varphi(g')| g'\in S(h)\}
$$
is a subset of $U$. Since its cardinality does not exceed
$\card(G)$, and hence does not exceed the cardinality of an element
of $U$, we have $\varphi(h)\in U$ (\hyperlink{sga4.exp.I.appendix.II.sec.1}{${\rm n^{\circ}}$ 1},
\hyperlink{sga4.exp.I.appendix.II.sec.1.proposition.7}{Proposition 7}). This contradiction shows that $\varphi$ takes its
values in $U$.

With this \pageoriginale %%
established, let us prove a) of \hyperlink{sga4.exp.I.appendix.II.sec.6.theorem.2}{Theorem 2}. We may confine ourselves to
the assertion without the \resp clause, since the other follows at
once. Let $\sheaf{\sheaf{c}}$ be an \emph{infinite} cardinal. If $A$
is a set of type $\sheaf{\sheaf{c}}$, then:
\hypertarget{sga4.exp.I.appendix.II.sec.6.theorem.2.proof.equation.2}{}
\begin{equation}
\card(G(A))\leq \sheaf{\sheaf{c}}{\quoi}%%
\ .\tag{2}
\end{equation}
Indeed, if $A_n$ denotes the set of components of order $n$ of $A$,
then $\card(A_1)\leq \sheaf{\sheaf{c}}$ and $\card(A_{n+1})\leq
\sheaf{\sheaf{c}}\cdot \card(A_n)$, whence $\card(A_n)\leq
\sheaf{\sheaf{c}}^n$ by induction. But
$\sheaf{\sheaf{c}}^n=\sheaf{\sheaf{c}}$, since
$\sheaf{\sheaf{c}}$ is infinite (Chapter III); hence
$\card(\bigcup^\infty_{n=0} A_n)\leq
\card(\NN)\cdot \sheaf{\sheaf{c}}=\sheaf{\sheaf{c}}$ (Chapter III).
Now $G(A)$ is a set of \emph{finite} sequences of elements of
$\bigcup_n A_n$, so inequality (2) indeed holds (Chapter III).
Now, if $E$ is a set of cardinality $\sheaf{\sheaf{c}}$, specifying
an order structure on a subset $E'$ of $E$ is equivalent to specifying
the subset of $E\times E$ consisting of the $(x,y)$ such that $x\in
E'$, $y\in E'$, and $x\leq y$. Consequently, by \hyperlink{sga4.exp.I.appendix.II.sec.6.theorem.2.proof.equation.2}{(2)}, the isomorphism
classes of the ordered sets $G(A)$ (where $A$ is of type
$\sheaf{\sheaf{c}}$) form a set
$\mathfrak{F}_{\sheaf{\sheaf{c}}}$, and
$\card(\mathfrak{F}_{\sheaf{\sheaf{c}}})\leq
\card(\mathcal{P}(E\times E))=2^{\sheaf{\sheaf{c}}}$. Let
$\mathfrak{F}'_{\sheaf{\sheaf{c}}}$ be the subset of
$\mathfrak{F}_{\sheaf{\sheaf{c}}}$ consisting of the classes of
Noetherian ordered sets having a least element. If, to every $G\in
\mathfrak{F}'_{\sheaf{\sheaf{c}}}$, one assigns the value of the
function $\varphi$ of \hyperlink{sga4.exp.I.appendix.II.sec.6.lemma.1}{Lemma 1} at the least element of $G$, one obtains
a map $\theta$ on $\mathfrak{F}'_{\sheaf{\sheaf{c}}}$ whose image
contains \emph{all} Artinian sets of type $\sheaf{\sheaf{c}}$. The
latter therefore do indeed form a set $A_{\sheaf{\sheaf{c}}}$, and
we have $\card(A_{\sheaf{\sheaf{c}}})\leq
\card(\mathfrak{F}'_{\sheaf{\sheaf{c}}})\leq
\card(\mathfrak{F}_{\sheaf{\sheaf{c}}})\leq
2^{\sheaf{\sheaf{c}}}$.

It remains %%
\pageoriginale to show that
$\card(A_{\sheaf{\sheaf{c}}})=2^{\sheaf{\sheaf{c}}}$. For this it
suffices to show that there exists an Artinian set $B$ of type
$\sheaf{\sheaf{c}}$ and cardinality $\sheaf{\sheaf{c}}$, since the
subsets of $B$ will then be elements of $A_{\sheaf{\sheaf{c}}}$.
This existence follows from the following lemma:

% END INCLUDED SGA4_English_current_buildable_source_20260906/source/01/drafts/SGA4_I_LINES_8514_8642_English.texfrag


% BEGIN INCLUDED SGA4_English_current_buildable_source_20260906/source/01/drafts/SGA4_I_LINES_8643_8705_English.texfrag
\hypertarget{sga4.exp.I.appendix.II.sec.6.lemma.2}{}
\begin{enonce*}{LEMMA 2}
For every cardinal $\sheaf{\sheaf{c}}$, there exists an Artinian set
$B$ of type $\sheaf{\sheaf{c}}$ and cardinality
$\sheaf{\sheaf{c}}$.
\end{enonce*}

Proceed by transfinite induction on $\sheaf{\sheaf{c}}$. For
$\sheaf{\sheaf{c}}=0$ there is no choice, and one takes $B=∅$. If
$\sheaf{\sheaf{c}}$ has a predecessor $\sheaf{\sheaf{c}}'$, let $B'$
be an Artinian set of type $\sheaf{\sheaf{c}}'$ and cardinality
$\sheaf{\sheaf{c}}'$; one then has
$\sheaf{\sheaf{c}}\leq2^{\sheaf{\sheaf{c}}'}$, so there exists a
subset $B$ of $\mathcal{P}(B')$ of cardinality
$\sheaf{\sheaf{c}}$; the elements of $B$ are subsets of $B'$ and
therefore have cardinalities
$\leq\sheaf{\sheaf{c}}'\leq\sheaf{\sheaf{c}}$; the higher-order
components of $B$ are components of $B'$, and hence also have
cardinalities $\leq\sheaf{\sheaf{c}}'\leq\sheaf{\sheaf{c}}$.
Finally, if $\sheaf{\sheaf{c}}$ has no predecessor, choose, for every
cardinal $\sheaf{\sheaf{c}}_\lambda<\sheaf{\sheaf{c}}$, an Artinian
set $B_\lambda$ of type $\sheaf{\sheaf{c}}_\lambda$ and cardinality
$\sheaf{\sheaf{c}}_\lambda$; then
$B=\bigcup_\lambda B_\lambda$ has the required property. This proves
\hyperlink{sga4.exp.I.appendix.II.sec.6.lemma.2}{Lemma 2} and completes the proof of part a).

We turn to b). Let $\sheaf{\sheaf{c}}$ be a strongly inaccessible
cardinal. By \hyperlink{sga4.exp.I.appendix.II.sec.6.theorem.2.part.a}{a)}, we already know that the Artinian sets of strict type
$\sheaf{\sheaf{c}}$ form a set $U_{\sheaf{\sheaf{c}}}$. The fact that
$U_{\sheaf{\sheaf{c}}}$ is a universe follows immediately from
properties \hyperlink{sga4.exp.I.appendix.II.sec.6.property.AR.I}{(AR.I)} to \hyperlink{sga4.exp.I.appendix.II.sec.6.property.AR.IV}{(AR.IV)} of Artinian sets (beginning of
\hyperlink{sga4.exp.I.appendix.II.sec.6}{${\rm n^{\circ}}$}), and\pageoriginale from cardinal estimates analogous
to those of \hyperlink{sga4.exp.I.appendix.II.sec.5.proposition.10}{Proposition 10} of \hyperlink{sga4.exp.I.appendix.II.sec.5}{${\rm n^{\circ}}$~5}. The relation
$\sup_{x\in U_{\sheaf{\sheaf{c}}}}\card(x)=\sheaf{\sheaf{c}}$ follows%%
 from \hyperlink{sga4.exp.I.appendix.II.sec.6.lemma.2}{Lemma 2}, applied to cardinals $<\sheaf{\sheaf{c}}$. Finally, to
show that
$\card(U_{\sheaf{\sheaf{c}}})=\sheaf{\sheaf{c}}$, first suppose that
$\sheaf{\sheaf{c}}$ is \emph{uncountable}; by%%
 the \hyperlink{sga4.exp.I.appendix.II.sec.5.lemma.of.section.5}{lemma of ${\rm n^{\circ}}$~5},
$U_{\sheaf{\sheaf{c}}}$ is the union
$\bigcup_{\sheaf{\sheaf{d}}<\sheaf{\sheaf{c}}}
A_{\sheaf{\sheaf{d}}}$, where $A_{\sheaf{\sheaf{d}}}$ denotes the
set of Artinian sets of type $\sheaf{\sheaf{d}}$; now
$\card(A_{\sheaf{\sheaf{d}}})=2^{\sheaf{\sheaf{d}}}
<\sheaf{\sheaf{c}}$ (by \hyperlink{sga4.exp.I.appendix.II.sec.6.theorem.2.part.a}{a)}); on the other hand, the set of cardinals
$\sheaf{\sheaf{d}}<\sheaf{\sheaf{c}}$ has cardinality
$\leq\sheaf{\sheaf{c}}$ (Chapter III); whence
$\card(U_{\sheaf{\sheaf{c}}})\leq
\sheaf{\sheaf{c}}\cdot\sheaf{\sheaf{c}}=\sheaf{\sheaf{c}}$, and also
$\card(U_{\sheaf{\sheaf{c}}})\geq\sheaf{\sheaf{c}}$, since
$\card(U_{\sheaf{\sheaf{c}}})\geq
\card(A_{\sheaf{\sheaf{d}}})=2^{\sheaf{\sheaf{d}}}$ for every
$\sheaf{\sheaf{d}}<\sheaf{\sheaf{c}}$.


The case $\sheaf{\sheaf{c}}=0$ being trivial, there remains the case
where $\sheaf{\sheaf{c}}$ is the \emph{countably infinite cardinal.}
In this case $U_{\sheaf{\sheaf{c}}}$ is the set of Artinian sets of
finite type (\ie finite together with all their components), and one
uses a pleasant result of a combinatorial nature.


% END INCLUDED SGA4_English_current_buildable_source_20260906/source/01/drafts/SGA4_I_LINES_8643_8705_English.texfrag


% BEGIN INCLUDED SGA4_English_current_buildable_source_20260906/source/01/drafts/SGA4_I_LINES_8706_8780_English.texfrag
\hypertarget{sga4.exp.I.appendix.II.sec.6.lemma.3}{}
\begin{enonce*}{LEMMA 3}
{\rm (D. K"onig?)}. Consider two infinite sequences
$(E_n)_{n\geq1}$, $(f_n)_{n\geq1}$ of finite sets $E_n$ and maps
$f_n:E_{n+1}\rightarrow E_n$. If there exists no infinite sequence
$(x_n)_{n%%
\geq1}$ such that $x_n\in E_n$ and $f_n(x_{n+1})=x_n$ for every
$n$, then $E_n$ is empty for all sufficiently large $n$.
\end{enonce*}

In other words, if all sequences $(x_n)$ such that
$f_n(x_{n+1})=x_n$ are \emph{finite}, their lengths are
\emph{bounded}. This can be expressed in terms of projective limits: a
projective limit of nonempty finite sets is nonempty (cf. General
Topology, Chapter I, 2nd ed., \S 9, ${\rm n^{\circ}}$ 6,
Proposition 8, $2^{\circ}$)).

Indeed,\pageoriginale argue by contradiction. If there exist nonempty
$E_n$ for arbitrarily large $n$, no $E_n$ is empty (since
$E_n=∅$ implies $E_{n+1}=∅$ in view of the existence of
$f_n:E_{n+1}\rightarrow E_n$). Call a finite sequence
$(x_j)_{1\leq j\leq n}$ \emph{coherent} if
$f_j(x_{j+1})=x_j$ for $j=1,\ldots,n-1$. We prove by induction on $n$
the existence of a coherent sequence $(a_1,\ldots,a_n)$
$(a_i\in E_i)$ that, for every $q\geq n$, can be extended to a
coherent sequence $(a_1,\ldots,a_n,x_{n+1},\ldots,x_q)$ of length
$q$. This is evident for $n=0$, since no $E_q$ is empty. Pass from $n$
to $n+1$. If, for every $x\in f_n^{-1}(\{a_n\})\subset E_{n+1}$,
all coherent sequences extending $(a_1,\ldots,a_n,x)$ had lengths
bounded by an integer $q(x)$, then all coherent sequences extending
$(a_1,\ldots,a_n)$ would have bounded lengths (by $\sup_xq(x)$), since
$E_{n+1}$ is finite; hence there exists
$a_{n+1}\in f_n^{-1}(\{a_n\})$ such that the coherent sequence
$(a_1,\ldots,a_n,a_{n+1})$ admits extensions of arbitrary length. We
then obtain an infinite sequence $(a_n)_{n%%
\geq1}$ contradicting the hypothesis.

It follows from \hyperlink{sga4.exp.I.appendix.II.sec.6.lemma.3}{Lemma 3} that if $A$ is an Artinian set of finite type,
then the ordered set $G(A)$ of its chains is \emph{finite}: indeed,
take for $E_n$ the set of chains
$(x_n\in x_{n-1}\in\ldots\in x_1\in A)$ with $n+1$ terms (which is
finite because the components of $A$ of order $\leq n+1$ are finite
in number and are all finite), and for $f_n$ the map
$(x_{n+1}\in x_n\in x_{n-1}\in\ldots)\mapsto
(x_n\in x_{n-1}\in\ldots)$. Now the set of isomorphism classes of
finite\pageoriginale ordered sets is \emph{countable}: indeed, giving
an order structure on a finite subset of $\NN$ is equivalent to giving
its graph, which is a finite subset of $\NN\times\NN$; on the other
hand, the set of finite subsets of a countable set is countable
(Chapter III). It follows from \hyperlink{sga4.exp.I.appendix.II.sec.6.theorem.2.proof.claim.star}{(*)} and \hyperlink{sga4.exp.I.appendix.II.sec.6.lemma.1}{Lemma 1} that the set
$\mathfrak{a}$ of Artinian sets of finite type is countable. It is
infinite by \hyperlink{sga4.exp.I.appendix.II.sec.6.lemma.2}{Lemma 2} (or, more simply, by using the sequence
$(z_n)_{n\geq0}$ defined by $z_0=∅$, $z_{n+1}=\{z_n\}$). This
completes the proof of b).

\begin{enonce*}[remark]{Remark}
We have just proved that if $A$ is an Artinian set of finite type, it
has only finitely many components. There therefore exists an integer
$n$ such that $A$ is of type $n$.
\end{enonce*}
We turn to the proof of c). Let $\sheaf{\sheaf{c}}$ be an infinite
cardinal, let $U$ be a universe admitting an element of cardinality
$\sheaf{\sheaf{c}}$, and let $A$ be an Artinian set of type
$\sheaf{\sheaf{c}}$. Then the ordered set $G(A)$ has cardinality
$\leq\sheaf{\sheaf{c}}$ (\hyperlink{sga4.exp.I.appendix.II.sec.6.theorem.2.proof.equation.2}{formula (2)} above). The second assertion of
\hyperlink{sga4.exp.I.appendix.II.sec.6.lemma.1}{Lemma 1}, applied to $G(A)$, shows that the map
$(x_0,\ldots,x_n)\rightsquigarrow x_n$ from $G(A)$ takes its values in
$U$. In other words, all the components of $A$ are elements of $U$.
This proves c). The proof of \hyperlink{sga4.exp.I.appendix.II.sec.6.theorem.2.part.c-prime}{${\rm c'}$)} is analogous: if $A$ is an
Artinian set of finite type, we have just seen that $G(A)$ is finite;
since $U$ is nonempty, it contains an element equipotent to $G(A)$ by
the \hyperlink{sga4.exp.I.appendix.II.sec.1.corollary.to.proposition.7}{corollary to Proposition 7} (\hyperlink{sga4.exp.I.appendix.II.sec.1}{${\rm n^{\circ}}$ 1}). \cqfd


% END INCLUDED SGA4_English_current_buildable_source_20260906/source/01/drafts/SGA4_I_LINES_8706_8780_English.texfrag


% BEGIN INCLUDED SGA4_English_current_buildable_source_20260906/source/01/drafts/SGA4_I_LINES_8781_8916_English.texfrag
\section{Vague metamathematical remarks}%
\hypertarget{sga4.exp.I.appendix.II.sec.7}{}\pageoriginale

\begin{enumerate}
\renewcommand{\theenumi}{\alph{enumi}}
\renewcommand{\labelenumi}{\theenumi)}
\item The axiom ``\emph{every set is Artinian}'' is
\emph{harmless}. Indeed, if one has a model $M$ of Set Theory, the set
$M'$ of the Artinian elements of $M$ is also a model (\cf \hyperlink{sga4.exp.I.appendix.II.sec.6.proposition.11}{Proposition
11}).

\item The universe axiom \hyperlink{sga4.exp.I.appendix.II.sec.4.axiom.A.6}{(A.6)} (and the equivalent axiom \hyperlink{sga4.exp.I.appendix.II.sec.5.axiom.A-prime.6}{$(A'.6)$} on
strongly inaccessible cardinals) is \emph{independent} of the rest of
Set Theory. Indeed, let $\sheaf{\sheaf{c}}$ be the first uncountable
strongly inaccessible cardinal. The universe
$U_{\sheaf{\sheaf{c}}}$ of Artinian sets of strict type
$\sheaf{\sheaf{c}}$ (\hyperlink{sga4.exp.I.appendix.II.sec.6.theorem.2}{Theorem 2}, \hyperlink{sga4.exp.I.appendix.II.sec.6.theorem.2.part.b}{b)}) is a model of Set Theory without
\hyperlink{sga4.exp.I.appendix.II.sec.4.axiom.A.6}{(A.6)}: the elements of $U_{\sheaf{\sheaf{c}}}$ are called ``sets,''
the ``membership relation'' is the restriction to
$U_{\sheaf{\sheaf{c}}}$ of the usual one, \etc The ``universes'' of
the model are therefore the ordinary universes that are elements of
$U_{\sheaf{\sheaf{c}}}$. But we have seen that the only universes
which are elements of $U_{\sheaf{\sheaf{c}}}$ are the two
peque${\rm\widetilde{n}}$os $U_0=∅$ and $U_1$. Thus $U_1$ is a
``set'' which is not an element of any ``universe.'' We therefore
have a model of Set Theory in which \hyperlink{sga4.exp.I.appendix.II.sec.4.axiom.A.6}{(A.6)} is false.

\item Bourbaki was too cautious in merely ``presuming'' that the
\emph{axiom of infinity} (A.5) is \emph{independent} of the preceding
axioms and schemes. It is indeed independent, since the countable
universe $U_1$ of Artinian sets of finite type is a model in which
(A.5) is false and the preceding axioms and schemes are true.

\item It would be very interesting to prove that the universe axiom
\hyperlink{sga4.exp.I.appendix.II.sec.4.axiom.A.6}{(A.6)} is harmless. That looks difficult---and in fact it is
unprovable, says Paul Cohen.
\end{enumerate}

The adjective\pageoriginale ``vague'' in the title means that, when
constructing models, we did not take the trouble to fiddle with the
symbol $\tau$ so that it could not take us outside the model. The
trick, if one is considering a model $M$, is to replace the ordinary
$\tau_x(ℛ(x))$ by:
$$
\tau_x(ℛ(x)\mbox{ and } x\in M){\quoi}.
$$
One still has to check that this really turns the ordinary
quantifiers into the quantifiers formerly called ``typical'':
$$
(\forall x \in M)\quav (\exists x \in M){\quoi}.
$$

\hypertarget{sga4.exp.I.appendix.II.sec.7.exercises}{}
\begin{enonce*}[remark]{Exercises}
\begin{enumerate}
\renewcommand{\labelenumi}{\theenumi)}
\hypertarget{sga4.exp.I.appendix.II.sec.7.exercise.1}{}
\item Let $n$ be an integer $\geq 1$. Show that the set of Artinian
sets of type $n$ is infinite (the sets $z_0=∅$, $z_1=\{∅\}$,
$z_{q+1}=\{z_q\}$ are of type 1).

\hypertarget{sga4.exp.I.appendix.II.sec.7.exercise.2}{}
\item Define the \emph{height} of a set $A$ to be the supremum
(finite or infinite) of the integers $n$ for which there exists a
sequence $x_n\in x_{n-1}\in \ldots \in x_0=A$. Show that the sets of
height $\leq n$ are finite and form a finite set, whose cardinality
$p_n$ can be calculated from $p_0=1$, $p_{n+1}=2^{p_n}$ (proceed by
induction on $n$, noting that the elements of a set $A$ of height
$\leq n$ are sets of height $\leq n-1$).

\hypertarget{sga4.exp.I.appendix.II.sec.7.exercise.3}{}
\item Let $G$ be a Noetherian ordered set having a least element
$g_0$, and suppose that, for every $h\in G$, the set of $g\in G$ such
that $g\leq h$ is finite and totally ordered.
\begin{enumerate}
\renewcommand{\theenumii}{\alph{enumii}}
\renewcommand{\labelenumii}{\theenumii)}
\hypertarget{sga4.exp.I.appendix.II.sec.7.exercise.3.part.a}{}
\item Show\pageoriginale that every element $h\neq g_0$ of $G$ has
a unique predecessor.

\hypertarget{sga4.exp.I.appendix.II.sec.7.exercise.3.part.b}{}
\item Let $\varphi$ be the map on $G$ defined in \hyperlink{sga4.exp.I.appendix.II.sec.6.lemma.1}{Lemma 1} (\ie
$\varphi(g)$ is the set of all $\varphi(g')$ as $g'$ ranges over the
successors of $g$). Put $A=\varphi(g_0)$. Show that $A$ is an Artinian
set. For $g\in G$, let $g_0<g_1<\ldots<g_n=g$ be the sequence of
elements $\leq g$; put
$f(g)=(\varphi(g_0),\varphi(g_1),\ldots,\varphi(g_n))$; show that $f$
is an increasing surjection from $G$ onto the ordered set $G(A)$ of
the text. Show that, if $f$ is \emph{injective}, it is an
\emph{isomorphism} from $G$ onto $G(A)$.

\hypertarget{sga4.exp.I.appendix.II.sec.7.exercise.3.part.c}{}
\item For $g\in G$, let $M_g$ be the set of upper bounds of $g$.
Suppose that, for every pair of distinct elements $g$, $g'$ having
the same predecessor $p$, the ordered sets $M_g$ and $M_{g'}$ are
not isomorphic. Show that the map $f$ of \hyperlink{sga4.exp.I.appendix.II.sec.7.exercise.3.part.b}{b)} is then an
\emph{isomorphism} from $G$ onto $G(A)$. [If $f(g)=f(g')$ with
$g\neq g'$, and if $g_0<g_1<\ldots<g_n=g$ and
$g'_0<g'_1<\ldots <g'_{n'}=g'$ are the sequences of elements
$\leq g$ and $\leq g'$, respectively, show that $n=n'$, and that one
may assume that $g$ and $g'$ have the same predecessor $p$; then
consider the set of $p\in G$ for which there exist two distinct
successors $g$, $g'$ of $p$ such that $f(g)=f(g')$, a maximal element
$q$ of this set, and two distinct successors $h$, $h'$ of $q$ such
that $f(h)=f(h')$; note that the restrictions of $f$ to $M_h$ and
$M_{h'}$ are injective and hence (by \hyperlink{sga4.exp.I.appendix.II.sec.7.exercise.3.part.b}{b)}) are isomorphisms from $M_h$
onto $G(\varphi(h))$ and from $M_{h'}$ onto $G(\varphi(h'))$; deduce
from the hypothesis that $M_h$ and $M_{h'}$ are not isomorphic that
$\varphi(h)\neq \varphi(h')$, which contradicts $f(h)=f(h')$].
\end{enumerate}
N.B.\pageoriginale \hyperlink{sga4.exp.I.appendix.II.sec.7.exercise.3}{Exercise 3}) gives very precise information on how
Artinian sets are ``made.'' We already knew that such a set $A$ is
determined by the isomorphism class of the ordered set $G(A)$ (\hyperlink{sga4.exp.I.appendix.II.sec.6.lemma.1}{Lemma
1}). We now know how to \emph{characterize} the ordered sets $G$
isomorphic to some $G(A)$:
\begin{itemize}
\item[$\alpha)$] $G$ is Noetherian and has a least element;

\item[$\beta)$] For every $g\in G$, the set of $h\in G$ such that
$h\leq g$ is totally ordered and finite (whence the existence and
uniqueness of the predecessor of every nonleast $g$);

\item[$\gamma)$] If $g$, $g'$ are distinct elements having the same
predecessor, the set $M_g$ of upper bounds of $g$ and the set
$M_{g'}$ of upper bounds of $g'$ are not isomorphic.
\end{itemize}

\hypertarget{sga4.exp.I.appendix.II.sec.7.exercise.4}{}
\item Let $(A_n)_{n\geq 0}$ be the sequence of finite \emph{ordinals}
defined by $A_0=∅$, $A_{n+1}=A_n\cup \{A_n\}$. Show that the
ordered set $G(A_n)$ has $2^n$ elements, and that the number of its
elements of height $q$ (in the sense of \hyperlink{sga4.exp.I.appendix.II.sec.7.exercise.2}{Exercise 2}) is
$\binom{n}{q}$.
\end{enumerate}
\end{enonce*}

% END INCLUDED SGA4_English_current_buildable_source_20260906/source/01/drafts/SGA4_I_LINES_8781_8916_English.texfrag


% BEGIN INCLUDED SGA4_English_current_buildable_source_20260906/source/01/drafts/SGA4_I_LINES_8917_8927_English.texfrag
\makeatletter
\let\toclevel@section\sgaFourSavedToclevelSection
\makeatother
\pdfbookmark[2]{Bibliography}{sga4.exp.I.appendix.II.bibliography}
\begingroup
\makeatletter
\def\Hy@writebookmark#1#2#3#4#5{}
\makeatother
\begin{thebibliography}{10}%%
\bibitem[1]{sga4.exp.I.bib.I.2} \sisi{B. \textsc{Mitchell}, Theory of categories, Academic Press
(1965).}{B. \textsc{Mitchell}, \textit{Theory of categories},
Academic Press, 1965.}
\end{thebibliography}
\endgroup
\hypersetup{bookmarksdepth=1}

\ifx\danslelivre\undefined
%\bibliography{../fichiers-configuration-latex/bibliographie}
%\bibliographystyle{../fichiers-configuration-latex/smfalpha}
\end{document}
\fi

% END INCLUDED SGA4_English_current_buildable_source_20260906/source/01/drafts/SGA4_I_LINES_8917_8927_English.texfrag


% END INCLUDED SGA4_English_current_buildable_source_20260906/source/01/01_English.tex


\makeatletter
\def\input@path{{../02/}}
\makeatother
% BEGIN INCLUDED SGA4_English_current_buildable_source_20260906/source/02/02_English.tex
\ifx\danslelivre\undefined
% Release engine: XeLaTeX (the LuaTeX compatibility layer breaks Xy-pic output).
\documentclass[10pt]{../fichiers-configuration-latex/smfart}
\usepackage{ifthen}
% BEGIN INCLUDED SGA4_English_current_buildable_source_20260906/source/fichiers-configuration-latex/paquets.tex
\usepackage{polyglossia}
%\usepackage[francais]{babel}
\usepackage{fontspec}
\usepackage[math-style=ISO,bold-style=upright,sans-style=upright]{unicode-math}
%\usepackage[mathletters]{ucs}
%\usepackage[utf8]{inputenc}
%\usepackage[T1]{fontenc}
\usepackage[unicode,hypertexnames=false,colorlinks=true,linkcolor=blue,citecolor=Sepia,urlcolor=magenta]{hyperref}
\setmainlanguage{french}
\setmainfont[Ligatures=TeX]{Linux Libertine O}
% variante possible : TeXGyrePagella ?
\setmathfont{XITS Math}

% maths
\usepackage{amsmath}
%\usepackage{amssymb}
%\usepackage{mathrsfs}
%\usepackage{stmaryrd}
%\usepackage{euscript}

% couleurs
\usepackage[usenames,dvipsnames]{xcolor}

% graphiques
\usepackage{graphicx}
\usepackage[all]{xy}
%\usepackage{tikz}
%\usetikzlibrary{matrix,arrows,calc}
%\tikzset{column sep=3em,row sep=2em,text height=1.5ex, text depth=0.25ex,>=angle 90}
%\tikzset{description/.style={fill=white,inner sep=1pt}}

% hyperliens, 'etiquettes
\usepackage{xr-hyper}
\usepackage{../fichiers-configuration-latex/smfthm}
%\ifx\cacheretiquettes\undefined
%\usepackage[inline]{showlabels} \else \fi % pour rendre visibles les 'etiquettes
%\usepackage{../fichiers-configuration-latex/makeidx} % pour faire un index
%\usepackage[pagebackref,unicode]{hyperref}  % pour voir o`u sont cit'ees les r'ef'erences
% hyperref is loaded once above with the cumulative-reader destination policy.
%\hypersetup{breaklinks=true} %permet le retour `a la ligne dans les liens trop longs. Incompatible avec pdflatex
%\renewcommand*{\backref}[1]{{\small [cit\'e page(s)~#1.]}}

\usepackage{xspace}
\usepackage{pifont} % pour expose XVII

% END INCLUDED SGA4_English_current_buildable_source_20260906/source/fichiers-configuration-latex/paquets.tex

% BEGIN INCLUDED SGA4_English_current_buildable_source_20260906/source/fichiers-configuration-latex/numerotation.tex
\renewcommand{\thefootnote}{\roman{footnote}} % pour mettre les notes en bas de page sous la forme (ii) et non (2), qui ressemble trop (?) a un exposant.

\newcounter{pageoriginal}
\newcommand{\pageoriginale}{\refstepcounter{pageoriginal}\marginpar{\footnotesize\textbf{\thepageoriginal}}\xspace}
\let\pageoriginaled\pageoriginale

\usepackage{ifthen}
\newboolean{original}
\def\sisi{\ifthenelse{\boolean{original}}}
\setboolean{original}{false}

\newcounter{cnote}
\def\nde#1{\ifthenelse{\boolean{original}}{}{\stepcounter{cnote}\def\thefootnote{\arabic{cnote}}\protect\footnote[\value{cnote}]{N.D.E.: #1}}}
\def\ndemark{\ifthenelse{\boolean{original}}{}{\stepcounter{cnote}\def\thefootnote{\arabic{cnote}}\protect\footnotemark[\value{cnote}]}}
\def\ndetext#1{\ifthenelse{\boolean{original}}{}{\def\thefootnote{\arabic{cnote}}\protect\footnotetext[\value{cnote}]{N.D.E.: #1}}}

% END INCLUDED SGA4_English_current_buildable_source_20260906/source/fichiers-configuration-latex/numerotation.tex

% BEGIN INCLUDED SGA4_English_current_buildable_source_20260906/source/fichiers-configuration-latex/macros.tex
\setmainlanguage{french}
\setotherlanguage{english}


%% Fixed, engine-independent release identifier.  The inherited dynamic Git
%% probe used Lua io.popen, exposed build-state-dependent output, and printed
%% a visible question mark outside a Git checkout.
\def\showgitstatus{English release r2, 2026-07-23}


% macros pour LaTeX
\renewcommand{\mathsf}{}

% Police différente pour U+2113 SCRIPT SMALL L
\setmathfont[range="2113]{latinmodern-math.otf}

% Give every inherited semantic label its own stable PDF destination instead
% of inheriting whichever automatic equation/page anchor happened to be live.
% The legacy label remains as a compatibility alias used by existing edges.
\makeatletter
\def\etiquette[#1]#2{%
  \Hy@raisedlink{\hyper@anchorstart{sga4.semantic.#2}\hyper@anchorend}%
  \def\@currentHref{sga4.semantic.#2}%
  \def\@currentlabel{#1}%
  \label{#2}%
}
\makeatother
\def\eqhyperref[#1]#2{(\hyperref[#1]{#2})}
\def\pagehyperref[#1]#2{(page \hyperref[#1]{#2})}
\def\hyperrefext[#1][#2]#3{\hyperref[#2]{#1 #3}}
\def\eqhyperrefext[#1][#2]#3{\hyperref[#2]{#1 #3}}     %cf. exposes 17
\def\neweqhyperrefext[#1][#2]#3{\hyperref[#2]{#1 #3}}  %cf. exposes 17

\def\shat#1{\protect\widehat{#1}}
\def\stilde#1{\protect\widetilde{#1}}
%%%%%%%%%%%%%%%%%%


\def\etcn{etc.\xspace}
\def\etco{etc...\xspace}
\def\etc{\sisi{\etco}{\etcn}}
\def\ldotsn{\sisi{...}{\etcn}}

\def\cf{{cf.}\kern.3em}
\def\Cf{{Cf.}\kern.3em}
\def\ie{{i.e.}\kern.3em}
\def\iev{{i.e.},\xspace}
\def\loccito{loc.\kern2pt cit.\xspace}
\def\loccitn{\textit{loc.\kern2pt cit.}\xspace}
\def\loccit{\sisi{\loccito}{\loccitn}}
\def\EGA{\textup{EGA}\kern.2em}
\def\SGA{\textup{SGA}\kern.3em}

\def\quoi{}

\def\quav{\quad,\quad}
\def\quaet{\quad \text{et}\quad}
\def\fibreproduct#1#2#3{#1 ×_{#3} #2}
\let\fprod\fibreproduct
\def\ccheck#1#2{\mbox{${\check{\check{\rm #1}}}^{\lower2pt\hbox{$\scriptstyle #2$}}$}}
\def\sheaf#1{\protect\mathcal{#1}}
\def\vieuxsheaf#1{{#1}} % cf. problème avec sheaf{a\Phi_i} par exemple...
\DeclareMathOperator{\Hom}{Hom}
\DeclareMathOperator{\can}{can}
\DeclareMathOperator{\Ind}{Ind}
\DeclareMathOperator{\ind}{ind}
%\DeclareMathOperator{\Hhomtop}{\mathcal{H}\!\mathit{omtop}}
\DeclareMathOperator{\Hhomtop}{\mathbf{Homtop}}
\DeclareMathOperator{\Point}{Point}
\DeclareMathOperator{\Pro}{Pro}
\DeclareMathOperator{\Top}{Top}
\DeclareMathOperator{\Mod}{Mod}
\DeclareMathOperator{\Cart}{Cart}
%\DeclareMathOperator{\Mmorsite}{\mathcal{M}\!\mathit{orsite}}
\DeclareMathOperator{\Mmorsite}{\mathbf{Morsite}}
\DeclareMathOperator{\loc}{loc}
\DeclareMathOperator{\Loc}{Loc}
\DeclareMathOperator{\Pic}{Pic}
\DeclareMathOperator{\Et}{Et}
\DeclareMathOperator{\card}{card}
\DeclareMathOperator{\Ens}{Ens}
\DeclareMathOperator{\h}{\underline{H}}
\DeclareMathOperator{\Fonct}{Fonct}
\DeclareMathOperator{\Ker}{Ker}
\DeclareMathOperator{\but}{but}
\DeclareMathOperator{\Ob}{Ob}
\DeclareMathOperator{\ob}{ob}
\DeclareMathOperator{\Lib}{Lib}
\DeclareMathOperator{\esj}{esj}
\DeclareMathOperator{\id}{id}
\DeclareMathOperator{\Id}{Id}
\DeclareMathOperator{\ad}{ad}
\DeclareMathOperator{\Iim}{Im}
\DeclareMathOperator{\pr}{pr}
\DeclareMathOperator{\Cat}{Cat}
\DeclareMathOperator{\car}{car}
\DeclareMathOperator{\SysInd}{SysInd}
\DeclareMathOperator{\Sup}{Sup}
\DeclareMathOperator{\Coker}{Coker}
\DeclareMathOperator{\Filt}{Filt}
\DeclareMathOperator{\acc}{acc}
\DeclareMathOperator{\sur}{sur}
\DeclareMathOperator{\cat}{cat}
\DeclareMathOperator{\Ord}{Ord}
\DeclareMathOperator{\Gr}{Gr}
\DeclareMathOperator{\Cov}{Cov}
\DeclareMathOperator{\Caf}{Caf}
\DeclareMathOperator{\coker}{coker}
\DeclareMathOperator{\coim}{coim}
\DeclareMathOperator{\im}{im}
\DeclareMathOperator{\Ouv}{Ouv}
\DeclareMathOperator{\diag}{diag}
\DeclareMathOperator{\Esp}{Esp}
\DeclareMathOperator{\TOP}{TOP}
\DeclareMathOperator{\Unif}{Unif}
\DeclareMathOperator{\Iind}{\mathbf{Ind}}
\DeclareMathOperator{\Hhomcart}{\mathbf{Homcart}}
\DeclareMathOperator{\Hhomtopess}{\mathbf{Homtopess}}
\DeclareMathOperator{\Ppoints}{\mathbf{Points}}
\DeclareMathOperator{\Ppoint}{\mathbf{Point}}
\DeclareMathOperator{\Ttors}{\mathbf{Tors}}
\DeclareMathOperator{\Ssex}{\mathbf{Sex}}
\DeclareMathOperator{\Ffib}{\mathbf{Fib}}
\DeclareMathOperator{\Vvois}{\mathbf{Vois}}
\DeclareMathOperator{\Ppointess}{\mathbf{Pointess}}
\DeclareMathOperator{\scrHom}{\mathbf{Hom}}
\DeclareMathOperator{\sheafPic}{\mathbf{Pic}}
\DeclareMathOperator{\sheafAut}{\mathbf{Aut}}
\DeclareMathOperator{\Hhomtoplocan}{\mathbf{Homtoplocan}}
\DeclareMathOperator{\Hhomtopan}{\mathbf{Homtopan}}
\DeclareMathOperator{\Pptgeom}{\mathbf{Ptgéom}}
\DeclareMathOperator{\Morsite}{Morsite}
\DeclareMathOperator{\Fib}{Fib}
\DeclareMathOperator{\Fais}{Fais}
\DeclareMathOperator{\kar}{kar}
\DeclareMathOperator{\Groupes}{Groupes}
\DeclareMathOperator{\Restr}{Restr}
\DeclareMathOperator{\Isom}{Isom}
\DeclareMathOperator{\Prol}{Prol}
\DeclareMathOperator{\Res}{Res}
\DeclareMathOperator{\Tors}{Tors}
\DeclareMathOperator{\esp}{esp}
\DeclareMathOperator{\es}{es}
\DeclareMathOperator{\od}{od}
\DeclareMathOperator{\gr}{gr}
\DeclareMathOperator{\sob}{sob}
\DeclareMathOperator{\dfn}{dfn}
\DeclareMathOperator{\Sch}{Sch}
\DeclareMathOperator{\ouv}{ouv}
\DeclareMathOperator{\cons}{cons}
\DeclareMathOperator{\Nerf}{Nerf}
\DeclareMathOperator{\Points}{Points}
\DeclareMathOperator{\Espsob}{Espsob}
\DeclareMathOperator{\Homtop}{Homtop}
\DeclareMathOperator{\Pt}{Pt}
\DeclareMathOperator{\Quot}{Quot}
\DeclareMathOperator{\glob}{glob}
\DeclareMathOperator{\Homfib}{Homfib}
\DeclareMathOperator{\Qouv}{Qouv}
\DeclareMathOperator{\Quotouv}{Quotouv}
\DeclareMathOperator{\ens}{ens}
\DeclareMathOperator{\ab}{ab}
\DeclareMathOperator{\an}{an}
\DeclareMathOperator{\fs}{fs}
\DeclareMathOperator{\red}{red}
\DeclareMathOperator{\Ab}{Ab}
\DeclareMathOperator{\cosq}{cosq}
\DeclareMathOperator{\Ann}{Ann}
\DeclareMathOperator{\Zar}{Zar}


%%%%%%%%%%%%
%%%%% SGA4-2
\DeclareMathOperator{\Sex}{Sex}
\DeclareMathOperator{\vois}{vois}
\DeclareMathOperator{\Vois}{Vois}
\DeclareMathOperator{\Ext}{Ext}
\DeclareMathOperator{\ext}{ext}
\DeclareMathOperator{\Fer}{Fer}
\DeclareMathOperator{\Add}{Add}
\DeclareMathOperator{\Cocart}{Cocart}
\DeclareMathOperator{\cosk}{cosk}
\DeclareMathOperator{\plat}{plat}

\DeclareMathOperator{\faisc}{faisc}
\DeclareMathOperator{\sk}{sk}
\DeclareMathOperator{\ch}{ch}

\def\HH{\mathbf{H}}
\def\qc{\mathrm{qc}}
\def\et{\mathrm{et}}
\def\coh{\mathrm{coh}}

\def\Hvr#1{%
\begin{array}[b]{l}
{\kern2pt\scriptstyle r}\\[-8pt]
{\kern1pt\scriptstyle \vee}\\[-5pt]
H^{#1}
\end{array}}

\def\Hvv#1{%
\begin{array}[b]{l}
\kern2pt\rotatebox{90}{${\scriptscriptstyle\ll}$}\\[-6pt]
H^{#1}
\end{array}}

\def\Hvi#1{%
\begin{array}[b]{l}
{\kern0.2pt\scriptscriptstyle \infty}\\[-8pt]
{\kern1pt\scriptstyle \vee}\\[-5pt]
H^{#1}
\end{array}}

\def\Hv#1{%
\begin{array}[b]{l}
{\kern1pt\scriptstyle \vee}\\[-5pt]
H^{#1}
\end{array}}

\DeclareMathOperator{\cart}{cart}
\DeclareMathOperator{\Ffaisc}{\mathbf{Faisc}}


\def\Ev#1{%
\begin{array}[b]{l}
{\kern1pt\scriptstyle \vee}\\[-5pt]
E'_{#1}
\end{array}}

\DeclareMathOperator{\Cont}{Cont}
\DeclareMathOperator{\Ccont}{\mathscr{C}\!\mathit{ont}}
\DeclareMathOperator{\Site}{Site}
\newcommand\MM{\mathbf{M}}
\def\limtop#1{{\displaystyle{\varprojlim_{#1}\!\mbox{\rm top }}}}
\def\Limtop{\mathop{\varprojlim\!\mbox{\rm top}}}
\DeclareMathOperator{\spec}{spec}
\def\ss{\scriptstyle{\mathscr{S}}}
\DeclareMathOperator{\Gal}{Gal}
\def\smallm{\scriptstyle{\mathscr{M}}}
\DeclareMathOperator{\End}{End}
\DeclareMathOperator{\ass}{ass}
\DeclareMathOperator{\zar}{zar}
\DeclareMathOperator{\Ogg}{Ogg}
\DeclareMathOperator{\res}{res}
\DeclareMathOperator{\tr}{tr}
\DeclareMathOperator{\Supp}{Supp}
\DeclareMathOperator{\codim}{codim}
\DeclareMathOperator{\rk}{rk}
\DeclareMathOperator{\cur}{cur}
\DeclareMathOperator{\cdqc}{cdqc}
\newcommand\DD{\mathbf{D}}
\DeclareMathOperator{\cl}{cl}
\def\JJ{\mathbf{J}}
\def\FF{\mathbf{F}}
\def\BB{\mathbf{B}}
\DeclareMathOperator{\Alget}{Alget}
\DeclareMathOperator{\Ass}{Ass}
\DeclareMathOperator{\supp}{supp}
\DeclareMathOperator{\rad}{rad}
\DeclareMathOperator{\prof}{prof}
\DeclareMathOperator{\cint}{int}
\DeclareMathOperator{\sic}{sic}
\DeclareMathOperator{\Tor}{Tor}
\DeclareMathOperator{\tors}{tors}
\DeclareMathOperator{\Sym}{Sym}
\DeclareMathOperator{\parf}{parf}
\DeclareMathOperator{\tor}{tor}
\DeclareMathOperator{\defn}{defn}
\DeclareMathOperator{\Div}{Div}
\DeclareMathOperator{\Pol}{Pol}
\DeclareMathOperator{\Fex}{Fex}
\DeclareMathOperator{\Lie}{Lie}
\DeclareMathOperator{\tot}{tot}
\DeclareMathOperator{\Aut}{Aut}
\DeclareMathOperator{\HOM}{HOM}
\DeclareMathOperator{\pch}{pch}
\DeclareMathOperator{\Ch}{Ch}
\DeclareMathOperator{\EXT}{EXT}
\DeclareMathOperator{\PIC}{PIC}
\DeclareMathOperator{\Rf}{Rf}
\DeclareMathOperator{\inj}{inj}
\newcommand\VV{\mathbf{V}}

\def\Xv{%
\begin{array}[b]{l}
{\kern1pt\scriptstyle \vee}\\[-5pt]
X
\end{array}}

%%%%%%%
\newcommand\TT{\mathbf{T}}
\newcommand\XX{\mathbf{X}}
\newcommand\EE{\mathbf{E}}
%%%%%%

\def\qedsymbol{\textup{C.Q.F.D.}}
\let\CQFD\qed
\let\cqfd\qed

%%%%%\bigdot
\font\bigsymb=cmsy10 at 4pt
\def\bigdot{{\kern1.2pt\raise 1.5pt\hbox{\bigsymb\char15}}}
%%%%%

\def\Lright{\mathop{\vcenter{\hbox{\hbox{$\longrightarrow$}\lower8pt \hbox{\kern-16pt\hbox{$\longrightarrow$}}}}}}

\def\Aa{{\mathcal A}}
\def\Bb{{\mathcal B}}
\def\Cc{{\mathcal C}}
\def\Dd{{\mathcal D}}
\def\Ee{{\mathcal E}}
\def\Ff{{\mathcal F}}
\def\Gg{{\mathcal G}}
\def\Hh{{\mathcal H}}
\def\Ii{{\mathcal I}}
\def\Jj{{\mathcal J}}
\def\Kk{{\mathcal K}}
\def\Ll{{\mathcal L}}
\def\Mm{{\mathcal M}}
\def\Nn{{\mathcal N}}
\def\Oo{{\mathcal O}}
\def\Pp{{\mathcal P}}
\def\Qq{{\mathcal Q}}
\def\Rr{{\mathcal R}}
\def\Ss{{\mathcal S}}
\def\Tt{{\mathcal T}}
\def\Uu{{\mathcal U}}
\def\Vv{{\mathcal V}}
\def\Ww{{\mathcal W}}
\def\Xx{{\mathcal X}}
\def\Yy{{\mathcal Y}}
\def\Zz{{\mathcal Z}}

\newcommand{\CC}{\mathbf{C}}
\newcommand{\GG}{\mathbf{G}}
\newcommand{\KK}{\mathbf{K}}
\newcommand{\LL}{\mathbf{L}}
\newcommand{\NN}{\mathbf{N}}
\newcommand{\PP}{\mathbf{P}}
\newcommand\QQ{\mathbf{Q}}
\newcommand\RR{\mathbf{R}}
\newcommand\ZZ{\mathbf{Z}}

\def\to{\mathchoice{\longrightarrow}{\rightarrow}{\rightarrow}{\rightarrow}}
\def\from{\mathchoice{\longleftarrow}{\leftarrow}{\leftarrow}{\leftarrow}}
\def\mto{\mathchoice{\longmapsto}{\mapsto}{\mapsto}{\mapsto}}
\def\mfrom{\mathrel{\from\joinrel\mapstochar\,}}
\def\hto{\mathrel{\lhook\joinrel\to}}
\def\hfrom{\mathrel{\from\joinrel\rhook}}

\def\implies{\mathchoice{\Longrightarrow}{\Rightarrow}{\Rightarrow}{\Rightarrow}}
\def\implied{\mathchoice{\Longleftarrow}{\Leftarrow}{\Leftarrow}{\Leftarrow}}
\def\iff{\mathchoice{\Longleftrightarrow}{\Leftrightarrow}{\Leftrightarrow}{\Leftrightarrow}}

\def\To#1{\mathchoice{\xrightarrow{\labelstyle~~#1~}}{\stackrel{#1}{\longrightarrow}}{}{}}
\def\Mto#1{\mathrel{\mapstochar\To{#1}}}
\def\From#1{\mathchoice{\xleftarrow{\labelstyle~~#1~}}{\stackrel{#1}{\longleftarrow}}{}{}}

\newcommand{\isomto}{\overset{\sim}{\longrightarrow}}
\newcommand{\isomfrom}{\overset{\sim}{\longleftarrow}}

\newcommand{\IFF}{\unskip~$\iff$~\ignorespaces}
\newcommand{\THEN}{\unskip~$\implies$~\ignorespaces}

\let\leq\leqslant
\let\geq\geqslant
\let\le\leqslant
\let\ge\geqslant

\def\cf{{cf.}\kern.3em}
\def\Cf{{Cf.}\kern.3em}
\def\ie{{i.e.}\kern.3em}
\def\iev{{i.e.},\xspace}
%\def\loccito{loc.\kern2pt cit.\xspace}
%\def\loccitn{\textit{loc.\kern2pt cit.}\xspace}
%\def\loccit{\sisi{\loccito}{\loccitn}}
\def\EGA{\textup{EGA}\kern.2em}
\def\SGA{\textup{SGA}\kern.3em}

\let\oldS\S
\def\S{\oldS\kern .15em}
\def\ptbl{.\kern .2em}
\def\resp{resp.\kern.3em}
\def\Exp{Exp.\kern.3em}

\let\bar\overline
\let\hat\widehat
\let\tilde\widetilde
\let\dpl\displaystyle
\let\ldots\dots
\let\emptyset\varnothing
\let\skpt\skippointrait
\let\epsilon\varepsilon

\def\sfrac#1#2{{#1}/{#2}}

\DeclareMathOperator{\cd}{cd}
%\DeclareMathOperator{\Hhom}{\mathcal{H}\!\mathit{om}}
\DeclareMathOperator{\Hhom}{\mathbf{Hom}}
%\DeclareMathOperator{\Eext}{\mathcal{E}\!\mathit{xt}}
\DeclareMathOperator{\Eext}{\mathbf{Ext}}
\DeclareMathOperator{\Spec}{Spec}
\DeclareMathOperator{\Tr}{Tr}

\let\emptyset=\varnothing

%%%%%%%%%%%%%%%
%%% expose XVII
%%%%%%%%%%%%%%%

\makeatletter
\let\old@optexpose\optexpose
\newcommand{\empty@optexpose}[1]{}
\newcommand{\paren@optexpose}[1]{\old@optexpose{#1}(}
\newcommand{\refnoexpose}[1]{\let\optexpose\empty@optexpose\ref{#1}\let\optexpose\old@optexpose}%\refnoexpose{X.cor5.2} engendre "5.2" au lieu de "X 5.2"
\newcommand{\neweqref}[1]{\let\optexpose\paren@optexpose\textup{\ref{#1})}\let\optexpose\old@optexpose}%\neweqref{XII.eq4.1} engendre "XII (4.1)" au lieu de "(XII 4.1)"
\makeatother
\newcommand{\refvb}[2][VI${}_{\text{B}}$ ]{\sisi{#1\refnoexpose{#2}}{\ref{#2}}}

\newcommand{\leftexp}[2]{{\vphantom{#2}}^{#1}{#2}}
\newcommand{\AAA}{\mathbf{A}}
\newcommand{\Sf}{\mathfrak{S}}
\DeclareMathOperator*{\quotprojlim}{\text{``$\varprojlim$''}}
\DeclareMathOperator*{\quotinjlim}{\text{``$\varinjlim$''}}
\DeclareMathOperator{\Fl}{F\sisi{\ell}{l}}
\DeclareMathOperator{\Pl}{P\sisi{\ell}{l}}
%\DeclareMathOperator{\cl}{c\sisi{\ell}{l}}
\DeclareMathOperator{\tdf}{tdf}
\DeclareMathOperator{\rg}{rg}
\newcommand{\iInd}{\sisi{I}{i}nd}
\newcommand{\iind}{ind}
\newcommand{\ppro}{pro}
\newcommand{\tTor}{\sisi{T}{t}or}
\newcommand{\ttor}{tor}
\newcommand{\chap}{Chap.}
\newcommand{\Ffex}{\Ff\mathit{ex}}
\newcommand{\sheafk}{\underline{k}}
\newcommand{\sheafn}{\underline{n}}
\newcommand{\sheafGamma}{\underline{\Gamma}}

% Repeated references to an already typeset footnote must not use
% \footnotemark without a matching \footnotetext: Hyperref would then create
% a live annotation to a destination that cannot exist.  This helper keeps the
% current printed footnote mark and links it to an explicit stable note target.
\makeatletter
\newcommand{\sgarepeatedfootnotemark}[1]{%
  \begingroup
  \edef\@thefnmark{\thefootnote}%
  \hyperlink{#1}{\@makefnmark}%
  \endgroup}
\makeatother

\sisi{
\DeclareMathOperator{\simarrow}{\rule{.15mm}{1.8mm}\kern -3.4pt\rightsquigarrow}
\newcommand{\sheafell}{\underline{\ell}}
\newcommand{\sheafZZ}{\underline{\ZZ}}
}{
\newcommand{\nderef}[1]{la N.D.E.~(\ref{#1}) page~\pageref{#1}}
}

%%%%%%%%%%%%%%%%%
%%%%% expose XVII
%%%%%%%%%%%%%%%%%

% END INCLUDED SGA4_English_current_buildable_source_20260906/source/fichiers-configuration-latex/macros.tex


\externaldocument{../01/01_English.tex}
%\externaldocument{../02/02_English.tex}
\externaldocument{../03/03.tex}
\externaldocument{../04/04.tex}

\begin{document}

\begin{center}
II. Topologies and sheaves\\
J.-L. Verdier
\end{center}

\selectlanguage{english}version: \showgitstatus
\tableofcontents

\else
\chapter{Topologies and sheaves}\MakeLinkTarget*{sga4.exp.II}\label{sga4.exp.II}
%\addtocontents{toc}{by J.-L. Verdier}
\begin{center}
J.-L. Verdier
\end{center}
\fi

% BEGIN INCLUDED SGA4_English_current_buildable_source_20260906/source/02/drafts/SGA4_II_LINES_31_212_English.texfrag
After\pageoriginale defining topologies and
pretopologies (\hyperref[II.sec1]{Section 1}) and sheaves of sets
(\hyperref[II.sec2]{Section 2}), in \hyperref[II.sec3]{Section 3} we turn to the
central theorem of this exposé (\hyperref[II.theo3.4]{3.4}): the existence
of the sheaf associated with a presheaf. In order to cover all the
cases encountered in practice, the existence of this functor is
proved for presheaves on a $𝒰$-site
(\hyperref[II.def3.0.2]{3.0.2}). The consequences of this theorem for
the behavior of inductive and projective limits in categories of
sheaves are drawn in \hyperref[II.sec4]{Sections 4} and 5. In \hyperref[II.sec6]{Section 6} we define and study
sheaves with values in arbitrary categories, immediately focusing
attention on sheaves of rings, groups,
modules, \etcn

\section{Topologies, covering families,
  pretopologies}\etiquette[1]{II.sec1}

\begin{enonce*}{Definition 1.1}\etiquette[1.1]{II.def1.1}
A topology on a category $C$ consists, for every
object $X$ of $C$, of a set $J(X)$ of sieves on $X$, subject to the
following axioms:
\begin{itemize}
\item[T 1)] {\rm{(Stability under base change).}} For every
  object $X$ of $C$, every sieve $R\in J(X)$, and every morphism
  $f:Y\rightarrow X$ $(Y\in \ob (C))$, the sieve $R\times_X Y$ on $Y$
  belongs to $J(Y)$.

\item[T 2)] {\rm(Local character).} If $R$ and $R'$ are two
  sieves on $X$, if $R\in J(X)$, and if, for every $Y\in \ob (C)$ and
  every morphism $Y\rightarrow R$, the sieve $R'\times_X Y$
  belongs\pageoriginale
  to $J(Y)$, then $R'$ belongs to $J(X)$.

\item[T 3)] For every object $X$ of $C$, $X$ belongs to $J(X)$.
\end{itemize}
\end{enonce*}

\setcounter{subsection}{1}
\subsubsection{}\etiquette[1.1.1]{II.subsubsec1.1.1}

The sieves belonging to $J(X)$ will be called the \emph{sieves
  covering $X$}, or the \emph{refinements of $X$}. It follows
  immediately from axioms (T1), (T2), and (T3) that the set of
  sieves covering $X$ is stable under finite intersections, and that
  every sieve containing a covering sieve is itself a covering
  sieve. Thus the set $J(X)$, ordered by inclusion, is
  cofiltered (\hyperrefext[I][I.sec8]{8}).

\subsubsection{}\etiquette[1.1.2]{II.subsubsec1.1.2}

Let $C$ be a category, and let $T$ and $T'$ be two topologies on $C$. The
topology $T$ is said to be \emph{finer than the topology} $T'$ if,
for every object $X$ of $C$, every refinement of $X$ for the
topology $T'$ is a refinement of $X$ for the topology $T$. This
defines an order structure on the set
of topologies.

\subsubsection{}\etiquette[1.1.3]{II.subsubsec1.1.3}

Let $(T_i)_{i\in I}$ be a family of topologies on $C$. For every
object $X$ of $C$, the set of sieves on $X$ that are
covering for all the topologies $T_i$ satisfies axioms
(T 1), (T 2), and (T 3), and hence defines a topology: the
\emph{intersection topology} of the $T_i$, \ie the infimum
of the $T_i$. It is the finest topology that is coarser than
all the topologies $T_i$. Consequently, the family $(T_i)_{i \in I}$
also admits a \emph{supremum}: the intersection topology of the
topologies finer than each of the~$T_i$.

\subsubsection{}\etiquette[1.1.4]{II.subsubsec1.1.4}
In\pageoriginale particular, taking $J(X)$ to be the set of all
sieves on $X$ gives a topology finer than every topology
on $C$; we call it the \emph{discrete topology}.

There is a topology coarser than every topology on $C$: the
topology for which $J(X)=\{X\}$ for every $X$ of $C$. This
topology is called the \emph{coarse} or
\emph{chaotic} topology.

\subsubsection{}\etiquette[1.1.5]{II.subsubsec1.1.5}
A category $C$ equipped with a topology is called a \emph{site.}
The category $C$ is called the \emph{category
  underlying the site.}

\begin{enonce*}{Definition 1.2}\etiquette[1.2]{II.def1.2}
Let $C$ be a site and $X$ an object of $C$. A family of morphisms
$(f_{\alpha}:X_{\alpha}\rightarrow X)$, $\alpha\in A$, is called a
covering family if the sieve generated by the family $f_{\alpha}$
{\rm(I~\hyperref[I.sec4.3.3]{4.3.3})} is a sieve covering $X$.
\end{enonce*}

\subsubsection{}\etiquette[1.1.6]{II.subsubsec1.1.6}
Let $C$ be a category. Suppose that, for each object $X$ of $C$,
we are given a set of families of morphisms with target $X$. There is
then a coarsest topology $T$ for which
the given families are covering, namely the
intersection (\hyperref[II.subsubsec1.1.3]{1.1.3}) of all the topologies in
question. This topology is called the \emph{topology generated
by the given sets of families
  of morphisms}. In general, it is difficult to describe all the
covering sieves of this topology.
The situation is more favorable, however, in the following case:

\begin{enonce*}{Definition 1.3}\etiquette[1.3]{II.def1.3}
Let $C$ be a category. A pretopology on $C$ consists, for
each object $X$ of $C$, of a set $\Cov(X)$ of families of
morphisms with target $X$, subject to the following
axioms:
\begin{itemize}
\item[PT 0)] For\pageoriginale every object $X$ of $C$, the morphisms
  in the families belonging to $\Cov(X)$ are pullbackable. (Recall
  that a morphism $Y\rightarrow X$ is called pullbackable if, for every
  morphism $Z\rightarrow X$, the fiber product $Z\times_X Y$ is
  representable.)

\item[PT 1)] For every object $X$ of $C$, every family
  $(X_{\alpha}\rightarrow X)_{\alpha \in A}$ belonging to
  $\Cov(X)$, and every morphism $Y\rightarrow X$ $(Y\in \ob (C))$, the
  family $(X_{\alpha}\times_X Y\rightarrow Y)_{\alpha \in A}$
  belongs to $\Cov(Y)$. (Stability under base change.)

\item[PT 2)] If $(X_{\alpha}\rightarrow X)_{\alpha \in A}$ belongs
  to $\Cov(X)$ and if, for every $\alpha\in A$,
  $(X_{\beta_{\alpha}}\rightarrow X_{\alpha})_{\beta_{\alpha}\in
  B_{\alpha}}$ belongs to $\Cov(X_{\alpha})$, then the family
  $(X_\gamma\rightarrow X)_{\gamma\in \amalg_{\alpha\in A}B_{\alpha}}$ (where,
  for every $\gamma=(\alpha,\beta_{\alpha})$, $\alpha\in A$,
  $\beta_{\alpha}\in B_{\alpha}$, the morphism $X_\gamma\rightarrow X$ is
  the composite morphism $X_\gamma=X_{\beta_{\alpha}}\rightarrow
  X_{\alpha}\rightarrow X$) belongs to $\Cov(X)$. (Stability under
  composition.)

\item[PT 3)] The family $X\xrightarrow{\id_X} X$ belongs to $\Cov(X)$.
\end{itemize}
\end{enonce*}

\setcounter{subsection}{3} \setcounter{subsubsection}{0}
\subsubsection{}\etiquette[1.3.1]{II.subsubsec1.3.1}

The definition in \hyperref[II.subsubsec1.1.6]{1.1.6} allows us to consider,
for a given \emph{pretopology} on $C$, the \emph{topology} on
$C$ generated by that pretopology. Note that if $C$ is a
category in which fiber products are representable, then every
topology $T$ on $C$ can be defined by a pretopology, namely
the one for which $\Cov(X)$ consists of \emph{all} the
families covering $X$ for the topology $T$.

\begin{enonce*}{Proposition 1.4}\etiquette[1.4]{II.prop1.4}
Let $C$ be a category, $E$ a pretopology on $C$, $T$ the
topology defined by the pretopology $E$
$(\hyperref[II.subsubsec1.1.6]{1.1.6})$, and $X$ an object of $C$. Denote by
$J_E(X)$ the set of sieves on $X$ generated by families
of the pretopology, and by $J_T(X)$ the set of refinements of
$X$ for the\pageoriginale topology $T$. Then $J_E(X)$ is cofinal
in $J_T(X)$. In other words, a sieve $R$ on $X$
belongs to $J_T(X)$ if and only if there exists a sieve
$R'\in J_E(X)$ such that $R'\subset R$.
\end{enonce*}

\begin{enonce*}[remark]{Proof}
For every object $X$ of $C$, let $J'(X)$ be the set of sieves on
$X$ that contain a sieve belonging to $J_E(X)$. Clearly,
$J'(X)\subset J_T(X)$. To prove that $J'(X)=J_T(X)$, it is enough to
show that the data $J'(X)$ define a topology on $C$. Now
the $J'(X)$ clearly satisfy axioms (T 1) and (T 3). It
therefore remains to verify (T 2). For this, it is enough to prove that if
$R'$ is a subsieve of $R\in J_E(X)$ such that, for every
$Y\rightarrow R$, the sieve $R'\times_X Y$ on $Y$ belongs to
$J'(Y)$, then the sieve $R'$ belongs to $J'(X)$. Now the sieve $R$ is
generated by a family $(X_{\alpha}\rightarrow X)$ belonging to
$\Cov(X)$, and for every $\alpha$ the sieve $R'\times_X X_{\alpha}$
contains a sieve generated by a family
$(X_{\beta_{\alpha}}\rightarrow X_{\alpha})$ belonging to
$\Cov(X_{\alpha})$. It follows that the sieve $R'$ contains a
sieve generated by the family $(X_{\beta_{\alpha}}\rightarrow X)$.
Thus, by axiom (PT 2), $R'$ contains a sieve belonging to $J_E(X)$
and consequently belongs to $J'(X)$, \cqfd
\end{enonce*}

% END INCLUDED SGA4_English_current_buildable_source_20260906/source/02/drafts/SGA4_II_LINES_31_212_English.texfrag


% BEGIN INCLUDED SGA4_English_current_buildable_source_20260906/source/02/drafts/SGA4_II_LINES_213_421_English.texfrag
\section{Sheaves of sets}\etiquette[2]{II.sec2}

\begin{enonce*}{Definition 2.1}\etiquette[2.1]{II.def2.1}
Let $C$ be a site whose underlying category is a
$𝒰$-category. A presheaf $F$ with values in
$𝒰-\Ens$ is said to be separated (\resp to be a sheaf) if, for
every sieve $R$ covering an object $X$ of $C$, the map
$$
\Hom_{C^{\widehat{\phantom{,}}}}(X,F)\longrightarrow
\Hom_{C^{\widehat{\phantom{,}}}}(R,F)
$$
is injective (\resp bijective). The full subcategory
of $C^{\widehat{\phantom{,}}}$ whose\pageoriginale objects are
the sheaves is called the category of sheaves
of sets on $C$, and is most often denoted by $\widehat{C}$
\footnote[1]{or
  $C^{\widetilde{\phantom{,}}}$, depending on the mood of the
  typewriter.}. When no ambiguity can result, we shall
simply say the category of sheaves on $C$; on the other hand,
we shall sometimes be more precise and say the category of sheaves with
values in $𝒰-\Ens$, or the category of
$𝒰$-sheaves.
\end{enonce*}

\begin{enonce*}{Proposition 2.2}\etiquette[2.2]{II.prop2.2}
Let $C$ be a $𝒰$-category, and let $\Ff=(F_i)_{i\in I}$ be a
family of $𝒰$-presheaves on $C$. For
each object $X$ of $C$, denote by $J_{\Ff}(X)$ the set of sieves
$R\rightarrow X$ such that, for every morphism $Y\rightarrow X$ in
$C$ with target $X$, the sieve $R\times_X Y$ has the following
property: the map
$$
\Hom_{C^{\widehat{\phantom{,}}}}(Y,F_i)\longrightarrow
\Hom_{C^{\widehat{\phantom{,}}}} (R\times_X Y, F_i)
$$
is bijective (\resp injective) for every $i\in I$. Then the
sets $J_{\Ff}(X)$ define a topology on $C$, which is
the finest topology for which each $F_i$ is a
sheaf (\resp separated presheaf).
\end{enonce*}

\begin{enonce*}[remark]{Proof}
The $J_{\Ff}(X)$ clearly satisfy axioms (T 1) and (T
3). It remains to show that they satisfy (T 2). For this it suffices
to show the following:
\begin{enumerate}
\renewcommand{\labelenumi}{\theenumi)}
\item If $R'\hookrightarrow R\hookrightarrow X$ are two sieves on
  $X$ such that $R\in J_{\Ff}(X)$ and such that, for every
  $Y\rightarrow R$ (where $Y$ is an object of $C$), the sieve
  $R'\times_X Y$ belongs to $J_{\Ff}(Y)$, then the sieve
  $R'$ belongs to $J_{\Ff}(X)$.

\item If $R'\hookrightarrow R\hookrightarrow X$ are two sieves on
  $X$ such that $R'\in J_{\Ff}(X)$, then the sieve $R$
  belongs to $J_{\Ff}(X)$.
\end{enumerate}
\end{enonce*}

Note that, in case 2), for every $Y\rightarrow R$ (where $Y$ is
an object of $C$), the sieve $R'\times_X Y$ belongs to $J_{\Ff}(Y)$.
Now, in cases 1) and 2), $R$ is the inductive limit of the objects $Y$ of
$C$ over $R$ (I \hyperref[I.prop3.4]{3.4}). The\pageoriginale inductive limits
in $C^{\widehat{\phantom{,}}}$ are universal
(\hyperrefext[I][I.cor3.3]{3.3}). Consequently, in cases 1) and 2), $R'$ is the inductive limit
of the $R'\times_X Y$. Now, in cases 1) and 2),
$R'\times_X Y$ belongs to $J_{\Ff}(Y)$. Passing to the
inductive limit over the objects $Y$ of $C$ over $R$, we deduce that
the map
$$
\Hom_{C^{\widehat{\phantom{,}}}}(R,F_i)\longrightarrow
\Hom_{C^{\widehat{\phantom{,}}}}(R',F_i)
$$
is a bijection (\resp an injection) for every $i\in I$. It
follows that the maps
$$
\Hom_{C^{\widehat{\phantom{,}}}}(X,F_i)\longrightarrow
\Hom_{C^{\widehat{\phantom{,}}}}(R',F_i)\quav
\Hom_{C^{\widehat{\phantom{,}}}}(X,F_i)\longrightarrow
\Hom_{C^{\widehat{\phantom{,}}}}(R,F_i)
$$
are, in cases 1) and 2), bijections (\resp injections).
Moreover, hypotheses 1) and 2) are plainly stable under
arbitrary base change $Y\rightarrow X$ ($Y$ an object of $C$).
It then follows that, for every object $Y$ of $C$ over $X$
and every $i\in I$, the maps
\begin{align*}
\Hom_{C^{\widehat{\phantom{,}}}}(Y,F_i)&\longrightarrow
\Hom_{C^{\widehat{\phantom{,}}}}(R\times_X Y, F_i)\\
\Hom_{C^{\widehat{\phantom{,}}}}(Y,F_i)&\longrightarrow
\Hom_{C^{\widehat{\phantom{,}}}}(R'\times_X Y, F_i)
\end{align*}
are in both cases bijections (\resp injections); this
shows that $R'$ and $R$ belong to $J_{\Ff}(X)$, \cqfd

\begin{enonce*}{Corollary 2.3}\etiquette[2.3]{II.cor2.3}
Let $C$ be a $𝒰$-category, and for every $X$ of $C$ let
$K(X)$ be a set of sieves on $X$. Suppose that the $K(X)$
satisfy axiom {\rm (T 1)} of {\rm(\hyperref[II.def1.1]{1.1})}. A
presheaf $F$ is a sheaf (\resp a separated presheaf)
for the topology generated (\hyperref[II.subsubsec1.1.6]{1.1.6})
by the $K(X)$ if and only if, for every object $X$ of $C$
and every sieve $R\in K(X)$, the map
$$
\Hom_{C^{\widehat{\phantom{,}}}}(X,F)\longrightarrow
\Hom_{C^{\widehat{\phantom{,}}}}(R,F)
$$
is a bijection (\resp an injection).
\end{enonce*}

\begin{enonce*}{Corollary 2.4}\etiquette[2.4]{II.cor2.4}
In\pageoriginale particular, let $C$ be a $𝒰$-category
equipped with a pretopology. A presheaf $F$ is a
sheaf (\resp a separated presheaf) if and only if,
for every object $X$ of $C$ and every covering family
$(X_{\alpha}\rightarrow X)$ belonging to $\Cov(X)$, the diagram
of sets
$$
F(X)\rightarrow \prod_{\alpha\in A}F(X_{\alpha})\rightrightarrows
\prod_{(\alpha,\beta)\in A\times A}F(X_{\alpha}\times_X X_{\beta})
$$
is exact (\resp the map
$$
F(X)\longrightarrow \prod_{\alpha\in A}F(X_{\alpha})
$$
is injective).
\end{enonce*}

\begin{enonce*}[remark]{Proof}
Apply \hyperref[II.cor2.3]{2.3}, then (\hyperrefext[I][I.cor3.5]{3.5}) and
(\hyperrefext[I][I.prop2.12]{2.12}).

Corollary \hyperref[II.cor2.4]{2.4} recovers the definition
given in [1].
\end{enonce*}

\begin{enonce*}{Definition 2.5}\etiquette[2.5]{II.def2.5}
Let $C$ be a $𝒰$-category. The canonical topology
of $C$ is the finest topology for which the representable functors
are sheaves (\hyperref[II.prop2.2]{2.2}). A sieve
covering $X$ for the canonical topology will be called a universally
strictly epimorphic sieve. A covering family for the
canonical topology will be called a universally strictly epimorphic
family. If, in addition, the morphisms of the covering family
are pullbackable, the family will be called universally effective
epimorphic $[2]$.
\end{enonce*}

\begin{enonce*}[remark]{Remark 2.5.1}\etiquette[2.5.1]{II.rem2.5.1}
For almost all the sites that have been used up to now,
the topology is \emph{coarser} than the canonical topology; in
other words, the representable functors on $C$ are
sheaves, \ie the covering families of $C$ are universally strictly
epimorphic. The\pageoriginale only exception to this
rule is the site $\widehat{C}$ studied in ${\rm n^{\circ}}$ 5
(whose topology is \emph{finer,} and most often strictly
finer, than the canonical topology).
\end{enonce*}

\begin{enonce*}{Proposition 2.6}\etiquette[2.6]{II.prop2.6}
A sieve $R\hookrightarrow X$ is universally strictly
epimorphic if and only if, for every object $Y\rightarrow X$
of $C$ over $X$, and for every object $Z$ of $C$,
the map
$$
\Hom_C(Y,Z)\longrightarrow \varprojlim_{C/(Y\times_X R)}\Hom_C (.,
Z)
$$
is a bijection.
\end{enonce*}

\begin{enonce*}[remark]{Proof}
Immediate upon applying \hyperref[II.prop2.2]{2.2} and then (\hyperrefext[I][I.sec5.3]{5.3}).
\end{enonce*}

\begin{enonce*}[remark]{Remarks 2.7}\etiquette[2.7]{II.rem2.7}
\begin{enumerate}
\renewcommand{\labelenumi}{\theenumi)}
\item Proposition \hyperref[II.prop2.6]{2.6} gives a characterization of the universally
  strictly epimorphic sieves of a category $C$
  that is independent of the universe in which the presheaves take
  their values, provided only that the morphism sets
  $\Hom(X,Y)$ of $C$ belong to that universe. It
  thus makes it possible to define the canonical topology for any
  category $C$, without having to specify the universes.

\item Let $C$ be a site whose underlying category is a
  $𝒰$-category, let $F$ be a $𝒰$-presheaf, and let
  $𝒱$ be a universe containing $𝒰$. The category
  underlying $C$ is a $𝒱$-category, and $F$ may
  be regarded as a $𝒱$-presheaf. The
  $𝒰$-presheaf $F$ is a sheaf (\resp a
  separated presheaf) if and only if the
  $𝒱$-presheaf $F$ is a sheaf (\resp a
  separated presheaf). In other words, the condition for a
  presheaf $F$ to be a sheaf (\resp a separated
  presheaf) does not depend (in the sense just\pageoriginale
  specified) on the universe in which the presheaf $F$ takes its
  values. In particular, let $C$ be a site and let $F$ be a
  $𝒰$-presheaf; we say that $F$ is a sheaf (\resp a
  separated presheaf) if there exists a universe $𝒱$
  containing $𝒰$ such that the category $C$ is a
  $𝒱$-category and such that $F$ is a $𝒱$-sheaf
  (\resp a separated $𝒱$-presheaf). This property
  does not depend on the universe $𝒱$.
\end{enumerate}
\end{enonce*}

% END INCLUDED SGA4_English_current_buildable_source_20260906/source/02/drafts/SGA4_II_LINES_213_421_English.texfrag


% BEGIN INCLUDED SGA4_English_current_buildable_source_20260906/source/02/drafts/SGA4_II_LINES_422_562_English.texfrag
\section{The sheaf associated with a presheaf}\etiquette[3]{II.sec3}

\begin{enonce*}{Definition 3.0.1}\etiquette[3.0.1]{II.def3.0.1}
Let $C$ be a site. A topologically generating family (or, when no
confusion can result, a generating family) of $C$ is a set $G$ of
objects of $C$ such that every object of $C$ is the target of a
covering family of morphisms of $C$ (\hyperref[II.def1.2]{1.2}) whose
sources are elements of $G$.
\end{enonce*}

\begin{enonce*}{Definition 3.0.2}\etiquette[3.0.2]{II.def3.0.2}
Let $𝒰$ be a universe. A $𝒰$-site is a site $C$ whose underlying
category is a $𝒰$-category
{\rm(\hyperrefext[I][I.def1.1]{1.1})}, and which possesses a small
topologically generating family. Let $C$ be a $𝒰$-category; a
$𝒰$-topology on $C$ is a topology on $C$ making $C$ a $𝒰$-site. A
site $C$ is called $𝒰$-small, or, by abuse of language, small, if its
underlying category is small {\rm(\hyperrefext[I][I.sec1.0]{1.0})}.
\end{enonce*}

\setcounter{subsubsection}{2}
\subsubsection{}\etiquette[3.0.3]{II.subsubsec3.0.3}
It follows immediately from the definitions that every topology finer
than a $𝒰$-topology is a $𝒰$-topology, and that a small site is a
$𝒰$-site.

\begin{enonce*}{Proposition 3.0.4}\etiquette[3.0.4]{II.prop3.0.4}
Let\pageoriginale $C$ be a $𝒰$-site and $G$ a small topologically
generating family of $C$. For every $X\in\ob C$, denote by $J_G(X)$
the set of sieves covering $X$ generated by a family of morphisms
$$
(Y_{\alpha}\xrightarrow{u_{\alpha}} X), \alpha\in A, \mbox{ where }
Y_{\alpha}\in G{\quoi}.
$$
\begin{enumerate}
\renewcommand{\labelenumi}{\theenumi)}
\item The set $J_G(X)$ is small.

\item The set $J_G(X)$ is cofinal in the set $J(X)$ of all sieves
covering $X$, ordered by inclusion.

\item For every $R\in J_G(X)$, there exists a small epimorphic family
{\rm(\hyperrefext[I][I.sec10.2]{10.2})}
$(u_\alpha:Y\rightarrow R)$, $\alpha\in A$, with $Y\in G$.
\end{enumerate}
\end{enonce*}

\begin{enonce*}[remark]{Proof}
\begin{enumerate}
\renewcommand{\labelenumi}{\theenumi)}
\item Put $A(X)=\amalg_{Y\in G}\hom(Y,X)$. The set $A(X)$ is small
(\hyperrefext[I][I.sec0]{0}) and
$\card(J_G(X))\leq2^{\card(A(X))}$. Consequently $J_G(X)$ is small.

\item Let $R\in J(X)$. Put $A(R)=\amalg_{Y\in G}\Hom(Y,R)$, and let
$R'$ be the sieve of $X$ generated by the family
$(Y\xrightarrow{u}R\hookrightarrow X)$, $u\in A(R)$. One has
$R'\subset R$, and it suffices to show that $R'$ is covering. By
axiom (T 2) of \hyperref[II.def1.1]{1.1}, it suffices to show that for
every morphism $Z\rightarrow R$, $Z\in\ob C$, the sieve
$R'\times_XZ$ of $Z$ is covering. Now the sieve $R'\times_XZ$
contains the sieve generated by the family of all morphisms
$Y\rightarrow Z$ with $Y\in G$, which is covering by hypothesis.
Therefore the sieve $R'\times_XZ$ of $Z$ is covering (axiom T2 of
\hyperref[II.def1.1]{1.1}).

\item Let $R\in J_G(X)$. The family $(Y\xrightarrow{u}R)$,
$u\in A(R)=\amalg_{Y\in G}\Hom(Y,R)$, is epimorphic by hypothesis.
Now, for every $Y\in\ob C$, $\Hom(Y,R)$ is contained in $\Hom(Y,X)$,
which is small. Consequently $A(R)$ is small.
\end{enumerate}
\end{enonce*}

\setcounter{subsubsection}{4}
\subsubsection{}\etiquette[3.0.5]{II.subsubsec3.0.5}
Let $C$ be a $𝒰$-site, $𝒱$ a universe such that $C\in𝒱$ and
$𝒰\subset𝒱$, and $G$ a $𝒰$-small topologically generating family of
$C$. The category $C^{\widehat{\phantom{,}}}$ of presheaves of
$𝒰$-sets on $C$ is a $𝒱$-category (I \hyperref[I.sec1.1.1]{1.1.1}). Let
$X$\pageoriginale be an object of $C$. The set $J(X)$ of sieves
covering $X$ is $𝒱$-small and, ordered by inclusion, is cofiltered
(\hyperref[II.subsubsec1.1.1]{1.1.1}). For every $𝒰$-presheaf $F$,
the inductive limit
$\varinjlim_{J(X)}\Hom_{C^{\widehat{\phantom{,}}}}(R,F)$ is therefore
representable by an element of $𝒱$
(\hyperrefext[I][I.def2.4.1]{2.4.1}). Moreover, it follows from
(\hyperref[II.prop3.0.4]{3.0.4} 3)) that, for every $R\in J_G(X)$,
$\Hom_{C^{\widehat{\phantom{,}}}}(R,F)$ is $𝒰$-small, and since
$J_G(X)$ is a $𝒰$-small set cofinal in $J(X)$ (\loccit), it follows
from (\hyperrefext[I][I.sec2.4.2]{2.4.2}) that
$\varinjlim_{J(X)}\Hom_{C^{\widehat{\phantom{,}}}}(R,F)$ is
$𝒰$-small. Choose, for every $F$ and every $X$, an element of $𝒰$
representing this inductive limit, and put
$$
LF(X)=\varinjlim_{J(X)}\Hom_{C^{\widehat{\phantom{,}}}}(.,F).
$$
Let $g:Y\rightarrow X$ be a morphism of $C$. The base-change functor
$g^\ast:J(X)\rightarrow J(Y)$ defines a map
$$
LF(g):LF(X)\longrightarrow LF(Y){\quoi},
$$
which makes $X\mapsto LF(X)$ a presheaf on $C$.

Since the morphism $\id_X:X\rightarrow X$ is an element of $J(X)$,
for every object $X$ of $C$ one has a map
$$
\ell(F)(X):F(X)\longrightarrow LF(X){\quoi},
$$
which visibly defines a morphism of functors
$$
\ell(F):F\longrightarrow LF{\quoi}.
$$
It is moreover clear that $F\mapsto LF$ is a functor in $F$ and that
the morphisms $\ell(F)$ define a morphism
$$
\ell:\Id\longrightarrow L{\quoi}.
$$
Finally, let $R\hookrightarrow X$ be a covering sieve of $X$. The
definition of $LF(X)$ and (\hyperrefext[I][I.prop1.4]{1.4}) provide a
map
$$
Z_R:\Hom_{C^{\widehat{\phantom{,}}}}(R,F)\longrightarrow
\Hom_{C^{\widehat{\phantom{,}}}}(X,LF){\quoi},
$$
and,\pageoriginale for every morphism $Y\xrightarrow{g}X$ of $C$, the
definition of the functor $LF$ shows that the following diagram is
commutative:
\[\vcenter{
\xymatrix{ \qquad Z_R :& \Hom_{C^{\widehat{\phantom{a}}}}(R,F)
\ar[r]\ar[d] &
\Hom_{C^{\widehat{\phantom{a}}}}(X,LF)\ar[d]\\
Z_{\fprod{R}{Y}{X}}:&  \Hom_{C^{\widehat{\phantom{a}}}}(\fprod{R}{Y}{X},F) \ar[r]&  \Hom_{C^{\widehat{\phantom{a}}}}(Y,LF)\\
}}\leqno{(*)}
\]
(the vertical arrows are the evident arrows).

\noindent Let us now copy [2].


% END INCLUDED SGA4_English_current_buildable_source_20260906/source/02/drafts/SGA4_II_LINES_422_562_English.texfrag


% BEGIN INCLUDED SGA4_English_current_buildable_source_20260906/source/02/drafts/SGA4_II_LINES_563_764_English.texfrag
\begin{enonce*}{Lemma 3.1}\etiquette[3.1]{II.lem3.1}
\begin{enumerate}
\renewcommand{\labelenumi}{\theenumi)}
\item For every covering sieve $i_R:R\hookrightarrow X$ and every
$u:R\rightarrow F$, the diagram
\[\vcenter{
\xymatrix@L=5pt@R=1.35cm@C=3.5cm@M=.2cm{
F\ar[r]^{\ell(F)} & LF\\
R\ar[u]^{u}\ar@{^{(}->}_{i_R}[r] & X \ar[u]_{Z_R(u)} }}\leqno{(**)}
\]
is commutative.

\item For every morphism $v:X\rightarrow LF$, there exists a covering
sieve $R$ of $X$ and a morphism $u:R\rightarrow F$ such that
$Z_R(u)=v$.

\item Let $Y$ be an object of $C$, and let
$u,v:Y\rightrightarrows F$ be two morphisms such that
$\ell(F)\circ u=\ell(F)\circ v$. Then the kernel of the pair $(u,v)$
is a covering sieve of $Y$.

\item Let $R,R'$ be two covering sieves of $X$, and let
$u:R\rightarrow F$, $u':R'\rightarrow F$ be two morphisms. For
$Z_R(u)=Z_{R'}(u')$, it is necessary and sufficient that $u$ and $u'$
coincide on a covering sieve $R''\hookrightarrow R\times_XR'$.
\end{enumerate}
\end{enonce*}

\begin{enonce*}[remark]{Proof}
The only nontrivial assertion is assertion 1). We must show that
$Z_R(u)\circ i_R=\ell(F)\circ u$. For this it suffices to show
(\hyperrefext[I][I.prop3.4]{3.4}) that the composites of these
morphisms with every morphism $Y\xrightarrow{g}R$ ($Y$ an object
of\pageoriginale $C$) are equal. Consider $f=i_R\circ g$ and the fiber
product $R\times_XY$:
\[\vcenter{
\xymatrix@L=5pt@R=1cm@C=3.5cm@M=.2cm{
F  \ar[r]^{\ell(F)}   & LF \\
R \ar[u]^{u}        \ar@{^{(}->}_{i_R}[r]  & X \ar[u]_{Z_R(u)} \\
\fprod{R}{Y}{X} \ar[u]^{g'} \ar[r]^(.35){\sim}_{i'} & Y
\ar[lu]^{g}\ar[u]_{f}. }}
\]

In the diagram above, $i'$ is a monomorphism admitting a section.
Consequently (\hyperref[II.def1.3]{1.3} 3) $i'$ is an isomorphism. By
the definition of the morphism $\ell(F)$, the morphism
$Z_{R\times_XY}(u\circ g')$ is equal to $\ell(F)\circ u\circ g$. On
the other hand, the commutativity of diagram (*) gives the equality
$Z_{R\times_XY}(u\circ g')=Z_R(u)\circ f$, and hence one indeed has
$\ell(F)\circ u\circ g=Z_R(u)\circ i_R\circ g$.
\end{enonce*}

\begin{enonce*}{Proposition 3.2}\etiquette[3.2]{II.prop3.2}
\begin{enumerate}
\renewcommand{\labelenumi}{\theenumi)}
\item The functor $L$ is left exact
{\rm(\hyperrefext[I][I.cor2.3.2]{2.3.2}).}

\item For every presheaf $F$, $LF$ is a separated presheaf.

\item The presheaf $F$ is separated if and only if the morphism
$\ell(F):F\rightarrow LF$ is a monomorphism. The presheaf $LF$ is then
a sheaf.

\item The following properties are equivalent:
\begin{enumerate}
\renewcommand{\theenumii}{\roman{enumii}}
\renewcommand{\labelenumii}{(\theenumii)}
\item $\ell(F):F\rightarrow LF$ is an isomorphism.

\item $F$ is a sheaf.
\end{enumerate}
\end{enumerate}
\end{enonce*}

\begin{enonce*}[remark]{Proof}
\begin{enumerate}
\renewcommand{\labelenumi}{\theenumi)}
\item It suffices to show (\hyperrefext[I][I.prop3.1]{3.1}) that, for
every object $X$ of $C$, the functor $F\mapsto LF(X)$ commutes with
finite projective limits. By the definition of a projective limit, the
functor
$$
F\longmapsto \Hom_{C^{\widehat{\phantom{,}}}}(R,F)\quad
R\hookrightarrow X\in \ob(J(X))
$$
commutes\pageoriginale with projective limits, and the inductive limit
$\varinjlim_{J(X)}$ commutes with finite projective limits because
$J(X)$ is a cofiltered set (\hyperrefext[I][I.def2.7]{2.7}).

\item Let $X$ be an object of $C$, and let
$f,g:X\rightrightarrows LF$ be two morphisms that coincide on a
covering sieve $R\hookrightarrow X$ of $X$. By
\hyperref[II.lem3.1]{3.1}\quad 2), there exists a covering sieve
$R'\hookrightarrow X$, which may always be supposed contained in $R$,
and two morphisms $u,v:R'\rightrightarrows F$ such that
$Z_{R'}(u)=f$ and $Z_{R'}(v)=g$. By
\hyperref[II.lem3.1]{3.1} 1), one then has
$\ell(F)\circ u=\ell(F)\circ v$. Consequently
(\hyperref[II.lem3.1]{3.1} 4)) $u$ and $v$ coincide on a covering
sieve $R''\hookrightarrow R'$. Let $w$ be the restriction of $u$ to
$R''$. One has $Z_{R''}(w)=Z_{R'}(u)=Z_{R'}(v)$, and consequently
$f=g$. Thus the presheaf $LF$ is separated.

\item If $F$ is separated, the morphism $\ell(F)$ is a monomorphism
because a filtered inductive limit of monomorphisms is a monomorphism.
If $\ell(F)$ is a monomorphism, the presheaf $F$ is a subpresheaf of a
separated presheaf, and therefore is separated. Let us show that $LF$
is then a sheaf. Let $i:R\hookrightarrow X$ be a sieve covering an
object $X$ of $C$, and let $u:R\rightarrow LF$ be a morphism. It
suffices to show that $u$ factors through $X$. Put
$R'=\fprod{F}{R}{LF}$ and $v=Z_{R'}(\pr_1)$:
\[
\xymatrix@L=5pt@R=1cm@C=1.5cm{
F \ar[rr]^{\ell(F)}               && LF                      &\\
R' \ar[u]^{\pr_1} \ar[rr]^{\pr_2}   && R \ar[r]^{i} \ar[u]_{u} & X \ar@{-}[lu]_{v}\\
R''\ar[u]^{p_1} \ar[rr]^{p_2}     && Y  \ar[u]_{m}               &
\quad . }
\]
To\pageoriginale show that $u=v\circ i$, it suffices to show (I \hyperref[I.prop3.4]{3.4})
that for every morphism $m:Y\rightarrow R$ ($Y$ an object of $C$),
$v\circ i\circ m=u\circ m$. Put $R''=Y\times_RR'$, and let $p_1,p_2$
be the projections. The projection
$\xymatrix@M=5pt{R''\ar@{^{(}->}[r]^{p_2}&Y}$ is a monomorphism and
makes $R''$ a covering sieve of $Y$ (\hyperref[II.lem3.1]{3.1} 2)).
One has $v\circ i\circ\pr_2=\ell(F)\circ\pr_1$
(\hyperref[II.lem3.1]{3.1} 1)), and hence
$v\circ i\circ\pr_2=u\circ m\circ p_2$. It follows that
$v\circ i\circ\pr_2\circ p_1=u\circ\pr_2\circ p_1$, \ie
$v\circ i\circ m\circ p_2=u\circ m\circ p_2$. Since the presheaf
$LF$ is separated, one has $v\circ i\circ m=u\circ m$, \cqfd

\item Clear.
\end{enumerate}
\end{enonce*}

\begin{enonce*}[remark]{Remark 3.3}\etiquette[3.3]{II.rem3.3}
Let $J'(X)$ be a cofinal subset of $J(X)$. One has
$$
LF(X)=\varinjlim_{J'(X)}\Hom_{C^{\widehat{\phantom{,}}}}(.,F){\quoi}.
$$

In particular, if the topology of $C$ is defined by a pretopology
$X\mapsto\Cov(X)$ (\hyperref[II.subsubsec1.1.3]{1.1.3}), the functor
$L$ can be described using the covering families of $\Cov(X)$
(\hyperrefext[I][I.prop2.12]{2.12} and
\hyperrefext[I][I.cor3.5]{3.5}). Writing the formulas explicitly, one
recovers the construction of [
\end{enonce*}

\begin{enonce*}{Theorem 3.4}\etiquette[3.4]{II.theo3.4}
Let $C$ be a $𝒰$-site. The inclusion functor
$i:C^{\widetilde{\phantom{,}}}\hookrightarrow
C^{\widehat{\phantom{,}}}$ from sheaves into presheaves admits a left
adjoint $\underline{a}$, which is left exact
(\hyperrefext[I][I.cor2.3.2]{2.3.2}):
\[
\xymatrix@=3cm@M=.2cm{
 \widetilde{C}  \ar @{^{(}->} @<-1ex>[r]_{i} & C^{\widehat{\phantom{,}}} \ar@<-1ex>[l]_{\underline{a}}
\qquad .}
\]

The functor $i\circ\underline{a}$ is canonically isomorphic to the
functor $L\circ L$ (cf. \hyperref[II.subsubsec3.0.5]{3.0.5}). For
every presheaf $F$, the adjunction morphism
$F\rightarrow i\circ\underline{a}(F)$ is deduced, via the preceding
isomorphism, from the morphism
$\ell(LF)\ell(F):F\rightarrow L\circ L(F)$.
\end{enonce*}

\begin{enonce*}{Definition 3.5}\etiquette[3.5]{II.def3.5}
The\pageoriginale sheaf $\underline{a}F$ is called the sheaf associated
with the presheaf $F$.
\end{enonce*}

Theorem \hyperref[II.theo3.4]{3.4} follows immediately from Proposition
\hyperref[II.prop3.2]{3.2}.

\begin{enonce*}{Proposition 3.6}\etiquette[3.6]{II.prop3.6}
Let $C$ be a $𝒰$-site and let $𝒱\supset𝒰$ be a universe. Denote by
$C^{\widehat{\phantom{,}}}_{𝒰}$ and
$C^{\widetilde{\phantom{,}}}_{𝒰}$ (respectively
$C^{\widehat{\phantom{,}}}_{𝒱}$ and
$C^{\widetilde{\phantom{,}}}_{𝒱}$) the categories of $𝒰$-presheaves
and $𝒰$-sheaves (respectively $𝒱$-presheaves and $𝒱$-sheaves), and
$\underline{a}_{𝒰}:C^{\widehat{\phantom{,}}}_{𝒰}\rightarrow
C^{\widetilde{\phantom{,}}}_{𝒰}$ (respectively
$\underline{a}_{𝒱}:C^{\widehat{\phantom{,}}}_{𝒱}\rightarrow
C^{\widetilde{\phantom{,}}}_{𝒱}$) the corresponding ``associated
sheaf'' functors. The diagram
\[
\xymatrix@C=3.5cm@R=1.25cm{ C^{\widehat{\phantom{a}}}_{U}
\ar[r]^{\underline{a}_{\underline{U}}}
\ar @{_{(}->}[d]& C^{\widetilde{\phantom{a}}}_{\underline{U}}\ar @{_{(}->}[d]\\
C^{\widehat{\phantom{a}}}_{\underline{V}}
\ar[r]^{\underline{a}_{\underline{V}}} &
C^{\widetilde{\phantom{a}}}_{\underline{V}} , }
\]
where the vertical functors are the canonical inclusions, is
commutative up to canonical isomorphism.
\end{enonce*}

\begin{enonce*}[remark]{Proof}
This follows from the construction of the functors $\underline{a}_U$
and $\underline{a}_V$ (\hyperref[II.theo3.4]{3.4} and
\hyperref[II.subsubsec3.0.5]{3.0.5}).
\end{enonce*}


% END INCLUDED SGA4_English_current_buildable_source_20260906/source/02/drafts/SGA4_II_LINES_563_764_English.texfrag


% BEGIN INCLUDED SGA4_English_current_buildable_source_20260906/source/02/drafts/SGA4_II_LINES_765_975_English.texfrag
\section{Exactness properties of the category of sheaves}\etiquette[4]{II.sec4}

The exactness properties of the category of sheaves follow from the
exactness properties of the category of presheaves via Theorem
\hyperref[II.theo3.4]{3.4}. This section makes this philosophy explicit
through some of the most useful standard statements.

\begin{enonce*}{Theorem 4.1}\etiquette[4.1]{II.theo4.1}
Let $C$ be a $𝒰$-site, let $\widetilde{C}$ be the category of
sheaves, let $\underline{a}:C^{\widehat{\phantom{,}}}\rightarrow
\widetilde{C}$ be the associated-sheaf functor, and let
$i:\widetilde{C}\rightarrow C^{\widehat{\phantom{,}}}$ be the
inclusion functor.
\begin{enumerate}
\renewcommand{\labelenumi}{\theenumi)}
\item The functor $\underline{a}$ commutes with inductive limits and is exact.

\item The\pageoriginale $𝒰$-inductive limits in
$\widetilde{C}$ are representable. For every small category $I$ and
every functor $E:I\rightarrow \widetilde{C}$, the canonical morphism
$$
\varinjlim_I E\longrightarrow \underline{a}(\varinjlim_I i\circ E)
$$
is an isomorphism.

\item The $𝒰$-projective limits in $\widetilde{C}$ are
representable. For every object $X$ of $C$, the functor on
$\widetilde{C}:F\mapsto F(X)$ commutes with projective limits,
\ie the inclusion functor $i:\widetilde{C}\rightarrow \widehat{C}$
commutes with projective limits.
\end{enumerate}
\end{enonce*}

\begin{enonce*}[remark]{Proof}
These properties follow essentially from Theorem
\hyperref[II.theo3.4]{3.4} and from
(\hyperrefext[I][I.prop2.11]{2.11}).
\end{enonce*}

Thus, in the category of sheaves, products indexed by an element of
$𝒰$, fiber products, sums indexed by an element of $𝒰$, amalgamated
sums, kernels, cokernels, images, and coimages are representable.

\begin{enonce*}{Corollary 4.1.1}\etiquette[4.1.1]{II.cor4.1.1}
Let $C$ be a $𝒰$-site and let $F$ be a sheaf of sets on $C$. The
canonical homomorphism
$$
\varinjlim_{C/F} \underline{a}X\longrightarrow F
$$
is an isomorphism.
\end{enonce*}

\begin{enonce*}[remark]{Proof}
This follows from (\hyperrefext[I][I.prop3.4]{3.4}) and the fact that
$\underline{a}$ commutes with inductive limits.
\end{enonce*}

\begin{enonce*}{Proposition 4.2}\etiquette[4.2]{II.prop4.2}
Every morphism of $\widetilde{C}$ that is both an epimorphism and a
monomorphism is an isomorphism.
\end{enonce*}

\begin{enonce*}[remark]{Proof}
Let $u:G\rightarrow H$ be a morphism of $\widetilde{C}$ that is both
an epimorphism and a monomorphism. First observe that the morphism $u$
is a monomorphism of\pageoriginale presheaves
(\hyperref[II.theo4.1]{4.1} 1)). Construct the amalgamated sum
$H\amalg_G H=K$ and the two canonical morphisms $i_1$,
$i_2:H\rightrightarrows K$ in the category of presheaves. Since $u$
is a monomorphism of presheaves, one immediately verifies that the
following diagram is cartesian and cocartesian
(\hyperrefext[I][I.sec10.1]{10.1}):
\[
\xymatrix@C=3cm@R=1.5cm{
G \ar[r]^{u} \ar[d]_{u}& H \ar[d]^{i_2}\\
H \ar[r]^{i_1}& K . }
\]
Applying the ``associated sheaf'' functor, one therefore obtains a
cartesian and cocartesian diagram in the category of sheaves
(\hyperref[II.theo4.1]{4.1} 1)):
\[\vcenter{
\xymatrix@C=3cm@R=1.5cm{
G \ar[r]^{u} \ar[d]_{u}& H \ar[d]^{\underline{a}i_2}\\
H \ar[r]^{\underline{a}i_1}& \underline{a}K . }}\leqno{(*)}
\]
Since $u$ is an epimorphism of sheaves, the morphism
$\underline{a}i_1$ is an isomorphism, and since diagram (*) is
cartesian, the morphism $u$ is an isomorphism.
\end{enonce*}

\begin{enonce*}{Proposition 4.3}\etiquette[4.3]{II.prop4.3}
\begin{enumerate}
\renewcommand{\labelenumi}{\theenumi)}
\item The inductive limits in $C^{\widetilde{\phantom{,}}}$ that are
representable are universal {\rm(\hyperrefext[I][I.def2.5]{2.5}).}

\item Every epimorphic family {\rm(\hyperrefext[I][I.sec10.2]{10.2})}
of morphisms is universally effective epimorphic
(\hyperref[II.prop2.6]{2.6}).

\item Every equivalence relation is universally effective {\rm(I
\hyperref[I.sec10.6]{10.6}).}

\item Filtered $𝒰$-inductive limits commute with finite projective
limits {\rm(\hyperrefext[I][I.prop2.6]{2.6})}.
\end{enumerate}
\end{enonce*}

\begin{enonce*}[remark]{Proof}
Assertions\pageoriginale 1) through 4) hold in the category of sets,
hence in the category of presheaves
(\hyperrefext[I][I.prop3.1]{3.1}). Assertions 1) and 4) then follow
immediately from the corresponding assertions for presheaves and from
\hyperref[II.theo4.1]{4.1}. Let us prove 3). Let $p_1$, $p_2:$
$\xymatrix{R\ar@<2pt>[r] \ar@<-2pt>[r]& X}$ be an equivalence
relation in $C^{\widetilde{\phantom{,}}}$. The diagram
$\xymatrix{R \ar@<2pt>[r]^{p_1} \ar@<-2pt>[r]_{p_2}&X}$ is then an
equivalence relation of presheaves
(\hyperref[II.cor4.1.1]{4.1.1}). There is therefore a morphism of
presheaves $u:X\rightarrow Y$ such that the diagram
\[
\xymatrix@C=3cm@R=1.25cm{
R \ar[r]^{p_2} \ar[d]_{p_1}& X\ar[d]^{u}\\
X \ar[r]^{u}& Y }
\]
is cartesian and cocartesian in the category of presheaves. Applying
the ``associated sheaf'' functor, one obtains a cartesian and
cocartesian diagram in the category of sheaves
(\hyperref[II.theo4.1]{4.1} 1)). It follows that
$\xymatrix{R \ar@<2pt>[r]^{p_1} \ar@<-2pt>[r]_{p_2}&X}$ is an
effective equivalence relation. Since every equivalence relation is
effective by what precedes, every equivalence relation is universally
effective. Let us prove 2). Let $(u_i:X_i\rightarrow X)$, $i\in I$,
be an epimorphic family in $C^{\widetilde{\phantom{,}}}$. By
(II\ \hyperref[II.prop2.6]{2.6}) and
(\hyperrefext[I][I.prop2.12]{2.12}), it suffices to prove the following
assertions:
\begin{enumerate}
\renewcommand{\theenumi}{\alph{enumi}}
\renewcommand{\labelenumi}{\theenumi)}
\item For every morphism of sheaves $v:Y\rightarrow X$, the family
of morphisms $\pr_{2,i}:X_i\times_X Y\rightarrow Y$, $i\in I$, is an
epimorphic family in $C^{\widetilde{\phantom{,}}}$.

\item For every sheaf $Z$, the following diagram of sets is exact:
$$
\Hom(X,Z)\rightarrow \prod_{i\in I} \Hom(X_i,Z)\rightrightarrows
\prod_{(i,k)\in I\times I}\Hom(X_i\times_X X_k,Z)
$$
\end{enumerate}
\end{enonce*}

Let\pageoriginale $X'$ be the union of the images of the morphisms
$u_i$, formed in the category of presheaves, and let
$j:X'\hookrightarrow X$ be the canonical injection. For every
$i\in I$, the morphism $u_i:X_i\rightarrow X$ factors as a morphism
$u'_i:X_i\rightarrow X'$ followed by $j$, and the family
$u'_i:X_i\rightarrow X'$, $i\in I$, is an epimorphic family in
$C^{\widehat{\phantom{,}}}$. Since $j$ is a monomorphism, the morphism
$\underline{a}j:\underline{a}X'\rightarrow X$ is a monomorphism of
sheaves (\hyperref[II.theo4.1]{4.1}). Since, for every $i\in I$, one
has $u_i=\underline{a}j\underline{a}u'_i$, $\underline{a}j$ is an
epimorphism of sheaves. It follows that $\underline{a}j$ is an
isomorphism (\hyperref[II.prop4.2]{4.2}). Let $v:Y\rightarrow X$ be a
morphism of sheaves. Base change then yields a morphism
$j_v:X'\times_X Y\rightarrow Y$ and morphisms
$u'_{i,v}:X_i\times_X Y\rightarrow X'\times_X Y$. Since
$\underline{a}$ commutes with fiber products, $\underline{a}j_v$ is
an isomorphism. Since epimorphic families in
$C^{\widehat{\phantom{,}}}$ remain epimorphic under base change, the
family $u'_{i,v}$, $i\in I$, is epimorphic in
$C^{\widehat{\phantom{,}}}$. Since $\underline{a}$ commutes with
inductive limits, the family
$\underline{a}u'_{i,v}:X_i\times_X Y\rightarrow
\underline{a}(X'\times_X Y)$, $i\in I$, is epimorphic in
$C^{\widetilde{\phantom{,}}}$, proving a). Let us prove b). Let $Z$
be a sheaf. Since $\underline{a}j$ is an isomorphism, the map
$$
\Hom(X,Z)\longrightarrow \Hom(X',Z)
$$
is a bijection. Since the $u'_i$, $i\in I$, form an epimorphic family
in $C^{\widehat{\phantom{,}}}$, the following diagram of sets is
exact:
$$
\Hom(X',Z)\rightarrow \prod_{i\in I} \Hom(X_i,Z)\rightrightarrows
\prod_{(i,k)\in I\times I}\Hom(X_i\times_{X'}X_k, Z)
$$
Finally, since $X'$ is a subpresheaf of $X$, the presheaf
$X_i\times_{X'}X_k$, $(i,k)\in I\times I$, is canonically isomorphic
to $X_i\times_X X_k$, proving b).

\begin{enonce*}[remark]{Remarks 4.3.1}\etiquette[4.3.1]{II.rem4.3.1}
\begin{enumerate}
\renewcommand{\labelenumi}{\theenumi)}
\item Proposition\pageoriginale \hyperref[II.prop4.3]{4.3} 2)
formally gives us a second proof of \hyperref[II.prop4.2]{4.2}: it is
clear that, in any category, a morphism that is both a universal
effective epimorphism and a monomorphism is in fact an isomorphism.

\item It follows from the preceding propositions that every morphism
in the category of sheaves factors uniquely as an effective
epimorphism followed by an effective monomorphism. This property
naturally generalizes, to nonadditive categories, axiom (AB2) for
abelian categories (coim $\simeq$ im) [Tohoku].
\end{enumerate}
\end{enonce*}

\setcounter{subsection}{4} \setcounter{subsubsection}{-1}

% END INCLUDED SGA4_English_current_buildable_source_20260906/source/02/drafts/SGA4_II_LINES_765_975_English.texfrag


% BEGIN INCLUDED SGA4_English_current_buildable_source_20260906/source/02/drafts/SGA4_II_LINES_976_1126_English.texfrag
\subsubsection{}\etiquette[4.4.0]{II.subsubsec4.4.0}
Let $C$ be a $𝒰$-site. Composing the functor $h:C\rightarrow
C^{\widehat{\phantom{,}}}$ (\hyperrefext[I][I.sec1.3.1]{1.3.1}) with
the associated sheaf functor yields a functor$$
\epsilon_C:C\longrightarrow C^{\widetilde{\phantom{,}}}{\quoi},
$$
called the \emph{canonical functor} from $C$ to
$C^{\widetilde{\phantom{,}}}$, which will be used constantly in what
follows. The functor $\epsilon_C$ commutes with finite projective
limits. When the topology on $C$ is coarser than the canonical
topology, $\epsilon_C$ is fully faithful and commutes with projective
limits; moreover, it is then defined even when $C$ is not assumed to
possess a small topologically generating family. We shall
not study in detail the behavior of $\epsilon_C$ with respect to
inductive limits (\cf [SGA 3 IV]). We shall, however, need the results
below.

\begin{enonce*}{Theorem 4.4}\etiquette[4.4]{II.theo4.4}
Let $C$ be a $𝒰$-site and let $(u_i:X_i\rightarrow X)$, $i\in
I$, be a family of morphisms of $C$ with target $X$. The following
conditions are equivalent \footnote[1]{When the site $C$ does not
necessarily possess a small topologically generating family and when
the topology on $C$ is coarser than the canonical topology, one has
ii) $\Rightarrow$ i) (\cf the proof of
\hyperref[II.theo4.4]{4.4}).}:
\begin{enumerate}
\renewcommand{\theenumi}{\roman{enumi}}
\renewcommand{\labelenumi}{\theenumi)}
\item The family $(\epsilon_C u_i:\epsilon_C X_i\rightarrow \epsilon_C
  X)$, $i\in I$, is an epimorphic family in
  $C^{\widetilde{\phantom{,}}}$ {\rm(\hyperrefext[I][I.sec10.2]{10.2}).}\pageoriginale

\item The family $(u_i:X_i\rightarrow X)$, $i\in I$, is a covering
  family of $C$ (\hyperref[II.def1.2]{1.2}).
\end{enumerate}
\end{enonce*}

\begin{enonce*}[remark]{Proof}
ii) $\Rightarrow$ i). Let $R\hookrightarrow X$ be the sieve generated
by the $u_i$, $i\in I$ (\hyperrefext[I][I.sec4.3.3]{4.3.3}), and let
$u'_i:X_i\rightarrow R$ be the morphisms induced by the $u_i$. The
morphisms $u'_i$, $i\in I$, form an epimorphic family in
$C^{\widehat{\phantom{,}}}$, and $R$ is a covering sieve of $X$.
Consequently, for every sheaf $F$, the map
$$
\Hom(X,F)\longrightarrow \Hom(R,F)
$$
is bijective, and the map
$$
\Hom(R,F)\longrightarrow \prod_i \Hom(X_i,F)
$$
is injective. Thus the map
$$
\Hom(X,F)\longrightarrow \prod_i\Hom(X_i,F)
$$
is injective, and consequently the map
$$
\Hom(\underline{a}X,F)\longrightarrow \prod_i \Hom(\underline{a}X_i,F)
$$
is injective, which proves i).

\noindent i) $\Rightarrow$ ii). With the notation introduced above,
let $J_R:R\hookrightarrow X$ be the canonical inclusion into $X$ of
the sieve generated by the $u_i$, $i\in I$. It follows from i) that
the morphism of sheaves
$\underline{a}J_R:\underline{a}R\rightarrow \underline{a}X$ is both an
epimorphism and a monomorphism in the category of sheaves. It is
therefore an isomorphism of sheaves (\hyperref[II.prop4.2]{4.2}). With the
notation of ${\rm n^{\circ}}$ 3, we have a commutative diagram:
\[
\xymatrix@C=2.5cm@R=1.5cm@M=.25cm{
X \ar[r]^{\ell(X)} & LX \ar[r]^{\ell(LX)} & \mathsf{\underline{a}}X\\
R \ar[r]^{\ell(R)} \ar@{^{(}->}[u]^{J_R}& LR \ar[r]^{\ell(LR)}
\ar@{^{(}->}[u]^{LJ_R}& \mathsf{\underline{a}}R
\ar[u]^{\wr}_{\mathsf{\underline{a}}J_R}. }
\]
First of\pageoriginale all, by \hyperref[II.lem3.1]{3.1} 2), there exist a covering sieve
$J_{S_1}:S_1\rightarrow X$ and a morphism
$u_1:S_1\rightarrow LR$ such that
\begin{equation}
\ell(LX)\circ \ell (X)\circ J_{S_1}=\underline{a}J_R\circ \ell (LR)\circ
u_1=\ell(LX)\circ LJ_R\circ u_1{\quoi}.\tag{4.6.1}
\end{equation}\etiquette[4.6.1]{II.eq4.6.1}
Put $m=\ell(X)\circ J_{S_1}$ and $n=LJ_R\circ u_1$. The morphisms
$m$ and $n$ are morphisms from $S_1$ to $LX$. Let
$S_2\hookrightarrow S_1$ be the kernel of the pair of arrows $(m,n)$.
For every object $Y$ of $C$ and every morphism $\alpha:Y\rightarrow S_1$,
one has, by \hyperref[II.eq4.6.1]{4.6.1}, $\ell(LX)\circ m\circ \alpha
=\ell(LX)\circ n\circ\alpha$. It then follows from
(\hyperref[II.lem3.1]{3.1} 3)) that the kernel $S_2\times_{S_1}Y$ of
the pair of arrows $(m\circ \alpha,n\circ \alpha)$ is a covering
sieve of $Y$. Consequently, by axiom (T 2) for topologies,
$J_{S_2}:S_2\hookrightarrow X$ is a covering sieve of $X$. Thus there
is a morphism $u_2:S_2\rightarrow LR$ and a commutative diagram:
\hypertarget{sga4.exp.II.equation.4.6.2}{}\[\vcenter{
\xymatrix@C=2.5cm@R=.8cm@M=.2cm{
S_2 \ar@{^{(}->}[rr]^{J_{S_{2}}} \ar[dd]_{u_2}&  &X \ar[dd]^{\ell(X)}\\
                  &R\ar@{^{(}->}[ru]_{J_R} \ar[ld]^{\ell(R)}&\\
LR \ar[rr]^{LJ_R}  &  &LX&. }}\leqno{(4.6.2)}
\]

Let $Y$ be an object of $C$ and let $\beta:Y\rightarrow S_2$ be a
morphism. By \hyperref[II.lem3.1]{3.1} 2) and 1), there exist a
covering sieve $J_{Q_1}:Q_1\hookrightarrow Y$ and a morphism
$v_1:Q_1\rightarrow R$ such that
$u_2\circ\beta\circ J_{Q_1}=\ell(R)\circ v_1$. Put
$f=J_{S_{2}}\circ \beta\circ J_{Q_1}$ and $g=J_R\circ v_1$. Let
$Q_2$ be the kernel of the pair $(f,g):Q_1\rightrightarrows X$, and
let $J_{Q_2}:Q_2\hookrightarrow Y$ be the canonical inclusion. For
every object $Z$ of $C$ and every morphism $\gamma:Z\rightarrow Q_1$,
the commutativity of (\hyperref[II.cor4.6.2]{4.6.2}) gives:
\begin{gather*}
\ell(X)\circ f\circ \gamma =\ell(X)\circ J_{S_2}\circ \beta\circ
J_{Q_1}\circ \gamma =LJ_R\circ u_2\circ \beta\circ J_{Q_1}\circ \gamma=\\
=LJ_R\circ \ell(R)\circ v_1\circ \gamma =\ell (X)\circ J_R\circ
v_1\circ \gamma = \ell(X)\circ g\circ \gamma{\quoi}.
\end{gather*}
It\pageoriginale then follows from \hyperref[II.lem3.1]{3.1} 3) that the
kernel $Q_2\times_{Q_1}Z$ of the pair $(f\circ\gamma,g\circ\gamma)$
is a covering sieve of $Z$. By axiom (T 2) for topologies,
$J_{Q_2}:Q_2\hookrightarrow Y$ is a covering sieve of $Y$. Thus, for
every morphism $\beta:Y\rightarrow S_2$, there exist a covering sieve
$J_{Q_2}:Q_2\hookrightarrow Y$ and a morphism $v_2:Q_2\rightarrow
R$ such that the following diagram is commutative:
\[
\xymatrix@C=2.5cm@R=1.5cm@M=.2cm{
Y\ar[r]^{\beta}   & S_2\ar[r]^{J_{S_2}} & X\\
Q_2 \ar[r]^{v_2} \ar@{^{(}->}[u]^{J_{Q_2}}& R
\ar@{^{(}->}[ru]_{J_{R}}  &}
\]

Denote by $S_2\cap R$ the fiber product of $S_2$ and $R$ over $X$.
The sieve (of $Y$) $(S_2\cap R)\times_{S_2}Y$ contains the sieve
$Q_2$, and hence $(S_2\cap R)\times_{S_2}Y$ is a covering sieve of
$Y$. It then follows from axiom (T 2) for topologies that
$S_2\cap R$ is a covering sieve of $X$, and consequently the sieve
$R$, which contains $S_2\cap R$, is a covering sieve of $X$, \cqfd
\end{enonce*}

\setcounter{subsubsection}{3}
\begin{enonce*}{Corollary 4.4.4}\etiquette[4.4.4]{II.cor4.4.4}
Let $T$ be the topology of the $\Uu$-site $C$, and let $T'$ (\resp
$T''$) be the finest topology on $C$ among those for which the
objects of $\widetilde{C}$ are sheaves \emph{(\resp
separated presheaves)} (\hyperref[II.prop2.2]{2.2}). Then $T=T'=T''$.
\end{enonce*}

Indeed, one trivially has $T\leq T'\leq T''$, and
\hyperref[II.theo4.4]{4.4} and \hyperref[II.prop2.2]{2.2} imply that
$T''\leq T$, \cqfd

% END INCLUDED SGA4_English_current_buildable_source_20260906/source/02/drafts/SGA4_II_LINES_976_1126_English.texfrag


% BEGIN INCLUDED SGA4_English_current_buildable_source_20260906/source/02/drafts/SGA4_II_LINES_1127_1298_English.texfrag
\begin{enonce*}{Definition 4.5}\etiquette[4.5]{II.def4.5}
An initial object of a category $A$ is an object $\phi_A$ that
represents the empty inductive limit, \ie such that, for every $X\in
\ob A$, there exists a unique arrow $\phi_A\rightarrow X$. An object
$\phi_A$ is called \emph{strict initial} if it is initial and every
morphism with target $\phi_A$ is an isomorphism. Let $(S_i)$, $i\in I$,
be a family of objects of a category $A$. Suppose that the sum
$S=\amalg_{i\in I} S_i$ is representable. The sum $S$ is called
disjoint if the structure morphisms $S_i\rightarrow S$ are
pullbackable, if they are monomorphisms, and if, for every pair $i$,
$j$ with $i\neq j$, the fiber product $S_i\times_S S_j$ is an initial
object of $A$. The sum $S$ is called disjoint and
universal\pageoriginale if it is disjoint and remains a disjoint sum
after every base change $T\rightarrow S$; it follows that, for every
pair $i$, $j$ with $i\neq j$, the objects $S_i\times_S S_j$ are strict
initial objects of $A$.
\end{enonce*}

\begin{enonce*}[remark]{Example 4.5.1}\etiquette[4.5.1]{II.exem4.5.1}
In the category of sets, sums are disjoint and universal. The same is
therefore true in every category of presheaves of sets
(\hyperrefext[I][I.prop3.1]{3.1}), and hence in every category of
sheaves of sets on a $𝒰$-site (\hyperref[II.theo4.1]{4.1}); in
particular, the initial object of $C^{\widetilde{\phantom{,}}}$ is
strict.
\end{enonce*}

\begin{enonce*}{Proposition 4.6}\etiquette[4.6]{II.prop4.6}
Let $C$ be a $𝒰$-category and let $(s_i:X_i\rightarrow X)$,
$i\in I$, be a family of morphisms of $C$. For every
$𝒰$-topology $\Tt$ on $C$ (\hyperref[II.def3.0.2]{3.0.2}), denote by
$C^{\widetilde{\phantom{,}}}_{\Tt}$ the corresponding category of
sheaves and by $\epsilon_{\Tt}:C\rightarrow
C^{\widetilde{\phantom{,}}}_{\Tt}$ the corresponding functor.
\begin{enumerate}
\renewcommand{\labelenumi}{\theenumi)}
\item Let $\Tt$ be a $𝒰$-topology such that
$$
\coprod_{i\in I}
\epsilon_{\Tt}(X_i)\xrightarrow{(\epsilon_{\Tt}(s_i))}\epsilon_{\Tt}(X)
$$
is an isomorphism. Then, for every topology $\Tt'$ finer than $\Tt$,
the morphism
$$
\coprod_{i\in
  I}\epsilon_{\Tt'}(X_i)\xrightarrow{(\epsilon_{\Tt'}(s_i))}\epsilon_{\Tt'}(X)
$$
is an isomorphism.

\item Let $\Tt$ be a $𝒰$-topology on $C$. The following
  properties {\rm(i)} and {\rm(ii)} are equivalent:
\begin{itemize}
\item[i)]
\begin{itemize}
\item[a)] The family $(s_i:X_i\rightarrow X)$, $i\in I$, is covering
  for $\Tt$.

\item[b)] For every $i\in I$, the diagonal morphism of presheaves
  $\Delta_i:X_i\hookrightarrow X_i\times_X X_i$\pageoriginale
  is sent by the associated sheaf functor (for $\Tt$) to an
  isomorphism (as is the case if the $s_i$ are monomorphisms).

\item[c)] For every pair $(i,j)$ of elements of $I$ with $i\neq j$,
  the presheaf $X_i\times_X X_j$ is sent by the associated sheaf
  functor (for $\Tt$) to the initial object of
  $C^{\widetilde{\phantom{,}}}_{\Tt}$.
\end{itemize}
\item[ii)] $\epsilon_{\Tt}(X)$ is the sum of the
  $\epsilon_{\Tt}(X_i)$.
\end{itemize}
\end{enumerate}
\end{enonce*}

\begin{enonce*}[remark]{Proof}
\begin{enumerate}
\renewcommand{\labelenumi}{\theenumi)}
\item A sheaf for $\Tt'$ is a sheaf for $\Tt$. Consequently, for
  every sheaf $F$ for $\Tt'$, the morphism
$$
\Hom_{C^{\widehat{\phantom{,}}}}(X,F)\longrightarrow \prod_{i\in I}
\Hom_{C^{\widehat{\phantom{,}}}}(X_i,F)
$$
is an isomorphism, which proves the assertion.

\item By definition, property ii) holds if and only if the morphism
  of presheaves
$$
\Phi=(h(s_i),i\in I):\coprod_i h(X_i)\longrightarrow h(X)
$$
is sent by the associated sheaf functor to an isomorphism, \ie
(\hyperref[II.prop4.2]{4.2}) if and only if $\underline{a}(\Phi)$ is
both an epimorphism and a monomorphism of sheaves. By
(\hyperref[II.theo4.4]{4.4}), the morphism $\underline{a}(\Phi)$ is an
epimorphism if and only if property a) holds. The diagonal morphism
$$
\Delta:\coprod_i h(X_i)\longrightarrow (\coprod_i
h(X_i))\times_{h(X)}(\coprod_i h(X_i))
$$
is the direct sum of a family $\Delta_{i,j}$, $(i,j)\in I\times I$,
of morphisms defined as follows:
\begin{itemize}
\item[$\beta)$] When $i=j$, $\Delta_{i,i}$ is the diagonal morphism
  $h(X_i)\rightarrow h(X_i)\times_{h(X)}h(X_i)$.

\item[$\gamma)$] When $i\neq j$, $\Delta_{i,j}$ is the morphism
  $\phi_{C^{\widehat{\phantom{,}}}}\rightarrow
  h(X_i)\times_{h(X)}h(X_j)$ ($\phi_{C^{\widehat{\phantom{,}}}}$
  denotes the initial object of $C^{\widehat{\phantom{,}}}$).
\end{itemize}
\end{enumerate}
\end{enonce*}

The\pageoriginale morphism $\underline{a}(\Phi)$ is a monomorphism if
and only if $\underline{a}(\Delta)$ is an isomorphism, and by
(\hyperref[II.theo4.1]{4.1}) $\underline{a}(\Delta)$ is an isomorphism
if and only if the $\underline{a}(\Delta_{i,j})$ are isomorphisms, i.e.
if and only if properties b) and c) hold.

\begin{enonce*}{Corollary 4.6.1}\etiquette[4.6.1]{II.cor4.6.1}
Let $C$ be a $𝒰$-category and $X$ an object of $C$.
\begin{enumerate}
\renewcommand{\labelenumi}{\theenumi)}
\item Let $\Tt$ be a $𝒰$-topology on $C$. The object $X$ of $C$ is
  sent by the associated sheaf functor (for the topology $\Tt$) to
  the initial object of $C^{\widetilde{\phantom{,}}}_{\Tt}$ if and
  only if the empty sieve covers $X$.

\item When $X$ is a strict initial object of $C$, the sheaf associated
  with $X$, for every $𝒰$-topology finer than the canonical topology,
  is an initial object.
\end{enumerate}
\end{enonce*}

\begin{enonce*}[remark]{Proof}
\begin{enumerate}
\renewcommand{\labelenumi}{\theenumi)}
\item Take the indexing set $I$ in \hyperref[II.prop4.6]{4.6} 2) to
be the empty set.

\item By 1) and \hyperref[II.prop4.6]{4.6} 1), it is enough to show
  that the empty sieve covers $X$ for the canonical topology; \ie by
  (\hyperref[II.prop2.6]{2.6}), it is enough to show that, for every
  object $Y\rightarrow X$ over $X$, $Y$ is an initial object of $C$,
  which follows from the definition (\hyperref[II.def4.5]{4.5}).
\end{enumerate}
\end{enonce*}

\begin{enonce*}{Corollary 4.6.2}\etiquette[4.6.2]{II.cor4.6.2}
Let $C$ be a $𝒰$-site and let $(s_i:X_i\rightarrow X)$, $i\in I$,
be a family of pullbackable morphisms of $C$ with the same target and
having the following properties:
\begin{itemize}
\item[$\alpha)$] The family of the $s_i$, $i\in I$, is covering.

\item[$\beta)$] For every $i\in I$, $s_i:X_i\rightarrow X$ is a
  monomorphism.

\item[$\gamma)$] For every pair $(i,j)$ of elements of $I$ with
  $i\neq j$, $X_i\times_X X_j$ is covered by the empty sieve, \ie for
  every sheaf $F$ on $C$, $F(X_i\times_X X_j)$ is a singleton set.
\end{itemize}
Then $\epsilon_C(X)$ is the sum of the $\epsilon_C(X_i)$, $i\in I$.
\end{enonce*}

\begin{enonce*}[remark]{Proof}
This follows\pageoriginale immediately from
\hyperref[II.prop4.6]{4.6} 2) and
\hyperref[II.cor4.6.1]{4.6.1} 1).
\end{enonce*}

% END INCLUDED SGA4_English_current_buildable_source_20260906/source/02/drafts/SGA4_II_LINES_1127_1298_English.texfrag


% BEGIN INCLUDED SGA4_English_current_buildable_source_20260906/source/02/drafts/SGA4_II_LINES_1299_1528_English.texfrag
\begin{enonce*}{Corollary 4.6.3}\etiquette[4.6.3]{II.cor4.6.3}
Let $C$ be a $𝒰$-category and let $(s_i:X_i\rightarrow X)$,
$i\in I$, be a family of pullbackable morphisms with the same target.
Suppose that the canonical topology of $C$ is a $𝒰$-topology. The
following conditions are equivalent:
\begin{enumerate}
\renewcommand{\theenumi}{\roman{enumi}}
\renewcommand{\labelenumi}{\theenumi)}
\item There exists a $𝒰$-topology $\Tt$ on $C$, coarser than the
canonical topology, such that $\epsilon_\Tt(X)$ is the sum of the
$\epsilon_\Tt(X_i)$, $i\in I$.

\item The object $X$ is the disjoint and universal sum
(\hyperref[II.def4.5]{4.5}) (in $C$) of the $X_i$.
\end{enumerate}
\end{enonce*}

\begin{enonce*}[remark]{Proof}
ii) $\Rightarrow$ i). Since $X$ is the universal sum of the $X_i$,
$i\in I$, the family of the $s_i$, $i\in I$, is covering for the
canonical topology of $C$ (\hyperref[II.prop2.6]{2.6}). Condition
$\alpha$) of \hyperref[II.cor4.6.2]{4.6.2} is therefore satisfied when
$C$ is endowed with the canonical topology. Condition $\beta$) is
evidently satisfied, and property $\gamma$) follows from
\hyperref[II.prop4.6]{4.6} 1) 2), whence i).

\noindent i) $\Rightarrow$ ii). Let $\Tt$ be a topology on $C$ coarser
than the canonical topology such that $\epsilon_\Tt(X)$ is the sum of
the $\epsilon_\Tt(X_i)$. For every sheaf $F$ for $\Tt$, the canonical
map $\Hom(X,F)\rightarrow\prod_i\Hom(X_i,F)$ is a bijection. Moreover,
every representable presheaf is a sheaf. It follows immediately that
$X$ is the sum in $C$ of the $X_i$, $i\in I$. Applying this to the
case where the set $I$ is empty, one sees that if an object of $C$ is
carried to the initial object of $C^{\widetilde{\phantom{,}}}$, that
object is an initial object of $C$. The functor
$\epsilon_\Tt:C\rightarrow C^{\widetilde{\phantom{,}}}$ is fully
faithful and commutes with finite projective limits. Consequently,
condition b) of \hyperref[II.cor4.6.2]{4.6.2} implies that the $s_i$,
$i\in I$, are monomorphisms of $C$, and condition c) implies, by what
precedes, that the $X_i\times_XX_j$, $i\neq j$, are initial objects of
$C$. Consequently $X$ is the disjoint sum of the\pageoriginale $X_i$.
Since $\epsilon_\Tt$ commutes with fiber products, this last property
is stable under every base change $Y\rightarrow X$, and therefore $X$
is the disjoint and universal sum of the $X_i$, $i\in I$.
\end{enonce*}

\begin{enonce*}{Corollary 4.6.4}\etiquette[4.6.4]{II.cor4.6.4}
Let $C$ be a $𝒰$-category and $X$ an object of $C$. Suppose that the
canonical topology on $C$ is a $𝒰$-topology. The following properties
are equivalent:
\begin{enumerate}
\renewcommand{\theenumi}{\roman{enumi}}
\renewcommand{\labelenumi}{\theenumi)}
\item There exists on $C$ a topology coarser than the canonical
topology such that $\epsilon_\Tt(X)$ is an initial object of
$C^{\widetilde{\phantom{,}}}_\Tt$.

\item $X$ is a strict initial object of $C$.
\end{enumerate}
\end{enonce*}

\begin{enonce*}[remark]{Proof}
Take the empty set for the set $I$ in \hyperref[II.cor4.6.3]{4.6.3}.
\end{enonce*}

\begin{enonce*}{Proposition 4.7}\etiquette[4.7]{II.prop4.7}
Let $C$ be a $𝒰$-category, let $C^{\widetilde{\phantom{,}}}$ be the
category of sheaves on $C$ for the canonical topology, let
$\epsilon_C:C\rightarrow C^{\widetilde{\phantom{,}}}$ be the canonical
functor (\hyperref[II.subsubsec4.4.0]{4.4.0}), and let
$R\rightrightarrows X$ be an equivalence relation in $C$ admitting a
cokernel $Y$ in $C$. Let $\Pi:X\rightarrow Y$ be the canonical
morphism and suppose it is pullbackable. Consider the following
properties:
\begin{enumerate}
\renewcommand{\theenumi}{\roman{enumi}}
\renewcommand{\labelenumi}{\theenumi)}
\item $\epsilon_C(Y)$ is the quotient of the equivalence relation
$\epsilon_C(R)\rightrightarrows\epsilon_C(X)$.

\item The equivalence relation $R\rightrightarrows X$ is universally
effective {\rm(cf. ${\rm n^{\circ}}$ 7).}
\end{enumerate}
One has {\rm{ii) $\Rightarrow$ i).}} When $C$ possesses a
$𝒰$-topology coarser than the canonical topology, one has
{\rm{i) $\Rightarrow$ ii).}}
\end{enonce*}

\begin{enonce*}[remark]{Proof}
ii) $\Rightarrow$ i). The canonical morphism
$R\rightarrow X\times_YX$ is an isomorphism. Let $S\rightarrow Y$ be
the sieve generated by $\pi:X\rightarrow Y$. For every presheaf $F$,
there is an exact diagram (\hyperrefext[I][I.prop2.12]{2.12}):
$$
\Hom(S,F)\rightarrow \Hom(X,F)\rightrightarrows \Hom(R,F){\quoi}.
$$
It\pageoriginale therefore suffices to show that, for every sheaf $F$
for the canonical topology, the map
$$
\Hom(X,F)\longrightarrow \Hom(S,F)
$$
is a bijection, \ie it suffices to show that $S$ is covering for the
canonical topology, or equivalently that the sieve generated by
$\pi:X\rightarrow Y$ is covering, which follows from \hyperref[II.def2.5]{Definition 2.5}.

\noindent i) $\Rightarrow$ ii). The functor
$\epsilon_C:C\rightarrow C^{\widetilde{\phantom{,}}}$ commutes with
finite projective limits and is fully faithful. The canonical morphism
$\epsilon_C(R)\rightarrow J_C(X)\times_{J_C(Y)}J_C(X)$ is an
isomorphism (\hyperref[II.prop4.3]{4.3}). Consequently the canonical
morphism $R\rightarrow X\times_YX$ is an isomorphism. The morphism
$\epsilon_C(X)\rightarrow\epsilon_C(Y)$ is an epimorphism, which
implies (\hyperref[II.theo4.4]{4.4}) that $\pi:X\rightarrow Y$ is a
covering morphism of $C$ for the canonical topology, \cqfd
\end{enonce*}

\begin{enonce*}{Proposition 4.8}\etiquette[4.8]{II.prop4.8}
Let $C$ be a $𝒰$-site. The category of sheaves on $C$ has the
following properties:
\begin{enumerate}
\renewcommand{\theenumi}{\alph{enumi}}
\renewcommand{\labelenumi}{\theenumi)}
\item Finite projective limits are representable.

\item Direct sums indexed by an element of $𝒰$ are representable.
They are disjoint and universal (\hyperref[II.def4.5]{4.5}).

\item Equivalence relations are universally effective.
\end{enumerate}
\end{enonce*}

This was proved in \hyperref[II.theo4.1]{4.1}, 2) and 3),
\hyperref[II.prop4.3]{4.3}, 3), and
\hyperref[II.exem4.5.1]{4.5.1}.

These properties have been singled out because they will later allow
us to characterize the $𝒰$-categories equivalent to categories of
sheaves on categories belonging to $𝒰$
(\hyperrefext[IV][IV.theo1.2]{1.2}).

\begin{enonce*}[remark]{Remark}
Recall that property a) is equivalent to the property:

There exists a terminal object, and fiber products are representable
(\hyperrefext[I][I.sec2.3.1]{2.3.1}).
\end{enonce*}

\setcounter{subsection}{8}
\subsection{}\etiquette[4.9]{II.subsec4.9}
Let\pageoriginale $C$ be a site whose underlying category is a
$𝒰$-category. Then the category $\widetilde{C}$ of $𝒰$-sheaves on
$C$ satisfies the conditions considered in
\hyperrefext[I][I.cor7.3]{7.3} ((i) or (ii), as desired), and therefore,
by \loccit, the various variants of the notion of a \emph{generating
family} in $\widetilde{C}$ considered in
\hyperrefext[I][I.def7.1]{7.1} coincide. This being understood:

\begin{enonce*}{Proposition 4.10}\etiquette[4.10]{II.prop4.10}
With the preceding notation, consider the canonical functor
$\epsilon:C\rightarrow\widetilde{C}$
(\hyperref[II.subsubsec4.4.0]{4.4.0}), and let $G\subset\ob C$ be a
topologically generating family in $C$
(\hyperref[II.def3.0.1]{3.0.1}). Then the family
$(\epsilon(X))_{X\in G}$ of objects of $\widetilde{C}$ is a generating
family. In particular, the family $(\epsilon(X))_{X\in\ob C}$ of
objects of $\widetilde{C}$ is generating.
\end{enonce*}

We must prove that every morphism $u:F\rightarrow F'$ in
$\widetilde{C}$ such that
\begin{equation}
\Hom(\epsilon(X),F)\longrightarrow \Hom(\epsilon(X),F')\tag{4.10.1}
\end{equation}\etiquette[4.10.1]{II.eq4.10.1}
is a bijection for every $X\in G$, is an isomorphism. The preceding
map identifies with the map
\begin{equation}
u(X):F(X)\longrightarrow F'(X){\quoi},\tag{4.10.2}
\end{equation}\etiquette[4.10.2]{II.eq4.10.2}
and we must show that if the latter is bijective for every $X\in G$,
then it is so for every $X\in\ob C$. Let
$(X_i\rightarrow X)_{i\in I}$ be a covering family of $X$ by objects
of $G$, and, for every pair of indices $(i,j)$ in $I$, consider the
set of all morphisms $X_{ijk}\rightarrow X_i\times_XX_j$ in
$\widehat{C}$, with $X_{ijk}\in G$ ($k$ varying in an index set
$I_{i,j}$). One then obtains a morphism of exact diagrams of sets
\[
\xymatrix{ F(X) \ar[r] \ar[d]& \displaystyle{{\prod_i}} F(X_i)
\ar[d]^{\wr} \ar@<3pt>[r] \ar@<-3pt>[r] & \displaystyle{\prod_{ijk}}
F(X_{ijk}) \ar[d]^{\wr}\\
F'(X) \ar[r] & \displaystyle{{\prod_i}} F'(X_i) \ar@<3pt>[r]
\ar@<-3pt>[r] & \displaystyle{\prod_{ijk}} F'(X_{ijk})&, }
\]
where\pageoriginale, by hypothesis, the last two vertical arrows are
bijections. The same is therefore true of the first vertical arrow,
\cqfd

\begin{enonce*}{Corollary 4.11}\etiquette[4.11]{II.cor4.11}
Suppose that $C$ is a $𝒰$-site (\hyperref[II.def3.0.2]{3.0.2}). Then
the category $\widetilde{C}$ is a $𝒰$-category and admits a
$𝒰$-small generating family.
\end{enonce*}

Indeed, by hypothesis one may take for $G$ in
\hyperref[II.prop4.10]{4.10} a small generating family, which proves
the existence of a small generating family. On the other hand, if
$F,F'\in\ob\widetilde{C}$, a morphism from $F$ to $F'$ is known once
one knows the morphism \eqhyperref[II.eq4.10.1]{4.10.1}, \ie
\eqhyperref[II.eq4.10.2]{4.10.2}, for every $X\in\ob G$ (I \hyperref[I.sec7.1.1]{7.1.1}),
whence the map
$$
\Hom(F,F')\longrightarrow \prod_{X\in G}\Hom(F(X),F'(X))
$$
is injective. Since the second member is $𝒰$-small, so is the first,
which proves that $\widetilde{C}$ is a $𝒰$-category.

\begin{enonce*}{Corollary 4.12}\etiquette[4.12]{II.cor4.12}
Let $C$ be a $𝒰$-site. Then, for every object $X$ of
$\widetilde{C}$, the set of subobjects of $X$ and the set of quotient
objects of $X$ are $𝒰$-small.
\end{enonce*}

This follows from \hyperrefext[I][I.prop7.4]{7.4}, respectively
\hyperrefext[I][I.cor7.5]{7.5}, which apply by
\hyperref[II.cor4.11]{4.11}.


% END INCLUDED SGA4_English_current_buildable_source_20260906/source/02/drafts/SGA4_II_LINES_1299_1528_English.texfrag


% BEGIN INCLUDED SGA4_English_current_buildable_source_20260906/source/02/drafts/SGA4_II_LINES_1529_1645_English.texfrag
\section{Extension of a topology from $C$ to
  $C^{\shat{\;\hbox{}}}$}\etiquette[5]{II.sec5}

\begin{enonce*}{Proposition 5.1}\etiquette[5.1]{II.prop5.1}
Let $C$ be a $𝒰$-site and let $f:H\rightarrow K$ be a morphism of
$C^{\widehat{\phantom{,}}}$. The following conditions are
equivalent:
\begin{itemize}
\item[i)] For every $X\rightarrow K$, with $X\in \ob C$, the image of
  the corresponding morphism $H\times_K X\rightarrow X$ is a
  covering sieve of $X$.

\item[ii)] The morphism $\underline{a}(f):\underline{a}H\rightarrow \underline{a}K$
  between the associated sheaves is an epimorphism\pageoriginale in
  $C^{\widetilde{\phantom{,}}}$.

\item[ii bis)] For every sheaf $F$ on $C$, the map
  $\Hom(K,F)\rightarrow \Hom(H,F)$ induced by $f$ is injective.
\end{itemize}
\end{enonce*}

\begin{enonce*}[remark]{Proof}
It is clear that ii) is equivalent to ii bis).
\end{enonce*}

Let us prove that (ii) $\Rightarrow$ (i). By
\hyperref[II.theo4.4]{4.4}, (i) means that
$u:\underline{a}(H\times_K X)\rightarrow \underline{a}(X)$ is an
epimorphism. Now, since $\underline{a}(H\times_K X)\simeq
\underline{a}(H)\times_{\underline{a}(K)}\underline{a}(X)$, because
$\underline{a}$ commutes with finite $\varprojlim$
(\hyperref[II.theo4.1]{4.1}), the morphism in question is obtained
from $\underline{a}(f):\underline{a}(H)\rightarrow
\underline{a}(K)$ by base change. Since the epimorphisms of
$\widetilde{C}$ are universal, this proves (ii) $\Rightarrow$ (i).
Conversely, since the family of all $X_i\rightarrow K$, with
$X_i\in\ob C$, is epimorphic, the same is true of the family of
$\underline{a}(X_i)\rightarrow \underline{a}(K)$ in $\widetilde{C}$
(\hyperref[II.theo4.1]{4.1}). Hence, to verify that
$\underline{a}(f):\underline{a}(H)\rightarrow \underline{a}(K)$ is
epimorphic, it suffices to do so after every base change of the
preceding kind $\underline{a}(X_i)\rightarrow \underline{a}(K)$,
which proves i) $\Rightarrow$ ii), \cqfd

\begin{enonce*}{Definition 5.2}\etiquette[5.2]{II.def5.2}
\begin{enumerate}
\renewcommand{\labelenumi}{\theenumi)}
\item A morphism $f:H\rightarrow K$ satisfying the three equivalent
  conditions of \hyperref[II.prop5.1]{5.1} is called a
  \emph{covering morphism}. A family of morphisms
  $f_i:H_i\rightarrow K$, $i\in I$, with the same target is called
  \emph{covering} if the corresponding morphism
  $f:\amalg_i H_i\rightarrow K$ is covering.

\item A morphism $f:H\rightarrow K$ is called \emph{bicovering} if it
  is covering and if the diagonal morphism
  $H\rightarrow H\times_K H$ is covering. A family
  $f_i:H_i\rightarrow K$, $i\in I$, with the same target is called
  bicovering if the corresponding morphism
  $f:\amalg_i H_i\rightarrow K$ is bicovering.
\end{enumerate}
\end{enonce*}

By condition i) of \hyperref[II.prop5.1]{5.1}, to say that a family
$(f_i:H_i\rightarrow K)$, $i\in I$, is covering means that, for every
morphism $X\rightarrow K$, with $X\in \ob C$, the image of the family
$H_i\times_K X\rightarrow X$, $i\in I$, is a covering sieve
of\pageoriginale $X$; or equivalently, by ii), that the family of
morphisms $\underline{a}(f_i):\underline{a}H_i\rightarrow
\underline{a}K$ in $C^{\widetilde{\phantom{,}}}$ is epimorphic
(taking into account that the functor $\underline{a}$ commutes with
direct sums (\hyperref[II.theo4.1]{4.1})).\footnote[1]{Note that the
property of a morphism or a family of morphisms of
$C^{\widehat{\phantom{,}}}$ being covering (\resp bicovering) is
stable under base change.}

\begin{enonce*}{Proposition 5.3}\etiquette[5.3]{II.prop5.3}
Let $C$ be a $𝒰$-site and let $f:H\rightarrow K$ be a morphism of
$C^{\widehat{\phantom{,}}}$. The following conditions are
equivalent:
\begin{itemize}
\item[i)] The morphism $f$ is bicovering
  (\hyperref[II.def5.2]{5.2}).

\item[i bis)] The morphism $f$ is covering and, for every object $X$
  of $C$ and every pair of morphisms $u$, $v:X\rightrightarrows H$
  such that $fu=fv$, the kernel of $(u,v)$ is a covering sieve of
  $X$.

\item[ii)] The morphism
  $\underline{a}(f):\underline{a}H\to \underline{a}K$ is an
  isomorphism in $C^{\widetilde{\phantom{,}}}$.

\item[ii bis)] For every sheaf $F$ on $C$, the map
  $\Hom(K,F)\rightarrow \Hom(H,F)$ is a bijection.
\end{itemize}
\end{enonce*}

\begin{enonce*}[remark]{Proof}
The equivalence i) $\Leftrightarrow$ i bis) follows from condition i)
of \hyperref[II.prop5.1]{5.1}, applied to the diagonal morphism
$H\rightarrow H\times_K H$; the equivalence ii) $\Leftrightarrow$ ii
bis) is trivial.

\noindent i) $\Rightarrow$ ii). The morphism
$\underline{a}(f):\underline{a}H\rightarrow \underline{a}K$ is an
epimorphism (\hyperref[II.prop5.1]{5.1}). Since the functor
$\underline{a}$ commutes with the formation of fiber products
(\hyperref[II.theo4.1]{4.1}), the diagonal morphism
$\underline{a}H\rightarrow
\underline{a}H\times_{\underline{a}K}\underline{a}H$ is an
epimorphism (\hyperref[II.prop5.1]{5.1}). Since a diagonal morphism is
always a monomorphism, it is an isomorphism
(\hyperref[II.prop4.2]{4.2}). It follows that the morphism
$\underline{a}(f)$ is a monomorphism. It is therefore an isomorphism
(\hyperref[II.prop4.2]{4.2}).

\noindent ii) $\Rightarrow$ i). Since the functor $\underline{a}$
commutes with the formation of fiber products, the morphism $f$ and
the diagonal morphism $H\rightarrow H\times_K H$ are carried by
$\underline{a}$ to isomorphisms. In particular, they are carried by
$\underline{a}$ to epimorphisms. They are therefore covering, \cqfd
\end{enonce*}

\setcounter{subsection}{3} \setcounter{subsubsection}{0}

% END INCLUDED SGA4_English_current_buildable_source_20260906/source/02/drafts/SGA4_II_LINES_1529_1645_English.texfrag


% BEGIN INCLUDED SGA4_English_current_buildable_source_20260906/source/02/drafts/SGA4_II_LINES_1646_1819_English.texfrag
\subsubsection{}\etiquette[5.3.1]{II.subsubsec5.3.1}
It\pageoriginale follows from \hyperref[II.prop5.3]{5.3}, and from
the fact that $\underline{a}$ commutes with direct sums, that a family
$(f_i:H_i\rightarrow K)$, $i\in I$, of morphisms of
$C^{\widehat{\phantom{,}}}$ is bicovering if and only if the $f_i$
induce an isomorphism from the direct sum
$\amalg_i \underline{a}H_i$ onto $\underline{a}K$ on the associated
sheaves, or equivalently if and only if, for every sheaf $F$, the map
$$
\Hom(K,F)\longrightarrow \prod_i \Hom(H_i,F)
$$
is bijective.

\begin{enonce*}{Proposition 5.4}\etiquette[5.4]{II.prop5.4}
Let $C$ be a $𝒰$-site. There exists a (necessarily unique) topology
on $C^{\widehat{\phantom{,}}}$ such that a family
$H_i\rightarrow K$ of arrows of $C^{\widehat{\phantom{,}}}$ with the
same target is covering for this topology if and only if it is
covering in the sense of \hyperref[II.def5.2]{5.2}. It is also the
coarsest topology on $C^{\widehat{\phantom{,}}}$ among the topologies
$T$ having the following properties:
\begin{enumerate}
\renewcommand{\theenumi}{\alph{enumi}}
\renewcommand{\labelenumi}{\theenumi)}
\item $T$ is finer than the canonical topology of
  $C^{\widehat{\phantom{,}}}$ (\ie every epimorphic family in
  $C^{\widehat{\phantom{,}}}$ is covering for $T$).

\item Every covering family in $C$ is covering in
  $C^{\widehat{\phantom{,}}}$.
\end{enumerate}
\end{enonce*}

\begin{enonce*}[remark]{Proof}
We shall confine ourselves to giving indications. We leave it to the
reader to show that the families which are covering in the sense of
\hyperref[II.def5.2]{5.2} are the covering families of a topology
$T_C$ on $C^{\widehat{\phantom{,}}}$ (one uses
\hyperref[II.theo4.1]{4.1}). The covering families of the canonical
topology on $C^{\widehat{\phantom{,}}}$ are the epimorphic families
in $C^{\widehat{\phantom{,}}}$ (\hyperref[II.prop2.6]{2.6} and I
\hyperref[I.prop3.1]{3.1}). Since $\underline{a}$ commutes with inductive limits
(\hyperref[II.theo4.1]{4.1}), the topology $T_C$ is finer than the
canonical topology of $C^{\widehat{\phantom{,}}}$. Moreover, the
covering families of $C$ are covering families for $T_C$
(\hyperref[II.prop5.1]{5.1}). Thus, if $T'$ denotes the coarsest
topology on $C^{\widehat{\phantom{,}}}$ having properties a) and b),
$T_C$ is finer than $T'$. Let $(f_i:H_i\rightarrow K)$, $i\in I$, be
a\pageoriginale covering family for $T_C$. We show that $(f_i,i\in
I)$ is a covering family for $T'$. Let
$s_i:H_i\rightarrow H=\amalg_i H_i$ be the canonical monomorphisms,
and let $f=(f_i,i\in I):H=\amalg_i H_i\rightarrow K$ be the morphism
defined by the $f_i$. The family of the $s_i$ is covering for $T'$.
Thus, by axiom (T~2) for topologies, to show that $(f_i,i\in I)$ is a
covering family for $T'$, it suffices to show that the morphism
$f:H\rightarrow K$ is covering for $T'$. There exists an epimorphic
family in $C^{\widehat{\phantom{,}}}$,
$u_{\lambda}:X_{\lambda}\rightarrow K$, $\lambda\in \Lambda$, with
$X_{\lambda}\in \ob C$ (\hyperrefext[I][I.prop3.4]{3.4}). Therefore,
again by axiom (T 2) for topologies, to show that $f:H\rightarrow K$
is covering for $T'$, it suffices to show that, for every
$\lambda\in \Lambda$, the morphism
$\pr_2:H\times_K X_{\lambda}\rightarrow X_{\lambda}$ is covering
for $T'$. Let $v_j:Y_j\rightarrow H\times_K X_{\lambda}$, $j\in J$,
with $Y_j\in \ob C$, be an epimorphic family in
$C^{\widehat{\phantom{,}}}$. The family $(\pr_2\circ v_j,j\in J)$ is
covering for $T_C$. It is therefore a covering family in $C$
(\hyperref[II.prop5.1]{5.1}), and hence a covering family for $T'$.
Consequently, the sieve generated by
$\pr_2:H\times_K X_{\lambda}\rightarrow X_{\lambda}$ contains a
covering sieve for $T'$. It is therefore covering for $T'$, and hence
$\pr_2:H\times_K X_{\lambda}\rightarrow X_{\lambda}$ is covering,
\cqfd
\end{enonce*}

\begin{enonce*}[remark]{Remark 5.4.1}\etiquette[5.4.1]{II.rem5.4.1}
The proof of \hyperref[II.prop5.4]{5.4} in fact shows that every
topology $T'$ on $C^{\widehat{\phantom{,}}}$ which is finer than the
canonical topology of $C^{\widehat{\phantom{,}}}$ is the coarsest of
the topologies $T$ on $C^{\widehat{\phantom{,}}}$ having the
following properties:
\begin{enumerate}
\renewcommand{\theenumi}{\alph{enumi}}
\renewcommand{\labelenumi}{\theenumi)}
\item $T$ is finer than the canonical topology of
  $C^{\widehat{\phantom{,}}}$.

\item Every family covering for $T'$ of the form
  $u_i:X_i\rightarrow X$, $i\in I$, where $X$ and the $X_i$ are
  objects of $C$, is covering for $T$.
\end{enumerate}
\end{enonce*}

In particular, every topology $T'$ on
$C^{\widehat{\phantom{,}}}$ which is finer than the canonical
topology is uniquely determined by the families
$u_i:X_i\rightarrow X$, $i\in I$, with $X_i$ and $X$ objects of $C$,
which are covering for $T'$.

\begin{enonce*}[remark]{Remark 5.4.2}\etiquette[5.4.2]{II.rem5.4.2}
It\pageoriginale is easy to show that, for every topology $T'$ on
$C^{\widehat{\phantom{,}}}$ which is finer than the canonical
topology, the set of families of morphisms $(X_i\rightarrow X)$,
$i\in I$, with the same target in $C$, which are covering for $T'$ is
the set of covering families of a topology on $C$. Thus
\hyperref[II.prop5.4]{5.4} and \hyperref[II.rem5.4.1]{5.4.1}
establish a one-to-one correspondence between the topologies on $C$
and the topologies on $C^{\widehat{\phantom{,}}}$ finer than the
canonical topology.
\end{enonce*}

\setcounter{subsection}{5} \setcounter{subsubsection}{-1}
\subsubsection{}\etiquette[5.5.0]{II.subsubsec5.5.0}
Let $C$ be a small category. Denote by $\Caf$ the set of strictly full
subcategories (every object isomorphic to an object of the
subcategory is an object of the subcategory) of
$C^{\widehat{\phantom{,}}}$ whose inclusion functor admits a left
adjoint which commutes with finite projective limits. Also denote by
$\Tt$ the set of topologies on $C$. Theorem
\hyperref[II.theo3.4]{3.4} defines a map
$$
\Phi:\Tt\longrightarrow \Caf{quoi}.
$$

We shall define a map in the opposite direction. For this purpose,
to every element
$\xymatrix{ e=i': C' \ar@<-4pt>[r]&
\ar@<-2pt>[l]_(.35){a'} C^{\widehat{\phantom{,}}} }$ of $\Caf$, with
$\underline{a}'$ left adjoint to $i'$, we must associate a topology
$T_e$ on $C$. For every object $X$ of $C$, define $J_e(X)$ to be the
set of subobjects of $X$ whose inclusion morphism is carried by
$\underline{a}'$ to an isomorphism. Using the hypotheses made on
$\underline{a}'$, one verifies immediately that this defines a
topology $T_e$ on $C$. We have therefore defined a map
$$
\Psi:\Caf\longrightarrow \Tt{quoi}.
$$

We then have the following result (due to J. GIRAUD):

\begin{enonce*}{Theorem 5.5}\etiquette[5.5]{II.theo5.5}
The\pageoriginale map $\Phi$ is a bijection, and $\Psi$ is its
inverse.
\end{enonce*}

\begin{enonce*}[remark]{Proof}
The map $\Psi\circ \Phi$ is the identity. Indeed, this follows
immediately from \hyperref[II.prop5.1]{5.1}. The map
$\Phi\circ \Psi$ is the identity. Indeed, let
$\xymatrix{ i': C' \ar@<-4pt>[r]&
\ar@<-2pt>[l]_(.35){{\underline{a}'}}
C^{\widehat{\phantom{,}}} }$ be an element of $\Caf$, let $T_e$ be
the topology corresponding to it under $\Psi$, and let $C_e$ be the
category of sheaves for $T_e$. Using the definition of the topology
$T_e$ and the definition of bicovering morphisms
(\hyperref[II.def5.2]{5.2}), one proves that the following assertions
are equivalent:
\begin{enumerate}
\renewcommand{\theenumi}{\roman{enumi}}
\renewcommand{\labelenumi}{\theenumi)}
\item The morphism $u$ of $C^{\widehat{\phantom{,}}}$ is bicovering
  for the topology $T_e$.

\item The morphism $u$ of $C^{\widehat{\phantom{,}}}$ is carried by
  $\underline{a}'$ to an isomorphism.
\end{enumerate}

It is clear that $C'$ is a full subcategory of $C_e$. It therefore
suffices to show that every sheaf $F$ for the topology $T_e$ is an
object of $C'$. But, by the equivalence above, the morphism
$F\rightarrow i'\circ a'(F)$ is bicovering, and since its source and
target are sheaves, it follows from \hyperref[II.prop5.3]{5.3} (ii)
that $F\rightarrow i'\circ a'(F)$ is an isomorphism, \cqfd
\end{enonce*}


% END INCLUDED SGA4_English_current_buildable_source_20260906/source/02/drafts/SGA4_II_LINES_1646_1819_English.texfrag


% BEGIN INCLUDED SGA4_English_current_buildable_source_20260906/source/02/drafts/SGA4_II_LINES_1820_2120_English.texfrag
\section{Sheaves with values in a category}\etiquette[6]{II.sec6}

\setcounter{subsection}{-1}
\subsection{}\etiquette[6.0]{II.subsec6.0}
Let $C$ and $D$ be two categories. A contravariant functor from
$C$ to $D$, $F:C^{\circ}\rightarrow D$, is called a
\emph{presheaf on $C$ with values in $D$.}

\begin{enonce*}{Definition 6.1}\etiquette[6.1]{II.def6.1}
Let $C$ be a site and $D$ a category. A presheaf
$F:C^{\circ}\rightarrow D$ is called a \emph{sheaf on $C$ with
values in} $D$ (or, more briefly, a sheaf with values in $D$) if, for
every object $S$ of $D$, the presheaf of sets
$$
X\longmapsto \Hom_D(S,F(X))\quav X\in \ob(C)
$$
is a sheaf. The full subcategory of $\Hhom(C^{\circ},D)$ formed by
the sheaves on $C$ with values in $D$ will be denoted by
$\Hhom^{\sim}(C^{\circ},D)$.
\end{enonce*}

\begin{enonce*}[remark]{Remarks 6.2}\etiquette[6.2]{II.rem6.2}
\begin{enumerate}
\renewcommand{\labelenumi}{\theenumi)}
\hypertarget{sga4.exp.II.remark.6.2.item.1}{}\item The\pageoriginale condition of \hyperref[II.def6.1]{6.1} means
  that for every object $X$ of $C$, every covering sieve $R$ of $X$,
  and every object $S$ of $D$, one has an isomorphism
  (\hyperrefext[I][I.cor3.5]{3.5} and \hyperref[II.def2.1]{2.1}):
$$
\Hom_D(S,F(X))\xrightarrow{\sim}\varprojlim_{C/R}F(.){\quoi}.
$$
Thus, when $D$ is the category of $𝒰$-sets, one recovers the
definition of sheaves of sets (\loccit).

\item Let $D'$ be a $𝒰$-category and let $G:D\rightarrow D'$ be a
  functor commuting with projective limits. By composition, the
  functor $G$ carries sheaves with values in $D$ to sheaves with
  values in $D'$.
\end{enumerate}
\end{enonce*}

\setcounter{subsection}{3} \setcounter{subsubsection}{-1}
\subsubsection{}\etiquette[6.3.0]{II.subsubsec6.3.0}
Let $\gamma$ be a type of algebraic structure defined by finite
projective limits (\hyperrefext[I][I.cor2.9]{2.9}), and let
$𝒰-\gamma$ be the category of $\gamma$-sets belonging to $𝒰$.
Denote by $\esj:𝒰-\gamma\rightarrow 𝒰-\Ens$ the ``underlying
set'' functor. The functor $\esj$ commutes with projective limits and
therefore, by the preceding remark, defines by composition the
functors
\begin{align*}
&\esj^{\widehat{\phantom{,}}} :
\Hhom(C^{\circ},𝒰-\gamma)\longrightarrow
C^{\widehat{\phantom{,}}}\\
&\esj^{\widetilde{\phantom{,}}} :
\Hhom^{\sim}(C^{\circ},𝒰-\gamma) \longrightarrow
C^{\widetilde{\phantom{,}}}{\quoi}.
\end{align*}

The functor $\esj^{\widetilde{\phantom{,}}}$ is called the
``underlying sheaf of sets'' functor. In fact, it factors canonically
through the category $\gamma-\widetilde{C}$ of $\gamma$-objects of
$C^{\widetilde{\phantom{,}}}$. Denoting again by
$\widetilde{\esj}$ the resulting functor
$$
\Hhom^{\sim}(C^{\circ}, 𝒰-\gamma)\longrightarrow
\gamma-\widetilde{C},
$$
we may state:

\begin{enonce*}{Proposition 6.3.1}\etiquette[6.3.1]{II.prop6.3.1}
The functor $\esj^{\widetilde{\phantom{,}}}$ establishes an
equivalence among the category of sheaves on $C$ with values in
$𝒰-\gamma$, the category of $\gamma$-objects in the category of
sheaves of sets, and the category of $\gamma$-objects in the category
of presheaves of sets whose underlying presheaf is a sheaf.
\end{enonce*}

\begin{enonce*}[remark]{Proof}
The\pageoriginale proof essentially uses
(\hyperrefext[I][I.cor3.2]{3.2}) and the fact that a finite projective
limit of sheaves in the category of presheaves is a sheaf
(\hyperref[II.theo4.1]{4.1}).
\end{enonce*}

\setcounter{subsection}{3} \setcounter{subsubsection}{1}
\subsubsection{}\etiquette[6.3.2]{II.subsubsec6.3.2}
The preceding proposition justifies the abuse of language whereby a
sheaf with values in $𝒰-\gamma$ is identified with the corresponding
$\gamma$-object of $C^{\widetilde{\phantom{,}}}$. We shall
henceforth make this abuse of language systematically.

We shall use the classical terminology: \emph{sheaves of groups,
sheaves of commutative groups} (which we shall most often call
\emph{abelian sheaves}), \emph{sheaves of rings, sheaves of
$A$-modules} \ldotsn

\subsubsection{}\etiquette[6.3.3]{II.subsubsec6.3.3}
We denote by
$$
C^{\widetilde{\phantom{,}}}_{\gamma}(\resp
C^{\widehat{\phantom{,}}}_{\gamma})
$$
the category of $\gamma$-objects of $C^{\widetilde{\phantom{,}}}$
(\resp of $C^{\widehat{\phantom{,}}}$). When $\gamma$ is the type of
structure ``$A$-module'', we also write
$C^{\widetilde{\phantom{,}}}_{A}$, or simply
$\widetilde{C}_{\ab}$ when $A=\ZZ$.

Suppose that $C$ is a $𝒰$-site. The ``associated sheaf'' functor is
left exact (\hyperref[II.theo4.1]{4.1}), and therefore:

\begin{enonce*}{Proposition 6.4}\etiquette[6.4]{II.prop6.4}
The inclusion functor
$i_{\gamma}:C^{\widetilde{\phantom{,}}}_{\gamma}\rightarrow
C^{\widehat{\phantom{,}}}_{\gamma}$ admits a left adjoint
$a_{\gamma}$ which is left exact; that is, one has an isomorphism
$$
\Hom_{C^{\widetilde{\phantom{,}}}_{\gamma}}(a_{\gamma},\;)\xrightarrow{\sim}\Hom_{C^{\widehat{\phantom{,}}}_{\gamma}}(\;,i_{\gamma}){\quoi}.
$$

Let $X$ be an object of $C^{\widehat{\phantom{,}}}_{\gamma}$. The
underlying sheaf of sets of $a_{\gamma}(X)$ is canonically isomorphic
to the sheaf of sets associated with the underlying presheaf of sets
of $X$. The morphism of underlying presheaves of sets of the adjunction
morphism $\id\rightarrow i_{\gamma}a_{\gamma}$ is identified with the
adjunction morphism applied to the underlying presheaf of sets.
\end{enonce*}

\begin{enonce*}{Proposition 6.5}\etiquette[6.5]{II.prop6.5}
Suppose\pageoriginale that the functor
$\esj:\gamma-𝒰\rightarrow 𝒰-\Ens$ admits a left adjoint
\footnote[1]{One proves that such an adjoint always exists [C.F.]:
  free group, free abelian group, free $A$-module \ldotsn}:
$$
\Hom_{𝒰-\Ens}(.,\esj(.))\xrightarrow{\sim}\Hom_{\gamma-𝒰}(\Lib(.),.){\quoi}.
$$

The functor
$\esj^{\widetilde{\phantom{,}}}:C^{\widetilde{\phantom{,}}}_{\gamma}\rightarrow
C^{\widetilde{\phantom{,}}}$ (\resp
$\esj^{\widehat{\phantom{,}}}:C^{\widehat{\phantom{,}}}_{\gamma}\rightarrow
C^{\widehat{\phantom{,}}}$) admits a left adjoint
$\Lib^{\widetilde{\phantom{,}}}$ (\resp
$\Lib^{\widehat{\phantom{,}}}$), and one has a canonical isomorphism
$$
\Lib^{\widetilde{\phantom{,}}}=a_{\gamma}\Lib^{\widehat{\phantom{,}}}
i{\quoi}.
$$
\end{enonce*}

\begin{enonce*}[remark]{Proof}
The proof is formal and is left to the reader.
\end{enonce*}

\begin{enonce*}{Corollary 6.6}\etiquette[6.6]{II.cor6.6}
Let $(X_i)$, $i\in I$, be a generating family of
$C^{\widetilde{\phantom{,}}}$ (\hyperref[II.subsec4.9]{4.9}). Then
the family $\Lib^{\widetilde{\phantom{,}}}(X_i)$ is a generating
family of $C^{\widetilde{\phantom{,}}}_{\gamma}$.
\end{enonce*}

\begin{enonce*}[remark]{Proof}
The proof is formal once one observes that the functor
$\esj^{\widetilde{\phantom{,}}}$ commutes with projective limits and
is conservative ($\esj^{\widetilde{\phantom{,}}}(u)$ is an
isomorphism $\Leftrightarrow$ $u$ is an isomorphism).
\end{enonce*}

\begin{enonce*}{Proposition 6.7}\etiquette[6.7]{II.prop6.7}
Let $A$ be a $𝒰$-sheaf of rings, or a small ring, and let
$C^{\widetilde{\phantom{,}}}_A$ (\resp
$C^{\widehat{\phantom{,}}}_A$) be the category of sheaves (\resp of
presheaves) of unital $A$-modules
$(\hyperref[II.subsubsec6.3.3]{6.3.3})$ on a $𝒰$-site $C$. Then
$C^{\widetilde{\phantom{,}}}_A$ is an abelian $𝒰$-category satisfying
the axioms {\rm (AB 5)} (``existence of left exact filtered inductive
limits'') and {\rm (AB 3)$^{\ast}$} (``existence of infinite
products'') of {\rm[\textsc{Tohoku}]}. It has a family of generators
indexed by an element of $𝒰$.
\end{enonce*}

\begin{enonce*}[remark]{Proof}
It is clear that the category $C^{\widehat{\phantom{,}}}_A$ is an
abelian category satisfying the axioms (AB 3), (AB 4), and (AB 5)
(\hyperrefext[I][I.prop2.8]{2.8} and
\hyperrefext[I][I.cor3.3]{3.3}). It then follows\pageoriginale from
\hyperref[II.prop6.4]{6.4} that the category
$C^{\widetilde{\phantom{,}}}_A$ is an additive category in which
kernels and cokernels are representable. More precisely, let $i_A$
and $a_A$ be the inclusion and associated-sheaf functors, respectively,
and let $u:X\rightarrow Y$ be a morphism of
$C^{\widetilde{\phantom{,}}}_A$. The functor $a_A$ is left exact and
commutes with inductive limits. The functor $i_A$ commutes with
projective limits. Consequently the functor
$i_Aa_A:C^{\widehat{\phantom{,}}}_A\rightarrow
C^{\widehat{\phantom{,}}}_A$ is an additive left exact functor. One
deduces canonical isomorphisms
\begin{gather*}
i_A(\Ker(u))\xrightarrow{\sim} \Ker(i_A(u)){\quoi},\tag{$\ast$}\\
\coker(u)\xrightarrow{\sim} a_A \coker
(i_A(u)){\quoi}.\tag{${\ast}{\ast}$}
\end{gather*}
Consequently there is a canonical isomorphism
\begin{equation}\coim(u)\xrightarrow{\sim} a_A \coim (i_A (u)){\quoi}.
\tag{${\ast}{\ast}{\ast}$}
\end{equation}
Let us show that the canonical morphism
$\coim(u)\rightarrow \im(u)$ is an isomorphism. For this, by (*), it
suffices to show that the sequence
$$
0\longrightarrow i_A(\coim(u))\longrightarrow i_A(Y)\longrightarrow
i_A(\coker(u))
$$
is exact. Using the isomorphisms $({\ast}{\ast})$ and
$({\ast}{\ast}{\ast})$, and observing that $a_Ai_A(Y)$ is isomorphic
to $Y$, one sees that this sequence is isomorphic to the sequence
$$
0\longrightarrow i_Aa_A(\coim(i_A(u)))\longrightarrow i_Aa_Ai_A
(Y)\longrightarrow i_Aa_A(\coker(i_A(u))){\quoi},
$$
which is precisely the image under the functor $i_Aa_A$ of the
sequence
$$
0\longrightarrow \coim(i_A(u))\longrightarrow i_A(Y)\longrightarrow
\coker(i_A(u)){\quoi}.
$$
This sequence is exact, and the functor $i_Aa_A$ is left exact. Thus
the category $C^{\widetilde{\phantom{,}}}_A$ is an abelian category.
The functor $a_A:C^{\widehat{\phantom{,}}}_A\rightarrow
C^{\widetilde{\phantom{,}}}_A$\pageoriginale is left exact and
commutes with arbitrary inductive limits. Consequently it is
additive, exact, and commutes with inductive limits. Since the category
$C^{\widehat{\phantom{,}}}_A$ satisfies axiom (AB 5), the same is true
of the category $C^{\widetilde{\phantom{,}}}_A$. Finally, it is clear
that in the category $C^{\widetilde{\phantom{,}}}_A$, products indexed
by an element of $𝒰$ are representable, and hence axiom
(AB 3)$^{\ast}$ is satisfied. This completes the proof of the first
assertion.
\end{enonce*}

For every ring $A$ which is an element of $𝒰$, the functor
$$
\esj:𝒰-A-\mbox{module}\longrightarrow 𝒰-\Ens
$$
admits a left adjoint $X\longmapsto X_A$ (the free $A$-module
generated by $X$). It
therefore suffices to apply \hyperref[II.prop4.10]{4.10} and
\hyperref[II.cor6.6]{6.6} to complete the proof.

\begin{enonce*}[remark]{Notation 6.8}\etiquette[6.8]{II.not6.8}
Let $H$ be a sheaf of sets. The sheaf of $A$-modules associated with
the presheaf (\cf notation \hyperref[II.prop6.5]{6.5})
$$
X\longmapsto \Lib_A H(X)\qquad X\in \ob C
$$
is denoted by $A_H$. When $H=\underline{a}(X)$ (the sheaf associated
with the presheaf represented by $X$), one sometimes simply writes
$A_X$ (by abuse of notation). The proof of \hyperref[II.prop6.7]{6.7}
shows that the family of sheaves of $A$-modules $A_X$,
$(X\in \ob(C))$, is a generating family, indexed by an element of
$𝒰$, of the category of sheaves of $A$-modules.
\end{enonce*}

\begin{enonce*}[remark]{Remark 6.9}\etiquette[6.9]{II.rem6.9}
According to [\textsc{Tôhoku}], \hyperref[II.prop6.7]{6.7} shows that
the category $C^{\widetilde{\phantom{,}}}_A$, when $A$ is an element
of $𝒰$, has enough injectives. It is also known that infinite products
are not necessarily exact in $C^{\widetilde{\phantom{,}}}_A$, and
therefore that, in general, the category
$C^{\widetilde{\phantom{,}}}_A$ does not have enough projectives
[\textsc{Roos}].
\end{enonce*}

\begin{thebibliography}{99}
\bibitem{sga4.exp.II.bib.1} M. \textsc{Artin};\pageoriginale Grothendieck's topologies.

\bibitem{sga4.exp.II.bib.2} M. \textsc{Demazure}: Séminaire de Géométrie Algébrique III,
  Exposé IV

Lecture Notes. Springer-Verlag.

\bibitem{sga4.exp.II.bib.3} A. \textsc{Grothendieck}: Sur quelques points d'Algèbre
Homologique.

 Tohoku Math. Journal.

\bibitem{sga4.exp.II.bib.4} J.E. \textsc{Roos}: CR
\end{thebibliography}

\ifx\danslelivre\undefined
%\bibliography{../fichiers-configuration-latex/bibliographie}
%\bibliographystyle{../fichiers-configuration-latex/smfalpha}
\end{document}
\fi

% END INCLUDED SGA4_English_current_buildable_source_20260906/source/02/drafts/SGA4_II_LINES_1820_2120_English.texfrag


% END INCLUDED SGA4_English_current_buildable_source_20260906/source/02/02_English.tex


\makeatletter
\def\input@path{{../03/}}
\makeatother
% BEGIN INCLUDED SGA4_English_current_buildable_source_20260906/source/03/03_English.tex
\ifx\danslelivre\undefined
% Release engine: XeLaTeX (the LuaTeX compatibility layer breaks Xy-pic output).
\documentclass[10pt]{../fichiers-configuration-latex/smfart}
% BEGIN INCLUDED SGA4_English_current_buildable_source_20260906/source/fichiers-configuration-latex/paquets.tex
\usepackage{polyglossia}
%\usepackage[francais]{babel}
\usepackage{fontspec}
\usepackage[math-style=ISO,bold-style=upright,sans-style=upright]{unicode-math}
%\usepackage[mathletters]{ucs}
%\usepackage[utf8]{inputenc}
%\usepackage[T1]{fontenc}
\usepackage[unicode,hypertexnames=false,colorlinks=true,linkcolor=blue,citecolor=Sepia,urlcolor=magenta]{hyperref}
\setmainlanguage{french}
\setmainfont[Ligatures=TeX]{Linux Libertine O}
% variante possible : TeXGyrePagella ?
\setmathfont{XITS Math}

% maths
\usepackage{amsmath}
%\usepackage{amssymb}
%\usepackage{mathrsfs}
%\usepackage{stmaryrd}
%\usepackage{euscript}

% couleurs
\usepackage[usenames,dvipsnames]{xcolor}

% graphiques
\usepackage{graphicx}
\usepackage[all]{xy}
%\usepackage{tikz}
%\usetikzlibrary{matrix,arrows,calc}
%\tikzset{column sep=3em,row sep=2em,text height=1.5ex, text depth=0.25ex,>=angle 90}
%\tikzset{description/.style={fill=white,inner sep=1pt}}

% hyperliens, 'etiquettes
\usepackage{xr-hyper}
\usepackage{../fichiers-configuration-latex/smfthm}
%\ifx\cacheretiquettes\undefined
%\usepackage[inline]{showlabels} \else \fi % pour rendre visibles les 'etiquettes
%\usepackage{../fichiers-configuration-latex/makeidx} % pour faire un index
%\usepackage[pagebackref,unicode]{hyperref}  % pour voir o`u sont cit'ees les r'ef'erences
% hyperref is loaded once above with the cumulative-reader destination policy.
%\hypersetup{breaklinks=true} %permet le retour `a la ligne dans les liens trop longs. Incompatible avec pdflatex
%\renewcommand*{\backref}[1]{{\small [cit\'e page(s)~#1.]}}

\usepackage{xspace}
\usepackage{pifont} % pour expose XVII

% END INCLUDED SGA4_English_current_buildable_source_20260906/source/fichiers-configuration-latex/paquets.tex

% BEGIN INCLUDED SGA4_English_current_buildable_source_20260906/source/fichiers-configuration-latex/numerotation.tex
\renewcommand{\thefootnote}{\roman{footnote}} % pour mettre les notes en bas de page sous la forme (ii) et non (2), qui ressemble trop (?) a un exposant.

\newcounter{pageoriginal}
\newcommand{\pageoriginale}{\refstepcounter{pageoriginal}\marginpar{\footnotesize\textbf{\thepageoriginal}}\xspace}
\let\pageoriginaled\pageoriginale

\usepackage{ifthen}
\newboolean{original}
\def\sisi{\ifthenelse{\boolean{original}}}
\setboolean{original}{false}

\newcounter{cnote}
\def\nde#1{\ifthenelse{\boolean{original}}{}{\stepcounter{cnote}\def\thefootnote{\arabic{cnote}}\protect\footnote[\value{cnote}]{N.D.E.: #1}}}
\def\ndemark{\ifthenelse{\boolean{original}}{}{\stepcounter{cnote}\def\thefootnote{\arabic{cnote}}\protect\footnotemark[\value{cnote}]}}
\def\ndetext#1{\ifthenelse{\boolean{original}}{}{\def\thefootnote{\arabic{cnote}}\protect\footnotetext[\value{cnote}]{N.D.E.: #1}}}

% END INCLUDED SGA4_English_current_buildable_source_20260906/source/fichiers-configuration-latex/numerotation.tex

% BEGIN INCLUDED SGA4_English_current_buildable_source_20260906/source/fichiers-configuration-latex/macros.tex
\setmainlanguage{french}
\setotherlanguage{english}


%% Fixed, engine-independent release identifier.  The inherited dynamic Git
%% probe used Lua io.popen, exposed build-state-dependent output, and printed
%% a visible question mark outside a Git checkout.
\def\showgitstatus{English release r2, 2026-07-23}


% macros pour LaTeX
\renewcommand{\mathsf}{}

% Police différente pour U+2113 SCRIPT SMALL L
\setmathfont[range="2113]{latinmodern-math.otf}

% Give every inherited semantic label its own stable PDF destination instead
% of inheriting whichever automatic equation/page anchor happened to be live.
% The legacy label remains as a compatibility alias used by existing edges.
\makeatletter
\def\etiquette[#1]#2{%
  \Hy@raisedlink{\hyper@anchorstart{sga4.semantic.#2}\hyper@anchorend}%
  \def\@currentHref{sga4.semantic.#2}%
  \def\@currentlabel{#1}%
  \label{#2}%
}
\makeatother
\def\eqhyperref[#1]#2{(\hyperref[#1]{#2})}
\def\pagehyperref[#1]#2{(page \hyperref[#1]{#2})}
\def\hyperrefext[#1][#2]#3{\hyperref[#2]{#1 #3}}
\def\eqhyperrefext[#1][#2]#3{\hyperref[#2]{#1 #3}}     %cf. exposes 17
\def\neweqhyperrefext[#1][#2]#3{\hyperref[#2]{#1 #3}}  %cf. exposes 17

\def\shat#1{\protect\widehat{#1}}
\def\stilde#1{\protect\widetilde{#1}}
%%%%%%%%%%%%%%%%%%


\def\etcn{etc.\xspace}
\def\etco{etc...\xspace}
\def\etc{\sisi{\etco}{\etcn}}
\def\ldotsn{\sisi{...}{\etcn}}

\def\cf{{cf.}\kern.3em}
\def\Cf{{Cf.}\kern.3em}
\def\ie{{i.e.}\kern.3em}
\def\iev{{i.e.},\xspace}
\def\loccito{loc.\kern2pt cit.\xspace}
\def\loccitn{\textit{loc.\kern2pt cit.}\xspace}
\def\loccit{\sisi{\loccito}{\loccitn}}
\def\EGA{\textup{EGA}\kern.2em}
\def\SGA{\textup{SGA}\kern.3em}

\def\quoi{}

\def\quav{\quad,\quad}
\def\quaet{\quad \text{et}\quad}
\def\fibreproduct#1#2#3{#1 ×_{#3} #2}
\let\fprod\fibreproduct
\def\ccheck#1#2{\mbox{${\check{\check{\rm #1}}}^{\lower2pt\hbox{$\scriptstyle #2$}}$}}
\def\sheaf#1{\protect\mathcal{#1}}
\def\vieuxsheaf#1{{#1}} % cf. problème avec sheaf{a\Phi_i} par exemple...
\DeclareMathOperator{\Hom}{Hom}
\DeclareMathOperator{\can}{can}
\DeclareMathOperator{\Ind}{Ind}
\DeclareMathOperator{\ind}{ind}
%\DeclareMathOperator{\Hhomtop}{\mathcal{H}\!\mathit{omtop}}
\DeclareMathOperator{\Hhomtop}{\mathbf{Homtop}}
\DeclareMathOperator{\Point}{Point}
\DeclareMathOperator{\Pro}{Pro}
\DeclareMathOperator{\Top}{Top}
\DeclareMathOperator{\Mod}{Mod}
\DeclareMathOperator{\Cart}{Cart}
%\DeclareMathOperator{\Mmorsite}{\mathcal{M}\!\mathit{orsite}}
\DeclareMathOperator{\Mmorsite}{\mathbf{Morsite}}
\DeclareMathOperator{\loc}{loc}
\DeclareMathOperator{\Loc}{Loc}
\DeclareMathOperator{\Pic}{Pic}
\DeclareMathOperator{\Et}{Et}
\DeclareMathOperator{\card}{card}
\DeclareMathOperator{\Ens}{Ens}
\DeclareMathOperator{\h}{\underline{H}}
\DeclareMathOperator{\Fonct}{Fonct}
\DeclareMathOperator{\Ker}{Ker}
\DeclareMathOperator{\but}{but}
\DeclareMathOperator{\Ob}{Ob}
\DeclareMathOperator{\ob}{ob}
\DeclareMathOperator{\Lib}{Lib}
\DeclareMathOperator{\esj}{esj}
\DeclareMathOperator{\id}{id}
\DeclareMathOperator{\Id}{Id}
\DeclareMathOperator{\ad}{ad}
\DeclareMathOperator{\Iim}{Im}
\DeclareMathOperator{\pr}{pr}
\DeclareMathOperator{\Cat}{Cat}
\DeclareMathOperator{\car}{car}
\DeclareMathOperator{\SysInd}{SysInd}
\DeclareMathOperator{\Sup}{Sup}
\DeclareMathOperator{\Coker}{Coker}
\DeclareMathOperator{\Filt}{Filt}
\DeclareMathOperator{\acc}{acc}
\DeclareMathOperator{\sur}{sur}
\DeclareMathOperator{\cat}{cat}
\DeclareMathOperator{\Ord}{Ord}
\DeclareMathOperator{\Gr}{Gr}
\DeclareMathOperator{\Cov}{Cov}
\DeclareMathOperator{\Caf}{Caf}
\DeclareMathOperator{\coker}{coker}
\DeclareMathOperator{\coim}{coim}
\DeclareMathOperator{\im}{im}
\DeclareMathOperator{\Ouv}{Ouv}
\DeclareMathOperator{\diag}{diag}
\DeclareMathOperator{\Esp}{Esp}
\DeclareMathOperator{\TOP}{TOP}
\DeclareMathOperator{\Unif}{Unif}
\DeclareMathOperator{\Iind}{\mathbf{Ind}}
\DeclareMathOperator{\Hhomcart}{\mathbf{Homcart}}
\DeclareMathOperator{\Hhomtopess}{\mathbf{Homtopess}}
\DeclareMathOperator{\Ppoints}{\mathbf{Points}}
\DeclareMathOperator{\Ppoint}{\mathbf{Point}}
\DeclareMathOperator{\Ttors}{\mathbf{Tors}}
\DeclareMathOperator{\Ssex}{\mathbf{Sex}}
\DeclareMathOperator{\Ffib}{\mathbf{Fib}}
\DeclareMathOperator{\Vvois}{\mathbf{Vois}}
\DeclareMathOperator{\Ppointess}{\mathbf{Pointess}}
\DeclareMathOperator{\scrHom}{\mathbf{Hom}}
\DeclareMathOperator{\sheafPic}{\mathbf{Pic}}
\DeclareMathOperator{\sheafAut}{\mathbf{Aut}}
\DeclareMathOperator{\Hhomtoplocan}{\mathbf{Homtoplocan}}
\DeclareMathOperator{\Hhomtopan}{\mathbf{Homtopan}}
\DeclareMathOperator{\Pptgeom}{\mathbf{Ptgéom}}
\DeclareMathOperator{\Morsite}{Morsite}
\DeclareMathOperator{\Fib}{Fib}
\DeclareMathOperator{\Fais}{Fais}
\DeclareMathOperator{\kar}{kar}
\DeclareMathOperator{\Groupes}{Groupes}
\DeclareMathOperator{\Restr}{Restr}
\DeclareMathOperator{\Isom}{Isom}
\DeclareMathOperator{\Prol}{Prol}
\DeclareMathOperator{\Res}{Res}
\DeclareMathOperator{\Tors}{Tors}
\DeclareMathOperator{\esp}{esp}
\DeclareMathOperator{\es}{es}
\DeclareMathOperator{\od}{od}
\DeclareMathOperator{\gr}{gr}
\DeclareMathOperator{\sob}{sob}
\DeclareMathOperator{\dfn}{dfn}
\DeclareMathOperator{\Sch}{Sch}
\DeclareMathOperator{\ouv}{ouv}
\DeclareMathOperator{\cons}{cons}
\DeclareMathOperator{\Nerf}{Nerf}
\DeclareMathOperator{\Points}{Points}
\DeclareMathOperator{\Espsob}{Espsob}
\DeclareMathOperator{\Homtop}{Homtop}
\DeclareMathOperator{\Pt}{Pt}
\DeclareMathOperator{\Quot}{Quot}
\DeclareMathOperator{\glob}{glob}
\DeclareMathOperator{\Homfib}{Homfib}
\DeclareMathOperator{\Qouv}{Qouv}
\DeclareMathOperator{\Quotouv}{Quotouv}
\DeclareMathOperator{\ens}{ens}
\DeclareMathOperator{\ab}{ab}
\DeclareMathOperator{\an}{an}
\DeclareMathOperator{\fs}{fs}
\DeclareMathOperator{\red}{red}
\DeclareMathOperator{\Ab}{Ab}
\DeclareMathOperator{\cosq}{cosq}
\DeclareMathOperator{\Ann}{Ann}
\DeclareMathOperator{\Zar}{Zar}


%%%%%%%%%%%%
%%%%% SGA4-2
\DeclareMathOperator{\Sex}{Sex}
\DeclareMathOperator{\vois}{vois}
\DeclareMathOperator{\Vois}{Vois}
\DeclareMathOperator{\Ext}{Ext}
\DeclareMathOperator{\ext}{ext}
\DeclareMathOperator{\Fer}{Fer}
\DeclareMathOperator{\Add}{Add}
\DeclareMathOperator{\Cocart}{Cocart}
\DeclareMathOperator{\cosk}{cosk}
\DeclareMathOperator{\plat}{plat}

\DeclareMathOperator{\faisc}{faisc}
\DeclareMathOperator{\sk}{sk}
\DeclareMathOperator{\ch}{ch}

\def\HH{\mathbf{H}}
\def\qc{\mathrm{qc}}
\def\et{\mathrm{et}}
\def\coh{\mathrm{coh}}

\def\Hvr#1{%
\begin{array}[b]{l}
{\kern2pt\scriptstyle r}\\[-8pt]
{\kern1pt\scriptstyle \vee}\\[-5pt]
H^{#1}
\end{array}}

\def\Hvv#1{%
\begin{array}[b]{l}
\kern2pt\rotatebox{90}{${\scriptscriptstyle\ll}$}\\[-6pt]
H^{#1}
\end{array}}

\def\Hvi#1{%
\begin{array}[b]{l}
{\kern0.2pt\scriptscriptstyle \infty}\\[-8pt]
{\kern1pt\scriptstyle \vee}\\[-5pt]
H^{#1}
\end{array}}

\def\Hv#1{%
\begin{array}[b]{l}
{\kern1pt\scriptstyle \vee}\\[-5pt]
H^{#1}
\end{array}}

\DeclareMathOperator{\cart}{cart}
\DeclareMathOperator{\Ffaisc}{\mathbf{Faisc}}


\def\Ev#1{%
\begin{array}[b]{l}
{\kern1pt\scriptstyle \vee}\\[-5pt]
E'_{#1}
\end{array}}

\DeclareMathOperator{\Cont}{Cont}
\DeclareMathOperator{\Ccont}{\mathscr{C}\!\mathit{ont}}
\DeclareMathOperator{\Site}{Site}
\newcommand\MM{\mathbf{M}}
\def\limtop#1{{\displaystyle{\varprojlim_{#1}\!\mbox{\rm top }}}}
\def\Limtop{\mathop{\varprojlim\!\mbox{\rm top}}}
\DeclareMathOperator{\spec}{spec}
\def\ss{\scriptstyle{\mathscr{S}}}
\DeclareMathOperator{\Gal}{Gal}
\def\smallm{\scriptstyle{\mathscr{M}}}
\DeclareMathOperator{\End}{End}
\DeclareMathOperator{\ass}{ass}
\DeclareMathOperator{\zar}{zar}
\DeclareMathOperator{\Ogg}{Ogg}
\DeclareMathOperator{\res}{res}
\DeclareMathOperator{\tr}{tr}
\DeclareMathOperator{\Supp}{Supp}
\DeclareMathOperator{\codim}{codim}
\DeclareMathOperator{\rk}{rk}
\DeclareMathOperator{\cur}{cur}
\DeclareMathOperator{\cdqc}{cdqc}
\newcommand\DD{\mathbf{D}}
\DeclareMathOperator{\cl}{cl}
\def\JJ{\mathbf{J}}
\def\FF{\mathbf{F}}
\def\BB{\mathbf{B}}
\DeclareMathOperator{\Alget}{Alget}
\DeclareMathOperator{\Ass}{Ass}
\DeclareMathOperator{\supp}{supp}
\DeclareMathOperator{\rad}{rad}
\DeclareMathOperator{\prof}{prof}
\DeclareMathOperator{\cint}{int}
\DeclareMathOperator{\sic}{sic}
\DeclareMathOperator{\Tor}{Tor}
\DeclareMathOperator{\tors}{tors}
\DeclareMathOperator{\Sym}{Sym}
\DeclareMathOperator{\parf}{parf}
\DeclareMathOperator{\tor}{tor}
\DeclareMathOperator{\defn}{defn}
\DeclareMathOperator{\Div}{Div}
\DeclareMathOperator{\Pol}{Pol}
\DeclareMathOperator{\Fex}{Fex}
\DeclareMathOperator{\Lie}{Lie}
\DeclareMathOperator{\tot}{tot}
\DeclareMathOperator{\Aut}{Aut}
\DeclareMathOperator{\HOM}{HOM}
\DeclareMathOperator{\pch}{pch}
\DeclareMathOperator{\Ch}{Ch}
\DeclareMathOperator{\EXT}{EXT}
\DeclareMathOperator{\PIC}{PIC}
\DeclareMathOperator{\Rf}{Rf}
\DeclareMathOperator{\inj}{inj}
\newcommand\VV{\mathbf{V}}

\def\Xv{%
\begin{array}[b]{l}
{\kern1pt\scriptstyle \vee}\\[-5pt]
X
\end{array}}

%%%%%%%
\newcommand\TT{\mathbf{T}}
\newcommand\XX{\mathbf{X}}
\newcommand\EE{\mathbf{E}}
%%%%%%

\def\qedsymbol{\textup{C.Q.F.D.}}
\let\CQFD\qed
\let\cqfd\qed

%%%%%\bigdot
\font\bigsymb=cmsy10 at 4pt
\def\bigdot{{\kern1.2pt\raise 1.5pt\hbox{\bigsymb\char15}}}
%%%%%

\def\Lright{\mathop{\vcenter{\hbox{\hbox{$\longrightarrow$}\lower8pt \hbox{\kern-16pt\hbox{$\longrightarrow$}}}}}}

\def\Aa{{\mathcal A}}
\def\Bb{{\mathcal B}}
\def\Cc{{\mathcal C}}
\def\Dd{{\mathcal D}}
\def\Ee{{\mathcal E}}
\def\Ff{{\mathcal F}}
\def\Gg{{\mathcal G}}
\def\Hh{{\mathcal H}}
\def\Ii{{\mathcal I}}
\def\Jj{{\mathcal J}}
\def\Kk{{\mathcal K}}
\def\Ll{{\mathcal L}}
\def\Mm{{\mathcal M}}
\def\Nn{{\mathcal N}}
\def\Oo{{\mathcal O}}
\def\Pp{{\mathcal P}}
\def\Qq{{\mathcal Q}}
\def\Rr{{\mathcal R}}
\def\Ss{{\mathcal S}}
\def\Tt{{\mathcal T}}
\def\Uu{{\mathcal U}}
\def\Vv{{\mathcal V}}
\def\Ww{{\mathcal W}}
\def\Xx{{\mathcal X}}
\def\Yy{{\mathcal Y}}
\def\Zz{{\mathcal Z}}

\newcommand{\CC}{\mathbf{C}}
\newcommand{\GG}{\mathbf{G}}
\newcommand{\KK}{\mathbf{K}}
\newcommand{\LL}{\mathbf{L}}
\newcommand{\NN}{\mathbf{N}}
\newcommand{\PP}{\mathbf{P}}
\newcommand\QQ{\mathbf{Q}}
\newcommand\RR{\mathbf{R}}
\newcommand\ZZ{\mathbf{Z}}

\def\to{\mathchoice{\longrightarrow}{\rightarrow}{\rightarrow}{\rightarrow}}
\def\from{\mathchoice{\longleftarrow}{\leftarrow}{\leftarrow}{\leftarrow}}
\def\mto{\mathchoice{\longmapsto}{\mapsto}{\mapsto}{\mapsto}}
\def\mfrom{\mathrel{\from\joinrel\mapstochar\,}}
\def\hto{\mathrel{\lhook\joinrel\to}}
\def\hfrom{\mathrel{\from\joinrel\rhook}}

\def\implies{\mathchoice{\Longrightarrow}{\Rightarrow}{\Rightarrow}{\Rightarrow}}
\def\implied{\mathchoice{\Longleftarrow}{\Leftarrow}{\Leftarrow}{\Leftarrow}}
\def\iff{\mathchoice{\Longleftrightarrow}{\Leftrightarrow}{\Leftrightarrow}{\Leftrightarrow}}

\def\To#1{\mathchoice{\xrightarrow{\labelstyle~~#1~}}{\stackrel{#1}{\longrightarrow}}{}{}}
\def\Mto#1{\mathrel{\mapstochar\To{#1}}}
\def\From#1{\mathchoice{\xleftarrow{\labelstyle~~#1~}}{\stackrel{#1}{\longleftarrow}}{}{}}

\newcommand{\isomto}{\overset{\sim}{\longrightarrow}}
\newcommand{\isomfrom}{\overset{\sim}{\longleftarrow}}

\newcommand{\IFF}{\unskip~$\iff$~\ignorespaces}
\newcommand{\THEN}{\unskip~$\implies$~\ignorespaces}

\let\leq\leqslant
\let\geq\geqslant
\let\le\leqslant
\let\ge\geqslant

\def\cf{{cf.}\kern.3em}
\def\Cf{{Cf.}\kern.3em}
\def\ie{{i.e.}\kern.3em}
\def\iev{{i.e.},\xspace}
%\def\loccito{loc.\kern2pt cit.\xspace}
%\def\loccitn{\textit{loc.\kern2pt cit.}\xspace}
%\def\loccit{\sisi{\loccito}{\loccitn}}
\def\EGA{\textup{EGA}\kern.2em}
\def\SGA{\textup{SGA}\kern.3em}

\let\oldS\S
\def\S{\oldS\kern .15em}
\def\ptbl{.\kern .2em}
\def\resp{resp.\kern.3em}
\def\Exp{Exp.\kern.3em}

\let\bar\overline
\let\hat\widehat
\let\tilde\widetilde
\let\dpl\displaystyle
\let\ldots\dots
\let\emptyset\varnothing
\let\skpt\skippointrait
\let\epsilon\varepsilon

\def\sfrac#1#2{{#1}/{#2}}

\DeclareMathOperator{\cd}{cd}
%\DeclareMathOperator{\Hhom}{\mathcal{H}\!\mathit{om}}
\DeclareMathOperator{\Hhom}{\mathbf{Hom}}
%\DeclareMathOperator{\Eext}{\mathcal{E}\!\mathit{xt}}
\DeclareMathOperator{\Eext}{\mathbf{Ext}}
\DeclareMathOperator{\Spec}{Spec}
\DeclareMathOperator{\Tr}{Tr}

\let\emptyset=\varnothing

%%%%%%%%%%%%%%%
%%% expose XVII
%%%%%%%%%%%%%%%

\makeatletter
\let\old@optexpose\optexpose
\newcommand{\empty@optexpose}[1]{}
\newcommand{\paren@optexpose}[1]{\old@optexpose{#1}(}
\newcommand{\refnoexpose}[1]{\let\optexpose\empty@optexpose\ref{#1}\let\optexpose\old@optexpose}%\refnoexpose{X.cor5.2} engendre "5.2" au lieu de "X 5.2"
\newcommand{\neweqref}[1]{\let\optexpose\paren@optexpose\textup{\ref{#1})}\let\optexpose\old@optexpose}%\neweqref{XII.eq4.1} engendre "XII (4.1)" au lieu de "(XII 4.1)"
\makeatother
\newcommand{\refvb}[2][VI${}_{\text{B}}$ ]{\sisi{#1\refnoexpose{#2}}{\ref{#2}}}

\newcommand{\leftexp}[2]{{\vphantom{#2}}^{#1}{#2}}
\newcommand{\AAA}{\mathbf{A}}
\newcommand{\Sf}{\mathfrak{S}}
\DeclareMathOperator*{\quotprojlim}{\text{``$\varprojlim$''}}
\DeclareMathOperator*{\quotinjlim}{\text{``$\varinjlim$''}}
\DeclareMathOperator{\Fl}{F\sisi{\ell}{l}}
\DeclareMathOperator{\Pl}{P\sisi{\ell}{l}}
%\DeclareMathOperator{\cl}{c\sisi{\ell}{l}}
\DeclareMathOperator{\tdf}{tdf}
\DeclareMathOperator{\rg}{rg}
\newcommand{\iInd}{\sisi{I}{i}nd}
\newcommand{\iind}{ind}
\newcommand{\ppro}{pro}
\newcommand{\tTor}{\sisi{T}{t}or}
\newcommand{\ttor}{tor}
\newcommand{\chap}{Chap.}
\newcommand{\Ffex}{\Ff\mathit{ex}}
\newcommand{\sheafk}{\underline{k}}
\newcommand{\sheafn}{\underline{n}}
\newcommand{\sheafGamma}{\underline{\Gamma}}

% Repeated references to an already typeset footnote must not use
% \footnotemark without a matching \footnotetext: Hyperref would then create
% a live annotation to a destination that cannot exist.  This helper keeps the
% current printed footnote mark and links it to an explicit stable note target.
\makeatletter
\newcommand{\sgarepeatedfootnotemark}[1]{%
  \begingroup
  \edef\@thefnmark{\thefootnote}%
  \hyperlink{#1}{\@makefnmark}%
  \endgroup}
\makeatother

\sisi{
\DeclareMathOperator{\simarrow}{\rule{.15mm}{1.8mm}\kern -3.4pt\rightsquigarrow}
\newcommand{\sheafell}{\underline{\ell}}
\newcommand{\sheafZZ}{\underline{\ZZ}}
}{
\newcommand{\nderef}[1]{la N.D.E.~(\ref{#1}) page~\pageref{#1}}
}

%%%%%%%%%%%%%%%%%
%%%%% expose XVII
%%%%%%%%%%%%%%%%%

% END INCLUDED SGA4_English_current_buildable_source_20260906/source/fichiers-configuration-latex/macros.tex


\externaldocument{../01/01_English.tex}
\externaldocument{../02/02_English.tex}
%\externaldocument{../03/03_English.tex}
\externaldocument{../04/04.tex}

\setboolean{original}{false}

\begin{document}

\begin{center}
III. Functoriality of categories of sheaves\\
J.-L. Verdier
\end{center}

\selectlanguage{english}version: \showgitstatus
\tableofcontents

\else
\chapter{Functoriality of categories of sheaves}
%\addtocontents{toc}{J.-L. Verdier}
\begin{center}
J.-L. Verdier
\end{center}
\fi
\phantomsection\MakeLinkTarget*{sga4.exp.III}\label{sga4.exp.III}

% BEGIN INCLUDED SGA4_English_current_buildable_source_20260906/source/03/drafts/SGA4_III_LINES_31_241_English.texfrag
In\pageoriginale \hyperrefext[I][I.sec5]{5}, we studied the behavior of categories of
presheaves with respect to functors between the categories of
arguments. In this exposé, we extend this study to the case of sites
and categories of sheaves (${\rm n^{\circ}}$ 1 and 2). After
introducing the induced topology (${\rm n^{\circ}}$ 3), we turn in
${\rm n^{\circ}}$ 4 to the comparison lemma, which will play an
important role in Exp. IV. In number 5, for the patient reader, we
study certain commutative diagrams related to localized categories
(of the form $C/X$). In Theorem \hyperrefext[I][I.prop5.1]{5.1}, we impose
smallness hypotheses on the categories under consideration. In
general, these hypotheses are not satisfied by the sites encountered
in practice ($𝒰$-sites). This necessitates a few contortions
(\hyperref[III.subsec4.2]{4.2}, \hyperref[III.subsec4.3]{4.3},
\hyperref[III.subsec4.4]{4.4}).

\section{Continuous functors}\etiquette[1]{III.sec1}

\begin{enonce*}{Definition 1.1}\etiquette[1.1]{III.def1.1}
Let $C$ and $C'$ be two $𝒰$-sites and let $u:C\rightarrow C'$ be a
functor between their underlying categories. We say that $u$ is
continuous if, for every sheaf of sets $F$ on $C'$, the presheaf
$X\mapsto F\circ u(X)$ on $C$ is a sheaf.
\end{enonce*}

This notion of continuity depends a priori on the universe $𝒰$
with respect to which the two sites are $𝒰$-sites. Proposition
\hyperref[III.prop1.5]{1.5} shows that it is independent of it.

%\setcounter{subsection}{1}
\setcounter{subsection}{10}
\subsection{}\etiquette[1.11]{III.subsec1.11}
%% original misprint: 1.11 should be 1.1.1!? \XXX
Denote\pageoriginale by $i:C^{\sim}\rightarrow
C^{\widehat{\phantom{,}}}$ (\resp by $i':{C'}^{\sim}\rightarrow
{C'}^{\widehat{\phantom{,}}}$) the canonical inclusion functor from
sheaves into presheaves. By Definition \hyperref[III.def1.1]{1.1}, the
functor $u$ is continuous if and only if there exists a functor
$u_s:{C'}^{\sim}\rightarrow C^{\sim}$ such that
$iu_s=u^\ast i'$ (where $u^\ast F=F\circ u$
(\hyperrefext[I][I.sec5.0]{5.0})).

\begin{enonce*}{Proposition 1.2}\etiquette[1.2]{III.prop1.2}
Let $C$ be a small site, $C'$ a $𝒰$-site, and
$u:C\rightarrow C'$ a functor between their underlying
categories. The following properties are equivalent:
\begin{enumerate}
\renewcommand{\theenumi}{\roman{enumi}}
\renewcommand{\labelenumi}{\theenumi)}
\item The functor $u$ is continuous.

\item For every object $X$ of $C$ and every sieve $R\hookrightarrow X$
  covering $X$, the morphism $u_{!}R\hookrightarrow u(X)$ is
  bicovering in ${C'}^{\widehat{\phantom{,}}}$ {\rm(\hyperrefext[II][II.def5.2]{5.2}
  and \hyperrefext[I][I.prop5.1]{5.1}).} \footnote[1]{``Covering'' in place of
  ``bicovering'' did not necessarily suffice; \cf Example
  \hyperref[III.exem1.9.3]{1.9.3} below.}

\item For every bicovering family $H_i\rightarrow K$, $i\in I$, in
  $C^{\widehat{\phantom{,}}}$, the family of morphisms
  $u_{!}H_i\rightarrow u_{!}K$, $i\in I$, is bicovering in
  ${C'}^{\widehat{\phantom{,}}}$.

\item There exists a functor
  $u^s:C^{\sim}\rightarrow
  {C'}^{\sim}$, commuting with inductive limits and
  ``extending $u$'', \ie such that the diagram of canonical functors
\[
\xymatrix@C=3cm@R=1.25cm{
  C \ar[r]^{u} \ar[d]_{\epsilon_C} & C' \ar[d]^{\epsilon_{C'}}\\
C{\widetilde{\phantom{A}}} \ar[r]^{u^s}&
  C'{\widetilde{\phantom{A}}},
}
\]
commutes up to isomorphism.
\end{enumerate}

Moreover, when $C$ is no longer necessarily a small site but is a
$𝒰$-site, one still has {\rm i) $\Leftrightarrow$ iv)}, and the
functor $u^s$ of {\rm iv)} is necessarily a left adjoint of the functor
$u_s:\widetilde{C}'\rightarrow \widetilde{C}$, and is consequently
determined up to a unique isomorphism.
\end{enonce*}

\begin{enonce*}[remark]{Proof}
The\pageoriginale proof of i) $\Leftrightarrow$ iv) when $C$ is a
$𝒰$-site will be given in \hyperref[III.subsec4.2]{4.2}. Assume that
$C$ is small. It is clear that iii) $\Rightarrow$ ii).
\begin{itemize}
\item[i) $\Rightarrow$ iii):] Let $H_i\rightarrow K$, $i\in I$, be a
  bicovering family in $C^{\widehat{\phantom{,}}}$. For every sheaf
  $F$ on $C'$, the presheaf $u^\ast F$
  (\hyperrefext[I][I.sec5.0]{5.0}) is a sheaf on $C$. Hence
  (\hyperrefext[II][II.prop5.3]{5.3}) the canonical map
$$
\Hom(K, u^\ast F)\longrightarrow \prod_i \Hom(H_i, u^\ast F){\quoi},
$$
is bijective. By adjunction (\hyperrefext[I][I.prop5.1]{5.1}), it follows
that the canonical map
$$
\Hom(u_{!}K,F)\longrightarrow \prod_i \Hom(u_{!}H_i, F)
$$
is bijective. Consequently (\hyperrefext[II][II.prop5.3]{5.3}) the family
$u_{!}H_i\rightarrow u_{!}K$, $i\in I$, is bicovering.

\item[ii) $\Rightarrow$ i):] Let $X$ be an object of $C$, let
  $R\hookrightarrow X$ be a covering sieve, and let $F$ be a sheaf on
  $C'$. By (\hyperrefext[II][II.prop5.3]{5.3}), the map
$$
\Hom(u_{!}X,F)\longrightarrow \Hom(u_{!}R,F)
$$
is bijective. By adjunction (\hyperrefext[I][I.prop5.1]{5.1}), it follows
that the map
$$
\Hom(X,u^\ast F)\longrightarrow \Hom(R, u^\ast F)
$$
is bijective. Consequently $u^\ast F$ is a sheaf.

\item[i) $\Rightarrow$ iv):] We use the customary notation
  $\underline{a}$ and $i$ (\resp $\underline{a}'$ and $i'$) for the
  ``associated sheaf'' and ``inclusion into presheaves'' functors.
  For every sheaf $G$ on $C$, set $u^sG=\underline{a}'u_{!}iG$.
  For every sheaf $F$ on $C'$, we have the following sequence of
  natural isomorphisms:
\begin{align*}
\Hom(G,u_sF) &\simeq \Hom(iG, iu_s F)\simeq \Hom(iG, u^\ast i'F)\\
             &\simeq \Hom(u_{!}iG, i'F)\simeq \Hom(\underline{a}'u_{!} iG,
             F){\quoi}.
\end{align*}
(The\pageoriginale first isomorphism follows from the full
faithfulness of $i$; the second follows from the definition of $u_s$;
the third and fourth are obtained by adjunction.) Consequently $u^s$
is left adjoint to $u_s$. In particular, $u^s$ commutes with inductive
limits. For every presheaf $K$ on $C$ and every sheaf $F$ on $C'$, we
have the following sequence of natural isomorphisms:
$\Hom(u^s\underline{a}K,F)\simeq
\Hom(\underline{a}K, u_s F)\simeq \Hom(K,iu_sF)\simeq \Hom(K,u^\ast
i'F)\simeq \Hom(u_{!}K, i'F)\simeq \Hom(\underline{a}'u_{!}K,F)$ (the
third isomorphism follows from the definition of $u_s$; the others are
obtained by adjunction); whence a functorial isomorphism
$u^s\underline{a}K\simeq \underline{a}'u_{!}K$. In particular, using
this isomorphism when $K$ is representable, and taking
\hyperrefext[I][I.prop5.4]{5.4} 3) into account, we obtain a diagram
commutative up to isomorphism:
\[
\xymatrix@C=2.75cm@R=1.25cm{
C \ar[d]_{\epsilon_C} \ar[r]^{u}& C'\ar[d]^{\epsilon_{C'}}\\
\widetilde{C} \ar[r]^{u^{S}}& \widetilde{C'} }
\]

\item[iv) $\Rightarrow$ ii):] Let $h:C\rightarrow
  C^{\widehat{\phantom{,}}}$ (\resp $h':C'\rightarrow
  {C'}^{\widehat{\phantom{,}}}$) be the functor that assigns to an
  object the presheaf it represents. Consider the diagram:
\[
\xymatrix@C=2.5cm@R=1cm{
C \ar[r]^{u}\ar[d]_{h}             &   C' \ar[d]^{h'}\\
C{\widehat{\phantom{i}}} \ar[d]_{\underline{a}} \ar[r]^{u_!}&
C'{\widehat{\phantom{i}}} \ar[d]^{\underline{a}'}\\
C{\widetilde{\phantom{A}}} \ar[r]^{u^{S}}&
C'{\widetilde{\phantom{A}}}. }
\]
The\pageoriginale upper square is commutative
(\hyperref[III.rem1.4]{1.4} 3)), and we have
$\underline{a}h=\epsilon_C$, $\underline{a}'h'=\epsilon_{C'}$. For
every object $K$ of $C^{\widehat{\phantom{,}}}$, we have a canonical
isomorphism (I \hyperref[I.prop3.4]{3.4}):
$$
\varinjlim_{(h(X)\rightarrow K)\in \ob (C/K)} h(X)\xrightarrow{\sim}
K{\quoi}.
$$
Since the functors $u_{!}$, $u^s$, $\underline{a}$, and
$\underline{a}'$ commute with inductive limits, we obtain
isomorphisms:
\begin{gather*}
\varinjlim_{(h(X)\rightarrow K)\in \ob (C/K)}
u^s\underline{a}h(X)\xrightarrow{\sim} u^s \underline{a}K{\quoi},\\
\varinjlim_{(h(X)\rightarrow K)\in \ob (C/K)}
\underline{a}'u_{!}h(X)\xrightarrow{\sim} \underline{a}'u_{!}K{\quoi}.
\end{gather*}
We have $u_{!}h=h'u$ and, by hypothesis, an isomorphism
$\underline{a}'h'u\simeq u^s \underline{a}h$, whence an isomorphism
$$
u^s\underline{a}K\simeq \underline{a}'u_{!}K{\quoi},
$$
which is immediately seen to be functorial in $K$. The diagram
\[
\xymatrix@C=2.5cm@R=1.35cm{
C\widehat{\phantom{a}} \ar[r]^{u_!} \ar[d]_{\underline{a}}&  C'\widehat{\phantom{a}}\ar[d]^{\underline{a}'}\\
C \widetilde{\phantom{A}} \ar[r]^{u^{S}}&  C',& }
\]
is therefore commutative up to isomorphism. Since $\underline{a}i$ is
isomorphic to the identity functor, we have an isomorphism
$u^s\simeq \underline{a}'u_{!}i$, whence the uniqueness of $u^s$. Let
$v:H\rightarrow K$ be a bicovering morphism in
$C^{\widehat{\phantom{,}}}$. Then $\underline{a}(v)$ is an isomorphism
(\hyperrefext[II][II.prop5.3]{5.3}), and consequently
$\underline{a}'u_{!}(v)$ (isomorphic to $u^s\underline{a}(v)$) is an
isomorphism. Thus $u_{!}(v)$ is carried by $\underline{a}'$ to an
isomorphism, and consequently $u_{!}(v)$ is bicovering
(\hyperrefext[II][II.prop5.3]{5.3}), whence ii).
\end{itemize}
\end{enonce*}

The additional assertion,\pageoriginale in the case where $C$ is
small, was proved in the course of the proof, \cqfd

\setcounter{subsection}{2} \setcounter{subsubsection}{0}

% END INCLUDED SGA4_English_current_buildable_source_20260906/source/03/drafts/SGA4_III_LINES_31_241_English.texfrag


% BEGIN INCLUDED SGA4_English_current_buildable_source_20260906/source/03/drafts/SGA4_III_LINES_242_461_English.texfrag
\subsubsection{}\etiquette[1.2.1]{III.subsubsec1.2.1}

The functor $u^s:C^{\sim}\rightarrow {C'}^{\sim}$ of
\hyperref[III.prop2.2]{2.2} iv) may often be interpreted as an
``inverse image'' functor for a ``morphism of topoi''
${C'}^{\sim}\rightarrow C^{\sim}$, \cf \hyperrefext[IV][IV.subsec4.9]{4.9}.
Its properties are summarized in the following proposition:

\begin{enonce*}{Proposition 1.3}\etiquette[1.3]{III.prop1.3}
Let $u:C\rightarrow C'$ be a continuous functor from a small site $C$
to a $U$-site $C'$.
\begin{enumerate}
\renewcommand{\labelenumi}{\theenumi)}
\item The functor $u^s$ is left adjoint to the functor $u_s$.

\item There is a canonical isomorphism $u^s\simeq \underline{a}'u_{!}i$.

\item There is a canonical isomorphism $u^s\underline{a}\simeq
  \underline{a}'u_{!}$.

\item The functor $u^s$ commutes with inductive limits.

\item When the functor $u_{!}$ is left exact (\cf \hyperrefext[I][I.prop5.4]{5.4} 4)),
  the functor $u^s$ is also left exact. More generally, the functor
  $u^s$ commutes with those types of finite projective limits with
  which the functor $u_{!}$ commutes.
\end{enumerate}
\end{enonce*}

\begin{enonce*}[remark]{Proof}
Assertions 1), 2), 3), and 4) were proved in the course of the proof
of \hyperref[III.prop1.2]{1.2}. Assertion 5) follows from the
isomorphism in 2), since $i$ commutes with projective limits and
$\underline{a}'$ commutes with finite projective limits
(\hyperrefext[II][II.theo4.1]{4.1}).
\end{enonce*}

\begin{enonce*}[remark]{Remark 1.4}\etiquette[1.4]{III.rem1.4}
Combining \hyperrefext[I][I.prop5.4]{5.4} 3) and
\hyperref[III.prop1.3]{1.3} 3), we obtain a diagram:
\[
\xymatrix@C=3.5cm@R=1cm{
C \ar[r]^{u} \ar[d]_{h}& C'\ar[d]_{h'}\\
C\widehat{\phantom{a}} \ar[r]^{u_!} \ar[d]_{\underline{a}}&
C'\widehat{\phantom{a}} \ar[d]_{\underline{a}'}\\
C\widetilde{\phantom{a}} \ar[r]& C'\widetilde{\phantom{a}}, }
\]
where\pageoriginale the upper square commutes and the lower square
commutes up to isomorphism. But the functors $\underline{a}$,
$\underline{a}'$, $u_{!}$, and $u^s$ are defined only up to
isomorphism. The reader may verify that the functors $\underline{a}$,
$\underline{a}'$, and $u^s$ can be chosen so that:
\begin{enumerate}
\renewcommand{\labelenumi}{\theenumi)}
\item The composite functors $\underline{a}h$ and
  $\underline{a}'h'$ are injective on the sets of objects;

\item the diagram
\[
\xymatrix@C=2cm@R=1.25cm{
C \ar[r]^{u} \ar[d]_{\underline{a}h}& C' \ar[d]_{\underline{a}'h'}\\
C{\widetilde{\phantom{a}}} \ar[r]^{u^{s}}&
C'{\widetilde{\phantom{a}}} }
\]
commutes.
\end{enumerate}

More precisely, the functors $\underline{a}$ and $\underline{a}'$ can
be chosen so as to satisfy condition 1). Once the functors
$\underline{a}$ and $\underline{a}'$ have been chosen, the functor
$u^s$ can be chosen so as to satisfy condition 2).
\end{enonce*}

\begin{enonce*}{Proposition 1.5}\etiquette[1.5]{III.prop1.5}
Let $𝒰\subset 𝒱$ be two universes, let $C$ and $C'$ be two
$𝒰$-sites, and let $u:C\rightarrow C'$ be a functor between the
underlying categories. Then $u$ is continuous relative to $𝒰$ if
and only if it is continuous relative to $𝒱$. Moreover, when $u$ is
continuous, denote by $u^s_{𝒰}$ (\resp $u^s_{𝒱}$) the functor between
the categories of $𝒰$-sheaves (\resp $𝒱$-sheaves) introduced in
\hyperref[III.prop1.2]{1.2} {\rm iv)}. Then the diagram
\[\vcenter{
\xymatrix@C=3cm@R=1.5cm{ \widetilde{C}_{𝒰}
\ar[r]^{u_{\underline{U}}^s} \ar[d]&
\widetilde{C'}_{\underline{U}} \ar[d]\\
 C\widetilde{\phantom{a}}_{\underline{V}} \ar[r]^{u_{\underline{V}}^s}
 &  C'\widetilde{\phantom{a}}_{\underline{V}} ,
}}\leqno{(1.5.1)}
\]\etiquette[1.5.1]{III.eq1.5.1}
where\pageoriginale the vertical functors are the canonical inclusion
functors, commutes up to canonical isomorphism.
\end{enonce*}

\begin{enonce*}[remark]{Proof}
We give the proof only in the case where $C$ is $𝒰$-small. The
general case will be proved in \hyperref[III.subsec4.3]{4.3}. It is
clear that if the functor $u^\ast$ carries every $𝒱$-sheaf on $C'$
to a $𝒱$-sheaf on $C$, then it carries every $𝒰$-sheaf on $C'$ to a
$𝒰$-sheaf on $C$. Suppose that $u^\ast$ carries every $𝒰$-sheaf on
$C'$ to a $𝒰$-sheaf on $C$. Denote by
$C^{\widehat{\phantom{,}}}_{𝒰}\hookrightarrow
C^{\widehat{\phantom{,}}}_{𝒱}$ (\resp
${C'}^{\widehat{\phantom{,}}}_{𝒰}\hookrightarrow
{C'}_{𝒱}^{\widehat{\phantom{,}}}$) the categories of
$𝒰$-presheaves and $𝒱$-presheaves, and by $u_{𝒰!}$ and $u_{𝒱!}$
the ``inverse image'' functors for $𝒰$-presheaves and
$𝒱$-presheaves, respectively. We have a commutative diagram
(\cf the explicit construction of $u_{!}$ in
\hyperrefext[I][I.prop5.1]{5.1}):
\[\vcenter{
  \xymatrix@C=3cm@R=1.25cm@M=.2cm{
C\widehat{\phantom{a}}_{\underline{U}} \ar[r]^{u_{{\underline{U}}!}}
\ar @{_{(}->}[d]&
C'\widehat{\phantom{a}}_{\underline{V}} \ar @{_{(}->}[d]\\
C\widehat{\phantom{a}}_{\underline{V}}
\ar[r]^{u_{{\underline{V}}!}}&
C'\widehat{\phantom{a}}_{\underline{V}}. }}\leqno{(*)}
\]
The inclusion functors from $𝒰$-presheaves into $𝒱$-presheaves are
fully faithful and commute with projective limits. It follows from
(\hyperrefext[II][II.prop5.1]{5.1} i)) that a bicovering morphism of
$𝒰$-presheaves is a bicovering morphism of $𝒱$-presheaves. Now let
$R\hookrightarrow X$ be a sieve covering an object $X$ of $C$. By
\hyperref[III.prop1.2]{1.2} ii), the morphism
$u_{𝒰!}R\rightarrow u_{𝒰!}X$ is bicovering. Using the commutativity of
$(\ast)$ and what precedes, we deduce that the morphism
$u_{𝒱!}R\rightarrow u_{𝒱!}X$ is bicovering. Consequently
(\hyperref[III.prop1.2]{1.2}), the functor $u^\ast$ carries
$𝒱$-sheaves on $C'$ to $𝒱$-sheaves on $C$. The commutativity of
\eqhyperref[III.eq1.5.1]{1.5.1} then follows from the commutativity of
$(\ast)$, from \hyperrefext[II][II.prop3.6]{3.6}, and from
\hyperref[III.prop1.3]{1.3} 2).
\end{enonce*}

\begin{enonce*}{Proposition 1.6}\etiquette[1.6]{III.prop1.6}
Let\pageoriginale $C$ and $C'$ be two $𝒰$-sites, and let
$u:C\rightarrow C'$ be a functor between the underlying categories.
Consider the following conditions:
\begin{enumerate}
\renewcommand{\theenumi}{\roman{enumi}}
\renewcommand{\labelenumi}{\theenumi)}
\item $u$ is continuous {\rm(\cf. \hyperref[III.def1.1]{1.1} and
  \hyperref[III.prop1.2]{1.2}).}

\item For every covering family $(X_i\rightarrow X)$, $i\in I$, in
  $C$, the family $(uX_i\rightarrow uX)$, $i\in I$, in $C'$ is
  covering.

We have the implication {\rm i) $\Rightarrow$ ii).}

Suppose that the topology of $C$ can be defined by a pretopology
{\rm $T$ (\hyperrefext[II][II.def1.3]{1.3})} and that the functor $u$
commutes with fiber products.\footnote[1]{In fact, it suffices that
$u$ commute with the fiber products occurring in base changes of
morphisms belonging to covering families for $T$.} Then conditions
{\rm i)} and {\rm ii)} are equivalent and are also equivalent to the
following condition:

\item The functor $u$ carries the covering families of $T$ to
  covering families in $C'$.
\end{enumerate}

In particular, suppose that fiber products in $C$ are representable
and that the functor $u$ commutes with fiber products. Then conditions
{\rm i)} and {\rm ii)} are equivalent.
\end{enonce*}

\begin{enonce*}[remark]{Proof}
Let us prove that i) $\Rightarrow$ ii). Let $X_i\rightarrow X$,
$i\in I$, be a covering family in $C$. For every sheaf $F$ on $C'$,
$u^\ast F$ is a sheaf. Consequently, the map
$$
\Hom(X,u^\ast F)\longrightarrow \prod_i\Hom(X_i,u^\ast F)
$$
is injective (\hyperrefext[II][II.prop5.1]{5.1}). Hence there is an
injective map
$$
\Hom(u_{!}X,F)\longrightarrow \prod_i \Hom(u_{!}X_i, F){\quoi}.
$$
Therefore the family $u_{!}X_i\rightarrow u_{!}X$, $i\in I$, is
covering in $C'$ (\hyperrefext[II][II.prop5.1]{5.1}).

Every covering family of the pretopology $T$ is a covering family in
$C$, whence ii) $\Rightarrow$ iii). Let us prove that iii)
$\Rightarrow$ i). Let $X$ be an object\pageoriginale of $C$ and let
$X_i\rightarrow X$, $i\in I$, be a covering family in $\Cov(X)$ (II
\hyperref[II.def1.3]{1.3}). The morphisms $X_i\rightarrow X$ are pullbackable, and the
functor $u$ commutes with fiber products. Consequently, the fiber
products $u(X_i)\times_{u(X)} u(X_j)$ are representable and
canonically isomorphic to $u(X_i\times_X X_j)$. Let $F$ be a sheaf on
$C'$. Then the diagram of sets
$$
F(u(X))\rightarrow \prod_i F(u(X_i))\rightrightarrows
\prod_{i,j}F(u(X_i\times_X X_j))
$$
is exact (\hyperrefext[II][II.def2.1]{2.1},
\hyperrefext[I][I.cor3.5]{3.5}, and \hyperrefext[I][I.prop2.12]{2.12}).
Consequently, the diagram of sets
$$
u^\ast F(X)\rightarrow \prod_i u^\ast F(X_i)\rightrightarrows
\prod_{i,j}u^\ast F(X_i\times_X X_j)
$$
is exact. Therefore $u^\ast F$ is a sheaf
(\hyperrefext[II][II.cor2.4]{2.4}), whence i). The final assertion of
\hyperref[III.prop1.6]{1.6} follows from what precedes and from the
fact that, when fiber products are representable in $C$, all covering
families in $C$\footnote[1]{indexed by sets belonging to a suitable
universe.} define a pretopology on $C$ whose associated topology is
the topology of $C$.
\end{enonce*}


% END INCLUDED SGA4_English_current_buildable_source_20260906/source/03/drafts/SGA4_III_LINES_242_461_English.texfrag


% BEGIN INCLUDED SGA4_English_current_buildable_source_20260906/source/03/drafts/SGA4_III_LINES_462_612_English.texfrag
\begin{enonce*}{Proposition 1.7}\etiquette[1.7]{III.prop1.7}
Let $C$ be a small site, $C'$ a $𝒰$-site, and
$u:C\rightarrow C'$ a continuous functor. Let $\gamma$ be an algebraic
structure defined by finite projective limits, such that in the category
of $\gamma$-objects of $𝒰-\Ens$ the $𝒰$-indexed inductive limits are
representable
\footnote[2]{In fact, one can show that this condition is always
  satisfied.}. We use the notation of Proposition {\rm II \hyperref[II.prop6.4]{6.4}.} The
functor $u_s$ commutes with projective limits and consequently defines a
functor $u^{\gamma}_s:C^{\sim}_{\gamma}\rightarrow
{C'}^{\sim}_{\gamma}$. The functor $u^{\gamma}_s$ has a left adjoint
$u^s_{\gamma}$ with the following properties:
\begin{enumerate}
\renewcommand{\labelenumi}{\theenumi)}
\item There\pageoriginale is a canonical isomorphism
  $u^s_{\gamma}\xrightarrow{\sim}\underline{a}'_{\gamma}u_{\gamma !}
  i_{\gamma}$, and $u^s_{\gamma}$ commutes with every finite projective
  limit with which $u_{\gamma !}$ commutes.

\item There is a canonical isomorphism
  $u^s_{\gamma}\underline{a}_{\gamma}\xrightarrow{\sim}
  \underline{a}'_{\gamma} u'_{\gamma !}$.

\item The functor $u^s_{\gamma}$ commutes with inductive limits.

\item If $u^s$ is left exact, the diagram {\rm(in the notation of II
  \hyperref[II.subsubsec6.3.0]{6.3.0})}
\[
\xymatrix@C=3.5cm@R=1.25cm{ C^{\widetilde{\phantom{a}}_\gamma}
\ar[r]^{u_{Y}^{s}}
\ar[d]_{\esj^{\widetilde{\phantom{a}}}}& C'\widetilde{\phantom{a}}_\gamma \ar[d]_{{\esj'}^{\widetilde{\phantom{a}}}}\\
C^{\widetilde{\phantom{a}}} \ar[r]^{u^{s}}  &
{C'}\widetilde{\phantom{a}} }
\]
is commutative up to isomorphism, and the functor $u^s_{\gamma}$ is
exact.

\item Suppose that the functor $\esj^{\sim}$
  (\resp $\esj^{\sim}$) has a left adjoint
  $\Lib^{\sim}$ (\resp
  ${\Lib'}^{\sim}$) {\rm (\hyperrefext[II][II.prop6.5]{6.5})}. There is a
  canonical isomorphism
  $u^s_{\gamma}\Lib^{\sim}\xrightarrow{\sim}
  {\Lib'}^{\sim} u^s$.
\end{enumerate}

Moreover, when $C$ is not necessarily a small site but is a $𝒰$-site,
the functor $u^{\gamma}_s$ has a left adjoint $u^s_{\gamma}$. Finally,
if $V$ is a universe containing $𝒰$ and if $u^s_{\gamma𝒰}$ (\resp
$u^s_{\gamma𝒱}$) denotes the left adjoint of $u_s$ relative to the
universe $𝒰$ (\resp $𝒱$), the following diagram, analogous to diagram
\eqhyperref[III.eq1.5.1]{1.5.1}:

\hypertarget{sga4.exp.III.equation.1.7.1}{}\[\vcenter{
\xymatrix@C=3.5cm@R=1.25cm{ C^{\widetilde{\phantom{a}}}_{\gamma
\underline{U}} \ar[r]^{u_{\gamma \underline{U}}^{s}}
\ar@{_{(}->}[d]&
\widetilde{C}'_{\gamma \underline{U}} \ar@{_{(}->}[d]\\
C^{\widetilde{\phantom{a}}}_{\gamma \underline{V}}
\ar[r]^{u^s_{r\underline{V}}}  &  \widetilde{C}'_{\gamma
\underline{V}} }}\leqno{(1.7.1)}
\]
is commutative up to canonical isomorphism.
\end{enonce*}

\begin{enonce*}[remark]{Proof}
The proof, in the case where $C$ is a small site, is left to the
reader. It will be completed in \hyperref[III.subsec4.3]{4.3} in the
case where $C$ is a $𝒰$-site.
\end{enonce*}

\begin{enonce*}[remark]{Notation 1.8}\etiquette[1.8]{III.not1.8}
Henceforth\pageoriginale we shall use the notation $u_s$ for the
functor $u^{\gamma}_s$. This notation is unlikely to cause confusion,
because the functor $u_s$ (defined on sheaves of sets) commutes with
finite projective limits.
\end{enonce*}

Likewise, when $u^s$ is left exact, we shall use the notation $u^s$ for
the functor $u^s_{\gamma}$. This abuse of notation is justified by
\hyperref[III.prop1.7]{1.7} 4).

\begin{enonce*}[remark]{Example 1.9.1}\etiquette[1.9.1]{III.exem1.9.1}
Let $X$ be a small topological space. Denote by $\Ouv(X)$ the category
of open subsets of $X$ (the objects of $\Ouv(X)$ are the open subsets
of $X$, and the morphisms of $\Ouv(X)$ are inclusions), equipped with
the canonical topology. A family $(U_i\subset U)$, $i\in I$, is
covering if and only if the open subsets $U_i$ cover $U$
(\hyperrefext[II][II.prop2.6]{2.6}). The sheaves of sets on $\Ouv(X)$
are therefore the sheaves of sets in the sense of [TF]. Let
$f:X\rightarrow Y$ be a continuous map. It induces a functor
$f^{-1}:\Ouv(Y)\rightarrow \Ouv(X)$: for every open subset $U$ of $Y$,
$f^{-1}(U)$ is the inverse image of $U$ under the map $f$. The functor
$f^{-1}$ is continuous (\hyperref[III.prop1.6]{1.6}). The functor
$f^{-1}_s:\Ouv(X)^{\sim}\rightarrow \Ouv(Y)^{\sim}$ is the direct-image
functor for sheaves, in the sense of [TF]. The functor
$(f^{-1})^s:\Ouv(Y)^{\sim}\rightarrow \Ouv(X)^{\sim}$ is the
inverse-image functor for sheaves in the sense of [TF]. The functor
$(f^{-1})^s$ commutes with finite projective limits, by
\hyperref[III.prop1.3]{1.3} 5).
\end{enonce*}

\begin{enonce*}[remark]{Example 1.9.2}\etiquette[1.9.2]{III.exem1.9.2}
Let $X$ be a topological space and let $Y\hookrightarrow X$ be an open
subset of $X$. Every open subset of $Y$ is an open subset of $X$, and
thus we obtain a functor $\Phi:\Ouv(Y)\rightarrow \Ouv(X)$. The functor
$\Phi$ is continuous (\hyperref[III.prop1.6]{1.6}). The functor
$\Phi_s:\Ouv(X)^{\sim}\rightarrow
\Ouv(Y)^{\sim}$ is the\pageoriginale ``restriction to $Y$'' functor.
The functor $\Phi^s_{\ab}:\Ouv(Y)^{\sim}_{\ab}\rightarrow
\Ouv(X)^{\sim}_{\ab}$ (\hyperrefext[II][II.subsubsec6.3.3]{6.3.3}, I \hyperref[III.prop1.7]{1.7})
is the ``extension by zero outside $Y$'' functor introduced in [TF].
The functor $\Phi^s:\Ouv(Y)^{\sim}\rightarrow
\Ouv(X)^{\sim}$ (\hyperref[III.prop1.3]{1.3}) is analogous to the
``extension by zero'' functor: for every sheaf $H$ on $Y$ and every
point $y\in Y$, the stalk of $\Phi^sH$ at $y$ is canonically
isomorphic to the stalk of $H$ at $y$; for every point $x\in X$,
$x\not\in Y$, the stalk of $\Phi^sH$ at $x$ is empty. The functor
$\Phi^s$ is called the ``extension by the empty set outside $Y$''
functor. The functor $\Phi^s$ commutes with fiber products and with
finite products over a nonempty index set, but it does not carry the
terminal object of $\Ouv(Y)^{\sim}$ to the terminal object of
$\Ouv(X)^{\sim}$.
\end{enonce*}

\begin{enonce*}[remark]{Example 1.9.3}\etiquette[1.9.3]{III.exem1.9.3}
Let $C$ and $C'$ be two small sites, and let $u:C\rightarrow C'$ be a
functor between the underlying categories that carries every covering
family of $C$ to a covering family of $C'$. The functor $u$ is not
continuous in general (\cf nevertheless \hyperref[III.prop1.6]{1.6}).
Here is a counterexample. Let $S$ be a set of infinite cardinality in
the universe $𝒰$, and let $C$ be the full subcategory of the category
of sets whose objects are the finite, nonempty subsets of $S$. Equip
$C$ with the canonical topology. Observe that the covering families of
$C$ are the surjective families of maps and that covering families are
nonempty. Observe also that, up to isomorphism, there are only two
constant sheaves on $C$: the sheaf with value the empty set and the
sheaf with value a one-element set. Set $C=C'$ and let
$u:C\rightarrow C'$ be a constant functor. The functor $u$ carries the
covering families of $C$ to covering families of $C'$. The functor $u$
is not continuous. Indeed, since the representable presheaves of $C'$
are sheaves, for every object $X$ of $C'$ there exists a
sheaf\pageoriginale $F$ on $C'$ such that the cardinality of $F(X)$ is
$\geq 2$, and consequently the presheaf $u^\ast F$ is not a sheaf.
\end{enonce*}


% END INCLUDED SGA4_English_current_buildable_source_20260906/source/03/drafts/SGA4_III_LINES_462_612_English.texfrag


% BEGIN INCLUDED SGA4_English_current_buildable_source_20260906/source/03/drafts/SGA4_III_LINES_613_861_English.texfrag
\section{Cocontinuous functors}\etiquette[2]{III.sec2}

\begin{enonce*}{Definition 2.1}\etiquette[2.1]{III.def2.1}
Let $C$ and $C'$ be two $𝒰$-sites, and let $u:C\rightarrow C'$ be a
functor between the underlying categories. We say that $u$ is
cocontinuous if it has the following property:

$(COC)$ For every object $Y$ of $C$ and every covering sieve
$R\hookrightarrow u(Y)$, the sieve on $Y$ generated by the arrows
$Z\rightarrow Y$ such that $u(Z)\rightarrow u(Y)$ factors through
$R$ covers~$Y$.
\end{enonce*}

\begin{enonce*}{Proposition 2.2}\etiquette[2.2]{III.prop2.2}
Let $C$ be a small site, let $C'$ be a $𝒰$-site, and let
$u:C\rightarrow C'$ be a functor between the underlying categories.
Denote by $\widehat{u}_{\ast}:C^{\widehat{\phantom{,}}}\rightarrow
{C'}^{\widehat{\phantom{,}}}$ the functor right adjoint to the
functor $F\mapsto F\circ u=\widehat{u}^{\ast}F$ {\rm
(\hyperrefext[I][I.prop5.1]{5.1}).} The functor $u$ is cocontinuous if and only if,
for every sheaf $G$ on $C$, $\widehat{u}_{\ast}G$ is a sheaf on $C'$.
\end{enonce*}

\begin{enonce*}[remark]{Proof}
\emph{The condition is sufficient.} Suppose that for every sheaf
$G$ on $C$, the presheaf $\widehat{u}_{\ast}G$ is a sheaf. Let $Y$
be an object of $C$, and let $R\hookrightarrow u(Y)$ be a covering
sieve. There is an isomorphism $\Hom(u(Y),
u_{\ast}G)\xrightarrow{\sim}\Hom(R,u_{\ast}G)$, hence, by adjunction
(\hyperrefext[I][I.prop5.1]{5.1}), an isomorphism
$\Hom(\widehat{u}^{\ast}\circ
u(Y),G)\xrightarrow{\sim}\Hom(\widehat{u}^{\ast}R,G)$. It follows
(\hyperrefext[II][II.prop5.3]{5.3}) that the morphism $\widehat{u}^{\ast}R\rightarrow
\widehat{u}^{\ast}u(Y)$ is bicovering. But the functor
$\widehat{u}^{\ast}$ commutes with projective limits, and therefore
the morphism $\widehat{u}^{\ast}R\rightarrow \widehat{u}^{\ast}u(Y)$
is a monomorphism. It is thus a covering monomorphism. On the other
hand, there is a canonical morphism
$Y\xrightarrow{p}\widehat{u}^{\ast}u(Y)$ deduced from the adjunction
morphism $\id_C\rightarrow
\widehat{u}^{\ast}u_{!}$\pageoriginale (\hyperrefext[I][I.prop5.4]{5.4} 3)). We then
make the base change $Y\xrightarrow{p}\widehat{u}^{\ast}u(Y)$. This
gives a covering sieve on $Y$:
$$
\widehat{u}^{\ast}R\times_{\widehat{u}^{\ast}u(Y)}Y\longrightarrow
Y{\quoi},
$$
and, on determining the arrows $Z\rightarrow Y$ that factor through
this sieve, one obtains the sieve described by property (COC).
\end{enonce*}

\emph{The condition is necessary.} Let $S$ be an object of $C'$
and let $R\hookrightarrow S$ be a covering sieve on $S$. We must
prove that, for every sheaf $G$ on $C$, there is an isomorphism
$$
\Hom(S,u_{\ast}G)\xrightarrow{\sim} \Hom(R,u_{\ast}G){\quoi}.
$$

Using Proposition \hyperrefext[II][II.prop5.3]{5.3} and the adjunction
isomorphisms, we see that it suffices to prove that the monomorphism
$\widehat{u}^{\ast}R\rightarrow \widehat{u}^{\ast}S$ is covering.
For this, it suffices to prove (\hyperrefext[II][II.prop5.1]{5.1}) that for every
base change $Y\rightarrow \widehat{u}^{\ast}S$, the sieve
$\widehat{u}^{\ast}R\times_{\widehat{u}^{\ast}S}Y$ is covering,
where $Y$ is an object of $C$. But, by the properties of adjoint
functors and \hyperrefext[I][I.prop5.4]{5.4} 3), the morphism
$Y\rightarrow \widehat{u}^{\ast}S$ factors through
$Y\xrightarrow{p}\widehat{u}^{\ast}u(Y)\rightarrow
\widehat{u}^{\ast}S$, and consequently the sieve
$\widehat{u}^{\ast}R\times_{\widehat{u}^{\ast}S}Y\hookrightarrow Y$
is obtained, by the base change
$Y\xrightarrow{p}\widehat{u}^{\ast}u(Y)$, from the monomorphism
$\xymatrix@1@M=.15cm{
\fprod{\hat{u}^{\ast}R}{\hat{u}^{\ast}u(Y)}{\hat{u}^{\ast}S}
\ar@{^{(}->}@<-1pt>[r]^(.6){q}&\hat{u}^{\ast}u(Y)}$.
Now the functor $\widehat{u}^{\ast}$ commutes with projective limits,
and consequently the monomorphism $q$ is the image under
$\widehat{u}^{\ast}$ of the sieve $R\times_S u(Y)\hookrightarrow u(Y)$,
which is covering. Hypothesis (COC) then allows us to conclude that
the sieve
$\widehat{u}^{\ast}R\times_{\widehat{u}^{\ast}S}Y\hookrightarrow Y$
is covering, \cqfd

\begin{enonce*}{Proposition 2.3}\etiquette[2.3]{III.prop2.3}
Let $C$ and $C'$ be two $𝒰$-sites, and let $u:C\rightarrow C'$ be a
cocontinuous functor. Denote by $u^{\ast}:\widehat{C}'\rightarrow
\widehat{C}$ the functor $u^{\ast}=\underline{a}\widehat{u}^{\ast}i'$,
where $\underline{a}$ denotes the ``associated sheaf'' functor for
$C^{\widehat{\phantom{,}}}$, $\widehat{u}^{\ast}$ the functor
$F\mapsto F\circ u$, and $i'$ the inclusion functor\pageoriginale
$\widetilde{C}'\rightarrow {C'}^{\widehat{\phantom{,}}}$.
\begin{enumerate}
\renewcommand{\labelenumi}{\theenumi)}
\item The functor $u^{\ast}$ commutes with inductive limits and is
  exact.

\item There is a canonical isomorphism $u^{\ast}\underline{a}'\simeq
  \underline{a}\widehat{u}^{\ast}$, where $\underline{a}'$ denotes
  the ``associated sheaf'' functor for
  ${C'}^{\widehat{\phantom{,}}}$.

\item The functor $u^\ast$ has a right adjoint
  $u_{\ast}:\widetilde{C}\rightarrow \widetilde{C}'$. When $C$ is
  small, the diagram
\[\vcenter{
\xymatrix@C=3.25cm@R=1.25cm@M=.25cm{
C\widetilde{\phantom{a}}\ar[r]^{u_\ast} \ar@{_{(}->}[d]_{i}&
\widetilde{C}' \ar@{_{(}->}[d]^{i'}\\
C\widehat{\phantom{a}} \ar[r]^{\widehat{u}_{\ast}}  &
C'\widehat{\phantom{a}}, }}\leqno{(2.3.1)}
\]\etiquette[2.3.1]{III.eq2.3.1}
in which the lower functor is right adjoint to the functor
$\widehat{u}^{\ast}$ (\hyperrefext[I][I.prop5.1]{5.1}), commutes up
to canonical isomorphism.

\item Let $𝒱\supset 𝒰$ be a universe. Denote by
  $u_{\ast_{𝒰}}:
  C^{\sim}_{𝒰}\rightarrow
  \widetilde{C}'_{𝒰}$ (\resp
  $u_{\ast_{𝒱}}:C^{\sim}_{𝒱}\rightarrow
  \widetilde{C}'_{𝒱}$) the right adjoint relative to the universe
  $𝒰$ (\resp $𝒱$) whose existence is asserted in $3$). The diagram
\[\vcenter{
\xymatrix@C=3.25cm@R=1.25cm@M=.2cm{
C\widetilde{\phantom{a}}_{\underline{U}}
\ar[r]^{u_{\ast_{\underline{U}}}} \ar@{_{(}->}[d]&
\widetilde{C}'_{\underline{U}} \ar@{_{(}->}[d]\\
C\widetilde{\phantom{a}}_{\underline{V}}
\ar[r]^{u_{\ast_{\underline{V}}}}& \widetilde{C}'_{\underline{V}}
}}\leqno{(2.3.2)}
\]\etiquette[2.3.2]{III.eq2.3.2}
commutes up to isomorphism.
\end{enumerate}
\end{enonce*}

\begin{enonce*}[remark]{Proof}
Assertions 1) and 2) follow immediately from
\hyperrefext[II][II.theo3.4]{3.4}. Assertion 3), when $C$ is small,
follows from \hyperref[III.prop2.2]{2.2}. When $C$ is a $𝒰$-site, it
will be proved in \hyperref[III.subsec4.4]{4.4}. Assertion 4), when
$C$ is small, follows from the analogous commutativity assertion when
the topologies on $C$ and $C'$ are chaotic and from
\hyperrefext[I][I.cor3.5]{3.5}. When the topologies on $C$ and $C'$
are chaotic,\pageoriginale the commutativity of
\eqhyperref[III.eq2.3.2]{2.3.2} follows immediately from the explicit
description of $u_{\ast}$ (\hyperrefext[I][I.prop5.1]{5.1}).
\end{enonce*}

\setcounter{subsection}{3}
\subsection{}\etiquette[2.4]{III.subsec2.4}
We shall not develop here the considerations concerning
$\gamma$-objects of categories of sheaves. It will suffice to note
that the functors $u^{\ast}$ and $u_{\ast}$ introduced in this
section always commute with finite projective limits (unlike what
happened in the preceding section for the functor $u^s$).
Consequently, they extend naturally to functors defined on
$\gamma$-objects; these functors are adjoint to one another and
commute with the ``underlying sheaf of sets'' functors. We shall make
the notational abuses mentioned in \hyperref[III.not1.8]{1.8}, using
$u^{\ast}$ and $u_{\ast}$ also for the extensions to $\gamma$-objects.

\begin{enonce*}{Proposition 2.5}\etiquette[2.5]{III.prop2.5}
Let $C$ and $C'$ be two $𝒰$-sites, and let $\xymatrix{ C
\ar@<2pt>[r]^{v}& C'\ar@<2pt>[l]^{u} }$ be a pair of functors, where
$v$ is left adjoint to $u$. The following properties are equivalent:
\begin{enumerate}
\renewcommand{\theenumi}{\roman{enumi}}
\renewcommand{\labelenumi}{\theenumi)}
\item The functor $u$ is continuous.

\item The functor $v$ is cocontinuous.
\end{enumerate}
Moreover, under these equivalent conditions, there are canonical
isomorphisms $v_{\ast}\simeq u_{s}$ and $v^{\ast}\simeq u^s$. In
particular, the functor $u^s$ commutes with finite projective limits.
\end{enonce*}

\begin{enonce*}[remark]{Proof}
The property of a functor being continuous or cocontinuous does not
depend on the universe $𝒰$ for which $C$ and $C'$ are $𝒰$-sites
(this is immediate for a cocontinuous functor, and follows from
\hyperref[III.prop1.5]{1.5} for a continuous functor). We may
therefore, after enlarging the universe, suppose that $C$ and $C'$
are small. The proposition then follows from
\hyperref[III.prop2.2]{2.2} and \hyperrefext[I][I.prop5.6]{5.6}.
\end{enonce*}

\begin{enonce*}{Proposition 2.6}\etiquette[2.6]{III.prop2.6}
Let\pageoriginale $u:C\rightarrow C'$ be a continuous and
cocontinuous functor between two $𝒰$-sites. Then the functor
$u^{\ast}:{C'}^{\sim}\rightarrow C^{\sim}$
$(\hyperref[III.prop2.3]{2.3})$ commutes with $𝒰$-inductive and
projective limits. Denote by $u_{!}:C^{\sim}\rightarrow {C'}^{\sim}$
and $u_{\ast}:C^{\sim}\rightarrow {C'}^{\sim}$ the functors left and
right adjoint to $u^{\ast}$, respectively. The functor $u_{!}$ is
fully faithful if and only if the functor $u_{\ast}$ is fully
faithful. When $u$ is fully faithful, $u_{!}$ is fully faithful, and
the converse holds when the topologies of $C$ and $C'$ are coarser
than the canonical topology.
\end{enonce*}

The functor $u^{\ast}$ has a right adjoint $u_{\ast}$
(\hyperref[III.prop2.3]{2.3}) and a left adjoint $u_{!}$
(\hyperref[III.prop1.2]{1.2}). The functor $u^{\ast}$ therefore
commutes with inductive and projective limits
(\hyperrefext[I][I.prop2.11]{2.11}). The fact that $u_{!}$ is fully
faithful if and only if $u_{\ast}$ is fully faithful is a general
property of adjoint functors (\hyperrefext[I][I.rem5.7]{5.7} 1).
When $u$ is fully faithful, the functor
$\widehat{u}_{\ast}:C^{\widehat{\phantom{,}}}_{𝒱}\rightarrow
{C'}^{\widehat{\phantom{,}}}_{𝒱}$ on $𝒱$-presheaves, for a
sufficiently large universe $𝒱$, is fully faithful
(\hyperrefext[I][I.rem5.7]{5.7}). Consequently, by the commutativity
of \eqhyperref[III.eq2.3.1]{2.3.1} and
\eqhyperref[III.eq2.3.2]{2.3.2}, the functor
$u_{\ast}:C^{\sim}\rightarrow \widetilde{C}'$ is fully faithful, and
hence so is $u_{!}$ by what precedes. The converse follows from the
existence of the commutative diagram in \hyperref[III.prop1.2]{1.2}
iv), taking into account the fact that the functors $\epsilon_C$ and
$\epsilon_{C'}$ are fully faithful when the topologies are coarser
than the canonical topology, \cqfd

\begin{enonce*}[remark]{Example 2.7}\etiquette[2.7]{III.exem2.7}
With the notation of \hyperref[III.exem1.9.2]{1.9.2}, the functor
$$
\Phi:\Ouv(Y)\longrightarrow \Ouv(X)
$$
is continuous (\hyperref[III.exem1.9.2]{1.9.2}) and cocontinuous
(\hyperref[III.prop2.2]{2.2}). Moreover, let $i:Y\rightarrow X$ be
the canonical inclusion\pageoriginale and let
$i^{-1}:\Ouv(X)\rightarrow \Ouv(Y)$ be the functor that it defines
(\hyperref[III.exem1.9.1]{1.9.1}). The functor $i^{-1}$ is right
adjoint to the functor $\Phi$. There is therefore a sequence of three
adjoint functors:
\begin{align*}
\Phi^s\quad ,\quad & \Phi_s\\
        & ‖\\
        &\Phi^{\ast}\quad,\quad \Phi_{\ast},
\end{align*}
and canonical isomorphisms $\Phi^{\ast}\simeq (i^{-1})^s$ and
$\Phi_{\ast}\simeq i^{-1}_s$.
\end{enonce*}

% END INCLUDED SGA4_English_current_buildable_source_20260906/source/03/drafts/SGA4_III_LINES_613_861_English.texfrag


% BEGIN INCLUDED SGA4_English_current_buildable_source_20260906/source/03/drafts/SGA4_III_LINES_862_1113_English.texfrag
\section{Induced topology}\etiquette[3]{III.sec3}

\subsection{}\etiquette[3.1]{III.subsec3.1}
Let $C'$ be a site, $C$ a category, and $u:C\rightarrow C'$ a
functor. For every universe $𝒰$ such that $C'$ is a
$𝒰$-site and $C$ is a $𝒰$-small category, denote by
$𝒯_{𝒰}$ the finest among the topologies $T$ on $C$
that make $u$ continuous (\hyperref[III.def1.1]{1.1}). (Such a topology
exists by \hyperrefext[II][II.prop2.2]{2.2}.) The topology $𝒯_{𝒰}$ does
not depend on the universe $𝒰$. Indeed, if $𝒱\supset
𝒰$ is a universe, then $𝒯_{𝒰}=𝒯_{𝒱}$
(\hyperref[III.def1.1]{1.1} and \hyperref[III.prop1.5]{1.5}). The topology
$𝒯_{𝒰}$ is called \emph{the topology induced on $C$ by the topology
  of $C'$ by means of the functor $u$}
  \footnote[1]{When no confusion can result, this topology is called
    the \emph{topology induced on $C$ by the topology of $C'$.}}.

\begin{enonce*}{Proposition 3.2}\etiquette[3.2]{III.prop3.2}
Let $C$ be a small category, $C'$ a $𝒰$-site,
$u:C\rightarrow C'$ a functor, and $𝒯$ the topology on $C$ induced
by $u$. Let $X$ be an object of $C$, and let $R\rightarrow X$ be a
sieve on $X$. The following properties are equivalent:
\begin{enumerate}
\renewcommand{\theenumi}{\roman{enumi}}
\renewcommand{\labelenumi}{\theenumi)}
\item The sieve $R\hookrightarrow X$ is covering for $𝒯$.

\item For every base change $Y\rightarrow X$, where $Y$ is an
  object of $C$, the morphism $u_{!}(R\times_X Y)\rightarrow u(Y)$ is
  bicovering in ${C'}^{\widehat{\phantom{,}}}$ {\rm (\hyperrefext[II][II.def5.2]{5.2}).}
\end{enumerate}
\end{enonce*}

\begin{enonce*}[remark]{Proof}
\begin{itemize} %%% \XXX display problem [2016-11-23]: crossed-out i in i) below
\item[i) $\Rightarrow$ ii):] This follows from axiom (T 1) for
  topologies and from \hyperref[III.prop1.2]{1.2}.

\item[ii) $\Rightarrow$ i):] For every sheaf $F$ on $C'$, the map
$$
\Hom(Y,u^{\ast}F)\longrightarrow \Hom(R\times_X Y, u^{\ast}F)
$$
is\pageoriginale isomorphic, by adjunction (\hyperrefext[I][I.prop5.1]{5.1}), to
the map
$$
\Hom(u(Y),F)\longrightarrow \Hom(u_{!}(R\times_X Y),F),
$$
which is bijective (\hyperrefext[II][II.prop5.3]{5.3}). Consequently
(\hyperrefext[II][II.prop2.2]{2.2}), the sieve $R\hookrightarrow X$ is covering
for $𝒯$.
\end{itemize}
\end{enonce*}

\begin{enonce*}{Corollary 3.3}\etiquette[3.3]{III.cor3.3}
Let $C'$ be a site, $C$ a category, $u:C\rightarrow C'$ a
functor, and $𝒯$ the topology on $C$ induced by the topology of
$C'$. Let $(X_i\rightarrow X)$, $i\in I$, be a family of pullbackable
morphisms in $C$, and suppose that $u$ commutes with fiber products
\footnote[2]{In fact, it suffices that $u$ commute with fiber products
  of the form $\fprod{X_i}{X'}{X}$.}. The following properties are
equivalent:
\begin{enumerate}
\renewcommand{\theenumi}{\roman{enumi}}
\renewcommand{\labelenumi}{\theenumi)}
\item The family $X_i\rightarrow X$, $i\in I$, is covering for
  $𝒯$.

\item The family $(u(X_i)\rightarrow u(X))$, $i\in I$, is covering.
\end{enumerate}
\end{enonce*}

\begin{enonce*}[remark]{Proof}
\begin{itemize}
\item[i) $\Rightarrow$ ii):] This follows from \hyperref[III.prop1.6]{1.6}.

\item[ii) $\Rightarrow$ i):] Let $𝒰$ be a universe such that $C$
  is $𝒰$-small and $C'$ is a $𝒰$-site. Let
  $R\hookrightarrow X$ be the sieve generated by the $X_i\rightarrow
  X$. The presheaf $R$ is the coequalizer of the pair of arrows
(\hyperrefext[I][I.prop2.12]{2.12} and \hyperrefext[I][I.cor3.5]{3.5})
$$
\coprod_{i,j} \fprod{X_i}{X_j}{X}\rightrightarrows \coprod_i X_i
$$
(the coproduct is taken here in $C^{\wedge}$).
Since the functor $u_{!}$ commutes with inductive limits
(\hyperrefext[I][I.prop5.4]{5.4}), the presheaf $u_{!}R$ is the coequalizer of the
pair of arrows
$$
\coprod_{i,j}u(\fprod{X_i}{X_j}{X})\rightrightarrows \coprod_i
u(X_i){\quoi}.
$$
Since the functor $u$ commutes with fiber products, the presheaf
$u_{!}R$ is the coequalizer of the pair of arrows
$$
\coprod_{i,j}\fprod{u(X_i)}{u(X_j)}{u(X)}\rightrightarrows \coprod_i
u(X_i){\quoi},
$$
and\pageoriginale consequently $u_{!}R\rightarrow u(X)$ is the sieve on
$u(X)$ generated by the $u(X_i)\rightarrow u(X)$, $i\in I$.
Using once again the fact that $u$ commutes with fiber products, one
shows that for every base change $Y\rightarrow X$, where $Y$ is an
object of $C$, $u_{!}(\fprod{R}{Y}{X})\rightarrow u(Y)$ is the sieve
on $u(Y)$ generated by the
$\fprod{u(X_i)}{u(Y)}{u(X)}\rightarrow u(Y)$, $i\in I$. Since the
family $u(X_i)\rightarrow u(X)$, $i\in I$, is covering, the family
$\fprod{u(X_i)}{u(Y)}{u(X)}\rightarrow u(Y)$, $i\in I$, is covering.
Thus $u_{!}(\fprod{R}{Y}{X})\rightarrow u(Y)$ is a covering sieve,
and consequently (\hyperref[III.prop3.2]{3.2}) $R\rightarrow X$ is
covering for $𝒯$, \cqfd
\end{itemize}
\end{enonce*}

\begin{enonce*}{Corollary 3.4}\etiquette[3.4]{III.cor3.4}
Let $C'$ be a $𝒰$-site, let $C$ be a full subcategory
of $C'$, and let $u:C\rightarrow C'$ be the inclusion functor. Suppose
that fiber products are representable in $C$ and that $u$
commutes with fiber products. The following conditions are
equivalent:
\begin{enumerate}
\renewcommand{\theenumi}{\roman{enumi}}
\renewcommand{\labelenumi}{\theenumi)}
\item
\begin{enumerate}
\renewcommand{\theenumii}{\alph{enumii}}
\renewcommand{\labelenumii}{\theenumii)}
\item For every object $X$ of $C$, every covering family
  $(Y_j\rightarrow X)$, $j\in J$, in $C'$ is refined by a covering
  family $X_i\rightarrow X$, $i\in I$, where the $X_i$ are objects
  of $C$.

\item There exists a small set $G$ of objects of $C$ such that every
  object of $C$ is the target of a covering family in $C'$ of
  morphisms whose sources belong to $G$.
\end{enumerate}

\item The topology induced by the topology of $C'$ is a
  $𝒰$-topology {\rm (I \hyperref[II.def3.0.2]{3.0.2})}, and the functor $u$ is continuous and
  cocontinuous for this topology.
\end{enumerate}
\end{enonce*}

\begin{enonce*}[remark]{Proof}
This follows immediately from \hyperref[III.cor3.3]{3.3} and
\hyperref[III.def2.1]{2.1}.
\end{enonce*}

\begin{enonce*}{Proposition 3.5}\etiquette[3.5]{III.prop3.5}
Let\pageoriginale $C$ be a $𝒰$-site and let
$\epsilon_C:C\rightarrow C^{\sim}$ be the canonical functor {\rm
(\hyperrefext[II][II.subsubsec4.4.0]{4.4.0})}. Equip $C^{\sim}$ with the canonical
topology. Then the topology of the site $C$ is the topology induced
by the topology of $C^{\sim}$.
\end{enonce*}

\begin{enonce*}[remark]{Proof}
The covering families of $C^{\sim}$ for the canonical topology are
the universally effective epimorphic families
(\hyperrefext[II][II.def2.5]{2.5}), that is (\hyperrefext[II][II.prop4.3]{4.3}), the
epimorphic families. Let $T$ be the topology of the site $C$, and let
$𝒯_{𝒰}$ be the finest among the topologies $T'$ on $C$ such that,
for every sheaf $F$ on $C^{\sim}$, $F\circ\epsilon_C$ is a sheaf for
$T'$. It suffices to show that $T=𝒯_{𝒰}$, for then
$𝒯_{𝒰}$ is a $𝒰$-topology and consequently $𝒯_{𝒰}$ is the induced
topology.
\begin{itemize}
\item[A)] \emph{The topology $T$ is finer than $𝒯_{𝒰}$}:
Let $(X_i\rightarrow X)$, $i\in I$, be a covering family for
$𝒯_{𝒰}$. Then, for every sheaf $F$ on $C^{\sim}$, the map
$F(\epsilon_C(X))\rightarrow\prod_i F(\epsilon_C(X_i))$ is injective,
and consequently (\hyperrefext[II][II.def5.2]{5.2})
$(\epsilon_C X_i\rightarrow \epsilon_CX)$, $i\in I$, is covering for
the canonical topology of $C^{\sim}$, hence
(\hyperrefext[II][II.def5.2]{5.2}) $(X_i\rightarrow X)$, $i\in I$, is covering
for $T$.

\item[B)] \emph{The topology $T$ is coarser than
$𝒯_{𝒰}$}: It suffices to show that
$\epsilon_C:C\rightarrow C^{\sim}$ is continuous
(\hyperref[III.def1.1]{1.1}). Now we shall prove in IV ${\rm n^{\circ}}$ 1
that every sheaf on $C^{\sim}$ is representable, and consequently,
for every sheaf $F$ on $C^{\sim}$, $F\circ\epsilon_C$ is a sheaf on
$C$. (We shall not use \hyperref[III.prop3.5]{3.5} before
\hyperrefext[IV][IV.sec1]{1}.)
\end{itemize}
We record two results that will be used in VI 7.
\end{enonce*}

\begin{enonce*}{Proposition 3.6}\etiquette[3.6]{III.prop3.6}
Let $(C_i)_{i\in I}$ be a family of sites, let $C$ be a category,
and, for every $i\in I$, let $u_i:C_i\rightarrow C$ be a functor.
Let $\sheaf{\Uu}$ be a universe such that the categories $C_i$ and
$C$ are $𝒰$-small. There exists a topology $𝒯_{𝒰}$ on $C$ that is
the coarsest among the topologies for which the $u_i$ are continuous.
The topology $𝒯_{𝒰}$ does not depend on the universe $𝒰$ relative
to which the categories under consideration are small.
\end{enonce*}

The last assertion of \hyperref[III.prop3.6]{3.6} follows from
\hyperref[III.prop1.5]{1.5}.

Let $T$ be a topology on $C$. The functors $u_i$ are continuous if
and only if, for every $i\in I$ and every covering sieve
$R\hookrightarrow X$ on an object $X$ of $C_i$,\pageoriginale the
morphism
$$
u_{i!}(R)\longrightarrow u_i(X)
$$
in $\widehat{C}$ is bicovering (\hyperref[III.prop1.2]{1.2} (ii)). Thus
\hyperref[III.prop3.6]{3.6} is a consequence of the following lemma.

\begin{enonce*}{Lemma 3.6.1}\etiquette[3.6.1]{III.lem3.6.1}
Let $C$ be a small category and let $(u_i:F_i\rightarrow G_i)_{i\in
I}$ be a family of arrows in $\widehat{C}$. There exists on $C$ a
topology that is coarsest among those making the morphisms $u_i$
covering (\resp bicovering) {\rm(\hyperrefext[II][II.def5.2]{5.2})}.
\end{enonce*}

To say that $u:F\rightarrow G$ is covering for a given topology $T$
means that, for every arrow $X\rightarrow G$, with $X\in\Ob C$, the
corresponding arrow $F\times_G X\rightarrow X$ is covering; equivalently,
the family of arrows $X'\rightarrow X$ in $C$ that factor through the
preceding arrow is covering. The existence of a coarsest topology
among those for which the $u_i:F_i\rightarrow G_i$ are covering thus
follows from I \hyperref[II.subsubsec1.1.6]{1.1.6}, proving the assertion of
\hyperref[III.lem3.6.1]{3.6.1} without the parenthetical alternative.
The parenthetical assertion follows by recalling that a morphism
$u:F\rightarrow G$ is bicovering if and only if the morphisms
$u:F\rightarrow G$ and $\diag_u:F\rightarrow F\times_G F$ are
covering.

\begin{enonce*}{Proposition 3.7}\etiquette[3.7]{III.prop3.7}
Let $(C_i)_{i\in I}$ be a family of sites, let $C$ be a category,
for every $i\in I$ let $u_i:C_i\rightarrow C$ be a functor, and let
$\sheaf{\Uu}$ be a universe such that the categories $C_i$ and $C$
are $𝒰$-small. There exists a topology $\mathfrak{T}_{𝒰}$ that is
the finest for which the $u_i$ are cocontinuous. The topology
$\mathfrak{T}_{𝒰}$ does not depend on the universe $𝒰$ relative to
which the categories under consideration are small.
\end{enonce*}

Let $𝒰$ be a universe relative to which the categories under
consideration are small. The functors $u_i$ are cocontinuous for a
topology $\mathfrak{T}$ on $C$ if and only if, for every $i\in I$ and
every sheaf $F$ on $C_i$, the presheaf
$\widehat{u}_{i^{\ast}}F$ is a sheaf for $\mathfrak{T}$
(\hyperref[III.prop2.2]{2.2}). It follows that the topology
$\mathfrak{T}_{𝒰}$ is the finest topology for which the presheaves
$\widehat{u}_{i^\ast}F$, $i\in I$, $F\in C_i^{\sim}$, are sheaves
(\hyperrefext[II][II.prop2.2]{2.2}). The last assertion follows from
\hyperref[III.prop2.2]{2.2}.


% END INCLUDED SGA4_English_current_buildable_source_20260906/source/03/drafts/SGA4_III_LINES_862_1113_English.texfrag


% BEGIN INCLUDED SGA4_English_current_buildable_source_20260906/source/03/drafts/SGA4_III_LINES_1114_1313_English.texfrag
\section{Comparison lemma}\etiquette[4]{III.sec4}\pageoriginale

\begin{enonce*}{Theorem 4.1}\etiquette[4.1]{III.theo4.1}
{\rm(comparison lemma)}. Let $C$ be a small category, let $C'$ be a
site whose underlying category is a $𝒰$-category, and let
$u:C\rightarrow C'$ be a fully faithful functor. Endow $C$ with the
topology induced by $u$ {\rm(\hyperref[III.subsec3.1]{3.1}).} Consider
the following properties:
\begin{enumerate}
\renewcommand{\theenumi}{\roman{enumi}}
\renewcommand{\labelenumi}{\theenumi)}
\item Every object of $C'$ can be covered by objects coming from $C$.

\item The functor $F\mapsto F\circ u$ induces an equivalence from the
category of sheaves on $C'$ to the category of sheaves on $C$.
\end{enumerate}

We always have {\rm i) $\Rightarrow$ ii).} When $C'$ is a $𝒰$-site
and the topology of $C'$ is coarser than the canonical topology, we
have {\rm ii) $\Rightarrow$ i).}
\end{enonce*}

\setcounter{subsection}{1}
\subsubsection{}\etiquette[4.1.1]{III.subsubsec4.1.1}
Let us first prove i) $\Rightarrow$ ii). The proof proceeds in two
steps.

%\emph{ %% causes compilation problems
First step. For every presheaf $H$ on $C'$, the adjunction morphism
$\varphi:u_{!}u^{\ast}H\rightarrow H$ {\rm(\hyperrefext[I][I.prop5.1]{5.1})} is
bicovering {\rm(\hyperrefext[II][II.prop5.3]{5.3})}, and the functor
$u_s:\widetilde{C'}\rightarrow \widetilde{C}$ $\hyperref[III.subsec1.11]{(1.1.1)}$ is fully
faithful.
%} %% cf. above

Let $C/H$ be the small category whose objects are the objects $Y$ of
$C$ equipped with a morphism $u(Y)\rightarrow H$, and whose
morphisms are the commutative diagrams
\[
\xymatrix@=1.25cm{
u(Y) \ar[rr]^{u(m)} \ar[rd]&     &u(Y')\ar[ld]\\
     &  H  & .
}
\]
We have $u^{\ast}H=\varinjlim_{Y\in \ob C/H} Y$
(\hyperrefext[I][I.prop3.4]{3.4}), and hence
(\hyperrefext[I][I.prop5.4]{5.4})
$u_{!}u^{\ast}H=\varinjlim_{Y\in \ob C/H} u(Y)$. The adjunction
morphism is the evident morphism arising from the description of
$u_{!}u^{\ast}H$ as an inductive limit. The morphism $\varphi$ has
the following property:\pageoriginale $(\ast)$ For every object $Y$
of $C$, every morphism $m:u(Y)\rightarrow H$ factors uniquely as the
canonical morphism $\alpha(m):u(Y)\rightarrow u_{!}u^{\ast}H$
followed by $\varphi$. It follows immediately from i) that
$\varphi$ is covering (\hyperrefext[II][II.prop5.1]{5.1}); let us show
that it is bicovering. Let
$p,q:Z\rightrightarrows u_{!}u^{\ast}H$ be two morphisms from an
object $Z$ of $C'$ such that $\varphi p=\varphi q$. For every object
$Y$ of $C$ and every morphism $n:u(Y)\rightarrow Z$, we have
$\varphi pn=\varphi qn$. Property $(\ast)$ then implies that
$pn=qn$, and since the $u(Y)$ cover $Z$, the kernel of $(p,q)$ is a
covering sieve of $Z$. Thus $\varphi$ is bicovering
(\hyperrefext[II][II.prop5.3]{5.3}). For every sheaf $H$ on $C'$,
$u^{\ast}H$ is a sheaf on $C$, denoted by $u_sH$ \hyperref[III.subsec1.11]{(1.1.1)}. Moreover,
$u^su_sH=\underline{a}'u_{!}u_sH$ (\hyperref[III.prop1.3]{1.3}), and
the adjunction morphism $u^su_sH\rightarrow H$ is obtained by
applying the ``associated sheaf'' functor to
$\varphi:u_{!}u^{\ast}H\rightarrow H$. Consequently
(\hyperrefext[II][II.prop5.3]{5.3}), the adjunction morphism
$u^su_sH\rightarrow H$ is an isomorphism. Hence
$u_s:\widetilde{C}'\rightarrow \widetilde{C}$ is fully faithful.

\noindent \emph{Second step. The functor $u$ is cocontinuous and the
functor $u^s:C^{\sim}\rightarrow {C'}^{\sim}$
$(\hyperref[III.prop1.2]{1.2})$ is fully faithful. Consequently
$u_s:{C'}^{\sim}\rightarrow C^{\sim}$ is an equivalence.} Let $Y$
be an object of $C$ and let $i:R\hookrightarrow u(Y)$ be a covering
sieve. Since the functor $u^{\ast}$ commutes with projective limits
(\hyperrefext[I][I.prop5.5]{5.5}), the morphism
$u^{\ast}(i):u^{\ast}(R)\hookrightarrow u^{\ast}u(Y)$ is a
monomorphism. Since $u$ is fully faithful, we have
$u^{\ast}u(Y)\simeq Y$, and thus obtain a sieve of $Y$,
$u^{\ast}(i):u^{\ast}(R)\rightarrow Y$. To show that $u$ is
cocontinuous, it suffices to show that
$u^{\ast}(i):u^{\ast}(R)\hookrightarrow Y$ is a covering sieve for
the induced topology on $C$ (\hyperref[III.def2.1]{2.1}). For this,
it suffices to show (\hyperref[III.prop3.2]{3.2}) that for every base
change $m:X\rightarrow Y$, the morphism
$u_{!}(u^{\ast}(R)\times_Y X)\rightarrow u(X)$ is bicovering. But
since $u^{\ast}$ commutes with\pageoriginale projective limits, we
have
$\fprod{u^{\ast}(R)}{X}{Y}=u^{\ast}(\fprod{R}{u(X)}{u(Y)})$, and
$\fprod{R}{u(X)}{u(Y)}$ is a covering sieve of $u(X)$. It therefore
suffices to show that for every object $Y$ of $C$ and every covering
sieve $i:R\hookrightarrow u(Y)$, the morphism in
${C'}^{\widehat{\phantom{,}}}$,
$u_{!}u^{\ast}(R)\xrightarrow{u_{!}u^{\ast}(i)}u(Y)$, is
bicovering. Now this morphism factors as the adjunction morphism
$u_{!}u^{\ast}(R)\rightarrow R$, which is bicovering (first step),
followed by the covering monomorphism $R\hookrightarrow u(Y)$. It is
therefore bicovering (\hyperrefext[II][II.prop5.3]{5.3}). This shows
that $u$ is cocontinuous. Since $u$ is continuous and fully
faithful, it follows from \hyperref[III.prop2.6]{2.6} (used in the
case where $C$ is small) that $u^s$ (denoted by $u_{!}$ in
\hyperref[III.prop2.6]{2.6}) is fully faithful. Since $u^s$ and
$u_s$ are adjoint to one another and are fully faithful, they are
quasi-inverse functors, and hence $u_s$ is an equivalence.

\setcounter{subsection}{1}
\subsubsection{}\etiquette[4.1.2]{III.subsubsec4.1.2}
Let us now prove ii) $\Rightarrow$ i). The objects $Y$ of $C$ form a
generating family of $C^{\sim}$ (\hyperrefext[II][II.prop4.10]{4.10}),
and consequently, for every object $X$ of $C'$, there exists an
epimorphic family (in $C^{\sim}$)
$v_i:Y_i\rightarrow u_s X$, $i\in I$. We have
$u_su(Y_i)=Y_i$, and, since $u_s$ is an equivalence of categories,
$v_i=u_s(w_i)$, where the $w_i:u(Y_i)\rightarrow X$ form an
epimorphic family in ${C'}^{\sim}$. It then follows from
\hyperrefext[II][II.prop5.1]{5.1} that the family
$w_i:u(Y_i)\rightarrow X$, $i\in I$, is covering.

\subsection{Completion of the proof of \hyperref[III.prop1.2]{1.2}}\etiquette[4.2]{III.subsec4.2}

Let $G$ be a full subcategory of $C$ whose objects form a small
topologically generating family of $C$ (II \hyperref[II.def3.0.1]{3.0.1}). Endow $G$ with
the topology induced by the inclusion functor $i:G\rightarrow C$
(\hyperref[III.subsec3.1]{3.1}). The functor
$(u\circ i)_s=i_s\circ u_s$ \hyperref[III.subsec1.11]{(1.1.1)} has a left adjoint by the first
part of the proof of \hyperref[III.prop1.2]{1.2}. Since $i_s$ is an
equivalence of\pageoriginale categories (\hyperref[III.theo4.1]{4.1}),
the functor $u_s$ has a left adjoint $u^s$, and there is a functorial
isomorphism $u^s\simeq (u\circ i)^s\circ i_s$. Let us show that the
diagram
\hypertarget{sga4.exp.III.equation.4.2.1}{}\[\vcenter{
\xymatrix@C=3cm@R=1.5cm{
C \ar[r]^{u} \ar[d]_{\epsilon_C}& C'\ar[d]^{\epsilon_{C'}}\\
C\widetilde{\phantom{a}} \ar[r]^{u^{s}}& \widetilde{C}'
}}\leqno{(4.2.1)}
\]
commutes up to isomorphism. For every sheaf $H$ on $C'$, we have
$\Hom_{\widetilde{C}'}(\epsilon_{C'}uX,H)\simeq u_sH(X)$ for every
$X\in \ob C$. Moreover,
$\Hom_{\widetilde{C}'}(u^s\epsilon_CX,H)\simeq
\Hom_{\widetilde{C}}(\epsilon_CX,u_sH)$ by adjunction, and then
$\Hom_{\widetilde{C}}(\epsilon_CX,u_sH)\simeq u_sH(X)$ by the
definition of $\epsilon_C$. Hence there is an isomorphism
$\epsilon_{C'}uX\simeq u^s\epsilon_CX$ for every $X\in \ob C$. This
proves i) $\Rightarrow$ iv). Let us prove iv) $\Rightarrow$ i). We
have a commutative diagram
\hypertarget{sga4.exp.III.equation.4.2.2}{}\[\vcenter{
\xymatrix@C=2.25cm@R=1.5cm{ G \ar[r]^{i} \ar[d]_{\epsilon_G} & C
\ar[r]^{u} \ar[d]_{\epsilon_C}&
C' \ar[d]_{\epsilon_{C'}}\\
\widetilde{G} \ar[r]^{i^{s}}& \widetilde{C} \ar[r]^{u^s}&
\widetilde{C'}, }}\leqno{(4.2.2)}
\]
and consequently, by the first part of the proof, the functor
$u\circ i:G\rightarrow C'$ is continuous. It follows at once from
\hyperref[III.theo4.1]{4.1} and \hyperref[III.def1.1]{1.1} that $u$
is continuous, whence iv) $\Rightarrow$ i). This also gives the
uniqueness of $u^s$, since the first part of the proof gives
uniqueness when $C$ is small.

\subsection{Completion of the proofs of \hyperref[III.prop1.5]{1.5} and \hyperref[III.prop1.7]{1.7}}\etiquette[4.3]{III.subsec4.3}
\pageoriginale

Let $u:C\rightarrow C'$ be a functor between two $𝒰$-sites, and let
$G$ be a full subcategory of $C$ whose set of objects is a small
topologically generating family of $C$ (II \hyperref[II.def3.0.1]{3.0.1}). Endow $G$ with
the topology induced by the inclusion functor $i:G\rightarrow C$.
It follows immediately from \hyperref[III.theo4.1]{4.1} and
\hyperref[III.def1.1]{1.1} that $u$ is continuous if and only if
$u\circ i$ is continuous. This observation and the first part of
the proof of \hyperref[III.prop1.5]{1.5} give the general case. The
same observation also reduces the proof of \hyperref[III.prop1.7]{1.7}
to the case where the category $C$ is small.

\subsection{Completion of the proof of \hyperref[III.prop2.3]{2.3}}\etiquette[4.4]{III.subsec4.4}

Let $u:C\rightarrow C'$ be a functor between two $𝒰$-sites, and let
$G$ be a full subcategory of $C$ whose set of objects is a small
topologically generating family of $C$
(\hyperrefext[II][II.def3.0.1]{3.0.1}), endowed with the topology
induced by the inclusion functor $i:G\rightarrow C$. It follows from
the proof of \hyperref[III.theo4.1]{4.1} (\hyperref[III.subsubsec4.1.1]{4.1.1}, second step) that
$i$ is cocontinuous. Consequently, when $u$ is cocontinuous, the
functor $u\circ i:G\rightarrow C'$ is cocontinuous. Moreover,
$i^{\ast}:\widetilde{C}\rightarrow \widetilde{G}$ is an equivalence
of categories (\hyperref[III.theo4.1]{4.1}). Therefore the functor
$u^{\ast}:\widetilde{C}'\rightarrow \widetilde{C}$ has a right
adjoint. This proves 3). To prove 4), it suffices to observe that
$(u\circ i)_{\ast}\simeq u_{\ast}\circ i_{\ast}$ and that
$i_{\ast}$ is an equivalence (\hyperref[III.theo4.1]{4.1}). The
commutativity of \eqhyperref[III.eq2.3.2]{2.3.2} then follows from
the commutativity of these diagrams when $C$ is small.


% END INCLUDED SGA4_English_current_buildable_source_20260906/source/03/drafts/SGA4_III_LINES_1114_1313_English.texfrag


% BEGIN INCLUDED SGA4_English_current_buildable_source_20260906/source/03/drafts/SGA4_III_LINES_1314_1551_English.texfrag
\section{Localization}\etiquette[5]{III.sec5}\pageoriginale

\subsection{}\etiquette[5.1]{III.subsec5.1}
Now let $C$ be a $𝒰$-site and $X$ an object of
$C^{\widehat{\phantom{,}}}$. Unless expressly stated otherwise, the
category $C/X$ will be endowed with the topology $𝒯$ induced by the
functor $j_X:C/X\rightarrow C$ (\hyperref[III.subsec3.1]{3.1}). The notation
$C/X$ will denote the category $C/X$ endowed with the topology $𝒯$.
Proposition \hyperrefext[I][I.prop5.11]{5.11} shows that the functor $j_{X!}$
commutes with fiber products and consequently sends every
monomorphism to a monomorphism. In particular, for every object
$(Y\rightarrow X)$ of $C/X$, the functor $j_{X!}$ establishes a
one-to-one correspondence between the sieves, in the category
$C/X$, on the object $(Y\rightarrow X)$ and the sieves, in $C$, on
the object $Y$.

\begin{enonce*}{Proposition 5.2}\etiquette[5.2]{III.prop5.2}
Let $C$ be a $𝒰$-site, $X$ an object of
$C^{\widehat{\phantom{,}}}$, and $j_X:C/X\rightarrow C$ the
corresponding continuous functor.
\begin{enumerate}
\renewcommand{\labelenumi}{\theenumi)}
\item A sieve $R$ on an object $(Z\rightarrow X)$ is covering in
$C/X$ if and only if the sieve $j_{X!}(R)\hookrightarrow Z$ is
covering in $C$.

\item The functor $j_X:C/X\rightarrow C$ is cocontinuous and
continuous.

\item Let $m:Y\rightarrow X$ be a morphism of
$C^{\widehat{\phantom{,}}}$. We then have the commutative diagram
\[
\xymatrix@C=1cm@R=1.25cm{
C/Y \ar[rr]^{j_m} \ar[rd]_{j_Y}&       &C/X \ar[ld]^{j_X}\\
    &  C    &.
}
\]
The topology induced by the functor $j_Y$ on $C/Y$ is equal to the
topology induced by $j_m$ on $C/Y$.

\item The topology induced by $j_X:C/X\rightarrow C$ is a
$𝒰$-topology (\hyperrefext[II][II.def3.0.2]{3.0.2}).
\end{enumerate}
\end{enonce*}

\begin{enonce*}[remark]{Proof}
\begin{enumerate}
\renewcommand{\labelenumi}{\theenumi)}
\item If\pageoriginale the sieve $R$ on $(Z\rightarrow X)$ is
covering, the sieve $j_{X!}(R)\hookrightarrow Z$ is covering
(\hyperref[III.prop1.6]{1.6}).

Conversely, if the sieve $j_{X!}(R)\rightarrow Z$ is covering in
$C$, one sees that the same is true for every sieve obtained by a
base change in $C/X$. The sieve $R$ is therefore covering
(\hyperref[III.prop3.2]{3.2}).

\item This follows immediately from 1) by applying
\hyperref[III.def2.1]{2.1}.

\item This follows immediately from the description of covering
sieves given in 1).

\item Let $(G_i)_{i\in I}$ be a small topologically generating
family of $C$. One verifies immediately that the small family
$(u:G_i\rightarrow X)$, $u\in \amalg_{i\in
I}\Hom_{C^{\widehat{\phantom{,}}}}(G_i, X)$, is topologically
generating in $C/X$.
\end{enumerate}
\end{enonce*}

\phantomsection\MakeLinkTarget*{sga4.exp.III.terminology.5.3}\label{sga4.exp.III.terminology.5.3}\label{III.subsec5.3}%
\noindent \emph{Terminology and notation} 5.3. By the preceding
proposition, the functor $j_X$ is both continuous and cocontinuous.
It therefore defines a sequence of three adjoint functors (4.3.2)
between the categories of sheaves of sets (\hyperref[III.prop1.3]{1.3}
and \hyperref[III.prop2.3]{2.3}):
\begin{gather*}
j^s_X:(C/X)^{\sim}\longrightarrow
C^{\sim}\\
j^{\ast}_X=j_{X,s}:C^{\sim}\longrightarrow
(C/X)^{\sim}\\
j_{X\ast}:(C/X)^{\sim}\longrightarrow C^{\sim}{\quoi}.
\end{gather*}

In the particular situation of Proposition \hyperref[III.prop5.2]{5.2},
we shall use the following terminology and notation:
\begin{enumerate}
\renewcommand{\labelenumi}{\theenumi)}
\item The functor $j_{X\ast}$ will be called the \emph{direct image
functor.}

\item The functor $j_{X,s}=j^{\ast}_X$ will be denoted by
$j^{\ast}_X$ and will be called the \emph{restriction functor to
$C/X$.}

\item The\pageoriginale functor $j^s_X$ on sheaves of sets will be
called the ``extension by the empty set to $C$''
functor and will be denoted by $j_{X!}$ (\cf 2.9.2).
\end{enumerate}

Thus we have a sequence of three adjoint functors between
$(C/X)^{\sim}$ and $C^{\sim}$:
$$
j_{X!},\; j^{\ast}_X,\; j_{X\ast}{\quoi}.
$$

\begin{enonce*}{Proposition 5.4}\etiquette[5.4]{III.prop5.4}
The functor $j_{X!}:(C/X)^{\sim}\rightarrow C^{\sim}$ factors
through the category $C^{\sim}/\underline{a}X$ ($\underline{a}$ is
the associated sheaf functor):

$(C/X)^{\sim}\xrightarrow{e^{\sim}_X} C^{\sim}/\underline{a}X\rightarrow
C^{\sim}$. The functor
$$
e^{\sim}_X:(C/X)^{\sim}\longrightarrow \widetilde{C}/\underline{a}X
$$
is an equivalence of categories. The restriction functor to $C/X$,
followed by the equivalence $e^{\sim}_X$, is isomorphic to the
``base change by $\underline{a}X\rightarrow \epsilon$'' functor
($\epsilon$ being the terminal object of $C^{\sim}$):
$F\mapsto (F\times
\underline{a}X\xrightarrow{\pr_2}\underline{a}X)$.
\end{enonce*}

\begin{enonce*}[remark]{Proof}
The image under $j_{X!}$ of the terminal object of $(C/X)^{\sim}$ is
the object $\underline{a}X$, which gives the factorization. To show
that the functor $e^{\sim}_X$ is an equivalence, we shall confine
ourselves to a few indications. By \hyperrefext[I][I.prop5.11]{5.11}, a
presheaf on $C/X$ is defined by a presheaf $F$ on $C$ equipped with
a morphism $F\rightarrow X$. One then proves that the following two
properties are equivalent:
\begin{enumerate}
\renewcommand{\theenumi}{\roman{enumi}}
\renewcommand{\labelenumi}{\theenumi)}
\item The presheaf on $C/X$ defined by $F\rightarrow X$ is a
sheaf.

\item The following diagram is cartesian (where $i$ and
$\underline{a}$ denote the inclusion in the category of presheaves
and the associated sheaf functor, respectively):
\[
\xymatrix@C=2.75cm@R=1cm{
F \ar[r]\ar[d]& i\mathsf{\underline{a}}F \ar[d]\\
X \ar[r]& i\mathsf{\underline{a}}X\quad . }
\]
\end{enumerate}

It\pageoriginale follows immediately that $e^{\sim}_X$ is an
equivalence. The last assertion is trivial.
\end{enonce*}

\begin{enonce*}{Proposition 5.5}\etiquette[5.5]{III.prop5.5}
\begin{enumerate}
\renewcommand{\labelenumi}{\theenumi)}
\item Let $C$ be a $𝒰$-site and $X$ an object of
$C^{\widehat{\phantom{,}}}$. The diagram below of categories and
functors commutes up to canonical isomorphisms:
\[
\xymatrix@C=1.5cm@R=1.5cm{ C/X \ar[r]^{h_X} \ar[d]&
(C/X)\widehat{\phantom{a}} \ar[d]^{e\widehat{\phantom{a}}_X}
\ar[r]^{\underline{a}_X}& (C/X)\widetilde{\phantom{a}} \ar[rr]^{i_X}
\ar[d]^{e\widetilde{\phantom{a}}_X}& &
(C/X)\widehat{\phantom{a}} \ar[d]^{e\widehat{\phantom{a}}/X}\\
C/X \ar[r]^{h/X}& C\widehat{\phantom{a}}/X \ar[r]^{\underline{a}/X}&
C\widetilde{\phantom{a}}/\underline{a}X \ar[r]^{i/\underline{a}X}&
C\widehat{\phantom{a}}/i\underline{a}X \ar[r]^{\pi}&
C\widehat{\phantom{a}}/X. }
\]

The notations $\underline{a}_X$ and $i_X$ denote the ``associated
sheaf'' and inclusion into presheaves functors for the site $C/X$.
The notation $\underline{a}/X$ denotes the natural extension of the
functor $\underline{a}$ (the associated sheaf functor for the site
$C$) to the category of arrows with target $X$; similarly for the
notation $i/\underline{a}X$. Finally, the notation $\pi$ denotes the
base-change functor along the canonical arrow
$X\rightarrow i\underline{a}X$.

\item Let, in addition, $m:Y\rightarrow X$ be a morphism of
$C^{\widehat{\phantom{,}}}$. The following diagram commutes up to
canonical isomorphism:
\[
\xymatrix@C=1.75cm@H=.5cm{ C\widetilde{\phantom{a}}
/\underline{a}X/\underline{a}Y \ar@{=}[dd]&
(C/X)\widetilde{\phantom{a}}/\underline{a}_X Y \ar[l]_{f}& \\
& & (C/X/Y)\widetilde{\phantom{a}} \ar[lu]_{g} \ar[ld]_{e\widetilde{\phantom{a}}_m}\\
C\widetilde{\phantom{a}}/\underline{a}Y & (C/Y)\widetilde{\phantom{a}}
\ar[l]_{e\widetilde{\phantom{a}}_Y}&. }
\]

The\pageoriginale arrow $f$ is the natural extension of the
equivalence $e^{\sim}/X$ to the category of arrows with target
$\underline{a}_XY$, and is consequently an equivalence of
categories. The arrow $g$ is precisely the equivalence of
\hyperref[III.prop5.4]{5.4} applied to the situation $(C/X)/[m]$.

\item Denote by
$j_{\underline{a}X!}:C^{\sim}/\underline{a}X\rightarrow C^{\sim}$ the
``forget $\underline{a}X$'' functor, and by
$j^{\ast}_{\underline{a}X}:C^{\sim}\rightarrow C^{\sim}/\underline{a}X$ the
``product with $\underline{a}X$'' functor. The following diagrams
commute up to canonical isomorphism:
\begin{displaymath}
\xymatrix@C=2.25cm@R=1.25cm{ (C/X)^{\widetilde{\phantom{a}}}
\ar[r]^{j_{X!}} \ar[d]   &
  C\widetilde{\phantom{a}} \ar@{=}[d]\\
C\widetilde{\phantom{a}}/\underline{a}X \ar[r]^{j_{\underline{a}X!}} &
C\widetilde{\phantom{a}} } \qquad {\raisebox{-1.75cm}{,}}\qquad
\xymatrix@C=2.5cm@R=1.25cm{
C{\widetilde{\phantom{a}}} \ar[r]^{j_X^{\ast}} \ar@{=}[d]&  (C/X)\widetilde{\phantom{a}}\ar[d]^{e\widetilde{\phantom{a}}/X}\\
C{\widetilde{\phantom{a}}} \ar[r]^{j_{\underline{a}X}^{\ast}} &
C{\widetilde{\phantom{a}}}/\underline{a}X. }
\end{displaymath}
\end{enumerate}
\end{enonce*}

\begin{enonce*}[remark]{Proof}
The proof of these assertions follows immediately from the
definition of the equivalences $e^{\sim}_X$
(\hyperref[III.prop5.4]{5.4}) and
$e^{\widehat{\phantom{\,}}}_X$ (\hyperrefext[I][I.prop5.11]{5.11}).
\end{enonce*}

\begin{thebibliography}{99}%
\bibitem[TF]{sga4.exp.III.bib.TF} R. Godement, théorie des faisceaux, Hermann,
1958. Act. Scient.

Ind ${\rm n^{\circ}}$ 1252, (Paris).
\end{thebibliography}


%I didn't get the label...
%---\widetilde{C}$ $(1.1.1)$ is fully faithful.}
%---\emph{Terminology and notation} 5.3. By the proposition
%---4.3.2
%---2.9.2

\ifx\danslelivre\undefined
\end{document}
\fi

% END INCLUDED SGA4_English_current_buildable_source_20260906/source/03/drafts/SGA4_III_LINES_1314_1551_English.texfrag


% END INCLUDED SGA4_English_current_buildable_source_20260906/source/03/03_English.tex


\makeatletter
\def\input@path{{../04/}}
\makeatother
% BEGIN INCLUDED SGA4_English_current_buildable_source_20260906/source/04/04_English.tex
\ifx\danslelivre\undefined
% Release engine: XeLaTeX (the LuaTeX compatibility layer breaks Xy-pic output).
\documentclass[10pt]{../fichiers-configuration-latex/smfart}
% BEGIN INCLUDED SGA4_English_current_buildable_source_20260906/source/fichiers-configuration-latex/paquets.tex
\usepackage{polyglossia}
%\usepackage[francais]{babel}
\usepackage{fontspec}
\usepackage[math-style=ISO,bold-style=upright,sans-style=upright]{unicode-math}
%\usepackage[mathletters]{ucs}
%\usepackage[utf8]{inputenc}
%\usepackage[T1]{fontenc}
\usepackage[unicode,hypertexnames=false,colorlinks=true,linkcolor=blue,citecolor=Sepia,urlcolor=magenta]{hyperref}
\setmainlanguage{french}
\setmainfont[Ligatures=TeX]{Linux Libertine O}
% variante possible : TeXGyrePagella ?
\setmathfont{XITS Math}

% maths
\usepackage{amsmath}
%\usepackage{amssymb}
%\usepackage{mathrsfs}
%\usepackage{stmaryrd}
%\usepackage{euscript}

% couleurs
\usepackage[usenames,dvipsnames]{xcolor}

% graphiques
\usepackage{graphicx}
\usepackage[all]{xy}
%\usepackage{tikz}
%\usetikzlibrary{matrix,arrows,calc}
%\tikzset{column sep=3em,row sep=2em,text height=1.5ex, text depth=0.25ex,>=angle 90}
%\tikzset{description/.style={fill=white,inner sep=1pt}}

% hyperliens, 'etiquettes
\usepackage{xr-hyper}
\usepackage{../fichiers-configuration-latex/smfthm}
%\ifx\cacheretiquettes\undefined
%\usepackage[inline]{showlabels} \else \fi % pour rendre visibles les 'etiquettes
%\usepackage{../fichiers-configuration-latex/makeidx} % pour faire un index
%\usepackage[pagebackref,unicode]{hyperref}  % pour voir o`u sont cit'ees les r'ef'erences
% hyperref is loaded once above with the cumulative-reader destination policy.
%\hypersetup{breaklinks=true} %permet le retour `a la ligne dans les liens trop longs. Incompatible avec pdflatex
%\renewcommand*{\backref}[1]{{\small [cit\'e page(s)~#1.]}}

\usepackage{xspace}
\usepackage{pifont} % pour expose XVII

% END INCLUDED SGA4_English_current_buildable_source_20260906/source/fichiers-configuration-latex/paquets.tex

% BEGIN INCLUDED SGA4_English_current_buildable_source_20260906/source/fichiers-configuration-latex/numerotation.tex
\renewcommand{\thefootnote}{\roman{footnote}} % pour mettre les notes en bas de page sous la forme (ii) et non (2), qui ressemble trop (?) a un exposant.

\newcounter{pageoriginal}
\newcommand{\pageoriginale}{\refstepcounter{pageoriginal}\marginpar{\footnotesize\textbf{\thepageoriginal}}\xspace}
\let\pageoriginaled\pageoriginale

\usepackage{ifthen}
\newboolean{original}
\def\sisi{\ifthenelse{\boolean{original}}}
\setboolean{original}{false}

\newcounter{cnote}
\def\nde#1{\ifthenelse{\boolean{original}}{}{\stepcounter{cnote}\def\thefootnote{\arabic{cnote}}\protect\footnote[\value{cnote}]{N.D.E.: #1}}}
\def\ndemark{\ifthenelse{\boolean{original}}{}{\stepcounter{cnote}\def\thefootnote{\arabic{cnote}}\protect\footnotemark[\value{cnote}]}}
\def\ndetext#1{\ifthenelse{\boolean{original}}{}{\def\thefootnote{\arabic{cnote}}\protect\footnotetext[\value{cnote}]{N.D.E.: #1}}}

% END INCLUDED SGA4_English_current_buildable_source_20260906/source/fichiers-configuration-latex/numerotation.tex

% BEGIN INCLUDED SGA4_English_current_buildable_source_20260906/source/fichiers-configuration-latex/macros.tex
\setmainlanguage{french}
\setotherlanguage{english}


%% Fixed, engine-independent release identifier.  The inherited dynamic Git
%% probe used Lua io.popen, exposed build-state-dependent output, and printed
%% a visible question mark outside a Git checkout.
\def\showgitstatus{English release r2, 2026-07-23}


% macros pour LaTeX
\renewcommand{\mathsf}{}

% Police différente pour U+2113 SCRIPT SMALL L
\setmathfont[range="2113]{latinmodern-math.otf}

% Give every inherited semantic label its own stable PDF destination instead
% of inheriting whichever automatic equation/page anchor happened to be live.
% The legacy label remains as a compatibility alias used by existing edges.
\makeatletter
\def\etiquette[#1]#2{%
  \Hy@raisedlink{\hyper@anchorstart{sga4.semantic.#2}\hyper@anchorend}%
  \def\@currentHref{sga4.semantic.#2}%
  \def\@currentlabel{#1}%
  \label{#2}%
}
\makeatother
\def\eqhyperref[#1]#2{(\hyperref[#1]{#2})}
\def\pagehyperref[#1]#2{(page \hyperref[#1]{#2})}
\def\hyperrefext[#1][#2]#3{\hyperref[#2]{#1 #3}}
\def\eqhyperrefext[#1][#2]#3{\hyperref[#2]{#1 #3}}     %cf. exposes 17
\def\neweqhyperrefext[#1][#2]#3{\hyperref[#2]{#1 #3}}  %cf. exposes 17

\def\shat#1{\protect\widehat{#1}}
\def\stilde#1{\protect\widetilde{#1}}
%%%%%%%%%%%%%%%%%%


\def\etcn{etc.\xspace}
\def\etco{etc...\xspace}
\def\etc{\sisi{\etco}{\etcn}}
\def\ldotsn{\sisi{...}{\etcn}}

\def\cf{{cf.}\kern.3em}
\def\Cf{{Cf.}\kern.3em}
\def\ie{{i.e.}\kern.3em}
\def\iev{{i.e.},\xspace}
\def\loccito{loc.\kern2pt cit.\xspace}
\def\loccitn{\textit{loc.\kern2pt cit.}\xspace}
\def\loccit{\sisi{\loccito}{\loccitn}}
\def\EGA{\textup{EGA}\kern.2em}
\def\SGA{\textup{SGA}\kern.3em}

\def\quoi{}

\def\quav{\quad,\quad}
\def\quaet{\quad \text{et}\quad}
\def\fibreproduct#1#2#3{#1 ×_{#3} #2}
\let\fprod\fibreproduct
\def\ccheck#1#2{\mbox{${\check{\check{\rm #1}}}^{\lower2pt\hbox{$\scriptstyle #2$}}$}}
\def\sheaf#1{\protect\mathcal{#1}}
\def\vieuxsheaf#1{{#1}} % cf. problème avec sheaf{a\Phi_i} par exemple...
\DeclareMathOperator{\Hom}{Hom}
\DeclareMathOperator{\can}{can}
\DeclareMathOperator{\Ind}{Ind}
\DeclareMathOperator{\ind}{ind}
%\DeclareMathOperator{\Hhomtop}{\mathcal{H}\!\mathit{omtop}}
\DeclareMathOperator{\Hhomtop}{\mathbf{Homtop}}
\DeclareMathOperator{\Point}{Point}
\DeclareMathOperator{\Pro}{Pro}
\DeclareMathOperator{\Top}{Top}
\DeclareMathOperator{\Mod}{Mod}
\DeclareMathOperator{\Cart}{Cart}
%\DeclareMathOperator{\Mmorsite}{\mathcal{M}\!\mathit{orsite}}
\DeclareMathOperator{\Mmorsite}{\mathbf{Morsite}}
\DeclareMathOperator{\loc}{loc}
\DeclareMathOperator{\Loc}{Loc}
\DeclareMathOperator{\Pic}{Pic}
\DeclareMathOperator{\Et}{Et}
\DeclareMathOperator{\card}{card}
\DeclareMathOperator{\Ens}{Ens}
\DeclareMathOperator{\h}{\underline{H}}
\DeclareMathOperator{\Fonct}{Fonct}
\DeclareMathOperator{\Ker}{Ker}
\DeclareMathOperator{\but}{but}
\DeclareMathOperator{\Ob}{Ob}
\DeclareMathOperator{\ob}{ob}
\DeclareMathOperator{\Lib}{Lib}
\DeclareMathOperator{\esj}{esj}
\DeclareMathOperator{\id}{id}
\DeclareMathOperator{\Id}{Id}
\DeclareMathOperator{\ad}{ad}
\DeclareMathOperator{\Iim}{Im}
\DeclareMathOperator{\pr}{pr}
\DeclareMathOperator{\Cat}{Cat}
\DeclareMathOperator{\car}{car}
\DeclareMathOperator{\SysInd}{SysInd}
\DeclareMathOperator{\Sup}{Sup}
\DeclareMathOperator{\Coker}{Coker}
\DeclareMathOperator{\Filt}{Filt}
\DeclareMathOperator{\acc}{acc}
\DeclareMathOperator{\sur}{sur}
\DeclareMathOperator{\cat}{cat}
\DeclareMathOperator{\Ord}{Ord}
\DeclareMathOperator{\Gr}{Gr}
\DeclareMathOperator{\Cov}{Cov}
\DeclareMathOperator{\Caf}{Caf}
\DeclareMathOperator{\coker}{coker}
\DeclareMathOperator{\coim}{coim}
\DeclareMathOperator{\im}{im}
\DeclareMathOperator{\Ouv}{Ouv}
\DeclareMathOperator{\diag}{diag}
\DeclareMathOperator{\Esp}{Esp}
\DeclareMathOperator{\TOP}{TOP}
\DeclareMathOperator{\Unif}{Unif}
\DeclareMathOperator{\Iind}{\mathbf{Ind}}
\DeclareMathOperator{\Hhomcart}{\mathbf{Homcart}}
\DeclareMathOperator{\Hhomtopess}{\mathbf{Homtopess}}
\DeclareMathOperator{\Ppoints}{\mathbf{Points}}
\DeclareMathOperator{\Ppoint}{\mathbf{Point}}
\DeclareMathOperator{\Ttors}{\mathbf{Tors}}
\DeclareMathOperator{\Ssex}{\mathbf{Sex}}
\DeclareMathOperator{\Ffib}{\mathbf{Fib}}
\DeclareMathOperator{\Vvois}{\mathbf{Vois}}
\DeclareMathOperator{\Ppointess}{\mathbf{Pointess}}
\DeclareMathOperator{\scrHom}{\mathbf{Hom}}
\DeclareMathOperator{\sheafPic}{\mathbf{Pic}}
\DeclareMathOperator{\sheafAut}{\mathbf{Aut}}
\DeclareMathOperator{\Hhomtoplocan}{\mathbf{Homtoplocan}}
\DeclareMathOperator{\Hhomtopan}{\mathbf{Homtopan}}
\DeclareMathOperator{\Pptgeom}{\mathbf{Ptgéom}}
\DeclareMathOperator{\Morsite}{Morsite}
\DeclareMathOperator{\Fib}{Fib}
\DeclareMathOperator{\Fais}{Fais}
\DeclareMathOperator{\kar}{kar}
\DeclareMathOperator{\Groupes}{Groupes}
\DeclareMathOperator{\Restr}{Restr}
\DeclareMathOperator{\Isom}{Isom}
\DeclareMathOperator{\Prol}{Prol}
\DeclareMathOperator{\Res}{Res}
\DeclareMathOperator{\Tors}{Tors}
\DeclareMathOperator{\esp}{esp}
\DeclareMathOperator{\es}{es}
\DeclareMathOperator{\od}{od}
\DeclareMathOperator{\gr}{gr}
\DeclareMathOperator{\sob}{sob}
\DeclareMathOperator{\dfn}{dfn}
\DeclareMathOperator{\Sch}{Sch}
\DeclareMathOperator{\ouv}{ouv}
\DeclareMathOperator{\cons}{cons}
\DeclareMathOperator{\Nerf}{Nerf}
\DeclareMathOperator{\Points}{Points}
\DeclareMathOperator{\Espsob}{Espsob}
\DeclareMathOperator{\Homtop}{Homtop}
\DeclareMathOperator{\Pt}{Pt}
\DeclareMathOperator{\Quot}{Quot}
\DeclareMathOperator{\glob}{glob}
\DeclareMathOperator{\Homfib}{Homfib}
\DeclareMathOperator{\Qouv}{Qouv}
\DeclareMathOperator{\Quotouv}{Quotouv}
\DeclareMathOperator{\ens}{ens}
\DeclareMathOperator{\ab}{ab}
\DeclareMathOperator{\an}{an}
\DeclareMathOperator{\fs}{fs}
\DeclareMathOperator{\red}{red}
\DeclareMathOperator{\Ab}{Ab}
\DeclareMathOperator{\cosq}{cosq}
\DeclareMathOperator{\Ann}{Ann}
\DeclareMathOperator{\Zar}{Zar}


%%%%%%%%%%%%
%%%%% SGA4-2
\DeclareMathOperator{\Sex}{Sex}
\DeclareMathOperator{\vois}{vois}
\DeclareMathOperator{\Vois}{Vois}
\DeclareMathOperator{\Ext}{Ext}
\DeclareMathOperator{\ext}{ext}
\DeclareMathOperator{\Fer}{Fer}
\DeclareMathOperator{\Add}{Add}
\DeclareMathOperator{\Cocart}{Cocart}
\DeclareMathOperator{\cosk}{cosk}
\DeclareMathOperator{\plat}{plat}

\DeclareMathOperator{\faisc}{faisc}
\DeclareMathOperator{\sk}{sk}
\DeclareMathOperator{\ch}{ch}

\def\HH{\mathbf{H}}
\def\qc{\mathrm{qc}}
\def\et{\mathrm{et}}
\def\coh{\mathrm{coh}}

\def\Hvr#1{%
\begin{array}[b]{l}
{\kern2pt\scriptstyle r}\\[-8pt]
{\kern1pt\scriptstyle \vee}\\[-5pt]
H^{#1}
\end{array}}

\def\Hvv#1{%
\begin{array}[b]{l}
\kern2pt\rotatebox{90}{${\scriptscriptstyle\ll}$}\\[-6pt]
H^{#1}
\end{array}}

\def\Hvi#1{%
\begin{array}[b]{l}
{\kern0.2pt\scriptscriptstyle \infty}\\[-8pt]
{\kern1pt\scriptstyle \vee}\\[-5pt]
H^{#1}
\end{array}}

\def\Hv#1{%
\begin{array}[b]{l}
{\kern1pt\scriptstyle \vee}\\[-5pt]
H^{#1}
\end{array}}

\DeclareMathOperator{\cart}{cart}
\DeclareMathOperator{\Ffaisc}{\mathbf{Faisc}}


\def\Ev#1{%
\begin{array}[b]{l}
{\kern1pt\scriptstyle \vee}\\[-5pt]
E'_{#1}
\end{array}}

\DeclareMathOperator{\Cont}{Cont}
\DeclareMathOperator{\Ccont}{\mathscr{C}\!\mathit{ont}}
\DeclareMathOperator{\Site}{Site}
\newcommand\MM{\mathbf{M}}
\def\limtop#1{{\displaystyle{\varprojlim_{#1}\!\mbox{\rm top }}}}
\def\Limtop{\mathop{\varprojlim\!\mbox{\rm top}}}
\DeclareMathOperator{\spec}{spec}
\def\ss{\scriptstyle{\mathscr{S}}}
\DeclareMathOperator{\Gal}{Gal}
\def\smallm{\scriptstyle{\mathscr{M}}}
\DeclareMathOperator{\End}{End}
\DeclareMathOperator{\ass}{ass}
\DeclareMathOperator{\zar}{zar}
\DeclareMathOperator{\Ogg}{Ogg}
\DeclareMathOperator{\res}{res}
\DeclareMathOperator{\tr}{tr}
\DeclareMathOperator{\Supp}{Supp}
\DeclareMathOperator{\codim}{codim}
\DeclareMathOperator{\rk}{rk}
\DeclareMathOperator{\cur}{cur}
\DeclareMathOperator{\cdqc}{cdqc}
\newcommand\DD{\mathbf{D}}
\DeclareMathOperator{\cl}{cl}
\def\JJ{\mathbf{J}}
\def\FF{\mathbf{F}}
\def\BB{\mathbf{B}}
\DeclareMathOperator{\Alget}{Alget}
\DeclareMathOperator{\Ass}{Ass}
\DeclareMathOperator{\supp}{supp}
\DeclareMathOperator{\rad}{rad}
\DeclareMathOperator{\prof}{prof}
\DeclareMathOperator{\cint}{int}
\DeclareMathOperator{\sic}{sic}
\DeclareMathOperator{\Tor}{Tor}
\DeclareMathOperator{\tors}{tors}
\DeclareMathOperator{\Sym}{Sym}
\DeclareMathOperator{\parf}{parf}
\DeclareMathOperator{\tor}{tor}
\DeclareMathOperator{\defn}{defn}
\DeclareMathOperator{\Div}{Div}
\DeclareMathOperator{\Pol}{Pol}
\DeclareMathOperator{\Fex}{Fex}
\DeclareMathOperator{\Lie}{Lie}
\DeclareMathOperator{\tot}{tot}
\DeclareMathOperator{\Aut}{Aut}
\DeclareMathOperator{\HOM}{HOM}
\DeclareMathOperator{\pch}{pch}
\DeclareMathOperator{\Ch}{Ch}
\DeclareMathOperator{\EXT}{EXT}
\DeclareMathOperator{\PIC}{PIC}
\DeclareMathOperator{\Rf}{Rf}
\DeclareMathOperator{\inj}{inj}
\newcommand\VV{\mathbf{V}}

\def\Xv{%
\begin{array}[b]{l}
{\kern1pt\scriptstyle \vee}\\[-5pt]
X
\end{array}}

%%%%%%%
\newcommand\TT{\mathbf{T}}
\newcommand\XX{\mathbf{X}}
\newcommand\EE{\mathbf{E}}
%%%%%%

\def\qedsymbol{\textup{C.Q.F.D.}}
\let\CQFD\qed
\let\cqfd\qed

%%%%%\bigdot
\font\bigsymb=cmsy10 at 4pt
\def\bigdot{{\kern1.2pt\raise 1.5pt\hbox{\bigsymb\char15}}}
%%%%%

\def\Lright{\mathop{\vcenter{\hbox{\hbox{$\longrightarrow$}\lower8pt \hbox{\kern-16pt\hbox{$\longrightarrow$}}}}}}

\def\Aa{{\mathcal A}}
\def\Bb{{\mathcal B}}
\def\Cc{{\mathcal C}}
\def\Dd{{\mathcal D}}
\def\Ee{{\mathcal E}}
\def\Ff{{\mathcal F}}
\def\Gg{{\mathcal G}}
\def\Hh{{\mathcal H}}
\def\Ii{{\mathcal I}}
\def\Jj{{\mathcal J}}
\def\Kk{{\mathcal K}}
\def\Ll{{\mathcal L}}
\def\Mm{{\mathcal M}}
\def\Nn{{\mathcal N}}
\def\Oo{{\mathcal O}}
\def\Pp{{\mathcal P}}
\def\Qq{{\mathcal Q}}
\def\Rr{{\mathcal R}}
\def\Ss{{\mathcal S}}
\def\Tt{{\mathcal T}}
\def\Uu{{\mathcal U}}
\def\Vv{{\mathcal V}}
\def\Ww{{\mathcal W}}
\def\Xx{{\mathcal X}}
\def\Yy{{\mathcal Y}}
\def\Zz{{\mathcal Z}}

\newcommand{\CC}{\mathbf{C}}
\newcommand{\GG}{\mathbf{G}}
\newcommand{\KK}{\mathbf{K}}
\newcommand{\LL}{\mathbf{L}}
\newcommand{\NN}{\mathbf{N}}
\newcommand{\PP}{\mathbf{P}}
\newcommand\QQ{\mathbf{Q}}
\newcommand\RR{\mathbf{R}}
\newcommand\ZZ{\mathbf{Z}}

\def\to{\mathchoice{\longrightarrow}{\rightarrow}{\rightarrow}{\rightarrow}}
\def\from{\mathchoice{\longleftarrow}{\leftarrow}{\leftarrow}{\leftarrow}}
\def\mto{\mathchoice{\longmapsto}{\mapsto}{\mapsto}{\mapsto}}
\def\mfrom{\mathrel{\from\joinrel\mapstochar\,}}
\def\hto{\mathrel{\lhook\joinrel\to}}
\def\hfrom{\mathrel{\from\joinrel\rhook}}

\def\implies{\mathchoice{\Longrightarrow}{\Rightarrow}{\Rightarrow}{\Rightarrow}}
\def\implied{\mathchoice{\Longleftarrow}{\Leftarrow}{\Leftarrow}{\Leftarrow}}
\def\iff{\mathchoice{\Longleftrightarrow}{\Leftrightarrow}{\Leftrightarrow}{\Leftrightarrow}}

\def\To#1{\mathchoice{\xrightarrow{\labelstyle~~#1~}}{\stackrel{#1}{\longrightarrow}}{}{}}
\def\Mto#1{\mathrel{\mapstochar\To{#1}}}
\def\From#1{\mathchoice{\xleftarrow{\labelstyle~~#1~}}{\stackrel{#1}{\longleftarrow}}{}{}}

\newcommand{\isomto}{\overset{\sim}{\longrightarrow}}
\newcommand{\isomfrom}{\overset{\sim}{\longleftarrow}}

\newcommand{\IFF}{\unskip~$\iff$~\ignorespaces}
\newcommand{\THEN}{\unskip~$\implies$~\ignorespaces}

\let\leq\leqslant
\let\geq\geqslant
\let\le\leqslant
\let\ge\geqslant

\def\cf{{cf.}\kern.3em}
\def\Cf{{Cf.}\kern.3em}
\def\ie{{i.e.}\kern.3em}
\def\iev{{i.e.},\xspace}
%\def\loccito{loc.\kern2pt cit.\xspace}
%\def\loccitn{\textit{loc.\kern2pt cit.}\xspace}
%\def\loccit{\sisi{\loccito}{\loccitn}}
\def\EGA{\textup{EGA}\kern.2em}
\def\SGA{\textup{SGA}\kern.3em}

\let\oldS\S
\def\S{\oldS\kern .15em}
\def\ptbl{.\kern .2em}
\def\resp{resp.\kern.3em}
\def\Exp{Exp.\kern.3em}

\let\bar\overline
\let\hat\widehat
\let\tilde\widetilde
\let\dpl\displaystyle
\let\ldots\dots
\let\emptyset\varnothing
\let\skpt\skippointrait
\let\epsilon\varepsilon

\def\sfrac#1#2{{#1}/{#2}}

\DeclareMathOperator{\cd}{cd}
%\DeclareMathOperator{\Hhom}{\mathcal{H}\!\mathit{om}}
\DeclareMathOperator{\Hhom}{\mathbf{Hom}}
%\DeclareMathOperator{\Eext}{\mathcal{E}\!\mathit{xt}}
\DeclareMathOperator{\Eext}{\mathbf{Ext}}
\DeclareMathOperator{\Spec}{Spec}
\DeclareMathOperator{\Tr}{Tr}

\let\emptyset=\varnothing

%%%%%%%%%%%%%%%
%%% expose XVII
%%%%%%%%%%%%%%%

\makeatletter
\let\old@optexpose\optexpose
\newcommand{\empty@optexpose}[1]{}
\newcommand{\paren@optexpose}[1]{\old@optexpose{#1}(}
\newcommand{\refnoexpose}[1]{\let\optexpose\empty@optexpose\ref{#1}\let\optexpose\old@optexpose}%\refnoexpose{X.cor5.2} engendre "5.2" au lieu de "X 5.2"
\newcommand{\neweqref}[1]{\let\optexpose\paren@optexpose\textup{\ref{#1})}\let\optexpose\old@optexpose}%\neweqref{XII.eq4.1} engendre "XII (4.1)" au lieu de "(XII 4.1)"
\makeatother
\newcommand{\refvb}[2][VI${}_{\text{B}}$ ]{\sisi{#1\refnoexpose{#2}}{\ref{#2}}}

\newcommand{\leftexp}[2]{{\vphantom{#2}}^{#1}{#2}}
\newcommand{\AAA}{\mathbf{A}}
\newcommand{\Sf}{\mathfrak{S}}
\DeclareMathOperator*{\quotprojlim}{\text{``$\varprojlim$''}}
\DeclareMathOperator*{\quotinjlim}{\text{``$\varinjlim$''}}
\DeclareMathOperator{\Fl}{F\sisi{\ell}{l}}
\DeclareMathOperator{\Pl}{P\sisi{\ell}{l}}
%\DeclareMathOperator{\cl}{c\sisi{\ell}{l}}
\DeclareMathOperator{\tdf}{tdf}
\DeclareMathOperator{\rg}{rg}
\newcommand{\iInd}{\sisi{I}{i}nd}
\newcommand{\iind}{ind}
\newcommand{\ppro}{pro}
\newcommand{\tTor}{\sisi{T}{t}or}
\newcommand{\ttor}{tor}
\newcommand{\chap}{Chap.}
\newcommand{\Ffex}{\Ff\mathit{ex}}
\newcommand{\sheafk}{\underline{k}}
\newcommand{\sheafn}{\underline{n}}
\newcommand{\sheafGamma}{\underline{\Gamma}}

% Repeated references to an already typeset footnote must not use
% \footnotemark without a matching \footnotetext: Hyperref would then create
% a live annotation to a destination that cannot exist.  This helper keeps the
% current printed footnote mark and links it to an explicit stable note target.
\makeatletter
\newcommand{\sgarepeatedfootnotemark}[1]{%
  \begingroup
  \edef\@thefnmark{\thefootnote}%
  \hyperlink{#1}{\@makefnmark}%
  \endgroup}
\makeatother

\sisi{
\DeclareMathOperator{\simarrow}{\rule{.15mm}{1.8mm}\kern -3.4pt\rightsquigarrow}
\newcommand{\sheafell}{\underline{\ell}}
\newcommand{\sheafZZ}{\underline{\ZZ}}
}{
\newcommand{\nderef}[1]{la N.D.E.~(\ref{#1}) page~\pageref{#1}}
}

%%%%%%%%%%%%%%%%%
%%%%% expose XVII
%%%%%%%%%%%%%%%%%

% END INCLUDED SGA4_English_current_buildable_source_20260906/source/fichiers-configuration-latex/macros.tex


\externaldocument{../01/01_English.tex}
\externaldocument{../02/02_English.tex}
\externaldocument{../03/03_English.tex}
%\externaldocument{../04/04_English.tex}

\setboolean{original}{false}

\begin{document}

\begin{center}
IV. Topoi\\
A. Grothendieck and J.-L. Verdier
\end{center}

\selectlanguage{english}version: \showgitstatus
\tableofcontents

\else
\chapter{Topoi}
%\addtocontents{toc}{A. Grothendieck and J.-L. Verdier}
\begin{center}
A. Grothendieck and J.-L. Verdier
\end{center}
\fi
\phantomsection\MakeLinkTarget*{sga4.exp.IV}\label{sga4.exp.IV}

% BEGIN INCLUDED SGA4_English_current_buildable_source_20260906/source/04/drafts/SGA4_IV_LINES_31_187_English.texfrag
\setcounter{pageoriginal}{298}
\setcounter{section}{-1}
\section{Introduction}\etiquette[0]{IV.sec0}\pageoriginale

\subsection{}\etiquette[0.1]{IV.subsec0.1}
We saw in \hyperref[sga4.exp.II]{II} various exactness properties of categories of the
form $\widetilde{𝒞}=$ the category of sheaves of sets on $𝒞$, where
$𝒞$ is a small site. These properties can be expressed by saying
that in many respects such categories (which we shall call
\emph{topoi}) inherit the familiar properties of the category
$(\Ens)$ of (small) sets. On the other hand, experience has shown
that various situations in Mathematics should be regarded
\emph{above all as a technical means of constructing the
  corresponding categories of sheaves} (of sets), \ie \emph{the
  corresponding ``topoi.''} It appears that all the genuinely
important notions attached to a site (for example its cohomological
invariants, studied in \hyperref[sga4.exp.V]{V}; various other ``topological''
invariants, such as its homotopy invariants recently studied by M.
\textsc{Artin} and B. \textsc{Mazur} [1]; and the notions studied in
J. \textsc{Giraud}'s book on noncommutative cohomology) are in fact
expressed directly in terms of the associated topos. From this point
of view, two sites should be regarded as essentially equivalent when
their associated topoi are equivalent categories. The datum of a
site (at least in the case, especially important in practice, where
its topology is coarser than its canonical topology) should be
regarded as amounting to the datum of a topos $ℰ$ (namely the
associated topos formed by the sheaves of sets on the site) together
with a generating family of objects of $ℰ$ (\cf\pageoriginale II
\hyperref[II.subsec4.9]{4.9}, and \hyperref[IV.cor1.2.1]{1.2.1} below). This viewpoint is
analogous to associating a group with a system of generators and a
system of relations among those generators, and attaching interest
rather to the structure of that group than to the system of
generators and relations used to generate it (which are regarded as
auxiliary data of the situation). Moreover, the ``comparison lemma''
III \hyperref[III.theo4.1]{4.1} provides numerous examples of pairs of sites $𝒞$, $𝒞'$ that
are not isomorphic, and are not even equivalent as categories, but
give rise to equivalent topoi; thus $𝒞$ and $𝒞'$ should be regarded
as essentially equivalent.

\subsection{}\etiquette[0.2]{IV.subsec0.2}
In the present exposé we give a characterization of topoi by simple
exactness properties (due to J. \textsc{Giraud}); we study the
natural notion of a morphism of topoi, inspired by the notion of a
continuous map from one topological space to another; and we develop
in the setting of topoi certain constructions familiar from ordinary
sheaf theory (sheaves $\Hhom$, tensor-product sheaves, supports).
Finally, following M. \textsc{Artin}, we show how a topos can be
reconstructed from an ``open'' part, its complementary ``closed''
part, and a certain left exact functor connecting them, which may in
fact be chosen almost arbitrarily.

\subsection{}\etiquette[0.3]{IV.subsec0.3}
This gives a procedure for gluing topoi which, when applied to topoi
arising from ordinary topological spaces, will in general produce a
topos no longer of the same type. This is a first indication of the
remarkable stability of the notion of a topos under various natural
constructions, a stability lacking for the notion of a topological
space (from which the notion of a\pageoriginale topos is inspired).
As a second remarkable example, let us also mention the relative
classifying topos of a group in a topos (\cf [1] or
\hyperref[IV.exer5.9]{5.9} below), inspired by the classical notion
of the classifying space of a topological group, and the notion of a
``moduli topos'' associated with various ``moduli problems'' in
Algebraic Geometry or Analytic Geometry [10] [13].

Other topoi, such as the \emph{étale topos} of a scheme
(systematically studied in the present Seminar beginning with \Exp
VII) or the \emph{crystalline topos} of a relative scheme [6], arise
naturally when one wishes to develop usable cohomology theories for
abstract algebraic varieties (and, more generally, schemes),
replacing the classical Betti cohomology of algebraic varieties over
the field of complex numbers.

\subsection{}\etiquette[0.4]{IV.subsec0.4}
Thus one may say that the notion of a topos, a natural outgrowth of
the \emph{sheaf-theoretic viewpoint} in Topology, in turn constitutes
a substantial enlargement of the notion of a topological
space\footnote[1]{\Cf [9], or \hyperref[IV.subsec4.1]{4.1} and
\hyperref[IV.subsec4.2]{4.2} below, for the precise relations between
  the notion of a topos and that of a topological space.}, embracing
a large number of situations which formerly were not regarded as
belonging to topological intuition. The characteristic feature of
such situations is the availability of a notion of ``localization,''
a notion formalized precisely by that of a site and, ultimately, by
that of a topos (through the topos associated with the site). As the
term ``topos'' itself is intended precisely to suggest, it seems
reasonable and legitimate to the authors of the present Seminar to
regard the subject of Topology as the study of \emph{topoi} (and not
merely of topological spaces).

\subsection{}\etiquette[0.5]{IV.subsec0.5}
It seemed useful to include in this general exposé on topoi a fairly
large number of examples, many of which are only remotely\pageoriginale related to
the original aim of this seminar (namely, the study of étale
cohomology). The reader in a hurry, interested exclusively in étale
cohomology, may of course omit these examples without inconvenience,
as well as most of the present exposé, consulting it only when the
need arises.

\section{Definition and characterization of topoi}\etiquette[1]{IV.sec1}

\begin{enonce*}{Definition 1.1}\etiquette[1.1]{IV.def1.1}
A category $E$ is called a $𝒰$-topos, or simply a topos when no
confusion is likely, if there exists a site $C\in 𝒰$ such that $E$ is
equivalent to the category $C^{\widetilde{\phantom{,}}}$ of
$𝒰$-sheaves of sets on $C$.
\end{enonce*}

\setcounter{subsection}{1}
\subsubsection{}\etiquette[1.1.1]{IV.subsubsec1.1.1}
Let $E$ be a $𝒰$-topos. We always regard $E$ as endowed with its
canonical topology (\hyperrefext[II][II.def2.5]{2.5}), which therefore
makes it a \emph{site}, and even, by \hyperref[IV.subsubsec1.1.2]{1.1.2}
d) below, a $𝒰$-site (\hyperrefext[II][II.def3.0.2]{3.0.2}). Unless
expressly stated otherwise, we shall consider no topology on $E$
other than the one just specified.

\subsubsection{}\etiquette[1.1.2]{IV.subsubsec1.1.2}
We saw in \hyperrefext[II][II.prop4.8]{4.8},
\hyperrefext[II][II.cor4.11]{4.11} that a $𝒰$-topos $E$ is a
$𝒰$-category (\hyperrefext[I][I.def1.1]{1.1}) satisfying the following
conditions:
\begin{enumerate}
\renewcommand{\theenumi}{\alph{enumi}}
\renewcommand{\labelenumi}{\theenumi)}
\item Finite projective limits are representable in $E$.

\item Direct sums indexed by an element of $𝒰$ are representable in
  $E$. They are disjoint and universal
  (\hyperrefext[II][II.def4.5]{4.5}).

\item Equivalence relations in $E$ are universally effective
  \sisi{(\hyperrefext[I][I.sec10.9]{10.9}).}{(\hyperrefext[I][I.sec10.10]{10.10}).}

\item $E$ admits a generating family
  (\hyperrefext[II][II.subsec4.9]{4.9}) indexed by an element of $𝒰$.
\end{enumerate}

In\pageoriginale fact, we shall see that these intrinsic properties
\emph{characterize} $𝒰$-topoi:

% END INCLUDED SGA4_English_current_buildable_source_20260906/source/04/drafts/SGA4_IV_LINES_31_187_English.texfrag


% BEGIN INCLUDED SGA4_English_current_buildable_source_20260906/source/04/drafts/SGA4_IV_LINES_188_382_English.texfrag
\begin{enonce*}{Theorem 1.2}\etiquette[1.2]{IV.theo1.2}
{\rm (J. Giraud).} Let $E$ be a $𝒰$-category. The following
properties are equivalent:
\begin{itemize}
\item[i)] $E$ is a $𝒰$-topos $(\hyperref[IV.def1.1]{1.1})$.

\item[ii)] $E$ satisfies conditions {\rm a), b), c),} and {\rm d)
  of \hyperref[IV.subsubsec1.1.2]{1.1.2}.}

\item[iii)] The $𝒰$-sheaves on $E$ for the canonical topology
  are representable, and $E$ has a small generating family
  (condition {\rm \hyperref[IV.subsubsec1.1.2]{1.1.2} d)}).

\item[iv)] There exist a category $C\in 𝒰$ and a fully faithful
  functor $i:E\rightarrow \widehat{C}$ (where $\widehat{C}$ denotes
  the category of $𝒰$-presheaves on $C$) having a left adjoint
  a which is left exact.

\item[${\rm i')}$] There exists a site $C\in 𝒰$ such that
  projective limits are representable in $C$ and the topology of
  $C$ is coarser than its canonical topology {\rm (II \hyperref[II.def2.5]{2.5}),} %% label!
  and such that $E$ is equivalent to the category
  $C^{\widetilde{\phantom{,}}}$ of $𝒰$-sheaves of sets on
  $C$.
\end{itemize}
\end{enonce*}

Moreover:

\begin{enonce*}{Corollary 1.2.1}\etiquette[1.2.1]{IV.cor1.2.1}
Let $E$ be a $𝒰$-topos, and let $C$ be a full subcategory of
$E$, endowed with the topology induced {\rm (\hyperrefext[III][III.subsec3.1]{3.1})}
by that of $E$ (\hyperref[IV.subsubsec1.1.1]{1.1.1}). Consider the
functor
$$
E\longrightarrow C^{\widetilde{\phantom{,}}}
$$
which associates to every $X\in \ob E$ the restriction to $C$ of
the functor represented by $X$. This functor is an equivalence of
categories if and only if $\ob C$ is a generating family of $E$.
\end{enonce*}

The equivalence\pageoriginale i) $\Leftrightarrow$ iv) follows
immediately from \hyperrefext[II][II.theo5.5]{5.5}, and we have already recalled
above that i) $\Rightarrow$ ii). Since ${\rm i'}$) $\Rightarrow$ i)
trivially, it remains to prove ii) $\Rightarrow$ iii) and iii)
$\Rightarrow$ ${\rm i'}$), which will be done in
\hyperref[IV.subsubsec1.2.4]{1.2.4} and \hyperref[IV.subsubsec1.2.3]{1.2.3} below. The
corollary is then obtained by observing that the functor in question
factors as $E\rightarrow E^{\widetilde{\phantom{,}}}\rightarrow
C^{\widetilde{\phantom{,}}}$; since $E\rightarrow
E^{\widetilde{\phantom{,}}}$ is an equivalence by virtue of
(iii), the question reduces to determining when
$E^{\widetilde{\phantom{,}}}\rightarrow C^{\widetilde{\phantom{,}}}$
is an equivalence. The conclusion then follows from the ``comparison
lemma'' \hyperrefext[III][III.theo4.1]{4.1}; \cf the proof in
\hyperref[IV.subsubsec1.2.3]{1.2.3} below.

\setcounter{subsection}{2} \setcounter{subsubsection}{2}
\subsubsection{}\etiquette[1.2.3]{IV.subsubsec1.2.3}
\emph{Proof of} iii) $\Rightarrow$ ${\rm i'}$). Let
$\mathfrak{X}=(X_i)_{i\in I}$, $I\in 𝒰$, be a small generating
family of $E$. Since $E$ is a $𝒰$-category, the set of
isomorphism classes of finite diagrams in $E$ whose objects are
members $X_i$ is $𝒰$-small. Consequently, the smallest set
$\mathfrak{X}'$ of objects of $E$ containing the finite projective
limits of objects of $\mathfrak{X}$ is a countable union of small
sets and is therefore small~\footnote[1]{Here we reason with classes
of objects of $E$ up to isomorphism.}. Setting
$\mathfrak{X}^{(n+1)}=\sisi{\mathfrak{X}^{(n)(1)}}{(\mathfrak{X}^{(n)})'}$, and
$\mathfrak{X}=\bigcup_n \mathfrak{X}^{(n)}$, we see that, after
enlarging the generating family, we may suppose that
$\mathfrak{X}$ is stable under finite projective limits. Let
$𝒱$ be a universe containing $𝒰$ such that $E$ is
$𝒱$-small. For every object $H$ of $E$, denote by $I(H)$
the set $\amalg_{i\in I}\Hom(X_i, H)$. Let $H$ be an object of $E$
and let $H'\hookrightarrow H$ be the subsheaf of $H$, for the
canonical topology, which is the ``union'' of the images of the
morphisms $u:X_i\rightarrow H$, $(u,i)\in I(H)$
(\hyperrefext[II][II.theo4.1]{4.1}). Since $H'$ is a subsheaf of a
$𝒰$-sheaf, $H'$ is a $𝒰$-sheaf. It is therefore
representable. Since the family $\mathfrak{X}$ is generating and,
for every $i\in I$, the map $\Hom(X_i, H')\rightarrow \Hom(X_i,H)$
is bijective, the morphism $H'\hookrightarrow H$ is an isomorphism
(\hyperrefext[II][II.subsec4.9]{4.9}). Consequently
(\sisi{\hyperrefext[II][II.def6.1]{6.1}}{\hyperrefext[II][II.def5.2]{5.2}}), the family
$(u:X_i\rightarrow H)$, $(u,i)\in I(H)$, is covering for the
canonical topology of $E$. Thus\pageoriginale every object of $E$
can be covered, for the canonical topology of $E$, by objects of
$\mathfrak{X}$. Let $C$ be the full subcategory of $E$ defined by
the objects of $\mathfrak{X}$, and let $u:C\hookrightarrow E$ be the
inclusion functor. The topology $\Cc$ induced on $C$ by the canonical
topology of $E$ is coarser than the canonical topology of $C$, and
when $𝒰$ contains an element of infinite cardinality, finite
projective limits are representable in $C$. It follows from
\hyperrefext[III][III.theo4.1]{4.1} that the functor $F\mapsto F\circ
u$ is an equivalence from $E$ to the category of $𝒰$-sheaves
on $C$ for the topology $\Cc$, \cqfd

\subsubsection{}\etiquette[1.2.4]{IV.subsubsec1.2.4}
\emph{Proof of} ii) $\Rightarrow$ iii). The proof consists of four
steps.

Let $𝒱$ be a universe containing $𝒰$, such that $E$ is
an element of $𝒱$. Let $\widetilde{E}$ be the category of
$𝒱$-sheaves on $E$ for the canonical topology, and let
$J_E:E\rightarrow \widetilde{E}$ be the canonical functor.

\paragraph{}\etiquette[1.2.4.1]{IV.par1.2.4.1}
Let $(g_i:G_i\rightarrow H)$, $i\in I\in 𝒰$, be an
epimorphic family in $\widetilde{E}$. If the $G_i$ and the fiber
products $\fprod{G_i}{G_j}{H}$ are representable, then the sheaf
$H$ is representable.

Indeed, the direct sum $\coprod_{i\in I}G_i$ is represented in
$\widetilde{E}$ (\hyperrefext[II][II.theo4.1]{4.1}) by a representable sheaf
$G$ (property b) and \hyperrefext[II][II.prop4.6]{4.6}). The same holds for the
direct sum $K=\coprod_{(i,j)\in I\times
I}\fprod{G_i}{G_j}{H}$. Moreover, the diagram $K\rightrightarrows
G\rightarrow H$ is exact and $K$ is the fiber square of $G$ over
$H$. (This latter property holds in the category of presheaves, and
hence in the category of sheaves (\hyperrefext[II][II.theo4.1]{4.1}).) It follows
from c) and \hyperrefext[II][II.prop4.7]{4.7} that the sheaf $H$ is
representable.

\paragraph{}\etiquette[1.2.4.2]{IV.par1.2.4.2}
Let $X_{\alpha}$, $\alpha\in A\in 𝒰$, be a generating
family of $E$. For every sheaf $H$, denote by $I(H)$ the set
$\coprod_{\alpha \in A}\Hom_E (X_{\alpha}, H)$. The family
$(u:X_{\alpha}\rightarrow H, (u,\alpha)\in
I(H))$\pageoriginale is epimorphic in $\widetilde{E}$. When $H$ is
a $𝒰$-sheaf, $I(H)$ is an element of $𝒰$.

Indeed, since every sheaf $H$ is the target of an epimorphic family
of morphisms whose sources are representable sheaves (
\hyperrefext[I][I.prop3.4]{3.4} and \hyperrefext[II][II.theo4.1]{4.1}), it is enough to prove the
first assertion when the sheaf $H$ is representable. Let $G$ then
be the image of the family $(u:X_{\alpha}\rightarrow H,
(u,\alpha)\in I(H))$. The morphism $G\rightarrow H$ is a
monomorphism. We are therefore in the situation of the first step,
since $\fprod{X_{\alpha}}{X_{\beta}}{G}$ is isomorphic to
$\fprod{X_{\alpha}}{X_{\beta}}{H}$, which is representable, and
moreover $I(H)$ is an element of $𝒰$. It follows that $G$ is
representable. But since $X_{\alpha}$ is a generating family,
$G\rightarrow H$ is an isomorphism. The last assertion is immediate
from the definition of $𝒰$-sheaves.

\paragraph{}\etiquette[1.2.4.3]{IV.par1.2.4.3}
Every subsheaf of a representable sheaf is representable.

Indeed, let $G\rightarrow H$ be a subsheaf of a representable
sheaf. Then $G$ is a $𝒰$-sheaf. The family
$(u:X_{\alpha}\rightarrow G, (u,\alpha)\in I(G))$ is therefore
epimorphic in $\widetilde{E}$ and indexed by an element of
$𝒰$. Moreover, the fiber products
$\fprod{X_{\alpha}}{X_{\beta}}{G}$ are isomorphic to the fiber
products $\fprod{X_{\alpha}}{X_{\beta}}{H}$, which are
representable. Thus we are in the situation of the first step and
$G$ is representable.

\paragraph{}\etiquette[1.2.4.4]{IV.par1.2.4.4}
Every $𝒰$-sheaf is representable.

Indeed, by the first and second steps, it is enough to show that the
fiber products $\fprod{X_{\alpha}}{X_{\beta}}{H}$ are representable.
But these fiber products are subobjects of the products
$\fprod{X_{\alpha}}{X_{\beta}}{}$. The conclusion therefore follows
from the third step. This completes the proof of Theorem
\hyperref[IV.theo1.2]{1.2}.

\begin{enonce*}[remark]{Remark 1.3}\etiquette[1.3]{IV.rem1.3}
Of\pageoriginale course, for a given $𝒰$-topos $E$, there is in
general no preferred way of representing it, up to equivalence, in
the form $C^{\widetilde{\phantom{,}}}$, with $C$ a \emph{small}
site; or, what amounts essentially to the same thing by virtue of
\hyperref[IV.cor1.2.1]{1.2.1} when one restricts to $C$ whose topology is
coarser than the canonical topology, there is no preferred
\emph{small} generating family in $E$. When one ceases to impose on
$C$ the condition that $C$ be small, there is, on the other hand (by virtue of
\hyperref[IV.theo1.2]{1.2} iii)), an entirely canonical choice of a
$𝒰$-site $C$ such that $E$ is equivalent to
$C^{\widetilde{\phantom{,}}}$, namely $E$ itself! This is one of the
technical reasons why in practice it is inconvenient to work only
with small sites: in fact, the most important sites of all, namely
$𝒰$-topoi, are not small! Moreover, the sites generating topoi
that arise in many questions in algebraic geometry (and even in
topology, \cf \hyperref[IV.subsec2.5]{2.5}) are not small either;
examples include the étale site of a scheme \sisi{(VIII.1).}{(VII.1).}
\end{enonce*}

% END INCLUDED SGA4_English_current_buildable_source_20260906/source/04/drafts/SGA4_IV_LINES_188_382_English.texfrag


% BEGIN INCLUDED SGA4_English_current_buildable_source_20260906/source/04/drafts/SGA4_IV_LINES_383_536_English.texfrag
\begin{enonce*}{Proposition 1.4}\etiquette[1.4]{IV.prop1.4}
Let $E$ be a $𝒰$-topos, let $E'$ be a category (not
necessarily a $𝒰$-category), and let $F$ be a presheaf
on $E$ with values in $E'$. In order that $F$ be a sheaf with
values in $E'$ {\rm(\hyperrefext[II][II.def6.1]{6.1}),} it is necessary and sufficient that
$F$ take $𝒰$-inductive limits in $E$ to projective
limits in $E'$.
\end{enonce*}

Let $𝒱$ be a universe containing $𝒰$ and such that $E'$ is
a $𝒱$-category. By composing $F$ with the functors
$\Hom(X',-):E'\rightarrow 𝒱\text{-}\Ens$ defined by the
objects $X'$ of $E'$, we are reduced to the case where
$E'=𝒱\text{-}\Ens$, with $𝒱$ a universe.
Suppose that $F$ takes $𝒰$-inductive limits to projective
limits. It then follows at once from the definitions that $F$ is
a sheaf, since a covering family\pageoriginale
$X_i\rightarrow X$ in $ℰ$, by the definition of the canonical
topology of $ℰ$, allows one to regard $X$ as an inductive limit
of the diagram $ \left(
\fprod{X_i}{X_j}{X}\raise6pt\hbox{$\mathop{\nearrow}^{\mbox{$X_i$}}_{\displaystyle{\kern-10pt\raise-8pt\hbox{$\searrow_{\mbox{$X_j$}}$}}}$}
\right) $. Suppose that $F$ is a sheaf, and let us prove that it
takes $𝒰$-inductive limits to projective limits. If
$𝒱\subset 𝒰$, we may suppose that $𝒱=𝒰$,
and it suffices to apply criterion \hyperref[IV.theo1.2]{1.2} iii). To
treat the general case, a little more work is needed. Let
$C$ be a full subcategory of $E$ generated by a generating family
indexed by an element of $𝒰$, and let
$u:C\rightarrow E$ be the inclusion functor. Endow $C$ with the
topology induced by the canonical topology of $E$ \sisi{(III 4.5).}{(III \hyperref[III.prop3.5]{3.5}).}
Denote by $C^{\widetilde{\phantom{,}}}_{𝒰}$ and
$C^{\widetilde{\phantom{,}}}_{𝒱}$ the categories of
sheaves on $C$, by $E^{\widetilde{\phantom{,}}}_{𝒱}$ the
category of $V$-sheaves on $E$ for the canonical topology,
by $J_E:E\rightarrow E_{𝒱}$ the canonical functor, and
by $i_{𝒰,𝒱}:C^{\widetilde{\phantom{,}}}_{𝒰}\rightarrow
C^{\widetilde{\phantom{,}}}_{𝒱}$ the inclusion functor. The
diagram
\[\vcenter{
\xymatrix@C=3.25cm@R=1.25cm{
C^{\widetilde{\phantom{a}}}_{\underline{U}} \ar[r]^{i_{\underline{U},\underline{V}}} & C^{\widetilde{\phantom{a}}}_{\underline{V}}\\
E  \ar[r]^{J_E} \ar[u]&
E^{\widetilde{\phantom{a}}}_{\underline{V}}\ar[u]\quad,
}}\leqno{(*)}
\]
where the vertical arrows are induced by the functor
$F\mapsto F\circ u$, is commutative. It follows from the explicit
construction of the associated sheaf functor (\hyperrefext[II][II.sec3]{3}) and from
\hyperrefext[II][II.theo4.1]{4.1} that the functor $i_{𝒰,𝒱}$ commutes
with $𝒰$-inductive limits. Moreover, it follows from
\hyperrefext[III][III.subsec5.1]{5.1} and \hyperref[IV.cor1.5]{1.5} that the vertical arrows
in $(\ast)$ are equivalences of categories. Consequently,
$J_E:E\rightarrow E^{\widetilde{\phantom{,}}}_{𝒱}$
commutes with $𝒰$-inductive limits. For every object $X$ of
$E$, we have
$$
F(X)\simeq
\Hom_{E^{\widetilde{\phantom{,}}}_{𝒱}}(J_E(X),F){\quoi}.
$$
Thus\pageoriginale $F$ takes $𝒰$-inductive limits
in $E$ to projective limits.

\begin{enonce*}{Corollary 1.5}\etiquette[1.5]{IV.cor1.5}
Let $E$ be a $𝒰$-topos, let $E'$ be a $𝒰$-category, and let
$f:E\rightarrow E'$ be a functor. In order that $f$ admit a right
adjoint, it is necessary and sufficient that $f$ commute with
$𝒰$-inductive limits.
\end{enonce*}

This is an immediate consequence of \hyperref[IV.prop1.4]{1.4} in the case
$E'=𝒰\text{-}\Ens$.

\begin{enonce*}{Corollary 1.6}\etiquette[1.6]{IV.cor1.6}
Let $E$ and $E'$ be two $𝒰$-topoi, and let $f:E\rightarrow E'$ be a
functor. The following conditions are equivalent:
\begin{enumerate}
\renewcommand{\theenumi}{\roman{enumi}}
\renewcommand{\labelenumi}{\theenumi)}
\item $f$ commutes with $𝒰$-inductive limits.

\item $f$ admits a right adjoint.

\item $f$ is continuous {\rm(\hyperrefext[III][III.def1.1]{1.1}).}
\end{enumerate}
\end{enonce*}

The equivalence i) $\Leftrightarrow$ ii) was proved in
\hyperref[IV.cor1.5]{1.5}. To prove i) $\Leftrightarrow$ iii), apply the
definition of continuous functors, choosing a universe
$𝒱$ such that $𝒰\in 𝒱$ (so $E$ and $E'$ are
$𝒱$-small sites). We must express that for every
$𝒱$-sheaf $F$ on $E'$, the composite $F\circ f$ is a
sheaf on $E$, that is (\hyperref[IV.prop1.4]{1.4}), that it takes
$𝒰$-inductive limits to projective limits. For this it is
sufficient that i) hold, since by \hyperref[IV.prop1.4]{1.4} $F$
itself takes $𝒰$-inductive limits to projective
limits; it is also necessary, as one sees by taking $F$ to be
a representable functor.

\begin{enonce*}{Corollary 1.7}\etiquette[1.7]{IV.cor1.7}
With the notation of \hyperref[IV.cor1.6]{1.6}, in order that $f$ be the
inverse-image functor $u^{\ast}$ for a morphism of topoi
$(\hyperref[IV.def3.1]{3.1})$ $u:E'\rightarrow E$, it is necessary and sufficient that
$f$ be left exact and take epimorphic families to epimorphic
families.
\end{enonce*}

The\pageoriginale necessity is trivial by definition (NB every
right exact functor takes epimorphisms to epimorphisms). The
sufficiency follows from \sisi{\hyperrefext[III][III.prop2.6]{2.6}}{\hyperrefext[III][III.prop1.6]{1.6}}, which implies
that $f$ is continuous under the stated conditions, and hence
commutes with $𝒰$-inductive limits by
\hyperref[IV.cor1.6]{1.6}.

\begin{enonce*}[remark]{Remark 1.8}\etiquette[1.8]{IV.rem1.8}
With the notation of \hyperref[IV.cor1.5]{1.5}, one can show that $f$ admits
a left adjoint if and only if it commutes with
$𝒰$-projective limits. In other words, a
\emph{covariant} functor $E\rightarrow (𝒰\text{-}\Ens)$ is
representable if and only if it commutes with $𝒰$-projective
limits \footnote[1]{A more general statement is found in
\hyperrefext[I][I.cor8.12.8]{8.12.8}, \hyperrefext[I][I.exer8.12.9]{8.12.9}; the outline of the proof
following a) through c) corresponds to the proof given in
\loccit}. We indicate only the principle of the proof, which
proceeds in three steps:
\begin{enumerate}
\renewcommand{\theenumi}{\alph{enumi}}
\renewcommand{\labelenumi}{\theenumi)}
\item Standard arguments [5, ${\rm n^{\circ}}$ 195, \S 3]
show that if $F$ commutes with $\varprojlim$, it is
prorepresentable by a strict projective system $(T_i)_{i\in
I}$, where $I$ is a filtered ordered set, \emph{not
necessarily small.} We may suppose that if $i>j$, then
$T_i\rightarrow T_j$ is not an isomorphism, and under this
hypothesis $F$ is representable if and only if $I$ is
small (which in fact implies that the projective system is
essentially constant).

\item To prove that $I$ is small, knowing that for every object $X$
of $E$ the set $F(X)=\varinjlim_i \Hom(T_i,X)$ is small, it suffices
to have a \emph{small cogenerating family}
$(X_j)_{j\in J}$ (\ie one that is generating for the opposite
category $E^{\circ}$). Now one shows that every $𝒰$-topos
$E$ has a small cogenerating family.

\item To prove this last point, one observes by standard
arguments [Toh] that every object $X$ of $E$ admits a monomorphism
into an ``injective'' object; and then that for every generating
family $(L_{\alpha})$ of $E$, if\pageoriginale each $L_{\alpha}$ is
embedded in this way into an injective object $I_{\alpha}$, the
family $(I_{\alpha})$ is cogenerating.
\end{enumerate}
\end{enonce*}

% END INCLUDED SGA4_English_current_buildable_source_20260906/source/04/drafts/SGA4_IV_LINES_383_536_English.texfrag


% BEGIN INCLUDED SGA4_English_current_buildable_source_20260906/source/04/drafts/SGA4_IV_LINES_537_947_English.texfrag
\section{Examples of topoi}\etiquette[2]{IV.sec2}

\setcounter{subsection}{-1}
\subsection{}\etiquette[2.0]{IV.subsec2.0}
We have brought together in this section a fairly large number of
typical examples of topoi, which the reader will already have had
occasion to encounter elsewhere, and which are intended to facilitate
his access to the ``yoga'' of topoi. For other examples (drawn from
algebraic geometry) of topologies on sites, giving rise to as many
topoi, he may consult \SGA 3 IV 6, and (for the étale topos) Exposé VII
of the present seminar. Since later on we shall hardly refer to the
present section except for questions of notation or terminology, we
leave it to the reader, as an exercise, to verify the statements with
which we have supplemented these examples for his general instruction.
All the examples of the present section will be made more precise in
4, where their functorial dependence on the data will be examined, and
in the following sections by way of illustration of the general notions
concerning topoi.

\subsection{Topos associated with a topological space}\etiquette[2.1]{IV.subsec2.1}

Let $X$ be a small topological space, and let $\Ouv(X)$ be the category
of open subsets of $X$, endowed with the canonical topology
\sisi{(III 2.9.1).}{(III \hyperref[III.exem1.9.1]{1.9.1}).} We shall denote by $\Top(X)$ the
topos of $𝒰$-sheaves on $\Ouv(X)$. This topos is equivalent to the
category of étale topological spaces over $X$, by associating with every
such space $X'$ over $X$ the sheaf
$U\mapsto \Gamma(X'/U)=\Hom_X(U,X')$ on $\Ouv(X)$
[TF]\nde{Here [TF] denotes reference [8].} By an abuse\pageoriginale
of language, we shall sometimes be unable to refrain from denoting by
the same letter $X$ the topos $\Top(X)$ defined by the topological space
$X$.

It is of course the preceding example that served chiefly as a guide
and as an intuitive basis for the development of the theory of topoi.
One should note, however, that the topoi obtained from topological
spaces are of a very special nature, owing in particular to the fact
that they have been described by sites $C=\Ouv(X)$ in which
\emph{all morphisms are monomorphisms} (and hence whose underlying
category reduces to a preordered set). It follows in particular that
the sheaves represented by the objects of $C$ are subsheaves of the
terminal sheaf, and consequently that \emph{the subsheaves of the
terminal sheaf form a generating family of the topos under
consideration}. This property is not shared by most topoi that arise
naturally in algebraic geometry or in algebra, \cf the examples below.
It is, up to minor qualifications, characteristic of topoi of the form
$\Top(X)$ (\cf \hyperref[IV.exer7.1.9]{7.1.9} below).

It is easily verified that the map $\Ouv(X)\rightarrow \Top(X)$,
which associates with every open subset of $X$ the sheaf it represents,
is a \emph{bijection} from $\Ouv(X)$ onto the set of subobjects of the
terminal object of $\Top(X)$; this bijection is even an isomorphism for
the natural order structures, \ie it induces an isomorphism of the
corresponding categories. This suggests that it should be possible to
reconstruct the topological space $X$, up to homeomorphism, from
$\Top(X)$ known up to equivalence. We shall see below
(\hyperref[IV.subsec4.2]{4.2}) that this is indeed the case, subject to
a slight restriction on $X$.

\subsection{Point topos or terminal topos, and empty topos or initial topos}\etiquette[2.2]{IV.subsec2.2}\pageoriginale

When $X$ is a topological space consisting of a single point, the
functor
$$
F\longrightarrow F(X):\Top(X)\longrightarrow 𝒰\mbox{-}\Ens
$$
is an equivalence of categories. This shows in particular that the
\emph{category} $𝒰$-$\Ens$ is a $𝒰$-\emph{topos.} We saw by
examples in \Exp II that this $𝒰$-topos is typical from the point
of view of exactness properties, in that the verification of many
properties (especially exactness properties) of general topoi reduces
to this particular topos. The interpretation given here in terms of
the one-point topological space justifies the abuse of language by
which a topos equivalent to the category $𝒰$-$\Ens$ is called a
\emph{point topos} (although, as a category, it is not at all
equivalent to the one-point category!). This is the terminology that
corresponds to the correct geometric intuition of the role played by
these topoi. By another abuse of language, a point topos is sometimes
called a \emph{terminal topos}, \cf \hyperref[IV.subsec4.3]{4.3}; we
shall say ``\emph{the} terminal topos'' for the topos
($𝒰$-$\Ens$).

When $X$ is the empty topological space, so that $\Ouv(X)$ is the
one-point category, a presheaf $F$ on $\Ouv(X)$ is a sheaf if and only
if its value at the unique object $\phi$ of $\Ouv(X)$ is a singleton
set. It follows that $\Top(∅)$ is isomorphic to the category of
singleton $𝒰$-sets, a category equivalent to the one-point category.
From this one concludes in particular that the one-point category (as
well as any $𝒰$-category equivalent to it) is a $𝒰$-topos. By an
abuse\pageoriginale of language, it is sometimes called the
\emph{empty topos} or the \emph{initial topos} (\cf
\hyperref[IV.subsec4.4]{4.4}); one should take care that it is not
equivalent to the empty category.

\subsection{Topos associated with a space with operators}\etiquette[2.3]{IV.subsec2.3}

Let $X$ be a topological space and $G$ a discrete group acting on $X$
by homeomorphisms. In \sisi{[Toh. \hyperref[IV.subsec5.1]{5.1}]}{[Toh]\nde{Here
[Toh] denotes reference [4].} \hyperref[IV.subsec5.1]{5.1}} the category
of $G$-\emph{sheaves on} $X$, or, as we shall also say, sheaves on
$(X,G)$, was defined; these are sheaves (of sets) on $X$
\emph{equipped} with actions of $G$ compatible with those of $G$ on
$X$. One sees at once, using Giraud's criterion
\hyperref[IV.theo1.2]{1.2} iii), that this %% \XXX ii) est probablement plus facile à vérifier
category is a topos (NB $𝒰$ is understood throughout), which we shall
simply denote by $\Top(X,G)$. When $G$ is the trivial group, we recover
the example of \hyperref[IV.subsec2.1]{2.1}; when $X$ is the one-point
space, we obtain the topos of sets on which $G$ acts on the left (or
$G$-\emph{sets}), also called the \emph{classifying topos of the
discrete group} $G$, and denoted by $B_G$. It is easily verified that
the only subobjects of the terminal object $e$ of the topos $B_G$ are
$e$ and the empty sheaf $\phi$; in particular, if $G$ is not the
trivial group, the subobjects of the terminal object of the classifying
topos $B_G$ do not form a generating family of $B_G$. Thus $B_G$ is
then not equivalent to a topos of the type $\Top(X)$ considered in
\hyperref[IV.subsec2.1]{2.1}.

The notion of a $G$-sheaf was introduced in \loccit in order to develop
the cohomological theory of abelian $G$-sheaves. Interpreting the latter
as the abelian sheaves of the topos $(X,G)$, that theory becomes part
of the theory of \Exp V, developed in the setting of general topoi.

More\pageoriginale generally, one may seek to associate an appropriate
topos with a topological space $X$ equipped with a topological group
$G$ (not necessarily discrete) of automorphisms, which would give rise
to an adequate cohomological theory. The same may be done in the
context of differentiable manifolds, real or complex analytic
manifolds, or schemes. This is indeed possible, \cf
\hyperref[IV.subsec2.5]{2.5} below.

\subsection{Classifying topos of a group object}\etiquette[2.4]{IV.subsec2.4}

Let $E$ be a topos, and let $G$ be a group object of $E$. Let $(E,G)$
be the category of objects of $E$ on which $G$ acts. It follows at
once from Giraud's criterion that this is a topos. It is called the
\emph{classifying topos} of the group object $G$, and is denoted by
$B_G$. When $E$ is the point topos (\hyperref[IV.subsec2.2]{2.2}),
\ie when $G$ is an ordinary group, we recover the classifying topos of
\hyperref[IV.subsec2.3]{2.3}.

The terminology adopted here is justified by the fact that the topos
$B_G$ plays a universal role in the classification of ``torsors'' (or
homogeneous principal bundles) under $G$, or more generally under
$G_{E'}=f^{\ast}(G)$, where $E'$ is a topos ``over $E$'', \ie endowed
with a morphism $f:E'\rightarrow E$ (\cf
\hyperref[IV.def3.1]{3.1} below). This role, made explicit in [3 Chap V]
or in \hyperref[IV.exer5.9]{5.9} below, shows that $B_G$ plays, in the
context of topoi, the same role as the classical classifying spaces of
topological groups in the homotopy theory of topological spaces. The
latter may be regarded (\cf \hyperref[IV.subsec2.5]{2.5}) as a weakened
version of the former, obtained by retaining from the classifying topos
only the ``homotopy type'' of that topos, in a suitable sense which
need not be specified here.

\subsection{``Big site'' and ``big topos'' of a topological space.
Classifying topos of a topological group}\footnote[1]{The introduction
of these sites and topoi is due to M. \textsc{Giraud}, who also brought
out their advantages over the traditional ``small'' site.}\etiquette[2.5]{IV.subsec2.5}\pageoriginale

Let $U-\Esp$, or simply $(\Esp)$, be the category of topological
spaces $\in 𝒰$. It is known that finite projective limits are
representable in $(\Esp)$. Consider on $(\Esp)$ the pretopology
(\hyperrefext[II][II.def1.3]{1.3}) for
which $\Cov(X)$ is the set of surjective families of open embeddings
$u_i:X_i\rightarrow X$. We regard $(\Esp)$ as a site by means of the
topology generated by the preceding pretopology. For every object $X$
of $(\Esp)$, regard the category
$$
(\Esp)_{/X}
$$
of objects of $(\Esp)$ over $X$, \ie topological spaces over $X$, as a
site, using the topology induced by that of $(\Esp)$ via the forgetful
functor $(\Esp)_{/X}\rightarrow (\Esp)$
(\hyperrefext[III][III.prop5.2]{5.2} 4)). This site is called the
\emph{big site} associated with $X$. One should note that it is not
$\in 𝒰$; nor is it a $𝒰$-site in the sense of
\hyperrefext[II][II.def3.0.2]{3.0.2}, so precautions are necessary when
applying the usual results to it. To remedy this inconvenience, one may
choose a universe $𝒱$ such that $𝒰\in 𝒱$, so that
$(\Esp)_{/X}$ becomes a $𝒱$-site, and work with the associated
$𝒱$-topos $(\Esp)_{/X}$, which may be denoted by $\TOP(X)$ and will
be called the \emph{big topos of} $X$. If one is reluctant to enlarge
$𝒰$, one may choose a cardinal $c$ bounding the cardinalities of $X$
and of all the topological spaces one intends to use in the arguments
(most often,
\sisi{$\Sup(\card X, \card ℛ)$}{$\Sup(\card X, \card \mathbb{R})$}
will suffice!), and replace $(\Esp)_{/X}$ by the subcategory
$(\Esp)'_{/X}$ consisting of those $X'$ over $X$ such that
$\card X'\leq c$, endowed with the induced topology; one then denotes
by $\TOP(X)$ the topos of sheaves on this\pageoriginale site. To fix
ideas, suppose that the first definition has been adopted.

The advantage of the big topos of $X$ over the small one is that the
site defining it contains $(\Esp)_{/X}$ as a full subcategory; since
the topology of this site is manifestly coarser than the canonical
topology, the canonical functor from $(\Esp)_{/X}$ to $\TOP(X)$,
associating with every space $X'$ over $X$ the sheaf it represents,
is seen to be \emph{fully faithful.} Consequently, \emph{a space $X'$
over $X$ is determined up to $X$-isomorphism by the sheaf
$(\in \Top(X))$ it defines}; thus the notion of a sheaf on (the big
site of) $X$ may be regarded as a \emph{generalization} of the notion
of a topological space over $X$, by means of which all constructions
of sheaf theory acquire a meaning for topological spaces over $X$.

Thus, when $G$ is a group object of the category $(\Esp)_{/X}$ of
topological spaces over $X$, one may associate with it the classifying
topos $B_G$ (\hyperref[IV.subsec2.4]{2.4}), hence classifying cohomology
groups, classifying homotopy groups, \etc (defined as the corresponding
invariants of the $𝒱$-topos $B_G$). In particular, when $X$ is a
one-point space, $G$ is identified with an ordinary topological group.
Subject to the usual local conditions ensuring that the singular
cohomology of the Cartesian products $G^n$ agrees with cohomology in
the sense of sheaves (for constant coefficients, say), for example if
$G$ is locally contractible, one may verify that the cohomology of the
classifying topos of $G$ is canonically isomorphic to that of the
classifying space of $G$ in the sense of topologists.

The introduction\pageoriginale of classifying topoi (via ``big
sites'') has over classifying spaces the advantage of providing a
richer theory, since in particular they provide useful cohomological
invariants for coefficients more general than constant or locally
constant coefficients. Moreover, the definition considered here
adapts in an evident way to the other customary contexts:
differentiable manifolds, real or complex analytic manifolds or
spaces, and schemes. This point of view makes it possible in particular
to connect the study of characteristic classes from the traditional
and the ``arithmetic'' points of view, by regarding the ``classical
groups'' as arising from schemes defined over the ring of integers;
\cf [\hyperref[IV.sec7]{7}] for indications in this direction.
Likewise, the general results of J. \textsc{Giraud}
[\hyperref[IV.sec3]{3}] on the classification of extensions of group
objects, developed in the very general and flexible setting of topoi,
may, by means of ``big topoi'', be specialized to results on the
classification of extensions of topological groups, or of real or
complex Lie groups, which seem to have been scarcely known to
topologists except in the case of extensions with abelian kernel
[\hyperref[IV.sec11]{11}].

\subsection{Topoi of the form $\widehat{C}$}\etiquette[2.6]{IV.subsec2.6}

Let $C$ be a small category. Then the category $\widehat{C}$ of
$U$-presheaves on $C$ is plainly a $𝒰$-topos, since it is of the
form $C^{\widetilde{\phantom{,}}}$, where $C$ is endowed with the
chaotic topology. We shall give below some details on the relations
between $C$ and $\widehat{C}$. Let us note here only that a topos $E$
equivalent to a topos of the form $\widehat{C}$ is of a rather special
nature, because \emph{it has a small generating family $(X_i)$
consisting of connected projective objects}, \ie objects $X$ such
that the functor $Y\mapsto \Hom(X,Y)$\pageoriginale takes epimorphisms
to epimorphisms and coproducts to coproducts: indeed, it suffices to
take in $\widehat{C}$ the generating family consisting of the functors
represented by the $X\in \ob C$. Moreover, note that if in a topos
$E$ one has a covering family $X_i\rightarrow X$, where the $X_i$ are
projective and connected, then every other covering family of $X$ is
refined by the former (\hyperrefext[I][I.sec4.3.2]{4.3.2},
\hyperrefext[I][I.sec4.3.3]{4.3.3}). Consequently, in a topos $E$ of
the form $\widehat{C}$ every object $X$ admits a covering family
refining all the others. A topos of the form $\Top(X)$
(\hyperref[IV.subsec2.1]{2.1}), where $X$ is a topological space whose
points are closed, has the preceding property only if $X$ is discrete.

When the category $C$ has a single object, $C$ is identified with a
monoid $G$. A presheaf on $C$ is then identified with a set on which
$G$ acts \emph{on the right} (since it is a functor
$G^{\circ}\rightarrow (\Ens)$), and $\widehat{C}$ is the topos of
sets with a monoid of right operators, which may also be denoted by
$B_{G^{\circ}}$ in view of \hyperref[IV.subsec2.3]{2.3}: it is the
\emph{topos of sets with monoid of operators $G^{\circ}$ (the monoid
opposite to $G$)}. When $G$ is a group, using the isomorphism
$g\mapsto g^{-1}$ from $G$ onto $G^{\circ}$, one recovers the
classifying topos $B_G$ of \hyperref[IV.subsec2.3]{2.3}.

\subsection{Classifying topos of a pro-group}\etiquette[2.7]{IV.subsec2.7}

\subsubsection{}\etiquette[2.7.1]{IV.subsubsec2.7.1}
Let $𝒢=(G_i)_{i\in I}$ be a projective system of groups, with
$G_i$, $I\in 𝒰$. Assume that the projective system is
\emph{strict}, \ie that the transition morphisms $G_j\rightarrow G_i$
are surjective. If $E$ is a set, by a (left, say) \emph{action of
$𝒢$ on $E$} we mean the following structure: a) a family
$(E_i)_{i\in I}$ of subsets of $E$, whose union is $E$; b) for every
$i\in I$, an action of the group $G_i$ on the set $E_i$; these data
are furthermore required to satisfy\pageoriginale the following
condition: for $j\geq i$, $E_i$ is the subset of $E_j$ consisting of
the elements fixed under the kernel of $G_j\rightarrow G_i$. We also
say that $𝒢$ \emph{acts on} $E$ (on the left) if a (left) action of
$𝒢$ on $E$ has been given. The sets $\in 𝒰$ endowed with an
action of $𝒢$ form a category in the evident way. By Giraud's
criterion, one sees at once that this category is a $𝒰$-topos. It
is denoted by $B_{𝒢}$ and called the \emph{classifying topos of}
$𝒢$. When $I$ has an initial object $i_{\circ}$, on putting
$G=G_{i_\circ}$ one recovers the classifying topos of
\hyperref[IV.subsec2.3]{2.3}.

\subsubsection{}\etiquette[2.7.2]{IV.subsubsec2.7.2}
Another important example is that in which the groups $G_i$ are
finite, so that
$$
G=\varprojlim G_i
$$
is a compact totally disconnected topological group, or
\emph{profinite group.} An action of $𝒢$ on $E$ then amounts to a
continuous action of $G$ on $E$, or, equivalently, to one for which
the stabilizer of every point of $E$ is an open subgroup of $G$.
The classifying topos $B_{𝒢}$ will also be denoted by $B_{G}$, where,
of course, $G$ is to be regarded as endowed with its profinite
topology.

\subsubsection{}\etiquette[2.7.3]{IV.subsubsec2.7.3}
Using the remarks of \hyperref[IV.subsec2.6]{2.6}, it is easy to verify
that the topos $B_{𝒢}$ defined by a strict projective system
$𝒢=(G_i)_{i\in I}$ of groups is equivalent to a topos of the form
$\widehat{C}$ only if this projective system is essentially constant;
in the case of a profinite group, this means that the group is in fact
finite.

\subsubsection{}\etiquette[2.7.4]{IV.subsubsec2.7.4}
The following geometric interpretation of the classifying topos $B_G$
of a discrete group $G$ is useful for giving a correct geometric
intuition for these topoi. (\Cf also, in the same vein,
\hyperref[IV.subsec4.5]{4.5}, \hyperref[IV.subsec5.8]{5.8},
\hyperref[IV.exer5.9]{5.9}, and \hyperref[IV.subsec7.2]{7.2} below.)
Let $X$ be a connected, locally connected, and\pageoriginale locally
simply connected topological space, let $x$ be a point of $X$, and let
$G$ be its fundamental group at $x$. (NB It is known that, up to
isomorphism, every discrete group $G$ may be obtained in this way.)
Then the Galois theory of coverings of $X$ gives an equivalence between
the category $B_G$ of $G$-sets and the category of \emph{étale
coverings} of $X$, \ie spaces $X'$ over $X$ that are locally
$X$-isomorphic to $X$-spaces of the form $X\times I$, where $I$ is a
discrete space. (Compare \SGA 1 V 4,5.) The latter can moreover be
interpreted as the category of locally constant sheaves on $X$, \ie
the category of locally constant objects (IX \hyperref[IX.subsec2.0]{2.0}) of the topos
$\Top(X)$.

When $G$ is a profinite group, there is an analogous geometric
interpretation of $B_G$ as the category of $X$-schemes that are
coproducts of finite étale coverings of a connected scheme $X$,
equipped with a geometric point $x$ and an isomorphism
$G\simeq \pi_1(X,x)$. It is again known that every profinite group may
be obtained as the fundamental group of a suitable connected scheme
(the spectrum of a field, if desired). Finally, pro-groups (not
necessarily profinite or essentially constant) also occur in the
classification of coverings of connected and locally connected spaces
that are not locally simply connected, or in the classification of
étale coverings, not necessarily finite or ind-finite, of connected
nonnormal schemes. For the latter case, \cf \SGA 3 X 6.

\begin{enonce*}[remark]{Exercise 2.7.5}\etiquette[2.7.5]{IV.exer2.7.5}
For a topos $E$, define the notions of connectedness, local
connectedness~\footnote[1]{\cf
\sisi{\hyperref[IV.exer8.7]{8.7} 1).}{\hyperref[IV.exer8.7]{8.7} l).}},
simple connectedness, and local simple connectedness. Define the notion
of a constant and locally constant object of $E$ (\cf IX \hyperref[IX.subsec2.0]{2.0}). Let
$f:E'\rightarrow E$\pageoriginale be a morphism of topoi
(\hyperref[IV.def3.1]{3.1}), with $E'$ simply connected (for example,
$E'$ the point topos (\hyperref[IV.subsec2.2]{2.2})), and $E$ connected
and locally connected. Define a strict pro-group
$\pi_1(E,f)=(G_i)_{i\in I}=𝒢$ (called the \emph{fundamental
pro-group of $E$ at $f$}) and an equivalence of categories between
$B_{𝒢}$ and the category of locally constant objects of
$E$.\footnote[2]{One may take inspiration from \SGA 3 X 6.} Show that,
when $E$ is locally simply connected, $\pi_1(E,f)$ is essentially
constant and hence is identified with an ordinary discrete group
$\pi_1(E,f)$, called the \emph{fundamental group of $E$ at $f$}.
Show that every strict pro-group $𝒢$ is isomorphic (as a pro-group)
to the fundamental pro-group of a suitable connected and locally
connected topos at a suitable $f$, with $E'$ the point topos (take
$E=B_{𝒢}$, and $f:E'\rightarrow E$ defined by the forgetful functor
$f^{\ast}:E\rightarrow E'=(\Ens)$). When $𝒢$ is essentially
constant, \ie isomorphic (as a pro-group) to an ordinary discrete
group, prove that one may take $E$ above to be locally simply connected
(again take $E=B_G$).
\end{enonce*}

\subsection{Example of a false topos}\etiquette[2.8]{IV.subsec2.8}

Let $𝒢=(G_i)_{i\in I}$ be a strict pro-group, where $I$ is a
filtered ordered set and where $i>j$ implies that
$G_i\rightarrow G_j$ is not an isomorphism. Suppose that
$\card(I)\not\in 𝒰$. Consider the category of sets $E\in 𝒰$
on which $𝒢$ acts on the left
(\hyperref[IV.subsubsec2.7.1]{2.7.1}). It is a $𝒰$-category, and one
sees as in \hyperref[IV.subsubsec2.7.1]{2.7.1} that this category
satisfies conditions a), b), and c) of
\hyperref[IV.subsubsec1.1.2]{1.1.2}. Nevertheless, it is not a
$𝒰$-topos, since it has no generating family that is
$𝒰$-small. Similarly, it is not a $𝒱$-topos for any universe
$𝒱$.

% END INCLUDED SGA4_English_current_buildable_source_20260906/source/04/drafts/SGA4_IV_LINES_537_947_English.texfrag


% BEGIN INCLUDED SGA4_English_current_buildable_source_20260906/source/04/drafts/SGA4_IV_LINES_948_1319_English.texfrag
\section{Morphisms of topoi}\etiquette[3]{IV.sec3}\pageoriginale

\begin{enonce*}{Definition 3.1}\etiquette[3.1]{IV.def3.1}
Let $E$ and $E'$ be two $𝒰$-topoi. A morphism\nde{Some authors
speak of a \emph{geometric morphism}.} from
$E$ to $E'$, or sometimes (by abuse of language) a continuous map
from $E$ to $E'$, is a triple $u=(u_{\ast}, u^{\ast},\varphi)$,
consisting of functors
$$
u_{\ast}:E\longrightarrow E'\quav u^{\ast}:E'\longrightarrow E
$$
and an ``adjunction'' isomorphism $\varphi$ of bifunctors
in $X'\in \Ob E'$, $Y\in \Ob E$:
$$
\varphi:\Hom_E (u^{\ast}(X'),Y)\xrightarrow{\sim} \Hom_{E'}(X',
u_{\ast}(Y)){\quoi},
$$
where the functor $u^{\ast}$ is moreover required to be left exact,
\ie to commute with finite projective limits. The functor $u_{\ast}$
is called the direct-image functor for the morphism of topoi $u$,
the functor $u^{\ast}$ is called the inverse-image functor for the
morphism of topoi $u$, and the isomorphism $\varphi$ is called the
adjunction isomorphism for $u$.
\end{enonce*}

\setcounter{subsection}{1}
\subsubsection{}\etiquette[3.1.1]{IV.subsubsec3.1.1}
Unless expressly stated otherwise, for a morphism of topoi
$u:E\rightarrow E'$, we shall henceforth denote the corresponding
direct- and inverse-image functors by $u_{\ast}$ and $u^{\ast}$
\footnote[1]{It is sometimes appropriate also to write $u^{-1}$
instead of $u^{\ast}$, \cf \hyperref[IV.subsubsec13.2.3]{13.2.3}.}.
Note that, since $u_{\ast}$ is right adjoint to $u^{\ast}$ and
$u^{\ast}$ is left adjoint to $u_{\ast}$ by the adjunction
isomorphism $\varphi$, each of the two functors $u_{\ast}$,
$u^{\ast}$ determines the other up to unique isomorphism, by the
well-known general facts about adjoint functors [\hyperref[IV.sec14]{14}].
In practice, depending on the circumstances, it may be more convenient
to define a morphism of topoi $u:E\rightarrow E'$ either by specifying
$u^{\ast}:E'\rightarrow E$, or by specifying
$u_{\ast}:E\rightarrow E'$; in the first case, one need only verify
that the given functor $u^{\ast}$ has\pageoriginale a right adjoint
and is left exact. In the second case, one must verify that the given functor
$u_{\ast}$ has a left adjoint which is left exact. In either case,
from the partial data, by \emph{choosing} an adjoint functor and an
adjunction isomorphism, one obtains a morphism of topoi
$u:E\rightarrow E'$, and the latter will be ``unique up to unique
isomorphism'' in terms of the datum $u^{\ast}$ \resp $u_{\ast}$, in
a sufficiently clear sense which will moreover be made completely
explicit below \hyperlink{sga4.exp.IV.subsubsection.3.2.1}{(3.2.1)}.

\subsubsection{}\etiquette[3.1.2]{IV.subsubsec3.1.2}
If $u:E\rightarrow E'$ is a morphism of topoi, it follows from the
properties of adjoint functors (\hyperrefext[I][I.prop2.11]{2.11}) that
the direct-image functor $u_{\ast}:E\rightarrow E'$ commutes with
projective limits, and the functor $u^{\ast}:E'\rightarrow E$ commutes
with inductive limits \footnote[2]{Moreover
(\hyperref[IV.cor1.5]{1.5} and \hyperref[IV.rem1.8]{1.8}), for a given functor
$u_{\ast}:E\rightarrow E'$ \resp $u^{\ast}:E'\rightarrow E$,
this functor has a left adjoint (\resp a right adjoint) if and only if
it commutes with $𝒰$-projective limits (\resp $𝒰$-inductive
limits).}\MakeLinkTarget*{sga4.exp.IV.footnote.3.1.1.2}\label{sga4[1]}. Since the latter is assumed in addition to
be left exact, \ie to commute with finite projective limits, we see in
particular that $u^{\ast}$ \emph{is exact}. We may therefore say that
\emph{it is the inverse-image functor $u^{\ast}$}, in the pair
$(u_{\ast},u^{\ast})$, \emph{which has the most remarkable exactness
properties.} These properties ensure that, for every type of
algebraic structure $\Sigma$ whose data can be described in terms of
arrows between the underlying sets and sets deduced from them by
repeated application of operations of finite projective limits and
arbitrary inductive limits, and for every ``object of $ℰ'$ endowed
with a structure of type $\Sigma$'' (a notion which makes sense by
virtue of the internal exactness properties of the topos $ℰ'$
(\hyperrefext[II][II.theo4.1]{4.1})), its image under $u^{\ast}$ is endowed
with the same structures. Rather than undertake the unrewarding task
of giving a precise meaning to this statement and formally justifying
it,\pageoriginale we advise the reader to spell it out and convince
himself of its validity for types of structure such as groups, rings,
modules over a ring, comodules over a coring, bialgebras over a ring,
and torsors under a group. (In these examples, the first three types
of structure are defined exclusively in terms of finite projective
limits, while the other notions implicitly involve constructions
which also use inductive limits.) Furthermore, the functor
$u^{\ast}$ ``\emph{commutes}'' with all the usual functorial
operations performed in terms of such structures, more precisely
with all operations which can be expressed in terms of
$\varinjlim$ and finite $\varprojlim$: constructions of free objects
(free groups or modules, for example) generated by an object, tensor
products (\cf \hyperref[IV.sec12]{\S 12} below), \etcn

As for the direct-image functor $u_{\ast}:E\rightarrow E'$, since it
commutes with projective limits, it consequently ``preserves'' every
algebraic structure on an object (or a family of objects) of $E$
which can be defined exclusively in terms of projective limits (such
as the structures of group, ring, or module over a ring among the
preceding examples). On the other hand, the functor $u_{\ast}$ is not
in general right exact, \ie it does not in general commute with finite
inductive limits, and it does not even in general take epimorphisms to
epimorphisms (it is, moreover, this failure of exactness of
$u_{\ast}$ which is the source of its cohomological properties; these
will be studied, from the viewpoint of commutative homological
algebra, in the following exposé). Consequently it does not, in
general, extend to a functor on objects of the type comodule,
\pageoriginale bialgebra, or torsor under a group, and does not in
general commute with operations such as ``free module generated by'',
tensor product of modules, \etc

\subsubsection{}\etiquette[3.1.3]{IV.subsubsec3.1.3}
In practice, when one is given a functor $f:E\rightarrow F$ from one
$𝒰$-topos to another, one should make all of its exactness properties
explicit, including the possible existence of left or right adjoint
functors (\cf the footnote on page \pageref{sga4[1]}), in order to
understand the ``geometric nature'' of $f$, an understanding which
will generally be an indispensable guide to a correct geometric
intuition of the situation. Thus, if it turns out that $f$ commutes
with arbitrary inductive limits and finite projective limits, one
should write $f$ in the form
$$
f=u^{\ast}{\quoi},
$$
where
$$
u:F\longrightarrow E
$$
is a morphism of topoi, \ie interpret $f$ as an ``inverse-image''
functor for a ``continuous map'' of topoi. When $f$ commutes with
arbitrary projective limits, so that it has a left adjoint, and if
the latter (which a priori commutes with arbitrary inductive limits)
\emph{also} commutes with finite projective limits, one should write
$f$ in the form
$$
f=v_{\ast}{\quoi},
$$
where
$$
v:E\longrightarrow F
$$
is a morphism of topoi. In certain cases it may happen that $f$
satisfies\pageoriginale both of the properties under consideration
(\cf \hyperref[IV.subsec4.10]{4.10} for an example). In that case one
should introduce simultaneously the morphisms of topoi
$$
u:F\longrightarrow E \quav v:E\longrightarrow F{\quoi},
$$
which give rise to a sequence of three adjoint functors
(\hyperrefext[I][I.sec5.3]{5.3}):
$$
v^{\ast} \quav v_{\ast}=u^{\ast}=f\quav u_{\ast}{\quoi}.
$$
One must then distinguish carefully between the morphisms of topoi
$u$ and $v$, lest one lose the geometric intuition of the situation.

In this connection, note that when between two topoi $E$, $F$ one has
a sequence of three adjoint functors
$$
e, f, g\; (e,g: F\rightrightarrows E, f:E\rightarrow F){\quoi},
$$
then $f$ commutes with arbitrary inductive and projective limits, and
therefore can always be written in the form $u^{\ast}$, where
$u:F\rightarrow E$ is a morphism of topoi. Thus $g$ can be written
$g=u_{\ast}$. On the other hand, of course, $e$ can be written in the
form $v^{\ast}$ (and then $f$ in the form $v_{\ast}$) only if it also
commutes with finite projective limits. This will evidently be the
case if $e$ is itself the right adjoint of a fourth functor $d$.
Without a condition of this kind, one will often write
$$
e=u_{!}
$$
for the left adjoint of an inverse-image functor $f=u^{\ast}$ when
such a left adjoint exists, this notation being suggested by the
example of an open immersion $u:X\rightarrow Y$ of topological
spaces. Similarly, if $e$ is of the form $v^{\ast}$, \ie if $f$ is
of the form $v_{\ast}$, one sometimes writes $g=v^{!}$ for the right
adjoint of a direct-image functor $v_{\ast}$, when this adjoint exists
\footnote[1]{\Cf below \hyperref[IV.subsec14.4]{14.4}, in the case of
abelian sheaves.}.

\subsection{}\etiquette[3.2]{IV.subsec3.2}
Let\pageoriginale $E$, $E'$ be two $𝒰$-topoi, and let
$$
u=(u_{\ast},u^{\ast},\varphi)\quav
v=(v_{\ast},v^{\ast},\Psi):E\rightrightarrows E'
$$
be two morphisms of topoi from $E$ to $E'$. A \emph{morphism from
$u$ to $v$} is any morphism from $u_{\ast}$ to $v_{\ast}$ (in the
sense of the category $\Hhom(E,E')$ of functors from $E$ to $E'$).
Morphisms of morphisms of topoi compose in the evident way, thereby
defining a category, which is in fact a
$\sheaf{\mathcal{U}}$-\emph{category} (\hyperrefext[I][I.prop7.8]{7.8}),
denoted
$$
\Hhomtop(E,E'){\quoi},
$$
and called the \emph{category of morphisms} (or \emph{continuous
maps}) \emph{from $E$ to $E'$}. One then has an evident functor
$$
u\longmapsto u_{\ast}:\Hhomtop(E,E')\longrightarrow
\Hhom(E,E'){\quoi},
$$
which is fully faithful by the definition of arrows in the first
member (but is not injective on objects).

\hypertarget{sga4.exp.IV.subsubsection.3.2.1}{}\subsubsection{}
Care must be taken that, if $u$ and $v$ are given as above, the theory
of adjoint functors supplies a canonical bijection
$$
\Hom(u_{\ast},v_{\ast})\xrightarrow{\sim}\Hom(v^{\ast},
u^{\ast}){\quoi};
$$
in particular, one obtains a \emph{contravariant functor} on
$\Hhomtop(E,E')$:
$$
u\longmapsto u^{\ast}:\Hhomtop(E,E')^{\circ}\longrightarrow
\Hhom(E,E'){\quoi}.
$$
In accordance with geometric intuition, according to which the
direct-image functor $u_{\ast}$ ``goes in the same direction'' as the
continuous map which gives rise to it, one should therefore define
the direction of arrows for morphisms between morphisms of topoi
\emph{in terms of the direct-image functors}, and not in terms of the
inverse-image functors (although the latter, as we have seen, possess
the exactness properties characteristic of the notion of morphism of
topoi).

\subsubsection{}\etiquette[3.2.2]{IV.subsubsec3.2.2}
Having\pageoriginale defined the category $\Hhomtop(E,E')$ of
morphisms from the topos $E$ to the topos $E'$, the notion of an
\emph{isomorphism} between two morphisms $u$, $v$ from $E$ to $E'$
is also defined. In practice there is no need to distinguish
essentially between two isomorphic morphisms of topoi, at least when
one has a \emph{canonical} isomorphism between the two (just as there
is often no need to distinguish between two objects of a category
when a canonical isomorphism between them is given). Let us point out
in this connection that, most often, when one deals with diagrams of
morphisms of topoi and questions of commutativity of such diagrams
(a notion which makes sense by \hyperlink{sga4.exp.IV.subsection.3.3}{(3.3)}), one means only commutativity up
to (``canonical'') isomorphism; by abuse of language, one then treats
these diagrams as actually commutative diagrams.

One might try to justify this abuse of language by introducing the
set \\ $\Homtop(E,E')/\Isom$ of morphisms from $E$ to $E'$ up to
isomorphism, and calling such an isomorphism class a morphism, instead
of following Definition \hyperref[IV.subsec3.2]{3.2}. But this runs
into the very serious difficulties which arise whenever one tries to
identify two isomorphic objects of a category without having a
canonical isomorphism between them. Experience shows that such a
viewpoint is impracticable, and that one must retain the ``fine''
notion \hyperref[IV.def3.1]{3.1} of a morphism between morphisms of
topoi, even if one is sometimes obliged to struggle with
compatibilities between canonical isomorphisms \footnote[1]{For
examples of such struggles (apparently successful), we refer the
reader to Mme M. \textsc{Hakim}'s book on relative schemes [9].}.

\hypertarget{sga4.exp.IV.subsection.3.3}{}\setcounter{subsection}{3} \setcounter{subsubsection}{0}
\subsubsection{}\etiquette[3.3.1]{IV.subsubsec3.3.1}
Let\pageoriginale $E$, $E'$, $E''$ be three $𝒰$-topoi, and consider
morphisms of topoi
$$
u:E\longrightarrow E'\quav v:E'\longrightarrow E''{\quoi}.
$$
The theory of adjoint functors then gives us an adjunction
isomorphism between the composite functors $v_{\ast}u_{\ast}$ and
$u^{\ast}v^{\ast}$, in terms of the adjunction isomorphisms for the
pairs $(u_{\ast},u^{\ast})$ and $(v_{\ast},v^{\ast})$. On the other
hand, the functor $u^{\ast}v^{\ast}$ is left exact, as the composite
of two left exact functors. Consequently one obtains a morphism from
$E$ to $E''$, called the \emph{composite of the morphisms $u$ and
$v$}, and denoted $vu$:
$$
vu:E\longrightarrow E''{\quoi}.
$$
One verifies at once that composition of morphisms is associative,
and that for every $𝒰$-topos $E$ there is a morphism from $E$ to
itself which is a two-sided identity for composition: it is the
morphism $(\id_E,\id_E,\varphi)$, where $\varphi$ is the evident
adjunction isomorphism of $\id_E$ with itself. Let $𝒱$ be a universe
such that $𝒰\in 𝒱$. One defines a category
$$
(𝒱\text{-}𝒰\text{-}\Top){\quoi},
$$
whose objects are the $𝒰$-topoi which belong to $𝒱$, whose arrows
are morphisms between such $𝒰$-topoi, and whose composition of arrows
is the one just described.

\subsubsection{}\etiquette[3.3.2]{IV.subsubsec3.3.2}
In fact, the composition map
$$
\Homtop(E, E')\times \Homtop(E',E'')\longrightarrow \Homtop(E,E'')
$$
is the map induced on objects by a ``functor of composition of
morphisms'':
$$
\Hhomtop(E,E')\times \Hhomtop(E',E'')\longrightarrow
\Hhomtop(E,E''){\quoi},
$$
whose\pageoriginale effect on arrows is the usual ``convolution
product'' operation for morphisms between functors (here,
direct-image functors). These composition functors satisfy a (strict)
associativity property, for four topoi $E$, $E'$, $E''$, $E'''$,
which refines the associativity of composition of morphisms of topoi.
One may also say, in language which is beginning to become familiar
[2], [9], that the $𝒰$-topoi are the objects (or 0-arrows) of a
2-\emph{category}, whose 1-arrows are the morphisms of topoi and
whose 2-arrows are morphisms of morphisms of topoi.

It is the fact that the $𝒰$-topoi (elements of a universe $𝒱$) form
a 2-category, rather than merely an ordinary category like ordinary
topological spaces, which from the technical viewpoint constitutes
the most important difference between the theory of topoi and that
of topological spaces. This fact is the source of certain technical
complications already alluded to, but also of facts essentially new
in comparison with traditional topology.

\subsection{}\etiquette[3.4]{IV.subsec3.4}
The fact that the $𝒰$-topoi (elements of a universe $𝒱$) form a
2-category (\hyperref[IV.subsubsec3.3.2]{3.3.2}) makes it possible in
particular to define the notion of \emph{equivalence of two
$𝒰$-topoi $E$, $E'$}: one says that $E$ and $E'$ are
\emph{equivalent} if there exist morphisms of topoi
$u:E\rightarrow E'$ and $v:E'\rightarrow E$ such that the composites
$uv$ and $vu$ are isomorphic respectively to the identity morphisms
of $E$ and $E'$; the morphisms $u$ and $v$ are then said to be
\pageoriginale \emph{quasi-inverse equivalences of one another.}

One sees at once that for the morphism of topoi $u:E\rightarrow E'$
to be an equivalence, it is necessary and sufficient that $u_{\ast}$
be an equivalence, or, what amounts to the same thing, that
$u^{\ast}$ be an equivalence. (Use the fact that a functor
$f:E\rightarrow E'$ between two topoi which is an equivalence is
both of the form $u_{\ast}$ and of the form $v^{\ast}$, with $u$ and
$v$ morphisms of topoi); and for $E$ and $E'$ to be equivalent in
the sense of the preceding paragraph, it is necessary and sufficient
that they be equivalent as categories (\ie as objects of the
2-category $(𝒱\text{-}\cat)$). As it should, this shows that the
notions of equivalence introduced do not depend on the choice of the
universe $𝒱$ in \hyperref[IV.subsubsec3.3.1]{3.3.1}.

\subsubsection{}\etiquette[3.4.1]{IV.subsubsec3.4.1}
In practice, there is most often no reason to distinguish essentially
between equivalent $𝒰$-topoi, just as there is often no reason to
distinguish essentially between two equivalent categories---provided,
however, that one has an explicit equivalence from one to the other,
or at least an equivalence defined up to unique isomorphism. Here it
is the notion of equivalence of topoi which replaces the traditional
notion of homeomorphism between two topological spaces. See the
example in \hyperref[IV.subsec4.2]{4.2} below for the precise relations
between these two notions.

% END INCLUDED SGA4_English_current_buildable_source_20260906/source/04/drafts/SGA4_IV_LINES_948_1319_English.texfrag


% BEGIN INCLUDED SGA4_English_current_buildable_source_20260906/source/04/drafts/SGA4_IV_LINES_1320_1659_English.texfrag
\section{Examples of morphisms of topoi}\etiquette[4]{IV.sec4}

Here we return to the examples of 2, using the notion of a
morphism of topoi. The comments of \hyperref[IV.subsec2.0]{2.0}
apply equally to the present section. The universe
$𝒰$ will generally be understood.

\subsection{The topos $\Top(X)$ for a varying topological space $X$}
\etiquette[4.1]{IV.subsec4.1}\pageoriginale

\subsubsection{}\etiquette[4.1.1]{IV.subsubsec4.1.1}
Let
$$
f:X\longrightarrow Y
$$
be a continuous map of topological spaces. We shall associate with it
canonically a morphism of topoi
$$
\Top(f)\mbox{ or } f:\Top(X)\rightarrow \Top(Y){\quoi},
$$
with the notation of \hyperref[IV.subsec2.1]{2.1}. When $\Top(X)$ is
defined as $\Ouv(X)^{\widetilde{\phantom{,}}}$, the most convenient
description of $\Top(f)$ is by means of the direct-image functor on
sheaves
$$
f_{\ast}:\Top(X)\longrightarrow \Top(Y){\quoi},
$$
defined by the formula
$$
f_{\ast}(F)=F\circ f^{-1}{\quoi},
$$
where
$$
f^{-1}:\Ouv(Y)\longrightarrow \Ouv(X)
$$
is the evident functor $U\rightarrow f^{-1}(U)$. We have already noted
that this functor is continuous and left exact, and therefore indeed
defines the functor $f_{\ast}$ above, which admits a left adjoint
$f^{\ast}$ that is left exact (\hyperrefext[III][III.exem1.9.1]{1.9.1}).
Of course, strictly speaking, the morphism of topoi $\Top(f)$ depends
on the choice of the left adjoint $f^{\ast}$ of $f_{\ast}$ and hence
is defined only up to canonical isomorphism. In what follows we shall
refrain from explicitly pointing out phenomena of this kind.

When one adopts the ``étale spaces'' point of view in defining
$\Top(X)$, it is the inverse-image functor
$$
f^{\ast}:\Top(Y)\longrightarrow \Top(X)
$$
that\pageoriginale is the most convenient for defining the morphism of
topoi $\Top(f)$, by simply setting
$$
f^{\ast}(Y')=\fprod{X}{Y'}{Y}
$$
for every étale space $Y'$ over $Y$; it is clear that the fiber product
is indeed an étale space over $X$, and that the functor $f^{\ast}$ so
obtained is left exact and commutes with arbitrary $\varinjlim$, and
therefore defines a morphism of topoi $\Top(f)$. For the compatibility
of the two definitions obtained, we refer to [TF].

For two composable continuous maps
$$
X\xrightarrow{f}Y\xrightarrow{g}Z{\quoi},
$$
there is a canonical isomorphism
$\Top(gf)\simeq\Top(g)\Top(f)$ of morphisms of topoi. These
\emph{transitivity isomorphisms}, for three composable continuous maps
$f$, $g$, $h$, satisfy a compatibility relation which we shall refrain
from writing here, and which is precisely the one considered in
\SGA 1 VI 7.4 B) (for $\epsilon = (\Esp)^{\circ}$). This may be
expressed by saying that, as $X$ varies in the category $(\Esp)$,
$$
X\longmapsto \Top(X)
$$
is a ``pseudofunctor''
\begin{equation*}
(𝒰\text{-}\esp)\longrightarrow
(𝒱\text{-}𝒰\text{-}\top){\quoi},\tag{4.1.1.1}
\end{equation*}\etiquette[4.1.1.1]{IV.eq4.1.1.1}
or, in the terminology of 2-categories, that one has a
\emph{non-strict} functor of 2-categories [9]. In practice, we shall
most often permit ourselves the abuse of language consisting in
identifying $\Top(gf)$ and $\Top(g)\Top(f)$, that is, reasoning as if
\eqhyperref[IV.eq4.1.1.1]{4.1.1.1} were an actual functor between
ordinary categories. We shall permit ourselves analogous abuses of
language in the other examples treated below.\pageoriginale

\subsubsection{}\etiquette[4.1.2]{IV.subsubsec4.1.2}
The preceding considerations extend immediately to topoi associated
with topological spaces equipped with operator groups
(\hyperref[IV.subsec2.3]{2.3}). If $f=(f^{\es},f^{\gr})$,
$$
f:(X,G)\longrightarrow (Y,H)
$$
is a morphism of spaces with operators (where
$$
f^{\es}:X\longrightarrow Y\quav f^{\gr}:G\longrightarrow H
$$
are respectively continuous maps and group homomorphisms, compatible
in the evident sense), one associates with it a morphism of topoi
$$
\Top(f)\mbox{ or } f:\Top(X,G)\longrightarrow \Top(Y,H){\quoi},
$$
whose definition is left to the reader. When the groups $G$ and $H$
are trivial, one recovers the definition of
\hyperref[IV.subsubsec4.1.1]{4.1.1}; when, on the other hand, the spaces
$X$ and $Y$ are each reduced to a point, the inverse-image functor
$f^{\ast}$ is the ``restriction of the operator group'' functor
$$
f^{\ast}:B_H\longrightarrow B_G{\quoi},
$$
which associates with every $H$-set the $G$-set that it defines by
means of $f:G\rightarrow H$. We shall encounter this example again in
other forms in \hyperref[IV.subsec4.5]{4.5} and
\hyperref[IV.subsubsec4.6.1]{4.6.1}.

\subsubsection{}\etiquette[4.1.3]{IV.subsubsec4.1.3}
Given a continuous map $f:X\rightarrow Y$ of topological spaces, one
also associates with it a morphism between the corresponding ``large
topoi'' (\hyperref[IV.subsec2.5]{2.5})
$$
\TOP(f)\mbox{ or } f:\TOP(X)\longrightarrow \TOP(Y){\quoi},
$$
most conveniently defined\pageoriginale by the inverse-image functor
$$
f^{\ast}:\TOP(Y)\longrightarrow \TOP(X){\quoi},
$$
which is simply the \emph{restriction functor}. This morphism
$\TOP(f)$ is a special case of the morphism called an ``inclusion''
for an induced topos, which will be studied in 5. Thus, subject to the
same qualifications as in \hyperref[IV.subsubsec4.1.1]{4.1.1}, the
topos $\TOP(X)$ may be regarded as a functor of $X$, as $X$ varies in
$(\Esp)$.

\subsection{Fidelity properties of $X\mapsto \Top(X)$}\etiquette[4.2]{IV.subsec4.2}

We propose to specify to what extent a topological space $X$ can be
reconstructed in terms of the topos $\Top(X)$. For this purpose it is
appropriate, for two spaces $X$ and $Y$, to describe the category of
morphisms from $\Top(X)$ to $\Top(Y)$
(\hyperref[IV.subsec3.2]{3.2}), so as to determine the fidelity
properties of the ``functor'' $X\mapsto \Top(X)$. We shall merely
state the results obtained, referring the reader to [9] for details.
The reader who wishes to carry out the verification himself may
consult Exercise \hyperref[IV.exer7.8]{7.8}.

\subsubsection{}\etiquette[4.2.1]{IV.subsubsec4.2.1}
Recall that a topological space $X$ is called \emph{sober} if every
irreducible closed subset of $X$ has exactly one generic point. Almost
all spaces used in practice are sober; this is true in particular of a
Hausdorff space, more generally of a space all of whose points are
closed, or of the underlying space of a scheme. With a topological
space $X$ one associates (\loccit or \EGA $0_I$, new edition) a sober
topological space $X_{\sob}$ and a continuous map\pageoriginale
\begin{equation*}
\varphi: X\longrightarrow X_{\sob}\tag{4.2.1.1}
\end{equation*}\etiquette[4.2.1.1]{IV.eq4.2.1.1}
which is \emph{universal} for continuous maps from $X$ to sober spaces;
in other words, one constructs a left adjoint $X\mapsto X_{\sob}$ to
the inclusion functor $(\Espsob)\rightarrow (\Esp)$ from the category
of sober spaces to that of ``arbitrary'' topological spaces (the
quotation marks serving as a reminder that a universe is involved!).
The explicit construction takes as the points of $X_{\sob}$ the
irreducible closed subsets of $X$, and as its open subsets the sets of
the form $U'$, where $U$ is an open subset of $X$ and
$U'\subset X_{\sob}$ denotes the set of irreducible closed subsets of
$X$ that meet $U$. The map
\eqhyperref[IV.eq4.2.1.1]{4.2.1.1} is obtained by associating with each
$x\in X$ the closure of $\{x\}$. The space $X$ is sober if and only if
the preceding map is bijective, and hence a homeomorphism.

One checks that the functor
$$
\varphi^{-1}:\Ouv(X_{\sob})\longrightarrow \Ouv(X)
$$
induced by $\varphi$ is an isomorphism, which implies that the morphism
of topoi
$$
\Top(\varphi):\Top(X)\longrightarrow \Top(X_{\sob})
$$
defined by $\varphi$ is also an isomorphism. This explains a priori
why $X_{\sob}$ must necessarily enter into the question of
reconstructing $X$ from $\Top(X)$: since the latter depends only on
$X_{\sob}$ up to isomorphism, the question can have an affirmative
answer only if $X$ is sober. We shall specify below
(\hyperref[IV.subsec7.1]{7.1}) how $X_{\sob}$ can in fact be
reconstructed\pageoriginale in terms of $\Top(X)$, by interpreting its
points as ``points'' of the topos $\Top(X)$ (or, equivalently, as fiber
functors).

\subsubsection{}\etiquette[4.2.2]{IV.subsubsec4.2.2}
On every topological space one should introduce the order relation
$\leq$ for which
$$
x\leq y\Longleftrightarrow \bar{\{x\}}\subset \bar{\{y\}}\mbox{ \ie
} x\in \bar{\{y\}}
$$
(one also expresses this by saying that $x$ is a
\emph{specialization} of $y$, or that $y$ is a \emph{generization} of
$x$). For a space of the form $X_{\sob}$, this is simply the inclusion
relation among the irreducible closed subsets of $X$.

With this understood, on the set of maps from a space $X$ to another
space $Y$ one should introduce the order relation called
``specialization'', deduced from that of $Y$, namely
$$
f\leq g\Longleftrightarrow f(x)\leq g(x)\mbox{ for every } x\in
X{\quoi}.
$$
With these conventions, one has the following result:

\subsubsection{}\etiquette[4.2.3]{IV.subsubsec4.2.3}
\emph{Let $X$, $Y$ be two topological spaces, with $Y$ sober, and let
$f$ and $g$ be two continuous maps from $X$ to $Y$. Then there is at
most one morphism from $\Top(f)$ to $\Top(g)$, and such a morphism
exists if and only if $g$ is a specialization of $f$. Finally, every
morphism of topoi $\Top(X)\rightarrow \Top(Y)$ is isomorphic to a
morphism of the form $\Top(f)$, where $f:X\rightarrow Y$ is a
continuous map (uniquely determined by the first assertion).}

If\pageoriginale, for an ordered set $I$, one defines the associated
category $\cat(I)$ by declaring that, for $i\geq j$, there is exactly
one arrow from $i$ to $j$, the preceding result may be summarized by
saying that there is a canonical equivalence of categories
$$
\cat(\Hom_{(\esp)}(X,Y))\xrightarrow{\approx}\Hhomtop(\Top(X),\Top(Y))
$$
(where the second member is defined in \hyperref[IV.subsec3.2]{3.2}).

The following is a formal consequence of these results:

\begin{enonce*}{Corollary 4.2.4}\etiquette[4.2.4]{IV.cor4.2.4}
\begin{enumerate}
\renewcommand{\theenumi}{\alph{enumi}}
\renewcommand{\labelenumi}{\theenumi)}
\item Let $f:X\rightarrow Y$ be a continuous map. Then
$\Top(f):\Top(X)\rightarrow \Top(Y)$ is an equivalence of topoi if
and only if $f_{\sob}:X_{\sob}\rightarrow Y_{\sob}$ is a
homeomorphism (and hence, when $X$ and $Y$ are sober, if and only if
$f$ is a homeomorphism).

\item Let $X$ and $Y$ be topological spaces. The topoi $\Top(X)$ and
$\Top(Y)$ are equivalent if and only if $X_{\sob}$ and $Y_{\sob}$ are
homeomorphic (and hence, if $X$ and $Y$ are sober, if and only if $X$
and $Y$ are homeomorphic).
\end{enumerate}
\end{enonce*}

\subsection{Morphisms into the terminal topos: constant objects of a topos,
section functors}\etiquette[4.3]{IV.subsec4.3}

Let $P$ (the initial letter of ``point'') denote the standard terminal
topos, \ie $P=(\Ens)$ (\hyperref[IV.subsec2.2]{2.2}). Let $E$ be any
topos. We shall see that \emph{up to unique isomorphism, there exists a
unique morphism of topoi}
\hypertarget{sga4.exp.IV.equation.4.3.1}{}$$
f:E\rightarrow P
$$
or, more precisely, that \emph{the category $\Hhomtop(E,P)$ is
equivalent to the one-point category}: between any two objects of this
category there is therefore a\pageoriginale unique arrow, and it is an
isomorphism. (This justifies to some extent the name ``terminal
topos''.) For this, recall \hyperlink{sga4.exp.IV.subsubsection.3.2.1}{(3.2.1)} that $\Hhomtop(E,P)$ is equivalent
to the full subcategory of $\Hhom(P,E)^{\circ}$ formed by the functors
$$
f^{\ast}:P=(\Ens)\longrightarrow E
$$
that commute with $\varinjlim$ and are left exact. Let $e$ be a
singleton set. Every set $X$ can then be written canonically as the
``coproduct of $X$ copies of $e$'', from which it follows that the
category of functors $g:(\Ens)\rightarrow E$ that commute with
$\varinjlim$ is equivalent to the category $E$, by associating with
each $g$ the object $g(e)$ of $E$. In terms of $T=g(e)$, the functor
$g$ is reconstructed, up to unique isomorphism, by
$g(I)=T\times I$, where one sets
$T\times I=\amalg_{i\in I}T_i$, with $T_i=T$ for every $i\in I$. The
fact that the functor of $I$ so defined by $T$ does indeed commute with
$\varinjlim$ follows from the fact that it is manifestly
\sisi{right}{left} adjoint to the functor
$X\mapsto \Hom(T,X)$ from $E$ to $(\Ens)$. This said, for $g$ to be
left exact, it is plainly necessary that $g(e)=T$ be a terminal object
of $E$ (since $e$ is a terminal object of $(\Ens)$), and it follows
easily from the fact that in $E$ ``coproducts are universal''
(\hyperref[IV.subsubsec1.1.2]{1.1.2} b)) that this condition is also
sufficient. Thus the category of inverse-image functors $f^{\ast}$ is
equivalent to the category of terminal objects of $E$, which is itself
plainly equivalent to the terminal category.

It follows from the foregoing that the choice of a morphism \hyperlink{sga4.exp.IV.equation.4.3.1}{(4.3.1)}
is essentially equivalent to the choice of a terminal object of $E$,
say $e_E$. In terms of this object, one then has canonical
isomorphisms of functors
\begin{equation*}
 f^{\ast}(I)\simeq e_E\times
 I=\mbox{ coproduct of } I \mbox{ copies of } e_E \quad \mbox{ for } I\in \ob
 (\Ens),\tag{4.3.2}
\end{equation*}\etiquette[4.3.2]{IV.eq4.3.2}
and\pageoriginale
\begin{equation*}
f_{\ast}(X)=\Hom(e_E,X) \quad \quad \mbox{ for } X\in \Ob E.\tag{4.3.3}
\end{equation*}\etiquette[4.3.3]{IV.eq4.3.3}


\setcounter{subsubsection}{3}
\subsubsection{}\etiquette[4.3.4]{IV.subsubsec4.3.4}
The preceding two functors will play an important role below. For a
set $I$, the object $f^{\ast}(I)=e_E\times I$ of
\eqhyperref[IV.eq4.3.2]{4.3.2} is called the \emph{constant object
with value} $I$ in $E$ (or, when $E$ is presented as a category
$C^{\widetilde{\phantom{,}}}$ in terms of a site $C$, the
\emph{constant sheaf with value $I$ on $C$}). It will also often be
denoted by $I_E$, or by $I_C$ when $E$ is defined by the site $C$. The
fact that $I\mapsto I_E$ is the inverse-image functor of a morphism of
topoi specifies its exactness properties, which imply in particular
that this functor respects all the usual kinds of algebraic structure,
transforming a group into a group object of $E$, and so on
(\hyperref[IV.subsubsec3.1.2]{3.1.2}). If, for example, one has a
group $G$ of $E$, one says that it is a \emph{constant group} (or,
where appropriate, a \emph{constant sheaf of groups}) if it is
isomorphic to a group of the form $G_{\circ E}$, where $G_{\circ}$ is
an ordinary group. The same terminology is used for every other kind
of ``algebraic'' structure, in the sense specified (more or less) in
\hyperref[IV.subsubsec3.1.2]{3.1.2}.

% END INCLUDED SGA4_English_current_buildable_source_20260906/source/04/drafts/SGA4_IV_LINES_1320_1659_English.texfrag


% BEGIN INCLUDED SGA4_English_current_buildable_source_20260906/source/04/drafts/SGA4_IV_LINES_1660_1806_English.texfrag
\subsubsection{}\etiquette[4.3.5]{IV.subsubsec4.3.5}
Notice that the functor $I\mapsto I_E$ is not necessarily fully
faithful (or even faithful: take $E$ to be the ``empty topos''
(\hyperref[IV.subsec2.2]{2.2})); thus a constant object of $E$ does
not in general determine, up to unique isomorphism, the set $I$ from
which it arises. We say that $E$ is $0$-\emph{acyclic}, or
\emph{connected and nonempty}, if the functor $I\mapsto I_E$ is fully
faithful. By the general properties of adjoint functors, this amounts
to saying that the adjunction morphism
$$
I\longrightarrow f_{\ast}(f^{\ast}(I))=f_{\ast}(I_E)=\Hom(e_E, I_E)
$$
is\pageoriginale an isomorphism (so that the functor
\eqhyperref[IV.eq4.3.3]{4.3.3} recovers the ``value'' of a constant
object of $E$). It is easily checked that this holds if and only if
$e_E$ is not the initial object $\phi_E$ of $E$, \ie if $E$ is not
an ``empty topos'' (which expresses the \emph{faithfulness} of the
functor $I\mapsto I_E$\footnote[1]{or equivalently the fact that this
functor is conservative (\hyperrefext[I][I.cor6.3]{6.3}).}), \emph{and}
if $e_E$ is a \emph{connected} object of $E$, \ie is not the sum of
two objects of $E$ which are not ``empty'' (\ie which are not initial
objects of $E$).

\subsubsection{}\etiquette[4.3.6]{IV.subsubsec4.3.6}
The functor \eqhyperref[IV.eq4.3.3]{4.3.3} is also often called the
\emph{sections functor} and is denoted by $\Gamma_E$, $\Gamma(E,-)$,
or simply $\Gamma$:
\hypertarget{sga4.exp.IV.equation.4.3.6.1}{}\begin{equation*}
\Hom(e_E,X)=\Gamma_E(X)=\Gamma(E,X)=\Gamma(X){\quoi}.\tag{4.3.6.1}
\end{equation*}
It commutes with arbitrary projective limits, but is not in general
right exact; its derived functors (on abelian group objects) will be
studied in the next exposé.

\subsection{Morphisms from the ``empty topos''}\etiquette[4.4]{IV.subsec4.4}

Let $\phi_{\text{top}}$ be an empty topos, corresponding to a category
of sheaves $\Phi$ equivalent to the terminal category
(\hyperref[IV.subsec2.2]{2.2}). Let $E$ be a topos. The category of
functors from $E$ to $\Phi$ is clearly equivalent to the one-point
category, and every such functor commutes with inductive and
projective limits (without any merit in this), and hence may be
regarded as an inverse-image functor for a morphism of topoi
$\phi_{\text{top}}\rightarrow E$. It follows that the \emph{category
$\Hhomtop(\phi_{\text{top}},E)$ is equivalent to the one-point
category}, and in particular that \emph{there exists, up to
unique\pageoriginale isomorphism, exactly one morphism of topoi}
\begin{equation*}
\phi_{\text{top}}\longrightarrow E{\quoi}.\tag{4.4.1}
\end{equation*}\etiquette[4.4.1]{IV.eq4.4.1}
This justifies to some extent the terminology ``initial topos''
introduced in \hyperref[IV.subsec2.2]{2.2}.

One may also determine the morphisms of topoi
\begin{equation*}
E\longrightarrow \phi_{\text{top}}{\quoi}\ ;\tag{4.4.2}
\end{equation*}\etiquette[4.4.2]{IV.eq4.4.2}
one checks at once that such a morphism exists if and only if the
initial object of $E$ is also a terminal object, \ie if and only if
$E$ itself is an ``empty topos,'' and that in this case the category
$\Hhomtop(E,\phi_{\text{top}})$ is again equivalent to the one-point
category. The unique morphism \eqhyperref[IV.eq4.4.2]{4.4.2} (up to
isomorphism) is then an equivalence of topoi.

\subsection{The classifying topos $B_G$ for a varying group $G$}\etiquette[4.5]{IV.subsec4.5}

\subsubsection{}\etiquette[4.5.1]{IV.subsubsec4.5.1}
Let $E$ be a topos, and let
$$
f:G\longrightarrow H
$$
be a morphism of group objects in $E$. It induces a ``restriction of
the group of operators'' functor
$$
f^{\ast}:B_H\longrightarrow B_G{\quoi},
$$
where the notation is that of \hyperref[IV.subsec2.4]{2.4}. It is
immediate that this functor commutes with inductive and projective
limits; a fortiori, it may be interpreted as an inverse-image
functor associated with a morphism of topoi
$$
B_f\mbox{ or } f:B_G\longrightarrow B_H{\quoi}.
$$
The\pageoriginale corresponding direct-image functor
$$
f_{\ast}:B_G\longrightarrow B_H
$$
is readily made explicit by the formula
$$
f_{\ast}(X)=\Hhom_G(H_s,X){\quoi},
$$
where $X$ is an object of \sisi{$X$}{$E$} with a left action of $G$,
where $H_s$ is $H$ regarded as endowed with the left actions by $G$
induced by \sisi{$F$}{$f$}, and where $\Hhom_G$ denotes the evident
subobject of the object $\Hhom$ defined below
(\hyperref[IV.subsec10.2]{10.2}); the left action of $H$ on
$\Hhom_G(H_s,X)$ is induced by the right actions of $H$ on $H_s$ by
\emph{right} translations.

Since the inverse-image functor $f^{\ast}$ commutes with arbitrary
\sisi{$\varinjlim$}{$\varprojlim$} (and not only with finite
\sisi{$\varinjlim$}{$\varprojlim$}), it is itself the right adjoint of
a functor
$$
f_{!}:B_G\longrightarrow B_H{\quoi},
$$
so that there is a sequence of three adjoint functors as in
\hyperref[IV.subsubsec3.1.3]{3.1.3}:
$$
f_{!},f^{\ast}, f_{\ast}{\quoi}.
$$
The functor $f_{!}$ is readily made explicit by the formula
$$
f_{!}(X)=H\times {}^GX{\quoi},
$$
where the right-hand side denotes the ``contracted product'' induced
by the actions of $G$ on $X$ (on the left) and on $H$ (on the right,
via right translations and $f$), defined as the quotient of
$H\times X$ by the action of $G$ given by
$$
g\circ (h,x)=(hg^{-1}, gx){\quoi}.
$$

The functor $f_{!}$, being a left adjoint, clearly commutes with
inductive limits, but is not in general left exact (\ie
it\pageoriginale cannot in turn be regarded as an inverse-image
functor for a morphism of topoi $B_H\rightarrow B_G$). In fact, one
checks easily that it can be left exact only if $f:G\rightarrow H$
is an isomorphism. Likewise, the functor $f_{\ast}$, which commutes
with projective limits, is not in general right exact and, a
fortiori, does not in general admit a right adjoint. At least when
$E$ is the punctual topos, \ie when $G$ and $H$ are ordinary groups,
$f_{\ast}$ is right exact only if $f$ is an isomorphism.

Given a second morphism of groups $g:H\rightarrow K$, one obtains,
as in \hyperref[IV.subsubsec4.1.1]{4.1.1}, transitivity up to
canonical isomorphism, so that, subject to the usual reservation, the
classifying topos $B_G$ may be regarded as depending
\emph{functorially} on the group $G$. We leave it to the reader to
generalize this functorial behavior to the case where both $G$
\emph{and} the topos $E$ vary.


% END INCLUDED SGA4_English_current_buildable_source_20260906/source/04/drafts/SGA4_IV_LINES_1660_1806_English.texfrag


% BEGIN INCLUDED SGA4_English_current_buildable_source_20260906/source/04/drafts/SGA4_IV_LINES_1807_1981_English.texfrag
\subsubsection{The topos $B_{𝒢}$ for a varying pro-group
  $𝒢$}\etiquette[4.5.2]{IV.subsubsec4.5.2}

We leave it to the reader to make precise the covariant character of
the topos $B_{𝒢}$ (\hyperref[IV.subsec2.7]{2.7}) with respect to $𝒢$,
by following the account given in \hyperref[IV.subsubsec4.5.1]{4.5.1}.
One should note, however, that in the case of a morphism
$f:𝒢\rightarrow ℋ$ of pro-groups which are not essentially
constant, the corresponding morphism of topoi
$B_{𝒢}\rightarrow B_{ℋ}$ does not in general permit the definition
of a functor $f_{!}$ (whose right adjoint would be the functor
$f^{\ast}$ restricting the pro-group of operators). Suppose, to
simplify the statement, that $𝒢$ and $ℋ$ are defined by
\emph{profinite} groups $G$ and $H$. It is then easily seen that
$f_{!}$ exists if and only if the image of the morphism in question
$f:G\rightarrow H$ has finite index in $H$, and in\pageoriginale
that case $f_{!}$ is given by the same formula as in
\hyperref[IV.subsec4.5]{4.5}.

\subsection{The topos $\widehat{C}$ for a varying category $C$}\etiquette[4.6]{IV.subsec4.6}

\subsubsection{}\etiquette[4.6.1]{IV.subsubsec4.6.1}
Let
$$
f:C\longrightarrow C'
$$
be a functor from a category $C\in 𝒰$ to another category $C'$, and
hence let
$$
f^{\ast}:\widehat{C}'\longrightarrow \widehat{C}\quav
f^{\ast}(F)=F\circ f{\quoi}.
$$
It is immediate that this functor commutes with projective and
inductive limits; a fortiori, it may be regarded as the inverse-image
functor for a morphism of topoi
$$
\widehat{f}\mbox{ or } f:\widehat{C}\rightarrow \widehat{C}'{\quoi}.
$$
The corresponding direct-image functor
$$
f_{\ast}:\widehat{C}\longrightarrow \widehat{C}'
$$
is none other than the functor likewise denoted $f_{\ast}$ in I \hyperref[I.prop5.1]{5.1}.
Moreover (as was to be expected from the fact that $f^{\ast}$ also
commutes with arbitrary $\varprojlim$), $f^{\ast}$ also admits a left
adjoint
$$
f_{!}:\widehat{C}\longrightarrow \widehat{C}'{\quoi},
$$
(which was likewise denoted $f_{!}$ in
\hyperrefext[I][I.prop5.1]{5.1}). Thus one obtains a sequence of three
adjoint functors
$$
f_{!}, f^{\ast}, f_{\ast}{\quoi},
$$
the first of which is moreover an extension of $f:C\rightarrow C'$
to the categories of presheaves (for the usual embeddings of $C$ and
$C'$ in the categories $\widehat{C}$ and $\widehat{C}'$). One should
note that $f_{!}$ is not in general left exact, nor is $f_{\ast}$ in
general right exact; this removes any ambiguity about the direction
of\pageoriginale variance of the topos $\widehat{C}$ associated with
the varying category $C$: this topos is a \emph{covariant} functor of
$C$, subject to the usual reservation arising from the transitivity
isomorphisms (\cf \hyperref[IV.subsubsec4.1.1]{4.1.1}). When the sets
of objects of $C$ and $C'$ each consist of a single element, so that
$C$ and $C'$ identify with monoids $G$ and $G'$, the corresponding
topoi are the classifying topoi $B_{G^{\circ}}$ and
$B_{{G'}^{\circ}}$, and one recovers their variance with respect to
the varying $G$, already encountered from other viewpoints in
\hyperref[IV.subsubsec4.1.2]{4.1.2} and
\hyperref[IV.subsec4.5]{4.5} (where the restriction to groups rather
than monoids was inessential).

\subsubsection{}\etiquette[4.6.2]{IV.subsubsec4.6.2}
The 2-functorial dependence of the topos $\widehat{C}$ on $C$ can be
made more precise by introducing, for two categories $C,C'\in 𝒰$, a
canonical functor
\begin{equation*}
\Hhom(C,C')\longrightarrow
\Hhomtop(\widehat{C},\widehat{C}'){\quoi}.\tag{4.6.2.1}
\end{equation*}\etiquette[4.6.2.1]{IV.eq4.6.2.1}
It remains to define this functor on arrows. For this, note that
$f\mapsto f^{\ast}$ identifies (up to equivalence of categories) the
right-hand side with a full subcategory of
$\Hhom(\widehat{C}',\widehat{C})^{\circ}$ \hyperlink{sga4.exp.IV.subsubsection.3.2.1}{(3.2.1)}. If
$f,g:C\rightarrow C'$ are two functors, every morphism of functors
$u:f\rightarrow g$ defines, for $F\in\ob\widehat{C}'$, a morphism of
functors $F\circ g\rightarrow F\circ f$ (since $F$ is a
\emph{contra}variant functor), hence a morphism
$g^{\ast}\rightarrow f^{\ast}$ and therefore a morphism
$\widehat{f}\rightarrow\widehat{g}$, as asserted.

When $C$ is the one-point category, $\widehat{C}$ is the punctual
topos (\hyperref[IV.subsec2.2]{2.2}), denoted $P$, and
\eqhyperref[IV.eq4.6.2.1]{4.6.2.1} is interpreted as a natural functor
\begin{equation*}
C'\longrightarrow \Hhomtop(P,\widehat{C}') \mathop{=}^{\dfn}
\Ppoints(\widehat{C}')\tag{4.6.2.2}
\end{equation*}\etiquette[4.6.2.2]{IV.eq4.6.2.2}
from $C'$ to the ``\emph{category of points}'' of $\widehat{C}'$,
which will be studied in 6. This\pageoriginale functor is not
necessarily an equivalence of categories \hyperlink{sga4.exp.IV.subsubsection.7.3.3}{(7.3.3)}; a fortiori,
\eqhyperref[IV.eq4.6.2.1]{4.6.2.1} is not necessarily an equivalence
of categories.

On the other hand, the functor \eqhyperref[IV.eq4.6.2.1]{4.6.2.1} is
always \emph{fully faithful}. To see this, note that if
$f,g:E\rightarrow E'$ are two morphisms of topoi for which $f_{!}$
and $g_{!}$ are defined (\hyperref[IV.subsubsec3.1.3]{3.1.3}), then
the theory of adjoint functors gives canonical isomorphisms
$$
\Hom(f_{!},g_{!})\simeq \Hom(g^{\ast}, f^{\ast})\simeq
\Hom(f_{\ast}, g_{\ast}) \mathop{=}^{\dfn} \Hom(f,g){\quoi}.
$$
Applying this to functors of the form $\widehat{f}$ and $\widehat{g}$
associated with $f,g:C\rightarrow C'$, one obtains the asserted
result, taking into account that the natural extension map
$\Hom(f,g)\rightarrow \Hom(f_{!},g_{!})$ is bijective (I \hyperref[I.prop7.8]{7.8}).

\subsubsection{}\etiquette[4.6.3]{IV.subsubsec4.6.3}
One may ask when the functor \eqhyperref[IV.eq4.6.2.1]{4.6.2.1} is an
equivalence of categories, \ie when it is essentially surjective.
This is a special case of the problem of determining all morphisms of
topoi $\widehat{C}\rightarrow\widehat{C}'$. More generally, if $C$ is
a category in $𝒰$ and $E$ is a topos, one may seek to determine the
morphisms of topoi
$$
f:\widehat{C}\longrightarrow E{\quoi}.
$$
The category of these morphisms is equivalent to the opposite of the
category of functors $f^{\ast}:E\rightarrow\widehat{C}$ commuting
with arbitrary inductive limits and finite projective limits. A
functor $E\rightarrow\widehat{C}$ may be interpreted as a functor
$E\times C^{\circ}\rightarrow(𝒰\text{-}\Ens)=(\Ens)$, or equivalently
as a functor $F:C^{\circ}\rightarrow\Hhom(E,(\Ens))$. The exactness
property in question is then expressed by requiring that, for every
object $X$ of $C$, the functor $F(X):E\rightarrow(\Ens)$ commute with
arbitrary inductive limits and finite projective limits (or, as we
shall say in 6, that $F(X)$ is a ``\emph{fiber\pageoriginale
functor}'' on $E$). This amounts to saying that $F(X)$ is the
inverse-image functor for a morphism of topoi $P\rightarrow E$, and
one therefore obtains finally a canonical equivalence of categories
\begin{equation*}
\Hhomtop(\widehat{C}, E) \xrightarrow{\approx} \Hhom(C,
\Ppoints(E)){\quoi},\tag{4.6.3.1}
\end{equation*}\etiquette[4.6.3.1]{IV.eq4.6.3.1}
where, for brevity, we write
$$
\Ppoints(E)=\Hhomtop(P,E)
$$
for the category of points of the topos $E$.

When $E$ is of the form $\widehat{C}'$, it follows immediately from
the definitions that the composite of
\eqhyperref[IV.eq4.6.2.1]{4.6.2.1} and the preceding equivalence
\eqhyperref[IV.eq4.6.3.1]{4.6.3.1} is the functor
\begin{equation*}
\Hhom(C,C')\longrightarrow
\Hhom(C,\Ppoints(\widehat{C}'))\tag{4.6.3.2}
\end{equation*}\etiquette[4.6.3.2]{IV.eq4.6.3.2}
defined by $F\mapsto i\circ F$, where
$i:C'\rightarrow\Ppoints(\widehat{C}')$ is the canonical embedding
\eqhyperref[IV.eq4.6.2.2]{4.6.2.2}. It follows at once that
\emph{for} \eqhyperref[IV.eq4.6.2.1]{4.6.2.1} \emph{to be an
equivalence, it is necessary and sufficient that $C$ be empty or
that} \eqhyperref[IV.eq4.6.2.2]{4.6.2.2} \emph{be essentially
surjective} (and hence an equivalence of categories). This last
condition on $C'$ is satisfied in certain interesting cases, notably
when $C'$ is the one-object category defined by a group $G$
(\hyperref[IV.subsubsec7.2.5]{7.2.5}).



% END INCLUDED SGA4_English_current_buildable_source_20260906/source/04/drafts/SGA4_IV_LINES_1807_1981_English.texfrag


% BEGIN INCLUDED SGA4_English_current_buildable_source_20260906/source/04/drafts/SGA4_IV_LINES_1982_2188_English.texfrag
\begin{enonce*}[remark]{Remark 4.6.4}\etiquette[4.6.4]{IV.rem4.6.4}\nde{For
a recent perspective on the notion of the homotopy type of a topos, see
[15] $\S$ \hyperref[IV.subsubsec7.1.6]{7.1.6}.}
The association of a topos $\widehat{C}$ with an arbitrary category
$C$ suggests that a category $C$ has invariants of a topological
nature (cohomology groups, homotopy groups, \etc), just as a topos
does. The cohomology groups of $\widehat{C}$ with coefficients in an
abelian group object $F$ (in the general sense studied in V) are
precisely the values of the derived functors
$\varprojlim^{(n)}$ of the functor $\varprojlim$, already familiar
to topologists.\pageoriginale J. L. Verdier and (independently)
D. G. Quillen verified that, when one restricts to constant, or more
generally locally constant, coefficients, these cohomology groups
identify with the cohomology groups of the semi-simplicial set
$\Nerf(C)$ canonically associated with $C$ [5, ${\rm n^{\circ}}$ 212,
prop. 4.1], and moreover that, up to isomorphism in the ``homotopy
category'' of [2], every semi-simplicial set can be obtained from a
category $C$. With a suitable notion of homotopy type for topoi,
which we do not specify here, one may say that the semi-simplicial
homotopy types of topologists are precisely the homotopy types of
topoi of the special form $\widehat{C}$, and more generally of topoi
$E$ in which every object admits a covering that refines every other
one (\hyperref[IV.subsec2.6]{2.6})
\footnote[1]{and, more generally still, of topoi that are ``locally
$\infty$-connected'' in an evident sense which we leave it to the
reader to make precise.}. In the absence of this condition on $E$,
its homotopy type can at best be expressed by a suitable
\emph{projective system} of semi-simplicial sets [1].
\end{enonce*}

\subsection{The topos $C^{\sim}$ for a variable site $C$
(cocontinuous functors)}\etiquette[4.7]{IV.subsec4.7}

Let
$$
f:C\longrightarrow C'
$$
be a cocontinuous functor (\hyperrefext[III][III.def2.1]{2.1}) between
sites in $𝒰$, \ie a functor such that the functor
$\widehat{f}_{\ast}$ from $\widehat{C}$ to $\widehat{C}'$
(\hyperref[IV.subsubsec4.6.1]{4.6.1}) maps $C^{\sim}$ into
$C'^{\sim}$, that is, induces a functor
\begin{equation*}
\widetilde{f}_{\ast}:\widetilde{C}\longrightarrow
{C'}^{\widetilde{\phantom{,}}}\tag{4.7.1}
\end{equation*}\etiquette[4.7.1]{IV.eq4.7.1}
making the following diagram of functors commutative:
\[\vcenter{
\xymatrix@C=2.5cm@R=1.25cm{ C^{\widetilde{\phantom{a}}}
\ar[r]^{\widetilde{f}_{\ast}} \ar[d]_{\sisi{i}{i_{\ast}}}&
C'^{\widetilde{\phantom{a}}} \ar[d]^{\sisi{i_{\ast}}{i_{\ast}'}}\\
\widehat{C}  \ar[r]^{\widehat{f}_{\ast}}       & \widehat{C}'\quad,
}}\leqno{(4.7.2)}
\]\etiquette[4.7.2]{IV.eq4.7.2}
where\pageoriginale \sisi{$i$}{$i_{\ast}$} and
\sisi{$i'$}{$i_{\ast}'$} are the inclusion functors. We saw
(\hyperrefext[III][III.prop2.3]{2.3}) that the functor
$\widetilde{f}_{\ast}$ then has a left adjoint
$\widetilde{f}^{\ast}$, and that the latter is left exact. In other
words, $\widetilde{f}_{\ast}$ is the direct-image functor associated
with a morphism of topoi
$$
\widetilde{f}\mbox{ or } f:C^{\sim}\longrightarrow
{C'}^{\widetilde{\phantom{,}}}{\quoi},
$$
the corresponding inverse-image functor being, of course,
$\widetilde{f}^{\ast}$. Taking left adjoints of the functors involved,
the commutative diagram \eqhyperref[IV.eq4.7.2]{4.7.2} also gives a
diagram commutative up to isomorphism:
\[\vcenter{
\xymatrix@C=2.5cm@R=1.5cm{
C^{\widetilde{\phantom{a}}} &  C'^{\widetilde{\phantom{a}}} \ar[l]_{\widetilde{f}^{\ast}}\\
\widehat{C}  \ar[u]^{\underline{a}}  & \widehat{C}'
\ar[l]_{\widehat{f}^{\ast}}\ar[u]_{\underline{a}'}\quad,
}}\leqno{(4.7.3)}
\]\etiquette[4.7.3]{IV.eq4.7.3}
where $\underline{a}$ and $\underline{a}'$ are the ``associated
sheaf'' functors. This diagram immediately recovers the formula
(\hyperrefext[III][III.prop2.3]{2.3})
$$
\widetilde{f}^{\ast}=\underline{a}\widehat{f}^{\ast} i'{\quoi}.
$$

The transitivity property for the morphisms of topoi
$\widehat{f}:\widehat{C}\rightarrow \widehat{C}'$ implies the same
property for the morphisms of topoi
$\widetilde{f}:C^{\widetilde{\phantom{,}}}\rightarrow
{C'}^{\widetilde{\phantom{,}}}$ associated with cocontinuous
functors. Thus one may say that \emph{the topos
$C^{\widetilde{\phantom{,}}}$ varies functorially with $C$} in a
\emph{covariant} manner, when cocontinuous functors are taken as the
``morphisms'' of sites.

In\pageoriginale all the cases encountered thus far, the
cocontinuous functor $f$ used is also continuous, that is
(\hyperrefext[III][III.def1.1]{1.1}), it ``extends'' to a functor
$$
\widetilde{f}_{!}:C^{\widetilde{\phantom{,}}}\longrightarrow
{C'}^{\widetilde{\phantom{,}}}
$$
commuting with inductive limits, which is left adjoint to
$\widetilde{f}^{\ast}$, so that there is a sequence of three adjoint
functors
$$
\widetilde{f}_{!}, \widetilde{f}^{\ast},
\widetilde{f}_{\ast}{\quoi}.
$$
One should note that the functor $\widetilde{f}_{!}$ is not in
general left exact, nor is $\widetilde{f}_{\ast}$ right exact. This
removes any ambiguity concerning the direction of the variance of
the topos $C^{\widetilde{\phantom{,}}}$, when $C$ varies through
functors assumed only to be cocontinuous, or even continuous and
cocontinuous.

\begin{enonce*}[remark]{Remark 4.7.4}\etiquette[4.7.4]{IV.rem4.7.4}
Given a morphism of topoi
$$
F:E\longrightarrow E'{\quoi},
$$
for there to exist a functor $F_{!}$ left adjoint to $F^{\ast}$, it
is necessary and sufficient that one be able to ``realize'' (up to
equivalence) $E$ and $E'$ in the form
$C^{\widetilde{\phantom{,}}}$ and ${C'}^{\widetilde{\phantom{,}}}$,
for two sites $C,C'\in 𝒰$, and that one be able to find a
\emph{continuous and cocontinuous} functor $f:C\rightarrow C'$ such
that $F$ identifies with $\widetilde{f}$. This is clearly sufficient;
for necessity, it is enough to take for $C$ and $C'$ small full
generating subcategories of $E$ and $E'$, respectively, such that
$$
F_{!}(\ob C)\subset \ob C'{\quoi},
$$
endowed with the topologies induced by those of $E$ and $E'$, and to
take for $f$ the functor induced by $F_{!}$ (\cf
\hyperref[IV.cor1.2.1]{1.2.1}). Recall
(\hyperref[IV.rem1.8]{1.8}) that the existence of $F_{!}$ also means
that $F^{\ast}$ commutes with projective limits, or,\pageoriginale
equivalently here since this functor is left exact, that it commutes
with products.
\end{enonce*}

\setcounter{subsubsection}{4}
\subsubsection{}\etiquette[4.7.5]{IV.subsubsec4.7.5}
If one asks which morphisms of topoi $F:E\rightarrow E'$ can be
realized by a cocontinuous (not necessarily continuous) functor of
sites, one likewise finds the following answer: the set of objects
$X$ of $E$ such that the functor
$X'\mapsto \Hom(X,F^{\ast}(X'))$ from $E'$ to $(\Ens)$ commutes with
projective limits (or, equivalently, with products) must be a
\emph{generating family of $E$.}\nde{This statement does not appear
to be correct: the condition expressed is equivalent to the assertion
that $F^{\ast}$ commutes with arbitrary products, or equivalently that
$F^{\ast}$ has a left adjoint $F_{!}$. For every site $C$, the
identity functor $F=id:C\longrightarrow C$ is cocontinuous when the
term on the right is endowed with the trivial topology. The functor
$F^{\ast}:\hat{C}\longrightarrow\widetilde{C}$ is then the associated
sheaf functor, and does not in general commute with arbitrary
products.}

\subsection{The morphism of topoi $\widetilde{C}\rightarrow
\widehat{C}$ for a site $C$}\etiquette[4.8]{IV.subsec4.8}

Let $C$ be a small site, with which the two topoi
$C^{\widetilde{\phantom{,}}}$ and $\widehat{C}$ are therefore
associated (the latter does not depend on the topology placed on
$C$). In II \hyperref[II.theo3.4]{3.4} we defined the ``associated sheaf'' functor
$$
\underline{a}:\widehat{C}\longrightarrow
C^{\widetilde{\phantom{,}}}{\quoi},
$$
and proved that it is left exact and commutes with inductive limits.
It is therefore the inverse-image functor associated with a morphism
of topoi
$$
p:\widetilde{C}\longrightarrow \widehat{C},\;
p^{\ast}=\underline{a}{\quoi},
$$
the corresponding direct-image functor being the canonical inclusion
$$
i=p_{\ast}:C^{\widetilde{\phantom{,}}}\longrightarrow
\widehat{C}{\quoi}.
$$
It is well known that the latter functor is not in general right
exact (its derived functors on abelian objects give rise to the
cohomology presheaves $\Xx^n(F)$ of V 2), and that $\underline{a}$
does not in general commute with arbitrary $\varprojlim$. This
removes any ambiguity concerning the direction of the ``natural''
morphism of topoi between $\widehat{C}$ and
$C^{\widetilde{\phantom{,}}}$.

Given a cocontinuous functor
$$
f:C\longrightarrow C'
$$
of\pageoriginale sites, one obtains a diagram of morphisms of topoi
\[
\xymatrix@C=2.25cm@R=1.25cm{ C^{\widetilde{\phantom{a}}} \ar[d]_{\sisi{p}{p_{\ast}}}
\ar[r]^{\widetilde{f}} &
C'^{\widetilde{\phantom{a}}} \ar[d]^{\sisi{p'}{p'_{\ast}}}\\
\widehat{C} \ar[r]^{\widehat{f}}&  \widehat{C}' }
\]
which is commutative up to canonical isomorphism: this is precisely
what the commutativity of the diagram of functors
\eqhyperref[IV.eq4.7.2]{4.7.2} expresses. Thus one may say that the
morphism of topoi
$p:C^{\widetilde{\phantom{,}}}\rightarrow \widehat{C}$ is functorial
in $C$, when $C$ is varied through \emph{cocontinuous} functors
between sites.


% END INCLUDED SGA4_English_current_buildable_source_20260906/source/04/drafts/SGA4_IV_LINES_1982_2188_English.texfrag


% BEGIN INCLUDED SGA4_English_current_buildable_source_20260906/source/04/drafts/SGA4_IV_LINES_2189_2354_English.texfrag
\subsection{Effect of a continuous functor of sites. Morphisms of sites}\etiquette[4.9]{IV.subsec4.9}

\subsubsection{}\etiquette[4.9.1]{IV.subsubsec4.9.1}
Let
$$
f:C\longrightarrow C'
$$
be a continuous functor from $C$ to $C'$, \ie such that the functor
$\widehat{f}^{\ast}:\widehat{C}'\rightarrow\widehat{C}$ maps
${C'}^{\widetilde{\phantom{,}}}$ into
$C^{\widetilde{\phantom{,}}}$, and hence induces a functor
$$
f_s:{C'}^{\widetilde{\phantom{,}}}\longrightarrow
C^{\widetilde{\phantom{,}}}{\quoi}.
$$
We saw (\hyperrefext[III][III.prop1.2]{1.2}) that the latter admits a
left adjoint
$$
f^s:C^{\widetilde{\phantom{,}}}\longrightarrow
{C'}^{\widetilde{\phantom{,}}}{\quoi},
$$
which moreover ``extends'' $f$ in the evident sense. One should note
that in general $f^s$ is not \sisi{right}{left} exact (even if $f$ is
also cocontinuous), nor does $f_s$ commute with inductive limits, so
that the datum of $f$ does not, without further hypotheses, describe
a morphism of topoi in either direction between
$C^{\widetilde{\phantom{,}}}$ and
${C'}^{\widetilde{\phantom{,}}}$. The case where $f$ is cocontinuous,
\ie where $f_s$ commutes with inductive limits and may\pageoriginale
therefore be regarded as an inverse-image functor for a morphism of
topoi $\widetilde{f}:C^{\widetilde{\phantom{,}}}\rightarrow
{C'}^{\widetilde{\phantom{,}}}$, was examined in
\hyperref[IV.subsec4.7]{4.7}. We shall examine the case where the
functor $f^s$ is left exact, and may therefore be regarded as an
inverse-image functor for a morphism of topoi in the \emph{opposite
direction}:
\begin{equation*}
\Top(f)=g:{C'}^{\sim}\longrightarrow C^{\sim}{\quoi}. \tag{4.9.1.1}
\end{equation*}\etiquette[4.9.1.1]{IV.eq4.9.1.1}
Notice that \emph{we have deliberately not denoted this morphism by
the letter $\widetilde{f}$ or $f$, so as to avoid confusion with the
situation of \hyperref[IV.subsec4.7]{4.7}, in accordance with the
general recommendations of \hyperref[IV.subsubsec3.1.3]{3.1.3}.} A
continuous functor $f:C\rightarrow C'$ is sometimes called a
\emph{morphism of sites from $C'$ to $C$} (note: \emph{not from $C$
to $C'$}) if the functor $f^s$ is left exact; in other words, if
there exists a morphism of topoi \eqhyperref[IV.eq4.9.1.1]{4.9.1.1}
such that the corresponding inverse-image functor
$$
g^{\ast}:C^{\widetilde{\phantom{,}}}\longrightarrow
{C'}^{\widetilde{\phantom{,}}}
$$
\emph{extends} the functor $f$, \ie makes the following diagram of
functors commute up to isomorphism:
\[
\xymatrix@C=2cm@R=1.5cm{
C \ar[d]_{\epsilon} \ar[r]^{f}& C' \ar[d]^{{\epsilon}'}\\
C^{\widetilde{\phantom{a}}} \ar[r]^{g^{\ast}} &
C'^{\widetilde{\phantom{a}}} &, }
\]
%% widetilde + \phantom: horrible hack \XXX
where $\epsilon$ and $\epsilon'$ are the canonical functors of II
\hyperref[II.subsubsec4.4.0]{4.4.0}.

\subsubsection{}\etiquette[4.9.2]{IV.subsubsec4.9.2}
In practice, one recognizes that a continuous functor
$f:C\rightarrow C'$ is a morphism of sites from $C'$ to $C$ from the
fact that finite projective limits are representable in $C$ and that
$f$ commutes with them (\hyperrefext[III][III.prop1.3]{1.3} 5)).
Subject\pageoriginale to the indicated condition on $C$ (which is
almost always satisfied in practice), and assuming in addition that
the topology of $C'$ is coarser than the canonical topology (also
almost always satisfied), the preceding sufficient condition ($f$
left exact) for $f$ to be a morphism of sites from $C'$ to $C$ is
also necessary.

\subsubsection{}\etiquette[4.9.3]{IV.subsubsec4.9.3}
In the spirit of the preceding discussion, if $C$ and $C'$ are two
$𝒰$-sites, the category of morphisms of sites
$\Mmorsite(C',C)$ from $C$ to $C'$ should be defined as the full
subcategory of the \emph{opposite category
$\Hhom(C,C')^{\circ}$} of the category of functors from $C$ to $C'$,
formed by those functors which are morphisms of sites (from $C'$ to
$C$). In this way one obtains a canonical functor (defined up to
unique isomorphism)
$$
\Mmorsite(C',C)\longrightarrow
\Hhomtop({C'}^{\widetilde{\phantom{,}}},
C^{\widetilde{\phantom{,}}}){\quoi}.
$$

\begin{enonce*}{Proposition 4.9.4}\etiquette[4.9.4]{IV.prop4.9.4}
Let $E$ be a $𝒰$-topos and $C$ a $𝒰$-site. The functor
$f\mapsto f^{\ast}| C=f^{\ast}\circ\epsilon_C$, associating with
every morphism of topoi
$f:E\rightarrow C^{\widetilde{\phantom{,}}}$ the ``restriction'' to
$C$ of the associated inverse-image functor
$f^{\ast}:C^{\widetilde{\phantom{,}}}\rightarrow E$, induces an
equivalence of categories
$$
\Hhomtop(E,C^{\widetilde{\phantom{,}}})\xrightarrow{\approx}\Mmorsite(E,C)(\hookrightarrow
\Hhom(C,E)^{\circ}){\quoi}.
$$
When finite $\varprojlim$ are representable in $C$, the corresponding
fully faithful functor
$$
\Hhomtop(E,C^{\widetilde{\phantom{,}}})\longrightarrow
\Hhom(C,E)^{\circ}
$$
has as its essential image the functors $g:C\rightarrow E$ which are
left exact and continuous, or equivalently those which are left exact
and carry\pageoriginale covering families to covering families.
\end{enonce*}

The last assertion follows from the first by
\hyperref[IV.subsubsec4.9.2]{4.9.2}. On the other hand, it follows
from \hyperrefext[III][III.prop1.2]{1.2} iv) and IV
\hyperref[IV.theo1.2]{1.2} iii) that the functor
$G\rightarrow G\circ\epsilon$ induces an equivalence between the
category of \emph{continuous} functors from
$C^{\widetilde{\phantom{,}}}$ to $E$ and the category of
\emph{continuous} functors from $C$ to $E$, a quasi-inverse being
given by $g\rightarrow g^s$ in the notation of \loccit On the other
hand, by definition, $g$ is a morphism of sites if and only if $g^s$
is the inverse-image functor associated with a morphism of topoi
$E\rightarrow C^{\widetilde{\phantom{,}}}$, whence the conclusion by
\hyperlink{sga4.exp.IV.subsubsection.3.2.1}{3.2.1}; the ``or equivalently'' follows from
\hyperrefext[III][III.prop1.6]{1.6}.

The main point to retain from \hyperref[IV.prop4.9.4]{4.9.4} is that
(when finite $\varprojlim$ are representable in $C$) ``it amounts to
the same thing'' to give a morphism of topoi
$f:E\rightarrow C^{\widetilde{\phantom{,}}}$ or to give a functor
$g:C\rightarrow E$ which is left exact and carries covering families
to covering families.

\setcounter{subsubsection}{4}
\subsubsection{}\etiquette[4.9.5]{IV.subsubsec4.9.5}
Using the developments of \hyperref[IV.subsubsec4.9.1]{4.9.1} and
\hyperref[IV.subsubsec4.9.3]{4.9.3}, one sees as usual that, for a
varying $𝒰$-site $C$, via the notions of a morphism of sites and a
morphism between morphisms of sites just made explicit, the topos
$C^{\widetilde{\phantom{,}}}$ depends functorially (or more exactly,
2-functorially) on the site $C$. Note that, with the terminology just
introduced, $C^{\widetilde{\phantom{,}}}$ depends
\emph{covariantly} on the site $C$.

\subsubsection{}\etiquette[4.9.6]{IV.subsubsec4.9.6}
It is immediate that (unlike what happens for the notion of a
cocontinuous functor, \cf \hyperref[IV.subsec4.7]{4.7} 4)
\emph{every} morphism of topoi $F:E\rightarrow E'$ can be realized
by means of a morphism of sites $f:C\rightarrow C'$ (\ie a continuous
functor $f:C'\rightarrow C$ such that \ldots): it suffices to choose
small full generating subcategories $C$\pageoriginale and $C'$ of
$E$ and $E'$, respectively, endowed with the topologies induced by
those of $E$ and $E'$, such that
$$
F^{\ast}(\ob C')\subset \ob C{\quoi},
$$
and to take for $f$ the functor induced by $F^{\ast}$. This explains
why most morphisms of topoi encountered in practice are indeed
described by morphisms of sites (rather than by cocontinuous functors
as in \hyperref[IV.subsec4.7]{4.7}).


% END INCLUDED SGA4_English_current_buildable_source_20260906/source/04/drafts/SGA4_IV_LINES_2189_2354_English.texfrag


% BEGIN INCLUDED SGA4_English_current_buildable_source_20260906/source/04/drafts/SGA4_IV_LINES_2355_2641_English.texfrag
\subsection{Relations between the small and large topoi associated with
a topological space $X$}\etiquette[4.10]{IV.subsec4.10}

We retain the notation of \hyperref[IV.subsec2.5]{2.5}. In particular,
$𝒱$ is a universe such that $𝒰\in 𝒱$. We depart
from the convention of \hyperref[IV.subsec2.1]{2.1} by using
$\Top(X)$ to denote the topos of $𝒱$-sheaves (rather than
$𝒰$-sheaves) on $X$. Thus, we shall work with
$𝒱$-topoi and not $𝒰$-topoi. (NB it would be
possible to retain $𝒰$-topoi by adopting the appropriate
convention for the definition of $\TOP(X)$, so that the latter is a
$𝒰$-topos, \cf \hyperref[IV.subsec2.5]{2.5}). With this understood,
we shall define TWO morphisms of topoi
\begin{equation*}
\begin{cases}
 f: \Top(X)\longrightarrow \TOP(X)\\
 g: \TOP(X)\longrightarrow \Top(X)\quav gf\simeq \id_{\Top(X)}{\quoi},
\end{cases}\tag{4.10.1}
\end{equation*}\etiquette[4.10.1]{IV.eq4.10.1}
giving rise to a sequence of three adjoint functors
\begin{equation*}
g^{\ast}=f_{!}\quav g_{\ast}=f^{\ast}\quav
g^{!}=f_{\ast}{\quoi},\tag{4.10.1.1}
\end{equation*}\etiquette[4.10.1.1]{IV.eq4.10.1.1}
with the notation of \hyperref[IV.subsubsec3.1.3]{3.1.3}. We shall
successively define the\pageoriginale first two functors in this
sequence, the third then being defined as the right adjoint of the
second.

\setcounter{subsubsection}{1}
\subsubsection{}\etiquette[4.10.2]{IV.subsubsec4.10.2}
The functor
$$
g_{\ast}=f^{\ast}:\TOP(X)\longrightarrow \Top(X)
$$
is defined simply as the \emph{restriction functor} from
$(\Esp)_{/X}$ to the site $\Ouv(X)$ of open subsets of $X$; it does
indeed take sheaves to sheaves, as follows immediately from the
definition. Using underlined letters for sheaves on the large site of
$X$, for such a sheaf $\underline{F}$ we denote by
$\underline{F}_X$ its restriction to the site $\Ouv(X)$, thereby
obtaining a functor
\begin{equation*}
\Restr:\underline{F}\longmapsto \underline{F}_X:\TOP(X)\longrightarrow
\Top(X){\quoi}.\tag{4.10.2.1}
\end{equation*}\etiquette[4.10.2.1]{IV.eq4.10.2.1}
It is clear that this functor commutes with $𝒱$-projective
limits, since these are computed pointwise (one could also invoke the
existence of the left adjoint constructed in
\hyperref[IV.subsubsec4.10.4]{4.10.4} below). I claim that it also
commutes with $𝒱$-inductive limits. To see this, we shall give a
very convenient interpretation of ``large'' sheaves on $X$, \ie of
objects of $\TOP(X)$, in terms of ordinary sheaves (NB these are
$𝒱$-sheaves) on topological spaces $X'$ over~$X$.

\subsubsection{}\etiquette[4.10.3]{IV.subsubsec4.10.3}
For a large sheaf $\underline{F}$ on $X$, and for every space $X'$
over $X$ (understood that $X'\in 𝒰$), one defines in the
obvious way, as in \hyperref[IV.subsubsec4.10.2]{4.10.2}, the
``small'' sheaf $\underline{F}_{X'}$, the restriction of
$\underline{F}$ to $X'$. If
$$
u:X''\longrightarrow X'
$$
is\pageoriginale a morphism in $(\Esp)_{/X}$, one then defines in
the obvious way a morphism
$u_{\ast}(\underline{F}_{X''})\rightarrow \underline{F}_{X'}$, or,
what amounts to the same thing, a ``\emph{transition morphism}''
\begin{equation*}
\varphi_u :u^{\ast}(\underline{F}_{X'})\longrightarrow
\underline{F}_{X''}{\quoi}. \tag{4.10.3.1}
\end{equation*}\etiquette[4.10.3.1]{IV.eq4.10.3.1}
As $u$ varies, these morphisms satisfy an evident transitivity
condition for a composite
$$
X'''\xrightarrow{u}X''\xrightarrow{v}X'
$$
of morphisms in $(\Esp)_{/X}$, whose explicit formulation is left to
the reader. In this way one obtains a natural functor from the
category $\TOP(X)$ of large sheaves on $X$ to the category of systems
$$
(\underline{F}_{X'})\quad(X'\in \ob (\Esp)_{/X}), (\varphi_u)\quad(u\in
Fl(\Esp)_{/X}){\quoi},
$$
satisfying the transitivity condition just described and, in
addition, such that for every morphism $u:X''\rightarrow X'$ which is
an \emph{open immersion} (or, more generally, a local homeomorphism),
the transition morphism $\varphi_u$ is an isomorphism. Since the
functors $u^{\ast}:\Top(X')\rightarrow \Top(X'')$ used to describe
the $\varphi_u$ commute with $𝒱$-inductive limits, it follows at
once that, in the preceding description of the objects of $\TOP(X)$
in terms of ``small'' sheaves $F_{X'}$, $𝒱$-inductive limits are
computed pointwise, \ie the functors of the form
$\underline{F}\rightarrow \underline{F}_{X'}$ commute with
$𝒱$-inductive limits.

In particular, this holds for the functor $\underline{F}\mapsto
\underline{F}_X$ considered in \hyperref[IV.subsubsec4.10.2]{4.10.2}.
It therefore does have a right adjoint (\hyperref[IV.cor1.5]{1.5}). It
remains to construct a left adjoint (whose existence also follows a
priori from \hyperref[IV.rem1.8]{1.8}) and to verify that the latter is
left exact. This\pageoriginale will complete the definition of the
morphisms of topoi announced in \eqhyperref[IV.eq4.10.1]{4.10.1} in
terms of the functors \eqhyperref[IV.eq4.10.1.1]{4.10.1.1}.

\subsubsection{}\etiquette[4.10.4]{IV.subsubsec4.10.4}
The functor
$$
g^{\ast}=f_{!}:\Top(X)\longrightarrow \TOP(X)
$$
is obtained by associating with every small sheaf $F$ on $X$ the
system of its inverse images $(F_{X'})$, $X'\in \ob (\Esp)_{/X}$,
under the structural morphisms $X'\rightarrow X$, the transition
morphism $\varphi_u$ for $u:X''\rightarrow X'$ being the transitivity
isomorphism for inverse images of sheaves
(\hyperref[IV.subsubsec4.1.1]{4.1.1}). One sees at once that this
gives a functor
$$
\Prol:\Top(X)\longrightarrow \TOP(X){\quoi},
$$
which is \emph{fully faithful}, and whose essential image consists of
the large sheaves $ℱ$ on $X$ for which all the transition
morphisms $\varphi_u$ of \sisi{4.11.3}{\eqhyperref[IV.eq4.10.3.1]{4.10.3.1}}
are isomorphisms. We leave to the reader the task of defining an
adjunction isomorphism between this canonical prolongation functor and
the restriction functor of \hyperref[IV.subsubsec4.10.2]{4.10.2},
proving that the latter is right adjoint to the former. This is
immediate from the description in \hyperref[IV.subsubsec4.10.3]{4.10.3}
of the category $\TOP(X)$.

\begin{enonce*}[remark]{Remarks 4.10.5}\etiquette[4.10.5]{IV.rem4.10.5}
\begin{enumerate}
\renewcommand{\theenumi}{\alph{enumi}}
\renewcommand{\labelenumi}{\theenumi)}
\item The construction in \hyperref[IV.subsubsec4.10.4]{4.10.4}
also shows that $g^{\ast}=f_{!}$ is fully faithful (or, what amounts
to the same thing, that $g_{\ast}=f^{\ast}$ is a localization
functor, or again that $g^{!}=f_{\ast}$ is fully faithful). The large
sheaves on $X$ belonging to the essential image of the functor
$g^{\ast}=f_{!}=\Prol$ deserve the name of large \emph{étale
sheaves} on $X$, since they form a category equivalent to the category
of ordinary sheaves on $X$, or again to that of étale spaces
over\pageoriginale $X$. It is moreover easy to see that if $X'$ is a
topological space over $X$, then the large sheaf on $X$ represented by
it is étale in the preceding sense if and only if $X'$ is an étale
space over $X$.

\item The full faithfulness of $f_{\ast}$ may also be expressed by
saying that the adjunction morphism
$f^{\ast}f_{\ast}\rightarrow \id_{\Top(X)}$ is an isomorphism, \ie
that $g_{\ast}f_{\ast}\simeq \id$, \ie that $gf\simeq \id$. Thus $g$
makes $\TOP(X)$ into a \emph{topos over} $\Top(X)$, admitting a
``section'' $f$ over $\Top(X)$.

\item The fact that the functor $g^{\ast}=f_{!}$ is fully faithful,
exact, and commutes with $𝒱$-inductive limits partially justifies
the very convenient point of view (due to J. \textsc{Giraud})
according to which, in practically all questions in the theory of
sheaves on $X$, it is harmless to replace the usual sheaves, or
``small'' sheaves, by the associated ``large'' sheaves. This applies
in particular to cohomological questions, since $g_{\ast}$ is exact,
the functors $R^i g_{\ast}$ (\sisi{V\quad}{\hyperrefext[V][V.sec5]{5}})
vanish for $i>0$, and hence (\sisi{V\quad}{\hyperrefext[V][V.prop5.4]{5.4}})
for every \emph{large} sheaf $ℱ$ on $X$ there are canonical
isomorphisms
$$
H^i(\TOP(X),ℱ)\simeq H^i(\Top(X),
ℱ_X)=H^i(X,F){\quoi}.
$$
Applying this to a sheaf of the form $\Prol(F)$, where $F$ is a small
sheaf on $X$, we obtain (since $\Prol(F)_X\simeq F$ canonically) a
canonical isomorphism
$$
H^i(X,F)=H^i(\TOP(X), \Prol(F)){\quoi}.
$$
Thus, \emph{the cohomological invariants of $X$, computed by means of
the small or the large topos of $X$, are essentially identical.} The
same result is moreover valid in noncommutative cohomology.
\end{enumerate}
\end{enonce*}

\begin{enonce*}[remark]{Exercise 4.10.6}\etiquette[4.10.6]{IV.exer4.10.6}
Let\pageoriginale $S$ be a $𝒰$-site whose topology is
coarser than the canonical topology, and let $\underline{M}$ be a
subset of $Fl$ $S$ satisfying the following conditions:
\begin{enumerate}
\renewcommand{\theenumi}{\alph{enumi}}
\renewcommand{\labelenumi}{\theenumi)}
\item The morphisms of $ℳ$ are pullbackable
(\sisi{\hyperrefext[I][I.sec10.3]{10.3}}{\hyperrefext[I][I.sec10.7]{10.7}}),
and $\underline{M}$ is stable under base change.

\item $\underline{M}$ contains the identity arrows and is stable under
composition.

\item An arrow $u:X\rightarrow Y$ such that there exists a covering
family $Y_i\rightarrow Y$, with the arrows
$\fprod{X}{Y_i}{Y}\rightarrow Y_i$ in $\underline{M}$, itself belongs
to $\underline{M}$.

\item For every $X\in \ob S$, every covering family of $X$ is refined
by a covering family $(f_i:X_i\rightarrow X)$, with
$f_i\in \underline{M}$.
\end{enumerate}

For every $X\in \ob S$, consider the site $S(X)$ (the ``small site of
$X$''), whose underlying category is the full subcategory of
$S_{/X}$ consisting of the objects whose structural morphism belongs
to $\underline{M}$, endowed with the topology induced
(\hyperrefext[III][III.subsec3.1]{3.1}) by that of $S$.
\begin{itemize}
\item[$1^{\circ})$] For every arrow $u:X\rightarrow Y$ of $S$, show
that base change by $u$ from $S(Y)$ to $S(X)$ is a continuous
functor, hence gives a functor commuting with inductive limits
(\sisi{III}{\hyperrefext[III][III.prop1.2]{1.2} iv)})
$$
S(u)^x:S(Y)^{\sim}\longrightarrow S(X)^{\sim}{\quoi}.
$$

\item[$2^{\circ})$] Define an equivalence between the topos
$S^{\sim}$ and the category of systems
$$
(F_X)\quad (X\in \ob S)\quav (\varphi_u)\quad (u\in Fl S)
$$
consisting\pageoriginale of objects $F_X\in \ob S(X)^{\sim}$ and,
for every arrow $u:X\rightarrow Y$ in $S$, of a morphism
$\varphi_u:S(u)^x(F_Y)\rightarrow F_X$, these systems being subject
to a transitivity condition for a composite $v\circ u$ of arrows of
$S$, and to the condition that $u\in \underline{M}$ imply that
$\varphi_u$ is an isomorphism.

\item[$3^{\circ})$] Define ``restriction'' and ``prolongation''
functors
$$
\Res_X:(S_{/X})^{\sim}\longrightarrow S(X)^{\sim},
\Prol_X:S(X)^{\sim}\longrightarrow (S_{/X})^{\sim}
$$

Show that $\Res_X$ commutes with small inductive and projective limits
and that $\Prol_X$ is \emph{fully faithful}, its essential image
consisting of the sheaves $F$ such that $\varphi_u$ is an isomorphism
for every arrow $u$ of $S_{/X}$.

\item[$4^{\circ})$] Define an adjunction morphism making $\Res_X$ the
right adjoint of $\Prol_X$. Deduce that there exists a morphism of
topoi
$$
f:S(X)^{\sim}\rightarrow (S_{/X})^{\sim}{\quoi}
$$
making $S(X)^{\sim}$ a subtopos of $(S_{/X})^{\sim}$ and such that
\begin{align*}
f_{!} &= \Prol_X\\
f^x &= \Res_X{\quoi}.
\end{align*}

Show that $\Prol_X$ takes abelian sheaves to abelian sheaves.

\item[$5^{\circ})$] Show that if $S(X)$ admits fiber products, then
$\Prol_X$ is exact. Deduce that there then exists a morphism of topoi
$g:(S_{/X})^{\sim}\rightarrow S(X)^{\sim}$ which is a left retraction
of $f$, \ie such that
\begin{align*}
g^x &= \Prol_X\\
g_x &= \Res_X{\quoi}.
\end{align*}
(For\pageoriginale an example in which $S(X)$ does not admit fiber
products, take for $S$ the category of schemes endowed with the étale
topology, for $M$ the smooth morphisms, and for $X$ a Noetherian
scheme of dimension $>0$.)

\item[$6^{\circ})$] Show that for every abelian sheaf $F$ of
$(S_{/X})^{\sim}$ there is an isomorphism
$$
H^q(X,F)\simeq H^q(X,\Res_X F)\qquad \forall q{\quoi}.
$$
Show, using hypercoverings for example, that for every abelian sheaf
$G$ of $S(X)^{\sim}$ there is an isomorphism
$$
H^q(X,G)\simeq H^q(X,\Prol_X G)\qquad \forall q{\quoi}.
$$

\item[$7^{\circ})$] Buy a chocolate medal for the editor.
\end{itemize}
\end{enonce*}

% END INCLUDED SGA4_English_current_buildable_source_20260906/source/04/drafts/SGA4_IV_LINES_2355_2641_English.texfrag


% BEGIN INCLUDED SGA4_English_current_buildable_source_20260906/source/04/drafts/SGA4_IV_LINES_2642_2987_English.texfrag
\section{Induced topos}\etiquette[5]{IV.sec5}

\subsection{}\etiquette[5.1]{IV.subsec5.1}
Let $E$ be a topos and $X$ an object of $E$. Then the category
$E_{/X}$ of objects of $E$ over $X$ is a topos, as follows, for
example, immediately from Giraud's criterion
\hyperref[IV.theo1.2]{1.2} ii). Alternatively, by \hyperref[IV.cor1.2.1]{1.2.1},
one may realize $E$ as a category of sheaves $C^{\sim}$, where
$C$ is a generating subcategory\pageoriginale of $E$ that may be
chosen so that $X\in \ob C$; one then knows that
$E_{/X}=C^{\sim}_{/X}$ is equivalent to $(C_{/X})^{\sim}$
(\hyperrefext[III][III.prop5.4]{5.4}), and hence is a topos.

The topos $E_{/X}$ is called the \emph{topos induced} on the object
$X$ of $E$.

\subsection{}\etiquette[5.2]{IV.subsec5.2}
We shall define a canonical morphism of topoi
\begin{equation*}
j_X:E_{/X}\longrightarrow E{\quoi},\tag{5.2.1}
\end{equation*}\etiquette[5.2.1]{IV.eq5.2.1}
which is called the \emph{inclusion morphism} of the induced topos
$E_{/X}$ into the ambient topos $E$, or, better (\cf
\hyperref[IV.subsec5.7]{5.7}), the \emph{localization morphism of $E$ at $X$.}
This morphism corresponds to a sequence of three adjoint functors
(\cf \hyperref[IV.subsubsec3.1.3]{3.1.3})
\begin{equation*}
j_{X!}\quav j^{\ast}_X\quav j_{X^{\ast}}{\quoi},\tag{5.2.2}
\end{equation*}\etiquette[5.2.2]{IV.eq5.2.2}
which may be described explicitly as follows:
\begin{enumerate}
\renewcommand{\theenumi}{\alph{enumi}}
\renewcommand{\labelenumi}{\theenumi)}
\item The functor
$$
j_{X!}:E_{/X}\longrightarrow E
$$
is the functor that \emph{forgets the structural arrow} from
$E_{/X}$ to $E$.

\item The functor
$$
j^{\ast}_X:E\longrightarrow E_{/X}
$$
is defined by
$$
j^{\ast}_X(Z)=(X\times Z, \pr_1){\quoi},
$$
where $\pr_1:X\times Z\rightarrow X$ is the first projection. It
may also be interpreted as the \emph{base-change functor} relative
to the morphism
$$
X\longrightarrow e_E{\quoi},
$$
where $e_E$ is the terminal object of $E$. By the definition of
the base-change functor, it is immediate that $j^{\ast}_X$ is
indeed right adjoint to $j_{X!}$. In particular, it commutes with
projective limits. It also commutes with inductive
limits,\pageoriginale by virtue of the special exactness properties
of topoi (\hyperrefext[II][II.sec4]{4}).

\item It follows from this last fact (\hyperref[IV.cor1.6]{1.6}) that $j^{\ast}_X$
has a right adjoint
$$
j_{X^{\ast}}:E_{/X}\longrightarrow E{\quoi}.
$$
For an object $X'$ over $X$, the object $j_{X^{\ast}}(X')$ is also
sometimes denoted by one of the symbols $\prod_{X/e_E}(X'/X)$,
$\Hhom_{X/e_E}(X,X')$, or $\Res_{X/e_E}(X'/X)$ (``Weil
restriction''), one or another of which is doubtless already
familiar to the learned reader of our modest work.
\end{enumerate}

\subsection{}\etiquette[5.3]{IV.subsec5.3}
One should note that the functor $j_{X!}$ commutes with fiber
products and takes monomorphisms to monomorphisms, but it is not in
general left exact for all that. It is left exact only if $X$ is a
terminal object of $E$ (since $X=j_{X!}(X,\id_X)$ and
$(X,\id_X)$ is a terminal object of $E_{/X}$), that is, if and only
if $j_X$ is in fact an equivalence of topoi (and hence all the
functors in \eqhyperref[IV.eq5.2.2]{5.2.2} are equivalences).
\footnote[1]{For a commutation property of $j_{X!}$ under change of
topos, \cf XVII \hyperref[XVII.lem5.1.2]{5.1.2}.}

Likewise, the functor $j_{X^{\ast}}$ is not in general right exact
and does not necessarily take epimorphisms to epimorphisms.

It follows from these observations that the direction of the
morphism of topoi relating $E_{/X}$ to $E$, for an object $X$ of the
topos $E$, is determined without any possible ambiguity.

\subsection{}\etiquette[5.4]{IV.subsec5.4}
The functor $j^{\ast}_X$ is often called the \emph{localization
functor} or the \emph{restriction functor}. The latter terminology
is justified by identifying the objects of $E$ \resp of $E_{/X}$
with the sheaves on $E$ \resp on $E_{/X}$ (\hyperref[IV.theo1.2]{1.2}
iii)), and observing that, under this identification, the adjunction
formula between $j_{X!}$ (the forgetful functor) and $j^{\ast}_X$
may be interpreted as saying that $j^{\ast}_X(F)$ \emph{is\pageoriginale
the restriction to $E_{/X}$ of the sheaf $F$ on $E$} (where
``restriction'' is understood in the generalized sense: composition
with the ``inclusion'' functor $j_{X!}:E_{/X}\rightarrow E$).
In accordance with notation familiar from other contexts, we shall
also often write, interchangeably:
\begin{equation*}
j^{\ast}_X(F)=F| X=F_X=\mbox{ restriction of } F\mbox{ to }
X{\quoi}.\tag{5.4.1}
\end{equation*}\etiquette[5.4.1]{IV.eq5.4.1}

Likewise, when $E$ is of the form $C^{\sim}$, where $C$ is a
(small) site, and $X$ is of the form $\epsilon(S)$ with $S\in
\ob C$, so that one has the equivalence of categories recalled in
\hyperref[IV.subsec5.1]{5.1}
$$
E_{/X}=C^{\sim}_{/\epsilon(S)}\xrightarrow{\approx}(C_{/S})^{\sim}
$$
(where $C_{/S}$ is endowed with the topology induced by that of
$C$), the functor $j^{\ast}_X$ is identified simply with the functor
that ``restricts'' a variable sheaf $F$ on $C$ to the category
$C_{/S}$.

These considerations also make it possible to foresee the important
role that induced topoi and the localization morphisms
\eqhyperref[IV.eq5.2.1]{5.2.1} will play in all questions in which one is led
to argue by ``localization on $E$'' (\cf 8), that is, in practically
all questions involving sites or topoi.

\subsection{}\etiquette[5.5]{IV.subsec5.5}
Let
$$
f:X\longrightarrow Y
$$
be an arrow of $E$, which thus allows us to regard $X$ (or, more
correctly, $(X,f)$) as an object of $E_{/Y}$. It is clear that there
is a canonical isomorphism
\begin{equation*}
(E_{/Y})_{/X}\xrightarrow{\sim}E_{/X}\tag{5.5.1}
\end{equation*}\etiquette[5.5.1]{IV.eq5.5.1}
(\emph{transitivity of induced topoi}). Applying the construction of
\hyperref[IV.subsec5.1]{5.1} to $E_{/Y}$ in place of $E$, one therefore
obtains a canonical morphism of topoi
\begin{equation*}
\loc(f)\mbox{ or } f:E_{/X}\longrightarrow E_{/Y}{\quoi},\tag{5.5.2}
\end{equation*}\etiquette[5.5.2]{IV.eq5.5.2}
also called the\pageoriginale \emph{localization morphism
(associated with $f$).} When $Y$ is the terminal object of $E$, one
essentially recovers \eqhyperref[IV.eq5.2.1]{5.2.1}. The localization morphism
for $f$ is associated with a sequence of three adjoint functors
\begin{equation*}
f_{!}\quav f^{\ast}\quav f_{\ast}{\quoi},\tag{5.5.3}
\end{equation*}\etiquette[5.5.3]{IV.eq5.5.3}
which may be interpreted respectively as a forgetful functor, a
restriction functor, and a functor denoted, according to preference
(with $X'$ as argument), by $\prod_{X/Y}(X'/X)$,
$\Hhom_{X/Y}(X,X')$, or $\Res_{X/Y}(X'/X)$.

\subsection{}\etiquette[5.6]{IV.subsec5.6}
The transitivity of the forgetful functors implies that, for two
composable morphisms
$$
X\xrightarrow{f} Y \xrightarrow{g} Z
$$
of $E$, there is a canonical isomorphism of morphisms of topoi
$$
\loc(gf)\simeq \loc(g)\loc(f):E_{/X}\longrightarrow
E_{/Y}\longrightarrow E_{/Z}{\quoi}.
$$
Moreover, for three composable morphisms, the preceding transitivity
isomorphisms satisfy the usual compatibility relation. Thus, subject
to the customary abuse of language, this allows one to regard
$$
X\longmapsto E_{/X}
$$
as a covariant functor from $E$ to the category
($𝒱$-$𝒰$-$\text{top}$) (\hyperref[IV.subsubsec3.3.1]{3.3.1}).
More precisely, one obtains a 2-functor (not strict in general) from
$E$ to the 2-category ($𝒱$-$𝒰$-$\text{top}$).

Before continuing with the general theory of induced topoi
(\hyperref[IV.subsec5.10]{5.10}), let us give a few instructive examples.

\subsection{}\etiquette[5.7]{IV.subsec5.7}
Let\pageoriginale $X$ be a topological space, and hence let
$\Top(X)$ be a topos (\hyperref[IV.subsec2.1]{2.1}). Let $X'$ be an object
of $\Top(X)$, which it will be convenient to interpret as an
\'etale space over $X$, $p:X'\rightarrow X$. The category
$\Top(X)_{/X'}$ is then identified with the category of \'etale
spaces $X''$ over $X$, equipped with an $X$-morphism
$X''\rightarrow X'$. It is known that such a morphism makes $X''$
an \'etale space over $X'$, and in this way one obtains a canonical
equivalence of topoi
$$
\Top(X)_{/X'}\xrightarrow{\approx} \Top(X'){\quoi}.
$$
The localization morphism is therefore identified with a morphism
of topoi $\Top(X')\rightarrow \Top(X)$, and one sees immediately
that the latter is none other than the morphism
$$
\Top(p):\Top(X')\longrightarrow \Top(X)
$$
associated with the structural continuous map $p:X'\rightarrow X$.
Intuitively, the localization morphism is thus simply the translation,
into the language of topoi, of the \'etale map $p:X'\rightarrow X$.
This both explains the appropriateness of the term ``localization
morphism'' and counsels caution in using the term ``inclusion
morphism'', which seems chiefly appropriate when $p$ is an open
immersion. It will consequently be prudent in
\hyperref[IV.subsec5.1]{5.1} to reserve the term ``inclusion morphism'' for
the case where the morphism $X\rightarrow e_E$ is a monomorphism,
\ie where $X$ is identified with a subobject of the terminal object
of $E$.

\subsection{}\etiquette[5.8]{IV.subsec5.8}
Let $E$ be a topos, $G$ a group in $E$, and $H$ a subgroup of $G$,
and let
$$
X=G/H
$$
be the quotient homogeneous space, regarded as an object of $E$
with a left action by $G$, \ie as an object of the classifying topos
$B_G$ (\hyperref[IV.subsec2.4]{2.4}). We\pageoriginale shall determine the
induced topos
$$
(B_G)_{/X}\quad (X=G/H){\quoi},
$$
by defining an equivalence of topoi
\begin{equation*}
c:B_H \xrightarrow{\approx} (B_G)_{/X}{\quoi},\tag{5.8.1}
\end{equation*}\etiquette[5.8.1]{IV.eq5.8.1}
such that the following diagram commutes up to canonical
isomorphism:
\[\vcenter{
\xymatrix@C=2cm{ B_H \ar[r]^{c} \ar[r] \ar
@{^{(}->}@<-9pt>[rr]_{B_i}& {(B_{G})_{/X}} \ar[r]^{j_X}& B_G
}}\leqno{(5.8.2)}
\]\etiquette[5.8.2]{IV.eq5.8.2}
where $j_X$ is the localization morphism in $B_G$, and where
$$
B_i:B_H\longrightarrow B_G
$$
is the morphism deduced from the inclusion $i:H\rightarrow G$
(\hyperref[IV.subsec4.5]{4.5}).

To define \eqhyperref[IV.eq5.8.1]{5.8.1}, we shall simply give the
descriptions of the inverse-image functor $c^{\ast}$ and the
direct-image functor $c_{\ast}$, leaving it to the reader to verify
that these are indeed quasi-inverse equivalences and give rise to
the commutative diagram \eqhyperref[IV.eq5.8.2]{5.8.2} of morphisms of
topoi. Let $e$ be the subobject of $E$ underlying $X$, the image of
the section of $X$ (over a chosen terminal object $e_E$ of $E$)
deduced from the unit section of $G$. If $X'$ is an object of $B_G$
over \sisi{$G$}{$X$}, then $c^{\ast}(X')$ denotes the inverse image
of $e$ in $X'$,
$$
c^{\ast}(X')=\fprod{X'}{e}{X}{\quoi},
$$
which is stable under the left action of $H$ on $X'$, obtained by
restricting the given action of $G$ on $X'$. This indeed defines a
functor
$$
c^{\ast}:(B_G)_{/X}\longrightarrow B_H{\quoi}.
$$
On the other\pageoriginale hand, if $X'$ is an object of $B_H$, then
the morphism from $X'$ to the terminal object $e_H$ of $B_H$ (\ie
the terminal object $e_E$ of $E$ with the trivial action of $H$)
gives, upon application of the functor $i_{!}$
(\hyperref[IV.subsec4.5]{4.5}), a morphism in $B_G$
$$
i_{!}(X')\longrightarrow i_{!}(e_H)\simeq X{\quoi},
$$
which allows $i_{!}(X')$ to be regarded as an object of
$(B_G)_{/X}$, and hence gives the desired functor
$$
c_{\ast}:X'\longmapsto (i_{!}(X')\longrightarrow
X):B_H\longrightarrow (B_G)_{/X}{\quoi}.
$$

\setcounter{subsubsection}{2}
\subsubsection{}\etiquette[5.8.3]{IV.subsubsec5.8.3}
Thus, by \eqhyperref[IV.eq5.8.2]{5.8.2}, one sees that if $G$ is a group
in a topos and $H$ a subgroup, \emph{the functor ``restriction of
the operators from $G$ to $H$'' may be interpreted as a localization
functor.} In particular, taking $H$ to be the unit subgroup, one sees
that the functor that \emph{forgets the actions of $G$} may be
interpreted as a localization functor. More generally, it follows
that if $(X,G)$ is an object of $E$ with an action of $G$ (\ie an
object of the classifying topos $B_G$), then the functor that
``forgets the actions of $G$'',
$$
E'=(B_G)_{/(X,G)}\longrightarrow E_{/X}
$$
may be interpreted as a localization functor with respect to a
suitable object $E_{G,(X,G)}$ of $E'$ covering the terminal object.
It suffices to take $E_{G,(X,G)}=E_G\times (X,G)$, giving the
cartesian diagram (where $E_G=G_S=G$ with the action of $G$ by left
translation)
\[\vcenter{
\xymatrix@C=2cm@R=1.5cm{
E_G \ar[d]& **[r]E_{G,(X,G)} = Z \ar[l] \ar[d]\\
e_G & (X,G) \ar[l]          \quad , }}\leqno{(5.8.3.1)}
\]\etiquette[5.8.3.1]{IV.eq5.8.3.1}
and\pageoriginale to observe that, by transitivity of induced topoi,
the induced topos $E'=((B_G)_{X/G)/Z}$ is equivalent to the topos
$((B_G)_{/E_G})_{/Z}$; since $(B_G)_{/E_G}$ is equivalent to $E$
by the functor $c^{\ast}$ of \eqhyperref[IV.eq5.8.1]{5.8.1} (with $H=e$),
it suffices to verify that this equivalence takes $Z$ to the object
$X$ of $E$ (which is immediate) in order to deduce the desired
equivalence of categories
$$
E'_{/Z}\xrightarrow{\approx}E_{/X}{\quoi}.
$$
We leave it to the reader to verify that the composite of this
equivalence with the localization functor
$$
E'=(B_G)_{/(X,G)}\longrightarrow E'_{/Z}
$$
is the functor that ``forgets the actions of $G$''. Thus one sees
that, upon passing to induced topoi, the cartesian diagram above gives
a diagram of morphisms of topoi, commutative up to canonical
isomorphism (\hyperref[IV.subsec5.6]{5.6}):
\[\vcenter{
  \xymatrix@C=1.25cm@R=1cm{
**[r](B_G)_{/E_G}\approx E \ar[d]_{j} & **[r](B_G)_{/Z} \approx
E'_{/Z} \ar[l]_(.3){i}
  \approx E_{/X} \ar[d]^{j'}\\
B_G & **[r](B_G)_{/(X,G)} =  E' \ar[l]_(.5){f}\quad,
}}\leqno{(5.8.3.2)}
\]\etiquette[5.8.3.2]{IV.eq5.8.3.2}
where the inverse-image functors associated with the vertical arrows
$j$, $j'$ are the functors that ``forget the actions of $G$'', and
where the inverse-image functor $i^{\ast}$ is identified with the
localization functor $E\rightarrow E_{/X}$. This diagram\pageoriginale is,
moreover, ``2-cartesian'' in the sense of \hyperref[IV.prop5.11]{5.11}
below.

\subsubsection{}\etiquette[5.8.4]{IV.subsubsec5.8.4}
We leave it to the reader to extend the preceding considerations to
the case of a pro-group $𝒢=(G_i)$, in the particular case where $H$
is defined as the ``kernel'' of $𝒢\rightarrow G_{i_{\circ}}$. (In
the case of profinite groups, this means restricting attention to
\emph{open} subgroups $H$ of $G$, \ie subgroups that are closed and
of \emph{finite index}.)

% END INCLUDED SGA4_English_current_buildable_source_20260906/source/04/drafts/SGA4_IV_LINES_2642_2987_English.texfrag


% BEGIN INCLUDED SGA4_English_current_buildable_source_20260906/source/04/drafts/SGA4_IV_LINES_2988_3412_English.texfrag
\begin{enonce*}[remark]{Exercise 5.9}\etiquette[5.9]{IV.exer5.9}
Let $E$ be a topos, $G$ a group of $E$, $B_G$ the classifying topos
of $G$ (\hyperref[IV.subsec2.4]{2.4}), and
$$
\pi:B_G\longrightarrow E
$$
the morphism of topoi deduced from the homomorphism from $G$ to the
unit group $e$ of $E$ (taking into account that $B_e\approx E$).
\begin{enumerate}
\renewcommand{\theenumi}{\alph{enumi}}
\renewcommand{\labelenumi}{\theenumi)}
\item Let
$$
E_G=G_s
$$
be the object of $B_G$ whose underlying $E$-object is $G$, with the
actions of $G$ defined by left translations. Regard $\pi^{\ast}(G)$
as a group of $B_G$ (it is the group $G$ of $E$ with the trivial
action of $G$ on it). Show that the morphism of $E$
$$
G_s\times G\longrightarrow G_s
$$
defined by the right translations of $G$ on $G_s$ is compatible
with the actions of $G$ on $G_s$ and on $G=\pi^{\ast}(G)$, and
defines a morphism of $B_G$
$$
E_G\times \pi^{\ast}(G)\longrightarrow E_G{\quoi}.
$$
Show that this morphism equips $E_G$ with a right action of the
$B_G$-group $\pi^{\ast}(G)$, and that this action makes $E_G$ a
\emph{torsor} under\pageoriginale $\pi^{\ast}(G)$ (that is, the
terminal object of $B_G$ can be covered by objects $X_i$ such that,
for every $i$, the restriction of $E_G$ to $X_i$ is isomorphic to
the trivial bundle with operators $(\pi^{\ast}(G)_{X_i})_d$).

\item Let
$$
q:E'\longrightarrow E
$$
be a topos over $E$. For every topos $B$ over $E$, define the
category $\Hhomtop_E(E',B)$ of morphisms of topoi from $E'$ to $B$
compatible, up to a specified isomorphism, with the structural
morphisms $E'\rightarrow E$ and $B\rightarrow E$. Taking $B=B_G$,
define, by the formula $f\mapsto f^{\ast}(E_G)$, an equivalence of
categories
$$
\Hhomtop_E(E',B_G)\longrightarrow \Ttors(E',q^{\ast}(G)){\quoi},
$$
where the right-hand side denotes the category of (right)
$q^{\ast}(G)$-torsors on $E'$.

\item Let $H$ be a subgroup of $G$, so that $\pi^{\ast}(H)$ is a
subgroup of $\pi^{\ast}(G)$ and hence acts on $E_G$ by restriction
of the operator group, giving a quotient object
$$
X=E_G/\pi^{\ast}(H){\quoi},
$$
which is simply the object $G/H$ endowed with the usual left action
of $G$. Let $e_G$ be the terminal object of $B_G$ (that is, the
terminal object $e_E$ of $E$ with the necessarily trivial action of
$G$), and consider in $B_G$ the diagram of morphisms
\[
\xymatrix@C=-1cm@R=.75cm{
E_G \ar[dd] \ar[rd]& \\
    & \qquad\qquad\quad X = E_{G}/\pi^{*}(H)\ar[ld]\\
e_G & }
\]
which,\pageoriginale upon passing to the topoi induced by $B_G$ on
these objects, gives a diagram of morphisms of topoi, commutative up
to canonical isomorphism (\hyperref[IV.subsec5.6]{5.6}). Show that
this diagram is equivalent to the diagram
\[
\xymatrix@C=.05cm@R=.85cm{
E & = & B_e \ar[dd] \ar[rrrd]& & &\\
  &   &                    &  & &B_H \ar[llld]\\
  &   & B_G                &  & & &,
}
\]
whose three arrows are the arrows associated by
\hyperref[IV.subsec4.5]{4.5} with the group morphisms
$e\rightarrow H\rightarrow G$ in $E$. (Thus, \emph{the classifying
topos $B_H$ may intuitively be interpreted as a homogeneous space
over $B_G$, with group $\pi^{\ast}(G)$, associated with the universal
torsor} (= homogeneous principal bundle) \emph{$E_G$ over $B_G$}.)

\item Reconsider the example of \hyperref[IV.subsec5.7]{5.7} in the
case where $X'$ is an \emph{\'etale covering} (locally trivial, \cf
\hyperref[IV.subsubsec2.7.4]{2.7.4}) of $X$, interpreting it in terms
of the considerations of the present exercise.
\end{enumerate}
\end{enonce*}

\setcounter{subsection}{9}
\subsection{Inverse images of induced topoi}\etiquette[5.10]{IV.subsec5.10}

Let
$$
f:E'\rightarrow E
$$
be a morphism of $𝒰$-topoi, let $X$ be an object of $E$, and let
$X'=f^{\ast}(X)$ be its inverse image in $E'$. We shall define a
diagram of morphisms of topoi
\[\vcenter{
\xymatrix@C=2.5cm@R=1.5cm{
E_{/X}  \ar[d]_{j_X} & E'_{/X'} \ar[l]_{f_{/X}} \ar[d]^{j_{X'}} \\
E      &  E'\ar[l]_{f}\quad , }}\leqno{(5.10.1)}
\]\etiquette[5.10.1]{IV.eq5.10.1}
where\pageoriginale $j_X$ and $j_{X'}$ are the localization
morphisms (\hyperref[IV.subsec5.2]{5.2}), and where $f_{/X}$ is
defined below; this diagram will be commutative up to canonical
isomorphism. We shall define $f_{/X}$ here by means of its inverse
image functor
\begin{equation*}
(f_{/X})^{\ast}:E_{/X}\longrightarrow E'_{/X'}{\quoi},\tag{5.10.2}
\end{equation*}\etiquette[5.10.2]{IV.eq5.10.2}
which we simply take to be the functor ``induced'' by $f^{\ast}$, in
the evident sense. Since the inclusion functors
$j_{X!}:E_{/X}\rightarrow E$ and
$j_{X'!}:E'_{/X'}\rightarrow E'$ commute with $𝒰$-inductive limits
and fiber products, and are conservative, it follows at once that
the functor $(f_{/X})^{\ast}$ commutes with $𝒰$-inductive limits and
fiber products, just as does the functor $f^{\ast}$ which induces
it. Moreover, since $(f_{/X})^{\ast}$ plainly takes the terminal
object $X$ to the terminal object $X'$, it is left exact and hence
is associated with a morphism of topoi
\begin{equation*}
f_{/X}:E'_{/X'}\longrightarrow E_{/X}{\quoi},\tag{5.10.3}
\end{equation*}\etiquette[5.10.3]{IV.eq5.10.3}
defined up to unique isomorphism
(\hyperref[IV.subsubsec3.1.1]{3.1.1}). On the other hand, the diagram
deduced from \eqhyperref[IV.eq5.10.1]{5.10.1} by passing to the
inverse-image functors is commutative (up to canonical isomorphism)
by construction and by the fact that $f^{\ast}$ commutes with the
products $X\times Y$. Thus \eqhyperref[IV.eq5.10.1]{5.10.1} itself
is commutative, up to a unique isomorphism inducing the canonical
isomorphism on the inverse-image functors of the two composites
$f\circ j_{X'}$ and $j_X\circ f_{/X}$ \hyperlink{sga4.exp.IV.subsubsection.3.2.1}{(3.2.1)}.

\setcounter{subsubsection}{3}
\subsubsection{}\etiquette[5.10.4]{IV.subsubsec5.10.4}
When $f:E'\rightarrow E$ is associated with a continuous map
$Y'\rightarrow Y$ of topological spaces
(\hyperref[IV.subsec4.1]{4.1}), so that $X$ is identified with an
\'etale space over $Y$, and $X'$ with the fiber product
$\fprod{X}{Y'}{Y}$, then, identifying the induced topoi
$\Top(Y)_{/X}$ and $\Top(Y')_{/X'}$ with $\Top(X)$ and $\Top(X')$
respectively (\hyperref[IV.subsec5.7]{5.7}), one sees that the
morphism $f_{/X}$ just constructed\pageoriginale is (up to unique
isomorphism) the morphism $\Top(X')\rightarrow\Top(X)$ deduced from
the canonical continuous map $X'\rightarrow X$. In the light of
this example, one may therefore say that the topos $E'_{/X'}$ plays
the role of a fiber product of $E_{/X}$ and $E'$ over $E$. This
intuition is made precise by the following result:

\begin{enonce*}{Proposition 5.11}\etiquette[5.11]{IV.prop5.11}
With the notation of \hyperref[IV.subsec5.10]{5.10}, the diagram
\eqhyperref[IV.eq5.10.1]{5.10.1}, equipped with the compatibility
isomorphism $\alpha:f\circ j_{X'}\rightarrow j_X\circ f_{/X}$, is
``$2$-cartesian.'' More precisely, for every $𝒰$-topos $F$, if to
each morphism of topoi $g:F\rightarrow E'_{/X'}$ one associates the
triple $(f_{/X}\circ g, j_{X'}\circ g,
\alpha\ast g)=(g_1,g_2,\beta)$, where
$g_1:F\rightarrow E_{/X}$ and $g_2:F\rightarrow E'$ are morphisms
of topoi and $\beta:j_X\circ g_1\rightarrow f\circ g_2$ is an
isomorphism of morphisms of topoi from $F$ to $E$, then one obtains
an equivalence between the category $\Hhomtop(F,E'_{/X'})$ and the
category
\sisi{$\fprod{\Hhomtop(F,E_{/X})^2}{\Hhomtop(F,E')}{\Hhomtop(F,E)}$}{$\Hhomtop(F,E_{/X})\times^{2}_{\Hhomtop(F,E)}\Hhomtop(F,E')$}
of all triples $(g_1,g_2,\beta)$ as above.
\end{enonce*}

(NB. A $2$ has been placed above the cartesian-product sign to
recall that this is not an ordinary fiber product of categories, but
a ``$2$-fiber product,'' in the sense made explicit in the
statement.)

To prove \hyperref[IV.prop5.11]{5.11}, we shall describe morphisms of
topoi explicitly by means of their associated inverse-image
functors. For this we shall need the auxiliary results given in
\hyperref[IV.subsubsec5.11.1]{5.11.1} through
\hyperref[IV.prop5.12]{5.12} below.

\setcounter{subsection}{11} \setcounter{subsubsection}{0}
\subsubsection{}\etiquette[5.11.1]{IV.subsubsec5.11.1}
Consider, in general, the situation in which two categories $E$ and
$E'$ are given, together with an object $X$ of $E$ which is assumed
to be pullbackable, that is, such that the functor
$$
j:A\longmapsto A\times X:E\longrightarrow E_{/X}
$$
is\pageoriginale defined, and finally a left exact functor
\begin{equation*}
\varphi:E_{/X}\longrightarrow E'{\quoi}.\tag{$\ast$}
\end{equation*}
Since $E_{/X}$ has a terminal object, namely $(X,\id_X)=X$, the same
is therefore true of $E'$, which has the terminal object
$$
e'\simeq \varphi(X){\quoi}.
$$
Having chosen the terminal object $e'$ of $E'$, we shall associate
with every $\varphi$ as above a pair
\begin{equation*}
(\Psi,u),\mbox{ with } \Psi=\varphi\circ j:E\rightarrow E',\mbox{ and
} u\in \Hom(e',\Psi(X)){\quoi}.\tag{${\ast}{\ast}$}
\end{equation*}
To define $u$, note that
$$
j(X)=X\times X
$$
has a canonical section $\delta_X$ over the terminal object $X$ of
$E_{/X}$, namely the diagonal section, and define
$$
u=\varphi(\delta_X):\varphi(X)=e'\longrightarrow \varphi(X\times
X)=\Psi(X){\quoi}.
$$
I claim that knowledge of the pair $({\ast}{\ast})$ makes it
possible to reconstruct the functor $\varphi$ of $(\ast)$ up to
unique isomorphism. For this, for every object
$$
p:X'\longrightarrow X
$$
of $E_{/X}$, consider the following cartesian diagram in $E_{/X}$:
\[
\xymatrix@C=1cm@R=1.5cm{
X' \ar[d] \ar[r] & **[r]j(X') = X'\times X \ar[d]^{j(p)}\\
X  \ar[r]^{\delta_X}& **[r]j(X) = X \times X   &, }
\]
where the first horizontal arrow is the graph morphism
$\Gamma_p=(\id_{X'},p)$. Applying the left exact functor $\varphi$ to
this diagram, and using the definition of $\Psi$ as $\varphi\circ
j$, gives a cartesian diagram
\[
\xymatrix@C=2.25cm@R=1.25cm{
\varphi(X',p) \ar[d] \ar[r]& \Psi (X') \ar[d]^{\Psi(p)}\\
e'  \ar[r]          & \Psi (X) \quad , }
\]
or,\pageoriginale in other words, a canonical isomorphism
\begin{equation*}
\varphi(X',p)\simeq
\fprod{\Psi(X')}{e'}{\Psi(X)}{\quoi}\tag{5.11.1.1}
\end{equation*}\etiquette[5.11.1.1]{IV.eq5.11.1.1}
It is immediately seen to be functorial in the object $(X',p)$ of
$E_{/X}$, which describes explicitly how the functor
$\varphi:E_{/X}\rightarrow E'$ is reconstructed, up to isomorphism,
from the pair $(\Psi,u)$ of $({\ast}{\ast})$. Notice also that the
functor $\Psi$ is left exact. More generally, $\Psi$ commutes with
every type of projective limit with which $\varphi$ commutes,
because $j$ commutes with projective limits.

Conversely, suppose now that finite projective limits are
representable in $E$ and $E'$, and start with a pair
$$
(\Psi,u),\Psi:E\longrightarrow E', u\in \Hom(e',\Psi(X)){\quoi},
$$
where the functor $\Psi$ is left exact. Formula
\eqhyperref[IV.eq5.11.1.1]{5.11.1.1} then defines a functor
$\varphi:E_{/X}\rightarrow E'$. One checks at once that this functor
is left exact and, more generally, that it commutes with every type
of projective limit representable in $E$ with which $\Psi$ commutes.
This follows immediately from the cartesian diagram
\[
\xymatrix@C=2.5cm@R=1cm{ \displaystyle{\varprojlim^{E}_{I} X'_i}
\ar[d] &
\displaystyle{\varprojlim^{E_{/X}}_{I} X'_i} \ar[d] \ar[l]\\
\displaystyle{\varprojlim^{E}_{I} X} & X \ar[l]_{\delta} }
\]
relating\pageoriginale the projective limits computed in $E$ and in
$E_{/X}$ (where $\delta$ denotes the diagonal morphism into the
projective limit of the constant functor with value $X$): apply the
functor $\Psi$ to it, obtaining a cartesian diagram in $E'$, and
compose the latter with the cartesian square deduced from
$\Psi(X)\leftarrow e'$. The resulting composite cartesian square
expresses the compatibility of $\varphi$ with the projective limit
under consideration.

The pair $(\Psi,u)$ is reconstructed up to isomorphism from
$\varphi$ by observing that there is a canonical isomorphism
\begin{equation*}
\Psi(X')=\varphi(j(X')){\quoi},\tag{5.11.1.2}
\end{equation*}\etiquette[5.11.1.2]{IV.eq5.11.1.2}
deduced from \eqhyperref[IV.eq5.11.1.1]{5.11.1.1} by replacing
$(X',p)$ there with $j(X')=(X'\times X,\pr_2)$ and observing that
$\Psi(X'\times X)\simeq\Psi(X')\times\Psi(X)$. The preceding
isomorphism is manifestly functorial in $X'$, that is, it gives an
isomorphism
$$
\Psi\simeq \varphi\circ j{\quoi},
$$
and one likewise checks that, under this isomorphism,
$u:e'\rightarrow\Psi(X)$ is identified with $\varphi(\delta_X)$.
We have thus obtained the essential part of the following result.

\begin{enonce*}{Lemma 5.11.2}\etiquette[5.11.2]{IV.lem5.11.2}
Let $E$ and $E'$ be two categories in which finite projective limits
are representable, let $X$ be an object of $E$, let $\Ssex(E,E')$
be the category of left exact functors $\Psi$ from $E$ to $E'$, and
let $\Ssex(E_{/X},E')$ be the analogous category of left exact
functors $\varphi$ from $E_{/X}$ to $E'$. There is then an
equivalence between the category $\Ssex(E_{/X},E')$ and the category
$\Ssex(E,E')_{/\Gamma}$ of pairs $(\Psi,u)$, where $\Psi\in\ob
\Ssex(E,E')$ and $u:e'\rightarrow\Psi(X)$ (with $e'$ the terminal
object of $E'$). This equivalence, as well as a quasi-inverse
functor, is described explicitly in
\hyperref[IV.subsubsec5.11.1]{5.11.1}. If $\varphi$ and $(\Psi,u)$
correspond, then $\varphi$ commutes with a given type of projective
limit if\pageoriginale and only if $\Psi$ does. If the base-change
functors in $E$ and $E'$ commute with a given type of inductive
limit, then $\varphi$ commutes with inductive limits of that type if
and only if $\Psi$ does.
\end{enonce*}

We leave it to the reader to make explicit the category structure on
$\Ssex(E,E')_{/\Gamma}$ corresponding to the pairs $(\Psi,u)$. If
the functor $\Gamma:\Ssex(E,E')\rightarrow(\Ens)$ is defined by
$\Gamma(\Psi)=\Hom(e',\Psi(X))$, then $u$ in the pair $(\Psi,u)$ may
be interpreted as an element of $\Gamma(\Psi)$, that is, as a
homomorphism $\Psi\rightarrow\Gamma$ in the category of presheaves
on $\Ssex(E,E')^{\circ}$. This justifies the notation
${}_{/\Gamma}$ used in the statement of
\hyperref[IV.subsubsec5.11.1]{5.11.1}. To finish the proof of
\hyperref[IV.lem5.11.2]{5.11.2}, it remains to make explicit in
\hyperref[IV.subsubsec5.11.1]{5.11.1} the functorial nature of the
constructions considered when $\varphi$, respectively $(\Psi,u)$,
varies. This is essentially trivial and is left to the reader. It
also remains to verify the assertion concerning commutation with
inductive limits; this too follows immediately from the explicit
formulas \eqhyperref[IV.eq5.11.1.1]{5.11.1.1} and
\eqhyperref[IV.eq5.11.1.2]{5.11.1.2} relating these functors.

In particular, when $E$ and $E'$ are $𝒰$-topoi, and morphisms of
topoi are interpreted by means of their associated inverse-image
functors, one obtains:

\begin{enonce*}{Proposition 5.12}\etiquette[5.12]{IV.prop5.12}
Let $E$ and $E'$ be two $𝒰$-topoi, let $X$ be an object of $E$, and
consider on $\Hhomtop(E',E)$ the contravariant functor
$$
\gamma:f\longmapsto
\Gamma(E',f^{\ast}(X))=\Hom(e',f^{\ast}(X)):\Hhomtop(E',E)^{\circ}\longrightarrow
(\Ens){\quoi}.
$$
There is then an equivalence of categories
\begin{equation*}
\Hhomtop(E',E_{/X})\xrightarrow{\approx}\Hhomtop(E',E)_{/\gamma}{\quoi},\tag{5.12.1}
\end{equation*}\etiquette[5.12.1]{IV.eq5.12.1}
where\pageoriginale the right-hand side is the full subcategory of
$\Hhomtop(E',E)^{\widehat{\phantom{,}}}_{/\gamma}$ consisting of
pairs $(f,u)$ with $f\in\Homtop(E',E)$ and $u\in\gamma(f)$, that is,
$u\in\Gamma(E',f^{\ast}(X))$. This functor is obtained by
associating with every morphism of topoi $h:E'\rightarrow E_{/X}$
the pair $(f,u)$ consisting of the composite
$$
f=j_X\circ h:E'\rightarrow E_{/X}\rightarrow E{\quoi},
$$
and the morphism
$$
u:e'=h^{\ast}(X,\id_X)\longrightarrow
f^{\ast}(X)=h^{\ast}(j^{\ast}_X(X))=h^{\ast}(X\times X)
$$
deduced from the diagonal morphism $\delta_X:X\rightarrow X\times
X$ by applying $h^{\ast}$ to it.
\end{enonce*}

\setcounter{subsection}{12}
\subsubsection{}\etiquette[5.12.2]{IV.subsubsec5.12.2}
Using \hyperref[IV.prop5.12]{5.12}, the proof of
\hyperref[IV.prop5.11]{5.11} becomes almost evident, and essentially
reduces to the following (which will serve in place of a formal
proof that would be no more instructive): to give a morphism of
topoi $g:F\rightarrow E'_{/X'}$ ``amounts to the same thing'' as to
give a morphism of topoi $g_2:F\rightarrow E'$ and a section $u$ of
$g^{\ast}_2(X')$. But
$g^{\ast}_2(X')=g^{\ast}_2(f^{\ast}(X))=(fg_2)^{\ast}(X)$, so giving
$u$ is also equivalent to giving a lifting of the morphism of topoi
$fg_2:F\rightarrow E$ to a morphism of topoi
$g_1:F\rightarrow E_{/X}$. This expresses precisely that giving a
morphism of topoi $g:F\rightarrow E'_{/X'}$ is essentially
equivalent to giving a triple $(g_1,g_2,\alpha)$ as in
\hyperref[IV.prop5.11]{5.11}.

\begin{enonce*}[remark]{Remark 5.13}\etiquette[5.13]{IV.rem5.13}
We shall see (\S 15\nde{This section does not exist. See, however,
[16], Proposition 4.47, for a general existence statement for fiber
products of topoi (noting that all geometric morphisms between
Grothendieck topoi are bounded).}) that for every diagram
\[
\xymatrix@C=2cm@R=1.25cm{
E_1 \ar[d]&\\
E   & E_2 \ar[l] }
\]
of\pageoriginale topoi, there exists a $2$-fiber product in the sense
of the $2$-category ($𝒱$-$𝒰$-$\Top$) of $𝒰$-topoi belonging to
$𝒱$ (independent, up to equivalence, of the choice of a universe
$𝒱$ such that $𝒰\in𝒱$). For natural examples of ``fiber products
of topoi'' other than \hyperref[IV.prop5.11]{5.11}, see the final
chapter of \textsc{Giraud}'s book [3].
\end{enonce*}

\begin{enonce*}[remark]{Exercise 5.14}\etiquette[5.14]{IV.exer5.14}
\begin{enumerate}
\renewcommand{\theenumi}{\alph{enumi}}
\renewcommand{\labelenumi}{\theenumi)}
\item Let $F$ and $G$ be two topoi over a topos $E$. Define the
category $\Hhomtop_E(F,G)$ of pairs $(f,\alpha)$, where
$f:F\rightarrow G$ is a morphism of topoi and $\alpha:p\rightarrow
q\circ f$ is an isomorphism of morphisms of topoi from $F$ to $E$
(where $p:F\rightarrow E$ and $q:G\rightarrow E$ are the structural
morphisms of topoi). Define pairing functors
$\Hhomtop_E(F,G)\times\Hhomtop_E(G,H)\rightarrow\Hhomtop_E(F,H)$,
and associativity isomorphisms.

\item Let $X$ and $Y$ be two objects of a topos $E$. Define a
natural functor from the discrete category defined by the set
$\Hom(X,Y)$ to the category
$\Hhomtop_E(E_{/X},E_{/Y})$. Establish compatibility with the
composition of morphisms in $E$ and with the pairings considered in
a).

\item Prove that the functor considered in b) is an equivalence of
categories. In particular, an object $X$ of a topos $E$ can be
reconstructed up to unique isomorphism when one knows, ``up to
$E$-equivalence,'' the induced topos $E_{/X}$ as a topos
\emph{over $E$.}~\footnote[1]{This result is due to P.
\textsc{Deligne}. The case where $E$ is the punctual topos had
previously been treated by Mme \textsc{Hakim}.}
\end{enumerate}
\end{enonce*}

% END INCLUDED SGA4_English_current_buildable_source_20260906/source/04/drafts/SGA4_IV_LINES_2988_3412_English.texfrag


% BEGIN INCLUDED SGA4_English_current_buildable_source_20260906/source/04/drafts/SGA4_IV_LINES_3413_3631_English.texfrag
\section{Points of a topos and fiber functors}\etiquette[6]{IV.sec6}

\setcounter{subsection}{-1}
\subsection{}\etiquette[6.0]{IV.subsec6.0}
Let $P$ be the standard $𝒰$-\emph{punctual topos}
(\hyperref[IV.subsec2.2]{2.2})
$$
P=(𝒰\text{-}\Ens){\quoi}.
$$
In view of the rather different intuitions attached on the one hand
to the symbol\pageoriginale $P$, representing a geometric object of
the nature of a point, and on the other hand to
$(𝒰\text{-}\Ens)$, representing a ``large category'' to be
interpreted as ``the category of sheaves on $P$,'' one should use one
or the other symbol exclusively according to context, although from
a strictly logical viewpoint they denote one and the same object.

\begin{enonce*}{Definition 6.1}\etiquette[6.1]{IV.def6.1}
Let $E$ be a $𝒰$-topos. A point of $E$ is a morphism of topoi
$P\rightarrow E$ from the punctual topos $P$
(\hyperref[IV.subsec6.0]{6.0}) to $E$. The category
$\Hhomtop(P,E)$ (\hyperref[IV.subsec3.2]{3.2}) is called the category
of points of $E$ and is denoted by $\Ppoints(E)$ or
$\sheaf{\Pt}(E)$. If $p:P\rightarrow E$ is a point of $E$, with
associated inverse-image functor
$p^{\ast}:E\rightarrow(𝒰\text{-}\Ens)$, and if $F$ is an object of
$E$, the set $p^{\ast}(F)$ is called the fiber of $F$ at $p$ and is
denoted by $F_p$.
\end{enonce*}

\setcounter{subsection}{1}
\subsubsection{}\etiquette[6.1.1]{IV.subsubsec6.1.1}
Of course, when $F$ is a group object (\resp ring object,
\resp\ldots) of $E$, its fiber $F_p$ at $p$ is a group (\resp a ring,
\resp\ldots) (\hyperref[IV.subsubsec3.1.2]{3.1.2}).

\setcounter{subsubsection}{1}

\subsubsection{}\etiquette[6.1.2]{IV.subsubsec6.1.2}
By \hyperlink{sga4.exp.IV.subsubsection.3.2.1}{3.2.1}, the category of points of $E$ is equivalent, via the
functor $p\mapsto p^{\ast}$, to the \emph{opposite} of the category
of functors
$$
\varphi:E\longrightarrow (𝒰\text{-}\Ens)
$$
which commute with $𝒰$-inductive limits and are left exact. Thus it
amounts essentially to the same thing to give a point of $E$ or to
give such a functor $\varphi$.

\begin{enonce*}{Definition 6.2}\etiquette[6.2]{IV.def6.2}
A fiber functor of the $𝒰$-topos $E$ is a functor
$\varphi:E\rightarrow(𝒰\text{-}\Ens)$ which commutes with
$𝒰$-inductive limits and is left exact. The category of fiber
functors on $E$, denoted by $\Ffib(E)$, is the\pageoriginale full
subcategory of $\Hhom(E,(𝒰\text{-}\Ens))$ formed by the fiber
functors on $E$.
\end{enonce*}

\setcounter{subsection}{2}
\subsubsection{}\etiquette[6.2.1]{IV.subsubsec6.2.1}
By \hyperref[IV.subsubsec6.1.2]{6.1.2}, the category $\Ffib(E)$ of
fiber functors on $E$ is therefore equivalent to the \emph{opposite}
of the category $\Ppoint(E)$ of points of $E$, via the functor
$p\mapsto p^{\ast}$ from the latter to the former:
\begin{equation*}
\Ppoint(E)\xrightarrow{\approx}\Ffib(E)^{\circ}{\quoi}.
\tag{6.2.1.1}
\end{equation*}\etiquette[6.2.1.1]{IV.eq6.2.1.1}

A functor
$$
\varphi:E\longrightarrow (𝒰\text{-}\Ens)
$$
is a fiber functor if and only if it is the fiber functor
$F\mapsto F_p$ associated with a point $p$ of $E$
(\hyperref[IV.def6.1]{6.1}). In other words, the functor
\eqhyperref[IV.eq6.2.1.1]{6.2.1.1} is surjective. Let us also note
that a functor $\varphi$ is a fiber functor if and only if it is left
exact and carries covering families to surjective families
(\hyperref[IV.cor1.7]{1.7}). This last condition may also be expressed
by saying that $\varphi$ commutes with $𝒰$-sums and carries
epimorphisms to epimorphisms.

\subsection{}\etiquette[6.3]{IV.subsec6.3}
Let $C\in 𝒰$ be a site. A \emph{point of the site} $C$ is a point of
the topos $C^{\sim}$, and the \emph{category of points of the site}
$C$ is the category
$$
\Ppoint(C)=\Ppoint(C^{\sim}){\quoi}.
$$
Consider, on the other hand, the functor
\begin{equation*}
\varphi\longrightarrow \varphi|
C=\varphi\circ\epsilon_C:\Ffib(C^{\sim})\longrightarrow
\Hhom(C,(𝒰\text{-}\Ens)){\quoi},\tag{6.3.1}
\end{equation*}\etiquette[6.3.1]{IV.eq6.3.1}
which is fully faithful and whose essential image is the full
subcategory $\Mmorsite((𝒰\text{-}\Ens),C)^{\circ}$
(\hyperref[IV.prop4.9.4]{4.9.4}), which we shall also denote by
$\Ffib(C)$. A functor
$$
\Psi:C\longrightarrow (𝒰\text{-}\Ens)
$$
is\pageoriginale called a \emph{fiber functor on the site} $C$ if it
belongs to the essential image of \eqhyperref[IV.eq6.3.1]{6.3.1},
just described. Such a functor is therefore \emph{continuous} (a
fortiori (III \hyperref[III.prop1.6]{1.6}), it carries covering families to covering
families, \ie to surjective families), and it is left exact. When
finite projective limits are representable in $C$ (a condition
satisfied in almost every case encountered in practice), the
preceding properties \emph{characterize} the fiber functors on $C$:
they are then precisely the functors
$\Psi:C\rightarrow(𝒰\text{-}\Ens)$ which are left exact and carry
covering families to surjective families
(\hyperref[IV.prop4.9.4]{4.9.4}).

One should retain that the functor $\varphi\rightarrow\varphi|C$ is
an equivalence of categories
$$
\Ffib(C^{\sim})\xrightarrow{\approx}\Ffib(C){\quoi},
$$
so that \emph{it amounts essentially to the same thing to give a
fiber functor on the topos $C^{\sim}$} (or equivalently a point of
that topos) \emph{or to give a fiber functor on the site $C$.} Given
a fiber functor $\Psi$ on $C$, arising up to isomorphism from a fiber
functor $\varphi$ on $C^{\sim}$, the latter is reconstructed from
$\Psi$, up to canonical isomorphism, by the formula
\begin{equation*}
\varphi(F)\simeq \varinjlim_{C_{/F}}\Psi (X){\quoi},\tag{6.3.2}
\end{equation*}\etiquette[6.3.2]{IV.eq6.3.2}
where $C_{/F}$ is the category of objects $X$ of $C$ endowed with a
morphism $X\rightarrow F$ (in $\widehat{C}$), \ie an element of
$F(X)$. Formula \eqhyperref[IV.eq6.3.2]{6.3.2} is an immediate
consequence of the fact that $\varphi$ commutes with inductive limits
and of the canonical isomorphism, functorial in $F$
(\hyperrefext[II][II.cor4.1.1]{4.1.1}):
$$
F\simeq \varinjlim_{C_{/F}}\epsilon_C(X){\quoi}.
$$
More\pageoriginale generally, when $G$ is a presheaf on $C$, there
is a canonical isomorphism, functorial in $G$:
\begin{equation*}
\varphi(\underline{a}G)\simeq
\varinjlim_{C_{/G}}\Psi(X){\quoi},\tag{6.3.3}
\end{equation*}\etiquette[6.3.3]{IV.eq6.3.3}
where $\underline{a}G$ is the sheaf associated with $G$
(\hyperrefext[II][II.cor4.1.1]{4.1.1}). Indeed, one has the formula
$$
\underline{a}G\simeq \varinjlim_{C_{/G}}\epsilon_C(X){\quoi}.
$$

\setcounter{subsection}{4} \setcounter{subsubsection}{-1}
\subsubsection{}\etiquette[6.4.0]{IV.subsubsec6.4.0}
We refer to \hyperrefext[I][I.def6.1]{6.1} for the notion of a
\emph{conservative family of functors}
$$
\varphi_i:E\longrightarrow E_i\qquad i\in I{\quoi}.
$$
Note that when $E$ and the $E_i$ are $𝒰$-topoi and the functors
$\varphi_i$ are inverse-image functors $f_i^{\ast}$ associated with
morphisms of topoi
$$
f_i:E_i\longrightarrow E{\quoi},
$$
all the exactness properties assumed in
\hyperrefext[I][I.prop6.2]{6.2}, \hyperrefext[I][I.cor6.3]{6.3}, and
\hyperrefext[I][I.prop6.4]{6.4} are satisfied, provided that in
\hyperrefext[I][I.prop6.2]{6.2} (v) $D$ is assumed finite. The
following conditions are then equivalent: $(f_i^{\ast})_{i\in I}$ is
conservative; it is conservative for \emph{monomorphisms}
(\hyperrefext[I][I.def6.1]{6.1}); it is conservative for
\emph{epimorphisms}; and it is \emph{faithful}.

When\pageoriginale the family of functors $f_i^{\ast}$ is
conservative, one also sometimes says that \emph{the family of
morphisms of topoi $(f_i)_{i\in I}$ is conservative}. In particular:

\begin{enonce*}{Definition 6.4.1}\etiquette[6.4.1]{IV.def6.4.1}
Let $E$ be a $𝒰$-topos and let $(p_i:P\rightarrow E)_{i\in I}$ be a
family of points of $E$. This family of points is called conservative
if the family of associated fiber functors $F\mapsto F_{p_i}$ from
$E$ to $(𝒰\text{-}\Ens)$ is conservative
(\hyperref[IV.subsubsec6.4.0]{6.4.0}). The topos $E$ is said to have
enough points if it admits a conservative family of points (\ie if
the family of all points of the topos is conservative).
\end{enonce*}

\setcounter{subsection}{4} \setcounter{subsubsection}{1}
\subsubsection{}\etiquette[6.4.2]{IV.subsubsec6.4.2}
Let us note that all $𝒰$-topoi used thus far have enough points. It
is nevertheless possible, ``on purpose,'' to construct topoi which
do not have enough points (\hyperref[IV.exer7.2.6]{7.2.6} e) and
\hyperref[IV.subsec7.4]{7.4}); such a topos is necessarily not
``empty'' in the geometric sense of \hyperref[IV.subsec2.2]{2.2}. A
topos having enough points admits a conservative family of points
indexed by some $I\in𝒰$ (\hyperref[IV.prop6.5]{6.5}). Note, however,
that if $E$ is a $𝒰$-topos, the set of isomorphism classes of points
of $E$ is not necessarily $𝒰$-small (\hyperref[IV.subsec7.3]{7.3}),
although it is so in many cases encountered in practice. Finally,
for an interesting existence theorem (due to P. Deligne) providing
enough points and covering all cases arising in algebraic
geometry\sisi{.}{, see \hyperrefext[VI][VI.sec9]{9}.}

\subsubsection{}\etiquette[6.4.3]{IV.subsubsec6.4.3}
When $(p_i)_{i\in I}$ is a conservative family of points of the
topos $E$, the remarks of \hyperrefext[I][I.prop6.2]{6.2} apply. In
particular, they imply that two arrows $u,v:F\rightrightarrows G$ in
$E$ are equal if and only if, for every $i\in I$, the induced arrows
on fibers $u_{p_i},v_{p_i}:F_{p_i}\rightrightarrows G_{p_i}$
are\pageoriginale equal; that an arrow $u:F\rightarrow G$ in $E$ is
a monomorphism (\resp an epimorphism) if and only if, for every
$i\in I$, the map $u_{p_i}:F_{p_i}\rightarrow G_{p_i}$ is so; that
an object $F$ of $E$ is initial (\resp terminal) if and only if, for
every $i\in I$, $F_{p_i}$ is empty (\resp consists of one point); and
that an object $F$ is a subobject of the terminal object $e_E$ if and
only if, for every $i\in I$, $F_{p_i}$ has at most one point.


% END INCLUDED SGA4_English_current_buildable_source_20260906/source/04/drafts/SGA4_IV_LINES_3413_3631_English.texfrag


% BEGIN INCLUDED SGA4_English_current_buildable_source_20260906/source/04/drafts/SGA4_IV_LINES_3632_3863_English.texfrag
\begin{enonce*}{Proposition 6.5}\etiquette[6.5]{IV.prop6.5}
\begin{enumerate}
\renewcommand{\theenumi}{\alph{enumi}}
\renewcommand{\labelenumi}{\theenumi)}
\item Let $C$ be a $𝒰$-site and let $(\varphi_i)_{i\in I}$ be a
family of fiber functors on $C$ (\hyperref[IV.subsec6.3]{6.3}). The
family $(p_i)_{i\in I}$ of corresponding points of the topos
$E=C^{\sim}$ is conservative if and only if every family
$(X_j\rightarrow X)_{j\in J}$ in $C$ for which, for every $i\in I$,
the corresponding family
$(\varphi_i(X_j)\rightarrow\varphi_i(X))_{j\in J}$ is surjective is
a covering family.

\item Let $E$ be a $𝒰$-topos. If $E$ has enough points
(\hyperref[IV.def6.4.1]{6.4.1}), then $E$ admits a $𝒰$-small
conservative family of points.
\end{enumerate}
\end{enonce*}



Assertion b) is a special case of \hyperrefext[I][I.prop7.7]{7.7}.
To prove a), apply \hyperrefext[I][I.prop7.7]{7.7} to the generating
family in $E$ formed by the $\epsilon(X)$, $X\in\ob C$. Recall
(\hyperrefext[II][II.theo4.4]{4.4}) that the family
$X_j\rightarrow X$ is covering if and only if the family
$\epsilon(X_j)\rightarrow\epsilon(X)$ is epimorphic. If the family of
points $(p_i)_{i\in I}$ is conservative, this amounts
(\hyperref[IV.subsubsec6.4.3]{6.4.3}) to saying that, for every
$i\in I$, the family $(X_j)_{p_i}\rightarrow(X)_{p_i}$ is
surjective, \ie that the family
$\varphi_i(X_j)\rightarrow\varphi_i(X)$ is surjective. This proves
the ``only if'' part of a). For the ``if'' part, apply the criterion
\hyperrefext[I][I.prop7.7]{7.7}, which reduces us to verifying that
every monomorphism $F\rightarrow\epsilon(X)$\pageoriginale for which
$F_{p_i}\rightarrow\epsilon(X)_{p_i}$ is an isomorphism for every
$i\in I$ is an isomorphism, or equivalently an epimorphism. Now one
can find a covering family of morphisms
$\epsilon(X_j)\rightarrow F$, and, after refining further, one may
suppose that the morphisms
$\epsilon(X_j)\rightarrow\epsilon(X)$ are induced by morphisms
$X_j\rightarrow X$. By hypothesis, for every $i\in I$, the composite
morphisms $(X_j)_{p_i}\rightarrow F_{p_i}\rightarrow(X)_{p_i}$ form
a surjective family (as the composite of two surjective families),
\ie the family
$\varphi_i(X_j)\rightarrow\varphi_i(X)$ is surjective. It follows
from the hypothesis that the family $X_j\to X$ is covering, hence
that the family $\epsilon(X_j)\rightarrow\epsilon(X)$ is
epimorphic, and a fortiori that $F\rightarrow\epsilon(X)$ is an
epimorphism.

\begin{enonce*}{Corollary 6.5.1}\etiquette[6.5.1]{IV.cor6.5.1}
Let $C$ be a category. A topology on $C$ making $C$ a $𝒰$-site with
enough fiber functors (\ie such that the associated topos has enough
points) is completely determined by the full subcategory $\Phi$ of
$\Hhom(C,\Ens)$ formed by the fiber functors on $C$.
\end{enonce*}

Indeed, by \hyperref[IV.prop6.5]{6.5} a), one can then describe the
covering families in terms of the family of fiber functors.

In fact, we shall see below
\eqhyperref[IV.subsubsec6.8.5]{6.8.5} that the
subcategory $\Phi$ is contained in the category of pro-representable
functors on $C$, so that
$\Phi^{\circ}\simeq\Ppoint(C^{\sim})$ identifies, up to equivalence,
with a full subcategory of $\Pro-C$.

Compare \hyperref[IV.cor6.5.1]{6.5.1} with \textsc{Giraud}'s theorem
\hyperrefext[II][II.theo5.5]{5.5}, which establishes a one-to-one
correspondence between the set of topologies on $C$ and a certain set
of full subcategories of $\Hhom(C^{\circ},\Ens)$.

\begin{enonce*}[remark]{Exercise 6.5.2}\etiquette[6.5.2]{IV.exer6.5.2}
Let\pageoriginale $C$ be a category equivalent to a category in $𝒰$,
and let $P$ be a subcategory of $\Pro(C)$. For every $p\in P$, denote
by $X\mapsto X_p$ the functor pro-represented by $p$. A family
$(X_i\rightarrow X)_{i\in I}$ in $C$ is called
$P$-\emph{covering} if, for every $p\in P$, the family
$(X_{ip}\rightarrow X_p)_{i\in I}$ is surjective. Show that there is
a topology $T_P$ on $C$ whose covering families are precisely the
$P$-covering families, and that $T_P$ is the finest topology on $C$
for which the functors $X\mapsto X_p$, $p\in P$, are fiber functors.
Show that the topology $T_P$ makes $C$ a site with \emph{enough}
fiber functors. Show that, for every topology $T$ on $C$, among the
topologies $T'$ finer than $T$ which have enough fiber functors there
is a coarsest one: it is the topology $T_P$, where $P$ is the
strictly full subcategory of $\Pro(C)$, equivalent to
$\Ppoint(C)$, described in
\eqhyperref[IV.subsubsec6.8.5]{6.8.5} below.
\end{enonce*}

\begin{enonce*}[remark]{Problem 6.5.3}\etiquette[6.5.3]{IV.prop6.5.3}
Characterize the subcategories (strictly full, stable under
$𝒰$-filtered projective limits, \ldots) $P$ of $\Pro(C)$ which can
arise from a topology on $C$ as the essential image of
$\Ppoint(C^{\sim})$ (\hyperref[IV.subsubsec6.8.5]{6.8.5}).
\footnote[1]{\cf \hyperref[IV.exer9.1.8]{9.1.8} a) and e) for
necessary conditions, and \hyperref[IV.exer9.1.12]{9.1.12} b) for
necessary and sufficient conditions in the cases of $\Top(X)$, with
$X$ a sober locally Noetherian space.}
\end{enonce*}

\setcounter{subsection}{5}
\subsection{}\etiquette[6.6]{IV.subsec6.6}
Let
$$
f:E\longrightarrow E'
$$
be a morphism of $𝒰$-topoi. It induces a canonical functor
$$
\Ppoint(f)=(p\longmapsto f\circ p):\Ppoint(E)\longrightarrow
\Point(E'){\quoi}.
$$
Given two composable morphisms of topoi
$$
E\xrightarrow{f} E'\xrightarrow{g} E''{\quoi},
$$
one\pageoriginale checks immediately that
$$
\Ppoint(gf)=\Ppoint(g)\circ\Ppoint(f):\Ppoint(E)\longrightarrow
\Ppoint(E')\longrightarrow \Ppoint(E''){\quoi}.
$$
This makes precise the functorial dependence (without any abuse of
language for once) of $\Ppoint(E)$ on the topos $E$: if $𝒱$ is a
universe such that $𝒰\in𝒱$, one obtains a (genuine!) functor
$$
\Ppoint:(𝒱\text{-}𝒰\text{-}\text{top})\longrightarrow
(𝒱\text{-}\cat)
$$
from the category of $𝒰$-topoi which are elements of $𝒱$ to the
category of categories which are elements of $𝒱$.

\subsection{}\etiquette[6.7]{IV.subsec6.7}
Let $E$ be a $𝒰$-topos. Every object $F$ of $E$ then defines a
contravariant functor
$$
p\longmapsto F_p=p^{*}(F):\Ppoint(E)^{\circ}\longrightarrow
(𝒰\text{-}\Ens){\quoi},
$$
and hence, for varying $F$, a canonical functor
\begin{equation*}
E\longrightarrow
\Ppoint(E)^{\widehat{\phantom{,}}}=\Hhom(\Ppoint(E)^{\circ},
(𝒰\text{-}\Ens)){\quoi},\tag{6.7.1}
\end{equation*}\etiquette[6.7.1]{IV.eq6.7.1}
which, by definition (\hyperref[IV.prop6.5]{6.5}), is faithful if and
only if $E$ has enough points. (The functor
\eqhyperref[IV.eq6.7.1]{6.7.1} bears a certain formal analogy to a
Fourier transform\ldots)

Let $X$ be an object of $E$, which therefore defines a presheaf
$\widehat{X}$ on
$$
C=\Ppoint(E){\quoi}.
$$
We may therefore define the category
$$
C_{/\widehat{X}}=\Ppoint(E)_{/\widehat{X}}
$$
of morphisms $p\rightarrow\widehat{X}$ in the category
$\widehat{C}$ of presheaves on $C$, \ie the category of pairs
$(p,\xi)$ where $p$ is a point of $E$ and $\xi$ is an element of
$\widehat{X}(p)=X_p$, the fiber of $X$ at $p$. With this understood,
I claim that \emph{there is a canonical equivalence of categories}
\begin{equation*}
C'=\Ppoint(E_{/X})\xrightarrow{\approx}C_{/\widehat{X}}\quav \mbox{
where }\quad C=\Ppoint(E){\quoi},\tag{6.7.2}
\end{equation*}\etiquette[6.7.2]{IV.eq6.7.2}
whose\pageoriginale composite with the inclusion functor
$C_{/\widehat{X}}\rightarrow C$ is none other than the functor
$$
\Ppoint(j_X):C'=\Ppoint(E_{/X})\longrightarrow C=\Ppoint(E)
$$
induced functorially (\hyperref[IV.subsec6.6]{6.6}) by the
localization morphism \eqhyperref[IV.eq5.2.1]{5.2.1}
$$
j_X:E_{/X}\longrightarrow E{\quoi}.
$$
This is simply the special case of \hyperref[IV.prop5.12]{5.12}
obtained by taking $E'=P=(\Ens)$.

\begin{enonce*}{Corollary 6.7.3}\etiquette[6.7.3]{IV.cor6.7.3}
Let $E$ be a topos. If $E$ has enough points, then so does every
induced topos $E_{/X}$, $X\in\ob E$. Conversely, if
$(X_i)_{i\in I}$ is a family of objects of $E$ covering the terminal
object and, for every $i\in I$, $E_{/X_i}$ has enough points, then
$E$ has enough points.
\end{enonce*}

For the first assertion, let $f:X'\rightarrow X''$ be a morphism in
$E_{/X}$ which is not an isomorphism. We prove that there exists a
point $q$ of $E_{/X}$ such that $f_q$ is not an isomorphism. Since
$E$ has enough points, there is a point $p$ of $E$ such that
$f_p:X'_p\rightarrow X''_p$ is not an isomorphism. This implies that
there exists $q\in X_p$ such that the map induced by $f_p$ on the
fibers $X'_q$ and $X''_q$ of $X'_p$ and $X''_p$ over $X_p$ at $q$
is not an isomorphism. But by the equivalence
\eqhyperref[IV.eq6.7.2]{6.7.2}, $q$ may be identified with a point of
$E_{/X}$, and the map $X'_q\rightarrow X''_q$ just considered is
none other than the map on fibers at $q$ induced by
$f:X'\rightarrow X''$, which proves the assertion.

Conversely, suppose that the $X_i$ cover the terminal object of $E$
and that the $E_{/X_i}$ have enough points. We prove the same for
$E$. Let\pageoriginale $f:X'\rightarrow X''$ be a morphism in $E$
which is not an isomorphism. There is then some $i\in I$ such that
$f_i:X'_i\rightarrow X''_i$ is not an isomorphism, where the
subscript $i$ denotes restriction to $X_i$. Hence there is a point
$p_i$ of $E_{/X_i}$ such that
$f_{ip_i}:X'_{ip_i}\rightarrow X''_{ip_i}$ is not an isomorphism.
If $p$ denotes the image of $p_i$ in $E$ under the localization
morphism $E_{/X_i}\rightarrow E$, this says that
$X'_p\rightarrow X''_p$ is not an isomorphism, which proves that $E$
has enough points.

\begin{enonce*}[remark]{Remark 6.7.4}\etiquette[6.7.4]{IV.rem6.7.4}
The preceding argument shows that if $(p_{\alpha})$ is a conservative
family of points of $E$, then the family of points of $E_{/X}$ lying
over one of the $p_{\alpha}$ (a family indexed by the disjoint union
of the $X_{p_{\alpha}}$) is conservative. Likewise, if the $X_i$
cover the terminal object of $E$ and, for every $i\in I$, $P_i$ is a
conservative set of points of $E_{/X_i}$, then the set of points of
$E$ which are the \emph{images} of the $p\in P_i$ under the
localization morphisms $E_{/X_i}\rightarrow E$ (with $i\in I$ and
$p$ varying) is conservative.
\end{enonce*}

\setcounter{subsection}{7}

% END INCLUDED SGA4_English_current_buildable_source_20260906/source/04/drafts/SGA4_IV_LINES_3632_3863_English.texfrag


% BEGIN INCLUDED SGA4_English_current_buildable_source_20260906/source/04/drafts/SGA4_IV_LINES_3864_4120_English.texfrag
\subsection{}\etiquette[6.8]{IV.subsec6.8}
Let $E$ be a topos and $p$ a point of $E$. A \emph{neighborhood
of the point $p$ of the topos $E$} is a pair $(X,u)$, where $X\in \ob E$ and
$u\in X_p$. By \eqhyperref[IV.eq6.7.2]{6.7.2}, one may also interpret $u$
as a lift of $p$ to a point of the induced topos $E_{/X}$.
Writing $\varphi_p$ for the fiber functor associated with the point $p$, one
may also interpret a neighborhood of the point $p$ as an object of
the category $E_{/\varphi_p}$ (the full subcategory of
$(E^\circ)^{\widehat{\phantom{,}}}_{/\varphi_p}$ formed by the objects
whose source lies in $E\subset \widehat{E}$). This latter
category, \ie the category of pairs $(X,u)$ as above,
will naturally be called the \emph{category of neighborhoods} of the
point $p$ of the topos $E$; it is also denoted by $\Vvois(p)$. Since
$\varphi_p$\pageoriginale is left exact and finite projective limits
are representable in $E$, finite projective limits are representable
in the category $\Vvois(p)$, and the canonical functor
$\Vvois(p)\rightarrow E$ commutes with them
(\hyperrefext[I][I.sec8.10.10]{8.10.10}). A fortiori, the opposite category
$\Vvois(p)^{\circ}$ is filtered. Moreover, it is clear
(\hyperrefext[I][I.prop3.4]{3.4}) that there is a canonical isomorphism, functorial in
$F\in \ob E$,
\begin{equation*}
\varphi_p(F)=F_p\simeq
\varinjlim_{\Vvois(p)^{\circ}}F(X){\quoi}.\tag{6.8.1}
\end{equation*}\etiquette[6.8.1]{IV.eq6.8.1}
In other words, \emph{the functor $\varphi_p$, ``fiber at the
point $p$,'' is isomorphic to the filtered inductive limit of the
functors represented in the topos $E$ by the neighborhoods of the point
$p$ of $E$.} In fact, \emph{the fiber functor is
pro-representable} (\hyperrefext[I][I.sec8]{8}), which also means that it is
represented by a pro-object $(X_i)_{i\in I}$ of $E$, where $I$ is
a filtered ordered set such that $I\in𝒰$. Indeed, by
\loccit, this is equivalent to saying that the category
$\Vvois(p)$ admits a small full cofinal subcategory. Now let
$C$ be a small full generating subcategory of $E$; I claim that
the full subcategory $C_{/\varphi_{p}}$ of
$\Vvois(p)=E_{/\varphi_p}$, formed by the neighborhoods $(X,u)$ of $p$
such that $X\in \ob C$ (a category that is obviously small), is
cofinal in $\Vvois(p)$. Indeed, let $(X,u)$ be a neighborhood of
$p$; one must find a neighborhood $(X',u')$ of $p$ lying over
$(X,u)$, with $X'\in \ob C$. There is a covering family
$X'_i\rightarrow X$, with the $X'_i$ in $C$, from which it follows that the
$X'_{ip}$ cover $X_p$; hence there exist an $X'=X'_{i_{\circ}}$ and
a $u'\in X'_p$ such that $(X',u')$ lies over $(X,u)$, as
claimed.

\setcounter{subsubsection}{1}
\subsubsection{}\etiquette[6.8.2]{IV.subsubsec6.8.2}
Let\pageoriginale $C$ be a $𝒰$-site and $p$ a point of
$C$, \ie a point of the topos $C^{\sim}$. By an abuse of
notation, we again denote by $\varphi_p$ the ``restriction'' of the fiber functor
$\varphi_p$ to $C$, \ie the composite $\varphi_p\circ \epsilon$; for
an object $X$ of $C$, we shall also often write
$$
X_p=\varphi_p(X)=\epsilon(X)_p{\quoi}.
$$
A \emph{neighborhood of the point $p$ in the site $C$} is a pair
$(X,u)$, where $X\in \ob C$ and $u\in X_p=\varphi_p(X)$. These neighborhoods
again form a category $C_{/\varphi_p}$, which may be denoted by
$\Vvois_C(p)$. Let us show that \emph{the category $\Vvois_C(p)^{\circ}$
is again filtered, that it admits a $𝒰$-small full cofinal
subcategory, and that there is an isomorphism, functorial
in the sheaf $F$ on $C$:}
\begin{equation*}
F_p\simeq \varinjlim_{\Vvois_C(p)^{\circ}} F(X){\quoi}.\tag{6.8.3}
\end{equation*}\etiquette[6.8.3]{IV.eq6.8.3}
First note, by repeating the argument of \eqhyperref[IV.eq6.8.1]{6.8.1}, that if
$C'$ is a full subcategory of $C$ that is topologically
generating (II \hyperref[II.def3.0.1]{3.0.1}), then $\Vvois_{C'}(p)^{\circ}$ is cofinal
in $\Vvois_C(p)^{\circ}$. Taking $C'\in𝒰$ to be small, we are
thus reduced, in proving the assertion, to the case where $C\in𝒰$.
In this case $\widehat{C}$ is a $𝒰$-topos, and there is an
associated-sheaf functor $\underline{a}:\widehat{C}\rightarrow
C^{\sim}$, which allows us to construct on $\widehat{C}$ the composite
fiber functor $\varphi_p\underline{a}$. Applying \eqhyperref[IV.eq6.8.1]{6.8.1} to the
topos $\widehat{C}$ and to its full generating
subcategory $C$, one finds that the category
$C_{/\varphi_p\underline{a}}^{\circ},$ which is manifestly isomorphic to
$\Vvois_C(p)^{\circ}$, is filtered, and that there is an isomorphism
functorial in the presheaf $P$ on $C$:
$$
\varphi_p\underline{a}(P)\simeq \varinjlim_X P(X){\quoi}.
$$
If $F$ is a sheaf on $C$, then $F=\underline{a}F$, and applying
the preceding isomorphism\pageoriginale to $F$ regarded as a
presheaf gives \eqhyperref[IV.eq6.8.3]{6.8.3}. At the same time, this argument
establishes the more general formula
\begin{equation*}
(\underline{a}P)_p\simeq \varinjlim_{\Vvois_C(p)^{\circ}}
P(X){\quoi},\tag{6.8.4}
\end{equation*}\etiquette[6.8.4]{IV.eq6.8.4}
an isomorphism functorial in an arbitrary presheaf $P$ on $C$, at
least when $C$ is small. The general case follows by
introducing a universe $𝒱$ such that $C\in𝒱$, by
regarding $\varphi_p$ on $C$ as a fiber functor with values
in $(𝒱\text{-}\Ens)$, and using the compatibility of the
formation of the associated sheaf with enlargement of universes
(\hyperrefext[II][II.prop3.6]{3.6}).

\setcounter{subsubsection}{4}
\subsubsection{}\etiquette[6.8.5]{IV.subsubsec6.8.5}
To every point $p$ of the topos $C^{\sim}$ we have associated a pro-object of
the $𝒰$-site $C$, defined by the canonical functor
$\Vvois_C(p)\rightarrow C$, taking into account the fact that the category of
neighborhoods of $p$ in $C$ is cofiltered and admits a small
full cofinal subcategory. One immediately verifies that, as $p$
varies, this gives a functor
\begin{equation*}
\Ppoint(C^{\sim})\longrightarrow \Pro(C){\quoi}.\tag{6.8.5.1}
\end{equation*}\etiquette[6.8.5.1]{IV.eq6.8.5.1}
This functor is \emph{fully faithful.} Via the associated
fiber functors, it may be interpreted as a functor
$\Fib(C)\rightarrow \Fib(C^{\sim})\rightarrow \Pro(C)^{\circ}$. The
first functor is an equivalence by \hyperref[IV.prop4.9.4]{4.9.4}
(recalled for fiber functors in \hyperref[IV.subsec6.3]{6.3}), and the
second is fully faithful by the general formalism
(\hyperrefext[I][I.sec8]{8}) of pro-objects.

\subsubsection{}\etiquette[6.8.6]{IV.subsubsec6.8.6}
Consider the special case where $C$ is $𝒰$-small and endowed with
the chaotic topology, so that $C^{\sim}=\widehat{C}$. I claim that
in this case the preceding functor is even an equivalence of
categories
\begin{equation*}
\Ppoint(\widehat{C})\xrightarrow{\approx}\Pro(C){\quoi}.\tag{6.8.6.1}
\end{equation*}\etiquette[6.8.6.1]{IV.eq6.8.6.1}
Indeed,\pageoriginale it remains to prove that this functor is
essentially surjective, which follows from the fact that the functors
of the form $P\mapsto P(X)$ on $\widehat{C}$ are fiber functors,
and from the evident fact that every filtered inductive limit of
fiber functors on a topos is again a fiber functor.

Identifying $C$ with a full subcategory of $\Pro(C)$
(\hyperrefext[I][I.sec8]{8}), the functor
$$
C\longrightarrow \Ppoint(\widehat{C})
$$
induced by a quasi-inverse to \sisi{(6.8.5.2)}{\eqhyperref[IV.eq6.8.6.1]{6.8.6.1}} is manifestly
isomorphic to the canonical functor already considered in
\eqhyperref[IV.eq4.6.2.2]{4.6.2.2};
the quasi-inverse in question may be regarded as its canonical
extension (\hyperrefext[I][I.sec8]{8}), taking into account the fact that
$\Ppoint(\widehat{C})$ is stable under small projective limits.

\subsubsection{}\etiquette[6.8.7]{IV.subsubsec6.8.7}
More generally, let $(X_i)_{i\in I}$ be a pro-object of the
$𝒰$-site $C$, and hence a functor on $C^{\sim}$
\begin{equation*}
\varphi(F)=\varinjlim_I F(X_i){\quoi},\tag{6.8.7.1}
\end{equation*}\etiquette[6.8.7.1]{IV.eq6.8.7.1}
and, in view of (\hyperref[IV.subsubsec6.8.5]{6.8.5}), it is natural to ask
when this functor is a fiber functor. The following
conditions are equivalent:
\begin{enumerate}
\renewcommand{\theenumi}{\roman{enumi}}
\renewcommand{\labelenumi}{(\theenumi)}
\item The preceding functor $\varphi$ is a fiber functor on
$C^{\sim}$.

\item The ``restriction'' of $\varphi$ to $C$ is a fiber functor
on the site $C$ (\hyperref[IV.subsec6.3]{6.3}).

\item For every covering family $Y_{\alpha}\rightarrow Y$ of an
object $Y$ of $C$, every $i_{\circ}\in I$, and every morphism
$X_{i_{\circ}}\rightarrow Y$, there exist an $i\geq i_{\circ}$, an
$\alpha$, and a morphism $X_i\rightarrow Y_{\alpha}$ making
the diagram
\[
\xymatrix@C=2cm@R=1.25cm{
X_i \ar[d] \ar[r]& Y_{\alpha} \ar[d]\\
X_{{i}_{o}} \ar[r]     &  Y \quad . }
\]
commutative.

Indeed,\pageoriginale condition (iii) simply expresses the fact
that $\varphi| C$ carries covering families to covering
families. The implications (i) $\Rightarrow$ (ii) $\Rightarrow$
(iii) are trivial, and it remains to prove (iii) $\Rightarrow$ (i).
Since the functor $\varphi$ is manifestly left exact, it
remains to prove (\hyperref[IV.prop4.9.4]{4.9.4}) that it carries covering
families $F_{\alpha}\rightarrow F$ in $C^{\sim}$ to covering
families, \ie that for every $i_{\circ}$ and every $x_{\circ}\in
F(X_{i_{\circ}})$, there exist an $i\geq i_{\circ}$, an $\alpha$, and a
$y\in F_{\alpha}(X_i)$ such that $x_{\circ}$ and $y$ have the same image
in $F(X_i)$. Since the family $F_{\alpha}\rightarrow F$ is covering, there
is a covering family $Z_{\gamma}\rightarrow X_{i_{\circ}}$
of $X_{i_{\circ}}$ such that, for every $\gamma$, the inverse image of
$x_{\circ}$ in $F(Z_{\gamma})$ lifts to an element of some
$F_{\alpha}(Z_{\gamma})$. By hypothesis (iii), there exist an $i\geq
i_{\circ}$ and an $X_{i_{\circ}}$-morphism $X_i\rightarrow
Z_{\gamma}$. The image of $x_{\circ}$ in $F(X_i)$ therefore lifts
to an element $y$ of some $F_{\alpha}(X_i)$, \cqfd
\end{enumerate}

\begin{enonce*}[remark]{Exercise 6.9}\etiquette[6.9]{IV.exer6.9}
Let $C\in𝒰$ be a site, and let $\varphi$ be a fiber functor on
$C$, pro-represented by an object $(X_i)_{i\in I}$ of $\Pro(C)$. For
every index $i\in I$, every object $Y$ of $C$, every morphism
$u:X_i\rightarrow Y$, and every covering sieve $R$ on $Y$, choose
an index $j\geq i$ such that the composite $X_j\rightarrow
X_i\xrightarrow{u}Y$ factors through $R$. For fixed $i$, let
$J(i)$ be the set formed by $i$ and by the $j$ obtained as $Y$, $u$,
and $R$ vary; for a subset $I'$ of $I$, likewise let $J(I')$ be the
union of the $J(i)$ for $i\in I'$. On the other hand, if $I'$ is a
finite subset of $I$, let $m(I')\in I$ be an upper bound of $I'$, and for
an arbitrary subset $I'$ of $I$, let $M(I')$ be the set of all $m(I'')$,
where $I''$ ranges over the set of finite subsets
of\pageoriginale $I'$. We assume that the function
$I'\mapsto m(I')$ is chosen so that, when $I'$ consists of a single element
$i$, one has $m(I')=i$; this implies that, for every subset $I'$ of
$I$, one has $I'\subset M(I')$. Consider the iterates $(MJ)^n$ of
the map $MJ$ from the set of subsets of $I$ to itself,
and, for every subset $I'$ of $I$, let $P(I')$ be the union of the
$(MJ)^n(I')$.
\begin{enumerate}
\renewcommand{\theenumi}{\alph{enumi}}
\renewcommand{\labelenumi}{\theenumi)}
\item Show that, for every subset $I'$ of $I$, $P(I')$ is a
filtered subset of $I$ for the induced order, and that $I$ is the
increasing filtered union of the $P(I')$ as $I'$ ranges over the set
of \emph{finite} subsets of $I$.

\item Show that there is a cardinal $c\in 𝒰$, depending
only on the cardinal $\card \text{fl}(C)=a$, such that, for every
finite subset $I'$ of $I$, one has $\card(P(I'))\leq c$. (If $a$ is infinite,
one may take $c=2^{a}$.)

\item Show, using \hyperref[IV.subsubsec6.8.7]{6.8.7}, that for every subset $I'$ of $I$
the functor $\varphi_{I'}$ on $C$, pro-represented by the pro-object
$(X_i)_{i\in P(I')}$, is a fiber functor on $C$, and that the fiber
functor $\varphi$ is the filtered inductive limit of the functors
$\varphi_{I'}$ as $I'$ ranges over the set of finite subsets
of $I$.

\item Deduce that there is a small full subcategory $\Phi$
of $\Ffib(C)\approxeq \Ffib(C^{\sim})$ such that every object of $\Ffib(C)$
is a filtered inductive limit of objects of $\Phi$. (If $a$ is
infinite, one may take $\Phi$ to have cardinal $\leq 2^{2^{a}}$; if $a$
is finite, one may of course take $\Phi=C$.)
\end{enumerate}
\end{enonce*}

\begin{enonce*}[remark]{Exercise 6.10}\etiquette[6.10]{IV.exer6.10}
Let $C$ be a $\mathcal{U}$-site, and let $D$ be a
$\sheaf{\mathcal{U}}$-category in which small filtered $\varinjlim$
are representable; let $\Fais(C,D)$ be the category of
sheaves on $C$ with values in $D$ (\hyperrefext[II][II.def6.1]{6.1}). By
extending \eqhyperref[IV.eq6.8.3]{6.8.3}, define a ``fiber functor''
\begin{equation*}
F\longrightarrow F_p:\Fais(C,D)\longrightarrow D{\quoi}.\tag{6.10.1}
\end{equation*}\etiquette[6.10.1]{IV.eq6.10.1}
If finite $\varprojlim$ are representable in $D$ and
commute with filtered $\varinjlim$, then the preceding functor
is left exact.
\end{enonce*}

% END INCLUDED SGA4_English_current_buildable_source_20260906/source/04/drafts/SGA4_IV_LINES_3864_4120_English.texfrag


% BEGIN INCLUDED SGA4_English_current_buildable_source_20260906/source/04/drafts/SGA4_IV_LINES_4121_4475_English.texfrag
\section{Examples of fiber functors and points of topoi}\etiquette[7]{IV.sec7}\pageoriginale

\subsection{The case of $\Top(X)$ for a topological space
  $X$}\etiquette[7.1]{IV.subsec7.1}

Let $x$ be a point of $X$. Regard the singleton $\{x\}$ as a
topological space, and consider the inclusion
$$
i_x:\{x\}\longrightarrow X{\quoi}.
$$
By \hyperref[IV.subsubsec4.1.1]{4.1.1}, it defines a morphism of topoi
(defined up to unique isomorphism)
\begin{equation*}
\Top(i_x):\Top(\{x\})\longrightarrow \Top(X){\quoi}.\tag{7.1.1}
\end{equation*}\etiquette[7.1.1]{IV.eq7.1.1}
On the other hand, one may identify $P$ with $\Top(\{x\})$
(\hyperref[IV.subsec2.2]{2.2}), and one thereby obtains a point $p_x$ of
$\Top(X)$:
\begin{equation*}
p_x:P\longrightarrow \Top(X){\quoi},\tag{7.1.2}
\end{equation*}\etiquette[7.1.2]{IV.eq7.1.2}
defined by $x$ up to unique isomorphism. The associated fiber functor
is given by the familiar formula $[TF]$, which moreover follows
immediately from \hyperrefext[I][I.prop5.1]{5.1}:
\begin{equation*}
p^{\ast}_x(F)=F_x=F_{p_x}=\varinjlim_{U\ni x}F(U){\quoi},\tag{7.1.3}
\end{equation*}\etiquette[7.1.3]{IV.eq7.1.3}
the limit being taken over the open neighborhoods $U$ of $x$ in
$X$.

Since $\Top(X)$ is known to be isomorphic to $\Top(X_{\sob})$
(\hyperref[IV.subsubsec4.2.1]{4.2.1}), one sees more generally that every point
of $X_{\sob}$, \ie every irreducible closed subset $Z$ of $X$,
defines a point of $\Top(X)$, or equivalently a fiber functor on
$\Top(X)$, which we shall still denote by $F\mapsto F_Z$. Returning to
its definition by formula \eqhyperref[IV.eq7.1.3]{7.1.3} on $X_{\sob}$, one obtains
\begin{equation*}
F_Z=\varinjlim_{U\cap Z\neq \phi} F(U){\quoi},\tag{7.1.4}
\end{equation*}\etiquette[7.1.4]{IV.eq7.1.4}
the\pageoriginale inductive limit being taken over the open subsets
$U$ of $X$ which meet $Z$.

\setcounter{subsubsection}{4}
\subsubsection{}\etiquette[7.1.5]{IV.subsubsec7.1.5}
It is well known that the fiber functors $F\mapsto F_x$ $(x\in X)$
already form a conservative family. This is especially clear from
the interpretation of sheaves on $X$ in terms of étale spaces, since
the fiber of $F$ at $x$ is then simply its fiber in the sense of
spaces fibered over $X$. Note that the index set of the conservative
family of fiber functors under consideration is $X$, and hence is
$𝒰$-small.

\subsubsection{}\etiquette[7.1.6]{IV.subsubsec7.1.6}
By \hyperref[IV.subsubsec4.2.3]{4.2.3}, the category $\Ppoint(\Top(X))$
is equivalent to the category associated with the set $X_{\sob}$
ordered by the specialization relation. In other words, every fiber
functor on $\Top(X)$ is isomorphic to a fiber functor
$F\mapsto F_Z$, where $Z$ is a uniquely determined element of
$X_{\sob}$, \ie a uniquely determined irreducible closed subset
of $X$; moreover, if $Z$, $Z'\in X_{\sob}$, then the set
$\Hom(F_Z,F_{Z'})$ is either empty or consists of a single point, the
latter case occurring if and only if $Z$ is a specialization of $Z'$
in $X_{\sob}$, \ie if and only if $Z\subset Z'$ as a subset of
$X$. When $X$ itself is sober, one may replace $X_{\sob}$ by $X$
itself in these statements.

Note that the automorphism group of a fiber functor of $\Top(X)$
(opposite to the automorphism group of the corresponding point of
$\Top(X)$) is always the trivial group (more generally,
\hyperref[IV.subsubsec4.2.3]{4.2.3} tells us the same thing about the
automorphism group of every morphism of topoi associated with
topological spaces). This is a phenomenon very special to the case
under consideration: see the example in \hyperref[IV.subsec7.2]{7.2}, as\pageoriginale
well as VIII 7. In the case of topoi associated with ``moduli
problems'' (\cf for example [10]), the automorphism groups of fiber
functors moreover have a remarkable interpretation as the
automorphism groups of the algebraic structures (over algebraically
closed fields) which one proposes to classify.

\subsubsection{}\etiquette[7.1.7]{IV.subsubsec7.1.7}
We have just seen how one can reconstruct a sober space $X$ (or the
sober space $X_{\sob}$ associated with an arbitrary topological
space), at least as a set ordered by the specialization relation,
from the topos $\Top(X)$ (which it even suffices to know up to
equivalence), as the set of isomorphism classes of ``points'' of
$\Top(X)$. From what was said in \hyperref[IV.subsec2.1]{2.1}, one also
reconstructs the topology of $X$, \ie the family of its open subsets,
as follows: for every subobject $U$ of the terminal object $e_E$ of
$E$, let $\Pt(U)$ be the set of $x\in X$ such that $U_x\neq\phi$
(which moreover identifies, by \eqhyperref[IV.eq6.7.2]{6.7.2}, with the set of
isomorphism classes of points of the induced topos $E_{/U}$). Then
$U\mapsto \Pt(U)$ is an isomorphism of ordered sets from the set of
subobjects of $e_E$ onto the set of open subsets of $X$.

\subsubsection{}\etiquette[7.1.8]{IV.subsubsec7.1.8}
The determination in \hyperref[IV.subsubsec7.1.6]{7.1.6} of the category of
points of the topos $\Top(X)$ defined by a topological space $X$
leads to the following terminology for points $p$, $p'$ of an
arbitrary topos $E$: one says that $p$ is a \emph{specialization}
of $p'$, or that $p'$ is a \emph{generization} of $p$, when there
exists a morphism from $p'$ to $p$ (in the sense of the category
$\Ppoint(E)$ of \hyperref[IV.def6.1]{6.1}). These relations are still
transitive, but, using \hyperref[IV.subsubsec6.8.6]{6.8.6}, one easily finds
examples in which $p$ and $p'$ are each a specialization of\pageoriginale the
other, without $p$ and $p'$ being isomorphic (much less equal). The
category of generizations of a point $p$ of the topos $E$, by which
we mean the category $\Ppoint(E)_{/p}$, plays in some respects a role
analogous to that of passing to the localization of a scheme at a
point. This role will be made precise in
\sisi{Chap. VI}{\hyperrefext[VI][VI.subsec8.4]{8.4}} by the construction of the topos obtained by
localizing $E$ at the point $p$, whose category of points is
canonically equivalent to the category of generizations of $p$.

\begin{enonce*}[remark]{Exercise 7.1.9}\etiquette[7.1.9]{IV.exer7.1.9}
Let $E$ be a $𝒰$-topos. Prove that $E$ is equivalent to a
topos of the form $\Top(X)$, where $X$ is a topological space $\in
𝒰$, if and only if it satisfies the following two conditions:
\begin{enumerate}
\renewcommand{\theenumi}{\alph{enumi}}
\renewcommand{\labelenumi}{\theenumi)}
\item The family of subobjects of the terminal object $e_E$ is
generating for the topos $E$.

\item $E$ has enough points (\hyperref[IV.prop6.5]{6.5}). (This condition is not
superfluous: \cf \hyperref[IV.subsec7.4]{7.4}.)
\end{enumerate}

Compare with \hyperref[IV.exer7.8]{7.8} c).
\end{enonce*}

\begin{enonce*}[remark]{Exercise 7.1.10}\etiquette[7.1.10]{IV.exer7.1.10}
Let $X$ be a topological space endowed with an operator group $G$
$(X,G\in𝒰)$, and let $E=\Top(X,G)$ (\hyperref[IV.subsec2.3]{2.3}).
\begin{enumerate}
\renewcommand{\theenumi}{\alph{enumi}}
\renewcommand{\labelenumi}{\theenumi)}
\item Show that, for every object $(X',G)$ of $E$, the induced topos
$E_{/(X',G)}$ is canonically equivalent to the topos $\Top(X',G)$
(compare \hyperref[IV.subsec5.7]{5.7}).

\item When $G$ acts properly and freely on $X'$ (Bourbaki,
Gen. Top. Chap. \hyperrefext[III][III.sec4]{4}), show that the morphism of
spaces with operators $(X',G)\rightarrow (X'/G,e)$ induces
(\hyperref[IV.subsubsec4.1.2]{4.1.2}) an equivalence of topoi
$$
\Top(X',G)\longrightarrow \Top(X'/G){\quoi}.
$$

\item Deduce\pageoriginale from b) that there exists an object $Z$ of
$E$ such that the induced topos $E_{/Z}$ is equivalent to $\Top(X)$,
the localization functor $E\rightarrow E_{/Z}$ being isomorphic to
the functor that ``forgets the actions of $G$''. (Imitate the
argument of \hyperref[IV.subsubsec5.8.3]{5.8.3}.)

\item Deduce from c) that every fiber functor on $E$ is induced, via
the functor that ``forgets the actions of $G$'', by a fiber functor
on $\Top(X)$, and can therefore be defined by a point of $X_{\sob}$.

\item Determine the structure of the category $\Ppoint(E)$ (up to
equivalence) in terms of the space with operators $(X_{\sob},G)$.
Conclude in particular that two points of $X_{\sob}$ define
isomorphic fiber functors on $E$ if and only if they are conjugate
under the action of $G$.
\end{enumerate}
\end{enonce*}

\begin{enonce*}[remark]{Exercise 7.1.11}\etiquette[7.1.11]{IV.exer7.1.11}
Let $X\in 𝒰$ be a topological space. Prove that the
$𝒱$-topos $\TOP(X)$ (\hyperref[IV.subsec2.5]{2.5}) has enough points.
More precisely, for every object $X'$ of $\TOP(X)$ and every $x'\in
X'$, define a fiber functor $F\mapsto F_{X',x'}$ on $\TOP(X)$ which
depends only on the ``germ'' of the space $(X',x')$ over $X$, and
prove that this family of fiber functors is conservative (and indexed
by a $𝒱$-small index set). Show, using a suitable filtered projective
system $(X'_i,x'_i)_{i\in I}$, with $I$ not $𝒰$-small, that one does
not obtain in this way all the fiber functors of the $𝒱$-topos
$\TOP(X)$. Show that the set of isomorphism classes of such fiber
functors does not have cardinality $\in 𝒱$.
\end{enonce*}

\subsection{Points of a classifying topos
  $B_G$}\etiquette[7.2]{IV.subsec7.2}\pageoriginale

Let $E$ be a $𝒰$-topos and $G$ a group of $E$, whence a classifying
topos $B_G$ (\hyperref[IV.subsec2.4]{2.4}). If $e$ is the unit group, the
morphisms $e\rightarrow G\rightarrow e$ define
(\hyperref[IV.subsec4.5]{4.5}) morphisms of topoi (taking into account that
$B_e\approx E$)
\begin{equation*}
E\longrightarrow B_G \longrightarrow E{\quoi},\tag{7.2.1}
\end{equation*}\etiquette[7.2.1]{IV.eq7.2.1}
whose composite is isomorphic to $\id_E$; passage to the categories
of points therefore gives functors
\begin{equation*}
\Ppoint(E)\longrightarrow \Ppoint(B_G)\longrightarrow
\Ppoint(E){\quoi},\tag{7.2.2}
\end{equation*}\etiquette[7.2.2]{IV.eq7.2.2}
whose composite is isomorphic to the identity. Use the fact that the
first morphism of topoi in \eqhyperref[IV.eq7.2.1]{7.2.1} identifies with a
localization morphism relative to the object $E_G=G$ of $B_G$
(\hyperref[IV.subsec5.8]{5.8}); by \eqhyperref[IV.eq6.7.2]{6.7.2}, it follows that the
first functor in \eqhyperref[IV.eq7.2.2]{7.2.2} identifies with the natural
``inclusion'' functor
\begin{equation*}
\Ppoint(B_G)_{/\widehat{E}_G}\longrightarrow
\Ppoint(B_G){\quoi},\tag{7.2.3}
\end{equation*}\etiquette[7.2.3]{IV.eq7.2.3}
where $\widehat{E}_G$ denotes the presheaf $p'\mapsto (E_G)_{p'}$ on
$\Ppoint(B_G)$. Since $E_G$ plainly covers the terminal object of
$B_G$, its fibers $(E_G)_p$ are nonempty, from which it follows that
\eqhyperref[IV.eq7.2.3]{7.2.3} is essentially surjective; this means that
\emph{every fiber functor on $B_G$ is induced (up to isomorphism) by
a fiber functor on $E$, via the functor that ``forgets the actions of
$G$'' $B_G\rightarrow E$.} One concludes in particular that
\emph{if $E$ has enough points (\resp a small conservative family of
points),} then the same is true of $B_G$.

\begin{enonce*}[remark]{Remark 7.2.4}\etiquette[7.2.4]{IV.rem7.2.4}
One can go further and determine (up to equivalence) the structure of
the category $C'=\Ppoint(B_G)$ in terms of the category
$C=\Ppoint(E)$ and the presheaf of groups
$$
\widehat{G}:p\longmapsto G_p: C^{\circ}\longrightarrow
(\text{Groups})
$$
on\pageoriginale it. Indeed, define a new category $C_1$, having the
same objects as $C$, and such that for two objects $p$, $q$ of $C$,
one has
$$
\Hom_{C_1}(p,q)=\Hom_C(p,q)\times \widehat{G}(p){\quoi},
$$
composition of arrows in $C_1$ being induced by that in $C$ and by
the group law of the presheaf $\widehat{G}$. The datum of a functor
$\varphi_1:C_1\rightarrow D$ into an arbitrary category $D$ is then
equivalent to the datum of a functor $\varphi:C\rightarrow D$
endowed with an action of the presheaf $\widehat{G}$ on the latter
(\ie the datum, for every $p\in \ob C$, of an action of
$\widehat{G}(p)$ on $\varphi(p)$, satisfying an evident functoriality
condition as $p$ varies). Using this observation, one defines a
canonical functor
\begin{equation*}
\varphi_1:C_1\longrightarrow C'{\quoi},\tag{7.2.4.1}
\end{equation*}\etiquette[7.2.4.1]{IV.eq7.2.4.1}
corresponding to the functor $\varphi:C\rightarrow C'$ of
\eqhyperref[IV.eq7.2.2]{7.2.2}, which associates with every point $p$ of $E$
the point $\varphi(p)$ induced on $B_G$, together with the natural
actions of $\widehat{G}(p)=G_p$ on the latter, arising from the
actions of $G_p$ on the $X_p$ as $X$ ranges over $B_G$. We leave to
the reader the proof that \eqhyperref[IV.eq7.2.4.1]{7.2.4.1} is an equivalence
of categories, using the structure \eqhyperref[IV.eq7.2.3]{7.2.3} of the first
functor in \eqhyperref[IV.eq7.2.2]{7.2.2}, and noting that the natural
inclusion functor $C\rightarrow C_1$ has an analogous structure,
relative to the inverse image $P_1$ of the presheaf
$P'=\widehat{E}_G$ on $C'$.
\end{enonce*}

\setcounter{subsubsection}{4}
\subsubsection{}\etiquette[7.2.5]{IV.subsubsec7.2.5}
One\pageoriginale concludes in particular that the \emph{isomorphism
classes of points of $B_G$ correspond exactly to the isomorphism
classes of points of $E$: every point of the classifying topos $B_G$
is induced, up to a} non-unique \emph{isomorphism, by a point of
$E$}. In particular, when $E$ is the punctual topos $P$, and hence
$G$ is an ordinary group, the category $\Ppoint(B_G)$ is a connected
groupoid with fundamental group $G$: every fiber functor on $B_G$ is
isomorphic (noncanonically) to the functor $\omega$ that ``forgets
the actions of $G$'', and the monoid of endomorphisms of the latter
functor is the group $G$ (hence every endomorphism of $\omega$ is an
automorphism).

\begin{enonce*}[remark]{Exercise 7.2.6}\etiquette[7.2.6]{IV.exer7.2.6}
\begin{enumerate}
\renewcommand{\theenumi}{\alph{enumi}}
\renewcommand{\labelenumi}{\theenumi)}
\item Let (Tors) be the category of pairs $(\Gamma,P)\in 𝒰$,
where $\Gamma$ is a group and $P$ is a right torsor under $\Gamma$.
The functor $(\Gamma,P)\mapsto \Gamma$
\begin{equation*}
(\Tors)\longrightarrow (\Groupes)\tag{7.2.6.1}
\end{equation*}\etiquette[7.2.6.1]{IV.eq7.2.6.1}
is a cofibered functor. With the notation of \hyperref[IV.rem7.2.4]{7.2.4},
consider the functor
$$
\widehat{G}:C^{\circ}\longrightarrow (\Groupes){\quoi},
$$
and let $D'$ be the cofibered category over $C^{\circ}$ obtained by
inverse image from the cofibered category \eqhyperref[IV.eq7.2.6.1]{7.2.6.1}, and
let $D$ be the corresponding category fibered over $C$, having the
same fiber categories as $D'$, so that for $p\in \ob C$ one has
$$
D_p=\mbox{ category of right torsors under $G_p$}{\quoi}.
$$
Construct canonical functors
\begin{equation*}
D\longrightarrow C'=\Ppoint(B_G)\quav C'\longrightarrow
D{\quoi},\tag{7.2.6.2}
\end{equation*}\etiquette[7.2.6.2]{IV.eq7.2.6.2}
\emph{quasi-inverse to one another.} In particular, conclude that
the category of points of $B_G$ above a given point $p$ of $E$ is
canonically equivalent to the category of right torsors under the
group $G_p$.

\item Let\pageoriginale $𝒢=(G_i)_{i\in I}$ be a strict
projective system of groups of the $𝒰$-topos $E$ ($I\in
𝒰$ a filtered preordered set), whence a classifying topos
$B_{𝒢}$, by imitating the definition of
\hyperref[IV.subsubsec2.7.1]{2.7.1}. Determine the category of points of
$B_{𝒢}$ by imitating a). (NB. Replace the categories
$(\Tors)$ and $(\Groupes)$ by the categories of projective systems
indexed by $I$ in these categories.) Conclude that if $I$ admits a
countable cofinal subset, then every fiber functor on $B_{𝒢}$ is
isomorphic to a functor induced by a fiber functor on $E$ (via the
functor that ``forgets the actions of $G$''), the latter being
determined up to non-unique isomorphism.

\item Taking $E$ to be the punctual topos, so that $𝒢$ is an
ordinary pro-group, show that the conclusion of b) is likewise valid
if $I$ is arbitrary but $𝒢$ is profinite (\ie the $G_i$ are
finite): every fiber functor is isomorphic to the functor that
``forgets the actions of $G$''.

\item Still taking the case where $E$ is the punctual topos, give an
example of a fiber functor on $B_{𝒢}$, for a suitable projective
system $𝒢=(G_i)_{i\in I}$, which is not isomorphic to the functor
that ``forgets the actions of $G$''.

\item Regard $𝒢=(G_i)$ as a group of the topos
$\widehat{I}$, and let $𝒯$ be a \emph{gerbe} on $I$ with
\emph{band} $G$ [\hyperref[IV.sec3]{3}]. Show how one can
``twist'' the classifying topos $B_G$ by means of the gerbe $T$, to
obtain a topos $B^T_{𝒢}$, an inductive limit of subcategories
$B^T(i)$ equivalent (noncanonically) to the classifying topoi
$B_{G_i}$. Show that the topos $B^T_{𝒢}$ admits a fiber functor if
and only if the gerbe $T$ is ``neutral'', by establishing an\pageoriginale
equivalence of categories between the category of fiber functors on
$B^T_{𝒢}$ and the category of sections of $T$ over $I$. Conclude,
if the $G_i$ are commutative, so that gerbes on $I$ with band $𝒢$
are classified by
$\displaystyle{\varprojlim_I}^{(2)}G_i=H^{2}(I,𝒢)$
(\loccit), that there is an example of a topos (not ``empty''
(\hyperref[IV.subsec2.2]{2.2})) of the form $B^T_{𝒢}$ which has no
points. (Take an example in which
$\displaystyle{\varprojlim_I}^{(2)}G_i\neq 0$.)
\end{enumerate}
\end{enonce*}

% END INCLUDED SGA4_English_current_buildable_source_20260906/source/04/drafts/SGA4_IV_LINES_4121_4475_English.texfrag


% BEGIN INCLUDED SGA4_English_current_buildable_source_20260906/source/04/drafts/SGA4_IV_LINES_4476_4830_English.texfrag
\subsection{Points of presheaf topoi $\widehat{C}$; examples of
$𝒰$-presheaf topoi $\widehat{C}$ whose category
of points is not equivalent to a small
category}\etiquette[7.3]{IV.subsec7.3}

Let $E$ be a $𝒰$-topos, and let $(\varphi_i)_{i\in I}$ be a filtered
family of fiber functors on $E$. It follows from the exactness
properties of the functors $\varinjlim$
(\sisi{I\quad and}{\hyperrefext[I][I.prop2.8]{2.8}}) that the functor
$\varphi=\varinjlim\varphi_i$ is also a fiber functor:
\emph{every filtered inductive limit of fiber functors is a fiber
functor.} Consequently, the category $\Ffib(E)$ admits $𝒰$-filtered
inductive limits (and the inclusion functor
$\Ffib(E)\rightarrow \Hhom(E,(𝒰\text{-}\Ens))$ commutes with them);
in other words, the category $\Ppoint(E)$ admits $𝒰$-filtered
projective limits. This fact was already used in Exercise
\sisi{\hyperref[IV.exer7.1.10]{7.1.10}}{\hyperref[IV.exer7.1.11]{7.1.11}}
to give an example of ``large'' categories of fiber functors.

\hypertarget{sga4.exp.IV.subsubsection.7.3.3}{}Considerably simpler examples can be constructed, with presheaf
topoi of the form $E=\widehat{C}$, $C\in𝒰$, by using
\eqhyperref[IV.eq6.8.6.1]{6.8.6.1}. This immediately gives examples
in which $\Ffib(\widehat{C})$ is not equivalent to a category
$\in𝒰$, \ie in which the cardinal of the set of isomorphism classes
of objects of this $𝒰$-category is not $\in𝒰$. This means
that the category of pro-objects $\Pro-(C)$ is not equivalent to a
category $\in 𝒰$. For example, it suffices to take
$D=C^{\circ}$ to be the category of finite sets of the form
$[0,n]$, where $n\geq 0$ is an\pageoriginale integer; in
this case $\Ind(D)\simeq \Pro(C)^{\circ}$ is equivalent to the
category of nonempty sets. The topos $E$ is then the familiar
category of \emph{cosimplicial sets}. One could also take $C$ to be
the category of nonempty finite sets; one then finds that the
category of points of the presheaf topos $\widehat{C}$ of simplicial
sets is equivalent to the category of profinite sets, or equivalently
to the category of totally disconnected compact spaces.

\subsection{Nonempty topoi without points}\etiquette[7.4]{IV.subsec7.4}

The following example is due to $P$. Deligne. (For another example,
\cf \hyperref[IV.exer7.2.6]{7.2.6} e).) Let $K$ be a compact space
equipped with a measure $\mu$, and let $U$ be the ordered set of
measurable subsets of $K$, modulo sets of measure zero. Make $U$
into a category $\underline{U}$ with $\ob \underline{U}=U$, the
morphisms of $\underline{U}$ being the ``inclusion morphisms''
between elements of $U$. Make $\underline{U}$ into a site by taking
the pretopology for which $\Cov(E)$ (for $E\in U$) consists of the
\emph{countable} families of elements $E_i$ of $U$ bounded above by
$E$, such that $E$ is the union of the $E_i$ modulo a set of measure
zero. This gives a topos $\Top(\mu)=\underline{U}^{\sim}$ whose set
of subobjects of the terminal object is a generating family (a topos
which appears to have escaped the attention of probabilists). This
topos is an ``empty topos'' (\hyperref[IV.subsec2.2]{2.2}) if and
only if $\mu=0$. On the other hand, the category of points of this
topos is equivalent to the discrete category defined by the set of
points $x\in K$ such that $\mu(\{x\})\neq 0$ (proof left to the
reader). It is therefore empty if $K$ has no such points, for
example if $K$ is the unit interval in the real line, with the
measure induced by Lebesgue measure.

\begin{enonce*}[remark]{Exercise 7.5}\etiquette[7.5]{IV.exer7.5}
(\emph{\sisi{Karoubian}{karoubian} categories\pageoriginale and
morphisms of topoi})\footnote[1]{Compare
\hyperrefext[I][I.exer8.7.8]{8.7.8}.}
\begin{enumerate}
\renewcommand{\theenumi}{\alph{enumi}}
\renewcommand{\labelenumi}{\theenumi)}
\item Let $C$ be a category, let $X$ be an object of $C$, and let
$p$ be an endomorphism of $X$. We say that $p$ is a
\emph{projector} if $p^2=p$. Prove that, if $p$ is a projector,
$\Ker(\id_X,p)$ is representable if and only if
$\Coker(\id_X,p)$ is representable, and that the two objects of $C$
so obtained are canonically isomorphic. We then say that the
projector $p$ \emph{admits an image,} and
$\Ker(\id_X,p)\simeq \Coker(\id_X,p)$ is called the \emph{image of
the projector $p$}; according to the context, it is identified with
either a subobject or a quotient object of $C$. An object isomorphic
to $\text{Im\,} p$, for a suitable projector $p$ in $X$, is called
a \emph{direct factor of the object $X$ of $C$.} We say that $C$ is
a \emph{category with direct factors}, or a
\emph{\sisi{Karoubian}{karoubian} category,} if every projector in
an object of $C$ admits an image. If $F:C\rightarrow C'$ is a
functor that commutes with kernels or cokernels, then $F$ sends a
projector admitting an image to a projector admitting an
image\nde{No condition on $F$ is needed here: every functor
preserves projectors admitting images.}.

\item Show that, for every category $C$, one can find a functor
$\varphi:C\rightarrow \kar (C)$ from $C$ to a Karoubian category
(determined up to equivalence) such that, for every Karoubian
category $C'$, the functor
$$
f\mapsto f\circ \varphi:\Hhom(\kar(C),C')\longrightarrow \Hhom(C,C')
$$
is an equivalence of categories; $\kar(C)$ is called the
\emph{Karoubi envelope} of the category $C$. (Hint: take
$\ob \kar (C)$ to be the set of pairs $(X,p)$, with $X\in \ob C$
and $p$ a projector in $X$, and take $\Hom((X,p),(Y,q))$ to be the
subset of $\Hom(X,Y)$ consisting of those $f:X\rightarrow Y$ such
that $f=qfp.$) Show that $\varphi$ is fully faithful.

\item Let $F:E\rightarrow E'$ be a functor from one category to
another. Show that if $F$ commutes with a given type of inductive or
projective limits,\pageoriginale then the same is true of every
direct factor of $F$. In particular, if $E$ and $E'$ are
$𝒰$-topoi and $F$ is an inverse-image functor for a morphism of
topoi $E'\rightarrow E$, then the same is true of every direct
factor of $F$; consequently, the category $\Hhomtop(E',E)$ is
Karoubian, and in particular the category $\Ppoint(E)$ is
Karoubian.

\item Show that, for every $𝒰$-category $C$, the category
$\Ind(C)$ is Karoubian (where $\Ind(C)$ is formed using inductive
systems in $C$ indexed by a filtered preordered set $I\in 𝒰$). (If
$C\in 𝒰$, use, for example, the fact
(\hyperref[IV.subsec7.3]{7.3}) that $\Ind(C)$ is equivalent to the
category
$\Ffib((C^{\circ})^{\widehat{\phantom{,}}})=\Ppoint((C^{\circ})^{\widehat{\phantom{,}}})^{\circ}.$)
Deduce another construction of $\kar(C)$ as the full subcategory of
$\Ind(C)$ formed by the images in $\Ind(C)$ of the projectors of
objects $X$ of $C$.
\end{enumerate}
\end{enonce*}

\begin{enonce*}[remark]{Exercise 7.6}\etiquette[7.6]{IV.exer7.6}
(\emph{Essential morphisms of topoi, essential points}).
\begin{enumerate}
\renewcommand{\theenumi}{\alph{enumi}}
\renewcommand{\labelenumi}{\theenumi)}
\item Let $f:E\rightarrow F$ be a morphism of topoi. Show that the
following conditions are equivalent:
\begin{itemize}
\item[i)] $f_{!}$ exists, \ie $f^{\ast}$ admits a left adjoint.

\item[ii)] $f^{\ast}$ commutes with $𝒰$-projective limits.

\item[${\rm ii')}$] $f^{\ast}$ commutes with $𝒰$-products.
\end{itemize}
We then say that $f$ is an \emph{essential morphism} from the topos
$E$ to the topos $F$.

Show that if $f$ satisfies these conditions, then the same is true
of every direct factor of $f$. (Use \hyperref[IV.rem1.8]{1.8} and
\hyperref[IV.exer7.5]{7.5} c).)

\item Let $E$ be a topos. An \emph{essential point of the topos}
$E$ is a point $p:P\rightarrow E$ of $E$ such that $p_{!}$ exists.
Show that if $E=\Top(X)$, where $X$ is a topological space, and if
$p$ is the point of $E$ defined by an $x\in X$, then $p$ is
essential\pageoriginale if and only if $x$ has a smallest open
neighborhood $U$ (which means, if the points of $X$ are closed, that
$x$ is an \emph{isolated point} of $X$). For a morphism of topoi
$f:E\rightarrow F$ to be essential (a)), it is necessary that $f$
send essential points to essential points, and this condition is
also sufficient when $E$ has enough essential points (\ie when the
family of these points is conservative) (\cf d)).

\item A point $p$ of the topos $E$ is essential if and only if the
associated fiber functor is representable by an object $X$ of $E$.
For the covariant functor $\varphi:E\rightarrow (\Ens)$ represented
by a given object $X$ of $E$ to be a fiber functor (\ie to define a
point, necessarily an essential point, of $E$), it is necessary and
sufficient that $X$ be nonempty and connected
(\hyperref[IV.subsubsec4.3.5]{4.3.5}), and projective (\ie that the
functor $\varphi$ send epimorphisms to epimorphisms). (Hint: first
show that $\varphi$ commutes with coproducts if and only if $X$ is
nonempty and connected, and then use the criterion
\hyperref[IV.prop4.9.4]{4.9.4}.) Deduce an equivalence between the
category $\Ppointess(E)$ and the full subcategory $P$ of $E$ formed
by the nonempty connected projective objects.

\item Show that the topology on $P$ induced by that of $E$ is the
chaotic topology. Deduce the equivalence of the following conditions
on $E$ (due to J. E. ROOS [12 c), prop. 1]): (i) The family of
essential points of $E$ is conservative; (ii) The full subcategory
$P$ of $E$ formed by the nonempty connected projective objects is
generating; (iii)~$E$ is equivalent to a presheaf topos of the form
$\widehat{C}$, where $C$ is a category equivalent to a category
$\in 𝒰$ (or, if one prefers, $C\in 𝒰$). (Use $c$) and the fact that,
if $P$ is generating, the natural functor $E\rightarrow P^{\sim}$
is an equivalence of\pageoriginale categories.)

\item The category of essential points of a topos $E$ is Karoubian
(\cf \hyperref[IV.exer7.5]{7.5} c)). Let $C$ be a category
equivalent to a category $\in 𝒰$. Prove that the canonical fully
faithful functor (\hyperref[IV.subsubsec4.6.2]{4.6.2})
$$
C\longrightarrow \Ppoint(\widehat{C})
$$
factors through a functor
$$
C\longrightarrow \Ppointess(\widehat{C}){\quoi},
$$
which makes $\Ppointess(\widehat{C})$ a \emph{Karoubi envelope}
(\hyperref[IV.exer7.5]{7.5} b)) \emph{of} $C$. In particular, this
functor is an equivalence of categories if and only if the category
$C$ is Karoubian. (Hint: using $c$), prove that every essential point
of $\widehat{C}$ is isomorphic to a direct factor of some
$X\in \ob C$.)

\item Let $C$ and $C'$ be two categories equivalent to categories
$\in 𝒰$. Prove that the canonical fully faithful functor
(\hyperref[IV.subsubsec4.6.2]{4.6.2})
$$
\Hhom(C,C')\longrightarrow \Hhomtop(\widehat{C},\widehat{C}')
$$
takes its values in the full subcategory
$\Hhomtopess(\widehat{C},\widehat{C}')$ of essential morphisms
\emph{of topoi} (a)), and that the induced functor
$$
\Hhom(C,C')\longrightarrow \Hhomtopess(\widehat{C},\widehat{C}')
$$
is an equivalence of categories if and only if $C$ is empty or $C'$
is Karoubian. (Use $g$) below.)

\item Let $C$ be a category equivalent to a category $\in 𝒰$, and
let $E$ be a topos. Define an equivalence between the category
$\Hhomtopess(\widehat{C},E)$ of essential morphisms of topoi
$\widehat{C}\rightarrow E$, and the category
$\Hhom(C, \Ppointess(E))$. (Use $b$)\pageoriginale and the fact that
$\widehat{C}$ has enough essential points.)

\item Let $C$ be a finite category. Prove that every point of the
presheaf topos $\widehat{C}$ is essential, and that
$\Ppoint(\widehat{C})$ is equivalent to a finite category. (Note
that the Karoubi envelope of $C$ is finite, and, using
\eqhyperref[IV.eq6.8.6.1]{6.8.6.1}, reduce to proving that if $C$ is
a finite Karoubian category, then the canonical functor
$C\rightarrow \Pro(C)$ is an equivalence of categories.) Using b),
deduce that \emph{every} morphism of topoi
$\widehat{C}'\rightarrow \widehat{C}$ is essential, and that
$\Hhom(C,C')\rightarrow \Hhomtop(\widehat{C},\widehat{C}')$ is an
equivalence if and only if $C$ is empty or $C'$ is Karoubian.

\item Let $X$ be a topological space, and let
$f:\Top(X)\rightarrow P$ be the morphism of topoi induced by the
continuous map from $X$ to the one-point space. Show that $f$ is
essential if and only if the space $X$ is \emph{locally connected},
\ie if it satisfies the following condition: for every $x\in X$ and
every open neighborhood $U$ of $x$, the connected component of $x$
in $U$ is a neighborhood of $x$; or equivalently: for every open
subset $U$ of $X$, the connected components of $U$ are open. (\cf
\hyperref[IV.exer8.7]{8.7} l) for a generalization).
\end{enumerate}
\end{enonce*}

\begin{enonce*}[remark]{Exercise 7.7}\etiquette[7.7]{IV.exer7.7}
(\emph{Unusual points of a classifying topos})\footnote[1]{\cf also
\hyperref[IV.exer7.2.6]{7.2.6} d).}.
\begin{enumerate}
\renewcommand{\theenumi}{\alph{enumi}}
\renewcommand{\labelenumi}{\theenumi)}
\item Let $G$ be a monoid. The monoid $G$ contains an element $g$
such that $g\neq e$ and $g^2=g$ if and only if the classifying topos
$B_G$ admits an essential point that is not isomorphic to the
trivial point (corresponding to the fiber functor ``forget the
$G$-action''). (Use \hyperref[IV.exer7.6]{7.6} e)).

\item Let $G$ be the additive monoid of the integers $\geq 0$.
Show that every essential point of the classifying topos $B_G$ is
isomorphic to the trivial point. Construct a point of $B_G$ that is
not isomorphic to the trivial point. Show that the category
$\Ppoint(B_G)$ is not equivalent to a category $\in 𝒰$.
\end{enumerate}
\end{enonce*}

\begin{enonce*}[remark]{Exercise 7.8}\etiquette[7.8]{IV.exer7.8}
(\emph{Topology on $\Point(E)$, and topoi associated with ordered
sets}).
\begin{enumerate}
\renewcommand{\theenumi}{\alph{enumi}}
\renewcommand{\labelenumi}{\theenumi)}
\item Let $E$ be a topos. Denote by $\Point(E)$ the set of
isomorphism classes of points of $E$. For every open $U$ of $E$,
let $U'$ be the set\pageoriginale of classes of points $p$ of $E$
such that $U_p\neq \phi$. Show that $U\mapsto U'$ is an
order-preserving map from the ordered set $\Ouv(E)$ of opens of $E$
to the set of subsets of $X=\Point(E)$, and that this map commutes
with finite Inf and arbitrary $\Sup$. Deduce that the subsets of
$X$ of the form $U'$ define a topology on $X$ (called the
\emph{canonical topology on the set of classes of points of the
topos $E$}). Show that if $E$ has enough points, the map
$U\mapsto U'$ is an isomorphism from the ordered set $\Ouv(E)$ onto
the ordered set $\Ouv(X)$.

\item Let $\Oo\in 𝒰$ be an ordered set admitting finite Inf and
arbitrary $\Sup$, and hence defining a category $\cat(\Oo)$ having
finite products and arbitrary coproducts. We say that a family of
morphisms $U_i\rightarrow U$ with common target $U$ is covering if
$U$ is the $\Sup$ of the $U_i$. Suppose that coproducts in
$\cat(\Oo)$ are universal, \ie that in $\Oo$ the arbitrary $\Sup$
commute with Inf. This defines a topology on $\cat(\Oo)$, making
it a site, and hence a topos $\cat(\Oo)^{\sim}$, also denoted simply
by $\Oo^{\sim}$. Show that $\Oo$ is canonically isomorphic to
$\Ouv(\Oo^{\sim})$.

Show that the category $\Hhomtop(E,\Oo^{\sim})$ is equivalent to the
category associated with the ordered set of morphisms of ordered
sets $\Oo\rightarrow \Ouv(E)$ that commute with finite Inf and
arbitrary $\Sup$. (Use the criterion
\hyperref[IV.prop4.9.4]{4.9.4}.) If $\Oo'$ is an ordered set
satisfying the same conditions as $\Oo$, deduce that
$\Hhomtop(\Oo^{\sim}, {\Oo'}^{\sim})$ is equivalent to the category
associated with the ordered set of all morphisms of ordered sets
${\Oo}'\rightarrow\Oo$ that commute with finite Inf and
arbitrary $\Sup$.

\item Show\pageoriginale that the topos $E$ is equivalent to a topos
of the form $\Oo^{\sim}$, with $\Oo$ as in b), if and only if the
family of opens of $E$ is generating in $E$. It is equivalent to
some $\Top(X)$ with $X\in 𝒰$ if and only if it satisfies the
preceding condition and has enough points. Assuming the first
condition, express the second in terms of the structure of the
ordered set $\Oo=\Ouv(E)$, using the last assertion of b).

\item Define a canonical morphism of topoi
$$
E\longrightarrow \Ouv(E)^{\sim}{\quoi},
$$
which is universal (up to equivalence) for morphisms from $E$ to
topoi generated by their opens (\cf c)).

\item Let $X'$ be a \emph{small} subset of $X=\Point(E)$, equipped
with the topology induced by that of $X$. Define a canonical
morphism of topoi
$$
\Ouv(E)^{\sim}\longrightarrow \Top(X'){\quoi},
$$
which is an equivalence when the family $X'$ of points of $E$ is
conservative. Then deduce a canonical morphism of topoi
$$
E\longrightarrow \Top(X'){\quoi}.
$$
Give a universal characterization (up to equivalence) of this
morphism of topoi with respect to morphisms from the topos $E$ to
topoi of the form $\Top(Y)$. (Note in this connection that, up to
equivalence, $\Top(X')$ does not depend on the chosen small
conservative family of fiber functors, or equivalently
(\hyperref[IV.subsec4.2]{4.2}), that $X'_{\sob}$ is independent, up
to homeomorphism, of the choice of such a family.)

\item Suppose\pageoriginale that the category $\Ppoint(E)$ is not
equivalent to a small category (\cf \hyperref[IV.subsec7.3]{7.3}),
and that $E$ admits a small conservative family of fiber functors.
Show that, if such a family is taken sufficiently large, the subset
$X'$ of $\Ppoint(E)$ that it defines is not a sober topological
space for the induced topology.

\item For the relations with Exercise \hyperref[IV.exer7.6]{7.6},
\cf \hyperref[IV.sec8]{8}.
\end{enumerate}
\end{enonce*}

% END INCLUDED SGA4_English_current_buildable_source_20260906/source/04/drafts/SGA4_IV_LINES_4476_4830_English.texfrag


% BEGIN INCLUDED SGA4_English_current_buildable_source_20260906/source/04/drafts/SGA4_IV_LINES_4831_5055_English.texfrag
\section{Localization. Opens of a topos}\etiquette[8]{IV.sec8}

\subsection{}\etiquette[8.1]{IV.subsec8.1}
Let $C$ be a category, and let $T(X)$ be a relation involving an
object $X$ of $C$. The relation $T(X)$ is said to be \emph{stable
under base change} (in $X$) if, for every morphism $X'\rightarrow X$
of $C$, the relation $T(X)$ implies $T(X')$. This amounts to saying
that the full subcategory $R$ of $C$ formed by the objects $X$ for
which $T(X)$ holds is a \emph{sieve}
(\hyperrefext[I][I.def4.1]{4.1}).

\begin{enonce*}[remark]{Remark 8.1.1}\etiquette[8.1.1]{IV.rem8.1.1}
Suppose that $C$ is a $𝒰$-category
(\hyperrefext[I][I.def1.1]{1.1}), so that $C$ may be identified with
a full subcategory of $\widehat{C}$
(\hyperrefext[I][I.prop1.4]{1.4}). It is sometimes convenient to
extend the definition of the relation $T$ on $\ob C$ to a relation
$\widehat{T}$ on an object $F$ of $\widehat{C}$, by defining
$\widehat{T}(F)$ to mean that $T(X)$ holds for every arrow
$X\rightarrow F$ in $\widehat{C}$ with $X\in\ob C$. The relation
$\widehat{T}(F)$ is clearly still stable under base change in $F$.
Again denoting by $R$ the subobject of the terminal object $e$ of
$\widehat{C}$ defined by the sieve $R$ of $C$
(\hyperrefext[I][I.rem4.2.1]{4.2.1}), the relation $\widehat{T}(F)$
may also be expressed as
$\Hom_C^{\widehat{\phantom{,}}}(F,R)\neq∅$, \ie it means that the
unique morphism $F\rightarrow e$ factors through the subobject $R$
of $e$.
\end{enonce*}

\subsection{}\etiquette[8.2]{IV.subsec8.2}
Now suppose that $C$ is a \emph{site}. A relation $T(X)$ in an
argument $X\in\ob C$ is said to be \emph{stable under descent} if,
for every covering\pageoriginale family
$(X_i:X_i\rightarrow X)_{i\in I}$ of an object $X$ of $C$, the
relation ``$T(X_i)$ for every $i\in I$'' implies $T(X)$. We say that
$T$ is a \emph{relation of local nature} if it is stable under base
change (\hyperref[IV.subsec8.1]{8.1}) and stable under descent. When
$C$ is a $𝒰$-category and $T$ is already assumed stable under base
change, \ie is defined by a sieve $R$ of $C$, which we identify with
a subpresheaf of the terminal presheaf on $C$, one sees immediately
that $T$ is of local nature if and only if $R$ is a \emph{sheaf}.
Thus \emph{relations of local nature in an argument $X\in\ob C$
correspond exactly to the subsheaves of the terminal sheaf of $C$,
\ie to the subobjects of the terminal object of the topos
$C^{\sim}$.} They therefore depend essentially only on the topos
$C^{\sim}$ defined by the site $C$ (like all important notions
associated with a site!). In terms of the subsheaf $R$ of the
terminal sheaf, the relation $T(X)$ is expressed as
$\Hom(\epsilon(X),R)\neq∅$. It extends canonically to the relation
$T^{\sim}(F):\Hom(F,R)\neq∅$ on the topos $C^{\sim}$.

\begin{enonce*}[remark]{Remark 8.2.1}\etiquette[8.2.1]{IV.rem8.2.1}
With the notation introduced in \hyperref[IV.rem8.1.1]{8.1.1}, for a
presheaf $F$ on $C$, since
$\Hom(F,R)\simeq\Hom(\underline{a}(F),R)$, the relation
$\widehat{T}(F)$ is equivalent to the relation
$\widehat{T}(\underline{a}(F))$.
\end{enonce*}

\begin{enonce*}{Definition 8.3}\etiquette[8.3]{IV.def8.3}
An open of a $𝒰$-topos $E$ is a subobject of the terminal object
$e_E$ of $E$; if $X$ is an object of $E$, an open of $X$ sometimes
means an open of the induced topos $E_{/X}$, \ie a subobject of $X$.

An open of a $𝒰$-site $C$ is an open of the associated $𝒰$-topos
$C^{\sim}$.
\end{enonce*}

Since the terminal objects of $E$ are canonically isomorphic, the
ambiguity introduced in \hyperref[IV.def8.3]{8.3} by the choice of
$e_E$ is harmless. One\pageoriginale may also, in a more intrinsic
but practically less convenient way, define the opens of $E$ as the
full subcategories $R$ of $E$ for which the relation
$F\in\ob R$, for an object $F$ of $E$, is of local nature. Likewise,
the opens of a $𝒰$-site $C$ correspond bijectively to the full
subcategories of $C$ for which the relation $F\in\ob R$, for an
object of $C$, is of local nature; this notion is therefore
essentially independent of the chosen universe $𝒰$. Finally, by what
was said in \hyperref[IV.subsec8.2]{8.2}, a relation of local nature
on a topos (or on a site) is essentially the same thing as an open of
that topos (or site).

\setcounter{subsection}{3}
\subsection{Examples of opens}\etiquette[8.4]{IV.subsec8.4}

\subsubsection{}\etiquette[8.4.1]{IV.subsubsec8.4.1}
Let $X$ be a topological space, and interpret $\Top(X)$
(\hyperref[IV.subsec2.1]{2.1}) as the category of étale spaces over
$X$. Since the monomorphisms in $\Top(X)$ are plainly the injective
maps, it follows that the opens of the topos $\Top(X)$ identify with
the opens of the space $X$. More generally, the opens of a topos
$\Top(X,G)$ (\hyperref[IV.subsec2.4]{2.4}) identify with the opens of
$X$ which are \emph{invariant} under the action of $G$.

\subsubsection{}\etiquette[8.4.2]{IV.subsubsec8.4.2}
The opens of a $𝒰$-topos $E$ form a $𝒰$-small set
(\hyperrefext[I][I.prop7.4]{7.4}), ordered by inclusion, admitting
arbitrary suprema and finite infima. Its least element is the initial
object (or ``empty object'') $\phi_E$, and its greatest element is
the terminal object $e_E$ of $E$. These two elements are identical
only if $E$ is an ``empty topos'' (\hyperref[IV.subsec2.2]{2.2}).

\subsubsection{}\etiquette[8.4.3]{IV.subsubsec8.4.3}
Let $G$ be a monoid object in $E$, giving a topos $B_G$
(\sisi{\hyperref[IV.subsec2.4]{2.4}}{\hyperref[IV.subsec2.6]{2.6}}),
the category of objects of $E$ on which $G$ acts on the left. The
terminal object $e_G$ of $B_G$ is the\pageoriginale terminal object
$e_E$ of $E$ with the trivial action of $G$. It follows that the
ordered set of opens of $B_G$ is canonically isomorphic to that of
the opens of $E$. In particular, if $E$ is the punctual topos, so
that $G$ is an ordinary monoid, the set of opens of $B_G$ consists
of exactly two elements, namely $\phi_{B_G}$ and $e_G$.

\subsubsection{}\etiquette[8.4.4]{IV.subsubsec8.4.4}
If $E$ is a topos of the form $\widehat{C}$, where $C$ is a
$𝒰$-category, then (\hyperref[IV.subsec8.1]{8.1}) the ordered set of
opens of $E$ is isomorphic to the ordered set of sieves of $C$.

\subsubsection{}\etiquette[8.4.5]{IV.subsubsec8.4.5}
Let $X$ be a sober nondiscrete topological space. Show that
$\TOP(X)$ (\hyperref[IV.subsec2.5]{2.5}) has opens which do not arise
from opens of $X$, \ie from opens of $\Top(X)$, by inverse image
under the canonical morphism \eqhyperref[IV.eq4.10.1]{4.10.1}
$\TOP(X)\rightarrow\Top(X)$.

\subsection{Examples of relations of local nature}\etiquette[8.5]{IV.subsec8.5}

For every object $X$ of the topos $E$, denote by
$F\rightarrow F_X$ the \emph{localization functor}
$j_X^{\ast}:E\rightarrow E_{/X}$, $Y\mapsto Y\times X$
(\hyperref[IV.subsec5.2]{5.2}).

\subsubsection{}\etiquette[8.5.1]{IV.subsubsec8.5.1}
Let $u:F\rightarrow G$ be a morphism of $E$. The following relation
in the argument $X\in\ob E$:

$u_X:F_X\rightarrow G_X$ is a monomorphism (\resp an epimorphism,
\resp an isomorphism)

\noindent is a relation of local nature. In particular, if $G$ is a
group object of $E$, the relation ``$G_X$ is the trivial group over
$X$'' is of local nature; the corresponding open of $E$ is called
the \emph{cosupport of the group $G$.}

\subsubsection{}\etiquette[8.5.2]{IV.subsubsec8.5.2}
Let\pageoriginale $u,v:F\rightrightarrows G$ be two morphisms of
$E$. The relation in the argument $X\in\ob E$
$$
u_X=v_X:F_X\longrightarrow G_X
$$
is a relation of local nature. In particular, if $G$ is a group
object of $E$, and $u$ is a \emph{section} of $G$
(\hyperref[IV.subsubsec4.3.6]{4.3.6}), the relation in $X$
$$
\mbox{the section $u_X$ of the group $G_X$ of $E_{/X}$ is zero}
$$
is of local nature; the corresponding open of $E$ is called the
\emph{cosupport of the section $u$ of the group $G$.}

\begin{enonce*}[remark]{Remark 8.5.3}\etiquette[8.5.3]{IV.rem8.5.3}
When $E=\Top(X)$, the cosupport of a sheaf of groups $G$, \resp of a
section of such a sheaf $G$, is simply the open of $X$ complementary
to the \emph{support of} $G$, \resp to the \emph{support of} $u$, in
the classical terminology [TF].
\end{enonce*}

\setcounter{subsubsection}{3}
\subsubsection{}\etiquette[8.5.4]{IV.subsubsec8.5.4}
Let $P(f)$ be a relation in the argument $f\in{\rm f\ell\,}C$, where
$C$ is a site. When fiber products are representable in $C$, the
relation $P(f)$ is said to be \emph{stable under base change (\resp
stable under descent, \resp of local nature) on the target} (or
\emph{on the base}, or ``\emph{downstairs}'') if, for every morphism
$f:X\rightarrow Y$ of $C$, the following relation in the argument
$Y'\in\ob C_{/Y}$:

``the morphism $f':X'=\fprod{X}{Y'}{Y}\rightarrow Y'$ obtained from
$f$ by base change along the structural morphism $Y'\rightarrow Y$
satisfies $P(f')$''

is stable under base change, \resp stable under descent, \resp is of
local nature. When the topology of $C$ is defined by means of a
pretopology, the relation $P(f)$ is said to be \emph{of local nature
on the source} (or \emph{``upstairs''}) if, for every morphism
$f:X\rightarrow Y$ of $C$ and every family
$(g_i:X_i\rightarrow X)_{i\in I}$\pageoriginale in $\Cov(X)$, the
relation $P(f)$ is equivalent to the relation ``$P(fg_i)$ for every
$i\in I$.'' When $\Cov(X)$ consists of all covering families of $C$,
this condition also means that the following relation in the object
$X'$ of the induced site $C_{/X}$
\begin{center}
$P(fg)$, where $g:X'\rightarrow X$ is the structural morphism
\end{center}
is of local nature. Note that if the relation $P(f)$ is of local
nature upstairs, it is a fortiori of local nature downstairs.

\subsubsection{}\etiquette[8.5.5]{IV.subsubsec8.5.5}
Take, for example, $C=(\Sch)$, the category of schemes in $𝒰$, with
the faithfully flat quasi-compact topology (\SGA 3 IV 6.3) or a
coarser topology. Each of the following properties of a morphism is
then of local nature downstairs:

surjective, radicial, universally open, universally closed, proper,
quasi-compact, quasi-compact and dominant, universal homeomorphism,
separated, quasi-separated, locally of finite type, locally of finite
presentation, of finite type, of finite presentation, an
isomorphism, a monomorphism, an open immersion, a closed immersion, a
quasi-compact immersion, affine, quasi-affine, finite, quasi-finite,
integral, flat, faithfully flat, neat, smooth, étale

\noindent (\EGA IV 2.6.4, 2.7.1 and \EGA IV 17.7.1). On the other
hand, endow $(\Sch)$ with the pretopology for which, for every scheme
$X\in𝒰$, $\Cov(X)$ consists of the families of morphisms
$(f_i:X_i\rightarrow X)_{i\in I}$ which are surjective and for which
the $f_i$ are flat and locally of finite presentation. Each of the
following properties is then \emph{of local nature upstairs}:

locally of finite type, locally of finite presentation, of finite
type, flat, neat, smooth, étale.

\noindent (\EGA IV 17.7.5, 17.7.7).

% END INCLUDED SGA4_English_current_buildable_source_20260906/source/04/drafts/SGA4_IV_LINES_4831_5055_English.texfrag


% BEGIN INCLUDED SGA4_English_current_buildable_source_20260906/source/04/drafts/SGA4_IV_LINES_5056_5253_English.texfrag
\begin{enonce*}[remark]{Exercise 8.6}\etiquette[8.6]{IV.exer8.6}
Let\pageoriginale $f:E\rightarrow F$ be a morphism of $𝒰$-topoi.
Consider the following relation in the argument $X\in\ob F$: the
induced morphism $E_{/f^{\ast}(X)}\rightarrow F_{/X}$ is an
equivalence of topoi. Prove that this relation is of local nature.
\end{enonce*}

\begin{enonce*}[remark]{Exercise 8.7}\etiquette[8.7]{IV.exer8.7}
(\emph{Partitions of a topos, sum of topoi}). Let $E$ be a
$𝒰$-topos. A family $(e_i)_{i\in I}$ of opens of $E$, \ie of
subobjects of the terminal object $e$ of $E$, is called a
\emph{partition of} $e$ (or of the topos $E$) if the canonical
morphism $\amalg_{i\in I}e_i\rightarrow e$ is an isomorphism, \ie
(\hyperrefext[II][II.cor4.6.2]{4.6.2}) if $e$ is the supremum of the
$e_i$ and $i\neq j$ implies $e_i\cap e_j=∅_E$.
\begin{enumerate}
\renewcommand{\theenumi}{\alph{enumi}}
\renewcommand{\labelenumi}{\theenumi)}
\item Show that a family $(e_i)_{i\in I}$ of objects of $E$ is a
partition of $E$ if and only if the functor
\begin{equation*}
E\longrightarrow \prod_i E_{/e_i}\tag{8.7.1}
\end{equation*}\etiquette[8.7.1]{IV.eq8.7.1}
defined by the localization functors $E\rightarrow E_{/e_i}$ is an
equivalence of categories.

\item Conversely, let $(E_i)_{i\in I}$ be a family of $𝒰$-topoi,
with $I\in𝒰$. Prove that $E=\Pi_{i\in I}E_i$ is a $𝒰$-topos (called
the \emph{sum topos} of the family $(E_i)_{i\in I}$). Define a
partition $(e_i)_{i\in I}$ of $E$ and equivalences of categories
$E_{/e_i}\approx E_i$ in such a way that the projection functors
$E\rightarrow E_i$ identify with the localization functors
$E\rightarrow E_{/e_i}$.

\item Let $E$ be the sum topos of the family $(E_i)_{i\in I}$. Show
that, for every $i\in I$, the projection functor $E\rightarrow E_i$
is of the form $u_i^{\ast}$, where
\begin{equation*}
u_i:E_i\longrightarrow E\tag{8.7.2}
\end{equation*}\etiquette[8.7.2]{IV.eq8.7.2}
is\pageoriginale a morphism of topoi. Prove that, for every
$𝒰$-topos $F$, the functor $v\mapsto(v\circ u_i)_{i\in I}$
\begin{equation*}
\Hhomtop(E,F)\longrightarrow \prod_i\Hhomtop(E_i,F)\tag{8.7.3}
\end{equation*}\etiquette[8.7.3]{IV.eq8.7.3}
is an equivalence of categories.

\item Let $(e_i)_{i\in I}$ be a partition of the topos $E$. For
every morphism of topoi $g:F\rightarrow E$, the family
$(f_i)_{i\in I}$, where $f_i=g^{\ast}(e_i)$, is a partition of $F$,
and the induced morphisms $g_i:F_{/f_i}\rightarrow E_{/e_i}$ make it
possible to reconstruct $g$ (up to unique isomorphism), the functor
$g^{\ast}$ identifying with the Cartesian product of the functors
$g_i^{\ast}$ (taking into account the equivalences of the type in
\hyperref[IV.exer8.7]{8.7} 1)). Deduce a complete description of the
category $\Hhomtop(F,E)$ in terms of categories of the form
$\Hhomtop(F',E_i)$, where $F'$ is a topos of the form $F_{/f}$ (with
$f$ a direct summand of the terminal object of $F$) and
$E_i=E_{/e_i}$.

\item In particular, if $F$ is connected and nonempty (\cf
\hyperref[IV.subsubsec4.3.5]{4.3.5}), \ie if, for every partition
$(f_i)_{i\in I}$ of $F$, there exists a unique $i\in I$ such that
$f_i\neq\phi_F$, prove that the canonical functor
\begin{equation*}
\coprod_{i\in I}\Hhomtop(F,E_i)\longrightarrow
\Hhomtop(F,E)\tag{8.7.4}
\end{equation*}\etiquette[8.7.4]{IV.eq8.7.4}
induced by the morphisms of topoi
\eqhyperref[IV.eq8.7.2]{8.7.2} is an equivalence of categories. More
particularly, the family of morphisms of topoi
\eqhyperref[IV.eq8.7.2]{8.7.2} induces an equivalence of categories
$$
\coprod_{i\in I}\Ppoint(E_i)\xrightarrow{\approx}\Ppoint(E){\quoi}.
$$

\item A\pageoriginale subset $I$ of the set of opens of $E$ is
called a \emph{reduced partition} of $E$ if the identity family
$(e_i)_{i\in I}$ indexed by $I$ is a partition and the $e_i$ are
different from $\phi_E$. Show that the refinement relation
(\hyperrefext[I][I.sec4.3.2]{4.3.2}) between reduced partitions makes
the set $P$ of reduced partitions into a $𝒰$-small ordered set in
which the supremum of any two elements exists. This gives a strict
projective system $(I_p)_{p\in P}$, regarded as a pro-set and denoted
by $\pi_{\circ}(E)$. It is called the \emph{pro-$\pi_{\circ}$ of the
topos} $E$.

\item Prove that $\pi_{\circ}(E)$ pro-represents the functor
$$
T\longmapsto f_{\ast}f^{\ast}(T)=\Gamma(E,
T_E):(𝒰\text{-}\Ens)\longrightarrow (𝒰\text{-}\Ens)
$$
associated with the canonical morphism $f:E\rightarrow P$ from $E$
to the punctual topos $P$ (\hyperref[IV.subsec4.3]{4.3}). In
particular, this functor is representable if and only if
$\pi_{\circ}(E)$ is essentially constant, \ie isomorphic (as a
pro-set) to an ``ordinary'' set, again denoted $\pi_{\circ}(E)$. The
pro-set $\pi_{\circ}(E)$ consists of one point if and only if $E$ is
``connected and nonempty.'' It is empty if and only if $E$ is the
``empty topos'' (\hyperref[IV.subsec2.2]{2.2}).

\item Suppose that $E$ is quasi-compact, \ie every covering of its
terminal object $e$ admits a finite subcovering. Show that
$\pi_{\circ}(E)$ is then a profinite set and hence may be identified
(using the familiar equivalence between pro-objects of the category
of finite sets and compact totally disconnected spaces) with a
compact totally disconnected space, also denoted $\pi_{\circ}(E)$.

\item Let $X$ be a topological space, and let $\pi_{\circ}(X)$ be the
set of connected components of $X$, regarded as a constant pro-set.
Define a canonical\pageoriginale morphism of pro-sets
\begin{equation*}
\pi_{\circ}(X)\longrightarrow
\pi_{\circ}(\Top(X)){\quoi},\tag{8.7.5}
\end{equation*}\etiquette[8.7.5]{IV.eq8.7.5}
or equivalently a canonical map
\begin{equation*}
\pi_{\circ}(X)\longrightarrow \varprojlim
\pi_{\circ}(\Top(X)){\quoi}.\tag{8.7.6}
\end{equation*}\etiquette[8.7.6]{IV.eq8.7.6}
Show that the first morphism is an isomorphism if and only if $X$ is
homeomorphic to a sum space of connected spaces (which is the case,
in particular, when $X$ is locally connected).

\item Suppose that $X$ is a quasi-compact scheme. Prove that
\eqhyperref[IV.eq8.7.6]{8.7.6} is bijective.

\item Establish the ``functorial character'' in $X$ of the morphisms
\eqhyperref[IV.eq8.7.5]{8.7.5} and
\eqhyperref[IV.eq8.7.6]{8.7.6}, first making precise the meaning of
this expression.


\item (Compare \hyperref[IV.exer7.6]{7.6} i)). Show that the
following two conditions on the $𝒰$-topos $E$ are equivalent: (i)
The canonical morphism from $E$ to the punctual topos is
``essential'' (Exercise \hyperref[IV.exer7.6]{7.6} a)), \ie the
functor $I\mapsto I_E$ from $(𝒰\text{-}\Ens)$ to $E$ commutes with
products indexed by sets in $𝒰$. (ii) For every object $X$ of $E$,
the induced topos $E_{/X}$ has a $\pi_{\circ}(E_{/X})$ which is an
ordinary set, \ie $X$ is isomorphic to the sum of a family of
connected objects of $E$. In this case $E$ is called a
\emph{locally connected topos} (compare
\sisi{IV~}{}\hyperref[IV.exer2.7.5]{2.7.5}).
\end{enumerate}
\end{enonce*}

\begin{enonce*}[remark]{Exercise 8.8}\etiquette[8.8]{IV.exer8.8}
(\emph{Dominant morphisms of topoi}).
\begin{enumerate}
\renewcommand{\theenumi}{\alph{enumi}}
\renewcommand{\labelenumi}{\theenumi)}
\item Let $f:E'\rightarrow E$ be a morphism of topoi. Show that the
following conditions are equivalent:

\begin{itemize}
\item[(i)] $f_{\ast}(\phi_{E'})=\phi_E$ (where $\phi_E$ and
$\phi_{E'}$ denote the initial objects of $E$ and $E'$, respectively).

\item[(ii)] For\pageoriginale every object $X$ of $E$ which is
``nonempty,'' \ie not isomorphic to $\phi_E$, the object
$f^{\ast}(X)$ is nonempty in $E'$.

\item[$\rm(ii')$] As in (ii), with $X$ a subobject of the terminal
object $e_E$ of $E$, \ie with $X$ an open of the topos $E$.
\end{itemize}

The morphism of topoi $f$ is then called \emph{dominant}.

\item Let $f:X'\rightarrow X$ be a continuous map of topological
spaces, giving a morphism of topoi
$\Top(f):\Top(X')\rightarrow\Top(X)$
(\hyperref[IV.subsec4.1]{4.1}). Show that $\Top(f)$ is dominant if
and only if $f$ is dominant, \ie if $f(X')$ is dense in $X$.

\item Show that if $f:E'\rightarrow E$ is a ``conservative''
morphism of topoi (\hyperref[IV.subsubsec6.4.0]{6.4.0}), \ie if
$f^{\ast}$ is faithful, then $f$ is dominant. Show that the converse
need not hold, even for a localization morphism
$E_{/U}\rightarrow E$, where $U$ is an open of the topos $E$.
(Show that in this case $f$ is conservative only if $U$ is the
terminal object of $E$ (and hence only if $f$ is an equivalence of
topoi), and use the example in b).)

\item Let $f:X'\rightarrow X$ be an arrow in a topos $E$. We call
$f$ a \emph{dominant morphism} in the topos $E$ if the corresponding
morphism of induced topoi $E_{/X'}\rightarrow E_{/X}$
\eqhyperref[IV.eq5.5.2]{5.5.2} is dominant. Show that this property
depends only on the image $f(X')$ of $X'$ in $X$, as an element of
the ordered set $\ouv(X)$ of subobjects of $X$. Show that the
morphism $E_{/X'}\rightarrow E_{/X}$ is conservative if and only if
$f$ is covering.
\end{enumerate}
\end{enonce*}

% END INCLUDED SGA4_English_current_buildable_source_20260906/source/04/drafts/SGA4_IV_LINES_5056_5253_English.texfrag


% BEGIN INCLUDED SGA4_English_current_buildable_source_20260906/source/04/drafts/SGA4_IV_LINES_5254_5490_English.texfrag
\section{Subtopoi and gluing of topoi}\etiquette[9]{IV.sec9}\pageoriginale

\subsection{Subtopoi and embeddings of topoi}\etiquette[9.1]{IV.subsec9.1}

\begin{enonce*}{Definition 9.1.1}\etiquette[9.1.1]{IV.def9.1.1}
\begin{enumerate}
\renewcommand{\theenumi}{\alph{enumi}}
\renewcommand{\labelenumi}{\theenumi)}
\item A morphism of topoi $f:F\rightarrow E$ is called an embedding
if the functor $f_{\ast}:F\rightarrow E$ is fully faithful (\ie if
the adjunction morphism $f^{\ast}f_{\ast}\rightarrow\id_F$ is an
isomorphism).

\item A subtopos of a topos $E$ is a strictly full subcategory $E'$
of $E$ which is a topos and for which the inclusion functor
$\alpha:E'\rightarrow E$ is of the form $i_{\ast}$, where
$i:E'\rightarrow E$ is a morphism of topoi (\ie
(\hyperref[IV.subsubsec3.1.1]{3.1.1}) $\alpha$ admits a left adjoint
$i^{\ast}$ which is left exact). The morphism $i$ (which is
determined up to unique isomorphism) is then called the inclusion
morphism of the subtopos $E'$ in $E$.
\end{enumerate}
\end{enonce*}

\begin{enonce*}[remark]{Remark 9.1.2}\etiquette[9.1.2]{IV.rem9.1.2}
\begin{enumerate}
\renewcommand{\theenumi}{\alph{enumi}}
\renewcommand{\labelenumi}{\theenumi)}
\item It is clear that the inclusion morphism of a subtopos is an
embedding, and that a subcategory $E'$ of a topos $E$ is a subtopos
if and only if it is the essential image of a direct-image functor
$f_{\ast}$ associated with a morphism of topoi $f:F\rightarrow E$
which is an embedding.

\item It follows immediately from the definitions that a morphism of
topoi $f:F\rightarrow E$ is an embedding if and only if it is
isomorphic to a composite morphism
$$
F\xrightarrow{g}E'\xrightarrow{i}E{\quoi},
$$
where $g$ is an equivalence of topoi and $i$ is the inclusion
morphism of a subtopos $E'$ of $E$. The latter is uniquely determined
by the preceding conditions as the essential image of $f_{\ast}$,
and the preceding factorization is likewise unique up to unique
isomorphism. In practice, of course, one should use $g$ to identify
$F$ with the subtopos $E'$ of $E$ (compare
\hyperref[IV.subsubsec3.4.1]{3.4.1}).

\item An equivalence of topoi (\hyperref[IV.subsec3.4]{3.4})
$u:E\xrightarrow{\approx}E_1$ defines in the evident way a bijection
between the set of subtopoi of $E$ and the set of subtopoi of $E_1$.
Indeed, since $u_{\ast}$ is an equivalence of categories, it
establishes\pageoriginale a bijection between the set of strictly
full subcategories of $E$ and the set of strictly full subcategories
of $E_1$ (to $E'\subset E$ corresponds the essential image $E'_1$ of
$u_{\ast}|E'$), and one sees immediately that $E'$ is a subtopos of
$E$ if and only if $E'_1$ is a subtopos of $E_1$. Thus, in studying
subtopoi, one may reduce to the case where $E$ is of the form
$\widetilde{C}$, with $C$ a small site. In this case,
\hyperrefext[II][II.theo5.5]{5.5} shows that a strictly full
subcategory $E'$ of $E$ is a subtopos if and only if its inclusion
functor $\alpha:E'\rightarrow E$ admits a left adjoint which is left
exact (which already implies that $E'$ is a topos), and that these
subtopoi are in one-to-one correspondence with the topologies $T'$
on $C$ finer than the given topology $T$ of $C$: to such a $T'$ one
associates the strictly full subcategory $(C,T')^{\sim}$ of
$C^{\sim}$ formed by the sheaves for the topology $T'$. This also
shows, for \emph{every} $𝒰$-topos $E$, that a strictly full
subcategory $E'$ of $E$ is a subtopos if and only if its inclusion
functor $\alpha:E'\rightarrow E$ admits a left adjoint which is left
exact; \emph{the ordered set $\mathfrak{E}(E)$ of subtopoi of $E$ is
$𝒰$-small, and every subset of $\mathfrak{E}(E)$ admits a supremum
and an infimum.}

\item If $E'$ and $E''$ are two subtopoi of a topos $E$, their
intersection is a subtopos of $E$. Indeed, since this intersection
is strictly full, one reduces to the case where $E$ is of the form
$C^{\sim}$. With the notation of \sisi{b)}{c)}, $E'$ and $E''$
then correspond to two topologies $T'$ and $T''$ on $C$ finer than
$T$, and it suffices to see that $E'\cap E''$ \emph{is the category
of sheaves for the topology $T'''$ which is the supremum of $T'$ and
$T''$}. We indicate the proof of this fact (which should have
appeared as a corollary after \hyperrefext[II][II.theo4.4]{4.4} or
\hyperrefext[II][II.theo5.5]{5.5}!). Let $E'''$ be the category of
sheaves for $T'''$, plainly contained in $E'$ and $E''$, and hence in
$E'\cap E''$. For every full subcategory $D$ of $E$, let $t(D)$ be
the finest topology among those for which the objects of $D$ are
sheaves\pageoriginale (\hyperrefext[II][II.prop2.2]{2.2}). The
inclusions $E'''\subset E'\cap E''\subset E',E''$ plainly imply
$t(E'),t(E'')\leq t(E'\cap E'')\leq t(E''')$. On the other hand,
(\sisi{\hyperrefext[II][II.theo4.4]{4.4} I}{\hyperrefext[II][II.cor4.4.4]{4.4.4}})
$t(E')=T'$, $t(E'')=T''$, and $t(E''')=T'''$, so the preceding
inequalities and the definition of $T'''$ as the supremum imply
$T'''=t(E'\cap E'')$, and hence $E'\cap E''\subset E'''$, \cqfd\
More generally, this argument shows that if $T'$ is the supremum of
an arbitrary family of topologies $(T_i)$, the category $E'$ of
sheaves for $T'$ is the intersection of the categories $E_i$ of
sheaves for the $T_i$. In particular, one obtains:
\end{enumerate}
\end{enonce*}

\begin{enonce*}[remark]{Proposition 9.1.3}\etiquette[9.1.3]{IV.prop9.1.3}
Let $E$ be a topos. For every family of subtopoi $(E_i)_{i\in I}$ of
$E$, the intersection $\bigcap_{i\in I}E_i$ is a subtopos of $E$.
\end{enonce*}

\begin{enonce*}{Proposition 9.1.4}\etiquette[9.1.4]{IV.prop9.1.4}
Let $i:E'\rightarrow E$ be an embedding of topoi, and let $F$ be a
topos. The functor
\begin{equation*}
f\longmapsto i\circ f:\Hhomtop(F,E')\longrightarrow
\Hhomtop(F,E)\tag{9.1.4.1}
\end{equation*}\etiquette[9.1.4.1]{IV.eq9.1.4.1}
is fully faithful, and its essential image consists of the morphisms
of topoi $g:F\rightarrow E$ for which the essential image of
$g_{\ast}$ is contained in that of $i_{\ast}$.
\end{enonce*}

Indeed, consider the evident commutative diagram
\[
\xymatrix@C=1.75cm@R=1cm@M=6pt{
\Hhomtop (F,E') \ar @{^{(}->}[d] \ar[r]& \Hhomtop (F,E) \ar @{^{(}->}[d]\\
\Hhom (F,E') \ar[r] &\Hhom (F,E) , }
\]
where the vertical arrows are the functors $h\mapsto h_{\ast}$,
which are fully faithful by the definition of $\Hhomtop$
(\sisi{IV}{}\hyperref[IV.subsec3.2]{3.2}). The second horizontal arrow
$G\mapsto i_{\ast}G$ is also fully faithful, since $i_{\ast}$ is, and
therefore so is the first horizontal arrow. It remains to prove the
characterization of the essential image of
\eqhyperref[IV.eq9.1.4.1]{9.1.4.1}. The necessity of the condition is
clear from the formula $(if)_{\ast}=i_{\ast}f_{\ast}$. Conversely,
suppose that the essential image of $g_{\ast}$ is contained in that
of $i_{\ast}$, \ie that one may write
$g_{\ast}=i_{\ast}f_{\ast}$, with $f_{\ast}:F\rightarrow E'$ a
functor. We must show that $g$ belongs to the essential image of
\eqhyperref[IV.eq9.1.4.1]{9.1.4.1}. This amounts to saying that
$f_{\ast}$ admits a left adjoint which is left exact.
Now\pageoriginale it is clear that $f_{\ast}$ admits
$g^{\ast}i_{\ast}$ as a left adjoint, and this functor is left exact,
being the composite of the left exact functors $g^{\ast}$ and
$i_{\ast}$, \cqfd

For a converse to \hyperref[IV.prop9.1.4]{9.1.4}, \cf Exercise
\hyperref[IV.exer9.1.6]{9.1.6}.

Taking $F$ to be the punctual topos
(\sisi{IV~}{}\hyperref[IV.subsec2.2]{2.2}), one obtains in particular
from \hyperref[IV.prop9.1.4]{9.1.4}:

\begin{enonce*}{Corollary 9.1.5}\etiquette[9.1.5]{IV.cor9.1.5}
If $f:E'\rightarrow E$ is an embedding morphism of topoi, the
corresponding functor
\begin{equation*}
\Point(f):\Ppoint(E')\longrightarrow \Ppoint(E)\tag{9.1.5.1}
\end{equation*}\etiquette[9.1.5.1]{IV.eq9.1.5.1}
is fully faithful.
\end{enonce*}

\begin{enonce*}[remark]{Exercise 9.1.6}\etiquette[9.1.6]{IV.exer9.1.6}
(\emph{Embeddings of topoi and 2-monomorphisms}).

Let $i:E'\rightarrow E$ be a morphism of topoi.

\begin{enumerate}
\renewcommand{\theenumi}{\alph{enumi}}
\renewcommand{\labelenumi}{\theenumi)}
\item If $F$ and $G$ are two topoi \emph{over} $E'$, define (with
the notation of \hyperref[IV.exer5.14]{5.14} a)) a canonical functor
\begin{equation*}
f\longmapsto f\ast i:\Hhomtop_{E'}(F,G)\longrightarrow
\Hhomtop_E(F,G){\quoi}.\tag{9.1.6.1}
\end{equation*}\etiquette[9.1.6.1]{IV.eq9.1.6.1}

\item Suppose that $i$, as a 1-arrow of the 2-category of
$𝒱\text{-}𝒰$-topoi (\hyperref[IV.subsubsec3.3.2]{3.3.2}), is a
$2$-monomorphism, \ie that for every $𝒰$-topos $F$ which is an
element of $𝒱$, the functor \hyperref[IV.eq9.1.4.1]{(9.1.4.1)} is fully % add a ref
faithful. Show that, for every pair of $𝒰$-topoi $F,G$ (not
necessarily in $𝒱$), the functor \hyperref[IV.eq9.1.6.1]{(9.1.6.1)} is %% add a ref
fully faithful.

\item Show that $i$ is an embedding if and only if $i$ is a
$2$-monomorphism. (For sufficiency, apply b) to induced topoi
$E'_{/X'}$ and $E'_{/Y'}$, and use
\hyperref[IV.exer5.14]{5.14} c)).

\item Let $g:F\rightarrow E$ be a morphism of topoi, and suppose that
the $2$-fiber product (\hyperref[IV.prop5.11]{5.11})
\sisi{$F'=\fprod{F^2}{E'}{E}$}{$F'=F\times_E^2E'$}
exists (a condition which is always satisfied, as P. DELIGNE has
shown\nde{See also [16], Proposition 4.47.}). Let
$j:F'\rightarrow F$ be the projection morphism. Prove that if $i$ is
an embedding, then so is $j$. (Use c), and note that the stability of
the notion of a 2-monomorphism under 2-base change is formal.) When
$E'$ is a subtopos of $E$ and $i:E'\rightarrow E$ is the canonical
inclusion morphism, the essential image of $F'$ under $j_{\ast}$
(which is a subtopos of $F$ independent\pageoriginale of the
indeterminacy in the construction of a 2-fiber product) is called the
\emph{subtopos of $F$ which is the inverse image under
$g:F\rightarrow E$ of the subtopos $E'$ of $E$.}

\item With the notation at the end of d), choose as 2-fiber product
the subtopos of $F$ which is the inverse image of the subtopos $E'$ of
$E$. Let $X$ be an object of $F$. Show that, for $X$ to belong to
$F'$, it is necessary that it satisfy the following condition: for
every object $U$ of $F$, writing
$g_U:F_{/U}\rightarrow E$ for the morphism of topoi induced by $g$,
one has $g_{U\ast}(X|U)\in\Ob E'$, where
$X|U=(X\times U\xrightarrow{\pr_2}U)$ is the ``restriction of $X$ to
$U$.'' (NB: Note that if $i_U:F_{/U}\rightarrow F$ is the canonical
inclusion, then $g_U=g\circ i_U$, whence
$g_{U\ast}(X|U)=g_{\ast}(i_{U\ast}(X|U))=g_{\ast}(\Hhom(U,X))$,
where $\Hhom(U,X)$ is the object defined in \hyperref[IV.prop10.1]{10.1} below.)

\noindent \emph{Problem.} Conversely, is the preceding condition on
$X$ sufficient for $X$ to belong to $F'$? This amounts to asking
whether the strictly full subcategory of $F$ formed by these objects
$X$ is a \emph{subtopos} of $F$.

\item Let $X$ be an object of $E$, let $X'=i^{\ast}(X)$, and let
$i_{/X}:E'_{/X'}\rightarrow E_{/X}$ be the morphism of topoi induced
by $i$ (\eqhyperref[IV.eq5.10.1]{5.10.1}). Show that if $i$ is an
embedding, then so is $i_{/X}$. (Use d) and
\hyperref[IV.prop5.11]{5.11}.) Conversely, if
$(X_{\alpha})_{\alpha}$ is a family of objects covering the terminal
object, and if $i_{/X_{\alpha}}$ is an embedding for every $\alpha$,
then so is $i$.
\end{enumerate}
\end{enonce*}

% END INCLUDED SGA4_English_current_buildable_source_20260906/source/04/drafts/SGA4_IV_LINES_5254_5490_English.texfrag


% BEGIN INCLUDED SGA4_English_current_buildable_source_20260906/source/04/drafts/SGA4_IV_LINES_5491_5778_English.texfrag
\begin{enonce*}[remark]{Exercise 9.1.7}\etiquette[9.1.7]{IV.exer9.1.7}
(\emph{Subtopoi and multiplicative systems of arrows}).

\begin{enumerate}
\renewcommand{\theenumi}{\alph{enumi}}
\renewcommand{\labelenumi}{\theenumi)}
\item Let $i:E'\rightarrow E$ be a morphism of topoi, let $S$ be the set of
arrows $u$ of $E$ such that $i^{\ast}(u)$ is an isomorphism, and
consider the canonical functor induced by $i^{\ast}$ (\cf [2,
Chap. I] or \sisi{VI}{the proof of \hyperrefext[VI][VI.prop6.2]{6.2}}):
\begin{equation*}
\rho:E[S^{-1}]\longrightarrow E'{\quoi}.\tag{9.1.7.1}
\end{equation*}\etiquette[9.1.7.1]{IV.eq9.1.7.1}

Show that $i$ is an embedding if and only if the preceding functor
$\rho$ is an equivalence of categories, and that in this case $S$
admits both a calculus of left fractions and a calculus of right
fractions. (Use [2, Chap. I, 1.3].) In this case, the essential image
of $i_{\ast}$ can be recovered from $S$ as consisting of those
$X\in \Ob \sisi{\not{E}}{(E)}$ such that for every\pageoriginale
$u:Y\rightarrow Y'$ in $S$, the map $\Hom(u,X):\Hom(Y',X)\rightarrow
\Hom(Y,X)$ is bijective.

\item We henceforth assume that i is an embedding. Show that, for
every topos $F$, the functor
$$
f^{\ast}\longmapsto f^{\ast}i^{\ast}:\Hhom(E',F)\longrightarrow
\Hhom(E,F)
$$
induces an equivalence between the category of the $f^{\ast}$ arising
from morphisms of topoi $f:F\rightarrow E'$ (\ie which are left exact
and admit a right adjoint) and the category of functors
$g^{\ast}:E\rightarrow F$ arising from morphisms of topoi
$g:F\rightarrow E$ and such that $g^{\ast}$ takes arrows of $S$ to
isomorphisms.

\item Let $S$ be a subset of $Fl$ $E$, where $E$ is a
$𝒰$-topos. Show that $S$ is associated with a subtopos $E'$
of $E$ by the procedure of a) if and only if $S$ satisfies the
following conditions:
\begin{itemize}
\item[ST 1)] $S$ contains the invertible arrows and is stable under
composition and base change.

\item[ST 2)] If a composite vu belongs to $S$, then $u\in S$ if
and only if $v\in S$.

\item[ST 3)] If $u:X\rightarrow Y$ is such that there exists a
covering family $Y_i\rightarrow Y$, $i\in I$, for which
$u_i:X_i=\fprod{X}{Y_i}{Y}\to Y_{i}$ belongs to $S$ for every $i\in
I$, then $u\in S$.
\end{itemize}

Show that this gives a one-to-one correspondence between the set of
subtopoi $E'$ of $E$ and the set of subsets $S$ of $Fl$ $E$
satisfying conditions ST 1), ST 2), and ST 3). (Given $S$, associate
with it a $𝒰$-topology $T'$, finer than the canonical topology
$T$ of $E$, whose covering families $(X_i\rightarrow X)_{i\in I}$
are those for which the inclusion $X'\hookrightarrow X$ of the image
of the $X_i$ is $\in S$. Show that the category of sheaves for
$T'$ is equivalent to a subtopos $E'$ of $E$, and use
\hyperrefext[II][II.theo4.4]{4.4} to prove that it recovers $S$.)

\item Show that $S$ is determined by the subset $S_0$ of $S$
consisting of the \emph{monomorphisms}. Show that the map
$S\mapsto S_0$ establishes a bijection between the set of subsets
$S$ of $Fl$ $E$ satisfying conditions ST 1) through ST 3) of c) and
the set of subsets of $Fl$ $E$ consisting of
monomorphisms\pageoriginale and satisfying the same conditions ST
1) through ST 3).

\item Let $(E'_i)_{i\in I}$ be a family of subtopoi of $E$, and let
$(S_i)_{i\in I}$ be the corresponding family of subsets of $Fl$ $E$.
Show that $\bigcap_{i\in I} S_i$ satisfies conditions ST 1) through
ST 3) of c), and hence corresponds to a subtopos $E'$ of $E$, and
that the latter is the \emph{subtopos generated by the} $E'_i$, \ie
the \emph{supremum} (\cf \hyperref[IV.rem9.1.2]{9.1.2} c)) of the
subtopoi $E'_i$.

\item Let $A$ be a subset of $Fl$ $E$. Show that there is a smallest
subset $S$ of $Fl$ $E$ among those containing $A$ and satisfying
conditions ST 1) through ST 3), and that $S$ corresponds to the
largest subtopos $E'$ of $E$ such that the corresponding
multiplicative system of $Fl$ $E$ contains $A$. Show that $S$ is
equal to the union of the subsets $A_n$ $(n\geq 0)$ of $Fl$ $E$,
constructed inductively as follows: $A_0=A$, and $A_{n+1}$ consists
of the arrows $u:X\rightarrow Y$ which are of one of the following
four types (i) through (iv):
\begin{enumerate}
\renewcommand{\theenumii}{\roman{enumii}}
\renewcommand{\labelenumii}{(\theenumii)}
\item $u$ is invertible, or is the composite of two arrows of $A_n$,
or is obtained from an arrow of $A_n$ by base change.

\item $u$ is one of the factors of a composite arrow whose other
factor and whose composite both belong to $A_n$.

\item $u$ is the coproduct of a small family of arrows $\in
A_n$.

\item There exists an epimorphism $Y'\rightarrow Y$ such that
$u'=X'=\fprod{X}{Y'}{Y}\rightarrow Y'$, obtained from $u$ by base
change, belongs to $A_n$.
\end{enumerate}

\item With the notation of e), let $A$ be the union of the $S_i$,
let $S$ be the subset of $Fl$ $E$ obtained from $A$ by the procedure
of f), and let $E'$ be the corresponding subtopos of $E$. Show that
$E'$ is the intersection (= infimum) of the subtopoi $E'_i$.
\end{enumerate}
\end{enonce*}



\begin{enonce*}[remark]{Exercise 9.1.7.2}\etiquette[9.1.7.2]{IV.exer9.1.7.2}
(\emph{Canonical factorization of a morphism of topoi}).
\begin{enumerate}
\renewcommand{\theenumi}{\alph{enumi}}
\renewcommand{\labelenumi}{\theenumi)}
\item Let $f:F\longrightarrow E$
be a morphism of topoi. Show that, up to isomorphism, one can find a
factorization of $f$ as
\begin{equation*}
F\xrightarrow{f'}E'\xrightarrow{i} E,\tag{9.1.7.3}
\end{equation*}\etiquette[9.1.7.3]{IV.eq9.1.7.3}
where $f'$ is \emph{conservative} (\ie (\hyperref[IV.subsubsec6.4.0]{6.4.0})
${f'}^{\ast}$ is conservative, or equivalently faithful)\pageoriginale
and $i$ is an embedding, and that this factorization is unique up to
equivalence (\cf \hyperref[IV.prop9.1.4]{9.1.4}). To do this, introduce
the set $S_f\subset Fl$ $E$ of arrows $u$ of $E$ such that
$f^{\ast}(u)$ is an isomorphism, prove that $S_f$ satisfies the
conditions of \hyperref[IV.exer9.1.7]{9.1.7} c), take for $E'$ the corresponding subtopos of
$E$, and define ${f'}^{\ast}$ by \hyperref[IV.eq9.1.7.1]{(9.1.7.1)}, which corresponds to a
morphism of topoi $f'$ by \hyperref[IV.exer9.1.7]{9.1.7} b).

\item The subtopos of $E$ defined by the embedding $i$ is called the
subtopos of $E$ which is the \emph{image of the morphism of topoi}
$f$. Show that it is the smallest of the subtopoi $E_1$ of $E$ such
that $f$ factors, up to isomorphism, through $E_1$, \ie
(\hyperref[IV.prop9.1.4]{9.1.4}) the smallest of the subtopoi of $E$
containing the image of $f_{\ast}$ (or equivalently the essential
image of $f_{\ast}$). For $f$ to be an embedding, it is necessary
and sufficient that $f$ induce an equivalence $f'$ of $F$ with its
image; for $f$ to be conservative, it is necessary and sufficient
that its image be all of $E$.

\item Suppose that $f=\Top(f_0)$ is associated with a continuous map
of sober topological spaces $f_0:Y\rightarrow X$
(\hyperref[IV.subsec4.1]{4.1}). Consider the canonical factorization of
$f_0$ as
$$
Y\xrightarrow{{f'}_0}X'=f_0(Y)\xrightarrow{i_0}X{\quoi},
$$
where $i_0$ is the inclusion. Show that the canonical factorization
of $f$ then identifies with the corresponding factorization
\[
\xymatrix@C=2cm@M=.2cm{ \Top(Y) \ar[r]^{\Top(f'_0)}& \Top(X')
\ar@{^{(}->}[r]^{\Top(i_0)}& \Top(X). }
\]
Let $Z$ be the smallest sober subspace of $X$ containing $X'=f_0(Y)$,
\ie the set of points $x$ of $X$ such that $\{\bar{x}\}\cap X'$
is dense in $\{\bar{x}\}$. (Note that $Z=X'$ if the points of $X$
are closed, or if $X'$ is a \emph{constructible} subset of the
locally Noetherian space $X$.) Then $f$ is conservative, \ie the
subtopos of $E=\Top(X)$ which is the image of $F=\Top(Y)$ under
$f=\Top(f_0)$ is equal to $E$ itself, if and only if $Z=X$, \ie if
and only if $f_0(Y)$ is \emph{very dense} (EGA IV 10.13) in $X$. This
makes precise the extent to which the notion of a conservative
morphism of topoi may legitimately be regarded as the natural
generalization of the notion of a \emph{surjective} continuous map
of topological spaces.

\item Let\pageoriginale
\[
\xymatrix{
Y \ar[d]& Y'\ar[l] \ar[d]\\
X & X'\ar[l] }
\]
be a Cartesian diagram of topological spaces. If $Y$ and $X'$ are
discrete spaces, and hence $Y'$ is discrete, then the corresponding
diagram of topoi is a 2-fiber product diagram
(\hyperref[IV.prop5.11]{5.11}). (NB. For a counterexample when $Y$
and $X'$ are no longer assumed discrete, with $Y$ and $X'$ subspaces
of the Hausdorff space $X$, \cf \hyperref[IV.exer9.1.10]{9.1.10} c)).
Deduce an example of a 2-fiber product diagram of topoi
\[
\xymatrix{
F \ar[d]_{f}& F'\ar[l] \ar[d]^{f'}\\
E & E'\ar[l] }
\]
such that $f$ is conservative but $f'$ is not conservative, more
precisely with $E'=$ the punctual topos and $F'=$ the empty topos
(\hyperref[IV.subsec2.2]{2.2}). (Take $Y$ such that the image of $Y$
in $X$ is very dense and $\neq X$, and take $X'$ to consist of a
point not belonging to the image.)

\item Let $(E_i)_{i\in I}$ be a family of subtopoi of a topos $E$.
Prove that $E$ is the supremum of the $E_i$ if and only if the family
of embedding morphisms $E_i\rightarrow E$ is conservative
(\hyperref[IV.subsubsec6.4.0]{6.4.0}).
\end{enumerate}
\end{enonce*}

\begin{enonce*}[remark]{Exercise 9.1.8}\etiquette[9.1.8]{IV.exer9.1.8}
(\emph{Subtopoi and points of $E$. The case of topological spaces}).

Let $E$ be a topos.
\begin{enumerate}
\renewcommand{\theenumi}{\alph{enumi}}
\renewcommand{\labelenumi}{\theenumi)}
\item Let $Z$ be a full subcategory of $\Ppoint(E)$ (or, what amounts
to the same thing, a subset of $\Ob \Ppoint(E)$). Show that the set
$S(Z)$ of arrows of $E$ which are taken to bijections by the fiber
functors corresponding to the $z\in \Ob Z$ satisfies the conditions
of \hyperref[IV.exer9.1.7]{9.1.7} c), and hence defines a subtopos $E_Z$ of $E$. Show that
this subtopos has enough points, and that $Z$ is equivalent to a full
subcategory of $\Ppoint(E_Z)$. Show that every subtopos $E'$ of $E$
which has enough points can be obtained by the preceding procedure.
(Take for $Z$ the essential image of the functor \hyperref[IV.eq9.1.5.1]{(9.1.5.1)}.) Show
that one obtains\pageoriginale a bijection between the set of
subtopoi $E'$ of $E$ having enough points and a subset of the set of
strictly full subcategories $Z$ of $\Ppoint(E)$ which are stable
under direct factors (\sisi{\hyperrefext[I][I.sec10]{10}}{\hyperrefext[I][I.sec10.6]{10.6}})
and under small filtered projective limits (\cf
\hyperref[IV.subsec7.3]{7.3}). (Problem: Which subcategories of
$\Ppoint(E)$ arise in this way? This is essentially, in another
form, the problem already raised in \hyperref[IV.prop6.5.3]{6.5.3}, in view of the dictionary
in \hyperref[IV.rem9.1.2]{9.1.2} c). Note that, besides the necessary
conditions just indicated, there are also subtler conditions on the
topological nature of $\Point(E)$, such as sobriety; \cf e)).

\item Let $(Z_i)_{i\in I}$ be a family of full subcategories of
$\Ppoint(E)$, and let $Z$ be the full subcategory which is the union
of the $Z_i$. Show that the corresponding subtopos $E_Z$ defined in
a) is the subtopos of $E$ generated by the $E_{Z_i}$. (Use \hyperref[IV.exer9.1.7]{9.1.7}
e).) Conclude that the subtopos of $E$ generated by a family of
subtopoi having enough points also has enough points.

\item Let $f:E'\rightarrow E$ be an embedding of topoi. Show that
the map $U\mapsto f^{\ast}(U)$
\begin{equation*}
\Ouv(E)\longrightarrow \Ouv(E')\tag{9.1.8.1}
\end{equation*}\etiquette[9.1.8.1]{IV.eq9.1.8.1}
is surjective, and conclude that the map
\begin{equation*}
\Point(E')\longrightarrow \Point(E)\tag{9.1.8.2}
\end{equation*}\etiquette[9.1.8.2]{IV.eq9.1.8.2}
induced by $f$ on the isomorphism classes of points is a
homeomorphism from $\Point(E')$ onto a subspace of $\Point(E)$, for
the topologies defined in \hyperref[IV.exer7.8]{7.8} a). (For
injectivity, use \hyperref[IV.cor9.1.5]{9.1.5}.)

\item Let $f:X'\rightarrow X$ be a continuous map, giving a morphism
of topoi (\hyperref[IV.subsec4.1]{4.1})
$$
\Top(f):\Top(X')\longrightarrow \Top(X){\quoi}.
$$
Prove that the latter is an embedding if and only if the topology of
$X'$ is the inverse image under $f$ of that of $X$ (\ie, when $X'$
is a Kolmogorov space, for example a sober space, if and only if $f$
induces a homeomorphism of $X'$ onto a subspace of $X$). (For the
necessity,\pageoriginale use c), and for the sufficiency use
\hyperref[IV.exer9.1.7.2]{9.1.7.2} c)).

\item Let $X$ be a sober topological space, and let $E=\Top(X)$. Using
a) and d), establish an isomorphism of ordered sets $X'\mapsto
E_{X'}$ between the set of sober subspaces $X'$ of $X$ (\ie, when
$X$ is Hausdorff, the set of \emph{all} subsets of $X$) and
the set of subtopoi of $E$ having enough points. Show that $E_{X'}$
is equivalent to $\Top(X')$. If $(X'_i)_{i\in I}$ is a family of
sober subspaces of $X$, show that there is a smallest sober subspace
$X'$ of $X$ containing the $X'_i$, equal to the union of the $X'_i$
if I is finite or $X$ is Hausdorff, and that $E_{X'}$ is the
subtopos of $E$ generated by the $E_{X'_i}$. (Use b)).

\item Give an example of a topos $E$ of the form $\Top(X)$ and a
subtopos $E'$ which does not have enough points, and which is even
nonempty (\hyperref[IV.subsec2.2]{2.2}) but has no points. (Take $E$
of the form $\widehat{U}$ and $E'=U^{\sim}$, where $U$ is the ordered
set defined in \hyperref[IV.subsec7.4]{7.4}. Or better, take the
example in \sisi{9.1. 10. b)}{\hyperref[IV.exer9.1.10]{9.1.10} b)},
which shows that one may take for $X$ any nonempty compact space
without isolated points.)
\end{enumerate}
\end{enonce*}

% END INCLUDED SGA4_English_current_buildable_source_20260906/source/04/drafts/SGA4_IV_LINES_5491_5778_English.texfrag


% BEGIN INCLUDED SGA4_English_current_buildable_source_20260906/source/04/drafts/SGA4_IV_LINES_5779_6107_English.texfrag
\begin{enonce*}[remark]{Exercise 9.1.9}\etiquette[9.1.9]{IV.exer9.1.9}
(\emph{The case of topoi defined by an ordered set}).

\begin{enumerate}
\renewcommand{\theenumi}{\alph{enumi}}
\renewcommand{\labelenumi}{\theenumi)}
\item We use the yoga of \hyperref[IV.exer7.8]{7.8} b) and c), which gives a dictionary
between, on the one hand, topoi generated by their opens and morphisms
of such topoi, and, on the other hand, ordered sets admitting
arbitrary $\Sup$, finite infima, and in which the infimum of two elements %% standardize notation sup/Sup/\Sup etc.
distributes over arbitrary $\Sup$; morphisms between such ordered
sets are increasing maps commuting with finite infima and arbitrary
$\Sup$.

\item Let $E$ be a topos and $f:E'\rightarrow E$ an embedding of
topoi. Show that, for every generating family $(X_i)_{i\in I}$ in
$E$, the family $(f^{\ast}(X_i))_{i \in I}$ is generating. Deduce
that if $E$ is generated by its opens, then the same is true of
$E'$. (Use \hyperref[IV.exer9.1.8]{9.1.8} c).)

\item Let $f:E'\rightarrow E$ be a morphism of topoi generated by
their opens, associated with a morphism $g:\Oo\rightarrow \Oo'$ of
the corresponding ordered sets. Show that $f$ is an embedding if
and only if $g$ is surjective, \emph{and} $g$\pageoriginale is
surjective on the graphs of the order relations, \ie the order on
$\Oo'$ is the quotient of that on $\Oo$.

\item Let $E$ be a topos of the form $\Oo^{\sim}$ as in a). Deduce
from b) and c) a bijective correspondence between the set of
subtopoi of $E$ and the set of equivalence relations $R$ on $\Oo$
having the following properties, analogous to conditions ST 1) to
ST 3) of \hyperref[IV.exer9.1.7]{9.1.7} c):
\begin{itemize}
\item[Q0 1)] The relation $R$ is compatible with finite infima; in
other words, if $U$ and $U'$ are equivalent under $R$, then so are
$U\cap V$ and $U'\cap V$ for every $V\in \Oo$.

\item[Q0 2)] The relation $R$ is compatible with arbitrary $\Sup$;
in other words, if $U_i$ is equivalent to $U'_i$ for every $i\in I$,
then $\Sup_i U_i$ is equivalent to $\Sup_i U'_i$.
\end{itemize}

Show that if subtopoi $E'_i$ of $E$ correspond to the equivalence
relations $R_i$, then the subtopos generated by the $E'_i$
corresponds to the intersection of the relations~$R_i$.
\end{enumerate}
\end{enonce*}

\begin{enonce*}[remark]{Exercise 9.1.10}\etiquette[9.1.10]{IV.exer9.1.10}
(\emph{Intersections of subspaces; new topoi without points; nonexistence
of a complement to a subtopos}).

Let $X$ be a topological space, $E=\Top(X)=\Ouv(X)^{\sim}$
(\hyperref[IV.subsec2.1]{2.1}).

\begin{enumerate}
\renewcommand{\theenumi}{\alph{enumi}}
\renewcommand{\labelenumi}{\theenumi)}
\item Let $R$ be the following relation on $\Ouv(X):(U,U')\in
R\Longleftrightarrow U\cap U'$ is dense in both $U$ and $U'$.
Show that $R$ is an equivalence relation satisfying conditions Q0
1) and Q0 2) of \hyperref[IV.exer9.1.9]{9.1.9} d), and therefore defines a
subtopos $E'$ of $E$. Show that, for a point $x\in X$, the
corresponding point belongs to the essential image of $\Ppoint(E')$
if and only if $x$ is \emph{thick}, \ie belongs to the closure of the
interior of $\bar{x}$. Show that $E'$ is the empty topos
(\hyperref[IV.subsec2.2]{2.2}) if and only if $X$ is empty.

\item If $X$ is sober, show that the category $\Ppoint(E')$ is
equivalent to the category defined by the ordered subset of $X$
(for the specialization relation) consisting of the thick points
of $X$. Show that if the points of $X$ are closed (\resp if $X$ is
Noetherian), then the thick points of $X$ are the isolated points of
$X$, \ie those for which $\{x\}$ is open (\resp are the maximal
points of $X$). Conclude that if $X$ is\pageoriginale a nonempty
Hausdorff topological space without isolated points, then the
subtopos $E'$ of $E=\Top(X)$ is a non-``empty'' topos having no
points.

\item Let $X'$ be a subset of $X$, and let $E_{X'}$ be the subtopos
of $E$ associated with $X'$ by the construction of
\hyperref[IV.exer9.1.8]{9.1.8} a). Show that if $X'$ is dense, then $E_{X'}$
contains $E'$. Show that if $X$ is sober and $X'$ and $X''$ are two
sober subsets of $X$, then their intersection $X'\cap X''$ is sober
and $E_{X'\cap X''}\subset E_{X'}\cap E_{X''}$, and that the
inclusion can be strict
\footnote[1]{See, however, \hyperref[IV.exer9.1.11]{9.1.11} b) below for the
case where $X$ is locally Noetherian.}. (Take a nonempty Hausdorff
space $X$ having two disjoint dense subsets $X'$, $X''$.) Show that
the preceding inclusion is an equality if and only if the subtopos
$E_{X'}\cap E_{X''}$ has enough points. Show that this is the case
if $X'$ or $X''$ is a \emph{locally closed} subset of~$X$.

\item Let $X'$ be a subset of $X$, and let $X''$ be the complement
of $X'$. A subtopos $F$ of $E$ contains the $E_{\{x\}}$ for $x\in
X''$ if and only if it contains $E_{X''}$ (\cf
\hyperref[IV.exer9.1.8]{9.1.8} b)). Conclude that if $X''$ is a union of
locally closed subsets of $X$ (for example, if the points of $X''$
are closed), then the set of subtopoi $F$ of $E$ whose intersection
with $E_{X'}$ is the empty subtopos $\Top(∅)$ has a greatest element
(which then deserves the name of \emph{weak complement of}
$E_{X'}$ in $E$; \cf \hyperref[IV.exer9.1.13]{9.1.13}) \emph{only if}
$E_{X'}\cap E_{X''}$ is the empty subtopos (a condition that is not
satisfied if $X\neq ∅$ and if $X'$ and $X''$ are both dense in $X$,
with $X$, $X'$, and $X''$ sober; \cf c)). Show that when this
condition fails, the infimum of two elements in the ordered set of
subtopoi of $E$ is not distributive with respect to arbitrary
$\Sup$. (The editor does not know whether it is distributive with
respect to finite $\Sup$; \cf, however, \hyperref[IV.exer9.1.11]{9.1.11} e)
and \hyperref[IV.exer9.1.12]{9.1.12} d)).
\end{enumerate}
\end{enonce*}

\begin{enonce*}[remark]{Exercise 9.1.11}\etiquette[9.1.11]{IV.exer9.1.11}
(\emph{Subtopoi of locally Noetherian topoi. The case of Noetherian
spaces}).
\begin{enumerate}
\renewcommand{\theenumi}{\alph{enumi}}
\renewcommand{\labelenumi}{\theenumi)}
\item Let $f:E'\rightarrow E$ be an embedding of topoi. Show that if
$X$ is a pre-Noetherian object of $E$ (\ie (\sisi{VI 1.}{\hyperrefext[VI][VI.def1.30]{1.30}}) every
increasing sequence of subobjects\pageoriginale of $X$ is stationary),
then $X'=f^{\ast}(X)$ is a pre-Noetherian object of $E'$.
(By introducing the induced topoi ${E'}_{/X'}$ and $E_{/X}$ and
using \hyperref[IV.exer9.1.6]{9.1.6} e), first reduce to the case where $X$
is the final object $e$ of $E$, so that $X'$ is the final object $e'$
of $E'$. Then observe that every increasing sequence of subobjects
of $e'$ defines, upon applying $f_{\ast}$, an increasing sequence of
subobjects of $f_{\ast}(e')=e$.) Conclude that if $E$ admits a
generating family consisting of pre-Noetherian objects (\resp if $E$
is a Noetherian topos, \resp locally Noetherian (\sisi{VI 2.)}{\hyperrefext[VI][VI.def2.11]{2.11}})), then
$E'$ satisfies the same condition. (Use \hyperref[IV.exer9.1.9]{9.1.9} b).)

\item Let $E$ be a locally Noetherian topos (\sisi{VI 2.}{\hyperrefext[VI][VI.def2.11]{2.11}}). Show that
every subtopos of $E$ is of the form $E_Z$ (\hyperref[IV.exer9.1.8]{9.1.8}
a)), for a suitable full subcategory $Z$ of $\Ppoint(E)$. (Use a)
and \textsc{Deligne}'s theorem that every locally Noetherian topos
has enough points.) Conclude that if $E$ is of the form $\Top(X)$,
where $X$ is a sober locally Noetherian topological space, then
$Z\mapsto E_Z$ is an isomorphism from the ordered set of sober
subspaces of $X$ to the ordered set of subtopoi of $E$. (Use
\hyperref[IV.exer9.1.8]{9.1.8} e).) The intersection of sober subspaces $Z_i$
is necessarily sober and defines the intersection of the subtopoi
$E_{Z_i}$ (in contrast with what can occur in the general case
(\hyperref[IV.exer9.1.10]{9.1.10} c)).

\item Let $X$ be a sober locally Noetherian space, let $X'$ be a
subset of $X$, and let $X''$ be its complement. Show that the
following conditions are equivalent:
\begin{enumerate}
\renewcommand{\theenumii}{\roman{enumii}}
\renewcommand{\labelenumii}{(\theenumii)}
\item $X'$ is constructible.

\item $X'$ and $X''$ are sober.

\item $X'$ is sober, and the subtopos $E_{X'}\approxeq \Top(X')$ of
$E=\Top(X)$ admits a ``complementary'' subtopos (\cf
\sisi{\hyperref[IV.exer9.1.13]{9.1.13} e)}{\hyperref[IV.exer9.1.13]{9.1.13}}).

\item The intersection of the subtopoi $E_{X'}\approx\Top(X')$ and
$E_{X''}\approx\Top(X'')$ of $E=\Top(X)$ is the ``empty'' subtopos
of $E$.
\end{enumerate}

Show that if $E'$ and $E''$ are two subtopoi of $E$, they are weak
complements if and only if they are complements
(\hyperref[IV.exer9.1.13]{9.1.13}).

\item Let\pageoriginale $X$ be a sober Noetherian space. Show that
the following conditions are equivalent:
\begin{itemize}
\item[(i)] Every sober subspace $X'$ of $X$ is constructible, \ie,
according to c), its complement is sober.

\item[(i bis)] For every maximal point $x$ of $X$, $X-\{x\}$ is
sober, \ie $\{x\}$ is constructible.

\item[(ii)] $X$ is finite.

\item[(iii)] Every subtopos of $\Top(X)$ admits a complementary
subtopos (\cf \hyperref[IV.exer9.1.13]{9.1.13}).

\item[(iv)] In the ordered set of subtopoi of $\Top(X)$, the
infimum of two elements is distributive with respect to arbitrary
$\Sup$.
\end{itemize}
(Use the argument of \hyperref[IV.exer9.1.10]{9.1.10} d), together with the
intersection formula proved in b).)

\item Let $X$ be a locally Noetherian space. Show that in the ordered
set of subtopoi of $\Top(X)$, the infimum of two elements is
distributive with respect to \emph{finite} $\Sup$. (Reduce to the
case where $X$ is sober, and use b).)

\item For every sober subspace $Z$ of the sober locally Noetherian
space $X$, let $T_Z$ be the topology on the category $\Ouv(X)$ of
open subsets of $X$ for which a family of inclusions
$U_i\rightarrow U$ is covering if and only if $U\cap
Z=\bigcup_i U_i \cap Z$. Show that $Z\mapsto T_Z$ establishes a
bijection (reversing inclusion relations) between the set of sober
subspaces $Z$ of $X$ and the set of topologies on $\Ouv(X)$ finer
than the canonical topology $T_X$. (Use b) and
\hyperref[IV.rem9.1.2]{9.1.2} c).)

\item Let $X$ be a sober topological space. A subset $X'$ of $X$ is
sober if and only if its complement is a union of locally closed
subsets; if $X$ is locally Noetherian, this also means that $X'$ is
``proconstructible'' (EGA IV 1.9.4). Show that there is a topology
$T_{\cons'}$ on $X$ such that the closed subsets of
$X_{\cons'}=(X,T_{\cons'})$ are the sober subsets of $X$. Assume
henceforth that $X$ is locally Noetherian, so that $X_{\cons'}$ is
none other than the space $X_{\cons}$\pageoriginale of EGA IV 9.1.13.
Deduce from b) and c) that the ordered set $\Sigma$ of subtopoi of
$\Top(X)$ is isomorphic to the ordered set of closed subsets of
$X_{\cons}$, and that, under this correspondence, the subtopoi
admitting a complement correspond to the subsets that are both open
and closed in $X_{\cons}$. Show that \sisi{$X$}{$X_{cons}$} is locally compact,
and that it is compact if and only if $X$ is Noetherian, \ie the
topos \sisi{$\Top(E)$}{$\Top(X)$} is Noetherian.

\item Let $E$ be a locally Noetherian topos. Suppose that it has the
property considered in \hyperref[IV.subsubsec9.1.14]{9.1.14} b) (a condition
that may perhaps always hold\ldots). Show that the essential images
of the canonical functors $\Ppoint(E')\rightarrow \Ppoint(E)$, where
$E'$ ranges over the set $\Sigma$ of subtopoi of $E$, define a
collection of subsets of the set $X=\Point(E)$ of isomorphism classes
of points of $E$ which is the collection of closed subsets for a
\emph{topology} $T_{\cons}$ \emph{on} $X$, and that in this way one
obtains an isomorphism of ordered sets between the set $\Sigma$ of
subtopoi of $E$ and the set of closed subsets of $X_{\cons}$. Conclude
that in $\Sigma$, the $\Sup$ of two elements is distributive with
respect to \emph{arbitrary} infima. Show that $X_{\cons}$, which is not
necessarily $𝒰$-small, has an associated sober space
(\hyperref[IV.subsubsec4.2.1]{4.2.1}) which is $𝒰$-small.
\end{enumerate}
\end{enonce*}

\begin{enonce*}[remark]{Exercise 9.1.12}\etiquette[9.1.12]{IV.exer9.1.12}
(\emph{Finite topoi}). A topos $E$ is said to be \emph{finite} if it
is equivalent to a topos of the form $\widehat{C}$, with $C$ a
finite category.
\begin{enumerate}
\renewcommand{\theenumi}{\alph{enumi}}
\renewcommand{\labelenumi}{\theenumi)}
\item (Dictionary). Let $𝒱$ be a universe such that $𝒰\in 𝒱$. Let
(Karfiness) be the 2-category defined as the full 2-subcategory of
the category $(\Cat)_{𝒱}$ consisting of categories that are elements
of $𝒱$, are \emph{Karoubian and equivalent to a finite category},
and let (Topfin) be the 2-category defined as the full 2-subcategory
of $(𝒱\text{-}𝒰\text{-}\Top)$
(\hyperref[IV.subsubsec3.3.1]{3.3.1}) consisting of the \emph{finite}
$𝒰$-topoi that are $\in 𝒱$. Show that there are mutually
quasi-inverse 2-equivalences:
\hypertarget{sga4.exp.IV.equation.9.1.12.1}{}\begin{align*}
C & \longmapsto \widehat{C}: (\text{Karfiness})\longrightarrow
(\text{Topfin})\tag{9.1.12.1}\\
E & \longmapsto \Ppoint(E): (\text{Topfin})\longrightarrow
(\text{Karfiness}).\tag{9.1.12.2.}
\end{align*}\etiquette[9.1.12.2]{IV.eq9.1.12.2}
(Use \hyperref[IV.exer7.6]{7.6} h), e).)

\item Show\pageoriginale that a finite topos is Noetherian
\sisi{(VI 2)}{\hyperrefext[VI][VI.def2.11]{2.11}}.

\item Show that every subtopos of a finite topos $E$ is a finite
topos. More precisely, if $E$ is equivalent to $\widehat{C}$, where
$C$ is a finite Karoubian category (or, more generally, is equivalent
to a finite category), show that one obtains an order isomorphism
between the set of subtopoi $E'$ of $E$ and the set of strictly full
subcategories $C'$ of $C$ that are Karoubian (or, equivalently,
stable under taking direct factors in $C$), by associating with every
$C'$ the essential image of
$f_{\ast}:\widehat{C}'\rightarrow \widehat{C}$, where
$f:C'\rightarrow C$ is the inclusion functor. (Use
\hyperref[IV.subsec5.6]{5.6} to verify that $f_{\ast}$ is fully faithful, and
b) and \hyperref[IV.exer9.1.11]{9.1.11} b) for the fact that every subtopos
$E'$ of $E$ is obtained in the indicated way.)

\item Let $E$ be a finite topos. Construct on the finite set
$X=\Point(E)$ of isomorphism classes of points of $E$ a topology
making it a sober topological space and such that the ordered set of
subtopoi of $E$ is canonically isomorphic to the set of closed subsets
of $X$. (Take the ``constructible topology'' $\Cc_{\cons}$ on $X$,
defined by the order relation $(x\leq y)\Leftrightarrow$ ($x$ is a
direct factor of $y$), for which the closure of a point $x$ consists
of the $y$ such that $y\leq x$.) In particular, the set $\Sigma$ of
subtopoi of $E$ is finite, and the binary infimum operation $\wedge$
is distributive with respect to arbitrary $\Sup$.

\item Let $C$ be a category equivalent to a finite category. Show
that, for every full subcategory $C'$ of $C$, there exists a topology
$T_{C'}$ on $C$ for which a family $(X_i\rightarrow X)_{i\in I}$ is
covering if and only if, for every $Y\in \Ob C'$, the family of maps
$\Hom(Y,X_i)\rightarrow \Hom(Y,X)$ $(i\in I)$ is surjective. Show
that the category of sheaves on $(C,T_{C'})$ is equivalent to
$\widehat{C}'$, and that $C'\mapsto T_{C'}$ is a bijection (reversing
the order relations) between the set of strictly full subcategories
$C'$ of $C$ that are Karoubian (or, equivalently, stable under taking
direct factors in $C$), and the set of topologies
(\hyperrefext[II][II.def1.1]{1.1}) on the category $C$. (Use $c$) and
\hyperref[IV.rem9.1.2]{9.1.2} c).)\pageoriginale In particular, if $C$ is a
site whose underlying category is equivalent to a finite category,
then the topos $C^{\sim}$ is a finite topos.

\item Let $f:C'\rightarrow C$ be a functor of $𝒰$-categories, with
$C$ equivalent to a finite category. Show that the morphism of topoi
$\widehat{f}:\widehat{C}'\rightarrow \widehat{C}$
(\hyperref[IV.subsec4.6]{4.6}) is conservative if and only if every object of
$C$ is isomorphic to a direct factor of an object in the image of
$f$. (Reduce to the case where $f$ is an inclusion
$f:C'\hookrightarrow C$, with $C'$ as in c).)

\item Let $C$ be the category having two objects $e$ (the final
object) and $a$, and, besides the identity arrows, three arrows
$f:a\rightarrow e$, $g:e\rightarrow a$, $p=gf$ (thus $p^2=p$).
Show that $E=\widehat{C}$ has, besides the subtopoi $E$ and
$\Top(∅)$, exactly one subtopos $E'$, namely
$E'=\widehat{\{e\}}$; consequently, for every subtopos $E''$ of $E$
with $E''\neq \Top(∅)$, one has $E'\cap E''\neq \Top(∅)$.
\end{enumerate}
\end{enonce*}

% END INCLUDED SGA4_English_current_buildable_source_20260906/source/04/drafts/SGA4_IV_LINES_5779_6107_English.texfrag


% BEGIN INCLUDED SGA4_English_current_buildable_source_20260906/source/04/drafts/SGA4_IV_LINES_6108_6401_English.texfrag
\begin{enonce*}[remark]{Exercise 9.1.13}\etiquette[9.1.13]{IV.exer9.1.13}
(\emph{Complementary subtopoi of a topos}).

Let $E$ be a topos and let $E'$ and $E''$ be two subtopoi. We say
that $E'$ and $E''$ are \emph{complementary} if
$E'\cap E''=\Top(∅)$ and $E'\vee E''=E$ (where $\vee$ denotes the
$\Sup$ in the ordered set of subtopoi of $E$).

Suppose that in the set $\Sigma$ of subtopoi of $E$, the infimum of
two elements is distributive over finite suprema.

\begin{enumerate}
\renewcommand{\theenumi}{\alph{enumi}}
\renewcommand{\labelenumi}{\theenumi)}
\item If $E'$ and $E''$ are complementary, then $E''$ is the largest
among the subtopoi $E''_1$ such that
$E'\cap E''_1=\Top(∅)$ (\ie $E''$ is a weak complement of $E'$, in
the terminology of \hyperref[IV.exer9.1.10]{9.1.10} d)).

\item Show that the following conditions on $E$ are equivalent:
\begin{itemize}
\item[(i)] If $E'$ and $E''$ are two subtopoi of $E$ such that $E''$
is a weak complement of $E'$, then $E''$ is a complement of $E'$.

\item[(ii)] For every subtopos $E'$ of $E$ such that $E'\neq E$,
there exists a subtopos $E''$ of $E$ such that
$E''\neq\Top(∅)$ and $E'\cap E''=\Top(∅)$.

\item[(i bis)] If $E'$ and $E''$ are two subtopoi of $E$ such that
$E''$ is maximal among the subtopoi $E''_1$ satisfying
$E'\cap E''_1=\Top(∅)$, then $E''$ is a complement of $E'$.
\end{itemize}

Show that even if $E$ is a finite topos
(\hyperref[IV.exer9.1.12]{9.1.12}), these\pageoriginale conditions
need not hold (\hyperref[IV.exer9.1.12]{9.1.12} g)). They do hold,
however, if $E$ is of the form $\Top(X)$, with $X$ a locally
Noetherian space. (Use \hyperref[IV.exer9.1.11]{9.1.11} b)).

\item A subtopos $E'$ of $E$ is called \emph{complemented} if it
admits a complement (compare \hyperref[IV.exer9.1.12]{9.1.12} c)).
Prove that the set of complemented subtopoi is stable under finite
infima and suprema and under complementation, and that
complementation transforms infima into suprema and suprema into
infima.

\item State the duals of observations a), b), and c) (which were in
fact trivial general statements about an abstract ordered set
$\Sigma$; note that the hypothesis made on $\Sigma$ is itself
self-dual).
\end{enumerate}
\end{enonce*}

\setcounter{subsubsection}{13}
\subsubsection{}\etiquette[9.1.14]{IV.subsubsec9.1.14}
\emph{Open questions}\nde{Positive answers to questions a) and d)
were given in [17].}. Let $E$ be a topos.

\begin{itemize}
\item[a)] In the set of subtopoi of $E$, is the infimum of two
elements distributive over finite suprema?

\item[b)] Let $E'$ and $E''$ be two subtopoi of $E$ such that
$E=E'\vee E''$. Is every point of $E$ then isomorphic to the image of
a point of $E'$ or a point of $E''$?

Note that if the answer to b) is affirmative for $E$ and all its
subtopoi, then the answer to a) is affirmative at least for subtopoi
having enough points (and hence for all subtopoi if $E$ is locally
Noetherian (\hyperref[IV.exer9.1.11]{9.1.11} b)). The answer to a)
(and a fortiori to b)) is nevertheless not known for all Noetherian
topoi. The answers to a) and b) are affirmative, however, if $E$ is
of the form $\Top(X)$, with $X$ a locally Noetherian space
(\hyperref[IV.exer9.1.11]{9.1.11}), or if $E$ is a finite topos
(\hyperref[IV.exer9.1.12]{9.1.12}).

Note that a) can be reformulated as follows (applied to all subtopoi
$F$ of $E$):

\item[$\rm a'$)] If $E'$ and $E''$ are two subtopoi of a topos $F$
such that the canonical morphism (\cf \hyperref[IV.exer8.7]{8.7} b))
\begin{equation*}
E'\amalg E''\longrightarrow F\tag{9.1.14.1}
\end{equation*}\etiquette[9.1.14.1]{IV.eq9.1.14.1}
is conservative (which also means that $F=E'\vee E''$
(\hyperref[IV.exer9.1.7.2]{9.1.7.2} e)), then the same is true after
2-base change by $F_1\hookrightarrow F$,\pageoriginale where $F_1$
is a subtopos of $F$.

There is an analogous reformulation $\rm b'$) of b), taking a base
change $F_1\rightarrow F$ by a \emph{punctual} topos $F_1$
(\hyperref[IV.subsec2.2]{2.2}). Thus affirmative answers to a) and b)
would follow from an affirmative answer to the following question
(where $F$ is any subtopos of $E$):

\item[c)] If $E'$ and $E''$ are two subtopoi of the topos $F$ whose
supremum is $F$, \ie such that the morphism of topoi
\eqhyperref[IV.eq9.1.14.1]{9.1.14.1} is conservative, is this
morphism even \emph{universally conservative}, \ie does it remain
conservative after every 2-base change $F_1\rightarrow F$?

This statement makes sense by a result of \textsc{Deligne} (see [16],
Proposition 4.47) asserting the existence of 2-fiber products of
topoi. See, however, the example in
\hyperref[IV.exer9.1.7.2]{9.1.7.2} d).

On the other hand, \hyperref[IV.exer9.1.11]{9.1.11} h) shows that an
affirmative answer to b) would, in the locally Noetherian case,
provide an affirmative answer to the following question:

\item[d)] In the ordered set $\Sigma$ of subtopoi of $E$, is the
infimum of two elements distributive over \emph{arbitrary} (not
necessarily finite) suprema?

This also means that $\Sigma$ may be interpreted as the ordered set
of ``closed subtopoi'' of a suitable topos $E_{\cons'}$ (namely
$\Sigma^{0^{\sim}}$, where the opposite ordered set $\Sigma^0$ of
$\Sigma$, regarded as a category, is endowed with its canonical
topology, \cf \hyperref[IV.exer7.8]{7.8} b)), which is generated by
the subobjects of the terminal object and is therefore very close to
an ordinary topological space. It would remain to study this topos
$E_{\cons'}$, taking inspiration from the examples
\hyperref[IV.exer9.1.11]{9.1.11} g) and
\hyperref[IV.exer9.1.12]{9.1.12} d), and in particular to determine
whether it is Noetherian when $E$ is. This amounts to the following
question for subtopoi $F$ of $E$, which makes sense independently of
d):

\item[e)] Let $F$ be a Noetherian topos, and let $(E_i)_{i\in I}$ be
a family of subtopoi of $F$ whose intersection is $\Top(∅)$. Is there
a finite subfamily with the same property?

In a rather different direction, recall also the
question\pageoriginale raised in \hyperref[IV.exer9.1.6]{9.1.6} e),
which deserves clarification.
\end{itemize}

\subsection{The case of open subtopoi}\etiquette[9.2]{IV.subsec9.2}

\subsubsection{}\etiquette[9.2.1]{IV.subsubsec9.2.1}

Let $U$ be an open of a topos $E$, \ie a subobject of the terminal
object $e$ of $E$ (\hyperref[IV.def8.3]{8.3}). Consider the
localization morphism (\hyperref[IV.subsec5.2]{5.2})
\begin{equation*}
j:E_{/U}\longrightarrow E{\quoi}.\tag{9.2.2}
\end{equation*}\etiquette[9.2.2]{IV.eq9.2.2}
Since $U\rightarrow e$ is a monomorphism, the functor $j_{!}$ (which
is interpreted as the forgetful functor) is fully faithful, and
hence so is its biadjoint $j_{\ast}$
(\sisi{1.5.7 a)}{\hyperrefext[I][I.rem5.7]{5.7}}). In other words,
\emph{the localization morphism $j$ in \hyperref[IV.eq9.2.2]{(9.2.2)} associated with an
open $U$ of the topos $E$ is an embedding of topoi}
(\hyperref[IV.def9.1.1]{9.1.1} a)).

An embedding of topoi $F\rightarrow E$ is called an \emph{open
embedding} if it is isomorphic to the embedding of topoi defined by
an open $U$ of $E$. A subtopos $E'$ (\hyperref[IV.def9.1.1]{9.1.1}
b)) of $E$ is called an \emph{open subtopos} if it is defined by an
open of $E$, \ie if the canonical inclusion morphism
$E'\rightarrow E$ is an open embedding. Note that the open $U$ of
$E$ associated with an open embedding $j:E'\rightarrow E$ is
uniquely determined as the subobject $j_{!}(e')$ of $e$, where $e'$
denotes the terminal object of $E'$. Thus one obtains a one-to-one
correspondence between opens $U$ of $E$ and open subtopoi of $E$.

\begin{enonce*}[remark]{Remark 9.2.3}\etiquette[9.2.3]{IV.rem9.2.3}
In accordance with the preceding definitions, the open subtopos of
$E$ associated with the open $U$ of $E$ is the strictly full
subcategory of $E$ which is the essential image of the direct-image
functor
$$
j_{\ast}:E_{/U}\longrightarrow E.
$$
One should not confuse this subcategory of $E$ with the strictly full
subcategory which is the essential image of $j_{!}$, formed by the
objects $X$ of $E$ for which the structural morphism $X\rightarrow e$
factors through $U$. These two subcategories plainly determine one
another and are canonically equivalent. This is why some authors may
have been tempted to call (and indeed have called) the strictly full
subcategory which is the essential image of $j_{!}$ the ``open
subtopos defined by $U$,'' since its description in terms of $U$ is
simpler than that\pageoriginale of the essential image of
$j_{\ast}$. This terminology causes no difficulty so long as one does
not intend to study kinds of subtopoi other than open subtopoi. Since
that is not our situation, we shall not follow the authors just
mentioned.

For the reader's convenience, we summarize the most important special
properties of open embeddings:
\end{enonce*}

\begin{enonce*}{Proposition 9.2.4}\etiquette[9.2.4]{IV.prop9.2.4}
Let
$$
j:\mathcal{U}\longrightarrow E
$$
be an open embedding of topoi
(\hyperref[IV.subsubsec9.2.1]{9.2.1}). Then the functor $j_{!}$ left
adjoint to $j^{\ast}$ exists, so that there is a sequence of three
adjoint functors
\[
\xymatrix@R=.35cm@C=.5cm{
           & j_!,\quad j^{\ast}, \quad j_{\ast} : &\\
&&\\
\mathcal{U} \ar @<-16pt>[rr]^{j_{\ast}} \ar
           @<14pt>[rr]^{j_{!}}&            &\ar @<1pt>[ll]_{j^{\ast}} E
}
\]

Moreover:
\begin{enumerate}
\renewcommand{\theenumi}{\alph{enumi}}
\renewcommand{\labelenumi}{\theenumi)}
\item The functors $j_{!}$ and $j_{\ast}$ are fully faithful.

\item The functor $j_{!}$ commutes with fiber products, with products
indexed by small sets $I\neq\phi$, and with projective limits indexed
by small cofiltered categories (\hyperrefext[I][I.def2.7]{2.7}).

\item For every object $X$ of $E$, the adjunction morphism
$$
j_{!}j^{\ast}(X)\longrightarrow X
$$
is a monomorphism.
\end{enumerate}
\end{enonce*}

\noindent \emph{Proof}. We may suppose that $j$ is the localization
morphism defined by an open $U$ of $E$. Then a) is recalled for
reference, b) follows from the interpretation of $j_{!}$ as a
forgetful functor and from the fact that $U\rightarrow e$ is a
monomorphism. Finally, c) is immediate on observing that the morphism
in question is obtained from the inclusion $U\rightarrow e$ by base
change along $X\rightarrow e$, since a monomorphism is carried to a
monomorphism by base change.

\begin{enonce*}[remark]{Remark 9.2.4.1}\etiquette[9.2.4.1]{IV.rem9.2.4.1}
There exist embeddings of topoi $j:E'\rightarrow E$ for which
$j_{!}$ exists but does not commute with products of two objects or
with cofiltered projective limits (Exercise
\hyperref[IV.exer9.5.9]{9.5.9} c)).
\end{enonce*}

\setcounter{subsubsection}{4}
\subsubsection{}\etiquette[9.2.5]{IV.subsubsec9.2.5}
Now\pageoriginale let
$$
g:E'\longrightarrow E
$$
be an arbitrary morphism of topoi, and let $U'=g^{\ast}(U)$ be the
inverse image of the open $U$ of $E$. We saw in
\hyperref[IV.prop5.11]{5.11} that the diagram of topoi
\[
\xymatrix@C=2cm@R=1.5cm{
E'/_{U'} \ar[d]_{g_U} \ar[r]^{j'}& E' \ar[d]^{g}\\
E/_{U}  \ar[r]^{j} & E }
\]
is 2-Cartesian. It follows in particular that the inverse image under
$g$ of an open subtopos of $E$ (\cf
\hyperref[IV.exer9.1.6]{9.1.6} d)) is an open subtopos of $E'$, or
equivalently that the notion of an open embedding of topoi
(\hyperref[IV.subsubsec9.2.1]{9.2.1}) is stable under 2-base changes
in the 2-category of $𝒰$-topoi (which are elements of a given
universe $𝒱$).

Also note that the morphism of topoi $g:E'\rightarrow E$ factors, up
to isomorphism, through $j:E_{/U}\rightarrow E$ if and only if
$j':E'_{/U'}\rightarrow E'$ is an equivalence of categories, which
also means that the inclusion $U'\hookrightarrow e'$ (where $e'$ is
the terminal object of $E'$) is an isomorphism. This makes
\hyperref[IV.prop9.1.4]{9.1.4} more precise by giving a simple
characterization of the essential image of the functor considered
in \loccit

\subsubsection{}\etiquette[9.2.6]{IV.subsubsec9.2.6}
Apply \hyperref[IV.subsubsec9.2.5]{9.2.5} when $E'$ is of the form
$E_{/V}$, where $V$ is a second open of $E$, and $g$ is the
localization morphism. Then $U'=U\cap V$, and the 2-Cartesian product
of $E_{/U}$ and $E_{/V}$ over $E$ identifies with
$E_{/U\cap V}$. It follows that \emph{a finite intersection of open
subtopoi of $E$ is an open subtopos of $E$}; more precisely, the map
\begin{equation*}
U\longmapsto \text{ the subtopos of $E$ defined by $U$}
\tag{9.2.6.1}
\end{equation*}\etiquette[9.2.6.1]{IV.eq9.2.6.1}
commutes with finite intersections.

\begin{enonce*}[remark]{Exercise 9.2.7}\etiquette[9.2.7]{IV.exer9.2.7}
Prove that the map \eqhyperref[IV.eq9.2.6.1]{9.2.6.1} also commutes with \emph{arbitrary}
$\Sup$. (Use \hyperref[IV.exer9.1.7.2]{9.1.7.2} e)).
\end{enonce*}

% END INCLUDED SGA4_English_current_buildable_source_20260906/source/04/drafts/SGA4_IV_LINES_6108_6401_English.texfrag


% BEGIN INCLUDED SGA4_English_current_buildable_source_20260906/source/04/drafts/SGA4_IV_LINES_6402_6597_English.texfrag
\subsection{Construction of the closed subtopos complementary to an
open subtopos}\etiquette[9.3]{IV.subsec9.3}\pageoriginale

\subsubsection{}\etiquette[9.3.1]{IV.subsubsec9.3.1}
Let $\mathsf{T}$ be a topological space and $U$ an open subset of
$\mathsf{T}$, so that $\Top(U)$ identifies with an open subtopos of
$\Top(\mathsf{T})$ (notation of \hyperref[IV.subsec2.1]{2.1}). Let
$Y$ be the closed topological subspace of $\mathsf{T}$ complementary
to $U$. Up to equivalence, $\Top(Y)$ may then be regarded as a
subtopos of $\Top(\mathsf{T})$, namely the subtopos formed by the
objects $F$ of $\Top(\mathsf{T})$ whose restriction to $U$ is
\sisi{the terminal topos of $U$}{the terminal object of $U$}. This
description of a subcategory of $E=\Top(\mathsf{T})$ makes sense
whenever one has a topos and an open $U$ of it, and we shall see that
it always gives a subtopos of $E$. Moreover, in the particular case
first considered, it is immediate (with the terminology introduced
in Exercise \hyperref[IV.exer9.1.13]{9.1.13}) that $\Top(Y)$ and
$\Top(U)$ are complementary subtopoi (use
\hyperref[IV.subsubsec9.2.5]{9.2.5} and
\hyperref[IV.exer9.1.7.2]{9.1.7.2} e)). The same will hold in the
general case, and we shall see that this property uniquely
characterizes the subtopos in question. It therefore deserves in
every respect the name of the closed subtopos complementary to the
given open $U$, or to the open subtopos defined by $U$. The details
of the construction of this topos $\mathcal{F}$ and of the
inverse-image functor $i^{\ast}:E\rightarrow\mathcal{F}$ will be
given in the present section, while Section
\hyperref[IV.subsec9.4]{9.4} develops its first properties.

\subsubsection{}\etiquette[9.3.2]{IV.subsubsec9.3.2}
Thus let $E$ be a topos and $U$ an open of $E$, as in
\hyperref[IV.subsubsec9.2.1]{9.2.1}, and consider the localization
morphism
\begin{equation*}
j:\mathcal{U}=E_{/U}\longrightarrow E{\quoi},\tag{9.3.2.1}
\end{equation*}\etiquette[9.3.2.1]{IV.eq9.3.2.1}
which is an open embedding.

For every object $X$ of $E$, put
\begin{equation*}
X_{\complement U}=U\coprod_{U\times X}
X{\quoi},\tag{9.3.2.2}
\end{equation*}\etiquette[9.3.2.2]{IV.eq9.3.2.2}
where the amalgamation is formed with respect to the canonical
projections
$$
\pr_1:U\times X\longrightarrow U,\quad \pr_2:U\times
X\longrightarrow X.
$$
Thus there is a \emph{cocartesian} diagram
(\hyperrefext[I][I.sec10]{10}) in $E$, depending functorially on
$X$:\pageoriginale
\[\vcenter{
\xymatrix@C=1.75cm@R=1.5cm{
U\times X \ar[r]^{\pr_2} \ar[d]_{\pr_1}& X \ar[d]\\
U \ar[r]        & X_{\complement U} }}\leqno{(9.3.2.3)}
\]\etiquette[9.3.2.3]{IV.eq9.3.2.3}
Note that in the example considered in
\hyperref[IV.subsubsec9.3.1]{9.3.1}, $X_{\complement U}$ is
canonically isomorphic to $i_{\ast}i^{\ast}(X)$ (where
$i:Y\rightarrow T$ is the inclusion), and
$X\rightarrow X_{\complement U}$ identifies with the adjunction
morphism.

\begin{enonce*}{Proposition 9.3.3}\etiquette[9.3.3]{IV.prop9.3.3}
With the preceding notation, the following statements hold:
\begin{enumerate}
\renewcommand{\theenumi}{\alph{enumi}}
\renewcommand{\labelenumi}{\theenumi)}
\item For every object $X$ of $E$, the following conditions are
equivalent:
\begin{itemize}
\item[(i)] There exists an object $Y$ of $E$ and an isomorphism
$X\simeq Y_{\complement U}$.

\item[(ii)] The canonical morphism
$X\xrightarrow{m}X_{\complement U}$ is an isomorphism.

\item[(iii)] The canonical morphism
$\pr_2:X\times U\rightarrow U$ is an isomorphism (\ie
$j^{\ast}(X)$ is a terminal object of $E_{/U}$).
\end{itemize}

\item For every morphism
$$
m:X\longrightarrow X'
$$
in $E$, the following conditions are equivalent:
\begin{itemize}
\item[${\rm (i')}$] The diagram
\[
\xymatrix@C=2cm@R=1.5cm{
X \ar[r]^{m} \ar[d]& X'\ar[d]\\
X_{\complement U} \ar[r]& X'_{\complement U} }
\]
is Cartesian.

\item[${\rm (ii')}$] The morphism
$m\times\id_U:X\times U\rightarrow X'\times U$ is an isomorphism,
\ie $j^{\ast}(m)$ is an isomorphism.
\end{itemize}

\item The functor
$$
X\longrightarrow X_{\complement U}
$$
from $E$ to $E$ is left exact, carries epimorphic families to
epimorphic families, and commutes with filtered inductive limits
and\pageoriginale amalgamated sums.
\end{enumerate}
\end{enonce*}

\noindent \emph{Proof}. First prove the proposition when $E$ is of
the form $\widehat{C}$, where $C$ is a small category. By \hyperref[I.prop3.1]{I.3.1} this
reduces to the case where $E$ is the ``punctual'' topos $(\Ens)$,
where the proposition is evident. One then passes to the general
case by the standard procedures of \hyperrefext[II][II.sec4]{4},
using the ``associated sheaf'' functor.

\begin{enonce*}{Proposition 9.3.4}\etiquette[9.3.4]{IV.prop9.3.4}
With the notation of \hyperref[IV.subsubsec9.3.2]{9.3.2}, the
strictly full subcategory $\mathcal{F}$ of $E$ formed by the objects
$X$ of $E$ such that $j^{\ast}(X)$ is a terminal object of
$\mathcal{U}$ \textup{(\cf \hyperref[IV.prop9.3.3]{9.3.3} a))} is a
subtopos of $E$: \ie (\hyperref[IV.def9.1.1]{9.1.1}) it is a topos,
and the inclusion functor $i_{\ast}:\mathcal{F}\rightarrow E$ is the
direct-image functor of a morphism of topoi
$i:\mathcal{F}\rightarrow E$, whose inverse-image functor is
$X\mapsto i^{\ast}X=X_{\complement U}$. The adjunction morphism
$X\rightarrow i_{\ast}i^{\ast}X$ is the canonical morphism
$X\rightarrow X_{\complement U}$ of
\eqhyperref[IV.eq9.3.2.3]{9.3.2.3}.
\end{enonce*}

\noindent \emph{Proof}: For every object $X$ of $E$, denote by
$u(X)$ the canonical morphism $X\rightarrow X_{\complement U}$. The
morphism $u(X)$ is functorial in $X$, and for every object $X$,
$u(X_{\complement U})$ is an isomorphism. It follows formally and
immediately from these two properties that the functor $i^{\ast}$ is
left adjoint to the functor $i_{\ast}$ and that the adjunction
morphism is $u(X)$. Since the functor
$X\rightarrow X_{\complement U}$ is left exact, the functor
$i^{\ast}$ is left exact. It then follows from
\hyperrefext[II][II.theo5.5]{5.5} that $\mathcal{F}$ is a topos, and
from Definition \hyperref[IV.def3.1]{3.1} that the pair
$(i^{\ast},i_{\ast})$ is a morphism of topoi.

\setcounter{subsubsection}{4}
\subsubsection{}\etiquette[9.3.5]{IV.subsubsec9.3.5}
The subtopos $\mathcal{F}$ of $E$ described in
\hyperref[IV.prop9.3.4]{9.3.4} is called the \emph{closed subtopos of
$E$ complementary to the open $U$ of $E$}, or, equivalently, to the
open subtopos of $E$ corresponding to $U$. In accordance with
\hyperref[IV.def9.1.1]{9.1.1} b), the morphism
$i:\mathcal{F}\rightarrow E$ constructed in
\hyperref[IV.prop9.3.4]{9.3.4} is called the \emph{inclusion
morphism}. A full subcategory $\mathcal{F}$ of $E$ is called a
\emph{closed subtopos} of $E$ if there exists an open $U$ of $E$ for
which $\mathcal{F}$ is the closed subtopos complementary to $U$.
Note that this $U$ is uniquely determined as
$i_{\ast}(\phi_{\mathcal{F}})$, where $\phi_{\mathcal{F}}$ is the
initial object of $\mathcal{F}$, equal to $i^{\ast}(\phi_E)$; it is
also called \emph{the open of $E$ complementary to the closed
subtopos $\mathcal{F}$}, and\pageoriginale the subtopos of $E$
defined by $U$ is also called the \emph{open subtopos of $E$
complementary to $\mathcal{F}$}. Finally, an embedding of topoi
(\hyperref[IV.def9.1.1]{9.1.1} a)) $i:F\rightarrow E$ is called a
\emph{closed embedding} if the corresponding subtopos of $E$ is
closed.

If $G$ is a group object on $E$ and $s\in\Hom(e_E,G)$, the
\emph{support} of $G$ (\resp of $s$) is the closed subtopos of $E$
complementary to the cosupport open of $G$ (\resp of $s$)
(\hyperref[IV.subsubsec8.5.1]{8.5.1},
\hyperref[IV.subsubsec8.5.2]{8.5.2}).

\subsubsection{}\etiquette[9.3.6]{IV.subsubsec9.3.6}
With the notation of \hyperref[IV.subsubsec9.3.2]{9.3.2} and
\hyperref[IV.prop9.3.4]{9.3.4}, one therefore has two morphisms of
topoi
\begin{equation*}
\mathcal{U}\xrightarrow{j}
E\xleftarrow{i}\mathcal{F}{\quoi},\tag{9.3.6.1}
\end{equation*}\etiquette[9.3.6.1]{IV.eq9.3.6.1}
defining five functors which form two sequences of adjoint functors,
$(j_{!},j^{\ast},j_{\ast})$ and $(i^{\ast},i_{\ast})$:
\[
\xymatrix@C=2cm@R=1.5cm@M=.2cm{ \Uu \ar@<6pt>[r]^{j_!}
\ar@<-13pt>[r]|{j_{\ast}} & E \ar@<3pt>[r]^{i^{\ast}}
\ar@<4pt>[l]|{j^{\ast}}& \Ff
\ar@<13pt>[l]|{i_{\ast}}.}\leqno{(9.3.6.2)}
\]\etiquette[9.3.6.2]{IV.eq9.3.6.2}

The functor
\begin{equation*}
\rho=i^{\ast}j_{\ast}:\mathcal{U}\longrightarrow
\mathcal{F}\tag{9.3.6.3}
\end{equation*}\etiquette[9.3.6.3]{IV.eq9.3.6.3}
is called the \emph{gluing functor} relative to the open $U$ of $E$.
More generally, if $\mathcal{U}$ and $\mathcal{F}$ are arbitrary
topoi, a functor $\rho:\mathcal{U}\rightarrow\mathcal{F}$ is called
a \emph{gluing functor} if it is left exact and accessible (\hyperref[I.def9.2]{I.9.2}),
or equivalently
(\sisi{\hyperref[IV.rem1.8]{1.8}\quad}{\hyperrefext[I][I.sec8]{8}})
if it admits a pro-adjoint. The reasons for this terminology will
appear in \hyperref[IV.prop9.4.1]{9.4.1} d) and
\hyperref[IV.theo9.5.4]{9.5.4} below.

% END INCLUDED SGA4_English_current_buildable_source_20260906/source/04/drafts/SGA4_IV_LINES_6402_6597_English.texfrag


% BEGIN INCLUDED SGA4_English_current_buildable_source_20260906/source/04/drafts/SGA4_IV_LINES_6598_6787_English.texfrag
\subsection{First properties of the closed subtopos
$\mathcal{F}$ and of $i:\mathcal{F}\rightarrow
E$}\etiquette[9.4]{IV.subsec9.4}

\begin{enonce*}{Proposition 9.4.1}\etiquette[9.4.1]{IV.prop9.4.1}
The notation is that of \hyperref[IV.subsubsec9.3.2]{9.3.2} and
\hyperref[IV.prop9.3.4]{9.3.4}.
\begin{enumerate}
\renewcommand{\theenumi}{\alph{enumi}}
\renewcommand{\labelenumi}{\theenumi)}
\item The functor $i_{\ast}:\mathcal{F}\rightarrow E$ carries
epimorphic families to epimorphic families. It commutes with the
formation of amalgamated sums and filtered inductive limits.

\item For every object $X$ of $E$ such that the projection morphism
$U\times X\rightarrow U$ is an epimorphism, the adjunction morphism
$X\rightarrow i_{\ast}i^{\ast}X$ is an epimorphism.

\hypertarget{sga4.exp.IV.proposition.9.4.1.item.c}{}\item The\pageoriginale pair of functors $(i^{\ast},j^{\ast})$ from
$E$ to $\mathcal{F}$ and $\mathcal{U}$, respectively, is
conservative {\rm(I \hyperref[I.def6.1]{6.1})}, \ie a morphism $m$ of $E$ is an
isomorphism if and only if $i^{\ast}(m)$ and $j^{\ast}(m)$ are
isomorphisms (equivalently, the pair $(i^{\ast},j^{\ast})$ is
faithful (\hyperref[IV.subsubsec6.4.0]{6.4.0})).

\item The gluing functor
$\rho=i^{\ast}j_{\ast}:\mathcal{U}\rightarrow\mathcal{F}$ is left
exact and accessible {\rm(\hyperrefext[I][I.def9.2]{9.2})}, or
equivalently {\rm(\hyperrefext[I][I.sec8]{8})} it admits a pro-adjoint
$\sigma:\Pro\mathcal{F}\rightarrow\Pro\mathcal{U}$.
\end{enumerate}
\end{enonce*}

\noindent \emph{Proof}: a) follows from \hyperref[IV.prop9.3.3]{9.3.3} c). For b), it suffices
to use the definition \eqhyperref[IV.eq9.3.2.2]{9.3.2.2} of
$X_{\complement U}$, taking into account that
$i_{\ast}i^{\ast}X=X_{\complement U}$
(\hyperref[IV.prop9.3.4]{9.3.4}). Assertion c) follows from
\hyperref[IV.prop9.3.3]{9.3.3} b). The functors $i^{\ast}$ and
$j_{\ast}$ are left exact, and hence so is $\rho$. The functor
$j_{\ast}$ is a direct-image functor for a morphism of topoi and is
therefore accessible (\hyperrefext[I][I.prop9.5]{9.5}); note that a
topos is an accessible category
(\hyperrefext[I][I.rem9.11.3]{9.11.3}). Since the functor $i^{\ast}$
commutes with inductive limits (it is an inverse-image functor), the
functor $\rho$, being a composite of two accessible functors, is
accessible.

\begin{enonce*}{Proposition 9.4.2}\etiquette[9.4.2]{IV.prop9.4.2}
Let $E'$ be a topos. Then the functor $h\mapsto i\circ h$,
$$
\Hhomtop(E',\mathcal{F})\longrightarrow \Hhomtop(E',E),
$$
is fully faithful, and its essential image consists of the morphisms
of topoi $g:E'\rightarrow E$ such that $g^{\ast}(U)$ is an initial
object $\phi_{E'}$ of $E'$.
\end{enonce*}

By \hyperref[IV.prop9.1.4]{9.1.4}, the functor in question is fully
faithful, and $g$ belongs to its essential image if and only if, for
every object $X'$ of $E'$, the object $g_{\ast}(X')$ of $E$ belongs
to $\mathcal{F}$, \ie if
\begin{equation*}
j^{\ast}(g_{\ast}(X'))\simeq e_{\mathcal{U}}\text{ for every }X'\in
\Ob E'{\quoi}.\tag{*}
\end{equation*}
Put $U'=g^{\ast}(U)$ and consider the diagram of topoi
\[
\xymatrix@C=1.5cm{
E'_{/U'}\ar[r]^{j'} \ar[d]_{g_U}& E'\ar[d]^{g}\\
E_{/U}   \ar[r]^{j}             & E\quad , }
\]
There\pageoriginale is plainly a canonical isomorphism
$$
j^{\ast}(g_{\ast}(X'))\simeq g_{U\ast}({j'}^{\ast}(X')){\quoi}.
$$
The functor ${j'}^{\ast}$ is also plainly essentially surjective, so
condition (*) is equivalent to
\begin{equation*}
g_{U\ast}(Y')\simeq e_{\mathcal{U}}\text{ for every }Y'\in\Ob
E'_{/U'}{\quoi}.\tag{**}
\end{equation*}
Using the adjunction between $g_{U\ast}$ and $g_U^{\ast}$, one sees
immediately that this means that $g_U^{\ast}(Y)$ is an initial object
of $E'_{/U'}$ for every object $Y$ of $E_{/U}$; or, using the fact
that the initial object of $E'_{/U'}$ is strict
(\hyperrefext[II][II.def4.5]{4.5}), that it suffices for this to hold
for $Y=e_{\mathcal{U}}$, the terminal object of $\mathcal{U}$. But
this also means that the terminal object $U'$ of $E'_{/U'}$ is
initial, which is manifestly equivalent to $U'$ being an initial
object $\phi_{E'}$ of $E'$, \cqfd

\begin{enonce*}{Corollary 9.4.3}\etiquette[9.4.3]{IV.cor9.4.3}
Let $g:E'\rightarrow E$ be a morphism of topoi, let
$U'=g^{\ast}(U)$, let $\mathcal{U}'=E'_{/U'}$, and let
$\mathcal{F}'$ be the closed subtopos of $E'$ complementary to the
open $U'$; finally, let $i':\mathcal{F}'\rightarrow E'$ be the
inclusion morphism. Then the morphism of topoi
$g\circ i':\mathcal{F}'\rightarrow E$ factors, up to isomorphism,
through a morphism of topoi
$g_{\mathcal{F}}:\mathcal{F}'\rightarrow\mathcal{F}$, and the
corresponding diagram of topoi
\[\vcenter{
\xymatrix@=.5cm{ \Ff' \ar[rr]^{i'}\ar[dd]_{g_{\Ff}}  &&
E'\ar[dd]^{g}
\\
                          &            &  \\
\Ff  \ar[rr]^(.4){i}  &  \ar@/^.5pc/[ru]  &   E }}\leqno{(9.4.3.1)}
\]\etiquette[9.4.3.1]{IV.eq9.4.3.1}
is $2$-Cartesian {\rm(\cf \hyperref[IV.prop5.11]{5.11}).}
\end{enonce*}

This follows formally from \hyperref[IV.prop9.4.2]{9.4.2} and the
definitions.

With the terminology introduced in \hyperref[IV.exer9.1.6]{9.1.6}
d), one may therefore say in particular that the inverse image under
a morphism of topoi $g:E'\rightarrow E$ of a \emph{closed} subtopos
$\mathcal{F}$ of $E$ is a \emph{closed} subtopos $\mathcal{F}'$ of
$E'$ (namely the closed subtopos complementary to the open
$U'=g^{\ast}(U)$, where $U$ is the open of $E$ complementary to
$\mathcal{F}$).

\begin{enonce*}{Corollary 9.4.4}\etiquette[9.4.4]{IV.cor9.4.4}
Let $E$ be a topos, $\mathcal{U}$ an open subtopos of $E$, and
$\mathcal{F}$ the complementary closed subtopos.
\begin{enumerate}
\renewcommand{\theenumi}{\alph{enumi}}
\renewcommand{\labelenumi}{\theenumi)}
\item One has\pageoriginale
$$
\mathcal{U}\cap\mathcal{F}\approxeq\Top(∅)\quav
\mathcal{U}\vee\mathcal{F}=E{\quoi},
$$
where $\vee$ denotes the $\Sup$ in the set of all subtopoi of $E$
{\rm(\hyperref[IV.rem9.1.2]{9.1.2} c))}. (In the terminology
introduced in \hyperref[IV.exer9.1.13]{9.1.13}, $\mathcal{U}$ and
$\mathcal{F}$ are complementary subtopoi of $E$.)

\item $\mathcal{U}$ (\resp $\mathcal{F}$) is the largest among the
subtopoi $E'$ of $E$ such that
$E'\cap\mathcal{F}\approxeq\Top(∅)$ (\resp
$\mathcal{U}\cap E'\approxeq\Top(∅)$).
\end{enumerate}
\end{enonce*}

\noindent \emph{Proof.}
\begin{enumerate}
\renewcommand{\theenumi}{\alph{enumi}}
\renewcommand{\labelenumi}{\theenumi)}
\item The first relation amounts to saying that if $E'$ is a topos
and $g:E'\rightarrow E$ is a morphism of topoi factoring, up to
isomorphism, through both $\mathcal{U}$ and $\mathcal{F}$, then
$E'\approxeq\Top(∅)$. By \hyperref[IV.subsubsec9.2.5]{9.2.5} and
\hyperref[IV.prop9.4.2]{9.4.2}, the hypothesis means that
$g^{\ast}(U)$ is both an initial and a terminal object of $E'$, which
proves the assertion. For the second relation, note that by
\hyperref[IV.exer9.1.7.2]{9.1.7.2} e) it is equivalent to
\hyperref[IV.prop9.4.1]{9.4.1} c).

\item Suppose that $E'\cap\mathcal{F}\approxeq\Top(∅)$. By \hyperref[IV.cor9.4.3]{9.4.3},
$E'\cap\mathcal{F}$ is equivalent to the subtopos of $E'$
complementary to $g^{\ast}(U)$, where $g:E'\rightarrow E$ is the
inclusion morphism. Thus the condition means that
$U'=g^{\ast}(U)$ is a terminal object of $E'$, \ie
(\hyperref[IV.subsubsec9.2.5]{9.2.5}) that
$E'\subset\mathcal{U}$. Suppose that
$\mathcal{U}\cap E'\approxeq\Top(∅)$. By
\hyperref[IV.prop5.11]{5.11}, $\mathcal{U}\cap E'$ is equivalent to
$E'_{/U'}$, so the condition means that $U'$ is an initial object of
$E'$, \ie (\hyperref[IV.prop9.4.2]{9.4.2}) that
$E'\subset\mathcal{F}$.
\end{enumerate}

\begin{enonce*}{Corollary 9.4.5}\etiquette[9.4.5]{IV.cor9.4.5}
Let $(\mathcal{F}_i)_{i\in I}$ be a family of closed subtopoi of the
topos $E$, and for every $i\in I$ let $U_i$ be the open of $E$
complementary to $\mathcal{F}_i$. Then the intersection of the
$\mathcal{F}_i$ as subtopoi of $E$
(\hyperref[IV.prop9.1.3]{9.1.3}) is a closed subtopos whose
complementary open is $U=\Sup_iU_i$.
\end{enonce*}

This follows formally from the universal property of the
intersection as a 2-product and from
\hyperref[IV.prop9.4.2]{9.4.2}, taking into account that for a
morphism of topoi $g:E'\rightarrow E$,
$g^{\ast}(U)=\Sup_i g^{\ast}(U_i)$ is an initial object if and only
if all the $g^{\ast}(U_i)$ are initial.

Similarly one proves:

\begin{enonce*}{Corollary 9.4.6}\etiquette[9.4.6]{IV.cor9.4.6}
With the notation of \hyperref[IV.cor9.4.5]{9.4.5}, suppose that $I$ is finite. Then the
supremum\pageoriginale of the subtopoi $\mathcal{F}_i$ of $E$
{\rm(\hyperref[IV.rem9.1.2]{9.1.2} c))} is a closed subtopos of $E$,
whose complementary open is $\bigcap_iU_i$.
\end{enonce*}


% END INCLUDED SGA4_English_current_buildable_source_20260906/source/04/drafts/SGA4_IV_LINES_6598_6787_English.texfrag


% BEGIN INCLUDED SGA4_English_current_buildable_source_20260906/source/04/drafts/SGA4_IV_LINES_6788_7070_English.texfrag
\begin{enonce*}[remark]{Exercise 9.4.7}\etiquette[9.4.7]{IV.exer9.4.7}
(\emph{Gluing of subtopoi}). Let $E$ be a topos.
\begin{enumerate}
\renewcommand{\theenumi}{\alph{enumi}}
\renewcommand{\labelenumi}{\theenumi)}
\item Let $E'$ be a subtopos of $E$, let $(S_i)_{i\in I}$ be a family
of objects of $E$ covering the terminal object, and, for every $i\in I$,
let $S'_i$ be the inverse image of $S_i$ in $E'$, and let $E'_i\approxeq
E'_{/S'_i}$ be the subtopos of $E_i=E_{/S_i}$ which is the inverse image,
under the localization morphism $E_i\rightarrow E$, of the subtopos $E'$ of
$E$ (\hyperref[IV.exer9.1.6]{9.1.6} d) and \hyperref[IV.prop5.11]{5.11}). Show that, for
an object $X$ of $E$, $X$ belongs to $E'$ if and only if, for
every $i\in I$, the inverse image $X| S_i$ of $X$ in $E_i$ belongs
to $E'_i$.

\item Suppose that $E$ is of the form $C^{\sim}$, where $C$ is a
$𝒰$-site. Show that the presheaf on $C$
\begin{center}
$S\mapsto$ set of subtopoi of the topos
$(C_{/S})^{\sim}\approxeq E_{/\varepsilon(S)}$
\end{center}
is a sheaf. For arbitrary $E$, deduce that there
exists an object of $E$ which represents the functor
$$
\text{$S\mapsto$ set of subtopoi of the topos $E_{/S}$;}
$$

\item With the notation of a), show that the subtopos
$E'$ of $E$ is an open subtopos (\resp a closed subtopos)
if and only if, for every $i\in I$, the same is true of the
induced subtopos $E'_i$ of $E_i$.
\end{enumerate}
\end{enonce*}

\begin{enonce*}[remark]{Exercise 9.4.8}\etiquette[9.4.8]{IV.exer9.4.8}
(\emph{Interior, exterior, closure, and boundary of a
subtopos.}) Let $E$ be a topos, $E'$ a subtopos, and
$i:E'\rightarrow E$ the inclusion morphism.
\begin{enumerate}
\renewcommand{\theenumi}{\alph{enumi}}
\renewcommand{\labelenumi}{\theenumi)}
\item Show that $i_{\ast}(\phi_{E'})$ is the largest open $U$
of $E$ such that $E'| \mathcal{U} \approxeq \Top(∅)$, where
$\mathcal{U}$ is the open subtopos of $E$ defined by $U$. This
$U$, or $\mathcal{U}$, is called the \emph{exterior} of the
subtopos $E'$ of $E$, and the closed subtopos of $E$ complementary
to $U$ is called the \emph{closure} of $E'$.

\item Show that there exists a largest open $V$ of $E$ such that the
corresponding open subtopos $\mathcal{V}$ of $E$ is contained
in $E'$. (Use \hyperref[IV.exer9.2.7]{9.2.7}.) This $V$, or
$\mathcal{V}$, is called the \emph{interior} of the subtopos $E'$ of $E$. The
opens $U$ and $V$ of $E$ are disjoint $(U\cap V=\phi_E)$; the
closed subtopos of $E$ complementary\pageoriginale to
$\Sup(U,V)=\mathcal{U}\amalg \mathcal{V}$ is called the
\emph{boundary} of the subtopos $E'$ of $E$.

\item The subtopos $E'$ is an open subtopos (\resp a closed subtopos) of $E$ if
and only if it is equal to its interior (\resp if it is
equal to its closure). The boundary of $E'$ is empty if and
only if $E'$ is both an open and a closed subtopos of
$E$, or equivalently, if it is the open subtopos of $E$ defined by an
open $V$ of $E$ which is a \emph{direct summand} of the terminal object
$e$, \ie such that there exists a subobject $U$ of $e$ with $e\simeq
U\amalg V$. The exterior of $E'$ is ``empty'' \ie the
closure is equal to $E$, if and only if
$i:E'\rightarrow E$ is dominant (\hyperref[IV.exer8.8]{8.8}).

\item Let $S$ be an object of $E$, let $F=E_{/S}$ be the induced topos, and let $F'\simeq
F_{/S'}$ be the subtopos of $F$ induced by the subtopos $E'$ of $E$.
Show that the exterior (\resp the closure, \resp
the interior, \resp
the boundary) of $F'$ is induced by the exterior (\resp\ldots) of
$E'$.
\end{enumerate}
\end{enonce*}

\begin{enonce*}[remark]{Exercise 9.4.9}\etiquette[9.4.9]{IV.exer9.4.9}
(\emph{Locally closed subtopoi of a topos.}) A subtopos $F$
of a topos $E$ is called a \emph{locally closed subtopos} of $E$
if it is the intersection of an open subtopos and a
closed subtopos.
\begin{enumerate}
\renewcommand{\theenumi}{\alph{enumi}}
\renewcommand{\labelenumi}{\theenumi)}
\item Prove that the intersection of a finite family of locally closed subtopoi
of $E$ is locally closed. If $g:E'\rightarrow
E$ is a morphism of topoi, prove that the inverse image
(\hyperref[IV.exer9.1.6]{9.1.6} d)) of a locally closed subtopos of $E$ is
a locally closed subtopos of $E'$.

\item Let $X$ be a topological space, and let $E=\Top(X)$. For every
subset $X'$ of $X$, consider the subtopos $T(X')$ of $E$ which is the essential
image of $\Top(X')$ in $\Top(X)$. Show that if $X'$ is
a locally closed subset, $T(X')$ is a locally closed subtopos
of $E$, the converse being true if, in addition, $X$ and
$X'$ are assumed sober. Prove that $X'\mapsto T(X')$ defines an order
isomorphism between the set of locally closed subsets
$X'$ of $X$ and the set of locally closed subtopoi of~$E$.

\item Let $F$ be a subtopos of $E$. Show that every open subtopos
of $F$\pageoriginale is induced by an open subtopos of $E$.
(Use \hyperref[IV.exer9.1.8]{9.1.8} c).) Show that $F$ is a locally closed
subtopos of $E$ if and only if it is an open subtopos
of its closure (\hyperref[IV.exer9.4.8]{9.4.8} a)). Deduce, using
\hyperref[IV.exer9.4.8]{9.4.8} d), that the property for $F$
of being a
locally closed subtopos of $E$ is a local property
on $E$.

\item Let $F$ be a locally closed subtopos of $E$. Show that,
among all the ways of writing $F$ as the intersection of an
open subtopos $\mathcal{U}$ and a closed subtopos
$\mathcal{F}$ of $E$, there is one for which $\mathcal{U}$ is the largest
possible, and $\mathcal{F}$ the smallest possible. (Take
$\mathcal{F} =$ the closure of $F$, and $\mathcal{U} =$ the largest
open subtopos inducing on $\mathcal{F}$ the subtopos $F$.)
Show that the formation of $\mathcal{U}$ and $\mathcal{F}$ is compatible
with localization on $E$.
\end{enumerate}
\end{enonce*}

\begin{enonce*}[remark]{Exercise 9.4.10}\etiquette[9.4.10]{IV.exer9.4.10}
Develop the notions of a \emph{constructible, locally
constructible, quasi-constructible, and locally quasi-constructible subtopos},
following the model of the familiar notions for topological
spaces (EGA ${\rm 0_{III}}$ \hyperref[IV.subsec9.1]{9.1}, EGA IV 10.1).
For a topos of the form $\Top(X)$, one recovers a
bijection between the set of constructible subsets (\resp\ldots)
of $X$, and the set of constructible subtopoi (\resp\ldots) of
$\Top(X)$.
\end{enonce*}

\subsection{The gluing theorem}\etiquette[9.5]{IV.subsec9.5}

\subsubsection{}\etiquette[9.5.1]{IV.subsubsec9.5.1}
Let $f:A\rightarrow B$ be a functor between two categories. Denote
by $(B,A,f)$ the following category: the objects of $(B,A,f)$ are
the triples
$$
(X,Y,u)
$$
where $X$ is an object of $B$, $Y$ an object of $A$, and $u$ a morphism
from $X$ to $f(Y)$; the morphisms between two objects $(X,Y,u)$ and
$(X',Y',u')$ are the pairs $(m,m')$, where $m$ is a morphism from
$X$ to $X'$ and $m'$ a morphism from $Y$ to $Y'$, such that the
following diagram is commutative:
\[
\xymatrix@C=1.5cm@R=1cm{
X \ar[r]^{u} \ar[d]_{m} & **[r]f(Y) \ar[d]^{f(m')}\\
X' \ar[r]^{u'}& **[r]f(Y') }
\]
the\pageoriginale composition of morphisms is defined in the
evident way.

\subsubsection{}\etiquette[9.5.2]{IV.subsubsec9.5.2}
Note that the category $(B,A,f)$ depends functorially
on the system of data $A$, $B$, $f$, in a sense whose precise
formulation we leave to the reader. In particular, if $f'$ is a
second functor from $A$ to $B$, then every morphism of functors
$$
f\xrightarrow{u} f'
$$
defines a functor $(B,A,f)\rightarrow (B,A,f')$, which is an
isomorphism if $u:f\rightarrow f'$ is one.

\subsubsection{}\etiquette[9.5.3]{IV.subsubsec9.5.3}
Return to the situation of \hyperref[IV.subsec9.3]{9.3}, and consider the
category $(\mathcal{F},\mathcal{U}, \rho)$ defined in
\hyperref[IV.subsubsec9.5.1]{9.5.1}. To each object $X$ of $E$ there corresponds an
object $(i^{\ast}X, j^{\ast}X, h(X):i^{\ast}X\rightarrow \rho
j^{\ast}X)$ of $(\mathcal{F}, \mathcal{U},\rho)$, where $h(X)$ is obtained
by applying the functor $i^{\ast}$ to the adjunction morphism
$X\rightarrow j_{\ast}j^{\ast}(X)$ (one has $\rho=i^{\ast}j_{\ast}$
\sisi{(9.10.)}{\hyperref[IV.eq9.3.6.3]{9.3.6.3}}). Thus this defines a functor
\begin{equation*}
\Phi:E\longrightarrow (\mathcal{F}, \mathcal{U},
\rho){\quoi}.\tag{9.5.3.1}
\end{equation*}\etiquette[9.5.3.1]{IV.eq9.5.3.1}

\begin{enonce*}{Theorem 9.5.4}\etiquette[9.5.4]{IV.theo9.5.4}
\begin{enumerate}
\renewcommand{\theenumi}{\alph{enumi}}
\renewcommand{\labelenumi}{\theenumi)}
\item Let $E$ be a topos, let $\mathcal{U}$ be an open subtopos of $E$, and let
$\mathcal{F}$ be the complementary closed subtopos. The functor
$$
\rho=i^{\ast}j_{\ast}:\mathcal{U} \longrightarrow \mathcal{F}
$$
is left exact and accessible, and the functor $\Phi$ $\hyperref[IV.eq9.5.3.1]{(9.5.3.1)}$
is an equivalence of categories.

\item Conversely, let $\mathcal{U}$ and $\mathcal{F}$ be two topoi, and let
$\rho:\mathcal{U}\rightarrow \mathcal{F}$ be a gluing functor
$(\hyperref[IV.subsubsec9.3.6]{9.3.6})$ (\ie accessible {\rm (\hyperrefext[I][I.def9.2]{9.2})}
and left exact). The category $E=(\mathcal{F},\mathcal{U},\rho)$
is a topos. Let $\phi_{\mathcal{F}}$ be an initial object of
$\mathcal{F}$, let $e_{\mathcal{U}}$ be a terminal object of $\mathcal{U}$, and
let $U=(\phi_{\mathcal{F}},e_{\mathcal{U}},
\phi_{\mathcal{F}}\rightarrow \rho(e_{\mathcal{U}}))\in \Ob E$. The
functor $X\mapsto (\phi_{\mathcal{F}},
X,\phi_{\mathcal{F}}\rightarrow \rho(X))$ induces an equivalence of
$\mathcal{U}$ with $E_{/U}$, hence with the open subtopos
$\mathcal{U}'$ of $E$ defined by $U$. The functor $Y\mapsto
(Y,e,Y\rightarrow \rho(e))$ is a closed embedding
$(\hyperref[IV.subsubsec9.3.5]{9.3.5})$, \ie it induces an equivalence of
$\mathcal{F}$ with a closed subtopos $\mathcal{F}'$ of $E$. The
subtopoi $\mathcal{U}'$ and $\mathcal{F}'$\pageoriginale are
complementary. Let $\rho':\mathcal{U}'\rightarrow \mathcal{F}'$
be the gluing functor. The functors $\rho$ and $\rho'$ are
compatible with the equivalences described above.
\end{enumerate}
\end{enonce*}

\begin{enonce*}{Definition 9.5.5}\etiquette[9.5.5]{IV.def9.5.5}
Let $\mathcal{U}$ and $\mathcal{F}$ be two topoi, and let
$\rho:\mathcal{U}\rightarrow \mathcal{F}$ be a gluing functor
$(\hyperref[IV.subsubsec9.3.6]{9.3.6})$. The topos $(\mathcal{F},
\mathcal{U},\rho)$ is called the topos constructed from
$\mathcal{U}$ and $\mathcal{F}$ by gluing by means of the
functor $\rho$.
\end{enonce*}

\setcounter{subsubsection}{5}
\subsubsection{}\etiquette[9.5.6]{IV.subsubsec9.5.6}
\emph{Proof of} \hyperref[IV.theo9.5.4]{9.5.4} a) \footnote[1]{For a
gluing statement more general than \hyperref[IV.def9.5.5]{9.5.5}, \cf
Exercise\stepcounter{cnote}\textsuperscript{\arabic{cnote}}%
\textup{ [N.D.E.: This exercise does not seem to exist.]}.} The first assertion merely recalls
\hyperref[IV.prop9.4.1]{9.4.1} d). The pair of functors $(i^{\ast},j^{\ast})$
is conservative, or equivalently, faithful (\hyperref[IV.prop9.4.1]{9.4.1} c)). A
fortiori, the functor $\Phi$ is faithful. Let us show that it is
fully faithful. It follows from \hyperref[IV.prop9.3.3]{9.3.3} b) that, for
every object $X$ of $E$, the following diagram is cartesian:
\[\vcenter{
\xymatrix@C=2.5cm{
X  \ar[r] \ar[d]               & j_{\ast}j^{\ast}X \ar[d]\\
i_{*}i^{*}X   \ar[r]     & i_{\ast}i^{\ast}j_{\ast}j^{\ast}X ;
}}\leqno{(*)}
\]
Let $X$ and $Y$ be two objects of $E$, and suppose that $(m,m')$:
\[\vcenter{
\xymatrix@C=2.75cm@R=1cm{
 i^{\ast}X  \ar[r] \ar[d]_{m}               & i^{\ast}j_{\ast}j^{\ast}X \ar[d]_{i^{\ast}j_{\ast}m'}\\
i^{*}Y   \ar[r]     & i^{\ast}j_{\ast}j^{\ast}Y ; }}
\]
is a morphism from $\Phi(X)$ to $\Phi(Y)$. Using the
functoriality of the fiber product and the diagrams $(\ast)$, one
obtains a morphism $w:X\rightarrow Y$ which makes the
following diagrams commutative:
\begin{displaymath}
\xymatrix@C=.5cm@R=1.25cm{
X \ar@{-}[rr]^{w} \ar[d]  &             &Y \ar[d]\\
 j_{\ast}j^{\ast}X \ar[rr]^{j_{\ast}m'}       &   &
j_{\ast}j^{\ast}Y , }\qquad\qquad \xymatrix@C=.5cm@R=1.25cm{
X \ar@{-}[rr]^{w} \ar[d]& &Y \ar[d]\\
 i_{\ast}i^{\ast}X \ar[rr]^{i_{\ast}m}&   & i_{\ast}i^{\ast}Y
}
\end{displaymath}
By\pageoriginale applying the functors $j^{\ast}$
and $i^{\ast}$ respectively, one obtains $j^{\ast}w=m'$ and $i^{\ast}w=m$. Let us
now show that the functor $\Phi$ is essentially surjective.
Let $(Y,X,u:Y\rightarrow i^{\ast}j_{\ast}(X))$ be an object of
$(\mathcal{F},\mathcal{U},i^{\ast} j_{\ast})$. From it one obtains a
diagram:
\[
\xymatrix@C=2cm{
                             & j_{\ast}X \ar[d]\\
i_{\ast}Y \ar[r]^{i_{\ast}u} & i_{\ast}i^{\ast}j_{\ast}X }
\]
Hence, on taking the fiber product, one obtains a cartesian diagram:
\[\vcenter{
  \xymatrix@C=2.75cm@R=1.25cm{
W \ar[r] \ar[d]& j_{\ast}X \ar[d]\\
i_{\ast}Y \ar[r]^{i_{\ast}u}& i_{\ast}i^{\ast}j_{\ast} X\\
}}\leqno{(**)}
\]

We leave it to the reader to show that $\Phi(W)$ is
isomorphic to $(Y,X,u:Y\rightarrow i^{\ast}j_{\ast}(X))$ and that the
diagram (${\ast}{\ast}$) is isomorphic to the diagram ($\ast$)
associated with $W$.

\hypertarget{sga4.exp.IV.subsubsection.9.5.7}{}\subsubsection{}

% END INCLUDED SGA4_English_current_buildable_source_20260906/source/04/drafts/SGA4_IV_LINES_6788_7070_English.texfrag


% BEGIN INCLUDED SGA4_English_current_buildable_source_20260906/source/04/drafts/SGA4_IV_LINES_7071_7243_English.texfrag
\emph{Proof} of \hyperref[IV.theo9.5.4]{9.5.4}
\sisi{1)}{b)}. We shall confine ourselves to showing that
$(\mathcal{F},\mathcal{U},\rho)$ is a topos. The proof of the other
assertions is left to the reader. We verify properties a), b), c),
d) of \hyperref[IV.subsubsec1.1.2]{1.1.2}.
\begin{enumerate}
\renewcommand{\theenumi}{\alph{enumi}}
\renewcommand{\labelenumi}{\theenumi)}
\item \emph{Finite projective limits are representable}: This is
clear because the categories $\mathcal{U}$ and $\mathcal{F}$ are
topoi and $\rho$ commutes with finite projective limits.

\item \emph{Direct sums are representable, disjoint, and universal}:
Let $(Y_i\xrightarrow{u_i}\rho(X_i))$, $i\in I$, be a family of
objects of $(\mathcal{F},\mathcal{U},\rho)$. By the definition of
inductive limits, there is a canonical morphism
$$
\coprod_{i\in I}\rho(X_i)\longrightarrow\rho(\coprod_{i\in I}X_i),
$$
and hence,\pageoriginale by composing with $\coprod u_i$, a morphism
$$
n:\coprod_{i\in I}Y_i\longrightarrow\rho(\coprod_{i\in I}X_i){\quoi}.
$$
One verifies immediately that this last object is the direct sum of
the given family and that this direct sum is disjoint and universal.

\item \emph{Equivalence relations are effective and universal}:

Let
\[
\xymatrix@C=2cm@R=1.5cm{
R_1 \ar@<3pt>[r] \ar@<-3pt>[r] \ar[d]& Y \ar[d]^{w}\\
\rho(R_2) \ar@<3pt>[r]^{f(u)} \ar@<-3pt>[r]_{f(v)} & \rho(X) }
\]
be an equivalence relation on the object
$Y\xrightarrow{w}\rho(X)$. Then
$\xymatrix@C=.6cm{R_1\ar@<-2pt>[r]\ar@<2pt>[r]&Y}$ is an equivalence
relation on $Y$ in $\mathcal{F}$, and
$\xymatrix{R_2\ar@<2pt>[r]^{u}\ar@<-2pt>[r]_{v}&X\\}$ is an
equivalence relation on $X$ in $\mathcal{U}$. The quotients
$Y/R_1$, $\rho(X)/\rho(R_2)$, and $X/R_2$ are representable in
$\mathcal{F}$ and $\mathcal{U}$ because these categories are topoi.
Moreover, by the definition of inductive limits, the canonical
morphism $\rho(X)\rightarrow\rho(X/R_2)$ factors through
$\rho(X)/\rho(R_2)$. By functoriality of inductive limits, one obtains
a canonical morphism
$Y/R_1\xrightarrow{t}\rho(X/R_2)$. One verifies that this last object
is the quotient of $Y\xrightarrow{w}\rho(X)$ by the equivalence
relation in question, that this quotient is effective (I~10), and
that all these properties are preserved under base change.

\item $(\mathcal{F},\mathcal{U},\rho)$ \emph{admits a small
generating family}. This property follows from
\sisi{(1.9.2 5)}{\hyperrefext[I][I.cor9.25]{9.25}} (taking into
account that $\rho$ is accessible), by taking
$B=\Delta_1=(0\xrightarrow{f}1)$, $E_0=\mathcal{F}$,
$E_1=\mathcal{U}$, and $f^{\ast}=\rho:E_1\rightarrow E_0$, whence
$\Hhom_B(B,E)=(\mathcal{F},\mathcal{U},\rho)$.

One may also avoid using
\sisi{(1.9.2 5)}{\hyperrefext[I][I.cor9.25]{9.25}} (whose proof as
given was rather arduous!) and use the hypothesis on $\rho$ in the
form that $\rho$ admits a pro-adjoint
$$
\sigma:\Pro(\mathcal{F})\longrightarrow\Pro(\mathcal{U}){\quoi},
$$
(this\pageoriginale interpretation (\hyperrefext[I][I.sec8]{8})
being based on \hyperrefext[I][I.sec9]{9}, whose proof is distinctly
more comprehensible). It suffices to note that if $C'$ (\resp $C''$)
is a generating subcategory of $\mathcal{U}$ (\resp $\mathcal{F}$),
and if, for every $X''\in\Ob C''$, one represents $\sigma(X'')$ as a
projective system $(\sigma(X'')_i)_{i\in I(X'')}$, where $I(X'')$ is
a small filtered ordered set, then the family of objects of $E$ which
are either of the form
$(\phi_{\mathcal{F}},X',\phi_{\mathcal{F}}\rightarrow\rho(X'))$ with
$X'\in\Ob C'$, or of the form
$(X'',\sigma(X'')_i,u_{X''}:X''\rightarrow
\rho(\sigma(X'')_i))$, where $X''\in\Ob C''$, $i\in I(X'')$, and
$u_{X''}$ corresponds under adjunction to the canonical morphism
$\sigma(X'')\rightarrow\sigma(X'')_i$, is a generating family
(plainly a small one) when $C'$ and $C''$ are chosen small. The
verification is indeed immediate and is left to the reader.
\end{enumerate}

\setcounter{subsubsection}{4}\setcounter{paragraph}{7}
\paragraph{}\etiquette[9.5.4.8]{IV.par9.5.4.8}
Use the notation of \hyperlink{sga4.exp.IV.subsubsection.9.5.7}{9.5.7} b), and let us explain how, up to canonical
isomorphisms of functors, the system of functors
\eqhyperref[IV.eq9.3.6.2]{9.3.6.2} can be made explicit when
$E=(\mathcal{F},\mathcal{U},\rho)$. The details verifying the
assertions below (which are essentially mechanical applications of
\hyperref[IV.theo9.5.4]{9.5.4}) are left to the reader.

\setcounter{subsubsection}{8}\setcounter{paragraph}{0}
\paragraph{}\etiquette[9.5.8.1]{IV.par9.5.8.1}
The functor
$$
i^{\ast}:(\mathcal{F},\mathcal{U},\rho)\longrightarrow\mathcal{F}
$$
is given by
$$
i^{\ast}(X,Y,u:X\longrightarrow\rho(Y))=X{\quoi}.
$$

\paragraph{}\etiquette[9.5.8.2]{IV.par9.5.8.2}
The functor
$$
i_{\ast}:\mathcal{F}\longrightarrow(\mathcal{F},\mathcal{U},\rho)
$$
is given by $i_{\ast}(X)=(X,e,X\rightarrow\rho(e))$, where $e$ is
the terminal object of $\mathcal{U}$.

\paragraph{}\etiquette[9.5.8.3]{IV.par9.5.8.3}
The adjunction morphism
$$
\id\longrightarrow i_{\ast}i^{\ast}
$$
is isomorphic to the functorial morphism
\[
\xymatrix@C=2cm{
X \ar[r]^{u} \ar[d]_{id}& \rho(Y) \ar[d]\\
X \ar[r]                & \rho(e) }
\]

\paragraph{}\etiquette[9.5.8.4]{IV.par9.5.8.4}
The\pageoriginale functor
$$
j^{\ast}:(\mathcal{F},\mathcal{U},\rho)\longrightarrow\mathcal{U}
$$
is given by\qquad $j^{\ast}(X,Y,u:X\longrightarrow\rho(Y))=Y$.

\paragraph{}\etiquette[9.5.8.5]{IV.par9.5.8.5}
The functor
$$
j_{!}:\mathcal{U}\longrightarrow(\mathcal{F},\mathcal{U},\rho)
$$
is given by\qquad
$j_{!}(Y)=(\phi_{\mathcal{F}},Y,
\phi_{\mathcal{F}}\longrightarrow\rho(Y))$

\noindent (where $\phi_{\mathcal{F}}$ is the initial object of
$\mathcal{F}$).

\paragraph{}\etiquette[9.5.8.6]{IV.par9.5.8.6}
The functor
$$
j_{\ast}:\mathcal{U}\longrightarrow(\mathcal{F},\mathcal{U},\rho)
$$
is given by\qquad
$j_{\ast}(Y)=(\rho(Y),Y,\id:\rho(Y)\longrightarrow\rho(Y))$.

\paragraph{}\etiquette[9.5.8.7]{IV.par9.5.8.7}
The adjunction morphism
$$
j_{!}j^{\ast}\longrightarrow\id
$$
is isomorphic to the functorial morphism
\[
\xymatrix{
\phi_{\Ff} \ar[r] \ar[d] & \rho(Y) \ar[d]^{\rho(\id)}\\
X \ar[r]^{u}                 & \rho(Y)     &. }
\]

\paragraph{}\etiquette[9.5.8.8]{IV.par9.5.8.8}
The adjunction morphism
$$
\id\longrightarrow j_{\ast}j^{\ast}
$$
is isomorphic to the functorial morphism
\[
\xymatrix{
X  \ar[r]^{u} \ar[d]_{u}& \rho(Y) \ar[d]^{\rho(\id)}\\
\rho(Y) \ar[r]^{\id}& \rho(Y) }
\]

% END INCLUDED SGA4_English_current_buildable_source_20260906/source/04/drafts/SGA4_IV_LINES_7071_7243_English.texfrag


% BEGIN INCLUDED SGA4_English_current_buildable_source_20260906/source/04/drafts/SGA4_IV_LINES_7244_7443_English.texfrag
\begin{enonce*}[remark]{Exercise 9.5.9}\etiquette[9.5.9]{IV.exer9.5.9}
(\emph{Exactness properties of a glued topos.}) Let
$\mathcal{U}$, $\mathcal{F}$ be two topoi, let
$\rho:\mathcal{U}\rightarrow \mathcal{F}$ be a gluing functor, let
$E=(\mathcal{F}, \mathcal{U},\rho)$ be the glued topos, and let
$j:\mathcal{U}\rightarrow E$ and $i:\mathcal{F}\rightarrow E$ be the
canonical morphisms of topoi, so that
$\rho=i^{\ast}j_{\ast}$.
\begin{enumerate}
\renewcommand{\theenumi}{\alph{enumi}}
\renewcommand{\labelenumi}{\theenumi)}
\item Prove that $i^{\ast}$ commutes with (small) products (\ie (\hyperref[IV.rem1.8]{1.8}
or \sisi{\hyperref[IV.rem1.8]{1.8}}{\hyperrefext[I][I.sec8]{8}}) the functor $i^{\ast}$ admits a left adjoint
$i_{!}$) if and only if the same holds for $\rho$, \ie
(\loccit) if and only if $\rho$ admits a left adjoint
$\sigma:\mathcal{F}\rightarrow \mathcal{U}$ (or equivalently if and
only if the pro-adjoint of $\rho$, restricted\pageoriginale to
$\mathcal{F}$ \eqhyperref[IV.eq9.6.3]{9.6.3}, factors, up to isomorphism, through
$\mathcal{U}$): one then has $\sigma=j^{\ast}i_{!}$. Thus the
gluing functors $\rho$ having this property correspond,
up to isomorphism, to the functors $\sigma:\mathcal{F}\rightarrow
\mathcal{U}$ which have a right adjoint, \ie (\hyperref[IV.cor1.6]{1.6}) which
commute with $𝒰$-small inductive limits, or equivalently which are
continuous. When $\mathcal{F}$ is equivalent to a topos
$C^{\sim}$, where $C$ is a $𝒰$-site, these functors
therefore correspond, up to isomorphism, to continuous functors
$C\rightarrow \mathcal{U}$ (\hyperrefext[III][III.prop1.2]{1.2} and \hyperrefext[III][III.prop1.7]{1.7}).

\item Suppose that the functor $i_{!}$ is defined. Show that
$i_{!}$ commutes with a specified type of small projective limits
if and only if $\sigma:\mathcal{F}\rightarrow
\mathcal{U}$ commutes with them. (For necessity, use the fact that $j^{\ast}$
commutes with small projective limits; for sufficiency,
use the equivalent description of $E$ in terms of $\sigma$
as consisting of the triples $(Y,X,u)$ with $Y\in \Ob \mathcal{F}$, $X\in
\Ob \mathcal{U}$, $u:\sigma (Y)\rightarrow X$.)

\item Deduce from a) and b) an example of a closed embedding of topoi
$i:\mathcal{F}\rightarrow E$ such that $i_{!}$ exists, but commutes
neither with fiber products nor with filtered projective limits. (It
suffices to find a continuous functor between two topoi which
has neither of these two exactness properties.)
\end{enumerate}
\end{enonce*}

\begin{enonce*}[remark]{Exercise 9.5.10}\etiquette[9.5.10]{IV.exer9.5.10}
(\emph{Gluing of topoi of the form $\widehat{C}$.}) Let $C$
be a small category, and let $E=\widehat{C}$, which is a
$𝒰$-topos.
\begin{enumerate}
\renewcommand{\theenumi}{\alph{enumi}}
\renewcommand{\labelenumi}{\theenumi)}
\item Recall that the opens $U$ of $E$ correspond to
\emph{sieves} $C'$ of $C$ \sisi{(1.4.1, \hyperref[IV.subsec4.2]{4.2})}{(\hyperrefext[I][I.def4.1]{4.1}, \hyperref[IV.subsubsec8.4.4]{8.4.4})}, with $E_{/U}$
then being equivalent to $\widehat{C'}$, and the functor
$j^{\ast}:E\rightarrow E_{/U}$ identifying with the restriction functor
\hyperref[IV.subsec5.4]{5.4}, \ie the inclusion morphism $j:E_{/U}\rightarrow E$ being
$\widehat{f}:\widehat{C'}\rightarrow \widehat{C}$
(\hyperref[IV.subsubsec4.6.1]{4.6.1}), where $f:C'\rightarrow C$ is the inclusion.
Let $C''$ be the full subcategory of $C$ complementary to $C'$
(\ie $\Ob C''$ is the complement of $\Ob C'$ in $\Ob C$),
and let $g:C''\rightarrow C$ be the inclusion. Then the closed subtopos
$\mathcal{F}$ of $E=\widehat{C}$, complementary to the open subtopos
defined by $C'$, is canonically equivalent to
$\widehat{C''}$, with the inclusion morphism $i:\mathcal{F}\rightarrow
E$ identifying with $\widehat{g}:\widehat{C''}\rightarrow
\widehat{C}$.

\item Consider\pageoriginale the functor
\begin{equation*}
h:{C'}^0\times C''\longrightarrow (\Ens),\quad
h(X',X'')=\Hom_C(X',X''){\quoi},\tag{9.5.10.1}
\end{equation*}\etiquette[9.5.10.1]{IV.eq9.5.10.1}
and note that $C$, together with its full subcategory $C'$, can be
reconstructed up to isomorphism from the categories $C'$,
$C''$ and the bifunctor $h$ (the latter may be chosen
arbitrarily). (Compare \hyperref[IV.subsec9.7.1]{9.7.1} below.) The datum
of $h$ is equivalent to that of a functor $C''\rightarrow
\widehat{C'}=\Hhom({C'}^0,(\Ens))$, or equivalently (\hyperrefext[III][III.prop1.2]{1.2} and
\hyperrefext[III][III.prop1.7]{1.7}) to that of a functor
\begin{equation*}
\sigma:\widehat{C''}\longrightarrow \widehat{C'}\tag{9.5.10.2}
\end{equation*}\etiquette[9.5.10.2]{IV.eq9.5.10.2}
which is continuous, \ie (\hyperref[IV.cor1.6]{1.6}) commutes with small inductive
limits, or equivalently, admits a right adjoint
\begin{equation*}
\rho:\widehat{C}'\longrightarrow \widehat{C''}{\quoi}.\tag{9.5.10.3}
\end{equation*}\etiquette[9.5.10.3]{IV.eq9.5.10.3}
Show that $\rho$ is precisely the gluing functor
associated with the open subtopos $\widehat{C}'$ of $E=\widehat{C}$.

\item Deduce from a) and b) that if one is given a
gluing functor \eqhyperref[IV.eq9.5.10.3]{9.5.10.3} between two presheaf categories,
and hence a glued topos $E$, the following conditions are
equivalent:
\begin{itemize}
\item[(i)] $E$ is equivalent to a topos of the form $\widehat{C}$.

\item[(ii)] The functor $\rho$ commutes with small products (or equivalently, it
admits a left adjoint \eqhyperref[IV.eq9.5.10.2]{9.5.10.2}).

\item[(iii)] The inclusion morphism $i:\widehat{C''}\rightarrow E$ is ``essential,'' \ie $i^{\ast}$ commutes with products, or equivalently
$i^{\ast}$ admits a left adjoint $i_{!}$.
\end{itemize}
When this is so, describe explicitly a category $C$ such that
$E\approx \widehat{C}$ using the bifunctor \eqhyperref[IV.eq9.5.10.1]{9.5.10.1}
defined by the restriction of \eqhyperref[IV.eq9.5.10.2]{9.5.10.2} to $C''$.

\item Show that a topos obtained by gluing two punctual topoi
need not be equivalent to a topos of the form
$\widehat{C}$.
\end{enumerate}
\end{enonce*}

\begin{enonce*}[remark]{Exercise 9.5.11}\etiquette[9.5.11]{IV.exer9.5.11}
(Morphism from a glued topos to a topos.)
\begin{enumerate}
\renewcommand{\theenumi}{\alph{enumi}}
\renewcommand{\labelenumi}{\theenumi)}
\item Let $\rho \sisi{}{: \mathcal{U}} \rightarrow \mathcal{F}$ be a functor between two
categories, let $E=(\mathcal{F},\mathcal{U},\rho)$\pageoriginale be the
``glued'' category of \hyperref[IV.subsubsec9.5.1]{9.5.1}, and let $F$ be an
arbitrary category. Consider the functor $f\mapsto \rho\circ f$:
$$
\bar{\rho}:\bar{\mathcal{U}}=\Hhom(F,\mathcal{U})\longrightarrow
\bar{\mathcal{F}}=\Hhom(F,\mathcal{F}){\quoi},
$$
and define an equivalence of categories
$$
\Hhom(F,E)\xrightarrow{\approx}(\bar{\mathcal{F}},\bar{\mathcal{U}},\bar{\rho}){\quoi}.
$$

\item Let $f:F\rightarrow E$ be a functor, and consider its composites
with the canonical functors $j^{\ast}:E\rightarrow \mathcal{U}$ and
$i^{\ast}:E\rightarrow \mathcal{F}$. Show that $f$ commutes with a
specified type of inductive limits (assumed to be representable in
$\mathcal{U}$ and $\mathcal{F}$, and hence in $E$) if and only if
$j^{\ast}f$ and $i^{\ast}f$ commute with them. Show that if $\rho$
commutes with a specified type of projective limits (which is then
representable in $E$), then $f$ commutes with them if and only if
the same holds for $j^{\ast}f$ and $i^{\ast}f$.

\item Now suppose that $\mathcal{U}$ and $\mathcal{F}$ are
topoi, and that
$\rho$ is a gluing functor. Let $F$ be a topos and
$f:F\rightarrow E$ a functor. Conclude from b) that $f$ is an
inverse-image functor associated with a morphism of topoi $E\rightarrow
F$ if and only if the same holds for the functors
$j^{\ast}f:F\rightarrow \mathcal{U}$ and $i^{\ast}f:F\rightarrow
\mathcal{F}$. Conclude from this and from a) that there is an equivalence
between the category $\Hhomtop(E,F)$ and the category of triples
$(f',f'',u)$, where $f'$ (\resp $f''$) is an object of
$\Hhomtop(\mathcal{U}, F)$ (\resp $\Hhomtop(\mathcal{F}, F)$), and where
$u:{f''}^{\ast}\rightarrow \rho\circ {f'}^{\ast}$ is a
homomorphism of functors $F\rightarrow \mathcal{F}$: if $f$ is an
object of $\Hhomtop(E,F)$ and $(f',f'',u)$ is the corresponding triple,
then $f'=f\circ j$, $f''=f\circ i$.

\item Specify in what sense c) may be said to provide a
characterization, up to equivalence (in a 2-category of topoi),
of the topos $E$ in terms of the triple $(\mathcal{U},\mathcal{F},\rho)$.
\end{enumerate}
\end{enonce*}

\begin{enonce*}[remark]{Exercise 9.5.12}\etiquette[9.5.12]{IV.exer9.5.12}
(\emph{Gluing of two topological spaces.})
\begin{enumerate}
\renewcommand{\theenumi}{\alph{enumi}}
\renewcommand{\labelenumi}{\theenumi)}
\item Let $X$ be a topological space, let $\rho:\Top(X)\rightarrow
(\Ens)=\Top(pt)$ be a gluing functor, \ie (\hyperrefext[I][I.sec8]{8}) a
pro-representable functor, let $A\in \Ob \Pro(\Top(X))$ be the object which
pro-represents it, and let $E$ be the topos obtained from $\rho$ by gluing.
Show that $E$ is equivalent to a topos of the form $\Top(Y)$ if
and only if $A$ is isomorphic to an object of $\Pro(\Ouv(X))$.
Conclude\pageoriginale that if $X$ is nonempty, one can find
$\rho$ such that $E$ is not equivalent to a topos of the form
$\Top(Y)$.

\item More generally, consider a gluing functor
$\rho:\Top(X)\rightarrow \Top(Y)$, where $X$, $Y$ are two topological
spaces. Deduce from it a map
$$
\phi:Y\longrightarrow \Ob \Pro(\Top(X))
$$
(the composite of $Y\rightarrow \Ob \Ppoint(\Top(Y))\rightarrow
\Ob\Pro(\Top(Y))$ (\hyperref[IV.subsubsec6.8.5]{6.8.5}) and the pro-adjoint
$\sigma$ of $\rho$). (N.B. for $y\in Y$, $\phi(y)$ identifies with the
\emph{fiber at $y$ of the gluing sheaf} \eqhyperref[IV.eq9.6.4]{9.6.4}
\emph{defined by} $\rho$.) Show that if $E$ is equivalent to a
topos of the form $\Top(Z)$, then $\phi$ takes its values in
the essential image of $\Pro(\Ouv(X))$.
\end{enumerate}
\end{enonce*}

\begin{enonce*}[remark]{Problem}
Is the converse true?
\end{enonce*}

% END INCLUDED SGA4_English_current_buildable_source_20260906/source/04/drafts/SGA4_IV_LINES_7244_7443_English.texfrag


% BEGIN INCLUDED SGA4_English_current_buildable_source_20260906/source/04/drafts/SGA4_IV_LINES_7444_7700_English.texfrag
\subsection{Gluing sheaf}\etiquette[9.6]{IV.subsec9.6}

We saw in \hyperref[IV.theo9.5.4]{9.5.4} that the datum of a topos $E$
equipped with an open $U$ amounts essentially to that of two topoi
$\mathcal{U}$ and $\mathcal{F}$ and a gluing functor
\begin{equation*}
\rho:\mathcal{U}\longrightarrow \mathcal{F}\tag{9.6.1}
\end{equation*}\etiquette[9.6.1]{IV.eq9.6.1}
\ie a functor admitting a pro-adjoint
\begin{equation*}
\sigma:\Pro(\mathcal{F})\longrightarrow
\Pro(\mathcal{U}){\quoi}.\tag{9.6.2}
\end{equation*}\etiquette[9.6.2]{IV.eq9.6.2}
By the sorites on pro-adjoint functors (\hyperrefext[I][I.sec8]{8}), the
functor $\rho$ is determined up to unique isomorphism (and consequently
the glued topos $E=(\mathcal{F},\mathcal{U},\rho)$ is determined up to
equivalence of topoi) once the functor $\sigma$ is known, or,
equivalently, once the functor induced by $\sigma$ is known:
\begin{equation*}
\sigma_{\circ}:\mathcal{F}\longrightarrow
\Pro(\mathcal{U}){\quoi}.\tag{9.6.3}
\end{equation*}\etiquette[9.6.3]{IV.eq9.6.3}
The functor $\sigma$ is determined, up to unique isomorphism, by the
property of commuting with small filtered projective limits and of
extending the functor $\sigma_\circ$. On the other hand, for a given
functor \eqhyperref[IV.eq9.6.3]{9.6.3} to be isomorphic to a pro-adjoint, it
is plainly necessary and sufficient that, for every object $X$ of
$\mathcal{U}$, the contravariant functor
$$
Y\longmapsto \Hom_{\Pro(\mathcal{U})}(\sigma_{\circ}(Y),X)
$$
on\pageoriginale $\mathcal{F}$ be representable, or equivalently be a
sheaf on $\mathcal{F}$ for the canonical topology
(\sisi{\hyperref[IV.subsubsec1.1.2]{1.1.2} (iii)}{\hyperref[IV.theo1.2]{1.2} iii)}). It follows at once that this
also means that the opposite functor
\begin{equation*}
\sigma_{\circ}^{\circ}:\mathcal{F}^{\circ}\longrightarrow
\Pro(\mathcal{U})^{\circ}\subset
\check{\mathcal{U}}=\Hhom(\mathcal{U},(\Ens))\tag{9.6.4}
\end{equation*}\etiquette[9.6.4]{IV.eq9.6.4}
is a \emph{sheaf} on $\mathcal{F}$ with values in
$\Pro(\mathcal{U})^{\circ}$ (\hyperrefext[II][II.def6.1]{6.1}). Thus, \emph{the
datum of a gluing functor} \eqhyperref[IV.eq9.6.1]{9.6.1} \emph{is also
equivalent} (up to isomorphism) \emph{to that of a sheaf on
$\mathcal{F}$ with values in} $\Pro(\mathcal{U})^{\circ}$. The sheaf
thus associated with a gluing functor (or with a situation
$(E,U,\mathcal{F})$ as in (\hyperref[IV.subsec9.3]{9.3})) is called the
\emph{gluing sheaf} associated with the gluing functor $\rho$ under
consideration (\resp with the open $U$ of the topos $E$).

\begin{enonce*}[remark]{Exercise 9.6.5}\etiquette[9.6.5]{IV.exer9.6.5}
Refine \hyperref[IV.theo9.5.4]{9.5.4} into a statement of 2-equivalence of
2-categories between the 2-category consisting of pairs $(E,U)$, where
$E$ is a topos and $U$ an open of the latter, and the 2-category consisting
of triples $(\mathcal{U},\mathcal{F},\rho)$, where $\mathcal{U}$ and
$\mathcal{F}$ are topoi and
$\rho:\mathcal{U}\rightarrow \mathcal{F}$ is a gluing functor (the
explicit description of the 2-categories in question being left to
the reader). Give a variant of this statement involving the
2-category of triples $(\mathcal{U},\mathcal{F}, \phi)$, where $\phi$
is now a sheaf on $\mathcal{F}$ with values in
$\Pro(\mathcal{U})^{\circ}$.
\end{enonce*}

\subsection{Points of a glued topos}\etiquette[9.7]{IV.subsec9.7}

\subsubsection{}\etiquette[9.7.1]{IV.subsec9.7.1}
Let $C$ be a category, and let
$$
u:D'\longrightarrow C\quav v:D''\longrightarrow C
$$
be two fully faithful functors such that the essential images $C'$ and
$C''$ of $u$ and $v$ satisfy
$$
\Ob C'\cap \Ob C''=\phi\quav \Ob C'\cup \Ob C''=\Ob C{\quoi}.
$$
Suppose, in addition, that
$$
X'\in \Ob C'\quav X''\in \Ob C''\Rightarrow
\Hom(X'',X')=\phi{\quoi},
$$
or, equivalently, that
$$
\Hom(v(A''), u(A'))=\phi\mbox{ for } A'\in \Ob D', A''\in \Ob
D''{\quoi}.
$$
It is then immediate that the category $C$ can be reconstructed, up
to canonical\pageoriginale isomorphism, from the categories $C'$ and
$C''$ and the functor
$$
(X',X'')\longmapsto \Hom(X',X''):{C'}^{\circ}\times
C''\longrightarrow (\Ens){\quoi};
$$
similarly, $C$ can be reconstructed up to equivalence (this
equivalence being defined up to unique isomorphism) from $D'$, $D''$,
and the functor
$$
(A',A'')\longmapsto \Hom(u(A'), v(A'')):{D'}^{\circ}\times
D''\longrightarrow (\Ens){\quoi}.
$$
This observation applies in particular to the situation described in
the following proposition, and shows that the category
$\Ppoint(E)$ is determined, up to equivalence, by the categories
$\Ppoint(\mathcal{U})$ and $\Ppoint(\mathcal{F})$ and by a certain
functor
$$
\Ppoint(\mathcal{U})^{\circ}\times
\Ppoint(\mathcal{F})\longrightarrow (\Ens)
$$
derived from the gluing functor $\rho:\mathcal{U}\rightarrow
\mathcal{F}$.

\begin{enonce*}{Proposition 9.7.2}\etiquette[9.7.2]{IV.prop9.7.2}
Let $E$ be a topos, let $U$ be an open of $E$, let
$\mathcal{U}=E_{/U}$, and let $\mathcal{F}$ be the closed subtopos of
$E$ complementary to $U$. Consider the functors $p\mapsto p\circ j$
and $q\mapsto q\circ i$ induced by the inclusion morphisms
$j:\mathcal{U}\rightarrow E$ and $i:\mathcal{F} \rightarrow E$:
\begin{equation*}
u:\Ppoint(\mathcal{U})\longrightarrow \Ppoint(E),\quad
v:\Ppoint(\mathcal{F})\longrightarrow
\Ppoint(E){\quoi}.\tag{9.7.2.1}
\end{equation*}\etiquette[9.7.2.1]{IV.eq9.7.2.1}
Then the following assertions hold:
\begin{enumerate}
\renewcommand{\theenumi}{\alph{enumi}}
\renewcommand{\labelenumi}{\theenumi)}
\item The preceding functors $u$ and $v$ are fully faithful, and every
point of $E$ belongs to the essential image of exactly one of the
functors $u$ and $v$.

\item Let $p$ and $q$ be two points of $E$ such that $p$ (\resp $q$)
belongs to the essential image of $u$ (\resp of $v$). Then
\begin{equation*}
\Hom(q,p)=\phi{\quoi}.\tag{9.7.2.2}
\end{equation*}\etiquette[9.7.2.2]{IV.eq9.7.2.2}

\item Let $p$ (\resp $q$) be a point of $\mathcal{U}$ (\resp of
$\mathcal{F}$). There is then an isomorphism, functorial in $(q,p)$:
\begin{equation*}
\Hom(u(p),v(q))\simeq \Hom(\Pro(\rho)(p),q)\simeq
\Hom(p,\sigma(q)){\quoi},\tag{9.7.2.3}
\end{equation*}\etiquette[9.7.2.3]{IV.eq9.7.2.3}
where
$$
\sigma:\Pro(\mathcal{F})\longrightarrow \Pro(\mathcal{U})
$$
is the functor $\eqhyperref[IV.eq9.6.2]{9.6.2}$, pro-adjoint to the gluing
functor $\rho:\mathcal{U}\rightarrow \mathcal{F}$, and where, in
writing
$$
\rho:\mathcal{U}\longrightarrow \mathcal{F}{\quoi},
$$
we have\pageoriginale identified (up to equivalence) the categories
$\Ppoint(\mathcal{U})$ and $\Ppoint(\mathcal{F})$ with full
subcategories of $\Pro(\mathcal{U})$ and $\Pro(\mathcal{F})$,
respectively (\hyperref[IV.subsubsec6.8.5]{6.8.5}).
\end{enumerate}
\end{enonce*}

\noindent \emph{Proof}
\begin{enumerate}
\renewcommand{\theenumi}{\alph{enumi}}
\renewcommand{\labelenumi}{\theenumi)}
\item The full faithfulness of $u$ and $v$ is known
(\hyperref[IV.prop9.1.4]{9.1.4}). The assertion concerning the essential
images of $u$ and $v$ follows from the criteria
\hyperref[IV.subsubsec9.2.5]{9.2.5} and \hyperref[IV.prop9.2.4]{9.2.4} for a
point $P\rightarrow E$ to factor, up to isomorphism, through
$\mathcal{U}$ \resp $\mathcal{F}$, taking into account the fact that
the punctual topos $P$ has only two opens, namely the initial object
and the terminal object of $P$.

\item In view of a), the assertion means that if the point
$p:P\rightarrow E$ factors through $\mathcal{U}$, then the same holds
for every point $p':P\rightarrow E$ such that
$\Hom(p',p)\neq \phi$. Now, if $\alpha :p'\rightarrow p$, it induces
$p^{\ast}(U)\rightarrow {p'}^{\ast}(U)$, and
then $p^{\ast}(U)=e_p$ implies
${p'}^{\ast}(U)=e_p$, \cqfd

\item By the lemma \hyperref[IV.lem9.7.2.4]{9.7.2.4} below, applied to the
inclusions $j:\mathcal{U}\rightarrow E$ and
$i:\mathcal{F}\rightarrow E$, there is a diagram of functors
commuting up to canonical isomorphisms:
\[
\xymatrix@C=2cm@R=.55cm@M=.18cm{ \Point(\Uu) \ar@{^{(}->}[d]
\ar[r]^{u}& \Point(E) \ar@{^{(}->}[d]&
\Point(\Gg) \ar@{^{(}->}[d] \ar[l]_{v}\\
\Pro(\Uu) \ar[r]^{\Pro(j_{\ast})}  & \Pro(E)   & \Pro(\Gg)
\ar[l]_{\Pro(i_{\ast})} , }
\]
where the vertical arrows denote the canonical fully faithful
functors (\hyperref[IV.subsubsec6.8.5]{6.8.5}). Moreover, since
$i_{\ast}$ admits a left adjoint $i^{\ast}$,
$\Pro(i_{\ast})$ admits the left adjoint $\Pro(i^{\ast})$
\hyperrefext[I][I.sec8.11.5]{8.11.5}, and one finally obtains the functorial isomorphisms
\begin{align*}
\Hom(u(p),v(q)) \simeq & \Hom_{\Pro(E)}(\Pro(j_{\ast})(p),
\Pro(i_{\ast})(q))\\
& \simeq \Hom_{\Pro(\mathcal{F})}(\Pro(i^{\ast})(\Pro(j_{\ast})(p)),q)\\
& \simeq \Hom_{\Pro(\mathcal{F})}(\Pro(i^{\ast}j_{\ast})(p),q),
\end{align*}
which is precisely the first formula in \hyperref[IV.eq9.7.2.3]{(9.7.2.3)}, the second then
being immediate from the definition of $\sigma$ as a pro-adjoint of
$\rho$. This proves \hyperref[IV.prop9.7.2]{9.7.2}, subject to the
following lemma.
\end{enumerate}

\begin{enonce*}{Lemma 9.7.2.4}\etiquette[9.7.2.4]{IV.lem9.7.2.4}
Let $f:F\rightarrow E$ be a morphism of topoi. Then the diagram of
functors
\[\vcenter{
\xymatrix@C=2cm@R=.5cm{
  \Point(F)\ar[r]^{p \mapsto f \circ p} \ar[d] & \Point(E)\ar[d]\\
\Pro(F) \ar[r]^{\Pro(f_{\ast})}     & \Pro(E) }}\leqno{(9.7.2.4.1)}
\]\etiquette[9.7.2.4.1]{IV.eq9.7.2.4.1}
commutes\pageoriginale up to canonical isomorphism (the vertical
arrows denoting the fully faithful functors
$\hyperref[IV.subsubsec6.8.5]{6.8.5}$).
\end{enonce*}

The verification is immediate from the definitions and is left to
the reader.

\begin{enonce*}{Corollary 9.7.3}\etiquette[9.7.3]{IV.cor9.7.3}
Let $E$ be a topos having enough points. Then every closed subtopos
$\mathcal{F}$ of $E$ also has enough points. More precisely, if
$(p_i)_{i\in I}$ is a conservative family of points of $E$, then the
family of points of $\mathcal{F}$ defined by the subfamily
$(p_i)_{i\in J}$ consisting of those $p_i$ which come from points of
$\mathcal{F}$ is conservative. (Compare
$\hyperref[IV.cor6.7.3]{6.7.3}$, $\hyperref[IV.rem6.7.4]{6.7.4}$.)
\end{enonce*}

Indeed, let $f:X\rightarrow Y$ be an arrow of $\mathcal{F}$ which is
carried to an isomorphism by the fiber functors associated with the
points under consideration of $\mathcal{F}$. Since
$i^{\ast}i_{\ast}\simeq \id_{\mathcal{F}}$, it follows that
$i_{\ast}(f):i_{\ast}(X)\rightarrow i_{\ast}(Y)$ is carried to an
isomorphism by the fiber functors associated with the $p_i$ for
$i\in J$; the same is true for the $p_i$ with $i\in I-J$, \ie those
coming from points of $U$, since $j^{\ast}i_{\ast}$ is the constant
functor with value the terminal object. Therefore $i_{\ast}(f)$ is an
isomorphism, and hence so is $f\simeq i^{\ast}i_{\ast}(f)$, \cqfd

\begin{enonce*}[remark]{Remark 9.7.4}\etiquette[9.7.4]{IV.rem9.7.4}
If $E$ is a topos having enough points, it is clear that an open $U$
of $E$ is determined once the full subcategory of $\Ppoint(E)$ which
is the essential image of $\Ppoint(E_{/U})$ is known; in fact, if
$C$ is a full subcategory of $\Ppoint(E)$ defining a conservative
family of points of $E$, it even suffices to know the subcategory
$C_{/U}$ consisting of those elements of $C$ which belong to the
essential image of $\Ppoint(U)$. It follows from this and from
\hyperref[IV.prop9.7.2]{9.7.2} a) that a closed subtopos
$\mathcal{F}$ of $E$ is determined once the essential image of
$\Ppoint(\mathcal{F})$ in $\Ppoint(E)$ is known, or even merely its
intersection with $C$.
\end{enonce*}

\begin{enonce*}[remark]{Exercise 9.7.5}\etiquette[9.7.5]{IV.exer9.7.5}

% END INCLUDED SGA4_English_current_buildable_source_20260906/source/04/drafts/SGA4_IV_LINES_7444_7700_English.texfrag


% BEGIN INCLUDED SGA4_English_current_buildable_source_20260906/source/04/drafts/SGA4_IV_LINES_7701_7780_English.texfrag
Let $E$ be a topos. For every subtopos $F$ of $E$, let $P(F)$ be the
subset of $\Point(E)$ formed by the isomorphism classes of points of
$E$ which factor through \sisi{$E$}{$F$}. Thus $P(F)$ is homeomorphic
to $\Point(F)$ (\hyperref[IV.exer9.1.8]{9.1.8} c)). Show that if $F$
is an open (\resp closed, \resp locally closed
(\hyperref[IV.exer9.4.9]{9.4.9})) subtopos of $E$, then $P(F)$ is an
open\pageoriginale (\resp closed, \resp locally closed) subset of
\sisi{$E$}{$\Point(E)$}. Show that when $E$ has enough points, the
same is true of every locally closed subtopos of $E$, and the map
$F\mapsto P(F)$ induces a bijection between the set of open (\resp
closed, \resp locally closed) subtopoi of $E$ and the set of open
(\resp closed, \resp locally closed) subsets of $\Point(E)$. Extend
the preceding considerations to the case where $\Point(E)$ is
replaced by any subspace $X$ of $\Point(E)$ corresponding to a
conservative family of points of $E$.
\end{enonce*}

\subsection{Complements concerning certain topoi associated with topological spaces}\etiquette[9.8]{IV.subsec9.8}

In this section we give a few complements in the form of exercises.
The reader is advised to look through them in order to become
accustomed to the viewpoint of topoi in various classical
situations.

\begin{enonce*}[remark]{Exercise 9.8.1}\etiquette[9.8.1]{IV.exer9.8.1}
(\emph{Descent of sheaves on a topological space.})

Let $f:X\rightarrow Y$ be a continuous map of topological spaces in
$𝒰$, giving a functor
$$
f^{\ast}:\Top(Y)\longrightarrow\Top(X){\quoi}.
$$
\begin{enumerate}
\renewcommand{\theenumi}{\alph{enumi}}
\renewcommand{\labelenumi}{\theenumi)}
\item The following conditions are equivalent:
\begin{enumerate}
\renewcommand{\theenumii}{\roman{enumii}}
\renewcommand{\labelenumii}{(\theenumii)}
\item $f^{\ast}$ is faithful.

\item $f^{\ast}$ is conservative.

\item $f(X)$ is a very dense subset (EGA IV 10.1.3) of
$Y$.
\end{enumerate}
It suffices for $f$ or $f_{\sob}$ to be surjective. Give an example
where $f^{\ast}$ is faithful and $f=f_{\sob}$ is not surjective
(take $X$ discrete, $f$ injective, $f(X)$ the set of closed points
of $Y$, and $Y$ a nondiscrete Noetherian topological space). Give an
example where $f_{\sob}$ is surjective but $f$ is not, and another
where $f$ is surjective but $f_{\sob}$ is not.

\item Suppose that $f(X)$ is very dense in $Y$ and that its
topology is the quotient topology induced from $X$. Prove that the
arrow $f$ of $(\Esp)$ is a descent morphism for the fibered category
of sheaves of sets on varying spaces, \ie that the functor
$f^{\ast}$ induces a \emph{fully faithful}\pageoriginale functor
from $\Top(Y)$ to the category of objects of $\Top(X)$ equipped with
descent data relative to $f:X\rightarrow Y$. Reduce to the case
where $f$ is surjective and follow the argument of VIII \hyperref[VIII.prop9.1]{9.1} given in
the setting of étale topologies of schemes. Give a converse.

\item Suppose that $f(X)$ is very dense in $Y$, that its
topology is the quotient topology induced from $X$, and that the
fibers of $f$ are connected. Prove that the functor $f^{\ast}$ is
fully faithful. (Follow the argument of \SGA 1 IX 3.4.) Give a
converse, at least if the points of $Y$ are closed.

\item Suppose that $Y$ is the union of open subsets $Y_i$ over which
$X$ admits sections. Prove that $f$ is then an \emph{effective
descent morphism} for the fibered category of sheaves of sets on
varying topological spaces, \ie that the functor considered in b) is
in fact an equivalence of categories.
\end{enumerate}
\end{enonce*}

% END INCLUDED SGA4_English_current_buildable_source_20260906/source/04/drafts/SGA4_IV_LINES_7701_7780_English.texfrag


% BEGIN INCLUDED SGA4_English_current_buildable_source_20260906/source/04/drafts/SGA4_IV_LINES_7781_8025_English.texfrag
\begin{enonce*}[remark]{Exercise 9.8.2}\etiquette[9.8.2]{IV.exer9.8.2}
(\emph{Quotient topoi of topological spaces. Topological
\'etendues.})

\begin{enumerate}
\renewcommand{\theenumi}{\alph{enumi}}
\renewcommand{\labelenumi}{\theenumi)}
\item Let
\hypertarget{sga4.exp.IV.equation.9.8.2.1}{}\[
\xymatrix@C=1.5cm@R=1.5cm@M=.2cm{
X_2 \ar@<6pt>[r] \ar@<-1.5pt>[r] \ar@<-8pt>[r]& X_1 \ar@<2pt>[r]^{p_1} \ar@<-4pt>[r]_{p_2}& X_0\\
}\leqno{(9.8.2.1)}
\]\etiquette[9.8.2.1]{IV.eq9.8.2.1}
be a semi-simplicial object truncated in degree $2$ in the category
$(\Esp)$ of topological spaces $\in 𝒰$; the inverse-image functors
then give a diagram of categories of sheaves
\[
\xymatrix@M=.2cm@C=1.25cm{
  \Top(X_0) \ar@<2pt>[r]^{p_1^{\ast}} \ar@<-2pt>[r]_{p_2^{\ast}}&
  \Top(X_1) \ar@<-5pt>[r] \ar[r] \ar@<5pt>[r]& \Top(X_2).
}
\]
(More precisely, one has a category cofibered over the category of
standard simplices $\Delta_n$, $0\leq n\leq 2$.) Show that the
category $\varprojlim$ of this diagram (\ie the category of sheaves
of sets on $X_0$ \emph{equipped} with a descent datum relative to
the diagram \sisi{(9.8.6.1)}{\hyperref[IV.eq9.8.2.1]{9.8.2.1}}, \ie with an
isomorphism $p^{\ast}_1(F)\simeq p^{\ast}_2(F)$ satisfying\ldots), is
a $𝒰$-topos $T$. (For a more general statement on the stability of
topoi under $\varprojlim$ operations, \cf \sisi{VI)}{\hyperrefext[VI][VI.sec8]{8}}.) Define a morphism
of topoi $q:\Top(X_0)\rightarrow T$, and prove that for every
$𝒰$-topos\pageoriginale $E$, $\Hhomtop(T,E)$ is equivalent, via
$q^{\ast}$, to the category $\varprojlim$ of the following diagram
of categories, deduced from \hyperref[IV.eq9.8.2.1]{(9.8.2.1)}:
\[
\xymatrix@M=.2cm{ \Hhomtop(\Top(X_0),E)\ar@<2pt>[r] \ar@<-4pt>[r]&
\Hhomtop(\Top(X_1),E)\ar[r]\ar@<6pt>[r]\ar@<-6pt>[r]&
\Hhomtop(\Top(X_2),E). }
\]
\item In particular, let $R$ be an equivalence relation in an object
$X$ of $(\Esp)$, \ie (\sisi{\hyperrefext[I][I.sec10]{10}}{\hyperrefext[I][I.sec10.9]{10.9}}) a subobject $R$ of
$X\times X$ (N.B. the inclusion $R\hookrightarrow X\times X$ is
continuous, but the topology of $R$ is not necessarily induced by
that of $X\times X$), such that, for every object $Y$ of $(\Esp)$,
$\Hom(Y,R)$ is the graph of an equivalence relation in $\Hom(Y,X)$.
This gives a diagram of the form \hyperref[IV.eq9.8.2.1]{(9.8.2.1)}, with $X_0=X$, $X_1=R$,
$X_2=\fprod{(R,p_2)}{R,p_1}{X}$, where $p_1$, $p_2$ are the two
projections from $R$ to $X$, and hence a topos, which we denote by
$\Top(X)/R$, or even $\Top(X/R)$, or indeed $X/R$, by abuse of
notation, and which plays the role of a quotient of $X$ by $R$ in
the sense made precise in a). Suppose that the topology of $R$ is
induced by that of $X\times X$ and that, denoting by $X/R$ the
ordinary quotient topological space, $X$ locally admits sections
over $X/R$. Prove that $T$ is then equivalent to $\Top(X/R)$.
(Use \hyperref[IV.exer9.8.1]{9.8.1} d).)

\item Let $H\rightarrow X$ be an injective morphism of topological
groups, \ie a monomorphism of group objects in $(\Esp)$, and let $R$
be the equivalence relation that it defines in $X$, \ie
$R=H\times X$, with $p_1=\pr_2$ and $p_2$ defined by the action of
$H$ on $X$ by left translations. Show that the topology of $R$ is
induced by that of $X\times X$ if and only if the topology of $H$ is
induced by that of $X$. Give examples in which this condition is not
satisfied and in which the ordinary quotient topology on $X/H$ is
the indiscrete topology, with a) $H=\mathbb{Z}\times\mathbb{Z}$,
$X=\mathbb{R}$, and b) $H=\mathbb{R}$,
$X=\mathbb{T}\times\mathbb{T}$, where
$\mathbb{T}=\mathbb{R}/\mathbb{Z}$ (``geodesics on the torus'').
Prove that the two topoi so obtained are equivalent and are not
equivalent to $\Top(X/H)$ (which is a final topos
(\hyperref[IV.subsec2.2]{2.2})), nor to any topos of the form $\Top(Y)$.

\item Let $G$ be a topological group acting on a topological space
$X$, whence, in the familiar way, a semi-simplicial object
\[
\xymatrix@M=.2cm{ ... G\times G \times X
\ar@<6pt>[r]\ar@<-1.5pt>[r]\ar@<-8pt>[r] &
  G\times X \ar@<3pt>[r] \ar@<-4pt>[r] & X.
}
\]
Show\pageoriginale that the topos $T$ deduced from it by a) identifies
with the category of spaces $X'$ over $X$ having $G$ as a group of
operators, such that $X'\rightarrow X$ is an \'{e}tale map
(compatible with the action of $G$). In particular, when $G$ is
discrete, one recovers the topos $\Top(X,G)$ of \hyperref[IV.subsec2.3]{2.3}.

The preceding considerations justify the notation $\Top(X)/G$, or
indeed $\Top(X/G)$ or even simply $X/G$, for the preceding topos $T$.
One should bear in mind, however, that when $G$ is a discrete group
(to fix ideas) acting \emph{properly} on $X$, so that the ordinary
quotient space $X/G$ has reasonably good properties, the natural
morphism of topoi $T\rightarrow\Top(X/G)$ deduced from the universal
characterization a) of $T$ is an equivalence of topoi only if $G$
acts freely, \ie without fixed points, \ie when
$G\times X\rightarrow X\times X$ is a monomorphism. Thus, in the
case of a ``pre-equivalence relation'' (or ``groupoid'' in the sense
of (\SGA 3 V 1)) which is not an equivalence relation, passage to
the quotient in the sense of topoi (or ``passage to the fine
quotient'') does not in general correspond (via the correspondence
$Y\mapsto\Top(Y)$) to the usual topological quotient.

\item Let $T$ be a $𝒰$-topos. Show that the following conditions are
equivalent:
\begin{itemize}
\item[(i)] There exists a family $(S_i)_{i\in I}$ of objects of $T$
covering the final object of $T$, such that for every $i\in I$ the
induced topos $T_{/S_i}$ is equivalent to a topos of the form
$\Top(X)$, with $X\in(\Esp)$.

\item[${\rm (i')}$] There exists an $S\in\Ob T$ covering the final
object such that the induced topos $T_{/S}$ is equivalent to a topos
of the form $\Top(X)$.

\item[(ii)] There exist a topological space $X$ and an
\emph{\'{e}tale} pre-equivalence relation $R$ in $X$ (in the sense of
the category $(\Esp)$) (\ie one for which $p_1:R\rightarrow X$ is an
\'{e}tale map), such that $T$ is equivalent to $\Top(X/R)$, with the
notation of b).
\end{itemize}

The\pageoriginale topos $T$ is then said to be \emph{locally a
topological space}, or alternatively a \emph{topological
\'{e}tendue} (or simply an \emph{\'{e}tendue}, when no confusion is
likely).

\item Show that the topos $T$ constructed in a) has enough points
and, more precisely, that it admits a conservative family of points
parametrized by the (small) set $X_0$. In particular, every
\'{e}tendue has a small conservative family of points, and hence has
enough points. When an \'{e}tendue $T$ is represented in the form
$\Top(X/R)$ as in e) ii), prove that every point of $T$ is isomorphic
to the image of a point of $\Top(X)$, and is therefore defined by an
ordinary point of $X_{\sob}$ (or of $X$ itself, when $X$ is sober).

\item Show that the topos $\Top(X/G)$ of \hyperref[IV.subsec2.3]{2.3} is an
\'{e}tendue. (Use its description in d) or \hyperref[IV.exer7.1.10]{7.1.10} c).)
Show that, for every $x\in X$, the monoid of endomorphisms of the
point of $\Top(X,G)=\Top(X)/G$ defined by $x$ is canonically
isomorphic to the stabilizer group $G_x$ of $x$. In particular, the
points of an \'{e}tendue may have nontrivial automorphism groups.
Determine $\Ppoint(T)$ for the \'{e}tendue $T$ defined by an \'{e}tale
pre-equivalence relation in a space $X$, and show that every
endomorphism of a point of $T$ is an automorphism.

\item Let $T$ be a topos. Show that the following conditions are
equivalent:
\begin{itemize}
\item[i)] For every point $p$ of $T$, every endomorphism of $p$
(resp. every automorphism of $p$) is the identity.

\item[ii)] For every topos $S$ having enough points and every
morphism of topoi $f:S\rightarrow T$, every endomorphism of $f$
(resp. every automorphism of $f$) is the identity.
\end{itemize}
Likewise, show that the following conditions are equivalent:
\begin{itemize}
\item[${\rm i')}$] For two points $p$, $q$ of $T$, there exists at
most one morphism from $p$ to $q$.

\item[${\rm ii')}$] For two morphisms $f$, $g:S\rightrightarrows T$
from a topos $S$ having enough\pageoriginale points to $T$, there
exists at most one morphism from $f$ to $g$.
\end{itemize}

When these latter conditions hold, the topos $T$ is said to be
\emph{pointwise rigid}, or simply \emph{rigid}. Show that if $X$ is
a topological space equipped with an \'{e}tale pre-equivalence
relation $R$, the quotient \'{e}tendue $T=\Top(X)/R$ is rigid if and
only if $R$ is an equivalence relation.

\item Let $(X,G)$ be a topological space with a discrete group of
operators, let $H$ be a normal subgroup of $G$ such that the induced
action of $H$ on $X$ is proper and free, and let $X'=X/H$,
$G'=G/H$, so that $G'$ acts on $X'$ in the evident way. Prove that
the morphism of topoi
$$
\Top(X,G)\longrightarrow \Top(X',G')
$$
induced by the canonical morphism $(X,G)\rightarrow(X',G')$ of
spaces with operators (\hyperref[IV.subsubsec4.1.2]{4.1.2}) is an equivalence of
topoi.

Show that the set of open subtopoi of $\Top(X/G)$ identifies with
the set of open subsets of $X$ stable under the action of $G$, or
equivalently with the set of open subsets of the quotient
topological space $X/G$. Deduce that $\Top(X/G)$ is connected if and
only if every subset of $X$ which is both open and closed and stable
under the action of $G$ is either empty or equal to $X$. When the
connected components of $X$ are open, $\Top(X/G)$ is connected if
and only if $G$ acts transitively on the set $\pi_0(X)$ of connected
components of $X$; more generally, $\pi_0(\Top(X/G))$
(\hyperref[IV.exer8.7]{8.7} g)) is an essentially constant pro-set
isomorphic to the quotient set $\pi_0(X)/G$.

Show that when $X$ is locally connected and locally simply
connected, and $T=\Top(X,G)$ is connected, \ie $G$ is transitive on
$\pi_0(X)$, then among all the ways of representing $T$ by means of
a space with operators $(X',G')$ (\cf above for some such
representations), there is one, unique up to nonunique isomorphism,
for which $X'$ is connected and simply connected. Show that choosing
such an $(X',G')$ amounts to choosing a ``universal covering'' of
the final object of $T$, and that $G'$ is isomorphic to the group
$\pi_1(T)$ relative to this universal covering (\cf
\hyperref[IV.exer2.7.5]{2.7.5}).

\noindent (Hint:\pageoriginale show that giving an equivalence of
$T$ with a $\Top(X,G)$ amounts to giving a $G_T$-torsor $P$ in $T$,
and that, if $(X,G)$ and $P$ correspond, there is a canonical
equivalence $\Top(X)\simeq T_{/P}$.)

Deduce from this a trivial proof of the last assertion of c).
Characterize the topoi $T$ equivalent to a $\Top(X,G)$, where $G$
is a discrete group acting on a locally connected and locally simply
connected topological space while transitively permuting its
connected components, as the connected, locally connected, and
locally simply connected \'{e}tendues whose universal covering $P$
of the final object $e_T$ of $T$ is such that the induced topos
$T_{/P}$ is equivalent to a topos of the form $\Top(X)$ (or, as we
shall simply say by abuse of language, $T_{/P}$ ``is'' a topological
space).

Give an example of an \'{e}tendue, locally isomorphic to $\mathbb{R}$,
which is connected and simply connected but is not a topological
space. (Take the universal covering of the line $\mathbb{R}$ with
doubled origin.)

\item A ringed topos $(T,{\Oo}_T)$ (\cf \hyperref[IV.sec11]{\S 11}) is called a
\emph{differentiable \'{e}tendue} (resp. a \emph{real-analytic
\'{e}tendue}, resp. a \emph{complex-analytic \'{e}tendue}) if there
exist objects $S_i$ of $T$ covering the final object such that the
induced ringed topos $(T_{/S_i},{\Oo}_T/S_i)$ is equivalent to the
ringed topos defined by a differentiable manifold with its sheaf of
real $C^{\infty}$ functions (resp. to the ringed topos defined by a
real-analytic space, resp. to the ringed topos defined by a
complex-analytic space). Give a constructive description of these
ringed topoi of the type in e) ii) (taking for $X$, respectively, a
differentiable manifold, a real-analytic space, or a
complex-analytic space). Give examples of such ringed topoi which are
not equivalent to ringed topoi associated with topological spaces,
by returning to the examples considered in c).
\end{enumerate}
\end{enonce*}

% END INCLUDED SGA4_English_current_buildable_source_20260906/source/04/drafts/SGA4_IV_LINES_7781_8025_English.texfrag


% BEGIN INCLUDED SGA4_English_current_buildable_source_20260906/source/04/drafts/SGA4_IV_LINES_8026_8245_English.texfrag

\begin{enonce*}[remark]{Exercise 9.8.3}\etiquette[9.8.3]{IV.exer9.8.3}
(\emph{Topoi\pageoriginale associated with local equivalence relations
and with foliations}) \footnote[1]{This exercise is given subject to
every reservation, since it was written hastily and insufficiently
checked. Mr. N. \textsc{Saavedra} checked parts a) through e). The
reader is requested to communicate any observations to us.}.

Let $X$ be a topological space.
\begin{enumerate}
\renewcommand{\theenumi}{\alph{enumi}}
\renewcommand{\labelenumi}{\theenumi)}
\item For every open subset $U$ of $X$, let $\Quot(U)$ be the set of
equivalence relations on $U$, and, for an open subset $V\subset U$,
consider the natural map $\Quot(U)\rightarrow \Quot(V)$. This gives a
presheaf $\Quot_X$ on $X$. Let $Q_X$, or simply $Q$, be the associated
sheaf, and let $r$ be a section of $Q$ (a ``\emph{local equivalence
relation on $X$}''). From the order relation on the set $\Quot(U)$ of
equivalence relations on an open subset $U$ of $X$, deduce on the
presheaf $\Quot$, and on the associated sheaf $Q$, an order structure,
\ie a structure of a presheaf \resp a sheaf with values in the category
of ordered sets.

\item For every equivalence relation $R$ on $X$, consider the section
$\loc(R)$ of $Q$ that it defines. Let $E(r)$ be the set of equivalence
relations on $X$ such that $\loc(R)$ is less fine than $r$, and let
$\glob(r)$ be the least upper bound of the equivalence relations in
$E(r)$ (its graph is the intersection of the graphs of the $R\in
E(r)$). Let $(U_x)_{x\in X}$ be a family of open neighborhoods of the
points $x\in X$, and, for every $x\in X$, let $R_x$ be an equivalence
relation on $U_x$ whose germ at $x$ is $r_x$. For every family $𝒱$
of open subsets $V_x$ $(x\in X, x\in V_x\subset U_x)$, let $R_{𝒱}$
be the equivalence relation on $X$ generated by the family of
relations $R_x| V_x$. Show that if $𝒱'\leq 𝒱$ (in the obvious
sense), then $R_{𝒱'}\geq R_{𝒱}$, and that $\glob(r)$ is the least
upper bound of the increasing filtered family of equivalence relations
$R_{𝒱}$. Give an example in which $\glob(r)$ does not belong to
$E(r)$, \ie in which there is an $a\in X$ such that, for\pageoriginale
every open neighborhood $W\subset U_a$ of $a$, there exist a
$𝒱=(V_x)$ and two points $y,z$ of $W$ which are equivalent under
$R_a$ but not under $R_{𝒱}$. Show that, for every $R\in\Quot(X)$,
one has $\glob\loc(R)\geq R$.

\item The local equivalence relation $r$ is said to be
\emph{precoherent (\resp coherent)} if $\glob(r)\in E(r)$, \ie if
$\loc(\glob(r))\leq r$ (\resp if, for every open subset $U$ of $X$,
the restriction of $r$ to $U$ is precoherent). It is said to be
globally coherent if it is coherent and if, in addition,
$r=\loc\glob(r)$. An equivalence relation $R$ on $X$ is said to be
\emph{locally coherent} if $r=\loc(R)$ is coherent, and
\emph{coherent} if, in addition, $R=\glob(r)$, \ie if
$R=\glob(\loc(R))$. Show that globally coherent local equivalence
relations $r$ on $X$ correspond bijectively to coherent equivalence
relations $R$ on $X$, by $r\mapsto\glob(r)$ and
$R\mapsto\loc(R)$.

\item Show that, in order for $r$ to be coherent, it is enough that,
for every $a\in X$ and every open neighborhood $W'\subset U_a$ of
$a$, there exist an open neighborhood $W\subset W'$ of $a$ such that
every equivalence class of $R_a|W$ is contained in a connected
component of an equivalence class of $R_a|W'$ (a fortiori, it is
enough that there exist a fundamental system of open neighborhoods
$W_a\subset U_a$ of $a$ such that the fibers of the induced
equivalence relation $R_a|W_a$ are connected). (Hint: reduce to
proving precoherence in the case where $r$ is defined by an
$R\in\Quot(X)$, and observe in that case that the fibers of the
equivalence relations $R_{𝒱}$ are relatively open---and hence also
relatively closed---subsets of the fibers of $R$.) Show that, in order
for a global equivalence relation $R$ to be coherent, it is enough
that it be locally coherent and that its fibers be connected; prove
that this condition is necessary if the fibers are closed. Give an
example of an equivalence relation with connected and locally
connected fibers which is not locally coherent (?).

\item Let\pageoriginale $f:X'\rightarrow X$ be a continuous map.
Define a natural homomorphism of presheaves
$f_{\ast}(\Quot_{X'})\leftarrow\Quot_X$, inducing a homomorphism of
sheaves $f_{\ast}(Q_{X'})\leftarrow Q_X$, and hence a homomorphism
$f^{\ast}(Q_X)\rightarrow Q_{X'}$. If $r$ is a local equivalence
relation on $X$, then $f$ is said to be a ``\emph{fiber map}'' for
the local equivalence relation $r$ on $X$ if the inverse image of $r$
is the coarse local equivalence relation on $X'$. For every
topological space $X'$, let $\Homfib_r(X',X)$ be the set of continuous
maps from $X'$ to $X$ which are fiber maps relative to $r$. Show that,
as $X'$ varies, this defines a contravariant functor in $X'$.

\item Suppose that $r$ is coherent. Show that the preceding functor is
representable by a topological space $X^r$ over $X$. Show that the
universal fiber map $\Phi:X^r\rightarrow X$ is bijective, and that
every $x\in X^r$ has an open neighborhood $U$ such that the map
$U\rightarrow X$ induced by $\Phi$ is a homeomorphism of $U$ onto its
image. Show that the connected components of $X^r$ are open and
correspond under the bijection $\Phi$ to the fibers of the equivalence
relation $R=\glob(r)$. Give an example in which the restriction of
$\Phi$ to the connected components of $X^r$ does not induce
homeomorphisms from those spaces onto their images in $X$.

\item The local equivalence relation $r$ is said to be \emph{open} if
it is a section of the subsheaf $\Qouv_X$ of $Q$ arising from the
subpresheaf $\Quotouv_X$ of $\Quot_X$ whose value on every open subset
$U$ of $X$ is the set of \emph{open} equivalence relations on $U$.
The local equivalence relation $r$ is said to be \emph{strictly open}
if it is coherent and open or, equivalently, if it is coherent and if,
for every open subset $U$ of $X$, $\glob(r|U)$ is an open equivalence
relation on $U$. An equivalence relation $R$ on $X$ is said to be
\emph{strictly open} if it is coherent and if $\loc(R)$ is a strictly
open local equivalence relation. Prove that, for the open relation $r$
to be strictly open, it is enough that it satisfy\pageoriginale the
condition considered in d); if $r$ is locally defined by an $R$ with
closed fibers, then that condition is also necessary.

\item Let $F$ be a sheaf of sets on $X$. For every open subset $U$ of
$X$, let $\Quot(U,F)$ be the set of pairs consisting of an equivalence
relation $R_U$ on $U$ and an equivalence relation $R_{F|U}$ on $F|U$
($F$ being regarded as an étale space over $X$), such that the
structure map $p:F|U\rightarrow U$ is compatible with the equivalence
relations $R_U$ and $R_{F|U}$, and such that the corresponding diagram
of topological spaces
\[
\xymatrix@C=2cm@R=1cm{
F|U \ar[r]\ar[d] \ar[d]_{p}& (F|U)/R_{F|U} \ar[d]_{q}\\
U   \ar[r]&  U/R_{U} }
\]
is Cartesian, with $q$ an étale map (so that $F|U$ is identified
with the inverse image of the sheaf $(F|U)/R_{F|U}$ on $U/R_U$).
Show that the $\Quot(U,F)$, for varying $U$, define a presheaf on $X$,
whose associated sheaf will be denoted by $Q(F/X)$. Define a
homomorphism of sheaves $Q(F/X)\rightarrow Q_X$. For a given section
$r$ of $Q_X$, an $r$-\emph{structure} on the sheaf $F$ is a section
of $Q(F/X)$ lying above the given section $r$ of $Q_X$. Define the
category $\Top(X/r)$ of $r$-sheaves on $X$, and define a conservative
and faithful functor, ``forgetting the $r$-structure'',
$\Top(X/r)\rightarrow\Top(X)$. Prove that finite projective limits and
finite inductive limits are representable in $\Top(X/r)$, and that the
preceding functor commutes with those limits. Deduce that finite sums
in $\Top(X/r)$ are disjoint and universal and that equivalence
relations are universally effective. Prove that $\Top(X/r)$ has a
small generating family. Give an example, with $r$ coherent, in which
$\Top(X/r)$ does not have infinite coproducts (?), and hence is not a
topos.

\item Let $f:X\rightarrow Y$ be a continuous map compatible with
$R=\glob(r)$. Define a functor
$$
f^{\ast}_r:\Top(Y)\longrightarrow\Top(X/r)
$$
whose\pageoriginale composite with the inclusion functor
$\Top(X/r)\rightarrow\Top(X)$ is $f^{\ast}$. Show that, if $r$ is
globally coherent and strictly open (\ie if $R$ is strictly open and
$r=\loc(R)$), and if $f$ is the canonical map
$X\rightarrow Y=X/R$, then $f^{\ast}_r$ is an equivalence of
categories. Consequently, $\Top(X/r)$ is a $𝒰$-topos, and
$\Top(X/r)\rightarrow\Top(X)$ is the inverse-image functor of a
morphism of topoi.

\item Deduce from i) that, if $r$ is strictly open, then $\Top(X/r)$
is a topos, and that the inclusion functor
$\Top(X/r)\rightarrow\Top(X)$ is the inverse-image functor associated
with a morphism of topoi
$$
p:\Top(X)\longrightarrow\Top(X/r){\quoi}.
$$

\item Define a functorial construction of the topos $\Top(X/r)$ and
of the preceding morphism of topoi $p$ with respect to the pair
$(X,r)$ consisting of a topological space $X$ equipped with a strictly
open local equivalence relation. Show that, if $X'\rightarrow X$ is an
étale map and $r'$ is the inverse image of $r$ on $X'$, then the
induced morphism $\Top(X'/r')\rightarrow\Top(X/r)$ is equivalent to a
localization morphism $\Top(X/r)_{/E}\rightarrow\Top(X/r)$, where $E$
is an object of $\Top(X/r)$ (determined up to a unique isomorphism).
The object $E$ covers the final object of $\Top(X/r)$ if and only if
the saturation under $R=\glob(r)$ of the image of $X'$ in $X$ is all
of $X$.

\item Deduce from k) that $\Top(X/r)$ is a \emph{rigid étendue}
(\hyperref[IV.exer9.8.2]{9.8.2} h)). Show that, if $X$ is sober, the
set of isomorphism classes of points of $\Top(X/r)$, endowed with its
canonical topology (\hyperref[IV.exer7.8]{7.8}), is homeomorphic to the
quotient topological space $X/\glob(r)$, and that, for any two points
of $\Top(X/r)$ arising from points $x$ and $x'$ of $X$, there is at
most one morphism from one to the other. Such a morphism exists if and
only if $x'$ and $x$ are equivalent modulo $\glob(r)$ to points $x'_1$
and $x_1$ such that $x_1$ is a generalization of $x'_1$.

\item Study the canonical morphism of topoi
$\Top(X)\rightarrow\Top(X/r)$, observing\pageoriginale that, for every
object $U$ of $\Top(X/r)$ such that the topos induced on $U$ ``is'' an
ordinary topological space, the induced morphism
$\Top(X)_{/p^{\ast}(U)}\rightarrow\Top(X/r)_{/U}$ ``is'' a continuous
map of ordinary topological spaces $X_U\rightarrow U$. Show that the
fibers of the map $X_U\rightarrow U$ are homeomorphic to connected
components of the space $X^r$ of f).

\item Let $X$ be a differentiable manifold equipped with a foliation,
\ie with a subsheaf $F$, locally a direct factor of the tangent sheaf
of $X$, which is stable under brackets. Define on $X$ a strictly open
local equivalence relation associated with the foliation, and define
on the quotient topos $T=\Top(X/r)$ a ring which makes it a
\emph{differentiable étendue} (\hyperref[IV.exer9.8.2]{9.8.2} j)).
Define an equivalence between the category of locally free modules on
the differentiable étendue $T$ (also called differentiable vector
bundles on $T$) and the category of locally free modules on $X$
equipped with a connection relative to the given subbundle $F$ of the
tangent bundle. Give variants in the real-analytic and complex-analytic
cases.
\end{enumerate}
\end{enonce*}

% END INCLUDED SGA4_English_current_buildable_source_20260906/source/04/drafts/SGA4_IV_LINES_8026_8245_English.texfrag


% BEGIN INCLUDED SGA4_English_current_buildable_source_20260906/source/04/drafts/SGA4_IV_LINES_8246_8573_English.texfrag
\section{Sheaves of morphisms}\etiquette[10]{IV.sec10}\pageoriginale

\begin{enonce*}{Proposition 10.1}\etiquette[10.1]{IV.prop10.1}
Let $E$ be a topos and let $X$ and $Y$ be two objects of $E$; the
functor
$Z\mapsto \Hom_E(Z\times X,Y)\simeq \Hom_{E_{/Z}}(X_Z,Y_Z)$ is
representable.
\end{enonce*}

Indeed, inductive limits are universal in $E$ (II \hyperref[II.prop4.3]{4.3}). Hence the
functor $Z\mapsto Z\times X$ commutes with inductive limits.
Consequently, the functor $Z\mapsto \Hom_E(Z\times X,Y)$ carries
inductive limits in the argument $Z$ to projective limits. It is
therefore representable \sisi{(IV~\hyperref[IV.prop1.4]{1.4} and
\hyperref[IV.theo1.2]{1.2})}{(\hyperref[IV.cor1.2.1]{1.2.1})}.

\setcounter{subsection}{1}
\subsection{}\etiquette[10.2]{IV.subsec10.2}
The object representing the functor $Z\mapsto \Hom_E(Z\times X,Y)$
is denoted by $\scrHom_E(X,Y)$ (or more simply by $\scrHom(X,Y)$) and
is called the \emph{sheaf of morphisms from $X$ to $Y$}. It is a
bifunctor in $X$ and $Y$. We thus have a trifunctorial isomorphism
\begin{equation*}
\Hom_E(Z,\scrHom_E(X,Y))\simeq \Hom_E(Z\times X,
Y){\quoi}.\tag{10.2.1}
\end{equation*}\etiquette[10.2.1]{IV.eq10.2.1}

It then follows from formula \eqhyperref[IV.eq10.2.1]{10.2.1} that the
bifunctor $(X,Y)\mapsto \scrHom(X,Y)$ carries inductive limits in the
argument $X$ (\resp projective limits in the argument $Y$) to
projective limits in $E$.

\begin{enonce*}{Proposition 10.3}\etiquette[10.3]{IV.prop10.3}
Let $v:E\rightarrow E'$ be a morphism of topoi. For a variable object
$X$ of $E$ and a variable object $Y$ of $E'$, there is a bifunctorial
isomorphism
\begin{equation*}
v_{\ast}\scrHom_E(v^{\ast}Y,X)\simeq
\scrHom_{E'}(Y,v_{\ast}X){\quoi}.\tag{10.3.1}
\end{equation*}\etiquette[10.3.1]{IV.eq10.3.1}
\end{enonce*}

\noindent \emph{Proof}. For every object $Z$ of $E'$, we have a
sequence of isomorphisms, functorial in all the arguments:
\begin{alignat*}{4}
&\Hom_{E'}(Z,v_{\ast}\scrHom_E(v^{\ast}Y,X)) &&\simeq
&& \Hom_E(v^{\ast}Z,\scrHom_E(v^{\ast}Y,X)) &&\quad \text{(adjunction)}\\
& \Hom_E(v^{\ast}Z,\scrHom_E(v^{\ast}Y,X)) && \simeq
&& \Hom_E(v^{\ast}Z\times v^{\ast}Y,X) &&\quad \eqhyperref[IV.eq10.2.1]{10.2.1}\\
&\Hom_E(v^{\ast}Z\times v^{\ast}Y,X) &&\simeq &&
\Hom_E(v^{\ast}(Z\times Y),X) &&\quad \text{($v^{\ast}$ is left
exact)}\\
&\Hom_E(v^{\ast}(Z\times Y),X) &&\simeq && \Hom_{E'}(Z\times Y,
v_{\ast}X) &&\quad \text{(adjunction)}\\
& \Hom_{E'}(Z\times Y,v_{\ast}X) &&\simeq &&
\Hom_{E'}(Z,\scrHom_{E'}(Y,v_{\ast}X)) &&\quad \eqhyperref[IV.eq10.2.1]{10.2.1}
\end{alignat*}
The two sides of \eqhyperref[IV.eq10.3.1]{10.3.1} therefore represent
isomorphic functors, \cqfd

\begin{enonce*}{Corollary 10.4}\etiquette[10.4]{IV.cor10.4}
Let\pageoriginale $C$ be a $𝒰$-site, $X$ a $𝒰$-presheaf on
$C$, $Y$ a $𝒰$-sheaf on $C$, and $𝒱$ a universe containing
$𝒰$ such that $C$ is $𝒱$-small. Let
$C^{\widehat{\phantom{,}}}_{𝒱}$ be the topos of $𝒱$-presheaves on
$C$, and let $C^{\sim}_{𝒰}$ be the topos of $𝒰$-sheaves on $C$.
The presheaf
$\scrHom_{C^{\widehat{\phantom{,}}}_{𝒱}}(X,Y)$ is a $𝒰$-sheaf. Let
$X\rightarrow \underline{a}X$ be the canonical morphism from $X$ to
its associated sheaf. The corresponding morphism
$\scrHom_{C^{\widehat{\phantom{,}}}_{𝒱}}(\underline{a}X,Y)\rightarrow
\scrHom_{C^{\widehat{\phantom{,}}}_{𝒱}}(X,Y)$ is an isomorphism. There
is a canonical isomorphism
$\scrHom_{C^{\widehat{\phantom{,}}}_{𝒱}}(\underline{a}X,Y)\simeq
\scrHom_{C^{\sim}_{𝒰}}(\underline{a}X,Y)$.
\end{enonce*}

\setcounter{subsection}{4} \setcounter{subsubsection}{0}
\subsubsection{}\etiquette[10.4.1]{IV.subsubsec10.4.1}
Suppose first that $C$ is a small site and that $𝒱=𝒰$. We then have
a morphism of topoi
$C^{\sim}_{𝒰}\rightarrow
C^{\widehat{\phantom{,}}}_{𝒰}$ (\sisi{IV~}{}\hyperref[IV.subsec4.8]{4.8}),
from which \eqhyperref[IV.eq10.3.1]{10.3.1} gives the assertions in
this case. To pass from this to the general case, note first that the
``associated sheaf'' functor is independent of the universe
(\hyperrefext[II][II.prop3.6]{3.6}), and that the topos of
$𝒱$-sheaves on $C$ is equivalent to the topos of $𝒱$-sheaves on
$C^{\sim}_{𝒰}$ (\hyperrefext[III][III.sec4]{4} and
\hyperref[IV.sec1]{1}). It therefore suffices to prove the following
lemma:

\begin{enonce*}{Lemma 10.4.2}\etiquette[10.4.2]{IV.lem10.4.2}
Let $E$ be a $𝒰$-topos, let $𝒱$ be a universe containing $𝒰$, let
$E^{\sim}_{𝒱}$ be the topos of $𝒱$-sheaves on $E$, and let $X$ and
$Y$ be two objects of $E$. There exists a canonical isomorphism
$\scrHom_E(X,Y)\simeq \scrHom_{E^{\sim}_{𝒱}}(X,Y)$.
\end{enonce*}

\setcounter{subsubsection}{2}
\subsubsection{}\etiquette[10.4.3]{IV.subsubsec10.4.3}
The functor $\epsilon_E:E\rightarrow E_{\widetilde{𝒱}}$ is fully
faithful and left exact. Consequently, for every object $Z$ of $E$,
there is an isomorphism
$$
\Hom_{E_{\widetilde{𝒱}}}(\epsilon_E Z,\epsilon_E
\scrHom_E(X,Y))\simeq \Hom_{E_{\widetilde{𝒱}}}(\epsilon_E
Z\times \epsilon_E X,\epsilon_E Y).
$$
But every object of $E_{\widetilde{𝒱}}$ is an inductive limit of
objects coming from $E$; the assertion therefore follows from
\eqhyperref[IV.eq10.2.1]{10.2.1}.

\subsection{}\etiquette[10.5]{IV.subsec10.5}
Let $v:E\rightarrow E'$ be a morphism of topoi. We shall define four
bifunctorial morphisms:
\begin{alignat*}{4}
\Phi_v & : v^{\ast}\scrHom(X,Y) &&\longrightarrow
\scrHom(v^{\ast}X,v^{\ast}Y) &\quav
& X \text{ and } Y \text{ objects of } E'{\quoi},\\
\Psi_v & : v_{\ast}\scrHom(X,Y) &&\longrightarrow \scrHom(v_{\ast}X,
v_{\ast}Y) &\quav & X \text{ and } Y \text{ objects of } E{\quoi},\\
\Xi_v & : \scrHom(v_{!}X,Y) &&\longrightarrow
v_{\ast}\scrHom(X,v^{\ast}Y) &\quav & X\in \ob E, Y\in \ob
E'{\quoi},\\
\Lambda_v & : v_{!}(X\times v^{\ast}Y) &&\longrightarrow
v_{!}(X)\times Y & & X\in \ob E, Y\in \ob E'{\quoi},
\end{alignat*}
where,\pageoriginale in order to define the morphisms $\Xi_v$ and
$\Lambda_v$, we suppose that the inverse-image functor $v^{\ast}$ of
$v$ admits a left adjoint $v_{!}$.

\subsubsection{}\etiquette[10.5.1]{IV.subsubsec10.5.1}
\emph{Definition of} $\Psi_v$. Let $X$ be an object of $E$. There is
an adjunction morphism $v^{\ast}v_{\ast}X\rightarrow X$, and hence,
for every object $Z$ of $E'$, a bifunctorial morphism
$v^{\ast}(Z\times v_{\ast}X)\rightarrow (v^{\ast}Z)\times X$. Thus,
for every object $Y$ of $E$, there is a trifunctorial morphism
$$
\Hom_E((v^{\ast}Z)\times X,Y)\longrightarrow
\Hom_E(v^{\ast}(Z\times v_{\ast}X),Y){\quoi}.
$$
By adjunction this gives a trifunctorial morphism
$$
\Hom_{E'}(Z,v_{\ast}\scrHom(X,Y))\longrightarrow \Hom_{E'}(Z,
\scrHom(v_{\ast}X,v_{\ast}Y)){\quoi};
$$
and hence the morphism $\Psi_v$.

\subsubsection{}\etiquette[10.5.2]{IV.subsubsec10.5.2}
\emph{Definition of} $\Lambda_v$. For every object $X$ of $E$, there
is an adjunction morphism $X\rightarrow v^{\ast}v_{!}X$, and hence,
for every object $Y$ of $E'$, a bifunctorial morphism
$X\times v^{\ast}Y\rightarrow v^{\ast}v_{!}(X)\times v^{\ast}Y
\simeq v^{\ast}(v_{!}(X)\times Y)$; by adjunction this gives the
morphism $\Lambda_v$.

\subsubsection{}\etiquette[10.5.3]{IV.subsubsec10.5.3}
\emph{Definition of} $\Xi_v$. Let $X$ be an object of $E$, and let
$Z$ and $Y$ be two objects of $E'$. The morphism
$\Lambda_v:v_{!}(X\times v^{\ast}Z)\rightarrow v_{!}(X)\times Z$
gives a trifunctorial morphism
$\Hom_{E'}(v_{!}(X)\times Z,Y)\rightarrow
\Hom_{E'}(v_{!}(X\times v^{\ast}Z),Y)$; by adjunction this gives a
trifunctorial morphism
$$
\Hom_{E'}(Z,\scrHom(v_{!}X,Y))\longrightarrow
\Hom_{E'}(Z,v_{\ast}\scrHom(X,v^{\ast}Y)){\quoi};
$$
and hence the morphism $\Xi_v$.

\subsubsection{}\etiquette[10.5.4]{IV.subsubsec10.5.4}
\emph{Definition of} $\Phi_v$. The trifunctorial morphism of
(\hyperref[IV.subsubsec10.5.3]{10.5.3})
$$
\Hom_{E'}(v_{!}(Z)\times X,Y)\longrightarrow
\Hom_{E'}(v_{!}(Z\times v^{\ast}X),Y){\quoi},
$$
where $Z$ is an object of $E$ and $X,Y$ are objects of $E'$, gives,
by adjunction, a trifunctorial morphism
$$
\Hom_E(Z,v^{\ast}\scrHom(X,Y))\longrightarrow
\Hom_E(Z,\scrHom(v^{\ast}X,v^{\ast}Y)){\quoi};
$$
and hence a definition of the morphism $\Phi_v$. There is a second
way to define $\Phi_v$. Let $X$, $Y$, and $Z$ be three objects of
$E'$. The functor $v^{\ast}:E'\rightarrow E$ gives a trifunctorial
morphism
$$
\Hom_{E'}(Z\times X,Y)\longrightarrow \Hom_E(v^{\ast}(Z)\times
v^{\ast}(X),v^{\ast}Y){\quoi};
$$
and hence, by adjunction, a trifunctorial morphism
$$
\Hom_{E'}(Z,\scrHom(X,Y))\longrightarrow
\Hom_{E'}(Z,v_{\ast}\scrHom(v^{\ast}X,v^{\ast}Y)){\quoi}.
$$
We\pageoriginale therefore have a bifunctorial morphism
$$
\scrHom(X,Y)\longrightarrow v_{\ast}\scrHom(v^{\ast}X,
v^{\ast}Y){\quoi};
$$
and hence, by adjunction, a morphism which the reader may verify is
indeed the morphism $\Phi_v$ defined above.

\begin{enonce*}{Proposition 10.6}\etiquette[10.6]{IV.prop10.6}
Let $v:E\rightarrow E'$ be a morphism of topoi.
\begin{enumerate}
\renewcommand{\labelenumi}{\theenumi)}
\item The morphism $\Psi_v$ (\hyperref[IV.subsubsec10.5.1]{10.5.1})
is an isomorphism for every object $Y$ of $E$ if and only if the
adjunction morphism $v^{\ast}v_{\ast}X\rightarrow X$ is an
isomorphism. In particular, $\Psi_v$ is an isomorphism for every pair
of objects $(X,Y)$ if and only if the functor
$v_{\ast}:E\rightarrow E'$ is fully faithful, \ie if $v$ is an
embedding of topoi.

\item The following conditions are equivalent:
\begin{enumerate}
\renewcommand{\theenumii}{\roman{enumii}}
\renewcommand{\labelenumii}{(\theenumii)}
\item For every pair $(X,Y)$ of objects of $E'$, the morphism
$\Phi_v$ (\hyperref[IV.subsubsec10.5.4]{10.5.4}) is an isomorphism.

\item The functor $v^{\ast}$ admits a left adjoint $v_{!}$, and for
every object $X$ of $E$ and every object $Y$ of $E'$, the morphism
$\Xi_v$ (\hyperref[IV.subsubsec10.5.3]{10.5.3}) is an isomorphism.

\item The functor $v^{\ast}$ admits a left adjoint $v_{!}$, and for
every object $X$ of $E$ and every object $Y$ of $E'$, the morphism
$\Lambda_v$ (\hyperref[IV.subsubsec10.5.2]{10.5.2}) is an isomorphism.
\end{enumerate}
\end{enumerate}
\end{enonce*}

\setcounter{subsection}{6} \setcounter{subsubsection}{0}
\subsubsection{}\etiquette[10.6.1]{IV.subsubsec10.6.1}
The first assertion follows immediately from
(\hyperref[IV.subsubsec10.5.1]{10.5.1}). It also follows immediately
from (\hyperref[IV.subsubsec10.5.2]{10.5.2}),
(\hyperref[IV.subsubsec10.5.3]{10.5.3}), and
(\hyperref[IV.subsubsec10.5.4]{10.5.4}) that
(iii) $\Leftrightarrow$ (ii) $\Leftrightarrow$ ((i) and the functor
$v^{\ast}$ admits a left adjoint $v_{!}$). It remains to show that
(i) implies that $v^{\ast}$ admits a left adjoint, \ie
(\hyperref[IV.rem1.8]{1.8}) that (i) implies that $v^{\ast}$ commutes
with small products. Let $e'$ be the final object of $E'$ and let $I$
be a small set. Since $v^{\ast}$ is left exact, $v^{\ast}(e')$ is a
final object of $E$, and since $v^{\ast}$ commutes with inductive
limits, $v^{\ast}(\amalg_I e')\simeq \amalg_I v^{\ast}(e')$. For
every object $Y$ of $E'$, we have
$\scrHom(\amalg_I e',Y)\simeq \Pi_I\scrHom(e',Y)\simeq \Pi_IY$ and
$\scrHom(\amalg_I v^{\ast}(e'),v^{\ast}Y)\simeq
\Pi_I\scrHom(v^{\ast}(e'),v^{\ast}Y)\simeq
\Pi_Iv^{\ast}Y$.\pageoriginale The functorial morphism $\Phi_v$
therefore induces an isomorphism
$v^{\ast}\Pi_IY\simeq \Pi_Iv^{\ast}Y$, which one checks is the
canonical isomorphism. \cqfd

\begin{enonce*}{Corollary 10.7}\etiquette[10.7]{IV.cor10.7}
Let $E$ be a topos, let $Z$ be an object of $E$, let
$j_Z:E_{/Z}\rightarrow E$ be the localization morphism
(\hyperref[IV.subsec5.2]{5.2}), and let $X,Y$ be two objects of $E$.
\begin{enumerate}
\renewcommand{\labelenumi}{\theenumi)}
\item There is a bifunctorial isomorphism
$$
j^{\ast}_Z\scrHom(X,Y)\simeq \scrHom(j^{\ast}_ZX,
j^{\ast}_ZY){\quoi}.
$$

\item Let $X'$ be an object of $E_{/Z}$. There is a bifunctorial
isomorphism
$$
\scrHom_E(j_{Z!}X',Y)\simeq
j_{Z\ast}\scrHom_{E/Z}(X',j^{\ast}_ZY){\quoi}.
$$

\item Let $Y'$ be an object of $E/Z$. When $Z$ is open in $E$, there
is a bifunctorial isomorphism
$$
j_{Z\ast}\scrHom(X',Y')\simeq \scrHom(j_{Z\ast}X',
j_{Z\ast}Y'){\quoi}.
$$
\end{enumerate}
\end{enonce*}

Assertions 1) and 2) are proved by noting that the morphism
$\Lambda_{j_Z}$ is an isomorphism. For assertion 3), it suffices to
note that $j_{Z\ast}$ is fully faithful, since $j_{Z!}$ is fully
faithful (\hyperrefext[I][I.rem5.7]{5.7}. a)).

\begin{enonce*}{Corollary 10.8}\etiquette[10.8]{IV.cor10.8}
Let $E$ be a topos, let $X$, $Y$, and $Z$ be three objects of $E$,
and let $j_Z:E_{/Z}\rightarrow E$ be the localization morphism.
\begin{enumerate}
\renewcommand{\labelenumi}{\theenumi)}
\item There is a canonical isomorphism
$$
j_{Z\ast}j^{\ast}_ZY\simeq \scrHom(Z,Y){\quoi}.
$$

\item There is a canonical isomorphism
$$
\Hom_E(Z,\scrHom(X,Y))\simeq
\Hom_{E_{/Z}}(j^{\ast}_ZX,j^{\ast}_ZY){\quoi}.
$$
\end{enumerate}
\end{enonce*}

Let $e_Z$ be the final object of $E_{/Z}$ (\ie the object
$\id_Z:Z\rightarrow Z$). It follows from
\eqhyperref[IV.eq10.2.1]{10.2.1} that there is a canonical
isomorphism
$$
\scrHom_{E/Z}(e_Z,j^{\ast}_ZY)\simeq j^{\ast}_ZY{\quoi};
$$
hence, by (\hyperref[IV.cor10.7]{10.7} 2), an isomorphism
$$
j_{Z\ast}j^{\ast}_ZY\simeq \scrHom(j_{Z!}e_Z,Y){\quoi.}
$$
Since $j_{Z!}e_Z=Z$, this proves 1). Let us prove 2). From
(\hyperref[IV.cor10.7]{10.7}) we obtain a canonical\pageoriginale
isomorphism
$$
\Hom_{E_{/Z}}(e_Z,j^{\ast}_Z\scrHom(X,Y))\simeq
\Hom_{E_{/Z}}(j^{\ast}_ZX,j^{\ast}_ZY){\quoi};
$$
hence, by adjunction on the first member,
$$
\Hom_E(Z,\scrHom(X,Y))\simeq
\Hom_{E_{/Z}}(j^{\ast}_ZX,j^{\ast}_ZY){\quoi}.
$$


% END INCLUDED SGA4_English_current_buildable_source_20260906/source/04/drafts/SGA4_IV_LINES_8246_8573_English.texfrag


% BEGIN INCLUDED SGA4_English_current_buildable_source_20260906/source/04/drafts/SGA4_IV_LINES_8574_8831_English.texfrag
\section{Ringed topoi, localization in ringed topoi}\etiquette[11]{IV.sec11}

\setcounter{subsection}{1} \setcounter{subsubsection}{0}
\subsubsection{}\etiquette[11.1.1]{IV.subsubsec11.1.1}
Let $𝒰$ be a universe. A $𝒰$-\emph{ringed topos} is a pair
$(E,A)$ where $E$ is a $𝒰$-topos and $A$ is an object endowed with
a ring structure. A $𝒰$-ringed site is a pair $(C,A)$ where $C$ is
a $𝒰$-site and $A$ is a $𝒰$-sheaf of rings on $C$. We omit mention
of the universe when the context leaves no room for confusion. To a
ringed site $(C,A)$ is associated the ringed topos $(C^{\sim},A)$.
To a ringed topos $(E,A)$ is associated the ringed site consisting
of the site $E$ and the sheaf of rings represented by $A$.

\subsubsection{}\etiquette[11.1.2]{IV.subsubsec11.1.2}
Let $(E,A)$ be a ringed topos. We denote by ${}_AE$ (\resp $E_A$)
the category of sheaves of unital left (\resp right) $A$-modules
(we will also write $A$-Modules). The category ${}_AE$ (\resp
$E_A$) is an abelian category
(\hyperrefext[II][II.prop6.7]{6.7}). Let $M$ and $N$ be two sheaves
of left (\resp right) $A$-modules. The commutative group of
$A$-Module morphisms from $M$ to $N$ is denoted by $\Hom_A(M,N)$.

\subsubsection{}\etiquette[11.1.3]{IV.subsubsec11.1.3}
Let $A$, $B$, and $C$ be three rings in a topos $E$, let $M$ be a
sheaf of $A-B$-bimodules, and let $N$ be a sheaf of
$A-C$-bimodules. Let $e$ be a final object of $E$, and put
$\Hom(e,B)=\Gamma(B)$ and $\Hom(e,C)=\Gamma(C)$. The
$A-B$-bimodule structure on $M$ gives a ring homomorphism from
$\Gamma(B)^{\circ}$ to the ring of endomorphisms of the $A$-module
$M$; similarly, the $A-C$-bimodule structure on $N$ gives a ring
homomorphism from $\Gamma(C)^{\circ}$ to the ring of endomorphisms
of the $A$-module $N$. By functoriality, we deduce a
$\Gamma(B)-\Gamma(C)$-bimodule structure on the commutative group
$\Hom_A(M,N)$. In particular, when $A$ is a sheaf of commutative
rings, the group $\Hom_A(M,N)$ is canonically endowed with a
$\Gamma(A)$-module structure.

\subsubsection{}\etiquette[11.1.4]{IV.subsubsec11.1.4}
Let\pageoriginale $E$ be a topos, let $A$ and $B$ be two rings in
$E$, let $u:A\rightarrow B$ be a morphism of sheaves of rings, and
let $M$ be a $B$-module (a left module, to fix ideas). The sheaf $M$
may be regarded as a sheaf of $A$-modules by means of $u$: for every
object $X$ of $E$, $M(X)$ is endowed with the $A(X)$-module
structure obtained from its $B(X)$-module structure and the homomorphism
$u(X):A(X)\rightarrow B(X)$ by restriction of scalars. We thus
obtain a functor ``\emph{restriction of scalars along} $u$''
$$
\Res(u):{}_BE\longrightarrow {}_AE{\quoi}.
$$

\subsubsection{}\etiquette[11.1.5]{IV.subsubsec11.1.5}
In particular, when $A=\ZZ$ (the constant sheaf
$\sisi{\underline{\underline{\ZZ}}}{\underline{\ZZ}}$, that is, the
sheaf associated to the constant presheaf $\ZZ$) and when
$u:\sisi{\underline{\underline{\ZZ}}}{\underline{\ZZ}}\rightarrow B$
is the unique canonical morphism, we obtain a functor from ${}_BE$
to the category
${}_{\sisi{\underline{\underline{\ZZ}}}{\underline{\ZZ}}}E$, which
is none other than the category of abelian sheaves, denoted by
$E_{Ab}$. This functor is called the \emph{underlying abelian-sheaf
functor}.

\begin{enonce*}{Proposition 11.1.6}\etiquette[11.1.6]{IV.prop11.1.6}
The restriction-of-scalars functor along $u$ commutes with inductive
and projective limits. It is conservative.
\end{enonce*}

The proposition is true for the restriction-of-scalars functor for
ordinary modules, that is, when $E$ is the punctual topos. It is
therefore true when $E$ is the topos of presheaves on a small site.
It then follows from the determination of inductive and projective
limits by means of the associated-sheaf functor
(\hyperrefext[II][II.prop6.4]{6.4}) that the proposition is true in
general.

\setcounter{subsection}{2} \setcounter{subsubsection}{0}
\subsubsection{}\etiquette[11.2.1]{IV.subsubsec11.2.1}
Let $(E,A)$ be a ringed topos, let $X$ be an object of $E$, and let
$j_X:E_{/X}\rightarrow E$ be the localization morphism
\sisi{(\hyperref[IV.sec8]{8})}{\hyperref[IV.subsec5.2]{5.2}}.
By \hyperrefext[III][III.prop1.7]{1.7} or
\hyperref[IV.subsubsec3.1.2]{3.1.2}, the sheaf $j^{\ast}_XA$ is
canonically endowed with a ring structure. The sheaf of rings
$j^{\ast}_XA$ is most often denoted by $A|X$, or, by abuse, simply
by $A$. Unless otherwise stated, the topos $E_{/X}$ will be ringed
by $A|X$. The sheaf $j_{X\ast}j^{\ast}_XA$ is canonically endowed
with a ring structure. The adjunction morphism
$A\rightarrow j_{X\ast}j^{\ast}_XA$ is a morphism of sheaves of
rings.

\subsubsection{}\etiquette[11.2.2]{IV.subsubsec11.2.2}
Let\pageoriginale, in addition, $M$ be an $A$-module (a left module,
to fix ideas). The sheaf $j^{\ast}_XM$ is endowed with an
$A|X$-module structure (\hyperrefext[III][III.prop1.7]{1.7}; or
\hyperref[IV.subsubsec3.1.2]{3.1.2}); hence there is a functor
$$
j^{\ast}_X:{}_AE\longrightarrow {}_{A|X}E_{/X}
$$
called the restriction functor to $E_{/X}$, or also the restriction
functor to $X$. The functor $j^{\ast}_X$ commutes with inductive and
projective limits (\loccit). In particular, it is exact. Now let $N$
be an $A|X$-Module. The sheaf $j_{X\ast}N$ is a sheaf of
$j_{X\ast}A|X$-Modules (\hyperrefext[III][III.prop1.7]{1.7} or
\hyperref[IV.subsubsec3.1.2]{3.1.2}); hence, by restriction of
scalars along the adjunction morphism
$A\rightarrow j_{X\ast}(A|X)$
(\hyperref[IV.subsubsec11.1.4]{11.1.4}), it is an $A$-Module, again
denoted by $j_{X\ast}N$. We have therefore defined a functor
$$
j_{X\ast}:{}_{A|X}E_{/X}\longrightarrow {}_AE
$$
which commutes with projective limits
(\hyperref[IV.prop11.1.6]{11.1.6} and
\hyperrefext[III][III.prop1.7]{1.7}). For every $A$-Module $M$, the
adjunction morphism $M\rightarrow j_{X\ast}j^{\ast}_XM$ is a
morphism of $A$-Modules. For every $A|X$-Module $N$ and every
$A$-Module $M$, the adjunction morphism
$M\rightarrow j_{X\ast}j^{\ast}_XM$ defines a morphism,
bifunctorial in $M$ and $N$,
\begin{equation*}
\Hom_{A|X}(j^{\ast}_XM,N)\longrightarrow
\Hom_A(M,j_{X\ast}N){\quoi}.\tag{11.2.2.1}
\end{equation*}\etiquette[11.2.2.1]{IV.eq11.2.2.1}
The latter morphism is an isomorphism; that is, the functors
$j^{\ast}_X$ and $j_{X\ast}$ for $A$-Modules are adjoint
(III \hyperref[III.prop1.7]{1.7}).

\subsubsection{}\etiquette[11.2.3]{IV.subsubsec11.2.3}
The functors $j^{\ast}_X$ and $j_{X\ast}$ for Modules commute with
the ``underlying set'' functor
(\hyperrefext[III][III.prop1.7]{1.7} 4)). They therefore commute with
the restriction-of-scalars functors and, in particular, with the
underlying abelian-sheaf functor.

\setcounter{subsection}{3} \setcounter{subsubsection}{0}
\begin{enonce*}{Proposition 11.3.1}\etiquette[11.3.1]{IV.prop11.3.1}
Let $(E,A)$ be a ringed topos and let $X$ be an object of $E$. The
functor
$$
j^{\ast}_X:{}_{A}E\longrightarrow {}_{A|X}E_{/X}
$$
admits a left adjoint, denoted by
$j_{X!}:{}_{A|X}E_{/X}\rightarrow {}_{A}E$ and called extension by
zero. The extension-by-zero functor is exact and faithful, and
commutes with inductive limits. The functors $j_{X!}$ for
Modules\pageoriginale commute with restriction-of-scalars functors
and, in particular, they commute with the underlying abelian-sheaf
functor.
\end{enonce*}

The existence of the functor $j_{X!}$ follows from
\hyperrefext[III][III.prop1.7]{1.7}. Since $j_{X!}$ is a left
adjoint, it commutes with inductive limits
(\hyperrefext[I][I.sec2]{2}). To prove the other assertions, first
suppose that $E$ is the topos of presheaves of sets on a small
category $C$ containing the object $X$. We then know that $E_{/X}$
is equivalent to the category of presheaves on $C_{/X}$ and that,
under this equivalence, the functor
$j^{\ast}_X:E\rightarrow E_{/X}$ is simply composition with the
forgetful functor $C_{/X}\rightarrow C$
(\hyperrefext[I][I.prop5.11]{5.11}). It then follows immediately
from the explicit construction of $j_{X!}$
(\hyperrefext[I][I.prop5.1]{5.1}) that, for every $A|X$-Module $N$
and every object $Y$ of $C$, one has
$$
j_{X!}N(Y)=\bigoplus_{u\in \Hom_C(Y,X)}N(u){\quoi};
$$
hence $j_{X!}$ is exact, and extension by zero commutes with
restriction-of-scalars functors in this case. In the general case,
we may suppose that $E$ is the topos of sheaves on a small site $C$
and that $X$ comes from an object of $C$
(\hyperref[IV.cor1.2.1]{1.2.1}). The topos $E_{/X}$ is then
equivalent to the topos of sheaves on $C_{/X}$ (endowed with the
induced topology) (\hyperrefext[III][III.prop5.4]{5.4}), and the
topos morphism $j_X:E_{/X}\rightarrow E$ comes from the forgetful
functor $C_{/X}\rightarrow C$, which is continuous and cocontinuous
(\hyperrefext[III][III.prop5.2]{5.2}). It then follows from
\hyperrefext[III][III.prop1.7]{1.7} that extension by zero for
sheaves is obtained by composing extension by zero for presheaves
with the associated-sheaf functor; hence the assertion of exactness
and the compatibility with restrictions of scalars.

To prove faithfulness, it is equivalent to show that, for every
$A|X$-Module $N$, the adjunction morphism
$$
\ad_N:N\longrightarrow j^{\ast}_X j_{X!}N
$$
is a monomorphism. Now, denoting by $j_{\widehat{X}!}$ the
extension-by-zero functor for presheaves and by $\underline{a}$ the
``associated sheaf'' functor, we have a commutative diagram
\[
\xymatrix@C=2.75cm@R=1.5cm{ N \ar[r]^{\ad_{N}}
\ar[rd]_{\ad_{\hat{N}}} & j^{\ast}_{X} \underline{a}
j_{\widehat{X}!} N \\
                   & \underline{a} j^{\ast}_X j_{\widehat X!} N,\ar[u]
}
\]
where\pageoriginale $\ad_{\widehat{N}}$ is the adjunction morphism
for presheaves. The morphism $\ad_{\widehat{N}}$ is a monomorphism
by what precedes, and therefore
$\underline{a}(\ad_{\widehat{N}})$ is a monomorphism. Moreover,
since the forgetful functor $C_{/X}\rightarrow C$ is continuous and
cocontinuous, the canonical morphism
$\underline{a}j^{\ast}_X\rightarrow j^{\ast}_X\underline{a}$ is an
isomorphism (\hyperrefext[III][III.prop2.3]{2.3}). Consequently,
$\ad_N$ is a monomorphism.

\begin{enonce*}[remark]{Remark 11.3.2}\etiquette[11.3.2]{IV.rem11.3.2}
Let $N$ be an $A|X$-Module. The sheaf of sets underlying $j_{X!}N$
is not, in general, isomorphic to the sheaf obtained by extending
the sheaf of sets underlying $N$ by the empty set
(\hyperref[sga4.exp.III.terminology.5.3]{III 5.3} and \hyperref[IV.subsec5.2]{5.2}). One must
therefore take care not to confuse the \emph{extension-by-zero}
functor denoted by
$j_{X!}:{}_{A|X}E_{/X}\rightarrow {}_{A}E$ in
\hyperref[IV.prop11.3.1]{11.3.1} with the extension-by-the-empty-set
functor, also denoted by $j_{X!}:E_{/X}\rightarrow E$, in III \hyperref[III.subsec5.3]{5.3}
and IV \hyperref[IV.subsec5.2]{5.2}. In most cases encountered in practice, the abuse of
notation just noted causes no confusion. When confusion is
nevertheless possible, we will use the notations $j^{\ab}_{X!}$ and
$j^{\ens}_{X!}$ for the extension-by-zero and
extension-by-the-empty-set functors, respectively.
\end{enonce*}

\begin{enonce*}{Proposition 11.3.3}\etiquette[11.3.3]{IV.prop11.3.3}
Let $(E,A)$ be a ringed topos and let $X$ be an object of $E$. The
$A$-Module $j_{X!}(A|X)$, most often denoted by $A_X$ or $A_{X,E}$,
is the free $A$-Module generated by $X$
{\rm (\hyperrefext[II][II.prop6.5]{6.5})}; that is, for every
$A$-Module $M$, there is a canonical isomorphism, functorial in $M$:
$$
\Hom_E(X,M)\xleftarrow{\sim} \Hom_A (A_X, M){\quoi}.
$$
Let $(X_i)_{i\in I}$ be a topologically generating family of $E$
{\rm (\hyperrefext[II][II.def3.0.1]{3.0.1}).} The family
$(A_{X_i})_{i\in I}$ is a generating family of the category of
$A$-Modules.
\end{enonce*}

Let $e_X$ be the final object of the topos $E_{/X}$
$(e_X=X\xrightarrow{\id}X)$. There is an isomorphism
$$
\Hom_{A|X}(A|X,j^{\ast}_XM)\xrightarrow{\sim}
\Hom_{E_{/X}}(e_X,j^{\ast}_XM){\quoi};
$$
deduced\pageoriginale from the unit section $e_X\rightarrow A|X$;
hence, by adjunction, there is an isomorphism
$$
\Hom_A(j^{\ab}_{X!}(A|X),
M)\xrightarrow{\sim}\Hom_E(j^{\ens}_{X!}(e_X),M){\quoi}.
$$
Since $j^{\ens}_{X!}(e_X)=X$, we obtain the asserted isomorphism.
The last assertion follows from \hyperrefext[II][II.cor6.6]{6.6}.

\begin{enonce*}[remark]{Remark 11.3.4}\etiquette[11.3.4]{IV.rem11.3.4}
Let $A$ and $B$ be two sheaves of rings on a topos $E$, let $X$ be
an object of $E$, and let $N$ be an $A|X$-$B| X$-bimodule. The
abelian sheaf obtained by extending $N$ by zero is canonically
endowed with an $A$-$B$-bimodule structure, as follows immediately
from its explicit description (\hyperref[IV.prop11.3.1]{11.3.1}).
In particular, the free $A$-module so obtained, $A_X$, is an
$A$-bimodule.
\end{enonce*}

\begin{enonce*}[remark]{Exercise 11.3.5}\etiquette[11.3.5]{IV.exer11.3.5}
Let $E$ be a topos, let $X$ be an object of $E$, let
$x:P\rightarrow E$ be a point of $E$ (\hyperref[IV.def6.1]{6.1}),
and let $(x_i)_{i\in I}$ be the family of points of $E_{/X}$ above
$x$ (in bijective correspondence with the fiber $X_x$
(\hyperref[IV.eq6.7.2]{6.7.2})). Show that, for every abelian sheaf
$M$ on $E_{/X}$, the fiber $(j_{X!}M)_x$ is canonically isomorphic
to $\oplus_i M_{x_i}$.
\end{enonce*}

% END INCLUDED SGA4_English_current_buildable_source_20260906/source/04/drafts/SGA4_IV_LINES_8574_8831_English.texfrag


% BEGIN INCLUDED SGA4_English_current_buildable_source_20260906/source/04/drafts/SGA4_IV_LINES_8832_9247_English.texfrag
\section{Operations on modules}\etiquette[12]{IV.sec12}

\begin{enonce*}{Proposition 12.1}\etiquette[12.1]{IV.prop12.1}
Let $(E,A)$ be a ringed topos, and let $M$ and $N$ be two left
$A$-modules (\resp right $A$-modules). The functor on $E$ which
assigns to every object $X$ of $E$ the commutative group
$\Hom_A(M,\scrHom_E(X,N))$ is representable by an abelian sheaf
denoted by $\scrHom_A(M,N)$ (or sometimes $\scrHom(M,N)$ when no
confusion can result). For every object $X$ of $E$, there is a
canonical isomorphism
\begin{equation*}
\scrHom_A(M,N)(X)\simeq \Hom_{A| X}(j^{\ast}_XM,
j^{\ast}_XN){\quoi}.\tag{12.1.1}
\end{equation*}\etiquette[12.1.1]{IV.eq12.1.1}
\end{enonce*}

There is a canonical isomorphism
$\scrHom_E(X,N)=j_{X\ast}j^{\ast}_XN$
(\hyperref[IV.cor10.8]{10.8}); consequently, $\scrHom_E(X,N)$ is
endowed, functorially in $X$, with a structure of left $A$-module
(\resp right $A$-module). The functor
$X\longmapsto\scrHom_E(X,N)$ transforms inductive limits in the
argument $X$ into projective limits of $A$-modules
(\hyperref[IV.subsec10.2]{10.2} and
\hyperref[IV.subsubsec11.2.3]{11.2.3}); consequently, the functor
$$
X\longmapsto \Hom_A(M,\scrHom_E(X,N))
$$
transforms inductive limits\pageoriginale in the argument $X$ into
projective limits of commutative groups, and hence into projective
limits of the underlying sets. It is therefore representable
(\hyperref[IV.prop1.4]{1.4} and \hyperref[IV.theo1.2]{1.2}), and the
object representing it is endowed with the structure of an abelian
sheaf. By definition, for every object $X$ of $E$, there is a
canonical isomorphism
\begin{equation*}
\scrHom_A(M,N)(X)=\Hom(X,\scrHom_A(M,N))\simeq
\Hom_A(M,\scrHom(X,N))\tag{12.1.2}
\end{equation*}\etiquette[12.1.2]{IV.eq12.1.2}
and the isomorphism (\hyperref[IV.eq12.1.1]{12.1.1}) follows from
\eqhyperref[IV.eq12.1.2]{12.1.2}, from the isomorphism
$\scrHom(X,N)\approx j_{X\ast}j^{\ast}_XN$
(\hyperref[IV.cor10.8]{10.8}), and from the adjunction formulas of 10.

\setcounter{subsection}{1}
\subsection{}\etiquette[12.2]{IV.subsec12.2}
The abelian sheaf $\scrHom_A(M,N)$ is called the \emph{sheaf of
$A$-module homomorphisms from $M$ to $N$}. It is a bifunctor in $M$
and $N$. It follows from its definition and from
(\hyperref[IV.subsec10.2]{10.2}) that it transforms projective limits
in the argument $N$ (\resp inductive limits in $M$) into projective
limits of abelian sheaves. In particular, it is left exact in both
arguments.

\begin{enonce*}{Proposition 12.3}\etiquette[12.3]{IV.prop12.3}
Let $(E,A)$ be a ringed topos, let $M$ and $N$ be two right
$A$-modules (\resp left $A$-modules), and let $X$ be an object of $E$:
\begin{enumerate}
\renewcommand{\theenumi}{\alph{enumi}}
\renewcommand{\labelenumi}{\theenumi)}
\item There are canonical isomorphisms:
$$
\scrHom_A(M,N)(X)\approx \Hom_A(M, j_{X\ast}j^{\ast}_XN)\simeq
\Hom_{A| X}(j^{\ast}_XM, j^{\ast}_XN)\simeq
\Hom_A(j_{X!}j^{\ast}_XM, N).
$$

\item There is a canonical isomorphism:
$$
\Phi_X:j^{\ast}_X\scrHom_A(M,N)\simeq
\scrHom_{A| X}(j^{\ast}_XM,j^{\ast}_XN).
$$
Let, in addition, $P$ be a right $A|X$-module (\resp a left
$A|X$-module):

\item There is a canonical isomorphism:
$$
j_{X\ast}\scrHom_{A| X}(j^{\ast}_XM,P)\simeq
\scrHom_A(M,j_{X\ast}P).
$$

\item There is a canonical isomorphism:
$$
\Xi_X:\scrHom_A(j_{X!}P,N)\simeq j_{X\ast}\scrHom_{A|X}
(P,j^{\ast}_XN).
$$
\end{enumerate}
\end{enonce*}

\noindent The isomorphisms in a) follow from
\eqhyperref[IV.eq12.1.2]{12.1.2}, from the isomorphism
$\scrHom_E(X,N)\simeq j_{X\ast}j^{\ast}_XN$
(\hyperref[IV.cor10.8]{10.8}), and from the adjunction formulas of 10.

\begin{itemize}
\item[Proof of b).] For every object $Y$ of $E_{/X}$, there is the
following sequence of isomorphisms:

$\Hom(Y,j^{\ast}_X\scrHom_A(M,N))\simeq
\Hom(j_{X!}Y,\scrHom_A(M,N))$ (adjunction for sheaves of sets)
\begin{alignat*}{2}
& \Hom(j_{X!}Y,\scrHom_A(M,N))\simeq
\Hom_A(M,\scrHom_E(j_{X!}Y,N))
&\quad& \hyperref[IV.eq12.1.2]{(12.1.2)}\\
& \Hom_A(M,\scrHom_E(j_{X!}Y,N))\simeq
\Hom_A(M,j_{X\ast}\scrHom_{E_{/X}}(Y,j^{\ast}_XN))
&& (\hyperref[IV.cor10.7]{10.7})\\
& \Hom_A(M,j_{X\ast}\scrHom_{E_{/X}}(Y,j^{\ast}_XN))\simeq
\Hom_{A|X}(j^{\ast}_XM,\scrHom_E(Y,j^{\ast}_XN))
&& \eqhyperref[IV.eq11.2.2.1]{11.2.2.1}\\
& \Hom_{A|X}(j^{\ast}_XM,\scrHom_E(Y,j^{\ast}_XN))\simeq
\Hom(Y,\scrHom_{A|X}(j^{\ast}_XM,j^{\ast}_XN))
&& \hyperref[IV.eq12.1.2]{(12.1.2)}
\end{alignat*}\pageoriginale

\item[Proof of c).] For every object $Y$ of $E$, there is the
following sequence of isomorphisms:
\begin{alignat*}{2}
& \Hom(Y,j_{X\ast}\scrHom_{A|X}(j^{\ast}_XM,P))\simeq
\Hom(j^{\ast}_XY,\scrHom_{A|X}(j^{\ast}_XM,P))
&\quad& \text{(adjunction)}\\
& \Hom(j^{\ast}_XY,\scrHom_{A|X}(j^{\ast}_XM,P))\simeq
\Hom_{A|X}(j^{\ast}_XM,\scrHom_{E_{/X}}(j^{\ast}_XY,P))
&& \hyperref[IV.eq12.1.2]{(12.1.2)}\\
& \Hom_{A|X}(j^{\ast}_XM,\scrHom_{E_{/X}}(j^{\ast}_XY,P))\simeq
\Hom_A(M,j_{X\ast}\scrHom_{E_{/X}}(j^{\ast}_XY,P))
&& \eqhyperref[IV.eq11.2.2.1]{11.2.2.1}\\
& \Hom_A(M,j_{X\ast}\scrHom_{E_{/X}}(j^{\ast}_XY,P))\simeq
\Hom_A(M,\scrHom_E(Y,j_{X\ast}P))
&& (\hyperref[IV.prop10.3]{10.3})\\
& \Hom_A(M,\scrHom_E(Y,j_{X\ast}P))\simeq
\Hom(Y,\scrHom_A(M,j_{X\ast}P))
&& \hyperref[IV.eq12.1.2]{(12.1.2)}
\end{alignat*}

\item[Proof of d).] For every object $Y$ of $E$, there is the
following sequence of isomorphisms:
\begin{alignat*}{2}
& \Hom(Y,\scrHom_A(j_{X!}P,N))\simeq
\Hom_A(j_{X!}P,\scrHom_E(Y,N))
&\quad& \hyperref[IV.eq12.1.2]{(12.1.2)}\\
& \Hom_A(j_{X!}P,\scrHom_E(Y,N))\simeq
\Hom_{A|X}(P,j^{\ast}_X\scrHom_E(Y,N))
&& (\hyperref[IV.prop11.3.1]{11.3.1})\\
& \Hom_{A|X}(P,j^{\ast}_X\scrHom_E(Y,N))\simeq
\Hom_{A|X}(P,\scrHom_{E_{/X}}(j^{\ast}_XY,j^{\ast}_XN))
&& (\hyperref[IV.cor10.7]{10.7})\\
& \Hom_{A|X}(P,\scrHom_{E_{/X}}(j^{\ast}_XY,j^{\ast}_XN))\simeq
\Hom(j^{\ast}_XY,\scrHom_{A|X}(P,j^{\ast}_XN))
&& \hyperref[IV.eq12.1.2]{(12.1.2)}\\
& \Hom(j^{\ast}_XY,\scrHom_{A|X}(P,j^{\ast}_XN))\simeq
\Hom(Y,j_{X\ast}\scrHom_{A|X}(P,j^{\ast}_XN))
&& \text{(adjunction)}
\end{alignat*}
\end{itemize}

\begin{enonce*}{Corollary 12.4}\etiquette[12.4]{IV.cor12.4}
Let $C$ be a $𝒰$-site, let $A'$ be a $𝒰$-presheaf of rings on
$C$, let $M$ be a $𝒰$-presheaf of $A'$-modules (left modules, to fix
ideas), and let $N$ be a $𝒰$-sheaf of $A'$-modules (left modules).
Denote by $A$ the sheaf associated with $A'$, so that the sheaves $N$
and $\underline{a}M$ (the sheaf associated with $M$) are $A$-modules.
The presheaf
$X\longmapsto\Hom_{A'|X}(j^{\ast}_XM,j^{\ast}_XN)$ is a
$𝒰$-abelian sheaf canonically isomorphic to
$\scrHom_A(\underline{a}M,N)$.
\end{enonce*}

Indeed, it follows from \hyperrefext[III][III.prop5.5]{5.5} and
\hyperrefext[III][III.prop2.3]{2.3} that the localization functor to
$C_{/X}$ commutes with the associated-sheaf functors (for $C$ and
$C_{/X}$). Consequently, there are isomorphisms, functorial in the
variable object $X$ of $C$:
\begin{align*}
\Hom_{A'|X}(j^{\ast}_XM,j^{\ast}_XN)
&\simeq \Hom_{A|X}(\underline{a}j^{\ast}_XM,j^{\ast}_XN)\\
&\simeq \Hom_{A|X}(j^{\ast}_X\underline{a}M,j^{\ast}_XN)
\simeq \scrHom_A(\underline{a}M,N)(X){\quoi}.
\end{align*}

\begin{enonce*}{Corollary 12.5}\etiquette[12.5]{IV.cor12.5}
Let\pageoriginale $E$ be a topos, let $A$, $B$, and $C$ be three
sheaves of rings on $E$, let $M$ be an $A$-$B$-bimodule, and let $N$
be an $A$-$C$-bimodule. The abelian sheaf $\scrHom_A(M,N)$ is
canonically endowed with a $B$-$C$-bimodule structure. In particular,
when $A$ is a sheaf of commutative rings and $M$ and $N$ are
$A$-modules, $\scrHom_A(M,N)$ is canonically endowed with an
$A$-module structure. The canonical isomorphisms b), c), and d) of
$\hyperref[IV.prop12.3]{12.3}$ are isomorphisms of bimodules.
\end{enonce*}

We shall confine ourselves to giving indications. For every object
$X$ of $E$, one has
$\scrHom_A(M,N)(X)\simeq\Hom_{A|X}(j^{\ast}_XM,j^{\ast}_XN)$
(\hyperref[IV.prop12.3]{12.3}). Consequently, the commutative group\\
$\scrHom_A(M,N)(X)$ is canonically endowed with a
$B(X)$-$C(X)$-bimodule structure
(\hyperref[IV.subsubsec11.1.3]{11.1.3}), which is readily checked to
be functorial in $X$. The other assertions are left to the reader.

\begin{enonce*}{Corollary 12.6}\etiquette[12.6]{IV.cor12.6}
Let $(E,A)$ be a ringed topos, let $X$ be an object of $E$, and let
$N$ be an $A$-module (a left module, to fix ideas). There are canonical
isomorphisms of $A$-modules:
$$
\scrHom_E(X,N)\simeq j_{X\ast}j^{\ast}_XN\simeq
\scrHom_A(A_X,N){\quoi};
$$
the $A$-module structure on $\scrHom_A(A_X,N)$ comes from the
bimodule structure on $A_X$ $(\hyperref[IV.rem11.3.4]{11.3.4})$.
\end{enonce*}

The first isomorphism follows from \hyperref[IV.cor10.7]{10.7}, and
the second from the isomorphism of $A$-modules
$N\simeq\scrHom_A(A,N)$ and from \hyperref[IV.prop12.3]{12.3}.

\begin{enonce*}{Proposition 12.7}\etiquette[12.7]{IV.prop12.7}
Let $(E,A)$ be a ringed topos, let $M$ be a right $A$-module, and let
$N$ be a left $A$-module. The functor which assigns to every abelian
sheaf $P$ the commutative group
$\Hom_A(M,\scrHom_{\sisi{\underline{\underline{\ZZ}}}{\underline{\ZZ}}}(N,P))$
is representable by an abelian sheaf denoted by $M\otimes_A N$ and
called the tensor product over $A$ of $M$ and $N$.
\end{enonce*}

There is a canonical injection of
$\Hom_A(M,\scrHom_{\sisi{\underline{\underline{\ZZ}}}{\underline{\ZZ}}}(N,P))$
into the set
$\Hom(M,\scrHom(N,P))\simeq\Hom(M\times N,P)$; returning to the
definitions, one sees immediately that a morphism $f$ of sheaves of
sets from $M\times N$ to $P$ comes from an element of
$\Hom_A(M,\scrHom_Z(N,P))$ if and only if, for every object $X$ of
$E$, the map $f(X):M(X)\times N(X)\rightarrow P(X)$ is an
$A(X)$-bilinear map\pageoriginale from $M(X)\times N(X)$ to $P(X)$.
We call morphisms of sheaves of sets from $M\times N$ to $P$ having
this property $A$-bilinear morphisms. Thus, the problem is to represent
the functor of $A$-bilinear morphisms from $M\times N$ to $P$. To do
this, one proceeds as in the ordinary case, \ie the case where $E$ is
the punctual topos. Consider the eight morphisms:
\begin{itemize}
\item[(i):] $M\times M\times N\longrightarrow M\times N$\qquad
$1\leq i\leq 3$

\item[(i):] $M\times N\times N\longrightarrow M\times N$\qquad
$4\leq i\leq 6$

\item[(i):] $M\times A\times N\longrightarrow M\times N$\qquad
$7\leq i\leq 8$
\end{itemize}
defined by the formulas:
\begin{tabbing}
\quad\= (1) \= $(m_1, m_2, n)$ \= $\longmapsto$ \= $(m_1+m_2, n)$\\
\> (2) \> $(m_1, m_2, n)$ \> $\longmapsto$ \> $(m_1, n)$\\
\> (3) \> $(m_1, m_2, n)$ \> $\longmapsto$ \> $(m_2, n)$\\
\> (4) \> $(m, n_1, n_2)$ \> $\longmapsto$ \> $(m, n_1+n_2)$\\
\> (5) \> $(m, n_1, n_2)$ \> $\longmapsto$ \> $(m, n_1)$\\
\> (6) \> $(m, n_1, n_2)$ \> $\longmapsto$ \> $(m, n_2)$\\
\> (7) \> $(m, a, n)$ \> $\longmapsto$ \> $(ma, n)$\\
\> (8) \> $(m, a, n)$ \> $\longmapsto$ \> $(m, an)$;
\end{tabbing}
here formula (i) describes the map (i)$(X)$ for the variable objects
$X$ of $E$. Let $L$ denote the ``free generated abelian sheaf''
functor. It is clear that the largest quotient (in the sense of
abelian sheaves) of $L(M\times N)$ which equalizes $L(1)$ with
$L(2)+L(3)$, $L(4)$ with $L(5)+L(6)$, and $L(7)$ with $L(8)$
represents the functor of $A$-bilinear morphisms from $M\times N$ to
$P$.

\setcounter{subsection}{7}





\subsection{}\etiquette[12.8]{IV.subsec12.8}
One observes that the functor
$P\longmapsto
\Hom_A(N,\scrHom_{\sisi{\underline{\underline{\ZZ}}}{\underline{\ZZ}}}(M,P))$
is also canonically isomorphic to the functor of $A$-bilinear maps
from $M\times N$ to $P$. Consequently, the tensor product
$M\otimes_A N$ also represents the functor
$\Hom_A(N,\scrHom_{\sisi{\underline{\underline{\ZZ}}}{\underline{\ZZ}}}(M,P))$.
Thus there are isomorphisms, functorial in all the arguments:
\hypertarget{sga4.exp.IV.equation.12.8.1}{}\begin{align*}%\etiquette[12.8.1]{IV.eq12.8.1}
\Hom_Z(M\otimes_A N,P)
&\simeq \Hom_A(M,\scrHom_{\sisi{\underline{\underline{\ZZ}}}{\underline{\ZZ}}}(N,P)){\quoi},\tag{12.8.1}\\
\Hom_Z(M\otimes_A N,P)
&\simeq \Hom_A(N,\scrHom_{\sisi{\underline{\underline{\ZZ}}}{\underline{\ZZ}}}(M,P)){\quoi}.\tag{12.8.2}
\end{align*}\hypertarget{sga4.exp.IV.equation.12.8.2}{}

\subsection{}\etiquette[12.9]{IV.subsec12.9}
It\pageoriginale follows from \hyperref[IV.subsec12.8]{(12.8)}, or alternatively from \hyperlink{sga4.exp.IV.equation.12.8.1}{(12.8.1)}
and (\hyperref[IV.subsec12.2]{12.2}), that the functor $\otimes_A$
commutes with inductive limits in both arguments.

\begin{enonce*}{Proposition 12.10}\etiquette[12.10]{IV.prop12.10}
Let $C$ be a $𝒰$-site, let $A'$ be a presheaf of rings on $C$, let
$M'$ (\resp $N'$) be a right (\resp left) $A'$-module, and let $A$,
$M$, and $N$ be the sheaves associated with $A'$, $M'$, and $N'$
respectively. The sheaf associated with the presheaf
$X\longmapsto M'(X)\otimes_{A'(X)}N'(X)$ $(X\in\ob C)$ is
canonically isomorphic to the sheaf $M\otimes_A N$.
\end{enonce*}

The presheaf $X\longmapsto M'(X)\otimes_{A'(X)}N'(X)$ can be
constructed from the presheaves $M'$, $N'$, and $A'$ by the operations
specified in the proof of \hyperref[IV.prop12.7]{12.7}. Since the
``associated sheaf'' functor commutes with finite projective and
inductive limits and with the ``free generated abelian object''
functors, formation of the tensor product commutes with the
``associated sheaf'' functor.

\begin{enonce*}{Proposition 12.11}\etiquette[12.11]{IV.prop12.11}
Let $(E,A)$ be a ringed topos, let $M$ be a right $A$-module, let $N$
be a left $A$-module, and let $X$ be an object of $E$.
\begin{enumerate}
\renewcommand{\theenumi}{\alph{enumi}}
\renewcommand{\labelenumi}{\theenumi)}
\item There is a canonical isomorphism
$$
j^{\ast}_X(M\otimes_A N)\simeq
(j^{\ast}_XM)\otimes_{A|X}(j^{\ast}_XN){\quoi}.
$$
Let, in addition, $P$ be a right $A|X$-module and $Q$ a left
$A|X$-module.

\item There are canonical isomorphisms (projection formulas):
\begin{align*}
j_{X!}(P\otimes_{A|X}j^{\ast}_XN)&\simeq(j_{X!}P)\otimes_A N\\
j_{X!}(j^{\ast}_XM\otimes_{A|X}Q)&\simeq
M\otimes_A(j_{X!}Q){\quoi}.
\end{align*}
\end{enumerate}
\end{enonce*}

\begin{itemize}
\item[Proof of a).] For every abelian sheaf $R$ on $E_{/X}$, there is
the following sequence of isomorphisms, functorial in $R$:
\begin{alignat*}{2}
&\Hom_{\sisi{\underline{\underline{\ZZ}}}{\underline{\ZZ}}}
(j^{\ast}_X(M\otimes_A N),R)\simeq
\Hom_{\sisi{\underline{\underline{\ZZ}}}{\underline{\ZZ}}}
(M\otimes_A N,j_{X\ast}R)
&\quad& \eqhyperref[IV.eq11.2.2.1]{11.2.2.1}\\
&\Hom_{\sisi{\underline{\underline{\ZZ}}}{\underline{\ZZ}}}
(M\otimes_A N,j_{X\ast}R)\simeq
\Hom_A(M,\scrHom_Z(N,j_{X\ast}R))
&& \hyperlink{sga4.exp.IV.equation.12.8.1}{(12.8.1)}\\
&\Hom_A(M,\scrHom_Z(N,j_{X\ast}R))\simeq
\Hom_A(M,j_{X\ast}\scrHom_{\sisi{\underline{\underline{\ZZ}}}{\underline{\ZZ}}}
(j^{\ast}_XN,R))
&& (\hyperref[IV.prop12.3]{12.3})\\
&\Hom_A(M,j_{X\ast}\scrHom_{\sisi{\underline{\underline{\ZZ}}}{\underline{\ZZ}}}
(j^{\ast}_XN,R))\simeq
\Hom_{A|X}(j^{\ast}_XM,
\scrHom_{\sisi{\underline{\underline{\ZZ}}}{\underline{\ZZ}}}
(j^{\ast}_XN,R))
&& \eqhyperref[IV.eq11.2.2.1]{11.2.2.1}\\
&\Hom_{A|X}(j^{\ast}_XM,
\scrHom_{\sisi{\underline{\underline{\ZZ}}}{\underline{\ZZ}}}
(j^{\ast}_XN,R))\simeq
\Hom_{\sisi{\underline{\underline{\ZZ}}}{\underline{\ZZ}}}
(j^{\ast}_XM\otimes_{A|X}j^{\ast}_XN,R)
&& \hyperlink{sga4.exp.IV.equation.12.8.1}{(12.8.1)}
\end{alignat*}

Let us exhibit\pageoriginale the first isomorphism in b). For every
abelian sheaf $R$, there is a sequence of isomorphisms, functorial in
$R$:
\begin{alignat*}{2}
&\Hom_Z(j_{X!}(P\otimes_{A|X}j^{\ast}_XN),R)\simeq
\Hom_Z(P\otimes_{A|X}j^{\ast}_XN,j^{\ast}_XR)
&\quad& (\hyperref[IV.prop11.3.1]{11.3.1})\\
&\Hom_Z(P\otimes_{A|X}j^{\ast}_XN,j^{\ast}_XR)\simeq
\Hom_{\sisi{A/X}{A|X}}(P,\scrHom_Z(j^{\ast}_XN,j^{\ast}_XR))
&& \hyperlink{sga4.exp.IV.equation.12.8.1}{(12.8.1)}\\
&\Hom_{A|X}(P,\scrHom_Z(j^{\ast}_XN,j^{\ast}_XR))\simeq
\Hom_{A|X}(P,j^{\ast}_X\scrHom_Z(N,R))
&& (\hyperref[IV.prop12.3]{12.3})\\
&\Hom_{A|X}(P,j^{\ast}_X\scrHom_Z(N,R))\simeq
\Hom_A(j_{X!}P,\scrHom_Z(N,R))
&& (\hyperref[IV.prop11.3.1]{11.3.1})\\
&\Hom_A(j_{X!}P,\scrHom_Z(N,R))\simeq
\Hom_Z(j_{X!}P\otimes_A N,R)
&& \hyperlink{sga4.exp.IV.equation.12.8.1}{(12.8.1)}
\end{alignat*}

The second isomorphism is obtained similarly using \hyperlink{sga4.exp.IV.equation.12.8.2}{(12.8.2)}.
\end{itemize}

\begin{enonce*}{Corollary 12.12}\etiquette[12.12]{IV.cor12.12}
Let $E$ be a topos, let $A$, $B$, and $C$ be three sheaves of rings
on $E$, let $M$ be a $B$-$A$-bimodule, and let $N$ be an
$A$-$C$-bimodule. The abelian sheaf $M\otimes_A N$ is canonically
endowed with a $B$-$C$-bimodule structure. In particular, when $A$ is
a sheaf of commutative rings and $M$ and $N$ are $A$-modules,
$M\otimes_AN$ is canonically endowed with an $A$-module structure.
The isomorphisms of \hyperref[IV.prop12.11]{12.11} are isomorphisms of
bimodules.
\end{enonce*}

For every object $X$ of $E$, $j^{\ast}_XM$ is a right $A|X$-module
on which the ring $B(X)$ acts on the left, and likewise $j^{\ast}_XN$
is a left $\sisi{A/X}{A|X}$-module on which $C(X)$ acts on the right. By
functoriality, the abelian sheaf
$j^{\ast}_XM\otimes_{\sisi{A/X}{A|X}}j^{\ast}_XN\simeq
j^{\ast}_X(M\otimes_A N)$ is endowed with the structure of a
$B(X)$-$C(X)$-bimodule, and this structure varies functorially in
$X$. Consequently, $M\otimes_A N$ is endowed with a $B$-$C$-bimodule
structure. This $B(X)$-$C(X)$-bimodule structure on
$j^{\ast}_X(M\otimes_A N)$ is reflected in the evident way on the
functor of ``$A$-bilinear morphisms''. One then observes that, when
$A$ is commutative and $M$ and $N$ are $A$-modules, the right and left
$A$-module structures so obtained are ``equal''; consequently,
$M\otimes_A N$ is an $A$-module in this case. The last assertion is
left to the reader.

\begin{enonce*}{Corollary 12.13}\etiquette[12.13]{IV.cor12.13}
Let $(E,A)$ be a ringed topos, let $M$ be a left\pageoriginale
$A$-module (\resp a right $A$-module), and let $X$ be an object of
$E$. With the notation $A_X$ of \hyperref[IV.prop11.3.3]{11.3.3},
there is a canonical isomorphism
$$
A_X\otimes_AM\simeq j_{X!}j^{\ast}_XM
\quad\text{(\resp }M\otimes_AA_X\simeq j_{X!}j^{\ast}_XM\text{)}{\quoi}.
$$
\end{enonce*}

\noindent This follows from the projection formulas
(\hyperref[IV.prop12.11]{12.11}).

\begin{enonce*}{Proposition 12.14}\etiquette[12.14]{IV.prop12.14}
Let $E$ be a topos, let $A$ and $B$ be two sheaves of rings on $E$,
let $M$ be a right $A$-module, let $N$ be an $A$-$B$-bimodule, and let
$P$ be a right $B$-module. There is a canonical isomorphism
$$
\Hom_B(M\otimes_A N,P)\simeq
\Hom_A(M,\scrHom_B(N,P)){\quoi}.
$$
\end{enonce*}

\noindent From \hyperlink{sga4.exp.IV.equation.12.8.1}{(12.8.1)} one obtains a canonical $A$-bilinear morphism
$\scrHom_Z(N,P)\times N\rightarrow P$; hence, by restricting it to
$\scrHom_B(N,P)\times N$ (which is a subsheaf), one obtains an
$A$-bilinear morphism $\scrHom_B(N,P)\times N\rightarrow P$, which is
immediately seen to be $B$-linear in the second factor. Thus there is
a canonical morphism of $B$-modules
$\scrHom_B(N,P)\otimes_A N\rightarrow P$; hence a canonical map,
functorial in $M$:
\begin{equation*}
\Hom_A(M,\scrHom_B(N,P))\longrightarrow
\Hom_B(M\otimes_A N,P){\quoi}.\tag{12.14.1}
\end{equation*}\etiquette[12.14.1]{IV.eq12.14.1}
Let us show that this morphism of functors is an isomorphism. Since
both sides transform inductive limits in $M$ into projective limits,
it is enough to show that \eqhyperref[IV.eq12.14.1]{12.14.1} is an
isomorphism when $M$ ranges over a generating family of the category
$E_A$. It is therefore enough to show that
\eqhyperref[IV.eq12.14.1]{12.14.1} is an isomorphism when $M=A_X$,
where $X$ is an object of $E$. The verification is then immediate.

% END INCLUDED SGA4_English_current_buildable_source_20260906/source/04/drafts/SGA4_IV_LINES_8832_9247_English.texfrag


% BEGIN INCLUDED SGA4_English_current_buildable_source_20260906/source/04/drafts/SGA4_IV_LINES_9248_9593_English.texfrag
\section{Morphisms of ringed topoi}\etiquette[13]{IV.sec13}

\begin{enonce*}{Definition 13.1}\etiquette[13.1]{IV.def13.1}
Let $(E,A)$ and $(E',A')$ be two ringed topoi. A morphism of ringed
topoi $u:(E,A)\rightarrow(E',A')$ is a pair $(m,\theta)$, where
$m:E\rightarrow E'$ is a morphism of topoi
\textup{(\hyperref[IV.def3.1]{3.1})} and $\theta:m^{\ast}A'\rightarrow A$ is
a morphism of Rings.
\end{enonce*}

\setcounter{subsection}{1}
\subsubsection{}\etiquette[13.1.1]{IV.subsubsec13.1.1}
Since the functor $m^{\ast}:E'\rightarrow E$ is left adjoint to the
functor $m_{\ast}:E\rightarrow E'$, giving a morphism of ringed
topoi $(m,\theta):(E,A)\rightarrow(E',A')$ amounts to giving a
morphism of topoi $m:E\rightarrow E'$ and a morphism of sheaves of
rings $\Theta:A'\rightarrow m_{\ast}A$.

\subsection{}\etiquette[13.2]{IV.subsec13.2}
To\pageoriginale a morphism of ringed topoi
$u=(m,\theta):(E,A)\rightarrow(E',A')$, we associate two notable
functors between the categories of Modules:

\subsubsection{The direct-image functor for
  Modules:}\etiquette[13.2.1]{IV.subsubsec13.2.1}
Let $M$ be a left (\resp right) $A$-Module. The object $m_{\ast}M$ is
canonically endowed with an $m_{\ast}A$-Module structure; restriction
of scalars along the canonical morphism
$\Theta:A'\rightarrow m_{\ast}A$ therefore gives an $A'$-Module,
denoted by $u_{\ast}(M)$ and called the \emph{direct image of $M$
under the morphism $u$}.

\subsubsection{The inverse-image functor for
  Modules:}\etiquette[13.2.2]{IV.subsubsec13.2.2}
Let $N$ be a left (\resp right) $A'$-Module. The object $m^{\ast}N$
of $E$ is canonically endowed with the structure of a left (\resp
right) $m^{\ast}A'$-Module. The left $A$-Module
$A\otimes_{m^{\ast}A'}m^{\ast}N$ (\resp the right $A$-Module
$m^{\ast}N\otimes_{m^{\ast}A'}A$), where $A$ is endowed with the
$m^{\ast}A'$-Module structure defined by
$\theta:m^{\ast}A'\rightarrow A$, is denoted by $u^{\ast}N$ and is
called the \emph{inverse image of the Module $N$ under the morphism
of ringed topoi} $u$.

\subsubsection{}\etiquette[13.2.3]{IV.subsubsec13.2.3}
Note that for every $A$-Module $M$, the sheaf of sets and the abelian
sheaf underlying $u_{\ast}M$ is the sheaf $m_{\ast}M$. Thus the
notation $u_{\ast}$ is most often also used for the functor
$m_{\ast}$, the direct image for sheaves of sets. On the other hand,
for an $A'$-Module $N$, the sheaf of sets or abelian sheaf underlying
$u^{\ast}N$ is in general neither equal nor isomorphic to
$m^{\ast}N$ (except when $\theta$ is an isomorphism). One must
therefore distinguish between the inverse image for Modules and the
inverse image for sheaves of sets or abelian sheaves. The notation
$u^{-1}:E'\rightarrow E$ is most often used for the functor
$m^{\ast}$, the inverse image for sheaves of sets. The functor
$u^{-1}$ is called the \emph{``set-theoretic'' inverse-image functor}
under the morphism of ringed topoi $u$, as opposed to the inverse
image ``in the sense of Modules'' $u^{\ast}$.

\begin{enonce*}{Definition 13.3}\etiquette[13.3]{IV.def13.3}
Let $(C,A)$ and $(C',A')$ be two $𝒰$-ringed sites. A morphism of
ringed sites $u:(C,A)\rightarrow(C',A')$ is a pair $(m,\theta)$,
where $m$ is a morphism from the site $C$ to the site $C'$
(\hyperref[IV.subsec4.9]{4.9}) and
$\theta:m^{\ast}A'\rightarrow A$ is\pageoriginale a morphism of
Rings.
\end{enonce*}

\setcounter{subsection}{3} \setcounter{subsubsection}{0}
\subsubsection{}\etiquette[13.3.1]{IV.subsubsec13.3.1}
A morphism of ringed topoi is a morphism of $𝒰$-ringed sites. A
morphism of ringed sites gives rise to a morphism between the
corresponding ringed topoi (\hyperref[IV.subsubsec11.1.1]{11.1.1}
and \hyperref[IV.subsubsec4.9.1]{4.9.1}).

\begin{enonce*}{Proposition 13.4}\etiquette[13.4]{IV.prop13.4}
Let $u:(E,A)\rightarrow(E',A')$ be a morphism of ringed topoi.
\begin{enumerate}
\renewcommand{\theenumi}{\alph{enumi}}
\renewcommand{\labelenumi}{\theenumi)}
\item For a variable left (\resp right) $A$-Module $M$ and a variable
left (\resp right) $A'$-Module $N$, there are canonical
bifunctorial isomorphisms (called adjunction isomorphisms):
\hypertarget{sga4.exp.IV.equation.13.4.1}{}\begin{align*}
\Hom_A(u^{\ast}N,M) &\simeq
\Hom_{A'}(N,u_{\ast}M){\quoi},\tag{13.4.1}\\
u_{\ast}\scrHom_A(u^{\ast}N,M) &\simeq
\scrHom_{A'}(N,u_{\ast}M){\quoi}.\tag{13.4.2}
\end{align*}\hypertarget{sga4.exp.IV.equation.13.4.2}{}

\item For a variable object $X$ of $E'$, there is a canonical
isomorphism, functorial in $X$:
\begin{equation*}
u^{\ast}A'_X\xrightarrow{\sim} A_{u^{-1}(X)}{\quoi},\tag{13.4.3}
\end{equation*}\etiquette[13.4.3]{IV.eq13.4.3}
where $A'_X$ and $A_{u^{-1}(X)}$ denote the generated free Modules
(\hyperref[IV.prop11.3.3]{11.3.3}).

\item When $A'$ is commutative and the canonical morphism
$u^{-1}A'\rightarrow A$ is central (\resp when the canonical
morphism $u^{-1}A'\rightarrow A$ is an isomorphism), there is, for a
variable right $A'$-Module $M$ and a variable left $A'$-Module $N$,
a canonical bifunctorial isomorphism:
\hypertarget{sga4.exp.IV.equation.13.4.4}{}\begin{align*}
u^{\ast}M\otimes_A u^{\ast}N &\simeq
u^{\ast}(M\otimes_{A'}N)\tag{13.4.4}\\
u^{\ast}M\otimes_A u^{\ast}N &\simeq
u^{-1}(M\otimes_{A'}N){\quoi}.\tag*{(\resp (13.4.5))}
\end{align*}\hypertarget{sga4.exp.IV.equation.13.4.5}{}
\end{enumerate}
\end{enonce*}

First, there is a canonical isomorphism
$\Hom_A(u^{\ast}N,M)\simeq \Hom_{u^{-1}A'}(u^{-1}N,M)$
(\sisi{\hyperref[IV.cor12.12]{12.12}}{\hyperref[IV.prop12.14]{12.14}}),
and then an isomorphism
$\Hom_{u^{-1}A'}(u^{-1}N,M)\simeq
\Hom_{A'}(N,u_{\ast}M)$
(\hyperrefext[III][III.prop1.7]{1.7}); this gives \hyperlink{sga4.exp.IV.equation.13.4.1}{(13.4.1)}. Let us
exhibit the isomorphism \hyperlink{sga4.exp.IV.equation.13.4.2}{(13.4.2)}. For every object $X$ of $E'$, there
is the following sequence of functorial isomorphisms:
\begin{alignat*}{3}
& \Hom_{E'}(X,u_{\ast}\scrHom_A(u^{\ast}N,M))\simeq && \Hom(
u^{-1}X,
 \scrHom_A(u^{\ast}N,M)) &\quad& \text{(adjunction){\quoi},}\\
&\Hom_A(u^{\ast}N,\scrHom_E(u^{-1}X,M)) \simeq &&\qquad ''\qquad ''\qquad
''\qquad '' && (\hyperref[IV.prop12.1]{12.1}){\quoi},\\
& \Hom_A(u^{\ast}N,\scrHom_E(u^{-1}X,M))\simeq &&\Hom_{A'}(N,
u_{\ast}\scrHom_E(u^{-1}X,M)) && \hyperlink{sga4.exp.IV.equation.13.4.1}{(13.4.1)}{\quoi},\\
& \Hom_{A'}(N,\scrHom_{E'}(X,u_{\ast}M))\simeq &&\qquad ''\qquad ''\qquad
 ''\qquad '' && \eqhyperref[IV.eq10.3.1]{10.3.1}{\quoi},\\
& \qquad ''\qquad ''\qquad ''\qquad '' \simeq &&
 \Hom_{E'}(X,\scrHom_{A'}(N,u_{\ast}M)) &&
 (\hyperref[IV.prop12.1]{12.1}){\quoi}.
\end{alignat*}
Formula \hyperlink{sga4.exp.IV.equation.13.4.5}{(13.4.5)}\pageoriginale then follows from formula \hyperlink{sga4.exp.IV.equation.13.4.2}{(13.4.2)}
by adjunction (the definition of the tensor product
(\hyperref[IV.prop12.7]{12.7})). Formula \hyperlink{sga4.exp.IV.equation.13.4.4}{(13.4.4)} follows from
\hyperlink{sga4.exp.IV.equation.13.4.5}{(13.4.5)}, using the commutativity of the tensor product when the
base ring is commutative (\hyperref[IV.subsec12.8]{12.8}). Finally,
to prove \eqhyperref[IV.eq13.4.3]{13.4.3}, consider the sequence of
isomorphisms
\begin{alignat*}{3}
& \Hom_A(u^{\ast}A'_X,M) && \simeq\Hom_{A'}(A'_X,u_{\ast}M)
&\quad& \hyperlink{sga4.exp.IV.equation.13.4.1}{(13.4.1)}{\quoi},\\
& \Hom_{E'}(X,u_{\ast}M) && \simeq\quad ''\qquad ''\qquad ''&&
(\hyperref[IV.prop11.3.3]{11.3.3}){\quoi},\\
&\quad ''\qquad ''\qquad ''&& \simeq\Hom_E(u^{-1}X,M) &&
\text{(adjunction)}{\quoi},\\
&\Hom_A(A_{u^{-1}X},M) &&\simeq\quad ''\qquad ''\qquad ''&&
(\hyperref[IV.prop11.3.3]{11.3.3}){\quoi}.
\end{alignat*}

\begin{enonce*}{Corollary 13.5}\etiquette[13.5]{IV.cor13.5}
Let $(E,A)$ be a ringed topos, let $x:P\rightarrow E$ be a point of
$E$, and let $F\mapsto F_x$ be the associated fiber functor
(\hyperref[IV.def6.1]{6.1}). The fiber functor at $x$ carries the
tensor product of $A$-Modules to the (ordinary) tensor product of
$A_x$-modules.
\end{enonce*}

The tensor product in $P$ is the ordinary tensor product of modules
($P$, the punctual topos, is the category of sets). The assertion
therefore follows from \hyperlink{sga4.exp.IV.equation.13.4.4}{(13.4.4)}.

\begin{enonce*}{Corollary 13.6}\etiquette[13.6]{IV.cor13.6}
Let $u:(E,A)\rightarrow(E',A')$ be a morphism of ringed topoi. The
direct-image functor $u_{\ast}$ for right or left Modules commutes
with projective limits and, in particular, is left exact. The
inverse-image functor $u^{\ast}$ for right or left Modules commutes
with inductive limits and, in particular, is right exact.
\end{enonce*}

This follows from formula \hyperlink{sga4.exp.IV.equation.13.4.1}{(13.4.1)}
(\hyperrefext[I][I.prop2.11]{2.11}).

\setcounter{subsection}{6}
\subsection{}\etiquette[13.7]{IV.subsec13.7}
Note that the inverse-image functor $u^{\ast}$ for Modules is not, in
general, exact, whereas the set-theoretic inverse-image functor
$u^{-1}$ is exact. The exactness of $u^{\ast}$ can nevertheless be
asserted when the canonical morphism $u^{-1}A'\rightarrow A$ is flat
(on the right\pageoriginale or on the left) (V \hyperref[V.def1.8]{1.8}). This is so, in
particular, when the canonical morphism $u^{-1}A'\rightarrow A$ is an
isomorphism.

\subsection{}\etiquette[13.8]{IV.subsec13.8}
Let $C$ be a $𝒰$-site and let $A'$ be a presheaf of rings on $C$;
denote its associated sheaf by $A$. The category of presheaves of
$A'$-modules which are sheaves is equivalent to the category of
$A$-Modules. The notation $\Hom_{A'}(M,N)$ is sometimes used for the
group $\Hom_A(M,N)$ (\hyperref[IV.subsubsec11.1.2]{11.1.2}).
Likewise, and by abuse, the notations $\scrHom_{A'}(M,N)$ and
$M\otimes_{A'}N$ are used for the sheaves $\scrHom_A(M,N)$ and
$M\otimes_A N$. When $A'=k_E$ is the constant presheaf defined by an
ordinary ring $k$, we also write $\scrHom_k$, $\otimes_k$ instead of
$\scrHom_{k_E}$, $\otimes_{k_E}$.

\begin{enonce*}[remark]{Exercise 13.9}\etiquette[13.9]{IV.exer13.9}
\emph{Locally ringed topoi} (\cf [\hyperref[IV.sec9]{9}] for further
information in the direction considered below).

Let $(E,{\Oo}_E)$ be a commutatively ringed topos. For $X\in \ob E$
and $f\in {\Oo}_E(X)$, let $X_f$ be the largest subobject of $X$ on
which $f$ is invertible.
\begin{enumerate}
\renewcommand{\theenumi}{\alph{enumi}}
\renewcommand{\labelenumi}{\theenumi)}
\item Show that, for $X'\in \ob E_{/X}$, $f_{X'}$ is invertible if
and only if the structural morphism $X'\rightarrow X$ factors through
$X_f$, and that, for $f,g\in {\Oo}_E(X)$, one has
$$
X_{fg}=X_f\cap X_g{\quoi}.
$$

\item Show that the following conditions (i) through (iii) are
equivalent:
\begin{enumerate}
\renewcommand{\theenumii}{\roman{enumii}}
\renewcommand{\labelenumii}{(\theenumii)}
\item For $X\in \ob E$ and $f,g\in {\Oo}_E(X)$, one has
$$
X_{f+g}\subset \Sup(X_f,X_g).
$$

\item For $X\in \ob E$ and $f,g\in {\Oo}_E(X)$ such that $f+g$ is
invertible, one has $X=\Sup(X_f,X_g)$.

\item For $X\in \ob E$ and $f\in {\Oo}_E(X)$, one has
$$
X=\Sup(X_f,X_{1-f}){\quoi}.
$$

Show that these conditions imply the following condition (iv),
and\pageoriginale are equivalent to it if $E$ has enough points:

\item For every point $p$ of $E$, the fiber ring ${\Oo}_{E,p}$ is a
\emph{local ring}.
\end{enumerate}

When the equivalent conditions (i) through (iii) hold, $(E,{\Oo}_E)$
is said to be a \emph{locally ringed topos}; likewise, a ringed site
is said to be a \emph{locally ringed site} if the corresponding
ringed topos is a locally ringed topos.

\item Let $(E,{\Oo}_E)$ and $(E',{\Oo}_{E'})$ be two locally ringed
topoi, and let $f$ be a morphism of ringed topoi from the first to the
second. Show that the following condition (i) implies condition (ii),
and is equivalent to it if $E$ has enough points:
\begin{enumerate}
\renewcommand{\theenumii}{\roman{enumii}}
\renewcommand{\labelenumii}{(\theenumii)}
\item For every $X'\in \ob E'$ and $s'\in {\Oo}_{E'}(X')$, setting
$X=f^{-1}(X')$ and $s=f^{\ast}(s')$, one has
$$
X_s=f^{-1}(X'_{s'}){\quoi}.
$$

\item For every point $p$ of $E$, setting $p'=f(p)$, the natural
homomorphism on the fibers
$$
{\Oo}_{E',p'}\longrightarrow {\Oo}_{E,p}
$$
is a \emph{local} homomorphism of local rings.
\end{enumerate}

When condition (i) holds, $f$ is said to be a \emph{morphism of
locally ringed topoi}, or also an \emph{admissible morphism of locally
ringed topoi} if there is a risk of confusion. We denote by
$\Hhomtoplocan(E,E')$ the full subcategory of the category
$\Hhomtopan(E,E')$ of all morphisms of ringed topoi from $E$ to $E'$
defined by the admissible morphisms from $E$ to $E'$.

\item Suppose that $(E,{\Oo}_E)$ is the ringed topos associated with
a ringed topological space $(X,{\Oo}_X)$
(\hyperref[IV.subsec2.1]{2.1}). Show that $(E,{\Oo}_E)$ is locally
ringed if and only if $(X,{\Oo}_X)$ is locally ringed, \ie if and
only if, for every $x\in X$, the fiber ${\Oo}_{X,x}$ is a local
ring. (Note that in criterion (iv) of b), it suffices to take $p$ in
a \emph{conservative} family of points of $E$.) Let\pageoriginale
$f:(X,{\Oo}_X)\rightarrow(X',{\Oo}_{X'})$ be a morphism of ringed
spaces, where ${\Oo}_X$ and ${\Oo}_{X'}$ are sheaves of local rings.
Show that the corresponding morphism of ringed topoi
$\Top(X,{\Oo}_X)\rightarrow \Top(X',{\Oo}_{X'})$ is admissible if and
only if the same is true of the morphism $f$, \ie if and only if,
for every $x\in X$,
$${\Oo}_{X',f(x)}\longrightarrow {\Oo}_{X,x}$$
is a \emph{local} homomorphism of local rings.

\item Let $(E,{\Oo}_E)$ and $(E',{\Oo}_{E'})$ be two locally ringed
topoi such that $E$ has enough points and $(E',{\Oo}_{E'})$ is
equivalent to the ringed topos defined by a \emph{scheme}
$(X',{\Oo}_{X'})$; prove that the category $\Hhomtoplocan(E,E')$ is
equivalent to a \emph{discrete} category, \ie that it is a rigid
groupoid (every morphism is an isomorphism, and the automorphism
groups of the objects are the trivial groups).

\noindent (Hint: reduce to the case where $(E,{\Oo}_E)$ is the
punctual topos ringed by a field.) Recall, on the other hand, that
$\Hhomtop(E,E')$ is not in general equivalent to a discrete
category, even if $E$ and $E'$ are topoi defined by schemes (and even
if $E$ is the punctual topos), \cf
\hyperref[IV.subsubsec4.2.3]{4.2.3}.

\item Let $E$ be a locally ringed topos, and let $P$ be the ringed
topos defined by the ringed space $\Spec k$, where $k$ is a field.
Show that the admissible morphisms of locally ringed topoi
$P\rightarrow E$ correspond to pairs $(p,u)$ consisting of a point
$p$ of $E$ and an injection of $k(p)$ into $k$, where $k(p)$ is the
residue field of the local ring ${\Oo}_{E,p}$. Generalize this to a
statement exhibiting the admissible morphisms from a punctual
locally ringed topos to $E$ (generalizing EGA I 2.4.4).

A \emph{geometric point of a locally ringed topos} $E$ is any
admissible morphism into $E$ from the locally ringed topos defined by
a ringed space of the form $\Spec(k)$, where $k$ is an
\emph{algebraically closed} field. Define the \emph{category of
geometric points of} $E$ (with the understanding that the
corresponding fields $k$ belong to $𝒰$), denoted by $\Pptgeom(E)$,
and a canonical functor $\Pptgeom(E)\rightarrow \Ppoint(E)$. Show
that this functor is faithful if\pageoriginale and only if $E$ has no
point, \ie $\Ppoint(E)$ is the empty category.

\item Let $E$ be a $𝒰$-topos and let $𝒱$ be a universe such that
$𝒰\in 𝒱$. Define a 2-category
$(𝒱\text{-}𝒰\text{-}\Top)_{/E}$ in terms of the object $E$ of the
2-category $(𝒱\text{-}𝒰\text{-}\Top)$
(\hyperref[IV.subsubsec3.3.1]{3.3.1}), following
\hyperref[IV.exer5.14]{5.14} a). When $E$ is ringed by a Ring $A$,
define similarly a 2-category
$(𝒱\text{-}𝒰\text{-}\text{Topan})_{/E}$, and when $E$ is locally
ringed, use the notion of admissible morphism (\cf c)) to define a
2-category
$(𝒱\text{-}𝒰\text{-}\text{Topanloc})_{/E}$ having $\Pptgeom(E)$
(\cf f)) as a full subcategory.
\end{enumerate}
\end{enonce*}

% END INCLUDED SGA4_English_current_buildable_source_20260906/source/04/drafts/SGA4_IV_LINES_9248_9593_English.texfrag


% BEGIN INCLUDED SGA4_English_current_buildable_source_20260906/source/04/drafts/SGA4_IV_LINES_9594_9852_English.texfrag
\section{Modules on a topos defined by gluing}\etiquette[14]{IV.sec14}

\subsection{}\etiquette[14.1]{IV.subsec14.1}
Let $(E,A)$ be a ringed topos, let $U$ be an open subtopos of $E$,
let $F$ be the complementary closed subtopos, and let
$j:U\rightarrow E$ and $i:F\rightarrow E$ be the canonical
morphisms (\hyperref[IV.subsec9.3]{9.3}). The topos $U$ is ringed by
$j^{\ast}A$ (\sisi{\hyperref[IV.subsubsec11.2.2]{11.2.2}}{\hyperref[IV.subsubsec11.2.1]{11.2.1}}),
and the morphism of topoi $j$, completed in the evident way, becomes
a morphism of ringed topoi. In this section, the topos $F$ will be
ringed by the sheaf $i^{\ast}A$, denoted by $A|F$, and the morphism
$i$, completed in the evident way, is then a morphism of ringed
topoi.


\subsection{}\etiquette[14.2]{IV.subsec14.2}
We therefore have two morphisms of ringed topoi:
$$
j:(U,A| U)\longrightarrow (E,A)\text{ and } i:(F,A| F)\longrightarrow
(E,A){\quoi};
$$
whence five functors between the corresponding categories of
modules (left modules, to fix ideas):
\[\vcenter{
\xymatrix@C=2cm@R=1.5cm@M=.25cm{ ({}_{A|U}U) \ar@<14pt>[r]^{j_!}
\ar@<-14pt>[r]^{j_{\ast}}&
  ({}_{A}E) \ar[l]_{j^{\ast}} \ar@<7pt>[r]^{i^{\ast}} &
  ({}_{A|F}F)\quad. \ar@<7pt>[l]_{i_{\ast}}
}}\leqno(14.2.1)
\]\etiquette[14.2.1]{IV.eq14.2.1}
Each functor in diagram \eqhyperref[IV.eq14.2.1]{14.2.1} is left
adjoint to the functor lying below it.

\subsection{}\etiquette[14.3]{IV.subsec14.3}
We\pageoriginale know that the functor $X\mapsto (j^{\ast}X,
i^{\ast}X;i^{\ast}X\rightarrow i^{\ast}j_{\ast}j^{\ast}X)$ is an
equivalence from $E$ to the topos $(U,F,i^{\ast}j_{\ast})$
(\hyperref[IV.theo9.5.4]{9.5.4}). This equivalence induces an
equivalence between the corresponding categories of modules, and
hence the functor $P\mapsto (j^{\ast}P,i^{\ast}P;
i^{\ast}P\rightarrow i^{\ast}j_{\ast}j^{\ast}P)$ is an equivalence,
denoted by $\Phi$, from the category ${}_AE$ to the category
$({}_{A| U}U,{}_{A| F}F,i^{\ast}j_{\ast})$. The functors in
\hyperref[IV.eq14.2.1]{(14.2.1)}, composed with $\Phi$ or $\Phi^{-1}$, are then the
functors:
\begin{align*}
\Phi\circ j_{\ast} &: M\mapsto (M,i^{\ast} j_{\ast}M;
i^{\ast}j_{\ast}M\xrightarrow{\id} i^{\ast}j_{\ast}M){\quoi},\\
j^{\ast}\circ \Phi^{-1} &: (M,N; N\rightarrow
i^{\ast}j_{\ast}M)\mapsto M{\quoi},\\
\Phi\circ j_{!} &: M\mapsto (M,0;0\rightarrow i^{\ast}
j_{\ast}M){\quoi},\\
\Phi\circ i_{\ast} &: N\mapsto (0,N;N\rightarrow 0){\quoi},\\
i^{\ast}\circ \Phi^{-1} &: (M,N;N\rightarrow
i^{\ast}j_{\ast}M)\mapsto N{\quoi}.
\end{align*}
The reader may, as an exercise, make the various adjunction
morphisms explicit.

\subsection{}\etiquette[14.4]{IV.subsec14.4}
Denote by
\begin{equation*}
i^{!}:{}_AE\longrightarrow {}_{A| F}F\tag{14.4.1}
\end{equation*}\etiquette[14.4.1]{IV.eq14.4.1}
the functor defined by the formula
\begin{equation*}
i^{!}\circ
\Phi^{-1}(M,N;N\xrightarrow{u}i^{\ast}j_{\ast}M)=\Ker(u){\quoi}.\tag{14.4.2}
\end{equation*}\etiquette[14.4.2]{IV.eq14.4.2}

\begin{enonce*}{Proposition 14.5}\etiquette[14.5]{IV.prop14.5}
The functor $i^{!}$ is right adjoint to the functor $i_{\ast}$. The
adjunction morphism $i_{\ast}i^{!}\rightarrow\id$ is a
monomorphism. The functors $i^{!}$ (for varying rings) commute with
restriction-of-scalars functors.
\end{enonce*}

It is clear that every morphism
$$
(0,N;0)\longrightarrow (M,N;N\xrightarrow{u}i^{\ast}j_{\ast}M)
$$
factors uniquely through $(0,\ker(u);0)$, which proves the
adjunction property. The other assertions are trivial.

\begin{enonce*}{Proposition 14.6}\etiquette[14.6]{IV.prop14.6}
For\pageoriginale every object $P$ of ${}_AE$, there are exact
sequences, functorial in $P$:
\hypertarget{sga4.exp.IV.equation.14.6.1}{}\begin{gather}
0\longrightarrow j_{!}j^{\ast}P\longrightarrow P\longrightarrow
i_{\ast}i^{\ast}P\longrightarrow 0{\quoi};\tag{14.6.1}\\
0\longrightarrow i_{\ast}i^{!}P\longrightarrow P\longrightarrow
j_{\ast}j^{\ast}P{\quoi},\tag{14.6.2}
\end{gather}\hypertarget{sga4.exp.IV.equation.14.6.2}{}
where the nontrivial arrows are the adjunction arrows.
\end{enonce*}

First note that a sequence
$$
(M',N';u')\longrightarrow(M,N;u)\longrightarrow(M'',N'';u'')
$$
in $({}_{A|U}U,{}_{A|F}F;i^{\ast}j_{\ast})$ is exact if and only
if the corresponding sequences
\begin{align*}
& M'\longrightarrow M\longrightarrow M''\\
& N'\longrightarrow N\longrightarrow N''
\end{align*}
are exact.

The sequences \hyperlink{sga4.exp.IV.equation.14.6.1}{(14.6.1)} and \hyperlink{sga4.exp.IV.equation.14.6.2}{(14.6.2)} are transformed by the
equivalence $\Phi$ into sequences of the type in
\eqhyperref[IV.subsec14.3]{14.3}:
\begin{align*}
& 0\longrightarrow (M,0;0)\longrightarrow(M,N;u)\longrightarrow
(0,N;0)\longrightarrow 0\\
& 0\longrightarrow(0,\ker(u);0)\longrightarrow
(M,N,u)\longrightarrow(M,i^{\ast}j_{\ast}M;\id){\quoi}.
\end{align*}
The exactness of these sequences is immediate.

\setcounter{subsection}{6}
\subsection{}\etiquette[14.7]{IV.subsec14.7}
The exact sequence \hyperlink{sga4.exp.IV.equation.14.6.2}{(14.6.2)} gives a new interpretation of the
functor $i_{\ast}i^{!}$. Indeed, let $X$ be an object of $E$.
Sequence \hyperlink{sga4.exp.IV.equation.14.6.2}{(14.6.2)} gives the exact sequence of abelian groups
$$
0\longrightarrow \Hom_E(X,i_{\ast}i^{!}P)\longrightarrow
\Hom_E(X,P)\longrightarrow \Hom_E(X,j_{\ast}j^{\ast}P){\quoi}.
$$
Also denote by $U$ the open subobject of $E$ which is the final
object of the open subtopos $U$. It follows from the adjunction
properties of the functors $j_{\ast}$, $j^{\ast}$, and
$j_{!}^{\ens}$ that the group $\Hom_E(X,j_{\ast}j^{\ast}P)$ is
canonically isomorphic to $\Hom_E(X\times U,P)$, and that the
morphism $\Hom_E(X,P)\rightarrow\Hom_E(X,j_{\ast}j^{\ast}P)$ is
precisely the morphism defined by the monomorphism
$X\times U\rightarrow X$ (\hyperref[IV.sec5]{5}). In other words
(\hyperref[IV.subsubsec8.5.2]{8.5.2}):

\begin{enonce*}{Proposition 14.8}\etiquette[14.8]{IV.prop14.8}
The\pageoriginale sections of $i_{\ast}i^{!}P$ on $E_{/X}$ are the
sections of $P$ on $E_{/X}$ whose support
$(\hyperref[IV.subsubsec9.3.5]{9.3.5})$ is contained in $F$
\sisi{(notation of $5.9.1$)}{}. In other words, $i_{\ast}i^{!}P$ is
the largest subsheaf of $P$ whose support is contained in $F$.
\end{enonce*}

\setcounter{subsection}{8}
\subsection{}\etiquette[14.9]{IV.subsec14.9}
By abuse of language, one sometimes says that $i_{\ast}i^{!}P$ is
the submodule of $P$ defined by the sections of $P$ supported in
$F$. It follows from \hyperref[IV.rem8.5.3]{8.5.3} that this
terminology merely extends to general topoi terminology used for
topoi of sheaves on topological spaces.

\begin{enonce*}{Proposition 14.10}\etiquette[14.10]{IV.prop14.10}
\begin{enumerate}
\renewcommand{\labelenumi}{\theenumi)}
\item The functor $P\mapsto i_{\ast}i^{\ast}P$ is isomorphic to the
functor $P\mapsto i_{\ast}i^{\ast}A\otimes_A P$.

\item The functor $P\mapsto i_{\ast}i^{!}P$ is isomorphic to the
functor $P\mapsto\scrHom_A(i_{\ast}i^{\ast}A,P)$.
\end{enumerate}
\end{enonce*}

The exact sequence \hyperlink{sga4.exp.IV.equation.14.6.1}{(14.6.1)}, in the particular case where $P$ is
the $A$-module $A$, reads
\begin{equation*}
0\longrightarrow A_U\longrightarrow A\longrightarrow
i_{\ast}i^{\ast}A\longrightarrow 0.\tag{14.10.1}
\end{equation*}\etiquette[14.10.1]{IV.eq14.10.1}
Here $A_U$ is the free $A$-module generated by the open subobject
$U$ corresponding to the open subtopos $U$ (9 and \hyperref[IV.prop11.3.3]{11.3.3}). For
every $A$-module $P$, the $A$-modules $j_{!}j^{\ast}P$ and
$j_{\ast}j^{\ast}P$ are canonically isomorphic respectively to
$A_U\otimes_A P$ and $\scrHom_A(A_U,P)$
(\hyperref[IV.prop12.3]{12.3} and \hyperref[IV.cor12.6]{12.6}).
Moreover, under these isomorphisms, the canonical morphisms
$j_{!}j^{\ast}P\rightarrow P$ and
$P\rightarrow j_{\ast}j^{\ast}P$ arise from the monomorphism
$A_U\rightarrow A$. The exact sequence
\hyperref[IV.eq14.10.1]{14.10.1} therefore gives two exact
sequences (\hyperref[IV.subsec12.2]{12.2} and
\hyperref[IV.prop12.11]{12.11}):
\begin{align*}
& j_{!}j^{\ast}P\longrightarrow P\longrightarrow
i_{\ast}i^{\ast}A\otimes_AP\longrightarrow 0{\quoi},\\
& 0\longrightarrow \scrHom_A(i_{\ast}i^{\ast}A,P)\longrightarrow
P\longrightarrow j_{\ast}j^{\ast}P{\quoi};
\end{align*}
and the stated isomorphisms follow by comparison with \hyperlink{sga4.exp.IV.equation.14.6.1}{(14.6.1)} and
\hyperlink{sga4.exp.IV.equation.14.6.2}{(14.6.2)}.

\begin{thebibliography}{99}\pageoriginale
\bibitem{sga4.exp.IV.bib.1} M. Artin and B. Mazur, Homotopy of varieties in the etale
topology, in: Proceedings of a Conference on Local Fields,
Driebergen 1966, Springer.
\bibitem{sga4.exp.IV.bib.2} P. Gabriel and G. Zisman, Calculus of fractions and
homotopy theory, Ergebnisse der Mathematik, Bd 35.

\bibitem{sga4.exp.IV.bib.3} J. Giraud, Algèbre homologique non commutative,
Grundlehren Springer 1971.

\bibitem{sga4.exp.IV.bib.4} A. Grothendieck, Sur quelques points d'Algèbre
homologique, Tohoku Math. Journal (cited as [Toh]).

\bibitem{sga4.exp.IV.bib.5} A. Grothendieck, Fondements de la Géométrie
Algébrique (collection of Bourbaki talks 1957/62), Secrétariat
mathématique, 11 rue P. Curie, Paris.

\bibitem{sga4.exp.IV.bib.6} A. Grothendieck, Crystals and the De Rham cohomology
of schemes (notes by I. Coates and O. Jussila), in: Dix exposés sur
la cohomologie des schémas, North Holland Pub. Cie., 1969.

\bibitem{sga4.exp.IV.bib.7} A. Grothendieck, Classes de Chern des
représentations linéaires des groupes discrets, in: Dix exposés sur
la cohomologie des schémas, North Holland Pub. Cie., 1969.

\bibitem{sga4.exp.IV.bib.8} R. Godement, Théorie des faisceaux,
Act. Scient. Ind. ${\rm n^{\circ}}$ 1 252 (1958), Hermann (Paris)
(cited as [TF]).

\bibitem{sga4.exp.IV.bib.9} M. Hakim, Topos annelés et schémas relatifs,
mimeographed thesis, Orsay 1967.

\bibitem{sga4.exp.IV.bib.10} D. Mumford, Picard groups of moduli problems, in
Arithmetic Algebraic Geometry, Harper's Series in Modern
Mathematics.

\bibitem{sga4.exp.IV.bib.11} Nguyen Dinh Ngoc, Thèse Sciences Mathématiques,
Paris 1963, ${\rm n^{\circ}}$ 4 995.

\bibitem{sga4.exp.IV.bib.12} J.E. Roos, Distributivité des $\varinjlim$ par
rapport aux $\varprojlim$ des topos

\begin{tabular}{@{}r@{\;}l@{\hspace{1cm}}r@{\;}l}
a) & CR t.259 p. 969-972 & b) & CR T. 259 p. 1605-1608\\
c) & CR t.259 p. 1801-1804 & & (August and September 1964).
\end{tabular}

\bibitem{sga4.exp.IV.bib.13} P. Deligne--D. Mumford, The irreducibility of
the space of curves of given genus, Pub. math. ${\rm n^{\circ}}$ 36
(1969).

\bibitem{sga4.exp.IV.bib.14} B. Mitchell, Theory of categories, Academic Press
(1965).

\bibitem{sga4.exp.IV.bib.15} J. Lurie, Higher Topos Theory, Annals of Mathematics Studies, 170. Princeton University Press,
Princeton, NJ, 2009. xviii+925 pp.

\bibitem{sga4.exp.IV.bib.16} P. Johnstone, Topos theory,
London Mathematical Society Monographs, Vol. 10.
Academic Press [Harcourt Brace Jovanovich, Publishers], London--New York, 1977. xxiii+367 pp.

\bibitem{sga4.exp.IV.bib.17} F. Borceux, G. Van den Bossche, Structure des topologies
d'un topos, Cahiers de topologie et géométrie différentielle catégoriques,
t. 25, n°1 (1984), p. 37--39.

\end{thebibliography}

\ifx\danslelivre\undefined
\end{document}
\fi

% END INCLUDED SGA4_English_current_buildable_source_20260906/source/04/drafts/SGA4_IV_LINES_9594_9852_English.texfrag


% END INCLUDED SGA4_English_current_buildable_source_20260906/source/04/04_English.tex


\makeatletter
\def\input@path{{../05/}}
\makeatother
% BEGIN INCLUDED SGA4_English_current_buildable_source_20260906/source/05/05_English.tex
\ifx\danslelivre\undefined
% Release engine: XeLaTeX (the LuaTeX compatibility layer breaks Xy-pic output).
\documentclass[10pt]{../fichiers-configuration-latex/smfart}
% BEGIN INCLUDED SGA4_English_current_buildable_source_20260906/source/fichiers-configuration-latex/paquets.tex
\usepackage{polyglossia}
%\usepackage[francais]{babel}
\usepackage{fontspec}
\usepackage[math-style=ISO,bold-style=upright,sans-style=upright]{unicode-math}
%\usepackage[mathletters]{ucs}
%\usepackage[utf8]{inputenc}
%\usepackage[T1]{fontenc}
\usepackage[unicode,hypertexnames=false,colorlinks=true,linkcolor=blue,citecolor=Sepia,urlcolor=magenta]{hyperref}
\setmainlanguage{french}
\setmainfont[Ligatures=TeX]{Linux Libertine O}
% variante possible : TeXGyrePagella ?
\setmathfont{XITS Math}

% maths
\usepackage{amsmath}
%\usepackage{amssymb}
%\usepackage{mathrsfs}
%\usepackage{stmaryrd}
%\usepackage{euscript}

% couleurs
\usepackage[usenames,dvipsnames]{xcolor}

% graphiques
\usepackage{graphicx}
\usepackage[all]{xy}
%\usepackage{tikz}
%\usetikzlibrary{matrix,arrows,calc}
%\tikzset{column sep=3em,row sep=2em,text height=1.5ex, text depth=0.25ex,>=angle 90}
%\tikzset{description/.style={fill=white,inner sep=1pt}}

% hyperliens, 'etiquettes
\usepackage{xr-hyper}
\usepackage{../fichiers-configuration-latex/smfthm}
%\ifx\cacheretiquettes\undefined
%\usepackage[inline]{showlabels} \else \fi % pour rendre visibles les 'etiquettes
%\usepackage{../fichiers-configuration-latex/makeidx} % pour faire un index
%\usepackage[pagebackref,unicode]{hyperref}  % pour voir o`u sont cit'ees les r'ef'erences
% hyperref is loaded once above with the cumulative-reader destination policy.
%\hypersetup{breaklinks=true} %permet le retour `a la ligne dans les liens trop longs. Incompatible avec pdflatex
%\renewcommand*{\backref}[1]{{\small [cit\'e page(s)~#1.]}}

\usepackage{xspace}
\usepackage{pifont} % pour expose XVII

% END INCLUDED SGA4_English_current_buildable_source_20260906/source/fichiers-configuration-latex/paquets.tex

% BEGIN INCLUDED SGA4_English_current_buildable_source_20260906/source/fichiers-configuration-latex/numerotation.tex
\renewcommand{\thefootnote}{\roman{footnote}} % pour mettre les notes en bas de page sous la forme (ii) et non (2), qui ressemble trop (?) a un exposant.

\newcounter{pageoriginal}
\newcommand{\pageoriginale}{\refstepcounter{pageoriginal}\marginpar{\footnotesize\textbf{\thepageoriginal}}\xspace}
\let\pageoriginaled\pageoriginale

\usepackage{ifthen}
\newboolean{original}
\def\sisi{\ifthenelse{\boolean{original}}}
\setboolean{original}{false}

\newcounter{cnote}
\def\nde#1{\ifthenelse{\boolean{original}}{}{\stepcounter{cnote}\def\thefootnote{\arabic{cnote}}\protect\footnote[\value{cnote}]{N.D.E.: #1}}}
\def\ndemark{\ifthenelse{\boolean{original}}{}{\stepcounter{cnote}\def\thefootnote{\arabic{cnote}}\protect\footnotemark[\value{cnote}]}}
\def\ndetext#1{\ifthenelse{\boolean{original}}{}{\def\thefootnote{\arabic{cnote}}\protect\footnotetext[\value{cnote}]{N.D.E.: #1}}}

% END INCLUDED SGA4_English_current_buildable_source_20260906/source/fichiers-configuration-latex/numerotation.tex

% BEGIN INCLUDED SGA4_English_current_buildable_source_20260906/source/fichiers-configuration-latex/macros.tex
\setmainlanguage{french}
\setotherlanguage{english}


%% Fixed, engine-independent release identifier.  The inherited dynamic Git
%% probe used Lua io.popen, exposed build-state-dependent output, and printed
%% a visible question mark outside a Git checkout.
\def\showgitstatus{English release r2, 2026-07-23}


% macros pour LaTeX
\renewcommand{\mathsf}{}

% Police différente pour U+2113 SCRIPT SMALL L
\setmathfont[range="2113]{latinmodern-math.otf}

% Give every inherited semantic label its own stable PDF destination instead
% of inheriting whichever automatic equation/page anchor happened to be live.
% The legacy label remains as a compatibility alias used by existing edges.
\makeatletter
\def\etiquette[#1]#2{%
  \Hy@raisedlink{\hyper@anchorstart{sga4.semantic.#2}\hyper@anchorend}%
  \def\@currentHref{sga4.semantic.#2}%
  \def\@currentlabel{#1}%
  \label{#2}%
}
\makeatother
\def\eqhyperref[#1]#2{(\hyperref[#1]{#2})}
\def\pagehyperref[#1]#2{(page \hyperref[#1]{#2})}
\def\hyperrefext[#1][#2]#3{\hyperref[#2]{#1 #3}}
\def\eqhyperrefext[#1][#2]#3{\hyperref[#2]{#1 #3}}     %cf. exposes 17
\def\neweqhyperrefext[#1][#2]#3{\hyperref[#2]{#1 #3}}  %cf. exposes 17

\def\shat#1{\protect\widehat{#1}}
\def\stilde#1{\protect\widetilde{#1}}
%%%%%%%%%%%%%%%%%%


\def\etcn{etc.\xspace}
\def\etco{etc...\xspace}
\def\etc{\sisi{\etco}{\etcn}}
\def\ldotsn{\sisi{...}{\etcn}}

\def\cf{{cf.}\kern.3em}
\def\Cf{{Cf.}\kern.3em}
\def\ie{{i.e.}\kern.3em}
\def\iev{{i.e.},\xspace}
\def\loccito{loc.\kern2pt cit.\xspace}
\def\loccitn{\textit{loc.\kern2pt cit.}\xspace}
\def\loccit{\sisi{\loccito}{\loccitn}}
\def\EGA{\textup{EGA}\kern.2em}
\def\SGA{\textup{SGA}\kern.3em}

\def\quoi{}

\def\quav{\quad,\quad}
\def\quaet{\quad \text{et}\quad}
\def\fibreproduct#1#2#3{#1 ×_{#3} #2}
\let\fprod\fibreproduct
\def\ccheck#1#2{\mbox{${\check{\check{\rm #1}}}^{\lower2pt\hbox{$\scriptstyle #2$}}$}}
\def\sheaf#1{\protect\mathcal{#1}}
\def\vieuxsheaf#1{{#1}} % cf. problème avec sheaf{a\Phi_i} par exemple...
\DeclareMathOperator{\Hom}{Hom}
\DeclareMathOperator{\can}{can}
\DeclareMathOperator{\Ind}{Ind}
\DeclareMathOperator{\ind}{ind}
%\DeclareMathOperator{\Hhomtop}{\mathcal{H}\!\mathit{omtop}}
\DeclareMathOperator{\Hhomtop}{\mathbf{Homtop}}
\DeclareMathOperator{\Point}{Point}
\DeclareMathOperator{\Pro}{Pro}
\DeclareMathOperator{\Top}{Top}
\DeclareMathOperator{\Mod}{Mod}
\DeclareMathOperator{\Cart}{Cart}
%\DeclareMathOperator{\Mmorsite}{\mathcal{M}\!\mathit{orsite}}
\DeclareMathOperator{\Mmorsite}{\mathbf{Morsite}}
\DeclareMathOperator{\loc}{loc}
\DeclareMathOperator{\Loc}{Loc}
\DeclareMathOperator{\Pic}{Pic}
\DeclareMathOperator{\Et}{Et}
\DeclareMathOperator{\card}{card}
\DeclareMathOperator{\Ens}{Ens}
\DeclareMathOperator{\h}{\underline{H}}
\DeclareMathOperator{\Fonct}{Fonct}
\DeclareMathOperator{\Ker}{Ker}
\DeclareMathOperator{\but}{but}
\DeclareMathOperator{\Ob}{Ob}
\DeclareMathOperator{\ob}{ob}
\DeclareMathOperator{\Lib}{Lib}
\DeclareMathOperator{\esj}{esj}
\DeclareMathOperator{\id}{id}
\DeclareMathOperator{\Id}{Id}
\DeclareMathOperator{\ad}{ad}
\DeclareMathOperator{\Iim}{Im}
\DeclareMathOperator{\pr}{pr}
\DeclareMathOperator{\Cat}{Cat}
\DeclareMathOperator{\car}{car}
\DeclareMathOperator{\SysInd}{SysInd}
\DeclareMathOperator{\Sup}{Sup}
\DeclareMathOperator{\Coker}{Coker}
\DeclareMathOperator{\Filt}{Filt}
\DeclareMathOperator{\acc}{acc}
\DeclareMathOperator{\sur}{sur}
\DeclareMathOperator{\cat}{cat}
\DeclareMathOperator{\Ord}{Ord}
\DeclareMathOperator{\Gr}{Gr}
\DeclareMathOperator{\Cov}{Cov}
\DeclareMathOperator{\Caf}{Caf}
\DeclareMathOperator{\coker}{coker}
\DeclareMathOperator{\coim}{coim}
\DeclareMathOperator{\im}{im}
\DeclareMathOperator{\Ouv}{Ouv}
\DeclareMathOperator{\diag}{diag}
\DeclareMathOperator{\Esp}{Esp}
\DeclareMathOperator{\TOP}{TOP}
\DeclareMathOperator{\Unif}{Unif}
\DeclareMathOperator{\Iind}{\mathbf{Ind}}
\DeclareMathOperator{\Hhomcart}{\mathbf{Homcart}}
\DeclareMathOperator{\Hhomtopess}{\mathbf{Homtopess}}
\DeclareMathOperator{\Ppoints}{\mathbf{Points}}
\DeclareMathOperator{\Ppoint}{\mathbf{Point}}
\DeclareMathOperator{\Ttors}{\mathbf{Tors}}
\DeclareMathOperator{\Ssex}{\mathbf{Sex}}
\DeclareMathOperator{\Ffib}{\mathbf{Fib}}
\DeclareMathOperator{\Vvois}{\mathbf{Vois}}
\DeclareMathOperator{\Ppointess}{\mathbf{Pointess}}
\DeclareMathOperator{\scrHom}{\mathbf{Hom}}
\DeclareMathOperator{\sheafPic}{\mathbf{Pic}}
\DeclareMathOperator{\sheafAut}{\mathbf{Aut}}
\DeclareMathOperator{\Hhomtoplocan}{\mathbf{Homtoplocan}}
\DeclareMathOperator{\Hhomtopan}{\mathbf{Homtopan}}
\DeclareMathOperator{\Pptgeom}{\mathbf{Ptgéom}}
\DeclareMathOperator{\Morsite}{Morsite}
\DeclareMathOperator{\Fib}{Fib}
\DeclareMathOperator{\Fais}{Fais}
\DeclareMathOperator{\kar}{kar}
\DeclareMathOperator{\Groupes}{Groupes}
\DeclareMathOperator{\Restr}{Restr}
\DeclareMathOperator{\Isom}{Isom}
\DeclareMathOperator{\Prol}{Prol}
\DeclareMathOperator{\Res}{Res}
\DeclareMathOperator{\Tors}{Tors}
\DeclareMathOperator{\esp}{esp}
\DeclareMathOperator{\es}{es}
\DeclareMathOperator{\od}{od}
\DeclareMathOperator{\gr}{gr}
\DeclareMathOperator{\sob}{sob}
\DeclareMathOperator{\dfn}{dfn}
\DeclareMathOperator{\Sch}{Sch}
\DeclareMathOperator{\ouv}{ouv}
\DeclareMathOperator{\cons}{cons}
\DeclareMathOperator{\Nerf}{Nerf}
\DeclareMathOperator{\Points}{Points}
\DeclareMathOperator{\Espsob}{Espsob}
\DeclareMathOperator{\Homtop}{Homtop}
\DeclareMathOperator{\Pt}{Pt}
\DeclareMathOperator{\Quot}{Quot}
\DeclareMathOperator{\glob}{glob}
\DeclareMathOperator{\Homfib}{Homfib}
\DeclareMathOperator{\Qouv}{Qouv}
\DeclareMathOperator{\Quotouv}{Quotouv}
\DeclareMathOperator{\ens}{ens}
\DeclareMathOperator{\ab}{ab}
\DeclareMathOperator{\an}{an}
\DeclareMathOperator{\fs}{fs}
\DeclareMathOperator{\red}{red}
\DeclareMathOperator{\Ab}{Ab}
\DeclareMathOperator{\cosq}{cosq}
\DeclareMathOperator{\Ann}{Ann}
\DeclareMathOperator{\Zar}{Zar}


%%%%%%%%%%%%
%%%%% SGA4-2
\DeclareMathOperator{\Sex}{Sex}
\DeclareMathOperator{\vois}{vois}
\DeclareMathOperator{\Vois}{Vois}
\DeclareMathOperator{\Ext}{Ext}
\DeclareMathOperator{\ext}{ext}
\DeclareMathOperator{\Fer}{Fer}
\DeclareMathOperator{\Add}{Add}
\DeclareMathOperator{\Cocart}{Cocart}
\DeclareMathOperator{\cosk}{cosk}
\DeclareMathOperator{\plat}{plat}

\DeclareMathOperator{\faisc}{faisc}
\DeclareMathOperator{\sk}{sk}
\DeclareMathOperator{\ch}{ch}

\def\HH{\mathbf{H}}
\def\qc{\mathrm{qc}}
\def\et{\mathrm{et}}
\def\coh{\mathrm{coh}}

\def\Hvr#1{%
\begin{array}[b]{l}
{\kern2pt\scriptstyle r}\\[-8pt]
{\kern1pt\scriptstyle \vee}\\[-5pt]
H^{#1}
\end{array}}

\def\Hvv#1{%
\begin{array}[b]{l}
\kern2pt\rotatebox{90}{${\scriptscriptstyle\ll}$}\\[-6pt]
H^{#1}
\end{array}}

\def\Hvi#1{%
\begin{array}[b]{l}
{\kern0.2pt\scriptscriptstyle \infty}\\[-8pt]
{\kern1pt\scriptstyle \vee}\\[-5pt]
H^{#1}
\end{array}}

\def\Hv#1{%
\begin{array}[b]{l}
{\kern1pt\scriptstyle \vee}\\[-5pt]
H^{#1}
\end{array}}

\DeclareMathOperator{\cart}{cart}
\DeclareMathOperator{\Ffaisc}{\mathbf{Faisc}}


\def\Ev#1{%
\begin{array}[b]{l}
{\kern1pt\scriptstyle \vee}\\[-5pt]
E'_{#1}
\end{array}}

\DeclareMathOperator{\Cont}{Cont}
\DeclareMathOperator{\Ccont}{\mathscr{C}\!\mathit{ont}}
\DeclareMathOperator{\Site}{Site}
\newcommand\MM{\mathbf{M}}
\def\limtop#1{{\displaystyle{\varprojlim_{#1}\!\mbox{\rm top }}}}
\def\Limtop{\mathop{\varprojlim\!\mbox{\rm top}}}
\DeclareMathOperator{\spec}{spec}
\def\ss{\scriptstyle{\mathscr{S}}}
\DeclareMathOperator{\Gal}{Gal}
\def\smallm{\scriptstyle{\mathscr{M}}}
\DeclareMathOperator{\End}{End}
\DeclareMathOperator{\ass}{ass}
\DeclareMathOperator{\zar}{zar}
\DeclareMathOperator{\Ogg}{Ogg}
\DeclareMathOperator{\res}{res}
\DeclareMathOperator{\tr}{tr}
\DeclareMathOperator{\Supp}{Supp}
\DeclareMathOperator{\codim}{codim}
\DeclareMathOperator{\rk}{rk}
\DeclareMathOperator{\cur}{cur}
\DeclareMathOperator{\cdqc}{cdqc}
\newcommand\DD{\mathbf{D}}
\DeclareMathOperator{\cl}{cl}
\def\JJ{\mathbf{J}}
\def\FF{\mathbf{F}}
\def\BB{\mathbf{B}}
\DeclareMathOperator{\Alget}{Alget}
\DeclareMathOperator{\Ass}{Ass}
\DeclareMathOperator{\supp}{supp}
\DeclareMathOperator{\rad}{rad}
\DeclareMathOperator{\prof}{prof}
\DeclareMathOperator{\cint}{int}
\DeclareMathOperator{\sic}{sic}
\DeclareMathOperator{\Tor}{Tor}
\DeclareMathOperator{\tors}{tors}
\DeclareMathOperator{\Sym}{Sym}
\DeclareMathOperator{\parf}{parf}
\DeclareMathOperator{\tor}{tor}
\DeclareMathOperator{\defn}{defn}
\DeclareMathOperator{\Div}{Div}
\DeclareMathOperator{\Pol}{Pol}
\DeclareMathOperator{\Fex}{Fex}
\DeclareMathOperator{\Lie}{Lie}
\DeclareMathOperator{\tot}{tot}
\DeclareMathOperator{\Aut}{Aut}
\DeclareMathOperator{\HOM}{HOM}
\DeclareMathOperator{\pch}{pch}
\DeclareMathOperator{\Ch}{Ch}
\DeclareMathOperator{\EXT}{EXT}
\DeclareMathOperator{\PIC}{PIC}
\DeclareMathOperator{\Rf}{Rf}
\DeclareMathOperator{\inj}{inj}
\newcommand\VV{\mathbf{V}}

\def\Xv{%
\begin{array}[b]{l}
{\kern1pt\scriptstyle \vee}\\[-5pt]
X
\end{array}}

%%%%%%%
\newcommand\TT{\mathbf{T}}
\newcommand\XX{\mathbf{X}}
\newcommand\EE{\mathbf{E}}
%%%%%%

\def\qedsymbol{\textup{C.Q.F.D.}}
\let\CQFD\qed
\let\cqfd\qed

%%%%%\bigdot
\font\bigsymb=cmsy10 at 4pt
\def\bigdot{{\kern1.2pt\raise 1.5pt\hbox{\bigsymb\char15}}}
%%%%%

\def\Lright{\mathop{\vcenter{\hbox{\hbox{$\longrightarrow$}\lower8pt \hbox{\kern-16pt\hbox{$\longrightarrow$}}}}}}

\def\Aa{{\mathcal A}}
\def\Bb{{\mathcal B}}
\def\Cc{{\mathcal C}}
\def\Dd{{\mathcal D}}
\def\Ee{{\mathcal E}}
\def\Ff{{\mathcal F}}
\def\Gg{{\mathcal G}}
\def\Hh{{\mathcal H}}
\def\Ii{{\mathcal I}}
\def\Jj{{\mathcal J}}
\def\Kk{{\mathcal K}}
\def\Ll{{\mathcal L}}
\def\Mm{{\mathcal M}}
\def\Nn{{\mathcal N}}
\def\Oo{{\mathcal O}}
\def\Pp{{\mathcal P}}
\def\Qq{{\mathcal Q}}
\def\Rr{{\mathcal R}}
\def\Ss{{\mathcal S}}
\def\Tt{{\mathcal T}}
\def\Uu{{\mathcal U}}
\def\Vv{{\mathcal V}}
\def\Ww{{\mathcal W}}
\def\Xx{{\mathcal X}}
\def\Yy{{\mathcal Y}}
\def\Zz{{\mathcal Z}}

\newcommand{\CC}{\mathbf{C}}
\newcommand{\GG}{\mathbf{G}}
\newcommand{\KK}{\mathbf{K}}
\newcommand{\LL}{\mathbf{L}}
\newcommand{\NN}{\mathbf{N}}
\newcommand{\PP}{\mathbf{P}}
\newcommand\QQ{\mathbf{Q}}
\newcommand\RR{\mathbf{R}}
\newcommand\ZZ{\mathbf{Z}}

\def\to{\mathchoice{\longrightarrow}{\rightarrow}{\rightarrow}{\rightarrow}}
\def\from{\mathchoice{\longleftarrow}{\leftarrow}{\leftarrow}{\leftarrow}}
\def\mto{\mathchoice{\longmapsto}{\mapsto}{\mapsto}{\mapsto}}
\def\mfrom{\mathrel{\from\joinrel\mapstochar\,}}
\def\hto{\mathrel{\lhook\joinrel\to}}
\def\hfrom{\mathrel{\from\joinrel\rhook}}

\def\implies{\mathchoice{\Longrightarrow}{\Rightarrow}{\Rightarrow}{\Rightarrow}}
\def\implied{\mathchoice{\Longleftarrow}{\Leftarrow}{\Leftarrow}{\Leftarrow}}
\def\iff{\mathchoice{\Longleftrightarrow}{\Leftrightarrow}{\Leftrightarrow}{\Leftrightarrow}}

\def\To#1{\mathchoice{\xrightarrow{\labelstyle~~#1~}}{\stackrel{#1}{\longrightarrow}}{}{}}
\def\Mto#1{\mathrel{\mapstochar\To{#1}}}
\def\From#1{\mathchoice{\xleftarrow{\labelstyle~~#1~}}{\stackrel{#1}{\longleftarrow}}{}{}}

\newcommand{\isomto}{\overset{\sim}{\longrightarrow}}
\newcommand{\isomfrom}{\overset{\sim}{\longleftarrow}}

\newcommand{\IFF}{\unskip~$\iff$~\ignorespaces}
\newcommand{\THEN}{\unskip~$\implies$~\ignorespaces}

\let\leq\leqslant
\let\geq\geqslant
\let\le\leqslant
\let\ge\geqslant

\def\cf{{cf.}\kern.3em}
\def\Cf{{Cf.}\kern.3em}
\def\ie{{i.e.}\kern.3em}
\def\iev{{i.e.},\xspace}
%\def\loccito{loc.\kern2pt cit.\xspace}
%\def\loccitn{\textit{loc.\kern2pt cit.}\xspace}
%\def\loccit{\sisi{\loccito}{\loccitn}}
\def\EGA{\textup{EGA}\kern.2em}
\def\SGA{\textup{SGA}\kern.3em}

\let\oldS\S
\def\S{\oldS\kern .15em}
\def\ptbl{.\kern .2em}
\def\resp{resp.\kern.3em}
\def\Exp{Exp.\kern.3em}

\let\bar\overline
\let\hat\widehat
\let\tilde\widetilde
\let\dpl\displaystyle
\let\ldots\dots
\let\emptyset\varnothing
\let\skpt\skippointrait
\let\epsilon\varepsilon

\def\sfrac#1#2{{#1}/{#2}}

\DeclareMathOperator{\cd}{cd}
%\DeclareMathOperator{\Hhom}{\mathcal{H}\!\mathit{om}}
\DeclareMathOperator{\Hhom}{\mathbf{Hom}}
%\DeclareMathOperator{\Eext}{\mathcal{E}\!\mathit{xt}}
\DeclareMathOperator{\Eext}{\mathbf{Ext}}
\DeclareMathOperator{\Spec}{Spec}
\DeclareMathOperator{\Tr}{Tr}

\let\emptyset=\varnothing

%%%%%%%%%%%%%%%
%%% expose XVII
%%%%%%%%%%%%%%%

\makeatletter
\let\old@optexpose\optexpose
\newcommand{\empty@optexpose}[1]{}
\newcommand{\paren@optexpose}[1]{\old@optexpose{#1}(}
\newcommand{\refnoexpose}[1]{\let\optexpose\empty@optexpose\ref{#1}\let\optexpose\old@optexpose}%\refnoexpose{X.cor5.2} engendre "5.2" au lieu de "X 5.2"
\newcommand{\neweqref}[1]{\let\optexpose\paren@optexpose\textup{\ref{#1})}\let\optexpose\old@optexpose}%\neweqref{XII.eq4.1} engendre "XII (4.1)" au lieu de "(XII 4.1)"
\makeatother
\newcommand{\refvb}[2][VI${}_{\text{B}}$ ]{\sisi{#1\refnoexpose{#2}}{\ref{#2}}}

\newcommand{\leftexp}[2]{{\vphantom{#2}}^{#1}{#2}}
\newcommand{\AAA}{\mathbf{A}}
\newcommand{\Sf}{\mathfrak{S}}
\DeclareMathOperator*{\quotprojlim}{\text{``$\varprojlim$''}}
\DeclareMathOperator*{\quotinjlim}{\text{``$\varinjlim$''}}
\DeclareMathOperator{\Fl}{F\sisi{\ell}{l}}
\DeclareMathOperator{\Pl}{P\sisi{\ell}{l}}
%\DeclareMathOperator{\cl}{c\sisi{\ell}{l}}
\DeclareMathOperator{\tdf}{tdf}
\DeclareMathOperator{\rg}{rg}
\newcommand{\iInd}{\sisi{I}{i}nd}
\newcommand{\iind}{ind}
\newcommand{\ppro}{pro}
\newcommand{\tTor}{\sisi{T}{t}or}
\newcommand{\ttor}{tor}
\newcommand{\chap}{Chap.}
\newcommand{\Ffex}{\Ff\mathit{ex}}
\newcommand{\sheafk}{\underline{k}}
\newcommand{\sheafn}{\underline{n}}
\newcommand{\sheafGamma}{\underline{\Gamma}}

% Repeated references to an already typeset footnote must not use
% \footnotemark without a matching \footnotetext: Hyperref would then create
% a live annotation to a destination that cannot exist.  This helper keeps the
% current printed footnote mark and links it to an explicit stable note target.
\makeatletter
\newcommand{\sgarepeatedfootnotemark}[1]{%
  \begingroup
  \edef\@thefnmark{\thefootnote}%
  \hyperlink{#1}{\@makefnmark}%
  \endgroup}
\makeatother

\sisi{
\DeclareMathOperator{\simarrow}{\rule{.15mm}{1.8mm}\kern -3.4pt\rightsquigarrow}
\newcommand{\sheafell}{\underline{\ell}}
\newcommand{\sheafZZ}{\underline{\ZZ}}
}{
\newcommand{\nderef}[1]{la N.D.E.~(\ref{#1}) page~\pageref{#1}}
}

%%%%%%%%%%%%%%%%%
%%%%% expose XVII
%%%%%%%%%%%%%%%%%

% END INCLUDED SGA4_English_current_buildable_source_20260906/source/fichiers-configuration-latex/macros.tex


\externaldocument{../01/01_English.tex}
\externaldocument{../02/02_English.tex}
\externaldocument{../03/03_English.tex}
\externaldocument{../04/04_English.tex}
%\externaldocument{../05/05_English.tex}

\setboolean{original}{false}

\begin{document}

\begin{center}
V. Cohomology in Topoi\\
J.-L. Verdier
\end{center}

\selectlanguage{english}version: \showgitstatus
\tableofcontents

\else
\chapter{Cohomology in Topoi}\MakeLinkTarget*{sga4.exp.V}\label{sga4.exp.V}
%\addtocontents{toc}{J.-L. Verdier}
\begin{center}
J.-L. Verdier
\end{center}
\fi

\setcounter{pageoriginal}{1}

% BEGIN INCLUDED SGA4_English_current_buildable_source_20260906/source/05/drafts/SGA4_V_LINES_57_369_English.texfrag
\section*{Introduction}

In\pageoriginale this exposé we present the commutative and
elementary cohomological invariants of topoi. In ${\rm n^{\circ}}$
1, we study flat modules and flat morphisms of ringed topoi. The
proofs are given using the hypothesis, most often satisfied in
practice, that the topoi have enough points (IV 6). These proofs are
repeated in the general case in Appendix nᵒ 8, where Deligne, by
means of the technique of local inductive limits, also generalizes
to topoi D. Lazard's theorem on the structure of flat modules. The
theorems in this exposé\sisi{,}{} are existence theorems for spectral
sequences relating the various cohomological invariants (nᵒ 3, 5,
6). It is known that, even for topological spaces, Čech cohomology
does not in general coincide with sheaf cohomology \cite{sga4.exp.V.bib.11}. In
Appendix nᵒ 7, we introduce a modified Čech computation which makes
it possible to obtain the sheaf cohomology of an arbitrary topos by
means of coverings. In that appendix we are led to use simplicial
coverings (hypercoverings), whose homotopy invariants were studied
in \cite{sga4.exp.V.bib.1} (\cf also \cite{sga4.exp.V.bib.17}).

The cohomological invariants introduced here are elementary in the
sense that we do not use derived categories \cite{sga4.exp.V.bib.12}. A reader
familiar with this language will immediately translate the various
statements of this exposé and will then be able to generalize them
to complexes and hypercohomology. The language of derived
categories is, moreover, used later in this seminar.

Here we restrict ourselves to commutative cohomology. For
noncommutative $H^1$, which is used in this seminar, and for
noncommutative $H^2$, we refer to \cite{sga4.exp.V.bib.9}. The functors that we
derive are additive; we say nothing about multiplicative structures
(\cf \cite{sga4.exp.V.bib.7}).

\setcounter{section}{-1}
\section{Generalities on abelian categories}\etiquette[0]{V.sec0}

In\pageoriginale this section we recall several lemmas, most of
which can be found in \cite{sga4.exp.V.bib.11}.

\begin{enonce*}{Proposition 0.1}\etiquette[0.1]{V.prop0.1}
Let $\mathfrak{A}$ be an abelian category having a generator. The
following conditions are equivalent:
\begin{enumerate}
\renewcommand{\theenumi}{\roman{enumi}}
\renewcommand{\labelenumi}{\theenumi)}
\item The category $\mathfrak{A}$ satisfies axiom $AB$ $5$): Small
  direct sums are representable, and if $(X_i)$, $i \in I$, is a
  small increasing filtered family of subobjects of an object $X$
  of $\mathfrak{A}$ and $Y$ is a subobject of $X$, then
$$
(\displaystyle{\sup_i}\ X_i) \cap Y = \displaystyle{\sup_i}\ (X_i \cap Y) {\quoi}.
$$
\item Small pseudo-filtered inductive limits \emph{(\hyperref[I.def2.7]{I.2.7})} %% label \XXX
are representable and commute with finite projective limits.
\end{enumerate}

Moreover, if the preceding conditions hold, small filtered
inductive limits are universal \emph{(I \hyperref[I.prop2.6]{2.6})}.
\end{enonce*}

\begin{enonce*}[remark]{Proof}
It is clear that (ii) $\Rightarrow$ (i). To show that
\sisi{i)}{(i)} $\Rightarrow$ (ii), and to prove the additional
assertion, it suffices to use the fact that $\mathfrak{A}$ is a full
subcategory of a category of modules $\Mm$ over a suitable ring,
such that the inclusion functor $u : \mathfrak{A} \rightarrow \Mm$
admits an exact left adjoint $v$ \cite{sga4.exp.V.bib.5}. The verification is then
immediate.
\end{enonce*}

\setcounter{subsection}{1}
\subsubsection{}\etiquette[0.1.1]{V.subsubsec0.1.1}

It is known \cite{sga4.exp.V.bib.11} that an abelian category $\mathfrak{A}$
having a generator and satisfying axiom $AB$ 5) has enough
injectives, that is, every object embeds in an injective object.
Moreover, by the result already cited \cite{sga4.exp.V.bib.5}, small products are
representable in $\mathfrak{A}$ (axiom $AB$ 3) $^{\ast}$).

\begin{enonce*}{Proposition 0.2}\etiquette[0.2]{V.prop0.2}
Let $\mathfrak{A}$ and $\mathfrak{B}$ be two abelian categories and
let
$\xymatrix{\mathfrak{A} \ar@<-2pt>[r]_{u}& \mathfrak{B} \ar@<-2pt>[l]_{v}
}
$
be two adjoint functors ($u$ is left adjoint to $v$). Consider the
following two properties:
\begin{enumerate}
\renewcommand{\theenumi}{\roman{enumi}}
\renewcommand{\labelenumi}{\theenumi)}
\item The functor $u$ is exact.

\item The\pageoriginale functor $v$ sends injective objects of
$\mathfrak{B}$ to injective objects of $\mathfrak{A}$.
\begin{enumerate}
\renewcommand{\theenumii}{\arabic {enumii}}
\renewcommand{\labelenumii}{\theenumii)}
\item The implication \emph{(i) $\Rightarrow$ (ii)} always holds.

If, in addition, every nonzero object of $\mathfrak{B}$ is the
source of a nonzero morphism into an injective object, then
\emph{(ii) $\Leftrightarrow$ (i)}.
\item Suppose that:
\begin{enumerate}
\renewcommand{\theenumiii}{\alph {enumiii}}
\renewcommand{\labelenumiii}{\theenumiii)}
\item the category $\mathfrak{B}$ has enough injectives;
\item one of the two equivalent conditions {\rm(i)} and {\rm(ii)}
  above holds;
\item the functor $u$ is faithful.
\end{enumerate}
\end{enumerate}
\end{enumerate}
Then the category $\mathfrak{A}$ has enough injectives.
\end{enonce*}

\begin{enonce*}[remark]{Proof}
The proof is left to the reader.
\end{enonce*}

\begin{enonce*}[remark]{Remark 0.2.1}\etiquette[0.2.1]{V.rem0.2.1}
The category of commutative groups has enough injectives. Applying
Proposition \hyperref[V.prop0.2]{0.2}, we deduce that every category
of unital modules over a ring with identity has enough injectives.
Applying the result of \cite{sga4.exp.V.bib.5} (used in the proof of
\hyperref[V.prop0.1]{0.1}) together with
\hyperref[V.prop0.2]{0.2}, we then deduce that every abelian category
having a generator and exact small filtered inductive limits has
enough injectives; this gives a new proof of that fact.
\end{enonce*}

\begin{enonce*}{Proposition 0.3}\etiquette[0.3]{V.prop0.3}
Let $\mathfrak{A}$, $\mathfrak{B}$, and $\mathfrak{C}$ be three
abelian categories, and let $u : \mathfrak{A} \rightarrow
\mathfrak{B}$ and $v : \mathfrak{B} \rightarrow \mathfrak{C}$ be
two left exact additive functors. Suppose that $\mathfrak{A}$ and
$\mathfrak{B}$ have enough injective objects. The following two
properties are equivalent:
\begin{enumerate}
\renewcommand{\theenumi}{\roman{enumi}}
\renewcommand{\labelenumi}{\theenumi)}
\item There exists a spectral functor whose $E^{p, q}_2$ term is
$$
R^p vR^qu
$$
and which converges to $R^{p+q} vu$ (suitably filtered).
\item The functor $u$ sends injective objects of $\mathfrak{A}$ to
  objects that are acyclic for the functor $v$.
\end{enumerate}
\end{enonce*}

\begin{enonce*}[remark]{Proof}
(i) $\Rightarrow$ (ii) is immediate, since it suffices to apply the
spectral functor to an\pageoriginale injective object. The
implication (ii) $\Rightarrow$ (i) is proved in \cite{sga4.exp.V.bib.11}.
\end{enonce*}

\begin{enonce*}{Proposition 0.4}\etiquette[0.4]{V.prop0.4}
Let $\mathfrak{A}$ and $\mathfrak{B}$ be two abelian categories and
let $u : \mathfrak{A} \rightarrow \mathfrak{B}$ be a %% correct word break?
left exact additive functor. Let $M$ be a subset of the set of
objects of $\mathfrak{A}$ having the following properties:
\begin{enumerate}
\renewcommand{\labelenumi}{\theenumi)}
\item Every object of $\mathfrak{A}$ embeds in an element of $M$.

\item If $X \oplus Y$ belongs to $M$, then the object $X$ belongs to
  $M$.

\item If
$$
0 \longrightarrow X' \longrightarrow X \longrightarrow X''
\longrightarrow 0
$$
is an exact sequence and $X'$ and $X$ belong to $M$, then $X''$
belongs to $M$, and the sequence
$$
0 \longrightarrow u(X') \longrightarrow u(X) \longrightarrow u(X'')
\longrightarrow 0
$$
is exact. Zero objects belong to $M$.
\end{enumerate}

Then every injective belongs to $M$, and the objects of $M$ are
acyclic for $u$, that is, for every $q \neq 0$ and every object $X$
of $M$, one has $R^q u(X) = 0$. (In particular, resolutions by
objects of $M$ can be used to compute the derived functors of $u$.)
\end{enonce*}

For the proof\sisi{}{,} see \cite{sga4.exp.V.bib.11} 3.3.1.

\begin{enonce*}{Proposition 0.5}\etiquette[0.5]{V.prop0.5}
Let $𝒰 \subset 𝒱$ be two universes, let $\mathfrak{A}$ (\resp
$\mathfrak{B}$) be a $𝒰$-category (\resp $𝒱$-category) which is
abelian, satisfies axiom $AB$ 5) relative to $𝒰$ (\resp to $𝒱$),
and has a $𝒰$-small (\resp $𝒱$-small) generating family. Let
$\epsilon : \mathfrak{A} \to \mathfrak{B}$ be an exact fully
faithful functor. The following conditions are equivalent:
\begin{itemize}
\item[1)] There exists a generating family $(X_i)_{i\in I}$ of
  $\mathfrak{A}$ such that the family $(\epsilon(X_i))_{i\in I}$ is
  generating in $\mathfrak{B}$.

\item[1′)] Every object of $\mathfrak{B}$ is isomorphic to a
  quotient of an object of the form
  $\displaystyle{⨁_{\alpha \in A}} \epsilon (Y_\alpha)$, where
  $A \in 𝒱$.

Under these conditions, the following property holds:

\item[2)] $\epsilon$ sends $𝒰$-small products to products (and thus
  commutes with $𝒰$-small projective limits).
\end{itemize}

Moreover, under the equivalent conditions 1) or 1′), the following
conditions are equivalent\pageoriginale:
\begin{itemize}
\item[a)] For every object $Y$ of $\mathfrak{A}$, every subobject of
  $\epsilon(Y)$ is isomorphic to the image under $\epsilon$ of a
  subobject of $Y$.

\item[a′)] There exists a generating family $(X_i)_{i \in I}$ of
  $\mathfrak{A}$ such that the family $(\epsilon(X_i))_{i \in I}$
  is generating in $\mathfrak{B}$ and such that, for every $i \in
  I$, every subobject of $\epsilon(X_i)$ is isomorphic to the image
  under $\epsilon$ of a subobject of $X_i$.

\item[b)] Every object of $\mathfrak{B}$ is isomorphic to a
  subobject of an object of the form
  $\prod_{\alpha \in A} \epsilon(Y_\alpha)$, where $A \in 𝒱$.

\item[c)] $\epsilon$ commutes with $𝒰$-small direct sums (and thus
  commutes with $𝒰$-small inductive limits).

Moreover, under conditions 1) and a′), one has:

\item[d)] $\epsilon$ sends injective objects to injective objects.
\end{itemize}
\end{enonce*}

\begin{enonce*}[remark]{Remark 0.5.1}\etiquette[0.5.1]{V.rem0.5.1}
When in \hyperref[V.prop0.5]{0.5} one has $𝒰 = 𝒱$, conditions 1) and
a′) imply that $\epsilon$ is an equivalence of categories (since b),
2), and a) then hold).
\end{enonce*}

\setcounter{subsection}{5}
\subsubsection{}\etiquette[0.5.2]{V.subsubsec0.5.2}
We will confine ourselves to giving indications for the proof. The
implications 1) $\Leftrightarrow$ $1'$) and 1) $\Rightarrow$ 2) are
left to the reader. It is clear that a) $\Rightarrow$ a′). We show
that a′) implies d). Let $M$ be injective in $\mathfrak{A}$. For
every $i \in I$ and every subobject $Y \hookrightarrow X_i$ of
$X_i$, the homomorphism $\Hom(X_i, M) \rightarrow \Hom(Y, M)$ is
surjective. Consequently, by a′) and the full faithfulness of
$\epsilon$, for every subobject $U \hookrightarrow \epsilon(X_i)$,
the homomorphism $\Hom(\epsilon (X_i), \epsilon(M)) \to \Hom(U,
\epsilon(M))$ is surjective. Since the $\epsilon(X_i)$ form a
generating family of $\mathfrak{B}$, $\epsilon(M)$ is injective
\cite{sga4.exp.V.bib.11}.

We show that a′) $\Rightarrow$ b). By enlarging the family of the
$X_i$ if necessary, we may suppose that, for every $i \in I$, every
quotient of $X_i$ is isomorphic to an $X_j$ for a suitable $j$. For
every $i \in I$, let $X_i \hookrightarrow M_i$ then be a
monomorphism into an injective object. The family
$(\epsilon (X_i))_{i\in I}$ is stable under quotients, and for $i
\in I$, the morphism $\epsilon(X_i) \hookrightarrow
\epsilon(M_i)$ is, by d), a monomorphism into an injective object.
It is then immediately verified that, since the family
$\epsilon(X_i)$ is generating, the family $(\epsilon(M_i))_{i \in
I}$ is cogenerating, whence b).

We show that b) $\Rightarrow$ c). Let $(Z_\alpha)_{\alpha \in A}$
be a $𝒰$-small family of objects of $\mathfrak{A}$, and let us show
that the canonical morphism
$$
⨁_{\alpha\in A} \epsilon(Z_\alpha) \to \epsilon(⨁_{\alpha \in A} Z_\alpha)
$$
is\pageoriginale an isomorphism. Since the family of objects
$\epsilon(Y)$ is cogenerating, it suffices to show that, for every
object $Y$ of $\mathfrak{A}$, the homomorphism
$$
\Hom(\epsilon(⨁_a Z_a), \epsilon(Y)) \to
\Hom(⨁_a \epsilon(Z_a), \epsilon(Y))
$$
is an isomorphism, which follows from the full faithfulness of
$\epsilon$.

It remains to show that c) $\Rightarrow$ a). Let $Z$ be a subobject
of $\epsilon(Y)$. By $1'$), there exists a $𝒰$-small family
$Y_\alpha$ of objects of $\mathfrak{A}$ and an epimorphism from
$\bigoplus_{\alpha} \epsilon(Y_\alpha)$ onto $Z$. Since $\epsilon$
commutes with $𝒰$-small direct sums, there is therefore an
epimorphism $\epsilon(Y') \to Z$. Denote by $u : \epsilon(Y') \to Z
\to \epsilon(Y)$ the composite morphism. We have $Z =
\text{Im}(u)$. Since \sisi{}{$\epsilon$} is fully faithful, we have
$u \sisi{:}{=} \epsilon(v)$, and consequently
$Z = \text{Im}(\epsilon(\sisi{V}{v})) = \epsilon(\text{Im}(v))$.


\begin{enonce*}[remark]{Exercise 0.5.2}\etiquette[0.5.2]{V.exe0.5.2}
Let $\Sex_{𝒱} (\mathfrak{A})$ be the category of contravariant
functors from $\mathfrak{A}$ to the category of $𝒱$-commutative
groups which commute with $𝒰$-small inductive limits. Show that,
under conditions 1) and a′), the canonical functor from
$\mathfrak{B}$ to $\Sex_{𝒱}(\mathfrak{A})$ is an equivalence.
\end{enonce*}


% END INCLUDED SGA4_English_current_buildable_source_20260906/source/05/drafts/SGA4_V_LINES_57_369_English.texfrag


% BEGIN INCLUDED SGA4_English_current_buildable_source_20260906/source/05/drafts/SGA4_V_LINES_370_602_English.texfrag
\section{Flat modules}\etiquette[1]{V.sec1}

\begin{enonce*}{Definition 1.1}\etiquette[1.1]{V.def1.1}
Let $(E,A)$ be a ringed topos \emph{(IV \hyperref[IV.subsubsec11.1.1]{11.1.1})}. %% label
A right (\resp left) $A$-module $M$ is said to be flat if the functor
$M\otimes_A⋅$ (\resp $⋅\otimes_A M$) from the category of left
(\resp right) $A$-modules to the category of abelian sheaves of
$E$\sisi{,}{} is exact.
\end{enonce*}

\begin{enonce*}{Proposition 1.2}\etiquette[1.2]{V.prop1.2}
Let $M$ be a $B$-$A$-bimodule.
\begin{enumerate}
\renewcommand{\labelenumi}{\theenumi)}
\item The following properties are equivalent:
\renewcommand{\theenumii}{\roman{enumii}}
\renewcommand{\labelenumii}{\theenumii)}
\begin{enumerate}
\item The module $M$ is flat as a right $A$-module.
\item For every injective $B$-module $I$, the left $A$-module
  $\Hhom_B(M,I)$ is injective.
\end{enumerate}
\item A module $M$ which is a pseudo-filtered inductive limit
\emph{(I \hyperref[I.exem2.7.1]{2.7.1})} %% label
of flat modules is flat.

\item Finally, if
$M_∙=\cdots\sisi{}{\to}M_{i+1}\to M_i\sisi{}{\to}\cdots$ is an
acyclic complex of flat modules $(M_i=0$ for $i<i_0)$, then, for every
module $F$, the complex
$$
M_∙\otimes_A F=\cdots\sisi{}{\to}M_{i+1}\otimes_A F\to
M_i\otimes_A F\sisi{}{\to}\cdots
$$
is\pageoriginale acyclic.
\end{enumerate}
\end{enonce*}

\begin{enonce*}[remark]{Proof}
By (IV \hyperref[IV.cor12.12]{12.12}), %% label
there is an adjunction isomorphism:
$$
\Hom_B(M\otimes_A\mathord\cdot,\mathord\cdot)\xrightarrow{\sim}
\Hom_A(\mathord\cdot,\Hhom_B(M,\mathord\cdot)){\quoi}.
$$
It is then enough to apply \hyperref[V.prop0.2]{0.2} to obtain the
equivalence (i) $\Leftrightarrow$ (ii). This adjunction isomorphism
also shows that the tensor product commutes with inductive limits.
The fact that pseudo-filtered inductive limits are exact
(\hyperref[V.prop0.1]{0.1}) implies the second assertion. To show that
the complex $M_∙\otimes_A F$ is acyclic, it is enough to show that,
for every injective abelian sheaf $I$, the complex
$\Hom^{∙}_{\ZZ}(M_∙\otimes_A F,I)$ is acyclic. By the adjunction
formulas, this complex is isomorphic to
$\Hom^{∙}_A(F,\Hhom_{\ZZ}(M_∙,I))$. But, by the equivalence
(i) $\Leftrightarrow$ (ii), the complex of sheaves
$$
\Hhom_{\ZZ}(M_∙,I)
$$
is an acyclic complex whose objects are injective, whence the
conclusion.
\end{enonce*}

\begin{enonce*}{Proposition 1.3.1}\etiquette[1.3.1]{V.prop1.3.1}
Let $(E,A)$ be a ringed topos, let $H$ be an object of $E$, and let
$M$ be a flat $A_{/H}$-module. Then $j_{H!}M$ is a flat $A$-module.
In particular, $A_H$ is flat on both the right and the left.
\end{enonce*}

Suppose, to fix ideas, that $M$ is a right $A_{/H}$-module. For every
left $A$-module $P$, there is a canonical isomorphism (IV 12) %% label
$$
j_{H!}M\otimes_A P\simeq
j_{H!}(M\otimes_{A_{/H}}j_H^{\ast}P){\quoi}.
$$
The functors $j_{H!}$ and $j_H^{\ast}$ are exact (IV \hyperref[IV.prop11.3.1]{11.3.1} and
\hyperref[IV.subsubsec11.2.2]{11.12.2}), and, by hypothesis, the functor
$M\otimes_{A_{/H}}-$ is exact. Consequently, the functor
$P\longmapsto j_{H!}M\otimes_A P$ is exact, and $j_{H!}M$ is flat.

\begin{enonce*}{Proposition 1.3.2}\etiquette[1.3.2]{V.prop1.3.2}
(Projection formula for closed immersions):

Let $(E,A)$ be a ringed topos, and let $i:F\to E$ be a closed subtopos
of $E$. Set $i^{\ast}A=A_{/F}$. For every right $A_{/F}$-module $M$
and every left $A$-module $P$, there\pageoriginale is a functorial
isomorphism
$$
i_\ast(M\otimes_{A_{/F}}i^\ast P)\simeq(i_\ast M)\otimes_A P.
$$
\end{enonce*}

Let $U$ be the open complement of $F$, and let $j:U\to E$ be the
canonical immersion morphism. One has
$j^{\ast}(i_\ast M\otimes_A P)\simeq
0\otimes_{A_{/U}}j^\ast(P)$ (IV 12); consequently,
$i_\ast M\otimes_A P$ has its support in $F$. Hence the adjunction
morphism
$i_\ast M\otimes_A P\to i_\ast i^\ast(i_\ast M\otimes_A P)$ is an
isomorphism (IV 14). One has
$i^{\ast}(i_\ast M\otimes_A P)\simeq
i^\ast i_\ast M\otimes_{A_{/F}}i^\ast P$ (IV 12), and, since
$i:F\to E$ is a closed immersion,
$i^{\ast}i_\ast M\simeq M$ (IV 14); this gives the asserted
isomorphism.

\begin{enonce*}{Corollary 1.3.3}\etiquette[1.3.3]{V.cor1.3.3}
For every flat $A_{/F}$-module $M$, $i_\ast M$ is flat.
\end{enonce*}

It follows from \hyperref[V.prop1.3.2]{1.3.2} and from (IV 14) that
the functor $P\longmapsto(i_\ast M)\otimes_A P$ is exact.

\setcounter{subsection}{3}
\subsection{}\etiquette[1.4]{V.subsec1.4}
Let $(E,A)$ be a ringed topos, let $x:P\to E$ be a point of $E$
(IV \hyperref[IV.def6.1]{6.1}), and let $\Vois(x)$ be the category of neighborhoods of $x$
(IV \hyperref[IV.subsec6.8]{6.8}). To every object $V$ of $\Vois(x)$ there corresponds an
object of $E$, also denoted by $V$, and a point
$x_V:P\to E/V$ of $E/V$. Moreover, to every morphism
$u:V\to W$ in $\Vois(x)$ there corresponds an essentially
commutative diagram of morphisms of topoi (IV \hyperref[IV.subsec6.7]{6.7})
\[\vcenter{
\xymatrix@C=2.5cm@R=1.5cm{
P \ar[r]^{x_V} \ar[rd]_{x_W}& E/V \ar[d]^{j_u}\\
& E/W\quad .
}}\leqno{(1.4.1)}
\]\etiquette[1.4.1]{V.eq1.4.1}
Let $N$ be an $A_x$-module (\resp a set). From the diagram
\eqhyperref[V.eq1.4.1]{1.4.1}, one obtains a diagram of $A_x$-modules
(\resp sets):
\[\vcenter{
\xymatrix@C=2cm@R=1.5cm{
x^{\ast}_W x_{W\ast}N \ar[r]^{\ad_W} \ar[d]_{\phi(u)}& N\\
x^{\ast}_V x_{V\ast}N \ar[ru]_{\ad_V}&
}}\leqno{(1.4.2)}
\]\etiquette[1.4.2]{V.1.4.2}
where\pageoriginale $\ad_W$ and $\ad_V$ are the adjunction morphisms.
Moreover, one immediately checks that
$\phi(uv)=\phi(v)\phi(u)$. Thus there is a functor from the filtered
category $\Vois(x)^{\circ}$ to the category of $A_x$-modules
(\resp sets), and a homomorphism
\begin{equation}
\Lambda(N):\varinjlim_{V\in\Vois(x)^{\circ}}
x^{\ast}_Vx_{V\ast}N\to N{\quoi}.\tag{1.4.3}
\end{equation}\etiquette[1.4.3]{V.1.4.3}

\begin{enonce*}{Proposition 1.5}\etiquette[1.5]{V.prop1.5}
For every $A_x$-module (\resp set) $N$, $\Lambda(N)$ is an
isomorphism.
\end{enonce*}

This proposition is a special case of a proposition due to Deligne
(\hyperref[V.prop8.2.6]{8.2.6}). Let us give a direct proof. By passing
to the underlying sets, it is enough to prove the proposition when
$N$ is a set (I \hyperref[I.prop2.8]{2.8}). %% label
The cofiltered category $\mathrm{Fl}(\Vois(x))$ of arrows of
$\Vois(x)$ is fibered over the category $\Vois(x)$ by the functor
$p:\mathrm{Fl}(\Vois(x))\to\Vois(x)$ which assigns to an arrow its
target. Let $D$ be a category and
$F:\mathrm{Fl}(\Vois(x))\to D$ a pro-object (I \hyperref[I.sec8.10]{8.10}). %% label
For every object $V$ of $\Vois(x)$, denote by $F_V$ the pro-object
obtained by restricting the functor $F$ to the fiber category
$\Vois(x)/V$. To every morphism $m:U\to V$ of $\Vois(x)$, the
base-change functor by $m$ associates a morphism from the pro-object
$F_U$ to the pro-object $F_V$, and the canonical morphisms of
pro-objects $F\to F_V$ determine a morphism of pro-objects
\begin{equation*}
F\to\varprojlim_{V\in\Vois(x)^{\circ}}F_V\tag{1.5.1}
\end{equation*}\etiquette[1.5.1]{V.1.5.1}
which is immediately checked to be an isomorphism. Apply this
observation to the pro-object of pointed sets
\begin{equation*}
(U,V,m)\longmapsto x^{\ast}_V(m)=F(U,V,m){\quoi}.\tag{1.5.2}
\end{equation*}\etiquette[1.5.2]{V.1.5.2}
One obtains an isomorphism of pro-objects
\begin{equation*}
F\xrightarrow{\sim}\varprojlim_{\Vois(x)}
\varprojlim_{\Vois(x)/V}x^{\ast}_V(m){\quoi};\tag{1.5.3}
\end{equation*}\etiquette[1.5.3]{V.1.5.3}
hence, for every set $N$, a bijection
\begin{equation*}
\varinjlim_{\mathrm{Fl}(\Vois(x))^{\circ}}
\Hom(x^{\ast}_V(m),N)\simeq
\varinjlim_{\Vois(x)^{\circ}}
\varinjlim_{(\Vois(x)/V)^{\circ}}
\Hom(x^{\ast}_V(m),N).\tag{1.5.4}
\end{equation*}\etiquette[1.5.4]{V.1.5.4}
But,\pageoriginale for fixed $V$,
$\varinjlim_{(\Vois(x)/V)^{\circ}}\Hom(x^{\ast}_V(m),N)$ is
identified with $x^{\ast}_Vx_{V\ast}N$. Indeed, by IV \hyperref[IV.eq6.8.1]{6.8.1} one has
$x^{\ast}_Vx_{V\ast}N\simeq
\varinjlim_{\Vois(x_V)^\circ}\Hom_{E/V}(W,x_{V\ast}N)$, and hence,
by adjunction,
$x^{\ast}_Vx_{V\ast}N\simeq
\varinjlim_{\Vois(x_V)^\circ}\Hom(x^{\ast}_V(W),N)$; moreover, the
category $\Vois(x_V)$ is equivalent to the category $\Vois(x)/V$
(IV \hyperref[IV.eq6.7.2]{6.7.2}). Finally, the canonical transition maps
\[
\varinjlim_{(\Vois(x)/U)^{\circ}}\Hom(x^{\ast}_U(m),N)\to
\varinjlim_{(\Vois(x)/V)^{\circ}}\Hom(x^{\ast}_V(n),N)
\]
in
\eqhyperref[V.1.5.4]{1.5.4} are identified with the canonical maps
$x^{\ast}_Ux_{U\ast}N\to x^{\ast}_Vx_{V\ast}N$ of
\eqhyperref[V.1.4.3]{1.4.3}, as the reader may verify. Thus there is a
bijection
\begin{equation*}
\varinjlim_{\mathrm{Fl}(\Vois(x))^{\circ}}
\Hom(x^{\ast}_V(m),N)\simeq
\varinjlim_{\Vois(x)^{\circ}}x^{\ast}_Vx_{V\ast}N\tag{1.5.5}
\end{equation*}\etiquette[1.5.5]{V.1.5.5}
and the map
$\Lambda(N):\varinjlim_{\Vois(x)^{\circ}}x^{\ast}_Vx_{V\ast}N\to N$
of \eqhyperref[V.1.4.3]{1.4.3}, composed with the bijection
\eqhyperref[V.1.5.5]{1.5.5}, comes from the maps
$\Hom(x^{\ast}_V(m),N)\to N$ which assign to a map
$r:x^{\ast}_V(m)\to N$ the image under $r$ of the distinguished point
of $x^{\ast}_V(m)$. To prove the proposition, it is then enough to
observe that every object $(U,V,m)$ of
$\mathrm{Fl}(\Vois(x))$ is bounded below by $(U,U,\id_U)$, and that
$x^{\ast}_U(\id_U)$ consists of a single element or, in other words,
that the canonical morphism from the constant pro-object consisting
of a single element to $F$ is an isomorphism of pro-objects.

\begin{enonce*}{Proposition 1.6}\etiquette[1.6]{V.prop1.6}
Let $(E,A)$ be a ringed topos, and let $M$ be an $A$-module.
\begin{enumerate}
\renewcommand{\labelenumi}{\theenumi)}
\item If $M$ is flat, then, for every point $x:P\to E$ of $E$, the
  $A_x$-module $M_x$ is flat.
\item Let $(x_i)_{i\in I}$ be a conservative family of points of $E$
  such that, for every $i\in I$, $M_{x_i}$ is a flat $A_{x_i}$-module.
\end{enumerate}
Then $M$ is flat.
\end{enonce*}

\begin{enonce*}{Lemma 1.6.1}\etiquette[1.6.1]{V.lem1.6.1}
Let $M$ be a flat module and let $V$ be an object of $E$. The
$A/V$-module $j^{\ast}_VM$ is flat.
\end{enonce*}

We must show that the functor
$P\longmapsto P\otimes_{A/V}j^{\ast}_VM$ is exact. Since the
extension-by-zero functor $j_{V!}$ is exact and faithful
(IV \hyperref[IV.prop11.3.1]{11.3.1}), it is enough to show that the functor
$P\longmapsto j_{V!}(P\otimes_{A/V}j^{\ast}_VM)$ is exact. There is a
functorial isomorphism\pageoriginale
$j_{V!}(P\otimes_{A/V}j^{\ast}_VM)\simeq
j_{V!}(P)\otimes_A M$ (IV \hyperref[IV.prop12.11]{12.11}); consequently, the functor
$P\longmapsto j_{V!}(P)\otimes_A M$ is exact (IV \hyperref[IV.prop11.3.1]{11.3.1}).

\setcounter{subsection}{6}
\setcounter{subsubsection}{1}

% END INCLUDED SGA4_English_current_buildable_source_20260906/source/05/drafts/SGA4_V_LINES_370_602_English.texfrag


% BEGIN INCLUDED SGA4_English_current_buildable_source_20260906/source/05/drafts/SGA4_V_LINES_603_925_English.texfrag
\subsubsection{}\etiquette[1.6.2]{V.subusbsec1.6.2}
Let us prove 1). To fix ideas, suppose that $M$ is a flat left
$A$-Module, and let $x:P\to E$ be a point of $E$. We shall show
that the functor $N\mapsto N\otimes_{A_x}M_x$ from the category of
right $A_x$-modules to the category of commutative groups is exact.
For this it suffices to show that this functor is left exact. With
the notation of \hyperref[V.subsec1.4]{1.4}, let $V$ be a
neighborhood of $x$, and let $j_V:E/V\to E$ be the localization
morphism. For every $A_x$-module $N$, we have
$x_V^\ast x_{V\ast}N\otimes_{A_x}M_x\simeq
x_V^\ast x_{V\ast}N\otimes_{A_x}x_V^\ast j_V^\ast M\simeq
x_V^\ast(x_{V\ast}N\otimes_{A/V}j_V^\ast M)$ (IV \hyperref[IV.prop12.11]{12.11}).
Consequently, the functor
$N\mapsto x_V^\ast x_{V\ast}N\otimes_{A_x}M_x$ is left exact. It
follows that the functor $N\mapsto N\otimes_{A_x}M_x$, being a
filtered inductive limit of left exact functors
(\hyperref[V.prop1.5]{1.5}), is left exact.

Let us prove 2). Let
$0\to P'\to P\to P''\to0$ be an exact sequence of $A$-Modules, and
let us show that the sequence
$0\to P'\otimes_A M\to P\otimes_A M\to P''\otimes_A M\to0$ is
exact. It suffices to show that, for every $i\in I$, the sequence
obtained by passing to the fibers at $x_i$ is exact (IV 6). But the
sequence of fibers at $x_i$ is the sequence (IV \hyperref[IV.cor13.5]{13.5})
$$
0\to P'_{x_i}\otimes_{A_{x_i}}M_{x_i}\to
P_{x_i}\otimes_{A_{x_i}}M_{x_i}\to
P''_{x_i}\otimes_{A_{x_i}}M_{x_i}\to0,
$$
and since $M_{x_i}$ is flat, this sequence is exact.

\begin{enonce*}{Proposition 1.7}\etiquette[1.7]{V.prop1.7}
Let $(E,A)$ be a ringed topos, let $u:E'\to E$ be a morphism of
topoi, and put $A'=u^{\ast}A$. Let $M$ be a right (\resp left)
$A'$-Module. The following conditions are equivalent:
\begin{enumerate}
\renewcommand{\theenumi}{\roman{enumi}}
\renewcommand{\labelenumi}{\theenumi)}
\item The functor $N\mapsto M\otimes_{A'}u^{\ast}N$ (\resp
$N\mapsto u^{\ast}N\otimes_{A'}M$) from the category of left
(\resp right) $A$-Modules to the category of abelian sheaves on $E'$
is exact.
\item $M$ is a flat $A'$-Module.
\end{enumerate}
\end{enonce*}

It is clear that ii) $\Rightarrow$ i).

Let us prove that i) $\Rightarrow$ ii). We give the proof only in the
case where $E'$ possesses a conservative family of points
$(x_i)_{i\in I}$. The general case is treated in
(\hyperref[V.prop8.2.7]{8.2.7}). In this particular case, which
covers most applications, we reduce\pageoriginale to the case where
$E'$ is the punctual topos and hence where the morphism $u$, which we
then denote by $x$, is a point of $E$
(\hyperref[V.prop1.6]{1.6} and IV \hyperref[IV.prop11.3.1]{11.3.1}). We restrict ourselves to
the case where $M$ is a right $A'$-Module. The case where $M$ is a
left $A'$-Module reduces to this one by passing to opposite rings.
Let $V$ be a neighborhood of $x$, and let $V_x$ be its fiber at $x$.
With the notation of \hyperref[V.subsec1.4]{1.4}, there is an
essentially commutative diagram of morphisms of topoi:
\[
\xymatrix@C=2.25cm@R=1.5cm{
P \ar[r]^{j} \ar[rd]_{x_V}& P/V_{x}\ar[r]^{j_{V_x}} \ar[d]^{x/V}& P\ar[d]^{x}\\
  & E/V \ar[r] & E\quad \sisi{,}{}
}
\]
where $x/V$ is the morphism deduced from $x$ by localization over
$V$ (IV \hyperref[IV.subsec5.10]{5.10}), and $j$ is the morphism deduced from the marked point
of $V_x$. Let us show that the functor
$N'\mapsto M\otimes_{A'}x_V^{\ast}N'$ is exact. We have
$x_V^{\ast}=j^{\ast}(x/V)^{\ast}$ and
$\id=j^{\ast}j_{V_x}^{\ast}$; consequently, there is a canonical
isomorphism
$$
M\otimes_{A'}x_V^{\ast}N'\simeq
j^{\ast}\bigl(j_{V_x}^{\ast}M\otimes_{A'/V_x}(x/V)^{\ast}N'\bigr).
$$
Since $j^{\ast}$ is exact and the functor $j_{V_x!}$ is exact and
faithful (IV \hyperref[IV.prop11.3.1]{11.3.1}), it suffices to show that the functor
$$
N'\longmapsto j_{V_x!}\Bigl(j_{V_x}^{\ast}M\otimes_{A'/V_x}
(x/V)^{\ast}N'\bigr)
$$
is exact; consequently (IV \hyperref[IV.prop12.11]{12.11}), it suffices to show that the
functor $N'\mapsto M\otimes_{A'}x^{\ast}j_{V!}N'$ is exact. But it
follows from the definition of the inverse image of an induced topos
(\hyperrefext[IV][IV.eq5.10.2]{5.10.2}) that
$j_{V_x!}(x/V)^{\ast}N'$ is equal to $x^{\ast}j_{V!}N'$, and since
the functor $j_{V!}$ is exact
(\hyperrefext[IV][IV.prop11.3.1]{11.3.1}), the functor
$N'\mapsto M\otimes_{A'}j_{V_x!}(x/V)^{\ast}N'$ is exact by
hypothesis.

For every $A'$-module $Q$, there is a functorial isomorphism
(\hyperref[V.prop1.5]{1.5}):
$$
Q\simeq \varinjlim_{\Vois(x)^{\circ}}x_V^{\ast}x_{V\ast}Q{\quoi}.
$$
By what precedes, the functor $Q\mapsto M\otimes_{A'}Q$ is a
filtered inductive limit of left exact functors. It is therefore
left exact, and hence $M$ is flat.

\begin{enonce*}{Corollary 1.7.1}\etiquette[1.7.1]{V.cor1.7.1}
Let $u:(E',A')\to(E,A)$ be a morphism of ringed topoi, and let $M$
be a flat $A$-Module. Then $u^{\ast}M$ is a flat $A'$-Module.
\end{enonce*}

This follows from the criterion given in \hyperref[V.prop1.7]{1.7}
and formula (IV \hyperlink{sga4.exp.IV.equation.13.4.5}{13.4.5}).

\begin{enonce*}{Corollary 1.7.2}\etiquette[1.7.2]{V.cor1.7.2}
Let\pageoriginale $u:(E',A')\to(E,A)$ be a morphism of ringed
topoi. The following conditions are equivalent:
\begin{enumerate}
\renewcommand{\theenumi}{\roman{enumi}}
\renewcommand{\labelenumi}{\theenumi)}
\item The right (\resp left) $u^{-1}(A)$-Module $A'$ is flat.
\item The functor $u^{\ast}$ from the category of left (\resp right)
$A$-Modules to the category of left (\resp right) $A'$-Modules is
exact.
\end{enumerate}
\end{enonce*}

The equivalence follows from \hyperref[V.prop1.7]{1.7} and the
definition of $u^{\ast}$ (IV \hyperref[IV.subsubsec13.2.1]{13.2.1}).

\begin{enonce*}{Definition 1.8}\etiquette[1.8]{V.def1.8}
A morphism of ringed topoi having the equivalent properties of
\hyperref[V.cor1.7.2]{1.7.2} is called a morphism of ringed topoi
which is left (\resp right) flat.
\end{enonce*}

\begin{enonce*}{Proposition 1.9}\etiquette[1.9]{V.prop1.9}
Let $𝒰\subset 𝒱$ be two universes, let $C$ be a $𝒰$-site, let
$C^{\sim}_{𝒰}$ and $C^{\sim}_{𝒱}$ be the categories of $𝒰$- and
$𝒱$-sheaves of sets, respectively, let
$\epsilon:C^{\sim}_{𝒰}\to C^{\sim}_{𝒱}$ be the canonical inclusion
functor, and let $A$ be a $𝒰$-sheaf of rings on $C$.
\begin{enumerate}
\renewcommand{\labelenumi}{\theenumi)}
\item The functor $\epsilon$ commutes with $𝒰$-small inductive and
projective limits of Modules. The functor $\epsilon$ is conservative
and fully faithful on the categories of Modules.

\item For every object $H$ of $C^{\sim}_{𝒰}$, there is a canonical
isomorphism $\epsilon(A_H)\simeq \epsilon(A)_{\epsilon(H)}$.

\item Every injective left or right $A$-Module is carried by
$\epsilon$ to an injective $\epsilon(A)$-Module.

\item Every flat right or left $A$-Module is carried by $\epsilon$
to a flat $\epsilon(A)$-Module.

\item The functor $\epsilon$ commutes with the formation of tensor
products and sheaves of homomorphisms.
\end{enumerate}
\end{enonce*}

The formation of the associated sheaf is independent of the
universe (II \hyperref[II.prop3.6]{3.6}). It then follows from the construction of
inductive and projective limits in categories of sheaves (II \hyperref[II.theo4.1]{4.1}
and \hyperref[II.prop6.4]{6.4}) that the functor $\epsilon$ commutes with $𝒰$-small
projective limits of sets or Modules, which gives 1). Let us prove
2). By taking a $𝒰$-small generating subcategory of
$C^{\sim}_{𝒰}$, we may reduce to the case where $C$ is
$𝒰$-small\pageoriginale (III \hyperref[III.theo4.1]{4.1}). It then follows from (II \hyperref[II.prop6.5]{6.5})
that the free $A$-Module generated by $H$ is obtained by forming the
presheaf of free $A$-modules generated by $H$, a construction that
commutes with enlargement of universes, and then forming the
associated sheaf, an operation which likewise commutes with
enlargement of universes (II \hyperref[II.prop3.6]{3.6}). To prove 3), it suffices, by
\hyperref[V.prop0.5]{0.5}, to show that every $𝒱$-subsheaf of a
$𝒰$-sheaf is a $𝒰$-sheaf, which is clear.

Let us prove 4). Let $M$ be a flat $A$-Module. It suffices to show
that, for every injective $\epsilon(A)$-Module $J$, the abelian sheaf
$\Hhom_{\epsilon(A)}(\epsilon(M),J)$ is injective
(\hyperref[V.prop1.2]{1.2}). By \hyperref[V.prop0.5]{0.5}, $J$ is a
subobject, hence a direct factor, of an object of the form
$\prod_{\alpha}\epsilon(I_{\alpha})$, where the $I_{\alpha}$ are
injective. We may therefore suppose that $J=\epsilon(I)$, where $I$
is an injective object. If the homomorphism
$$
\epsilon(\Hhom_A(M,I))\longrightarrow
\Hhom_{\epsilon(A)}(\epsilon(M),\epsilon(I))
$$
is an isomorphism, it follows from 3) that
$\Hhom_{\epsilon(A)}(\epsilon(M),\epsilon(I))$ is injective. It
therefore remains to prove 5). The fact that formation of the tensor
product commutes with the functor $\epsilon$ follows from IV \hyperref[IV.cor12.6]{12.6} and
the fact that formation of the associated sheaf commutes with
enlargement of the universe (II \hyperref[II.prop3.6]{3.6}). For every object $X$ of
$C^{\sim}_{𝒰}$ and every pair of $A$-Modules $M,N$ of
$C^{\sim}_{𝒰}$, we therefore have
$$
\Hom_{C^{\sim}_{𝒱}}(\epsilon(X),\epsilon(\Hhom_A(M,N)))\simeq
\Hom_{C^{\sim}_{𝒰}}(X,\Hhom_A(M,N))\simeq
\Hom_A(\ZZ_X\otimes_{\ZZ}M,N)
$$
by the full faithfulness of $\epsilon$, and then
$$
\Hom_A(\ZZ_X\otimes_{\ZZ}M,N)\simeq
\Hom_{\epsilon(A)}(\epsilon(\ZZ_X\otimes_{\ZZ}M),\epsilon(N))
\simeq
\Hom_{\epsilon(A)}(\epsilon(\ZZ_X)\otimes_{\ZZ}\epsilon(M),
\epsilon(N))
$$
by the full faithfulness of $\epsilon$ and what precedes. Using 2)
and IV \hyperref[IV.prop12.14]{12.14}, we finally obtain an isomorphism
$$
\Hom_{C^{\sim}_{𝒱}}(\epsilon(X),\epsilon(\Hhom_A(M,N)))\simeq
\Hom_{C^{\sim}_{𝒱}}(\epsilon(X),
\Hhom_{\epsilon(A)}(\epsilon(M),\epsilon(N)));
$$
hence the asserted isomorphism.

\setcounter{subsection}{9}
\subsection{}\etiquette[1.10]{V.subsec1.10}
Let $E$ be a topos and let $\Uu=(U_i\to X)_{i\in I}$ be a small
family of morphisms with the same target. For every finite set
$\Delta_n=\sisi{[0,\ldots,n]}{\{0,\ldots,n\}}$, put
$$
S_n(\Uu)=\coprod_{f:\Delta_n\to I}
U_{f(0)}\times_X\cdots\times_X U_{f(n)}{\quoi},
$$
the coproduct being taken over the set of maps from $\Delta_n$ to
$I$. Let $m,n$ be two integers, and let\pageoriginale
$g:\Delta_m\to\Delta_n$ be a map. We define a morphism
\begin{equation*}
s(g):S_n(\Uu)\to S_m(\Uu){\quoi}\sisi{,}{}\tag{1.10.1}
\end{equation*}\etiquette[1.10.1]{V.1.10.1}
as follows: for every $f:\Delta_n\to I$, the restriction of $s(g)$
to
\mbox{$U_{f(0)}\times_X\cdots\times_X U_{f(n)}$}
is the composite of the canonical inclusion
$$
U_{f(g(0))}\times_X\cdots\times_X U_{f(g(m))}
\hookrightarrow S_m(\Uu){\quoi},
$$
and the unique morphism
$$
s_f(g):U_{f(0)}\times_X\cdots\times_X U_{f(n)}\to
U_{f(g(0))}\times_X\cdots\times_X U_{f(g(m))}
$$
such that, for every $i\in\Delta_m$,
$$
\pr_{f(g(i))}s_f(g)=\pr_{f(g(i))}
$$
($\pr_a$ denotes the $a$-th projection). We thus obtain a
contravariant functor $\Delta_n\mapsto S_n$ from the category of
finite sets to $E$; in other words, a semi-simplicial complex
$S_∙(\Uu)$ of objects of $E$. Note that this complex is canonically
augmented toward $X$. Every functor from $E$ to a category $C$
carries $S_∙(\Uu)$ to a simplicial object of $C$. In particular, if
$A$ is a Ring of $E$, the ``free $A$-Module generated by'' functor
carries $S_∙(\Uu)$ to a simplicial complex of $A$-biModules,
augmented toward $A_X$ and denoted by $A_∙(\Uu)$. We have
\begin{equation*}
A_n(\Uu)=\bigoplus_{f:\Delta_n\to I}
A_{U_{f(0)}\times_X\cdots\times_X U_{f(n)}}{\quoi}.\tag{1.10.2}
\end{equation*}\etiquette[1.10.2]{V.1.10.2}
This complex will often be written
\[
\xymatrix@C=1cm@M=.5cm{
  \cdots  \ar@<12pt>[r]^(.4){s_0} \ar[r]^(.4){s_1} \ar@<-12pt>[r]^(.4){s_2}&
\displaystyle{\mathop{\bigoplus}_{i,j}} \quad A_{U_i ×_X U_j}
 \ar@<6pt>[r]^(.55){s_0}\ar@<-6pt>[r]^(.55){s_1} &
\displaystyle{\mathop{\bigoplus}_{i}}\quad A_{U_i} \ar@<2pt>[r]& A_X \\
}\tag{1.10.3} %% problem (numbering in the wrong place) \XXX
\]\etiquette[1.10.3]{V.1.10.3}
where the arrows in \eqhyperref[V.1.10.3]{1.10.3} (except the last,
which is the augmentation) represent the face operators of the
simplicial complex $A_∙(\Uu)$, that is, the morphisms corresponding
to the increasing injective maps from $\Delta_n$ into
$\Delta_{n+1}$ ($s_i$ omits the integer $i$). To the complex
$A_∙(\Uu)$ we associate a differential complex augmented toward
$A_X$:
\begin{equation}
\cdots\xrightarrow{d}\bigoplus_{i,j}A_{U_i\times_XU_j}
\xrightarrow{d}\bigoplus_iA_{U_i}\to A_X{\quoi},
\sisi{\tag*{1.10.4}}{\tag{1.10.4}}
\end{equation}\etiquette[1.10.4]{V.1.10.4}
by\pageoriginale putting
\begin{equation}
d=\sum(-1)^is_i.\sisi{\tag*}{\tag}{1.10.5}
\end{equation}\etiquette[1.10.5]{V.1.10.5}

\begin{enonce*}{Proposition 1.11}\etiquette[1.11]{V.prop1.11}
When the family $\Uu=(U_i\to X)_{i\in I}$ is epimorphic, the
differential complex \eqhyperref[V.1.10.4]{1.10.4} is acyclic and
gives a resolution of $A_X$.
\end{enonce*}

Denote by $\ZZ$ the constant sheaf with value $\ZZ$. By the
definition of the ``free $A$-Module generated by'' functor, for every
object $H$ of $E$ we have
$$
A_H\simeq {\ZZ}_H\otimes_{\ZZ}A{\quoi},
$$
hence
$$
A_∙(\Uu)\simeq \ZZ_∙(\Uu)\otimes_{\ZZ}A{\quoi}.
$$
Since the components of $\ZZ_∙(\Uu)$ are flat $\ZZ$-Modules
(\hyperref[V.prop1.3.1]{1.3.1}), it suffices to prove the proposition
when $A=\ZZ$.

Suppose first that $E$ is the topos of sets. The augmented complex
$S_∙(\Uu)$ is then a direct sum of augmented complexes of the form
\[
\xymatrix@C=.6cm@M=.18cm@R=.3cm{
\cdots \ar@<9pt>[r] \ar@<3pt>[r] \ar@<-3pt>[r] \ar@<-9pt>[r]&
S\times S\times S \ar@<5pt>[r]\ar[r]\ar@<-5pt>[r]&S\times S
\ar@<2pt>[r]\ar@<-2pt>[r]&S\ar[r]&\{e\}&
\txt{($\{e\}$ a one-element set).}
}
\]
Each of these augmented complexes is homotopically trivial. Thus
$S_∙(\Uu)$ is a homotopically trivial augmented complex, and its
homology is therefore trivial, proving the proposition in this case.

Now let $p:\text{Ens}\to E$ be a point of $E$. Since formation of
the complex $\ZZ_∙(\Uu)$ commutes with the inverse-image functors of
morphisms of topoi,
$p^{\ast}(\ZZ_∙(\Uu))\simeq\ZZ_∙(p^{\ast}(\Uu))$ is a resolution of
${\ZZ}_{p^{\ast}X}\simeq p^{\ast}({\ZZ}_X)$; this proves the
proposition when $E$ has enough fiber functors (IV \hyperref[IV.subsec4.6]{4.6}). This holds,
in particular, when $E$ is a topos of presheaves $C^{\wedge}$,
since for every object $X$ of $C$, $\Gamma(X,-)$ is a fiber functor.
In the general case, $E$ is equivalent to a topos of sheaves on a
small site $C$ (IV 1), and the epimorphic family
$\Uu=(U_i\to X)_{i\in I}$ is the image, under the ``associated
sheaf'' functor, of an epimorphic family
$\Uu'=(U'_i\to X')_{i\in I}$. Consequently,\pageoriginale
${\ZZ}_∙(\Uu)\simeq\underline{a}{\ZZ}_∙(\Uu')$ is a resolution of
$\underline{a}{\ZZ}_{X'}\simeq{\ZZ}_X$.


% END INCLUDED SGA4_English_current_buildable_source_20260906/source/05/drafts/SGA4_V_LINES_603_925_English.texfrag


% BEGIN INCLUDED SGA4_English_current_buildable_source_20260906/source/05/drafts/SGA4_V_LINES_926_1361_English.texfrag
\section{Čech cohomology. Cohomological notation}\etiquette[2]{V.sec2}

\subsection{General notation}\etiquette[2.1]{V.subsec2.1}

\subsubsection{}\etiquette[2.1.1]{V.subsubsec2.1.1}
Let $(E, A)$ be a ringed topos, and let $M$, $N$ be two $A$-Modules
(left modules, to fix ideas). We denote by $\Ext^q_A (E; M, N)$ the
value at $N$ of the $q$-th right derived functor of the functor
$\Hom_A(M,⋅)$ \cite{sga4.exp.V.bib.11}. The functors $\Ext_A^q(E; M, N)$, $q \in
\NN$, form a $\delta$-functor in the variable $N$. They also form a
contravariant $\delta$-functor in the variable $M$. By definition,
$\Ext_A^0(E; M, N) = \Hom_A(M, N)$.

\subsubsection{}\etiquette[2.1.2]{V.subsubsec2.1.2}

Let $X$ be an object of $E$. When $M = A_X$, the free $A$-Module
generated by $X$ (IV 12), we put $\Ext_A^q(E; A_X, N) = H^q(X,
N)$. Notice that the ring $A$ no longer appears in this new
notation. This cannot cause confusion, since we will show
(\hyperref[V.cor3.5]{3.5}) that the formation of $H^q(X,⋅)$
\sisi{commutes}{commutes} with restriction of scalars. The functor
$H^q(X,⋅)$ is the $q$-th right derived functor of
$\Hom_A(A_X,⋅) = \Hom_E(X,⋅)$, also denoted by $\Gamma(X,⋅)$. When
$M = A$, we put $\Ext^q_A(E; A, N) = H^q(E, N)$.

\subsection{Localization}\etiquette[2.2]{V.subsec2.2}

Let $X$ be an object of $E$, and let $j : E_{/X} \to E$ be the
localization morphism (IV 8). The functor $j^\ast_X$ for Modules
\sisi{,}{} is exact (\hyperref[V.prop4.11]{4.11}) and admits an exact
left adjoint $j_{X!}$ (\hyperref[V.prop4.11]{4.11}). Consequently,
it sends injective Modules to injective Modules. Thus, for a
variable $A$-Module $N$ of $E$ and a variable $A | X$-Module $M$ of
$E_{/X}$, there are functorial isomorphisms
\begin{equation*}
\Ext^q_{A | X}(E_{/X} ; M, j^\ast_X N) \simeq \Ext^q_A(E; j_{X!}M,
N){\quoi}.\tag{2.2.1}
\end{equation*}\etiquette[2.2.1]{V.2.2.1}
In particular, there are canonical isomorphisms
$$
H^q(E_{/X}, j^\ast_X N) \simeq H^q (X, N){\quoi}.
$$

For every object $X$ of $E$ and every pair $M$, $N$ of $A$-Modules,
we put
\begin{equation*}
\Ext^q_A(X; M, N) \simeq \Ext^q_{A|X}(E_{/X}; M | X, N|X){\quoi}.\tag{2.2.2}
\end{equation*}\etiquette[2.2.2]{V.2.2.2}
By\pageoriginale what precedes, the functors $\Ext^q_A(X; M, ⋅)$
are the derived functors of the functors $N \mapsto
\Hom_{A|X}(M|X, N|X)$. \sisi{The functors}{The functors}
$(M, N) \mapsto \Ext^q_A(X; M, N)$, $q > 0$, form a
$\delta$-functor with respect to each variable.

\subsection{The case of presheaf topoi. Cohomology of a
  covering}\etiquette[2.3]{V.subsec2.3}

\subsubsection{}\etiquette[2.3.1]{V.subsubsec2.3.1}
Let $C$ be a small category endowed with a presheaf of rings $A$,
let $C^{\wedge}$ be the topos of presheaves on $C$, and let $X$ be
a representable object of $C^{\wedge}$. The functor which associates
to an $A$-Module $M$ the group $\Gamma(X, M) = M(X)$ is exact (I 3).
Consequently, $H^q(X, M)=0$ for every $q > 0$ and every $A$-Module
$M$, or, equivalently, $A_X$ is a projective $A$-Module.

\subsubsection{}\etiquette[2.3.2]{V.subsubsec2.3.2}
Let $S$ be a presheaf on $C$. For every $A$-Module $M$, there is a
canonical isomorphism (I 2):
$$
\Gamma(S, M) \simeq \varprojlim_{C/S} M|S
$$
Moreover, for every injective $A$-Module $M$, the $A | S$-Module
$M | S$ is injective \hyperref[V.subsec2.2]{(2.2)}. Consequently, the groups $H^{q}(S, M)$
are the values at $M | S$ of the right derived functors of
$\varprojlim_{C/S}$. Denoting these derived functors by
$\varprojlim_{C/S}^{q}$, there are canonical isomorphisms
$$
H^q (S, M) \simeq {\varprojlim_{C/S}}^{q} M | S{\quoi}.
$$
In particular, there are canonical isomorphisms
$$
H^q (C^{\wedge}, M) \simeq {\varprojlim_{C}}^{q} M{\quoi}.
$$

\subsubsection{}\etiquette[2.3.3]{V.subsubsec2.3.3}
Let $X$ be an object of $C$, and let $\Uu = (U_i \to X)$, $i \in
I$, be a family of morphisms of $C$ such that, for every $i \in I$,
$U_i \to X$ is pullbackable (I 10). Denote by $A_∙$ the simplicial
complex studied in \hyperref[V.subsec1.10]{1.10}:
\[
\xymatrix@M=.25cm@C=1cm{
A_∙ = \cdots \ar@<7pt>[r] \ar@<1pt>[r]  \ar@<-5pt>[r] &
\displaystyle{\mathop{\coprod}_{(i,j)\in I \times I}}\quad
A_{U_i \times_X U_j} \ar@<3pt>[r]\ar@<-3pt>[r]&
\displaystyle{\mathop{\coprod}_{i\in I}} \quad A_{U_i}.
}
\]
% Why ∐ rather than ⨁ as above?
For every $A$-Module $M$, put $C^{∙}(\Uu, M) = \Hom_A (A_∙, M)$:
\[\vcenter{
\xymatrix@C=.75cm@M=.25cm{
C^{∙}(\mathcal{U},M): \displaystyle{\mathop{\prod}_{i\in I}}
M(U_i) \ar@<6pt>[r]\ar@<-2pt>[r]& \displaystyle{\mathop{\prod}_{(i,j)\in
    I\times I}} M(U_i ×_X U_j) \ar@<10pt>[r] \ar@<3pt>[r]
\ar@<-5pt>[r]& \cdots .
}}\leqno{(2.3.3.1)}
\]\etiquette[2.3.3.1]{V.2.3.3.1}
Put $H^q(\Uu, M) = H^q(C^{∙} (\Uu, M))$.


\begin{enonce*}{Proposition 2.3.4}\etiquette[2.3.4]{V.prop2.3.4}
\mbox{ }
\begin{enumerate}
\renewcommand{\labelenumi}{\theenumi)}
\item With\pageoriginale the notation of
  \hyperref[V.subsubsec2.3.3]{2.3.3}, let $R \hookrightarrow X$ be
  the sieve generated by the family $(U_i \to X)$, $i \in I$. There
  is a canonical isomorphism
$$
H^{q}(\Uu, M) \simeq H^q(R, M){\quoi}.
$$

\item The functors $H^q(\Uu,⋅)$ commute with restrictions of
  scalars.
\end{enumerate}
\end{enonce*}

Since $R$ is a subobject of $X$ in $C^{\wedge}$, the fiber products
$U_{i_1} \times_R U_{i_2} \sisi{\times}{\times_R} \cdots
\times_{R} U_{i_n}$ and $U_{i_1} \times_X U_{i_2}
\sisi{\times}{\times_X} \cdots \times_{X} U_{i_n}$ are canonically
isomorphic. It follows from \hyperref[V.prop1.11]{1.11} that the
complex $A_{𝒰 ∙}$ is a resolution of $A_R$, and from
\hyperref[V.subsubsec2.3.1]{2.3.1} that the components of $A_{𝒰 ∙}$
are projective Modules. The \sisi{cohomology}{cohomology} groups of
$C^{∙}(\mathcal{U}, M)$ are therefore canonically isomorphic to
$\Ext^{q}_A(C^{\wedge}; A_R, M)$, whence the isomorphism. Assertion
2) follows immediately from the description of
$C^{∙}(\mathcal{U}, M)$ \eqhyperref[V.2.3.3.1]{2.3.3.1}.

\begin{enonce*}{Corollary 2.3.5}\etiquette[2.3.5]{V.cor2.3.5}
Let $\mathcal{U} = (U_i \to X, i \in I)$ and
$\mathcal{U}' = (U'_j \to X, j \in J)$ be two families of morphisms
with target $X$; let $\phi = (\phi : I \to J, m_i : U_i \to
U'_{\phi(i)})$ and $\phi' = (\phi' : I \to J, m'_i : U_i \to
U'_{\phi(i)})$ be two morphisms (over $X$) from $\mathcal{U}$ to
$\mathcal{U}'$. The morphisms $\phi$ and $\phi'$ induce the
morphisms $\phi^{q}$ and ${\phi'}^{q} :
\sisi{{\phi'}^{q} :}{} H^q(\mathcal{U}, M) \to H^q(\mathcal{U}',
M)$. The morphisms $\phi^q$ and ${\phi'}^{q}$ are equal. In
particular, if the families $\mathcal{U}$ and $\mathcal{U}'$ are
equivalent (\ie if there is a morphism from $\mathcal{U}$ to
$\mathcal{U}'$ and a morphism from $\mathcal{U}'$ to $\mathcal{U}$),
the $A$-Modules $H^{q} (\mathcal{U}, M)$ and $H^q(\mathcal{U}', M)$
are canonically isomorphic.
\end{enonce*}

\begin{enonce*}[remark]{Proof}
Indeed, let $R$ and $R'$ be the sieves of $X$ generated by the
families $\mathcal{U}$ and $\mathcal{U}'$. The morphisms $\phi$ and
$\phi'$ define the same morphism from $R$ to $R'$, and therefore
induce two homotopic morphisms between the projective resolutions
$A_{\mathcal{U}∙}$ and $A_{\mathcal{U}'∙}$.
\end{enonce*}

\begin{enonce*}[remark]{Exercise 2.3.6}\etiquette[2.3.6]{V.exer2.3.6}
(Standard resolution)
\begin{enumerate}
\renewcommand{\theenumi}{\alph{enumi}}
\renewcommand{\labelenumi}{\theenumi)}
\item Let $C$ be a small category. For every integer $n > 0$, denote
  by $\mathrm{Fl}^{n}(C)$ the set of sequences of morphisms of $C :
  (u_1, \ldots, u_n)$ such that, for every $i$, $0 < i < n$, the
  target of $u_i$ equals the source of $u_{i + 1}$, so that the
  morphisms $u_i$ and $u_{i +1}$ are\pageoriginale composable. Define
  a semisimplicial set
\[
\xymatrix@M=.25cm@C=.75cm{
ES(C) = \ob \;C  & \mathrm{Fl}^{1}(C) \ar@<3pt>[l]\ar@<-3pt>[l]
 & \mathrm{Fl}^{2}(C) \ar@<6pt>[l] \ar[l] \ar@<-6pt>[l] & \mathrm{Fl}^{3}(C)\ar@<9pt>[l]
\ar@<3pt>[l] \ar@<-3pt>[l] \ar@<-9pt>[l] \quad \cdots ,
}
\]
whose face operators $s_i : \mathrm{Fl}^n(C) \to
\mathrm{Fl}^{n-1}(C)$ are as follows:
\begin{align*}
s_{0}(u_1, \ldots, u_n) &= (u_2, \ldots, u_n),\\
s_{i}(u_1, \ldots, u_n) &= (u_1, \ldots, u_{i+1} u_i, \ldots, u_n), \hskip1em 0
< i < n ,\\
s_{n}(u_1, \ldots, u_n) &= (u_1, \ldots, u_{n-1}).
\end{align*}

\item Show that when $C$ has an initial object or a final object,
  the complex $ES(C)$ is homotopically trivial.

\item For every object $X$ of $C$, denote by ${}_{X\backslash}C$ the
  category of arrows with source $X$. Define a semisimplicial
  presheaf $ES(C)$ whose value at every object $X$ of $C$ is
  $ES({}_{X\backslash}C)$. Denote by $\ZZ_{ES(C)}$ the free abelian
  semisimplicial presheaf generated by $ES(C)$. It has a canonical
  augmentation $\ZZ_{ES(C)}\to \ZZ$ to the constant presheaf with
  value $\ZZ$. Show that, upon passing to the associated differential
  complex, the complex $\ZZ_{ES(C)}$ is a resolution of $\ZZ$ and
  that its components are projective abelian presheaves.

\item For every abelian presheaf $M$, put $ST^{∙}(M) =
  \Hom(\ZZ_{ES(C)}, M)$. Describe the components of $ST^{∙}(M)$
  explicitly. Show that, for every integer $q \geq 0$,
  $H^q(ST^{∙}(M)) = H^q(C^{\wedge}, M)$.

\item Show that, for every sheaf $S$ on $C$, the functors $H^q(S,⋅)$
  commute with restrictions of scalars.

\item Observe that when $C$ is a group, $\ZZ_{ES(C)}$ is the
  standard resolution of the trivial module $\ZZ$, \cite{sga4.exp.V.bib.3}.

\item Show that, for every abelian presheaf $M$ on $C$,
  $\ZZ_{ES(C)} \otimes M$ is a (left) resolution of $M$. Define the
  functor $\varinjlim_{C}$ on presheaves. Show that it is right exact.
  Denote its left derived functors by $H_q(C, M)$. Show that the
  components of the complex $\ZZ_{ES(C)} \otimes M$ are acyclic for
  $\varinjlim_{C}$. Deduce, denoting by $ST_{∙}(M)$ the complex
  $\varinjlim_{C}(\ZZ_{ES(C)} \otimes M)$, isomorphisms
  $H_q(ST_{∙}(M)) \simeq H_q(C, M)$.

\item Let\pageoriginale $u : C^{\wedge} \to \Ens$ be the unique
  morphism from the topos $C^{\wedge}$ to the punctual topos. Denote
  by $u_{!} : C^{\wedge}_{\ZZ} \to \Ab$ the left adjoint of the
  functor $u^\ast : \Ab \to C^{\wedge}_{\ZZ}$. Show that for every
  abelian presheaf $M$, $u_{!}M =\varinjlim_{C} M$.

\item Henceforth denote by $ST^{∙}(C, M)$ and $ST_{∙}(C, M)$ the
  complexes denoted above by $ST^{∙}(M)$ and $ST_{∙}(M)$. A
  presheaf $M$ on $C$ is called locally constant if it sends all
  morphisms of $C$ to isomorphisms. Associate to every locally
  constant presheaf $M$ a locally constant presheaf $\check{M}$ on
  the category $C^{\circ}$ (the category opposite to $C$) such that,
  for every morphism $u$ of $C$, $\check{M}(u) = M(u)^{-1}$. Find a
  canonical isomorphism between the complexes $ST^{∙}(C, M)$ and
  $ST^{∙}(C^{\circ}, \check{M})$ (\resp $ST_{∙}(C, M)$ and
  $ST_{∙}(C^{\circ}, \check{M})$).

\item Let $C$ be a small category having an initial object (\resp a
  small filtered category). Show that, for every constant presheaf
  $M$, $H^{q} (C^{\wedge}, M) = 0$ for $q > 0$ (\resp $H_q(C, M) =
  0$ for $q > 0$).
\end{enumerate}
\end{enonce*}

\subsection{The case of small sites. Čech cohomology}\etiquette[2.4]{V.subsec2.4}

\subsubsection{}\etiquette[2.4.1]{V.subsubsec2.4.1}
Let $(C, A)$ be a ringed $𝒰$-site, let $C^{\sim}$ be the topos of
sheaves on $C$, and let $\epsilon: C \to C^{\sim}$ be the canonical
functor which associates to an object of $C$ the sheaf associated to
the presheaf represented by that object. By abuse of notation, for
every object $X$ of $C$ and every sheaf of $A$-modules $M$, we put
$H^q(X, M) = H^q(\epsilon(X), M)$ (\cf \hyperref[V.subsubsec2.3.1]{2.3.1}). Recall that when the
topology of $C$ is coarser than the canonical topology, as is always
the case in practice, the functor $\epsilon$ is fully faithful and
allows us to identify $C$ with its image under $\epsilon$.

\subsubsection{}\etiquette[2.4.2]{V.subsubsec2.4.2}
Denote by $\Hh^0 : C^{\sim{\phantom{a}}}_A \to C^{\wedge}_A$ the
inclusion functor from sheaves of $A$-modules to the category of
presheaves of $A$-modules. Thus, for every sheaf of $A$-modules $M$,
one has by definition
\begin{equation*}
\Hh^{0}(M)(X) = H^{0}(X, M) = M(X){\quoi},\tag{2.4.2.1}
\end{equation*}\etiquette[2.4.2.1]{V.2.4.2.1}
for every object $X$ of $C$. The functor $\Hh^{0}$ is left exact.
Its right derived functors are denoted by $\Hh^{q}$. Since, for
every object $X$ of $C$, the ``sections over $X$'' functor is exact
in the category of presheaves, one has
\begin{equation*}
\mathcal{H}^q(M)(X) = H^q(X, M){\quoi}, \tag{2.4.2.2}
\end{equation*}\etiquette[2.4.2.2]{V.2.4.2.2}
for\pageoriginale every object $X$ of $C$ and every sheaf of
$A$-modules $M$, so that the presheaf $\mathcal{H}^q(M)$ is simply
the presheaf $X \mapsto H^q(X, M)$.

\subsubsection{}\etiquette[2.4.3]{V.subsubsec2.4.3}
Suppose that $(C, A)$ is a ringed \emph{small site}, so that
$C^{\wedge}$ is a \emph{topos} to which the results of
\hyperref[V.subsec2.3]{2.3} may be applied. Let $X$ be an object of
$C$, and let $R \hookrightarrow X$ be a covering sieve. For every
\emph{presheaf} of $A$-modules $G$, the groups $H^{q}(R, G)$ (which
are thus computed in the topos $C^{\wedge}$) are called the
\emph{Čech cohomology groups of the presheaf $G$ relative to the
  covering sieve $R$.} When $R \hookrightarrow X$ is the sieve
generated by a covering family $\mathcal{U} = (U_i \to X)_{i \in I}$
of pullbackable morphisms, these groups can be computed by means of
the complex $C^{∙}(\mathcal{U}, G)$
(\hyperref[V.subsubsec2.3.3]{2.3.3}), called the \emph{Čech complex
  of $G$ relative to the covering family $\mathcal{U}$} (or relative
to the covering $\Uu$). The groups $H^{q}(\Uu,G)=H^{q}(R,G)$
(\hyperref[V.prop2.3.4]{2.3.4}) are then called the \emph{Čech
cohomology groups of $G$ relative to the covering family $\Uu$}.

\subsubsection{}\etiquette[2.4.4]{V.subsubsec2.4.4}
Let $M$ be a sheaf of $A$-modules on $C$. The groups
$H^q(\mathcal{U}, \Hh^{0}(M))$ are most often denoted, by abuse, by
$H^q(\mathcal{U}, M)$ and called the \sisi{Čech}{Čech} cohomology groups of the
sheaf $M$ relative to the covering family $\mathcal{U}$.

\subsubsection{}\etiquette[2.4.5]{V.subsubsec2.4.5}
Denote by $\check{\Hh}^0 : C^{\wedge}_A \to C^{\wedge}_A$ the
natural extension to presheaves of $A$-modules of the functor $L$
described in II. Thus, by definition, for a presheaf $G$ and an
object $X$ of $C$,
\begin{equation}
\check{\Hh}^0 (G) (X) = \varinjlim_{R \hookrightarrow X}
G(R){\quoi},\tag{2.4.5.1}
\end{equation}\etiquette[2.4.5.1]{V.2.4.5.1}
where the inductive limit is taken over the covering sieves of $X$.
It follows from \eqhyperref[V.2.4.5.1]{2.4.5.1} that the functor
$\check{\Hh}^{0}$ is left exact. The right derived functors of
$\check{\Hh}^{0}$ are denoted by $\check{\Hh}^{q}$. Since the
``sections over $X$'' and ``filtered inductive limit'' functors are
exact, it follows from \eqhyperref[V.2.4.5.1]{2.4.5.1} that
\begin{equation}
\check{\Hh}^{q} (G) (X) = \varinjlim_{R \hookrightarrow
  X} H^q(R, G){\quoi},\tag{2.4.5.2}
\end{equation}\etiquette[2.4.5.2]{V.2.4.5.2}
where the limit is taken over the covering sieves of $X$.

The\pageoriginale presheaves $\check{\Hh}^{q}$ are called the
\emph{Čech cohomology presheaves}. For every object $X$ of $C$, put
\begin{equation*}
{\check{H}}^q (X, G) = \check{\Hh}^{q} (G) (X){\quoi}.\tag{2.4.5.3}
\end{equation*}\etiquette[2.4.5.3]{V.2.4.5.3}
The groups ${\check{H}}^q(X, G)$ are called the \emph{Čech
cohomology groups}. When the topology of $C$ is defined by a
pretopology, as is most often the case in practice, it follows from
\hyperref[V.prop2.3.4]{2.3.4} that
\begin{equation*}
{\check{H}}^q (X, G) \simeq \varinjlim_{\mathcal{U}} {H}^{q}
(\mathcal{U}, G){\quoi},\tag{2.4.5.4}
\end{equation*}\etiquette[2.4.5.4]{V.2.4.5.4}
where the inductive limit is taken over the pullbackable covering
families preordered by the natural order relation on their
corresponding sieves (\hyperref[V.cor2.3.5]{2.3.5}).

\subsubsection{}\etiquette[2.4.6]{V.subsubsec2.4.6}
Let $M$ be a sheaf of $A$-modules. By abuse of notation, put
\begin{equation}
{\check{H}}^q (X, M) = {\check{H}}^q (X, \check{\Hh}^{0}(M)),
\qquad \check{\Hh}^{q} (M) = \check{\Hh}^{q}(\check{\Hh}^{0} (M)){\quoi},\tag{2.4.6.1}
\end{equation}\etiquette[2.4.6.1]{V.2.4.6.1}
and the groups ${\check{H}}^q(X, M)$ are called the \sisi{Čech}{Čech} cohomology
groups of the sheaf $M$. Note that, although the functors
${\check{H}}^q$ are derived functors on the category of presheaves,
they do not in general form a $\delta$-functor on the category of
sheaves.

\subsection{Change of universe. Čech cohomology for $𝒰$-sites}\etiquette[2.5]{V.subsec2.5}

\subsubsection{}\etiquette[2.5.1]{V.subsubsec2.5.1}
Let $(C, A)$ be a ringed $𝒰$-site, and let $𝒱$ be a universe
containing $𝒰$. The site $(C, A)$ is then a $𝒱$-site, and there are
a $𝒰$-topos $C^{\sim}_{𝒰}$, a $𝒱$-topos $C^{\sim}_{𝒱}$, and a
canonical inclusion functor $\epsilon : C^{\sim}_{𝒰}
\hookrightarrow C^{\sim}_{𝒱}$. The functor $\epsilon$ is exact and
fully faithful on the categories of Modules, and sends injective
Modules to injective Modules (\hyperref[V.prop1.9]{1.9}). Thus, for
every pair of $𝒰$-sheaves of $A$-modules, there are canonical
isomorphisms
\begin{equation*}
\Ext^q_A(C^{\sim}_{𝒰}; M, N) \simeq
\Ext^q_{\epsilon A}(C^{\sim}_{𝒱}; \epsilon M,
\epsilon N),\quad q \geq 0{\quoi}.\tag{2.5.1.1}
\end{equation*}\etiquette[2.5.1.1]{V.2.5.1.1}
In particular, for every $𝒰$-sheaf of sets $R$ on $C$, there are
canonical isomorphisms (\hyperref[V.subsubsec2.1.1]{2.1.1})
\begin{equation*}
H^q(R, M) \simeq H^q(\epsilon R, \epsilon M),\quad q \geq
0{\quoi},\tag{2.5.1.2}
\end{equation*}\etiquette[2.5.1.2]{V.2.5.1.2}
and, even more particularly, for\pageoriginale every object $X$ of
$C$, there are canonical isomorphisms
(\hyperref[V.subsubsec2.4.1]{2.4.1})
\begin{equation*}
H^q(X, M) \simeq H^q(X, \epsilon M){\quoi}.\tag{2.5.1.3}
\end{equation*}\etiquette[2.5.1.3]{V.2.5.1.3}
Loosely speaking, we may therefore say that sheaf cohomology does
not depend on the choice of universes, and, when required for a
proof \sisi{or}{or} a construction, one may always enlarge the
universe in order to compute the cohomology of a sheaf.

\subsubsection{}\etiquette[2.5.2]{V.subsubsec2.5.2}
Let $(C, A)$ be a ringed $𝒰$-site, and let $𝒱$ be a universe
containing $𝒰$. Denote by $C^{\sim}_{A}$ the category of
$𝒰$-sheaves of modules, and let
$\widehat{\epsilon} : C^{\wedge}_{A, 𝒰} \to
C^{\wedge}_{A, 𝒱}$ be the inclusion functor from $𝒰$-presheaves to
$𝒱$-presheaves of $A$-modules. The functor $\widehat{\epsilon}$ is
exact, and therefore the derived functors of the functor
$\widehat{\epsilon} \Hh^{0} : C^{\sim}_A \to C^{\wedge}_{A,
𝒱} (\phantom{ABCD})$ are the functors $\widehat{\epsilon} \Hh^q$,
$q \geq 0$. By abuse of notation, we will again denote by
$\Hh^q :C^{\sim}_A \to C^{\wedge}_{A, 𝒱}$, $q \geq 0$, the
functors $\widehat{\epsilon} \Hh^q$. When $C$ is $𝒱$-small, this
enlargement of the universe has the following advantage: the
category $C^{\wedge}_{𝒰}$ is not in general a $𝒰$-topos, and
$𝒰$-presheaves of $A$-modules are not necessarily submodules of
injective $𝒰$-presheaves, whereas the category of $𝒱$-presheaves is
a topos (a $𝒱$-topos), and consequently every $𝒱$-presheaf of
$A$-modules is a subobject of an injective $𝒱$-presheaf
(\hyperref[V.subsubsec0.1.1]{0.1.1}). Thus, for every $𝒱$-presheaf
of sets $R$ (and in particular when $R$ is a $𝒰$-presheaf) and every
$𝒰$-sheaf of $A$-modules $M$, the groups $H^q(R, \Hh^q (M))$ are
defined by \hyperref[V.subsubsec2.1.1]{2.1.1}, and it follows from
\eqhyperref[V.2.5.1.2]{2.5.1.2} that these groups do not depend on
the universe $𝒱$ under consideration. Likewise, for every pair of
integers $\geq 0$, denoted by $p$ and $q$, the presheaves
$\check{\Hh}^p (\Hh^q (M))$ are defined by
\hyperref[V.subsubsec2.4.5]{2.4.5}, and it follows from
\eqhyperref[V.2.5.1.2]{2.5.1.2} and
\eqhyperref[V.2.4.5.4]{2.4.5.4} that these presheaves do not depend
on the universe $𝒱$ used to define them.

% END INCLUDED SGA4_English_current_buildable_source_20260906/source/05/drafts/SGA4_V_LINES_926_1361_English.texfrag


% BEGIN INCLUDED SGA4_English_current_buildable_source_20260906/source/05/drafts/SGA4_V_LINES_1362_1552_English.texfrag
\section{The Cartan--Leray spectral sequence associated with a
covering}\etiquette[3]{V.sec3}\pageoriginale

\begin{enonce*}{Proposition 3.1}\etiquette[3.1]{V.prop3.1}
Let $(C,A)$ be a ringed $𝒰$-site, and let $𝒱$ be a universe
containing $𝒰$. The functor
$\Hh^{0}:C^{\sim}_A\to C^{\wedge}_{A,𝒱}$
(\hyperref[V.subsubsec2.5.2]{2.5.2}) sends injective $A$-modules to
injective presheaves. For every integer $q>0$ and every $A$-module
$M$, the sheaf associated with the presheaf $\Hh^q(M)$ is zero.
\end{enonce*}

Denote by $\underline{a}_{𝒱}$ the ``associated sheaf'' functor for
$𝒱$-presheaves, and by
$\epsilon:C^{\sim}_A\hookrightarrow C^{\sim}_{A,𝒱}$ the inclusion
functor from $𝒰$-sheaves to $𝒱$-sheaves. One has
$\underline{a}_{𝒱}\Hh^{0}=\epsilon$ (II \hyperref[II.prop3.6]{3.6}), and, since the functors
$\underline{a}_{𝒱}$ and $\epsilon$ are exact, one has
$\underline{a}_{𝒱}\Hh^q=0$ for every integer $q>0$. For every
$𝒰$-sheaf $M$ and every $𝒱$-presheaf $N$, there is a functorial
isomorphism
$\Hom_{C^{\wedge}_{A,𝒱}}(N,\Hh^{0}(M))\simeq
\Hom_{C^{\sim}_{A,𝒱}}(\underline{a}_{𝒱}N,\epsilon M)$. When $M$ is
injective, $\epsilon M$ is injective (\hyperref[V.prop1.9]{1.9}), and,
since the functor $\underline{a}_{𝒱}$ is exact, the functor
$N\longmapsto\Hom_{C^{\wedge}_{A,𝒱}}(N,\Hh^{0}(M))$ is exact.
Consequently, $\Hh^{0}(M)$ is injective.

\begin{enonce*}{Theorem 3.2}\etiquette[3.2]{V.thm3.2}
Let $(C,A)$ be a ringed $𝒰$-site, let $R$ be a $𝒰$-presheaf of
sets on $C$, and let $M$ be a sheaf of $A$-modules. There is a
spectral sequence, functorial in $R$ and $M$:
\begin{equation*}
H^{p+q}(\underline{a}R,M)\Longleftarrow E^{p,q}_2=
H^p(R,\Hh^q(M)){\quoi}.
\tag{3.2.1}
\end{equation*}\etiquette[3.2.1]{V.3.2.1}
(When $C$ is not $𝒰$-small, the term $H^p(R,\Hh^q(M))$ must be
interpreted as the cohomology of the presheaf $\Hh^q(M)$ in the topos
$C^{\wedge}_{𝒱}$, where $𝒱$ is a universe containing $𝒰$ such
that $C$ is $𝒱$-small (\hyperref[V.subsubsec2.5.2]{2.5.2}).)
\end{enonce*}

By the definition of the ``associated sheaf'' functor (\cf II), there
is an isomorphism of functors
$H^{0}(\underline{a}R,M)\simeq H^0(R,\Hh^{0}(M))$. The functor
$\Hh^{0}$ sends injective objects to injective objects. The spectral
sequence of composite functors (\hyperref[V.prop0.3]{0.3}) is the
desired spectral sequence.

\begin{enonce*}{Corollary 3.3}\etiquette[3.3]{V.cor3.3}
Let $X$ be an object of $C$, and let
$\mathcal{U}=(U_i\to X)$, $i\in I$, be a covering family such that,
for every $i\in I$, the morphism $U_i\to X$ is pullbackable. Then
there is a spectral sequence (called the Cartan--Leray spectral
sequence):
\begin{equation}
H^{p+q}(X,M)\Longleftarrow E^{p,q}_2=
H^p(\mathcal{U},\Hh^q(M)){\quoi}.\tag{3.3.1}
\end{equation}\etiquette[3.3.1]{V.3.3.1}
\end{enonce*}

Let\pageoriginale $R\hookrightarrow X$ be the sieve generated by
$\mathcal{U}$. Since this sieve is covering, the sheaf associated with
$R$ is the sheaf associated with $X$ (II \hyperref[II.def5.2]{5.2}). Taking
\hyperref[V.subsubsec2.4.1]{2.4.1} into account, one has
$H^{p+q}(\underline{a}R,M)=H^{p+q}(X,M)$. The corollary then follows
from \hyperref[V.prop2.3.4]{2.3.4}.

\begin{enonce*}{Corollary 3.4}\etiquette[3.4]{V.cor3.4}
There is a spectral sequence, functorial in the sheaf $M$ and in the
object $X$ of $C$,
\begin{equation}
H^{p+q}(X,M)\Longleftarrow E^{p,q}_2=
\check{H}^p(X,\Hh^q(M)){\quoi}.\tag{3.4.1}
\end{equation}\etiquette[3.4.1]{V.3.4.1}
As $X$ varies in $C$, this spectral sequence gives the spectral
sequence of presheaves
\begin{equation}
\Hh^{p+q}(M)\Longleftarrow E^{p,q}_2=
\check{\Hh}^p(\Hh^q(M)){\quoi}.
\tag{3.4.2}
\end{equation}\etiquette[3.4.2]{V.3.4.2}
These spectral sequences give functorial morphisms (edge morphisms):
\hypertarget{sga4.exp.V.equation.3.4.3}{}\begin{align*}%\etiquette[3.4.3]{V.3.4.3}
\Phi^q(M)&:\check{\Hh}^q(M)\to\Hh^q(M){\quoi},\tag{3.4.3}\\
\Phi^q_X(M)&:\check{H}^q(X,M)\to H^q(X,M){\quoi}.\tag{3.4.4}
\end{align*}\hypertarget{sga4.exp.V.equation.3.4.4}{}
The morphisms $\Phi^q$ and $\Phi^q_X$ are isomorphisms when $q$ is
$0$ or $1$. The morphisms $\Phi^2$ and $\Phi^2_X$ are monomorphisms.
More generally, when the presheaves $\Hh^i(M)$ vanish for
$0<i<n$, the morphisms $\Phi^q(M)$ and $\Phi^q_X(M)$ are isomorphisms
for $0\leq q\leq n$ and monomorphisms for $q=n+1$.
\end{enonce*}

The first spectral sequence is obtained by taking the inductive limit
in the spectral sequence \eqhyperref[V.3.2.1]{3.2.1} over the covering
sieves $R\hookrightarrow X$ of $X$ \eqhyperref[V.2.4.5.2]{2.4.5.2}.
The second spectral sequence follows immediately from it
\eqhyperref[V.2.4.5.3]{2.4.5.3}. The sheaves associated with the
presheaves $\Hh^q(M)$ are zero when $q>0$
(\hyperref[V.prop3.1]{3.1}). Thus
$\check{\Hh}^{0}\check{\Hh}^{0}\check{\Hh}^{q}(M)=0$ for every
$q>0$ (\hyperref[II.theo3.4]{II.3.4}). Consequently,
$\check{\Hh}^{0}\check{\Hh}^{q}(M)=0$ for $q>0$ (II \hyperref[II.prop3.2]{3.2}). The
assertions concerning the morphisms $\Phi^q$ and $\Phi^q_X$ follow
immediately.

\begin{enonce*}{Corollary 3.5}\etiquette[3.5]{V.cor3.5}
Let $(E,A)$ be a ringed topos, and let $M$ be an $A$-module (a left
module, to fix ideas). Denote by $ℳ$ the underlying abelian sheaf of
$M$. The functor $M\longmapstoℳ$ is exact; consequently, for every
object $X$ of $E$, the canonical isomorphism
$$
H^{0}(X,M)\xrightarrow{\sim}H^{0}(X,ℳ)
$$
extends\pageoriginale to morphisms
$$
H^q(X,M)\to H^q(X,ℳ)\qquad q\geq0{\quoi}.
$$
These morphisms are isomorphisms.
\end{enonce*}

For every object $Y$ of $E$, one has \eqhyperref[V.2.4.5.4]{2.4.5.4}:
$$
H^p(Y,M)=\varinjlim_{\mathcal{U}}H^p(\mathcal{U},M)\quav
H^p(Y,ℳ)=\varinjlim_{\mathcal{U}}H^p(\mathcal{U},ℳ),
$$
where the inductive limit is taken over the epimorphic families
$\mathcal{U}=(Y_i\to Y)$, $i\in I$.

Consequently, the canonical homomorphism
$\check{H}^p(Y,M)\to\check{H}^p(Y,ℳ)$ is an isomorphism
(\hyperref[V.prop2.3.4]{2.3.4}). Hence, by
\eqhyperref[V.2.4.5.3]{2.4.5.3}, the canonical homomorphism
$\check{\Hh}^p(M)\to\check{\Hh}^p(ℳ)$ is an isomorphism. If $M$ is
an injective $A$-module, then $\check{\Hh}^p(ℳ)=0$ for $p>0$; hence
$\Hh^1(ℳ)=0$ and, by induction on $p$, $\Hh^p(ℳ)=0$ for $p>0$
(\hyperref[V.cor3.4]{3.4}). Thus $H^p(X,ℳ)=0$ for $p>0$
\eqhyperref[V.2.4.2.2]{2.4.2.2} and \hyperref[V.subsec2.5]{2.5}.
Consequently, the functor $M\longmapstoℳ$ sends injective objects to
objects acyclic for the functor $H^{0}(X,⋅)$, which gives the asserted
isomorphism.

\begin{enonce*}[remark]{Exercise 3.6}\etiquette[3.6]{V.exer3.6}
Let $G$ be a topological group and let $B_G$ be its classifying topos
(IV \hyperref[IV.subsec2.5]{2.5}). Denote by $E_G$ the object of $B_G$ consisting of the
underlying topological space of $G$, endowed with the action of $G$ by
left translations. The canonical morphism from $E_G$ to the final
object $e_G$ of $B_G$ is an epimorphism; hence there is a covering
$\mathcal{U}=(E_G\to e_G)$ and, for every abelian sheaf $F$ of $B_G$,
a spectral sequence
\begin{equation*}
E^{p,q}_2=H^p(\mathcal{U},\Hh^q(F))\Longrightarrow
H^{p+q}(B_G,F){\quoi},\tag{3.6.1}
\end{equation*}\etiquette[3.6.1]{V.3.6.1}
which we propose to study.
\end{enonce*}

\begin{enumerate}
\renewcommand{\labelenumi}{\theenumi)}
\item Show that, for every integer $n$, the topos
$$
B_{G/E_G\times E_G\times\ldots\times E_G}\quad
(n\quad\mbox{factors})
$$
is canonically equivalent to the topos (IV \hyperref[IV.subsec2.5]{2.5})
$$
\TOP(G\times G\times G\times\ldots\times G)\quad
(n-1\quad\mbox{factors}).
$$
For\pageoriginale every integer $n$, denote by $F_n$ the sheaf on the
topological space $G\times\ldots\times G$ ($n$ factors) induced by
$F$.

\item Deduce that the term $E^{p,q}_2$ of
\eqhyperref[V.3.6.1]{3.6.1} is canonically isomorphic to the $p$-th
cohomology group of a complex of the form
$$
H^{q}(e,F_0)\to H^{q}(G,F_1)\to H^{q}(G\times G,F_2)\to
\cdots{\quoi},
$$
which is to be made explicit ($e$ denotes the one-point topological
space).

\item Deduce that the cohomology of the classifying topos $B_G$ with
values in locally constant sheaves is isomorphic to the corresponding
singular cohomology of a classifying space, provided that, on the
spaces of the form $G\times G\times\ldots\times G$, cohomology with
coefficients in locally constant sheaves coincides with the
corresponding singular cohomology (as is the case, for example, when
$G$ is locally contractible). (For classifying spaces, one may consult
\cite{sga4.exp.V.bib.15}.)
\end{enumerate}


% END INCLUDED SGA4_English_current_buildable_source_20260906/source/05/drafts/SGA4_V_LINES_1362_1552_English.texfrag


% BEGIN INCLUDED SGA4_English_current_buildable_source_20260906/source/05/drafts/SGA4_V_LINES_1553_1946_English.texfrag
\section{Acyclic sheaves}\etiquette[4]{V.sec4}

\begin{enonce*}{Definition 4.1}\etiquette[4.1]{V.def4.1}
\emph{Let $(E,A)$ be a ringed topos, let $F$ be an $A$-Module, and
let $S$ be a topologically generating family of $E$. The sheaf $F$
is said to be $S$-acyclic if, for every object $X$ of $S$ and every
integer $q>0$, one has $H^q(X,F)=0$. When $S$ is equal to $\ob E$,
the $S$-acyclic sheaves are called} flasque sheaves.
\end{enonce*}

\setcounter{subsection}{1}
\subsection{}\etiquette[4.2]{V.subsec4.2}
Let $(C,A)$ be a $𝒰$-ringed site, and let $F$ be a sheaf of
$A$-modules on $C$. The sheaf $F$ is said to be $C$-\emph{acyclic}
if it is $S$-acyclic, where $S$ is the family of sheaves associated
with the objects of $C$.

\begin{enonce*}{Proposition 4.3}\etiquette[4.3]{V.prop4.3}
Let $(C,A)$ be a $𝒰$-ringed site, and let $F$ be a sheaf of
$A$-modules. Denote by
$\mathcal{H}^{0}:C^{\sim}_{A}\to C^{\wedge}_{A}$ the inclusion
functor from the $A$-Modules into the category of presheaves of
$A$-modules (here these are $𝒱$-presheaves, where $𝒱$ is a universe
such that $C$ is $𝒱$-small). The following conditions are equivalent:
\begin{enumerate}
\renewcommand{\theenumi}{\roman{enumi}}
\renewcommand{\labelenumi}{\theenumi)}
\item $F$\pageoriginale is $C$-acyclic.

\item For every object $X$ of $C$ and every covering sieve
$R\hookrightarrow X$, one has
$$
H^q(R,\Hh^{0}(F))=0\quad\mbox{for every }q>0{\quoi}.
$$

\item For every object $X$ of $C$, one has
$$
\check{H}^q(X,F)=0\quad\mbox{for every }q>0{\quoi}.
$$
\end{enumerate}
\end{enonce*}

\begin{itemize}
\item[i $\Rightarrow$ ii)]: Since $F$ is $C$-acyclic, the presheaves
$\Hh^q(F)$ vanish for $q>0$. The spectral sequence
\eqhyperref[V.3.2.1]{3.2.1} then gives an isomorphism
$H^q(R,\Hh^{0}(F))\simeq H^q(X,F)$ for every $q$. Consequently,
$H^q(R,\Hh^{0}(F))=0$ for $q>0$.
\item[ii $\Rightarrow$ iii)]: This is clear upon passing to the
inductive limit.
\item[iii $\Rightarrow$ i)]: By induction on $q$, one proves that
the presheaves $\Hh^q(F)$ vanish for $q>0$. For this, use
\eqhyperref[V.3.4.2]{3.4.2}.
\end{itemize}

\setcounter{subsection}{3}
\subsection{}\etiquette[4.4]{V.subsec4.4}
It follows from criterion \hyperref[V.prop4.3]{4.3} ii) and
\hyperref[V.cor3.5]{3.5} that the property of being $S$-acyclic
depends only on the underlying abelian sheaf. In particular, a sheaf
of $A$-modules is flasque if and only if its underlying abelian sheaf
is flasque.

\begin{enonce*}{Corollary 4.5}\etiquette[4.5]{V.cor4.5}
Let $(E,A)$ be a ringed topos, and let $F$ be an $A$-Module. The
following properties are equivalent:
\begin{enumerate}
\renewcommand{\theenumi}{\roman{enumi}}
\renewcommand{\labelenumi}{\theenumi)}
\item $F$ is flasque;

\item For every epimorphic family
$𝒳=(X_i\to X)_{i\in I}$,
$$
H^q(𝒳,F)=0\quad\mbox{for every }q>0{\quoi}.
$$
\end{enumerate}
\end{enonce*}

In \hyperref[V.subsec4.2]{4.2}, take $C$ to be the topos $E$ itself.
The corollary then follows from the equivalence i)
$\Leftrightarrow$ ii) of \hyperref[V.prop4.3]{4.3}.

\setcounter{subsection}{5}
\subsection{}\etiquette[4.6]{V.subsec4.6}
Injective sheaves are flasque. Flasque sheaves are $S$-acyclic for
every topologically generating family $S$. Flasque sheaves are not
necessarily injective (take for $E$ the topos of sets). $S$-acyclic
sheaves are not necessarily flasque (Exercise \hyperref[V.exer4.13]{4.1.3}).

\begin{enonce*}{Proposition 4.7}\etiquette[4.7]{V.prop4.7}
Let $(E,A)$ be a ringed topos, let $F$ be a flasque $A$-Module, and
let $X$ be an object of $E$. For every subobject $Y$ of $X$, the
canonical homomorphism
$H^{0}(X,F)\rightarrow H^{0}(Y,F)$ is\pageoriginale surjective.
\end{enonce*}

Suppose that $Y$ is a subobject of $X$ such that the morphism
$H^{0}(X,F)\to H^{0}(Y,F)$ is not surjective. Let $Z$ be the object
obtained by gluing two copies $X_1$ and $X_2$ of $X$ along $Y$. The
object $Z$ is covered by the two subobjects $X_1$ and $X_2$, and one
has $X_1\times_ZX_2=Y$. Denote the resulting covering by
$𝒳=(X_1,X_2)$. One sees immediately that
$H^1(𝒳,F)\neq0$, a contradiction.

\setcounter{subsection}{7}
\subsection{}\etiquette[4.8]{V.subsec4.8}
For general topoi, the property described in
\hyperref[V.prop4.7]{4.7} does not characterize flasque sheaves
(Exercise \hyperref[V.exer4.15]{4.15}). It does, however,
characterize them for topoi generated by their open subobjects, and
in particular for topoi associated with topological spaces (Exercise
\hyperref[V.exer4.16]{4.16}). In the case of topological spaces, the
terminology ``flasque'' adopted here agrees with that of \cite{sga4.exp.V.bib.7}
(Exercise \hyperref[V.exer4.16]{4.16}).

\begin{enonce*}{Proposition 4.9}\etiquette[4.9]{V.prop4.9}
Let $u:(E,A)\to(E',A')$ be a morphism of ringed topoi.
\begin{enumerate}
\renewcommand{\labelenumi}{\theenumi)}
\item The functor $u_\ast$ carries flasque $A$-Modules to flasque
$A'$-Modules.

\item Let $S$ and $S'$ be topologically generating families of $E$
and $E'$, respectively, such that $u^\ast(S')\subset S$. The functor
$u_\ast$ carries $S$-acyclic $A$-Modules to $S'$-acyclic
$A'$-Modules.

\item When $u$ is a flat morphism (\hyperref[V.def1.8]{1.8}), the
functor $u_\ast$ carries injective $A$-Modules to injective
$A'$-Modules.
\end{enumerate}
\end{enonce*}

Let $X$ be an object of $E'$, let
$𝒳=(X_i\to X)_{i\in I}$ be an epimorphic family, let $F$ be a
flasque $A$-Module, and let $C^{∙}(𝒳,u_\ast F)$ be the Čech complex
of the covering $𝒳$. Using the adjunction between $u_\ast$ and
$u^\ast$, and the fact that $u^\ast$ commutes with fiber products,
one obtains a canonical isomorphism
$$
C^{∙}(𝒳,u_\ast F)\simeq C^{∙}(u^\ast(𝒳),F).
$$
Moreover, $u^\ast$ commutes with inductive limits, and hence
$u^\ast(𝒳)=(u^\ast X_i\to u^\ast X)_{i\in I}$ is an epimorphic
family. Since $F$ is flasque,
$H^q(u^\ast(𝒳),F)=0$ for $q>0$, and consequently
$$
H^q(𝒳,u_\ast F)=H^q(C^{∙}(𝒳,u_\ast F))
=H^q(C^{∙}(u^\ast(𝒳),F))=H^q(u^\ast(𝒳),F)=0
$$
for $q>0$. Thus $u_\ast F$ is flasque.

Let us prove 2). When $S'$ is stable under fiber products, an
argument analogous to the preceding one proves 2). In the general
case, we use the spectral sequence of the morphism $u$
\eqhyperref[V.5.3.2]{5.3.2}. Assertion 2) will not be used
before\pageoriginale \eqhyperref[V.5.3.2]{5.3.2}. Let $F$ be an
$S$-acyclic sheaf. The sheaves $R^qu_\ast F$ are the sheaves
associated with the presheaves $X\mapsto H^q(u^\ast X,F)$. Since
$S'$ is a topologically generating family and $F$ is $S$-acyclic,
one has $R^qu_\ast F=0$ for every $q>0$. The spectral sequence
\eqhyperref[V.5.3.2]{5.3.2} then gives, for every object $X$ of $E'$,
an isomorphism
$H^q(X,u_\ast F)\simeq H^q(u^\ast X,F)$; consequently, for every
$X$ in $S'$, $H^q(X,u_\ast F)=0$.

Let us prove 3). When $u$ is flat, the inverse-image functor
$u^\ast$ for Modules is exact (\hyperref[V.def1.8]{1.8}). Therefore
its right adjoint $u_\ast$ carries injective objects to injective
objects (\hyperref[V.prop0.2]{0.2}).

\begin{enonce*}{Proposition 4.10}\etiquette[4.10]{V.prop4.10}
Let $F$ be an $A$-Module of the ringed topos $(E,A)$, and let $G$ be
an injective $A$-Module.
\begin{enumerate}
\renewcommand{\labelenumi}{\theenumi)}
\item The functor $F\mapsto\Hhom_A(F,G)$ is exact.

\item The abelian sheaf $\Hhom_A(F,G)$ is flasque.
\end{enumerate}
\end{enonce*}

\begin{enonce*}[remark]{Proof}
Let us prove the assertions in turn.
\begin{enumerate}
\renewcommand{\labelenumi}{\theenumi)}
\item Let
$$
0\to F'\to F\to F''\to0
$$
be an exact sequence. We must show that the sequence
$$
0\to\Hhom_A(F'',G)\to\Hhom_A(F,G)\to
\Hhom_A(F',G)\to0
$$
is exact; for this it suffices to show that, for every object $H$ of
$E$, the sequence
\begin{align*}
0\to \Hom_E(H,\Hhom_A(F'',G))\to \Hom_E(H,
\Hhom_A(F,G))\hspace{3cm} \\
\to \Hom_E(H,\Hhom_A(F',G))\to0
\end{align*}
is exact. By (IV \hyperref[IV.cor12.6]{6.12}), the latter sequence is isomorphic to
$$
0\to\Hom_A(A_H\otimes_AF'',G)\to
\Hom_A(A_H\otimes_AF,G)\to
\Hom_A(A_H\otimes_AF',G)\to0,
$$
and the $A$-Module $A_H$ is flat
(\hyperref[V.prop1.3.1]{1.3.1}).

\item Let us show that $\Hhom_A(F,G)$ is flasque. Let
$$
𝒳=(X_i\to X)_{i\in I}{\quoi}
$$
be an epimorphic family of $E$, and let
\[
\xymatrix@M=.25cm@C=1cm{
C_∙(𝒳):\cdots \ar@<5pt>[r] \ar[r] \ar@<-5pt>[r] &
\displaystyle{\mathop{\coprod}_{i,j}}\ZZ_{X_i\times_X X_j}
\ar@<3pt>[r]\ar@<-3pt>[r]&
\displaystyle{\mathop{\coprod}_{i}}\ZZ_{X_i}
}
\]
be the\pageoriginale complex defined in \hyperref[V.prop1.5]{1.5}.
This complex is a flat resolution of the object $\ZZ_X$, which is
itself flat (\hyperref[V.subsec1.4]{1.4} and \hyperref[V.prop1.3.1]{1.3}).

\noindent
We then compute the groups:
$$
H^q(𝒳,\Hhom_A(F,G))\xrightarrow{\sim}
H^q\bigl(\Hom_{\ZZ}^{∙}(C_∙(𝒳),\Hhom_A(F,G))\bigr)
\xrightarrow{\sim}
H^q\bigl(\Hom_A(C_∙(𝒳)\otimes_{\ZZ}F,G)\bigr){\quoi}.
$$
The complex
$$
C_∙(𝒳)\otimes_{\ZZ}F
$$
is acyclic in degrees different from $0$, because $C_∙(𝒳)$ is a
flat resolution of a flat module. Consequently,
$\Hom_A(C_∙(𝒳)\otimes_{\ZZ}F,G)$ is acyclic in degrees different
from $0$.
\end{enumerate}
\end{enonce*}

\begin{enonce*}{Proposition 4.11}\etiquette[4.11]{V.prop4.11}
Let $(E,A)$ be a ringed topos, and let $F$ be a flasque (\resp
injective) sheaf of $A$-modules.
\begin{enumerate}
\renewcommand{\labelenumi}{\theenumi)}
\item For every object $X$ of $E$, the $A|X$-Module $j_X^\ast F$ is
flasque (\resp injective).

\item For every closed subtopos $Z$ of $E$, the sheaf of sections of
$F$ with support in $Z$ \emph{(IV 14)} is flasque (\resp injective).
\end{enumerate}
\end{enonce*}

Assertion 1) follows from \hyperref[V.cor4.5]{4.5} when $F$ is
flasque. The functor $j_X^\ast$ admits an exact left adjoint
$j_{X!}$ (IV 11). Consequently, when $F$ is injective, $j_X^\ast F$
is injective (\hyperref[V.prop0.2]{0.2}). Let us prove 2). Denote the
inclusion morphism by $i:Z\to E$. The sheaf of sections of $F$ with
support in $Z$ is the sheaf $i_\ast i^{!}F$ (IV 14). The functor
$i_\ast i^{!}$ is right adjoint to the exact functor
$i_\ast i^\ast$ (IV 14). It therefore carries injective sheaves to
injective sheaves. Let $U$ be the open complement of $Z$, let
$j:U\to E$ be the inclusion morphism, and let $F$ be a flasque
sheaf. There is an exact sequence (IV 14):
\begin{equation*}
0\to i_\ast i^{!}F\to F\to j_\ast j^\ast F{\quoi}.\tag{4.11.1}
\end{equation*}\etiquette[4.11.1]{V.4.11.1}
For every object $X$ of $E$, one has
$j_\ast j^\ast F(X)=F(X\times U)$, and the morphism
$F(X)\to j_\ast j^\ast F(X)$ induced by the last arrow of
\eqhyperref[V.4.11.1]{4.11.1} comes from the canonical injection
$X\times U\hookrightarrow X$. Since $F$ is flasque, this morphism is
surjective (\hyperref[V.prop4.7]{4.7}), and hence the last arrow of
\eqhyperref[V.4.11.1]{4.11.1} is an epimorphism of presheaves. For
every object $X$ of $E$, the exact cohomology sequence deduced from
\eqhyperref[V.4.11.1]{4.11.1} gives
$H^q(X,i_\ast i^{!}F)=0$ for $q>0$, and hence $i_\ast i^{!}F$ is
flasque.

\setcounter{subsection}{11}
\subsection{}\etiquette[4.12]{V.subsec4.12}
The\pageoriginale property of a sheaf being flasque (\resp
injective) localizes \eqhyperref[V.4.11.1]{4.11.1}. It is not, in
general, a property of local nature (Exercise
\hyperref[V.exer4.15]{4.15}). It is, however, a property of local
nature for topoi generated by their open subobjects, and in
particular for topoi associated with topological spaces (Exercise
\hyperref[V.exer4.16]{4.16}).

\begin{enonce*}[remark]{Exercise 4.13}\etiquette[4.13]{V.exer4.13}
Let $X$ be a locally compact space, and let $F$ be a $c$-soft sheaf
\cite{sga4.exp.V.bib.7}. Using the local nature of softness (\loccit), show that
the restriction of $F$ to every paracompact open subset of $X$ is a
soft sheaf. Show that the paracompact open subsets form a basis for
the topology of $X$. Deduce that, if $S$ denotes the family of
paracompact open subsets of $X$, then $F$ is $S$-acyclic. Show that
the sheaf of continuous functions on $X$ is $S$-acyclic but is not
flasque when $X$ is not discrete.
\end{enonce*}

\begin{enonce*}[remark]{Problem 4.14}\etiquette[4.14]{V.prob4.14}
Study totally acyclic topoi, \ie topoi such that
$H^q(X,F)=0$ for every $q>0$, every object $X$, and every abelian
sheaf $F$.
\end{enonce*}

\begin{enonce*}[remark]{Exercise 4.15}\etiquette[4.15]{V.exer4.15}
Let $G$ be a discrete group, and let $B_G$ be its classifying topos
(IV \hyperref[IV.subsec2.4]{2.4}). Show that for every abelian sheaf $F$ and every
monomorphism $X\hookrightarrow Y$, the homomorphism $F(Y)\to F(X)$
is surjective. Let $E(G)$ be the group $G$ regarded as a homogeneous
space under itself. It is an object of $B_G$. The topos
$B_{G/E(G)}$ is equivalent to the punctual topos (IV 8). The
morphism $E(G)\to e$ ($e$ the final object of $B_G$) is an
epimorphism. For every abelian sheaf $F$ on $B_G$, the sheaf
$F|E(G)$ is flasque. Deduce that the property of being flasque or
injective is not of local nature.
\end{enonce*}

\begin{enonce*}[remark]{Exercise 4.16}\etiquette[4.16]{V.exer4.16}
A topos $E$ is said to be generated by its open subobjects if the
open subobjects of $E$ (\ie the subobjects of the final object $e$ of
$E$) form a generating family (I 7). Such a topos has the following
property:

\noindent
(P) Every epimorphic family $X_i\to X$, $i\in I$, admits a
refinement by an epimorphic family $U_j\to X$, $j\in J$, in which
the $U_j\to X$ are monomorphisms.

\begin{enumerate}
\renewcommand{\theenumi}{\alph{enumi}}
\renewcommand{\labelenumi}{\theenumi)}
\item Are there\pageoriginale topoi which have property (P) but are
not generated by their open subobjects? Topoi associated with
topological spaces are generated by their open subobjects. If $E$
is generated by its open subobjects (\resp has (P)), then for every
object $X$ of $E$, $E_{/X}$ is generated by its open subobjects
(\resp has (P)). Every subtopos of a topos generated by its open
subobjects is generated by its open subobjects. Property (P) is not
a property of local nature.

\item Let $E$ be a topos, and let $\Ouv(E)$ be the category of open
subobjects of $E$, endowed with the induced topology. The inclusion
functor $\Ouv(E)\to E$ is a morphism of sites, and hence gives a
morphism of topoi $\Pi:E\to\Ouv(E)^{\sim}$. The functor $\Pi^{\ast}$
is fully faithful. With respect to the 2-category of topoi generated
by their open subobjects, the morphism $\Pi$ has a universal property
which the reader is invited to make explicit.

\item Let $E$ be a topos generated by its open subobjects, and let
$\Point(E)$ be the set of points of $E$ up to isomorphism. Show that
$\Point(E)$ is small. Put a topology on $\Point(E)$ and define a
morphism $\Top(\Point(E))\to\Ouv(E)^{\sim}$ making
$\Top(\Point(E))$ a subtopos of $\Ouv(E)^{\sim}$. For every
topological space $X$, show that every morphism from $\Top(X)$ to
$E$ gives a morphism from $\Top(X)$ to $\Top(\Point(E))$. Show that
a topos generated by its open subobjects is equivalent to a topos
$\Top(X)$ ($X$ a topological space) if and only if it has enough
points. Show that there exist topoi generated by their open
subobjects which are not equivalent to any topos $\Top(X)$, where
$X$ is a topological space.

\item Let $E$ be a topos having property (P). An abelian sheaf $F$
on $E$ is flasque if and only if, for every monomorphism
$X\hookrightarrow Y$, the homomorphism $F(Y)\to F(X)$ is
surjective.

\item Let $E$ be a topos having property (P). Show that the property
of a sheaf $F$ being flasque is a property of local nature.

\item Let $E$ be a topos generated by its open subobjects, and let
$e$ be the final object of $E$. An abelian sheaf $F$ is flasque if
and only if, for every open subobject $U$ of $e$, the canonical
morphism $F(e)\to F(U)$ is surjective.
\end{enumerate}
\end{enonce*}

\begin{enonce*}[remark]{Exercise 4.17}\etiquette[4.17]{V.exer4.17}
(Flasque\pageoriginale sheaves and change of universe)

Let $C$ be a $𝒰$-site (for example, a $𝒰$-topos), and let $𝒱$ be a
universe containing $𝒰$. Denote the corresponding categories of
sheaves by $C^{\sim}_{𝒰}$ and $C^{\sim}_{𝒱}$, and denote the
canonical inclusion by
$\epsilon:C^{\sim}_{𝒰}\to C^{\sim}_{𝒱}$. Let $F$ be a flasque
abelian $𝒰$-sheaf on $C$. We shall show that $\epsilon F$ is
flasque. First note that, for every object $X$ of
$C^{\sim}_{𝒰}$, one has $H^q(\epsilon X,\epsilon F)=0$ for $q>0$.
Every object $Y$ of $C^{\sim}_{𝒱}$ admits an epimorphic family
$\epsilon X_i\to Y$, $i\in I$, in which the $\epsilon X_i\to Y$
are monomorphisms. It follows that
$$
H^q(Y,\epsilon F)=displaystyle
\varprojlim_{\epsilon X\hookrightarrow Y}^{q}F(X).
$$
To show that these $\varprojlim^q$ vanish for $q\neq0$, one may
reduce to the topos $\Ouv(C_{𝒱/Y})^{\sim}$ and use the fact that, for
this topos, a sheaf is flasque if it is so locally.
\end{enonce*}


% END INCLUDED SGA4_English_current_buildable_source_20260906/source/05/drafts/SGA4_V_LINES_1553_1946_English.texfrag


% BEGIN INCLUDED SGA4_English_current_buildable_source_20260906/source/05/drafts/SGA4_V_LINES_1947_2047_English.texfrag
\section{The $R^qu_*$ and the Cartan--Leray spectral sequence for a morphism of topoi}\etiquette[5]{V.sec5}

\setcounter{subsection}{-1}
\subsection{}\etiquette[5.0]{V.subsec5.0}
Let $u:(E,A)\to(E',A')$ be a morphism of ringed topoi, and let
$u_\ast:E\to E'$ be the direct-image functor. The notation $u_\ast$
will also denote the extension of the direct-image functor to
modules. The functor $u_\ast$ on modules (left modules, to fix ideas)
is left exact. Its right derived functors are denoted by
$R^qu_\ast$, $q\geq 0$.

\begin{enonce*}{Proposition 5.1}\etiquette[5.1]{V.prop5.1}
\begin{enumerate}
\renewcommand{\labelenumi}{\theenumi)}
\item For every $A$-module $M$, the sheaf $R^qu_\ast M$ is the
sheaf associated with the presheaf $X'\mapsto H^q(u^\ast X',M)$
$(X'\in\ob E')$.

\item Formation of the functors $R^qu_\ast$ commutes with
restriction of scalars.

\item Formation of the functors $R^qu_\ast$ commutes with
localization. More precisely, for every object $X'$ of $E'$, if
$u_{/X'}:E_{/u^\ast X'}\to E'_{/X'}$ denotes the morphism deduced
from $u$ by localization \emph{(IV 8)}, then, for every $A$-module
$M$, there is a canonical isomorphism
\begin{equation*}
R^q(u_{/X'})_\ast(M|u^\ast X')\simeq R^qu_\ast(M)|X'\qquad q\geq
0{\quoi}.\tag{5.1.1}
\end{equation*}\etiquette[5.1.1]{V.5.1.1}
\end{enumerate}
\end{enonce*}

Let $u^{\wedge}_\ast:E^{\wedge}\to E^{\prime\wedge}$ denote the
direct-image functor for $𝒰$-presheaves
$(u^{\wedge}_\ast M=M\circ u^\ast)$. Since $u^\ast$ and $u_\ast$
are adjoint, there is an isomorphism
$$
u_\ast\simeq\underline{a}u^{\wedge}_\ast{}^{0}{\quoi},
$$
where\pageoriginale $\underline{a}$ is the associated-sheaf functor
for $E'$. Since the functors $\underline{a}$ and $u^{\wedge}_\ast$
are exact, we have
$$
R^qu_\ast\simeq\underline{a}u^{\wedge}_\ast\Hh^q{\quoi},
$$
which is another way of stating 1). Assertion 2) then follows from
1) and \hyperref[V.cor3.5]{3.5}. By definition of the morphism
$u_{/X'}$, there is a canonical isomorphism
\eqhyperref[V.5.1.1]{5.1.1} for $q=0$. The general case follows by
noting that the localization functors (IV 8) are exact and take
injective objects to injective objects (\hyperref[V.prop4.11]{4.11}).

\begin{enonce*}{Proposition 5.2}\etiquette[5.2]{V.prop5.2}
Let $u:E\to E'$ be a morphism of topoi and let $S'$ be a generating
family of $E'$. The sheaves $M$ which are acyclic for the functors
$H^0(u^\ast X',\mathord\cdot)$, $X'\in S'$, are acyclic for the
functor $u_\ast$. In particular, flabby sheaves are acyclic for
$u_\ast$.
\end{enonce*}

This follows from \hyperref[V.prop5.1]{5.1}, 1).

\begin{enonce*}{Proposition 5.3}\etiquette[5.3]{V.prop5.3}
Let $u:E\to E'$ be a morphism of topoi and let $M$ be an abelian
group of $E$. There is a spectral sequence
\begin{equation*}
E^{p,q}_2=H^p(E',R^qu_\ast M)\Longrightarrow H^{p+q}(E,M){\quoi}.
\tag{5.3.1}
\end{equation*}\etiquette[5.3.1]{V.5.3.1}
More generally, for every object $X'$ of $E'$, there is a spectral
sequence
\begin{equation}
E^{p,q}_2=H^p(X',R^qu_\ast M)\Rightarrow H^{p+q}(u^\ast X',M){\quoi}.
\tag{5.3.2}
\end{equation}\etiquette[5.3.2]{V.5.3.2}
\end{enonce*}

By definition of the direct- and inverse-image functors, there is an
isomorphism $H^0(X',u_\ast M)\simeq H^0(u^\ast X',M)$. The functor
$u_\ast$ takes injective objects to flabby sheaves
(\hyperref[V.subsec4.6]{4.6} and \hyperref[V.prop4.9]{4.9}). The
asserted spectral sequences are therefore spectral sequences of
composite functors (\hyperref[V.prop0.3]{0.3}).

\begin{enonce*}{Proposition 5.4}\etiquette[5.4]{V.prop5.4}
Let $u:E\to E'$ and $v:E'\to E''$ be two morphisms of topoi. There
is a spectral sequence
$$
R^pv_\ast R^qu_\ast\Longrightarrow R^{p+q}(v\circ u)_\ast{\quoi}.
$$
\end{enonce*}

We have $v_\ast u_\ast\simeq(vu)_\ast$, and the functor $u_\ast$
takes injective objects to\pageoriginale flabby sheaves
(\hyperref[V.prop4.9]{4.9}), hence to objects acyclic for $v_\ast$
(\hyperref[V.prop5.2]{5.2}). Thus there is a spectral sequence of
composite functors (\hyperref[V.prop0.3]{0.3}).


% END INCLUDED SGA4_English_current_buildable_source_20260906/source/05/drafts/SGA4_V_LINES_1947_2047_English.texfrag


% BEGIN INCLUDED SGA4_English_current_buildable_source_20260906/source/05/drafts/SGA4_V_LINES_2048_2300_English.texfrag
\section{Local Ext and cohomology with supports}\etiquette[6]{V.sec6}

\setcounter{subsection}{-1}
\subsection{}\etiquette[6.0]{V.subsec6.0}

Let $(E,A)$ be a ringed topos, and let $M$ be an $A$-module, a left
module to fix ideas. The functor
$N\longmapsto\Hhom_A(M,N)$, from the category of left $A$-modules to
the category of abelian groups, is left exact (IV 12). Its right
derived functors are denoted by
\begin{equation*}
\Eext_A^q(M,N){\quoi}.\tag{6.0.1}
\end{equation*}\etiquette[6.0.1]{V.6.0.1}
In particular, one has
\begin{equation*}
\Eext_A^{0}(M,N)=\Hhom_A(M,N){\quoi}.\tag{6.0.2}
\end{equation*}\etiquette[6.0.2]{V.6.0.2}
By definition, there are canonical isomorphisms
(\hyperref[V.subsec2.1]{2.1} and IV 12)
\begin{equation*}
H^{0}(E,\Eext_A^{0}(M,N))=\Ext^{0}_A(E;M,N)
=\Hom_A(M,N){\quoi}.\tag{6.0.3}
\end{equation*}\etiquette[6.0.3]{V.6.0.3}
More generally, for every object $X$ of $E$, one has
(\hyperref[V.subsec2.2]{2.2})
\begin{equation*}
H^{0}(X,\Eext_A^{0}(M,N))=\Ext^{0}_A(X;M,N)
=\Hom_{A|X}(M|X,N|X){\quoi}.\tag{6.0.4}
\end{equation*}\etiquette[6.0.4]{V.6.0.4}

\begin{enonce*}{Proposition 6.1}\etiquette[6.1]{V.prop6.1}
\begin{enumerate}
\renewcommand{\labelenumi}{\theenumi)}
\item Formation of the functors $\Eext^q_A$ commutes with
localization. More precisely, for every object $X$ of $E$, there are
functorial isomorphisms
\begin{equation*}
\Eext^q_A(M,N)|X\simeq
\Eext^q_{A|X}(M|X,N|X){\quoi}.\tag{6.1.1}
\end{equation*}\etiquette[6.1.1]{V.6.1.1}

\item The sheaf $\Eext^q_A(M,N)$ is isomorphic to the sheaf associated
with the presheaf $X\longmapsto\Ext^q_A(X;M,N)$.

\item There is a spectral sequence
\begin{equation*}
\Ext^{p+q}_A(E;M,N)\Longleftarrow E^{p,q}_2
=H^p(E,\Eext^q_A(M,N)){\quoi}.\tag{6.1.2}
\end{equation*}\etiquette[6.1.2]{V.6.1.2}
\end{enumerate}
More\pageoriginale generally, for every object $X$ of $E$, there is a
spectral sequence functorial in $X$ and in the arguments $M$ and $N$:
\begin{equation*}
\Ext^{p+q}_A(X;M,N)\Longleftarrow E^{p,q}_2
=H^p(X,\Eext^q_A(M,N)){\quoi}.\tag{6.1.3}
\end{equation*}\etiquette[6.1.3]{V.6.1.3}
\end{enonce*}

By definition, \eqhyperref[V.6.1.1]{6.1.1} is an isomorphism when
$q=0$ (IV 12). The general case follows, taking into account the fact
that the localization functors are exact and send injective modules to
injective modules (\hyperref[V.subsec2.2]{2.2}). The functor
$N\longmapsto\Hhom_A(M,N)=\Eext^0_A(M,N)$ sends injective modules to
flasque sheaves, and hence to sheaves acyclic for $H^{0}(X,⋅)$
(\hyperref[V.prop4.10]{4.10}). The spectral sequences
\eqhyperref[V.6.1.2]{6.1.2} and \eqhyperref[V.6.1.3]{6.1.3} are the
spectral sequences of composite functors, in view of
\eqhyperref[V.6.0.3]{6.0.3} and \eqhyperref[V.6.0.4]{6.0.4}. As $X$
varies in $E$, the spectral sequence \eqhyperref[V.6.1.3]{6.1.3}
gives a spectral sequence of presheaves and hence, upon passing to the
associated sheaves, a spectral sequence of sheaves. Since the sheaves
associated with the presheaves $X\longmapsto H^p(X,⋅)$ are zero when
$p\neq0$ (\hyperref[V.prop3.1]{3.1}), this spectral sequence
degenerates and gives the isomorphism asserted in 2).

\begin{enonce*}{Proposition 6.2}\etiquette[6.2]{V.prop6.2}
The functors $(M,N)\longmapsto\Eext^q_A(M,N)$, $q\geq0$, form a
$\delta$-functor in the variable $M$ and in the variable $N$.
Explicitly, for every exact sequence
$0\to N'\to N\to N''\to0$
(\resp $0\to M'\to M\to M''\to0$), there are long exact sequences
\begin{equation*}
\cdots\to\Eext^q_A(M,N')\to\Eext^q_A(M,N)\to\Eext^q_A(M,N'')
\xrightarrow{\delta}\Eext^{q+1}_A(M,N')\to\cdots\tag{6.2.1}
\end{equation*}\etiquette[6.2.1]{V.6.2.1}
(\resp
\begin{equation*}
\cdots\to\Eext^q_A(M'',N)\to\Eext^q_A(M,N)\to\Eext^q_A(M',N)
\xrightarrow{\delta}\Eext^{q+1}_A(M'',N)\to\cdots){\quoi}.
\tag{6.2.2}
\end{equation*}\etiquette[6.2.2]{V.6.2.2}
\end{enonce*}

It follows from the general properties of the functors $\Ext^q$ that,
for every object $X$ of $E$, the functors
$(M,N)\longmapsto\Ext^q_A(X;M,N)$, $q\geq0$, form a
$\delta$-functor in each argument. The assertion follows by allowing
$X$ to vary and taking the associated sheaf
(\hyperref[V.prop6.1]{6.1}).

\setcounter{subsection}{2}
\subsection{}\etiquette[6.3]{V.subsec6.3}
Let $(E,A)$ be a ringed topos, let $M$ be an $A$-module, let $Z$ be a
closed subtopos of $E$ (IV 9), and let $U$ be its open complement.
Denote by $H^{0}_Z(E,M)$ the group of sections of $M$ whose
support\pageoriginale is contained in $Z$ (IV 14), and by
$ℋ^{0}_Z(M)$ the subsheaf of $M$ defined by the sections of $M$
``with supports in $Z$'' (IV 14). The functors $H^{0}_Z(E,⋅)$ and
$ℋ^{0}_Z(⋅)$ are left exact (IV 14). Their derived functors are denoted
by $H^q_Z(E,⋅)$ and $ℋ^q_Z(⋅)$ respectively, and are called the
\emph{cohomology groups (\emph{\resp} sheaves) of $M$ with supports in
$Z$}\footnote{Compare with \SGA 2 I for the case of ordinary
topological spaces, as well as with the exposé of \textsc{Hartshorne}
cited on p.~80 below.}.

There are canonical isomorphisms (IV 14)
\hypertarget{sga4.exp.V.equation.6.3.1}{}\begin{align*}
H^{0}_Z(E,M)&\simeq\Hom_A(A_Z,M)\simeq
\Ext^{0}_A(E;A_Z,M){\quoi},\tag{6.3.1}\\
ℋ^{0}_Z(M)&\simeq\Hhom_A(A_Z,M)\simeq
\Eext^{0}_A(A_Z,M){\quoi},\tag{6.3.2}
\end{align*}\hypertarget{sga4.exp.V.equation.6.3.2}{}
and hence, for every $q\geq0$, there are isomorphisms
\hypertarget{sga4.exp.V.equation.6.3.3}{}\begin{align*}%\etiquette[6.3.3]{V.6.3.3}
H^q_Z(E,M)&\simeq\Ext^q_A(E;A_Z,M){\quoi},\tag{6.3.3}\\
ℋ^q_Z(M)&\simeq\Eext^q_A(A_Z,M){\quoi}.\tag{6.3.4}
\end{align*}\hypertarget{sga4.exp.V.equation.6.3.4}{}
Note that, since $A_Z$ is a bimodule, the sheaves
$\Eext^q_A(A_Z,M)\simeqℋ^q_Z(M)$ are canonically endowed with
$A$-module structures.

For every object $X$ of $E$, denote by $Z_{/X}$ the closed subtopos of
$E_{/X}$ obtained from $Z$ by localization (it is the complement of
$U\times X$). By definition, there are canonical isomorphisms
\begin{equation*}
H^{0}(X,ℋ^{0}_Z(M))\simeq
H^{0}_{Z_{/X}}(E_{/X},M|X){\quoi}.\tag{6.3.5}
\end{equation*}\etiquette[6.3.5]{V.6.3.5}
Set
\begin{equation*}
H^q_Z(X,M)\simeq H^q_{Z_{/X}}(E_{/X},M|X){\quoi}.\tag{6.3.6}
\end{equation*}\etiquette[6.3.6]{V.6.3.6}
Taking \hyperlink{sga4.exp.V.equation.6.3.2}{(6.3.2)} into account, there are canonical isomorphisms
\begin{equation*}
H^q_Z(X,M)\simeq\Ext^q_A(X;A_Z,M){\quoi}.\tag{6.3.7}
\end{equation*}\etiquette[6.3.7]{V.6.3.7}

\begin{enonce*}{Proposition 6.4}\etiquette[6.4]{V.prop6.4}
\begin{enumerate}
\renewcommand{\labelenumi}{\theenumi)}
\item Formation\pageoriginale of the functors $ℋ^q_Z$ commutes with
localization. More precisely, for every object $X$ of $E$ and every
$A$-module $M$, there are canonical isomorphisms
\begin{equation*}
ℋ^{q}_Z(M)|X\simeqℋ^q_{Z_{/X}}(M|X){\quoi}.\tag{6.4.1}
\end{equation*}\etiquette[6.4.1]{V.6.4.1}

\item The sheaf $ℋ^q_Z(M)$ is the sheaf associated with the presheaf
$X\longmapsto H^q_Z(X,M)$.

\item There is a spectral sequence
\begin{equation*}
H^{p+q}_Z(E,M)\Longleftarrow H^p(E,ℋ^q_Z(M)){\quoi}.\tag{6.4.2}
\end{equation*}\etiquette[6.4.2]{V.6.4.2}
More generally, for every object $X$ of $E$, there is a spectral
sequence
\begin{equation*}
H^{p+q}_Z(X,M)\Longleftarrow H^p(X,ℋ^q_Z(M)){\quoi}.\tag{6.4.3}
\end{equation*}\etiquette[6.4.3]{V.6.4.3}
\end{enumerate}
\end{enonce*}

This is simply Proposition \hyperref[V.prop6.1]{6.1} translated by
means of the dictionary in \hyperref[V.subsec6.3]{6.3}.

\begin{enonce*}{Proposition 6.5}\etiquette[6.5]{V.prop6.5}
With the notation of \hyperref[V.subsec6.3]{6.3}, denote by
$j:U\to E$ the canonical morphism. For every $A$-module $M$, there is
an exact sequence of sheaves
\begin{equation*}
0\toℋ^{0}_Z(M)\to M\to j_\ast(M|U)\to
ℋ^1_Z(M)\to0{\quoi},\tag{6.5.1}
\end{equation*}\etiquette[6.5.1]{V.6.5.1}
and, for $q\geq2$, there are isomorphisms
\begin{equation*}
ℋ^q_Z(M)\simeq\Eext^{q-1}_A(A_U,M)\simeq
R^{q-1}j_\ast(M|U){\quoi}.\tag{6.5.2}
\end{equation*}\etiquette[6.5.2]{V.6.5.2}
There is, in addition, a long exact sequence
\begin{equation*}
\cdots\to H^q_Z(E,M)\to H^q(E,M)\to H^q(U,M)\to
H^{q+1}_Z(E,M)\to\cdots
\footnote{This exact sequence makes precise the role of the global
cohomological invariants $H^q_Z(E,M)$ as ``cohomology groups of $E$
\emph{modulo} the open subtopos $U$, with coefficients in $M$.''}
\tag{6.5.3}
\end{equation*}\etiquette[6.5.3]{V.6.5.3}
and, more generally, for every object $X$ of $E$, there is a long
exact sequence
\begin{equation*}
\cdots\to H^q_Z(X,M)\to H^q(X,M)\to H^q(X\times U,M)\to
H^{q+1}_Z(X,M)\to\cdots.\tag{6.5.4}
\end{equation*}\etiquette[6.5.4]{V.6.5.4}
\end{enonce*}

By the definition of $A_Z$, there is an exact sequence (IV 14)
\begin{equation*}
0\to A_U\to A\to A_Z\to0{\quoi}.\tag{6.5.5}
\end{equation*}\etiquette[6.5.5]{V.6.5.5}
The\pageoriginale functors $\Eext^q_A(A,⋅)$ vanish for $q>0$. One has
$\Hhom_A(A_U,M)\simeq j_\ast(M|U)$ (IV 14), and hence
(\hyperref[V.subsec2.2]{2.2})
$\Eext^q_A(A_U,M)\simeq R^qj_\ast(M|U)$. Finally,
$\Eext^q_A(A_Z,M)\simeqℋ^q_Z(M)$ \hyperlink{sga4.exp.V.equation.6.3.4}{(6.3.4)}. The long exact sequence
\eqhyperref[V.6.2.2]{6.2.2} gives, in this case,
\eqhyperref[V.6.5.1]{6.5.1} and \eqhyperref[V.6.5.2]{6.5.2}. The exact
sequences \eqhyperref[V.6.5.3]{6.5.3} and
\eqhyperref[V.6.5.4]{6.5.4} follow from the dictionary in
\hyperref[V.subsec6.3]{6.3} and from the long exact sequence of the
$\delta$-functor $\Ext^q_A(X;\mathord\cdot,M)$, $q\geq0$, associated
with \eqhyperref[V.6.5.5]{6.5.5}.

\begin{enonce*}{Proposition 6.6}\etiquette[6.6]{V.prop6.6}
\begin{enumerate}
\renewcommand{\labelenumi}{\theenumi)}
\item Flasque sheaves are acyclic for the functors $ℋ^{0}_Z$ and
$H^{0}_Z(X,⋅)$.

\item For every $q\geq0$, the functors $ℋ^{q}_Z$ and $H^{q}_Z(X,⋅)$
commute with restriction of scalars.
\end{enumerate}
\end{enonce*}

Let $M$ be a flasque sheaf. The exact sequence
\eqhyperref[V.6.5.4]{6.5.4} gives the equalities
$H^q_Z(X,M)=0$ for every $X$ and every $q\geq2$, together with an
exact sequence
$$
0\to H^{0}_Z(X,M)\to H^{0}(X,M)\to H^{0}(X\times U,M)
\to H^1_Z(X,M)\to0{\quoi}.
$$
But, since the sheaf $M$ is flasque, the morphism
$H^{0}(X,M)\to H^{0}(X\times U,M)$ is surjective
(\hyperref[V.sec4]{4}). Consequently, $H^1_Z(X,M)=0$, and $M$ is
acyclic for $H^{0}_Z(X,⋅)$. Passing to the associated sheaves, one
deduces that $M$ is acyclic for the functor $ℋ^{0}_Z$
(\hyperref[V.prop6.4]{6.4}). It is clear that the functors $ℋ^{0}_Z$
and $H^{0}_Z(X,⋅)$ commute with restriction of scalars; since flasque
sheaves are acyclic for both functors, assertion 2) follows.

\begin{enonce*}{Proposition 6.7}\etiquette[6.7]{V.prop6.7}
Let $(E,A)$ be a ringed topos, let $Z$ be a closed subtopos of $E$,
let $U$ be its open complement, and let $M$ and $N$ be two
$A$-modules. There are isomorphisms, functorial in $M$ and $N$ and
compatible with changes of closed subtopos:
$$
ℋ^{0}_Z(\Hhom_A(M,N))\simeq\Hhom_A(M,ℋ^0_Z(N))
\simeq\Hhom_A(M\otimes_AA_Z,N){\quoi}.
$$
\end{enonce*}

This follows from (IV 14), taking \hyperlink{sga4.exp.V.equation.6.3.2}{(6.3.2)} into account.

\setcounter{subsection}{7}
\subsection{}\etiquette[6.8]{V.subsec6.8}

% END INCLUDED SGA4_English_current_buildable_source_20260906/source/05/drafts/SGA4_V_LINES_2048_2300_English.texfrag


% BEGIN INCLUDED SGA4_English_current_buildable_source_20260906/source/05/drafts/SGA4_V_LINES_2301_2629_English.texfrag
Set
\begin{equation*}
\Hhom_{A,Z}(M,N)=ℋ^{0}_Z(\Hhom_A(M,N)){\quoi}.\tag{6.8.1}
\end{equation*}\etiquette[6.8.1]{V.6.8.1}
The functor $N\longmapsto\Hhom_{A,Z}(M,N)$ is left exact. Its derived
functors are denoted by $\Eext^q_{A,Z}(M,N)$: these are the
\emph{$\Eext$ sheaves with supports in $Z$}. Thus
\begin{equation*}
\Eext^{0}_{A,Z}(M,N)=\Hhom_{A,Z}(M,N){\quoi}.\tag{6.8.2}
\end{equation*}\etiquette[6.8.2]{V.6.8.2}
Set\pageoriginale $M_Z=M\otimes_AA_Z$. It follows from IV 14 that
$M_Z$ is the direct image on $E$ of the inverse image of $M$ on $Z$.
Thus, taking \hyperref[V.prop6.7]{6.7} into account, there are
isomorphisms
\begin{equation*}
\Eext^{q}_{A,Z}(M,N)\simeq\Eext^{q}_{A}(M_Z,N){\quoi}.\tag{6.8.3}
\end{equation*}\etiquette[6.8.3]{V.6.8.3}
We now pass to the global invariants. Set
\begin{equation*}
\Hom_{A,Z}(M,N)=\Ext^{0}_{A,Z}(E;M,N)
=H^{0}_Z(E,\Hhom_A(M,N)){\quoi}.\tag{6.8.4}
\end{equation*}\etiquette[6.8.4]{V.6.8.4}
The group $\Hom_{A,Z}(M,N)$ is the subgroup of the group of
morphisms from $M$ to $N$ whose support is in $Z$ (IV 14), \ie those
which vanish on $U$. More generally, for every object $X$ of $E$, set
\begin{equation*}
\Ext^{0}_{A,Z}(X;M,N)=H^{0}_Z(X,\Hhom_A(M,N)){\quoi}.\tag{6.8.5}
\end{equation*}\etiquette[6.8.5]{V.6.8.5}
The functors $N\longmapsto\Ext^{0}_{A,Z}(X;M,N)$ are left exact.
Their derived functors are denoted by $\Ext^q_{A,Z}(X;M,N)$. These
are the \emph{$\Ext$ groups with supports in $Z$}. Definitions
\eqhyperref[V.6.3.6]{6.3.6}, \eqhyperref[V.6.8.1]{6.8.1}, and
\eqhyperref[V.6.8.5]{6.8.5}, together with the isomorphisms of
\hyperref[V.prop6.7]{6.7}, give isomorphisms
\begin{equation*}
\Ext^{0}_{A,Z}(X;M,N)\simeq
\begin{cases}
H^{0}(X,\Hhom_{A,Z}(M,N)){\quoi},\\[3pt]
\Ext^{0}_A(X;M,ℋ^{0}_Z(N)){\quoi},\\[3pt]
\Ext^{0}_A(X;M_Z,N){\quoi}.
\end{cases}\tag{6.8.6}
\end{equation*}\etiquette[6.8.6]{V.6.8.6}
The last isomorphism in \eqhyperref[V.6.8.6]{6.8.6} gives
isomorphisms
\begin{equation*}
\Ext^{q}_{A,Z}(X;M,N)\simeq
\Ext^{q}_{A}(X;M_Z,N){\quoi}.\tag{6.8.7}
\end{equation*}\etiquette[6.8.7]{V.6.8.7}

\begin{enonce*}{Proposition 6.9}\etiquette[6.9]{V.prop6.9}
\begin{enumerate}
\renewcommand{\labelenumi}{\theenumi)}
\item There are two spectral sequences, functorial in $M$ and $N$ and
compatible with changes of closed subtopos:
\begin{equation*}
\Eext^{p+q}_{A,Z}(M,N)\Longleftarrow\tag{6.9.1}
\begin{cases}
E^{p,q}_2=ℋ^p_Z(\Eext^q_A(M,N)){\quoi},\\[3pt]
'E^{p,q}_2=\Eext^p_A(M,ℋ^q_Z(N)){\quoi}.
\end{cases}
\end{equation*}\etiquette[6.9.1]{V.6.9.1}

\item There\pageoriginale are three spectral sequences, functorial in
$X$, $M$, and $N$ and compatible with changes of closed subtopos:
\begin{equation*}
\Ext^{p+q}_{A,Z}(X;M,N)\Longleftarrow\tag{6.9.2}
\begin{cases}
E^{p,q}_2=H^p_Z(X,\Eext^q_A(M,N)){\quoi},\\[3pt]
'E^{p,q}_2=H^p(X,\Eext^q_{A,Z}(M,N)){\quoi},\\[3pt]
``E^{p,q}_2=\Ext^p_A(X;M,ℋ^q_Z(N)).
\end{cases}
\end{equation*}\etiquette[6.9.2]{V.6.9.2}

\item The sheaves $\Eext^q_{A,Z}(M,N)$ are canonically isomorphic to
the sheaves associated with the presheaves
$X\longmapsto\Ext^q_{A,Z}(X;M,N)$.
\end{enumerate}
\end{enonce*}

When $N$ is injective, the sheaf $\Hhom_A(M,N)$ is flasque
(\hyperref[V.prop4.10]{4.10}), and hence acyclic for $ℋ^{0}_Z$
(\hyperref[V.prop6.6]{6.6}). The first spectral sequence in
\eqhyperref[V.6.9.1]{6.9.1} is a spectral sequence of composite
functors \eqhyperref[V.6.8.1]{6.8.1}. Likewise, when $N$ is injective,
$ℋ^{0}_Z(N)$ is injective (\hyperref[V.prop4.11]{4.11}), and the
second spectral sequence in \eqhyperref[V.6.9.1]{6.9.1} is a spectral
sequence of composite functors deduced from
\hyperref[V.prop6.7]{6.7}. The spectral sequences in
\eqhyperref[V.6.9.2]{6.9.2} are spectral sequences of composite
functors deduced from definition \eqhyperref[V.6.8.5]{6.8.5} and the
first two isomorphisms in \eqhyperref[V.6.8.6]{6.8.6}. Finally, by
allowing $X$ to vary in the second spectral sequence in
\eqhyperref[V.6.9.2]{6.9.2} and taking the associated sheaves, one
obtains a spectral sequence of sheaves which degenerates by
\hyperref[V.prop3.1]{3.1} and gives the isomorphisms in 3).

\begin{enonce*}{Proposition 6.10}\etiquette[6.10]{V.prop6.10}
The functors $(M,N)\longmapsto\Eext^q_{A,Z}(M,N)$, $q\geq0$, and
$(M,N)\longmapsto\Ext^q_{A,Z}(X;M,N)$, $q\geq0$, are
$\delta$-functors in each of the variables $M$ and $N$. Denoting by
$M_U$ the sheaf $M\otimes_AA_U$ (\cf IV 11), there is a long exact
sequence
\begin{equation}
\ldots\xrightarrow{\delta}\Eext^q_{A,Z}(M,N)\to\Eext^q_A(M,N)
\to\Eext^q_A(M_U,N)\xrightarrow{\delta}
\Eext^{q+1}_{A,Z}(M,N)\to\ldots\tag{6.10.1}
\end{equation}\etiquette[6.10.1]{V.6.10.1}
and an analogous exact sequence for the $\Ext$ groups.
\end{enonce*}

The functors $\Eext^q_A(⋅,⋅)$ and $\Ext^q_A(X;⋅,⋅)$ form
$\delta$-functors with respect to each variable
(\hyperref[V.prop6.2]{6.2}), and the functor
$M\longmapsto M_Z$ is exact because $A_Z$ is flat
(\hyperref[V.cor1.3.3]{1.3.3}). The first assertion therefore follows
from the isomorphisms \eqhyperref[V.6.8.3]{6.8.3} and
\eqhyperref[V.6.8.7]{6.8.7}. The exact sequence
\eqhyperref[V.6.10.1]{6.10.1}, and the analogous exact sequence for
the $\Ext$ groups, is the long exact sequence of the $\delta$-functor
$$
\Eext^q_A(⋅,N),\ q\geq0
(\resp\ \Ext^q_A(X;⋅,N),\ q\geq0)
$$
associated\pageoriginale with the exact sequence
$0\to M_U\to M\to M_Z\to0$, taking
\eqhyperref[V.6.8.3]{6.8.3} and \eqhyperref[V.6.8.7]{6.8.7} into
account.

\setcounter{subsection}{10}
\subsection{}\etiquette[6.11]{V.subsec6.11}
Let us indicate briefly how these results may be extended to the case
of families of supports. Let $E$ be a topos. For every object $X$ of
$E$, denote by $\Fer(X)$ the set of closed subtopoi of the topos
$E_{/X}$. By passage to complements, this set is in bijective
correspondence with the set of open subtopoi of $E_{/X}$, \ie with the
set of subobjects of $X$ (IV 8). The presheaf
$\Fer:X\longmapsto\Fer(X)$ is in fact a $𝒰$-sheaf, as is immediately
verified, and is therefore representable. In other words, there is an
object of $E$, again denoted by $\Fer$, and a closed subtopos
$Z_{\Fer}$ of $E_{/\Fer}$ such that, for every $X$, every element of
$\Fer(X)$ is obtained from $Z_{\Fer}$ by base change along a morphism
$u_Z:X\to\Fer$ uniquely determined by $Z$.

\begin{enonce*}{Definition 6.12}\etiquette[6.12]{V.def6.12}
A family of supports of $E$ is a subset $\Phi$ of the set of closed
subtopoi of $E$ having the following properties:
\begin{itemize}
\item[{\rm (S1)}] The union of a finite family of elements of $\Phi$
belongs to $\Phi$.
\item[{\rm (S2)}] Every closed subtopos of $E$ contained in an element
of $\Phi$ is an element of $\Phi$.
\end{itemize}
\end{enonce*}

\setcounter{subsection}{12}
\setcounter{subsubsection}{0}
\subsubsection{}\etiquette[6.12.1]{V.subsubsec6.12.1}
Let $E$ be a topos, let $\Phi$ be a family of supports of $E$, and let
$X$ be an object of $E$. Denote by $\Phi(X)$ the smallest family of
supports of $E_{/X}$ containing the closed subtopoi of $E_{/X}$
obtained from closed subtopoi in $\Phi$ by the base change $X\to e$
($e$ being the final object of $E$). The functor
$X\longmapsto\Phi(X)$ is a subpresheaf of the sheaf $\Fer$. It is
therefore separated.

\subsubsection{}\etiquette[6.12.2]{V.subsubsec6.12.2}
A family $\Phi$ of supports of $E$ is said to be \emph{of local
character} if it has the following property:

\noindent
(LC) For every epimorphic family $(X_i\to e)$, $i\in I$ (where $e$ is
the final object of $E$), the sequence of sets
$$
\Phi\to\prod_i\Phi(X_i)\rightrightarrows
\prod_{i,j}\Phi(X_i\times X_j)
$$
is exact.

Let $\vieuxsheaf{a\Phi}$ be the sheaf associated with the presheaf
$X\longmapsto\Phi(X)$ (\hyperref[V.subsubsec6.12.1]{6.12.1}).
Condition (LC) is equivalent to the condition that the canonical
morphism $\Phi\to\vieuxsheaf{a\Phi}(e)$ be\pageoriginale a bijection
(\cf the construction of the associated sheaf in II in the case of a
separated presheaf). To verify (LC), one may therefore restrict to a
final family of epimorphic families $(X_i\to e)$, $i\in I$.

\begin{enonce*}[remark]{Example 6.12.3}\etiquette[6.12.3]{V.exem6.12.3}
\begin{enumerate}
\renewcommand{\labelenumi}{\theenumi)}
\hypertarget{sga4.exp.V.example.6.12.3.item.1}{}\item Let $Z$ be a closed subtopos of $E$. The set of closed subtopoi
of $E$ contained in $Z$ is a family of supports of $E$. It is of local
character.

\item Let $T$ be a topological space and let $\Phi$ be a
paracompactifying family of supports of $T$ \cite{sga4.exp.V.bib.7}. In general, the
family $\Phi$ is not of local character.

\item Let $T$ be a topological space and let $p$ be an integer. The
family $\Phi_p$ of closed subsets of $T$ of Krull codimension
$\geq p$ is a family of supports of $T$. It is of local character.

\item Let $E$ be a topos with the following property: every epimorphic
family $(X_i\to e)$, $i\in I$, is dominated by a finite epimorphic
family. Then every family of supports of $E$ is of local character.
Indeed, using \hyperlink{sga4.exp.V.example.6.12.3.item.1}{example 1}), it is enough to show that a filtered
inductive limit $\Phi_\lambda$ of families of local character is a
family of local character. This follows immediately by taking the
inductive limit in the sequence of sets
$$
\Phi_\lambda\to\prod_i\Phi_\lambda(X_i)
\rightrightarrows\prod_{i,j}\Phi_\lambda(X_i\times X_j){\quoi},
$$
where $(X_i\to e)$, $i\in I$, is a finite epimorphic family. Indeed,
by \hyperref[V.subsubsec6.12.2]{6.12.2} and the hypothesis on $E$, it
is enough to consider finite epimorphic families in verifying
condition (LC); and, since filtered inductive limits commute with
finite projective limits (I 2), exactness of these sequences of sets
is preserved upon passage to the inductive limit.

\item Let $E$ be a coherent topos (\hyperrefext[VI][VI.def2.3]{2.3}).
Then, by the preceding example, every family of supports of $E$ is of
local character.
\end{enumerate}
\end{enonce*}

\setcounter{subsection}{12}
\subsection{}\etiquette[6.13]{V.subsec6.13}
Let $(E,A)$ be a ringed topos, let $\Phi$ be a family of supports of
$E$, let $N$ be an $A$-module (a left module, to fix ideas), and let
$X$ be an object of $E$. Set
\hypertarget{sga4.exp.V.equation.6.13.1}{}\begin{align*}
ℋ^{0}_{\Phi}(N)&=\varinjlim_{Z\in\Phi}ℋ^{0}_Z(N){\quoi},\tag{6.13.1}\\
H^{0}_{\Phi}(E,N)&=\varinjlim_{Z\in\Phi}H^{0}_Z(E,N){\quoi}.
\tag{6.13.2}
\end{align*}\hypertarget{sga4.exp.V.equation.6.13.2}{}
Let\pageoriginale $M$ be a left $A$-module, and set
\hypertarget{sga4.exp.V.equation.6.13.3}{}\begin{gather*}
\Eext^{0}_{A,\Phi}(M,N)=\Hhom_{A,\Phi}(M,N)
=\varinjlim_{Z\in\Phi}\Eext^{0}_{A,Z}(M,N){\quoi},\tag{6.13.3}\\
\Ext^{0}_{A,\Phi}(X;M,N)=
\Hom_{A|X,\Phi(X)}(M|X,N|X)
=\varinjlim_{Z\in\Phi}\Ext^{0}_{A,Z}(X;M,N){\quoi}.
\tag{6.13.4}
\end{gather*}\hypertarget{sga4.exp.V.equation.6.13.4}{}
There are canonical isomorphisms
\begin{equation*}%
\Eext^{0}_{A,\Phi}(M,N)\simeq
ℋ^{0}_{\Phi}(\Eext^{0}_A(M,N)){\quoi},\tag{6.13.5}
\end{equation*}\etiquette[6.13.5]{V.6.13.5}

\begin{equation*}%
\Ext^{0}_{A,\Phi}(X;M,N)\simeq
H^{0}_{\Phi}(X,\Eext^{0}_A(M,N)){\quoi}.\tag{6.13.6}
\end{equation*}\etiquette[6.13.6]{V.6.13.6}

When $\Phi$ is the family of closed subtopoi contained in a fixed
closed subtopos $Z$, there are isomorphisms
\begin{equation*}
\begin{cases}
ℋ^{0}_{\Phi}\simeqℋ^{0}_Z,\\
H^{0}_{\Phi}\simeq H^{0}_Z,\\
\Eext^{0}_{A,\Phi}(⋅,⋅)\simeq\Eext^{0}_{A,Z}(⋅,⋅),\\
\Ext^{0}_{A,\Phi}(X;⋅,⋅)\simeq\Ext^{0}_{A,Z}(X;⋅,⋅){\quoi}.
\end{cases}\tag{6.13.7}
\end{equation*}\etiquette[6.13.7]{V.6.13.7}
The inductive limits defining the preceding functors are filtered.
Consequently, these functors are left exact. Their right derived
functors are denoted by
\begin{equation}
ℋ^q_{\Phi},\ H^q_{\Phi},\ \Eext^q_{A,\Phi}(⋅,⋅),\
\Ext^q_{A,\Phi}(X;⋅,⋅){\quoi}.\tag{6.13.8}
\end{equation}\etiquette[6.13.8]{V.6.13.8}
They too are computed by inductive limits of the functors $ℋ^q_Z$,
$H^q_Z$, \etcn as $Z$ ranges through $\Phi$.

From the isomorphisms \eqhyperref[V.6.13.5]{6.13.5} and \hyperref[V.6.13.6]{(6.13.6)}, one
obtains two spectral sequences by taking inductive limits in the first
spectral sequence in \eqhyperref[V.6.9.1]{6.9.1} and the first spectral
sequence in \eqhyperref[V.6.9.2]{6.9.2}:
\hypertarget{sga4.exp.V.equation.6.13.9}{}\begin{gather*}
\Eext^{p+q}_{A,\Phi}(M,N)\Leftarrow E^{p,q}_2=
ℋ^p_{\Phi}(\Eext^q_A(M,N)){\quoi},\tag{6.13.9}\\
\Ext^{p+q}_{A,\Phi}(X;M,N)\Leftarrow E^{p,q}_2=
H^p_{\Phi}(X,\Eext^q_A(M,N)){\quoi}.\tag{6.13.10}
\end{gather*}\hypertarget{sga4.exp.V.equation.6.13.10}{}

\subsection{}\etiquette[6.14]{V.subsec6.14}\pageoriginale
With the notation of \hyperref[V.subsec6.13]{6.13}, and without any
hypothesis on $\Phi$, there is a functorial morphism
\begin{equation*}
\theta:H^{0}_{\Phi}(N)\to H^{0}(E,ℋ^{0}_{\Phi}(N)){\quoi}.
\tag{6.14.1}
\end{equation*}\etiquette[6.14.1]{V.6.14.1}
This morphism is always injective, but is not in general an
isomorphism. However, if $\Phi$ is \emph{of local character}, the
morphism $\theta$ is an isomorphism, as is immediately verified.
Likewise, if $\Phi$ is \emph{of local character}, there is an
isomorphism
\begin{equation}
\theta':\Ext^{0}_{A,\Phi}(E;M,N)\simeq
H^{0}(E,\Eext^{0}_{A,\Phi}(M,N)){\quoi}.\tag{6.14.2}
\end{equation}\etiquette[6.14.2]{V.6.14.2}

\begin{enonce*}{Proposition 6.15}\etiquette[6.15]{V.prop6.15}
Let $(E,A)$ be a ringed topos such that $E$ is coherent \emph{(VI)},
let $\Phi$ be a family of supports of $E$, and let $M$ and $N$ be two
$A$-modules. There are two spectral sequences:
\hypertarget{sga4.exp.V.equation.6.15.1}{}\begin{align*}%\etiquette[6.15.1]{V.6.15.1}
H^{p+q}_{\Phi}(E,N)\Leftarrow E^{p,q}_2
&=H^p(E,ℋ^q_{\Phi}(N)){\quoi},\tag{6.15.1}\\
\Ext^{p+q}_{A,\Phi}(E;M,N)\Leftarrow E^{p,q}_2
&=H^p(E,\Eext^q_{A,\Phi}(M,N)){\quoi}.\tag{6.15.2}
\end{align*}\hypertarget{sga4.exp.V.equation.6.15.2}{}
\end{enonce*}

These spectral sequences are deduced from the spectral sequence
\eqhyperref[V.6.4.2]{6.4.2} and the second spectral sequence in
\eqhyperref[V.6.9.2]{6.9.2} by taking the inductive limit over the
closed subtopoi $Z$ in $\Phi$, taking into account that the cohomology
of a coherent topos commutes with inductive limits of sheaves (IV 5).

\setcounter{subsection}{15}
\subsection{}\etiquette[6.16]{V.subsec6.16}
We mention, without proof, that the second spectral sequence in
\eqhyperref[V.6.9.1]{6.9.1} and the third spectral sequence in
\eqhyperref[V.6.9.2]{6.9.2} (with $X$ equal to the final object of
$E$) can be extended to the case of families of supports, provided
that the topos $E$ is coherent and the module $M$ is perfect
\cite{sga4.exp.V.bib.13}.

Finally, one can also generalize the notion of families of supports
by introducing presheaves of families of supports, sheaves of families
of supports, and the corresponding cohomology groups\pageoriginale
and sheaves\footnote{\cf \cite{sga4.exp.V.bib.12}, Chap.~IV, \S 1 (Lecture Notes 20,
Springer), for the development of these notions in the form of a fugue
with variations.}.

% END INCLUDED SGA4_English_current_buildable_source_20260906/source/05/drafts/SGA4_V_LINES_2301_2629_English.texfrag


% BEGIN INCLUDED SGA4_English_current_buildable_source_20260906/source/05/drafts/SGA4_V_LINES_2630_2810_English.texfrag
\section{Appendix: Čech cohomology}\etiquette[7]{V.sec7}

Developing an idea due to P. Cartier, we show in this section how,
in an arbitrary topos, one can compute the cohomology of a sheaf by
means of coverings. For the classical theory of paracompact spaces,
we refer to \cite{sga4.exp.V.bib.7}; for another method which applies to certain
topoi, and in particular to the étale topos, see \cite{sga4.exp.V.bib.14}.

\subsection{Skeleton and coskeleton}\etiquette[7.1]{V.subsec7.1}

\setcounter{subsubsection}{-1}
\subsubsection{}\etiquette[7.1.0]{V.subsubsec7.1.0}

Let $𝚫$ be the category of standard simplices (the objects of
$𝚫$ are the ordered sets $\Delta_n=[0,\ldots,n]$, and the
morphisms are the increasing maps). Let $E$ be a $𝒰$-topos, and
let $𝚫(E)$ be the category of presheaves on $𝚫$ with values in
$E$; in other words, the category of \emph{semi-simplicial objects}
of $E$. Denote by $𝚫[n]$ the full subcategory of $𝚫$ defined by
the objects $\Delta_p$, $p\leq n$, by
$i_n:𝚫[n]\to𝚫$ the inclusion functor, and by $𝚫 E[n]$ the
category of presheaves on $𝚫[n]$ with values in $E$; in other
words, the category of \emph{semi-simplicial objects of $E$
truncated at order} $n$. By I \hyperref[I.prop5.1]{5.1}, there is a sequence of three
adjoint functors (I \hyperref[I.sec5.3]{5.3})
\[
\xymatrix@C=3cm@M=.5cm{
𝚫 E[n] \ar@<17pt>[r]^{i_{n!}} \ar@<-19pt>[r]^{i_{n\ast}}&
𝚫 E \ar@<1pt>[l]_{i^{\ast}_n}\\
}
\]
where $i_n^\ast$ is the restriction functor to the category
$𝚫[n]$, and $i_{n\ast}$ and $i_{n!}$ are its right and left
adjoints, respectively. Note that $𝚫 E$ and $𝚫 E[n]$ are
$𝒰$-topoi (IV \hyperref[IV.theo1.2]{1.2}), and that $(i_n^\ast,i_{n\ast})$ is a morphism
of topoi (IV \hyperref[IV.def3.1]{3.1}) which is an embedding of $𝚫 E[n]$ into
$𝚫 E$ (I \hyperref[I.prop5.6]{5.6} and IV \hyperref[IV.def9.1.1]{9.1.1}).

\begin{enonce*}{Definition 7.1.1}\etiquette[7.1.1]{V.def7.1.1}
The\pageoriginale functor $i_{n!}i_n^\ast$ (\resp
$i_{n\ast}i_n^\ast$) from
$𝚫_E$ to $𝚫_E$ is denoted by $\sk_n$ (\resp $\cosk_n$) and is
called the \emph{skeleton functor of order $n$} (\resp the
\emph{coskeleton functor of order $n$}).
\end{enonce*}

The skeleton and coskeleton functors are used constantly in the
theory of semi-simplicial sets \cite{sga4.exp.V.bib.3}.

\begin{enonce*}{Proposition 7.1.2}\etiquette[7.1.2]{V.prop7.1.2}
\begin{enumerate}
\renewcommand{\labelenumi}{\theenumi)}
\item The adjunction morphism $\sk_n\to\id$ is a monomorphism. The
canonical morphisms $\sk_n\sk_m\to\sk_n$ (\resp
$\cosk_n\to\cosk_n\cosk_m$) are isomorphisms when $n\leq m$
(\resp $n\geq m$).

\item Let $u:E\to E'$ be a morphism of topoi. The functor $u^\ast$,
extended to semi-simplicial and truncated semi-simplicial objects,
commutes with the functors $i_{n!}$, $i_n^\ast$, $\sk_n$, and
$\cosk_n$.
\end{enumerate}
\end{enonce*}

The first assertion follows immediately from the definitions. To
prove 2), it suffices to observe that, by I \hyperref[I.prop5.1]{5.1}, the functors
$i_{n!},\ldots,\cosk_n$ are computed by means of inductive limits
and finite projective limits.

\subsection{An acyclicity lemma}\etiquette[7.2]{V.subsec7.2}

\setcounter{subsubsection}{-1}
\subsubsection{}\etiquette[7.2.0]{V.subsubsec7.2.0}

Let $M$ be an abelian group of $E$. The \emph{homology} of a
semi-simplicial object $K$ of $E$ with coefficients in $M$ is
defined in the usual way. First consider $A(K)$, the free abelian
semi-simplicial group generated by $K$, and then form the tensor
product $M\otimes_{\ZZ}A(K)$; this gives an abelian semi-simplicial
group of $E$. One then considers the associated complex of abelian
groups (formed by taking the alternating sum of the boundary
operators) and takes its homology. The homology objects are denoted
by $H_i(K,M)$. They are abelian groups of $E$. When $M$ is the group
$\mathbb{Z}_E$ (the free abelian group generated by the final object
of $E$), we write simply $H_i(K)$. The groups $H_i(K)$ are called
the \emph{homology groups} of $K$. A semi-simplicial object $K$ is
said to be \emph{acyclic} if, for every abelian group $M$ of $E$,
the groups $H_i(K,M)$ vanish for $i>0$.

Formation of homology commutes with inverse images under morphisms
of topoi\pageoriginale.

The purpose of this section is to prove the following lemma:

\begin{enonce*}{Lemma 7.2.1}\etiquette[7.2.1]{V.lem7.2.1}
Let $E$ be a topos, let $K_{∙}$ and $L_{∙}$ be two semi-simplicial
objects of $E$, let $v_{∙}:K_{∙}\to L_{∙}$ be a morphism of
semi-simplicial objects, and let $n$ be an integer $\geq0$. Suppose
that $v$ has the following properties:
\begin{enumerate}
\renewcommand{\labelenumi}{\theenumi)}
\item For every integer $p<n$ $(p\geq0)$, the morphism
$v_p:K_p\to L_p$ is an isomorphism.

\item The morphism $K_n\xrightarrow{v_n}L_n$ is an epimorphism.

\item The canonical morphisms
$K_{∙}\to\cosk_n(K_{∙})$ and $L_{∙}\to\cosk_n(L_{∙})$ are
isomorphisms.
\end{enumerate}
Then, for every integer $p$, the morphism $v_p$ is an epimorphism,
and the morphism $v_{∙}$ induces an isomorphism on homology objects.
\end{enonce*}

\begin{enonce*}[remark]{Remark}
The proof in fact shows that, after adopting a suitable definition
of homotopy in topoi, the morphism $v_{∙}$ is a homotopy equivalence.
\end{enonce*}

\begin{enonce*}[remark]{Proof}
We may suppose that $E$ is the topos of sheaves on a small site $C$
(IV 1), and hence that $E$ is a subtopos of a topos $E'$ having
enough fiber functors (for example, the topos $C^{\wedge}$). Denote
by $a^\ast:E'\to E$ the inverse-image functor of the embedding
morphism $E\hookrightarrow E'$. Let
$$
v_{∙}[n]:K_{∙}[n]\to L_{∙}[n]
$$
be the morphism of truncated semi-simplicial objects obtained by
truncating $v_{∙}$ at order $n$. Denote by $L_{∙}[n]'$ the image (in
$E'$) of $K_{∙}[n]$ under $v_{∙}[n]$, and by
$v_{∙}[n]':K_{∙}[n]\to L_{∙}[n]'$ the morphism induced by
$v_{∙}[n]$. Put
$L!=i_{n\ast}L_{∙}[n]'$, $K!=i_{n\ast}K_{∙}[n]$, and
$v!=i_{n\ast}v_{∙}[n]'$ (here $i_{n\ast}$ is relative to $E'$).
By 1), the morphism $v!:K!\to L!$ has properties 1), 2), and 3) of
the lemma. Moreover, by 2), 3), and \hyperref[V.prop7.1.2]{7.1.2},
the morphism $a^\ast v!$ is none other than $v_{∙}$. It therefore
suffices to prove the lemma for $v!$, and hence we may suppose that
$E$ has enough fiber functors. Using these fiber functors, we may
then reduce to the case where $E$ is the category of sets.
\end{enonce*}

First observe that the hypotheses of the lemma on the morphism
$v_{∙}$ are\pageoriginale stable under every base change
$M_{∙}\to L_{∙}$ for which $M_{∙}$ is a coskeleton of order $n$.
Consequently, the fiber $P_{∙}$ of $v_{∙}$ at any base point of
$L_{∙}$ is of the form $i_{n\ast}P_{∙}[n]$, where $P_{∙}[n]$ is a
nonempty complex truncated at order $n$ of the form
\begin{equation*}
P_n\to e\to e\to e\to\cdots\to e{\quoi}.\tag{7.2.2}
\end{equation*}\etiquette[7.2.2]{V.7.2.2}
For every integer $i$, denote by $\delta\sheaf{\Delta}_i$ the
semi-simplicial set forming the boundary of the simplex
$\sheaf{\Delta}_i$. Every morphism from
$\delta\sheaf{\Delta}_i$ to $P_{∙}$ extends to a morphism from
$\sheaf{\Delta}_i$ to $P_{∙}$. Indeed, this property is immediate
for $i\geq n$, because $P_{∙}$ is a coskeleton of order $n$, and it
is clear for $i<n$ from the description \eqhyperref[V.7.2.2]{7.2.2}.
Since $P_{∙}$ is a Kan complex, it is homotopically trivial
\cite{sga4.exp.V.bib.6}. Note that the morphism $v_{∙}$ is a fibration in the sense
of Kan \cite{sga4.exp.V.bib.6}. Since the fibers of this fibration are
homotopically trivial, $v_{∙}$ induces an isomorphism on the homotopy
groups of $K_{∙}$ and $L_{∙}$ at all base points of $K_{∙}$. By the
Hurewicz theorem, this implies that $v_{∙}$ induces an isomorphism on
homology \cite{sga4.exp.V.bib.6}. This proves the second assertion of the lemma. To
prove the first, observe that every morphism from a complex truncated
at order $n$ to $L_{∙}[n]$ lifts to a morphism to $K_{∙}[n]$.
Consequently, by 3), every morphism from a complex to $L_{∙}$ lifts
to a morphism to $K_{∙}$. Thus $v_{∙}$ is surjective.

% END INCLUDED SGA4_English_current_buildable_source_20260906/source/05/drafts/SGA4_V_LINES_2630_2810_English.texfrag


% BEGIN INCLUDED SGA4_English_current_buildable_source_20260906/source/05/drafts/SGA4_V_LINES_2811_3070_English.texfrag
\subsection{Hypercoverings}\etiquette[7.3]{V.subsec7.3}

\setcounter{subsubsection}{-1}
\subsubsection{}\etiquette[7.3.0]{V.subsubsec7.3.0}
Let $C$ be a site belonging to the universe $𝒰$, in which fiber
products and products of two objects are representable. Denote by
$C^{\wedge}$ the topos of $𝒰$-presheaves on $C$, and by
$C^{\sim}$ the topos of $𝒰$-sheaves on $C$. An object $K$ of
$C^{\wedge}$ is called \emph{semi-representable} if it is isomorphic
to a coproduct of representable presheaves.

Let $SR(C)$ be the full subcategory of $C^{\wedge}$ defined by the
semi-representable objects. In the category $SR(C)$, projective
limits indexed by finite nonempty categories are representable, and
the inclusion functor $SR(C) \hookrightarrow C^{\wedge}$ commutes
with these finite projective limits. By a remark already made, it
follows that if $K_∙$ is a semi-simplicial object truncated in degree
$n$ in $C^{\wedge}$, all of whose objects\pageoriginale are
semi-representable, then the extension $i^+_n(K_∙)$ has
the same property. Thus, if $K_∙$ is a semi-simplicial object of
$C^{\wedge}$ all of whose objects are semi-representable, then, for
every $n$, the object $\cosk_n (K_∙)$ has the same property.

\subsubsection{Definitions and notation}\etiquette[7.3.1]{V.subsubsec7.3.1}

\paragraph{}\etiquette[7.3.1.1]{V.par7.3.1.1}

Let $p > 0$ be an integer. A semi-simplicial object $K_∙$ of
$C^{\wedge}$ is called a \emph{hypercovering of type} $p$ of $C$,
or, when no confusion can result, a hypercovering of type $p$, if it
has the following properties:
\begin{itemize}
\item[HR1)] For every integer $n \geq 0$, $K_n$ is semi-representable.

\item[HR2)] The canonical morphism $K_∙ \to \cosk_p (K_∙)$ is an
isomorphism.

\item[HR3)] For every integer $n \geq 0$, the canonical morphism of
presheaves $K_{n + 1} \to (\cosk_n(K_∙))_{n + 1}$ is a covering
morphism of presheaves (II \hyperref[II.rem6.2]{6.2}). The canonical morphism of
presheaves $K_{0} \to e$, where $e$ is the final object of
$C^{\wedge}$, is a covering morphism of presheaves.
\end{itemize}

\paragraph{}\etiquette[7.3.1.2]{V.par7.3.1.2}
A hypercovering of type $p$ is a hypercovering of type $q$ for every
$q \geq p$. A semi-simplicial object $K_∙$ will be called a
\emph{hypercovering} (of $C$) if it has properties HR 1) and HR 3).

\paragraph{}\etiquette[7.3.1.3]{V.par7.3.1.3}
Denote by $HR_p$ (\resp $HR$), and call the category of
hypercoverings of type $p$ (\resp the category of hypercoverings),
the following category:
\begin{enumerate}
\renewcommand{\theenumi}{\alph {enumi}}
\renewcommand{\labelenumi}{\theenumi)}
\item The objects of $HR_p$ (\resp $HR$) are the hypercoverings of
  type $p$ (\resp the hypercoverings).

\item Let $K_∙$ and $L_∙$ be two objects of $HR_p$ (\resp $HR$). A
  morphism of $HR_p$ (\resp $HR$) with source $K_∙$ and target $L_∙$
  is a morphism $v_∙ : K_∙ \to L_∙$ of semi-simplicial objects of
  $C^{\wedge}$.
\end{enumerate}

\paragraph{}\etiquette[7.3.1.4]{V.par7.3.1.4}
Let $C$ be a site in which fiber products are representable, and let
$X$ be an object of $C$. A hypercovering of $X$ (\resp a
hypercovering of type $p$ of $X$) means a hypercovering of the site
$C/X$ (\resp a hypercovering of type $p$ of $C/X$). The category of
hypercoverings of $X$ (\resp of hypercoverings of type $p$ of $X$)
is defined similarly\pageoriginale. There is a sequence of inclusion
functors
$$
\cdots HR_p \hookrightarrow HR_{p + 1} \hookrightarrow \cdots
\hookrightarrow HR{\quoi}.
$$
These functors are fully faithful. We denote by $HR_{\infty}$ the
inductive limit of the categories $HR_p$.

\setcounter{paragraph}{5}
\paragraph{}\etiquette[7.3.1.6]{V.par7.3.1.6}
For every integer $n \geq 0$, denote by $\sheaf{\Delta}_n^c$ the
standard semi-simplicial object of dimension $n$ with values in the
category of sets, \ie the functor on $𝚫$ represented by the
ordered set $[0, n]$. Also denote by $\sheaf{\Delta}_n^c$ the
presheaf on $C$, with values in the category of semi-simplicial sets,
that is constant with value $\sheaf{\Delta}_n$. Thus the presheaf
$\sheaf{\Delta}_n^c$ is a presheaf of semi-simplicial sets. The two
canonical injections from $\Delta_0$ into $\Delta_1$ define two
morphisms of semi-simplicial presheaves:
\[
\xymatrix@C=2cm@M=.25cm{
\underline{\Delta}_0^{c} \ar@<3pt>[r]^{e_{0}} \ar@<-3pt>[r]_{e_{1}} &
\underline{\Delta}_1^c ,
}
\]
and consequently define, for every semi-simplicial presheaf $K_∙$,
two canonical injections:
\[
\xymatrix@C=2cm@M=.25cm{
K_{∙} \ar@<3pt>[r]^(.45){e_{0}} \ar@<-3pt>[r]_(.45){e_{1}}&  K_{∙}
\times \underline{\Delta}_1^c .
}
\]
Two morphisms
$
\xymatrix@C=1cm@M=.25cm{
K_{∙} \ar@<3pt>[r]^(.45){u_{0}} \ar@<-3pt>[r]_(.45){u_{1}}&  L_{∙}
}
$
of semi-simplicial presheaves are called \emph{homotopic morphisms}
if there exists a morphism
$$
v : K_∙ \times \sheaf{\Delta}_1^c \to L_∙
$$
such that the diagrams
\[
\xymatrix@R=.5cm@C=1.25cm{
K_∙ \ar[rr]^{e_i} \ar[rdd]_{u_i}&   &K_{∙}\times \sheaf{\Delta}_1^c
\ar[ldd]^v\\
                                    &   & &\txt{i = 0,1}\\
                       & L_{∙} &
}
\]
are\pageoriginale commutative. The morphism $v$ is called a
\emph{homotopy joining $u_0$ to $u_1$}. The relation ``$u_0$ and
$u_1$ are two homotopic morphisms from $K_∙$ to $L_∙$'' is not, in
general, an equivalence relation. Nevertheless, the equivalence
relation generated by this relation is compatible with composition
of morphisms. This therefore allows us to define the category of
semi-simplicial presheaves up to homotopy by passing to the quotient
by the equivalence relation generated by the homotopy relation.

\paragraph{}\etiquette[7.3.1.7]{V.par7.3.1.7}
In this way we define the categories $\sheaf{HR}_p$,
$\sheaf{HR}_{\infty}$, and $\sheaf{HR}$, the \emph{categories of
hypercoverings up to homotopy}.

\begin{enonce*}{Theorem 7.3.2}\etiquette[7.3.2]{V.thm7.3.2}
Let $C$ be a site satisfying the conditions of
\hyperref[V.subsubsec7.3.0]{7.3.0}.
\begin{enumerate}
\renewcommand{\labelenumi}{\theenumi)}
\item The category $\sheaf{HR}^{0}_p$ (\resp $\sheaf{HR}^{0}$) is
  filtered \emph{(I \hyperref[I.def2.7]{2.7})}.

\item Let $K_∙$ be an object of $HR_p$ (\resp of $HR$), let $n$ be
  an integer such that $0 \leq n \leq p$ (\resp $0 \leq n$), and
let
$$
u : X \to K_n
$$
be a covering morphism of presheaves. There exists an object $L_∙$
of $HR_p$ (\resp of $HR$) and a morphism $f : L_∙ \to K_∙$ such that
the morphism
$$
f_n : L_n \to K_n
$$
factors through $u$.

\item The associated semi-simplicial sheaves \emph{(II \hyperref[II.def3.5]{3.5})} of
hypercoverings are acyclic (\hyperref[V.subsubsec7.2.0]{7.2.0}) in
strictly positive degrees. The $0$-th homology sheaf is isomorphic
to the sheaf associated to the constant presheaf $𝒵$.
\end{enumerate}
\end{enonce*}

\begin{enonce*}[remark]{Proof}
We prove 3). Let ${K}^{\sim}_{∙}$ be the semi-simplicial sheaf
associated to $K_∙$. The ``associated sheaf'' functor commutes with
the functor $\cosk_n$ (\hyperref[V.prop7.1.2]{7.1.2}. 2)). It follows
that the semi-simplicial sheaf ${K}^{\sim}_{∙}$ has the following
properties:
\begin{enumerate}
\renewcommand{\theenumi}{\alph {enumi}}
\renewcommand{\labelenumi}{\theenumi)}
\item The canonical morphism $K_0 \to e$ is an epimorphism of
  sheaves.

\item For every integer $n \geq 0$, the canonical morphism
  $K_{n + 1} \to (\cosk_n(K_∙))_{n + 1}$ is an epimorphism of
  sheaves.
\end{enumerate}
\end{enonce*}

For $n \geq 0$, put ${L_∙}_n = \cosk_n {K}^{\sim}_{∙}$. Then, for
every $n$, there is a sequence of morphisms of semi-simplicial
sheaves
$$
{K_∙}^{\sim} \xrightarrow{u_n}{L_∙}_n\xrightarrow{v_n} {L_∙}_{n - 1}
\cdots \to {L_∙}_{0} \xrightarrow{v_0} e_∙{\quoi},
$$
and\pageoriginale, for every $j$, $0 \leq j \leq n$, $v_j$ has the
properties of \hyperref[V.lem7.2.1]{7.2.1}. Moreover, the morphism
$u_n$ induces an isomorphism on the components of degree $\leq n$.
It then follows from \hyperref[V.lem7.2.1]{7.2.1}, by induction on
$n$, that $K_∙$ is acyclic.

To prove 1) and 2), we introduce some terminology.

\begin{enonce*}{Definition 7.3.3}\etiquette[7.3.3]{V.def7.3.3}
Let $f : X_{∙} \to Y_∙$ be a morphism of semi-simplicial presheaves.
We say that $f$ is \emph{special of type} $p$ (\resp
\emph{special}) if:
\begin{enumerate}
\renewcommand{\labelenumi}{\theenumi)}
\item The objects $X_{∙}$ and $Y_∙$ are canonically isomorphic to
  their coskeleta of order $p$, and $\cosk_p (f) = f$ (\resp no
  conditions are imposed on $f$, \sisi{$X^{∙}$}{$X_{∙}$}, or
  \sisi{$Y^{∙}$}{$Y_{∙}$}).

\item For every integer $n$ such that $0 \leq n \leq p$ (\resp
  $0 \leq n$), the morphism $\Phi_{n + 1}$ appearing in the diagram
  below is covering:
\[
\xymatrix@C=1cm@R=1.25cm{
X_{n+1} \ar[rrr]^{f_{n+1}} \ar[rd]^{\Phi_{n+1}} \ar[dd]       & & & Y_{n+1}\ar[dd]\\
                                    &P_{n+1} \ar[ld] \ar[rru]& &\\
(\cosk_n X_{∙})_{n+1}\ar[rrr]_{(\cosk_n f)_{n+1}}& & &
(\cosk_n Y_{∙})_{n+1} .
}
\]
(The vertical arrows are defined by the canonical morphisms
$X_{∙} \to \cosk_n X_{∙}$ and $Y_∙ \to \cosk_n Y_∙$. The object
$P_{n+1}$ is the fiber product, and $\Phi_{n+1}$ is the unique arrow
making the diagram commutative.)

\item The morphism $f_0 : X_0 \to Y_0$ is covering.
\end{enumerate}

A semi-simplicial presheaf $K_∙$ is called special of type $p$
(\resp special) if the canonical morphism
$$
K_∙ \to e_{∙}\quad (e_{∙}\quad \mbox{\rm is the final semi-simplicial object})
$$
is special of type $p$ (\resp special), \ie if $K_∙$ satisfies
conditions \emph{HR 2)} and \emph{HR 3)} (\resp \emph{HR 3)}) of
\hyperref[V.par7.3.1.1]{7.3.1.1}.
\end{enonce*}

\begin{enonce*}{Lemma 7.3.4}\etiquette[7.3.4]{V.lem7.3.4}
\begin{enumerate}
\renewcommand{\labelenumi}{\theenumi)}
\item The\pageoriginale composite of two special morphisms (\resp
  special morphisms of type $p$) is a special morphism (\resp a
  special morphism of type $p$).

\item Let $K_∙$ be a special semi-simplicial presheaf (\resp special
  of type $p$), let $X_{∙}\xrightarrow{f}Y_∙$ be a special morphism
  (\resp special of type $p$), and let $u : K_∙ \to Y_∙$ be a
  morphism of complexes. Then the fiber product $P_{∙} =
  K_∙ ×_{Y_∙} X_∙$ is special (\resp special of type $p$).
\end{enumerate}
\end{enonce*}

The proof of this lemma is left to the reader.

\setcounter{subsubsection}{4}

% END INCLUDED SGA4_English_current_buildable_source_20260906/source/05/drafts/SGA4_V_LINES_2811_3070_English.texfrag


% BEGIN INCLUDED SGA4_English_current_buildable_source_20260906/source/05/drafts/SGA4_V_LINES_3071_3430_English.texfrag
\subsubsection{}\etiquette[7.3.5]{V.subsubsec7.3.5}
\emph{Proof of assertion} 2) \emph{of} \hyperref[V.thm7.3.2]{7.3.2}.
We may first suppose that $X$ is semi-representable. For every
semi-simplicial object $L_∙$, denote by $\Hom_{(X)}(L_∙,K_∙)$ the
set of morphisms of semi-simplicial objects equipped with a
factorization $L_n\to X\xrightarrow{u}K_n$. The functor
$\Hom_{(X)}(⋅,K_∙)$ is representable. Indeed, this functor carries
inductive limits to projective limits, and the category
$𝚫(C^{\wedge})$ is a topos. Denote by $P_{∙}$ an object
representing this functor. For an object $Z$ of $C^{\wedge}$,
likewise denote by $j_n^{+}(Z)$ the object representing the functor
$L_∙\mapsto\Hom(L_n,Z)$ on $𝚫 C^{\wedge}$, whose construction by
$\varprojlim$ is made explicit in III \hyperref[III.def1.1]{1.1}. Observe here that this
construction involves only finite projective limits. It follows
immediately that $j_n^{+}$ carries semi-representable presheaves to
semi-simplicial presheaves having semi-representable components, and
covering morphisms $Z'\to Z$ to special morphisms
(\hyperref[V.def7.3.3]{7.3.3}).

By the definition of $P_{∙}$, there is a cartesian square:
\[
\xymatrix@C=2cm@R=1.5cm{
P_{∙} \ar[r]\ar[d]& j_n^{+}(X) \ar[d]^{j_n^{+}(u)}\\
K_{∙} \ar[r]      &  j_n^{+}(K_n)\quad.
}
\]
By what has just been observed, the semi-simplicial objects
$j_n^{+}(X)$ and $j_n^{+}(K_n)$ are semi-representable, and the
morphism $j_n^{+}(u)$ is a special morphism of type $n$. The result
now follows by applying \hyperref[V.cor3.4]{3.4} and the sorites
\hyperref[V.subsubsec7.3.0]{7.3.0}.

\setcounter{subsubsection}{6}
\setcounter{paragraph}{-1}
\paragraph{}\etiquette[7.3.6.0]{V.par7.3.6.0}
Let\pageoriginale $M_∙$ and $N_∙$ be two semi-simplicial
presheaves. The functor
$$
L_∙\longmapsto\Hom(L_∙\times M_∙,N_∙)
$$
is representable. The semi-simplicial presheaf representing it will
be denoted by
$$
\Hhom_∙(M_∙,N_∙){\quoi}.
$$
The degree-$n$ component presheaf of this object will be denoted by
$\Hhom_n(M_∙,N_∙)$. There is a canonical isomorphism
$$
\Hhom_n(M_∙,N_∙)\xrightarrow{\sim}
\Hhom_0(M_∙\times\sheaf{\Delta}^c_n,N_∙){\quoi}.
$$

\begin{enonce*}{Lemma 7.3.6}\etiquette[7.3.6]{V.lem7.3.6}
Let
\[
\xymatrix@C=2cm@R=1.5cm{
M_{∙} \ar@{^{(}->}[r]& N_{∙}\\
\sk_n\sheaf{\Delta}_{n+1}^c \ar[u] \ar@{^{(}->}[r]^{u}&
\sheaf{\Delta}_{n+1}^c \ar[u] &
\kern -1.5cm \txt{\rm ($u$ the canonical injection)}
}
\]
be a cocartesian diagram of semi-simplicial presheaves.

For every hypercovering $L_∙$, the morphism
$$
\Hhom_0(N_∙,L_∙)\to\Hhom_0(M_∙,L_∙)
$$
is covering.
\end{enonce*}

\begin{enonce*}[remark]{Proof}
There is a cartesian diagram:
\[
\xymatrix{
\Hhom_0(N_{∙},L_{∙})\ar[d] \ar[r]&
\Hhom_0(M_{∙},L_{∙}) \ar[d]\\
\Hhom_0(\sheaf{\Delta}^c_{n+1},L_{∙})\ar[r]^{(\ast)} &
\Hhom_0(\sk_n\sheaf{\Delta}^c_{n+1},L_{∙}).
}
\]
It therefore suffices to show that the morphism $(\ast)$ is
covering. But $(\ast)$ is isomorphic to the morphism
$$
L_{n+1}\to(\cosk_nL_∙)_{n+1},
$$
which is covering by hypothesis. \cqfd
\end{enonce*}

\subsubsection{}\etiquette[7.3.7]{V.subsubsec7.3.7}
Proof of assertion 1) of \hyperref[V.thm7.3.2]{7.3.2}. The product
of two objects of $HR_p$ (\resp of $HR$) is again an object of
$HR_p$ (\resp of $HR$) (\hyperref[V.lem7.3.4]{7.3.4}). To prove
that the category $\sheaf{HR}^{0}_p$ (\resp $\sheaf{HR}^{0}$) is
filtered, it suffices to show that, given two morphisms in $HR_p$
(\resp $HR$)
\[
\xymatrix@C=2cm@M=.2cm{
K_{∙} \ar@<3pt>[r]^{u_0} \ar@<-3pt>[r]_{u_1}&L_∙
}
\]
there exists a morphism $v:M_∙\to K_∙$ in $HR_p$ (\resp in $HR$)
such that the morphisms
\[
\xymatrix@C=2cm@M=.2cm{
M_{∙} \ar@<3pt>[r]^{u_0v} \ar@<-3pt>[r]_{u_1v}&L_∙
}
\]
are homotopic,\pageoriginale \ie such that there exists
$w:M_∙\times\sheaf{\Delta}_1^c\to L_∙$ making the diagrams
\[
\vcenter{
\xymatrix@C=2cm@R=1.5cm{
M_{∙} \ar[d]_{v} \ar[r]^(.4){e_i}&
M_{∙}\times\underline{\Delta}^c_1 \ar[d]^{w}\\
K_{∙} \ar[r]^{u_i}&L_{∙}&\kern -1.5cm\txt{$i=0,1$.}
}}\leqno({\rm x})
\]
commute. For every semi-simplicial object $M_∙$ of $C^{\wedge}$
(not necessarily a hypercovering), let $F(M_∙)$ be the set of pairs
$(v,w)$, where $v:M_∙\to K_∙$ and
$w:M_∙\times\sheaf{\Delta}_1^c\to L_∙$, such that the diagrams (x)
commute. The functor $M_∙\mapsto F(M_∙)$ is contravariant in $M_∙$,
carries inductive limits to projective limits, and is consequently
represented by a semi-simplicial object $F_{∙}$. The object $F_{∙}$
is clearly the upper-left vertex of a cartesian diagram
\[
\xymatrix@=2cm{
F_{∙}\ar[r]\ar[d]&
\Hhom_{∙}(\sheaf{\Delta}^{c}_{1},L_{∙})\ar[d]^{\pi}\\
K_{∙}\ar[r]_{u_0,u_1}&L_{∙}\times L_{∙},
}
\]
where $\pi$ is defined by the two inclusions of
$\sheaf{\Delta}^c_0$ into $\sheaf{\Delta}^c_1$. To prove the
assertion, it therefore suffices to show that $F_{∙}$ is an object
of $HR_p$ (\resp of $HR$). For this, by
\hyperref[V.lem7.3.4]{7.3.4} 2) and the sorites
\hyperref[V.subsubsec7.3.0]{7.3.0}, it suffices to show that

\begin{enumerate}
\renewcommand{\theenumi}{\alph{enumi}}
\renewcommand{\labelenumi}{\theenumi)}
\item $\Hhom_∙(\sheaf{\Delta}^c_1,L_∙)$\pageoriginale is
semi-representable,

\item the morphism $\pi$ is special of type $p$ (\resp special).
\end{enumerate}

First, it follows immediately that
$\Hhom_∙(\sheaf{\Delta}^c_1,L_∙)$ is a semi-representable
semi-simplicial object. Indeed, its components are computed by
finite projective limits from the $L_n$, and are therefore
semi-representable. Let us now verify b). First one shows that, when
$L_∙$ is special of type $p$, the morphism
$$
\Hhom_∙(\sheaf{\Delta}_1^c,L_∙)\to
\cosk_p\Hhom_∙(\sheaf{\Delta}^c_1,L_∙)
$$
is an isomorphism; this verifies property 1) of
\hyperref[V.def7.3.3]{7.3.3}. Next, the morphism
$$
\pi_0:\Hhom_0(\sheaf{\Delta}^c_1,L_∙)\to L_0\times L_0
$$
is isomorphic to the canonical morphism
$$
L_1\xrightarrow{(d_0,d_1)}L_0\times L_0,
$$
which is covering by hypothesis ($L_∙$ is a hypercovering).

It remains to verify property 2) of \hyperref[V.def7.3.3]{7.3.3},
\ie to verify that, for every $n$, in the diagram
\[
\xymatrix@C=.75cm@R=1.25cm{
\Hhom_{n+1}(\sheaf{\Delta}_1^c,L_∙)
\ar[rr]^{\pi_{n+1}} \ar[dd] \ar[rd]_{\Phi_{n+1}}&&
L_{n+1}\times L_{n+1}\ar[dd]\\
&P_{n+1}\ar[ru]\ar[ld]&\\
(\cosk_n\Hhom_∙(\sheaf{\Delta}_1^c,L_∙))_{n+1}
\ar[rr]^{(\cosk_n\pi)_{n+1}}&&
(\cosk_nL_∙)_{n+1}\times(\cosk_nL_∙)_{n+1}
}
\]
($P_{n+1}$ is the fiber product), the morphism $\Phi_{n+1}$ is
covering. A simple ``adjoint-functor chase'' shows that
\begin{align*}
P_{n+1}&\xrightarrow{\sim}
\Hhom_0\bigl(\sk_{n+1}(\sheaf{\Delta}_{n+1}^c\times
\sheaf{\Delta}_1^c),L_∙\bigr),\\
\Hhom_{n+1}(\sheaf{\Delta}_1^c,L_∙)&\xrightarrow{\sim}
\Hhom_0(\sheaf{\Delta}_{n+1}^c\times
\sheaf{\Delta}_1^c,L_∙),
\end{align*}
and\pageoriginale that the morphism $\Phi_{n+1}$ is isomorphic to
the morphism
$$
\Hhom_0(\sheaf{\Delta}_{n+1}^c\times\sheaf{\Delta}_1^c,L_∙)
\to
\Hhom_0\bigl(\sk_{n+1}(\sheaf{\Delta}_{n+1}^c\times
\sheaf{\Delta}_1^c),L_∙\bigr)
$$
arising from the injection
$$
\sk_{n+1}(\sheaf{\Delta}_{n+1}^c\times\sheaf{\Delta}_1^c)
\hookrightarrow
\sheaf{\Delta}_{n+1}^c\times\sheaf{\Delta}_1^c{\quoi}.
$$
There is a sequence of subobjects of
$\sheaf{\Delta}_{n+1}^c\times\sheaf{\Delta}_1^c$:
$$
\sk_{n+1}(\sheaf{\Delta}_{n+1}^c\times\sheaf{\Delta}_1^c)
=M_0\hookrightarrow M_1\hookrightarrow\cdots\hookrightarrow M_k
=\sheaf{\Delta}_{n+1}^c\times\sheaf{\Delta}_1^c{\quoi},
$$
such that, for every $0\leq i<k$, the morphism
$$
M_i\hookrightarrow M_{i+1}
$$
fits into a cocartesian diagram:
\[
\xymatrix@C=1.5cm@R=1.5cm{
M_i \ar@{^{(}->}[r]&M_{i+1}&\\
\sk_{n+1}\underline{\Delta}{}^c_{n+2}
\ar@{^{(}->}[r]\ar[u]&
\underline{\Delta}{}^c_{n+2}\ar[u]\;.
}
\]
(The nondegenerate simplices of dimension $n+2$ in
$\sheaf{\Delta}_{n+1}\times\sheaf{\Delta}_1$ are added one after
another.) By \hyperref[V.lem7.3.6]{7.3.6}, the morphism
$\Phi_{n+1}$ is a composite of covering morphisms and is therefore
itself covering. This completes the proof of
\hyperref[V.thm7.3.2]{7.3.2}.

\subsection{The isomorphism theorem}\etiquette[7.4]{V.subsec7.4}

\setcounter{subsubsection}{-1}
\subsubsection{}\etiquette[7.4.0]{V.subsubsec7.4.0}

Let $F$ be a presheaf of abelian groups on a site $C$ satisfying
condition \hyperref[V.subsubsec7.3.0]{7.3.0}. For every
hypercovering $K_∙$, put
$$
H^q(K_∙,F)=H^q(\Hom_{C^{\wedge}}(K_∙,F)){\quoi}.
$$
The $H^q(K_∙,F)$ form a $\delta$-functor on the category of abelian
presheaves on $C$, but they are not, in general, the derived
functors of the functor $H^0(K_∙,F)$.

Now\pageoriginale put
\begin{equation*}
\Hvr{q}(C,F)=\varinjlim_{HR_r}H^q(⋅,F)
\text{ }(r\mbox{ may be infinite}){\quoi},\tag{7.4.0.1}
\end{equation*}\etiquette[7.4.0.1]{V.7.4.0.1}
and
\begin{equation*}
\Hvv{q}(C,F)=\varinjlim_{HR}H^q(⋅,F){\quoi}.\tag{7.4.0.2}
\end{equation*}\etiquette[7.4.0.2]{V.7.4.0.2}
The $\Hvr{q}(C,F)$ (\resp $\Hvv{q}(C,F)$) again form a
$\delta$-functor, because the category $\sheaf{HR}^0_r$ (\resp
$\sheaf{HR}^0$) is filtered (\hyperref[V.thm7.3.2]{7.3.2}). Since
the categories $\sheaf{HR}_r$ are not necessarily $𝒰$-categories,
these $\delta$-functors take values in the category of abelian
$𝒱$-groups, for a suitable universe $𝒱$.

Suppose now that the presheaf $F$ is a sheaf. Since the
semi-simplicial sheaf associated with a hypercovering is acyclic
(\hyperref[V.thm7.3.2]{7.3.2}), there is a spectral sequence,
functorial in $F$ and in $K_∙$:
\begin{equation*}
H^p(K_∙,\Hh^q(F))\Rightarrow
H^{p+q}(C^{\sim},F){\quoi}.\tag{7.4.0.3}
\end{equation*}\etiquette[7.4.0.3]{V.7.4.0.3}
Passing to the limit gives the spectral sequences
\begin{equation*}
\begin{split}
\Hvr{p}(C,\Hh^q(F))&\Rightarrow H^{p+q}(C^{\sim},F){\quoi},\\
\Hvv{p}(C,\Hh^q(F))&\Rightarrow H^{p+q}(C^{\sim},F){\quoi}.
\end{split}\tag{7.4.0.4}
\end{equation*}\etiquette[7.4.0.4]{V.7.4.0.4}

\begin{enonce*}{Theorem 7.4.1}\etiquette[7.4.1]{V.thm7.4.1}
Let $C$ be a site satisfying condition
\hyperref[V.subsubsec7.3.0]{7.3.0}, and let $F$ be an abelian sheaf
on $C$.
\begin{enumerate}
\renewcommand{\labelenumi}{\theenumi)}
\item The spectral sequences \eqhyperref[V.7.4.0.4]{7.4.0.4} define
isomorphisms
$$
\Hvr{q}(C,F)\xrightarrow{\sim}H^q(C^{\sim},F),\quad q\leq r+1,
$$
and a monomorphism
$$
\Hvr{r+2}(C,F)\longrightarrow H^{r+2}(C^{\sim},F){\quoi}.
$$

\item There is\pageoriginale an isomorphism of $\delta$-functors:
$$
\Hvi{q}(C,F)\xrightarrow{\sim}\Hvv{q}(C,F)
\xrightarrow{\sim}H^q(C^{\sim},F)\quad(\text{for every }q){\quoi}.
$$

\item Let $G$ be a presheaf of abelian groups, and let $F$ be its
associated sheaf. There are functorial isomorphisms
$$
\Hvr{q}(C,G)\xrightarrow{\sim}H^q(C^{\sim},F),\quad q\leq r-1,
$$
and a monomorphism
$$
\Hvr{r}(C,G)\longrightarrow H^r(C^{\sim},F){\quoi}.
$$
When $G$ is a separated presheaf, the latter morphism is an
isomorphism, and there is a monomorphism
$$
\Hvr{r+1}(C,G)\to H^{r+1}(C^{\sim},F){\quoi}.
$$

\item There are isomorphisms
$$
\Hvi{q}(C,G)\xrightarrow{\sim}\Hvv{q}(C,G)
\xrightarrow{\sim}H^q(C^{\sim},F)\quad(\text{for every }q){\quoi}.
$$
\end{enumerate}
\end{enonce*}

\begin{enonce*}[remark]{Proof}
It suffices to show that, if $N$ is an abelian presheaf whose
associated sheaf is zero, then
$\Hvr{q}(C,N)=0$ for $q\leq r$ (\resp
$\Hvv{q}(C,N)=0$ for every $q$); this follows immediately by using
\hyperref[V.thm3.2]{3.2}.
\end{enonce*}

\begin{enonce*}[remark]{Remark 7.4.2}\etiquette[7.4.2]{V.rem7.4.2}
\begin{enumerate}
\renewcommand{\labelenumi}{\theenumi)}
\item Note that, for sites satisfying condition
\hyperref[V.subsubsec7.3.0]{7.3.0}, hypercoverings of type $0$ are
ordinary coverings, and the functors $\Hvr{q}$ for $r=0$ are none
other than the functors $\Hv{q}$ introduced in
\eqhyperref[V.2.4.5.1]{2.4.5.1}.

\item Let $C$ be a site in which finite projective limits are
representable, and suppose that for every object $X$ of $C$ and
every covering family $(X_i\to X)$, $i\in I$, there exists an
$i_0\in I$ such that $X_{i_0}\to X$ is covering. One can then show
that hypercoverings of type $p$ whose components are representable
are cofinal in $\sheaf{HR}_p$. Likewise, hypercoverings $K_∙$ for
which there exists an integer $p$ such that $K_∙$ is of type $p$,
and whose components are representable, are cofinal
in\pageoriginale $\sheaf{HR}$. (This may be proved by following the
argument of \hyperref[V.cor3.5]{3.5}.) Thus, in the case under
consideration, these hypercoverings with representable components
suffice to compute the cohomology of the sheaves associated with
presheaves.
\end{enumerate}
\end{enonce*}


% END INCLUDED SGA4_English_current_buildable_source_20260906/source/05/drafts/SGA4_V_LINES_3071_3430_English.texfrag


% BEGIN INCLUDED SGA4_English_current_buildable_source_20260906/source/05/drafts/SGA4_V_LINES_3431_3814_English.texfrag
\section{Appendix. Local direct limits (by P. Deligne)}\etiquette[8]{V.sec8}

The editor recommends that the reader should, as a rule, avoid
reading this appendix. It presents a technique which sometimes
makes it possible to extend to topoi without enough points
assertions which the existence of points makes trivial. This
technique yields a sheaf-theoretic variant of D. Lazard's theorem
asserting that flat modules over a ring are direct limits of finite
free modules.

\subsection{Locally filtering categories}\etiquette[8.1]{V.subsec8.1}

\setcounter{subsubsection}{-1}
\subsubsection{}\etiquette[8.1.0]{V.subsubsec8.1.0}
If $\Bb$ is a category and $p:\Aa\to\Bb$ is a category over
$\Bb$, we shall use the following notation:
\begin{itemize}
\item For $U\in\Ob\Bb$, $\Aa_U$ is the fiber category $p^{-1}(U)$.

\item For $f:U\to V$ in $\Bb$ and $\lambda,\mu\in\Ob\Aa$ such
that $p(\lambda)=U$ and $p(\mu)=V$, put
$$
\Hom_f(\lambda,\mu)=p^{-1}(f),\quad\text{where }p:\Hom(\lambda,\mu)
\to\Hom(U,V){\quoi}.
$$
\end{itemize}

\begin{enonce*}{Definition 8.1.1}\etiquette[8.1.1]{V.def8.1.1}
Let $\Ss$ be a site. A \emph{locally cofiltering category (or
locally left filtering category)} over $\Ss$ is a category
$p:\Ll\to\Ss$ over $\Ss$ such that:
\begin{itemize}
\item[($\Ll0$)] For every $f:U\to V$ in $\Ss$ and every
$\mu\in\Ob\Ll_V$, there exists $\lambda\in\Ob\Ll_U$ such that
$\Hom_f(\lambda,\mu)\neq\varnothing$;

\item[($\Ll1$)] For every $U\in\Ob\Ss$ and every finite family
$(\lambda_i)$ of objects of $\Ll_U$, there exist a covering
$f_j:U_j\to U$ of $U$ and objects $\mu_j\in\Ob\Ll_{U_j}$ such that
$\Hom_{f_j}(\mu_j,\lambda_i)\neq\varnothing$ for all $i$ and $j$;

\item[($\Ll2$)] For every $f:U\to V$ in $\Ss$ and every pair of
parallel arrows $(\phi_0,\phi_1):\lambda\rightrightarrows\mu$ over
$f$, there exist a covering $f_j:U_j\to U$ of $U$ and arrows
$\psi_j$ with target $\lambda$ such that $p(\psi_j)=f_j$ and
$\phi_0\psi_j=\phi_1\psi_j$.
\end{itemize}
\end{enonce*}

\setcounter{subsubsection}{1}
\paragraph{}\etiquette[8.1.1.1]{V.par8.1.1.1}
This definition simplifies when $\Ll$ is fibered over $\Ss^{\ast}$:
the axiom
\begin{itemize}
\item[$(\Ll_\ast0)$] is then satisfied, and $(\Ll1)$, $(\Ll2)$ may
be stated as follows:

\item[$(\Ll'1)$] For every $U\in\Ob\Ss$ and every finite family
$\lambda_i$ of objects of $\Ll_U$, locally on $U$ there exists a
$\mu$ mapping to all the $\lambda_i$.

\item[$(\Ll'2)$] For every $U\in\Ob\Ss$, every pair of parallel
arrows $(\phi_0,\phi_1):\lambda\rightrightarrows\mu$ in $\Ll_U$ can,
locally on $U$, be equalized by an arrow with target $\lambda$.
\end{itemize}

\subsubsection{}\etiquette[8.1.2]{V.subsubsec8.1.2}
Let $\Ss$ be a site, let $\Ll$ be a category over $\Ss$, and let
$\mathsf{e}$ be a stack over $\Ss$. Denote by
$\sheaf{\Gamma}\mathsf{e}$ the category of global sections
$\Hhom^{\cart}_{\Ss}(\Ss,\mathsf{e})$ of $\mathsf{e}$. If $\Ss$
has a final object $S$, then $\sheaf{\Gamma}\mathsf{e}$ ``is none
other than'' $\mathsf{e}_S$.

\begin{enonce*}{Definition 8.1.2.1}\etiquette[8.1.2.1]{V.def8.1.2.1}
The category of \emph{local inverse systems of objects of
$\mathsf{e}$ indexed by $\Ll$} is the category
$\Hhom_{\Ss}(\Ll,\mathsf{e})$.
\end{enonce*}

Denote by $c$ (``associated constant local inverse system'') the
composite functor
$$
\sheaf{\Gamma}\mathsf{e}=\Hhom^{\cart}_{\Ss}(\Ss,\mathsf{e})
\hookrightarrow\Hhom_{\Ss}(\Ss,\mathsf{e})
\xrightarrow{u\mapsto u\circ p}\Hhom_{\Ss}(\Ll,\mathsf{e}).
$$
When the dependence on $\Ll$ must be made explicit, we shall write
$c_{\Ll}$ instead.

\begin{enonce*}{Definition 8.1.3}\etiquette[8.1.3]{V.def8.1.3}
The local inverse-limit functor, denoted by
$\varprojlim_{\lambda\in\Ll}$, is the partially defined functor
right adjoint to the functor $c$.
\end{enonce*}

Thus the functor
$$
\varprojlim:\Hhom_{\Ss}(\Ll,\mathsf{e})\to
\sheaf{\Gamma}\mathsf{e}
$$
satisfies
$$
\Hom(X,\varprojlim_{\lambda\in\Ll}X_\lambda)\simeq
\Hom(c(X),(X_\lambda)_{\lambda\in\Ll}){\quoi}.
$$

\begin{enonce*}{Definition 8.1.4}\etiquette[8.1.4]{V.def8.1.4}
An $\Ss$-functor $F:\Ll\to\Mm$ between locally cofiltering
categories over $\Ss$ is said to be cofinal if it satisfies the
following two conditions:
\begin{itemize}
\item[$(C\mathfrak{F}1)$] For every\pageoriginale $U\in\Ob\Ss$ and
$\mu\in\Ob\Mm_U$, there exist a covering $f_j:U_j\to U$ of $U$ and
objects $\lambda_j\in\Ob\Ll_{U_j}$ such that
$\Hom_{f_j}(F(\lambda_j),\mu)\neq\varnothing$.

\item[$(C\mathfrak{F}2)$] For every $f:U\to V$ in $\Ss$ and every
pair of parallel arrows $(\phi_0,\phi_1):F(\lambda)
\rightrightarrows\mu$ over $f$, there exist a covering
$f_j:U_j\to U$ of $U$ and arrows $\psi_j$ with target $\lambda$ in
$\Ll$, such that $p(\psi_j)=f_j$ and
$\phi_0F(\psi_j)=\phi_1F(\psi_j)$.
\end{itemize}
\end{enonce*}

The composite of two cofinal functors is cofinal; equivalences of
locally filtering categories are cofinal (the proof is left to the
reader).

\begin{enonce*}{Proposition 8.1.5}\etiquette[8.1.5]{V.prop8.1.5}
Let $F:\Ll\to\Mm$ be a cofinal functor between locally cofiltering
categories over $\Ss$, let $\mathsf{e}$ be a stack over $\Ss$, and
let $(X_\mu)_{\mu\in\Mm}$ be a local inverse system indexed by
$\Mm$. Then the morphism, functorial in $X\in\sheaf{\Gamma}\mathsf{e}$,
\[
F^{\ast}:\Hom(c_{\Mm}(X),(X_\mu)_{\mu\in\Mm})\longrightarrow
\Hom(c_{\Ll}(X),(X_{F(\lambda)})_{\lambda\in\Ll})
\]
is an isomorphism, so that
$$
F^{\ast}:\varprojlim_{\Mm}X_\mu\longrightarrow
\varprojlim_{\Ll}X_{F(\lambda)}
$$
is an isomorphism, the two sides being simultaneously defined or
undefined.
\end{enonce*}

The proof uses the following lemma, which is left to the reader:

\begin{enonce*}{Lemma 8.1.6}\etiquette[8.1.6]{V.lem8.1.6}
If $F:\Ll\to\Mm$ is cofinal, then for every $f:U\to V$ in $\Ss$
and arrows over $f$, $\phi_i:F(\lambda_i)\to\mu$ $(i=1,2)$, there
exist a covering $f_j:U_j\to U$ of $U$, objects
$\lambda'_j\in\Ob\Ll_{U_j}$, and commutative diagrams
\[
\xymatrix@=1.25cm{
              & F(\lambda_1)   \ar[rd]^{\phi_1}&\\
F(\lambda'_j) \ar@{--}[ru] \ar@{--}[rd]&                & \mu\\
              & F(\lambda_2) \ar[ru]_{\phi_2}  &
}
\]
$$
U_j\xrightarrow{f_j}U\to V
$$
whose projection in $\Ss$ is $(f_j,f)$.
\end{enonce*}

For every $U\in\Ob\Ss$ and $\mu\in\Mm_U$, condition
$(C\mathfrak{F}1)$ gives a covering $f_j:U_j\to U$ of $U$ and
arrows $\phi_j:F(\lambda_j)\to\mu$ over the $f_j$. If
$$
\psi=(\psi_\lambda)\in\Hom(c_{\Ll}X,
(X_{F(\lambda)})_{\lambda\in\Ll}):\quad
\psi_\lambda:X|p(\lambda)\to X_{F(\lambda)}
$$
is the image of
$$
\phi=(\phi_\mu)\in\Hom(c_{\Mm}X,(X_\mu)_{\mu\in\Mm}):\quad
\phi_\mu:X|p(\mu)\to X_\mu,
$$
then the diagrams
\[
\xymatrix@=1.5cm{
                               & X_{F(\lambda_j)}
                               \ar[ld]_{\psi_{\lambda_j}} \ar[rd]^{\phi_j}&        \\
X|U_j \ar[rr]_{\phi_\mu|U_j} &                  & X_\mu|U_j
}
\]
commute. Thus the $\phi_\mu$ are locally determined by $\psi$; since
$\mathsf{e}$ is a stack, it follows that $F^{\ast}$ is injective.
To prove that $F^{\ast}$ is surjective, start with $\psi$ and
construct a $\phi$ such that $\psi=F^{\ast}\phi$. The preceding
construction gives the composite arrows
$$
\phi_{\mu,j}=\phi_j\circ\psi_{\lambda_j}:X|U_j\to X_\mu|U_j{\quoi}.
$$
Let us verify that these arrows glue, so that there exists
$\phi_\mu:X|U\to X_\mu$ such that
$\phi_\mu|U_j=\phi_{\mu,j}$. Consider a commutative diagram in
\[
\xymatrix@=1cm{
                 & U_{ij}\ar[ld]_{g_i} \ar[rd]^{g_j}     &            \\
U_i  \ar[rd]_{f_i}            &            & U_j \ar[ld]^{f_j}\\
                 & U     &
}
\]
To\pageoriginale verify that
$\phi_{\mu,j}|U_{ij}=\phi_{\mu,i}|U_{ij}$, it is enough to do so
locally on $U_{ij}$; by \hyperref[V.lem8.1.6]{8.1.6} and $(\Ll0)$ this reduces us to the
case where there is a commutative diagram
\[
\xymatrix@=1.5cm{
             & F(\lambda_{ij}) \ar[ld]_{F(a_{ij})} \ar[rd]^{F(a_{ji})} &\\
F(\lambda_i) \ar[rd]_{\phi_i}&                 & F(\lambda_j) \ar[ld]^{\phi_j}\\
             & \mu             &
}
\]
over the preceding one. The diagram
\[
\xymatrix@R=1.5cm@C=2cm{
  & X_{F(\lambda_{ij})} \ar[ld] \ar[rd] \ar@/_1.68pc/[dd]&\\
X_{F(\lambda_i)}\ar[rd] & X|U_{ij}\ar[u]\ar[l]\ar[r] & X_{F(\lambda_j)} \ar[ld]\\
                 & X_\mu &
}
\]
is then commutative, and there exists a family $(\phi_\mu)$ such
that, for every arrow $\phi:F(\lambda)\to\mu$ over
$f:V\to U$, the diagram
\[
\xymatrix@=1.5cm{
                              & X_{F(\lambda)} \ar[rd]^{\phi}      &    \\
X|V \ar[rr]_{\phi_\mu|V} \ar[ru]^{\psi_\lambda} &
                              & X_\mu|V
}
\]
commutes. If $\sigma:\mu\to\mu'$ is an arrow of $\Mm$, then
$(\sigma|V)(\phi_\mu|V)=(\sigma|V)\phi\cdot\psi_\lambda
=(\sigma\phi|V)\cdot\psi_\lambda=\phi_{\mu'}|V$; hence
$\phi=(\phi_\mu)$ is a morphism of functors such that
$\psi=F^{\ast}\phi$, as required.

If\pageoriginale $E$ is an ordered set, we shall also denote by $E$
the category whose set of objects is $E$ and in which $\Hom(i,j)$
is a singleton or empty according as $i\leq j$ or not.

\begin{enonce*}{Proposition 8.1.7}\etiquette[8.1.7]{V.prop8.1.7}
Let $\Ll$ be a locally cofiltering category over a site $\Ss$.
There exist an ordered set $E$ and a functor $F$ from $E$ to $\Ll$
such that:
\begin{enumerate}
\renewcommand{\theenumi}{\roman{enumi}}
\renewcommand{\labelenumi}{(\theenumi)}
\item $E$, regarded as a category over $\Ss$ by means of
$pF:E\to\Ss$, is locally cofiltering.
\item The functor $F$ is cofinal.
\end{enumerate}
\end{enonce*}

Let $\Ll'$ be the locally cofiltering category which is the product
of $\Ll$ with the category defined by the ordered set $\ZZ$. The
projection from $\Ll'$ to $\Ll$ is cofinal, so it suffices to prove
\hyperref[V.prop8.1.7]{8.1.7} for $\Ll'$.

If $X$ is a finite set of arrows in $\Ll'$, denote by $X^0$ the set
of endpoints of the arrows in $X$. Let $E$ be the set of nonempty
finite subsets $X$ of $\mathrm{Fl}(\Ll')$ for which there exists
$\lambda_X\in\Ob\Ll'$ satisfying:
\begin{enumerate}
\renewcommand{\theenumi}{\alph{enumi}}
\renewcommand{\labelenumi}{\theenumi)}
\item no arrow of $X$, except $1_{\lambda_X}$, has target
$\lambda_X$;

\item if $\mu\in X^0$, then $1_\mu\in X$ and
$\Hom(\lambda_X,\mu)\cap X\neq\varnothing$;

\item every diagram of the form
\[
\xymatrix@=1.75cm{
\lambda_X \ar[r] \ar[rd]& \mu \ar[d]\\
            & \mu'
}
\]
consisting of arrows of $X$ is commutative.
\end{enumerate}

Order $E$ by the relation opposite to inclusion. If $X\in E$ and
$\mu\in X^0$, there is a unique arrow from $\lambda_X$ to $\mu$
which belongs to $X$. If $X,Y\in E$ and $X\supset Y$, let
$\lambda_{X,Y}$ be the unique arrow in $X$ from $\lambda_X$ to
$\lambda_Y$. The assignments $X\mapsto\lambda_X$ and
$(X\supset Y)\mapsto\lambda_{X,Y}$ form a functor from $E$ to
$\Ll'$.

\begin{enonce*}{Lemma 8.1.8}\etiquette[8.1.8]{V.lem8.1.8}
The category\pageoriginale $E$ is locally cofiltering over $\Ss$,
and the functor $\lambda:E\to\Ll'$ is cofinal.
\end{enonce*}

We must verify the following axioms successively:
\begin{itemize}
\item[(i) \emph{axiom} ($\Ll0$)] Let $f:V\to U$ and $X\in E_U$.
There exist $\lambda\in\Ob\Ll'_V$ and
$\phi\in\Hom_f(\lambda,\lambda_X)$. Moreover, $\lambda$ may be
chosen so that $\lambda\notin X^0$ (this is where
$\Ll'=\Ll\times\ZZ$ is used). Let
$X'=X\cup\{1_\lambda\}\cup Y$, where $Y$ is the set of composite
arrows
$\lambda\xrightarrow{\phi}\lambda_X\xrightarrow{\psi}\mu$
$(\psi\in X)$. Then $X'\in E$, $\lambda_{X'}=\lambda$, and
$\Hom(X',X)\neq\varnothing$.

\item[(ii) \emph{axiom} ($\Ll1$)] Let $(X_i)$ be a finite family of
objects of $E_U$. There exist a covering $f_j:U_j\to U$, objects
$\lambda_j\in\Ob\Ll'_{U_j}$, and arrows
$\phi_{ji}\in\Hom_{f_j}(\lambda_j,\lambda_{X_i})$. If
$x\in X_a^0\cap X_b^0$, then for each $j$ there exist a covering
$g_{jk}:U_{jk}\to U_j$ and arrows
$\psi_{jk}:\lambda_{jk}\to\lambda_j$ such that
$p(\psi_{jk})=g_{jk}$ and which equalize the pair of arrows
$\lambda_j\rightrightarrows x$, one belonging to $X_a$ and the
other to $X_b$. We may arrange that the $\lambda_{jk}$ belong to
none of the $X_i^0$. Replace the $(f_j,\lambda_j,\phi_{ji})$ by the
$(f_jg_{jk},\lambda_{jk},\phi_{ji}\psi_{jk})$, and repeat this
construction for every $x$ in an intersection of two of the $X_i$.
We obtain a new system $(f_j,\lambda_j,\phi_{ji})$; put
$$
X_j=\bigcup_iX_i\cup\{1_{\lambda_j}\}\cup Y_j,
$$
where $Y_j$ is the set of composite arrows
$$
\lambda_j\xrightarrow{\phi_{ji}}\lambda_{X_i}
\xrightarrow{\psi}x\quad(\psi\in X_i){\quoi}.
$$
Then $X_j\in E_{U_j}$, and these $X_j$ satisfy $(\Ll1)$.

\item[(iii) \emph{axiom} ($\Ll2$)] Trivial, since there are no
nontrivial parallel arrows!

\item[(iv) \emph{axiom} $(C\mathfrak{F}1)$] Trivial, since the
functor $\lambda$ is surjective.

\item[(v) \emph{axiom} $(C\mathfrak{F}2)$] Let $f:U\to V$ and let
$(\phi_0,\phi_1):\lambda_X\rightrightarrows\lambda$ be a pair of
parallel arrows over $f$. There exist a covering $f_j:U_j\to U$ of
$U$ and arrows $\psi_j:\lambda_j\to\lambda_X$ such that
$p(\psi_j)=f_j$, $\phi_0\psi_j=\phi_1\psi_j$, and
$\lambda_j\notin X^0$. Let
$X_j=X\cup\{1_{\lambda_j}\}\cup Y_j$, where $Y_j$ is the set of
composite arrows
$$
\lambda_j\xrightarrow[]{\psi_j}\lambda_X\xrightarrow{\psi}\mu
\quad(\psi\in X){\quoi}.
$$
Then\pageoriginale $X_j\in E_{U_j}$, and
$F(\lambda_{X_j,X})=\psi_j$ equalizes $(\phi_0,\phi_1)$.
\end{itemize}

\setcounter{subsubsection}{9}
\setcounter{paragraph}{-1}
\paragraph{}\etiquette[8.1.9.0]{V.par8.1.9.0}

Suppose that $\Ss$ is (the site of nonempty open subsets of) the
one-point topological space. A locally cofiltering category over
$\Ss$ is then simply a left filtering category (\ie one for which
$\Ss^0$ is filtering in the sense of I \hyperref[I.def2.7]{2.7}), and cofinal functors
correspond to cofinal functors between left filtering categories
(I 8). Proposition \hyperref[V.prop8.1.7]{8.1.7} may then be restated
as follows.

\begin{enonce*}{Proposition 8.1.9}\etiquette[8.1.9]{V.prop8.1.9}
For every cofiltering category $\Ll$, there exist a cofiltering
ordered set $E$ and a cofinal functor from $E$ to $\Ll$.
\end{enonce*}

In this special case, the proof of \hyperref[V.prop8.1.7]{8.1.7}
reduces essentially to the proof of \hyperref[V.prop8.1.9]{8.1.9}
given in I 8.

% END INCLUDED SGA4_English_current_buildable_source_20260906/source/05/drafts/SGA4_V_LINES_3431_3814_English.texfrag


% BEGIN INCLUDED SGA4_English_current_buildable_source_20260906/source/05/drafts/SGA4_V_LINES_3815_4065_English.texfrag
\subsubsection{}\etiquette[8.1.10]{V.subsubsec8.1.10}
Let $q:\Tt\to\Ss$ be a morphism of sites, and let
$p:\Ll\to\Ss$ be a locally cofiltered category over $\Ss$. Define
the category $\Ll^{\ast}$ as follows:
\begin{enumerate}
\renewcommand{\theenumi}{\roman{enumi}}
\renewcommand{\labelenumi}{(\theenumi)}
\item An object of $\Ll^{\ast}$ is a quadruple $(V,U,\phi,\lambda)$
such that $V\in\Ob\Tt$, $U\in\Ob\Ss$,
$\phi\in\Hom(V,q^{\ast}U)$, and $\lambda\in\Ob\Ll_U$.

\item An arrow
$f:(V,U,\phi,\lambda)\to(V',U',\phi',\lambda')$ is a triple
$(f_1,f_2,f_3)$, where $f_1\in\Hom(V,V')$,
$f_2\in\Hom(U,U')$, and $f_3\in\Hom(\lambda,\lambda')$, such that
$q^{\ast}f_2\circ\phi=\phi'\circ f_1$ and $p(f_3)=f_2$.

The functor $f\mapsto f_1$ makes $\Ll^{\ast}$ a category over
$\Tt$.
\end{enumerate}

\begin{enonce*}{Lemma 8.1.11}\etiquette[8.1.11]{V.lem8.1.11}
The category $\Ll^{\ast}$ is locally cofiltered over $\Tt$.
\end{enonce*}

Axiom $(\Ll0)$ is trivial. Let us verify $(\Ll1)$.

Let $L_i=(V,U_i,\phi_i,\lambda_i)$ be a finite family. The $\phi_i$
define a morphism from $V$ to the sheaf
$q^{\ast}a\prod U_i=a\prod q^{\ast}U_i$. Consequently, there are
commutative diagrams
\[
\xymatrix@C=2.5cm@R=1.5cm{
V_j \ar[r]^{f_j}\ar[d]^{\psi_j}&V\ar[d]^{\phi_i}\\
q^{\ast}U_j\ar[r]_{q^{\ast}(p_{ij})}&q^{\ast}U_i
}
\]
such\pageoriginale that the $f_j$ cover $V$.

Choose $\lambda_{ji}$ such that
$\Hom_{p_{ij}}(\lambda_{ji},\lambda_i)\neq\varnothing$. By
$(\Ll1)$, after refining $U_j$ (and $V_j$), we may take
$\lambda_{ji}=\lambda_j$ independently of $i$, and the
$L_j=(V_j,U_j,\psi_j,\lambda_j)$ satisfy $(\Ll1)$.

Let us verify $(\Ll2)$. Let
$(f_1,f_2):L^1\rightrightarrows L^2$, with
$L^i=(V^i,U^i,\phi^i,\lambda^i)$ and
$f_i=(f,g_i,h_i)$. The $\phi^i$ define a morphism from $V^1$ to the
sheaf
$q^{\ast}a\Ker(U^1\rightrightarrows U^2)
=a\Ker(q^{\ast}U^1\rightrightarrows q^{\ast}U^2)$, so there are
commutative diagrams
\[
\xymatrix@C=2cm@R=1.5cm{
V_j \ar[r]^{f_j}\ar@{-}[d]^{\psi_j}&
V^1\ar[r]\ar@{-}[d]^{\phi^1}&V^2\ar@{-}[d]^{\phi^2}\\
q^{\ast}U_j\ar[r]_{q^{\ast}p_j}&q^{\ast}U^1
\ar@<2pt>[r]^{q^{\ast}g_1}\ar@<-2pt>[r]_{q^{\ast}g_2}&
q^{\ast}U^2&\kern-1cm(g_1p_j=g_2p_j),
}
\]
such that the $f_j$ cover $V^1$. Let
$\chi_j:\lambda_j\to\lambda^1$ satisfy $p(\chi_j)=p_j$. Applying
$(\Ll2)$ to $(h_1\chi_j,h_2\chi_j)$, and refining $U_j$ (and
$V_j$) if necessary, we may arrange that
$h_1\chi_j=h_2\chi_j$; the
$L_j=(V_j,U_j,\psi_j,\lambda_j)$ then satisfy $(\Ll2)$.

\begin{enonce*}{Proposition 8.1.12}\etiquette[8.1.12]{V.prop8.1.12}
Let $t^{\ast}:\Ss^{\ast}\to\Ss$ be a locally cofiltered category
over $\Ss$, and endow $\Ss^{\ast}$ with the topology induced from
that of $\Ss$ by $t^{\ast}$. Then $t^{\ast}$ is an equivalence of
sites $t:\Ss\to\Ss^{\ast}$.
\end{enonce*}

By construction, the functor
$t^\ast:\Ff\mapsto\Ff\circ t^{\ast}$ carries sheaves on $\Ss$ to
sheaves on $\Ss^{\ast}$. By $(\Ll0)$ and $(\Ll1)$ applied to the
empty family, every ``sufficiently small'' object of $\Ss$ belongs
to the image of $t^{\ast}$ (that is, the image of $t^{\ast}$ is a
\emph{covering sieve} in $\Ss$). Hence, by the ``comparison lemma''
(III 4), the inclusion into $\Ss$ of the full subcategory defined by
$t^{\ast}(\Ob\Ss^{\ast})$ is an equivalence of sites. We may
therefore suppose that $t^{\ast}$ is surjective on objects.

Let $\Ff$ be a sheaf on $\Ss^{\ast}$.

\begin{enumerate}
\renewcommand{\theenumi}{\roman{enumi}}
\renewcommand{\labelenumi}{(\theenumi)}
\item Let $(f_j:U_j\to U)$ be a covering in $\Ss$, let
$\lambda^1,\lambda^2\in\Ob\Ss^{\ast}_U$, and let
$f_j^i\in\Hom_{f_j}(\mu_j,\lambda^i)$ for $i=1,2$. Then
$\Ff(\lambda^1)$ and $\Ff(\lambda^2)$ have the same image in
$\displaystyle\mathop{\prod}_j\Ff(\mu_j)$.

Indeed, let\pageoriginale $x=(x_j)$ belong to the image of
$\Ff(\lambda^1)$. For $x$ to belong to the image of
$\Ff(\lambda^2)$, it suffices that, for every commutative diagram
\[
\xymatrix@C=2.5cm@R=1.5cm{
\mu\ar[r]\ar[d]&\mu_j\ar[d]\\
\mu_k\ar[r]&\lambda^2,
}
\]
$x_j$ and $x_k$ have the same image in $\Ff(\mu)$. Since $\Ff$ is
a sheaf, it suffices to verify this after replacing $\mu$ by the
objects of one of its coverings. By $(\Ll2)$, we may then suppose
that the diagrams analogous to the preceding one, with $\lambda^2$
replaced by $\lambda^1$, also commute, in which case the assertion is
clear.

\item Let $f:U\to V$ be an arrow of $\Ss$, let
$\lambda\in\Ob\Ss^{\ast}_U$, and let
$\mu\in\Ob\Ss^{\ast}_V$. There then exist a covering
$f_j:U_j\to U$, an object $\mu'\in\Ob\Ss^{\ast}_U$, and arrows in
the diagram
\[
\xymatrix@C=2.5cm@R=1cm{
&\lambda&\\
\lambda_j\ar[ru]^{\phi_j}\ar[rd]_{\psi_j}&&\mu\\
&\mu'\ar[ru]_{\psi}&
}
\]
lying over
$$
U_j\xrightarrow{f_j}U\xrightarrow{f}V{\quoi}.
$$
By (i), there is a unique arrow $\overline f$ making the following
diagram commute:
\[
\xymatrix@C=2.5cm@R=1.5cm{
\Ff(\mu)\ar[r]^{\overline f}\ar[d]_{\psi}&
\Ff(\lambda)\ar[d]^{\phi_j}\\
\Ff(\mu')\ar[r]_{\psi_j}&\Ff(\lambda_j).
}
\]
The arrow $\overline f$ does not change if $(f_j)$ is replaced by a
finer covering $f_k:U_k\to U_{j(k)}$, if $\lambda_j$ is replaced by
$\lambda_k$ equipped with
$\phi_k\in\Hom_{f_k}(\lambda_k,\lambda_{j(k)})$, and if
$\phi_j$ (respectively $\psi_j$) is replaced by
$\phi_{j(k)}\circ\phi_k$ (respectively
$\psi_{j(k)}\circ\phi_k$).\pageoriginale

The reader may deduce that $\overline f$ depends only on $f$.

\item Let us prove that $\overline f$ is functorial in $f$. Thus let
$$
W\xrightarrow{h}V\xrightarrow{g}U
$$
be arrows in $\Ss$, and let $\lambda^0$, $\mu^0$, and $\nu^0$
belong respectively to $\Ob\Ss^{\ast}_U$, $\Ob\Ss^{\ast}_V$, and
$\Ob\Ss^{\ast}_W$.
There exists a commutative diagram
\[
\xymatrix@C=.5cm@R=.5cm{
\nu^{0}& &&&&\mu^{0}&&\lambda^{0}\\
&&\nu^{1}\ar[rrru]&&&&\mu^{1}\ar[ru]&\\
&&&{\nu^{2}}\ar[rrru]&&&&\\
&&&&&&&\\
&\nu_{jk}\ar[luuuu]^(.7){e^{0}_{jk}}
\ar[uuur]^(.7){e^{1}_{jk}}\ar[rruu]_{e^{2}_{jk}}\ar[rrrr]&&&&
\mu_j\ar[uuur]_{f^{1}_j}\ar[uuuu]^(.3){f^{0}_j}&&
}
\]
such that

\begin{enumerate}
\renewcommand{\theenumii}{\alph{enumii}}
\renewcommand{\labelenumii}{\theenumii)}
\item $t^{\ast}(\nu^i)$, $t^{\ast}(\mu^i)$,
$t^{\ast}(e^i_{jk})$, and $t^{\ast}(f^i_j)$ do not depend on $i$;

\item for fixed $i$, the $e^i_{jk}$ cover $W$, and the $f^i_j$
cover $V$.
\end{enumerate}

First construct $\mu^1$, $\nu^1$, and $\nu^2$ by $(\Ll0)$, and
the $\mu_j$ as in (ii). Let $(W'_{jk})$ be a covering of $W$ mapping
to the covering $V_j=t^{\ast}(\mu_j)$ as required by the diagram.
After refining $(W'_{jk})$ and applying $(\Ll0)$ and $(\Ll1)$, one
obtains a diagram of the required type, satisfying a) and b), but
not necessarily commutative. Refining once more and applying
$(\Ll2)$ makes it commutative.

It remains only to contemplate the following commutative diagram:
\[
\xymatrix@C=1.5cm@R=.25cm{
\Ff(\lambda^{0})\ar[r]^{\bar g}&\Ff(\mu^0)\ar[rd]\ar[rr]^{\bar h}&&
\Ff(\nu^{0})\ar[dddd]\\
&&\Ff(\nu^1)\ar[rddd]&\\
\Ff(\mu^{1})\ar[r]&\prod\Ff(\mu_j)\ar[rrdd]&&&\\
&&&\\
\Ff(\nu^2)\ar[rrr]&&&\prod\Ff(\nu_{jk})
}
\]

\item It\pageoriginale is clear that $\overline f$ is functorial
in $\Ff$. Thus every sheaf $\Ff$ on $\Ss^{\ast}$ is the direct
image under $t$ of a uniquely determined presheaf $\Ff_0$ on
$\Ss$, and every morphism of sheaves $f:\Ff\to\mathfrak g$ is the
image of a unique morphism of presheaves
$f_0:\Ff_0\to\mathfrak g_0$.
\end{enumerate}

It already follows that $t_{\ast}$ is fully faithful; it remains to
prove that $\Ff_0$ is a sheaf. If $(f_j:U_j\to U)$ is a covering,
then this covering is the image of a covering in $\Ss^{\ast}$, so
$\Ff_0(U)$ injects into $\prod_j\Ff_0(U_j)$. If elements
$x_j\in\Ff_0(U_j)$ glue, then a fortiori they glue in $\Ss^{\ast}$,
and hence come from an element of $\Ff_0(U)$. This completes the
proof of \hyperref[V.prop8.1.12]{8.1.12}.

\setcounter{subsubsection}{13}
\setcounter{paragraph}{-1}
\paragraph{}\etiquette[8.1.13.0]{V.par8.1.13.0}

Let $q:\Tt\to\Ss$ and $p:\Ll\to\Ss$ be as in
(\hyperref[V.subsubsec8.1.10]{8.1.10}), let $\Ll^{\ast}$ be the
category defined there, and let $\mathsf e$ be a stack on $\Tt$.
The stack $q_{\ast}\mathsf e$ on $\Ss$ is the stack defined by
$(q_{\ast}\mathsf e)_U=\mathsf e_{q^{\ast}U}$.

If $(X_{\lambda})_{\lambda\in\Ll}$ is a local projective system
indexed by $\Ll$ and taking values in $q_{\ast}\mathsf e$, it
defines a local projective system
$(X_L)_{L\in\Ll^{\ast}}$ taking values in $\mathsf e$ by the formula
$$
X_L=\phi^{\ast}X_{\lambda}quad\text{for }
L=(V,U,\phi,\lambda){\quoi}.
$$

The reader may verify the following:

\begin{enonce*}{Proposition 8.1.13.1}\etiquette[8.1.13.1]{V.prop8.1.13.1}
With the preceding notation, one has
$$
\varprojlim_{\lambda\in\Ll}X_{\lambda}
=\varprojlim_{L\in\Ll^{\ast}}X_L{\quoi},
$$
the two sides being simultaneously defined or undefined.
\end{enonce*}

Let $\Ss$ be a site, and let $p^{\ast}:\Ss^{\ast}\to\Ss$ be a
locally filtered category over $\Ss$. If $\mathsf e$ is a stack on
$\Ss$, then $p_{\ast}\mathsf e$ is a stack on $\Ss^{\ast}$
(\hyperref[V.par8.1.13.0]{8.1.13.0}). Moreover, the identity functor
$\Ss^{\ast}\to\Ss^{\ast}$ is a locally filtered category over
$\Ss^{\ast}$, endowed with the induced topology, and every local
inductive system $(X_{\lambda})$ indexed by $\Ss^{\ast}$ defines a
local inductive system over $\Ss^{\ast}$ indexed by $\Ss^{\ast}$.

One verifies easily:

\begin{enonce*}{Proposition 8.1.14}\etiquette[8.1.14]{V.prop8.1.14}
With\pageoriginale the preceding notation, one has
$$
\varprojlim_{\lambda\in\Ss^{\ast}/\Ss}X_{\lambda}
=\varprojlim_{\lambda\in\Ss^{\ast}/\Ss^{\ast}}X_{\lambda}{\quoi},
$$
the two sides being simultaneously defined or undefined.
\end{enonce*}

% END INCLUDED SGA4_English_current_buildable_source_20260906/source/05/drafts/SGA4_V_LINES_3815_4065_English.texfrag


% BEGIN INCLUDED SGA4_English_current_buildable_source_20260906/source/05/drafts/SGA4_V_LINES_4066_4380_English.texfrag
\subsection{Local inductive limits in categories of
sheaves}\etiquette[8.2]{V.subsec8.2}

\setcounter{subsubsection}{-1}
\subsubsection{}\etiquette[8.2.0]{V.subsubsec8.2.0}

Let $p:\Tt\to\Ss$ be a morphism of sites. Let $\mathsf{e}$ be the
following stack on $\Ss$:
\begin{enumerate}
\renewcommand{\theenumi}{\roman{enumi}}
\renewcommand{\labelenumi}{(\theenumi)}
\item An object of $\mathsf{e}$ is a pair $(U,\Ff)$ consisting of an
object $U$ of $\Ss$ and a sheaf $\Ff$ on $p^{\ast}U$.

\item An arrow $f:(U,\Ff)\to(V,\Gg)$ is a pair consisting of an arrow
$f_0:U\to V$ and a morphism of sheaves on $p^{\ast}U$:
$$
f_1:p^{\ast}(f_0)^{\ast}\Gg\to\Ff.
$$

\item The functor from $\mathsf{e}$ to $\Ss$ is given by
$f\longmapsto f_0$.
\end{enumerate}

Care must be taken that $\mathsf{e}_U$ is the \emph{opposite category}
of the category of sheaves on $p^{\ast}U$.

Let $\Ll$ be a locally cofiltered category (on the left) over $\Ss$.

\begin{enonce*}{Definition 8.2.1}\etiquette[8.2.1]{V.def8.2.1}
With the preceding notation, the category of local inductive systems,
indexed by $\Ll$, of sheaves on $\Tt$ is the category
$\Hom_{\Ss}(\Ll,\mathsf{e})^{\circ}$.
\end{enonce*}

The ``projective limit'' functor of \hyperref[V.def8.1.3]{8.1.3} is
called here the \emph{``local inductive limit'' functor}. It is a
functor from $\Hom_{\Ss}(\Ll,\mathsf{e})^{\circ}$ to
$(\Gamma\mathsf{e})^{\circ}=(\Tt)^{\sim}$.

\begin{enonce*}{Theorem 8.2.2}\etiquette[8.2.2]{V.thm8.2.2}
The ``local inductive limit'' functor from the category of inductive
systems of sheaves on $\Tt$, indexed by $\Ll$, to the category of
sheaves on $\Tt$ is defined everywhere, commutes with finite
projective limits, and commutes with arbitrary inductive limits.
\end{enonce*}

The\pageoriginale \hyperref[V.prop8.1.13.1]{propositions 8.1.13},
\hyperref[V.prop8.1.14]{8.1.14}, and \hyperref[V.prop8.1.12]{8.1.12}
reduce the assertion to the case $\Ss=\Tt=\Ll$. A local inductive
system then consists of a sheaf $\Ff_U$ on $\Ss/U$ for every
$U\in\Ob\Ss$, together, for every arrow $f:V\to U$ in $\Ss$, with an
arrow $f^{\ast}\Ff_U\to\Ff_V$, functorial in $f$.

For such inductive systems, the $\varinjlim$ are computed as follows:

\begin{enonce*}{Lemma 8.2.3}\etiquette[8.2.3]{V.lem8.2.3}
Let $U\longmapsto\Ff_U$ be a section of the category of presheaves on
the objects of $\Ss$. Then
$$
a(U\longmapsto\Gamma(U,\Ff_U))\simeq
\varinjlim_U a\Ff_U{\quoi}.
$$
\end{enonce*}

By definition, the right-hand side represents the functor which
assigns to every sheaf $\Ff$ on $\Ss$ the set of coherent systems of
arrows $\phi_U:\Ff_U\to\Ff|U$. In turn, an arrow $\phi_U$ is a
coherent system of arrows $\phi_g$, one for every $g:V\to U$, where
$\phi_g:\Ff_U(V)\to\Ff(V)$. Apply \hyperref[V.lem8.1.11]{8.1.11} to
the identity functor of $\Ss$, thereby obtaining $\Ss^{\ast}$ locally
filtered over $\Ss$. We have seen that
$$
\varinjlim_{g\in\Ss^{\ast}}\Ff_U(V_g)^{\sim}
\simeq\varinjlim_U a\Ff_U{\quoi},
$$
where $V_g$ is the source of $g$ and ${}^{\sim}$ denotes the
``generated constant sheaf'' functor. The functor
$U\longmapsto1_U$ from $\Ss$ to $\Ss^{\ast}$ is cofinal, and hence,
again by \hyperref[V.prop8.1.5]{8.1.5},
$$
\varinjlim_U\Ff_U(U)^{\sim}\simeq
\varinjlim_U a\Ff_U{\quoi}.
$$
Returning to the definitions, one finally finds
$$
\varinjlim_U\Ff_U(U)^{\sim}
=a(U\longmapsto\Gamma(U,\Ff_U)){\quoi}.
$$

This shows that, in the case to which we have reduced, the functor
$\varinjlim$ is defined everywhere and is computed as the composite
of the ``sheaf generated by a presheaf'' functor and the functor
$(\Ff_U)\longmapsto(U\longmapsto\Gamma(U,\Ff_U))$. Both of these
functors commute with finite projective limits, and the functor
$\varinjlim$ also commutes with arbitrary inductive limits by its
definition as an adjoint functor.

Lemma\pageoriginale \hyperref[V.lem8.2.3]{8.2.3} provides a systematic
method for proving properties of the local-inductive-limit functor
from the analogous properties of the ``sheaf associated with a
presheaf'' functor. One obtains in this way:

\begin{enonce*}{Proposition 8.2.4}\etiquette[8.2.4]{V.prop8.2.4}
Let $\Tt$, $\Ss$, and $\Ll$ be as above, let $\Aa_\lambda$ be a local
inductive system of sheaves of rings on $\Tt$, and let $\Ll_\lambda$
(\resp $\Mm_\lambda$) be a local inductive system of sheaves of right
(\resp left) modules over $\Aa_\lambda$, all three indexed by $\Ll$.

Then
$$
\varinjlim(\Ll_\lambda\otimes_{\Aa_\lambda}\Mm_\lambda)
\simeq\varinjlim\Ll_\lambda
\mathop{\otimes}_{\varinjlim\Aa_\lambda}
\varinjlim\Mm_\lambda{\quoi}.
$$
\end{enonce*}

(Reduce to the case $\Tt=\Ss=\Ll$, and note that both sides coincide
with the sheaf generated by the presheaf
$$
\Ll_U(U)\mathop{\otimes}_{\Aa_U(U)}\Mm_U(U){\quoi}.
$$
)

\begin{enonce*}{Proposition 8.2.5}\etiquette[8.2.5]{V.prop8.2.5}
Let $\Tt_1\xrightarrow{q}\Tt_2\to\Ss$ be two morphisms of sites, and
let $(F_\lambda)$ be a local inductive system of sheaves on $\Tt_2$
indexed by a locally cofiltered category $\Ll$ over $\Ss$. Then
$$
q^{\ast}\varinjlim_{\lambda\in\Ll}F_\lambda\simeq
\varinjlim_{\lambda\in\Ll}q^{\ast}F_\lambda{\quoi}.
$$
\end{enonce*}

This follows immediately from the fact that $q^{\ast}$ has a right
adjoint $q_\ast$.

Let $\Ss$ and $\Tt$ be two sites in which finite $\varprojlim$ are
representable, let $p:\Tt\to\Ss$ be a morphism of sites such that
$p^{\ast}:\Ss\to\Tt$ commutes with finite $\varprojlim$, and let
$\Ll$ be the locally cofiltered category over $\Tt$ whose objects are
triples $(V,U,\phi)$, where $V\in\Ob\Tt$, $U\in\Ob\Ss$, and
$\phi\in\Hom(V,p^{\ast}U)$. It is obtained by applying
\hyperref[V.lem8.1.11]{8.1.11} to $p:\Tt\to\Ss$.

\begin{enonce*}{Proposition 8.2.6}\etiquette[8.2.6]{V.prop8.2.6}
For every sheaf $\Ff$ on $\Tt$, one has
$$
\Ff=\varinjlim_{\lambda\in\Ll}\phi^{\ast}\phi_\ast\Ff,
$$
where, by abuse of notation, $\phi$ denotes the morphism induced by
$p$ from $\Tt/V$ to $\Ss/U$.
\end{enonce*}

Let\pageoriginale $\phi'$ be the inverse-image functor in the sense of
presheaves. From \hyperref[V.lem8.2.3]{8.2.3} one obtains
$$
\varinjlim_{\lambda\in\Ll}{\phi'}^{\ast}\phi_\ast\Ff
=\varinjlim_{\lambda\in\Ll}a{\phi'}^{\ast}\phi_\ast\Ff
=\varinjlim_{\lambda\in\Ll}\Ff(V)^{\sim}
=\varinjlim_V\Ff(V)^{\sim}=\Ff{\quoi}.
$$

\begin{enonce*}{Proposition 8.2.7}\etiquette[8.2.7]{V.prop8.2.7}
Let $p:T\to S$ be a morphism of topoi, let $\Oo_S$ be a sheaf of
rings on $S$, and let $\Oo_T$ be its inverse image on $T$. A right
$\Oo_T$-module $\Mm$ is flat over $\Oo_T$ if and only if the functor
\begin{equation*}
\Nn\longmapsto\Mm\mathop{\otimes}_{\Oo_T}p^{\ast}\Nn
\tag{8.2.7.1}
\end{equation*}\etiquette[8.2.7.1]{V.8.2.7.1}
from the category of left $\Oo_S$-modules to the category of abelian
sheaves on $T$ is exact.
\end{enonce*}

The ``only if'' assertion is trivial; suppose, therefore, that the
functor \eqhyperref[V.8.2.7.1]{8.2.7.1} is exact, and let us prove that
$\Mm$ is flat.

We may suppose that $p$ is defined by a morphism of sites
$p:\Tt\to\Ss$ satisfying the hypotheses made in \hyperref[V.prop8.2.6]{8.2.6}. Let
$V\in\Ob\Tt$, $U\in\Ob\Ss$, and
$\phi:V\to p^{\ast}U$. We again denote by $\phi$ the induced morphism
of sites $\phi:\Tt/V\to\Ss/U$.

\begin{enonce*}{Lemma 8.2.8}\etiquette[8.2.8]{V.lem8.2.8}
The functor
$\Nn\longmapsto\Mm|V\mathop{\otimes}_{\Oo_V}\phi^{\ast}\Nn$,
from the category of $\Oo_S|U$-modules on $\Ss/U$ to that of
$\Oo_T|V$-modules on $\Tt/V$, is exact.
\end{enonce*}

We reduce to taking $V=p^{\ast}U$. Let $j$ and $j'$ be the
``inclusion morphisms'' $j:\Ss/U\to\Ss$ and
$j':\Tt/p^{\ast}U\to\Tt$. One has
$$
j'_!(\Mm|V\mathop{\otimes}_{\Oo_V}\phi^{\ast}\Nn)
=\Mm\mathop{\otimes}_{\Oo_T}p^{\ast}j_!\Nn{\quoi},
$$
which gives \hyperref[V.lem8.2.8]{8.2.8}, since $j_!$ and $j'_!$ are exact and faithful.

Let now $Q$ be a left $\Oo_T$-module. By \hyperref[V.prop8.2.6]{(8.2.6)}, one has
$$
Q=\varinjlim_{\phi:V\to p^{\ast}U}
\phi^{\ast}\phi_\ast(Q|V),
$$
so that (\hyperref[V.prop8.2.4]{8.2.4})
$$
\Mm\mathop{\otimes}_{\Oo_T}Q
=\varinjlim_{\phi:V\to p^{\ast}U}
\Mm|V\mathop{\otimes}_{\Oo_V}
\phi^{\ast}\phi_\ast(Q|V){\quoi}.
$$
By\pageoriginale \hyperref[V.lem8.2.8]{8.2.8}, the functors
$Q\longmapsto\Mm|V\mathop{\otimes}_{\Oo_V}
\phi^{\ast}\phi_\ast(Q|V)$ are left exact; therefore
$\Mm\mathop{\otimes}_{\Oo_T}Q$ is left exact in $Q$, and hence exact.

\begin{enonce*}{Corollary 8.2.9}\etiquette[8.2.9]{V.cor8.2.9}
Let $f:(T,\Oo_T)\to(S,\Oo_S)$ be a morphism of ringed topoi. Then the
inverse image by $f$ of a sheaf of flat right modules is a sheaf of
flat modules.
\end{enonce*}

Factoring $f$ through $(T,f^{\ast}\Oo_S)$ reduces us to the case
$\Oo_T=f^{\ast}\Oo_S$. Let $\Mm$ then be a sheaf of flat modules on
$S$. By \hyperref[V.prop8.2.7]{8.2.7}, to verify that $f^{\ast}\Mm$
is flat, it is enough to verify the exactness of the functor
$$
\Nn\longmapsto f^{\ast}\Mm\mathop{\otimes}_{\Oo_T}f^{\ast}\Nn
=f^{\ast}(\Mm\mathop{\otimes}_{\Oo_S}\Nn),
$$
which proves the assertion, since $f^{\ast}$ is exact.

\begin{enonce*}{Lemma 8.2.10}\etiquette[8.2.10]{V.lem8.2.10}
Let $(S,\Oo_S)$ be a ringed topos, let $\Ff$ be a sheaf of left
modules locally of finite presentation, let $\Gg$ be a sheaf of
bimodules, and let $\Hh$ be a flat sheaf of left modules. Then the
canonical arrow
$$
\Hhom(\Ff,\Gg)\otimes\Hh\to
\Hhom(\Ff,\Gg\otimes\Hh)
$$
is an isomorphism.
\end{enonce*}

The question is local on $S$, which allows us to suppose that $\Ff$
has a finite presentation
$$
\Ff_1\to\Ff_0\to\Ff\to0{\quoi}.
$$
The first row of the following diagram is exact, since $\Hh$ is flat:
\[
\xymatrix{
0 \to \Hhom(\Ff,\Gg) \otimes \Hh \ar[r] \ar[d]&
\Hhom(\Ff_0,\Gg) \otimes \Hh \ar[r]\ar[d]^{S}&
\Hhom(\Ff_1,\Gg) \otimes \Hh \ar[d]^{S}\\
0 \to \Hhom(\Ff,\Gg \otimes \Hh) \ar[r]&
\Hhom(\Ff_0,\Gg \otimes \Hh) \ar[r]&
\Hhom(\Ff_1,\Gg \otimes \Hh)
}
\]
which proves the assertion. Taking $\Gg=\Oo_S$, one deduces from
\hyperref[V.lem8.2.10]{8.2.10}:

\begin{enonce*}{Lemma 8.2.11}\etiquette[8.2.11]{V.lem8.2.11}
Every\pageoriginale morphism from a sheaf of modules locally of finite
presentation to a flat sheaf of modules factors locally on $S$ through
a finite free sheaf $\Oo_S^n$ $(n\in\NN)$.
\end{enonce*}

\begin{enonce*}{Theorem 8.2.12}\etiquette[8.2.12]{V.thm8.2.12}
\emph{(D. Lazard).} A sheaf of modules $\Mm$ on a ringed site
$(S,\Oo_S)$ is flat if and only if it is a local inductive limit of
finite free modules.
\end{enonce*}

The ``if'' assertion follows from \hyperref[V.thm8.2.2]{8.2.2} and
\hyperref[V.prop8.2.4]{8.2.4}.

For every $U\in\Ob\Ss$, let $\Ll_0(U)$ (\resp $\Ll_1(U)$) be the
category of finite free sheaves of modules (\resp sheaves of modules
of finite presentation) on $U$, equipped with a linear map to
$\Mm|U$. For every arrow $f:V\to U$ in $\Ss$, let $f^{\ast}$ be the
restriction functor from $\Ll_i(U)$ to $\Ll_i(V)$. These functors
define a category $\Ll_i$, fibered over $\Ss$, whose fibers are the
opposite categories of the categories $\Ll_i(U)$.

The reader will verify that the category $\Ll_1$ is locally filtered
over $\Ss$ and that
$$
\Mm=\varinjlim_{(M,f)\in\Ll_1}M.
$$

If $\Mm$ is flat, the inclusion of $\Ll_0$ in $\Ll_1$ is fully
faithful and, by \hyperref[V.lem8.2.11]{8.2.11}, satisfies
$(C\mathfrak{F}1)$ of \hyperref[V.def8.1.4]{8.1.4}. It then
automatically satisfies $(C\mathfrak{F}2)$, and $\Ll_0$ is locally
filtered. By \hyperref[V.prop8.1.5]{8.1.5}, one has
$$
\Mm=\varinjlim_{(M,f)\in\Ll_0}M{\quoi},
$$
which proves \hyperref[V.thm8.2.12]{8.2.12}.

% END INCLUDED SGA4_English_current_buildable_source_20260906/source/05/drafts/SGA4_V_LINES_4066_4380_English.texfrag


% BEGIN INCLUDED SGA4_English_current_buildable_source_20260906/source/05/drafts/SGA4_V_LINES_4381_4410_English.texfrag
\begin{thebibliography}{99}\pageoriginale
\bibitem{sga4.exp.V.bib.1} M. Artin and B. Mazur: Etale Homotopy Theory, Lecture Notes no.~100, Springer-Verlag.
\bibitem{sga4.exp.V.bib.2} Blum and Herrera: Article forthcoming in Inventiones.
\bibitem{sga4.exp.V.bib.3} H. Cartan and S. Eilenberg: Homological Algebra.
\bibitem{sga4.exp.V.bib.4} P. Deligne: Théorie de Hodge (Publications de l'I.H.E.S.).
\bibitem{sga4.exp.V.bib.5} P. Gabriel and N. Popescu: CRAS.
\bibitem{sga4.exp.V.bib.6} P. Gabriel and M. Zisman: Homotopy Theory and Calculus of
  Fractions, Ergebnisse der Mathematik, Bd 35.
\bibitem{sga4.exp.V.bib.7} R. Godement: Théorie des faisceaux. Hermann.
\bibitem{sga4.exp.V.bib.8} J. Giraud: Méthode de la descente. Mémoire de la S.M.F.
\bibitem{sga4.exp.V.bib.9} J. Giraud: Algèbre homologique non commutative (forthcoming).
\bibitem{sga4.exp.V.bib.10} A. Grothendieck: On the De Rham Cohomology of Algebraic
  Varieties, I.H.E.S. no.~29.
\bibitem{sga4.exp.V.bib.11} A. Grothendieck: Sur quelques points d'Algèbre Homologique,
  Tohoku Math. Journal.
\bibitem{sga4.exp.V.bib.12} R. Hartshorne: Residues and Duality. Lecture Notes
  no.~20, Springer-Verlag.
\bibitem{sga4.exp.V.bib.13} L. Illusie: \SGA 6 I.
\bibitem{sga4.exp.V.bib.14} S. Lubkin: On a Conjecture of A. Weil. Am. J. of
  Math., p.~456, 1967.
\bibitem{sga4.exp.V.bib.15} G. Segal: Classifying Spaces and Spectral Sequences,
  I.H.E.S. no.~34.
\bibitem{sga4.exp.V.bib.16} Séminaire Cartan 1957--1958.
\bibitem{sga4.exp.V.bib.17} D. Sullivan: Geometric Topology, Part I, mimeographed
  notes, M.I.T., 1970.
\end{thebibliography}

\ifx\danslelivre\undefined
\end{document}
\fi

% END INCLUDED SGA4_English_current_buildable_source_20260906/source/05/drafts/SGA4_V_LINES_4381_4410_English.texfrag


% END INCLUDED SGA4_English_current_buildable_source_20260906/source/05/05_English.tex


\renewcommand*{\thechapter}{Vbis}
\makeatletter
\def\input@path{{../05bis/}}
\makeatother
% BEGIN INCLUDED SGA4_English_current_buildable_source_20260906/source/05bis/05bis_English.tex
\ifx\danslelivre\undefined
% Release engine: XeLaTeX (the LuaTeX compatibility layer breaks Xy-pic output).
\documentclass[10pt]{../fichiers-configuration-latex/smfart}
% BEGIN INCLUDED SGA4_English_current_buildable_source_20260906/source/fichiers-configuration-latex/paquets.tex
\usepackage{polyglossia}
%\usepackage[francais]{babel}
\usepackage{fontspec}
\usepackage[math-style=ISO,bold-style=upright,sans-style=upright]{unicode-math}
%\usepackage[mathletters]{ucs}
%\usepackage[utf8]{inputenc}
%\usepackage[T1]{fontenc}
\usepackage[unicode,hypertexnames=false,colorlinks=true,linkcolor=blue,citecolor=Sepia,urlcolor=magenta]{hyperref}
\setmainlanguage{french}
\setmainfont[Ligatures=TeX]{Linux Libertine O}
% variante possible : TeXGyrePagella ?
\setmathfont{XITS Math}

% maths
\usepackage{amsmath}
%\usepackage{amssymb}
%\usepackage{mathrsfs}
%\usepackage{stmaryrd}
%\usepackage{euscript}

% couleurs
\usepackage[usenames,dvipsnames]{xcolor}

% graphiques
\usepackage{graphicx}
\usepackage[all]{xy}
%\usepackage{tikz}
%\usetikzlibrary{matrix,arrows,calc}
%\tikzset{column sep=3em,row sep=2em,text height=1.5ex, text depth=0.25ex,>=angle 90}
%\tikzset{description/.style={fill=white,inner sep=1pt}}

% hyperliens, 'etiquettes
\usepackage{xr-hyper}
\usepackage{../fichiers-configuration-latex/smfthm}
%\ifx\cacheretiquettes\undefined
%\usepackage[inline]{showlabels} \else \fi % pour rendre visibles les 'etiquettes
%\usepackage{../fichiers-configuration-latex/makeidx} % pour faire un index
%\usepackage[pagebackref,unicode]{hyperref}  % pour voir o`u sont cit'ees les r'ef'erences
% hyperref is loaded once above with the cumulative-reader destination policy.
%\hypersetup{breaklinks=true} %permet le retour `a la ligne dans les liens trop longs. Incompatible avec pdflatex
%\renewcommand*{\backref}[1]{{\small [cit\'e page(s)~#1.]}}

\usepackage{xspace}
\usepackage{pifont} % pour expose XVII

% END INCLUDED SGA4_English_current_buildable_source_20260906/source/fichiers-configuration-latex/paquets.tex

% BEGIN INCLUDED SGA4_English_current_buildable_source_20260906/source/fichiers-configuration-latex/numerotation.tex
\renewcommand{\thefootnote}{\roman{footnote}} % pour mettre les notes en bas de page sous la forme (ii) et non (2), qui ressemble trop (?) a un exposant.

\newcounter{pageoriginal}
\newcommand{\pageoriginale}{\refstepcounter{pageoriginal}\marginpar{\footnotesize\textbf{\thepageoriginal}}\xspace}
\let\pageoriginaled\pageoriginale

\usepackage{ifthen}
\newboolean{original}
\def\sisi{\ifthenelse{\boolean{original}}}
\setboolean{original}{false}

\newcounter{cnote}
\def\nde#1{\ifthenelse{\boolean{original}}{}{\stepcounter{cnote}\def\thefootnote{\arabic{cnote}}\protect\footnote[\value{cnote}]{N.D.E.: #1}}}
\def\ndemark{\ifthenelse{\boolean{original}}{}{\stepcounter{cnote}\def\thefootnote{\arabic{cnote}}\protect\footnotemark[\value{cnote}]}}
\def\ndetext#1{\ifthenelse{\boolean{original}}{}{\def\thefootnote{\arabic{cnote}}\protect\footnotetext[\value{cnote}]{N.D.E.: #1}}}

% END INCLUDED SGA4_English_current_buildable_source_20260906/source/fichiers-configuration-latex/numerotation.tex

% BEGIN INCLUDED SGA4_English_current_buildable_source_20260906/source/fichiers-configuration-latex/macros.tex
\setmainlanguage{french}
\setotherlanguage{english}


%% Fixed, engine-independent release identifier.  The inherited dynamic Git
%% probe used Lua io.popen, exposed build-state-dependent output, and printed
%% a visible question mark outside a Git checkout.
\def\showgitstatus{English release r2, 2026-07-23}


% macros pour LaTeX
\renewcommand{\mathsf}{}

% Police différente pour U+2113 SCRIPT SMALL L
\setmathfont[range="2113]{latinmodern-math.otf}

% Give every inherited semantic label its own stable PDF destination instead
% of inheriting whichever automatic equation/page anchor happened to be live.
% The legacy label remains as a compatibility alias used by existing edges.
\makeatletter
\def\etiquette[#1]#2{%
  \Hy@raisedlink{\hyper@anchorstart{sga4.semantic.#2}\hyper@anchorend}%
  \def\@currentHref{sga4.semantic.#2}%
  \def\@currentlabel{#1}%
  \label{#2}%
}
\makeatother
\def\eqhyperref[#1]#2{(\hyperref[#1]{#2})}
\def\pagehyperref[#1]#2{(page \hyperref[#1]{#2})}
\def\hyperrefext[#1][#2]#3{\hyperref[#2]{#1 #3}}
\def\eqhyperrefext[#1][#2]#3{\hyperref[#2]{#1 #3}}     %cf. exposes 17
\def\neweqhyperrefext[#1][#2]#3{\hyperref[#2]{#1 #3}}  %cf. exposes 17

\def\shat#1{\protect\widehat{#1}}
\def\stilde#1{\protect\widetilde{#1}}
%%%%%%%%%%%%%%%%%%


\def\etcn{etc.\xspace}
\def\etco{etc...\xspace}
\def\etc{\sisi{\etco}{\etcn}}
\def\ldotsn{\sisi{...}{\etcn}}

\def\cf{{cf.}\kern.3em}
\def\Cf{{Cf.}\kern.3em}
\def\ie{{i.e.}\kern.3em}
\def\iev{{i.e.},\xspace}
\def\loccito{loc.\kern2pt cit.\xspace}
\def\loccitn{\textit{loc.\kern2pt cit.}\xspace}
\def\loccit{\sisi{\loccito}{\loccitn}}
\def\EGA{\textup{EGA}\kern.2em}
\def\SGA{\textup{SGA}\kern.3em}

\def\quoi{}

\def\quav{\quad,\quad}
\def\quaet{\quad \text{et}\quad}
\def\fibreproduct#1#2#3{#1 ×_{#3} #2}
\let\fprod\fibreproduct
\def\ccheck#1#2{\mbox{${\check{\check{\rm #1}}}^{\lower2pt\hbox{$\scriptstyle #2$}}$}}
\def\sheaf#1{\protect\mathcal{#1}}
\def\vieuxsheaf#1{{#1}} % cf. problème avec sheaf{a\Phi_i} par exemple...
\DeclareMathOperator{\Hom}{Hom}
\DeclareMathOperator{\can}{can}
\DeclareMathOperator{\Ind}{Ind}
\DeclareMathOperator{\ind}{ind}
%\DeclareMathOperator{\Hhomtop}{\mathcal{H}\!\mathit{omtop}}
\DeclareMathOperator{\Hhomtop}{\mathbf{Homtop}}
\DeclareMathOperator{\Point}{Point}
\DeclareMathOperator{\Pro}{Pro}
\DeclareMathOperator{\Top}{Top}
\DeclareMathOperator{\Mod}{Mod}
\DeclareMathOperator{\Cart}{Cart}
%\DeclareMathOperator{\Mmorsite}{\mathcal{M}\!\mathit{orsite}}
\DeclareMathOperator{\Mmorsite}{\mathbf{Morsite}}
\DeclareMathOperator{\loc}{loc}
\DeclareMathOperator{\Loc}{Loc}
\DeclareMathOperator{\Pic}{Pic}
\DeclareMathOperator{\Et}{Et}
\DeclareMathOperator{\card}{card}
\DeclareMathOperator{\Ens}{Ens}
\DeclareMathOperator{\h}{\underline{H}}
\DeclareMathOperator{\Fonct}{Fonct}
\DeclareMathOperator{\Ker}{Ker}
\DeclareMathOperator{\but}{but}
\DeclareMathOperator{\Ob}{Ob}
\DeclareMathOperator{\ob}{ob}
\DeclareMathOperator{\Lib}{Lib}
\DeclareMathOperator{\esj}{esj}
\DeclareMathOperator{\id}{id}
\DeclareMathOperator{\Id}{Id}
\DeclareMathOperator{\ad}{ad}
\DeclareMathOperator{\Iim}{Im}
\DeclareMathOperator{\pr}{pr}
\DeclareMathOperator{\Cat}{Cat}
\DeclareMathOperator{\car}{car}
\DeclareMathOperator{\SysInd}{SysInd}
\DeclareMathOperator{\Sup}{Sup}
\DeclareMathOperator{\Coker}{Coker}
\DeclareMathOperator{\Filt}{Filt}
\DeclareMathOperator{\acc}{acc}
\DeclareMathOperator{\sur}{sur}
\DeclareMathOperator{\cat}{cat}
\DeclareMathOperator{\Ord}{Ord}
\DeclareMathOperator{\Gr}{Gr}
\DeclareMathOperator{\Cov}{Cov}
\DeclareMathOperator{\Caf}{Caf}
\DeclareMathOperator{\coker}{coker}
\DeclareMathOperator{\coim}{coim}
\DeclareMathOperator{\im}{im}
\DeclareMathOperator{\Ouv}{Ouv}
\DeclareMathOperator{\diag}{diag}
\DeclareMathOperator{\Esp}{Esp}
\DeclareMathOperator{\TOP}{TOP}
\DeclareMathOperator{\Unif}{Unif}
\DeclareMathOperator{\Iind}{\mathbf{Ind}}
\DeclareMathOperator{\Hhomcart}{\mathbf{Homcart}}
\DeclareMathOperator{\Hhomtopess}{\mathbf{Homtopess}}
\DeclareMathOperator{\Ppoints}{\mathbf{Points}}
\DeclareMathOperator{\Ppoint}{\mathbf{Point}}
\DeclareMathOperator{\Ttors}{\mathbf{Tors}}
\DeclareMathOperator{\Ssex}{\mathbf{Sex}}
\DeclareMathOperator{\Ffib}{\mathbf{Fib}}
\DeclareMathOperator{\Vvois}{\mathbf{Vois}}
\DeclareMathOperator{\Ppointess}{\mathbf{Pointess}}
\DeclareMathOperator{\scrHom}{\mathbf{Hom}}
\DeclareMathOperator{\sheafPic}{\mathbf{Pic}}
\DeclareMathOperator{\sheafAut}{\mathbf{Aut}}
\DeclareMathOperator{\Hhomtoplocan}{\mathbf{Homtoplocan}}
\DeclareMathOperator{\Hhomtopan}{\mathbf{Homtopan}}
\DeclareMathOperator{\Pptgeom}{\mathbf{Ptgéom}}
\DeclareMathOperator{\Morsite}{Morsite}
\DeclareMathOperator{\Fib}{Fib}
\DeclareMathOperator{\Fais}{Fais}
\DeclareMathOperator{\kar}{kar}
\DeclareMathOperator{\Groupes}{Groupes}
\DeclareMathOperator{\Restr}{Restr}
\DeclareMathOperator{\Isom}{Isom}
\DeclareMathOperator{\Prol}{Prol}
\DeclareMathOperator{\Res}{Res}
\DeclareMathOperator{\Tors}{Tors}
\DeclareMathOperator{\esp}{esp}
\DeclareMathOperator{\es}{es}
\DeclareMathOperator{\od}{od}
\DeclareMathOperator{\gr}{gr}
\DeclareMathOperator{\sob}{sob}
\DeclareMathOperator{\dfn}{dfn}
\DeclareMathOperator{\Sch}{Sch}
\DeclareMathOperator{\ouv}{ouv}
\DeclareMathOperator{\cons}{cons}
\DeclareMathOperator{\Nerf}{Nerf}
\DeclareMathOperator{\Points}{Points}
\DeclareMathOperator{\Espsob}{Espsob}
\DeclareMathOperator{\Homtop}{Homtop}
\DeclareMathOperator{\Pt}{Pt}
\DeclareMathOperator{\Quot}{Quot}
\DeclareMathOperator{\glob}{glob}
\DeclareMathOperator{\Homfib}{Homfib}
\DeclareMathOperator{\Qouv}{Qouv}
\DeclareMathOperator{\Quotouv}{Quotouv}
\DeclareMathOperator{\ens}{ens}
\DeclareMathOperator{\ab}{ab}
\DeclareMathOperator{\an}{an}
\DeclareMathOperator{\fs}{fs}
\DeclareMathOperator{\red}{red}
\DeclareMathOperator{\Ab}{Ab}
\DeclareMathOperator{\cosq}{cosq}
\DeclareMathOperator{\Ann}{Ann}
\DeclareMathOperator{\Zar}{Zar}


%%%%%%%%%%%%
%%%%% SGA4-2
\DeclareMathOperator{\Sex}{Sex}
\DeclareMathOperator{\vois}{vois}
\DeclareMathOperator{\Vois}{Vois}
\DeclareMathOperator{\Ext}{Ext}
\DeclareMathOperator{\ext}{ext}
\DeclareMathOperator{\Fer}{Fer}
\DeclareMathOperator{\Add}{Add}
\DeclareMathOperator{\Cocart}{Cocart}
\DeclareMathOperator{\cosk}{cosk}
\DeclareMathOperator{\plat}{plat}

\DeclareMathOperator{\faisc}{faisc}
\DeclareMathOperator{\sk}{sk}
\DeclareMathOperator{\ch}{ch}

\def\HH{\mathbf{H}}
\def\qc{\mathrm{qc}}
\def\et{\mathrm{et}}
\def\coh{\mathrm{coh}}

\def\Hvr#1{%
\begin{array}[b]{l}
{\kern2pt\scriptstyle r}\\[-8pt]
{\kern1pt\scriptstyle \vee}\\[-5pt]
H^{#1}
\end{array}}

\def\Hvv#1{%
\begin{array}[b]{l}
\kern2pt\rotatebox{90}{${\scriptscriptstyle\ll}$}\\[-6pt]
H^{#1}
\end{array}}

\def\Hvi#1{%
\begin{array}[b]{l}
{\kern0.2pt\scriptscriptstyle \infty}\\[-8pt]
{\kern1pt\scriptstyle \vee}\\[-5pt]
H^{#1}
\end{array}}

\def\Hv#1{%
\begin{array}[b]{l}
{\kern1pt\scriptstyle \vee}\\[-5pt]
H^{#1}
\end{array}}

\DeclareMathOperator{\cart}{cart}
\DeclareMathOperator{\Ffaisc}{\mathbf{Faisc}}


\def\Ev#1{%
\begin{array}[b]{l}
{\kern1pt\scriptstyle \vee}\\[-5pt]
E'_{#1}
\end{array}}

\DeclareMathOperator{\Cont}{Cont}
\DeclareMathOperator{\Ccont}{\mathscr{C}\!\mathit{ont}}
\DeclareMathOperator{\Site}{Site}
\newcommand\MM{\mathbf{M}}
\def\limtop#1{{\displaystyle{\varprojlim_{#1}\!\mbox{\rm top }}}}
\def\Limtop{\mathop{\varprojlim\!\mbox{\rm top}}}
\DeclareMathOperator{\spec}{spec}
\def\ss{\scriptstyle{\mathscr{S}}}
\DeclareMathOperator{\Gal}{Gal}
\def\smallm{\scriptstyle{\mathscr{M}}}
\DeclareMathOperator{\End}{End}
\DeclareMathOperator{\ass}{ass}
\DeclareMathOperator{\zar}{zar}
\DeclareMathOperator{\Ogg}{Ogg}
\DeclareMathOperator{\res}{res}
\DeclareMathOperator{\tr}{tr}
\DeclareMathOperator{\Supp}{Supp}
\DeclareMathOperator{\codim}{codim}
\DeclareMathOperator{\rk}{rk}
\DeclareMathOperator{\cur}{cur}
\DeclareMathOperator{\cdqc}{cdqc}
\newcommand\DD{\mathbf{D}}
\DeclareMathOperator{\cl}{cl}
\def\JJ{\mathbf{J}}
\def\FF{\mathbf{F}}
\def\BB{\mathbf{B}}
\DeclareMathOperator{\Alget}{Alget}
\DeclareMathOperator{\Ass}{Ass}
\DeclareMathOperator{\supp}{supp}
\DeclareMathOperator{\rad}{rad}
\DeclareMathOperator{\prof}{prof}
\DeclareMathOperator{\cint}{int}
\DeclareMathOperator{\sic}{sic}
\DeclareMathOperator{\Tor}{Tor}
\DeclareMathOperator{\tors}{tors}
\DeclareMathOperator{\Sym}{Sym}
\DeclareMathOperator{\parf}{parf}
\DeclareMathOperator{\tor}{tor}
\DeclareMathOperator{\defn}{defn}
\DeclareMathOperator{\Div}{Div}
\DeclareMathOperator{\Pol}{Pol}
\DeclareMathOperator{\Fex}{Fex}
\DeclareMathOperator{\Lie}{Lie}
\DeclareMathOperator{\tot}{tot}
\DeclareMathOperator{\Aut}{Aut}
\DeclareMathOperator{\HOM}{HOM}
\DeclareMathOperator{\pch}{pch}
\DeclareMathOperator{\Ch}{Ch}
\DeclareMathOperator{\EXT}{EXT}
\DeclareMathOperator{\PIC}{PIC}
\DeclareMathOperator{\Rf}{Rf}
\DeclareMathOperator{\inj}{inj}
\newcommand\VV{\mathbf{V}}

\def\Xv{%
\begin{array}[b]{l}
{\kern1pt\scriptstyle \vee}\\[-5pt]
X
\end{array}}

%%%%%%%
\newcommand\TT{\mathbf{T}}
\newcommand\XX{\mathbf{X}}
\newcommand\EE{\mathbf{E}}
%%%%%%

\def\qedsymbol{\textup{C.Q.F.D.}}
\let\CQFD\qed
\let\cqfd\qed

%%%%%\bigdot
\font\bigsymb=cmsy10 at 4pt
\def\bigdot{{\kern1.2pt\raise 1.5pt\hbox{\bigsymb\char15}}}
%%%%%

\def\Lright{\mathop{\vcenter{\hbox{\hbox{$\longrightarrow$}\lower8pt \hbox{\kern-16pt\hbox{$\longrightarrow$}}}}}}

\def\Aa{{\mathcal A}}
\def\Bb{{\mathcal B}}
\def\Cc{{\mathcal C}}
\def\Dd{{\mathcal D}}
\def\Ee{{\mathcal E}}
\def\Ff{{\mathcal F}}
\def\Gg{{\mathcal G}}
\def\Hh{{\mathcal H}}
\def\Ii{{\mathcal I}}
\def\Jj{{\mathcal J}}
\def\Kk{{\mathcal K}}
\def\Ll{{\mathcal L}}
\def\Mm{{\mathcal M}}
\def\Nn{{\mathcal N}}
\def\Oo{{\mathcal O}}
\def\Pp{{\mathcal P}}
\def\Qq{{\mathcal Q}}
\def\Rr{{\mathcal R}}
\def\Ss{{\mathcal S}}
\def\Tt{{\mathcal T}}
\def\Uu{{\mathcal U}}
\def\Vv{{\mathcal V}}
\def\Ww{{\mathcal W}}
\def\Xx{{\mathcal X}}
\def\Yy{{\mathcal Y}}
\def\Zz{{\mathcal Z}}

\newcommand{\CC}{\mathbf{C}}
\newcommand{\GG}{\mathbf{G}}
\newcommand{\KK}{\mathbf{K}}
\newcommand{\LL}{\mathbf{L}}
\newcommand{\NN}{\mathbf{N}}
\newcommand{\PP}{\mathbf{P}}
\newcommand\QQ{\mathbf{Q}}
\newcommand\RR{\mathbf{R}}
\newcommand\ZZ{\mathbf{Z}}

\def\to{\mathchoice{\longrightarrow}{\rightarrow}{\rightarrow}{\rightarrow}}
\def\from{\mathchoice{\longleftarrow}{\leftarrow}{\leftarrow}{\leftarrow}}
\def\mto{\mathchoice{\longmapsto}{\mapsto}{\mapsto}{\mapsto}}
\def\mfrom{\mathrel{\from\joinrel\mapstochar\,}}
\def\hto{\mathrel{\lhook\joinrel\to}}
\def\hfrom{\mathrel{\from\joinrel\rhook}}

\def\implies{\mathchoice{\Longrightarrow}{\Rightarrow}{\Rightarrow}{\Rightarrow}}
\def\implied{\mathchoice{\Longleftarrow}{\Leftarrow}{\Leftarrow}{\Leftarrow}}
\def\iff{\mathchoice{\Longleftrightarrow}{\Leftrightarrow}{\Leftrightarrow}{\Leftrightarrow}}

\def\To#1{\mathchoice{\xrightarrow{\labelstyle~~#1~}}{\stackrel{#1}{\longrightarrow}}{}{}}
\def\Mto#1{\mathrel{\mapstochar\To{#1}}}
\def\From#1{\mathchoice{\xleftarrow{\labelstyle~~#1~}}{\stackrel{#1}{\longleftarrow}}{}{}}

\newcommand{\isomto}{\overset{\sim}{\longrightarrow}}
\newcommand{\isomfrom}{\overset{\sim}{\longleftarrow}}

\newcommand{\IFF}{\unskip~$\iff$~\ignorespaces}
\newcommand{\THEN}{\unskip~$\implies$~\ignorespaces}

\let\leq\leqslant
\let\geq\geqslant
\let\le\leqslant
\let\ge\geqslant

\def\cf{{cf.}\kern.3em}
\def\Cf{{Cf.}\kern.3em}
\def\ie{{i.e.}\kern.3em}
\def\iev{{i.e.},\xspace}
%\def\loccito{loc.\kern2pt cit.\xspace}
%\def\loccitn{\textit{loc.\kern2pt cit.}\xspace}
%\def\loccit{\sisi{\loccito}{\loccitn}}
\def\EGA{\textup{EGA}\kern.2em}
\def\SGA{\textup{SGA}\kern.3em}

\let\oldS\S
\def\S{\oldS\kern .15em}
\def\ptbl{.\kern .2em}
\def\resp{resp.\kern.3em}
\def\Exp{Exp.\kern.3em}

\let\bar\overline
\let\hat\widehat
\let\tilde\widetilde
\let\dpl\displaystyle
\let\ldots\dots
\let\emptyset\varnothing
\let\skpt\skippointrait
\let\epsilon\varepsilon

\def\sfrac#1#2{{#1}/{#2}}

\DeclareMathOperator{\cd}{cd}
%\DeclareMathOperator{\Hhom}{\mathcal{H}\!\mathit{om}}
\DeclareMathOperator{\Hhom}{\mathbf{Hom}}
%\DeclareMathOperator{\Eext}{\mathcal{E}\!\mathit{xt}}
\DeclareMathOperator{\Eext}{\mathbf{Ext}}
\DeclareMathOperator{\Spec}{Spec}
\DeclareMathOperator{\Tr}{Tr}

\let\emptyset=\varnothing

%%%%%%%%%%%%%%%
%%% expose XVII
%%%%%%%%%%%%%%%

\makeatletter
\let\old@optexpose\optexpose
\newcommand{\empty@optexpose}[1]{}
\newcommand{\paren@optexpose}[1]{\old@optexpose{#1}(}
\newcommand{\refnoexpose}[1]{\let\optexpose\empty@optexpose\ref{#1}\let\optexpose\old@optexpose}%\refnoexpose{X.cor5.2} engendre "5.2" au lieu de "X 5.2"
\newcommand{\neweqref}[1]{\let\optexpose\paren@optexpose\textup{\ref{#1})}\let\optexpose\old@optexpose}%\neweqref{XII.eq4.1} engendre "XII (4.1)" au lieu de "(XII 4.1)"
\makeatother
\newcommand{\refvb}[2][VI${}_{\text{B}}$ ]{\sisi{#1\refnoexpose{#2}}{\ref{#2}}}

\newcommand{\leftexp}[2]{{\vphantom{#2}}^{#1}{#2}}
\newcommand{\AAA}{\mathbf{A}}
\newcommand{\Sf}{\mathfrak{S}}
\DeclareMathOperator*{\quotprojlim}{\text{``$\varprojlim$''}}
\DeclareMathOperator*{\quotinjlim}{\text{``$\varinjlim$''}}
\DeclareMathOperator{\Fl}{F\sisi{\ell}{l}}
\DeclareMathOperator{\Pl}{P\sisi{\ell}{l}}
%\DeclareMathOperator{\cl}{c\sisi{\ell}{l}}
\DeclareMathOperator{\tdf}{tdf}
\DeclareMathOperator{\rg}{rg}
\newcommand{\iInd}{\sisi{I}{i}nd}
\newcommand{\iind}{ind}
\newcommand{\ppro}{pro}
\newcommand{\tTor}{\sisi{T}{t}or}
\newcommand{\ttor}{tor}
\newcommand{\chap}{Chap.}
\newcommand{\Ffex}{\Ff\mathit{ex}}
\newcommand{\sheafk}{\underline{k}}
\newcommand{\sheafn}{\underline{n}}
\newcommand{\sheafGamma}{\underline{\Gamma}}

% Repeated references to an already typeset footnote must not use
% \footnotemark without a matching \footnotetext: Hyperref would then create
% a live annotation to a destination that cannot exist.  This helper keeps the
% current printed footnote mark and links it to an explicit stable note target.
\makeatletter
\newcommand{\sgarepeatedfootnotemark}[1]{%
  \begingroup
  \edef\@thefnmark{\thefootnote}%
  \hyperlink{#1}{\@makefnmark}%
  \endgroup}
\makeatother

\sisi{
\DeclareMathOperator{\simarrow}{\rule{.15mm}{1.8mm}\kern -3.4pt\rightsquigarrow}
\newcommand{\sheafell}{\underline{\ell}}
\newcommand{\sheafZZ}{\underline{\ZZ}}
}{
\newcommand{\nderef}[1]{la N.D.E.~(\ref{#1}) page~\pageref{#1}}
}

%%%%%%%%%%%%%%%%%
%%%%% expose XVII
%%%%%%%%%%%%%%%%%

% END INCLUDED SGA4_English_current_buildable_source_20260906/source/fichiers-configuration-latex/macros.tex


\externaldocument{../01/01_English.tex}
\externaldocument{../02/02_English.tex}
\externaldocument{../03/03_English.tex}
\externaldocument{../04/04_English.tex}
\externaldocument{../05/05_English.tex}
%\externaldocument{../05bis/05bis_English.tex}

\setboolean{original}{false}

\begin{document}

\begin{center}
Vbis. Cohomological Descent Techniques\\
B. Saint-Donat
\end{center}

\selectlanguage{english}version: \showgitstatus
\tableofcontents

\else
\chapter{Cohomological Descent Techniques}
%\addtocontents{toc}{B. Saint-Donat}
\begin{center}
B. Saint-Donat
\end{center}
\fi
\phantomsection\MakeLinkTarget*{sga4.exp.Vbis}\label{sga4.exp.Vbis}

\setcounter{pageoriginal}{83}

%\footnotetext[1]{From concise notes by P. Deligne}

\let\cosk=\cosq
\let\Ysheaf=\sheaf

% BEGIN INCLUDED SGA4_English_current_buildable_source_20260906/source/05bis/drafts/SGA4_VBIS_LINES_59_205_English.texfrag
\section*{Introduction}

\begin{enumerate}
\renewcommand{\labelenumi}{\theenumi.}
\item Let\pageoriginale $X$ be a topological space and let
$\Uu=(U_i)_{i\in I}$ be a covering of $X$, assumed either open or
closed and locally finite. If $\mathfrak{F}$ is an abelian sheaf on
$X$, the Leray spectral sequence
\begin{equation*}
\Hv{p}(\Uu,\mathcal{H}^q(\mathfrak{F}))\Rightarrow H^{\ast}
(X,\mathfrak{F})\tag{1.1}
\end{equation*}\etiquette[1.1]{Vbis.eq1.1}
defined by $\Uu$ [\cite{sga4.exp.Vbis.bib.5} II (5.2.4) and (5.4.1)] may be described
as follows.

The covering $\Uu$ defines a ``\v{C}ech'' resolution
$\mathcal{C}^{\ast}(\Uu,\mathfrak{F})$, functorial in
$\mathfrak{F}$ (ibid. (\hyperref[Vbis.subsubsec5.2.1]{5.2.1})). On
the other hand, for every $\mathfrak{F}$ there is a ``canonical
flasque'' resolution $C^{\ast}(\mathfrak{F})$, functorial in
$\mathfrak{F}$ (ibid. (\hyperref[Vbis.subsec4.3]{4.3})). With this
notation, when $\Uu$ is open the spectral sequence
(\hyperref[Vbis.subsec1.1]{1.1}) is obtained from the double complex
\begin{equation*}
\Gamma(X,\mathcal{C}^{\ast}(\Uu,C^{\ast}(\mathfrak{F}))){\quoi}.
\tag{1.2}
\end{equation*}\etiquette[1.2]{Vbis.eq1.2}
When $\Uu$ is closed and locally finite, one considers the double
complex
\begin{equation*}
\Gamma(X,C^{\ast}(\mathcal{C}^{\ast}(\Uu,\mathfrak{F}))){\quoi}.
\tag{1.3}
\end{equation*}\etiquette[1.3]{Vbis.eq1.3}

\item Let us seek a unified description of these double complexes.
Let $X_0$ be the topological disjoint union of the $U_i$, and for
$n\geq0$ let $X_n$ be the $(n+1)$-fold fiber product of $X_0$ with
itself over $X$:
\begin{equation*}
X_n=\mathop{\coprod}_{i_0,\ldots,i_n\in I}U_{i_0}\cap\cdots\cap U_{i_n}
=\mathop{\coprod}_{\sigma\in\Hom([n],I)}
\mathop{\bigcap}_{i=0}^{n}U_{\sigma(i)}{\quoi}.\tag{2.1}
\end{equation*}\etiquette[2.1]{Vbis.eq2.1}
The $X_n$ form a simplicial system of topological spaces, and if
$j_n$ denotes the projection from $X_n$ to $X$, then
\begin{equation*}
\mathcal{C}^n(\Uu,\mathfrak{F})=j_{n\ast}j_n^{\ast}(\mathfrak{F}){\quoi}.
\tag{2.2}
\end{equation*}\etiquette[2.2]{Vbis.eq2.2}

Note that formation of the canonical flasque resolutions commutes
with\pageoriginale restriction to an open subset and with direct
image by a closed immersion. Consequently:
\begin{enumerate}
\renewcommand{\theenumi}{\alph{enumi}}
\renewcommand{\labelenumi}{\theenumi)}
\item if $\Uu$ is open,
\begin{equation*}
\mathcal{C}^q(\Uu,C^p(\mathfrak{F}))=j_{q\ast}j_q^{\ast}
C^p(\mathfrak{F})=j_{q\ast}C^p(j_q^{\ast}\mathfrak{F});\tag{2.3}
\end{equation*}\etiquette[2.3]{Vbis.eq2.3}
\item if $\Uu$ is closed and locally finite,
\begin{equation*}
C^p(\mathcal{C}^q(\Uu,\mathfrak{F}))=C^p(j_{q\ast}j_q^{\ast}
\mathfrak{F})=j_{q\ast}C^p(j_q^{\ast}\mathfrak{F}){\quoi}.\tag{2.4}
\end{equation*}\etiquette[2.4]{Vbis.eq2.4}
\end{enumerate}
Thus, to obtain a unified description of
\eqhyperref[Vbis.eq1.2]{1.2} and \eqhyperref[Vbis.eq1.3]{1.3}, it is
enough to take the canonical flasque resolution of
$j_q^{\ast}(\mathfrak{F})$ on $X_q$ for every $q$, and then apply
the functor $j_{q\ast}$ to this resolution.

\item The preceding description remains meaningful for any
simplicial system of topological spaces over $X$:
\hypertarget{sga4.exp.Vbis.equation.3.1}{}\begin{gather*}
\Delta^0\longrightarrow\Top/X\tag{3.1}\\
[n]\longmapsto X_n.\notag
\end{gather*}
Such a system need not be of the form \eqhyperref[Vbis.eq2.1]{2.1}.
The double complex, coaugmented by $\mathfrak{F}$,
\begin{equation*}
\mathfrak{F}\longrightarrow
(j_{q\ast}C^p(j_q^{\ast}(\mathfrak{F})))_{p,q}\tag{3.2}
\end{equation*}\etiquette[3.2]{Vbis.eq3.2}
will not, however, in general define a resolution of
$\mathfrak{F}$.

This work is devoted to finding sufficient conditions for
\eqhyperref[Vbis.eq3.2]{3.2} to define a resolution of
$\mathfrak{F}$. In that case, the spectral sequence
(\hyperref[Vbis.subsec1.1]{1.1}) generalizes to a spectral sequence
\begin{equation*}
\Hv{p}(H^q(X_p,j_p^{\ast}(\mathfrak{F})))\Rightarrow
H^{p+q}(X,\mathfrak{F}),\tag{3.3}
\end{equation*}\etiquette[3.3]{Vbis.eq3.3}
called the ``descent spectral sequence''.

In\pageoriginale the case of ``constant coefficients'', analogous
spectral sequences were obtained by Segal (\cf \cite{sga4.exp.Vbis.bib.8}), by other
methods and for other ``cohomology theories'', such as $K$-theory.
Segal works in the category of CW-complexes: he uses a ``geometric
realization'' functor which associates a new topological space with
a semi-simplicial complex of topological spaces; this new space is
to be compared with the topos associated with a simplicial topos
[\cf \hyperref[Vbis.prop1.2.12]{(1.2.12)}].

\item In \hyperref[Vbis.sec5]{Section 5} we illustrate the criteria obtained by using
resolution of singularities to construct, for every analytic space
$X$ over $\CC$, a simplicial system of nonsingular analytic spaces
over $X$,
$$
[n]\longmapsto X_n,
$$
such that \eqhyperref[Vbis.eq3.2]{3.2} defines a resolution of
$\mathfrak{F}$. Taking the constant sheaf $\CC$ for
$\mathfrak{F}$ gives in particular a spectral sequence
\begin{equation*}
H^q(X_p,\CC)\Longrightarrow H^{p+q}(X,\CC)\tag{4.1}
\end{equation*}\etiquette[4.1]{Vbis.eq4.1}
which expresses the complex cohomology of $X$ in terms of the
complex cohomology of nonsingular analytic spaces. Moreover, if $X$
is projective, all the $X_p$ may be assumed projective; this is the
essential ingredient which yields a kind of ``Hodge theory'' for
$X$ (\cf \cite{sga4.exp.Vbis.bib.2}).

\item The preceding constructions extend verbatim when the sheaf
$\mathfrak{F}$ is replaced by a bounded-below complex of sheaves.
They lead to ``localization'' techniques in derived categories.

It is known that, for $X_0$ given by \eqhyperref[Vbis.eq2.1]{2.1},
the arrow
$$
j_0^{\ast}:D^b(X)\longrightarrow D^b(X_0)
$$
is not faithful in general. We shall show that data more precise
than $j_0^{\ast}(K^{\bigdot})$\pageoriginale (for
$K^{\bigdot}\in D^+(X)$), involving the $X_n$, sometimes make it
possible to reconstruct the complex $K^{\bigdot}$.

The resulting statements will be used in the appendix to Exposé
XVII to extend the definition of the functor $Rf_!$ (where $f$ is a
separated morphism of finite type between schemes) to the case in
which $f$ is not assumed compactifiable.

For this application it is not possible to consider only
topological spaces, which are replaced here by the étale sites of
schemes. Moreover, to carry out the proofs it will be necessary to
consider both simplicial systems of spaces and multisimplicial
systems. This explains, justifies, or excuses the degree of
hypergenerality from which we shall begin.
\end{enumerate}


% END INCLUDED SGA4_English_current_buildable_source_20260906/source/05bis/drafts/SGA4_VBIS_LINES_59_205_English.texfrag


% BEGIN INCLUDED SGA4_English_current_buildable_source_20260906/source/05bis/drafts/SGA4_VBIS_LINES_206_516_English.texfrag
\section{Preliminaries}\etiquette[1]{Vbis.sec1}

\subsection{Notation}\etiquette[1.1]{Vbis.subsec1.1}

\subsubsection{}\etiquette[1.1.1]{Vbis.subsubsec1.1.1}
Throughout what follows, $\Uu$ is a universe such that $\ZZ \in
\Uu$: \emph{all the topoi considered will be
$\Uu\text{-}topoi$.}

Let $T$ and $T'$ be two topoi: a morphism $\varphi : T \to T'$
consists of a pair of functors $\varphi_{\ast} : T \to T'$ and
$\varphi^{\ast} : T'\to T$, endowed with an adjunction
$$\Hom_{T'}(.,\varphi_{\ast}.) \isomto{\Hom}_{T}(\varphi^{\ast} . ,
.),$$
such that $\varphi^{\ast}$ is left exact (\ie preserves finite
projective limits).

Let $(T, \mathcal{O}_{T})$ and $(T', \mathcal{O}_{T'})$ be two
ringed topoi: a morphism from $(T, \mathcal{O}_{T})$ to
$(T', \mathcal{O}_{T'})$ is a pair $(\varphi, \theta)$, where
$\varphi : T \to T'$ is a morphism of topoi and
$\theta : \mathcal{O}_{T'} \to \varphi_{\ast} (\mathcal{O}_T)$ is a
morphism of rings.

\subsubsection{}\etiquette[1.1.2]{Vbis.subsubsec1.1.2}
We will make constant use of the language of fibered categories as
presented in [\SGA 1 VI]; the reader may also consult \cite{sga4.exp.Vbis.bib.4}. Let
us merely fix some notation: if $E \to B$ is a fibered (\resp
cofibered) functor, then, for a morphism $m : i \to j$ in $B$, we
denote by $m^{\ast} : E_{j} \to E_i$ (\resp
$m_{\ast} : E_i \to E_j$) the associated ``inverse image'' (\resp
\pageoriginale ``direct image'') functor; each of these functors is
defined up to a unique functorial isomorphism. If
$\varphi : E \to E'$ is a $B$-functor, then, for every object $i$ of
$B$, we denote by $\varphi_i : E_i \to E'_i$ the restriction of
$\varphi$ to $E_i$.

\subsubsection{}\etiquette[1.1.3]{Vbis.subsubsec1.1.3}

Finally, we denote by $\Delta$ the following category: the objects
of $\Delta$ are the ordered sets $[n] = \{ 0, 1, \ldots, n \}$, and
the morphisms of $\Delta$ are all weakly increasing maps.
$\Delta^{+}$ (\resp $\Delta^{-}$) denotes the category whose objects
are those of $\Delta$ and whose arrows are the monomorphisms (\resp
the epimorphisms) of $\Delta$. We will freely use, as needed, the
classical notation introduced in connection with $\Delta$ [\cf
\cite{sga4.exp.Vbis.bib.3} II 2].

\subsection{D-topoi}\etiquette[1.2]{Vbis.subsec1.2}

\begin{enonce*}{Definition 1.2.1}\etiquette[1.2.1]{Vbis.def1.2.1}
Let $D$ be a category, and let $E \to D$ be a fibered and cofibered
functor. We say that $E$ is bifibered in topoi over $D$, or that $E$
is a $D$-topos, if the following conditions hold:
\begin{enumerate}
\renewcommand{\theenumi}{\alph{enumi}}
\renewcommand{\labelenumi}{(\theenumi)}
\item For every object $i$ of $D$, the fiber category $E_i$ is a
  topos.

\item For every morphism $m : i \to j$ in $D$, there exists a
  morphism of topoi $f : E_j \to E_i$ such that
  $f_{\ast} = m^{\ast}$ and $f^{\ast} = m_{\ast}$.
\end{enumerate}
\end{enonce*}

\begin{enonce*}[remark]{Remark}
In view of (a), condition (b) can equivalently be expressed by saying
that the functor $m_{\ast}$ is left exact\nde{Care must be taken here:
in a fibered category, $m_\star$ is left adjoint to $m^\star$ and is
therefore tautologically right exact; cf. [\SGA 1 VI.10]. This
explains the formulas $f_{\ast} = m^{\ast}$ and
$f^{\ast}=m_{\ast}$.}.
\end{enonce*}

When $D = \Delta$ (\resp $\Delta \times \Delta$), we speak of a
\emph{simplicial topos} (\resp a \emph{double simplicial topos}) to
mean a $\Delta$-topos (\resp a $\Delta \times \Delta$-topos).

In practice, we will encounter $D$-topoi by means of the following
considerations.

\begin{enonce*}{Definition 1.2.2}\etiquette[1.2.2]{Vbis.def1.2.2}
Let\pageoriginale $\mathcal{E}$ be a fibered and cofibered category
over $D$. We say that $\mathcal{E}$ is \emph{bifibered in duals of
topoi over $D$} if $\mathcal{E}^{\circ}$ is a $D^{\circ}$-topos.
\end{enonce*}

\hypertarget{sga4.exp.Vbis.remark.1.2.3}{}\begin{enonce*}[remark]{Remark 1.2.3}
Explicitly, $\mathcal{E}$ is bifibered in duals of topoi over $D$ if
and only if the following conditions hold:
\begin{enumerate}
\renewcommand{\theenumi}{\alph{enumi}}
\renewcommand{\labelenumi}{(\theenumi)}
\item For every object $i$ of $D$, the opposite category
  \sisi{}{${\mathcal{E}_i}^\circ$} of the fiber category
  $\mathcal{E}_i$ is a topos.

\item For every morphism $m : i \to j$ in $D$, the functor
  $m^{\ast}$ is left exact\nde{Cf. the preceding note. In other
  words, the family of functors
  $((m_{\ast})^\circ,(m^{\ast})^\circ)$ defines a compatible system
  of morphisms of topoi
  $\mathcal{E}_i^\circ\rightarrow \mathcal{E}_j^\circ$.}.
\end{enumerate}
\end{enonce*}

\setcounter{subsection}{2}
\setcounter{subsubsection}{3}
\subsubsection{}\etiquette[1.2.4]{Vbis.subsubsec1.2.4}
The reader will find examples of categories bifibered in duals of
topoi in \hyperref[Vbis.sec4]{paragraph 4}.

\subsubsection{}\etiquette[1.2.5]{Vbis.subsubsec1.2.5}
Let $\mathcal{E} \to B$ be a category bifibered in duals of topoi
over $B$, and let $X : D^{\circ} \to B$ be a functor. We leave it to
the reader to verify that
$(\fprod{D^{\circ}}{\mathcal{E}}{B})^{\circ}$ is a $D$-topos, which
we denote by $\bar{X}$.

\subsubsection{}\etiquette[1.2.6]{Vbis.subsubsec1.2.6}
We will often say that a functor $X : D^{\circ} \to B$ is a
$D$-\emph{object of} $B$, and we denote by $X_i$ the image under
this functor of an object $i$ of $D$; the $D$-objects of $B$ form a
category denoted by $D^{\circ}B$. If $S$ is an object of $B$, a
$D$-object of $B/S$ will be called a \emph{$D$-object of $B$
augmented toward} $S$. The datum of a $D$-object of $B$ augmented
toward $S$ is trivially equivalent to that of a functorial morphism
$X \to C^{D}_S$, where $X$ is a $D$-object of $B$ and $C^D_S$ is
the constant $D$-object of $B$ with value $S$.

When $D = \Delta$, we speak of a \emph{simplicial object} (\resp an
\emph{augmented simplicial object}).

We will also use \emph{double simplicial objects} (taking
$D = \Delta \times \Delta$).

\subsubsection{}\etiquette[1.2.7]{Vbis.subsubsec1.2.7}
Now suppose\pageoriginale that the category $B$ has finite fiber
products.

Let $f : R \to S$ be an arrow in $B$: the bifunctor
$$
([n], X) \rightsquigarrow \Hom_{\Ens} ([n], \Hom_S (X, R))
$$
from $\Delta^{\circ} \times (B / S)^{\circ}$ to $\Ens$ defines a
functor
$$
[n] \rightsquigarrow X_{[n]} = X_n
$$
by taking as $X_n$ a representative of the functor
$$
Z \rightsquigarrow \Hom_{\Ens} ([n], \Hom_S (Z, R))
$$
$(X_n \simeq \displaystyle\mathop{\overbrace{R \times_S R
    \times \ldots \times_S R}}^{n+1 \text{ times }})$.

We denote by $[R |_{f} S]$ or $[R | S]$ the semi-simplicial object
augmented toward $S$ thus constructed.

Finally, if $X$ and $X'$ are two semi-simplicial objects of $B$ (or
of $B/S$), and if $$u : X \to X'$$ is a functorial morphism, then,
for technical reasons, we introduce the object $$[X |_{u} X']$$
computed in the category $\Delta^{\circ}B$. This object can be
interpreted as a double simplicial object\nde{Explicitly, with the
usual ordering of the indices, one has
$$[[X |_u X']]_{(n,m)}=\displaystyle\mathop{\overbrace{X_m \times_{X'_m} X_m
    \times \ldots \times_{X'_m} X_m}}^{n+1 \text{ times }}$$
with the functorialities deduced from those on the factors.}
of $B$, which we then denote by $[[X |_u X']]$: indeed, there is a
canonical isomorphism of categories
$$
\Delta^{\circ} (\Delta^{\circ} B) \xrightarrow{\sim} (\Delta
\times \Delta)^{\circ} B{\quoi}.
$$

We now return to the general notion of a $D$-topos.

\subsubsection{}\etiquette[1.2.8]{Vbis.subsubsec1.2.8}
Let $E$ be a $D$-topos: we denote by $\Ysheaf{\Gamma}(E)$ the
category $\Hom_D (D, E)$. Let $f : D'\to D$ be a functor: the
category $\fprod{D'}{E}{D}$ is a $D'$-topos and, by composition with
$f$, we obtain a functor
$$
f^{\ast} : \Ysheaf{\Gamma} (E) \longrightarrow \Ysheaf{\Gamma} (\fprod{D'}{E}{D}){\quoi}.
$$

In\pageoriginale the case where $D'$ is reduced to an object $i$ of
$D$ (with the identity of $i$ as its only morphism), and where $f$
is the canonical inclusion denoted by $e_i$, we may take
$\fprod{D'}{E}{D}$ to be the fiber category $E_i$, and $e^{\ast}_i$ is then
identified with the ``\emph{evaluation at $i$}'' functor.

\begin{enonce*}{Proposition 1.2.9}\etiquette[1.2.9]{Vbis.prop1.2.9}
If $D'$ is a $\Uu$-small category, the functor $f^{\ast}$ has a
right adjoint and a left adjoint (denoted by $f_{\ast}$ and $f_{!}$,
respectively).
\end{enonce*}

This follows from a slight generalization of Kan's lemma [III 1.1],
which we will use repeatedly.

\begin{enonce*}{Lemma 1.2.10}\etiquette[1.2.10]{Vbis.lem1.2.10}
Let $I$, $J$, and $A$ be three categories over the same category
$B$. Suppose that $I$ is $\Uu$-small and that $A$ is fibered and
cofibered over $B$. Let a $B$-functor $f : I \to J$ be given, and
denote by $f^{\ast}$ the functor
$$
\Hom_B (J, A) \longrightarrow \Hom_B (I, A)
$$
defined by composition with $f$. If, in every fiber of $A$ over
$B$, $\Uu$-inductive limits (\resp projective limits) exist, then
$f^{\ast}$ has a left adjoint (\resp a right adjoint).
\end{enonce*}

\begin{enonce*}[remark]{Proof}
We give only the proof of the existence of the left adjoint, the
\resp part of the lemma following by duality. We use the following
fact, whose verification is left to the reader.
\end{enonce*}

\begin{enonce*}{Lemma 1.2.10.1}\etiquette[1.2.10.1]{Vbis.lem1.2.10.1}
Let $A$ be a bifibered category over a category $B$, let $b$ be an
object of $B$, and let $(F_{\lambda})_{\lambda \in \Lambda}$ be a
functor from a category $\Lambda$ to the fiber $A_b$. For every $G$
in $A_{b'}$, and every $u : b'\to b$ (\resp $u : b \to b'$) in $B$,
there is a bijection
$$
\Hom_u(G, \varprojlim F_{\lambda}) \xrightarrow{\sim}
\varprojlim \Hom_u (G, F_{\lambda})
$$
(\resp $\Hom_u (\varinjlim F_{\lambda}, G)
\xrightarrow{\sim}\varprojlim \Hom_u (F_{\lambda}, G)$)
whenever the first member is defined.
\end{enonce*}

\sisi{$+$}{} With this understood\pageoriginale, let
$I \amalg_f J$ be the category over $B$ defined by
$$
\ob (I \textstyle{\amalg}_f J) = \ob (I) \amalg \ob (J)
$$
\begin{equation*}
\Hom(x, y)=
\begin{cases}
\Hom_I(x, y) & \text{ if } x, y \in \ob (I)\\
\Hom_J(x, y) & \text{ if } x, y \in \ob (J)\\
\Hom_J(f(x), y) & \text{ if } x \in \ob(I) \text{ and } y \in \ob(J)\\
∅ & \text{ if } x \in \ob(J) \text{ and } y \in
\ob(I){\quoi}.\\
\end{cases}
\end{equation*}

\sisi{$+$}{} The category $\Hom_B(I \amalg_f J, A)$ is equivalent
to the category of triples consisting of a B-functor $F$ from $I$
to $A$, a $B$-functor $G$ from $J$ to $A$, and a morphism
$\varphi$ of $B$-functors from $F$ to
$G \circ f = f(G)$.

\sisi{$+$}{} For every object $j$ of $J$, denote by $I /j$ the
category of objects of $I$ ``placed by $f$ over $j$'': the objects
of $I /j$ are the pairs $(i, \alpha)$, where $i$ is an object of $I$
and $\alpha : i \to j$ is an arrow in $I \amalg_f J/j$; the
morphisms of $I/j$ are those of $I\amalg_{f}J/j$. If $p_I$ and $p_J$
are the projections of $I$ and $J$ onto $B$, and if $(i, \alpha)$ is
an object of $I/j$, denote by
$\alpha_{\ast} : A_{p_{I}(i)} \to A_{p_{J}(j)}$ the functor
$p_{J}(\alpha)_{\ast}$.

\sisi{$+$}{} To give a $B$-functor from $I \amalg_f J$ to $A$ is
equivalently to give $F \in \Hom_B (I, A)$,
$G \in \Hom_B (J, A)$, and a morphism, functorial in $j$,
$$
\psi_j : \varinjlim_{(i, \alpha)\in I/j} \alpha_{\ast}(F(i))
\longrightarrow G(j){\quoi}.
$$
(The functoriality in $j$ of the left-hand member follows from
\hyperref[Vbis.lem1.2.10.1]{1.2.10.1}.) The left adjoint of
$f^{\ast}$ is therefore given by the formula
$$
f_{!}(F) (j) = \varinjlim_{(i, \alpha)\in I/j}
\alpha_{\ast}(F(i)){\quoi}.
$$

\begin{enonce*}{Corollary 1.2.11}\etiquette[1.2.11]{Vbis.cor1.2.11}
Let\pageoriginale $E$ be a $D$-topos and let $i$ be an object of
$D$; the functor $e^{\ast}_i$ has a left adjoint (\resp a right
adjoint), given by
$$
e_{i!}(a)(j) = \coprod_{\alpha \in \Hom(i, j)} \alpha_{\ast}(a)
\left( \emph{\resp} e_{i\ast} (a) (j) = \prod_{\alpha \in \Hom(j,
  i)} \alpha^{\ast}(a)\right)
$$
where a is an object of $E$ over $i$.
\end{enonce*}

\begin{enonce*}{Proposition 1.2.12}\etiquette[1.2.12]{Vbis.prop1.2.12}
Let $D$ be a $\Uu$-small category and let $E$ be a $D$-topos. Then
the category $\Ysheaf{\Gamma}(E)$ is a $\Uu$-topos.
\end{enonce*}

\sisi{$+$}{} One verifies ``fiber by fiber,'' using
\hyperref[Vbis.lem1.2.10.1]{1.2.10.1}, that the category
$\Ysheaf{\Gamma}(E)$ has the following properties:
\begin{enumerate}
\renewcommand{\theenumi}{\alph{enumi}}
\renewcommand{\labelenumi}{\theenumi)}
\item Finite projective limits are representable.
\item Coproducts indexed by an element of $\Uu$ are representable.
  They are disjoint and universal.
\item Equivalence relations are universally effective.
\end{enumerate}

\sisi{$+$}{} It remains to show that $\Ysheaf{\Gamma}(E)$ has a
system of generators indexed by an element of $\Uu$. But if, for
every object $i$ of $D$, $(G_{i\lambda})_{\lambda \in \Lambda_{i}}$
is a system of generators of $E_i$ (where $\Lambda_i$ is a
$\Uu$-small set), then the family
$(e_{i!}(G_{i\lambda}))_{i, \lambda}$ is a system of generators of
$\Ysheaf{\Gamma}(E)$.

\setcounter{subsection}{2}
\setcounter{subsubsection}{12}

% END INCLUDED SGA4_English_current_buildable_source_20260906/source/05bis/drafts/SGA4_VBIS_LINES_206_516_English.texfrag


% BEGIN INCLUDED SGA4_English_current_buildable_source_20260906/source/05bis/drafts/SGA4_VBIS_LINES_517_829_English.texfrag
\subsubsection{}\etiquette[1.2.13]{Vbis.subsubsec1.2.13}
We now introduce the notion of a morphism between $D$-topoi. First,
let us specify that if $F$ and $F'$ are two categories over the same
category $B$ and $T:F\to F'$ is a $B$-functor, a
\emph{$B$-left adjoint} to $T$ means a functor $S:F'\to F$ left
adjoint to $T$ such that the canonical morphisms
$1\to T\circ S$ and $S\circ T\to1$ are $B$-morphisms of functors.
Under these conditions, one verifies immediately that $T$ is
cartesian and $S$ is cocartesian.

\begin{enonce*}{Definition 1.2.14}\etiquette[1.2.14]{Vbis.def1.2.14}
Let\pageoriginale $E$ and $E'$ be two $D$-topoi. A morphism from
$E$ to $E'$ is a pair of $D$-functors
$\Phi_{\ast}:E\to E'$ and $\Phi^{\ast}:E'\to E$, equipped with a
$D$-adjunction
$$
\Hom(\Phi^{\ast}(-),-)\xrightarrow[\xi]{\sim}
\Hom(-,\Phi_{\ast}(-)),
$$
such that, for every object $i$ of $D$, the pair
$(\Phi_{\ast i},\Phi_i^{\ast})$, equipped with the adjunction induced
by $\xi$, is a morphism of topoi from $E_i$ to $E'_i$.
\end{enonce*}

\begin{enonce*}{Proposition 1.2.15}\etiquette[1.2.15]{Vbis.prop1.2.15}
Let $E$ and $E'$ be two $D$-topoi, and let
$(\Phi_{\ast},\Phi^{\ast}):E\to E'$ be a morphism. If $D$ is a
$\Uu$-small category, then the pair
$(\Gamma(\Phi_{\ast}),\Gamma(\Phi^{\ast})):
\Ysheaf{\Gamma}(E)\to\Ysheaf{\Gamma}(E')$ is a morphism of topoi.
\end{enonce*}

This follows immediately from the preceding definition.

The following lemma, whose proof is left to the reader, provides a
way of constructing morphisms of $D$-topoi.

\begin{enonce*}{Lemma 1.2.16}\etiquette[1.2.16]{Vbis.lem1.2.16}
Let $E$ and $E'$ be two bifibered categories over the same category
$D$, and let $\Phi$ be a \emph{cartesian} $D$-functor from $E$ to
$E'$ such that, for every object $i$ of $D$,
$\Phi_i:E_i\to E'_i$ has a left adjoint. Then the choice, for every
$i$, of a left adjoint to $\Phi_i$ canonically determines a
$D$-functor $\Psi:E'\to E$ which is $D$-left adjoint to $\Phi$.
\end{enonce*}

\begin{enonce*}[remark]{Scholium 1.2.17}\etiquette[1.2.17]{Vbis.sch1.2.17}
Under the conditions of (\hyperref[Vbis.lem1.2.16]{1.2.16}), suppose
that $E$ and $E'$ are two $D$-topoi and that, for every object $i$
of $D$, every left adjoint to $\Phi_i$ is left exact. Then, if
$\Psi:E'\to E$ is a $D$-left adjoint to $\Phi$, the pair
$(\Phi,\Psi):E\to E'$ is a morphism of $D$-topoi.
\end{enonce*}

Now let $\mathcal E\to B$ be a category bifibered in duals of topoi
over $B$, let $X$ and $X'$ be two $D$-objects of $B$, and let
$\alpha:X\to X'$ be a functorial morphism. A choice of normalized
cleavages for $\mathcal E$ and $\mathcal E^{\circ}$ then canonically
determines a morphism
$(\alpha_{\ast},\alpha^{\ast}):\bar X\to\bar X'$ of $D$-topoi such
that, for every object\pageoriginale $i$ of $D$,
$(\alpha_{\ast i},\alpha_i^{\ast}):\bar X_i\to\bar X'_i$ is equal
to $(\alpha_{i\ast}^{\circ},\alpha_i^{\ast\circ})$, where
$\alpha_i:X_i\to X'_i$ is the arrow of $B$ given by $\alpha$. For
two different choices of cleavages for $\mathcal E$ and
$\mathcal E^{\circ}$, there is a unique $D$-isomorphism between the
resulting morphisms. (For the verification of these facts, the
reader may consult \cite{sga4.exp.Vbis.bib.4}, (1.17).)

\subsection{Ringed $D$-topoi}\etiquette[1.3]{Vbis.subsec1.3}

\begin{enonce*}{Definition 1.3.1}\etiquette[1.3.1]{Vbis.def1.3.1}
A ringed $D$-topos is a pair $(E,A)$, where $E$ is a $D$-topos and
$A$ is a ring of $\Ysheaf{\Gamma}(E)$.
\end{enonce*}

One then verifies that, for every object $i$ of $D$, $A_i$ is a ring
of the topos $E_i$, and that, for every morphism $m:i\to j$, the
canonical arrow $A_i\to m^{\ast}(A_j)$ is a ring homomorphism.

\begin{enonce*}{Definition 1.3.2}\etiquette[1.3.2]{Vbis.def1.3.2}
Let $(E,A)$ and $(E',A')$ be two ringed $D$-topoi. A morphism from
$(E,A)$ to $(E',A')$ is a pair $(\Phi,\theta)$, where
$\Phi:E\to E'$ is a morphism of $D$-topoi and
$\theta:A'\to\Gamma(\Phi_{\ast})(A)$ is a ring homomorphism.
\end{enonce*}

\begin{enonce*}[remark]{Remark 1.3.3}\etiquette[1.3.3]{Vbis.rem1.3.3}
A morphism $\Phi:(E,A)\to(E',A')$ of ringed $D$-topoi induces a
morphism
$(\Gamma(\Phi),\theta):(\Ysheaf{\Gamma}(E),A)\to
(\Ysheaf{\Gamma}(E'),A')$ of ringed $\Uu$-topoi when $D$ is a
$\Uu$-small category (\cf \hyperref[Vbis.prop1.2.15]{1.2.15}).
\end{enonce*}

\setcounter{subsubsection}{3}
\subsubsection{}\etiquette[1.3.4]{Vbis.subsubsec1.3.4}
Let $\mathcal E\to B$ be a category bifibered in duals of topoi over
$B$, and let $\mathcal O$ be a ring of
$\Ysheaf{\Gamma}(\mathcal E^{\circ})=
\Hom_{B^{\circ}}(B^{\circ},\mathcal E^{\circ})$. If
$X:D^{\circ}\to B$ is a $D$-object of $B$, the $D$-topos $\bar X$
(\cf \hyperref[Vbis.subsubsec1.2.5]{(1.2.5)}) is naturally ringed by
$(\mathcal O\mathbin{.}X^{\circ}):D\to\mathcal E$; the resulting
ringed $D$-topos is denoted by $(\bar X,\mathcal O)$. If
$\alpha:X\to X'$ is a functorial morphism, the morphism
$(\alpha_{\ast},\alpha^{\ast}):\bar X\to\bar X'$ (\cf
\hyperref[Vbis.sch1.2.17]{1.2.17}) canonically induces a morphism
$(\bar X,\mathcal O)\to(\bar X',\mathcal O)$ of ringed $D$-topoi,
again denoted by $\alpha$.

\subsubsection{}\etiquette[1.3.5]{Vbis.subsubsec1.3.5}
A ringed $D$-topos $(E,A)$ canonically defines a category
$\Ysheaf{\Mod}(E,A)$ bifibered in abelian categories over $D$, whose
fiber at an object $i$ of\pageoriginale $D$ is the category
$\Mod(E_i,A_i)$ of modules on the ringed topos $(E_i,A_i)$. With
this notation, the category of modules of $\Ysheaf{\Gamma}(E)$ over
$A$, denoted by $\Mod(\Ysheaf{\Gamma}(E),A)$, identifies with the
category $\Hom_D(D,\Ysheaf{\Mod}(E,A))$.

\subsubsection{}\etiquette[1.3.6]{Vbis.subsubsec1.3.6}
Let $\varphi=(\Phi,\theta):(E,A)\to(E',A')$ be a morphism of ringed
$D$-topoi. It defines two functors
$\varphi_{\ast}:\underline{\Mod}(E,A)\to
\underline{\Mod}(E',A')$ and
$\varphi^{\ast}:\underline{\Mod}(E',A')\to\Mod(E,A)$ between the
corresponding categories of modules:

\begin{itemize}
\item Let $M$ be an object of $\Ysheaf{\Mod}(E,A)$ over an object
$i$ of $D$. Then $\Phi_{\ast}(M)$ is a module over
$\Phi_{\ast}(A_i)$, and the morphism
$\theta_i:A'_i\to\Phi_{\ast}(A_i)$ makes it an $A'_i$-module,
denoted by $\varphi_{\ast}(M)$. The functor $\varphi_{\ast}$ is
called the \emph{direct-image functor under the morphism} $\varphi$.

\item Let $M'$ be an object of $\Ysheaf{\Mod}(E',A')$ over an object
$i$ of $D$. Then $\Phi^{\ast}(M')$ is a module over
$\Phi^{\ast}(A'_i)$, and
$\varphi^{\ast}(M')=\Phi^{\ast}(M')
\otimes_{\Phi^{\ast}(A'_i)}A_i$ is canonically endowed with an
$A_i$-module structure. By means of
(\hyperref[Vbis.lem1.2.16]{1.2.16}), this defines a functor
$\varphi^{\ast}$ left adjoint to $\varphi_{\ast}$, called the
\emph{inverse-image functor under the morphism} $\varphi$.
\end{itemize}

\paragraph{}\etiquette[1.3.6.1]{Vbis.par1.3.6.1}
We say that $\varphi$ is \emph{flat} if the functor
$\Gamma(\varphi^{\ast})$ is exact.

\begin{enonce*}{Proposition 1.3.7}\etiquette[1.3.7]{Vbis.prop1.3.7}
Let $f:D'\to D$ be a functor, and let $(E,A)$ be a ringed $D$-topos.
Then the canonical functor
$$
f^{\ast}:\Mod(\Ysheaf{\Gamma}(E),A)\longrightarrow
\Mod(\Ysheaf{\Gamma}(\fprod{E}{D'}{D}),A\circ f)
$$
has both a right and a left adjoint if $D'$ is a $\Uu$-small
category. In particular, it is exact.
\end{enonce*}

This follows immediately from (\hyperref[Vbis.lem1.2.10]{1.2.10})
and the identification
$\Hom_D(D,\Ysheaf{\Mod}(E,A))\simeq
\Mod(\Ysheaf{\Gamma}(E),A)$.

In accordance\pageoriginale with the general notation, we denote the
left (\resp right) adjoint to $f^{\ast}$ by $f_!$ (\resp
$f_{\ast}$).

\setcounter{subsubsection}{7}
\subsubsection{}\etiquette[1.3.8]{Vbis.subsubsec1.3.8}
The considerations below provide a convenient computational method
for derived functors of the form
$$
\sisi{\RR^{+}}{R}\Gamma(\varphi_{\ast}):
D^{+}(\Ysheaf{\Gamma}(E),A)\to
D^{+}(\Ysheaf{\Gamma}(E'),A'),
$$
where $\varphi:(E,A)\to(E',A')$ is a morphism of ringed $D$-topoi
(\cf (\hyperref[Vbis.subsubsec1.3.6]{1.3.6})). [If
$(S,\mathcal O_S)$ is a ringed topos, we denote by
$D^{+}(S,\mathcal O_S)$ the bounded-below derived
category\nde{If $f:(S,\mathcal O_S)\rightarrow(T,\mathcal O_T)$ is
a flat morphism of ringed topoi, we simply write $f^*$ (and not
$L^+(f^*)$) for the derived functor
$$f^*:D^+(T,\mathcal O_T)\rightarrow D^+(S,\mathcal O_S),$$
as in the original version.} of the category of
$\mathcal O_S$-modules on $S$.]

\subsubsection{}\etiquette[1.3.9]{Vbis.subsubsec1.3.9}
This formal calculation may moreover be applied in other contexts,
such as ``descent in $\ell$-adic cohomology'' (\cf \SGA 5).

\begin{enonce*}{Proposition 1.3.10}\etiquette[1.3.10]{Vbis.prop1.3.10}
Let $D$ be a $\Uu$-small category. If $(E,A)$ is a ringed
$D$-topos, denote by $I_{(E,A)}$ the class of objects of
$\Mod(\Ysheaf{\Gamma}(E),A)$ isomorphic to an object of the form
$\prod_{i\in\ob(D)}e_{i\ast}(Q_i)$, where, for every object $i$ of
$D$, $Q_i$ is totally acyclic \emph{(\cf V \hyperref[V.def4.1]{4.1})}\footnote{``Flasque''
in the terminology of V.}; $I_{(E,A)}$ has the following properties:
\begin{enumerate}
\renewcommand{\theenumi}{\roman{enumi}}
\renewcommand{\labelenumi}{(\theenumi)}
\item For every object $F$ of $I_{(E,A)}$ and every object $i$ of
$D$, $e_i^{\ast}(F)$ is totally acyclic.

\item Every object of $\Mod(\Ysheaf{\Gamma}(E),A)$ embeds into an
object of $I_{(E,A)}$.

\item $I_{(E,A)}$ is stable under finite direct sums.

\item For every morphism
$\varphi:(E,A)\to(E',A')$ of ringed $D$-topoi, the functor
$\Gamma(\varphi_{\ast})$ carries every acyclic complex in
$C^{+}(\Ysheaf{\Gamma}(E),A)$ formed of objects of $I_{(E,A)}$ to
an acyclic complex formed of objects of $I_{(E',A')}$.
\end{enumerate}
\end{enonce*}

\noindent{\underline{\emph{Proof}}}

\begin{enumerate}
\renewcommand{\theenumi}{\roman{enumi}}
\renewcommand{\labelenumi}{(\theenumi)}
\item In view of the explicit expression for $e_i^{\ast}$ (\cf
\hyperref[Vbis.cor1.2.11]{(1.2.1.1)}), it suffices to prove the following lemma:

\begin{enonce*}{Lemma 1.3.10.1}\etiquette[1.3.10.1]{Vbis.lem1.3.10.1}
Let\pageoriginale $(S,\mathcal O_S)$ be a ringed topos, and let
$(F_t)_{t\in T}$ be a family of totally acyclic
$\mathcal O_S$-modules indexed by a $\Uu$-small set $T$. Then
$\prod_{t\in T}F_t$ is totally acyclic.
\end{enonce*}

Let $X$ be an object of $S$. There is a spectral sequence
\begin{equation*}
H^p\left(X,{\mathop{\prod}_{t\in T}}^{(q)}F_t\right)Rightarrow
\mathop{\prod}_{t\in T}H^{p+q}(X,F_t)\tag{1.3.10.2}
\end{equation*}\etiquette[1.3.10.2]{Vbis.eq1.3.10.2}
where $\prod_{t\in T}^{(q)}$ denotes the $q$-th derived functor of
the ``product indexed by $T$'' functor. Since
$\prod_{t\in T}^{(q)}F_t$ is the sheaf associated with the presheaf
$U\mapsto\prod_{t\in T}H^q(U,F_t)$,
\eqhyperref[Vbis.eq1.3.10.2]{1.3.10.2} degenerates and gives
isomorphisms
\begin{equation*}
H^n\left(X,\prod_{t\in T}F_t\right)\xrightarrow{\sim}
\prod_{t\in T}H^n(X,F_t)\tag{1.3.10.3}
\end{equation*}\etiquette[1.3.10.3]{Vbis.eq1.3.10.3}
and hence, finally,
$H^n(X,\prod_{t\in T}F_t)=0$ for every $n>0$.

\item Let $F$ be an object of $\Mod(\Ysheaf{\Gamma}(E),A)$, and put
$F_i=e_i^{\ast}(F)$. For every $i$, choose a monomorphism
$F_i\to Q_i$, where $Q_i$ is
totally acyclic in $\Mod(E_i,A_i)$. Then
$e_{i\ast}(F_i)\to e_{i\ast}(Q_i)$ is a monomorphism (because
$e_{i\ast}$ has a left adjoint), and the canonical arrow
\begin{equation*}
F\to\prod_{i\in\Ob(D)}e_{i\ast}(F_i)\to
\prod_{i\in\Ob(D)}e_{i\ast}(Q_i)\tag{1.3.10.4}
\end{equation*}\etiquette[1.3.10.4]{Vbis.eq1.3.10.4}
is again a monomorphism, as is immediately verified.

\item The proof is left to the reader.

\item The reader may also verify that $\Gamma(\varphi_{\ast})$
carries an object of $I_{(E,A)}$ to an object of $I_{(E',A')}$. Once
this point is established, in view of (i), it suffices to verify the
following lemma:
\end{enumerate}

\begin{enonce*}{Lemma 1.3.10.5}\etiquette[1.3.10.5]{Vbis.lem1.3.10.5}
If\pageoriginale $0\to F\to G\to H\to0$ is a short exact sequence
in $\Mod(\Ysheaf{\Gamma}(E),A)$, with $F,G\in I_{(E,A)}$, then the
sequence
$$
0\to\Gamma(\varphi_{\ast})(F)\to
\Gamma(\varphi_{\ast})(G)\to
\Gamma(\varphi_{\ast})(H)\to0
$$
is exact in $\Mod(\Ysheaf{\Gamma}(E'),A')$.
\end{enonce*}

It suffices to observe that $e_i^{\ast}(H)$ is acyclic for every
object $i$ of $D$, and that $\Gamma(\varphi_{\ast})$ is computed
fiber by fiber. This completes the proof of
(\hyperref[Vbis.prop1.3.10]{1.3.10}).

\begin{enonce*}{Corollary 1.3.11}\etiquette[1.3.11]{Vbis.cor1.3.11}
Let $\varphi:(E,A)\to(E',A')$ be a morphism of ringed $D$-topoi.
The functor $\sisi{\RR^+}{R}\Gamma(\varphi_{\ast})$ can be computed
by means of resolutions formed of objects of $I_{(E,A)}$.
\end{enonce*}

Let $L\subset K^{+}(\Ysheaf{\Gamma}(E),A)$ be the full subcategory
of closed complexes of objects of $I_{(E,A)}$. By
((\hyperref[Vbis.prop1.3.10]{1.3.10}), (ii)), $L$ is a triangulated
subcategory, and the theorem (\hyperref[Vbis.subsec5.1]{5.1}) of
[\cite{sga4.exp.Vbis.bib.7}, Chapter 1] applies to it.

\begin{enonce*}{Corollary 1.3.12}\etiquette[1.3.12]{Vbis.cor1.3.12}
Let $\varphi:(E,A)\to(E',A')$ be a morphism of ringed $D$-topoi.
For every object $i$ of $D$, denote by
$\varphi_i:(E_i,A_i)\to(E'_i,A'_i)$ the induced morphism of ringed
topoi over $i$. Then the diagram
\[
\xymatrix@C=2.75cm@R=1.75cm{
D^{+}(\Gamma(E),A) \ar[r]^{\sisi{\RR^{+}(e_i^{\ast})}{e_i^{\ast}}}
\ar[d]_{\sisi{\RR^{+}}{R}\Gamma(\varphi_{\ast})}&
D^{+}(E_i,A_i) \ar[d]^{\sisi{\RR^{+}}{R}\varphi_{i\ast}}\\
D^{+}(\Gamma(E'),A')
\ar[r]_{\sisi{\RR^{+}(e_i^{\ast})}{e_i^{\ast}}}&D^{+}(E'_i,A'_i)
}
\]
is essentially commutative.
\end{enonce*}

There is a morphism
$$
\sisi{\RR^{+}(e_i^{\ast})}{e_i^{\ast}}\circ
\sisi{\RR^{+}}{R}\Gamma(\varphi_{\ast})\to
\sisi{\RR^{+}}{R}\varphi_{i\ast}\circ
\sisi{\RR^{+}(e_i^{\ast})}{e_i^{\ast}}
$$
which is verified to be an isomorphism by
(\hyperref[Vbis.prop1.3.10]{1.3.10}).

\begin{enonce*}[remark]{Remark 1.3.13}\etiquette[1.3.13]{Vbis.rem1.3.13}
Denote\pageoriginale by $\mathcal A_{(E,A)}$ the class of objects
$F$ of $\Mod(\Ysheaf{\Gamma}(E),A)$ such that $e_i^{\ast}(F)$ is
totally acyclic for every object $i$ of $D$. It follows from the
proof of (\hyperref[Vbis.prop1.3.10]{1.3.10}) that
$\sisi{\RR^{+}}{R}\Gamma(\varphi_{\ast})$ can be computed by means
of resolutions formed of objects of $\mathcal A_{(E,A)}$; this is
what is done, in particular, in the introduction by using the
``canonical flasque resolution''\footnote{See Exposé XVII for the
generalization of this notion.}. By
[(\hyperref[Vbis.prop1.3.10]{1.3.10})(i)], one has
$I_{(E,A)}\subset\mathcal A_{(E,A)}$, but we will use $I_{(E,A)}$
explicitly in \hyperref[Vbis.sec2]{Section 2}.
\end{enonce*}

% END INCLUDED SGA4_English_current_buildable_source_20260906/source/05bis/drafts/SGA4_VBIS_LINES_517_829_English.texfrag


% BEGIN INCLUDED SGA4_English_current_buildable_source_20260906/source/05bis/drafts/SGA4_VBIS_LINES_830_1073_English.texfrag
\section{The method of cohomological descent}\etiquette[2]{Vbis.sec2}

In this section, $D$ is a $\Uu$-small category.

\subsection{Generalities. Notation}\etiquette[2.1]{Vbis.subsec2.1}

\subsubsection{}\etiquette[2.1.1]{Vbis.subsubsec2.1.1}
Let $(S,\mathcal{O}_S)$ be a ringed topos. The category $S\times D$,
equipped with the canonical projection $S\times D\to D$, is a
$D$-topos; moreover, the constant section with value $\mathcal{O}_S$
defines a ringed $D$-topos $(S\times D,\mathcal{O}_S)$, called the
\emph{constant ringed $D$-topos with value} $(S,\mathcal{O}_S)$. With
this notation, the category
$\Mod(\Ysheaf{\Gamma}(S\times D),\mathcal{O}_S)$ is identified with
the category of covariant functors from $D$ to
$\Mod(S,\mathcal{O}_S)$.

Define an exact functor
$$
\epsilon^{\ast}:\Mod(S,\mathcal{O}_S)\to
\Mod(\Ysheaf{\Gamma}(S\times D),\mathcal{O}_S)
$$
by associating with every module $F$ on $S$ the constant functor with
value $F$. The functor $\epsilon^{\ast}$ has a right adjoint
$\epsilon_\ast$, which associates with every functor
$H:D\to\Mod(\mathcal{O}_S)$ its projective limit; the adjunction
morphism $F\to\epsilon_\ast\epsilon^{\ast}(F)$ is the morphism from
$F$ to the projective limit of the constant functor with value $F$.
Thus $\epsilon^{\ast}$ is fully faithful if and only if $D$ is
connected.

\begin{enonce*}{Definition 2.1.2}\etiquette[2.1.2]{Vbis.def2.1.2}
Let\pageoriginale $(E,A)$ be a ringed $D$-topos. An augmentation of
$(E,A)$ is a morphism (of ringed $D$-topoi) from $(E,A)$ to a constant
ringed $D$-topos. A ringed $D$-topos equipped with an augmentation will
be called an augmented ringed $D$-topos.
\end{enonce*}

\setcounter{subsubsection}{2}
\subsubsection{}\etiquette[2.1.3]{Vbis.subsubsec2.1.3}
Let $\mathcal{E}\to B$ be a category bifibered in duals of topoi, and
let $\mathcal{O}$ be a ring in
$\Ysheaf{\Gamma}(\mathcal{E}^{\circ})=
\Hom_{B^{\circ}}(B^{\circ},\mathcal{E}^{\circ})$. Let $S$ be an
object of $B$. A functor $D^{\circ}\to B/S$, that is, a functorial
morphism $X\xrightarrow{\theta}C^D_S$, where $X$ is a functor
$D^{\circ}\to B$ (\cf (\hyperref[Vbis.subsubsec1.2.6]{1.2.6})),
induces an augmentation
$$
\displaylines{\hfill
\theta:(\bar{X},\mathcal{O})\longrightarrow
(\mathcal{E}^{\circ}_S\times D,\mathcal{O}_S)
\hfill\llap{(\cf \hyperref[Vbis.subsubsec1.3.4]{(1.3.4)}){\quoi}.}}
$$
which we shall denote by the same letter as the functorial morphism
from which it arises.


\subsection{Cohomological descent}\etiquette[2.2]{Vbis.subsec2.2}

\subsubsection{}\etiquette[2.2.1]{Vbis.subsubsec2.2.1}
Let $\theta:(E,A)\to(S\times D,\mathcal{O}_S)$ be an augmented ringed
$D$-topos. With the notation of (\hyperref[Vbis.subsubsec1.3.6]{1.3.6})
and (\hyperref[Vbis.subsubsec2.1.1]{2.1.1}), set
$$
\bar{\theta}^{\ast}=\Gamma(\theta^{\ast})\circ\epsilon^{\ast}:
\Mod(S,\mathcal{O}_S)\to\Mod(\Ysheaf{\Gamma}(E),A)
$$
and
$\bar{\theta}_\ast=\epsilon_\ast\circ\Gamma(\theta_\ast):
\Mod(\Ysheaf{\Gamma}(E),A)\to\Mod(S,\mathcal{O}_S)$.

The image of $\bar{\theta}^{\ast}$ lies in the full subcategory of
$\Mod(\Ysheaf{\Gamma}(E),A)$ consisting of the cocartesian sections
of $\Ysheaf{\Gamma}(D,\Ysheaf{\Mod}(E,A))$. We shall denote this
category by
$\Ysheaf{\Gamma}^{\text{cocart}}(D,\Ysheaf{\Mod}(E,A))$
\nde{In other words, an augmented ringed $D$-topos is the datum of a
compatible system of morphisms of $S$-topoi
$$(m^\ast,m_\ast):E_j\longrightarrow E_i,$$
where $m\in\Hom_D(i,j)$. The category of cocartesian sections is the
category of descent data; its objects are families
$F_i\in\text{Ob}(E_i)$ equipped with cocycles of isomorphisms
$$F_j\isomto m^\ast F_i.$$
The functor which associates with every object $F$ on $S$ the constant
functor with value $F$ is the inverse image of a morphism of topoi
$$\epsilon:\Gamma(S\times D)\longrightarrow S,$$
inducing the morphism of ringed topoi of
\hyperref[Vbis.subsubsec2.1.3]{2.1.3}. The pair
$(\bar\theta_\ast,\bar\theta^\ast)$ is therefore simply induced by the
composite of morphisms of topoi
$\bar\theta=\epsilon\circ\Gamma(\theta):\Gamma(E)\to S$
(cf. \hyperref[Vbis.prop1.2.15]{1.2.15}).}

\begin{enonce*}{Definition 2.2.2}\etiquette[2.2.2]{Vbis.def2.2.2}
The morphism $\theta$ is said to be an \emph{effective descent
augmentation} if
$$
\sisi{}{\bar\theta^{\ast}=}
\Gamma(\theta^{\ast})\circ\epsilon^{\ast}:
\Mod(S,\mathcal{O}_S)\to
\Ysheaf{\Gamma}^{\text{\rm cocart}}(D,\Ysheaf{\Mod}(E,A))
$$
is an equivalence of categories.
\end{enonce*}

\setcounter{subsubsection}{2}
\subsubsection{}\etiquette[2.2.3]{Vbis.subsubsec2.2.3}
With\pageoriginale the notation of
(\hyperref[Vbis.subsubsec2.2.1]{2.2.1}), suppose that $\theta$ is
flat, so that $\bar\theta^{\ast}$ is exact and passes trivially to the
derived categories; \sisi{let $\LL^{+}(\bar\theta^{\ast})$}{we again
denote by $\bar\theta^{\ast}$} the functor thereby obtained. With this
notation:

\begin{enonce*}{Lemma 2.2.3.1}\etiquette[2.2.3.1]{Vbis.lem2.2.3.1}
There are two functorial morphisms
$$
\alpha:\id_{D^{+}(S,\mathcal{O}_S)}\to
\sisi{\RR^{+}}{R}\bar\theta_\ast\circ
\sisi{\LL^{+}(\bar\theta^{\ast})}{\bar\theta^{\ast}}
\quad\mbox{and}\quad
\beta:\sisi{\LL^{+}(\bar\theta^{\ast})}{\bar\theta^{\ast}}\circ
\sisi{\RR^{+}}{R}\bar\theta_\ast\to
\id_{D^{+}(\Ysheaf{\Gamma}(E),A)}
$$
which make these two functors adjoint.
\end{enonce*}

(\cf \cite{sga4.exp.Vbis.bib.9} (3.3)).

\begin{enonce*}{Definition 2.2.4}\etiquette[2.2.4]{Vbis.def2.2.4}
The morphism $\theta$ is said to be a \emph{1-cohomological-descent
augmentation} if
\begin{enumerate}
\renewcommand{\theenumi}{\arabic{enumi}}
\renewcommand{\labelenumi}{\theenumi$^{\circ}$)}
\item $\theta$ is flat;
\item
\sisi{$\LL^{+}(\bar\theta^{\ast})$}
{$\bar\theta^{\ast}:D^{+}(S,\mathcal{O}_S)\to
D^{+}(\Ysheaf{\Gamma}(E),A)$} is fully faithful.
\end{enumerate}
\end{enonce*}

\begin{enonce*}[remark]{Remark 2.2.5}\etiquette[2.2.5]{Vbis.rem2.2.5}
Condition $2^{\circ}$) of the preceding definition can also be
expressed by saying that the morphism
$\alpha\sisi{}{:\id_{D^{+}(S,\mathcal{O}_S)}\to
R\bar\theta_\ast\circ\bar\theta^{\ast}}$ in
(\hyperref[Vbis.lem2.2.3.1]{2.2.3.1}) is an isomorphism.
\sisi{}{This is what one proves in practice and what is useful.}
\end{enonce*}

\begin{enonce*}{Definition 2.2.6}\etiquette[2.2.6]{Vbis.def2.2.6}
The morphism $\theta$ is said to be a \emph{2-cohomological-descent
augmentation} (or an \emph{effective cohomological-descent
augmentation}) if $\theta$ is both an effective descent augmentation
and a 1-cohomological-descent augmentation.
\end{enonce*}

The preceding terminology is justified by the following result:

\begin{enonce*}{Proposition 2.2.7}\etiquette[2.2.7]{Vbis.prop2.2.7}
Let $\theta$ be an effective cohomological-descent augmentation. Then
the essential image of
$\sisi{\LL^{+}(\bar\theta^{\ast})}{\bar\theta^{\ast}}$ is the full
subcategory of $D^{+}(\Ysheaf{\Gamma}(E),A)$ consisting of the
complexes $F^{\bigdot}$ such that, for every $i$,
$H^{i}(F^{\bigdot})$ is a cocartesian section of
$\Ysheaf{\Mod}(E,A)$.
\end{enonce*}

For brevity, we shall call a complex $F^{\bigdot}$ satisfying the
preceding conditions a \emph{cohomological descent datum}.

\sisi{}{\begin{proof}}
It is clear that, if $K^{\bigdot}\in D^{+}(S,\mathcal{O}_S)$,
\sisi{}{then}
$\sisi{\LL^+(\bar\theta^{\ast})}
{\bar\theta^{\ast}(K^{\bigdot})}$ is a cohomological descent datum. It
is therefore enough to show that, for every cohomological descent
datum $F^{\bigdot}$, the canonical morphism
$$
\beta(F^{\bigdot}):
\sisi{\LL^{+}(\bar\theta^{\ast})}{\bar\theta^{\ast}}\circ
\sisi{\RR^{+}(\bar\theta_\ast)}{R\bar\theta_\ast}(F^{\bigdot})
\to F^{\bigdot}
$$
is\pageoriginale an isomorphism in
$D^{+}(\Ysheaf{\Gamma}(E),A)$.

\begin{enumerate}
\renewcommand{\theenumi}{\alph{enumi}}
\renewcommand{\labelenumi}{\theenumi)}
\item \emph{The case where $F^{\bigdot}$ is bounded}: we argue by
induction on the length $\ell$ of the interval of integers $i$ for
which $H^{i}(F^{\bigdot})\neq0$:
\begin{itemize}
\item for $\ell\leq1$, the assertion is true because $\theta$ is of
effective descent.

\item suppose $\ell>1$, and let $n$ be the greatest integer such that
$H^n(F^{\bigdot})\neq0$. There is a distinguished
triangle\footnote{For every complex $K^{\bigdot}$,
$K^{\bigdot}[-n]$ denotes the complex $T^{-n}(K^{\bigdot})$, where
$T$ is the translation functor.}
\[
\xymatrix@R=1.5cm{
          & H^{n}(F^{\bigdot})[-n] \ar[ld]&\\
F'^{\bigdot} \ar[rr] & & F^{\bigdot} \ar[lu]
}
\]
such that the induction hypothesis applies to $F'^{\bigdot}$. Since
$\beta(H^{n}(F^{\bigdot})[-n])$ and $\beta(F'^{\bigdot})$ are
isomorphisms, the same is true of $\beta(F^{\bigdot})$.
\end{itemize}

\item \emph{The general case}: denote by
$\sisi{\sigma}{\tau}_{\leq n}(F^{\bigdot})$ the complex
$$
\ldots\to F^{n-2}\to F^{n-1}\to\Ker d^n\to0\to0\ldots{\quoi},
$$
(\cf \cite{sga4.exp.Vbis.bib.7} (7.1)), so that, for every $n$, there is a commutative
diagram
\[
\xymatrix@R=1.5cm@C=2cm{
\bar{\theta}^{\ast}\circ R\bar{\theta}_{\ast}
(\tau_{\leq n}(F^{\bigdot}))
\ar[r]^{\sim}\ar[d]&
\tau_{\leq n}(F^{\bigdot})\ar[d]\\
\bar{\theta}^{\ast}\circ R\bar{\theta}_{\ast}
(F^{\bigdot})\ar[r]&F^{\bigdot}
}
\]
and it is enough to see that the morphism
$$
\sisi{\RR^{+}(\bar\theta_\ast)}{R\bar\theta_\ast}
(\sisi{\sigma}{\tau}_{\leq n}(F^{\bigdot}))
\to
\sisi{\RR^{+}(\bar\theta_\ast)}{R\bar\theta_\ast}(F^{\bigdot})
$$
induces\pageoriginale an isomorphism on $H^{i}$ for $i\leq n$. But
there is a triangle
\[
\xymatrix{
&F''^{\bigdot}\ar[ld]&\\
\tau_{\leq n}(F^{\bigdot})\ar[rr]&&
F^{\bigdot}\ar[lu]
}
\]
with $H^i({F''}^{\bigdot})=0$ for $i\leq n$. It follows that
$H^{i}(\sisi{\RR^{+}(\bar\theta_\ast)}{R\bar\theta_\ast}
(F''^{\bigdot}))=0$ for $i\leq n$, and then that
$H^{i}(\sisi{\RR^{+}(\bar\theta_\ast)}{R\bar\theta_\ast}
(\sisi{\sigma}{\tau}_{\leq n}(F^{\bigdot})))\to
H^{i}(\sisi{\RR^{+}(\bar\theta_\ast)}{R\bar\theta_\ast}
(F^{\bigdot}))$ is an isomorphism for $i\leq n$, which completes the
proof.
\end{enumerate}

\sisi{}{\end{proof}}

We shall now present an explicit method for computing
$$
\sisi{\RR^{+}}{R}\bar\theta_\ast
\sisi{}{=R\epsilon_\ast\circ R\Gamma(\theta)_\ast},
$$
which is very useful in proving certain criteria for cohomological
descent and in exploiting that notion.

\sisi{\subsection{A method for computing
$R^+\epsilon_\ast$}\etiquette[2.3]{Vbis.subsec2.3}}
{\subsection{A method for computing
$R\epsilon_\ast$}\etiquette[2.3]{Vbis.subsec2.3}}

We shall begin with two sorites.

% END INCLUDED SGA4_English_current_buildable_source_20260906/source/05bis/drafts/SGA4_VBIS_LINES_830_1073_English.texfrag


% BEGIN INCLUDED SGA4_English_current_buildable_source_20260906/source/05bis/drafts/SGA4_VBIS_LINES_1074_1378_English.texfrag
\subsubsection{}\etiquette[2.3.1]{Vbis.subsubsec2.3.1}
Let $D$ be a $\Uu$-small category, and let $\Ab$ be the category of
abelian groups belonging to $\Uu$. Denote by $r_D$ the functor
$D\to(\Hom(D,\Ab))^{\circ}$ defined by
$$
r_D(i)(j)=\ZZ^{\Hom(i,j)}
$$
for every pair $(i,j)$ of objects of $D$. For every pair $(i,X)$
consisting of an object of $D$ and an object of $\Hom(D,\Ab)$, one
then constructs a canonical and functorial isomorphism
$$
\Hom(r_D(i),X)\xrightarrow{\sim}X(i).
$$

\paragraph{}\etiquette[2.3.1.1]{Vbis.par2.3.1.1}
Let $A$ be an additive category (\cf \sisi{Tohokû}{Tôhoku}), and let
$F$ be a functor $D\to A$. Define a functor
$\bar F:(\Hom(D,\Ab))^{\circ}\to\Hom(A^{\circ},\Ens)$ by
$$
\bar F(X)(a)=\Hom_{\Hom(D,\Ab)}
(X,\Hom_A(a,F(.))),
$$
so\pageoriginale that the following diagram, in which $h_A$ denotes
the canonical functor, is essentially commutative:
\[
\xymatrix@C=3cm@R=1.5cm{
D \ar[r]^{r_D}\ar[d]_{F}&(\Hom(D,\Ab))^{\circ}\ar[d]^{\bar F}\\
A \ar[r]_{h_A}&\Hom(A^{\circ},\Ens).
}
\]

We leave it to the reader to verify that $\bar F$ sends finite direct
sums in $\Hom(D,\Ab)$ to products and that, if $\ZZ$ denotes the
constant functor $D\to\Ab$ with value $\ZZ$, then
$\bar F(\ZZ)=\varprojlim F$.

Finally, the correspondence $F\rightsquigarrow\bar F$ defines a
covariant functor
$$
\Hom(D,A)\to\Hom((\Hom(D,\Ab))^{\circ},
\Hom(A^{\circ},\Ens)){\quoi}.
$$

\paragraph{}\etiquette[2.3.1.2]{Vbis.par2.3.1.2}
Denote by $\Add(D)$ the full subcategory of
$(\Hom(D,\Ab))^{\circ}$ defined by the objects of the form $r_D(i)$,
where $i$ is an object of $D$, and their finite direct sums. The
category $\Add(D)$ is additive, and the functor
$r_D:D\to\Add(D)$ has the following universal property:

\paragraph{}\etiquette[2.3.1.3]{Vbis.par2.3.1.3}
\emph{For every additive category $A$, the functor
$G\longmapsto G\circ r_D$ induces an equivalence from the full
subcategory of $\Hom(\Add(D),A)$ consisting of the additive functors}
onto \emph{the category $\Hom(D,A)$.}

The category $\Add(D)$, characterized up to equivalence by
(\hyperref[Vbis.par2.3.1.3]{2.3.1.3}), will be called the
\emph{additive category generated by $D$}.

%% compilation problem!

%\nde{In more concrete terms, used implicitly below, the objects of
%$\Add(D)$ are the finite sequences
%$$[i_1]\oplus\cdots\oplus[i_n]:=(i_1,\ldots,i_n)$$
%of objects of $D$, with morphisms in $\Add(D)$ defined by the rules
%\begin{itemize}
%  \item the empty sequence is the zero object;
%  \item for all $i,j\in\text{ob}(D)$,
%      $\Hom([i],[j])=\ZZ^{\Hom_D(i,j)}$;
%  \item for all finite families $(i),(j)$ of objects of $D$,
%      $\Hom(\oplus[i],\oplus[j])=\prod\Hom([i],[j])$,
%\end{itemize}
%while the composition rules are deduced from those of $D$ by
%bilinearity. With this presentation, $r_D$ is the tautological
%functor $i\rightsquigarrow[i]$; and, if $F$ is a functor from $D$ to
%an additive category $A$, it extends uniquely to an additive functor
%$$\widetilde F:\Add(D)\longrightarrow A$$
%defined by $\widetilde F(\oplus[i])=\oplus F_i$, with the evident
%functoriality rules.}

\subsubsection{}\etiquette[2.3.2]{Vbis.subsubsec2.3.2}
Let $A$ be an abelian category, and let $K^{\bigdot\bigdot}$ be a
double complex, regarded as an object of $C^{+}(C^+(A))$, with the
first index corresponding to the first\pageoriginale symbol $C^+$.
Denote by $(K^{\bigdot\bigdot})_s$ the \emph{associated simple
complex} (\cf \sisi{Tohokû}{Tôhoku}
(\hyperref[Vbis.subsec2.4]{2.4}), whose notation we shall retain).
This defines a functor
\begin{equation*}
(\quad)_s:C^+(C^+(A))\to C^+(A)\tag{2.3.2.1}
\end{equation*}\etiquette[2.3.2.1]{Vbis.eq2.3.2.1}
and we leave it to the reader to verify the following lemma:

\begin{enonce*}{Lemma 2.3.2.2}\etiquette[2.3.2.2]{Vbis.lem2.3.2.2}
The functor $(\quad)_s$ defines a triangulated functor
$$
K^+(C^+(A))\to K^+(A)
$$
which preserves quasi-isomorphisms; it therefore defines a
triangulated functor
$$
D^+(C^+(A))\to D^+(A){\quoi},
$$
again denoted by $(\quad)_s$.
\end{enonce*}

\begin{enonce*}[remark]{Remark}
The same result holds for the functor
$(\quad)_s:C(C^b(A))\to C(A)$, since the spectral sequence under
consideration is biregular by (\EGA III (11.3.3)).
\end{enonce*}

\subsubsection{}\etiquette[2.3.3]{Vbis.subsubsec2.3.3}
Now let $(S,\mathcal{O}_S)$ be a ringed topos, let $D$ be a
$\Uu$-small category, and let
$Z^{\bigdot}\in C^{+}(\Add(D))$ be a cochain complex such that
$Z^n=0$ for $n<0$.

\sisi{Every object $F$ of
$\Mod(\Ysheaf{\Gamma}(S\times D),\mathcal{O}_S)$ (which is identified
with a functor $D\to\Mod(S,\mathcal{O}_S)$ by
(\hyperref[Vbis.subsubsec2.1.1]{2.1.1})) defines, by the universal
property of $\Add(D)$, an object
$\varepsilon_{Z^{\bigdot}\ast}(F)^{\bigdot}$ of
$C^{+}(S,\mathcal{O}_S)$ which varies functorially with $F$ by
(\hyperref[Vbis.par2.3.1.3]{2.3.1.3})}
{Every object $F$ of
$\Mod(\Ysheaf{\Gamma}(S\times D),\mathcal{O}_S)$ (which is identified
with a functor $D\to\Mod(S,\mathcal{O}_S)$ by
(\hyperref[Vbis.subsubsec2.1.1]{2.1.1})) defines, by the universal
property of $\Add(D)$ (\hyperref[Vbis.par2.3.1.3]{2.3.1.3}), an
additive functor
$$\widetilde F:\Add(D)\longrightarrow\Mod(S,\mathcal{O}_S).$$
The image complex
$$\varepsilon_{Z^{\bigdot}\ast}(F):=\widetilde F(Z^{\bigdot})$$
is an object of $C^{+}(S,\mathcal{O}_S)$ which varies functorially
with $F$. We have thus defined an additive functor
$$
\varepsilon_{Z^{\bigdot}\ast}:
\Mod(\Ysheaf{\Gamma}(S\times D),\mathcal{O}_S)
\longrightarrow C^{+}(S,\mathcal{O}_S),
$$
which is visibly exact}.

This defines a triangulated functor
\nde{$K^+$ denotes the category of bounded-below complexes up to
homotopy.}:
\begin{equation*}
K^{+}(\varepsilon_{Z^{\bigdot}\ast}):
K^+(\Ysheaf{\Gamma}(S\times D),\mathcal{O}_S)
\to K^+(C^+(S,\mathcal{O}_S)){\quoi}.\tag{2.3.3.1}
\end{equation*}\etiquette[2.3.3.1]{Vbis.eq2.3.3.1}
which preserves quasi-isomorphisms. Moreover, it follows from
(\hyperref[Vbis.lem2.3.2.2]{2.3.2.2}) that the composite functor
$$
(\quad)_s\circ K^{+}(\varepsilon_{Z^{\bigdot}\ast}):
K^+(\Ysheaf{\Gamma}(S\times D),\mathcal{O}_S)
\to K^{+}(S,\mathcal{O}_S)
$$
sends acyclic objects to acyclic objects. It therefore defines a
triangulated\pageoriginale functor
\begin{equation*}
\sisi{\RR^+}{R}\varepsilon_{Z^{\bigdot}\ast}:
D^+(\Ysheaf{\Gamma}(S\times D),\mathcal{O}_S)
\to D^{+}(S,\mathcal{O}_S)\tag{2.3.3.2}
\end{equation*}\etiquette[2.3.3.2]{Vbis.eq2.3.3.2}
satisfying the relation
\nde{In other words, $\varepsilon_{Z^{\bigdot}\ast}$ derives
trivially as
$F^{\bigdot}\rightsquigarrow
\varepsilon_{Z^{\bigdot}\ast}(F^{\bigdot})$.}
$Q\circ(\quad)_s\circ K^+(\varepsilon_{Z^{\bigdot}\ast})
=\sisi{\RR^{+}}{R}\varepsilon_{Z^{\bigdot}\ast}\circ Q$
\footnote{For every abelian category $A$, the letter $Q$ denotes the
canonical functor $K^{+}(A)\to D^{+}(A)$.}

\subsubsection{}\etiquette[2.3.4]{Vbis.subsubsec2.3.4}
Retain the preceding data\sisi{.}{ and set
$A=C^+(S,\mathcal{O}_S)$.
%%%%%%%
%%%%%%
First, $\widetilde F$ and $\bar F$ are compatible in the sense of
commutativity of the Yoneda-type diagram
$$
\xymatrix{
&A\ar[r]^h&\Hom(A^\circ,\Ens)\\
D\ar[rd]_{{r_D}}\ar[r]\ar[ru]^{F}&
\Add(D)\ar[d]_{\widetilde{r_D}}\ar[u]_{\widetilde F}\\
&\Hom(D,\Ab)\ar[ruu]_{\bar F}
}
$$
where $h$ and $r_D$ are respectively the Yoneda-type arrows defined
by
$$
h(a)=(\alpha\longmapsto\Hom_A(\alpha,a))
\quad\text{and}\quad
r_D(d)\stackrel{\hyperref[Vbis.subsubsec2.3.1]{2.3.1}}{=}
(\delta\longmapsto\Hom_D(d,\delta)),
$$
that is, $h\circ\widetilde F=\bar F\circ\widetilde{r_D}$. It follows
that
$$
h(\varepsilon_{Z^{\bigdot}\ast}(F))
=\bar F(\widetilde{r_D}(Z^{\bigdot})).
$$

By the definition of a projective limit,
$h(\epsilon_\ast F)=\varprojlim F$, and, as already observed, one has
tautologically $\bar F(\ZZ)=\varprojlim F$; hence
$$\bar F(\ZZ)=h(\epsilon_\ast(F)).$$
%%%%%%
%%%%%%%
}


Regard $Z^{\bigdot}$ as a chain complex in $\Hom(D,\Ab)$
\sisi{.}{, that is, as a complex of the abelian category
$\Hom(D,\Ab)^{\circ}$ by means of the functor $\widetilde r_D$. We
may therefore speak of exact sequences and resolutions. Explicitly, a
sequence in $\Add(D)$
$$Z\longrightarrow Z'\longrightarrow Z''$$
is exact if, for every object $j$ of $D$, the sequence
$$
\Hom_{\Add(D)}(Z'',[j])\longrightarrow
\Hom_{\Add(D)}(Z',[j])\longrightarrow
\Hom_{\Add(D)}(Z,[j])
$$
is exact in $\Ab$.}
Let $t:Z^0\to\ZZ$ be a morphism from $Z^0$ to the constant functor
with value $\ZZ$ such that the composite morphism
\begin{equation*}
Z^1\to Z^0\xrightarrow{t}\ZZ\tag{2.3.4.1}
\end{equation*}\etiquette[2.3.4.1]{Vbis.eq2.3.4.1}
is zero.

By functoriality (\cf (\hyperref[Vbis.par2.3.1.1]{2.3.1.1})), one
deduces a morphism
$\varepsilon_\ast(F)\to\varepsilon_{Z^{\bigdot}\ast}(F)^0$ such that
the composite morphism
\begin{equation*}
\varepsilon_\ast(F)\to\varepsilon_{Z^{\bigdot}\ast}(F)^0\to
\varepsilon_{Z^{\bigdot}\ast}(F)^1\tag{2.3.4.2}
\end{equation*}\etiquette[2.3.4.2]{Vbis.eq2.3.4.2}
is zero.

There is then a canonical morphism of triangulated functors
\begin{equation*}
Q\circ K^+(\varepsilon_\ast)\xrightarrow{j}
Q\circ(\quad)_s\circ K^+(\varepsilon_{Z^{\bigdot}\ast})
=\sisi{\RR^{+}}{R}\varepsilon_{Z^{\bigdot}\ast}\circ Q
\tag{2.3.4.3}
\end{equation*}\etiquette[2.3.4.3]{Vbis.eq2.3.4.3}

\begin{enonce*}{Proposition 2.3.5}\etiquette[2.3.5]{Vbis.prop2.3.5}
Retain the conditions and notation of
(\hyperref[Vbis.subsubsec2.3.3]{2.3.3}) and
{(\hyperref[Vbis.subsubsec2.3.4]{2.3.4})}; suppose, in addition, that
the augmented complex
$$
\ldots\to Z^n\to Z^{n-1}\to\ldots\to Z^1\to Z^0
\xrightarrow{t}\ZZ
$$
defines a resolution of $\ZZ$.

Then, for every complex $F^{\bigdot}$ consisting of objects of
$I_{(S\times D,\mathcal{O}_S)}$
(\cf \hyperref[Vbis.prop1.3.10]{1.3.10}),
$$
j(F^{\bigdot}):
Q\circ K^+(\varepsilon_\ast)(F^{\bigdot})\to
(\sisi{\RR^+}{R}\varepsilon_{Z^{\bigdot}\ast}\circ Q)(F^{\bigdot})
$$
is\pageoriginale an isomorphism.
\end{enonce*}

\sisi{$+$}{} Let $i$ be an object of $D$, and let $Q_i$ be a module on
$(S,\mathcal{O}_S)$. We show that $j(e_{i\ast}(Q_i))$ is an
isomorphism. It is enough to see that the sequence
\begin{equation*}
0\to Q_i\to\varepsilon_{Z^{\bigdot}\ast}(e_{i\ast}(Q_i))^0\to
\varepsilon_{Z^{\bigdot}\ast}(e_{i\ast}(Q_i))^1\to\ldots
\tag{2.3.5.1}
\end{equation*}\etiquette[2.3.5.1]{Vbis.eq2.3.5.1}
is exact.

For this, it is enough to show that, for every object $X$ of
$\Mod(S,\mathcal{O}_S)$, the sequence
\begin{equation*}
0\to\Hom(X,Q_i)\to
\Hom(X,\varepsilon_{Z^{\bigdot}\ast}(e_{i\ast}(Q_i))^0)\to
\Hom(X,\varepsilon_{Z^{\bigdot}\ast}(e_{i\ast}(Q_i))^1)\to\ldots
\tag{2.3.5.2}
\end{equation*}\etiquette[2.3.5.2]{Vbis.eq2.3.5.2}
is exact. An immediate calculation shows that
\eqhyperref[Vbis.eq2.3.5.2]{2.3.5.2} reduces to
\hypertarget{sga4.exp.Vbis.equation.2.3.5.3}{}\begin{gather*}
0\to\Hom(\ZZ,\Hom(X,Q_i))\to
\Hom(\Hom_{\Add(D)}(Z^0,[i]),\Hom(X,Q_i))\to\tag{2.3.5.3}\\
\Hom(\Hom_{\Add(D)}(Z^1,[i]),\Hom(X,Q_i))\to\ldots{\quoi}.
\end{gather*}
\sisi{}{By hypothesis, the sequence
$$
0\leftarrow\ZZ\leftarrow\Hom_{\Add(D)}(Z^0,[i])
\leftarrow\Hom_{\Add(D)}(Z^1,[i])\leftarrow\ldots
$$
is an exact sequence of free modules and is therefore null-homotopic.
It follows that \hyperlink{sga4.exp.Vbis.equation.2.3.5.3}{(2.3.5.3)} is exact.}

\sisi{$+$}{} If, for every object $i$ of $D$, $Q_i$ is a totally
acyclic module, then, by applying the functor
$\prod_{i\in\Ob(D)}$ to the complexes
\eqhyperref[Vbis.eq2.3.5.1]{2.3.5.1}, one sees that the resulting
complex is still acyclic by (\hyperref[Vbis.lem1.3.10.1]{1.3.10.1}).
Thus $j(\prod_{i\in\Ob(D)}e_{i\ast}(Q_i))$ is an isomorphism.

\sisi{$+$}{} We leave it to the reader to deduce that
$j(F^{\bigdot})$ is an isomorphism when
$F^n\in I_{(S\times D,\mathcal{O}_S)}$ for every $n$ (for example,
by spectral sequences).

\begin{enonce*}{Corollary 2.3.6}\etiquette[2.3.6]{Vbis.cor2.3.6}
Under\pageoriginale the conditions of
(\hyperref[Vbis.prop2.3.5]{2.3.5}), the morphism $j$ identifies
$\sisi{\RR^+}{R}\varepsilon_{Z^{\bigdot}\ast}$ with the derived
functor of $\varepsilon_\ast$.
\end{enonce*}

The subcategory of
$K^{+}(\Ysheaf{\Gamma}(S\times D),\mathcal{O}_S)$ defined by the
complexes consisting of objects of
$I_{(S\times D,\mathcal{O}_S)}$ satisfies the conditions of theorem
(\hyperref[Vbis.subsec5.1]{5.1}) of [\cite{sga4.exp.Vbis.bib.7}, I] for the functor
$K^{+}(\varepsilon_\ast)$, by
(\hyperref[Vbis.prop1.3.10]{1.3.10}) and
(\hyperref[Vbis.prop2.3.5]{2.3.5}). The corollary follows immediately
from this observation.

\begin{enonce*}{Corollary 2.3.7}\etiquette[2.3.7]{Vbis.cor2.3.7}
Under the conditions of (\hyperref[Vbis.prop2.3.5]{2.3.5}), let $(E,A)$
be a ringed $D$-topos and let
$\theta:(E,A)\to(S\times D,\mathcal{O}_S)$ be an augmentation. There
is a canonical isomorphism
$$
\sisi{\RR^+}{R}\bar\theta_\ast\xrightarrow{\sim}
\sisi{\RR^+}{R}\varepsilon_{Z^{\bigdot}\ast}\circ
\sisi{\RR^+}{R}\Gamma(\theta_\ast){\quoi}.
$$
\end{enonce*}

This follows from [\cite{sga4.exp.Vbis.bib.7}, I, 5.4].

\begin{enonce*}[remark]{Remark 2.3.8}\etiquette[2.3.8]{Vbis.rem2.3.8}
We knew a priori that $\sisi{\RR^+}{R}\varepsilon_\ast$ exists and
that there is an isomorphism
$\sisi{\RR^+}{R}\bar\theta_\ast\xrightarrow{\sim}
\sisi{\RR^+}{R}\varepsilon_\ast\circ
\sisi{\RR^+}{R}\Gamma(\theta_\ast)$, since
$\Gamma(\theta_\ast)$ and $\varepsilon_\ast$ are induced by morphisms
of ringed topoi (\cf (\hyperref[Vbis.subsubsec1.3.6]{1.3.6}) and
(\hyperref[Vbis.subsubsec2.1.1]{2.1.1})). Nevertheless, the preceding
calculation will prove very useful and applicable in other contexts,
since in practice the conditions of
(\hyperref[Vbis.prop2.3.5]{2.3.5}) will always be satisfied.
\end{enonce*}

We shall now give the fundamental examples to which
(\hyperref[Vbis.prop2.3.5]{2.3.5}) can be applied.

\setcounter{subsubsection}{8}

% END INCLUDED SGA4_English_current_buildable_source_20260906/source/05bis/drafts/SGA4_VBIS_LINES_1074_1378_English.texfrag


% BEGIN INCLUDED SGA4_English_current_buildable_source_20260906/source/05bis/drafts/SGA4_VBIS_LINES_1379_1421_English.texfrag
\subsubsection{}\etiquette[2.3.9]{Vbis.subsubsec2.3.9}
In $\Add(\Delta)$, put $Z^n=[n]$ and define
$d^n:Z^n\to Z^{n+1}$ by the formula
\begin{equation*}
d^n=\sum_{i=0}^{n+1}(-1)^i\partial^i.\tag{2.3.9.1}
\end{equation*}\etiquette[2.3.9.1]{Vbis.eq2.3.9.1}
Here $\partial^i:[n]\to[n+1]$ is the $i$th face
[\cf \cite{sga4.exp.Vbis.bib.3} II 2]. For $t:Z^0\to\ZZ$, take the evident natural
augmentation\nde{Explicitly, if $F$ is a simplicial complex with
differential $\mathrm d_F$, then $\varepsilon_{Z^{\bigdot}*}(F)$ is
the simple complex associated with the double complex $F^j(i)$ whose
horizontal differential is the simplicial differential $\mathrm d^i$
and whose vertical differential is $\mathrm d_F^j$.}.

The condition of \hyperref[Vbis.subsubsec2.3.4]{2.3.4} is trivially
satisfied. The condition of \hyperref[Vbis.prop2.3.5]{2.3.5} follows
from \cite{sga4.exp.Vbis.bib.5} I (3.7.4).

\subsubsection{}\etiquette[2.3.10]{Vbis.subsubsec2.3.10}
In\pageoriginale $\Add(\Delta\times\Delta)$, consider the simple
complex associated with the double complex
\hypertarget{sga4.exp.Vbis.equation.2.3.10.1}{}\[\vcenter{
\xymatrix@C=2cm@R=1.75cm{
\big([n],[m+1]\big)  \ar[r]^{s_1^{n,m+1}} & \big([n+1],[m+1]\big)\\
\big([n],[m]\big) \ar[r]^{s_1^{n,m}} \ar[u]^{s_2^{n,m}} &
\big([n+1],[m]\big)\ar[u]_{s_2^{n+1,m}}
}}\leqno{(2.3.10.1)}
\]
with
\begin{equation*}
\tag{2.3.10.2}
\begin{cases}
S^{n,m}_1=\displaystyle\mathop{\sum}_{i=0}^{n+1}(-1)^i
(\partial^i,\id_{[m]})\,,\\[15pt]
S^{n,m}_2=\displaystyle\mathop{\sum}_{j=0}^{m+1}(-1)^j
(\id_{[n]},\partial^j){\quoi}.
\end{cases}
\end{equation*}\etiquette[2.3.10.2]{Vbis.eq2.3.10.2}

The reader may verify (using \hyperref[Vbis.subsubsec2.3.9]{2.3.9}) that the complex so defined,
equipped with the evident natural augmentation, satisfies the
conditions of \hyperref[Vbis.prop2.3.5]{2.3.5}.

\subsubsection{}\etiquette[2.3.11]{Vbis.subsubsec2.3.11}
The multisimplicial case is treated similarly (with
$\underbrace{\Delta\times\Delta\times\cdots\times\Delta}_{p\text{ times}}$).


% END INCLUDED SGA4_English_current_buildable_source_20260906/source/05bis/drafts/SGA4_VBIS_LINES_1379_1421_English.texfrag


% BEGIN INCLUDED SGA4_English_current_buildable_source_20260906/source/05bis/drafts/SGA4_VBIS_LINES_1422_1690_English.texfrag
\subsection{Relative cohomological
  descent}\etiquette[2.4]{Vbis.subsec2.4}\pageoriginale

\setcounter{subsubsection}{-1}
\subsubsection{}\etiquette[2.4.0]{Vbis.subsubsec2.4.0}
Let $\Ee\to B$ be a category bifibered in duals of topoi, and let
$\mathcal O$ be a ring of $\Ysheaf{\Gamma}(\Ee^{\circ})$.

These data canonically define a category fibered and cofibered in
abelian categories over $B^{\circ}$, denoted by
$\Ysheaf{\Mod}(\Ee^{\circ},\mathcal O)$ (\cf
(\hyperref[Vbis.rem1.3.3]{1.3.3}) and
(\hyperref[Vbis.subsubsec1.3.4]{1.3.4})). Throughout what follows,
the ``direct-image'' and ``inverse-image'' functors will always be
taken with respect to the fibered and cofibered functor
$\Ysheaf{\Mod}(\mathcal E^{\circ},\mathcal O)\to B^{\circ}$.

\begin{enonce*}{Definition 2.4.0.1}\etiquette[2.4.0.1]{Vbis.def2.4.0.1}
A morphism $f:T\to S$ in $B$ is said to be \emph{flat} (relative to
$(\mathcal E,\mathcal O)$) if
$f_{\ast}:\Mod(\mathcal E^{\circ}_S,\mathcal O_S)\to
\Mod(\mathcal E^{\circ}_T,\mathcal O_T)$ is an exact
functor\nde{This notation may cause confusion: one should probably
write $(f^\circ)_\star$, which formally coincides with
$(f^\star)^\circ$, where $f^\star$ is the fibered functor associated
with $\Ee\rightarrow B$. For example, for the category of sheaves on
topological spaces fibered over the category of topological spaces
(cf. \hyperref[Vbis.subsec4.1]{4.1}), one has
$(f^\circ)_\star(S,F)=(T,f^{-1}F)=(f^\star)^\circ(S,F)$, where
$f:T\rightarrow S$ is a continuous map of topological spaces, $F$
is a sheaf on $S$, and $f^{-1}$ is the usual inverse image. One then
recognizes the usual flatness conditions.}.
\end{enonce*}

\subsubsection{}\etiquette[2.4.1]{Vbis.subsubsec2.4.1}
Let $G$ be a fibered and cofibered subcategory of
$\Ysheaf{\Mod}(\mathcal E^{\circ},\mathcal O)$ such that, for every
object $S$ of $B$, $G_S$ is a \emph{thick} subcategory (\cf
\sisi{Tohokû}{Tôhoku} (1.11)) of
$(\Ysheaf{\Mod}(\mathcal E^{\circ},\mathcal O))_S$. The reader will
find examples of such situations in \hyperref[Vbis.sec4]{Section 4}.

If $S$ is an object of $B$, denote by $D^+(S)$ the derived category
of $\Mod(\mathcal E^{\circ}_S,\mathcal O_S)$. Following
[\cite{sga4.exp.Vbis.bib.7} (I \S\ 4)], introduce the category $D^+_{G_S}(S)$, which
identifies with the full subcategory of $D^+(S)$ formed by the
objects $X^{\bigdot}$ such that $H^n(X^{\bigdot})\in G_S$ for every
$n$.

\emph{We assume throughout that the following condition holds.}

\paragraph{}\etiquette[2.4.1.1]{Vbis.par2.4.1.1}
\emph{For every morphism $h:T\to S$ in $B$, the functor
$\sisi{\RR^+}{Rh^{\ast}}:D^+(T)\to D^+(S)$ has the property that
$\sisi{\RR^+}{Rh^{\ast}}(F)\in D^+_{G_S}(S)$ whenever
$F\in G_T$}\nde{As above, it would be more correct to write
$R(h^\circ)^\ast=R(h_\ast)^\circ$ rather than $Rh^\ast$, which would
make the notation more consistent with the usual geometric
intuition. One could equally well have taken the fibered and
cofibered functors with respect to
$\Mod(\mathcal E,\mathcal O)\rightarrow B$ rather than the opposite
categories. The fibers, however, would then have been only duals of
topoi.}.

By [\cite{sga4.exp.Vbis.bib.7} (I. (7.3). (ii))], this condition implies that
$\sisi{\RR^+}{Rh^{\ast}}(D^+_{G_T}(T))\subset D^+_{G_S}(S)$.

\subsubsection{}\etiquette[2.4.2]{Vbis.susubsec2.4.2}
Let\pageoriginale $D$ be a $\Uu$-small category. \emph{Throughout
what follows, we suppose given a complex $Z^{\bigdot}$ of
$\Add(D)$ satisfying the conditions of}
(\hyperref[Vbis.prop2.3.5]{2.3.5}) (in applications, one will in fact
have $D=\Delta$ or $D=\Delta\times\Delta$).

If $X:D^{\circ}\to B$ is a $D$-object of $B$, denote by
$\Mod_G(\Ysheaf{\Gamma}(\bar X),\mathcal O)$ the full subcategory of
$\Mod(\Ysheaf{\Gamma}(\bar X),\mathcal O)$ formed by the objects $F$
such that $e_i^{\ast}(F)\in G_i$ $(=G_{X_i})$ for every object $i$
of $D$. The category $\Mod_G(\Ysheaf{\Gamma}(\bar X),\mathcal O)$
is a thick subcategory of
$\Mod(\Ysheaf{\Gamma}(\bar X),\mathcal O)$. Denote by
$K_G^+(\Ysheaf{\Gamma}(\bar X),\Oo)$ (\resp
$D_G^+(\Ysheaf{\Gamma}(\bar X),\Oo)$) the full subcategory of
$K^+(\Ysheaf{\Gamma}(\bar X),\Oo)$ (\resp
$D^+(\Ysheaf{\Gamma}(\bar X),\Oo)$) formed by the objects
$X^{\bigdot}$ such that
$H^n(X^{\bigdot})\in\Mod_G(\Ysheaf{\Gamma}(\bar X),\Oo)$ for every
$n$.

\subsubsection{}\etiquette[2.4.3]{Vbis.subsubsec2.4.3}
Let $X$ and $X'$ be two $D$-objects of $B$, and let
$\alpha:X\to X'$ be a functorial morphism. Recall (\cf
(\hyperref[Vbis.subsubsec1.3.4]{1.3.4})) that the corresponding
morphism of ringed $D$-topoi is again denoted by
$\alpha:(\bar X,\mathcal O)\to(\bar X',\mathcal O)$.

\begin{enonce*}{Definition 2.4.3.1}\etiquette[2.4.3.1]{Vbis.def2.4.3.1}
The morphism $\alpha:X\to X'$ is said to be \emph{flat} if the
corresponding morphism of ringed $D$-topoi is flat in the sense of
(\hyperref[Vbis.par1.3.6.1]{1.3.6.1}).
\end{enonce*}

We then have the evident lemma:

\begin{enonce*}{Lemma 2.4.3.2}\etiquette[2.4.3.2]{Vbis.lem2.4.3.2}
The morphism $\alpha:X\to X'$ is flat if and only if, for every
object $i$ of $D$, $\alpha_i:X_i\to X'_i$ is a flat morphism
\emph{(\cf (\hyperref[Vbis.def2.4.0.1]{2.4.0.1}))}.
\end{enonce*}

\subsubsection{}\etiquette[2.4.4]{Vbis.subsubsec2.4.4}
A morphism $\alpha:X\to X'$ defines a functor
\begin{equation*}
K^+(\Gamma(\alpha_*)):K^+(\Ysheaf{\Gamma}(\bar X),\Oo)
\longrightarrow K^+(\Ysheaf{\Gamma}(\bar X'),\Oo){\quoi}.
\tag{2.4.4.1}
\end{equation*}\etiquette[2.4.4.1]{Vbis.eq2.4.4.1}

\begin{enonce*}{Lemma 2.4.4.2}\etiquette[2.4.4.2]{Vbis.lem2.4.4.2}
The functor
$$
K^+(\Gamma(\alpha_*))|_{K_G^+(\Ysheaf{\Gamma}(\bar X),\Oo)}:
K_G^+(\Ysheaf{\Gamma}(\bar X),\Oo)\longrightarrow
K^+(\Ysheaf{\Gamma}(\bar X'),\Oo)
$$
has a right derived functor, denoted by
$\sisi{\RR^+}{R}_G\Gamma(\alpha_{\ast})$, and the canonical
morphism
$$
\sisi{\RR^+}{R}_G\Gamma(\alpha_{\ast})\longrightarrow
\sisi{\RR^+}{R}\Gamma(\alpha_{\ast})|_{D_G^+
(\Ysheaf{\Gamma}(\bar X),\Oo)}
$$
is\pageoriginale an isomorphism.

Moreover,
$$
\sisi{\RR^+}{R}_G\Gamma(\alpha_{\ast})
(D_G^+(\Ysheaf{\Gamma}(\bar X),\Oo))
\subset D_G^+(\Ysheaf{\Gamma}(\bar X'),\Oo).
$$
\end{enonce*}

The first part follows from (\hyperref[Vbis.prop1.3.10]{1.3.10}) and
[\cite{sga4.exp.Vbis.bib.7}, (I. (5.2))]. To verify the last assertion, by
[\cite{sga4.exp.Vbis.bib.7}, (I. (7.3). (ii))], it suffices to show that
$\sisi{\RR^+}{R}\Gamma(\alpha_{\ast})(F)\in
D_G^+(\Ysheaf{\Gamma}(\bar X'),\Oo)$ whenever
$F\in\Mod_G(\Ysheaf{\Gamma}(\bar X),\Oo)$. This follows from the
explicit calculation of
$\sisi{\RR^+}{R}\Gamma(\alpha_\ast)$ given by \hyperref[Vbis.cor1.3.11]{(1.3.1.1)} and from
hypothesis (\hyperref[Vbis.par2.4.1.1]{2.4.1.1}).

\subsubsection{}\etiquette[2.4.5]{Vbis.subsubsec2.4.5}
Let $S$ be an object of $B$. The ringed $D$-topos
$(\bar{C^D_S},\mathcal O)$ (\cf
(\hyperref[Vbis.subsubsec1.2.6]{1.2.6})) identifies with the constant
ringed $D$-topos
$(\mathcal E^{\circ}_S\times D,\mathcal O_S)$, and we have the
following lemma, which follows from [\cite{sga4.exp.Vbis.bib.7} (I. (7.3). (ii))]:

\begin{enonce*}{Lemma 2.4.5.1}\etiquette[2.4.5.1]{Vbis.lem2.4.5.1}
The functor
$\sisi{\RR^+}{R}\varepsilon_{Z^{\bigdot}\ast}|_{
D_G^+(\Ysheaf{\Gamma}(\bar{C^D_S}),\mathcal O)}$, denoted by
$\sisi{\RR^+}{R}_G\varepsilon_{Z^{\bigdot}\ast}$, takes values in
$D^+_{G_S}(S)$ and identifies with a right derived functor of
$$
K^+(\varepsilon_{\ast})|_{K_G^+
(\Ysheaf{\Gamma}(\bar{C^D_S}),\mathcal O)}:
K_G^+(\Ysheaf{\Gamma}(\bar{C^D_S}),\mathcal O)
\longrightarrow K^+(S){\quoi}.
$$
\end{enonce*}

\begin{enonce*}{Proposition 2.4.6}\etiquette[2.4.6]{Vbis.prop2.4.6}
Let $\theta:X\to C^D_S$ be an augmented $D$-object of $B$. The
functor
$K^+(\bar\theta_{\ast})|_{K_G^+
(\Ysheaf{\Gamma}(\bar X),\mathcal O)}$ has a right derived functor,
denoted by $\sisi{\RR^+}{R}_G\bar\theta_{\ast}$, taking values in
$D^+_{G_S}(S)$. Moreover, there is a canonical isomorphism
$$
\sisi{\RR^+}{R}_G\bar\theta_{\ast}\xrightarrow{\sim}
\sisi{\RR^+}{R}_G\varepsilon_{Z^{\bigdot}\ast}\circ
\sisi{\RR^+}{R}_G\Gamma(\theta_{\ast}){\quoi}.
$$
\end{enonce*}

The proof, from what precedes, is left to the reader.

We are now in a position to give the definitions concerning
cohomological descent.

\setcounter{subsubsection}{6}
\subsubsection{}\etiquette[2.4.7]{Vbis.subsubsec2.4.7}
Let\pageoriginale $\theta:X\to C^D_S$ be an augmented $D$-object of
$B$. Suppose that $\theta$ is flat, so that the restriction to
$D^+_{G_S}(S)$ of the functor $\bar\theta^{\ast}$ takes values in
$D_G^+(\Ysheaf{\Gamma}(\bar X),\mathcal O)$; denote this restriction
by $\bar\theta_G^{\ast}$. With this notation:

\begin{enonce*}{Lemma 2.4.7.1}\etiquette[2.4.7.1]{Vbis.lem2.4.7.1}
There are two functorial morphisms
$\alpha:\id_{D^+_{G_S}(S)}\to
\sisi{\RR^+}{R}_G\bar\theta_{\ast}\circ\bar\theta_G^{\ast}$ and
$\beta:\bar\theta_G^{\ast}\circ
\sisi{\RR^+}{R}_G\bar\theta_{\ast}\to
\id_{D_G^+(\Ysheaf{\Gamma}(\bar X),\mathcal O)}$ which make these
two functors adjoint.
\end{enonce*}

This follows immediately from
(\hyperref[Vbis.lem2.2.3.1]{2.2.3.1}) and
(\hyperref[Vbis.prop2.4.6]{2.4.6}).

\begin{enonce*}{Definition 2.4.8}\etiquette[2.4.8]{Vbis.def2.4.8}
The morphism $\theta$ is called an augmentation of cohomological
$1$-descent relative to $G$ if
\begin{enumerate}
\renewcommand{\theenumi}{\arabic{enumi}}
\renewcommand{\labelenumi}{\theenumi$^{\circ}$)}
\item $\theta$ is flat.

\item The functor $\bar\theta_G^{\ast}$ is fully faithful.
\end{enumerate}
\end{enonce*}

\begin{enonce*}{Remark 2.4.9}\etiquette[2.4.9]{Vbis.rem2.4.9}
Condition $2^{\circ})$ of the preceding definition may equivalently
be expressed by saying that the morphism $\alpha$ in
(\hyperref[Vbis.lem2.4.7.1]{2.4.7.1}) is an isomorphism.
\end{enonce*}

\begin{enonce*}{Definition 2.4.10}\etiquette[2.4.10]{Vbis.def2.4.10}
Let $\theta:X\to C^D_S$ be an augmented $D$-object of $B$. The
morphism $\theta$ is called an \emph{augmentation of effective
descent relative to} $G$ if the functor
$$
\Gamma(\theta^{\ast})\circ\varepsilon^{\ast}:G_S\longrightarrow
\Ysheaf{\Gamma}^{\Cocart}(D,G)
$$
is an equivalence of categories.
\end{enonce*}

\begin{enonce*}{Definition 2.4.11}\etiquette[2.4.11]{Vbis.def2.4.11}
The morphism $\theta$ is called an \emph{augmentation of
cohomological $2$-descent} (or of \emph{effective cohomological
descent}) \emph{relative to} $G$ if $\theta$ is both an augmentation
of cohomological $1$-descent and an augmentation of effective
descent relative to $G$.
\end{enonce*}

The proof of (\hyperref[Vbis.prop2.2.7]{2.2.7}) then gives the
following result:

\begin{enonce*}{Theorem 2.4.12}\etiquette[2.4.12]{Vbis.thm2.4.12}
Let\pageoriginale $\theta:X\to C^D_S$ be an augmentation of
cohomological $2$-descent relative to $G$. Then the essential image
of $\bar\theta_G^{\ast}$ is the full subcategory of
$D_G^+(\Ysheaf{\Gamma}(\bar X),\mathcal O)$ formed by the complexes
$F^{\bigdot}$ such that, for every $n$, $H^n(F^{\bigdot})$ is a
cocartesian section of $G$.
\end{enonce*}

The following proposition, whose verification is left to the reader,
translates Definition (\hyperref[Vbis.def2.4.10]{2.4.10}) into the
language of (\cite{sga4.exp.Vbis.bib.4}, 6):

\begin{enonce*}{Proposition 2.4.13}\etiquette[2.4.13]{Vbis.prop2.4.13}
The morphism $\theta:X\to C^D_S$ is an augmentation of effective
descent relative to $G$ if and only if the functor
$D^{\circ}\to B/S$ over $B$ defined by $\theta$ induces an
equivalence between the category of cartesian sections of
$G^{\circ}$ over $B/S$ and the category of cartesian sections of
$G^{\circ}$ over $D^{\circ}$.
\end{enonce*}

\begin{enonce*}{Corollary 2.4.14}\etiquette[2.4.14]{Vbis.cor2.4.14}
Suppose that finite fiber products are representable in $B$, and let
$f:R\to S$ be a morphism. The augmented semi-simplicial object
$[R|_fS]$ \emph{(\cf
(\hyperref[Vbis.subsubsec1.2.7]{1.2.7}))} is of effective descent
relative to $G$ if and only if $f$ is an effective descent morphism
for the fibered category $G^{\circ}$ over $B$.
\end{enonce*}

This follows from \hyperref[Vbis.prop2.4.13]{(2.4.13)} and (\cite{sga4.exp.Vbis.bib.4}, 9).

\begin{enonce*}{Definition 2.4.15}\etiquette[2.4.15]{Vbis.def2.4.15}
Suppose that finite fiber products are representable in $B$, and let
$f:R\to S$ be a morphism. We say that $f$ is a morphism of
cohomological $i$-descent relative to $G$ $(i=1,2)$ if the augmented
semi-simplicial object $[R|_fS]$ is of cohomological $i$-descent
relative to $G$.
\end{enonce*}

\begin{enonce*}{Definition 2.4.16}\etiquette[2.4.16]{Vbis.def2.4.16}
Suppose that finite fiber products are representable in $B$, and let
$\theta:X\to C^D_S$ be an augmentation. We say that $\theta$ is an
\emph{augmentation of universal cohomological $i$-descent
$(i=1,2)$ relative to $G$} if, for every morphism $T\to S$, the
augmentation $X_T\xrightarrow{\theta_T}C^D_T$ obtained by base
change\pageoriginale is an augmentation of cohomological
$i$-descent relative to $G$.
\end{enonce*}

The notion of a \emph{morphism of universal cohomological
$i$-descent relative to} $G$ follows immediately from the two
preceding definitions.

We will often use the terminology ``augmentation of
$G$-cohomological $i$-descent'' in place of ``augmentation of
cohomological $i$-descent relative to $G$''.

\subsection{The descent spectral sequence}\etiquette[2.5]{Vbis.subsec2.5}


% END INCLUDED SGA4_English_current_buildable_source_20260906/source/05bis/drafts/SGA4_VBIS_LINES_1422_1690_English.texfrag


% BEGIN INCLUDED SGA4_English_current_buildable_source_20260906/source/05bis/drafts/SGA4_VBIS_LINES_1691_1929_English.texfrag

\subsubsection{}\etiquette[2.5.1]{Vbis.subsubsec2.5.1}
Let $A$ be an abelian category. Denote by $F(A)$ the category whose
objects are the objects of $A$ equipped with a discrete and codiscrete
filtration, and whose morphisms are the filtered morphisms of $A$:
the category $F(A)$ is an additive category (but is not abelian).

The category $K(F(A))$ of filtered complexes of $A$, whose filtration
is discrete and codiscrete degree by degree, is a triangulated
category, and the canonical functors
\begin{equation*}
\Gr_n : K(F(A)) \longrightarrow K(A)\tag{2.5.1.1}
\end{equation*}\etiquette[2.5.1.1]{Vbis.eq2.5.1.1}
are triangulated. The set $\Sigma$ of morphisms $f$ of $K(F(A))$
such that $\Gr_n(f)$ is a quasi-isomorphism for every $n$ is therefore
a saturated multiplicative system [\cf \cite{sga4.exp.Vbis.bib.9}. \S~2. ${\rm n^\circ}$ 1].

By inverting the arrows of this system, one obtains a new
triangulated category, denoted by $D F(A)$, and the functors $\Gr_n$
extend to functors
\begin{equation*}
\Gr_n : D F(A) \longrightarrow D(A){\quoi}.\tag{2.5.1.2}
\end{equation*}\etiquette[2.5.1.2]{Vbis.eq2.5.1.2}

In what follows, we will consider only complexes bounded below: this
naturally gives rise to the notation $K^+(F(A))$ and $D^+ F(A)$.

\subsubsection{}\etiquette[2.5.2]{Vbis.subsubsec2.5.2}
Let\pageoriginale $B$ be a thick subcategory of $A$, and denote by
$K^+_B(F(A))$ the full subcategory of $K^+(F(A))$ whose objects are
the complexes $X^{\bigdot}$ such that $H^{i}(\Gr_j(X^{\bigdot})) \in B$
for every pair of integers $(i, j): K_B^+(F(A))$ is a
\emph{localizing triangulated subcategory} of $K(F(A))$ [\cf
  \cite{sga4.exp.Vbis.bib.7}. (I. \S\ 5)]. We denote by $D_B^+ F(A)$ the full subcategory
of $D F(A)$ whose objects are the complexes bounded below
$X^{\bigdot}$ such that $H^i(
\Gr_j(X^{\bigdot})) \in B$ for every pair $(i, j)$: by (\loccit),
$D^+_B F(A)$ identifies with the category of fractions
$K^+_B(F(A))_{\Sigma \cap K^{+}_{B} (F(A))}$.

The ``forgetful functor for filtrations''
$$
\iota : K^+(F(A)) \longrightarrow K^+(A)
$$
is triangulated. Moreover, it follows from [\cite{sga4.exp.Vbis.bib.5}. I. (4.7)] that
$\iota(f)$ is a quasi-isomorphism if $f \in \Sigma$ and that
$\iota(K^+_B(F(A))) \subset K^+_B(A)$. The functor $\iota$ thus
extends to a triangulated functor
\begin{equation*}
\iota : D^+ F(A) \longrightarrow D^+ (A) \tag{2.5.2.1}
\end{equation*}\etiquette[2.5.2.1]{Vbis.eq2.5.2.1}
such that $\iota (D^+_B F(A)) \subset D^+_B (A)$.

Finally, note that there is a canonical \emph{spectral functor}
$D^+_B F(A) \to B$, converging to $H^n \circ \iota$, whose $E_1$ term
is given by $H^{p+q} \circ \Gr_{p}$ [\cf \cite{sga4.exp.Vbis.bib.5} I. (4.2)].

\emph{We now return to the notation of} (2.3), \emph{assuming that}
$D = \Delta$: $Z^{\bigdot}$ will denote the complex of $\Add(\Delta)$
defined by \eqhyperref[Vbis.eq2.3.9.1]{2.3.9.1}.

\begin{enonce*}{Proposition 2.5.3}\etiquette[2.5.3]{Ii.prop2.5.3}
Let $(S, \mathcal{O}_S)$ be a ringed topos. The functor
$$\sisi{\RR^+}{R}
\varepsilon_{Z^{\bigdot} \ast} : D^{+} (\Ysheaf{\Gamma} (S \times D),
\mathcal{O}_S) \to D^{+} (S, \mathcal{O}_S)$$ admits a canonical
factorization
$$
D^{+} (\Gamma (S \times D), \mathcal{O}_S) \xrightarrow{F^{\sisi{\RR^+}{R}
\varepsilon} Z^{\bigdot}\ast} D^{+} F(S, \mathcal{O}_S)
\xrightarrow{\iota} D^{+} (S, \mathcal{O}_S){\quoi}.
$$
Moreover, for every integer $i$, the following diagram is commutative:
\[
\xymatrix@=2.5cm{
D^{+}(\Ysheaf{\Gamma}(S\times D),\Oo_S)
\ar[r]^(.55){\sisi{\RR^{+}}{R}e^{\ast}_i} \ar[d]_{F^{\sisi{\RR^{+}}{R}\epsilon_{Z^{\bigdot}\ast}}}& D^{+}(S,\Oo_S)\\
D^{+}F(S,\Oo_S) \ar[ru]_{\Gr_i}& .
}
\]
\end{enonce*}

Consider\pageoriginale the functor $C^+ (C^+(S, \mathcal{O}_S))
\xrightarrow{(\quad)_s} C^+(S, \mathcal{O}_S)$: if $X^{\cdot\cdot}$ is
an object of $C^+ (C^+(S, \mathcal{O}_S))$, we equip
$(X^{\cdot\cdot})_s$ with its second canonical filtration [\cf
  \cite{sga4.exp.Vbis.bib.5}. I.4.8]; this gives a factorization
\begin{equation*}
(\quad)_s : C^+ (C^+(S, \mathcal{O}_S)) \longrightarrow C^+ (F(S,
  \mathcal{O}_S)) \xrightarrow{\iota} C^+(S, \mathcal{O}_S)\tag{2.5.3.1}
\end{equation*}\etiquette[2.5.3.1]{Vbis.eq2.5.3.1}
which passes to the categories $K^+$:
\begin{equation*}
(\quad)_s : K^+ (C^+(S, \mathcal{O}_S)) \longrightarrow K^+ (F(S,
  \mathcal{O}_S)) \xrightarrow{\iota} K^+(S, \mathcal{O}_S)\tag{2.5.3.2}
\end{equation*}\etiquette[2.5.3.2]{Vbis.eq2.5.3.2}
because a homotopy in $C^+ (C^+(S, \mathcal{O}_S))$ induces a
filtered homotopy.

We then have a diagram
\[\vcenter{
\xymatrix{
D^{+}(\Ysheaf{\Gamma}(S\times\Delta),\Oo_S) \ar[rr]^{R\epsilon_{Z^{\bigdot}}\ast} \ar@{-->}[rd]&                   &D^{+}(S,\Oo_S)\\
                                              & D^{+}F(S,\Oo_S)
\ar[ru]_{\imath}&\\
                                  & K^{+}(F(S,\Oo_S))
\ar[u] \ar[rd]^{\imath}&\\
K^{+}(\Ysheaf{\Gamma}(S\times \Delta),\Oo_S) \ar[uuu]
\ar[rr]_{(\;)_s\circ K^{+}(\epsilon_{Z^{\bigdot}}\ast)} \ar[ru]& & K^{+}(S,\Oo_S)\ar[uuu]
}}\leqno{(2.5.3.3)}
\]\etiquette[2.5.3.3]{Vbis.eq2.5.3.3}
and one verifies that there is a unique functor
$$
{}_F \sisi{\RR^{+}}{R} \epsilon_{Z^{\bigdot} \ast} : D^{+}(\Ysheaf{\Gamma}(S \times
\Delta), \mathcal{O}_S) \longrightarrow D^{+}F(S, \mathcal{O}_S)
$$
making\pageoriginale \eqhyperref[Vbis.eq2.5.3.3]{2.5.3.3} commutative.

The rest of the proposition is evident.

\begin{enonce*}{Proposition 2.5.4}\etiquette[2.5.4]{Vbis.prop2.5.4}
Let $(E, A)$ be a semi-simplicial ringed topos and
$$\theta : (E, A) \to (S \times \Delta,
\mathcal{O}_S)$$ an augmentation. For every integer $i$, denote by
$$\theta_i : (E_i, A_i) \to (S, \mathcal{O}_S)$$ the morphism of
ringed topoi induced by $\theta$ over $i$.

The functor $\sisi{\RR^+}{R} \bar{\theta}_\ast : D^+ (\Ysheaf{\Gamma}(E), A)
\to D^+ (S, \mathcal{O}_S)$ admits a canonical factorization
$$
D^+ (\Ysheaf{\Gamma}(E), A) \xrightarrow{{}_F\sisi{\RR^{+}}{R} \bar{\theta}_{\ast}}
D^{+} F(S, \mathcal{O}_S) \xrightarrow{\iota} D^+ (S, \mathcal{O}_S)
$$
such that, for every integer $i$, the canonical diagram
\[
\xymatrix@C=2.5cm@R=1.75cm{
D^{+}(\Ysheaf{\Gamma}(E),A)  \ar[r]^{\sisi{\RR^{+}}{R}e_i^{\ast}}
\ar[d]_{{}_F \sisi{\RR^{+}}{R}
\bar{\theta}_\ast}&
D^{+}(E_i,A_i) \ar[d]^{\sisi{\RR^{+}}{R}\theta_{i\ast}}\\
D^{+}F(S,\Oo_S)  \ar[r]_{\Gr_i} & D^{+}(S,\Oo_S)\\
}
\]
is essentially commutative.
\end{enonce*}

This follows from what precedes and from \hyperref[Vbis.cor1.3.12]{(1.3.1.2)}.

\begin{enonce*}{Proposition 2.5.5}\etiquette[2.5.5]{Vbis.prop2.5.5}
Let $(E, A)$ be a semi-simplicial ringed topos and let
$$\theta : (E, A) \to (S \times \Delta, \mathcal{O}_S)$$ be an
augmentation of cohomological $1$-descent. Let
$$H : (S, \mathcal{O}_S) \to (R, \mathcal{O}_R)$$ be a morphism of
ringed topoi. Then there is a spectral functor from
$D^+(S, \mathcal{O}_S)$ to $\Mod(R, \mathcal{O}_R)$:
$$
E^{pq}_1 = R^q(H \circ \theta_p)_{\ast} \circ
\theta^{\ast}_p \longrightarrow R^{p+q}H_{\ast}
$$
called\pageoriginale the \emph{descent spectral functor}.
\end{enonce*}

Let $H \circ \theta : (E, A) \to (R \times \Delta, \mathcal{O}_R)$
be the augmentation canonically deduced from $\theta$ by composition
with $H$; we have
\begin{equation*}
\bar{(H \circ \theta)}_{\ast} = H_{\ast} \circ \bar{\theta}_{\ast}\tag{2.5.5.1}
\end{equation*}\etiquette[2.5.5.1]{Vbis.eq2.5.5.1}
so that the diagram
\[\vcenter{
\xymatrix@C=2.25cm@R=2.2cm{
D^{+}(S,\Oo_S) \ar[r]^{\sisi{\LL^{+}(\bar{\theta}^{\ast})}{\bar{\theta}^{\ast}}} \ar[rd]_{\sisi{\RR^+}{R}H_{\ast}} &
D^{+}(\Ysheaf{\Gamma}(E),A)
\ar[d]^{\sisi{\RR^{+}}{R}(\overline{H\circ
\theta})_\ast}\\
                  & D^{+}(R,\Oo_R)
}}\leqno{(2.5.5.2)}
\]\etiquette[2.5.5.2]{Vbis.eq2.5.5.2}
is essentially commutative, since
$$
\sisi{\RR^+}{R} \bar{(H \circ \theta)}_{\ast} \xrightarrow{\sim} \sisi{\RR^+}{R}
H_\ast \circ \sisi{\RR^+}{R} \bar{\theta}_{\ast} \mbox{ and } \id_{D^{+}(S,
  \mathcal{O}_S)} \xrightarrow{\sim} \sisi{\RR^+}{R} \bar{\theta}_\ast \circ
\bar{\theta}^{\ast}{\quoi}.
$$
By (\hyperref[Vbis.prop2.5.4]{2.5.4}), for every $i$ we have an
essentially commutative diagram
\[\vcenter{
\xymatrix@C=2.25cm@R=1.75cm{
D^{+}(S,\Oo_S) \ar[r]^{\bar{\theta}^{\ast}} \ar[rdd]_{\sisi{\RR^{+}}{R}H_{\ast}}&
D^{+}(\Ysheaf{\Gamma}(E),A)
\ar[r]^{\sisi{\RR^{+}}{R}e^{\ast}_i} \ar[d]_{{}_F\sisi{\RR^{+}}{R}(\bar{H\circ\theta})_\ast}&
D^{+}(E_i,A_i) \ar[d]^{\sisi{\RR^{+}}{R}(H\circ\theta_{i})_{\ast}}\\
                  & D^{+}F(R,\Oo_R)   \ar[d]_{\imath}\ar[r]_{\Gr_i}      & D^{+}(R,\Oo_R)\\
                  & D^{+}(R,\Oo_R)          &   .
}}\leqno{(2.5.5.3)}
\]\etiquette[2.5.5.3]{Vbis.eq2.5.5.3}
The\pageoriginale asserted spectral functor is obtained by taking the
canonical spectral functor on $D^+ F(R, \mathcal{O}_R)$ and composing
it with ${}_F R(\bar{H \circ \theta})_{\ast}\circ\bar{\theta}^{\ast}$.

We leave it to the reader to verify the following proposition:

\begin{enonce*}{Proposition 2.5.6}\etiquette[2.5.6]{Vbis.prop2.5.6}
With the preceding notation, the $E^{pq}_2$ term of the preceding
spectral functor is
$$
\check{H}^p \circ R^q \Gamma ((H \circ \theta)_{\ast}) \circ
\bar{\theta}^\ast
$$
where $\check{H}^p : \Mod(\Ysheaf{\Gamma}(R \times \Delta), \mathcal{O}_R) \to
\Mod(R, \mathcal{O}_R)$ is the functor that assigns to every functor
$\Delta \xrightarrow{F} \Mod(R, \mathcal{O}_R)$ the homology object
in degree $p$ of the associated complex (denoted by
$\varepsilon_{Z^{\bigdot}\ast} (F)^{\bigdot}$ in
(\hyperref[Vbis.subsubsec2.3.3]{2.3.3})).
\end{enonce*}

We now return to the terminology introduced in
(\hyperref[Vbis.subsec2.4]{2.4}): by the sorites
(\hyperref[Vbis.subsubsec2.5.2]{2.5.2}), all the preceding
considerations carry over word for word by placing the letter $G$ as
a subscript wherever this makes sense. The details are left to the
reader; in particular, one obtains:

\begin{enonce*}{Proposition 2.5.7}\etiquette[2.5.7]{Vbis.prop2.5.7}
Let\pageoriginale $\theta : X \to C^D_S$ be an augmentation of
cohomological $1$-descent relative to $G$, and let $h : S \to R$ be a
morphism of $B$; there is a spectral functor from $D^+_{G_S}(S)$ to
$G_R$:
$$
E^{pq}_1 = R^q_G(h \circ \theta)_\ast \circ \theta_{p,G}^\ast
\to R^{p + q}_G h_{\ast}
$$
with $E^{pq}_2 = \check{H}^p \circ R^q_G \Gamma((h \circ \theta)_\ast)
\circ \bar\theta_G^{\star}$.
\end{enonce*}

% END INCLUDED SGA4_English_current_buildable_source_20260906/source/05bis/drafts/SGA4_VBIS_LINES_1691_1929_English.texfrag


% BEGIN INCLUDED SGA4_English_current_buildable_source_20260906/source/05bis/drafts/SGA4_VBIS_LINES_1930_2298_English.texfrag
\section{Descent criteria}\etiquette[3]{Vbis.sec3}

\setcounter{subsection}{-1}
\subsection{Notation}\etiquette[3.0]{Vbis.subsec3.0}

Throughout what follows, we retain the notation of
(\hyperref[Vbis.subsec2.4]{2.4}). We shall always assume \emph{that
finite fiber products are representable in} $B$.

We shall also assume that $B$ satisfies the following condition:

\setcounter{subsubsection}{-1}
\subsubsection{}\etiquette[3.0.0]{Vbis.subsubsec3.0.0}
\emph{Coproducts exist in $B$, are disjoint and universal, and are
families of effective $\mathcal{E}$-descent and effective
$G^{\circ}$-descent} [\cf \cite{sga4.exp.Vbis.bib.4} (9.23), (9.25), and (9.27)].

\subsubsection{}\etiquette[3.0.1]{Vbis.subsubsec3.0.1}
We now recall some classical notation concerning simplicial objects
of a category [\cf \cite{sga4.exp.Vbis.bib.3}, Chap.~II, and (V, Appendix)].

Let $E$ be a category \emph{having finite projective limits}.
$\Delta^{\circ}E$ denotes the category of simplicial objects of $E$
(\cf (\hyperref[Vbis.subsubsec1.2.6]{1.2.6})).

For an integer $n$, denote by $\mathop{\Delta}_n$ (\resp
${\mathop{\Delta}_n}^{+}$, ${\mathop{\Delta}_n}^{-}$) the full
subcategory of $\Delta$ (\resp $\Delta^+$, $\Delta^-$) consisting of
the objects $[p]$ such that $p\leq n$.

The restriction functor
\begin{equation*}
i_n^{\ast}:\Delta^{\circ}E\longrightarrow
{\mathop{\Delta}_n}^{\circ}E\tag{3.0.1.1}
\end{equation*}\etiquette[3.0.1.1]{Vbis.eq3.0.1.1}
has a right adjoint $i_{n\ast}$, since finite projective limits exist
in\pageoriginale $E$ (\cf (\hyperref[Vbis.lem1.2.10]{1.2.10})). The
functor $i_{n\ast}\circ i_n^{\ast}$ is denoted by $\cosk_n$ and is
called the \emph{coskeleton functor of order $n$}; by abuse of
notation, we shall also use $\cosk_n$ for the functor $i_{n\ast}$.

For every integer $p$, denote by $\mathop{\Delta}_{n[p]}$ (\resp
${\mathop{\Delta}_{n[p]}}^+$) the full subcategory of
$\Delta_{[p]}$ (\resp $\Delta_{[p]}^+$), the category of objects of
$\Delta$ (\resp $\Delta^+$) over $[p]$, defined by the objects $[q]$
over $[p]$ such that $q\leq n$. We leave it to the reader to verify
that
\begin{equation*}
(\cosk_n(X))_p=\varprojlim_{\mathop{\Delta}_{n[p]}}X_q
\xrightarrow{\sim}
\varprojlim_{\mathop{\Delta}_{n[p]}^{+}}X_q{\quoi}.
\tag{3.0.1.2}
\end{equation*}\etiquette[3.0.1.2]{Vbis.eq3.0.1.2}

\subsubsection{}\etiquette[3.0.2]{Vbis.subsubsec3.0.2}
Let $X$ and $X'$ be two objects of $\Delta^{\circ}E$, and let
$f,g:X\rightrightarrows X'$ be two morphisms. A \emph{homotopy from
$f$ to $g$} consists of a map
$h_n:\Hom_{\Delta}([n],[1])\to\Hom_E(X_n,X'_n)$ for every $n$,
satisfying the following two conditions:
\begin{equation*}
h_n(\partial^1)=f_n\quad\mbox{and}\quad
h_n(\partial^\circ)=g_n\quad\mbox{for every }n{\quoi}.
\tag{3.0.2.1}
\end{equation*}\etiquette[3.0.2.1]{Vbis.eq3.0.2.1}

\setcounter{paragraph}{1}
\paragraph{}\etiquette[3.0.2.2]{Vbis.par3.0.2.2}
For every arrow $[n]\to[p]$ in $\Delta$, the following diagram is
commutative:
\[
\xymatrix@C=2.25cm{
\Hom_{\Delta}\big([p],[1]\big) \ar[r] \ar[dd]&
\Hom_E(X_p,X'_p) \ar[rd]&\\
& & \Hom_E(X_p,X'_p)\\
\Hom_{\Delta}\big([n],[1]\big) \ar[r]&
\Hom_E(X_n,X'_n) \ar[ru]&.
}
\]
N.B.: The relation thus introduced on the set of morphisms from $X$
to $X'$ is not an equivalence relation. This inconvenience can be
remedied by introducing the notion of \emph{composite homotopy}
[\cf \cite{sga4.exp.Vbis.bib.3}, Chap.~III], but we shall not use that notion.

\begin{enonce*}{Lemma 3.0.2.3}\etiquette[3.0.2.3]{Vbis.lem3.0.2.3}
Let\pageoriginale $X$ and $X'$ be two objects of $\Delta^{\circ}E$,
let $f,g:X\rightrightarrows X'$ be two morphisms, and let
$F:E\to C$ be a functor, where $C$ is an abelian category. A
simplicial homotopy from $f$ to $g$ induces a homotopy between the
morphisms of cochain complexes canonically associated with $F(f)$ and
$F(g)$.
\end{enonce*}

This reduces to the case where $C$ is the category of modules over a
ring, and one applies [\cite{sga4.exp.Vbis.bib.5}, I, (3.7.1)].

\begin{enonce*}{Lemma 3.0.2.4}\etiquette[3.0.2.4]{Vbis.lem3.0.2.4}
Let $n$ be an integer, and let $X$ and $X'$ be two objects of
$\Delta^{\circ}E$. Let $f$ and $g$ be two morphisms from $X$ to $X'$
such that $f_p=g_p$ for $p<n$. Then there exists a simplicial homotopy
from $\cosk_n(f)$ to $\cosk_n(g)$.
\end{enonce*}

For $p<n$, define
$$h_p:\Hom_{\Delta}([p],[1])\to\Hom(X_p,X'_p)$$
to be the constant map with value $f_p=g_p$. Next define
$h_n:\Hom_{\Delta}([n],[1])\to\Hom(X_n,X'_n)$ by sending every
element of $\Hom_{\Delta}([n],[1])$, except $\partial^{\circ}$, to
$f_n$, and $\partial^{\circ}$ to $g_n$. Finally, observe that
$$
\Hom_{\Delta}([k],[1])=
\varprojlim_{\substack{[q]\to[k]\\q\leq n}}\Hom([q],[1])
$$
for $k>n$, which makes it possible to define $h_k$ canonically.

\subsection{Comparison of two augmentations from the viewpoint of
1-cohomological descent}\etiquette[3.1]{Vbis.subsec3.1}

Let $S$ be an object of $B$. Denote by
$\Hom_{\plat}(D^{\circ},B/S)$ the full subcategory of
$\Hom(D^{\circ},B/S)$ whose objects are the flat augmentations
[\cf (\hyperref[Vbis.subsubsec1.2.6]{1.2.6}) and
(\hyperref[Vbis.def2.4.3.1]{2.4.3.1})].

\begin{enonce*}{Proposition 3.1.1}\etiquette[3.1.1]{Vbis.prop3.1.1}
Let $X\xrightarrow{u}C^D_S$ and $X'\xrightarrow{u'}C^D_S$ be two
objects of $\Hom_{\plat}(D^{\circ},B/S)$. Let $f:X\to X'$ be a
morphism over $S$ (\ie a morphism in $\Hom(D^{\circ},B/S)$). There is
naturally associated with it a functorial morphism
$$
\eta^f_G:\sisi{\RR^+}{R}_G\bar u'_\ast\circ{\bar u'}_G^\ast
\longrightarrow
\sisi{\RR^+}{R}_G\bar u_\ast\circ\bar u_G^\ast
$$
such that the following diagram is commutative:
\[
\xymatrix@=1.5cm{
&R_G\bar u'_\ast\circ{\bar u'}_G^{\ast}
\ar[dd]^{\eta^f_G}\\
\id_{D_{G_S}^{+}(S)}\ar[ru]\ar[rd]&\\
&R_G\bar u_\ast\circ\bar u_G^{\ast}.
}
\]

Moreover,\pageoriginale if $D=\Delta$, then $\eta^f_G$ depends only on the homotopy
class \emph{(\cf (\hyperref[Vbis.subsubsec3.0.2]{3.0.2}))} of $f$ in
$\Hom(\Delta^{\circ},B/S)$.
\end{enonce*}

It is clear that, for the construction of $\eta_G^f$, we may suppose
that $G=\Ysheaf{\Mod}(\mathcal{E}^{\circ},\mathcal{O})$.

Let $I$ be the category defined by the graph
\begin{equation*}
{0}\longrightarrow{1}{\quoi}.\tag{3.1.1.1}
\end{equation*}\etiquette[3.1.1.1]{Vbis.eq3.1.1.1}
Denote by $r_0:D\to I\times D$ (\resp $r_1$) the fully faithful
functor defined by $r_0(i)=(0,i)$ (\resp $r_1(i)=(1,i)$).

By virtue of the canonical isomorphism
\begin{equation*}
\Hom((I\times D)^{\circ},B/S)\xrightarrow{\sim}
\Hom(I^{\circ},\Hom(D^{\circ},B/S))\tag{3.1.1.2}
\end{equation*}\etiquette[3.1.1.2]{Vbis.eq3.1.1.2}
the data of (\hyperref[Vbis.prop3.1.1]{3.1.1}) define a flat
augmentation
\nde{Explicitly, one has $(XfX')_{(0,i)}=X_i$ and
$(XfX')_{(1,j)}=X'_j$; to
$\alpha\in\Hom_D(i,j)=\Hom((0,i),(0,j))$ is associated
$\alpha\circ f_i:(XfX')_{(0,i)}\to(XfX')_{(1,j)}$, and the
endomorphisms of $(\varepsilon,i)$, $\varepsilon=0,1$, act in the
evident way.}
\begin{equation*}
XfX'\xrightarrow{ufu'}C^{I\times D}_S{\quoi}.\tag{3.1.1.3}
\end{equation*}\etiquette[3.1.1.3]{Vbis.eq3.1.1.3}

Let $F$ be an object of $\Mod(\Ee^\circ_S,\Oo_S)$. The canonical
morphism
\begin{equation*}
\Ee_{I\times D}^{\ast}(F)\longrightarrow
\Gamma((ufu')_\ast)\circ\Gamma((ufu')^{\ast})
(\Ee_{I\times D}^{\ast}(F)){\quoi}\tag{3.1.1.4}
\end{equation*}\etiquette[3.1.1.4]{Vbis.eq3.1.1.4}
may be interpreted as a commutative triangle in
$\Mod(\Ysheaf{\Gamma}(\bar{C^D_S}),\Oo)$:
\[\vcenter{
\xymatrix@=1.5cm{
&\Gamma(u'_\ast)\circ\Gamma({u'}^{\ast})(\Ee_D^{\ast}(F))\ar[dd]\\
\Ee_D^{\ast}(F)\ar[ru]\ar[rd]&\\
&\Gamma(u_\ast)\circ\Gamma(u^{\ast})(\Ee_D^{\ast}(F))
}}\leqno{(3.1.1.5)}
\]\etiquette[3.1.1.5]{Vbis.eq3.1.1.5}
and hence\pageoriginale gives a functorial morphism
$\bar u'_\ast\circ{\bar u'}^{\ast}
\xrightarrow{\alpha_f}\bar u_\ast\circ\bar u^{\ast}$ such that the
diagram
\[\vcenter{
\xymatrix@=1.5cm{
&\bar u'_\ast\circ{\bar u'}^{\ast}\ar[dd]^{\alpha_f}\\
\id_{\Mod(\Ee^{\circ}_S,\Oo_S)}\ar[ru]\ar[rd]&\\
&\bar u_\ast\circ\bar u^{\ast}
}}\leqno{(3.1.1.6)}
\]\etiquette[3.1.1.6]{Vbis.eq3.1.1.6}
is commutative.

One consequently obtains a functorial morphism such that the diagram
\[\vcenter{
\xymatrix@=1.5cm{
&K^{+}(\bar u'_\ast)\circ K^{+}({\bar u'}^{\ast})
\ar[dd]^{K^{+}(\alpha_f)}\\
\id_{K^{+}(\Ee^{\circ}_S,\Oo_S)}\ar[rd]\ar[ru]&\\
&K^{+}(\bar u_\ast)\circ K^{+}(\bar u^{\ast})
}}\leqno{(3.1.1.7)}
\]\etiquette[3.1.1.7]{Vbis.eq3.1.1.7}
is commutative.

\noindent
(Observe that $K^+(\alpha_f)$ can also be obtained formally by
adjunction, using the isomorphism
$K^+(\Gamma(f^{\ast}))\circ K^+({\bar u'}^{\ast})
\xrightarrow{\sim}K^+(\bar u^{\ast})$.)

Let $F^{\bigdot}$ be a complex in $K^+(S)$, and let
$\xi:K^+(\overline{ufu'}^{\ast})(F^{\bigdot})\to J^{\bigdot}$ be a
quasi-isomorphism such that, for every integer $n$, $J^n$ is totally
acyclic objectwise (\cf (\hyperref[Vbis.rem1.3.13]{1.3.13})). By
(\loccit), the canonical arrows
$$
K^+({\bar u'}^\ast)(F^{\bigdot})\to r_1^{\ast}(J^{\bigdot})
\quad\mbox{and}\quad
K^+(\bar u^\ast)(F^{\bigdot})\to r_0^{\ast}(J^{\bigdot})
$$
are\pageoriginale quasi-isomorphisms (\cf
(\hyperref[Vbis.subsubsec1.2.8]{1.2.8}) for the notation). One
therefore obtains a morphism
$$
\sisi{\RR^+}{R}\bar u'_\ast\circ{\bar u'}^\ast(F^{\bigdot})
\xrightarrow{\eta^f(F^{\bigdot})}
\sisi{\RR^+}{R}\bar u_\ast\circ\bar u^\ast(F^{\bigdot})
$$
such that the following diagram is commutative in $D^+(S)$:
{\def\objectstyle{\scriptstyle}
\def\labelstyle{\scriptstyle}

\[\vcenter{
\xymatrix@C=0cm@R=1cm{
&Q\circ K^{+}(\bar u'_\ast)(r^{\ast}_{0}(J^{\bigdot}))
\ar[rr]^{\sim}\ar@{->}'[d][ddd]&&
R\bar u'_\ast\circ
Q(r^{\ast}_0(J^{\bigdot}))\ar[ddd]\\
Q\circ K^{+}(\bar u'_\ast)\circ K^{+}({\bar u'}^{\ast})(F^{\bigdot})
\ar[rr]\ar[ru]\ar[ddd]_{Q(K^{+}(\alpha_f)(F^{\bigdot}))}&&
R\bar u'_\ast\circ
{\bar u'}^{\ast}(F^{\bigdot})\ar[ru]^{\sim}
\ar@{-->}[ddd]^(.35){\eta^{f}(F^{\bigdot})}&\\
&&&\\
&Q\circ K^{+}(\bar u_\ast)(r^{\ast}_1(J^{\bigdot}))
\ar@{->}'[r][rr]^(.25){\sim}&&
R\bar u_\ast\circ Q(r_1^{\ast}(J^{\bigdot}))\\
Q\circ K^{+}(\bar u_\ast)\circ K^{+}(\bar u^{\ast})(F^{\bigdot})
\ar[ru]\ar[rr]&&
R\bar u_\ast\circ\bar u^{\ast}(F^{\bigdot})
\ar[ru]&
}}\leqno{(3.1.1.8)}
\]\etiquette[3.1.1.8]{Vbis.eq3.1.1.8}
}
\noindent
It is clear that $\eta^f(F^{\bigdot})$ is independent of the chosen
resolution, and one verifies that it is functorial in $F^{\bigdot}$
as the latter varies in $D^{+}(S)$.

To complete the proof of \hyperref[Vbis.prop3.1.1]{(3.1.1)}, it is enough to prove the following
lemma:

\begin{enonce*}{Lemma 3.1.1.9}\etiquette[3.1.1.9]{Vbis.lem3.1.1.9}
Let\pageoriginale $X\xrightarrow{u}C^{\Delta}_S$ and
$X'\xrightarrow{u'}C^{\Delta}_S$ be two objects of
$\Hom_{\plat}(\Delta^{\circ},B/S)$. Let $f,g:X\to X'$ be two
morphisms in $\Hom(\Delta^{\circ},B/S)$. If there exists a simplicial
homotopy from $f$ to $g$, then $\eta^f=\eta^g$.
\end{enonce*}

Let
$h_n:\Hom_{\Delta}([n],[1])\to\Hom_{B/S}(X_n,{X'}_n)$ be a
simplicial homotopy from $f$ to $g$. Let $\overline{I\times D}$ be the
category obtained from $I\times D$ by adjoining the arrows
$(1,n)\to(0,n)$ in bijective correspondence with
$h_n(\Hom([n],[1]))$, together with the relations imposed by the
compatibility of the $h_n$ as $n$ varies. The data of the lemma define
a flat augmentation
$\overline{XfX'}\xrightarrow{t}C^{\overline{I\times D}}_S$. Let
$F^{\bigdot}$ be a complex in $C^+(S)$, and let $J^{\bigdot}$ be a
resolution of $C^+(t^{\ast})(F^{\bigdot})$ such that, for every
integer $n$, $J^n$ is totally acyclic objectwise. There exists a
simplicial homotopy between
$$
C^{+}(\Gamma(u'_\ast))(r^{\ast}_0(J^{\bigdot}))
\rightrightarrows
C^{+}(\Gamma(u_\ast))(r^\ast_1(J^{\bigdot})),
$$
and hence a homotopy after passage to the total complexes
(\cf (\hyperref[Vbis.lem3.0.2.3]{3.0.2.3})).

\begin{enonce*}{Definition 3.1.2}\etiquette[3.1.2]{Vbis.def3.1.2}
With the notation of \emph{\hyperref[Vbis.prop3.1.1]{(3.1.1)}}, we say that $f$ is an
\emph{equivalence for $G$-1-cohomological descent} if $\eta^f_G$ is an
isomorphism.
\end{enonce*}

\setcounter{subsubsection}{2}
\subsubsection{}\etiquette[3.1.3]{Vbis.subsubsec3.1.3}
In what follows, denote by $L_i:\Delta\to\Delta\times\Delta$
(\resp $C_i:\Delta\to\Delta\times\Delta$) the canonical functor
defined by $[n]\longmapsto[n]\times[i]$ (\resp
$[n]\longmapsto[i]\times[n]$).

\begin{enonce*}{Proposition 3.1.4}\etiquette[3.1.4]{Vbis.prop3.1.4}
Let $X\xrightarrow{u}C^{\Delta\times\Delta}_S$ and
$X'\xrightarrow{u'}C^{\Delta\times\Delta}_S$ be two objects of
$\Hom_{\plat}((\Delta\times\Delta)^{\circ},B/S)$. Let $f:X\to X'$ be
a morphism over $S$. Suppose that
$f\ast\id_{L_i}:X\circ L_i\to X'\circ L_i$ (\resp
$f\ast\id_{C_i}:X\circ C_i\to X'\circ C_i$) is an equivalence for
$G$-1-cohomological descent for every integer $i$. Then $f$ is an
equivalence for $G$-1-cohomological descent.
\end{enonce*}

By the description of $\eta^f$ (\cf the proof of
(\hyperref[Vbis.prop3.1.1]{3.1.1})), it is enough to show that a
certain morphism of triple complexes induces an isomorphism on the
cohomology of the total complexes. A standard spectral-sequence
argument then gives the conclusion.

\begin{enonce*}{Proposition 3.1.5}\etiquette[3.1.5]{Vbis.prop3.1.5}
Let\pageoriginale $X$ and $X'$ be two objects of
$\Hom_{\plat}(\Delta^{\circ},B/S)$. Let $f:X\to X'$ be a functorial
morphism such that $f_n:X_n\to X'_n$ is a morphism of
$G$-1-cohomological descent. Then the canonical morphism
$$
[[X|_fX']]\to[[X'|_{\id}X']]
$$
\emph{(\cf (\hyperref[Vbis.subsubsec1.2.7]{1.2.7}))} is an equivalence
for $G$-1-cohomological descent.
\end{enonce*}

This follows from (\hyperref[Vbis.prop3.1.4]{3.1.4}) and the following
lemma, whose proof is left to the reader:

\begin{enonce*}{Lemma 3.1.5.1}\etiquette[3.1.5.1]{Vbis.lem3.1.5.1}
Let $f:R\to S$ be a flat morphism, and let
$Y\xrightarrow{u}C^{\Delta}_S$ and
$Y\xrightarrow{v}C^{\Delta}_R$ be two \emph{flat} augmentations such
that the following diagram is commutative in
$\Hom(\Delta^{\circ},B)$:
\[
\xymatrix@=1.5cm{
Y\ar[rr]^{v}\ar[rd]_{u}&&C^{\Delta}_R\ar[ld]^f\\
&C^{\Delta}_S&
}
\]
If $v$ is an augmentation of $G$-1-cohomological descent, then it is
an equivalence for $G$-1-cohomological descent.
\end{enonce*}

In the applications, we shall combine
(\hyperref[Vbis.prop3.1.4]{3.1.4}) and
(\hyperref[Vbis.prop3.1.5]{3.1.5}) with the following result:

\begin{enonce*}{Proposition 3.1.6}\etiquette[3.1.6]{Vbis.prop3.1.6}
Let $X\xrightarrow{u}C^{\Delta\times\Delta}_S$ be a flat
augmentation. For $u$ to be an augmentation of
$G$-1-cohomological descent, it is enough that
$$
X\circ L_i\xrightarrow{u\ast\id_{L_i}}C^{\Delta}_S
\quad(\resp\ X\circ C_i\xrightarrow{u\ast\id_{C_i}}C^{\Delta}_S)
$$
be so for every integer $i$.
\end{enonce*}

This follows from the explicit construction of
$\sisi{\RR^+}{R}_G\bar u_\ast$
(\cf (\hyperref[Vbis.subsubsec2.3.10]{2.3.10})).

\begin{enonce*}{Corollary 3.1.7}\etiquette[3.1.7]{Vbis.cor3.1.7}
Let $X\xrightarrow{u}C^{\Delta}_S$ be an augmentation. Then $u$ is of
$G$-1-cohomological descent if and only if
$[[X|_{\id}X]]\to C^{\Delta\times\Delta}_S$ is so.
\end{enonce*}

% END INCLUDED SGA4_English_current_buildable_source_20260906/source/05bis/drafts/SGA4_VBIS_LINES_1930_2298_English.texfrag


% BEGIN INCLUDED SGA4_English_current_buildable_source_20260906/source/05bis/drafts/SGA4_VBIS_LINES_2299_2585_English.texfrag
\subsection{Localization criteria}\etiquette[3.2]{Vbis.subsec3.2}\pageoriginale

\begin{enonce*}{Proposition 3.2.1}\etiquette[3.2.1]{Vbis.prop3.2.1}
Let $Y\xrightarrow{v}C^{\Delta}_S$ be an augmentation of
$G$-cohomological $1$-descent. For a flat augmentation
$X\xrightarrow{u}C^{\Delta}_S$ to be of $G$-cohomological
$1$-descent, it suffices that it become so after all base changes
$Y_n\to S$.
\end{enonce*}

By means of the base changes $Y_n\to S$, one constructs a double
semi-simplicial object $\fprod{X}{Y}{S}$ augmented toward $S$. By
(\hyperref[Vbis.prop3.1.4]{3.1.4}) and
(\hyperref[Vbis.lem3.1.5.1]{3.1.5.1}), the canonical morphism
$$
\fprod{X}{Y}{S}\to[[X|_{\id}X]]
$$
is an equivalence for $G$-cohomological $1$-descent. By
(\hyperref[Vbis.cor3.1.7]{3.1.7}), it therefore suffices to show that
the augmentation of $\fprod{X}{Y}{S}$ toward $S$ is of
$G$-cohomological $1$-descent. This follows from the fact that the
canonical morphism
$$
\fprod{X}{Y}{S}\to[[Y|_{\id}Y]]
$$
is an equivalence for $G$-cohomological $1$-descent; again one uses
(\hyperref[Vbis.prop3.1.4]{3.1.4}),
(\hyperref[Vbis.lem3.1.5.1]{3.1.5.1}), and
(\hyperref[Vbis.cor3.1.7]{3.1.7}).

\setcounter{subsubsection}{1}
\subsubsection{}\etiquette[3.2.2]{Vbis.subsubsec3.2.2}
Let
\[
\vcenter{
\xymatrix@C=3cm@R=1.5cm{
Y'\ar[r]^{\mathsf g'}\ar[d]_{f'}&Y\ar[d]^{f}\\
X'\ar[r]_{\mathsf g}&X
}}\leqno{(3.2.2.1)}
\]\etiquette[3.2.2.1]{Vbis.eq3.2.2.1}
be a commutative diagram in $B$, and let $F$ be an object of
$\Mod(\mathcal E^{\circ}_Y,\mathcal O_Y)$. There is a unique
morphism
$$
\varphi(F):g_{\ast}(f^{\ast}(F))\to
f'^{\ast}(g'_{\ast}(F))
$$
such that the following diagram (in
$\Ysheaf{\Mod}(\mathcal E^{\circ},\mathcal O)$) commutes\nde{As
above, one must regard the commutative diagram as lying in
$B^\circ$ and use the fibered and cofibered structure of
$\Mod(\mathcal E^\circ,\mathcal O)\rightarrow B^\circ$, which
explains the unusual appearance of the base-change morphism. Under
the (adjunction) arrows in the following diagram we have written
their images in $B^\circ$}:
\[
\vcenter{
\xymatrix@R=1.85cm@C=1.75cm{
&f^{\ast}(F)\ar[r]^{f^\circ}\ar[ld]_{g^\circ}&F\ar[d]^{g'^\circ}\\
g_{\ast}(f^{\ast}(F))\ar[rd]_{\Id}&&g'_{\ast}(F)\\
&f'^{\ast}(g'_{\ast}(F))\ar[ru]_{f'^\circ}&
}}\leqno{(3.2.2.2)}
\]\etiquette[3.2.2.2]{Vbis.eq3.2.2.2}

If\pageoriginale all the arrows of
\eqhyperref[Vbis.eq3.2.2.1]{3.2.2.1} are now assumed flat, one
defines in the same way a morphism, called the \emph{base-change}
morphism\nde{Here again, it would be more correct to write
$f^\circ$, $g^\circ$, $f'^\circ$, $g'^\circ$ instead of $f$, $g$,
$f'$, $g'$. If one considers the nonopposite fibered and cofibered
structure, $\xi_G$ becomes
$$g_G^\ast\circ R_Gf_{\ast}\longrightarrow
R_Gf'_{\ast}\circ g_G'^\ast.$$}
\begin{equation*}
\xi_G:g_{G\ast}\circ\sisi{\RR^+}{R}_Gf^{\ast}
\longrightarrow\sisi{\RR^+}{R}_Gf'^{\ast}\circ g'_{G\ast}{\quoi}.
\tag{3.2.2.3}
\end{equation*}\etiquette[3.2.2.3]{Vbis.eq3.2.2.3}

\begin{enonce*}{Definition 3.2.2.4}\etiquette[3.2.2.4]{Vbis.def3.2.2.4}
The diagram \emph{\hyperref[Vbis.eq3.2.2.1]{(3.2.2.1)} is said to satisfy the base-change
theorem relative to $G$ if $f$, $g$, $f'$, and $g'$} are flat
morphisms and $\xi_G$ is an isomorphism.
\end{enonce*}

\setcounter{subsubsection}{2}
\subsubsection{}\etiquette[3.2.3]{Vbis.subsubsec3.2.3}
Let $h:S\to S'$ be a \emph{flat} morphism in $B$, let
$X\xrightarrow{u}C^{\Delta}_S$ and
$X'\xrightarrow{u'}C^{\Delta}_{S'}$ be two \emph{flat}
augmentations, and let $f:X\to X'$ be a functorial morphism such that
the diagram (in $\Hom(\Delta^0,B)$)
\[
\vcenter{
\xymatrix@C=3cm@R=1.5cm{
X\ar[r]^{f}\ar[d]_{u}&X'\ar[d]^{u'}\\
C^{\Delta}_S\ar[r]^{h_c}&C^{\Delta}_{S'}
}}\leqno{(3.2.3.1)}
\]\etiquette[3.2.3.1]{Vbis.eq3.2.3.1}
commutes. (Here $h_c$ denotes the constant functorial morphism
defined by $h$.)

Let $K^{\bigdot}$ be a complex in $D_G^+(S')$. The morphism
$$
\eta_G^f(K^{\bigdot}):
\sisi{\RR^+}{R}\bar u'_{\ast}\circ\bar u'^{\ast}(K^{\bigdot})
\longrightarrow
\sisi{\RR^+}{R}(\overline{h_c\circ u})_{\ast}\circ
(\overline{h_c\circ u})^{\ast}(K^{\bigdot})
$$
induces, in view of the identities
$$
\sisi{\RR^+}{R}(\overline{h_c\circ u})_{\ast}\simeq
\sisi{\RR^+}{R}h_c^{\ast}\circ
\sisi{\RR^+}{R}\bar u_{\ast}
\quad\mbox{and}\quad
(\overline{h_c\circ u})^{\ast}\simeq
\bar u^{\ast}\circ h_{c\ast},
$$
a morphism $\ch_G^{f,h}(K^{\bigdot})$, functorial in $K^{\bigdot}$,
such that the following diagram commutes\nde{For readability, we
have suppressed the reference to the thick category $G$ in the
derived functors. The same unusual derivations $h_{c\ast}$ reappear,
owing to the choice of the cofibered structure.}:
\[
\vcenter{
\xymatrix@C=1cm@R=2.25cm{
{}_{}h_{c\ast}\circ
R_{}\bar u'_{\ast}{}_{}\circ
\bar u'^{\ast}(K^{\bigdot})
\ar[rr]^{\ch_G^{f,h}{(K^{\bigdot})}}&&
R_{}\bar u_{\ast}\circ
{}_{}\bar u^{\ast}\circ{}_{}h_{c\ast}(K^{\bigdot})\\
&h_{c\ast}(K^{\bigdot})
\ar[lu]^{{}_{}h_{c\ast}(\alpha(K^{\bigdot}))}\ar[ru]&.
}}\leqno{(3.2.3.2)}
\]\etiquette[3.2.3.2]{Vbis.eq3.2.3.2}

\begin{enonce*}{Lemma 3.2.3.3}\etiquette[3.2.3.3]{Vbis.lem3.2.3.3}
Suppose\pageoriginale that, for every integer $i$, the diagram
\[
\xymatrix@C=3cm@R=1.5cm{
X_i\ar[d]_{u_i}\ar[r]^{f_i}&X'_i\ar[d]^{u'_i}\\
S\ar[r]^{h}&S'
}
\]
satisfies the base-change theorem relative to $G$. Then
$\ch_G^{f,h}$ is an isomorphism.
\end{enonce*}

The morphism $\ch_G^{f,h}(K^{\bigdot})$ may be computed as follows.
Consider the augmentation
$$
X f X'\xrightarrow{(h_c\circ u)f u'}C^{I\times\Delta}_S
$$
(\cf \eqhyperref[Vbis.eq3.1.1.2]{3.1.1.2}), and choose a resolution
$J^{\bigdot}$ of
$K^+(\overline{(h_c\circ u)f u'}^{\ast})(K^{\bigdot})$ such that,
for every integer $n$, $J^n$ is totally acyclic object by object.

One obtains a quasi-isomorphism
$$
\theta:K^+(\Gamma(f^{\ast}))(r_0^{\ast}J^{\bigdot})
\longrightarrow r_1^{\ast}(J^{\bigdot})
$$
and a morphism
$$
t:K^+(\Gamma(h_c^{\ast}))\circ K^+(\Gamma(u'_{\ast}))
(r_0^{\ast}(J^{\bigdot}))\longrightarrow
K^+(\Gamma(u_{\ast}))(r_1^{\ast}(J^{\bigdot}))
$$
which has the factorization
\begin{align*}
&K^+(\Gamma(h_c^{\ast}))\circ K^+(\Gamma(u'_{\ast}))
(r_0^{\ast}(J^{\bigdot}))\xrightarrow{\varphi}
K^+(\Gamma(u_{\ast}))\circ K^+(\Gamma(f^{\ast}))
(r_0^{\ast}(J^{\bigdot}))
\xrightarrow{K^+(\Gamma(u_{\ast}))(\theta)}\\
&\longrightarrow K^+(\Gamma(u_{\ast}))(r_1^{\ast}(J^{\bigdot})),
\end{align*}
and\pageoriginale $\ch^{f,h}(K^{\bigdot})$ is obtained by taking the
image of $t$ under the functor
$\sisi{\RR^+}{R}\epsilon_{Z^{\bigdot\ast}}$.

The preceding factorization shows that $t$ identifies ``object by
object'' with the base-change morphisms associated with the diagrams
\[
\xymatrix@C=3cm@R=1.5cm{
X_i\ar[d]_{u_i}\ar[r]^{f_i}&X'_i\ar[d]^{u'_i}\\
S\ar[r]^{h}&S'
}
\]
It follows, by spectral sequences, that
$\ch_G^{f,h}(K^{\bigdot})$ is an isomorphism, completing the proof.

This leads to a \emph{second localization criterion}:

\begin{enonce*}{Proposition 3.2.4}\etiquette[3.2.4]{Vbis.prop3.2.4}
Let $X\xrightarrow{u}C^{\Delta}_S$ be an object of
$\Hom_{\plat}(\Delta^{\circ},B/S)$. For $u$ to be an augmentation of
$G$-cohomological $1$-descent, it suffices that it become so after
``enough for $G$''\footnote{The expression ``enough for $G$'' means
that there exists a family
$(S_{\alpha}\xrightarrow{h_{\alpha}}S)_{\alpha\in A}$ of flat
morphisms satisfying the preceding conditions and such that the
family of functors
$(h_{\alpha\ast}:G_S\to G_{S_\alpha})_{\alpha\in A}$ is
\emph{conservative}.} flat base changes $h:S'\to S$ for which the
cartesian diagrams
\[
\xymatrix@C=3cm@R=1.5cm{
X'_i\ar[d]\ar[r]&X_i\ar[d]\\
S'\ar[r]^{h}&S
}
\]
satisfy the base-change theorem relative to $G$.
\end{enonce*}

Proposition\pageoriginale (\hyperref[Vbis.prop3.2.4]{3.2.4}) follows
immediately from (\hyperref[Vbis.lem3.2.3.3]{3.2.3.3}), in view of
diagram \eqhyperref[Vbis.eq3.2.3.2]{3.2.3.2}.

\subsection{Properties of morphisms of cohomological descent}
\etiquette[3.3]{Vbis.subsec3.3}

In this section, we assume that the \emph{class of flat morphisms in
$B$ is stable under base change}.

\begin{enonce*}{Proposition 3.3.1}\etiquette[3.3.1]{Vbis.prop3.3.1}
Let $i=1,2$.
\begin{enumerate}
\renewcommand{\theenumi}{\alph{enumi}}
\renewcommand{\labelenumi}{\theenumi)}
\item Every flat morphism which has a section is a morphism of
universal $G$-cohomological $2$-descent.

\item Let $f:X\to S$ be a flat morphism, let $g:S'\to S$ be a
morphism of $G$-cohomological $i$-descent, put
$X'=\fprod{X}{S'}{S}$, and let $f'=f_{(S')}:X'\to S'$. Then $f$ is
of universal $G$-cohomological $i$-descent if and only if $f'$ is.

\item If the composite of two morphisms $f:X\to Y$ and $g:Y\to Z$
is of universal $G$-cohomological $i$-descent, then $g$ is of
universal $G$-cohomological $i$-descent.

\item The composite of two morphisms of universal
$G$-cohomological $i$-descent is a morphism of universal
$G$-cohomological $i$-descent.

\item If $f:X\to X'$ and $g:Y\to Y'$ are two $S$-morphisms of
universal $G$-cohomological $i$-descent, then $\fprod{f}{g}{S}$ is
of universal $G$-cohomological $i$-descent.

\item Let $(u_{\alpha}:X_{\alpha}\to Y_{\alpha})_{\alpha\in A}$
be a family of morphisms of $G$-cohomological $i$-descent. Then
$\amalg u_{\alpha}:\amalg_{\alpha\in A}X_{\alpha}\to
\amalg_{\alpha\in A}Y_{\alpha}$ is a morphism of
$G$-cohomological $i$-descent.
\end{enumerate}
\end{enonce*}

\noindent\emph{Proof.}
For\pageoriginale effective descent relative to $G$, we refer to
\cite{sga4.exp.Vbis.bib.4}.
\begin{enumerate}
\renewcommand{\theenumi}{\alph{enumi}}
\renewcommand{\labelenumi}{\theenumi)}
\item Using a section $s$ of $f$, compare the semi-simplicial objects
augmented toward $S$, $[X|_fS]$ and $[S|_{\id}S]$. By
(\hyperref[Vbis.lem3.0.2.4]{3.0.2.4}), one obtains an equivalence for
$G$-cohomological $1$-descent\nde{For this, note that $[X|_fS]$ and
$[S|_{\id}S]$ are simply the coskeleta of order $0$ of $f$ and
$\id$, respectively, so that $s$ and $f$ induce inverse
isomorphisms in cohomology by
\hyperref[Vbis.prop3.1.1]{3.1.1}.}.

\item This follows from (\hyperref[Vbis.prop3.2.1]{3.2.1}).

\item This follows from a) and b).

\item Let $f:X\to Y$ and $g:Y\to S$ be two morphisms of universal
$G$-cohomological $1$-descent. Put $h=g\circ f$ and consider the
fiber product
\[
\xymatrix@C=3cm@R=1.5cm{
Y\ar[d]_{g}&R\ar[d]^{g'}\ar[l]_{h'}\\
S&X\ar[l]^{h}
}
\]
The morphism $g'$ has a section $s:X\to R$ such that
$h'\circ s=f$. By c), $h'$ is of universal $G$-cohomological
$1$-descent, and the same holds for $h$ by b).

\item This follows formally from b) and d).

\item Let $(Z_{\lambda})_{\lambda\in\Lambda}$ be a family of
objects of $B$. It follows from hypotheses
(\hyperref[Vbis.subsubsec3.0.0]{3.0.0}) that there is a canonical
equivalence
$$
\Mod\left(\mathcal E^{\circ}_{\coprod_{\lambda}Z_{\lambda}},
\mathcal O_{\coprod_{\lambda}Z_{\lambda}}\right)
\xrightarrow{\sim}\displaystyle\mathop{\prod}_{\lambda}
\Mod(\mathcal E^{\circ}_{Z_\lambda},\mathcal O_{Z_\lambda})
$$
\end{enumerate}
which induces an equivalence
$$
G_{\coprod_{\lambda}Z_{\lambda}}
\xrightarrow{\sim}\prod_{\lambda}G_{Z_{\lambda}}{\quoi}.
$$
Moreover, a family
$(Q_{\lambda})_{\lambda}\in
\prod_{\lambda}\Mod(\mathcal E^{\circ}_{Z_\lambda},
\mathcal O_{Z_\lambda})$ corresponds to a totally acyclic object if
and only if $Q_{\lambda}$ is totally acyclic for every $\lambda$.

Now let
$(u_{\lambda}:X_{\lambda}\to Y_{\lambda})_{\lambda\in\Lambda}$ be a
family of morphisms of $G$-cohomological $1$-descent. Since direct
sums in $B$ are disjoint and\pageoriginale universal, for every $n$
there is an identification
$$
\left[\displaystyle\mathop{\coprod}_{\lambda}X_{\lambda}
\mathrel{|}_{\coprod u_{\lambda}}\coprod Y_{\lambda}\right]_n
\simeq\displaystyle\mathop{\coprod}_{\lambda}
[X_{\lambda}|_{u_{\lambda}}Y_{\lambda}]_n
$$
which is functorial in $n$. It then follows from the preceding
remarks and the explicit calculation of
$\sisi{\RR^+}{R}_G\theta_{\ast}$ (\cf \hyperref[Vbis.subsec2.3]{(2.3)}) that
$\coprod_{\lambda}u_{\lambda}$ is a morphism of
$G$-cohomological $1$-descent; the details are left to the reader.
The reader may also verify that
$\coprod_{\lambda}u_{\lambda}$ \emph{remains of $G$-cohomological
$1$-descent after every base change if the same is true of every
$u_{\lambda}$}.

\setcounter{subsubsection}{1}

% END INCLUDED SGA4_English_current_buildable_source_20260906/source/05bis/drafts/SGA4_VBIS_LINES_2299_2585_English.texfrag


% BEGIN INCLUDED SGA4_English_current_buildable_source_20260906/source/05bis/drafts/SGA4_VBIS_LINES_2586_2816_English.texfrag
\subsubsection{}\etiquette[3.3.2]{Vbis.subsubsec3.3.2}
Associate with each object $X$ of $B$ the set of families of
morphisms $(X_\alpha\to X)_{\alpha\in A}$ such that
$\coprod_\alpha X_\alpha\to X$ is a morphism of universal
$G$-$i$-cohomological descent. By \hyperref[Vbis.prop3.3.1]{3.3.1},
this defines a pretopology on $B$.

The topology generated by this pretopology is called the
\emph{topology of $G$-$i$-cohomological descent}. It follows from
\hyperref[Vbis.prop3.3.1]{3.3.1} c) that the covering morphisms for
this topology are precisely the morphisms of universal
$G$-$i$-cohomological descent. Finally, note that the topology of
$G$-$2$-cohomological descent is coarser than the topology of
$G^{\circ}$-descent [\cf \cite{sga4.exp.Vbis.bib.4} (6.23)].

\begin{enonce*}{Theorem 3.3.3}\etiquette[3.3.3]{Vbis.thm3.3.3}
Suppose that \emph{all arrows of $B$ are flat}. Let $S$ be an object
of $B$. Every hypercovering of $S$ for the topology of
$G$-$1$-cohomological descent whose objects are all representable
\emph{(\cf V, Appendix)} defines an augmentation of universal
$G$-$1$-cohomological descent.
\end{enonce*}

The proof proceeds in two steps; we specify that all coskeleta are
to be computed in the category $B/S$.

\begin{enonce*}{Lemma 3.3.3.1}\etiquette[3.3.3.1]{Vbis.lem3.3.3.1}
Let\pageoriginale $X$ be an object of
$\Hom(\Delta^{\circ},B/S)$. For $X$ to be of
$1$-cohomological descent (resp. universal $1$-cohomological
descent), it is enough that $\cosk_n(X)$ have this property for all
sufficiently large $n$.
\end{enonce*}

Indeed, let $\alpha_n:X\to\cosk_n(X)$ be the canonical morphism. One
checks that if $K^{\bigdot}$ is a complex of $D^+_{G_S}(S)$ such
that $H^j(K^{\bigdot})=0$ for $j<N$, then
$H^i(\eta_G^{\alpha_n})$ is an isomorphism for $i<N+n$, whence the
assertion.

\begin{enonce*}{Lemma 3.3.3.2}\etiquette[3.3.3.2]{Vbis.lem3.3.3.2}
Let $n\geq0$, let $X$ and $X'$ be two simplicial objects of $B/S$,
and let $f:X\to X'$ be a morphism. Suppose that:
\begin{enumerate}
\renewcommand{\theenumi}{\roman{enumi}}
\renewcommand{\labelenumi}{(\theenumi)}
\item $X\to\cosk_{n+1}(X)$ is an isomorphism;

\item $X'\to\cosk_{n+1}(X')$ is an isomorphism;

\item $f_i:X_i\to X'_i$ is an isomorphism for $i\leq n$;

\item $f_{n+1}$ is a morphism of universal
$G$-$1$-cohomological descent.
\end{enumerate}
Then, if $X$ is of universal $G$-$1$-cohomological descent, the same
is true of $X'$.
\end{enonce*}

Lemmas \hyperref[Vbis.lem3.3.3.1]{3.3.3.1} and
\hyperref[Vbis.lem3.3.3.2]{3.3.3.2} clearly prove Theorem
\hyperref[Vbis.thm3.3.3]{3.3.3} by induction.

\begin{enonce*}{Lemma 3.3.3.3}\etiquette[3.3.3.3]{Vbis.lem3.3.3.3}
Under the hypotheses of \hyperref[Vbis.lem3.3.3.2]{3.3.3.2}, all
the morphisms $f_p$ are morphisms of universal
$G$-$1$-cohomological descent.
\end{enonce*}

This is trivial for $p\leq n+1$. For $p>n+1$, $X_p$ (resp. $X'_p$)
may be written as
$\varprojlim_{\substack{\Delta^+\\ n+1[p]}}X_q$
(resp. $\varprojlim_{\substack{\Delta^+\\ n+1[p]}}X'_q$).

We may suppose that $p>n+1$. Write
$$
\cosq_{n+1}(X)_p=
\varprojlim_{\substack{[q]\rightarrow[p]\\q\leq n+1}}X_q
$$
as the kernel of the pair of arrows
$$
\Pi X=\prod_{\substack{[q]\rightarrow[p]\\q\leq n+1}}X_q
\rightrightarrows
\prod_{\substack{[i]\stackrel{\alpha}{\hookrightarrow}[j]\\
\searrow\swarrow\\{}[p]\\j\leq n+1}}X_i.
$$
The component $\alpha_X$ indexed by
$\alpha\in\Hom_{[p]}([i],[j])$ of this pair consists, on the one
hand, of the projection $\Pi X\to X_i$ indexed by
$[i]\to[p]$, and, on the other hand, of the composite
$$
\Pi X\longrightarrow X_j\longrightarrow X_i
$$
of the projection indexed by $[j]\to[p]$ with
$\alpha\in\Hom(X_j,X_i)$.

It follows that the diagram
$$
\xymatrix{\cosq_{n+1}(X)_p\ar[r]\ar[d]&\Pi X\ar[d]\\
\cosq_{n+1}(X')_p\ar[r]&\Pi X'}
$$
is Cartesian. Indeed, by the preceding description, it is enough to
prove that for every \textit{increasing injection}
$\alpha\in\Hom_{[p]}([i],[j])$, the diagram
$$
\xymatrix{\Ker(\alpha_X:\Pi X\rightrightarrows X_i)\ar[r]\ar[d]&
\Pi X\ar[d]\\
\Ker(\alpha_{X'}:\Pi X'\rightrightarrows X'_i)\ar[r]&\Pi X'}
$$
is Cartesian. If $i=n+1$, then
$\alpha:[i]\hookrightarrow[j]$ is the identity of $[n+1]$, so the
pair of arrows imposes no condition. The respective kernels are
$\Pi X$ and $\Pi X'$, which settles this case. Otherwise $i<n+1$
and $X_i\to X'_i$ is an isomorphism, so the diagram is Cartesian in
this case as well (testing an equality in $X_i$ or $X'_i$ amounts to
the same thing)\nde{The proof of Lemma
\hyperref[Vbis.lem3.3.3.3]{3.3.3.3} has been slightly rewritten.}.

Using \hyperref[Vbis.prop3.3.1]{3.3.1}, the relevant morphisms are
therefore morphisms of universal $G$-$1$-cohomological descent. Let
$[X'/X]^p$ be the $(p+1)$-fold fiber product of $X'$ over $X$. By
\hyperref[Vbis.lem3.3.3.3]{3.3.3.3}, it is enough to verify
\hyperref[Vbis.lem3.3.3.2]{3.3.3.2} after a base change
$[X'/X]^p\to X$. After such a base change, the hypotheses of
\hyperref[Vbis.lem3.3.3.2]{3.3.3.2} still hold and, moreover,
$f_{n+1}$ admits a section. Lemma
\hyperref[Vbis.lem3.0.2.4]{3.0.2.4} then completes the proof.

For effective descent, we have the following result.

\begin{enonce*}{Proposition 3.3.4}\etiquette[3.3.4]{Vbis.prop3.3.4}
Let $X:\Delta^{\circ}\to B/S$ be a functor. Suppose that
$X_0\to S$ is a morphism of universal $G^{\circ}$-$2$-descent, as
are the morphisms $X_{n+1}\to(\cosk_n(X))_{n+1}$ for $n=0,1$. Then
the functor $X$ is $G^{\circ}$-$2$-faithful and remains so after
every base change (in other words, $X$ defines an augmentation of
universal effective descent in the sense of
\hyperref[Vbis.def2.4.11]{2.4.11}).
\end{enonce*}

By \cite{sga4.exp.Vbis.bib.4} (7.12), it is enough to show that $i_2^{\ast}(X)$ is
$G^{\circ}$-$2$-faithful; for this we use the following lemmas.

\begin{enonce*}{Lemma 3.3.4.1}\etiquette[3.3.4.1]{Vbis.lem3.3.4.1}
Let\pageoriginale $X$ and $X'$ be two functors
$\mathop{\Delta}_2^{\circ}\to B/S$, and consider a commutative
diagram
\[
\xymatrix@R=1.5cm@C=2cm{
X'_2 \ar[r]\ar@<8pt>[r]\ar@<-8pt>[r]\ar[d]_{k}&X'_1\ar@<4pt>[l]
\ar@<-4pt>[l]\ar@<4pt>[r]\ar@<-4pt>[r]\ar@{=}[d]&X'_0\ar[l]
\ar[r]\ar@{=}[d]&S\ar@{=}[d]\\
X_2 \ar[r]\ar@<8pt>[r]\ar@<-8pt>[r]&X_1\ar@<4pt>[l]
\ar@<-4pt>[l]\ar@<4pt>[r]\ar@<-4pt>[r]&X_0\ar[l]\ar[r]&S
}
\]
such that $k$ is a morphism of $G^{\circ}$-$0$-descent. If $X$ is
$G^{\circ}$-$2$-faithful, then so is $X'$.
\end{enonce*}

This is evident.

\begin{enonce*}{Lemma 3.3.4.2}\etiquette[3.3.4.2]{Vbis.lem3.3.4.2}
Let $X$ and $X'$ be two functors
$\mathop{\Delta}_2^{\circ}\to B/S$, and consider a commutative
diagram
\def\labelstyle{\scriptstyle}
\[
\xymatrix@R=2cm@C=2cm{
X'_2 \ar[r]^{\partial^{1}} \ar@<17pt>[r]^{\partial^{0}}
\ar@<-14pt>[r]_{\partial^{2}} \ar[d]_{f_2}&X'_1\ar@<8pt>[l]
\ar@<-11pt>[l]\ar@<11pt>[r]^{\partial^{0}}\ar@<-8pt>[r]_{\partial^{1}}
\ar[d]^{f_1}&X'_0\ar[l]_{\sigma^{0}}\ar[r]\ar@{=}[d]&S\ar@{=}[d]\\
X_2 \ar[r]\ar@<12pt>[r]\ar@<-12pt>[r]&X_1\ar@<6pt>[l]
\ar@<-6pt>[l]\ar@<6pt>[r]\ar@<-6pt>[r]&X_0\ar[l]\ar[r]&S
}
\]
such that $f_1$ is of $G^{\circ}$-$1$-descent and $f_2$ is of
$G^{\circ}$-$0$-descent. Suppose in addition that
$X'_2\xrightarrow{\sim}(\cosk_1(X'))_2$. If $X$ is
$G^{\circ}$-$2$-faithful, then so is $X'$.
\end{enonce*}

Let $X''_1=\fprod{X'_1}{X'_1}{X_1}$, and let
$\Ysheaf{\partial}_0:X''_1\to X'_1$ and
$\Ysheaf{\partial}_1:X''_1\to X'_1$ be the two projections.

The definition of a coskeleton gives an arrow
$\varphi:X''_1\to X'_2$ such that
\begin{equation*}
\begin{cases}
\partial_0\circ\varphi=\Ysheaf{\partial}_0,\\
\partial_1\circ\varphi=\Ysheaf{\partial}_1,\\
\partial_2\circ\varphi=\sigma_0\partial_1\Ysheaf{\partial}_0
=\sigma_0\partial_1\Ysheaf{\partial}_1{\quoi}.
\end{cases}
\end{equation*}

Now take a descent datum on $X'$: it gives an object on $X_0$, two
objects on $X_1$, and an isomorphism between them on $X'_1$. The
arrow $\varphi$ shows that the two inverse images of these
isomorphisms on $X''_1$ are equal. Since $f_1$ is of
$1$-descent, we obtain an isomorphism between the two objects on
$X_1$; this isomorphism is a descent datum (because
$X'_2\to X_2$ is of $0$-descent), and the result follows.

\begin{enonce*}{Corollary 3.3.5}\etiquette[3.3.5]{Vbis.cor3.3.5}
Suppose that \emph{all arrows of $B$ are flat}. Every hypercovering
of $S$ for the topology of $G$-$2$-cohomological descent whose
objects are representable defines an augmentation of universal
$G$-$2$-cohomological descent.
\end{enonce*}

% END INCLUDED SGA4_English_current_buildable_source_20260906/source/05bis/drafts/SGA4_VBIS_LINES_2586_2816_English.texfrag


% BEGIN INCLUDED SGA4_English_current_buildable_source_20260906/source/05bis/drafts/SGA4_VBIS_LINES_2817_2992_English.texfrag
\section{Examples}\etiquette[4]{Vbis.sec4}

\subsection{Sheaves of abelian groups on topological
  spaces}\etiquette[4.1]{Vbis.subsec4.1}

\setcounter{subsubsection}{-1}
\subsubsection{}\etiquette[4.1.0]{Vbis.subsubsec4.1.0}
In this section, Top denotes the category of topological spaces
[elements of the fixed universe $\Uu$]. We define a category
$\mathcal{E}$ bifibered in duals of topoi over Top by taking as
objects the pairs $(X, F)$, where $X$ is a topological space and $F$
is a sheaf of sets on $X$, and as morphisms the pairs
$(f, \varphi) : (X, F) \to (Y, H)$, where $f : X \to Y$ is a
continuous map and $\varphi : H \to f_{\ast} (F)$ is a morphism of
sheaves. We take for $\mathcal{O}$ the section of $\mathcal{E}^0$
over $\Top^0$ that associates with each topological space the
constant sheaf $\ZZ_X$: we set
$G = \Ysheaf{\Mod}(\mathcal{E}^0, \mathcal{O})$, so that, for every
topological space $X$, $G_X$ identifies with the category of sheaves
of abelian groups on $X$.

Recall that a morphism $f : X \to Y$ is \emph{separated} if the
diagonal in $\fprod{X}{X}{Y}$ is closed. A \emph{proper} morphism is
a separated and universally closed morphism (note that this
definition is more restrictive than Bourbaki's).

The\pageoriginale proof of the ``base-change theorem'' below is
inspired by (\cite{sga4.exp.Vbis.bib.5} II(4.11.1)).

\begin{enonce*}{Theorem 4.1.1}\etiquette[4.1.1]{Vbis.thm4.1.1}
Let $f : X \to Y$ be a proper morphism, let $F$ be an abelian sheaf
on $X$, and consider a Cartesian diagram
\[
\xymatrix@C=2.5cm@R=1.5cm{
X' \ar[r]^{g'} \ar[d]^{f'}& X\ar[d]^{f}\\
Y' \ar[r]^{g}& Y.
}
\]
Then the canonical map \emph{(XII \hyperref[XII.eq4.2]{4.2})} or
\emph{\eqhyperref[Vbis.eq3.2.2.3]{3.2.2.3}}
$$
g^{\ast} R^i f_{\ast} F \xrightarrow{\sim} R^i f'_{\ast} {g'}^{\ast}F
$$
is an isomorphism\nde{Unlike above, these are the derived functors
of the usual functors, and not those arising from the fibered-category
structure $G^\circ\rightarrow\Top^\circ$.}.
\end{enonce*}

This theorem is equivalent to the following corollary (the corollary
is obtained by taking $Y'= (\Point)$; the theorem is obtained by
applying the corollary to $f$ and $f'$).

\begin{enonce*}{Corollary 4.1.2}\etiquette[4.1.2]{Vbis.cor4.1.2}
Let $f : X \to Y$ be a proper morphism, let $F$ be an abelian sheaf
on $X$, and let $y \in Y$. The canonical map
$$
(R^i f_{\ast} F)_y \longrightarrow H^i(f^{-1} (y), F)
$$
is bijective.
\end{enonce*}

By definition, as $U$ ranges over the neighborhoods of $y$, we have
$$
(R^if_{\ast} F)_y = \varinjlim H^i(f^{-1} (U), F){\quoi}.
$$
Since $f$ is closed, the $f^{-1}(U)$ form a fundamental system of
neighborhoods of $f^{-1}(y)$, and the result follows from the
following lemma, (\hyperref[Vbis.lem4.1.3]{4.1.3}).

\begin{enonce*}{Lemma 4.1.3}\etiquette[4.1.3]{Vbis.lem4.1.3}
Let\pageoriginale $X$ be a topological space, let $K \subset X$, and
let $F$ be an abelian sheaf on $X$. Suppose that $K$ is compact and
that any two distinct points of $K$ have disjoint neighborhoods in
$X$. Then, as $U$ ranges over the neighborhoods of $K$, we have
\begin{equation*}
\varinjlim H^i(U, F) \xrightarrow{\sim} H^i(K, F){\quoi}.\tag{4.1.3.1}
\end{equation*}\etiquette[4.1.3.1]{Vbis.eq4.1.3.1}
\end{enonce*}

We first treat the case $i = 0$. In this case, it is clear that
\eqhyperref[Vbis.eq4.1.3.1]{4.1.3.1} is injective. For surjectivity,
we use the following result.

\begin{enonce*}{Lemma 4.1.4}\etiquette[4.1.4]{Vbis.lem4.1.4}
Under the hypotheses of \emph{(\hyperref[Vbis.lem4.1.3]{4.1.3})}, if
$A$ and $B$ are two closed subsets of $K$ and $W$ is a neighborhood
of $A \cap B$ in $X$, then there are neighborhoods $U$ and $V$ of
$A$ and $B$ in $X$ such that $U \cap V \subset W$.
\end{enonce*}

We first treat the case where $A$ consists of a single point $a$.
The assertion is trivial if $a \in B$ (take $U = W$). If
$a \not \in B$, then for every $b \in B$ there are disjoint open
neighborhoods $U_b$ and $V_b$ of $a$ and $b$ in $X$. We take for
$V$ a finite union of the $V_b$ that covers $B$, and for $U$ the
intersection of the corresponding $U_b$.

In the general case, for every $a \in A$, there are open
neighborhoods $U_a$ and $V_a$ of $a$ and $B$ in $X$ such that
$U_a \cap V_a \subset W$. We take for $U$ a finite union of the
$U_a$ that covers $A$, and for $V$ the corresponding intersection
of the $V_a$.

Return to (\hyperref[Vbis.lem4.1.3]{4.1.3}). If
$s \in H^{\circ}(K, F)$, there is an open covering $U_i$ of $K$ in
$X$ and there are sections $s_i \in H^{\circ}(U_i, F)$ such that
$s_i = s$ on $U_i \cap K$. We may suppose that there are only
finitely many $U_i$. Let $(K_i)$ be a closed covering of $K$, with
$K_i \subset K \cap U_i$. Let $W_y$ be the open subset of
$U_i \cap U_j$ on which $s_i = s_j$; we have
$K_i \cap K_j \subset W_{ij}$. Apply
(\hyperref[Vbis.lem4.1.4]{4.1.4}) to all pairs $(K_i, K_j)$ and to
the $W_{ij}$: we find open subsets $V_{ij}$ with
$K_i \subset V_{ij} \subset U_i$ and
$V_{ij} \cap V_{ji} \subset W_{ij}$. Let
$U'_i = \displaystyle\mathop{\cap}_j V_{ij}$: then
$K_i \subset U'_i \subset U_i$, and $s_i = s_j$ on
$U'_i \cap U'_j$. The $s_i$ therefore glue\pageoriginale on the
neighborhood of $K$ given by the union of the $U'_i$, and
\eqhyperref[Vbis.eq4.1.3.1]{4.1.3.1} is surjective (hence bijective)
for $i = 0$\nde{The two meanings of the symbol $i$ should not be
confused.}.

\begin{enonce*}{Lemma 4.1.5}\etiquette[4.1.5]{Vbis.lem4.1.5}
If $F$ is flabby, the sheaf $F | K$ on $K$ is soft.
\end{enonce*}

Let $A \subset K$ be closed in $K$. Every section of $F$ on $A$
extends to a neighborhood, by what precedes. Since $F$ is flabby, it
extends to $X$, and hence in particular to $K$.

We now prove (\hyperref[Vbis.lem4.1.3]{4.1.3}). Let $F^{\ast}$ be a
flabby resolution of $F$. We have
\begin{eqnarray*}
\varinjlim H^i (U, F) &\stackrel{\ndemark}{=}& \varinjlim H^i (\Gamma(U, F^{\ast})) = H^i\varinjlim \Gamma(U,F^{\ast})
\stackrel{\ndemark}{=} H^i \Gamma (K, F^{\ast})\\
& \stackrel{\text{\hyperref[Vbis.lem4.1.5]{(4.1.5)}}}{=} & H^i (K, F){\quoi}.
\end{eqnarray*}
\addtocounter{cnote}{-1}\ndetext{$H^i (\Gamma(U, F^{\ast}))$ denotes here the $i$-th cohomology group of the complex $\Gamma(U,F^\ast)$.}\addtocounter{cnote}{1}
\ndetext{By \hyperref[Vbis.lem4.1.3]{(4.1.3)}, applied with $\Gamma=H^0$}
\begin{enonce*}{Corollary 4.1.6}\etiquette[4.1.6]{Vbis.cor4.1.6}
Let $f : X \to Y$ be a proper surjective morphism of topological
spaces. Then $f$ is of $G$-$1$-cohomological descent.
\end{enonce*}

This follows from (\hyperref[Vbis.cor4.1.2]{4.1.2}) and
(\hyperref[Vbis.prop3.2.4]{3.2.4}).

\begin{enonce*}{Corollary 4.1.7}\etiquette[4.1.7]{Vbis.cor4.1.7}
Let $\Uu = (U_i)_{i \in I}$ be a locally finite closed covering of a
topological space $X$. Then the canonical morphism
$f : \coprod_{i\in I} U_i \to X$ is a morphism of
$G$-$1$-cohomological descent.
\end{enonce*}

Indeed, $f$ is proper and surjective. Applying
(\hyperref[Vbis.prop2.5.6]{2.5.6}), we recover the spectral sequence
of [\cite{sga4.exp.Vbis.bib.5}. II (5.2.4)] (\cf the Introduction).

\begin{enonce*}{Proposition 4.1.8}\etiquette[4.1.8]{Vbis.prop4.1.8}
Let $f : X \to Y$ be a surjective local homeomorphism of topological
spaces. Then $f$ is a morphism of universal
$G$-$1$-cohomological descent.
\end{enonce*}

By localization on $Y$ (\cf (\hyperref[Vbis.prop3.2.4]{3.2.4})), we
are reduced to the case where $f$ has a section.

\begin{enonce*}{Corollary 4.1.9}\etiquette[4.1.9]{Vbis.cor4.1.9}
Let\pageoriginale $\Uu = (U_i)_{i \in I}$ be an open covering of a
topological space $X$. Then the canonical morphism
$f : \coprod_{i \in I} U_i \to X$ is a morphism of universal
$G$-$1$-cohomological descent.
\end{enonce*}

The descent spectral sequence then gives [\cite{sga4.exp.Vbis.bib.5} II (5.4.1)]
(\cf the Introduction).

\begin{enonce*}[remark]{Remark 4.1.10}\etiquette[4.1.10]{Vbis.rem4.1.10}
%%%%%%%%%%%%%%%%%%%
%%%%%%%%%%%%Check that this compiles
%%%%%%%%%%%%%%%%%%%%
The proofs of (\hyperrefext[VIII][VIII.sec9]{9}) show that
``1-cohomological descent'' may be replaced by
``2-cohomological descent'' in all the preceding statements.
\end{enonce*}

\subsection{Quasi-coherent modules on schemes}\etiquette[4.2]{Vbis.subsec4.2}

\setcounter{subsubsection}{-1}

% END INCLUDED SGA4_English_current_buildable_source_20260906/source/05bis/drafts/SGA4_VBIS_LINES_2817_2992_English.texfrag


% BEGIN INCLUDED SGA4_English_current_buildable_source_20260906/source/05bis/drafts/SGA4_VBIS_LINES_2993_3150_English.texfrag
\subsubsection{}\etiquette[4.2.0]{Vbis.subsubsec4.2.0}
Let $(\Sch)_s$ be the category whose objects are the schemes that are
elements of $\Uu$ and whose arrows are the quasi-compact and quasi-separated
morphisms. We define a category $\mathcal{E}$ bifibered in duals of topoi over
$(\Sch)_s$ by taking as objects the pairs $(X,F)$, where $X$ is a scheme and
$F$ a sheaf on $X$ (for the Zariski topology), and as morphisms the pairs
$(f,\varphi):(X,F)\to(Y,G)$, where $f:X\to Y$ is a continuous map and
$\varphi:G\to f_{\ast}(F)$ is a morphism of sheaves. For $\mathcal{O}$ we
take the section of $\mathcal{E}^0$ over $(\Sch)^0_s$ that associates with
every scheme its structure sheaf; for $G$ we take the subcategory of
$\Ysheaf{\Mod}(\mathcal{E}^0,\mathcal{O})$ such that, for every scheme $X$,
$G_X$ is the category of quasi-coherent modules on $X$.

\begin{enonce*}{Proposition 4.2.1}\etiquette[4.2.1]{Vbis.prop4.2.1}
Every faithfully flat, quasi-compact, and quasi-separated morphism
$f:X\to Y$ is a morphism of universal $G$-$2$-cohomological descent.
\end{enonce*}

By [\SGA I (VIII.5.2)], $f$ is a morphism of effective
$G^0$-descent. By [\EGA IV (1.7.21)], we may apply
(\hyperref[Vbis.prop3.2.4]{3.2.4}) to the Cartesian square
\[
\xymatrix@C=3cm@R=1.5cm{
  \fprod{X}{X}{Y} \ar[r] \ar[d]& X \ar[d]^{f}\\
X   \ar[r]_{f}            & Y
}
\]

\subsection{Étale sheaves on schemes}\etiquette[4.3]{Vbis.subsec4.3}
\pageoriginale

\setcounter{subsubsection}{-1}
\subsubsection{}\etiquette[4.3.0]{Vbis.subsubsec4.3.0}
Let $(\Sch)$ be the category of all schemes belonging to $\Uu$. We thus
define a category $\mathcal{E}$ bifibered in duals of topoi over $(\Sch)$ by
taking as objects the pairs $(X,\mathcal{F})$, where $X$ is a scheme and
$\mathcal{F}$ an étale sheaf on $X$, and as morphisms the pairs
$(f,\varphi):(X,\mathcal{F})\to(Y,\mathcal{G})$, where $f:X\to Y$ is a
morphism of schemes and $\varphi:\mathcal{G}\to f_{\ast}(\mathcal{F})$ is a
morphism of sheaves.

Let $n$ be an integer $\geq0$. We denote by $\mathcal{O}_n$ the section of
$\mathcal{E}^0$ over $(\Sch)^0$ that associates with every scheme $X$ the
constant sheaf $(\ZZ/n\ZZ)_X$. Put
$F_n=\Ysheaf{\Mod}(\mathcal{E}^0,\mathcal{O}_n)$, so that, for every scheme
$X$, $F_{n_X}$ identifies with the category of sheaves of $\ZZ/n\ZZ$-modules
on $X$.

\subsubsection{}\etiquette[4.3.1]{Vbis.subsubsec4.3.1}
Now let $(\Sch)_s$ be the subcategory of $(\Sch)$ whose objects are the
schemes that are elements of $\Uu$ and whose arrows are the quasi-compact and
quasi-separated morphisms. Denote by $\mathcal{E}_s$ the category bifibered
in duals of topoi over $(\Sch)_s$ deduced from $\mathcal{E}$ by the base
change $(\Sch)_s\to(\Sch)$; we again denote by $\mathcal{O}_0$ the
restriction of $\mathcal{O}_0$ to $(\Sch)_s$. Let
$G\subset\Ysheaf{\Mod}(\mathcal{E}^0_s,\mathcal{O}_0)$ be the fibered and
cofibered subcategory of
$\Ysheaf{\Mod}(\mathcal{E}^0_s,\mathcal{O}_0)$ such that, for every scheme
$X$, $G_X$ is the category of torsion abelian sheaves on $X$. The reader
will note that condition (\hyperref[Vbis.par2.4.1.1]{2.4.1.1}) is satisfied.

\begin{enonce*}{Proposition 4.3.2}\etiquette[4.3.2]{Vbis.prop4.3.2}
Let $f:X\to Y$ be a \emph{proper surjective} morphism. Then $f$ is a
morphism of universal $F_n$-$2$-cohomological descent for $n\geq1$.
Considered as a morphism of $(\Sch)_s$, $f$ is a morphism of universal
$G$-$2$-cohomological descent.
\end{enonce*}

By\pageoriginale (\hyperrefext[VIII][VIII.thm9.4]{9.4}), $f$ is a morphism
of effective $F^0$-descent, and by (XII.\hyperref[XII.thm5.1]{(5.1)}) we may apply
(\hyperref[Vbis.prop3.2.4]{3.2.4}) to the Cartesian diagram
\[
\xymatrix@C=3cm@R=1.5cm{
  \fprod{X}{X}{Y} \ar[r] \ar[d]& X \ar[d]^{f}\\
X   \ar[r]_{f}            & Y.
}
\]

\begin{enonce*}{Proposition 4.3.3}\etiquette[4.3.3]{Vbis.prop4.3.3}
Let $f:X\to Y$ be a \emph{surjective} morphism of schemes. Suppose that one
of the following hypotheses is satisfied:
\begin{enumerate}
\renewcommand{\theenumi}{\alph {enumi}}
\renewcommand{\labelenumi}{\theenumi)}
\item $f$ is integral.

\item $Y$ is locally Noetherian, and $f$ is universally open and locally of
  finite presentation.

\item $f$ is flat and locally of finite presentation.
\end{enumerate}
Then $f$ is a morphism of $F_n$-$2$-cohomological descent for every
$n\geq0$.
\end{enonce*}

\begin{enumerate}
\renewcommand{\theenumi}{\alph {enumi}}
\renewcommand{\labelenumi}{\theenumi)}
\item Repeat the proof of (\hyperref[Vbis.prop4.3.2]{4.3.2}), replacing
  (XII.\hyperref[XII.thm5.1]{(5.1)}) by ((\hyperrefext[VIII][VIII.cor5.6]{5.6})).

\item Since the question is local on $Y$, we may suppose that $Y$ is
  Noetherian. Apply (\EGA IV (14.5.10)). By
  (\hyperref[Vbis.prop3.3.1]{3.3.1}) and a), we reduce to the case in which
  every point $y$ of $Y$ has a neighborhood $U$ such that
  $f^{-1}(U)\to U$ has a section, and the result then follows from
  (\hyperref[Vbis.prop3.2.4]{3.2.4}).

\item We must show that $f$ is a morphism of
  \emph{$F$-$1$-cohomological descent} (we know from
  ((\hyperrefext[VIII][VIII.thm9.4]{9.4})) that $f$ is a morphism of
  effective $F^0$-descent).
\end{enumerate}

By localization on $Y$ (\cf (\hyperref[Vbis.prop3.2.4]{3.2.4})), we may
suppose that $Y$ is affine, with ring $A$. We may then write
$\Spec A=\varprojlim_{i\in I}\Spec A_i$, where $I$ is a filtered ordered set
and $A_i$ is a Noetherian ring for every $i$, with
$X=\varprojlim X_i$, where, for every $i$, $X_i$ is a scheme \emph{flat} and
\emph{locally of finite presentation} over $\Spec A_i$ such that
$f_i:X_i\to\Spec A_i$ is \emph{surjective}; moreover, the diagrams
\[
\xymatrix@C=3cm@R=1.5cm{
X \ar[r]\ar[d]_{f}& X_i \ar[d]^{f_i}\\
Y \ar[r]^{u_i}& \Spec A_i
}
\]
are\pageoriginale Cartesian.

Let $\mathfrak{F}$ be an abelian sheaf on $Y$. We have
$\mathfrak{F}={\varinjlim}_i u^{\ast}_i u_{i\ast}(\mathfrak{F})$, so that,
by a standard passage to the limit, we reduce to the case in which
$\mathfrak{F}$ comes from a sheaf on $\Spec A_i$ for some $i$, and b) gives
the conclusion.

\begin{enonce*}[remark]{Remark 4.3.4}\etiquette[4.3.4]{Vbis.rem4.3.4}
If $f:X\to Y$ is integral, then $R^qf_{\ast}=0$ for $q>0$. It follows
readily that it is unnecessary to restrict to complexes bounded below when
performing descent by means of $f$.
\end{enonce*}

\begin{enonce*}{Corollary 4.3.5}\etiquette[4.3.5]{Vbis.cor4.3.5}
Let $(X_{\alpha}\xrightarrow{u_\alpha}X)_{\alpha\in A}$ be a covering family
for the étale topology. Then
$\coprod_{\alpha\in A}X_{\alpha}\to X$ is a morphism of universal
$F_n$-$2$-cohomological descent for every $n\geq0$.
\end{enonce*}

Indeed, there exists a factorization
\[
\xymatrix@=1.5cm{
  &\displaystyle{\mathop{\coprod}_{a\in A} X_{\alpha}} \ar[rd]&\\
X \ar[ru]         \ar[rr]^{h} &                              & X
}
\]
where $h$ is a surjective étale morphism.


% END INCLUDED SGA4_English_current_buildable_source_20260906/source/05bis/drafts/SGA4_VBIS_LINES_2993_3150_English.texfrag


% BEGIN INCLUDED SGA4_English_current_buildable_source_20260906/source/05bis/drafts/SGA4_VBIS_LINES_3151_3452_English.texfrag
\section{Applications}\etiquette[5]{Vbis.sec5}\pageoriginale

\subsection{Construction of simplicial objects}\etiquette[5.1]{Vbis.subsec5.1}

\setcounter{subsubsection}{-1}
\subsubsection{}\etiquette[5.1.0]{Vbis.subsubsec5.1.0}

Up to and including \hyperref[Vbis.prop5.1.3]{5.1.3}, $H$ is a
category in which nonempty finite inverse limits exist and in which
finite coproducts exist and are disjoint and universal. If
$X:\Delta^0\to H$ is a simplicial object of $H$, we denote by
$d^k_i:X_k\to X_{k-1}$ (\resp by $s^k_i:X_k\to X_{k+1}$), for
$i\in[k]$ and $0\leq k$, the arrows corresponding to the face
operators $\partial^i_k:[k-1]\to[k]$ (\resp to the degeneracy
operators $\sigma^i_k:[k]\to[k+1]$) (\cf \cite{sga4.exp.Vbis.bib.3} II 2).

\begin{enonce*}{Definition 5.1.1}\etiquette[5.1.1]{Vbis.def5.1.1}
A simplicial object $X:\Delta^0\to H$ of $H$ is
$\sigma$-split if there exists a family of subobjects $NX_j$ of
$X_j$ $(j\geq0)$ such that the morphisms
$$
h_i:\coprod_{\substack{\Hom_{\Delta^{-}}([i], [j])\\
j\leq i}} NX_j \longrightarrow X_i
$$
are isomorphisms.
\end{enonce*}

Likewise, a $k$-truncated simplicial object
$X:\Delta^0_k\to H$ is $\sigma$-split if there exist $NX_j$
$(0\leq j\leq k)$ satisfying the preceding condition for
$i\leq k$.

The $NX_j$ are uniquely determined by $X$. Indeed, if $X$ is
$\sigma$-split, the degeneracy operators
$s^{k-1}_\ell:X_{k-1}\to X_k$ are isomorphisms from $X_{k-1}$
onto direct summands of $X_k$, and
$$
NX_k = \mathop{\cap}_\ell\quad (\text{complement of } s^{k-1}_\ell(X_{k-1})) {\quoi}.
$$
For $k=0$, $NX_0=X_0$.

\setcounter{subsubsection}{1}
\subsubsection{}\etiquette[5.1.2]{Vbis.subsubsec5.1.2}
For a $\sigma$-split $(n+1)$-truncated simplicial object $X$ of
$H$, we denote by $\alpha_n(X)$ the triple consisting of
\begin{enumerate}
\renewcommand{\theenumi}{\alph {enumi}}
\renewcommand{\labelenumi}{\theenumi)}
\item the restriction $i^{\ast}_n(X)$ of $X$ to $\Delta^0_n$;

\item $NX_{n+1}$;\pageoriginale

\item the evident map from $NX_{n+1}$ to $(\cosq_n(X))_{n+1}$.
\end{enumerate}

This triple $(Y,N,\beta)$ satisfies the following condition:

\noindent
$(\ast)$ $Y$ is a $\sigma$-split $n$-truncated simplicial object of
$H$, and $\beta$ is a map from $N$ to $(\cosq Y)_{n+1}$.

\begin{enonce*}{Proposition 5.1.3}\etiquette[5.1.3]{Vbis.prop5.1.3}
Let $(Y,N,\beta)$ be a triple satisfying $(\ast)$ above.
\begin{enumerate}
\renewcommand{\theenumi}{\roman {enumi}}
\renewcommand{\labelenumi}{(\theenumi)}
\item Up to unique isomorphism, there exists a unique $\sigma$-split
  $X:\Delta^0_{n+1}\to H$ such that
  $\alpha_n(X)\simeq(Y,N,\beta)$.

\item To give a morphism $f$ from $X$ to a truncated simplicial
  object $Z:\Delta^0_{n+1}\to H$ is equivalent to giving
\begin{enumerate}
\renewcommand{\labelenumii}{\theenumii)}
\item a morphism $f':Y\to i^{\ast}_n(Z)$;

\item a morphism $f'':N\to Z_{n+1}$ such that the diagram
\[
\xymatrix@C=4cm@R=2cm{
N  \ar[r] \ar[d]^{f''}& (\cosq_n Y)_{n+1} \ar[d]^{f'}\\
Z_{n+1} \ar[r]& (\cosq_n Z)_{n+1}
}
\]
commutes.
\end{enumerate}
\end{enumerate}
\end{enonce*}

Let us construct $X$.

Set $Y'_{n+1}=(\cosk_nY)_{n+1}$ and denote by
$$
{d'}^{n+1}_i:Y'_{n+1}\to Y_n
\quad(\textup{resp. }{s'}^n_i:Y_n\to Y'_{n+1})
$$
the face operators (\resp the degeneracy operators) obtained by
functoriality.

\noindent \sisi{$+$}{} Set
$Y_{n+1}=N\coprod(\coprod_{\Hom_{\Delta^{-}_{\ell\leq n}}([n+1],
  [\ell])}NY_{\ell})$, and let $\alpha:Y_{n+1}\to Y'_{n+1}$ be
the arrow whose components are $\beta$ and the arrows
$NY_{\ell}\to Y_{\ell}\to Y'_{n+1}$ for every epimorphism
$[n+1]\to[\ell]$ with $\ell\leq n$.

\noindent \sisi{$+$}{} Let\pageoriginale
$s^n_i:Y_n\simeq
\coprod_{\Hom_{\Delta^{-}_{\ell'\leq n}}([n],[{\ell}'])}NY_{{\ell}'}
\to Y_{n+1}$ be the arrow induced by $\sigma^i_n$, so that
\begin{enumerate}
\renewcommand{\theenumi}{\Roman{enumi}}
\renewcommand{\labelenumi}{(\theenumi)}
\item $\alpha\circ s^n_i={{s}'}^n_i$.

Set

\item $d^{n+1}_i={d'}^{n+1}_i\circ\alpha$.
\end{enumerate}

\noindent \sisi{$+$}{} To show that this defines an object of
$\Hom(\Delta^0_{n+1},H)$, it is enough, by [\cite{sga4.exp.Vbis.bib.3}. II (2.4)], to
verify the formulas
\begin{enumerate}
\renewcommand{\theenumi}{\alph{enumi}}
\renewcommand{\labelenumi}{\theenumi)}
\item $d^n_i\cdot d^{n+1}_j=d^n_{j-1}\cdot d^{n+1}_i$ \qquad $i<j$;

\item $s^n_i\cdot s^{n-1}_j=s^n_{j+1}\cdot s^{n-1}_i$ \qquad $i\leq j$;

\item $d^{n+1}_i\cdot s^n_j=
\begin{cases}
s^{n-1}_{j-1}\cdot d^n_i & i<j,\\
\id & i=j \mbox{ or } i=j+1,\\
s^{n-1}_{j}\cdot d^n_{i-1} & i>j+1.
\end{cases}$
\end{enumerate}

Formulas a) and c) are evident by (I) and (II), since they are known
to hold if $d^{n+1}_i$ is replaced by ${d'}^{n+1}_i$ and $s^n_i$
by ${s'}^n_i$.

Formula b) holds by construction.

We leave it to the reader to verify that the object $X$ thus obtained
satisfies \hyperref[Vbis.prop5.1.3]{5.1.3}.

\noindent \sisi{$+$}{} Note that the functor constructed in the proof
of \hyperref[Vbis.prop5.1.3]{5.1.3}, from the category of triples
satisfying (\hyperref[Vbis.subsubsec5.1.2]{5.1.2}) $(\ast)$ to the
category of $(n+1)$-truncated simplicial objects, is faithful but not
full.

Proposition \hyperref[Vbis.prop5.1.3]{5.1.3} provides a very convenient
method for constructing simplicial objects of $H$ by induction. In
the remainder of this section, we formalize this observation in the
form that will be used in \hyperref[Vbis.subsec5.2]{5.2} and
\hyperref[Vbis.subsec5.3]{5.3}.

\setcounter{subsubsection}{3}
\subsubsection{}\etiquette[5.1.4]{Vbis.subsubsec5.1.4}

Let\pageoriginale $E\xrightarrow{\pi}B$ be a functor such that, for
every object $b$ of $B$, the fiber category $E_b$ is nonempty and
satisfies the conditions of \hyperref[Vbis.subsubsec5.1.0]{5.1.0}.

Let $Q$ be a property of an object of $E$, stable under isomorphism.
Suppose that the following conditions hold:
\begin{enumerate}
\renewcommand{\theenumi}{\alph{enumi}}
\renewcommand{\labelenumi}{\theenumi)}
\item \emph{For every object $b$ of $B$, there exists an object $x$ of
  $E_b$ satisfying $Q$.}

\item \emph{For every object $b$ of $B$ and every pair $(x,y)$ of
  objects of $E_b$ satisfying $Q$, there exists an object $z$ of $E_b$
  satisfying $Q$ and morphisms $z\to x$, $z\to y$ in $E_b$.}

\item \emph{For every morphism $b'\xrightarrow{t}b$ of $B$ and every
  object $x$ of $E_b$ satisfying $Q$, there exists an object $x'$ of
  $E_{b'}$ satisfying $Q$ and a morphism $x'\to x$ over $t$.}

Let, on the other hand, $P$ be a property of an arrow of $E$ over an
identity, stable under isomorphism. Suppose that the following
conditions hold:

\item \emph{For every object $b$ of $B$ and every object $x$ of $E_b$,
  there exists an arrow $y\to x$ in $E_b$ satisfying $P$.}

\item \emph{For every object $b$ of $B$, every diagram in $E_b$ of
  the form}
\[
\xymatrix{
u \ar[rrrr]^{f} & & &                 &x\\
                & & &z\ar[ru]\ar[rd] &\\
v \ar[rrrr]_{g} & & &               & y
}
\]
\emph{in which $f$ and $g$ satisfy $P$, there is a completion to a
commutative diagram in $E_b$}
\[
\xymatrix@R=1.25cm{
  & u\ar[rrrr]^{f} & & &  & x\\
w\ar[rrrr]^{h}\ar[ru]\ar[rd] &   & & &z\ar[ru]\ar[rd] & \\
  & v\ar[rrrr]_{g} & & &  & y
}
\]
\emph{in which}\pageoriginale $h$ \emph{satisfies $P$.}

\item \emph{For every morphism $b'\xrightarrow{t}b$ of $B$ and every
  morphism $x\xrightarrow{f}y$ of $E_b$ satisfying $P$, there exists
  a morphism $x'\xrightarrow{f'}y'$ of $E_{b'}$ satisfying $P$ and a
  commutative diagram}
\[
\xymatrix@C=3cm@R=1.5cm{
x' \ar[r]^{\alpha} \ar[d]_{f'}& x\ar[d]^{f}\\
y' \ar[r]_{\beta}& y
}
\]
\emph{in which $\alpha$ and $\beta$ are morphisms over $t$.}
\end{enumerate}

\begin{enonce*}{Definition 5.1.5}\etiquette[5.1.5]{Vbis.def5.1.5}
With the notation of \hyperref[Vbis.subsubsec5.1.4]{5.1.4}, we say
that an object $X$ of $\Hom(\Delta^0,E)$ satisfies condition
\emph{(APQ)} if the following conditions hold:
\begin{enumerate}
\item[(i)] The image of $\pi\circ X$ is the constant simplicial object
  defined by an object $b$ of $B$.

\item[(ii)] As an object of $\Hom(\Delta^0,E_b)$, $X$ is
  $\sigma$-split.

\item[(iii)] $X_0$ satisfies condition $Q$.

\item[(iv)] For every integer $n$, the arrow
  $NX_{n+1}\to(\cosq_nX)_{n+1}$ satisfies condition $P$.
\end{enumerate}
\end{enonce*}

\setcounter{subsubsection}{5}
\subsubsection{}\etiquette[5.1.6]{Vbis.subsubsec5.1.6}
Let $E_{(APQ)}$ denote the subcategory of $\Hom(\Delta^0,E)$ whose
objects are the simplicial objects satisfying $(APQ)$ and whose
morphisms are the morphisms $X\xrightarrow{\alpha}Y$ of
$\Hom(\Delta^0,E)$ such that
$\pi(\alpha_n):\pi(X_n)\to\pi(Y_n)$ is the same morphism of $B$ for
every $n$. There is then a canonical functor
$\bar{\pi}:E_{(APQ)}\to B$.

\begin{enonce*}{Proposition 5.1.7}\etiquette[5.1.7]{Vbis.prop5.1.7}
The\pageoriginale functor $\bar{\pi}$ is surjective. Moreover, for
every object $b$ of $B$ and every pair $(X,Y)$ of objects of
$E_{(APQ),b}$, there exists a diagram
\[
\xymatrix@=1.5cm{
  & Z \ar[ld]\ar[rd]&\\
X &   & Y
}
\]
in $E_{(APQ),b}$.
\end{enonce*}

The proof follows immediately from \hyperref[Vbis.prop5.1.3]{5.1.3},
by virtue of the hypotheses made in
\hyperref[Vbis.subsubsec5.1.4]{5.1.4}.

\begin{enonce*}{Proposition 5.1.8}\etiquette[5.1.8]{Vbis.prop5.1.8}
Let $C$ be a category and $F$ a functor $E_{(APQ)}\to C$. We introduce
the following conditions on $F$:
\begin{enumerate}
\renewcommand{\theenumi}{\roman{enumi}}
\renewcommand{\labelenumi}{(\theenumi)}
\item For every object $b$ of $B$, the restriction of $F$ to
  $E_{(APQ),b}$ factors through the \emph{coarse} category
  \footnote{A category is called \emph{coarse} if, for every pair
  $(x,y)$ of objects, $\Hom(x,y)$ consists of exactly one element.}
  defined by the objects of $E_{(APQ),b}$.

\item For every morphism $t:b'\to b$ of $B$ and every pair
  $(X'\to X,Y'\to Y)$ of morphisms over $t$, the following diagram
  commutes:
\[
\xymatrix@C=3cm@R=2cm{
F(X')         \ar[r]  \ar[d]_{\alpha'}& F(X)\ar[d]^{\alpha}\\
F(Y')          \ar[r] & F(Y)
}
\]
where $\alpha$ and $\alpha'$ are the arrows determined by condition
\emph{(i)}.
\end{enumerate}

Then the functor $G\mapsto G\circ\bar{\pi}$ induces an equivalence
between the category $\Hom(B,C)$ and the full subcategory of
$\Hom(E_{(APQ)},C)$ defined by the functors satisfying conditions
\emph{(i) and (ii)}.
\end{enonce*}

The\pageoriginale details of the proof are left to the reader. We
indicate only how to prove that the preceding functor is essentially
surjective.

Let $F:E_{(APQ)}\to C$ be a functor satisfying conditions (i) and
(ii). For every object $b$ of $B$, choose an object $X$ of
$E_{(APQ),b}$ and set $\bar{F}(b)=F(X)$; one then verifies that
$\bar{F}(b)$ varies functorially with $b$.


% END INCLUDED SGA4_English_current_buildable_source_20260906/source/05bis/drafts/SGA4_VBIS_LINES_3151_3452_English.texfrag


% BEGIN INCLUDED SGA4_English_current_buildable_source_20260906/source/05bis/drafts/SGA4_VBIS_LINES_3453_3612_English.texfrag
\subsection{Singular cohomology of a scheme of finite type over a
  field of characteristic zero}\etiquette[5.2]{Vbis.subsec5.2}

In this section, $k$ will denote a field of characteristic zero.

\setcounter{subsubsection}{-1}
\subsubsection{}\etiquette[5.2.0]{Vbis.subsubsec5.2.0}
We shall denote by $\Sch_f /k$ (\resp Ann) the category of schemes
\emph{locally of finite type} over $k$ (\resp complex analytic
spaces). All simplicial objects of $\Sch_f /k$ (\resp Ann) will be
regarded as simplicial objects of Top by means of the canonical
functor $\Sch_f /k \to  \Top$ (\resp $\Ann \to \Top$). We retain the
data of (\hyperref[Vbis.subsubsec4.1.0]{4.1.0}) concerning cohomological descent.

Following the usual notation, we write $S \to S^{\an}$ for the
canonical functor $\Sch_f / \CC \to$ Ann defined in G.A.G.A.

\subsubsection{}\etiquette[5.2.1]{Vbis.subsubsec5.2.1}
Let $X$ be a simplicial object of $\Sch / k$ (\resp $\Ann$): the
usual compatibilities in the construction of the de Rham complex
of a morphism of schemes (\resp of analytic spaces) make it possible
to construct a complex $\Omega_{X/k}^{\Zar.}$ (\resp $\Omega_{X /
  \CC}^{\bigdot}$) of $D^+(\Ysheaf{\Gamma} (\bar{X},
\ZZ))$ such that, for every integer $i$, $\sisi{\RR^+}{R} e^{\ast}_i
(\Omega_{X/k}^{\Zar.})$ (\resp $\sisi{\RR^+}{R} e^{\ast}_i$
$(\Omega^{\bigdot}_{X /\CC})$) is the de Rham complex
$\Omega_{X/k}^{\Zar.}$ (\resp $\Omega^{\bigdot}_{X_i/ \CC}$).

Let $X \xrightarrow{u} C^{\Delta}_S$ be an augmented simplicial
object of $\Sch_f / k$ (\resp Ann). We define $H^{\Zar.}_{DR}(u)$
(\resp $H^{\bigdot}_{DR}(u)$ and $H^{\bigdot}(u, \CC)$) to be the
hypercohomology groups of the complex $\sisi{\RR^+}{R} \bar{u}_+
(\Omega_{X/k}^{\Zar.})$ (\resp $\sisi{\RR^+}{R} \bar{u}_{\ast}
(\Omega_{X/\CC}^{\bigdot})$ and $\sisi{\RR^+}{R} \bar{u} (\CC)$).

\subsubsection{}\etiquette[5.2.2]{Vbis.subsubsec5.2.2}
If\pageoriginale $X \xrightarrow{u} C^{\Delta}_S$ is a semisimplicial
object of $\Sch_f / \CC$, there are three canonical morphisms:
\begin{align*}
H^{\Zar.}_{DR} (u) &\xrightarrow{\alpha} H^{\bigdot}_{DR}(u^{\an})\\
H^{\bigdot} (u^{\an}, \CC) &\xrightarrow{\beta}
H^{\bigdot}_{DR}(u^{\an})\\
H^{\bigdot} (S^{\an}, \CC) &\xrightarrow{\gamma}
H^{\bigdot}(u^{\an}, \CC){\quoi}.
\end{align*}

\begin{enonce*}{Proposition 5.2.3}\etiquette[5.2.3]{Vbis.prop5.2.3}
With the notation of (\hyperref[Vbis.subsubsec5.2.2]{5.2.2}), suppose
that $X_n$ is smooth for every $n$: then $\alpha$ and $\beta$ are
isomorphisms. If, moreover, $X^{\an} \xrightarrow{u^{\an}}
C^{\Delta}_{S^{\an}}$ is a hypercovering of $S^{\an}$ for the
topology of universal $1$-cohomological descent, then $\gamma$ is
an isomorphism.
\end{enonce*}

The assertion that $\alpha$ is an isomorphism follows from
[\cite{sga4.exp.Vbis.bib.6}, Th. $1'$]. The Poincar\'e lemma shows that $\beta$ is an
isomorphism. The final part of the proposition follows from
(\hyperref[Vbis.thm3.3.3]{3.3.3}).

\setcounter{subsubsection}{3}
\subsubsection{}\etiquette[5.2.4]{Vbis.subsubsec5.2.4}
We shall now apply (\hyperref[Vbis.subsubsec5.1.4]{5.1.4}) in the case
where $B = \Sch_f /k$ and $E$ is the category over $B$ such that, for
every scheme $S$, the fiber category $E_S$ is the category of schemes
of $\Sch_f /k$ over $S$, the arrows over an arrow of $B$ being the
obvious ones. We shall say that a morphism $h : X \to Y$ over $S$
satisfies property $P$ if $X$ is smooth over $k$ and if $h$ is a
composite of finitely many morphisms satisfying one of the following
conditions:
\begin{enumerate}
\renewcommand{\theenumi}{\roman{enumi}}
\renewcommand{\labelenumi}{(\theenumi)}
\item it is the coproduct of a family of \'etale surjective morphisms;
\item it is the coproduct of a family of proper surjective morphisms.
\end{enumerate}
We shall say that a scheme $X$ of $E_S$ satisfies property $Q$ if the
structural morphism $X \to S$ satisfies property $P$.

We leave it to the reader to verify, using resolution of singularities,
that all the conditions of (\hyperref[Vbis.subsubsec5.1.4]{5.1.4}) are satisfied.

The\pageoriginale groups $H^{\Zar.}_{DR}(u)$ define a contravariant
functor from the category $E_{(APQ)}$ to the category of graded
abelian groups, denoted by $H^{\Zar.}_{DR}$.

\begin{enonce*}{Proposition 5.2.5}\etiquette[5.2.5]{Vbis.prop5.2.5}
The functor $H^{\Zar.}_{DR}$ satisfies the conditions of
(\hyperref[Vbis.prop5.1.8]{5.1.8}).
\end{enonce*}

By the Lefschetz principle, we are reduced to the case $k = \CC$.
In that case, for every morphism $S'\xrightarrow{t} S$ and every pair
$(u'\to u$, $v'\to v)$ of morphisms over $t$, there is a commutative
diagram:
\[
\xymatrix@C=3.5cm@R=2cm{
H^{\Zar.}_{DR}(u') \ar[r] \ar[d]_{\gamma^{-1}.\beta^{-1}.\alpha}&  H_{DR}^{\Zar.}(u)\ar[d]^{\gamma^{-1}.\beta^{-1}.\alpha}\\
H^{\bigdot}(S'^{\an},\CC)
\ar[r]^{H^{\bigdot}(t^{\an},\CC)}
& H^{\bigdot}(S^{\an},
\CC) \\
H^{\Zar.}_{DR}(v') \ar[r] \ar[u]^{\gamma^{-1}.\beta^{-1}.\alpha}& H_{DR}^{\Zar.}(v) \ar[u]_{\gamma^{-1}.\beta^{-1}.\alpha}
}
\]
which completes the proof.

\begin{enonce*}{Definition 5.2.6}\etiquette[5.2.6]{Vbis.def5.2.6}
The contravariant functor defined by
(\hyperref[Vbis.prop5.1.8]{5.1.8}) from $H^{\Zar.}_{DR}$ will be
denoted by $H^{\bigdot} (\ast, k)$. For every scheme $S$ locally of
finite type over $k$, the graded abelian group $H^{\bigdot}(S, k)$
will be called the \emph{singular cohomology of} $S$.
\end{enonce*}

\setcounter{subsubsection}{6}
\subsubsection{}\etiquette[5.2.7]{Vbis.subsubsec5.2.7}
We conclude this section with an illustration of the general
principles introduced in it. Recall first that if $S$ is a complex
analytic space, the \emph{de Rham cohomology} of $S$, denoted by
$\HH^{\bigdot}_{DR}(S)$, is by definition the hypercohomology of the
de Rham complex $\Omega^{\bigdot}_{S / \CC}$.

\begin{enonce*}{Proposition 5.2.8}\etiquette[5.2.8]{Vbis.prop5.2.8}
\emph{(Blum--Herrera)}.\pageoriginale Let $S$ be a scheme locally of
finite type over $\CC$. The canonical arrow
$$
H^{\bigdot} (S^{\an}, \CC) \longrightarrow
\HH^{\bigdot}_{DR}(S^{\an})
$$
has a retraction (in the category of graded abelian groups) that is
functorial in $S$.
\end{enonce*}

By (\hyperref[Vbis.prop5.1.8]{5.1.8}), it is enough to consider the
morphism $H^{\bigdot}(+, \CC) \circ \bar{\pi} \to
\HH^{\bigdot}_{DR} \circ \bar{\pi}$ obtained from the preceding one
by composition with $\bar{\pi}$.

Using, for every augmentation $X \xrightarrow{u} C^{\Delta}_S$, the
canonical morphism $\LL \bar{u}^{+}(\Omega^{\bigdot}_{S/C}) \to
\Omega^{\bigdot}_{X/C}$, one defines a functorial morphism
$\HH^{\bigdot}_{DR} \circ \bar{\pi} \to \HH^{\bigdot}_{DR}$ so that
the following diagram is commutative:
\[
\xymatrix@=2cm{
H^{\bigdot}_{DR}  & H^{\bigdot}_{DR} \circ \bar{\pi} \ar[l]\\
H^{\bigdot}(\ast,\CC)_{\circ}\bar{\pi}
\ar[u]_{\text{\rotatebox{90}{$\backsim$}}}^{\beta \circ \gamma}
\ar[ru]&
}
\]
which gives the desired retraction.

\begin{enonce*}[remark]{Remark 5.2.9}\etiquette[5.2.9]{Vbis.rem5.2.9}
Statement (\hyperref[Vbis.prop5.2.8]{5.2.8}) remains valid for every
analytic space once resolution of singularities is available in the
analytic setting.
\end{enonce*}


% END INCLUDED SGA4_English_current_buildable_source_20260906/source/05bis/drafts/SGA4_VBIS_LINES_3453_3612_English.texfrag


% BEGIN INCLUDED SGA4_English_current_buildable_source_20260906/source/05bis/drafts/SGA4_VBIS_LINES_3613_3730_English.texfrag
\subsection{Mixed Hodge theories}\etiquette[5.3]{Vbis.subsec5.3}

We give here a technical result that plays a key role in \cite{sga4.exp.Vbis.bib.2}.

$k$ is a field of characteristic zero.

\subsubsection{}\etiquette[5.3.1]{Vbis.subsubsec5.3.1}
Let $\Sch_{\fs} /k$ be the category of separated schemes of finite
type over $k$; denote by $E \xrightarrow{\pi} \Sch_{\fs}/k$ the
category over $\Sch_{\fs}/k$ defined as follows:

\noindent \sisi{$+$}{} If\pageoriginale $S$ is an object of
$\Sch_{\fs} /k$, an object of $E_S$ is a triple $(\bar{X}, X, i)$,
where $\bar{X}$ is a reduced projective scheme over $k$, $X$ is a
proper scheme over $S$, and $i : X \to \bar{X}$ is an open immersion
over $k$ such that $i(X)$ is dense in $\bar{X}$.

\noindent \sisi{$+$}{} If $t : S' \to S$ is a morphism of schemes,
$(\bar{X}, X, i)$ is an object over $S$, and $(\bar{X}', X', i')$ is
an object over $S'$, a morphism from $(\bar{X}, X, i)$ to
$(\bar{X}', X', i')$ over $t$ is a pair $(\bar{h}, h)$, where
$\bar{h} : \bar{X}'\to \bar{X}$ is a morphism over $k$ and
$h : X'\to X$ is a morphism over $S$, such that the following diagram
is commutative.
\[
\xymatrix@C=3.5cm@R=2cm{
X' \ar[r]^{i'} \ar[d]_{h}& \bar{X}'\ar[d]^{\bar{h}}\\
X  \ar[r]^{i}& \bar{X}.
}
\]
N.B.: the properness hypotheses imply that
$\bar{h}^{-1}(i(X)) = i(X')$.

\begin{enonce*}{Lemma 5.3.2}\etiquette[5.3.2]{Vbis.lem5.3.2}
For every object $S$ of $\Sch_{\fs} / k$, the category $E_S$ is
nonempty, has nonempty projective limits, and has finite coproducts
that are disjoint and universal.
\end{enonce*}

The fact that the category $E_S$ is nonempty follows from Chow's
lemma.

We shall describe only the construction of the direct product of two
objects $(\bar{X}, X, i)$ and $(\bar{X}', X', i')$ of $E_S$.

The canonical arrow $j : \fprod{X}{X'}{S} \to
\fprod{\bar{X}}{\bar{X}'}{k}$ is an immersion. Let $T$ be the reduced
closed subscheme of $\fprod{\bar{X}}{\bar{X}'}{k}$ whose underlying
space is the closure of the image of $j$. Then
$(T, {\fprod{X}{X'_{\red}}{S}}, j_{\red})$ has the required universal
property.

\setcounter{subsubsection}{2}
\subsubsection{}\etiquette[5.3.3]{Vbis.subsubsec5.3.3}
Let\pageoriginale $S$ be an object of $\Sch_{\fs} /k$. We shall say
that an object $(\bar{X}, X, i)$ of $E_S$ satisfies property $Q$ if
$\bar{X}$ is smooth over $k$ and $i(X)$ is the complement of a normal
crossings divisor. We shall say that a morphism
$(\bar{h}, h) : (\bar{X}', X', i') \to (\bar{X}, X, i)$ satisfies
property $P$ if $\bar{h}$ is surjective and $(\bar{X}', X', i')$
satisfies property $Q$.

\begin{enonce*}{Proposition 5.3.4}\etiquette[5.3.4]{Vbis.prop5.3.4}
With the notation of (\hyperref[Vbis.subsubsec5.3.3]{5.3.3}), all the
conditions of (\hyperref[Vbis.subsubsec5.1.4]{5.1.4}) are satisfied.
\end{enonce*}

Conditions a), b), d), and e) are readily verified using resolution
of singularities. Moreover, c) follows from f), so it remains to
verify f).

We must show that every commutative diagram
\[
\xymatrix@C=3cm@R=1cm{
              &Z         \ar[r]^{i_Z} \ar[d]_{h}& \bar{Z}\ar[d]^{\bar{h}}\\
              &X \ar[r]^{i_X} \ar[d]& \bar{X}\ar[d]\\
 S'\ar[r]^{t} &S  \ar[r]            & k
}
\]
where $(\bar{X}, X, i_X)$ and $(\bar{Z}, Z, i_Z)$ are objects of
$E_S$ and $\bar{h}$ is surjective can be completed to a commutative
diagram
\[
\xymatrix@C=1.7cm@R=.45cm{
Z' \ar[rr] \ar[rd]_{i_{Z'}} \ar[dd]_{h'}&    & Z \ar[rd]^{i_Z} \ar@{->} '[d][dd]_{h} &\\
         & \bar{Z}' \ar[rr] \ar@{->} '[d][dd]&    & \bar{Z}\ar[dd]^{\bar{h}}\\
X'  \ar[rr]^(.65){h'} \ar[dd]\ar[rd]_{i_{X'}}   &    & X \ar[rd]_{i_X} \ar@{->} '[d][dd]&\\
         & \bar{X}'  \ar@{->} '[d][dd] \ar[rr]&    & \bar{X} \ar[dd]\\
S'\ar[rr]^(.65){t}\ar[rd]        &    & S \ar[rd] &\\
         & k  \ar@{=}[rr]&    & k
}
\]
where\pageoriginale $(\bar{X}', X', i_{X'})$ and
$(\bar{Z}', Z', i_{Z'})$ are objects of $E_{S'}$ and $\bar{h'}$ is
surjective.

Consider a diagram
\[
\xymatrix@C=2cm@R=1.5cm@M=.2cm{
T \ar[r]^{\varphi} \ar@{_{(}->}[d]_{i_T}& S'\ar[r]^{t} & S\ar[d]\\
\bar{T} \ar[rr]_{\psi}&  & k\\
}
\]
where $\varphi$ is proper and $\psi$ is projective. By base change we
obtain a diagram
\[
\xymatrix@C=1cm@R=.6cm{
\fprod{Z}{T_{\red}}{S} \ar[rrr] \ar[rd]_{p} \ar[dd]&&    & Z \ar[rd] \ar@{->} '[d][dd] &\\
         & \fprod{\bar{Z}}{\bar{T}_{\red}}{k} \ar[rrr] \ar@{->} '[d][dd]& &   & \bar{Z}\ar[dd]\\
\fprod{X}{T_{\red}}{S}  \ar[rrr] \ar[dd]\ar[rd]_{q} &  &    & X \ar[rd] \ar@{->} '[d][dd]&\\
         &\fprod{\bar{X}}{\bar{T}_{\red}}{k}  \ar@{->} '[d][dd] \ar[rrr]& &   & \bar{X} \ar[dd]\\
T \ar[rd]\ar[r]& S'\ar[rr]        &    & S \ar[rd] &\\
         & \bar{T}  \ar[rrr]&  &  & k
}
\]
where\pageoriginale $p$ and $q$ are immersions. It then suffices to
replace $\fprod{\bar{Z}}{\bar{T}_{\red}}{k}$ and
$\fprod{\bar{X}}{\bar{T}_{\red}}{k}$ by the closed images of $p$ and
$q$ to obtain the required diagram. This completes the proof of
(\hyperref[Vbis.prop5.3.4]{5.3.4}).


% END INCLUDED SGA4_English_current_buildable_source_20260906/source/05bis/drafts/SGA4_VBIS_LINES_3613_3730_English.texfrag


% BEGIN INCLUDED SGA4_English_current_buildable_source_20260906/source/05bis/drafts/SGA4_VBIS_LINES_3731_3755_English.texfrag
\begin{thebibliography}{99}\pageoriginale
\bibitem{sga4.exp.Vbis.bib.1} T. Blum -- M. Herrera : De Rham Cohomology of an
  analytic Space. Inv. Math. vol 7. Fas 4-1969. p.275-296.
\bibitem{sga4.exp.Vbis.bib.2} P. Deligne : Théorie de Hodge III (I.H.E.S).
\bibitem{sga4.exp.Vbis.bib.3} P. Gabriel -- M. Zisman : Homotopy theory and calculus of
  fractions Ergebnisse der Mathematik -- Band 35 -- Springer 1967.
\bibitem{sga4.exp.Vbis.bib.4} J. Giraud : Méthode de la descente : Mémoires de la
  S.M.F. 2. 1964.
\bibitem{sga4.exp.Vbis.bib.5} R. Godement : Théorie des faisceaux.
                                Publications de l'Institut de
                                Mathématique de l'Université de
                                Strasbourg. Hermann 1964.
\bibitem{sga4.exp.Vbis.bib.6} A. Grothendieck : On the De Rham Cohomology of Algebraic
  Varieties, Publications de 1' I.H.E.S ${\rm n^{\circ}}$ 29.
\bibitem{sga4.exp.Vbis.bib.7} R. Hartshorne : Residues and Duality. Springer
  Lecture Notes ${\rm n^{\circ}}$ 20.
\bibitem{sga4.exp.Vbis.bib.8} G. Segal : Classifying spaces and Spectral sequences
  Publications de 1'I.H.E.S. ${\rm n^{\circ}}$ 34.
\bibitem{sga4.exp.Vbis.bib.9} J.L.Verdier : Catégories dérivées - Etat 0 -
  (I.H.E.S)
\end{thebibliography}

\ifx\danslelivre\undefined
\end{document}
\fi

% END INCLUDED SGA4_English_current_buildable_source_20260906/source/05bis/drafts/SGA4_VBIS_LINES_3731_3755_English.texfrag


% END INCLUDED SGA4_English_current_buildable_source_20260906/source/05bis/05bis_English.tex

\setcounter{chapter}{5}
\def\thechapter{\Roman{chapter}}

\makeatletter
\def\input@path{{../06/}}
\makeatother
% BEGIN INCLUDED SGA4_English_current_buildable_source_20260906/source/06/06_English.tex
\ifx\danslelivre\undefined
% Release engine: XeLaTeX (the LuaTeX compatibility layer breaks Xy-pic output).
\documentclass[10pt]{../fichiers-configuration-latex/smfart}
% BEGIN INCLUDED SGA4_English_current_buildable_source_20260906/source/fichiers-configuration-latex/paquets.tex
\usepackage{polyglossia}
%\usepackage[francais]{babel}
\usepackage{fontspec}
\usepackage[math-style=ISO,bold-style=upright,sans-style=upright]{unicode-math}
%\usepackage[mathletters]{ucs}
%\usepackage[utf8]{inputenc}
%\usepackage[T1]{fontenc}
\usepackage[unicode,hypertexnames=false,colorlinks=true,linkcolor=blue,citecolor=Sepia,urlcolor=magenta]{hyperref}
\setmainlanguage{french}
\setmainfont[Ligatures=TeX]{Linux Libertine O}
% variante possible : TeXGyrePagella ?
\setmathfont{XITS Math}

% maths
\usepackage{amsmath}
%\usepackage{amssymb}
%\usepackage{mathrsfs}
%\usepackage{stmaryrd}
%\usepackage{euscript}

% couleurs
\usepackage[usenames,dvipsnames]{xcolor}

% graphiques
\usepackage{graphicx}
\usepackage[all]{xy}
%\usepackage{tikz}
%\usetikzlibrary{matrix,arrows,calc}
%\tikzset{column sep=3em,row sep=2em,text height=1.5ex, text depth=0.25ex,>=angle 90}
%\tikzset{description/.style={fill=white,inner sep=1pt}}

% hyperliens, 'etiquettes
\usepackage{xr-hyper}
\usepackage{../fichiers-configuration-latex/smfthm}
%\ifx\cacheretiquettes\undefined
%\usepackage[inline]{showlabels} \else \fi % pour rendre visibles les 'etiquettes
%\usepackage{../fichiers-configuration-latex/makeidx} % pour faire un index
%\usepackage[pagebackref,unicode]{hyperref}  % pour voir o`u sont cit'ees les r'ef'erences
% hyperref is loaded once above with the cumulative-reader destination policy.
%\hypersetup{breaklinks=true} %permet le retour `a la ligne dans les liens trop longs. Incompatible avec pdflatex
%\renewcommand*{\backref}[1]{{\small [cit\'e page(s)~#1.]}}

\usepackage{xspace}
\usepackage{pifont} % pour expose XVII

% END INCLUDED SGA4_English_current_buildable_source_20260906/source/fichiers-configuration-latex/paquets.tex

% BEGIN INCLUDED SGA4_English_current_buildable_source_20260906/source/fichiers-configuration-latex/numerotation.tex
\renewcommand{\thefootnote}{\roman{footnote}} % pour mettre les notes en bas de page sous la forme (ii) et non (2), qui ressemble trop (?) a un exposant.

\newcounter{pageoriginal}
\newcommand{\pageoriginale}{\refstepcounter{pageoriginal}\marginpar{\footnotesize\textbf{\thepageoriginal}}\xspace}
\let\pageoriginaled\pageoriginale

\usepackage{ifthen}
\newboolean{original}
\def\sisi{\ifthenelse{\boolean{original}}}
\setboolean{original}{false}

\newcounter{cnote}
\def\nde#1{\ifthenelse{\boolean{original}}{}{\stepcounter{cnote}\def\thefootnote{\arabic{cnote}}\protect\footnote[\value{cnote}]{N.D.E.: #1}}}
\def\ndemark{\ifthenelse{\boolean{original}}{}{\stepcounter{cnote}\def\thefootnote{\arabic{cnote}}\protect\footnotemark[\value{cnote}]}}
\def\ndetext#1{\ifthenelse{\boolean{original}}{}{\def\thefootnote{\arabic{cnote}}\protect\footnotetext[\value{cnote}]{N.D.E.: #1}}}

% END INCLUDED SGA4_English_current_buildable_source_20260906/source/fichiers-configuration-latex/numerotation.tex

% BEGIN INCLUDED SGA4_English_current_buildable_source_20260906/source/fichiers-configuration-latex/macros.tex
\setmainlanguage{french}
\setotherlanguage{english}


%% Fixed, engine-independent release identifier.  The inherited dynamic Git
%% probe used Lua io.popen, exposed build-state-dependent output, and printed
%% a visible question mark outside a Git checkout.
\def\showgitstatus{English release r2, 2026-07-23}


% macros pour LaTeX
\renewcommand{\mathsf}{}

% Police différente pour U+2113 SCRIPT SMALL L
\setmathfont[range="2113]{latinmodern-math.otf}

% Give every inherited semantic label its own stable PDF destination instead
% of inheriting whichever automatic equation/page anchor happened to be live.
% The legacy label remains as a compatibility alias used by existing edges.
\makeatletter
\def\etiquette[#1]#2{%
  \Hy@raisedlink{\hyper@anchorstart{sga4.semantic.#2}\hyper@anchorend}%
  \def\@currentHref{sga4.semantic.#2}%
  \def\@currentlabel{#1}%
  \label{#2}%
}
\makeatother
\def\eqhyperref[#1]#2{(\hyperref[#1]{#2})}
\def\pagehyperref[#1]#2{(page \hyperref[#1]{#2})}
\def\hyperrefext[#1][#2]#3{\hyperref[#2]{#1 #3}}
\def\eqhyperrefext[#1][#2]#3{\hyperref[#2]{#1 #3}}     %cf. exposes 17
\def\neweqhyperrefext[#1][#2]#3{\hyperref[#2]{#1 #3}}  %cf. exposes 17

\def\shat#1{\protect\widehat{#1}}
\def\stilde#1{\protect\widetilde{#1}}
%%%%%%%%%%%%%%%%%%


\def\etcn{etc.\xspace}
\def\etco{etc...\xspace}
\def\etc{\sisi{\etco}{\etcn}}
\def\ldotsn{\sisi{...}{\etcn}}

\def\cf{{cf.}\kern.3em}
\def\Cf{{Cf.}\kern.3em}
\def\ie{{i.e.}\kern.3em}
\def\iev{{i.e.},\xspace}
\def\loccito{loc.\kern2pt cit.\xspace}
\def\loccitn{\textit{loc.\kern2pt cit.}\xspace}
\def\loccit{\sisi{\loccito}{\loccitn}}
\def\EGA{\textup{EGA}\kern.2em}
\def\SGA{\textup{SGA}\kern.3em}

\def\quoi{}

\def\quav{\quad,\quad}
\def\quaet{\quad \text{et}\quad}
\def\fibreproduct#1#2#3{#1 ×_{#3} #2}
\let\fprod\fibreproduct
\def\ccheck#1#2{\mbox{${\check{\check{\rm #1}}}^{\lower2pt\hbox{$\scriptstyle #2$}}$}}
\def\sheaf#1{\protect\mathcal{#1}}
\def\vieuxsheaf#1{{#1}} % cf. problème avec sheaf{a\Phi_i} par exemple...
\DeclareMathOperator{\Hom}{Hom}
\DeclareMathOperator{\can}{can}
\DeclareMathOperator{\Ind}{Ind}
\DeclareMathOperator{\ind}{ind}
%\DeclareMathOperator{\Hhomtop}{\mathcal{H}\!\mathit{omtop}}
\DeclareMathOperator{\Hhomtop}{\mathbf{Homtop}}
\DeclareMathOperator{\Point}{Point}
\DeclareMathOperator{\Pro}{Pro}
\DeclareMathOperator{\Top}{Top}
\DeclareMathOperator{\Mod}{Mod}
\DeclareMathOperator{\Cart}{Cart}
%\DeclareMathOperator{\Mmorsite}{\mathcal{M}\!\mathit{orsite}}
\DeclareMathOperator{\Mmorsite}{\mathbf{Morsite}}
\DeclareMathOperator{\loc}{loc}
\DeclareMathOperator{\Loc}{Loc}
\DeclareMathOperator{\Pic}{Pic}
\DeclareMathOperator{\Et}{Et}
\DeclareMathOperator{\card}{card}
\DeclareMathOperator{\Ens}{Ens}
\DeclareMathOperator{\h}{\underline{H}}
\DeclareMathOperator{\Fonct}{Fonct}
\DeclareMathOperator{\Ker}{Ker}
\DeclareMathOperator{\but}{but}
\DeclareMathOperator{\Ob}{Ob}
\DeclareMathOperator{\ob}{ob}
\DeclareMathOperator{\Lib}{Lib}
\DeclareMathOperator{\esj}{esj}
\DeclareMathOperator{\id}{id}
\DeclareMathOperator{\Id}{Id}
\DeclareMathOperator{\ad}{ad}
\DeclareMathOperator{\Iim}{Im}
\DeclareMathOperator{\pr}{pr}
\DeclareMathOperator{\Cat}{Cat}
\DeclareMathOperator{\car}{car}
\DeclareMathOperator{\SysInd}{SysInd}
\DeclareMathOperator{\Sup}{Sup}
\DeclareMathOperator{\Coker}{Coker}
\DeclareMathOperator{\Filt}{Filt}
\DeclareMathOperator{\acc}{acc}
\DeclareMathOperator{\sur}{sur}
\DeclareMathOperator{\cat}{cat}
\DeclareMathOperator{\Ord}{Ord}
\DeclareMathOperator{\Gr}{Gr}
\DeclareMathOperator{\Cov}{Cov}
\DeclareMathOperator{\Caf}{Caf}
\DeclareMathOperator{\coker}{coker}
\DeclareMathOperator{\coim}{coim}
\DeclareMathOperator{\im}{im}
\DeclareMathOperator{\Ouv}{Ouv}
\DeclareMathOperator{\diag}{diag}
\DeclareMathOperator{\Esp}{Esp}
\DeclareMathOperator{\TOP}{TOP}
\DeclareMathOperator{\Unif}{Unif}
\DeclareMathOperator{\Iind}{\mathbf{Ind}}
\DeclareMathOperator{\Hhomcart}{\mathbf{Homcart}}
\DeclareMathOperator{\Hhomtopess}{\mathbf{Homtopess}}
\DeclareMathOperator{\Ppoints}{\mathbf{Points}}
\DeclareMathOperator{\Ppoint}{\mathbf{Point}}
\DeclareMathOperator{\Ttors}{\mathbf{Tors}}
\DeclareMathOperator{\Ssex}{\mathbf{Sex}}
\DeclareMathOperator{\Ffib}{\mathbf{Fib}}
\DeclareMathOperator{\Vvois}{\mathbf{Vois}}
\DeclareMathOperator{\Ppointess}{\mathbf{Pointess}}
\DeclareMathOperator{\scrHom}{\mathbf{Hom}}
\DeclareMathOperator{\sheafPic}{\mathbf{Pic}}
\DeclareMathOperator{\sheafAut}{\mathbf{Aut}}
\DeclareMathOperator{\Hhomtoplocan}{\mathbf{Homtoplocan}}
\DeclareMathOperator{\Hhomtopan}{\mathbf{Homtopan}}
\DeclareMathOperator{\Pptgeom}{\mathbf{Ptgéom}}
\DeclareMathOperator{\Morsite}{Morsite}
\DeclareMathOperator{\Fib}{Fib}
\DeclareMathOperator{\Fais}{Fais}
\DeclareMathOperator{\kar}{kar}
\DeclareMathOperator{\Groupes}{Groupes}
\DeclareMathOperator{\Restr}{Restr}
\DeclareMathOperator{\Isom}{Isom}
\DeclareMathOperator{\Prol}{Prol}
\DeclareMathOperator{\Res}{Res}
\DeclareMathOperator{\Tors}{Tors}
\DeclareMathOperator{\esp}{esp}
\DeclareMathOperator{\es}{es}
\DeclareMathOperator{\od}{od}
\DeclareMathOperator{\gr}{gr}
\DeclareMathOperator{\sob}{sob}
\DeclareMathOperator{\dfn}{dfn}
\DeclareMathOperator{\Sch}{Sch}
\DeclareMathOperator{\ouv}{ouv}
\DeclareMathOperator{\cons}{cons}
\DeclareMathOperator{\Nerf}{Nerf}
\DeclareMathOperator{\Points}{Points}
\DeclareMathOperator{\Espsob}{Espsob}
\DeclareMathOperator{\Homtop}{Homtop}
\DeclareMathOperator{\Pt}{Pt}
\DeclareMathOperator{\Quot}{Quot}
\DeclareMathOperator{\glob}{glob}
\DeclareMathOperator{\Homfib}{Homfib}
\DeclareMathOperator{\Qouv}{Qouv}
\DeclareMathOperator{\Quotouv}{Quotouv}
\DeclareMathOperator{\ens}{ens}
\DeclareMathOperator{\ab}{ab}
\DeclareMathOperator{\an}{an}
\DeclareMathOperator{\fs}{fs}
\DeclareMathOperator{\red}{red}
\DeclareMathOperator{\Ab}{Ab}
\DeclareMathOperator{\cosq}{cosq}
\DeclareMathOperator{\Ann}{Ann}
\DeclareMathOperator{\Zar}{Zar}


%%%%%%%%%%%%
%%%%% SGA4-2
\DeclareMathOperator{\Sex}{Sex}
\DeclareMathOperator{\vois}{vois}
\DeclareMathOperator{\Vois}{Vois}
\DeclareMathOperator{\Ext}{Ext}
\DeclareMathOperator{\ext}{ext}
\DeclareMathOperator{\Fer}{Fer}
\DeclareMathOperator{\Add}{Add}
\DeclareMathOperator{\Cocart}{Cocart}
\DeclareMathOperator{\cosk}{cosk}
\DeclareMathOperator{\plat}{plat}

\DeclareMathOperator{\faisc}{faisc}
\DeclareMathOperator{\sk}{sk}
\DeclareMathOperator{\ch}{ch}

\def\HH{\mathbf{H}}
\def\qc{\mathrm{qc}}
\def\et{\mathrm{et}}
\def\coh{\mathrm{coh}}

\def\Hvr#1{%
\begin{array}[b]{l}
{\kern2pt\scriptstyle r}\\[-8pt]
{\kern1pt\scriptstyle \vee}\\[-5pt]
H^{#1}
\end{array}}

\def\Hvv#1{%
\begin{array}[b]{l}
\kern2pt\rotatebox{90}{${\scriptscriptstyle\ll}$}\\[-6pt]
H^{#1}
\end{array}}

\def\Hvi#1{%
\begin{array}[b]{l}
{\kern0.2pt\scriptscriptstyle \infty}\\[-8pt]
{\kern1pt\scriptstyle \vee}\\[-5pt]
H^{#1}
\end{array}}

\def\Hv#1{%
\begin{array}[b]{l}
{\kern1pt\scriptstyle \vee}\\[-5pt]
H^{#1}
\end{array}}

\DeclareMathOperator{\cart}{cart}
\DeclareMathOperator{\Ffaisc}{\mathbf{Faisc}}


\def\Ev#1{%
\begin{array}[b]{l}
{\kern1pt\scriptstyle \vee}\\[-5pt]
E'_{#1}
\end{array}}

\DeclareMathOperator{\Cont}{Cont}
\DeclareMathOperator{\Ccont}{\mathscr{C}\!\mathit{ont}}
\DeclareMathOperator{\Site}{Site}
\newcommand\MM{\mathbf{M}}
\def\limtop#1{{\displaystyle{\varprojlim_{#1}\!\mbox{\rm top }}}}
\def\Limtop{\mathop{\varprojlim\!\mbox{\rm top}}}
\DeclareMathOperator{\spec}{spec}
\def\ss{\scriptstyle{\mathscr{S}}}
\DeclareMathOperator{\Gal}{Gal}
\def\smallm{\scriptstyle{\mathscr{M}}}
\DeclareMathOperator{\End}{End}
\DeclareMathOperator{\ass}{ass}
\DeclareMathOperator{\zar}{zar}
\DeclareMathOperator{\Ogg}{Ogg}
\DeclareMathOperator{\res}{res}
\DeclareMathOperator{\tr}{tr}
\DeclareMathOperator{\Supp}{Supp}
\DeclareMathOperator{\codim}{codim}
\DeclareMathOperator{\rk}{rk}
\DeclareMathOperator{\cur}{cur}
\DeclareMathOperator{\cdqc}{cdqc}
\newcommand\DD{\mathbf{D}}
\DeclareMathOperator{\cl}{cl}
\def\JJ{\mathbf{J}}
\def\FF{\mathbf{F}}
\def\BB{\mathbf{B}}
\DeclareMathOperator{\Alget}{Alget}
\DeclareMathOperator{\Ass}{Ass}
\DeclareMathOperator{\supp}{supp}
\DeclareMathOperator{\rad}{rad}
\DeclareMathOperator{\prof}{prof}
\DeclareMathOperator{\cint}{int}
\DeclareMathOperator{\sic}{sic}
\DeclareMathOperator{\Tor}{Tor}
\DeclareMathOperator{\tors}{tors}
\DeclareMathOperator{\Sym}{Sym}
\DeclareMathOperator{\parf}{parf}
\DeclareMathOperator{\tor}{tor}
\DeclareMathOperator{\defn}{defn}
\DeclareMathOperator{\Div}{Div}
\DeclareMathOperator{\Pol}{Pol}
\DeclareMathOperator{\Fex}{Fex}
\DeclareMathOperator{\Lie}{Lie}
\DeclareMathOperator{\tot}{tot}
\DeclareMathOperator{\Aut}{Aut}
\DeclareMathOperator{\HOM}{HOM}
\DeclareMathOperator{\pch}{pch}
\DeclareMathOperator{\Ch}{Ch}
\DeclareMathOperator{\EXT}{EXT}
\DeclareMathOperator{\PIC}{PIC}
\DeclareMathOperator{\Rf}{Rf}
\DeclareMathOperator{\inj}{inj}
\newcommand\VV{\mathbf{V}}

\def\Xv{%
\begin{array}[b]{l}
{\kern1pt\scriptstyle \vee}\\[-5pt]
X
\end{array}}

%%%%%%%
\newcommand\TT{\mathbf{T}}
\newcommand\XX{\mathbf{X}}
\newcommand\EE{\mathbf{E}}
%%%%%%

\def\qedsymbol{\textup{C.Q.F.D.}}
\let\CQFD\qed
\let\cqfd\qed

%%%%%\bigdot
\font\bigsymb=cmsy10 at 4pt
\def\bigdot{{\kern1.2pt\raise 1.5pt\hbox{\bigsymb\char15}}}
%%%%%

\def\Lright{\mathop{\vcenter{\hbox{\hbox{$\longrightarrow$}\lower8pt \hbox{\kern-16pt\hbox{$\longrightarrow$}}}}}}

\def\Aa{{\mathcal A}}
\def\Bb{{\mathcal B}}
\def\Cc{{\mathcal C}}
\def\Dd{{\mathcal D}}
\def\Ee{{\mathcal E}}
\def\Ff{{\mathcal F}}
\def\Gg{{\mathcal G}}
\def\Hh{{\mathcal H}}
\def\Ii{{\mathcal I}}
\def\Jj{{\mathcal J}}
\def\Kk{{\mathcal K}}
\def\Ll{{\mathcal L}}
\def\Mm{{\mathcal M}}
\def\Nn{{\mathcal N}}
\def\Oo{{\mathcal O}}
\def\Pp{{\mathcal P}}
\def\Qq{{\mathcal Q}}
\def\Rr{{\mathcal R}}
\def\Ss{{\mathcal S}}
\def\Tt{{\mathcal T}}
\def\Uu{{\mathcal U}}
\def\Vv{{\mathcal V}}
\def\Ww{{\mathcal W}}
\def\Xx{{\mathcal X}}
\def\Yy{{\mathcal Y}}
\def\Zz{{\mathcal Z}}

\newcommand{\CC}{\mathbf{C}}
\newcommand{\GG}{\mathbf{G}}
\newcommand{\KK}{\mathbf{K}}
\newcommand{\LL}{\mathbf{L}}
\newcommand{\NN}{\mathbf{N}}
\newcommand{\PP}{\mathbf{P}}
\newcommand\QQ{\mathbf{Q}}
\newcommand\RR{\mathbf{R}}
\newcommand\ZZ{\mathbf{Z}}

\def\to{\mathchoice{\longrightarrow}{\rightarrow}{\rightarrow}{\rightarrow}}
\def\from{\mathchoice{\longleftarrow}{\leftarrow}{\leftarrow}{\leftarrow}}
\def\mto{\mathchoice{\longmapsto}{\mapsto}{\mapsto}{\mapsto}}
\def\mfrom{\mathrel{\from\joinrel\mapstochar\,}}
\def\hto{\mathrel{\lhook\joinrel\to}}
\def\hfrom{\mathrel{\from\joinrel\rhook}}

\def\implies{\mathchoice{\Longrightarrow}{\Rightarrow}{\Rightarrow}{\Rightarrow}}
\def\implied{\mathchoice{\Longleftarrow}{\Leftarrow}{\Leftarrow}{\Leftarrow}}
\def\iff{\mathchoice{\Longleftrightarrow}{\Leftrightarrow}{\Leftrightarrow}{\Leftrightarrow}}

\def\To#1{\mathchoice{\xrightarrow{\labelstyle~~#1~}}{\stackrel{#1}{\longrightarrow}}{}{}}
\def\Mto#1{\mathrel{\mapstochar\To{#1}}}
\def\From#1{\mathchoice{\xleftarrow{\labelstyle~~#1~}}{\stackrel{#1}{\longleftarrow}}{}{}}

\newcommand{\isomto}{\overset{\sim}{\longrightarrow}}
\newcommand{\isomfrom}{\overset{\sim}{\longleftarrow}}

\newcommand{\IFF}{\unskip~$\iff$~\ignorespaces}
\newcommand{\THEN}{\unskip~$\implies$~\ignorespaces}

\let\leq\leqslant
\let\geq\geqslant
\let\le\leqslant
\let\ge\geqslant

\def\cf{{cf.}\kern.3em}
\def\Cf{{Cf.}\kern.3em}
\def\ie{{i.e.}\kern.3em}
\def\iev{{i.e.},\xspace}
%\def\loccito{loc.\kern2pt cit.\xspace}
%\def\loccitn{\textit{loc.\kern2pt cit.}\xspace}
%\def\loccit{\sisi{\loccito}{\loccitn}}
\def\EGA{\textup{EGA}\kern.2em}
\def\SGA{\textup{SGA}\kern.3em}

\let\oldS\S
\def\S{\oldS\kern .15em}
\def\ptbl{.\kern .2em}
\def\resp{resp.\kern.3em}
\def\Exp{Exp.\kern.3em}

\let\bar\overline
\let\hat\widehat
\let\tilde\widetilde
\let\dpl\displaystyle
\let\ldots\dots
\let\emptyset\varnothing
\let\skpt\skippointrait
\let\epsilon\varepsilon

\def\sfrac#1#2{{#1}/{#2}}

\DeclareMathOperator{\cd}{cd}
%\DeclareMathOperator{\Hhom}{\mathcal{H}\!\mathit{om}}
\DeclareMathOperator{\Hhom}{\mathbf{Hom}}
%\DeclareMathOperator{\Eext}{\mathcal{E}\!\mathit{xt}}
\DeclareMathOperator{\Eext}{\mathbf{Ext}}
\DeclareMathOperator{\Spec}{Spec}
\DeclareMathOperator{\Tr}{Tr}

\let\emptyset=\varnothing

%%%%%%%%%%%%%%%
%%% expose XVII
%%%%%%%%%%%%%%%

\makeatletter
\let\old@optexpose\optexpose
\newcommand{\empty@optexpose}[1]{}
\newcommand{\paren@optexpose}[1]{\old@optexpose{#1}(}
\newcommand{\refnoexpose}[1]{\let\optexpose\empty@optexpose\ref{#1}\let\optexpose\old@optexpose}%\refnoexpose{X.cor5.2} engendre "5.2" au lieu de "X 5.2"
\newcommand{\neweqref}[1]{\let\optexpose\paren@optexpose\textup{\ref{#1})}\let\optexpose\old@optexpose}%\neweqref{XII.eq4.1} engendre "XII (4.1)" au lieu de "(XII 4.1)"
\makeatother
\newcommand{\refvb}[2][VI${}_{\text{B}}$ ]{\sisi{#1\refnoexpose{#2}}{\ref{#2}}}

\newcommand{\leftexp}[2]{{\vphantom{#2}}^{#1}{#2}}
\newcommand{\AAA}{\mathbf{A}}
\newcommand{\Sf}{\mathfrak{S}}
\DeclareMathOperator*{\quotprojlim}{\text{``$\varprojlim$''}}
\DeclareMathOperator*{\quotinjlim}{\text{``$\varinjlim$''}}
\DeclareMathOperator{\Fl}{F\sisi{\ell}{l}}
\DeclareMathOperator{\Pl}{P\sisi{\ell}{l}}
%\DeclareMathOperator{\cl}{c\sisi{\ell}{l}}
\DeclareMathOperator{\tdf}{tdf}
\DeclareMathOperator{\rg}{rg}
\newcommand{\iInd}{\sisi{I}{i}nd}
\newcommand{\iind}{ind}
\newcommand{\ppro}{pro}
\newcommand{\tTor}{\sisi{T}{t}or}
\newcommand{\ttor}{tor}
\newcommand{\chap}{Chap.}
\newcommand{\Ffex}{\Ff\mathit{ex}}
\newcommand{\sheafk}{\underline{k}}
\newcommand{\sheafn}{\underline{n}}
\newcommand{\sheafGamma}{\underline{\Gamma}}

% Repeated references to an already typeset footnote must not use
% \footnotemark without a matching \footnotetext: Hyperref would then create
% a live annotation to a destination that cannot exist.  This helper keeps the
% current printed footnote mark and links it to an explicit stable note target.
\makeatletter
\newcommand{\sgarepeatedfootnotemark}[1]{%
  \begingroup
  \edef\@thefnmark{\thefootnote}%
  \hyperlink{#1}{\@makefnmark}%
  \endgroup}
\makeatother

\sisi{
\DeclareMathOperator{\simarrow}{\rule{.15mm}{1.8mm}\kern -3.4pt\rightsquigarrow}
\newcommand{\sheafell}{\underline{\ell}}
\newcommand{\sheafZZ}{\underline{\ZZ}}
}{
\newcommand{\nderef}[1]{la N.D.E.~(\ref{#1}) page~\pageref{#1}}
}

%%%%%%%%%%%%%%%%%
%%%%% expose XVII
%%%%%%%%%%%%%%%%%

% END INCLUDED SGA4_English_current_buildable_source_20260906/source/fichiers-configuration-latex/macros.tex


\externaldocument{../01/01_English.tex}
\externaldocument{../02/02_English.tex}
\externaldocument{../03/03_English.tex}
\externaldocument{../04/04_English.tex}
\externaldocument{../05/05_English.tex}
\externaldocument{../05bis/05bis_English.tex}

\setboolean{original}{false}

\begin{document}

\begin{center}
VI. Finiteness Conditions. Fibered Topoi and Sites. Applications
to Questions of Passage to the Limit.\\
A. Grothendieck and J.-L. Verdier
\end{center}

\selectlanguage{english}version: \showgitstatus
\tableofcontents

\else
\chapter[Finiteness conditions and passage to the limit]{Finiteness Conditions. Fibered Topoi and Sites. Applications
to Questions of Passage to the Limit.}
%\addtocontents{toc}{A. Grothendieck and J.-L. Verdier}
\begin{center}
A. Grothendieck and J.-L. Verdier
\end{center}
\fi
\phantomsection\MakeLinkTarget*{sga4.exp.VI}\label{sga4.exp.VI}

% BEGIN INCLUDED SGA4_English_current_buildable_source_20260906/source/06/drafts/SGA4_VI_LINES_44_390_English.texfrag
\setcounter{pageoriginal}{164}

\setcounter{section}{-1}
\section{Introduction}\etiquette[0]{VI.sec0}

In\pageoriginale the present exposé, we study two closely related
types of questions.

First, we undertake a systematic study of finiteness conditions for
topoi, successively examining the notions of \emph{quasi-compactness},
\emph{quasi-separatedness}, and \emph{coherence}
(= quasi-compactness $+$ quasi-separatedness). These are inspired by
the familiar analogous notions from scheme theory. We study them
first for an object $X$ of a topos and for an arrow $u : X \to Y$ of
a topos (\hyperref[VI.sec1]{\S 1}), then for the topos $E$ itself (\hyperref[VI.sec2]{\S 2}), and finally for
a morphism of topoi $f : E \to E'$ (\hyperref[VI.sec3]{\S 3}). \hyperref[VI.sec4]{Section 4} studies these
notions in the special case of a topos obtained by gluing two topoi
(IV 9). For the subsequent sections, the essential facts to retain
from these somewhat technical developments are the following: a
\emph{coherent topos} is a topos equivalent to one of the form
$C^{\sim}$, where $C$ is a small site in which finite projective
limits are representable and every covering family admits a finite
covering subfamily; and a \emph{coherent morphism}
$f : E \to E'$ of coherent topoi is a morphism that can be described
(up to equivalence) as the morphism associated with a morphism of
sites $f : C \to C'$, where $C$ and $C'$ are as above and the
corresponding functor $f^{\ast} : C'\to C$ is left exact.

It is no doubt useful to note that, with the finiteness notions used
here, and especially that of a coherent topos, we resolutely depart
(and for the first time) from the topological spaces familiar to
analysts and ``topologists''\footnote{The quotation marks indicate
that this is a ``\ldots'' unfamiliar with the notion of a topos
(\cf IV \hyperref[IV.subsec0.4]{0.4} for the meaning of the word ``topology'').}. Thus the
topos associated (IV \hyperref[IV.subsec2.1]{2.1}) with a separated topological space $X$ is
coherent only if $X$ is a finite set! In other words, the present
exposé is directed entirely toward the types of topoi arising from
algebraic geometry and algebra. In fact, all the topoi used thus far
in these disciplines (except, of course, in algebraic geometry by
transcendental methods over the field of complex numbers!) turn out
to be locally coherent.

Second, we develop two \emph{theorems on passage to the inductive
limit} for the cohomology of topoi. One
(\hyperref[VI.cor5.2]{5.2}) tells us that the functors
$H^{i}(E, -)$\pageoriginale on a coherent topos commute with
inductive limits of abelian sheaves. The other
(\hyperref[VI.subsec8.7]{8.7}) says essentially that if a topos $X$
is represented as a filtered projective limit of coherent topoi
$X_i$, with coherent transition morphisms
$f_{ij} : X_j \to X_i$, then the cohomology of $X$ (with coefficients
in an arbitrary abelian sheaf) is computed as the inductive limit of
the cohomologies of the $X_i$. To give this statement a precise
meaning, the main point is to specify the notion of a ``projective
limit of topoi,'' which is done in \S\hyperref[VI.sec7]{\S 7}, 8. This is a very useful
geometric notion, undoubtedly also in contexts other than that of
projective limits of schemes (which had served as the model for
M.\textsc{Artin} in his initial definition of this notion). We show,
for example, in (\hyperref[VI.subsec8.4]{8.4}) how the various notions
of being ``localized at a point'' (for a Noetherian space, for a
scheme with its usual topology, or even with its étale topology) can
be unified by defining the localization as the projective limit of
the ``neighborhoods'' of that point---in perfect agreement with the
intuitive idea of localization at a point.

The two limit-passage theorems will be stated again explicitly in the
next exposé in the special setting of étale cohomology of schemes
(\hyperrefext[VII][VII.prop3.3]{3.3} and
\hyperref[VI.prop5.7]{5.7}); in that case, they will moreover be used
constantly throughout the remainder of this Seminar, and practically
everywhere étale cohomology has been used up to the present. The
reader willing to accept these two theorems in the special case in
which they are stated in Exposé VII, and interested exclusively in
applications of étale cohomology, may therefore omit the present
exposé. It is true that the notion of a \emph{constructible object} of
a topos (\ie one whose structural morphism $X \to e$ to the final
object $e$ is coherent) will likewise play an important role
throughout the rest of the Seminar (whereas it plays only a very
secondary role in the present exposé, if only because, for a coherent
topos, it coincides with the notion of a coherent object, which takes
center stage here). But this notion is developed independently in
Exposé IX, in the special case of étale sheaves on schemes, without
using the general study made in the present exposé; in that case it
also exhibits very useful special phenomena that served as the basis
for its treatment in IX (beginning with \hyperref[IX.def2.3]{Definition IX 2.3}). Since the
text of IX predates the present\pageoriginale text of Exposé VI by
five years, yet remains quite usable for the reader, we have confined
ourselves at the end of Exposé IX to proving the equivalence between
the notion of constructibility used in IX and the one introduced
here.

\section{Finiteness conditions for the objects and arrows of a
  topos}\etiquette[1]{VI.sec1}

\begin{enonce*}{Definition 1.1}\etiquette[1.1]{VI.def1.1}
Let $C$ be a site. An object $X$ of $C$ is called quasi-compact if,
for every covering family $(X_i \to X)_{i\in I}$, there is a finite
subset $J$ of the index set $I$ such that the subfamily
$(X_i \to X)_{i \in J}$ is still covering.
\end{enonce*}

Since a $𝒰$-topos $E$ is always regarded as a site (IV \hyperref[IV.subsubsec1.1.1]{1.1.1}), the
preceding definition applies in particular to an object of $E$. This
case is moreover the general one, by the following proposition.

\begin{enonce*}{Proposition 1.2}\etiquette[1.2]{VI.prop1.2}
Let $C$ be a $𝒰$-site, let $\epsilon : C \to C^{\sim}$ be the
canonical functor, and let $X$ be an object of $C$. Then $X$ is a
quasi-compact object of the site $C$ if and only if $\epsilon(X)$ is
a quasi-compact object of the topos $C^{\sim}$.
\end{enonce*}

Suppose that $X$ is quasi-compact; we prove that $\epsilon(X)$ is as
well. Let $(F_i \to \epsilon(X))_{i \in I}$ be an epimorphic family.
For every $i \in I$, choose an epimorphic family
$\epsilon(X_{ij}) \to F_i$, $j \in J_i$. The composite
$\epsilon(X_{ij}) \to \epsilon(X)$ identifies with an element
$\alpha_{ij}$ of $\epsilon(X)(X_{ij})$. In general one cannot assert
that it comes from a morphism $X_{ij} \to X$, but, by the construction
of the associated sheaf $\epsilon(X)$, one knows that, for fixed $i$
and $j$, a covering family $Y_k \to X_{ij}$ can be found such that
the images of $\alpha_{ij}$ in the $\epsilon(X)(Y_k)$ come from
elements of $X(Y_k)$, \ie from morphisms $Y_k \to X$. Thus, after
refining the epimorphic family $\epsilon(X_{ij}) \to F_i$, we may
suppose that the composites
$\epsilon(X_{ij}) \to \epsilon(X)$ come from morphisms
$X_{ij} \to X$.\pageoriginale As $i$ and $j$ vary, the family of
$\epsilon(X_{ij}) \to \epsilon(X)$ is epimorphic, so (II
\hyperref[VI.cor4.4]{4.4}) the family of $X_{ij} \to X$ is covering.
By the hypothesis on $X$, it admits a finite covering subfamily, and
the corresponding finite family of morphisms
$\epsilon(X_{ij}) \to \epsilon(X)$ is again epimorphic. Since these
factor through finitely many of the morphisms
$F_i \to \epsilon(X)$, that finite family is already epimorphic.
This proves that $\epsilon(X)$ is quasi-compact.

Conversely, the fact that quasi-compactness of $\epsilon(X)$ implies
quasi-compactness of $X$ is trivial, by \loccit

The preceding result therefore shows that the notion of
quasi-compactness on sites reduces to the corresponding notion in
topoi.

\begin{enonce*}{Proposition 1.3}\etiquette[1.3]{VI.prop1.3}
Let $C$ be a site and let $(X_i \to X)_{i \in I}$ be a finite
covering family, with the $X_i$ quasi-compact. Then $X$ is
quasi-compact.
\end{enonce*}

To see this, we may assume, by \hyperref[VI.prop1.2]{1.2}, that $C$
is a topos. Let $(X'_j \to X)$ be a covering, \ie epimorphic, family
in $C$. Then, for every $i \in I$, the family obtained by the base
change $X_{i} \to X$ is epimorphic and hence (since $X_i$ is
quasi-compact) admits a finite epimorphic subfamily, corresponding to
a finite subset $J_i \subset J$ of the index set $J$. Set
$J'= \bigcup_{i\in I} J_i$. Then $J'$ is a finite subset of $J$
because $I$ is finite, and the family
$(X'_j \to X)_{j \in {J'}}$ is already epimorphic, since it is so
locally, \cqfd

\begin{enonce*}{Corollary 1.4}\etiquette[1.4]{VI.cor1.4}
Let\pageoriginale $(X_i)_{i \in I}$ be a family of objects of a topos
$E$, and let $X$ be their sum. Then $X$ is quasi-compact if and only
if all the $X_i$ are quasi-compact and all but finitely many of the
$X_i$ are isomorphic to the ``empty object'' $∅_{E}$ of $E$.
\end{enonce*}

Sufficiency follows immediately from \hyperref[VI.prop1.3]{1.3}
(since the sum does not change when the empty summands are omitted).
For necessity, note that a finite subfamily $(X_i)_{i \in J}$ must
cover $X$, which implies that, for $i \in I-J$, one has
$X_i = \fprod{X_{i}}{\amalg_{j=J}}{X}X_{j} =
\fprod{\amalg_{j=J}X_{i}}{X_j}{X} = ∅$ (II
\hyperref[VI.lem4.5.1]{4.5.1}). On the other hand, to show that every
$X_i$ is quasi-compact, note that $X \simeq X_i \amalg {X'}_i$, and
that every covering family $Y_j \to X_i$ of $X_i$ defines a covering
family $Y_j \amalg {X'}_i \to X_i \amalg {X'}_i = X$ of $X$. The
latter admits a finite covering subfamily, which plainly implies that
the corresponding finite family of $Y_j \to X_i$ is also covering;
thus $X_i$ is quasi-compact.

\begin{enonce*}[remark]{Remark 1.5.1}\etiquette[1.5.1]{VI.rem1.5.1}
The property for an object $X$ of a topos $E$ of being quasi-compact
depends only on the induced topos $E_{/X}$, and means that the final
object of the latter is quasi-compact. Likewise, an object $X$ of a
site $C$ is quasi-compact if and only if the final object of the
induced site $C_{/X}$(III \hyperref[III.subsec5.1]{5.1}) is quasi-compact.
\end{enonce*}

\setcounter{subsection}{5}
\setcounter{subsubsection}{1}
\subsubsection{}\etiquette[1.5.2]{VI.subsubsec1.5.2}
Let $X$ be an object of a topos $E$. Then $X$ is quasi-compact if and
only if every filtered increasing family $(X_i)_{i \in I}$ of
subobjects of $X$ that covers $X$ contains some $X_i$ equal to $X$.
This condition therefore depends\pageoriginale only on the ordered
set of subobjects of $X$ (\ie the open subsets of the induced topos
$E_{/X}$).

\subsubsection{}\etiquette[1.5.3]{VI.subsubsec1.5.3}
Consider a generating family $(X_i)_{i \in I}$ of $E$. It follows
from the definitions and from \hyperref[VI.prop1.3]{1.3} that, for an
object $X$ of $E$ to be quasi-compact, it is necessary that it admit
a covering family consisting of finitely many of the $X_i$, \ie that
$X$ be isomorphic to a quotient object of a finite sum of objects of
$E$; this condition is sufficient if the $X_i$ are quasi-compact.
Note also that, under the hypothesis made on $E$, after replacing the
generating family under consideration by a subfamily if necessary,
we may clearly suppose that $I \in 𝒰$. Using the fact that the set of
quotients of an object of $E$ is small (I \hyperref[I.cor7.5]{7.5}), we then find that the
full subcategory $E_{\qc}$ of $E$ formed by the quasi-compact objects
of $E$ is equivalent to a category $\in 𝒰$, \ie the cardinal of the set of
isomorphism classes of objects of $E_{\qc}$ is $\in 𝒰$.

\begin{enonce*}[remark]{Example 1.6.1}\etiquette[1.6.1]{VI.exem1.6.1}
Let $X$ be a topological space. An open subset $U$ of $X$ is a
quasi-compact object of the site $\Ouv(X)$ of open subsets of $X$ if
and only if $U$ is a quasi-compact space for the topology induced by
$X$. More generally, a sheaf $F$ on $X$ is a quasi-compact object of
the topos $\Top(X)$ if and only if the étalé space $X'$ over $X$
associated with $F$ is quasi-compact. (This follows from the preceding
assertion, from \hyperref[VI.rem1.5.1]{1.5}, and from the fact that $\Top (X)_{/F}$ is
equivalent to $\Top(X')$ (IV \hyperref[IV.subsec5.7]{5.7})).
\end{enonce*}

\setcounter{subsection}{6}
\setcounter{subsubsection}{1}
\subsubsection{}\etiquette[1.6.2]{VI.subsubsec1.6.2}
Consider\pageoriginale on the category ($\Sch$) of schemes (elements
of $𝒰$) one of the topologies $T_i$ of \cite{sga4.exp.VI.bib.6}
\hyperref[VI.def6.3]{6.3} (for example, the Zariski topology, the
étale topology, the fppf topology, or the fpqc topology). A scheme is
quasi-compact for this topology if and only if it is a quasi-compact
scheme in the usual sense. (With the notation of \loccit, only the
fact that every family $\in P'$ is finite is used.)

\begin{enonce*}{Definition 1.7}\etiquette[1.7]{VI.def1.7}
Let $E$ be a topos and let $f : X \to Y$ be an arrow of $E$. We say
that $f$ is a quasi-compact morphism if, for every arrow $Y'\to Y$
with $Y'$ quasi-compact, the object
$X'= \fprod{X}{Y'}{Y}$ is also quasi-compact. We say that $f$ is
quasi-separated if the diagonal morphism
$X \to \fprod{X}{X}{Y}$ is quasi-compact. We say that $f$ is coherent
if $f$ is quasi-compact and quasi-separated.
\end{enonce*}

\setcounter{subsection}{7}
\setcounter{subsubsection}{0}
\subsubsection{}\etiquette[1.7.1]{VI.subsubsec1.7.1}
Unwinding the definition of a quasi-separated morphism
$f : X \to Y$, we see that it also means that, for every pair of
arrows $g_1, g_2 : X'\rightrightarrows X$ ``over $Y$,'' \ie such that
$fg_1 = fg_2$, quasi-compactness of $X'$ implies quasi-compactness of
$\Ker(g_1, g_2)$.

\begin{enonce*}{Proposition 1.8}\etiquette[1.8]{VI.prop1.8}
Let $E$ be a topos.
\begin{enumerate}
\renewcommand{\theenumi}{\roman{enumi}}
\renewcommand{\labelenumi}{(\theenumi)}
\item Every isomorphism in $E$ is a coherent morphism, \ie is
  quasi-compact and quasi-separated. The composite of two
  quasi-compact (\resp quasi-separated, \resp coherent) morphisms is
  quasi-compact (\resp quasi-separated, \resp coherent).

\item Let\pageoriginale $f : X \to Y$ and $g : Y' \to Y$ be
  morphisms in $E$, and let
  $f': X'= \fprod{X}{Y'}{Y} \to Y'$ be deduced from $f$ by the base
  change $g$. If $f$ is quasi-compact (\resp quasi-separated, \resp
  coherent), then so is $f'$.

\item Let $f : X \to Y$ and $g : Y \to Z$ be
  morphisms in $E$. If
  $gf$ is quasi-separated, then $f$ is quasi-separated.

\item Let $f : X \to Y$ and $g : Y \to Z$ be morphisms in $E$, with
  $g$ quasi-separated. If $gf$ is quasi-compact (\resp coherent),
  then $f$ is quasi-compact (\resp coherent).
\end{enumerate}
\end{enonce*}

The proofs are well known (\Cf \EGA IV 1 or EGA I, 2nd edition).
Statements (i) and (ii) in the quasi-compact case follow trivially
from the definitions. In the quasi-separated case, stability under
base change in (ii) follows immediately from the definition and the
preceding case, as does the fact that every isomorphism is
quasi-separated. Stability under composition $X \to Y \to Z$ follows
from the diagram
\[\vcenter{
\xymatrix@R=1.5cm@C=1.5cm{
X \ar[d] \ar@/_3.5pc/[dd]^(.2){\Delta_{X/Y}}_{\Delta_{X/Z}}&\\
\fprod{X}{X}{Y} \ar[d]^{u} \ar[rr]&  &Y \ar[d]^{\Delta_{Y/Z}}\\
\fprod{X}{X}{Z}  \ar[rr]&        & \fprod{Y}{Y}{Z}
}}\leqno{(*)}
\]
whose square\pageoriginale is Cartesian. By hypothesis, $\Delta_{Y/Z}$ is
quasi-compact, and hence so is $u$, which is obtained from it by base
change; moreover, $\Delta_{X/Y}$ is quasi-compact, and therefore so is
the composite $u\Delta_{X/Y} = \Delta_{X/Z}$. This proves (i) and
(ii), the coherent case following by combining the two preceding
cases. To prove the first case considered in (iv), consider the
following diagram with Cartesian squares:
\[\vcenter{
\xymatrix@C=2.75cm@R=1.5cm{
  &     X \ar[r]^{gf}          & Z\\
  X \ar[r]^{\Gamma_f=(\id_X,f)}  \ar[d]
  \ar@/_2pc/[rr]_(.55){f\qquad\qquad \fprod{f}{\id_{Y}}{Z}}  &\fprod{X}{Y}{Z} \ar[r]^{\pr_2}\ar[d]\ar[u]^{\pr_1}& Y\ar[u]_{g}\\
Y \ar[r]^{\Delta_{Y/Z}}&  \fprod{Y}{Y}{Z}& \quad .
}}\leqno{(**)}
\]
By hypothesis, $g$ is quasi-separated, so $\Delta_{Y/Z}$ is
quasi-compact and hence $\Gamma_f$ is quasi-compact by base change.
Likewise, by hypothesis $gf$ is quasi-compact, so $\pr_2$ is also
quasi-compact by base change. Consequently,
$f = \pr_2 \Gamma_f$ is quasi-compact as a composite of two
quasi-compact morphisms. The corresponding case of (iv) follows
immediately from the preceding case and (iii), which remains to be
proved. For this, return to diagram $(\ast)$. The hypothesis that
$gf$ is quasi-separated means that
$\Delta_{X/Z} = u \Delta_{X/Y}$ is quasi-compact, while the conclusion
that $f$ is quasi-separated means that $\Delta_{X/Y}$ is
quasi-compact. This follows from (iv) once we have proved that $u$ is
quasi-separated. But $u$ is\pageoriginale plainly a monomorphism, and indeed we
have the following trivial result:

\begin{enonce*}{Corollary 1.8.1}\etiquette[1.8.1]{VI.cor1.8.1}
Every monomorphism is quasi-separated.
\end{enonce*}

Indeed, if $u : X \to Y$ is a monomorphism, then $\Delta_{X/Y}$ is an
isomorphism and is therefore quasi-compact, \cqfd

\begin{enonce*}{Corollary 1.8.2}\etiquette[1.8.2]{VI.cor1.8.2}
Let $f:X \to Y$ and $f':X' \to Y'$ be two $S$-morphisms in $E$. If
$f$ and $f'$ are quasi-compact (\resp quasi-separated, \resp
coherent), then so is
$\fprod{f}{f'}{S} : \fprod{X}{X'}{S} \to \fprod{Y}{Y'}{S}$.
\end{enonce*}

This follows in the usual way by combining (i) and (ii).

\begin{enonce*}{Remark 1.9.1}\etiquette[1.9.1]{VI.rem1.9.1}

% END INCLUDED SGA4_English_current_buildable_source_20260906/source/06/drafts/SGA4_VI_LINES_44_390_English.texfrag


% BEGIN INCLUDED SGA4_English_current_buildable_source_20260906/source/06/drafts/SGA4_VI_LINES_391_677_English.texfrag
Let $E$ be any category, and let $Q$ be a subset of $\ob E$ stable
under isomorphism (thus, giving $Q$ amounts to giving a strictly full
subcategory of $E$), playing the role of the ``quasi-compact''
objects. Suppose that fiber products are representable in $E$. One
can then paraphrase Definition \hyperref[VI.def1.7]{1.7} to define the
notions of $Q$-quasi-compact, $Q$-quasi-separated, and $Q$-coherent
morphisms. Proposition \hyperref[VI.prop1.8]{1.8} and Corollary
\hyperref[VI.cor1.8.1]{1.8.1} then hold, as does
\hyperref[VI.cor1.8.2]{1.8.2} when products of two objects are
representable in $E$.
\end{enonce*}

\setcounter{subsection}{9}
\setcounter{subsubsection}{1}
\subsubsection{}\etiquette[1.9.2]{VI.subsubsec1.9.2}
Let $S$ be an object of $E$, and let $f:X\to Y$ be a morphism of
$E_{/S}$. It follows immediately from \hyperref[VI.rem1.5.1]{Remark 1.5} that $f$ is
quasi-compact (\resp quasi-separated, \resp coherent) if and only if
the corresponding morphism\pageoriginale in $E$ (obtained from $f$
by the ``forget $S$'' functor) has the same property. In particular,
if $f:X\to Y$ is a morphism in $E$, then, by taking $S=Y$ above, we
see that $f$ is quasi-compact (\resp quasi-separated, \resp coherent)
if and only if the structural morphism of the object $(X,f)$ of the
induced topos $E_{/Y}$, with target the final object of that topos, is
quasi-compact (\resp quasi-separated, \resp coherent). One also says
that (by means of the structural morphism $f$) $X$ \emph{is
quasi-compact (\emph{\resp} quasi-separated, \emph{\resp} coherent)
over} $Y$, an expression that may equally well be interpreted as
referring to the topos $E$ or to the induced topos $E_{/Y}$ (in which
$Y$ is regarded as a final object).

\subsubsection{}\etiquette[1.9.3]{VI.subsubsec1.9.3}
An object $X$ of a topos $E$ is called \emph{constructible} if it is
coherent over the final object. Every morphism from one constructible
object to another is coherent (\hyperref[VI.prop1.8]{1.8} (iv)). The
full subcategory $E_{\cons}$ of $E$ formed by the constructible
objects of $E$ is stable under finite inverse limits, as follows
immediately from \hyperref[VI.cor1.8.2]{1.8.2} and
\hyperref[VI.prop1.8]{1.8} (ii).

\subsubsection{}\etiquette[1.9.4]{VI.subsubsec1.9.4}
The notions of \hyperref[VI.def1.7]{1.7} are of real interest only in
topoi $E$ that admit a generating family formed by quasi-compact
objects. In that case, these notions are ``local on the base'' (IV
\hyperref[IV.subsubsec8.5.4]{8.5.4}), as we now show.

\begin{enonce*}{Proposition 1.10}\etiquette[1.10]{VI.prop1.10}
Suppose\pageoriginale that the topos $E$ admits a generating
subcategory $C$ formed by quasi-compact objects, and let $f:X\to Y$
be a morphism in $E$.
\begin{enumerate}
\renewcommand{\theenumi}{\roman{enumi}}
\renewcommand{\labelenumi}{(\theenumi)}
\item The morphism $f$ is quasi-compact if and only if, for every
  morphism $S\to Y$ with $S\in\Ob C$, the object
  $\fprod{X}{S}{Y}$ is quasi-compact.

\item Let $Y_i\to Y$, $i\in I$, be a covering family. The morphism
  $f$ is quasi-compact (\emph{\resp} quasi-separated,
  \emph{\resp} coherent) if and only if, for every $i\in I$, the
  morphism $f_i:X_i=\fprod{X}{Y_i}{Y}\to Y_i$ has the same property.
\end{enumerate}
\end{enonce*}

Let us prove (i). Necessity is immediate from the definition. For
sufficiency, we must prove that, for every quasi-compact object $Y'$
and every morphism $Y'\to Y$, the object
$X'=\fprod{X}{Y'}{Y}$ is quasi-compact. There is a covering family
$S_i\to Y'$ with the $S_i$ in $\Ob C$, and since $Y'$ is
quasi-compact, this covering family may be taken to be \emph{finite}.
By hypothesis the objects
${X'}_i=\fprod{X}{S_i}{Y}\simeq\fprod{X'}{S_i}{Y'}$ are
quasi-compact, and since they cover $X'$ (the $S_i$ covering $Y'$),
$X'$ is quasi-compact by \hyperref[VI.prop1.3]{1.3}.

Let us prove (ii). It is enough to treat the quasi-compact case,
since the quasi-separated case follows from Definition
\hyperref[VI.def1.7]{1.7}, and the coherent case likewise follows by
combining the two preceding cases. The necessity has already been
proved in \hyperref[VI.prop1.8]{1.8} (ii). It remains to show that if
the $f_i$ are quasi-compact, then $f$ is quasi-compact, \ie that, for
every quasi-compact $Y'$ and every morphism $Y'\to Y$, the object
$X'=\fprod{X}{Y'}{Y}$ is quasi-compact. Since $(Y_i\to Y)$ is a
covering family,\pageoriginale there exist objects $S_{\alpha}$ over
$Y$ that cover $Y'$ and such that each composite
$S_{\alpha}\to Y'\to Y$ factors through some $Y_i$. Since $C$ is
generating, the $S_{\alpha}$ may be taken in $\Ob C$, and since $Y'$
is quasi-compact, the covering family may be taken finite. The
$\fprod{X'}{S_{\alpha}}{Y'}$ then form a finite covering family of
$X'$, and by \hyperref[VI.prop1.3]{1.3} we are reduced to proving
that each
$\fprod{X'}{S_{\alpha}}{Y'}=\fprod{X}{S_{\alpha}}{Y}$ is
quasi-compact. Choosing a factorization of $S_{\alpha}\to Y$ through
some $Y_i$, the object in question is also isomorphic to
$\fprod{X_i}{S_{\alpha}}{Y_i}$, which is indeed quasi-compact since
$X_i$ is quasi-compact over $Y_i$ and $S_{\alpha}$ is quasi-compact.

\begin{enonce*}{Corollary 1.11}\etiquette[1.11]{VI.cor1.11}
Let $g:Y\to Z$ be a morphism in the topos $E$, and let
$(f_i:X_i\to Y)_{i\in I}$ be a covering family of morphisms with
target $Y$. Then:
\begin{enumerate}
\renewcommand{\theenumi}{\roman{enumi}}
\renewcommand{\labelenumi}{(\theenumi)}
\item Suppose that $I$ is finite. If the $gf_i$ are quasi-compact,
  then so is $g$.

\item Suppose that the $f_i$ are quasi-compact. If the $gf_i$ are
  quasi-separated, then so is $g$. If $I$ is finite and the $gf_i$
  are coherent, then $g$ is coherent.
\end{enumerate}
\end{enonce*}

Case (i) follows immediately from the definitions and
\hyperref[VI.prop1.3]{1.3}, without any hypothesis on $E$. To prove
(ii), it is enough, by (i), to treat the nonparenthetical case.
Consider diagram \hyperref[VI.prop1.8]{1.8} $(\ast)$ with $X$
replaced by some $X_i$. Since the family of the $f_i$ is covering,
so is the family of the $\fprod{f_i}{\id_Y}{Z}$ obtained from it by
base change. Hence,\pageoriginale by \hyperref[VI.prop1.10]{1.10}
(ii), to prove that $\Delta_{Y/Z}$ is quasi-compact, it is enough to
prove that the $\Gamma_{f_i}$ are quasi-compact. Since $gf_i$ is
quasi-separated by hypothesis, so is
$\pr_2:\fprod{X_i}{Y}{Z}\to Y$, which is obtained from it by base
change. On the other hand, $\pr_2\Gamma_{f_i}=f_i$ is quasi-compact
by hypothesis. Thus $\Gamma_{f_i}$ is quasi-compact by
\hyperref[VI.prop1.8]{1.8} (iv), applied to the composite of
$\Gamma_{f_i}$ and $\pr_2$, \cqfd

\begin{enonce*}{Corollary 1.12}\etiquette[1.12]{VI.cor1.12}
Let $(X_i)_{i\in I}$ be a finite family of objects of $E$, let $X$ be
the coproduct of the $X_i$, let $f_i:X_i\to Y$ be morphisms in $E$,
and let $f:X\to Y$ be the corresponding morphism. The morphism $f$
is quasi-compact (\resp quasi-separated, \resp coherent) if and only
if the $f_i$ have the same property; in the quasi-separated case, the
hypothesis that $I$ be finite may be omitted.
\end{enonce*}

Sufficiency is a special case of \hyperref[VI.cor1.11]{1.11}, taking
into account that the canonical morphisms $X_i\to X$ form a finite
covering family and are quasi-compact morphisms, as follows
immediately from the definition and \hyperref[VI.cor1.4]{1.4}.
Necessity follows from transitivity, \hyperref[VI.prop1.8]{1.8} (i),
taking into account that the morphisms $X_i\to X$ are even coherent
(for they are quasi-separated by \hyperref[VI.cor1.8.1]{1.8.1}).

\begin{enonce*}{Definition 1.13}\etiquette[1.13]{VI.def1.13}
Let $X$ be an object of a topos $E$. We say that $X$ is a
\emph{quasi-separated object of the topos} $E$ if, for every
quasi-compact object $S$ of $E$, every morphism from $S$ to $X$ is
quasi-compact, \ie (\hyperref[VI.def1.7]{1.7}) if, for any two
quasi-compact objects $S,T$ of $E$ and morphisms $S\to X$, $T\to X$,
the fiber product $\fprod{S}{T}{X}$ is always quasi-compact. We say
that $X$ is a coherent object of the topos $E$ if $X$ is
quasi-compact and quasi-separated.
\end{enonce*}

\begin{enonce*}{Proposition 1.14}\etiquette[1.14]{VI.prop1.14}
Let\pageoriginale $f:X\to Y$ be a morphism in a topos $E$.
\begin{enumerate}
\renewcommand{\theenumi}{\roman{enumi}}
\renewcommand{\labelenumi}{(\theenumi)}
\item If $Y$ is quasi-compact, then $f$ quasi-compact implies that
  $X$ is quasi-compact. If $Y$ is quasi-separated, then $X$
  quasi-compact implies that $f$ is quasi-compact. If $Y$ is
  coherent, $f$ is quasi-compact if and only if $X$ is quasi-compact.

\item If $Y$ is quasi-compact (\emph{\resp} quasi-separated,
  \emph{\resp} coherent) and $f$ is quasi-compact
  (\emph{\resp} quasi-separated, \emph{\resp} coherent), then $X$
  is quasi-compact (\emph{\resp} quasi-separated,
  \emph{\resp} coherent).
\end{enumerate}
\end{enonce*}

The first two assertions of (i) are contained immediately in the
definitions of a quasi-compact morphism
(\hyperref[VI.def1.7]{1.7}) and, respectively, a quasi-separated
object (\hyperref[VI.def1.13]{1.13}); the third follows by combining
the first two. The first case of (ii) is merely a restatement, and the
third again follows by combining the first two. It remains to treat
the quasi-separated case. Let $S$ be a quasi-compact object of $E$
and $g:S\to X$ a morphism; we prove that $g$ is quasi-compact. Since
$Y$ is quasi-separated, $fg$ is quasi-compact, and since $f$ is
quasi-separated, it follows (\hyperref[VI.prop1.8]{1.8} (iii)) that
$g$ is quasi-compact, \cqfd

\begin{enonce*}{Corollary 1.15}\etiquette[1.15]{VI.cor1.15}
Every subobject of a quasi-separated object of $E$ is quasi-separated.
Let $(X_i)_{i\in I}$ be a family of objects of $E$, with $I\in 𝒰$;
then $X=\amalg_iX_i$ is quasi-separated (\emph{\resp} coherent) if
and only if the $X_i$ are.
\end{enonce*}

The\pageoriginale first assertion follows from
\hyperref[VI.prop1.14]{1.14} (ii) and
\hyperref[VI.cor1.8.1]{1.8.1}, but is also immediate from the
definition. For the second assertion in the nonparenthetical case,
since the $X_i\to X$ are monomorphisms, it remains only to prove
sufficiency, which follows easily from the definitions and
\hyperref[VI.prop1.3]{1.3}; the parenthetical case follows from the
case just treated and \hyperref[VI.cor1.4]{1.4}.

\begin{enonce*}{Corollary 1.16}\etiquette[1.16]{VI.cor1.16}
Under the conditions of \hyperref[VI.prop1.10]{1.10} (ii), if the
$Y_i$ are coherent, then $f$ is quasi-compact if and only if the
$X_i=\fprod{X}{Y_i}{Y}$ are quasi-compact objects of $E$.
\end{enonce*}

Indeed, by the criterion \hyperref[VI.prop1.10]{1.10} (ii), one must
express that the $f_i:X_i\to Y_i$ are quasi-compact; by
\hyperref[VI.prop1.14]{1.14} (i), this is equivalent to the $X_i$
being quasi-compact.

\begin{enonce*}{Corollary 1.17}\etiquette[1.17]{VI.cor1.17}
Let $E$ be a topos admitting a generating family formed by
quasi-compact objects, and let $(g_i:Y_i\to Y)_{i\in I}$ be a
covering (\resp finite covering) family in $E$, with the $Y_i$
coherent. The object $Y$ is quasi-separated (\resp coherent) if and
only if, for every $i\in I$, $g_i$ is quasi-compact, or equivalently
if, for every pair of indices $i,j\in I$,
$\fprod{Y_i}{Y_j}{Y}$ is a quasi-compact object of $E$.
\end{enonce*}

The parenthetical case follows immediately from the nonparenthetical
case and \hyperref[VI.prop1.3]{1.3}. Necessity of the first condition
is immediate from the definition; let us prove sufficiency. Let
$f:X\to Y$ be a morphism, with\pageoriginale $X$ quasi-compact; we
prove that $f$ is quasi-compact. Since the $g_i$ are quasi-compact,
the $X_i=\fprod{X}{Y_i}{Y}$ are quasi-compact, which implies by
\hyperref[VI.cor1.16]{1.16} that $f$ is quasi-compact. Finally,
expressing the quasi-compactness of the $g_i$ by the same criterion
\hyperref[VI.cor1.16]{1.16} gives the second criterion of
\hyperref[VI.cor1.17]{1.17}. In particular:

\begin{enonce*}{Corollary 1.17.1}\etiquette[1.17.1]{VI.cor1.17.1}
Let $E$ be as in \hyperref[VI.cor1.17]{1.17}, let $X$ be a coherent
object of $E$, let $\Rr\rightrightarrows X$ be an equivalence relation
in $X$, and let $Y=X/\Rr$. Then the following conditions are
equivalent:
\begin{enumerate}
\renewcommand{\theenumi}{\roman{enumi}}
\renewcommand{\labelenumi}{(\theenumi)}
\item $Y$ is coherent;
\item $f:X\to Y$ is quasi-compact;
\item $\Rr$ is quasi-compact.
\end{enumerate}
\end{enonce*}

\begin{enonce*}{Corollary 1.18}\etiquette[1.18]{VI.cor1.18}
Let $E$ be a topos admitting a generating subcategory $C$ formed by
quasi-compact objects, and let $Y$ be an object of $E$. The object
$Y$ is quasi-separated if and only if every morphism $X\to Y$ with
$X\in\ob C$ is quasi-compact.
\end{enonce*}

Necessity is immediate from the definition, since the objects
$X\in\ob C$ are quasi-compact by hypothesis. To prove sufficiency,
let $Z$ be a quasi-compact object of $E$; we prove that every
morphism $Z\to Y$ is quasi-compact. By \hyperref[VI.prop1.10]{1.10}
(i), it is enough to verify that, for every $X\to Y$ with
$X\in\ob C$, the object $\fprod{Z}{X}{Y}$ is quasi-compact. This
follows from the fact that the $X\to Y$ are quasi-compact by
hypothesis and that $Z$ is quasi-compact.

\begin{enonce*}{Corollary 1.19}\etiquette[1.19]{VI.cor1.19}
Let\pageoriginale $E$ be a topos admitting a generating subcategory
formed by quasi-compact objects, and let
$(g_i:Y_i\to Y)_{i\in I}$ be a covering family in $E$, with the $Y_i$
quasi-separated and the $g_i$ quasi-compact. Then $Y$ is
quasi-separated.
\end{enonce*}

Indeed, let $X$ be a quasi-compact object of $E$; we show that every
morphism $f:X\to Y$ is quasi-compact. By
\hyperref[VI.prop1.10]{1.10} (ii), it is enough to prove that the
induced morphisms $f_i:X_i=\fprod{X}{Y_i}{Y}\to Y_i$ are
quasi-compact. Since $g_i$ is quasi-compact, $X_i$ is quasi-compact
(because $X$ is), and since $Y_i$ is quasi-separated, it follows that
$f_i$ is quasi-compact.

\begin{enonce*}[remark]{Remark 1.20.1}\etiquette[1.20.1]{VI.rem1.20.1}
The reader accustomed to the use of the terms ``quasi-compact'' and
``quasi-separated'' in algebraic geometry will expect certain other
relations between the ``absolute'' notions (concerning objects of the
topos $E$) and the ``relative'' notions (concerning arrows of $E$)
considered here; for example, if $f:X\to Y$ is a morphism in $E$ such
that $X$ is quasi-compact and quasi-separated, then $f$ is
quasi-separated. Such properties hold only under restrictive
finiteness conditions on the topos $E$, stronger than those
considered in \hyperref[VI.prop1.10]{1.10}; we shall study these
conditions in the next section.
\end{enonce*}

\setcounter{subsection}{20}
\setcounter{subsubsection}{1}
\subsubsection{}\etiquette[1.20.2]{VI.subsubsec1.20.2}
It follows again from \hyperref[VI.rem1.5.1]{1.5} that whether an object $X$ of the topos $E$
is quasi-separated depends only on the induced topos $E_{/X}$, and
means that the final object of that topos is quasi-separated, \ie
that in $E_{/X}$ the product of two quasi-compact objects is
quasi-compact.


% END INCLUDED SGA4_English_current_buildable_source_20260906/source/06/drafts/SGA4_VI_LINES_391_677_English.texfrag


% BEGIN INCLUDED SGA4_English_current_buildable_source_20260906/source/06/drafts/SGA4_VI_LINES_678_1095_English.texfrag
\subsection{Case of sites}\etiquette[1.21]{VI.subsec1.21}\pageoriginale

Let $C$ be a $𝒰$-site, $E$ the topos of $𝒰$-sheaves on $C$, $f$ an
arrow of $C$, and $X$ an object of $C$. We say that $f$ is a
\emph{quasi-compact morphism} (\resp quasi-separated morphism) \emph{of the
site} $C$ if $\epsilon(f)$ is a quasi-compact morphism (\resp
quasi-separated morphism) of the topos $E$; likewise, we say that $X$ is a
\emph{quasi-separated object of the site} $C$ if $\epsilon(X)$ is a
quasi-separated object of the topos $E$ (compare \hyperref[VI.prop1.2]{1.2}). When $C$ is a
$𝒰$-topos endowed with its canonical topology, the canonical functor
$\epsilon:C\to E$ is an equivalence of categories, and consequently the
preceding definition is compatible with the earlier \hyperref[VI.def1.7]{definitions 1.7} and
\hyperref[VI.def1.13]{1.13}. Note also that the definition under
consideration does not change when the universe $𝒰$ is replaced by another
universe $𝒱$ such that $C$ is also a $𝒱$-site, at least when $C$ admits a
topologically generating family consisting of quasi-compact objects. To see
this, we are immediately reduced to the case $𝒰\subset 𝒱$, so that $E$ is
a $𝒱$-site, and to proving that, for every arrow $f$ of the $𝒰$-topos $E$,
$f$ is quasi-compact (\resp quasi-separated) if and only if it defines a
quasi-compact arrow in the topos $E'$ of $𝒱$-sheaves on $E$, and likewise
that, for every object $X$ of $E$, $X$ is quasi-separated if and only if the
corresponding object of $E'$ is quasi-separated. We already know (without
any hypothesis on $E$) the analogous statement for a quasi-compact object
(\hyperref[VI.prop1.2]{1.2}). The case ``$f$ quasi-compact'' follows from
this, by the criterion \hyperref[VI.prop1.10]{1.10} (i), applied successively
in $E$ and\pageoriginale in $E'$ to the same generating family
$(X_i)_{i\in I}$ of $E$ consisting of quasi-compact objects. This also gives
the case ``$f$ quasi-separated'', which means that the corresponding
diagonal morphism is quasi-compact. Finally, in the case ``$X$
quasi-separated'', we are again reduced to the case ``$X$ quasi-compact''
by the criterion \hyperref[VI.cor1.18]{1.18}.

\subsection{}\etiquette[1.22]{VI.subsec1.22}
{\textit{\textbf{Examples.}}} Return to the example
\hyperref[VI.subsubsec1.6.2]{1.6.2} of the category $(\Sch)$ with a topology
$T_i$. This is a $𝒱$-site for a suitable $𝒱$, and it admits the subcategory
of affine schemes as a topologically generating category consisting of
quasi-compact objects. We leave it to the reader to verify that an arrow
$f:X\to Y$ of this site is quasi-compact (\resp quasi-separated) if and only
if it is a quasi-compact (\resp quasi-separated) morphism of schemes in the
usual sense (\cite{sga4.exp.VI.bib.3} or \cite{sga4.exp.VI.bib.4}); likewise, an object of the site is
quasi-compact (\resp quasi-separated) if and only if it is a quasi-compact
(\resp quasi-separated) scheme in the usual sense of \loccit. It is in fact
appropriate to say, as we do here, \emph{coherent morphism of schemes,
coherent scheme}, instead of ``quasi-compact and quasi-separated morphism of
schemes'', ``quasi-compact and quasi-separated scheme'', as is also done in
\EGA I, $2^{\text{nd}}$ edition.

\setcounter{subsection}{22}
\subsubsection{}\etiquette[1.22.1]{VI.subsubsec1.22.1}

For a given scheme $X\in 𝒰$, one may also consider the topos $\Top(X)$
associated with the underlying topological space, or the topos
$\Top(X_{\et})$ (\resp $\Top(X_{\text{fppf}})$) associated with the site of
étale $X$-schemes (\resp $X$-schemes locally\pageoriginale of finite
presentation) that are elements of $𝒰$, endowed with the étale topology
(\resp fppf topology). It is a $𝒰$-topos admitting a generating family
consisting of coherent objects, corresponding to the affine objects of the
site of definition.

With this understood, the final object of this topos is quasi-compact
(\resp quasi-separated, \resp coherent) if and only if the scheme $X$ is
quasi-compact (\resp quasi-separated, \resp coherent). Likewise, a morphism
of the site of étalé spaces over $X$ (\resp of the site $X_{\text{et}}$ of
étale schemes over $X$, \resp of the site $X_{\text{fppf}}$ of schemes
locally of finite presentation over $X$) is quasi-compact (\resp
quasi-separated, \resp coherent) if and only if it is a morphism of schemes
that is, in the usual sense of \loccit, quasi-compact (\resp quasi-separated,
\resp coherent).

\begin{enonce*}{Theorem 1.23}\etiquette[1.23]{VI.thm1.23}
Let $E$ be a $𝒰$-topos and $X$ an object of $E$.
\begin{enumerate}
\renewcommand{\theenumi}{\roman{enumi}}
\renewcommand{\labelenumi}{(\theenumi)}
\item For $X$ to be quasi-compact, it is necessary and sufficient that, for
  every filtered inductive system $(Y_i)_{i\in I}$ of $E$, with $I\in 𝒰$,
  the natural map
\begin{equation*}
\varinjlim  \Hom(X, Y_i) \to \Hom(X, \varinjlim Y_i)\tag{1.23.1}
\end{equation*}\etiquette[1.23.1]{VI.eq1.23.1}
be injective.

\item Suppose that $E$ admits a generating subcategory $C$ consisting of
  quasi-compact objects. A necessary condition for $X$ to be coherent is
  that, for every datum as in \emph{(i)}, the map \emph{\hyperref[VI.eq1.23.1]{(1.23.1)}} be
  bijective.
\end{enumerate}
\end{enonce*}

\begin{enumerate}
\renewcommand{\theenumi}{\roman{enumi}}
\renewcommand{\labelenumi}{(\theenumi)}
\item \underline{Necessity}.
Suppose\pageoriginale that $X$ is quasi-compact. We must prove that if
  $i\in I$ and $f_i,g_i:X\rightrightarrows Y_i$ are such that the composites
  $f,g$ of $f_i,g_i$ with $Y_i\to Y=\varinjlim Y_j$ are equal, then there
  exists $j\geq i$ such that the composites $f_j,g_j$ with $Y_i\to Y_j$ are
  already equal. Since filtered inductive limits in $E$ commute with finite
  projective limits, and in particular with kernels, we have
$$
\Ker (f, g) = \mathop{\lim}_{j\geq l} \Ker (f_j, g_j){\quoi}.
$$
By hypothesis $\Ker(f,g)=X$, so $X$ is the limit of the filtered increasing
family of its subobjects $\Ker(f_j,g_j)$. Since $X$ is quasi-compact, it is
equal to one of these subobjects; hence there exists $j\geq i$ such that
$f_j=g_j$.

\underline{Sufficiency}. To prove that $X$ is quasi-compact, it is equivalent
(\hyperref[VI.subsubsec1.5.2]{1.5.2}) to prove that, for every filtered
increasing family $(X_i)_{i\in I}$ of subobjects of $X$ whose $\varinjlim$ is
$X$, there exists $i\in I$ such that $X_i=X$. Since every monomorphism in
$E$ is effective (II \hyperref[II.sec4]{4.0}), it follows that
$$
X_i = \Ker (f_i, g_i : X \rightrightarrows X \amalg_{X_{i}} X){\quoi}.
$$
Let $Y_i=X\amalg_{X_i}X$, and consider the inductive system formed by the
$Y_i$. If $Y$ is its inductive limit, then
$Y\simeq X\amalg_{\varinjlim X_i}X=X\amalg_X X=X$.
Consequently, for a fixed $i\in I$, $f_i,g_i$ give the same element of
$\Hom(X,Y)$, so by hypothesis there exists $j\geq i$ such that $f_j=g_j$,
\ie $X_j=X$, \cqfd

\item It remains\pageoriginale to prove that, under the stated conditions,
  every morphism $X\to Y$ comes from a morphism $X\to Y_i$ for a suitable
  $i\in I$. Since $C$ generates $E$ and the $Y_i$ cover $Y$, we may cover
  $X$ by objects $X_{\alpha}$ of $C$ in such a way that each composite
  $X_{\alpha}\to X\to Y$ factors through one of the $Y_i$. Since $X$ is
  quasi-compact, we may suppose that the family of the $X_{\alpha}$ is
  finite, and hence that the morphisms $X_{\alpha}\to Y$ factor through a
  fixed $Y_i$ as morphisms $g_{\alpha i}:X_{\alpha}\to Y_i$. Since $X$ is
  quasi-separated, the $\fprod{X_{\alpha}}{X_{\beta}}{X}$ are quasi-compact.
  Moreover, for every $(\alpha,\beta)$, the composites of $\pr_1$ and
  $\pr_2$ on $\fprod{X_{\alpha}}{X_{\beta}}{X}$ with
  $X_{\alpha},X_{\beta}\to Y_i$ are
\[
\xymatrix{
\fprod{X_{\alpha}}{X_{\beta}}{X} \ar[rd] \ar[dd]& &\\
                            &X_{\alpha} \ar[ld] \ar[rrrr]&   & & & Y_i
\ar[lld]\\
X    \ar[rrr]                       &          & & Y  &
}
\]
such that their composites with $Y_i\to Y$ are equal. By (i), there
  therefore exists $j\geq i$ such that the composites with $Y_i\to Y_j$ are
  already equal. Since the set of pairs $(\alpha,\beta)$ is finite, $j$ may
  be chosen independently of the pair. Consequently, the morphisms
  $g_{\alpha j}$ arise by composition from a morphism $g_j:X\to Y_j$, whose
  composite with $Y_j\to Y$ is moreover the given morphism $X\to Y$, since
  this is so after composition with the covering maps $X_{\alpha}\to X$.
  This completes the proof.
\end{enumerate}

\begin{enonce*}[remark]{Remarks 1.23.2}\etiquette[1.23.2]{VI.rem1.23.2}
The hypothesis on $E$ in (ii) may be replaced by an additional hypothesis
\emph{on} $X$, namely that $X$ admits a fundamental system of covering
families by quasi-compact objects.
\end{enonce*}

\begin{enonce*}{Corollary 1.24}\etiquette[1.24]{VI.cor1.24}
Let\pageoriginale $E$ be a topos, and let $C$ be a strictly full subcategory
of $E$ consisting of coherent objects. Using the fact that the
$𝒰$-inductive limits in $E$ are representable, one obtains \emph{(I \hyperref[I.sec8.6.1]{8.6.1})}
a canonical functor
\begin{equation*}
\Ind (C) \to E{\quoi}.\tag{1.24.1}
\end{equation*}\etiquette[1.24.1]{VI.eq1.24.1}
This functor is fully faithful. If $E$ admits a generating subcategory
consisting of coherent objects, and if $C$ is stable under direct factors,
then the following conditions are equivalent:
\begin{enumerate}
\renewcommand{\theenumi}{\roman{enumi}}
\renewcommand{\labelenumi}{(\theenumi)}
\item The functor \emph{\hyperref[VI.eq1.24.1]{(1.24.1)}} is essentially surjective, \ie it is an
  equivalence of categories.

\item $C$ is a generating subcategory of $E$, stable under finite inductive
  limits.

\item The category $C$ is equal to the full subcategory $E_{\coh}$ of $E$
  consisting of all coherent objects of $E$, and the necessary condition
  \emph{\hyperref[VI.thm1.23]{1.23} (ii)} for an object of $E$ to be coherent
  is also sufficient.
\end{enumerate}
\end{enonce*}

Let $E_{PF}$ be the full subcategory of $E$ defined in I \hyperref[I.cor8.7.6]{8.7.6}. By I \hyperref[I.prop8.7.5]{8.7.5}
a), statement \hyperref[VI.thm1.23]{1.23} (ii) is equivalent to the second
inclusion in
$$
C \subset E_{\coh} \subset E_{PF}{\quoi},
$$
and the composite inclusion $C\supset E_{PF}$ means that \hyperref[VI.eq1.24.1]{(1.24.1)} is fully
faithful. Since\pageoriginale, by hypothesis, the subcategory $E_{\coh}$ of
$E$ is generating, the same is a fortiori true of $E_{PF}$; moreover, finite
projective limits are representable in $E$, so that we may apply I \hyperref[VI.cor8.7.7]{8.7.7}.
Condition (iii) of \hyperref[VI.cor1.24]{1.24} may also be written
$C=E_{\coh}=E_{PF}$ and is equivalent to the condition $C=E_{PF}$, which is
precisely condition (ii bis) of \loccit (taking into account that, by
hypothesis, $E_{\coh}$ and hence $E_{PF}$ is a generating subcategory of
$E$). Likewise, condition (ii) of \hyperref[VI.cor1.24]{1.24} is precisely
condition (iii bis) of \loccit. Thus the equivalence of conditions (i),
(ii bis), and (iii bis) in \loccit gives the equivalence of conditions (i),
(ii), and (iii) of \hyperref[VI.cor1.24]{1.24}, \cqfd

The proof (or, more precisely, the invocation of I \hyperref[I.cor8.7.7]{8.7.7} (ii bis)) also gives
part a) of the

\begin{enonce*}{Corollary 1.24.2}\etiquette[1.24.2]{VI.cor1.24.2}
Let $E$ be a topos admitting a generating subfamily consisting of coherent
objects, and let $E_{PF}$ be the strictly full subcategory of objects $X$ of
$E$ such that the covariant functor $\Hom(X,-)$ represented by $X$ commutes
with filtered inductive limits. Then:
\begin{enumerate}
\renewcommand{\theenumi}{\alph{enumi}}
\renewcommand{\labelenumi}{\theenumi)}
\item The natural functor
$$
\Ind (E_{PF}) \to E
$$
is an equivalence of categories.

\item For an object $X$ of $E$ to belong to $E_{PF}$, it is necessary and
  sufficient that it be isomorphic to a direct factor ($=$ image of a
  projector (\hyperref[I.sec10.6]{I, 10.6})) of the cokernel\pageoriginale of a double arrow
  $X_1\rightrightarrows X_0$, with $X_1$ and $X_0$ coherent.
\item The subcategory $E_{PF}$ of $E$ is stable under finite inductive
  limits, and is equivalent to a small category.
\end{enumerate}
\end{enonce*}

It remains to prove b) and c). The first assertion of c) is immediate from
I \hyperref[I.prop2.8]{2.8}, and implies the ``if'' part of b) by \hyperref[VI.thm1.23]{1.23} (ii).
For the ``only if'' part, use \hyperref[VI.thm1.23]{1.23} (i), which shows
that $X$ is quasi-compact, whence the existence of an epimorphism
$X_0\to X$, with $X_0$ coherent, in view of the fact that $E_{\coh}$
generates $E$ and is stable under finite sums. Let
$R=\fprod{X_0}{X_0}{X}$, so that $X$ identifies with the cokernel of
$R\rightrightarrows X_0$. By a), we have $R=\varinjlim R_i$, the filtered
inductive limit of the coker $(R_i\rightrightarrows X_0)=X_i$. By the
hypothesis $X\in\Ob E_{PF}$, for sufficiently large $i$ the morphism
$u_i:X'_i\to X$ admits a section $v:X\to X'_i$. Let
$w=vu_i:X'_i\to X'_i$, so that $w^2=w$ and
$X\simeq\Coker(w,\id_{X'_i}:X'_i\rightrightarrows X'_i)$. Since $R_i$ is
again quasi-compact, cover it by $X_1\to R_i$, with $X_i$ coherent, whence
$X_i\simeq\Coker(X_1\rightrightarrows X_0)$, which proves b).

Finally, since $E_{PF}\subset E_{\text{qu.cpct}}$, the last assertion of c)
follows from \hyperref[VI.subsubsec1.5.3]{1.5.3}.

\begin{enonce*}{Corollary 1.25}\etiquette[1.25]{VI.cor1.25}
Let $E$ be a topos such that the full subcategory $C$ of $E$ consisting of
coherent objects is generating. Consider the canonical functor
$$
\varphi : \Ind(c) \to E{\quoi}.
$$
The following\pageoriginale conditions are equivalent:
\begin{itemize}
\item[(i)] The functor $\varphi$ is essentially surjective, \ie every object
  of $E$ is a filtered inductive limit of coherent objects.

\item[(i bis)] The functor $\varphi$ is an equivalence of categories.

\item[(ii)] Every finite inductive limit in $E$ of coherent objects is a
  coherent object.

\item[(ii bis)] The cokernel in $E$ of a double arrow of coherent objects is
  a coherent object.

\item[(iii)] The converse of \emph{\hyperref[VI.thm1.23]{1.23} (ii)} holds.
\end{itemize}
\end{enonce*}

This is a special case of \hyperref[VI.cor1.24]{1.24}, taking into account
that $C$ is stable under finite sums (\hyperref[VI.cor1.4]{1.4}), which
implies the equivalence of (ii) and (ii bis), and taking into account the

\begin{enonce*}{Lemma 1.25.1}\etiquette[1.25.1]{VI.lem1.25.1}
Let $E$ be a topos. Every direct factor $X'$ of a coherent object $X$ of $E$
is coherent.
\end{enonce*}

Indeed, $X'$ is quasi-compact as a quotient of $X$
(\hyperref[VI.prop1.3]{1.3}), and quasi-separated as a subobject of $X$
(\hyperref[VI.cor1.15]{1.15}).

\begin{enonce*}{Corollary 1.26}\etiquette[1.26]{VI.cor1.26}
Let $E$ be a topos satisfying the equivalent conditions
\emph{\hyperref[VI.cor1.25]{1.25}.} Then every induced topos does likewise.
\end{enonce*}

This follows immediately from criterion (ii), taking into account
\hyperref[VI.subsubsec1.20.2]{1.20.2} and the fact that the existence of a
generating subcategory of $E$ consisting of coherent objects plainly
implies the same property for an induced topos.

\begin{enonce*}[remark]{Remark 1.27}\etiquette[1.27]{VI.rem1.27}
Among\pageoriginale the important examples of topoi satisfying the
conditions of \hyperref[VI.cor1.25]{1.25}, let us mention Noetherian topoi
(\hyperref[VI.subsec2.14]{2.14}), and the Zariski or étale topos
(\hyperrefext[VII][VII.sec1]{1}) of a scheme $X$ (\cf IX, note p.42). But
even when the topos $E$ is \emph{coherent} (\hyperref[VI.def2.3]{2.3}) (\ie
the full subcategory $C$ of coherent objects is generating and stable under
finite projective limits), it is not always true that $C$ is \emph{perfect}
(\hyperref[VI.def2.9.1]{2.9.1}), \ie satisfies the equivalent conditions of
\hyperref[VI.cor1.25]{1.25}; \cf \hyperref[VI.exer1.28]{1.28} below.
\end{enonce*}

\begin{enonce*}[remark]{Exercise 1.28}\etiquette[1.28]{VI.exer1.28}
Let $E$ be a topos and $p_1,p_2:X_1\rightrightarrows X_0$ a double arrow in
$E$. Define an increasing sequence of subobjects $R_n$ $(n\geq0)$ of
$X_0\times X_0$ by the condition
$R_0=\Sup(\text{Im}(p_1,p_2),\text{Im}(p_2,p_1))$, and
$R_n=\pr_{13}(\pr_{12}^{-1}(R_{n-1}),\pr_{23}^{-1}(R_{n-1}))$ for every
$n\geq1$, where $\pr_{ij}$ $(1\leq i<j\leq3)$ denote the evident
projections from the triple product of $X_0$ to the double product.
\begin{enumerate}
\renewcommand{\theenumi}{\alph{enumi}}
\renewcommand{\labelenumi}{\theenumi)}
\item Show that one thus obtains an increasing sequence of subobjects of
  $X_0\times X_0$, that $R=\varinjlim R_n$ is an equivalence graph (I \hyperref[I.sec10.4]{10.4})
  in $X_0$, and that the canonical morphism
$$
X = \Coker (p_1, p_2) \to X_0 /R
$$
is an isomorphism.

\item Deduce that if $X_0$ is quasi-compact and $X=\Coker(p_1,p_2)$ is
  quasi-separated (hence coherent), then the sequence $(R_n)$ is stationary.
  Conversely,\pageoriginale if this sequence is stationary, and if $X_0$ is
  assumed coherent and $X_1$ quasi-compact, then $X$ is coherent. (For this
  last point, write
  $R_n\simeq\fprod{(R_{n-1},\pr_2)}{(R_{n-1},\pr_1)}{X_0}$ and deduce that
  $R_n$ is quasi-compact over $X_0$ for every $n$, then use \hyperref[VI.cor1.19]{1.19}.)

\item Construct an example of two maps of sets
  $p_1,p_2:X_1\rightrightarrows X_0$ such that the sequence $(R_n)$ is not
  stationary.

\item Let $C$ be a small category in which finite inductive limits and
  finite projective limits are representable, and in which every morphism
  factors as an effective epimorphism followed by a monomorphism. Repeat the
  construction of the $R_n$ for a double arrow $(p_1,p_2)$ of $C$. Endowing
  $C$ with the canonical topology, show that the canonical functor
  $\epsilon:C\to C^{\sim}$ is exact, and consequently commutes with the
  formation of the $R_n$.

\item Take for $C$ a small full subcategory of $(\Ens)$, stable up to
  isomorphism under finite projective limits and finite inductive limits, and
  containing the sets $X_0,X_1$ of c). Deduce that the topos $C^{\sim}$
  (which is a \emph{coherent} topos (\hyperref[VI.def2.3]{2.3}), \ie is
  generated by the subcategory of its coherent objects, which is stable
  under finite projective limits) does not satisfy the equivalent conditions
  of \hyperref[VI.cor1.25]{1.25}.

\item Let $k$ be a field, and let $C$ be the fppf site of schemes locally of
  finite presentation over $k$ (\SGA 3 IV 6.3). Show that there exists a
  reduced affine scheme $X$ over $k$ admitting an automorphism $f$ that is
  not of finite order (take, for example, the affine plane $E^2$ endowed with
  the automorphism $f(x)=x$, $f(y)=x+y$).\pageoriginale Show that if
  $p_1,p_2:X_1\rightrightarrows X_0$ is a double arrow in $C$ such that the
  corresponding double arrow in $(\Ens) p_1(\bar{k}):p_2(\bar{k}):
  X_1(\bar{k})\rightrightarrows X_0(\bar{k})$ gives a nonstationary sequence
  $(R_n)$, then the same is true of the double arrow of $C^{\sim}$ defined by
  $p_1,p_2$. Thus, if $X_0$ is of finite type over $k$, the cokernel in
  $C^{\sim}$ of $(p_1,p_2)$ is not coherent. Deduce that the cokernel in
  $C^{\sim}$ of the pair $(f,\id_X)$ is not coherent, and hence that the
  topos $C^{\sim}$ does not satisfy the equivalent conditions of
  \hyperref[VI.cor1.25]{1.25}. Deduce more generally that if $S$ is a
  nonempty scheme and $C$ is the fppf site of schemes locally of finite
  presentation over $S$, then the topos $C^{\sim}$ does not satisfy the
  conditions of \hyperref[VI.cor1.25]{1.25}. (Care must be taken that,
  contrary to appearances, $C^{\sim}$ is therefore Noetherian
  (\hyperref[VI.def2.11]{2.11}) only if $S$ is empty, as follows from the
  preceding discussion and \hyperref[VI.subsec2.14]{2.14}.)
\end{enumerate}
\end{enonce*}

\begin{enonce*}{Definition 1.30}\etiquette[1.30]{VI.def1.30}
Let $E$ be a topos. An object $X$ of $E$ is called a \emph{pre-Noetherian
object}\footnote{The stronger notion of a \emph{Noetherian} object is given
  in \hyperref[VI.def2.11]{2.11} below.} of the topos $E$ if it satisfies the
following two equivalent conditions:
\begin{enumerate}
\renewcommand{\theenumi}{\roman{enumi}}
\renewcommand{\labelenumi}{(\theenumi)}

\item Every subobject of $X$ is quasi-compact.

\item Every increasing sequence of subobjects of $X$ is stationary.
\end{enumerate}
\end{enonce*}

\setcounter{subsection}{30}
\setcounter{subsubsection}{0}
\subsubsection{}\etiquette[1.30.1]{VI.subsubsec1.30.1}
The equivalence of conditions (i) and (ii) is clear. Note that these
conditions again depend only on the induced topos $E_{/X}$, and indeed only
on the ordered set of opens of $E_{/X}$, which is isomorphic to the ordered
set of subobjects of $X$. They are stable under passage to a subobject of
$X$.



% END INCLUDED SGA4_English_current_buildable_source_20260906/source/06/drafts/SGA4_VI_LINES_678_1095_English.texfrag


% BEGIN INCLUDED SGA4_English_current_buildable_source_20260906/source/06/drafts/SGA4_VI_LINES_1096_1382_English.texfrag
\subsection{}\etiquette[1.31]{VI.subsec1.31}
One proves,\pageoriginale as in \hyperref[VI.prop1.3]{1.3} and
\hyperref[VI.cor1.4]{1.4}, the following facts: if $(X_i\to X)$ is a
finite covering family, with the $X_i$ pre-Noetherian, then $X$ is
pre-Noetherian; in particular, a quotient of a pre-Noetherian object
of $E$ is pre-Noetherian. If $(X_i)_{i\in I}$ is a family of objects
of $E$, then their coproduct $X$ is a pre-Noetherian object of $E$ if
and only if all but finitely many of the $X_i$ are isomorphic to
$∅_E$, and all the $X_i$ are pre-Noetherian.

\subsection{}\etiquette[1.32]{VI.subsec1.32}
Clearly, a pre-Noetherian object of a topos $E$ is quasi-compact. When
$E$ admits a generating family formed by pre-Noetherian objects, the
converse holds (as follows immediately from
\hyperref[VI.subsec1.31]{1.31}): the pre-Noetherian objects of $E$ are
then precisely its quasi-compact objects.

Suppose that $E$ admits a generating family $(X_i)_{i\in I}$ formed
by quasi-compact objects. Then the following three conditions are
equivalent:
\begin{enumerate}
\renewcommand{\theenumi}{\roman{enumi}}
\renewcommand{\labelenumi}{(\theenumi)}
\item The $X_i$ are pre-Noetherian.

\item Every quasi-compact object of $E$ is pre-Noetherian.

\item Every monomorphism in $E$ is quasi-compact.
\end{enumerate}
\noindent
(We have just proved the equivalence of (i) and (ii); (ii)
$\Rightarrow$ (iii) is immediate from the definitions, and (iii)
$\Rightarrow$ (ii) follows from Definition
\hyperref[VI.def1.30]{1.30} (i).)

Note\pageoriginale that (iii) implies that \emph{every morphism of
$E$ is quasi-separated}; hence (\hyperref[VI.prop1.14]{1.14} (ii))
every object of $E$ lying over a quasi-separated object of $E$ is
itself quasi-separated.

\begin{enonce*}[remark]{Example 1.33}\etiquette[1.33]{VI.exem1.33}
(\emph{Classifying topos of a group}). Let $G$ be a group
$\in 𝒰$, and let $B_G$ be its classifying topos (IV \hyperref[IV.subsec2.4]{2.4}), \ie the
category of sets $\in 𝒰$ on which $G$ acts on the left. The nonempty
connected objects of $B_G$ (IV \hyperref[IV.subsubsec4.3.5]{4.3.5}) are the nonempty $G$-sets on
which $G$ acts transitively, and every object $E$ of $B_G$ is the
coproduct of its connected components, which are the minimal nonempty
subobjects of $E$ (namely, the orbits of $G$ in $E$). Consequently,
$E$ is a quasi-compact object of $B_G$ if and only if the set
$\pi_0(E)$ of orbits of $G$ in $E$ is finite, and in that case $E$ is
even a pre-Noetherian object of $B_G$. These objects (and even the
single object $G_s$, which is nonempty and connected) form a
generating subcategory of $B_G$. The final object $e$ of $B_G$ is
quasi-separated, \ie the product of two quasi-compact objects of
$B_G$ is quasi-compact, if and only if the product $G_s\times G_s$ is
quasi-compact, which is plainly equivalent to $G$ being finite. More
generally, one concludes that an object $E$ of $B_G$ is
quasi-separated if and only if the stabilizers of its points are
\emph{finite} subgroups of $G$, \ie if it is isomorphic to a coproduct
of homogeneous spaces $G/H$, with $H$ \emph{finite}. Indeed, by
\hyperref[VI.cor1.15]{1.15}, one may assume in this criterion that
$E$ is connected; but if $x\in E$ has stabilizer group $G_x$, we know
(IV \hyperref[IV.subsec5.8]{5.8})\pageoriginale that $B_G/E$ is equivalent to the topos
$B_{G_x}$, and hence its final object is quasi-separated if and only
if $G_x$ is finite. This criterion shows that every object of $B_G$
lying over a quasi-separated object is quasi-separated; a fortiori,
the product of two quasi-separated objects of $B_G$ is
quasi-separated. It also shows that the coherent objects of $B_G$ are
the objects isomorphic to finite coproducts of homogeneous spaces
$G/H$, with $H$ \emph{finite}. Since $G_s$ is coherent, the coherent
objects of $B_G$ already form a generating subcategory of $B_G$.
Moreover, this subcategory is stable under fiber products (as follows
from the fact that every object lying over a quasi-separated object
is quasi-separated). (In the terminology of
\hyperref[VI.def2.11]{2.11} below, $B_G$ is a \emph{locally
Noetherian} topos, but it is quasi-separated only when $G$ is finite,
in which case $B_G$ is even Noetherian.)

It is immediate that a morphism $f:E'\to E$ is quasi-compact if and
only if the induced map $\pi_0(f):\pi_0(E')\to\pi_0(E)$ on the sets
of orbits is proper, \ie has finite fibers. We recover the fact that
a monomorphism is quasi-compact (\hyperref[VI.subsec1.32]{1.32}), and
hence that every morphism is quasi-separated. Thus the coherent
morphisms in $B_G$ are simply the quasi-compact morphisms, which have
just been characterized.

Let $E$ be an object of the topos $B_G$. One verifies easily that $E$
is a filtered inductive limit of coherent objects of $B_G$ if and
only if the stabilizers of the points of $E$ are ind-finite groups
(\ie filtered\pageoriginale inductive limits of their finite
subgroups). In particular, if $G$ is not ind-finite (for example, if
$G=\ZZ$), the final object of the topos $B_G$ is not an inductive
limit of coherent objects. More precisely, the topos $B_G$ satisfies
the equivalent conditions of \hyperref[VI.cor1.25]{1.25} if and only
if its final object is a filtered inductive limit of coherent objects,
or equivalently if the group $G$ is ind-finite.
\end{enonce*}


\section{Finiteness conditions for a topos}\etiquette[2]{VI.sec2}

\begin{enonce*}{Proposition 2.1}\etiquette[2.1]{VI.prop2.1}
Let $E$ be a $𝒰$-topos. The following conditions are equivalent:
\begin{enumerate}
\renewcommand{\theenumi}{\roman{enumi}}
\renewcommand{\labelenumi}{(\theenumi)}
\item There exists a full generating subcategory $C$ of $E$, formed
  by quasi-compact objects and stable under fiber products.

\item There exists a $𝒰$-site $C$ such that every object of $C$ is
  quasi-compact, fiber products are representable in $C$, and $E$ is
  equivalent to the category of sheaves $C^{\sim}$.
\end{enumerate}

Suppose that these conditions hold. Then in \emph{(i) (\resp (ii))}
the category $C$ may even be chosen to be $𝒰$-small. Moreover, for
every $X\in\ob C$, the object $X$ \emph{(\resp $\epsilon(X)$)} of the
topos $E$ is \emph{coherent}.
\end{enonce*}

The implication (i) $\Rightarrow$ (ii) follows from IV \hyperref[IV.cor1.2.1]{1.2.1}, and the
fact that in (i) (\resp (ii)) one may take $C$ to be $𝒰$-small
follows immediately from the argument of IV \hyperref[IV.subsubsec1.2.3]{1.2.3}, applied to a small
full generating subcategory in the topological sense (II \hyperref[II.def3.0.1]{3.0.1}) of
$C$. It remains to prove (ii) $\Rightarrow$ (i) and the last assertion
of \hyperref[VI.prop2.1]{2.1}; these follow immediately from the
following result.

\begin{enonce*}{Corollary 2.1.1}\etiquette[2.1.1]{VI.cor2.1.1}
Let\pageoriginale $C$ be a $𝒰$-site in which fiber products are
representable and every object is quasi-compact, let $E=C^{\sim}$,
and let $\varepsilon:C\to E$ be the canonical functor. Then, for every
$X\in\ob C$, $\varepsilon(X)$ is a coherent object of $E$. The full
subcategory of $E$ formed by the coherent objects of $E$ lying over
some $\varepsilon(X)$ is stable under fiber products (and is plainly
generating in $E$).
\end{enonce*}

We already know that $\epsilon(X)$ is quasi-compact
(\hyperref[VI.prop1.2]{1.2}); let us prove that it is quasi-separated.
Let $F,G$ be two quasi-compact objects of $E$, and let
$F\to\varepsilon(X)$, $G\to\epsilon(X)$ be two morphisms. We must
prove that the fiber product $\fprod{F}{G}{\epsilon(X)}$ is
quasi-compact. Covering $F$ and $G$ by finitely many objects of the
form $\epsilon(Y_i)$ and, respectively, $\epsilon(Z_j)$, and
proceeding as in \hyperref[VI.prop1.2]{1.2} to reduce to the case in
which the composites $\epsilon(Y_i)\to\epsilon(X)$ and
$\epsilon(Z_j)\to\epsilon(X)$ come from morphisms $Y_i\to X$ and,
respectively, $Z_j\to X$, we are reduced by
\hyperref[VI.prop1.3]{1.3} to the case in which the morphisms
$F\to\epsilon(X)$ and $G\to\epsilon(X)$ are obtained by applying
$\epsilon$ to morphisms $Y\to X$ and $Z\to X$. But since fiber
products are representable in $C$ and $\epsilon$ commutes with them,
we are reduced to proving that $\epsilon(\fprod{Y}{Z}{X})$ is
quasi-compact, which was already noted above for every object of the
form $\epsilon(U)$.

Now let $G\to F$, $H\to F$ be morphisms between coherent objects of
$E$, and suppose that there exists a morphism $F\to\epsilon(X)$ for
some $X\in\ob C$. We prove that $\fprod{G}{H}{F}$ is coherent. Since
$G$ and $H$ are quasi-compact and $F$ is quasi-separated, it already
follows from the definitions that the fiber product is quasi-compact.
It remains\pageoriginale to see that it is quasi-separated, or,
equivalently (\hyperref[VI.cor1.15]{1.15}), that
$\fprod{G}{H}{\epsilon(X)}$ is quasi-separated. This follows from the
more general fact:

\begin{enonce*}{Corollary 2.1.2}\etiquette[2.1.2]{VI.cor2.1.2}
Under the conditions of \hyperref[VI.cor2.1.1]{2.1.1}, for every
quasi-separated object $G$ of $E$, every morphism
$G\to\epsilon(X)$ is quasi-separated. If $H\to\epsilon(X)$ is a
second morphism, with $H$ quasi-separated, then
$\fprod{G}{H}{\epsilon(X)}$ is quasi-separated.
\end{enonce*}

The first assertion implies the second: by base change, it implies
that $\fprod{G}{H}{\epsilon(X)}\to H$ is quasi-separated, and since
$H$ is quasi-separated, $\fprod{G}{H}{\epsilon(X)}$ is likewise
quasi-separated (\hyperref[VI.prop1.14]{1.14} (ii)). To prove that
$G\to\epsilon(X)$ is quasi-separated, we shall use the following
lemma.

\begin{enonce*}{Lemma 2.1.3}\etiquette[2.1.3]{VI.lem2.1.3}
Let $E$ be a topos and $f:X\to Y$ a morphism in $E$. Suppose that $E$
admits a generating family of quasi-compact objects, and that there
exists a covering family $(g_i:X_i\to X)_{i\in I}$ of $X$ such that
the $X_i$ are quasi-compact and, for every pair of indices
$i,j\in I$, $\fprod{X_i}{X_j}{Y}$ is quasi-separated. Then, if $X$
is a quasi-separated object of $E$, $f$ is a quasi-separated
morphism.
\end{enonce*}

Indeed, the family
$(\fprod{g_i}{g_j}{Y}:\fprod{X_i}{X_j}{Y}\to
\fprod{X}{X}{Y})_{(i,j)\in I\times I}$ is covering. Therefore, to see
that the morphism $\Delta_{X/Y}:X\to\fprod{X}{X}{Y}$ is
quasi-compact, it is enough to see that, for every
$(i,j)\in I\times I$, the morphism obtained from it\pageoriginale by
base change along $\fprod{g_i}{g_j}{Y}$ is quasi-compact
(\hyperref[VI.prop1.10]{1.10} (ii)). The latter is precisely the
canonical morphism
$\fprod{X_i}{X_j}{X}\to\fprod{X_i}{X_j}{Y}$. By hypothesis $X$ is
quasi-separated and the $X_i$ are quasi-compact, so
$\fprod{X_i}{X_j}{X}$ is quasi-compact; and since
$\fprod{X_i}{X_j}{Y}$ is quasi-separated by hypothesis, the morphism
$\fprod{X_i}{X_j}{X}\to\fprod{X_i}{X_j}{Y}$ is indeed quasi-compact,
\cqfd

We can now complete the proof of \hyperref[VI.cor2.1.2]{2.1.2} by
proving that every morphism $f:G\to\epsilon(X)$, with $G$
quasi-separated, is quasi-separated. To do this, cover $G$ by objects
$\epsilon(X_i)$ such that the composite morphisms
$\epsilon(X_i)\to\epsilon(X)$ come from morphisms $X_i\to X$. By
\hyperref[VI.lem2.1.3]{2.1.3}, since the $\epsilon(X_i)$ are
quasi-compact, it is enough to verify that the fiber products
$\fprod{\epsilon(X_i)}{\epsilon(X_j)}{\epsilon(X)}=
\epsilon(\fprod{X_i}{X_j}{X})$ are quasi-separated. We have already
seen that every object of $E$ of the form $\epsilon(U)$ is coherent.
This proves \hyperref[VI.cor2.1.2]{2.1.2} and completes the proof of
\hyperref[VI.prop2.1]{2.1}.

\begin{enonce*}{Proposition 2.2}\etiquette[2.2]{VI.prop2.2}
Let $E$ be a topos containing a full generating subcategory $C$
formed by coherent objects. The following conditions on $E$ are
equivalent:
\begin{itemize}
\item[(i)] Every quasi-separated object of $E$ is quasi-separated over
  the final object.

\item[(i bis)] For every morphism $f:X\to Y$ in $E$, if $X$ is
  quasi-separated, then $f$ is quasi-separated.

\item[(i ter)] Every object $X$ of $C$ is quasi-separated over the
  final object of $E$, \ie the diagonal morphism
  $X\to X\times X$ is quasi-compact.

\item[(ii)] The product\pageoriginale of two quasi-separated objects
  of $E$ is quasi-separated.

\item[(ii bis)] The product in $E$ of two objects of $C$ is
  quasi-separated.
\end{itemize}

These conditions imply that the full subcategory $E_{\coh}$ of $E$
formed by the coherent objects of $E$ is stable under fiber products
(and, a fortiori, that $E$ satisfies the equivalent conditions of
\hyperref[VI.prop2.1]{2.1}).
\end{enonce*}

Clearly (i bis) $\Rightarrow$ (i), and the converse follows from
\hyperref[VI.prop1.8]{1.8} (iii). We have (i) $\Rightarrow$ (ii) by
an argument already used: if $X$ and $Y$ are quasi-separated, then by
hypothesis (i), $X$ is quasi-separated over $e$, so $X\times Y$ is
quasi-separated over $Y$ by base change, and hence $X\times Y$ is
quasi-separated by \hyperref[VI.prop1.14]{1.14} (ii). The same
argument shows that (i ter) $\Rightarrow$ (ii bis), while (i)
$\Rightarrow$ (i ter) and (ii) $\Rightarrow$ (ii bis) are immediate.
It therefore remains to prove (ii bis) $\Rightarrow$ (i). If $X$ is a
quasi-separated object of $E$, cover $X$ by objects $X_i$ of $C$
(which is possible since $C$ is generating). Since the $X_i$ are
quasi-compact by the hypothesis on $C$, we may apply
\hyperref[VI.lem2.1.3]{2.1.3} to the morphism $X\to e$; thus, to
prove that this morphism is quasi-separated, we are reduced precisely
to verifying that the $X_i\times X_j$ are quasi-separated, which is
hypothesis (ii bis).

\begin{enonce*}[remark]{Remark 2.2.1}\etiquette[2.2.1]{VI.rem2.2.1}
Condition (i ter) is also equivalent to the following condition (i
quater): for every double arrow $g_1,g_2:Y\rightrightarrows X$ in
$C$, its kernel in $E$ is quasi-compact (or equivalently
(\hyperref[VI.cor1.15]{1.15}) coherent).
\end{enonce*}

To see\pageoriginale this, it is enough to apply criterion
\hyperref[VI.prop1.10]{1.10} (i) to the diagonal morphism
$X\to X\times X$ considered in (i ter).

\begin{enonce*}{Definition 2.3}\etiquette[2.3]{VI.def2.3}
A topos $E$ satisfying the equivalent conditions of
\emph{\hyperref[VI.prop2.2]{2.2} (\resp
\hyperref[VI.prop2.1]{2.1})} is called an \emph{algebraic topos
(\resp} a \emph{locally algebraic topos}, or \emph{locally
coherent)}. An object $X$ of a topos $E$ is called an
\emph{algebraic object of the topos} $E$ if the induced topos $E_{/X}$
is algebraic. A topos $E$ is called \emph{quasi-separated
(\resp coherent)} if it is algebraic and its final object is
quasi-separated (\resp coherent).
\end{enonce*}

\setcounter{subsection}{3}
\subsubsection{}\etiquette[2.3.1]{VI.subsubsec2.3.1}
It is in algebraic topoi that the finiteness notions developed in
${\rm n^{\circ}}$~1 have fully satisfactory properties, analogous to
those of the notions bearing the same names in the category of
schemes. All locally algebraic (= locally coherent) topoi encountered
in practice are in fact algebraic, so the practical interest of this
notion seems rather limited for the present. Nevertheless, one can
construct locally algebraic topoi that are not algebraic
(\hyperref[VI.exer2.17]{2.17} f).


% END INCLUDED SGA4_English_current_buildable_source_20260906/source/06/drafts/SGA4_VI_LINES_1096_1382_English.texfrag


% BEGIN INCLUDED SGA4_English_current_buildable_source_20260906/source/06/drafts/SGA4_VI_LINES_1383_1678_English.texfrag
\subsection{}\etiquette[2.4]{VI.subsec2.4}
Here are some general remarks concerning the notions introduced in
\hyperref[VI.def2.3]{2.3}, which will convince the reader (or so we hope)
that the terminology introduced there is coherent.

\subsubsection{}\etiquette[2.4.1]{VI.subsubsec2.4.1}
Under the conditions of \hyperref[VI.prop2.1]{2.1} (i) \resp (ii), for every
$X\in C$, $X$ (\resp $\epsilon(X)$) is an algebraic (coherent) object of $E$
(\hyperref[VI.cor2.1.2]{2.1.2} and 1.19.2); in other\pageoriginale words, the
induced topos $E_{/X}$ (\resp
$E_{\epsilon(X)}\approxeq(C_{/X})^{\sim}$) is algebraic, hence
\emph{coherent}. Equivalently, \emph{when $C$ admits a final object, $E$
itself is algebraic, hence coherent.}

\subsubsection{}\etiquette[2.4.2]{VI.subsubsec2.4.2}
If $E$ is an algebraic topos, then every induced topos $E_{/X}$ is also
algebraic (\hyperref[VI.prop2.2]{2.2} (i bis)). For a topos $E$ to be locally
algebraic $=$ locally coherent, it is necessary and sufficient that there
exist a family $(X_i)_{i\in I}$ of objects of $E$ covering the final object,
such that the $X_i$ are algebraic (\resp coherent), \ie such that the
induced topoi $E_{/X_i}$ are algebraic (\resp coherent). Sufficiency follows
immediately from the definitions, and necessity from the observations in
\hyperref[VI.subsubsec2.4.1]{2.4.1}. Of course, an algebraic topos is locally
algebraic, and if $E$ is a locally algebraic topos, every induced topos
$E_{/X}$ is locally algebraic.

\subsubsection{}\etiquette[2.4.3]{VI.subsubsec2.4.3}
If $X$ is an algebraic object of a topos $E$, then every object $X'$ of $E$
lying over $X$ is again algebraic (\cf
\hyperref[VI.subsubsec2.4.2]{2.4.2}). In particular, every object of an
algebraic topos is algebraic, and conversely, of course.

\subsubsection{}\etiquette[2.4.4]{VI.subsubsec2.4.4}
The full subcategory of $E$ consisting of the \emph{quasi-separated} objects
that are \emph{algebraic} is stable under fiber products and subobjects
(\hyperref[VI.cor1.15]{1.15}), and hence also under kernels. If $E$ is
algebraic, this category (which is then simply\pageoriginale the category of
quasi-separated objects of $E$) is also stable under products of two
objects. If $E$ is quasi-separated (and only then), this category is stable
under arbitrary finite projective limits, or equivalently contains the
final object of $E$.

The full subcategory $C$ of $E$ consisting of the \emph{coherent} objects
that are \emph{algebraic} is stable under fiber products. If $E$ is locally
algebraic, \ie locally coherent, saying that $E$ is algebraic amounts to
saying that $C$ is also stable under kernels of double arrows
(\hyperref[VI.rem2.2.1]{2.2.1}); and saying that $E$ is quasi-separated
amounts to saying that, in addition, $C$ is stable under products of two
objects (\cf \hyperref[VI.cor1.17]{1.17}). Finally, saying that $E$ is
coherent amounts to saying that $C$ is stable under arbitrary finite
projective limits, or equivalently that it contains the final object of
$E$.

\subsubsection{}\etiquette[2.4.5]{VI.subsubsec2.4.5}
Let $E$ be a topos. For $E$ to be locally coherent (\resp algebraic, \resp
quasi-separated, \resp coherent), it is necessary and sufficient that it
admit a full generating subcategory $C$ consisting of quasi-compact objects
(which will automatically be coherent (\hyperref[VI.prop2.1]{2.1})) and
\emph{stable} under fiber products (\resp under fiber products and kernels,
\resp under fiber products and products of two objects, \resp under
arbitrary finite projective limits); or equivalently that $E$ be equivalent
to a topos $C^{\sim}$, where $C$ is a $𝒰$-site every object of which is
quasi-compact, and in which fiber products and products of two objects
(\resp all finite projective limits) are representable. (For necessity,
take\pageoriginale for $C$ the category of objects of $E$ that are coherent
and \emph{algebraic}, and apply \hyperref[VI.subsec2.4]{2.4}; for
sufficiency, apply \hyperref[VI.prop2.2]{2.2} (i ter).) In this statement,
one may of course also suppose that $C$ is $𝒰$-small.

\subsubsection{}\etiquette[2.4.6]{VI.subsubsec2.4.6}
Let $E$ be a topos. If $E$ is algebraic, an object $X$ of $E$ is
quasi-separated (\resp coherent) if and only if the induced topos $E_{/X}$
is so. When $E$ is no longer assumed algebraic, one can only say that
$E_{/X}$ is quasi-separated (\resp coherent) if and only if $X$ is
quasi-separated (\resp coherent) and \emph{algebraic}. This is of little
consequence, since we will probably not have to use these notions outside
the case in which $E$ is algebraic.

\subsubsection{}\etiquette[2.4.7]{VI.subsubsec2.4.7}
For the notion of a coherent (\resp quasi-separated) object of a topos $E$
to have satisfactory properties, one must restrict to objects that are also
algebraic, a condition automatically satisfied if $E$ is algebraic. Since
we will probably never have to work with locally algebraic topoi $E$ that
are not algebraic, the question of revising the terminology introduced in
\hyperref[VI.def1.13]{1.13}, possibly reserving it for algebraic objects
alone, therefore does not arise with any great urgency!

Moreover, we have refrained from defining the notion of a quasi-compact
topos. In the spirit of the present paragraph, this name would naturally be
given to \emph{algebraic} topoi whose final object is quasi-compact.

We\pageoriginale hesitated to introduce such a usage, since if $X$ is a
topological space, it would not be true that $X$ is quasi-compact if and
only if the associated topos $\Top(X)$ is so. In this connection, note that
$\Top(X)$ is locally coherent if and only if $X$ admits a basis of opens
consisting of quasi-compact opens $U$ such that, for
$U_j,U_k\subset U_i$, $U_j\cap U_k$ is again quasi-compact (and then
$\Top(X)$ is even algebraic). This condition is satisfied for the
topological spaces underlying schemes, but rarely for spaces encountered by
analysts, even when the spaces themselves are compact. For ease of
reference, let us state explicitly:

\begin{enonce*}{Proposition 2.5}\etiquette[2.5]{VI.prop2.5}
Let $E$ be an algebraic topos and $f:X\to Y$ a morphism in $E$. If $Y$ is a
quasi-separated (\resp coherent) object of $E$, then, for $X$ to be a
quasi-separated (\resp coherent) object of $E$, it is necessary and
sufficient that the morphism $f$ be quasi-separated (\resp coherent).
\end{enonce*}

Sufficiency was already seen in \hyperref[VI.prop1.14]{1.14} (ii), and holds
without any hypothesis on $E$. Necessity in the non-parenthetical case holds
even without supposing that $Y$ is quasi-separated, and is simply Definition
\hyperref[VI.def2.3]{2.3}, via \hyperref[VI.prop2.2]{2.2} i bis). It remains
to prove that if $Y$ is coherent and $X$ is coherent, then $f$ is coherent.
Since we already know that $f$ is quasi-separated, it is enough to see that
$f$ is quasi-compact, which follows from the fact that $X$ is quasi-compact
and $Y$ quasi-separated (\hyperref[VI.def1.13]{1.13}).

\begin{enonce*}{Corollary 2.6}\etiquette[2.6]{VI.cor2.6}
Let\pageoriginale $E$ be an algebraic topos, $f:X\to Y$ a morphism in $E$,
and $(Y_i\to Y)_{i\in I}$ a covering family, with the $Y_i$ quasi-separated
(\resp coherent). For $f$ to be quasi-separated (\resp quasi-compact,
\resp coherent), it is necessary and sufficient that, for every $i\in I$,
$X_i=\fprod{X}{Y_i}{Y}$ be a quasi-separated (\resp quasi-compact, \resp
coherent) object of $E$.
\end{enonce*}

It is enough to combine \hyperref[VI.prop1.10]{1.10} (ii) and
\hyperref[VI.prop2.5]{2.5} in the ``quasi-separated'' or ``coherent'' case;
the ``quasi-compact'' case was already seen in
\hyperref[VI.cor1.16]{1.16}. Note that the necessity of the condition stated
in \hyperref[VI.cor2.6]{2.6} also holds without any hypothesis on the topos
$E$.

As a special case of \hyperref[VI.cor2.6]{2.6}, one obtains:

\begin{enonce*}{Corollary 2.7}\etiquette[2.7]{VI.cor2.7}
Let $E$ be a coherent topos and $X$ an object of $E$. For $X$ to be a
quasi-compact (\resp quasi-separated, \resp coherent) object of $E$, it is
necessary and sufficient that it be quasi-compact (\resp quasi-separated,
\resp coherent) over the final object
(\hyperref[VI.subsubsec1.9.2]{1.9.2}).
\end{enonce*}

In particular, $X$ is constructible (\hyperref[VI.subsubsec1.9.2]{1.9.2})
if and only if it is coherent.

\begin{enonce*}{Corollary 2.8}\etiquette[2.8]{VI.cor2.8}
Let $E$ be a locally coherent topos and $f:X\to Y$ a morphism in $E$. For
$f$ to be quasi-separated, it is necessary and sufficient that, for every
quasi-separated object $Y'$ of $E_{/Y}$,
$f^{\ast}(Y')=\fprod{X}{Y'}{Y}$ be a quasi-separated object of $E_{/X}$ (or
of $E$, which amounts to the same thing
(\hyperref[VI.subsubsec1.20.2]{1.20.2})).
\end{enonce*}

The\pageoriginale ``only if'' part holds without any condition on $E$,
since if $Y'$ is quasi-separated, then
$f':X'=\fprod{X}{Y'}{Y}\to Y'$ is quasi-separated by base change, and hence
$X'$ is quasi-separated (\hyperref[VI.prop1.14]{1.14} (ii)). For the
converse, note that there exists a covering family $Y_i\to Y$ of $Y$ by
quasi-separated algebraic objects. By \hyperref[VI.prop1.10]{1.10} (ii), to
verify that $f$ is quasi-separated it suffices to verify that the
$f_i:X_i=\fprod{X}{X_i}{Y}\to Y_i$ are so. But the hypothesis on $f^{\ast}$
implies that the $X_i$ are quasi-separated, and therefore so are the $f_i$,
since $E_{/Y_i}$ is an algebraic topos.

\begin{enonce*}[remark]{Remark 2.8.1}\etiquette[2.8.1]{VI.rem2.8.1}
Statements \hyperref[VI.prop2.5]{2.5} and \hyperref[VI.cor2.6]{2.6} remain
valid if the hypothesis ``$E$ algebraic'' is replaced by the more general
hypothesis ``$Y$ is algebraic'', as is seen immediately by applying the
original statement to the induced topos $E_{/Y}$, which is algebraic.
Likewise, in \hyperref[VI.cor2.8]{2.8} it is enough to suppose that $E_{/Y}$
(rather than $E$) is locally coherent.
\end{enonce*}

\begin{enonce*}{Proposition 2.9}\etiquette[2.9]{VI.prop2.9}
Let $E$ be a topos and $C$ a subcategory of $E$. The following conditions
are equivalent:
\begin{enumerate}
\renewcommand{\theenumi}{\roman{enumi}}
\renewcommand{\labelenumi}{(\theenumi)}
\item $E$ is locally coherent (\resp coherent), satisfies the equivalent
  conditions of \hyperref[VI.cor1.25]{1.25}, and $C$ is the full subcategory
  $E_{\coh}$ of $E$ consisting of the coherent objects.

\item $C$ is a strictly full generating subcategory of $E$, stable under
  finite inductive limits and fiber products (\resp and under finite
  projective limits), and consists of quasi-compact objects.
\end{enumerate}
\end{enonce*}

The implication\pageoriginale (i) $\Rightarrow$ (ii) follows immediately
from the definitions and from formulation \hyperref[VI.cor1.25]{1.25} (ii)
of the conditions under consideration in \hyperref[VI.cor1.25]{1.25}. Let
us prove the converse. The fact that $C$ is a full generating subcategory
consisting of quasi-compact objects and stable under fiber products (\resp
under finite projective limits) implies, by definition, that $E$ is locally
coherent (\resp coherent). The additional fact that $C$ is strictly full
and stable under finite inductive limits then implies that $C=E_{\coh}$ (by
\hyperref[VI.cor1.24]{1.24} (ii) $\Rightarrow$ (iii)), and that condition
\hyperref[VI.cor1.25]{1.25} (ii) is satisfied, \cqfd

\begin{enonce*}{Definition 2.9.1}\etiquette[2.9.1]{VI.def2.9.1}
A $𝒰$-topos $E$ is called perfect if it is coherent
(\hyperref[VI.def2.3]{2.3}) and satisfies the equivalent conditions of
\hyperref[VI.cor1.25]{1.25}, \ie if the subcategory $E_{\coh}$ of $E$
consisting of the coherent objects of $E$ is stable under finite inductive
limits. We say that $E$ is locally perfect if there exists a family
$(X_i)_{i\in I}$ of objects of $E$ covering the final object such that, for
every $i\in I$, the induced topos $E_{/X_i}$ is perfect.
\end{enonce*}

\setcounter{subsection}{9}
\setcounter{subsubsection}{1}
\subsubsection{}\etiquette[2.9.2]{VI.subsubsec2.9.2}

It follows immediately from this definition that a perfect (\resp locally
perfect) topos is coherent (\resp locally coherent). As examples of perfect
(\resp locally perfect) topoi, let us mention the Noetherian (\resp locally
Noetherian) topoi introduced in \hyperref[VI.def2.11]{2.11} below (\cf
\hyperref[VI.subsec2.14]{2.14}), and the Zariski and étale topoi of a
coherent scheme (\resp an arbitrary scheme) (IX, note page 42). As a
counterexample,\pageoriginale let us mention the fppf topos of a scheme $S$,
which is locally coherent (and even coherent if $S$ is a coherent scheme,
\ie quasi-compact and quasi-separated)
(\hyperrefext[VII][VII.lem5.6]{5.6}), but is not locally perfect if
$S\neq∅$ (\hyperref[VI.exer1.28]{1.28} f)).

In a locally perfect topos, the notion of constructible object is
particularly stable (\hyperref[VI.prop2.9.3]{2.9.3} below), which is one
reason why the notion seems useful. Another reason is that, like the notion
of a coherent or locally coherent topos, the notion of a perfect or locally
perfect topos is stable under formation of the closed subtopos complementary
to an open (\hyperref[VI.prop4.6]{4.6}), and behaves particularly simply
under the operation of gluing topoi (\hyperref[VI.cor4.11]{4.11},
\hyperref[VI.cor4.14]{4.14}).

\begin{enonce*}{Proposition 2.9.3}\etiquette[2.9.3]{VI.prop2.9.3}
Let $E$ be a locally perfect topos (\hyperref[VI.def2.9.1]{2.9.1}). Then the
full subcategory $E_{\cons}$ of $E$ consisting of the constructible objects
of $E$ (\hyperref[VI.subsubsec1.9.3]{1.9.3}) is stable under finite
inductive limits (and also, of course, under finite projective limits, by
(\hyperref[VI.subsubsec1.9.3]{1.9.3})).
\end{enonce*}

Since the property of an object of $E$ being constructible is local on $E$
(\hyperref[VI.prop1.10]{1.10} (ii)), we are reduced to the case in which $E$
is a perfect topos. But then $E_{\cons}=E_{\coh}$
(\hyperref[VI.cor2.7]{2.7}), and the conclusion follows from Definition
\hyperref[VI.def2.9.1]{2.9.1}.

\begin{enonce*}{Corollary 2.9.4}\etiquette[2.9.4]{VI.cor2.9.4}
Let\pageoriginale $E$ be a locally perfect topos. Then $E$ is perfect if and
only if $E$ is coherent. For every object $X$ of $E$, the induced topos
$E_{/X}$ is locally perfect.
\end{enonce*}

\begin{enonce*}[remark]{Problem 2.9.5}\etiquette[2.9.5]{VI.prob2.9.5}
It is conceivable that perfect topoi are sufficiently special to admit a
\emph{structure theory} that is equally explicit (following a model
proposed by M.\textsc{Artin} at the time of the oral seminar). Let us call a
topos \emph{finite} if it is equivalent to a topos of the form $\widehat{C}$,
where $C$ is a finite category (\cf Exercise
\hyperref[VI.exer3.11]{3.11}), and call a topos \emph{profinite} if it is a
filtered projective limit of finite topoi (in the sense of 6 below, which
applies by \hyperref[VI.exer3.12]{3.12} c)). Since a finite topos is
Noetherian (\hyperref[VI.exer2.17]{2.17} g)), hence perfect, and a filtered
projective limit of perfect topoi is perfect (6), it follows that a profinite
topos is perfect. The question is whether, conversely, every perfect topos
is profinite. This is equivalent to asking whether every finite subcategory
of $E_{\coh}$ is contained in a full subcategory $F_0$ of $E_{\coh}$ that is
equivalent to a category $F_{\coh}$, where $F$ is a finite topos (\cf
\hyperref[VI.exer3.11]{3.11}). (If the answer were negative, one would need
to find manageable additional intrinsic conditions on a perfect topos that
ensure that it is profinite.) Finally, it would remain to study the
structure of profinite topoi in terms of a suitable notion of a
``Karoubian profinite category'', inspired by IV \hyperref[IV.exer7.6]{7.6} h). This\pageoriginale
study has not yet been carried out even in the special case in which $E$ is
the Zariski or étale topos of a good coherent scheme $X$, and should then
give invariants finer than the compact space $X_{\cons}$ (\EGA IV 1.)
associated with $X$.

In the same vein, let us mention the following question: Let $E$ be a
perfect topos, $E'$ the topos induced on an open of $E$, and $E''$ the
corresponding closed subtopos (IV \hyperref[IV.subsubsec9.3.5]{9.9}). If $E'$ and $E''$ are finite topoi,
is the same true of $E$?
\end{enonce*}

% END INCLUDED SGA4_English_current_buildable_source_20260906/source/06/drafts/SGA4_VI_LINES_1383_1678_English.texfrag


% BEGIN INCLUDED SGA4_English_current_buildable_source_20260906/source/06/drafts/SGA4_VI_LINES_1679_2073_English.texfrag
\begin{enonce*}{Proposition 2.10}\etiquette[2.10]{VI.prop2.10}
Let $E$ be a topos. The following conditions are equivalent:
\begin{itemize}
\item[(i)] $E$ admits a full generating subcategory, stable under
  fiber products, consisting of pre-Noetherian objects of $E$.

\item[(ii)] $E$ admits a full generating subcategory consisting of
  quasi-separated pre-Noetherian objects (\ie coherent
  pre-Noetherian objects).

\item[(iii)] $E$ is locally coherent (\hyperref[VI.def2.3]{2.3}), and
  every quasi-compact object of $E$ is pre-Noetherian.

\item[(iii bis)] $E$ is locally coherent and admits a generating
  family consisting of pre-Noetherian objects of $E$.
\end{itemize}
\end{enonce*}

The equivalence of (iii) and (iii bis) follows from
\hyperref[VI.subsec1.32]{1.32}, and it is trivial that
(i) $\Rightarrow$ (iii bis) and (iii) $\Rightarrow$ (i); hence (i) is
equivalent to (iii) and (iii bis). Finally, (i) $\Rightarrow$ (ii)
follows from the last assertion of \hyperref[VI.prop2.1]{2.1},
\pageoriginale and it remains to prove (ii) $\Rightarrow$ (i). This
implication follows from the fact that the full subcategory of $E$
formed by the quasi-separated Noetherian objects, if it is generating,
is stable under fiber products: such a fiber product is quasi-compact
by the definitions, hence Noetherian by the first assertion of
\hyperref[VI.subsec1.32]{1.32}; and it is quasi-separated by the last
assertion of \hyperref[VI.subsec1.32]{1.32}.

\begin{enonce*}{Definition 2.11}\etiquette[2.11]{VI.def2.11}
A topos $E$ is called locally Noetherian if it satisfies the
equivalent conditions of \hyperref[VI.prop2.10]{2.10}, and Noetherian
if, in addition, its final object is coherent (or, equivalently
(\hyperref[VI.subsec1.32]{1.32}), pre-Noetherian and quasi-separated).
An object $X$ of a topos $E$ is called a Noetherian object of $E$ if
the induced topos $E_{/X}$ is a Noetherian topos.
\end{enonce*}

\setcounter{subsection}{11}
\subsection{}\etiquette[2.12]{VI.subsec2.12}
If $E$ is a locally Noetherian topos, then so is every induced topos
$E_{/X}$. Conversely, if one can find a family $(X_i)$ covering the
final object such that the $E_{/X_i}$ are locally Noetherian, then so
is $E$. If $C$ is a full generating subcategory of $E$ consisting of
pre-Noetherian objects and stable under fiber products
(\hyperref[VI.def2.11]{2.11} (i)), then for every $X \in C$,
$E_{/X}$ is a Noetherian topos, \ie the objects of $C$ are even
\emph{Noetherian} (for, if $E$ is locally Noetherian and
$X \in \Ob E$, then $E_{/X}$ is Noetherian if and only if $X$ is
coherent). It follows that a topos $E$ is locally Noetherian if and
only if the final object of $E$ can be covered by Noetherian objects
$X_i$.

\subsection{}\etiquette[2.13]{VI.subsec2.13}
In a\pageoriginale locally Noetherian topos $E$, every monomorphism is
quasi-compact and every morphism is quasi-separated
(\hyperref[VI.subsec1.32]{1.32}). A fortiori, $E$ is an
\emph{algebraic} topos (\hyperref[VI.def2.3]{2.3}), and not merely
locally coherent, \ie locally algebraic. Thus a topos $E$ is
Noetherian if and only if it is locally Noetherian and \emph{coherent}
(\hyperref[VI.def2.3]{2.3}). In other words, if $E$ is a locally
Noetherian topos, then the Noetherian objects of $E$ are precisely the
\emph{coherent} objects of $E$.

\subsection{}\etiquette[2.14]{VI.subsec2.14}
Let $E$ be a locally Noetherian topos. If $E$ is quasi-separated, \ie
if its final object is quasi-separated, then \emph{every} object of
$E$ is quasi-separated (\hyperref[VI.subsec1.32]{1.32}); hence the
pre-Noetherian objects of $E$ are the same as the Noetherian (\ie
coherent) objects of $E$. Since a quotient object of a pre-Noetherian object
is pre-Noetherian, it follows that $E$ satisfies the equivalent
conditions of \hyperref[VI.cor1.25]{1.25} (because it satisfies
condition \hyperref[VI.cor1.25]{1.25} (ii bis)). In particular, if
$C$ is the full subcategory of $E$ consisting of the Noetherian
(coherent) objects of $E$, then the natural functor
$$
\Ind(C) \to E
$$
is an equivalence of categories. We conclude that \emph{a Noetherian
topos is perfect} (\hyperref[VI.def2.9.1]{2.9.1}), and hence that a
locally Noetherian topos is locally perfect.

If one assumes only that $E$ is locally Noetherian, it remains true
that every object of $E$ is the inductive limit of its pre-Noetherian
subobjects (for, in a topos admitting a generating subcategory
consisting\pageoriginale of quasi-compact objects, every object is the
filtered inductive limit of its quasi-compact subobjects). But since a
pre-Noetherian object of $E$ is no longer necessarily coherent, this
statement then has only very limited usefulness, and we know from
(\hyperref[VI.exem1.33]{1.33}) that the conditions of
\hyperref[VI.cor1.25]{1.25} are no longer necessarily satisfied.

\begin{enonce*}[remark]{Example 2.15}\etiquette[2.15]{VI.exem2.15}
(\emph{Topological spaces}.) Let $X$ be a topological space. Then the
following conditions are equivalent: (i) the topological space $X$ is
Noetherian (\ie every open subset of $X$ is quasi-compact, \ie every
increasing sequence of open subsets of $X$ is stationary); (ii) the
topos $\Top(X)$ is Noetherian; (iii) the final object $X$ of
$\Top(X)$ is Noetherian; (iv) the final object $X$ of $\Top(X)$ is
pre-Noetherian.

Indeed, trivially (ii) $\Leftrightarrow$ (iii) $\Rightarrow$ (iv)
$\Leftrightarrow$ (i). On the other hand, (i) implies that the
generating subcategory Ouv(X) of $\Top(X)$, which is stable under
finite projective limits, consists of pre-Noetherian objects, whence
(ii).

Likewise, one sees that the topological space $X$ is locally
Noetherian (\ie is a union of Noetherian open subsets) if and only if
the topos $\Top(X)$ is locally Noetherian.
\end{enonce*}

\setcounter{subsection}{15}
\setcounter{subsubsection}{0}
\subsubsection{}\etiquette[2.15.1]{VI.subsubsec2.15.1}
Thus our Definitions \hyperref[VI.def1.30]{1.30} and
\hyperref[VI.def2.11]{2.11} agree with the accepted terminology for
topological spaces. On the other hand, in the general case we have
chosen to give the term ``Noetherian object'' a\pageoriginale stronger
meaning than the more naive notion of \hyperref[VI.def1.30]{1.30}, so
that all the properties associated with this notion (rather than
merely the grammatical structure of the definition) agree fully with
the intuition attached to \emph{Noetherian} topological spaces and
schemes.

The reader will find further examples in the following exercises.

\begin{enonce*}[remark]{Exercise 2.16}\etiquette[2.16]{VI.exer2.16}
(\emph{Spaces with operators and classifying topoi}.)

Let $X$ be a topological space on which a discrete group $G$ acts,
giving (IV \hyperref[IV.subsec2.3]{2.3}) a topos $E = \Top(X, G)$. The objects $X'$ of this
topos are interpreted as étalé spaces over $X$ with operator group
$G$, and we note that the induced topos $E_{/X'}$ is canonically
equivalent to $\Top(X', G)$.
\begin{enumerate}
\renewcommand{\theenumi}{\alph{enumi}}
\renewcommand{\labelenumi}{\theenumi)}
\item The final object of the topos $E$ is quasi-compact if and only
  if the quotient topological space $X/G$ is quasi-compact (use IV
  \hyperref[IV.subsubsec8.4.1]{8.4.1}).

\item Let $P$ be the object of $E$ defined by the trivial étalé space
  $X \times G$ with the diagonal action of $G$. Show that the induced
  topos $E_{/P}$ is equivalent to the topos $\Top(X)$ (compare IV
  \hyperref[IV.subsubsec5.8.3]{5.8.3}). Conclude that $E$ is locally coherent (\resp locally
  Noetherian) if and only if $\Top(X)$ is locally coherent (\cf
  \hyperref[VI.subsubsec2.4.7]{2.4.7}) (\resp locally Noetherian (\cf
  \hyperref[VI.exem2.15]{2.15})). Show that if $E$ is locally coherent,
  \ie locally algebraic, it is in fact algebraic, using the generating
  family formed\pageoriginale by the objects $T_U$, where $U$ is a
  coherent open subset of $X$, $T_U = G \times U$ with the action
  $g \cdot (u, g') = (u, gg')$ of $G$, and the structural morphism
  $p_U : T_U \to X$ defined by $p_U(u, g) = g\cdot u$.

\item Suppose that $E$ is algebraic, \ie that $\Top(X)$ is algebraic
  (\hyperref[VI.subsubsec2.4.7]{2.4.7}). Show that the structural
  morphism $P \to e$ is quasi-separated. Show that $E$ is
  quasi-separated if and only if $\Top(X)$ is quasi-separated (or,
  equivalently, the topological space $X$ is quasi-separated, \ie the
  intersection of two quasi-compact open subsets of $X$ is a
  quasi-compact open subset), \emph{and} $G$ is finite or $X$ is empty.
  (Compare \hyperref[VI.exem1.33]{1.33}.)

\item Show that the structural morphism $P \to e$ is quasi-compact if
  and only if $G$ is finite. Show that $E$ is coherent (\resp
  Noetherian) if and only if $\Top(X)$ is so, \emph{and} $G$ is finite
  or $X$ is empty.

\item Let $E$ be a topos and $G$ a group of $E$, giving a classifying
  topos $B_G$ (IV \hyperref[IV.subsec2.4]{2.4}). Study the finiteness conditions in $B_G$, using
  what precedes as a guide. Ask the same question starting from a
  strict pro-group $𝒢 = (G_i)_{i \in I}$ of $E$. In particular, show
  that if $E$ is coherent (\resp Noetherian) and the $G_i$ are
  coherent, then $B_{𝒢}$ is coherent (\resp Noetherian).
\end{enumerate}
\end{enonce*}

\begin{enonce*}[remark]{Exercise 2.17}\etiquette[2.17]{VI.exer2.17}
(\emph{Topoi of the form $\widehat{C}$}.)
Let $C$ be a category equivalent to a category $\in 𝒰$, and let
$E = \widehat{C}$ be the category of presheaves on $C$. By means of
the canonical functor $\varepsilon : C \to \widehat{C}$, we identify
$C$ with a full subcategory of $E$.
\begin{enumerate}
\renewcommand{\theenumi}{\alph {enumi}}
\renewcommand{\labelenumi}{(\theenumi)}
\item $C$ is\pageoriginale a generating subcategory of $E$ consisting
  of quasi-compact objects. These objects are pre-Noetherian (\cf
  \hyperref[VI.subsec1.32]{1.32}) if and only if, for every object $X$
  of $C$, every sieve on $X$ is generated by a finite family of
  morphisms $X_i \to X$ of $C$. The final object of $E$ is
  quasi-compact (\resp pre-Noetherian) if and only if there is a final
  subcategory of $C$ whose set of objects is finite (\resp if every
  sieve on $C$ is generated by a subcategory of $C$ whose set of
  objects is finite).

\item The topos $E$ admits a generating subcategory consisting of
  coherent (or, equivalently, quasi-separated) objects if and only if
  the objects of $C$ are coherent in $E$, or equivalently if, for two
  arrows $Y \xrightarrow{f} X$, $Z \xrightarrow{g} X$ in $C$, the
  object $\fprod{Y}{Z}{X}$ of $E$ is quasi-compact, \ie if one can find
  a finite family of commutative squares in $C$
\[
\xymatrix@=1cm{
 & T_i \ar[ld] \ar[rd]& \\
Y\ar[rd]_{f}&    & Z\ar[ld]^{g}\\
 & X  &,
}
\]
such that every other commutative square in $C$
\[
\xymatrix@=1cm{
 & T \ar[ld] \ar[rd]& \\
Y\ar[rd]_{f}&    & Z\ar[ld]^{g}\\
 & X  &,
}
\]
comes\pageoriginale from a morphism $T \to T_i$. (Note that if
$(F_i)_{i \in I}$ is a generating family of $E$, then every
$X \in \ob C$ is isomorphic to a direct factor of one of the $F_i$.)

\item The topos $E$ is locally coherent, \ie is generated by a family
  of coherent and \emph{algebraic} objects, if and only if the objects
  $X$ of $C$ are coherent and algebraic objects of $E$, \ie if
  condition b) and the following condition hold: for every pair of
  arrows $f$, $g : Y \rightrightarrows Z$ in $C$ \emph{over an object
  $X$ of} $C$, the kernel of this pair of arrows in $E$ is a
  quasi-compact object of $E$, \ie there is a finite family of
  commutative diagrams in $C$
\[
\xymatrix@C=1.5cm{
T_i \ar[r]& Y \ar@<2pt>[r]^{f,g}\ar@<-2pt>[r]&Z
}
\]
such that every other commutative diagram in $C$
\[
\xymatrix@C=1.5cm{
T \ar[r]& Y \ar@<2pt>[r]^{f,g}\ar@<-2pt>[r]&Z
}
\]
comes from a morphism $T \to T_i$. The topos $E$ is algebraic if and
only if it satisfies condition b) and the preceding condition, but
with an arbitrary pair of arrows $(f, g)$ of $C$ (not necessarily
over an object $X$ of $C$). The topos $E$ is quasi-separated if and
only if it satisfies the conditions of b), and if, in addition, for
two objects $X$, $Y\in\ob C$, the product $X\times Y$ in $E$ is
quasi-compact. The topos $E$ is coherent if and only if it satisfies
the two preceding conditions and there is a final subcategory of $C$
whose underlying set is finite.

\item Give\pageoriginale an example in which the objects of the
  generating subcategory $C$ of $E$ are pre-Noetherian, but $E$ admits
  no generating family consisting of coherent objects (\ie b) fails,
  or the objects of $C$ are not all coherent). For this, take for $C$
  the category having as objects $X$, $T_i$ $(i = 0, 1, \ldots)$, and
  $e$ (the final object), the only morphisms between these objects,
  besides the structural morphisms to $e$ and the identity morphisms,
  being morphisms $u_i : T_i \to X$ and $v_i : X \to T_i$ subject to
  the conditions $u_i v_i = \id_X$, and finally set
  $p_i = v_i u_i$ (which necessarily satisfies
  $p_i^2 = p_i$). One verifies that the only sieves on $X$ or on any
  $T_i$ are the two trivial sieves, and that $e$ admits exactly one
  nontrivial sieve. Thus, by a), the objects of $C$ are
  pre-Noetherian objects of $E$. Nevertheless, the product
  $X \times X$ in $E$ is not quasi-compact, since it does not satisfy
  the criterion of b).

\item Give an example in which the generating subcategory $C$ of $E$
  consists of coherent objects of $E$, but $E$ is not locally
  coherent. For this, take for $C$ the category whose set of objects
  consists of the distinct objects $e$ (the final object), $Y$, $Z$,
  and $T_i$ $(i=0, 1, \ldots)$, with the only morphisms between these
  objects, besides the morphisms to the final object $e$ and the
  identity morphisms, being morphisms $f$, $g : Y \rightrightarrows Y$
  subject to the conditions $fu_i = gu_i$. One verifies that $C$
  satisfies condition b) (with a little patience; one finds that all
  fiber products of objects of $C$ are isomorphic to $∅_E$ or belong
  to $C$,\pageoriginale with the exception of
  $\fprod{Z}{Z}{e}$ and $\fprod{Y}{Z}{e}$, which are covered by two
  elements of $C$), but clearly $\Ker(f, g)$ does not satisfy
  condition c).

\item Give an example in which $C$ is stable under fiber products (a
  fortiori, $E$ is a locally coherent, \ie locally algebraic, topos),
  but $E$ is not an algebraic topos. (Take the example given in d),
  and the full subcategory $C'$ of $E$ consisting of the objects $Y$,
  $Z$, $T_{i}$ $(i > 0)$, and $∅_{E}$.)

\item Suppose that the category $C$ is finite. Prove that the topos
  $E$ is Noetherian; that its Noetherian (\ie quasi-compact) objects
  are the contravariant functors $F : C^{\circ} \to (\Ens)$ such that,
  for every object $X$ of $C$, $F(X)$ is a \emph{finite} set; and that
  the fiber functors on $E$ take Noetherian objects to finite sets
  (\cf IV \hyperref[IV.exer7.6]{7.6}. h)).

\item Suppose that every morphism $f: X \to X$ of $C$ factors as a
  composite $ip$, where $i$ is a monomorphism and $p$ has a right
  inverse. Let $X$ be an object of $C$, let $I(X)$ be the set of its
  subobjects in the sense of $C$, regarded as an ordered set, and let
  $S(I(X))$ be the set of subsets $U$ of $I(X)$ such that, for two
  elements $X'$, $X'' \in I(X)$, the conditions $X'\in U$ and
  $X''\leq X'$ imply $X''\in U$. Show that the ordered set of
  subobjects in $E$ of $X$ is isomorphic to the set $S(I(X))$ (ordered
  by inclusion). Conclude that if $I(X)$ is finite, then $X$ is a
  pre-Noetherian object of $E$. In particular, if (in addition to the
  factorization condition above) $C$ is stable under fiber products
  (\resp under finite projective limits), and\pageoriginale if, for
  every $X \in \ob C$, the set $I(X)$ of subobjects of $X$ in $C$ is
  finite, then the topos $E$ is locally Noetherian (\resp Noetherian).

\item Let $C$ be the category of finite sets (or, alternatively, of
  nonempty finite sets) $\in 𝒰$, so that $E = \widehat{C}$ is the
  category of augmented simplicial sets (\resp of simplicial sets
  simpliciter). Show that $E$ is a Noetherian topos. (Use h).) Taking
  a nonessentially constant pro-object $(X_i)_{i \in I}$ of $C$, show
  that $E$ admits fiber functors that do not take Noetherian objects
  to Noetherian objects (\ie to finite sets).

\item Consider the following diagram of implications among properties
  of a topos $E$:
\[
\xymatrix@R=.75cm{
**[r]\txt{E is Noetherian} \ar@{=>}[d] \ar@{=>}[r]& **[r]\txt{E is
  locally Noetherian} \ar@{=>}[d] & \ar@{=>}[rdd]\\
**[r]\txt{E is coherent} \ar@{=>}[r] & **[r]\txt{E
is algebraic} \ar@{=>}[d]\\
 & **[r]\txt<9.5pc>{E is loc. coherent, i.e.\\E \raisebox{-1ex}{cohalg}
  generates E} \ar@{=>}[d]& & \txt{E \raisebox{-1ex}{pre-Noeth}
  generates E} \ar@{=>}[ldd] \\
 & **[r]\txt{E \raisebox{-1ex}{coh} generates E} \ar@{=>}[d]&\\
 & **[r]\txt{E \raisebox{-1ex}{qucpct} generates E}  &&.\\
}
\]
Show that every implication in this diagram is strict, and that among
the notions considered there are no implications other than the
composites of the implications in the preceding diagram. Here
$E_{\coh}, E_{\text{qucpct}}$, $E_{\text{prénoeth}}$,
$E_{\text{cohalg}}$\pageoriginale denote respectively the full
subcategories of $E$ consisting of the coherent, \resp quasi-compact,
\resp pre-Noetherian, \resp coherent and algebraic objects of $E$.
(One may restrict attention to topoi of the form $\widehat{C}$, using
the results stated in d), e), and f).)

\item Answer the following question (to which the author of these
  lines does not know the answer): can $E = \widehat{C}$ be locally
  Noetherian without the $\Hom(X, Y)$ (for $X$, $Y \in \ob C$) being
  finite?
\end{enumerate}
\end{enonce*}

\begin{enonce*}[remark]{Exercise 2.18}\etiquette[2.18]{VI.exer2.18}
Let $E$ be a topos.
\begin{enumerate}
\renewcommand{\theenumi}{\alph{enumi}}
\renewcommand{\labelenumi}{\theenumi)}
\item Let $X$ be a pre-Noetherian object of $E$. Show that $X$ is
  isomorphic to a finite sum of connected objects (IV \hyperref[IV.subsubsec4.3.5]{4.3.5}) of $E$.
  (Prove by contradiction that an increasing sequence of partitions
  of $X$ (IV \hyperref[IV.exer8.7]{8.7}) is stationary.)

\item Let $X$ be an object of $E$ and let
  $(f_i : X_i \to X)_{i \in I}$ be a covering family of $X$. Show that
  if each $X_i$ is isomorphic to a sum of connected objects of $E$,
  then so is $X$. (Reduce to the case in which the $X_i$ are connected,
  then to the case in which they are subobjects of $X$, and consider
  on the index set $I$ the equivalence relation $R$ generated by the
  relation $X_i \cap X_j \neq \phi_E$. Show that if $J = I/R$, then
  $X$ is the sum of the
  $X(i) = \displaystyle{\mathop{\Sup}_{i \in I}} X_i .)$

\item Conclude from b) that if $(X_i)_{i \in I}$ denotes a generating
  family of $E$, then $E$ is locally connected (IV \hyperref[IV.exer8.7]{8.7}. $\ell$)) if
  and\pageoriginale only if every $X_i$ $(i \in I)$ is isomorphic to a
  sum of connected objects of $E$.

\item Conclude from a) and c) that if $E$ admits a generating
  subcategory consisting of pre-Noetherian objects, in particular if
  $E$ is locally Noetherian, then $E$ is locally connected, and a
  fortiori is isomorphic to the sum topos (IV \hyperref[IV.exer8.7]{8.7} b)) of a family of
  connected topoi (IV \hyperref[IV.exer8.7]{8.7} e)), \ie topoi whose final object is
  connected. (In particular, one can associate with $E$, for every
  morphism $f : P \to E$, where $P$ is a ``nonempty connected and
  simply connected'' topos (IV \hyperref[IV.exer2.7.5]{2.7.5}), a fundamental pro-group
  $\pi_1(E, f)$. Note that if $E$ is locally Noetherian and ``nonempty,''
  such an $f$ can always be found, with $P$ the punctual topos, by a
  result of \textsc{Deligne} ([9]).
\end{enumerate}
\end{enonce*}

% END INCLUDED SGA4_English_current_buildable_source_20260906/source/06/drafts/SGA4_VI_LINES_1679_2073_English.texfrag


% BEGIN INCLUDED SGA4_English_current_buildable_source_20260906/source/06/drafts/SGA4_VI_LINES_2074_2418_English.texfrag
\section{Finiteness conditions for a morphism of topoi}\etiquette[3]{VI.sec3}

\begin{enonce*}{Definition 3.1}\etiquette[3.1]{VI.def3.1}
Let $f:E'\to E$ be a morphism of topoi. We say that $f$ is
\emph{quasi-compact (\resp quasi-separated)} if, for every
quasi-compact (\resp quasi-separated) object $X$ of $E$,
$f^{\ast}(X)$ is quasi-compact (\resp quasi-separated). We say that
$f$ is coherent if $f$ is quasi-compact and quasi-separated.
\end{enonce*}

\setcounter{subsection}{1}
\setcounter{subsubsection}{0}
\subsubsection{}\etiquette[3.1.1]{VI.subsubsec3.1.1}
Note that if $f$ has any one of the preceding properties, then every
morphism of topoi isomorphic to $f$ has the same property. It is also
immediate that the composite of two quasi-compact (\resp
quasi-separated, \resp coherent) morphisms of topoi is again
quasi-compact (\resp quasi-separated, \resp coherent).

\begin{enonce*}{Proposition 3.2}\etiquette[3.2]{VI.prop3.2}
With the\pageoriginale notation of \hyperref[VI.def3.1]{3.1}, let
$(X_i)_{i\in I}$ be a generating family of quasi-compact (\resp
coherent) objects of $E$. In the parenthetical case, suppose that
$E'$ admits a generating family formed by quasi-compact objects. Then
$f$ is quasi-compact (\resp coherent) if and only if, for every
$i\in I$, $f^{\ast}(X_i)$ is quasi-compact (\resp coherent).
\end{enonce*}

Necessity is immediate. For sufficiency in the first case, note that,
for every quasi-compact object $X$ of $E$, there is a finite
epimorphic family of morphisms $X_i\to X$. Hence there is a finite
epimorphic family $f^{\ast}(X_i)\to f^{\ast}(X)$, with the
$f^{\ast}(X_i)$ quasi-compact by hypothesis, and therefore
$f^{\ast}(X)$ is quasi-compact (\hyperref[VI.prop1.3]{1.3}). In the
second case, it remains to see that if $X$ is a quasi-separated
object of $E$, then $f^{\ast}(X)$ is quasi-separated. The hypothesis
on $X$ may be expressed (\hyperref[VI.cor1.17]{1.17}) by the
existence of an epimorphic family of morphisms $X_i\to X$ such that
the $\fprod{X_i}{X_j}{X}$ are quasi-compact. We then have a covering
family $f^{\ast}(X_i)\to f^{\ast}(X)$ such that the fiber products
$\fprod{f^{\ast}(X_i)}{f^{\ast}(X_j)}{f^{\ast}(X)}$ are
quasi-compact (since, by the first case, $f^{\ast}$ carries
quasi-compact objects to quasi-compact objects). As the
$f^{\ast}(X_i)$ are coherent, the conclusion follows again from
\hyperref[VI.cor1.17]{1.17}.

\begin{enonce*}{Corollary 3.3}\etiquette[3.3]{VI.cor3.3}
Let $C$ and $C'$ be two $𝒰$-sites in which fiber products are
representable and every covering family admits a finite covering
subfamily, and let $g:C\to C'$ be a morphism of sites \emph{(IV
\hyperref[IV.subsubsec4.9.1]{4.9.1})}. Then the morphism of topoi
$f:C^{\sim}\to{C'}^{\sim}$ defined by $g$ is coherent.
\end{enonce*}

Indeed,\pageoriginale $\epsilon_C(C)$ (\resp
$\epsilon_{C'}(C')$) is a generating family of $C^{\sim}$ (\resp
${C'}^{\sim}$) satisfying the corresponding conditions of
\hyperref[VI.prop3.2]{3.2} (\hyperref[VI.cor2.1.1]{2.1.1}).

\begin{enonce*}{Proposition 3.4}\etiquette[3.4]{VI.prop3.4}
Let $E$ be a locally coherent topos (\hyperref[VI.def2.3]{2.3}), and
let $f:X\to Y$ be an arrow of $E$, whence a morphism of induced topoi
\emph{(IV \hyperref[IV.eq5.5.2]{(5.5.2)})} $E_{/X}\to E_{/Y}$. This morphism of topoi is
quasi-compact (\resp quasi-separated, \resp coherent) if and only if
$f$ is a quasi-compact (\resp quasi-separated, \resp coherent)
morphism.
\end{enonce*}

The quasi-compact case follows immediately from the definitions
(without any condition on $E$); the quasi-separated case is precisely
\hyperref[VI.cor2.8]{2.8}; and the coherent case follows by combining
the two preceding cases.

\begin{enonce*}{Proposition 3.5}\etiquette[3.5]{VI.prop3.5}
Let $f:E'\to E$ be a morphism of locally coherent topoi. Let
$(X_i)_{i\in I}$ be a generating family in $E$ formed by coherent
algebraic objects. The following conditions are equivalent:
\begin{itemize}
\item[(i)] For every quasi-compact arrow $u$ of $E$, $f^{\ast}(u)$ is
  quasi-compact.

\item[(i bis)] For every arrow $u:X_i\to X_j$, $f^{\ast}(u)$ is
  quasi-compact.

\item[(ii)] For every coherent object $Y'$ of $E'$, every coherent
  algebraic object $Y$ of $E$, and every arrow
  $v:Y'\to f^{\ast}(Y)$, the corresponding morphism of topoi
$$
f_v:{E'}_{/Y'}\to E_{/Y}
$$
obtained by composing\pageoriginale the localization morphism
${E'}_{/Y'}\to{E'}_{/f^{\ast}(Y)}$ deduced from $v$
\emph{(IV \hyperref[IV.eq5.5.2]{(5.5.2)})} with the morphism
${E'}_{/f^{\ast}(Y)}\to E_{/Y}$ induced by $f$
\emph{(IV \hyperref[IV.eq5.10.1]{(5.10.1)})} is coherent.

\item[(ii bis)] The same condition as in \emph{(ii)}, but restricted
  to a collection of data
  $v_{\alpha}:{Y'}_{\alpha}\to f^{\ast}(Y_{\alpha})$ such that the
  ${Y'}_{\alpha}$ are algebraic and cover the final object of $E'$.

\item[(ii ter)] (When a generating family
  $(X'_{i'})_{i'\in I'}$ in $E'$ formed by coherent objects is given,
  together with, for every $i'\in I'$, an $i\in I$ and an arrow
  $v_{i'}:X'_{i'}\to f^{\ast}(X_i)$.) For every $i'\in I'$, the
  morphism of topoi $E'_{/{X'}_{i'}}\to E_{/X_i}$ defined by
  $v_{i'}$ is coherent.
\end{itemize}
\end{enonce*}

Since every arrow $X_i\to X_j$ is quasi-compact, it is clear that (i)
$\Rightarrow$ (i bis). Conversely, suppose (i bis) holds and let us
prove (i). Let $u:X\to Y$ be a quasi-compact arrow of $E$. In terms
of an epimorphic family $X_i\to Y$, the hypothesis on $u$ says that
the $\fprod{X}{X_i}{Y}$ are quasi-compact
(\hyperref[VI.cor1.16]{1.16}), \ie that for every $i$ there exists a
finite epimorphic family of morphisms
$X_j\to\fprod{X}{X_i}{Y}$ (\hyperref[VI.prop1.3]{1.3}). We then have
the corresponding epimorphic family
$f^{\ast}(X_i)\to f^{\ast}(Y)$, such that for every $i$ there is a
finite epimorphic family
$f^{\ast}(X_j)\to f^{\ast}(\fprod{X}{X_i}{Y})\simeq
\fprod{f^{\ast}(X)}{f^{\ast}(X_i)}{f^{\ast}(Y)}$, whose composite
with the projection $\pr_2$ to $f^{\ast}(X_i)$ is a quasi-compact
morphism $f^{\ast}(X_j)\to f^{\ast}(X_i)$ by hypothesis (i bis). It
follows (\hyperref[VI.cor1.11]{1.11} (i)) that, for every $i$,
$\pr_2:\fprod{f^{\ast}(X)}{f^{\ast}(X_i)}{f^{\ast}(Y)}$ is
quasi-compact, and hence (\hyperref[VI.prop1.10]{1.10} (ii)) that
$f^{\ast}(u):f^{\ast}(X)\to f^{\ast}(Y)$ is quasi-compact.

It remains\pageoriginale to prove (i) $\Rightarrow$ (ii) and
(ii bis) $\Rightarrow$ (i); the implication (ii) $\Rightarrow$
(ii bis) is immediate, and (ii ter) is a special case of (ii bis).
The implication (i) $\Rightarrow$ (ii) is immediate: indeed, an
object $X$ of $E_{/Y}$ is quasi-compact (\resp quasi-separated) if
and only if its structural morphism $X\to Y$ (\resp its diagonal
morphism $X\to\fprod{X}{X}{Y}$) is quasi-compact
(\hyperref[VI.cor2.7]{2.7}), and the same criterion of
quasi-compactness and quasi-separatedness applies to
$f^{\ast}_Y(X')$. Finally, let us prove (ii bis) $\Rightarrow$ (i).
Let $u:X\to Y$ be a quasi-compact arrow in $E$; we must show that
$f^{\ast}(u)$ is quasi-compact. Since the $Y'_{\alpha}$ cover the
final object of $E'$, it is enough to prove that the
$f^{\ast}(u)|Y'_{\alpha}$ are quasi-compact
(\hyperref[VI.prop1.10]{1.10} (ii)). These morphisms are precisely
the $f^{\ast}_{v_\alpha}(u|Y_{\alpha})$, so we are reduced to proving
the following result.

\begin{enonce*}{Corollary 3.6}\etiquette[3.6]{VI.cor3.6}
Let $f:E'\to E$ be a coherent morphism of locally coherent topoi.
Then, for every quasi-compact arrow $u$ of $E$, $f^{\ast}(u)$ is
quasi-compact.
\end{enonce*}

By the implication (i bis) $\Rightarrow$ (i) of
\hyperref[VI.prop3.5]{3.5} already proved, it is enough to show that,
if $u:X\to Y$ is an arrow of $E$ with $X,Y$ coherent, then
$f^{\ast}(u)$ is a quasi-compact morphism. This follows immediately
from the hypothesis on $f$, which implies that $f^{\ast}(X)$ and
$f^{\ast}(Y)$ are coherent as well.

\begin{enonce*}{Definition 3.7}\etiquette[3.7]{VI.def3.7}
Let $f:E'\to E$\pageoriginale be a morphism of locally coherent
topoi. We say that $f$ is \emph{locally coherent} if it satisfies the
equivalent conditions of \hyperref[VI.prop3.5]{3.5}.
\end{enonce*}

Observe that ``locally'' means locally on the source.

\setcounter{subsection}{7}
\setcounter{subsubsection}{0}
\subsubsection{}\etiquette[3.7.1]{VI.subsubsec3.7.1}
This notion again depends only on the isomorphism class of the
morphism of topoi $f$, and it is plainly stable under composition of
morphisms of topoi. Moreover, by \hyperref[VI.prop3.5]{3.5}, if $f$
is coherent then it is locally coherent; the converse holds if $E'$
and $E$ are coherent (by \hyperref[VI.prop3.5]{3.5} (ii bis)).
Criterion \hyperref[VI.prop3.5]{3.5} (ii bis) immediately gives the
following criterion:

\begin{enonce*}{Corollary 3.7}\etiquette[3.7]{VI.cor3.7}
Let $f:E'\to E$ be a morphism of locally coherent topoi, let
$(Y_i)_{i\in I}$ and $({Y'}_i)_{i\in I}$ be two families of objects
of $E$ and $E'$, and, for every $i\in I$, let
$v_i:Y'_i\to f^{\ast}(Y_i)$ be a morphism, whence a morphism of topoi
$$
f_{v_i}:E'_{/Y'_i}\to E_{/Y_i}{\quoi}.
$$
Suppose that the ${Y'}_i$ cover the final object of $E'$. Then $f$ is
locally coherent if and only if, for every $i\in I$, $f_{v_i}$ is
locally coherent.
\end{enonce*}

In particular, applying this to the identity morphism of an induced
topos gives:

\begin{enonce*}{Corollary 3.8}\etiquette[3.8]{VI.cor3.8}
Let $E$ be a locally coherent topos, and let $v:X\to Y$ be an arrow of
$E$. Then the morphism of induced topoi $E_{/X}\to E_{/Y}$ defined by
$v$ is locally coherent.
\end{enonce*}

\begin{enonce*}{Remark 3.9}\etiquette[3.9]{VI.rem3.9}
It seems\pageoriginale that all morphisms of algebraic topoi
encountered in practice are locally coherent (which is why the term
``locally coherent'' is probably not destined to be used very often).
Equivalently, the morphisms between coherent topoi encountered in
practice are coherent. Nevertheless, one can construct noncoherent
morphisms between coherent topoi and, more particularly, a
noncoherent morphism from the punctual topos \emph{(IV \hyperref[IV.subsec2.2]{2.2})} to a
Noetherian topos $E$, \emph{\cf \hyperref[VI.exer2.17]{2.17} i)}.
\end{enonce*}

\begin{enonce*}[remark]{Examples 3.10}\etiquette[3.10]{VI.exem3.10}
(\emph{Topoi associated with schemes}.)
Return to the example \hyperref[VI.subsec1.22]{1.22} of the category
$(\Sch)$ with the topologies $T_i$ of \SGA IV 6.3. For every scheme
$X$, let $X_{T_i}$ be the $𝒱$-topos defined by the induced site
$(\Sch)_{/X}$. A morphism of schemes $f:X\to Y$ then defines a
morphism of topoi $f_{T_i}:X_{T_i}\to Y_{T_i}$, and it follows from
\hyperref[VI.prop3.4]{3.4} and \hyperref[VI.subsec1.22]{1.22} that
this morphism of topoi is quasi-compact (\resp quasi-separated,
\resp coherent) if and only if the morphism of schemes $f$ is
quasi-compact (\resp quasi-separated, \resp coherent) in the usual
sense of \EGA IV 1. In any event, $f_{T_i}$ is locally coherent by
\hyperref[VI.cor3.8]{3.8}.

One may also associate with a scheme $X$ the $𝒰$-topoi $\Top(X)$ and
$\Top(X_{\text{ét}})$ as in \hyperref[VI.subsubsec1.22.1]{1.22.1},
and to every morphism of schemes $f:X\to Y$ there are then associated
morphisms $\Top(f)$ and $\Top(f_{\text{ét}})$ of the corresponding
topoi. Let $T(f)$ denote either of these two morphisms of
topoi.\pageoriginale One verifies immediately by
\hyperref[VI.prop3.2]{3.2} that this morphism of topoi is always
locally coherent, and that it is quasi-compact (\resp quasi-separated,
\resp coherent) if and only if the morphism of schemes $f$ is
quasi-compact (\resp quasi-separated, \resp coherent) in the usual
sense.

These observations show once again that the terminology introduced
in the present section is compatible with the accepted terminology
of scheme theory, and therefore deserves to be accepted by even the
most recalcitrant reader.
\end{enonce*}

\begin{enonce*}[remark]{Exercise 3.11}\etiquette[3.11]{VI.exer3.11}
(\emph{Coherent topoi and pretopoi}.)
\begin{enumerate}
\renewcommand{\theenumi}{\alph {enumi}}
\renewcommand{\labelenumi}{\theenumi)}
\item A $𝒰$-pretopos (or simply a pretopos) is a category $C$
  satisfying the following conditions:
\begin{enumerate}
\renewcommand{\theenumii}{\arabic{enumii}}
\renewcommand{\labelenumii}{\theenumii)}
\item Finite projective limits are representable in $C$.

\item Finite coproducts are representable in $C$, and they are
  disjoint and universal.

\item Equivalence relations in $C$ are effective, and every
  epimorphism in $C$ is universally effective.

\item $C$ is equivalent to a category $\in 𝒰$.
\end{enumerate}
Show that if $E$ is a coherent $𝒰$-topos, then the full subcategory
$E_{\coh}$ of $E$ formed by its coherent objects is a $𝒰$-pretopos,
and that the inclusion functor $E_{\coh}\to E$ is left exact,
commutes with\pageoriginale finite coproducts, and commutes with
passage to the quotient by an equivalence relation. (Use
\hyperref[VI.subsubsec1.5.3]{1.5.3}, \hyperref[VI.cor1.15]{1.15},
and \hyperref[VI.cor1.17.1]{1.17.1}.) Show that the topology induced
by $E$ on $C$ is the ``precanonical'' topology, \ie the topology whose
covering families $X_i\to X$ are those that admit a \emph{finite}
subfamily covering for the canonical topology of $C$. Consequently,
$E$ can be recovered up to equivalence from the pretopos
$C=E_{\coh}$ as the topos $C^{\sim}$ (with $C$ endowed with its
precanonical topology).

\item Let $C$ be a $𝒰$-pretopos. Endow $C$ with the precanonical
  topology. Show that every morphism $f$ of $C$ factors as $f''f'$,
  with $f'$ an epimorphism and $f''$ a monomorphism. Show that the
  topos $E=C^{\sim}$ is coherent and that the canonical functor
  $\epsilon:C\to E$ induces an equivalence from $C$ to the
  subcategory $E_{\coh}$ of $E$. (First show that $\varepsilon$ is a
  fully faithful left exact functor commuting with finite coproducts
  and with passage to the quotient by an equivalence relation, and
  then that a coherent subobject in $E$ of an object of $C$ lies in
  the essential image of $C$.) Consequently, the $𝒰$-pretopos $C$
  can be recovered up to equivalence of categories from its associated
  $𝒰$-topos $E=C^{\sim}$.

\item Let $C$ and $C'$ be two $𝒰$-pretopoi. Show that a functor
  $\varphi:C'\to C$ is a morphism of sites from $C$ to $C'$ (for the
  precanonical topologies) if and only if $\varphi$ is left exact and
  commutes with finite coproducts and with passage to a quotient by an
  equivalence relation.

\item With the\pageoriginale notation of c), suppose that $C$ and
  $C'$ are associated with two coherent $𝒰$-topoi $E,E'$ as in a).
  Show that the canonical functor (IV \hyperref[IV.subsubsec4.9.3]{4.9.3})
$$
\Mmorsite(C,C')\to\Hhomtop(E,E')
$$
induces an equivalence from its source (described explicitly in c))
to the full subcategory of its target formed by the \emph{coherent}
morphisms of topoi $E\to E'$. A quasi-inverse is obtained by
associating with every coherent morphism of topoi $f:E\to E'$ the
functor $E'_{\coh}=C'\to E_{\coh}=C$ induced by $f^{\ast}$.

\item Let $C$ be a $𝒰$-pretopos and $E=C^{\sim}$. Express directly in
  terms of exactness properties of $C$ the equivalent conditions of
  \hyperref[VI.cor1.25]{1.25} (consult
  \hyperref[VI.exer1.28]{1.28} b). Show that an object $X$ of $C$ is
  Noetherian in $E$ if and only if every increasing sequence of
  subobjects of $X$ in $C$ is stationary, and hence that $E$ is
  Noetherian if and only if every object of $C$ satisfies the
  preceding condition.
\end{enumerate}
\end{enonce*}

\begin{enonce*}[remark]{Exercise 3.12}\etiquette[3.12]{VI.exer3.12}
A topos $E$ is called a \emph{finite} topos if there exists a finite
category $C$ such that $E$ is equivalent to $\widehat{C}$.
\begin{enumerate}
\renewcommand{\theenumi}{\alph{enumi}}
\renewcommand{\labelenumi}{\theenumi)}
\item The topos $E$ is finite if and only if $E$ has enough essential
  points (IV \hyperref[IV.exer7.6]{7.6} b)) and the category $\Ppointess(E)$ of essential
  points of $E$ is equivalent to a finite category. (Use IV \hyperref[IV.exer7.6]{7.6} d)
  and h).)

\item Suppose\pageoriginale that $E$ is finite. Prove that every point
  of $E$ is essential. Prove that the 2-functors
  $E\mapsto\Ppoint(E)$ and $C\mapsto\widehat{C}$ establish
  2-equivalences between the 2-category of finite topoi $E$ and the
  2-category of Karoubian categories (IV \hyperref[IV.exer7.5]{7.5} a)) $C$ that are
  equivalent to finite categories, and that every morphism of finite
  topoi is essential (IV \hyperref[IV.exer7.6]{7.6} a)). (Use a) and IV \hyperref[IV.exer7.6]{7.6} h).)

\item Suppose that $E=\widehat{C}$, with $C$ equivalent to a finite
  category, and let $F$ be a coherent topos. Recall (IV \hyperref[IV.eq4.6.3.1]{(4.6.3.1)})
  that morphisms of topoi $f:E\to F$ correspond to functors
  $u:C\to\Ppoint(F)$. Show that $f$ is coherent if and only if $f$
  carries every coherent point of $E$ (\hyperref[VI.def3.1]{3.1}) to
  a coherent point of $F$, or equivalently if, for every $X\in\ob C$,
  $u(X)$ is a \emph{coherent} point of $F$. (Use
  \hyperref[VI.prop3.2]{3.2} and \hyperref[VI.exer2.17]{2.17} g).) Conclude that every
  morphism $f:E\to F$ from a finite topos to a finite topos is
  coherent.

\item Let $X$ be a sober topological space (IV \hyperref[IV.subsubsec4.2.1]{4.2.1}). The topos
  $\Top(X)$ (IV \hyperref[IV.subsec2.1]{2.1}) is finite if and only if $X$ is a finite set.
  (Use IV \hyperref[IV.subsubsec7.1.6]{7.1.6}.)
\end{enumerate}
\end{enonce*}


% END INCLUDED SGA4_English_current_buildable_source_20260906/source/06/drafts/SGA4_VI_LINES_2074_2418_English.texfrag


% BEGIN INCLUDED SGA4_English_current_buildable_source_20260906/source/06/drafts/SGA4_VI_LINES_2419_2936_English.texfrag
\section{Finiteness conditions in a topos obtained by
  gluing}\etiquette[4]{VI.sec4}

The present section will not be used later in the Seminar.

\subsection{}\etiquette[4.1]{VI.subsec4.1}
Let $E$ be a topos and $U$ an open of $E$ (\ie a subobject of the final
object of $E$), and consider the corresponding open and closed subtopoi
of\pageoriginale $E$ (IV 9)
$$
E'= E_{/U}, E'' = E_{/U}{\quoi}.
$$
We denote by
$$
j : E'\to E, i : E''\to E
$$
the canonical morphisms of topoi. Recall (\hyperref[VI.prop3.4]{3.4}) that
for $j$ to be a quasi-compact (\resp coherent) morphism of topoi, it is
necessary and sufficient that the inclusion
$$
j : U \to e
$$
of $U$ into the final object of $E$ (an inclusion that we shall also denote
by $j$) be a quasi-compact (\resp coherent) morphism in $E$; moreover,
$j:E'\to E$ is always quasi-separated (because $j:U\to e$ is so
(\hyperref[VI.cor1.8.1]{1.8.1})). We propose to give criteria for the topos
$E''$ to be coherent and for the morphism of topoi $i:E''\to E$ to be
coherent. First note:

\begin{enonce*}{Proposition 4.2}\etiquette[4.2]{VI.prop4.2}
With the notation of \hyperref[VI.subsec4.1]{4.1}, the inclusion morphism
$i:E''\to E$ is quasi-compact, \ie for every quasi-compact object $X$ of
$E$, the object $i^{\ast}(X)$ of $E''$ is quasi-compact; moreover, if $X$ is
pre-Noetherian (\hyperref[VI.def1.30]{1.30}), then $i^{\ast}(X)$ is
pre-Noetherian.
\end{enonce*}

This follows from \hyperref[VI.def1.1]{Definition 1.1} and the

\begin{enonce*}{Lemma 4.2.1}\etiquette[4.2.1]{VI.lem4.2.1}
The map\pageoriginale $Y\mapsto i^{\ast}(Y)$ induces an isomorphism of
ordered sets between the set of subobjects of $X$ containing the subobject
$X_U=X\times U$ of $X$, and the set of subobjects of $i^{\ast}(X)$.
\end{enonce*}

Note that since $i_{\ast}$ is conservative (because it is fully faithful)
and left exact (in particular, it commutes with fiber products), it follows
immediately that a morphism $u:Y\to Z$ in $E''$ is a monomorphism if and only
if $i_{\ast}(u):i_{\ast}(Y)\to i_{\ast}(Z)$ is one. Consequently, on
identifying (by $i_{\ast}$) $E''$ with the full subcategory of $E$ consisting
of objects $Z$ satisfying the equivalent conditions of IV \hyperref[IV.subsec9.7]{9.7} 1), the
subobjects in $E''$ of the object $Z$ of $E''$ identify with the subobjects
$Y$ of $Z$ in $E$ that belong to $E''$, or, equivalently, that contain
$Z_U\xrightarrow{\sim}U$. For $Z$ of the form
$i_{\ast}i^{\ast}(X)=X_{\complement U}$ given by the amalgamated sum
IV \hyperref[IV.subsubsec9.5.1]{(9.5.1)},
$$
i_{\ast} i^{\ast} (X) = X_{\complement U} = X \amalg_{X_U} U {\quoi},
$$
the datum of a subobject $Y$ of the latter in $E$ is equivalent to the datum
of a pair consisting of a subobject $Y_1$ of $X$ and a subobject $Y_2$ of
$U$, subject to the condition that the inverse images of these subobjects
in $X_U$ coincide. The condition that $Y$ contain $U$ then means that
$Y_2=U$, and the remaining condition on $Y_1$ is that it contain $X_U$.
This therefore establishes a bijection between the set of subobjects of
$i^{\ast}(X)$ in $E''$ and the set of subobjects of $X$ that contain $X_U$.
It is clear\pageoriginale that this is an isomorphism of ordered sets and
that its inverse map is indeed the one asserted in \hyperref[VI.lem4.2.1]{4.2.1}.

\begin{enonce*}{Corollary 4.3}\etiquette[4.3]{VI.cor4.3}
\begin{enumerate}
\renewcommand{\theenumi}{\alph{enumi}}
\renewcommand{\labelenumi}{\theenumi)}
\item Let $X$ be an object of $E$. For $X$ to be quasi-compact, it is
  sufficient that $i^{\ast}(X)$ and $j^{\ast}(X)$ be so, and this condition
  is also necessary if $j:U\to e$ is quasi-compact.

\item Let $Y$ be an object of $E''$. For $Y$ to be quasi-compact, it is
  sufficient that $i_{\ast}(Y)$ be so, and this condition is also necessary
  if $U$ is quasi-compact.
\end{enumerate}
\end{enonce*}

\noindent
\underline{Proof}
\begin{enumerate}
\renewcommand{\theenumi}{\alph{enumi}}
\renewcommand{\labelenumi}{\theenumi)}
\item Sufficiency follows from \hyperref[VI.def1.1]{Definition 1.1} and the fact that the pair
  $(i^{\ast},j^{\ast})$ is conservative (IV \hyperlink{sga4.exp.IV.proposition.9.4.1.item.c}{9.11} 3)). Necessity follows from
  the fact that $i$ and $j$ are quasi-compact (by \hyperref[VI.prop4.2]{3.14} and the hypothesis
  that $j$ is quasi-compact).

\item Since $Y\simeq i^{\ast}i_{\ast}(Y)$, sufficiency follows from
  \hyperref[VI.prop4.2]{4.2}. Necessity follows from the sufficiency in a),
  taking into account that $i^{\ast}i_{\ast}(Y)\simeq Y$ and
  $j^{\ast}i_{\ast}(Y)\simeq U$.
\end{enumerate}

\begin{enonce*}{Corollary 4.4}\etiquette[4.4]{VI.cor4.4}
Let $Y$ be an object of $E''$. If $U$ is quasi-compact or if $E$ admits a
generating family consisting of quasi-compact objects, then for $Y$ to be
quasi-separated (\resp coherent), it is \emph{sufficient} that
$i_{\ast}(Y)$ be so. If $U$ is quasi-separated (\resp coherent) and if
$j:U\to e$ is a quasi-compact morphism, then for $Y$ to be quasi-separated
(\resp coherent), it is necessary that $i_{\ast}(Y)$ be so.
\end{enonce*}

The parenthetical\pageoriginale case follows from the non-parenthetical
case, taking \hyperref[VI.cor4.3]{4.3} b) into account. Suppose that
$i_{\ast}(Y)$ is quasi-separated, and let us prove that $Y$ is so under one
of the two stated hypotheses. Thus we must prove that, for two quasi-compact
objects $Y'$ and $Y''$ over $Y$, the fiber product
$\fprod{Y'}{Y''}{Y}$ is quasi-compact. When $U$ is assumed quasi-compact,
note that it is enough to prove that
$i_{\ast}(\fprod{Y'}{Y''}{Y})$ is quasi-compact
(\hyperref[VI.cor4.3]{4.3} b)); this follows from the fact that
$i_{\ast}(Y')$ and $i_{\ast}(Y'')$ are so
(\hyperref[VI.cor4.3]{4.3} b)), using here the quasi-compactness of $U$, the
fact that $i_{\ast}$ commutes with fiber products, and the hypothesis that
$i_{\ast}(Y)$ is quasi-separated. When $E$ is assumed to admit a generating
family consisting of quasi-compact objects, $i_{\ast}(Y')$ is a filtered
inductive limit of its quasi-compact subobjects ${X'}_{\alpha}$; hence
$Y'\simeq i^{\ast}i_{\ast}(Y')$ is a filtered inductive limit of the
subobjects $i^{\ast}({X'}_{\alpha})$, and since $Y'$ is quasi-compact, it is
equal to one of the $i^{\ast}({X'}_{\alpha})$. Likewise, $Y''$ is of the
form $i^{\ast}({X''}_{\beta})$, where ${X''}_{\beta}$ is a quasi-compact
subobject of $i_{\ast}(Y'')$. But then
$\fprod{Y'}{Y''}{Y}=
i^{\ast}(\fprod{X'_{\alpha}}{X''_{\beta}}{i_{\ast}Y})$, and since
$\fprod{X'_{\alpha}}{X''_{\beta}}{i_{\ast}(Y)}$ is quasi-compact by the
hypothesis that $i_{\ast}(Y)$ is quasi-separated, it follows by
\hyperref[VI.prop4.2]{4.2} that $\fprod{Y'}{Y''}{Y}$ is also quasi-compact.

Conversely, suppose that $U$ is quasi-separated and $j:U\to e$ is
quasi-compact, and assume that $Y$ is quasi-separated. Let us show that
$i_{\ast}(Y)$ is also quasi-separated, \ie that for two quasi-compact
objects $X',X''$ over $i_{\ast}(Y)$, the fiber product
$\fprod{X'}{X''}{i_{\ast}(Y)}$ is quasi-compact. By
\hyperref[VI.cor4.3]{4.3} a), it is enough to prove that its image under
$i^{\ast}$ is quasi-compact\pageoriginale (which follows from
\hyperref[VI.prop4.2]{4.2} and the hypothesis that $Y$ is quasi-separated)
\emph{and} that its image under $j^{\ast}$ is quasi-compact. The latter is
$j^{\ast}(X')\times j^{\ast}(X'')$, and is indeed quasi-compact because
$j^{\ast}(X')$ and $j^{\ast}(X'')$ are so ($j$ being quasi-compact) and $U$
is quasi-separated.

\begin{enonce*}{Corollary 4.5}\etiquette[4.5]{VI.cor4.5}
Suppose that $E$ admits a generating subcategory consisting of quasi-compact
(\emph{\resp} pre-Noetherian) objects. Then the same is true of $E''$.
\end{enonce*}

This follows from \hyperref[VI.prop4.2]{4.2} and the

\begin{enonce*}{Lemma 4.5.1}\etiquette[4.5.1]{VI.lem4.5.1}
If $(X_{\alpha})_{\alpha}$ is a generating family in $E$, then
$(i^{\ast}(X_{\alpha}))_{\alpha}$ is a generating family in $E''$.
\end{enonce*}

Indeed, this means that the family of functors
$$
Y \to \Hom (i^{\ast}(X_{\alpha}), Y)
$$
on $E''$ is conservative. But
$\Hom(i^{\ast}(X_{\alpha}),Y)\simeq\Hom(X_{\alpha},i_{\ast}(Y))$, and it is
enough to use the fact that the functor $i_{\ast}$ is conservative.

\begin{enonce*}{Proposition 4.6}\etiquette[4.6]{VI.prop4.6}
The notation is that of \hyperref[VI.subsec4.1]{4.1}.
\begin{enumerate}
\renewcommand{\theenumi}{\alph{enumi}}
\renewcommand{\labelenumi}{\theenumi)}
\item If $E$ admits a generating family consisting of coherent objects,
  then the same is true of the closed subtopos $E''$, and the inclusion
  morphism $i:E''\to E$ is coherent.

\item Suppose that the morphism $j:U\to e$ in $E$ is quasi-compact.
If\pageoriginale $E$ is locally coherent (\emph{\resp} coherent,
\emph{\resp} quasi-separated, \emph{\resp} locally Noetherian,
\emph{\resp} Noetherian, \emph{\resp} locally perfect, \emph{\resp}
perfect), then the same is true of $E''$, and the inclusion morphism
$i:E''\to E$ is coherent.
\end{enumerate}
\end{enonce*}

\begin{enumerate}
\renewcommand{\theenumi}{\alph{enumi}}
\renewcommand{\labelenumi}{\theenumi)}
\item The first assertion will follow from the second and
  \hyperref[VI.lem4.5.1]{4.5.1}. In any event,
  \hyperref[VI.cor4.5]{4.5} implies that $E''$ admits a generating family
  consisting of quasi-compact objects, so \hyperref[VI.prop3.2]{3.2}
  applies and shows that it is enough to verify that, for every coherent
  object $X$ of $E$, the object $i^{\ast}(X)$ of $E''$ is coherent. Using
  the compatibility of the formation of the topos complementary to an open
  with localization (IV \hyperref[IV.cor9.4.3]{9.13} b)), we reduce to the case in which $X$ is the
  final object of $E$, and hence to proving that the final object of $E''$
  is coherent. Now \hyperref[VI.cor4.4]{4.4} applies, so it is enough to
  prove that $i_{\ast}(e_{E''})=e_E$ is coherent, which is indeed the case.

\item Under each of the hypotheses made in b), $E$ admits a generating
  family consisting of coherent objects, so it already follows from a) that
  $i:E''\to E$ is coherent. Moreover, \hyperref[VI.lem4.5.1]{4.5.1} then
  implies that $E''$ admits as a generating subcategory the full subcategory
  ${E''}_{\coh}$ consisting of its coherent objects. Suppose that $E$ is
  coherent, and let us show that $E''$ is also coherent, \ie
  (\hyperref[VI.subsubsec2.4.5]{2.4.5}) that the subcategory
  ${E''}_{\coh}$ of $E''$ is stable under finite projective limits, knowing
  that this is so for the subcategory $E_{\coh}$ of $E$. This follows
  immediately from the parenthetical criterion
  \hyperref[VI.cor4.4]{4.4} (which applies because $U$ is now coherent, $e$
  being\pageoriginale coherent and $j:U\to e$ being quasi-compact), taking
  into account that $i_{\ast}$ is left exact. By localization, again using
  IV \hyperref[IV.cor9.4.3]{9.13} b), we conclude that if $E$ is locally coherent, then the same is
  true of $E''$. If $E$ is quasi-separated, then $E''$ is also
  quasi-separated: indeed, its final object is quasi-separated by
  \hyperref[VI.cor4.4]{4.4}, since that of $E$ is, and it remains to prove
  that the product of two quasi-separated objects of $E''$ is
  quasi-separated (knowing that this is so in $E$); this again follows from
  the criterion \hyperref[VI.cor4.4]{4.4} and the fact that $i_{\ast}$
  commutes with products (taking into account that the subobject $U$ of $e$
  is quasi-separated, since $e$ is, so that \hyperref[VI.cor4.4]{4.4}
  applies).
\end{enumerate}

Since a topos is locally Noetherian (\resp Noetherian) if and only if it is
locally coherent (\resp coherent) and admits a generating family consisting
of pre-Noetherian objects (\hyperref[VI.prop2.10]{2.10} (iii bis)), it
follows from the preceding discussion and \hyperref[VI.cor4.5]{4.5} that if
$E$ is locally Noetherian (\resp Noetherian), then the same is true of
$E''$. Now suppose that $E$ is perfect, and let us prove that $E''$ is
perfect. Since we already know that it is coherent, it remains to verify
that it satisfies criterion \hyperref[VI.cor1.25]{1.25} (i), \ie that every
object $Y$ of $E''$ is a filtered inductive limit of coherent objects,
knowing that the analogous statement is true in $E$. Since $i_{\ast}(Y)$ is
a filtered inductive limit of coherent objects $X_{\alpha}$,
$Y\simeq i^{\ast}i_{\ast}(Y)$ is a filtered inductive limit of the coherent
objects $i^{\ast}(X_{\alpha})$. By localization (IV \hyperref[IV.cor9.4.3]{9.13} b)), we conclude
that if $E$ is locally perfect, then the same is true of $E''$. This
completes the proof of \hyperref[VI.prop4.6]{4.6}.

\begin{enonce*}[remark]{Remark 4.6.1}\etiquette[4.6.1]{VI.rem4.6.1}
In fact,\pageoriginale in \hyperref[VI.prop4.6]{4.6} b), the hypothesis that
$j:U\to e$ is quasi-compact is unnecessary, except perhaps in the case
``$E$ quasi-separated''. Indeed, by the preceding proof it is enough to see
that if $E$ is coherent, then $E''$ is coherent. Now $U$ is the filtered
inductive limit of its quasi-compact subobjects $U_{\alpha}$, and in
\hyperref[VI.sec7]{7} we shall verify that $E''$ then identifies with the
projective-limit topos (in the sense of \hyperref[VI.sec7]{7}) of the closed
subtopoi $E''_i$ complementary to the $E_{/U_i}$, which by
\hyperref[VI.prop4.6]{4.6} b) are coherent topoi with coherent transition
morphisms. It then follows that $E''$ is a coherent topos
(\hyperref[VI.sec7]{7}).
\end{enonce*}

\begin{enonce*}{Corollary 4.7}\etiquette[4.7]{VI.cor4.7}
Suppose that the subcategory $E_{\coh}$ of $E$ consisting of coherent
objects is generating, and let $X$ be an object of $E$.
\begin{enumerate}
\renewcommand{\theenumi}{\alph{enumi}}
\renewcommand{\labelenumi}{\theenumi)}
\item If $X$ is quasi-separated, then $i^{\ast}(X)$ and $j^{\ast}(X)$ are
  quasi-separated; this condition is also sufficient if $j:U\to e$ is
  quasi-compact.

\item Suppose that $j:U\to e$ is quasi-compact. Then $X$ is coherent if and
  only if $i^{\ast}(X)$ and $j^{\ast}(X)$ are coherent.

\item Suppose that $E$ is coherent and $j:U\to e$ quasi-compact. Then $E$
  is perfect (\hyperref[VI.def2.9.1]{2.9.1}) if and only if $E'$ and $E''$
  are perfect.
\end{enumerate}
\end{enonce*}

\begin{enumerate}
\renewcommand{\theenumi}{\alph{enumi}}
\renewcommand{\labelenumi}{\theenumi)}
\item We noted in \hyperref[VI.subsec4.1]{4.1} that $j:E'\to E$ is
  quasi-separated, and we know from \hyperref[VI.prop4.6]{4.6} a) that the
  same is true of $i$, which gives necessity. For sufficiency, let $X'$ and
  $X''$ be quasi-compact objects over $X$. We must prove that
  $\fprod{X'}{X''}{X}$ is quasi-compact, and for this it\pageoriginale is
  enough to prove that its images under $i^{\ast}$ and $j^{\ast}$ are so
  (\hyperref[VI.cor4.3]{4.3} a)). Since $i$ and $j$ are quasi-compact, by
  \hyperref[VI.prop4.2]{4.2} and the hypothesis that $j:U\to e$ is
  quasi-compact, the conclusion follows from the fact that $i^{\ast}$ and
  $j^{\ast}$ commute with fiber products and from the hypothesis that
  $i^{\ast}(X)$ and $j^{\ast}(X)$ are quasi-separated.

\item This follows by combining a) with \hyperref[VI.cor4.3]{4.3} a).

\item Necessity was already seen (\hyperref[VI.prop4.6]{4.6} b)).
  Sufficiency follows from the fact that, since $E$ is coherent, $E'$ and
  $E''$ are coherent (\hyperref[VI.prop4.6]{4.6} b)), so that for any of the
  topoi under consideration, being perfect means that the category of its
  coherent objects is stable under finite $\varinjlim$. We therefore
  conclude by b).
\end{enumerate}

\begin{enonce*}[remark]{Exercise 4.8}\etiquette[4.8]{VI.exer4.8}
\begin{enumerate}
\renewcommand{\theenumi}{\alph{enumi}}
\renewcommand{\labelenumi}{\theenumi)}
\item Solve the following question (to which the writer confesses, to his
  embarrassment, that he does not know the answer): with the notation of
  \hyperref[VI.subsec4.1]{4.1}, if $E$ is quasi-separated
  (\resp algebraic), is the same true of $E''$?

\item Suppose that $E$ is equivalent to a topos of the form $\widehat{C}$,
  where $C$ is a category equivalent to a category $\in 𝒰$. Prove directly,
  using \hyperref[VI.exer2.17]{2.17} c) and IV 9.24, that if $E$ is locally
  coherent (\resp coherent, \resp algebraic, \resp quasi-separated, \resp
  locally Noetherian, \resp Noetherian), then the same is true of $E''$.
\end{enumerate}
\end{enonce*}

\setcounter{subsection}{8}
\subsection{}\etiquette[4.9]{VI.subsec4.9}
Keep the notation of \hyperref[VI.subsec4.1]{4.1}, and recall (IV \hyperref[IV.theo9.5.4]{9.16}) that
the datum of a situation $(E,U)$ consisting of a topos $E$ and an open $U$
of $E$ is essentially equivalent to that of a triple $(E',E'',f)$, where
$E'$ and $E''$ are\pageoriginale topoi and
\begin{equation*}
f : E'\to E''\tag{4.9.1}
\end{equation*}\etiquette[4.9.1]{VI.eq4.9.1}
is a ``gluing functor'', \ie a functor that is \emph{left exact and
accessible}; starting from $(E,U)$ as in \hyperref[VI.subsec4.1]{4.1}, the
associated gluing functor is given by
\begin{equation*}
  f = i^{\ast} j_{\ast}{\quoi}.\tag{4.9.2}
\end{equation*}\etiquette[4.9.2]{VI.eq4.9.2}

We propose to express, in terms of the data $E'$, $E''$, and $f$, the fact
that $E$ is a coherent (\resp perfect) topos and that $U$ is a quasi-compact
object, or equivalently, that $E$ and $E'$ are coherent. We already know
that this implies that $E''$ is also coherent
(\hyperref[VI.prop4.6]{4.6} b)), so the question becomes the following:
given two \emph{coherent} topoi $E'$ and $E''$ and a gluing functor $f$
\eqhyperref[VI.eq4.9.1]{4.9.1}, under what conditions on $f$ is the glued
topos $E$ coherent (\resp perfect)? We also know
(\hyperref[VI.cor4.7]{4.7} c)) that $E$ is perfect if and only if $E$ is
coherent and $E'$ and $E''$ are perfect; thus the parenthetical case of the
problem reduces to the non-parenthetical case. Let us nevertheless note that
the solution (\hyperref[VI.prop4.10]{4.10}) of the latter is more elegant if
$E'$ is already assumed perfect.

\setcounter{subsection}{9}
\setcounter{subsubsection}{2}
\subsubsection{}\etiquette[4.9.3]{VI.subsubsec4.9.3}

First note that if $E$ is coherent, then the full subcategory $E_{\coh}$ of
$E$ consisting of the objects $X=(X',X'',u:X''\to f(X'))$ of $E$ that are
coherent is reconstructed as the subcategory $E_0$ of $E$ consisting of the
$X$ for\pageoriginale which $X'=j^{\ast}(X)$ and $X''=i^{\ast}(X)$ are
coherent (\hyperref[VI.cor4.7]{4.7} b)). A necessary condition for $E$ to be
coherent is therefore that the preceding subcategory $E_0$ be generating.
This condition is also sufficient, because by \hyperref[VI.cor4.3]{4.3} a),
$E_0$ consists of quasi-compact objects; moreover, it is clear that $E_0$ is
stable under finite projective limits, and we conclude by \hyperref[VI.subsubsec2.4.5]{3.4.5}.

Recall now (IV \hyperref[IV.subsec9.6]{9.18}) that the datum of a gluing functor
\eqhyperref[VI.eq4.9.1]{4.9.1} is equivalent to that of a \emph{sheaf} on
$E''$ with values in $\Pro(E')^0$:
\begin{equation*}
G \in \Ffaisc (E'', \Pro(E')^{0}), G : {E''}^{0} \to \Pro({E'}^{0})
\tag{4.9.4}{\quoi},
\end{equation*}\etiquette[4.9.4]{VI.eq4.9.4}
or equivalently to that of a functor $g=G^0$,
\begin{equation*}
g:E''\to \Pro (E')\tag{4.9.5}
\end{equation*}\etiquette[4.9.5]{VI.eq4.9.5}
that \emph{commutes with inductive limits}. If $C$ is a generating
subcategory of $E''$, endowed with the induced topology, we know
(II \hyperref[VI.subsec6.10]{6.10}) that the datum of $G$ is also equivalent
(up to unique isomorphism) to that of a sheaf on $C$ with values in
$\Pro(E')^0$
\begin{equation*}
G_C \in \Ffaisc (C, \Pro(E')^0), G_C : C^{\circ} \to \Pro (E')^{0}{\quoi},\tag{4.9.6}
\end{equation*}\etiquette[4.9.6]{VI.eq4.9.6}
or, equivalently, to that of a functor $g_C=G_C^0$,
\begin{equation*}
g_C : C \to \Pro (E'){\quoi},\tag{4.9.7}
\end{equation*}\etiquette[4.9.7]{VI.eq4.9.7}
satisfying\pageoriginale the familiar left-exactness conditions
(II \hyperlink{sga4.exp.II.remark.6.2.item.1}{6.2.1})). Of course, $g_C$ is simply the \emph{restriction of $g$
$\hyperref[VI.eq4.9.5]{(4.9.5)}$ to $C$}. The case of greatest interest to us is that in which we
take $C=E''_{\coh}$, giving the objects
\hypertarget{sga4.exp.VI.equation.4.9.8}{}\begin{align*}
& G_0 \in \Ffaisc (E''_{\coh}, \Pro (E')^{0}), G_0 : {E''}_{\coh}^0 \to
 \Pro (E')^0 {\quoi},\tag{4.9.8}\\
& g_0 = G^0_0 : E''_{\coh} \to \Pro(E'){\quoi},\tag{4.9.9}
\end{align*}\hypertarget{sga4.exp.VI.equation.4.9.9}{}
each of which is again equivalent to the datum of $f$.

\setcounter{subsubsection}{9}
\subsubsection{}\etiquette[4.9.10]{VI.subsubsec4.9.10}

Suppose that the full generating subcategory $C$ of $E''$ consists of
quasi-compact objects and is stable in $E''$ under finite sums, passage to
the quotient by equivalence relations, and fiber products; this is the case,
for example, for $C=E''_{\coh}$
(\hyperref[VI.cor1.15]{1.15} and \hyperref[VI.cor1.17.1]{1.17.1}). It is then
immediate, if $P$ is a category in which finite projective limits are
representable (for example $P=\Pro(E')^0$), that a functor
$G_C:C^{\circ}\to P$ is a sheaf with values in $P$ if and only if the functor
$g_C=G_C^0:C\to P^0$ commutes with finite sums and with passage to the
quotient by an equivalence relation. This specifies, in particular, which
objects are the sheaves \eqhyperref[VI.eq4.9.6]{4.9.6} (and thus expresses
the gluing functors $f:E'\to E''$).

With these reminders in place, we can give the solution to the problem posed
in \hyperref[VI.subsec4.9]{4.9}:

\begin{enonce*}{Proposition 4.10}\etiquette[4.10]{VI.prop4.10}
Let $E'$ and $E''$\pageoriginale be two coherent topoi, and let
$f:E'\to E''$ be a gluing functor \emph{(IV \hyperref[IV.eq9.3.6.3]{9.10})}, \ie a left exact and
accessible functor. Let $E$ be the topos deduced from it by gluing
\emph{(IV \hyperref[IV.theo9.5.4]{9.16})}, and denote by $E'_{\coh}$ (\resp $E'_{PF}$) the strictly
full subcategory of $E'$ consisting of the coherent objects (\resp the
objects $X$ such that the covariant functor $\Hom(X,-)$ represented by $X$
commutes with small filtered inductive limits). Consider the following
conditions:
\begin{itemize}
\item[(i)] $E$ is coherent.

\item[(ii)] For every object $X'$ of $E'$, there exists a family of
  morphisms
$$
v_{\alpha} : Y'_{\alpha} \to X'
$$
with target $X'$ and with coherent objects of $E'$ as sources, such that the
family of the
$$
f(v_{\alpha}) : f(Y'_{\alpha}) \to f(X')
$$
in $E''$ is covering.

\item[(ii bis)] $f$ commutes with (small) filtered inductive limits, and
  \emph{(ii)} holds for every $X'\in\Ob E'_{PF}$.

\item[(ii ter)] The canonical functor \emph{(IV \hyperref[IV.eq9.7.2.3]{9.20.1})}
\begin{equation*}
\pi : \Ppoint(E'') \to \Pro(E')\tag{4.10.1}
\end{equation*}\etiquette[4.10.1]{VI.eq4.10.1}
defined by the gluing functor $f$ factors (up to isomorphism) through
$\Pro(E'_{\coh})$ (via the fully faithful functor
$\Pro(E'_{\coh})\to\Pro(E')$ arising from the inclusion
$E'_{\coh}\to E'$):
\begin{equation*}
\pi_0 : \Ppoint (E'') \to \Pro(E'_{\coh}){\quoi}.\tag{4.10.2}
\end{equation*}\etiquette[4.10.2]{VI.eq4.10.2}

\item[(iii)] The\pageoriginale functor $f$ commutes with (small) filtered
  inductive limits.

\item[(iii bis)] The functor $g_0$ \hyperlink{sga4.exp.VI.equation.4.9.9}{(4.9.9)} factors (up to isomorphism)
  through a functor
\begin{equation*}
E''_{\coh} \to \Pro(E'_{PF}){\quoi},\tag{4.10.3}
\end{equation*}\etiquette[4.10.3]{VI.eq4.10.3}
via the fully faithful functor $\Pro(E'_{PF})\to\Pro(E')$ deduced from the
inclusion $E'_{PF}\to E'$.

\item[(iii ter)] The canonical functor $\pi$
  \eqhyperref[VI.eq4.10.1]{4.10.1} factors (up to isomorphism) as
\begin{equation*}
\pi_1 : \Ppoint(E'') \to \Pro(E'_{PF}){\quoi}.\tag{4.10.4}
\end{equation*}\etiquette[4.10.4]{VI.eq4.10.4}

\item[(iv)] The functor $g_0$ \hyperlink{sga4.exp.VI.equation.4.9.9}{(4.9.9)} factors (up to isomorphism) through a
  functor
\begin{equation*}
E''_{\coh} \to \Pro(E'_{\coh}){\quoi}.\tag{4.10.5}
\end{equation*}\etiquette[4.10.5]{VI.eq4.10.5}
We then have the following diagram of implications:
\begin{equation*}
\text{\rm (iv) $\Rightarrow$ (i) $\Leftrightarrow$ (ii) $\Leftrightarrow$ (ii bis)
  $\Leftrightarrow$ (ii ter) $\Rightarrow$ (iii) $\Leftrightarrow$ (iii bis)
  $\Leftrightarrow$ (iii ter)}{\quoi}.\tag{4.10.6}
\end{equation*}\etiquette[4.10.6]{VI.eq4.10.6}
\end{itemize}
\end{enonce*}

Let us immediately record the

\begin{enonce*}{Corollary 4.11}\etiquette[4.11]{VI.cor4.11}
\begin{enumerate}
\renewcommand{\theenumi}{\alph{enumi}}
\renewcommand{\labelenumi}{\theenumi)}
\item Suppose that $E'$ is perfect. Then all the conditions considered in
  \hyperref[VI.prop4.10]{4.10} are equivalent; in particular, $E$ is
  coherent if and only if $f$ commutes with filtered inductive limits, or
  equivalently if and only if the functor $g_0$ \hyperlink{sga4.exp.VI.equation.4.9.9}{(4.9.9)} factors (up to
  isomorphism)\pageoriginale through $\Pro(E'_{\coh})$.

\item For $E$ to be perfect, it is necessary and sufficient that $E'$ and
  $E''$ be perfect and that $f$ (\resp $g_0$) satisfy the condition stated
  in {\rm a)}.
\end{enumerate}
\end{enonce*}

Indeed, a) follows from the fact that $E'$ perfect means
$E'_{\coh}=E'_{PF}$, so that (iii bis) implies (iv). Assertion b) follows,
as is clear from the preliminary remarks of \hyperref[VI.subsec4.9]{4.9}.


\setcounter{subsection}{11}

% END INCLUDED SGA4_English_current_buildable_source_20260906/source/06/drafts/SGA4_VI_LINES_2419_2936_English.texfrag


% BEGIN INCLUDED SGA4_English_current_buildable_source_20260906/source/06/drafts/SGA4_VI_LINES_2937_3239_English.texfrag
\subsection{Proof of 4.10}\etiquette[4.12]{VI.subsec4.12}

\begin{enumerate}
\renewcommand{\theenumi}{\alph{enumi}}
\renewcommand{\labelenumi}{\theenumi)}
\item Let us spell out the condition equivalent to (i) obtained in
  \hyperref[VI.subsubsec4.9.3]{4.9.3}, namely that the full subcategory
  $E_0$ of $E$ is generating, \ie that for every object
  $X = (X', X'', u : X''\to f(X'))$ of $E$, the family of all
  morphisms
$$
v = (v', v'') : Y = (Y', Y'', v : Y''\to f(Y')) \to X {\quoi},
$$
with $Y \in \Ob E_0$, \ie with $Y'$ and $Y''$ coherent, is
epimorphic. Since the pair of functors $(i^{\ast}, j^{\ast})$ is
conservative, this also means that the family of corresponding
morphisms $Y' \to X'$ is epimorphic in $E'$, and that the family of
corresponding morphisms $Y''\to X''$ is epimorphic in $E''$. This is
clear for the first family, as is seen by taking $Y \in \Ob E_0$
\emph{over} $U$, \ie such that $Y'' = \varnothing_{E''}$: one obtains
the family of all arrows in $E'$ with target $X'$ and coherent source,
which is indeed epimorphic because $E'$ is coherent and hence the
subcategory ${E'}_{\coh}$ is generating. It therefore remains to
express the condition that the family $Y''\to X''$ be epimorphic for
every given object $X$ of $E$.

\item Now suppose\pageoriginale that condition (ii) is satisfied for
  the object $X'$ under consideration. It follows that one can find
  an epimorphic family of morphisms $Y''_{\beta} \to X''$ whose
  composites with $u : X''\to f(X')$ each lift to a morphism
  $Y''_{\beta} \to f({Y'_{\alpha}})$ for a suitable
  $\alpha = \phi(\beta)$. Since $E''_{\coh}$ is a generating
  subcategory of $E''$, we may moreover suppose the $Y''_{\beta}$
  coherent. The commutative diagrams
\[
\xymatrix@C=2.25cm@R=1.25cm{
Y''_{\beta}  \ar[r] \ar[d]& f(Y'_{\varphi(\beta)}) \ar[d]\\
X'' \ar[r] & f(X')
}
\]
then provide the required family of morphisms in $E$, giving an
epimorphic family $Y''_{\beta} \to X''$. Thus (ii) $\Rightarrow$ (i).

Conversely, suppose (i), and let us prove (ii). Apply the condition
made explicit in a) to the object
$$
X = j_{\ast}(X') = (X', f(X'), \id = f(X') \xrightarrow{\sim} f(X')){\quoi}.
$$
Since the $Y''\to X''\simeq f(X')$ form an epimorphic family, a
fortiori so does the family of $f(v') : f(Y') \to f(X')$. The
conclusion follows, since the $Y'$ are coherent by hypothesis. Thus
$$
{\rm(i)} \Leftrightarrow {\rm(ii)}{\quoi}.
$$

\item Note\pageoriginale now that condition (iv) also means that for
  every object $X''$ of $E''$, the pro-representable functor
$$
g(X'') : X'\longmapsto \Hom (X'', f(X'))
$$
on $E'$ is pro-represented by a pro-object whose components belong to
${E'}_{\coh}$. Equivalently, in the category $E'_{/g(X'')}$ formed by
pairs consisting of an object $X'$ of $E'$ and an element
$u \in g(X'')(X') =\Hom(X'', f(X'))$, $u : X''\to f(X')$, there is a
cofinal set of objects $(X'_{\alpha},u_{\alpha})$, with
$u_{\alpha} : X'' \to f(X'_{\alpha})$ and
$X'_{\alpha} \in \Ob E'_{\coh}$. Since the category
$E'_{/g(X'')}$ is filtered, the cofinality condition simply means
that for every object $(X',u)$, $u : X''\to f(X')$, of
$E'_{/g(X'')}$, there is a morphism from one of the
$(X'_{\alpha},u_{\alpha})$ to this object, \ie a morphism
$v'_{\alpha} : X'_{\alpha} \to X'$ making the triangle
\[\vcenter{
\xymatrix@C=2cm@R=1.5cm{
                      &f(X'_{\alpha}) \ar[d]^{f(v'_\alpha)}\\
X''\ar@{-}[ru]^{u_{\alpha}} \ar[r]& f(X').
}}\leqno{(*)}
\]
commutative. Thus condition (iv) is equivalent to the following:

\begin{itemize}
\item[(iv bis)] For every arrow $u : X''\to f(X')$, with
  $X'\in \Ob E'$ and $X''\in \Ob E''_{\coh}$, there is an arrow
  $v' : Y'\to X'$ in $E'$, with $Y'$ coherent, such that $u$ factors
  through $f(v') : f(Y') \to f(X')$.

It is clear that condition (iv bis) implies (ii), since for fixed
$X'$ the family of $u : X''\to f(X')$ whose source is coherent is
epimorphic. A fortiori, the same is therefore true of the family of
all $f(v'_{\alpha})$ in the corresponding diagrams $(\ast)$. Thus
$$
{\rm(iv)} \Rightarrow {\rm(i)}{\quoi}.
$$
\end{itemize}

\item Let us spell out condition (ii ter). Let $p$ be a point of
  $E''$. By definition, $\pi(p)$ is the pro-object of $E'$ which
  pro-represents the functor $h : X'\longmapsto f(X')_{p}$, or,
  equivalently, the pro-object defined by the canonical functor
  $E'_{/h} \to E'$, where $E'_{/h}$ denotes the category of pairs
  $(X',u)$ with $X'\in \Ob E'$ and $u\in h(X')$, \ie
  $u\in f(X')_p$. To say that this pro-object is isomorphic to a
  pro-object coming from $\Pro(E'_{\coh})$ means that the subcategory
  of $E'_{/h}$ formed by the $(X',u)$ with
  $X'\in\Ob E'_{\coh}$ is cofinal. In other words, for every object
  $(X',u)$, $u\in f(X')_p$, of $E'_{/h}$, there is an object
  $(X'_0,u_0)$ mapping to it, with $X'_0$ coherent; that is, there is
  a morphism $X'_0\to X'$ ($X'_0$ coherent) such that
  $u\in\text{Im}(f(X'_0)\to f(X'))$. To say that this holds for every
  point $p$ means that, for every object $X'$ of $E'$, the family of
  maps $f(X'_0)\to f(X')$, with $X'_0$ coherent, is surjective on
  every stalk. As we shall see in the appendix (10) that a coherent
  topos $E''$ has enough points, the resulting condition also says
  that, for every $X'\in\Ob E'$, the family
  $f(X'_0)\to f(X')$ is epimorphic. This is precisely condition (ii),
  and hence
$$
{\rm(ii)} \Leftrightarrow \text{(ii ter)}{\quoi}.
$$

\item The\pageoriginale condition that $f$ commute with filtered
  inductive limits (condition (iii)) is equivalent, since the family
  of fiber functors of $E''$ is conservative, to the condition that
  the composite functors $h : X'\longmapsto f(X')_p$ do so. Let us
  prove that this means that the pro-object $\pi(p)$ which
  pro-represents this functor lies in the essential image of
  $\Pro(E'_{PF})$. This will prove the implications
$$
{\rm(iii)} \Leftrightarrow \text{(iii ter)} \Leftarrow \text{(ii ter)}
\Rightarrow \text{(iii bis)}{\quoi}.
$$
We are therefore led to prove the following lemma.

\begin{enonce*}{Lemma 4.12.1}\etiquette[4.12.1]{VI.lem4.12.1}
Let $E'$ be a topos in which the subcategory $E'_{\coh}$ is
generating, and let $X' \in \Ob \Pro(E')$. Then $X'$ belongs to the
essential image of $\Pro(E'_{PF})$ if and only if the functor
$h : E'\to (\Ens)$ which it pro-represents commutes with filtered
inductive limits.
\end{enonce*}

Since a filtered inductive limit of functors commuting with filtered
inductive limits has the same property, sufficiency follows from the
definition of $E'_{PF}$. For necessity, we must prove that if the
functor $h$ commutes with filtered inductive limits, then in the
filtered category $E'_{/h}$ of pairs $(X',u)$ with $u\in h(X')$, the
subcategory formed by the pairs $(X',u)$ with
$X'\in\Ob E'_{PF}$ is cofinal. That is, for every $(X',u)$ in
$E'_{/h}$, there is a morphism $X'_0\to X'$, with
$X'_0\in\Ob E'_{PF}$, such that
$u\in\text{Im}(h(X'_0)\to h(X'))$.\pageoriginale But $X'$ is known to
be a filtered inductive limit of objects $X'_i$ of $E'_{PF}$
(\hyperref[VI.cor1.24.2]{1.24.2} a)), so $h(X')$ is the filtered
inductive limit of the $h(X'_i)$, which proves the assertion.

\item It remains only to prove the implications
$$
\text{(ii bis)} \Rightarrow \text{(ii)\quad and\quad (iii)}
\Leftrightarrow \text{(iii bis)}{\quoi}.
$$
The first implication follows trivially from the fact that every
object $X'$ of $E'$ is a filtered inductive limit of objects $X'_i$
of $E'_{PF}$. For the second, it is enough to note that the family of
functors
$$
Y'' \longmapsto \Hom(X'', Y'') {\quoi},
$$
for $X''\in\Ob {E''}_{\coh}$, is conservative and consists of
functors commuting with filtered inductive limits
(\hyperref[VI.thm1.23]{1.23} (ii)). Consequently, the functor
$f : E'\to E''$ commutes with filtered inductive limits if and only
if the composite functors
$X'\longmapsto \Hom(X'',f(X'))=
\Hom_{\Pro(E')}(g_0(X''),X')$ do so. By
\hyperref[VI.lem4.12.1]{4.12.1}, this is equivalent to the fact that
the $g_0(X'')$ lie in the essential image of $\Pro(E'_{PF})$. \cqfd
\end{enumerate}

\begin{enonce*}[remark]{Remark 4.13}\etiquette[4.13]{VI.rem4.13}
It is useful to state in sheaf-theoretic language the meaning of
condition \hyperref[VI.prop4.10]{4.10} (iv). Consider the commutative
diagram of functors
\[\vcenter{
\xymatrix@C=2.5cm@R=1.5cm{
\Pro(E'_{\coh})^{\circ}  \ar@{^{(}->}[r]^{\alpha}  \ar[d]^{\beta_0}&
\Pro(E')^{\circ}_{\beta} \ar[d]^{\beta}\\
\text{$\Ev{\coh}$} & \text{$\Ev{}$}\ar[l]^{\gamma} ,
}}\leqno{(4.13.1)}
\]\etiquette[4.13.1]{VI.eq4.13.1}
where\pageoriginale the superscript circles denote opposite
categories, and the second horizontal arrow $\gamma$ is restriction
of functors, all the other arrows denoting the familiar fully
faithful functors. Note that the inclusion arrow $\alpha$ \emph{does
not in general commute with projective limits}, even finite ones.
Thus, in general, a sheaf $F$ on a site $C$ (such as $E''$) defines a
presheaf $\alpha\circ F$ on the same site, with values in
$\Pro(E'')^{0}$, which need not be a sheaf. Nevertheless, if a
\emph{presheaf} $F$ on $C$ with values in
$\Pro(E'_{\coh})^{0}$ is such that the corresponding presheaf
$\alpha F$ with values in $\Pro(E')^{0}$ is a sheaf, then $F$ itself
is a sheaf. Indeed, the sheaf condition is a left-exactness property
(II \hyperlink{sga4.exp.II.remark.6.2.item.1}{6.2.1})), and the functors $\beta$, $\beta_0$, and $\gamma$ commute
with small projective limits.
\end{enonce*}

On the other hand, the characterization in
\hyperref[VI.subsubsec4.9.10]{4.9.10} of sheaves on $E''_{\coh}$ with
values in an arbitrary category $P$ shows that, for every functor
$\alpha:P\to Q$ such that $\alpha^{0}:P^0\to Q^0$ commutes with
\emph{finite coproducts} and with \pageoriginale
\emph{passage to quotients by equivalence relations}, the functor
$g\mapsto\alpha\circ g$ from $\Hom({E''}^0_{\coh},P)$ to
$\Hom({E''}^0_{\coh},Q)$ takes sheaves to sheaves. We also know that
in $\Pro(E'_{\coh})$ and in $\Pro(E')$, finite coproducts and
quotients by equivalence relations are representable, and that the
inclusion $E'_{\coh}\to E'$ commutes with them
(\hyperref[VI.cor1.15]{1.15} and
\hyperref[VI.cor1.17.1]{1.17.1}). The same is therefore true for the
categories $\Pro(E'_{\coh})$ and $\Pro(E')$ and for the inclusion
between them (I \hyperref[I.prop8.9.5]{8.9.5} b) and \hyperref[I.prop8.9.7]{8.9.7}), which is precisely $\alpha^0$,
where $\alpha$ is the functor occurring in
\eqhyperref[VI.eq4.13.1]{4.13.1}. It follows that if
\begin{equation*}
\varphi:{E''}^0_{\coh}\to\Pro(E'_{\coh})^0\tag{4.13.2}
\end{equation*}\etiquette[4.13.2]{VI.eq4.13.2}
is a functor, then $\varphi$ \emph{is a sheaf on $E''_{\coh}$ with
values in $\Pro(E'_{\coh})^0$ if and only if $\alpha\circ\varphi$ is
a sheaf on $E''_{\coh}$ with values in $\Pro(E')^0$}. Thus the
category of gluing functors $f:E'\to E''$ satisfying condition
\hyperref[VI.prop4.10]{4.10} (iv) is canonically equivalent to the
category of sheaves $\varphi$ on $E''_{\coh}$ with values in
$\Pro(E'_{\coh})^0$. Taking \hyperref[VI.prop4.10]{4.10} into account,
we conclude that such a sheaf $\varphi$ canonically defines a gluing
functor (up to unique isomorphism), and that the topos $E$ obtained
from the latter is \emph{coherent}. Moreover
(\hyperref[VI.cor4.11]{4.11}), \emph{if $E'$ is perfect}, one thereby
obtains \emph{all} coherent topoi that can be obtained from $E'$ and
$E''$ by gluing.

Note\pageoriginale also that $E'_{\coh}$ is equivalent to a small
category and is stable under finite projective limits. Consequently
(I \hyperref[I.sec8.10.14]{8.10.14}), by the fully faithful functor $\beta_0$ of
\eqhyperref[VI.eq4.13.1]{4.13.1}, $\Pro(E'_{\coh})^0$ is equivalent
to the full subcategory of
$\widehat{E'_{\coh}}=\Hom(E'_{\coh},(\Ens))$ formed by the left exact
functors. Thus the data of a sheaf
\eqhyperref[VI.eq4.13.2]{4.13.2} are equivalent to those of a functor
\begin{equation*}
F:{E''}_{\coh}^0\times E'_{\coh}\to(\Ens)\tag{4.13.3}
\end{equation*}\etiquette[4.13.3]{VI.eq4.13.3}
which is a sheaf in the first argument and left exact in the second;
or, equivalently, to those of a functor
\begin{equation*}
f_0:E'_{\coh}\to E''\tag{4.13.4}
\end{equation*}\etiquette[4.13.4]{VI.eq4.13.4}
which is left exact. If $f$ is the gluing functor satisfying condition
\hyperref[VI.prop4.10]{4.10} (iv) associated with $\varphi$, one
immediately verifies that $f_0$ is canonically isomorphic to the
restriction of $f:E'\to E''$. Thus \emph{under the conditions being
considered, $f_0$ determines $f$} (up to canonical isomorphism). To
summarize the content of \hyperref[VI.cor4.11]{4.11}, we finally state:

\begin{enonce*}{Corollary 4.14}\etiquette[4.14]{VI.cor4.14}
When $E'$ is a perfect topos, the category of gluing functors
$f:E'\to E''$ which give a coherent topos $E$ by gluing is equivalent,
under $f\mapsto f_0=f|E'_{\coh}$, to the category of left exact
functors $f_0:E'_{\coh}\to E''$, or equivalently to the\pageoriginale
category of sheaves $\varphi$ on the site $E''_{\coh}$ (or, what
amounts to the same thing, on the topos $E''$) with values in the
category $\Pro(E'_{\coh})^0$ \emph{(\cf
\hyperref[VI.subsubsec4.9.10]{4.9.10})}.
\end{enonce*}

The first assertion can also be deduced from
\hyperref[VI.prop4.10]{4.10} by observing that
$E'\approxeq\Ind(E'_{\coh})$ (\hyperref[VI.cor1.25]{1.25}). Indeed,
restriction gives an equivalence between the category of functors
$f:E'\to P$ commuting with filtered inductive limits (where $P$ is a
category in which small filtered inductive limits are representable)
and the category of arbitrary functors $f:E'_{\coh}\to P$. It is
immediate that, under these conditions, if finite projective limits
are representable in $P$, then $f$ is left exact if and only if $f_0$
is.

\begin{enonce*}[remark]{Remarks 4.15}\etiquette[4.15]{VI.rem4.15}
\begin{enumerate}
\renewcommand{\theenumi}{\alph{enumi}}
\renewcommand{\labelenumi}{\theenumi)}
\item The proof given for \hyperref[VI.prop4.10]{4.10} shows that
  conditions equivalent to (ii ter) \resp (iii ter) are obtained by
  replacing the category $\Ppoint(E'')$ by a subcategory which
  defines a \emph{conservative} family of fiber functors.

\item Taking for $E''$ the point topos (IV \hyperref[IV.subsec2.2]{2.2}), one sees immediately
  that conditions (ii ter) and (iii ter) are equivalent \emph{only if}
  $E'_{\coh}=E'_{PF}$, \ie if $E'$ is perfect. Thus the condition
  ``$E$ is coherent'' is equivalent to the condition ``$f$ commutes
  with filtered inductive limits'' only if $E'$ is perfect. On the
  other hand, taking more generally for $E''$ a coherent topos of the
  form $\widehat{C}$, one sees that \emph{in this case condition} (i)
  \emph{is equivalent to} (iv) \emph{for every $E'$}. In fact, in the
  general case the editor has been unable to construct an example in
  which $E$ is coherent without condition (iv) being satisfied. This
  is shameful!
\end{enumerate}
\end{enonce*}


% END INCLUDED SGA4_English_current_buildable_source_20260906/source/06/drafts/SGA4_VI_LINES_2937_3239_English.texfrag


% BEGIN INCLUDED SGA4_English_current_buildable_source_20260906/source/06/drafts/SGA4_VI_LINES_3240_3413_English.texfrag
\section{Commutation of the functors $H^i(X,-)$ with filtered
  inductive limits}\etiquette[5]{VI.sec5}\pageoriginale

\begin{enonce*}{Theorem 5.1}\etiquette[5.1]{VI.thm5.1}
Let $E'$ be an algebraic topos (\hyperref[VI.def2.3]{2.3}), let $E$ be
a locally coherent topos (\hyperref[VI.def2.3]{2.3}), and let
$f:E'\to E$ be a coherent morphism of topoi
(\hyperref[VI.def3.1]{3.1}). Then, for every integer $q$, the functors
$R^qf_{\ast}$ \emph{(V \hyperref[V.subsec5.0]{5.0})} commute with filtered inductive limits
of abelian sheaves.
\end{enonce*}

\begin{enonce*}{Corollary 5.2}\etiquette[5.2]{VI.cor5.2}
Let $E'$ be a coherent topos (\hyperref[VI.def2.3]{2.3}). For every
integer $q$, the functor $H^q(E',-)$ commutes with filtered inductive
limits of abelian sheaves.
\end{enonce*}

\begin{enonce*}{Corollary 5.3}\etiquette[5.3]{VI.cor5.3}
Let $E$ be a topos (\resp an algebraic topos)
(\hyperref[VI.def2.3]{2.3}), and let $X$ be an algebraic and coherent
(\resp coherent) object of $E$ (\hyperref[VI.def2.3]{2.3}). For every
integer $q$, the functor $H^q(X,-)$ commutes with filtered inductive
limits of abelian sheaves.
\end{enonce*}

Corollary \hyperref[VI.cor5.2]{5.2} follows from
\hyperref[VI.thm5.1]{5.1} by taking for $E$ the point topos and for
$f:E'\to E$ the unique morphism (IV \hyperref[IV.subsec2.2]{2.2}). Let us show that
\hyperref[VI.cor5.2]{5.2} implies \hyperref[VI.cor5.3]{5.3}. Let
$j_X:E_{/X}\to E$ be the localization morphism. There is a canonical
isomorphism, functorial in the abelian sheaf $F$ (V \hyperref[V.2.2.1]{2.2.1}),
$H^q(X,F)\simeq H^q(E_{/X},j_X^{\ast}F)$. The functor $j_X^{\ast}$
commutes with inductive limits, and the topos $E_{/X}$ is coherent
(\hyperref[VI.subsubsec2.4.6]{2.4.6}), whence
\hyperref[VI.cor5.3]{5.3}. Let us show that
\hyperref[VI.cor5.3]{5.3} implies \hyperref[VI.thm5.1]{5.1}. Let
$E_{\text{cohalg}}$ be the full subcategory of $E$ defined by the
algebraic and coherent objects of $E$. It is a generating category
(\hyperref[VI.subsubsec2.4.5]{2.4.5}). For every abelian sheaf $F$,
$R^qf_{\ast}F$ is the sheaf associated with the presheaf
$X\mapsto H^q(f^{\ast}X,F)$ $(X\in\ob C)$
(\hyperrefext[V][V.prop5.1]{5.1}). Since $f$ is coherent,
$f^{\ast}X$ is coherent for every object $X$ of $C$
(\hyperref[VI.def3.1]{3.1}). By \hyperref[VI.cor5.3]{5.3}, the
presheaf $X\mapsto H^q(f^{\ast}X,F)$ commutes with filtered inductive
limits in the argument $F$, and \hyperref[VI.thm5.1]{5.1} follows,
since the associated-sheaf functor commutes with inductive limits.
(Only the following property of the morphism $f$ is used: there is a
generating family $S$ of $E$ such that, for every $X$ of $S$,
$f^{\ast}(X)$ is algebraic and coherent.) It therefore remains to
prove \hyperref[VI.cor5.2]{5.2}. We already know that the functor
$H^{\circ}(E',-)$ commutes with filtered inductive limits of abelian
sheaves \hyperref[VI.thm1.23]{(1.2.3)}. By a classical result of \cite{sga4.exp.VI.bib.7}, to prove
\hyperref[VI.cor5.2]{5.2} it is enough to show that a filtered
inductive limit of sheaves acyclic for the functor $H^{\circ}(E',-)$
is acyclic for $H^{\circ}(E',-)$. This follows from the next lemma,
applied with $C=E'_{\coh}$.

\begin{enonce*}{Lemma 5.4}\etiquette[5.4]{VI.lem5.4}
Let\pageoriginale $E$ be a locally coherent topos
(\hyperref[VI.def2.3]{2.3}), and let $C$ be a full generating
subcategory of $E$, stable under fiber products, all of whose objects
are quasi-compact (\hyperref[VI.def1.1]{1.1}). A filtered inductive
limit of $C$-acyclic sheaves \emph{(\hyperrefext[V][V.subsec4.2]{4.2})}
is a $C$-acyclic sheaf.
\end{enonce*}

Let $(F_i)_{i\in I}$ be a filtered inductive system of $C$-acyclic
sheaves, and denote its inductive limit by $F$. Every object $Y$ of
$C$ is coherent (\hyperref[VI.prop2.1]{2.1}), and consequently
$F(Y)=\displaystyle\mathop{\varinjlim}_iF_i(Y)$ \hyperref[VI.thm1.23]{(1.2.3)}. To show that
$F$ is $C$-acyclic, it is enough to show that
$\check H^q(Y,F)=0$ for every $q>0$ and every $Y$ of $C$
(\hyperrefext[V][V.prop4.3]{4.3}). Let $X$ be an object of $C$, and
let $\Xx=(X_{\alpha}\to X)_{\alpha\in A}$ be a finite covering
family by objects of $C$. We have
$H^q(\Xx,F)=H^q(C^{\bigdot}(\Xx,F))$
(\hyperrefext[V][V.subsubsec2.4.3]{2.4.3}). By the preceding and the
fact that $C^{\bigdot}(\Xx,F)$ involves only finite products of groups
(\hyperrefext[V][V.subsubsec2.3.3]{2.3.3}), we have
$C^{\bigdot}(\Xx,F)=\varinjlim_iC^{\bigdot}(\Xx,F_i)$. Hence
$$
H^q(\Xx,F)
=H^q\bigl(\varinjlim_iC^{\bigdot}(\Xx,F_i)\bigr)
=\varinjlim_iH^q(C^{\bigdot}(\Xx,F_i))=0
$$
for $q>0$ (\hyperrefext[V][V.prop4.3]{4.3}). Consequently
$\check H^q(X,F)=\varinjlim_{\mathfrak X}H^q(\mathfrak X,F)=0$.

\begin{enonce*}{Corollary 5.5}\etiquette[5.5]{VI.cor5.5}
Let $E$ be a coherent topos (\hyperref[VI.def2.3]{2.3}), and let
${(X_i)_{i\in I}}$ be a decreasing filtered family of coherent
subobjects of the final object. Denote by $\Phi$ the family of closed
subtopoi of $E$ \emph{(IV 9)} contained in the closed complement of
one of the $X_i$. Then $\Phi$ is a family of supports of $E$
\emph{(\hyperrefext[V][V.def6.12]{6.12})}, and for every integer $q$
the functors $H^q_{\Phi}(E,-)$ and $H^q_{\Phi}$
\emph{(\hyperrefext[V][V.subsec6.13]{6.13})} commute with filtered
inductive limits.
\end{enonce*}

Let $Z_i$ be the closed complement of $X_i$. For every abelian sheaf
$F$, there is an exact sequence (\hyperrefext[V][V.prop6.5]{6.5}):
$$
0 \to H^{0}_{Z_i} (E, F) \to H^{0} (E, F) \to H^{0} (X_i, F) \to
H^{1}_{Z_i} (E, F) \to \ldots\quad {\quoi}.
$$
The functors $H^q(E,F)$ and $H^q(X_i,F)$ commute with filtered
inductive limits in $F$ (\hyperref[VI.cor5.2]{5.2},
\hyperref[VI.cor5.3]{5.3}). Consequently $H^q_{Z_i}(E,F)$ does so as
well, whence the analogous property for $H^q_{\Phi}(E,F)$ after
taking the inductive limit over the $Z_i$
(\hyperrefext[V][V.subsec6.13]{6.13}). The assertion concerning the
functors $\mathcal H^q_{\Phi}$ follows by localization and passage to
the associated sheaf.

\begin{enonce*}[remark]{Definition 5.6}\etiquette[5.6]{VI.def5.6}
Let\pageoriginale $E$ be a ringed topos with ring $A$, and let $q$ be
an integer. An $A$-module $G$ is said to be of finite
$q$-presentation if there is a family $X_i$ of objects of $E$ covering
the final object of $E$ such that, for every $i$, there is an exact
sequence of modules on $E_{/X_i}$:
\begin{equation*}
L_q \to L_{q-1} \to \cdots L_1 \to L_0 \to G_{/X_i} \to 0\tag{5.5.1}
\end{equation*}\etiquette[5.5.1]{VI.eq5.5.1}
where, for every $p$, $L_p=A^{n_p}_{/X_i}$, with $n_p$ an integer.
\end{enonce*}

\begin{enonce*}{Proposition 5.7}\etiquette[5.7]{VI.prop5.7}
Let $G$ be a sheaf of finite $q$-presentation. For every $i\leq q-1$,
the functor $\Eext^i_A(G,-)$ commutes with filtered inductive limits
of $A$-modules.
\end{enonce*}

The problem is local on $E$ (\hyperrefext[V][V.prop6.1]{6.1}). Thus,
after localizing to the $E_{/X_i}$, we may suppose that there is an
exact sequence of the form \eqhyperref[VI.eq5.5.1]{5.5.1}. Put
$L^{\bigdot}=L_q\to\cdots\to L_1\to L_0$. We have
$\Eext^i_A(G,F)\simeq\Hh^i(\Hhom_A(L^{\bigdot},F))$, and the functor
$F\mapsto\Hhom_A(L^{\bigdot},F)$ commutes with inductive limits,
which proves the proposition.

\begin{enonce*}{Corollary 5.8}\etiquette[5.8]{VI.cor5.8}
Let $(E,A)$ be a ringed topos, let $G$ be a sheaf of $A$-modules of
finite $q$-presentation, and let $X$ be an algebraic and coherent
object of $E$. The functors $F\mapsto\Ext^i_A(X;G,F)$, for
$\frac{i(i-1)}{2}\leq q-1$, commute with filtered inductive limits
in $F$.
\end{enonce*}

This follows from \hyperref[VI.prop5.7]{5.7} and
\hyperref[VI.cor5.3]{5.3} by the spectral sequence of V \hyperref[V.6.1.3]{6.1.3}.

\setcounter{subsection}{8}
\subsection{}\etiquette[5.9]{VI.subsec5.9}

Let $(E,A)$ be a ringed topos, let $\Phi$ be a family of supports of
$E$ (\hyperrefext[V][V.def6.12]{6.12}), let $G$ be an $A$-module of
finite $r$-presentation (\hyperref[VI.def5.6]{5.6}), and let $X$ be
an algebraic and coherent object of $E$ (\hyperref[VI.def2.3]{2.3}).
Taking the inductive limit over the closed subtopoi in $\Phi$ in the
last spectral sequences of V \hyperref[V.6.9.1]{6.9.1} and V \hyperref[V.6.9.2]{6.9.2}, respectively, gives
two spectral sequences (\hyperrefext[V][V.subsec6.13]{6.13}):
\begin{equation*}
\begin{cases}
'E^{pq}_2 = \varinjlim_{Z \in \Phi} \Ext^p_A (X; G, H^q_Z F)
\Rightarrow \Ext^{p+q}_{A, \Phi} (X; G, F){\quoi},\\
''E^{pq}_2 = \varinjlim_{Z \in \Phi} \Eext^p_A (G, H^q_Z F)
\Rightarrow \Eext^{p+q}_{A, \Phi} (G, F){\quoi}.
\end{cases}\tag{5.9.1}
\end{equation*}\etiquette[5.9.1]{VI.eq5.9.1}

It follows from \hyperref[VI.prop5.7]{5.7} and
\hyperref[VI.cor5.8]{5.8} that there are canonical isomorphisms
\begin{equation*}
\begin{cases}
'E^{pq}_2 \simeq \Ext^p_A (X; G, H^q_\Phi F)&\quav
\frac{p(p-1)}{2} \leq r - 1{\quoi},\\
''E^{pq}_2 \simeq \Eext^p_A (G, H^q_\Phi F) &\quav p \leq r-1{\quoi}.
\end{cases}\tag{5.9.2}
\end{equation*}\etiquette[5.9.2]{VI.eq5.9.2}

\begin{enonce*}{Corollary 5.10}\etiquette[5.10]{VI.cor5.10}
Use\pageoriginale the hypotheses and notation of
\hyperref[VI.cor5.5]{5.5}. Let $A$ be a ring of $E$, and let $G$ be an
$A$-module of finite $q$-presentation. For every integer $i$ such
that $\dfrac{i(i-1)}{2}\leq q-1$, the functors
$F\mapsto\Ext^i_{A,\Phi}(E;G,F)$ and
$F\mapsto\Eext^i_{A,\Phi}(G,F)$
\emph{(\hyperrefext[V][V.subsec6.13]{6.13})} commute with filtered
inductive limits.
\end{enonce*}

This follows from \hyperref[VI.cor5.5]{5.5} and
\hyperref[VI.prop5.7]{5.7} by the first spectral sequences of V \hyperref[V.6.9.1]{6.9.1}
and V \hyperref[V.6.9.2]{6.9.2}, respectively.


% END INCLUDED SGA4_English_current_buildable_source_20260906/source/06/drafts/SGA4_VI_LINES_3240_3413_English.texfrag


% BEGIN INCLUDED SGA4_English_current_buildable_source_20260906/source/06/drafts/SGA4_VI_LINES_3414_3872_English.texfrag
\section{Inductive and projective limits of a fibered category}\etiquette[6]{VI.sec6}

\setcounter{subsection}{-1}
\subsection{}\etiquette[6.0]{VI.subsec6.0}

We assume that the reader is familiar with the theory of fibered
categories \cite{sga4.exp.VI.bib.5}. The essential purpose of this section is to fix
the terminology and notation concerning fibered categories. Unless
otherwise stated, all categories considered in this section belong to
a fixed universe $\Uu$.

Let $\Ff$, $\Gg$, and $\Ee$ be three categories, and let
$\pi:\Ff\to\Ee$ and ${\pi}':\Gg\to\Ee$ be two functors. We denote by
$\Hhom_{\Ee}(\Ff,\Gg)$ the category whose objects are the functors
$u:\Ff\to\Gg$ such that the diagram
\[
\xymatrix@=1.75cm{
\Ff \ar@{-}[rr]^{u} \ar[rd]_{\pi}&             &\Gg \ar[ld]^{\pi'}\\
           & \Ee &
}\]
commutes, and whose morphisms are the $\Ee$-morphisms of functors,
\ie the morphisms of functors carried by $\pi'$ to identity
morphisms.

Let $\xi$ be an object of $\Ee$. The objects $X$ of $\Ff$ such that
$\pi(X)=\xi$ are called \emph{the objects of $\Ff$ over} $\xi$. Let
$f:\xi\to\eta$ be a morphism of $\Ee$. The morphisms $m$ of $\Ff$
such that $\pi(m)=f$ are called the morphisms of $\Ff$
\emph{over}\pageoriginale $f$. Let $f:\xi\to\eta$ be a morphism of
$\Ee$, let $X$ be an object of $\Ff$ over $\xi$, and let $Y$ be an
object of $\Ff$ over $\eta$; we denote by $\Hom_f(X,Y)$ the set of
morphisms from $X$ to $Y$ over $f$.

\begin{enonce*}{Definition 6.1}\etiquette[6.1]{VI.def6.1}
\begin{enumerate}
\renewcommand{\labelenumi}{\theenumi)}
\item A morphism $m:X\to Y$ of $\Ff$ is called \emph{Cartesian} if,
  for every morphism $p:Z\to Y$ over $\pi(m)$, there exists a unique
  morphism $q:Z\to X$ over the identity of $\pi(X)$ such that
  $mq=p$.

\item The category $\Ff$ is called \emph{prefibered} by $\pi$ over
  $\Ee$ if, for every morphism $f:\xi\to\eta$ of $\Ee$, every object
  of $\Ff$ over $\eta$ is the target of a Cartesian morphism over
  $f$.

\item The category $\Ff$ is called \emph{fibered} by $\pi$ over
  $\Ee$ if it is prefibered and the composite of any two composable
  Cartesian morphisms of $\Ff$ is Cartesian.
\end{enumerate}
\end{enonce*}

\setcounter{subsection}{1}
\subsubsection{}\etiquette[6.1.1]{VI.subsubsec6.1.1}

Let $\xi$ be an object of $\Ee$. The \emph{fiber category} of $\Ff$
at $\xi$, denoted by $\Ff_{\xi}$, is the subcategory of $\Ff$ whose
objects are the objects of $\Ff$ over $\xi$ and whose morphisms are
the morphisms of $\Ff$ over $\id_{\xi}$. For two objects $X$ and $Y$
of $\Ff_{\xi}$, we denote by $\Hom_{\xi}(X,Y)$ the set of morphisms
from $X$ to $Y$ in $\Ff_{\xi}$.

\subsubsection{}\etiquette[6.1.2]{VI.subsubsec6.1.2}

Let $f:\xi\to\eta$ be a morphism of $\Ee$. An object $Y$ over
$\eta$ is the target of a Cartesian morphism over $f$ if and only if
the set-valued functor on the fiber $\Ff_{\xi}$
$$
Z \mapsto \Hom_f(Z, Y)\quad Z \in \ob (\Ff_{\xi}){\quoi},
$$
is representable. When the category $\Ff$ is prefibered, the functor
$$
Y \mapsto \Hom_f(., Y)\quad Y \in \ob (\Ff_{\eta}){\quoi},
$$
with values in the category of presheaves of sets on $\Ff_{\xi}$,
in fact takes its values in the category of representable functors on
$\Ff_{\xi}$, and therefore defines, up to unique isomorphism, a
functor $f^{\ast}:\Ff_{\eta}\to\Ff_{\xi}$, called the \emph{inverse
image\pageoriginale functor for} $f$. Suppose that $\Ff$ is
prefibered and choose, for every morphism $f$ of $\Ee$, a base-change
functor $f^{\ast}$. The category $\Ff$ is then fibered over $\Ee$ if
and only if, for every pair $f,g$ of composable morphisms of $\Ee$,
the composite functor $f^{\ast}g^{\ast}$ is an inverse image functor.
Let
\begin{equation*}
C_{f, g} : f^{\ast} g^{\ast} \to (gf)^{\ast}\tag{6.1.2.1}
\end{equation*}\etiquette[6.1.2.1]{VI.eq6.1.2.1}
be the canonical isomorphism. The isomorphisms $C_{f,g}$ satisfy a
cocycle condition arising from the associativity of composition of
morphisms in $\Ee$:
\begin{equation*}
C_{f, hg}(f^{\ast} \circ C_{g, h}) = C_{gh, f}(C_{f, g} \circ h^{\ast})\tag{6.1.2.2}
\end{equation*}\etiquette[6.1.2.2]{VI.eq6.1.2.2}

\subsubsection{}\etiquette[6.1.3]{VI.subsubsec6.1.3}

Conversely, suppose that for every object $\xi$ of $\Ee$ one is given
a category $\Ff_{\xi}$, for every morphism $f:\xi\to\eta$ a functor
$f^{\ast}:\Ff_\eta\to\Ff_{\xi}$, and for every pair $(f,g)$ of
morphisms of $\Ee$ an isomorphism of functors
$C_{f,g}:f^{\ast}g^{\ast}\to(fg)^{\ast}$, such that the $C_{f,g}$
satisfy condition \eqhyperref[VI.eq6.1.2.2]{6.1.2.2}. One can then
construct, in an essentially unique way, a category $\Ff$ fibered over
$\Ee$ whose fibers ``are'' the categories $\Ff_{\xi}$ and whose
base-change functors may be chosen to be ``equal'' to the functors
given in advance \cite{sga4.exp.VI.bib.5}.

\subsubsection{}\etiquette[6.1.4]{VI.subsubsec6.1.4}
Let $\Ff$ and $\Gg$ be two categories over $\Ee$. We denote by
$$
\Hhom_{\cart/\Ee}(\Ff, \Gg)
$$
the full subcategory of $\Hhom_{\Ee}(\Ff,\Gg)$ defined by the
functors that carry Cartesian morphisms of $\Ff$ to Cartesian
morphisms of $\Gg$. The objects of \\
$\Hhom_{\cart/\Ee}(\Ff,\Gg)$ are called \emph{Cartesian functors}.

Let $C$ and $C'$ be two categories, and let $S$ be a set of morphisms
of the category $C$. We denote by $\Hom_{S^{-1}}(C,C')$ the set of
functors from $C$ to $C'$ that carry the morphisms of $S$ to
isomorphisms.

Let $(\Cat)$ denote the category whose objects are the categories
belonging to the universe and whose morphisms are the functors between
these categories (the category $(\Cat)$ does not belong to the
universe).

Let $\Ff$ be a category over $\Ee$, and let $S$ be the set of
Cartesian morphisms of $\Ff$. The category $\Ff$ defines two functors
on $(\Cat)$ with values in the category\pageoriginale of sets
belonging to the universe:
$$
C \mapsto \Hom_{S^{-1}}(\Ff,C)\quad C \in \ob (\Cat)
$$

$C \mapsto \Hom_{\Cart /\Ee}(\Ff,C\times\Ee) =$ \{Cartesian
functors from $\Ff$ to $C\times\Ee$\}

(the category $C\times\Ee$ is regarded as a category over $\Ee$ by
the second projection functor).

\begin{enonce*}{Proposition 6.2}\etiquette[6.2]{VI.prop6.2}
The functors $(\Cat)\to(\Ens):$
\begin{align*}
C &\mapsto \Hom_{S^{-1}} (\Ff, C)\\
C &\mapsto \Hom_{\Cart/\Ee} (\Ff, C \times \Ee)
\end{align*}
are canonically isomorphic. They are representable.
\end{enonce*}

\begin{enonce*}[remark]{Proof}
To prove the first assertion, it is enough to observe that the
Cartesian morphisms of $C\times\Ee$ are the morphisms of the form
$m\times f$, where $m$ is an isomorphism of $C$. To prove the second
assertion, it is enough to prove that the functor
$\Hom_{S^{-1}}(\Ff,.)$ is representable. This follows from \cite{sga4.exp.VI.bib.1}.
We indicate only the idea of the proof. One formally adjoins to the
morphisms of $\Ff$ inverses of the morphisms of $S$. One considers
the free category generated by the objects of $\Ff$, the morphisms of
$\Ff$, and the formal inverses of the morphisms of $S$ (the path
category). One then takes the quotient by the relations arising from
the relations among morphisms of $\Ff$ and relations of the form
$$
s^{-1}s=\id\quav \text{s}s^{-1}=\id\quad s\in S{\quoi}.
$$
The category thus obtained is denoted by $\Ff(S^{-1})$; equipped with
the canonical functor $Q:\Ff\to\Ff(S^{-1})$, it represents the
functor $\Hom_{S^{-1}}(\Ff,.)$.
\end{enonce*}

\begin{enonce*}{Definition 6.3}\etiquette[6.3]{VI.def6.3}
When $\Ff$ is fibered over $\Ee$, the functor
$\Hom_{S^{-1}}(\Ff,.)$ is denoted by
$\varinjlim_{\mathcal{E}^{\circ}}\Ff$ and is called the inductive-limit
functor of $\Ff$ over $\mathcal{E}^{\circ}$. The category representing
it is likewise denoted by
$\varinjlim_{\mathcal{E}^{\circ}}\Ff$ and is called the inductive-limit
category of $\Ff$ over $\Ee^{\circ}$.
\end{enonce*}

\setcounter{subsection}{4}
\setcounter{subsubsection}{-1}
\subsubsection{}\etiquette[6.4.0]{VI.subsubsec6.4.0}

Suppose that $\Ff$ is fibered over $\Ee$. For every morphism $f$ of
$\Ee$, choose a base-change functor over $f$, and let
$$
Q :\Ff \to \varinjlim_{\mathcal{E}^{\circ}} \Ff
$$
be the\pageoriginale canonical functor. Composing the inclusion
functors of the fibers of $\Ff$ into $\Ff$ with $Q$ gives, for every
object $\xi$ of $\Ee$, a functor
$$
u_{\xi} : \Ff_{\xi} \to \varinjlim_{\mathcal{E}} \Ff
$$
and, for every morphism $f:\xi\to\eta$ of $\Ee$, a diagram that
commutes up to canonical isomorphism:
\[
\xymatrix@R=1cm@C=3cm{
\Ff_{\eta} \ar[rd]^{u_{\eta}} \ar[dd]_{f^\ast}&\\
                   &\displaystyle\mathop{\varinjlim}_{\Ee^{0}}\Ff\\
\Ff_{\xi} \ar[ru]_{u_{\xi}} &
}
\]
The category $\varinjlim_{\mathcal{E}^{\circ}}\Ff$ then appears as
the inductive limit \emph{in the sense of pseudofunctors} \cite{sga4.exp.VI.bib.5} of
the pseudofunctor $\Ee^{\circ}\to\Cat$ that associates
$\Ff_{\xi}$ with every $\xi\in\ob$, and the functor
$f^{\ast}:\Ff_{\eta}\to\Ff_{\xi}$ with every
$f:\xi\to\eta$. Note, however, that even when
$\xi\mapsto\Ff_{\xi}$ is an actual functor, the category
$\varinjlim_{\mathcal{E}^{\circ}}\Ff$ is not in general the
inductive limit in the sense of I 2 of the functor
$\xi\mapsto\Ff_{\xi}$ (\cf \hyperref[VI.exe6.8]{6.8}).

\begin{enonce*}{Proposition 6.4}\etiquette[6.4]{VI.prop6.4}
Let $\Ff$ be a category fibered over $\Ee$. Suppose that the category
$\Ee$ has the following properties:
\begin{itemize}
\item[L1)] Every diagram
\[
\xymatrix@C=1.5cm@R=.65cm{
\bigdot \ar[rd]& \\
        & \bigdot\\
\bigdot \ar[ru]&
}
\]
can be embedded in a commutative diagram
\[
\xymatrix@C=1.25cm@R=1cm{
        & \bigdot \ar[rd]&\\
\bigdot \ar[ru]\ar[rd]&         & \bigdot\\
        & \bigdot \ar[ru]&
}
\]

\item[L2)] For\pageoriginale every pair of morphisms
$u,v:\cdot\rightrightarrows\cdot$ such that there exists a morphism
$t$ satisfying $tu=tv$, there exists a morphism $w$ satisfying
$uw=vw$.

(Note that if the category $\Ee^{\circ}$ is pseudofiltered
\emph{(I \hyperref[I.def2.7]{2.7})}, it has properties \emph{L1)} and \emph{L2)}.) The set
$S$ of Cartesian morphisms of $\Ff$ then has the following
properties:

\item[Fr1)] The set $S$ is stable under composition. Isomorphisms
  belong to~$S$.

\item[Fr2)] Every diagram
\[
\xymatrix@C=1.5cm@R=.65cm{
\bigdot \ar[rd]& \\
        & \bigdot\\
\bigdot \ar[ru]^{s}&,
}
\]
where $s$ belongs to $S$, can be completed to a commutative diagram
\[
\xymatrix@C=1.25cm@R=1cm{
        & \bigdot \ar[rd]&\\
\bigdot \ar[ru]^{t}\ar[rd]&         & \bigdot\\
        & \bigdot \ar[ru]^{s}&
}
\]
where $t$ belongs to $S$.

\item[Fr3)] For every pair of morphisms
$u,v:\cdot\rightrightarrows\cdot$ such that there exists a morphism
$s\in S$ satisfying $su=sv$, there exists a morphism $t\in S$ such
that $ut=vt$.
\end{itemize}
\end{enonce*}

\begin{enonce*}[remark]{Proof}
Left to the reader as an exercise.
\end{enonce*}

\begin{enonce*}{Proposition 6.5}\etiquette[6.5]{VI.prop6.5}
Let $\Ff$\pageoriginale be a category, and let $S$ be a set of
morphisms of $\Ff$ having properties \emph{Fr1), Fr2), and Fr3)
(\hyperref[VI.prop6.4]{6.4})}. For every object $X$ of $\Ff$, let
$S(X)$ denote the category of morphisms of $S$ with target $X$. The
category $S(X)$ is cofiltered \emph{(I \hyperref[I.def2.7]{2.7})}. The category
$\Ff(S^{-1})$ (which represents the functor
$\Hom_{S^{-1}}(\Ff,\cdot)$) may then be described as follows:
\begin{enumerate}
\renewcommand{\theenumi}{\alph{enumi}}
\renewcommand{\labelenumi}{\theenumi)}
\item The objects of $\Ff(S^{-1})$ are the objects of $\Ff$.

\item Let $X$ and $Y$ be two objects of $\Ff(S^{-1})$. We have
$$
\Hom_{\Ff(S^{-1})} (X, Y)= \varinjlim_{S(X)}
\Hom_{\Ff} (., Y){\quoi}.
$$

\item Let $f:X\to Y$ and $g:Y\to Z$ be two morphisms of
$\Ff(S^{-1})$. Let
\[
\xymatrix{
\bigdot  \ar[d]_{t} \ar[rrd]^{f'}&\\
X        & &Y & & t\in S
}
\]
(resp.
\[
\xymatrix{
\bigdot \ar[d]_{t'} \ar[rrd]^{g'}&\\
Y                          & & Z & & t'\in S)
}
\]
be a diagram in $\Ff$ whose image is $f$ (\resp $g$). Let
\[
\xymatrix@C=2.5cm@R=1cm{
\bigdot \ar[rd]^{f''} \ar[d]_{t''}& \\
\bigdot \ar[rd]^{f'}        & \bigdot \ar[d]^{t'}\\
                            & Y      & t''\in S
}
\]
be a commutative diagram whose existence is ensured by Fr2). The image
in $\Hom_{\Ff(S^{-1})}(X,Z)$ of the diagram
\[
\xymatrix@C=2cm{
\bigdot \ar[d]_{tt''} \ar[rd]^{g'f''}& \\
X       & Z
}
\]
is the composite morphism $gf$.

Let\pageoriginale
$$
Q : \Ff \to \Ff (S^{-1})
$$
be the canonical functor.

\item The functor $Q^{\ast}:\Ff(S^{-1})^{\widehat{\phantom{,}}}\to
  \Ff^{\widehat{\phantom{,}}}$ ($F\mapsto F\circ Q$ \emph{I \hyperref[I.sec5.0]{5.0}})
  is fully faithful and injective on objects.

\item The functor $Q_{!}:\Ff^{\widehat{\phantom{,}}}\to
  \Ff(S^{-1})^{\widehat{\phantom{,}}}$ (left adjoint to the functor
  $Q^{\ast}$ \emph{(I \hyperref[I.prop5.1]{5.1})}) is left exact. (Consequently, the same is
  true of $Q$ when finite projective limits are representable in
  $\Ff$.)
\end{enumerate}
\end{enonce*}

\begin{enonce*}[remark]{Proof}
For the proof, we refer to [1].
\end{enonce*}

\begin{enonce*}[remark]{Exercise 6.6}\etiquette[6.6]{VI.exe6.6}
Let $\Ff$ be a category, and let $S$ be a set of morphisms of $\Ff$
having properties Fr1), Fr2), and Fr3). For every object $X$ of
$\Ff$, let $S(X)$ denote the category of morphisms with target $X$.
\begin{enumerate}
\renewcommand{\theenumi}{\alph{enumi}}
\renewcommand{\labelenumi}{\theenumi)}
\item The category $S(X)$ is cofiltered (I 2).

\item For every object $X$ of $\Ff$, let $J_S(X)$ denote the set of
  sieves of $X$ that contain a sieve generated by a morphism of $S$.
  The $J_S(X)$ define a topology $T$ on $\Ff$.

\item For the topology $T$, the morphisms of $S$ are bicovering
  (II \hyperref[VI.cor5.2]{5.2}).

\item For every object $X$ of $\Ff$, let $𝒥^2(X)$ denote the category
  of bicovering morphisms with target $X$. The category $S(X)$ is
  cofinal in $𝒥^2(X)$ (I 8).

\item Let $G$ be a presheaf of sets on $\Ff$, and let $\sheaf{a}$ be
  the associated-sheaf functor for the topology $T$. For every object
  $X$ of $\Ff$, there is a canonical isomorphism
$$
\sheaf{a} G(X) \xrightarrow{\sim} \frac{\lim}{S(X)} G( \cdot ){\quoi}.
$$
A presheaf is a sheaf if and only if it carries every morphism of $S$
to an isomorphism.

\item Choose the associated-sheaf functor so that its restriction to
$\Ff$ is injective on objects. Let $\underrightarrow{\Ff}$ denote the
full subcategory of the category of sheaves of sets on $\Ff$ defined
by the sheaves associated with representable presheaves.
\pageoriginale The category $\underrightarrow{\Ff}$ is isomorphic to
the category $\Ff(S^{-1})$ described in \hyperref[VI.prop6.5]{6.5}.
The objects of $\underrightarrow{\Ff}$ form a generating family of
the topos $\Ff^{\sim}$ of sheaves of sets on $\Ff$. The topology on
$\underrightarrow{\Ff}$ induced by the canonical topology of
$\Ff^{\sim}$ is the coarse topology (II \hyperref[II.subsubsec1.1.4]{1.1.4}).

\item Let $Q:\Ff\to\underrightarrow{\Ff}$ be the canonical functor.
The functor $Q$ is a morphism from the site $\underrightarrow{\Ff}$
(with the coarse topology) to the site $\Ff$ (with the topology $T$
of b)) (IV \hyperref[IV.exer5.9]{5.9}), and induces an isomorphism on the corresponding
topoi.

\item Let $u:\Ff\to C$ be a functor. The functor $u$ is continuous if
and only if it carries the morphisms of $S$ to isomorphisms (III \hyperref[III.def1.1]{1.1}).

\item Let $u:\Ff\to C$ be a functor carrying the morphisms of $S$ to
isomorphisms. The functor $u$ factors uniquely as
\[
\xymatrix@=1.5cm{
  \Ff\ar[rr]^{u} \ar[rd]_{Q}&  & C\\
            & \underrightarrow{\Ff}{} \ar[ru]_{v}&
}
\]
(Use III \hyperref[III.rem1.4]{1.4}.)

\item Let $\Pro_S\Ff$ be the full subcategory of $\Pro\Ff$ (I 8)
defined by the pro-objects whose transition morphisms belong to $S$.
Show that the functor $Q$ factors as
\[
\xymatrix@M=.25cm@C=2cm{
\Ff\ar@{^{(}->}[r]^{i}& \Pro_S\Ff\ar[r]^{j}&
\Ff_{S^{-1}}
}
\]
where $j$ is the restriction of $\Pro Q$ to $\Pro_S\Ff$. Show that
the functor $j$ admits a fully faithful left adjoint
$$
A : \Ff_{S^{-1}} \to \Pro_S \Ff{\quoi}.
$$
Show that, for every object $X$ of $\Ff$, $AQ(X)$ is the pro-object
$$
Y \to X \in S(X) \mapsto Y{\quoi}.
$$
\end{enumerate}
\end{enonce*}


% END INCLUDED SGA4_English_current_buildable_source_20260906/source/06/drafts/SGA4_VI_LINES_3414_3872_English.texfrag


% BEGIN INCLUDED SGA4_English_current_buildable_source_20260906/source/06/drafts/SGA4_VI_LINES_3873_4230_English.texfrag
\begin{enonce*}{Proposition 6.7}\etiquette[6.7]{VI.prop6.7}
Let\pageoriginale $\Ee$ be a cofiltered category (I \hyperref[I.sec8.7]{8.7}) and
$\Ff\to\Ee$ a fibered category \emph{(\hyperref[VI.subsubsec6.1.3]{6.1.3})}.
\begin{enumerate}
\renewcommand{\labelenumi}{\theenumi)}
\item Let $\xi$ be an object of $\Ee$, and let $\Ff/\xi$ be the category
  fibered over $\Ee/\xi$ obtained by the base change
  $\Ee/\xi\to\Ee$. The canonical functor
  $\varinjlim_{\Ee/\xi}\Ff/\xi\to\varinjlim_{\Ee}\Ff$ is an equivalence of
  categories.

\item More generally, let $\Ee'\to\Ee$ be a cofinal functor \emph{(I 8)}.
  Let $\Ff'\to\Ee'$ be the fibered category deduced from
  $\Ff\to\Ee$ by the base change $\Ee'\hookrightarrow\Ee$. The canonical
  functor
$$
\varinjlim_{\Ee'} \Ff' \to \varinjlim_{\Ee} \Ff
$$
is an equivalence of categories.
\end{enumerate}
\end{enonce*}

\begin{enonce*}[remark]{Proof}
Exercise.
\end{enonce*}

\begin{enonce*}[remark]{Exercise 6.8}\etiquette[6.8]{VI.exe6.8}
\begin{enumerate}
\renewcommand{\labelenumi}{\theenumi)}
\item Let $g:{\Ee}^{\circ}\to\Cat$ be a functor and $\Gg\to\Ee$ the
  corresponding fibered category \cite{sga4.exp.VI.bib.5}. Show that there exists a
  canonical functor
  $\varinjlim_{\Ee^{\circ}}\Gg\to\varinjlim_{\Ee^{\circ}}g$. Show that
  this functor is not necessarily an equivalence of categories. (One may
  take for $g$ the constant functor whose value is the final object of Cat,
  and for $\Ee$ the category associated with a group.)

\item Suppose that the category $\Ee^{\circ}$ is pseudo-filtered. Show that
  the canonical functor
  $\varinjlim_{\Ee^{\circ}}\Gg\to\varinjlim_{\Ee^{\circ}}g$ is then an
  equivalence of categories.
\end{enumerate}
\end{enonce*}

\begin{enonce*}{Proposition 6.9}\etiquette[6.9]{VI.prop6.9}
Let $\Ff\to\Ee$ be a fibered category (\hyperref[VI.def6.1]{6.1}). The
functor $(\Cat)\to(\Ens)$
$$
C \mapsto \Hom_{\Cart/\Ee} (C \times \Ee, \Ff)
$$
is represented by the category $\Hhom_{\Cart/\Ee}(\Ee,\Ff)$.
\end{enonce*}

Let $C$ be a category, let $F:C\times\Ee\to\Ff$ be a Cartesian
$\Ee$-functor (\hyperref[VI.subsubsec6.1.4]{6.1.4}), and let $X$ be an
object of $C$. The functor $\xi\mapsto F((X,\xi))$ is a Cartesian functor
from $\Ee$ to $\Ff$, denoted $F'(X)$, that depends functorially on $X$.
Hence there is a map $F\mapsto F'$ from
$\Hom_{\Cart/\Ee}(C\Ee,\Ff)$ to
$\Hom(C,\Hhom_{\Cart/\Ee}(\Ee,\Ff))$, functorial in $C$, which is
immediately verified to be a bijection.

\setcounter{subsection}{9}
\subsection{}\etiquette[6.10]{VI.subsec6.10}

The category\pageoriginale $\Hhom_{\Cart/}(\Ee,\Ff)$ is called the category
of \emph{Cartesian sections} of the fibered category $\Ff$. It is also
sometimes called the \emph{projective-limit} category of $\Ff$ over
$\Ee^{\circ}$. It is then denoted
$\varprojlim_{\Ee^{\circ}}\Ff$ \cite{sga4.exp.VI.bib.2}.

\subsection{}\etiquette[6.11]{VI.subsec6.11}

Choose a splitting of the fibered category $\Ff\to\Ee$, \ie for every arrow
$f:\xi\to\eta$ of $\Ee$, choose a base-change functor
$f^{\ast}:\Ff_{\eta}\to\Ff_{\xi}$. For every object $\xi$ of $\Ee$, denote
by
$$
v_{\xi} : \varprojlim_{\Ee^{\circ}} \Ff \to \Ff_{\xi}
$$
the evaluation functor at $\xi$. For every morphism $f:\xi\to\eta$, there
is a diagram of functors commutative up to canonical isomorphism:
\[
\xymatrix@C=2.25cm@R=1.25cm{
\displaystyle\mathop{\varprojlim}_{\Ee^{\circ}} \Ff
\ar@{-}[r]^{v_{\xi}} \ar[rd]^{v_\eta}& \Ff_{\xi} \ar[d]^{f^\ast}\\
                                 & \Ff_\eta.
}
\]
The category $\varprojlim_{\Ee^{\circ}}\Ff$ thus appears as the projective
limit, \emph{in the sense of pseudofunctors}, of
$\xi\mapsto\Ff_{\xi}$. Even when $\xi\mapsto\Ff_{\xi}$ is an actual
functor (\ie even when the chosen splitting is a cleavage \cite{sga4.exp.VI.bib.5}), the
categories $\varprojlim_{\Ee^{\circ}}\Ff$
(\hyperref[VI.subsec6.10]{6.10}) and
$\varprojlim_{\Ee^{\circ}}\Ff_{\xi}$ (I 2) are not equivalent in general.

\section{Fibered topoi and sites}\etiquette[7]{VI.sec7}

\subsection{Fibered topoi}\etiquette[7.1]{VI.subsec7.1}

\begin{enonce*}{Definition 7.1.1}\etiquette[7.1.1]{VI.def7.1.1}
Let $𝒰$ be a universe. A $𝒰$-topos fibered over a category $I$ is a category
fibered over $I$ (\hyperref[VI.def6.1]{6.1})
$$
p : F \to I
$$
whose fibers are $𝒰$-topoi and whose inverse-image
functors\footnote{For the fibered structure
  (\hyperref[VI.subsubsec6.1.2]{6.1.2}).} relative to the morphisms
$f:i\to j$ of $I$ are functors $f^{\ast}:F_j\to F_i$ that are the
inverse-image functors of morphisms of topoi $f\bigdot:F_i\to F_j$
\emph{(IV \hyperref[IV.subsubsec3.1.2]{3.1.2})}.
\end{enonce*}

\setcounter{subsubsection}{1}
\subsubsection{}\etiquette[7.1.2]{VI.subsubsec7.1.2}

Equivalently, a $𝒰$-topos fibered over a category $I$ is a fibered category
$p:F\to I$ whose fibers are $𝒰$-topoi and whose inverse-image
functors\pageoriginale are exact and commute with inductive limits
(IV \hyperref[IV.cor1.6]{1.6}).


\subsubsection{}\etiquette[7.1.3]{VI.subsubsec7.1.3}

For every morphism $f:i\to j$ of $I$, choose an inverse-image functor
$f^{\ast}:F_j\to F_i$. For every pair
$i\xrightarrow{f}j\xrightarrow{g}k$ of composable morphisms of $I$, there is
then a canonical isomorphism
$c_{f,g}:f^{\ast}g^{\ast}\to(gf)^{\ast}$
\eqhyperref[VI.eq6.1.2.1]{6.1.2.1}, and the $c_{f,g}$ have the cocycle
property \eqhyperref[VI.eq6.1.2.2]{6.1.2.2}. In addition, for every
$f:i\to j$, choose a right adjoint $f_{\ast}:F_i\to F_j$ to the functor
$f^{\ast}:F_j\to F_i$. The functor $f_{\ast}:F_i\to F_j$ is defined up to
canonical isomorphism by the property of being right adjoint to $f^{\ast}$.
The general results on adjoint functors then give canonical isomorphisms
\begin{equation*}
{c'}_{g, f} : g_{\ast} f_{\ast} \to (gf)_{\ast}\tag{7.1.3.1}{\quoi}.
\end{equation*}\etiquette[7.1.3.1]{VI.eq7.1.3.1}
These have a cocycle property analogous to property \hyperref[VI.eq6.1.2.2]{6.2.2.2}, which we leave
to the reader to make explicit. Applying the results of
\hyperref[VI.subsubsec6.1.3]{6.1.3}, we then obtain a category fibered over
$I^{\circ}$:
\begin{equation*}
p': F'\to I^{\circ}{\quoi}.\tag{7.1.3.2}
\end{equation*}\etiquette[7.1.3.2]{VI.eq7.1.3.2}
It is readily verified that the fibered category $p':F'\to I^{\circ}$ is
independent, up to unique $I^{\circ}$-isomorphism, of the various choices
used to construct it\footnote{\etiquette[112]{page112}The writer apologizes
  for the nonintrinsic character of this construction.}. The fiber at every
object $i$ of $I^{\circ}$ of the fibered category $p':F'\to I^{\circ}$ is
canonically isomorphic to the $𝒰$-topos $F_i$ and is identified with it. The
inverse-image functors of the fibered category $p':F'\to I^{\circ}$ are
\emph{direct-image} functors of morphisms of topoi.

\subsubsection{}\etiquette[7.1.4]{VI.subsubsec7.1.4}

Let $p:F\to I$ be a fibered $𝒰$-topos. For every $f:i\to j$, choose a
morphism of topoi (IV \hyperref[IV.subsubsec3.1.2]{3.1.2}) $f\bigdot:F_i\to F_j$ such that the
inverse-image functor $f^{\ast}$ (under $f\bigdot$) is the inverse-image
functor for the fibered structure (\hyperref[VI.def6.1]{6.1}). We then
obtain cocycles $c_{f,g}$ relating these morphisms of topoi
\eqhyperref[VI.eq7.1.3.1]{7.1.3.1} and
\eqhyperref[VI.eq7.1.3.2]{7.1.3.2}. The collection of the morphisms
$f.$, $f\in\text{Fl}(I)$, and the cocycles $c_{f,g}$ is called a
\emph{bisplitting} of the fibered topos $F\to I$. Choosing a bisplitting
makes it possible to associate with the fibered topos $F\to I$ a
pseudofunctor $i\mapsto F_i$ from the category $I$ to the 2-category
of\pageoriginale $𝒰$-topoi. Conversely, one can associate with it in an
essentially unique manner (\cf \cite{sga4.exp.VI.bib.5}) a fibered $𝒰$-topos $F\to I$
endowed with a bisplitting such that the corresponding pseudofunctor is
canonically isomorphic to $i\mapsto F_i$. In practice, a fibered $𝒰$-topos
$F\to I$ will most often be denoted by $(F_i)_{i\in I}$, with the implicit
assumption that a bisplitting has been chosen. Note, however, that the
constructions carried out on fibered topoi (such as the associated total
topos (\hyperref[VI.subsec7.4]{7.4}) and the projective limit
(\hyperref[VI.subsec8.1]{8.1})) depend only on the fibered topoi and not on
any bisplittings that may have been chosen.

\begin{enonce*}{Definition 7.1.5}\etiquette[7.1.5]{VI.def7.1.5}
Let $p:F\to I$ and $q:G\to I$ be two $𝒰$-topoi fibered over $I$. A morphism
$m$ of fibered topoi from $(F,p)$ to $(G,q)$ consists of the datum of an
$I$-functor $m^{\ast}:G\to F$ \emph{(\hyperref[VI.subsubsec6.1.4]{6.1.4})}
and, for every object $i$ of $I$, the datum of a right adjoint
$m_{i\ast}:F_i\to G_i$ to the functor $m^{\ast}_i:G_i\to F_i$ obtained by
restricting $m^{\ast}$ to the fibers at $i$. These data must moreover be
subject to the following condition: for every object $i$ of $I$, the pair
of adjoint functors $(m^{\ast}_i,m_{i\ast})$ is a morphism from the topos
$F_i$ to the topos $G_i$.
\end{enonce*}

\setcounter{subsubsection}{5}
\subsubsection{}\etiquette[7.1.6]{VI.subsubsec7.1.6}

Let $m:(F,p)\to(G,q)$ be a morphism of $𝒰$-topoi fibered over $I$, and let
$p':F'\to I^{\circ}$ and $q':G'\to I^{\circ}$ be the fibered categories
associated by the construction of \hyperref[VI.subsubsec7.1.3]{7.1.3} with
the fibered topoi $(F,p)$ and $(G,q)$, respectively. Choose bisplittings of
$F$ and $G$ (\hyperref[VI.subsubsec7.1.4]{7.1.4}), denoted in both cases by
$(f:i\to j)\mapsto f.:(f^{\ast},f_{\ast})$, from $F_i$ to $F_j$ and from
$G_i$ to $G_j$, respectively. Since $f^{\ast}:G\to F$ is an $I$-functor,
for every morphism $f:i\to j$ there is a morphism of functors (transition
morphism)
\begin{equation*}
b_f : f^{\ast}m^{\ast}_j \to m^{\ast}_i f^{\ast}\tag{7.1.6.1}{\quoi},
\end{equation*}\etiquette[7.1.6.1]{VI.eq7.1.6.1}
and for every pair $i\xrightarrow{f}j\xrightarrow{g}k$ of composable
morphisms, we have
\begin{equation*}
m^{\ast}_i (c_{f, g}) b^{\ast}_f f^{\ast} (b_g) = b_{gf} c_{f, g}\tag{7.1.6.2}{\quoi}.
\end{equation*}\etiquette[7.1.6.2]{VI.eq7.1.6.2}
Passing to adjoint functors gives morphisms
\begin{equation*}
b'_f : f_{\ast}m_{i \ast} \to m_{j\ast} f_{\ast}\tag{7.1.6.3}
\end{equation*}\etiquette[7.1.6.3]{VI.eq7.1.6.3}
that satisfy relations analogous to \eqhyperref[VI.eq7.1.6.2]{7.1.6.2}.
Since the functors $f_{\ast}$ are\pageoriginale the inverse-image functors
of the fibered categories $(F',p')$ and $(G',q')$
(\hyperref[VI.subsubsec7.1.3]{7.1.3}), one can construct an
$I^{\circ}$-functor
\begin{equation*}
m_{\ast} : F'\to G'\tag{7.1.6.4}{\quoi},
\end{equation*}\etiquette[7.1.6.4]{VI.eq7.1.6.4}
whose restrictions to the fibers are the functors $m_{i\ast}$ and whose
transition isomorphisms are the $b'_f$. It is verified, with a little
patience, that the functor $m_{\ast}$ so constructed does not depend on the
choice of the morphisms of topoi $f$.\footnote{See the note at the bottom of
  page \pagehyperref[page112]{112}.}

\subsubsection{}\etiquette[7.1.7]{VI.subsubsec7.1.7}

Denote by $\Homtop_I(F,G)$ the set of morphisms of fibered topoi between the
fibered topos $(F,p)$ and the fibered topos $(G,q)$. The set
$\Homtop_I(F,G)$ is the set of objects of a category denoted
$\Hhomtop_I(F,G)$: if $m$ and $n$ are two morphisms of fibered topoi, a
morphism from $m$ to $n$ is an $I$-morphism of functors from $n^{\ast}$ to
$m^{\ast}$. By adjunction, one also obtains an $I^{\circ}$-morphism from the
functor $m_{\ast}$ to the functor $n_{\ast}$
(\hyperref[VI.subsubsec7.1.6]{7.1.6}). Thus, in the end, there are two
functors
\hypertarget{sga4.exp.VI.equation.7.1.7.1}{}\begin{align*}
\Hhomtop_I (F, G) & \to \Hhom_I (G, F)^{\circ}\tag{7.1.7.1}\\
\Hhomtop_I (F, G) &\to \Hhom_{I^{\circ}} (F', G')\tag{7.1.7.2}
\end{align*}\hypertarget{sga4.exp.VI.equation.7.1.7.2}{}
that associate with every morphism
$(m^{\ast},(m_{i\ast})_{i\in\Ob I})$, on the one hand the functor
$m^{\ast}\in\Hom_I(G,F)$ and, on the other hand, the functor
$m_{\ast}\in\Hom_{I^\circ}(F',G')$ obtained by gluing the $m_{i\ast}$ as in
(\hyperref[VI.subsubsec7.1.6]{7.1.6}). The functor $m^{\ast}:G\to F$
\hyperlink{sga4.exp.VI.equation.7.1.7.1}{(7.1.7.1)} is called the \emph{inverse-image functor} under the morphism of
fibered topoi $m:F\to G$; the functor $m_{\ast}:F'\to G'$ \hyperlink{sga4.exp.VI.equation.7.1.7.2}{(7.1.7.2)} is
called the \emph{direct-image functor} under the morphism of fibered topoi
$m:F\to G$. It follows immediately from the definitions that the functors
$m\mapsto m^{\ast}$ and $m\mapsto m_{\ast}$ are \emph{fully faithful}, and
that their essential images are, on the one hand, the set of $I$-functors
from $G$ to $F$ that, fiber by fiber, are exact and commute with inductive
limits and, on the other hand, the set of $I^0$-functors from $F'$ to $G'$
that, fiber by fiber, are exact and commute with inductive limits and, on the
other hand, the set of $I^0$-functors from $F'$ to $G'$ that, fiber by fiber,
possess an exact left adjoint. This justifies the abuses of language in
which a morphism of fibered topoi is defined by its direct image or its
inverse image. In fact, such a morphism of fibered\pageoriginale topoi is
then defined only up to unique isomorphism.

Denote by $\Hhomtop_{\Cart/I}(F,G)$ the full subcategory of
$\Hhomtop_I(F,G)$ generated by the morphisms $m$ such that $m_{\ast}$ (or,
equivalently, $m^{\ast}$) is a Cartesian functor
(\hyperref[VI.subsubsec6.1.4]{6.1.4}). Such morphisms are called
\emph{Cartesian morphisms}. Thus there are fully faithful functors induced
by \hyperlink{sga4.exp.VI.equation.7.1.7.1}{(7.1.7.1)} and \hyperlink{sga4.exp.VI.equation.7.1.7.2}{(7.1.7.2)}:
\hypertarget{sga4.exp.VI.equation.7.1.7.3}{}\begin{align*}
\Hhomtop_{\Cart/I} (F, G) & \to \Hhom_{\cart /I} (G,
  F)^{\circ}{\quoi},\tag{7.1.7.3}\\
\Hhomtop_{\Cart/I} (F, G) &\to \Hhom_{\cart /I^{\circ}} (F',
  G'){\quoi}.\tag{7.1.7.4}
\end{align*}\hypertarget{sga4.exp.VI.equation.7.1.7.4}{}

\subsubsection{}\etiquette[7.1.8]{VI.subsubsec7.1.8}

Let $p:F\to I$, $q:G\to I$, and $r:H\to I$ be three $𝒰$-topoi fibered over
$I$, and let $m:F\to G$ and $n:G\to H$ be two morphisms of fibered
$𝒰$-topoi. As in IV \hyperlink{sga4.exp.IV.subsection.3.3}{3.3}, define the composite of the two morphisms $m$ and
$n$, denoted $nm:F\to H$, and if $𝒱$ is a universe such that $𝒰\in 𝒱$,
define a category ($𝒱$-$𝒰$-$\Top_{/I}$) whose objects are the $𝒰$-topoi
fibered over $I$ with fiber categories that are elements of $𝒱$, and whose
morphisms are the morphisms of $𝒰$-topoi fibered over $I$. Moreover, as in
IV \hyperref[IV.subsubsec3.3.2]{3.3.2}, the composition map
$$
\Homtop_{/I} (F, G) \times \Homtop_{/I} (G, F) \to \Homtop_{/I} (F, H)
$$
is the map induced on objects by a ``functor of composition of morphisms'':
$$
\Hhomtop_{/I} (F, G) \Hhomtop_{/I} (G, H) \to \Hhomtop_{/I} (F, H)
$$
and the considerations of IV \hyperlink{sga4.exp.IV.subsection.3.3}{3.3} and IV \hyperref[IV.subsec3.4]{3.4} extend without change to the
case of fibered $U$-topoi. The reader will moreover have noticed that when
$I$ is a one-point category (one object and one identity morphism), a
$𝒰$-topos fibered over $I$ is ``nothing but'' a $𝒰$-topos, and a morphism
of fibered $𝒰$-topoi is ``nothing but'' a morphism of $𝒰$-topoi; the
preceding constructions then reduce to those of IV 3.

\subsubsection{}\etiquette[7.1.9]{VI.subsubsec7.1.9}

Let $p:F\to I$ be a fibered $𝒰$-topos and $\phi:J\to I$ a functor. The
fibered category $p_J:F_J\to J$ deduced from $(F,p)$ by the base change
$\phi:J\to I$ is a $𝒰$-topos fibered over $J$. Let $p':F'\to I^{\circ}$ and
$(p_J)':(F_J)'\to J^{\circ}$ be the fibered categories deduced respectively
from $(F,p)$ and $(F_J,p_J)$ by the construction of
\hyperref[VI.subsubsec7.1.3]{7.1.3}. It is immediately verified that the
category $((F_J)',(p_J)')$ is canonically isomorphic to the category
$(p')_J:(F')_J\to J^{\circ}$ deduced\pageoriginale from $p':F'\to I^{\circ}$
by the base change $\phi:J^{\circ}\to I^{\circ}$. These two categories
fibered over $J^{\circ}$ will henceforth be identified and denoted
$p'_j:F'_j\to J^{\circ}$. The base-change operation is functorial, and even
2-functorial, in its argument: for every pair $p:F\to I$ and $G\to I$ of
$𝒰$-topoi fibered over $I$, there are two commutative diagrams of categories
and functors:
\[\vcenter{
\xymatrix{
\Hhomtop_I(F,G) \ar[r] \ar[d]& \Hhom_I(G,F)^{\circ} \ar[d]\\
\Hhomtop_J(F_J,G_J) \ar[r]& \Hhom_J(G_J,F_J)^{\circ},
}}\leqno(7.1.9.1)
\]\etiquette[7.1.9.1]{VI.eq7.1.9.1}
\[\vcenter{
\xymatrix{
\Hhomtop_I(F,G) \ar[r] \ar[d]& \Hhom_{I^{\circ}}(F',G') \ar[d]\\
\Hhomtop_J(F_J,G_J) \ar[r]&
\Hhom_{J^{\circ}}(F'_{J^{\circ}},G_{J^{\circ}}),
}}\leqno{(7.1.9.2)}
\]\etiquette[7.1.9.2]{VI.eq7.1.9.2}
where the vertical functors are base-change functors, and the horizontal
functors in \eqhyperref[VI.eq7.1.9.1]{7.1.9.1} are respectively the functors
``morphism $\mapsto$ inverse image'' \hyperlink{sga4.exp.VI.equation.7.1.7.1}{(7.1.7.1)}, while those in
\eqhyperref[VI.eq7.1.9.2]{7.1.9.2} are the functors ``morphism $\to$ direct
image'' \hyperlink{sga4.exp.VI.equation.7.1.7.2}{(7.1.7.2)}.


% END INCLUDED SGA4_English_current_buildable_source_20260906/source/06/drafts/SGA4_VI_LINES_3873_4230_English.texfrag


% BEGIN INCLUDED SGA4_English_current_buildable_source_20260906/source/06/drafts/SGA4_VI_LINES_4231_4582_English.texfrag
\subsection{Fibered sites}\etiquette[7.2]{VI.subsec7.2}

\subsubsection{}\etiquette[7.2.1]{VI.subsubsec7.2.1}

A \emph{fibered $𝒰$-site} $C$ over a category $I$ is a fibered
category $p : C \to I$ whose fibers are endowed with topologies making
them $𝒰$-sites (IV \hyperref[II.def3.0.2]{3.0.2}), such that for every morphism
$f : i \to j$ of $I$, the inverse-image functor
$f^{\ast} : C_j \to C_i$ is a morphism from the site $C_i$ to the site
$C_j$ (IV \hyperref[IV.subsubsec4.9.1]{4.9.1}).

\subsubsection{}\etiquette[7.2.2]{VI.subsubsec7.2.2}

Let $p : C \to I$ and $q : D \to I$ be two fibered $𝒰$-sites over
$I$. A \emph{morphism of fibered $𝒰$-sites from $C$ to} $D$ is an
$I$-functor (\hyperref[VI.subsec6.0]{6.0}) $m : D \to C$ such that,
for every object $i$ of $I$, the functor $m_i : D_i \to C_i$ is a
morphism from the site $C_i$ to the site $D_i$. Denote by
$\Morsite_I(C, D)$ the set of morphisms of fibered sites from $C$ to
$D$. A \emph{continuous functor from the fibered site $D$ to the
fibered site} $C$ is an $I$-functor from $D$ to $C$ that induces
continuous functors on the fibers (III 1). Denote by
$\Cont_{I}(D, C)$ the set of continuous functors from the fibered site
$D$\pageoriginale to the fibered site $C$. Denote by
$\Mmorsite_I (C, D)$ the full subcategory of
$\Hhom_I (D, C)^{\circ}$ defined by the set of objects
$\Morsite_I (C, D) \subset \Hom_I (D, C)$. Likewise, denote by
$\Ccont_I(D, C)$ the full subcategory of $\Hhom_I (D, C)$ defined by
the subset of objects
$\Ccont_I (D, C) \subset \Hom_I (D, C)$. Thus we have two fully
faithful functors, injective on objects,
\begin{equation*}
\Mmorsite_I (C, D) \hookrightarrow \Ccont_I (D, C)^{\circ} \hookrightarrow
\Hhom_I (D, C)^{\circ}{\quoi}.\tag{7.2.2.1}
\end{equation*}\etiquette[7.2.2.1]{VI.eq7.2.2.1}
As in \hyperref[VI.subsubsec7.1.7]{7.1.7}, one defines
\emph{cartesian} morphisms from the fibered site $C$ to the fibered
site $D$. We have a diagram of full subcategories
\[\vcenter{
\xymatrix@M=.2cm{
\Mmorsite_{\Cart/I}(C,D)\ar@{^{(}->}[r] \ar@{_{(}->}[d]&
\Hhom_{\Cart/I}(D,C)^{\circ} \ar@{_{(}->}[d]\\
\Mmorsite_I(C,D)\ar@{^{(}->}[r]& \Hhom_I(D,C)^{\circ}.
}}\leqno{(7.2.2.2)}
\]\etiquette[7.2.2.2]{VI.eq7.2.2.2}

\subsubsection{}\etiquette[7.2.3]{VI.subsubsec7.2.3}

Morphisms of fibered $𝒰$-sites can be composed. This makes it
possible to define the category ($𝒱$-$𝒰$-Site/$I$), whose objects are
the $𝒰$-sites whose fibers belong to a universe $𝒱$ (of which $𝒰$ is
an element). The category $𝒱$-$𝒰$-site is moreover endowed with a
2-category structure; in particular, composition of morphisms of
fibered $𝒰$-sites is functorial with respect to morphisms between
morphisms. For every functor $\phi : J \to I$ and every fibered site
$p : C \to I$, the category $p_J : C_J \to J$ deduced from $(C, p)$
by the base change $\phi : J \to I$ is canonically endowed with a
fibered $𝒰$-site structure over $J$. The base-change operation is
2-functorial.

\subsubsection{}\etiquette[7.2.4]{VI.subsubsec7.2.4}

Let $p : C \to I$ be a fibered category whose fibers are endowed with
topologies making them $𝒰$-sites. Suppose that projective limits are
representable in the fiber categories. For $p : C \to I$ to be a
fibered $𝒰$-site, it is necessary that the inverse-image functors be
left exact and take covering families to covering families; this
condition is also sufficient when the topologies of the fibers are
coarser than the canonical topology (IV \hyperref[IV.subsubsec4.9.2]{4.9.2}). All the fibered sites
used in this Seminar will be of the type described above. Likewise,
let $p : C \to I$ and $q : D \to I$ be two fibered sites such that
finite projective limits are representable in the fiber categories.
\pageoriginale A cartesian functor $m : D \to C$ that induces left
exact functors on the fiber categories and takes covering families of
the fiber categories to covering families is a morphism of fibered
sites from $C$ to $D$; the converse holds if the fiber topologies of
$C$ are coarser than the canonical topology (IV \hyperref[IV.subsubsec4.9.2]{4.9.2}). All the
morphisms of fibered sites used in this Seminar will be of the type
described above.

\subsubsection{}\etiquette[7.2.5]{VI.subsubsec7.2.5}

A fibered $𝒰$-topos $p : F \to I$ is a fibered $𝒰$-site when its
fibers are endowed with the canonical topology, as we shall always do
from now on. Let $m: F \to G$ be a morphism of fibered $𝒰$-topoi. The
inverse-image functor $m^{\ast} : G \to F$
(\hyperref[VI.subsubsec7.1.7]{7.1.7}) is a morphism from the fibered
$𝒰$-site $F$ to the fibered $𝒰$-site $G$. Thus we have a functor
$(𝒱\text{-}𝒰\text{-}\Top_{/I}) \to
(𝒱\text{-}𝒰\text{-}\Site_{/I})$, which is not faithful in general.
But this functor extends naturally to a 2-functor which, for two
fibered $𝒰$-topoi $F$ and $G$ over $I$, induces a functor
\begin{equation*}
\Hhomtop_{I}(F, G) \to \Mmorsite_I (F, G)\tag{7.2.5.1}
\end{equation*}\etiquette[7.2.5.1]{VI.eq7.2.5.1}
that is an \emph{equivalence of categories}, as follows immediately
from the definitions. Moreover, the functor
\eqhyperref[VI.eq7.2.5.1]{7.2.5.1}, composed with the canonical
inclusion $\Mmorsite_I (F, G) \hookrightarrow
\Hhom_{I}(G, F)^{\circ}$, is precisely the functor \hyperlink{sga4.exp.VI.equation.7.1.7.1}{(7.1.7.1)} that
associates its inverse image with a morphism of fibered topoi.

\subsubsection{}\etiquette[7.2.6]{VI.subsubsec7.2.6}

Let $p : C \to I$ be a fibered $𝒰$-site. For every morphism
$f : i \to j$ of $I$, choose a base-change functor
$f^{\ast} : C_j \to C_{i}$. Passing to the categories of
$\sheaf{\Uu}$-sheaves, the functor $f^{\ast}$ extends to a functor
$\Top(f)^{\ast} : C^{\sim}_{j} \to C^{\sim}_{i}$, where
$\Top(f) : C^{\sim}_{i} \to C^{\sim}_{j}$ is a morphism of topoi (IV
\hyperref[IV.eq4.9.1.1]{4.9.1.1}), making the following diagram commutative up to isomorphism:
\[\vcenter{
\xymatrix@C=2.5cm@R=1cm{
C_j \ar[r]^{f^{\ast}} \ar[d]_{\varepsilon_j}& C_i\ar[d]^{\varepsilon_i}\\
C^{\widetilde{\phantom{a}}}_j \ar[r]^{\Top(f)^{\ast}}&
C^{\widetilde{\phantom{a}}}_i
}}\leqno{(7.2.6.1)}
\]\etiquette[7.2.6.1]{VI.eq7.2.6.1}
Let
\begin{equation}
c_{f, g} : f^{\ast} g^{\ast} \to (gf)^{\ast}\tag{7.2.6.2}
\end{equation}\etiquette[7.2.6.2]{VI.eq7.2.6.2}
be the\pageoriginale canonical isomorphisms
\eqhyperref[VI.eq6.1.2.1]{6.1.2.1}, and let
\begin{equation*}
b_f : \Top(f)^{\ast} \epsilon_j \to \epsilon_i f^{\ast}\tag{7.2.6.3}
\end{equation*}\etiquette[7.2.6.3]{VI.eq7.2.6.3}
be the isomorphisms that ``make'' the diagrams
\eqhyperref[VI.eq7.2.6.1]{7.2.6.1} commutative. Using the fact that the
functors $\Top(f)^{\ast}$ commute with inductive limits, one verifies
immediately that, for every pair of composable morphisms
$i \xrightarrow{f} j \xrightarrow{g} k$, there is a unique
isomorphism
\begin{equation*}
\Top(c_{f, g}) : \Top(f)^{\ast} \Top(g)^{\ast} \to \Top(gf)^{\ast}\tag{7.2.6.4}
\end{equation*}\etiquette[7.2.6.4]{VI.eq7.2.6.4}
such that
\begin{equation*}
\epsilon_i (c_{f, g})b_f \Top(f)^{\ast}(b_g) = b_{gf} \Top(c_{f,g})\tag{7.2.6.5}{\quoi}.
\end{equation*}\etiquette[7.2.6.5]{VI.eq7.2.6.5}
Moreover, the isomorphisms \eqhyperref[VI.eq7.2.6.4]{7.2.6.4}
automatically satisfy the cocycle condition of \hyperref[VI.eq6.1.2.2]{6.2.2.2}. We can
therefore construct a fibered category
\begin{equation*}
p^{\sim} : C^{\sim / I} \to I\tag{7.2.6.6}
\end{equation*}\etiquette[7.2.6.6]{VI.eq7.2.6.6}
whose fibers are canonically isomorphic to the $𝒰$-topoi
$C^{\sim}_i$ (\hyperref[VI.subsubsec6.1.3]{6.1.3}), and a cartesian
functor
\begin{equation*}
\epsilon_{C/I} : C \to C^{\sim /I}\tag{7.2.6.7}
\end{equation*}\etiquette[7.2.6.7]{VI.eq7.2.6.7}
that induces on the fibers the functors
$\epsilon_i : C_{i} \to C^{\sim}_i$. It follows immediately from the
definitions that $p^{\sim} : C^{\sim/I} \to I$ is a fibered
$𝒰$-topos over $I$, that $\epsilon_{C/I}$ is a cartesian morphism of
fibered $𝒰$-sites from $C^{\sim/I}$ to $C$, and that the fibered
$𝒰$-topos $C^{\sim /I} \to I$ and the cartesian functor
$\epsilon_{C/I}$ are independent, up to canonical isomorphism, of the
various choices used to construct them (the choices of the base-change
functors $f^{\ast}$ and of the extensions $\Top(f)^{\ast}$). The
fibered $𝒰$-topos $C^{\sim/I}$ is called the \emph{associated
fibered $𝒰$-topos} of the fibered $𝒰$-site.

\setcounter{paragraph}{7}
\paragraph{}\etiquette[7.2.6.8]{VI.par7.2.6.8}

One verifies that formation of the fibered $𝒰$-topos associated with
a fibered site ``commutes'' with base-category change operations.

\begin{enonce*}{Proposition 7.2.7}\etiquette[7.2.7]{VI.prop7.2.7}
Let\pageoriginale $E$ be a fibered $𝒰$-topos over $I$, and let $C$
be a fibered $𝒰$-site over $I$. The functor
$m \mapsto m^{\ast} \epsilon_{C/I}$, which associates with every
morphism of fibered topoi $m : E \to C^{\sim /I}$ the composite with
$\epsilon_{C/I}$ of the associated inverse-image functor
$m^{\ast} : C^{\sim / I} \to E$, induces an equivalence of categories
preserving cartesian objects
$$
\Hhomtop_I (E, C^{\sim / I}) \xrightarrow{\approx} \Mmorsite_I(E,
C){\quoi}.
$$
When finite projective limits are representable in the fiber
categories of $C$, the corresponding fully faithful functor
$$
\Hhomtop_{I} (E, C^{\sim/ I}) \to \Hhom_{I} (C, E)^{\circ}
$$
has as its essential image the set of $I$-functors $g : C \to E$ that
are left exact fiber by fiber and take covering families of the fiber
categories to epimorphic families of the corresponding fiber
categories.
\end{enonce*}

This proposition merely generalizes \hyperref[IV.prop4.9.4]{Proposition IV 4.9.4} to the
fibered context, and the proof given in IV \hyperref[IV.prop4.9.4]{4.9.4} generalizes
immediately.

\subsection{Example}\etiquette[7.3]{VI.subsec7.3}

\begin{enonce*}[remark]{Example 7.3.1}\etiquette[7.3.1]{VI.exem7.3.1}
(\emph{The fibered topos of arrows of a topos})

Let $E$ be a $𝒰$-topos, let $\text{Fl}(E)$ be the category of arrows
of $E$, and let $p : \text{Fl} (E) \to E$ be the functor that assigns
to an arrow its target. The category $(\text{Fl} (E), p)$ is a
fibered $𝒰$-topos over $E$ (\hyperref[VI.def7.1.1]{7.1.1} and IV
\hyperref[VI.thm5.1]{5.1}). The fiber over every object $X$ of $E$ is
the topos $E_{/X}$. The inverse-image functor
$f^{\ast} : E_{/Y} \to E_{/X}$ for a morphism $f : X \to Y$ is the
fiber-product functor. Thus, with every functor $\phi : I \to E$, one
can associate the fibered $𝒰$-topos over $I$
\begin{equation*}
p_I : \text{Fl}(E)_I \to I{\quoi},\tag{7.3.1.1}
\end{equation*}\etiquette[7.3.1.1]{VI.eq7.3.1.1}
obtained by the base change $\phi : I \to E$, whose fiber over every
object $i$ of $I$ is $E_{/\phi (i)}$
(\hyperref[VI.subsubsec7.1.9]{7.1.9}). If $𝒱$ is a universe of which
$E$ is an element, then, with the notation of
\hyperref[VI.subsubsec7.1.8]{7.1.8}, we have a map
\begin{equation*}
\Hom(I, E) \to \ob (𝒱\text{-}𝒰\text{-}\text{\rm top}_{/I}){\quoi},\tag{7.3.1.2}
\end{equation*}\etiquette[7.3.1.2]{VI.eq7.3.1.2}
which associates with every functor $\phi \in \Hom (I, E)$ the
fibered $𝒰$-topos over $I$ deduced from $(\text{Fl}(E), p)$ by the
\pageoriginale base change $\phi : I \to E$. Let $\phi_1$ and
$\phi_2$ be two functors from $I$ to $E$, and let
$m : \phi_1 \to \phi_2$ be a morphism of functors. Denote by
$p_1 : \text{Fl}(E)_1 \to I$ and $p_2 : \text{Fl} (E)_2 \to I$ the
fibered $𝒰$-topoi over $I$ deduced from
$p : \text{Fl}(E) \to E$ by the base changes $\phi_1 : I \to E$ and
$\phi_2 : I \to E$, respectively. For every object $i$ of $I$,
choose a morphism of topoi (IV \hyperref[IV.eq5.5.2]{5.5.2})
\begin{equation*}
\loc(m(i)) : E_{/{\phi}_{1}(i)} \to E_{/{\phi}_{2}(i)}{\quoi}.\tag{7.3.1.3}
\end{equation*}\etiquette[7.3.1.3]{VI.eq7.3.1.3}
One then verifies easily that there is a unique\footnote{Uniqueness
comes from requiring, with the notation of (IV \hyperref[IV.eq5.5.3]{5.5.3}), the formula
$(gf)_! = g_! f_!$.} cartesian morphism of fibered topoi over $I$,
\begin{equation*}
\loc(m) : \text{Fl}(E)_1 \to \text{Fl} (E)_2{\quoi},\tag{7.3.1.4}
\end{equation*}\etiquette[7.3.1.4]{VI.eq7.3.1.4}
that induces the morphisms $\loc(m(i))$ on the fiber topoi. The
morphism $\loc(m) : \text{Fl}(E)_1 \to \text{Fl} (E)_2$ is determined
up to canonical isomorphism by the morphism of functors
$m : \phi_1 \to \phi_2$. If $n : \phi_2 \to \phi_3$ is a morphism of
functors, we have an isomorphism
\begin{equation*}
c(m, n) : \loc (n) \loc(m) \simeq \loc(nm){\quoi},\tag{7.3.1.5}
\end{equation*}\etiquette[7.3.1.5]{VI.eq7.3.1.5}
and the morphisms $c(m, n)$ satisfy a cocycle property analogous to
\eqhyperref[VI.eq6.1.2.2]{6.1.2.2}.
\end{enonce*}

\begin{enonce*}[remark]{Example 7.3.2}\etiquette[7.3.2]{VI.exem7.3.2}
(The fibered site of arrows of a site).

Let $C$ be a $𝒰$-site in which \emph{fiber products are
representable}, let $\text{Fl}(C)$ be the category of arrows of $C$,
and let $p : \text{Fl}(C) \to C$ be the functor that assigns to an
arrow its target. The fibered category $(\text{Fl}(C), p)$ is a
fibered $𝒰$-site over $C$ when the fiber categories are endowed with
the induced topologies (II \hyperref[VI.cor5.2]{5.2}). Let
$\epsilon_C : C \to C^{\sim}$ be the canonical functor from $C$ to
the category of $𝒰$-sheaves on $C$. Denoting by
$\text{Fl}(\epsilon_C)$ the natural extension of the latter functor to
the arrow categories, we have a commutative diagram of categories and
functors:
\[\vcenter{
\xymatrix@C=2cm{
\text{Fl}(C) \ar[r]^{\text{Fl}(\varepsilon_C)}\ar[d]_{p}& \text{Fl}(C^{\widetilde{\phantom{a}}}) \ar[d]^{p}\\
C       \ar[r]^{\varepsilon_C} & C^{\widetilde{\phantom{a}}}
}}\leqno{(7.3.2.1)}
\]\etiquette[7.3.2.1]{VI.eq7.3.2.1}
and hence a\pageoriginale cartesian functor between categories
fibered over $C$:
\[\vcenter{
\xymatrix{
\text{Fl}(C) \ar[rr] \ar[rd]_{p}&     &
\text{Fl}(C^{\widetilde{\phantom{a}}})_C \ar[ld]^{p_{/C}}\\
                   & C     &\\
}}\leqno{(7.3.2.2)}
\]\etiquette[7.3.2.2]{VI.eq7.3.2.2}
where $p_{/C} : \text{Fl}(C^{\sim})_{C} \to C$ is deduced from
$\text{Fl}(C^{\sim}) \to C^{\sim}$ by the base change
$\epsilon_{C} : C \to C^{\sim}$. The functor
$\text{Fl}(C) \to \text{Fl}(C^{\sim})_{C}$ of
\eqhyperref[VI.eq7.3.2.2]{7.3.2.2} is a morphism from the fibered site
$\text{Fl}(C^{\sim})_{C}$ to the fibered site $\text{Fl}(C)$
(\hyperref[VI.subsubsec7.2.2]{7.2.2}); hence, by
\hyperref[VI.prop7.2.7]{7.2.7}, there is a morphism of fibered topoi
over $C$,
\begin{equation*}
\can : \text{Fl}(C^{\sim})_{C} \to \text{Fl} (C)^{\sim/C}{\quoi},\tag{7.3.2.3}
\end{equation*}\etiquette[7.3.2.3]{VI.eq7.3.2.3}
where $\text{Fl}(C)^{\sim/C}$ is the fibered topos over $C$
associated with the fibered site $\text{Fl}(C) \to C$
(\hyperref[VI.subsubsec7.2.6]{7.2.6}). It follows from the definition
of the associated fibered topos
(\hyperref[VI.subsubsec7.2.6]{7.2.6}) and from II
\hyperref[VI.cor5.5]{5.5} that the morphism can of
\eqhyperref[VI.eq7.3.2.3]{7.3.2.3} is a $C$-equivalence of fibered
$𝒰$-topoi over $C$, \ie there is a morphism
$\Theta : \text{Fl}(C)^{\sim/C} \to \text{Fl} (C^{\sim})_C$ of
fibered $𝒰$-topoi over $C$ such that the composites
$\Theta \circ \can$ and $\can \circ \Theta$ are $C$-isomorphic to the
identity.

Let $\phi : I \to C$ be a functor. By base change, we deduce a
fibered site over $I$,
\begin{equation*}
p_I : \text{Fl} (C)_I \to I\tag{7.3.2.4}
\end{equation*}\etiquette[7.3.2.4]{VI.eq7.3.2.4}
and, since formation of the associated fibered topos commutes with
base change (\hyperref[VI.subsubsec7.2.6]{7.2.6}), we have an
$I$-equivalence of fibered $𝒰$-topoi
\begin{equation*}
\can_I : \text{Fl} (C^{\sim})_{I} \to \text{Fl}(C)^{\sim/I}_{I}\tag{7.3.2.5}
\end{equation*}\etiquette[7.3.2.5]{VI.eq7.3.2.5}
where $\text{Fl}(C^{\sim})_{I} \to I$ is the fibered $𝒰$-topos
obtained by the base change
$\epsilon_C \circ \phi : I \to C^{\sim}$.

Finally, let $\phi_1 : I \to C$ and $\phi_2 : I \to C$ be two
functors, and let $m : \phi_1 \to \phi_2$ be a morphism of functors.
Denote by $\text{Fl}(C)_1$ and $\text{Fl}(C)_2$ the fibered sites over
$I$ deduced from $p : \text{Fl}(C) \to C$ by the base changes
$\phi_1$ and $\phi_2$, respectively. As in
\hyperref[VI.exem7.3.1]{7.3.1}, one constructs a morphism of fibered
sites over $I$,
\begin{equation*}
\loc(m) : \text{Fl} (C)_1 \to \text{Fl}(C)_2{\quoi}.\tag{7.3.2.6}
\end{equation*}\etiquette[7.3.2.6]{VI.eq7.3.2.6}
Denote\pageoriginale by $\text{Fl}(C^{\sim})_1 \to I$ and
$\text{Fl}(C^{\sim})_2 \to I$ the fibered $𝒰$-topoi over $I$
obtained by the base changes $\epsilon_C \circ \phi_1$ and
$\epsilon_C \circ \phi_2$, respectively. We have a diagram of
morphisms over $I$ that is commutative up to isomorphism:
\[\vcenter{
\xymatrix@C=2.5cm{
\text{Fl}(C^{\sim})_1 \ar[d]\ar[r]^{\loc(\varepsilon_C\ast m)}&
\text{Fl}(C^{\sim})_2 \ar[d]_{\bigdot}\\
\text{Fl}(C)_1       \ar[r]^{\loc(m)} & \text{Fl}(C)_2
}}\leqno{(7.3.2.7)}
\]\etiquette[7.3.2.7]{VI.eq7.3.2.7}
where the vertical morphisms are deduced by base change from the
morphisms of \eqhyperref[VI.eq7.3.2.2]{7.3.2.2}.
\end{enonce*}

\begin{enonce*}[remark]{Exercise 7.3.3}\etiquette[7.3.3]{VI.exe7.3.3}
(\emph{Small and big fibered sites}).

This exercise continues Exercise IV \hyperref[IV.exer4.10.6]{4.10.6}, whose hypotheses and
notation we use. Denote by $\MM$ the full subcategory of
$\text{Fl}(S)$ defined by the morphisms of $ℳ$, and let
$p : \MM \to S$ be the functor that assigns to a morphism its target.

$6^{\circ}$) Show that $(\MM, p)$ is a fibered site over $S$ and that,
when fiber products are representable in $S$, the inclusion functor
$\MM \to \text{Fl}(S)$ is a cartesian morphism from the fibered site
$\text{Fl} (S) \to S$ to the fibered site $\MM \to S$.

$7^{\circ}$) Let $\MM^{\sim/S} \to S$ be the fibered topos associated
with $\MM \to S$. Show that, with the notation of
\hyperref[VI.exem7.3.2]{7.3.2}, there is a canonical morphism
$g : \text{Fl} (S^{\sim})_{/S} \to \MM^{\sim/S}$ of fibered topoi
over $S$ which induces on each fiber over $X \in \ob (S)$ the
morphism $(\Prol_{X}, \Res_X)$ (IV \hyperref[IV.exer4.10.6]{4.10.6} $4^{\circ})$). Show that,
although for every $X \in \ob (S)$ the morphism
$g_X : S^{\sim}_{/X} \to S(X)^{\sim}$ admits a section (\loccit), the
morphism $g$ does not in general admit a section.
\end{enonce*}

% END INCLUDED SGA4_English_current_buildable_source_20260906/source/06/drafts/SGA4_VI_LINES_4231_4582_English.texfrag


% BEGIN INCLUDED SGA4_English_current_buildable_source_20260906/source/06/drafts/SGA4_VI_LINES_4583_4929_English.texfrag
\subsection{The total topology of a fibered topos or site}\etiquette[7.4]{VI.subsec7.4}

\begin{enonce*}{Definition 7.4.1}\etiquette[7.4.1]{VI.def7.4.1}
Let $p:C\to I$ be a fibered $𝒰$-site and, for every
$i\in\ob(I)$, denote by $\alpha_{i!}:C_i\to C$ the inclusion functor
of the fiber category $C_i$ into $C$. The \emph{total topology} on
$C$ is the coarsest topology on $C$ for which the functors
$\alpha_{i!}$, $i\in\ob(I)$, are continuous \emph{(III \hyperref[III.prop3.6]{3.6})}. The
category $C$ equipped with the total topology is called the
\emph{total site}.
\end{enonce*}

\begin{enonce*}{Proposition 7.4.2}\etiquette[7.4.2]{VI.prop7.4.2}
Let\pageoriginale $p:C\to I$ be a fibered $𝒰$-site, let $T$ be the
total topology on $C$, and, for every $i\in\ob(I)$, let $T_i$ be the
topology of the fiber $C_i$.
\begin{enumerate}
\renewcommand{\labelenumi}{\theenumi)}
\item Let $X$ be an object of $C$ over an object $i$ of $I$. A family
  $(X_\beta\to X)_{\beta\in B}$ is covering for $T$ if and only if it
  is refined by a covering family for $T_i$,
  $(Y_{\gamma}\to X)_{\gamma\in\Gamma}$, in $C_i$.

\item The topology $T$ induces the topology $T_i$ on each category
  $C_i$.

\item For every $i\in\ob(I)$, the functor $\alpha_{i!}:C_i\to C$ is
  cocontinuous.
\end{enumerate}
\end{enonce*}

\setcounter{subsubsection}{2}
\setcounter{paragraph}{0}
\paragraph{}\etiquette[7.4.2.1]{VI.par7.4.2.1}

Property 3) follows from 1) and III \hyperref[III.def2.1]{2.1}. Let us prove 1). For every
object $X$ of $C$, let $J(X)$ be the set of sieves of $C/X$ described
in 1). One verifies that the $J(X)$ have properties $T$ 1), $T$ 2),
and $T$ 3) of II \hyperref[II.def1.1]{1.1}, and hence that they define a topology $T'$ on
$C$. The topology $T'$ is the coarsest topology on $C$ for which the
covering families for the topologies $T_i$ are covering. Thus $T'$ is
coarser than $T$ (III \hyperref[III.prop1.6]{1.6}). To prove that $T'=T$, it is therefore
enough to show that, for every object $i$ of $I$, the functor
$\alpha_{i!}:C_i\to C$ is continuous when $C$ is equipped with the
topology $T'$; equivalently, it is enough to show that the topology
induced by $T'$ on the category $C_i$ is $T_i$, which will prove
property 2) at the same time. Let $T'_i$ be the topology on $C_i$
induced by $T'$, and let $u:R\hookrightarrow X$ be a covering sieve
for $T'_i$. By III \hyperref[III.prop3.2]{3.2}, the morphism
$\alpha_{i!}(u):\alpha_{i!}(R)\to\alpha_{i!}(X)$ is bicovering and in
particular covering. By the definition of $T'$, this means that the
sieve $R\hookrightarrow X$ is covering for $T_i$. Thus $T'_i$ is
coarser than $T_i$. To prove that $T'_i$ is finer than $T_i$, by the
definition of the induced topology (III \hyperref[III.subsec3.1]{3.1}), it is enough to show
that the functor $\alpha_{i!}:C_i\to C$ is continuous when $C_i$ and
$C$ are equipped with $T_i$ and $T'$, respectively. In summary, by
III \hyperref[III.prop1.2]{1.2} and \hyperref[III.prop1.5]{1.5}, to prove properties 1) and 2), it is enough to prove
the following lemma:

\begin{enonce*}{Lemma 7.4.2.2}\etiquette[7.4.2.2]{VI.lem7.4.2.2}
Under the hypotheses and with the notation of
\hyperref[VI.prop7.4.2]{7.4.2} and
\hyperref[VI.par7.4.2.1]{7.4.2.1}, let $X$ be an object of $C$ over
$i$, and let $R\xrightarrow{u}X$ be a morphism of
$C^{\widehat{\phantom{,}}}_i$. The following properties are
equivalent:
\begin{enumerate}
\renewcommand{\theenumi}{\roman{enumi}}
\renewcommand{\labelenumi}{\theenumi)}
\item The morphism\pageoriginale $R\to X$ is bicovering for $T_i$.

\item The morphism
  $\alpha_{i!}(u):\alpha_{i!}(R)\to\alpha_{i!}(X)$ is bicovering for
  $T'$.
\end{enumerate}
\end{enonce*}

i) $\Rightarrow$ ii): It is clear that if $R\to X$ is bicovering for
$T_i$, then the morphism $\alpha_{i!}(u)$ is covering for $T'$. It
remains to show that it is bicovering (II
\hyperref[VI.cor5.2]{5.2}). Let $Y$ be an object of $C$ over an object
$j$ of $I$, let $m,n:Y\rightrightarrows\alpha_{i!}(R)$ be two
morphisms of $C^{\widehat{\phantom{,}}}$ such that
$\alpha_{i!}(u)m=\alpha_{i!}(u)n$, and set
$f=p(\alpha_{i!}(u)m)$. Since the functor $\alpha_{i!}$ commutes with
inductive limits (I \hyperlink{sga4.exp.I.proposition.5.4.item.3}{5.4.3}), we have (I \hyperref[I.prop3.4]{3.4}):
$$
\Hom_{C^{\widehat{\phantom{,}}}} (Y, \alpha_{i!} (R)) \simeq
\varinjlim_{C_{i}/R} \Hom_C (Y, \alpha_{i!} (.)){\quoi}.
$$
For every object $Z$ of $C_i$, with the notation of
\hyperref[VI.subsubsec6.1.1]{6.1.1}, we have
$$
\Hom_{C} (Y, Z) \simeq \coprod_{w \in \Hom(j, i)} \Hom_w (Y,
Z){\quoi},
$$
and therefore
$$
\varinjlim_{C_{i}/R} \Hom (Y, \alpha_{i!} (.)) \simeq
\coprod_{w \in \Hom(j, i)} \varinjlim_{C_{i}/R}
\Hom_{w} (Y, \alpha_{i!} (.)){\quoi}.
$$
For every $w\in\Hom(j,i)$, however, there is an inverse image functor
$w^{\ast}:C_i\to C_j$, and
$\Hom_w(Y,\alpha_{i!}(.))\simeq\Hom_{C_j}(Y,w^{\ast}(.))$. Thus,
still denoting by $w^{\ast}:C^{\widehat{\phantom{,}}}_i\to
C^{\widehat{\phantom{,}}}_j$ the extension of $w^{\ast}$ to
presheaves (which, according to I \hyperref[VI.lem5.4]{5.4}, should be denoted
$(w^{\ast})_{!}$), and observing that the functor
$w^{\ast}:C^{\widehat{\phantom{,}}}_i\to
C^{\widehat{\phantom{,}}}_j$ commutes with inductive limits (I
\hyperlink{sga4.exp.I.proposition.5.4.item.3}{5.4.3}), we obtain an isomorphism
$$
\Hom_{C^{\widehat{\phantom{,}}}} (Y, \alpha_{i!} (R)) =
\coprod_{w \in \Hom(j, i)}
\Hom_{C^{\widehat{\phantom{,}}}_{j}}(Y,w^{\ast}(R)){\quoi}.
$$
Under this isomorphism, the morphisms
$m,n:Y\rightrightarrows\alpha_{i!}(R)$ correspond to two morphisms
$m',n':Y\rightrightarrows f^{\ast}R$ such that
$f^{\ast}(u)m'=f^{\ast}(u)n'$. Since the inverse image functor is a
morphism of topoi, it is in particular continuous, and consequently
$f^{\ast}(u)$ is bicovering (III \hyperref[III.prop3.2]{3.2}). There is therefore a covering
family $(Y_{\delta}\xrightarrow{v_{\delta}}Y)_{\delta}$ in $C_j$
such that $m'v_{\delta}=n'v_{\delta}$ for every $\delta$. It follows
that $mv_{\delta}=nv_{\delta}$ for every $\delta$, and hence that
$\alpha_{i!}(u)$ is bicovering for $T'$.

ii) $\Rightarrow$ i):\pageoriginale Suppose that
$\alpha_{i!}(R)\xrightarrow{\alpha_{i!}(u)}\alpha_{i!}(X)$ is
bicovering. The functor $\alpha_{i!}:C_i\to C$ is cocontinuous (III
\hyperref[III.def2.1]{2.1}). Hence $\alpha^{\ast}_{i}\alpha_{i!}(u)$ is bicovering (III
\hyperref[III.eq2.3.2]{2.3.2} and II \hyperref[VI.cor5.3]{5.3} ii)). The functor
$\alpha_{i!}:C_i\to C$ has property (PPF) of I \hyperref[I.prop5.14]{5.14}. We therefore
have a Cartesian diagram in $C^{\widehat{\phantom{,}}}_i$ (I \hyperref[I.prop5.14]{5.14}
3)):
\[
\xymatrix@C=2.25cm@R=1.25cm{
R  \ar[r] \ar[d]_{u}& \alpha_i^{\ast}\alpha_{i!}(R) \ar[d]^{\alpha_i^{\ast}\alpha_{i!}(u)}\\
X     \ar[r]          & \alpha_i^{\ast}\alpha_{i!}(X).
}
\]
It follows that $u$ is bicovering (II \hyperref[VI.cor5.2]{5.2}).

\begin{enonce*}[remark]{Remark 7.4.3}\etiquette[7.4.3]{VI.rem7.4.3}
\begin{enumerate}
\renewcommand{\labelenumi}{\theenumi)}
\item Proposition \hyperref[VI.prop7.4.2]{7.4.2} can be stated and
  proved in a slightly more general setting than that of fibered
  sites, but apparently without useful application: instead of a
  fibered site $p:C\to I$, one considers a fibered category
  $p:C\to I$ whose fibers are equipped with topologies and whose
  inverse image functors are continuous for these topologies (whereas
  in the case of fibered sites one requires in addition that these
  inverse image functors be morphisms of sites (IV \hyperref[IV.subsec4.9]{4.9})).

\item One can prove that the total topology of a fibered site is the
  finest topology making the functors $\alpha_{i!}:C_i\to C$
  cocontinuous (III 2).

\hypertarget{sga4.exp.VI.remark.7.4.3.item.3}{}\item Let $p:C\to I$ be a fibered $𝒰$-site over a category
  equivalent to a $𝒰$-small category. Then the total site $C$ is a
  $𝒰$-site (II \hyperref[II.def3.0.2]{3.0.2}). Indeed, if, for every object $i$ of $I$,
  $(X_{\beta})_{\beta\in B_i}$ is a small topologically generating
  family of $C_i$, then the family
  $(X_{\beta})_{\beta\in\amalg_iB_i}$ is topologically generating for
  $C$ and is $𝒰$-small. The topos of sheaves on the total site is
  called the \emph{total topos} and is denoted by $\Top(C)$.

\hypertarget{sga4.exp.VI.remark.7.4.3.item.4}{}\item Let $p:C\to I$ be a fibered $𝒰$-site over a small category.
  Since the functor $\alpha_{i!}:C_i\to C$ is cocontinuous, it gives
  rise to a morphism of topoi
  $\alpha_i=(\alpha^{\ast}_i,\alpha^{i\ast})$ from $C^{\sim}_i$ to
  $\Top(C)$ (IV \hyperref[IV.subsec4.7]{4.7}). Since, moreover, the functor
  $\alpha_{i!}:C_i\to C$ is continuous, the functor
  $\alpha^{\ast}_i:\Top(C)\to C^{\sim}_i$ admits a left adjoint,
  again denoted by\pageoriginale
  $\alpha_{i!}:C^{\sim}_i\to\Top(C)$, which extends the functor
  $\alpha_{i!}:C_i\to C$ (III \hyperref[III.prop1.2]{1.2} iv)).
\end{enumerate}
\end{enonce*}

\begin{enonce*}{Proposition 7.4.4}\etiquette[7.4.4]{VI.prop7.4.4}
Let $p:C\to I$ be a fibered $𝒰$-site over a small category $I$.
\begin{enumerate}
\renewcommand{\labelenumi}{\theenumi)}
\item A presheaf $F$ on $C$ is a sheaf for the total topology if and
  only if, for every object $i$ of $I$, the presheaf
  $F\circ\alpha_{i!}=F|C_i$ is a sheaf on $C_i$.

\item A functor $m$ from $C$ to a site $D$ is a continuous functor
  from the total site $C$ to the site $D$ if and only if, for every
  object $i$ of $I$, the functor
  $m|C_i=m\circ\alpha_{i!}:C_i\to D$ is continuous.
\end{enumerate}
\end{enonce*}

If $F$ is a sheaf for the total topology, the presheaf
$F\circ\alpha_{i!}$ is a sheaf on $C_i$ because
$\alpha_{i!}:C_i\to C$ is continuous. Conversely, suppose that
$F\circ\alpha_{i!}$ is a sheaf on $C_i$ for every $i$. Let $X$ be an
object of $C$ over an object $i$ of $I$. For every covering sieve $R$
of $C_i/X$, we have $F(X)\simeq F\circ\alpha_{i!}(R)$ and hence
$F(X)\simeq F(\alpha_{i!}(R))$. Let $R'\hookrightarrow X$ be the
sieve of $C/X$ that is the image of the canonical morphism of
presheaves $\alpha_{i!}(R)\to X$. The morphism
$\alpha_{i!}(R)\to R'$ is an epimorphism of presheaves, and therefore
$F(R')\to F(\alpha_{i!}(R))$ is injective. Since
$F(R')\to F(\alpha_{i!}(R))$ is also surjective, it is bijective, and
hence $F(X)\to F(R')$ is bijective. Since the covering sieves of $X$
are described by \eqhyperref[VI.par7.4.2.1]{7.4.2.1}, the presheaf
$F$ is a sheaf: returning to the construction of the associated
sheaf (II 3), one sees that $F$ is isomorphic to its associated sheaf
(II \hyperref[II.rem3.3]{3.3}). The second assertion follows immediately from the first and
the definition of continuity (III \hyperref[III.def1.1]{1.1}).

\setcounter{subsubsection}{4}
\subsubsection{}\etiquette[7.4.5]{VI.subsubsec7.4.5}

Let $p:C\to I$ be a fibered $𝒰$-site over a category $I$ equivalent
to a small category, and let
$$
f : i \to j
$$
be an arrow of the index category $I$. Again denoting by
$f:C^{\sim}_i\to C^{\sim}_j$ the morphism of topoi determined by the
fibered-site structure, we have (with the notation of \hyperlink{sga4.exp.VI.remark.7.4.3.item.4}{7.4.3.4}) a
diagram of morphisms of $𝒰$-topoi
\[\vcenter{
\xymatrix@C=2.25cm@R=1.25cm{
C^{\sim}_i \ar[d]_{f} \ar[r]^{\alpha_i} & \Top(C)\\
C^{\sim}_j  \ar[ru]_{\alpha_j}&
}}\leqno{(7.4.5.1)}
\]\etiquette[7.4.5.1]{VI.eq7.4.5.1}
where $\Top(C)$\pageoriginale is the total topos \hyperlink{sga4.exp.VI.remark.7.4.3.item.3}{(7.4.3.3)}. The
diagram \eqhyperref[VI.eq7.4.5.1]{7.4.5.1} \emph{is not commutative in
general}. We shall define a canonical morphism of morphisms of topoi:
\begin{equation*}
\rho_f : \alpha_i \to \alpha_j \circ f{\quoi}.\tag{7.4.5.2}
\end{equation*}\etiquette[7.4.5.2]{VI.eq7.4.5.2}
It is enough to define such a morphism at the level of inverse images,
\ie to define a morphism of functors
\begin{equation*}
\beta_f^{\ast} : f^{\ast}\circ \alpha_j^{\ast} \to
\alpha_i^{\ast}{\quoi}, \tag{7.4.5.3}
\end{equation*}\etiquette[7.4.5.3]{VI.eq7.4.5.3}
or, equivalently, using the adjunction between $f^{\ast}$ and
$f_{\ast}$, it is enough to define a morphism of functors
\begin{equation*}
\gamma_f : \alpha^{\ast}_j \to f_{\ast} \alpha^{\ast}_i{\quoi}.\tag{7.4.5.4}
\end{equation*}\etiquette[7.4.5.4]{VI.eq7.4.5.4}
Let $Y$ be an object of $C_j$ and $F$ a sheaf on the total site $C$.
By definition,
$\alpha^{\ast}_jF(Y)=F(\alpha_{j!}(Y))$ and
$f_{\ast}\alpha^{\ast}_iF(Y)=F(\alpha_{i!}f^{\ast}(Y))$, where in the
second member $f^{\ast}:C_j\to C_i$ denotes the inverse image functor
for the fibered structure. By the definition of this inverse image
functor, however, there is a canonical Cartesian morphism
\begin{equation*}
m_Y : \alpha_{i!} f^{\ast} \to \alpha_{j!}{\quoi};\tag{7.4.5.5}
\end{equation*}\etiquette[7.4.5.5]{VI.eq7.4.5.5}
hence, on applying the functor $F$, a map functorial in $Y$:
\begin{equation*}
\gamma_f (F) (Y) = F(m_Y) : \alpha_{j}^{\ast} F(Y) \to f_{\ast}
\alpha^{\ast}_i F(Y){\quoi},\tag{7.4.5.6}
\end{equation*}\etiquette[7.4.5.6]{VI.eq7.4.5.6}
which defines a morphism of sheaves on $C_j$:
\begin{equation*}
\gamma_f (F) : \alpha_{j}^{\ast} F \to f_{\ast} \alpha^{\ast}_i
F{\quoi};\tag{7.4.5.7}
\end{equation*}\etiquette[7.4.5.7]{VI.eq7.4.5.7}
whence the morphism of functors $\gamma_f$ of
\eqhyperref[VI.eq7.4.5.4]{7.4.5.4}.

Now let\pageoriginale $f:i\to j$ and $g:j\to k$ be two composable
morphisms of $I$. There is a canonical isomorphism
$$
c'_{g, f} : g_{\ast} f_{\ast} \to (gf)_{\ast}
$$
and one immediately verifies the compatibility formula
\begin{equation*}
c'_{g, f} \circ g_{\ast} (\gamma_f) \circ \gamma_g =
\gamma_{gf}{\quoi}.\tag{7.4.5.8}
\end{equation*}\etiquette[7.4.5.8]{VI.eq7.4.5.8}

\subsubsection{}\etiquette[7.4.6]{VI.subsubsec7.4.6}

Introduce the fibered $𝒰$-topos
$p^{\sim/I}:C^{\sim/I}\to I$ associated with $p:C\to I$
(\hyperref[VI.subsubsec7.2.6]{7.2.6}) and the corresponding category
fibered over $I^{\circ}$,
$(p^{\sim/I})':(C^{\sim/I})'\to I^{\circ}$
\eqhyperref[VI.eq7.1.3.2]{7.1.3.2}. The preceding considerations
associate with every object $F$ of $\Top(C)$ a family
$F_i=(\alpha^{\ast}_iF)_{i\in\ob(I)}$ of objects of the
$C^{\sim}_i$, together with a family of morphisms
$$
\gamma_f F : F_j \to f_{\ast} F_i,\quad f\in F1(I),\quad f : i \to j{\quoi},
$$
subject to the compatibility conditions of
\eqhyperref[VI.eq7.4.5.8]{7.4.5.8}. Thus with $F$ we have associated
an $I^{\circ}$-functor from the category $I^{\circ}$ to the category
$(C^{\sim/I})'$, \ie an object $\Theta_C(F)$ of the category
$\Hhom_{I^{\circ}}(I^{\circ},(C^{\sim/I})')$. Since this construction
is functorial in $F$, we obtain a functor
\begin{equation*}
\Theta_C : \Top(C) \to \Hhom_{I^{\circ}} (I^{\circ},
(C^{\sim/I})'){\quoi}.\tag{7.4.6.1}
\end{equation*}\etiquette[7.4.6.1]{VI.eq7.4.6.1}

\begin{enonce*}{Proposition 7.4.7}\etiquette[7.4.7]{VI.prop7.4.7}
Let $C\to I$ be a fibered $\mathcal{U}$-site, with $I$ equivalent to
a small category. The functor $\Theta_C$
\eqhyperref[VI.eq7.4.6.1]{7.4.6.1} is an equivalence of categories.
\end{enonce*}

We shall content ourselves with describing a quasi-inverse to
$\Theta_C$. Let \\
$\sigma\in\Hom_{I^{\circ}}(I^{\circ},(C^{\sim/I})')$, \ie an
$I^{\circ}$-functor $i\mapsto\sigma(i)$ from $I^{\circ}$ to
$(C^{\sim/I})'$. For every object $X$ of $C$ over
$p(X)=i\in\ob(I)$, there is a sheaf
$\sigma(p(X))\in\ob C^{\sim}_i$, and hence a set
$\sigma(p(X))(X)$. Let $m:X\to Y$ be a morphism of $C$ over the
morphism $f:i\to j$ of $I$. The morphism $m$ factors uniquely as a
morphism $m':X\to f^{\ast}Y$ of $C_i$ followed by the canonical
Cartesian morphism $f^{\ast}Y\to Y$. Likewise, the morphism
$\sigma(f)$ factors uniquely as a morphism
$\gamma_f(\sigma):\sigma(p(Y))\to f_{\ast}\sigma(p(X))$ followed by
the canonical $I^{\circ}$-Cartesian morphism
$f_{\ast}\sigma(p(X))\to\sigma(p(X))$. We therefore obtain a map
from $\sigma(p(Y))(Y)$ to $\sigma(p(X))(X)$,\pageoriginale by
composing the maps
$$
\sigma(p(Y))(Y) \xrightarrow{\gamma_f(\sigma)(Y)}
f_{\ast}\sigma(p(X))(Y) =
\sigma(p(X))(f^{\ast}Y) \xrightarrow{\sigma(p(X))(m')}
\sigma(p(X))(X){\quoi}.
$$
Thus with every object $X$ of $C$ we have associated a set
$\sigma(p(X))(X)$, and with every morphism $m:X\to Y$ a map
$\sigma(p(Y))(Y)\to\sigma(p(X))(X)$. One verifies that this indeed
defines a presheaf on $C$, and that this presheaf does not depend on
the choice of inverse image functors used in its construction. It
follows from \hyperref[VI.prop7.4.4]{7.4.4} that the presheaf
$X\mapsto\sigma(p(X))(X)$ is a sheaf on the total site $C$, and it is
clear that this sheaf depends functorially on the object $\sigma$ of
$\Hhom_{I^{\circ}}(I^{\circ},(C^{\sim/I})')$. Returning to the
definition of $\Theta_C$, it is also clear that the functor from
$\Hhom_{I^{\circ}}(I^{\circ},(C^{\sim/I})')$ to $\Top(C)$ just
constructed is a quasi-inverse to $\Theta_C$.

\begin{enonce*}[remark]{Remark 7.4.8}\etiquette[7.4.8]{VI.rem7.4.8}
It follows in particular from \hyperref[VI.prop7.4.7]{7.4.7} that the
category $\Hhom_{I^{\circ}}(I^{\circ},(C^{\sim/I})')$ is a
$𝒰$-topos. This latter fact can be seen directly. Using I \hyperref[I.lem9.21.10]{9.21.10},
one sees immediately that conditions a), b), and c) of IV \hyperref[IV.subsubsec1.1.2]{1.1.2} are
satisfied, and the existence of a small generating category follows
from I \hyperref[I.cor9.25]{9.25}.
\end{enonce*}

\setcounter{subsubsection}{8}

% END INCLUDED SGA4_English_current_buildable_source_20260906/source/06/drafts/SGA4_VI_LINES_4583_4929_English.texfrag


% BEGIN INCLUDED SGA4_English_current_buildable_source_20260906/source/06/drafts/SGA4_VI_LINES_4930_5251_English.texfrag
\subsubsection{}\etiquette[7.4.9]{VI.subsubsec7.4.9}

Let $p:C\to I$ and $q:D\to I$ be two $𝒰$-sites fibered over a small
category, and let $m$ be a morphism from the fibered site $C$ to the fibered
site $D$ (\hyperref[VI.subsubsec7.2.2]{7.2.2}). Denote by
$m^{\sim/I}:C^{\sim/I}\to D^{\sim/I}$ the corresponding morphism between
the associated fibered $𝒰$-topoi. The morphism of fibered topoi
$m^{\sim/I}$ is defined up to canonical isomorphism
(\hyperref[VI.prop7.2.7]{7.2.7}). It has a corresponding direct-image
functor $m^{\sim/I}_{\ast}:(C^{\sim/I})'\to(D^{\sim/I})'$
(\hyperref[VI.subsubsec7.1.7]{7.1.7}), which, on passing to sections over
$I^{\circ}$, gives a functor
\begin{equation*}
\Hom_{I^{\circ}} (I^{\circ}, m^{\sim/I}_{\ast}) : \Hhom_{I^{\circ}}
(I^{\circ}, (C^{\sim/I})') \to \Hhom_{I^{\circ}} (I^{\circ},
(D^{\sim/I})'){\quoi}.\tag{7.4.9.1}
\end{equation*}\etiquette[7.4.9.1]{VI.eq7.4.9.1}

Moreover, the functor $m$ is a continuous functor between the total sites
(\hyperref[VI.prop7.4.4]{7.4.4}), and hence, by composition, gives a functor
$m_{\ast}:\Top(C)\to\Top(F)$.

In summary, there is a diagram of functors
\[\vcenter{
\xymatrix@C=2.75cm@R=1.5cm{
\Top(C) \ar[r]^{\theta_C} \ar[d]_{m_{\ast}}& \Hhom_I^{\circ}(I^{\circ},(C^{\sim/I})') \ar[d]_{\Hom_{I^\circ}(I^{\circ},m_{\ast}^{\sim/I})}\\
\Top(D)  \ar[r]^{\theta_D} & \Hhom_I^{\circ}(I^{\circ},(D^{\sim/I})')
}}\leqno{(7.4.9.2)}
\]\etiquette[7.4.9.2]{VI.eq7.4.9.2}

\begin{enonce*}{Proposition 7.4.10}\etiquette[7.4.10]{VI.prop7.4.10}
The diagram\pageoriginale \eqhyperref[VI.eq7.4.9.2]{7.4.9.2} is commutative
up to isomorphism. When $m$ is Cartesian, $m$ is a morphism from the total
site $C$ to the total site $D$.
\end{enonce*}

We leave it to the reader to verify that diagram
\eqhyperref[VI.eq7.4.9.2]{7.4.9.2} is commutative up to isomorphism. To show
that $m$ is a morphism between the total sites, it is enough, in view of
\hyperref[VI.prop7.4.7]{7.4.7}, to show that the functor left adjoint to
$\Hom_{I^{\circ}}(I^{\circ},m^{\sim/I}_{\ast})$ is left exact. We first
describe the inverse-image functor of the morphism $m^{\sim/I}$
(\hyperref[VI.prop7.2.7]{7.2.7}); for every object $i$ of $I$, denote by
$m^{\ast}_i:D^{\sim}_i\to C^{\sim}_i$ the functor induced by
$(m^{\sim/i})^{\ast}$ on the fibers at $i$. For every morphism $f:i\to j$
of $I$, choose morphisms of topoi, denoted in both cases by
$(f^{\ast},f_{\ast}):C^{\sim}_i\to C^{\sim}_j$ and
$(f^{\ast},f_{\ast}):D^{\sim}_i\to D^{\sim}_j$. Since
$(m^{\sim/I})^{\ast}$ is a Cartesian $I^{\circ}$-functor, for every
$f:i\to j$ there is a canonical isomorphism
\begin{equation*}
m^{\ast}_i f_{\ast} \to f^{\ast} m^{\ast}_j{\quoi},\tag{7.4.10.1}
\end{equation*}\etiquette[7.4.10.1]{VI.eq7.4.10.1}
and since the functors $f^{\ast}$ and $f_{\ast}$ are adjoint, there are
canonical morphisms (adjunction morphisms)
\begin{align*}
& \phi : \id \to f_{\ast} f^{\ast}\\
& \psi : f^{\ast} f_{\ast} \to \id{\quoi}. \tag{7.4.10.2}
\end{align*}\etiquette[7.4.10.2]{VI.eq7.4.10.2}
Now consider the following sequence of functorial morphisms:
\begin{equation*}
m^{\ast}_j f^{\ast} \xrightarrow{(1)} f_{\ast} f^{\ast} m^{\ast}_j
f_{\ast} \xrightarrow{(2)} f_{\ast} m^{\ast}_i f^{\ast} f_{\ast}
\xrightarrow{(3)} f_{\ast} m^{\ast}_i{\quoi},\tag{7.4.10.3}
\end{equation*}\etiquette[7.4.10.3]{VI.eq7.4.10.3}
where morphism (1) is deduced from $\phi$, morphism (2) is deduced from the
isomorphism \eqhyperref[VI.eq7.4.10.1]{7.4.10.1}), and morphism (3 is
deduced from $\psi$. By composing the various morphisms\pageoriginale in
the sequence \eqhyperref[VI.eq7.4.10.3]{7.4.10.3}, we obtain a functorial
morphism
\begin{equation*}
\can_f : m^{\ast}_j f_{\ast} \to f_{\ast} m^{\ast}_i{\quoi}.\tag{7.4.10.4}
\end{equation*}\etiquette[7.4.10.4]{VI.eq7.4.10.4}
Now let $\tau$ be an object of
$\Hhom_{I^{\circ}}(I^{\circ}(D^{\sim/I})')$. Thus, for every object $i$ of
$I$, there is an object $\tau(i)$ of $D^{\sim}_i$, and for every morphism
$f:i\to j$, a morphism
\begin{equation*}
\gamma_{f} (\tau) : \tau(j) \to f_{\ast} \tau(i){\quoi},\tag{7.4.10.5}
\end{equation*}\etiquette[7.4.10.5]{VI.eq7.4.10.5}
where the family of the $\gamma_f(\tau)$ has a compatibility property
analogous to \eqhyperref[VI.eq7.4.5.8]{7.4.5.8}). Put
$\sigma(i=m^{\ast}_i(\tau(i))$ and denote by
\begin{equation*}
\gamma_{f} (\sigma) : \sigma(j) \to f_{\ast} \sigma(i)\tag{7.4.10.6}
\end{equation*}\etiquette[7.4.10.6]{VI.eq7.4.10.6}
the composite morphism
$m^{\ast}_j\tau(j)\xrightarrow{m_j^{\ast}\gamma_f(\tau)}
m_j^{\ast}f_{\ast}\tau(i)\xrightarrow{\can_f(\tau(i))}
f_{\ast}m^{\ast}_i\tau(i)$. One verifies that the $\gamma_f(\sigma)$ have
the compatibility property of \eqhyperref[VI.eq7.4.5.8]{7.4.5.8}, and hence
that an object of $\Hhom_{I^{\circ}}(I^{\circ},(C^{\sim/I})')$ has thus been
defined. An ``adjoint functor chasing'' then shows that the functor from
$\Hhom_{I^{\circ}}(I^{\circ},(D^{\sim/I})')$ to
$\Hhom_{I^{\circ}}(I^{\circ},(C^{\sim/I})')$ just constructed is left
adjoint to the functor
$\Hom_{I^{\circ}}(I^{\circ},m^{\sim/I}_{\ast})$. It remains to show that
this left adjoint commutes with finite projective limits, which follows
immediately from the fact that projective limits in the categories of
sections are computed fiber by fiber, and that all the functors used to
construct this left adjoint commute with finite projective limits.

\begin{enonce*}{Corollary 7.4.11}\etiquette[7.4.11]{VI.cor7.4.11}
Let $p:C\to I$ be a $𝒰$-site fibered over a small category, let
$p^{\sim/I}:C^{\sim/I}\to I$ be the associated fibered $𝒰$-topos
(\hyperref[VI.subsubsec7.2.6]{7.2.6}), and let
$\epsilon_I:C\to C^{\sim/I}$ be the canonical morphism from the fibered
$𝒰$-site $C^{\sim/I}$ to the fibered $𝒰$-site $C$. The morphism
$\epsilon_I$ induces an equivalence
$\Top(\epsilon_I):\Top(C)\to\Top(C^{\sim/I})$ between the corresponding
total topoi.
\end{enonce*}

This follows from the commutativity of diagram
\eqhyperref[VI.eq7.4.9.2]{7.4.9.2} and from the fact that the morphism
$\epsilon_I$ gives an equivalence between the associated fibered
$𝒰$-topoi.

\begin{enonce*}[remark]{Remark 7.4.11.1}\etiquette[7.4.11.1]{VI.rem7.4.11.1}
It follows from \hyperref[VI.cor7.4.11]{7.4.11} that in questions concerning
the total topos one may replace fibered sites by the corresponding fibered
topoi.
\end{enonce*}

\begin{enonce*}{Lemma 7.4.12}\etiquette[7.4.12]{VI.lem7.4.12}
Let\pageoriginale $p:C\to I$ be a site fibered over a small category $I$,
and let $i$ be a final object of $I$ \emph{(I 10)}. The functor
$\alpha_{i!}:C^{\sim}_i\to\Top(C)$ \hyperlink{sga4.exp.VI.remark.7.4.3.item.4}{(7.4.3.4)} is exact and fully faithful.
The pairs of functors
$\beta_i=(\alpha_{i!},\alpha^{\ast}_i):\Top(C)\to C^{\sim}_i$ and
$\alpha_i=(\alpha_i^{\ast},\alpha_{i\ast}):C^{\sim}_i\to\Top(C)$ are
morphisms of topoi. The composite morphism
$\beta_i\alpha_i:C^{\sim}_i\to C^{\sim}_i$ is isomorphic to the identity
morphism.
\end{enonce*}

For every object $j$ of $I$, denote by $f_j:j\to i$ the unique morphism from
$j$ to $i$. Using the equivalence
$\Theta_C:\Top(C)\xrightarrow{\sim}\Hhom_{I^{\circ}}(I^{\circ},
(C^{\sim/I})')$ (\hyperref[VI.prop7.4.7]{7.4.7}), one sees that for every
object $X$ of $C_i$, $\alpha_{i!(X)}$ is the section
$j\to f^{\ast}_i(X)$, and that for every section
$j\mapsto\sigma(j)\in\Hom_{I^{\circ}}(I^{\circ},(C^{\sim/I})')$,
$\alpha^{\ast}_i(j\mapsto\sigma(j))$ is the object $\sigma(i)$. Hence
$\alpha_i^{\ast}\alpha_{i!}$ is isomorphic to the identity. Therefore
$\alpha_{i!}$ is fully faithful. Since the functors $f_j^{\ast}$ are left
exact, the functor $\alpha_{i!}$ is left exact. Thus $\beta_i$ is a morphism
of topoi, and since $\alpha_i^{\ast}\alpha_{i!}$ is isomorphic to the
identity, the morphism $\beta_i\alpha_i$ is isomorphic to the identity
morphism.


\begin{enonce*}{Theorem 7.4.13.1}\etiquette[7.4.13.1]{VI.thm7.4.13.1}
Let $p:C\to I$ and $q:D\to I$ be two $𝒰$-sites fibered over a small
category $I$, and let $m:D\to C$ be an $I$-functor. For $m$ to be a morphism
from the fibered site $C$ to the fibered site $D$
\emph{(\hyperref[VI.subsubsec7.2.2]{7.2.2})}, it is sufficient that $m$ be a
morphism from the total site $C$ to the total site $D$.
\end{enonce*}

\setcounter{subsubsection}{13}
\setcounter{paragraph}{1}
\paragraph{}\etiquette[7.4.13.2]{VI.par7.4.13.2}

We must show that if an $I$-functor $m:D\to C$ is a morphism from the total
site $C$ to the total site $D$, then $m$ is a morphism from the fibered site
$C$ to the fibered site $D$, \ie for every object $i$ of $I$, the functor
$m_i:D_i\to C_i$ induced by $m$ on the fibers at $i$ is a morphism from the
site $C_i$ to the site $D_i$. The proof of this last assertion occupies
paragraphs \hyperref[VI.par7.4.13.3]{7.4.13.3} through
\hyperref[VI.par7.4.13.7]{7.4.13.7}.

\paragraph{}\etiquette[7.4.13.3]{VI.par7.4.13.3}
Let us show that, for every object $i$ of $I$, the functor $m_i:D_i\to C_i$
is continuous. Indeed, there is a commutative diagram:\etiquette[133]{page133}
\[
\xymatrix@C=2cm@R=1.5cm{
D_i \ar[r]^{m_i} \ar[d]_{\alpha_{i!}}& C_i \ar[d]^{\alpha_{i!}}\\
D   \ar[r]^{m}& C
}
\]
and it\pageoriginale is enough to show that, for every covering sieve
$R\hookrightarrow X$ of $C_i/X$, the morphism
$m_{i!}(R)\to m_i(X)$ is bicovering (III \hyperref[III.prop1.2]{1.2} and \hyperref[III.prop1.5]{1.5}). But by the
commutativity of the diagram on page \pagehyperref[page133]{133} and the
fact that the functors $\alpha_{i!}$ and $m$ are continuous, the morphism
$\alpha_{i!}m_{i!}(R)\to\alpha_{i!}m_i(X)$ is bicovering. The assertion
therefore follows from \hyperref[VI.lem7.4.2.2]{7.4.2.2}.

\paragraph{}\etiquette[7.4.13.4]{VI.par7.4.13.4}
\emph{Reduction to the case in which $C$ and $D$ are fibered $𝒰$-topoi.}
Since the functors $m_i$ are continuous, they admit natural extensions to
the categories of sheaves, which we denote by
$m_i^{\ast}:D_i^{\sim}\to C_i^{\sim}$ (III \hyperref[III.prop1.2]{1.2} iv)). Let $D^{\sim/I}$ and
$C^{\sim/I}$ be the fibered $𝒰$-topoi associated with the fibered sites $D$
and $C$, respectively, and for every morphism $f$ of $I$, denote by
$f^{\ast}$ the inverse-image functors for the four fibered categories $C$,
$D$, $C^{\sim/I}$, and $D^{\sim/I}$. Since the functor $m$ is an
$I$-functor, for every $f:i\to j$ there is a canonical morphism
$$
f^{\ast} mj \leftarrow m_i f^{\ast}{\quoi}.
$$
Since the functors $f^{\ast}$, $m_i$, and $m_j$ are continuous, by passing
to the categories of sheaves we deduce canonical isomorphisms
$$
f^{\ast} m_j^{\ast} \longleftarrow m_i^{\ast} f^{\ast}{\quoi}.
$$
These isomorphisms have compatibility properties that make it possible to
construct a Cartesian functor
$m^{\ast}:D^{\sim/I}\to C^{\sim/I}$ inducing the functors $m_i^{\ast}$ on
the fiber topoi. Moreover, the diagram
\[
\xymatrix@C=2.5cm@R=1.5cm{
D \ar[r]^{m} \ar[d]_{\varepsilon_I}& C \ar[d]^{\varepsilon_I}\\
D^{\sim/I}   \ar[r]^{m^{\ast}}& C^{\sim/I}
}
\]
is commutative up to isomorphism. Since the functors $\epsilon_I$ induce
equivalences on the total topoi (\hyperref[VI.cor7.4.11]{7.4.11}), and since
$m$ is a morphism between the total sites, the functor
$m^{\ast}:D^{\sim/I}\to C^{\sim/I}$ is a morphism between the total sites.
We are therefore reduced to proving the assertion of
\hyperref[VI.par7.4.13.2]{7.4.13.2} when $D$ and $C$ are fibered $𝒰$-topoi,
which we assume henceforth.

\paragraph{}\etiquette[7.4.13.5]{VI.par7.4.13.5}
\emph{For every\pageoriginale object $i$ of $I$, the functor $m_i:D_i\to C_i$
takes the final object of $D_i$ to the final object of $C_i$}. Denote by
$m:D^{\sim}\to C^{\sim}$ the natural extension of $m$ to the total topoi
(III \hyperref[III.prop1.2]{1.2}). Using the equivalences
$\Theta_D:D^{\sim}\to\Hhom_{I^{\circ}}(I^{\circ},D')$
(\hyperref[VI.prop7.4.7]{7.4.7}), one immediately verifies that $m$
associates with every section
$\sigma\in\Hom_{I^{\circ}}({I^{\circ}},D')$ the section
$(i\mapsto m_i\sigma(i))\in\Hom_{I^{\circ}}({I^{\circ}},C')$. Since $m$ is
a morphism of sites, the functor $m$ is left exact and in particular takes
the final object of $\Hhom_{I^{\circ}}({I^{\circ}},D)$ to the final object
of $\Hhom_{I^{\circ}}({I^{\circ}},C')$. For every object $i$ of $I$, denote
by $e_i$ a final object of $D_i$. A final object of
$\Hhom_{I^{\circ}}({I^{\circ}},D')$ is the section $i\mapsto e_i$. Thus the
section $i\mapsto m_i(e_i)$ is a final object of
$\Hhom_{I^{\circ}}({I^{\circ}},C')$, and consequently, for every object
$i$ of $I$, $m_i(e_i)$ is a final object of $C_i$.

\paragraph{}\etiquette[7.4.13.6]{VI.par7.4.13.6}
\emph{Reduction to the case in which $i$ is a final object of} $I$. Let $i$
be an object of $I$. The functor $\alpha_{i!}:D_i\to D$ factors as
$$
D_i \xrightarrow{\bar{\alpha}_{i!}} D_{/e_i} \xrightarrow{\loc(e_i)} D{\quoi},
$$
where $\bar{\alpha}_i$ is the evident functor and $\loc(e_i)$ the forgetful
functor. Denote by $m/i:D_{/e_i}\to D_{m(e_i)}$ the functor that associates
with every object $u:X\to e_i$ of $D_{/e_i}$ the object
$m(u):m(X)\to m(e_i)$ of $C_{/m(e_i)}$. There is a commutative diagram
\[
\xymatrix@C=2.25cm@R=1.5cm{
D_i \ar[r]^{m_i} \ar[d]_{\bar{\alpha}_{i!}}& C_i\ar[d]^{\bar{\alpha}_{i!}}\\
D_{/e_i}  \ar[r]^{m/i} & C_{/m(e_i)}
}
\]
Since $e_i$ and $m(e_i)$ are final objects of $D_i$ and $C_i$, respectively
(\hyperref[VI.par7.4.13.5]{7.4.13.5}), the categories $D_{/e_i}$ and
$C_{/m(e_i)}$ are $𝒰$-topoi fibered over $I_{/i}$, and the functor $m_{/i}$
is an $I_{/i}$-functor. Since the property of a functor being a morphism of
sites is local (IV \hyperref[IV.subsec4.9]{4.9}), the functor $m_{/i}$ is a morphism from the total
site $C_{/m(e_i)}$ to the total site $D_{/(e_i)}$. We are therefore reduced
to proving that $m_i$ is a morphism of topoi when $i$ is a final object of
$I$, which we assume henceforth.\pageoriginale

\paragraph{}\etiquette[7.4.13.7]{VI.par7.4.13.7}
\emph{End of the proof}. Denote by $\alpha_{i!}:D_i\to D^{\sim}$,
$m^{\sim}:D^{\sim}\to C^{\sim}$, and $\alpha_{i!}:C_i\to C^{\sim}$ the
natural extensions to the categories of sheaves of the functors
$\alpha_{i!}$ and $m$. There is a diagram commutative up to isomorphism:
\[
\xymatrix@C=2.5cm@R=1.5cm{
D_i \ar[r]^{m_i} \ar[d]_{\alpha_{i!}}& C_i\ar[d]^{\alpha_{i!}}\\
D^{\sim}  \ar[r]^{m^{\sim}} & C .
}
\]
Since $i$ is a final object, the functor
$\alpha_i^{\ast}\alpha_{i!}:C_i\to C_i$ is isomorphic to the identity
(\hyperref[VI.lem7.4.12]{7.4.12}), and consequently $m_i$ is isomorphic to
$\alpha_i^{\ast}m^{\sim}\alpha_{i!}$. Moreover, the functors
$\alpha_i^{\ast}$, $m^{\sim}$, and $\alpha_i^{\ast}$ are left exact
(\hyperref[VI.lem7.4.12]{7.4.12}). Consequently, $m_i$ is left exact and is
therefore a morphism from the site $C_i$ to the site $D_i$.

\begin{enonce*}[remark]{Exercise 7.4.14}\etiquette[7.4.14]{VI.exe7.4.14}
Let $X=(X_i)_{i\in I}$ be a topos fibered over $I$ and let $T$ be a topos.
Show that the categories $\Hhomtop(X_i,T)$ are the fibers of a fibered
category denoted $\Hhomtop_I(X,T\times I)$. Show that the category of
sections over $I$ of $\Hhomtop_I(X,T\times I)$ is canonically equivalent to
$\Hhomtop_I(\Top(X),T)$, where $\Top(X)$ is the total topos of $X$.
\end{enonce*}

\begin{enonce*}[remark]{Exercise 7.4.15}\etiquette[7.4.15]{VI.exe7.4.15}
Let $X=(X_i)_{i\in I}$ be a topos fibered over $I$. Define a Cartesian
morphism $m$ from $X$ to the constant fibered topos with fiber $(\Ens)$.
Show that the total topos of the constant fibered topos with fiber $(\Ens)$
is canonically equivalent to $I^{\widehat{\phantom{,}}}$. Deduce a morphism
of topoi
$$
\Top(m) :\Top(X) \to I^{\widehat{\phantom{,}}}{\quoi}.
$$
Let $F=(F_i)_{i\in I}$ (where $F_i=F|X_i$) be an object of $\Top(X)$
(\hyperref[VI.prop7.4.7]{7.4.7}). Show that
$$
\Top(m)(F) = (i \mapsto \Gamma (X_i, F_i)){\quoi}.
$$
Deduce from \hyperref[VI.lem7.4.12]{7.4.12} that if $F$ is injective
abelian, then $F_i$ is injective abelian for every $i\in\ob I$. Deduce that
$$
R^q \Top(m)(F) = (i \mapsto H^{q} (X_i, F_i)){\quoi}.
$$
Show\pageoriginale that the Cartan--Leray spectral sequence of the morphism
$m$ (\hyperrefext[V][V.sec5]{5}) is
$$
H^{p+q} (\Top(X), F) \Leftarrow {\varprojlim_{I}}^{(p)} H^{q}
(X_i, F|X_i){\quoi}.
$$
\end{enonce*}


% END INCLUDED SGA4_English_current_buildable_source_20260906/source/06/drafts/SGA4_VI_LINES_4930_5251_English.texfrag


% BEGIN INCLUDED SGA4_English_current_buildable_source_20260906/source/06/drafts/SGA4_VI_LINES_5252_5495_English.texfrag
\section{Projective limits of fibered topoi}\etiquette[8]{VI.sec8}

\subsection{Generalities}\etiquette[8.1]{VI.subsec8.1}

\begin{enonce*}{Definition 8.1.1}\etiquette[8.1.1]{VI.def8.1.1}
Let $𝒰$ be a universe and let $p:F \to I$ be a fibered $𝒰$-topos. A
pair $(C, m)$ consisting of a $𝒰$-topos $C$ and a cartesian morphism
$m : C \times I \to F$ of fibered $𝒰$-topoi over $I$ is called a
projective limit of the fibered topos $F$ if, for every $𝒰$-topos
$D$, the functor
\begin{equation*}
\Hhomtop(D, C) \to \Hhomtop_{\cart/I}(D\times I, F)\tag{8.1.1.1}
\end{equation*}\etiquette[8.1.1.1]{VI.eq8.1.1.1}
obtained by composing the base-change functor
$$
\Hhomtop(D, C) \to \Hhomtop_{I}(D\times I, C \times I)
$$
with the functor of composition with $m$
(\hyperref[VI.subsubsec7.1.8]{7.1.8}) is an equivalence of categories.
\emph{(\cf \hyperref[VI.par8.1.3.3]{8.1.3.3} for a more
``geometric'' formulation)}.
\end{enonce*}

\setcounter{subsubsection}{1}
\subsubsection{}\etiquette[8.1.2]{VI.subsubsec8.1.2}

We use the notation of \hyperref[VI.def8.1.1]{8.1.1}. Let $D$ be a
$𝒰$-topos and let $m_1 : D \times I \to F$ be a morphism of fibered
$𝒰$-topoi. Translating Definition \hyperref[VI.def8.1.1]{8.1.1}, one
immediately sees that there is a morphism of topoi $n : D \to C$ and
an isomorphism of morphisms of fibered topoi over $I$,
$\gamma : m_{1} \xrightarrow{\sim} m \circ (n \times \id_I)$, making
the diagram commutative:
\[\vcenter{
\xymatrix@C=.75cm@R=.5cm{
D \times I  \ar[rrrr]^{m_1} \ar[dddd]_{n\times \id_I} &&{\phantom{.}}\ar@/_.5pc/@{-->}[dr]_{\gamma}  &&F\\
                         & &{\phantom{a}}& &\\
                         & && & \\
                         & && & \\
C \times I \ar[rrrruuuu]_{m}&& &&{\phantom{a}}\\
}}\leqno{(8.1.2.1)}
\]\etiquette[8.1.2.1]{VI.eq8.1.2.1}
Moreover, if $(n', \gamma')$ is another pair with the property above,
there is a \emph{unique} isomorphism
$\eta : n \xrightarrow{\sim} n'$ such that
\begin{equation*}
\gamma \circ m (\eta \times \id) = \gamma'{\quoi}.\tag{8.1.2.2}
\end{equation*}\etiquette[8.1.2.2]{VI.eq8.1.2.2}
It\pageoriginale follows that if $(D, m_1)$ is a projective limit of
$F$, then $n$ is an equivalence of $𝒰$-topoi (IV \hyperref[IV.subsec3.4]{3.4}), and the
additional datum of the isomorphism $\gamma$ making diagram
\eqhyperref[VI.eq8.1.2.1]{8.1.2.1} commutative determines the pair
$(n, \gamma)$ up to canonical isomorphism. Existence of the
projective limit will be proved in \hyperref[VI.thm8.2.3]{8.2.3} when
$I$ is cofiltered.

\subsubsection{}\etiquette[8.1.3]{VI.subsubsec8.1.3}
We denote by
$$
\limtop{I} F \mbox{ or } \limtop{I}
F_i \mbox{ or also } \underleftarrow{F}
$$
a projective limit of the fibered topos $F$, by
$\mu : (\limtop{I} F) \times I \to F$ the canonical morphism, and,
for every object $i$ of $I$, by
$$
\mu_i : \underleftarrow{F} \simeq \underleftarrow{F} \times {i} \to F_{i}
$$
the morphism of topoi induced by $\mu$ on the fiber over $i$.

Choose a bisplitting of $F$ (\hyperref[VI.subsubsec7.1.4]{7.1.4}),
\ie for every $f : i \to j$ a morphism of topoi
$f\bigdot = (f^{\ast}, f_{\ast}) : F_i \to F_j$ such that $f^{\ast}$
is an inverse-image functor for the fibered structure. Then, for every
pair $i \xrightarrow{f} j \xrightarrow{g} k$ of composable morphisms
of $I$, we have an isomorphism of morphisms of topoi
(\hyperref[VI.subsubsec7.1.3]{7.1.3})
$$
c_{f, g} : g.f. \simeq (gf).\quad {\quoi}.
$$
Since $\mu$ is a cartesian morphism of fibered topoi, for every
morphism $f : i \to j$ of $I$ we have an isomorphism of morphisms of
topoi (\hyperref[VI.subsubsec7.1.6]{7.1.6})
\begin{equation*}
b_f  : f. \mu_i \xrightarrow{\sim} \mu_j{\quoi},\tag{8.1.3.1}
\end{equation*}\etiquette[8.1.3.1.]{VI.eq8.1.3.1}
and the family of the $b_f$ satisfies relations analogous to
\eqhyperref[VI.eq7.1.6.2]{7.1.6.2}.

\setcounter{paragraph}{1}
\paragraph{}\etiquette[8.1.3.2]{VI.par8.1.3.2}

Let $D$ be a $𝒰$-topos, let $m : D \times I \to F$ be a morphism of
fibered topoi, let $m_i : D = D \times \{i\} \to F_i$ be the
morphisms induced by $m$ on the fibers, and, for every $f : i \to j$,
let
$$
b'_f : f . m_i \xrightarrow{\sim} m_j{\quoi},
$$
be the transition isomorphism. It follows from
\hyperref[VI.subsubsec8.1.2]{8.1.2} that there is a morphism of topoi
$n : D \to \underleftarrow{F}$\pageoriginale and a family of
isomorphisms
$$
\gamma_i : m_i \xrightarrow{\sim} \mu_i n \quad i \in \ob I {\quoi},
$$
such that, for every $f : i \to j$,
$$
b_f \circ f. (\gamma_i) = \gamma_j \circ b'_f{\quoi}.
$$
Moreover, the morphism $n$ and the family of the $\gamma_i$,
$i \in \ob (I)$, are determined up to unique isomorphism
(\hyperref[VI.subsubsec8.1.2]{8.1.2}).

Conversely, suppose we are given a topos $\underleftarrow{F}$,
morphisms $\mu_i : \underleftarrow{F} \to F_i$, and isomorphisms $b_f$
as in \eqhyperref[VI.eq8.1.3.1]{8.1.3.1.}, satisfying the usual
compatibility conditions, and suppose all these data have the
universal property described above. Then there is a unique morphism
$\mu : \underleftarrow{F} \times I \to F$ of fibered topoi giving
rise to the $\mu_{i}$ and the $b_f$, and $(\underleftarrow{F}, \mu)$
is a projective limit of the fibered topos $F$.

\paragraph{}\etiquette[8.1.3.3]{VI.par8.1.3.3}

Let $𝒱$ be a universe such that $\sheaf{\Uu} \in 𝒱$. The preceding
considerations therefore allow projective limits of fibered topoi to
be interpreted as ``projective limits,'' in the sense of
2-categories, of the pseudofunctors with values in the 2-category of
fibered $𝒰$-topoi that are elements of $𝒱$, deduced from the fibered
structure by choosing the morphisms $f.$. The reader should take care,
however, not to confuse this notion of ``projective limit'' in the
sense of 2-categories (with commutation up to isomorphism) with the
strict (or naive) notion of projective limit (strict commutation of
diagrams). Even when the strict notion makes sense (when the
pseudofunctor is a genuine functor), generalized ``projective
limits'' and strict projective limits do not coincide in general.

\subsubsection{}\etiquette[8.1.4]{VI.subsubsec8.1.4}

Let $p : F \to I$ and $q : G \to I$ be two fibered $𝒰$-topoi, and
let $m : F \to G$ be a morphism of fibered topoi. Suppose that the
fibered topoi $F$ and $G$ have projective limits
(\hyperref[VI.def8.1.1]{8.1.1}). It follows immediately from
Definition \hyperref[VI.def8.1.1]{8.1.1} that there is a morphism of
topoi
\begin{equation*}
\underleftarrow{m} : \underleftarrow{F} \to
\underleftarrow{G}{\quoi},\tag{8.1.4.1}
\end{equation*}\etiquette[8.1.4.1]{VI.eq8.1.4.1}
and an isomorphism of morphisms of fibered topoi
$\gamma : m \circ \mu \xrightarrow{\sim}
\mu \circ \underleftarrow{m} \times \id_{I}$, making the diagram
\pageoriginale commutative:
\[\vcenter{
  \xymatrix@C=.6cm@R=.45cm{
\underleftarrow{F}\times I \ar[rrr]^{\mu}
\ar[ddd]_{\underleftarrow{m}\times \id_I}   && & F\ar[ddd]^{m}\\
                          &            & &\\
                          &            & &\ar@/_.55pc/@{-->}[ld]_{\gamma} \\
\underleftarrow{G}\times I \ar[rrr]^{\mu} &       & &  G&.
}}\leqno{(8.1.4.2)}
\]\etiquette[8.1.4.2]{VI.eq8.1.4.2}
Moreover, if $(\underleftarrow{m}', \gamma')$ is another pair with the
property above, there is a \emph{unique} isomorphism
$\eta: \underleftarrow{m} \xrightarrow{\sim} \underleftarrow{m}'$ of
morphisms of topoi such that
$$
\gamma \circ \mu (\eta \times \id) = \gamma'{\quoi}.
$$

\subsubsection{}\etiquette[8.1.5]{VI.subsubsec8.1.5}

Let $p : F \to I$ be a fibered $𝒰$-topos that is an element of $𝒱$,
let $\phi : I'\to I$ be a functor that is an element of $𝒱$, and let
$p_{\phi} : F_{\phi} \to I'$ be the fibered $𝒰$-topos deduced from
$(F, p)$ by the base change $\phi$
(\hyperref[VI.subsubsec7.1.9]{7.1.9}). Let
$(\underleftarrow{F}, \mu)$ be a projective limit of $F$. By base
change, we obtain a morphism of fibered topoi over $I'$:
$$
\mu' : \underleftarrow{F} \times I' \to F_{\phi}{\quoi}.
$$
Let $(\underleftarrow{F}_{\phi}, \mu_{\phi})$ be a projective limit of
$F_{\phi}$. By the universal property of the projective limit
(\hyperref[VI.subsubsec8.1.2]{8.1.2}), we then have a morphism of
topoi
$$
m_{\phi} : \underleftarrow{F} \to \underleftarrow{F}_{\phi}{\quoi},
$$
and an isomorphism of morphisms of fibered topoi over $I'$:
$$
\gamma : \mu'\xrightarrow{\sim} \mu_{\phi} \circ (m_{\phi} \times \id_{I'}){\quoi}.
$$
The pair $(m_{\phi}, \gamma)$ is determined up to unique isomorphism.

\subsubsection{}\etiquette[8.1.6]{VI.subsubsec8.1.6}

Let $p : F \to I$ be a fibered $𝒰$-site over a small category $I$,
let $F^{\sim/I} \to I$ be the associated fibered $𝒰$-topos
(\hyperref[VI.subsubsec7.2.6]{7.2.6}), and let $D$ be a $𝒰$-topos.
First, we have a functor
\begin{equation*}
\Hhomtop_{\cart / I} (D \times I, F^{\sim/I}) \to
\Mmorsite_{\cart/I}(D \times I, F){\quoi},\tag{8.1.6.1}
\end{equation*}\etiquette[8.1.6.1]{VI.eq8.1.6.1}
which associates with every $m$ the composite of $m^{\ast}$ with
$\epsilon_{F /I}$ (\hyperref[VI.prop7.2.7]{7.2.7}). Moreover, by the
definition of the category $\Mmorsite_{\cart /I} (D \times I, F)$
\eqhyperref[VI.eq7.2.2.2]{7.2.2.2}, we have a functor
\begin{equation*}
\Mmorsite_{\cart/I}(D \times I, F) \to \Hhom_{\cart / I} (F, D \times
I)^{\circ}{\quoi}.\tag{8.1.6.2}
\end{equation*}\etiquette[8.1.6.2]{VI.eq8.1.6.2}
Finally,\pageoriginale by the definition of the inductive limit
(\hyperref[VI.def6.3]{6.3}), we have a functor
\begin{equation*}
\Hhom_{\cart / I} (F, D \times I)^{\circ} \to
\Hhom(\varinjlim_{I^{\circ}}F, D)^{\circ}{\quoi}.\tag{8.1.6.3}
\end{equation*}\etiquette[8.1.6.3]{VI.eq8.1.6.3}
Denote by
\begin{equation*}
\Theta : \Hhomtop_{\cart / I}(D \times I, F^{\sim / I}) \to
\Hom(\varinjlim_{I^{\circ}}F, D)^{\circ}\tag{8.1.6.4}
\end{equation*}\etiquette[8.1.6.4]{VI.eq8.1.6.4}
the composite of these last three functors, and by
\begin{equation*}
\Pi : F \to \varinjlim_{I^{\circ}}F{\quoi},\tag{8.1.6.5}
\end{equation*}\etiquette[8.1.6.5]{VI.eq8.1.6.5}
the canonical functor.

\begin{enonce*}{Proposition 8.1.7}\etiquette[8.1.7]{VI.prop8.1.7}
The functor $\Theta$ is fully faithful. Its essential image is the set
of functors $m : \varinjlim_{I^{\circ}}F \to D$ such that the functor
$$
(m \circ \Pi, p) : F \to D \times I
$$
is a morphism from the total site $D \times I$ to the total site $F$.
\end{enonce*}

It follows from \hyperref[VI.prop7.2.7]{7.2.7} that the functor
\eqhyperref[VI.eq8.1.6.1]{8.1.6.1} is an equivalence of categories,
from \hyperref[VI.subsubsec7.2.2]{7.2.2} that the functor \hyperref[VI.eq8.1.6.2]{(8.1.6.2)}
is fully faithful, and from \hyperref[VI.def6.3]{6.3} that the functor
\eqhyperref[VI.eq8.1.6.3]{8.1.6.3} is an equivalence of categories.
Consequently, $\Theta$ is fully faithful. A functor
$m : \varinjlim_{I^{\circ}}F \to D$ is isomorphic to the image, under
$\Hhom_{\cart / I} (F, D \times I)^{\circ} \to
\Hhom(\varinjlim_{I^{\circ}}F, D)^{\circ}$, of the functor
$(m \circ \Pi, p)$, and such a functor belongs
to the essential image of \eqhyperref[VI.eq8.1.6.2]{8.1.6.2} if and
only if it is a morphism from the total site $D \times I$ to the total
site $F$ (\hyperref[VI.thm7.4.13.1]{7.4.13.1}); hence the proposition.

% END INCLUDED SGA4_English_current_buildable_source_20260906/source/06/drafts/SGA4_VI_LINES_5252_5495_English.texfrag


% BEGIN INCLUDED SGA4_English_current_buildable_source_20260906/source/06/drafts/SGA4_VI_LINES_5496_5856_English.texfrag
\subsection{Construction of the projective limit when the index
  category is cofiltered.}\etiquette[8.2]{VI.subsec8.2}

\begin{enonce*}{Lemma 8.2.1}\etiquette[8.2.1]{VI.lem8.2.1}
Let $I$ and $I'$ be two cofiltered categories \emph{(I \hyperref[I.sec8.1]{8.1})}, and let
$\phi:I'\to I$ be a functor such that
$\phi^{\circ}:{I'}^{\circ}\to I^{\circ}$ is cofinal \emph{(I \hyperref[I.sec8.1]{8.1})}.
Let, moreover, $p:F\to I$ be a fibered $𝒰$-topos, let $D$ be a
$𝒰$-topos, let $m:D\times I\to F$ be a Cartesian morphism of fibered
topoi, and let $m':D\times I'\to F'$ be the morphism of fibered topoi
deduced from $m$ by the base change $\phi$
(\hyperref[VI.subsubsec7.1.9]{7.1.9}). The pair $(D,m)$ is a
projective limit of $F$ if and only if $(D,m')$ is a projective limit
of $F'$.
\end{enonce*}

The proof\pageoriginale is merely a rather lengthy routine
verification. It is left to the patient reader, who may use the
viewpoint of projective limits of pseudofunctors
(\hyperref[VI.par8.1.3.3]{8.1.3.3}).

\begin{enonce*}{Lemma 8.2.2}\etiquette[8.2.2]{VI.lem8.2.2}
Let $p:F\to I$ be a fibered $𝒰$-site over a small cofiltered
category, let $\underrightarrow{F}$ be the inductive limit of the
fibered category $F$ \hyperref[VI.def6.3]{(6.3.9)}, and let
$\pi:F\to\underrightarrow{F}$ be the canonical functor. Equip
$\underrightarrow{F}$ with the coarsest topology making $\pi$
continuous \emph{(III 1)}. Let $D$ be a $𝒰$-site and
$n:\underrightarrow{F}\to D$ a functor. The following properties are
equivalent:
\begin{enumerate}
\renewcommand{\theenumi}{\roman{enumi}}
\renewcommand{\labelenumi}{\theenumi)}
\item $n$ is a morphism of sites $D\to\underrightarrow{F}$.

\item $n\circ\pi$ is a morphism of sites (where $F$ is equipped with
  the total topology).

\item $(n\circ\pi,p):F\to D\times I$ is a morphism from the total
  site $D\times I$ to the total site $F$.

\item For every object $i$ of $I$, the composite functor
  $F_i\xrightarrow{\alpha_{i!}}F\xrightarrow{n\circ\pi}D$ is a
  morphism of sites.
\end{enumerate}
\end{enonce*}

The equivalence of i) and ii) follows from
\hyperref[VI.exe6.6]{6.6}, and that of iii) and iv) follows from
\hyperref[VI.thm7.4.13.1]{7.4.13}. Let us prove the equivalence of iii) and ii). For this it is
enough to show that the first projection functor
$\pr_1:D\times I\to D$ is a morphism of sites. Denote by $D^{\sim}$
the topos of sheaves on $D$. It follows from
\hyperref[VI.prop7.4.7]{7.4.7} that $(D\times I)^{\sim}$ is equivalent
to the topos
$\Hhom_{I^{\circ}}(I^{\circ},D^{\sim}\times I^{\circ})=
\Hhom(I^{\circ},D^{\sim})$. Moreover, one sees immediately that the
natural extension of $\pr_1:D\times I\to D$ to sheaves (III \hyperref[III.prop1.2]{1.2}) is
the functor from $\Hhom(I^{\circ},D^{\sim})$ to $D^{\sim}$ that
associates with the functor $i\mapsto\sigma(i)$ the object
$\varinjlim_{I^{\circ}}\sigma(i)$. Since $I^{\circ}$ is filtered,
inductive limits over $I^{\circ}$ are exact (II 4), and hence
$\pr_1$ is a morphism of sites (IV \hyperref[IV.subsec4.9]{4.9}). It remains to prove the
equivalence of ii) and iv). First treat the case in which $D$ is a
$𝒰$-topos and $F$ a fibered $𝒰$-topos over $I$.

Let $i$ be an object of $I$. Arguing as in
\hyperref[VI.par7.4.13.6]{7.4.13.6}, one sees that the functor
$\alpha_{i!}:F_i\to F$ factors as a morphism
$\bar{\alpha}_{i!}:F_i\to F/\alpha_{i!}(e_i)$ followed by the
localization functor $F/\alpha_{i!}(e_i)\to F$, where $e_i$ is a
final object of $F_i$. It is therefore enough to show that the
composite functor
$F/\alpha_{i!}(e_i)\xrightarrow{n\circ\pi}D$ is a morphism of sites.
But there is a diagram commutative up to isomorphism:
\[
\xymatrix@C=2.5cm@R=1.25cm{
F/\alpha_{i!}(e_i) \ar[d]\ar[r]^{n\circ
  \pi/\alpha_{i!}(e_i)}& D/n\circ \pi \alpha_{i!}(e_i) \ar[d]\\
F \ar[r]^{n\circ \pi}& D
}
\]
Since\pageoriginale the property of being a morphism of sites is
local (IV \hyperref[IV.subsec5.10]{5.10}), it is enough to show that
$n\circ\pi\circ\alpha_{i!}(e_i)$ is a final object of $D$. For every
morphism $f:j\to k$ of $I$, there exists a unique morphism
$e_f:\alpha_{j!}(e_j)\to\alpha_{k!}(e_k)$ over $f$, and this
morphism is Cartesian. Thus $j\mapsto\alpha_{j!}(e_j)$ is a Cartesian
section of $F$ over $I$. Consequently, the canonical morphism
$n\circ\pi\circ\alpha_{i!}(e_i)\to
\varinjlim_I n\circ\pi\circ\alpha_{j!}(e_j)$ is an isomorphism,
because $\pi$ carries Cartesian morphisms to isomorphisms. Moreover,
$\varinjlim_I\alpha_{j!}(e_j)$, where the inductive limit is taken in
$F^{\widehat{\phantom{,}}}$, is a final object of
$F^{\widehat{\phantom{,}}}$. Since $n\circ\pi$ is a morphism of
sites, it carries the final object of $F^{\sim}$ to a final object of
$D$. Hence $\varinjlim_I n\circ\pi\circ\alpha_{j!}(e_j)$ is a final
object of $D$, completing the proof in this case.

In the general case, the functor $(n\pi,p):F\to D\times I$ is
Cartesian and continuous (\hyperref[VI.prop7.4.4]{7.4.4}). Moreover,
$n\circ\pi$ is the composite functor
$$
F \xrightarrow{(n \pi, p)} D \times I \xrightarrow{\pr_1} D{\quoi}.
$$
Passing to the associated fibered $𝒰$-topoi gives a diagram
commutative up to isomorphism:
\[
\xymatrix@C=2cm@R=1.75cm{
F       \ar@{-}[r]^{(n\pi,p)} \ar[d]_{\varepsilon_{F/I}}& D\times I
\ar[r]^{\pr_1}\ar[d]^{\varepsilon_D\times \id} & D \ar[d]^{\varepsilon_D}\\
F^{\sim/I}  \ar[r]^{(n\pi,p)^{\sim/I}} & D^{\sim}\times I \ar[r]^{\pr_1} & D^{\sim}&.
}
\]
Since $(n\pi,p)^{\sim/I}$ is Cartesian, it is of the form
$(n'\pi',p')$, where $p':F^{\sim/I}\to I$ and
$\pi':F^{\sim/I}\to\varinjlim_{I^{\circ}}F^{\sim}_i$ are the
canonical functors. Thus
$\pr_{1^{\circ}}(n\pi,p)^{\sim/I}=n'$. Since the functors
$\epsilon_{F/I}$ and $\epsilon_D$ induce equivalences on the topoi of
sheaves (\hyperref[VI.cor7.4.11]{7.4.11}), the functors $n$ and $n'$
induce isomorphic functors after passage to associated sheaves (III
\hyperref[III.prop1.2]{1.2}), and hence $n'$ is a morphism of sites. By what precedes,\pageoriginale
the functor $(n\pi,p)^{\sim/I}$ induces morphisms of topoi on the
fibers. It is therefore a morphism from the total site $F^{\sim/I}$
to the total site $D^{\sim}\times I$ \hyperref[VI.thm7.4.13.1]{(7.4.13)}. Since the functors
$\epsilon_{F/I}$ and $\epsilon_D\times\id$ induce equivalences on the
total topoi (\hyperref[VI.cor7.4.11]{7.4.11}), the functor $(n\pi,p)$
is a morphism of fibered sites \hyperref[VI.thm7.4.13.1]{(7.4.13)}, which proves the assertion.

\begin{enonce*}{Theorem 8.2.3}\etiquette[8.2.3]{VI.thm8.2.3}
Let $p:F\to I$ be a fibered $𝒰$-site over an essentially small
cofiltered category \emph{(I \hyperref[I.sec8.1]{8.1})}.
\begin{enumerate}
\renewcommand{\labelenumi}{\theenumi)}
\item Let $\underrightarrow{F}$ denote the site whose underlying
  category is the inductive limit of the fibered category $F\to I$
  \emph{(\hyperref[VI.def6.3]{6.3})} and whose topology is the
  coarsest topology making the canonical functor
  $\pi:F\to\underrightarrow{F}$ continuous. Then
  $\underrightarrow{F}$ is a $𝒰$-site and
$$
(\pi, p) : F \to \underrightarrow{F} \times I
$$
is a morphism from the fibered site $\underrightarrow{F}\times I$ to
the fibered site $F$.

\item Let $\mu:\underrightarrow{F}^{\sim}\times I\to F^{\sim/I}$ be
  the morphism of fibered $𝒰$-topoi deduced from $(\pi,p)$ by passing
  to the associated fibered $𝒰$-topoi. The pair
  $(\underrightarrow{F}^{\sim},\mu)$ is a projective limit of
  $F^{\sim/I}$.
\end{enumerate}
\end{enonce*}

\begin{enonce*}[remark]{Remark 8.2.4}\etiquette[8.2.4]{VI.rem8.2.4}
\begin{enumerate}
\renewcommand{\labelenumi}{\theenumi)}
\item In view of the notation introduced in
  \hyperref[VI.subsubsec8.1.3]{8.1.3}, one may set
$$
\limtop{I} F^{\sim} = \underrightarrow{F}^{\sim}.
$$

\item It follows from \hyperref[VI.prop7.4.4]{7.4.4} that the topology
  on $\underrightarrow{F}$ is also the coarsest topology making the
  family of functors
  $\pi_j:F_j\to F\xrightarrow{\pi}\underrightarrow{F}$,
  $j\in\ob I$, continuous.
\end{enumerate}
\end{enonce*}

\begin{enonce*}{Definition 8.2.5}\etiquette[8.2.5]{VI.def8.2.5}
The site $\underrightarrow{F}$ introduced in
\hyperref[VI.thm8.2.3]{8.2.3} is called the projective-limit site of
the fibered site $F$.
\end{enonce*}

As an exercise, the reader may, as in \hyperref[VI.lem8.2.1]{8.2.1},
define the notion of a projective limit of a fibered site and verify
that the site $\underrightarrow{F}$ is indeed a projective limit of
the site $F$. The second assertion of the theorem may therefore be
paraphrased as follows: the topos of sheaves on the projective-limit
site is equivalent to the projective limit of the associated fibered
topos.

\setcounter{subsubsection}{5}
\subsubsection{Proof of the theorem: Reduction to the case of a small
  index category}\etiquette[8.2.6]{VI.subsubsec8.2.6}

We\pageoriginale shall give only indications. Let $I'$ be a full
subcategory of $I$ such that ${I'}^{\circ}$ is cofinal in
$I^{\circ}$. Denote by $F'\to I'$ the fibered category deduced from
$F$ by the base change $I'\to I$, and by
$\pi':F'\to\underrightarrow{F}'$ the canonical functor. One first
observes that the categories $\underrightarrow{F}$ and
$\underrightarrow{F}'$ are equivalent and that the topology on
$\underrightarrow{F}'$ deduced from the topology of
$\underrightarrow{F}$ by this equivalence is the coarsest topology
making the functor $\pi':F'\to\underrightarrow{F}'$ continuous.
Assertion 1) of the theorem follows when it has been proved in the
case where $I$ is small. Assertion 2) in the general case then
follows from assertion 2) in the case where $I$ is small by Lemma
\hyperref[VI.lem8.2.1]{8.2.1}.

\subsubsection{End of the proof of the theorem}\etiquette[8.2.7]{VI.subsubsec8.2.7}
The first assertion follows from Lemma \hyperref[VI.lem8.2.2]{8.2.2},
applied with $D=\underrightarrow{F}$ and $n=\id$. Let $D$ be a
$𝒰$-topos. There is a diagram commutative up to isomorphism:
{\def\objectstyle{\scriptstyle}
\def\labelstyle{\scriptstyle}
\[
\xymatrix@C=.8cm@R=1.15cm{
\Hhomtop(D,{(\underrightarrow{F})}^{\sim}) \ar[r]^{(1)}\ar[d]_{(3)}&
\Mmorsite(D,\underrightarrow{F}) \ar@{^{(}->}[r]^{(2)} \ar[d]_{(3)}&
\Hhom(\underrightarrow{F},D)^{\circ}\ar[dd]_{(5)}\\
\Hhomtop_{\Cart/I}(D\times I,(\underrightarrow{F})^{\sim}\times
I)  \ar[r]^{(1)} \ar[d]_{(4)}& \Mmorsite_{\Cart/I}(D\times I,
(\underrightarrow{F})\times I) \ar[d]&\\
\Hhomtop_{\Cart/I}(D\times I,F^{\sim/I})
\ar[r]^{(1)} & \Mmorsite_{\Cart/I}(D\times I,F)  \ar@{^{(}->}[r]^{(2)} &
\Hhom_{\Cart/I}(F,D\times I)^{\circ},
}
\]}

\noindent where the functors of type (1) are equivalences (IV \hyperref[IV.prop4.9.4]{4.9.4}
and \hyperref[VI.prop7.2.7]{7.2.7}), the functors of type (2) are fully faithful (IV \hyperref[IV.subsubsec4.9.1]{4.9.1} and
\hyperref[IV.eq7.2.2]{7.2.2}), the functors of type (3) are base-change functors, the
functors of type (4) are composition functors (with $\mu$, \resp
$(\pi,p)$), and functor (5) is the equivalence deduced from the
universal property of the inductive limit
(\hyperref[VI.def6.3]{6.3}). To prove that the functor
$\Hhomtop(D,(\underrightarrow{F})^{\sim})\xrightarrow{(4)\circ(3)}
\Hhomtop_{/I}(D\times I,F^{\sim/I})$ is an equivalence, it is enough
to show that the functor
$$
\Mmorsite(D, \underrightarrow{F})
\xrightarrow{(4) \circ (3)} \Mmorsite_{/I}(D \times I,F)
$$
is an\pageoriginale equivalence of categories. Equivalently, it is
enough to show that the two fully faithful functors
\begin{align*}
& \Mmorsite(D, \underrightarrow{F}) \xrightarrow{(5) \circ (2)} \Hhomcart_I(F, D \times I)^{\circ}\\
& \Mmorsite_{/I}(D \times I, F) \xrightarrow{(2)} \Hhomcart_I(F, D \times
I)^{\circ}
\end{align*}
have the same essential images, which follows immediately from Lemma
\hyperref[VI.lem8.2.2]{8.2.2}.

\subsubsection{}\etiquette[8.2.8]{VI.subsubsec8.2.8}

Let $p:F\to I$ be a fibered $𝒰$-topos over a category $I$, and let
$(\limtop{I}F_i,\mu)$ be a projective limit of $F$. Thus there is a
Cartesian morphism of fibered $𝒰$-topoi over $I$:
$$
\mu:\limtop{I} F_i \times I \to F{\quoi},
$$
whence, on passing to the categories fibered over $I^{\circ}$ by the
direct image functors (\hyperref[VI.subsubsec7.1.3]{7.1.3}), a
Cartesian functor
$$
\mu_{\ast}: \limtop{I} F_i \times I^{\circ} \to F'
$$
which is the direct image functor for $\mu$
(\hyperref[VI.subsubsec7.1.7]{7.1.7}). By the definition of the
projective limit of a fibered category (\hyperref[VI.subsec6.10]{6.10}),
there is therefore a canonical functor
\begin{equation*}
\Theta : \limtop{I} F_i  \to
\varprojlim_{I, f_{\ast}} F_i = \Hhom_{\cart/I^{\circ}} (I^{\circ},
F'){\quoi}.\tag{8.2.8.1}
\end{equation*}\etiquette[8.2.8.1]{VI.eq8.2.8.1}

\begin{enonce*}{Theorem 8.2.9}\etiquette[8.2.9]{VI.thm8.2.9}
Use the notation of \hyperref[VI.subsubsec8.2.8]{8.2.8}, and suppose
that $I$ is an essentially small cofiltered category. Then the functor
$\Theta$ is an equivalence of categories. Suppose that $I$ is small.
Set $\limtop{I}F_i=(\underrightarrow{F})^{\sim}$ (which is permitted
by \hyperref[VI.thm8.2.3]{8.2.3} and
\hyperref[VI.rem8.2.4]{8.2.4}), and interpret the category
$\Hhom_{I^{\circ}}(I^{\circ},F')$ as the category of sheaves on the
total site $F$ (\hyperref[VI.prop7.4.7]{7.4.7}). Then the composite
functor
$$
(\underrightarrow{F})^{\sim} = \limtop{I} F_i
\xrightarrow{\Theta} \Hhom_{\cart/I^{\circ}} (I^{\circ}, F') \hookrightarrow \Hhom_{I^{\circ}}
(I^{\circ}, F') \approxeq F^{\sim}{\quoi},
$$
is precisely the direct image functor for the morphism of sites
defined by the functor
$$
\pi : F \to \underrightarrow{F}{\quoi}.
$$
\end{enonce*}

\setcounter{subsubsection}{9}
\subsubsection{}\etiquette[8.2.10]{VI.subsubsec8.2.10}

Some\pageoriginale comments before the proof of Theorem
\hyperref[VI.thm8.2.9]{8.2.9}. Besides solving the problem of the
existence of projective limits when the index category is cofiltered,
a solution which is curiously nontrivial in this theory, Theorems
\hyperref[VI.thm8.2.3]{8.2.3} and \hyperref[VI.thm8.2.9]{8.2.9}
provide two further pieces of information essential for working with
projective-limit topoi in practice. First, Theorem
\hyperref[VI.thm8.2.3]{8.2.3} interprets the projective-limit topos as
the topos of sheaves on a site whose underlying category is the
inductive-limit category of the categories $F_i$ under the system of
inverse image functors. In applications, one will be able to
interpret this category and its topology concretely. Second, Theorem
\hyperref[VI.thm8.2.9]{8.2.9} asserts that a sheaf on the projective
limit is determined by the system of its direct images in the topoi
$F_i$ (which gives a Cartesian section over $I^{\circ}$ of the
category $(F^{\sim/I})'$). The theorem further asserts conversely
that, when one is given for every object $i$ of $I$ a sheaf $X_i$ on
$F_i$ and, for every morphism $f:i\to j$ of $I$, an isomorphism
$X_j\simeq f_{\ast}X_i$, these isomorphisms being subject to
compatibility conditions, there exists essentially a unique sheaf on
the projective limit whose direct images give the system of the
$X_i$.

\subsubsection{}\etiquette[8.2.11]{VI.subsubsec8.2.11}

Assertion 2) of \hyperref[VI.thm8.2.3]{8.2.3} and the first assertion
of \hyperref[VI.thm8.2.9]{8.2.9} may be summarized by the striking
formula
\begin{equation*}
\varprojlim_{I, f_{\ast}}F_i^{\sim} \approxeq
\limtop{I, f.} F_i^{\sim} \approxeq
(\varinjlim_{I^{\circ}, f_{\ast}} F_i)^{\sim}{\quoi}. \tag{8.2.11.1}
\end{equation*}\etiquette[8.2.11.1]{VI.eq8.2.11.1}

\subsubsection{}\etiquette[8.2.12]{VI.subsubsec8.2.12}
\emph{Proof of Theorem} \hyperref[VI.thm8.2.9]{8.2.9}.
To prove the first assertion, reduce to the case where the category
$I$ is small by using \hyperref[VI.lem8.2.1]{8.2.1} and an analogous
lemma on projective limits of fibered categories. We leave the details
to the reader. Suppose henceforth that $I$ is a small category. We
know that the functor
$$
\pi : F \to \varinjlim_{I^{\circ}} F_i
$$
is a morphism of sites (\hyperref[VI.lem8.2.2]{8.2.2}). The direct
image functor
$$
\pi_{\ast} : (\varinjlim_{I^{\circ}} F_i)^{\sim} \to F^{\sim}
$$
that it\pageoriginale defines on the topoi is then fully faithful,
and its essential image is the set of sheaves on $F$ that carry the
Cartesian morphisms of $F$ to isomorphisms
(\hyperref[VI.exe6.6]{6.6}). These sheaves correspond, under the
equivalence
$F^{\sim}\simeq\Hhom_{I^{\circ}}(I^{\circ},F')$, to the Cartesian
sections of $F'$ over $I^{\circ}$
(\hyperref[VI.prop7.4.7]{7.4.7}). Moreover, for every
$j\in\ob(I)$, denote by $\alpha_{j!}:F_j\to F$ the inclusion functor,
which gives rise to the three functors
$\xymatrix{(\alpha_{j!},\alpha^{\ast}_j,\alpha_{j\ast}):F_j
\ar@<4pt>[r]\ar@<-4pt>[r]&F\ar[l]}$
(\hyperref[VI.rem7.4.3]{7.4.3}). It follows from the definition of the
equivalence $F^{\sim}\simeq\Hhom_{I^{\circ}}(I^{\circ},F')$
(\hyperref[VI.prop7.4.7]{7.4.7}) that a sheaf $X$ on
$\varinjlim_{I^{\circ}}F_i$ is carried by the functor
$(\varinjlim_{I^{\circ}}F_i)^{\sim}\xrightarrow{\pi_{\ast}}F^{\sim}
\simeq\Hhom_{I^{\circ}}(I^{\circ},F')$ to the Cartesian section
$j\mapsto\alpha^{\ast}_j\pi_{\ast}(X)$. But the description in
\hyperref[VI.thm8.2.3]{8.2.3} of the functor
$\mu:(\varinjlim_{I^{\circ}}F_i)^{\sim}\times I\to F$ shows that,
for every object $j$ of $I$, the direct image functor
$\mu_{j\ast}:(\varinjlim_{I^{\circ}}F_i)^{\sim}\to F_j$ for the
morphism $\mu_j:(\varinjlim_{I^{\circ}}F_i)^{\sim}\to F_j$ obtained
from $\mu$ by passing to the fiber at $j$ is precisely the functor
$\alpha^{\ast}_j\pi_{\ast}$. Consequently, the functor
$(\varinjlim_{I^{\circ}}F_i)^{\sim}\xrightarrow{\pi_{\ast}}F^{\sim}
\simeq\Hhom_{I^{\circ}}(I^{\circ},F')$ is precisely the functor
$(\varinjlim_{I^{\circ}}F_i)^{\sim}\xrightarrow{\Theta}
\Hhom_{\cart/I^{\circ}}(I^{\circ},F')\hookrightarrow
\Hhom_{I^{\circ}}(I^{\circ},F')$, which proves the theorem.


% END INCLUDED SGA4_English_current_buildable_source_20260906/source/06/drafts/SGA4_VI_LINES_5496_5856_English.texfrag


% BEGIN INCLUDED SGA4_English_current_buildable_source_20260906/source/06/drafts/SGA4_VI_LINES_5857_6042_English.texfrag
\subsection{Topology of the projective-limit site: the case of
  coherent topoi}\etiquette[8.3]{VI.subsec8.3}

\subsubsection{}\etiquette[8.3.1]{VI.subsubsec8.3.1}

In this section, $p:F\to I$ is a fibered $\Uu$-site over an
essentially small cofiltered category, whose inverse-image functors
$f^{\ast}:F_i\to F_j$ commute with fiber products. In addition, for
every object $X$ of a fiber category $F_i$, we are given a set
$\Cov(X)$ of covering families for the topology of $F_i$ which have
properties (PTO) and (PT1) of II \hyperref[II.def1.3]{1.3} and generate the topology of
$F_i$. Finally, suppose that for every $f:i\to j$, every
$Y\in\ob F_j$, and every family
$(Y_{\alpha}\to Y)_{\alpha\in A}$ in $\Cov(Y)$, the family
$(f^{\ast}(Y_{\alpha})\to f^{\ast}(Y))_{\alpha\in A}$ belongs to
$\Cov(f^{\ast}(Y))$. These conditions on $\Cov(X)$ are satisfied, for
example, if one takes $\Cov(X)=$ \{all covering families of $X$ in
$F_i$\}. In\pageoriginale this section we shall study the
projective-limit site of $F$ under suitable finiteness hypotheses
(\hyperref[VI.thm8.3.13]{8.3.13}).

\subsubsection{}\etiquette[8.3.2]{VI.subsubsec8.3.2}

Let $\pi:F\to\underrightarrow{F}$ be the canonical functor, and let
$Y$ be an object of $\underrightarrow{F}$. Denote by $\Cov(Y)$ the set
of families $(Y_{\alpha}\xrightarrow{n_{\alpha}}Y)_{\alpha\in A}$
of the following form:

There exist $i\in I$, an object $X$ of $F_i$, a family
$(X_{\alpha}\xrightarrow{{n'}_\alpha}X)_{\alpha\in A}\in\Cov(X)$,
an isomorphism $\phi:\pi(X)\simeq Y$, and a family of isomorphisms
$\phi_{\alpha}:\pi(X_{\alpha})\simeq Y_{\alpha}$, $\alpha\in A$,
such that for every $\alpha\in A$ the diagram
\[
\xymatrix@C=2.5cm@R=2cm{
\pi(X_{\alpha})
\ar@{-}[r]^{\displaystyle{\mathop{\sim}^{\phi_{\alpha}}}} \ar[d]^{\pi(n'_{\alpha})}& Y_{\alpha}\ar[d]^{n_{\alpha}} \\
\pi(X)  \ar@{-}[r]^{\displaystyle{\mathop{\sim}^{\phi}}} & Y
}
\]
is commutative.

\begin{enonce*}{Proposition 8.3.3}\etiquette[8.3.3]{VI.prop8.3.3}
\begin{enumerate}
\renewcommand{\labelenumi}{\theenumi)}
\item The families belonging to $\Cov(Y)$, for
  $Y\in\ob\underrightarrow{F}$, generate the topology of the
  projective-limit site (\hyperref[VI.def8.2.5]{8.2.5}).

\item The family $Y\mapsto\Cov(Y)$, $Y\in\ob\underrightarrow{F}$,
  has properties \emph{(PTO)} and \emph{(PT1)} of \emph{II \hyperref[II.def1.3]{1.3}}.
\end{enumerate}
\end{enonce*}

We shall need the following lemma.

\begin{enonce*}{Lemma 8.3.4}\etiquette[8.3.4]{VI.lem8.3.4}
Let $i$ be an object of $I$, and let $n:X'\to X$ be a base-changeable
morphism of $F_i$ such that $f^{\ast}(n)$ is base-changeable for every
$f:j\to i$. Then $\pi(n)$ is base-changeable. The functors
$$
\pi_j:F_j\hookrightarrow F\xrightarrow{\pi}\underrightarrow{F},
\quad j\in\ob I{\quoi},
$$
commute with fiber products.
\end{enonce*}

Recall that every object $Y$ of $\underrightarrow{F}$ is equal to an
object $\pi(Z)$, $Z\in\ob F$, and that every morphism
$m:Y\to\pi(X)$ is of the form
$$
\pi(Z)\xrightarrow{\pi(s)^{-1}}\pi(Z')
\xrightarrow{\pi(m')}\pi(X){\quoi},
$$
where $s$ is a cartesian morphism of $F$. Consequently, to show that
$\pi(n)$ is base-changeable,\pageoriginale it is enough to show that
the fiber product of every diagram
\[
\xymatrix@C=2.5cm@R=1.25cm{
                           &\pi(X')\ar[d]^{\pi(n)}\\
\pi(Z')\ar[r]^{\pi(m')}&\pi(X)
}
\]
is representable. The morphism $m':Z'\to X$ factors as
$Z'\xrightarrow{m''}X''\xrightarrow{s'}X$, where $p(m'')$ is the
identity of $p(Z')=j$ and $s'$ is cartesian over $p(s')=f$. Put
$X'''=f^{\ast}(X')$ and $n'=f^{\ast}(n)$. We then have a commutative
diagram
\[
\xymatrix@C=2.5cm@R=1.5cm{
X'''\ar[r]^{s''}\ar[d]^{n'}&X'\ar[d]^{n}\\
X''\ar[r]^{s'}&X,
}
\]
and hence a commutative diagram in $\underrightarrow{F}$:
\[
\xymatrix@C=2.5cm@R=1.5cm{
        &\pi(X''')\ar[r]^{\sim}\ar[d]_{\pi(n')}&\pi(X')\ar[d]^{\pi(n)}\\
\pi(Z')\ar@{-}[r]^{\pi(m'')}&\pi(X'')\ar[r]^{\sim}&\pi(X)
}
\]
where $n'$ is a base-changeable morphism of $F_j$. Thus, to show that
$\pi(n)$ is base-changeable, it is enough to show that the product
$\pi(Z')\times_{\pi(X'')}\pi(X''')$ is representable, and therefore
it is enough to show that
$\pi_j:F_j\hookrightarrow F\to\underrightarrow{F}$ commutes with
fiber products.

Let
\[
\xymatrix@C=2.5cm@R=1.25cm{
Y'\ar[r]\ar[d]&X'\ar[d]\\
Y\ar[r]&X
}
\]
be a cartesian diagram in $F_j$, and let $W$ be an object of $F$ over
$k\in\ob I$. We have (\hyperref[VI.prop6.5]{6.5}):
\def\labelstyle{\scriptstyle}
$$\displaystyle{\Hhom(\pi(W),\pi(Y'))\simeq\mathop{\varinjlim_{
\xymatrix@R=.3cm@C=.3cm@M=.05cm{
            &j\\
        l\ar[ru]^{f}\ar[rd]_{g}&\\
            &k
}}}\Hhom_{F_{\ell}}(g^{\ast}(W),f^{\ast}(Y'))}$$

Since\pageoriginale $f^{\ast}:F_j\to F_{\ell}$ commutes with fiber
products, we have
$$
\Hom(g^{\ast}(W),f^{\ast}(Y'))\simeq
\Hom(g^{\ast}(W),f^{\ast}(Y))
\times_{\Hom(g^{\ast}(W),f^{\ast}(X))}
\Hom(g^{\ast}(W),f^{\ast}(X')).
$$
Using the commutation of filtered limits with fiber products (I 2)
and formula (\hyperref[VI.prop6.5]{6.5}), we obtain an isomorphism
$$
\Hom(\pi(W),\pi(Y'))\simeq\Hom(\pi(W),\pi(Y))
\times_{\Hom(\pi(W),\pi(X))}\Hom(\pi(W),\pi(X')){\quoi},
$$
which shows that the diagram
\[
\xymatrix@C=2.5cm@R=1.5cm{
\pi(Y')\ar[r]\ar[d]&\pi(X')\ar[d]\\
\pi(Y)\ar[r]&\pi(X)
}
\]
is cartesian in $\underrightarrow{F}$.

\setcounter{subsubsection}{4}
\subsubsection{}\etiquette[8.3.5]{VI.subsubsec8.3.5}
\emph{Proof of} \hyperref[VI.prop8.3.3]{8.3.3}. We first prove the
second assertion. It follows from \hyperref[VI.lem8.3.4]{8.3.4} that
the families in $\Cov(Y)$, $Y\in\underrightarrow{F}$, consist of
base-changeable morphisms, whence property (PTO). To show that these
families have property (PT1) (stability under base change), a series
of reductions like those in the proof of
\hyperref[VI.prop8.3.3]{8.3.3} reduces us to proving the following
assertion:

For every $i\in\ob I$, every $Y\in\ob F_i$, every
$(Y_{\alpha}\to Y)_{\alpha\in A}\in\Cov(Y)$, and every morphism
$u:Z\to Y$ of $F_i$, the family
$$
(\pi(Z)\times_{\pi(Y)}\pi(Y_{\alpha})\to\pi(Z))_{\alpha\in A}
$$
belongs to $\Cov(\pi(Z))$.

This assertion follows from the fact that the
$\pi_i:F_i\hookrightarrow F\xrightarrow{\pi}\underrightarrow{F}$
commute with fiber products (\hyperref[VI.lem8.3.4]{8.3.4}) and that
the families $X\mapsto\Cov(X)$ of $F_i$ are stable under base change.
Let us prove the first assertion. It is clear that the families
belonging to the $\Cov(Y)$, $Y\in\underrightarrow{F}$, are covering
for the coarsest topology $T$ making the functor $\pi$ continuous
(III 1). Thus the topology $T'$ generated by these families is
coarser than $T$. To show that it is finer than\pageoriginale $T$
(and hence equal to $T$), it is enough to show that every sheaf $M$
for $T'$ is such that $M\circ\pi$ is a sheaf on $F$, or equivalently
(\hyperref[VI.rem8.2.4]{8.2.4}) that $M\circ\pi_i$ is a sheaf on
$F_i$ for every $i\in\ob I$. Since the functors
$\pi_i:F_i\to\underrightarrow{F}$ commute with fiber products
(\hyperref[VI.lem8.3.4]{8.3.4}), $M$ is a sheaf for $T'$ if and only
if (II \hyperref[II.cor2.3]{2.3} and I \hyperref[I.prop2.12]{2.12}), for every $i\in\ob I$, every $Y\in\ob F_i$,
and every $(Y_{\alpha}\to Y)_{\alpha\in A}\in\Cov(Y)$, the sequence
of sets
$$
M(\pi_i(Y))\to\prod_{\alpha\in A}M(\pi_i(Y_{\alpha}))
\rightrightarrows\prod_{(\alpha,\beta)\in A\times A}
M\bigl(\pi_i(\fprod{Y_{\alpha}}{Y_{\beta}}{Y})\bigr)
$$
is exact, \ie (\loccit) if and only if $M\circ\pi_i$ is a sheaf on
$F_i$ for every $i\in\ob I$.


% END INCLUDED SGA4_English_current_buildable_source_20260906/source/06/drafts/SGA4_VI_LINES_5857_6042_English.texfrag


% BEGIN INCLUDED SGA4_English_current_buildable_source_20260906/source/06/drafts/SGA4_VI_LINES_6043_6336_English.texfrag
\begin{enonce*}{Proposition 8.3.6}\etiquette[8.3.6]{VI.prop8.3.6}
We use the notation and hypotheses of \hyperref[VI.subsubsec8.3.1]{8.3.1}
and \hyperref[VI.subsubsec8.3.2]{8.3.2}. If, for every $i\in\ob I$,
$X\mapsto\Cov(X)$ is a pretopology on $F_i$ \emph{(II \hyperref[II.def1.3]{1.3})}, and if for
every $X\in\ob F_i$ the covering families in $\Cov(X)$ are finite, then for
every $Y\in\underrightarrow{F}$ the covering families in $\Cov(Y)$ are
finite and $Y\mapsto\Cov(Y)$ is a pretopology on $\underrightarrow{F}$.
\end{enonce*}

\setcounter{subsubsection}{6}
\subsubsection{}\etiquette[8.3.7]{VI.subsubsec8.3.7}
The only thing to prove is that the families $Y\mapsto\Cov(Y)$ have property
(PT2) of (II \hyperref[II.def1.3]{1.3}) (stability under composition). Let $i\in\ob I$, let $Y$
be an object of $F_i$, and let
$(Y_\alpha\to Y)_{\alpha\in A}\in\Cov(Y)$. In addition, let
$i_\alpha$, $\alpha\in A$, be a family of objects of $I$; for every
$\alpha$, let $Z_\alpha$ be an object of $F_i$; and let
$(Z_{\alpha\beta}\xrightarrow{n_{\alpha\beta}}Z_\alpha)_{\alpha\beta\in
B_\alpha}$ be a family in $\Cov(Z_\alpha)$. Finally, for every
$\alpha\in A$, let an isomorphism
$\phi:\pi(Z_\alpha)\xrightarrow{\sim}\pi(Y_\alpha)$ be given. Returning to
the definition of the families in $\Cov(X)$ for
$X\in\ob\underrightarrow{F}$
(\hyperref[VI.subsubsec8.3.2]{8.3.2}), we see that we must prove that the
family
$$
(\pi (Z_{\alpha \beta}) \xrightarrow{\pi(n_{\alpha}) \circ \phi_{\alpha} \circ
  \pi(n_{\alpha \beta})} \pi(Y))_{\alpha \beta \in \mathop{\coprod}_A
  B_{\alpha}}
$$
belongs to $\Cov(\pi(Y))$. We shall do this after five reductions.

\setcounter{subsubsection}{7}
\subsubsection{}\etiquette[8.3.8]{VI.subsubsec8.3.8}
\emph{Reduction to the case in which, for every $\alpha\in A$, we have}
$\phi_\alpha=\pi(m_\alpha)$, $m_\alpha:Z_\alpha\to Y_\alpha$.
\emph{For every $\alpha\in A$, we have
$\phi_\alpha=\pi(m_\alpha)\pi(s_\alpha)^{-1}$
$(\hyperref[VI.prop6.5]{6.5})$, where $s_\alpha$ is a Cartesian morphism
over $f_\alpha:i'_\alpha\to i_\alpha$. Thus there are commutative diagrams}
\[
\xymatrix@C=2.5cm@R=1.25cm{
Z_{\alpha\beta} \ar[d]_{n_{\alpha\beta}} & f^{\ast}(Z_{\alpha\beta})
\ar[l]_{s_{\alpha\beta}} \ar[d]^{f^{\ast}_{\alpha}(n_{\alpha\beta})}\\
Z_{\alpha} &   f^{\ast}_{\alpha}(Z_{\alpha}) \ar[l]_{s_{\alpha}},  &
**[l]\alpha\beta\in \displaystyle{\coprod_A} B_{\alpha},
}
\]
which\pageoriginale give commutative diagrams
\[
\xymatrix@C=2.5cm@R=1.25cm@H=.2cm{
\pi(Z_{\alpha\beta})  \ar[r]^{\sim} \ar[d]^{\pi(n_{\alpha\beta})}&
\pi(f^{\ast}_{\alpha}(Z_{\alpha\beta})) \ar[d]^{\pi(f^{\ast}_{\alpha}(m_{\alpha\beta}))}\\
\pi(Z_{\alpha})  \ar[r]^{\displaystyle\mathop{\sim}^{\pi(S_\alpha)^{-1}}} & \pi(f^{\ast}_{\alpha}(Z_{\alpha}))
\ar[r]^{\pi(m_{\alpha})} & \pi(Y_\alpha),\quad
\alpha\beta\in \displaystyle{\coprod_A} B_{\alpha}.}
\]
Since the families
$(f^{\ast}_{\alpha}(Z_{\alpha\beta})\xrightarrow{f^{\ast}(n_{\alpha})}
f^{\ast}(Z_{\alpha}))$ belong to
$\Cov(f^{\ast}_{\alpha}(Z_{\alpha}))$, we may reduce to the case in which
$\phi_\alpha=\pi(m_\alpha)$ for every $\alpha$.

\subsubsection{}\etiquette[8.3.9]{VI.subsubsec8.3.9}

\emph{Reduction to the case in which, for every $\alpha\in A$,
$i_\alpha=j$ and} $p(m_\alpha)=k:j\to i$. Put
$g_\alpha=p(m_\alpha)$. Since $I$ is cofiltered and $A$ is finite, there
exist an object $j$ of $I$ and morphisms $h_\alpha:j\to i_\alpha$ such that,
for every pair $(\alpha,\alpha')$, we have
$g_\alpha h_\alpha=g_{\alpha'}h_{\alpha'}=k$. Using the inverse-image
functors $h_\alpha^{\ast}$ as in \hyperref[VI.subsubsec8.3.8]{8.3.8}, we
reduce to the stated case.

\subsubsection{}\etiquette[8.3.10]{VI.subsubsec8.3.10}

\emph{Reduction to the case in which} $i=j$ and $k:j\to i$ \emph{is the
identity}. For every $\alpha$, there is a commutative diagram
\[
\xymatrix@C=2.5cm@R=1.5cm{
Z_{\alpha}  \ar[r]^{m'_\alpha} &  k^{\ast}(Y_{\alpha}) \ar[d]_{k^{\ast}(n_{\alpha})}
\ar[r]^{t_{\alpha}} & Y_{\alpha} \ar[d]^{n_{\alpha}}\\
                    & k^{\ast}(Y)  \ar[r]^{t} & Y,
}
\]
where the morphisms $t$ and $t_\alpha$, $\alpha\in A$, are Cartesian,
$m'_\alpha$ lies over the identity of $j$, and
$t_\alpha\circ m'_\alpha=m_\alpha$. Replacing the family
$(Y_\alpha\to Y)\alpha\in A$ by the\pageoriginale family
$(k^{\ast}(Y_\alpha)\to k^{\ast}(Y))_{\alpha\in A}\in
\Cov(k^{\ast}(Y_\alpha))$, and the morphisms $m_\alpha$ by the morphisms
$m'_\alpha$, we reduce to the stated case.

\begin{enonce*}{Lemma 8.3.11}\etiquette[8.3.11]{VI.lem8.3.11}
Let $j$ be an object of $I$ and let $m:X'\to X$ be a morphism of $F_j$ such
that $\pi(m)$ is an isomorphism. There exists a morphism $\ell:j'\to j$ such
that $\ell^{\ast}(m)$ is an isomorphism.
\end{enonce*}

Let us show that there exist a morphism $f:i\to j$ and a morphism
$q:f^{\ast}(X)\to f^{\ast}(X')$ such that
$f^{\ast}(m)q=\id_{f^{\ast}(X)}$ and such that $\pi(q)$ is an isomorphism.
Indeed, there exists an isomorphism $n:\pi(X)\to\pi(X')$ such that
$\pi(m)\circ n=\id_{\pi(X)}$. There consequently exist
(\hyperref[VI.prop6.5]{6.5}) two morphisms
$X\xleftarrow{s_1}Y_1\xrightarrow{q_1}X'$, where $s_1$ is Cartesian, such
that $n=\pi(q_1)\pi(s_1)^{-1}$. Thus
$\pi(m)\circ\pi(q_1)=\pi(s_1)$. The morphisms $mq_1$ and $s_1$ from $Y_1$
to $X$ have the same image in $\underrightarrow{F}$. Consequently (\cf the
description of the morphisms in the inductive limit
(\hyperref[VI.prop6.5]{6.5})), there exists a Cartesian morphism
$s_2:Y_2\to Y_1$ such that $mq_1s_2=s_1s_2$. Put $q_1s_2=q_2$ and denote by
$s_3$ the Cartesian morphism $s_1s_2$. Thus $mq_2=s_3$. The morphism
$q_2:Y_2\to X'$ factors in an essentially unique manner as
$Y_2\xrightarrow{q}Y_3\xrightarrow{s_4}X'$, where $s_4$ is Cartesian and
$q$ lies over the identity of $p(Y_2)$. Thus there is a commutative diagram
\[\vcenter{
\xymatrix@C=2.25cm@R=1.5cm{
Y_3 \ar[d]_{s_4} & Y_2 \ar[l]_{q}\ar[d]^{s_3}\\
X'  \ar[r]^{m} & X.          .
}}\leqno{(8.3.11.1)}
\]\etiquette[8.3.11.1]{VI.eq8.3.11.1}
Put $p(s_4)=p(s_3)=f:i\to j$. We have
$f^{\ast}(X')\simeq Y_3$, $f^{\ast}(X')\simeq Y_2$, and
$f^{\ast}(m):Y_2\to Y_3$ makes the following diagram commutative:
\[
\xymatrix@C=2.25cm@R=1.5cm{
Y_3  \ar[r]^{f^{\ast}(m)} \ar[d]_{s_4} & Y_2 \ar[d]^{s_3}\\
X'  \ar[r]^{m} & X.
}
\]
Thus\pageoriginale $s_3=s_3(f^{\ast}(m)\circ q)$; since $s_3$ is Cartesian
and $f^{\ast}(m)\circ q$ lies over the identity of $p(Y_2)$, we have
$f^{\ast}(m)\circ q=\id_{Y_2}$. Moreover, the commutativity of diagram
\eqhyperref[VI.eq8.3.11.1]{8.3.11.1} shows that $\pi(q)$ is an isomorphism.

Now apply the preceding result to the morphism
$q:f^{\ast}(X)\to f^{\ast}(X')$: there exist a morphism $f':j'\to i$ and a
morphism $q':f'^{\ast}(X')\to f'^{\ast}(X)$ such that
$f'^{\ast}(q)q'=\id_{f'^{\ast}(f'^{\ast}(X'))}$. But on putting
$\ell^{\ast}=ff'$, there is an isomorphism
$f'^{\ast}f^{\ast}\simeq\ell^{\ast}$. Thus there is a morphism
$f'^{\ast}(q):\ell^{\ast}(X)\to\ell^{\ast}(X')$ and an isomorphism
$f'^{\ast}f^{\ast}\simeq\ell^{\ast}$. Thus there is a morphism
$f'^{\ast}(q):\ell^{\ast}(X)\to\ell^{\ast}(X')$ and a morphism
$q':\ell^{\ast}(X')\to\ell^{\ast}(X)$ such that
$\ell^{\ast}(m){f'}^{\ast}(q)=\id_{\ell^{\ast}(X)}$ and
${f'}^{\ast}(q)q'=\id_{\ell^{\ast}(X')}$. Consequently
$\ell^{\ast}(m)$ is an isomorphism.

\setcounter{subsubsection}{11}
\subsubsection{}\etiquette[8.3.12]{VI.subsubsec8.3.12}
\emph{End of the proof}. Since $I$ is cofiltered and $A$ finite, it follows
from \hyperref[VI.lem8.3.11]{8.3.11} that there exists an inverse-image
functor $\ell^{\ast}$ that takes all the morphisms $m_\alpha$,
$\alpha\in A$, to isomorphisms. We are thus reduced to the case in which the
$m_\alpha$ are isomorphisms. But then the covering families
$(Z_{\alpha\beta}\to Z)_{\alpha\beta\in B_\alpha}$ are deduced by the base
change $m_\alpha:Z_\alpha\to Y_\alpha$ from covering families
$(Y_{\alpha\beta}\to Y_\alpha)_{\alpha\beta\in B_\alpha}\in
\Cov(Y_\alpha)$. We are therefore reduced to the case in which
$Z_\alpha=Y_\alpha$ and $\Phi_\alpha=\id_{\pi(Y_\alpha)}$. But in this case
the family
$(Y_{\alpha\beta}\xrightarrow{n_\alpha\circ n_{\alpha\beta}}Y)_{
\alpha\beta\in\mathop{\coprod}_A B_\alpha}$ belongs to $\Cov(Y)$ (property
(PT2)). Hence its image under $\pi$ belongs to $\Cov(\pi(Y))$
(\hyperref[VI.subsubsec8.3.2]{8.3.2}).

\begin{enonce*}{Theorem 8.3.13}\etiquette[8.3.13]{VI.thm8.3.13}
Let $p:F\to I$ be a $𝒰$-topos fibered over an essentially small cofiltered
category. If, for every object $i$ of $I$, the fiber topos $F_i$ is coherent
$(\hyperref[VI.def2.3]{2.3})$, and if, for every morphism $f:i\to j$, the
morphism of topoi
$f\bigdot=(f^{\ast},f_{\ast}):F_i\to F_j$ is coherent
(\hyperref[VI.def3.1]{3.1}), then the topos $\limtop{I}F_i$ is coherent and,
for every object $i$ of $I$, the canonical morphism of topoi
$$
\mu_i : \limtop{I} F_i \to F_i
$$
is coherent.

Moreover, if $F_{\coh}$ denotes the full subcategory of $F$ defined by the
objects of $F$ that are coherent in their fiber
(\hyperref[VI.def1.13]{1.13}), then the category $F_{\coh}$ is fibered over
$I$, and $\underrightarrow{F}_{\coh}$ is canonically equivalent to the
category of coherent objects of $\limtop{I}F_i$.
\end{enonce*}

For every\pageoriginale $i\in\ob I$, denote by $(F_i)_{\coh}$ the $𝒰$-site
(endowed with the induced topology) of coherent objects of $F_i$
(\hyperref[VI.def1.13]{1.13}). Since $F_i$ is coherent, $F_i$ is the topos
of sheaves on $(F_i)_{\coh}$. For every morphism $f:i\to j$, the morphisms
$f\bigdot=(f^{\ast},f_{\ast})$ are coherent, and consequently
$f^{\ast}(F_j)_{\coh}\subset(F_i)_{\coh}$. Thus there is a fibered $𝒰$-site
$p_{\coh}:F_{\coh}\to I$ whose fiber sites are the sites $(F_i)_{\coh}$ and
whose associated fibered topos is equivalent to $p:F\to I$. Consequently
(\hyperref[VI.thm8.2.3]{8.2.3}),
\begin{equation*}
\limtop{I, f.} F_i \approxeq
(\underrightarrow{F}_{\coh})^{\sim}\tag{8.3.13.1}
\end{equation*}\etiquette[8.3.13.1]{VI.eq8.3.13.1}
Since the topoi $F_i$ are coherent, finite products and fiber products are
representable in $(F_i)_{\coh}$ (\hyperref[VI.prop2.2]{2.2}). Moreover, the
functors $f^{\ast}:(F_i)_{\coh}\to(F_j)_{\coh}$ are left exact. Finally,
since the objects of $(F_i)_{\coh}$ are quasi-compact, the finite covering
families form a pretopology. It then follows from
(\hyperref[VI.lem8.3.4]{8.3.4}) that finite products and fiber products are
representable in $\varprojlim_{I^{\circ},F}(F_i)_{\coh}$ (addendum to Lemma
\hyperref[VI.lem8.3.4]{8.3.4}: show that the canonical functors
$(F_i)_{\coh}\to\underrightarrow{F}_{\coh}$ take the final object to the
final object), and from (\hyperref[VI.prop8.3.6]{8.3.6}) that the objects of
$\underrightarrow{F}_{\coh}$ are quasi-compact. Consequently
(\hyperref[VI.subsubsec2.4.5]{2.4.5}), $\limtop{I}(F_i)$ is coherent.
Finally, the canonical morphisms
$\mu_i:\limtop{I}(F_i)\to F_i$ are deduced, by passage to the topoi of
sheaves, from the morphisms of sites (\hyperref[VI.thm8.2.3]{8.2.3})
$$
(F_i)_{\coh} \to \underrightarrow{F}_{\coh}
$$
and hence (\hyperref[VI.cor3.3]{3.3}) the morphisms $\mu_i$ are coherent.

Let us prove the last assertion. We have \eqhyperref[VI.eq8.3.13.1]{8.3.13.1}
$$
\limtop{ℐ} (F_i) \simeq
(\underrightarrow{F}_{\coh})^{\sim}{\quoi}.
$$
Let $X$ be a coherent object of $(F_{\coh})^{\sim}$. Since $X$ is
quasi-compact, there exists a finite covering family
$(Y_\alpha\to X)$, where
$Y_\alpha\in\ob\underrightarrow{F}_{\coh}$
(\hyperref[VI.def1.1]{1.1}). On taking an index $i\in\ob I$ sufficiently
small, we may suppose that the $Y_\alpha$ are the images under
$\mu_i:F_{i\coh}\to\underrightarrow{F}_{\coh}$ of a family $Y'_\alpha$;
hence, by taking the direct sum of the $Y'_\alpha$
(\hyperref[VI.cor1.15]{1.15}), we see that there exist an index $i\in\ob I$,
an object $Y'\in F_{i\coh}$, and a surjective morphism
$\mu_i(Y')\to X$. The equivalence relation
$R=\fprod{\mu_i(Y')}{\mu_i(Y')}{X}$ is then a coherent
object\pageoriginale because $(\underrightarrow{F}_{\coh})^{\sim}$ is a
coherent topos (\hyperref[VI.prop2.2]{2.2}). Let us first show that $R$ is an
object of $\underrightarrow{F}_{\coh}$. By definition, $R$ is a subobject of
$\mu_i(Y'\times Y')$, and since $R$ is coherent, we may, after changing the
index $i$, find an arrow $m:Z\to(Y'\times Y')$ of $F_{i\coh}$ such that
$\text{Im}(\mu_i(m))=R$. It follows that, since $F_i$ is coherent,
$\text{Im}(m)$ is coherent in $F_i$ (\hyperref[VI.cor1.17.1]{1.17.1}), and
hence belongs to $F_{i\coh}$. There consequently exist an index $i$ and a
diagram
$\def\labelstyle{\scriptstyle} \xymatrix@C=1cm{
R \ar@<3pt>[r]^{p_1}\ar@<-3pt>[r]_{p_2} & Y'
}
$
of $F_{i\coh}$ such that
$\def\labelstyle{\scriptstyle}
\xymatrix@C=1cm{
  \mu_i(R)\ar@<3pt>[r]^{\mu_i(p_1)} \ar@<-3pt>[r]_{\mu_i(p_2)}& \mu_i(Y')
}
$
is an equivalence relation and such that
$X=\mu_i(Y')/\mu_i(R)$. Put $X_i=\coker(p_1,p_2)$ and
$R'=\fprod{Y'}{Y'}{X'_i}$. There is a canonical morphism in $F_i$,
$u:R\to R'$. Moreover, $\mu_i^{\ast}(u)$ is an isomorphism. Thus, after
changing the index $i$, we may suppose that $R=R'$, \ie that $R$ is an
equivalence relation. But then $X_i$ belongs to $F_{i\coh}$
(\hyperref[VI.cor1.17.1]{1.17.1}), and consequently
$X=\mu_i(X_i)$ belongs to $\underrightarrow{F}_{\coh}$.

\begin{enonce*}{Corollary 8.3.14}\etiquette[8.3.14]{VI.cor8.3.14}
Let $I$ be an essentially small cofiltered category, and let $p:F\to I$ and
$q:G\to I$ be two $𝒰$-topoi fibered over $I$ such that, for every object
$i$ of $I$, the topoi $F_i$ and $G_i$ are coherent, and such that, for every
morphism $f:i\to j$, the morphisms $f.:F_i\to F_j$ and $G_i\to G_j$ are
coherent. In addition, let $m:F\to G$ be a Cartesian morphism of fibered
$𝒰$-topoi such that, for every $i\in\ob(I)$,
$m_i:F_i\to G_i$ is coherent. Then the morphism $\underleftarrow{m}$ deduced
from $m$ by passage to the projective limit is coherent.
\end{enonce*}

For every object $i$ of $I$, the morphism $m_i$ induces a functor
$m^{\ast}_{i,\coh}:G_{i,\coh}\to F_{i,\coh}$, whence a Cartesian functor
$m^{\ast}_{\coh}=G_{\coh}\to F_{\coh}$ that is a morphism from the fibered
site $F_{\coh}$ to the fibered site $G_{\coh}$ \hyperref[VI.thm7.4.13.1]{(7.4.13)}. The functor
$\varinjlim_{I^{\circ}}m^{\ast}_{\coh}:
\varinjlim_{I^{\circ}}G_{\coh}\to\varinjlim_{I^{\circ}}F_{\coh}$ is a
morphism of sites that, on passage to the corresponding topoi, gives the
morphism $\underleftarrow{m}$. The assertion then follows from
\hyperref[VI.thm8.3.13]{8.3.13} and \hyperref[VI.cor3.3]{3.3}.

\begin{enonce*}[remark]{Exercise 8.3.15}\etiquette[8.3.15]{VI.exe8.3.15}
With the notation of \hyperref[VI.thm8.3.13]{8.3.13}, show that if the $F_i$,
$i\in I$, are algebraic (\hyperref[VI.def2.3]{2.3}) and if the morphisms
$f.:F_i\to F_j$, $f\in\text{\rm Fl}(I)$, are coherent,\pageoriginale then
$\limtop{I}F_i$ is algebraic and
$(\limtop{I}F_i)_{\coh}\xrightarrow{\sim}\underrightarrow{F}_{\coh}$.
\end{enonce*}


% END INCLUDED SGA4_English_current_buildable_source_20260906/source/06/drafts/SGA4_VI_LINES_6043_6336_English.texfrag


% BEGIN INCLUDED SGA4_English_current_buildable_source_20260906/source/06/drafts/SGA4_VI_LINES_6337_6444_English.texfrag
\subsection{Example: local topoi}\etiquette[8.4]{VI.subsec8.4}

\subsubsection{}\etiquette[8.4.1]{VI.subsubsec8.4.1}

Let $X=(X_i)_{i\in I}$ be a topos fibered over a category $I$. One
deduces from it a topos fibered over the category $\Pro(I)$ (I \hyperref[I.sec8.10]{8.10})
by the formula
\begin{equation*}
X_{\sheaf{\alpha}}=\limtop{I_{\alpha}}X_{i\alpha}\tag{8.4.1.1}
\end{equation*}\etiquette[8.4.1.1]{VI.eq8.4.1.1}
for every pro-object $\sheaf{\alpha}=(i_{\alpha})$ of $I$ (to see
this, one may generalize Exercise I
\hyperref[VI.subsubsec8.2.8]{8.2.8} to pseudofunctors).

\subsubsection{}\etiquette[8.4.2]{VI.subsubsec8.4.2}

Let $E$ be a topos, and let $p:\text{Fl}(E)\to E$ be the fibered topos
considered in \hyperref[VI.exem7.3.1]{7.3.1}. Extending it as in
\hyperref[VI.subsubsec8.4.1]{8.4.1}, one obtains a fibered topos
$\bar{\text{Fl}(E)}\to\Pro(E)$. The category $\Point(E)$ maps by a
fully faithful functor into $\Pro(E)$ (IV \hyperref[IV.subsubsec6.8.5]{6.8.5}), whence, by base
change (\hyperref[VI.subsubsec7.1.9]{7.1.9}), a fibered topos
$\Loc(E)\to\Point(E)$. The fiber $\Loc_p(E)$ at a point $p$ of $E$
is called the \emph{localization of $E$ at the point} $p$. By the
preceding, this topos depends covariantly on the point $p$. By
definition,
\begin{equation*}
\Loc_p(E)=\Limtop_{X\in\Vois(p)}E_{/X}{\quoi},\tag{8.4.2.1}
\end{equation*}\etiquette[8.4.2.1]{VI.eq8.4.2.1}
where $\Vois(p)$ is the category of neighborhoods of $p$ (IV \hyperref[IV.subsec6.8]{6.8}).

\subsubsection{}\etiquette[8.4.3]{VI.subsubsec8.4.3}

Let $m:p'\to p$ be a morphism of $\Point(E)$ (IV 6). For every
$X\in\Vois(p)$ one thereby obtains a point $p'_X$ of $E_{/X}$ and
hence, by the universal property of $\limtop{}$
\hyperref[VI.subsec8.1]{8.1}, a point $\theta_p(m)$ of $\Loc_p(E)$.
This defines a functor
\begin{equation*}
\theta_p:\Point(E)_{/p}\to\Point(\Loc_p(E)){\quoi},\tag{8.4.3.1}
\end{equation*}\etiquette[8.4.3.1]{VI.eq8.4.3.1}
which is immediately seen to be an equivalence of categories. Thus
$\Point(\Loc_p(E))$ is canonically equivalent to the category of
generalizations of $p$ (IV \hyperref[IV.subsubsec7.1.8]{7.1.8}). There is moreover a canonical
equivalence of topoi
$$
\can\bigdot:\Loc_{\theta_p(m)}(\Loc_p(E))\to\Loc_{p'}(E){\quoi},
$$
which fits into a diagram commutative up to isomorphism
\[
\xymatrix@R=1.5cm{
\Loc_{\theta_p(m)}(\Loc_p(E))\ar[rr]^{\can\bigdot}\ar[rd]_{\phi}
&&\Loc_{p'}(E)\ar[ld]^{m}\\
&\Loc_p(E)&
}
\]
where $\phi$\pageoriginale is the canonical morphism of
\eqhyperref[VI.eq8.4.2.1]{8.4.2.1}.

\subsubsection{}\etiquette[8.4.4]{VI.subsubsec8.4.4}

Localized topoi seem to be of interest only for topoi arising from
algebraic geometry, or at least for topoi having suitable finiteness
properties. Thus, when $E$ is the topos of sheaves on a separated
topological space, one sees immediately (using, for example,
\hyperref[VI.thm8.2.9]{8.2.9}) that the localized topoi are point
topoi. Moreover, in this case the category $\Point(E)$ is discrete
and the situation described in \hyperref[VI.subsubsec8.4.2]{8.4.2}
is trivial. On the other hand, let $E=\Top(X)$ be the topos of sheaves
for the Zariski topology on a scheme $X$, let $x$ be a point of $X$,
let $O_x$ be the local ring of $X$ at $x$, and let
$Y=\spec(O_x)$. One finds that
$$
\Loc_x(E)=\Top(Y){\quoi}.
$$
For other examples arising from the \'etale topology, we refer to
Expos\'e VIII of this seminar.

\subsubsection{}\etiquette[8.4.5]{VI.subsubsec8.4.5}

When $E$ is locally coherent (\hyperref[VI.def2.3]{2.3}), the
localized topoi $\Loc_p(E)$ are coherent (in
\eqhyperref[VI.eq8.4.2.1]{8.4.2.1} one may restrict to the
$X\in\Vois(p)$ which are algebraic and coherent, and then apply
\hyperref[VI.thm8.3.13]{8.3.13}). For every morphism of points
$m:p'\to p$, the morphism of topoi
$m\bigdot:\Loc_{p'}(E)\to\Loc_p(E)$ is coherent
(\hyperref[VI.cor8.3.14]{8.3.14}).

\subsubsection{}\etiquette[8.4.6]{VI.subsubsec8.4.6}

A \emph{local topos} is a topos $X$ such that the functor
$\Gamma(X,-)=$ ``sections over $X$'' (IV \hyperref[IV.subsec4.3]{4.3}) is a fiber functor
(IV 6). The point corresponding to this fiber functor is called the
\emph{center} of the local topos. When $E$ is locally coherent, its
localized topoi are local topoi (\hyperref[VI.thm8.2.3]{1.2.3},
\hyperref[VI.prop8.5.2]{8.5.2}, and \hyperref[VI.cor8.5.7]{8.5.7}).
The center of $\Loc_p(E)$ is canonically isomorphic to
$\theta_p(\id_p)$ (\eqhyperref[VI.eq8.4.3.1]{8.4.3.1}). The image in
$E$ of the center of $\Loc_p(E)$ is isomorphic to $p$. The localized
topos $\Loc_p(E)$, equipped with the canonical morphism\pageoriginale
$\Loc_p(E)\to E$, is the solution of the universal (2-universal!)
problem of mapping local topoi into $E$ in such a way that the center
maps ``to'' $p$.

\subsubsection{}\etiquette[8.4.7]{VI.subsubsec8.4.7}

A number of questions concerning local topoi have not been addressed
by the editors. Thus, if $X$ is a local topos, the final object $e$ of
$X$ has a maximal open subtopos $U\ne e$. The complement $Y$ of $U$
is a local topos having no nontrivial open subtopos. Is such a topos a
point topos? Let $U$ be a topos and let
$X=(U,\Phi:U\to(\Ens))$ be a topos obtained by gluing. What conditions
must the gluing functor $\Phi$ satisfy for $X$ to be a local topos?


% END INCLUDED SGA4_English_current_buildable_source_20260906/source/06/drafts/SGA4_VI_LINES_6337_6444_English.texfrag


% BEGIN INCLUDED SGA4_English_current_buildable_source_20260906/source/06/drafts/SGA4_VI_LINES_6445_6610_English.texfrag
\subsection{Structure of the sheaves of a filtered projective limit
  of topoi}\etiquette[8.5]{VI.subsec8.5}

\subsubsection{}\etiquette[8.5.1]{VI.subsubsec8.5.1}

In this section, $F=(F_i)_{i\in I}$ is a $\Uu$-topos fibered over a
small cofiltered category $I$. Choose a bicleavage
(\hyperref[VI.subsubsec7.1.4]{7.1.4}) of $F$, \ie for every
$f\in\text{Fl}(I):i\to j$, choose morphisms of topoi
$f\bigdot:F_i\to F_j$ such that $f^{\ast}:F_j\to F_i$ is the
inverse-image functor for the fibered structure. There are then
canonical isomorphisms $c_{f,g}$ satisfying a cocycle condition
(\hyperref[VI.subsubsec7.1.3]{7.1.3}). Denote by $\underleftarrow{F}$
the projective limit of the fibered topos $F$
(\hyperref[VI.thm8.2.3]{8.2.3}), and for every $i\in\ob I$, denote by
\begin{equation*}
\mu_i:\underleftarrow{F}\to F_i{\quoi},\tag{8.5.1.1}
\end{equation*}\etiquette[8.5.1.1]{VI.eq8.5.1.1}
the canonical morphism (\hyperref[VI.subsubsec8.1.3]{8.1.3}). Denote
by $\Top(F)$ the total topos of $F$ \hyperlink{sga4.exp.VI.remark.7.4.3.item.3}{(7.4.3.3)}. By
\hyperref[VI.thm8.2.9]{8.2.9}, the canonical functor
$F\to\underrightarrow{F}$ defines a morphism of topoi
\begin{equation*}
Q:\underleftarrow{F}\to\Top(F){\quoi}.\tag{8.5.1.2}
\end{equation*}\etiquette[8.5.1.2]{VI.eq8.5.1.2}
The topos $\underleftarrow{F}$ is identified with
$\Hhom_{\Cart/I^{\circ}}(I^{\circ},F')$
(\hyperref[VI.thm8.2.9]{8.2.9}), and the topos $\Top(F)$ is identified
with $\Hhom_{I^{\circ}}(I^{\circ},F')$
(\hyperref[VI.prop7.4.7]{7.4.7}). Under these identifications, $Q$ is
the embedding morphism of
$\Hhom_{\cart/I^{\circ}}(I^{\circ},F')$ into
$\Hhom_{I^{\circ}}(I^{\circ},F')$
(\hyperref[VI.thm8.2.9]{8.2.9}), and for every object $M$ of
$\underleftarrow{F}$ we have
\begin{equation*}
Q_{\ast}(M)=(i\longmapsto\mu_{i\ast}(M)){\quoi}.\tag{8.5.1.3}
\end{equation*}\etiquette[8.5.1.3]{VI.eq8.5.1.3}
By \hyperlink{sga4.exp.VI.remark.7.4.3.item.4}{7.4.3.4}, for every $i\in\ob I$ there is a morphism of topoi
\begin{equation*}
\alpha_i:F_i\to\Top(F){\quoi}.\tag{8.5.1.4}
\end{equation*}\etiquette[8.5.1.4]{VI.eq8.5.1.4}
The diagram\pageoriginale
\[\vcenter{
\xymatrix@R=1.25cm{
{\underleftarrow{F}}\ar[rr]^{\mu_i}\ar[rd]_{Q}&&F_i\ar[ld]^{\alpha_i}\\
&\Top(F)&
}}\leqno{(8.5.1.5)}
\]\etiquette[8.5.1.5]{VI.eq8.5.1.5}
is not commutative in general (even up to isomorphism). But the
definition of the morphisms involved gives a canonical isomorphism
\begin{equation*}
\alpha_i^{\ast}Q_{\ast}\simeq\mu_{i\ast}.\tag{8.5.1.6}{\quoi}
\end{equation*}\etiquette[8.5.1.6]{VI.eq8.5.1.6}
When $i$ is a \emph{final object} of $I$, $\alpha_i$ is an embedding
having a retraction $\beta_i:\Top(F)\to F_i$
(\hyperref[VI.lem7.4.12]{7.4.12}), and the diagram
\[\vcenter{
\xymatrix@R=1.25cm{
\underleftarrow{F}\ar[rr]^{\mu_i}\ar[rd]_{Q}&&F_i\ar[ld]^{\beta_i}\\
&\Top(F)&,
}}\leqno{(8.5.1.7)}
\]\etiquette[8.5.1.7]{VI.eq8.5.1.7}
is commutative up to canonical isomorphism.

\begin{enonce*}{Proposition 8.5.2}\etiquette[8.5.2]{VI.prop8.5.2}
Let $j\mapsto M_j$ be an object of $\Top(F)$. There is a functorial
isomorphism
\begin{equation*}
Q^{\ast}(j\mapsto M_j)\simeq\mathop{\varinjlim}_{I^{\circ}}
\mu_j^{\ast}M_j{\quoi}.\tag{8.5.2.1}
\end{equation*}\etiquette[8.5.2.1]{VI.eq8.5.2.1}
\end{enonce*}
\noindent
An object of $\Top(F)$ (\hyperref[VI.subsubsec7.1.3]{7.1.3}) consists
of the data of an assignment $j\mapsto M_j$, $M_j\in\ob F_j$, and,
for every morphism $f:i\to j$, a morphism
\begin{equation*}
\beta_f:M_j\to f_{\ast}M_i\tag{8.5.2.2}
\end{equation*}\etiquette[8.5.2.2]{VI.eq8.5.2.2}
or, equivalently by adjunction, a morphism
\begin{equation*}
\beta'_f:f^{\ast}M_j\to M_i{\quoi},\tag{8.5.2.3}
\end{equation*}\etiquette[8.5.2.3]{VI.eq8.5.2.3}
where the morphisms $\beta'_f$ are subject to the condition that, for
every pair of composable morphisms
$i\xrightarrow{f}j\xrightarrow{g}k$, the following diagram be
commutative:
\[\vcenter{
\xymatrix@C=2.25cm@R=1.5cm{
f^{\ast}g^{\ast}M_k\ar[r]^{f^{\ast}\beta'_g}\ar@{-}[d]_{c^{\ast}_{g,f}}&
f^{\ast}M_j\ar@{->}[d]^{\beta'_f}\\
(gf)^{\ast}M_k\ar[r]^{\beta'_{gf}}&M_i
}}\leqno{(8.5.2.4)}
\]\etiquette[8.5.2.4]{VI.eq8.5.2.4}
Recall\pageoriginale moreover from
\eqhyperref[VI.eq8.1.3.1]{8.1.3.1} that, for every morphism
$f:i\to j$, there is an isomorphism
$b_f:f\bigdot\mu_i\xrightarrow{\sim}\mu_j$, and that for every pair
of composable morphisms $i\xrightarrow{f}j\xrightarrow{g}k$ there is
a commutative diagram
\[
\xymatrix@C=2cm@R=1.25cm{
g\bigdot f\bigdot\mu_i\ar[r]^{c_{g,f}}\ar[d]^{g\bigdot(b_f)}&
(gf)\bigdot\mu_i\ar[d]^{b_{gf}}\\
g\bigdot\mu_j\ar[r]^{b_g}&\mu_k
}\leqno{(8.5.2.5)}
\]\etiquette[8.5.2.5]{VI.eq8.5.2.5}
Let $f:i\to j$ be a morphism of $I$. Denote by
\begin{equation*}
t_f:\mu_j^{\ast}(M_j)\to\mu_i^{\ast}(M_i)\tag{8.5.2.6}
\end{equation*}\etiquette[8.5.2.6]{VI.eq8.5.2.6}
the composite
$\mu_j^{\ast}(M_j)\xrightarrow{b_f^{\ast}}
\mu_i^{\ast}(f^{\ast}(M_j))
\xrightarrow{\mu_i^{\ast}(\beta'_f)}\mu_i^{\ast}(M_i)$.
We leave it to the reader to verify that the commutativity of diagrams
\eqhyperref[VI.eq8.5.2.4]{8.5.2.4} and
\eqhyperref[VI.eq8.5.2.5]{8.5.2.5} implies that for every pair of
composable morphisms $i\to j\to k$, one has
\begin{equation*}
t_ft_g=t_{gf}{\quoi}.\tag{8.5.2.7}
\end{equation*}\etiquette[8.5.2.7]{VI.eq8.5.2.7}
We have therefore defined a functor
$I^{\circ}\to\underleftarrow{F}$, whose inductive limit
$\varinjlim_{I^{\circ}}\mu_i^{\ast}(M_i)$ can be considered. Let $N$
be an object of $\underleftarrow{F}$. We have
\begin{equation*}
\Hom(\varinjlim_{I^{\circ}}\mu_i^{\ast}(M_i),N)\simeq
\varprojlim_{I^{\circ}}\Hom(\mu_i^{\ast}(M_i),N){\quoi},\tag{8.5.2.8}
\end{equation*}\etiquette[8.5.2.8]{VI.eq8.5.2.8}
and hence, by adjunction, an isomorphism
\begin{equation*}
\Hom(\varinjlim_{I^{\circ}}\mu_i^{\ast}(M_i),N)\simeq
\varprojlim_{I^{\circ}}\Hom_{F_i}(M_i,\mu_{i\ast}(N)){\quoi}.
\tag{8.5.2.9}
\end{equation*}\etiquette[8.5.2.9]{VI.eq8.5.2.9}
We\pageoriginale now interpret the right-hand side of
\eqhyperref[VI.eq8.5.2.9]{8.5.2.9}. For every morphism $f:i\to j$,
denote by
$$
d_f:\Hom_{F_i}(M_i,\mu_{i\ast}(N))\to
\Hom_{F_j}(M_j,\mu_{j\ast}(N)){\quoi}
$$
the transition map of the projective system occurring in
\eqhyperref[VI.eq8.5.2.9]{8.5.2.9}. It follows from the definition of
the morphisms $t_f$ (\eqhyperref[VI.eq8.5.2.6]{8.5.2.6}) that $d_f$
associates with a morphism
$$
u_i:M_i\to\mu_{i\ast}(N)
$$
the morphism
$$
d_f(u_i):M_j\to\mu_{j\ast}(N){\quoi},
$$
obtained by composing the morphisms
$$
M_j\xrightarrow{\beta_f}f_{\ast}(M_i)
\xrightarrow{f_{\ast}(u_i)}f_{\ast}\mu_{i\ast}(N)
\xrightarrow{b_f}\mu_{j\ast}(N){\quoi}.
$$
Thus an element of the right-hand side of
\eqhyperref[VI.eq8.5.2.9]{8.5.2.9} is a family of morphisms
$u_i:M_i\to\mu_{i\ast}(N)$, $i\in\ob I$, such that for every
$f:i\to j$ the diagram
\[
\xymatrix@C=2cm@R=1.25cm{
M_j\ar[r]\ar[d]_{u_j}&f_{\ast}(M_i)\ar[d]^{f_{\ast}(u_i)}\\
\mu_{j\ast}(N)\ar[r]^{\sim}&f_{\ast}\mu_{i\ast}(N)
}
\]
is commutative, that is, a morphism from the section
$(i\mapsto M_i)\in\Hhom_{I^{\circ}}(I^{\circ},F')$ to the section
$Q_{\ast}(N)$. We therefore have a functorial isomorphism
$$
\Hom(\varinjlim_{I^{\circ}}\mu_i^{\ast}(M_i),N)
\simeq\Hom((i\mapsto M_i),Q_{\ast}(N)){\quoi};
$$
the proposition follows by adjunction.


% END INCLUDED SGA4_English_current_buildable_source_20260906/source/06/drafts/SGA4_VI_LINES_6445_6610_English.texfrag


% BEGIN INCLUDED SGA4_English_current_buildable_source_20260906/source/06/drafts/SGA4_VI_LINES_6611_6839_English.texfrag
\begin{enonce*}{Proposition 8.5.3}\etiquette[8.5.3]{VI.prop8.5.3}
If, for every $f:i\to j\in\text{\rm Fl}(I)$, the functors
$f_{\ast}:F_i\to F_j$ commute with small filtered inductive limits, then for
every section $(i\mapsto M_i)\in\Hom_I^{\circ}(I^{\circ},F')$ we have
\begin{equation*}
Q_{\ast} Q^{\ast} (i \mapsto M_i) \simeq i \mapsto
\varinjlim_{f : j \to i} f_{\ast} M_j{\quoi}.\tag{8.5.3.1}
\end{equation*}\etiquette[8.5.3.1]{VI.eq8.5.3.1}
\end{enonce*}

\noindent We\pageoriginale leave it to the reader to describe explicitly,
in terms of the $\beta_f$ \eqhyperref[VI.eq8.5.2.2]{8.5.2.2} and the
$c_{f,g}$ (\hyperref[VI.subsubsec8.5.1]{8.5.1}), the transition morphisms of
the inductive systems $(f:j\to i)\mapsto f_{\ast}(M_j)$ and the transition
morphisms of the section
$i\mapsto\varinjlim_{f:j\to i}f_{\ast}M_j$. Since the functors $f_{\ast}$
commute with filtered inductive limits, the section
$i\mapsto\varinjlim_{f:j\to i}f_{\ast}M_j$ is Cartesian. Moreover, there is
a natural morphism $u$ from $(i\mapsto M_i)$ to the section
$(i\mapsto\varinjlim_{f:j\to i}f_{\ast}M_j)$, and it is clear that every
morphism from $(i\mapsto M_i)$ to a Cartesian section factors uniquely
through $u$, giving the isomorphism \eqhyperref[VI.eq8.5.3.1]{8.5.3.1}.

\begin{enonce*}{Corollary 8.5.4}\etiquette[8.5.4]{VI.cor8.5.4}
Under the hypotheses of \hyperref[VI.prop8.5.3]{8.5.3}, the functor
$Q_{\ast}$ commutes with small filtered inductive limits.
\end{enonce*}

Let $\alpha\mapsto N^{\alpha}$ be a small filtered inductive system of
$\underleftarrow{F}$. Put $M^{\alpha}=Q_{\ast}(N^{\alpha})$, so that
$M^{\alpha}$ is a section $i\mapsto M_i^{\alpha}$. We have
$N^{\alpha}\simeq Q^{\ast}Q_{\ast}(N^{\alpha})$, and since $Q^{\ast}$
commutes with inductive limits, we have
$\varinjlim_{\alpha}N^{\alpha}\simeq
Q^{\ast}(\varinjlim_{\alpha}M^{\alpha})$. Inductive limits in
$\Top(F)=\Hom_{I^{\circ}}(I^{\circ},F')$ are computed fiber by fiber. Thus
$\varinjlim_{\alpha}M^{\alpha}\simeq
(i\mapsto\varinjlim_{\alpha}M_i^{\alpha})$. By
\eqhyperref[VI.eq8.5.3.1]{8.5.3.1}, we have
$$
Q_{\ast} (\varinjlim_{\alpha} {N}^{\alpha}) \simeq Q_{\ast}
Q^{\ast} (\varinjlim_{\alpha} {M}^{\alpha}) \simeq i \mapsto
\varinjlim_{f : j \to i} f_{\ast}
(\varinjlim_{\alpha} {M}^{\alpha}_j){\quoi}.
$$
Since the functors $f_{\ast}$ commute with filtered inductive limits, it
follows that
$$
Q_{\ast} (\varinjlim_{\alpha} {N}^{\alpha}) \simeq i \mapsto
\varinjlim_{f: j\to i} \varinjlim_{\alpha} f_{\ast}
(M^{\alpha}_j) \simeq i \mapsto \varinjlim_{\alpha} \varinjlim_{f:
  j\to i} f_{\ast}
(M^{\alpha}_j) \simeq \varinjlim_{\alpha} Q_{\ast} (N^{\alpha}){\quoi}.
$$

\begin{enonce*}{Corollary 8.5.5}\etiquette[8.5.5]{VI.cor8.5.5}
If the functors $f_{\ast}:F_i\to F_j$ commute with small filtered inductive
limits, then for every object $i$ of $I$ and every $M_i\in\ob F_i$ we have
\begin{equation*}
\mu_{i\ast} \mu_{i}^{\ast} (M_i) \simeq \varinjlim_{f : j \to
      i} f_{\ast} f^{\ast}(M_i){\quoi}.\tag{8.5.5.1}
\end{equation*}\etiquette[8.5.5.1]{VI.eq8.5.5.1}
\end{enonce*}
\noindent
After making the base change $I/i\to I$, we may suppose that $i$ is a final
object of $I$ (\hyperref[VI.lem8.2.1]{8.2.1}). The morphism
$\mu_i:\underleftarrow{F}\to F_i$ is then the composite of
the\pageoriginale morphisms (\hyperref[VI.thm8.2.9]{8.2.9}) and
(\hyperref[VI.lem7.4.12]{7.4.12}):
\begin{align*}
Q &: \underleftarrow{F} \to F^{\sim}{\quoi},\\
\beta_i &: (\alpha_{i!}, \alpha^{\ast}_i) : F^{\sim} \to F_i{\quoi}.
\end{align*}
Thus $\mu_i^{\ast}(M_i)\simeq Q^{\ast}\beta_i^{\ast}(M_i)$, where
$\beta_i^{\ast}(M_i)$ is the section
$(f:i\to j)\mapsto f^{\ast}(M_i)$ whose transition morphisms are deduced
from the $c_{f,g}$. The assertion then follows from
\hyperref[VI.prop8.5.3]{8.5.3}.

\begin{enonce*}{Corollary 8.5.6}\etiquette[8.5.6]{VI.cor8.5.6}
Under the hypotheses of \hyperref[VI.cor8.5.5]{8.5.5}, the functors
$$
\mu_{i\ast} : \underleftarrow{F} \to F_i
$$
commute with filtered inductive limits.
\end{enonce*}

The functor $\mu_{i\ast}$ is the composite of the functor $Q_{\ast}$, which
commutes with filtered inductive limits, and the functor ``restriction to
the fiber at $i$'', which commutes with inductive limits.

\begin{enonce*}{Corollary 8.5.7}\etiquette[8.5.7]{VI.cor8.5.7}
Under the hypotheses of \hyperref[VI.prop8.5.3]{8.5.3}, let $i$ be an object
of $I$ and $X$ an object of $F_i$ such that the functor $\Hom(X,-)$ on $F_i$
commutes with filtered inductive limits. Then the functor
$\Hom(\mu_i^{\ast}(X),-)$ commutes with filtered inductive limits, and for
every object $M_i$ of $F_i$ we have
\begin{equation*}
\Hom(\mu^{\ast}_i(X), \mu^{\ast}_i(M_i)) \simeq \varinjlim_{f
: j \to i} \Hom(f^{\ast}(X), f^{\ast}(M_i))\tag{8.5.7.1}{\quoi}.
\end{equation*}\etiquette[8.5.7.1]{VI.eq8.5.7.1}
\end{enonce*}
\noindent
For every object $(j\mapsto M_j)$ of $\Top(F)$, we have
\begin{equation*}
\Hom(\mu^{\ast}_i(X), Q^{\ast}(j \mapsto M_j)) \simeq \varinjlim_{f
: j \to i} \Hom(f^{\ast}(X), M_{j}){\quoi}.\tag{8.5.7.2}
\end{equation*}\etiquette[8.5.7.2]{VI.eq8.5.7.2}

The functor $\Hom(\mu_i^{\ast}(X),-)$ is isomorphic to the functor
$\Hom(X,\mu_{i\ast}(-))$, and the latter commutes with filtered inductive
limits (\hyperref[VI.cor8.5.4]{8.5.4}). By
\eqhyperref[VI.eq8.5.5.1]{8.5.5.1}, we have
$$
\Hom(\mu^{\ast}_i(X), \mu^{\ast}_i (M_i)) \simeq \Hom (X, \mu_{i\ast}\mu^{\ast}_i (M_i)) \simeq  \Hom(X, \varinjlim_{f : j \to
  i} f_{\ast} f^{\ast} (M_i)){\quoi},
$$
which gives formula \eqhyperref[VI.eq8.5.7.1]{8.5.7.1}. Likewise, by
\eqhyperref[VI.eq8.5.3.1]{8.5.3.1}, we have
$$
\Hom(\mu^{\ast}_i (X), Q^{\ast}(j \mapsto M_j)) \simeq \Hom(X,
\mu_{i\ast} Q^{\ast} (j \mapsto M_j)) \simeq \Hom(X,
\varinjlim_{f : j \to i} f_{\ast} (M_j)){\quoi},
$$
which\pageoriginale gives formula \eqhyperref[VI.eq8.5.7.2]{8.5.7.2}.

\begin{enonce*}{Proposition 8.5.8}\etiquette[8.5.8]{VI.prop8.5.8}
Let $G\to I$ be a fibered $𝒰$-topos and $m:F\to G$ a Cartesian morphism of
fibered topoi \hyperref[VI.def7.1.5]{(7.1.5)}. For $G$ we use the notation introduced in
\hyperref[VI.def8.1.1]{8.1.1}. Let
$\underleftarrow{m}:\underleftarrow{F}\to\underleftarrow{G}$ be the
morphism deduced from $m$ by passage to the projective limit
(\hyperref[VI.subsubsec8.1.4]{8.1.4}). The diagram of topoi and morphisms of
topoi
\[\vcenter{
\xymatrix@C=2.25cm@R=1.25cm{
\underset{\leftarrow}{F}  \ar[r]^{Q} \ar[d]^{\underset{\leftarrow}{m}}& \Top(F)\ar[d]^{m^{\sim}}\\
\underset{\leftarrow}{G}  \ar[r]^{Q} & \Top(G)
}}\leqno{(8.5.8.1)}
\]\etiquette[8.5.8.1]{VI.eq8.5.8.1}
is commutative up to isomorphism. For every object $N$ of
$\underleftarrow{F}$, we have
\begin{equation*}
\underleftarrow{m}_{\ast} (N) \simeq \varinjlim_{I^{\circ}} \mu_j^{\ast}
m_{j\ast} \mu_{j\ast} (N){\quoi}.\tag{8.5.8.2}
\end{equation*}\etiquette[8.5.8.2]{VI.eq8.5.8.2}
\end{enonce*}

The commutativity of diagram \eqhyperref[VI.eq8.5.8.1]{8.5.8.1} follows
immediately from the definitions (\hyperref[VI.subsubsec8.2.8]{8.2.8} and
\hyperref[VI.subsubsec8.5.1]{8.5.1}). We have
$\underleftarrow{m}_{\ast}(N)\simeq
Q^{\ast}Q_{\ast}\underleftarrow{m}_{\ast}(N)$ and, by the commutativity of
\eqhyperref[VI.eq8.5.8.1]{8.5.8.1},
$\underleftarrow{m}_{\ast}(N)\simeq
Q^{\ast}m^{\sim}_{\ast}Q_{\ast}(N)$. Consequently, there is an isomorphism
$\underleftarrow{m}_{\ast}(N)\simeq
Q^{\ast}m^{\sim}_{\ast}(j\mapsto\mu_{j\ast}(N))$. The functor
$m_{\ast}^{\sim}$ is simply the functor
$\Hom_{I^{\circ}}(I^{\circ},m_{\ast})$
(\hyperref[VI.prop7.4.10]{7.4.10}), and consequently there is a canonical
isomorphism
$\underleftarrow{m}_{\ast}(N)\simeq
Q^{\ast}(j\mapsto m_{j\ast}\mu_{j\ast}(N))$. Formula
\eqhyperref[VI.eq8.5.8.2]{8.5.8.2} then follows from
\hyperref[VI.prop8.5.2]{8.5.2}.

\begin{enonce*}{Proposition 8.5.9}\etiquette[8.5.9]{VI.prop8.5.9}
Under the hypotheses and with the notation of \hyperref[VI.prop8.5.8]{8.5.8},
suppose in addition that, for every object $i$ of $I$, the functor
$m_{i\ast}:F_i\to G_i$ commutes with filtered inductive limits, and that for
every morphism $f:i\to j$ of $I$, the functor $f_{\ast}:F_i\to F_j$ commutes
with filtered inductive limits. Then the functor
$\underleftarrow{m}_{\ast}:\underleftarrow{F}\to\underleftarrow{G}$ commutes
with filtered inductive limits, and for every object $(i\mapsto M_i)$ of
$\Top(F)$ there is a canonical isomorphism
\begin{equation*}
\underleftarrow{m}_{\ast} Q^{\ast} (i \mapsto M_i) \simeq
\varinjlim_{I^{\circ}} \mu^{\ast}_j m_{j\ast} (M_j){\quoi}.\tag{8.5.9.1}
\end{equation*}\etiquette[8.5.9.1]{VI.eq8.5.9.1}
\end{enonce*}

The first\pageoriginale assertion follows from
\eqhyperref[VI.eq8.5.8.2]{8.5.8.2} and
\hyperref[VI.cor8.5.6]{8.5.6}. By \eqhyperref[VI.eq8.5.8.2]{8.5.8.2} and
\eqhyperref[VI.eq8.5.3.1]{8.5.3.1}, there is a functorial isomorphism
$$
\underleftarrow{m}_{\ast} Q^{\ast} (i \mapsto M_i) \simeq
\varinjlim_i \mu^{\ast}_{i} m_{i\ast} (\varinjlim_{f
  : j \to i} f_{\ast} (M_{j})){\quoi}.
$$
Using the commutation of the $m_{i\ast}$ with filtered inductive limits and
the isomorphisms $m_{i\ast}f_{\ast}\simeq f_{\ast}m_{j\ast}$
(\hyperref[VI.subsubsec7.1.6]{7.1.6}), we obtain
$$
\underleftarrow{m}_{\ast} Q^{\ast} (i \mapsto M_i) \simeq
\varinjlim_i \mu^{\ast}_{i} (\varinjlim_{f : j \to
  i} f_{\ast} m_{j\ast} (M_j)){\quoi}.
$$
Since the functors $\mu_i^{\ast}$ commute with inductive limits, it follows
that
$$
\underleftarrow{m}_{\ast}Q^{\ast}(i\mapsto M_{i})\simeq \varinjlim_{i}
\varinjlim_{f:j\to i} \mu^{\ast}_{i}f_{\ast}m_{j\ast}(M_{j}){\quoi}.
$$
This last object can be interpreted as an inductive limit over the category
$\text{Fl}(I)^{\circ}$, where $\text{Fl}(I)$ is the category of morphisms of
$I$. Let $\phi:I\to\text{Fl}(I)$ be the functor that associates with every
object $i$ of $I$ the identity morphism of $i$. The functor
$\phi^{\circ}:I^{\circ}\to\text{Fl}(I^{\circ})$ is cofinal, and consequently
(I \hyperref[I.sec8.1]{8.1}) there is a canonical isomorphism
$$
\underleftarrow{m}_{\ast} Q^{\ast} (i \mapsto M_i) \simeq
\varinjlim_{I^{\circ}} \mu^{\ast}_{j} m_{j\ast}(M_j){\quoi}.
$$

\begin{enonce*}{Corollary 8.5.10}\etiquette[8.5.10]{VI.cor8.5.10}
Under the conditions of \hyperref[VI.prop8.5.9]{8.5.9}, for every object $i$
of $I$ and every object $M_i$ of $F_i$, there is a canonical isomorphism
\begin{equation*}
\underleftarrow{m}_{\ast} \mu^{\ast}_i (M_i) \simeq
\varinjlim_{f:j\to i} \mu^{\ast}_{j} m_{j\ast}f^{\ast}
(M_{i}){\quoi}.\tag{8.5.10.1}
\end{equation*}\etiquette[8.5.10.1]{VI.eq8.5.10.1}
\end{enonce*}

The proof is analogous to that of Corollary \hyperref[VI.cor8.5.5]{8.5.5}.

\begin{enonce*}[remark]{Remark 8.5.11}\etiquette[8.5.11]{VI.rem8.5.11}
Recall that the direct-image functor under a coherent morphism between
coherent topoi commutes with filtered inductive limits
(\hyperref[VI.thm5.1]{5.1}). Consequently, Propositions
\hyperref[VI.prop8.5.3]{8.5.3} through \hyperref[VI.cor8.5.7]{8.5.7}, and
\hyperref[VI.prop8.5.9]{8.5.9}, \hyperref[VI.cor8.5.10]{8.5.10}, apply when
the fibered topoi under consideration are coherent and the morphisms of
fibered topoi under consideration are coherent.
\end{enonce*}


% END INCLUDED SGA4_English_current_buildable_source_20260906/source/06/drafts/SGA4_VI_LINES_6611_6839_English.texfrag


% BEGIN INCLUDED SGA4_English_current_buildable_source_20260906/source/06/drafts/SGA4_VI_LINES_6840_6915_English.texfrag
\subsection{Ringed fibered $𝒰$-topoi}\etiquette[8.6]{VI.subsec8.6}

\subsubsection{}\etiquette[8.6.1]{VI.subsubsec8.6.1}

A fibered $𝒰$-topos is called \emph{ringed} if it is equipped with a
sheaf of rings on the total site (\hyperref[VI.def7.4.1]{7.4.1}). Let
$p : F \to I$ be a fibered $𝒰$-topos. Choose a bisplitting
\pageoriginale (\hyperref[VI.subsubsec7.1.4]{7.1.4}). It follows from
\hyperref[VI.prop7.4.7]{7.4.7} that giving a sheaf of rings on $F$ is
equivalent to giving, for every object $i$ of $I$, a ring $A_i$ of
$F_i$, and, for every morphism $f : i \to j$, a morphism
$\phi_f : A_j \to f_{\ast} A_i$, the family of the $\phi_f$ being
subject to the compatibility conditions made explicit in \loccit.
The morphisms of topoi $f. : F_i \to F_j$ are therefore morphisms of
topoi ringed respectively by $A_i$ and $A_j$ (IV 11), and the ringed
structure on $F$, together with the choice of the morphisms $f$,
gives a pseudofunctor from $I$ to the 2-category of ringed topoi.
Conversely, given such a pseudofunctor, one can reconstruct the ringed
fibered $𝒰$-topos from which it arises.

\subsubsection{}\etiquette[8.6.2]{VI.subsubsec8.6.2}

As in \hyperref[VI.subsec8.1]{8.1}, one can define the notion of the
projective limit of a ringed $𝒰$-topos $(F_i, A_i, i \in I)$. We
shall suppose that this immediate generalization has been made, and
we shall freely apply to this situation the language introduced in
\hyperref[VI.subsec8.1]{8.1}. When it exists, this projective limit is
a ringed $𝒰$-topos $(H, B)$, and, for every object $i$ of $I$, we
have morphisms of ringed topoi
$\mu_i : (H, B) \to (F_i, A_i)$. If the underlying (unringed) fibered
topos $F$ admits a projective limit $\limtop{I} F_i$, and if the
inductive system of rings $i \mapsto \mu_i^{\ast}(A_i)$ admits an
inductive limit $\varinjlim_{I^{\circ}}\mu_i^{\ast} A_i$ in the
category of rings of $\limtop{I} F_i$ (as is always the case if $I$
is small), then the ringed topos
($\limtop{I} F_i$, $\varinjlim_{I^{\circ}} \mu^{\ast}_i A_i$) is
evidently a projective limit of the ringed fibered topos under
consideration.

\subsubsection{}\etiquette[8.6.3]{VI.subsubsec8.6.3}

The formulas of \hyperref[VI.subsec8.5]{8.5}, established for sheaves
of sets, are valid for abelian sheaves. They are also valid for
sheaves of modules, provided the inverse-image functors occurring in
them are interpreted as inverse-image functors in the sense of
sheaves of modules (IV \hyperref[IV.def1.1]{1.1}). We leave the verification to the reader.

\subsubsection{}\etiquette[8.6.4]{VI.subsubsec8.6.4}

Let $(p : F \to I, A)$ be a ringed fibered $𝒰$-topos over a
cofiltered category. We call $(p : F \to I, A)$ a ringed fibered
$𝒰$-topos flat on the right (\resp on the left) if, for every morphism
$f : i \to j$, the
morphism $f. : (F_i, A_i) \to (F_j, A_j)$ is a morphism of ringed
topoi flat on the right (\resp on the left)
(\hyperrefext[V][V.def1.8]{1.8}). Suppose that $I$ is essentially
small. Then the morphism of ringed topoi
$Q : (\limtop{I} F_i,
\varinjlim_{I^{\circ}} \mu^{\ast}_i (A_i)) \to
(\Hom_{I^{\circ}}(I^{\circ}, F'), A)$\pageoriginale
(\hyperref[VI.subsubsec8.5.1]{8.5.1}) is flat on the right (\resp on
the left), as follows from the definition and from
\hyperref[VI.prop8.5.2]{8.5.2}; this holds without any flatness
hypothesis on the transition morphisms $f.$. If, in addition,
$(p : F \to I, A)$ is a ringed fibered $𝒰$-topos flat on the right
(\resp on the left), the morphisms
$\mu_i : (\limtop{I} F_i,
\varinjlim_{I^{\circ}} \mu^{\ast}_i A_i) \to (F_i, A_i)$ are flat on
the right (\resp on the left). This follows from the fact that the
inverse image of a flat module is flat
(\hyperrefext[V][V.cor1.7.1]{1.7.1}), and that a filtered inductive
limit of flat modules is flat. The same arguments prove the following
fact: if $m : (F_i, A_i, i \in I) \to (G_i, A_i, i \in I)$ is a
cartesian morphism of ringed fibered $𝒰$-topoi over an essentially
small cofiltered category, and if, for every object $i \in \ob (I)$,
the morphism $m_i$ is flat on the right (\resp on the left), then the
morphism $\underleftarrow{m}$ deduced from $m$ by passage to the
projective limit is flat on the right (\resp on the left).

% END INCLUDED SGA4_English_current_buildable_source_20260906/source/06/drafts/SGA4_VI_LINES_6840_6915_English.texfrag


% BEGIN INCLUDED SGA4_English_current_buildable_source_20260906/source/06/drafts/SGA4_VI_LINES_6916_7181_English.texfrag
\subsection{Cohomology of sheaves on a projective limit of
  topoi}\etiquette[8.7]{VI.subsec8.7}

\subsubsection{}\etiquette[8.7.1]{VI.subsubsec8.7.1}

In this section, $I$ denotes a small cofiltered category,
$(p:F\to I,A)$ and $(q:G\to I,B)$ two ringed fibered $𝒰$-topoi, and
$m:(p:F\to I,A)\to(q:G\to I,B)$ a morphism of ringed fibered
$𝒰$-topoi. We suppose that, for every morphism $f:i\to j$ of $I$, the
derived functors $R^nf_{\ast}$, $n\in\NN$, of the functor
$f_{\ast}:\Mod(F_i,A_i)\to\Mod(F_j,A_j)$ on modules commute with
filtered inductive limits, and that, for every object $i\in I$, the
derived functors $R^nm_{i\ast}$ of the functor
$m_{i\ast}:(F_i,A_i)\to(G_i,B_i)$ on modules\footnote{To fix ideas,
we shall take left modules.} commute with filtered inductive limits.
Recall that these hypotheses hold when the categories $F_i,G_i$,
$(i\in\ob(I))$, are \emph{coherent} and when the morphisms
$f:F_i\to F_j$, $(f:i\to j\in \text{Fl}(I))$, and
$m_i:F_i\to G_i$, $(i\in I)$, are \emph{coherent}
(\hyperref[VI.cor5.2]{5.2}). We use the notation of
\hyperref[VI.subsubsec8.5.1]{8.5.1}. In addition, we denote by
$\underrightarrow{A}$ (\resp $\underrightarrow{B}$) the sheaf
$\varinjlim\mu^{\ast}_i(A_i)$ (\resp
$\varinjlim\mu^{\ast}_i(B_i)$). Finally, for modules, we use the
notation $\mu^{-1}_i$ for inverse image in the sense of abelian
sheaves, reserving $\mu_i^{\ast}$ for inverse image in the sense of
modules (IV 11).

\begin{enonce*}{Lemma 8.7.2}\etiquette[8.7.2]{VI.lem8.7.2}
Let\pageoriginale $j\mapsto M_j$ be an injective $A$-Module of
$\Top(F)$. Then $Q^{\ast}(j\mapsto M_j)$ is an $A$-Module acyclic for
$\underleftarrow{m}_{\ast}$, and, for every $i\in\ob(I)$, $M_i$ is
flasque.
\end{enonce*}

Let $i\in\ob(I)$, and let $e_i$ be a final object of $F_i$. The
fibered $𝒰$-topos $F_{/e_i}\to I_{/i}$ is deduced from $F\to I$ by
the base change $I_{/i}\to I$. Since the sheaf $j\mapsto M_j$ is
injective, it is flasque (\hyperrefext[V][V.subsec4.6]{4.6}). Its
restriction to the localized topos $F_{/e_i}$ is flasque
(\hyperrefext[V][V.subsec4.12]{4.12}). The restriction functor from
$F_{/e_i}$ to $F_i$ is a direct image functor for a morphism of topoi
(\hyperref[VI.lem7.4.12]{7.4.12}). It therefore carries flasque
Modules to flasque Modules (\hyperrefext[V][V.prop5.2]{5.2}). Thus
$M_i$ is flasque.

Let us prove the first assertion of the lemma. Set
$N'=Q^{\ast}(j\mapsto M_j)$. By
\eqhyperref[VI.eq8.5.3.1]{8.5.3.1}, for every $i\in\ob(I)$ we have
$$
\mu_{i\ast}(N') = \mu_{i\ast} Q^{\ast} (j \mapsto M_j) \simeq
\varinjlim_{f : j \to i} f_{\ast} (M_j){\quoi}.
$$
Using the hypothesis of \hyperref[VI.subsubsec8.7.1]{8.7.1}, we see
that the sheaf $N'$ has the following property:

(P)~For every object $i\in\ob(I)$, $\mu_{i\ast}(N')$ is
$m_{i\ast}$-acyclic, and, for every morphism $f:i\to j$ of $I$,
$\mu_{i\ast}(N')$ is $f_{\ast}$-acyclic.

It is enough to show that every object $N'$ having property (P) is
$\underleftarrow{m}_{\ast}$-acyclic. Since injectives have property
(P), it is enough to verify the following two properties
(\hyperrefext[V][V.prop0.4]{0.4}): for every exact sequence
$0\to N'\to N\to N''\to0$, where $N'$ and $N$ have property (P),
(a) $N''$ has property (P), and (b)
$m_{\ast}(N)\to m_{\ast}(N'')$ is an epimorphism. Verification of (a)
is immediate. Let us verify (b). Set
$K=\coker(Q_{\ast}(N')\to Q_{\ast}(N))$, so that $K$ is the section
$(i\mapsto K_i=\coker(\mu_{i\ast}(N')\to\mu_{i\ast}(N)))$. For every
$f:i\to j$, there is a commutative diagram
\[
\xymatrix{
0  \ar[r] & f_{\ast}\mu_{i\ast}(N') \ar[r] & f_{\ast}\mu_{i\ast}(N)\ar[r]
 & f_{\ast}(K_i) \ar[r]& 0\\
0  \ar[r] & \mu_{j\ast}(N')\ar[u] \ar[r] & \mu_{j\ast}(N) \ar[u] \ar[r] & K_j
\ar[r] \ar[u]&0.
}
\]
Since\pageoriginale $\mu_{i\ast}(N')$ is $f_{\ast}$-acyclic, the
top row is exact. Since $i\mapsto\mu_{i\ast}(N')$ and
$i\mapsto\mu_{i\ast}(N)$ are Cartesian sections, the first two
vertical morphisms are isomorphisms. Consequently, the canonical
morphism $K_j\to f_{\ast}(K_i)$ is an isomorphism, and
$i\mapsto K_i$ is a Cartesian section. Thus $K$ is of the form
$Q_{\ast}(K')$, and the canonical morphisms $Q_{\ast}(N)\to K$ and
$K\to Q_{\ast}(N'')$ are of the form $Q_{\ast}(u)$ and
$Q_{\ast}(v)$, respectively, because $Q_{\ast}$ is fully faithful.
Since $Q^{\ast}$ is exact and $Q^{\ast}Q_{\ast}$ is isomorphic to the
identity, $v$ is an isomorphism, and hence $Q_{\ast}(v)$ is an
isomorphism. Thus the sequence
$0\to Q_{\ast}(N')\to Q_{\ast}(N)\to Q_{\ast}(N'')\to0$ is exact,
and consequently, for every object $i$ of $I$, the sequence
$0\to\mu_{i\ast}(N')\to\mu_{i\ast}(N)\to
\mu_{i\ast}(N'')\to0$ is exact. Since $\mu_{i\ast}(N')$ is
$m_{i\ast}$-acyclic, the sequence
$0\to m_{i\ast}\mu_{i\ast}(N')\to
m_{i\ast}\mu_{i\ast}(N)\to
m_{i\ast}\mu_{i\ast}(N'')\to0$ is exact. Moreover, there is an
isomorphism
$\underleftarrow{m}_{\ast}\simeq
\varinjlim_{I^{\circ}}\mu^{-1}_jm_{j\ast}\mu_{j\ast}$
\eqhyperref[VI.eq8.5.8.2]{8.5.8.2}. The functor $\mu^{-1}_j$ is exact,
and filtered inductive limits are exact. Therefore the sequence
$0\to\underleftarrow{m}_{\ast}(N')\to
\underleftarrow{m}_{\ast}(N)\to
\underleftarrow{m}_{\ast}(N'')\to0$ is exact.

\begin{enonce*}{Theorem 8.7.3}\etiquette[8.7.3]{VI.thm8.7.3}
The notation and hypotheses are those of
\hyperref[VI.subsubsec8.7.1]{8.7.1}. For every integer $n\in\NN$,
the derived functor $R^n\underleftarrow{m}_{\ast}$ commutes with
filtered inductive limits, and for every $A$-Module
$j\mapsto M_j$ there is a functorial isomorphism
\begin{equation*}
R^{n} \underleftarrow{m}_{\ast} Q^{\ast} (j \mapsto M_j) \simeq
\varinjlim_{I^{\circ}}
\mu^{\ast}_j R^n m_{j\ast} (M_j)\tag{8.7.3.1}{\quoi}.
\end{equation*}\etiquette[8.7.3.1]{VI.eq8.7.3.1}
\end{enonce*}

We have $Q^{-1}(A)=\underrightarrow{A}$, and hence inverse image for
Modules is isomorphic to inverse image for abelian sheaves. It follows
from \hyperref[VI.eq8.5.8.2]{8.5.8.2} that, for every section
$j\mapsto M_j$, there is a canonical isomorphism
$\varinjlim_{I^{\circ}}\mu^{-1}_j(M_j)\simeq
\varinjlim_{I^{\circ}}\mu^{\ast}_j(M_j)$. Moreover, under the
hypotheses of \hyperref[VI.subsubsec8.7.1]{8.7.1}, there is a
canonical isomorphism \eqhyperref[VI.eq8.5.9.1]{8.5.9.1}
$\underleftarrow{m}_{\ast}Q^{\ast}(j\mapsto M_j)\simeq
\varinjlim_{I^{\circ}}\mu^{-1}_jm_{j\ast}(M_j)$. Since, for every
injective $j\mapsto M_j$, the $M_j$ are flasque and
$Q^{\ast}(j\mapsto M_j)$ is acyclic for
$\underleftarrow{m}_{\ast}$ (\hyperref[VI.lem8.7.2]{8.7.2}), there is
an isomorphism
$R^n\underleftarrow{m}_{\ast}(Q^{\ast}(j\mapsto M_j))\simeq
\varinjlim_{I^{\circ}}\mu^{-1}_jR^nm_{j\ast}(M_j)$, which gives
formula \eqhyperref[VI.eq8.7.3.1]{8.7.3.1} by what precedes. The first
assertion follows immediately from formula
\eqhyperref[VI.eq8.7.3.1]{8.7.3.1}.

\begin{enonce*}{Corollary 8.7.4}\etiquette[8.7.4]{VI.cor8.7.4}
For\pageoriginale $n\in\ZZ$, there are functorial isomorphisms
\begin{equation*}
R^n \underleftarrow{m}_{\ast} (N) \simeq \varinjlim_{I^{\circ}}
\mu^{\ast}_j R^n m_{j\ast} \mu_{j\ast} (N){\quoi}.\tag{8.7.4.1}
\end{equation*}\etiquette[8.7.4.1]{VI.eq8.7.4.1}
\end{enonce*}

Indeed, $N$ is isomorphic to
$Q^{\ast}(j\mapsto\mu_{j\ast}(N))$.

\begin{enonce*}{Corollary 8.7.5}\etiquette[8.7.5]{VI.cor8.7.5}
For every object $i$ of $I$, every integer $n$, and every
$A_i$-Module $M_i$, there is a functorial isomorphism
\begin{equation*}
R^n \underleftarrow{m}_{\ast} \mu^{\ast}_i (M_i) \simeq
\varinjlim_{f : j \to i} \mu^{\ast}_j R^n m_{j\ast} f^{\ast}
(M_i){\quoi}.\tag{8.7.5.1}
\end{equation*}\etiquette[8.7.5.1]{VI.eq8.7.5.1}
\end{enonce*}

After making the base change $I_{/i}\to I$, we may assume that $i$
is a final object of $I$ (\hyperref[VI.lem8.2.1]{8.2.1}). The object
$\mu_i^{\ast}(M_i)$ is then isomorphic to
$Q^{\ast}((f:j\to i)\mapsto f^{\ast}(M_i))$
(\cf \hyperref[VI.cor8.5.5]{8.5.5}).

\begin{enonce*}{Corollary 8.7.6}\etiquette[8.7.6]{VI.cor8.7.6}
Let $i$ be an object of $I$ and $n$ an integer. There is an
isomorphism, functorial in the section $j\mapsto M_j$:
\begin{equation*}
R^n \mu_{i\ast} Q^{\ast} (j \mapsto M_j) \simeq \varinjlim_{f :
  j \to i} R^n f_{\ast} (M_j){\quoi}.\tag{8.7.6.1}
\end{equation*}\etiquette[8.7.6.1]{VI.eq8.7.6.1}
There is an isomorphism functorial in the $\underrightarrow{A}$-Module
$N$:
\begin{equation*}
R^n \mu_{i\ast} (N) \simeq \varinjlim_{f:j\to i}
R^n f_{\ast} \mu_{j\ast} (N){\quoi}.\tag{8.7.6.2}
\end{equation*}\etiquette[8.7.6.2]{VI.eq8.7.6.2}
There is an isomorphism functorial in the $A_i$-Module $M_i$:
\begin{equation*}
R^n \mu_{i\ast} \mu^{\ast}_i (M_i) \simeq \varinjlim_{f:j\to i}
R^n f_{\ast} f^{\ast} (M_i){\quoi}.\tag{8.7.6.3}
\end{equation*}\etiquette[8.7.6.3]{VI.eq8.7.6.3}
The functors $R^n\mu_{i\ast}$ commute with filtered inductive limits.
\end{enonce*}

Making the base change $I_{/i}\to I$ reduces us to the case where
$i$ is a final object of $I$. Formulas
\eqhyperref[VI.eq8.7.6.1]{8.7.6.1} and
\eqhyperref[VI.eq8.7.6.2]{8.7.6.2} are then special cases of formulas
\eqhyperref[VI.eq8.7.3.1]{8.7.3.1},
\eqhyperref[VI.eq8.7.4.1]{8.7.4.1}, and
\eqhyperref[VI.eq8.7.5.1]{8.7.5.1}, respectively, obtained by taking
for $G$ the constant fibered topos with fiber $F_i$ and for $m$ the
morphism $(f:F_j\to F_i, f\in\ob(I_{/i}))$. The last assertion
follows from \eqhyperref[VI.eq8.7.6.2]{8.7.6.2}, taking
\hyperref[VI.cor8.5.4]{8.5.4} into account.

\begin{enonce*}{Corollary 8.7.7}\etiquette[8.7.7]{VI.cor8.7.7}
Let\pageoriginale $i$ be an object of $I$, and let $X$ be an object
of $F_i$ such that the functors $H^n(X,-)$, $n\in\NN$, commute with
filtered inductive limits. For every $n$, there is a canonical
isomorphism, functorial in the $A_i$-Module $M_i$:
\begin{equation*}
H^n (\mu_{i}^{\ast}(X), \mu^{\ast}_i (M_i)) \simeq \varinjlim_{f:j\to i}
H^n (f^{\ast} (X), f^{\ast} (M_i))\tag{8.7.7.1}
\end{equation*}\etiquette[8.7.7.1]{VI.eq8.7.7.1}
and the functors $H^n(\mu_i^{\ast}(X),-)$ commute with filtered
inductive limits. In particular, if the functors $H^n(F_i,-)$ commute
with filtered inductive limits, then the functors
$H^n(\underleftarrow{F},-)$ commute with filtered inductive limits,
and there are canonical isomorphisms, functorial in the $A_i$-Modules
$M_i$:
\begin{equation*}
H^n (\underleftarrow{F}, \mu_{i}^{\ast} (M_i)) \simeq \varinjlim_{f:j\to i}
H^n (F_j, f^{\ast} (M_i)){\quoi}.\tag{8.7.7.2}
\end{equation*}\etiquette[8.7.7.2]{VI.eq8.7.7.2}
\end{enonce*}

\noindent
It follows from \hyperrefext[V][V.prop5.3]{5.3} that there are two
spectral sequences
\begin{align*}
'E_{2}^{p, q} &= \varinjlim_{f:j\to i} H^p(X, R^q f_{\ast}
f^{\ast} (M_i)) \Rightarrow \varinjlim_{f:j\to i}
H^{p+q}(f^{\ast} (X), f^{\ast} (M_i))\\
''E_{2}^{p, q} &=  H^p(X, R^q \mu_{i\ast} \mu^{\ast}_i (M_i))
\Rightarrow H^{p+q} (\mu^{\ast}_i (X), \mu^{\ast}_i (M_i)){\quoi},
\end{align*}
and from \eqhyperref[VI.eq8.5.7.1]{8.5.7.1} that there is a morphism
between these two spectral sequences. Since the $H^{p+q}(X,-)$
commute with filtered inductive limits, it follows from
\eqhyperref[VI.eq8.7.6.2]{8.7.6.2} that this morphism of spectral
sequences is an isomorphism at the $E_2^{p,q}$ level. It therefore
induces an isomorphism on the abutments, giving
\eqhyperref[VI.eq8.7.7.1]{8.7.7.1}. For every
$\underrightarrow{A}$-Module $N$, there is a spectral sequence
(\hyperrefext[V][V.prop5.3]{5.3})
$$
E^{p, q}_2 = H^{p} (X, R^q \mu_{i\ast} (N)) \Rightarrow H^{p+q}
(\mu^{\ast}_i (X), N){\quoi}.
$$
The assertion that the functors $H^n(\mu_i^{\ast}(X),-)$ commute with
filtered inductive limits then follows from the analogous properties
of the functors $H^n(X,-)$ and $R^n\mu_{i\ast}$
(\hyperref[VI.cor8.7.6]{8.7.6}).

\begin{enonce*}[remark]{Remark 8.7.8}\etiquette[8.7.8]{VI.rem8.7.8}
When the topoi $F_j$, $j\in\ob(I)$, are coherent and the transition
morphisms are coherent, the hypothesis made on $X$ (\resp $F_i$) in
\hyperref[VI.cor8.7.7]{8.7.7} holds when $X$ (\resp $F_i$) is
coherent (\hyperref[VI.cor5.2]{5.2}). Moreover, we know that in this
case $\mu_i^{\ast}(X)$ (\resp $\underleftarrow{F}$) is coherent
(\hyperref[VI.thm8.3.13]{8.3.13}).
\end{enonce*}

\begin{enonce*}{Corollary 8.7.9}\etiquette[8.7.9]{VI.cor8.7.9}
Let\pageoriginale $i$ be an object of $I$, let $M_i$ be a left
$A_i$-Module, and let $L_i$ be a right $A_i$-Module admitting a
resolution of the form
$$
P_{i, k} \to P_{i, k-1} \to \cdots \to P_{i, 0} \to L_i \to 0
$$
where, for every integer $k$, $P_{i,k}$ is isomorphic to a finite
direct sum of objects of the form $A_{i|X}$ \emph{(IV \hyperref[IV.prop11.3.3]{11.3.3})}, with
$X$ satisfying the hypotheses of \hyperref[VI.cor8.7.7]{8.7.7}. If,
in addition to the hypotheses of \hyperref[VI.subsubsec8.7.1]{8.7.1},
the fibered topos $F$ is right flat (\hyperref[VI.subsec8.6]{8.6}),
then for $n\leq k-1$ there are canonical isomorphisms
\begin{equation*}
\Ext^n_{\underrightarrow{A}}(\underleftarrow{F}, \mu^{\ast}_i(L_i),
\mu^{\ast}_i (M_i)) \simeq \varinjlim_{f : j \to i}
\Ext^{n}_{A_j} (F_j, f^{\ast}(L_i), f^{\ast} (M_i)){\quoi}.\tag{8.7.9.1}
\end{equation*}\etiquette[8.7.9.1]{VI.eq8.7.9.1}
\end{enonce*}

\noindent
Denote by $P_{i,.}$ the resolution of $L_i$. Since $F$ is right flat,
for every $f:j\to i$, $f^{\ast}(P_{i,.})$ is a resolution of
$f^{\ast}(L_i)$, and $\mu_i^{\ast}(P_{i,.})$ is a resolution of
$\mu_i^{\ast}(L_i)$ (\hyperref[VI.subsec8.6]{8.6}). These resolutions
give two spectral sequences converging respectively to the two sides
of \eqhyperref[VI.eq8.7.9.1]{8.7.9.1}, together with a morphism between
the two spectral sequences. At the $E_1^{p,q}$ level this morphism is
an isomorphism by \eqhyperref[VI.eq8.7.7.1]{8.7.7.1}. It therefore
induces an isomorphism on the abutments.


% END INCLUDED SGA4_English_current_buildable_source_20260906/source/06/drafts/SGA4_VI_LINES_6916_7181_English.texfrag


% BEGIN INCLUDED SGA4_English_current_buildable_source_20260906/source/06/drafts/SGA4_VI_LINES_7182_7367_English.texfrag
\section{Appendix. Criterion for the existence of points}\etiquette[9]{VI.sec9}

\centerline{\textit{by} P. \textsc{Deligne}}

\medskip
\begin{enonce*}{Proposition 9.0}\etiquette[9.0]{VI.prop9.0}
Every locally coherent topos $S$ has enough points.
\end{enonce*}

Since the question is local on $S$, we may suppose that the topos $S$
is defined by a site $\Ss$ in which finite projective limits are
representable and such that every covering $f_i:U_i\to U$ admits a
finite subcovering. It is enough to prove that if
$f:\Ff\to\Gg$ is not a monomorphism, then there is a point $x$ of
$\Ss$ such that $f_x$ is not injective. By hypothesis, there are
$U\in\ob\Ss$ and $\ss,\ss'\in\Ff(U)$ such that $\ss\ne\ss'$ and
$f(\ss)=f(\ss')$. Replacing $\Ss$ by $\Ss/U$, we reduce to the
following lemma.

\begin{enonce*}{Lemma 9.1}\etiquette[9.1]{VI.lem9.1}
If\pageoriginale $\ss$ and $\ss'$ are two distinct global sections of
a sheaf $\mathfrak F$ on a site $\Ss$ satisfying the preceding
hypotheses, then there is a point $x$ of $\Ss$ such that
$\ss_x\ne\ss'_x$.
\end{enonce*}

Let $P=(U_i)_{i\in I}$ be a projective system in $\Ss$, indexed by a
filtered ordered set $I$. For every sheaf $\Hh$ on $\Ss$, put
$P(\Hh)=\varinjlim\Hh(U_i)$, and for every $V\in\Ob\Ss$, put
$P(V)=\varinjlim\Hom(U_i,V)$. The functors $P(\Hh)$ and $P(V)$
commute with fiber products. If the functor $P(V)$ takes coverings to
surjective families of maps, it is a morphism of sites from the point
topos to $\Ss$, and $P(\Hh)$ is the corresponding fiber functor.

If $P=(U_i)_{i\in I}$ and $Q=(V_j)_{j\in J}$ are two projective
systems in $\Ss$, we say that $Q$ refines $P$ if $I$ is a subset of
$J$ with the induced order and $P$ is the restriction of $Q$ to $I$.
There are then morphisms of functors from $P(\Hh)$ to $Q(\Hh)$ and
from $P(V)$ to $Q(V)$.

For every $P$ as above, denote by $\ss_P$ and $\ss'_P$ the images of
$\ss$ and $\ss'$ in $P(\Ff)$. Lemma \hyperref[VI.lem9.1]{9.1}
follows from the next lemma, in which $P$ is taken to be the
projective system indexed by $\{0\}$ and consisting of the final
object of $\Ss$.

\begin{enonce*}{Lemma 9.2}\etiquette[9.2]{VI.lem9.2}
For every $P$ as above satisfying $\ss_P\ne\ss'_P$, there is a
refinement $Q$ of $P$ satisfying $\ss_Q\ne\ss'_Q$ and such that, for
every covering $f_i:V_i\to V$ and every $f\in Q(V)$, $f$ belongs to
the image of one of the $Q(V_i)$.
\end{enonce*}

We first prove:

\begin{enonce*}{Lemma 9.3}\etiquette[9.3]{VI.lem9.3}
Let $P=(U_i)_{i\in I}$ satisfy the hypotheses of Lemma
\hyperref[VI.lem9.2]{9.2}, let $f_i:V_i\to V$ be a finite covering in
$\Ss$, and let $f\in P(V)$. There is a refinement $Q$ of $P$
satisfying $\ss_Q\ne\ss'_Q$ and such that the image of $f$ in $Q(V)$
belongs to the image of one of the $Q(V_i)$.
\end{enonce*}

There are $i_{\circ}\in I$ and
$f'\in\Hom(U_{i_{\circ}},V)$ which define $f$. For
$i\geq i_{\circ}$, put
$U_{i,k}=\fprod{U_i}{V_k}{V}$, this fiber product being defined by
$f'$. Since $(V_k)$ is a covering of $V$, $(U_{i,k})$ is a covering
of $U_i$; the following arrow is injective (for $i>i_{\circ}$):
$$
\Ff(U_i)\to\prod_k\Ff(U_{i,k}).
$$
Passing\pageoriginale to the limit, and using the commutation of
\emph{finite} products with filtered inductive limits, we see that
$$
P(\Ff)\to\prod_k\varinjlim\Ff(U_{i,k})
$$
is injective. Hence there is a $k$ such that $\ss_P$ and $\ss'_P$
have distinct images in
$\varinjlim_{i>i_{\circ}}\Ff(U_{i,k})$. Put
$I_{\circ}=\{i\mid i\geq i_{\circ}\text{ and }i\in I\}$ and
$J=I\coprod I_{\circ}$. Order $J$ by declaring $j'\leq j''$ if the
images of $j'$ and $j''$ in $I$ satisfy $j'\leq j''$ and it is not
the case that $j'\in I_{\circ}$ and $j''\in I$. The $U_i$ and
$U_{i,k}$, indexed by $J$, form a refinement $Q$ of $P$ such that
$\ss_Q\ne\ss'_Q$ and the image of $f$ in $Q(V)$ belongs to the image
of $Q(V_k)$.

\begin{enonce*}{Lemma 9.4}\etiquette[9.4]{VI.lem9.4}
Let $P=(U_i)_{i\in I}$ satisfy the hypotheses of Lemma
\hyperref[VI.lem9.2]{9.2}. There is a refinement $Q$ of $P$ such that
$\ss_Q\ne\ss'_Q$ and such that, for every finite covering $(V_k)$ of
an object $V$ of $\Ss$ and every $f\in P(V)$, the image of $f$ in
$Q(V)$ belongs to the image of one of the $Q(V_k)$.
\end{enonce*}

Let $E$ be the set of triples consisting of an object $V$ of $\Ss$, a
finite covering $(V_k)$ of $V$, and an element $f\in P(V)$. Let
$\leq$ be a well-ordering of $E$, and let $\bar E$ be the set obtained
from $E$ by adjoining a greatest element, denoted by $\infty$. By
transfinite induction on $e\in\bar E$, we shall define refinements
$Q_e$ of $P$ satisfying $\ss_{Q_e}\ne\ss'_{Q_e}$ and such that
\begin{enumerate}
\renewcommand{\theenumi}{\roman{enumi}}
\renewcommand{\labelenumi}{(\theenumi)}
\item if $e'\leq e$, then $Q_e$ refines $Q_{e'}$;

\item if $e=(V,(V_k),f)$, then the image of $f$ in $Q_e(V)$ belongs
  to the image of one of the $Q_e(V_k)$.
\end{enumerate}

Suppose the $Q_{e'}$ already defined for $e'<e$. If $e$ is the first
element of $E$, put $Q'_e=P$. If $e$ has a predecessor $e-1$, put
$Q'_e=Q_{e-1}$. Otherwise, let $Q'_e$ be the projective system with
index set $\bigcup_{e'<e}I_{e'}$ which refines the $Q_{e'}$ for
$e'<e$. In this case,
$$
Q'_e(\Ff)=\varinjlim_{e'<e}Q_{e'}(\Ff),
$$
so that, in every case,\pageoriginale $Q'_e$ refines the $Q_{e'}$ for
$e'<e$ and satisfies $\ss_{Q'_e}\ne\ss'_{Q'_e}$.

The required projective system $Q_e$ is obtained by applying Lemma
\hyperref[VI.lem9.3]{9.3} to $Q'_e$ and
$e=(V,(V_k),f)$ (\resp by taking $Q_e=Q'_e$ if $e=\infty$).

The projective system $Q_{\infty}$ satisfies Lemma
\hyperref[VI.lem9.4]{9.4}.

Lemma \hyperref[VI.lem9.4]{9.4} makes it possible to define,
inductively on $n$, a sequence $Q_n=(U_i)_{i\in I_n}$ of projective
systems in $\Ss$ such that
\begin{enumerate}
\renewcommand{\theenumi}{\roman{enumi}}
\renewcommand{\labelenumi}{(\theenumi)}
\item $Q_0=P$;

\item $Q_{n+1}$ refines $Q_n$;

\item $\ss_{Q_n}\ne\ss'_{Q_n}$;

\item for every covering $f_k:V_k\to V$ and every $f\in Q_n(V)$,
  the image of $f$ in $Q_{n+1}(V)$ belongs to the image of one of the
  $Q_{n+1}(V_k)$.
\end{enumerate}

The projective system $Q$, whose index set is the union of the $I_n$
and which extends the various $Q_n$, then satisfies Lemma
\hyperref[VI.lem9.2]{9.2}. The proof also shows:

\begin{enonce*}{Corollary 9.5}\etiquette[9.5]{VI.cor9.5}
Let $\Ss$ be a site in which fiber products are representable and
such that every covering in $\Ss$ admits a finite subcovering. Put
$c=\sup(\aleph_0,\card(\text{Fl }\Ss))$. Then there is a very dense
set $X$ of points of $\Ss$ such that
\begin{enumerate}
\renewcommand{\theenumi}{\roman{enumi}}
\renewcommand{\labelenumi}{(\theenumi)}
\item $\card(X)\leq2^c$;

\item if $x\in X$ and $U\in\Ob\Ss$, then $\card(U_x)\leq c$.
\end{enumerate}
\end{enonce*}

\begin{thebibliography}{99}\pageoriginale
\bibitem{sga4.exp.VI.bib.1} P.\textsc{Gabriel} and M.\textsc{Zisman}: Homotopy Theory
  and Calculus of Fractions, Ergebnisse der Mathematik, Bd 35.

\bibitem{sga4.exp.VI.bib.2} J. \textsc{Giraud}: M\'ethode de la descente. M\'emoire de la S.M.F.

\bibitem{sga4.exp.VI.bib.3} A. \textsc{Grothendieck} and J. \textsc{Dieudonn\'e}: \'El\'ements de
  G\'eom\'etrie Alg\'ebrique I, 2nd edition.

\bibitem{sga4.exp.VI.bib.4} A. \textsc{Grothendieck} and J. \textsc{Dieudonn\'e}: \'El\'ements de
  G\'eom\'etrie Alg\'ebrique IV I, I.H.E.S.~${\rm n^{\circ}}$

\bibitem{sga4.exp.VI.bib.5} \SGA. 1 VI, by A. \textsc{Grothendieck}, I.H.E.S.

\bibitem{sga4.exp.VI.bib.6} \SGA. 3 IV 6.3 by M. \textsc{Demazure}, Lecture Notes
  ${\rm n^{\circ}}$, Springer Verlag.

\bibitem{sga4.exp.VI.bib.7} A. \textsc{Grothendieck}: Sur quelques points d'alg\`ebre
  homologique, Tohoku Math. J.
\end{thebibliography}

\ifx\danslelivre\undefined
\end{document}
\fi

% END INCLUDED SGA4_English_current_buildable_source_20260906/source/06/drafts/SGA4_VI_LINES_7182_7367_English.texfrag


% END INCLUDED SGA4_English_current_buildable_source_20260906/source/06/06_English.tex


\makeatletter
\def\input@path{{../07/}}
\makeatother
% BEGIN INCLUDED SGA4_English_current_buildable_source_20260906/source/07/07_English.tex
\ifx\danslelivre\undefined
% Release engine: XeLaTeX (the LuaTeX compatibility layer breaks Xy-pic output).
\documentclass[10pt]{../fichiers-configuration-latex/smfart}
% BEGIN INCLUDED SGA4_English_current_buildable_source_20260906/source/fichiers-configuration-latex/paquets.tex
\usepackage{polyglossia}
%\usepackage[francais]{babel}
\usepackage{fontspec}
\usepackage[math-style=ISO,bold-style=upright,sans-style=upright]{unicode-math}
%\usepackage[mathletters]{ucs}
%\usepackage[utf8]{inputenc}
%\usepackage[T1]{fontenc}
\usepackage[unicode,hypertexnames=false,colorlinks=true,linkcolor=blue,citecolor=Sepia,urlcolor=magenta]{hyperref}
\setmainlanguage{french}
\setmainfont[Ligatures=TeX]{Linux Libertine O}
% variante possible : TeXGyrePagella ?
\setmathfont{XITS Math}

% maths
\usepackage{amsmath}
%\usepackage{amssymb}
%\usepackage{mathrsfs}
%\usepackage{stmaryrd}
%\usepackage{euscript}

% couleurs
\usepackage[usenames,dvipsnames]{xcolor}

% graphiques
\usepackage{graphicx}
\usepackage[all]{xy}
%\usepackage{tikz}
%\usetikzlibrary{matrix,arrows,calc}
%\tikzset{column sep=3em,row sep=2em,text height=1.5ex, text depth=0.25ex,>=angle 90}
%\tikzset{description/.style={fill=white,inner sep=1pt}}

% hyperliens, 'etiquettes
\usepackage{xr-hyper}
\usepackage{../fichiers-configuration-latex/smfthm}
%\ifx\cacheretiquettes\undefined
%\usepackage[inline]{showlabels} \else \fi % pour rendre visibles les 'etiquettes
%\usepackage{../fichiers-configuration-latex/makeidx} % pour faire un index
%\usepackage[pagebackref,unicode]{hyperref}  % pour voir o`u sont cit'ees les r'ef'erences
% hyperref is loaded once above with the cumulative-reader destination policy.
%\hypersetup{breaklinks=true} %permet le retour `a la ligne dans les liens trop longs. Incompatible avec pdflatex
%\renewcommand*{\backref}[1]{{\small [cit\'e page(s)~#1.]}}

\usepackage{xspace}
\usepackage{pifont} % pour expose XVII

% END INCLUDED SGA4_English_current_buildable_source_20260906/source/fichiers-configuration-latex/paquets.tex

% BEGIN INCLUDED SGA4_English_current_buildable_source_20260906/source/fichiers-configuration-latex/numerotation.tex
\renewcommand{\thefootnote}{\roman{footnote}} % pour mettre les notes en bas de page sous la forme (ii) et non (2), qui ressemble trop (?) a un exposant.

\newcounter{pageoriginal}
\newcommand{\pageoriginale}{\refstepcounter{pageoriginal}\marginpar{\footnotesize\textbf{\thepageoriginal}}\xspace}
\let\pageoriginaled\pageoriginale

\usepackage{ifthen}
\newboolean{original}
\def\sisi{\ifthenelse{\boolean{original}}}
\setboolean{original}{false}

\newcounter{cnote}
\def\nde#1{\ifthenelse{\boolean{original}}{}{\stepcounter{cnote}\def\thefootnote{\arabic{cnote}}\protect\footnote[\value{cnote}]{N.D.E.: #1}}}
\def\ndemark{\ifthenelse{\boolean{original}}{}{\stepcounter{cnote}\def\thefootnote{\arabic{cnote}}\protect\footnotemark[\value{cnote}]}}
\def\ndetext#1{\ifthenelse{\boolean{original}}{}{\def\thefootnote{\arabic{cnote}}\protect\footnotetext[\value{cnote}]{N.D.E.: #1}}}

% END INCLUDED SGA4_English_current_buildable_source_20260906/source/fichiers-configuration-latex/numerotation.tex

% BEGIN INCLUDED SGA4_English_current_buildable_source_20260906/source/fichiers-configuration-latex/macros.tex
\setmainlanguage{french}
\setotherlanguage{english}


%% Fixed, engine-independent release identifier.  The inherited dynamic Git
%% probe used Lua io.popen, exposed build-state-dependent output, and printed
%% a visible question mark outside a Git checkout.
\def\showgitstatus{English release r2, 2026-07-23}


% macros pour LaTeX
\renewcommand{\mathsf}{}

% Police différente pour U+2113 SCRIPT SMALL L
\setmathfont[range="2113]{latinmodern-math.otf}

% Give every inherited semantic label its own stable PDF destination instead
% of inheriting whichever automatic equation/page anchor happened to be live.
% The legacy label remains as a compatibility alias used by existing edges.
\makeatletter
\def\etiquette[#1]#2{%
  \Hy@raisedlink{\hyper@anchorstart{sga4.semantic.#2}\hyper@anchorend}%
  \def\@currentHref{sga4.semantic.#2}%
  \def\@currentlabel{#1}%
  \label{#2}%
}
\makeatother
\def\eqhyperref[#1]#2{(\hyperref[#1]{#2})}
\def\pagehyperref[#1]#2{(page \hyperref[#1]{#2})}
\def\hyperrefext[#1][#2]#3{\hyperref[#2]{#1 #3}}
\def\eqhyperrefext[#1][#2]#3{\hyperref[#2]{#1 #3}}     %cf. exposes 17
\def\neweqhyperrefext[#1][#2]#3{\hyperref[#2]{#1 #3}}  %cf. exposes 17

\def\shat#1{\protect\widehat{#1}}
\def\stilde#1{\protect\widetilde{#1}}
%%%%%%%%%%%%%%%%%%


\def\etcn{etc.\xspace}
\def\etco{etc...\xspace}
\def\etc{\sisi{\etco}{\etcn}}
\def\ldotsn{\sisi{...}{\etcn}}

\def\cf{{cf.}\kern.3em}
\def\Cf{{Cf.}\kern.3em}
\def\ie{{i.e.}\kern.3em}
\def\iev{{i.e.},\xspace}
\def\loccito{loc.\kern2pt cit.\xspace}
\def\loccitn{\textit{loc.\kern2pt cit.}\xspace}
\def\loccit{\sisi{\loccito}{\loccitn}}
\def\EGA{\textup{EGA}\kern.2em}
\def\SGA{\textup{SGA}\kern.3em}

\def\quoi{}

\def\quav{\quad,\quad}
\def\quaet{\quad \text{et}\quad}
\def\fibreproduct#1#2#3{#1 ×_{#3} #2}
\let\fprod\fibreproduct
\def\ccheck#1#2{\mbox{${\check{\check{\rm #1}}}^{\lower2pt\hbox{$\scriptstyle #2$}}$}}
\def\sheaf#1{\protect\mathcal{#1}}
\def\vieuxsheaf#1{{#1}} % cf. problème avec sheaf{a\Phi_i} par exemple...
\DeclareMathOperator{\Hom}{Hom}
\DeclareMathOperator{\can}{can}
\DeclareMathOperator{\Ind}{Ind}
\DeclareMathOperator{\ind}{ind}
%\DeclareMathOperator{\Hhomtop}{\mathcal{H}\!\mathit{omtop}}
\DeclareMathOperator{\Hhomtop}{\mathbf{Homtop}}
\DeclareMathOperator{\Point}{Point}
\DeclareMathOperator{\Pro}{Pro}
\DeclareMathOperator{\Top}{Top}
\DeclareMathOperator{\Mod}{Mod}
\DeclareMathOperator{\Cart}{Cart}
%\DeclareMathOperator{\Mmorsite}{\mathcal{M}\!\mathit{orsite}}
\DeclareMathOperator{\Mmorsite}{\mathbf{Morsite}}
\DeclareMathOperator{\loc}{loc}
\DeclareMathOperator{\Loc}{Loc}
\DeclareMathOperator{\Pic}{Pic}
\DeclareMathOperator{\Et}{Et}
\DeclareMathOperator{\card}{card}
\DeclareMathOperator{\Ens}{Ens}
\DeclareMathOperator{\h}{\underline{H}}
\DeclareMathOperator{\Fonct}{Fonct}
\DeclareMathOperator{\Ker}{Ker}
\DeclareMathOperator{\but}{but}
\DeclareMathOperator{\Ob}{Ob}
\DeclareMathOperator{\ob}{ob}
\DeclareMathOperator{\Lib}{Lib}
\DeclareMathOperator{\esj}{esj}
\DeclareMathOperator{\id}{id}
\DeclareMathOperator{\Id}{Id}
\DeclareMathOperator{\ad}{ad}
\DeclareMathOperator{\Iim}{Im}
\DeclareMathOperator{\pr}{pr}
\DeclareMathOperator{\Cat}{Cat}
\DeclareMathOperator{\car}{car}
\DeclareMathOperator{\SysInd}{SysInd}
\DeclareMathOperator{\Sup}{Sup}
\DeclareMathOperator{\Coker}{Coker}
\DeclareMathOperator{\Filt}{Filt}
\DeclareMathOperator{\acc}{acc}
\DeclareMathOperator{\sur}{sur}
\DeclareMathOperator{\cat}{cat}
\DeclareMathOperator{\Ord}{Ord}
\DeclareMathOperator{\Gr}{Gr}
\DeclareMathOperator{\Cov}{Cov}
\DeclareMathOperator{\Caf}{Caf}
\DeclareMathOperator{\coker}{coker}
\DeclareMathOperator{\coim}{coim}
\DeclareMathOperator{\im}{im}
\DeclareMathOperator{\Ouv}{Ouv}
\DeclareMathOperator{\diag}{diag}
\DeclareMathOperator{\Esp}{Esp}
\DeclareMathOperator{\TOP}{TOP}
\DeclareMathOperator{\Unif}{Unif}
\DeclareMathOperator{\Iind}{\mathbf{Ind}}
\DeclareMathOperator{\Hhomcart}{\mathbf{Homcart}}
\DeclareMathOperator{\Hhomtopess}{\mathbf{Homtopess}}
\DeclareMathOperator{\Ppoints}{\mathbf{Points}}
\DeclareMathOperator{\Ppoint}{\mathbf{Point}}
\DeclareMathOperator{\Ttors}{\mathbf{Tors}}
\DeclareMathOperator{\Ssex}{\mathbf{Sex}}
\DeclareMathOperator{\Ffib}{\mathbf{Fib}}
\DeclareMathOperator{\Vvois}{\mathbf{Vois}}
\DeclareMathOperator{\Ppointess}{\mathbf{Pointess}}
\DeclareMathOperator{\scrHom}{\mathbf{Hom}}
\DeclareMathOperator{\sheafPic}{\mathbf{Pic}}
\DeclareMathOperator{\sheafAut}{\mathbf{Aut}}
\DeclareMathOperator{\Hhomtoplocan}{\mathbf{Homtoplocan}}
\DeclareMathOperator{\Hhomtopan}{\mathbf{Homtopan}}
\DeclareMathOperator{\Pptgeom}{\mathbf{Ptgéom}}
\DeclareMathOperator{\Morsite}{Morsite}
\DeclareMathOperator{\Fib}{Fib}
\DeclareMathOperator{\Fais}{Fais}
\DeclareMathOperator{\kar}{kar}
\DeclareMathOperator{\Groupes}{Groupes}
\DeclareMathOperator{\Restr}{Restr}
\DeclareMathOperator{\Isom}{Isom}
\DeclareMathOperator{\Prol}{Prol}
\DeclareMathOperator{\Res}{Res}
\DeclareMathOperator{\Tors}{Tors}
\DeclareMathOperator{\esp}{esp}
\DeclareMathOperator{\es}{es}
\DeclareMathOperator{\od}{od}
\DeclareMathOperator{\gr}{gr}
\DeclareMathOperator{\sob}{sob}
\DeclareMathOperator{\dfn}{dfn}
\DeclareMathOperator{\Sch}{Sch}
\DeclareMathOperator{\ouv}{ouv}
\DeclareMathOperator{\cons}{cons}
\DeclareMathOperator{\Nerf}{Nerf}
\DeclareMathOperator{\Points}{Points}
\DeclareMathOperator{\Espsob}{Espsob}
\DeclareMathOperator{\Homtop}{Homtop}
\DeclareMathOperator{\Pt}{Pt}
\DeclareMathOperator{\Quot}{Quot}
\DeclareMathOperator{\glob}{glob}
\DeclareMathOperator{\Homfib}{Homfib}
\DeclareMathOperator{\Qouv}{Qouv}
\DeclareMathOperator{\Quotouv}{Quotouv}
\DeclareMathOperator{\ens}{ens}
\DeclareMathOperator{\ab}{ab}
\DeclareMathOperator{\an}{an}
\DeclareMathOperator{\fs}{fs}
\DeclareMathOperator{\red}{red}
\DeclareMathOperator{\Ab}{Ab}
\DeclareMathOperator{\cosq}{cosq}
\DeclareMathOperator{\Ann}{Ann}
\DeclareMathOperator{\Zar}{Zar}


%%%%%%%%%%%%
%%%%% SGA4-2
\DeclareMathOperator{\Sex}{Sex}
\DeclareMathOperator{\vois}{vois}
\DeclareMathOperator{\Vois}{Vois}
\DeclareMathOperator{\Ext}{Ext}
\DeclareMathOperator{\ext}{ext}
\DeclareMathOperator{\Fer}{Fer}
\DeclareMathOperator{\Add}{Add}
\DeclareMathOperator{\Cocart}{Cocart}
\DeclareMathOperator{\cosk}{cosk}
\DeclareMathOperator{\plat}{plat}

\DeclareMathOperator{\faisc}{faisc}
\DeclareMathOperator{\sk}{sk}
\DeclareMathOperator{\ch}{ch}

\def\HH{\mathbf{H}}
\def\qc{\mathrm{qc}}
\def\et{\mathrm{et}}
\def\coh{\mathrm{coh}}

\def\Hvr#1{%
\begin{array}[b]{l}
{\kern2pt\scriptstyle r}\\[-8pt]
{\kern1pt\scriptstyle \vee}\\[-5pt]
H^{#1}
\end{array}}

\def\Hvv#1{%
\begin{array}[b]{l}
\kern2pt\rotatebox{90}{${\scriptscriptstyle\ll}$}\\[-6pt]
H^{#1}
\end{array}}

\def\Hvi#1{%
\begin{array}[b]{l}
{\kern0.2pt\scriptscriptstyle \infty}\\[-8pt]
{\kern1pt\scriptstyle \vee}\\[-5pt]
H^{#1}
\end{array}}

\def\Hv#1{%
\begin{array}[b]{l}
{\kern1pt\scriptstyle \vee}\\[-5pt]
H^{#1}
\end{array}}

\DeclareMathOperator{\cart}{cart}
\DeclareMathOperator{\Ffaisc}{\mathbf{Faisc}}


\def\Ev#1{%
\begin{array}[b]{l}
{\kern1pt\scriptstyle \vee}\\[-5pt]
E'_{#1}
\end{array}}

\DeclareMathOperator{\Cont}{Cont}
\DeclareMathOperator{\Ccont}{\mathscr{C}\!\mathit{ont}}
\DeclareMathOperator{\Site}{Site}
\newcommand\MM{\mathbf{M}}
\def\limtop#1{{\displaystyle{\varprojlim_{#1}\!\mbox{\rm top }}}}
\def\Limtop{\mathop{\varprojlim\!\mbox{\rm top}}}
\DeclareMathOperator{\spec}{spec}
\def\ss{\scriptstyle{\mathscr{S}}}
\DeclareMathOperator{\Gal}{Gal}
\def\smallm{\scriptstyle{\mathscr{M}}}
\DeclareMathOperator{\End}{End}
\DeclareMathOperator{\ass}{ass}
\DeclareMathOperator{\zar}{zar}
\DeclareMathOperator{\Ogg}{Ogg}
\DeclareMathOperator{\res}{res}
\DeclareMathOperator{\tr}{tr}
\DeclareMathOperator{\Supp}{Supp}
\DeclareMathOperator{\codim}{codim}
\DeclareMathOperator{\rk}{rk}
\DeclareMathOperator{\cur}{cur}
\DeclareMathOperator{\cdqc}{cdqc}
\newcommand\DD{\mathbf{D}}
\DeclareMathOperator{\cl}{cl}
\def\JJ{\mathbf{J}}
\def\FF{\mathbf{F}}
\def\BB{\mathbf{B}}
\DeclareMathOperator{\Alget}{Alget}
\DeclareMathOperator{\Ass}{Ass}
\DeclareMathOperator{\supp}{supp}
\DeclareMathOperator{\rad}{rad}
\DeclareMathOperator{\prof}{prof}
\DeclareMathOperator{\cint}{int}
\DeclareMathOperator{\sic}{sic}
\DeclareMathOperator{\Tor}{Tor}
\DeclareMathOperator{\tors}{tors}
\DeclareMathOperator{\Sym}{Sym}
\DeclareMathOperator{\parf}{parf}
\DeclareMathOperator{\tor}{tor}
\DeclareMathOperator{\defn}{defn}
\DeclareMathOperator{\Div}{Div}
\DeclareMathOperator{\Pol}{Pol}
\DeclareMathOperator{\Fex}{Fex}
\DeclareMathOperator{\Lie}{Lie}
\DeclareMathOperator{\tot}{tot}
\DeclareMathOperator{\Aut}{Aut}
\DeclareMathOperator{\HOM}{HOM}
\DeclareMathOperator{\pch}{pch}
\DeclareMathOperator{\Ch}{Ch}
\DeclareMathOperator{\EXT}{EXT}
\DeclareMathOperator{\PIC}{PIC}
\DeclareMathOperator{\Rf}{Rf}
\DeclareMathOperator{\inj}{inj}
\newcommand\VV{\mathbf{V}}

\def\Xv{%
\begin{array}[b]{l}
{\kern1pt\scriptstyle \vee}\\[-5pt]
X
\end{array}}

%%%%%%%
\newcommand\TT{\mathbf{T}}
\newcommand\XX{\mathbf{X}}
\newcommand\EE{\mathbf{E}}
%%%%%%

\def\qedsymbol{\textup{C.Q.F.D.}}
\let\CQFD\qed
\let\cqfd\qed

%%%%%\bigdot
\font\bigsymb=cmsy10 at 4pt
\def\bigdot{{\kern1.2pt\raise 1.5pt\hbox{\bigsymb\char15}}}
%%%%%

\def\Lright{\mathop{\vcenter{\hbox{\hbox{$\longrightarrow$}\lower8pt \hbox{\kern-16pt\hbox{$\longrightarrow$}}}}}}

\def\Aa{{\mathcal A}}
\def\Bb{{\mathcal B}}
\def\Cc{{\mathcal C}}
\def\Dd{{\mathcal D}}
\def\Ee{{\mathcal E}}
\def\Ff{{\mathcal F}}
\def\Gg{{\mathcal G}}
\def\Hh{{\mathcal H}}
\def\Ii{{\mathcal I}}
\def\Jj{{\mathcal J}}
\def\Kk{{\mathcal K}}
\def\Ll{{\mathcal L}}
\def\Mm{{\mathcal M}}
\def\Nn{{\mathcal N}}
\def\Oo{{\mathcal O}}
\def\Pp{{\mathcal P}}
\def\Qq{{\mathcal Q}}
\def\Rr{{\mathcal R}}
\def\Ss{{\mathcal S}}
\def\Tt{{\mathcal T}}
\def\Uu{{\mathcal U}}
\def\Vv{{\mathcal V}}
\def\Ww{{\mathcal W}}
\def\Xx{{\mathcal X}}
\def\Yy{{\mathcal Y}}
\def\Zz{{\mathcal Z}}

\newcommand{\CC}{\mathbf{C}}
\newcommand{\GG}{\mathbf{G}}
\newcommand{\KK}{\mathbf{K}}
\newcommand{\LL}{\mathbf{L}}
\newcommand{\NN}{\mathbf{N}}
\newcommand{\PP}{\mathbf{P}}
\newcommand\QQ{\mathbf{Q}}
\newcommand\RR{\mathbf{R}}
\newcommand\ZZ{\mathbf{Z}}

\def\to{\mathchoice{\longrightarrow}{\rightarrow}{\rightarrow}{\rightarrow}}
\def\from{\mathchoice{\longleftarrow}{\leftarrow}{\leftarrow}{\leftarrow}}
\def\mto{\mathchoice{\longmapsto}{\mapsto}{\mapsto}{\mapsto}}
\def\mfrom{\mathrel{\from\joinrel\mapstochar\,}}
\def\hto{\mathrel{\lhook\joinrel\to}}
\def\hfrom{\mathrel{\from\joinrel\rhook}}

\def\implies{\mathchoice{\Longrightarrow}{\Rightarrow}{\Rightarrow}{\Rightarrow}}
\def\implied{\mathchoice{\Longleftarrow}{\Leftarrow}{\Leftarrow}{\Leftarrow}}
\def\iff{\mathchoice{\Longleftrightarrow}{\Leftrightarrow}{\Leftrightarrow}{\Leftrightarrow}}

\def\To#1{\mathchoice{\xrightarrow{\labelstyle~~#1~}}{\stackrel{#1}{\longrightarrow}}{}{}}
\def\Mto#1{\mathrel{\mapstochar\To{#1}}}
\def\From#1{\mathchoice{\xleftarrow{\labelstyle~~#1~}}{\stackrel{#1}{\longleftarrow}}{}{}}

\newcommand{\isomto}{\overset{\sim}{\longrightarrow}}
\newcommand{\isomfrom}{\overset{\sim}{\longleftarrow}}

\newcommand{\IFF}{\unskip~$\iff$~\ignorespaces}
\newcommand{\THEN}{\unskip~$\implies$~\ignorespaces}

\let\leq\leqslant
\let\geq\geqslant
\let\le\leqslant
\let\ge\geqslant

\def\cf{{cf.}\kern.3em}
\def\Cf{{Cf.}\kern.3em}
\def\ie{{i.e.}\kern.3em}
\def\iev{{i.e.},\xspace}
%\def\loccito{loc.\kern2pt cit.\xspace}
%\def\loccitn{\textit{loc.\kern2pt cit.}\xspace}
%\def\loccit{\sisi{\loccito}{\loccitn}}
\def\EGA{\textup{EGA}\kern.2em}
\def\SGA{\textup{SGA}\kern.3em}

\let\oldS\S
\def\S{\oldS\kern .15em}
\def\ptbl{.\kern .2em}
\def\resp{resp.\kern.3em}
\def\Exp{Exp.\kern.3em}

\let\bar\overline
\let\hat\widehat
\let\tilde\widetilde
\let\dpl\displaystyle
\let\ldots\dots
\let\emptyset\varnothing
\let\skpt\skippointrait
\let\epsilon\varepsilon

\def\sfrac#1#2{{#1}/{#2}}

\DeclareMathOperator{\cd}{cd}
%\DeclareMathOperator{\Hhom}{\mathcal{H}\!\mathit{om}}
\DeclareMathOperator{\Hhom}{\mathbf{Hom}}
%\DeclareMathOperator{\Eext}{\mathcal{E}\!\mathit{xt}}
\DeclareMathOperator{\Eext}{\mathbf{Ext}}
\DeclareMathOperator{\Spec}{Spec}
\DeclareMathOperator{\Tr}{Tr}

\let\emptyset=\varnothing

%%%%%%%%%%%%%%%
%%% expose XVII
%%%%%%%%%%%%%%%

\makeatletter
\let\old@optexpose\optexpose
\newcommand{\empty@optexpose}[1]{}
\newcommand{\paren@optexpose}[1]{\old@optexpose{#1}(}
\newcommand{\refnoexpose}[1]{\let\optexpose\empty@optexpose\ref{#1}\let\optexpose\old@optexpose}%\refnoexpose{X.cor5.2} engendre "5.2" au lieu de "X 5.2"
\newcommand{\neweqref}[1]{\let\optexpose\paren@optexpose\textup{\ref{#1})}\let\optexpose\old@optexpose}%\neweqref{XII.eq4.1} engendre "XII (4.1)" au lieu de "(XII 4.1)"
\makeatother
\newcommand{\refvb}[2][VI${}_{\text{B}}$ ]{\sisi{#1\refnoexpose{#2}}{\ref{#2}}}

\newcommand{\leftexp}[2]{{\vphantom{#2}}^{#1}{#2}}
\newcommand{\AAA}{\mathbf{A}}
\newcommand{\Sf}{\mathfrak{S}}
\DeclareMathOperator*{\quotprojlim}{\text{``$\varprojlim$''}}
\DeclareMathOperator*{\quotinjlim}{\text{``$\varinjlim$''}}
\DeclareMathOperator{\Fl}{F\sisi{\ell}{l}}
\DeclareMathOperator{\Pl}{P\sisi{\ell}{l}}
%\DeclareMathOperator{\cl}{c\sisi{\ell}{l}}
\DeclareMathOperator{\tdf}{tdf}
\DeclareMathOperator{\rg}{rg}
\newcommand{\iInd}{\sisi{I}{i}nd}
\newcommand{\iind}{ind}
\newcommand{\ppro}{pro}
\newcommand{\tTor}{\sisi{T}{t}or}
\newcommand{\ttor}{tor}
\newcommand{\chap}{Chap.}
\newcommand{\Ffex}{\Ff\mathit{ex}}
\newcommand{\sheafk}{\underline{k}}
\newcommand{\sheafn}{\underline{n}}
\newcommand{\sheafGamma}{\underline{\Gamma}}

% Repeated references to an already typeset footnote must not use
% \footnotemark without a matching \footnotetext: Hyperref would then create
% a live annotation to a destination that cannot exist.  This helper keeps the
% current printed footnote mark and links it to an explicit stable note target.
\makeatletter
\newcommand{\sgarepeatedfootnotemark}[1]{%
  \begingroup
  \edef\@thefnmark{\thefootnote}%
  \hyperlink{#1}{\@makefnmark}%
  \endgroup}
\makeatother

\sisi{
\DeclareMathOperator{\simarrow}{\rule{.15mm}{1.8mm}\kern -3.4pt\rightsquigarrow}
\newcommand{\sheafell}{\underline{\ell}}
\newcommand{\sheafZZ}{\underline{\ZZ}}
}{
\newcommand{\nderef}[1]{la N.D.E.~(\ref{#1}) page~\pageref{#1}}
}

%%%%%%%%%%%%%%%%%
%%%%% expose XVII
%%%%%%%%%%%%%%%%%

% END INCLUDED SGA4_English_current_buildable_source_20260906/source/fichiers-configuration-latex/macros.tex


\externaldocument{../01/01_English.tex}
\externaldocument{../02/02_English.tex}
\externaldocument{../03/03_English.tex}
\externaldocument{../04/04_English.tex}
\externaldocument{../05/05_English.tex}
\externaldocument{../05bis/05bis_English.tex}
\externaldocument{../06/06_English.tex}

\setboolean{original}{false}

\begin{document}

\begin{center}
VII. The étale site and topos of a scheme\\
A. Grothendieck
\end{center}

\selectlanguage{english}version: \showgitstatus
\tableofcontents

\else
\chapter{The étale site and topos of a scheme}
%\addtocontents{toc}{A. Grothendieck}
\begin{center}
A. Grothendieck
\end{center}
\fi
\phantomsection\MakeLinkTarget*{sga4.exp.VII}\label{sga4.exp.VII}

% BEGIN INCLUDED SGA4_English_current_buildable_source_20260906/source/07/drafts/SGA4_VII_LINES_43_396_English.texfrag
\setcounter{pageoriginal}{341}

In\pageoriginale the present exposé and the next, we develop the most
elementary properties concerning étale topology and cohomology. The
developments of the present exposé concern certain properties that are, for
the most part, valid for other, very different topologies, such as the
``fppf topology''. In the next exposé, we shall develop properties rather
special to the étale topology, arising from the very particular nature of
étale morphisms.

Here we partly follow three oral exposés by J.E. \textsc{Roos} (which he had
not been able to write up), notably in the proof of
\hyperrefext[VIII][VIII.thm6.3]{6.3}.

Unless expressly stated otherwise, it will be understood that the schemes
considered in the present exposé and the following ones are elements of the
fixed universe $𝒰$.

\section{The étale topology}\etiquette[1]{VII.sec1}

\subsection{}\etiquette[1.1]{VII.subsec1.1}

We denote by $(\Sch)$ the category of schemes (elements of the fixed
universe $𝒰$). Recall that a morphism $f:X\to Y$ in $(\Sch)$ is
called\pageoriginale \emph{étale} if it is locally of finite presentation
and flat, and if for every $y\in Y$, the fiber $X_y$ is discrete and its
local rings are finite \emph{separable} extensions of $k(y)$. Equivalently,
$f$ is locally of finite presentation and, for every affine $Y$-scheme $Y'$
(in the absolute sense) and every subscheme $Y'_0$ defined by a nilpotent
ideal, the map
$$
\Hom_Y (Y', X) \to \Hom_Y (Y'_{0}, X)
$$
is \emph{bijective}. For the most important properties of this notion, I
refer to \EGA IV \S\S 17 and 18, and, in the meantime, to \SGA 1 I and IV.


\subsection{}\etiquette[1.2]{VII.subsec1.2}

The \emph{étale topology} on $(\Sch)$ is the topology generated by the
pretopology for which, for every $X\in(\Sch)$, the set $\Cov X$ consists of
the families $(X_i\xrightarrow{u_i}X)_{i\in I}$ (indexed by an $I\in U$)
such that the $u_i$ are étale and
$X=\bigcup_i u_i(X_i)$ (in the set-theoretic sense).

It is convenient to associate with every $X$ the subcategory $\Et_{/X}$ of
$(\Sch)_{/X}$ consisting of the arrows $X'\to X$ that are étale. We endow it
with the ``induced topology'' (III) from the étale topology of $(\Sch)$
(also called the ``étale topology on $X$''), and denote the resulting site
by $X_{\text{ét}}$ (the ``étale site'' of $X$)
\footnote{Note, however, that it would be preferable to denote by
  $X_{\text{ét}}$ the topos $\widetilde{X_{\text{ét}}}$ defined by the
  étale site of $X$. For practical reasons, in this Seminar we shall retain
  the notation introduced here (which is that of the original seminar).}.
Note that every morphism of $\Et_{/X}$ is étale (\SGA 1 I 4.8), whence a
family $(X'_i\xrightarrow{u_i}X')_{i\in I}$ in $X_{\text{ét}}$ is covering
if and only if it is \emph{surjective}, \ie
$\bigcup_i u_i(X'_i)=X'$. One sees immediately (\cf
\hyperref[VII.prop3.1]{3.1}) that $X_{\text{ét}}$ is a $\Uu$-site, so the
results of Exposés I through VI apply. The topos
$\widetilde{X_{\text{ét}}}$ of $𝒰$-sheaves on $X_{\text{ét}}$ is the
\emph{étale topos} of $X$.

\subsection{}\etiquette[1.3]{VII.subsec1.3}

In what\pageoriginale follows, unless expressly stated otherwise, all the
topological notions considered in $\Et_{/X}$ or $(\Sch)$ are understood
with respect to the étale topology. Moreover, in language and notation we
shall commonly write $X$ in place of $X_{\text{ét}}$ or $Et_{/X}$. Thus a
``sheaf on $X$'' (understood: for the étale topology) means a sheaf on the
site $X_{\text{ét}}$. We denote by $\widetilde{X_{\text{ét}}}$ the category
of these sheaves, which is a $\Uu$-topos. If $F$ is an abelian sheaf on $X$,
we denote its cohomology groups by $H^i(X,F)$; these would be denoted
$H^i(X_{\text{ét}},F)$ in Exposé V.

\subsection{}\etiquette[1.4]{VII.subsec1.4}

Let $f:X\to Y$ be a morphism of schemes. The inverse-image functor
$$
f^{\ast} : Y' \mapsto \fprod{Y'}{X}{Y}
$$
induces a functor
\begin{equation*}
f^{\ast}_{\text{ét}} : \Et_{/Y} \to \Et_{/X}\tag{$\ast$}
\end{equation*}
that commutes with finite $\varprojlim$ and takes covering families to
covering families; it is therefore a morphism of sites
$$
f^{\ast}_{\text{ét}} : X_{\text{ét}} \to Y_{\text{ét}}{\quoi},
$$
and thus induces a functor on the categories of sheaves:
$$
f^{\text{ét}}_{\ast} : \widetilde{X_{\text{ét}}} \to \widetilde{Y_{\text{ét}}}
$$
given by
$$
f^{\text{ét}}_{\ast}  (F) = F \circ f^{\ast}_{\text{ét}}{\quoi}.
$$
Moreover, $f^{\text{ét}}_{\ast}$ admits a left adjoint
$$
f^{\ast}_{\text{ét}} : \widetilde{Y_{\text{ét}}} \to
\widetilde{X_{\text{ét}}}
$$
extending\pageoriginale the functor considered in $(\ast)$ and commuting
with arbitrary $\varinjlim$ and finite $\varprojlim$, \ie
$f^{\ast}_{\text{ét}}$ defines a \emph{morphism of topoi}
$$
f_{\text{ét}} : \widetilde{X_{\text{ét}}} \to \widetilde{Y_{\text{ét}}}{\quoi}.
$$

Clearly, for a composite of two morphisms of schemes
$$
X \xrightarrow{f} Y \xrightarrow{g} Z
$$
there are canonical isomorphisms
\begin{equation*}
\begin{cases}
(gf)^{\text{ét}}_{\ast} \simeq g_{\ast}^{\text{ét}}
  f_{\ast}^{\text{ét}} &\quav \mbox{ whence}\\
(gf)^{\ast}_{\text{ét}} \simeq f^{\ast}_{\text{ét}}
  g^{\ast}_{\text{ét}} &\quav \mbox{ \ie there is an
    isomorphism}\\
(gf)_{\text{ét}} \simeq g_{\text{ét}} f_{\text{ét}}&\quav
\end{cases}
\end{equation*}
(so that one obtains a ``pseudofunctor'' (\SGA 1 VI 8) from the category
$(\Sch)$ to the category $(\Top)$ of topoi $\in𝒰'$, where $𝒰'$ is the
smallest universe such that $𝒰\in𝒰'$; compare IV).

\subsection{Notation}\etiquette[1.5]{VII.subsec1.5}

In what follows, we shall often write $f^{\ast}$ and $f_{\ast}$ in place of
$f^{\ast}_{\text{ét}}$ and $f^{\text{ét}}_{\ast}$. If $f:X\to Y$ is a
morphism of schemes, the functors $R^if_{\ast}^{\text{ét}}$ will therefore
simply be denoted $R^if_{\ast}$. In this notation, recall the Leray spectral
sequence for $f$ and an abelian sheaf $F$ on $X$
(\hyperrefext[V][V.sec5]{5}):
$$
H^{\ast} (X, F) \Longleftarrow E^{p. q}_2 = H^{p}(Y, R^q f_\ast (F)){\quoi}.
$$
Likewise, if there are two morphisms $f:X\to Y$ and $g:Y\to Z$, there is a
Leray spectral sequence (the spectral sequence of composite functors)
$$
R^{\ast} (gf)_{\ast} (F) \Longleftarrow E^{p. q}_2 = R^{p} g_{\ast} (R^q
f_\ast (F)){\quoi}.
$$

\subsection{}\etiquette[1.6]{VII.subsec1.6}

When\pageoriginale $f:X\to Y$ is itself an étale morphism, $X$ may be
regarded as an object of $Y_{\text{ét}}$, and there is a canonical
\emph{isomorphism of sites}
$$
X_{\text{ét}} \xrightarrow{\sim} (Y_{\text{ét}})_{/X}{\quoi}.
$$
Under this isomorphism, the functor
$f^{\ast}_{\text{ét}}:Y'\mapsto\fprod{Y'}{X}{Y}$ of $(\ast)$ identifies with
the internal base-change functor in the category $\Et_{/Y}$. It follows that
the functor
$$
f^{\ast} : \widetilde{Y_{\text{ét}}} \to \widetilde{X_{\text{ét}}}
$$
is then isomorphic to the functor ``\emph{restriction} to $Y$''
$$
f^{\ast} (F) \simeq F \circ i_{X/Y}{\quoi},
$$
where $F$ is a sheaf on $Y_{\text{ét}}$ and
$i_{X/Y}=X_{\text{ét}}\to Y_{\text{ét}}$ is the evident functor (which
regards an étale scheme over $X$, $u:X'\to X$, as an étale scheme over $Y$
via $fu:X'\to Y$).

\subsection{Questions of universes}\etiquette[1.7]{VII.subsec1.7}

Note that if $X\neq∅$, then $X_{\text{ét}}$ \emph{is not} an element of the
chosen universe $𝒰$. Nevertheless, as we have noted, it is easy to see
(\hyperref[VII.prop3.1]{3.1}) that one can find a full subcategory $C$ of
$X_{\text{ét}}$, \emph{an element of} $𝒰$, satisfying the conditions of
the ``comparison lemma'' (III), so that the restriction functor induces an
equivalence $\widetilde{C}\to\widetilde{X_{\text{ét}}}$. Thus the
\emph{étale topos} is equivalent to a topos defined in terms of a site
$C\in𝒰$. In other words, $X_{\text{ét}}$ is a $U$-site, and hence by
\loccit the results of Exposés I through VI apply to this site.

\section{Examples of sheaves}\etiquette[2]{VII.sec2}\pageoriginale

\begin{enumerate}
\renewcommand{\theenumi}{\alph{enumi}}
\renewcommand{\labelenumi}{\theenumi)}
\item Let $F\in\Ob(\Sch)_{/X}$ be a scheme over $X$, and for every $X'$ over
  $X$ put
$$
F(X') = \Hom_X (X', F){\quoi}.
$$
The functor $(\Sch)^{\circ}_{/X}\to(\Ens)$ thus defined is a sheaf for the
étale topology (and even for the finer fpqc topology studied in \SGA 3 IV).
In other words, the étale topology on $(\Sch)$ is coarser than the canonical
topology. This is, essentially, the content of \SGA 1 VIII 5.1 (\cf also
\SGA 3 IV 6.3.1). A fortiori, the restriction of this sheaf to
$X_{\text{ét}}$ is a sheaf. It will again be denoted by $F$ when there is no
risk of confusion\footnote{\etiquette[6]{page6}Care must be taken, however,
  that if $X\neq∅$, the functor $\varphi$ that associates with
  $F\in\Ob(\Sch)_{/X}$ the corresponding sheaf $\varphi(F)$ is \emph{not}
  fully faithful, nor even faithful, and that $F$ cannot be reconstructed
  (up to isomorphism) from $\varphi(F)$. For example, if $S=\Spec k$, with
  $k$ an algebraically closed field, knowledge of $\varphi(F)$ is equivalent
  merely to knowledge of the underlying set of $F$ (by
  \hyperrefext[VIII][VIII.cor2.4]{2.4}). Compare IV}. Note also that the
functor thus obtained
\begin{equation*}
(\Sch)_{/X} \to \widetilde{X_{\text{ét}}}\tag{$\ast$}
\end{equation*}
commutes with finite $\varprojlim$ (this is immediate). This implies, for
example, that when $F$ is a group scheme (\resp\ldots) over $X$, the sheaf
it defines is a sheaf of groups (\resp\ldots). Note that the functor
$\Et_{/X}\to\widetilde{X_{\text{ét}}}$ induced by (*) is simply the
canonical functor, associating with every
$X'\in\Ob X_{\text{ét}}$ the functor on $X_{\text{ét}}$ that it
represents. It is therefore an isomorphism of the category $\Et_{/X}$ onto a
full subcategory of the étale topos $\widetilde{X_{\text{ét}}}$, by which we
generally identify an $X\in\Ob X_{\text{ét}}$ with the corresponding sheaf,
denoted $\widetilde{X}$ or simply $X$.

Clearly, $H^{\circ}(X,F)=\Hom_{(\Sch)_{/X}}(X,F)$.

Let us also give an interpretation of $H^1(X,F)$ when $F$ is a group
prescheme over $X$ (commutative, if desired, so that the definition of
$H^1(X,F)$ falls within Exposé V;\pageoriginale for the general case one may
consult J. Giraud's thesis\footnote{\etiquette[7]{page7}}). Familiar
arguments, which I shall not repeat here, then show that $H^1(X,F)$ is
canonically isomorphic to the group of classes (up to isomorphism) of
\emph{sheaves of sets on $X_{\text{ét}}$ that are principal homogeneous
under $F$\footnote{or $F$-torsors in the terminology of \loccit.}}. When
$F$ is \emph{affine} over $X$, \SGA 1 VIII 2.1 implies (again by standard
arguments, \cf J. \textsc{Giraud}'s thesis) that $H^1(X,F)$ is also the
group of classes of \emph{schemes} $P$ over $X$ on which $F$ acts on the
right and which are ``principal homogeneous bundles over $X$ for the étale
topology'', \ie locally trivial for the étale topology\footnote{or
  representable $F$-\emph{torsors} in the now accepted terminology}.
\end{enumerate}

\begin{enonce*}[remark]{Remarks 2.1}\etiquette[2.1]{VII.rem2.1}
Care must be taken that if, for every scheme $Z$ over $X$, the associated
sheaf on $X_{\text{ét}}$ is denoted by $\varphi_X(Z)$, and if $f:Y\to X$ is
a morphism of schemes, there is an evident homomorphism (functorial in $Z$)
\begin{equation*}
f^{\ast} (\varphi_X (Z)) \to \varphi_Y (Z_Y),\quad Z_Y =
\fprod{Z}{Y}{X}{\quoi},\tag{$\ast$}
\end{equation*}
\ie an evident homomorphism
$$
\varphi_X (Z) \to f_{\ast} (\varphi_Y (Z_Y)){\quoi},
$$
namely the homomorphism, functorial in $X'\in\Ob X_{\text{ét}}$,
$$
\Hom_X (X', Z) \to \Hom_Y (X'_{Y}, Z_{Y}){\quoi}.
$$
But note that \emph{in general $(\ast)$ is not an isomorphism, \ie the
formation of the étale sheaf associated with a relative scheme does not
commute with inverse-image functors}. Nevertheless, $(\ast)$ is an
isomorphism in the special case in which $Z$ is étale over $X$. Likewise,
the functor\pageoriginale
$$
\varphi_{X} : \Sch_{/X} \to \widetilde{X_{\text{ét}}}
$$
\emph{is not faithful} if $X\neq∅$\footnote{\Cf the note at the bottom of
  page \pagehyperref[page6]{6}}; nevertheless, its restriction to
$X_{\text{ét}}$, which is simply the canonical functor
$X_{\text{ét}}\to\widetilde{X_{\text{ét}}}$, is \emph{fully faithful},
because by a) the topology of $X_{\text{ét}}$ is coarser than its canonical
topology.

\begin{enumerate}
\renewcommand{\theenumi}{\alph{enumi}}
\renewcommand{\labelenumi}{\theenumi)}
\setcounter{enumi}{1}
\item Note that $(\ast)$ also commutes with direct sums indexed by the
  $I\in𝒰$, as is readily verified. In particular, if for every set $I$ we
  denote by $I_X$ the constant $X$-scheme defined by $I$ (\SGA 3 I 1.8), the
  direct sum of $I$ copies of $X$, then the associated sheaf is simply the
  constant sheaf $I_{X_{\text{ét}}}$ (the direct sum of $I$ copies of the
  final sheaf on $X_{\text{ét}}$). Since the functor $I\mapsto I_X$ also
  commutes with finite $\varprojlim$, we see that it takes a group to a
  group, a commutative group to a commutative group, \etc If $G$ is an
  ordinary commutative group, we shall simply write $H^i(X,G)$ in place of
  $H^i(X,G_X)$. Suppose, for example, that $G$ is a finite group. Then $G_X$
  is finite, hence affine over $X$, and thus, using the final remark of a),
  we obtain an interpretation of $H^1(X,G)$ as the set \emph{of classes of
  principal Galois coverings with group $G$} (\SGA 1 V 2.7). When $X$ is
  connected and endowed with a geometric point $a$, the ``fundamental
  group'' $\pi_1(X,a)$ (\SGA 1 V 7) gives the canonical isomorphism
\begin{equation*}
H^1(X, G) \simeq \Hom(\pi_1(X, a), G){\quoi},\tag{${\ast}{\ast}$}
\end{equation*}
(where $G$ has been assumed commutative).

One may say that in passing from Zariski cohomology to the étale topology,
``we have done what was needed'' to obtain the ``right'' $H^1$ (which occurs
in the second\pageoriginale member of $({\ast}{\ast})$) for a finite constant
coefficient group $G$. \emph{It is a remarkable fact, to be proved in the
remainder of this seminar, that this also suffices to obtain the ``right''
$H^i(X,G)$ for every torsion coefficient group (at least if $G$ is prime to
the residual characteristics of $X$)}.

\item Let $F$ be an $\Oo_X$-module on $X$, for the Zariski topology. Then
  (with the notation of \SGA 3 I 4.6) define a functor
$$
W(F) : (\Sch)_{/X}^\circ \to (\Ens)
$$
by
$$
W(F)(X') = \Gamma (X', F \otimes_{\Oo_{X}} \Oo_{X'}){\quoi}.
$$
This functor is again a sheaf for the étale topology (and even for the fpqc
topology), as follows again from \SGA 1 VIII 1.6 and \SGA 3 IV 6.3.1. A
fortiori, the restriction of $W(F)$ to $X_{\text{ét}}$ is again a sheaf,
which we also denote by $W(F)$. By definition, therefore,
$H^{\circ}(X,W(F))=\Gamma(F)=H^{\circ}(X,F)$. But there is a stronger result
if $F$ is quasi-coherent; \cf \hyperref[VII.prop4.3]{4.3}.
\end{enumerate}
\end{enonce*}

% END INCLUDED SGA4_English_current_buildable_source_20260906/source/07/drafts/SGA4_VII_LINES_43_396_English.texfrag


% BEGIN INCLUDED SGA4_English_current_buildable_source_20260906/source/07/drafts/SGA4_VII_LINES_397_665_English.texfrag
\section{Generators of the étale topos. Cohomology of an
  $\protect\varinjlim$ of sheaves}\etiquette[3]{VII.sec3}

\begin{enonce*}{Proposition 3.1}\etiquette[3.1]{VII.prop3.1}
Let $X$ be a scheme, and let $C$ be a full subcategory of
$X_{\text{\rm ét}}$ such that, for every affine scheme $X'$ étale
over $X$, $C$ contains an object isomorphic to $X'$. Then $C$ is a
``family of topological generators'' of the site $X_{\text{\rm ét}}$,
and hence a family of generators of the topos
$\widetilde{X_{\text{\rm ét}}}$. The restriction functor
$\widetilde{X_{\text{\rm ét}}}\to\widetilde{C}$ is an equivalence of
categories (where $C$ is equipped with the topology induced by that
of $X_{\text{\rm ét}}$).
\end{enonce*}

Immediate\pageoriginale from the ``comparison lemma.''

\begin{enonce*}{Corollary 3.2}\etiquette[3.2]{VII.cor3.2}
Suppose that $X$ is quasi-separated. Call the full subcategory $C$ of
$X_{\text{\rm ét}}$ formed by the schemes étale over $X$ that are of
finite presentation over $X$, equipped with the topology induced by
$X_{\text{\rm ét}}$, the restricted étale site of $X$. Then:
\begin{enumerate}
\renewcommand{\theenumi}{\roman{enumi}}
\renewcommand{\labelenumi}{(\theenumi)}
\item $C$ is stable under fiber products and is a site of finite type
  if $X$ is quasi-compact.

\item The restriction functor
  $\widetilde{X_{\text{\rm ét}}}\to\widetilde{C}$ is an equivalence of
  categories.
\end{enumerate}
\end{enonce*}

Assertion (i) is immediate, and (ii) follows from
\hyperref[VII.prop3.1]{3.1}, since $C$ satisfies the condition of
\hyperref[VII.prop3.1]{3.1} by the quasi-separatedness of $X$. Indeed,
this implies that, if $X'$ is quasi-compact, for example affine, then
it is quasi-compact \emph{over} $X$ (\EGA IV 1.2.4, with
$Z=\Spec\ZZ$), and hence is of finite presentation over $X$ if $X'$
is locally of finite type over $X$.

\begin{enonce*}{Proposition 3.3}\etiquette[3.3]{VII.prop3.3}
Suppose that $X$ is quasi-compact and quasi-separated. Then the
functors $H^q(X_{\text{\rm ét}},F)$ on the category of abelian sheaves
on $X_{\text{\rm ét}}$ commute with $\varinjlim$.
\end{enonce*}

Indeed, by \hyperref[VII.cor3.2]{3.2} (ii), one may replace
$X_{\text{ét}}$ by $C$. Since $X$ is quasi-compact, $X\in\Ob C$, and
the conclusion follows from \hyperref[VII.cor3.2]{3.2} (i) and
\hyperrefext[VI][VI.subsubsec6.1.2]{6.1.2} (3).

\section{Comparison with other topologies}\etiquette[4]{VII.sec4}

\setcounter{subsection}{-1}
\subsection{}\etiquette[4.0]{VII.subsec4.0}

First note that the examples of sheaves on $X_{\text{ét}}$ considered
in ${\rm n^{\circ}}$~2 are in fact naturally \emph{restrictions} of
sheaves\pageoriginale defined on $(\Sch)_{/X}$ equipped with its étale
topology (or even with the fpqc topology). More generally, let
$$
u^{\ast} : X_{\text{ét}} \to (\Sch)_{/X}
$$
be the inclusion functor. It is continuous (III) and commutes with
finite $\varprojlim$, and hence defines a \emph{morphism of sites}
$$
u : (\Sch)_{/X} \to X_{\text{ét}}{\quoi},
$$
whence a functor
$$
u_{\ast} : \widetilde{(\Sch)_{/X}} \to \widetilde{X_{\text{ét}}}
$$
on the associated categories of sheaves:
$$
u_{\ast}(F)=F\circ u{\quoi},
$$
and a left adjoint\footnote{Strictly speaking, since $(\Sch)_{/X}$ is
a $\Uu$-category but \emph{not} a $\Uu$-site, one cannot invoke III
for the existence of $u^{\ast}$. Nevertheless, one constructs
$u^{\ast}$ easily by
$$
u^{\ast} (F) (X') = (p_{X'})^{\ast}_{\text{ét}} (F) (X')\quad (X'\in \Ob \Sch_{/X}){\quoi},
$$
where $p_{X'}$ denotes the structural morphism $X'\to X$. Moreover,
to give meaning to \hyperref[VII.prop4.1]{4.1} (ii) and justify the
indicated proof of \hyperref[VII.prop4.1]{4.1}, one must introduce a
universe $\Vv$ such that $\Uu\in\Vv$, and regard the sheaves occurring
in the formulas of \hyperref[VII.prop4.1]{4.1} (i) as sheaves with
values in $\Vv\text{-}(\Ens)$. Alternatively, if one wishes to work
only with $\Uu$-sites, one may replace $(\Sch)_{/X}$ by a full
subcategory of $(\Sch)_{/X}$ stable under finite $\varprojlim$,
containing the $C$ of \hyperref[VII.prop3.1]{3.1}, and
$\Uu$-small.} to the latter:
$$
u^{\ast} : \widetilde{X_{\text{ét}}} \to \widetilde{\Sch_{/X}}.
$$
With this understood:

\begin{enonce*}{Proposition 4.1}\etiquette[4.1]{VII.prop4.1}
\begin{enumerate}
\renewcommand{\theenumi}{\roman{enumi}}
\renewcommand{\labelenumi}{(\theenumi)}
\item The functor $u^{\ast}$ is fully faithful; hence, for every sheaf
  $G$ on $X_{\text{\rm ét}}$, the canonical homomorphism
  $G\to u_{\ast}u^{\ast}(G)$ is an isomorphism.

\item Let $G$ (\resp $F$) be an abelian sheaf on
  $X_{\text{\rm ét}}$ (\resp $(\Sch)_{/X}$). Then the following
  canonical homomorphisms (defined, for example, as edge
  homomorphisms of the corresponding Leray spectral sequences) are
  isomorphisms:\pageoriginale
\begin{gather*}
H^{\ast} (X_{\text{\rm ét}}, F \circ u) \xrightarrow{\sim} H^{\ast}((\Sch)_{/X}, F)\\
H^{\ast} (X_{\text{\rm ét}}, G) \xrightarrow{\sim} H^{\ast}((\Sch)_{/X}, u^{\ast}G){\quoi}.
\end{gather*}
\end{enumerate}

\end{enonce*}

Indeed, the functor $u$ is a comorphism, and one can apply (III).

Thus (i) shows that, for the study of the cohomology of sheaves, it is
essentially equivalent to work with the ``small'' étale site
$X_{\text{ét}}$ or with the ``big'' étale site $(\Sch)_{/X}$.

\setcounter{subsection}{1}
\subsection{}\etiquette[4.2]{VII.subsec4.2}

One should also introduce on $(\Sch)$ (and hence on the
$(\Sch)_{/X}$) various topologies other than the étale topology, the
most useful of which are defined in \SGA 3 IV 6.3. The coarsest among
them is the \emph{Zariski topology} $(\Zar)$, defined by the
pretopology whose covering families are the surjective families of
open immersions; it is coarser than the étale topology. The finest of
these topologies is the \emph{``faithfully flat quasi-compact''
topology}, abbreviated (fpqc), which is the coarsest topology for
which the covering families in the Zariski sense, as well as
faithfully flat quasi-compact morphisms, are covering; the fpqc
topology is finer than the étale topology. As already observed, the
various examples of sheaves on $(\Sch)_{/X}$ considered in
${\rm n^{\circ}}$~2 are in fact already sheaves for the fpqc topology.

\subsubsection{}\etiquette[4.2.1]{VII.subsubsec4.2.1}

One must bear in mind, however, that for an abelian sheaf $F$ on
$(\Sch)_{/X}$ for the fpqc topology (or for a topology such as fppf,
considered in \SGA 3), the cohomology groups of $F$ for the fpqc
topology (\resp\ldots) are not always isomorphic to the cohomology
groups for\pageoriginale the étale topology. This is so even if $X$ is the spectrum
of a field $k$ and $F$ is representable by an algebraic group over
$k$, even for $H^1$. In general, one can show that the étale topology
gives the ``correct'' cohomology groups for coefficient groups that
are \emph{étale}, or more generally \emph{smooth}, group schemes over
$X$\footnote{\Cf A. Grothendieck, Le groupe de Brauer III, Th.~11.7,
in ``Dix exposés sur la topologie des schémas,'' North-Holland
Publishing Co.}, but this is no longer true for group schemes such as
\emph{radicial} groups over $S$, for which one must replace the étale
topology, still too coarse, by the fpqc or fppf topology.

\subsubsection{}\etiquette[4.2.2]{VII.subsubsec4.2.2}

As an example of the relations between cohomologies for different
topologies, consider the Zariski and étale topologies. Denote by
$X_{\Zar}$ the site of Zariski open subsets of $X$, so that there is a
canonical inclusion functor
$$
u^{\ast} : X_{\Zar} \to X_{\text{ét}}
$$
which defines a morphism of sites
$$
f : X_{\text{ét}} \to X_{\Zar}{\quoi},
$$
whence corresponding direct image functors
$$
f_{\ast} : \widetilde{X_{\text{ét}}} \to \widetilde{X_{\Zar}}\quad
(f_{\ast} (F) = F \circ u){\quoi},
$$
and inverse image functors
$$
f^{\ast} : \widetilde{X_{\Zar}} \to \widetilde{X_{\text{ét}}}{\quoi},
$$
adjoint to one another. Geometrically, the pair $(f_{\ast},f^{\ast})$
should be interpreted as a morphism of topoi
$$
f : \widetilde{X_{\text{ét}}} \to \widetilde{X_{\Zar}}{\quoi}.
$$
One obtains a homomorphism of cohomological functors
$$
H^{\ast} (X_{\Zar}, f_{\ast} (F)) \to H^{\ast} (X_{\text{ét}}, F)
$$
and\pageoriginale a Leray spectral sequence
$$
H^{\ast} (X_{\text{ét}}, F) \Longleftarrow H^{p} (X_{\Zar}, R^q
f_{\ast} (F)){\quoi},
$$
where $F$ is an abelian sheaf on $X_{\text{ét}}$. This spectral
sequence summarizes the general relations between étale and Zariski
cohomology. Of course, for coefficient sheaves $F$ such as constant
coefficient sheaves, this spectral sequence is in general far from
trivial, \ie in general $R^qf_{\ast}(F)\neq0$ (notably, take the case
where $X$ is the spectrum of a field). Nevertheless:

\begin{enonce*}{Proposition 4.3}\etiquette[4.3]{VII.prop4.3}
Let $F$ be a sheaf of varying modules on the scheme $X$ (a sheaf for
the Zariski topology), whence a sheaf $F_{\text{\rm ét}}$ on
$X_{\text{\rm ét}}$ (\cf 2, c) and a homomorphism of cohomological
functors
$$
H^{\ast} (X, F) \to H^{\ast} (X_{\text{\rm ét}}, F_{\text{\rm ét}}){\quoi}.
$$
When $F$ is quasi-coherent, the preceding homomorphism is an
isomorphism.
\end{enonce*}

Indeed, the first member is precisely
$H^{\ast}(X_{\Zar},f_{\ast}(F_{\text{ét}}))$, and with the preceding
notation it is enough to prove that
$$
R^q f_{\ast} (F_{\text{ét}}) = 0\quad \mbox{for}\quad q > 0{\quoi}.
$$
By the method for computing the $R^qf_{\ast}$ and the fact that the
affine open subsets form a basis for the topology of $X$, this follows
immediately from:

\begin{enonce*}{Corollary 4.4}\etiquette[4.4]{VII.cor4.4}
If $X$ is affine, then
$$
H^{q} (X_{\text{\rm ét}}, F_{\text{\rm ét}}) = 0\quad \mbox{for}\quad q > 0{\quoi}.
$$
\end{enonce*}

To see this, let $C=X_{\text{ét}}$, and let $C'$ be the full
subcategory formed by the \emph{affine} $X'$. Then, by
\hyperref[VII.prop3.1]{3.1},
$$
H^{q} (C, F_{\text{ét}}) \simeq H^{q} (C', F | C'){\quoi}.
$$

It is therefore\pageoriginale enough to prove that $F|C'$ is
$C'$-acyclic (\hyperrefext[V][V.def4.1]{4.1}), or equivalently that it
satisfies the condition of \hyperrefext[V][V.prop4.3]{4.3}, \ie that,
for every $X'\in\Ob C'$ and every covering family
$ℛ=(X'_i\to X')_i$ in $C'$, one has $H^q(ℛ,F)=0$ for $q>0$. One may
suppose the family $(X'_i)_i$ finite and then, after replacing the
$X'_i$ by their coproduct, that the covering family consists of a
single covering morphism $X'_i\to X'$, \ie an (étale and) surjective
morphism. We are thus reduced to proving that, if $f:X\to Y$ is a
surjective étale morphism of affine schemes and $F$ a quasi-coherent
module on $Y$, then
\begin{equation*}
H^q (X /Y, F) = 0\quad \mbox{for}\quad q > 0{\quoi}.\tag{$\ast$}
\end{equation*}
This is proved in TDTE I, B, \hyperref[VII.subsec1.1]{1.1} (in FGA,
\cf reference [4] of IX).

\begin{enonce*}[remark]{Remarks 4.5}\etiquette[4.5]{VII.rem4.5}
In fact, by \loccit, for $(\ast)$ it is enough that $X\to Y$ be
surjective and flat (rather than étale). The preceding proof then
shows that one still has isomorphisms analogous to that of
\hyperref[VII.prop4.3]{4.3}, with the étale site $X_{\text{ét}}$
replaced by $(\Sch)_{/X}$ equipped with any of the topologies finer
than the Zariski topology considered in \SGA 3 IV 6.3, for example
the fpqc topology.
\end{enonce*}

% END INCLUDED SGA4_English_current_buildable_source_20260906/source/07/drafts/SGA4_VII_LINES_397_665_English.texfrag


% BEGIN INCLUDED SGA4_English_current_buildable_source_20260906/source/07/drafts/SGA4_VII_LINES_666_1032_English.texfrag
\section{Cohomology of a projective limit of schemes}\etiquette[5]{VII.sec5}

\subsection{}\etiquette[5.1]{VII.subsec5.1}

Let $I$ be an increasing filtered preordered set,
$$
\mathfrak{X} = (X_i)_{i \in I}
$$
a projective system of schemes, the transition morphisms $u_{ji} :
X_i \to X_j$ $(i \geq j)$ being \emph{affine}. Recall (\EGA IV
8\ldots) that under these conditions,\pageoriginale the projective limit
$$
X  = \varprojlim X_i
$$
exists in the category of schemes (and even in the category of locally
ringed spaces) and may be constructed as follows: one chooses
$i_{0} \in I$, so that for $i \geq i_0$ one has an
$X_{i_0}$-isomorphism (\EGA II 1.2 and 1.3)
$$
X_i = \Spec 𝒜_{i}{\quoi},
$$
where $𝒜_i$ is a quasi-coherent Algebra on $X_{i_0}$, the
$𝒜_i$ $(i \geq i_0)$ thus forming an inductive system of
such Algebras. Setting
$$
𝒜 = \varinjlim_{i} 𝒜_i{\quoi},
$$
one may take
$$
X = \Spec 𝒜{\quoi}.
$$
Moreover, denoting by esp $S$ the underlying topological space of a
scheme $S$, one proves (\EGA IV 8\ldots) that the canonical map
$$
\esp X \to \varprojlim \esp (X_i)
$$
is a \emph{homeomorphism} and that the canonical homomorphism of
sheaves of rings
$$
\varinjlim_{i} \gamma^{-1}_{i} (\Oo_{X_i}) \to \Oo_X
$$
is bijective, where $\gamma_i : X \to X_i$ is the canonical continuous
map.

\subsection{}\etiquette[5.2]{VII.subsec5.2}

Informally, the content of \EGA IV 8, 9 may be summarized by saying
that, \emph{when $X_{i_0}$ is quasi-compact and quasi-separated},
then ``every datum of a schematic nature on $X$, of finite presentation
over $X$'', is equivalent\pageoriginale to a datum of the same nature
on one of the $X_i$, ``for sufficiently large $i$''. Thus one proves
(\EGA IV 8 $\ldots$):

\begin{enumerate}
\renewcommand{\theenumi}{\alph{enumi}}
\renewcommand{\labelenumi}{(\theenumi)}
\item if $i_1 \in I$ and $X'_{i_1}$, $X''_{i_1}$ are two
schemes of finite presentation over $X_{i_1}$, then, setting for every
$i \geq i_1$
$$
X'_{i} = \fprod{X'_{i_1}}{X_i}{X_{i_1}},\quad X' = \fprod{X'_{i_1}}{X}{X_{i_1}}{\quoi},
$$
and defining $X''_i$, $X''$ in the same way, the canonical map
$$
\varinjlim_{i \geq i_1} \Hom_{X_i} (X'_i, X''_i) \to
\Hom_X (X', X'')
$$
is bijective.

\item For every scheme $X'$ of finite presentation over $X$, there
exist an index $i_1 \in I$, a scheme $X'_{i_1}$ of finite presentation
over $X_{i_1}$, and an $X$-isomorphism
$$
X'\simeq \fprod{X'_{i_1}}{X}{X_{i_1}}{\quoi}.
$$
\end{enumerate}

\subsection{}\etiquette[5.3]{VII.subsec5.3}

The preceding result, and analogous results contained in EGA IV 8,
may be expressed in the language of the $\varinjlim$ of fibered
categories introduced in VI. For this purpose, consider the category
$\Ff$ of morphisms $f : T \to S$ of finite presentation of schemes,
and the ``target functor''
$$
\pi : \Ff \to (\Sch){\quoi},
$$
which is plainly a \emph{fibered functor}, its fiber categories being,
moreover, equivalent to $\Uu$-small categories. Consider the inverse
image of the fibered category $\Ff/(\Sch)$ by the functor
$$
i \mapsto X_i : \cat (I) \to (\Sch)
$$
(where\pageoriginale $\cat(I)$ is the category associated with $I$, with
$\Hom(i, j) \neq ∅$ if and only if $i \geq j$), giving a fibered
category
$$
\pi_{\mathfrak{X}} : \Ff_{\mathfrak{X}} \to \cat(I){\quoi},
$$
whose fiber at each $i \in I$ is canonically isomorphic to the full
subcategory $\Ff_{X_i}$ of $(\Sch)_{/X_i}$ consisting of schemes of
finite presentation over $X_i$, the base-change morphism
$\Ff_i \to \Ff_{j}$ relative to $j \geq i$ being none other than
$X'_i \mapsto X'_j = \fprod{X'_i}{X_j}{X_i}$. This being so, by
associating to every $X'_i$ of finite presentation over $X_i$ its
inverse image on $X$
$$
\varphi(X'_i) = \fprod{X'_i}{X}{X_i}{\quoi},
$$
one obtains a natural functor
$$
\varphi : \Ff_{\mathfrak{X}} \to \Ff_{X}{\quoi},
$$
which plainly takes every cartesian morphism of
$\Ff_{\mathfrak{X}}$ to an isomorphism, and hence, by definition,
factors uniquely through a canonical functor
\begin{equation*}
\psi : \varinjlim_{\Ff_X /\Cat I}
\Ff_{\mathfrak{X}} \to \Ff_{X}{\quoi}.\tag{$\ast$}
\end{equation*}
In view of the description of the first member, the result of \EGA IV
8 recalled above may then be stated by saying that $(\ast)$ is an
\emph{equivalence of categories}.

\subsection{}\etiquette[5.4]{VII.subsec5.4}

Statement a) of \hyperref[VII.subsec5.2]{5.2}. above may be refined in various ways,
by introducing some class $(\Mm)$ of morphisms of schemes, stable under
base change, and stating that for a given $X_{i_1}$-morphism\pageoriginale
$u_{i_1} : X'_{i_1} \to X''_{i_1}$, the corresponding morphism
$u': X'\to X''$ is $\Ee (\Mm)$ if and only if there exists
$i\geq i_{1}$ such that the $X_{i}$-morphism
$u_{i}:X'_{i}\to X''_{i}$ is $\Ee (\Mm)$. The statement obtained in
this way is true, for example, when $(\Mm)$ is one of the following
classes of arrows: proper morphisms (respectively projective, \resp
quasi-projective, \resp affine, \resp quasi-affine, \resp finite, \resp
quasi-finite, \resp radicial, \resp surjective, \resp flat, \resp
smooth, \resp étale, \resp unramified). The reader will find the
corresponding statements in \EGA IV, §8 (for the first nine), §11 (for
the flat case), and §17 (for the last three cases).

\subsection{}\etiquette[5.5]{VII.subsec5.5}

Now replace $\Ff \to (\Sch)$ by the fibered subcategory $\Gg$
consisting of \emph{étale morphisms} of finite presentation, whose
fiber category $\Gg_{S}$, for every scheme $S$, is none other than the
category of schemes étale and of finite presentation over $S$, and
similarly form the fibered category $\Gg_{\mathfrak{X}}$ over
$\cat(I)$, whose fiber category at every $i \in I$ is the category
$\Gg_{X_i}$. By \hyperref[VII.cor3.2]{3.2} (ii), the étale topoi
$\widetilde{X_{i_{\text{ét}}}}$, $\widetilde{X_{\text{ét}}}$ are also
canonically equivalent to $\widetilde{\Gg_{X_{i}}}$,
$\widetilde{\Gg_{X}}$, (where each $\Gg_S$ is endowed with the
topology induced by $S_{\text{ét}}$). Moreover, by
\hyperref[VII.cor3.2]{3.2} (i), each $\Gg_{X_i}$ is a site of finite
type, and is moreover equivalent to a $\Uu$-small site. In view of the
topologies on the $\Gg_{X_i}$, $\Gg$ becomes a \emph{fibered site}
over $\cat(I)$. This being so, we may state the following

\begin{enonce*}{Lemma 5.6}\etiquette[5.6]{VII.lem5.6}
The canonical functor
$$
\varinjlim_{\Gg_{\mathfrak{X}}/\cat(I)}
\Gg_{\mathfrak{X}} \to \Gg_X
$$
defines\pageoriginale an equivalence, respecting the topologies, between the
$\varinjlim$ site of the restricted étale sites
(\hyperref[VII.cor3.2]{3.2}) of the $X_i$ and the restricted étale site
of $X = \varprojlim_i X_i$.
\end{enonce*}

The fact that one obtains an equivalence of categories is one of the
statements recalled in \hyperref[VII.subsec5.4]{5.4} (the one in which
$(\Mm)$ is the class of étale morphisms in $(\Sch)$). The assertion
concerning the topologies is obtained in the same way by taking
$(\Mm)=$ the class of surjective morphisms in $(\Sch)$.

The preceding result allows us to apply the results of VI to the
present situation. In particular, we obtain

\begin{enonce*}{Theorem 5.7}\etiquette[5.7]{VII.thm5.7}
Let $\mathfrak{X} = (X_i)_{i \in I}$ be a projective system of
schemes, with $I$ an increasing filtered preordered set. Suppose that
the transition morphisms $u_{ij} : X_j \to X_i$ are affine (so that the
scheme $X = \varprojlim X_i$ is defined
(\hyperref[VII.subsec5.1]{5.1})), and that the $X_i$ are quasi-compact
and quasi-separated. Consider the fibered site $\Gg$ over $\cat(I)$
described in \hyperref[VII.subsec5.5]{5.5}, whose fiber $\Gg_i$ at
$i \in I$ is canonically isomorphic to the ``restricted étale site''
of $X_i$ (consisting of schemes étale and of finite presentation over
$X_i$). Let $F$ be an abelian sheaf on $\Gg$ (\ie $(3)$ a functor
$\Gg^{0} \to (\Ab)$ whose restriction $F_i$ to each $\Gg_i$ is a
sheaf, \ie a sheaf on $X_i$). Let $F_{\infty} =
\varinjlim_{i} u^{\ast}_i (F_i)$ be the induced sheaf on $X$, where
$u_i : X \to X_i$ is the canonical morphism. Then the canonical
homomorphisms
$$
\varinjlim_i H^n (X_i, F_i) \to H^n (X, F_{\infty})
$$
are isomorphisms.
\end{enonce*}

In\pageoriginale view of \hyperref[VII.lem5.6]{5.6}, this follows by
combining \hyperrefext[VI][VI.thm8.7.3]{8.7.3} and
\hyperrefext[VI][VI.prop8.5.2]{8.5.2}. We shall use
\hyperref[VII.thm5.7]{5.7} chiefly in the following particular case,
which is also made explicit in

\begin{enonce*}{Corollary 5.8}\etiquette[5.8]{VII.cor5.8}
With the preceding hypotheses and notation for $(X_i)_{i \in I}$,
$X$, $u_{ij}$, $u_i$, let $i_0 \in I$, and let $F_{i_0}$ be an
abelian sheaf on $X_{i_0}$. For every $i \geq i_0$, let
$F_i = u_{i_0 i} {}^{\ast}(F_{i_0})$, and let moreover
$F_{\infty} = u^{\ast}_{i_0} (F_{i_0})$. Under these conditions, the
canonical homomorphisms
$$
\varinjlim_i H^n (X_i, F_i) \to H^n (X, F_{\infty})
$$
are isomorphisms.
\end{enonce*}

Here is, however, a sometimes useful case that follows from
\hyperref[VII.thm5.7]{5.7} and not from \hyperref[VII.cor5.8]{5.8}:

\begin{enonce*}{Corollary 5.9}\etiquette[5.9]{VII.cor5.9}
With the hypotheses and notation of \hyperref[VII.thm5.7]{5.7} for
$(X_i)_{i \in I}$, $X$, $u_{ij}$, $u_i$, let $G_{i_0}$ be a
commutative group scheme locally of finite presentation over
$X_{i_0}$ (where $i_0 \in I$ is given). For every $i \geq i_0$, let
$G_i = \fprod{G_{i_0}}{X_i}{X_{i_0}}$, and similarly let
$G_{\infty} = \fprod{G_{i_0}}{X}{X_{i_0}}$. Then the canonical
homomorphisms
$$
\varinjlim_i H^n (X_i, G_i) \to H^n (X, G_{\infty})
$$
are isomorphisms.
\end{enonce*}

Indeed, we know that $G_{i_0}$ defines a sheaf on
$(\Sch)_{/X_{i_0}}$ (\cf 2 (a)), and, supposing that $i_0$ is an
initial object of $I$ (which we may do), hence obtaining a canonical
functor $\Gg \to (\Sch)_{/X_{i_0}}$, composition gives a contravariant
functor $F$ on $\Gg$, whose restriction to $\Gg_i$ is\pageoriginale
canonically isomorphic to the sheaf on $X_i$ defined by $G_i$. In
particular, $F$ is a sheaf on $\Gg$, and we may apply
\hyperref[VII.thm5.7]{5.7} to it. The hypothesis that $G_{i_0}$ is
locally of finite presentation over $X_{i_0}$ serves to ensure that the
sheaf $F_{\infty}$ of \hyperref[VII.thm5.7]{5.7} is isomorphic, by the natural homomorphism
$F_{\infty} \to \faisc(G_{\infty})$, to the sheaf
$\faisc(G_{\infty})$ defined by $G_{\infty}$ (as follows immediately
from \EGA IV 8.8.2 (i), by considering the values of the sheaves in
question on the objects of the restricted étale site of $X$
and applying \hyperref[VII.lem5.6]{5.6}). This completes the proof of
\hyperref[VII.cor5.9]{5.9}.

\begin{enonce*}{Corollary 5.10}\etiquette[5.10]{VII.cor5.10}
With the hypotheses and notation of \hyperref[VII.thm5.7]{5.7} for
$(X_i)_{i \in I}$, $X$, $u_{ij}$, $u_i$, let $F$ be a sheaf of
commutative groups on $X$; then the canonical homomorphisms
$$
\varinjlim_{i} H^n (X_i, u_{i\ast} (F)) \to H^n (X, F)
$$
are isomorphisms.
\end{enonce*}

The statements \hyperref[VII.thm5.7]{5.7} through
\hyperref[VII.cor5.10]{5.10} generalize to corresponding statements
for the $\Ext^i$ of sheaves of Modules, by
\hyperrefext[VI][VI.cor8.7.9]{8.7.9}; we shall not repeat them in the
present particular case. On the other hand, for later reference we
shall make explicit the variants of \hyperref[VII.cor5.8]{5.8} and \hyperref[VII.cor5.9]{5.9}
in terms of the functors $R^n f_{\ast}$.

\begin{enonce*}{Corollary 5.11}\etiquette[5.11]{VII.cor5.11}
Let $I$ be an increasing filtered preordered set, let $(X_i)_{i \in I}$
and $(Y_i)_{i \in I}$ be two projective systems of schemes with affine
transition morphisms, $X = \varprojlim_{i} X_i$,
$Y = \varprojlim_{i} Y_i$, $(f_i)_{i \in I}$ a projective system of
quasi-compact and quasi-separated morphisms $f_i : X_i \to Y_i$,
$i_0 \in I$, and $F_{i_0}$ a sheaf on $X_{i_0}$. For every
$i \geq i_0$, let $F_i$ be the sheaf\pageoriginale on $X_i$ that is
the inverse image of $F_{i_0}$. Similarly, let $f : X \to Y$ be
deduced from $(f_i)$ and let $F$ be the sheaf on $X$ that is the inverse
image of $F_{i_0}$. Under these conditions, the canonical homomorphism
$$
\varinjlim_{i} v^{\ast}_i (R^n f_{i\ast} (F_i)) \to R^n
f_{\ast} (F)
$$
is an isomorphism (where $v_i : Y \to Y_i$ is the canonical morphism).
\end{enonce*}

We may suppose that $i_0$ is an initial object of $I$, and, since the
question is local on $Y_{i_0}$, that $Y_{i_0}$ is affine, so that the
$Y_i$ and the $X_i$ are quasi-compact and quasi-separated. We are then
under the hypotheses for applying
\hyperrefext[VI][VI.thm8.7.3]{8.7.3} (where one reduces to
\hyperref[VII.cor5.8]{5.8}, in view of the computation of the
$R^{n} f_{\ast}$ as sheaves associated with presheaves). In the same
way one proves:

\begin{enonce*}{Corollary 5.12}\etiquette[5.12]{VII.cor5.12}
The same statement as \hyperref[VII.cor5.11]{5.11}, except that now
$F_{i_0}$ denotes a commutative group scheme locally of finite
presentation over $X_{i_0}$, and $F_i$, $F$ its inverse images (in the
sense of base change for schemes).
\end{enonce*}

Note that \hyperref[VII.cor5.9]{5.9} (\resp
\hyperref[VII.cor5.12]{5.12}) is not a particular case of
\hyperref[VII.cor5.8]{5.8} (\resp \hyperref[VII.cor5.11]{5.11}),
\cf \hyperref[VII.sec2]{2} a) remark \hyperref[VII.rem2.1]{2.1}.

\begin{enonce*}{Corollary 5.13}\etiquette[5.13]{VII.cor5.13}
With the notation of \hyperref[VII.cor5.11]{5.11}, let $F$ be a sheaf
of abelian groups on $X$. Then the canonical homomorphisms
$$
\varinjlim_i v_i^{\ast} (R^n f_{i\ast} (u_{i\ast} (F))) \to
R^n f_{\ast} (F)
$$
are isomorphisms (where $u_i : X \to X_i$ and $v_i : Y \to Y_i$ are
the canonical homomorphisms).
\end{enonce*}

The proof is essentially the same as before.

\begin{enonce*}[remark]{Remarks 5.14}\etiquette[5.14]{VII.rem5.14}
\begin{enumerate}
\renewcommand{\theenumi}{\alph{enumi}}
\renewcommand{\labelenumi}{\theenumi)}
\item The\pageoriginale preceding limit-passage results remain valid
for $H^{0}$ \resp $f_{\ast}$, for sheaves \emph{of sets} (instead of
abelian sheaves), and for $H^{1}$ \resp $R^1 f_{\ast}$ for sheaves of
groups (not necessarily commutative), where the noncommutative $H^1$
and $R^1 f_{\ast}$ are defined as usual in terms of torsors (or
principal homogeneous bundles) (\cf Giraud's thesis
\footnote{cited p.\pagehyperref[page7]{7} footnote $(\ast)$}). This can
be checked directly in the general context of VI. It is certainly
possible (and doubtless useful) also to give a variant for the
noncommutative $H^2$ of ``links'' studied by Giraud.

\item The limit-passage results for fibered sites developed in VI
assumed only that the base category was a filtered category, without it
being necessary to suppose that it was associated with a filtered
preordered set. The case of an arbitrary filtered index category is
also the natural framework for the statements developed in the present
${\rm n^{\circ}}$. If we have worked in an overly restrictive
framework, it was in order to be able to give correct references to
\EGA IV, where it is unfortunately assumed, in questions of passage to
the limit (from §8 onward), that the index category is defined in terms
of a filtered preordered set. We shall nevertheless assume henceforth
that all the results used from \EGA IV are valid for arbitrary filtered
index categories (the proofs given in \loccit being valid, essentially
without change, in this more general case). Thus we shall also use,
without further comment, the results of the present ${\rm n^{\circ}}$
in the case where $I$ is replaced by an arbitrary filtered index
category.
\end{enumerate}
\end{enonce*}

\ifx\danslelivre\undefined
\end{document}
\fi

% END INCLUDED SGA4_English_current_buildable_source_20260906/source/07/drafts/SGA4_VII_LINES_666_1032_English.texfrag


% END INCLUDED SGA4_English_current_buildable_source_20260906/source/07/07_English.tex


\makeatletter
\def\input@path{{../08/}}
\makeatother
% BEGIN INCLUDED SGA4_English_current_buildable_source_20260906/source/08/08_English.tex
\ifx\danslelivre\undefined
% Release engine: XeLaTeX (the LuaTeX compatibility layer breaks Xy-pic output).
\documentclass[10pt]{../fichiers-configuration-latex/smfart}
% BEGIN INCLUDED SGA4_English_current_buildable_source_20260906/source/fichiers-configuration-latex/paquets.tex
\usepackage{polyglossia}
%\usepackage[francais]{babel}
\usepackage{fontspec}
\usepackage[math-style=ISO,bold-style=upright,sans-style=upright]{unicode-math}
%\usepackage[mathletters]{ucs}
%\usepackage[utf8]{inputenc}
%\usepackage[T1]{fontenc}
\usepackage[unicode,hypertexnames=false,colorlinks=true,linkcolor=blue,citecolor=Sepia,urlcolor=magenta]{hyperref}
\setmainlanguage{french}
\setmainfont[Ligatures=TeX]{Linux Libertine O}
% variante possible : TeXGyrePagella ?
\setmathfont{XITS Math}

% maths
\usepackage{amsmath}
%\usepackage{amssymb}
%\usepackage{mathrsfs}
%\usepackage{stmaryrd}
%\usepackage{euscript}

% couleurs
\usepackage[usenames,dvipsnames]{xcolor}

% graphiques
\usepackage{graphicx}
\usepackage[all]{xy}
%\usepackage{tikz}
%\usetikzlibrary{matrix,arrows,calc}
%\tikzset{column sep=3em,row sep=2em,text height=1.5ex, text depth=0.25ex,>=angle 90}
%\tikzset{description/.style={fill=white,inner sep=1pt}}

% hyperliens, 'etiquettes
\usepackage{xr-hyper}
\usepackage{../fichiers-configuration-latex/smfthm}
%\ifx\cacheretiquettes\undefined
%\usepackage[inline]{showlabels} \else \fi % pour rendre visibles les 'etiquettes
%\usepackage{../fichiers-configuration-latex/makeidx} % pour faire un index
%\usepackage[pagebackref,unicode]{hyperref}  % pour voir o`u sont cit'ees les r'ef'erences
% hyperref is loaded once above with the cumulative-reader destination policy.
%\hypersetup{breaklinks=true} %permet le retour `a la ligne dans les liens trop longs. Incompatible avec pdflatex
%\renewcommand*{\backref}[1]{{\small [cit\'e page(s)~#1.]}}

\usepackage{xspace}
\usepackage{pifont} % pour expose XVII

% END INCLUDED SGA4_English_current_buildable_source_20260906/source/fichiers-configuration-latex/paquets.tex

% BEGIN INCLUDED SGA4_English_current_buildable_source_20260906/source/fichiers-configuration-latex/numerotation.tex
\renewcommand{\thefootnote}{\roman{footnote}} % pour mettre les notes en bas de page sous la forme (ii) et non (2), qui ressemble trop (?) a un exposant.

\newcounter{pageoriginal}
\newcommand{\pageoriginale}{\refstepcounter{pageoriginal}\marginpar{\footnotesize\textbf{\thepageoriginal}}\xspace}
\let\pageoriginaled\pageoriginale

\usepackage{ifthen}
\newboolean{original}
\def\sisi{\ifthenelse{\boolean{original}}}
\setboolean{original}{false}

\newcounter{cnote}
\def\nde#1{\ifthenelse{\boolean{original}}{}{\stepcounter{cnote}\def\thefootnote{\arabic{cnote}}\protect\footnote[\value{cnote}]{N.D.E.: #1}}}
\def\ndemark{\ifthenelse{\boolean{original}}{}{\stepcounter{cnote}\def\thefootnote{\arabic{cnote}}\protect\footnotemark[\value{cnote}]}}
\def\ndetext#1{\ifthenelse{\boolean{original}}{}{\def\thefootnote{\arabic{cnote}}\protect\footnotetext[\value{cnote}]{N.D.E.: #1}}}

% END INCLUDED SGA4_English_current_buildable_source_20260906/source/fichiers-configuration-latex/numerotation.tex

% BEGIN INCLUDED SGA4_English_current_buildable_source_20260906/source/fichiers-configuration-latex/macros.tex
\setmainlanguage{french}
\setotherlanguage{english}


%% Fixed, engine-independent release identifier.  The inherited dynamic Git
%% probe used Lua io.popen, exposed build-state-dependent output, and printed
%% a visible question mark outside a Git checkout.
\def\showgitstatus{English release r2, 2026-07-23}


% macros pour LaTeX
\renewcommand{\mathsf}{}

% Police différente pour U+2113 SCRIPT SMALL L
\setmathfont[range="2113]{latinmodern-math.otf}

% Give every inherited semantic label its own stable PDF destination instead
% of inheriting whichever automatic equation/page anchor happened to be live.
% The legacy label remains as a compatibility alias used by existing edges.
\makeatletter
\def\etiquette[#1]#2{%
  \Hy@raisedlink{\hyper@anchorstart{sga4.semantic.#2}\hyper@anchorend}%
  \def\@currentHref{sga4.semantic.#2}%
  \def\@currentlabel{#1}%
  \label{#2}%
}
\makeatother
\def\eqhyperref[#1]#2{(\hyperref[#1]{#2})}
\def\pagehyperref[#1]#2{(page \hyperref[#1]{#2})}
\def\hyperrefext[#1][#2]#3{\hyperref[#2]{#1 #3}}
\def\eqhyperrefext[#1][#2]#3{\hyperref[#2]{#1 #3}}     %cf. exposes 17
\def\neweqhyperrefext[#1][#2]#3{\hyperref[#2]{#1 #3}}  %cf. exposes 17

\def\shat#1{\protect\widehat{#1}}
\def\stilde#1{\protect\widetilde{#1}}
%%%%%%%%%%%%%%%%%%


\def\etcn{etc.\xspace}
\def\etco{etc...\xspace}
\def\etc{\sisi{\etco}{\etcn}}
\def\ldotsn{\sisi{...}{\etcn}}

\def\cf{{cf.}\kern.3em}
\def\Cf{{Cf.}\kern.3em}
\def\ie{{i.e.}\kern.3em}
\def\iev{{i.e.},\xspace}
\def\loccito{loc.\kern2pt cit.\xspace}
\def\loccitn{\textit{loc.\kern2pt cit.}\xspace}
\def\loccit{\sisi{\loccito}{\loccitn}}
\def\EGA{\textup{EGA}\kern.2em}
\def\SGA{\textup{SGA}\kern.3em}

\def\quoi{}

\def\quav{\quad,\quad}
\def\quaet{\quad \text{et}\quad}
\def\fibreproduct#1#2#3{#1 ×_{#3} #2}
\let\fprod\fibreproduct
\def\ccheck#1#2{\mbox{${\check{\check{\rm #1}}}^{\lower2pt\hbox{$\scriptstyle #2$}}$}}
\def\sheaf#1{\protect\mathcal{#1}}
\def\vieuxsheaf#1{{#1}} % cf. problème avec sheaf{a\Phi_i} par exemple...
\DeclareMathOperator{\Hom}{Hom}
\DeclareMathOperator{\can}{can}
\DeclareMathOperator{\Ind}{Ind}
\DeclareMathOperator{\ind}{ind}
%\DeclareMathOperator{\Hhomtop}{\mathcal{H}\!\mathit{omtop}}
\DeclareMathOperator{\Hhomtop}{\mathbf{Homtop}}
\DeclareMathOperator{\Point}{Point}
\DeclareMathOperator{\Pro}{Pro}
\DeclareMathOperator{\Top}{Top}
\DeclareMathOperator{\Mod}{Mod}
\DeclareMathOperator{\Cart}{Cart}
%\DeclareMathOperator{\Mmorsite}{\mathcal{M}\!\mathit{orsite}}
\DeclareMathOperator{\Mmorsite}{\mathbf{Morsite}}
\DeclareMathOperator{\loc}{loc}
\DeclareMathOperator{\Loc}{Loc}
\DeclareMathOperator{\Pic}{Pic}
\DeclareMathOperator{\Et}{Et}
\DeclareMathOperator{\card}{card}
\DeclareMathOperator{\Ens}{Ens}
\DeclareMathOperator{\h}{\underline{H}}
\DeclareMathOperator{\Fonct}{Fonct}
\DeclareMathOperator{\Ker}{Ker}
\DeclareMathOperator{\but}{but}
\DeclareMathOperator{\Ob}{Ob}
\DeclareMathOperator{\ob}{ob}
\DeclareMathOperator{\Lib}{Lib}
\DeclareMathOperator{\esj}{esj}
\DeclareMathOperator{\id}{id}
\DeclareMathOperator{\Id}{Id}
\DeclareMathOperator{\ad}{ad}
\DeclareMathOperator{\Iim}{Im}
\DeclareMathOperator{\pr}{pr}
\DeclareMathOperator{\Cat}{Cat}
\DeclareMathOperator{\car}{car}
\DeclareMathOperator{\SysInd}{SysInd}
\DeclareMathOperator{\Sup}{Sup}
\DeclareMathOperator{\Coker}{Coker}
\DeclareMathOperator{\Filt}{Filt}
\DeclareMathOperator{\acc}{acc}
\DeclareMathOperator{\sur}{sur}
\DeclareMathOperator{\cat}{cat}
\DeclareMathOperator{\Ord}{Ord}
\DeclareMathOperator{\Gr}{Gr}
\DeclareMathOperator{\Cov}{Cov}
\DeclareMathOperator{\Caf}{Caf}
\DeclareMathOperator{\coker}{coker}
\DeclareMathOperator{\coim}{coim}
\DeclareMathOperator{\im}{im}
\DeclareMathOperator{\Ouv}{Ouv}
\DeclareMathOperator{\diag}{diag}
\DeclareMathOperator{\Esp}{Esp}
\DeclareMathOperator{\TOP}{TOP}
\DeclareMathOperator{\Unif}{Unif}
\DeclareMathOperator{\Iind}{\mathbf{Ind}}
\DeclareMathOperator{\Hhomcart}{\mathbf{Homcart}}
\DeclareMathOperator{\Hhomtopess}{\mathbf{Homtopess}}
\DeclareMathOperator{\Ppoints}{\mathbf{Points}}
\DeclareMathOperator{\Ppoint}{\mathbf{Point}}
\DeclareMathOperator{\Ttors}{\mathbf{Tors}}
\DeclareMathOperator{\Ssex}{\mathbf{Sex}}
\DeclareMathOperator{\Ffib}{\mathbf{Fib}}
\DeclareMathOperator{\Vvois}{\mathbf{Vois}}
\DeclareMathOperator{\Ppointess}{\mathbf{Pointess}}
\DeclareMathOperator{\scrHom}{\mathbf{Hom}}
\DeclareMathOperator{\sheafPic}{\mathbf{Pic}}
\DeclareMathOperator{\sheafAut}{\mathbf{Aut}}
\DeclareMathOperator{\Hhomtoplocan}{\mathbf{Homtoplocan}}
\DeclareMathOperator{\Hhomtopan}{\mathbf{Homtopan}}
\DeclareMathOperator{\Pptgeom}{\mathbf{Ptgéom}}
\DeclareMathOperator{\Morsite}{Morsite}
\DeclareMathOperator{\Fib}{Fib}
\DeclareMathOperator{\Fais}{Fais}
\DeclareMathOperator{\kar}{kar}
\DeclareMathOperator{\Groupes}{Groupes}
\DeclareMathOperator{\Restr}{Restr}
\DeclareMathOperator{\Isom}{Isom}
\DeclareMathOperator{\Prol}{Prol}
\DeclareMathOperator{\Res}{Res}
\DeclareMathOperator{\Tors}{Tors}
\DeclareMathOperator{\esp}{esp}
\DeclareMathOperator{\es}{es}
\DeclareMathOperator{\od}{od}
\DeclareMathOperator{\gr}{gr}
\DeclareMathOperator{\sob}{sob}
\DeclareMathOperator{\dfn}{dfn}
\DeclareMathOperator{\Sch}{Sch}
\DeclareMathOperator{\ouv}{ouv}
\DeclareMathOperator{\cons}{cons}
\DeclareMathOperator{\Nerf}{Nerf}
\DeclareMathOperator{\Points}{Points}
\DeclareMathOperator{\Espsob}{Espsob}
\DeclareMathOperator{\Homtop}{Homtop}
\DeclareMathOperator{\Pt}{Pt}
\DeclareMathOperator{\Quot}{Quot}
\DeclareMathOperator{\glob}{glob}
\DeclareMathOperator{\Homfib}{Homfib}
\DeclareMathOperator{\Qouv}{Qouv}
\DeclareMathOperator{\Quotouv}{Quotouv}
\DeclareMathOperator{\ens}{ens}
\DeclareMathOperator{\ab}{ab}
\DeclareMathOperator{\an}{an}
\DeclareMathOperator{\fs}{fs}
\DeclareMathOperator{\red}{red}
\DeclareMathOperator{\Ab}{Ab}
\DeclareMathOperator{\cosq}{cosq}
\DeclareMathOperator{\Ann}{Ann}
\DeclareMathOperator{\Zar}{Zar}


%%%%%%%%%%%%
%%%%% SGA4-2
\DeclareMathOperator{\Sex}{Sex}
\DeclareMathOperator{\vois}{vois}
\DeclareMathOperator{\Vois}{Vois}
\DeclareMathOperator{\Ext}{Ext}
\DeclareMathOperator{\ext}{ext}
\DeclareMathOperator{\Fer}{Fer}
\DeclareMathOperator{\Add}{Add}
\DeclareMathOperator{\Cocart}{Cocart}
\DeclareMathOperator{\cosk}{cosk}
\DeclareMathOperator{\plat}{plat}

\DeclareMathOperator{\faisc}{faisc}
\DeclareMathOperator{\sk}{sk}
\DeclareMathOperator{\ch}{ch}

\def\HH{\mathbf{H}}
\def\qc{\mathrm{qc}}
\def\et{\mathrm{et}}
\def\coh{\mathrm{coh}}

\def\Hvr#1{%
\begin{array}[b]{l}
{\kern2pt\scriptstyle r}\\[-8pt]
{\kern1pt\scriptstyle \vee}\\[-5pt]
H^{#1}
\end{array}}

\def\Hvv#1{%
\begin{array}[b]{l}
\kern2pt\rotatebox{90}{${\scriptscriptstyle\ll}$}\\[-6pt]
H^{#1}
\end{array}}

\def\Hvi#1{%
\begin{array}[b]{l}
{\kern0.2pt\scriptscriptstyle \infty}\\[-8pt]
{\kern1pt\scriptstyle \vee}\\[-5pt]
H^{#1}
\end{array}}

\def\Hv#1{%
\begin{array}[b]{l}
{\kern1pt\scriptstyle \vee}\\[-5pt]
H^{#1}
\end{array}}

\DeclareMathOperator{\cart}{cart}
\DeclareMathOperator{\Ffaisc}{\mathbf{Faisc}}


\def\Ev#1{%
\begin{array}[b]{l}
{\kern1pt\scriptstyle \vee}\\[-5pt]
E'_{#1}
\end{array}}

\DeclareMathOperator{\Cont}{Cont}
\DeclareMathOperator{\Ccont}{\mathscr{C}\!\mathit{ont}}
\DeclareMathOperator{\Site}{Site}
\newcommand\MM{\mathbf{M}}
\def\limtop#1{{\displaystyle{\varprojlim_{#1}\!\mbox{\rm top }}}}
\def\Limtop{\mathop{\varprojlim\!\mbox{\rm top}}}
\DeclareMathOperator{\spec}{spec}
\def\ss{\scriptstyle{\mathscr{S}}}
\DeclareMathOperator{\Gal}{Gal}
\def\smallm{\scriptstyle{\mathscr{M}}}
\DeclareMathOperator{\End}{End}
\DeclareMathOperator{\ass}{ass}
\DeclareMathOperator{\zar}{zar}
\DeclareMathOperator{\Ogg}{Ogg}
\DeclareMathOperator{\res}{res}
\DeclareMathOperator{\tr}{tr}
\DeclareMathOperator{\Supp}{Supp}
\DeclareMathOperator{\codim}{codim}
\DeclareMathOperator{\rk}{rk}
\DeclareMathOperator{\cur}{cur}
\DeclareMathOperator{\cdqc}{cdqc}
\newcommand\DD{\mathbf{D}}
\DeclareMathOperator{\cl}{cl}
\def\JJ{\mathbf{J}}
\def\FF{\mathbf{F}}
\def\BB{\mathbf{B}}
\DeclareMathOperator{\Alget}{Alget}
\DeclareMathOperator{\Ass}{Ass}
\DeclareMathOperator{\supp}{supp}
\DeclareMathOperator{\rad}{rad}
\DeclareMathOperator{\prof}{prof}
\DeclareMathOperator{\cint}{int}
\DeclareMathOperator{\sic}{sic}
\DeclareMathOperator{\Tor}{Tor}
\DeclareMathOperator{\tors}{tors}
\DeclareMathOperator{\Sym}{Sym}
\DeclareMathOperator{\parf}{parf}
\DeclareMathOperator{\tor}{tor}
\DeclareMathOperator{\defn}{defn}
\DeclareMathOperator{\Div}{Div}
\DeclareMathOperator{\Pol}{Pol}
\DeclareMathOperator{\Fex}{Fex}
\DeclareMathOperator{\Lie}{Lie}
\DeclareMathOperator{\tot}{tot}
\DeclareMathOperator{\Aut}{Aut}
\DeclareMathOperator{\HOM}{HOM}
\DeclareMathOperator{\pch}{pch}
\DeclareMathOperator{\Ch}{Ch}
\DeclareMathOperator{\EXT}{EXT}
\DeclareMathOperator{\PIC}{PIC}
\DeclareMathOperator{\Rf}{Rf}
\DeclareMathOperator{\inj}{inj}
\newcommand\VV{\mathbf{V}}

\def\Xv{%
\begin{array}[b]{l}
{\kern1pt\scriptstyle \vee}\\[-5pt]
X
\end{array}}

%%%%%%%
\newcommand\TT{\mathbf{T}}
\newcommand\XX{\mathbf{X}}
\newcommand\EE{\mathbf{E}}
%%%%%%

\def\qedsymbol{\textup{C.Q.F.D.}}
\let\CQFD\qed
\let\cqfd\qed

%%%%%\bigdot
\font\bigsymb=cmsy10 at 4pt
\def\bigdot{{\kern1.2pt\raise 1.5pt\hbox{\bigsymb\char15}}}
%%%%%

\def\Lright{\mathop{\vcenter{\hbox{\hbox{$\longrightarrow$}\lower8pt \hbox{\kern-16pt\hbox{$\longrightarrow$}}}}}}

\def\Aa{{\mathcal A}}
\def\Bb{{\mathcal B}}
\def\Cc{{\mathcal C}}
\def\Dd{{\mathcal D}}
\def\Ee{{\mathcal E}}
\def\Ff{{\mathcal F}}
\def\Gg{{\mathcal G}}
\def\Hh{{\mathcal H}}
\def\Ii{{\mathcal I}}
\def\Jj{{\mathcal J}}
\def\Kk{{\mathcal K}}
\def\Ll{{\mathcal L}}
\def\Mm{{\mathcal M}}
\def\Nn{{\mathcal N}}
\def\Oo{{\mathcal O}}
\def\Pp{{\mathcal P}}
\def\Qq{{\mathcal Q}}
\def\Rr{{\mathcal R}}
\def\Ss{{\mathcal S}}
\def\Tt{{\mathcal T}}
\def\Uu{{\mathcal U}}
\def\Vv{{\mathcal V}}
\def\Ww{{\mathcal W}}
\def\Xx{{\mathcal X}}
\def\Yy{{\mathcal Y}}
\def\Zz{{\mathcal Z}}

\newcommand{\CC}{\mathbf{C}}
\newcommand{\GG}{\mathbf{G}}
\newcommand{\KK}{\mathbf{K}}
\newcommand{\LL}{\mathbf{L}}
\newcommand{\NN}{\mathbf{N}}
\newcommand{\PP}{\mathbf{P}}
\newcommand\QQ{\mathbf{Q}}
\newcommand\RR{\mathbf{R}}
\newcommand\ZZ{\mathbf{Z}}

\def\to{\mathchoice{\longrightarrow}{\rightarrow}{\rightarrow}{\rightarrow}}
\def\from{\mathchoice{\longleftarrow}{\leftarrow}{\leftarrow}{\leftarrow}}
\def\mto{\mathchoice{\longmapsto}{\mapsto}{\mapsto}{\mapsto}}
\def\mfrom{\mathrel{\from\joinrel\mapstochar\,}}
\def\hto{\mathrel{\lhook\joinrel\to}}
\def\hfrom{\mathrel{\from\joinrel\rhook}}

\def\implies{\mathchoice{\Longrightarrow}{\Rightarrow}{\Rightarrow}{\Rightarrow}}
\def\implied{\mathchoice{\Longleftarrow}{\Leftarrow}{\Leftarrow}{\Leftarrow}}
\def\iff{\mathchoice{\Longleftrightarrow}{\Leftrightarrow}{\Leftrightarrow}{\Leftrightarrow}}

\def\To#1{\mathchoice{\xrightarrow{\labelstyle~~#1~}}{\stackrel{#1}{\longrightarrow}}{}{}}
\def\Mto#1{\mathrel{\mapstochar\To{#1}}}
\def\From#1{\mathchoice{\xleftarrow{\labelstyle~~#1~}}{\stackrel{#1}{\longleftarrow}}{}{}}

\newcommand{\isomto}{\overset{\sim}{\longrightarrow}}
\newcommand{\isomfrom}{\overset{\sim}{\longleftarrow}}

\newcommand{\IFF}{\unskip~$\iff$~\ignorespaces}
\newcommand{\THEN}{\unskip~$\implies$~\ignorespaces}

\let\leq\leqslant
\let\geq\geqslant
\let\le\leqslant
\let\ge\geqslant

\def\cf{{cf.}\kern.3em}
\def\Cf{{Cf.}\kern.3em}
\def\ie{{i.e.}\kern.3em}
\def\iev{{i.e.},\xspace}
%\def\loccito{loc.\kern2pt cit.\xspace}
%\def\loccitn{\textit{loc.\kern2pt cit.}\xspace}
%\def\loccit{\sisi{\loccito}{\loccitn}}
\def\EGA{\textup{EGA}\kern.2em}
\def\SGA{\textup{SGA}\kern.3em}

\let\oldS\S
\def\S{\oldS\kern .15em}
\def\ptbl{.\kern .2em}
\def\resp{resp.\kern.3em}
\def\Exp{Exp.\kern.3em}

\let\bar\overline
\let\hat\widehat
\let\tilde\widetilde
\let\dpl\displaystyle
\let\ldots\dots
\let\emptyset\varnothing
\let\skpt\skippointrait
\let\epsilon\varepsilon

\def\sfrac#1#2{{#1}/{#2}}

\DeclareMathOperator{\cd}{cd}
%\DeclareMathOperator{\Hhom}{\mathcal{H}\!\mathit{om}}
\DeclareMathOperator{\Hhom}{\mathbf{Hom}}
%\DeclareMathOperator{\Eext}{\mathcal{E}\!\mathit{xt}}
\DeclareMathOperator{\Eext}{\mathbf{Ext}}
\DeclareMathOperator{\Spec}{Spec}
\DeclareMathOperator{\Tr}{Tr}

\let\emptyset=\varnothing

%%%%%%%%%%%%%%%
%%% expose XVII
%%%%%%%%%%%%%%%

\makeatletter
\let\old@optexpose\optexpose
\newcommand{\empty@optexpose}[1]{}
\newcommand{\paren@optexpose}[1]{\old@optexpose{#1}(}
\newcommand{\refnoexpose}[1]{\let\optexpose\empty@optexpose\ref{#1}\let\optexpose\old@optexpose}%\refnoexpose{X.cor5.2} engendre "5.2" au lieu de "X 5.2"
\newcommand{\neweqref}[1]{\let\optexpose\paren@optexpose\textup{\ref{#1})}\let\optexpose\old@optexpose}%\neweqref{XII.eq4.1} engendre "XII (4.1)" au lieu de "(XII 4.1)"
\makeatother
\newcommand{\refvb}[2][VI${}_{\text{B}}$ ]{\sisi{#1\refnoexpose{#2}}{\ref{#2}}}

\newcommand{\leftexp}[2]{{\vphantom{#2}}^{#1}{#2}}
\newcommand{\AAA}{\mathbf{A}}
\newcommand{\Sf}{\mathfrak{S}}
\DeclareMathOperator*{\quotprojlim}{\text{``$\varprojlim$''}}
\DeclareMathOperator*{\quotinjlim}{\text{``$\varinjlim$''}}
\DeclareMathOperator{\Fl}{F\sisi{\ell}{l}}
\DeclareMathOperator{\Pl}{P\sisi{\ell}{l}}
%\DeclareMathOperator{\cl}{c\sisi{\ell}{l}}
\DeclareMathOperator{\tdf}{tdf}
\DeclareMathOperator{\rg}{rg}
\newcommand{\iInd}{\sisi{I}{i}nd}
\newcommand{\iind}{ind}
\newcommand{\ppro}{pro}
\newcommand{\tTor}{\sisi{T}{t}or}
\newcommand{\ttor}{tor}
\newcommand{\chap}{Chap.}
\newcommand{\Ffex}{\Ff\mathit{ex}}
\newcommand{\sheafk}{\underline{k}}
\newcommand{\sheafn}{\underline{n}}
\newcommand{\sheafGamma}{\underline{\Gamma}}

% Repeated references to an already typeset footnote must not use
% \footnotemark without a matching \footnotetext: Hyperref would then create
% a live annotation to a destination that cannot exist.  This helper keeps the
% current printed footnote mark and links it to an explicit stable note target.
\makeatletter
\newcommand{\sgarepeatedfootnotemark}[1]{%
  \begingroup
  \edef\@thefnmark{\thefootnote}%
  \hyperlink{#1}{\@makefnmark}%
  \endgroup}
\makeatother

\sisi{
\DeclareMathOperator{\simarrow}{\rule{.15mm}{1.8mm}\kern -3.4pt\rightsquigarrow}
\newcommand{\sheafell}{\underline{\ell}}
\newcommand{\sheafZZ}{\underline{\ZZ}}
}{
\newcommand{\nderef}[1]{la N.D.E.~(\ref{#1}) page~\pageref{#1}}
}

%%%%%%%%%%%%%%%%%
%%%%% expose XVII
%%%%%%%%%%%%%%%%%

% END INCLUDED SGA4_English_current_buildable_source_20260906/source/fichiers-configuration-latex/macros.tex


\externaldocument{../01/01_English.tex}
\externaldocument{../02/02_English.tex}
\externaldocument{../03/03_English.tex}
\externaldocument{../04/04_English.tex}
\externaldocument{../05/05_English.tex}
\externaldocument{../05bis/05bis_English.tex}
\externaldocument{../06/06_English.tex}
\externaldocument{../07/07_English.tex}

\setboolean{original}{false}

\begin{document}

\begin{center}
VIII. Fiber functors, supports, cohomological study of finite morphisms\\
A. Grothendieck
\end{center}

\selectlanguage{english}version: \showgitstatus
\tableofcontents

\else
\chapter{Fiber functors, supports, cohomological study of finite morphisms}
%\addtocontents{toc}{A. Grothendieck}
\begin{center}
A. Grothendieck
\end{center}
\fi
\phantomsection\MakeLinkTarget*{sga4.exp.VIII}\label{sga4.exp.VIII}

% BEGIN INCLUDED SGA4_English_current_buildable_source_20260906/source/08/drafts/SGA4_VIII_LINES_43_576_English.texfrag
\setcounter{pageoriginal}{366}

% remarks:
% -- title too long at the top of the page

\section{Topological invariance of the étale
  topos}\etiquette[1]{VIII.sec1}

\begin{enonce*}{Theorem 1.1}\etiquette[1.1]{VIII.thm1.1}
Let \pageoriginale  $f : X'\to X$ be an \emph{integral surjective
radicial} morphism. Then the base-change functor
$$
f^{\ast} : X_{\text{\rm ét}} \to X'_{\text{\rm ét}}
$$
is an equivalence of sites (\ie an equivalence of categories that is
continuous, as is every quasi-inverse functor). Consequently, the
functors
\begin{align*}
f_{\ast} : \widetilde{X'_{\text{\rm ét}}} &\to \widetilde{X_{\text{\rm ét}}}\\
f^{\ast} : \widetilde{X_{\text{\rm ét}}} &\to \widetilde{X'_{\text{\rm
      ét}}}
\end{align*}
are equivalences of categories, quasi-inverse to one another.
\end{enonce*}

The first assertion is well known (\SGA\sisi{}{ 1} IX 4.10 and 4.11) when
$f$ is of finite presentation or when $f$ is flat; the general case
reduces to the case in which $f$ is of finite presentation. Indeed, it
is already known that the functor $f^{\ast}$ is fully faithful (\SGA 1 IX
3.4). It is enough to prove that $f^{\ast}$\pageoriginale  is
essentially surjective, \ie that every $Z'$ étale over $X'$ ``comes
from'' a $Z$ étale over $X$. Using the fact that $f$ is a universal
homeomorphism and the full faithfulness of $f^{\ast}$, we are reduced
to the case where $X$, $X'$, and $Z'$ are affine, the spectra of rings
$A$, $A'$, and $B'$. Writing $A'$ as the inductive limit of its
finitely generated $A$-subalgebras $A'_{i}$, $B'$ comes from an étale
algebra $B'_{i}$ over some $A'_{i}$ (\cf \EGA IV 8), which reduces us
to the case in which $A'$ is integral and of \emph{finite type} over
$A$, hence finite over $A$. We then have an $A$-isomorphism $A'\simeq
A''/J$, where $A'' \simeq A [t_1, \ldots t_n]$ and $J$ is an ideal of
$A''$. Writing $J$ as the inductive limit of its finitely generated
subideals $J_i$, and setting $A'_i = A'' /J_i$, we again have $A' =
\varinjlim_i A'_{i}$, so $B'$ comes from an étale algebra $B'_i$ over
some $A'_i$ (\emph{loco citato}). On the other hand, since $\Spec
A'\to \Spec A$ is surjective, the same holds for the $\Spec A'_i \to
\Spec A$; moreover, one easily verifies that since $\Spec A'\to
\Spec A$ is radicial, the same holds for $\Spec A'_i \to \Spec A$
for sufficiently large $i$ (apply \EGA IV 1.9.9 to $\Spec
A$\footnote{Let $S_i \subset S = \Spec A$ be the set of $s \in S$
such that the fiber of $X_i = \Spec A'_i$ \sisi{as sets}{at $s$} is
of separable rank 1. Then the $S_i$ are \emph{constructible} subsets
of $S$ (\EGA IV 9.7.9) forming an increasing family; their union is
$S$ by the hypothesis on $S'=\Spec A'\to S$ and the fact that the
fibers of $S'_i$ are \emph{noetherian} schemes. Thus $S=S_i$ for
sufficiently large $i$, by \EGA IV 1.9.9, \cqfd.}). This reduces us
to the case in which $A'$ is one of the $A'_i$, hence is of finite
presentation over $A$.

This proves the first assertion of
\sisi{\hyperref[VIII.prop2.1]{2.1}}{\hyperref[VIII.thm1.1]{1.1}}.
The assertion that $f_{\ast}$ and $f^{\ast}$ are mutually
quasi-inverse equivalences follows at once.

\begin{enonce*}{Corollary 1.2}\etiquette[1.2]{VIII.cor1.2}
Let $F$ be an abelian sheaf on $X$, and let $F'$ be its inverse image
on $X'$. Then the canonical homomorphisms
\begin{equation*}
H^{i} (X, F) \to H^{i} (X', F')\tag{$\ast$}
\end{equation*}
are isomorphisms. Likewise, if $F'$ is an abelian sheaf on $X'$ and
$F$ its direct image on $X$, the canonical homomorphisms $(\ast)$ are
isomorphisms.
\end{enonce*}

\begin{enonce*}[remark]{Examples 1.3}\etiquette[1.3]{VIII.exem1.3}
Here\pageoriginale  are some frequently used applications of
\hyperref[VIII.thm1.1]{1.1}:
\begin{enumerate}
\renewcommand{\theenumi}{\alph{enumi}}
\renewcommand{\labelenumi}{\theenumi)}
\item $X'$ is a subscheme of $X$ having the same underlying set as
  $X$, \ie defined by a nil-ideal $I$.

\item $X$ is a scheme over a separably closed field $k$, $k'$ is an
  algebraic closure of $k$, and $X'=X\otimes_k k'$.

\item Let $X$ be a geometrically unibranch scheme (for example, an
  algebraic curve over a field $k$ whose only singularities are
  ``cuspidal'' singularities, thus excluding ordinary double points in
  particular). Then, if $X'$ is the normalization of $X_{\red}$, by
  definition $X'\to X$ is radicial, so
  \sisi{\hyperref[VIII.prop2.1]{2.1}}{\hyperref[VIII.thm1.1]{1.1}}
  applies.
\end{enumerate}
\end{enonce*}

\begin{enonce*}[remark]{Remarks 1.4}\etiquette[1.4]{VIII.rem1.4}
A morphism $f$ as in \hyperref[VIII.thm1.1]{1.1} is a \emph{universal
  homeomorphism}, \ie it is a homeomorphism and remains so after every
  base extension. Conversely, if $f$ is a universal homeomorphism, one
  proves that $f$ satisfies the hypotheses of
  \sisi{\hyperref[VIII.prop2.1]{2.1}}{\hyperref[VIII.thm1.1]{1.1}}
  \footnote{\Cf \EGA IV 8.11.6 if $f$ is of finite presentation, and
  \EGA IV 18.12.11 in the general case.}, which explains the title of
  this section.

Let us point out explicitly, in connection with Example
\sisi{\hyperref[VIII.cor2.3]{2.3} c)}{\hyperref[VIII.exem1.3]{1.3} c)},
that if $X$ is an algebraic curve (over an algebraically closed field
$k$, to fix ideas) having at least one singular point that is not
``unibranch'' (for example, an ordinary double point), then (with $X'$
denoting the normalization of $X_{\red}$) the conclusion of
\sisi{\hyperref[VIII.cor2.2]{2.2}}{\hyperref[VIII.cor1.2]{1.2}}
already fails for $H^1$ with constant coefficients (compare \SGA 1 I
11 a) and \SGA 1 IX 5).
\end{enonce*}

\section{Sheaves on the spectrum of a field}\etiquette[2]{VIII.sec2}\pageoriginale

\begin{enonce*}{Proposition 2.1}\etiquette[2.1]{VIII.prop2.1}
Let $k$ be a field, $\bar{k}$ a separable closure of $k$, $\pi =
\Gal(\bar{k}/k)$ its topological Galois group, $X=\Spec k$, and
$\bar{X}=\Spec\bar{k}$ (on which $\pi$ acts on the left). Consider the
canonical functor
$$
i : X_{\text{\rm ét}} \to \widetilde{X_{\text{\rm ét}}}
$$
(defined by $i(X')(X'')=\Hom_X(X'',X')$), and the canonical functor
$$
j : \widetilde{X_{\text{\rm ét}}} \to \Bb_{\pi}
$$
from the category of sheaves on $X_{\text{\rm ét}}$ to the category
$\Bb_{\pi}$ \emph{(\cf IV \hyperref[IV.subsec2.7]{2.7})} of (discrete) sets on which $\pi$ acts
continuously on the left, defined by
$$
j(F) = \varinjlim_{\alpha} F(\Spec(k_{\alpha})) {\quoi},
$$
where $k_{\alpha}$ ranges over the finite subextensions of $\bar{k}$.
Then $i$ and $j$ are equivalences of categories. (Consequently, the
topos $\widetilde{X_{\text{\rm ét}}}$ is equivalent \emph{(IV \hyperref[IV.subsec3.4]{3.4})}
to the ``classifying topos'' $\Bb_{\pi^{\bigdot}}$.)
\end{enonce*}

The composite functor
$$
ji : X_{\text{\rm ét}} \to \widetilde{X_{\text{ét}}} \to (\pi\text{-}\Ens)
$$
is an equivalence of categories by the fundamental theorem of Galois
theory (in the form of \SGA I V). Since $\Bb_{\pi}$ is plainly a topos
(II 4.14), the same is therefore true of $X_{\text{ét}}$. Moreover, a
family $(X_i\to X)$ of morphisms in $X_{\text{ét}}$ is covering if and
only if it is surjective, \ie its image in $\Bb_{\pi}$ is surjective,
or, equivalently, covering for the canonical topology of $\Bb_{\pi}$.
This shows that the topology of $X_{\text{ét}}$ is indeed its canonical
topology. Consequently $i$ is an equivalence, and hence so is $j$.

\begin{enonce*}{Corollary 2.2}\etiquette[2.2]{VIII.cor2.2}
The functor\pageoriginale  $j$ of \hyperref[VIII.prop2.1]{2.1}
induces an equivalence \sisi{of}{between} the category of abelian
sheaves on $X=\Spec k$ and the category of Galois $\pi$-modules.
\end{enonce*}

Indeed, the latter are precisely the ``abelian groups'' of the topos
$\Bb_{\pi}$. The same argument, by ``transport of structure,'' gives:

\begin{enonce*}{Corollary 2.3}\etiquette[2.3]{VIII.cor2.3}
Let $F$ be an abelian sheaf on $X=\Spec k$, and let $M=j(F)$ be the
associated Galois $\pi$-module. Then there is a canonical isomorphism
of $\partial$-functors in $F$:
$$
H^{\ast} (X, F) \simeq H^{\ast} (\pi, M){\quoi},
$$
(where the right-hand side denotes Galois cohomology, studied, for
example, in Serre's course \emph{C.G.}).
\end{enonce*}

\begin{enonce*}{Corollary 2.4}\etiquette[2.4]{VIII.cor2.4}
Suppose that $k$ is separably closed. Then the functor
$$
\Gamma : \widetilde{X_{\text{\rm ét}}} \to (\Ens)
$$
is an equivalence of categories. If $F$ is an abelian sheaf on
$X=\Spec k$, then $H^{i}(X,F)=0$ for $i\neq0$. (In other words, the
topos $\widetilde{X_{\text{\rm ét}}}$ is a ``punctual topos''
\emph{(IV \hyperref[IV.subsec2.2]{2.2})}.)
\end{enonce*}

\begin{enonce*}{Corollary 2.5}\etiquette[2.5]{VIII.cor2.5}
Let $k'$ be a separably closed extension of a separably closed field
$k$, let $X=\Spec k$, $X'=\Spec k'$, and let $u:X'\to X$ be the
canonical morphism. Then the functor
$$
u^{\ast} : \widetilde{X_{\text{\rm ét}}} \to \widetilde{X'_{\text{\rm ét}}}
$$
is an equivalence of categories.
\end{enonce*}

\section{Fiber functors relative to geometric points of a
  scheme}\etiquette[3]{VIII.sec3}\pageoriginale

\subsection{}\etiquette[3.1]{VIII.subsec3.1}

Let $X$ be a scheme. We shall call a \emph{geometric point} of $X$
any $X$-scheme $\xi$ that is the spectrum of a \emph{separably closed}
field $\Omega$. The datum of such a point is therefore equivalent to
that of an ordinary point $x\in X$ together with a \emph{separably
closed} extension $\Omega$ of the residue field $k(x)$ of $x$. The
case most frequent for us will be that in which $\Omega$ is a
\emph{separable closure} of $k(x)$, which we shall then generally
denote by $\bar{k(x)}$; the corresponding geometric point of $X$ is
then denoted by $\bar{x}$. For a given $x\in X$, the $\bar{k(x)}$
(resp. $\bar{x}$) are determined up to \emph{nonunique}
$k$-isomorphism (resp. $X$-isomorphism).

\begin{enonce*}[remark]{Remarks 3.2}\etiquette[3.2]{VIII.rem3.2}
In most geometric questions, it is more natural to restrict to the
case in which $\Omega$ is \emph{algebraically closed}. The different
convention used here is specific to the study of the étale topology.
It is the convention we adopted in the definition of the fundamental
group in \SGA 1 V 7. Note, however, that the developments in
\emph{loco citato} are equally valid with the convention adopted here
(because the property of $\Omega$ used there is that every étale
covering of the spectrum of $\Omega$ is trivial, \ie precisely that
$\Omega$ is separably closed).
\end{enonce*}

\begin{enonce*}{Definition 3.3}\etiquette[3.3]{VIII.def3.3}
Let $X$ be a scheme, let $\xi$ be a geometric point of $X$, and let
$u:\xi\to X$ be its structural morphism. The (geometric) \emph{fiber
functor} relative to $\xi$, denoted $F\mapsto F_{\xi}$, is the
composite functor
$$
\widetilde{X_{\text{\rm ét}}} \to (\Ens)
$$
of
$$
\widetilde{X_{\text{\rm ét}}} \xrightarrow{u^{\ast}}
\widetilde{\xi_{\text{\rm ét}}} \xrightarrow{\Gamma_{\xi}} (\Ens){\quoi}.
$$
\end{enonce*}

Equivalently, by \hyperref[VIII.cor2.4]{2.4}, one may say that a
geometric point $\xi$ of $X$ defines a point of the topos
$X^{\sim}_{\text{ét}}$ (IV \hyperref[IV.def6.1]{6.1}), whose associated fiber functor
(\loccit) is simply being taken.

\setcounter{subsection}{3}
\subsection{}\etiquette[3.4]{VIII.subsec3.4}

Since the\pageoriginale  functor $\Gamma_{\xi}$ (the global-sections
functor on $\xi$) is an \emph{equivalence} of categories \hyperref[VIII.cor2.4]{(2.4)},
knowing the fiber functors $F\mapsto F_{\xi}$ is essentially
equivalent to knowing the inverse-image functors $u^{\ast}$. Moreover,
it follows at once from \hyperref[VIII.cor2.5]{2.5} that if $\xi$ is a
geometric point of $X$ for which there exists an $X$-morphism
$\xi'\to\xi$ (\ie $\xi'$ corresponds to a separably closed extension
$\Omega'$ of $\Omega$), then the canonical homomorphism
$F_{\xi}\to F_{\xi'}$ is a \emph{functorial isomorphism}
$$
F_{\xi} \simeq F_{\xi'}{\quoi}.
$$
This shows that, for the study of fiber functors, one may (after
replacing $\Omega$ by the separable algebraic closure of $k$ in
$\Omega$) restrict oneself to the case where $\Omega = k(\xi)$ is a
separable closure $\bar{k(x)}$ of $k(x)$, and that the fiber functors
corresponding to geometric points of $X$ lying over the same $x
\in X$ are isomorphic to one another (noncanonically).

In accordance with the conventions of \hyperref[VIII.subsec3.1]{3.1},
we shall denote by $F_{\bar{x}}$
a fiber functor corresponding to the choice of a separable closure
$\bar{k(x)}$ of $k(x)$.

Let us point out an evident transitivity property (which shows the
technical usefulness of not restricting oneself exclusively to geometric
points defined by separable closures of residue fields): Let $f :
X'\to X$ be a morphism of schemes, and let $u' : \xi'\to X'$ be a
geometric point of $X'$. Then $u = f u': \xi'\to X$ is a geometric
point of $X$, which we shall denote by $\xi$; relative to these geometric
points, there is an isomorphism, functorial in $F \in \Ob
\widetilde{X_{\text{ét}}}$:
$$
f^{\ast} (F)_{\xi'} \simeq F_{\xi}{\quoi}.
$$
Indeed, this\pageoriginale  follows from the transitivity formula for
inverse-image
functors
$$
(fu')^{\ast} \xrightarrow{\sim} u'^{\ast} f^{\ast}{\quoi}.
$$

\begin{enonce*}{Theorem 3.5}\etiquette[3.5]{VIII.thm3.5}
Let $X$ be a scheme.
\begin{enumerate}
\renewcommand{\theenumi}{\alph{enumi}}
\renewcommand{\labelenumi}{\theenumi)}
\item For every geometric point $\xi$ of $X$, the fiber functor $F
  \mapsto F_{\xi}$, $\widetilde{X_{\text{\rm ét}}} \to (\Ens)$ commutes with
  arbitrary $\varinjlim$ (indexed by categories $\in
  \Uu$) and with finite $\varprojlim$\footnote{Thus it is indeed
  a ``fiber functor'' in the sense of Exposé IV.}. In particular, it
  takes sheaves of groups (\resp abelian sheaves, \etc)
  to groups (\resp abelian groups, \etc).

\item As $x$ ranges over the points of $X$, the fiber functors
  $F_{\bar{x}}$ form a ``conservative'' family of functors, \ie if $u : F \to G$ is a homomorphism of
  sheaves on $X$, then $u$ is an isomorphism if and only if, for
  every $x \in X$, the corresponding homomorphism $u_{\bar{x}} :
  F_{\bar{x}} \to G_{\bar{x}}$ is an isomorphism\footnote{Thus there are
  ``enough fiber functors'' of the form $F \mapsto
  F_{\bar{x}}$, in the terminology of IV \hyperref[IV.prop6.5]{6.5}.}.
\end{enumerate}
\end{enonce*}

Taking into account the exactness properties in a), let us immediately
note the following formal consequences of b):

\begin{enonce*}{Corollary 3.6}\etiquette[3.6]{VIII.cor3.6}
Let $u : F \to G$ be a homomorphism of sheaves. Then $u$ is a
monomorphism (\resp an epimorphism) if and only if, for every $x
\in X$, $u_{\bar{x}}:F_{\bar{x}}\to G_{\bar{x}}$ is injective
(\resp surjective).
\end{enonce*}

\begin{enonce*}{Corollary 3.7}\etiquette[3.7]{VIII.cor3.7}
Let $u$, $v : F \to G$ be two morphisms of sheaves on $X$. Then
$u = v$ if and only if, for every $x \in X$, one has $u_{\bar{x}} =
v_{\bar{x}} : F_{\bar{x}} \to G_{\bar{x}}$. In particular, if $u$,
$v$ are two sections of $F$, one has\pageoriginale  $u = v$ if and
only if
$u_{\bar{x}} = v_{\bar{x}}$ in $F_{\bar{x}}$ for every $x \in X$.
\end{enonce*}

\begin{enonce*}{Corollary 3.8}\etiquette[3.8]{VIII.cor3.8}
Let $F \to G \to H$ be a sequence of two homomorphisms of sheaves on
$X$. Then the sequence is exact if and only if, for every $x \in X$,
the corresponding sequence $F_{\bar{x}} \to G_{\bar{x}} \to H_{\bar{x}}$
is exact.
\end{enonce*}

We leave to the reader the task of deducing the corollaries
\footnote{\Cf IV \hyperref[IV.exer6.5.2]{6.5.2}.}, %% issue: check that (**) is indeed present upon compilation
which will moreover be proved in part
in the proof that follows.

Assertion a) follows from the exactness properties already noted
in (\hyperrefext[VII][VII.subsec1.4]{1.4}) for every inverse-image functor, such as
$u^{\ast} : \widetilde{X_{\text{ét}}} \to
\widetilde{\xi_{\text{ét}}}$, and from the fact that
$\widetilde{\xi_{\text{ét}}} \to (\Ens)$ is an equivalence. To
prove b), let us note the following formula, which gives a procedure for
\emph{computing} fiber functors:

\begin{enonce*}{Proposition 3.9\footnotemark}\etiquette[3.9]{VIII.prop3.9}
\footnotetext{\Cf IV \hyperref[IV.eq6.8.3]{6.8.3}} %% issue: check the number of * upon compilation
Let $\xi \xrightarrow{u} X$ be a geometric
point of $X$, and let $C_{\xi}$ be the category of ``pointed étale
$X$-schemes at $\xi$'', \ie of pairs $(X', u')$, where
$X'\in \Ob X_{\text{\rm ét}}$ is a scheme étale over $X$, and $u'$ is an
$X$-morphism $\xi \to X'$. Then the opposite category $C^{0}_{\xi}$
is filtered and, for a varying sheaf $F$ (\resp presheaf $P$)
on $X$, there is a canonical functorial isomorphism
$$
F_{\xi} \simeq \varinjlim_{C^{0}_{\xi}} F(X'), (\emph{\resp}
(aP)_{\xi} \simeq \varinjlim_{C^{0}_{\xi}} P(X')){\quoi},
$$
(where $aP$ = the sheaf associated with $P$).
\end{enonce*}

Indeed, let us first note that for every presheaf $P$ on
$\xi_{\text{ét}}$,
if $aP$ is the associated sheaf, the natural homomorphism
$$
P(\xi) \to aP(\xi)
$$
is an isomorphism, as follows, for example, immediately from the
explicit computation of\pageoriginale  $aP$ as $L(LP)$ (II \hyperref[II.theo3.4]{3.19}). Let
$\varphi :
\widetilde{X_{\text{ét}}} \to \widetilde{\xi_{\text{ét}}}$ be the
``inverse-image'' functor. Then $u^{\ast} = \varphi^{\ast}$ is nothing
other (III \hyperref[III.prop2.3]{2.3}. (d)) than the composite $a\varphi^{\bigdot} i$, where $i :
\widetilde{X_{\text{ét}}} \to \widehat{X_{\text{ét}}}$
is the inclusion functor, and $a :
\widehat{\xi_{\text{ét}}} \to
\widetilde{\xi_{\text{ét}}}$ is the ``associated
sheaf'' functor. By the preceding remark we find
$$
u^{\ast} (F) (\xi) = (\varphi^{\bigdot} i(F)) (\xi){\quoi},
$$
so it is enough to apply the explicit expression for
$\varphi^{\bigdot}$ given in the relevant proof.\SGAArchiveOnly{\nde{The frozen source
leaves the exact internal reference unspecified.}}
The corresponding case, where one starts with a presheaf $P$ on $X$, is proved in the same way, by
using the functorial isomorphism $\varphi^{\ast} a(P) \simeq a
\varphi^{\ast} (P)$.\SGAArchiveOnly{\nde{The frozen source leaves the exact internal
reference unspecified.}}
The fact that the category $C_{\xi}$ is filtered (which remains true for \emph{every} $X$-scheme $\xi$, not
necessarily reduced to a geometric point)
follows immediately from the
fact that finite projective limits exist in $C_{\xi}$, which
follows from the fact that $X_{\text{ét}}$ is a subcategory of
$(\Sch) /X$ stable under finite projective limits.

Now let $u : F \to G$ be a homomorphism of sheaves on $X$
such that, for every $x \in X$, $u_{\bar{x}} : F_{\bar{x}} \to
G_{\bar{x}}$ is a monomorphism; let us prove that $u$ is a
monomorphism. For this, we must prove that if $X'\in \Ob
X_{\text{ét}}$, $\varphi$, $\psi \in F(X')$ are such that $u
(\varphi) = u (\psi)$, then $\varphi = \psi$. Replacing $X$ by
$X'$ and using the transitivity of fiber functors, we may
suppose $X = X'$. For every $x \in X$, we then have $u
(\varphi)_{\bar{x}} = u (\psi)_{\bar{x}}$, \ie $u_{\bar{x}}
(\varphi_{\bar{x}}) = u_{\bar{x}} (\psi_{\bar{x}})$, hence
$\varphi_{\bar{x}} = \psi_{\bar{x}}$ since $u_{\bar{x}}$ is a
monomorphism. Using \hyperref[VIII.prop3.9]{3.9}, %% XXX Reference to correct}
it follows that there exists an $X'_x \in \Ob X_{\text{ét}}$, whose image contains $x$, such
that $\varphi$ and $\psi$ have the same inverse image on $X'_x$. Since,
as $x$ varies in $X$, the $X'_x$ form a covering family of
$X$, it follows that $\varphi = \psi$.

This\pageoriginale proves that if the $u_{\bar{x}}$ are
monomorphisms, then so
is $u$. Suppose in addition that the $u_{\bar{x}}$ are
epimorphisms; let us prove that the same holds for $u$, and hence that $u$ is
an isomorphism. It remains to prove that for every $X'\in \Ob X_{\text{ét}}$
and every $\psi \in G (X')$, there exists $\varphi \in F(X')$ such that
$u(X') (\varphi) = \psi$. We may again suppose $X'= X$, and
using \hyperref[VIII.prop3.9]{3.9}, %% XXX Reference to correct
we see that, for every $x \in X$, \sisi{}{there} exists $X'_x \in \Ob
X_{\text{ét}}$ such that the image of $X'_x$ in $X$ contains $x$, and a
$\varphi_x \in F(X'_x)$ whose image in $G (X'_x)$ is
the inverse image of $\psi$. Using the fact that $u$ is a
monomorphism, we see that the $\varphi_{x}$ agree on the
$\fprod{X'_x}{X'_y}{X}$ $(x, y \in X)$ and therefore come from a
uniquely determined $\varphi \in F (X)$; one then has $u(\varphi) =
\psi$, since the inverse images on the $X'_{x}$
agree. \cqfd.

\begin{enonce*}[remark]{Scholium 3.10}\etiquette[3.10]{VIII.sch3.10}
Theorem \hyperref[VIII.thm3.5]{3.5} implies that the collection of
``fiber functors''
for the étale topology enjoys the same essential properties as
in the usual theory of sheaves on a topological
space. We shall use \hyperref[VIII.thm3.5]{3.5}
and its corollaries very frequently,
and for this reason we \sisi{}{shall} generally dispense with referring to them explicitly.
\end{enonce*}

\setcounter{subsection}{10}
\subsection{}\etiquette[3.11]{VIII.subsec3.11}

For every $x \in X$ we shall also, by abuse of language, denote by
$x$ the $X$-scheme $\Spec k(x)$, for the usual structural morphism
$$
i_x : x = \Spec k(x) \to X{\quoi}.
$$
It gives rise to a canonical functor
$$
i_x^{\ast} : \widetilde{X_{\text{ét}}} \to \widetilde{x_{\text{ét}}}{\quoi}.
$$
If $F$ is a sheaf on $X$, we shall put $F_x = i^{\ast}_x
(F)$; thus it is a sheaf on $x = \Spec k(x)$ (and as such
may also be identified with an\pageoriginale étale scheme over $x$, by
virtue of \hyperref[VIII.prop2.1]{2.1}). It depends functorially on $F$, and
the functor $i^{\ast}_x : F \mapsto F_x$ still commutes with
arbitrary $\varinjlim$ and finite $\varprojlim$, by virtue of
\hyperrefext[VII][VII.subsec1.4]{1.4}.

If $\bar{x} = \Spec \bar{k(x)}$, the fiber functor $F \mapsto
F_{\bar{x}}$ is canonically isomorphic to the composite of the functor $F
\mapsto F_x$ and the functor $j$ of \sisi{\hyperref[VIII.thm1.1]{1.1}}{\hyperref[VIII.prop2.1]{2.1}} (taking into account that the latter is
isomorphic to the fiber functor relative to the geometric point $\Spec
\bar{k}$ of $\Spec k$). It follows immediately from this and from
\hyperref[VIII.thm3.5]{3.5} that the
system of functors $(F \mapsto F_x)$ $(x \in X)$ is likewise
conservative. Moreover, if $k(x)$ is separably closed, \ie $x =
\bar{x}$, the notations $F_x$ and $F_{\bar{x}}$ are, strictly speaking,
in conflict (the first denoting an object of
$\widetilde{x_{\text{ét}}}$, the second a set), but by virtue of
\sisi{\hyperref[VIII.thm1.1]{1.1}}{\hyperref[VIII.prop2.1]{2.1}} this conflict is inessential.

\subsection{}\etiquette[3.12]{VIII.subsec3.12}

The preceding remarks show that if $\bar{x} = \Spec
\bar{k(x)}$, then, for every sheaf $F$ on $X$, the group $\pi_x = \Gal
(\bar{k(x)}/ k(x))$ acts naturally (on the left) on
the fiber set $F_{\bar{x}}$, so that $F \mapsto F_{\bar{x}}$
may in fact be regarded as a functor
$$
\widetilde{X_{\text{ét}}} \to \Bb_{\pi_x}
$$
whose knowledge is essentially equivalent (thanks to
\sisi{\hyperref[VIII.thm1.1]{1.1}}{\hyperref[VIII.prop2.1]{2.1}}) to that of the
functor $F \mapsto F_x$.

This shows, in an important special case, that contrary to what
happens for sheaves on topological spaces, ``fiber
functors'' generally admit nontrivial automorphisms.
In general, we shall determine in \hyperref[VIII.thm7.9]{7.9} all the
homomorphisms from one fiber functor to another.

\begin{enonce*}[remark]{Remark 3.13}\etiquette[3.13]{VIII.rem3.13}
Here is\pageoriginale a sometimes useful generalization
of \hyperref[VIII.thm3.5]{3.5} b). Consider a subset $E$ of $X$ such that,
for every étale morphism $f : X' \to X$, $E' = f^{-1} (E)$ is ``very dense'' in $X'$ (EGA IV 10.1.3), \ie for every $f$
as above, every open subset $U$ of $X'$ which contains $E'$ is equal
to $X'$. Then the fiber functors $F_{\bar{x}}$ on
$\widetilde{X_{\text{\rm ét}}}$, for $x \in E$, already form a
``conservative family'' (and consequently Corollaries
\hyperref[VIII.cor3.6]{3.6}, \hyperref[VIII.cor3.7]{3.7}, and \hyperref[VIII.cor3.8]{3.8} remain
valid when one restricts oneself to choosing $x \in E$ in them). This follows
immediately by repeating the proof given for
\hyperref[VIII.thm3.5]{3.5} b). The preceding result applies in particular
in the following situations:
\begin{enumerate}
\renewcommand{\theenumi}{\alph{enumi}}
\renewcommand{\labelenumi}{\theenumi)}
\item $X$ is a Jacobson scheme (EGA IV 10.4.1.), for example
  locally of finite type over a field, or over $\Spec(𝒵)$, and $E$
  is the set of closed points of $X$.

\item $X$ is locally of finite type over a scheme $S$, and $E$ is
  the set of points of $X$ which are closed in their fiber.

\item $X$ is locally noetherian, and $E$ is the set of $x \in
  X$ such that $\bar{x}$ is a finite set (\ie a
  semilocal scheme of dimension $\leq 1$) \cf \EGA IV 10.5.3 and 10.5.5.
\end{enumerate}
\end{enonce*}

% END INCLUDED SGA4_English_current_buildable_source_20260906/source/08/drafts/SGA4_VIII_LINES_43_576_English.texfrag


% BEGIN INCLUDED SGA4_English_current_buildable_source_20260906/source/08/drafts/SGA4_VIII_LINES_577_804_English.texfrag
\section{Strictly local rings and schemes}\etiquette[4]{VIII.sec4}

\subsection{}\etiquette[4.1]{VIII.subsec4.1}

Recall that, following Azumaya, a local ring $A$ is said to be
\emph{henselian} if it satisfies the following equivalent conditions:
\begin{enumerate}
\renewcommand{\theenumi}{\roman{enumi}}
\renewcommand{\labelenumi}{(\theenumi)}
\item Every finite $A$-algebra $B$ is a direct product
$\prod_i B_i$ of local rings (the $B_i$ are then necessarily
isomorphic to the $B_{\mathfrak m_i}$, where\pageoriginale
$\mathfrak m_i$ runs through the maximal ideals of $B$).

\item Every $A$-algebra $B$ which is the localization at a prime ideal
$\mathfrak p$ lying over the maximal ideal $\mathfrak m$ of $A$ of
an $A$-algebra $C$ of finite type, and such that
$B/\mathfrak mB$ is finite over $k=A/\mathfrak m$, is finite over
$A$.

\item The same as (ii), but with $C$ assumed étale over $A$ at
$\mathfrak p$ and $k\to B/\mathfrak mB$ an isomorphism.

\item \lq\lq Hensel's lemma'' in its classical form holds for $A$, that is,
for every monic polynomial $F\in A[T]$, writing
$\overset{\bigdot}{F}\in k[T]$ for the corresponding reduced
polynomial, and for every factorization
$$
\overset{\bigdot}{F}
=\overset{\bigdot}{F}_1\overset{\bigdot}{F}_2
$$
of $\overset{\bigdot}{F}$ as a product of two relatively prime monic
polynomials, $\overset{\bigdot}{F}_1$ and
$\overset{\bigdot}{F}_2$ lift (necessarily uniquely) to
$F_1,F_2\in A[T]$ such that $F=F_1F_2$.

For the proof of these equivalences and the general properties of
henselian rings, see \EGA IV 18.5.11. Recall that if $A$ is any local
ring, one proves (Nagata) that there exists a local homomorphism
$$
A\to A^h
$$
with $A^h$ \emph{henselian}, which is \emph{universal} for local
homomorphisms from $A$ to local
\sisi{Noetherian}{henselian} rings. The ring $A^h$ is called the
henselization of $A$. The construction given in \EGA IV 18.6
(different from Nagata's) consists in setting
$$
A^h=\varinjlim_i A_i{\quoi},
$$
where $A_i$ runs through the $A$-algebras essentially of finite type
and essentially étale\pageoriginale over $A$ (that is, localizations
of an étale $A$-algebra) such that $A\to A_i$ is a local
homomorphism and the residue extension $k(A_i)/k(A)$ is trivial. This
construction shows in particular that $A^h$ is flat over $A$, and
that if $A$ is Noetherian, then so is its henselization.
\end{enumerate}

\begin{enonce*}{Definition 4.2}\etiquette[4.2]{VIII.def4.2}
A local ring $A$ is said to be strictly local \emph{if it satisfies
one of the following equivalent conditions}:
\begin{enumerate}
\renewcommand{\theenumi}{\roman{enumi}}
\renewcommand{\labelenumi}{(\theenumi)}
\item $A$ is henselian (\hyperref[VIII.subsec4.1]{4.1}) and its
residue field is separably closed.

\item Every local homomorphism $A\to B$, where $B$ is a localization
of an étale algebra over $A$, is an isomorphism.
\end{enumerate}

A scheme is called a \emph{strictly local scheme} if it is isomorphic
to the spectrum of a strictly local ring.
\end{enonce*}

The equivalence of (i) and (ii) follows at once from
\hyperref[VIII.subsec4.1]{4.1}. Note that, in geometric form,
\hyperref[VIII.def4.2]{4.2} (ii) (respectively
\hyperref[VIII.subsec4.1]{4.1} (iii)) says that for every étale
$Y$-scheme $X$ and every point $x$ of the fiber $X_y$ (respectively
every point of $X_y$ \emph{rational over} $k(y)$), where $y$ is the
closed point of $Y$, there exists a section of $X$ over $Y$ passing
through $x$ (and this section is necessarily unique).

\begin{enonce*}{Definition 4.3}\etiquette[4.3]{VIII.def4.3}
Let $X$ be a scheme, let $\Oo_{X_{\text{\rm ét}}}$ be the sheaf on
$X_{\text{\rm ét}}$ defined by
$$
\Oo_{X_{\text{\rm ét}}}(X')=\Gamma(X',\Oo_{X'})
$$
(compare \emph{\hyperrefext[VII][VII.sec2]{2}. c)}), and let
$u:\xi=\Spec(\Omega)\to X$ be a geometric point of $X$. The
\emph{strictly local ring of $X$ at $\xi$}, denoted
$\Oo_{{X,}_{\xi}}$ or simply\pageoriginale $\Oo_\xi$, is the
fiber of the sheaf $\Oo_{X_{\text{\rm ét}}}$ at the point $\xi$.
Its spectrum is called the \emph{strict localization of $X$ at
$\xi$}.
\end{enonce*}

By \hyperref[VIII.prop3.9]{3.9}, the fiber $\Oo_{{X,}_{\xi}}$ has
the description
\begin{equation*}
\Oo_{{X,}_{\xi}}=\varinjlim\Gamma(X',\Oo_{X'}){\quoi},
\tag{$\ast$}
\end{equation*}
where $X'$ runs through the opposite filtered category of étale
$X$-schemes pointed by $\xi$, that is, equipped with an $X$-morphism
$\xi\to X'$. By transitivity of inductive limits, one may replace
$\Gamma(X',\Oo_{X'})$ in the right-hand side by $\Oo_{X',x'}$,
where $x'$ is the image of $\xi$ in $X'$, and one obtains
\begin{equation*}
\Oo_{{X,}_{\xi}}=\varinjlim_i A_i
\tag{\ast\ast}
\end{equation*}
where $A_i$ runs through the algebras essentially of finite type and
étale over $A=\Oo_{X,x}$, equipped with a $k$-homomorphism
$k(A_i)\to\Omega$. (Here $x$ denotes the image of $\xi$ in $X$.)
The transition homomorphisms between the $A_i$ are, of course, the
local homomorphisms of $A$-algebras $A_i\to A_j$ that make the
corresponding triangle commute:
\[
\xymatrix@R=1.25cm@C=1.75cm{
k(A_i)  \ar[r] \ar[rd]& k(A_j)\ar[d]    \\
                      &  \Omega.
}
\]
This implies that if there exists an admissible homomorphism from
$A_i$ to $A_j$, it is unique. Thus the inductive limit
$({\ast}{\ast})$ is a limit over an upward-filtered ordered index
set.

\setcounter{subsection}{3}
\subsection{}\etiquette[4.4]{VIII.subsec4.4}

The description\pageoriginale $({\ast}{\ast})$ shows that the
strictly local ring of $X$ at $\xi$ depends essentially only on the
usual local ring $A=\Oo_{X,x}$ ($x=u(\xi)$) and on the extension
$\Omega$ of $k=k(A)$. It is also called the \emph{strict
henselization} of $A$ (relative to the given separably closed
extension $\Omega$ of $k$), and is denoted $A^{hs}$. We refer to
\EGA IV 18.8 for the general properties of $A^{hs}$. We merely point
out that the construction again shows that $A^{hs}$ is flat over
$A$, and that it is Noetherian if $A$ is. Moreover, one proves easily
(\loccit) that the local homomorphism
$$
A\to A^{hs}
$$
together with the $k$-homomorphism
$$
k(A^{hs})\to\Omega
$$
solves the universal problem, relative to the data of $A$ and the
extension $\Omega$ of $k$, of finding local homomorphisms
$$
A\to A'{\quoi},
$$
where $A'$ is \emph{strictly local}, equipped with a
$k$-homomorphism
$$
k(A')\to\Omega{\quoi}.
$$
In particular, $A^{hs}$ \emph{is indeed a strictly local ring}
(which justifies the terminology of \hyperref[VIII.def4.3]{4.3}).

\subsection{}\etiquette[4.5]{VIII.subsec4.5}

Keep the notation of \hyperref[VIII.def4.3]{4.3}, and choose an
affine open neighborhood $U$ of $x$. Note that in
\hyperref[VIII.prop3.9]{3.9} one may plainly restrict to those $X'$
which are affine over $U$. These then form a pseudo-filtered inverse
system of affine schemes, so that we are\pageoriginale in the general
situation of \hyperrefext[VII][VII.sec5]{5} (see VI 5.12 b).
By the expression $(\ast)$ for $\Oo_{X,\xi}$, there is an
$X$-isomorphism
\begin{equation*}
\varprojlim_{\substack{X'\ \mbox{\rm étale}\\ \mbox{affine over }U,\\
X'_{\xi}\mbox{\rm -pointed}}}X'\simeq\Spec\Oo_{X,\xi}
\tag{\ast\ast\ast}
\end{equation*}
which makes precise the intuitive geometric meaning of
$\Spec(\Oo_{X,\xi})$ as the \lq\lq\emph{limit of the étale
neighborhoods of $\xi$}''.

\begin{enonce*}{Proposition 4.6}\etiquette[4.6]{VIII.prop4.6}
Let $X$ be a strictly local scheme, and let $x$ be its closed point,
which is therefore a geometric point of $X$. Then the functor
$F\mapsto\Gamma(X,F):\widetilde{X_{\text{\rm ét}}}\to(\Ens)$
is canonically isomorphic to the fiber functor $F\mapsto F_x$, and
hence commutes with arbitrary $\varinjlim$\footnote{not necessarily
filtered!}.
\end{enonce*}

Indeed, by Definition \hyperref[VIII.def4.2]{4.2} (ii), the final
object $X$ of the category of étale $X$-schemes pointed by $x$
considered in \hyperref[VIII.prop3.9]{3.9} dominates every other
object; hence $\varinjlim F(X')$ is canonically isomorphic to $F(X)$,
as required.

In particular, the functor induced by $\Gamma$ on the category of
abelian sheaves is exact, and therefore:

\begin{enonce*}{Corollary 4.7}\etiquette[4.7]{VIII.cor4.7}
Under the hypotheses of \hyperref[VIII.prop4.6]{4.6}, for every
abelian sheaf $F$ on $X$ one has
$$
H^q(X,F)=0\quad\mbox{\rm for}\quad q\neq0{\quoi}.
$$
\end{enonce*}

\begin{enonce*}{Corollary 4.8}\etiquette[4.8]{VIII.cor4.8}
Let $X$ be a scheme, let $\xi$ be a geometric point of $X$, let
$\bar X=\Spec\Oo_{X,\xi}$ be the corresponding strict
localization, let $F$ be a sheaf on $X$, and let $\bar F$ be its
inverse image on $\bar X$. Then there is a functorial isomorphism
$$
F_\xi\simeq\Gamma(\bar X,\bar F){\quoi}.
$$
\end{enonce*}

Indeed,\pageoriginale $\xi$ may also be regarded as a geometric point
of $\bar X$, and it suffices to apply
\hyperref[VIII.prop4.6]{4.6} and the transitivity of fibers
(\hyperref[VIII.subsec3.4]{3.4}).


% END INCLUDED SGA4_English_current_buildable_source_20260906/source/08/drafts/SGA4_VIII_LINES_577_804_English.texfrag


% BEGIN INCLUDED SGA4_English_current_buildable_source_20260906/source/08/drafts/SGA4_VIII_LINES_805_1011_English.texfrag
\section{Application to the computation of the fibers of the
$R^qf_*$}\etiquette[5]{VIII.sec5}

\subsection{}\etiquette[5.1]{VIII.subsec5.1}

Let
$$
f:X\to Y
$$
be a morphism of schemes and let $F$ be an abelian sheaf on $X$. We
wish to determine the $R^qf_*(F)$. By
\hyperref[VIII.thm3.5]{3.5}, this amounts, up to very little, to
knowing the geometric fibers $R^qf_*(F)_{\bar y}$ for $y\in Y$. Now
(\hyperrefext[V][V.prop5.1]{5.1}) $R^qf_*(F)$ is the sheaf associated
with the presheaf
$$
\Hh^q:Y'\longmapsto H^q(\fprod{X}{Y'}{Y},F)
$$
on $Y$, and hence (\hyperref[VIII.prop3.9]{3.9})
$$
R^qf_*(F)_{\bar y}=\varinjlim_{Y'}\Hh^q(Y'){\quoi},
$$
where the inductive limit is taken over the opposite filtered
category of the étale $Y$-schemes pointed by $\bar y$. After choosing
an affine open neighborhood $U$ of $y$, one may restrict the limit on
the right to the cofinal category of those $Y'$ which are affine and
lie over $U$. Thus one obtains
\begin{equation*}
R^qf_*(F)_{\bar y}\simeq
\varinjlim H^q(\fprod{X}{Y'}{Y},F){\quoi},
\tag{5.1.1}
\end{equation*}\etiquette[5.1.1]{VIII.eq5.1.1}
where $Y'$ varies in the preceding category.

Put
$$
\bar Y=\Spec(\Oo_{Y,\bar y})=\varprojlim Y'
$$
(see \hyperref[VIII.subsec4.5]{4.5}), and
$$
\bar X=\fprod{X}{\bar Y}{Y}{\quoi}.
$$
Observe that\pageoriginale in the inverse system
$\fprod{X}{Y'}{Y}$ (deduced from that of the $Y'$ by the base change
$X\to Y$), the transition morphisms are affine. We are therefore in
the general situation of \hyperrefext[VII][VII.subsec5.1]{5.1}, and
the construction of $\varprojlim$ in \loccit immediately gives a
canonical isomorphism
$$
\bar X=\varprojlim\fprod{X}{Y'}{Y}{\quoi}.
$$
Let $\bar F$ denote the sheaf on $\bar X$ that is the inverse image of
$F$. We then obtain a canonical homomorphism
$$
\varinjlim_{Y'}H^q(\fprod{X}{Y'}{Y},F)
\longrightarrow H^q(\bar X,\bar F),
$$
and hence, by comparison with \eqhyperref[VIII.eq5.1.1]{5.1.1}, a
canonical homomorphism
\begin{equation*}
R^qf_*(F)_{\bar y}\longrightarrow H^q(\bar X,\bar F)
\tag{5.1.2}{\quoi},
\end{equation*}\etiquette[5.1.2]{VIII.eq5.1.2}
which is plainly functorial in $F$.

Suppose now that $f:X\to Y$ is quasi-compact and quasi-separated.
Then the same is true of $\fprod{X}{Y'}{Y}\to Y'$, and since $Y'$ is
affine, the schemes $\fprod{X}{Y'}{Y}$ are quasi-compact and
quasi-separated. Applying \eqhyperref[VIII.eq5.1.1]{5.1.1} and the
limit-passage theorem \hyperrefext[VII][VII.cor5.8]{5.8}, one obtains:

\begin{enonce*}{Theorem 5.2}\etiquette[5.2]{VIII.thm5.2}
Let $f:X\to Y$ be a quasi-compact and quasi-separated morphism of
schemes, let $F$ be an abelian sheaf on $X$, let $y$ be a point of
$Y$, and let $\bar y$ be the geometric point over $y$ associated with
a separable closure $k(\bar y)$ of $k(y)$. Let
$\bar Y=\Spec(\Oo_{Y,\bar y})$ be the corresponding strict
localization, let $\bar X=\fprod{X}{\bar Y}{Y}$, and let $\bar F$ be
the inverse image of $F$ on $\bar X$. Then the canonical
homomorphism \eqhyperref[VIII.eq5.1.2]{5.1.2} is an isomorphism:
$$
R^qf_*(F)_{\bar y}\simeq H^q(\bar X,\bar F){\quoi}.
$$
\end{enonce*}

This statement practically reduces the determination of the fibers
of a sheaf\pageoriginale $R^qf_*(F)$ to the determination of the
cohomology of a scheme over a strictly local scheme, and it will be
used constantly below. Technically, \hyperref[VIII.thm5.2]{5.2}
explains the important role of henselian and strictly local rings in
the study of étale cohomology.

\begin{enonce*}[remark]{Remark 5.3}\etiquette[5.3]{VIII.rem5.3}
The statement remains valid for $R^0f_*(F)=f_*(F)$, for an arbitrary
sheaf of sets. Moreover, \hyperref[VIII.thm5.2]{5.2} also remains
valid for $R^1f_*(F)$ when $F$ is a sheaf of groups (not necessarily
commutative), with the usual definition of $R^1f_*(F)$ (see Giraud's
thesis).
\end{enonce*}

\setcounter{subsection}{3}
\subsection{}\etiquette[5.4]{VIII.subsec5.4}

Suppose now that $f$ is a \emph{finite} morphism. Then
$\bar f:\bar X\to\bar Y$ is finite, and since $\bar Y$ is strictly
local, it follows that $\bar X$ is a finite disjoint union of strictly
local schemes $\bar X_i$. Hence, using
\hyperref[VIII.cor4.7]{4.7}, one finds
\begin{equation*}
H^q(\bar X,\bar F)=0\quad\mbox{for}\quad q\neq0{\quoi}.
\tag{5.4.1}
\end{equation*}\etiquette[5.4.1]{VIII.eq5.4.1}
Furthermore, the components $\bar X_i$ of $\bar X$ correspond to the
points $\bar x_i$ of $\bar X$ above the closed point $\bar y$ of
$\bar Y$, that is, to the points of
$\bar X_{\bar y}=X_y\otimes_{k(y)}k(\bar y)$. These points may also
be regarded as geometric points of $X$, and $\bar X_i$ is then
precisely the strict localization of $X$ at $\bar x_i$. Thus
(\hyperref[VIII.cor4.8]{4.8})
$$
H^0(\bar X_i,\bar F)\simeq F_{\bar x_i}{\quoi},
$$
whence
\begin{equation*}
H^0(\bar X,\bar F)=\mathop{\prod}_iF_{\bar x_i}{\quoi}.
\tag{5.4.2}
\end{equation*}\etiquette[5.4.2]{VIII.eq5.4.2}
This is\pageoriginale an isomorphism functorial in the sheaf of sets
$F$ (there is no need here to restrict to abelian sheaves). Taking
\hyperref[VIII.thm5.2]{5.2} and \hyperref[VIII.rem5.3]{5.3} into
account, formulas \eqhyperref[VIII.eq5.4.1]{5.4.1} and
\eqhyperref[VIII.eq5.4.2]{5.4.2} give:

\begin{enonce*}{Proposition 5.5}\etiquette[5.5]{VIII.prop5.5}
Let $f:X\to Y$ be a finite morphism of schemes, and let $y$ be a
point of $Y$. Then for every sheaf $F$ on $X$, there is an
isomorphism, functorial in $F$,
$$
f_*(F)_{\bar y}\simeq
\prod_{\bar x\in X_y\otimes_{k(y)}k(\bar y)}F_{\bar x}{\quoi}.
$$
Consequently, the formation of $f_*(F)$ commutes with every base
change $Y'\to Y$; and if $F$ is an abelian sheaf, then
$$
R^qf_*(F)=0\quad\mbox{if}\quad q>0{\quoi}.
$$
\end{enonce*}

Note that the first formula of \hyperref[VIII.prop5.5]{5.5} is in
fact independent of \hyperref[VIII.thm5.2]{5.2} and of the general
limit-passage theorem \hyperrefext[VII][VII.thm5.7]{5.7}; together
with \hyperref[VIII.thm3.5]{3.5}, it implies that
$F\mapsto f_*(F)$ is an \emph{exact functor} on the category of
abelian sheaves, and hence again that $R^qf_*(F)=0$ for $q\neq0$.
Here is a slight variant of \hyperref[VIII.prop5.5]{5.5}:

\begin{enonce*}{Corollary 5.6}\etiquette[5.6]{VIII.cor5.6}
Let $f:X\to Y$ be an integral morphism. Then for every abelian sheaf
$F$ on $X$ one has
$$
R^qf_*(F)=0\quad\mbox{for}\quad q\neq0{\quoi}.
$$
Moreover, the functor $f_*$ on sheaves of sets commutes with every
base change $Y'\to Y$.
\end{enonce*}

Indeed, by \hyperref[VIII.thm5.2]{5.2} it is enough to prove that when
$Y$ is strictly local, $H^q(X,F)=0$ for every abelian sheaf $F$ on
$X$. We then have $Y=\Spec A$, $X=\Spec B$, where $B$ is an integral
$A$-algebra; writing\pageoriginale $B=\varinjlim B_i$, where $B_i$
runs through the finite-type subalgebras of $B$ (which are in fact
finite over $A$), one has $X=\varprojlim_iX_i$, where
$X_i=\Spec B_i$. By \hyperrefext[VII][VII.cor5.13]{5.13}, it is
enough to prove that $R^qf_*(F_i)=0$ for $q\neq0$, where $F_i$ is an
abelian sheaf on $X_i$; this follows from
\hyperref[VIII.prop5.5]{5.5}. The last assertion of
\hyperref[VIII.cor5.6]{5.6} is proved by the same reduction to
\hyperref[VIII.prop5.5]{5.5}.

\begin{enonce*}{Corollary 5.7}\etiquette[5.7]{VIII.cor5.7}
Let $f:X\to Y$ be an immersion. Then for every sheaf $F$ on $X$, the
canonical morphism
$$
f^*f_*(F)\to F
$$
is an isomorphism.
\end{enonce*}

Since $f$ factors as the composite of an open immersion and a closed
immersion, and \hyperref[VIII.cor5.7]{5.7} is trivial for an open
immersion (IV no.~3), it remains to consider the case where $f$ is a
closed immersion. We must prove that for every $x\in X$, the
corresponding homomorphism
$$
(f^*f_*(F))_{\bar x}\to F_{\bar x}
$$
is bijective. By transitivity of fibers
(\hyperref[VIII.subsec3.4]{3.4}), the left-hand side is precisely the
fiber $f_*(F)_{\bar x}$; it is therefore enough to verify that the
canonical homomorphism
$$
f_*(F)_{\bar x}\to F_{\bar x}
$$
is bijective, which follows at once from
\hyperref[VIII.prop5.5]{5.5}.

\begin{enonce*}[remark]{Remark 5.8}\etiquette[5.8]{VIII.rem5.8}
Using \hyperref[VIII.prop5.5]{5.5} and arguing again as in
\hyperref[VIII.cor5.6]{5.6}, one easily proves that if
$f:X\to Y$ is an integral morphism and $F$ is a sheaf of groups on
$X$ (not necessarily commutative), then $R^1f_*(F)$ is the terminal
sheaf on $Y$ (compare Remark \hyperref[VIII.rem5.3]{5.3}).
\end{enonce*}


% END INCLUDED SGA4_English_current_buildable_source_20260906/source/08/drafts/SGA4_VIII_LINES_805_1011_English.texfrag


% BEGIN INCLUDED SGA4_English_current_buildable_source_20260906/source/08/drafts/SGA4_VIII_LINES_1012_1556_English.texfrag
\section{Supports}\pageoriginale\etiquette[6]{VIII.sec6}

Let $U$ be a Zariski open subset of the scheme $X$. Then $U \in \Ob
X_{\text{\rm ét}}$,
and in fact $U$ is a subobject of the final object $X$ of $X_{\text{\rm
    ét}}$,
hence defines a subobject $\widetilde{U}$ of the final object of
$\widetilde{X_{\text{\rm ét}}}$, \ie an ``open subset'' of the étale
topos $X_{\text{\rm ét}}$
of $X$ (IV \hyperref[IV.def8.3]{8.3}).

\begin{enonce*}{Proposition 6.1}\etiquette[6.1]{VIII.prop6.1}
The preceding map $U \mapsto \overset{\sim}{U}$ is an isomorphism
from the ordered set of open subsets (in the Zariski sense) of $X$
onto the set of open subsets of the étale topos
$\widetilde{X_{\text{\rm ét}}}$.
\end{enonce*}

Since $X_{\text{\rm ét}} \to \widetilde{X_{\text{\rm ét}}}$ is fully
faithful, it follows immediately that the map $U \mapsto \widetilde{U}$
preserves the order structures, \ie $(U \subset V) \Leftrightarrow
(\widetilde{U} \subset \widetilde{V})$; in particular, the preceding
map is injective. To prove that it is surjective, consider a subsheaf
$F$ of the final sheaf $\widetilde{X}$, and consider the objects of
$X_{\text{ét}/ F}$, \ie the schemes $X'$ étale over $X$ such that
$F(X') \neq ∅$. Since $X' \to X$ is étale, it is an open map
(\EGA IV 2.4.6); in particular, its image (in the set-theoretic sense)
$\text{Im}(X')$ is open. Let $U$ be the open subset which is the union
of the $\text{Im}(X')$ $(X'\in \Ob X_{\text{ét} /  F})$. Since the
family of the $X'\to U$ $(X'\in \Ob X_{\text{ét}/F})$ is
surjective, hence covering, we conclude that $U \in \Ob
X_{\text{ét}/F}$, hence
$ X_{\text{ét}/F} =  X_{\text{ét}/U}$, hence $F = U$, \cqfd.

In view of \hyperref[VIII.prop6.1]{6.1}, we may therefore speak
unambiguously of an ``open subset'' of $X$, without specifying whether
we mean this notion in the usual Zariski sense or in the sense of the
étale topology.

\begin{enonce*}{Corollary 6.2}\etiquette[6.2]{VIII.cor6.2}
Let $U$ be an open subset of $X$, and let $j : U \to X$ be the
canonical immersion. Then the functor
$$
j^{\ast} :  \widetilde{X_{\text{\rm ét}}} \to \widetilde{U_{\text{\rm ét}}}
$$
induces\pageoriginale an equivalence of categories
$$
\widetilde{X_{\text{\rm ét}}} / \widetilde{U} \to
\widetilde{U_{\text{\rm ét}}}{\quoi}.
$$
\end{enonce*}

Indeed, one checks immediately that for every site $C$ in which finite
$\varprojlim$ exist, and for every subobject $U$ of the final object
$e$ of $C$, if one considers the functor $j : C \to C_{/U}$ defined by
$j(S) = S \times U$, then $j$ is a morphism of sites and $j^{\ast} :
\widetilde{C} \to \widetilde{C/U}$ induces an equivalence
$\widetilde{C}/\widetilde{U} \to \widetilde{C / U}$. It then suffices
to combine this general fact with the fact that $U_{\text{ét}}$ is
canonically isomorphic to $X_{\text{ét}/U}$.

In pictorial terms, one can express \hyperref[VIII.cor6.2]{6.2} by
saying that the operations of ``restriction to an open subset,'' in
the usual sense for schemes on the one hand and in the sense of topoi
on the other, are compatible. Here is an analogous compatibility for
the operations of ``restriction to a closed subset'':

\begin{enonce*}{Theorem 6.3}\etiquette[6.3]{VIII.thm6.3}
Let $X$ be a scheme, let $Y$ be a closed subscheme of $X$, let $U
= X - Y$ be endowed with the induced structure, and let $i : Y \to X$
and $j : U \to X$ be the canonical immersions. Then the functor
$$
i_{\ast} : \widetilde{Y_{\text{\rm ét}}} \to \widetilde{X_{\text{\rm ét}}}
$$
is fully faithful, and if $F \in \Ob \widetilde{X_{\text{\rm ét}}}$,
then $F$ is isomorphic to a sheaf of the form $i_{\ast} (G)$ if and
only if $j^{\ast}(F)$ is isomorphic to the final sheaf
$\widetilde{U}$ on $U$.
\end{enonce*}

Proof. Since $i_{\ast}$ and $i^{\ast}$ are adjoint to one another, the
fact that $i_{\ast}$ is fully faithful is also equivalent to the fact
that the functorial homomorphism
$$
i^{\ast} (i_{\ast}(G)) \to G
$$
is an\pageoriginale isomorphism, which is precisely
\hyperref[VIII.cor5.7]{5.7}. On the other hand, if $G \in \Ob
\widetilde{Y_{\text{ét}}}$, then one checks trivially, using
\hyperref[VIII.cor6.2]{6.2}, that $j^{\ast} (i_{\ast}(G))$ is the final
sheaf on $U$. Conversely, if $F \in \Ob
\widetilde{X_{\text{ét}}}$ is such that $j^{\ast}(F)$ is the final
sheaf, let us prove that $F$ is isomorphic to a sheaf of the form
$i_{\ast}(G)$, or, what amounts to the same thing, that the canonical
homomorphism
$$
F \to i_{\ast} i^{\ast} F
$$
is an isomorphism. It is enough once again to check that for every $x
\in X$, the homomorphism induced on the geometric fibers at $\bar{x}$
is an isomorphism. When $x \in U$, this is precisely the hypothesis
made on $F$. When $x \in Y$, by transitivity of fibers we are reduced
to checking that the homomorphism on the fibers at $\bar{x}$ induced
by
$$
i^{\ast} (F) \to i^{\ast} (i_{\ast} i^{\ast} (F))
$$
is an isomorphism; but the preceding homomorphism is an isomorphism by
\hyperref[VIII.cor5.7]{5.7}, applied to $F = i^{\ast} F$ and to $i$.
This completes the proof of \hyperref[VIII.thm6.3]{6.3}.

\begin{enonce*}{Corollary 6.4}\etiquette[6.4]{VIII.cor6.4}
The functor $i_{\ast}$ induces an equivalence from the category of
abelian sheaves on $Y$ to the category of abelian sheaves on $X$ whose
restriction to $U = X - Y$ is zero.
\end{enonce*}

In accordance with common usage, we shall therefore often identify an
abelian sheaf on $Y$ with an abelian sheaf on $X$ which is zero on $U$.

\setcounter{subsection}{4}
\subsection{}\etiquette[6.5]{VIII.subsec6.5}

By virtue of \hyperref[VIII.prop6.1]{6.1}, we thus know (if $X$ is a
scheme) how to interpret the ``open subsets'' of the topos $ℰ =
\widetilde{X_{\text{ét}}}$ as open subsets $U$ of $X$ in the usual
sense, and by virtue of \hyperref[VIII.thm6.3]{6.3}, under this
identification the ``residual topos''\pageoriginale $ℰ_\complement U$
of IV 3 is canonically equivalent to the topos
$\widetilde{Y_{\text{ét}}}$, where $Y = X - U$ is endowed with any
induced structure making $Y$ a scheme. We may consequently apply to
this situation the results of IV 3, in particular the description IV
\hyperref[IV.theo9.5.4]{3.3} of sheaves $F$ on $X$ in terms of triples
$(F', F'', u)$, where $F'$ is a sheaf on $Y$, $F''$ is a sheaf on $U$,
and $u$ is a homomorphism from $F'$ to $i^{\ast} j_{\ast} (F'')$. In
what follows we shall freely use the notation $i^{!}, j_{!}$ of IV.3,
which denotes functors
\begin{align*}
j_{!} : \widetilde{(U_{\text{ét}})}_{\ab} &\to
\widetilde{(X_{\text{ét}})}_{\ab}{\quoi},\\
i^{!} : \widetilde{(X_{\text{ét}})}_{\ab} &\to
\widetilde{(Y_{\text{ét}})}_{\ab}{\quoi},
\end{align*}
where the subscript ``$\ab$'' denotes the category of abelian sheaves.
These functors give rise to the two exact sequences IV \hyperref[IV.prop14.6]{3.7}.

\subsection{}\etiquette[6.6]{VIII.subsec6.6}

In view of the preceding developments and of the general terminology
introduced in IV \hyperref[IV.subsec8.5]{8.5}, it is appropriate, for an abelian sheaf $F$ on
$X$, or a section $\varphi$ of such a sheaf, to introduce the notion of
the \emph{support of} $F$, \resp of $\varphi$, as the closed subset (in
the usual, \ie Zariski, sense) complementary to the largest open subset
on which the object in question vanishes.

In the present situation, one should also replace the general notation
\hyperrefext[V][V.prop4.3]{4.3} $H^q (ℰ_{\complement U} F)$,
$ℋ^q_{\complement U} (F)$ by the notation $H^{q}_Y (X, F)$,
$ℋ^{q}_Y (F)$, the first denoting an abelian group and the second an
abelian sheaf on $Y$ (or equivalently an abelian sheaf on $X$ which is
zero on $U$). Thus, one has $\mathcal H^q_Y(F)=R^q i^!(F)$, \etc

We refer to \hyperrefext[V][V.sec6]{6} for the general properties of
the preceding functors.

\section{Specialization morphisms of fiber
  functors}\pageoriginale
\etiquette[7]{VIII.sec7}

\subsection{}\etiquette[7.1]{VIII.subsec7.1}

We saw in ${\rm N^{\circ}}$ 4 how one associates, with every geometric
point $\xi$ of the scheme $X$, a strictly local $X$-scheme
$$
\bar{X}(\xi) = \Spec \Oo_{X, \xi}{\quoi},
$$
which in fact depends only on the geometric point above the image $x$
of $\xi$ defined by the separable closure $\bar{k(x)}$ of $k(x)$ in
$\Omega = k(\xi)$. In what follows we shall often restrict ourselves
to geometric points $\xi = \Spec \Omega$ which are \emph{algebraic
  separable} over $X$, \ie such that $\Omega$ is a separable closure
of $k(x)$, \ie $\xi = \bar{x}$. Let us call an $X$-scheme $Z$ a
\emph{strict localization of} $X$ if it is $X$-isomorphic to a scheme
of the form $\Spec \Oo_{X, \bar{x}}$; thus one sees that
$$
\xi \mapsto \Spec \sheaf{\Oo}_{X, \xi} = \bar{X}(\xi)
$$
is an equivalence from the category of geometric points
\emph{algebraic separable} over $X$ to the category of $X$-schemes
which are strict localizations of $X$, when only isomorphisms are taken
as morphisms in both categories.

The last assertion is no longer correct when all $X$-morphisms are
taken as morphisms, since in the second category there may be
$X$-morphisms which are not isomorphisms. We then make the following

\begin{enonce*}{Definition 7.2}\etiquette[7.2]{VIII.def7.2}
Let $\xi$, $\xi'$ be two geometric points of the scheme $X$. A
specialization arrow from $\xi'$ to $\xi$ is any $X$-morphism between
the corresponding strict localizations
$$
\bar{X} (\xi') \to \bar{X} (\xi).
$$
We\pageoriginale say that $\xi$ is a specialization of $\xi'$, or
that $\xi'$ is a generization of $\xi$, if there exists a
specialization arrow from $\xi'$ to $\xi$.
\end{enonce*}

Note that specialization arrows compose in the evident way, so that,
with specialization arrows as morphisms, \emph{the geometric points of
  $X$ form a category}, equivalent (as is the full subcategory formed
by the geometric points algebraic and separable over $X$) to the full
subcategory of $(\Sch)_{/X}$ formed by the strict localizations of $X$.

\begin{enonce*}{Lemma 7.3}\etiquette[7.3]{VIII.lem7.3}
Let $X$ be a scheme, let $Z$ be an $X$-scheme which is isomorphic to a
pseudo-filtered projective limit of $X$-schemes $X_{i}$ which are étale,
with affine transition morphisms
\emph{(\hyperrefext[VII][VII.subsec5.1]{5.1})}, let $\xi'$ be a
geometric point of $X$, and let $Z' = \bar{X} (\xi')$ be the
corresponding strict localization. Then:
\begin{enumerate}
\renewcommand{\theenumi}{\alph{enumi}}
\renewcommand{\labelenumi}{\theenumi)}
\item The restriction map
$$
\Hom_X (Z', Z) \to \Hom_X(\xi', Z)
$$
is bijective.

\item In order that the two sides be nonempty, it is necessary and
  sufficient that the image $x'$ of $\xi'$ in $X$ lie in the image of
  $Z$.

\item Let $T$ be a $Z$-scheme. In order that $T$ be a strict
  localization of $Z$, it is necessary and sufficient that it be a
  strict localization of $X$.
\end{enumerate}
\end{enonce*}

Proof.
\begin{enumerate}
\renewcommand{\theenumi}{\alph{enumi}}
\renewcommand{\labelenumi}{\theenumi)}
\item We are immediately reduced to the case where $Z$ is one of the
  $X_i$, \ie where $Z$ is étale over $X$, and by the base change
  $Z'\to X$ we may suppose that $Z'\xrightarrow{\sim} X$, whence the
  conclusion follows from \hyperref[VIII.def4.2]{4.2} (ii).

\item Note that if $x'$ is the image of a point $z$ of $Z$, then
  $k(z)$ is necessarily a separable algebraic extension of $k(x')$;
  hence there exists a $k(x')$-homomorphism from the latter into
  $k(\xi')$, whence the conclusion.

\item The proof\pageoriginale is an easy exercise left to the reader.
\end{enumerate}

From \hyperref[VIII.def7.2]{7.2} and \hyperref[VIII.lem7.3]{7.3} we
conclude:

\begin{enonce*}{Proposition 7.4}\etiquette[7.4]{VIII.prop7.4}
Let $\xi$, $\xi'$ be two geometric points of the scheme $X$. Then the
restriction map defines a bijection from the set
$\Hom_{X} (\bar{X}(\xi'), \bar{X}(\xi))$ of specialization arrows from
$\xi'$ to $\xi$ to the set of $X$-morphisms from $\xi'$ to
$\bar{X}(\xi)$.
\end{enonce*}

\begin{enonce*}{Corollary 7.5}\etiquette[7.5]{VIII.cor7.5}
In order that $\xi$ be a specialization of $\xi'$, it is necessary and
sufficient that the same hold for the images $x$, $x'$ of $\xi$,
$\xi'$ in $X$, \ie that $x$ belong to the closure of $\{x'\}$.
\end{enonce*}

Since $\Spec \Oo_{X, \xi} \to \Spec \Oo_{X, x}$ is faithfully flat
(\hyperref[VIII.subsec4.4]{4.4}), it is surjective, and it is therefore
enough to apply the second assertion of \hyperref[VIII.lem7.3]{7.3}.

\begin{enonce*}{Corollary 7.6}\etiquette[7.6]{VIII.cor7.6}
For every scheme $X$, let $\sheaf{\Pt}(X)$ be the category of geometric
points over $X$ (or equivalently of geometric points algebraic
separable over $X$), the morphisms being specialization arrows. Let
$\xi$ be a geometric point of a scheme $X$, and let $\bar{X}(\xi)$ be
the corresponding strict localization. Then there is an equivalence of
categories
$$
\sheaf{\Pt}(\bar{X}(\xi)) \xrightarrow{\sim} \sheaf{\Pt}(X)_{/\xi}{\quoi},
$$
obtained by associating with every geometric point $\xi'$ of
$\bar{X}(\xi)$ the corresponding geometric point over $X$, together
with the specialization morphism to $\xi$ deduced from the structural
morphism $\xi' \to \bar{X}(\xi)$ by \hyperref[VIII.prop7.4]{7.4}.
\end{enonce*}

In other words, \emph{the datum of a generization $\xi'$ of the
geometric point $\xi$ is essentially equivalent to the datum of a
geometric point of $\bar{X}(\xi)$}. Note moreover that, by
\hyperref[VIII.lem7.3]{7.3} c), the corresponding
homomorphism\pageoriginale $\bar{X}(\xi') \to \bar{X}(\xi)$ makes
$\bar{X}(\xi')$ the strict localization of $\bar{X}(\xi)$ relative to
$\xi'$.

\setcounter{subsection}{6}
\subsection{}\etiquette[7.7]{VIII.subsec7.7}

Now interpret, for every geometric point $\xi$ of $X$ and every sheaf
$F$ on $X$, the fiber $F_{\xi}$ as $\Gamma (\bar{X} (\xi),
\bar{F}(\xi))$, where $\bar{F}(\xi)$ is the inverse image of $F$ on
$\bar{X}(\xi)$ (\hyperref[VIII.cor4.8]{4.8}). Then one sees that every
specialization arrow
$$
u : \xi'\to \xi
$$
induces a homomorphism, functorial in $F$:
$$
u^{\ast} : F_{\xi} \to F_{\xi'}{\quoi},
$$
called the \emph{specialization homomorphism} associated with the
specialization arrow $u$. It is evident from the transitivity of
inverse images of sheaves that, for a composite specialization arrow,
one has:
$$
(wu)^{\ast} = u^{\ast} w^{\ast}{\quoi}.
$$

\subsection{}\etiquette[7.8]{VIII.subsec7.8}

% SGA-SRC-000768: the scan supplies no exact citation; reconstructed from the mathematics.
If $ℰ$ is a topos, recall from Exposé IV, 6.1 and 6.2,
that a ``fiber functor,'' or (by abuse of language) a ``point'' of the
topos $ℰ$, is any morphism from the ``final topos'' $(\Ens)$
(isomorphic to the category of sheaves on a space reduced to one
point!) to $ℰ$, \ie any functor
$$
\varphi : ℰ \to (\Ens)
$$
which commutes with arbitrary inductive \sisi{lim.}{limits}, and with
finite projective \sisi{lim.}{limits}. The set of fiber functors of
$ℰ$ should be regarded as the set of objects of a full subcategory of
$\Hhom(ℰ, (\Ens))$, called the \emph{category of fiber functors} of the
topos $ℰ$, whose opposite is called the \emph{category of points of
  $ℰ$}, and is denoted $\sheaf{\Pt}(ℰ)$, \cf (IV \hyperref[IV.def6.1]{6.1}).

When $ℰ$ is of the form $\widetilde{X_{\text{ét}}}$, where $X$ is
a scheme, we have\pageoriginale defined in
\hyperref[VIII.def3.3]{3.3} and \hyperref[VIII.subsec7.7]{7.7} a
functor
\begin{equation*}
\sheaf{\Pt} (X) \to \sheaf{\Pt} (\widetilde{X_{\text{ét}}}){\quoi},\tag{$\ast$}
\end{equation*}
where the first member is defined in \hyperref[VIII.cor7.6]{7.6}. With
this understood, we have the following

\begin{enonce*}{Theorem 7.9}\etiquette[7.9]{VIII.thm7.9}
Let $X$ be a scheme. The preceding functor $(\ast)$ is an equivalence
from the category of geometric points over $X$ (with specialization
arrows as morphisms) to the category of points of the étale topos
$\widetilde{X_{\text{\rm ét}}}$ (the \emph{opposite} of that of fiber
functors on $\widetilde{X_{\text{\rm ét}}}$).
\end{enonce*}

Since this theorem will not be used later in the seminar, we confine
ourselves here to a sketch of the proof, in which we shall allow
ourselves certain liberties with questions of universes.
\begin{enumerate}
\renewcommand{\theenumi}{\alph{enumi}}
\renewcommand{\labelenumi}{\theenumi)}
\item \emph{The functor under consideration is fully faithful}. For
  every $X \in \Ob X_{\text{ét}}$, let $\Xv :
  \widetilde{X_{\text{ét}}} \to (\Ens)$ be defined by
$$
\Xv (F) = \Hom(X, F) = F(X){\quoi}.
$$
Note that the fiber functor
$$
\mathcal{E}_{\xi} : F \mapsto F_{\xi}
$$
can be written
$$
\mathcal{E}_{\xi} \simeq \varinjlim_i X^{\vee}_i{\quoi},
$$
where the $X_i$ are affine schemes étale over $X$, indexed by a
certain filtered category (\cf \hyperref[VIII.def4.3]{4.3} and
\hyperref[VIII.subsec4.5]{4.5}). Hence, for every functor $\varphi :
\widetilde{X_{\text{ét}}} \to (\Ens)$, one has:
$$
\Hom(\mathcal{E}_{\xi}, \varphi) \simeq \varprojlim_{i}
\Hom(X^{\vee}_i, \varphi){\quoi},
$$
and on the other hand it is well known that there is a bifunctorial
isomorphism
$$
\Hom(X^{\vee}_i, \varphi) \simeq \varphi (X_i){\quoi}.
$$

When\pageoriginale $\varphi$ is of the form $\Ee_{\xi'}$, hence
$$
\varphi \simeq \varinjlim_j {X'}_j^{\vee}{\quoi},
$$
there is therefore a natural bijection
\begin{equation*}
\Hom (\mathcal{E}_{\xi}, \mathcal{E}_{\xi'}) \simeq
\varprojlim_{i} \varinjlim_{j} \Hom_X (X'_j, X_i){\quoi}.\tag{$\ast$}
\end{equation*}
On the other hand, one has (in the category of schemes)
$$
\bar{X}(\xi) = \varprojlim_{i} X_i,\quad \bar{X} (\xi') =
\varprojlim_{j} X'_j{\quoi},
$$
whence
$$
\Hom_{X} (\bar{X} (\xi'), \bar{X}(\xi)) \simeq \varprojlim_{i}
\Hom_X (\bar{X} (\xi'), X_i){\quoi},
$$
and on the other hand, since $X_i$ is locally of finite presentation
over $X$, one has
$$
\Hom_X (\bar{X} (\xi'), X_i) \simeq \varinjlim_{j}
\Hom_X (X'_j, X_i){\quoi},
$$
whence
\begin{equation*}
\Hom_X (\bar{X} (\xi'), \bar{X} (\xi) ) \simeq
\varprojlim_{i} \varinjlim_{j} \Hom_X (X'_j,
X_i){\quoi}.\tag{${\ast}{\ast}$}
\end{equation*}
Comparison of $(\ast)$ and $({\ast}{\ast})$ gives the desired
conclusion (subject to checking compatibilities, which is left to the
reader).

\item \emph{The functor under consideration is essentially
  surjective}.
\end{enumerate}

\setcounter{subsection}{9}
\subsubsection{}\etiquette[7.9.1]{VIII.subsubsec7.9.1}

First note that if $C$ is a site in which finite $\varprojlim$ are
representable, and $\widetilde{C} = ℰ$ is the corresponding topos, then
every fiber functor $\varphi$ on $ℰ$ can be represented as a ``filtered
inductive limit'' of functors of the form
$$
\check{Y} (F) = F(Y) = \Hom(\widetilde{Y}, F), Y \in \Ob C{\quoi},
$$
(where $\widetilde{Y}$ is the sheaf associated with $Y$), as follows
\footnote{This is a special case of IV \hyperref[IV.eq6.8.3]{6.8.3}.}. Consider the
category\pageoriginale $C_{/\varphi}$ of pairs $(Y, \xi)$, with
$Y \in \Ob C$, $\xi \in \varphi (\widetilde{Y})$ (the morphisms being
defined in the evident way), and note that there is an evident
homomorphism
\begin{equation*}
\varinjlim_{(Y, \xi)\in C^0_{/\varphi}} \check{Y} \to
\varphi\tag{$\ast$}
\end{equation*}
by observing that
$$
\Hom (\varinjlim_{C^0_{/\varphi}} \check{Y},\varphi) \simeq
\varprojlim_{C_{/ \varphi}} \Hom(\check{Y}, \varphi) \simeq
\varprojlim_{C_{/ \varphi}} \varphi (\widetilde{Y}){\quoi}.
$$
There is a canonical element in $\varprojlim_{C_{/
    \varphi}} \varphi(\widetilde{Y}),$ obtained by associating with
every $(Y, \xi)$ the element $\xi$ of $\varphi (\widetilde{Y})$. Using the fact
that finite $\varprojlim$ exist in $C$, and that the functor $Y
\mapsto \varphi(\widetilde{Y})$ commutes with them, one sees that in
$C_{/\varphi}$ finite $\varprojlim$ exist; a fortiori
$C^0_{/\varphi}$ is filtered. Using the fact that every sheaf $F$ is
an inductive limit of sheaves of the form $\widetilde{Y}$, and using
the fact that $\varphi$ commutes with inductive limits, one easily
concludes that the homomorphism of functors $(\ast)$ is bijective and
thus gives the asserted representation.

\subsubsection{}\etiquette[7.9.2]{VIII.subsubsec7.9.2}

Note also that if $ℰ$ is a topos and $Y$ is an object of $ℰ$, hence
defines a topos $ℰ_{/Y}$, then the datum of a fiber functor $\varphi$
for $ℰ_{/Y}$ is equivalent to the datum of a pair $(\varphi, \xi)$,
where $\varphi$ is a fiber functor for $ℰ$, and $\xi \in \varphi(Y)$. With
the fiber functor $\psi : ℰ_{/Y} \to (\Ens)$ one associates the pair
$(\varphi, \xi)$, where $\varphi$ is the composite $Z \mapsto
\varphi(Z) = \psi (Z \times Y)$, and where $\xi \in \varphi(Y) =
\psi(Y \times Y)$ is the image of the unique element of $\psi(Y)$
under the diagonal morphism of $Y \times Y$.

\subsubsection{}\etiquette[7.9.3]{VIII.subsubsec7.9.3}

We now return to the case of the site $C =
X_{\text{ét}}$, and let $\varphi :
\widetilde{X_{\text{ét}}} \to (\Ens)$ be a fiber functor on the
étale topos of $X$. Considering\pageoriginale the restriction of
$\varphi$ to the full subcategory of
$\widetilde{X_{\text{ét}}}$ formed by the open subsets of this
topos, or, what amounts to the same thing
(\hyperref[VIII.prop6.1]{6.1}), by the open subsets $U$ of $X$, one
obtains a map
$$
\varphi_0 : \Uu (X) \to \{0, 1\}
$$
which commutes with arbitrary suprema and finite infima. (Note that for
$U \in \Uu (X)$, $\varphi(U)$ is either empty or consists of a single
point, and one takes $\varphi_{0} (U) = 0$ or 1 according as one is in
the former or the latter case.)

One easily sees, using the fact that every irreducible closed subset
of $X$ has one and only one generic point \footnote{\ie $X$ is
\emph{sober} in the terminology of IV \hyperref[IV.subsubsec4.2.1]{4.2.1}. The assertion made is a
special case of IV \hyperref[IV.subsubsec4.2.3]{4.2.3}.}, that $\varphi_0$ is defined by means of a
unique $x \in X$, by the condition
$$
\varphi_0 (U) = 1\quad \Longleftrightarrow\quad x \in U {\quoi},
$$
\ie
$$
\varphi (U) \ne\varnothing \Leftrightarrow  x \in U {\quoi}.
$$
Let $F\in \Ob \widetilde{X_{\text{ét}}}$. Then the image of $F$ in
the final sheaf $X$ is an open subset; and since
$$
\varphi (F) \to \varphi(U)
$$
is surjective, one sees that $\varphi(F) \neq ∅$ if and only if
$\varphi(U) \neq ∅$, \ie if and only if $x \in U$. Let
$C_{/\varphi}$ be the category of pairs $(X', \xi)$ where $X'\in \Ob
X_{\text{ét}}$, and $\xi \in \varphi (X')$. It was pointed out in
\hyperref[VIII.subsubsec7.9.1]{7.9.1} that $C^{0}_{/ \varphi}$ is
filtered and connected. Moreover, by
\hyperref[VIII.subsubsec7.9.2]{7.9.2} and by the preceding discussion
applied to $X'$ instead of $X$, if $(X', \xi) \in \Ob
C_{/\varphi}$, there is associated with it a unique point $x' \in
X'$, such that for an étale morphism $X'' \to X'$, $\xi$ lies in the
image of $\varphi (X'') \to \varphi(X')$ if and only if $x'$ lies in the
image of $X''$. It follows that the scheme $\varprojlim$ of the $X'$
over the category $C$ of the $(X', \xi)$ exists and is a strict
localization $Z$ of $X$, thus corresponding to a geometric point $\xi$
and a fiber functor $\mathcal{E}_{\xi}$. For\pageoriginale every sheaf
$F$, one then has
$$
\mathcal{E}_{\xi} (F) \simeq \varinjlim_{C^0_{/\varphi}} F(X'){\quoi}.
$$
In view of \hyperref[VIII.subsubsec7.9.1]{7.9.1}, we conclude that
$\varphi$ is isomorphic to $\mathcal{E}_{\xi}$, which completes the
proof of \hyperref[VIII.thm7.9]{7.9}.

% END INCLUDED SGA4_English_current_buildable_source_20260906/source/08/drafts/SGA4_VIII_LINES_1012_1556_English.texfrag


% BEGIN INCLUDED SGA4_English_current_buildable_source_20260906/source/08/drafts/SGA4_VIII_LINES_1557_1976_English.texfrag
\section{Two spectral sequences for integral morphisms}\etiquette[8]{VIII.sec8}

\begin{enonce*}{Proposition 8.1}\etiquette[8.1]{VIII.prop8.1}
Let $f : X \to Y$ be a surjective integral morphism, and let $F$ be an
abelian sheaf on $Y$. For every integer $i$, let $\Hh^{i} (F)$ be the
presheaf on $(\Sch)_{/Y}$ defined by
$$
\Hh^{i} (F) (Z) = H^i (Z, F_Z){\quoi},
$$
where $F_Z$ is the inverse image of $F$ on $Z$. Then there exists a
spectral sequence (functorial in $F$)
$$
H^{\ast} (X, F) \Longleftarrow E^{pq}_2 = H^p (X /Y, \Hh^q (F))
$$
(``descent spectral sequence'').
\end{enonce*}

Of course, the symbol $H^p (X / Y, G)$, for a presheaf $G$ on
$(\Sch)_{/Y}$, denotes the $p$-th relative \v{C}ech cohomology group,
defined by means of the complex of the $C^n(X / Y, G)
= G(X^{n+1})$, where $X^m$ denotes the $m$-th power in
$(\Sch)_{/Y}$. To prove \hyperref[VIII.prop8.1]{8.1}, let $p_n :
X^{n+1} \to Y$ be the projection, and set
$$
A^n = p_{n\ast} (𝒵_{X^{n + 1}}){\quoi},
$$
where $𝒵_{X^{n + 1}}$ denotes the constant sheaf $𝒵$ on
$X^{n+1}$. Then the $A^n$ are the components of a simplicial abelian
sheaf on $Y$, hence of a complex of abelian sheaves $A^{\ast}$ on $Y$.
Note that there is an evident homomorphism
\begin{equation*}
𝒵_Y \to A^{0}{\quoi}.\tag{$\ast$}
\end{equation*}
With this\pageoriginale understood, we have the following

\begin{enonce*}{Lemma 8.2}\etiquette[8.2]{VIII.lem8.2}
The complex $(A^{\ast})$, endowed with the homomorphism $(\ast)$, is a
resolution of $𝒵_Y$. More generally, for every abelian sheaf $F$ on
$Y$, $A^{\ast} \otimes_{𝒵}F$ is a resolution of $F$.
\end{enonce*}

\noindent
Proof: We may suppose that $Y$ is affine, hence $Y = \Spec A$, $X
= \Spec B$, where $B$ is an integral $A$-algebra. We have $B =
\varinjlim_i B_i$, where $B_i$ ranges over the finitely generated, hence
finite, subalgebras of $B$, whence $X=\varprojlim X_{i}$, where
$X_{i}=\Spec B_{i}$. Using \hyperrefext[VII][VII.cor5.11]{5.11}, one
sees that the augmented complex $A^{\ast}(X/Y)$ is then the inductive
limit of the augmented complexes $A^{\ast}(X_{i}/Y)$. This reduces us
to the case where $f$ is \emph{finite}. It suffices to prove that for
every geometric point $\bar{y}$ of $Y$, the complex $A^{\ast}_{\bar y}$ is
a resolution of $𝒵_{\bar{y}}$, and likewise after tensoring with an
$F$. But if $X_{\bar{y}}$ is the fiber of $X$ at $\bar{y}$, it follows
from \hyperref[VIII.prop5.5]{5.5} that the complex
$A^{\ast}_{\bar{y}}$ is none other than the analogous complex
$A^{\ast}(X_{\bar{y}}/\bar{y})$. This reduces us to the case where $Y$
is the spectrum of a separably closed field $k$. Using
\hyperref[VIII.exem1.3]{1.3} a) and b), we may even suppose that $k$ is
algebraically closed and $X$ is reduced, hence that $X$ is of the form
$I_{Y}$, where $I$ is a finite set. But then the complex
$A^{\ast}\otimes F$ is identified with the trivial cochain complex of
the index set $I$ with coefficients in $F_{\bar{y}}$, and is therefore
indeed a resolution of $F_{\bar{y}}$.

Thus, as a consequence of \hyperref[VIII.lem8.2]{8.2}, we obtain a
spectral sequence
$$
H^{\ast}(Y,F)\Longleftarrow E^{pq}_{2}=H^{p}(H^{q}(Y,A^{\ast}\otimes
F)){\quoi},
$$
and it remains to make the initial term explicit. There is a canonical
homomorphism\pageoriginale
$$
A^{n}\otimes F\longrightarrow p_{n_{\ast}}p^{\ast}_{n}(F){\quoi},
$$
and the latter is an isomorphism, as one sees once again by reduction
to the case where $f$ is finite and passage to fibers. Using
\hyperref[VIII.prop5.5]{5.5}, one concludes that
$$
H^{q}(Y,A^{n}\otimes F)\simeq
H^{q}(X^{n+1},p^{\ast}_{n}(F))=
\Hh^{q}(F)(X^{n+1}){\quoi},
$$
which indeed gives the initial term stated in
\hyperref[VIII.prop8.1]{8.1}.

\begin{enonce*}[remark]{Remark 8.3}\etiquette[8.3]{VIII.rem8.3}
\begin{enumerate}
\renewcommand{\theenumi}{\alph{enumi}}
\renewcommand{\labelenumi}{\theenumi)}
\item When one is given a \emph{locally finite} covering of $Y$ by
  closed subsets $Y_{i}$, and one sets $X=\coprod_{i}Y_{i}$, the
  canonical morphism $f:X\rightarrow Y$ is finite, and the initial
  term of the spectral sequence \hyperref[VIII.prop8.1]{8.1} can be
  made explicit as the cohomology of the complex defined by the
  $H^{q}(Y_{i_{0}\ldots i_{p}},F|Y_{i_{0}\ldots i_{p}})$, where
  $Y_{i_{0}\ldots i_{p}}=Y_{i_{0}}\cap\ldots\cap Y_{i_{p}}$. This is
  therefore the analogue of the Leray spectral sequence for a locally
  finite closed covering of an ordinary topological space (TF Chapter
  II 5.2.4). The latter can moreover also be generalized to a spectral
  sequence relative to a ``finite'' morphism, \ie one which is proper
  with finite fibers, by proceeding as in
  \hyperref[VIII.prop8.1]{8.1}.

\item Note the analogy between the spectral sequence of
  \hyperref[VIII.prop8.1]{8.1} and the Leray spectral sequence of a
  covering $(X_{i})$ of $Y$ (in this case by $X_{i}$ étale over $Y$);
  the latter would formally be obtained by writing the spectral
  sequence \hyperref[VIII.prop8.1]{8.1} for
  $X=\coprod_{i}X_{i}$. It is in fact likely that these two spectral
  sequences admit a common generalization, valid whenever one has a
  family of morphisms $(X_{i}\rightarrow Y)$ which is a ``universally
  effective descent family'' for the fibered category of étale sheaves
  over a variable base scheme\pageoriginale (\cf ${\rm n^{\circ}}$ 9
  below). The analogous question also arises in ordinary topology,
  concerning a common generalization of the two kinds of Leray
  spectral sequences of a covering, assumed either open or locally
  finite closed \footnote[1]{Since these lines were written, P.~Deligne
  has developed a general theory of ``descent spectral sequences,''
  \cf \Exp V bis.}.
\end{enumerate}
\end{enonce*}

\begin{enonce*}{Proposition 8.4}\etiquette[8.4]{VIII.prop8.4}
Let $Y$ be a scheme, let $\pi$ be a profinite group, let
$(\pi_{i})_{i}$ be the projective system of finite discrete quotient
groups of $\pi$, let $(X_{i})_{i\in I}$ be a projective system of
principal coverings of $Y$, with groups $\pi_{i}$, the transition
homomorphisms $X_{j}\rightarrow X_{i}$ being compatible with the
homomorphisms $\pi_{j}\rightarrow \pi_{i}$ on the groups of operators,
and let $X=\varprojlim X_{i}$
\emph{(\cf \hyperrefext[VII][VII.subsec5.1]{5.1})}, so that the group
$\pi$ acts on the $Y$-scheme $X$. Let $F$ be an abelian sheaf on $Y$.
Then there is a ``\emph{Hochschild--Serre}'' spectral sequence
(functorial in $F$)
$$
H^{\ast}(Y,F)\Longleftarrow
E^{pq}_{2}=H^{p}(\pi,\varinjlim_{i}H^{q}(X_{i},F_{X_{i}})){\quoi},
$$
where $F_{X_{i}}$ is the inverse image of $F$ on $X_{i}$, and
$H^{p}(\pi,-)$ denotes Galois cohomology.
\end{enonce*}

The term $E^{pq}_{2}$ written here is also (by the definition of
$H^{p}(\pi,-)$ and the transitivity of inductive limits) isomorphic to
$$
E^{pq}_{2}=\varinjlim_{i}
H^{p}(\pi_{i},H^{q}(X_{i},F_{X_{i}})){\quoi},
$$
which shows that it is enough to find an inductive system of spectral
sequences (depending on the index $i$)
$$
H^{\ast}(Y,F)\Longleftarrow
{}^{i}E^{pq}_{2}=H^{p}(\pi_{i},H^{q}(X_{i},F_{X_{i}})){\quoi}.
$$
This reduces us to defining the spectral sequence in the case where
$\pi$ is finite. Then $f:X\rightarrow Y$ is a covering morphism in
$Y_{\text{\rm ét}}$, so one may write\pageoriginale the Leray
spectral sequence of this morphism. In view of the canonical
isomorphisms
$$
X^{n}\simeq X\mathop×^{\pi}\pi^{n}{\quoi},
$$
a well-known calculation then shows that for every presheaf $\Hh$ on
$X_{\text{\rm ét}}$ which transforms sums into products,
$C^{\ast}(X/Y,\Hh)$ is none other than the complex of homogeneous
cochains of $\pi$ with coefficients in $\Hh(X)$, whence the stated form
of the initial term.

A variant of this proof, avoiding any calculation, consists in
% SGA-SRC-000769: corrupt source notation reconstructed as a contracted product.
regarding $E\mapsto X\mathop{\times}^{\pi}E$ as a morphism from the site of
finite $\pi$-sets to the étale site of $Y$, and writing the Leray
spectral sequence of this morphism. Finally, when $\pi$ is finite, one
may also regard the Hochschild--Serre spectral sequence as a special
case of \hyperref[VIII.prop8.1]{8.1} relative to the morphism
$f:X\rightarrow Y$.

\begin{enonce*}{Corollary 8.5}\etiquette[8.5]{VIII.cor8.5}
Suppose that $Y$ is quasi-compact and quasi-separated. Then the spectral
sequence of \hyperref[VIII.prop8.4]{8.4} takes the form
$$
H^{\ast}(Y,F)\Longleftarrow E^{pq}_{2}=H^{p}(\pi,
H^{q}(X,F_{X})){\quoi}.
$$
\end{enonce*}

Indeed, by \hyperrefext[VII][VII.cor5.8]{5.8}, there are then canonical
isomorphisms
$$
\varinjlim_{i} H^{q}(X_{i},F_{X_{i}})\simeq H^{q}(X,F_{X}){\quoi}.
$$

\begin{enonce*}{Corollary 8.6}\etiquette[8.6]{VIII.cor8.6}
Let $Y$ be a local henselian scheme with closed point $y$, let $F$ be
an abelian sheaf on $Y$, and let $F_{0}$ be the induced sheaf on
$Y_{0}=\Spec (k(y))$. Then the canonical homomorphisms
$$
H^{n}(Y,F)\longrightarrow H^{n}(Y_{0},F_{0})
$$
are\pageoriginale isomorphisms.
\end{enonce*}

Indeed, let $\bar{y}$ be a geometric point over $y$, corresponding to
a separable closure $\bar{k(y)}=k(\bar{y})$ of $k(y)$; since $Y$ is
henselian, the strict localization $X$ of $Y$ at $\bar{y}$ is the
projective limit of the connected étale Galois coverings pointed at
$\bar{y}$, denoted $X_{i}$, so that the hypotheses of
\hyperref[VIII.cor8.5]{8.5} apply. Since $H^{q}(X,F_{X})=0$ for $q>0$
by \hyperref[VIII.cor4.7]{4.7}, one obtains isomorphisms
$$
H^{n}(Y,F)\xleftarrow{\sim} H^{n}(\pi, F(X)){\quoi},
$$
where $\pi$ is the ``Galois group'' of $X$ over $Y$, isomorphic to that
of $\bar{k(y)}$ over $k=k(y)$. Similarly (or by
\hyperref[VIII.cor2.3]{2.3}) one finds
$$
H^{n}(Y_{0},F_{0})\xleftarrow{\sim} H^{n}(\pi,F_{0}(X_{0})){\quoi},
$$
where $X_{0}=\fprod{X}{Y}{Y_{0}}\simeq \Spec (k(\bar{y}))$. But the
restriction homomorphism $F(X)\rightarrow F_{0}(X_{0})$ is an
isomorphism by \hyperref[VIII.cor4.8]{4.8}, from which the conclusion
\hyperref[VIII.cor8.6]{8.6} follows immediately.

\section{Descent of étale sheaves}\etiquette[9]{VIII.sec9}

This section will not be used later in this Seminar, and may be omitted
on a first reading.

\begin{enonce*}{Proposition 9.1}\etiquette[9.1]{VIII.prop9.1}
Let\pageoriginale $f:X'\rightarrow X$ be a surjective (\resp
universally submersive \emph{(\SGA I IX 2.1)}) morphism of schemes.
Then the functor $f^{\ast}:X_{\text{\rm ét}}\rightarrow
X'_{\text{\rm ét}}$ is faithful and ``conservative''
\emph{(\cf \hyperref[VIII.cor3.6]{3.6})} (\resp induces a
\emph{fully faithful} functor from the category of sheaves on
$X_{\text{\rm ét}}$ to the category of sheaves on
$X'_{\text{\rm ét}}$ endowed with descent data relative to
$f:X'\rightarrow X$).
\end{enonce*}

The first assertion follows immediately from the transitivity of fiber
functors (\hyperref[VIII.subsec3.4]{3.4}) and from
(\hyperref[VIII.cor3.7]{3.7}); the second (which can also be stated by
saying that $f$ is a \emph{descent morphism} relative to the fibered
category of étale sheaves on variable schemes) also means that for two
sheaves $F$, $G$ on $X$, the natural diagram
\begin{equation*}
\Hom(F,G)\rightarrow \Hom (F',G')\rightrightarrows \Hom (F'',G'')\tag{$\ast$}
\end{equation*}
is exact, where $F'$ and $G'$ (\resp $F''$ and $G''$) are the inverse
images of $F$ and $G$ on $X'$ (\resp on
$X''=\fprod{X'}{X'}{X}$). Taking the final sheaf for $F$, the assertion
gives the following

\begin{enonce*}{Corollary 9.2}\etiquette[9.2]{VIII.cor9.2}
Let $f:X'\rightarrow X$ be a universally submersive morphism. Then for
every sheaf $F$ on $X$, the diagram
$$
\Gamma(X,F)\rightarrow \Gamma(X',F')\rightrightarrows \Gamma(X'',F'')
$$
is exact, where $X''=\fprod{X'}{X'}{X}$, and where $F'$, $F''$ are the
inverse images of $F$ on $X'$, $X''$.
\end{enonce*}

We shall prove \hyperref[VIII.prop9.1]{9.1} using the following

\begin{enonce*}{Lemma 9.3}\etiquette[9.3]{VIII.lem9.3}
Suppose that $f$ is universally submersive. Let $F$ be a sheaf on $X$,
denote by $S(F)$ the set of subsheaves of \sisi{$X$}{$F$},\pageoriginale
and define $S(F')$, $S(F'')$ analogously. Then the diagram of natural
maps
$$
S(F)\rightarrow S(F')\rightrightarrows S(F'')
$$
is exact.
\end{enonce*}

\noindent
Proof. The fact that $S(F)\rightarrow S(F')$ is injective follows
immediately from the fact that the functor $f^{\ast}$ is conservative;
for if $F_{i}$ $(i=1,2)$ are two subsheaves of $F$ such that
$f^{\ast}(F_{1})=f^{\ast}(F_{2})$, then the inclusions
$F_{3}=F_{1}\cap F_{2}\rightarrow F_{i}$ $(i=1,2)$ are isomorphisms,
since they become so after applying the functor $f^{\ast}$, hence
$F_{1}=F_{2}$. It remains to prove that if $G'$ is a subsheaf of $F'$
such that $\pr^{\ast}_{1}(G')=\pr^{\ast}_{2}(G')$, then $G'$ is the
inverse image of a subsheaf of $F$. Note that one can find an
epimorphism $F_{1}\rightarrow F$ in $X_{\text{\rm ét}}$, where
$F_{1}$ is \emph{representable} by a scheme étale over $X$.
Introduce $F_{2}=\fprod{F_{1}}{F_{1}}{F}$, and likewise $F'_{1}$,
$F''_{1}$, $F'_{2}$, $F''_{2}$, sheaves giving rise to a diagram of
sets
\[
\xymatrix@R=.75cm{
  S(F)  \ar[r] \ar[d]& S(F') \ar[r]\ar[d] & S(F'') \ar[d]\\
S(F_1) \ar[r] \ar@<3pt>[d] \ar@<-3pt>[d]& S(F'_1) \ar@<3pt>[r]
  \ar@<-3pt>[r] \ar@<3pt>[d] \ar@<-3pt>[d]& S(F''_{1})\ar@<3pt>[d] \ar@<-3pt>[d] \\
S(F_2)  \ar[r] & S(F'_2)\ar@<3pt>[r]
  \ar@<-3pt>[r] & S(F''_2),
}
\]
in which the columns are exact, by the exactness properties in the
category of sheaves on $X$, $X'$, $X''$, respectively. A standard
diagram chase then shows that, to prove that the first row is exact,
it is enough to prove this for rows 2 and 3. For row 2, this follows
from the fact that $F_{1}$ is representable, that $X'\rightarrow X$ is
universally\pageoriginale submersive, and from
\SGA\sisi{}{ 1} IX 2.3. On the other hand, $F_{2}$ is also
representable, since it is a subsheaf of
$\fprod{F_{1}}{F_{1}}{X}$, which is representable, and one applies
\hyperref[VIII.prop6.1]{6.1}. Thus row 3 is also exact, \cqfd

We now prove \hyperref[VIII.prop9.1]{9.1}, \ie that every morphism
$u':F'\rightarrow G'$ compatible with the descent data on $F'$, $G'$,
the morphism $u'$ comes from a morphism $u:F\rightarrow G$. Let
$H=F\times G$, hence
$H'=F'\times G'$, $H''=F''\times G''$; then the \emph{graph} of $u'$
is a subsheaf $\Gamma'$ of $H'$, whose two inverse images are equal to
the graph of the same morphism $u'':F''\rightarrow G''$. Hence, by
\hyperref[VIII.lem9.3]{9.3}, $\Gamma'$ comes from a subsheaf $\Gamma$
of $H$. I claim that $\Gamma$ is the graph of a morphism
$u:F\rightarrow G$, \ie that the morphism $p:\Gamma\rightarrow F$
induced by $\pr_{1}$ is an isomorphism: indeed, it becomes an
isomorphism after the base change $X'\rightarrow X$, and we apply the
already proved part of \hyperref[VIII.prop9.1]{9.1}. The morphism
$u:F\rightarrow G$ then has the required property, \cqfd

\begin{enonce*}{Theorem 9.4}\etiquette[9.4]{VIII.thm9.4}
Let $f:X'\rightarrow X$ be a surjective morphism of schemes. Suppose
that one of the following hypotheses holds:
\begin{enumerate}
\renewcommand{\theenumi}{\alph{enumi}}
\renewcommand{\labelenumi}{\theenumi)}
\item $f$ is integral.

\item $f$ is proper.

\item $f$ is flat and locally of finite presentation
\footnote[1]{It is likely that this second hypothesis is in fact
superfluous.}.

\item $X$ is discrete (e.g. the spectrum of a field).
\end{enumerate}

Then $f$ is an \emph{effective descent morphism} for the fibered
category $\Ff$ over $(\Sch)$ of étale sheaves on variable schemes,
\ie the functor $f^{\ast}$ induces an \emph{equivalence} between the
category of\pageoriginale sheaves on $X$ and the category of sheaves on
$X'$ endowed with descent data relative to $f:X'\rightarrow X$.
\end{enonce*}

\setcounter{subsection}{4}
\setcounter{subsubsection}{0}
\subsubsection{}\etiquette[9.4.1]{VIII.subsubsec9.4.1}

Let $F'$ be a sheaf on $X'$ endowed with descent data relative to
$f:X'\rightarrow X$. Define a functor
$$
G:X^{0}_{\text{\rm ét}}\longrightarrow (\Ens)
$$
by the formula
$$
G(Y)=\Ker (F'(Y')\rightrightarrows F''(Y'')){\quoi}.
$$
One sees immediately that $G$ is a \emph{sheaf} for the étale
topology. On the other hand, there is a canonical injective
homomorphism $G\rightarrow f_{\ast}(F')$, hence a homomorphism
$f^{\ast}(G)\rightarrow F'$, and the latter is evidently compatible
with the descent data. It remains to determine whether this
homomorphism is an \emph{isomorphism}. Note that $G$ can also be defined
by the exactness of
$$
G\rightarrow f_{\ast}(F')\rightrightarrows g_{\ast}(F''){\quoi},
$$
so for every base change $Y\xrightarrow{h}X$ \emph{which commutes with
  the formation of} $f_{\ast}(F')$ and $g_{\ast}(F'')$,
$h^{\ast}(G)=G_{Y}$ is identified with the sheaf ``descended'' from
$F'_{Y'}$ by $Y'\rightarrow Y$, and of course the homomorphism
$f_{Y}{}^{\ast}(G_{Y})\rightarrow (F_{Y})'$ is the one deduced from
$f^{\ast}(G)\rightarrow F'$ by base change. Now \emph{suppose that the
  base change} $f:Y=X'\rightarrow X$ \emph{commutes with the formation
  of} $f_{\ast}(F')$ and $g_{\ast}(F'')$, and note that since
$Y'=\fprod{X'}{X'}{X}$ has a section over $Y=X'$, hence $Y'\rightarrow
Y$ is an \emph{effective} descent morphism for every fibered category,
in particular for $\mathcal{F}$, it follows that
$f^{\ast}_{Y}(G_{Y})\rightarrow (F_{Y})'$ is an isomorphism. Thus
$f^{\ast}(G)\rightarrow F'$ becomes an isomorphism after the base
change $Y'=\fprod{X'}{X'}{X}\rightarrow X'$. Since the latter has a
section, $f^{\ast}(G)\rightarrow F'$ is an isomorphism, and hence the
descent datum\pageoriginale under consideration on $F'$ is effective.

\subsubsection{}\etiquette[9.4.2]{VIII.subsubsec9.4.2}

This proves Theorem \hyperref[VIII.thm9.4]{9.4} in case a), by
\hyperref[VIII.cor5.6]{5.6}. Case b) also follows from the fact that if
$f$ is proper, then $f_{\ast}$ commutes with \emph{every} base change
$Y'\rightarrow Y$ (which will be proved in a later exposé (XII \hyperref[XII.thm5.1]{5.1}
(i))).

\subsubsection{}\etiquette[9.4.3]{VIII.subsubsec9.4.3}

In case c), the question being local on $X$, we may suppose that $X$
is affine and that there exists an affine scheme $X'_{1}$, of finite
presentation, quasi-finite and faithfully flat over $X$, together with
an $X$-morphism $X'_{1}\rightarrow X'$, which reduces us to the case
where $f$ is quasi-finite. After localizing further on $X$, this time
in the sense of the étale topology, one sees that there will exist an
open subset $X'_{1}$ of $X'$ such that $X'_{1}\rightarrow X$ is finite
and surjective, which reduces us to the case where $f$ is finite and
surjective, already treated in a).

\subsubsection{}\etiquette[9.4.4]{VIII.subsubsec9.4.4}
In case d), we may suppose (by localizing on $X$) that $X$ consists of
a single point, and even, in view of \hyperref[VIII.thm1.1]{1.1}, that
it is the spectrum of a field $k$. Moreover, $f$ is universally open
(\EGA IV 2.4.9), so the hypotheses of \hyperref[VIII.prop9.1]{9.1}
hold, and the argument of \hyperref[VIII.subsubsec9.4.1]{9.4.1} applies
again, making a base change with $Y=\Spec k$, where $k=k(x)$, for an
$x\in X$. But then it is still true that $f_{\ast}$ commutes with the
base change $Y\rightarrow X$, as we shall see later (XVI \hyperref[XVI.compl1.4]{1.4} and \hyperref[XVI.cor1.5]{1.5}).

This completes the proof of \hyperref[VIII.thm9.4]{9.4}.


\ifx\danslelivre\undefined
\end{document}
\fi

% END INCLUDED SGA4_English_current_buildable_source_20260906/source/08/drafts/SGA4_VIII_LINES_1557_1976_English.texfrag


% END INCLUDED SGA4_English_current_buildable_source_20260906/source/08/08_English.tex


\makeatletter
\def\input@path{{../09/}}
\makeatother
% BEGIN INCLUDED SGA4_English_current_buildable_source_20260906/source/09/09_English.tex
\ifx\danslelivre\undefined
% Release engine: XeLaTeX (the LuaTeX compatibility layer breaks Xy-pic output).
\documentclass[10pt]{../fichiers-configuration-latex/smfart}
% BEGIN INCLUDED SGA4_English_current_buildable_source_20260906/source/fichiers-configuration-latex/paquets.tex
\usepackage{polyglossia}
%\usepackage[francais]{babel}
\usepackage{fontspec}
\usepackage[math-style=ISO,bold-style=upright,sans-style=upright]{unicode-math}
%\usepackage[mathletters]{ucs}
%\usepackage[utf8]{inputenc}
%\usepackage[T1]{fontenc}
\usepackage[unicode,hypertexnames=false,colorlinks=true,linkcolor=blue,citecolor=Sepia,urlcolor=magenta]{hyperref}
\setmainlanguage{french}
\setmainfont[Ligatures=TeX]{Linux Libertine O}
% variante possible : TeXGyrePagella ?
\setmathfont{XITS Math}

% maths
\usepackage{amsmath}
%\usepackage{amssymb}
%\usepackage{mathrsfs}
%\usepackage{stmaryrd}
%\usepackage{euscript}

% couleurs
\usepackage[usenames,dvipsnames]{xcolor}

% graphiques
\usepackage{graphicx}
\usepackage[all]{xy}
%\usepackage{tikz}
%\usetikzlibrary{matrix,arrows,calc}
%\tikzset{column sep=3em,row sep=2em,text height=1.5ex, text depth=0.25ex,>=angle 90}
%\tikzset{description/.style={fill=white,inner sep=1pt}}

% hyperliens, 'etiquettes
\usepackage{xr-hyper}
\usepackage{../fichiers-configuration-latex/smfthm}
%\ifx\cacheretiquettes\undefined
%\usepackage[inline]{showlabels} \else \fi % pour rendre visibles les 'etiquettes
%\usepackage{../fichiers-configuration-latex/makeidx} % pour faire un index
%\usepackage[pagebackref,unicode]{hyperref}  % pour voir o`u sont cit'ees les r'ef'erences
% hyperref is loaded once above with the cumulative-reader destination policy.
%\hypersetup{breaklinks=true} %permet le retour `a la ligne dans les liens trop longs. Incompatible avec pdflatex
%\renewcommand*{\backref}[1]{{\small [cit\'e page(s)~#1.]}}

\usepackage{xspace}
\usepackage{pifont} % pour expose XVII

% END INCLUDED SGA4_English_current_buildable_source_20260906/source/fichiers-configuration-latex/paquets.tex

% BEGIN INCLUDED SGA4_English_current_buildable_source_20260906/source/fichiers-configuration-latex/numerotation.tex
\renewcommand{\thefootnote}{\roman{footnote}} % pour mettre les notes en bas de page sous la forme (ii) et non (2), qui ressemble trop (?) a un exposant.

\newcounter{pageoriginal}
\newcommand{\pageoriginale}{\refstepcounter{pageoriginal}\marginpar{\footnotesize\textbf{\thepageoriginal}}\xspace}
\let\pageoriginaled\pageoriginale

\usepackage{ifthen}
\newboolean{original}
\def\sisi{\ifthenelse{\boolean{original}}}
\setboolean{original}{false}

\newcounter{cnote}
\def\nde#1{\ifthenelse{\boolean{original}}{}{\stepcounter{cnote}\def\thefootnote{\arabic{cnote}}\protect\footnote[\value{cnote}]{N.D.E.: #1}}}
\def\ndemark{\ifthenelse{\boolean{original}}{}{\stepcounter{cnote}\def\thefootnote{\arabic{cnote}}\protect\footnotemark[\value{cnote}]}}
\def\ndetext#1{\ifthenelse{\boolean{original}}{}{\def\thefootnote{\arabic{cnote}}\protect\footnotetext[\value{cnote}]{N.D.E.: #1}}}

% END INCLUDED SGA4_English_current_buildable_source_20260906/source/fichiers-configuration-latex/numerotation.tex

% BEGIN INCLUDED SGA4_English_current_buildable_source_20260906/source/fichiers-configuration-latex/macros.tex
\setmainlanguage{french}
\setotherlanguage{english}


%% Fixed, engine-independent release identifier.  The inherited dynamic Git
%% probe used Lua io.popen, exposed build-state-dependent output, and printed
%% a visible question mark outside a Git checkout.
\def\showgitstatus{English release r2, 2026-07-23}


% macros pour LaTeX
\renewcommand{\mathsf}{}

% Police différente pour U+2113 SCRIPT SMALL L
\setmathfont[range="2113]{latinmodern-math.otf}

% Give every inherited semantic label its own stable PDF destination instead
% of inheriting whichever automatic equation/page anchor happened to be live.
% The legacy label remains as a compatibility alias used by existing edges.
\makeatletter
\def\etiquette[#1]#2{%
  \Hy@raisedlink{\hyper@anchorstart{sga4.semantic.#2}\hyper@anchorend}%
  \def\@currentHref{sga4.semantic.#2}%
  \def\@currentlabel{#1}%
  \label{#2}%
}
\makeatother
\def\eqhyperref[#1]#2{(\hyperref[#1]{#2})}
\def\pagehyperref[#1]#2{(page \hyperref[#1]{#2})}
\def\hyperrefext[#1][#2]#3{\hyperref[#2]{#1 #3}}
\def\eqhyperrefext[#1][#2]#3{\hyperref[#2]{#1 #3}}     %cf. exposes 17
\def\neweqhyperrefext[#1][#2]#3{\hyperref[#2]{#1 #3}}  %cf. exposes 17

\def\shat#1{\protect\widehat{#1}}
\def\stilde#1{\protect\widetilde{#1}}
%%%%%%%%%%%%%%%%%%


\def\etcn{etc.\xspace}
\def\etco{etc...\xspace}
\def\etc{\sisi{\etco}{\etcn}}
\def\ldotsn{\sisi{...}{\etcn}}

\def\cf{{cf.}\kern.3em}
\def\Cf{{Cf.}\kern.3em}
\def\ie{{i.e.}\kern.3em}
\def\iev{{i.e.},\xspace}
\def\loccito{loc.\kern2pt cit.\xspace}
\def\loccitn{\textit{loc.\kern2pt cit.}\xspace}
\def\loccit{\sisi{\loccito}{\loccitn}}
\def\EGA{\textup{EGA}\kern.2em}
\def\SGA{\textup{SGA}\kern.3em}

\def\quoi{}

\def\quav{\quad,\quad}
\def\quaet{\quad \text{et}\quad}
\def\fibreproduct#1#2#3{#1 ×_{#3} #2}
\let\fprod\fibreproduct
\def\ccheck#1#2{\mbox{${\check{\check{\rm #1}}}^{\lower2pt\hbox{$\scriptstyle #2$}}$}}
\def\sheaf#1{\protect\mathcal{#1}}
\def\vieuxsheaf#1{{#1}} % cf. problème avec sheaf{a\Phi_i} par exemple...
\DeclareMathOperator{\Hom}{Hom}
\DeclareMathOperator{\can}{can}
\DeclareMathOperator{\Ind}{Ind}
\DeclareMathOperator{\ind}{ind}
%\DeclareMathOperator{\Hhomtop}{\mathcal{H}\!\mathit{omtop}}
\DeclareMathOperator{\Hhomtop}{\mathbf{Homtop}}
\DeclareMathOperator{\Point}{Point}
\DeclareMathOperator{\Pro}{Pro}
\DeclareMathOperator{\Top}{Top}
\DeclareMathOperator{\Mod}{Mod}
\DeclareMathOperator{\Cart}{Cart}
%\DeclareMathOperator{\Mmorsite}{\mathcal{M}\!\mathit{orsite}}
\DeclareMathOperator{\Mmorsite}{\mathbf{Morsite}}
\DeclareMathOperator{\loc}{loc}
\DeclareMathOperator{\Loc}{Loc}
\DeclareMathOperator{\Pic}{Pic}
\DeclareMathOperator{\Et}{Et}
\DeclareMathOperator{\card}{card}
\DeclareMathOperator{\Ens}{Ens}
\DeclareMathOperator{\h}{\underline{H}}
\DeclareMathOperator{\Fonct}{Fonct}
\DeclareMathOperator{\Ker}{Ker}
\DeclareMathOperator{\but}{but}
\DeclareMathOperator{\Ob}{Ob}
\DeclareMathOperator{\ob}{ob}
\DeclareMathOperator{\Lib}{Lib}
\DeclareMathOperator{\esj}{esj}
\DeclareMathOperator{\id}{id}
\DeclareMathOperator{\Id}{Id}
\DeclareMathOperator{\ad}{ad}
\DeclareMathOperator{\Iim}{Im}
\DeclareMathOperator{\pr}{pr}
\DeclareMathOperator{\Cat}{Cat}
\DeclareMathOperator{\car}{car}
\DeclareMathOperator{\SysInd}{SysInd}
\DeclareMathOperator{\Sup}{Sup}
\DeclareMathOperator{\Coker}{Coker}
\DeclareMathOperator{\Filt}{Filt}
\DeclareMathOperator{\acc}{acc}
\DeclareMathOperator{\sur}{sur}
\DeclareMathOperator{\cat}{cat}
\DeclareMathOperator{\Ord}{Ord}
\DeclareMathOperator{\Gr}{Gr}
\DeclareMathOperator{\Cov}{Cov}
\DeclareMathOperator{\Caf}{Caf}
\DeclareMathOperator{\coker}{coker}
\DeclareMathOperator{\coim}{coim}
\DeclareMathOperator{\im}{im}
\DeclareMathOperator{\Ouv}{Ouv}
\DeclareMathOperator{\diag}{diag}
\DeclareMathOperator{\Esp}{Esp}
\DeclareMathOperator{\TOP}{TOP}
\DeclareMathOperator{\Unif}{Unif}
\DeclareMathOperator{\Iind}{\mathbf{Ind}}
\DeclareMathOperator{\Hhomcart}{\mathbf{Homcart}}
\DeclareMathOperator{\Hhomtopess}{\mathbf{Homtopess}}
\DeclareMathOperator{\Ppoints}{\mathbf{Points}}
\DeclareMathOperator{\Ppoint}{\mathbf{Point}}
\DeclareMathOperator{\Ttors}{\mathbf{Tors}}
\DeclareMathOperator{\Ssex}{\mathbf{Sex}}
\DeclareMathOperator{\Ffib}{\mathbf{Fib}}
\DeclareMathOperator{\Vvois}{\mathbf{Vois}}
\DeclareMathOperator{\Ppointess}{\mathbf{Pointess}}
\DeclareMathOperator{\scrHom}{\mathbf{Hom}}
\DeclareMathOperator{\sheafPic}{\mathbf{Pic}}
\DeclareMathOperator{\sheafAut}{\mathbf{Aut}}
\DeclareMathOperator{\Hhomtoplocan}{\mathbf{Homtoplocan}}
\DeclareMathOperator{\Hhomtopan}{\mathbf{Homtopan}}
\DeclareMathOperator{\Pptgeom}{\mathbf{Ptgéom}}
\DeclareMathOperator{\Morsite}{Morsite}
\DeclareMathOperator{\Fib}{Fib}
\DeclareMathOperator{\Fais}{Fais}
\DeclareMathOperator{\kar}{kar}
\DeclareMathOperator{\Groupes}{Groupes}
\DeclareMathOperator{\Restr}{Restr}
\DeclareMathOperator{\Isom}{Isom}
\DeclareMathOperator{\Prol}{Prol}
\DeclareMathOperator{\Res}{Res}
\DeclareMathOperator{\Tors}{Tors}
\DeclareMathOperator{\esp}{esp}
\DeclareMathOperator{\es}{es}
\DeclareMathOperator{\od}{od}
\DeclareMathOperator{\gr}{gr}
\DeclareMathOperator{\sob}{sob}
\DeclareMathOperator{\dfn}{dfn}
\DeclareMathOperator{\Sch}{Sch}
\DeclareMathOperator{\ouv}{ouv}
\DeclareMathOperator{\cons}{cons}
\DeclareMathOperator{\Nerf}{Nerf}
\DeclareMathOperator{\Points}{Points}
\DeclareMathOperator{\Espsob}{Espsob}
\DeclareMathOperator{\Homtop}{Homtop}
\DeclareMathOperator{\Pt}{Pt}
\DeclareMathOperator{\Quot}{Quot}
\DeclareMathOperator{\glob}{glob}
\DeclareMathOperator{\Homfib}{Homfib}
\DeclareMathOperator{\Qouv}{Qouv}
\DeclareMathOperator{\Quotouv}{Quotouv}
\DeclareMathOperator{\ens}{ens}
\DeclareMathOperator{\ab}{ab}
\DeclareMathOperator{\an}{an}
\DeclareMathOperator{\fs}{fs}
\DeclareMathOperator{\red}{red}
\DeclareMathOperator{\Ab}{Ab}
\DeclareMathOperator{\cosq}{cosq}
\DeclareMathOperator{\Ann}{Ann}
\DeclareMathOperator{\Zar}{Zar}


%%%%%%%%%%%%
%%%%% SGA4-2
\DeclareMathOperator{\Sex}{Sex}
\DeclareMathOperator{\vois}{vois}
\DeclareMathOperator{\Vois}{Vois}
\DeclareMathOperator{\Ext}{Ext}
\DeclareMathOperator{\ext}{ext}
\DeclareMathOperator{\Fer}{Fer}
\DeclareMathOperator{\Add}{Add}
\DeclareMathOperator{\Cocart}{Cocart}
\DeclareMathOperator{\cosk}{cosk}
\DeclareMathOperator{\plat}{plat}

\DeclareMathOperator{\faisc}{faisc}
\DeclareMathOperator{\sk}{sk}
\DeclareMathOperator{\ch}{ch}

\def\HH{\mathbf{H}}
\def\qc{\mathrm{qc}}
\def\et{\mathrm{et}}
\def\coh{\mathrm{coh}}

\def\Hvr#1{%
\begin{array}[b]{l}
{\kern2pt\scriptstyle r}\\[-8pt]
{\kern1pt\scriptstyle \vee}\\[-5pt]
H^{#1}
\end{array}}

\def\Hvv#1{%
\begin{array}[b]{l}
\kern2pt\rotatebox{90}{${\scriptscriptstyle\ll}$}\\[-6pt]
H^{#1}
\end{array}}

\def\Hvi#1{%
\begin{array}[b]{l}
{\kern0.2pt\scriptscriptstyle \infty}\\[-8pt]
{\kern1pt\scriptstyle \vee}\\[-5pt]
H^{#1}
\end{array}}

\def\Hv#1{%
\begin{array}[b]{l}
{\kern1pt\scriptstyle \vee}\\[-5pt]
H^{#1}
\end{array}}

\DeclareMathOperator{\cart}{cart}
\DeclareMathOperator{\Ffaisc}{\mathbf{Faisc}}


\def\Ev#1{%
\begin{array}[b]{l}
{\kern1pt\scriptstyle \vee}\\[-5pt]
E'_{#1}
\end{array}}

\DeclareMathOperator{\Cont}{Cont}
\DeclareMathOperator{\Ccont}{\mathscr{C}\!\mathit{ont}}
\DeclareMathOperator{\Site}{Site}
\newcommand\MM{\mathbf{M}}
\def\limtop#1{{\displaystyle{\varprojlim_{#1}\!\mbox{\rm top }}}}
\def\Limtop{\mathop{\varprojlim\!\mbox{\rm top}}}
\DeclareMathOperator{\spec}{spec}
\def\ss{\scriptstyle{\mathscr{S}}}
\DeclareMathOperator{\Gal}{Gal}
\def\smallm{\scriptstyle{\mathscr{M}}}
\DeclareMathOperator{\End}{End}
\DeclareMathOperator{\ass}{ass}
\DeclareMathOperator{\zar}{zar}
\DeclareMathOperator{\Ogg}{Ogg}
\DeclareMathOperator{\res}{res}
\DeclareMathOperator{\tr}{tr}
\DeclareMathOperator{\Supp}{Supp}
\DeclareMathOperator{\codim}{codim}
\DeclareMathOperator{\rk}{rk}
\DeclareMathOperator{\cur}{cur}
\DeclareMathOperator{\cdqc}{cdqc}
\newcommand\DD{\mathbf{D}}
\DeclareMathOperator{\cl}{cl}
\def\JJ{\mathbf{J}}
\def\FF{\mathbf{F}}
\def\BB{\mathbf{B}}
\DeclareMathOperator{\Alget}{Alget}
\DeclareMathOperator{\Ass}{Ass}
\DeclareMathOperator{\supp}{supp}
\DeclareMathOperator{\rad}{rad}
\DeclareMathOperator{\prof}{prof}
\DeclareMathOperator{\cint}{int}
\DeclareMathOperator{\sic}{sic}
\DeclareMathOperator{\Tor}{Tor}
\DeclareMathOperator{\tors}{tors}
\DeclareMathOperator{\Sym}{Sym}
\DeclareMathOperator{\parf}{parf}
\DeclareMathOperator{\tor}{tor}
\DeclareMathOperator{\defn}{defn}
\DeclareMathOperator{\Div}{Div}
\DeclareMathOperator{\Pol}{Pol}
\DeclareMathOperator{\Fex}{Fex}
\DeclareMathOperator{\Lie}{Lie}
\DeclareMathOperator{\tot}{tot}
\DeclareMathOperator{\Aut}{Aut}
\DeclareMathOperator{\HOM}{HOM}
\DeclareMathOperator{\pch}{pch}
\DeclareMathOperator{\Ch}{Ch}
\DeclareMathOperator{\EXT}{EXT}
\DeclareMathOperator{\PIC}{PIC}
\DeclareMathOperator{\Rf}{Rf}
\DeclareMathOperator{\inj}{inj}
\newcommand\VV{\mathbf{V}}

\def\Xv{%
\begin{array}[b]{l}
{\kern1pt\scriptstyle \vee}\\[-5pt]
X
\end{array}}

%%%%%%%
\newcommand\TT{\mathbf{T}}
\newcommand\XX{\mathbf{X}}
\newcommand\EE{\mathbf{E}}
%%%%%%

\def\qedsymbol{\textup{C.Q.F.D.}}
\let\CQFD\qed
\let\cqfd\qed

%%%%%\bigdot
\font\bigsymb=cmsy10 at 4pt
\def\bigdot{{\kern1.2pt\raise 1.5pt\hbox{\bigsymb\char15}}}
%%%%%

\def\Lright{\mathop{\vcenter{\hbox{\hbox{$\longrightarrow$}\lower8pt \hbox{\kern-16pt\hbox{$\longrightarrow$}}}}}}

\def\Aa{{\mathcal A}}
\def\Bb{{\mathcal B}}
\def\Cc{{\mathcal C}}
\def\Dd{{\mathcal D}}
\def\Ee{{\mathcal E}}
\def\Ff{{\mathcal F}}
\def\Gg{{\mathcal G}}
\def\Hh{{\mathcal H}}
\def\Ii{{\mathcal I}}
\def\Jj{{\mathcal J}}
\def\Kk{{\mathcal K}}
\def\Ll{{\mathcal L}}
\def\Mm{{\mathcal M}}
\def\Nn{{\mathcal N}}
\def\Oo{{\mathcal O}}
\def\Pp{{\mathcal P}}
\def\Qq{{\mathcal Q}}
\def\Rr{{\mathcal R}}
\def\Ss{{\mathcal S}}
\def\Tt{{\mathcal T}}
\def\Uu{{\mathcal U}}
\def\Vv{{\mathcal V}}
\def\Ww{{\mathcal W}}
\def\Xx{{\mathcal X}}
\def\Yy{{\mathcal Y}}
\def\Zz{{\mathcal Z}}

\newcommand{\CC}{\mathbf{C}}
\newcommand{\GG}{\mathbf{G}}
\newcommand{\KK}{\mathbf{K}}
\newcommand{\LL}{\mathbf{L}}
\newcommand{\NN}{\mathbf{N}}
\newcommand{\PP}{\mathbf{P}}
\newcommand\QQ{\mathbf{Q}}
\newcommand\RR{\mathbf{R}}
\newcommand\ZZ{\mathbf{Z}}

\def\to{\mathchoice{\longrightarrow}{\rightarrow}{\rightarrow}{\rightarrow}}
\def\from{\mathchoice{\longleftarrow}{\leftarrow}{\leftarrow}{\leftarrow}}
\def\mto{\mathchoice{\longmapsto}{\mapsto}{\mapsto}{\mapsto}}
\def\mfrom{\mathrel{\from\joinrel\mapstochar\,}}
\def\hto{\mathrel{\lhook\joinrel\to}}
\def\hfrom{\mathrel{\from\joinrel\rhook}}

\def\implies{\mathchoice{\Longrightarrow}{\Rightarrow}{\Rightarrow}{\Rightarrow}}
\def\implied{\mathchoice{\Longleftarrow}{\Leftarrow}{\Leftarrow}{\Leftarrow}}
\def\iff{\mathchoice{\Longleftrightarrow}{\Leftrightarrow}{\Leftrightarrow}{\Leftrightarrow}}

\def\To#1{\mathchoice{\xrightarrow{\labelstyle~~#1~}}{\stackrel{#1}{\longrightarrow}}{}{}}
\def\Mto#1{\mathrel{\mapstochar\To{#1}}}
\def\From#1{\mathchoice{\xleftarrow{\labelstyle~~#1~}}{\stackrel{#1}{\longleftarrow}}{}{}}

\newcommand{\isomto}{\overset{\sim}{\longrightarrow}}
\newcommand{\isomfrom}{\overset{\sim}{\longleftarrow}}

\newcommand{\IFF}{\unskip~$\iff$~\ignorespaces}
\newcommand{\THEN}{\unskip~$\implies$~\ignorespaces}

\let\leq\leqslant
\let\geq\geqslant
\let\le\leqslant
\let\ge\geqslant

\def\cf{{cf.}\kern.3em}
\def\Cf{{Cf.}\kern.3em}
\def\ie{{i.e.}\kern.3em}
\def\iev{{i.e.},\xspace}
%\def\loccito{loc.\kern2pt cit.\xspace}
%\def\loccitn{\textit{loc.\kern2pt cit.}\xspace}
%\def\loccit{\sisi{\loccito}{\loccitn}}
\def\EGA{\textup{EGA}\kern.2em}
\def\SGA{\textup{SGA}\kern.3em}

\let\oldS\S
\def\S{\oldS\kern .15em}
\def\ptbl{.\kern .2em}
\def\resp{resp.\kern.3em}
\def\Exp{Exp.\kern.3em}

\let\bar\overline
\let\hat\widehat
\let\tilde\widetilde
\let\dpl\displaystyle
\let\ldots\dots
\let\emptyset\varnothing
\let\skpt\skippointrait
\let\epsilon\varepsilon

\def\sfrac#1#2{{#1}/{#2}}

\DeclareMathOperator{\cd}{cd}
%\DeclareMathOperator{\Hhom}{\mathcal{H}\!\mathit{om}}
\DeclareMathOperator{\Hhom}{\mathbf{Hom}}
%\DeclareMathOperator{\Eext}{\mathcal{E}\!\mathit{xt}}
\DeclareMathOperator{\Eext}{\mathbf{Ext}}
\DeclareMathOperator{\Spec}{Spec}
\DeclareMathOperator{\Tr}{Tr}

\let\emptyset=\varnothing

%%%%%%%%%%%%%%%
%%% expose XVII
%%%%%%%%%%%%%%%

\makeatletter
\let\old@optexpose\optexpose
\newcommand{\empty@optexpose}[1]{}
\newcommand{\paren@optexpose}[1]{\old@optexpose{#1}(}
\newcommand{\refnoexpose}[1]{\let\optexpose\empty@optexpose\ref{#1}\let\optexpose\old@optexpose}%\refnoexpose{X.cor5.2} engendre "5.2" au lieu de "X 5.2"
\newcommand{\neweqref}[1]{\let\optexpose\paren@optexpose\textup{\ref{#1})}\let\optexpose\old@optexpose}%\neweqref{XII.eq4.1} engendre "XII (4.1)" au lieu de "(XII 4.1)"
\makeatother
\newcommand{\refvb}[2][VI${}_{\text{B}}$ ]{\sisi{#1\refnoexpose{#2}}{\ref{#2}}}

\newcommand{\leftexp}[2]{{\vphantom{#2}}^{#1}{#2}}
\newcommand{\AAA}{\mathbf{A}}
\newcommand{\Sf}{\mathfrak{S}}
\DeclareMathOperator*{\quotprojlim}{\text{``$\varprojlim$''}}
\DeclareMathOperator*{\quotinjlim}{\text{``$\varinjlim$''}}
\DeclareMathOperator{\Fl}{F\sisi{\ell}{l}}
\DeclareMathOperator{\Pl}{P\sisi{\ell}{l}}
%\DeclareMathOperator{\cl}{c\sisi{\ell}{l}}
\DeclareMathOperator{\tdf}{tdf}
\DeclareMathOperator{\rg}{rg}
\newcommand{\iInd}{\sisi{I}{i}nd}
\newcommand{\iind}{ind}
\newcommand{\ppro}{pro}
\newcommand{\tTor}{\sisi{T}{t}or}
\newcommand{\ttor}{tor}
\newcommand{\chap}{Chap.}
\newcommand{\Ffex}{\Ff\mathit{ex}}
\newcommand{\sheafk}{\underline{k}}
\newcommand{\sheafn}{\underline{n}}
\newcommand{\sheafGamma}{\underline{\Gamma}}

% Repeated references to an already typeset footnote must not use
% \footnotemark without a matching \footnotetext: Hyperref would then create
% a live annotation to a destination that cannot exist.  This helper keeps the
% current printed footnote mark and links it to an explicit stable note target.
\makeatletter
\newcommand{\sgarepeatedfootnotemark}[1]{%
  \begingroup
  \edef\@thefnmark{\thefootnote}%
  \hyperlink{#1}{\@makefnmark}%
  \endgroup}
\makeatother

\sisi{
\DeclareMathOperator{\simarrow}{\rule{.15mm}{1.8mm}\kern -3.4pt\rightsquigarrow}
\newcommand{\sheafell}{\underline{\ell}}
\newcommand{\sheafZZ}{\underline{\ZZ}}
}{
\newcommand{\nderef}[1]{la N.D.E.~(\ref{#1}) page~\pageref{#1}}
}

%%%%%%%%%%%%%%%%%
%%%%% expose XVII
%%%%%%%%%%%%%%%%%

% END INCLUDED SGA4_English_current_buildable_source_20260906/source/fichiers-configuration-latex/macros.tex


\externaldocument{../01/01_English.tex}
\externaldocument{../02/02_English.tex}
\externaldocument{../03/03_English.tex}
\externaldocument{../04/04_English.tex}
\externaldocument{../05/05_English.tex}
\externaldocument{../05bis/05bis_English.tex}
\externaldocument{../06/06_English.tex}
\externaldocument{../07/07_English.tex}
\externaldocument{../08/08_English.tex}

\setboolean{original}{false}

\begin{document}

\begin{center}
IX. Constructible sheaves; cohomology of an algebraic curve\\
M. Artin
\end{center}

\selectlanguage{english}version: \showgitstatus
\tableofcontents

\else
\chapter{Constructible sheaves; cohomology of an algebraic curve}
%\addtocontents{toc}{M. Artin}
\begin{center}
M. Artin
\end{center}
\fi
\phantomsection\MakeLinkTarget*{sga4.exp.IX}\label{sga4.exp.IX}

% BEGIN INCLUDED SGA4_English_current_buildable_source_20260906/source/09/drafts/SGA4_IX_LINES_43_278_English.texfrag
\setcounter{pageoriginal}{0}

\setcounter{section}{-1}
\section{Introduction}\etiquette[0]{IX.sec0}\pageoriginale

Note that, beginning with the present exposé, in contrast to the
preceding exposés and to the classical cohomological theory of
topological spaces, we are obliged, for the validity of the essential
statements, to restrict ourselves to \emph{torsion} sheaves (and often,
more precisely, to torsion sheaves prime to the residue
characteristics).

Let us point out that in fact there is no ``reasonable'' theory of
cohomology with integral (or even real) coefficients for varieties over
an algebraically closed field of characteristic $p\neq 0$, analogous
to the classical theory for the case $k=\CC$ (as Serre observed). More
precisely, there is no contravariant functor $H^1$, defined, say, on
the category of smooth projective schemes over an algebraically closed
field $K$ of characteristic $p>0$, with values in the category of
finite-dimensional vector spaces over $\RR$ (or a subfield of $\RR$),
``commuting with products,'' \ie such that $H^1(X\times
Y)\xleftarrow{\sim}H^1(X)\times H^1(Y)$, and such that $\dim H^1(X)=2$
if $X$ is a connected curve of genus 1. Indeed, it would follow that
if $X$ is an abelian variety over $K$, then the map $u\mapsto u^1$
from $\End(X)$ to $\End(H^1(X))$ is \emph{additive}, hence a
\emph{representation} of the opposite ring of $\End(X)$. But in
characteristic $p$ there exist elliptic\pageoriginale curves $X$ whose
endomorphism ring is a maximal order $A$ in a definite quaternion
algebra over $\ZZ$ ([3], p.198), and such an algebra plainly has no
representation in a vector space of dimension 2 over $\RR$, since
$A\otimes_{\ZZ}\RR$ is a division algebra (the quaternion algebra).

\section{The formalism of torsion sheaves}\etiquette[1]{IX.sec1}

Let $T$ be a topos and $F$ an abelian sheaf on $T$. One may multiply
by $n$ in $F$ ($n\in \ZZ$ an integer), this multiplication being
induced by the analogous operation in (Ab). Denote by ${}_nF$ the
kernel of this multiplication; thus ${}_nF(X)$, for $X\in \Ob T$, is
the group of sections of $F(X)$ whose order divides $n$.

We denote by $\PP$ the set of all prime numbers.

\begin{enonce*}{Definition 1.1}\etiquette[1.1]{IX.def1.1}
Let $p$ be a set of prime numbers. We say that $F$ is a $p$-torsion
sheaf, \sisi{where}{or} a $p$-sheaf, if the canonical morphism
$$
\varinjlim_n {}_nF\longrightarrow F\text{ is bijective}{\quoi},
$$
where $n$ ranges over the set of integers such that $\ass n\subset p$,
that is, such that the prime numbers dividing $n$ belong to $p$. If
$p=\PP$ is the set of all prime numbers, we simply say that $F$ is a
torsion sheaf.
\end{enonce*}

\begin{enonce*}{Proposition 1.2}\etiquette[1.2]{IX.prop1.2}
\begin{enumerate}
\renewcommand{\theenumi}{\roman{enumi}}
\renewcommand{\labelenumi}{(\theenumi)}
\item $F$ is a $p$-torsion sheaf if and only if $F$ is the sheaf
  associated with a presheaf taking values in $p$-torsion abelian
  groups; one may take
$$
P=\frac{\lim(\text{\rm presheaf})}{\ass n\subset p} {}_nF{\quoi}.
$$

\item If\pageoriginale $F$ is a $p$-torsion sheaf and if $X\in \Ob T$
  is a quasi-compact object\footnote[1]{An object $X$ of a topos is
  called quasi-compact if every covering family $\{X_i\rightarrow X\}$
  can be refined by a finite covering family. A morphism $u$ of
  topoi is quasi-compact if the inverse image of every quasi-compact
  object is again quasi-compact.}, then $F(X)$ is a $p$-torsion group.

\item If $T$ is locally of finite type {\rm (VI \hyperref[VI.def1.1]{1.1})}, then $F$ is a
  $p$-torsion sheaf if and only if $F(X)$ is a $p$-torsion group for
  every quasi-compact $X\in \Ob T$. In this case, one has
$$
\varinjlim_{\ass n\subset p} H^q(T/X; {}_nF)\xrightarrow{\sim}
H^q(T/X;F)
$$
for every quasi-compact $X$ and every $q$.

\item If $u:T\rightarrow T'$ is a morphism of topoi and if $F'$ is a
  $p$-torsion sheaf on $T'$, then the inverse image $u^{\ast}F'$ is a
  $p$-torsion sheaf on $T$.

\item If $u:T\rightarrow T'$ is a morphism of topoi, with $T$ and
  $T'$ locally of finite type and $u$ quasi-compact${}^\ast$, and if
  $F$ is a $p$-torsion sheaf on $T$, then the $R^qu_{\ast}F$ are
  $p$-torsion sheaves on $T'$ for every $q$.
\end{enumerate}
\end{enonce*}

\noindent
\emph{Proof}:
\begin{enumerate}
\renewcommand{\theenumi}{\roman{enumi}}
\renewcommand{\labelenumi}{(\theenumi)}
\item Clearly, if a sheaf $F$ is $p$-torsion, then $F$ is the sheaf
  associated with the presheaf $P$ whose group of sections $P(X)$ is
  the subgroup of $p$-torsion elements of $F(X)$. Conversely, let
  $F=\underline{a}P$ (II) be the sheaf associated with $P$,
  where $P$ takes values in the category of abelian
  $p$-torsion groups. From the exact sequence
$$
0\rightarrow {}_nP \rightarrow P\xrightarrow{n}P
$$
one deduces an exact sequence
$$
0\rightarrow \underline{a}({}_nP)\rightarrow F\xrightarrow{n}F{\quoi},
$$
whence\pageoriginale $\underline{a}({}_nP)={}_nF$. Since the functor $\underline{a}$
commutes with inductive limits, it follows from
$$
\frac{\lim(\text{\rm presheaf})}{\ass n\subset p}{}_nP\xrightarrow{\sim}P
$$
that one also has
$$
\varinjlim_{\ass n\subset p} {}_nF\xrightarrow{\sim}F{\quoi}.
$$

\item Put $P=\underrightarrow{\lim(\text{\rm presheaf})}{}_n F$. Then one has
  $P\subset F$, so $P$ is a separated presheaf, and consequently
$$
\underline{a}P(X)=\varinjlim (\ker(\prod_i P(X_i)\rightrightarrows)){\quoi},
$$
where, as usual, the limit is taken over the covering families
$\{X_i\rightarrow X\}$. But since $X$ is quasi-compact,
it is enough to take finite covering families. For these families,
$\prod_iP(X_i)$ is $p$-torsion; therefore so is the $\ker$. It follows
that $\underline{a}P(X)=F(X)$ is $p$-torsion.

\item To verify that $\varinjlim {}_nF\rightarrow F$ is bijective,
  it is enough to do so on the values at $X\in \Ob T$ for every
  $X$ in a family of generators. We are therefore reduced to (ii) for the
  first assertion. The second assertion is a consequence of

\item A consequence of (i), since $u^{\ast}(\underline{a}P')=\underline{a}(u\cdot P')$

\item An easy consequence of (iii).
\end{enumerate}

We shall also need a nonabelian variant. But there are technical
problems in the definition (I do not know whether they are serious),
and we shall therefore content ourselves with giving a definition for the
étale topology of a scheme.

\begin{enonce*}{Definition 1.3}\etiquette[1.3]{IX.def1.3}
Let\pageoriginale $p$ be a set of prime numbers. A group $G$ is called
an \emph{ind-p-group} if every finite subset of $G$
generates a finite subgroup of order $n$, with $\ass n\subset p$. This
is equivalent to saying that the finite subgroups $G_{\alpha}$
of order $n_{\alpha}$, with $\ass n_{\alpha}\subset p$, form a
filtered set, and that $G=\varinjlim G_{\alpha}$. If
$p=\PP$, one says \emph{ind-finite group} instead of ind-p-group.
\end{enonce*}

One immediately verifies the following
\begin{enonce*}{Proposition 1.4}\etiquette[1.4]{IX.prop1.4}
\begin{enumerate}
\renewcommand{\theenumi}{\roman{enumi}}
\renewcommand{\labelenumi}{(\theenumi)}
\item A subgroup as well as a quotient of an ind $p$-group is
  again an ind $p$-group.

\item A filtered inductive limit of ind-$p$-groups is an
  ind-$p$-group.

\item A finite projective limit of ind-p-groups is an ind-p-group.
\end{enumerate}
\end{enonce*}

\begin{enonce*}{Definition 1.5}\etiquette[1.5]{IX.def1.5}
Let $X$ be a scheme. A sheaf $F$ of groups on $X$ is called
a \emph{sheaf of ind $p$-groups} (\resp of \emph{ind-finite
  groups})
if, for every étale $U\rightarrow X$, with $U$ quasi-compact, $F(U)$ is an
ind $p$-group (\resp an ind-finite group).
\end{enonce*}

\begin{enonce*}{Proposition 1.6}\etiquette[1.6]{IX.prop1.6}
Let $F$ be a sheaf of groups on the scheme $X$.
\begin{enumerate}
\renewcommand{\theenumi}{\roman{enumi}}
\renewcommand{\labelenumi}{(\theenumi)}
\item $F$ is a sheaf of ind-$p$-groups if and only if, for
  every geometric point $\xi$ of $X$, the fiber $F_{\xi}$ is an
  ind-p-group.

\item If $f:X\rightarrow X'$ is a morphism and $F'$ is a
  sheaf of ind-p-groups on $X'$, then $f^{\ast}F'$ is a sheaf of
  ind-p-groups.

\item If $f:X\rightarrow X'$ is a quasi-compact morphism and $F$ is
  a sheaf of ind-p-groups on $X$, then $f_{\ast}F$ is a
  sheaf of ind-p-groups.
\end{enumerate}
\end{enonce*}

\noindent
\emph{Proof.}
(i) Suppose that $F$ is a sheaf of ind-p-groups, and let
  $\xi$ be a geometric point of $X$. Then $F_{\xi}=\varinjlim F(X')$, where
  $X'$ ranges over a\pageoriginale pseudo-filtered system of schemes
étale over $X$ (VIII \hyperref[VIII.prop3.9]{3.9}). It is evident that the affine $X'$ (hence
quasi-compact) form a cofinal set, and consequently that $F_{\xi}$ is an
inductive limit of ind $p$-groups, and hence is an ind-p-group by
\hyperref[IX.prop1.4]{1.4} (ii). Conversely, suppose that for every $\xi$ the fiber
$F_{\xi}$ is an ind-p-group, and let $U\rightarrow X$ be étale, with $U$
quasi-compact. Let $S\subset F(U)$ be a finite subset. For every
geometric point $\xi$ of $U$, the image of $S$ in $F_{\xi}$ generates
a finite group. Since a finite group is finitely presented, one
easily concludes that there exists a $U_{\xi}$ étale over $U$,
pointed at $\xi$, such that $S$ generates a finite group in $F(U_{\xi})$
(\cf VIII 4). A finite number of such $U_{\xi}$ form a covering
of $U$, say $\{U_1,\ldots,U_n\}$. Then the image of $S$ in $\prod
F(U_i)$ generates a finite subgroup, and since $F(U)\subset \prod
F(U_i)$, the result follows.

Assertion (ii) is trivial from (i), and (iii) is immediate.

% END INCLUDED SGA4_English_current_buildable_source_20260906/source/09/drafts/SGA4_IX_LINES_43_278_English.texfrag


% BEGIN INCLUDED SGA4_English_current_buildable_source_20260906/source/09/drafts/SGA4_IX_LINES_279_650_English.texfrag
\section{Constructible sheaves}\etiquette[2]{IX.sec2}

\setcounter{subsection}{-1}
\subsection{}\etiquette[2.0]{IX.subsec2.0}
Let $T$ be a topos. Recall that with every element $S\in \Ob (\Ens)$
one associates the sheaf $S_T$ associated with the presheaf $P$ such
that $P(X)=S$ for every $X\in \Ob (T)$. The object of $T$ representing
$S_T$ is
$$
\coprod_{s\in S}e= ``S\times e \mbox{''}{\quoi},
$$
where $e$ is the final object of $T$. The sheaf $S_T$ is called the
\emph{constant sheaf} with value $S$ (IV). In the evident way one
defines the notion of a \emph{constant sheaf of groups, constant sheaf
of $A$-modules} ($A$ a ring), and \emph{constant abelian sheaf}, with
value $M$. A morphism $f_T:S_T\rightarrow S'_T$ of constant sheaves is
said to be constant if it comes from a morphism of sets
$f:S\rightarrow S'$.

A sheaf $F$ is said to be \emph{locally constant} if there exists a
covering $\{e_i\rightarrow e\}$ of the final object such that $F$
becomes constant on each $e_i$. One\pageoriginale defines analogously
the notion of a locally constant \emph{sheaf of groups} or sheaf of
$A$-\emph{modules}, and of a \emph{locally constant morphism}
$f:F\rightarrow G$, when $F$ and $G$ are locally constant sheaves.
Finally, a locally constant sheaf of groups or of $A$-modules is said
to be \emph{of finite type} (\resp \emph{of finite presentation}) if
its local values are groups or $A$-modules of finite type (\resp of
finite presentation).

\begin{enonce*}{Lemma 2.1}\etiquette[2.1]{IX.lem2.1}
\begin{enumerate}
\renewcommand{\theenumi}{\roman{enumi}}
\renewcommand{\labelenumi}{(\theenumi)}
\item Let $f:F\rightarrow G$ be a morphism of locally constant sheaves
of sets on $T$, and suppose that the local values of $F$ are finite
sets. Then $f$ is locally constant.

\item Let $f:F\rightarrow G$ be a morphism of locally constant sheaves
of $A$-modules, with $F$ of finite type. Then $f$ is locally constant,
and the kernel and cokernel of $f$ are locally constant.

\item Let $0\rightarrow F'\rightarrow F\rightarrow F''\rightarrow 0$
be an exact sequence of sheaves of $A$-modules (\resp of groups), with
$F'$ and $F''$ locally constant, and $F''$ of finite presentation.
Then $F$ is locally constant.
\end{enumerate}
\end{enonce*}

Proof of (i) and (ii). We reduce without difficulty to the case in
which $F$ and $G$ are constant. Consider (ii): let $F=M_T$, $G=N_T$
($M$, $N$ being $A$-modules), and let $S\subset M$ be a finite set of
generators of $M$. Then $f$ is determined by its restriction to $S_T$,
and $f$ will be locally constant if $f|S_T$ is. Thus we are reduced to
(i) for the first assertion of (ii), and the second assertion is a
trivial consequence of the first.

It therefore remains to prove (i). Suppose that $F=S_T$, $G={S'}_T$
are constant, and also denote by $f$ the morphism of objects
$f:S\times e\rightarrow S'\times e$ (notation as above). Let
$f_s:e\rightarrow S'\times e$ $(s\in S)$ be the components of $f$:
since $S$ is finite, it plainly suffices to treat the case $f=f_s$,
that is, the case in which $S$ is a one-element set. Let
$X_{s'}\rightarrow e$ $(s'\in S')$ be the family of morphisms that make
the diagrams
$$
\xymatrix@C=2cm@R=1.5cm{
X_{s'} \ar[r] \ar[d]& e\ar[d]^{i_{s'}}\\
e   \ar[r]^{f}   & S'\times e
}
$$
cartesian, where\pageoriginale $i_{s'}:e\rightarrow S'\times e$ is
the ``inclusion'' into the $s'$-th component. The family $\{i_{s'}\}$
is trivially covering, and hence so is $\{X_{s'}\rightarrow e\}$. One
checks immediately (since coproducts in a topos are disjoint)
%% missing reference
that $f$ becomes constant on $X_{s'}$, and consequently $f$ is locally
constant.

Proof of (iii). Consider the case of sheaves of $A$-modules. We may
suppose $F'$ and $F''$ constant, with $F''$ having value $M''$, an
$A$-module of finite presentation. If $M''$ is free, and hence has a
finite basis $(e_i)_{1\leq i\leq n}$, then the $e_i$ define sections of
$F''$ and hence lift locally to sections of $F$; this shows that the
extension $F$ of $F''$ by $F'$ splits locally, which proves that $F$
is locally isomorphic to $F'\times F''$, and is therefore locally
constant. In the general case, choose a surjective homomorphism
$\varphi:L''\rightarrow M''$, with $L''$ free of finite type; the
kernel $R$ of $\varphi$ is of finite type, and setting
$L=\fprod{L''_T}{F}{M''_T}=\fprod{L''_T}{F}{F''}$, one sees on the one
hand that $L$ is an extension of $L''_T$ by $F'$, and is therefore
locally constant by what precedes, and on the other hand that there is
an exact sequence $0\rightarrow R_T\rightarrow L\rightarrow
F\rightarrow 0$, which, by (ii), implies that $F$ is locally constant.

From now on we restrict ourselves to the étale topoi of schemes.

\begin{enonce*}{Lemma 2.2}\etiquette[2.2]{IX.lem2.2}
Let $X$ be a scheme and let $F$ be a locally constant sheaf on $X$.
Then $F$ is represented by an étale $Y/X$. If, moreover, the fibers of
$F$ are finite (\resp and nonempty), then $Y$ is an étale covering
(\resp and is surjective).
\end{enonce*}

Proof. Descent (\SGA 1 IX 4.1 if $F$ has finite fibers, \SGA 3 X 5.4
in the general case).

\begin{enonce*}{Definition 2.3}\etiquette[2.3]{IX.def2.3}
A sheaf of sets (\resp of groups, \resp of $A$-modules) is said to be
constructible if, for every affine open $U\subset X$, there exists a
decomposition of $U$ as a union of finitely many constructible
{\rm (\EGA $0_{{\rm III}}$ 9.1.2)} reduced locally closed subschemes
$U_i$ such that the sheaf induced by $F$\pageoriginale on each $U_i$
is locally constant, with finite value (\resp finite, \resp of finite
presentation\footnote[1]{When $A$ is not assumed noetherian, for the
purposes of Homological Algebra it is preferable to require here,
instead of finite presentation alone for the $A$-module $M$, that it
have a left resolution by free $A$-modules of finite type. The notion
introduced in (\hyperref[IX.def2.3]{2.3}) should then be called:
$F$ is 1-constructible (more generally, one would define
$n$-constructibility in the evident way, for every integer $n\geq 0$).
For the present Exposé we shall provisionally retain the terminology in
the form (\hyperref[IX.def2.3]{2.3}), which is chiefly reasonable when
$A$ is noetherian, in which case it coincides with that just indicated.
Compare \SGA 6 I for finiteness notions in nonnoetherian cases.}).
\end{enonce*}

\setcounter{subsection}{3}
\subsubsection{}\etiquette[2.3.1]{IX.subsubsec2.3.1}
The finiteness hypothesis on the local values amounts to saying that
for every geometric point $\xi$, the fiber $F_{\xi}$ is finite (\resp
finite, \resp of finite presentation). It follows immediately from
\hyperref[IX.lem2.1]{2.1} (i) that a sheaf of groups $F$ is
constructible if and only if the underlying sheaf of sets of $F$ is
constructible.

Care must be taken that if $F$ is an abelian sheaf, it is constructible
as a sheaf of groups (or equivalently, as a sheaf of sets) if and only
if it is constructible as a sheaf of $\ZZ$-modules, and if
\emph{in addition} its fibers are finite. Thus, $\ZZ_X$ is
constructible as a $\ZZ$-module, but not as a sheaf of groups.

\begin{enonce*}{Proposition 2.4}\etiquette[2.4]{IX.prop2.4}
\begin{enumerate}
\renewcommand{\theenumi}{\roman{enumi}}
\renewcommand{\labelenumi}{(\theenumi)}
\item Let $X$ be quasi-compact and quasi-separated. Then a sheaf \\
(\resp sheaf of groups, \resp sheaf of $A$-modules) is constructible if
and only if there exists a decomposition of $X$ as a union of
constructible locally closed subsets such that $F$ becomes locally
constant and finite (\resp finite, \resp of finite presentation) on
each $X_i$.

\item To verify that $F$ is constructible, it is enough to verify that
$F| U_i$ is constructible for the $U_i$ of a given open covering of
$X$.

\item Let $f:Y\rightarrow X$ be a morphism. Then $f^{\ast}F$ is
constructible if $F$ is.

\item The\pageoriginale set of points at which the fiber of a
constructible sheaf is nonempty (\resp nonzero) is a locally
constructible subset {\rm (\EGA $0_{\rm III}$ 9.1.2)}.

\item Let $X$ be locally noetherian. Then a sheaf of sets
(\resp\ldots) $F$ on $X$ is constructible if and only if for every
$x\in X$ there exists a nonempty open subset $U$ of the closure $\bar{x}$
of $x$ such that $F$ induces a locally constant sheaf with finite value
(\resp\ldots) on $U$.
\end{enumerate}
\end{enonce*}

Proof. Plainly, if $f:Y\rightarrow X$ is a morphism and if there exists
a decomposition of $X$ as a union of constructible locally closed
subsets $X_i$ such that $F$ becomes locally constant on each $X_i$,
the same is true of $Y$ and $f^{\ast}F$ (\cf \EGA IV 1.8.2). This
proves the implication $\Leftarrow$ of (i). Now suppose that $X$ is
quasi-compact and quasi-separated, and let $\{U_1,\ldots,U_n\}$ be an
affine open covering of $X$ such that $F|U_i$ is constructible for
every $i$. To find a decomposition of $X$ as in (i), it is enough to
do so for $U_n$ and for $Y=X-U_n$ (with the induced reduced structure),
since $U_n$ and $Y$ are constructible in $X$ by virtue of the
hypothesis that ``$X$ is quasi-separated''. Now a decomposition exists
for
$U_n$ by the definition of constructibility, and $Y$ is the union of
$n-1$ affine opens $V_i=Y\cap U_i$ $(i=1,\ldots,n-1)$ with
$F| V_i$ constructible, whence result (i) follows by induction on $n$.
The same reasoning proves (ii); indeed, the hypothesis implies that
$F| U$ is constructible for every ``sufficiently small'' affine open
$U$, and an arbitrary affine open $V$ can be covered by finitely many
such opens. Assertion (ii) shows that the notion of constructibility is
local on $X$, and (iii), (iv) follow immediately. The implication
$\Rightarrow$ of (v) is immediate and does not depend on the
noetherian hypothesis. If $X$ is noetherian, one proves the implication
$\Leftarrow$ by the usual ``noetherian induction''.

\begin{enonce*}{Proposition 2.5}\etiquette[2.5]{IX.prop2.5}
Let $X$ be quasi-compact and quasi-separated, and let $F$ be a
constructible sheaf of groups or of $A$-modules on $X$. Then there
exists a finite filtration of $F$ whose successive quotients are of the
form $i_{!}G$, where $i:U\rightarrow X$ is the inclusion of a
constructible locally closed subset and where $G$ is locally constant
and constructible on $U$. If $X$ is noetherian, there exists such a
filtration with the $U$ irreducible.
\end{enonce*}

Proof.\pageoriginale By \hyperref[IX.prop2.4]{2.4} (i), there exists a
decomposition of $X$ as a finite union of constructible locally closed
subsets $X_i$ such that each $F|X_i$ is locally constant and finite
(\resp of finite presentation). Write $X_i=U_i\cap \complement V_i$,
where $U_i$ and $V_i$ are constructible opens, and let $N$ be the
number of opens of $X$ in the subtopology $T$ generated by
$\{U_i,V_i\}$. We argue by induction on $N$: if $W$ is a minimal
nonempty element of $T$, it is plain that $F|W$ is locally constant and
of finite presentation. Let $i:W\rightarrow X$ be the inclusion
morphism and let $Y=X-W$. We have the exact sequence
$$
0\longrightarrow i_{!}i^{\ast}F\longrightarrow F\longrightarrow F|
Y\longrightarrow 0{\quoi},
$$
and are thereby reduced to the same assertion for $Y$ and the sheaf
$F| Y$, where the induction hypothesis applies.

\begin{enonce*}{Proposition 2.6}\etiquette[2.6]{IX.prop2.6}
Suppose that $A$ is a noetherian ring.
\begin{itemize}
\item[(i)] A finite projective (\resp inductive) limit of constructible
sheaves of $A$-modules or sets is constructible. In particular, the
kernel, cokernel, and image of a morphism of constructible sheaves of
$A$-modules are constructible.

\item[(i bis)] A finite projective limit of constructible sheaves of
sets (\resp of groups) is constructible, and if $f:F\rightarrow G$ is
a morphism of constructible sheaves of groups, the kernel, image, and
cokernel (if $\Iim(f)$ is a normal subgroup) are constructible. A finite
inductive limit of constructible sheaves of sets is constructible.

\item[(ii)] Let $0\rightarrow M'\rightarrow M\rightarrow
M''\rightarrow 0$ be an exact sequence of sheaves of $A$-modules
(\resp of groups), with $M'$, $M''$ constructible. Then $M$ is
constructible.

\item[(iii)] Let $M_1\rightarrow M_2\rightarrow M_3\rightarrow
M_4\rightarrow M_5$ be an exact sequence of sheaves of $A$-modules.
Then if $M_i$ is constructible for $i=1,2,4,5$, the same holds for
$i=3$.
\end{itemize}
\end{enonce*}

\emph{Proof.} (i) We\pageoriginale leave the assertions for sets and
groups to the reader. For sheaves of $A$-modules, it is enough to prove
that the kernel and cokernel of a morphism $f:F\rightarrow G$ of
constructible sheaves of $A$-modules are constructible. The assertion
is local on $X$, and we may therefore suppose that $X$ is affine, hence
quasi-compact and quasi-separated. By \hyperref[IX.prop2.4]{2.4} (i),
we then reduce to the case in which $F$ and $G$ are locally constant,
and conclude by using \hyperref[IX.lem2.1]{2.1} (ii) and the noetherian
hypothesis on $A$. Assertion (ii) is proved analogously using
\hyperref[IX.lem2.1]{2.1} (iii). Assertion (iii) follows immediately
from (i) and (ii).

\begin{enonce*}{Proposition 2.7}\etiquette[2.7]{IX.prop2.7}
Let $X$ be a quasi-compact and quasi-separated scheme, and let $F$ be a
sheaf of sets (\resp of $A$-modules). In order that $F$ be
constructible, it is necessary and sufficient that it be isomorphic to
the cokernel of a pair of morphisms $H\rightrightarrows G$, where $H$
and $G$ are sheaves of sets representable by schemes étale and of
finite presentation over $X$ (\resp that $F$ be isomorphic to the
cokernel of a homomorphism $𝒜_{V,X}\rightarrow 𝒜_{U,X}$, where $U$ and
$V$ are two schemes étale and of finite presentation over $X$).
\end{enonce*}

%% A_V [and not 𝒜] in 2.9 \XXX

Sufficiency follows easily from \hyperref[IX.prop2.6]{2.6} (i bis).

Suppose that $F$ is a constructible sheaf of sets. Using the existence
of infinite coproducts in the étale site over $X$, one sees that one can
find an epimorphism $G\rightarrow F$, with $G$ representable by a
scheme étale over $X$, which may moreover be supposed to be a
coproduct of affine schemes $G_i$ $(i\in I)$, and hence separated over
$X$. For every finite subset $J$ of the index set $I$, let $G_J$ be
the coproduct of the $G_i$ for $i\in J$, and let $F_J$ be its image in
$F$. Since $G_J$ and $F$ are constructible, so is $F_J$
(\hyperref[IX.prop2.6]{2.6} (i bis)); hence the set $X_J$ of $x\in X$
such that the fibers of $F$ and $F_J$ at a geometric point of $X$ over
$x$ are equal is locally constructible (\hyperref[IX.prop2.4]{2.4}
(iv). Since the family of $X_J$ is increasing and has union $X$, and
since $X$ is quasi-compact, one of the $X_J$ equals $X$
(\EGA IV 1.9.9). This shows, after replacing $G$ by $G_J$, that we may
suppose $G$ affine, and hence separated over $X$, and quasi-compact over
$X$ since $X$ is quasi-separated (\EGA IV 1.2.4). Let then
$H=\fprod{G}{G}{F}$. It is a subsheaf of $\fprod{G}{G}{X}$, which is
representable, and is therefore itself representable (VIII \hyperref[VIII.prop6.1]{6.1}).
Since $H$ is moreover constructible (\hyperref[IX.prop2.6]{2.6}
(i bis)),\pageoriginale it follows again from the preceding argument
that it is quasi-compact, and therefore (being separated over $X$) of
finite presentation over $X$. This proves the first assertion of
\hyperref[IX.prop2.7]{2.7}, and the second is proved by the same method
(NB it is not necessary that $A$ be noetherian). Let us note in passing
that the proof also proves the following corollary:

\begin{enonce*}{Corollary 2.7.1}\etiquette[2.7.1]{IX.cor2.7.1}
Let $X$ be a scheme and let $F$ be a scheme étale over $X$. In order
that the corresponding étale sheaf on $X$ be constructible, it is
necessary and sufficient that $F$ be of finite presentation over $X$.
\end{enonce*}

Proposition \hyperref[IX.prop2.7]{2.7} will be useful here chiefly for
deducing certain limit-passage properties for constructible sheaves:

\begin{enonce*}{Corollary 2.7.2}\etiquette[2.7.2]{IX.cor2.7.2}
Let $X$ be a quasi-compact and quasi-separated scheme. Then every
sheaf of sets (\resp of ind-p-groups, \resp of $A$-modules) on $X$ is
the inductive limit of a filtered inductive system of constructible
sheaves of sets (\resp\ldots).
\end{enonce*}

If $F$ is a sheaf of sets, retain the notation of the proof of
\hyperref[IX.prop2.7]{2.7}, and let $L$ be the set of pairs $(J,H')$,
where $J$ is a finite subset of $I$ and $H'$ is a quasi-compact open
subset of $H_J=\fprod{G_J}{G_J}{F}$. Order $L$ in the evident way, and
observe that, since $F$ is the limit of the $F_J$, and each $F_J$ is
the limit of the $\Coker(H'\rightarrow G_J)$ for the quasi-compact
opens $H'$ of $H_J$, $F$ is the inductive limit of the inductive system
of the $F_{\ell}=\Coker(H'\rightarrow G_J)$ for
$\ell=(H',J)\in L$, proving the case of sets. The case of sheaves of
$A$-modules is treated in essentially the same way, beginning by
representing $F$ as the cokernel of a homomorphism
$𝒜_{V,X}\rightarrow 𝒜_{U,X}$, with $U$, $V$ schemes étale and
separated over $X$. The case of ind-$\sheaf{p}$-groups requires a little
more attention. Retaining the preceding notation, denote by $L(G_J)$
the free group sheaf generated by the sheaf of sets $G_J$, by $N_J$
the kernel of the homomorphism $L(G_J)\rightarrow F$ deduced from
$G_J\rightarrow F$, and by $L$ the set of pairs $(J,N')$, where $N'$
is a subsheaf of \emph{sets} of $N_J$ which is ``quasi-compact'', \ie
such that there exists an epimorphism $P\rightarrow N'$, with $P$
representable by a quasi-compact $X$-scheme. Order $L$ by declaring
that $(J,N')$ is smaller than $(J_1, N'_1)$ if $J\subset J_1$ and if
the canonical\pageoriginale homomorphism
$L(G_J)\rightarrow L(G_{J_1})$ maps $N'$ into $N'_1$. It is immediate
that $L$ is filtered, and that one obtains an inductive system of
sheaves of groups $F_{\ell}$, indexed by $L$, by taking, for
$\ell=(J,N')$, $F_{\ell}=$ the quotient sheaf of groups of $L(G_J)$ by
the invariant subsheaf of groups generated by $N'$. It is moreover
immediate that $F$ is isomorphic to the inductive limit of the
$F_{\ell}$, and it remains to prove that there exists a subset $L'$ of
$L$, cofinal in $L$, such that for $\ell\in L'$, $F_{\ell}$ is a
constructible sheaf of $\sheaf{p}$-groups. For this, it is enough to
prove that for every $J$ one can find a subsheaf of sets $N'$ of $N_J$
which is quasi-compact, contains a given quasi-compact subsheaf of sets
$N'_0$, and is such that the corresponding $F_{\ell}$ is a
constructible sheaf of $\sheaf{p}$-groups. Note first that it is enough
for this that the fibers of $F_{\ell}$ be finite $\sheaf{p}$-groups:
then $F_{\ell}$ will automatically be constructible, as the reader
will readily verify. On the other hand, the reader may also verify that
the set $X_{N'}$ of points $x$ of $X$ such that the fiber of
$F_{\ell}$ at a geometric point over $x$ is a finite
$\sheaf{p}$-group is locally constructible. Moreover, the subsets
$X_{N'}$ of $X$ are increasing functions of $N'$, and their union is
all of $X$: this last point follows from the fact that the inductive
limit of the $F_{\ell}$, for $\ell=(J,N')$ with $J$ fixed, is the
subsheaf of groups $\bar{F_J}$ of $F$ generated by the subsheaf of sets
$F_J$, which is quasi-compact, and hence $\bar{F_J}$ has fibers that
are finite $\sheaf{p}$-groups. On the other hand, in a free group on a
finite number of generators, every subgroup of finite index is a free group on a
finite number of generators; hence, if it is expressed as a filtered union of its
finite subsets, one of these generates the whole subgroup. It
follows readily that the union of the

% END INCLUDED SGA4_English_current_buildable_source_20260906/source/09/drafts/SGA4_IX_LINES_279_650_English.texfrag


% BEGIN INCLUDED SGA4_English_current_buildable_source_20260906/source/09/drafts/SGA4_IX_LINES_651_1136_English.texfrag
$X_{N'}$ is indeed $X$. Using (\EGA IV 1.9.9) once again, one finds
that one of the $X_{N'}$ is equal to $X$, which completes the proof of
\hyperref[IX.cor2.7.2]{2.7.2}.

\setcounter{subsection}{7}
\setcounter{subsubsection}{2}
\paragraph{}\etiquette[2.7.2.1]{IX.par2.7.2.1}
Note that the somewhat laborious proof in the case of group schemes
brings to light the lack of a convenient sorites for constructibility
conditions on sheaves of groups with fibers not necessarily finite; we
apologize to the reader, inviting them to fill this gap by their own
means if necessary, and point out by way of consolation that in
\hyperref[IX.prop2.9]{2.9} (iii) we obtain a simpler proof when $X$ is
Noetherian.

\begin{enonce*}{Corollary 2.7.3}\etiquette[2.7.3]{IX.cor2.7.3}
Let\pageoriginale $X$ be a quasi-compact and quasi-separated scheme,
and let $F$ be a constructible sheaf of sets (\resp of finite
ind-groups, \resp of $A$-modules). Then the functor $\Hom(F,G)$, where
$G$ is a variable sheaf of sets (\resp\ldots), commutes with inductive
limits in $G$. When $F$ is a sheaf of $A$-modules, with $A$
Noetherian, the functors $\Ext^{i}(X;F,G)$ in the $A$-module $G$ also
commute with inductive limits.
\end{enonce*}

This follows from \hyperref[IX.prop2.7]{2.7} and (VII \hyperref[VII.prop3.3]{3.3}) by the usual
arguments, which are left to the reader. Similarly, the conjunction of
\hyperref[IX.prop2.7]{2.7} and the results of (VII 5) easily gives:

\begin{enonce*}{Corollary 2.7.4}\etiquette[2.7.4]{IX.cor2.7.4}
Let $I$ be a filtered category, and let $i\mapsto
X_i:I^0\rightarrow (\Sch)$ be a functor which transforms arrows into
affine morphisms and objects into quasi-compact and quasi-separated
schemes. For every scheme $Y$, let $F_{c}(Y)$ be the category of
constructible sheaves of sets (\resp of ind-$\sheaf{p}$-groups, \resp
of $A$-modules) on $Y$. Then there is an equivalence of categories
$$
F_{c}(X)\xleftarrow{\approx}\varinjlim_i F_{c}(X_i){\quoi},
$$
where $X=\varprojlim_i X_i$ {\rm (VII 5)} and where the second member
denotes the $\varinjlim$ category of fibered categories, defined
in \hyperref[VI.def6.3]{VI 6.3}. In particular, every constructible sheaf of sets
(\resp\ldots) on $X$ is isomorphic to the inverse image of a sheaf of
the same kind on one of the $X_i$.
\end{enonce*}

\begin{enonce*}{Proposition 2.8}\etiquette[2.8]{IX.prop2.8}
Let $f:X\rightarrow Y$ be a surjective morphism locally of finite
presentation, and let $F$ be a sheaf of sets (\resp\ldots) on $Y$.
Then $F$ is constructible if and only if $f^{\ast}(F)$ is
constructible.
\end{enonce*}

We shall use the following

\begin{enonce*}{Lemma 2.8.1}\etiquette[2.8.1]{IX.lem2.8.1}
Let $f:X\rightarrow Y$ be a surjective morphism locally of finite
presentation, with $Y$ quasi-compact and quasi-separated. Then there
exists a finite partition of $Y$ into subschemes $Y_i$ of finite
presentation over $Y$ (in particular,\pageoriginale each $Y_i$ is
constructible in $Y$), and, for every $i$, finite surjective morphisms
$Y''_i\xrightarrow{g_i}Y'_i\xrightarrow{h_i}Y_i$, with $h_i$ étale
and $g_i$ locally free (\ie flat and of finite presentation, in
addition to the condition that $g_i$ be finite) and radicial, and
finally a $Y$-morphism $Y''_i\rightarrow X$
\footnote[1]{This assertion is also given in \EGA IV 17.16.4.}.
\end{enonce*}

Using the fact that $Y$ is quasi-compact and quasi-separated, we are
easily reduced to the case where $Y$ is quasi-affine, hence separated
(by using a finite affine open covering $(U_i)_{1\leq i\leq n}$ of $Y$
and considering the $Y_i=U_i-U_i\cap \bigcup_{j< i} U_j$), and then to
the case where $Y$ is affine (by the same argument). The scheme $X$ is
a union of affine open subsets $X_i$, and the image $T_i$ of $X_i$ in
$Y$ is constructible (\EGA IV 1.8.4); the union of the $T_i$ is $Y$
since $f$ is surjective, whence it follows that $Y$ is the union of a
finite number of the $T_i$ (\EGA IV 1.9.9), so that, replacing $X$ by
the sum scheme of the corresponding $X_i$, we may suppose that $X$ is
affine. The standard limit procedure (\EGA IV 8.1.2 c)) then allows
us to suppose in addition that $Y$ is Noetherian. The usual Noetherian
induction, together with the limit procedure (\EGA IV 8.1.2 a)),
reduces us to the case where $Y$ is the spectrum of a field $k$. Since
then $X\neq ∅$, there exists a closed point $x$ in $X$, corresponding
to a finite extension $k'$ of $k$, itself a radicial extension of a
separable extension $k'_s$ of $k$. We then take
$Y'=\Spec(k'_s)$, $Y''=\Spec(k')$, which completes the proof of
\hyperref[IX.lem2.8.1]{2.8.1}.

Necessity in \hyperref[IX.prop2.8]{2.8} being trivial
(\hyperref[IX.prop2.4]{2.4} (iii)), let us prove sufficiency. By the
preceding lemma we are reduced to the case where $f$ is of the form
$gh$, where $g$ and $h$ satisfy the conditions stated for $g_i$, $h_i$
in the lemma. We may therefore suppose successively either that $f$ is
finite radicial, a trivial case by (VIII \hyperref[VIII.thm1.1]{1.1}), or that $f$ is finite
étale. Note moreover that, using the hypothesis that $f^{\ast}(F)$ is
constructible and the definition of constructibility, we could suppose
(after replacing $X$ by an $X$-prescheme
$X'=\coprod_i X_i$ of finite presentation and with surjective
structural morphism) that $f^{\ast}(F)$ is already locally constant.
But when $f$ is étale and surjective, this implies that $F$ is locally
constant. Since in addition the hypotheses made on the fibers of
$f^{\ast}(F)$ remain valid for those of $F$, with $f$ surjective, it
indeed follows that $F$ is constructible, which completes the proof of
\hyperref[IX.prop2.8]{2.8}.

\begin{enonce*}{Proposition 2.9}\etiquette[2.9]{IX.prop2.9}
Let\pageoriginale $X$ be a Noetherian scheme, let $A$ be a Noetherian
ring, and let $p$ be a set of prime numbers.
\begin{enumerate}
\renewcommand{\theenumi}{\roman{enumi}}
\renewcommand{\labelenumi}{(\theenumi)}
\item The category of sheaves of sets (\resp of ind $p$-groups, \resp
of $A$-modules) on $X$ is locally Noetherian, that is, it has a set of
generators formed of Noetherian objects.

\item A sheaf $F$ of sets (\resp of ind $p$-groups, \resp of
$A$-modules) on $X$ is constructible if and only if it is Noetherian.
If $F$ is a sheaf of $A$-modules, then $F$ is constructible if and only
if it is a quotient of a finite sum of sheaves of the form $A_{U/X}$,
where $U\rightarrow X$ is étale and of finite type (notation of
{\rm \hyperref[IV.subsec2.4]{IV.2.4}}: $A_{U/X}=A_U$).

%% check the reference IV.2.4 [seems incorrect to me] \XXX

\item The constructible subsheaves of a sheaf $F$ of sets
(\resp\ldots) form an inductive system, and $F$ is the inductive limit
of these subsheaves.
\end{enumerate}
\end{enonce*}

Proof. We leave the assertion for sets to the reader. First consider
the case of a sheaf of $A$-modules. For (i), one must find a set of
generators whose elements are Noetherian objects. But the sheaves of
the form $A_{U/X}$, with $U\rightarrow X$ étale and of finite type,
form a set of generators which are easily seen to be constructible,
and it therefore suffices to prove the following lemma:

\begin{enonce*}{Lemma 2.10}\etiquette[2.10]{IX.lem2.10}
Let $X$ be Noetherian. Then every constructible $A$-module is
Noetherian.
\end{enonce*}

Proof of the lemma.

By Noetherian induction, we suppose the lemma true for every subscheme
$X'$ of $X$ distinct from $X$. Choose a nonempty open subset $X_0$ of
$X$ such that $F$ induces a locally constant sheaf $F_0$ on $X_0$.

Now let $\{F_i\}$, $i\in \NN$, be an increasing sequence of subsheaves
of $F_0$. Since, for a geometric point $P$ which is generic for an
irreducible component of $X_0$, the fiber of $F_0$ at $P$ is of finite
type, the sequence of fibers $(F_i)_P$\pageoriginale is stationary,
and we may suppose that it is constant. Now let $Q$ be a geometric
point of $X_0$ which is a \emph{specialization} of $P$ (VIII \hyperref[VIII.def7.2]{7.2}).
Since $F_0$ is locally constant on $X_0$, one sees immediately that the
specialization morphism $F_{0_Q}\rightarrow F_{0_P}$ (VIII.7) is
bijective, hence $(F_i)_Q\rightarrow (F_i)_P$ is \emph{injective} for
every $i$.

Let $V/X$ be étale and pointed at $P$, and let $s_1,\ldots,s_n$ be
sections $\in F_1(V)$ which generate $(F_1)_P$. Then the $s_i$ generate
$(F_i)_Q$ for every $i$ and every geometric point $Q$ which is a
specialization of $P$ above $V$. The sequence of sheaves $F_i$ is thus
constant in a neighborhood of the image of $P$, say in $X'\neq ∅$.
Let $Y=X-X'$; it therefore suffices to prove that the sequence of the
$F_i| Y$ is stationary, which follows from the Noetherian induction
hypothesis.

The assertions (ii) for $A$-modules are now immediate: since the
$A_{U/X}$ are Noetherian generators, one plainly has $F$ Noetherian
$\Leftrightarrow F$ is a quotient of a finite sum of sheaves
$A_{U/X}$. But $A_{U/X}$ is constructible, and consequently, in view
of \hyperref[IX.prop2.6]{2.6}, $F$ Noetherian $\Rightarrow F$
constructible; the converse implication having already been proved
(\hyperref[IX.lem2.10]{2.10}), this establishes (ii). Finally,
assertion (iii) for $A$-modules is a property of locally Noetherian
categories (\cf [1], Ch. IV).

Now consider the case of ind-$p$-groups. First note that a
constructible sheaf of groups $F$ is certainly Noetherian. Indeed,
since the fibers are finite, we may use the same argument as for
$A$-modules. We shall now prove assertion (iii) directly, which will
imply (i)---indeed, it will show that constructible sheaves form a set
of generators. Moreover, any Noetherian sheaf will then necessarily be
constructible, whence (ii).

To prove (iii), it plainly suffices to prove that a subsheaf (of
groups) $S$ of a sheaf $F$ of ind $p$-groups generated by a finite
number of sections $s_i\in F(U_i)$, with $U_i\rightarrow X$ étale and
of finite type, is constructible \footnote[1]{Recall the definition of
this sheaf: let $L_{U/X}$ be the sheaf representing the functor
$F\mapsto F(U)$ in the category of groups. The section
$s_i\in F(U_i)$ induces a morphism $L_{U_i/X}\rightarrow F$, and $S$ is
the image of the sum morphism $\amalg_i L_{U_i/X}\rightarrow F$. It is
the smallest subsheaf which ``contains'' the sections $s_i$.}. (Since
a\pageoriginale constructible sheaf is Noetherian, it is generated by
a finite number of sections.) By Noetherian induction, we may suppose
that this is true for every closed subscheme $Y$ of $X$ distinct from
$X$. Now the notion of a subsheaf generated by sections commutes with
the inverse image of sheaves; indeed, let $L_{U/X}$ be the sheaf which
represents the functor $F\mapsto F(U)$ in the category of groups. Then
if $f:Y\rightarrow X$, one has
$f^{\ast}(L_{U/X})=L_{\fprod{U}{Y}{X}/Y}$. Since $S$ is the image of
the sum of the $L_{U_i/X}$, and since $f^{\ast}$ commutes with sums, it
is clear that $f^{\ast}S$ is the subsheaf of $f^{\ast}F$ generated by
the sections $f^{\ast}(s_i)$.

We may therefore suppose that $S$ induces a constructible sheaf on
every closed subset $Y$ distinct from $X$, and it thus suffices
(\hyperref[IX.prop2.4]{2.4} (i)) to prove that $S$ induces a
constructible sheaf on a suitable nonempty open subset of $X$.
Replacing $X$ by a sufficiently small open subset, we may suppose that
every $U_i$ is an étale covering of $X$. It suffices to prove that $S$
is then locally constant, with finite fibers, on a suitable nonempty
open subset. This assertion is local on $X$ for the étale topology, so
we may suppose that every $U_i$ is completely split, say
$U_i=\amalg_{n_i}X$ (the sum of $n_i$ copies of $X$). Let $s_{ij}$
$(j=1,\ldots,n_i)$ be the sections of $F(X)$ such that
$s_i$ $(s_{i1},\ldots,s_{in_i})$. Plainly, $S$ is none other than the
subsheaf of $F$ generated by the $s_{ij}\in F(X)$. Since $F(X)$ is an
ind $p$-group (\hyperref[IX.def1.5]{1.5}), the $s_{ij}$ generate a
finite subgroup \hyperref[IX.def1.3]{1.3}. Thus we may suppose that the
set of the $s_i$ is a finite subgroup of $F(X)$; but then $S$ is
identical to the sheaf \emph{of sets} generated by the $s_i$, which is
constructible by \hyperref[IX.prop2.6]{2.6} (i) as a sheaf of sets,
hence as a sheaf of groups, \cqfd

In the following two propositions, we give criteria for a sheaf to be
locally constant or constructible, using specialization homomorphisms
(VIII \hyperref[VIII.subsec7.7]{7.7}).

\begin{enonce*}{Proposition 2.11}\etiquette[2.11]{IX.prop2.11}
Let $X$ be a scheme, and let $F$ be a constructible sheaf of sets
(\resp of ind-finite groups, \resp of $A$-modules); in the case of
sets,\pageoriginale suppose in addition that $F$ has finite fibers. Let
$x\in X$, and let $\bar{x}$ be a geometric point over $x$. In order that $F$
be locally constant in a neighborhood of $x$, it is necessary and
sufficient that, for every geometric point $\bar{x}'$ which is a
generization of $\bar{x}$, and every specialization arrow
$\bar{x'}\rightarrow \bar{x}$ {\rm (VIII \hyperref[VIII.def7.2]{7.2})}, the specialization
homomorphism {\rm (VIII \hyperref[VIII.subsec7.7]{7.7})} $F_{\bar{x}}\rightarrow F_{\bar{x'}}$ be
an isomorphism.
\end{enonce*}

The case of a sheaf of groups being a special case of that of a sheaf
of sets, we confine ourselves to the latter; the case of a sheaf of
modules is treated in essentially the same way. Necessity of the
condition being trivial, we need only establish sufficiency. Since the
fiber $I=F_{\bar{x}}$ is finite, after replacing $X$ by a suitable
scheme $X'$ étale over $X$, we may suppose that one can find a
homomorphism $u:I_X\rightarrow F$ from the constant sheaf $I_X$ to
$F$, inducing an isomorphism on the fibers at $\bar{x}$. Using the
constructibility of both sheaves and \hyperref[IX.prop2.6]{2.6}, one
easily sees that the set $Z$ of points $z$ of $X$ such that the
homomorphism induced on the fibers at a geometric point $\bar{z}$ over
$z$ is bijective is a locally constructible subset of $X$. The
hypothesis immediately implies that it contains the generizations of
$x$, hence (\EGA IV 1.10.1) it is a neighborhood of $x$, which proves
that $F$ is locally constant (indeed, even constant in this case) in a
neighborhood of $x$,\cqfd

\begin{enonce*}{Definition 2.12}\etiquette[2.12]{IX.def2.12}
\footnote[1]{The terminology is due to S. Lubkin.}: A scheme $X$ is
said to be path-connected if for every pair $P$, $Q$ of geometric
points of $X$, there exist geometric points $P=P_0,\ldots,P_n$;
$Q_1,\ldots,Q_n=Q$ and specializations {\rm (VIII \hyperref[VIII.def7.2]{7.2})}
$P_i\rightarrow Q_i$ and $P_{i-1}\rightarrow Q_i$
$(i=1,\ldots,n)$. One says that $X$ is locally path-connected (for the
étale topology) if for every étale $U\rightarrow X$ there exists an
étale covering $\{U_i\rightarrow U\}$ of $U$ such that $U_i$ is
path-connected for every $i$.
\end{enonce*}

One checks immediately that a locally Noetherian prescheme is locally
path-connected.

\begin{enonce*}{Proposition 2.13}\etiquette[2.13]{IX.prop2.13}
\begin{enumerate}
\renewcommand{\theenumi}{\roman{enumi}}
\renewcommand{\labelenumi}{(\theenumi)}
\item Let\pageoriginale $X$ be locally path-connected and let $F$ be a
sheaf of sets (\resp of groups, \resp of $A$-modules) on $X$. Suppose
that the fibers of $F$ are finite (\resp finite, \resp of finite
presentation). Then $F$ is locally constant if and only if, for every
specialization $P\rightarrow Q$ of geometric points, the specialization
morphism {\rm (VIII \hyperref[VIII.subsec7.7]{7.7})} $F_Q\rightarrow F_P$ is bijective.

\item Suppose that $X$ is locally Noetherian and that the fibers of $F$
are finite (\resp finite, \resp of finite presentation). Then, if for
every specialization $P\rightarrow Q$ the specialization morphism
$F_Q\rightarrow F_P$ is injective, $F$ is constructible.

\item Suppose that $X$ is locally Noetherian and that the fibers of the
sheaf of sets $F$ are finite. Let $c:X\rightarrow \NN$ be the function
which associates with $x\in X$ the number of elements of the fiber
$F_{\bar{x}}$ of $F$ at a geometric point $\bar{x}$ above $x$. Then
$F$ is constructible if and only if $c$ is a constructible function,
\ie if and only if $c^{-1}(n)$ is constructible for every
$n\in \NN$.
\end{enumerate}
\end{enonce*}

\emph{Proof.}
\begin{enumerate}
\renewcommand{\theenumi}{\roman{enumi}}
\renewcommand{\labelenumi}{(\theenumi)}
\item Let $P$ be a geometric point of $X$. We shall find an étale
neighborhood of $P$, that is, an étale $U\rightarrow X$ pointed at
$P$, such that $F$ becomes constant on $U$. Let $V'$ be an étale
neighborhood of $P$ such that $F(V')\rightarrow F_P$ is surjective.
Such a $V'$ exists by the finiteness hypothesis. Let $U\rightarrow V'$
be an étale neighborhood of $P$ in $V'$ such that $U$ is
path-connected. Then $F(U)\rightarrow F_P$ is still surjective. Choose
a subset $S\subset F(U)$ which maps bijectively onto $F_P$. If $F$ is
a sheaf of groups (\resp of $A$-modules), one easily sees that, after
changing $U$ if necessary, one may find such a set which is moreover a
subgroup (\resp submodule) of $F(U)$. This gives a morphism
$S_U\rightarrow F| U$, where $S_U$ is the constant sheaf with value
$S$. One has
$$
(S_U)_P\xrightarrow{\sim}F_P{\quoi}.
$$
Since\pageoriginale every specialization morphism on the fibers of $F$
(\resp $S_U$) is bijective, and since $U$ is path-connected, it follows
that $S_U\xrightarrow{\sim}F| U$, whence the result.

\item The assertion is local, and we may therefore suppose that $X$ is
Noetherian. Let $P$ be a geometric point above a maximal point of $X$.
There exists an irreducible étale neighborhood, hence a path-connected
one, $V\rightarrow X$ of $P$ such that $F(V)$ maps surjectively onto
$F_P$. It then maps surjectively onto $F_Q$ for every geometric point
$Q$ of $V$, or, what amounts to the same thing, for every geometric
point of the image $U$ of $V$ in $X$, since such a point is a
specialization of $P$ and $F_Q\rightarrow F_P$ is injective. Thus every
specialization morphism in $U$ is bijective, hence $F| U$ is locally
constant. By Noetherian induction we may suppose that $F| X-U$ is
constructible, whence the result.

\item It plainly suffices to treat the case where the function $c$ is
constant. Let $V\rightarrow X$ be an étale neighborhood of a geometric
point $P$ above a maximal point of $X$, such that $F(V)$ maps
surjectively onto $F_P$. Let $S$ be a subset of $F(V)$ such that
$S\xrightarrow{\sim}F_P$. Then, for the constant sheaf $S_V$ and every
specialization $P\rightarrow Q$ of geometric points, one has a
commutative diagram
$$
\xymatrix@C=2cm@R=1.25cm{
(S_V)_P   \ar[r]^{\sim}& F_P\\
(S_V)_Q   \ar[r] \ar[u]^{\text{\rotatebox{90}{$\backsim$}}}  & F_Q. \ar[u]
}
$$
Since $\lvert F_Q\rvert=\lvert F_P\rvert$, the arrow $F_Q\rightarrow F_P$ is bijective, so
$F$ is locally constant in a neighborhood of $P$ by (i). By
Noetherian induction, we are done.
\end{enumerate}

The following result (ii) will be of technical use later; it will allow
us to reduce certain verifications to the case of a constant sheaf.

\begin{enonce*}{Proposition 2.14}\etiquette[2.14]{IX.prop2.14}
Let\pageoriginale $X$ be a scheme.
\begin{enumerate}
\renewcommand{\theenumi}{\roman{enumi}}
\renewcommand{\labelenumi}{(\theenumi)}
\item For every finite morphism of finite presentation
$f:X'\rightarrow X$, and every constructible sheaf of sets (\resp of
ind-finite groups, \resp of $A$-modules) on $X'$, $f_{\ast}(F)$ is
constructible.

\item Suppose that $X$ is quasi-compact and quasi-separated, and let
$F$ be a constructible sheaf of sets (\resp\ldots) on $X$. Then one can
find a finite family $(p_i:X'_i\rightarrow X_i)$ of finite morphisms,
for every $i$ a constant constructible sheaf of sets (\resp\ldots)
$C_i$ on $X'_i$, and a monomorphism
$$
F\hookrightarrow \prod_i p_{i\ast}(C_i){\quoi}.
$$
\end{enumerate}
\end{enonce*}

First prove (i). By definition, we may suppose that $X$ is affine, and
then, using \hyperref[IX.cor2.7.4]{2.7.4}, the standard limit procedure
reduces us to the case where $X$ is Noetherian. A further passage to
the limit, again using \hyperref[IX.cor2.7.4]{2.7.4} and also
\hyperref[IX.prop2.4]{2.4}. (v), reduces us to the case where $X$ is the
spectrum of a field. The conclusion then reduces to the corresponding
assertion concerning the fiber of $f_{\ast}(F)$ at a geometric point
over $X$, which follows immediately from (VIII \hyperref[VIII.prop5.5]{5.5}, \hyperref[VIII.rem5.8]{5.8}).

Let us prove (ii). By \hyperref[IX.prop2.4]{2.4} (i), there exists a
partition of $X$ into subschemes $Y_i$, with inclusion morphisms
$u_i:Y_i\rightarrow X$ of finite presentation, such that the
restriction $F_i$ of $F$ to every $Y_i$ is a locally constant sheaf. It
is then immediate that the canonical homomorphisms $F\rightarrow
u_{i\ast}(F_i)$ define a homomorphism $F\rightarrow \prod_i
u_{i\ast}(F_i)$ which is a monomorphism.

In the preceding construction, one could even arrange that every $F_i$
be not only locally constant, but also \emph{isotrivial,} \ie that there
exist a surjective \emph{finite} étale morphism
$q_i:Y'_i\rightarrow Y_i$ such that $q^{\ast}_i(F_i)$ is a constant
sheaf on $Y'_i$, say $C_{iY'_i}$. This follows easily, for example, from
the definition of \hyperref[IX.lem2.8.1]{2.8.1} and from (VIII \hyperref[VIII.thm1.1]{1.1}). On
the other hand, it\pageoriginale follows from the ``Main Theorem'' in
the form (\EGA IV 8.12.6) that there exists a finite morphism
$p_i:Z_i\rightarrow X$ and an open immersion
$v_i:Y'_i\rightarrow Z_i$ of $X$-schemes.

Since $F_i$ embeds into
$q_{i\ast}(q^{\ast}_i(F_i))=q_{i\ast}(C_{iY'_i})$, $F$ embeds into the
product of the
$u_{i\ast}(q_{i\ast}(C_{iY'_i}))=p_{i\ast}(v_{i\ast}(C_{iY'_i}))$. If
every $Z_i$ is normal, then Lemma (\hyperref[IX.lem2.14.1]{2.14.1})
below implies that $v_{i\ast}(C_{iY'_i})\simeq C_{iZ_i}$, hence $F$
embeds into the product of the $p_{i\ast}(C_{iZ_i})$, and we are done.
In the general case, introduce the normalization $\bar{Z}_i$ of $Z_i$,
the inverse image $\bar{Y}'_i$ of $Y'_i$ in $\bar{Z}_i$, and the
immersion $\bar{v}_i:\bar{Y}'_i\rightarrow \bar{Z}_i$. One sees
immediately that $\bar{Y}'_i$ contains the maximal points of
$\bar{Z}_i$, and is therefore dense in $\bar{Z}_i$. If
$r_i:\bar{Z}_i\rightarrow Z_i$ and $s_i:\bar{Y}'_i\rightarrow Y'_i$
are the projections, since the latter is surjective, it follows that
$C_{iY'_i}$ embeds into $s_{i\ast}(C_{i\bar{Y}'_i})$, hence
$v_{i\ast}(C_{iY'_i})$ embeds into
$v_{i\ast}(s_{i\ast}(C_{i\bar{Y}'_i}))=r_{i\ast}(\bar{v}_{i\ast}(C_{i\bar{Y}'_i}))$,
which is isomorphic to $r_{i\ast}(C_{i\bar{Z}_i})$ by
(\hyperref[IX.lem2.14.1]{2.14.1}) below. Thus $F$ embeds into the
product of the
$p_{i\ast}(r_{i\ast}(C_{\bar{Z}_i}))=t_{i\ast}(C_{i\bar{Z}_i})$, where
$t_i=p_ir_i$ is the canonical projection $\bar{Z}_i\rightarrow X$.
When this morphism is finite for every $i$, we are done. In the general
case, one can only say that $\bar{Z}_i$ is \emph{integral} over $X$,
hence of the form $\sheaf{\Spec}(𝒜_i)$, where $𝒜_i$ is a
quasi-coherent sheaf of $𝒪_X$-algebras which is integral over $𝒪_X$.
Using (\EGA IV 1.7.7), one sees that $𝒜_i$ is the inductive limit of
its subalgebras $𝒜_{i,j}$ finite over $X$, hence $\bar{Z}_i$ is the
projective limit (VII 5) of the preschemes
$Z_{ij}=\Spec(A_{ij})$, finite over $X$. Applying (VII \hyperref[VII.cor5.11]{5.11}), one finds
that $t_{i\ast}(C_{i\bar{Z}_i})$ is the inductive limit of the
$t_{ij^{\ast}}(C_{iZ_{ij}})$, where $t_{ij}:Z_{ij}\rightarrow X$ is the
canonical projection. By (\hyperref[IX.cor2.7.3]{2.7.3}), the
homomorphism $F\rightarrow t_{i\ast}(C_{i\bar{Z}_i})$ therefore factors
through a homomorphism $F\rightarrow t_{ij\ast}(C_{iZ_{ij}})$ for a
suitable $j=j(i)$. We now set $X'_i=Z_{i,j(i)}$, and obtain an embedding
as stated in \hyperref[IX.prop2.14]{2.14}. It remains to prove the
following

\begin{enonce*}{Lemma 2.14.1}\etiquette[2.14.1]{IX.lem2.14.1}
Let $X$ be a geometrically unibranch scheme (for example, a normal
scheme), and let $i:U\rightarrow X$ be a quasi-compact open immersion.
Then the closure $Z$\pageoriginale of $U$, endowed with the induced
reduced structure, is also geometrically unibranch; more precisely, for
every $z\in Z$, one has
$𝒪_{Z,z}\xleftarrow{\sim}𝒪_{X_{\red},z}$. Suppose on the other hand
that $i$ is dominant, and let $C$ be a set, with $C_U$ the constant
sheaf on $U$ defined by $C$. Then the natural homomorphism
$C_X\rightarrow i_{\ast}(C_U)$ is an isomorphism. The same conclusions
hold if the open immersion $i$ is replaced by an inclusion
$i:x\rightarrow X$, where $x$ is a maximal point of $X$.
\end{enonce*}

Using (\EGA IV 2.3.11), we are reduced for the first assertion to the
case where $X$ is local, with closed point $z$; but then
$X_{\text{ réd}}$ is integral and $U$ is nonempty, hence
$Z=X_{\text{ réd}}$, and the assertion is evident (thus, instead of
assuming $X$ geometrically unibranch, it would have sufficed here to
suppose that the local rings of $X_{\text{réd}}$ are integral
domains). For the second assertion, looking fiber by fiber, we are
reduced by (VIII \hyperref[VIII.rem5.3]{5.3}) \loccit to the case where $X$ is strictly local,
and then to proving that $H^0(X,C_X)\rightarrow H^0(U,C_U)$ is
bijective. Here again, $X$ is irreducible, and $U$ is a nonempty open
subset of it and hence irreducible, and the two sides of the preceding
map are identified with $C$, whence the desired conclusion. (Here the
notion ``geometrically unibranch'' has been used in the form of
(\EGA IV 18.8.14), as meaning that the strict localizations of the
rings of $X_{\text{ réd}}$ are integral domains.) The case of the
inclusion of a maximal point (stated here for ease of future reference)
is treated in exactly the same way.

\begin{enonce*}[remark]{Remarks 2.14.2}\etiquette[2.14.2]{IX.rems2.14.2}
When $X$ is Noetherian, the proof given (or an immediate deduction)
shows that in \hyperref[IX.prop2.14]{2.14} (ii), the $X'_i$ may be
assumed integral; if in addition $X$ is assumed universally Japanese,
\ie the rings of its affine open subsets are universally Japanese
(\EGA ${\rm 0_{IV}}$ 23.1.1, IV 7.7.2), the $X'_i$ may moreover be
assumed normal.
\end{enonce*}

% END INCLUDED SGA4_English_current_buildable_source_20260906/source/09/drafts/SGA4_IX_LINES_651_1136_English.texfrag


% BEGIN INCLUDED SGA4_English_current_buildable_source_20260906/source/09/drafts/SGA4_IX_LINES_1137_1339_English.texfrag
\section{Kummer and Artin--Schreier theories}\etiquette[3]{IX.sec3}\pageoriginale

\setcounter{subsection}{-1}
\subsection{}\etiquette[3.0]{IX.subsec3.0}
We denote by $(\GG_m)_X$ the sheaf represented by the multiplicative
group $\Spec\Oo_X[t,t^{-1}]$ over $X$. Thus, for $U\to X$ étale, the
sections of $(\GG_m)_X$ over $U$ are the invertible elements
$\Gamma(U,\Oo_U^*)$ of $\Gamma(U,\Oo_U)$. Let $n\in\NN$. The kernel
of the $n$-th power map on $(\GG_m)_X$ is the \lq\lq sheaf of $n$-th
roots of unity'', denoted $(𝛍_n)_X$. Thus there is an exact sequence
\begin{equation*}
0\to(𝛍_n)_X\to(\GG_m)_X\xrightarrow{n}(\GG_m)_X{\quoi}.
\tag{3.1}
\end{equation*}\etiquette[3.1]{IX.eq3.1}
If $n$ is invertible on $X$ (that is, if $n$ is prime to every
residue characteristic of $X$), the $n$-th power map is surjective
\emph{as a morphism of sheaves}:

\begin{enonce*}{Kummer theory 3.2}\etiquette[3.2]{IX.Thr3.2}
If $n\in\NN$ is invertible on $X$, there is an exact sequence
$$
0\to(𝛍_n)_X\to(\GG_m)_X\xrightarrow{n}(\GG_m)_X\to0.
$$
\end{enonce*}

Indeed, let $u\in(\GG_m)_X(U)$, with $U\to X$ étale. We must find an
étale covering $\{U_i\to U\}$ such that $u$ becomes an $n$-th power
on every $U_i$. Since $n$ is invertible on $U$, the equation
$$
T^n-u=0
$$
is \lq\lq separable'' over $U$, that is,
$$
U'=\Spec\Oo_U[T]/(T^n-u)
$$
is étale over $U$. Since $U'\to U$ is surjective, it is a covering,
and $u$ becomes an $n$-th power on $U'$.

We point out\pageoriginale that when $n$ is invertible on $X$,
$(𝛍_n)_X$ is a locally constant sheaf, in fact locally isomorphic to
the sheaf $(\ZZ/n\ZZ)_X$, and is constant in important cases. Kummer
theory therefore gives information about the cohomology of $X$ with
values in certain constant sheaves. By definition,
$$
H^0(X,(\GG_m)_X)=\Gamma(X,\Oo_X^*){\quoi},
$$
and in dimension one we have:

\begin{enonce*}{Theorem 3.3}\etiquette[3.3]{IX.thm3.3}
(\lq\lq Hilbert's Theorem 90'') There is an isomorphism
$$
H^1(X,(\GG_m)_X)\xrightarrow{\sim}\Pic X{\quoi},
$$
where $\Pic X$ is the group of isomorphism classes of invertible
sheaves on $X$.
\end{enonce*}

\emph{Proof.} This says that the canonical morphism
$$
H^1(X_{\zar},(\GG_m)_X)\longrightarrow
H^1(X_{\text{\rm ét}},(\GG_m)_X)
$$
is bijective. Equivalently,
$$
R^1\epsilon_*(\GG_m)_X=0,
$$
where $\epsilon_*$ is direct image for the evident morphism of sites
$\epsilon:X_{\et}\to X_{\zar}$ (VIII 4). We are therefore reduced to
the case where $X$ is affine. Since $H^1$ may be computed by
\v{C}ech cohomology (V \hyperref[V.subsec2.5]{2.5}), and since for quasi-compact $X$ it is
enough to use quasi-compact coverings, this follows immediately from
descent theory for sheaves (see FGA 190 or \SGA VIII 1).

\setcounter{subsection}{3}
\setcounter{subsubsection}{0}
\subsubsection{}\etiquette[3.3.1]{IX.subsubsec3.3.1}
Suppose that $X$ is Noetherian and has no embedded component, and let
$x_j$ be the maximal points of $X$. Let $R_j$ be the local Artinian
ring of $X$ at $x_j$, and let $i_j:\Spec R_j\to X$ be the inclusion.
There is\pageoriginale a canonical injection
$$
(\GG_m)_X\longrightarrow
\prod_j i_{j^{\ast}}(\GG_m)_{\Spec R_j}
$$
of sheaves on $X_{\text{\rm ét}}$ (respectively $X_{\zar}$). Its
cokernel $D$ is called the \emph{sheaf of Cartier divisors} on
$X_{\text{\rm ét}}$ (respectively $X_{\zar}$)\footnote[1]{In
\EGA IV 21, the word \lq\lq divisors'' is used instead of
\lq\lq Cartier divisors''.}.

\begin{enonce*}{Corollary 3.4}\etiquette[3.4]{IX.cor3.4}
Let $X$ be Noetherian and without embedded components, let
$\epsilon:X_{\text{\rm ét}}\to X_{\zar}$ be the canonical morphism of
sites, and let $D$ be the sheaf of Cartier divisors on
$X_{\text{\rm ét}}$. Then $\epsilon_*D$ is the sheaf $D_{\zar}$ of
Cartier divisors on $X_{\zar}$, and in particular
$$
H^0(X_{\text{\rm ét}},D)\simeq H^0(X_{\zar},D_{\zar}){\quoi}.
$$
\end{enonce*}

\emph{Proof.} Applying $\epsilon_*$ to the exact sequence
$$
0\to(\GG_m)_X\longrightarrow
\prod_j i_{j^{\ast}}(\GG_m)_{\Spec R_j}
\longrightarrow D\to0
$$
gives an exact sequence, since
$R^1\epsilon_*(\GG_m)_X=0$ (\hyperref[IX.thm3.3]{3.3}); the first
assertion follows.

When $X$ has characteristic $p\neq0$, the following result sometimes
replaces \hyperref[IX.Thr3.2]{3.2} in the study of $p$-torsion
coefficients:

\begin{enonce*}{Artin--Schreier theory 3.5}\etiquette[3.5]{IX.Thr3.5}
If $X$ has characteristic $p>0$, there is an exact sequence
$$
0\to(\ZZ/p\ZZ)_X\to\Oo_X\xrightarrow{\mathfrak p}\Oo_X\to0{\quoi},
$$
where $\mathfrak p$ is the morphism of additive groups defined by
$\mathfrak p(f)=f^p-f$.
\end{enonce*}

The kernel of $\mathfrak p$ is the sheaf of sections of $\Oo_X$ which
locally belong to the prime field, and is therefore a constant sheaf.
The morphism $\mathfrak p$ is\pageoriginale surjective for the étale
topology because the equation $T^p-T-a$ ($a\in\Gamma(U,\Oo_U)$) is
always \lq\lq separable'' in characteristic $p$, that is, defines an
étale covering of $U$. In applying this exact sequence, we shall of
course use VII \hyperref[VII.subsec4.2]{4.2}.

\begin{enonce*}{Proposition 3.6}\etiquette[3.6]{IX.prop3.6}
\begin{enumerate}
\renewcommand{\theenumi}{\roman{enumi}}
\renewcommand{\labelenumi}{(\theenumi)}
\item Let $X$ be a scheme and let $i:x\to X$ be the inclusion of a
point. Then
$$
H^1(X,i_*(\ZZ)_x)=0,
$$
where $(\ZZ)_x$ is the \lq\lq constant sheaf'' of value $\ZZ$
(see \hyperref[IX.subsec2.0]{2.0}).

\item If $X$ is irreducible and if, for every geometric point
$\bar x$ of $X$, the strict localization $X_{\bar x}$ of $X$ at
$\bar x$ is irreducible (one should say that \lq\lq $X$ is
geometrically unibranch'', see \EGA IV 18.8.14), then
$$
H^1(X,(\ZZ)_X)=0{\quoi}.
$$
\end{enumerate}
\end{enonce*}

\emph{Proof.}
\begin{enumerate}
\renewcommand{\theenumi}{\roman{enumi}}
\renewcommand{\labelenumi}{(\theenumi)}
\item The Leray spectral sequence
$$
E^{pq}_2=H^p(X,R^qi_*(\ZZ)_x)\Longrightarrow
H^{p+q}(x,(\ZZ)_x)
$$
gives an injection
$H^1(X,i_*(\ZZ)_x)\to H^1(x,(\ZZ)_x)$. Since
$x=\Spec k(x)$ is the spectrum of a field, it is well known that
$H^1(x,(\ZZ)_x)=0$. (This follows from the fact that
$H^q(x,F)$ is torsion for $q>0$ in any case (see VII \hyperref[VII.rem2.1]{2.2} and CG),
and from the exact sequences
$0\to\ZZ\xrightarrow{n}\ZZ\to\ZZ/n\to0$.) This proves (i).

\item Let $i:x\to X$ be the inclusion of the generic point of $X$.
Under the hypotheses of (ii), the canonical morphism
$(\ZZ)_X\to i_*(\ZZ)_x$ is immediately seen to be bijective (and the
converse is also true); thus (ii) is reduced to (i).
\end{enumerate}

\begin{enonce*}[remark]{Remarks 3.7}\etiquette[3.7]{IX.rems3.7}
In \hyperref[IX.prop3.6]{3.6} (i) and (ii), $\ZZ$ may be replaced by
any torsion-free abelian group. On the other hand, (ii) becomes false
if the hypothesis that $X$ is geometrically unibranch is omitted, as
one sees, for example, by taking $X$ to be an algebraic curve over a
field $K$ having an ordinary double point (compute $H^1(X,\ZZ)$ in
this case!).
\end{enonce*}


% END INCLUDED SGA4_English_current_buildable_source_20260906/source/09/drafts/SGA4_IX_LINES_1137_1339_English.texfrag


% BEGIN INCLUDED SGA4_English_current_buildable_source_20260906/source/09/drafts/SGA4_IX_LINES_1340_1459_English.texfrag
\section{The case of an algebraic curve}\etiquette[4]{IX.sec4}

\setcounter{subsection}{-1}
\subsection{}\etiquette[4.0]{IX.subsec4.0}
Let $X$ be a Noetherian scheme of dimension $1$. Let $R(X)$ be the
ring of rational functions on $X$, and let
$i:\Spec R(X)\to X$ be the inclusion. The ring $R(X)$ is Artinian,
the product of the local rings of $X$ at its maximal points, and the
results of VIII 2 may be applied.

Let $F$ be an abelian sheaf on $\Spec R(X)$, and consider the sheaves
$R^qi_*F$, $q>0$, on $X$. It follows at once from VIII \hyperref[VIII.thm5.2]{5.2} that the
fiber of $R^qi_*F$ at a geometric point lying over a maximal point of
$X$ vanishes. Since $X$ has dimension $1$, such a sheaf is of a very
special kind:

\begin{enonce*}{Lemma 4.1}\etiquette[4.1]{IX.lem4.1}
Let $X$ be Noetherian and let $F$ be a sheaf on $X$. The following
conditions are equivalent (a sheaf satisfying them will be called a
\lq\lq skyscraper sheaf''):
\begin{enumerate}
\renewcommand{\theenumi}{\roman{enumi}}
\renewcommand{\labelenumi}{(\theenumi)}
\item The fiber of $F$ at every geometric point $\bar x$ of $X$
lying over a nonclosed point $x\in X$ vanishes.

\item Every section $s\in F(X')$, where $X'\to X$ is étale of finite
type, vanishes away from finitely many closed points of $X$.

\item There is an isomorphism
$F\simeq\bigoplus_x i_{x*}i_x^*F$, where $x$ runs through the closed
points of $X$ and $i_x:x\to X$ is the inclusion.
\end{enumerate}
\end{enonce*}

\emph{Proof.}\pageoriginale Suppose (i) holds and let
$z\in F(X')$, where $X'\to X$ is étale of finite type. The support of
$z$ (VIII \hyperref[VIII.subsec6.6]{6.6}) is then a closed subset of $X'$ consisting of
\emph{closed} points, and is therefore finite because $X'$ is
Noetherian. Thus (i) implies (ii). The implication (ii) to (i)
follows from the computation of fibers (VIII \hyperref[VIII.prop3.9]{3.9}), and (iii) implies
(i), since for every family $(G_x)$ of sheaves on the closed points
$x$ of $X$, the sheaf $F=\bigoplus i_{x*}(G_x)$ on $X$ satisfies
(i): fiber functors commute with direct sums, and the support of
$i_{x*}(G_x)$ is plainly contained in $\{x\}$.

It remains to prove that (ii) implies (iii). Consider the canonical
morphism
$$
F\longrightarrow\prod_x i_{x*}i_x^*F{\quoi},
$$
where $x$ runs through the closed points of $X$. By (ii), it factors
through
$$
F\longrightarrow\bigoplus_x i_{x*}i_x^*F{\quoi},
$$
which gives the desired morphism. We leave to the reader the
verification that it is bijective.

\setcounter{subsection}{1}
\subsubsection{}\etiquette[4.1.1]{IX.subsubsec4.1.1}
Apply the lemma to the sheaf $R^qi_*F$ introduced above. One has
$$
R^qi_*F\xrightarrow{\sim}
\bigoplus_x i_{x*}i_x^*(R^qi_*F),
$$
where $x$ runs through the closed points of $X$. By VIII 5, the fiber
of $R^qi_*F$ at a geometric point $\bar x$ of $X$ is
$$
(R^qi_*F)_{\bar x}\simeq H^q(\Spec R(X_{\bar x}),F),
$$
where
$R(X_{\bar x})\simeq R(X)\otimes_{\Oo_X}\Oo_{X_{\bar x}}$,
$X_{\bar x}$ is the strict localization of $X$ at $\bar x$ (VIII 4),
and $F$ also denotes the sheaf induced on
$\Spec R(X_{\bar x})$. Suppose now that the residue field of every
closed point $x$ of $X$ is separably closed. Then (see also VI 1)
\begin{align*}
H^p(X,R^qi_*F)
&\simeq H^p\left(X,\bigoplus_x i_{x*}i_x^*(R^qi_*F)\right)\\
&\simeq \bigoplus_x H^p\left(X,i_{x*}i_x^*(R^qi_*F)\right)\\
&\simeq \bigoplus_x H^p\left(x,i_x^*(R^qi_*F)\right)\\
&\simeq
\begin{cases}
\displaystyle\bigoplus_x H^q(\Spec R(X_x),F)
  &\mbox{if }p=0,\\
0 &\mbox{if }p>0{\quoi}.
\end{cases}
\end{align*}
Thus:

\begin{enonce*}{Corollary 4.2}\etiquette[4.2]{IX.cor4.2}
Let $X$ be a Noetherian scheme of dimension $1$ such that, for every
closed point $x$ of $X$, the residue field $k(x)$ is separably
closed. Let $F$ be an abelian sheaf on $\Spec R(X)$, where $R(X)$ is
the ring of rational functions on $X$. Consider the Leray spectral
sequence
$$
E^{pq}_2=H^p(X,R^qi_*F)\Longrightarrow
H^{p+q}(\Spec R(X),F){\quoi}.
$$
One has
\begin{align*}
E^{pq}_2&=0\quad\mbox{if }p>0\mbox{ and }q>0{\quoi},\\
E^{0q}_2&=\bigoplus_x H^q(\Spec R(X_x),F)
  \quad\mbox{if }q>0{\quoi},
\end{align*}
where $x$ runs through the closed points of $X$, and $X_x$ is the
strict localization of $X$ at $x$.

The\pageoriginale spectral sequence therefore gives an exact sequence
\begin{align*}
0\longrightarrow H^1(X,i_*F)&\longrightarrow H^1(\Spec R(X),F)
\longrightarrow\bigoplus_xH^1(\Spec R(X_x),F)\\
&\longrightarrow H^2(X,i_*F)\longrightarrow H^2(\Spec R(X),F)
\longrightarrow\bigoplus_xH^2(\Spec R(X_x),F)\\
&\longrightarrow\cdots .
\end{align*}
\end{enonce*}


% END INCLUDED SGA4_English_current_buildable_source_20260906/source/09/drafts/SGA4_IX_LINES_1340_1459_English.texfrag


% BEGIN INCLUDED SGA4_English_current_buildable_source_20260906/source/09/drafts/SGA4_IX_LINES_1460_1606_English.texfrag
\setcounter{subsection}{2}
\subsection{}\etiquette[4.3]{IX.subsec4.3}
Suppose moreover that the $\ell$-cohomological dimension of
$\Spec R(X)$, for a given prime number $\ell$, is at most $1$ (that
is, $H^q(\Spec R(X),F)=0$ for $q>1$ and every $\ell$-torsion sheaf
$F$ on $\Spec R(X)$). Let $K^1,\ldots,K^m$ be the residue fields of
the Artinian ring $R(X)$, that is, the residue fields of the maximal
points of $X$. By the dictionary of VIII 2, the
$\ell$-cohomological dimension of $\Spec R(X)$ is
$$
\cd_\ell(\Spec R(X))
=\mathop{\sup}_i\{\cd_\ell(G(\overline K^i/K^i))\}{\quoi},
$$
where $\cd_\ell(G)$ denotes the $\ell$-cohomological dimension of the
profinite group $G$ (see CG).

Every residue field $K_x$ of $R(\bar X(x))$, for a strict
localization $\bar X(x)$ of $X$, is plainly a separable (possibly
infinite) extension of one of the $K^i$, and hence
$$
\cd_\ell(G(\bar K_x/K_x))
\leq\cd_\ell(G(\bar K^i/K^i))
$$
by CG II 4.1, Proposition 10. Thus
$\cd_\ell(\Spec R(\bar X(x)))\leq1$ as well.

Applying the exact sequence of \hyperref[IX.cor4.2]{4.2} to an
$\ell$-torsion sheaf $F$, one obtains
\begin{equation*}
\begin{cases}
0\to H^1(X,i_*F)\to H^1(\Spec R(X),F)\to\\
\qquad\displaystyle\bigoplus_xH^1(\Spec R(\bar X(x)),F)
\to H^2(X,i_*F)\to0{\quoi},\\
H^q(X,i_*F)=0\quad\mbox{for }q>2{\quoi}.
\end{cases}
\tag{4.3.1}
\end{equation*}\etiquette[4.3.1]{IX.eq4.3.1}

Note\pageoriginale that the hypothesis
$\cd_\ell(\Spec R(X))\leq1$ holds if $X$ is an algebraic curve over a
separably closed field (Tsen's theorem [6, Theorem
\hyperref[IX.subsec4.3]{4.3}]). Analogous exact sequences in this case
were studied by $\Ogg[2]$ and \v{S}afarevi\v{c}.

\subsection{}\etiquette[4.4]{IX.subsec4.4}
Suppose that $\ell$ is invertible on $X$ and take
$F=(\GG_m)_{R(X)}$. Applying \hyperref[IX.thm3.3]{3.3} and
[CG, Chapter I, Proposition 12], still under the hypothesis
$\cd_\ell(\Spec R(X))\leq1$, one finds
\begin{equation*}
\begin{cases}
H^1(\Spec R(X),\GG_m)=0,\\
H^q(\Spec R(X),\GG_m)\;\displaystyle{\mathop{=}_{\ell}}\;0
\quad(q\geq2){\quoi},
\end{cases}
\tag{4.4.1}
\end{equation*}\etiquette[4.4.1]{IX.eq4.4.1}
where the $\ell$ under the equality sign means that the group on the
left, which is torsion in any case, has no $\ell$-torsion.

The same result holds with $R(X)$ replaced by $R(\bar X(x))$. Since
the sheaf $\GG_m$ on $\Spec R(X)$ plainly induces the sheaf $\GG_m$
on $\Spec R(\bar X(x))$, \hyperref[IX.cor4.2]{4.2} gives
\begin{equation*}
\begin{cases}
H^1(X,i_*(\GG_m)_{R(X)})=0,\\
H^q(X,i_*(\GG_m)_{R(X)})
\;\displaystyle{\mathop{=}_{\ell}}\;0
\quad\mbox{if }q>1{\quoi}.
\end{cases}
\tag{4.4.2}
\end{equation*}\etiquette[4.4.2]{IX.eq4.4.2}
For every sheaf $F$ on $X$, \hyperref[IX.lem4.1]{4.1} shows that the
kernel $P$ and cokernel $Q$ of the canonical homomorphism
$F\to i_*i^*F$ are skyscraper sheaves. Since the residue fields at
the closed points of $X$ are assumed separably closed, one therefore
has $H^j(X,P)=H^j(X,Q)=0$ for $j>0$. The long exact cohomology
sequence immediately implies that
$$
H^j(X,F)\longrightarrow H^j(X,i_*i^*F)
$$
is\pageoriginale an isomorphism for $j\geq2$. Taking
\eqhyperref[IX.eq4.4.2]{4.4.2} into account, one obtains
\begin{equation*}
H^j(X,\GG_m)\displaystyle{\mathop{\;=\;}_{\ell}}0
\quad\mbox{for }j\geq2{\quoi},
\tag{4.5}
\end{equation*}\etiquette[4.5]{IX.eq4.5}
valid when $X$ is a Noetherian scheme of dimension $1$ satisfying
$\cd_\ell(R(X))\leq1$ and whose residue fields at closed points are
separably closed. Kummer theory (\hyperref[IX.Thr3.2]{3.2}) and
\hyperref[IX.thm3.3]{3.3} now give:

\begin{enonce*}{Theorem 4.6}\etiquette[4.6]{IX.thm4.6}
Let $X$ be Noetherian of dimension $1$, and let $n$ be an integer
invertible on $X$. Suppose that $\cd_\ell(R(X))\leq1$ for every prime
number $\ell$ dividing $n$, and that the residue field $k(x)$ is
separably closed for every closed point $x$ of $X$. Then
$$
H^q(X,(𝛍_n)_X)=0\quad\mbox{if }q>2{\quoi},
$$
and the cohomology for $q\leq2$ is given by the exact sequence
\begin{align*}
0\to H^0(X,(𝛍_n)_X)&\to\Gamma(X,\Oo_X^*)
\xrightarrow{n}\Gamma(X,\Oo_X^*)\to\\
H^1(X,(𝛍_n)_X)&\to\Pic X\xrightarrow{n}\Pic X
\to H^2(X,(𝛍_n)_X)\to0{\quoi}.
\end{align*}
\end{enonce*}

\begin{enonce*}{Corollary 4.7}\etiquette[4.7]{IX.cor4.7}
Let $X$ be a complete connected curve over a separably closed field
$k$ whose characteristic is prime to $n$. Then
\begin{align*}
H^0(X,(𝛍_n)_X)&\simeq𝛍_n(k){\quoi},\\
H^1(X,(𝛍_n)_X)&\simeq{}_n(\Pic X)
  =\{\text{elements of $\Pic X$ whose order divides $n$}\},\\
H^2(X,(𝛍_n)_X)&\simeq(\Pic X)/n
  \simeq(\ZZ/n\ZZ)^c{\quoi},
\end{align*}
where\pageoriginale $c$ is the number of irreducible components of
$X$.
\end{enonce*}

For the last assertion, recall that the Picard scheme [TDTE V] of
such an $X/\Spec k$ has a decomposition
\begin{equation*}
0\longrightarrow A\longrightarrow\sheaf{\Pic}_{X/k}
\longrightarrow\ZZ^c\longrightarrow0
\tag{4.8}
\end{equation*}\etiquette[4.8]{IX.eq4.8}
where $A$ is a connected algebraic group over $\Spec k$. The
\lq\lq multiplication by $n$'' map on such an $A$ is étale, hence
surjective on $k$-valued points. It follows that
$$
{}_n(\Pic X)\simeq{}_n(A(k))
$$
and
$$
(\Pic X)/n\simeq(\ZZ/n\ZZ)^c{\quoi}.
$$

If $X$ is smooth over $\Spec k$, so that its Jacobian $A$ is an
abelian variety, the group \({}_n(\Pic X)\) is isomorphic to
$(\ZZ/n)^{2g}$, where $g$ is the genus of $X$. Of course,
$𝛍_n(k)\simeq\ZZ/n\ZZ$ as well. Thus the ranks of the cohomology
groups as $\ZZ/n$-modules are $1$, $2g$, and $1$, as expected.


% END INCLUDED SGA4_English_current_buildable_source_20260906/source/09/drafts/SGA4_IX_LINES_1460_1606_English.texfrag


% BEGIN INCLUDED SGA4_English_current_buildable_source_20260906/source/09/drafts/SGA4_IX_LINES_1607_1883_English.texfrag
\section{The trace method}\etiquette[5]{IX.sec5}

\subsection{}\etiquette[5.1]{IX.subsec5.1}
Let $X$ be a scheme, let $f:X'\rightarrow X$ be an étale covering, and
let $F$ be a sheaf on $X$. Then the adjunction formulas give a
\emph{restriction} morphism
$$
F\xrightarrow{\res}f_{\ast}f^{\ast}F{\quoi}.
$$
Since $f$ is finite, one has $H^q(X,f_{\ast}f^{\ast}F)\simeq H^q(X',
f^{\ast}F)$, whence a morphism, also called restriction:
\begin{equation*}
H^q(X,F)\xrightarrow{\res}H^q(X',f^{\ast} F){\quoi},\tag{5.1.1}
\end{equation*}\etiquette[5.1.1]{IX.eq5.1.1}
which\pageoriginale is moreover the evident morphism. But when $f$ is
étale, there is also another morphism, called the \emph{trace
morphism},
\begin{equation*}
f_{\ast}f^{\ast} F\xrightarrow{\tr} F{\quoi},\tag{5.1.2}
\end{equation*}\etiquette[5.1.2]{IX.eq5.1.2}
which induces morphisms on cohomology
\begin{equation*}
H^q(X', f^{\ast} F)\xrightarrow{\tr}H^q
(X,F){\quoi}.\tag{5.1.3}
\end{equation*}\etiquette[5.1.3]{IX.eq5.1.3}
The trace morphism is characterized by the following two properties:
\begin{enumerate}
\renewcommand{\theenumi}{\roman{enumi}}
\renewcommand{\labelenumi}{(\theenumi)}
\item It is local for the étale topology.

\item If $X'$ is the disjoint sum of $d$ copies of $X$, so that
$f_{\ast}f^{\ast} F\simeq F^d$, the trace is the summation map
$F^d\rightarrow F$.
\end{enumerate}

Uniqueness follows trivially from (i), (ii), since every étale
covering is locally constant. Once uniqueness is known, existence
follows by the usual descent argument.

From (ii) one deduces the following formula:
\begin{equation*}
\tr\circ \res:F\longrightarrow F\mbox{ is multiplication by the
local degree of $f$.}\tag{5.1.4}
\end{equation*}\etiquette[5.1.4]{IX.eq5.1.4}
If the degree is a constant $d$, it follows that the morphism
$\tr\circ \res: H^q(X,F)\rightarrow H^q(X,F)$ is multiplication by
$d$.

\begin{enonce*}{Corollary 5.2}\etiquette[5.2]{IX.cor5.2}
Let $f:X'\rightarrow X$ be an étale covering of constant degree $d$.
If $F$ is an $r$-torsion sheaf on $X$ with $(r,d)=1$, then
$H^q(X',f^{\ast} F)=0$ implies that $H^q(X,F)=0$.
\end{enonce*}

Indeed, multiplication by $d$ induces an isomorphism of $F$, hence
an\pageoriginale isomorphism on cohomology. If the latter isomorphism
factors through zero, then $H^q(X,F)=0$.

\setcounter{subsection}{2}
\subsubsection{}\etiquette[5.2.1]{IX.subsubsec5.2.1}
Now let $i:U\rightarrow X$ be an immersion, let $f:U'\rightarrow U$ be
an étale covering of degree $d$, and suppose that $f$ is induced by a
finite morphism $g:X'\rightarrow X$ by base change, \ie that one has
the cartesian diagram
$$
\xymatrix@C=2cm@R=1.25cm{
U' \ar[r]^{i'}\ar[d]_{f}& X'\ar[d]^{g}\\
U  \ar[r]^{i}& X
}
$$
Let $F$ be a sheaf on $U$. One has
$$
\xymatrix@C=2.25cm@M=.25cm{
F \ar[r]^{\txt{res}} \ar@/_1.5pc/[rr]_{d}& f_{\ast}f^{\ast}F
\ar[r]^{\txt{tr}}& F,
}
$$
hence, since $i_{\ast}f_{\ast}=g_{\ast}i'_{\ast}$,
\hypertarget{sga4.exp.IX.equation.5.3}{}$$
\xymatrix@C=2.25cm{
i_{\ast}F \ar@/_2pc/[rr]_{d} \ar[r]^{\txt{res}}& g_{\ast}i' {}_{\ast}f^{\ast}F \ar[r]^{\txt{tr}}& i_{\ast}F,
}\leqno{(5.3)}
$$
and since $i_{!}f_{\ast}\simeq g_{\ast}i'_{!}$ (which follows, for
example, from (VIII \hyperref[VIII.prop5.5]{5.5})),
$$
\xymatrix@C=2cm{
i_{!}F \ar@/_2pc/[rr]_{d} \ar[r]^{\txt{res}}& g_{\ast}i' {}_{!}f^{\ast}F \ar[r]^{\txt{tr}}& i_{!}F,
}\leqno{(5.4)}
$$\etiquette[5.4]{IX.eq5.4}
whence
$$
\xymatrix@C=2cm{
H^{q}(X,i_{\ast} F) \ar@/_2pc/[rr]_{d'} \ar[r]^{\txt{res}}& H^{q}(X', i' _{\ast}f^{\ast}F) \ar[r]^{\txt{tr}}& H^q(X,i_{\ast}F)
}\leqno{(5.3')}
$$\MakeLinkTarget*{sga4.exp.IX.equation.5.3-prime}\label{5.3'}
and\pageoriginale
$$
\xymatrix@C=2cm{
H^{q}(X,i_{!} F) \ar[r]^{\txt{res}} \ar@/_2.5pc/[rr]_{d}& H^{q}(X', i' _{!}f^{\ast}F)
\ar[r]^{\txt{tr}}& H^q(X,i_{!}F) ;
}\leqno{(5.4')}
$$\MakeLinkTarget*{sga4.exp.IX.equation.5.4-prime}\label{5.4'}
From \eqhyperref[IX.eq5.4]{5.4} one deduces the following corollaries:

\begin{enonce*}{Proposition 5.5}\etiquette[5.5]{IX.prop5.5}
Let $X$ be a quasi-compact and quasi-separated scheme, let $\ell \in
\PP$, and let $H$ be a half-exact functor with values
in {\rm (Ab)}, defined on the category of $\ell$-sheaves on $X$, and
compatible with filtered inductive limits. The following assertions
are equivalent:
\begin{enumerate}
\renewcommand{\theenumi}{\roman{enumi}}
\renewcommand{\labelenumi}{(\theenumi)}
\item $H=0$.

\item $H(f_{\ast}i_{!}(\ZZ/\ell))=0$ for every finite
$f:Y\rightarrow X$, and every open subscheme $i:U\rightarrow Y$ of
$Y$, such that $U$ is of finite presentation over $X$. If $X$ is
noetherian, one may even restrict to integral $Y$.
\end{enumerate}
\end{enonce*}

Proof. By \hyperref[IX.cor2.7.2]{2.7.2} and \hyperref[IX.prop2.5]{2.5}
it is enough to verify $H(F)=0$ for $F$ of the form $F=i$, $G$, where
$i:U\rightarrow X$ is the inclusion of a subscheme of $X$ of finite
presentation over $X$, and where $G$ is a locally constant
``isotrivial'' $\ell$-sheaf on $U$. Let $U''\rightarrow U$ be a
principal covering with group $g$, such that $g$ becomes constant on
$U''$. Let $A$ be the $\ell$-group of values of $G$ on $U''$. The
group $g$ acts on $A$, and this action determines the structure of
$G$. Let $\sheaf{h}$ be an $\ell$-Sylow subgroup of $g$, and let
$U'=U''/\sheaf{h}$ be the covering of $U$ corresponding to
$\sheaf{h}$. The degree $d$ of $U'$ over $U$ is prime to $\ell$.
Choose an extension of $f:U'\rightarrow U$ to a finite morphism
$g:X'\rightarrow X$, which always exists (\EGA IV 8.12.6), and let
$i':U'\rightarrow X'$ be the inclusion. Applying
\eqhyperref[IX.eq5.4]{5.4} and the fact that $d$ is prime to $\ell$,
we are reduced to proving
that $H(g_{\ast}i'_{!}f^{\ast} G)=0$.

Now\pageoriginale it is well known that the only abelian
$\ell$-torsion group that is simple under an action of an $\ell$-group
$\sheaf{h}$ is $\ZZ/\ell$ with trivial action, and it follows that
there exists a filtration of $A$ as an $\sheaf{h}$-module whose
successive quotients are isomorphic to $\ZZ/\ell$ with trivial action.
This filtration gives a corresponding filtration of $f^{\ast} G$, and
therefore, in order to prove that $H(g_{\ast}i'_{!} f^{\ast} G)=0$, it
is enough to prove that $H(g_{\ast}i'_{!}\ZZ/r)=0$, whence the result.

\begin{enonce*}{Proposition 5.6}\etiquette[5.6]{IX.prop5.6}
Let $X$ be a noetherian scheme, let $\ell \in \PP$, and
let $\varphi$ be a function with values in $\NN$, defined on the set
of integral schemes finite over $X$. Suppose that
$\varphi(Y_1)\leq \varphi(Y_2)$ whenever there exists an $X$-morphism
$Y_1\rightarrow Y_2$, and that the inequality is strict if $Y_1$ is a
closed subscheme of $Y_2$, distinct from $Y_2$. For every $Y$ finite
over $X$, put $(Y)=\Sup (Y_i)$, where the $Y_i$ are the irreducible
components of $Y$, endowed with the induced reduced structure. Then
the following assertions are equivalent:
\begin{enumerate}
\renewcommand{\theenumi}{\roman{enumi}}
\renewcommand{\labelenumi}{(\theenumi)}
\item For every $Y$ finite over $X$ and every $\ell$-sheaf $F$ on
$Y$, one has $H^q(Y,F)=0$ if $q>\varphi(\Supp F)$.

\item For every integral $Y$ finite over $X$, one has
$H^q(Y,\ZZ/\ell)=0$ if $q>\varphi(Y)$.
\end{enumerate}
\end{enonce*}

Proof. Suppose that (ii) is true. We may apply the preceding
proposition to each $H^q$, taking into account that cohomology commutes
with inductive limits and with direct images under finite morphisms;
we are thus reduced, for the verification of (i), to the case where
$Y$ is integral and finite over $X$, and $F$ is of the form
$i_{!}(\ZZ/\ell)_U$, where $i:U\rightarrow X$ is a nonempty open
subscheme.

Now there is an exact sequence
$$
0\longrightarrow i_{!}(\ZZ/\ell)_U\longrightarrow
(\ZZ/\ell)_Y\longrightarrow
(\ZZ/\ell)_Z\longrightarrow 0
$$
where $Z=Y-U$, hence $Z\neq Y$. Arguing by induction on the value of
$\varphi$, we may\pageoriginale suppose that
$H^q(\ZZ, \ZZ/\ell)=0$ for $q>\varphi (Y)-1$, whence
$H^q(Y,i_{!}(\ZZ/\ell))=0$ for $q>\varphi(Y)$ by the exact
cohomology sequence.

\begin{enonce*}{Corollary 5.7}\etiquette[5.7]{IX.cor5.7}
Let $X$ be an algebraic curve over a separably closed field $k$ of
characteristic different from $\ell$, with $\ell\in \PP$.
Then the cohomology of $X$ with coefficients in a sheaf $F$ of
$\ell$-torsion vanishes in degrees $>2$. If $X$ is affine, the cohomology
vanishes in degrees $>1$.
\end{enonce*}

For integral $Y$ finite over $X$, put $\varphi(Y)=2\dim Y$ (\resp
$\varphi(Y)=\dim Y$ if $X$, hence $Y$, is affine). By \hyperref[IX.prop5.6]{7.6} it is enough
to prove that $H^q(Y,\ZZ/\ell)=0$ if $q>\varphi(Y)$, which is trivial
if $\dim Y=0$. It therefore remains to prove that
$H^q(Y,\ZZ/\ell)=0$ if $q>2$ (\resp if $q>1$ in the affine case) for
an integral curve $Y$. Since $(\ZZ/\ell)_Y\simeq (𝛍_{\ell})_Y$, this
follows from \hyperref[IX.thm4.6]{4.6} for $q>2$.

Suppose that $X$ is affine. Since radicial extensions do not affect
cohomology (VIII \hyperref[VIII.thm1.1]{1.1}), we are reduced to the case where $k$ is
algebraically closed. Let $f:X'\rightarrow X$ be the normalization of
$X$, which is a finite morphism and an isomorphism outside a finite set
of points, say $S$. There is an exact sequence of sheaves on $X$
$$
0\longrightarrow (\ZZ/n)_X\longrightarrow
f_{\ast}(\ZZ/n)_{X'}\longrightarrow \epsilon
\longrightarrow 0{\quoi},
$$
where $\epsilon$ is concentrated on $S$, so that
$H^q(X,\epsilon)=0$ for $q>0$. By (VIII \hyperref[VIII.cor5.6]{5.6}), we are reduced to proving
that
$$
H^q(X,f_{\ast}(\ZZ/n)_{X'})\simeq
H^q(X',\ZZ/n)=0
$$
if $q>1$, that is, to the case where $X$ is normal. We may also plainly
suppose that $X$ is connected. Let $i:X\rightarrow \bar{X}$ be an open
immersion with $\bar{X}$ complete and nonsingular. Since $X$ is
affine, at least one point is missing from $\bar{X}$, and it follows
immediately that the morphism
$$
A(k)\longrightarrow \Pic X
$$
is\pageoriginale surjective, where $A$ is the Jacobian of $\bar{X}$
(\cf \eqhyperref[IX.eq4.8]{4.8}). Since $A(k)$ is divisible by
$\ell\neq \car k$, the same is true of $\Pic X$, whence
$H^2(X,𝛍_{\ell})=0$ by \hyperref[IX.thm4.6]{4.6}, \cqfd

\begin{enonce*}[remark]{Remark 5.8}\etiquette[5.8]{IX.rem5.8}
Let us record, for later reference, that the argument of
\hyperref[IX.prop5.5]{5.5} in fact establishes the following result:
Let $X$ be a quasi-compact and quasi-separated scheme, let $\ell$ be a
prime number, and let $C$ be a strictly full subcategory of the
category of constructible abelian $\ell$-torsion sheaves, stable under
direct factors and extensions. Suppose that for every finite morphism
$f:Y\rightarrow X$ and every open subscheme $i:U\rightarrow Y$ of $Y$
that is of finite presentation over $X$, one has
$f_{\ast}(i_{!}(𝐙/\ell𝐙)_U)\in C$. Then $C$ contains all constructible
%% 𝒵 here is ℤ [check the macros]
$\ell$-torsion sheaves.
\end{enonce*}

\begin{thebibliography}{99}
\bibitem{sga4.exp.IX.bib.1} Gabriel, P. Des Catégories Abéliennes,
Thèse, Paris (1962).

\bibitem{sga4.exp.IX.bib.2} Ogg, A. Cohomology of abelian varieties over
functions fields, Annals of Math. Vol. 76, No. 2 (1962).

\bibitem{sga4.exp.IX.bib.3} Deuring, M. Die Typen der Multiplikatorenringe elliptischer
Funktionenkörper, Hamb. Abh. Bol. 13 (1940) p. 197.

\bibitem{sga4.exp.IX.bib.4} Grothendieck, A. Fondements de la Géométrie
Algébrique, Extraits du Séminaire Bourbaki 1957-1962), Secrétariat
mathématique, 11, rue Pierre Curie, Paris 5ème. (\emph{Cited as} FGA).

\bibitem{sga4.exp.IX.bib.5} Serre, J.P. Cohomologie Galoisienne, Lecture Notes
nᵒ 5, (Springer) (\emph{Cited as} CG).

\bibitem{sga4.exp.IX.bib.6} Douady, A. Cohomologie des groupes compacts
totalement discontinus, Séminaire Bourbaki, exp. 189.
\end{thebibliography}

\ifx\danslelivre\undefined
\end{document}
\fi

% END INCLUDED SGA4_English_current_buildable_source_20260906/source/09/drafts/SGA4_IX_LINES_1607_1883_English.texfrag


% END INCLUDED SGA4_English_current_buildable_source_20260906/source/09/09_English.tex


\makeatletter
\def\input@path{{../10/}}
\makeatother
% BEGIN INCLUDED SGA4_English_current_buildable_source_20260906/source/10/10_English.tex
\ifx\danslelivre\undefined
% Release engine: XeLaTeX (the LuaTeX compatibility layer breaks Xy-pic output).
\documentclass[10pt]{../fichiers-configuration-latex/smfart}
% BEGIN INCLUDED SGA4_English_current_buildable_source_20260906/source/fichiers-configuration-latex/paquets.tex
\usepackage{polyglossia}
%\usepackage[francais]{babel}
\usepackage{fontspec}
\usepackage[math-style=ISO,bold-style=upright,sans-style=upright]{unicode-math}
%\usepackage[mathletters]{ucs}
%\usepackage[utf8]{inputenc}
%\usepackage[T1]{fontenc}
\usepackage[unicode,hypertexnames=false,colorlinks=true,linkcolor=blue,citecolor=Sepia,urlcolor=magenta]{hyperref}
\setmainlanguage{french}
\setmainfont[Ligatures=TeX]{Linux Libertine O}
% variante possible : TeXGyrePagella ?
\setmathfont{XITS Math}

% maths
\usepackage{amsmath}
%\usepackage{amssymb}
%\usepackage{mathrsfs}
%\usepackage{stmaryrd}
%\usepackage{euscript}

% couleurs
\usepackage[usenames,dvipsnames]{xcolor}

% graphiques
\usepackage{graphicx}
\usepackage[all]{xy}
%\usepackage{tikz}
%\usetikzlibrary{matrix,arrows,calc}
%\tikzset{column sep=3em,row sep=2em,text height=1.5ex, text depth=0.25ex,>=angle 90}
%\tikzset{description/.style={fill=white,inner sep=1pt}}

% hyperliens, 'etiquettes
\usepackage{xr-hyper}
\usepackage{../fichiers-configuration-latex/smfthm}
%\ifx\cacheretiquettes\undefined
%\usepackage[inline]{showlabels} \else \fi % pour rendre visibles les 'etiquettes
%\usepackage{../fichiers-configuration-latex/makeidx} % pour faire un index
%\usepackage[pagebackref,unicode]{hyperref}  % pour voir o`u sont cit'ees les r'ef'erences
% hyperref is loaded once above with the cumulative-reader destination policy.
%\hypersetup{breaklinks=true} %permet le retour `a la ligne dans les liens trop longs. Incompatible avec pdflatex
%\renewcommand*{\backref}[1]{{\small [cit\'e page(s)~#1.]}}

\usepackage{xspace}
\usepackage{pifont} % pour expose XVII

% END INCLUDED SGA4_English_current_buildable_source_20260906/source/fichiers-configuration-latex/paquets.tex

% BEGIN INCLUDED SGA4_English_current_buildable_source_20260906/source/fichiers-configuration-latex/numerotation.tex
\renewcommand{\thefootnote}{\roman{footnote}} % pour mettre les notes en bas de page sous la forme (ii) et non (2), qui ressemble trop (?) a un exposant.

\newcounter{pageoriginal}
\newcommand{\pageoriginale}{\refstepcounter{pageoriginal}\marginpar{\footnotesize\textbf{\thepageoriginal}}\xspace}
\let\pageoriginaled\pageoriginale

\usepackage{ifthen}
\newboolean{original}
\def\sisi{\ifthenelse{\boolean{original}}}
\setboolean{original}{false}

\newcounter{cnote}
\def\nde#1{\ifthenelse{\boolean{original}}{}{\stepcounter{cnote}\def\thefootnote{\arabic{cnote}}\protect\footnote[\value{cnote}]{N.D.E.: #1}}}
\def\ndemark{\ifthenelse{\boolean{original}}{}{\stepcounter{cnote}\def\thefootnote{\arabic{cnote}}\protect\footnotemark[\value{cnote}]}}
\def\ndetext#1{\ifthenelse{\boolean{original}}{}{\def\thefootnote{\arabic{cnote}}\protect\footnotetext[\value{cnote}]{N.D.E.: #1}}}

% END INCLUDED SGA4_English_current_buildable_source_20260906/source/fichiers-configuration-latex/numerotation.tex

% BEGIN INCLUDED SGA4_English_current_buildable_source_20260906/source/fichiers-configuration-latex/macros.tex
\setmainlanguage{french}
\setotherlanguage{english}


%% Fixed, engine-independent release identifier.  The inherited dynamic Git
%% probe used Lua io.popen, exposed build-state-dependent output, and printed
%% a visible question mark outside a Git checkout.
\def\showgitstatus{English release r2, 2026-07-23}


% macros pour LaTeX
\renewcommand{\mathsf}{}

% Police différente pour U+2113 SCRIPT SMALL L
\setmathfont[range="2113]{latinmodern-math.otf}

% Give every inherited semantic label its own stable PDF destination instead
% of inheriting whichever automatic equation/page anchor happened to be live.
% The legacy label remains as a compatibility alias used by existing edges.
\makeatletter
\def\etiquette[#1]#2{%
  \Hy@raisedlink{\hyper@anchorstart{sga4.semantic.#2}\hyper@anchorend}%
  \def\@currentHref{sga4.semantic.#2}%
  \def\@currentlabel{#1}%
  \label{#2}%
}
\makeatother
\def\eqhyperref[#1]#2{(\hyperref[#1]{#2})}
\def\pagehyperref[#1]#2{(page \hyperref[#1]{#2})}
\def\hyperrefext[#1][#2]#3{\hyperref[#2]{#1 #3}}
\def\eqhyperrefext[#1][#2]#3{\hyperref[#2]{#1 #3}}     %cf. exposes 17
\def\neweqhyperrefext[#1][#2]#3{\hyperref[#2]{#1 #3}}  %cf. exposes 17

\def\shat#1{\protect\widehat{#1}}
\def\stilde#1{\protect\widetilde{#1}}
%%%%%%%%%%%%%%%%%%


\def\etcn{etc.\xspace}
\def\etco{etc...\xspace}
\def\etc{\sisi{\etco}{\etcn}}
\def\ldotsn{\sisi{...}{\etcn}}

\def\cf{{cf.}\kern.3em}
\def\Cf{{Cf.}\kern.3em}
\def\ie{{i.e.}\kern.3em}
\def\iev{{i.e.},\xspace}
\def\loccito{loc.\kern2pt cit.\xspace}
\def\loccitn{\textit{loc.\kern2pt cit.}\xspace}
\def\loccit{\sisi{\loccito}{\loccitn}}
\def\EGA{\textup{EGA}\kern.2em}
\def\SGA{\textup{SGA}\kern.3em}

\def\quoi{}

\def\quav{\quad,\quad}
\def\quaet{\quad \text{et}\quad}
\def\fibreproduct#1#2#3{#1 ×_{#3} #2}
\let\fprod\fibreproduct
\def\ccheck#1#2{\mbox{${\check{\check{\rm #1}}}^{\lower2pt\hbox{$\scriptstyle #2$}}$}}
\def\sheaf#1{\protect\mathcal{#1}}
\def\vieuxsheaf#1{{#1}} % cf. problème avec sheaf{a\Phi_i} par exemple...
\DeclareMathOperator{\Hom}{Hom}
\DeclareMathOperator{\can}{can}
\DeclareMathOperator{\Ind}{Ind}
\DeclareMathOperator{\ind}{ind}
%\DeclareMathOperator{\Hhomtop}{\mathcal{H}\!\mathit{omtop}}
\DeclareMathOperator{\Hhomtop}{\mathbf{Homtop}}
\DeclareMathOperator{\Point}{Point}
\DeclareMathOperator{\Pro}{Pro}
\DeclareMathOperator{\Top}{Top}
\DeclareMathOperator{\Mod}{Mod}
\DeclareMathOperator{\Cart}{Cart}
%\DeclareMathOperator{\Mmorsite}{\mathcal{M}\!\mathit{orsite}}
\DeclareMathOperator{\Mmorsite}{\mathbf{Morsite}}
\DeclareMathOperator{\loc}{loc}
\DeclareMathOperator{\Loc}{Loc}
\DeclareMathOperator{\Pic}{Pic}
\DeclareMathOperator{\Et}{Et}
\DeclareMathOperator{\card}{card}
\DeclareMathOperator{\Ens}{Ens}
\DeclareMathOperator{\h}{\underline{H}}
\DeclareMathOperator{\Fonct}{Fonct}
\DeclareMathOperator{\Ker}{Ker}
\DeclareMathOperator{\but}{but}
\DeclareMathOperator{\Ob}{Ob}
\DeclareMathOperator{\ob}{ob}
\DeclareMathOperator{\Lib}{Lib}
\DeclareMathOperator{\esj}{esj}
\DeclareMathOperator{\id}{id}
\DeclareMathOperator{\Id}{Id}
\DeclareMathOperator{\ad}{ad}
\DeclareMathOperator{\Iim}{Im}
\DeclareMathOperator{\pr}{pr}
\DeclareMathOperator{\Cat}{Cat}
\DeclareMathOperator{\car}{car}
\DeclareMathOperator{\SysInd}{SysInd}
\DeclareMathOperator{\Sup}{Sup}
\DeclareMathOperator{\Coker}{Coker}
\DeclareMathOperator{\Filt}{Filt}
\DeclareMathOperator{\acc}{acc}
\DeclareMathOperator{\sur}{sur}
\DeclareMathOperator{\cat}{cat}
\DeclareMathOperator{\Ord}{Ord}
\DeclareMathOperator{\Gr}{Gr}
\DeclareMathOperator{\Cov}{Cov}
\DeclareMathOperator{\Caf}{Caf}
\DeclareMathOperator{\coker}{coker}
\DeclareMathOperator{\coim}{coim}
\DeclareMathOperator{\im}{im}
\DeclareMathOperator{\Ouv}{Ouv}
\DeclareMathOperator{\diag}{diag}
\DeclareMathOperator{\Esp}{Esp}
\DeclareMathOperator{\TOP}{TOP}
\DeclareMathOperator{\Unif}{Unif}
\DeclareMathOperator{\Iind}{\mathbf{Ind}}
\DeclareMathOperator{\Hhomcart}{\mathbf{Homcart}}
\DeclareMathOperator{\Hhomtopess}{\mathbf{Homtopess}}
\DeclareMathOperator{\Ppoints}{\mathbf{Points}}
\DeclareMathOperator{\Ppoint}{\mathbf{Point}}
\DeclareMathOperator{\Ttors}{\mathbf{Tors}}
\DeclareMathOperator{\Ssex}{\mathbf{Sex}}
\DeclareMathOperator{\Ffib}{\mathbf{Fib}}
\DeclareMathOperator{\Vvois}{\mathbf{Vois}}
\DeclareMathOperator{\Ppointess}{\mathbf{Pointess}}
\DeclareMathOperator{\scrHom}{\mathbf{Hom}}
\DeclareMathOperator{\sheafPic}{\mathbf{Pic}}
\DeclareMathOperator{\sheafAut}{\mathbf{Aut}}
\DeclareMathOperator{\Hhomtoplocan}{\mathbf{Homtoplocan}}
\DeclareMathOperator{\Hhomtopan}{\mathbf{Homtopan}}
\DeclareMathOperator{\Pptgeom}{\mathbf{Ptgéom}}
\DeclareMathOperator{\Morsite}{Morsite}
\DeclareMathOperator{\Fib}{Fib}
\DeclareMathOperator{\Fais}{Fais}
\DeclareMathOperator{\kar}{kar}
\DeclareMathOperator{\Groupes}{Groupes}
\DeclareMathOperator{\Restr}{Restr}
\DeclareMathOperator{\Isom}{Isom}
\DeclareMathOperator{\Prol}{Prol}
\DeclareMathOperator{\Res}{Res}
\DeclareMathOperator{\Tors}{Tors}
\DeclareMathOperator{\esp}{esp}
\DeclareMathOperator{\es}{es}
\DeclareMathOperator{\od}{od}
\DeclareMathOperator{\gr}{gr}
\DeclareMathOperator{\sob}{sob}
\DeclareMathOperator{\dfn}{dfn}
\DeclareMathOperator{\Sch}{Sch}
\DeclareMathOperator{\ouv}{ouv}
\DeclareMathOperator{\cons}{cons}
\DeclareMathOperator{\Nerf}{Nerf}
\DeclareMathOperator{\Points}{Points}
\DeclareMathOperator{\Espsob}{Espsob}
\DeclareMathOperator{\Homtop}{Homtop}
\DeclareMathOperator{\Pt}{Pt}
\DeclareMathOperator{\Quot}{Quot}
\DeclareMathOperator{\glob}{glob}
\DeclareMathOperator{\Homfib}{Homfib}
\DeclareMathOperator{\Qouv}{Qouv}
\DeclareMathOperator{\Quotouv}{Quotouv}
\DeclareMathOperator{\ens}{ens}
\DeclareMathOperator{\ab}{ab}
\DeclareMathOperator{\an}{an}
\DeclareMathOperator{\fs}{fs}
\DeclareMathOperator{\red}{red}
\DeclareMathOperator{\Ab}{Ab}
\DeclareMathOperator{\cosq}{cosq}
\DeclareMathOperator{\Ann}{Ann}
\DeclareMathOperator{\Zar}{Zar}


%%%%%%%%%%%%
%%%%% SGA4-2
\DeclareMathOperator{\Sex}{Sex}
\DeclareMathOperator{\vois}{vois}
\DeclareMathOperator{\Vois}{Vois}
\DeclareMathOperator{\Ext}{Ext}
\DeclareMathOperator{\ext}{ext}
\DeclareMathOperator{\Fer}{Fer}
\DeclareMathOperator{\Add}{Add}
\DeclareMathOperator{\Cocart}{Cocart}
\DeclareMathOperator{\cosk}{cosk}
\DeclareMathOperator{\plat}{plat}

\DeclareMathOperator{\faisc}{faisc}
\DeclareMathOperator{\sk}{sk}
\DeclareMathOperator{\ch}{ch}

\def\HH{\mathbf{H}}
\def\qc{\mathrm{qc}}
\def\et{\mathrm{et}}
\def\coh{\mathrm{coh}}

\def\Hvr#1{%
\begin{array}[b]{l}
{\kern2pt\scriptstyle r}\\[-8pt]
{\kern1pt\scriptstyle \vee}\\[-5pt]
H^{#1}
\end{array}}

\def\Hvv#1{%
\begin{array}[b]{l}
\kern2pt\rotatebox{90}{${\scriptscriptstyle\ll}$}\\[-6pt]
H^{#1}
\end{array}}

\def\Hvi#1{%
\begin{array}[b]{l}
{\kern0.2pt\scriptscriptstyle \infty}\\[-8pt]
{\kern1pt\scriptstyle \vee}\\[-5pt]
H^{#1}
\end{array}}

\def\Hv#1{%
\begin{array}[b]{l}
{\kern1pt\scriptstyle \vee}\\[-5pt]
H^{#1}
\end{array}}

\DeclareMathOperator{\cart}{cart}
\DeclareMathOperator{\Ffaisc}{\mathbf{Faisc}}


\def\Ev#1{%
\begin{array}[b]{l}
{\kern1pt\scriptstyle \vee}\\[-5pt]
E'_{#1}
\end{array}}

\DeclareMathOperator{\Cont}{Cont}
\DeclareMathOperator{\Ccont}{\mathscr{C}\!\mathit{ont}}
\DeclareMathOperator{\Site}{Site}
\newcommand\MM{\mathbf{M}}
\def\limtop#1{{\displaystyle{\varprojlim_{#1}\!\mbox{\rm top }}}}
\def\Limtop{\mathop{\varprojlim\!\mbox{\rm top}}}
\DeclareMathOperator{\spec}{spec}
\def\ss{\scriptstyle{\mathscr{S}}}
\DeclareMathOperator{\Gal}{Gal}
\def\smallm{\scriptstyle{\mathscr{M}}}
\DeclareMathOperator{\End}{End}
\DeclareMathOperator{\ass}{ass}
\DeclareMathOperator{\zar}{zar}
\DeclareMathOperator{\Ogg}{Ogg}
\DeclareMathOperator{\res}{res}
\DeclareMathOperator{\tr}{tr}
\DeclareMathOperator{\Supp}{Supp}
\DeclareMathOperator{\codim}{codim}
\DeclareMathOperator{\rk}{rk}
\DeclareMathOperator{\cur}{cur}
\DeclareMathOperator{\cdqc}{cdqc}
\newcommand\DD{\mathbf{D}}
\DeclareMathOperator{\cl}{cl}
\def\JJ{\mathbf{J}}
\def\FF{\mathbf{F}}
\def\BB{\mathbf{B}}
\DeclareMathOperator{\Alget}{Alget}
\DeclareMathOperator{\Ass}{Ass}
\DeclareMathOperator{\supp}{supp}
\DeclareMathOperator{\rad}{rad}
\DeclareMathOperator{\prof}{prof}
\DeclareMathOperator{\cint}{int}
\DeclareMathOperator{\sic}{sic}
\DeclareMathOperator{\Tor}{Tor}
\DeclareMathOperator{\tors}{tors}
\DeclareMathOperator{\Sym}{Sym}
\DeclareMathOperator{\parf}{parf}
\DeclareMathOperator{\tor}{tor}
\DeclareMathOperator{\defn}{defn}
\DeclareMathOperator{\Div}{Div}
\DeclareMathOperator{\Pol}{Pol}
\DeclareMathOperator{\Fex}{Fex}
\DeclareMathOperator{\Lie}{Lie}
\DeclareMathOperator{\tot}{tot}
\DeclareMathOperator{\Aut}{Aut}
\DeclareMathOperator{\HOM}{HOM}
\DeclareMathOperator{\pch}{pch}
\DeclareMathOperator{\Ch}{Ch}
\DeclareMathOperator{\EXT}{EXT}
\DeclareMathOperator{\PIC}{PIC}
\DeclareMathOperator{\Rf}{Rf}
\DeclareMathOperator{\inj}{inj}
\newcommand\VV{\mathbf{V}}

\def\Xv{%
\begin{array}[b]{l}
{\kern1pt\scriptstyle \vee}\\[-5pt]
X
\end{array}}

%%%%%%%
\newcommand\TT{\mathbf{T}}
\newcommand\XX{\mathbf{X}}
\newcommand\EE{\mathbf{E}}
%%%%%%

\def\qedsymbol{\textup{C.Q.F.D.}}
\let\CQFD\qed
\let\cqfd\qed

%%%%%\bigdot
\font\bigsymb=cmsy10 at 4pt
\def\bigdot{{\kern1.2pt\raise 1.5pt\hbox{\bigsymb\char15}}}
%%%%%

\def\Lright{\mathop{\vcenter{\hbox{\hbox{$\longrightarrow$}\lower8pt \hbox{\kern-16pt\hbox{$\longrightarrow$}}}}}}

\def\Aa{{\mathcal A}}
\def\Bb{{\mathcal B}}
\def\Cc{{\mathcal C}}
\def\Dd{{\mathcal D}}
\def\Ee{{\mathcal E}}
\def\Ff{{\mathcal F}}
\def\Gg{{\mathcal G}}
\def\Hh{{\mathcal H}}
\def\Ii{{\mathcal I}}
\def\Jj{{\mathcal J}}
\def\Kk{{\mathcal K}}
\def\Ll{{\mathcal L}}
\def\Mm{{\mathcal M}}
\def\Nn{{\mathcal N}}
\def\Oo{{\mathcal O}}
\def\Pp{{\mathcal P}}
\def\Qq{{\mathcal Q}}
\def\Rr{{\mathcal R}}
\def\Ss{{\mathcal S}}
\def\Tt{{\mathcal T}}
\def\Uu{{\mathcal U}}
\def\Vv{{\mathcal V}}
\def\Ww{{\mathcal W}}
\def\Xx{{\mathcal X}}
\def\Yy{{\mathcal Y}}
\def\Zz{{\mathcal Z}}

\newcommand{\CC}{\mathbf{C}}
\newcommand{\GG}{\mathbf{G}}
\newcommand{\KK}{\mathbf{K}}
\newcommand{\LL}{\mathbf{L}}
\newcommand{\NN}{\mathbf{N}}
\newcommand{\PP}{\mathbf{P}}
\newcommand\QQ{\mathbf{Q}}
\newcommand\RR{\mathbf{R}}
\newcommand\ZZ{\mathbf{Z}}

\def\to{\mathchoice{\longrightarrow}{\rightarrow}{\rightarrow}{\rightarrow}}
\def\from{\mathchoice{\longleftarrow}{\leftarrow}{\leftarrow}{\leftarrow}}
\def\mto{\mathchoice{\longmapsto}{\mapsto}{\mapsto}{\mapsto}}
\def\mfrom{\mathrel{\from\joinrel\mapstochar\,}}
\def\hto{\mathrel{\lhook\joinrel\to}}
\def\hfrom{\mathrel{\from\joinrel\rhook}}

\def\implies{\mathchoice{\Longrightarrow}{\Rightarrow}{\Rightarrow}{\Rightarrow}}
\def\implied{\mathchoice{\Longleftarrow}{\Leftarrow}{\Leftarrow}{\Leftarrow}}
\def\iff{\mathchoice{\Longleftrightarrow}{\Leftrightarrow}{\Leftrightarrow}{\Leftrightarrow}}

\def\To#1{\mathchoice{\xrightarrow{\labelstyle~~#1~}}{\stackrel{#1}{\longrightarrow}}{}{}}
\def\Mto#1{\mathrel{\mapstochar\To{#1}}}
\def\From#1{\mathchoice{\xleftarrow{\labelstyle~~#1~}}{\stackrel{#1}{\longleftarrow}}{}{}}

\newcommand{\isomto}{\overset{\sim}{\longrightarrow}}
\newcommand{\isomfrom}{\overset{\sim}{\longleftarrow}}

\newcommand{\IFF}{\unskip~$\iff$~\ignorespaces}
\newcommand{\THEN}{\unskip~$\implies$~\ignorespaces}

\let\leq\leqslant
\let\geq\geqslant
\let\le\leqslant
\let\ge\geqslant

\def\cf{{cf.}\kern.3em}
\def\Cf{{Cf.}\kern.3em}
\def\ie{{i.e.}\kern.3em}
\def\iev{{i.e.},\xspace}
%\def\loccito{loc.\kern2pt cit.\xspace}
%\def\loccitn{\textit{loc.\kern2pt cit.}\xspace}
%\def\loccit{\sisi{\loccito}{\loccitn}}
\def\EGA{\textup{EGA}\kern.2em}
\def\SGA{\textup{SGA}\kern.3em}

\let\oldS\S
\def\S{\oldS\kern .15em}
\def\ptbl{.\kern .2em}
\def\resp{resp.\kern.3em}
\def\Exp{Exp.\kern.3em}

\let\bar\overline
\let\hat\widehat
\let\tilde\widetilde
\let\dpl\displaystyle
\let\ldots\dots
\let\emptyset\varnothing
\let\skpt\skippointrait
\let\epsilon\varepsilon

\def\sfrac#1#2{{#1}/{#2}}

\DeclareMathOperator{\cd}{cd}
%\DeclareMathOperator{\Hhom}{\mathcal{H}\!\mathit{om}}
\DeclareMathOperator{\Hhom}{\mathbf{Hom}}
%\DeclareMathOperator{\Eext}{\mathcal{E}\!\mathit{xt}}
\DeclareMathOperator{\Eext}{\mathbf{Ext}}
\DeclareMathOperator{\Spec}{Spec}
\DeclareMathOperator{\Tr}{Tr}

\let\emptyset=\varnothing

%%%%%%%%%%%%%%%
%%% expose XVII
%%%%%%%%%%%%%%%

\makeatletter
\let\old@optexpose\optexpose
\newcommand{\empty@optexpose}[1]{}
\newcommand{\paren@optexpose}[1]{\old@optexpose{#1}(}
\newcommand{\refnoexpose}[1]{\let\optexpose\empty@optexpose\ref{#1}\let\optexpose\old@optexpose}%\refnoexpose{X.cor5.2} engendre "5.2" au lieu de "X 5.2"
\newcommand{\neweqref}[1]{\let\optexpose\paren@optexpose\textup{\ref{#1})}\let\optexpose\old@optexpose}%\neweqref{XII.eq4.1} engendre "XII (4.1)" au lieu de "(XII 4.1)"
\makeatother
\newcommand{\refvb}[2][VI${}_{\text{B}}$ ]{\sisi{#1\refnoexpose{#2}}{\ref{#2}}}

\newcommand{\leftexp}[2]{{\vphantom{#2}}^{#1}{#2}}
\newcommand{\AAA}{\mathbf{A}}
\newcommand{\Sf}{\mathfrak{S}}
\DeclareMathOperator*{\quotprojlim}{\text{``$\varprojlim$''}}
\DeclareMathOperator*{\quotinjlim}{\text{``$\varinjlim$''}}
\DeclareMathOperator{\Fl}{F\sisi{\ell}{l}}
\DeclareMathOperator{\Pl}{P\sisi{\ell}{l}}
%\DeclareMathOperator{\cl}{c\sisi{\ell}{l}}
\DeclareMathOperator{\tdf}{tdf}
\DeclareMathOperator{\rg}{rg}
\newcommand{\iInd}{\sisi{I}{i}nd}
\newcommand{\iind}{ind}
\newcommand{\ppro}{pro}
\newcommand{\tTor}{\sisi{T}{t}or}
\newcommand{\ttor}{tor}
\newcommand{\chap}{Chap.}
\newcommand{\Ffex}{\Ff\mathit{ex}}
\newcommand{\sheafk}{\underline{k}}
\newcommand{\sheafn}{\underline{n}}
\newcommand{\sheafGamma}{\underline{\Gamma}}

% Repeated references to an already typeset footnote must not use
% \footnotemark without a matching \footnotetext: Hyperref would then create
% a live annotation to a destination that cannot exist.  This helper keeps the
% current printed footnote mark and links it to an explicit stable note target.
\makeatletter
\newcommand{\sgarepeatedfootnotemark}[1]{%
  \begingroup
  \edef\@thefnmark{\thefootnote}%
  \hyperlink{#1}{\@makefnmark}%
  \endgroup}
\makeatother

\sisi{
\DeclareMathOperator{\simarrow}{\rule{.15mm}{1.8mm}\kern -3.4pt\rightsquigarrow}
\newcommand{\sheafell}{\underline{\ell}}
\newcommand{\sheafZZ}{\underline{\ZZ}}
}{
\newcommand{\nderef}[1]{la N.D.E.~(\ref{#1}) page~\pageref{#1}}
}

%%%%%%%%%%%%%%%%%
%%%%% expose XVII
%%%%%%%%%%%%%%%%%

% END INCLUDED SGA4_English_current_buildable_source_20260906/source/fichiers-configuration-latex/macros.tex


\externaldocument{../01/01_English.tex}
\externaldocument{../02/02_English.tex}
\externaldocument{../03/03_English.tex}
\externaldocument{../04/04_English.tex}
\externaldocument{../05/05_English.tex}
\externaldocument{../05bis/05bis_English.tex}
\externaldocument{../06/06_English.tex}
\externaldocument{../07/07_English.tex}
\externaldocument{../08/08_English.tex}
\externaldocument{../09/09_English.tex}

\setboolean{original}{false}

\begin{document}

\begin{center}
X. Cohomological dimension: first results\\
M. Artin
\end{center}

\selectlanguage{english}version: \showgitstatus
\tableofcontents

\else
\chapter{Cohomological dimension: first results}
%\addtocontents{toc}{M. Artin}
\begin{center}
M. Artin
\end{center}
\fi
\phantomsection\MakeLinkTarget*{sga4.exp.X}\label{sga4.exp.X}

% BEGIN INCLUDED SGA4_English_current_buildable_source_20260906/source/10/drafts/SGA4_X_LINES_43_382_English.texfrag
\setcounter{pageoriginal}{42}

\section{Introduction}\etiquette[1]{X.sec1}\pageoriginale

In this Exposé, unless expressly stated otherwise, ``sheaf'' means
``abelian sheaf''.

Let $X$ be a scheme and let $\ell\in \PP$. ``Recall'' that the
\emph{$\ell$-cohomological dimension} of $X$ is the largest integer
$\cd_{\ell}X=n$ (or $=\infty$ if no such number exists) such that
$H^n(X,F)\neq 0$ for at least one $\ell$-torsion sheaf $F$ on $X$.
If $X=\Spec A$, we write $\cd_{\ell} A=\cd_{\ell}\Spec A$. Taking
into account the dictionary (VIII 2), %% label
we thus recover the ``classical'' notion of [CG] when $A=k$ is a field.
Since we plainly have the formula
$$
\cd_{\ell}(\coprod_i X_i)=\displaystyle{\mathop{\sup}_i} \cd_{\ell}X_i
$$
the classical theory may, in the evident way, be extended to an
arbitrary artinian ring.\pageoriginale This theory is set out in
$(CG)=$ ``\sisi{Galois cohomology}{Galois cohomology}'' %% Galois?
by Serre. We shall add a few refinements in ${\rm n^{\circ}}$
\hyperref[X.sec2]{2}.

Consider, for example, the case of an irreducible algebraic variety
$X$, of dimension $n$, over a separably closed field $k$, and fix
$\ell\neq \car k$. A theorem to be found in ${\rm n^{\circ}}$
\hyperref[X.sec4]{4} gives the upper bound $2n$ for its cohomological
dimension; this is the ``true'' dimension of $X$ (suggested by the
topological analogy when $k=\CC$). This result is fairly elementary
once the theory for fields is known, and in fact the upper bound $2n$
can be improved except when $X$ is complete: there is another upper
bound, $2n-1$, which holds whenever $X$ is ``missing at least one
point''. The latter is the correct ``stable upper bound'', in the sense
that every $X$ of dimension $n$ contains an open subset of
cohomological dimension $2n-1$. A third upper bound is obtained for an
affine variety, where one obtains $n$ as the cohomological dimension.
This is a generalization of one of the Lefschetz theorems for
hyperplane sections. These last two assertions are deeper and can only
be proved later (XIV).

Note that cohomological dimension is not a local notion (in contrast
to what happens for paracompact spaces). For example, a henselian ring
with separably closed residue field has cohomological dimension zero,
but if the closed point is removed, one may obtain an arbitrarily large
number, namely $2n-1$ if $n=\dim A$ in the most important cases; this
also agrees with the topological analogy supplied by the cohomology of
$U-\{x\}$, where $U$ is a small open neighborhood of a point $x$ of a
complex analytic space.

\section{Auxiliary results over a field}\etiquette[2]{X.sec2}\pageoriginale

We have the following

\begin{enonce*}{Theorem 2.1}\etiquette[2.1]{X.thm2.1}
Let $k'/k$ be an extension of finite type of a field $k$, and let
$\ell\in \PP$. Then
$$
\cd_{\ell}k'\leqq \cd_{\ell}k+\deg.\tr.(k'/k){\quoi},
$$
If the inequality is strict and if $\ell\neq \car k$, then $\ell=2$,
$k$ is orderable, and $k'$ is not orderable (\ie $-1$ is a sum of
squares in $k'$).
\end{enonce*}

Recall that when $\ell=\car k$ one has $\cd_{\ell}k'\leqq 1$ (CG II
\hyperref[X.thm2.2]{2.2}), so the question of equality in
\hyperref[X.thm2.1]{2.1} scarcely arises. Moreover,
$\cd_2k=\infty$ if $k$ is orderable. Indeed, the algebraic closure
$\bar{k}$ of $k$ contains a subextension $\widetilde{k}$ (a maximal
ordered subfield of $\bar{k}$) with
$[\bar{k}:\widetilde{k}]=2$. Plainly,
$\cd_2\widetilde{k}=\infty$, and
$\cd_2\widetilde{k}\leqq \cd_2 k$ by (CG I Prop.~14).

Proof. Almost all of this is well known. The inequality\sisi{}{,} as
well as equality if $\cd_{\ell}k<\infty$, are results of Tate (CG II
4.1.2). It remains to prove that $\cd_{\ell}k'<\infty$ and
$\cd_{\ell}k=\infty$ \sisi{implies}{imply} that $k$ is orderable and
$\ell=2$, and it is enough to treat the following two cases:
\begin{itemize}
\item[a)] $[k':k]<\infty$,

\item[b)] $k'=k(x)$ with $x$ transcendental over $k$.
\end{itemize}
In case a) we may suppose $k'/k$ separable. Then the Galois group
$H=G(\bar{k}/k')$ is an open subgroup of $G=G(\bar{k}/k)$, and by a
result of Serre [5] one has $\cd_{\ell}H=\cd_{\ell}G$ unless $G$
contains an element of order $\ell$. Since an algebraically closed
field $\bar{k}$ cannot have a subfield $\widetilde{k}$ of finite index
unless $\ell=2$ and $\widetilde{k}$ is orderable [7], the result
follows.

Consider\pageoriginale case b). Let $R$ be the henselization of
$\Spec k[x]$ at the point $x=0$, which is a discrete valuation ring,
and let $K$ be the fraction field of $R$. Then
$\cd_{\ell}k(x)<\infty$ implies $\cd_{\ell}K<\infty$. It is therefore
enough to prove that $\cd_{\ell}K<\infty$ implies
$\cd_{\ell}k<\infty$, which is a consequence of the following theorem:

\begin{enonce*}{Theorem 2.2}\etiquette[2.2]{X.thm2.2}
\footnote{This theorem was proved recently by $Ax$
    [1].}. Let $R$ be a henselian discrete valuation ring, with
    fraction field $K$ and residue field $k$. Let $\ell\in
    \PP$. We have
$$
\cd_{\ell}K=\cd_{\ell}k+1
$$
in the following two situations:
\begin{enumerate}
\renewcommand{\theenumi}{\roman{enumi}}
\renewcommand{\labelenumi}{(\theenumi)}
\item $\ell\neq \car k$.

\item $\ell \neq \car K$, $k$ perfect, and $\cd_{\ell}k<\infty$.
\end{enumerate}
\end{enonce*}

\begin{enonce*}[remark]{Remark}
Naturally, if $\ell=\car K$ one has $\cd_{\ell}K\leq 1$ in every case
(CG II \hyperref[X.thm2.2]{2.2}). It therefore remains to study the
relation of $\cd_{\ell}K$ with, perhaps, the module of absolute
differentials $\Omega^1_k$ (\EGA IV 20.4.3), in the unequal
characteristic case when $\car k=\ell$, $k$ imperfect.
\SGAArchiveOnly{\nde{The frozen source calls for a citation to work of Kato but supplies
no bibliographic target.}}
\end{enonce*}

Proof of the theorem. As for Theorem (\hyperref[X.thm2.1]{2.1}), most
of it is due to Tate. Case (ii) is treated explicitly in (CG II
Prop.~12) if $R$ is complete, and we reduce to that case by means of the
following lemma:

\begin{enonce*}{Lemma 2.2.1}\etiquette[2.2.1]{X.lem2.2.1}
Let $R$ be a henselian discrete valuation ring with fraction field $K$,
and let $\widehat{K}$ be the fraction field of the completion
$\widehat{R}$ of $R$. Then
$$
G(\bar{\widehat{K}}/\widehat{K})\simeq G(\bar{K}/K){\quoi}.
$$
\end{enonce*}

Proof. This means that the functor $L\mapsto
\widehat{L}=L\otimes_K \widehat{K}$ from\pageoriginale the category of
étale $K$-algebras to the category of étale $\widehat{K}$-algebras is
an equivalence. Now every étale algebra $\widehat{K}'$ over
$\widehat{K}$ is induced by a suitable étale algebra $K'$ over $K$.
Indeed, if $\widehat{K}'=\widehat{K}[\alpha]$ where $\alpha$ satisfies
the separable irreducible equation $f(T)=0$ with coefficients in
$\widehat{R}$, it is enough to take as $K'$ the algebra
$K[\alpha_0]$, where $\alpha_0$ is a root of an equation $f_0(T)=0$
with coefficients in $R$, with $f_0\equiv f(\widehat{m}^N)$, $N$ large
and $\widehat{m}$ the maximal ideal of $\widehat{R}$. This is an
application of the usual ``Newton method'' (Bourbaki,
Alg.~Comm.~III 4.5). It remains to prove that the functor is fully
faithful (it is plainly faithful), and this amounts to proving that if
$K'$ is a \emph{field}, then so is $\widehat{K}'$. Let $R'$ be the
normalization of $R$ in $K'$, which is finite over $R$ since $K'/K$ is
separable. The ring $R'$ is a discrete valuation ring, and the
completion of $R'$ is $\widehat{R}'=R'\otimes_R \widehat{R}$. Hence
$\widehat{R}'$ is a discrete valuation ring, and since $K'$ is its
total ring of fractions, it is indeed a field.

Consider assertion (i) of
\sisi{\eqhyperref[X.eq3.1]{3.1}}{\hyperref[X.thm2.2]{2.2}}.
If $\cd_{\ell}k<\infty$, we are again reduced by
\hyperref[X.lem2.2.1]{2.2.1}
to (CG II Prop.~12) (the hypothesis in the statement of that
proposition that $k$ be perfect is \sisi{used}{used} in the proof only
in the case of residue characteristic $\ell$). It therefore remains to
prove that $\cd_{\ell}K<\infty$ implies $\cd_{\ell}k<\infty$.

Let $\bar{G}=G(\bar{k}/k)=G(\widetilde{K}/K)$, where $\widetilde{K}$ is
the maximal unramified extension of $K$ (\ie the fraction field of the
discrete valuation ring $\widetilde{R}$ with separably closed residue
field, the strict henselization of $R$), let
$G=G(\bar{K}/K)$, $H=G(\bar{K}/\widetilde{K})$, so that
$\bar{G}=G/H$. Let $\bar{G}'\subset \bar{G}$ be an $\ell$-Sylow
subgroup, and let $k'/k$, $K'/K$ be the extensions corresponding to
$\bar{G}'$. We have $\cd_{\ell}k<\infty$ if and only if
$\cd_{\ell}k'<\infty$ (CG I Prop.~14), so by replacing $R$ with its
normalization in $K'$, we reduce to the case in which $\bar{G}$ is an
$\ell$-group. It is then enough to prove that if
$\cd_{\ell}G<\infty$, then $H^N(\bar{G},\ZZ/\ell)=0$ for sufficiently
large $N$ (CG I Prop.~20).

We\pageoriginale have $\cd_{\ell}H=1$ by the part of i) already proved,
applied to $\widetilde{R}$. Moreover, we have
\begin{equation*}
H^1(H, 𝛍_{\ell})=H^1(\Spec \widetilde{K},
𝛍_{\ell})=\widetilde{K}^{\ast}/\widetilde{K}^{\ast\ell}=\ZZ/\ell{\quoi}.\tag{2.2.2}
\end{equation*}\etiquette[2.2.2]{X.eq2.2.2}
Indeed, $H^1(\Spec \widetilde{K},
𝛍_{\ell})=\widetilde{K}^{\ast}/\widetilde{K}^{\ast\ell}$
by Kummer theory \hyperrefext[IX][IX.Thr3.2]{3.2}. We have an exact
sequence
$$
0\longrightarrow \widetilde{R}^{\ast}\longrightarrow
\widetilde{K}^{\ast}\longrightarrow
\ZZ\longrightarrow 0{\quoi},
$$
and $\widetilde{R}^{\ast}$ is divisible by $\ell$ since
$\widetilde{R}$ is strictly local and $\ell$ is invertible in
$\widetilde{R}$, whence
\sisi{\hyperref[X.lem2.2.1]{2.2.1}}{\eqhyperref[X.eq2.2.2]{2.2.2}}.

Consider the Hochschild--Serre spectral sequence
$$
E^{pq}_2=H^p(\bar{G}, H^q(H,
𝛍_{\ell}))\Longrightarrow
H^{p+q}(G,𝛍_\ell){\quoi}.
$$
We have $E^{pq}_2=0$ if $q>1$. If, moreover,
$\cd_{\ell}G<\infty$, the ``boundary'' morphism
$$
H^{n-1}(\bar{G}, H^1(H,𝛍_{\ell}))\longrightarrow
H^{n+1}(\sisi{\bar{G}}{G}, 𝛍_{\ell})
$$
must be bijective for sufficiently large $n$, and it is necessary to
prove that then
$$
H^{n-1}
(\bar{G},H^1(H,𝛍_{\ell}))=H^{n-1}(\bar{G},\ZZ/\ell)=0{\quoi}.
$$

Set $R'=R[x']$ with ${x'}^{\ell}=x$, where $x$ is the local parameter
of $R$, and let $K'$ be the fraction field of $R'$,
$G'=G(\bar{K}/\sisi{k}{K}')$, \etcn The residue field of $R'$ is $k$,
since we have supposed that $\bar{G}$ is an $\ell$-group, and hence that
the $\ell$\sisi{-th}{-th} roots of unity are in $K$. We have a diagram
$$
\xymatrix@M=.25cm@C=1.5cm{
H' \ar[r] \ar@{^{(}->}[d]& G' \ar[r] \ar@{^{(}->}[d]&\bar{G}\\
H  \ar[r]& G \ar[r] & \bar{G},
}
$$
whence\pageoriginale a morphism of Hochschild--Serre spectral
sequences, which gives a commutative diagram
$$
\xymatrix{
H^{n-1}(\bar{G},H^{1}(H',𝛍_l)) \ar[rd]& &\\
                                              &
H^{n+1}(\sisi{\bar{G}}{G},𝛍_l) &\\
H^{n-1}(\bar{G},H^{1}(H,𝛍_l))\ar[uu]^{\varepsilon}
\ar[ru] & &.
}
$$
Since $\cd_{\ell}G<\infty$ implies $\cd_{\ell}G'<\infty$, all the
morphisms in this diagram are isomorphisms for large $n$. But the
morphism
$$
\ZZ/r=H^1(H,𝛍_{\ell})\longrightarrow
H^1(H',𝛍_{\ell})=\ZZ/\ell
$$
which induces the arrow $\epsilon$ is plainly zero, and hence
$H^{n-1}(\bar{G},\ZZ/\ell)=0$,\cqfd

\begin{enonce*}{Theorem 2.3}\etiquette[2.3]{X.thm2.3}
Let $A$ be a noetherian henselian ring of dimension $1$, and let
$\ell\in \PP$ be invertible in $A$. Let $k$ be the residue field of
$A$ and let $R(A)$ be the ring of rational functions
%% say that this is the total ring of fractions? [with an EGA reference]?
on $\Spec A$. We have
$$
\cd_{\ell}R(A)\leq \cd_{\ell}k+1{\quoi},
$$
and equality holds if $\ell\neq 2$ or if $k$ is not orderable.
\end{enonce*}

Proof. Replacing $A$ by $A/P$ for a minimal prime ideal $P$ of $A$ and
applying VIII 1, %% missing reference \XXX
we see that we may suppose $A$ integral. Let $K$ be its fraction field,
and let $R$ be its normalization, which is local and henselian, and
hence a discrete valuation ring (it is noetherian by Krull--Akizuki
(Bourbaki, Alg. Comm. Chap.~7, \S 2, prop.~5)), and let $k'$ be the
residue field of $R$. It is known that $[k':k]<\infty$ (\loccit). Now
$K$ is the fraction field of $R$, and therefore
$$
\cd_{\ell}K=\cd_{\ell}k'+1\leq \cd_{\ell}k+1
$$
by\pageoriginale \hyperref[X.thm2.2]{2.2}, with equality under the last
hypothesis of \hyperref[X.thm2.3]{2.3}, by
\hyperref[X.thm2.1]{2.1}.

\begin{enonce*}[remark]{Example}
Equality does not hold for $\ell=2$, when $A$ is the henselization at
the origin of $\RR[x,y]/(x^2+y^2)$.
\end{enonce*}

\begin{enonce*}{Corollary 2.4}\etiquette[2.4]{X.cor2.4}
Let $A$ be a noetherian local ring with residue field $k$, and let
$\ell\in \PP$ be invertible in $A$. Suppose that $\ell\neq 2$
\sisi{or}{or} that no residue field of $\Spec A$ is orderable. Then
$$
\cd_{\ell}R(A)\geq \cd_{\ell}k+\dim A{\quoi},
$$
where $R(A)$ is the ring of rational functions on $\Spec A$.
\end{enonce*}

Proof. We may suppose $A$ integral, \sisi{with}{with} fraction field
$R(A)=K$. Let $x\in 𝔪_A$\sisi{ be a non-zero-divisor}{}.
%% harmonize maximal/prime ideals \XXX
Then $\dim A/(x)=\dim A-1$. By induction on $\dim A$, we have
$$
\cd_{\ell}R(A/(x))\geq \cd_{\ell}k+\dim A-1{\quoi}.
$$
Let $k'$ be a residue field of $R(A/(x))$ such that
$\cd_{\ell}k'=\cd_{\ell}R(A/(x))$ (\cf VIII 1), and let $P$ %% 𝔭
be the prime ideal that is the kernel of $A\rightarrow k'$. The ring
$A_P$ has dimension $1$. It is therefore enough to prove the
corollary for the local ring $A_P$; that is, we are reduced to the case
in which $\dim A=1$. Let $\widetilde{A}$ be the henselization of $A$.
Every residue field of $R(\widetilde{A})$ identifies with a separable
extension of $R(\widetilde{A})=K$, and therefore
$$
\cd_{\ell}R(A)=\cd_{\ell}K\geq \cd_{\ell}R(\widetilde{A}){\quoi}.
$$
We are thus reduced to the case in which $A$ is henselian of dimension
1, which follows from \hyperref[X.thm2.3]{2.3}.

\begin{enonce*}{Corollary 2.5}\etiquette[2.5]{X.cor2.5}
Let\pageoriginale $X$ be a noetherian scheme and let $\ell\in
\PP$ be invertible on $X$. Then
$$
\cd_{\ell}R(X)\geq \dim X{\quoi}.
$$
\end{enonce*}

Indeed, let $x\in X$ be a point such that $\dim X=\dim \Oo_{X,x}$ (or
such that $\dim \Oo_{X,x}$ is very large if $\dim X=\infty$). Let $A$
be the strict henselization of $\Oo_{X,x}$. Then
$\cd_{\ell}R(X)\geq \cd_{\ell}R(\Oo_{X,x})\geq \cd_{\ell}R(A)$ by
the usual argument. Moreover, $\dim A=\dim X$ (\resp is very large).
If $\ell=2$, and if $\ell$ is invertible on $X$, then the equation
$T^2+1=0$ is separable over $A$, so $-1$ is a square in $A$ because
$A$ is strictly local; hence the residue fields of $A$ are not
orderable. Thus the hypotheses of \hyperref[X.cor2.4]{2.4} are
satisfied for $A$ in every case, and
$\cd_{\ell}R(A)\geq \dim A$, whence the result.

% END INCLUDED SGA4_English_current_buildable_source_20260906/source/10/drafts/SGA4_X_LINES_43_382_English.texfrag


% BEGIN INCLUDED SGA4_English_current_buildable_source_20260906/source/10/drafts/SGA4_X_LINES_383_678_English.texfrag
\section{Field of fractions of a strictly local ring}\etiquette[3]{X.sec3}

\setcounter{subsection}{-1}
\subsection{}\etiquette[3.0]{X.subsec3.0}
In the following sections we shall have essential need of knowing
$\cd_{\ell}R(A)$, where $A$ is a strictly local ring, that is, a
henselian ring with separably closed residue field. Unfortunately,
$\cd_{\ell}R(A)$ is not known in general, for example if
$A=\widehat{\ZZ}_p[[x]](p\neq \ell)$.

It is nevertheless conjectured\SGAArchiveOnly{\nde{The frozen source states that this
was proved by O.~Gabber but supplies no bibliographic target.}} that if
$A$ is noetherian and $\ell$ is invertible in $A$, then one
\sisi{will have}{has}
\begin{equation*}
\cd_{\ell}R(A)=\dim A\quad (\ell\mbox{ invertible in } A){\quoi},\tag{3.1}
\end{equation*}\etiquette[3.1]{X.eq3.1}
at least in the most important cases, for example if $A$ is an
\emph{excellent} ring (\EGA IV 7.8.2).

Note that if $\dim A=1$, \eqhyperref[X.eq3.1]{3.1} is true by
\hyperref[X.thm2.3]{2.3}. Moreover, the inequality $\geq$ holds in any
case by \hyperref[X.cor2.5]{2.5}. We shall prove later (XIX), using
resolution of singularities \cite{sga4.exp.X.bib.3}, that \eqhyperref[X.eq3.1]{3.1}
is also true if $A$ is excellent of residual\pageoriginale
characteristic zero. Here we shall prove \eqhyperref[X.eq3.1]{3.1}
only in the easy case of \hyperref[X.prop3.2]{3.2} below.

Naturally one has $\cd_{\ell}R(A)\leq 1$ if $A$ has characteristic
$\ell$, so this case is trivial. It remains, moreover, to analyze the
case of unequal characteristics with residual characteristic $\ell$,
and one does not expect \eqhyperref[X.eq3.1]{3.1} to be true in this
\sisi{case there}{case}. It will certainly be necessary to take into
account the rank of the module $\Omega^1_k$ ($k$ the residue field),
\ie the degree $[k:k^{\ell}]$.\SGAArchiveOnly{\nde{The frozen source points to work of
Gabber and Orgogozo but supplies no bibliographic target.}}

\begin{enonce*}{Proposition 3.2}\etiquette[3.2]{X.prop3.2}
Let $X$ be a prescheme of finite type over the spectrum of a field $k$
or over the spectrum of a discrete valuation ring $R$, and let
$\ell\in \PP$ be invertible in $k$ (\resp~$R$). Then $\hyperref[X.eq3.1]{3.1}$ is true for
every ring $A$ that is a strict localization of $X$.
\end{enonce*}

\emph{Proof}. Take $X$ integral (one reduces to this case as usual),
and let $n=\dim X$.

Suppose that $X$ is defined over a \emph{separably closed} field $k$.
Then for every ring $A$ that is a strict localization of $X$ at a
\emph{closed} point, one has $\dim A=n$. Since every residue field $K$
of $R(A)$ is identified with an algebraic extension of $R(X)$, and
since
$$
\deg.\tr.(R(X)/k)=\dim X=n{\quoi},
$$
%% \tr. etc. ugly: there is one space too many (tr .) \XXX
one has $\cd_{\ell}K\leq n$ by \hyperref[X.thm2.1]{2.1}, whence
$\cd_{\ell}R(A)\leq n$. The opposite inequality follows from
\hyperref[X.cor2.5]{2.5}.

Now suppose that $X$ is of finite type over $\Spec R$, where $R$ is a
discrete valuation ring, and let $A$ be the strict localization of $X$
at a point $x$. If $x$ does not lie above the closed point of $\Spec
R$, one reduces immediately to the case where $X$ is of finite type
over a field (by localizing $R$). For the same reason, one may also
suppose that the morphism $X\rightarrow \Spec R$ is dominant.

Suppose\pageoriginale for the moment, in addition, that $R$ is a
\emph{strictly local} discrete valuation ring and that the point $x$
is \emph{closed} in $X$. Then one has $\dim A=\dim X=n$ (\EGA IV
5.6.5). Since every residue field $K$ of $R(A)$ is identified with a
separable extension of $R(X)$, one has
$$
\cd_{\ell}R(A)\leq \cd_{\ell}R(X)\leq
\cd_{\ell}L+\deg.\tr.(R(X)/L)=1+(n-1)=n
$$
by \hyperref[X.thm2.1]{2.1} and \hyperref[X.thm2.2]{2.2}, where $L$ is
the field of fractions of $R$; the result follows in this case, taking
\hyperref[X.cor2.5]{2.5} into account.

The general case is deduced from these special cases by means of the
following lemma:

\begin{enonce*}{Lemma 3.3}\etiquette[3.3]{X.lem3.3}
\begin{enumerate}
\renewcommand{\theenumi}{\roman{enumi}}
\renewcommand{\labelenumi}{(\theenumi)}
\item Let $X$ be of finite type over the spectrum of a field $k$, and
  let $A$ be a strict localization ring of $X$ at a point $x$. Then
  there exist a separably closed field $k'$, a scheme $X'$ of finite
  type over $\Spec k'$, and a point $x'$ closed in $X'$ such that $A$
  is isomorphic to a strict localization ring of $X'$ at $x'$.

\item Let $X$ be of finite type over $\Spec R$, where $R$ is a
  discrete valuation ring, let $x$ be a point of $X$ above the closed
  point of $\Spec R$, and let $A$ be a strict localization of $X$ at
  $x$. There exist a strictly local discrete valuation ring $R'$, a
  scheme $X'$ of finite type over $\Spec R'$, and a closed point $x'$
  of $X'$ above the closed point of $\Spec R'$, such that $A$ is
  isomorphic to a strict localization of $X'$ at $x'$.
\end{enumerate}
\end{enonce*}

Proof.
\begin{enumerate}
\renewcommand{\theenumi}{\roman{enumi}}
\renewcommand{\labelenumi}{(\theenumi)}
\item We may take $X$ affine.

Let $Y$ be the closure of $x$ in $X$, with the induced reduced
structure, and let $d=\dim Y$. By the normalization lemma (\cite{sga4.exp.X.bib.2}
Chap. V \S 3 ${\rm n^{\circ}}$ 1), there exists a morphism
$$
X\longrightarrow \Spec k[y_1,\ldots, y_d]=T
$$
inducing\pageoriginale a \emph{finite} morphism $Y\rightarrow T$.
Consequently, the strict localization ring of $\Oo_{X,x}$ is
isomorphic to the strict localization of $\Oo_{X',x'}$, where
$X'=\fprod{X}{\Spec\bar{R(T)}}{T}$, and where $x'$ is a closed point
of $X'$.

\item This case is treated similarly. We may suppose that $X$ is
affine, with ring $A$. Let $Y$ be the closure of $x$ in $X$ with the
induced reduced structure; it is a closed subscheme of the closed
fiber $X_0$ of $X$ over $\Spec R$, and let $d=\dim Y$. Then, lifting
the elements $y_1,\ldots,y_d$ of $A\otimes_R k$ considered in (i) to
elements of $A$, one obtains a morphism
$$
X\longrightarrow \Spec R[y_1,\ldots,y_d]=T
$$
inducing a \emph{finite} morphism $Y\rightarrow T_0$. Let $R'$ be the
strict localization of $R[y_1,\ldots,y_d]=B$ at the prime ideal
$𝔪B$. %% specify that 𝔪 is the maximal ideal of R? \XXX
Then $R'$ is a strictly local discrete valuation ring, and one may
take $X'=\fprod{X}{\Spec R'}{T}$, and $x'$ a closed point of $X'$ above
$x$.
\end{enumerate}

\section{Cohomological dimension: the case where $ℓ$ is invertible in $𝒪_X$}\etiquette[4]{X.sec4}

\begin{enonce*}{Theorem 4.1}\etiquette[4.1]{X.thm4.1}
Let $X$ be a noetherian scheme, let $\ell\in \PP$, and let $\varphi$
be a function on $X$ with values in $\NN$ and satisfying the following
conditions:
\begin{enumerate}
\renewcommand{\theenumi}{\roman{enumi}}
\renewcommand{\labelenumi}{(\theenumi)}
\item For every $x\in X$, $\cd_{\ell}k(x)\leq \varphi (x)$.

\item Let $x\in X$ and let $y\in Y=$ the closure of $x$, with $y\neq
x$. Let $K$ be the ring of rational functions of a strict localization
of $Y$ above $y$. One has $\cd_{\ell}K<\varphi(x)-\varphi(y)$.
\end{enumerate}
Then for every $\ell$-torsion sheaf $F$ on $X$, one has
\begin{gather*}
H^q(X,F)=0\\
q>\varphi(F)\mathop{=}_{\text{\rm def}}\Sup\{\varphi(x)\mid x\in X,
F_{\bar{x}}\neq 0\}{\quoi}.\tag*{if}
\end{gather*}
\end{enonce*}

Proof.\pageoriginale
The hypotheses will also be satisfied by the restriction of
$\varphi$\sisi{,}{} to any subprescheme of $X$, so by noetherian
induction we may suppose the theorem true for every closed subprescheme
of $X$ distinct from $X$. Let $F$ be given, and write
$F=\varinjlim F_i$, where the $F_i$ are the constructible subsheaves
of $F$ (\hyperrefext[IX][IX.prop2.9]{2.9} (iii)).

For every $F_i$, the set of $x\in X$ such that $F_x\neq 0$ is
constructible in $X$, and we may apply the induction hypothesis to
every $F_i$ whose support is not dense. We thus see that the induction
hypothesis in fact implies that the theorem is true for every $F$ that
vanishes at at least one of the maximal points of $X$. We shall
therefore suppose that $F$ does not satisfy this condition.

Now let $i:\Spec R(X)\rightarrow X$ be the inclusion, where $R(X)$ is
the ring of rational functions of $X$, and consider the following
exact sequence, where $K$ and $C$ are defined as kernel and cokernel,
respectively:
$$
0\longrightarrow K\longrightarrow F\longrightarrow i_{\ast}
i^{\ast}F\longrightarrow C\longrightarrow 0{\quoi}.
$$
Here $K$ and $C$ vanish at every maximal point of $X$, so by (ii),
$\varphi(K)$ (\resp $\varphi(C)$) is strictly less than
$n=\varphi(F)$, whence $H^q(X,K)=H^q(X,C)=0$ for $q\geq n$. Thus,
\sisi{}{,} in order to prove $H^q(X,F)=0$ for $q>n$, it is enough, by
the cohomology exact sequences, to prove
$H^q(X,i_{\ast}i^{\ast}F)=0$; hence we may suppose that $F$ is of the
form $i_{\ast}G$, for a suitable $\ell$-sheaf $G$ on $\Spec R(X)$.

Consider the Leray spectral sequence\sisi{.}{}
$$
E^{p,q}_2=H^p(X,R^qi_{\ast}G)\Longrightarrow H^{p+q}(\Spec
R(X),G){\quoi}.
$$
Now $R(X)_{\text{\rm réd}}$ is the product of the fields $k(x)$ for
$x$ maximal in $X$, so $\cd_{\ell}R(X)\leq n$ by (i); therefore the
abutment of the spectral sequence vanishes in
\sisi{$\dim > n$}{dimension strictly greater than $n$}.
We\pageoriginale wish to prove that $E^{p0}_2=0$ for $p>n$. Examining
the spectral sequence, one sees that it is enough to prove that
$E^{pq}_2=0$ if $q>0$ and if $p>n-q-1$. Thus, by induction, it is
enough to prove that $\varphi(R^qi_{\ast}G)<n-q$ for $q>0$. But by
VIII \hyperref[VIII.thm5.2]{5.2}, the fiber of $R^qi_{\ast}G$ at a geometric point $\bar{y}$
above a point $y\in X$ is precisely $H^q(\Spec K, G_K)$, where $K$ is
the ring of rational functions of the strict localization of $X$ at
$\bar{y}$, and $G_K$ is the induced sheaf on $\Spec K$. By (ii), one
has $\cd_{\ell}K<n-\varphi(y)$, whence
$(R^qi_{\ast}G)_{\bar{y}}=0$ if $q\geq n-\varphi(y)$, which gives the
desired result.

\begin{enonce*}{Corollary 4.2}\etiquette[4.2]{X.cor4.2}
Let $X$ be a noetherian scheme, let $\ell\in \PP$, let
$n\in \NN$, and suppose that the following conditions are satisfied:
\begin{enumerate}
\renewcommand{\theenumi}{\roman{enumi}}
\renewcommand{\labelenumi}{(\theenumi)}
\item For every point $y$ of $X$ of codimension $c$, one has
  $\cd_{\ell}k(y)\leq n-2c$.

\item For every ring $A$ that is a strict localization of an
  irreducible closed subscheme $Y$ of $X$, one has
  $\cd_{\ell}R(A)\leq \dim A$.

Then
$$
\cd_{\ell}X\leq n{\quoi}.
$$
More precisely, if $F$ is an $\ell$-sheaf that vanishes at every point
$y$ of codimension $<s$, one has
$$
H^q(X,F)=0\mbox{ if } q>n-2s{\quoi}.
$$
\end{enumerate}
\end{enonce*}

Indeed, one may take $\varphi(y)=n-2\codim y$ and apply
\hyperref[X.thm4.1]{4.1}, whose condition (i) (\resp (ii)) follows
from the corresponding condition (i) (\resp (ii)) of
(\hyperref[X.cor4.2]{4.2}), taking into account for (ii) the inequality
$\dim \Oo_{Y,y}+\dim \Oo_{X,x}\leq \dim \Oo_{X,y}$.

\begin{enonce*}{Corollary 4.3}\etiquette[4.3]{X.cor4.3}
Let $X$ be of finite type and dimension $n$ over a field $k$, and let
$\ell\neq \car k$. Then $\cd_{\ell}X\leq 2n+\cd_{\ell}k$.
\end{enonce*}

One\pageoriginale applies the preceding corollary, replacing $n$ there
by $2n+cd_{\ell}k$, and using \hyperref[X.thm2.1]{2.1} and
\hyperref[X.prop3.2]{3.2} to verify conditions (i) and (ii) of
\hyperref[X.cor4.2]{4.2}, taking into account for (i) the inequality
$$
\deg.\tr.K(y)/k+\dim \Oo_{X,y}\leq \dim X{\quoi}.
$$

\begin{enonce*}{Corollary 4.4}\etiquette[4.4]{X.cor4.4}
Let $X$ be noetherian, and let $\ell\in \PP$ be invertible on $X$.
Suppose that condition {\rm (ii)} of $\hyperref[X.cor4.2]{4.2}$ is satisfied, and moreover
that $\ell\neq 2$ \sisi{where}{or} that no residue field of $X$ is
orderable. Then one has
$$
\cd_{\ell}X\leq \cd_{\ell}R(X)+\dim X\leq 2\cd_{\ell}R(X){\quoi}.
$$
\end{enonce*}

Put $n=\cd_{\ell}R(X)+\dim X$ in \hyperref[X.cor4.2]{4.2}, and apply
\hyperref[X.cor2.4]{2.4}, which implies condition (i) of
\hyperref[X.cor4.2]{4.2}, since
\[\cd_{\ell}k(y)\leq \cd_{\ell}R(\Oo_{X,y})-\dim \Oo_{X,y}\leq
\cd_{\ell}R(X)-\dim\Oo_{X,y}\leq n-2\dim \Oo_{X,y}.
\]
The second inequality in \hyperref[X.cor4.4]{4.4} follows from
\hyperref[X.cor2.5]{2.5}.

\begin{enonce*}[remark]{Example 4.5}\etiquette[4.5]{X.exem4.5}
The inelegant hypothesis on residue fields made in
\hyperref[X.cor4.4]{4.4} is necessary, as is seen by taking
$X=\Spec \RR[x,y,z]/(x^2+y^2+z^2)$, where all the other conditions are
satisfied, with $\dim X=\cd_{\ell}R(X)=\ell=2$, but the sheaf $\ZZ/2$
concentrated at the point $(0,0,0)$ has nonzero cohomology in every
degree, so $\cd_{\ell}X=+\infty$.
\end{enonce*}

\setcounter{subsection}{5}
\subsection{}\etiquette[4.6]{X.subsec4.6}
One is also tempted to try to impose hypotheses in
\hyperref[X.cor4.2]{4.2} only on the closed points of $X$, but this is
not possible. For example, there exist discrete valuation rings with
algebraically closed residue field and with function field of
arbitrarily large cohomological dimension\footnote[1]{To see this, it
  is enough to note that if $X$ is a smooth algebraic scheme over a
  field $k$, admitting a point $x\in X(k)$, then one can find a
  ``formal curve arc'' passing through $x$ whose ``algebraic envelope''
  is $X$, and whose function field therefore contains that $K$ of $X$
  and induces on it a discrete valuation whose residue field is $k$.}.
The spectrum of such a ring will likewise have very large
cohomological dimension.
%% to do (for the note): i -> (i) \XXX

% END INCLUDED SGA4_English_current_buildable_source_20260906/source/10/drafts/SGA4_X_LINES_383_678_English.texfrag


% BEGIN INCLUDED SGA4_English_current_buildable_source_20260906/source/10/drafts/SGA4_X_LINES_679_937_English.texfrag
\section{Cohomological dimension: the case $ℓ=p$}\etiquette[5]{X.sec5}\pageoriginale 

\setcounter{subsection}{-1}
\subsection{}\etiquette[5.0]{X.subsec5.0}
Using Artin\sisi{ - }{-}Schreier theory
(\hyperrefext[IX][IX.Thr3.5]{3.5}), one obtains a better upper bound
than \hyperref[X.cor4.3]{4.3} for the $p$-cohomological dimension of a
scheme $X$ of characteristic~$p$. Let $\cdqc X$ be the largest number
$n$ such that $H^n(X,F)\neq 0$ for at least one \emph{quasi-coherent}
sheaf of modules $F$ on $X$ (in the Zariski sense). We have the
following

\begin{enonce*}{Theorem 5.1}\etiquette[5.1]{X.thm5.1}
Let $X$ be a Noetherian scheme of characteristic $p>0$. Then
$$
\cd_pX\leq \cdqc X+1{\quoi}.
$$
In particular, the $p$-cohomological dimension of a Noetherian affine
scheme $X$ of characteristic $p>0$ is at most $1$.
\end{enonce*}

Proof. Apply \hyperrefext[IX][IX.prop5.5]{5.5}:

Taking account of the fact that the $R^q f_{\ast}F$, $q>0$, vanish for
a finite morphism $f$ and an abelian sheaf $F$ (VIII \hyperref[VIII.prop5.5]{5.5})
%% reference \XXX
(\resp a quasi-coherent sheaf $F$), we are reduced to proving that
$H^q(X,\sisi{i_1}{i_!}(\ZZ/p)_U)=0$ for $i:U\rightarrow X$ an open
immersion and for $q>\cdqc X$.

Let $Y$ be a closed subscheme of $X$ with support $X-U$, and let $J$
be the coherent sheaf of ideals defining $Y$. The morphism
$p:f\mapsto f^p-f$ (\hyperrefext[IX][IX.Thr3.5]{3.5}) induces a
morphism $J\rightarrow J$. Examining the exact sequence
\hyperrefext[IX][IX.Thr3.5]{3.5} and the exact sequence
$$
0\longrightarrow i_{!}(\ZZ/p)_U\longrightarrow
(\ZZ/p)_X\longrightarrow
(\ZZ/p)_Y\longrightarrow 0{\quoi},
$$
one sees that there is in fact an exact sequence
\begin{equation*}
0\longrightarrow i_{!}(\ZZ/p)_U\longrightarrow
J\longrightarrow J\longrightarrow 0{\quoi}.\tag{$\ast$}
\end{equation*}
Since $J$ is coherent, one has $H^q(X,J)=0$ for $q>\cdqc X$, whence the
result\pageoriginale by the long exact cohomology sequence.

For a scheme quasi-projective over a separably closed field $k$, this
upper bound can in certain cases be reduced \sisi{by}{by}~$1$. In fact,
we have

\begin{enonce*}{Corollary 5.2}\etiquette[5.2]{X.cor5.2}
Let $X$ be a scheme of finite type over a field $k$ of characteristic
$p>0$. One has
$$
\cd_pX\leq \dim X+1{\quoi},
$$
and if $k$ is separably closed and $X$ is quasi-projective
\footnote[1]{Using Chow's lemma, it is easy to replace the condition
  ``$X$ quasi-projective'' by ``$X$ separated.''},
one has
$$
\cd_pX\leq \dim X{\quoi}.
$$
\end{enonce*}

Indeed, the first assertion is clear, since $\cdqc X\leq \dim X$
(\cite{sga4.exp.X.bib.6} 4.15.2). Let us treat the second. For such an $X$, the
cohomology groups $H^n(X,F)$ $(n=\dim X)$, for a coherent sheaf $F$,
are finite-dimensional vector spaces over $k$. Let $V=H^n(X,J)$,
where $J$ is the sheaf of ideals occurring in the exact sequence
$(\ast)$. The morphism $\varphi:V\rightarrow V$ deduced from $(\ast)$
is evidently of the form $\varphi=\epsilon-\id$ \sisi{or}{where}
$\epsilon$ satisfies $\epsilon(rv)=r^p\epsilon(V)$ for $r\in k$ and
$v\in V$. The Jacobian of the map $\Phi$ of the affine scheme deduced
from $\Phi$ is therefore $-I$, \sisi{$I$}{where~$I$ is} the identity
matrix. Consequently $\Phi$ is étale, hence surjective because it is
additive, and it follows that $\Phi$ is surjective since $k$ is
separably closed, whence the conclusion immediately.

\section{Cohomological dimension for a prescheme of finite type over $\mathrm{Spec} ℤ$}\etiquette[6]{X.sec6}\pageoriginale

\begin{enonce*}{Proposition 6.1}\etiquette[6.1]{X.prop6.1}
Let $X$ be quasi-finite over $\Spec \ZZ$. Suppose that
$\ell\neq 2$ \sisi{where}{or} that every residue field of $X$ of
characteristic zero is totally imaginary. Then $\cd_{\ell}X\leq 3$.
\end{enonce*}

This is a consequence \sisi{of}{of} \sisi{Theorem}{Theorem}
\hyperref[X.thm2.2]{2.2} (ii), of \sisi{Corollary}{Corollary}
\hyperref[X.cor4.2]{4.2}, and of the fact that a number field (totally
imaginary if $\ell =2$) has $\ell$-cohomological dimension $\leq 2$
(CG II 4.4).

\sisi{It follows in fact from class field theory}{A consequence of
class field theory is} that equality holds if the morphism
$X\rightarrow \Spec \ZZ$ is proper and surjective.

Unfortunately, the results of ${\rm n^{\circ}}$ 4 do not apply as they
stand if $X$ is only of finite type over $\Spec \ZZ$, because
hypothesis (ii) of (\hyperref[X.cor4.2]{4.2}) is not in general
satisfied for local rings of \emph{residual characteristic} $\ell$.
One must therefore use result (\hyperref[X.thm5.1]{5.1}) to obtain the
correct upper bound; we shall now reap the benefit of Theorem
(\hyperref[X.thm5.1]{5.1}).

\begin{enonce*}{Theorem 6.2}\etiquette[6.2]{X.thm6.2}
Let $X$ be of finite type over $\Spec \ZZ$, and let $\ell\in \PP$.
Suppose, if $\ell=2$, that no residue field of $X$ is orderable. Then
$$
\cd_{\ell}X\leq 2\dim X+1{\quoi}.
$$
More precisely, let $F$ be a sheaf of $\ell$-torsion such that for
every point $x$ of $X$ such that $F_{\bar{x}}\neq 0$, one has
$$
\dim\{\bar{x}\}\leq
\begin{cases}
(N-1)/2 &\mbox{if\quad } \ell\neq \car k(x)\\
N-1     &\mbox{if\quad } \ell=\car k(x)
\end{cases}
$$
($\{\bar{x}\}$ is the closure of $\{x\}$); then
$$
H^p(X,F)=0\mbox{ if } p>N{\quoi}.
$$
\end{enonce*}

Proof.\pageoriginale By induction on $X$, we may suppose the theorem
true for every closed subscheme distinct from $X$. If $F$ is as in the
statement and we write $F=\varinjlim F_i$, where the $F_i$ are the
constructible subsheaves of $F$
(\hyperrefext[IX][IX.prop2.9]{2.9} (iii)), then every $F_i$ also
satisfies the same hypotheses. It therefore suffices to prove the
theorem when $F$ is constructible.

In this case, the set $E$ of $x\in X$ such that $F_{\bar{x}}\neq 0$ is
a constructible subset, and the induction hypothesis implies that the
theorem is true for every $F$ such that $E$ is not dense. Moreover, the
theorem is a consequence of \hyperref[X.thm5.1]{5.1} for a sheaf
concentrated on the fiber of $X$ over the point $(\ell)$ of
$\Spec \ZZ$, and of \hyperref[X.cor4.3]{4.3} for a sheaf concentrated
on the fiber over $(p)$, where $p\neq \ell$, taking account of the fact
that $\cd_{\ell}\ZZ/p \ZZ=1$ (CG II 2.2).

Take $X$ reduced, and let $X=X_{0}\cup X_1$, where $X_1$ is the closed
subset which is the union of the irreducible components of $X$ of
characteristic $\ell$, and where $X_{0}$ is the union of the other
irreducible components, $X_{0}$ and $X_1$ being endowed with their
induced reduced structures. Let $j_{\nu}:X_{\nu}\rightarrow X$ be the
inclusions, let $d_{\nu}=\dim X_{\nu}$, and let
$N=\sup \{2d_{0}+1, d_1+1\}$. For every constructible $F$ on $X$, there
is an exact sequence\sisi{.}{}
\begin{equation*}
0\longrightarrow F\longrightarrow
{j_{0}}_{\ast} j_{0}{}^{\ast}F\oplus {j_1}_{\ast} j_1
{}^{\ast}F\longrightarrow C\longrightarrow 0{\quoi},\tag{$\ast$}
\end{equation*}
where $C$ is concentrated on $X_{0}\cap X_1$, which has characteristic
$\ell$ and dimension $\leq d_1-1$. It follows that
$$
H^q(X,C)=0\mbox{ if } q>N-1{\quoi},
$$
because by \hyperref[X.thm5.1]{5.1} one has
$\cd_{\ell}(X_{0}\cap X_1)\leq \dim
\sisi{}{(}X_{0}\cap X_1\sisi{}{)}+1\leq d_1\leq N-1$. Examining the
relative long exact cohomology sequence of $(\ast)$, we are reduced to
the case $X=X_{0}$ \sisi{where}{or} $X=X_1$, and the latter is a
consequence of \hyperref[X.thm5.1]{5.1}, as already noted.

Suppose\pageoriginale therefore that every maximal point of $X$ has
characteristic different from $\ell$, and let
$i_{\nu}:x_{\nu}\rightarrow X$ be the inclusions of the maximal points.
Examining the exact sequence
$$
0\longrightarrow K\longrightarrow F\longrightarrow \bigoplus_{\nu}
i_{\nu\ast}i_{\nu}{}^{\ast}F\longrightarrow C\longrightarrow 0
$$
we are reduced by induction to the case
$F=i_{\nu\ast}i_{\nu}{}^{\ast}F$, that is, where $F$ is of the form
$i_{\ast}G$ for a suitable sheaf $G$ on $\Spec k(x)$, with $x$ a point
of $X$. We may therefore suppose that $X$ is irreducible and reduced,
with generic point $x$, and that the characteristic of $k(x)$ is zero
by \hyperref[X.cor4.3]{4.3}.

Consider the spectral sequence
\begin{equation*}
H^p(X,R^qi_{\ast}G)\Longrightarrow H^{p+q}(\Spec k(x),
G){\quoi}.\tag{${\ast}{\ast}$} 
\end{equation*}
Let $s=\dim X$. We must prove that
\begin{equation*}
E^{p0}_2=0\mbox{ for } p> 2s+1{\quoi}.\tag{$+$}
\end{equation*}
Now $k(x)[i]$ has transcendence degree $s-1$ over $\QQ[i]$
%% say that i = √(-1)? \XXX
which is a field of cohomological dimension~$2$. Since $k(x)$ is not
orderable if $\ell=2$, it follows from \hyperref[X.thm2.1]{2.1} that
$\cd_{\ell}k(x)=s+1$. This implies that the target of the spectral
sequence vanishes for $n>s+1$. It therefore suffices, for $(+)$, to
prove that in $({\ast}{\ast})$ one has
\begin{equation*}
E^{pq}_2=0\mbox{ if } p>2s-q\mbox{ and } q>0{\quoi}.\tag{${+}{+}$}
\end{equation*}

Let $y$ be a point of $X$, and let $K$ be the ring of fractions of a
strict localization at a geometric point $\widetilde{y}$ of $X$ above
$y$. Then $K$ is a direct product of fields, each of transcendence
degree $s-1$ over a strict localization of $\ZZ$, and hence, by
\hyperref[X.thm2.1]{2.1} and \hyperref[X.thm2.2]{2.2}, one has
$\cd_{\ell}K\leq s$. Moreover, one may apply
\hyperref[X.prop3.2]{3.2} if the residual characteristic of $y$ is not
$\ell$. Thus one has
\begin{equation*}
(R^qi_{\ast}G)_{\widetilde{y}}=0
\begin{cases}
\mbox{ if $q>\codim y$, \ie $\dim \bar{y}>s-q$, and $\ell\neq \car
k(y)$,}\\
\mbox{or if $q>s$ in all cases.}
\end{cases}
\tag{$''$}
\end{equation*}
In\pageoriginale particular, one has $E^{pq}_2=0$ if $q>s$. To verify
$({+}{+})$ in the case $0<q\leq s$, apply the induction hypothesis to
the sheaves $R^qi_{\ast}(G)$.

According to the statement and $('')$, one must set $N'=2s-q$ and
verify the two inequalities
\begin{align*}
& s-q\leq (N'-1)_{/2}\\
& s-1=\dim X_{\ell}\leq N'-1{\quoi},
\end{align*}
where $X_{\ell}$ is the \sisi{fibered space}{fiber} of $X$ over the
\sisi{point}{point} $(\ell)$ of $\Spec \ZZ$. They do indeed hold if
$0<q\leq s$\sisi{,}{.} \cqfd

\begin{thebibliography}{99}
\bibitem{sga4.exp.X.bib.1} 
Ax, J.  Proof of some conjectures of Serre on
cohomological dimension, Proc. Amer. Math. Society 16 (1965)
1214-1221.

\bibitem{sga4.exp.X.bib.2} Bourbaki, N. Algèbre \sisi{Commutative}{commutative}, Paris (1960)

\bibitem{sga4.exp.X.bib.3} 
Hironaka, H.  Resolution of singularities of an algebraic surface
over a field of characteristic zero, Annals of Math. Vol. 79, No. 1, 2
(1964)

 
\bibitem{sga4.exp.X.bib.4} Nagata, M. Local Rings, New York (1962)

\bibitem{sga4.exp.X.bib.5} 
Serre, J-P.  Sur la dimension cohomologique \sisi{de}{des} groupes profinis,
Topology Vol. 3, p. 413 (1965)


\bibitem{sga4.exp.X.bib.6} Godement, R. \sisi{Théorie des Faisceaux}{Topologie algébrique et théorie des faisceaux}, Paris (1958).

\bibitem{sga4.exp.X.bib.7} 
E. Artin et A.  Schreier - Eine Kennzeichnung der \sisi{reel}{reell}
abgeschlossenen Körper, Hamb. Abh. 5 (1927) pp. 225-231.


\bibitem{sga4.exp.X.bib.8} Serre, J.P. Cohomologie \sisi{Galoisienne}{galoisienne}, Lecture Notes
${\rm n^{\circ}}$ 5 (Springer).
\end{thebibliography}

\ifx\danslelivre\undefined
\end{document}
\fi

% END INCLUDED SGA4_English_current_buildable_source_20260906/source/10/drafts/SGA4_X_LINES_679_937_English.texfrag


% END INCLUDED SGA4_English_current_buildable_source_20260906/source/10/10_English.tex


\makeatletter
\def\input@path{{../11/}}
\makeatother
% BEGIN INCLUDED SGA4_English_current_buildable_source_20260906/source/11/11_English.tex
\ifx\danslelivre\undefined
% Release engine: XeLaTeX (the LuaTeX compatibility layer breaks Xy-pic output).
\documentclass[10pt]{../fichiers-configuration-latex/smfart}
% BEGIN INCLUDED SGA4_English_current_buildable_source_20260906/source/fichiers-configuration-latex/paquets.tex
\usepackage{polyglossia}
%\usepackage[francais]{babel}
\usepackage{fontspec}
\usepackage[math-style=ISO,bold-style=upright,sans-style=upright]{unicode-math}
%\usepackage[mathletters]{ucs}
%\usepackage[utf8]{inputenc}
%\usepackage[T1]{fontenc}
\usepackage[unicode,hypertexnames=false,colorlinks=true,linkcolor=blue,citecolor=Sepia,urlcolor=magenta]{hyperref}
\setmainlanguage{french}
\setmainfont[Ligatures=TeX]{Linux Libertine O}
% variante possible : TeXGyrePagella ?
\setmathfont{XITS Math}

% maths
\usepackage{amsmath}
%\usepackage{amssymb}
%\usepackage{mathrsfs}
%\usepackage{stmaryrd}
%\usepackage{euscript}

% couleurs
\usepackage[usenames,dvipsnames]{xcolor}

% graphiques
\usepackage{graphicx}
\usepackage[all]{xy}
%\usepackage{tikz}
%\usetikzlibrary{matrix,arrows,calc}
%\tikzset{column sep=3em,row sep=2em,text height=1.5ex, text depth=0.25ex,>=angle 90}
%\tikzset{description/.style={fill=white,inner sep=1pt}}

% hyperliens, 'etiquettes
\usepackage{xr-hyper}
\usepackage{../fichiers-configuration-latex/smfthm}
%\ifx\cacheretiquettes\undefined
%\usepackage[inline]{showlabels} \else \fi % pour rendre visibles les 'etiquettes
%\usepackage{../fichiers-configuration-latex/makeidx} % pour faire un index
%\usepackage[pagebackref,unicode]{hyperref}  % pour voir o`u sont cit'ees les r'ef'erences
% hyperref is loaded once above with the cumulative-reader destination policy.
%\hypersetup{breaklinks=true} %permet le retour `a la ligne dans les liens trop longs. Incompatible avec pdflatex
%\renewcommand*{\backref}[1]{{\small [cit\'e page(s)~#1.]}}

\usepackage{xspace}
\usepackage{pifont} % pour expose XVII

% END INCLUDED SGA4_English_current_buildable_source_20260906/source/fichiers-configuration-latex/paquets.tex

% BEGIN INCLUDED SGA4_English_current_buildable_source_20260906/source/fichiers-configuration-latex/numerotation.tex
\renewcommand{\thefootnote}{\roman{footnote}} % pour mettre les notes en bas de page sous la forme (ii) et non (2), qui ressemble trop (?) a un exposant.

\newcounter{pageoriginal}
\newcommand{\pageoriginale}{\refstepcounter{pageoriginal}\marginpar{\footnotesize\textbf{\thepageoriginal}}\xspace}
\let\pageoriginaled\pageoriginale

\usepackage{ifthen}
\newboolean{original}
\def\sisi{\ifthenelse{\boolean{original}}}
\setboolean{original}{false}

\newcounter{cnote}
\def\nde#1{\ifthenelse{\boolean{original}}{}{\stepcounter{cnote}\def\thefootnote{\arabic{cnote}}\protect\footnote[\value{cnote}]{N.D.E.: #1}}}
\def\ndemark{\ifthenelse{\boolean{original}}{}{\stepcounter{cnote}\def\thefootnote{\arabic{cnote}}\protect\footnotemark[\value{cnote}]}}
\def\ndetext#1{\ifthenelse{\boolean{original}}{}{\def\thefootnote{\arabic{cnote}}\protect\footnotetext[\value{cnote}]{N.D.E.: #1}}}

% END INCLUDED SGA4_English_current_buildable_source_20260906/source/fichiers-configuration-latex/numerotation.tex

% BEGIN INCLUDED SGA4_English_current_buildable_source_20260906/source/fichiers-configuration-latex/macros.tex
\setmainlanguage{french}
\setotherlanguage{english}


%% Fixed, engine-independent release identifier.  The inherited dynamic Git
%% probe used Lua io.popen, exposed build-state-dependent output, and printed
%% a visible question mark outside a Git checkout.
\def\showgitstatus{English release r2, 2026-07-23}


% macros pour LaTeX
\renewcommand{\mathsf}{}

% Police différente pour U+2113 SCRIPT SMALL L
\setmathfont[range="2113]{latinmodern-math.otf}

% Give every inherited semantic label its own stable PDF destination instead
% of inheriting whichever automatic equation/page anchor happened to be live.
% The legacy label remains as a compatibility alias used by existing edges.
\makeatletter
\def\etiquette[#1]#2{%
  \Hy@raisedlink{\hyper@anchorstart{sga4.semantic.#2}\hyper@anchorend}%
  \def\@currentHref{sga4.semantic.#2}%
  \def\@currentlabel{#1}%
  \label{#2}%
}
\makeatother
\def\eqhyperref[#1]#2{(\hyperref[#1]{#2})}
\def\pagehyperref[#1]#2{(page \hyperref[#1]{#2})}
\def\hyperrefext[#1][#2]#3{\hyperref[#2]{#1 #3}}
\def\eqhyperrefext[#1][#2]#3{\hyperref[#2]{#1 #3}}     %cf. exposes 17
\def\neweqhyperrefext[#1][#2]#3{\hyperref[#2]{#1 #3}}  %cf. exposes 17

\def\shat#1{\protect\widehat{#1}}
\def\stilde#1{\protect\widetilde{#1}}
%%%%%%%%%%%%%%%%%%


\def\etcn{etc.\xspace}
\def\etco{etc...\xspace}
\def\etc{\sisi{\etco}{\etcn}}
\def\ldotsn{\sisi{...}{\etcn}}

\def\cf{{cf.}\kern.3em}
\def\Cf{{Cf.}\kern.3em}
\def\ie{{i.e.}\kern.3em}
\def\iev{{i.e.},\xspace}
\def\loccito{loc.\kern2pt cit.\xspace}
\def\loccitn{\textit{loc.\kern2pt cit.}\xspace}
\def\loccit{\sisi{\loccito}{\loccitn}}
\def\EGA{\textup{EGA}\kern.2em}
\def\SGA{\textup{SGA}\kern.3em}

\def\quoi{}

\def\quav{\quad,\quad}
\def\quaet{\quad \text{et}\quad}
\def\fibreproduct#1#2#3{#1 ×_{#3} #2}
\let\fprod\fibreproduct
\def\ccheck#1#2{\mbox{${\check{\check{\rm #1}}}^{\lower2pt\hbox{$\scriptstyle #2$}}$}}
\def\sheaf#1{\protect\mathcal{#1}}
\def\vieuxsheaf#1{{#1}} % cf. problème avec sheaf{a\Phi_i} par exemple...
\DeclareMathOperator{\Hom}{Hom}
\DeclareMathOperator{\can}{can}
\DeclareMathOperator{\Ind}{Ind}
\DeclareMathOperator{\ind}{ind}
%\DeclareMathOperator{\Hhomtop}{\mathcal{H}\!\mathit{omtop}}
\DeclareMathOperator{\Hhomtop}{\mathbf{Homtop}}
\DeclareMathOperator{\Point}{Point}
\DeclareMathOperator{\Pro}{Pro}
\DeclareMathOperator{\Top}{Top}
\DeclareMathOperator{\Mod}{Mod}
\DeclareMathOperator{\Cart}{Cart}
%\DeclareMathOperator{\Mmorsite}{\mathcal{M}\!\mathit{orsite}}
\DeclareMathOperator{\Mmorsite}{\mathbf{Morsite}}
\DeclareMathOperator{\loc}{loc}
\DeclareMathOperator{\Loc}{Loc}
\DeclareMathOperator{\Pic}{Pic}
\DeclareMathOperator{\Et}{Et}
\DeclareMathOperator{\card}{card}
\DeclareMathOperator{\Ens}{Ens}
\DeclareMathOperator{\h}{\underline{H}}
\DeclareMathOperator{\Fonct}{Fonct}
\DeclareMathOperator{\Ker}{Ker}
\DeclareMathOperator{\but}{but}
\DeclareMathOperator{\Ob}{Ob}
\DeclareMathOperator{\ob}{ob}
\DeclareMathOperator{\Lib}{Lib}
\DeclareMathOperator{\esj}{esj}
\DeclareMathOperator{\id}{id}
\DeclareMathOperator{\Id}{Id}
\DeclareMathOperator{\ad}{ad}
\DeclareMathOperator{\Iim}{Im}
\DeclareMathOperator{\pr}{pr}
\DeclareMathOperator{\Cat}{Cat}
\DeclareMathOperator{\car}{car}
\DeclareMathOperator{\SysInd}{SysInd}
\DeclareMathOperator{\Sup}{Sup}
\DeclareMathOperator{\Coker}{Coker}
\DeclareMathOperator{\Filt}{Filt}
\DeclareMathOperator{\acc}{acc}
\DeclareMathOperator{\sur}{sur}
\DeclareMathOperator{\cat}{cat}
\DeclareMathOperator{\Ord}{Ord}
\DeclareMathOperator{\Gr}{Gr}
\DeclareMathOperator{\Cov}{Cov}
\DeclareMathOperator{\Caf}{Caf}
\DeclareMathOperator{\coker}{coker}
\DeclareMathOperator{\coim}{coim}
\DeclareMathOperator{\im}{im}
\DeclareMathOperator{\Ouv}{Ouv}
\DeclareMathOperator{\diag}{diag}
\DeclareMathOperator{\Esp}{Esp}
\DeclareMathOperator{\TOP}{TOP}
\DeclareMathOperator{\Unif}{Unif}
\DeclareMathOperator{\Iind}{\mathbf{Ind}}
\DeclareMathOperator{\Hhomcart}{\mathbf{Homcart}}
\DeclareMathOperator{\Hhomtopess}{\mathbf{Homtopess}}
\DeclareMathOperator{\Ppoints}{\mathbf{Points}}
\DeclareMathOperator{\Ppoint}{\mathbf{Point}}
\DeclareMathOperator{\Ttors}{\mathbf{Tors}}
\DeclareMathOperator{\Ssex}{\mathbf{Sex}}
\DeclareMathOperator{\Ffib}{\mathbf{Fib}}
\DeclareMathOperator{\Vvois}{\mathbf{Vois}}
\DeclareMathOperator{\Ppointess}{\mathbf{Pointess}}
\DeclareMathOperator{\scrHom}{\mathbf{Hom}}
\DeclareMathOperator{\sheafPic}{\mathbf{Pic}}
\DeclareMathOperator{\sheafAut}{\mathbf{Aut}}
\DeclareMathOperator{\Hhomtoplocan}{\mathbf{Homtoplocan}}
\DeclareMathOperator{\Hhomtopan}{\mathbf{Homtopan}}
\DeclareMathOperator{\Pptgeom}{\mathbf{Ptgéom}}
\DeclareMathOperator{\Morsite}{Morsite}
\DeclareMathOperator{\Fib}{Fib}
\DeclareMathOperator{\Fais}{Fais}
\DeclareMathOperator{\kar}{kar}
\DeclareMathOperator{\Groupes}{Groupes}
\DeclareMathOperator{\Restr}{Restr}
\DeclareMathOperator{\Isom}{Isom}
\DeclareMathOperator{\Prol}{Prol}
\DeclareMathOperator{\Res}{Res}
\DeclareMathOperator{\Tors}{Tors}
\DeclareMathOperator{\esp}{esp}
\DeclareMathOperator{\es}{es}
\DeclareMathOperator{\od}{od}
\DeclareMathOperator{\gr}{gr}
\DeclareMathOperator{\sob}{sob}
\DeclareMathOperator{\dfn}{dfn}
\DeclareMathOperator{\Sch}{Sch}
\DeclareMathOperator{\ouv}{ouv}
\DeclareMathOperator{\cons}{cons}
\DeclareMathOperator{\Nerf}{Nerf}
\DeclareMathOperator{\Points}{Points}
\DeclareMathOperator{\Espsob}{Espsob}
\DeclareMathOperator{\Homtop}{Homtop}
\DeclareMathOperator{\Pt}{Pt}
\DeclareMathOperator{\Quot}{Quot}
\DeclareMathOperator{\glob}{glob}
\DeclareMathOperator{\Homfib}{Homfib}
\DeclareMathOperator{\Qouv}{Qouv}
\DeclareMathOperator{\Quotouv}{Quotouv}
\DeclareMathOperator{\ens}{ens}
\DeclareMathOperator{\ab}{ab}
\DeclareMathOperator{\an}{an}
\DeclareMathOperator{\fs}{fs}
\DeclareMathOperator{\red}{red}
\DeclareMathOperator{\Ab}{Ab}
\DeclareMathOperator{\cosq}{cosq}
\DeclareMathOperator{\Ann}{Ann}
\DeclareMathOperator{\Zar}{Zar}


%%%%%%%%%%%%
%%%%% SGA4-2
\DeclareMathOperator{\Sex}{Sex}
\DeclareMathOperator{\vois}{vois}
\DeclareMathOperator{\Vois}{Vois}
\DeclareMathOperator{\Ext}{Ext}
\DeclareMathOperator{\ext}{ext}
\DeclareMathOperator{\Fer}{Fer}
\DeclareMathOperator{\Add}{Add}
\DeclareMathOperator{\Cocart}{Cocart}
\DeclareMathOperator{\cosk}{cosk}
\DeclareMathOperator{\plat}{plat}

\DeclareMathOperator{\faisc}{faisc}
\DeclareMathOperator{\sk}{sk}
\DeclareMathOperator{\ch}{ch}

\def\HH{\mathbf{H}}
\def\qc{\mathrm{qc}}
\def\et{\mathrm{et}}
\def\coh{\mathrm{coh}}

\def\Hvr#1{%
\begin{array}[b]{l}
{\kern2pt\scriptstyle r}\\[-8pt]
{\kern1pt\scriptstyle \vee}\\[-5pt]
H^{#1}
\end{array}}

\def\Hvv#1{%
\begin{array}[b]{l}
\kern2pt\rotatebox{90}{${\scriptscriptstyle\ll}$}\\[-6pt]
H^{#1}
\end{array}}

\def\Hvi#1{%
\begin{array}[b]{l}
{\kern0.2pt\scriptscriptstyle \infty}\\[-8pt]
{\kern1pt\scriptstyle \vee}\\[-5pt]
H^{#1}
\end{array}}

\def\Hv#1{%
\begin{array}[b]{l}
{\kern1pt\scriptstyle \vee}\\[-5pt]
H^{#1}
\end{array}}

\DeclareMathOperator{\cart}{cart}
\DeclareMathOperator{\Ffaisc}{\mathbf{Faisc}}


\def\Ev#1{%
\begin{array}[b]{l}
{\kern1pt\scriptstyle \vee}\\[-5pt]
E'_{#1}
\end{array}}

\DeclareMathOperator{\Cont}{Cont}
\DeclareMathOperator{\Ccont}{\mathscr{C}\!\mathit{ont}}
\DeclareMathOperator{\Site}{Site}
\newcommand\MM{\mathbf{M}}
\def\limtop#1{{\displaystyle{\varprojlim_{#1}\!\mbox{\rm top }}}}
\def\Limtop{\mathop{\varprojlim\!\mbox{\rm top}}}
\DeclareMathOperator{\spec}{spec}
\def\ss{\scriptstyle{\mathscr{S}}}
\DeclareMathOperator{\Gal}{Gal}
\def\smallm{\scriptstyle{\mathscr{M}}}
\DeclareMathOperator{\End}{End}
\DeclareMathOperator{\ass}{ass}
\DeclareMathOperator{\zar}{zar}
\DeclareMathOperator{\Ogg}{Ogg}
\DeclareMathOperator{\res}{res}
\DeclareMathOperator{\tr}{tr}
\DeclareMathOperator{\Supp}{Supp}
\DeclareMathOperator{\codim}{codim}
\DeclareMathOperator{\rk}{rk}
\DeclareMathOperator{\cur}{cur}
\DeclareMathOperator{\cdqc}{cdqc}
\newcommand\DD{\mathbf{D}}
\DeclareMathOperator{\cl}{cl}
\def\JJ{\mathbf{J}}
\def\FF{\mathbf{F}}
\def\BB{\mathbf{B}}
\DeclareMathOperator{\Alget}{Alget}
\DeclareMathOperator{\Ass}{Ass}
\DeclareMathOperator{\supp}{supp}
\DeclareMathOperator{\rad}{rad}
\DeclareMathOperator{\prof}{prof}
\DeclareMathOperator{\cint}{int}
\DeclareMathOperator{\sic}{sic}
\DeclareMathOperator{\Tor}{Tor}
\DeclareMathOperator{\tors}{tors}
\DeclareMathOperator{\Sym}{Sym}
\DeclareMathOperator{\parf}{parf}
\DeclareMathOperator{\tor}{tor}
\DeclareMathOperator{\defn}{defn}
\DeclareMathOperator{\Div}{Div}
\DeclareMathOperator{\Pol}{Pol}
\DeclareMathOperator{\Fex}{Fex}
\DeclareMathOperator{\Lie}{Lie}
\DeclareMathOperator{\tot}{tot}
\DeclareMathOperator{\Aut}{Aut}
\DeclareMathOperator{\HOM}{HOM}
\DeclareMathOperator{\pch}{pch}
\DeclareMathOperator{\Ch}{Ch}
\DeclareMathOperator{\EXT}{EXT}
\DeclareMathOperator{\PIC}{PIC}
\DeclareMathOperator{\Rf}{Rf}
\DeclareMathOperator{\inj}{inj}
\newcommand\VV{\mathbf{V}}

\def\Xv{%
\begin{array}[b]{l}
{\kern1pt\scriptstyle \vee}\\[-5pt]
X
\end{array}}

%%%%%%%
\newcommand\TT{\mathbf{T}}
\newcommand\XX{\mathbf{X}}
\newcommand\EE{\mathbf{E}}
%%%%%%

\def\qedsymbol{\textup{C.Q.F.D.}}
\let\CQFD\qed
\let\cqfd\qed

%%%%%\bigdot
\font\bigsymb=cmsy10 at 4pt
\def\bigdot{{\kern1.2pt\raise 1.5pt\hbox{\bigsymb\char15}}}
%%%%%

\def\Lright{\mathop{\vcenter{\hbox{\hbox{$\longrightarrow$}\lower8pt \hbox{\kern-16pt\hbox{$\longrightarrow$}}}}}}

\def\Aa{{\mathcal A}}
\def\Bb{{\mathcal B}}
\def\Cc{{\mathcal C}}
\def\Dd{{\mathcal D}}
\def\Ee{{\mathcal E}}
\def\Ff{{\mathcal F}}
\def\Gg{{\mathcal G}}
\def\Hh{{\mathcal H}}
\def\Ii{{\mathcal I}}
\def\Jj{{\mathcal J}}
\def\Kk{{\mathcal K}}
\def\Ll{{\mathcal L}}
\def\Mm{{\mathcal M}}
\def\Nn{{\mathcal N}}
\def\Oo{{\mathcal O}}
\def\Pp{{\mathcal P}}
\def\Qq{{\mathcal Q}}
\def\Rr{{\mathcal R}}
\def\Ss{{\mathcal S}}
\def\Tt{{\mathcal T}}
\def\Uu{{\mathcal U}}
\def\Vv{{\mathcal V}}
\def\Ww{{\mathcal W}}
\def\Xx{{\mathcal X}}
\def\Yy{{\mathcal Y}}
\def\Zz{{\mathcal Z}}

\newcommand{\CC}{\mathbf{C}}
\newcommand{\GG}{\mathbf{G}}
\newcommand{\KK}{\mathbf{K}}
\newcommand{\LL}{\mathbf{L}}
\newcommand{\NN}{\mathbf{N}}
\newcommand{\PP}{\mathbf{P}}
\newcommand\QQ{\mathbf{Q}}
\newcommand\RR{\mathbf{R}}
\newcommand\ZZ{\mathbf{Z}}

\def\to{\mathchoice{\longrightarrow}{\rightarrow}{\rightarrow}{\rightarrow}}
\def\from{\mathchoice{\longleftarrow}{\leftarrow}{\leftarrow}{\leftarrow}}
\def\mto{\mathchoice{\longmapsto}{\mapsto}{\mapsto}{\mapsto}}
\def\mfrom{\mathrel{\from\joinrel\mapstochar\,}}
\def\hto{\mathrel{\lhook\joinrel\to}}
\def\hfrom{\mathrel{\from\joinrel\rhook}}

\def\implies{\mathchoice{\Longrightarrow}{\Rightarrow}{\Rightarrow}{\Rightarrow}}
\def\implied{\mathchoice{\Longleftarrow}{\Leftarrow}{\Leftarrow}{\Leftarrow}}
\def\iff{\mathchoice{\Longleftrightarrow}{\Leftrightarrow}{\Leftrightarrow}{\Leftrightarrow}}

\def\To#1{\mathchoice{\xrightarrow{\labelstyle~~#1~}}{\stackrel{#1}{\longrightarrow}}{}{}}
\def\Mto#1{\mathrel{\mapstochar\To{#1}}}
\def\From#1{\mathchoice{\xleftarrow{\labelstyle~~#1~}}{\stackrel{#1}{\longleftarrow}}{}{}}

\newcommand{\isomto}{\overset{\sim}{\longrightarrow}}
\newcommand{\isomfrom}{\overset{\sim}{\longleftarrow}}

\newcommand{\IFF}{\unskip~$\iff$~\ignorespaces}
\newcommand{\THEN}{\unskip~$\implies$~\ignorespaces}

\let\leq\leqslant
\let\geq\geqslant
\let\le\leqslant
\let\ge\geqslant

\def\cf{{cf.}\kern.3em}
\def\Cf{{Cf.}\kern.3em}
\def\ie{{i.e.}\kern.3em}
\def\iev{{i.e.},\xspace}
%\def\loccito{loc.\kern2pt cit.\xspace}
%\def\loccitn{\textit{loc.\kern2pt cit.}\xspace}
%\def\loccit{\sisi{\loccito}{\loccitn}}
\def\EGA{\textup{EGA}\kern.2em}
\def\SGA{\textup{SGA}\kern.3em}

\let\oldS\S
\def\S{\oldS\kern .15em}
\def\ptbl{.\kern .2em}
\def\resp{resp.\kern.3em}
\def\Exp{Exp.\kern.3em}

\let\bar\overline
\let\hat\widehat
\let\tilde\widetilde
\let\dpl\displaystyle
\let\ldots\dots
\let\emptyset\varnothing
\let\skpt\skippointrait
\let\epsilon\varepsilon

\def\sfrac#1#2{{#1}/{#2}}

\DeclareMathOperator{\cd}{cd}
%\DeclareMathOperator{\Hhom}{\mathcal{H}\!\mathit{om}}
\DeclareMathOperator{\Hhom}{\mathbf{Hom}}
%\DeclareMathOperator{\Eext}{\mathcal{E}\!\mathit{xt}}
\DeclareMathOperator{\Eext}{\mathbf{Ext}}
\DeclareMathOperator{\Spec}{Spec}
\DeclareMathOperator{\Tr}{Tr}

\let\emptyset=\varnothing

%%%%%%%%%%%%%%%
%%% expose XVII
%%%%%%%%%%%%%%%

\makeatletter
\let\old@optexpose\optexpose
\newcommand{\empty@optexpose}[1]{}
\newcommand{\paren@optexpose}[1]{\old@optexpose{#1}(}
\newcommand{\refnoexpose}[1]{\let\optexpose\empty@optexpose\ref{#1}\let\optexpose\old@optexpose}%\refnoexpose{X.cor5.2} engendre "5.2" au lieu de "X 5.2"
\newcommand{\neweqref}[1]{\let\optexpose\paren@optexpose\textup{\ref{#1})}\let\optexpose\old@optexpose}%\neweqref{XII.eq4.1} engendre "XII (4.1)" au lieu de "(XII 4.1)"
\makeatother
\newcommand{\refvb}[2][VI${}_{\text{B}}$ ]{\sisi{#1\refnoexpose{#2}}{\ref{#2}}}

\newcommand{\leftexp}[2]{{\vphantom{#2}}^{#1}{#2}}
\newcommand{\AAA}{\mathbf{A}}
\newcommand{\Sf}{\mathfrak{S}}
\DeclareMathOperator*{\quotprojlim}{\text{``$\varprojlim$''}}
\DeclareMathOperator*{\quotinjlim}{\text{``$\varinjlim$''}}
\DeclareMathOperator{\Fl}{F\sisi{\ell}{l}}
\DeclareMathOperator{\Pl}{P\sisi{\ell}{l}}
%\DeclareMathOperator{\cl}{c\sisi{\ell}{l}}
\DeclareMathOperator{\tdf}{tdf}
\DeclareMathOperator{\rg}{rg}
\newcommand{\iInd}{\sisi{I}{i}nd}
\newcommand{\iind}{ind}
\newcommand{\ppro}{pro}
\newcommand{\tTor}{\sisi{T}{t}or}
\newcommand{\ttor}{tor}
\newcommand{\chap}{Chap.}
\newcommand{\Ffex}{\Ff\mathit{ex}}
\newcommand{\sheafk}{\underline{k}}
\newcommand{\sheafn}{\underline{n}}
\newcommand{\sheafGamma}{\underline{\Gamma}}

% Repeated references to an already typeset footnote must not use
% \footnotemark without a matching \footnotetext: Hyperref would then create
% a live annotation to a destination that cannot exist.  This helper keeps the
% current printed footnote mark and links it to an explicit stable note target.
\makeatletter
\newcommand{\sgarepeatedfootnotemark}[1]{%
  \begingroup
  \edef\@thefnmark{\thefootnote}%
  \hyperlink{#1}{\@makefnmark}%
  \endgroup}
\makeatother

\sisi{
\DeclareMathOperator{\simarrow}{\rule{.15mm}{1.8mm}\kern -3.4pt\rightsquigarrow}
\newcommand{\sheafell}{\underline{\ell}}
\newcommand{\sheafZZ}{\underline{\ZZ}}
}{
\newcommand{\nderef}[1]{la N.D.E.~(\ref{#1}) page~\pageref{#1}}
}

%%%%%%%%%%%%%%%%%
%%%%% expose XVII
%%%%%%%%%%%%%%%%%

% END INCLUDED SGA4_English_current_buildable_source_20260906/source/fichiers-configuration-latex/macros.tex


\externaldocument{../01/01_English.tex}
\externaldocument{../02/02_English.tex}
\externaldocument{../03/03_English.tex}
\externaldocument{../04/04_English.tex}
\externaldocument{../05/05_English.tex}
\externaldocument{../05bis/05bis_English.tex}
\externaldocument{../06/06_English.tex}
\externaldocument{../07/07_English.tex}
\externaldocument{../08/08_English.tex}
\externaldocument{../09/09_English.tex}
\externaldocument{../10/10_English.tex}

\setboolean{original}{false}

\begin{document}

\begin{center}
XI. Comparison with classical cohomology: the case of a smooth scheme\\
M. Artin
\end{center}

\selectlanguage{english}version: \showgitstatus
\tableofcontents

\else
\chapter{Comparison with classical cohomology: the case of a smooth scheme}
%\addtocontents{toc}{M. Artin}
\begin{center}
M. Artin
\end{center}
\fi
\phantomsection\MakeLinkTarget*{sga4.exp.XI}\label{sga4.exp.XI}

% BEGIN INCLUDED SGA4_English_current_buildable_source_20260906/source/11/drafts/SGA4_XI_LINES_43_199_English.texfrag
\setcounter{pageoriginal}{63}

\section{Introduction}\etiquette[1]{XI.sec1}\pageoriginale

In Section~\hyperref[XI.sec4]{4} we shall prove that étale cohomology
with values in a locally constant torsion sheaf agrees with classical
cohomology for a scheme $X$ smooth over $\Spec\CC$. The proof, which is
essentially the one given in \cite{sga4.exp.XI.bib.1}, is elementary once one has the
theorem of Grauert recalled below (\hyperref[XI.thm4.3]{4.3} (iii)),
which says that every finite covering of $X(\CC)$ comes from an étale
covering of $X$; this is why we give it here, although we shall prove a
more complete result later, using Hironaka's resolution of
singularities theorem \cite{sga4.exp.XI.bib.3}\footnote[1]{Note, moreover, that
  Hironaka's resolution of singularities also gives a very simple
  proof of Grauert's result, and in this capacity may in any case be
  regarded as underlying the comparison theorems. See in this
  connection the exposé by Mme M. Raynaud in \SGA 1 XII (North Holland
  Pub. Cie).}. Furthermore, the good neighborhoods constructed in
\hyperref[XI.sec3]{Section~3} could be useful in other\pageoriginale situations.

Note that étale cohomology and classical cohomology do not agree for
sheaves that are not torsion. For example, one has
$H^1(X_{\et},\ZZ)=0$ for every normal scheme $X$
(\hyperrefext[IX][IX.prop3.6]{3.6} (ii)), but classical cohomology does
not necessarily vanish.

\section{Existence of sufficiently general hyperplane sections}\etiquette[2]{XI.sec2}

The following result is classical, but a reference is lacking (pending
\EGA V):

\begin{enonce*}{Theorem 2.1}\etiquette[2.1]{XI.thm2.1} %Why is the statement
                                %not formulated for a linear subspace
                                %varying around a fixed point?
                                %That is, (i) with $P\in L$ or
                                %$(ii)$ with $d=1$
Let $k$ be a field and let $V\subset \PP^n_k$ be a projective
algebraic scheme purely of dimension $r$. Let $V'$ be the open
subscheme of points of $V$ smooth over $\Spec k$, and let $P$ be a
rational point of $\PP^n$. Then:
\begin{enumerate}
\renewcommand{\theenumi}{\roman{enumi}}
\renewcommand{\labelenumi}{(\theenumi)}
\item A ``sufficiently general'' linear subspace $L$ of a given
  dimension $d$ in $\PP^n$ meets $V$ transversely at every point of
  $L\cap V'$.

\item Let $\{H_1,\ldots,H_s\}$ $(s\leq r)$ be a system of %a
                                %'hyperplane' is not linear here
  hypersurfaces of degree $d\geq 2$ passing through $P$, and let
  $Y=H_1\cap\ldots \cap H_s$. If $\{H_1,\ldots,H_s\}$ is sufficiently
  general, then $Y$ meets $V$ transversely at every point of $Y\cap
  V'$.
\end{enumerate}
\end{enonce*}

The phrase ``for a sufficiently general linear subspace'' means that
there exists a dense open subset $U$ of the Grassmannian parametrizing
the linear subspaces of dimension % 'which parametrizes'
$d$, such that the assertion under consideration about $L$ is true
whenever the parameter of $L$ lies in $U$. Thus this does not imply
that such an $L$ (defined over $k$) exists if $k$ is not infinite.
Likewise,
% mention the Veronese embedding?
the hypersurfaces of degree $d$ passing through $P$ are parametrized
\pageoriginale by a projective space $\PP^N$, and the phrase ``for
$\{H_1,\ldots,H_s\}$ sufficiently general'' means that there exists a
dense open subset $U$ of $(\PP^N)^s$ such that the assertion under
consideration is true whenever the parameter of
$(H_1,\ldots,H_s)$ lies in $U$.

Proof of (i): By induction one is reduced immediately to the case
where $L$ is a hyperplane. Let $(V'_{\nu})$ be a finite open covering
of $V'$. If assertion (i) is true when the symbol $V'$ is replaced by
$V'_{\nu}$ for every $\nu$, then it is also true for $V'$, because the
intersection of finitely many dense open subsets is again open and
dense. We may therefore replace the situation by an affine algebraic
scheme, say $V\subset \EE^n_k=\Spec k[x_1,\ldots,x_n]$. The hyperplane
$L$ will be the zero set of a linear equation
$$
\ell(x)=a_0+\sum^n_{i=1}a_ix_i{\quoi}.
$$

Recall that a smooth scheme is locally a complete intersection. We may
therefore replace $V'$ by an open subset of points of $V$ such that
there exist polynomials $f_1(x),\ldots,f_{n-r}(x)$ vanishing on $V$ and
indices $1\leq j_{\nu}\leq n$ $(\nu=1,\ldots,n-r)$ such that the
determinant
$$
\left|\frac{\partial f_i}{\partial x_{j_{\nu}}}\right|\quad 
(i,\nu=1,\ldots,n-r)
$$
is invertible. Say $j_{\nu}=\nu$ $(\nu=1,\ldots,n-r)$. We may then
solve the equations
$$
y_j=\sum^{n-r}_{i=1}\quad c_i\quad \frac{\partial f_i}{\partial x_j}
$$
for\pageoriginale $c_i$ as rational functions of the variables $x_i$,
$y_j$ $(i=1,\ldots, n-r; j=1,\ldots,n)$, and these functions $c_i(x,y)$
are defined at every point of $V'$.

Now to say that $L$ meets $V$ transversely at a point $Q$ of $V'$ is
equivalent to saying that the value of the matrix
$$
\begin{bmatrix}
\frac{\partial f_1}{\partial x_1} & \cdots & \frac{\partial
  f_1}{\partial x_n}\\
\vdots & &\vdots\\
\frac{\partial f_{n-r}}{\partial x_1} & & \frac{\partial
  f_{n-r}}{\partial x_n}\\
a_1 & \cdots & a_n
\end{bmatrix}
$$
at $Q$ has rank $n-r+1$. Thus, for $L$ to meet $V$ transversely at every
point of $V'$, it is necessary and sufficient that the system of
$r+1$ equations
\begin{align*}
a_j &= \sum^{n-r}_{i=1} c_i (x,a)\frac{\partial f_i}{\partial x_j}\qquad
(j=n-r+1,\ldots, n)\\
a_0 &+ \sum^n_{i=1} a_i x_i =0
\end{align*}
have no solution on $V'$. Such a condition is plainly constructible in
$(a)$: Let $X=V'\times A$, where $A$ is the affine space with
coordinates $a_0,\ldots,a_n$, and let $Y\subset X$ be the closed subset
of solutions of these equations. The condition is that the fiber of
$Y/A$ be empty, which is indeed a constructible\pageoriginale
condition (\EGA IV 9.5.1). It is therefore enough % \ref
to verify that the fiber above the generic point of $A$ is empty,
which follows immediately from the form of the equations. This
completes the proof of (i). (ii): Consider the rational map
$\PP^n\rightarrow \PP^N$ given by the linear system of hypersurfaces
of degree $d$ ($d\geq 2$) passing through $P$. It is well known (and
easily verified) that this is a ``blowing up'' of $P$, and therefore
gives a locally closed immersion
$$
f:\PP^n-\{P\}\longrightarrow \PP^N{\quoi}.
$$
Let $V''=V'-(\{P\}\cap V')$, let
$\bar{V}{}''=\sisi{f}{F}(V'')$, and let $L_i$ be the hyperplane of
$\PP^N$ corresponding to the hypersurface $H_i$
$(i=1,\ldots,s)$. Then $V''$ is isomorphic to $\bar{V}{}''$, and by
the definition of $f$, to say that $Y$ meets $V''$ transversely is
equivalent to saying that $L_1\cap\ldots \cap L_s$ meets
$\bar{V}{}''$ transversely. Putting $\bar{V}=$ the closure of
$\bar{V}{}''$, we are reduced to assertion (i) for $\bar{V}$, which
implies (ii) with $V''$ in place of $V'$. It remains to consider the
point $P$, if it is a smooth point of $V$. But it is evident that if
$\{H_1,\ldots,H_s\}$ is sufficiently general and $P$ is smooth on
$V$, then $Y$ meets $V$ transversely at $P$, whence the result.

% END INCLUDED SGA4_English_current_buildable_source_20260906/source/11/drafts/SGA4_XI_LINES_43_199_English.texfrag


% BEGIN INCLUDED SGA4_English_current_buildable_source_20260906/source/11/drafts/SGA4_XI_LINES_200_353_English.texfrag
\section{Construction of good neighborhoods}\etiquette[3]{XI.sec3}

\begin{enonce*}{Definition 3.1}\etiquette[3.1]{XI.def3.1}
An elementary fibration is a morphism of schemes
${f:X\rightarrow S}$ which can be embedded in a commutative diagram
$$
\xymatrix@=2cm{
X  \ar[r]^{j} \ar[rd]_{f}& \bar{X} \ar[d]^{\bar{f}}& Y \ar[l]_{i}  \ar[ld]^{g}\\
               &    S   &      
}
$$
satisfying\pageoriginale the following conditions:
\begin{enumerate}
\renewcommand{\theenumi}{\roman{enumi}}
\renewcommand{\labelenumi}{(\theenumi)} 
% \item $i$ is a closed immersion.
\item $j$ is an open immersion, dense in every fiber, and
  $X=\bar{X}-Y$. 

\item $\bar{f}$ is smooth and projective, with geometrically
  irreducible fibers of dimension $1$.

\item $g$ is an étale covering, and every fiber of $g$ is nonempty.
\end{enumerate}
\end{enonce*}


One easily checks that such an embedding of $X$, if it exists, is
unique up to canonical isomorphism.

\begin{enonce*}{Definition 3.2}\etiquette[3.2]{XI.def3.2}
A good neighborhood relative to $S$ is an $S$-scheme $X$ for which
there exist $S$-schemes
$$
X=X_n,\ldots,X_0=S
$$
and elementary fibrations $f_i:X_i\rightarrow X_{i-1}$,
% $f_i$ is an $S$-morphism
$i=1,\ldots, n$.
\end{enonce*}

\begin{enonce*}{Proposition 3.3}\etiquette[3.3]{XI.prop3.3}
\footnote{This result extends without difficulty to a semilocal
    assertion.}
Let $k$ be an algebraically closed field, let $X/\Spec k$ be a smooth
scheme, and let $x\in X$ be a rational point. There exists an open
subset of $X$ containing $x$ which is a good neighborhood (relative to
$\Spec k$).
\end{enonce*}

Proof. Take $X$ irreducible. By induction on $\dim X=n$, it suffices to
find a neighborhood $U$ of $x$ in $X$ and an elementary fibration
$f:U\rightarrow V$, with $V$ smooth and of dimension $n-1$. Indeed,
there will exist a neighborhood $V'$ of $v=f(x)$ which is a good
neighborhood, and we may take
$U'=U\cap \sisi{f^{-1}(V')}{V'}$ as a good neighborhood of $x$.

We may suppose that $X\subset \EE^r$ is affine. Let $X_0$ be the
closure of $X$ in $\PP^r$. Let $\bar{X}$ be the normalization of $X_0$
and let $Y=\bar{X}-X$, endowed with the induced reduced structure. In
addition, let\pageoriginale $S\subset \bar{X}$ be the closed subset of
singular points. One has $S\subset Y$ and
\begin{align*}
\dim \bar{X} &= \dim X=n{\quoi},\\
\dim Y &= n-1{\quoi},\\
\dim S &\leq n-2{\quoi}.
\end{align*}

Embed $\bar{X}$ in a projective space $\PP^N$ by means of an invertible
sheaf $L$ of the form $M^{\otimes r}$, with $M$ very ample and
$r\geq 2$. There then exist hyperplanes % linear
$H_1,\ldots, H_{n-1}$ of $\PP^N$, where $H_i$ is the zero set of
$$
\sum^N_{\nu =0}a_{i\nu} x_{\nu} =0{\quoi},
$$
which contain $x$ and such that the intersection
$L=H_1\cap\ldots \cap H_{n-1}$ has dimension $N-n+1$ and meets
$\bar{X}$ and $Y$ transversely (\hyperref[XI.thm2.1]{2.1}).
% see the comment on Theorem 2.1
 %For the intersection with $\bar{X}$: $S\cap L=∅$ for sufficiently
 %general $L$, hence $\bar{X}\cap L=(\bar{X} - Y)\cap L$ and
 %$\bar{X}- Y$ is smooth.
 %For the intersection with $Y$: $Y$ is reduced and $k$ algebraically closed,
 %hence every irreducible component is generically smooth, and sufficiently
 %general $L$ does not meet the singular locus of $Y$.
The intersection $\bar{X}\cap L$ is a smooth connected curve (this
follows from Bertini's theorem), and $Y\cap L$ has dimension $0$.
Choose in addition another hyperplane
$$
H_0:\sum^N_{\nu=0}a_{0 \nu}x_{\nu}=0
$$
such that $H_0$ meets $\sisi{\bar{X}}{X}\cap L$ transversely and
$H_0\cap Y\cap L=∅$. % One must also choose $x\notin H_0$ so that
                    % the map $f'$ below is well defined at $x$

Consider the projection $\PP^N\rightarrow \PP^{n-1}$ obtained by
``introducing the projective coordinates''
$$
y_i=\sum^N_{\nu =0}a_{i_{\nu}}x_{\nu}{\quoi}.
$$
It is a rational map defined outside the center of projection
$C=H_0\cap\ldots \cap H_{n-1}$. Let $\epsilon :P'\rightarrow \PP^N$ be
the blowup of $C$, so that there is a diagram
$$
\xymatrix@=1.5cm{
\PP^{N}  \ar[rd]& &P'\ar[ld]^{\pi}\ar[ll]_{\epsilon}\\
               &\PP^{n-1} &
}
$$
where\pageoriginale $\pi$ is a morphism. Let $\bar{X}{}'\subset P'$ be
the ``proper'' inverse image of $\bar{X}$, that is, the closure of
$\epsilon^{-1}(\bar{X}-(\bar{X}\cap C))$. Since, by hypothesis, $C$
meets $\bar{X}$ transversely, the morphism
$\bar{X}{}'\rightarrow \bar{X}$ is a blowup of the finite set
$\bar{X}\cap C$. Let $X'=X-\bar{X}\cap C$, which is also identified
with an open subscheme of $\bar{X}{}'$, and let
$Y'=\bar{X}{}'-X'$ be the closed subscheme with the induced reduced
structure. There is a diagram $\DD$ of morphisms
$$
\xymatrix@=2cm{
  X'  \ar[r]^{j} \ar[rd]_{f'}& \bar{X}{}'  \ar[d]^{\bar{f}{}'} & Y'\ar[l]_{i} \ar[ld]^{g'}     &\\
                           & \PP^{n-1} &\qquad ,   
}
$$
and I claim that there exists a neighborhood $V$ of $v=f'(x)$ such
that the restriction of $\DD$ to $V$ satisfies conditions (i), (ii),
(iii) of (\hyperref[XI.def3.1]{3.1}), hence that $f'_{\mid V}$ is an
elementary fibration. This will complete the proof.

Condition (i) is trivial. For (ii), note that $\bar{X}\cap L$ is a
smooth curve by hypothesis, and one sees immediately that there is a
one-to-one morphism ${\bar{f}{}'}^{-1}(v)\rightarrow \bar{X}\cap L$
% one-to-one morphism?
induced by $\epsilon$. Hence
$\left({\bar{f}{}'}^{-1}(v)\right)_{\red}\rightarrow
\bar{X}\cap L$ is bijective. To check that $\bar{f}{}'$ is smooth
% remove the \left and \right?
over a neighborhood of $v$, it suffices, by Hironaka's lemma
(\SGA 1 II 2.6 %\ref
\footnote{and \EGA IV 5.12.10.}), to check that it is smooth at the
generic point of ${\bar{f}{}'}^{-1}(v)$. At this point, $\bar{X}{}'$ is
isomorphic to $\bar{X}$, and the morphism is smooth because $L$ meets
$\bar{X}$ transversely.

It remains to prove that $g'$ is étale in a neighborhood of $v$ (since
$Y$ has dimension $n-1$, it is evident that every fiber of $g'$ is
nonempty). One has $Y'=\epsilon^{-1}(Y)\amalg D_1
\amalg\ldots \amalg D_r$, where $\bar{X}\cap C=\{P_1,\ldots,P_r\}$ and
$D_i$ is the blowup\pageoriginale of $P_i$ in $\bar{X}$. Every $D_i$
is identified with $\epsilon^{-1}(P_i)$, and maps isomorphically onto
$\PP^{n-1}$\sisi{. Moreover, the projection}{} is defined at every
point of $Y$%since $Y\cap C=∅$
, and the induced morphism on $Y$ is none other than
$g'_{\mid \epsilon^{-1}(Y)}$. It is étale over $v$, because $L$ meets
$Y$ transversely, hence $g'$ is étale over~$v$.

% END INCLUDED SGA4_English_current_buildable_source_20260906/source/11/drafts/SGA4_XI_LINES_200_353_English.texfrag


% BEGIN INCLUDED SGA4_English_current_buildable_source_20260906/source/11/drafts/SGA4_XI_LINES_354_643_English.texfrag
\section{The comparison theorem}\etiquette[4]{XI.sec4}

\setcounter{subsection}{-1}
\subsection{}\etiquette[4.0]{XI.subsec4.0}
Let $X$ be a scheme locally of finite type over $\Spec \CC$, and
consider the usual topology on the space $X(\CC)$ of rational points
of $X$. We shall denote by $X_{\cl}$ the site of \emph{local
isomorphisms} $f:U\rightarrow X(\CC)$, that is, of morphisms $f$ of
topological spaces having the following property:

\sisi{for}{For} every $x\in U$ there exists an open neighborhood $U_x$
such that $f_{\mid U_x}$ is a homeomorphism of $U_x$ onto an open
neighborhood of $f(x)$.

A family $\{U_{\nu}\rightarrow U\}$ of morphisms of $X_{\cl}$ is said
to be \emph{covering} if $U$ is the union of the images of the
$U_{\nu}$.

Since an open immersion is a local isomorphism, we have a morphism
(inclusion) of sites (IV\quad ) % a reference is missing
\footnote{More precisely, the inclusion of categories
  $\Ouv(X(\CC))\rightarrow X_{\cl}$ defines a morphism of
  sites \hyperref[XI.eq4.1]{(4.1)} in the opposite direction.}
\begin{equation*}
\delta:X_{\cl}\longrightarrow X(\CC)\qquad
(\ast){\quoi}.\tag{4.1}
\end{equation*}\etiquette[4.1]{XI.eq4.1}
Now for every local isomorphism $U\rightarrow X(\CC)$ there exists, by
definition, a ``covering'' of $U$ by open subsets of $X(\CC)$, and one
deduces from (\quad) that the topoi % a reference is missing
associated with the two sites are equivalent (by $\delta_{\ast}$). In
particular, one may replace $X(\CC)$ by $X_{\cl}$ when computing
ordinary cohomology.

Now let $f:X'\rightarrow X$ be an étale morphism of schemes. Then
$f(\CC):X'(\CC)\rightarrow X(\CC)$ is a local isomorphism. This follows
immediately from the\pageoriginale Jacobian criterion
(\SGA 1 II 4.10) and the implicit function theorem. The functor
$X'\mapsto X'(\CC)$ therefore gives a morphism of sites
\begin{equation*}
\epsilon:X_{\cl}\longrightarrow X_{\et}{\quoi},\tag{4.2}
\end{equation*}\etiquette[4.2]{XI.eq4.2}
where $X_{\et}$ is the étale site.

Recall the following results\footnote{Which appear in the Exposé cited
  in the footnote on p.~2.}: % identify the Exposés
                                % above to identify the Exposé
                                % mentioned in the note?

\begin{enonce*}{Theorem 4.3}\etiquette[4.3]{XI.thm4.3}
\begin{enumerate}
\renewcommand{\theenumi}{\roman{enumi}}
\renewcommand{\labelenumi}{(\theenumi)}
\item $X$ is connected and nonempty if and only if $X(\CC)$ is
  connected and nonempty.

\item The functor $\epsilon$ is fully faithful.

\item (``Grauert--Remmert theorem''). The functor $\epsilon$ induces an
  equivalence between the category of finite coverings of $X(\CC)$,
  endowed with its usual topology as a locally compact space, and the
  category of étale coverings of $X$.
\end{enumerate}
\end{enonce*}

We shall accept (i) as known. It follows, for example, from (GAGA)
%\ref
\cite{sga4.exp.XI.bib.4}. Assertion (ii) follows easily from (i). Indeed, to verify the
surjectivity assertion contained in (ii), let
$\varphi:X'(\CC)\rightarrow X''(\CC)$ be an $X(\CC)$-morphism. We may
suppose that all the schemes are connected and affine, and hence
separated. Then the graph $\Gamma$ of $\varphi$ is a connected
component of $\fprod{X'}{X''}{X}(\CC)$, whence
$\Gamma=Y(\CC)$ for some connected component $Y$ of
$\fprod{X'}{X''}{X}$, by (i). Then $Y$ is the graph of a morphism
$f:X'\rightarrow X''$ inducing $\varphi$, \ie the projection
$Y\rightarrow X'$ is an isomorphism: indeed,
$Y(\CC)\rightarrow X(\CC)$ is an isomorphism, from which it follows
easily that the same is true of $Y\rightarrow X'$ (which is étale,
radicial, and surjective!). For (iii), in view of (ii), everything
amounts to proving that every finite covering of $X(\CC)$ is of the
form $X'(\CC)$, where $X'\rightarrow X$\pageoriginale is étale.

For this, descent (\SGA 1 IX 4.7) immediately reduces us to
%\ref
the case in which $X$ is connected and normal. The problem is local on
$X$ (for the Zariski topology), and we may therefore suppose $X$
affine. Let $f:X\rightarrow \EE^n$ be a finite morphism, and let
$T'\rightarrow X(\CC)$ be a connected finite covering. One checks
immediately that the composite morphism
$\eta':T'\rightarrow \EE^n(\CC)$ is an ``analytic covering'' in the
sense of (\cite{sga4.exp.XI.bib.2} \S \hyperref[XI.sec2]{2}, Def.~\hyperref[XI.sec3]{3}).
By (\cite{sga4.exp.XI.bib.2} \S \hyperref[XI.sec2]{2}, Satz 8),
\sisi{$\eta'$}{} extends to an ``analytic covering''
$\eta:T\rightarrow \PP^n(\CC)$, and it is enough to prove that $\eta$
is \sisi{induced}{induced} by a suitable morphism of schemes
$f:Y\rightarrow \PP^n$, with $Y$ normal. Now the fundamental theorem
(\cite{sga4.exp.XI.bib.2} \S 13 Satz 42) asserts that there is a canonical normal
analytic structure on $T$ such that $\eta$ induces a morphism of
analytic spaces. Since $\eta$ is then a finite morphism, the direct
image of the structure sheaf on $T$ is a sheaf $\mathfrak{a}$ of
coherent analytic algebras on $\PP^n$. It follows from
\sisi{(GAGA)}{GAGA} \cite{sga4.exp.XI.bib.4} % \ref
that $\mathfrak{a}$ is the analytic sheaf associated with an algebraic
sheaf $A$ on $\PP^n$, and we take $Y=\sheaf{\Spec}A$.

We can now prove the following result:

\begin{enonce*}{Theorem 4.4}\etiquette[4.4]{XI.thm4.4}
Let $X$ be a smooth scheme over $\Spec \CC$.
\begin{enumerate}
\renewcommand{\theenumi}{\roman{enumi}}
\renewcommand{\labelenumi}{(\theenumi)}
\item There is an equivalence between the category of constructible
  locally constant torsion sheaves on $X_{\et}$ and the category of
  locally constant torsion sheaves with finite fibers on $X_{\cl}$,
  the equivalence being given by the quasi-inverse functors
  ${}_{\ast}$ and ${}^{\ast}$.

\item Let $F$ be a locally constant torsion sheaf with finite fibers on
  $X_{\cl}$, and also denote by $F$ the induced sheaf
  $\epsilon_{\ast}F$ on $X_{\et}$. Then
$$
R^q\epsilon_{\ast}F=0\quad \mbox{ if } q>0{\quoi}.
$$

\item With\pageoriginale the notation of {\rm(i)}, the canonical
  morphism
$$
H^q(X_{\et},F)\longrightarrow H^q(X_{\cl},F)
$$
is bijective for every $q$. In particular,
$$
H^q(X_{\et}, \ZZ/\sisi{n\ZZ}{n})\xrightarrow{\sim} H^q(X_{\cl},\ZZ/\sisi{n\ZZ}{n})
$$
for every $q$.
\end{enumerate}
\end{enonce*}

Proof. For (i), note that the sheaves in question are those that are
representable by étale coverings of $X$ (\resp by finite coverings of
$X(\CC)$). Since there is an equivalence between the categories of
these coverings (\hyperref[XI.thm4.3]{4.3}), assertion (i) follows
immediately. The isomorphism of (iii) follows from (ii) and from the
Leray spectral sequence V \hyperref[V.prop5.2]{5.2} % \ref + 5.3 instead of 5.2?
for $\epsilon$. It remains to prove (ii).

Now $R^q\epsilon_{\ast}F$ is the sheaf associated with the presheaf
$\Rr^qF,$, where $(\Rr^qF)(X')=H^q(X'_{\cl},F)$ for
$X'\rightarrow X$ étale. To prove that $R^q\epsilon_{\ast}F=0$ if
$q>0$, we must therefore prove the following \sisi{lemma}{}:

\begin{enonce*}{Lemma 4.5}\etiquette[4.5]{XI.lem4.5}
Let $\xi\in H^q(X_{\cl},F)$ and $x\in X_{\cl}=X(\CC)$. There exists an
étale morphism $X'\rightarrow X$ whose image contains $x$ and such that
the image of $\xi$ in $H^q(X'_{\cl},F)$ is zero.
\end{enonce*}

Proof of the lemma. Induction on $n=\dim X$; the problem is local on
$X$ for the étale topology, \sisi{}{and} we may therefore suppose $F$
constant, and by dévissage we may even suppose that
$F=\ZZ/\sisi{n\ZZ}{n}$. Moreover, replacing $X$ by a Zariski open
neighborhood of $x$, we may suppose (\hyperref[XI.prop3.3]{3.3}) that
$X$ admits an elementary fibration (\hyperref[XI.def3.1]{3.1})
$f:X\rightarrow S$. Retain the notation of
(\hyperref[XI.def3.1]{3.1}): by\pageoriginale direct calculation, one
sees that $R^q{j}_{\ast}F=0$ if $q>1$, that
$R^1j_{\cl\!\ast}F$ is the extension by $0$ of a constant sheaf on
$Y_{\cl}$, and finally that $j_{\cl\!\ast}F$ is the constant sheaf
$\ZZ/\sisi{n\ZZ}{n}$ on $\bar{X}$. Moreover, one may compute
$R^q \bar{f}_{\cl\!\ast}(j_{\cl\!\ast}F)$ fiber by fiber. From the
Leray spectral sequence
$$
E^{pq}_2=R^p\bar{f}_{\cl\!\ast}(R^qj_{\cl\!\ast}F)\Longrightarrow
R^{p+q}f_{\cl\!\ast}F
$$
one deduces that $R^{p+q}f_{\cl\!\ast}F$ may likewise be computed fiber
by fiber, and hence that

$R^0f_{\cl\!\ast}\ZZ/\sisi{n\ZZ}{n}=\ZZ/\sisi{n\ZZ}{n}$,

$R^1f_{\cl\!\ast}\ZZ/\sisi{n\ZZ}{n}$ is a locally constant torsion
  sheaf on $S_{\cl}$,

$R^qf_{\cl\!\ast}\ZZ/\sisi{n\ZZ}{n}=0$ if $q>1$.

The Leray spectral sequence

$E^{pq}_2=H^p(S_{\cl}, R^qf_{\cl\!\ast}F)\Longrightarrow
H^{p+q}(X_{\cl},F)$

\noindent
thus reduces to the exact sequence
\begin{equation*}
\ldots\longrightarrow H^q(S_{\cl},f_{\cl\!\ast}F)\longrightarrow
H^q(X_{\cl},F)\longrightarrow H^{q-1}(S_{\cl},
R^1f_{\cl\!\ast}F)\longrightarrow\tag{$\ast$}
\end{equation*}
Let $s$ be the image of $x$ in $S$. By the induction hypothesis and by
(i), for every class $\eta\in H^q(S_{\cl},G)$ ($G$ locally constant
torsion, with finite fibers), there exists an ``étale neighborhood''
$S'\rightarrow S$ of $s$ such that the image of $\eta$ in
$H^q(S'_{\cl},G)$ is zero. Now the functors $R^qf_{\cl\!\ast}$ plainly
commute with étale base changes, and if $S'\rightarrow S$ is an étale
neighborhood of $s$, then
$X'=\fprod{X}{S'}{S}\rightarrow X$ is an étale neighborhood of $x$.
When $q\geq 1$, we therefore conclude by means of the exact sequence
$(\ast)$, first taking $G=R^1f_{\cl\!\ast}F$ and
then\pageoriginale $G=f_{\cl\!\ast}F$. % The proof ends here; the
                                % following is a remark for the case $q=1$
When $q=1$, note that $H^1(X_{\cl},\ZZ/n \ZZ)$ (\resp
$H^1(X,\ZZ/n\ZZ)$) classifies the principal étale coverings
% should one instead speak of finite coverings for $X_\cl$?
of $X_{\cl}$ (\resp $X$), with group $\ZZ/n\ZZ$; hence, by
\hyperref[XI.thm4.3]{4.3} (iii), the map
$H^1(X,\ZZ/n\ZZ)\rightarrow H^1(X_{\cl},\ZZ/n\ZZ)$ is bijective,
from which the effacement result \hyperref[XI.lem4.5]{4.5} for $q=1$
follows at once.

\begin{enonce*}[remark]{Variant 4.6}\etiquette[4.6]{XI.var4.6}
One may also prove (\hyperref[XI.lem4.5]{4.5}) in the following elegant
way, due to Serre: one notes that a connected good neighborhood
(\hyperref[XI.def3.2]{3.2}) $X$ satisfies the following two conditions:
\begin{enumerate}
\renewcommand{\theenumi}{\roman{enumi}}
\renewcommand{\labelenumi}{(\theenumi)}
\item $\pi_n(X(\CC))=0$ if $n>1$.

\item $\pi_1(X(\CC))$ is an iterated extension of free groups of finite
type.
\end{enumerate}
\end{enonce*}

Indeed, if $X\rightarrow S$ is an elementary fibration where $S$ is a
good neighborhood, we may suppose (i) and (ii) true for $S$ by
induction. Now $X(\CC)\rightarrow S(\CC)$ is a locally trivial fiber
space, as is readily seen, with fibers isomorphic to $X_0(\CC)$, where
$X_0$ is an arbitrary fiber of $X/S$. Now $X_0$ is a connected
nonsingular \emph{noncomplete} curve, and it is well known that this
implies that $\pi_n(X_0(\CC))=0$ if $n>1$, and that
$\pi_1(X_0(\CC))$ is a free group. % \ref
Then (i) and (ii) follow from the homotopy exact sequence
$$
\longrightarrow \pi_0(X_0(\CC))\longrightarrow
\pi_n (X(\CC))\longrightarrow \pi_n
(S(\CC))\longrightarrow\ldots
$$

Thus $X(\CC)$ is a $K(\pi,1)$ space, and this implies that every
cohomology class $\xi\in H^n(X(\CC),F)$ becomes zero on the universal
covering of $X(\CC)$. It is enough to prove that $\xi$ already becomes
zero on a finite covering $X'$ of $X$, which amounts to proving that
$\pi_1=\pi_1(X(\CC))$ is\pageoriginale a \emph{good group}
(\cf CG 1 16), %\ref
that is, that the morphism from $\pi_1$ to its completion
$\sisi{\widehat}{}{\pi}_1=\varprojlim_N \pi_1/N$ (where $N$ runs
through the set of invariant subgroups % 'invariant' = 'normal'?
of $\pi_1$ of finite index) induces a bijection
$$
H^q(\widehat{\pi}_1,M)\xrightarrow{\sim} H^q(\pi_1,M)
$$
for every continuous finite $\pi_1$-module $M$ (where the symbol on the
left denotes the cohomology of $\widehat{\pi}_1$ as a profinite group).
Now an iterated extension of free groups of finite type is a good group
(\cf CG I p.~15-16, exc.~1,2), whence the result. % \ref

\begin{thebibliography}{99}
\bibitem{sga4.exp.XI.bib.1} GIRAUD J., ``Analysis Situs'', Séminaire
Bourbaki, Exposé nᵒ~256.

\bibitem{sga4.exp.XI.bib.2} GRAUERT H\sisi{.}{} and REMMERT R., ``Komplexe Räume'',
Math, Ann. Bd. 136, p. 245 (1958).

\bibitem{sga4.exp.XI.bib.3} HIRONAKA H., ``Resolution of singularities of an
algebraic variety over a field of characteristic zero'' Annals of
Math. Vol. 39, nᵒ~1, 2 (1964).

\bibitem{sga4.exp.XI.bib.4} SERRE J.P., ``Géométrie Algébrique et Géométrie
Analytique'', Annales de l'Institut Fourier, t. VI, p. 1-12 (1956).
\end{thebibliography}


\ifx\danslelivre\undefined
\end{document}
\fi

% END INCLUDED SGA4_English_current_buildable_source_20260906/source/11/drafts/SGA4_XI_LINES_354_643_English.texfrag


% END INCLUDED SGA4_English_current_buildable_source_20260906/source/11/11_English.tex


\makeatletter
\def\input@path{{../12/}}
\makeatother
% BEGIN INCLUDED SGA4_English_current_buildable_source_20260906/source/12/12_English.tex
\ifx\danslelivre\undefined
% Release engine: XeLaTeX (the LuaTeX compatibility layer breaks Xy-pic output).
\documentclass[10pt]{../fichiers-configuration-latex/smfart}
% BEGIN INCLUDED SGA4_English_current_buildable_source_20260906/source/fichiers-configuration-latex/paquets.tex
\usepackage{polyglossia}
%\usepackage[francais]{babel}
\usepackage{fontspec}
\usepackage[math-style=ISO,bold-style=upright,sans-style=upright]{unicode-math}
%\usepackage[mathletters]{ucs}
%\usepackage[utf8]{inputenc}
%\usepackage[T1]{fontenc}
\usepackage[unicode,hypertexnames=false,colorlinks=true,linkcolor=blue,citecolor=Sepia,urlcolor=magenta]{hyperref}
\setmainlanguage{french}
\setmainfont[Ligatures=TeX]{Linux Libertine O}
% variante possible : TeXGyrePagella ?
\setmathfont{XITS Math}

% maths
\usepackage{amsmath}
%\usepackage{amssymb}
%\usepackage{mathrsfs}
%\usepackage{stmaryrd}
%\usepackage{euscript}

% couleurs
\usepackage[usenames,dvipsnames]{xcolor}

% graphiques
\usepackage{graphicx}
\usepackage[all]{xy}
%\usepackage{tikz}
%\usetikzlibrary{matrix,arrows,calc}
%\tikzset{column sep=3em,row sep=2em,text height=1.5ex, text depth=0.25ex,>=angle 90}
%\tikzset{description/.style={fill=white,inner sep=1pt}}

% hyperliens, 'etiquettes
\usepackage{xr-hyper}
\usepackage{../fichiers-configuration-latex/smfthm}
%\ifx\cacheretiquettes\undefined
%\usepackage[inline]{showlabels} \else \fi % pour rendre visibles les 'etiquettes
%\usepackage{../fichiers-configuration-latex/makeidx} % pour faire un index
%\usepackage[pagebackref,unicode]{hyperref}  % pour voir o`u sont cit'ees les r'ef'erences
% hyperref is loaded once above with the cumulative-reader destination policy.
%\hypersetup{breaklinks=true} %permet le retour `a la ligne dans les liens trop longs. Incompatible avec pdflatex
%\renewcommand*{\backref}[1]{{\small [cit\'e page(s)~#1.]}}

\usepackage{xspace}
\usepackage{pifont} % pour expose XVII

% END INCLUDED SGA4_English_current_buildable_source_20260906/source/fichiers-configuration-latex/paquets.tex

% BEGIN INCLUDED SGA4_English_current_buildable_source_20260906/source/fichiers-configuration-latex/numerotation.tex
\renewcommand{\thefootnote}{\roman{footnote}} % pour mettre les notes en bas de page sous la forme (ii) et non (2), qui ressemble trop (?) a un exposant.

\newcounter{pageoriginal}
\newcommand{\pageoriginale}{\refstepcounter{pageoriginal}\marginpar{\footnotesize\textbf{\thepageoriginal}}\xspace}
\let\pageoriginaled\pageoriginale

\usepackage{ifthen}
\newboolean{original}
\def\sisi{\ifthenelse{\boolean{original}}}
\setboolean{original}{false}

\newcounter{cnote}
\def\nde#1{\ifthenelse{\boolean{original}}{}{\stepcounter{cnote}\def\thefootnote{\arabic{cnote}}\protect\footnote[\value{cnote}]{N.D.E.: #1}}}
\def\ndemark{\ifthenelse{\boolean{original}}{}{\stepcounter{cnote}\def\thefootnote{\arabic{cnote}}\protect\footnotemark[\value{cnote}]}}
\def\ndetext#1{\ifthenelse{\boolean{original}}{}{\def\thefootnote{\arabic{cnote}}\protect\footnotetext[\value{cnote}]{N.D.E.: #1}}}

% END INCLUDED SGA4_English_current_buildable_source_20260906/source/fichiers-configuration-latex/numerotation.tex

% BEGIN INCLUDED SGA4_English_current_buildable_source_20260906/source/fichiers-configuration-latex/macros.tex
\setmainlanguage{french}
\setotherlanguage{english}


%% Fixed, engine-independent release identifier.  The inherited dynamic Git
%% probe used Lua io.popen, exposed build-state-dependent output, and printed
%% a visible question mark outside a Git checkout.
\def\showgitstatus{English release r2, 2026-07-23}


% macros pour LaTeX
\renewcommand{\mathsf}{}

% Police différente pour U+2113 SCRIPT SMALL L
\setmathfont[range="2113]{latinmodern-math.otf}

% Give every inherited semantic label its own stable PDF destination instead
% of inheriting whichever automatic equation/page anchor happened to be live.
% The legacy label remains as a compatibility alias used by existing edges.
\makeatletter
\def\etiquette[#1]#2{%
  \Hy@raisedlink{\hyper@anchorstart{sga4.semantic.#2}\hyper@anchorend}%
  \def\@currentHref{sga4.semantic.#2}%
  \def\@currentlabel{#1}%
  \label{#2}%
}
\makeatother
\def\eqhyperref[#1]#2{(\hyperref[#1]{#2})}
\def\pagehyperref[#1]#2{(page \hyperref[#1]{#2})}
\def\hyperrefext[#1][#2]#3{\hyperref[#2]{#1 #3}}
\def\eqhyperrefext[#1][#2]#3{\hyperref[#2]{#1 #3}}     %cf. exposes 17
\def\neweqhyperrefext[#1][#2]#3{\hyperref[#2]{#1 #3}}  %cf. exposes 17

\def\shat#1{\protect\widehat{#1}}
\def\stilde#1{\protect\widetilde{#1}}
%%%%%%%%%%%%%%%%%%


\def\etcn{etc.\xspace}
\def\etco{etc...\xspace}
\def\etc{\sisi{\etco}{\etcn}}
\def\ldotsn{\sisi{...}{\etcn}}

\def\cf{{cf.}\kern.3em}
\def\Cf{{Cf.}\kern.3em}
\def\ie{{i.e.}\kern.3em}
\def\iev{{i.e.},\xspace}
\def\loccito{loc.\kern2pt cit.\xspace}
\def\loccitn{\textit{loc.\kern2pt cit.}\xspace}
\def\loccit{\sisi{\loccito}{\loccitn}}
\def\EGA{\textup{EGA}\kern.2em}
\def\SGA{\textup{SGA}\kern.3em}

\def\quoi{}

\def\quav{\quad,\quad}
\def\quaet{\quad \text{et}\quad}
\def\fibreproduct#1#2#3{#1 ×_{#3} #2}
\let\fprod\fibreproduct
\def\ccheck#1#2{\mbox{${\check{\check{\rm #1}}}^{\lower2pt\hbox{$\scriptstyle #2$}}$}}
\def\sheaf#1{\protect\mathcal{#1}}
\def\vieuxsheaf#1{{#1}} % cf. problème avec sheaf{a\Phi_i} par exemple...
\DeclareMathOperator{\Hom}{Hom}
\DeclareMathOperator{\can}{can}
\DeclareMathOperator{\Ind}{Ind}
\DeclareMathOperator{\ind}{ind}
%\DeclareMathOperator{\Hhomtop}{\mathcal{H}\!\mathit{omtop}}
\DeclareMathOperator{\Hhomtop}{\mathbf{Homtop}}
\DeclareMathOperator{\Point}{Point}
\DeclareMathOperator{\Pro}{Pro}
\DeclareMathOperator{\Top}{Top}
\DeclareMathOperator{\Mod}{Mod}
\DeclareMathOperator{\Cart}{Cart}
%\DeclareMathOperator{\Mmorsite}{\mathcal{M}\!\mathit{orsite}}
\DeclareMathOperator{\Mmorsite}{\mathbf{Morsite}}
\DeclareMathOperator{\loc}{loc}
\DeclareMathOperator{\Loc}{Loc}
\DeclareMathOperator{\Pic}{Pic}
\DeclareMathOperator{\Et}{Et}
\DeclareMathOperator{\card}{card}
\DeclareMathOperator{\Ens}{Ens}
\DeclareMathOperator{\h}{\underline{H}}
\DeclareMathOperator{\Fonct}{Fonct}
\DeclareMathOperator{\Ker}{Ker}
\DeclareMathOperator{\but}{but}
\DeclareMathOperator{\Ob}{Ob}
\DeclareMathOperator{\ob}{ob}
\DeclareMathOperator{\Lib}{Lib}
\DeclareMathOperator{\esj}{esj}
\DeclareMathOperator{\id}{id}
\DeclareMathOperator{\Id}{Id}
\DeclareMathOperator{\ad}{ad}
\DeclareMathOperator{\Iim}{Im}
\DeclareMathOperator{\pr}{pr}
\DeclareMathOperator{\Cat}{Cat}
\DeclareMathOperator{\car}{car}
\DeclareMathOperator{\SysInd}{SysInd}
\DeclareMathOperator{\Sup}{Sup}
\DeclareMathOperator{\Coker}{Coker}
\DeclareMathOperator{\Filt}{Filt}
\DeclareMathOperator{\acc}{acc}
\DeclareMathOperator{\sur}{sur}
\DeclareMathOperator{\cat}{cat}
\DeclareMathOperator{\Ord}{Ord}
\DeclareMathOperator{\Gr}{Gr}
\DeclareMathOperator{\Cov}{Cov}
\DeclareMathOperator{\Caf}{Caf}
\DeclareMathOperator{\coker}{coker}
\DeclareMathOperator{\coim}{coim}
\DeclareMathOperator{\im}{im}
\DeclareMathOperator{\Ouv}{Ouv}
\DeclareMathOperator{\diag}{diag}
\DeclareMathOperator{\Esp}{Esp}
\DeclareMathOperator{\TOP}{TOP}
\DeclareMathOperator{\Unif}{Unif}
\DeclareMathOperator{\Iind}{\mathbf{Ind}}
\DeclareMathOperator{\Hhomcart}{\mathbf{Homcart}}
\DeclareMathOperator{\Hhomtopess}{\mathbf{Homtopess}}
\DeclareMathOperator{\Ppoints}{\mathbf{Points}}
\DeclareMathOperator{\Ppoint}{\mathbf{Point}}
\DeclareMathOperator{\Ttors}{\mathbf{Tors}}
\DeclareMathOperator{\Ssex}{\mathbf{Sex}}
\DeclareMathOperator{\Ffib}{\mathbf{Fib}}
\DeclareMathOperator{\Vvois}{\mathbf{Vois}}
\DeclareMathOperator{\Ppointess}{\mathbf{Pointess}}
\DeclareMathOperator{\scrHom}{\mathbf{Hom}}
\DeclareMathOperator{\sheafPic}{\mathbf{Pic}}
\DeclareMathOperator{\sheafAut}{\mathbf{Aut}}
\DeclareMathOperator{\Hhomtoplocan}{\mathbf{Homtoplocan}}
\DeclareMathOperator{\Hhomtopan}{\mathbf{Homtopan}}
\DeclareMathOperator{\Pptgeom}{\mathbf{Ptgéom}}
\DeclareMathOperator{\Morsite}{Morsite}
\DeclareMathOperator{\Fib}{Fib}
\DeclareMathOperator{\Fais}{Fais}
\DeclareMathOperator{\kar}{kar}
\DeclareMathOperator{\Groupes}{Groupes}
\DeclareMathOperator{\Restr}{Restr}
\DeclareMathOperator{\Isom}{Isom}
\DeclareMathOperator{\Prol}{Prol}
\DeclareMathOperator{\Res}{Res}
\DeclareMathOperator{\Tors}{Tors}
\DeclareMathOperator{\esp}{esp}
\DeclareMathOperator{\es}{es}
\DeclareMathOperator{\od}{od}
\DeclareMathOperator{\gr}{gr}
\DeclareMathOperator{\sob}{sob}
\DeclareMathOperator{\dfn}{dfn}
\DeclareMathOperator{\Sch}{Sch}
\DeclareMathOperator{\ouv}{ouv}
\DeclareMathOperator{\cons}{cons}
\DeclareMathOperator{\Nerf}{Nerf}
\DeclareMathOperator{\Points}{Points}
\DeclareMathOperator{\Espsob}{Espsob}
\DeclareMathOperator{\Homtop}{Homtop}
\DeclareMathOperator{\Pt}{Pt}
\DeclareMathOperator{\Quot}{Quot}
\DeclareMathOperator{\glob}{glob}
\DeclareMathOperator{\Homfib}{Homfib}
\DeclareMathOperator{\Qouv}{Qouv}
\DeclareMathOperator{\Quotouv}{Quotouv}
\DeclareMathOperator{\ens}{ens}
\DeclareMathOperator{\ab}{ab}
\DeclareMathOperator{\an}{an}
\DeclareMathOperator{\fs}{fs}
\DeclareMathOperator{\red}{red}
\DeclareMathOperator{\Ab}{Ab}
\DeclareMathOperator{\cosq}{cosq}
\DeclareMathOperator{\Ann}{Ann}
\DeclareMathOperator{\Zar}{Zar}


%%%%%%%%%%%%
%%%%% SGA4-2
\DeclareMathOperator{\Sex}{Sex}
\DeclareMathOperator{\vois}{vois}
\DeclareMathOperator{\Vois}{Vois}
\DeclareMathOperator{\Ext}{Ext}
\DeclareMathOperator{\ext}{ext}
\DeclareMathOperator{\Fer}{Fer}
\DeclareMathOperator{\Add}{Add}
\DeclareMathOperator{\Cocart}{Cocart}
\DeclareMathOperator{\cosk}{cosk}
\DeclareMathOperator{\plat}{plat}

\DeclareMathOperator{\faisc}{faisc}
\DeclareMathOperator{\sk}{sk}
\DeclareMathOperator{\ch}{ch}

\def\HH{\mathbf{H}}
\def\qc{\mathrm{qc}}
\def\et{\mathrm{et}}
\def\coh{\mathrm{coh}}

\def\Hvr#1{%
\begin{array}[b]{l}
{\kern2pt\scriptstyle r}\\[-8pt]
{\kern1pt\scriptstyle \vee}\\[-5pt]
H^{#1}
\end{array}}

\def\Hvv#1{%
\begin{array}[b]{l}
\kern2pt\rotatebox{90}{${\scriptscriptstyle\ll}$}\\[-6pt]
H^{#1}
\end{array}}

\def\Hvi#1{%
\begin{array}[b]{l}
{\kern0.2pt\scriptscriptstyle \infty}\\[-8pt]
{\kern1pt\scriptstyle \vee}\\[-5pt]
H^{#1}
\end{array}}

\def\Hv#1{%
\begin{array}[b]{l}
{\kern1pt\scriptstyle \vee}\\[-5pt]
H^{#1}
\end{array}}

\DeclareMathOperator{\cart}{cart}
\DeclareMathOperator{\Ffaisc}{\mathbf{Faisc}}


\def\Ev#1{%
\begin{array}[b]{l}
{\kern1pt\scriptstyle \vee}\\[-5pt]
E'_{#1}
\end{array}}

\DeclareMathOperator{\Cont}{Cont}
\DeclareMathOperator{\Ccont}{\mathscr{C}\!\mathit{ont}}
\DeclareMathOperator{\Site}{Site}
\newcommand\MM{\mathbf{M}}
\def\limtop#1{{\displaystyle{\varprojlim_{#1}\!\mbox{\rm top }}}}
\def\Limtop{\mathop{\varprojlim\!\mbox{\rm top}}}
\DeclareMathOperator{\spec}{spec}
\def\ss{\scriptstyle{\mathscr{S}}}
\DeclareMathOperator{\Gal}{Gal}
\def\smallm{\scriptstyle{\mathscr{M}}}
\DeclareMathOperator{\End}{End}
\DeclareMathOperator{\ass}{ass}
\DeclareMathOperator{\zar}{zar}
\DeclareMathOperator{\Ogg}{Ogg}
\DeclareMathOperator{\res}{res}
\DeclareMathOperator{\tr}{tr}
\DeclareMathOperator{\Supp}{Supp}
\DeclareMathOperator{\codim}{codim}
\DeclareMathOperator{\rk}{rk}
\DeclareMathOperator{\cur}{cur}
\DeclareMathOperator{\cdqc}{cdqc}
\newcommand\DD{\mathbf{D}}
\DeclareMathOperator{\cl}{cl}
\def\JJ{\mathbf{J}}
\def\FF{\mathbf{F}}
\def\BB{\mathbf{B}}
\DeclareMathOperator{\Alget}{Alget}
\DeclareMathOperator{\Ass}{Ass}
\DeclareMathOperator{\supp}{supp}
\DeclareMathOperator{\rad}{rad}
\DeclareMathOperator{\prof}{prof}
\DeclareMathOperator{\cint}{int}
\DeclareMathOperator{\sic}{sic}
\DeclareMathOperator{\Tor}{Tor}
\DeclareMathOperator{\tors}{tors}
\DeclareMathOperator{\Sym}{Sym}
\DeclareMathOperator{\parf}{parf}
\DeclareMathOperator{\tor}{tor}
\DeclareMathOperator{\defn}{defn}
\DeclareMathOperator{\Div}{Div}
\DeclareMathOperator{\Pol}{Pol}
\DeclareMathOperator{\Fex}{Fex}
\DeclareMathOperator{\Lie}{Lie}
\DeclareMathOperator{\tot}{tot}
\DeclareMathOperator{\Aut}{Aut}
\DeclareMathOperator{\HOM}{HOM}
\DeclareMathOperator{\pch}{pch}
\DeclareMathOperator{\Ch}{Ch}
\DeclareMathOperator{\EXT}{EXT}
\DeclareMathOperator{\PIC}{PIC}
\DeclareMathOperator{\Rf}{Rf}
\DeclareMathOperator{\inj}{inj}
\newcommand\VV{\mathbf{V}}

\def\Xv{%
\begin{array}[b]{l}
{\kern1pt\scriptstyle \vee}\\[-5pt]
X
\end{array}}

%%%%%%%
\newcommand\TT{\mathbf{T}}
\newcommand\XX{\mathbf{X}}
\newcommand\EE{\mathbf{E}}
%%%%%%

\def\qedsymbol{\textup{C.Q.F.D.}}
\let\CQFD\qed
\let\cqfd\qed

%%%%%\bigdot
\font\bigsymb=cmsy10 at 4pt
\def\bigdot{{\kern1.2pt\raise 1.5pt\hbox{\bigsymb\char15}}}
%%%%%

\def\Lright{\mathop{\vcenter{\hbox{\hbox{$\longrightarrow$}\lower8pt \hbox{\kern-16pt\hbox{$\longrightarrow$}}}}}}

\def\Aa{{\mathcal A}}
\def\Bb{{\mathcal B}}
\def\Cc{{\mathcal C}}
\def\Dd{{\mathcal D}}
\def\Ee{{\mathcal E}}
\def\Ff{{\mathcal F}}
\def\Gg{{\mathcal G}}
\def\Hh{{\mathcal H}}
\def\Ii{{\mathcal I}}
\def\Jj{{\mathcal J}}
\def\Kk{{\mathcal K}}
\def\Ll{{\mathcal L}}
\def\Mm{{\mathcal M}}
\def\Nn{{\mathcal N}}
\def\Oo{{\mathcal O}}
\def\Pp{{\mathcal P}}
\def\Qq{{\mathcal Q}}
\def\Rr{{\mathcal R}}
\def\Ss{{\mathcal S}}
\def\Tt{{\mathcal T}}
\def\Uu{{\mathcal U}}
\def\Vv{{\mathcal V}}
\def\Ww{{\mathcal W}}
\def\Xx{{\mathcal X}}
\def\Yy{{\mathcal Y}}
\def\Zz{{\mathcal Z}}

\newcommand{\CC}{\mathbf{C}}
\newcommand{\GG}{\mathbf{G}}
\newcommand{\KK}{\mathbf{K}}
\newcommand{\LL}{\mathbf{L}}
\newcommand{\NN}{\mathbf{N}}
\newcommand{\PP}{\mathbf{P}}
\newcommand\QQ{\mathbf{Q}}
\newcommand\RR{\mathbf{R}}
\newcommand\ZZ{\mathbf{Z}}

\def\to{\mathchoice{\longrightarrow}{\rightarrow}{\rightarrow}{\rightarrow}}
\def\from{\mathchoice{\longleftarrow}{\leftarrow}{\leftarrow}{\leftarrow}}
\def\mto{\mathchoice{\longmapsto}{\mapsto}{\mapsto}{\mapsto}}
\def\mfrom{\mathrel{\from\joinrel\mapstochar\,}}
\def\hto{\mathrel{\lhook\joinrel\to}}
\def\hfrom{\mathrel{\from\joinrel\rhook}}

\def\implies{\mathchoice{\Longrightarrow}{\Rightarrow}{\Rightarrow}{\Rightarrow}}
\def\implied{\mathchoice{\Longleftarrow}{\Leftarrow}{\Leftarrow}{\Leftarrow}}
\def\iff{\mathchoice{\Longleftrightarrow}{\Leftrightarrow}{\Leftrightarrow}{\Leftrightarrow}}

\def\To#1{\mathchoice{\xrightarrow{\labelstyle~~#1~}}{\stackrel{#1}{\longrightarrow}}{}{}}
\def\Mto#1{\mathrel{\mapstochar\To{#1}}}
\def\From#1{\mathchoice{\xleftarrow{\labelstyle~~#1~}}{\stackrel{#1}{\longleftarrow}}{}{}}

\newcommand{\isomto}{\overset{\sim}{\longrightarrow}}
\newcommand{\isomfrom}{\overset{\sim}{\longleftarrow}}

\newcommand{\IFF}{\unskip~$\iff$~\ignorespaces}
\newcommand{\THEN}{\unskip~$\implies$~\ignorespaces}

\let\leq\leqslant
\let\geq\geqslant
\let\le\leqslant
\let\ge\geqslant

\def\cf{{cf.}\kern.3em}
\def\Cf{{Cf.}\kern.3em}
\def\ie{{i.e.}\kern.3em}
\def\iev{{i.e.},\xspace}
%\def\loccito{loc.\kern2pt cit.\xspace}
%\def\loccitn{\textit{loc.\kern2pt cit.}\xspace}
%\def\loccit{\sisi{\loccito}{\loccitn}}
\def\EGA{\textup{EGA}\kern.2em}
\def\SGA{\textup{SGA}\kern.3em}

\let\oldS\S
\def\S{\oldS\kern .15em}
\def\ptbl{.\kern .2em}
\def\resp{resp.\kern.3em}
\def\Exp{Exp.\kern.3em}

\let\bar\overline
\let\hat\widehat
\let\tilde\widetilde
\let\dpl\displaystyle
\let\ldots\dots
\let\emptyset\varnothing
\let\skpt\skippointrait
\let\epsilon\varepsilon

\def\sfrac#1#2{{#1}/{#2}}

\DeclareMathOperator{\cd}{cd}
%\DeclareMathOperator{\Hhom}{\mathcal{H}\!\mathit{om}}
\DeclareMathOperator{\Hhom}{\mathbf{Hom}}
%\DeclareMathOperator{\Eext}{\mathcal{E}\!\mathit{xt}}
\DeclareMathOperator{\Eext}{\mathbf{Ext}}
\DeclareMathOperator{\Spec}{Spec}
\DeclareMathOperator{\Tr}{Tr}

\let\emptyset=\varnothing

%%%%%%%%%%%%%%%
%%% expose XVII
%%%%%%%%%%%%%%%

\makeatletter
\let\old@optexpose\optexpose
\newcommand{\empty@optexpose}[1]{}
\newcommand{\paren@optexpose}[1]{\old@optexpose{#1}(}
\newcommand{\refnoexpose}[1]{\let\optexpose\empty@optexpose\ref{#1}\let\optexpose\old@optexpose}%\refnoexpose{X.cor5.2} engendre "5.2" au lieu de "X 5.2"
\newcommand{\neweqref}[1]{\let\optexpose\paren@optexpose\textup{\ref{#1})}\let\optexpose\old@optexpose}%\neweqref{XII.eq4.1} engendre "XII (4.1)" au lieu de "(XII 4.1)"
\makeatother
\newcommand{\refvb}[2][VI${}_{\text{B}}$ ]{\sisi{#1\refnoexpose{#2}}{\ref{#2}}}

\newcommand{\leftexp}[2]{{\vphantom{#2}}^{#1}{#2}}
\newcommand{\AAA}{\mathbf{A}}
\newcommand{\Sf}{\mathfrak{S}}
\DeclareMathOperator*{\quotprojlim}{\text{``$\varprojlim$''}}
\DeclareMathOperator*{\quotinjlim}{\text{``$\varinjlim$''}}
\DeclareMathOperator{\Fl}{F\sisi{\ell}{l}}
\DeclareMathOperator{\Pl}{P\sisi{\ell}{l}}
%\DeclareMathOperator{\cl}{c\sisi{\ell}{l}}
\DeclareMathOperator{\tdf}{tdf}
\DeclareMathOperator{\rg}{rg}
\newcommand{\iInd}{\sisi{I}{i}nd}
\newcommand{\iind}{ind}
\newcommand{\ppro}{pro}
\newcommand{\tTor}{\sisi{T}{t}or}
\newcommand{\ttor}{tor}
\newcommand{\chap}{Chap.}
\newcommand{\Ffex}{\Ff\mathit{ex}}
\newcommand{\sheafk}{\underline{k}}
\newcommand{\sheafn}{\underline{n}}
\newcommand{\sheafGamma}{\underline{\Gamma}}

% Repeated references to an already typeset footnote must not use
% \footnotemark without a matching \footnotetext: Hyperref would then create
% a live annotation to a destination that cannot exist.  This helper keeps the
% current printed footnote mark and links it to an explicit stable note target.
\makeatletter
\newcommand{\sgarepeatedfootnotemark}[1]{%
  \begingroup
  \edef\@thefnmark{\thefootnote}%
  \hyperlink{#1}{\@makefnmark}%
  \endgroup}
\makeatother

\sisi{
\DeclareMathOperator{\simarrow}{\rule{.15mm}{1.8mm}\kern -3.4pt\rightsquigarrow}
\newcommand{\sheafell}{\underline{\ell}}
\newcommand{\sheafZZ}{\underline{\ZZ}}
}{
\newcommand{\nderef}[1]{la N.D.E.~(\ref{#1}) page~\pageref{#1}}
}

%%%%%%%%%%%%%%%%%
%%%%% expose XVII
%%%%%%%%%%%%%%%%%

% END INCLUDED SGA4_English_current_buildable_source_20260906/source/fichiers-configuration-latex/macros.tex


\externaldocument{../01/01_English.tex}
\externaldocument{../02/02_English.tex}
\externaldocument{../03/03_English.tex}
\externaldocument{../04/04_English.tex}
\externaldocument{../05/05_English.tex}
\externaldocument{../05bis/05bis_English.tex}
\externaldocument{../06/06_English.tex}
\externaldocument{../07/07_English.tex}
\externaldocument{../08/08_English.tex}
\externaldocument{../09/09_English.tex}
\externaldocument{../10/10_English.tex}
\externaldocument{../11/11_English.tex}

\setboolean{original}{false}

\begin{document}

\begin{center}
XII. The base-change theorem for a proper morphism\\
M. Artin
\end{center}

\selectlanguage{english}version: \showgitstatus
\tableofcontents

\else
\chapter{The base-change theorem for a proper morphism}
%\addtocontents{toc}{M. Artin}
\begin{center}
M. Artin
\end{center}
\fi
\phantomsection\MakeLinkTarget*{sga4.exp.XII}\label{sga4.exp.XII}

% BEGIN INCLUDED SGA4_English_current_buildable_source_20260906/source/12/drafts/SGA4_XII_LINES_43_233_English.texfrag
\setcounter{pageoriginal}{78}

\section{Introduction}\etiquette[1]{XII.sec1}\pageoriginale

In this exposé we prove the fundamental theorem which says, in a
slightly more precise form, that for a \emph{proper} morphism
$f:X\to Y$ of schemes and a torsion abelian sheaf $F$ on $X$, the
fiber of $R^qf_*F$ at a geometric point $\xi$ of $Y$ is isomorphic to
the cohomology $H^q(X_\xi,F|X_\xi)$ of the fiber $X_\xi$.

Recall that the analogous result for a proper morphism of paracompact
spaces is well known (Godement II 4.11) and easy. It is considerably
more delicate for schemes. For example, in contrast with the case of
paracompact spaces, the hypothesis that $F$ be torsion is essential.

The proof proceeds in several stages. By more or less formal
dévissage arguments, one reduces to the case where $Y$ is
Noetherian, the relative dimension of $f$ is at most $1$, and the
sheaf $F$ is constant. To handle this special case, one uses the
specialization theorem for the fundamental group
(\hyperref[XII.thm5.9]{5.9}), which is a nonabelian variant of the
theorem stated above and can be proved directly in the case under
consideration (\hyperrefext[XIII][XIII.sec2]{2}).

\section{An example}\etiquette[2]{XII.sec2}\pageoriginale

We show that the hypothesis that $F$ be torsion is necessary. Let
$X$ be a nonsingular surface over an algebraically closed field $k$,
let $Y$ be a nonsingular curve over $k$, and let $f:X\to Y$ be a
proper morphism whose generic fiber is smooth and whose geometric
fiber at a geometric point $\bar y$ of $Y$ is irreducible and has an
ordinary double point. Let $F$ be the constant sheaf $\ZZ_X$.

Recall the following fact. For any point $z$ of a scheme $Z$, let
$i:z\to Z$ be the inclusion. Then
$H^1(Z,i_*\ZZ_z)=0$ (\hyperrefext[IX][IX.prop3.6]{3.6} (i)). Since
$X$ is normal, hence unibranch, it follows that
$\ZZ_X\simeq i_*\ZZ_x$, where $i:x\to X$ is the inclusion of the
generic point. Thus $H^1(X,\ZZ_X)=0$ and likewise
$R^1f_*\ZZ_X=0$. But for the geometric fiber $X_\circ$ of $X/Y$
having the ordinary double point, I claim that
$H^1(X_\circ,\ZZ_{X_\circ})\simeq\ZZ$\footnote[1]{Compare, from the
point of view of $\Pi_1$, with \SGA 3 X (6, Examples) and (1.6).};
this gives the counterexample, since
$\ZZ_{X_\circ}\simeq\ZZ_X|X_\circ$.

Let $i_\circ:x_\circ\to X_\circ$ be the inclusion of the generic
point of $X_\circ$, and let $Q$ be the singular point of $X_\circ$.
The curve $X_\circ$ has two \lq\lq branches'' at $Q$, and it follows
that the fiber of $i_{\circ*}(\ZZ_{x_\circ})$ at $Q$ is isomorphic to
$\ZZ\times\ZZ$. Away from $Q$, the curve $X_\circ$ is normal, hence
$i_{\circ*}(\ZZ_{x_\circ})\simeq\ZZ_{X_\circ}$ on
$X_\circ-Q$. Therefore there is an exact sequence
$$
0\longrightarrow\ZZ_{X_\circ}\longrightarrow
i_{\circ*}(\ZZ_{x_\circ})\longrightarrow(\ZZ)_Q
\longrightarrow0{\quoi},
$$
where the last term is extension by zero of the sheaf $\ZZ$ at the
point $Q$. All three terms have $H^0$ isomorphic to $\ZZ$, and
$H^1(X_\circ,i_{\circ*}(\ZZ_{x_\circ}))=0$. The long exact
cohomology sequence therefore gives
$H^1(X_\circ,\ZZ_{X_\circ})=\ZZ$\pageoriginale.

\section{Review of nonabelian $H^1$}\etiquette[3]{XII.sec3}

We briefly recall some results on nonabelian cohomology that will be
needed below. We omit most proofs, which are well known in various
special cases and may be found in Giraud's thesis [1].

Let $X$ be a scheme and $F$ a sheaf of groups on $X$. For every
étale $X'/X$, one defines a pointed set $H^1(X',F)$, covariantly
functorial in $F$ and contravariantly functorial in $X'$. The functor
$H^1$ commutes with finite products of sheaves. If $f:X\to Y$ is a
morphism, one defines a sheaf of pointed sets (that is, with a
distinguished section) $R^1f_*F$. It is the sheaf associated with the
presheaf of pointed sets $\Rr^1F$ which assigns to every étale
$Y'/Y$ the set $H^1(\fprod{X}{Y'}{Y},F)$.

There are the usual exact cohomology sequences; we record one of
them:

\begin{enonce*}{Proposition 3.1}\etiquette[3.1]{XII.prop3.1}
Let $0\to F\xrightarrow{u}G$ be an injection of sheaves of groups,
and let $C=G/F$ be the sheaf of homogeneous spaces under $G$. There
is an exact sequence of sheaves of pointed sets
$$
0\to f_*F\to f_*G\to f_*C\xrightarrow{\partial}
R^1f_*F\xrightarrow{u^1}R^1f_*G{\quoi}.
$$
\end{enonce*}

This exact sequence is obtained, of course, by passing to the
associated sheaves in the exact sequence of presheaves whose value
on an étale $Y'/Y$ is
$$
0\to F(X')\to G(X')\to C(X')\xrightarrow{\partial}
H^1(X',F)\to H^1(X',G){\quoi},
$$
where\pageoriginale $X'=\fprod{X}{Y'}{Y}$. The equivalence relation
induced by $\partial$ is the one given by the action of $f_*G$ on
$f_*C$. As usual, the equivalence relation induced by $u^1$ may be
made explicit using the notion of \lq\lq twisting'' [1].

\begin{enonce*}{Proposition 3.2}\etiquette[3.2]{XII.prop3.2}
Let $X\xrightarrow{f}Y\xrightarrow{g}Z$ be morphisms of schemes. For
every étale $Z'/Z$, putting
$Y'=\fprod{Y}{Z'}{Z}$ and $X'=\fprod{X}{Z'}{Z}$, there is an exact
sequence
$$
0\longrightarrow H^1(Y',f_*F)\longrightarrow H^1(X',F)
\longrightarrow H^0(Y',R^1f_*F){\quoi}.
$$
Passing to the associated sheaves gives an exact sequence of pointed
sheaves on $Z$:
$$
0\longrightarrow (R^1g_*)f_*F\longrightarrow R^1(gf)_*F
\longrightarrow g_*(R^1f_*)F{\quoi}.
$$
\end{enonce*}

\begin{enonce*}{Proposition 3.3}\etiquette[3.3]{XII.prop3.3}
The family of functors of the form $H^1(X',F)$, for $X'$ étale over
$X$, and of the form $R^1f_*F$, for an arbitrary morphism
$f:X\to Y$, is effaceable: there exists an injection
$F\xrightarrow{u}G$ such that all the induced morphisms
$H^1(X',F)\to H^1(X',G)$ and $R^1f_*F\to R^1f_*G$ are zero. More
precisely, there exists such an injection for which every
$H^1(X',G)$, with $X'$ étale over $X$, and every pointed sheaf
$R^1f_*G$ vanish. If $X$ is Noetherian and $F$ is a constructible
sheaf of ind-$\LL$-groups, then such an injection exists with $G$
also a sheaf of ind-$\LL$-groups.
\end{enonce*}

\emph{Proof.} The first assertion is presumably true in an arbitrary
topos, but for the étale topology it is convenient to perform the
effacement using the \lq\lq Godement resolution''
\begin{equation*}
F\longrightarrow G=\prod_{x\in X}i_{x*}i_x^*F{\quoi}.
\tag{3.4}
\end{equation*}\etiquette[3.4]{XII.eq3.4}
Here\pageoriginale $x$ runs through the points of $X$ and
$i_x:\bar x\to X$ is a geometric point over $x$. For $X'/X$ étale,
one has
$$
G(X')=\prod_{x\in X}(i_x^*F)(\fprod{X'}{\bar x}{X})
=\prod_{x\in X}\left(\prod_{\bar x'/\bar x}F_{\bar x'}\right){\quoi},
$$
where the $\bar x'$ are geometric points of $X'$. For each point of
$X'$ there is plainly at least one $\bar x'$ above it; hence
$F(X')\subset G(X')$ (VIII 3), and therefore $F\subset G$.
Moreover, every étale $\bar u/\bar x$ is a disjoint union of spectra
of separably closed fields, so $H^1(\bar u,\bar G)=0$ for every sheaf
$\bar G$ on $\bar x$. It follows from
\hyperref[XII.prop3.2]{3.2} that
$H^1(X',i_{x*}i_x^*F)=0$ for every étale $X'/X$. Since $H^1$
commutes with products, $H^1(X',G)=0$, and consequently
$R^1f_*G=0$ for every $f:X\to Y$, by the description of $R^1f_*G$
as the sheaf associated with the presheaf
$Y'\mapsto H^1(\fprod{X}{Y'}{Y},G)$.

Suppose in addition that $F$ is constructible and $X$ is Noetherian.
It is immediate that there is a finite set of points
$\{x_1,\ldots,x_n\}$ of $X$ such that
$$
F\longrightarrow\prod_{1\leq i\leq n}i_{x_i*}i_{x_i}^*F
$$
is already injective. The sheaf $i_{x_i*}i_{x_i}^*F$ is a sheaf of
ind-finite groups (IX \hyperref[IX.def1.5]{1.5}), and hence so is the finite product
(\hyperrefext[IX][IX.prop1.4]{1.4} and
\hyperrefext[IX][IX.prop1.6]{1.6} (iii)). Replacing $G$ by this sheaf
gives the last assertion.

\setcounter{subsection}{4}
\subsection{}\etiquette[3.5]{XII.subsec3.5}
As in the case of abelian cohomology (treated in V \hyperref[V.subsec2.4]{2.4}), $H^1$ may be
computed by the \v{C}ech procedure [1].

% END INCLUDED SGA4_English_current_buildable_source_20260906/source/12/drafts/SGA4_XII_LINES_43_233_English.texfrag


% BEGIN INCLUDED SGA4_English_current_buildable_source_20260906/source/12/drafts/SGA4_XII_LINES_234_382_English.texfrag
\section{The base-change morphism}\etiquette[4]{XII.sec4}\pageoriginale

Let
$$
\xymatrix@C=2.75cm@R=1.5cm{
X \ar[d]_{f}&X'\ar[l]_{g'}\ar[d]^{f'}\\
S&S'\ar[l]_{g}
}
$$
be a Cartesian square of morphisms of schemes. From it one obtains a
morphism of functors $\varphi=\varphi_g$
\begin{equation*}
g^{\ast}f_{\ast}\xrightarrow{\varphi}f'_{\ast}{g'}^{\ast}\tag{4.1}
\end{equation*}\etiquette[4.1]{XII.eq4.1}
as follows. Giving such a morphism is equivalent, by the adjunction
between $g^{\ast}$ and $g_{\ast}$, to giving a morphism of functors
\begin{equation*}
f_{\ast}\longrightarrow g_{\ast}{f'}_{\ast}{g'}^{\ast}{\quoi}.
\tag{4.1 adj}
\end{equation*}\MakeLinkTarget*{sga4.exp.XII.equation.4.1-adj}\label{4.1 adj}
We take this to be the morphism
$$
f_{\ast}\xrightarrow{(\id\rightarrow {g'}_{\ast}{g'}^{\ast})}
f_{\ast}{g'}_{\ast}{g'}^{\ast}
=g_{\ast}{f'}_{\ast}{g'}^{\ast}{\quoi},
$$
using
$$
g_{\ast}{f'}_{\ast}=(gf')_{\ast}=(fg')_{\ast}
=f_{\ast}{g'}_{\ast}{\quoi}.
$$

There is another candidate for \eqhyperref[XII.eq4.1]{4.1}. Indeed,
giving such a morphism is also equivalent to giving
$$
{f'}^{\ast}g^{\ast}f_{\ast}\longrightarrow {g'}^{\ast},
$$
and one may take
$$
{f'}^{\ast}g^{\ast}f_{\ast}\simeq
{g'}^{\ast}f^{\ast}f_{\ast}
\xrightarrow{(f^{\ast}f_{\ast}\rightarrow\id)}{g'}^{\ast}{\quoi}.
$$

I do not know\pageoriginale whether these two choices always coincide
\footnote[1]{Mr. \textsc{Deligne} resolved this uncertainty in the
affirmative, in a considerably more general context; \cf Exposé XVII.};
in any event, we shall use only the first. We call it the
\emph{base-change morphism.}

Restricting \eqhyperref[XII.eq4.1]{4.1} to abelian sheaves
(respectively, to sheaves of groups) gives a morphism
$\varphi^q=\varphi^q_g$ for every $q\geq0$ (respectively, for
$q=0,1$):
\begin{equation*}
g^{\ast}(R^qf_{\ast})\xrightarrow{\varphi^q}
(R^q{f'}_{\ast}){g'}^{\ast}{\quoi}.
\tag{4.2}
\end{equation*}\etiquette[4.2]{XII.eq4.2}
Indeed, this is equivalent to giving
\begin{equation*}
(R^qf_{\ast})\longrightarrow
g_{\ast}(R^q{f'}_{\ast}){g'}^{\ast}{\quoi},
\tag{4.2 adj}
\end{equation*}\MakeLinkTarget*{sga4.exp.XII.equation.4.2-adj}\label{4.2 adj}
and there are the following \lq\lq evident'' morphisms:
\begin{gather*}
(R^qf_{\ast})\xrightarrow{a}
(R^qf_{\ast}){g'}_{\ast}{g'}^{\ast}
\xrightarrow{b}R^q(fg')_{\ast}{g'}^{\ast}\\
=R^q(gf')_{\ast}{g'}^{\ast}
\xrightarrow{c}g_{\ast}(R^q{f'}_{\ast}){g'}^{\ast}{\quoi}.
\end{gather*}

\begin{enonce*}[remark]{Remark 4.3}\etiquette[4.3]{XII.rem4.3}
Suppose that $g$ (and hence $g'$) is a closed immersion. Then the
morphisms $b$ and $c$ above are isomorphisms (VIII \hyperref[VIII.cor5.7]{5.7}). It follows
that \eqhyperref[XII.eq4.2]{4.2} is essentially just the morphism $a$
above, obtained by applying $R^qf_{\ast}$ to the canonical morphism
$F\rightarrow {g'}_{\ast}{g'}^{\ast}F$.
\end{enonce*}

\begin{enonce*}{Proposition 4.4}\etiquette[4.4]{XII.prop4.4}
(Composition of base-change morphisms).
\begin{enumerate}
\renewcommand{\theenumi}{\roman{enumi}}
\renewcommand{\labelenumi}{(\theenumi)}
\item Let
$$
\xymatrix@C=2.5cm@R=1.5cm{
X\ar[d]_{f}&X'\ar[l]_{g'}\ar[d]_{f'}&
X''\ar[l]_{h'}\ar[d]_{f''}\\
S&S'\ar[l]_{g}&S''\ar[l]_{h}
}
$$
be a\pageoriginale Cartesian diagram of morphisms of schemes. Then
the following diagram is commutative:
$$
\xymatrix@C=2cm@R=1.5cm{
(gh)^{\ast}f_{\ast}\ar[rr]^{\varphi_{gh}}
\ar[d]^{\text{\rotatebox{90}{$\backsim$}}}&&
f''_{\ast}(g'h')^{\ast}
\ar[d]^{\text{\rotatebox{90}{$\backsim$}}}\\
h^{\ast}g^{\ast}f_{\ast}
\ar[r]^{h^{\ast}\varphi_g}&
h^{\ast}f'_{\ast}g'^{\ast}
\ar[r]^{\varphi_h g'^{\ast}}&
f''_{\ast}h'^{\ast}g'^{\ast}.}
$$
The same is true of the diagram obtained by replacing
$f_{\ast}$, $f'_{\ast}$, and $f''_{\ast}$ by
$R^qf_{\ast}$, $R^qf'_{\ast}$, and $R^qf''_{\ast}$, in the natural
domains of definition of the functors in question
\textup{(\cf $\textup{n}^{\circ}$ 3)}.

\item Let
$$
\xymatrix@C=3.25cm@R=1.25cm{
Y\ar[d]_{e}&Y'\ar[l]_{g''}\ar[d]^{e'}\\
X\ar[d]_{f}&X'\ar[l]_{g'}\ar[d]^{f'}\\
S&S'\ar[l]_{g}
}
$$
be a Cartesian diagram of morphisms of schemes. Then the following
diagram is commutative:
$$
\xymatrix@C=2.75cm@R=1.5cm{
g^{\ast}(fe)_{\ast}\ar[rr]^{\varphi_g^{fe}}\ar@{=}[d]&&
(f'e')_{\ast}g''^{\ast}\ar@{=}[d]\\
g^{\ast}f_{\ast}e_{\ast}
\ar[r]^{\varphi_g^f e_{\ast}}&
f'_{\ast}g'{}^{\ast}e_{\ast}
\ar[r]^{f'_{\ast}\varphi_{g'}^e}&
f'_{\ast}e'_{\ast}g''^{\ast},
}
$$
where the three horizontal arrows arise respectively from the
base-change homomorphisms for $fe$, $f$, and $e$.

In the abelian case, the base-change morphisms for $fe$ are induced by
a morphism of spectral sequences
$$
\xymatrix{
E^{p,q}_2=g^{\ast}(R^pf_{\ast})(R^qe_{\ast})F
\ar@{=>}[r]\ar[d]_{\varphi}&
g^{\ast}R^{p+q}(fe)_{\ast}F\ar[d]_{\varphi}\\
(R^pf'_{\ast})(R^qe'_{\ast})g''^{\ast}F
\ar@{=>}[r]&R^{p+q}(f'e')_{\ast}g''^{\ast}F,
}
$$
which, on the initial terms, is induced by the base-change morphisms
for $e$ and $f$. For a sheaf of groups, there is a morphism of exact
sequences
$$
\xymatrix{
0\ar[r]&g^{\ast}(R^1f_{\ast})e_{\ast}F
\ar[r]\ar[d]^{\varphi}&
g^{\ast}R^1(fe)_{\ast}F\ar[r]\ar[d]^{\varphi}&
g^{\ast}f_{\ast}(R^1e_{\ast})F\ar[d]^{\varphi}\\
0\ar[r]&(R^1f'_{\ast})e'_{\ast}g''^{\ast}F\ar[r]&
R^1(f'e')_{\ast}g''^{\ast}F\ar[r]&
f'_{\ast}(R^1e'_{\ast})g''^{\ast}F.
}
$$
\end{enumerate}
\end{enonce*}

We leave the proof to the reader.

% END INCLUDED SGA4_English_current_buildable_source_20260906/source/12/drafts/SGA4_XII_LINES_234_382_English.texfrag


% BEGIN INCLUDED SGA4_English_current_buildable_source_20260906/source/12/drafts/SGA4_XII_LINES_383_784_English.texfrag
\section{Statement of the main theorem and some variants}\etiquette[5]{XII.sec5}

The theorem is as follows:

\begin{enonce*}{Theorem 5.1}\etiquette[5.1]{XII.thm5.1}
(Proper base-change theorem). Let
$$
\xymatrix@C=3.25cm@R=1.5cm{
X \ar[d]_{f} & X' \ar[l]_{g'} \ar[d]^{f'}\\
S & S' \ar[l]_{g}
}
$$
be a Cartesian diagram of schemes, with $f$ proper.
\begin{enumerate}
\renewcommand{\theenumi}{\roman{enumi}}
\renewcommand{\labelenumi}{(\theenumi)}
\item Let $F$ be a sheaf of sets on $X$. Then the morphism
  \eqhyperref[XII.eq4.1]{4.1}
$$
\varphi:g^{\ast}f_{\ast}F\longrightarrow {f'}_{\ast}{g'}^{\ast}F
$$
is bijective.

\item Let\pageoriginale $F$ be a sheaf of groups (resp. of
  ind-finite groups) on $X$. Then the morphism
  \eqhyperref[XII.eq4.2]{4.2}
$$
\varphi^1:g^{\ast}(R^1f_{\ast})F\longrightarrow
(R^1{f'}_{\ast}){g'}^{\ast}F
$$
is injective (resp. bijective).

\item Let $F$ be a torsion abelian sheaf on $X$. Then the morphism
  \eqhyperref[XII.eq4.2]{4.2}
$$
\varphi^q:g^{\ast}(R^qf_{\ast})F\longrightarrow
(R^q{f'}_{\ast}){g'}^{\ast}F
$$
is bijective for every $q\geq0$.
\end{enumerate}
\end{enonce*}

We first state some immediate consequences of the theorem. If
$S'=\xi$ is a geometric point of $S$, so that a sheaf on $\xi$ is
determined by its global sections, then
$H^0(\xi,R^q{f'}_{\ast}G)=H^q(X',G)$ for every sheaf $G$ on $X'$.
This gives the following corollary (which is, moreover, essentially
equivalent to \hyperref[XII.thm5.1]{5.1}):

\begin{enonce*}{Corollary 5.2}\etiquette[5.2]{XII.cor5.2}
Let $f:X\longrightarrow S$ be proper, let $\xi\longrightarrow S$ be
a geometric point, and let $X_{\xi}$ be the fiber of $f$ at $\xi$.
\begin{enumerate}
\renewcommand{\theenumi}{\roman{enumi}}
\renewcommand{\labelenumi}{(\theenumi)}
\item If $F$ is a sheaf of sets on $X$, the canonical base-change
  morphism
$$
(f_{\ast}F)_{\xi}\longrightarrow H^0(X_{\xi},F|X_{\xi})
$$
is bijective.

\item If $F$ is a sheaf of groups (resp. of ind-finite groups), the
  base-change morphism
$$
(R^1f_{\ast}F)_{\xi}\longrightarrow H^1(X_{\xi},F|X_{\xi})
$$
is injective (resp. bijective).

\item If\pageoriginale $F$ is a torsion abelian sheaf, then for every
  $q\geq0$ the base-change morphism
$$
(R^qf_{\ast}F)_{\xi}\longrightarrow H^q(X_{\xi},F|X_{\xi})
$$
is bijective.
\end{enumerate}
\end{enonce*}

Taking \hyperrefext[X][X.cor4.3]{4.3} and
\hyperrefext[X][X.cor5.2]{5.2} into account, we obtain:

\begin{enonce*}{Corollary 5.3}\etiquette[5.3]{XII.cor5.3}
Under the hypotheses of \hyperrefext[X][X.cor5.2]{5.2}, let $n$ be
the dimension of the fiber $X_{\xi}$. Then, for every torsion
abelian sheaf $F$ on $X$, the fiber $(R^qf_{\ast}F)_{\xi}$ vanishes
for $q>2n$. If $F$ is a sheaf of $p$-torsion groups, where $p$ is
the characteristic of $\xi$, then the fiber of $R^qf_{\ast}F$ at
$\xi$ vanishes for $q>n$.
\end{enonce*}

Recall that the \emph{relative dimension of $X/S$ is} $\leq n$ if
every fiber has dimension $\leq n$. We therefore have:

\begin{enonce*}{Corollary 5.3 bis}\etiquette[5.3 bis]{XII.corbis5.3}
Let $f:X\longrightarrow S$ be proper, with $X/S$ of relative
dimension $\leq n$. Then $R^qf_{\ast}F=0$ for $q>2n$ and every
torsion abelian sheaf $F$ on $X$. If $S$ has characteristic $p$ and
$F$ is a sheaf of $p$-torsion abelian groups, then
$R^qf_{\ast}F=0$ for $q>n$.
\end{enonce*}

Applying \hyperref[XII.thm5.1]{5.1} to a morphism between spectra of
separably closed fields gives:

\begin{enonce*}{Corollary 5.4}\etiquette[5.4]{XII.cor5.4}
(Invariance of cohomology under extension of the base field in the
  proper case). Let $k\subset K$ be separably closed fields and let
  $X$ be a proper scheme over $\Spec k$. Put $X'=X_K$. Let $F$ be a
  sheaf on $X$ and let $F'$ be its inverse-image sheaf on $X'$. Then:
\begin{enumerate}
\renewcommand{\theenumi}{\roman{enumi}}
\renewcommand{\labelenumi}{(\theenumi)}
\item $H^0(X,F)\xrightarrow{\sim}H^0(X',F')$.

\item If\pageoriginale $F$ is a sheaf of groups (resp. of
  ind-finite groups), the morphism
$$
H^1(X,F)\longrightarrow H^1(X',F')
$$
is injective (resp. bijective).

\item If $F$ is a torsion abelian sheaf, then
$$
H^q(X,F)\xrightarrow{\sim}H^q(X',F')
$$
for every $q$.
\end{enumerate}
\end{enonce*}

A slightly different variant of \hyperref[XII.cor5.2]{5.2} is the
following:

\begin{enonce*}{Corollary 5.5}\etiquette[5.5]{XII.cor5.5}
Let $S$ be the spectrum of a henselian ring, and let
$S'=s_{\circ}$ be its closed point. Let $f:X\longrightarrow S$ be
proper. Put $X_{\circ}=X'$ and $F_{\circ}={g'}^{\ast}F$. Then:
\begin{enumerate}
\renewcommand{\theenumi}{\roman{enumi}}
\renewcommand{\labelenumi}{(\theenumi)}
\item $H^0(X,F)\xrightarrow{\sim}H^0(X_{\circ},F_{\circ})$

for every sheaf of sets $F$.

\item The morphism
$H^1(X,F)\longrightarrow H^1(X_{\circ},F_{\circ})$

is injective (resp. bijective) for every sheaf of groups (resp. of
ind-finite groups) $F$.

\item $H^q(X,F)\xrightarrow{\sim}H^q(X_{\circ},F_{\circ})$

for every torsion abelian sheaf $F$ and every $q$.
\end{enumerate}
\end{enonce*}

Indeed, this statement is just a special case of
\hyperref[XII.cor5.2]{5.2} when $S$ is strictly local. Suppose that
\hyperref[XII.cor5.2]{5.2} is known. Let $\bar S$ be the strict
localization of $S$, let $\bar s$ be the closed point of $\bar S$,
and let $G$ be the Galois group of $\bar S/S$ (which is also the
Galois group of $\bar s/s_{\circ}$). Then $G$ acts on the fibers of
the $R^q$, and, with the evident notation, one has
\begin{align*}
H^0(X,F)&\simeq H^0(\bar X,\bar F)^G,\\
H^0(X_{\circ},F_{\circ})&\simeq
H^0(\bar X_{\circ},\bar F_{\circ})^G
\end{align*}
for\pageoriginale a sheaf of sets $F$. Assertion (i) follows, since
$H^0(\bar X,\bar F)\xrightarrow{\sim}
H^0(\bar X_{\circ},\bar F_{\circ})$ and this isomorphism plainly
commutes with the action of $G$.

If $F$ is a sheaf of groups, one readily defines a morphism of exact
sequences of pointed sets
$$\vcenter{\etiquette[5.5.1]{XII.eq5.5.1}
\xymatrix{
0 \ar[r] & H^1(G,H^0(\bar X,\bar F)) \ar[r]\ar[d] &
H^1(X,F) \ar[r]\ar[d] & H^1(\bar X,\bar F)^G \ar[d]\\
0 \ar[r] & H^1(G,H^0(\bar X_{\circ},\bar F_{\circ})) \ar[r] &
H^1(X_{\circ},F_{\circ}) \ar[r] &
H^1(\bar X_{\circ},\bar F_{\circ})^G ,
}}\leqno{(5.5.1)}
$$
in which the left-hand vertical arrow is an isomorphism and the
right-hand one is injective by \hyperref[XII.cor5.2]{5.2}. Let
$\alpha,\beta\in H^1(X,F)$ be
classes having the same image in $H^1(X_{\circ},F_{\circ})$, and let
$F^{\alpha}$ be the sheaf obtained by twisting $F$ by $\alpha$
\cite{sga4.exp.XII.bib.1}. Let $\beta'$ be the class in $H^1(X,F^{\alpha})$
corresponding to $\beta$. Since the class $\alpha'$ corresponding to
$\alpha$ is neutral, the image of $\beta'$ in
$H^1(X_{\circ},(F^{\alpha})_{\circ})$ is neutral.

Consider the diagram deduced from \eqhyperref[XII.eq5.5.1]{5.5.1}
by replacing $F$ with $F^{\alpha}$. It shows that the image of
$\beta'$ in $H^1(\bar X,\bar F^{\alpha})^G$ is neutral, and hence
that $\beta'$ comes from an element of
$H^1(G,H^0(\bar X,\bar F^{\alpha}))$; this element must itself be
neutral. Thus $\beta'$ is neutral, that is,
$\beta'=\alpha'$, whence $\beta=\alpha$. This proves the injectivity
of the morphism in (ii).

Now suppose that $F$ is a sheaf of ind-finite groups, and let us
prove that the arrow in (ii) is surjective. It is enough to show that
the functor $F\mapsto H^1(X_{\circ},F_{\circ})$ is effaceable.
Indeed, let $0\longrightarrow F\xrightarrow{u}G$ be an injection
which effaces this functor, and put $C=G/F$. There is a morphism of
exact sequences (\hyperref[XII.prop4.4]{4.4})
$$
\xymatrix{
H^0(X,C) \ar[r]\ar[d] & H^1(X,F) \ar[d]_{\epsilon}\ar[r] &
H^1(X,G) \ar[d]\\
H^0(X_{\circ},C_{\circ}) \ar[r] &
H^1(X_{\circ},F_{\circ}) \ar[r]^{u_{\circ}^1} &
H^1(X_{\circ},G_{\circ})
}
$$
in which $u_{\circ}^1$ is zero; this indeed implies the surjectivity
of $\epsilon$.

Let us prove\pageoriginale that the functor
$F\mapsto H^1(X_{\circ},F_{\circ})$ is effaceable for sheaves $F$ of
ind-finite groups. Let $F\xrightarrow{u}C$ be the \lq\lq Godement
resolution'' \eqhyperref[XII.eq3.4]{3.4}. We have
$$
H^1(G,H^0(\bar X_{\circ},\bar C_{\circ}))\simeq
H^1(G,H^0(\bar X,\bar C))\subset H^1(X,C)=0{\quoi}.
$$

Moreover, there is a commutative square
$$
\xymatrix{
H^1(\bar X,\bar F) \ar@{-}[r]^{\sim}\ar[d] &
H^1(\bar X_{\circ},\bar F_{\circ})\ar[d]\\
0=H^1(\bar X,\bar C) \ar[r] &
H^1(\bar X_{\circ},\bar C_{\circ})
}
$$
(indeed,
$H^1(\bar X,\bar C)=\varinjlim H^1(X',C)=0$, where $X'/X$ runs
through the étale morphisms), so its right-hand vertical arrow is
zero. We therefore have a morphism of exact sequences
$$
\xymatrix{
0 \ar[r] & H^1(G,H^0(\bar X_{\circ},\bar F_{\circ}))
\ar[r]\ar[d] & H^1(X_{\circ},F_{\circ}) \ar[r]\ar[d]^{\beta} &
H^1(\bar X_{\circ},\bar F_{\circ})^G \ar[d]^{\alpha^G}\\
&0 \ar[r] & H^1(X_{\circ},C_{\circ}) \ar[r] &
H^1(\bar X_{\circ},\bar C_{\circ})^G
}
$$
in which $\alpha^G$ is zero. It follows that $\beta$ is also zero,
and hence that the functor is effaceable.

Finally, let us prove assertion (iii) of
\hyperref[XII.cor5.5]{5.5}. There is a morphism of spectral
sequences
$$
\xymatrix{
**[l]E_2^{p,q}=H^p(G,H^q(\bar X,\bar F))
\ar@{=>}[r]\ar[d] & H^n(X,F)\ar[d]^{\varphi}\\
H^p(G,H^q(\bar X_{\circ},\bar F_{\circ}))
\ar@{=>}[r] & H^n(X_{\circ},F_{\circ}).
}
$$
By \hyperref[XII.cor5.2]{5.2}, the morphisms
$H^q(\bar X,\bar F)\longrightarrow
H^q(\bar X_{\circ},\bar F_{\circ})$ are isomorphisms. Consequently
the morphism of $E_2^{p,q}$ terms is an isomorphism, and therefore so
is the morphism of abutments. This proves the result, \cqfd

\setcounter{subsection}{5}
\subsection{}\etiquette[5.6]{XII.subsec5.6}
Let us examine\pageoriginale \hyperref[XII.cor5.5]{5.5} (i) and (ii)
when $F$ is constant. Let $F=T_X$ be the constant sheaf of sets on
$X$ with value $T$. Then $F(X)$ is simply the set of locally constant
maps from $X$ to $T$. If $X_{\circ}$ is a subscheme of $X$ and $T$
has at least two elements, it follows that
$F(X)\longrightarrow F(X_{\circ})$ is bijective if and only if the
map
$$
U\longmapsto U_{\circ}=U\cap X_{\circ}
$$
from the set $\Theta(X)$ of subsets of $X$ that are both open and
closed to the analogous set $\Theta(X_{\circ})$ is bijective. If the
set $\Pi_{\circ}(X_{\circ})$ of connected components of $X_{\circ}$
is finite, this is equivalent to saying that the map
$\Pi_{\circ}(X_{\circ})\longrightarrow\Pi_{\circ}(X)$ is bijective.
Thus \hyperref[XII.cor5.5]{5.5} (i) may also be restated as follows
(bearing in mind that in this case $X_{\circ}$ is of finite type over
a field, hence Noetherian, so that $\Pi_{\circ}(X_{\circ})$ is
finite):

\begin{enonce*}{Corollary 5.7}\etiquette[5.7]{XII.cor5.7}
Let $S$ be the spectrum of a henselian ring, let $s_{\circ}$ be its
closed point, let $f:X\longrightarrow S$ be a proper morphism, and
let $X_{\circ}$ be the closed fiber of $X/S$. Then the map
$$
\Pi_{\circ}(X_{\circ})\longrightarrow\Pi_{\circ}(X)
$$
induced by $X_{\circ}\longrightarrow X$ is bijective.
\end{enonce*}

(In particular, this implies that $\Pi_{\circ}(X)$ is finite.) We
shall prove \hyperref[XII.cor5.7]{5.7} directly in the Noetherian
case:

\begin{enonce*}{Lemma 5.8}\etiquette[5.8]{XII.lem5.8}
Statement \hyperref[XII.cor5.7]{5.7} is true if $S$ is Noetherian
\footnote[1]{\Cf \EGA IV 18.5.19.}.
\end{enonce*}

\emph{Proof.} In the Noetherian case, $\Pi_{\circ}(X)$ is certainly
finite. Replacing $X$ by a connected component, we may suppose that
$\Pi_{\circ}(X)$ has one element, that is, that $X$ is connected and
nonempty. We must show that $X_{\circ}$ is also connected and
nonempty. The image of $X$ in $S$ is closed because $f$ is proper,
and is nonempty because $X$ is nonempty; it therefore contains
$s_{\circ}$, which shows that $X_{\circ}$ is nonempty. Let
$$
\xymatrix@R=1cm@C=1.5cm{
X \ar[rd]^{f'} \ar[dd]_{f}& \\
  & S'\ar[ld]^{\pi}\\
S &
}
$$
be the\pageoriginale Stein factorization of $f$ (\EGA III 4.3.3), so
that $S'/S$ is finite and $S'$ is connected. Every fiber of $f'$ is
known to be connected and nonempty, and the underlying space of
$X_{\circ}$ is the union of the underlying spaces of the fibers over
the closed points of $S'$. It remains only to show that $S'$ is the
spectrum of a \emph{local} ring. Since $S$ is henselian, the ring
$A'$ of $S'$ is a product of local rings (VII \hyperref[VII.prop4.1]{4.1} (i)); because $S'$
is connected, $A'$ is local. This proves \hyperref[XII.lem5.8]{5.8}.

\setcounter{subsection}{9}
\setcounter{subsubsection}{-1}
\subsubsection{}\etiquette[5.9.0]{XII.subsubsec5.9.0}
Suppose now that $F=G_X$ is a constant sheaf of groups with value a
finite group $G$. Then $H^1(X,G_X)$ classifies the principal Galois
coverings of $X$ with group $G$ (VII \hyperref[VII.rem2.1]{2.1}). Corollary
\hyperref[XII.cor5.5]{5.5} (ii) says that the morphism
$X_{\circ}\longrightarrow X$ induces a bijection between the sets of
isomorphism classes of these coverings. By the definition of the
$\Pi_1$, it follows that \hyperref[XII.cor5.5]{5.5} (ii) for constant
sheaves of finite groups (in the presence of (i)) is equivalent to:

\begin{enonce*}{Theorem 5.9}\etiquette[5.9]{XII.thm5.9}
(Specialization theorem for the fundamental group). With the notation
  of \hyperref[XII.cor5.5]{5.5}, suppose that $X$ is connected and
  nonempty. Then $X_{\circ}$ is connected and nonempty; and, if
  $\bar x_{\circ}$ is a geometric point of $X_{\circ}$, the morphism
  of fundamental groups
$$
\Pi_1(X_{\circ},\bar x_{\circ})\longrightarrow
\Pi_1(X,\bar x_{\circ})
$$
is an isomorphism.
\end{enonce*}

On the other\pageoriginale hand, by Galois theory (\SGA V), statement
\hyperref[XII.thm5.9]{5.9} (in the case where $X_{\circ}$ is
connected, as we may suppose by \hyperref[XII.cor5.7]{5.7}) is
equivalent to:

\begin{enonce*}{Theorem 5.9 bis}\etiquette[5.9 bis]{XII.thmbis5.9}
Let $S$ be the spectrum of a henselian ring, let $s_{\circ}$ be the
closed point of $S$, let $f:X\longrightarrow S$ be a proper
morphism, and let $X_{\circ}$ be the closed fiber of $X/S$. For a
scheme $Z$, denote by $\Et(Z)$ the category of étale coverings of
$Z$. Then the restriction functor
$$
\Et(X)\longrightarrow\Et(X_{\circ})
$$
is an equivalence of categories.
\end{enonce*}

\begin{enonce*}[remark]{Remark 5.10}\etiquette[5.10]{XII.rem5.10}
Recall that \hyperref[XII.thm5.9]{5.9} was already proved
(\SGA 1 X 2.1) when $S$ is the spectrum of a complete Noetherian
local ring. Unfortunately, that result does not seem to yield a
proof of \hyperref[XII.thm5.9]{5.9} in general
\footnote[1]{The situation has changed since these lines were
  written; \cf Th.~\hyperref[XII.prop3.1]{3.1} in M. Artin,
  Algebraic approximation of structures over complete local rings
  (to appear).}. The proof of \hyperref[XII.thm5.9]{5.9} given in this exposé and the next is
very different from the one given in \SGA 1 X.
\end{enonce*}

\setcounter{subsection}{10}
\subsection{}\etiquette[5.11]{XII.subsec5.11}
One can give another \lq\lq noncommutative'' variant of
\hyperref[XII.cor5.5]{5.5} (ii), using J. \textsc{Giraud}'s theory of
2-cohomology (\cf J. \textsc{Giraud}'s thesis, in preparation).
Suppose that $S$ is Noetherian. Let $L$ be a \emph{band} on $X$, and
suppose that $L$ is represented locally (for the étale topology) by
a sheaf of ind-finite groups. Consider the map
$$
\varphi^{(2)}:H^2(X,L)\longrightarrow H^2(X_{\circ},L_{\circ}){\quoi},
$$
where $L_{\circ}$ is the restriction of $L$ to $X_{\circ}$. Then
\begin{equation*}
(\varphi^{(2)})^{-1}(H^2(X_{\circ},L_{\circ})')
=H^2(X,L)'{\quoi},\tag{5.11.1}
\end{equation*}\etiquette[5.11.1]{XII.eq5.11.1}
where\pageoriginale the \lq\lq primes'' denote the subsets formed by
the neutral elements.

In view of \hyperref[XII.cor5.5]{5.5} (iii), this statement is also
equivalent to the following:

\emph{For every gerbe $\mathcal G$ on $X$, with band $L$ as above,
  let $\mathcal G_{\circ}$ denote its restriction to $X_{\circ}$.
  Then the restriction functor on the categories of sections}
$$
\mathcal G(X)\longrightarrow\mathcal G_{\circ}(X_{\circ})
$$
\emph{is an equivalence of categories.}

Indeed, the fact that this functor is fully faithful follows from
\hyperref[XII.cor5.5]{5.5} (i), and the fact that it is essentially
surjective when $\mathcal G(X)\neq\varnothing$, that is, when the
class $\xi\in H^2(X,\sisi{\mathcal G}{L})$ of $\mathcal G$ is
neutral, follows from \hyperref[XII.cor5.5]{5.5} (ii). It remains to
express the fact that
$\mathcal G_{\circ}(X_{\circ})\neq\varnothing$ implies
$\mathcal G(X)\neq\varnothing$, and this is precisely
\eqhyperref[XII.eq5.11.1]{5.11.1}.

We shall not give the proof of \eqhyperref[XII.eq5.11.1]{5.11.1}
here. We merely point out that, for every Noetherian scheme $X$ and
every closed subscheme of $X$, one can in fact show that the validity of
statements \hyperref[XII.cor5.5]{5.5} (i), (ii) implies that of
\eqhyperref[XII.eq5.11.1]{5.11.1}. The proof is fairly elementary
and uses the descent result VIII \hyperref[VIII.thm9.4]{9.4} a).

% END INCLUDED SGA4_English_current_buildable_source_20260906/source/12/drafts/SGA4_XII_LINES_383_784_English.texfrag


% BEGIN INCLUDED SGA4_English_current_buildable_source_20260906/source/12/drafts/SGA4_XII_LINES_785_1200_English.texfrag
\section{First reductions}\etiquette[6]{XII.sec6}

We use the abbreviation $(f,F,g)$ for the data of
\hyperref[XII.thm5.1]{5.1} in any one of cases (i), (ii), or (iii) of
that theorem.

\begin{enonce*}{Lemma 6.1}\etiquette[6.1]{XII.lem6.1}
For \hyperref[XII.thm5.1]{5.1} to hold for the data $(f,F,g)$, where
$f$ and $F$ are fixed and $g$ is arbitrary, it is enough that the
theorem hold for $(\bar f,\bar F,\bar g)$ for every $S$-scheme
$\bar S$ that is the strict localization of a scheme $S'$ locally\pageoriginale of
finite type over $S$ at a point closed in its fiber. Here $\bar f$
and $\bar F$ are obtained from $f$ and $F$ by the base change
$\bar S\to S$, and $\bar g:\bar s\to\bar S$ is the inclusion of the
closed point of $\bar S$.
\end{enonce*}

\noindent
\emph{Proof.} We may plainly assume that $S$ and $S'$ are affine,
say the spectra of rings $A$ and $A'$. Then $A'$ is the filtered
direct limit of its finite-type $A$-subalgebras $A'_i$. Using the
theory of passage to the limit (VII \hyperref[VII.cor5.11]{5.11} and \hyperref[VII.rem5.14]{5.14}), we are reduced to
proving \hyperref[XII.thm5.1]{5.1} for the triples $(f,F,g_i)$, where
$g_i:\Spec(A'_i)=S'_i\to S$; this reduces us to the case in which
$g$ is of finite type. To prove that the base-change homomorphisms
are isomorphisms, it is enough to prove that they induce isomorphisms
on the fibers of the sheaves in question, and it is enough to
consider geometric fibers at points $s'$ of $S'$ that are closed in
their fibers over $S$ (VIII \hyperref[VIII.rem3.13]{3.13}). Let $\bar s'$ be a geometric
point of $S'$ corresponding to a separable closure of $k(s')$.
Since $s'$ is algebraic over the corresponding point $s$ of $S$,
$k(\bar s')$ is also a separable closure of $k(s)$. Denote by
$\bar s$ the prescheme $\bar s'$, regarded as lying over $S$, that
is, as a geometric point of $S$, and consider the strict
localizations $\bar S$ and $\bar S'$ of $S$ and $S'$ at $\bar s$ and
$\bar s'$, respectively. We then have a commutative square
$$
\vcenter{
\xymatrix@C=3cm@R=1.25cm{
\bar{s}'  \ar[r] \ar[d]& \bar{s}\ar[d]\\
\bar{S}'  \ar[r] &  \bar{S},
}}
$$
in which the first horizontal arrow is an isomorphism. In view of
VIII \hyperref[VIII.thm5.2]{5.2} and \hyperref[VIII.rem5.3]{5.3}, the homomorphism induced on the fibers at
$\bar s'$ by the base-change homomorphism associated with $(f,F,g)$
identifies with the homomorphism
$$
H^{i}(\bar{X},\bar{F})\longrightarrow H^{i}(\bar{X}',\bar{F}')
$$
induced\pageoriginale by $\bar X'\to\bar X$, where $\bar X$ and
$\bar X'$ denote the schemes obtained from $X$ by the base changes
$\bar S\to S$ and $\bar S'\to S$, respectively. The commutative
square above plainly gives a commutative square
$$
\xymatrix@C=3cm@R=1.5cm{
H^{i}(\bar{X}_{\bar{s}}, \bar{F}) \ar[r]&
H^{i}(\bar{X}'_{\bar{s}'},\bar{F}')\\
H^{i}(\bar{X},\bar{F})  \ar[r] \ar[u]&
H^{i}(\bar{X}',\bar{F}') \ar[u],
}
$$
in which the first horizontal arrow is again an isomorphism.
Consequently the second horizontal arrow is an isomorphism as soon
as the two vertical arrows are. This is precisely what the hypothesis
of \hyperref[XII.lem6.1]{6.1} guarantees, \cqfd

\setcounter{subsection}{1}
\subsubsection{}\etiquette[6.1.1]{XII.subsubsec6.1.1}
Thus, to prove \hyperref[XII.thm5.1]{5.1} in a situation $(F,f,-)$,
we are reduced to studying the special case
\hyperref[XII.cor5.5]{5.5}, with $S$ strictly local. This leads us to
study in general a morphism $h:Y\to X$ of schemes and to seek general
conditions under which the corresponding homomorphisms
$$
H^{i}(X,F)\longrightarrow H^{i}(Y,h^{\ast}(F))
$$
are bijective. This subsection gives some easy auxiliary results on
this situation. For convenience in later references, we state them
in greater generality and precision than is needed for the proof of
\hyperref[XII.thm5.1]{5.1}.

\begin{enonce*}{Lemma 6.2}\etiquette[6.2]{XII.lem6.2}
Let $C$ and $C'$ be two abelian categories, let
$T^{\bigdot}=(T^{i})$ and $T'^{\bigdot}=({T'}^{i})$ be two
cohomological functors from $C$ to $C'$, and let
$\varphi=(\varphi^{i})$ be a morphism of\pageoriginale cohomological functors from
$T^{\bigdot}$ to $T'^{\bigdot}$. Assume that $T^{\bigdot}$ and
$T'^{\bigdot}$ vanish in degrees $<0$, and that $T^{i}$ is
effaceable for $i>0$. Let $n$ be an integer. The following conditions
are equivalent:
\begin{enumerate}
\renewcommand{\theenumi}{\alph{enumi}}
\renewcommand{\labelenumi}{\theenumi)}
\item $\varphi^{i}$ is an isomorphism for $i\leq n$ and a
  monomorphism for $i=n+1$.

\item (If $n\geq -1$.) The homomorphism $\varphi^{i}$ is a
  monomorphism for $i=0$ and an epimorphism for $i\leq n$.

\item (If $n\geq 0$.) The homomorphism $\varphi^{\circ}$ is an
  isomorphism, and the ${T'}^{i}$ are effaceable for $0<i\leq n$.
\end{enumerate}

If these conditions hold and $A\in\Ob(C)$, then
$\varphi^{n+1}(A):T^{n+1}(A)\to {T'}^{n+1}(A)$ is an epimorphism
(hence an isomorphism) if and only if ${T'}^{n+1}(A)$ is effaceable,
that is, if there exists a monomorphism $A\to B$ such that the
corresponding homomorphism
${T'}^{n+1}(A)\to {T'}^{n+1}(B)$ is zero. If $C'=(\Ab)$, the category
of abelian groups, and $\xi\in {T'}^{n+1}(A)$, then $\xi$ lies in the
image of $T^{n+1}(A)$ if and only if $\xi$ is effaceable, that is, if
there exists a monomorphism $A\to B$ such that the image of $\xi$ in
${T'}^{n+1}(B)$ is zero.
\end{enonce*}

\noindent
\emph{Proof.} It is clear that (a) implies (b) and (c). Since the
hypotheses on $T^{\bigdot}$ imply that the $T^{i}$ for $i>0$ are the
right satellites of $T^{\circ}$, the implication (c) $\Rightarrow$
(a) follows from the axiomatic characterization of satellite
functors. It also follows, by induction on $n$, from the last
assertion of \hyperref[XII.lem6.2]{6.2}, proved by the familiar
argument that we do not repeat here (compare the proof of
\hyperref[XII.cor6.6]{6.6} below). The implication (b)
$\Rightarrow$ (a) is immediate if $n=-1$. If $n\geq0$, we proceed by
induction on $n$: the induction hypothesis allows us to suppose that
$\varphi^{i}$ is an isomorphism for $i\leq n-1$ and\pageoriginale a monomorphism
for $i=n$. The hypothesis then makes it an isomorphism also for
$i=n$, and it remains to show, using the fact that $\varphi^n$ is an
isomorphism, that it is a monomorphism for $i=n+1$. Let $A\to B$ be
a monomorphism that effaces $T^{n+1}$, giving an exact sequence
$0\to A\to B\to C\to0$ in $C$ and a morphism of exact sequences in
$C'$:
$$
\xymatrix@R=1.25cm{
T^{n}(B)  \ar[r] \ar[d]& T^{n}(C)  \ar[r]\ar[d] & T^{n+1}(A) \ar[d] \ar[r] & T^{n+1}(B)\ar[d]\\
T'^{n}(B)  \ar[r] & T'^{n}(C) \ar[r] & T'^{n+1}(A)  \ar[r] &
T'^{n+1}(B),
}
$$
where the first two vertical arrows are isomorphisms by what we have
just proved, and $T^{n+1}(A)\to T^{n+1}(B)$ is zero by construction.
The five lemma then implies that
$T^{n+1}(A)\to {T'}^{n+1}(A)$ is a monomorphism, \cqfd

\begin{enonce*}{Corollary 6.3}\etiquette[6.3]{XII.cor6.3}
Suppose that $T^{\bigdot}$ and $T'^{\bigdot}$ arise from
cohomological functors (denoted by the same symbols) on the right
derived category $D^{+}(C)$ \cite{sga4.exp.XII.bib.3}, with values in $C'$; the
$T^{i}$ for $i>0$ are not necessarily assumed effaceable. Suppose
that, for $K^{\bigdot}\in D^{+}(C)$, the condition
$H^{i}(K^{\bigdot})=0$ for $i<0$ implies
$T^{i}(K^{\bigdot})={T'}^{i}(K^{\bigdot})=0$ for $i<0$. Let
$K^{\bigdot}\in\Ob D^{+}(C)$ be a bounded-below complex in $C$ such
that $H^{i}(K^{\bigdot})=0$ for $i<0$.
\begin{enumerate}
\renewcommand{\theenumi}{\roman{enumi}}
\renewcommand{\labelenumi}{(\theenumi)}
\item If condition {\rm (a)} of \hyperref[XII.lem6.2]{6.2} is
  satisfied, then the homomorphism
  $\varphi^{i}(K^{\bigdot})=T^{i}(K^{\bigdot})\to
  {T'}^{i}(K^{\bigdot})$ is an isomorphism for $i\leq n$ and a
  monomorphism for $i=n+1$.

\item Let, in addition, $L^{\bigdot}\in\Ob D^{+}(C)$ satisfy
  $H^{i}(L^{\bigdot})=0$ for $i<0$, and let
  $u:L^{\bigdot}\to K^{\bigdot}$ be a morphism in $D^{+}(C)$ that
  induces an injective homomorphism
  $H^{\circ}(u):H^{\circ}(L^{\bigdot})\to
  H^{\circ}(K^{\bigdot})$. Then the homomorphism
  $T^{n+1}(L^{\bigdot})\to {T'}^{n+1}(L^{\bigdot})$ is surjective\pageoriginale if
  and only if the image of ${T'}^{n+1}(L^{\bigdot})$ in
  ${T'}^{n+1}(K^{\bigdot})$ is contained in that of
  $T^{n+1}(K^{\bigdot})$. If $C'=(\Ab)$ and
  $\xi\in {T'}^{n+1}(L^{\bigdot})$, then $\xi$ lies in the image of
  $T^{n+1}(L^{\bigdot})$ if and only if its image in
  ${T'}^{n+1}(K^{\bigdot})$ lies in that of
  $T^{n+1}(K^{\bigdot})$.
\end{enumerate}
\end{enonce*}

Recall that the hypothesis that $T^{\bigdot}$ and $T'^{\bigdot}$
arise from cohomological functors on $D^{+}(C)$ may be spelled out as
follows: they arise from cohomological functors on the abelian
category of bounded-below complexes of $C$ that commute with degree
shifts, send morphisms of complexes homotopic to zero to zero
morphisms in $C'$, and send every morphism
$u:K^{\bigdot}\to L^{\bigdot}$ of complexes that is a
quasi-isomorphism (that is, induces isomorphisms
$H^{i}(K^{\bigdot})\to H^{i}(L^{\bigdot})$) to an isomorphism in
$C'$. In practice, all the cohomological functors from $C$ to $C'$
that one encounters arise in this way.

\medskip
\noindent
\emph{Proof of} \hyperref[XII.cor6.3]{6.3}. We first prove (i). The
assertion is clear if $H^{i}(K^{\bigdot})=0$ for $i\leq n+1$, since
the degree properties imposed on $T^{\bigdot}$ and $T'^{\bigdot}$
then give
$T^{i}(K^{\bigdot})={T'}^{i}(K^{\bigdot})=0$ for $i\leq n+1$.
On the other hand, let $m\geq0$ be an integer such that
$H^{i}(K^{\bigdot})=0$ for $i<m$; such an integer exists, and for
example $m=0$ will do. We proceed by descending induction on $m$, the
assertion having been proved when $m=n+2$. Thus assume the assertion
proved for $m'>m$, and prove it for $m$. By the hypothesis on the
degrees of $K^{\bigdot}$, the complex $K^{\bigdot}$ is
quasi-isomorphic to a complex concentrated in degrees $\geq m$, so
we may suppose that it is itself concentrated in those degrees.
There follows an exact sequence of complexes
$$
0\longrightarrow H^m[-m]\longrightarrow K^{\bigdot}\longrightarrow
K'^{\bigdot}\longrightarrow 0{\quoi},
$$
where $H^m=H^m(K^{\bigdot})$, the notation $[-m]$ denotes a shift of
$-m$ in the degrees of a complex,\pageoriginale and $K'^{\bigdot}$ is defined as
the cokernel. The cohomology objects of $K'^{\bigdot}$ are therefore
those of $K^{\bigdot}$ in degrees $\geq m+1$ and are zero in degrees
$\leq m$. In particular, the induction hypothesis applies to
$K'^{\bigdot}$. The preceding exact sequence gives rise to a
morphism of exact sequences
$$
\xymatrix{
T^{i-1} (K'^{\;\bigdot}) \ar[d] \ar[r] & T^{i-m}(H^m)  \ar[r]\ar[d] &
T^{i}(K^{\bigdot}) \ar[r] \ar[d]& T^{i} (K'^{\;\bigdot})  \ar[r] \ar[d]&
T^{i+1-m}(H^m) \ar[d]\\
T'^{i-1}(K'^{\;\bigdot})  \ar[r] & T'^{i-m}(H^m)  \ar[r] &
T'^{i}(K^{\bigdot})  \ar[r] & T'^{i}(K'^{\;\bigdot})  \ar[r] & T'^{i+1-m}(H^m).
}
$$
Taking $i\leq n$ in this diagram and using the induction hypothesis
and $m\geq0$, we find that the vertical arrows surrounding the
middle vertical arrow are isomorphisms, hence in particular
epimorphisms, and that the rightmost vertical arrow is a
monomorphism. The five lemma therefore shows that the middle
vertical arrow is an epimorphism. Taking $i\leq n+1$, we similarly
find that the vertical arrows surrounding the middle one are
monomorphisms, while the leftmost vertical arrow is an isomorphism,
hence in particular an epimorphism. The five lemma then shows that
the middle vertical arrow is a monomorphism. This proves
\hyperref[XII.cor6.3]{6.3} (i).

To prove the second assertion, \hyperref[XII.cor6.3]{6.3} (ii), we
similarly use the exact sequence
$$
0\longrightarrow L^{\bigdot}\longrightarrow K^{\bigdot}\longrightarrow
K'^{\bigdot}\longrightarrow 0{\quoi},
$$
where $K'^{\bigdot}$ is defined as the mapping cylinder of
$L^{\bigdot}\to K^{\bigdot}$ and again has cohomology in
nonnegative degrees. The usual diagram chase in
$$
\xymatrix{
T^{n}(K^{\bigdot})  \ar[r] \ar[d]& T^{n}(K'^{\;\bigdot})  \ar[r] \ar[d]& T^{n+1}
(L^{\bigdot})  \ar[r]\ar[d] & {T^{n+1}(K^{\bigdot})}  \ar[r]
\ar[d]&{T^{n+1}(K'^{\;\bigdot})} \ar[d]\\
T'^{n}(K^{\bigdot}) \ar[r] & T'^{n}(K'^{\;\bigdot})  \ar[r] &
T'^{n+1}(L^{\bigdot}) \ar[r] & T'^{n+1}(K^{\bigdot}) \ar[r] & T'^{n+1}(K'^{\bigdot})
}
$$
gives\pageoriginale the last assertion of
\hyperref[XII.cor6.3]{6.3}, since, by the first part already proved,
the first two vertical arrows are isomorphisms and the last two are
monomorphisms. This completes the proof of
\hyperref[XII.cor6.3]{6.3}.

\begin{enonce*}{Lemma 6.4}\etiquette[6.4]{XII.lem6.4}
Theorem \hyperref[XII.thm5.1]{5.1} holds if $f$ is finite.
\end{enonce*}

Indeed, this follows at once from VIII \hyperref[VIII.prop5.5]{5.5} and \hyperref[VIII.rem5.8]{5.8}.

\begin{enonce*}{Proposition 6.5}\etiquette[6.5]{XII.prop6.5}
Let $h:Y\to X$ be a quasi-compact and quasi-separated morphism of
schemes. In {\rm (i)} below, $I$ denotes a fixed set having at least
two elements; in {\rm (ii)} and {\rm (iii)}, $\LL$ denotes a
nonempty subset of the set $\PP$ of prime numbers; and in {\rm (iii)},
$n$ denotes an integer.
\begin{enumerate}
\renewcommand{\theenumi}{\roman{enumi}}
\renewcommand{\labelenumi}{(\theenumi)}
\item The following conditions are equivalent:
\begin{itemize}
\item[a)] For every sheaf of sets $F$ on $X$, the canonical map
$$
H^{\circ}(X,F)\longrightarrow H^{\circ}(Y,h^{\ast}(F))
$$
is injective (respectively, bijective).

\item[b)] For every finite morphism $X'\to X$, let
$h':Y'\to X'$ be the morphism obtained from $h$ by the base change
$X'\to X$. The canonical map
$$
H^{\circ}(X', I_{X'})\longrightarrow H^{\circ}(Y',I_{Y'})
$$
induced by $h'$ is injective (respectively, bijective).

\item[${\rm b')}$] With the notation of {\rm (b)}, the map
$U\mapsto {h'}^{-1}(U)$ from the set of subsets of $X'$ that are both
open and closed to the set of subsets of $Y'$ that are both open and
closed is injective (respectively, bijective).

\item[${\rm b'')}$] (If\pageoriginale $X$ is locally Noetherian.)
With the notation of {\rm (b)}, the induced map on sets of connected
components
$$
\Pi_{\circ}(h'):\Pi_{\circ}(Y')\longrightarrow \Pi_{\circ}(X')
$$
is surjective (respectively, bijective). Equivalently, if $X'$ is
nonempty (respectively, connected and nonempty), then the same is true
of $Y'$.
\end{itemize}

Moreover, if these conditions hold, then for every sheaf of groups
$F$ on $X$, the functor $P\mapsto h^{\ast}(P)$ from the category of
torsors (homogeneous principal bundles) under $F$ to the category of
torsors under $h^{\ast}(F)$ is faithful (respectively, fully
faithful). In particular, in the latter case, the canonical map
$$
H^1(X,F)\longrightarrow H^1(Y,h^{\ast}(F))
$$
is injective. Finally, under the same hypotheses, the functor
$X'\mapsto Y'=\fprod{X'}{Y}{X}$ from the category of étale coverings
of $X$ to the category of étale coverings of $Y$ is faithful
(respectively, fully faithful).

\item The following conditions are equivalent:
\begin{itemize}
\item[a)] For every sheaf $F$ of ind-$\LL$-finite groups on $X$, the
canonical map
$$
H^i(X,F)\longrightarrow H^i(Y,h^{\ast}(F))
$$
is bijective for $i=0,1$.

\item[b)] For every finite morphism $X'\to X$, let
$h':Y'\to X'$ be the morphism obtained from $h$ by base change, and
let $G$ be any ordinary finite $\LL$-group. The canonical map
$$
H^i(X',G_{X'})\longrightarrow H^i(Y',G_{Y'})
$$
is bijective for $i=0,1$.

\item[${\rm b')}$] (If\pageoriginale $\LL=\PP$.) With the notation
of {\rm (b)}, the inverse-image functor under $h'$ induces an
equivalence from the category of étale coverings of $X'$ to the
category of étale coverings of $Y'$.

\item[${\rm b'')}$] (If $X$ is Noetherian and $\LL=\PP$.) With the
notation of {\rm (b)}, if $X'$ is nonempty then so is $Y'$. If $y'$ is
a geometric point of $Y'$ and $x'$ is its image in $X'$, the canonical
map
$$
\Pi_i(Y',y')\longrightarrow \Pi_i(X',x')\quav\mbox{ for }i=0,1{\quoi},
$$
is bijective.
\end{itemize}

\item The following conditions are equivalent:
\begin{itemize}
\item[a)] For every sheaf $F$ of $\LL$-torsion groups on $X$, the
homomorphism
$$
H^{i}(X,F)\longrightarrow H^{i}(Y,h^{\ast}(F))
$$
is an isomorphism for $i\leq n$ and a monomorphism for $i=n+1$.

\item[${\rm a')}$] (If $n\geq -1$.) For every $F$ as above, the
homomorphism in question is injective for $i=0$ and surjective for
$i\leq n$.

\item[b)] (If $n\geq -1$.) For every finite morphism $X'\to X$, let
$h':Y'\to X'$ be the morphism obtained from $h$ by the base change
$X'\to X$. Then for every $\ell\in\LL$ and $\nu>0$, the canonical
homomorphism
$$
H^{i}(X',(\ZZ/\ell^{\nu}\ZZ)_{X'})\longrightarrow
H^{i}(Y',(\ZZ/\ell^{\nu}\ZZ)_{Y'})
$$
is injective for $i=0$ and surjective for $i\leq n$.
\end{itemize}
\end{enumerate}
\end{enonce*}

In addition, we have the following supplement to the preceding
statement:

\begin{enonce*}{Corollary 6.6}\etiquette[6.6]{XII.cor6.6}
\begin{enumerate}
\renewcommand{\theenumi}{\roman{enumi}}
\renewcommand{\labelenumi}{(\theenumi)}
\item Suppose that the non-parenthesized condition in
\hyperref[XII.prop6.5]{6.5} {\rm (i) (a)} holds. Let $F$ be a sheaf of
sets on $X$, and let $\xi\in H^{\circ}(Y,h^{\ast}(F))$. Suppose\pageoriginale that
there exists a monomorphism of sheaves of sets $F\to G$ such that the
image of $\xi$ in $H^{\circ}(Y,h^{\ast}(G))$ lies in the image of
$H^{\circ}(X,G)$. Then $\xi$ lies in the image of
$H^{\circ}(X,F)$.

\item Suppose that the parenthesized condition in
\hyperref[XII.prop6.5]{6.5} {\rm (i) (a)} holds. Let $F$ be a sheaf of
groups on $X$, and let $\xi\in H^1(Y,h^{\ast}(F))$. Suppose that there
exists a monomorphism of sheaves of groups $F\to G$ such that the
image of $\xi$ in $H^1(Y,h^{\ast}(G))$ lies in the image of
$H^1(X,G)$. Then $\xi$ lies in the image of $H^1(X,F)$.

\item Suppose that condition
\hyperref[XII.prop6.5]{6.5} {\rm (iii) (a)} holds for a fixed value of
$n$. Let $F$ be a sheaf of $\LL$-torsion groups on $X$, and let
$\xi\in H^{n+1}(Y,h^{\ast}(F))$. Suppose that there exists a
monomorphism $F\to G$ of sheaves of $\LL$-torsion groups such that
the image of $\xi$ in $H^{n+1}(Y,h^{\ast}(G))$ lies in the image of
$H^{n+1}(X,G)$. Then $\xi$ lies in the image of $H^{n+1}(X,F)$.
\end{enumerate}
\end{enonce*}

\noindent
\emph{Proof of \hyperref[XII.prop6.5]{6.5} and
\hyperref[XII.cor6.6]{6.6}.}
We first prove \hyperref[XII.cor6.6]{6.6}. The case

% END INCLUDED SGA4_English_current_buildable_source_20260906/source/12/drafts/SGA4_XII_LINES_785_1200_English.texfrag


% BEGIN INCLUDED SGA4_English_current_buildable_source_20260906/source/12/drafts/SGA4_XII_LINES_1201_1600_English.texfrag
of \hyperref[XII.cor6.6]{6.6} (iii) is a special case of the last assertion of
\hyperref[XII.cor6.3]{6.3}, applied when $C$ is the category of abelian
$\LL$-torsion sheaves on $X$, $C'$ is the category $(\Ab)$, and
$T^{i}(F)=H^{i}(X,F)$, ${T'}^{i}(F)=H^{i}(Y,h^{\ast}(F))$ (note that
$T'^{\bigdot}$ is indeed a cohomological functor in $F$, since the functor
$h^{\ast}$ is exact). Of course, the use of complexes as in
\hyperref[XII.cor6.3]{6.3} is unnecessary here (it will be convenient later,
in \hyperref[XII.prop6.8]{6.8}); we use only the last assertion of
\hyperref[XII.cor6.3]{6.3}.

To prove (i), consider the amalgamated sum $H=G\amalg_F G$, the colimit in
the category of sheaves of sets on $X$ of the diagram
$$
\xymatrix@R=.5cm{
  & G\\
F \ar[ru] \ar[rd]&\\
  & G
}
$$
The\pageoriginale usual exactness properties of the topos of sheaves of sets
show, since $F\rightarrow G$ is a monomorphism, that there is an
\emph{exact} diagram of sheaves of sets on $X$
$$
F\to G\rightrightarrows H
$$
(this may be expressed by saying that every monomorphism in a topos is
effective, and may be proved by reducing, as usual, to the category of
sets). Since the functors $H^{0}$ and $h^{\ast}$ are left exact, we obtain a
morphism of exact diagrams of sets
$$
\xymatrix@R=1.25cm{
H^{0}(X,F) \ar[r]\ar[d]& H^{0}(X,G)  \ar[d]\ar@<2pt>[r]\ar@<-2pt>[r]& H^{0}(X,H)\ar[d]\\
H^{0}(Y,h^{\ast}(F)) \ar[r]& H^{0}(Y,h^{\ast}(G)) \ar@<2pt>[r]\ar@<-2pt>[r]& H^{0}(Y,h^{\ast}(H)),
}
$$
in which the vertical arrows are injective. An immediate diagram chase then
proves assertion \hyperref[XII.cor6.6]{6.6} (i).

To prove \hyperref[XII.cor6.6]{6.6} (ii), we first prove the second assertion
of \hyperref[XII.prop6.5]{6.5} (i), namely that, under condition a) of
\hyperref[XII.prop6.5]{6.5} (i), for every sheaf of groups $F$ on $X$, the
functor $P\mapsto h^{\ast}(P)$ from the category of torsors under $F$ to the
category of torsors under $h^{\ast}(F)$ is faithful (\resp fully faithful),
and hence, in the respective case, induces an injection on $H^1$. For this,
if $P$ and $P'$ are two torsors under $F$, denote by
$\sheaf{\Isom}_F(P,P')$ the sheaf of $F$-isomorphisms from $P$ to $P'$.
It is immediate that the formation of this sheaf commutes with every base
extension, and in particular with base extension by $h$. (This assertion is
valid for every morphism of sites.) Applying the hypothesis on the effect of
$h^{\ast}$ on $H^{0}$ to the sheaf $\sheaf{\Isom}_F(P,P')$, we find that the
map
$$
\Isom_F(P,P')\longrightarrow
\Isom_F(h^{\ast}(P),h^{\ast}(P'))
$$
is injective (\resp bijective), which says that the functor
$P\mapsto h^{\ast}(P)$ is\pageoriginale faithful (\resp fully faithful).
(Note that every morphism in the category of torsors under a sheaf of groups
is an isomorphism.)

On the other hand, note that if $F$ is a subsheaf of groups of a sheaf of
groups $G$, then giving a torsor $P$ under $F$ amounts to giving a torsor
$Q$ under $G$ (namely the one obtained from $P$ by extending the structure
group along $F\rightarrow G$), \emph{together with} a section of $Q/F$
(which is indeed interpreted as a reduction of the structure group from
$G$ to $F$). With the notation of \hyperref[XII.cor6.6]{6.6} (ii), the image
of $\xi$ in $H^1(Y,h^{\ast}(G))$ comes from an element of $H^1(X,G)$; the
latter is represented by a torsor $Q$ under $G$, so that the image of $\xi$
is the class of the torsor $h^{\ast}(Q)$ under $h^{\ast}(G)$. The fact that
this class comes from $\xi\in H^1(Y,h^{\ast}(F))$ means explicitly that
$\xi$ is the class of the $h^{\ast}(F)$-torsor defined by $h^{\ast}(Q)$ and
a suitable section of $h^{\ast}(Q)/h^{\ast}(F)$. Since the functor
$h^{\ast}$ is right exact, this last sheaf is none other than
$h^{\ast}(Q/F)$; and since $h^{\ast}$ induces a bijection on the $H^{0}$ of
sheaves of sets, the section in question of $h^{\ast}(Q)/h^{\ast}(F)$ comes
from a section of $Q/F$. This last section defines an $F$-torsor by reduction
of the structure group, and the class of this torsor in $H^1(X,F)$ plainly
has $\xi$ as its image in $H^1(Y,h^{\ast}(F))$. This proves
\hyperref[XII.cor6.6]{6.6} (ii) and completes the proof of
\hyperref[XII.cor6.6]{6.6}.

We now prove \hyperref[XII.prop6.5]{6.5}, beginning with
\hyperref[XII.prop6.5]{6.5} (iii). Clearly a) implies ${\rm a'}$), and the
converse implication is a special case of \hyperref[XII.lem6.2]{6.2}
(again applied when $C$ is the category of $\LL$-torsion sheaves on $X$ and
$C'$ the category of abelian groups), provided that one proves that the
functors $H^{i}(X,F)$, for $i>0$, are effaceable in $C$. This presents no
difficulty; alternatively, we can avoid proving that result by\pageoriginale
noting that $F$ is the filtered colimit of its subsheaves $F_i$ for which
there exists an integer $n>0$ annihilating $F_i$
(\hyperrefext[IX][IX.def1.1]{1.1}). Then $h^{\ast}(F)$ is the colimit of the
$h^{\ast}(F_i)$, and, in view of the commutation of cohomology with filtered
colimits (VII \hyperref[VII.prop3.3]{3.3}), which follows from the quasi-compactness and
quasi-separatedness hypotheses on $X$ and $Y$ (hypotheses not used for
\hyperref[XII.cor6.6]{6.6}), we are reduced to proving condition a) for the
$F_i$. Thus we are reduced to the case in which there exists an integer
$m>0$ such that $mF=0$. We may plainly suppose that every prime divisor of
$m$ belongs to $\LL$, and can then take for $C$ the category of sheaves of
$\ZZ/m\ZZ$-modules on $X$. In this category it is known that the functors
$H^{i}(X,-)$ are effaceable for $i\geq1$. This proves the equivalence of
conditions a) and ${\rm a'}$) in \hyperref[XII.prop6.5]{6.5} (iii).

On the other hand, apply a) to the direct image of a $\LL$-torsion sheaf
$F'$ on a scheme $X'$ finite over $X$; this direct image is a
$\LL$-torsion sheaf by \hyperrefext[IX][IX.prop1.2]{1.2} (v) and
\hyperrefext[IX][IX.prop1.6]{1.6} (iii). Using
\hyperref[XII.lem6.4]{6.4} and VIII \hyperref[VIII.prop5.5]{5.5}, we find that the homomorphism
$$
H^{i}(X',F')\longrightarrow H^{i}(Y',{h'}^{\ast}(F'))
$$
is an isomorphism for $i\leq n$ and a monomorphism for $i=n+1$, whence, a
fortiori, condition b). It remains to prove that this condition implies
${\rm a'}$). For this, use the fact (\hyperrefext[IX][IX.cor2.7.2]{2.7.2})
that every $\LL$-torsion sheaf $F$ on $X$ is the filtered colimit of
constructible $\LL$-torsion sheaves. The preceding colimit argument then
shows that, in ${\rm a'}$), it suffices to consider constructible $F$. By
\hyperrefext[IX][IX.prop2.14]{2.14}, there then exist finitely many finite
morphisms $p_i:X'_i\rightarrow X$, constant $\LL$-torsion sheaves
$G_{X'_i}$ on the $X'_i$, and an injective homomorphism
$$
F\longrightarrow \prod_{i}p_{i\ast}(G_{X'_i}).
$$
Now\pageoriginale if $n=-1$, that is, if the desired conclusion reduces to
the assertion that
$$
H^{0}(X,F)\longrightarrow H^{0}(Y,h^{\ast}(F))
$$
is injective, then in this question we may at once replace $F$ by a sheaf
$G$ into which $F$ embeds, and hence may replace $F$ by the product of the
$p_{i\ast}(G_{X'_i})$. In view of \hyperref[XII.lem6.4]{6.4}, we are thus
reduced to proving the injectivity assertion after replacing
$h:Y\rightarrow X$ by $h':Y'\rightarrow X'$ and $F$ by a constant sheaf on
$X'$ whose value is an ordinary finite $\LL$-torsion group $G$. But $G$ is
isomorphic to a direct sum of groups of the form
$\ZZ/\ell^{\nu}\ZZ$, with $\ell\in\LL$. We are therefore indeed reduced
to verifying b).

When $n\geq0$, we proceed similarly. By induction on $n$, we may suppose
that the desired conclusion (and hence a)) has been proved for $n-1$ in
place of $n$. Applying \hyperref[XII.cor6.6]{6.6} (iii), we see that we may
again replace $F$ by a sheaf $G$ into which $F$ embeds, and we again take
$G=p_{\ast}(G_{X'})$. In view of \hyperref[XII.lem6.4]{6.4} and the
vanishing of $R^{i}p_{\ast}$ for $i\geq1$ when $p$ is finite (VIII \hyperref[VIII.prop5.5]{5.5}),
we are again reduced to proving surjectivity for $H^n$ after replacing
$h:Y\rightarrow X$ by $h':Y'\rightarrow X'$ and $F$ by a constant sheaf
on $X'$ whose value is an ordinary finite $\LL$-torsion group. Injectivity
is already known by the induction hypothesis and the implication
${\rm a')}\Rightarrow{\rm a)}$; by devissage on $G$, we are thus again
reduced to the case in which $G$ is of the form
$\ZZ/\ell^{\nu}\ZZ$, namely the case considered in condition b) of
\hyperref[XII.prop6.5]{6.5} (iii). This proves case (iii) of
\hyperref[XII.prop6.5]{6.5}.

Cases (i) and (ii) are proved in exactly the same way as regards the
equivalence of conditions a) and b) in \hyperref[XII.prop6.5]{6.5} (i)
\resp \hyperref[XII.prop6.5]{6.5} (ii): one uses\pageoriginale
\hyperrefext[IX][IX.cor2.7.2]{2.7.2} and the commutation of $H^{0}$ and
$H^1$ with filtered colimits of sheaves (VII \hyperref[VII.rem5.14]{5.14}) to reduce to the case in
which $F$ is constructible, and then uses \hyperrefext[IX][IX.prop2.14]{2.14},
\hyperref[XII.cor6.6]{6.6} (i), (ii), and VIII \hyperref[VIII.rem5.8]{5.8} to reduce to the special
case considered in b). Moreover, in (i), the equivalence of conditions b),
${\rm b'}$), and ${\rm b''}$) is trivial and is recorded only for
reference, as is the fact that they imply the faithfulness (\resp full
faithfulness) of the functor $X'\mapsto Y'$ from the category of etale
coverings of $X$ to the category of etale coverings of $Y$. The same is true
of the equivalence of conditions b), ${\rm b'}$), and ${\rm b''}$) in
\hyperref[XII.prop6.5]{6.5} (ii), which were added for good measure. The
proof of \hyperref[XII.prop6.5]{6.5} is now complete.

\begin{enonce*}[remark]{Remark 6.7}\etiquette[6.7]{XII.rem6.7}
\footnote{The editor recommends omitting these remarks, as well as
\hyperref[XII.rems6.13]{6.13}; they were surreptitiously introduced
(together with various other more or less felicitous changes to the
original text) by an irreverent collaborator.}
\begin{enumerate}
\renewcommand{\labelenumi}{\theenumi)}
\item The arguments proving \hyperref[XII.prop6.5]{6.5} and
  \hyperref[XII.cor6.6]{6.6} are plainly very general and essentially
trivial, and it would be useful to extract them as abstract lemmas of the
same kind as \hyperref[XII.lem6.2]{6.2} (with the sheaves
$p_{\ast}(G_{X'})$ playing the role of cogenerators of the category $C$ in
which one works). We leave this pleasant exercise to the reader, merely
pointing out that essentially the same statement could be given for
functors such as $R^{i}f_{\ast}$ in place of the functors $H^{i}$. The same
remarks apply to \hyperref[XII.prop6.8]{6.8} and
\hyperref[XII.prop6.11]{6.11} below.

\item When $X$ is noetherian, the arguments given show that in conditions
b) and their variants in \hyperref[XII.prop6.5]{6.5} (i), (ii), and (iii),
one may restrict to \emph{integral} $X'$; when $X$ is universally Japanese,
one may even take them to be integral and normal. This is also possible\pageoriginale
without the Japanese hypothesis on $X$, provided that \emph{integral}
morphisms $p:X'\rightarrow X$ are used in place of finite morphisms. None of
these variants will be used later in the seminar.

\item The proof of \hyperref[XII.cor6.6]{6.6} does not use the
quasi-compactness and quasi-separatedness hypothesis on $X$. The same is
true of the fact that in (i), (ii), and (iii), condition a) implies the
other conditions, and of the mutual equivalence of conditions b),
${\rm b'}$), and ${\rm b''}$). Note also that the form b) of the conditions
in \hyperref[XII.prop6.5]{6.5} (i) shows that, for $\LL\neq\phi$, they are
special cases (when $X$ is quasi-compact and quasi-separated) of
\hyperref[XII.prop6.5]{6.5} (iii), for $n=-1$ \resp $n=0$. This remains
true, moreover (if one restricts to the conditions in form a)), without the
stated restriction on $X$. Indeed, for every sheaf of sets $F$ and every
prime number $\ell$, one can find two abelian $\ell$-torsion sheaves $G$
and $H$, and a homomorphism $u:G\rightarrow H$ of the underlying sheaves of
sets, such that $F$ is isomorphic to the inverse image of the zero section
(one takes for $G$ the free $\ZZ/\ell\ZZ$-module generated by $F$, and for
$H$ the constant sheaf $\ZZ/\ell\ZZ$ on $X$). The same argument shows that
(without any quasi-compactness or quasi-separatedness restriction)
\hyperref[XII.prop6.5]{6.5} (ii) a) (where it even suffices to take $i=0$)
implies the respective condition \hyperref[XII.prop6.5]{6.5} (i) a). Since
the latter already implies the injectivity of
$H^1(X,F)\rightarrow H^1(Y,h^{\ast}(F))$ for every sheaf of groups $F$ on
$X$, in the statement of \hyperref[XII.prop6.5]{6.5} (ii) a) it therefore
suffices, for $i=1$, to require the \emph{surjectivity} of
$H^1(X,F)\rightarrow H^1(Y,h^{\ast}(F))$. Unfortunately,
\hyperref[XII.prop6.5]{6.5} (ii) cannot be regarded as a special case of
\hyperref[XII.prop6.5]{6.5} (iii), which often forces us to repeat in the
noncommutative case an argument already carried out in essence in the
commutative case. Note also that, apart from \hyperref[XII.cor6.10]{6.10},
the results that follow are stated so as to remain valid
without\pageoriginale any quasi-compactness or quasi-separatedness
hypothesis.
\end{enumerate}
\end{enonce*}

\begin{enonce*}{Proposition 6.8}\etiquette[6.8]{XII.prop6.8}
(Descent lemma). With the general notation of \hyperref[XII.prop6.5]{6.5},
suppose in addition that a morphism $p:\bar{X}\rightarrow X$ is given. Let
$\bar{h}:\bar{Y}\rightarrow\bar{X}$ denote the morphism obtained from $h$
by the base change $p$, and let $p_Y:\bar{Y}\rightarrow Y$ denote the
canonical morphism. We suppose either that $p$ is surjective, or that $h$
is a closed immersion and $X=p(\bar{X})\cup h(Y)$.
\begin{enumerate}
\renewcommand{\theenumi}{\roman{enumi}}
\renewcommand{\labelenumi}{(\theenumi)}
\item Suppose that $\bar{h}:\bar{Y}\rightarrow\bar{X}$ satisfies the
injectivity (\resp bijectivity) condition of
\hyperref[XII.prop6.5]{6.5} {\rm(i) a)}. In the respective case, suppose in
addition that the base-change morphism
\begin{equation*}
h^{\ast}(p_{\ast}(\bar{F}))\longrightarrow
p_{Y\ast}(\bar{h}^{\ast}(\bar{F}))\tag{$\ast$}
\end{equation*}
is injective for every sheaf of sets $\bar{F}$ on $\bar{X}$. Then the
morphism $h:Y\rightarrow X$ also satisfies the injectivity (\resp
bijectivity) condition of \hyperref[XII.prop6.5]{6.5} {\rm(i) a)}.

\item Suppose that $\bar{h}:\bar{Y}\rightarrow\bar{X}$ satisfies condition
\hyperref[XII.prop6.5]{6.5} {\rm(ii) a)}, and that for every sheaf of
ind-$\LL$-groups $\bar{F}$ on $\bar{X}$ the base-change morphism $(\ast)$
is bijective. Then $h:Y\rightarrow X$ also satisfies condition
\hyperref[XII.prop6.5]{6.5} {\rm(ii) a)}.

\item Suppose that $\bar{h}:\bar{Y}\rightarrow\bar{X}$ satisfies condition
\hyperref[XII.prop6.5]{6.5} {\rm(iii) a)}, and that for every abelian
$\LL$-torsion sheaf $\bar{F}$ on $\bar{X}$ the base-change homomorphism
$$
h^{\ast}(R^{i}p_{\ast}(\bar{F}))\longrightarrow
R^{i}p_{Y\ast}(\bar{h}^{\ast}(\bar{F}))
$$
is bijective for $i\leq n-1$ and injective for $i=n$, where $n\geq1$ is a
given integer. Then $h:Y\rightarrow X$ also satisfies condition
\hyperref[XII.prop6.5]{6.5} {\rm(iii) a)}.
\end{enumerate}
\end{enonce*}

\noindent
\emph{Proof of} \hyperref[XII.prop6.8]{6.8}.\pageoriginale In the case where
$p$ is not assumed surjective, but $h$ is a closed immersion and
$X=p(\bar{X})\cup h(Y)$, consider the disjoint-union scheme $\bar{X}'$ of
$\bar{X}$ and $Y$, and the morphism $p':\bar{X}'\rightarrow X$ induced by
$p$ and $h$. Then $p'$ is surjective; moreover, in each of the three cases
(i), (ii), and (iii), it is immediate that the hypotheses on the pair
$(h,p)$ remain satisfied for the pair $(h,p')$. We are therefore reduced to
the case in which $p$ is surjective, which we assume from now on. In this
case, for every sheaf $F$ on $X$, the canonical homomorphism
$$
F\longrightarrow p_{\ast}(p^{\ast}(F))
$$
is injective. In case (i), to verify the injectivity of
$H^{0}(X,F)\rightarrow H^{0}(Y,h^{\ast}(F))$, we are immediately reduced
to the case in which $F$ is replaced by $p_{\ast}(\bar{F})$, where
$\bar{F}=p^{\ast}(F)$. Consider the commutative diagram
$$
\xymatrix@C=2.75cm@R=.75cm{
**[r]H^{0}(X,p_{\ast}(\bar{F})) \ar[dd]^{\text{\rotatebox{90}{$\backsim$}}} \ar[r]& **[r]H^{0}(Y,h^{\ast}(p_{\ast}(\bar{F}))) \ar[d]\\
                           &**[r]H^{0}(Y,p_{Y\ast}(\bar{h}^{\ast}(\bar{F})))\ar[d]_{\text{\rotatebox{90}{$\backsim$}}}\\
**[r]H^{0}(\bar{X},\bar{F})  \ar[r] &
                           **[r]H^{0}(\bar{Y},\bar{h}^{\ast}(\bar{F})),
}
$$
where the left vertical arrow and the second right vertical arrow are the
canonical isomorphisms, and the first right vertical arrow is induced by
the base-change homomorphism for $h$. By hypothesis the second horizontal
arrow is injective, and it follows immediately that the first is as well;
this proves the non-respective assertion of (i). For the respective
assertion, suppose that the second horizontal arrow is surjective. Since the
base-change homomorphism is a monomorphism by hypothesis, the first right
vertical arrow is a monomorphism, and hence so is the composite of the
two\pageoriginale right vertical arrows. It follows immediately that the
first horizontal arrow is also surjective, completing the proof of (i).

To prove (ii), what precedes and \hyperref[XII.prop6.5]{6.5} (i) reduce us
to proving the surjectivity of
$H^1(X,F)\rightarrow H^1(Y,h^{\ast}(F))$. Applying
\hyperref[XII.cor6.6]{6.6} (ii), we are further reduced to proving it for
$F$ of the form $p_{\ast}(\bar{F})$, where $\bar{F}$ is a sheaf of
ind-$\LL$-groups on $\bar{X}$ (\hyperrefext[IX][IX.prop1.6]{1.6} (ii)).
Writing the exact sequences for $H^1$ (variants of
\hyperref[XII.prop3.2]{3.2}) for the morphisms $p$ and $p_Y$, we obtain a
morphism of noncommutative exact sequences:
$$
\xymatrix@R=1.5cm{
H^{1}(X,p_{\ast}(\bar{F})) \ar[r] \ar[d] & H^{1}(\bar{X},\bar{F}) \ar[r]\ar[d] &
H^{0}(X,R^{1}p_{\ast}(\bar{F})) \ar[d]&\\
H^{1}(Y,p_{Y\ast}(\bar{h}^{\ast}(\bar{F}))) \ar[r] &
H^{1}(\bar{Y},\bar{h}^{\ast}(\bar{F})) \ar[r] &
H^{0}(Y,R^{1}p_{Y\ast}(\bar{h}^{\ast}(\bar{F}))),
}
$$
where in each row the first arrow is injective and identifies the first term
with the inverse image, under the second arrow, of the distinguished point
of the third term. By the hypothesis on $\bar{h}$, the middle vertical arrow
is bijective. The two outer vertical arrows factor, respectively, through
$$
H^1(Y,h^{\ast}(p_{\ast}(\bar{F})))
\quad\hbox{and}\quad
H^{0}(Y,h^{\ast}(R^1p_{\ast}(\bar{F}))),
$$
by composing the homomorphisms $H^{i}(X,-)\rightarrow H^{i}(Y,h^{\ast}(-))$
(for $i=1$ and $i=0$, respectively) with the homomorphisms obtained by
applying $H^{i}(Y,-)$ (for those same values of $i$) to the base-change
homomorphisms
$$
h^{\ast}(p_{\ast}(\bar{F}))\longrightarrow
p_{Y\ast}(\bar{h}^{\ast}(\bar{F}))
$$
and, respectively,
$$
h^{\ast}(R^1p_{\ast}(\bar{F}))\longrightarrow
R^1p_{Y\ast}(\bar{h}^{\ast}(\bar{F})).
$$
Since the first of these arrows is an isomorphism by hypothesis, the first
vertical arrow in the diagram above identifies with the homomorphism
$H^1(X,-)\rightarrow H^1(Y,h^{\ast}(-))$ that we wish to study. To prove
that it is bijective, it therefore remains to prove that the last vertical
arrow in the diagram is injective. By what we have just said, it is the
composite of two maps,\pageoriginale the first of which is a canonical
homomorphism $H^{0}(X,-)\rightarrow H^{0}(Y,h^{\ast}(-))$, and hence is
bijective, while the second is injective because it is obtained by applying
the left exact functor $H^{0}(Y,-)$ to the monomorphism
$$
h^{\ast}(R^1p_{\ast}(\bar{F}))\longrightarrow
R^1p_{Y\ast}(\bar{h}^{\ast}(\bar{F})).
$$
The latter fact follows readily from the bijectivity hypothesis for the
base-change homomorphism in the preceding dimension $0$, for example by
using the second assertion of \hyperref[XII.prop6.5]{6.5} (i) and the usual
calculation of the fibers of the higher direct images (VIII \hyperref[VIII.rem5.3]{5.3}). This
proves \hyperref[XII.prop6.8]{6.8} (ii).

It remains to prove \hyperref[XII.prop6.8]{6.8} (iii). For this, rather than
using the injection $F\rightarrow p_{\ast}(p^{\ast}(F))$ directly as above,
it is more convenient to use the homomorphism
$$
F\longrightarrow Rp_{\ast}(\bar{F})
$$
in the right derived category $D^{+}(C)$, where $C$ denotes the category of
abelian sheaves on $X$. Recall that, by definition,
$Rp_{\ast}(\bar{F})$ is the complex $p_{\ast}(C(\bar{F}))$, where
$C(\bar{F})$ is an injective resolution of $\bar{F}$ in the category of
abelian sheaves on $\bar{X}$. If $K^{\bigdot}$ is a complex on $X$, we
denote by $H^{i}(X,K^{\bigdot})$ the hypercohomology groups of $X$ with
coefficients in $K^{\bigdot}$, and use analogous notation on $Y$. By
\hyperref[XII.cor6.3]{6.3}, we are reduced to proving the surjectivity of
\begin{equation*}
H^n(X,Rp_{\ast}(\bar{F}))\longrightarrow
H^n(Y,h^{\ast}(Rp_{\ast}(\bar{F}))).\tag{1}
\end{equation*}\etiquette[1]{XII.eq1}
By hypothesis, the base-change homomorphism
\begin{equation*}
h^{\ast}(Rp_{\ast}(\bar{F}))\longrightarrow
Rp_{Y\ast}(\bar{h}^{\ast}(\bar{F}))\tag{2}
\end{equation*}\etiquette[2]{XII.eq2}
induces an isomorphism on the cohomology sheaves in degrees $\leq n-1$ and
a monomorphism in degree $n$. Equivalently, the ``mapping cylinder''
\pageoriginale $K^{\bigdot}$ of the preceding homomorphism has nonzero
cohomology sheaves only in degrees $\geq n$. Hence
$H^{i}(Y,K^{\bigdot})=0$ for $i<n$, and the cohomology exact sequence shows
that the homomorphisms
\begin{equation*}
H^{i}(Y,h^{\ast}(Rp_{\ast}(\bar{F})))\longrightarrow
H^{i}(Y,Rp_{Y\ast}(\bar{h}^{\ast}(\bar{F})))\tag{3}
\end{equation*}\etiquette[3]{XII.eq3}
induced by the homomorphism above are isomorphisms for $i\leq n-1$ and a
monomorphism for $i=n$. Consider the composite of (1) and (3):
\begin{equation*}
H^n(X,Rp_{\ast}(\bar{F}))\longrightarrow
H^n(Y,h^{\ast}(Rp_{\ast}(\bar{F})))\longrightarrow
H^n(Y,Rp_{Y\ast}(\bar{h}^{\ast}(\bar{F})))\tag{4}
\end{equation*}\etiquette[4]{XII.eq4}

% END INCLUDED SGA4_English_current_buildable_source_20260906/source/12/drafts/SGA4_XII_LINES_1201_1600_English.texfrag


% BEGIN INCLUDED SGA4_English_current_buildable_source_20260906/source/12/drafts/SGA4_XII_LINES_1601_1924_English.texfrag
The two outer terms are respectively isomorphic, by definition, to
$H^n(\bar{X},\bar{F})$ and
$H^n(\bar{Y},\bar{h}^{\ast}(\bar{F}))$, and the composite in (4) is
precisely the homomorphism induced by $\bar{h}$, which is bijective by the
hypothesis on $\bar{h}$. Since the second arrow in (4) is injective by what
we have just seen, the first is also bijective. This completes the proof of
\hyperref[XII.prop6.8]{6.8}.

Moreover, the preceding proof immediately gives the following slightly
more precise and more general result:

\begin{enonce*}{Corollary 6.9}\etiquette[6.9]{XII.cor6.9}
The notation is that of \hyperref[XII.prop6.8]{6.8}. Suppose in addition
that $h$ satisfies the unparenthesized condition in
\hyperref[XII.prop6.5]{6.5} \textup{(i)} (respectively, the parenthesized
condition in \hyperref[XII.prop6.5]{6.5} \textup{(i)}, respectively, the
condition in \hyperref[XII.prop6.5]{6.5} \textup{(iii)}, with $n\geq -1$), and
that, for every sheaf of sets $\bar{F}$ on $\bar{X}$, the base-change
homomorphism
$$
h^{\ast}(p_{\ast}(\bar{F}))\longrightarrow
p_{Y\ast}(\bar{h}^{\ast}(\bar{F}))
$$
is injective (respectively, bijective; respectively, that for every sheaf
of $\LL$-torsion groups $\bar{F}$ on $\bar{X}$, the base-change
homomorphism
$$
h^{\ast}R^{i}p_{\ast}(\bar{F})\longrightarrow
R^{i}p_{Y\ast}(\bar{h}^{\ast}(\bar{F}))
$$
is bijective for $i\leq n$ and injective for $i=n+1$). Let $F$ be a sheaf
of sets (respectively, of ind-$\LL$-groups, respectively, a sheaf of
$\LL$-torsion groups) on $X$, and let
$\xi\in H^{n+1}(Y,h^{\ast}(F))$ (where $n=-1$ in case\pageoriginale \textup{(i)} and $n=0$
in case \textup{(ii)}). Put $\bar{F}=p^{\ast}(F)$. Then $\xi$ lies in the image
of $H^{n+1}(X,F)$ if and only if its inverse image in
$$
H^{n+1}(\bar{Y},p_Y^{\ast}(h^{\ast}(F)))
=H^{n+1}(\bar{Y},\bar{h}^{\ast}(\bar{F}))
$$
lies in the image of $H^{n+1}(\bar{X},\bar{F})$.
\end{enonce*}

We shall use this sharpened form of \hyperref[XII.prop6.8]{6.8} to prove
the following result.

\begin{enonce*}{Corollary 6.10}\etiquette[6.10]{XII.cor6.10}
The notation is that of \hyperref[XII.prop6.5]{6.5}. For simplicity,
suppose that the unparenthesized condition in
\hyperref[XII.prop6.5]{6.5} \textup{(i)} holds.
\begin{enumerate}
\renewcommand{\theenumi}{\roman{enumi}}
\renewcommand{\labelenumi}{(\theenumi)}
\item The parenthesized condition in \hyperref[XII.prop6.5]{6.5}
\textup{(i)} holds (respectively, condition
\hyperref[XII.prop6.5]{6.5} \textup{(ii) a)} holds, respectively, condition
\hyperref[XII.prop6.5]{6.5} \textup{(iii) a)} holds) if and only if the
following condition is satisfied. For every sheaf $F$ of sets
(respectively, of ind-$\LL$-groups, respectively, of abelian
$\LL$-torsion groups) on $X$, every
$\xi\in H^{\circ}(Y,h^{\ast}(F))$ (respectively, every
$\xi\in H^{i}(Y,h^{\ast}(F))$ with $i=0,1$, respectively, every
$\xi\in H^{i}(Y,h^{\ast}(F))$ with $i\leq n$), and every nonempty closed
subset $X'$ of $X$, if $Y'$ denotes its inverse image in $Y$, there exists
a morphism $p:\bar{X}\rightarrow X$ satisfying the following conditions:
\begin{itemize}
\item[$1^{\circ}$)] The image $p(\bar{X})$ is contained in $X'$ and
contains a nonempty open subset of $X'$.

\item[$2^{\circ}$)] For every sheaf of sets $\bar{F}$ on $\bar{X}$, the
base-change homomorphism
$$
h^{\ast}(p_{\ast}(\bar{F}))\longrightarrow
p_{Y\ast}(\bar{h}^{\ast}(\bar{F}))
$$
is an isomorphism (respectively, for every sheaf of ind-$\LL$-groups
$\bar{F}$ on $\bar{X}$, the base-change homomorphism
$$
h^{\ast}(R^{i}p_{\ast}(\bar{F}))\longrightarrow
R^{i}p_{Y\ast}(\bar{h}^{\ast}(\bar{F}))
$$
is bijective for $i=0,1$; respectively, for every sheaf of abelian
$\LL$-torsion groups $\bar{F}$ on $\bar{X}$, the base-change homomorphism
is bijective for $i\leq n$).

\item[$3^{\circ}$)] The inverse image of $\xi$ in
$$
H^{i}(\bar{Y},p_Y^{\ast}(h^{\ast}(F)))
=H^{i}(\bar{Y},\bar{h}^{\ast}(p^{\ast}(F)))
$$
lies in the image of $H^{i}(\bar{X},p^{\ast}(F))$.
\end{itemize}

\item Let\pageoriginale $n\geq -1$, and suppose that, for every sheaf of
sets $F$ (respectively, every sheaf of ind-$\LL$-groups, respectively,
every sheaf of abelian $\LL$-torsion groups) on $X$, the homomorphism
$$
H^{i}(X,F)\longrightarrow H^{i}(Y,h^{\ast}(F))
$$
is bijective for $i\leq n$ and injective for $i=n+1$ (in the first case
we take $n=-1$, and in the second case $n=-1$ or $0$). Let
$\xi\in H^{n+1}(Y,h^{\ast}(F))$. Then $\xi$ lies in the image of
$H^{n+1}(X,F)$ if and only if it satisfies the condition stated in
\textup{(i)} above, with $i=n+1$.
\end{enumerate}
\end{enonce*}

\noindent
\emph{Proof of} \hyperref[XII.cor6.10]{6.10}. Using
\hyperref[XII.prop6.5]{6.5} and an immediate induction on $n$, one sees
that \hyperref[XII.cor6.10]{6.10} (i) follows from
\hyperref[XII.cor6.10]{6.10} (ii), which we now prove.

Suppose that $\xi$ does not lie in the image of $H^{n+1}(X,F)$. Let
$\Phi$ be the set of closed subsets $X'$ of $X$ such that the image of
$\xi$ in $H^{n+1}(Y',h^{\ast}(F)|Y')$ does not lie in the image of
$H^{n+1}(X',F|X')$. By hypothesis, $X\in\Phi$, so $\Phi$ is nonempty.
Order $\Phi$ by the relation $\supset$, and let us show that $\Phi$ is
inductive. It is enough to prove that, if $(X'_{\lambda})$ is a totally
ordered family of elements of $\Phi$ and $X'$ is their intersection, then
$X'\in\Phi$. Give the $X'_{\lambda}$ and $X'$ their induced reduced
structures. The morphisms $X'_{\lambda}\rightarrow X$ are affine (being
closed immersions), and hence the inverse system of the $X'_{\lambda}$
satisfies the hypotheses considered in VII 5. Moreover, one immediately
sees that $X'$ is the inverse limit of the $X'_{\lambda}$. Likewise, $Y'$
is the inverse limit of the inverse system of the $Y'_{\lambda}$, which
are affine over $Y$. By VII \hyperref[VII.cor5.8]{5.8}, therefore,
$$
\varinjlim H^{n+1}(X'_{\lambda},F|X'_{\lambda})\xrightarrow{\sim}
H^{n+1}(X',F|X')\quav
\varinjlim H^{n+1}(Y'_{\lambda},h^{\ast}(F)|Y'_{\lambda})
\xrightarrow{\sim}H^{n+1}(Y',h^{\ast}(F)|Y'){\quoi}.
$$
It follows at once that, if
$(\xi|Y')\in\Iim H^{n+1}(X',F|X')$, then there\pageoriginale would be an index $\lambda$
such that
$(\xi|Y'_{\lambda})\in\Iim H^{n+1}(X'_{\lambda},F|X'_{\lambda})$, which
is absurd. Thus $X'\in\Phi$, as asserted. Consequently $\Phi$ is
inductive and therefore contains an inclusion-minimal element, say $X'$.
Since $\varnothing\notin\Phi$, the set $X'$ is nonempty. Apply the
hypothesis of \hyperref[XII.cor6.10]{6.10} to $\xi$ and $X'$. This gives a
morphism $p:\bar{X}\rightarrow X$ satisfying conditions $1^{\circ}$)--
$3^{\circ}$) in \hyperref[XII.cor6.10]{6.10}. Observe that we may even
suppose that $p$ is \emph{surjective onto $X'$}. Indeed, let $U$ be a
nonempty open subset of $X'$ contained in $p(\bar{X})$, and let $X'_1$ be
its complement in $X'$. By construction, $X'_1\notin\Phi$. Put
$\bar{X}_1=\bar{X}\amalg X'_1$, and let
$p_1:\bar{X}_1\rightarrow X$ be the morphism defined by $p$ and the
inclusion $X'_1\hookrightarrow X$. Using the relation
$X'_1\notin\Phi$, one sees immediately that $p_1$ still satisfies the same
hypotheses as $p$; moreover, it is surjective onto $X'$. We may therefore
suppose that $p$ is surjective onto $X'$. Now apply
\hyperref[XII.cor6.9]{6.9}, replacing $h:Y\rightarrow X$ there by
$h':Y'\rightarrow X'$ (with $X'$ endowed, say, with its induced reduced
structure), and replacing $\xi$ by $\xi|Y'$. It follows from
\hyperref[XII.cor6.9]{6.9} that
$(\xi|Y')\in\Iim H^{n+1}(X',F|X')$, contrary to $X'\in\Phi$. This
completes the proof of \hyperref[XII.cor6.10]{6.10}.

We shall use \hyperref[XII.cor6.10]{6.10} to reduce
\hyperref[XII.cor5.5]{5.5} to the case in which $f$ is projective. In the
non-Noetherian case, however, this requires a non-Noetherian variant of
Chow's lemma, which will be given in the next section.

\begin{enonce*}{Proposition 6.11}\etiquette[6.11]{XII.prop6.11}
(Transitivity lemma). The notation is that of
\hyperref[XII.prop6.8]{6.8}, but no surjectivity hypothesis is imposed on
$p$. Suppose that $h$ satisfies condition \textup{(i) a)} (respectively,
\textup{(ii) a)}, respectively, \textup{(iii) a)}) of
\hyperref[XII.prop6.5]{6.5}. In case \textup{(i)}, suppose in addition that,
for every sheaf of sets $\bar{F}$ on $\bar{X}$, the base-change
homomorphism of \hyperref[XII.prop6.8]{6.8} is a monomorphism
(respectively, an isomorphism); in case \textup{(ii)}, suppose that $\bar{X}$
is Noetherian and that, for every sheaf of ind-$\LL$-groups $\bar{F}$, the
base-change homomorphism
$$
h^{\ast}(R^ip_{\ast}(\bar{F}))\longrightarrow
R^ip_{Y\ast}(\bar{h}^{\ast}(\bar{F}))
$$
is\pageoriginale bijective for $i=0,1$; and finally, in case \textup{(iii)}, suppose that the
base-change homomorphism is bijective for $i\leq n$ and injective for
$i=n+1$. Then $\bar{h}:\bar{Y}\rightarrow\bar{X}$ also satisfies
condition \textup{(i)} (respectively, \textup{(ii)}, respectively, \textup{(iii)}) of
\hyperref[XII.prop6.5]{6.5}.
\end{enonce*}

\noindent
\emph{Proof of} \hyperref[XII.prop6.11]{6.11}. In case (i), consider the
commutative diagram in the proof of \hyperref[XII.prop6.8]{6.8} (i). Here
the hypotheses imply that the first horizontal arrow and the upper right
vertical arrow are injective (respectively, bijective), and hence so is the
second horizontal arrow, as required. In case (ii), let $\bar{F}$ be a
sheaf of ind-$\LL$-groups on $\bar{X}$. We must show that every element of
$H^1(\bar{Y},\bar{h}^{\ast}(\bar{F}))$ comes from an element of
$H^1(\bar{X},\bar{F})$. By \hyperref[XII.cor6.6]{6.6} (ii), applied to
$\bar{h}$, it is enough to find a monomorphism
$\bar{F}\rightarrow\bar{G}$ of sheaves of groups on $\bar{X}$ that
\emph{effaces} $H^1(\bar{Y},\bar{h}^{\ast}(\bar{F}))$. Consider the
second row of the diagram used in the proof of
\hyperref[XII.prop6.8]{6.8} (ii). It is enough, successively, to find a
monomorphism $\bar{F}\rightarrow\bar{G}$, with $\bar{G}$ a sheaf of
ind-$\LL$-finite groups, that effaces the last term
$H^{\circ}(Y,R^1p_{Y\ast}(\bar{h}^{\ast}(\bar{F})))$, and then a
monomorphism of sheaves of groups $\bar{G}\rightarrow\bar{H}$ that
effaces $H^1(Y,p_{Y\ast}(\bar{h}^{\ast}(\bar{G})))$. As observed in the
proof of \hyperref[XII.prop6.8]{6.8} (i), the two terms to be effaced are
respectively isomorphic, by the outer vertical arrows of the diagram in
question, to $H^{\circ}(X,R^1p_{\ast}(\bar{F}))$ and
$H^1(X,p_{\ast}(\bar{F}))$, in view of the hypotheses made here on the
base-change homomorphisms for $p$ and $h$. These isomorphisms are plainly
functorial in $\bar{F}$, so it is enough to efface the latter two terms.
By \hyperref[XII.prop3.3]{3.3}, one can efface
$R^1p_{\ast}(\bar{F})$ (and, a fortiori,
$H^{\circ}(X,R^1p_{\ast}(\bar{F}))$) by embedding $\bar{F}$ into an
ind-$\LL$-group $\GG$. It remains to efface
$H^1(X,p_{\ast}(\bar{F}))$. This is possible because, by\pageoriginale the diagram in
question, the latter term maps into $H^1(\bar{X},\bar{F})$ by a
monomorphism functorial in $\bar{F}$, while
$H^1(\bar{X},\bar{F})$ is again effaceable by
\hyperref[XII.prop3.3]{3.3}.

It remains to treat case (iii), and here again we read the proof of
\hyperref[XII.prop6.8]{6.8} backwards\footnote[1]{Instead of
\hyperref[XII.cor6.3]{6.3}, one could also use the Leray spectral sequence
for $p$ and $p_Y$.}. We must prove that the homomorphism
$$
H^{i}(\bar{X},\bar{F})\longrightarrow
H^{i}(\bar{Y},\bar{h}^{\ast}(\bar{F}))
$$
is an isomorphism for $i\leq n$ and a monomorphism for $i=n+1$. These
homomorphisms are induced by a homomorphism of complexes
$$
\RR\Gamma_{\bar{X}}(\bar{F})\longrightarrow
\RR\Gamma_{\bar{Y}}(\bar{h}^{\ast}(\bar{F})){\quoi},
$$
obtained from an injective resolution $C(\bar{F})$ and a homomorphism from
$\bar{h}^{\ast}(C(\bar{F}))$ (which is a resolution of
$\bar{h}^{\ast}(\bar{F})$) to an injective resolution of
$\bar{h}^{\ast}(\bar{F})$. The homomorphism above also identifies with
the composite
$$
\RR\Gamma_X(Rp_{\ast}(\bar{F}))\longrightarrow
\RR\Gamma_Y(h^{\ast}(Rp_{\ast}(\bar{F})))\longrightarrow
\RR\Gamma_Y(Rp_{Y\ast}(\bar{h}^{\ast}(\bar{F}))){\quoi},
$$
where the first homomorphism is the homomorphism induced by $h^{\ast}$ for
$Rp_{\ast}(\bar{F})$, and the second is obtained by applying
$\RR\Gamma_Y$ to the base-change homomorphism
$$
h^{\ast}(Rp_{\ast}(\bar{F}))\longrightarrow
Rp_{Y\ast}(\bar{h}^{\ast}(\bar{F})).
$$
By the hypothesis on $h$ and \hyperref[XII.cor6.3]{6.3} (i), the first
homomorphism therefore induces isomorphisms on the cohomology groups of
the complexes in question in the required range. The same is true of the
second by our hypothesis on the base-change homomorphisms, using the
argument given in the proof of \hyperref[XII.prop6.8]{6.8} (iii). The same
conclusion thus holds for their composite, completing the proof of
\hyperref[XII.prop6.11]{6.11}.

\begin{enonce*}[remark]{Remark 6.12}\etiquette[6.12]{XII.rem6.12}
The argument\pageoriginale given for \hyperref[XII.prop6.11]{6.11} in cases (i) and
(iii) gives the following more precise result. In order to conclude, for a
\emph{given} sheaf $\bar{F}$ on $\bar{X}$, that
$$
H^{i}(\bar{X},\bar{F})\longrightarrow
H^{i}(\bar{Y},\bar{h}^{\ast}(\bar{F}))
$$
is bijective for $i\leq n$ and injective for $i=n+1$, it is enough to
assume that
\begin{itemize}
\item[$1^{\circ}$)] the base-change homomorphism
$$
h^{\ast}(R^{i}p_{\ast}(\bar{F}))\longrightarrow
R^{i}p_{Y\ast}(\bar{h}^{\ast}(\bar{F}))
$$
is an isomorphism for $i\leq n$ and a monomorphism for $i=n+1$, and

\item[$2^{\circ}$)] for every pair of integers $i,j\geq0$, the
homomorphism
$$
H^{i}(X,R^{j}p_{\ast}(\bar{F}))\longrightarrow
H^{i}(Y,h^{\ast}(R^{j}p_{\ast}(\bar{F})))
$$
induced by $h$ is an isomorphism if $i+j\leq n$ and a monomorphism if
$i+j=n+1$.
\end{itemize}

There is plausibly an analogous result in case (ii), but its statement
would necessarily involve noncommutative $2$-cohomology. To avoid having
to invoke that theory, in \hyperref[XII.prop6.11]{6.11} (ii) we used a
different argument, based on the effaceability properties of
noncommutative $H^1$; unlike the argument above, however, it yields no
result for a fixed $\bar{F}$.
\end{enonce*}

\begin{enonce*}[remark]{Remarks 6.13}\etiquette[6.13]{XII.rems6.13}
The most important case in which the conditions of
\hyperref[XII.prop6.5]{6.5} (i), (ii), and (iii) hold is that considered
in \hyperref[XII.cor5.5]{5.5}; it will be proved in the next exposé using
the reductions made in the present one. Another interesting case (in
which $h$ is no longer a closed immersion as in
\hyperref[XII.cor5.5]{5.5}) will be studied in XV. Let us also mention the
closely related case in which $X$ is the spectrum of a Noetherian ring
$A$, separated and complete for the topology defined by an ideal $J$, and
$Y=\Spec(A/J)$ is the closed subscheme of $X$ defined by $J$. One then
verifies condition \hyperref[XII.prop6.5]{6.5} (ii) directly, in the form\pageoriginale
$\textup{ b'')}$ (with $\LL=\PP$), by \EGA IV 18.3.2. It is very plausible
that the conditions of \hyperref[XII.prop6.5]{6.5} (iii) also hold for
every $n$; that is, for every torsion sheaf $F$ on $X$, the homomorphisms
$$
H^n(X,F)\longrightarrow H^n(Y,F|Y)
$$
are bijective. One can formulate considerably more general conjectures,
related to a generalization of the construction of henselization to
nonlocal schemes. First, no example is known in which $h$ is a closed
immersion and the conditions of \hyperref[XII.prop6.5]{6.5} (ii) and (iii)
(with $\LL=\PP$) fail when the parenthesized condition in
\hyperref[XII.prop6.5]{6.5} (i) holds, that is, when $(X,Y)$ is a
\emph{henselian pair} in the terminology of \EGA IV 18.5.5. One may
therefore conjecture that every such pair satisfies the conditions of
\hyperref[XII.prop6.5]{6.5} (ii) and (iii) (with $\LL=\PP$, for every
$n$).

On the other hand, one may ask under what conditions the following
property, for a pair $(X,Y)$ consisting of a quasi-compact and
quasi-separated scheme $X$ and a closed subscheme $Y$, implies that the
pair is henselian:

$(\ast)$ For every scheme $X'$ étale over $X$, put
$Y'=\fprod{X'}{Y}{X}$. The canonical map
$$
\Gamma(X'/X)\longrightarrow\Gamma(Y'/Y)
$$
is bijective (\cf \EGA IV 18.5.4 b)). This condition is plainly weaker
\emph{a priori} than the henselian condition
\hyperref[XII.prop6.5]{6.5} (i), and is in fact strictly weaker: it holds,
for example, when $X$ is a normal irreducible projective scheme of
dimension $\geq2$ over a field $k$ and $Y$ is a hyperplane section of $X$
(as follows readily from \SGA 2 XII 2.1, taking $S=\Spec(k)$), whereas
the pair $(X,Y)$ is plainly not henselian. It is possible, however, that
when $X$ is affine, condition $(\ast)$ implies that $(X,Y)$ is a henselian
pair. If so, for every affine scheme $X$ and every closed subset $Y$ of
$X$, a\pageoriginale process of étale localization along $Y$, analogous to the
henselization of local rings, ought to associate to it a \emph{henselian}
pair $(X',Y)$, where $X'$ is ``pro-étale'' over $X$. Another hypothesis on
$(X,Y)$ (suggested by results of Hironaka and Rossi on the contractibility
of certain subvarieties of algebraic varieties) under which the henselian
property might perhaps be deduced from the weaker property $(\ast)$ is the
following: the immersion $Y\rightarrow X$ is regular of codimension $1$,
and the conormal sheaf $\NN_{Y/X}=\JJ/\JJ^2$ (where $\JJ$ is the ideal
defining $Y$ in $X$) is an \emph{ample} invertible sheaf on $Y$. (In the
counterexample considered above, by contrast, this conormal sheaf was
anti-ample, that is, its inverse was ample.)
\end{enonce*}

\footnotetext{\hypertarget{sga4.exp.XII.note.section7}{}A reader interested in the statements of \hyperref[XII.sec5]{Section 5} only in
the Noetherian case may omit the present section.}

% END INCLUDED SGA4_English_current_buildable_source_20260906/source/12/drafts/SGA4_XII_LINES_1601_1924_English.texfrag


% BEGIN INCLUDED SGA4_English_current_buildable_source_20260906/source/12/drafts/SGA4_XII_LINES_1925_2169_English.texfrag
\section[A variant of Chow's lemma ${}^{(*)}$.]{A variant of Chow's lemma
\texorpdfstring{\protect\hyperlink{sga4.exp.XII.note.section7}{${}^{(*)}$}}{(*)}.}\etiquette[7]{XII.sec7}

We shall need the following variant of (\EGA II 5.6.1), stripped of
all Noetherian and irreducibility hypotheses.

\begin{enonce*}{Lemma 7.1}\etiquette[7.1]{XII.lem7.1}
Let $S$ be a quasi-compact and quasi-separated scheme, and let
$f:X\rightarrow S$ be a separated morphism of finite type, with $X$
nonempty. There exists a commutative diagram
$$
\xymatrix@=1.25cm{
X \ar[rd]_{f}&    &\bar{X}\ar[ll]_{\pi}\ar[ld]^{\bar{f}}\\
  &  S &
}
$$
with\pageoriginale $\pi$ projective and $\bar{f}$ quasi-projective,
and a nonempty open subset $U$ of $X$, such that $\pi$ induces an
isomorphism from
$\bar{U}=\fprod{U}{\bar{X}}{X}=\pi^{-1}(U)$ onto $U$.
\end{enonce*}

\noindent
\emph{Proof.} We first make the following observation. Let
$i:X'\rightarrow X$ be a closed $S$-subscheme of $X$, and let
$V\subset X'$ be a subscheme that is open in $X$ and dense in $X'$.
If the lemma holds for $X'$, then it holds for $X$. Indeed, let
$$
\xymatrix@=1.25cm{
X' \ar[rd]_{f'}&    &\bar{X}'\ar[ll]_{\pi'}\ar[ld]^{\bar{f}'}\\
  &  S &
}
$$
be a diagram as in the statement, and let $U'\neq\varnothing$ be an
open subset of $X'$ such that $\pi'$ induces an isomorphism
$\bar{U}'={\pi'}^{-1}(U')\rightarrow U'$. Since $V$ is dense in
$X'$, we have $U'\cap V\neq\varnothing$, and we may therefore replace
$U'$ by this open subset. Then $U'$ is open in $X$, and we set
$\pi=i\pi'$, $\bar{f}=\bar{f}'$, and $U=U'$.

Now $X$ admits a finite covering by nonempty affine open subsets
$U_k$, $k=1,\ldots,n$. The hypotheses on $S$ imply that every
$U_k\rightarrow S$ is quasi-affine, and hence, a fortiori,
quasi-compact. We prove the lemma for every $X$ by induction on $n$.
If the intersection $U$ of the $U_k$ is dense in $X$, we may copy the
proof of (\EGA II 5.6, B, C, D), with the necessary changes in
notation. If the intersection is nonempty, we reduce to the case in
which it is dense by replacing $X$ by the closure $X'$ of $U$ (as in
the observation above), and $U_k$ by $U'_k=U_k\cap X'$. Finally,
suppose that $U$ is empty, and choose $r$, $1\leq r<n$, such that
$$
V=\bigcap_{k=1}^{r}U_k\neq\varnothing
\qquad\hbox{but}\qquad
V\cap U_{r+1}=\varnothing;
$$
this is plainly possible since $U_1\neq\varnothing$. We may now
replace $X$ by the closure $X'$ of $V$ and each $U_k$ by
$U_k\cap X'$, discarding the empty members. Since
$U_{r+1}\cap X'=\varnothing$ (it is open in $X'$ and disjoint from
the dense open $V$), the result follows from the induction
hypothesis.

\section{Final reductions}\etiquette[8]{XII.sec8}\pageoriginale

\begin{enonce*}{Lemma 8.1}\etiquette[8.1]{XII.lem8.1}
To prove any one of the assertions
{\rm \hyperref[XII.thm5.1]{5.1} (i), (ii), (iii)}, it is enough to do
so under the assumptions that $f$ is \emph{projective} and $S$ is
Noetherian.
\end{enonce*}

Suppose first that \hyperref[XII.thm5.1]{5.1} has been proved when $f$
is projective; we show how the same assertion follows for arbitrary
$f$. Using \hyperref[XII.lem6.1]{6.1}, we are reduced to the situation
of \hyperref[XII.cor5.5]{5.5}, with $S$ strictly local. We now apply
\hyperref[XII.cor6.10]{6.10}, taking $Y=X_{\circ}$. Given
$(\xi,X')$ as in \hyperref[XII.cor6.10]{6.10}, we apply Chow's lemma
\hyperref[XII.lem7.1]{7.1} to find an $S$-morphism
$p:\bar{X}\rightarrow X'$, with $\bar{X}$ projective over $S$, such
that $p(\bar{X})$ contains a nonempty open subset of $X'$ (condition
$1^{\circ}$ of \hyperref[XII.cor6.10]{6.10}). Since $p$ is
projective, the base-change homomorphisms for $(p,h)$ are isomorphisms
by hypothesis (condition $2^{\circ}$ of
\hyperref[XII.cor6.10]{6.10}). Finally, since $\bar{X}$ is projective
over $S$, the inclusion
$\bar{Y}=\bar{X}_{\circ}\rightarrow\bar{X}$ induces, by hypothesis,
an isomorphism on cohomology groups
$$
H^{i}(\bar{X},\bar{F})\longrightarrow
H^{i}(\bar{X}_{\circ},\bar{h}^{\ast}(\bar{F})),
$$
which is condition $3^{\circ}$ of
\hyperref[XII.cor6.10]{6.10}. Thus all the conditions of
\hyperref[XII.cor6.10]{6.10} are satisfied, and the conclusion
follows.

We next show that, in order to verify any one of the three assertions
\hyperref[XII.thm5.1]{5.1} for every situation $(f,F,g)$ with $f$
projective, we may further assume that $S$ is Noetherian. By
\hyperref[XII.lem6.1]{6.1}, we may suppose that $S$ is strictly local
and that $S'\rightarrow S$ is the inclusion of the closed point; we
are thus reduced to verifying \hyperref[XII.cor5.5]{5.5} with $S$
strictly local. Moreover, since $X$ embeds into the projective scheme
$\PP^r_S$, compatibility \hyperref[XII.prop4.4]{4.4} (ii) allows us
to replace $X$ by $\PP^r_S$ and $F$ by its direct image on
$\PP^r_S$. We are therefore reduced to the case $X=\PP^r_S$.

Write $S=\Spec(A)$. The ring $A$ is the filtered direct limit of
Noetherian Henselian local subrings $A_i$ for which the inclusions
$A_i\rightarrow A$ are local: one may take the strict henselizations,
at the prime ideals induced by the maximal ideal of $A$, of the
finitely generated $\ZZ$-subalgebras of $A$. Under these conditions,
the residue field $k$ of $A$ is likewise the filtered direct limit of
the residue fields $k_i$ of the $A_i$ (\EGA IV 5.13.1). Set
$S_i=\Spec(A_i)$ and $X_i=\PP^r_{S_i}$. Suppose that a direct system
$(F_i)$ of sheaves of the type occurring in the assertion under
consideration is given on the $X_i$, as in VII \hyperref[VII.thm5.7]{5.7}, and let $F$ be
the sheaf it defines on $X$. Let $F_{i\circ}$ (respectively,
$F_{\circ}$) be the sheaf induced by $F_i$ on the fiber
$X_{i\circ}$ over the closed point $s_i$ of $S_i$ (respectively, by
$F$ on the fiber $X_{\circ}$ over the closed point $s$ of $S$). We
then have a commutative diagram
$$
\xymatrix{
\displaystyle\mathop{\varinjlim}_{i} H^{n}(X_i,F_i) \ar[r] \ar[d]&
H^{n}(X,F)\ar[d]\\
\displaystyle\mathop{\varinjlim}_{i} H^{n}(X_{i\circ},F_{i\circ}) \ar[r] &
H^{n}(X_{\circ},F_{\circ}),
}
$$
whose two horizontal arrows are isomorphisms by VII \hyperref[VII.thm5.7]{5.7}. Consequently,
to prove that the second vertical arrow is an isomorphism, it is
enough to prove this for the first one; hence it is enough to prove
that the homomorphisms
$$
H^n(X_i,F_i)\longrightarrow H^n(X_{i\circ},F_{i\circ})
$$
are isomorphisms. On the other hand, every sheaf $F$ of the type under
consideration on $X$ is isomorphic to the limit of a direct system of
the preceding kind; for example, one may take $F_i=u_{i*}(F)$, where
$u_i:X\rightarrow X_i$ is the canonical morphism, by
\hyperref[XII.cor5.2]{5.2} and
\hyperrefext[IX][IX.prop1.6]{1.6} (iii). This completes the proof of
\hyperref[XII.lem8.1]{8.1}.

\begin{enonce*}{Corollary 8.2}\etiquette[8.2]{XII.cor8.2}
Theorem {\rm \hyperref[XII.thm5.1]{5.1} (i)} is true.
\end{enonce*}

Indeed, it is enough to combine \hyperref[XII.lem8.1]{8.1}, criterion
\hyperref[XII.prop6.5]{6.5} (i) ${\rm b'')}$, and
\hyperref[XII.lem5.8]{5.8}.

\begin{enonce*}{Lemma 8.3}\etiquette[8.3]{XII.lem8.3}
To\pageoriginale prove any one of the assertions
{\rm \hyperref[XII.thm5.1]{5.1} (i), (ii), (iii)}, it is enough to do
so when $f$ is projective of relative dimension $\leq 1$ and $S$ is
Noetherian.
\end{enonce*}

Indeed, \hyperref[XII.lem8.1]{8.1} already shows that we may restrict
ourselves to the case in which $f$ is projective and $S$ is
Noetherian, and an argument used above further allows us to suppose
that $X=\PP^r_S$. We proceed by induction on $r$, noting that the
theorem holds by hypothesis for $r=1$. We may therefore suppose that
$r\geq2$ and that the theorem has been proved for
$\PP^{r-1}_S\rightarrow S$ (over every base $S$). Embed
$Z=\PP^{r-2}_S$ in $X=\PP^r_S$ in the usual way, and let
$p:\bar{X}\rightarrow X$ be the blowup of $X$ along $Z$. A direct
calculation (\cf \EGA V) gives a natural morphism
$p_1:\bar{X}\rightarrow X_1=\PP^1_S$ that makes $\bar{X}$ an
$X_1$-scheme locally isomorphic to $\PP^{r-1}_{X_1}$ (it is the
projective bundle over $X_1$ associated with a suitable locally free
sheaf of rank $r$ on $X_1$). On the other hand, the fibers of $p$ have
dimension $\leq1$, as do those of $f_1:X_1\rightarrow S$. We may
therefore apply the relative-dimension-one hypothesis to $p$ and
$f_1$, and the induction hypothesis to $p_1$, in the commutative
diagram
$$
\xymatrix{
   & \bar{X} \ar[ld]_{p} \ar[rd]^{p_1} &          \\
 X \ar[rd]_{f}&           & X_1\ar[ld]^{f_1}\\
   &    S      &
}$$
Since $p$ is surjective, \hyperref[XII.prop6.8]{6.8} reduces us to
proving \hyperref[XII.thm5.1]{5.1} for the composite
$\bar{f}=fp=f_1p_1$. But we may already apply
\hyperref[XII.thm5.1]{5.1} to both $f_1$ and $p_1$. Taking
\hyperref[XII.prop6.11]{6.11} into account, we conclude that
\hyperref[XII.thm5.1]{5.1} also holds for their composite $f_1p_1$.
This completes the proof of \hyperref[XII.lem8.3]{8.3}.

\begin{enonce*}{Lemma 8.4}\etiquette[8.4]{XII.lem8.4}
To\pageoriginale prove {\rm \hyperref[XII.thm5.1]{5.1} (ii)}, it is
enough to prove the ``specialization theorem for the fundamental
group'' {\rm \hyperref[XII.thmbis5.9]{5.9 bis}} in the case where $f$
is projective of relative dimension $\leq1$ and $S$ is Noetherian and
strictly local.
\end{enonce*}

This follows at once from \hyperref[XII.lem8.3]{8.3},
\hyperref[XII.lem6.1]{6.1}, and \hyperref[XII.prop6.5]{6.5} (ii)
(criterion b)).

The specialization theorem in the case of \hyperref[XII.lem8.4]{8.4}
will be proved directly in the next exposé
(\hyperrefext[XIII][XIII.sec2]{2}).

\begin{enonce*}{Lemma 8.5}\etiquette[8.5]{XII.lem8.5}
To prove {\rm \hyperref[XII.thm5.1]{5.1} (iii)}, it is enough to prove
{\rm \hyperref[XII.thm5.1]{5.1} (ii)} and, in addition, the following
proposition (which will be proved in the next exposé
{\rm (\hyperrefext[XIII][XIII.sec3]{3})}, or in
{\rm (\EGA IV 21.9.12)}):
\end{enonce*}

\begin{enonce*}{Proposition 8.6}\etiquette[8.6]{XII.prop8.6}
Let $S$ be Henselian and Noetherian, let $f:X\rightarrow S$ be
projective of relative dimension $\leq1$, and let $X_{\circ}$ be the
closed fiber of $X/S$. For every scheme $Z$, let
$\Pic Z=H^1(Z,(G_m)_Z)$ denote the Picard group of $Z$. The restriction
homomorphism
$$
\Pic X\longrightarrow \Pic X_{\circ}
$$
is surjective.
\end{enonce*}

\noindent
\emph{Proof of \hyperref[XII.lem8.5]{8.5}.} By
\hyperref[XII.prop6.5]{6.5} (iii) and
\hyperref[XII.cor8.2]{8.2}, it is enough to prove that
$$
\varphi^q:H^q(X,\ZZ/\ell^\nu \ZZ)\longrightarrow
H^q(X_{\circ},\ZZ/\ell^{\nu}\ZZ)
$$
is surjective for every $q\geq1$ and every prime $\ell$. For $q=1$,
surjectivity is already known from \hyperref[XII.thm5.1]{5.1} (ii).
Moreover, since $X_{\circ}$ is a projective scheme of dimension
$\leq1$ over a separably closed field $k$, we have
$$
H^q(X_{\circ},\ZZ/\ell^{\nu}\ZZ)=0
$$
for $q>2$ when $\ell\neq\operatorname{char} k$
(\hyperrefext[X][X.cor4.3]{4.3}), and for $q>1$ when
$\ell=\operatorname{char} k$ (\hyperrefext[X][X.cor5.2]{5.2}). Thus
surjectivity is automatic for these values of $q$.

It\pageoriginale remains to consider $q=2$ for a prime
$\ell\neq\operatorname{char} k$. Put $n=\ell^{\nu}$. The sheaf
$(\mu_n)_S$ of $n$-th roots of unity is locally isomorphic to
$(\ZZ/n\ZZ)_S$, and, since $S$ is strictly local, there is a
noncanonical isomorphism
$$
(\mu_n)_S\simeq(\ZZ/n\ZZ)_S,
$$
and hence $(\mu_n)_X\simeq(\ZZ/n\ZZ)_X$.

Kummer theory \hyperrefext[IX][IX.Thr3.2]{3.2}, applied to $X$ and
$X_{\circ}$, gives a morphism of exact sequences
$$
\xymatrix@R=1.5cm{
\Pic X \ar[r] \ar[d]^{\epsilon}&
H^{2}(X,(\mu_n)_X) \ar[d]^{\varphi^2}\ar[r]&
H^{2}(X,(G_m)_X)\ar[d]\\
\Pic X_{\circ} \ar[r]&
H^2(X_{\circ},(\mu_n)_{X_{\circ}}) \ar[r]&
H^2(X_{\circ},(G_m)_{X_{\circ}}).
}
$$
Now $H^2(X_{\circ},(G_m)_{X_{\circ}})=0$ by
\hyperrefext[IX][IX.thm4.6]{4.6}. Consequently, the surjectivity of
$\epsilon$ implies that of $\varphi^2$, which proves the lemma.

\begin{thebibliography}{99}
\bibitem{sga4.exp.XII.bib.1} \textsc{Giraud} J., \emph{Cohomologie non-abélienne},
thesis (forthcoming), and two notes in C.R., vol. 260, pp. 2392 and
2656 (1965).

\bibitem{sga4.exp.XII.bib.2} \textsc{Godement} R., \emph{Théorie des faisceaux}, Paris
(1958).

\bibitem{sga4.exp.XII.bib.3} \textsc{Verdier} J.L., \emph{Catégories dérivées},
Institut des Hautes Études Scientifiques (1964).
\end{thebibliography}

\ifx\danslelivre\undefined
\end{document}
\fi

% END INCLUDED SGA4_English_current_buildable_source_20260906/source/12/drafts/SGA4_XII_LINES_1925_2169_English.texfrag


% END INCLUDED SGA4_English_current_buildable_source_20260906/source/12/12_English.tex


\makeatletter
\def\input@path{{../13/}}
\makeatother
% BEGIN INCLUDED SGA4_English_current_buildable_source_20260906/source/13/13_English.tex
\ifx\danslelivre\undefined
% Release engine: XeLaTeX (the LuaTeX compatibility layer breaks Xy-pic output).
\documentclass[10pt]{../fichiers-configuration-latex/smfart}
% BEGIN INCLUDED SGA4_English_current_buildable_source_20260906/source/fichiers-configuration-latex/paquets.tex
\usepackage{polyglossia}
%\usepackage[francais]{babel}
\usepackage{fontspec}
\usepackage[math-style=ISO,bold-style=upright,sans-style=upright]{unicode-math}
%\usepackage[mathletters]{ucs}
%\usepackage[utf8]{inputenc}
%\usepackage[T1]{fontenc}
\usepackage[unicode,hypertexnames=false,colorlinks=true,linkcolor=blue,citecolor=Sepia,urlcolor=magenta]{hyperref}
\setmainlanguage{french}
\setmainfont[Ligatures=TeX]{Linux Libertine O}
% variante possible : TeXGyrePagella ?
\setmathfont{XITS Math}

% maths
\usepackage{amsmath}
%\usepackage{amssymb}
%\usepackage{mathrsfs}
%\usepackage{stmaryrd}
%\usepackage{euscript}

% couleurs
\usepackage[usenames,dvipsnames]{xcolor}

% graphiques
\usepackage{graphicx}
\usepackage[all]{xy}
%\usepackage{tikz}
%\usetikzlibrary{matrix,arrows,calc}
%\tikzset{column sep=3em,row sep=2em,text height=1.5ex, text depth=0.25ex,>=angle 90}
%\tikzset{description/.style={fill=white,inner sep=1pt}}

% hyperliens, 'etiquettes
\usepackage{xr-hyper}
\usepackage{../fichiers-configuration-latex/smfthm}
%\ifx\cacheretiquettes\undefined
%\usepackage[inline]{showlabels} \else \fi % pour rendre visibles les 'etiquettes
%\usepackage{../fichiers-configuration-latex/makeidx} % pour faire un index
%\usepackage[pagebackref,unicode]{hyperref}  % pour voir o`u sont cit'ees les r'ef'erences
% hyperref is loaded once above with the cumulative-reader destination policy.
%\hypersetup{breaklinks=true} %permet le retour `a la ligne dans les liens trop longs. Incompatible avec pdflatex
%\renewcommand*{\backref}[1]{{\small [cit\'e page(s)~#1.]}}

\usepackage{xspace}
\usepackage{pifont} % pour expose XVII

% END INCLUDED SGA4_English_current_buildable_source_20260906/source/fichiers-configuration-latex/paquets.tex

% BEGIN INCLUDED SGA4_English_current_buildable_source_20260906/source/fichiers-configuration-latex/numerotation.tex
\renewcommand{\thefootnote}{\roman{footnote}} % pour mettre les notes en bas de page sous la forme (ii) et non (2), qui ressemble trop (?) a un exposant.

\newcounter{pageoriginal}
\newcommand{\pageoriginale}{\refstepcounter{pageoriginal}\marginpar{\footnotesize\textbf{\thepageoriginal}}\xspace}
\let\pageoriginaled\pageoriginale

\usepackage{ifthen}
\newboolean{original}
\def\sisi{\ifthenelse{\boolean{original}}}
\setboolean{original}{false}

\newcounter{cnote}
\def\nde#1{\ifthenelse{\boolean{original}}{}{\stepcounter{cnote}\def\thefootnote{\arabic{cnote}}\protect\footnote[\value{cnote}]{N.D.E.: #1}}}
\def\ndemark{\ifthenelse{\boolean{original}}{}{\stepcounter{cnote}\def\thefootnote{\arabic{cnote}}\protect\footnotemark[\value{cnote}]}}
\def\ndetext#1{\ifthenelse{\boolean{original}}{}{\def\thefootnote{\arabic{cnote}}\protect\footnotetext[\value{cnote}]{N.D.E.: #1}}}

% END INCLUDED SGA4_English_current_buildable_source_20260906/source/fichiers-configuration-latex/numerotation.tex

% BEGIN INCLUDED SGA4_English_current_buildable_source_20260906/source/fichiers-configuration-latex/macros.tex
\setmainlanguage{french}
\setotherlanguage{english}


%% Fixed, engine-independent release identifier.  The inherited dynamic Git
%% probe used Lua io.popen, exposed build-state-dependent output, and printed
%% a visible question mark outside a Git checkout.
\def\showgitstatus{English release r2, 2026-07-23}


% macros pour LaTeX
\renewcommand{\mathsf}{}

% Police différente pour U+2113 SCRIPT SMALL L
\setmathfont[range="2113]{latinmodern-math.otf}

% Give every inherited semantic label its own stable PDF destination instead
% of inheriting whichever automatic equation/page anchor happened to be live.
% The legacy label remains as a compatibility alias used by existing edges.
\makeatletter
\def\etiquette[#1]#2{%
  \Hy@raisedlink{\hyper@anchorstart{sga4.semantic.#2}\hyper@anchorend}%
  \def\@currentHref{sga4.semantic.#2}%
  \def\@currentlabel{#1}%
  \label{#2}%
}
\makeatother
\def\eqhyperref[#1]#2{(\hyperref[#1]{#2})}
\def\pagehyperref[#1]#2{(page \hyperref[#1]{#2})}
\def\hyperrefext[#1][#2]#3{\hyperref[#2]{#1 #3}}
\def\eqhyperrefext[#1][#2]#3{\hyperref[#2]{#1 #3}}     %cf. exposes 17
\def\neweqhyperrefext[#1][#2]#3{\hyperref[#2]{#1 #3}}  %cf. exposes 17

\def\shat#1{\protect\widehat{#1}}
\def\stilde#1{\protect\widetilde{#1}}
%%%%%%%%%%%%%%%%%%


\def\etcn{etc.\xspace}
\def\etco{etc...\xspace}
\def\etc{\sisi{\etco}{\etcn}}
\def\ldotsn{\sisi{...}{\etcn}}

\def\cf{{cf.}\kern.3em}
\def\Cf{{Cf.}\kern.3em}
\def\ie{{i.e.}\kern.3em}
\def\iev{{i.e.},\xspace}
\def\loccito{loc.\kern2pt cit.\xspace}
\def\loccitn{\textit{loc.\kern2pt cit.}\xspace}
\def\loccit{\sisi{\loccito}{\loccitn}}
\def\EGA{\textup{EGA}\kern.2em}
\def\SGA{\textup{SGA}\kern.3em}

\def\quoi{}

\def\quav{\quad,\quad}
\def\quaet{\quad \text{et}\quad}
\def\fibreproduct#1#2#3{#1 ×_{#3} #2}
\let\fprod\fibreproduct
\def\ccheck#1#2{\mbox{${\check{\check{\rm #1}}}^{\lower2pt\hbox{$\scriptstyle #2$}}$}}
\def\sheaf#1{\protect\mathcal{#1}}
\def\vieuxsheaf#1{{#1}} % cf. problème avec sheaf{a\Phi_i} par exemple...
\DeclareMathOperator{\Hom}{Hom}
\DeclareMathOperator{\can}{can}
\DeclareMathOperator{\Ind}{Ind}
\DeclareMathOperator{\ind}{ind}
%\DeclareMathOperator{\Hhomtop}{\mathcal{H}\!\mathit{omtop}}
\DeclareMathOperator{\Hhomtop}{\mathbf{Homtop}}
\DeclareMathOperator{\Point}{Point}
\DeclareMathOperator{\Pro}{Pro}
\DeclareMathOperator{\Top}{Top}
\DeclareMathOperator{\Mod}{Mod}
\DeclareMathOperator{\Cart}{Cart}
%\DeclareMathOperator{\Mmorsite}{\mathcal{M}\!\mathit{orsite}}
\DeclareMathOperator{\Mmorsite}{\mathbf{Morsite}}
\DeclareMathOperator{\loc}{loc}
\DeclareMathOperator{\Loc}{Loc}
\DeclareMathOperator{\Pic}{Pic}
\DeclareMathOperator{\Et}{Et}
\DeclareMathOperator{\card}{card}
\DeclareMathOperator{\Ens}{Ens}
\DeclareMathOperator{\h}{\underline{H}}
\DeclareMathOperator{\Fonct}{Fonct}
\DeclareMathOperator{\Ker}{Ker}
\DeclareMathOperator{\but}{but}
\DeclareMathOperator{\Ob}{Ob}
\DeclareMathOperator{\ob}{ob}
\DeclareMathOperator{\Lib}{Lib}
\DeclareMathOperator{\esj}{esj}
\DeclareMathOperator{\id}{id}
\DeclareMathOperator{\Id}{Id}
\DeclareMathOperator{\ad}{ad}
\DeclareMathOperator{\Iim}{Im}
\DeclareMathOperator{\pr}{pr}
\DeclareMathOperator{\Cat}{Cat}
\DeclareMathOperator{\car}{car}
\DeclareMathOperator{\SysInd}{SysInd}
\DeclareMathOperator{\Sup}{Sup}
\DeclareMathOperator{\Coker}{Coker}
\DeclareMathOperator{\Filt}{Filt}
\DeclareMathOperator{\acc}{acc}
\DeclareMathOperator{\sur}{sur}
\DeclareMathOperator{\cat}{cat}
\DeclareMathOperator{\Ord}{Ord}
\DeclareMathOperator{\Gr}{Gr}
\DeclareMathOperator{\Cov}{Cov}
\DeclareMathOperator{\Caf}{Caf}
\DeclareMathOperator{\coker}{coker}
\DeclareMathOperator{\coim}{coim}
\DeclareMathOperator{\im}{im}
\DeclareMathOperator{\Ouv}{Ouv}
\DeclareMathOperator{\diag}{diag}
\DeclareMathOperator{\Esp}{Esp}
\DeclareMathOperator{\TOP}{TOP}
\DeclareMathOperator{\Unif}{Unif}
\DeclareMathOperator{\Iind}{\mathbf{Ind}}
\DeclareMathOperator{\Hhomcart}{\mathbf{Homcart}}
\DeclareMathOperator{\Hhomtopess}{\mathbf{Homtopess}}
\DeclareMathOperator{\Ppoints}{\mathbf{Points}}
\DeclareMathOperator{\Ppoint}{\mathbf{Point}}
\DeclareMathOperator{\Ttors}{\mathbf{Tors}}
\DeclareMathOperator{\Ssex}{\mathbf{Sex}}
\DeclareMathOperator{\Ffib}{\mathbf{Fib}}
\DeclareMathOperator{\Vvois}{\mathbf{Vois}}
\DeclareMathOperator{\Ppointess}{\mathbf{Pointess}}
\DeclareMathOperator{\scrHom}{\mathbf{Hom}}
\DeclareMathOperator{\sheafPic}{\mathbf{Pic}}
\DeclareMathOperator{\sheafAut}{\mathbf{Aut}}
\DeclareMathOperator{\Hhomtoplocan}{\mathbf{Homtoplocan}}
\DeclareMathOperator{\Hhomtopan}{\mathbf{Homtopan}}
\DeclareMathOperator{\Pptgeom}{\mathbf{Ptgéom}}
\DeclareMathOperator{\Morsite}{Morsite}
\DeclareMathOperator{\Fib}{Fib}
\DeclareMathOperator{\Fais}{Fais}
\DeclareMathOperator{\kar}{kar}
\DeclareMathOperator{\Groupes}{Groupes}
\DeclareMathOperator{\Restr}{Restr}
\DeclareMathOperator{\Isom}{Isom}
\DeclareMathOperator{\Prol}{Prol}
\DeclareMathOperator{\Res}{Res}
\DeclareMathOperator{\Tors}{Tors}
\DeclareMathOperator{\esp}{esp}
\DeclareMathOperator{\es}{es}
\DeclareMathOperator{\od}{od}
\DeclareMathOperator{\gr}{gr}
\DeclareMathOperator{\sob}{sob}
\DeclareMathOperator{\dfn}{dfn}
\DeclareMathOperator{\Sch}{Sch}
\DeclareMathOperator{\ouv}{ouv}
\DeclareMathOperator{\cons}{cons}
\DeclareMathOperator{\Nerf}{Nerf}
\DeclareMathOperator{\Points}{Points}
\DeclareMathOperator{\Espsob}{Espsob}
\DeclareMathOperator{\Homtop}{Homtop}
\DeclareMathOperator{\Pt}{Pt}
\DeclareMathOperator{\Quot}{Quot}
\DeclareMathOperator{\glob}{glob}
\DeclareMathOperator{\Homfib}{Homfib}
\DeclareMathOperator{\Qouv}{Qouv}
\DeclareMathOperator{\Quotouv}{Quotouv}
\DeclareMathOperator{\ens}{ens}
\DeclareMathOperator{\ab}{ab}
\DeclareMathOperator{\an}{an}
\DeclareMathOperator{\fs}{fs}
\DeclareMathOperator{\red}{red}
\DeclareMathOperator{\Ab}{Ab}
\DeclareMathOperator{\cosq}{cosq}
\DeclareMathOperator{\Ann}{Ann}
\DeclareMathOperator{\Zar}{Zar}


%%%%%%%%%%%%
%%%%% SGA4-2
\DeclareMathOperator{\Sex}{Sex}
\DeclareMathOperator{\vois}{vois}
\DeclareMathOperator{\Vois}{Vois}
\DeclareMathOperator{\Ext}{Ext}
\DeclareMathOperator{\ext}{ext}
\DeclareMathOperator{\Fer}{Fer}
\DeclareMathOperator{\Add}{Add}
\DeclareMathOperator{\Cocart}{Cocart}
\DeclareMathOperator{\cosk}{cosk}
\DeclareMathOperator{\plat}{plat}

\DeclareMathOperator{\faisc}{faisc}
\DeclareMathOperator{\sk}{sk}
\DeclareMathOperator{\ch}{ch}

\def\HH{\mathbf{H}}
\def\qc{\mathrm{qc}}
\def\et{\mathrm{et}}
\def\coh{\mathrm{coh}}

\def\Hvr#1{%
\begin{array}[b]{l}
{\kern2pt\scriptstyle r}\\[-8pt]
{\kern1pt\scriptstyle \vee}\\[-5pt]
H^{#1}
\end{array}}

\def\Hvv#1{%
\begin{array}[b]{l}
\kern2pt\rotatebox{90}{${\scriptscriptstyle\ll}$}\\[-6pt]
H^{#1}
\end{array}}

\def\Hvi#1{%
\begin{array}[b]{l}
{\kern0.2pt\scriptscriptstyle \infty}\\[-8pt]
{\kern1pt\scriptstyle \vee}\\[-5pt]
H^{#1}
\end{array}}

\def\Hv#1{%
\begin{array}[b]{l}
{\kern1pt\scriptstyle \vee}\\[-5pt]
H^{#1}
\end{array}}

\DeclareMathOperator{\cart}{cart}
\DeclareMathOperator{\Ffaisc}{\mathbf{Faisc}}


\def\Ev#1{%
\begin{array}[b]{l}
{\kern1pt\scriptstyle \vee}\\[-5pt]
E'_{#1}
\end{array}}

\DeclareMathOperator{\Cont}{Cont}
\DeclareMathOperator{\Ccont}{\mathscr{C}\!\mathit{ont}}
\DeclareMathOperator{\Site}{Site}
\newcommand\MM{\mathbf{M}}
\def\limtop#1{{\displaystyle{\varprojlim_{#1}\!\mbox{\rm top }}}}
\def\Limtop{\mathop{\varprojlim\!\mbox{\rm top}}}
\DeclareMathOperator{\spec}{spec}
\def\ss{\scriptstyle{\mathscr{S}}}
\DeclareMathOperator{\Gal}{Gal}
\def\smallm{\scriptstyle{\mathscr{M}}}
\DeclareMathOperator{\End}{End}
\DeclareMathOperator{\ass}{ass}
\DeclareMathOperator{\zar}{zar}
\DeclareMathOperator{\Ogg}{Ogg}
\DeclareMathOperator{\res}{res}
\DeclareMathOperator{\tr}{tr}
\DeclareMathOperator{\Supp}{Supp}
\DeclareMathOperator{\codim}{codim}
\DeclareMathOperator{\rk}{rk}
\DeclareMathOperator{\cur}{cur}
\DeclareMathOperator{\cdqc}{cdqc}
\newcommand\DD{\mathbf{D}}
\DeclareMathOperator{\cl}{cl}
\def\JJ{\mathbf{J}}
\def\FF{\mathbf{F}}
\def\BB{\mathbf{B}}
\DeclareMathOperator{\Alget}{Alget}
\DeclareMathOperator{\Ass}{Ass}
\DeclareMathOperator{\supp}{supp}
\DeclareMathOperator{\rad}{rad}
\DeclareMathOperator{\prof}{prof}
\DeclareMathOperator{\cint}{int}
\DeclareMathOperator{\sic}{sic}
\DeclareMathOperator{\Tor}{Tor}
\DeclareMathOperator{\tors}{tors}
\DeclareMathOperator{\Sym}{Sym}
\DeclareMathOperator{\parf}{parf}
\DeclareMathOperator{\tor}{tor}
\DeclareMathOperator{\defn}{defn}
\DeclareMathOperator{\Div}{Div}
\DeclareMathOperator{\Pol}{Pol}
\DeclareMathOperator{\Fex}{Fex}
\DeclareMathOperator{\Lie}{Lie}
\DeclareMathOperator{\tot}{tot}
\DeclareMathOperator{\Aut}{Aut}
\DeclareMathOperator{\HOM}{HOM}
\DeclareMathOperator{\pch}{pch}
\DeclareMathOperator{\Ch}{Ch}
\DeclareMathOperator{\EXT}{EXT}
\DeclareMathOperator{\PIC}{PIC}
\DeclareMathOperator{\Rf}{Rf}
\DeclareMathOperator{\inj}{inj}
\newcommand\VV{\mathbf{V}}

\def\Xv{%
\begin{array}[b]{l}
{\kern1pt\scriptstyle \vee}\\[-5pt]
X
\end{array}}

%%%%%%%
\newcommand\TT{\mathbf{T}}
\newcommand\XX{\mathbf{X}}
\newcommand\EE{\mathbf{E}}
%%%%%%

\def\qedsymbol{\textup{C.Q.F.D.}}
\let\CQFD\qed
\let\cqfd\qed

%%%%%\bigdot
\font\bigsymb=cmsy10 at 4pt
\def\bigdot{{\kern1.2pt\raise 1.5pt\hbox{\bigsymb\char15}}}
%%%%%

\def\Lright{\mathop{\vcenter{\hbox{\hbox{$\longrightarrow$}\lower8pt \hbox{\kern-16pt\hbox{$\longrightarrow$}}}}}}

\def\Aa{{\mathcal A}}
\def\Bb{{\mathcal B}}
\def\Cc{{\mathcal C}}
\def\Dd{{\mathcal D}}
\def\Ee{{\mathcal E}}
\def\Ff{{\mathcal F}}
\def\Gg{{\mathcal G}}
\def\Hh{{\mathcal H}}
\def\Ii{{\mathcal I}}
\def\Jj{{\mathcal J}}
\def\Kk{{\mathcal K}}
\def\Ll{{\mathcal L}}
\def\Mm{{\mathcal M}}
\def\Nn{{\mathcal N}}
\def\Oo{{\mathcal O}}
\def\Pp{{\mathcal P}}
\def\Qq{{\mathcal Q}}
\def\Rr{{\mathcal R}}
\def\Ss{{\mathcal S}}
\def\Tt{{\mathcal T}}
\def\Uu{{\mathcal U}}
\def\Vv{{\mathcal V}}
\def\Ww{{\mathcal W}}
\def\Xx{{\mathcal X}}
\def\Yy{{\mathcal Y}}
\def\Zz{{\mathcal Z}}

\newcommand{\CC}{\mathbf{C}}
\newcommand{\GG}{\mathbf{G}}
\newcommand{\KK}{\mathbf{K}}
\newcommand{\LL}{\mathbf{L}}
\newcommand{\NN}{\mathbf{N}}
\newcommand{\PP}{\mathbf{P}}
\newcommand\QQ{\mathbf{Q}}
\newcommand\RR{\mathbf{R}}
\newcommand\ZZ{\mathbf{Z}}

\def\to{\mathchoice{\longrightarrow}{\rightarrow}{\rightarrow}{\rightarrow}}
\def\from{\mathchoice{\longleftarrow}{\leftarrow}{\leftarrow}{\leftarrow}}
\def\mto{\mathchoice{\longmapsto}{\mapsto}{\mapsto}{\mapsto}}
\def\mfrom{\mathrel{\from\joinrel\mapstochar\,}}
\def\hto{\mathrel{\lhook\joinrel\to}}
\def\hfrom{\mathrel{\from\joinrel\rhook}}

\def\implies{\mathchoice{\Longrightarrow}{\Rightarrow}{\Rightarrow}{\Rightarrow}}
\def\implied{\mathchoice{\Longleftarrow}{\Leftarrow}{\Leftarrow}{\Leftarrow}}
\def\iff{\mathchoice{\Longleftrightarrow}{\Leftrightarrow}{\Leftrightarrow}{\Leftrightarrow}}

\def\To#1{\mathchoice{\xrightarrow{\labelstyle~~#1~}}{\stackrel{#1}{\longrightarrow}}{}{}}
\def\Mto#1{\mathrel{\mapstochar\To{#1}}}
\def\From#1{\mathchoice{\xleftarrow{\labelstyle~~#1~}}{\stackrel{#1}{\longleftarrow}}{}{}}

\newcommand{\isomto}{\overset{\sim}{\longrightarrow}}
\newcommand{\isomfrom}{\overset{\sim}{\longleftarrow}}

\newcommand{\IFF}{\unskip~$\iff$~\ignorespaces}
\newcommand{\THEN}{\unskip~$\implies$~\ignorespaces}

\let\leq\leqslant
\let\geq\geqslant
\let\le\leqslant
\let\ge\geqslant

\def\cf{{cf.}\kern.3em}
\def\Cf{{Cf.}\kern.3em}
\def\ie{{i.e.}\kern.3em}
\def\iev{{i.e.},\xspace}
%\def\loccito{loc.\kern2pt cit.\xspace}
%\def\loccitn{\textit{loc.\kern2pt cit.}\xspace}
%\def\loccit{\sisi{\loccito}{\loccitn}}
\def\EGA{\textup{EGA}\kern.2em}
\def\SGA{\textup{SGA}\kern.3em}

\let\oldS\S
\def\S{\oldS\kern .15em}
\def\ptbl{.\kern .2em}
\def\resp{resp.\kern.3em}
\def\Exp{Exp.\kern.3em}

\let\bar\overline
\let\hat\widehat
\let\tilde\widetilde
\let\dpl\displaystyle
\let\ldots\dots
\let\emptyset\varnothing
\let\skpt\skippointrait
\let\epsilon\varepsilon

\def\sfrac#1#2{{#1}/{#2}}

\DeclareMathOperator{\cd}{cd}
%\DeclareMathOperator{\Hhom}{\mathcal{H}\!\mathit{om}}
\DeclareMathOperator{\Hhom}{\mathbf{Hom}}
%\DeclareMathOperator{\Eext}{\mathcal{E}\!\mathit{xt}}
\DeclareMathOperator{\Eext}{\mathbf{Ext}}
\DeclareMathOperator{\Spec}{Spec}
\DeclareMathOperator{\Tr}{Tr}

\let\emptyset=\varnothing

%%%%%%%%%%%%%%%
%%% expose XVII
%%%%%%%%%%%%%%%

\makeatletter
\let\old@optexpose\optexpose
\newcommand{\empty@optexpose}[1]{}
\newcommand{\paren@optexpose}[1]{\old@optexpose{#1}(}
\newcommand{\refnoexpose}[1]{\let\optexpose\empty@optexpose\ref{#1}\let\optexpose\old@optexpose}%\refnoexpose{X.cor5.2} engendre "5.2" au lieu de "X 5.2"
\newcommand{\neweqref}[1]{\let\optexpose\paren@optexpose\textup{\ref{#1})}\let\optexpose\old@optexpose}%\neweqref{XII.eq4.1} engendre "XII (4.1)" au lieu de "(XII 4.1)"
\makeatother
\newcommand{\refvb}[2][VI${}_{\text{B}}$ ]{\sisi{#1\refnoexpose{#2}}{\ref{#2}}}

\newcommand{\leftexp}[2]{{\vphantom{#2}}^{#1}{#2}}
\newcommand{\AAA}{\mathbf{A}}
\newcommand{\Sf}{\mathfrak{S}}
\DeclareMathOperator*{\quotprojlim}{\text{``$\varprojlim$''}}
\DeclareMathOperator*{\quotinjlim}{\text{``$\varinjlim$''}}
\DeclareMathOperator{\Fl}{F\sisi{\ell}{l}}
\DeclareMathOperator{\Pl}{P\sisi{\ell}{l}}
%\DeclareMathOperator{\cl}{c\sisi{\ell}{l}}
\DeclareMathOperator{\tdf}{tdf}
\DeclareMathOperator{\rg}{rg}
\newcommand{\iInd}{\sisi{I}{i}nd}
\newcommand{\iind}{ind}
\newcommand{\ppro}{pro}
\newcommand{\tTor}{\sisi{T}{t}or}
\newcommand{\ttor}{tor}
\newcommand{\chap}{Chap.}
\newcommand{\Ffex}{\Ff\mathit{ex}}
\newcommand{\sheafk}{\underline{k}}
\newcommand{\sheafn}{\underline{n}}
\newcommand{\sheafGamma}{\underline{\Gamma}}

% Repeated references to an already typeset footnote must not use
% \footnotemark without a matching \footnotetext: Hyperref would then create
% a live annotation to a destination that cannot exist.  This helper keeps the
% current printed footnote mark and links it to an explicit stable note target.
\makeatletter
\newcommand{\sgarepeatedfootnotemark}[1]{%
  \begingroup
  \edef\@thefnmark{\thefootnote}%
  \hyperlink{#1}{\@makefnmark}%
  \endgroup}
\makeatother

\sisi{
\DeclareMathOperator{\simarrow}{\rule{.15mm}{1.8mm}\kern -3.4pt\rightsquigarrow}
\newcommand{\sheafell}{\underline{\ell}}
\newcommand{\sheafZZ}{\underline{\ZZ}}
}{
\newcommand{\nderef}[1]{la N.D.E.~(\ref{#1}) page~\pageref{#1}}
}

%%%%%%%%%%%%%%%%%
%%%%% expose XVII
%%%%%%%%%%%%%%%%%

% END INCLUDED SGA4_English_current_buildable_source_20260906/source/fichiers-configuration-latex/macros.tex


\externaldocument{../01/01_English.tex}
\externaldocument{../02/02_English.tex}
\externaldocument{../03/03_English.tex}
\externaldocument{../04/04_English.tex}
\externaldocument{../05/05_English.tex}
\externaldocument{../05bis/05bis_English.tex}
\externaldocument{../06/06_English.tex}
\externaldocument{../07/07_English.tex}
\externaldocument{../08/08_English.tex}
\externaldocument{../09/09_English.tex}
\externaldocument{../10/10_English.tex}
\externaldocument{../11/11_English.tex}
\externaldocument{../12/12_English.tex}

\setboolean{original}{false}

\begin{document}

\begin{center}
XIII. Base-change theorem for a proper morphism: conclusion of the proof\\
M. Artin
\end{center}

\selectlanguage{english}version: \showgitstatus
\tableofcontents

\else
\chapter{Base-change theorem for a proper morphism: conclusion of the proof}
%\addtocontents{toc}{M. Artin}
\begin{center}
M. Artin
\end{center}
\fi
\phantomsection\MakeLinkTarget*{sga4.exp.XIII}\label{sga4.exp.XIII}

% BEGIN INCLUDED SGA4_English_current_buildable_source_20260906/source/13/drafts/SGA4_XIII_LINES_43_316_English.texfrag
\setcounter{pageoriginal}{131}

\section{The projective and flat case}\etiquette[1]{XIII.sec1}\pageoriginale

Recall the statement of \hyperrefext[XII][XII.thm5.9]{5.9 bis} in the case
under consideration here:

\begin{enonce*}{Proposition 1.1}\etiquette[1.1]{XIII.prop1.1}
Let $S$ be the spectrum of a Noetherian henselian ring, and let
$f:X\rightarrow S$ be a projective and flat morphism. Let $X_{\circ}$ be
the closed fiber of $X/S$. The restriction functor
$$
\Et(X)\longrightarrow \Et(X_{\circ})
$$
is an equivalence of categories.
\end{enonce*}

As we have already observed, it follows from
\hyperrefext[XII][XII.prop6.5]{6.5} (i) that the functor is fully faithful.
It remains to prove that every étale covering $Y_{\circ}$ of $X_{\circ}$ is
induced by an étale covering $Y$ of $X$.

\begin{enonce*}{Lemma 1.2}\etiquette[1.2]{XIII.lem1.2}
Let $S$ be locally Noetherian, let $f:X\rightarrow S$ be projective and
flat, and let $Y\subset X$ be a subscheme. The subfunctor
$\epsilon=\prod_{X/S}Y/X$ of the terminal functor, \lq\lq which expresses
the condition $Y=X$\rq\rq, that is, such that for $S'\rightarrow S$ one has
$$
\epsilon(S')=
\begin{cases}
\phi & \mbox{if }\quad Y'\neq X',\\
\{\phi\} &\mbox{if }\quad Y'=X',
\end{cases}
$$
(where\pageoriginale a prime denotes the effect of the base change
$S'\rightarrow S$), is representable by a subscheme $Z$ of $S$. If
$Y$ is open (\resp closed), then so is $Z$.
\end{enonce*}

\noindent
\emph{Proof}. Let $Y$ be closed in an open subset $U$ of $X$, and let
$C$ be the closed complement of $U$. Since $f$ is proper, the image of
$C$ in $S$ is a closed subset $D$, and it is clear that the condition
$U=X$ is represented by the open subset $S-D$ of $S$. We may therefore
suppose that $U=X$ already, that is, that $Y$ is a closed subscheme.
Consider the surjective morphism
$\Oo_X\xrightarrow{\alpha}\Oo_Y$ of $\Oo_X$-modules; saying that $Y=X$
amounts to saying that there exists a morphism
$\Oo_Y\xrightarrow{\beta}\Oo_X$ such that the composite $\beta\alpha$ is
the identity, that is, that there exists a section of
$\Hhom(\Oo_Y,\Oo_X)$ lifting the identity section of
$\Hhom(\Oo_X,\Oo_X)$. Now it follows from \EGA III 7.7.8 that the functor
which assigns to $S'/S$ the set of sections of
$\Hhom(\Oo_Y,\Oo_X)$ (\resp $\Hhom(\Oo_X,\Oo_X)$) is represented by a
vector bundle $V_1$ (\resp $V_2$) over $S$. Let
$\varphi:V_1\rightarrow V_2$ be the morphism $u\mapsto u\alpha$ induced
by $\Oo_X\rightarrow\Oo_Y$, and let $s:S\rightarrow V_2$ be the section
corresponding to the identity $\Oo_X\rightarrow\Oo_X$. The inverse image
of the section $s$ under $\varphi$ is identified with a closed subscheme
$Z$ of $S$, and it is clear that $Z$ represents the functor $\epsilon$,
which proves the lemma.

\begin{enonce*}{Lemma 1.3}\etiquette[1.3]{XIII.lem1.3}
Let $S$ be a locally Noetherian scheme, let $f:X\rightarrow S$ be a
projective and flat morphism, and let $E$ be a locally free sheaf on $X$.
Let
$$
QE:(\Sch)/S\longrightarrow(\Ens)
$$
be the functor which assigns to $S'/S$ the set $QE(S')$ of quotient
sheaves $A'$ of $E'=E\otimes_{\Oo_S}\Oo_{S'}$, endowed with the structure
of an étale $\Oo_{X'}$-algebra. This functor is representable by a scheme
locally of finite type over $S$.
\end{enonce*}

\noindent
\emph{Proof}.\pageoriginale The functor $\sheaf{\Quot}(E)$, with
$\sheaf{\Quot}(E)(S')$ equal to the set of quotients of $E'$ that are flat
over $S'$, is representable by a scheme locally of finite type (TDTE IV,
exp.~221), say by $Z/S$. Let $M_u$ be the \lq\lq universal\rq\rq quotient
of $E\otimes_{\Oo_S}\Oo_Z$ over $X_Z$. The subfunctor of
$\sheaf{\Quot}(E)$ consisting of the quotients which are flat
\emph{over} $X'$ is represented by an open subset $U$ of $Z$. Indeed,
let $M'$ be a quotient of $E'$, given by an $S$-morphism $S'\rightarrow
Z$. Since $M'$ is flat over $S'$, saying that it is flat over $X'$ amounts
to saying that, for every point $s'\in S'$, the module $M'_{s'}$ is flat
over the fiber $X'_{s'}$ (\EGA IV 5.9), which is equivalent to
$M_u|_{X_z}$ being flat over the fiber $X_z$ at the point $z$ which is
the image of $s'$. But the set $U$ of $z\in Z$ such that
$M_u|_{X_z}$ is flat over $X_z$, that is, such that $M_u$ is flat over
$X_Z$ at every point of $X_z$, is open in $Z$ by \EGA IV 11.1.5,
proving the assertion.

It now suffices to represent the relative functor $QE/U[f]$. We are thus
reduced to the following lemma, after replacing $S$ by $U$ and $X$ by
$\fprod{X}{U}{S}$:

\begin{enonce*}{Lemma 1.4}\etiquette[1.4]{XIII.lem1.4}
Let $S$ be a locally Noetherian scheme, let $f:X\rightarrow S$ be a
projective and flat morphism, and let $M$ be a locally free
$\Oo_X$-module. The functor
$\Alget:(\Sch)/S\rightarrow(\Ens)$, which assigns to $S'/S$ the set
$\Alget(S')$ of algebra structures on $M'$ that are étale over
$\Oo_{X'}$, is representable by a scheme of finite type over $S$.
\end{enonce*}

\noindent
\emph{Proof}. For brevity, we shall say \lq\lq representable\rq\rq,
understanding \lq\lq by a relative scheme of \emph{finite type}\rq\rq.
All the vector bundles encountered below will be defined by
\emph{coherent} sheaves, and hence will be of finite type. A bilinear
composition law on $M'$ is given by a section of the locally free sheaf
$\Hhom(M'\otimes M',M')$. Since
$\Hhom(M'\otimes M',M')$ is flat over $S$, it follows from\pageoriginale
\EGA III 7.7.6 that the functor \lq\lq bilinear composition law\rq\rq is
represented by a vector bundle $Z$ over $S$. We must check that the
additional conditions that the law be associative, commutative, and
possess a unit section are represented by a subscheme of $Z$. By
considering the relative functor, we may replace $S$ by $Z$, and may
therefore suppose that a composition law on $M$ has been given and prove
that the subfunctor of $S$ expressing the preceding conditions is
representable by a subscheme of $S$.

For associativity, the condition is that two sections of
$\Hhom(M\otimes M\otimes M,M)$ become equal after the base change
$S'\rightarrow S$. Let $Y$ be the fiber product of these sections over
the vector bundle corresponding to
$\Hhom(M\otimes M\otimes M,M)$. The scheme $Y$ is identified with a
closed subscheme of $X$, and associativity after $S'\rightarrow S$ is
expressed by the condition $Y'=X'$. It therefore defines a subfunctor of
$S$ represented by a closed subscheme of $S$, by
\hyperref[XIII.lem1.2]{1.2}. For commutativity, the condition is that two
sections of $\Hhom(M\otimes M,M)$ become equal, which again defines a
subfunctor of $S$ represented by a closed subscheme, by
\hyperref[XIII.lem1.2]{1.2}. Finally, suppose, as we now may, that the law
is already associative and commutative, and consider the condition that a
unit section exist; such a section will then be unique. The functor whose
value on $S'/S$ is the set of sections of $M'$ is represented by a vector
bundle, by \EGA III 7.7.6, and it remains only to represent the subfunctor
of this bundle corresponding to the \emph{unit sections} (taking account
of the uniqueness just recalled). Passing to the relative situation as
usual, we may suppose that a section $\epsilon$ of $M$ has already been
given, and prove that the subfunctor of $S$ expressing the condition that
it become a\pageoriginale unit section is representable. But the condition
that $\epsilon$ be a unit is that the composite of the evident morphisms
$M'\rightarrow M'\otimes M'\rightarrow M'$ be the identity, that is, that
two specified sections of $\Hom(M,M)$ become equal. We conclude once
again by \hyperref[XIII.lem1.2]{1.2}.

Finally, suppose that a commutative unital algebra structure has been
given on the locally free module $M$. The locus where it is étale is open
in $X$ (\SGA I 4.5), and is therefore represented over $S$ by an (open)
subscheme. This completes the proof of \hyperref[XIII.lem1.4]{1.4}, and
hence of \hyperref[XIII.lem1.3]{1.3}.

\setcounter{subsection}{4}
\subsection{}\etiquette[1.5]{XIII.subsec1.5}
Let $Z/S$ now be the object representing the functor considered in
\hyperref[XIII.lem1.3]{1.3}, and consider the condition that $Z$ be smooth
over $S$ at a point $z$. Translating \SGA I III 3.1 into the language of
the functor $QE$, one obtains the following condition:

Let $B$ be a local Artinian ring, let $\Mm=\rad B$, let $J$ be an ideal of
$B$ such that $J\Mm=0$, put $S'=\Spec B$ and $S''=\Spec B/J$, and let
$S'\rightarrow S$ be a morphism. We use the evident notation for the
effects of the base changes $S'\rightarrow S$ and $S''\rightarrow S$.

\begin{itemize}
\item[$(\ast)$] For every quotient $A''$ of $E''$, endowed with the
structure of an étale $\Oo_{X''}$-algebra and corresponding to an
$S$-morphism $S''\rightarrow Z$ whose image is $z$, there exists a
quotient $A'$ of $E'$, endowed with the structure of an étale
$\Oo_{X'}$-algebra, which induces it.
\end{itemize}

Now, since $X'$ and $X''$ have the same underlying space, an algebra $A'$
inducing $A''$ is already determined up to unique isomorphism
(\SGA I 8.3). We\pageoriginale may therefore express condition
$({\ast}{\ast})$ as follows:

\begin{itemize}
\item[$({\ast}{\ast})$] Let $A'$ be an étale algebra over $\Oo_{X'}$,
and let $A''$ be the algebra induced by $A'$ on $X''$. Every surjective
homomorphism
$$
E''\xrightarrow{\varphi''}A''
$$
corresponding to a morphism $S''\rightarrow Z$ whose image is $z$ lifts
to a homomorphism (necessarily surjective by Nakayama's lemma)
$$
E'\xrightarrow{\varphi'}A'{\quoi}.
$$
\end{itemize}

From the exact sequence
$$
0\longrightarrow J\longrightarrow B\longrightarrow B/J\longrightarrow
0
$$
we obtain an exact sequence of sheaves on $X'$:
$$
0\longrightarrow J\otimes\Hhom(E',A')\longrightarrow
\Hhom(E',A')\longrightarrow\Hhom(E'',A'')\longrightarrow
0{\quoi}.
$$
The obstruction to lifting $\varphi''$ lies in
$H^1(X',J\otimes\Hhom(E',A'))$. Let
$S_{\circ}=\Spec B/\Mm$. Since $J\Mm=0$, one has
$$
J\otimes\Hhom(E',A')\simeq
J\otimes\Hhom(E_{\circ},A_{\circ}){\quoi},
$$
and it is clear that
$$
H^1(X',J\otimes\Hhom(E_{\circ},A_{\circ}))\simeq
J\otimes H^1(X_{\circ},\Hhom(E_{\circ},A_{\circ})){\quoi}.
$$
Taking account of the fact that cohomology commutes with extension of the
base fields, we obtain:

\begin{enonce*}{Corollary 1.6}\etiquette[1.6]{XIII.cor1.6}
Let $Z/S$ be the scheme representing the functor of
\hyperref[XIII.lem1.3]{1.3}, let
$X_Z=\fprod{X}{Z}{S}$, and let $A_Z$ be the (\lq\lq general\rq\rq)
quotient étale algebra of $E_Z$. Then $Z/S$ is smooth at every point
$z_{\circ}$ such that
$$
H^1(X_{\circ},\Hhom(E_{\circ},A_{\circ}))=0{\quoi},
$$
where\pageoriginale the subscript zero denotes restriction to the fiber
$X_{\circ}$ of $X_Z/Z$ at the point $z_{\circ}$.
\end{enonce*}

We can now complete the proof of \hyperref[XIII.prop1.1]{1.1}. With the
notation of the proposition, let $Y_{\circ}/X_{\circ}$ be an étale
covering, given by a sheaf of algebras $A_{\circ}$ that is étale over
$\Oo_{X_{\circ}}$. For $n$ sufficiently large, $A_{\circ}(n)$ is
generated by its sections (\EGA III 2.2.1). We thus have a surjection
$$
\Oo^N_{X_{\circ}}\longrightarrow A_{\circ}(n)\longrightarrow
0{\quoi},
$$
and hence a surjection
$$
\Oo^N_{X_{\circ}}(-n)=E_{\circ}\longrightarrow
A_{\circ}\longrightarrow 0{\quoi}.
$$
Choose $n$ sufficiently large that we also have
\begin{equation*}
H^1(X_{\circ},\Hhom(E_{\circ},A_{\circ}))=0{\quoi},\tag{$+$}
\end{equation*}
which is possible by \EGA III 2.2.1, since
$$
\Hhom(E_{\circ},A_{\circ})
=[\Hhom(\Oo^N_{X_{\circ}},A_{\circ})](n)
=[A_{\circ}(n)]^N{\quoi}.
$$
Put $E=\Oo^N_X(-n)$, and let $Z/S$ be the scheme representing the functor
considered in \hyperref[XIII.lem1.3]{1.3}. The quotient $A_{\circ}$ of
$E_{\circ}$ corresponds to a point $z_{\circ}\in Z$ in the closed fiber
of $Z/S$, rational over $k(s)$; by \hyperref[XIII.subsec1.5]{1.5} and
$(+)$, $Z/S$ is smooth at $z_{\circ}$. Since $S$ is the spectrum of a
\emph{henselian} ring, there is a section of $Z/S$ through $z_{\circ}$;
this follows, for example, immediately from the original definition
(\SGA II 1.1) of smoothness. It follows that $A_{\circ}$ is induced by an
algebra $A$ étale over $X$ (a quotient of $E$), which completes the proof
of \hyperref[XIII.prop1.1]{1.1}.

% END INCLUDED SGA4_English_current_buildable_source_20260906/source/13/drafts/SGA4_XIII_LINES_43_316_English.texfrag


% BEGIN INCLUDED SGA4_English_current_buildable_source_20260906/source/13/drafts/SGA4_XIII_LINES_317_549_English.texfrag
\section{The case of relative dimension at most one}\etiquette[2]{XIII.sec2}\pageoriginale

The assertion is the following.

\begin{enonce*}{Proposition 2.1}\etiquette[2.1]{XIII.prop2.1}
With the notation of \textup{\hyperrefext[XII][XII.thm5.9]{5.9 bis}}, the
functor $\Et(X)\rightarrow\Et(X_{\circ})$ is an equivalence of
categories if $S$ is Noetherian and strictly local and if $f$ is
projective of relative dimension $\leq1$.

By \textup{\hyperrefext[XII][XII.lem8.5]{8.5}}, this will complete the
proof of \textup{\hyperrefext[XII][XII.thm5.1]{5.1} (ii)}, which
generalizes both \hyperref[XIII.prop1.1]{1.1} and
\hyperref[XIII.prop2.1]{2.1}.
\end{enonce*}

\begin{enonce*}{Lemma 2.2}\etiquette[2.2]{XIII.lem2.2}
Let $S$ be the spectrum of a Noetherian local ring $A$, and let $X/S$
be projective of relative dimension $\leq n$. There exist a projective
and flat $S$-scheme $Z/S$ of relative dimension $\leq n$ and a closed
immersion $X\hookrightarrow Z$ over $S$.
\end{enonce*}

\noindent
\emph{Proof.} Let $X\rightarrow\PP^N_S$ be a projective embedding over
$S$, and let $J\subset A[x_{\circ},\ldots,x_N]$ be the homogeneous
ideal defining $X$. Put $k=A/\mathfrak m_A$, $S_{\circ}=\Spec k$, and
$X_{\circ}=X\otimes_Sk$. Then $X_{\circ}$ is defined by the
homogeneous ideal
$J_{\circ}=\Iim(J\rightarrow k[x_{\circ},\ldots,x_N])$. Since
$\dim X_{\circ}\leq n$, it is well known that there exist elements
$f^{\circ}_1,\ldots,f^{\circ}_{N-n}$ of $J_{\circ}$ defining a
subscheme $Z_{\circ}$ of $\PP^N_k$ of dimension $n$. Let
$f_i$ $(1\leq i\leq N-n)$ be elements of $J$ lifting the
$f^{\circ}_i$, and let $Z$ be the subscheme of $\PP^N_S$ defined by
$f_1,\ldots,f_{N-n}$. By \EGA IV 11.3, $Z$ is flat over $S$; and by
\EGA IV 13.1.5 it has relative dimension
$\dim Z_{\circ}=n$. This proves \hyperref[XIII.lem2.2]{2.2}.

\noindent
\emph{Proof of} \hyperref[XIII.prop2.1]{2.1}. Let $X/S$ be projective
of relative dimension $\leq1$, with $S$ Noetherian and strictly local.
We may assume that $X$ (and hence $X_{\circ}$) is connected. Let
$X\rightarrow Z$ be a closed immersion over $S$, where $Z$ is
projective, flat, and of relative dimension $\leq1$; see
\hyperref[XIII.lem2.2]{2.2}. Let $s_{\circ}$ be the closed point of
$S$. We have a corresponding\pageoriginale diagram of categories
$$
\xymatrix@C=2.5cm@R=1.25cm{
\Et(Z)\ar[r]\ar[d]^{a}&\Et(X)\ar[d]^{b}\\
\Et(Z_{\circ})\ar[r]^{c}&\Et(X_{\circ}).
}
$$
We wish to prove that $b$ is an equivalence. By
\hyperrefext[XII][XII.lem5.8]{5.8}, $b$ is fully faithful, so it remains
to prove that it is essentially surjective. The functor $a$ is an
equivalence by \hyperref[XIII.prop1.1]{1.1}; it therefore suffices to
prove that $c$ is essentially surjective. Dropping the subscript
${}_{\circ}$, we are reduced to the following proposition.

\begin{enonce*}{Proposition 2.3}\etiquette[2.3]{XIII.prop2.3}
Let $k$ be a separably closed field, let $Z$ be a $k$-scheme locally
of finite type and of dimension $\leq1$, and let $X\hookrightarrow Z$
be a connected closed subscheme. Then every étale covering of $X$ is
induced by an étale covering of $Z$.
\end{enonce*}

\noindent
\emph{Proof.} We may suppose that $Z$ and $X$ have dimension $1$, the
case in which $X$ is zero-dimensional being immediate from the
hypothesis on $k$. Then $Z=X\cup Y$, where $Y$ is the closed
subscheme obtained as the closure of $Z-X$. Put $V=X\cap Y$; this is
a discrete, that is, locally Artinian, scheme. The morphism
$X\amalg Y\rightarrow Z$ is finite and surjective, hence is a strict
descent morphism for the category of étale coverings
(\SGA \hyperrefext[IX][IX.cor4.7]{4.7}). Clearly
\begin{equation*}
\fprod{(X\amalg Y)}{(X\amalg Y)}{Z}
=X\amalg Y\amalg V^{(1)}\amalg V^{(2)}
\quad(V^{(i)}=V){\quoi},\tag{$\ast$}
\end{equation*}
where $\pr_i$ sends $V^{(i)}$ to the component $X$ of $X\amalg Y$
and sends $V^{(j)}$ to the component $Y$ when $j\neq i$.

Let $X'/X$ be an étale covering of degree $n$, and let $Y''/Y$ be a
covering of degree $n$ (for example, the trivial covering of degree
$n$), so that $X'\amalg Y''$ is\pageoriginale a covering of
$X\amalg Y$ of degree $n$. It is enough to find descent data for this
covering relative to $X\amalg Y\rightarrow Z$. Inspection of
$(\ast)$ shows that such descent data amount simply to an isomorphism
$$
\fprod{X'}{V}{X}=V'\xrightarrow{\sim}
V''=\fprod{Y''}{V}{Y}{\quoi}.
$$
Since $k$ is separably closed, so are the residue fields of the points
of $V$; consequently every étale covering of $V$ is completely split.
Thus $V'\simeq V''$, which proves the result.

\begin{enonce*}[remark]{Remark 2.4}\etiquette[2.4]{XIII.rem2.4}
One may avoid the delicate result from \SGA IX and descent theory by a
direct argument showing that the category of étale coverings of $Z$
is equivalent to the category of triples $(X',Y',\theta)$, where
$X'$ (respectively $Y'$) is an étale covering of $X$ (respectively
$Y$), and $\theta$ is a $V$-isomorphism
$\fprod{X'}{V}{X}\simeq\fprod{Y'}{V}{Y}$. This result, valid for any
scheme $Z$ that is the union of two closed subschemes $X$ and $Y$, is
proved directly from the analogous description of quasi-coherent
modules on $Z$ (when $X=V(\mathfrak J)$, $Y=V(\mathfrak K)$, and
$\mathfrak J\cap\mathfrak K=0$).
\end{enonce*}

\footnotetext[1]{This also appears in \EGA IV 21.9.11 and 21.9.12 in
a slightly more general form; in particular, in
\hyperref[XIII.cor3.2]{3.2} it is enough for $f$ to be \emph{separated}
rather than proper.}
\section{An auxiliary result on the Picard group ${}^{(*)}$}\etiquette[3]{XIII.sec3}

\begin{enonce*}{Proposition 3.1}\etiquette[3.1]{XIII.prop3.1}
Let $X$ be a Noetherian scheme and let $X_{\circ}$ be a closed
subscheme of $X$ such that \textup{(a)} $\dim X_{\circ}\leq1$, and
\textup{(b)} for every nonempty connected closed subset $Y$ of $X$,
$Y_{\circ}=Y\cap X_{\circ}$ is nonempty and connected.
\begin{enumerate}
\renewcommand{\theenumi}{\roman{enumi}}
\renewcommand{\labelenumi}{(\theenumi)}
\item Let $Z_{\circ}$ be the union of the sets
$\overline{\{x\}}\cap X_{\circ}$, for $x\in\Ass\Oo_X$, that consist
of a single point. For every Cartier divisor $D_{\circ}$ on
$X_{\circ}$ such that $\supp(D_{\circ})$ does not meet $Z_{\circ}$,
there exists a Cartier divisor $D$ on $X$ inducing $D_{\circ}$ (and if
$D_{\circ}\geq0$, one may take $D\geq0$).

\item Suppose there is an affine open subset $U_{\circ}$ of
$X_{\circ}$ containing the finite set
$\Ass\Oo_{X_{\circ}}\cup Z_{\circ}$ (this condition is automatically
satisfied if $X_{\circ}$ admits an ample invertible sheaf). Then
$\Pic(X)\rightarrow\Pic(X_{\circ})$ is surjective.
\end{enumerate}
\end{enonce*}

\noindent
\emph{Proof.}
\begin{enumerate}
\renewcommand{\theenumi}{\roman{enumi}}
\renewcommand{\labelenumi}{(\theenumi)}
\item Since $\dim X_{\circ}\leq1$, every Cartier divisor
$D_{\circ}$ on $X_{\circ}$ plainly has discrete support, and is
therefore the difference of two effective Cartier divisors whose
supports are contained in that of $D_{\circ}$. We are thus reduced in
(i) to the case where $D_{\circ}$ is effective. For every
$x\in\Supp D_{\circ}$ it is then defined at $x$ by a non-zero-divisor
$f_x\in\rad(\Oo_{X_{\circ},x})$. The element $f_x$ is induced by a
section $g_x$ of $\Oo_X$ on an open neighborhood $U_x$ of $x$. I claim
that, after shrinking $U_x$, the section $g_x$ is a non-zero-divisor
in $\Oo_{U_x}$, that is,
$V(g_x)\cap\Ass\Oo_{U_x}=\varnothing$.

Indeed, conditions (a) and (b) imply that, for every
$z\in\Ass\Oo_X$, the set $\overline{\{z\}}\cap X_{\circ}$ is a
connected closed subset of $X_{\circ}$ and hence is either a single
closed point or a connected union of one-dimensional irreducible
components of $X_{\circ}$. Since $x\notin Z_{\circ}$, in the first
case $x\notin\overline{\{z\}}$. In the second case, one cannot have
both $x\in\overline{\{z\}}$ and $z\in V(g_x)$, for then $g_x$, and
therefore $f_x$, would vanish on an irreducible component of
$\overline{\{z\}}\cap X_{\circ}$, which is also an irreducible
component of $X_{\circ}$; this is incompatible with $f_x$ being a
non-zero-divisor in $\Oo_{X_{\circ},x}$. After shrinking $U_x$, we may
therefore suppose\pageoriginale that
$V(g_x)\cap\Ass\Oo_{U_x}=\varnothing$, so $g_x$ is a
non-zero-divisor.

Let $Y_x$ be the closure of $V(g_x)$ in $X$. The point $x$ is isolated
in $Y_x\cap X_{\circ}\cap U_x$, hence also in
$Y_x\cap X_{\circ}$. By condition (b), $Y_x$ has a connected
component $Y'_x$ such that $Y'_x\cap X_{\circ}=\{x\}$. Moreover,
$Y'_x\subset U_x$: indeed,
$T_x=Y'_x-(Y'_x\cap U_x)$ is closed in $X$ and satisfies
$T_x\cap X_{\circ}=\varnothing$, so condition (b) gives
$T_x=\varnothing$. Let $D'_x$ be the Cartier divisor on $X$ induced
by $Y'_x$ on $U_x$ and by zero on $X-Y'_x$. Then
$D=\sum_xD'_x$ is an effective Cartier divisor on $X$ inducing
$D_{\circ}$ on $X_{\circ}$.

\item The conclusion follows from (i) and the fact that every
invertible module $\LL_{\circ}$ on $X_{\circ}$ is isomorphic to the
module defined by a Cartier divisor $D_{\circ}$ whose support does not
meet $Z_{\circ}$. This last assertion means that $\LL_{\circ}$ admits
a section on an open set $V_{\circ}$ containing
$\Ass\Oo_{X_{\circ}}\cup Z_{\circ}$ which vanishes at no point of
that set. It is enough to prove this after replacing $X_{\circ}$ by
the affine open $U_{\circ}$; there it follows from the familiar fact
that an invertible module on a semilocal scheme is isomorphic to the
structure sheaf. Finally, if $X_{\circ}$ admits an ample invertible
sheaf, every finite subset of $X_{\circ}$ is contained in an affine
open subset. This proves (ii).
\end{enumerate}

We have the following corollary, which generalizes
\hyperrefext[XII][XII.prop8.6]{8.6} and completes the proof of
\hyperrefext[XII][XII.thm5.1]{5.1} (iii).

\begin{enonce*}{Corollary 3.2}\etiquette[3.2]{XIII.cor3.2}
Let $S$ be the spectrum of a Noetherian Henselian ring, and let
$f:X\rightarrow S$ be a proper morphism of relative dimension
$\leq1$. For every closed subscheme $X'_{\circ}$ of $X$ having the
same underlying space as the closed fiber\pageoriginale $X_{\circ}$
of $X/S$, the canonical map
$\Pic(X)\rightarrow\Pic(X'_{\circ})$ is surjective.
\end{enonce*}

Indeed, apply the proposition to $(X,X'_{\circ})$. Condition (a) holds
by hypothesis, and condition (b) follows from
\hyperrefext[XII][XII.cor5.7]{5.7}. It remains to verify the existence
of $U_{\circ}$ required in (ii); this follows because $X'_{\circ}$ is
necessarily projective, as is every proper curve over a field.

\begin{enonce*}[remark]{Remark}
A reader who does not wish to assume or reconstruct the fact that
proper curves over a field are projective may use the corollary only
under the additional hypothesis that $f$ is projective; this suffices
for the application intended here. Conversely, admitting the
projectivity of proper algebraic curves over a field, one concludes
that under the hypotheses of \hyperref[XIII.cor3.2]{3.2}, $X$ is
necessarily projective over $S$. Indeed, choose an ample
$\xi_{\circ}\in\Pic(X_{\circ})$, lift it to some
$\xi\in\Pic(X)$, and apply \EGA III 4.7.1.
\end{enonce*}

\begin{thebibliography}{99}
\bibitem{sga4.exp.XIII.bib.TDTE} Grothendieck A., Technique de descente et
théorèmes d'existence en Géométrie Algébrique, in Fondements de la
Géométrie Algébrique (Extraits du Sém. Bourbaki 1957--62), Secrétariat
mathématique, 11, rue Pierre Curie.
\end{thebibliography}

\ifx\danslelivre\undefined
\end{document}
\fi

% END INCLUDED SGA4_English_current_buildable_source_20260906/source/13/drafts/SGA4_XIII_LINES_317_549_English.texfrag


% END INCLUDED SGA4_English_current_buildable_source_20260906/source/13/13_English.tex


\makeatletter
\def\input@path{{../14/}}
\makeatother
% BEGIN INCLUDED SGA4_English_current_buildable_source_20260906/source/14/14_English.tex
\ifx\danslelivre\undefined
% Release engine: XeLaTeX (the LuaTeX compatibility layer breaks Xy-pic output).
\documentclass[10pt]{../fichiers-configuration-latex/smfart}
% BEGIN INCLUDED SGA4_English_current_buildable_source_20260906/source/fichiers-configuration-latex/paquets.tex
\usepackage{polyglossia}
%\usepackage[francais]{babel}
\usepackage{fontspec}
\usepackage[math-style=ISO,bold-style=upright,sans-style=upright]{unicode-math}
%\usepackage[mathletters]{ucs}
%\usepackage[utf8]{inputenc}
%\usepackage[T1]{fontenc}
\usepackage[unicode,hypertexnames=false,colorlinks=true,linkcolor=blue,citecolor=Sepia,urlcolor=magenta]{hyperref}
\setmainlanguage{french}
\setmainfont[Ligatures=TeX]{Linux Libertine O}
% variante possible : TeXGyrePagella ?
\setmathfont{XITS Math}

% maths
\usepackage{amsmath}
%\usepackage{amssymb}
%\usepackage{mathrsfs}
%\usepackage{stmaryrd}
%\usepackage{euscript}

% couleurs
\usepackage[usenames,dvipsnames]{xcolor}

% graphiques
\usepackage{graphicx}
\usepackage[all]{xy}
%\usepackage{tikz}
%\usetikzlibrary{matrix,arrows,calc}
%\tikzset{column sep=3em,row sep=2em,text height=1.5ex, text depth=0.25ex,>=angle 90}
%\tikzset{description/.style={fill=white,inner sep=1pt}}

% hyperliens, 'etiquettes
\usepackage{xr-hyper}
\usepackage{../fichiers-configuration-latex/smfthm}
%\ifx\cacheretiquettes\undefined
%\usepackage[inline]{showlabels} \else \fi % pour rendre visibles les 'etiquettes
%\usepackage{../fichiers-configuration-latex/makeidx} % pour faire un index
%\usepackage[pagebackref,unicode]{hyperref}  % pour voir o`u sont cit'ees les r'ef'erences
% hyperref is loaded once above with the cumulative-reader destination policy.
%\hypersetup{breaklinks=true} %permet le retour `a la ligne dans les liens trop longs. Incompatible avec pdflatex
%\renewcommand*{\backref}[1]{{\small [cit\'e page(s)~#1.]}}

\usepackage{xspace}
\usepackage{pifont} % pour expose XVII

% END INCLUDED SGA4_English_current_buildable_source_20260906/source/fichiers-configuration-latex/paquets.tex

% BEGIN INCLUDED SGA4_English_current_buildable_source_20260906/source/fichiers-configuration-latex/numerotation.tex
\renewcommand{\thefootnote}{\roman{footnote}} % pour mettre les notes en bas de page sous la forme (ii) et non (2), qui ressemble trop (?) a un exposant.

\newcounter{pageoriginal}
\newcommand{\pageoriginale}{\refstepcounter{pageoriginal}\marginpar{\footnotesize\textbf{\thepageoriginal}}\xspace}
\let\pageoriginaled\pageoriginale

\usepackage{ifthen}
\newboolean{original}
\def\sisi{\ifthenelse{\boolean{original}}}
\setboolean{original}{false}

\newcounter{cnote}
\def\nde#1{\ifthenelse{\boolean{original}}{}{\stepcounter{cnote}\def\thefootnote{\arabic{cnote}}\protect\footnote[\value{cnote}]{N.D.E.: #1}}}
\def\ndemark{\ifthenelse{\boolean{original}}{}{\stepcounter{cnote}\def\thefootnote{\arabic{cnote}}\protect\footnotemark[\value{cnote}]}}
\def\ndetext#1{\ifthenelse{\boolean{original}}{}{\def\thefootnote{\arabic{cnote}}\protect\footnotetext[\value{cnote}]{N.D.E.: #1}}}

% END INCLUDED SGA4_English_current_buildable_source_20260906/source/fichiers-configuration-latex/numerotation.tex

% BEGIN INCLUDED SGA4_English_current_buildable_source_20260906/source/fichiers-configuration-latex/macros.tex
\setmainlanguage{french}
\setotherlanguage{english}


%% Fixed, engine-independent release identifier.  The inherited dynamic Git
%% probe used Lua io.popen, exposed build-state-dependent output, and printed
%% a visible question mark outside a Git checkout.
\def\showgitstatus{English release r2, 2026-07-23}


% macros pour LaTeX
\renewcommand{\mathsf}{}

% Police différente pour U+2113 SCRIPT SMALL L
\setmathfont[range="2113]{latinmodern-math.otf}

% Give every inherited semantic label its own stable PDF destination instead
% of inheriting whichever automatic equation/page anchor happened to be live.
% The legacy label remains as a compatibility alias used by existing edges.
\makeatletter
\def\etiquette[#1]#2{%
  \Hy@raisedlink{\hyper@anchorstart{sga4.semantic.#2}\hyper@anchorend}%
  \def\@currentHref{sga4.semantic.#2}%
  \def\@currentlabel{#1}%
  \label{#2}%
}
\makeatother
\def\eqhyperref[#1]#2{(\hyperref[#1]{#2})}
\def\pagehyperref[#1]#2{(page \hyperref[#1]{#2})}
\def\hyperrefext[#1][#2]#3{\hyperref[#2]{#1 #3}}
\def\eqhyperrefext[#1][#2]#3{\hyperref[#2]{#1 #3}}     %cf. exposes 17
\def\neweqhyperrefext[#1][#2]#3{\hyperref[#2]{#1 #3}}  %cf. exposes 17

\def\shat#1{\protect\widehat{#1}}
\def\stilde#1{\protect\widetilde{#1}}
%%%%%%%%%%%%%%%%%%


\def\etcn{etc.\xspace}
\def\etco{etc...\xspace}
\def\etc{\sisi{\etco}{\etcn}}
\def\ldotsn{\sisi{...}{\etcn}}

\def\cf{{cf.}\kern.3em}
\def\Cf{{Cf.}\kern.3em}
\def\ie{{i.e.}\kern.3em}
\def\iev{{i.e.},\xspace}
\def\loccito{loc.\kern2pt cit.\xspace}
\def\loccitn{\textit{loc.\kern2pt cit.}\xspace}
\def\loccit{\sisi{\loccito}{\loccitn}}
\def\EGA{\textup{EGA}\kern.2em}
\def\SGA{\textup{SGA}\kern.3em}

\def\quoi{}

\def\quav{\quad,\quad}
\def\quaet{\quad \text{et}\quad}
\def\fibreproduct#1#2#3{#1 ×_{#3} #2}
\let\fprod\fibreproduct
\def\ccheck#1#2{\mbox{${\check{\check{\rm #1}}}^{\lower2pt\hbox{$\scriptstyle #2$}}$}}
\def\sheaf#1{\protect\mathcal{#1}}
\def\vieuxsheaf#1{{#1}} % cf. problème avec sheaf{a\Phi_i} par exemple...
\DeclareMathOperator{\Hom}{Hom}
\DeclareMathOperator{\can}{can}
\DeclareMathOperator{\Ind}{Ind}
\DeclareMathOperator{\ind}{ind}
%\DeclareMathOperator{\Hhomtop}{\mathcal{H}\!\mathit{omtop}}
\DeclareMathOperator{\Hhomtop}{\mathbf{Homtop}}
\DeclareMathOperator{\Point}{Point}
\DeclareMathOperator{\Pro}{Pro}
\DeclareMathOperator{\Top}{Top}
\DeclareMathOperator{\Mod}{Mod}
\DeclareMathOperator{\Cart}{Cart}
%\DeclareMathOperator{\Mmorsite}{\mathcal{M}\!\mathit{orsite}}
\DeclareMathOperator{\Mmorsite}{\mathbf{Morsite}}
\DeclareMathOperator{\loc}{loc}
\DeclareMathOperator{\Loc}{Loc}
\DeclareMathOperator{\Pic}{Pic}
\DeclareMathOperator{\Et}{Et}
\DeclareMathOperator{\card}{card}
\DeclareMathOperator{\Ens}{Ens}
\DeclareMathOperator{\h}{\underline{H}}
\DeclareMathOperator{\Fonct}{Fonct}
\DeclareMathOperator{\Ker}{Ker}
\DeclareMathOperator{\but}{but}
\DeclareMathOperator{\Ob}{Ob}
\DeclareMathOperator{\ob}{ob}
\DeclareMathOperator{\Lib}{Lib}
\DeclareMathOperator{\esj}{esj}
\DeclareMathOperator{\id}{id}
\DeclareMathOperator{\Id}{Id}
\DeclareMathOperator{\ad}{ad}
\DeclareMathOperator{\Iim}{Im}
\DeclareMathOperator{\pr}{pr}
\DeclareMathOperator{\Cat}{Cat}
\DeclareMathOperator{\car}{car}
\DeclareMathOperator{\SysInd}{SysInd}
\DeclareMathOperator{\Sup}{Sup}
\DeclareMathOperator{\Coker}{Coker}
\DeclareMathOperator{\Filt}{Filt}
\DeclareMathOperator{\acc}{acc}
\DeclareMathOperator{\sur}{sur}
\DeclareMathOperator{\cat}{cat}
\DeclareMathOperator{\Ord}{Ord}
\DeclareMathOperator{\Gr}{Gr}
\DeclareMathOperator{\Cov}{Cov}
\DeclareMathOperator{\Caf}{Caf}
\DeclareMathOperator{\coker}{coker}
\DeclareMathOperator{\coim}{coim}
\DeclareMathOperator{\im}{im}
\DeclareMathOperator{\Ouv}{Ouv}
\DeclareMathOperator{\diag}{diag}
\DeclareMathOperator{\Esp}{Esp}
\DeclareMathOperator{\TOP}{TOP}
\DeclareMathOperator{\Unif}{Unif}
\DeclareMathOperator{\Iind}{\mathbf{Ind}}
\DeclareMathOperator{\Hhomcart}{\mathbf{Homcart}}
\DeclareMathOperator{\Hhomtopess}{\mathbf{Homtopess}}
\DeclareMathOperator{\Ppoints}{\mathbf{Points}}
\DeclareMathOperator{\Ppoint}{\mathbf{Point}}
\DeclareMathOperator{\Ttors}{\mathbf{Tors}}
\DeclareMathOperator{\Ssex}{\mathbf{Sex}}
\DeclareMathOperator{\Ffib}{\mathbf{Fib}}
\DeclareMathOperator{\Vvois}{\mathbf{Vois}}
\DeclareMathOperator{\Ppointess}{\mathbf{Pointess}}
\DeclareMathOperator{\scrHom}{\mathbf{Hom}}
\DeclareMathOperator{\sheafPic}{\mathbf{Pic}}
\DeclareMathOperator{\sheafAut}{\mathbf{Aut}}
\DeclareMathOperator{\Hhomtoplocan}{\mathbf{Homtoplocan}}
\DeclareMathOperator{\Hhomtopan}{\mathbf{Homtopan}}
\DeclareMathOperator{\Pptgeom}{\mathbf{Ptgéom}}
\DeclareMathOperator{\Morsite}{Morsite}
\DeclareMathOperator{\Fib}{Fib}
\DeclareMathOperator{\Fais}{Fais}
\DeclareMathOperator{\kar}{kar}
\DeclareMathOperator{\Groupes}{Groupes}
\DeclareMathOperator{\Restr}{Restr}
\DeclareMathOperator{\Isom}{Isom}
\DeclareMathOperator{\Prol}{Prol}
\DeclareMathOperator{\Res}{Res}
\DeclareMathOperator{\Tors}{Tors}
\DeclareMathOperator{\esp}{esp}
\DeclareMathOperator{\es}{es}
\DeclareMathOperator{\od}{od}
\DeclareMathOperator{\gr}{gr}
\DeclareMathOperator{\sob}{sob}
\DeclareMathOperator{\dfn}{dfn}
\DeclareMathOperator{\Sch}{Sch}
\DeclareMathOperator{\ouv}{ouv}
\DeclareMathOperator{\cons}{cons}
\DeclareMathOperator{\Nerf}{Nerf}
\DeclareMathOperator{\Points}{Points}
\DeclareMathOperator{\Espsob}{Espsob}
\DeclareMathOperator{\Homtop}{Homtop}
\DeclareMathOperator{\Pt}{Pt}
\DeclareMathOperator{\Quot}{Quot}
\DeclareMathOperator{\glob}{glob}
\DeclareMathOperator{\Homfib}{Homfib}
\DeclareMathOperator{\Qouv}{Qouv}
\DeclareMathOperator{\Quotouv}{Quotouv}
\DeclareMathOperator{\ens}{ens}
\DeclareMathOperator{\ab}{ab}
\DeclareMathOperator{\an}{an}
\DeclareMathOperator{\fs}{fs}
\DeclareMathOperator{\red}{red}
\DeclareMathOperator{\Ab}{Ab}
\DeclareMathOperator{\cosq}{cosq}
\DeclareMathOperator{\Ann}{Ann}
\DeclareMathOperator{\Zar}{Zar}


%%%%%%%%%%%%
%%%%% SGA4-2
\DeclareMathOperator{\Sex}{Sex}
\DeclareMathOperator{\vois}{vois}
\DeclareMathOperator{\Vois}{Vois}
\DeclareMathOperator{\Ext}{Ext}
\DeclareMathOperator{\ext}{ext}
\DeclareMathOperator{\Fer}{Fer}
\DeclareMathOperator{\Add}{Add}
\DeclareMathOperator{\Cocart}{Cocart}
\DeclareMathOperator{\cosk}{cosk}
\DeclareMathOperator{\plat}{plat}

\DeclareMathOperator{\faisc}{faisc}
\DeclareMathOperator{\sk}{sk}
\DeclareMathOperator{\ch}{ch}

\def\HH{\mathbf{H}}
\def\qc{\mathrm{qc}}
\def\et{\mathrm{et}}
\def\coh{\mathrm{coh}}

\def\Hvr#1{%
\begin{array}[b]{l}
{\kern2pt\scriptstyle r}\\[-8pt]
{\kern1pt\scriptstyle \vee}\\[-5pt]
H^{#1}
\end{array}}

\def\Hvv#1{%
\begin{array}[b]{l}
\kern2pt\rotatebox{90}{${\scriptscriptstyle\ll}$}\\[-6pt]
H^{#1}
\end{array}}

\def\Hvi#1{%
\begin{array}[b]{l}
{\kern0.2pt\scriptscriptstyle \infty}\\[-8pt]
{\kern1pt\scriptstyle \vee}\\[-5pt]
H^{#1}
\end{array}}

\def\Hv#1{%
\begin{array}[b]{l}
{\kern1pt\scriptstyle \vee}\\[-5pt]
H^{#1}
\end{array}}

\DeclareMathOperator{\cart}{cart}
\DeclareMathOperator{\Ffaisc}{\mathbf{Faisc}}


\def\Ev#1{%
\begin{array}[b]{l}
{\kern1pt\scriptstyle \vee}\\[-5pt]
E'_{#1}
\end{array}}

\DeclareMathOperator{\Cont}{Cont}
\DeclareMathOperator{\Ccont}{\mathscr{C}\!\mathit{ont}}
\DeclareMathOperator{\Site}{Site}
\newcommand\MM{\mathbf{M}}
\def\limtop#1{{\displaystyle{\varprojlim_{#1}\!\mbox{\rm top }}}}
\def\Limtop{\mathop{\varprojlim\!\mbox{\rm top}}}
\DeclareMathOperator{\spec}{spec}
\def\ss{\scriptstyle{\mathscr{S}}}
\DeclareMathOperator{\Gal}{Gal}
\def\smallm{\scriptstyle{\mathscr{M}}}
\DeclareMathOperator{\End}{End}
\DeclareMathOperator{\ass}{ass}
\DeclareMathOperator{\zar}{zar}
\DeclareMathOperator{\Ogg}{Ogg}
\DeclareMathOperator{\res}{res}
\DeclareMathOperator{\tr}{tr}
\DeclareMathOperator{\Supp}{Supp}
\DeclareMathOperator{\codim}{codim}
\DeclareMathOperator{\rk}{rk}
\DeclareMathOperator{\cur}{cur}
\DeclareMathOperator{\cdqc}{cdqc}
\newcommand\DD{\mathbf{D}}
\DeclareMathOperator{\cl}{cl}
\def\JJ{\mathbf{J}}
\def\FF{\mathbf{F}}
\def\BB{\mathbf{B}}
\DeclareMathOperator{\Alget}{Alget}
\DeclareMathOperator{\Ass}{Ass}
\DeclareMathOperator{\supp}{supp}
\DeclareMathOperator{\rad}{rad}
\DeclareMathOperator{\prof}{prof}
\DeclareMathOperator{\cint}{int}
\DeclareMathOperator{\sic}{sic}
\DeclareMathOperator{\Tor}{Tor}
\DeclareMathOperator{\tors}{tors}
\DeclareMathOperator{\Sym}{Sym}
\DeclareMathOperator{\parf}{parf}
\DeclareMathOperator{\tor}{tor}
\DeclareMathOperator{\defn}{defn}
\DeclareMathOperator{\Div}{Div}
\DeclareMathOperator{\Pol}{Pol}
\DeclareMathOperator{\Fex}{Fex}
\DeclareMathOperator{\Lie}{Lie}
\DeclareMathOperator{\tot}{tot}
\DeclareMathOperator{\Aut}{Aut}
\DeclareMathOperator{\HOM}{HOM}
\DeclareMathOperator{\pch}{pch}
\DeclareMathOperator{\Ch}{Ch}
\DeclareMathOperator{\EXT}{EXT}
\DeclareMathOperator{\PIC}{PIC}
\DeclareMathOperator{\Rf}{Rf}
\DeclareMathOperator{\inj}{inj}
\newcommand\VV{\mathbf{V}}

\def\Xv{%
\begin{array}[b]{l}
{\kern1pt\scriptstyle \vee}\\[-5pt]
X
\end{array}}

%%%%%%%
\newcommand\TT{\mathbf{T}}
\newcommand\XX{\mathbf{X}}
\newcommand\EE{\mathbf{E}}
%%%%%%

\def\qedsymbol{\textup{C.Q.F.D.}}
\let\CQFD\qed
\let\cqfd\qed

%%%%%\bigdot
\font\bigsymb=cmsy10 at 4pt
\def\bigdot{{\kern1.2pt\raise 1.5pt\hbox{\bigsymb\char15}}}
%%%%%

\def\Lright{\mathop{\vcenter{\hbox{\hbox{$\longrightarrow$}\lower8pt \hbox{\kern-16pt\hbox{$\longrightarrow$}}}}}}

\def\Aa{{\mathcal A}}
\def\Bb{{\mathcal B}}
\def\Cc{{\mathcal C}}
\def\Dd{{\mathcal D}}
\def\Ee{{\mathcal E}}
\def\Ff{{\mathcal F}}
\def\Gg{{\mathcal G}}
\def\Hh{{\mathcal H}}
\def\Ii{{\mathcal I}}
\def\Jj{{\mathcal J}}
\def\Kk{{\mathcal K}}
\def\Ll{{\mathcal L}}
\def\Mm{{\mathcal M}}
\def\Nn{{\mathcal N}}
\def\Oo{{\mathcal O}}
\def\Pp{{\mathcal P}}
\def\Qq{{\mathcal Q}}
\def\Rr{{\mathcal R}}
\def\Ss{{\mathcal S}}
\def\Tt{{\mathcal T}}
\def\Uu{{\mathcal U}}
\def\Vv{{\mathcal V}}
\def\Ww{{\mathcal W}}
\def\Xx{{\mathcal X}}
\def\Yy{{\mathcal Y}}
\def\Zz{{\mathcal Z}}

\newcommand{\CC}{\mathbf{C}}
\newcommand{\GG}{\mathbf{G}}
\newcommand{\KK}{\mathbf{K}}
\newcommand{\LL}{\mathbf{L}}
\newcommand{\NN}{\mathbf{N}}
\newcommand{\PP}{\mathbf{P}}
\newcommand\QQ{\mathbf{Q}}
\newcommand\RR{\mathbf{R}}
\newcommand\ZZ{\mathbf{Z}}

\def\to{\mathchoice{\longrightarrow}{\rightarrow}{\rightarrow}{\rightarrow}}
\def\from{\mathchoice{\longleftarrow}{\leftarrow}{\leftarrow}{\leftarrow}}
\def\mto{\mathchoice{\longmapsto}{\mapsto}{\mapsto}{\mapsto}}
\def\mfrom{\mathrel{\from\joinrel\mapstochar\,}}
\def\hto{\mathrel{\lhook\joinrel\to}}
\def\hfrom{\mathrel{\from\joinrel\rhook}}

\def\implies{\mathchoice{\Longrightarrow}{\Rightarrow}{\Rightarrow}{\Rightarrow}}
\def\implied{\mathchoice{\Longleftarrow}{\Leftarrow}{\Leftarrow}{\Leftarrow}}
\def\iff{\mathchoice{\Longleftrightarrow}{\Leftrightarrow}{\Leftrightarrow}{\Leftrightarrow}}

\def\To#1{\mathchoice{\xrightarrow{\labelstyle~~#1~}}{\stackrel{#1}{\longrightarrow}}{}{}}
\def\Mto#1{\mathrel{\mapstochar\To{#1}}}
\def\From#1{\mathchoice{\xleftarrow{\labelstyle~~#1~}}{\stackrel{#1}{\longleftarrow}}{}{}}

\newcommand{\isomto}{\overset{\sim}{\longrightarrow}}
\newcommand{\isomfrom}{\overset{\sim}{\longleftarrow}}

\newcommand{\IFF}{\unskip~$\iff$~\ignorespaces}
\newcommand{\THEN}{\unskip~$\implies$~\ignorespaces}

\let\leq\leqslant
\let\geq\geqslant
\let\le\leqslant
\let\ge\geqslant

\def\cf{{cf.}\kern.3em}
\def\Cf{{Cf.}\kern.3em}
\def\ie{{i.e.}\kern.3em}
\def\iev{{i.e.},\xspace}
%\def\loccito{loc.\kern2pt cit.\xspace}
%\def\loccitn{\textit{loc.\kern2pt cit.}\xspace}
%\def\loccit{\sisi{\loccito}{\loccitn}}
\def\EGA{\textup{EGA}\kern.2em}
\def\SGA{\textup{SGA}\kern.3em}

\let\oldS\S
\def\S{\oldS\kern .15em}
\def\ptbl{.\kern .2em}
\def\resp{resp.\kern.3em}
\def\Exp{Exp.\kern.3em}

\let\bar\overline
\let\hat\widehat
\let\tilde\widetilde
\let\dpl\displaystyle
\let\ldots\dots
\let\emptyset\varnothing
\let\skpt\skippointrait
\let\epsilon\varepsilon

\def\sfrac#1#2{{#1}/{#2}}

\DeclareMathOperator{\cd}{cd}
%\DeclareMathOperator{\Hhom}{\mathcal{H}\!\mathit{om}}
\DeclareMathOperator{\Hhom}{\mathbf{Hom}}
%\DeclareMathOperator{\Eext}{\mathcal{E}\!\mathit{xt}}
\DeclareMathOperator{\Eext}{\mathbf{Ext}}
\DeclareMathOperator{\Spec}{Spec}
\DeclareMathOperator{\Tr}{Tr}

\let\emptyset=\varnothing

%%%%%%%%%%%%%%%
%%% expose XVII
%%%%%%%%%%%%%%%

\makeatletter
\let\old@optexpose\optexpose
\newcommand{\empty@optexpose}[1]{}
\newcommand{\paren@optexpose}[1]{\old@optexpose{#1}(}
\newcommand{\refnoexpose}[1]{\let\optexpose\empty@optexpose\ref{#1}\let\optexpose\old@optexpose}%\refnoexpose{X.cor5.2} engendre "5.2" au lieu de "X 5.2"
\newcommand{\neweqref}[1]{\let\optexpose\paren@optexpose\textup{\ref{#1})}\let\optexpose\old@optexpose}%\neweqref{XII.eq4.1} engendre "XII (4.1)" au lieu de "(XII 4.1)"
\makeatother
\newcommand{\refvb}[2][VI${}_{\text{B}}$ ]{\sisi{#1\refnoexpose{#2}}{\ref{#2}}}

\newcommand{\leftexp}[2]{{\vphantom{#2}}^{#1}{#2}}
\newcommand{\AAA}{\mathbf{A}}
\newcommand{\Sf}{\mathfrak{S}}
\DeclareMathOperator*{\quotprojlim}{\text{``$\varprojlim$''}}
\DeclareMathOperator*{\quotinjlim}{\text{``$\varinjlim$''}}
\DeclareMathOperator{\Fl}{F\sisi{\ell}{l}}
\DeclareMathOperator{\Pl}{P\sisi{\ell}{l}}
%\DeclareMathOperator{\cl}{c\sisi{\ell}{l}}
\DeclareMathOperator{\tdf}{tdf}
\DeclareMathOperator{\rg}{rg}
\newcommand{\iInd}{\sisi{I}{i}nd}
\newcommand{\iind}{ind}
\newcommand{\ppro}{pro}
\newcommand{\tTor}{\sisi{T}{t}or}
\newcommand{\ttor}{tor}
\newcommand{\chap}{Chap.}
\newcommand{\Ffex}{\Ff\mathit{ex}}
\newcommand{\sheafk}{\underline{k}}
\newcommand{\sheafn}{\underline{n}}
\newcommand{\sheafGamma}{\underline{\Gamma}}

% Repeated references to an already typeset footnote must not use
% \footnotemark without a matching \footnotetext: Hyperref would then create
% a live annotation to a destination that cannot exist.  This helper keeps the
% current printed footnote mark and links it to an explicit stable note target.
\makeatletter
\newcommand{\sgarepeatedfootnotemark}[1]{%
  \begingroup
  \edef\@thefnmark{\thefootnote}%
  \hyperlink{#1}{\@makefnmark}%
  \endgroup}
\makeatother

\sisi{
\DeclareMathOperator{\simarrow}{\rule{.15mm}{1.8mm}\kern -3.4pt\rightsquigarrow}
\newcommand{\sheafell}{\underline{\ell}}
\newcommand{\sheafZZ}{\underline{\ZZ}}
}{
\newcommand{\nderef}[1]{la N.D.E.~(\ref{#1}) page~\pageref{#1}}
}

%%%%%%%%%%%%%%%%%
%%%%% expose XVII
%%%%%%%%%%%%%%%%%

% END INCLUDED SGA4_English_current_buildable_source_20260906/source/fichiers-configuration-latex/macros.tex


\externaldocument{../01/01_English.tex}
\externaldocument{../02/02_English.tex}
\externaldocument{../03/03_English.tex}
\externaldocument{../04/04_English.tex}
\externaldocument{../05/05_English.tex}
\externaldocument{../05bis/05bis_English.tex}
\externaldocument{../06/06_English.tex}
\externaldocument{../07/07_English.tex}
\externaldocument{../08/08_English.tex}
\externaldocument{../09/09_English.tex}
\externaldocument{../10/10_English.tex}
\externaldocument{../11/11_English.tex}
\externaldocument{../12/12_English.tex}
\externaldocument{../13/13_English.tex}

\setboolean{original}{false}

\begin{document}

\begin{center}
XIV. Finiteness theorem for a proper morphism; cohomological dimension of affine algebraic schemes\\
M. Artin
\end{center}

\selectlanguage{english}version: \showgitstatus
\tableofcontents

\else
\chapter[Proper finiteness and affine cohomological dimension]{Finiteness theorem for a proper morphism; cohomological dimension of affine algebraic schemes}
%\addtocontents{toc}{M. Artin}
\begin{center}
M. Artin
\end{center}
\fi
\phantomsection\MakeLinkTarget*{sga4.exp.XIV}\label{sga4.exp.XIV}

% BEGIN INCLUDED SGA4_English_current_buildable_source_20260906/source/14/drafts/SGA4_XIV_LINES_43_300_English.texfrag
\setcounter{pageoriginal}{144}

\newcommand{\degtr}{{\rm deg.tr.\ }}

This\pageoriginale exposé contains two important theorems whose proofs
use the base-change theorem for a proper morphism.

\section{Finiteness theorem for a proper morphism}\etiquette[1]{XIV.sec1}

\begin{enonce*}{Theorem 1.1}\etiquette[1.1]{XIV.thm1.1}
Let $f:X\rightarrow S$ be a proper morphism of finite presentation. Let
$A$ be a left Noetherian ring which is a $\ZZ/n$-algebra for some
suitable $n\in\NN$. Let $F$ be a \emph{constructible} sheaf of sets
(\resp groups, \resp $A$-modules) on $X$. Then the sheaf $f_*F$
(\resp $R^1f_*F$, \resp $R^qf_*F$ for every $q\geq0$) is also
constructible.
\end{enonce*}

In particular, we have:

\begin{enonce*}{Corollary 1.2}\etiquette[1.2]{XIV.cor1.2}
Let $X$ be a proper scheme over the spectrum of a separably closed
field $k$, and let $F$ be a constructible torsion abelian sheaf on
$X$. Then the groups $H^q(X,F)$ are finite for every $q\geq0$.
\end{enonce*}

\begin{enonce*}[remark]{Remark 1.3}\etiquette[1.3]{XIV.rem1.3}
The theorem is false for abelian sheaves if the condition that $f$ be
proper is omitted, as is seen by taking $F=\ZZ/p$, $p=\car k$,
$X=\EE^1_k$, the affine line, and $f$ either the inclusion of $X$ in
$\PP^1_k$ or the structural morphism $X\rightarrow\Spec k$.
Nevertheless, it is conjectured\nde{Ofer Gabber presented the following
generalization of \hyperref[XIV.thm1.1]{1.1} at the conferences held in
honor of Luc Illusie and Pierre Deligne in 2005. Let $f:X\to Y$ be a
morphism of finite type between quasi-excellent Noetherian schemes
(\EGA IV 7.8), and let $F$ be a constructible sheaf of
$\mathbf Z/n$-modules on $X_{\text{ét}}$, with $n$ invertible on $Y$.
Then the sheaves $R^if_*F$ on $Y_{\text{ét}}$ are constructible and
vanish for all but finitely many $i$. There are analogous assertions
for constructible sheaves of sets \resp groups of order invertible on
$Y$. \newline
Before Gabber's result, \hyperref[XIV.thm1.1]{Theorem 1.1} was known in
the following case. Let $S$ be a regular scheme of dimension $\leq1$,
let $f:X\to Y$ be a morphism of finite type between $S$-schemes, and
let $F$ be a sheaf of left $A$-modules on $X_{\text{ét}}$. Then the
sheaves $R^if_*F$ are constructible (\SGA 4 1/2, Finitude 1.1).}
that it is true if $F$ has order prime to the\pageoriginale residue
characteristics of $S$ and $S$ is \lq\lq excellent\rq\rq\ (\EGA IV
7.8). When resolution of singularities is available (for example, in
characteristic zero), this can be proved for an excellent
equicharacteristic scheme $S$
(\hyperrefext[XIX][XIX.thm5.1]{5.1}).
\end{enonce*}

\noindent
\emph{Proof of} \hyperref[XIV.thm1.1]{1.1}. The assertion is local on
$S$, so we may take $S$ affine. Then $S$ is a limit of spectra of rings
of finite type over $\ZZ$, and the data of the theorem can be obtained,
by base change $S\rightarrow S_0$, from a proper morphism
$f_0:X_0\rightarrow S_0$, with $S_0$ of finite type over $\ZZ$, and a
constructible sheaf $F_0$ on $X_0$; see \EGA IV 8 and
\hyperrefext[IX][IX.cor2.7.4]{2.7.4}. Since the inverse image of a
constructible sheaf is again constructible
(\hyperrefext[IX][IX.prop2.4]{2.4} (iii)), application of
\hyperrefext[XII][XII.thm5.1]{5.1} to the base-change morphism
$S_0\leftarrow S$ reduces us to proving the constructibility of
$f_{0*}F_0$ (\resp $R^qf_{0*}F_0$), whence the following lemma.

\begin{enonce*}{Lemma 1.4}\etiquette[1.4]{XIV.lem1.4}
It is enough to verify \hyperref[XIV.thm1.1]{1.1} when $S$ is of finite
type over $\Spec\ZZ$.
\end{enonce*}

\begin{enonce*}{Lemma 1.5}\etiquette[1.5]{XIV.lem1.5}
The set-theoretic assertion of \hyperref[XIV.thm1.1]{1.1} is true: that
is, $f_*F$ is a constructible sheaf of sets if $F$ is one.
\end{enonce*}

\noindent
\emph{Proof}. We may suppose that $S$ is Noetherian
(\hyperref[XIV.lem1.4]{1.4}). Let
$F\rightarrow G=\prod_i\pi_{i*}(C_i)$ be an injection, with
$\pi_i:X_i\rightarrow X$ finite and $C_i$ constant and constructible on
$X_i$ (\hyperrefext[IX][IX.prop2.14]{2.14}). We have
$f_*F\subset f_*G$; hence, by \hyperrefext[IX][IX.prop2.9]{2.9} (ii),
it is enough to treat the case $F=G$. Now
$f_*G=\prod_i(f\pi_i)_*C_i$, and a finite product of constructible
sheaves on $S$ is constructible
(\hyperrefext[IX][IX.prop2.6]{2.6} (ibis)). Replacing $X$ by $X_i$, we
are reduced to the case where $F=D_X$ is a constructible constant
sheaf with value $D$. By \hyperrefext[XII][XII.thm5.1]{5.1} (i), the
fiber $(f_*F)_{\bar s}$ is isomorphic to $D^{C(s)}$, where
$C(s)=\pi_0(X_{\bar s})$;\pageoriginale it is therefore finite. By
\hyperrefext[IX][IX.prop2.13]{2.13} (iii), it is enough to prove that
the function \lq\lq number of elements in the fiber of
$f_*F$\rq\rq\ is constructible on $S$. This function is
$\card(D^C)$, where $C$ is the function \lq\lq number of connected
components in the geometric fiber\rq\rq, and $C$ is constructible
(\EGA IV 9.7.9), which proves the result.

\begin{enonce*}{Lemma 1.6}\etiquette[1.6]{XIV.lem1.6}
Let $S$ be a Noetherian scheme, let $f:X\rightarrow S$ be proper, let
$F$ be a constructible sheaf of groups (\resp $A$-modules) on $X$, and
let $q$ be an integer equal to $0$ or $1$ (\resp an integer $\geq0$).
Let $S_i$ $(1\leq i\leq n)$ be subschemes of $S$ such that
$S=\bigcup_iS_i$, and let $F_i$ and $f_i:X_i\rightarrow S_i$ be the
\lq\lq restrictions\rq\rq\ of the data to $S_i$. If
$R^qf_{i*}F_i$ is constructible for every $i$, then $R^qf_*F$ is
constructible as well.
\end{enonce*}

\noindent
\emph{Proof}. The restriction of $R^qf_*F$ to $S_i$ is isomorphic to
$R^qf_{i*}F_i$ by \hyperrefext[XII][XII.thm5.1]{5.1}, and the lemma
follows from \hyperrefext[IX][IX.prop2.8]{2.8}.

We now proceed by arguments analogous to those of
\hyperrefext[XII][XII.sec8]{8}:

\begin{enonce*}{Lemma 1.7}\etiquette[1.7]{XIV.lem1.7}
Let $S$ be Noetherian, and let
$$
\xymatrix@=1.25cm{
X        \ar[rd]_{f} &           &\bar{X} \ar[ll]_{\pi}
 \ar[ld]^{\bar{f}}\\
                 &    S      &
}
$$
be a commutative diagram of proper morphisms, and let
$j:U\rightarrow X$ be an open subscheme such that the open subscheme
$\bar U=\fprod{U}{\bar X}{X}$ maps isomorphically onto $U$. Let
$Y=X-U$ be the reduced closed subscheme, and let $f_0:Y\rightarrow S$
and $i:Y\rightarrow X$ be the canonical morphisms. Suppose that
\hyperref[XIV.thm1.1]{Theorem 1.1}
holds\nde{The hypothesis concerning $R^q\pi_*$ occurs in the proof only
when $q=0$. In that case constructibility is already known from
\hyperref[XIV.lem1.5]{1.5}.}
for $\bar f$, $f_0$, and $\pi$. Then it also holds for $f$.
\end{enonce*}

\noindent
\emph{Proof}.\pageoriginale First consider a sheaf of $A$-modules $F$
on $X$. We have the exact sequence
$$
0\longrightarrow j_!j^*F\longrightarrow F\longrightarrow
i_*i^*F\longrightarrow0{\quoi}.
$$
In view of the corresponding cohomology exact sequence for the
$R^qf_*$ and \hyperrefext[IX][IX.prop2.6]{2.6}, it is enough to verify
constructibility for the two outer terms. Notice that $i_*i^*F$ is
constructible by \hyperrefext[IX][IX.prop2.4]{2.4} (iii) and
\hyperrefext[IX][IX.prop2.14]{2.14} (i), and hence so is $j_!j^*F$
(\hyperrefext[IX][IX.prop2.6]{2.6}). Since $i$ is a closed immersion,
$$
(R^qf_*)i_*(i^*F)\simeq(R^qf_{0*})(i^*F){\quoi},
$$
by \hyperrefext[VIII][VIII.cor5.6]{5.6}; constructibility in this case
therefore follows from the hypothesis on $f_0$. Moreover, if
$\bar j:\bar U\rightarrow\bar X$ denotes the inclusion morphism and
$\bar F=\pi^*F$, then
\begin{equation*}
j_!j^*F\simeq\pi_*\bar j_!\bar j^*\bar F{\quoi},\tag{$*$}
\end{equation*}
as is seen by first defining a morphism from left to right and then
proving it to be invertible fiber by fiber using
\hyperrefext[XII][XII.thm5.1]{5.1}. At every geometric point $\bar y$
of $Y$, the restriction of $\bar j_!\bar j^*\bar F$ to the fiber
$\bar X_{\bar y}$ of $\bar X$ over $X$ is zero; hence, by
\hyperrefext[XII][XII.thm5.1]{5.1} (iii), the fiber of
$(R^q\pi_*)(\bar j_!\bar j^*\bar F)$ at $\bar y$ is zero for every
$q\geq0$. On the other hand, $\pi$ induces an isomorphism over $X-Y$,
so
$$
(R^q\pi_*)(\bar j_!\bar j^*\bar F)=0
$$
for $q>0$. The Leray spectral sequence
$$
E^{p,q}_2=(R^pf_*)(R^q\pi_*)(\bar j_!\bar j^*\bar F)
\Longrightarrow R^n\bar f_*(\bar j_!\bar j^*\bar F)
$$
together with $(*)$ therefore gives isomorphisms
$$
(R^q\bar f_*)(\bar j_!\bar j^*\bar F)\simeq
(R^qf_*)(j_!j^*F)
$$
for $q\geq0$. Since $\bar j_!\bar j^*\bar F$ is constructible (by the
same argument as for $j_!j^*F$), the left-hand side is constructible by
the hypothesis on $\bar f$. Thus $(R^qf_*)(j_!j^*F)$ is constructible,
and consequently so is $R^qf_*F$.\pageoriginale

Now suppose that $F$ is a constructible sheaf of groups on $X$. Replace
$\pi:\bar X\rightarrow X$ by
$\pi\amalg i:(\bar X\amalg Y)\rightarrow X$. It is then enough to
prove that if $\pi:\bar X\rightarrow X$ is \emph{surjective} and
\hyperref[XIV.thm1.1]{1.1} holds for $\pi$ and $\bar f$, then it also
holds for $f$. First let $\bar F$ be a constructible sheaf of groups on
$\bar X$. Then $\pi_*\bar F$ is constructible on $X$
(\hyperref[XIV.lem1.5]{1.5}), and there is an injection
$$
(R^1f_*)\pi_*\bar F\hookrightarrow(R^1\bar f_*)\bar F.
$$
The right-hand side is constructible by hypothesis, and hence so is the
left-hand side (\hyperrefext[IX][IX.prop2.9]{2.9} (ii)). Since $\pi$ is
surjective, there is an injection
$F\rightarrow\pi_*\pi^*F=G$. The sheaf $\bar F=\pi^*F$ is
constructible, so $G$ and $R^1f_*G$ are constructible as well, by what
we have just proved. We are thus reduced to the following lemma.

\begin{enonce*}{Lemma 1.8}\etiquette[1.8]{XIV.lem1.8}
Let $S$ be Noetherian, let $f:X\rightarrow S$ be a proper morphism, and
let
$\displaystyle F\mathop{\hookrightarrow}^{u}G$
be an injection of constructible sheaves of groups on $X$. If
$R^1f_*G$ is constructible, then $R^1f_*F$ is constructible as well.
\end{enonce*}

\noindent
\emph{Proof}. Let $C=G/F$, which is a constructible sheaf of
homogeneous spaces under $G$ (apply
\hyperrefext[IX][IX.prop2.6]{2.6}). Consider the exact sequence
\begin{equation*}
f_*C\xrightarrow{\delta}R^1f_*F\xrightarrow{u^1}R^1f_*G{\quoi}.
\tag{$*$}
\end{equation*}
To prove that $R^1f_*F$ is constructible, it is enough, by Noetherian
induction and \hyperref[XIV.lem1.6]{1.6}, to prove that if
$S\neq\varnothing$, there exists a nonempty open subset of $S$ on which
$R^1f_*F$ is constructible. We may therefore suppose that $R^1f_*G$ is
locally constant. Since constructibility is local for the étale
topology, we may further suppose that $S$ is irreducible and that
$R^1f_*G$ is constant (and constructible), with value
$E=\{e_1,\ldots,e_n\}$. For $i=1,\ldots,n$, let $D_i$ be the
subsheaf\pageoriginale of $R^1f_*F$ which is the inverse image of the section $e_i$
under $u^1$. Since \lq\lq coproducts are universal\rq\rq\ in a topos,
we have $R^1f_*F=\amalg_iD_i$; it is therefore enough to prove that each
$D_i$ is constructible over a suitable nonempty scheme $S'$ étale over
$S$.

Recall that $R^1f_*F$ is the sheaf associated with the presheaf
$\Rr^1F$, where
$\Rr^1F(S')=H^1(X\times_SS',F)$. Consequently, we may suppose (after
replacing $S$ by a nonempty scheme $S'$ étale over

% END INCLUDED SGA4_English_current_buildable_source_20260906/source/14/drafts/SGA4_XIV_LINES_43_300_English.texfrag


% BEGIN INCLUDED SGA4_English_current_buildable_source_20260906/source/14/drafts/SGA4_XIV_LINES_301_601_English.texfrag
$S$) that $D_i$ is the “empty sheaf” (hence constructible),
\emph{or} that there exists an element $\alpha\in H^{1}(X,F)$ inducing
a section $\sisi{\alpha}{\bar{\alpha}}$ of $D_i$. In the latter case, let
$$
f_{\ast}C^{\alpha}\xrightarrow{\delta}R^{1}f_{\ast}F^{\alpha}\longrightarrow
R^{1}f_{\ast}G^{\alpha}
$$
be the exact sequence deduced from $(\ast)$ by twisting with
$\alpha$. We have
$$
R^{1}f_{\ast}F^{\alpha}\simeq R^{1}f_{\ast}F\quav
R^{1}f_{\ast}G^{\alpha}\simeq R^{1}f_{\ast}G{\quoi},
$$
and the section of $R^{1}f_{\ast}F^{\alpha}$ corresponding to
$\bar{\alpha}$ is the unit section. Thus the subsheaf
$D_i^{\alpha}$ of $R^{1}f_{\ast}F^{\alpha}$ corresponding to $D_i$
is the inverse image of the unit section of $R^{1}f_{\ast}G^{\alpha}$;
that is, $D_i^{\alpha}$ is the image of $f_{\ast}C^{\alpha}$ under
$\delta$. Since $f_{\ast}C^{\alpha}$ is constructible
(\hyperref[XIV.lem1.5]{1.5}), so is $D_i^{\alpha}$, and hence so is
$D_i$ (apply \hyperrefext[IX][IX.prop2.6]{2.6}). This proves the lemma.

\begin{enonce*}{Lemma 1.9}\etiquette[1.9]{XIV.lem1.9}
Let $S$ be noetherian, and let
$$
\xymatrix@=1.25cm{
X       \ar[rd]_{f}      &      &\bar{X} \ar[ll]_{\pi} \ar[ld]^{\bar{f}}\\
              & S    &
}
$$
be a commutative diagram of proper morphisms. If Theorem
\hyperref[XIV.thm1.1]{1.1} holds for\pageoriginale $f$ and $\pi$, then it also holds
for $\bar{f}$.
\end{enonce*}

\noindent
\emph{Proof}. In the case of a constructible sheaf $\bar{F}$ of
$A$-modules, we use the Leray spectral sequence
$$
E^{p,q}_2=(R^{p}f_{\ast})(R^{q}\pi_{\ast})\bar{F}\Longrightarrow
(R^{n}\bar{f}_{\ast})\bar{F}{\quoi}.
$$
By hypothesis, $E^{p,q}_2$ is constructible for every $p,q$; hence
the abutment is constructible as well, as follows from
\hyperrefext[IX][IX.prop2.6]{2.6}.

Let $\bar{F}$ be a constructible sheaf of groups. We have the exact
sequence (\hyperrefext[XII][XII.prop3.2]{3.2})
\begin{equation*}
0\longrightarrow (R^{1}f_{\ast})\pi_{\ast}\bar{F}\longrightarrow
(R^{1}\bar{f}_{\ast})\bar{F}\longrightarrow
f_{\ast}(R^{1}\pi_{\ast})\bar{F}{\quoi},\tag{$\ast$}
\end{equation*}
whose outer terms are constructible by the hypothesis and
\hyperref[XIV.lem1.5]{1.5}. We proceed as in the proof of
\hyperref[XIV.lem1.8]{1.8}: by noetherian induction and
\hyperref[XIV.lem1.6]{1.6}, it suffices to assume $S\neq ∅$ and prove
constructibility over a suitable nonempty open subset. We may thus
assume that $f_{\ast}(R^{1}\pi_{\ast})\bar{F}$ is locally constant and,
after an étale localization, even constant and constructible
(\sisi{\hyperrefext[IX][IX.prop2.8]{2.8}}{\hyperrefext[IX][IX.prop2.4]{2.4}~(i)}),
with value $E=\{e_1,\ldots,e_n\}$. Let $D_i$ be the subsheaf of
$(R^{1}\bar{f}_{\ast})\bar{F}$ that is the inverse image of $e_i$, so
that $(R^{1}\bar{f}_{\ast})\bar{F}=\coprod_iD_i$. It is enough to show
that each sheaf $D_i$ induces a constructible sheaf over some suitable
nonempty connected étale $S'/S$. If $D_i$ is not the “empty sheaf”, we
reduce to the case in which there exists an element
$\alpha\in H^{1}(\bar{X},\bar{F})$ inducing a section $\bar{\alpha}$ of
$D_i$. Let
$$
\sisi{
0\longrightarrow (R^{1}f_{\ast})\pi_{\ast}\bar{F}\longrightarrow
(R^{1}\bar{f}_{\ast})\bar{F}\longrightarrow
f_{\ast}(R^{1}\pi_{\ast})\bar{F}
}
{0\longrightarrow (R^{1}f_{\ast})\pi_{\ast}\bar{F}^{\alpha}\longrightarrow
(R^{1}\bar{f}_{\ast})\bar{F}^{\alpha}\longrightarrow
f_{\ast}(R^{1}\pi_{\ast})\bar{F}^{\alpha}}
$$
be the exact sequence deduced from $(\ast)$ by twisting $\bar{F}$ with
$\alpha$. The section of $(R^{1}\bar{f}_{\ast})\bar{F}^{\alpha}$
corresponding to $\bar{\alpha}$ is the unit section. Consequently the
subsheaf $D_i^{\alpha}$ of \hbox{$(R^{1}\bar{f}_{\ast})\bar{F}^{\alpha}$}
corresponding to $D_i$ is the inverse image of the unit section of
$f_{\ast}(R^{1}\pi_{\ast})\bar{F}^{\alpha}$,
\pageoriginale whence
$D_i^{\alpha}\simeq (R^{1}f_{\ast})\pi_{\ast}\bar{F}^{\alpha}$,
which is indeed a constructible sheaf, \cqfd

\begin{enonce*}{Lemma 1.10}\etiquette[1.10]{XIV.lem1.10}
It suffices to verify \hyperref[XIV.thm1.1]{1.1} when $S$ is of finite
type over $\Spec\ZZ$, and $f$ is projective and of relative dimension
$\leq 1$.
\end{enonce*}

\noindent
\emph{Proof}. Suppose first that $f$ is projective, and argue by
induction on the relative dimension of $f$, which may be assumed finite
after replacing $S$ by an affine open subset. If $f$ has relative
dimension $\leq n$, with $n\geq 2$, then locally on $S$ one can find a
finite morphism $\pi:X\rightarrow\PP^n_S$. By
\hyperref[XIV.lem1.9]{1.9} and
\sisi{IX 2.15}{\hyperrefext[IX][IX.prop2.14]{2.14}} (i), we reduce to
the case in which $f$ is the structural morphism of $\PP^n_S$.
Consider the commutative diagram of morphisms
$$
\xymatrix@C=2cm{
                   & \bar{X}\ar[ld]_{\sisi{}{\pi}} \ar[rd]^{a} \ar@{-}[dd]^{\bar{f}} &       \\
\PP^{n}_S \ar[rd]_{f}  &                  &\PP^{1}_S \ar[ld]^{b}\\
                   &  S                          &
}
$$
already considered in the proof of
\hyperrefext[XII][XII.lem8.3]{8.3}, where $\bar{X}$ is obtained from
$\PP^n_S$ by blowing up the closed subscheme
$Y\simeq\PP^{n-2}_S$. The relative dimension of $b$ is $\leq1$, and
that of $a$ is $\leq n-1$; therefore the theorem holds for $a$ and $b$
by the induction hypothesis, and hence for $\bar{f}$ by
\hyperref[XIV.lem1.9]{1.9}. Moreover, the relative dimension of $\pi$
is $\leq1$, while that of $f_0:Y\rightarrow S$ is $\leq n-1$. It
therefore follows from \hyperref[XIV.lem1.7]{1.7} that
\hyperref[XIV.thm1.1]{1.1} holds for $f$, and thus for every projective
morphism. Now let $f$ be an arbitrary proper morphism, and argue by
noetherian induction on $X$: we may assume the theorem holds for every
proper closed subscheme of $X$ and that $X\neq∅$. Let
$$
\xymatrix@=1.5cm{
X  \ar[rd]_{f}          &           &\bar{X} \ar[ll]_{\pi} \ar[ld]^{\bar{f}}\\
                        &    S       &
}
$$
be a\pageoriginale diagram of the type supplied by Chow's lemma
\hyperrefext[XII][XII.lem7.1]{7.1}. Applying
\hyperref[XIV.lem1.7]{1.7} to this diagram yields
\hyperref[XIV.thm1.1]{1.1} for $f$, and proves the lemma.

\setcounter{subsection}{10}
\subsection{}\etiquette[1.11]{XIV.subsec1.11}
We can now complete the proof of \hyperref[XIV.thm1.1]{1.1}. If $F$ is
a constructible sheaf of groups, one can find an injection
$F\hookrightarrow G=\prod_i\pi_{i\sisi{}{,}\ast}C_i$, where each
$\pi_i$ is finite and each $C_i$ is constant and constructible
(\sisi{IX 2.15}{\hyperrefext[IX][IX.prop2.14]{2.14}}); by
\hyperref[XIV.lem1.8]{1.8}, we may therefore assume that $F$ is
constant. Likewise, if $F$ is a constructible sheaf of $A$-modules,
then by applying \sisi{IX 2.15}{\hyperrefext[IX][IX.prop2.14]{2.14}}
one can find a resolution $F\rightarrow G^{\bigdot}$,
$G^{\bigdot}=\{0\rightarrow G^0\rightarrow G^1\rightarrow
G^2\rightarrow\ldots\}$, of $F$ such that every $G^{\nu}$ is of the
form
$G^{\nu}=\prod_i\sisi{\eta}{\pi}_{i\sisi{}{,}\ast}C_i$.
The spectral sequence of the resolution $G^{\bigdot}$,
$$
E^{p,q}_2=H^{\sisi{q}{p}}((R^{q}f_{\ast})G^{\bigdot})\Longrightarrow
(R^{n}f_{\ast})F{\quoi},
$$
then reduces us to the case in which $F$ is constant and constructible,
with value $M$.

Recall that $A$ is a $\ZZ/n$-algebra. Let $\ell$ be a prime number
dividing $n$. We have the exact sequence
$$
0\longrightarrow \ell M\longrightarrow M\longrightarrow M/\ell
M\longrightarrow 0{\quoi},
$$
and it is again enough to prove \hyperref[XIV.thm1.1]{1.1} for
$\ell M$ and $M/\ell M$.

An easy induction then reduces us to the case in which $A$ is a
$\ZZ/\ell$-algebra, $\ell\in\PP$, and $M$ is a finite-type
$A$-module. Then $M$ is a vector space over $\ZZ/\ell$, and the
canonical morphism
$$
H^q(Y,\ZZ/\ell)\otimes_{\ZZ/\ell}M\longrightarrow H^q(Y,M)
$$
is bijective for every quasi-compact and quasi-separated $Y$ (indeed,
this is trivial if $M$ has finite rank, and both sides commute with
filtered direct limits\pageoriginale
(\sisi{VI 1.6}{\hyperrefext[VI][VI.thm5.1]{5.1} and
\hyperrefext[VII][VII.prop3.3]{3.3}})). Hence
$$
[R^qf_{\ast}(\ZZ/\ell)]\otimes_{\ZZ/\ell}M\xrightarrow{\sim}
R^qf_{\ast}(M){\quoi}.
$$
It plainly suffices, therefore, to show that $R^qf_{\ast}(\ZZ/\ell)$
is constructible as a $\ZZ/\ell$-module; in other words, we are reduced
to the case $A=\ZZ/\ell$, which we assume from now on.

Observe that the case in which $f$ has relative dimension $\leq0$
(that is, $f$ is finite), with $F$ arbitrary, now follows from
\hyperrefext[VIII][VIII.cor5.6]{5.6} and
\hyperref[XIV.lem1.5]{1.5}.

We proceed by noetherian induction on $S$. By
\hyperref[XIV.lem1.6]{1.6}, we may assume $S\neq∅$ and may replace $S$
by any nonempty open subset. Moreover, we may assume that $S$ is
integral, with generic point $s$. The fiber $X_s$ is an algebraic
scheme of dimension $\leq1$; consequently there is a purely inseparable extension
$K'$ of $k(s)=K$ such that the normalization $Y$ of
$(X_s\otimes_KK')_{\text{\rm red}}$ is smooth over $\Spec K'$
\footnote[1]{\Cf \EGA IV 17.15.14.}. After replacing $S$ by a nonempty
open subset, we may assume that there exists a finite, radicial,
surjective morphism $S'\rightarrow S$ such that $S'_s$ is
$K$-isomorphic to $\Spec K'$, a smooth projective morphism
$\bar{X}'\rightarrow S'$, and a finite surjective $S'$-morphism
$\pi:\bar{X}'\rightarrow X'=\fprod{X}{S'}{S}$
inducing $Y\rightarrow X_s\otimes_KK'$ on the fibers over the generic
point $s$ of $S$ (\EGA IV 9.6.1). Now\sisi{}{,} by
\hyperrefext[VIII][VIII.thm1.1]{1.1}\sisi{}{,} making a purely inseparable
extension is harmless. We may therefore replace $S$ by $S'$ and $X$ by
$X'$; that is, we may assume that $f$ fits into a diagram
$$
\xymatrix@=1.5cm{
X   \ar[rd]_{f}& &\bar{X} \ar[ll]_{\pi} \ar[ld]^{\bar{f}}\\
         &         S        &           ,
}
$$
in which\pageoriginale $\pi$ is finite and surjective and $\bar f$ is
smooth. Moreover, by construction, we may assume that there is a closed
subscheme $Y$ of $X$ whose generic fiber is finite and such that $\pi$
induces an isomorphism over $X-Y$. Replacing $S$ once more by a
nonempty open subset, we reduce to the case in which
$f_0:Y\rightarrow S$ is also finite.

Since $F$ is constant, $\pi^{\ast}F$ is constant as well, and
$F\rightarrow\pi_{\ast}\pi^{\ast}F$ is injective because $\pi$ is
surjective. If $F$ is a sheaf of groups, it suffices to prove that
$(R^1f_{\ast})(\pi_{\ast}\pi^{\ast}F)\simeq
(R^1\bar f_{\ast})(\pi^{\ast}F)$ is constructible
(\hyperref[XIV.lem1.8]{1.8}). If $F$ is a sheaf of $\ZZ/\ell$-modules,
we have an exact sequence
$$
0\longrightarrow F\longrightarrow\pi_{\ast}\pi^{\ast}F\longrightarrow
C\longrightarrow0{\quoi},
$$
where $C$ is supported on $Y$, because $\pi$ induces an isomorphism over
$X-Y$. It follows that $R^qf_{\ast}C=R^qf_{0\sisi{}{,}\ast}C$ is
constructible (and vanishes if $q>0$), so it is once again enough to
show that
$(R^qf_{\ast})(\pi_{\ast}\pi^{\ast}F)=
(R^q\bar f_{\ast})(\pi^{\ast}F)$ is constructible. Replacing $f$ by
$\bar f$ and $F$ by $\pi^{\ast}F$, we are reduced to the case in which
$F$ is constant and constructible and $f$ is smooth.

Let us first treat a constant sheaf of groups, $\sisi{f}{F}=G_X$. All the fibers
of $(R^1f_{\ast})G_X$ are finite. Indeed, by applying
\hyperrefext[XII][XII.cor5.2]{5.2} (ii), we reduce, in order to prove
this, to the case in which $S$ is the spectrum of a separably closed
field and $X$ is smooth of dimension $\leq1$ over $S$. The result is
known in this case (\SGA 1 X 2.6). It therefore remains only to prove
that the specialization morphisms are injective
(\hyperrefext[IX][IX.prop2.13]{2.13} (ii)). For this we reduce to the
case in which $S$ is strictly local; by the base-change theorem
\hyperrefext[XII][XII.thm5.1]{5.1} (ii), we may even replace $S$ by the
spectrum of a strictly local discrete valuation ring. Moreover, we may
assume $X$ connected. Since\pageoriginale $H^1(X,G)$ classifies the
principal Galois coverings of $X$ with group $G$, it is enough to check
that if $X'\rightarrow X$ is a connected étale covering, then its
geometric fiber is connected as well. This follows readily from the
fact that $f$ is smooth and proper with geometrically connected fibers.

We now treat the case $F=(\ZZ/\ell)_X$. For $q=0$ and $q=1$, the
result is already known (\hyperref[XIV.lem1.5]{1.5} and the preceding
argument). \sisi{By \hyperrefext[X][X.cor4.3]{4.3} and}{By}
\hyperrefext[XII][XII.corbis5.3]{5.3 bis}, only the case $q=2$ remains;
furthermore, if $S$ has characteristic $\ell$, then
$R^2f_{\ast}(\ZZ/\ell)=0$. Replacing $S$ by a suitable nonempty open
subset, we therefore reduce to the case in which $\ell$ is invertible
on $S$. Then, by \hyperrefext[XII][XII.cor5.2]{5.2}
\sisi{(ii)}{(iii)} and
\sisi{X 4.7}{\hyperrefext[X][X.cor5.2]{5.2}}, the fiber of
$R^2f_{\ast}(\ZZ/\ell)$ at a geometric point $\bar s$ above $s\in S$
is isomorphic to $(\ZZ/\ell)^{c(s)}$, where $c(s)$ is the number of
one-dimensional irreducible components of the geometric fiber
$X_{\bar s}$. The function $c$ is constructible on $S$
(\EGA IV 9.7.9), and the result follows
(\hyperrefext[IX][IX.prop2.13]{2.13} (iii)).

\begin{enonce*}[remark]{Remark 1.11}\etiquette[1.11]{XIV.rem1.11}
The proof of \hyperref[XIV.thm1.1]{1.1} becomes much simpler if one
uses the formalism of cohomology with proper supports, which will be
developed in \hyperref[sga4.exp.XVII]{XVII} as a direct consequence of the base-change
theorem for a proper morphism \hyperrefext[XII][XII.thm5.1]{5.1}.
Moreover, the statement of \hyperref[XIV.thm1.1]{1.1} generalizes to
the case in which one is given an arbitrary sheaf $𝒜$ of torsion rings
on $\sisi{Y}{S}$ and a complex $K^{\bigdot}$ of $f^{\ast}(𝒜)$-modules
on $X$ satisfying a “constructibility” condition which, in the special
case of \hyperref[XIV.thm1.1]{1.1}, would simply say that the cohomology
sheaves $\Hh^i(K^{\bigdot})$ of $K^{\bigdot}$ are constructible
$A$-modules and vanish for sufficiently large $i$.
\end{enonce*}

% END INCLUDED SGA4_English_current_buildable_source_20260906/source/14/drafts/SGA4_XIV_LINES_301_601_English.texfrag


% BEGIN INCLUDED SGA4_English_current_buildable_source_20260906/source/14/drafts/SGA4_XIV_LINES_602_722_English.texfrag
\section{A variant of dimension}\etiquette[2]{XIV.sec2}\pageoriginale

\subsection{}\etiquette[2.1]{XIV.subsec2.1}
Let $X$ be a scheme. We shall write
\begin{equation*}
d(x)=\dim\bar{\{x\}}\tag{2.1.1}
\end{equation*}\etiquette[2.1.1]{XIV.eq2.1.1}
for $x\in X$, where $\bar{\{x\}}$ is the closure of $\{x\}$. If $F$
is an abelian sheaf on $X$, we put
\begin{equation*}
d(F)=\sup \{d(x)\mid x\in X \quaet F_{\bar{x}}\neq 0\}{\quoi}.\tag{2.1.2}
\end{equation*}\etiquette[2.1.2]{XIV.eq2.1.2}
We shall need a corrected notion in the local case:

\begin{enonce*}{Definition 2.2}\etiquette[2.2]{XIV.def2.2}
Let $Y$ be the spectrum of a universally catenary Noetherian local
ring (\EGA IV 5.6.2), and let $f:X\rightarrow Y$ be a morphism of
finite type. Let $x\in X$ and $y=f(x)$. We put
$$
\delta(x)=\dim\bar{\{y\}}+\degtr k(x)/k(y){\quoi}.
$$
If $F$ is an abelian sheaf on $X$, we put
$$
\delta(F)=\sup \{\delta(x)\mid x\in X\quaet F_{\bar{x}}\neq 0\}{\quoi}.
$$
\end{enonce*}

\begin{enonce*}{Proposition 2.3}\etiquette[2.3]{XIV.prop2.3}
\begin{enumerate}
%%%
%%% The typesetting of this enumerate environment is terrible.
%%% The spacing between the dash and (i) is too big.
%%% Also, the items (ii), ... should better be without indent.
%%% If you use the package enum.sty, then you can build the necessary list environment.
%%%
\renewcommand{\theenumi}{\roman{enumi}}
\renewcommand{\labelenumi}{(\theenumi)}
\item Let $i:X'\rightarrow X$ be an immersion and let $x'\in X'$.
Then $\delta(i(x'))=\delta(x')$.
%%% new line break within formula in text: spend extra line to induce new line breaks?

\item Let $X_{0}$ be the closed fiber of $X/Y$. We have
$$
\delta(x)=
\begin{cases}
d(x) & \mbox{if } \bar{\{x\}}\cap X_{0}\neq ∅\\
d(x)+1 & \mbox{if } \bar{\{x\}}\cap X_{0}=∅{\quoi};
\end{cases}
$$

\item $\delta(x)=d(x)$ if $f$ is proper.
\end{enumerate}
\end{enonce*}

\noindent
\emph{Proof}.\pageoriginale Assertion (i) follows immediately from the
definition. Since $\bar{\{x\}}\cap X_{0}$ is nonempty when $f$ is
proper, (iii) follows from (ii). For (ii), we may replace $X$ by
$\bar{\{x\}}$ and $Y$ by $\bar{\{y\}}$ (with their induced reduced
structures), as is permitted by (i). Let $x'\in X$ and $y'=f(x')$.
We have the formula (\EGA IV 5.6.1.1)
$$
\dim \Oo_{X,x'}=\dim
\Oo_{Y,y'}+\degtr k(x)/k(y) - \degtr k(x')/k(y'){\quoi}.
$$

Suppose that $X_{0}\neq ∅$. It is then clear that
$\displaystyle{\mathop{\sup}_{x'}}$
$\{\dim \Oo_{X,x'}\}=\dim X$ is attained by taking $x'$ to be a closed
point of $X_{0}$, whence
$$
\dim X=\dim \Oo_{X,x'}=\dim Y + \degtr k(x)/k(y)=\delta(x){\quoi}.
$$
If $X_{0}=∅$, the supremum is plainly attained by taking $x'$ to
be a closed point in its fiber $f^{-1}(y')$, with
$\dim\bar{\{y'\}}=1$. Such points exist because the $y'\in Y$ with
$\dim\bar{\{y'\}}=1$ form a very dense set (\EGA IV 10.1.3,
\sisi{10.5.9, 1.8.4}{10.3.1, 10.4.1, 10.5.9}). Thus
\begin{align*}
\dim X=\dim \Oo_{X,x'} &= \dim \Oo_{Y,y'}+\degtr k(x)/k(y)\\
&= \dim Y-1 +\degtr k(x)/k(y)=\delta(x) \sisi{=}{-} 1{\quoi},
\end{align*}
which proves the result.

\begin{enonce*}{Proposition 2.4}\etiquette[2.4]{XIV.prop2.4}
Let $f:X\rightarrow Y$ be a morphism of schemes of finite type over
the spectrum of a field $k$. Let $x\in X$, let $y=f(x)$, and let
$y_{0}$ be a specialization of $y$. Let $f':X'\rightarrow Y'$ be the
strict localization of $f$ at a geometric point $\bar{y}_{0}$ above
$y_{0}$, and let $x'\in X'$ be a point above $x$. Then
$$
\delta(x')=d(x)-d(y_{0}){\quoi}.
$$
\end{enonce*}

\noindent
\emph{Proof}.\pageoriginale We have
\begin{align*}
d(x) &= d(y)+\degtr k(x)/k(y)\\
\delta(x') &= d(y')+\degtr k(x')/k(y'){\quoi},
\end{align*}
where $y'=f'(x')$. Moreover,
$$
\degtr k(x)/k(y)=\degtr k(x')/k(y'){\quoi},
$$
because the extensions $k(x')/k(x)$ and $k(y')/k(y)$ are algebraic.
It therefore remains to prove that
$$
d(y')=d(y)-d(y_{0}){\quoi}.
$$

We may assume that $y$ is maximal in $Y$, and hence that $y'$ is
maximal in $Y'$. Then
\begin{align*}
d(y') &= \dim \Oo_{Y',\bar{y}_{0}}\\
      &= \dim \Oo_{Y,y_{0}}\quad \text{\rm (\EGA IV 6.1.3)}\\
      &= \sisi{\dim Y}{\dim_{y_0} Y} -d(y_{0})\quad \text{\rm (\EGA IV 5.2.3.1)}\\
      &= d(y)-d(y_{0}){\quoi},
\end{align*}
which proves the result.

% END INCLUDED SGA4_English_current_buildable_source_20260906/source/14/drafts/SGA4_XIV_LINES_602_722_English.texfrag


% BEGIN INCLUDED SGA4_English_current_buildable_source_20260906/source/14/drafts/SGA4_XIV_LINES_723_863_English.texfrag
\section{Cohomological dimension of affine algebraic schemes}\etiquette[3]{XIV.sec3}

\begin{enonce*}{Theorem 3.1~(d)}\etiquette[3.1]{XIV.thm3.1}\nde{Theorem \ref{XIV.thm3.1} is now known as the \textit{affine Lefschetz theorem}.}
Let $f:X\rightarrow Y$ be an affine morphism of schemes of finite type
over a field $k$. Let $F$ be a torsion sheaf on $X$ such that
$d(F)\leq d$ \emph{(\cf \hyperref[XIV.subsec2.1]{2.1})}.
Then $d(R^{q}f_{\ast}F)\leq d-q$.
\end{enonce*}

In particular, taking $Y=\Spec(k)$, we obtain:

\begin{enonce*}{Corollary 3.2}\etiquette[3.2]{XIV.cor3.2}
Let\pageoriginale $X$ be an affine scheme of finite type over a separably
closed field $k$.
Then\nde{Consider the analogous assertion in complex analytic geometry:
let $X$ be an affine variety over $\mathbf{C}$ and let $F$ be a constructible sheaf on $X$. Then $H^m(X^{\rm an},F)=0$ for every $m>\dim X$, where $X^{\rm an}$ is the analytic space associated with $X$. Although this assertion follows from \hyperref[XIV.cor3.2]{3.2} by the comparison theorem \hyperrefext[XVI][XVI.thm4.1]{4.1}, there are also two entirely analytic proofs, due to Pierre Deligne and Hélène Esnault; see Esnault, H., Variation on Artin's vanishing theorem, \textit{Advances in Math.} \textbf{198} (2005), 435--438.}
$$
\cd X\leq \dim X{\quoi}.
$$
\end{enonce*}

\begin{enonce*}{Corollary 3.3}\etiquette[3.3]{XIV.cor3.3}
Let $X\subset \PP^N_k$ be a projective scheme of dimension $n$ over a
separably closed field $k$, and let $Y=H\cap X$ be a ``hyperplane
section'' of $X$. Then, for every torsion sheaf $F$ on $X$, the canonical
homomorphism
$$
H^{q}_Y(X,F)\longrightarrow H^{q}(X,F)
$$
is bijective for $q>n+1$ and surjective for $q=n+1$.
\end{enonce*}

The corollary follows immediately from
\sisi{\hyperref[XIV.def2.2]{2.2}}{\hyperref[XIV.cor3.2]{3.2}} and the exact
sequence for a closed subset
(\sisi{IV 3.7.1}{\hyperrefext[V][V.6.5.4]{6.5.4}}),
%%% This last reference should read \ref{V.eq6.5.4}, but the TeX-file of chapter V contains the
%%% label without "eq".
since $X-Y=\sisi{\mathcal{U}}{U}$ is affine of dimension $\leq n$.
Notice that \hyperref[XIV.cor3.3]{3.3} is a generalization of one of
Lefschetz's theorems on hyperplane sections. Indeed, if $X$ is smooth
over $k$ and $F$ is locally constant of order prime to the characteristic,
the cohomology with support $H^{q}_Y(X,F)$ can be computed explicitly
using the relative cohomological purity theorem
(\hyperrefext[XVI][XVI.sec3]{3}) and the spectral sequence
(\sisi{\quad}{\hyperrefext[V][V.6.4.3]{6.4.3}}),
%%% This last reference should read \ref{V.eq6.4.3}, but the TeX-file of chapter V contains the
%%% label without "eq".
and one finds that
$$
H^{q}_Y(X,\ZZ/n)\simeq H^{q-2}(Y,\mu_n^{-1}) \mbox{\rm\quad for
all $q$;}
$$
the homomorphism $H^{q-2}(Y,\mu_n^{-1})\rightarrow
H^{q}(X,\ZZ/n)$ furnished by \hyperref[XIV.cor3.3]{3.3} is precisely the
``Gysin homomorphism'', which will be studied later
\footnote[1]{\Cf \SGA 2 XIV for a systematic study of Lefschetz-type
theorems.}.

Theorem \hyperref[XIV.thm3.1]{3.1}~(d) is equivalent\nde{More precisely,
the argument in the text gives the following result. Assertion
\hyperref[XIV.cor3.4]{3.4}~(d') for $d'\leq d$ implies
\hyperref[XIV.thm3.1]{3.1}~(d), and assertion
\hyperref[XIV.thm3.1]{3.1}~(d) implies
\hyperref[XIV.cor3.4]{3.4}~(d). This suffices for the conclusion below.}
to the following statement:

\begin{enonce*}{Corollary 3.4~(d)}\etiquette[3.4]{XIV.cor3.4}
Let $Y'$ be a strictly local scheme that is a strict localization of an
algebraic scheme over a field. Let $f':X'\rightarrow Y'$ be an affine
morphism\pageoriginale of finite type, and let $F'$ be a torsion sheaf on
$X'$ such that $\delta(F')\leq d$ \emph{(\cf
(\hyperref[XIV.def2.2]{2.2}))}. Then $H^{q}(X',F')=0$ for $q>d$.
\end{enonce*}

Indeed, suppose that \hyperref[XIV.cor3.4]{3.4}~(d) is true. Let
$f:X\rightarrow Y$ be an affine morphism of algebraic schemes over $k$,
and let $F$ be a torsion sheaf on $X$ such that $d(F)\leq d$. Let $y$ be
a point of $Y$, and let $\xi$ be a geometric point over $y$. To prove
\sisi{\ref{XIV.subsec3.1}}{\hyperref[XIV.thm3.1]{3.1}}~(d), it is enough
to prove that the stalk $(R^{q}f_{\ast}F)_{\xi}$ vanishes when
$d(y)>d-q$.

Let $f':X'\rightarrow Y'$ be the strict localization of $f$ at the point
$\xi$; that is, let $Y'$ be the strict localization of $Y$ at $\xi$, and
let $X'=\fprod{X}{Y'}{Y}$.

Let $F'$ be the sheaf induced by $F$ on $X'$, so that
$(R^{q}f_{\ast}F)_{\xi}=H^{q}(X',F')$
(\hyperrefext[VIII][VIII.thm5.2]{5.2}).
\sisi{We have
$$
\delta(F')\leq d(F)-d(y)\leq d-d(y)\quad (\hyperref[XIV.prop2.4]{2.4}){\quoi},
$$}
{By (\hyperref[XIV.prop2.4]{2.4}), we have
$$
\delta(F')\leq d(F)-d(y)\leq d-d(y){\quoi},
$$}
%%% I find it irritating that the reference to the proposition from which the above follows,
%%% is referred to in a style that might be confused with an equation number
%%% (although the correct equation number is 3.4.1)
and applying \hyperref[XIV.cor3.4]{3.4}~(d), we obtain
$H^{q}(X',F')=0$ if $q>d-d(y)$, which proves the result.

Conversely, suppose that
\sisi{\hyperref[XIV.subsec2.1]{2.1}}{\hyperref[XIV.thm3.1]{3.1}}~(d) is
true, and let $F'$ and $f':X'\rightarrow Y'$ be as in the statement of
\hyperref[XIV.cor3.4]{3.4}~(d). Since cohomology commutes with filtered
inductive limits (\hyperrefext[VII][VII.prop3.3]{3.3}) and $F'$ is the
inductive limit of its constructible subsheaves
(\hyperrefext[IX][IX.prop2.9]{2.9}), we may suppose that $F'$ is
constructible. Now $Y'$ is the strict localization of an algebraic
scheme $Y$ at a geometric point $\xi$ over a \emph{closed} point $y_0$
(\hyperrefext[X][X.lem3.3]{3.3}). Write
$Y'=\varprojlim_{\alpha}Y_{\alpha}$, where $(Y_{\alpha})$ is a filtered
inverse system of affine schemes étale over $Y$
(\hyperrefext[VIII][VIII.subsec4.5]{4.5}). By
\hyperrefext[IX][IX.cor2.7.4]{2.7.4}, $F'$ is induced by a constructible
torsion sheaf \sisi{}{$F_\alpha$} on
$\sisi{X=\fprod{X}{Y_K}{Y}}{X_\alpha=\fprod{X}{Y_\alpha}{Y}}$ for
sufficiently large $\alpha$. Using \hyperref[XIV.prop2.4]{2.4} and\pageoriginale \EGA IV 9.5.5,
we see immediately that we may further suppose that
$d(F_{\alpha})=\delta(F')\leq d$. Replacing $Y$ by $Y_{\alpha}$, we are
reduced to the case in which $F'$ is induced by a sheaf $F$ on $X$, with
$d(F)\leq d$. Then $R^{q}f_{\ast}F=0$ for $q>\sisi{D}{d}$ by
\sisi{\hyperref[XIV.subsec2.1]{2.1}}{\hyperref[XIV.thm3.1]{3.1}}~(d),
and hence in particular
$$
(R^{q}f_{\ast}F)_{\xi}=H^{q}(X',F')=0\quad{\mbox{\rm if}}\quad
q>d{\quoi},
$$
which proves \hyperref[XIV.cor3.4]{3.4}~(d).

As a special case of \hyperref[XIV.cor3.4]{3.4}~(d), we have:

\begin{enonce*}{Corollary 3.5}\etiquette[3.5]{XIV.cor3.5}
Let $Y'$ be a strictly local scheme of dimension $\leq d$, which is a
strict localization of an algebraic scheme over a field $k$. Let
$U\subset Y'$ be an affine open subset. Then $\cd U\leq d$.
\end{enonce*}

\begin{enonce*}[remark]{Remark 3.6}\etiquette[3.6]{XIV.rem3.6}
It seems very plausible\nde{Ofer Gabber has succeeded in proving the following analogue of \ref{XIV.thm3.1}: let $f:X\to Y$ be an affine morphism of finite type, with $Y$ quasi-excellent (\EGA IV 7.8) and endowed with a dimension function $\delta_Y$. Let $\delta_X$ be the associated rectified dimension function on $X$. If $n$ is invertible on $Y$, then for every sheaf $F$ of $\mathbf{Z}/n$-modules on $X$ one has $\delta_Y(R^q f_\ast F)\leq\delta_X(F)-q$ for every $q\geq0$. Here, as above, for a sheaf $G$ we define $\delta(G)=\sup\{\delta(x)\mid x\in X\quaet G_{\bar{x}}\neq0\}$.
}
that \hyperref[XIV.cor3.4]{3.4}, and hence
\hyperref[XIV.cor3.5]{3.5}, remains valid for \emph{every} local
Noetherian scheme $Y'$ (not necessarily a strict localization of an
algebraic scheme), at least when $Y'$ is \emph{excellent}
(\EGA IV 7.8). This will be proved in
(\sisi{XX 6}{\hyperrefext[XIX][XIX.sec6]{6}}) under the further
assumption that $Y$ has characteristic zero.
\end{enonce*}

% END INCLUDED SGA4_English_current_buildable_source_20260906/source/14/drafts/SGA4_XIV_LINES_723_863_English.texfrag


% BEGIN INCLUDED SGA4_English_current_buildable_source_20260906/source/14/drafts/SGA4_XIV_LINES_864_1076_English.texfrag
\section[Proof of Theorem 3.1]{Proof of Theorem \texorpdfstring{\hyperref[XIV.thm3.1]{3.1}}{3.1}}\etiquette[4]{XIV.sec4}

The proof proceeds by induction on $d$.

\begin{enonce*}{Lemma 4.1}\etiquette[4.1]{XIV.lem4.1}
\emph{\hyperref[XIV.thm3.1]{3.1}~(0)} and
\emph{\hyperref[XIV.thm3.1]{3.1}~(1)} hold.
\end{enonce*}

\noindent
\emph{Proof}. This is trivial for $d=0$, in which case $\supp F$ is
finite over $k$, and we apply
(\hyperrefext[VIII][VIII.prop5.5]{5.5}). Let $F$ and
$f:X\rightarrow Y$ be as in the statement of
\hyperref[XIV.thm3.1]{3.1}, and suppose that $d(F)\leq1$. Replacing
$F$ by a constructible subsheaf
(\hyperrefext[IX][IX.prop2.9]{2.9} and
\hyperrefext[VII][VII.prop3.3]{3.3}), we reduce to the case in which
$F$ is constructible. The support\pageoriginale of $F$ is then
contained in a closed subscheme of $X$ of dimension $\leq1$
(\hyperrefext[IX][IX.def2.3]{2.3}\sisi{ (ii)}{}), and we may replace
$X$ by this subscheme; that is, we may assume that $\dim X\leq1$.
By \hyperrefext[IX][IX.cor5.7]{5.7}, the cohomological dimension of an
affine algebraic scheme of dimension $1$ over a separably closed field
is $\leq1$. It is enough, moreover, to treat the case in which $k$ is
separably closed. Indeed, let $\bar f:\bar X\rightarrow\bar Y$ be the
morphism obtained from $f$ by the base change
$\Spec\bar k\rightarrow\Spec k$, where $\bar k$ is the separable
closure of $k$. For a geometric point $\xi$ of $Y$, the strict
localization $f':X'\rightarrow Y'$ of $f$ at $\xi$ identifies with a
strict localization of $\bar f$, and consequently the fiber
$(R^{q}f_{\ast}F)_\xi$ identifies with a fiber of
$R^{q}\bar f_{\ast}\bar F$. Since $f$ is affine, it follows that
$R^{q}f_{\ast}F=0$ for $q>1$. It is also clear that
$d(f_{\ast}F)\leq1$. For
$q=1$, note that the hypothesis $\dim X\leq1$ implies that there is a
finite closed subset $Z$ of $Y$ such that $f$ is finite over $Y-Z$.
It follows that $R^{1}f_{\ast}F$ vanishes on $Y-Z$
(\hyperrefext[VIII][VIII.cor5.6]{5.6}), and hence
$d(R^{1}f_{\ast}F)\leq1$. This proves the lemma.

Let $d\geq2$, and now suppose that
\hyperref[XIV.thm3.1]{3.1}~(d') (or, equivalently,
\hyperref[XIV.cor3.4]{3.4}~(d')) holds for $d'<d$.

\begin{enonce*}{Lemma 4.2}\etiquette[4.2]{XIV.lem4.2}
To prove \emph{\hyperref[XIV.cor3.4]{3.4}~(d)} in general, it is
enough to treat the case in which $X'=\EE^1_{Y'}$ is affine
one-space over $Y'$ and $f$ is the structural morphism.
\end{enonce*}

\noindent
\emph{Proof}. We may plainly assume that $X'=\EE^N_{Y'}$ for a
suitable $N$, since $X'$ can be embedded in some $\EE^N_{Y'}$.
Suppose that $N>1$ and that the result is known, for every $Y'$, for
$\EE^r_{Y'}$ with $r<N$. Let $F'$ be a torsion sheaf on
$\EE^N_{Y'}$ with\pageoriginale $d(F')\leq d$. Consider the diagram
$$
\xymatrix@=1.5cm{
\EE^{N}_{Y'} \ar[rd]_{f'}       \ar[rr]^{g'}&      &
\EE^{N-1}_{Y'} \ar[ld]^{h'}\\
                       & Y' & &,
}
$$
where $g'$ and $h'$ are the canonical projections, so that
$\EE^N_{Y'}$ is isomorphic to affine one-space over
$\EE^{N-1}_{Y'}$. Since every strict localization of
$\EE^{N-1}_{Y'}$ is plainly also the strict localization of an
algebraic scheme, we may apply the induction hypothesis to the fibers
of $R^{q}g'_{\ast}F'$, and obtain
\begin{equation*}
\delta(R^{q}g'_{\ast}F')\leq d-q{\quoi}.\tag{$\ast$}
\end{equation*}
Moreover, \hyperref[XIV.cor3.4]{3.4}~(d') holds for $h'$ and
$d'\leq d$, by the induction hypothesis on $N$.

Consider the Leray spectral sequence
$$
E^{p,q}_2=H^{p}(\EE^{N-1}_{Y'},R^{q}g'_{\ast}(F'))
\Longrightarrow H^{n}(\EE^N_{Y'},F'){\quoi}.
$$
We have $E^{p,q}_2=0$ if $p>d-q$, as follows by applying $(\ast)$ and
\hyperref[XIV.cor3.4]{3.4}~(d') to the morphism $h'$. Consequently the
abutment vanishes for $n>d$, proving the lemma.

Consider the following additional statement, which is a special case
of \hyperref[XIV.cor3.5]{3.5}:

\begin{enonce*}{Statement 4.3~(d)}\etiquette[4.3]{XIV.enon4.3}
Let $X'$ be strictly local of dimension $\leq d$, and the strict
localization of an algebraic scheme over a field. Let
$f\in\Gamma(X',\Oo_{X'})$, and let $U\subset X'$ be the open subset
$U=X'-V(f)$. Then $\cd U\leq d$.
\end{enonce*}

\begin{enonce*}{Lemma 4.4}\etiquette[4.4]{XIV.lem4.4}
\emph{\hyperref[XIV.enon4.3]{4.3}~(d)} implies
\emph{\hyperref[XIV.cor3.4]{3.4}~(d)}.
\end{enonce*}

\noindent
\emph{Proof}.\pageoriginale Consider the diagram
$$
\xymatrix@=1.5cm{
\EE^1_{Y'} \ar[rr]^i \ar[rd]_{f'}&            &
\PP^1_{Y'} \ar[ld]^{g'}\\
                    &       Y'   &
}
$$
where $Y'$ is as in the statement of
\sisi{2.6}{\hyperref[XIV.cor3.4]{3.4}}, and where $\PP^1_{Y'}$ is the
projective line obtained from $\EE^1$ ``by adjoining the section
$Y'_\infty$ at infinity.'' Let $F'$ be a torsion sheaf on $\EE^1$
with $d(\supp F')\leq d$. We have the spectral sequence
$$
H^{p}(\PP^1,R^{q}i_{\ast}F')\Longrightarrow
H^{n}(\EE^1,F'){\quoi}.
$$

Now the $R^{q}i_{\ast}F'$, for $q>0$, are supported on $Y'_\infty$, which is
$Y'$-isomorphic to $Y'$. Since $Y'$ is strictly local, it follows that
$H^{p}(\PP^1,R^{q}i_{\ast}F')=0$ if $p,q>0$. Moreover, by the \emph{proper
base-change theorem} (\hyperrefext[XII][XII.thm5.1]{5.1}), the
cohomology of a torsion sheaf on $\PP^1$ can be computed on the closed
fiber, which is a scheme of cohomological dimension $\leq2$
(\hyperrefext[X][X.cor4.3]{4.3}). Thus
$H^{p}(\PP^1,i_{\ast}F')=0$ if $p>2$. Since we wish to prove that
$H^{n}(\EE^1,F')=0$ for $n>d$, and since $d\geq2$, we are reduced to
proving that $H^{0}(\PP^1,R^{q}i_{\ast}F')=0$ for $q>d$. This group, which is
also isomorphic to $H^{0}(Y'_\infty,R^{\sisi{1}{q}}i_{\ast}F')$, is, by
(\hyperrefext[VIII][VIII.prop4.6]{4.6}), isomorphic to the fiber of
$R^{q}i_{\ast}F'$ at the point $Q$, where $Q$ is the point at infinity of
$\PP^1$ in the closed fiber.

It is enough (\hyperrefext[VII][VII.prop3.3]{3.3} and
\hyperrefext[IX][IX.prop2.9]{2.9}) to treat the case in which $F'$ is
constructible. Since $\delta(\supp F')\leq d$, there is then a closed
subscheme $X'$ of $\PP^1$, with $\dim X'\leq d$, which contains the
support of $F'$. Let $\widetilde X'$ be the localization of $X'$ at
the geometric point $Q$, put
$$
\widetilde U=\widetilde X'-
\fprod{\widetilde X'}{Y'_\infty}{\PP^1},
$$
and let $\widetilde F$ be the sheaf\pageoriginale on $\widetilde U$ induced by $F'$
via the morphism $\widetilde U\rightarrow\PP^1$. The fiber of
$R^{q}i_{\ast}F'$ at $Q$ is then precisely
$H^{\sisi{p}{q}}(\widetilde U,\widetilde F)$
(\hyperrefext[VIII][VIII.thm5.2]{5.2}), and we apply
\hyperref[XIV.enon4.3]{4.3}~(d) to the strictly local scheme
$\widetilde X'$ and the open subset $\widetilde U$.

The following lemma, together with \hyperref[XIV.lem4.4]{4.4}, will
complete the proof.

\begin{enonce*}{Lemma 4.5}\etiquette[4.5]{XIV.lem4.5}
\emph{\hyperref[XIV.cor3.4]{3.4}~(d')} for $d'<d$ implies
\emph{\hyperref[XIV.enon4.3]{4.3}~(d')} for $d'\leq d$.
\end{enonce*}

\noindent
\emph{Proof}. Suppose that \hyperref[XIV.cor3.4]{3.4}~(d') holds for
$d'<d$. By induction, we may suppose that
\hyperref[XIV.enon4.3]{4.3}~(d') holds for $d'<d$. Let $X'$, $f$, and
$U$ be as in the statement \hyperref[XIV.enon4.3]{4.3}~(d). Then
$X'$ is the strict localization, at a \emph{closed} geometric point,
of an algebraic scheme $X_0$ over a field $k$. After replacing $X_0$
by a neighborhood of $x_0$, we may take $X_0$ to be affine and of
dimension $\leq d$. Then $X'$ is the inverse limit of affine schemes
$X_i$ which are étale over $X_0$
(\hyperrefext[VIII][VIII.subsec4.5]{4.5}).

Let $F$ be a torsion sheaf on $U$. We must prove that $H^q(U,F)=0$ for
$q>d$, and once again it is enough to treat the case in which $F$ is
constructible. The data $f$, $F$, and $U$ then arise from analogous
data on one of the $X_i$
(\sisi{\hyperrefext[IX][IX.prop2.7]{2.7}}{\hyperrefext[IX][IX.cor2.7.4]{2.7.4}}).
Thus there is a sheaf $F_i$ on $U_i=X_i-V(f_i)$ which induces $F$ on
$U$. Since $U=\varprojlim U_i$, we have
(\hyperrefext[VII][VII.cor5.8]{5.8})
$$
H^q(U,F)=\varinjlim H^q(U_i,F_i){\quoi}.
$$
It is therefore enough to prove the following fact: for every $i$,
there is a commutative diagram
$$
\xymatrix@=1.5cm{
 U \ar[rr] \ar[rd]&        & U''\ar[ld]\\
   &  U_i   &
}
$$
in which $U\rightarrow U_i$ is the given arrow and
$\cd(U'')\sisi{<}{\leq}d$. Indeed, this implies that\pageoriginale every element of
$H^q(U_i,F_i)$, for $q>d$, has zero image in $H^q(U,F)$.

Suppose that $i=0$, and consider the morphism
$g:X_0\rightarrow\EE^1_k=Y_0$ defined by
$f_0\in\Gamma(X_0,\Oo_{X_0})$. Let $Y''$ be the strict localization
of $\EE^1_k$ at a geometric point above the origin. Thus $Y''$ is the
spectrum of a discrete valuation ring. Put
$X''=\fprod{X_0}{Y''}{Y_0}$.

The morphism $X'\rightarrow X_0$ factors through $X''$, since $X'$ is
strictly local. Moreover, the open subset
$U''=\fprod{X''}{U_0}{X_0}$ is precisely the generic fiber of
$X''/Y''$. It is therefore an affine algebraic scheme over the generic
point of $Y''$, and plainly $\dim U''\leq d-1$.

Let us verify that $\cd U''\leq d$, which will complete the proof. Let
$K$ be the residue field of $Y''$ at its generic point. This field has
cohomological dimension $1$; that is, $\cd(G)=1$, where $G$ is the
Galois group of the separable closure $\bar K$ of $K$
(\hyperrefext[X][X.thm2.2]{2.2}). Put
$\bar U''=\fprod{U''}{(\Spec\bar K)}{\Spec K}$. By the induction
hypothesis on $d$, we have $\cd\bar U''\leq d-1$ (this is
\sisi{\hyperref[XIV.def2.2]{2.2} for a scheme of dimension $\leq d$}
{\hyperref[XIV.cor3.2]{3.2} for a scheme of dimension $<d$}). It then
follows from the Hochschild--Serre spectral sequence
(\hyperrefext[VIII][VIII.prop8.4]{8.4}) that $\cd U''\leq d$:
$$
H^p(G,H^q(\bar U\sisi{'}{''},\bar F))\Longrightarrow H^n(U'',F){\quoi},
$$
which proves Lemma \hyperref[XIV.lem4.5]{4.5}.

\ifx\danslelivre\undefined
\end{document}
\fi

% END INCLUDED SGA4_English_current_buildable_source_20260906/source/14/drafts/SGA4_XIV_LINES_864_1076_English.texfrag


% END INCLUDED SGA4_English_current_buildable_source_20260906/source/14/14_English.tex


\makeatletter
\def\input@path{{../15/}}
\makeatother
% BEGIN INCLUDED SGA4_English_current_buildable_source_20260906/source/15/15_English.tex
\ifx\danslelivre\undefined
% Release engine: XeLaTeX (the LuaTeX compatibility layer breaks Xy-pic output).
\documentclass[10pt]{../fichiers-configuration-latex/smfart}
% BEGIN INCLUDED SGA4_English_current_buildable_source_20260906/source/fichiers-configuration-latex/paquets.tex
\usepackage{polyglossia}
%\usepackage[francais]{babel}
\usepackage{fontspec}
\usepackage[math-style=ISO,bold-style=upright,sans-style=upright]{unicode-math}
%\usepackage[mathletters]{ucs}
%\usepackage[utf8]{inputenc}
%\usepackage[T1]{fontenc}
\usepackage[unicode,hypertexnames=false,colorlinks=true,linkcolor=blue,citecolor=Sepia,urlcolor=magenta]{hyperref}
\setmainlanguage{french}
\setmainfont[Ligatures=TeX]{Linux Libertine O}
% variante possible : TeXGyrePagella ?
\setmathfont{XITS Math}

% maths
\usepackage{amsmath}
%\usepackage{amssymb}
%\usepackage{mathrsfs}
%\usepackage{stmaryrd}
%\usepackage{euscript}

% couleurs
\usepackage[usenames,dvipsnames]{xcolor}

% graphiques
\usepackage{graphicx}
\usepackage[all]{xy}
%\usepackage{tikz}
%\usetikzlibrary{matrix,arrows,calc}
%\tikzset{column sep=3em,row sep=2em,text height=1.5ex, text depth=0.25ex,>=angle 90}
%\tikzset{description/.style={fill=white,inner sep=1pt}}

% hyperliens, 'etiquettes
\usepackage{xr-hyper}
\usepackage{../fichiers-configuration-latex/smfthm}
%\ifx\cacheretiquettes\undefined
%\usepackage[inline]{showlabels} \else \fi % pour rendre visibles les 'etiquettes
%\usepackage{../fichiers-configuration-latex/makeidx} % pour faire un index
%\usepackage[pagebackref,unicode]{hyperref}  % pour voir o`u sont cit'ees les r'ef'erences
% hyperref is loaded once above with the cumulative-reader destination policy.
%\hypersetup{breaklinks=true} %permet le retour `a la ligne dans les liens trop longs. Incompatible avec pdflatex
%\renewcommand*{\backref}[1]{{\small [cit\'e page(s)~#1.]}}

\usepackage{xspace}
\usepackage{pifont} % pour expose XVII

% END INCLUDED SGA4_English_current_buildable_source_20260906/source/fichiers-configuration-latex/paquets.tex

% BEGIN INCLUDED SGA4_English_current_buildable_source_20260906/source/fichiers-configuration-latex/numerotation.tex
\renewcommand{\thefootnote}{\roman{footnote}} % pour mettre les notes en bas de page sous la forme (ii) et non (2), qui ressemble trop (?) a un exposant.

\newcounter{pageoriginal}
\newcommand{\pageoriginale}{\refstepcounter{pageoriginal}\marginpar{\footnotesize\textbf{\thepageoriginal}}\xspace}
\let\pageoriginaled\pageoriginale

\usepackage{ifthen}
\newboolean{original}
\def\sisi{\ifthenelse{\boolean{original}}}
\setboolean{original}{false}

\newcounter{cnote}
\def\nde#1{\ifthenelse{\boolean{original}}{}{\stepcounter{cnote}\def\thefootnote{\arabic{cnote}}\protect\footnote[\value{cnote}]{N.D.E.: #1}}}
\def\ndemark{\ifthenelse{\boolean{original}}{}{\stepcounter{cnote}\def\thefootnote{\arabic{cnote}}\protect\footnotemark[\value{cnote}]}}
\def\ndetext#1{\ifthenelse{\boolean{original}}{}{\def\thefootnote{\arabic{cnote}}\protect\footnotetext[\value{cnote}]{N.D.E.: #1}}}

% END INCLUDED SGA4_English_current_buildable_source_20260906/source/fichiers-configuration-latex/numerotation.tex

% BEGIN INCLUDED SGA4_English_current_buildable_source_20260906/source/fichiers-configuration-latex/macros.tex
\setmainlanguage{french}
\setotherlanguage{english}


%% Fixed, engine-independent release identifier.  The inherited dynamic Git
%% probe used Lua io.popen, exposed build-state-dependent output, and printed
%% a visible question mark outside a Git checkout.
\def\showgitstatus{English release r2, 2026-07-23}


% macros pour LaTeX
\renewcommand{\mathsf}{}

% Police différente pour U+2113 SCRIPT SMALL L
\setmathfont[range="2113]{latinmodern-math.otf}

% Give every inherited semantic label its own stable PDF destination instead
% of inheriting whichever automatic equation/page anchor happened to be live.
% The legacy label remains as a compatibility alias used by existing edges.
\makeatletter
\def\etiquette[#1]#2{%
  \Hy@raisedlink{\hyper@anchorstart{sga4.semantic.#2}\hyper@anchorend}%
  \def\@currentHref{sga4.semantic.#2}%
  \def\@currentlabel{#1}%
  \label{#2}%
}
\makeatother
\def\eqhyperref[#1]#2{(\hyperref[#1]{#2})}
\def\pagehyperref[#1]#2{(page \hyperref[#1]{#2})}
\def\hyperrefext[#1][#2]#3{\hyperref[#2]{#1 #3}}
\def\eqhyperrefext[#1][#2]#3{\hyperref[#2]{#1 #3}}     %cf. exposes 17
\def\neweqhyperrefext[#1][#2]#3{\hyperref[#2]{#1 #3}}  %cf. exposes 17

\def\shat#1{\protect\widehat{#1}}
\def\stilde#1{\protect\widetilde{#1}}
%%%%%%%%%%%%%%%%%%


\def\etcn{etc.\xspace}
\def\etco{etc...\xspace}
\def\etc{\sisi{\etco}{\etcn}}
\def\ldotsn{\sisi{...}{\etcn}}

\def\cf{{cf.}\kern.3em}
\def\Cf{{Cf.}\kern.3em}
\def\ie{{i.e.}\kern.3em}
\def\iev{{i.e.},\xspace}
\def\loccito{loc.\kern2pt cit.\xspace}
\def\loccitn{\textit{loc.\kern2pt cit.}\xspace}
\def\loccit{\sisi{\loccito}{\loccitn}}
\def\EGA{\textup{EGA}\kern.2em}
\def\SGA{\textup{SGA}\kern.3em}

\def\quoi{}

\def\quav{\quad,\quad}
\def\quaet{\quad \text{et}\quad}
\def\fibreproduct#1#2#3{#1 ×_{#3} #2}
\let\fprod\fibreproduct
\def\ccheck#1#2{\mbox{${\check{\check{\rm #1}}}^{\lower2pt\hbox{$\scriptstyle #2$}}$}}
\def\sheaf#1{\protect\mathcal{#1}}
\def\vieuxsheaf#1{{#1}} % cf. problème avec sheaf{a\Phi_i} par exemple...
\DeclareMathOperator{\Hom}{Hom}
\DeclareMathOperator{\can}{can}
\DeclareMathOperator{\Ind}{Ind}
\DeclareMathOperator{\ind}{ind}
%\DeclareMathOperator{\Hhomtop}{\mathcal{H}\!\mathit{omtop}}
\DeclareMathOperator{\Hhomtop}{\mathbf{Homtop}}
\DeclareMathOperator{\Point}{Point}
\DeclareMathOperator{\Pro}{Pro}
\DeclareMathOperator{\Top}{Top}
\DeclareMathOperator{\Mod}{Mod}
\DeclareMathOperator{\Cart}{Cart}
%\DeclareMathOperator{\Mmorsite}{\mathcal{M}\!\mathit{orsite}}
\DeclareMathOperator{\Mmorsite}{\mathbf{Morsite}}
\DeclareMathOperator{\loc}{loc}
\DeclareMathOperator{\Loc}{Loc}
\DeclareMathOperator{\Pic}{Pic}
\DeclareMathOperator{\Et}{Et}
\DeclareMathOperator{\card}{card}
\DeclareMathOperator{\Ens}{Ens}
\DeclareMathOperator{\h}{\underline{H}}
\DeclareMathOperator{\Fonct}{Fonct}
\DeclareMathOperator{\Ker}{Ker}
\DeclareMathOperator{\but}{but}
\DeclareMathOperator{\Ob}{Ob}
\DeclareMathOperator{\ob}{ob}
\DeclareMathOperator{\Lib}{Lib}
\DeclareMathOperator{\esj}{esj}
\DeclareMathOperator{\id}{id}
\DeclareMathOperator{\Id}{Id}
\DeclareMathOperator{\ad}{ad}
\DeclareMathOperator{\Iim}{Im}
\DeclareMathOperator{\pr}{pr}
\DeclareMathOperator{\Cat}{Cat}
\DeclareMathOperator{\car}{car}
\DeclareMathOperator{\SysInd}{SysInd}
\DeclareMathOperator{\Sup}{Sup}
\DeclareMathOperator{\Coker}{Coker}
\DeclareMathOperator{\Filt}{Filt}
\DeclareMathOperator{\acc}{acc}
\DeclareMathOperator{\sur}{sur}
\DeclareMathOperator{\cat}{cat}
\DeclareMathOperator{\Ord}{Ord}
\DeclareMathOperator{\Gr}{Gr}
\DeclareMathOperator{\Cov}{Cov}
\DeclareMathOperator{\Caf}{Caf}
\DeclareMathOperator{\coker}{coker}
\DeclareMathOperator{\coim}{coim}
\DeclareMathOperator{\im}{im}
\DeclareMathOperator{\Ouv}{Ouv}
\DeclareMathOperator{\diag}{diag}
\DeclareMathOperator{\Esp}{Esp}
\DeclareMathOperator{\TOP}{TOP}
\DeclareMathOperator{\Unif}{Unif}
\DeclareMathOperator{\Iind}{\mathbf{Ind}}
\DeclareMathOperator{\Hhomcart}{\mathbf{Homcart}}
\DeclareMathOperator{\Hhomtopess}{\mathbf{Homtopess}}
\DeclareMathOperator{\Ppoints}{\mathbf{Points}}
\DeclareMathOperator{\Ppoint}{\mathbf{Point}}
\DeclareMathOperator{\Ttors}{\mathbf{Tors}}
\DeclareMathOperator{\Ssex}{\mathbf{Sex}}
\DeclareMathOperator{\Ffib}{\mathbf{Fib}}
\DeclareMathOperator{\Vvois}{\mathbf{Vois}}
\DeclareMathOperator{\Ppointess}{\mathbf{Pointess}}
\DeclareMathOperator{\scrHom}{\mathbf{Hom}}
\DeclareMathOperator{\sheafPic}{\mathbf{Pic}}
\DeclareMathOperator{\sheafAut}{\mathbf{Aut}}
\DeclareMathOperator{\Hhomtoplocan}{\mathbf{Homtoplocan}}
\DeclareMathOperator{\Hhomtopan}{\mathbf{Homtopan}}
\DeclareMathOperator{\Pptgeom}{\mathbf{Ptgéom}}
\DeclareMathOperator{\Morsite}{Morsite}
\DeclareMathOperator{\Fib}{Fib}
\DeclareMathOperator{\Fais}{Fais}
\DeclareMathOperator{\kar}{kar}
\DeclareMathOperator{\Groupes}{Groupes}
\DeclareMathOperator{\Restr}{Restr}
\DeclareMathOperator{\Isom}{Isom}
\DeclareMathOperator{\Prol}{Prol}
\DeclareMathOperator{\Res}{Res}
\DeclareMathOperator{\Tors}{Tors}
\DeclareMathOperator{\esp}{esp}
\DeclareMathOperator{\es}{es}
\DeclareMathOperator{\od}{od}
\DeclareMathOperator{\gr}{gr}
\DeclareMathOperator{\sob}{sob}
\DeclareMathOperator{\dfn}{dfn}
\DeclareMathOperator{\Sch}{Sch}
\DeclareMathOperator{\ouv}{ouv}
\DeclareMathOperator{\cons}{cons}
\DeclareMathOperator{\Nerf}{Nerf}
\DeclareMathOperator{\Points}{Points}
\DeclareMathOperator{\Espsob}{Espsob}
\DeclareMathOperator{\Homtop}{Homtop}
\DeclareMathOperator{\Pt}{Pt}
\DeclareMathOperator{\Quot}{Quot}
\DeclareMathOperator{\glob}{glob}
\DeclareMathOperator{\Homfib}{Homfib}
\DeclareMathOperator{\Qouv}{Qouv}
\DeclareMathOperator{\Quotouv}{Quotouv}
\DeclareMathOperator{\ens}{ens}
\DeclareMathOperator{\ab}{ab}
\DeclareMathOperator{\an}{an}
\DeclareMathOperator{\fs}{fs}
\DeclareMathOperator{\red}{red}
\DeclareMathOperator{\Ab}{Ab}
\DeclareMathOperator{\cosq}{cosq}
\DeclareMathOperator{\Ann}{Ann}
\DeclareMathOperator{\Zar}{Zar}


%%%%%%%%%%%%
%%%%% SGA4-2
\DeclareMathOperator{\Sex}{Sex}
\DeclareMathOperator{\vois}{vois}
\DeclareMathOperator{\Vois}{Vois}
\DeclareMathOperator{\Ext}{Ext}
\DeclareMathOperator{\ext}{ext}
\DeclareMathOperator{\Fer}{Fer}
\DeclareMathOperator{\Add}{Add}
\DeclareMathOperator{\Cocart}{Cocart}
\DeclareMathOperator{\cosk}{cosk}
\DeclareMathOperator{\plat}{plat}

\DeclareMathOperator{\faisc}{faisc}
\DeclareMathOperator{\sk}{sk}
\DeclareMathOperator{\ch}{ch}

\def\HH{\mathbf{H}}
\def\qc{\mathrm{qc}}
\def\et{\mathrm{et}}
\def\coh{\mathrm{coh}}

\def\Hvr#1{%
\begin{array}[b]{l}
{\kern2pt\scriptstyle r}\\[-8pt]
{\kern1pt\scriptstyle \vee}\\[-5pt]
H^{#1}
\end{array}}

\def\Hvv#1{%
\begin{array}[b]{l}
\kern2pt\rotatebox{90}{${\scriptscriptstyle\ll}$}\\[-6pt]
H^{#1}
\end{array}}

\def\Hvi#1{%
\begin{array}[b]{l}
{\kern0.2pt\scriptscriptstyle \infty}\\[-8pt]
{\kern1pt\scriptstyle \vee}\\[-5pt]
H^{#1}
\end{array}}

\def\Hv#1{%
\begin{array}[b]{l}
{\kern1pt\scriptstyle \vee}\\[-5pt]
H^{#1}
\end{array}}

\DeclareMathOperator{\cart}{cart}
\DeclareMathOperator{\Ffaisc}{\mathbf{Faisc}}


\def\Ev#1{%
\begin{array}[b]{l}
{\kern1pt\scriptstyle \vee}\\[-5pt]
E'_{#1}
\end{array}}

\DeclareMathOperator{\Cont}{Cont}
\DeclareMathOperator{\Ccont}{\mathscr{C}\!\mathit{ont}}
\DeclareMathOperator{\Site}{Site}
\newcommand\MM{\mathbf{M}}
\def\limtop#1{{\displaystyle{\varprojlim_{#1}\!\mbox{\rm top }}}}
\def\Limtop{\mathop{\varprojlim\!\mbox{\rm top}}}
\DeclareMathOperator{\spec}{spec}
\def\ss{\scriptstyle{\mathscr{S}}}
\DeclareMathOperator{\Gal}{Gal}
\def\smallm{\scriptstyle{\mathscr{M}}}
\DeclareMathOperator{\End}{End}
\DeclareMathOperator{\ass}{ass}
\DeclareMathOperator{\zar}{zar}
\DeclareMathOperator{\Ogg}{Ogg}
\DeclareMathOperator{\res}{res}
\DeclareMathOperator{\tr}{tr}
\DeclareMathOperator{\Supp}{Supp}
\DeclareMathOperator{\codim}{codim}
\DeclareMathOperator{\rk}{rk}
\DeclareMathOperator{\cur}{cur}
\DeclareMathOperator{\cdqc}{cdqc}
\newcommand\DD{\mathbf{D}}
\DeclareMathOperator{\cl}{cl}
\def\JJ{\mathbf{J}}
\def\FF{\mathbf{F}}
\def\BB{\mathbf{B}}
\DeclareMathOperator{\Alget}{Alget}
\DeclareMathOperator{\Ass}{Ass}
\DeclareMathOperator{\supp}{supp}
\DeclareMathOperator{\rad}{rad}
\DeclareMathOperator{\prof}{prof}
\DeclareMathOperator{\cint}{int}
\DeclareMathOperator{\sic}{sic}
\DeclareMathOperator{\Tor}{Tor}
\DeclareMathOperator{\tors}{tors}
\DeclareMathOperator{\Sym}{Sym}
\DeclareMathOperator{\parf}{parf}
\DeclareMathOperator{\tor}{tor}
\DeclareMathOperator{\defn}{defn}
\DeclareMathOperator{\Div}{Div}
\DeclareMathOperator{\Pol}{Pol}
\DeclareMathOperator{\Fex}{Fex}
\DeclareMathOperator{\Lie}{Lie}
\DeclareMathOperator{\tot}{tot}
\DeclareMathOperator{\Aut}{Aut}
\DeclareMathOperator{\HOM}{HOM}
\DeclareMathOperator{\pch}{pch}
\DeclareMathOperator{\Ch}{Ch}
\DeclareMathOperator{\EXT}{EXT}
\DeclareMathOperator{\PIC}{PIC}
\DeclareMathOperator{\Rf}{Rf}
\DeclareMathOperator{\inj}{inj}
\newcommand\VV{\mathbf{V}}

\def\Xv{%
\begin{array}[b]{l}
{\kern1pt\scriptstyle \vee}\\[-5pt]
X
\end{array}}

%%%%%%%
\newcommand\TT{\mathbf{T}}
\newcommand\XX{\mathbf{X}}
\newcommand\EE{\mathbf{E}}
%%%%%%

\def\qedsymbol{\textup{C.Q.F.D.}}
\let\CQFD\qed
\let\cqfd\qed

%%%%%\bigdot
\font\bigsymb=cmsy10 at 4pt
\def\bigdot{{\kern1.2pt\raise 1.5pt\hbox{\bigsymb\char15}}}
%%%%%

\def\Lright{\mathop{\vcenter{\hbox{\hbox{$\longrightarrow$}\lower8pt \hbox{\kern-16pt\hbox{$\longrightarrow$}}}}}}

\def\Aa{{\mathcal A}}
\def\Bb{{\mathcal B}}
\def\Cc{{\mathcal C}}
\def\Dd{{\mathcal D}}
\def\Ee{{\mathcal E}}
\def\Ff{{\mathcal F}}
\def\Gg{{\mathcal G}}
\def\Hh{{\mathcal H}}
\def\Ii{{\mathcal I}}
\def\Jj{{\mathcal J}}
\def\Kk{{\mathcal K}}
\def\Ll{{\mathcal L}}
\def\Mm{{\mathcal M}}
\def\Nn{{\mathcal N}}
\def\Oo{{\mathcal O}}
\def\Pp{{\mathcal P}}
\def\Qq{{\mathcal Q}}
\def\Rr{{\mathcal R}}
\def\Ss{{\mathcal S}}
\def\Tt{{\mathcal T}}
\def\Uu{{\mathcal U}}
\def\Vv{{\mathcal V}}
\def\Ww{{\mathcal W}}
\def\Xx{{\mathcal X}}
\def\Yy{{\mathcal Y}}
\def\Zz{{\mathcal Z}}

\newcommand{\CC}{\mathbf{C}}
\newcommand{\GG}{\mathbf{G}}
\newcommand{\KK}{\mathbf{K}}
\newcommand{\LL}{\mathbf{L}}
\newcommand{\NN}{\mathbf{N}}
\newcommand{\PP}{\mathbf{P}}
\newcommand\QQ{\mathbf{Q}}
\newcommand\RR{\mathbf{R}}
\newcommand\ZZ{\mathbf{Z}}

\def\to{\mathchoice{\longrightarrow}{\rightarrow}{\rightarrow}{\rightarrow}}
\def\from{\mathchoice{\longleftarrow}{\leftarrow}{\leftarrow}{\leftarrow}}
\def\mto{\mathchoice{\longmapsto}{\mapsto}{\mapsto}{\mapsto}}
\def\mfrom{\mathrel{\from\joinrel\mapstochar\,}}
\def\hto{\mathrel{\lhook\joinrel\to}}
\def\hfrom{\mathrel{\from\joinrel\rhook}}

\def\implies{\mathchoice{\Longrightarrow}{\Rightarrow}{\Rightarrow}{\Rightarrow}}
\def\implied{\mathchoice{\Longleftarrow}{\Leftarrow}{\Leftarrow}{\Leftarrow}}
\def\iff{\mathchoice{\Longleftrightarrow}{\Leftrightarrow}{\Leftrightarrow}{\Leftrightarrow}}

\def\To#1{\mathchoice{\xrightarrow{\labelstyle~~#1~}}{\stackrel{#1}{\longrightarrow}}{}{}}
\def\Mto#1{\mathrel{\mapstochar\To{#1}}}
\def\From#1{\mathchoice{\xleftarrow{\labelstyle~~#1~}}{\stackrel{#1}{\longleftarrow}}{}{}}

\newcommand{\isomto}{\overset{\sim}{\longrightarrow}}
\newcommand{\isomfrom}{\overset{\sim}{\longleftarrow}}

\newcommand{\IFF}{\unskip~$\iff$~\ignorespaces}
\newcommand{\THEN}{\unskip~$\implies$~\ignorespaces}

\let\leq\leqslant
\let\geq\geqslant
\let\le\leqslant
\let\ge\geqslant

\def\cf{{cf.}\kern.3em}
\def\Cf{{Cf.}\kern.3em}
\def\ie{{i.e.}\kern.3em}
\def\iev{{i.e.},\xspace}
%\def\loccito{loc.\kern2pt cit.\xspace}
%\def\loccitn{\textit{loc.\kern2pt cit.}\xspace}
%\def\loccit{\sisi{\loccito}{\loccitn}}
\def\EGA{\textup{EGA}\kern.2em}
\def\SGA{\textup{SGA}\kern.3em}

\let\oldS\S
\def\S{\oldS\kern .15em}
\def\ptbl{.\kern .2em}
\def\resp{resp.\kern.3em}
\def\Exp{Exp.\kern.3em}

\let\bar\overline
\let\hat\widehat
\let\tilde\widetilde
\let\dpl\displaystyle
\let\ldots\dots
\let\emptyset\varnothing
\let\skpt\skippointrait
\let\epsilon\varepsilon

\def\sfrac#1#2{{#1}/{#2}}

\DeclareMathOperator{\cd}{cd}
%\DeclareMathOperator{\Hhom}{\mathcal{H}\!\mathit{om}}
\DeclareMathOperator{\Hhom}{\mathbf{Hom}}
%\DeclareMathOperator{\Eext}{\mathcal{E}\!\mathit{xt}}
\DeclareMathOperator{\Eext}{\mathbf{Ext}}
\DeclareMathOperator{\Spec}{Spec}
\DeclareMathOperator{\Tr}{Tr}

\let\emptyset=\varnothing

%%%%%%%%%%%%%%%
%%% expose XVII
%%%%%%%%%%%%%%%

\makeatletter
\let\old@optexpose\optexpose
\newcommand{\empty@optexpose}[1]{}
\newcommand{\paren@optexpose}[1]{\old@optexpose{#1}(}
\newcommand{\refnoexpose}[1]{\let\optexpose\empty@optexpose\ref{#1}\let\optexpose\old@optexpose}%\refnoexpose{X.cor5.2} engendre "5.2" au lieu de "X 5.2"
\newcommand{\neweqref}[1]{\let\optexpose\paren@optexpose\textup{\ref{#1})}\let\optexpose\old@optexpose}%\neweqref{XII.eq4.1} engendre "XII (4.1)" au lieu de "(XII 4.1)"
\makeatother
\newcommand{\refvb}[2][VI${}_{\text{B}}$ ]{\sisi{#1\refnoexpose{#2}}{\ref{#2}}}

\newcommand{\leftexp}[2]{{\vphantom{#2}}^{#1}{#2}}
\newcommand{\AAA}{\mathbf{A}}
\newcommand{\Sf}{\mathfrak{S}}
\DeclareMathOperator*{\quotprojlim}{\text{``$\varprojlim$''}}
\DeclareMathOperator*{\quotinjlim}{\text{``$\varinjlim$''}}
\DeclareMathOperator{\Fl}{F\sisi{\ell}{l}}
\DeclareMathOperator{\Pl}{P\sisi{\ell}{l}}
%\DeclareMathOperator{\cl}{c\sisi{\ell}{l}}
\DeclareMathOperator{\tdf}{tdf}
\DeclareMathOperator{\rg}{rg}
\newcommand{\iInd}{\sisi{I}{i}nd}
\newcommand{\iind}{ind}
\newcommand{\ppro}{pro}
\newcommand{\tTor}{\sisi{T}{t}or}
\newcommand{\ttor}{tor}
\newcommand{\chap}{Chap.}
\newcommand{\Ffex}{\Ff\mathit{ex}}
\newcommand{\sheafk}{\underline{k}}
\newcommand{\sheafn}{\underline{n}}
\newcommand{\sheafGamma}{\underline{\Gamma}}

% Repeated references to an already typeset footnote must not use
% \footnotemark without a matching \footnotetext: Hyperref would then create
% a live annotation to a destination that cannot exist.  This helper keeps the
% current printed footnote mark and links it to an explicit stable note target.
\makeatletter
\newcommand{\sgarepeatedfootnotemark}[1]{%
  \begingroup
  \edef\@thefnmark{\thefootnote}%
  \hyperlink{#1}{\@makefnmark}%
  \endgroup}
\makeatother

\sisi{
\DeclareMathOperator{\simarrow}{\rule{.15mm}{1.8mm}\kern -3.4pt\rightsquigarrow}
\newcommand{\sheafell}{\underline{\ell}}
\newcommand{\sheafZZ}{\underline{\ZZ}}
}{
\newcommand{\nderef}[1]{la N.D.E.~(\ref{#1}) page~\pageref{#1}}
}

%%%%%%%%%%%%%%%%%
%%%%% expose XVII
%%%%%%%%%%%%%%%%%

% END INCLUDED SGA4_English_current_buildable_source_20260906/source/fichiers-configuration-latex/macros.tex


\externaldocument{../01/01_English.tex}
\externaldocument{../02/02_English.tex}
\externaldocument{../03/03_English.tex}
\externaldocument{../04/04_English.tex}
\externaldocument{../05/05_English.tex}
\externaldocument{../05bis/05bis_English.tex}
\externaldocument{../06/06_English.tex}
\externaldocument{../07/07_English.tex}
\externaldocument{../08/08_English.tex}
\externaldocument{../09/09_English.tex}
\externaldocument{../10/10_English.tex}
\externaldocument{../11/11_English.tex}
\externaldocument{../12/12_English.tex}
\externaldocument{../13/13_English.tex}
\externaldocument{../14/14_English.tex}

\setboolean{original}{false}

\begin{document}

\begin{center}
XV. Acyclic morphisms\\
M. Artin
\end{center}

\selectlanguage{english}version: \showgitstatus
\tableofcontents

\else
\chapter{Acyclic morphisms}
%\addtocontents{toc}{M. Artin}
\begin{center}
M. Artin
\end{center}
\fi
\phantomsection\MakeLinkTarget*{sga4.exp.XV}\label{sga4.exp.XV}

% BEGIN INCLUDED SGA4_English_current_buildable_source_20260906/source/15/drafts/SGA4_XV_LINES_43_400_English.texfrag
\setcounter{pageoriginal}{167}

% To do:
% standardize notation for restrictions, e.g. F|X -> F_{|X}.
% add ~ to avoid short mathematical terms standing alone at the start of a line.
% replace \ref{...} by (\ref{}); cf. for example the statement of 1.15.
% do not set remarks in italics (cf. Remark 1.18).

\section*{Introduction}\pageoriginale

Let $g:X\rightarrow Y$ be a morphism of schemes. In the first section
we study conditions under which $g$ is acyclic, that is, such that for
every torsion sheaf $F$ on $Y$ one has
$H^{q}(Y,F)\simeq H^{q}(X,g^{\ast}F)$, and such that this remains true
when $Y$ is replaced by a $Y'$ étale over $Y$. A natural necessary
condition for acyclicity is that the geometric fibers be acyclic, that
is, have trivial cohomology for constant torsion sheaves. We shall see
that this condition, together with a local condition called local
acyclicity of $g$ (\hyperref[XV.def1.11]{1.11}), implies that $g$ is
acyclic.

The fundamental theorem is that a smooth morphism is locally acyclic
for torsion sheaves of order prime to the residue characteristics
\hyperref[XV.thm2.1]{2.1}. This will imply the base-change theorem for
a smooth morphism \hyperrefext[XVI][XVI.thm1.1]{1.1}, which will be
developed, together with its first consequences, in the next exposé.

The present exposé is independent of Exposés XI to XIV, and in
particular of the ``base-change theorem for a proper morphism.'' The
latter will again enter in an essential way, but in conjunction with
the fundamental theorem of the present exposé, beginning with the next
exposé.

\section{Generalities on globally and locally acyclic
  morphisms}\etiquette[1]{XV.sec1}\pageoriginale

% g' is not defined as g x_Y Y' here, whereas it is in 1.3.

\begin{enonce*}{Proposition 1.1}\etiquette[1.1]{XV.prop1.1}
Let $g:X\rightarrow Y$ be a morphism of schemes. The following
assertions {\rm(i)} to {\rm(iii)} are equivalent:
\begin{itemize}
\item[(i)] For every étale morphism $Y'\rightarrow Y$ (where $Y'$ may
  be taken affine over an affine open subset of $Y$) and every sheaf of
  sets $F$ on $Y'$, the canonical morphism $F\rightarrow
  g'_{\ast}{g'}^{\ast}F$ is injective (\resp bijective).

\item[(ii)] For every étale morphism $Y'\rightarrow Y$ (where $Y'$ may
  be taken affine over an affine open subset of $Y$) and every sheaf of
  sets $F$ on $Y'$, the canonical morphism $H^{0}(Y',F)\rightarrow
  H^{0}(X',{g'}^{\ast}F)$ is injective (\resp bijective).

\item[(iii)] For every locally quasi-finite morphism $Y'\rightarrow Y$
  and every sheaf of sets $F$ on $Y'$, the canonical morphism
  $H^{0}(Y',F)\rightarrow H^{0}(X',{g'}^{\ast}F)$ is injective
  (\resp bijective).

\item[(ii bis)] (If $g$ is quasi-compact and quasi-separated.) As in
  {\rm(ii)}, but with $F$ a constant sheaf of the form $I_{Y'}$, where
  $I$ is a prescribed set with $\card I\geq 2$.
\end{itemize}
\end{enonce*}

\noindent
\emph{Proof}. The equivalence of (i) and (ii) follows immediately from
the definitions, while (iii) $\Rightarrow$ (ii) $\Rightarrow$ (ii bis)
are trivial. We have (ii bis) $\Rightarrow$ (ii), since we may suppose
$Y$ affine, so that $X$ and $Y$ are quasi-compact and quasi-separated,
and may then apply \hyperrefext[XII][XII.prop6.5]{6.5} (i). For the
implication (i) $\Rightarrow$ (iii), we first record the following
lemma:

\begin{enonce*}{Lemma 1.2}\etiquette[1.2]{XV.lem1.2}
Let\pageoriginale $f:Y'\rightarrow Y$ be a locally quasi-finite
morphism. Then $f$ is ``locally finite on $Y'$ for the étale
topology,'' that is, there are commutative diagrams
$$\vcenter{
\xymatrix@=1.75cm{
Y'_i        \ar[r]^{a_i}   \ar[d]^{b_i}& Y'\ar[d]^{f}\\
Y_i          \ar[r]^{c_i}              & Y
}}\leqno{(*)}
$$
in which the $a_i$ and $c_i$ are étale, the $b_i$ are finite, and the
$a_i$ form a covering of $Y'$ \footnote[1]{\Cf \EGA IV 18.12.1.}.
\end{enonce*}

Indeed, let $y'$ be a point of $Y'$ and let $y=f(y')$. It is enough to
find a diagram $({\ast})$ such that $y'$ lies in the image of $Y'_i$.
Let $\widetilde{Y}$ be the spectrum of the henselization of the ring
$\Oo_{Y,y}$ \hyperrefext[VIII][VIII.subsec4.1]{4.1}, and let $Z$ be
the (ordinary) localization of
$\widetilde{Y}'=\fprod{Y'}{\widetilde{Y}}{Y}$ at the unique point
lying over $y'$ and the closed point $\widetilde{y}$ of
$\widetilde{Y}$. By \hyperrefext[VIII][VIII.subsec4.1]{4.1} (ii),
% label for (ii)?
$Z\rightarrow \widetilde{Y}$ is finite, and $Z$ is an induced open
subscheme of $\widetilde{Y}'$. Since $\widetilde{Y}$ is a limit of
schemes $Y_i$ étale over $Y$, one sees immediately from \EGA IV 9 that,
for a suitable $Y_i$, there is an open subscheme $Z_i$ of
$\fprod{Y'}{Y_i}{Y}$ such that
$Z=\fprod{Z_i}{\widetilde{Y}}{Y_i}$; in $({\ast})$ it is then enough to
take $Y'_i=Z_i$. \cqfd


Now suppose that (i) of \hyperref[XV.prop1.1]{1.1} holds. It is clear
that assertion (iii) of \hyperref[XV.prop1.1]{1.1} is equivalent to the
assertion obtained from (i) by replacing the étale morphism
$Y'\rightarrow Y$ by a locally quasi-finite morphism. In this form,
the assertion (for given $Y$ and $Y'$) is local on $Y$ and $Y'$ for
the étale topology; hence the lemma reduces us to the case of a
\emph{finite} morphism $f:Y'\rightarrow Y$.\pageoriginale Let
$f':X'\rightarrow X$ be the morphism obtained from $f$ by base change
along $g$. Then the base-change morphism
$g^{\ast}f_{\ast}F\rightarrow {f'}_{\ast}{g'}^{\ast}F$ is bijective
\hyperrefext[VIII][VIII.cor5.6]{5.6}. By (i),
\sisi{the adjunction morphism
$f_{\ast}F\rightarrow g_{\ast}g^{\ast}f_{\ast}F$ is injective
(resp.~bijective)}{the adjunction morphism
$f_{\ast}F\rightarrow g_{\ast}g^{\ast}f_{\ast}F$ is injective
(resp.~bijective)}.
Consequently, the morphism $H^{0}(Y',F)\xrightarrow{\epsilon}
H^{0}(Y',g'_{\ast}{g'}^{\ast}F)$ in (iii) is the composite
$H^{0}(Y',F)=H^{0}(Y,f_{\ast}F)\xrightarrow{\epsilon'}
H^{0}(Y,g_{\ast}g^{\ast}f_{\ast}F)=H^{0}(X,g^{\ast}f_{\ast}F)\simeq
H^{0}(X,{f'}_{\ast}{g'}^{\ast}F)=H^{0}(X',{g'}^{\ast}F)$.
Since $\epsilon'$ is injective (\resp bijective), the same is true of
$\epsilon$. This completes the proof of \hyperref[XV.prop1.1]{1.1}.

\begin{enonce*}{Definition 1.3}\etiquette[1.3]{XV.def1.3}
A morphism $g:X\rightarrow Y$ having one of the equivalent properties
{\rm(i)} to {\rm(iii)} (\resp the corresponding parenthesized
properties) of \hyperref[XV.prop1.1]{1.1} is called a
\emph{$(-1)$-acyclic morphism} (\resp a \emph{$0$-acyclic morphism}).
The morphism is said to be \emph{universally} $(-1)$-acyclic
(\resp $0$-acyclic) if, for every base-change morphism
$Y'\rightarrow Y$, the morphism
$g'=\fprod{g}{Y'}{Y}:\fprod{X}{Y'}{Y}\rightarrow Y'$ is
$(-1)$-acyclic (\resp $0$-acyclic). Finally, a scheme $X$ is said to
be $(-1)$-acyclic (\resp $0$-acyclic) if it is nonempty (\resp
connected and nonempty).
\end{enonce*}

\begin{enonce*}{Corollary 1.4}\etiquette[1.4]{XV.cor1.4}
Let $g:X\rightarrow Y$ be a $0$-acyclic morphism. Then for every
locally quasi-finite morphism $Y'\rightarrow Y$ and every sheaf of
groups $F$ on $Y'$, the canonical morphism
$$
H^{1}(Y',F)\longrightarrow H^{1}(X',{g'}^{\ast}F)
$$
is injective.
\end{enonce*}

This follows immediately from \hyperrefext[XII][XII.prop6.5]{6.5} (i).

\begin{enonce*}{Corollary 1.5}\etiquette[1.5]{XV.cor1.5}
Let\pageoriginale $g:X\rightarrow Y$ be a morphism. If $g$ is
surjective (\ie its fibers are nonempty, \ie its fibers are
$(-1)$-acyclic), then $g$ is $(-1)$-acyclic. The converse holds if $g$
is quasi-compact and quasi-separated.
\end{enonce*}

The first assertion is clear in the form \hyperref[XV.prop1.1]{1.1}
(ii), by considering the fibers of the sheaves in question. The second
follows from the definition in the form \hyperref[XV.prop1.1]{1.1}
(iii) and the limit argument of \hyperrefext[VII][VII.sec5]{5}, taking
into account that, for every $y\in Y$, $y=\Spec k(y)$ is the inverse
limit of affine schemes locally quasi-finite over $Y$, namely the
affine neighborhoods of $y$ in $\bar{\{y\}}$ (endowed with the induced
reduced structure). \sisi{}{--} Note that \hyperref[XV.cor1.5]{1.5}
implies that, when $g$ is quasi-compact and quasi-separated, the
$(-1)$-acyclicity of $g$ already implies that $g$ is
\emph{universally} $(-1)$-acyclic.

\begin{enonce*}{Proposition 1.6}\etiquette[1.6]{XV.prop1.6}
Let $g:X\rightarrow Y$ be a morphism of schemes, let
$\LL\subset\PP$ be a set of prime numbers, and let $n\in\NN$ be a
natural number (\resp let $n=1$). The following assertions {\rm(i)} to
{\rm(iii)} (\resp the corresponding parenthesized assertions) are
equivalent. If, moreover, $g$ is quasi-compact and quasi-separated,
then {\rm(i)} to {\rm(iii)} (\resp the corresponding parenthesized
assertions) are also equivalent to {\rm(iv)} (\resp the corresponding
parenthesized assertion).
\begin{enumerate}
\renewcommand{\theenumi}{\roman{enumi}}
\renewcommand{\labelenumi}{(\theenumi)}
\item For every étale morphism $Y'\rightarrow Y$ and every abelian
  $\LL$-torsion sheaf (\resp sheaf of
  $\text{\rm ind}$-$\LL$-groups) \emph{\hyperrefext[IX][IX.def1.5]{1.5}}
  $F$ on $Y'$, one has
$$
\begin{cases}
F\xrightarrow{\sim} {g'}_{\ast}{g'}^{\ast}F, &\\
(R^{q}{g'}_{\ast}){g'}^{\ast}F=0 & \mbox{if}\quad 1\leq q\leq n\quad
(\resp \mbox{if\; } q=1){\quoi}.
\end{cases}
$$

\item For\pageoriginale every étale morphism $Y'\rightarrow Y$ and
every abelian $\LL$-torsion sheaf (\resp sheaf of
ind-$\LL$-groups) $F$ on $Y'$, the canonical morphism
$$
H^{q}(Y',F)\longrightarrow H^{q}(X',{g'}^{\ast}F)
$$
is bijective if $q\leq n$ and injective if $q=n+1$ (\resp is
bijective if $q\leq1$).

\item For every locally quasi-finite morphism $Y'\rightarrow Y$ and
every abelian $\LL$-torsion sheaf (\resp sheaf of
ind-$\LL$-groups) $F$ on $Y'$, the morphism
$$
H^{q}(Y',F)\longrightarrow H^{q}(X', {g'}^{\ast}F)
$$
is bijective if $q\leq n$ and injective if $q=n+1$ (\resp is
bijective if $q\leq1$).

\item The morphism $g$ is surjective, and for every locally
quasi-finite morphism $Y'\rightarrow Y$, every $\nu>0$, and every
$\ell\in\LL$, the canonical morphism
$$
H^{q}(Y',\ZZ/\ell^{\nu})\longrightarrow H^{q}(X',\ZZ/\ell^{\nu})
$$
is \emph{surjective} for $0\leq q\leq n$ (\resp for every locally
quasi-finite morphism $Y'\rightarrow Y$ and every ordinary finite
$\LL$-group $G$, the canonical morphism
$$
H^{i}(Y',G)\longrightarrow H^{i}(X',G)
$$
is \emph{surjective} for $i=0,1$).
\end{enumerate}
\end{enonce*}

\noindent
\emph{Proof of}\pageoriginale \hyperref[XV.prop1.6]{1.6}. The
implications (ii) $\Leftarrow$ (iii) $\Rightarrow$ (iv) are trivial.
% trivial surjectivity? (g is not necessarily coherent)
Now $(R^{q}{g'}_{\ast}){g'}^{\ast}F$ is the sheaf associated with the
presheaf $R^{q}(Y'')=H^{q}(X'',{g''}^{\ast}F)$
($Y''\rightarrow Y'$ étale). If (ii) holds, we have
$H^{q}(Y'',F)\xrightarrow{\sim}H^{q}(X'',{g''}^{\ast}F)$, so the
associated sheaf is zero\sisi{ if $1\leq q\leq n$}{}, \ie
(ii) $\Rightarrow$ (i). The implication (i) $\Rightarrow$ (ii) follows
from \hyperref[XV.prop1.1]{1.1} (i) (which corresponds to the case
$n=0$), together with the Leray spectral sequence
$$
H^{p}(Y',(R^{q}{g'}_{\ast}){g'}^{\ast}F)\Longrightarrow
H^{p+q}(X',{g'}^{\ast}F)
$$
in the abelian case, and the exact sequence
$$
0\longrightarrow H^{1}(Y',{g'}_{\ast}{g'}^{\ast}F)\longrightarrow
H^{1}(X',{g'}^{\ast}F)\longrightarrow
H^{0}(Y',(R^{1}{g'}_{\ast}){g'}^{\ast}F){\quoi}
$$
in the case of sheaves of groups.

(i) $\Rightarrow$ (iii): Since (i) and (ii) are equivalent, it is
clear that (iii) is equivalent to the statement obtained from (i) by
replacing the word ``étale'' by the words ``locally quasi-finite.'' In
this form, for given $Y$ and $Y'$, it is an assertion local on $Y'$ and
on $Y$ for the étale topology, so \hyperref[XV.lem1.2]{1.2} reduces us
to the case in which $f:Y'\rightarrow Y$ is a \emph{finite} morphism.
Let $f':X'\rightarrow X$ be the morphism
\sisi{obtained from~$f$ by base change along~$g$}{obtained from~$f$ by base change along~$g$}.
By \hyperrefext[VIII][VIII.prop5.5]{5.5}
(\resp \hyperrefext[VIII][VIII.rem5.8]{5.8}) (applied three times), we
have
$$
f_{\ast}(R^{q}{g'}_{\ast}){g'}^{\ast}F\simeq
R^{q}(fg')_{\ast}{g'}^{\ast}F\simeq
R^{q}(gf')_{\ast}{g'}^{\ast}F\simeq
(R^{q}g_{\ast}){f'}_{\ast}{g'}^{\ast}F\simeq
(R^{q}g_{\ast})g^{\ast}(f_{\ast}F)
$$
and the last of these is zero, for the values of $q$ in question, if
(i) holds. Thus (i) $\Rightarrow$ (iii).

It remains to prove (iv) $\Rightarrow$ (iii), and this follows
immediately\pageoriginale from \hyperrefext[XII][XII.prop6.5]{6.5}.

\begin{enonce*}{Definition 1.7}\etiquette[1.7]{XV.def1.7}
Let $n\in\NN$ and $\LL\subset\PP$. A morphism
$g:X\rightarrow Y$ satisfying one of the equivalent conditions
{\rm(i)} to {\rm(iii)} (\resp the corresponding parenthesized
conditions) of \hyperref[XV.prop1.6]{1.6} is called an
\emph{$n$-acyclic morphism} (\resp a \emph{$1$-aspherical morphism})
for $\LL$. We say that $g$ is \emph{acyclic} for $\LL$ if it is
$n$-acyclic for $\LL$ for every $n$. We say that $g$ is
\emph{universally $n$-acyclic} (\resp \emph{universally
$1$-aspherical}) for $\LL$ if, for every $Y'\rightarrow Y$, the
morphism $g'=\fprod{g}{Y'}{Y}$ is $n$-acyclic (\resp $1$-aspherical)
for $\LL$. Finally, a scheme $X$ is said to be $n$-\emph{acyclic}
(\resp $1$-\emph{aspherical}) for $\LL$ if it is $0$-acyclic in the
sense of \hyperref[XV.def1.3]{1.3} and if, for every $\ell\in\LL$ and
every $1\leq q\leq n$, one has $H^{q}(X,\ZZ/\ell)=0$ (\resp and if,
for every ordinary finite $\LL$-group $G$, one has $H^{1}(X,G)=0$).
\end{enonce*}

\begin{enonce*}[remark]{Remarks 1.8}\etiquette[1.8]{XV.rem1.8}
\begin{enumerate}
\renewcommand{\theenumi}{\alph{enumi}}
\renewcommand{\labelenumi}{\theenumi)}
\item If $g$ is quasi-compact and quasi-separated, then, in order for
  $g$ to be universally $n$-acyclic, it is enough that $g'$ be
  $n$-acyclic whenever $Y'\rightarrow Y$ is locally of finite
  presentation. We do not need this fact and leave its proof, which is
  an exercise in passage to the limit, to the reader.

\item We do not know whether a morphism that is $n$-acyclic for
  $\LL$ is necessarily universally $n$-acyclic for $\LL$. The same
  question arises for the local variants below
  \hyperref[XV.def1.11]{1.11}.

\item When $n=0$, we recover the notion of
  \hyperref[XV.def1.3]{1.3}. Of course,
  \hyperref[XV.prop1.6]{1.6} and \hyperref[XV.def1.7]{1.7} could have
  been formulated so as to include the case $n=-1$. The following
  proposition is also valid in that case.\pageoriginale
\end{enumerate}
\end{enonce*}

\begin{enonce*}{Proposition 1.9}\etiquette[1.9]{XV.prop1.9}
\begin{enumerate}
\renewcommand{\theenumi}{\roman{enumi}}
\renewcommand{\labelenumi}{(\theenumi)}
\item Let $X\xrightarrow{g}Y\xrightarrow{h}Z$ be morphisms of schemes,
  let $\LL\subset\PP$, and let $n\in\NN$. Suppose that $g$ is
  $0$-acyclic (\resp $1$-aspherical for $\LL$, \resp $n$-acyclic for
  $\LL$). Put $f=hg$. Then, for every sheaf of sets (\resp of
  ind\sisi{-}{ }$\LL$-groups, \resp abelian $\LL$-torsion sheaf) $F$
  on $Z$, there is an isomorphism
  $f_{\ast}f^{\ast}F\xrightarrow{\sim}h_{\ast}h^{\ast}F$ (\resp
  isomorphisms $(R^{p}f_{\ast})f^{\ast}F\xrightarrow{\sim}
  (R^{p}h_{\ast})h^{\ast}F$ for $p\leq1$, \resp for $p\leq n$). In
  particular, $h$ is $0$-acyclic (\resp $1$-aspherical for $\LL$,
  \resp $n$-acyclic for $\LL$) if and only if $f$ is.

\item Let $\{g_{\alpha}:X_{\alpha}\rightarrow Y_{\alpha}\}$ be an
  inverse system of morphisms of schemes such that the morphisms
  $X_{\alpha}\rightarrow X_{\beta}$ and
  $Y_{\alpha}\rightarrow Y_{\beta}$ are affine and the $g_{\alpha}$
  are quasi-compact and quasi-separated. Let $g:X\rightarrow Y$ be the
  inverse limit of the $g_{\alpha}$
  \hyperrefext[VII][VII.subsec5.1]{5.1}. Then, if all the $g_{\alpha}$
  are $0$-acyclic (\resp $1$-aspherical for $\LL$, \resp $n$-acyclic
  for $\LL$), the same is true of $g$.
\end{enumerate}
\end{enonce*}

\noindent
\emph{Proof}. For (i), first observe that
$$
h_{\ast}h^{\ast}F\xrightarrow{\sim}
h_{\ast}(g_{\ast}g^{\ast})h^{\ast}F\simeq f_{\ast}f^{\ast}F{\quoi}.
$$
The second assertion of (i) is an immediate consequence of the Leray
spectral sequence
$$
(R^{p}h_{\ast})(R^{q}g_{\ast})f^{\ast}F\Longrightarrow
(R^{p+q}f_{\ast})f^{\ast}F{\quoi}.
$$
Finally, the nonabelian assertion follows from the exact sequence
\hyperrefext[XII][XII.prop3.2]{3.2}.

% END INCLUDED SGA4_English_current_buildable_source_20260906/source/15/drafts/SGA4_XV_LINES_43_400_English.texfrag


% BEGIN INCLUDED SGA4_English_current_buildable_source_20260906/source/15/drafts/SGA4_XV_LINES_401_700_English.texfrag
Assertion\pageoriginale (ii) follows readily from the theory of passage
to the limit \hyperrefext[VII][VII.sec5]{5} and VI. % label for chapter?

\begin{enonce*}{Corollary 1.9.1}\etiquette[1.9.1]{XV.cor1.9.1}
Let $g:X\rightarrow Y$ be a quasi-compact and quasi-separated morphism,
and let $n\in \NN$, $\LL\subset \PP$. If $g$ is
$n$-acyclic for $\LL$ (\resp $1$-aspherical for
$\LL$), then, for every geometric point $\bar{y}$ of $Y$
algebraic over a point $y$ of $Y$, the geometric fiber $X_{\bar{y}}$
is $n$-acyclic for $\LL$ (\resp $1$-aspherical for $\LL$).
\end{enonce*}

Indeed, $\bar{y}$ is the inverse limit of affine schemes $Y_i$ that
are locally quasi-finite over $Y$, and since the
$X_i=\fprod{X}{Y_i}{Y}\rightarrow Y_i$ are $n$-acyclic for
$\LL$ (\resp\ldots), the same holds for
$X_{\bar{y}}\rightarrow \bar{y}$ by \hyperref[XV.prop1.9]{1.9} (ii), which gives the
asserted conclusion.

\begin{enonce*}{Proposition 1.10}\etiquette[1.10]{XV.prop1.10}
Let $g:X\rightarrow Y$ be a morphism of schemes, let $\LL\subset
\PP$, and let $n\in \NN$. Then conditions {\rm(i)} through
{\rm(iv)} below are equivalent.
\begin{enumerate}
\renewcommand{\theenumi}{\roman{enumi}}
\renewcommand{\labelenumi}{(\theenumi)}
\item Let $\bar{y}$ be a geometric point of $Y$ and $\bar{x}$ a geometric
point of $X$ lying over $\bar{y}$. Let $\widetilde{X}$ and
$\widetilde{Y}$ be the strict localizations of $X$ and $Y$ at
$\bar{x}$ and $\bar{y}$, respectively, and let
$\widetilde{g}:\widetilde{X}\rightarrow
\widetilde{Y}$ be the morphism induced by $g$. Then $\widetilde{g}$ is
$0$-acyclic (\resp $1$-aspherical for $\LL$, \resp
$n$-acyclic for $\LL$).

\item For every diagram with Cartesian squares
$$
\vcenter{
\xymatrix@C=2cm@R=1.25cm{
X   \ar[d]^{g}& X'\ar[l]_{j}  \ar[d]^{g'} & X'' \ar[l]_{j'} \ar[d]^{g''}\\
Y    & Y'\ar[l]_{i} &   Y'' \ar[l]_{i'} 
}}\leqno{(\ast)}
$$
in which\pageoriginale $i$ is etale and $i'$ is etale and of finite
presentation, and for every sheaf of sets (\resp of ind-$\LL$-groups,
\resp an abelian $\LL$-torsion sheaf) $F$ on $Y''$, the base-change
morphism \hyperrefext[XII][XII.eq4.1]{4.1}, \hyperrefext[XII][XII.eq4.2]{4.2}
$$
{g'}^{\ast}i'_{\ast}F\longrightarrow {j'}_{\ast}{g''}^{\ast}F
$$
is bijective (\resp the morphisms
$$
{g'}^{\ast}(R^{q}i'_{\ast})F\longrightarrow
(R^{q}j'_{\ast}){g''}^{\ast}F
$$
are bijective for $q=0$, $1$, \resp are bijective for $0\leq q\leq
n$ and injective for $q=n+1$).

\item The same statement as {\rm(ii)}, except that the morphism $i'$ is
assumed only quasi-finite and quasi-separated.

\item The same statement as {\rm(ii)}, but with $i$ assumed only
locally quasi-finite and, on the other hand, $i'$ a quasi-compact open
immersion.
\end{enumerate}
\end{enonce*}

\noindent
\emph{Proof}. The implication (iii) $\Rightarrow$ (ii) is
trivial. Let us verify that (i) $\Rightarrow$ (iii). It is enough to check
the morphism in (iii) fiber by fiber. Let $\bar{y}$ be a geometric point
of $Y'$ and $\bar{x}$ a geometric point of $X'$ lying over
$\bar{y}$. Let
$$
\xymatrix@=1.75cm{
\widetilde{X}'   \ar[d]^{\widetilde{g}'}& \widetilde{X}''
\ar[l]_{\widetilde{j}'} \ar[d]^{\widetilde{g}''}\\
 \widetilde{Y}'           & \widetilde{Y}'' \ar[l]_{\widetilde{i}'}
}
$$
be the Cartesian diagram deduced from $({\ast})$, where
$\widetilde{X}'$ and $\widetilde{Y}'$ are the strict localizations of
$X'$ and $Y'$ at $\bar{x}$ and $\bar{y}$, respectively, and where
$\widetilde{Y}''=\fprod{Y''}{\widetilde{Y}'}{Y'}$,
$\widetilde{X}''=\fprod{X''}{\widetilde{X}'}{X'}$.\pageoriginale By
\hyperrefext[VIII][VIII.thm5.2]{5.2} and \hyperrefext[VIII][VIII.rem5.3]{5.3}, we have
$$
H^{q}(\widetilde{Y}'',\widetilde{F})\simeq
(R^{q}i'{}_{\ast}F)_{\bar{y}}\simeq
({g'}^{\ast}(R^{q}i'_{\ast})F)_{\bar{x}}{\quoi},
$$
where $\widetilde{F}$ is the inverse image of $F$ on
$\widetilde{Y}''$, and
$$
H^{q}(\widetilde{X}'',{\widetilde{g}''}{}^{\ast}\widetilde{F})\simeq
((R^{q}j'_{\ast}){g''}{}^{\ast}F)_{\bar{x}}{\quoi}.
$$
Since $\widetilde{i}'$ is quasi-finite, we may apply
\hyperref[XV.prop1.6]{1.6} (iii) (\resp \hyperref[XV.cor1.5]{1.5} (iii)) to the morphism
$$
H^{q}(\widetilde{Y}'',\widetilde{F})\longrightarrow
H^{q}(\widetilde{X}'',{\widetilde{g}''}{}^{\ast}\widetilde{F})
$$
for the values of $q$ under consideration, which proves (i)
$\Rightarrow$ (iii).

(ii) $\Rightarrow$ (i). Assume (ii) and let us verify the acyclicity
criterion \hyperref[XV.prop1.6]{1.6} (ii) for a morphism
$\widetilde{g}:\widetilde{X}\rightarrow \widetilde{Y}$. Let
$\widetilde{Y}'\rightarrow \widetilde{Y}$ be etale, with
$\widetilde{Y}'$ affine, and let $\widetilde{F}$ be a sheaf of sets
(\resp an abelian $\LL$-torsion sheaf, \resp a sheaf of
ind-$\LL$-groups) on $\widetilde{Y}'$. We must prove that
$$
H^{q}(\widetilde{Y}',\widetilde{F})\xrightarrow{\sim}
H^{q}(\widetilde{X}', \widetilde{g}'{}^{\ast}\widetilde{F}){\quoi},
$$
where $\widetilde{g}':\widetilde{X}'\rightarrow \widetilde{Y}'$ is the
morphism obtained from $\widetilde{g}$ by base change. Now write
$\widetilde{Y}=\varprojlim Y_{\alpha}$, where the $Y_{\alpha}$ are
affine and etale over $Y$ (assuming $Y$ affine, as we may). We can
``descend'' the etale morphism
$\widetilde{Y}'\rightarrow \widetilde{Y}$ to one of the $Y_{\alpha}$,
say to $Y'_{0}\rightarrow Y_{0}$, and we have
$\widetilde{Y}'=\varprojlim
Y'_{\alpha}=\fprod{Y}{Y'_{0}}{Y_{0}}$. Every sheaf
$\widetilde{F}$ on $\widetilde{Y}'$ is the direct limit of the sheaves
${f'_{\alpha}}^{\ast}{f'_{\alpha}}_{\ast}\widetilde{F}$, where
$f'_{\alpha}:\widetilde{Y}'\rightarrow Y'_{\alpha}$ is the canonical
morphism. By \sisi{XII 1.2 (v) and 1.6 (iii)}{\hyperrefext[IX][IX.prop1.2]{1.2} and \hyperrefext[IX][IX.prop1.6]{1.6}}, the sheaves
$f'_{\alpha}{}_{\ast}\widetilde{F}$ are sheaves of sets (\resp abelian
$\LL$-torsion sheaves, \resp sheaves of
ind-$\LL$-groups). In view\pageoriginale of \hyperrefext[VII][VII.prop3.3]{3.3},
it is therefore enough to treat the case where $\widetilde{F}$ descends
to one of the $Y'_{\alpha}$, say to a sheaf $F_{0}$ on $Y'_{0}$. Applying
(ii) to the diagram
$$
\xymatrix@C=2.25cm@R=1.5cm{
X \ar[d]_{g}& X_0 \ar[l] \ar[d]_{g_0}& X'_0 \ar[l]_{j'_0} \ar[d]_{g'_0}\\
Y & Y_0 \ar[l]& Y'_0 \ar[l]_{i'_0},
}
$$
we deduce that
${g_{0}}^{\ast}(R^{q}{i'_0}_{\ast})F_{0}\xrightarrow{\sim}
(R^{q}{j'_0}_{\ast}){g'_0}^{\ast}F_{0}$ for the applicable
values of $q$. The geometric points $\bar{y}$ and $\bar{x}$ lift
canonically to geometric points of $Y_{0}$ and $X_{0}$ (which we denote
by the same symbols), and by \hyperrefext[VIII][VIII.thm5.2]{5.2} and
\hyperrefext[VIII][VIII.rem5.3]{5.3} we have
$$
H^{q}(\widetilde{Y}',\widetilde{F})=((R^{q}{i'_0}_{\ast})F_{0})_{\bar{y}}=({g_0}^{\ast}
(R^{q}{i'_0}_{\ast})F_0)_{\bar{x}}
$$
and
$$
H^{q}(\widetilde{X}',{\widetilde{g}'}{}^{\ast}\widetilde{F})=((R^{q}j'_{0\ast}){g'_0}^{\ast}F_{0})_{\bar{x}}
$$
for the values of $q$ under consideration, which proves the result.

It remains to show that conditions (i) through (iii) are equivalent to
(iv). Using form (i), we see that these conditions are stable under a
locally quasi-finite base change $Y'\rightarrow Y$; in form (ii), this
immediately shows that they imply (iv). Conversely, (iv) implies (iii),
since, in order to verify (iii), it follows at once---using, for example,
the Leray spectral sequence or the descent result
\hyperrefext[XII][XII.prop6.8]{6.8}---that we may assume $Y'$ and $Y''$ affine. But then, by
the ``Main Theorem'' \EGA IV 8.12.8, $Y''\rightarrow Y'$ factors as
$Y''\xrightarrow{i'_1}Y'_1\xrightarrow{i'_2}Y'$, with $i'_1$ an open
immersion\pageoriginale and $i'_2$ a \emph{finite} morphism. Using
\hyperrefext[VIII][VIII.prop5.5]{5.5}, \hyperrefext[VIII][VIII.rem5.8]{5.8} for $i'_2$, we are reduced to
verifying (ii) for $i'_1$, which follows from (iv). This completes the
proof of \hyperref[XV.prop1.10]{1.10}.

\begin{enonce*}{Definition 1.11}\etiquette[1.11]{XV.def1.11}
A morphism $g:X\rightarrow Y$ of preschemes satisfying any one of the
equivalent conditions {\rm(i)} through {\rm(iv)} of $\hyperref[XV.prop1.10]{1.10}$ is called a
\emph{locally $0$-acyclic} morphism (\resp \emph{locally
$1$-aspherical} for $\LL$, \resp \emph{locally
$n$-acyclic} for $\LL$). In the obvious way (compare
\hyperref[XV.def1.7]{1.7}), one defines the notions of a \emph{locally acyclic} morphism
and of a \emph{universally locally $0$-acyclic} morphism
(\resp\ldots) for $\LL$.
\end{enonce*}

The statements (i) through (iii) of \hyperref[XV.prop1.10]{1.10} can of course be varied.
Thus, in (iii) and (iv), it is enough to restrict to a constant sheaf $F$
of the form $G_{I}$,
% what is this notation?
with $G$ a finite set (\resp a finite $\LL$-group, \resp $G$ of the
form $\ZZ/\ell^{\nu}\ZZ$, with $\ell\in\LL$, $\nu\in \NN$); this
follows readily from \hyperrefext[VIII][VIII.thm5.2]{5.2},
\hyperrefext[VIII][VIII.rem5.3]{5.3}, and \hyperrefext[XII][XII.prop6.5]{6.5}. We also record the following
variant:

\begin{enonce*}{Corollary 1.12}\etiquette[1.12]{XV.cor1.12}
Suppose that $g:X\rightarrow Y$ is locally of finite type. To verify
local $0$-acyclicity (\resp\ldots), it is enough to verify $\hyperref[XV.prop1.10]{1.10}$
{\rm(i)} in the case where the geometric point $\bar{x}$ is
\emph{closed} in the geometric fiber $X_{\bar{y}}$.
\end{enonce*}

Indeed, in the proof of (i) $\Rightarrow$ (iii) of \hyperref[XV.prop1.10]{1.10}, we
verified (iii) fiber by fiber, and it was enough to do so for the fibers
at geometric points $\bar{x}$ that are closed in their geometric
fiber,\pageoriginale since $g$ is locally of finite type
\hyperrefext[VIII][VIII.rem3.13]{3.13} b).

\begin{enonce*}{Proposition 1.13}\etiquette[1.13]{XV.prop1.13}
\begin{enumerate}
\renewcommand{\theenumi}{\roman{enumi}}
\renewcommand{\labelenumi}{(\theenumi)}
\item Let $X\xrightarrow{g}Y\xrightarrow{h}Z$ be morphisms of schemes,
let $\LL\subset \PP$, and let $n\in \NN$. Suppose that $g$ is
surjective and locally $0$-acyclic (\resp locally $1$-aspherical for
$\LL$, \resp locally $n$-acyclic for $\LL$). Then $h$ is
locally $0$-acyclic (\resp\ldots) if and only if $f=hg$ is.

\item Let $\{g_{\alpha}:X_{\alpha}\rightarrow Y_{\alpha}\}$ be an
inverse system of morphisms of preschemes such that the transition
morphisms $X_{\alpha}\rightarrow X_{\beta}$ and
$Y_{\alpha}\rightarrow Y_{\beta}$ are affine. Let $g$ be the inverse
limit of the $\{g_{\alpha}\}$ \hyperrefext[VII][VII.subsec5.1]{5.1}.
If every $g_{\alpha}$ is locally $0$-acyclic (\resp\ldots), then so is
$g$.
\end{enumerate}
\end{enonce*}

This follows by a chain of reasoning analogous to \hyperref[XV.prop1.9]{1.9}.

\begin{enonce*}[remark]{Remarks 1.14}\etiquette[1.14]{XV.rems1.14}
To simplify his task, the author omitted from \hyperref[XV.prop1.10]{1.10} through
\hyperref[XV.prop1.13]{1.13} a treatment of the case $n=-1$, and leaves it to the
reader to verify that the statements, definitions, and proofs apply to
this case essentially without change. Thus, with the notation of
\hyperref[XV.prop1.10]{1.10}, condition (i) says that the induced local morphisms
$\widetilde{X}\rightarrow \widetilde{Y}$ are $(-1)$-acyclic, i.e., by
\hyperref[XV.cor1.5]{1.5}, \emph{surjective}, and conditions (ii) through (iv) say
that the base-change morphism
${g'}^{\ast}({i'}_{\ast}(F))\rightarrow {j'}_{\ast}({g''}^{\ast}(F))$
is a \emph{monomorphism}. Under these conditions, we shall therefore
say that $g$ is \emph{locally $(-1)$-acyclic}. This condition is thus
stable\pageoriginale under locally quasi-finite base change. As an
example, a flat morphism is locally $(-1)$-acyclic, since for such a
morphism the surjectivity of the maps
$\widetilde{X}\rightarrow \widetilde{Y}$ is a well-known consequence
of their flatness. We also note that, when $f$ is locally of finite
presentation, $f$ is locally $(-1)$-acyclic if and only if it is
\emph{universally open}, at least when $Y$ is locally Noetherian; this
follows readily from \EGA IV 14.4.9, 1.10.3. For such a morphism, local
$(-1)$-acyclicity consequently implies \emph{universal} local
$(-1)$-acyclicity. (It is moreover very likely\footnote[1]{This was in
fact verified by M. Raynaud. \Cf the new edition of \EGA IV 14!} that,
in these last statements, the hypothesis that $Y$ be locally
Noetherian is unnecessary; this is at any rate so if the set of
irreducible components of $Y$ is locally finite \EGA IV 14.4.10.)
\end{enonce*}

\begin{enonce*}{Theorem 1.15}\etiquette[1.15]{XV.thm1.15}
Let $g:X\rightarrow Y$ be a quasi-compact and quasi-separated morphism,
let $\LL\subset \PP$ be a set of prime numbers, and let $n$ be a
natural number (\resp $n=1$). Suppose that $g$ is $(n-1)$-acyclic for
$\LL$ and locally $(n-1)$-acyclic for $\LL$
$\hyperref[XV.def1.3]{1.3}$, $\hyperref[XV.def1.7]{1.7}$, $\hyperref[XV.def1.11]{1.11}$, $\hyperref[XV.rems1.14]{1.14}$. Let $F$ be a sheaf on $Y$, which, for
$n\geq 1$, is assumed to have the structure of an abelian
$\LL$-torsion sheaf (\resp the structure of a sheaf of groups), and
let $\xi$ be an element of $H^{n}(X,g^{\ast}(F))$. In order that
$\xi$ come from an element of $H^{n}(Y,F)$ (which is then uniquely
determined, thanks to the $(n-1)$-acyclicity hypothesis on $g$), it is
necessary and sufficient that, for every geometric point $\bar{y}$ of
$Y$ algebraic over a point $y$ of $Y$, the image $\xi_{\bar{y}}$ of
$\xi$ in $H^{n}(X_{\bar{y}},F|X_{\bar{y}})$ lie in the image of
$H^{n}(\bar{y},F_{\bar{y}})$ (which, for $n\geq 1$,
means\pageoriginale that $\xi_{\bar{y}}=0$).
\end{enonce*}

We note at once that \hyperref[XV.thm1.15]{1.15} shows that, under the preliminary
conditions imposed on $g$, the morphism $g$ is $n$-acyclic for
$\LL$ (\resp $1$-aspherical for $\LL$) if and only if, for
every $\bar{y}$ as above, $X_{\bar{y}}$ is $n$-acyclic for
$\LL$ (\resp is $1$-aspherical for $\LL$): the ``if'' direction
follows directly from \hyperref[XV.thm1.15]{1.15} and the definitions, while the ``only
if'' direction is in any case immediate from \hyperref[XV.cor1.9.1]{1.9.1}. With this
said, an immediate induction on $n$ gives the following result.

\begin{enonce*}{Corollary 1.16}\etiquette[1.16]{XV.cor1.16}
Let $g$ be a quasi-compact and quasi-separated morphism that is locally
$(n-1)$-acyclic for $\LL$, where $n$ is a given natural number
(\resp where $n=1$). Then $g$ is $n$-acyclic for $\LL$
(\resp is $1$-aspherical for $\LL$) if and only if, for every
geometric point $\bar{y}$ of $Y$ algebraic over a point $y$ of $Y$,
the fiber $X_{\bar{y}}$ is $n$-acyclic for $\LL$ (\resp
$1$-aspherical for $\LL$).
\end{enonce*}

\begin{enonce*}{Corollary 1.17}\etiquette[1.17]{XV.cor1.17}
Let $g:X\rightarrow Y$ be a morphism of preschemes, let
$\LL\subset \PP$ be a set of prime numbers, and let $n$ be a
natural number (\resp $n=1$). In order that $g$ be locally

% END INCLUDED SGA4_English_current_buildable_source_20260906/source/15/drafts/SGA4_XV_LINES_401_700_English.texfrag


% BEGIN INCLUDED SGA4_English_current_buildable_source_20260906/source/15/drafts/SGA4_XV_LINES_701_969_English.texfrag
$n$-acyclic for $\LL$ (\resp locally $1$-aspherical for
$\LL$), it is necessary and sufficient that, for every geometric
point $\bar{x}$ of $X$, if $\widetilde{g}:\widetilde{X}\rightarrow
\widetilde{Y}$ is the morphism between the strict localizations of $X$
and $Y$ (at $\bar{x}$ and the corresponding geometric point $\bar{y}$
of $Y$) induced by $g$, every geometric fiber $X_{\bar{z}}$ of $g$,
with respect to a geometric point $\bar{z}$ of $Y$ algebraic over a
point $z$ of $Y$, is $n$-acyclic for $\LL$ (\resp
$1$-aspherical for $\LL$).
\end{enonce*}

Necessity follows from criterion \hyperref[XV.prop1.10]{1.10} (i) and
\hyperref[XV.cor1.9.1]{1.9.1}. For\pageoriginale sufficiency, we proceed by induction
on $n$, which allows us, when $n\geq 1$, to suppose that $g$ is already
locally $(n-1)$-acyclic for $\LL$. The same then holds for the
morphisms $\widetilde{g}$, and \hyperref[XV.cor1.16]{1.16} implies that
$\widetilde{g}$ is $n$-acyclic (\resp $1$-aspherical for
$\LL$); by criterion \hyperref[XV.prop1.10]{1.10} (i), this implies that
$g$ is locally $n$-acyclic (\resp locally $1$-aspherical) for
$\LL$.

\begin{enonce*}{Corollary 1.18}\etiquette[1.18]{XV.cor1.18}
With the notation of \hyperref[XV.cor1.17]{1.17}, if $g$ is locally of
finite type, then in the stated criterion it is enough to consider
geometric points $\bar{x}$ whose underlying point $x$ is closed in its
fiber.
\end{enonce*}

This follows from the proof just given and from \hyperref[XV.cor1.12]{1.12}.

\noindent
\emph{Proof of} \hyperref[XV.thm1.15]{1.15}\footnote[1]{The reader
interested only in the Noetherian case is invited to consult
\sisi{Remark }{}\hyperref[XV.rem1.18]{1.18}.}.
% There are two statements numbered 1.18 in this exposé.
We first show that we may suppose $Y$ strictly local. In the abelian
case, with $n\geq 1$, the Leray spectral sequence for $g$ and
$g^{\ast}(F)$, together with
$R^{i}g_{\ast}(g^{\ast}(F))=0$ for $0<i\leq n-1$ and
$F\xrightarrow{\sim}g_{\ast}g^{\ast}(F)$, which express the
$(n-1)$-acyclicity hypothesis on $g$, yields a canonical exact
sequence
$$
0\longrightarrow H^{n}(Y,F)\longrightarrow
H^{n}(X,g^{\ast}(F))\longrightarrow
H^{0}(Y,R^{n}g_{\ast}(g^{\ast}(F))){\quoi},
$$
which shows that it is enough to prove that the image of $\xi$ in the
third term is zero. This can be checked fiber by fiber, and in view of
\hyperrefext[VIII][VIII.thm5.2]{5.2} we are reduced to the case where
$Y$ is strictly local. The case $n=0$, with $F$ a sheaf of sets, is
handled in the same way: the hypothesis on $g$ implies that
$F\rightarrow g_{\ast}g^{\ast}(F)$ is injective, and $\xi$ is
identified with a section of the second sheaf that must be shown to
come from a section of the first. This too can be checked
fiber by\pageoriginale fiber, and one applies
\hyperrefext[VIII][VIII.rem5.3]{5.3}. Essentially the same argument
works in the noncommutative case, with $n=1$, using the exact sequence
\hyperrefext[XII][XII.prop3.2]{3.2}
$$
H^{1}(Y,F)\longrightarrow H^{1}(X,g^{\ast}(F))\longrightarrow
H^{0}(Y,R^{1}g_{\ast}(g^{\ast}(F))){\quoi},
$$
and using \hyperrefext[VIII][VIII.rem5.3]{5.3}.

We therefore suppose that $Y$ is strictly local, and hence that $X$ is
quasi-compact and quasi-separated. The morphism $g:X\rightarrow Y$
satisfies the conditions of \hyperrefext[XII][XII.cor6.10]{6.10},
which reduces us to the case in which, for every closed subset $Y_1$
of $Y$ distinct from $Y$, the restriction of $\xi$ to
$X_1=\fprod{X}{Y_1}{Y}$ lies in the image of
$H^{n}(Y_1,F|Y_1)\rightarrow H^{n}(X_1,g^{\ast}(F)|X_1)$, and to
checking in this case that one can find a finite morphism
$f:Y'\rightarrow Y$ whose image in $Y$ contains a nonempty open subset
of $Y$, and such that the inverse image $\xi'$ of $\xi$ on
$X'=\fprod{X}{Y'}{Y}$ lies in the image of
$H^{n}(Y',F')\rightarrow H^{n}(X',{g'}^{\ast}(F'))$. We note that, in
view of \hyperrefext[VIII][VIII.prop5.5]{5.5} and
\hyperrefext[VIII][VIII.rem5.8]{5.8}, condition $2^{\circ}$) of
\loccit is automatically satisfied for this morphism
$f:Y'\rightarrow Y$.

To find such an $f$, we may suppose that $Y$ is reduced, and we
consider a maximal point $y$ of $Y$. Let $\bar{y}$ be the geometric
point corresponding to a separable closure of $k(y)$, so that
$\bar{y}$ is the filtered inverse limit of open subsets $U_i$ of
finite integral $Y$-schemes $Y_i$, each of which has an image in $Y$
containing a neighborhood of $y$. Using the hypothesis on $\xi$, we
see that there is an index $i$ such that the inverse image of $\xi$ on
$\fprod{X}{U_i}{Y}$ is zero if $n\geq 1$, \resp lies in the image of
$H^{0}(U_i,F_{U_i})$ if $n=0$.

We\pageoriginale take $Y'$ to be $Y_i$, and note that we may suppose
that $U'=U_i$ is the inverse image of an affine open neighborhood $U$
of $y$ in $Y$. We then see that the inverse image of $\xi$ on
$\fprod{X}{Z'}{Y}$ (where $Z'=Y'-U'$ is the inverse image in $Y'$ of
$Z=Y-U$) lies in the image of $H^{n}(Z',F_{Z'})$. After replacing $Y$
by $Y'$, we are therefore reduced to the following case: there is an
affine open subset $U$ of $Y$, with complement $Z$, such that the
restrictions of $\xi$ to $X_U$ and $X_Z$ lie respectively in the image
of $H^{n}(U,F_U)$ and in the image of $H^{n}(Z,F_Z)$; moreover, if
$n\geq 1$, these restrictions may be assumed to be \emph{zero}.

When $n\geq 1$, we shall show that under these conditions $\xi=0$. To
simplify notation, we write $Y'$ in place of $X$, denote by $U'$ and
$Z'$ the inverse images of $U$ and $Z$ in $Y'$, and denote the inverse
image $g^{\ast}(F)$ by $F'$. Let
$i:U\rightarrow Y$, $j:Z\rightarrow Y$,
$i':U'\rightarrow Y'$, and $j':Z'\rightarrow Y'$ be the inclusions,
$$
\xymatrix@C=2cm@R=1.5cm{
U'  \ar[r]^{i'} \ar[d]^{g_U}& Y' = X  \ar[d]^{g} & Z'  \ar[l]_{j'} \ar[d]^{g_Z}\\
U   \ar[r]^{i} & Y & Z \ar[l]_{j},
}
$$
and consider the exact sequence of sheaves
$$
0\longrightarrow F'_{U',Y'}\longrightarrow F'\longrightarrow
F'_{Z',Y'}\longrightarrow 0{\quoi}.
$$
% Reference for this notation?
Since the image of $\xi$ in
$H^{n}(Y',F'_{Z',Y'})=H^{n}(Z',F'_{Z'})$ is zero, $\xi$ comes from an
element of $H^{n}(Y',F'_{U',Y'})$, and the restriction of the latter to
$U'$ is equal to that of $\xi$, and is therefore zero. Since
$F'_{U',Y'}$ is isomorphic to the inverse image of $F_{U,Y}$, we are
thus reduced\pageoriginale to the case of a sheaf of the form
$F_{U,Y}$; hence we may suppose that $F_Z=0$ (which takes account of
the hypothesis $\xi|Z'=0$). Note that if $F\rightarrow G$ is a
monomorphism of abelian $\LL$-torsion sheaves (\resp of sheaves of
groups), it is enough to prove that the image $\eta$ of $\xi$ in
$H^{n}(Y',G')$ lies in the image of $H^{n}(Y,G)$, \ie is zero
\hyperrefext[XII][XII.cor6.6]{6.6}. This allows us to replace $F$ by
$i_{\ast}(F_U)$, since the hypothesis $F_Z=0$ implies that the
canonical homomorphism $F\rightarrow i_{\ast}(F_U)$ is injective. Now,
since $g$ is locally $(n-1)$-acyclic for $\LL$, and a fortiori
$0$-acyclic for $\LL$ (because $n\geq 1$), the inverse image of
$i_{\ast}(F_U)$ on $Y'$ is identified with
$i'_{\ast}(F'_{U'})$. We are therefore reduced to showing that the
image $\eta$ of $\xi$ in $H^{n}(Y',i'_{\ast}(F'_{U'}))$ is zero. The
restriction of $\eta$ to $U'$ is zero, so it is enough to check that
the homomorphism
\begin{equation*}
H^{n}(Y',i'_{\ast}(F'_{U'}))\longrightarrow H^{n}(U',F'_{U'})\tag{$\ast$}
\end{equation*}
is a monomorphism. In the abelian case, we use the Leray spectral
sequence for $i'$ and $F'_{U'}$:
$$
E^{p,q}_2=H^p(Y',R^qi'_{\ast}(F'_{U'}))\Longrightarrow
H^{\ast}(U',F'_{U'}){\quoi},
$$
and the relations
\begin{equation*}
E^{p,q}_2=0\; \mbox{\rm for } 0<p\leq n-1\quav q\leq n-1{\quoi}.\tag{${\ast}{\ast}$}
\end{equation*}
To verify the latter, note that the local $(n-1)$-acyclicity
hypothesis on $g$ for $\LL$ implies that there are isomorphisms
$$
R^{q}i'_{\ast}(F'_{U'})\simeq g^{\ast}(R^{q}i_{\ast}(F_U))\;\mbox{\rm for } q\leq n-1{\quoi},
$$
and,\pageoriginale in view of the fact that $g$ is $(n-1)$-acyclic
for $\LL$, so that
$H^{p}(Y,G)\xrightarrow{\sim}H^{p}(Y',G')$ for $p\leq n-1$ and every
abelian $\LL$-torsion sheaf $G$ on $Y$, we obtain
$({\ast}{\ast})$, since $H^{p}(Y,G)=0$ for $p>0$. In the same way we
obtain
\begin{equation*}
E^{0,q}_2\simeq H^{0}(Y,R{}^qi_{\ast}(F_U))\;\mbox{\rm
for}\quad q\leq n-1{\quoi}.\tag{${\ast}{\ast}{\ast}$}
\end{equation*}
The Leray spectral sequence under consideration then gives rise to an
exact sequence that forms the second row of the following commutative
diagram:
$$
\xymatrix{
**[r]H^{0}(Y,R^{n-1} i_{\ast}(F_U))  \ar[r] \ar[d]^{\text{\rotatebox{90}{$\sim$}}}&
**[r]H^n(Y,i_{\ast} (F_U)) = 0 \ar[d]\\
**[r]H^{0}(Y', R^{n-1} i'_{\ast}(F'_{U'}))  \ar[r] &
**[r]H^{n}(Y',i'_{\ast}(F'_{U'}))  \ar[r] & **[r]H^{n}(U',F'_{U'}),
}
$$
where the first vertical arrow is bijective by
$({\ast}{\ast}{\ast})$. The diagram shows that $(\ast)$ is injective,
which completes the proof in this case. In the nonabelian case,
$n=1$, we use the same diagram with $n=1$:
$$
\xymatrix{
**[r]H^{0}(Y, i_{\ast}(F_U))  \ar[r] \ar[d]^{\text{\rotatebox{90}{$\sim$}}}&
**[r]H^1(Y,i_{\ast} (F_U)) = 0 \ar[d]\\
**[r]H^{0}(Y', i'_{\ast}(F'_{U'}))  \ar[r] &
**[r]H^{1}(Y',i'_{\ast}(F'_{U'}))  \ar[r] & **[r]H^{1}(U',F'_{U'}),&
}
$$
\hyperrefext[XII][XII.prop3.2]{3.2}, where the first vertical arrow is
again bijective by the hypothesis of local and global $0$-acyclicity
on $g$; once again we conclude that $(\ast)$ is injective in this
case.

Finally, consider the case $n=0$. By virtue of
\sisi{\qquad}{\qquad}, specifying a section $\xi$ of $F'$ amounts to
% Missing reference
specifying a section $\xi_{U'}$ of $F'_{U'}$ and a section
$\xi_{Z'}$ of $F'_{Z'}$ such that the image of $\xi_{Z'}$ under the
natural homomorphism
$$
F'_{Z'}\longrightarrow {j'}^{\ast}(i'_{\ast}(F'_{U'}))
$$
is\pageoriginale the section of the right-hand side induced by
$\xi_{U'}$. By hypothesis, there are sections $\eta_Z$ and $\eta_U$ of
$F_Z$ and $F_U$, respectively, that induce $\xi_{Z'}$ and
$\xi_{U'}$. It remains to see that they satisfy the analogous
compatibility condition, which says that they are the restrictions of
a single section $\eta$ of $F$. We must verify the equality of two
sections of $j^{\ast}(i_{\ast}(F_U))$. Since $Z'\rightarrow Z$ is
surjective ($g$ being surjective because it is $(-1)$-acyclic by
hypothesis), it is enough to check that the two corresponding sections
of
$$
g_Z{}^{\ast}(j^{\ast}(i_{\ast}(F_U)))
=j'{}^{\ast}(g^{\ast}(i_{\ast}(F_U)))
$$
are equal. Through the base-change homomorphism
\begin{equation*}
g^{\ast}(i_{\ast}(F_U))\longrightarrow
i'_{\ast}(g_U{}^{\ast}(F_U))=i'_{\ast}(F'_{U'}){\quoi},\tag{$\ast$}
\end{equation*}
the sheaf under consideration maps to
${j'}^{\ast}(i'_{\ast}(F'_{U'}))$, and this latter homomorphism is
injective because $(\ast)$ is injective, by the hypothesis that $g$ is
locally $(-1)$-acyclic. It therefore suffices to verify that the two
sections of ${j'}^{\ast}(i'_{\ast}(F'_{U'}))$ that are the images of
the preceding two sections are equal. But these are immediately seen
to be precisely the sections considered above, which are equal by
hypothesis. This completes the proof.

\begin{enonce*}{Remark 1.18}\etiquette[1.18]{XV.rem1.18}
There is a refinement of \hyperref[XV.thm1.15]{1.15} that gives a much
simpler proof of \hyperref[XV.thm1.15]{1.15} when $Y$ is assumed to be
locally Noetherian. In addition to the preliminary hypotheses on $g$
and $F$, suppose that there is a finite set of geometric points
$f_i:\bar{y}_i\rightarrow Y$ of $Y$ such that the morphism
$$
F\longrightarrow \prod_if_{i\ast}(f_i{}^{\ast}(F))
$$
is\pageoriginale injective (a hypothesis automatically satisfied if
$Y$ is Noetherian and $F$ is constructible
\hyperrefext[IX][IX.prop2.14]{2.14}). I claim that under these
conditions criterion \hyperref[XV.thm1.15]{1.15} remains valid when
one restricts to the geometric points $\bar{y}_i$ alone. Indeed, by
\hyperrefext[XII][XII.cor6.6]{6.6}, we may suppose that $F$ itself is
of the form $f_{\ast}(M_{\bar{y}})$, where $M$ is a set \resp a group
\resp an abelian $\LL$-torsion group, and
$f:\bar{y}\rightarrow Y$ is a geometric point. Let
$f':X'=X_{\bar{y}}\rightarrow X$ be the canonical morphism. Assuming,
for definiteness, that $n\geq 1$, the local hypothesis on $g$ implies
that
$$
g^{\ast}(f_{\ast}(M_{\bar{y}}))\simeq f'_{\ast}(M_{X'})\sisi{,\;}{\quav}
R^{i}f'_{\ast}(M_{X'})=0\quad\mbox{\rm for}\quad 1\leq i\leq
n-1{\quoi},
$$
the second relation following from the fact that its left-hand side is
isomorphic to $g^{\ast}(R^{i}f_{\ast}(M_{\bar{y}}))$, which is zero
since manifestly $R^{i}f_{\ast}(M_{\bar{y}})=0$ for $i\geq 1$. From
these relations we conclude that
$$
H^{n}(X,g^{\ast}(F))\simeq H^{n}(X,f'_{\ast}(M_{X'}))\longrightarrow
H^{n}(X',M_{X'})
$$
is injective, using the Leray spectral sequence for $f'$. This proves
that an element $\xi$ of $H^{n}(X,g^{\ast}(F))$ whose induced element
in $H^{n}(X_{\bar{y}},M_{X_{\bar{y}}})$ is zero is itself zero. To
deduce from this a simplified proof of \hyperref[XV.thm1.15]{1.15} in
the case where $Y$ is locally Noetherian, note that we may restrict to
the case where $Y$ is Noetherian; an easy passage to the limit, using
\hyperrefext[IX][IX.cor2.7.2]{2.7.2} and
\hyperrefext[VII][VII.sec5]{5}, then reduces us to the case where $F$
is also constructible, to which the statement just proved applies.
\end{enonce*}

% END INCLUDED SGA4_English_current_buildable_source_20260906/source/15/drafts/SGA4_XV_LINES_701_969_English.texfrag


% BEGIN INCLUDED SGA4_English_current_buildable_source_20260906/source/15/drafts/SGA4_XV_LINES_970_1136_English.texfrag
\section{Local acyclicity of a smooth morphism}\etiquette[2]{XV.sec2}\pageoriginale

The theorem is the following:

\begin{enonce*}{Theorem 2.1}\etiquette[2.1]{XV.thm2.1}
Let $g:X\rightarrow Y$ be a smooth morphism, and let $\LL\subset
\PP$ be the complement of the set of residual characteristics of
$X$. Then $g$ is locally acyclic and locally $1$-aspherical for
$\LL$.
\end{enonce*}

Let us first record the corollary:

\begin{enonce*}{Corollary 2.2}\etiquette[2.2]{XV.cor2.2}
Let $g:\EE^m_Y\rightarrow Y$ be the structural morphism of affine
$m$-space. Then $g$ is acyclic and $1$-aspherical for the set $\LL$
complementary to the set of residual characteristics of $Y$.
\end{enonce*}

Indeed, induction on $m$ and \hyperref[XV.prop1.9]{1.9} (i) reduce us
to the case $m=1$. By \hyperref[XV.thm2.1]{2.1}, $g$ is locally
$n$-acyclic and locally $1$-aspherical for $\LL$, so the \emph{local}
condition of \hyperref[XV.cor1.16]{1.16} is satisfied. It therefore
remains only to prove that the geometric fibers of $\EE^1_Y$ are
acyclic and $1$-aspherical for $\LL$; that is, that $\Spec k[t]$ is
acyclic and $1$-aspherical for $\LL$ when $k$ is a separably closed
field of characteristic $p\notin\LL$. The acyclicity assertion follows
from \hyperrefext[IX][IX.thm4.6]{4.6}, and the $1$-asphericity assertion
is well known. More generally, let $X'$ be a normal irreducible Galois
covering of $X=\PP^1_k$, unramified away from the point $\infty$ and
tamely ramified at that point. We prove that the degree $n$ of the
covering is $1$. The Hurwitz formula gives
$$
2(g'-1)=2n(g-1)+\deg \mathfrak{d}_{X'/X},
$$
where $g$ and $g'$ are the genera of $X$ and $X'$, and
$\mathfrak{d}_{X'/X}$ is the different divisor. We have $g=0$, while
tame ramification implies $\deg \mathfrak{d}_{X'/X}\leq n-1$; hence
$2(g'-1)\leq-(n+1)$, so $n\leq2(1-g')-1$. Since $g'\geq0$, it follows
that $n\leq1$, and therefore $n=1$, \cqfd

\medskip
\noindent
\emph{Proof of} \hyperref[XV.thm2.1]{2.1}. The assertion is local on
$X$ and $Y$ for the étale topology, so we may assume
\sisi{(}{}\SGA 1 II 1.1\sisi{)}{} that $g$ is the structural morphism
of affine $m$-space $\EE^m_Y$ and that $Y$ is affine. By the
transitivity of local acyclicity, \hyperref[XV.prop1.13]{1.13} (i),
and the identity $\EE^m_Y=\EE^{m-1}_{\EE^1_Y}$, we are immediately
reduced to $m=1$, that is, to $X=\Spec\Oo_Y[t]$.

Apply \hyperref[XV.cor1.18]{1.18}. It remains to prove that the
geometric fibers of the morphism of strictly local schemes
$\widetilde g:\widetilde X\rightarrow\widetilde Y$ associated with a
pair of geometric points $\bar x,\bar y$, where $\bar x$ is closed in
the fiber $X_{\bar y}$, are acyclic for $\LL$ and $1$-aspherical for
$\LL$. We may plainly assume that $Y=\widetilde Y$. Let
$Y'\rightarrow Y$ be a finite, local, surjective morphism, so that
$Y'$ is strictly local, and let $\bar y'$ be the point of $Y'$ above
$\bar y$ and $\bar x'$ the point of $X'$ above both $\bar x$ and
$\bar y'$. Since $X'\rightarrow X$ is finite, the strict localization
of $X'$ at $\bar x'$ is precisely
$\widetilde X'=\fprod{\widetilde X}{Y'}{Y}$. Consequently every
geometric fiber of $\widetilde X'/Y'$ is obtained from a geometric
fiber of $\widetilde X/Y$ by an algebraic (necessarily radicial)
extension of the base field; moreover, at least one geometric fiber of
$\widetilde X'/Y'$ lies above each geometric fiber of
$\widetilde X/Y$. Thus, by \hyperrefext[VIII][VIII.thm1.1]{1.1}, we
may replace $Y$ by $Y'$. Choosing $Y'$ suitably reduces us to the case
$k(\bar x)=k(\bar y)$. Then $\bar x$ is a rational point of the fiber
$X_{\bar y}=\Spec k(\bar y)[t]$, say the point $t=a^0$. Lift $a^0$ to
a section $a$ of $\Oo_{\sisi{Y}{\widetilde Y}}$ and make the
``translation'' $t_1=t-a$; we are reduced to the case in which
$\bar x$ is the point $\{t=0\}$ of $\Spec k(\bar y)[t]$. Put
$A=\Gamma(\sisi{Y}{\widetilde Y},
\Oo_{\sisi{Y}{\widetilde Y}})$, so that $A$ is a strictly local
ring.

\begin{enonce*}[remark]{Notation 2.3}\etiquette[2.3]{XV.not2.3}
Let $A$ be a Henselian local ring. We denote by $A\{t\}$ the
Henselization of $A[t]$ at the prime ideal generated by $\rad A$ and
$t$.
\end{enonce*}

\begin{enonce*}{Corollary 2.4}\etiquette[2.4]{XV.cor2.4}
Let $A$ be Henselian and let $A'$ be a finite $A$-algebra. Then
$$
A'\{t\}\simeq A'\otimes_A A\{t\}{\quoi}.
$$
\end{enonce*}

This follows immediately from \hyperrefext[VIII][VIII.subsec4.1]{4.1}.
Note that if $A$ is strictly local, then so is $A\{t\}$.

The proof of the theorem is reduced to the following lemma:

\begin{enonce*}{Lemma 2.5}\etiquette[2.5]{XV.lem2.5}
Let $A$ be a Henselian ring of residual characteristic $p$. Then every
geometric fiber of $\Spec A\{t\}$ over $\Spec A$ is acyclic and
$1$-aspherical for $\LL=\PP-\{p\}$.
\end{enonce*}

Since $A\{t\}$ is the Henselization of $A[t]$, a fiber
$\Spec(A\{t\}\otimes_A K)$, where $K$ is a residue field of $A$, is
an inverse limit\pageoriginale of affine schemes $X_i$ of finite type
over the ``line'' $\Spec K[t]$, that is, a limit of affine algebraic
curves. Hence every geometric fiber of $\Spec A\{t\}$ is a limit of
affine curves $(X_i)_{\bar K}$ over the separably closed field
$\bar K$. It follows from \hyperrefext[IX][IX.cor5.7]{5.7} and
\hyperrefext[VII][VII.thm5.7]{5.7} that the cohomological dimension of
the geometric fibers is at most $1$. It is therefore enough to verify
that the geometric fibers are $1$-\emph{aspherical} for $\LL$; that
is, that such a geometric fiber $\bar Z$ is connected and nonempty and
that $H^1(\bar Z,G)=0$ for every finite $\LL$-group $G$.

Write $A=\varinjlim B_\alpha$, where the $B_\alpha$ are rings of
finite type over $\ZZ$. The morphism $B_\alpha\rightarrow A$ factors
through the Henselization $A_\alpha$ of $B_\alpha$ at the prime ideal
induced by $\rad A$, whence $A=\varinjlim A_\alpha$. One readily
checks that $A\{t\}=\varinjlim A_\alpha\{t\}$ as well. Let $\bar y$ be
a geometric point of $\Spec A$, let $X_{\bar y}$ be the geometric
fiber of $S=\Spec A\{t\}$ over $\Spec A$ at $\bar y$, let
$\bar y_\alpha$ be the point of $\Spec A_\alpha$ induced by $\bar y$,
and let $X_{\alpha,\bar y_\alpha}$ be the geometric fiber of
$\Spec A_\alpha\{t\}$ over $\Spec A_\alpha$ at $\bar y_\alpha$. We
have $X_{\bar y}=\varprojlim X_{\alpha,\bar y_\alpha}$. By
\hyperref[XV.prop1.9]{1.9} (ii), applied to the morphisms
$X_{\alpha,\bar y_\alpha}\rightarrow\bar y_\alpha$, the proof of
\hyperref[XV.lem2.5]{2.5} is therefore reduced to the case
$A=A_\alpha$, that is, to the case in which $A$ is the Henselization
of an algebra of finite type over $\ZZ$ at a given prime ideal. We
assume this from now on. Recall that such an algebra is
\emph{excellent} \sisi{(}{}\EGA IV 7.8.3 (iii), 7.8.6 (i)\sisi{)}{},
and hence that $A_\alpha$ is excellent as well (as follows readily
from the definition; \cf \EGA IV 18.7.6).

Let us verify $0$-acyclicity. (For a more general result, see the
Appendix.) Since $A\{t\}/A$ has the section ``$t=0$'', a geometric
fiber\pageoriginale $\bar Z$ is nonempty. To see that it is connected,
we reduce, by replacing $A$ by a quotient
\sisi{(}{}\hyperref[XV.cor2.4]{2.4}\sisi{)}{}, to the case in which
$A$ is integral and $\bar Z$ is the geometric fiber over the generic
point of $\Spec A$. It is enough to prove that for every finite
separable extension $K'$ of the fraction field $K$ of $A$, the scheme
$Z_{K'}$, where $Z=\Spec(A\{t\}\otimes_A K)$, is connected. Let $A'$
be the normalization of $A$ in $K'$ (a finite $A$-algebra,
\sisi{(}{}\EGA IV 7.8.2\sisi{)}{}) and let $Z'$ be the generic fiber
of $\Spec A'\{t\}$ over $A'$. By \hyperref[XV.cor2.4]{2.4},
$A'\{t\}\simeq A'\otimes_A A\{t\}$, and hence
$Z_{K'}\simeq Z'$. Replace $A$ by $A'$. It is then enough to prove
that $Z$ is integral, and for this it is enough to prove that
$A\{t\}$ is integral. Since $A$ is normal, $A[t]$ is normal, and so
is $A\{t\}$ \sisi{(}{}\SGA 1 I 9.5 (i)\sisi{)}{}. This proves the
assertion.

The proof of \hyperref[XV.lem2.5]{2.5}, and hence of
\hyperref[XV.thm2.1]{2.1}, is thus reduced to the following lemma,
which will be proved in the next section:

\begin{enonce*}{Lemma 2.6}\etiquette[2.6]{XV.lem2.6}
Let $A$ be a Henselian excellent ring, and let $\bar Z$ be a geometric
fiber of $\Spec A\{t\}/\Spec A$. Then every principal Galois covering
$\bar Z'$ of $\bar Z$ with group $G$ of order prime to the residual
characteristic $p$ of $A$ is trivial.
\end{enonce*}

% END INCLUDED SGA4_English_current_buildable_source_20260906/source/15/drafts/SGA4_XV_LINES_970_1136_English.texfrag


% BEGIN INCLUDED SGA4_English_current_buildable_source_20260906/source/15/drafts/SGA4_XV_LINES_1137_1315_English.texfrag
\section{Proof of the main lemma}\etiquette[3]{XV.sec3}\pageoriginale

Replacing $A$ by a finite $A$-algebra $A'$ (\cf
\hyperref[XV.cor2.4]{2.4}), assertion \hyperref[XV.lem2.6]{2.6} is reduced
to the case where $A$ is integral and $\bar{Z}$ is the geometric generic
fiber of $\Spec A\{t\}$ over $\Spec A$. Let $K$ be the fraction field of
$A$, and let $Z=\Spec A\{t\}\otimes_A K$. If $\bar{Z}'/\bar{Z}$ is a
principal Galois covering, it can be descended to some $Z_{K'}$, where
$K'$ is a suitable finite separable extension of $K$. Replacing $A$
(\hyperref[XV.cor2.4]{2.4}) by its normalization in $K'$, which is a
finite $A$-algebra \sisi{(}{}\EGA IV 7.8.3 (vi)\sisi{)}{}, we are reduced
to proving the lemma in the following form:

\begin{enonce*}{Lemma 3.1}\etiquette[3.1]{XV.lem3.1}
With the notation of \hyperref[XV.lem2.6]{2.6}, suppose that $A$ is
normal, and let $Z'/Z$ be a principal étale Galois covering, with group
$G$ of order prime to the residual characteristic $p$ of $A$. Then $Z'$
is induced by a Galois extension $K'$ of $K$; that is, \emph{(\SGA 1 IX
6.2)}, the covering $\bar{Z}'$ of $\bar{Z}$ deduced from $Z'$ is trivial.
\end{enonce*}

\noindent
\emph{Proof}. Let $B'$ be the normalization of $A\{t\}$ in the ring of
rational functions of $Z'$, and consider the height-$1$ prime ideals $P$
of $A\{t\}$ at which the $A\{t\}$-algebra $B'$ is ramified. Since $Z'/Z$
is étale, $P$ does not lie over the generic point of $\Spec A$; that is,
$\mathfrak{p}=P\cap A\neq 0$. Thus
$A\{t\}/\mathfrak{p}A\{t\}=(A/\mathfrak{p})\{t\}$ is integral and has
dimension at most $\dim A$. Since $A\{t\}$ is excellent
\sisi{(}{}\EGA IV 7.8.6 (i) and 18.7.6\sisi{)}{}, hence catenary, we have
$\dim \Spec A\{t\}/P=(\dim A+1)-1=\dim A\geq
\dim A\{t\}/\mathfrak{p}A\{t\}$, whence\pageoriginale
$P=\mathfrak{p}A\{t\}$, since $P\supset \mathfrak{p}A\{t\}$.

Let $\mathfrak{p}_1,\ldots,\mathfrak{p}_m$ therefore be the height-$1$
prime ideals of $A$ such that $B'/A\{t\}$ is ramified at the ideals
$\mathfrak{p}_iA\{t\}$, and let $s$ be the least common multiple of
the orders of the inertia groups \sisi{(}{}\SGA IV 2\sisi{)}{}
$G_i\subset G$ for $B'$ at the prime ideals lying above the
$\mathfrak{p}_iA\{t\}$. This is an integer dividing the order of $G$,
and is therefore prime to $p$. Let $f\in A$ be an element that vanishes
to order~$1$ at every $\mathfrak{p}_i$ (such an $f$ exists, as follows,
for example, from Bourbaki, Alg. Comm., Chap. II, \S 3, cor. to prop.
17). Replacing $A$ \sisi{(}{}\hyperref[XV.cor2.4]{2.4}\sisi{)}{} by its
normalization $A_1$ in the finite extension
$K_1=K[x]/(x^s-f)$ of $K$, and replacing $B'$ by the normalization of
$A_1\otimes_A B'$, Abhyankar's lemma \sisi{(}{}\SGA 1 X
3.6\sisi{)}{} reduces us to the case where $B'$ is ramified at no
height-$1$ prime ideal of $A\{t\}$.

The algebra $B'/A\{t\}$ is then also étale over the fiber of
$\Spec A\{t\}$ at every codimension-$1$ point of $\Spec A$. Indeed, let
$\mathfrak{p}$ be a height-$1$ prime ideal of $A$. Since $A$ is normal,
$A_{\mathfrak{p}}$ is a discrete valuation ring, hence regular, and so
$A\{t\}\otimes_A A_{\mathfrak{p}}$ is regular as well---indeed, it is a
limit of schemes étale over $A_{\mathfrak{p}}[t]$. By hypothesis, the
algebra $B'/A\{t\}$ is étale at every point of codimension $\leq 1$ of
$\Spec A\{t\}\otimes_A A_{\mathfrak{p}}$, hence it is étale everywhere
by the Zariski--Nagata purity theorem \sisi{(}{}\SGA 1 X
3.1\sisi{)}{}; for a proof, see for example \SGA 2, X 3.4.

Replacing $A$ \sisi{(}{}\hyperref[XV.cor2.4]{2.4}\sisi{)}{} by a
suitable finite $A$-algebra, we are reduced to the case where, in
addition, the algebra $B'/tB'$ over
$A\{t\}/tA\{t\}\simeq A$ is completely split over the generic point of
$\Spec A$. It therefore suffices to prove that, under these conditions,
the algebra $B'$ over $A\{t\}$\pageoriginale is ``trivial,'' and, since
$A\{t\}$ is Henselian, it suffices to prove that $B'$ is étale over
$A\{t\}$. Indeed, the fact that $B'/tB'$ splits over the generic point
of $\Spec A$ implies that it splits over $\Spec(A)$
\sisi{(}{}\SGA 1 I 10.1\sisi{)}{}; thus the residual extensions of $B'$
over the maximal ideal of $A\{t\}$ are trivial, and we apply \EGA IV
18.5.14. We are therefore reduced to the following lemma:

\begin{enonce*}{Lemma 3.2}\etiquette[3.2]{XV.lem3.2}
Let $B$ be an excellent normal local ring, and let $0\neq t\in \rad B$
be such that $B/tB$ is normal. Let $B'$ be a finite integral normal
$B$-algebra containing $B$, such that:
\begin{enumerate}
\renewcommand{\theenumi}{\roman{enumi}}
\renewcommand{\labelenumi}{(\theenumi)}
\item The algebra $B'/tB'$ over $B/tB$ is completely split at the
maximal points of $\Spec(B/tB)$.

\item The algebra $B'$ over $B$ is étale at the codimension-$1$ points
of $\Spec(B/tB)$.
\end{enumerate}
Then $B'$ is étale over $B$.
\end{enonce*}

\noindent
\emph{Proof}. By (ii), the assertion is trivial for $\dim B\leq 2$,
that is, for $\dim B/tB\leq 1$. Suppose that $B'/B$ is not étale; let
$x$ be a maximal point of the set of points of $\Spec B/tB$ over which
$B'/B$ is not étale, and let $C$ be the local ring of $\Spec B$ at $x$.
Then $C$ is again an excellent normal ring (\EGA IV 7.8.2), we have
$0\neq t\in\rad C$, and (i) and (ii) hold for the $C$-algebra
$C'=C\otimes {}_BB'$. We may therefore suppose that $B=C$; that is,
$B'$ is étale over $B$ above every point of $\Spec B/tB$, except
possibly at the origin. Since $B/tB$ is normal\pageoriginale (in fact,
unibranchness would suffice), condition (i) implies that the algebra $B'/tB'$ is
completely split away from the origin. Moreover, we may suppose that
$\dim B\geq 3$. We are thus reduced to the following lemma:

\begin{enonce*}{Lemma 3.3}\etiquette[3.3]{XV.lem3.3}
Let $B$ be an excellent normal local ring, let $0\neq t\in \rad B$,
and suppose that $\dim B\geq 3$. Let $B'$ be a finite integral normal
$B$-algebra containing $B$, such that the $B/tB$-algebra $B'/tB'$ is
completely split away from the origin. Then $B'$ is étale over $B$.
\end{enonce*}

\noindent
\emph{Proof}. Let $X=\Spec B$, let $X_1=X-\{x_0\}$, where $x_0$ is the
closed point, and let $i:U\rightarrow X$ be the open immersion whose
image consists of the points $x\in X$ at which the algebra $B'/B$ is
étale. Then
$(X-U)\cap V(t)=\{x_0\}$ has dimension $0$. Since $V(t)$ is defined by
\emph{one} equation, it follows that $\dim(X-U)\leq 1$ (\EGA
$0_{\text{\rm IV}}$ 16.2.5); consequently (since $X$ is catenary),
$x\in X-U\Rightarrow \dim \mathcal{O}_{X,x}\geq \dim B-1\geq 2$; that
is, $\codim(X-U,X)\geq 2$.

Let $X'=\Spec(B')$, and let $U'$ be the inverse image of $U$ in $X'$.
Since $X$ is universally catenary, the relation
$\codim(X-U,X)\geq 2$ implies
$\codim(X'-U',X')\geq 2$ (\EGA IV 5.6.10). Since $X'$ is normal, we
therefore have
$$
B'=\Gamma(X',\Oo_{X'})\xrightarrow{\sim}
\Gamma(U',\Oo_{U'}){\quoi}.
$$
Let $\mathcal{B}'=\widetilde{B}'$ be the sheaf of $\Oo_X$-algebras that
defines $X'$ over $X$, and let $\mathcal{C}$ be its restriction to $U$.
The preceding formula is then equivalent to
$$
\mathcal{B}'\xrightarrow{\sim}i_{\ast}(\mathcal{C}){\quoi},
$$
and\pageoriginale \hyperref[XV.lem3.3]{3.3} is consequently reduced to
the following statement (in which $B'$ no longer occurs): Let $B$, $t$ be
as in \hyperref[XV.lem3.3]{3.3}, let $X=\Spec(B)$, let $Y=\Spec(B/tB)$, let $U$ be an open
neighborhood of $Y_1=Y-\{x_0\}$ in $X_1=X-\{x_0\}$, and let $U'$ be an
étale covering of $U$ whose restriction to $Y_1$ is completely split.
If $\mathcal{C}$ denotes the sheaf of $\Oo_U$-algebras defining $U'$,
then (with $i:U\rightarrow X$ denoting the canonical immersion)
$i_{\ast}(\mathcal{C})$ is a coherent \emph{étale} algebra on $X$ (or,
equivalently, in view of the relation $\prof \Oo_{X,x}\geq 2$ for
$x\in X-U$ noted above and \EGA IV 5.10.5, $U'$ extends to an
\emph{étale} covering $X'$ of $X$).

Let $\widehat{X}=\Spec\widehat{B}$, where $\widehat{B}$ is the
completion of $B$. Then $\widehat{B}$ is normal (\EGA IV 7.8.3 (v)) and
$\widehat{B}/B$ is faithfully flat. Let
$$
\xymatrix@C=1.5cm@R=1cm{
\hat{U}      \ar[r]^{\sisi{\hat{i}}{\hat{i}}} \ar[d]^{f_1}& \hat{X} \ar[d]^{f}\\
U            \ar[r]^{i} &  X
}
$$
be the Cartesian diagram deduced from $\widehat{X}\rightarrow X$.
Since $f$ is faithfully flat, the functor $i_{\ast}$ commutes with base
change by $f$ \sisi{(}{}\EGA IV 2.3.1\sisi{)}{}, and we are reduced
\sisi{(}{}\EGA IV 17.7.1\sisi{)}{} to proving that
$\hat{i}_{\ast}(\widehat{\mathcal{C}})$ is an étale algebra on
$\widehat{X}$. We may therefore replace $B$ by $\widehat{B}$; that is,
we may suppose that $B$ is \emph{complete}. In this case we shall prove
that $U'$ is a completely split covering of $U$ (which will plainly
imply that it extends to an étale covering of $X$, and will complete the
proof). This assertion is now an immediate consequence of ``local
Lefschetz theory'' \sisi{(}{}\SGA\pageoriginale 2 X 2.1 (i) and
2.5\sisi{)}{}, as is seen at once, taking account once again of the
relation $\prof \Oo_{X,x}\geq 2$ for $\dim \{x\}=1$.

\begin{enonce*}[remark]{Remark 3.4}\etiquette[3.4]{XV.rem3.4}
One can prove \hyperref[XV.lem3.2]{3.2} and
\hyperref[XV.lem3.3]{3.3} under more general conditions, by following
essentially the same proof: instead of assuming that $B$ is excellent
and normal, it is enough to suppose that the completion $\widehat{B}$
is normal (a condition stable under localization, by the results of
\EGA IV 7), or merely that the completion of $B$ for the $(tB)$-adic
topology is normal and is a quotient of a regular ring.
\end{enonce*}

% END INCLUDED SGA4_English_current_buildable_source_20260906/source/15/drafts/SGA4_XV_LINES_1137_1315_English.texfrag


% BEGIN INCLUDED SGA4_English_current_buildable_source_20260906/source/15/drafts/SGA4_XV_LINES_1316_1417_English.texfrag
\section{Appendix: A criterion for local $0$-acyclicity}\etiquette[4]{XV.sec4}\pageoriginale 

The results of this section will not be used later in the seminar.

\begin{enonce*}{Theorem 4.1}\etiquette[4.1]{XV.thm4.1}
Let $g:X\rightarrow Y$ be a morphism of schemes. Assume that $g$ is flat
with separable fibers, and that $X$ and $Y$ are locally noetherian
\sisi{(}{}respectively, that $g$ is locally of finite presentation\sisi{)}{}.
Then $g$ is universally locally $0$-acyclic.
\end{enonce*}

We may clearly assume that $X$ and $Y$ are affine; in the second case,
the standard method of \EGA IV 8, using \EGA IV 9.7.7, then reduces us to
the case in which $Y$, and hence $X$, is noetherian. We may therefore
confine ourselves to proving the first assertion. Since the hypothesis
is stable under base change of finite type, the passage to the limit in
\hyperref[XV.prop1.13]{1.13}~(ii) reduces us to proving that $g$ is locally
$0$-acyclic. For this, \hyperref[XV.cor1.17]{1.17} reduces us to proving that,
if $g$ is moreover assumed to be strictly local, then the geometric
fibers of $g$ are connected. This is precisely \EGA IV 18.9.8.

\begin{enonce*}{Corollary 4.2}\etiquette[4.2]{XV.cor4.2}
Let $g:X\rightarrow Y$ be a morphism of schemes, with $Y$ discrete.
Then $g$ is universally locally $0$-acyclic.
\end{enonce*}

By \hyperrefext[VIII][VIII.thm1.1]{1.1}, we may assume that $Y$ is the
spectrum of a perfect field $k$. We may also assume that $X$ is affine,
and hence is a projective limit of affine schemes of finite type over
$k$; in view of \hyperref[XV.prop1.13]{1.13}~(ii), this reduces us to the
case in which $X$ is of finite type over~$k$. Finally,\pageoriginale we may
clearly assume that $X$ is reduced, and it is then separable over $k$
because $k$ is perfect. The conclusion now follows from
\hyperref[XV.thm4.1]{4.1}.

\begin{enonce*}{Corollary 4.3}\etiquette[4.3]{XV.cor4.3}
Assume that $g$ is surjective and satisfies the hypotheses of
\hyperref[XV.thm4.1]{4.1} or \hyperref[XV.cor4.2]{4.2}; in the
non-parenthetical case of \hyperref[XV.thm4.1]{4.1}, assume in addition
that $g$ is either quasi-compact and quasi-separated or universally
open. Then $g$ is a morphism of universal effective descent for the
fibered category of étale sheaves on varying preschemes.
\end{enonce*}

In the cases under consideration, $g$ is universally submersive, and it
follows from \hyperrefext[VIII][VIII.prop9.1]{9.1} that $g$ is a morphism
of universal descent for the fibered category in question. Thus, for
the question of effectivity, we may confine ourselves to the case in
which $Y$ is affine. Moreover, in each of the cases under consideration,
$X$ can be covered by a \emph{finite} number of affine open subsets
$X_i$ whose images cover $Y$. Replacing $X$ by the disjoint union of the
$X_i$, we may therefore assume that $X$ is affine, and hence a fortiori
quasi-compact and quasi-separated over $Y$. We may then conclude by the
argument of \hyperrefext[VIII][VIII.subsubsec9.4.1]{9.4.1}, using the fact
\hyperrefext[XVI][XVI.thm1.1]{1.1} that, for a quasi-compact and
quasi-separated morphism $g$, the formation of $g_{\ast}(F)$ (where $F$
is an étale sheaf on $X$) commutes with every base change
$Y'\rightarrow Y$ that is locally $0$-acyclic.

\begin{enonce*}[remark]{Remarks 4.4}\etiquette[4.4]{XV.rems4.4}
\begin{enumerate}
\renewcommand{\theenumi}{\alph{enumi}}
\renewcommand{\labelenumi}{\theenumi)}
\item This gives the proof, which had been left pending, of
\hyperrefext[VIII][VIII.thm9.4]{9.4}~d), as a special case of
\hyperref[XV.cor4.3]{4.3}. % label for d) ?
Another, easier proof, not using the rather delicate result
\pageoriginale of \EGA IV 18.9.8, is obtained by observing that, by
\hyperrefext[VIII][VIII.prop9.1]{9.1}, we may confine ourselves to the
case of the morphism $\Spec(K)\rightarrow \Spec(k)$ induced by a field
extension $K/k$. But $\Spec(K)$ is universally locally $0$-acyclic over
$\Spec(k)$, since it is even universally acyclic for
$\LL=\PP-\{p\}$, $p=\car k$, as we shall see in
\hyperrefext[XVI][XVI.cor1.5]{1.5}.

\item It is not known whether every quasi-compact and quasi-separated,
surjective, locally $(-1)$-acyclic (or universally locally
$(-1)$-acyclic) morphism is a morphism of effective descent for the
fibered category of étale sheaves on varying preschemes. This statement
seems quite plausible and would strengthen \hyperref[XV.cor4.3]{4.3} and
\hyperrefext[VIII][VIII.thm9.4]{9.4}~c), d). % labels for c) and d) ?
Compare the effectivity statement in
\textsc{Murre}, Sém. Bourbaki ${\rm n^{\circ}}$~293, p.~17.

\item It is plausible that, under the hypotheses of
\hyperref[XV.cor4.2]{4.2}, $g$ is even locally acyclic for $\LL$ and
locally $1$-aspherical for $\LL$, where $\LL$ is the set of prime numbers
distinct from the residue characteristics of $Y$. For this question we
may clearly assume that $Y$ is the spectrum of an algebraically closed
field $k$, and it can be shown (\SGA 1964/65) that the answer is
affirmative whenever resolution of singularities is available for
schemes of finite type over $k$. Thus the answer is affirmative when
$Y$ has characteristic zero, as follows by using the results of
\textsc{Hironaka}~\footnote[1]{\Cf \SGA 1 XIII.}.
\end{enumerate}
\end{enonce*}


%XII 4.1, 4.2 equation ref.
%XII 1.2, 1.6 no ref. at 501

\ifx\danslelivre\undefined
\end{document}
\fi

% END INCLUDED SGA4_English_current_buildable_source_20260906/source/15/drafts/SGA4_XV_LINES_1316_1417_English.texfrag


% END INCLUDED SGA4_English_current_buildable_source_20260906/source/15/15_English.tex


\makeatletter
\def\input@path{{../16/}}
\makeatother
% BEGIN INCLUDED SGA4_English_current_buildable_source_20260906/source/16/16_English.tex
\ifx\danslelivre\undefined
% Release engine: XeLaTeX (the LuaTeX compatibility layer breaks Xy-pic output).
\documentclass[10pt]{../fichiers-configuration-latex/smfart}
% BEGIN INCLUDED SGA4_English_current_buildable_source_20260906/source/fichiers-configuration-latex/paquets.tex
\usepackage{polyglossia}
%\usepackage[francais]{babel}
\usepackage{fontspec}
\usepackage[math-style=ISO,bold-style=upright,sans-style=upright]{unicode-math}
%\usepackage[mathletters]{ucs}
%\usepackage[utf8]{inputenc}
%\usepackage[T1]{fontenc}
\usepackage[unicode,hypertexnames=false,colorlinks=true,linkcolor=blue,citecolor=Sepia,urlcolor=magenta]{hyperref}
\setmainlanguage{french}
\setmainfont[Ligatures=TeX]{Linux Libertine O}
% variante possible : TeXGyrePagella ?
\setmathfont{XITS Math}

% maths
\usepackage{amsmath}
%\usepackage{amssymb}
%\usepackage{mathrsfs}
%\usepackage{stmaryrd}
%\usepackage{euscript}

% couleurs
\usepackage[usenames,dvipsnames]{xcolor}

% graphiques
\usepackage{graphicx}
\usepackage[all]{xy}
%\usepackage{tikz}
%\usetikzlibrary{matrix,arrows,calc}
%\tikzset{column sep=3em,row sep=2em,text height=1.5ex, text depth=0.25ex,>=angle 90}
%\tikzset{description/.style={fill=white,inner sep=1pt}}

% hyperliens, 'etiquettes
\usepackage{xr-hyper}
\usepackage{../fichiers-configuration-latex/smfthm}
%\ifx\cacheretiquettes\undefined
%\usepackage[inline]{showlabels} \else \fi % pour rendre visibles les 'etiquettes
%\usepackage{../fichiers-configuration-latex/makeidx} % pour faire un index
%\usepackage[pagebackref,unicode]{hyperref}  % pour voir o`u sont cit'ees les r'ef'erences
% hyperref is loaded once above with the cumulative-reader destination policy.
%\hypersetup{breaklinks=true} %permet le retour `a la ligne dans les liens trop longs. Incompatible avec pdflatex
%\renewcommand*{\backref}[1]{{\small [cit\'e page(s)~#1.]}}

\usepackage{xspace}
\usepackage{pifont} % pour expose XVII

% END INCLUDED SGA4_English_current_buildable_source_20260906/source/fichiers-configuration-latex/paquets.tex

% BEGIN INCLUDED SGA4_English_current_buildable_source_20260906/source/fichiers-configuration-latex/numerotation.tex
\renewcommand{\thefootnote}{\roman{footnote}} % pour mettre les notes en bas de page sous la forme (ii) et non (2), qui ressemble trop (?) a un exposant.

\newcounter{pageoriginal}
\newcommand{\pageoriginale}{\refstepcounter{pageoriginal}\marginpar{\footnotesize\textbf{\thepageoriginal}}\xspace}
\let\pageoriginaled\pageoriginale

\usepackage{ifthen}
\newboolean{original}
\def\sisi{\ifthenelse{\boolean{original}}}
\setboolean{original}{false}

\newcounter{cnote}
\def\nde#1{\ifthenelse{\boolean{original}}{}{\stepcounter{cnote}\def\thefootnote{\arabic{cnote}}\protect\footnote[\value{cnote}]{N.D.E.: #1}}}
\def\ndemark{\ifthenelse{\boolean{original}}{}{\stepcounter{cnote}\def\thefootnote{\arabic{cnote}}\protect\footnotemark[\value{cnote}]}}
\def\ndetext#1{\ifthenelse{\boolean{original}}{}{\def\thefootnote{\arabic{cnote}}\protect\footnotetext[\value{cnote}]{N.D.E.: #1}}}

% END INCLUDED SGA4_English_current_buildable_source_20260906/source/fichiers-configuration-latex/numerotation.tex

% BEGIN INCLUDED SGA4_English_current_buildable_source_20260906/source/fichiers-configuration-latex/macros.tex
\setmainlanguage{french}
\setotherlanguage{english}


%% Fixed, engine-independent release identifier.  The inherited dynamic Git
%% probe used Lua io.popen, exposed build-state-dependent output, and printed
%% a visible question mark outside a Git checkout.
\def\showgitstatus{English release r2, 2026-07-23}


% macros pour LaTeX
\renewcommand{\mathsf}{}

% Police différente pour U+2113 SCRIPT SMALL L
\setmathfont[range="2113]{latinmodern-math.otf}

% Give every inherited semantic label its own stable PDF destination instead
% of inheriting whichever automatic equation/page anchor happened to be live.
% The legacy label remains as a compatibility alias used by existing edges.
\makeatletter
\def\etiquette[#1]#2{%
  \Hy@raisedlink{\hyper@anchorstart{sga4.semantic.#2}\hyper@anchorend}%
  \def\@currentHref{sga4.semantic.#2}%
  \def\@currentlabel{#1}%
  \label{#2}%
}
\makeatother
\def\eqhyperref[#1]#2{(\hyperref[#1]{#2})}
\def\pagehyperref[#1]#2{(page \hyperref[#1]{#2})}
\def\hyperrefext[#1][#2]#3{\hyperref[#2]{#1 #3}}
\def\eqhyperrefext[#1][#2]#3{\hyperref[#2]{#1 #3}}     %cf. exposes 17
\def\neweqhyperrefext[#1][#2]#3{\hyperref[#2]{#1 #3}}  %cf. exposes 17

\def\shat#1{\protect\widehat{#1}}
\def\stilde#1{\protect\widetilde{#1}}
%%%%%%%%%%%%%%%%%%


\def\etcn{etc.\xspace}
\def\etco{etc...\xspace}
\def\etc{\sisi{\etco}{\etcn}}
\def\ldotsn{\sisi{...}{\etcn}}

\def\cf{{cf.}\kern.3em}
\def\Cf{{Cf.}\kern.3em}
\def\ie{{i.e.}\kern.3em}
\def\iev{{i.e.},\xspace}
\def\loccito{loc.\kern2pt cit.\xspace}
\def\loccitn{\textit{loc.\kern2pt cit.}\xspace}
\def\loccit{\sisi{\loccito}{\loccitn}}
\def\EGA{\textup{EGA}\kern.2em}
\def\SGA{\textup{SGA}\kern.3em}

\def\quoi{}

\def\quav{\quad,\quad}
\def\quaet{\quad \text{et}\quad}
\def\fibreproduct#1#2#3{#1 ×_{#3} #2}
\let\fprod\fibreproduct
\def\ccheck#1#2{\mbox{${\check{\check{\rm #1}}}^{\lower2pt\hbox{$\scriptstyle #2$}}$}}
\def\sheaf#1{\protect\mathcal{#1}}
\def\vieuxsheaf#1{{#1}} % cf. problème avec sheaf{a\Phi_i} par exemple...
\DeclareMathOperator{\Hom}{Hom}
\DeclareMathOperator{\can}{can}
\DeclareMathOperator{\Ind}{Ind}
\DeclareMathOperator{\ind}{ind}
%\DeclareMathOperator{\Hhomtop}{\mathcal{H}\!\mathit{omtop}}
\DeclareMathOperator{\Hhomtop}{\mathbf{Homtop}}
\DeclareMathOperator{\Point}{Point}
\DeclareMathOperator{\Pro}{Pro}
\DeclareMathOperator{\Top}{Top}
\DeclareMathOperator{\Mod}{Mod}
\DeclareMathOperator{\Cart}{Cart}
%\DeclareMathOperator{\Mmorsite}{\mathcal{M}\!\mathit{orsite}}
\DeclareMathOperator{\Mmorsite}{\mathbf{Morsite}}
\DeclareMathOperator{\loc}{loc}
\DeclareMathOperator{\Loc}{Loc}
\DeclareMathOperator{\Pic}{Pic}
\DeclareMathOperator{\Et}{Et}
\DeclareMathOperator{\card}{card}
\DeclareMathOperator{\Ens}{Ens}
\DeclareMathOperator{\h}{\underline{H}}
\DeclareMathOperator{\Fonct}{Fonct}
\DeclareMathOperator{\Ker}{Ker}
\DeclareMathOperator{\but}{but}
\DeclareMathOperator{\Ob}{Ob}
\DeclareMathOperator{\ob}{ob}
\DeclareMathOperator{\Lib}{Lib}
\DeclareMathOperator{\esj}{esj}
\DeclareMathOperator{\id}{id}
\DeclareMathOperator{\Id}{Id}
\DeclareMathOperator{\ad}{ad}
\DeclareMathOperator{\Iim}{Im}
\DeclareMathOperator{\pr}{pr}
\DeclareMathOperator{\Cat}{Cat}
\DeclareMathOperator{\car}{car}
\DeclareMathOperator{\SysInd}{SysInd}
\DeclareMathOperator{\Sup}{Sup}
\DeclareMathOperator{\Coker}{Coker}
\DeclareMathOperator{\Filt}{Filt}
\DeclareMathOperator{\acc}{acc}
\DeclareMathOperator{\sur}{sur}
\DeclareMathOperator{\cat}{cat}
\DeclareMathOperator{\Ord}{Ord}
\DeclareMathOperator{\Gr}{Gr}
\DeclareMathOperator{\Cov}{Cov}
\DeclareMathOperator{\Caf}{Caf}
\DeclareMathOperator{\coker}{coker}
\DeclareMathOperator{\coim}{coim}
\DeclareMathOperator{\im}{im}
\DeclareMathOperator{\Ouv}{Ouv}
\DeclareMathOperator{\diag}{diag}
\DeclareMathOperator{\Esp}{Esp}
\DeclareMathOperator{\TOP}{TOP}
\DeclareMathOperator{\Unif}{Unif}
\DeclareMathOperator{\Iind}{\mathbf{Ind}}
\DeclareMathOperator{\Hhomcart}{\mathbf{Homcart}}
\DeclareMathOperator{\Hhomtopess}{\mathbf{Homtopess}}
\DeclareMathOperator{\Ppoints}{\mathbf{Points}}
\DeclareMathOperator{\Ppoint}{\mathbf{Point}}
\DeclareMathOperator{\Ttors}{\mathbf{Tors}}
\DeclareMathOperator{\Ssex}{\mathbf{Sex}}
\DeclareMathOperator{\Ffib}{\mathbf{Fib}}
\DeclareMathOperator{\Vvois}{\mathbf{Vois}}
\DeclareMathOperator{\Ppointess}{\mathbf{Pointess}}
\DeclareMathOperator{\scrHom}{\mathbf{Hom}}
\DeclareMathOperator{\sheafPic}{\mathbf{Pic}}
\DeclareMathOperator{\sheafAut}{\mathbf{Aut}}
\DeclareMathOperator{\Hhomtoplocan}{\mathbf{Homtoplocan}}
\DeclareMathOperator{\Hhomtopan}{\mathbf{Homtopan}}
\DeclareMathOperator{\Pptgeom}{\mathbf{Ptgéom}}
\DeclareMathOperator{\Morsite}{Morsite}
\DeclareMathOperator{\Fib}{Fib}
\DeclareMathOperator{\Fais}{Fais}
\DeclareMathOperator{\kar}{kar}
\DeclareMathOperator{\Groupes}{Groupes}
\DeclareMathOperator{\Restr}{Restr}
\DeclareMathOperator{\Isom}{Isom}
\DeclareMathOperator{\Prol}{Prol}
\DeclareMathOperator{\Res}{Res}
\DeclareMathOperator{\Tors}{Tors}
\DeclareMathOperator{\esp}{esp}
\DeclareMathOperator{\es}{es}
\DeclareMathOperator{\od}{od}
\DeclareMathOperator{\gr}{gr}
\DeclareMathOperator{\sob}{sob}
\DeclareMathOperator{\dfn}{dfn}
\DeclareMathOperator{\Sch}{Sch}
\DeclareMathOperator{\ouv}{ouv}
\DeclareMathOperator{\cons}{cons}
\DeclareMathOperator{\Nerf}{Nerf}
\DeclareMathOperator{\Points}{Points}
\DeclareMathOperator{\Espsob}{Espsob}
\DeclareMathOperator{\Homtop}{Homtop}
\DeclareMathOperator{\Pt}{Pt}
\DeclareMathOperator{\Quot}{Quot}
\DeclareMathOperator{\glob}{glob}
\DeclareMathOperator{\Homfib}{Homfib}
\DeclareMathOperator{\Qouv}{Qouv}
\DeclareMathOperator{\Quotouv}{Quotouv}
\DeclareMathOperator{\ens}{ens}
\DeclareMathOperator{\ab}{ab}
\DeclareMathOperator{\an}{an}
\DeclareMathOperator{\fs}{fs}
\DeclareMathOperator{\red}{red}
\DeclareMathOperator{\Ab}{Ab}
\DeclareMathOperator{\cosq}{cosq}
\DeclareMathOperator{\Ann}{Ann}
\DeclareMathOperator{\Zar}{Zar}


%%%%%%%%%%%%
%%%%% SGA4-2
\DeclareMathOperator{\Sex}{Sex}
\DeclareMathOperator{\vois}{vois}
\DeclareMathOperator{\Vois}{Vois}
\DeclareMathOperator{\Ext}{Ext}
\DeclareMathOperator{\ext}{ext}
\DeclareMathOperator{\Fer}{Fer}
\DeclareMathOperator{\Add}{Add}
\DeclareMathOperator{\Cocart}{Cocart}
\DeclareMathOperator{\cosk}{cosk}
\DeclareMathOperator{\plat}{plat}

\DeclareMathOperator{\faisc}{faisc}
\DeclareMathOperator{\sk}{sk}
\DeclareMathOperator{\ch}{ch}

\def\HH{\mathbf{H}}
\def\qc{\mathrm{qc}}
\def\et{\mathrm{et}}
\def\coh{\mathrm{coh}}

\def\Hvr#1{%
\begin{array}[b]{l}
{\kern2pt\scriptstyle r}\\[-8pt]
{\kern1pt\scriptstyle \vee}\\[-5pt]
H^{#1}
\end{array}}

\def\Hvv#1{%
\begin{array}[b]{l}
\kern2pt\rotatebox{90}{${\scriptscriptstyle\ll}$}\\[-6pt]
H^{#1}
\end{array}}

\def\Hvi#1{%
\begin{array}[b]{l}
{\kern0.2pt\scriptscriptstyle \infty}\\[-8pt]
{\kern1pt\scriptstyle \vee}\\[-5pt]
H^{#1}
\end{array}}

\def\Hv#1{%
\begin{array}[b]{l}
{\kern1pt\scriptstyle \vee}\\[-5pt]
H^{#1}
\end{array}}

\DeclareMathOperator{\cart}{cart}
\DeclareMathOperator{\Ffaisc}{\mathbf{Faisc}}


\def\Ev#1{%
\begin{array}[b]{l}
{\kern1pt\scriptstyle \vee}\\[-5pt]
E'_{#1}
\end{array}}

\DeclareMathOperator{\Cont}{Cont}
\DeclareMathOperator{\Ccont}{\mathscr{C}\!\mathit{ont}}
\DeclareMathOperator{\Site}{Site}
\newcommand\MM{\mathbf{M}}
\def\limtop#1{{\displaystyle{\varprojlim_{#1}\!\mbox{\rm top }}}}
\def\Limtop{\mathop{\varprojlim\!\mbox{\rm top}}}
\DeclareMathOperator{\spec}{spec}
\def\ss{\scriptstyle{\mathscr{S}}}
\DeclareMathOperator{\Gal}{Gal}
\def\smallm{\scriptstyle{\mathscr{M}}}
\DeclareMathOperator{\End}{End}
\DeclareMathOperator{\ass}{ass}
\DeclareMathOperator{\zar}{zar}
\DeclareMathOperator{\Ogg}{Ogg}
\DeclareMathOperator{\res}{res}
\DeclareMathOperator{\tr}{tr}
\DeclareMathOperator{\Supp}{Supp}
\DeclareMathOperator{\codim}{codim}
\DeclareMathOperator{\rk}{rk}
\DeclareMathOperator{\cur}{cur}
\DeclareMathOperator{\cdqc}{cdqc}
\newcommand\DD{\mathbf{D}}
\DeclareMathOperator{\cl}{cl}
\def\JJ{\mathbf{J}}
\def\FF{\mathbf{F}}
\def\BB{\mathbf{B}}
\DeclareMathOperator{\Alget}{Alget}
\DeclareMathOperator{\Ass}{Ass}
\DeclareMathOperator{\supp}{supp}
\DeclareMathOperator{\rad}{rad}
\DeclareMathOperator{\prof}{prof}
\DeclareMathOperator{\cint}{int}
\DeclareMathOperator{\sic}{sic}
\DeclareMathOperator{\Tor}{Tor}
\DeclareMathOperator{\tors}{tors}
\DeclareMathOperator{\Sym}{Sym}
\DeclareMathOperator{\parf}{parf}
\DeclareMathOperator{\tor}{tor}
\DeclareMathOperator{\defn}{defn}
\DeclareMathOperator{\Div}{Div}
\DeclareMathOperator{\Pol}{Pol}
\DeclareMathOperator{\Fex}{Fex}
\DeclareMathOperator{\Lie}{Lie}
\DeclareMathOperator{\tot}{tot}
\DeclareMathOperator{\Aut}{Aut}
\DeclareMathOperator{\HOM}{HOM}
\DeclareMathOperator{\pch}{pch}
\DeclareMathOperator{\Ch}{Ch}
\DeclareMathOperator{\EXT}{EXT}
\DeclareMathOperator{\PIC}{PIC}
\DeclareMathOperator{\Rf}{Rf}
\DeclareMathOperator{\inj}{inj}
\newcommand\VV{\mathbf{V}}

\def\Xv{%
\begin{array}[b]{l}
{\kern1pt\scriptstyle \vee}\\[-5pt]
X
\end{array}}

%%%%%%%
\newcommand\TT{\mathbf{T}}
\newcommand\XX{\mathbf{X}}
\newcommand\EE{\mathbf{E}}
%%%%%%

\def\qedsymbol{\textup{C.Q.F.D.}}
\let\CQFD\qed
\let\cqfd\qed

%%%%%\bigdot
\font\bigsymb=cmsy10 at 4pt
\def\bigdot{{\kern1.2pt\raise 1.5pt\hbox{\bigsymb\char15}}}
%%%%%

\def\Lright{\mathop{\vcenter{\hbox{\hbox{$\longrightarrow$}\lower8pt \hbox{\kern-16pt\hbox{$\longrightarrow$}}}}}}

\def\Aa{{\mathcal A}}
\def\Bb{{\mathcal B}}
\def\Cc{{\mathcal C}}
\def\Dd{{\mathcal D}}
\def\Ee{{\mathcal E}}
\def\Ff{{\mathcal F}}
\def\Gg{{\mathcal G}}
\def\Hh{{\mathcal H}}
\def\Ii{{\mathcal I}}
\def\Jj{{\mathcal J}}
\def\Kk{{\mathcal K}}
\def\Ll{{\mathcal L}}
\def\Mm{{\mathcal M}}
\def\Nn{{\mathcal N}}
\def\Oo{{\mathcal O}}
\def\Pp{{\mathcal P}}
\def\Qq{{\mathcal Q}}
\def\Rr{{\mathcal R}}
\def\Ss{{\mathcal S}}
\def\Tt{{\mathcal T}}
\def\Uu{{\mathcal U}}
\def\Vv{{\mathcal V}}
\def\Ww{{\mathcal W}}
\def\Xx{{\mathcal X}}
\def\Yy{{\mathcal Y}}
\def\Zz{{\mathcal Z}}

\newcommand{\CC}{\mathbf{C}}
\newcommand{\GG}{\mathbf{G}}
\newcommand{\KK}{\mathbf{K}}
\newcommand{\LL}{\mathbf{L}}
\newcommand{\NN}{\mathbf{N}}
\newcommand{\PP}{\mathbf{P}}
\newcommand\QQ{\mathbf{Q}}
\newcommand\RR{\mathbf{R}}
\newcommand\ZZ{\mathbf{Z}}

\def\to{\mathchoice{\longrightarrow}{\rightarrow}{\rightarrow}{\rightarrow}}
\def\from{\mathchoice{\longleftarrow}{\leftarrow}{\leftarrow}{\leftarrow}}
\def\mto{\mathchoice{\longmapsto}{\mapsto}{\mapsto}{\mapsto}}
\def\mfrom{\mathrel{\from\joinrel\mapstochar\,}}
\def\hto{\mathrel{\lhook\joinrel\to}}
\def\hfrom{\mathrel{\from\joinrel\rhook}}

\def\implies{\mathchoice{\Longrightarrow}{\Rightarrow}{\Rightarrow}{\Rightarrow}}
\def\implied{\mathchoice{\Longleftarrow}{\Leftarrow}{\Leftarrow}{\Leftarrow}}
\def\iff{\mathchoice{\Longleftrightarrow}{\Leftrightarrow}{\Leftrightarrow}{\Leftrightarrow}}

\def\To#1{\mathchoice{\xrightarrow{\labelstyle~~#1~}}{\stackrel{#1}{\longrightarrow}}{}{}}
\def\Mto#1{\mathrel{\mapstochar\To{#1}}}
\def\From#1{\mathchoice{\xleftarrow{\labelstyle~~#1~}}{\stackrel{#1}{\longleftarrow}}{}{}}

\newcommand{\isomto}{\overset{\sim}{\longrightarrow}}
\newcommand{\isomfrom}{\overset{\sim}{\longleftarrow}}

\newcommand{\IFF}{\unskip~$\iff$~\ignorespaces}
\newcommand{\THEN}{\unskip~$\implies$~\ignorespaces}

\let\leq\leqslant
\let\geq\geqslant
\let\le\leqslant
\let\ge\geqslant

\def\cf{{cf.}\kern.3em}
\def\Cf{{Cf.}\kern.3em}
\def\ie{{i.e.}\kern.3em}
\def\iev{{i.e.},\xspace}
%\def\loccito{loc.\kern2pt cit.\xspace}
%\def\loccitn{\textit{loc.\kern2pt cit.}\xspace}
%\def\loccit{\sisi{\loccito}{\loccitn}}
\def\EGA{\textup{EGA}\kern.2em}
\def\SGA{\textup{SGA}\kern.3em}

\let\oldS\S
\def\S{\oldS\kern .15em}
\def\ptbl{.\kern .2em}
\def\resp{resp.\kern.3em}
\def\Exp{Exp.\kern.3em}

\let\bar\overline
\let\hat\widehat
\let\tilde\widetilde
\let\dpl\displaystyle
\let\ldots\dots
\let\emptyset\varnothing
\let\skpt\skippointrait
\let\epsilon\varepsilon

\def\sfrac#1#2{{#1}/{#2}}

\DeclareMathOperator{\cd}{cd}
%\DeclareMathOperator{\Hhom}{\mathcal{H}\!\mathit{om}}
\DeclareMathOperator{\Hhom}{\mathbf{Hom}}
%\DeclareMathOperator{\Eext}{\mathcal{E}\!\mathit{xt}}
\DeclareMathOperator{\Eext}{\mathbf{Ext}}
\DeclareMathOperator{\Spec}{Spec}
\DeclareMathOperator{\Tr}{Tr}

\let\emptyset=\varnothing

%%%%%%%%%%%%%%%
%%% expose XVII
%%%%%%%%%%%%%%%

\makeatletter
\let\old@optexpose\optexpose
\newcommand{\empty@optexpose}[1]{}
\newcommand{\paren@optexpose}[1]{\old@optexpose{#1}(}
\newcommand{\refnoexpose}[1]{\let\optexpose\empty@optexpose\ref{#1}\let\optexpose\old@optexpose}%\refnoexpose{X.cor5.2} engendre "5.2" au lieu de "X 5.2"
\newcommand{\neweqref}[1]{\let\optexpose\paren@optexpose\textup{\ref{#1})}\let\optexpose\old@optexpose}%\neweqref{XII.eq4.1} engendre "XII (4.1)" au lieu de "(XII 4.1)"
\makeatother
\newcommand{\refvb}[2][VI${}_{\text{B}}$ ]{\sisi{#1\refnoexpose{#2}}{\ref{#2}}}

\newcommand{\leftexp}[2]{{\vphantom{#2}}^{#1}{#2}}
\newcommand{\AAA}{\mathbf{A}}
\newcommand{\Sf}{\mathfrak{S}}
\DeclareMathOperator*{\quotprojlim}{\text{``$\varprojlim$''}}
\DeclareMathOperator*{\quotinjlim}{\text{``$\varinjlim$''}}
\DeclareMathOperator{\Fl}{F\sisi{\ell}{l}}
\DeclareMathOperator{\Pl}{P\sisi{\ell}{l}}
%\DeclareMathOperator{\cl}{c\sisi{\ell}{l}}
\DeclareMathOperator{\tdf}{tdf}
\DeclareMathOperator{\rg}{rg}
\newcommand{\iInd}{\sisi{I}{i}nd}
\newcommand{\iind}{ind}
\newcommand{\ppro}{pro}
\newcommand{\tTor}{\sisi{T}{t}or}
\newcommand{\ttor}{tor}
\newcommand{\chap}{Chap.}
\newcommand{\Ffex}{\Ff\mathit{ex}}
\newcommand{\sheafk}{\underline{k}}
\newcommand{\sheafn}{\underline{n}}
\newcommand{\sheafGamma}{\underline{\Gamma}}

% Repeated references to an already typeset footnote must not use
% \footnotemark without a matching \footnotetext: Hyperref would then create
% a live annotation to a destination that cannot exist.  This helper keeps the
% current printed footnote mark and links it to an explicit stable note target.
\makeatletter
\newcommand{\sgarepeatedfootnotemark}[1]{%
  \begingroup
  \edef\@thefnmark{\thefootnote}%
  \hyperlink{#1}{\@makefnmark}%
  \endgroup}
\makeatother

\sisi{
\DeclareMathOperator{\simarrow}{\rule{.15mm}{1.8mm}\kern -3.4pt\rightsquigarrow}
\newcommand{\sheafell}{\underline{\ell}}
\newcommand{\sheafZZ}{\underline{\ZZ}}
}{
\newcommand{\nderef}[1]{la N.D.E.~(\ref{#1}) page~\pageref{#1}}
}

%%%%%%%%%%%%%%%%%
%%%%% expose XVII
%%%%%%%%%%%%%%%%%

% END INCLUDED SGA4_English_current_buildable_source_20260906/source/fichiers-configuration-latex/macros.tex


\externaldocument{../01/01_English.tex}
\externaldocument{../02/02_English.tex}
\externaldocument{../03/03_English.tex}
\externaldocument{../04/04_English.tex}
\externaldocument{../05/05_English.tex}
\externaldocument{../05bis/05bis_English.tex}
\externaldocument{../06/06_English.tex}
\externaldocument{../07/07_English.tex}
\externaldocument{../08/08_English.tex}
\externaldocument{../09/09_English.tex}
\externaldocument{../10/10_English.tex}
\externaldocument{../11/11_English.tex}
\externaldocument{../12/12_English.tex}
\externaldocument{../13/13_English.tex}
\externaldocument{../14/14_English.tex}
\externaldocument{../15/15_English.tex}

\setboolean{original}{false}

\begin{document}

\begin{center}
XVI. Smooth base change theorem and applications\\
M. Artin
\end{center}

\selectlanguage{english}version: \showgitstatus
\tableofcontents

\else
\chapter{Smooth base change theorem and applications}
%\addtocontents{toc}{M. Artin}
\begin{center}
M. Artin
\end{center}
\fi
\phantomsection\MakeLinkTarget*{sga4.exp.XVI}\label{sga4.exp.XVI}

% BEGIN INCLUDED SGA4_English_current_buildable_source_20260906/source/16/drafts/SGA4_XVI_LINES_43_274_English.texfrag
\setcounter{pageoriginal}{206}

\section{The smooth base change theorem}\etiquette[1]{XVI.sec1}\pageoriginale

\begin{enonce*}{Theorem 1.1}\etiquette[1.1]{XVI.thm1.1}
Let $\LL\subset \PP$, let $n\in \NN$, and let
$$\vcenter{
\xymatrix@C=1.5cm@R=1cm{
X   \ar[d]^{f}& X' \ar[l]_{g'} \ar[d]^{f'}\\
Y     & \ar[l]_{g} Y'
}}\leqno{(*)}
$$
be a cartesian diagram, with $f$ quasi-compact and quasi-separated\nde{That is,
coherent in the terminology of the source.}, and $g$
universally locally $0$-acyclic (\resp
\sisi{univ. loc.}{universally locally} $1$-aspherical for $\LL$, \resp
\sisi{univ. loc.}{universally locally} $n$-acyclic for $\LL$) \emph{(\hyperrefext[XV][XV.def1.11]{1.11}).} Then,
for every sheaf $F$ of sets (\resp of
\emph{ind}-$\LL$-groups, \resp an abelian sheaf of
$\LL$-torsion) on $X$, the base-change morphism
\emph{(XII \hyperref[XII.eq4.2]{4.1.2})}
$$
\varphi^{q}:g^{\ast}(R^{q}f_{\ast})F\longrightarrow
(R^{q}f'_{\ast}){g'}^{\ast}F
$$
is bijective for $q=0$ (\resp $q\leq 1$, \resp $q\leq n$).
\end{enonce*}

In particular, applying \hyperrefext[XV][XV.thm2.1]{2.1}, we obtain the following.

\begin{enonce*}{Corollary 1.2}\etiquette[1.2]{XVI.cor1.2}
(Smooth base change theorem). Suppose that
  the morphism $g$ in $(\ast)$ is smooth, that $f$ is quasi-compact
  and quasi-separated, and that $\LL\subset \PP$ is
  the complement of the set of residue characteristics
  of $Y$. Then, for every sheaf $F$ of sets
  (\resp of \emph{ind}-$\LL$-groups, \resp an abelian sheaf of
  $\LL$-torsion) on $X$, the base-
  change\pageoriginale morphism $\varphi^{q}$ above is
  bijective for $q=0$ (\resp for $q=0$, $1$, \resp for every $q$).
\end{enonce*}

\noindent
\emph{Proof of} \hyperref[XVI.thm1.1]{1.1}. We first treat the case in which
$f:X\rightarrow Y$ is quasi-projective: this is simply a combination of the
proper base change theorem \hyperrefext[XII][XII.thm5.1]{5.1} and
Definition \hyperrefext[XV][XV.prop1.9]{1.9}. Indeed, we have a commutative diagram
$$
\xymatrix@=1.25cm{
 X \ar[rr]^{i} \ar[rd]_{f}&        & \bar{X}\ar[ld]^{\bar{f}}\\
                         &   Y    &
}
$$
where $\bar{f}$ is projective, hence proper, and $i$ is an open
immersion. We know from \hyperrefext[XII][XII.thm5.1]{5.1} that the base change theorem holds
for $\bar{f}$, whatever the morphism $g:Y'\rightarrow Y$. Since
$\fprod{g}{\bar{X}}{Y}:\fprod{Y'}{\bar{X}}{Y}\rightarrow \bar{X}$ is
again universally $0$-acyclic (\resp\ldots), it follows immediately
from (\hyperrefext[XII][XII.prop4.4]{4.4} (ii)) that it suffices to prove the theorem for the
morphism $i$; that is, we are reduced to the case in which $f$ is an
open immersion. The theorem then follows from Definition
\hyperrefext[XV][XV.prop1.10]{1.10} (ii).

We now treat the general case. The assertion is local on $Y$, and
we may therefore suppose that $Y$ is affine.

\begin{enonce*}{Lemma 1.3}\etiquette[1.3]{XVI.lem1.3}
We may moreover suppose that $X$ is affine.
\end{enonce*}

\noindent
\emph{Proof}. Arguing as in \hyperrefext[XII][XII.lem6.1]{6.1}, in order to prove
\hyperref[XVI.thm1.1]{1.1} for a given $f$, we are reduced to proving it in the situation
deduced\pageoriginale from the given situation by the base changes
$\widetilde{Y}\rightarrow Y$, $\widetilde{Y}'\rightarrow Y'$
for a local morphism $\widetilde{g}:\widetilde{Y}'\rightarrow
\widetilde{Y}$ between strict localizations of $Y'$ and $Y$ induced by
$g$, and to proving that the homomorphisms
\begin{equation*}
H^{q}(\widetilde{X},F)\longrightarrow H^{q}(\widetilde{X}',F')\quad
(\mbox{\rm where } F'={g'}^{\ast}(F))\tag{$\ast$}
\end{equation*}
are bijective for $q\leq n$, where $\widetilde{X}=\fprod{X}{\widetilde{Y}}{Y}$,
$\widetilde{X}'=\fprod{X'}{\widetilde{Y}'}{Y'}$. We are thus reduced to the case in which
$Y$, $Y'$, and $g$ are strictly local, and to proving in this case the
bijectivity of the preceding maps, assuming that
\hyperref[XVI.thm1.1]{1.1} has been proved
for affine $f$. Let $(X_i)$ be a finite covering of $X$ by affine open
subschemes $X_i$, and let $Z$ be the disjoint union of the $X_i$, which is
therefore affine and is equipped with a surjective morphism $h:Z\rightarrow X$, whence
a morphism $h':Z'\rightarrow X'$. Note that $h$ is separated, and
that it is affine when $f$ is separated. Applying the descent lemma
\hyperrefext[XII][XII.prop6.8]{6.8}, we see that, in order to prove the bijectivity of the
maps $(\ast)$ for $q\leq n$ and every $F$, it suffices to
prove the bijectivity of the corresponding maps
$H^{q}(Z,h^{\ast}(F))\rightarrow H^{q}(Z',{h'}^{\ast}(F'))$,
provided that the base-change homomorphisms for
$Z\rightarrow X$ and the base change $X'\rightarrow X$ are
isomorphisms in dimensions $\leq n$. If $f$ is separated, then
$h$ is affine, and this is so by hypothesis; this proves \hyperref[XVI.thm1.1]{1.1}
when $f$ is assumed separated. In the general case, we may then
apply the preceding result to $h$, which is always separated, and again
conclude that \hyperref[XVI.thm1.1]{1.1} holds for $f$, proving \hyperref[XVI.lem1.3]{1.3}.

\begin{enonce*}[remark]{Remark}
One could also invoke the Leray spectral sequence for the
open covering $(X_i)$ of $X$, and its noncommutative
variants,\pageoriginale but this method, which does not differ
essentially from using \hyperrefext[XII][XII.prop6.8]{6.8}, has the disadvantage of
forcing us once again to distinguish the usual three cases, work
that has already been done in \loccit
\end{enonce*}

To complete the proof of the theorem in the general case, it
suffices to reduce the proof to the case in which $f:X\rightarrow Y$ is
affine and of finite type, hence quasi-projective. By \hyperref[XVI.lem1.3]{1.3}, we may
suppose that $X$ is affine. Then let $X=\varprojlim X_{\alpha}$, where the
morphisms $f_{\alpha}:X_{\alpha}\rightarrow Y$ are affine and of finite
type. The usual limit argument
(\hyperrefext[VII][VII.sec5]{5}) completes the proof.

\begin{enonce*}[remark]{Complements 1.4}\etiquette[1.4]{XVI.compl1.4}
\begin{enumerate}
\renewcommand{\theenumi}{\alph{enumi}}
\renewcommand{\labelenumi}{\theenumi)}
\item We leave it to the reader to check that the same
  proof still yields a conclusion when $g$ is assumed
  universally locally $(-1)$-acyclic (\hyperrefext[XV][XV.rems1.14]{1.14}): in this
  case, the base-change homomorphism
$$
\phi^{0}:g^{\ast}(f_{\ast}(F))\longrightarrow f'_{\ast}({g'}^{\ast}(F))
$$
is \emph{injective} for every sheaf of sets $F$ on
$X$. On the other hand, suppose again that $n\geq 0$. One may supplement
\hyperref[XVI.thm1.1]{1.1} with the following injectivity statement, which is also contained in the
preceding proof: if $g$ is universally locally
$0$-acyclic, then for every sheaf of groups $F$ on $X$,
the base-change homomorphism $\phi^{1}$ of \hyperref[XVI.thm1.1]{1.1} is a
monomorphism; if $g$ is universally locally $n$-acyclic
for $\LL$, then the homomorphism $\phi^{n+1}$ of \hyperref[XVI.thm1.1]{1.1} is a
monomorphism for every abelian sheaf of $\LL$-torsion $F$
on $X$.

\item It\pageoriginale is probably also possible to obtain an
  analogous injectivity conclusion from the hypothesis that $g$ is
  locally $1$-aspherical for $\LL$, by introducing the
  nonabelian $2$-cohomology invariants from \textsc{Giraud}'s thesis;
  compare \hyperrefext[XII][XII.subsec5.11]{5.11} \footnote[1]{This is indeed what is
  established in J. Giraud's book.}; the same question also arises
  for (\hyperref[XVI.cor1.6]{1.6}) and (\hyperref[XVI.cor2.3]{2.3}) below.
\end{enumerate}
\end{enonce*}

%raghu
\begin{enonce*}{Corollary 1.5}\etiquette[1.5]{XVI.cor1.5}
Let $K/k$ be a field extension, and let
$\LL=\PP-\{p\}$, where $p= \car k$. Then $\Spec
(K)\rightarrow \Spec (k)$ is universally locally acyclic
for $\LL$ and universally locally $1$-aspherical for
$\LL$; if $k$ and $K$ are separably closed, then the
preceding morphism is also universally acyclic for
$\LL$ and universally $1$-aspherical for $\LL$.
\end{enonce*}

We immediately note the following consequence of (\hyperref[XVI.cor1.5]{1.5}).

\begin{enonce*}{Corollary 1.6}\etiquette[1.6]{XVI.cor1.6}
Let $K/k$ be an extension of separably closed fields, let $X$ be a scheme,
and let $F$ be a sheaf of \sisi{set}{sets} (\resp of ind-$\LL$-groups, \resp
of abelian groups of $\LL$-torsion) on $X$, where
$\LL=\PP-\{p\}$, $p=\car k$. Then the map
$$
H^{q}(X,F)\longrightarrow H^{q}(X_{K},F_{K})
$$
is bijective for $q=0$ (\resp for $q\leq 1$, \resp for every $q$).
\end{enonce*}

Indeed, this follows at once from the fact that $\Spec(K)\rightarrow
\Spec(k)$ is universally acyclic for $\LL$ and
universally $1$-aspherical for $\LL$, and from the
definitions.\sisi{---~}{~}Note that, by \hyperrefext[VIII][VIII.cor2.4]{2.4}, the preceding homomorphism
on the $H^{q}$ may also be interpreted as being identified\sisi{}{,}
with the base-change homomorphism
$$
g^{\ast}(R^{q}f_{\ast}(F))\longrightarrow
R^{q}{f'}_{\ast}({g'}^{\ast}(F)){\quoi},
$$
where\pageoriginale $f:X\rightarrow \Spec(k)$, $f':X_K\rightarrow
\Spec(K)$, $g:\Spec(K)\rightarrow \Spec(k)$, and $g':X_K\rightarrow X$
are the evident morphisms. Thus, when $X$ is quasi-compact and
quasi-separated, the bijectivity of these homomorphisms may also be
regarded, by \hyperref[XVI.thm1.1]{1.1}, as a consequence of the fact that
$\Spec(K)\rightarrow \Spec(k)$ is universally \emph{locally}
acyclic for $\LL$ and universally locally
$1$-aspherical for $\LL$. This latter fact therefore already implies
\hyperref[XVI.cor1.6]{1.6} for $X$ quasi-compact and quasi-separated, which plainly suffices
to imply that $\Spec (K)\rightarrow \Spec(k)$ is
universally acyclic for $\LL$ and universally
$1$-aspherical for $\LL$ (\hyperrefext[XV][XV.def1.7]{1.7}, \hyperrefext[XV][XV.prop1.6]{1.6} (i)).

This shows that, to prove \hyperref[XVI.cor1.5]{1.5}, it suffices to prove the
first assertion of \hyperref[XVI.cor1.5]{1.5}. After passing to the perfect closure of
$k$, which is permissible by \hyperrefext[VIII][VIII.thm1.1]{1.1},
we may then suppose that $k$ is \emph{perfect}. Every extension $K$ of $k$ is the filtered direct limit of its
finitely generated $k$-subalgebras $A$, and since $k$ is perfect, after
localizing such an $A$, we may suppose it smooth, so that $K$
appears as a filtered direct limit of smooth subalgebras. Since the
latter are universally locally acyclic for
$\LL$ and universally locally $1$-aspherical for
$\LL$ by \hyperrefext[XV][XV.thm2.1]{2.1}, the conclusion follows from \hyperrefext[XV][XV.prop1.13]{1.13} (ii),
\cqfd

\begin{enonce*}[remark]{Remark 1.7}\etiquette[1.7]{XVI.rem1.7}
The invariance theorem \hyperref[XVI.cor1.5]{1.5} should be compared with
\hyperrefext[XII][XII.cor5.4]{5.4}, where it was not necessary to suppose that the torsion sheaf
under consideration is prime to the residue characteristics, but where, on the other hand,
$X$ must be assumed \emph{proper} over $k$. Already in the case of
$H^{1}$ $(X,\ZZ/p\ZZ)$, where $X$\pageoriginale is an affine algebraic
curve (for example, the affine line) over an algebraically closed
field $k$ of characteristic $p>0$, the analogue of the preceding statements becomes false,
because of the phenomena of ``wild ramification\nde{That is, ``wild.''}'' at infinity,
which imply that, in such a case, the classification of principal étale covers
with group $\ZZ/p\ZZ$ is essentially ``continuous.''
\end{enonce*}

% END INCLUDED SGA4_English_current_buildable_source_20260906/source/16/drafts/SGA4_XVI_LINES_43_274_English.texfrag


% BEGIN INCLUDED SGA4_English_current_buildable_source_20260906/source/16/drafts/SGA4_XVI_LINES_275_440_English.texfrag
\section{Specialization theorem for cohomology groups}\etiquette[2]{XVI.sec2}

\begin{enonce*}{Theorem 2.1}\etiquette[2.1]{XVI.thm2.1}
Let $\LL\subset \PP$, let $n\in \NN$, and let
$f:X\rightarrow Y$ be a proper morphism of finite presentation
\footnote[1]{For a less restrictive hypothesis than properness that yields
  the same conclusions, \cf \SGA 5 II 3 and \SGA 1 XIII.} that is locally
$0$-acyclic (\resp locally $1$-aspherical for $\LL$, \resp locally
$n$-acyclic for $\LL$). Let $F$ be a sheaf of sets (\resp of
\emph{ind}-$\LL$-groups, \resp an abelian sheaf of $\LL$-torsion) that is
constructible and locally constant on $X$. Then the $R^{q}f_{\ast}F$ are
constructible and \emph{locally constant} for $q=0$ (\resp for $q=0$, $1$,
\resp for $q\leq n$), and for every geometric point $\bar{y}$ of $Y$ one has
$$
(R^{q}f_{\ast}F)_{\bar{y}}\simeq
H^{q}(X_{\bar{y}},F_{\bar{y}}){\quoi},
$$
for the same values of $q$.
\end{enonce*}

\begin{enonce*}[remark]{Remark}
Note that the last assertion is simply the proper base-change theorem
\hyperrefext[XII][XII.cor5.2]{5.2}.
\end{enonce*}

In view of \hyperrefext[XV][XV.thm2.1]{2.1}, we immediately obtain the
following corollary:\pageoriginale

\begin{enonce*}{Corollary 2.2}\etiquette[2.2]{XVI.cor2.2}
(Specialization theorem for cohomology groups). Let
$f:X\rightarrow Y$ be a proper smooth morphism, and let
$\LL\subset \PP$ be the complement of the set of residual characteristics
of $Y$. Let $F$ be a sheaf of sets (\resp of ind-$\LL$-groups, \resp an
abelian sheaf of $\LL$-torsion) that is constructible on $X$. Then the
$R^{q}f_{\ast}(F)$ are constructible and \emph{locally constant} for $q=0$
(\resp for $q=0$, $1$, \resp for every $q$), and for every geometric point
$\bar{y}$ of $Y$ one has
$$
R^{q}f_{\ast}(F)_{\bar{y}}\simeq H^{q}(X_{\bar{y}},F_{\bar{y}})
$$
for the same values of $q$.
\end{enonce*}

\begin{enonce*}[remark]{Remark}
Taking $n=1$ and translating into the language of fundamental groups, we
recover (\SGA 1 X 3.8), which compares the fundamental groups of the fibers
of a proper smooth morphism.
\end{enonce*}

\noindent
\emph{Proof of} \hyperref[XVI.thm2.1]{2.1}. Since we already know that the
$R^{q}f_{\ast}(F)$ are constructible \hyperrefext[XIV][XIV.thm1.1]{1.1}, it
suffices, by \hyperrefext[IX][IX.prop2.11]{2.11}, to prove that for every
specialization morphism $\bar{y}_1\rightarrow \bar{y}_{0}$ of geometric
points of $Y$, the corresponding specialization morphism
(\hyperrefext[VIII][VIII.subsec7.7]{7.7})
\begin{equation*}
R^{q}f_{\ast}(F)_{\bar{y}_{0}}\longrightarrow
R^{q}f_{\ast}(F)_{\bar{y}_1}\tag{$\ast$}
\end{equation*}
is bijective for the values of $q$ under consideration. For this we have the
following slightly more general result (in which it is unnecessary to
assume\pageoriginale that $f$ is of finite presentation or that $F$ is
constructible):

\begin{enonce*}{Corollary 2.3}\etiquette[2.3]{XVI.cor2.3}
Let $\LL\subset \PP$, let $n\in \NN$, let
$f:X\rightarrow Y$ be a proper morphism that is locally $0$-acyclic (\resp
locally $1$-aspherical for $\LL$, \resp locally $n$-acyclic for $\LL$), and
let $F$ be a sheaf of sets (\resp of ind-$\LL$-groups, \resp an abelian
sheaf of $\LL$-torsion) on $X$ that is locally constant. Then for every
specialization morphism $\bar{y}_1\rightarrow \bar{y}_0$ of geometric points
of $Y$, the corresponding specialization homomorphism $(\ast)$ above is
bijective for $q=0$ (\resp for $q=0$, $1$, \resp for $q\leq n$). Moreover,
when $n=0$, \ie when $f$ is locally $0$-acyclic, for every locally constant
sheaf of groups $F$ on $X$, the homomorphism
$R^{1}f_{\ast}(F)_{\bar{y}_0}\rightarrow R^{1}f_{\ast}(F)_{\bar{y}_1}$ is
injective; for arbitrary $n$, and for every locally constant abelian sheaf
of $\LL$-torsion $F$ on $X$, the specialization homomorphism $(\ast)$ is
injective for $q=n+1$.
\end{enonce*}

There exists an integral strictly local scheme $Y'$ and a morphism
$g:Y'\rightarrow Y$ that maps the closed point $y'_0$ to $y_{0}$ and the
generic point $y'_1$ to $y_1$: for example, it is enough first to take the
strict localization of $Y$ at $\bar{y}_0$, and then its integral closed
subscheme defined by a point lying above $y_1$. Moreover, after replacing
$Y'$ by its normalization in an algebraic closure of its function field, we
may assume that $Y'$ is normal and that $k(y'_1)$ is algebraically closed,
so that $y'_0$ and $y'_1$ are geometric points of $Y'$. Using the
base-change theorem (\hyperrefext[XII][XII.thm5.1]{5.1}) for $(f,F,g)$, the
proof of \hyperref[XVI.cor2.3]{2.3} is reduced, after replacing\pageoriginale
$Y$, $\bar{y}_0$, $\bar{y}_1$ by $Y'$, $y'_0$, $y'_1$, to proving \hyperref[XVI.cor2.3]{2.3} when
$Y$ is normal and strictly local and when $\bar{y}_0$ and $\bar{y}_1$ are,
respectively, its closed point $y_0$ and its generic point $y_1$. Moreover,
using \hyperrefext[XII][XII.sec4]{4}, we see that the specialization
homomorphism $(\ast)$ is then identified with the homomorphism
$$
H^{q}(X,F)\longrightarrow H^{q}(X_1,g^{\ast}_1(F)){\quoi},
$$
where $X_1=X_{y_1}$ and where $g:y_1\rightarrow Y$ and
$g_1:X_1\rightarrow X$ are the canonical morphisms. Our assertion now
follows from the following lemma (in which $f$ is no longer assumed proper):

\begin{enonce*}{Lemma 2.4}\etiquette[2.4]{XVI.lem2.4}
With the preliminary notation of $\hyperref[XVI.cor2.3]{2.3}$, drop the properness assumption on
$f$, but assume instead that $Y$ is normal and integral, with generic point
$y_1$ such that $k(y_1)$ is separably closed, and let
$X_1=X_{y_1}$ and $g_1:X_1\rightarrow X$ be the canonical morphism. Then the
homomorphism
\begin{equation*}
H^{q}(X,F)\longrightarrow H^{q}(X_1,g^{\ast}_1(F))\tag{${\ast}{\ast}$}
\end{equation*}
is bijective for $q=0$ (\resp bijective for $q=0$, $1$, \resp bijective for
$q\leq n$ and injective for $q=n+1$). Moreover, when $n=0$, \ie when $f$ is
locally $0$-acyclic, for every sheaf of groups $F$ on $X$, the homomorphism
$({\ast}{\ast})$ is injective for $q=1$.
\end{enonce*}

An immediate argument (\hyperrefext[XV][XV.prop1.6]{1.6} (i) $\Rightarrow$
(ii)) shows that it is enough to prove the relation
$$
F\xrightarrow{\sim} g_{1\ast}g_1{}^{\ast}(F){\quoi},
$$
and,\pageoriginale for $n\geq 1$, the additional relations
$$
R^{q}g_{1\ast}g_1{}^{\ast}(F)=0\quad\mbox{\rm for}\quad 1\leq
q\leq n{\quoi}.
$$
These relations are local on $X$ for the etale topology, and $F$ is locally
constant for the etale topology, so we may assume that $F$ is constant and
hence of the form $G_X$, where $G$ is a set (\resp an ind-$\LL$-group,
\resp an abelian group of $\LL$-torsion). Note also that the assumption that
$Y$ is normal implies
$g_{\ast}(G_{y_1})\xleftarrow{\sim}G_Y$
(\hyperrefext[IX][IX.lem2.14.1]{2.14.1}), and the assumption that $k(y_1)$
is separably closed implies that the $R^{q}g_{\ast}(G_{y_1})$ vanish for
$q\geq 1$. On the other hand, the assumption that $f$ is locally
$0$-acyclic (\resp\ldots) implies that the formation of the
$R^{q}g_{\ast}(G_{y_1})$ commutes with the base change $f:X\rightarrow Y$
for $q\leq n$. Here we use the fact that we may assume $Y$ affine and that
$y_1$ then appears as the inverse limit of its affine open neighborhoods,
which are quasi-compact in $Y$, so that Definition
\hyperrefext[XV][XV.def1.11]{1.11} and the theory of passage to the limit
\hyperrefext[VII][VII.sec5]{5} apply. We conclude that indeed
$$
G_X\xrightarrow{\sim} g_{1\ast}(G_{X_1})\quav
R^{q}g_{1\ast}(G_{X_{1}})=0\quad\mbox{\rm for}\quad 1\leq q\leq
n{\quoi},
$$
which completes the proof of \hyperref[XVI.lem2.4]{2.4}, and hence that of
\hyperref[XVI.thm2.1]{2.1}.

\begin{enonce*}[remark]{Remark}
The preceding proof, via \hyperref[XVI.lem2.4]{2.4}, which is simpler than
our original proof, is due to M. Lubkin.
\end{enonce*}

\begin{enonce*}{Corollary 2.5}\etiquette[2.5]{XVI.cor2.5}
Under the hypotheses of $\hyperref[XVI.thm2.1]{2.1}$, assume that $Y$ is connected, and let
$q\leq n$. Then the groups $H^{q}(X_{\bar{y}},F)$, for geometric
points\pageoriginale $\bar{y}$ of $Y$, are all isomorphic to one another.
\end{enonce*}

% END INCLUDED SGA4_English_current_buildable_source_20260906/source/16/drafts/SGA4_XVI_LINES_275_440_English.texfrag


% BEGIN INCLUDED SGA4_English_current_buildable_source_20260906/source/16/drafts/SGA4_XVI_LINES_441_700_English.texfrag
\section{The relative cohomological purity theorem}\etiquette[3]{XVI.sec3}

\begin{enonce*}{Definition 3.1}\etiquette[3.1]{XVI.def3.1}
Let $S$ be a scheme. A \emph{smooth $S$-pair} $(Y,X)$ is a closed
$S$-immersion $i:Y\rightarrow X$ of smooth $S$-preschemes, that is, a
commutative diagram of preschemes
$$
\xymatrix@=1.25cm{
Y \ar[rr]^{i} \ar[rd]_{h}&     &X\ar[ld]^{f}\\
  & S   &
}
$$
such that $f$ and $h$ are smooth and $i$ is a closed immersion. We
write $U=X-Y$, denote by $j:U\rightarrow X$ the open immersion, and
by $g:U\rightarrow S$ the structural morphism. The \emph{codimension}
of $(Y,X)$ at a point $y\in Y$ is the codimension at $y$ of the fiber
$Y_s$ in $X_s$, where $s=h(y)$.
\end{enonce*}

Thus the codimension is a locally constant function on $Y$. We shall
express the fact that this function is in fact constant,
\sisi{with}{of} value~$c$, by saying that ``$(Y,X)$ has
codimension~$c$.''

Let $(Y,X)$ be a smooth $S$-pair and let $y\in Y$. Recall (\SGA II
4.10) that there exist integers $m$ and $n$, a neighborhood $X'$ of
$y$ in $X$, and an étale morphism
$$
\phi:X'\longrightarrow \EE^n_S=\Spec \Oo_S[t_1,\ldots,t_n]
$$
such\pageoriginale that $Y'=Y\cap X'$ is the inverse image of the
closed subscheme $\EE^m_S$ defined by the equations
$$
t_{m+1}=\ldots=t_n=0{\quoi}.
$$
This result can be expressed by saying that \emph{locally for the
  étale topology, every smooth $S$-pair $(Y,X)$ is isomorphic to the
  standard pair $(\EE^m_S, \EE^n_S)$ for suitable $m$ and $n$,} where
$$
\EE^r_S=\Spec \Oo_S[t_1,\ldots,t_r]{\quoi}.
$$

In the present section, we propose to calculate the local cohomology
sheaves $\HH^q_Y(X,F)=(R^{q}i^{!})F$
(\sisi{\hyperrefext[V][V.sec4]{4}}{\hyperrefext[V][V.sec6]{6}} and
\hyperrefext[VIII][VIII.subsec6.6]{6.6}) for a smooth $S$-pair $(Y,X)$,
with coefficients in a locally constant abelian torsion sheaf $F$
prime to the residue characteristics. We begin with some preliminary
considerations:

\begin{enonce*}{Proposition 3.2}\etiquette[3.2]{XVI.prop3.2}
Let
$$
\xymatrix@=1.25cm{
U \ar[rr]^{j} \ar[rd]_{g}&     &X\ar[ld]^{f}\\
  & S   &
}
$$
be a commutative diagram in which $f$ is smooth and $j$ is an open
immersion such that the fiber $U_s$ is dense in $X_s$ for every
$s\in S$. Let $F$ be a sheaf of sets on $S$. Then the canonical
morphism
$$
f^{\ast}F\longrightarrow j_{\ast}g^{\ast}F=j_{\ast}j^{\ast}f^{\ast}F
$$
is bijective.
\end{enonce*}

\noindent
\emph{Proof}.\pageoriginale We readily reduce to the case where $U$ is
quasi-compact in $X$. Let $\bar{s}$ be a geometric point of $S$, let
$\bar{x}$ be a geometric point of $X$ lying over $\bar{s}$, let
$\widetilde{f}:\widetilde{X}\rightarrow \widetilde{S}$ be the morphism
of the strict localizations of $X$ and $S$ at the corresponding
geometric points, and let
$\widetilde{U}=\fprod{U}{\widetilde{X}}{X}$. With the evident notation,
we have
$$
(f^{\ast}F)_{\bar{x}}=H^{0}(\widetilde{X},\widetilde{f}^{\ast}\widetilde{F})\simeq
H^{0}(\widetilde{S},\widetilde{F}){\quoi}.
$$
Consequently, it suffices to prove that
$\widetilde{g}:\widetilde{U}\rightarrow \widetilde{S}$ is
$0$-acyclic, which will imply
$$
(j_{\ast}g^{\ast}F)_{\bar{x}}\simeq
H^{0}(\widetilde{U},\widetilde{g}^{\ast}\widetilde{F})\simeq
H^{0}(\widetilde{S},\widetilde{F}){\quoi}.
$$
Since $\widetilde{g}$ is a limit of smooth morphisms, it is locally
$0$-acyclic (\hyperrefext[XV][XV.thm2.1]{2.1},
\hyperrefext[XV][XV.def1.11]{1.11}), and by
\hyperrefext[XV][XV.cor1.12]{1.12} it suffices to prove that the
geometric fibers of $\widetilde{U}/\widetilde{S}$ are $0$-acyclic,
that is, connected and nonempty. Let $\bar{s}'$ be a geometric point
of $\widetilde{S}$. Then the fiber $\widetilde{X}_{\bar{s}'}$ is
regular, connected, and nonempty, and the hypothesis implies that
$\widetilde{U}_{\bar{s}'}$ is a dense open subset of
$\widetilde{X}_{\bar{s}'}$, hence connected and nonempty, which proves
the result.

\begin{enonce*}{Corollary 3.2.1}\etiquette[3.2.1]{XVI.cor3.2.1}
Under the hypotheses of \hyperref[XVI.prop3.2]{3.2}, the functor
$X'\mapsto X'|U$ from the category of étale coverings of $X$ to the
category of étale coverings of $U$ is fully faithful. In particular,
for every finite group $G$, the restriction map
$$
H^{1}(X,G)\longrightarrow H^{1}(U,G)
$$
is injective.
\end{enonce*}

This\pageoriginale follows from \hyperref[XVI.prop3.2]{3.2}, by a
standard argument which, for ease of reference, we shall spell out in
a

\begin{enonce*}{Lemma 3.2.2}\etiquette[3.2.2]{XVI.lem3.2.2}
Let $f:X\rightarrow Y$ be a morphism of schemes, and let
$f^{\ast}:\Et(Y)\rightarrow \Et(X)$ be the inverse-image functor
$Y'\mapsto \fprod{Y'}{X}{Y}$ from the category $\Et(Y)$ of étale
coverings of $Y$ to $\Et(X)$. The following conditions are equivalent:
\begin{itemize}
\item[(i)] $f^{\ast}$ is faithful (\resp fully faithful).

\item[(ii)] For every étale covering $Y'$ of $Y$, denoting by
  $f':X'\rightarrow Y'$ the morphism deduced from $f$ by the base
  change $Y'\rightarrow Y$, the map $U\mapsto {f'}^{-1}(U)$ from the
  set of clopen subsets of $Y'$ to the set of clopen subsets of $X'$
  is injective (\resp bijective).

\item[(ii bis)] With the notation of {\rm(ii)}, for every constant
  sheaf of sets $C$ on $Y'$, the canonical map
$$
H^{0}(Y',C)\longrightarrow H^{0}(X',{f'}^{\ast}C)
$$
is injective (\resp bijective).

\item[(iii)] With the notation of {\rm(ii)}, the canonical map
$$
\Gamma(Y'/Y)\longrightarrow \Gamma(X'/X)
$$
is injective (\resp bijective).
\end{itemize}
\end{enonce*}

The proof is left to the reader (\cf \SGA 2, IX 3.1 and 3.2).

\begin{enonce*}[remark]{Remarks 3.2.3}\etiquette[3.2.3]{XVI.rems3.2.3}
One\pageoriginale can generalize Proposition
\hyperref[XVI.prop3.2]{3.2} considerably, both for sheaves of sets, by
weakening the hypotheses on
$f$, and for sheaves of groups, with conclusions concerning the
$R^{q}g_{\ast}(f^{\ast}F)$ (\cf \SGA 2 XIV 1 and \SGA 1 XIII). The
analogous remark applies to \hyperref[XVI.cor3.4]{3.4} through
\hyperref[XVI.cor3.6]{3.6}.
\end{enonce*}

Let us also ``recall'' the Zariski--Nagata purity theorem in its
relative form:

\begin{enonce*}{Theorem 3.3}\etiquette[3.3]{XVI.thm3.3}
(Relative purity theorem). Let
$$
\xymatrix@=1.25cm{
U \ar[rr]^{j} \ar[rd]_{g}&     &X\ar[ld]^{f}\\
  & S   &
}
$$
be a commutative diagram in which $f$ is smooth and $j$ is an open
immersion such that, for every $s\in S$, $X_s-U_s$ has codimension
$\geq 2$ everywhere in $X_s$. Then the functor $X'\mapsto X'|U$ from
the category of étale coverings of $X$ to the category of étale
coverings of $U$ is an \emph{equivalence of categories.}
\end{enonce*}

\noindent
\emph{Proof}. We saw in \hyperref[XVI.cor3.2.1]{3.2.1} that the
functor is fully faithful. The fact that it is
\sisi{ess.}{essentially} surjective will follow from the classical
theorem. Given an étale covering $U'$ of $U$, we must
show that it extends to an étale covering $X'$ of $X$. There is an
immediate and easy reduction first to the case where $U$ is
quasi-compact in $X$, and then to the case where $S$ is affine and of
finite type over $\Spec \ZZ$, $X$ is connected, and $f$ is of finite
type.

We\pageoriginale proceed by induction on $n=\dim S$, the assertion
being a consequence of the classical theorem (\SGA 2 X 3.4) when
$n=0$. Since uniqueness is available, we can apply the descent
technique of \sisi{\SGA IX 4.7}{\SGA 1 IX 4.7} to the normalization
$S\leftarrow \bar{S}$ and thereby reduce to the case where $S$ (and
hence $X$) is normal ($\bar{S}\rightarrow S$ being finite (\EGA IV
7.8.3 (iii) (vi)). Moreover, after enlarging $U$, we may suppose that
$U$ is maximal among the open subsets over which $U'$ can be extended,
and we must then prove that $U=X$. Otherwise, let $x$ be a maximal
point of $X-U$, and put $s=f(x)$. It is enough to prove that $U'$ can
be extended over a neighborhood of $x$. We may suppose that
$\dim \Oo_{S,s}=n\geq 1$ and that the theorem has already been proved
in dimension $<n$. Put $B=\Oo_{X,x}$, and let $B'$ be the normalization
of $B$ in the ring of rational functions on the covering $U'$ of $U$;
thus $B'$ is a finite $B$-algebra (\loccit). It plainly suffices to
prove that $B'$ is an étale $B$-algebra.

Let $A=\Oo_{S,s}$, $0\neq t\in \max A$, $A_{0}=A/tA$, and
$B_{0}=B/tB$. By the induction hypothesis, the restriction of the
covering to $U_{0}=U\otimes_A A_{0}$ extends to an étale covering of
$\Spec B_{0}$, which we denote by $\Spec B^{\ast}_{0}$. After replacing
$X$ by a suitable étale covering of a neighborhood of $x$ (which is
permissible), we may suppose that $\Spec B^{\ast}_{0}$ is completely
split over $\Spec B_{0}$. We are thereby reduced to the case where the
$B_{0}$-algebra $B'_{0}$ is completely split away from the closed
point of $\Spec B_{0}$. The hypotheses of
\hyperrefext[XV][XV.lem3.3]{3.3} are then satisfied, which proves the
result.

\begin{enonce*}{Corollary 3.4}\etiquette[3.4]{XVI.cor3.4}
Under the same hypotheses as in \hyperref[XVI.thm3.3]{3.3}, we have
$(R^{1}j_{\ast})g^{\ast}F=0$ for every sheaf $F$ of ind-finite groups
on $S$.
\end{enonce*}

We reduce\pageoriginale as usual to the case where $S$ is noetherian,
$f$ is of finite type, and $F$ is constructible. Let
$0\rightarrow F\rightarrow G$ be an injection of sheaves of
ind-finite groups on $S$, and put $C=G/F$. It follows from
\hyperref[XVI.prop3.2]{3.2} that the functor $j_{\ast}g^{\ast}$ is
exact, so that $j_{\ast}g^{\ast}G\rightarrow j_{\ast}g^{\ast}C$ is
surjective, and hence (\hyperrefext[XII][XII.prop3.1]{3.1}) that
$(R^{1}j_{\ast})g^{\ast}F\rightarrow
(R^{1}j_{\ast})g^{\ast}G$ is injective. We can therefore replace $F$
by $G$, and thus reduce, by \hyperrefext[IX][IX.prop2.14]{2.14} and
\hyperrefext[VIII][VIII.prop5.5]{5.5}, to the case where $F=G$ is
constant, with value an ordinary finite group $G$.
%% This reduction is not perfectly clear in the nonabelian case:
%% give a reference for R^1 j_* p_*G = p_* R^1 j_*G
Let $\bar{x}$ be a geometric point of $X$, let $\widetilde{X}$ be the
strict localization of $X$ at $\bar{x}$, and let
$\widetilde{U}=\fprod{\widetilde{X}}{U}{X}$. It then follows from
\hyperref[XVI.thm3.3]{3.3} that
$$
(R^{q}j_{\ast}g^{\ast}F)_{\bar{x}}=H^{q}(\widetilde{U},G_{\widetilde{U}})\simeq
H^{q}(\widetilde{X},G_{\widetilde{X}}){\quoi},
$$
and the latter is zero, which proves the result.

\begin{enonce*}{Lemma 3.5}\etiquette[3.5]{XVI.lem3.5}
(Relative Abhyankar's lemma). Let $(Y,X)$ be a smooth $S$-pair of
  codimension $1$, such that $Y$ is defined in $X$ by an equation
  $t=0$. Let $U=X-Y$, and let $V/U$ be a principal Galois covering of
  order $n$ prime to the residue characteristics of $Y$. Put
  $X'=\Spec \Oo_X[z]/(z^n-t)$, $U'=\fprod{U}{X'}{X}$, and
  $V'=\fprod{V}{X'}{X}$. Then the covering $V'$ extends uniquely to an
  étale covering of $X'$.
\end{enonce*}

\noindent
\emph{Proof}. We immediately reduce to the case where $S$, $X$, and
$Y$ are affine and of finite type over $\Spec \ZZ$. Uniqueness is a
special case of \hyperref[XVI.cor3.2.1]{3.2.1}. We shall deduce
existence from \hyperrefext[XV][XV.cor1.12]{1.12}: let $\bar{s}$ be a
geometric point of $S$ and let $\bar{x}$ be a geometric point of $Y$
lying over $\bar{s}$. Let
$\widetilde{f}:\widetilde{X}\rightarrow \widetilde{S}$ be the morphism
of the corresponding strict localizations\pageoriginale induced by
$f$, let $\widetilde{U}=\fprod{U}{\widetilde{X}}{X}$, etc\ldots By
descent, in view of uniqueness, it plainly suffices to prove that
$\widetilde{V}'$ extends to an étale covering of $\widetilde{X}'$, that
is (since $\widetilde{X}'$ admits no nontrivial covering), that
$\widetilde{V}'$ is a trivial covering of $\widetilde{U}'$.

Let

% END INCLUDED SGA4_English_current_buildable_source_20260906/source/16/drafts/SGA4_XVI_LINES_441_700_English.texfrag


% BEGIN INCLUDED SGA4_English_current_buildable_source_20260906/source/16/drafts/SGA4_XVI_LINES_701_956_English.texfrag
$\bar{U}=\varprojlim_m\widetilde{U}_{m}=\Spec \Oo_{\widetilde{U}}[t^{1/m}]$,
where $m$ runs through the set of integers prime to the residue
characteristic of $\widetilde{X}$. Then
$\bar{U}/\widetilde{U}$ is an “infinite Galois covering” with abelian
group $G=\varprojlim_m \ZZ/m\ZZ$, and one sees immediately that it is
enough to prove that $\bar{U}$ is $1$-aspherical for the set
$\LL\subset \PP$ complementary to the residue characteristic of
$\widetilde{X}$.

Consider the morphism $\bar{f}:\bar{U}\rightarrow
\widetilde{S}$. It is a limit of smooth morphisms, and is therefore
locally $1$-aspherical for $\LL$ (\hyperrefext[XV][XV.thm2.1]{2.1}, \hyperrefext[XV][XV.def1.11]{1.11}
(ii)). By \hyperrefext[XV][XV.cor1.12]{1.12}, it is enough to prove that the geometric
fibres of $\bar{U}/\widetilde{S}$ are $1$-aspherical for $\LL$. Let
$\bar{s}$ be a geometric point of $\widetilde{S}$. The geometric fibre
$\bar{U}_{\bar{s}}$ is the limit of the fibres
$(\widetilde{U}_m)_{\bar{s}}$. Consequently, it is enough to prove that
every principal Galois covering of some
$(\widetilde{U}_m)_{\bar{s}}$, of order prime to the
\sisi{res.\ char.}{residue characteristics}, induces a trivial covering
of $\bar{U}_{\bar{s}}$. Replacing $\widetilde{X}$ by
$\widetilde{X}_m=\Spec \Oo_{\widetilde{X}}[t^{1/m}]$ (which is still
smooth over $\widetilde{S}$), we reduce to the case $m=1$, that is,
$\widetilde{U}_m=\widetilde{U}$. Let $V$ be such a covering of
$\widetilde{U}_{\bar{s}}$. Now $\widetilde{X}_{\bar{s}}=Z$ is a
regular strictly local scheme, and $W=\widetilde{U}_{\bar{s}}$ is the
open complement of a regular subset of codimension $1$. We may
therefore apply the usual form of Abhyankar's lemma
(\SGA 1 X 3.6) to conclude that $V$ induces a
covering\pageoriginale $V_m$ of $W_m$ which extends to an étale
covering of $Z_m$, for a suitable $m$. But since $Z$, and hence $Z_m$,
is strictly local, this covering is trivial; therefore $V_m$ is a
trivial covering of $W_m$, \cqfd

\begin{enonce*}[remark]{Remark 3.5.1}\etiquette[3.5.1]{XVI.rem3.5.1}
When the order $n$ of $G$ is not assumed prime to the residue
characteristics of $X$, one can still prove the following: let $m$ be
the largest integer prime to the residue characteristics of $Y$ which
divides the order of $G$, and suppose that for every point $s\in S$,
the covering $V_s$ of $U_s=X_s-Y_s$ has, at the maximal points of
$Y_s$, inertia groups whose orders are prime to the residue
characteristics of $Y$. Then the conclusion of
\hyperref[XVI.lem3.5]{3.5} remains valid on taking
$X'=\Spec \Oo_X[z]/(z^m-t)$. More generally, one may also take for $m$
any common multiple of the orders of the inertia groups just
considered, provided these orders are prime to the corresponding
residue characteristics.
\end{enonce*}

\begin{enonce*}{Corollary 3.6}\etiquette[3.6]{XVI.cor3.6}
(Cohomological purity in dimension $1$). Let $(Y,X)$ be a smooth
  $S$-pair of codimension $1$, and let $F$ be a constructible locally
  constant sheaf of groups on $X$, whose orders are prime to the
  residue characteristics of $X$. Then, with the notation of $\hyperref[XVI.def3.1]{3.1}$,
  $(R^{1}j_{\ast})j^{\ast}F$ is locally isomorphic, as a sheaf of
  pointed sets, to $i_{\ast}(F_Y/\cint(F_Y))$, where $F_Y=F|Y$ and
  where $F_Y/\cint(F_Y)$ denotes the quotient of $F_Y$ by its action
  on itself through inner automorphisms.
\end{enonce*}

\begin{enonce*}[remark]{Remark}
The isomorphism of Corollary \hyperref[XVI.cor3.6]{3.6} is not
canonical.
\end{enonce*}

\noindent
\emph{Proof}.\pageoriginale We immediately reduce to the case where
$X$ is noetherian and connected. Since the assertion is local on $X$
for the étale topology, we may suppose that $F$ is a constant sheaf,
with value an ordinary finite group $G$ of order $m$ prime to the
residue characteristics of $X$, and moreover that
$𝛍_{mX}\simeq (\ZZ/m\ZZ)_X$. Let $U=X-Y$ and let
$U'=\Spec \Oo_U[t^{1/m}]$, which is an “essential” principal Galois
covering of $U$ with group $\ZZ/m\ZZ$, since
$𝛍_{mX}\simeq (\ZZ/m\ZZ)_X$. Principal coverings of $U$ with group
$G$ which are trivialized over $U'$ are classified locally on $Y$ by
$H=\Hom(\ZZ/m\ZZ,G)$ modulo $\cint(G)$, which is a pointed set
isomorphic to the underlying set of $G/\cint(G)$. This gives a
morphism from the constant sheaf $H_Y$ on $Y$ to
$(R^{1}j_{\ast})G_U$, and it is enough to prove that this morphism is
bijective, which follows immediately from \hyperref[XVI.lem3.5]{3.5}.

\begin{enonce*}{Theorem 3.7}\etiquette[3.7]{XVI.thm3.7}
\emph{(Cohomological purity)}. Let $(Y,X)$ be a smooth $S$-pair of
codimension $c>0$. Let $F$ be an abelian sheaf on $X$ locally
isomorphic (for the étale topology) to a sheaf of the form
$f^{\ast}(G)$, where $G$ is a torsion sheaf on $S$, prime to the
residue characteristics (for example, $F$ may be a locally constant
sheaf of finite groups whose orders are prime to the residue
characteristics of $X$). Then, with the notation of $\hyperref[XVI.def3.1]{3.1}$, one has
\begin{align*}
&\HH^{q}_Y(X,F)=(R^{q}i^{!})F=0\qquad\mbox{if}\quad q\neq 2c,\\
&\begin{cases}
 F\xrightarrow{\sim} j_{\ast}j^{\ast}F\quad\text{and} &\tag{\ie}\\
(R^{q}j_{\ast})j^{\ast}F=0 &\hspace{1.3cm} \mbox{if}\quad q\neq 0,\quad 2c-1,
\end{cases}
\end{align*}
and\pageoriginale
$$
\HH^{2c}_Y(X,F)=i^{\ast}[(R^{2c-1}j_{\ast})j^{\ast}F]
$$
is a sheaf locally isomorphic to $i^{\ast}(F)$.
\end{enonce*}

\noindent
\emph{Proof}. The assertion is local on $X$ for the étale topology;
we may therefore \footnote[1]{At least if $F$ were assumed locally
  constant. The reader will convince himself that the proof which
  follows also applies, essentially, to the general case.}
suppose that $F$ is constant, and hence, if desired, that
$F\simeq f^{\ast}(\ZZ/n\ZZ)_S$ for the constant sheaf
$(\ZZ/n\ZZ)_S$ on $S$.

First treat the case $c=1$. (The assertion concerning
$(R^{1}j_{\ast})j^{\ast}F$, at least in the case where $F$ is locally
constant, would follow immediately from \hyperref[XVI.cor3.6]{3.6},
but we shall derive it again.) Recall that $(Y,X)$ is locally
isomorphic to the pair $(\EE^{n-1}_S,\EE^n_S)$, and that, upon
replacing $S$ by $\EE^{n-1}_S$, we reduce to the case where
$X=\EE^{1}_S=\Spec \Oo_S[t]$ and $Y$ is the section $t=0$ of
$X/S$. This situation is locally isomorphic to the inclusion of the
section at infinity in the projective line $\PP^{1}_S$. We may
therefore take $X=\PP^{1}_S$, take $Y$ to be the section at infinity,
and put $U=\EE^{1}_S=X-Y$.

Consider the Leray spectral sequence
$$
(R^{p}f_{\ast})(R^{q}j_{\ast})(\ZZ/n)\Longrightarrow
(R^{p+q}g_{\ast})(\ZZ/n){\quoi}.
$$
Its abutment is zero for $p+q>0$, by \hyperrefext[XV][XV.cor2.2]{2.2}.
Moreover, the $R^{q}j_{\ast}$ are supported on $Y$ when $q>0$, and
$Y$ maps isomorphically onto $S$ under $f$; on the other hand,
$j_{\ast}(F|U)\xleftarrow{\sim}F$ by \hyperref[XVI.prop3.2]{3.2}. It
follows that $(R^{p}f_{\ast})(R^{q}j_{\ast})=0$ if $p,q>0$, that
$f_{\ast}(R^{q}j_{\ast})$ determines the sheaf $(R^{q}j_{\ast})$ when
$q>0$, and finally that
$\sisi{\RR}{R}^{p}f_{\ast}j_{\ast}(F|U)=
\sisi{\RR}{R}^{p}f_{\ast}(F)$. The spectral sequence therefore
reduces to isomorphisms\pageoriginale
$$
f_{\ast}(R^{q}j_{\ast})(F|U)\xrightarrow{\sim}
(R^{q+1}f_{\ast})(F)\quad\mbox{\rm if}\quad q>0{\quoi}.
$$
But $(R^{q}f_{\ast})(F)$ can be computed fibre by fibre
(\hyperrefext[XII][XII.cor5.2]{5.2}); it is therefore zero unless
$q=0,2$, and is a sheaf locally isomorphic to $G$ if $q=2$, as follows
easily from \hyperrefext[XII][XII.cor5.2]{5.2} and
\hyperrefext[IX][IX.cor4.7]{4.7}. Thus $(R^{q}j_{\ast})(F|U)=0$ if
$q\neq 1$, and it is a sheaf locally isomorphic to $i^{\ast}(F)$ if
$q=1$.

The case $c>1$ is now handled easily by induction: let $(Y,X)$ be a
smooth $S$-pair of codimension $c>1$. It is clear that, locally on
$X$, one can find a subscheme $Z$ of $X$ of codimension $1$ which
contains $Y$ and is smooth over $S$, so that there is a commutative
diagram
$$
\xymatrix@=1.25cm{
Y       \ar[rr]^{u} \ar[rd]_{i}&    & Z\ar[ld]^v            \\
                  & X &,
}
$$
in which each pair is a smooth $S$-pair (one ought to call
$(Y,Z,X)$ a smooth $S$-triple), the codimension of $(Z,X)$ is $1$,
and that of $(Y,Z)$ is $c-1$.

Now there is a spectral sequence of composite functors
$$
(R^{p}u^{!})(R^{q}v^{!})F\Longrightarrow (R^{p+q}i^{!})F{\quoi}.
$$
By the case already proved,
$(R^{q}v^{!})F=0$ if $q\neq 2$ and is locally isomorphic to
$v^{\ast}(F)$ if $q=2$. The induction hypothesis, applied to
$(R^{q}v^{!})F$, then shows that
$(R^{p}u^{!})(R^{q}v^{!})F$ is zero if
$(p,q)\neq(2c-2,2)$\pageoriginale and is locally isomorphic to
$u^{\ast}v^{\ast}(F)$ if $(p,q)=(2c-2,2)$. The result for
$(R^{p+q}i^{!})F$ follows immediately, \cqfd

\begin{enonce*}{Corollary 3.8}\etiquette[3.8]{XVI.cor3.8}
Let $(Y,X)$ be a smooth $S$-pair of codimension $c$, let $n$ be a
natural number prime to the residue characteristics of $X$, and put
$$
T_{Y/X}=\HH^{2c}_Y(\ZZ/n\ZZ){\quoi},
$$
which, by $\hyperref[XVI.thm3.7]{(3.7)}$, is a sheaf on $Y$ locally isomorphic (for the
étale topology) to the sheaf $(\ZZ/n\ZZ)_Y$. Let $F$ be an abelian
sheaf on $X$, annihilated by $n$, satisfying the condition stated in
$\hyperref[XVI.thm3.7]{3.7}$. Then (for the record) $\HH^{i}_Y(F)=0$ for $i\neq 2c$, and
moreover there is a canonical isomorphism
$$
\HH^{2c}_Y(F)\simeq i^{\ast}(F)\otimes T_{Y/X}{\quoi}.
$$
\end{enonce*}

\noindent
\emph{Proof}. One easily defines a canonical homomorphism
$$
\varphi:i^{\ast}(F)\otimes T_{Y/X}\longrightarrow \HH^{2c}_Y(F)
$$
that is, a homomorphism
$$
i^{\ast}(F)\longrightarrow
\sheaf{\Hom}(\HH^{2c}_Y(\ZZ/n\ZZ),\HH^{2c}_Y(F)){\quoi},
$$
using the canonical homomorphism
$F\rightarrow\sheaf{\Hom}(\ZZ/n\ZZ,F)$. It remains to prove that
$\varphi$ is an isomorphism, which is easily checked by following the
proof given for \hyperref[XVI.thm3.7]{3.7}.

\begin{enonce*}{Corollary 3.8.1}\etiquette[3.8.1]{XVI.cor3.8.1}
Under the conditions of \hyperref[XVI.cor3.8]{3.8}, the formation of
the $\HH^{i}_Y(F)$\pageoriginale commutes with every base change
$S'\rightarrow S$.
\end{enonce*}

This is clear.

\begin{enonce*}{Corollary 3.9}\etiquette[3.9]{XVI.cor3.9}
(Absolute cohomological purity theorem). Let $X$ be a regular scheme,
and let $Y$ be a regular closed subprescheme of codimension $c$ at
every point. Suppose moreover that $X$ is locally of finite type over
a perfect field $k$. Then, if $n$ is prime to the residue
characteristics of $X$, one has
$$
\HH^{i}_Y(\ZZ/n\ZZ)=0\mbox{ if } i\neq 2c{\quoi},
$$
and $\HH^{2c}_Y(\ZZ/n\ZZ)$ is a sheaf on $Y$ locally isomorphic (for
the étale topology) to the constant sheaf $\ZZ/n\ZZ$.
\end{enonce*}

This is a special case of \hyperref[XVI.thm3.7]{3.7}, since the
hypothesis implies that $X$ and $Y$ are smooth over $k$.

\begin{enonce*}[remark]{Remarks 3.10}\etiquette[3.10]{XVI.rems3.10}
\begin{enumerate}
\renewcommand{\theenumi}{\alph{enumi}}
\renewcommand{\labelenumi}{\theenumi)}
\item In \hyperref[XVI.cor3.8]{3.8} one can make the structure of the
sheaf $T_{Y/X}$ explicit: one finds that it is canonically isomorphic
to the sheaf $(𝛍_n)_Y^{\otimes -c}$, where $𝛍_n$ denotes the sheaf of
$n$th roots of unity (locally isomorphic to the constant sheaf
$\ZZ/n\ZZ$). The isomorphism in question is studied, under the name
of the local \emph{fundamental class} of $Y$ in $X$, in the slightly
more general setting of relative complete intersections, in Exposé
XVIII, which is devoted to duality.

\item It\pageoriginale is very plausible that the conclusions of
\hyperref[XVI.cor3.9]{3.9} and \hyperref[XVI.rems3.10]{3.10} a) remain
valid without assuming the existence of a base field $k$, at least
when $X$ is an “excellent” prescheme\SGAArchiveOnly{\nde{The frozen source attributes
this later development to work of Gabber but supplies no bibliographic
target.}}. This
will in any case be established in Exposé XIX when $X$ is of
characteristic zero, using Hironaka's resolution of singularities. As
we shall see in \SGA 5, the “absolute purity theorem” (as well as
resolution of singularities) is needed in particular to establish the
“local duality formula”, which itself is one of the principal
ingredients of the \textsc{Lefschetz--Verdier} formula.
\end{enumerate}
\end{enonce*}

% END INCLUDED SGA4_English_current_buildable_source_20260906/source/16/drafts/SGA4_XVI_LINES_701_956_English.texfrag


% BEGIN INCLUDED SGA4_English_current_buildable_source_20260906/source/16/drafts/SGA4_XVI_LINES_957_1200_English.texfrag
\section{Comparison theorem for the cohomology of algebraic schemes
over $\CC$.}\etiquette[4]{XVI.sec4}

Let $X$ be an algebraic scheme over $\Spec \CC$. We shall compare
the étale and classical sites of $X$, and we retain the notation of
\hyperrefext[XI][XI.sec4]{4}. The comparison theorem, whose general
form is \hyperrefext[XI][XI.thm4.4]{4.4}, is as follows:

\begin{enonce*}{Theorem 4.1}\etiquette[4.1]{XVI.thm4.1}
Let $f:X\rightarrow S$ be a morphism of finite type between schemes
locally of finite type over $\Spec \CC$, so that there is a
commutative diagram of morphisms of sites
$$
\xymatrix@C=2cm@R=1.5cm{
X_{\et}  \ar[d]_{f_{\et}}                   &   X_{\cl}
\ar[l]_{\varepsilon} \ar[d]^{f_{\cl}} \\
S_{\et}                    &   S_{\cl}\ar[l]_{\varepsilon}
}
$$
Let\pageoriginale $F$ be a sheaf of sets (\resp of ind-finite
groups, \resp a torsion abelian sheaf), and suppose that one of the
following cases holds:
\begin{enumerate}
\renewcommand{\theenumi}{\roman{enumi}}
\renewcommand{\labelenumi}{(\theenumi)}
\item $f$ is proper,

\item $F$ is constructible.
\end{enumerate}
Then the base-change morphisms \sisi{\emph{XII 4.1.2}}{\hyperrefext[XII][XII.eq4.2]{4.2}}
\footnote[1]{In fact, we use the base-change homomorphism under more
general conditions than those considered in Exposé XII, where we
restricted ourselves to a \emph{Cartesian} square of morphisms of
\emph{preschemes}. The reader will readily verify that the definition
extends verbatim to an \emph{essentially commutative} square of
morphisms of \emph{sites}. See the following exposé for further
developments on this point.}
$$
\varphi:\epsilon^{\ast}(R^{q}f_{\et\ast})F\longrightarrow
(R^{q}f_{\cl\ast})\epsilon {}^{\ast}F
$$
are bijective for $q=0$ (\resp for $q=0$, $1$, \resp for $q\geq
0$).
\end{enonce*}

Observe that the Grauert--Remmert theorem
(\hyperrefext[XI][XI.thm4.3]{4.3} (iii)), which we used in the proof
of \hyperrefext[XI][XI.thm4.4]{4.4}, is a special case of
\hyperref[XVI.thm4.1]{4.1}.

For\pageoriginale the proof, we first treat the case in which $f$ is
proper. In this case, the result will follow from GAGA and the proper
base-change theorem, by an elementary argument that could have
appeared in Exposé XIV. To treat the general case, we shall also use
resolution of singularities [1] and \hyperref[XVI.thm3.7]{3.7}.

We may clearly assume that $S$, and hence also $X$, is of finite type
over $\CC$.

\emph{The case in which $f$ is proper.} We recall the following
elementary facts:
\begin{enumerate}
\renewcommand{\theenumi}{\alph{enumi}}
\renewcommand{\labelenumi}{(\theenumi)}
\item The topological space $X_{\cl}$ is locally compact.

\item If $f:X\rightarrow S$ is proper, then
$f_{\cl}:X_{\cl}\rightarrow S_{\cl}$ is a \emph{proper} map of
topological spaces.
\end{enumerate}

Indeed, (a) is trivial, since locally $X_{\cl}$ is a closed subset of
some $\CC^n$. For (b), Chow's lemma reduces us to the case in which
$f$ is projective, and this case is trivial.

Let $F$ be a sheaf of sets (\resp\ldots) on $X_{\et}$.
For both morphisms $f_{\et}$ and $f_{\cl}$, formation of the $R^{q}$
commutes with formation of the fibers (\hyperrefext[XII][XII.cor5.2]{5.2}
and \cite{sga4.exp.XVI.bib.2}, 4.11.1). It is therefore enough to verify the theorem
fiber by fiber, at the closed points of $S_{\et}$.

In the following lemma, we do not assume that $g$ is proper:

\begin{enonce*}{Lemma 4.2}\etiquette[4.2]{XVI.lem4.2}
Let\pageoriginale
$$
\xymatrix@=1.5cm{
X  \ar[rd]_{g}& &\bar{X} \ar[ll]_{h}\ar[ld]^{\bar{g}}\\
& S &
}
$$
be a commutative diagram with $h$ proper, and let $U$ be an open
subset of $X$ such that $h$ induces an isomorphism from the open
subset $\bar{U}=\fprod{U}{\bar{X}}{X}$ of $\bar{X}$ onto $U$. Let
$j:U\rightarrow X$ and $\bar{j}:\bar{U}\rightarrow \bar{X}$ be the
inclusions, let $F$ be a sheaf of sets (\resp\ldots) on $U$, and let
$\bar{F}$ be the induced sheaf on $\bar{U}$. If Theorem $\hyperref[XVI.thm4.1]{4.1}$ is true
for $(\bar{g},\bar{j}_{!}\bar{F})$, then it is also true for
$(g,j_{!}F)$.
\end{enonce*}

\noindent
\emph{Proof}. The theorem is true for
$(h,\bar{j}_{!}\bar{F})$, because we are reduced to checking it fiber
by fiber, and there are then two cases: either the restriction of the
sheaf to the fiber is zero, or the morphism $h$ is an isomorphism. In
both cases the result is trivial. We find that
$$
\epsilon^{\ast}j_{!}F\simeq
\epsilon^{\ast}h_{\text{\rm et\,}\ast}(\bar{j}_{!}\bar{F})\simeq
h_{\cl\ast}\epsilon^{\ast}(\bar{j}_{!}\bar{F})
$$
and
$$
(R^{q}h_{\text{\rm et\,}\ast})\bar{j}_{!}\bar{F}
=(R^{q}h_{\cl\ast})\epsilon^{\ast}(\bar{j}_{!}\bar{F})=0
\text{ if }q>0
$$
for the values of $q$ under consideration, and hence that
$$
(R^{q}g_{\text{\rm et\,}\ast})j_{!}F\simeq
(R^{q}\bar{g}_{\text{\rm et\,}\ast})\bar{j}_{!}\bar{F}
$$
and
$$
(R^{q}g_{\cl \ast})j_{!}F\simeq (R^{q}\bar{g}_{\cl
\ast})\bar{j}_{!}\bar{F}{\quoi},
$$
which proves \hyperref[XVI.lem4.2]{Lemma 4.2}.\pageoriginale

\begin{enonce*}{Lemma 4.3}\etiquette[4.3]{XVI.lem4.3}
To verify \hyperref[XVI.thm4.1]{4.1} for all data $(f,F)$ with $f$
proper, it is enough to do so when, in addition, $f$ has relative
dimension $\leq 1$ and $S=\Spec \CC$ is a point.
\end{enonce*}

\noindent
\emph{Proof}. We proceed by induction on the relative dimension.
Suppose the result known in relative dimension $<n$, where $n>1$.
As we have observed, we may check it fiber by fiber; that is, we may
suppose that $S=\Spec \CC$ is a point, and hence that $X$ has
dimension $\leq n$. Moreover, we may assume that $X$ is reduced. Let
$\varphi$ be a rational function on $X$ that is nonconstant on
\sisi{every irreducible component}{every irreducible component} of
$X$, let $U$ be a dense open subset of $X$
on which $\varphi$ is defined, and let
$\bar{X}\subset X\times \PP^{1}$ be the closure of the graph of the
morphism $\varphi:U\rightarrow \EE^{1}$. We have a commutative
diagram
$$
\xymatrix@R=1.5cm@C=1.5cm{
&\bar{X}\ar[ld]_{h} \ar[rd]^{a} \ar[dd]^{\bar{g}}   &\\
X  \ar[rd]_{g}  &                                  &\PP^1 \ar[ld]_{b}\\
            & **[r]Y              =\Spec \CC&
}
$$
in which the relative dimensions of $a$ and $b$ are $<n$. By the
induction hypothesis, the theorem is true for $(a,.)$ and
$(b,.)$, and it follows that it is also true for
$(\bar{g},.)$. Indeed, this is trivial for $q=0$; if $F$
is a torsion abelian sheaf on $\bar{X}$, it follows\pageoriginale from
the morphism of Leray spectral sequences
$$
\xymatrix{
\varepsilon^{\ast}(R^{p}b_{\text{et}\ast})(R^{q}a_{\text{et}\ast})F \ar@{=>}[r] \ar[d]&
\varepsilon^{\ast}(R^{p+q}\bar{g}_{\text{et}\ast})F  \ar[d]\\
(R^{p}b_{\text{cl}\ast}(R^{q}{a_{\text{cl}\ast}}))
\varepsilon^{\ast}F  \ar@{=>}[r] & (R^{p+q}\bar{g}_{\text{cl}\ast})\varepsilon^{\ast}F;
}
$$
finally, if $F$ is a sheaf of ind-finite groups on $\bar{X}$, note
that since $\epsilon^{\ast}(R^{1}\bar{g}_{\text{\rm et\,}\ast})$ is
an effaceable functor, it is enough to prove (in view of the result for
$q=0$) that $(R^{1}\bar{g}_{\cl\ast})\epsilon^{\ast}$ is also
effaceable (\cf \hyperrefext[XII][XII.cor8.2]{8.2}). This follows
from the exact sequence
$$
0\longrightarrow (R^{1}b_{\cl\ast})a_{\cl
\ast}\epsilon^{\ast}F\longrightarrow (R^{1}\bar{g}_{\cl
\ast})\epsilon^{\ast}F\longrightarrow b_{\cl \ast}(R^{1}a_{\cl
\ast})\epsilon^{\ast}F{\quoi},
$$
and from the fact that $(R^{1}b_{\cl\ast})\epsilon^{\ast}$ and
$(R^{1}a_{\cl\ast})\epsilon^{\ast}$ are effaceable by the induction
hypothesis.

Now let $F$ be a sheaf of sets on $X$. We have
$F\hookrightarrow h_{\text{\rm et\,}\ast}h_{\et}{}^{\ast}F=G$.
The exact sequence
\begin{equation*}
F\rightarrow G\rightrightarrows G\amalg_F G{\quoi},\tag{4.3.1}
\end{equation*}\etiquette[4.3.1]{XVI.eq4.3.1}
reduces us (for $q=0$) to proving \hyperref[XVI.thm4.1]{4.1} for $G$,
and hence for a sheaf of the form
$h_{\text{\rm et\,}\ast}\bar{F}$. For such a sheaf, however, the
result follows from the theorem for $(\bar{g},.)$ and for
$(h,.)$; and $h$ has relative dimension $<n$, as is
immediate.

Let $F$ be a torsion abelian sheaf on $X$. We have the exact sequence
\begin{equation*}
0\longrightarrow j_{!}j^{\ast}F\longrightarrow F\longrightarrow
i_{\ast}i^{\ast}F\longrightarrow 0{\quoi}.\tag{$\ast$}
\end{equation*}
Since $U$ is dense in $X$, we have $\dim Y<n$, so
\hyperref[XVI.thm4.1]{4.1} is true\pageoriginale for
$(g|Y,i^{\ast}F)$, and hence for $(g,i_{\ast}i^{\ast}F)$. It is
therefore enough, by the five lemma, to verify
\hyperref[XVI.thm4.1]{4.1} for $(g,j_{!}j^{\ast}F)$, and this follows
from the theorem for $\bar{g}$ (\hyperref[XVI.lem4.2]{4.2}).

Let $F$ be a sheaf of ind-finite groups on $X$. It is enough to prove
that the functor $(R^{1}g_{\cl\ast})\epsilon^{\ast}$ is effaceable for
the sheaf $F$. This follows from ($\ast$) and from effaceability for
$j_{!}j^{\ast}F$ and for $i_{\ast}i^{\ast}F$, that is, from the
theorem for $(g,j_{!}j^{\ast}F)$ (\hyperref[XVI.lem4.2]{4.2}) and for
$(g,i_{\ast}i^{\ast}F)$ (the induction hypothesis), \cqfd

\begin{enonce*}{Lemma 4.4}\etiquette[4.4]{XVI.lem4.4}
To prove $\hyperref[XVI.thm4.1]{(4.1)}$ when $f$ is proper, it is enough
to prove the following assertion: let $X$ be a complete regular curve
over $\Spec \CC$, and let $F$ be a constant constructible sheaf of
sets (\resp\ldots) on $X$. Then the morphisms
$$
H^{q}(X_{\et},F)\longrightarrow H^{q}(X_{\cl},F)
$$
are bijective for the values of $q$ under consideration.
\end{enonce*}

\noindent
\emph{Proof}. Since it is enough to verify the theorem fiber by fiber,
\hyperref[XVI.lem4.3]{4.3} reduces us to the case in which $X$ is
proper and has dimension $\leq 1$. The case of dimension $0$ is trivial,
and it follows that the theorem is true for a finite morphism. Since
$f$ is proper, the cohomology of $X_{\et}$ and of $X_{\cl}$ commutes
with filtered direct limits (\hyperrefext[VII][VII.prop3.3]{3.3} and
[2], 4.12.1), and we are therefore reduced to the case in which $F$
is constructible (\hyperrefext[IX][IX.prop2.9]{2.9} (iii)).

We apply \hyperrefext[IX][IX.prop2.14]{2.14}. For a sheaf of sets, the
exact sequence \eqhyperref[XVI.eq4.3.1]{4.3.1} and
\hyperrefext[IX][IX.prop2.14]{2.14} show that we may take
$F=\pi{}_{\ast}C$, where\pageoriginale $\pi:X'\rightarrow X$ is
finite, $X'$ is normal and hence regular, and $C$ is a constant sheaf
on $X'$. For a torsion abelian sheaf, we take a resolution of $F$ by
finite products of sheaves of the form $\pi_{\ast}C$ and apply the
spectral sequence of a resolution, again reducing to the case
$F=\pi_{\ast}C$. For a sheaf of ind-finite groups, recall that it is
enough to prove effaceability of
$(R^{1}g_{\cl\ast})\epsilon^{\ast}$ for the sheaf $F$, hence for a
sheaf $G$ into which $F$ can be embedded; thus we are again reduced to
the case $F=\pi_{\ast}C$. It is now clear from
(\hyperrefext[VIII][VIII.prop5.5]{5.5} and \hyperrefext[VIII][VIII.rem5.8]{5.8}) that

% END INCLUDED SGA4_English_current_buildable_source_20260906/source/16/drafts/SGA4_XVI_LINES_957_1200_English.texfrag


% BEGIN INCLUDED SGA4_English_current_buildable_source_20260906/source/16/drafts/SGA4_XVI_LINES_1201_1426_English.texfrag
we may now replace $X$ by $X'$, which proves the lemma.

To treat this last case, it is plainly enough (in view of the
cohomological-dimension results
(\sisi{IX 7.7}{\hyperrefext[IX][IX.cor5.7]{5.7}}) and
\cite{sga4.exp.XVI.bib.2}, 4.14.1) to prove that for a complete regular curve one has
\begin{itemize}
\item[$(q=0)$:] $X$ is connected and nonempty if and only if
$X_{\cl}$ is connected and nonempty.

\item[$(q=1)$:] The base-change morphism induces an equivalence between
the category of étale coverings of $X$ and the category of finite
analytic étale coverings of $X_{\cl}$.

\item[$(q=2)$:] There is an isomorphism
$\epsilon:H^2(X_{\et},\mu_n)\xrightarrow{\sim}
H^2(X_{\cl},\mu_n)$.
\end{itemize}

The first two assertions follow immediately from GAGA. For the last,
note that $H^2(X_{\et},\GG_m)=0$
(\hyperrefext[IX][IX.eq4.5]{4.5}). Likewise, if $\Oo^*$ denotes the
sheaf of invertible holomorphic functions on $X_{\cl}$, then
$H^2(X_{\cl},\Oo^*)=0$. This follows from the exact sequence
$$
0\longrightarrow\ZZ\longrightarrow\Oo\xrightarrow{\exp}
\Oo^*\longrightarrow0
$$
and\pageoriginale from the fact that
$H^2(X_{\cl},\Oo)=H^2(X_{\et},\Oo_X)=0$ by GAGA. Consequently the
Kummer exact sequence \hyperrefext[IX][IX.Thr3.2]{3.2} (respectively,
the exact sequence on $X_{\cl}$
$$
0\rightarrow\mu_n\rightarrow
\Oo^*\xrightarrow{n}\Oo^*\rightarrow0)
$$
gives
\begin{align*}
&H^1(X_{\et},\GG_m)/n\xrightarrow{\sim}
H^2(X_{\et},\mu_n)\\
(\resp\quad &H^1(X_{\cl},\Oo^*)/n\xrightarrow{\sim}
H^2(X_{\cl},\mu_n)){\quoi}.
\end{align*}
By GAGA, the canonical morphism
$$
H^1(X_{\et},\GG_m)=\Pic X\longrightarrow
\Pic X_{\cl}=H^1(X_{\cl},\Oo^*)
$$
is bijective. This proves the result, and hence proves
\hyperref[XVI.thm4.1]{4.1} when $f$ is proper.

\medskip
\noindent
\emph{The case where $F$ is constructible}. The proof proceeds by
induction on $\dim X$. Suppose the result known in dimensions $<n$,
and prove it when $\dim X\leq n$.

\begin{enonce*}{Lemma 4.5}\etiquette[4.5]{XVI.lem4.5}
We may suppose that $f:X\rightarrow S$ is a dense open immersion and
that $F$ is constant.
\end{enonce*}

\noindent
\emph{Proof}. As in the proof of \hyperref[XVI.lem1.3]{1.3}, we may
suppose that $X$ and $S$ are affine, and hence that $f$ is
quasi-projective. We then apply the argument of
\hyperref[XVI.lem4.4]{4.4} to reduce to $F=\pi_*C$, where
$\pi:X'\rightarrow X$ is finite and $C$ is constant. Replacing $X$ by
$X'$ therefore reduces us\pageoriginale to the case where $F$ is
constant. We may replace $X$ by $X_{\red}$. Let
$$
\xymatrix@=1.25cm{
  X  \ar[rr]^{h}\ar[rd]_{f} &      &\bar{X} \ar[ld]^{\bar{f}}\\
             &     S  &
}
$$
be a commutative diagram in which $\bar X$ is reduced, $\bar f$ is
projective, and $h$ is a dense open immersion. It is enough to verify
the theorem for the two morphisms $h$ and $\bar f$ (compare the proof
of \hyperref[XVI.lem4.3]{4.3}); hence it is enough to verify it for
$h$, since $\bar f$ is proper. We are thus reduced to a dense open
immersion.

\begin{enonce*}{Lemma 4.6}\etiquette[4.6]{XVI.lem4.6}
We may further suppose that $S$ is regular.
\end{enonce*}

\noindent
\emph{Proof}. Let $S'\xrightarrow{h}S$ be a ``resolution of the
singularities of $S$'', that is, a surjective, proper, birational
morphism with $S'$ regular, and let
$$
\xymatrix@C=2cm@R=1.25cm{
X'     \ar[r]^{f'} \ar[d]^{h'}&     S'\ar[d]^{h}\\
X   \ar[r]^{f}  & S
}
$$
be the Cartesian diagram obtained from $h$.

Let $F$ be a constant sheaf of sets on $X$. We have
$F\hookrightarrow G=h'_*{h'}^*F$, and it follows that, for $q=0$, it
is enough to verify \hyperref[XVI.thm4.1]{4.1} for $G$ (compare
\eqhyperref[XVI.eq4.3.1]{4.3.1}), that is, for a sheaf of the form
${h'}_*F'$, where $F'$ is constant on $X'$. We have a
commutative\pageoriginale diagram (in which the symbol ``ét'' is
suppressed)
$$
\xymatrix@C=2cm@R=1.5cm{
\varepsilon^*f_* (h'_*F') = \varepsilon^*h_*f'_*F'  \ar[d]^{d}\ar[r]^(.6){a} & h_{\cl *}
	    \varepsilon^*f'_*F' \ar[d]^{b}\\
	    f_{\cl *}\varepsilon^*(h'_*F')  \ar[r]^(.4){c} &
	    f_{\cl *} h'_{\cl *}\varepsilon^*F' = h_{\cl *}f'_{\cl *}
	    \varepsilon^*F',
}
$$
and $a$ and $c$ are bijective by the result for a proper morphism. To
prove that $d$ is bijective, it is enough to prove that $b$ is, that
is, that
$$
\epsilon^*f'_*F'\xrightarrow{\sim}
f'_{\cl *}\epsilon^*F'.
$$
We are therefore reduced to $S=S'$, with $S$ regular.

\noindent
(Remark: \sisi{We}{we} have used, without mentioning it explicitly,
the finiteness theorem for a proper morphism
\hyperrefext[XIV][XIV.thm1.1]{1.1}.)

Let $F$ be a constant torsion abelian sheaf on $X$. It is enough to
treat $F=(\ZZ/n)_X$. The morphism $h'$ is an isomorphism above a dense
open subset $j:U\rightarrow X$; let $j':U'\rightarrow X'$ be the open
subset ${h'}^{-1}(U)\xrightarrow{\sim}U$, and put $Y=X-U$. Since $U$
is dense, $\dim Y<n$, so the theorem holds for
$(f,(\ZZ/n)_Y)$. The exact sequence
$$
0\longrightarrow j_!(\ZZ/n)_U\longrightarrow
(\ZZ/n)_X\longrightarrow(\ZZ/n)_Y\longrightarrow0
$$
and the five lemma reduce us to proving the theorem for
$j_!(\ZZ/n)_U$ and $f$, and hence, by
\hyperref[XVI.lem4.2]{4.2}, to proving it for
$(hf',{j'_!}(\ZZ/n)_{U'})$. We have an exact sequence
$$
0\longrightarrow {j'_!}(\ZZ/n)_{U'}\longrightarrow
(\ZZ/n)_{X'}\longrightarrow(\ZZ/n)_{Y'}\longrightarrow0{\quoi},
$$
and\pageoriginale the theorem holds for $(hf',(\ZZ/n)_{Y'})$ by the
induction hypothesis. It is therefore enough to prove it for
$(hf',(\ZZ/n)_{X'})$. But the proper case gives the theorem for
$(h,\cdot)$, and the Leray spectral sequence for the pair of morphisms
$h,f'$ shows that it is enough to prove it for $(f',\ZZ/n)$. This
proves the lemma in the abelian case.

Finally, suppose that $F$ is a constant sheaf of finite groups. Recall
that it is enough to prove that the functor
$(R^1f_{\cl *})\epsilon^*(\mathord\cdot)$ is effaceable for $F$, hence
for a sheaf $G$ into which $F$ embeds, and therefore for
$G={h'}_*{h'}^*F$, that is, for a sheaf of the form ${h'}_*F'$, where
$F'$ is constant on $X'$. We have
$$
\epsilon^*({h'}_*F')\xrightarrow{\sim}
h'_{\cl *}(\epsilon^*F')
$$
and
$$
0\longrightarrow
(R^1f_{\cl *})h'_{\cl *}\epsilon^*(\mathord\cdot)\longrightarrow
(R^1(f_{\cl}h'_{\cl})_*)\epsilon^*(\mathord\cdot)
=(R^1(h_{\cl}f'_{\cl})_*)\epsilon^*(\mathord\cdot){\quoi}.
$$
Thus it is enough to prove that the last functor is effaceable for a
constant sheaf. But the exact sequence
$$
0\longrightarrow (R^1h_{\cl *})f'_{\cl *}
\epsilon^*(\mathord\cdot)\longrightarrow
(R^1(h_{\cl}f'_{\cl})_*)\epsilon^*(\mathord\cdot)\longrightarrow
h_{\cl *}(R^1f'_{\cl *})\epsilon^*(\mathord\cdot){\quoi}
$$
shows that the first term is effaceable, for example by the Godement
resolution (\hyperrefext[XII][XII.prop3.3]{3.3}), because $h$ is
proper. It is therefore enough to prove that the right-hand term is
effaced by the Godement resolution, hence that
$(R^1f'_{\cl *})\epsilon^*(\mathord\cdot)$ is so effaced, and thus
that the theorem holds for $(f',F')$, where $F'$ is constant. This
proves the lemma.

\noindent
\emph{End of the proof}.\pageoriginale We now suppose that $F$ is a
constant sheaf and that $f:X\rightarrow S$ is a dense open immersion,
with $S$ regular. The assertion for sheaves of sets now follows
immediately from \hyperref[XVI.prop3.2]{3.2} and the analogous trivial
lemma for the classical topology\sisi{---indeed}{. Indeed}, these show
that $\epsilon^*f_*F$ and $f_{\cl *}\epsilon^*F$ are constant with the
same value and hence are isomorphic under $\varphi$.

For $q>0$, we first reduce further to the case $X=S-Y$ with $Y$
regular, that is, to the case in which $(Y,S)$ is a smooth
$\Spec\CC$-pair. Indeed, write $X=S-C_1$, where $C_1$ is given its
induced reduced structure, and define recursively

$Y_\nu=$ the dense open subset of regular points of $C_\nu$,

$C_{\nu+1}=C_\nu-Y_\nu$, with its induced reduced structure.

\noindent
Since $Y_\nu$ is dense in $C_\nu$, we have
$\dim C_{\nu+1}<\dim C_\nu$, and hence a sequence
$$
X=X_1\subset X_2=S-C_2\subset X_3=S-C_3\subset\cdots\subset X_r=S{\quoi},
$$
where the inclusions $i_\nu:X_\nu\rightarrow X_{\nu+1}$ are open
immersions and
$$
X_\nu=X_{\nu+1}-Y_\nu,\qquad
Y_\nu\ \mbox{closed in $X_{\nu+1}$ and regular}{\quoi}.
$$
The direct image $i_{\nu *}F$ of a constant sheaf on $X$ is constant
(\hyperref[XVI.prop3.2]{3.2}), and the $R^qi_{\nu *}F$ are
constructible and concentrated on the subvarieties of $Y$ of
dimension at most $n-1$ (\hyperref[XVI.cor3.4]{3.4},
\hyperref[XVI.cor3.6]{3.6}, \hyperref[XVI.thm3.7]{3.7}) for the values
of $q$ under consideration. The induction hypothesis and the Leray
spectral sequence (respectively, the exact sequence
\hyperrefext[XII][XII.prop3.2]{3.2}) therefore reduce us\pageoriginale
to the case of $i_\nu$, which proves the asserted reduction.

For a smooth pair $(Y,S)$, however, the values of $(R^qf_*)F$ have
been computed explicitly for the étale topology
(\hyperref[XVI.cor3.4]{3.4}, \hyperref[XVI.cor3.6]{3.6},
\hyperref[XVI.thm3.7]{3.7}); an analogous and easy computation, which
we leave to the reader, gives the corresponding results for the
classical topology. Direct comparison proves the theorem for $f$.


% END INCLUDED SGA4_English_current_buildable_source_20260906/source/16/drafts/SGA4_XVI_LINES_1201_1426_English.texfrag


% BEGIN INCLUDED SGA4_English_current_buildable_source_20260906/source/16/drafts/SGA4_XVI_LINES_1427_1560_English.texfrag
\section{The finiteness theorem for algebraic schemes in
characteristic zero}\etiquette[5]{XVI.sec5}

\begin{enonce*}{Theorem 5.1}\etiquette[5.1]{XVI.thm5.1}
Let $k$ be a field of characteristic zero, and let $f:X\rightarrow S$
be a morphism of finite type between schemes locally of finite type
over $k$. Let $F$ be a constructible sheaf of sets (\resp of finite
groups, \resp a torsion abelian sheaf) on $X$. Then the
$R^qf_*F$ are constructible for $q=0$ (\resp for $q=0,1$, \resp for
every $q$).
\end{enonce*}

The hypothesis that $k$ have characteristic zero is imposed because
we shall use the resolution-of-singularities theorem \cite{sga4.exp.XVI.bib.1}. The
proof is also valid in characteristic $p>0$ if resolution of
singularities is assumed, \emph{provided that the orders of $F$ are
prime to $p$}. Thus in nonzero characteristic one obtains a result for
$\dim X\leq2$ by applying Abhyankar's results \cite{sga4.exp.XVI.bib.3}. In addition to
the ever-present\pageoriginale proper base-change theorem, another
ingredient of the proof is the ``absolute cohomological purity
theorem'' \hyperref[XVI.cor3.9]{3.9}\SGAArchiveOnly{\nde{The frozen source attributes
this later development to work of Deligne and Gabber but supplies no
bibliographic target.}}.

Recall that the theorem has already been proved for a proper morphism
and arbitrary $S$ (\hyperrefext[XIV][XIV.thm1.1]{1.1}). Moreover, in
Exposé XIX we shall obtain (again in characteristic zero) the finiteness
theorem under the much weaker hypotheses that $S$ be an
\emph{excellent} scheme and $f$ be of finite type, by proving for such
schemes the absolute-purity theorem in the form indicated in
\hyperref[XVI.rems3.10]{3.10}.

In nonzero characteristic and dimension at least $3$, for the time
being we have only the following easy consequence of
\hyperrefext[XI][XI.prop3.3]{3.3}:

\begin{enonce*}{Theorem 5.2}\etiquette[5.2]{XVI.thm5.2}
Let $X$ be a smooth algebraic scheme over $\Spec k$, where $k$ is a
separably closed field, and let $F$ be a locally constant sheaf of
finite abelian groups (\resp finite groups), of orders prime to the
characteristic of $k$. Then $H^q(X,F)$ is a finite group (\resp a
finite pointed set) for every $q$ (\resp for $q=0,1$).
\end{enonce*}

\noindent
\emph{Proof of} \hyperref[XVI.thm5.2]{5.2}. We may suppose that $k$ is
algebraically closed (\hyperrefext[VIII][VIII.thm1.1]{1.1}). We argue
by induction on $n=\dim X$. By \hyperrefext[XI][XI.prop3.3]{3.3}, there
is a hypercovering of $X$ by open subsets that are ``good
neighborhoods''. The spectral sequence of a hypercovering immediately
reduces us to the case in which $X$ is a good neighborhood and hence
admits an elementary fibration (\hyperrefext[XI][XI.def3.1]{3.1})
$$
\xymatrix@=1.5cm{
X \ar[rr]^{j} \ar[rd]_{f}&         &
\bar{X}\ar[ld]^{\bar{f}} \\
             &      S  & .
}
$$
By\pageoriginale \hyperref[XVI.prop3.2]{3.2},
\hyperref[XVI.cor3.6]{3.6}, and \hyperref[XVI.thm3.7]{3.7}, the sheaf
$j_*F$ is locally constant on $\bar X$, $(R^1j_*)F$ is locally
constant on $Y=\bar X-X$, and $(R^qj_*)F=0$ for $q>1$. It follows from
\hyperref[XVI.thm2.1]{2.1} that
$(R^p\bar f_*)(R^qj_*)F$ is a locally constant sheaf on $S$ for the
values of $q$ under consideration. The result follows from the
induction hypothesis and the Leray spectral sequence (\resp the exact
sequence \hyperrefext[XII][XII.prop3.2]{3.2}).

The proof of \hyperref[XVI.thm5.1]{5.1} is very close to the proof of
\hyperref[XVI.thm4.1]{4.1} in the constructible case. We merely
indicate the main steps. Argue by induction on
$\dim X=n$. First reduce to $X$ and $S$ affine, and hence $f$
quasi-projective, by applying the method used in the proof of \hyperref[XVI.lem1.3]{1.3}.
Next apply \sisi{IX 2.15}{\hyperrefext[IX][IX.prop2.14]{2.14}} and
\hyperrefext[VIII][VIII.prop5.5]{5.5},
\hyperrefext[VIII][VIII.rem5.8]{5.8} to reduce to the case in which
$F$ is constant. There is a commutative diagram
$$
\xymatrix@=1.5cm{
X \ar[rr]^{h} \ar[rd]_{f}&         & \bar{X}\ar[ld]^{\bar{f}}\\
             &      S  & ,
}
$$
in which $h$ is a dense open immersion and $\bar f$ is proper, so the
theorem holds for $\bar f$
(\hyperrefext[XIV][XIV.thm1.1]{1.1}). It is therefore enough to prove
the theorem for $(h,F)$, that is, for a dense open
immersion\pageoriginale and a constant sheaf. Let
$h:S'\rightarrow S$ be a resolution of the singularities of $S$, and
let
$$
\xymatrix@C=1.75cm@R=1cm{
X  \ar[d]^{f} & X' \ar[l]_{h'} \ar[d]^{f'}\\
S          & S' \ar[l]_{h}
}
$$
be the Cartesian diagram obtained from $h$. We reduce to proving the
theorem for $f'$, and hence to the case in which $S=S'$ is regular.
The inductive method used at the end of the proof of
\hyperref[XVI.thm4.1]{4.1} reduces us to the case in which $Y=S-X$ is
nonsingular, that is, $(Y,S)$ is a smooth pair. The proof is completed
by applying \hyperref[XVI.prop3.2]{3.2},
\hyperref[XVI.cor3.4]{3.4}, \hyperref[XVI.cor3.6]{3.6}, and
\hyperref[XVI.thm3.7]{3.7}.

\begin{enonce*}[remark]{Remark 5.3}\etiquette[5.3]{XVI.rem5.3}
The proof of \hyperref[XVI.thm5.1]{5.1} just outlined is also valid
when one is given a left Noetherian ring $A$ annihilated by an integer
$n>0$ (which must be assumed prime to the characteristic when the
latter is not assumed to be zero; compare the remarks following
\hyperref[XVI.thm5.1]{5.1}) and considers constructible sheaves of left
$A$-modules.
\end{enonce*}

\begin{thebibliography}{99}
\bibitem{sga4.exp.XVI.bib.1} Hironaka, H. Resolution of singularities of an
algebraic variety over a field of characteristic zero, Annals of
Math., vol. 79 (1964), p. 109.

\bibitem{sga4.exp.XVI.bib.2} Godement,\pageoriginale P. Théorie des
\sisi{Faisceaux}{faisceaux}, Paris, 1958\sisi{}{.}

\bibitem{sga4.exp.XVI.bib.3} Abhyankar, S. Local uniformization on algebraic
surfaces over ground fields of characteristic $p\neq0$, Ann. of
Math., vol. 63 (1956), p. 491.

\bibitem{sga4.exp.XVI.bib.4} J.-P. Serre, \sisi{Géométrie Algébrique et Géométrie
Analytique}{Géométrie algébrique et géométrie analytique},
Annales Inst. Fourier 1956, pp. 1--42 (\sisi{Cited}{cited} GAGA).
\end{thebibliography}

% XII 4.1 no ref. 78
% XII 4.1.2 no ref. 992
% 4.11.1 no ref. 1040
% XI 7.7 no ref. 1210
% IX 4.5 is an equation ref.
% IX 2.15 no ref. 1504

\ifx\danslelivre\undefined
\end{document}
\fi

% END INCLUDED SGA4_English_current_buildable_source_20260906/source/16/drafts/SGA4_XVI_LINES_1427_1560_English.texfrag


% END INCLUDED SGA4_English_current_buildable_source_20260906/source/16/16_English.tex


\makeatletter
\def\input@path{{../17/}}
\makeatother
% BEGIN INCLUDED SGA4_English_current_buildable_source_20260906/source/17/17_English.tex
\ifx\danslelivre\undefined
% Release engine: XeLaTeX (the LuaTeX compatibility layer breaks Xy-pic output).
\documentclass[10pt]{../fichiers-configuration-latex/smfart}
% BEGIN INCLUDED SGA4_English_current_buildable_source_20260906/source/fichiers-configuration-latex/paquets.tex
\usepackage{polyglossia}
%\usepackage[francais]{babel}
\usepackage{fontspec}
\usepackage[math-style=ISO,bold-style=upright,sans-style=upright]{unicode-math}
%\usepackage[mathletters]{ucs}
%\usepackage[utf8]{inputenc}
%\usepackage[T1]{fontenc}
\usepackage[unicode,hypertexnames=false,colorlinks=true,linkcolor=blue,citecolor=Sepia,urlcolor=magenta]{hyperref}
\setmainlanguage{french}
\setmainfont[Ligatures=TeX]{Linux Libertine O}
% variante possible : TeXGyrePagella ?
\setmathfont{XITS Math}

% maths
\usepackage{amsmath}
%\usepackage{amssymb}
%\usepackage{mathrsfs}
%\usepackage{stmaryrd}
%\usepackage{euscript}

% couleurs
\usepackage[usenames,dvipsnames]{xcolor}

% graphiques
\usepackage{graphicx}
\usepackage[all]{xy}
%\usepackage{tikz}
%\usetikzlibrary{matrix,arrows,calc}
%\tikzset{column sep=3em,row sep=2em,text height=1.5ex, text depth=0.25ex,>=angle 90}
%\tikzset{description/.style={fill=white,inner sep=1pt}}

% hyperliens, 'etiquettes
\usepackage{xr-hyper}
\usepackage{../fichiers-configuration-latex/smfthm}
%\ifx\cacheretiquettes\undefined
%\usepackage[inline]{showlabels} \else \fi % pour rendre visibles les 'etiquettes
%\usepackage{../fichiers-configuration-latex/makeidx} % pour faire un index
%\usepackage[pagebackref,unicode]{hyperref}  % pour voir o`u sont cit'ees les r'ef'erences
% hyperref is loaded once above with the cumulative-reader destination policy.
%\hypersetup{breaklinks=true} %permet le retour `a la ligne dans les liens trop longs. Incompatible avec pdflatex
%\renewcommand*{\backref}[1]{{\small [cit\'e page(s)~#1.]}}

\usepackage{xspace}
\usepackage{pifont} % pour expose XVII

% END INCLUDED SGA4_English_current_buildable_source_20260906/source/fichiers-configuration-latex/paquets.tex

% BEGIN INCLUDED SGA4_English_current_buildable_source_20260906/source/fichiers-configuration-latex/numerotation.tex
\renewcommand{\thefootnote}{\roman{footnote}} % pour mettre les notes en bas de page sous la forme (ii) et non (2), qui ressemble trop (?) a un exposant.

\newcounter{pageoriginal}
\newcommand{\pageoriginale}{\refstepcounter{pageoriginal}\marginpar{\footnotesize\textbf{\thepageoriginal}}\xspace}
\let\pageoriginaled\pageoriginale

\usepackage{ifthen}
\newboolean{original}
\def\sisi{\ifthenelse{\boolean{original}}}
\setboolean{original}{false}

\newcounter{cnote}
\def\nde#1{\ifthenelse{\boolean{original}}{}{\stepcounter{cnote}\def\thefootnote{\arabic{cnote}}\protect\footnote[\value{cnote}]{N.D.E.: #1}}}
\def\ndemark{\ifthenelse{\boolean{original}}{}{\stepcounter{cnote}\def\thefootnote{\arabic{cnote}}\protect\footnotemark[\value{cnote}]}}
\def\ndetext#1{\ifthenelse{\boolean{original}}{}{\def\thefootnote{\arabic{cnote}}\protect\footnotetext[\value{cnote}]{N.D.E.: #1}}}

% END INCLUDED SGA4_English_current_buildable_source_20260906/source/fichiers-configuration-latex/numerotation.tex

% BEGIN INCLUDED SGA4_English_current_buildable_source_20260906/source/fichiers-configuration-latex/macros.tex
\setmainlanguage{french}
\setotherlanguage{english}


%% Fixed, engine-independent release identifier.  The inherited dynamic Git
%% probe used Lua io.popen, exposed build-state-dependent output, and printed
%% a visible question mark outside a Git checkout.
\def\showgitstatus{English release r2, 2026-07-23}


% macros pour LaTeX
\renewcommand{\mathsf}{}

% Police différente pour U+2113 SCRIPT SMALL L
\setmathfont[range="2113]{latinmodern-math.otf}

% Give every inherited semantic label its own stable PDF destination instead
% of inheriting whichever automatic equation/page anchor happened to be live.
% The legacy label remains as a compatibility alias used by existing edges.
\makeatletter
\def\etiquette[#1]#2{%
  \Hy@raisedlink{\hyper@anchorstart{sga4.semantic.#2}\hyper@anchorend}%
  \def\@currentHref{sga4.semantic.#2}%
  \def\@currentlabel{#1}%
  \label{#2}%
}
\makeatother
\def\eqhyperref[#1]#2{(\hyperref[#1]{#2})}
\def\pagehyperref[#1]#2{(page \hyperref[#1]{#2})}
\def\hyperrefext[#1][#2]#3{\hyperref[#2]{#1 #3}}
\def\eqhyperrefext[#1][#2]#3{\hyperref[#2]{#1 #3}}     %cf. exposes 17
\def\neweqhyperrefext[#1][#2]#3{\hyperref[#2]{#1 #3}}  %cf. exposes 17

\def\shat#1{\protect\widehat{#1}}
\def\stilde#1{\protect\widetilde{#1}}
%%%%%%%%%%%%%%%%%%


\def\etcn{etc.\xspace}
\def\etco{etc...\xspace}
\def\etc{\sisi{\etco}{\etcn}}
\def\ldotsn{\sisi{...}{\etcn}}

\def\cf{{cf.}\kern.3em}
\def\Cf{{Cf.}\kern.3em}
\def\ie{{i.e.}\kern.3em}
\def\iev{{i.e.},\xspace}
\def\loccito{loc.\kern2pt cit.\xspace}
\def\loccitn{\textit{loc.\kern2pt cit.}\xspace}
\def\loccit{\sisi{\loccito}{\loccitn}}
\def\EGA{\textup{EGA}\kern.2em}
\def\SGA{\textup{SGA}\kern.3em}

\def\quoi{}

\def\quav{\quad,\quad}
\def\quaet{\quad \text{et}\quad}
\def\fibreproduct#1#2#3{#1 ×_{#3} #2}
\let\fprod\fibreproduct
\def\ccheck#1#2{\mbox{${\check{\check{\rm #1}}}^{\lower2pt\hbox{$\scriptstyle #2$}}$}}
\def\sheaf#1{\protect\mathcal{#1}}
\def\vieuxsheaf#1{{#1}} % cf. problème avec sheaf{a\Phi_i} par exemple...
\DeclareMathOperator{\Hom}{Hom}
\DeclareMathOperator{\can}{can}
\DeclareMathOperator{\Ind}{Ind}
\DeclareMathOperator{\ind}{ind}
%\DeclareMathOperator{\Hhomtop}{\mathcal{H}\!\mathit{omtop}}
\DeclareMathOperator{\Hhomtop}{\mathbf{Homtop}}
\DeclareMathOperator{\Point}{Point}
\DeclareMathOperator{\Pro}{Pro}
\DeclareMathOperator{\Top}{Top}
\DeclareMathOperator{\Mod}{Mod}
\DeclareMathOperator{\Cart}{Cart}
%\DeclareMathOperator{\Mmorsite}{\mathcal{M}\!\mathit{orsite}}
\DeclareMathOperator{\Mmorsite}{\mathbf{Morsite}}
\DeclareMathOperator{\loc}{loc}
\DeclareMathOperator{\Loc}{Loc}
\DeclareMathOperator{\Pic}{Pic}
\DeclareMathOperator{\Et}{Et}
\DeclareMathOperator{\card}{card}
\DeclareMathOperator{\Ens}{Ens}
\DeclareMathOperator{\h}{\underline{H}}
\DeclareMathOperator{\Fonct}{Fonct}
\DeclareMathOperator{\Ker}{Ker}
\DeclareMathOperator{\but}{but}
\DeclareMathOperator{\Ob}{Ob}
\DeclareMathOperator{\ob}{ob}
\DeclareMathOperator{\Lib}{Lib}
\DeclareMathOperator{\esj}{esj}
\DeclareMathOperator{\id}{id}
\DeclareMathOperator{\Id}{Id}
\DeclareMathOperator{\ad}{ad}
\DeclareMathOperator{\Iim}{Im}
\DeclareMathOperator{\pr}{pr}
\DeclareMathOperator{\Cat}{Cat}
\DeclareMathOperator{\car}{car}
\DeclareMathOperator{\SysInd}{SysInd}
\DeclareMathOperator{\Sup}{Sup}
\DeclareMathOperator{\Coker}{Coker}
\DeclareMathOperator{\Filt}{Filt}
\DeclareMathOperator{\acc}{acc}
\DeclareMathOperator{\sur}{sur}
\DeclareMathOperator{\cat}{cat}
\DeclareMathOperator{\Ord}{Ord}
\DeclareMathOperator{\Gr}{Gr}
\DeclareMathOperator{\Cov}{Cov}
\DeclareMathOperator{\Caf}{Caf}
\DeclareMathOperator{\coker}{coker}
\DeclareMathOperator{\coim}{coim}
\DeclareMathOperator{\im}{im}
\DeclareMathOperator{\Ouv}{Ouv}
\DeclareMathOperator{\diag}{diag}
\DeclareMathOperator{\Esp}{Esp}
\DeclareMathOperator{\TOP}{TOP}
\DeclareMathOperator{\Unif}{Unif}
\DeclareMathOperator{\Iind}{\mathbf{Ind}}
\DeclareMathOperator{\Hhomcart}{\mathbf{Homcart}}
\DeclareMathOperator{\Hhomtopess}{\mathbf{Homtopess}}
\DeclareMathOperator{\Ppoints}{\mathbf{Points}}
\DeclareMathOperator{\Ppoint}{\mathbf{Point}}
\DeclareMathOperator{\Ttors}{\mathbf{Tors}}
\DeclareMathOperator{\Ssex}{\mathbf{Sex}}
\DeclareMathOperator{\Ffib}{\mathbf{Fib}}
\DeclareMathOperator{\Vvois}{\mathbf{Vois}}
\DeclareMathOperator{\Ppointess}{\mathbf{Pointess}}
\DeclareMathOperator{\scrHom}{\mathbf{Hom}}
\DeclareMathOperator{\sheafPic}{\mathbf{Pic}}
\DeclareMathOperator{\sheafAut}{\mathbf{Aut}}
\DeclareMathOperator{\Hhomtoplocan}{\mathbf{Homtoplocan}}
\DeclareMathOperator{\Hhomtopan}{\mathbf{Homtopan}}
\DeclareMathOperator{\Pptgeom}{\mathbf{Ptgéom}}
\DeclareMathOperator{\Morsite}{Morsite}
\DeclareMathOperator{\Fib}{Fib}
\DeclareMathOperator{\Fais}{Fais}
\DeclareMathOperator{\kar}{kar}
\DeclareMathOperator{\Groupes}{Groupes}
\DeclareMathOperator{\Restr}{Restr}
\DeclareMathOperator{\Isom}{Isom}
\DeclareMathOperator{\Prol}{Prol}
\DeclareMathOperator{\Res}{Res}
\DeclareMathOperator{\Tors}{Tors}
\DeclareMathOperator{\esp}{esp}
\DeclareMathOperator{\es}{es}
\DeclareMathOperator{\od}{od}
\DeclareMathOperator{\gr}{gr}
\DeclareMathOperator{\sob}{sob}
\DeclareMathOperator{\dfn}{dfn}
\DeclareMathOperator{\Sch}{Sch}
\DeclareMathOperator{\ouv}{ouv}
\DeclareMathOperator{\cons}{cons}
\DeclareMathOperator{\Nerf}{Nerf}
\DeclareMathOperator{\Points}{Points}
\DeclareMathOperator{\Espsob}{Espsob}
\DeclareMathOperator{\Homtop}{Homtop}
\DeclareMathOperator{\Pt}{Pt}
\DeclareMathOperator{\Quot}{Quot}
\DeclareMathOperator{\glob}{glob}
\DeclareMathOperator{\Homfib}{Homfib}
\DeclareMathOperator{\Qouv}{Qouv}
\DeclareMathOperator{\Quotouv}{Quotouv}
\DeclareMathOperator{\ens}{ens}
\DeclareMathOperator{\ab}{ab}
\DeclareMathOperator{\an}{an}
\DeclareMathOperator{\fs}{fs}
\DeclareMathOperator{\red}{red}
\DeclareMathOperator{\Ab}{Ab}
\DeclareMathOperator{\cosq}{cosq}
\DeclareMathOperator{\Ann}{Ann}
\DeclareMathOperator{\Zar}{Zar}


%%%%%%%%%%%%
%%%%% SGA4-2
\DeclareMathOperator{\Sex}{Sex}
\DeclareMathOperator{\vois}{vois}
\DeclareMathOperator{\Vois}{Vois}
\DeclareMathOperator{\Ext}{Ext}
\DeclareMathOperator{\ext}{ext}
\DeclareMathOperator{\Fer}{Fer}
\DeclareMathOperator{\Add}{Add}
\DeclareMathOperator{\Cocart}{Cocart}
\DeclareMathOperator{\cosk}{cosk}
\DeclareMathOperator{\plat}{plat}

\DeclareMathOperator{\faisc}{faisc}
\DeclareMathOperator{\sk}{sk}
\DeclareMathOperator{\ch}{ch}

\def\HH{\mathbf{H}}
\def\qc{\mathrm{qc}}
\def\et{\mathrm{et}}
\def\coh{\mathrm{coh}}

\def\Hvr#1{%
\begin{array}[b]{l}
{\kern2pt\scriptstyle r}\\[-8pt]
{\kern1pt\scriptstyle \vee}\\[-5pt]
H^{#1}
\end{array}}

\def\Hvv#1{%
\begin{array}[b]{l}
\kern2pt\rotatebox{90}{${\scriptscriptstyle\ll}$}\\[-6pt]
H^{#1}
\end{array}}

\def\Hvi#1{%
\begin{array}[b]{l}
{\kern0.2pt\scriptscriptstyle \infty}\\[-8pt]
{\kern1pt\scriptstyle \vee}\\[-5pt]
H^{#1}
\end{array}}

\def\Hv#1{%
\begin{array}[b]{l}
{\kern1pt\scriptstyle \vee}\\[-5pt]
H^{#1}
\end{array}}

\DeclareMathOperator{\cart}{cart}
\DeclareMathOperator{\Ffaisc}{\mathbf{Faisc}}


\def\Ev#1{%
\begin{array}[b]{l}
{\kern1pt\scriptstyle \vee}\\[-5pt]
E'_{#1}
\end{array}}

\DeclareMathOperator{\Cont}{Cont}
\DeclareMathOperator{\Ccont}{\mathscr{C}\!\mathit{ont}}
\DeclareMathOperator{\Site}{Site}
\newcommand\MM{\mathbf{M}}
\def\limtop#1{{\displaystyle{\varprojlim_{#1}\!\mbox{\rm top }}}}
\def\Limtop{\mathop{\varprojlim\!\mbox{\rm top}}}
\DeclareMathOperator{\spec}{spec}
\def\ss{\scriptstyle{\mathscr{S}}}
\DeclareMathOperator{\Gal}{Gal}
\def\smallm{\scriptstyle{\mathscr{M}}}
\DeclareMathOperator{\End}{End}
\DeclareMathOperator{\ass}{ass}
\DeclareMathOperator{\zar}{zar}
\DeclareMathOperator{\Ogg}{Ogg}
\DeclareMathOperator{\res}{res}
\DeclareMathOperator{\tr}{tr}
\DeclareMathOperator{\Supp}{Supp}
\DeclareMathOperator{\codim}{codim}
\DeclareMathOperator{\rk}{rk}
\DeclareMathOperator{\cur}{cur}
\DeclareMathOperator{\cdqc}{cdqc}
\newcommand\DD{\mathbf{D}}
\DeclareMathOperator{\cl}{cl}
\def\JJ{\mathbf{J}}
\def\FF{\mathbf{F}}
\def\BB{\mathbf{B}}
\DeclareMathOperator{\Alget}{Alget}
\DeclareMathOperator{\Ass}{Ass}
\DeclareMathOperator{\supp}{supp}
\DeclareMathOperator{\rad}{rad}
\DeclareMathOperator{\prof}{prof}
\DeclareMathOperator{\cint}{int}
\DeclareMathOperator{\sic}{sic}
\DeclareMathOperator{\Tor}{Tor}
\DeclareMathOperator{\tors}{tors}
\DeclareMathOperator{\Sym}{Sym}
\DeclareMathOperator{\parf}{parf}
\DeclareMathOperator{\tor}{tor}
\DeclareMathOperator{\defn}{defn}
\DeclareMathOperator{\Div}{Div}
\DeclareMathOperator{\Pol}{Pol}
\DeclareMathOperator{\Fex}{Fex}
\DeclareMathOperator{\Lie}{Lie}
\DeclareMathOperator{\tot}{tot}
\DeclareMathOperator{\Aut}{Aut}
\DeclareMathOperator{\HOM}{HOM}
\DeclareMathOperator{\pch}{pch}
\DeclareMathOperator{\Ch}{Ch}
\DeclareMathOperator{\EXT}{EXT}
\DeclareMathOperator{\PIC}{PIC}
\DeclareMathOperator{\Rf}{Rf}
\DeclareMathOperator{\inj}{inj}
\newcommand\VV{\mathbf{V}}

\def\Xv{%
\begin{array}[b]{l}
{\kern1pt\scriptstyle \vee}\\[-5pt]
X
\end{array}}

%%%%%%%
\newcommand\TT{\mathbf{T}}
\newcommand\XX{\mathbf{X}}
\newcommand\EE{\mathbf{E}}
%%%%%%

\def\qedsymbol{\textup{C.Q.F.D.}}
\let\CQFD\qed
\let\cqfd\qed

%%%%%\bigdot
\font\bigsymb=cmsy10 at 4pt
\def\bigdot{{\kern1.2pt\raise 1.5pt\hbox{\bigsymb\char15}}}
%%%%%

\def\Lright{\mathop{\vcenter{\hbox{\hbox{$\longrightarrow$}\lower8pt \hbox{\kern-16pt\hbox{$\longrightarrow$}}}}}}

\def\Aa{{\mathcal A}}
\def\Bb{{\mathcal B}}
\def\Cc{{\mathcal C}}
\def\Dd{{\mathcal D}}
\def\Ee{{\mathcal E}}
\def\Ff{{\mathcal F}}
\def\Gg{{\mathcal G}}
\def\Hh{{\mathcal H}}
\def\Ii{{\mathcal I}}
\def\Jj{{\mathcal J}}
\def\Kk{{\mathcal K}}
\def\Ll{{\mathcal L}}
\def\Mm{{\mathcal M}}
\def\Nn{{\mathcal N}}
\def\Oo{{\mathcal O}}
\def\Pp{{\mathcal P}}
\def\Qq{{\mathcal Q}}
\def\Rr{{\mathcal R}}
\def\Ss{{\mathcal S}}
\def\Tt{{\mathcal T}}
\def\Uu{{\mathcal U}}
\def\Vv{{\mathcal V}}
\def\Ww{{\mathcal W}}
\def\Xx{{\mathcal X}}
\def\Yy{{\mathcal Y}}
\def\Zz{{\mathcal Z}}

\newcommand{\CC}{\mathbf{C}}
\newcommand{\GG}{\mathbf{G}}
\newcommand{\KK}{\mathbf{K}}
\newcommand{\LL}{\mathbf{L}}
\newcommand{\NN}{\mathbf{N}}
\newcommand{\PP}{\mathbf{P}}
\newcommand\QQ{\mathbf{Q}}
\newcommand\RR{\mathbf{R}}
\newcommand\ZZ{\mathbf{Z}}

\def\to{\mathchoice{\longrightarrow}{\rightarrow}{\rightarrow}{\rightarrow}}
\def\from{\mathchoice{\longleftarrow}{\leftarrow}{\leftarrow}{\leftarrow}}
\def\mto{\mathchoice{\longmapsto}{\mapsto}{\mapsto}{\mapsto}}
\def\mfrom{\mathrel{\from\joinrel\mapstochar\,}}
\def\hto{\mathrel{\lhook\joinrel\to}}
\def\hfrom{\mathrel{\from\joinrel\rhook}}

\def\implies{\mathchoice{\Longrightarrow}{\Rightarrow}{\Rightarrow}{\Rightarrow}}
\def\implied{\mathchoice{\Longleftarrow}{\Leftarrow}{\Leftarrow}{\Leftarrow}}
\def\iff{\mathchoice{\Longleftrightarrow}{\Leftrightarrow}{\Leftrightarrow}{\Leftrightarrow}}

\def\To#1{\mathchoice{\xrightarrow{\labelstyle~~#1~}}{\stackrel{#1}{\longrightarrow}}{}{}}
\def\Mto#1{\mathrel{\mapstochar\To{#1}}}
\def\From#1{\mathchoice{\xleftarrow{\labelstyle~~#1~}}{\stackrel{#1}{\longleftarrow}}{}{}}

\newcommand{\isomto}{\overset{\sim}{\longrightarrow}}
\newcommand{\isomfrom}{\overset{\sim}{\longleftarrow}}

\newcommand{\IFF}{\unskip~$\iff$~\ignorespaces}
\newcommand{\THEN}{\unskip~$\implies$~\ignorespaces}

\let\leq\leqslant
\let\geq\geqslant
\let\le\leqslant
\let\ge\geqslant

\def\cf{{cf.}\kern.3em}
\def\Cf{{Cf.}\kern.3em}
\def\ie{{i.e.}\kern.3em}
\def\iev{{i.e.},\xspace}
%\def\loccito{loc.\kern2pt cit.\xspace}
%\def\loccitn{\textit{loc.\kern2pt cit.}\xspace}
%\def\loccit{\sisi{\loccito}{\loccitn}}
\def\EGA{\textup{EGA}\kern.2em}
\def\SGA{\textup{SGA}\kern.3em}

\let\oldS\S
\def\S{\oldS\kern .15em}
\def\ptbl{.\kern .2em}
\def\resp{resp.\kern.3em}
\def\Exp{Exp.\kern.3em}

\let\bar\overline
\let\hat\widehat
\let\tilde\widetilde
\let\dpl\displaystyle
\let\ldots\dots
\let\emptyset\varnothing
\let\skpt\skippointrait
\let\epsilon\varepsilon

\def\sfrac#1#2{{#1}/{#2}}

\DeclareMathOperator{\cd}{cd}
%\DeclareMathOperator{\Hhom}{\mathcal{H}\!\mathit{om}}
\DeclareMathOperator{\Hhom}{\mathbf{Hom}}
%\DeclareMathOperator{\Eext}{\mathcal{E}\!\mathit{xt}}
\DeclareMathOperator{\Eext}{\mathbf{Ext}}
\DeclareMathOperator{\Spec}{Spec}
\DeclareMathOperator{\Tr}{Tr}

\let\emptyset=\varnothing

%%%%%%%%%%%%%%%
%%% expose XVII
%%%%%%%%%%%%%%%

\makeatletter
\let\old@optexpose\optexpose
\newcommand{\empty@optexpose}[1]{}
\newcommand{\paren@optexpose}[1]{\old@optexpose{#1}(}
\newcommand{\refnoexpose}[1]{\let\optexpose\empty@optexpose\ref{#1}\let\optexpose\old@optexpose}%\refnoexpose{X.cor5.2} engendre "5.2" au lieu de "X 5.2"
\newcommand{\neweqref}[1]{\let\optexpose\paren@optexpose\textup{\ref{#1})}\let\optexpose\old@optexpose}%\neweqref{XII.eq4.1} engendre "XII (4.1)" au lieu de "(XII 4.1)"
\makeatother
\newcommand{\refvb}[2][VI${}_{\text{B}}$ ]{\sisi{#1\refnoexpose{#2}}{\ref{#2}}}

\newcommand{\leftexp}[2]{{\vphantom{#2}}^{#1}{#2}}
\newcommand{\AAA}{\mathbf{A}}
\newcommand{\Sf}{\mathfrak{S}}
\DeclareMathOperator*{\quotprojlim}{\text{``$\varprojlim$''}}
\DeclareMathOperator*{\quotinjlim}{\text{``$\varinjlim$''}}
\DeclareMathOperator{\Fl}{F\sisi{\ell}{l}}
\DeclareMathOperator{\Pl}{P\sisi{\ell}{l}}
%\DeclareMathOperator{\cl}{c\sisi{\ell}{l}}
\DeclareMathOperator{\tdf}{tdf}
\DeclareMathOperator{\rg}{rg}
\newcommand{\iInd}{\sisi{I}{i}nd}
\newcommand{\iind}{ind}
\newcommand{\ppro}{pro}
\newcommand{\tTor}{\sisi{T}{t}or}
\newcommand{\ttor}{tor}
\newcommand{\chap}{Chap.}
\newcommand{\Ffex}{\Ff\mathit{ex}}
\newcommand{\sheafk}{\underline{k}}
\newcommand{\sheafn}{\underline{n}}
\newcommand{\sheafGamma}{\underline{\Gamma}}

% Repeated references to an already typeset footnote must not use
% \footnotemark without a matching \footnotetext: Hyperref would then create
% a live annotation to a destination that cannot exist.  This helper keeps the
% current printed footnote mark and links it to an explicit stable note target.
\makeatletter
\newcommand{\sgarepeatedfootnotemark}[1]{%
  \begingroup
  \edef\@thefnmark{\thefootnote}%
  \hyperlink{#1}{\@makefnmark}%
  \endgroup}
\makeatother

\sisi{
\DeclareMathOperator{\simarrow}{\rule{.15mm}{1.8mm}\kern -3.4pt\rightsquigarrow}
\newcommand{\sheafell}{\underline{\ell}}
\newcommand{\sheafZZ}{\underline{\ZZ}}
}{
\newcommand{\nderef}[1]{la N.D.E.~(\ref{#1}) page~\pageref{#1}}
}

%%%%%%%%%%%%%%%%%
%%%%% expose XVII
%%%%%%%%%%%%%%%%%

% END INCLUDED SGA4_English_current_buildable_source_20260906/source/fichiers-configuration-latex/macros.tex


\externaldocument{../01/01_English.tex}
\externaldocument{../02/02_English.tex}
\externaldocument{../03/03_English.tex}
\externaldocument{../04/04_English.tex}
\externaldocument{../05/05_English.tex}
\externaldocument{../05bis/05bis_English.tex}
\externaldocument{../06/06_English.tex}
\externaldocument{../07/07_English.tex}
\externaldocument{../08/08_English.tex}
\externaldocument{../09/09_English.tex}
%\externaldocument{../10/10_English.tex}
\externaldocument{../11/11_English.tex}
\externaldocument{../12/12_English.tex}
\externaldocument{../13/13_English.tex}
\externaldocument{../14/14_English.tex}
\externaldocument{../15/15_English.tex}
\externaldocument{../16/16_English.tex}

\setboolean{original}{false}

\begin{document}

\begin{center}
XVII. Cohomology with proper supports\\
P. Deligne
\end{center}

\selectlanguage{english}version: \showgitstatus
\tableofcontents
\mainmatter

\else
\chapter{Cohomology with proper supports}\MakeLinkTarget*{sga4.exp.XVII}\label{sga4.exp.XVII}
%\addtocontents{toc}{P. Deligne}
\begin{center}
P. Deligne
\end{center}
\fi

% The edition is fixed to the corrected (second) reading.  Keeping the
% generic ifthen-based selector inside Xy-pic matrices breaks current
% Xy-pic; this equivalent expandable selector avoids that toolchain bug.
\renewcommand{\sisi}[2]{#2}

% BEGIN INCLUDED SGA4_English_current_buildable_source_20260906/source/17/drafts/SGA4_XVII_LINES_44_260_English.texfrag
\setcounter{pageoriginal}{251}

\footnotetext[1]{The present exposé and the following one, written in
  1968 and 1969, reproduce and supplement the oral lectures of
  A.~\textsc{Grothendieck} (spring 1964). The author of this written
  version, who did not attend the oral seminar, drew in part on notes
  by A.~\textsc{Grothendieck}.}


\section*{Introduction}\pageoriginale

In this exposé we develop the formalism of cohomology with proper
support. The questions of variance have been treated with considerable
care, whereas they were only made plausible in the oral seminar. This
explains the length of the exposé, whose Sections
\hyperref[XVII.sec1]{1} to \hyperref[XVII.sec4]{4} are devoted to
categories and general topology. The reader is very strongly advised
to read only Sections \hyperref[XVII.sec5]{5} and
\hyperref[XVII.sec6]{6}, where the geometric substance of the exposé
is concentrated (\S \hyperref[XVII.sec5]{5}: construction and
variance of cohomology with proper supports, base-change, finiteness,
finite \tTor-dimension, and comparison theorems; \S
\hyperref[XVII.sec6]{6}: theory of the \sisi{trace}{trace} morphism for
a quasi-finite flat morphism).

In Sections \hyperref[XVII.sec1]{1} and \hyperref[XVII.sec2]{2}, we
``recall'' some results on derived categories. The \S
\hyperref[XVII.sec1]{1} ${\rm n^{\circ}}$ 1 is devoted to the sign
formalism. In the \S \hyperref[XVII.sec3]{3}, we treat the problem of
gluing two variance formalisms. In the \S \hyperref[XVII.sec4]{4}, we
introduce flat resolutions and study the special properties of
\textsc{Godement}'s flasque resolutions.

In\pageoriginale this exposé, the higher direct-image functors with
proper supports $R^{q}f_{!}$, and the functor $Rf_{!}$ from which they
arise, are defined only when $f$ is a compactifiable morphism
(\hyperref[XVII.def3.2.1]{3.2.1}). In the appendix, written by
\textsc{B. Saint-Donat}, it is shown how to extend the definition to
separated morphisms of finite type with quasi-compact quasi-separated
target.\nde{\MakeLinkTarget*{sga4.exp.XVII.editor-note.introduction.nagata-compactification}\label{XVII.ndeConrad}The Nagata compactification theorem
ensures that every separated morphism of finite type with
quasi-compact quasi-separated target is compactifiable. See B.~Conrad,
``Deligne's notes on Nagata compactifications'', \emph{J. Ramanujan
Math.\ Soc.} \textbf{22} (2007), 205--257, cited as [C] below, \hyperref[XVII.subsec4.1]{4.1}, as
well as ``Erratum for `Deligne's notes on Nagata compactifications''',
\emph{J. Ramanujan Math.\ Soc.} \textbf{24} (2009), 427--428.}

The \S \hyperref[XVII.sec5]{5} ${\rm n^{\circ}}$
\hyperref[XVII.sec5]{5} (cohomology of a symmetric product) will not be
used again in this seminar. It will be used in \SGA 5 to refine the
rationality theorem for $L$-functions.\nde{See \SGA
4$\frac{1}{2}$, $L$-functions modulo $\ell^n$ and modulo $p$.}

The \S \hyperref[XVII.sec6]{6} ${\rm n^{\circ}}$
\hyperref[XVII.sec3]{3} (the theory of the trace for \emph{continuous}
coefficients) will be used in Exposé XVIII only in the relatively easy
case of smooth groups; a reader interested in the Poincaré duality
theorem (global duality), and willing to accept a transcendental
argument, may even dispense entirely with reading this \S
\hyperref[XVII.sec6]{6} ${\rm n^{\circ}}$
\hyperref[XVII.sec3]{3}, as well as most of the \S
\refnoexpose{XVIII.sec1} of XVIII.

\setcounter{section}{-1}
\section{Terminological preliminaries}\etiquette[0]{XVII.sec0}

\subsection{}\etiquette[0.1]{XVII.subsec0.1}
The sign $=$ placed between two groups of symbols denoting objects of a
category will sometimes mean (by abuse of notation) that these objects
are canonically isomorphic. The category and the canonical isomorphism
should in principle have been defined beforehand. In a diagram, the
sign $=$ will then denote the isomorphism itself.

\hypertarget{sga4.exp.XVII.subsection.0.2}{}\setcounter{subsection}{2}
\setcounter{subsubsection}{-1}
\subsubsection{}\etiquette[0.2.0]{XVII.subsubsec0.2.0}
Let $f:S\rightarrow S'$ be a morphism of sites
(\hyperrefext[IV][IV.subsubsec4.9.3]{4.9.3}). If $U$ and $U'$ are
objects of $S$ and $S'$, an $f$-morphism from $U$ to $U'$ will by
definition be a morphism from $U$ to $f^{\ast}U'$. For the étale sites
of schemes, this recovers the usual notion.

\subsubsection{}\etiquette[0.2.1]{XVII.subsubsec0.2.1}
If\pageoriginale $\mathcal{F}$ and $\mathcal{F}'$ are sheaves on $S$
and $S'$, an $f$-morphism from $\mathcal{F}$ to $\mathcal{F}'$ will
mean, equivalently,
\begin{enumerate}
\renewcommand{\theenumi}{\roman{enumi}}
\renewcommand{\labelenumi}{(\theenumi)}
\item a morphism from $f^{\ast}\mathcal{F}'$ to $\mathcal{F}$;

\item a morphism from $\mathcal{F}'$ to $f_{\ast}\mathcal{F}$;

\item a function which, to every $f$-morphism $\varphi$ from an object
  $U$ of $S$ to an object $U'$ of $S'$, assigns a function from
  $\mathcal{F}'(U')$ to $\mathcal{F}(U)$, compatibly with the
  composition of $\varphi$ with an arrow of $S$ or of $S'$
  \footnote{These terminologies are not compatible with identifying
  the objects of $S$ with the associated sheaves.}.
\end{enumerate}

It follows from (i), (ii), and (iii) that sheaves of sets form a
fibered and cofibered category over the category of sites.

Likewise for sheaves of left modules on ringed sites. Likewise for
étale sheaves on schemes; this is not immediately a special case of
the preceding statement, since sites in fact form a 2-category and
``the étale site of $X$'' is only a pseudofunctor in $X$
(\hyperrefext[VII][VII.subsec1.4]{1.4}).

The preceding definition makes sheaves of sets (\resp\ldots) on
varying sites into a fibered category over the category of sites,
whose fiber categories are the \emph{opposite categories} of the
usual categories of sheaves of sets (\resp\ldots).

\subsection{}\etiquette[0.3]{XVII.subsec0.3}
Let $I$ be a finite set, let $\epsilon$ be a function on $I$ with
values in $\{+1,-1\}$, let $(\Aa_i)_{i\in I}$ be a family of additive
categories graded by translation functors $T_i$ (\cite{sga4.exp.XVII.bib.CD} I
\sisi{1.10}{1.1.0}), and let $F$ be a multiadditive multifunctor from
the categories $\Aa_i$ to an additive category $\Aa$ graded by the
translation functor $T$. We suppose that $F$ is covariant (\resp
contravariant) in the variables $i$ for which $\epsilon_i=1$ (\resp
$\epsilon_i=-1$). Modifying and supplementing\pageoriginale
(\cite{sga4.exp.XVII.bib.CD} \sisi{I 1.1}{II 2.1.1}), we shall say that $F$ is a
\emph{graded functor} if we are given a family
$(\varphi_i)_{i\in I}$ of isomorphisms of functors between $T\circ F$
and $F\circ T_i^{\varepsilon_{i}}$, such that the following diagrams are
anticommutative for $i\neq j$ in $I$:
\[
\xymatrix{
F \circ (T_i^{\varepsilon_i}, T_j^{\varepsilon_j})  \ar[rr]^{\varphi_j
  \ast T_i^{\varepsilon_i}} \ar[dd]^{\varphi_i\ast T_j^{\varepsilon_j}}& &T\circ F \circ T_i^{\varepsilon_i} \ar[dd]^{T\ast
  \varphi_i}\\
 & - & \\
 T\circ F \circ T_j^{\varepsilon_j}  \ar[rr]^{T\ast \varphi_j} & &T\circ T
  \circ F.
}
\]
We leave it to the reader to define the composite of two graded
functors and verify that it is again a graded functor.

\subsection{}\etiquette[0.4]{XVII.subsec0.4}
Let $\Aa$ be an additive category and let $I$ be a finite set. For
$i\in I$, denote by $1_i$ the $i$\textsuperscript{th} basis vector of
$\ZZ^{I}$. An \emph{$I$-fold complex} (\resp a
\sisi{naive}{\emph{naive}} \emph{$I$-fold complex}) of $\Aa$ consists of
\begin{enumerate}
\renewcommand{\theenumi}{\alph{enumi}}
\renewcommand{\labelenumi}{(\theenumi)}
\item a family $(K^{k})_{k\in \ZZ^{I}}$ of objects of $\Aa$;

\item for every $i\in I$ and every $k\in \ZZ^{I}$, an arrow
  $d^{k}_i:K^{k}\rightarrow K^{k+1_i}$, these arrows satisfying
  $d^{k+1_i}_i d^{k}_i=0$ and, for $i\neq j$,
  $d^{k+1_j}_i d^{k}_j+d^{k+1_i}_j d^{k}_i=0$ (\resp
  $d^{k+1_j}_i d^{k}_j=d^{k+1_i}_j d^{k}_i$).
\end{enumerate}

\subsection{}\etiquette[0.5]{XVII.subsec0.5}
A\pageoriginale \emph{\ppro-object} of a category $\sisi{C}{\Cc}$ is a functor from
$\sisi{C}{\Cc}$ to $(\Ens)$ which is a filtered direct limit (over a
small filtered category) of representable functors (\cf
\hyperrefext[I][I.sec8.10]{8.10}). Every \ppro-object is a direct limit,
over a small filtered ordered set, of representable functors. If
$X_{i}$ is an inverse system of objects of $\sisi{C}{\Cc}$, indexed by a
small filtered category, we denote by $\quotprojlim X_{i}$ the
\ppro-object $\varinjlim h^{X_i}$.

\subsection{}\etiquette[0.6]{XVII.subsec0.6}
The reader may dualize \hyperref[XVII.subsec0.5]{0.5} to the case of
\iind-objects (\Cf \hyperrefext[I][I.sec8.2]{8.2}).

\subsection{}\etiquette[0.7]{XVII.subsec0.7}
In accordance with the new terminology, what was formerly called a
prescheme is now called a \emph{scheme}, and what was formerly called
a scheme is now called a \emph{separated scheme}.

\subsection{}\etiquette[0.8]{XVII.subsec0.8}
The category of \emph{ringed schemes} is the category whose objects are
schemes whose étale site is endowed with a sheaf of rings, an arrow
from $(S,\Aa)$ to $(T,\Bb)$ being a pair consisting of a morphism of
schemes $f$ from $S$ to $T$ and a homomorphism $\varphi$ of sheaves of
rings from $f^{\ast}\Bb$ to $\Aa$. The morphism of schemes $f$ is
said to be \emph{induced by} $(f,\varphi)$.

\subsection{}\etiquette[0.9]{XVII.subsec0.9}
If $X$ is a scheme over $Y$, we denote by $(X/Y)^n$ the $n$-fold fiber
product of $X$ over $Y$.

\subsection{}\etiquette[0.10]{XVII.subsec0.10}
Let $S$ be a scheme. The \emph{big étale site} (\resp the \emph{big
fppf site}, \resp the \emph{big fpqc site}) of $S$ is the site
$\Sch/{S}$, endowed with the étale topology (\resp the fppf topology,
\resp the fpqc topology) (\SGA 3 IV 6.3). Note carefully that this is
not a $\sisi{U}{\Uu}$-site ($\sisi{U}{\Uu}$ being the
\sisi{finite}{fixed} universe, and $\Sch/{S}$ consisting of the schemes
$\in \sisi{U}{\Uu}$).

The\pageoriginale \emph{small fpqc site} (\resp the \emph{small fppf
site}) of $S$ is the site of schemes flat over $S$ (\resp flat and of
finite presentation), endowed with the fpqc topology (\resp the fppf
topology). When it is necessary to avoid ambiguity, the \emph{small
étale site} of $S$ will mean the site $S_{\et}$
(\hyperrefext[VII][VII.subsec1.2]{1.2}).

\subsection{}\etiquette[0.11]{XVII.subsec0.11}
In many cases, the author has allowed himself to speak of commutative
diagrams of morphisms of sites where one should have spoken of
essentially commutative diagrams, \ie diagrams commutative up to
isomorphism (\cf IV \hyperref[IV.subsubsec3.2.2]{3.2.2}). The reader may verify that the arguments we
give also apply to the general situation.

\subsection{}\etiquette[0.12]{XVII.subsec0.12}
The terminology ``\emph{coherent scheme}'' for ``quasi-compact
quasi-separated scheme'' has been surreptitiously introduced here and
there by A.~Grothendieck.

\subsection{}\etiquette[0.13]{XVII.subsec0.13}
A torsion sheaf $F$ on a scheme $S$ will be said to be \emph{prime to
the residual characteristics of} $S$ if it is the direct limit of its
subsheaves annihilated by integers $n$ invertible on $S$.

% END INCLUDED SGA4_English_current_buildable_source_20260906/source/17/drafts/SGA4_XVII_LINES_44_260_English.texfrag

% BEGIN INCLUDED SGA4_English_current_buildable_source_20260906/source/17/drafts/SGA4_XVII_LINES_261_600_English.texfrag
\section{Derived categories}\etiquette[1]{XVII.sec1}

\subsection{Exact functors (the sign rules)}\etiquette[1.1]{XVII.subsec1.1}

\subsubsection{}\etiquette[1.1.1]{XVII.subsubsec1.1.1}
For the fundamental theorems concerning derived categories, I refer
to \textsc{Verdier} \cite{sga4.exp.XVII.bib.CD} and \cite{sga4.exp.XVII.bib.V2}. In this subsection,
particularly careful attention has been paid to questions of signs.

Recall that if $f:X\rightarrow Y$ is a morphism of complexes
(in an additive category, which will henceforth be understood), its
\emph{cone} $C(f)$\pageoriginale is defined by
$$\etiquette[1.1.1.1]{XVII.eq1.1.1.1}
C(f)^{n}=X^{n+1}\oplus Y^{n}\quad d^{n}=-d^{n+1}_X+f^{n+1}+d^{n}_Y{\quoi}.\leqno(1.1.1.1)
$$
When $X=0$ (\resp $Y=0$), we have $C(f)=Y$ (\resp $C(f)=X[1]$), so
that the commutative diagram
\[
\xymatrix{
0        \ar@{=}[r] \ar[d]&   0   \sisi{\ar@{=}[r]}{\ar[r]} \ar[d]&  X
\ar@{=}[r]\ar[d]^{f}  & X \ar[d]\\
X        \ar[r]^{f}   &  Y    \ar@{=}[r] & Y     \ar[r]        & 0
}
\]
defines a ``triangle''
\[
X\xrightarrow{f}Y\rightarrow C(f)\xrightarrow{j}X[1]{\quoi}.
\]

In the category of complexes up to homotopy, a triangle is called
\emph{distinguished} if it is isomorphic to a triangle of this type,
and \emph{antidistinguished} if it becomes distinguished when the
signs of its arrows are changed. One checks that the triangle defined
by an exact sequence of complexes which is split in each degree
(\cite{sga4.exp.XVII.bib.CD} p.~10, \cf the proof of
\hyperref[XVII.lem1.1.5.4]{1.1.5.4}) is antidistinguished and that, up
to isomorphism, every antidistinguished triangle is obtained in this
way.

\subsubsection*{${}^{\ast}$1.1.2.}\etiquette[*1.1.2]{XVII.subsubsec*1.1.2}
Sign calculations are most easily carried out using the following
``formal rules''. Write an element of $C(f)^{n}$ in the form
$1\otimes x+y$ $(x\in X^{n+1}, y\in Y^{n})$, where $1$ is a
``cell'' of dimension $1$, and set $d(1)\otimes x=f(x)$. Thus
\[
d(1\otimes x+y)=d1\otimes x-1\otimes dx+dy=-1\otimes
dx+f(x)+dy{\quoi}._{\ast}
\]

\begin{enonce*}{Definition 1.1.3}\etiquette[1.1.3]{XVII.def1.1.3}
\begin{enumerate}
\renewcommand{\theenumi}{\roman{enumi}}
\renewcommand{\labelenumi}{(\theenumi)}
\item A\pageoriginale graded functor (\hyperref[XVII.subsec0.3]{0.3})
  from a triangulated category $\Aa$ to a triangulated category $\Bb$
  is said to be exact if it sends distinguished triangles to
  distinguished triangles.

\item A contravariant graded functor from a triangulated category
  $\Aa$ to a triangulated category $\Bb$ is said to be
  \emph{exact} if, for every distinguished triangle
  $(X,Y,Z,u,v,w)$ of $\Aa$, the triangle
  $(F(Z),F(Y), F(X), F(v),F(u),TF(w))$ is distinguished (in this
  definition, $FX$ is identified with $TFTX$ by means of the grading
  of $F$).

\item Let $(\Aa_i)_{i\in I}$ be a finite family of triangulated
  categories, let $\epsilon$ be a function from $I$ to
  $\{+1,-1\}$, and let $\Aa$ be a triangulated category. A graded
  functor $F$ from the product of the $\Aa_i$ to $\Aa$, covariant
  in those $i$ for which $\epsilon(i)=1$ and contravariant in those
  for which $\epsilon(i)=-1$, is said to be \emph{exact} if, for every
  $i\in I$, the functors from $\Aa_i$ to $\Aa$ obtained from $F$ by
  fixing all variables except the $i$-th are exact functors.
\end{enumerate}
\end{enonce*}

A composite of exact functors is again an exact functor.

\setcounter{subsubsection}{3}
\subsubsection{}\etiquette[1.1.4]{XVII.subsubsec1.1.4}
If $(K^{n_1\ldots n_p}, d_1,\ldots,d_p)$ is a multicomplex
(\hyperref[XVII.subsec0.4]{0.4}), its associated simple complex is
defined by
$$\etiquette[1.1.4.0]{XVII.eq1.1.4.0}
K^{n}=\prod_{\sum n_i=n}K^{(n_i)}\quav d=\sum d_i{\quoi}.\leqno(1.1.4.0)
$$

When, exceptionally, it is defined using a direct sum rather than a
product, this fact will be stated explicitly.

Let $I$ be a finite set. Following \textsc{Cartan--Eilenberg} and
\textsc{J.~P. Serre}, we shall explain how every naive $I$-fold
complex $K$ has canonically\pageoriginale associated with it an
$I$-fold complex (\hyperref[XVII.subsec0.4]{0.4}).

For every total order $<$ on $I$, let $K(<)$ be the following
$I$-fold complex:
$$\etiquette[1.1.4.1]{XVII.eq1.1.4.1}
\begin{cases}
K(<)^{\sheafk }=K^{\sheafk }\\
d(<)_i^{\sheafk }=(-1)^{\sum\limits_{j<i} k_j} d_i{\quoi}.
\end{cases}\leqno{(1.1.4.1)}
$$

If $<_1$ and $<_2$ are two \sisi{total}{total} orders on $I$, the
\emph{canonical isomorphism $\tau$ between $K(<_1)$ and $K(<_2)$} is,
by definition, the one given by
$$\etiquette[1.1.4.2]{XVII.eq1.1.4.2}
\tau^{\sheafk }=(-1)^{\epsilon(<_1,<_2)},\quad\mbox{\rm where}\quad
\epsilon(<_1,<_2)=\sum_{\substack{i<_1j\\ \sisi{i<_2 j}{j<_2 i}}}k_ik_j{\quoi}.\leqno{(1.1.4.2)}
$$

\subsubsection*{1.1.4.2.1}\etiquette[1.1.4.2.1]{XVII.par1.1.4.2.1}
One checks that these canonical isomorphisms form a transitive system
of isomorphisms among the $K(<)$, as $<$ ranges over the total orders
on $I$; these $K(<)$ therefore determine an $I$-fold complex, unique
up to a unique isomorphism, which in this subsection will be denoted by
$K(\ast)$ and is called the \emph{$I$-fold complex associated with the
naive $I$-fold complex $K$}. Note that $K(\ast)^{\sheafk}$ is
\emph{not canonically isomorphic to $K^{\sheafk}$: only the choice of
a total order on $I$ makes it possible to identify these two objects.}

If $K$ is a naive $I$-fold complex and $i\in I$, let $K[1_i]$ denote
the following naive complex, whose differentials are denoted by
$d_j[1_i]$:
$$\etiquette[1.1.4.3]{XVII.eq1.1.4.3}
\begin{cases}
K[1_i]^{\sheafk } = K^{\sheafk +1_i}\\
d_j[1_i]^{\sheafk } = d_j^{\sheafk +1_{i}}\quad{\mbox{\rm if }} i\neq
j\\
d_i[1_i]^{\sheafk }= -d_i^{\sheafk +1_i}{\quoi}.
\end{cases}\leqno{(1.1.4.3)}
$$

If\pageoriginale $K$ is an $I$-fold complex, set
$$\etiquette[1.1.4.4]{XVII.eq1.1.4.4}
\begin{cases}
K[1_i]^{\sheafk } = K^{\sheafk +1_i}\\
d_j[1_i]^{\sheafk } = -d_j[1_i]^{\sheafk +1_i}\sisi{}{.}
\end{cases}\leqno{(1.1.4.4)}
$$

If $K$ is a naive $I$-fold complex and $<$ is a total order on $I$,
so that $K(\ast)$ is given by \eqhyperref[XVII.eq1.1.4.1]{1.1.4.1},
define an isomorphism $\sigma$ between $K[1_i](\ast)$ and
$K(\ast)[1_i]$ by the formula
$$\etiquette[1.1.4.5]{XVII.eq1.1.4.5}
\sigma^{\sheafk }=(-1)^{\sum\limits_{j<i}k_j} \id_{K^{\sheafk +1_i}}{\quoi}.\leqno{(1.1.4.5)}
$$

\setcounter{subsubsection}{4}
\setcounter{paragraph}{5}
\paragraph{}\etiquette[1.1.4.6]{XVII.par1.1.4.6}
One checks that this isomorphism is independent of the choice of the
total order $<$; it therefore defines an \emph{isomorphism, called the
canonical isomorphism, between} $K[1_i](\ast)$ and $K(\ast)[1_i]$.

\paragraph{}\etiquette[1.1.4.7]{XVII.par1.1.4.7}
For any $i,j\in I$, if $K$ is an $I$-fold complex or a naive $I$-fold
complex, we identify $K[1_i][1_j]$ and $K[1_j][1_i]$ in the evident
way.

For $i\neq j$ in $I$, the following diagram of canonical
isomorphisms of type \hyperref[XVII.par1.1.4.6]{1.1.4.6} or
\hyperref[XVII.par1.1.4.7]{1.1.4.7} is then \emph{anticommutative}:
$$\etiquette[1.1.4.7.1]{XVII.eq1.1.4.7.1}
\vcenter{
\xymatrix{
K[1_i][1_j](\ast) \ar@{=}[r] \ar@{<->}[d]& K[1_j][1_i](\ast)
\ar@{<->}[r] & K[1_j](\ast)[1_i]\ar@{<->}[d] \\
K[1_i](\ast)[1_j] \ar@{<->}[r] & K(\ast)[1_i][1_j]  \ar@{=}[r]& K(\ast)[1_j][1_i]\rlap{.}
}}\leqno{(1.1.4.7.1)}
$$


\subsubsection{}\etiquette[1.1.5]{XVII.subsubsec1.1.5}
Let\pageoriginale $F:\Aa\rightarrow \Bb$ be a contravariant additive
functor between additive categories. Its \emph{extension to complexes}
$F:C(\Aa)\rightarrow C(\Bb)$ is defined as follows. If $K$ is a
complex of $\Aa$, set
$$\etiquette[1.1.5.1]{XVII.eq1.1.5.1}
\begin{cases}
F(K)^{k}=F(K^{-k})\\
d^{k}_{F(K)}=(-1)^{k+1}F(d^{-k-1})\sisi{}{.}
\end{cases}\leqno(1.1.5.1)
$$

The \emph{canonical isomorphism $\sigma$ between $F(K)[1]$ and
$F(K[-1])$} is defined by
$$\etiquette[1.1.5.2]{XVII.eq1.1.5.2}
\sigma^{\sisi{\sheafk }{k}}=(-1)^{k+1}\id_{F(K^{-k-1})}{\quoi}.\leqno(1.1.5.2)
$$

If $F:\Aa\rightarrow \Bb$ and $G:\Bb\rightarrow
\sisi{C}{\Cc}$ are two contravariant additive functors, define the
canonical isomorphism $\rho$ between $G{\circ}F(K)$ \emph{and}
$G(F(K))$ by
$$\etiquette[1.1.5.3]{XVII.eq1.1.5.3}
\rho^{k}=(-1)^{k}\sisi{}{\id_{G\circ F(K^k)}}{\quoi}.\leqno(1.1.5.3)
$$

\begin{enonce*}{Lemma 1.1.5.4}\etiquette[1.1.5.4]{XVII.lem1.1.5.4}
\begin{enumerate}
\renewcommand{\theenumi}{\roman{enumi}}
\renewcommand{\labelenumi}{(\theenumi)}
\item For every contravariant additive functor
  $F:\Aa\rightarrow \Bb$, formulas
  \eqhyperref[XVII.eq1.1.5.1]{1.1.5.1} and
  \eqhyperref[XVII.eq1.1.5.2]{1.1.5.2} define an exact contravariant
  functor from $K(\Aa)$ to $K(\Bb)$.

\item For every pair $(G,F)$ of composable contravariant additive
  functors, formula \eqhyperref[XVII.eq1.1.5.3]{1.1.5.3} defines an
  isomorphism of graded functors between the extension of $GF$ to
  complexes and the composite of the extensions of $G$ and $F$ to
  complexes.
\end{enumerate}
\end{enonce*}

\begin{enonce*}[remark]{Proof}
A\pageoriginale morphism of degree $i$, $f:K\rightarrow L$, may be
regarded as any one of the following:
\begin{enumerate}
\renewcommand{\theenumi}{\alph{enumi}}
\renewcommand{\labelenumi}{(\theenumi)}
\item a morphism from $K$ to $L[i]$;

\item a morphism from $K[-i]$ to $L$;

\item a family $f^{n}:K^{n}\rightarrow L^{n+i}$ which commutes or
  anticommutes with the differentials according to the parity of $i$.
\end{enumerate}
\end{enonce*}

If $f:K\rightarrow L[i]$ is a morphism of degree $i$, then
$F(f):F(L)[-i]\rightarrow F(K)$ is again a morphism of degree $i$;
for $i=1$, we have
$$\etiquette[1.1.5.5]{XVII.eq1.1.5.5}
F(f)^{n}=(-1)^{n+1}F(f^{-n-1}){\quoi}.\leqno(1.1.5.5)
$$

If $0\rightarrow K\xrightarrow{u} L\xrightarrow{v} M\rightarrow 0$
is an exact sequence of complexes split degree by degree, the choice
of a splitting allows the differential of $L$ to be written in the
form
\[
d_L=d_K+d_M+f{\quoi},
\]
where $f$ is a morphism of degree one from $M$ to $K$. Likewise, the
dual exact sequence
\[
0\longrightarrow F(M)\longrightarrow F(L)\longrightarrow
F(K)\longrightarrow 0
\]
defines a morphism of degree one $f^{\ast}$ from $F(K)$ to
$F(\sisi{L}{M})$. We have $f^{\ast}=F(f)$. It follows that the image
under $F$ of a triangle $(K,L,M,u,v,f)$ defined by a short exact
sequence of complexes is again a triangle of the same type and, by
\hyperref[XVII.subsubsec1.1.1]{1.1.1}, this proves (i).

Under the hypotheses of (ii), for $f$ of degree one, we have
$G(F(f))^{n}=(-1)^{n+1}(-1)^{-n}G\circ F(f)^{n}=-G\circ F(f)^{n}$; in
particular, the grading isomorphism $G(F(K[1]))\simeq
G(F(K))[1]$ is the isomorphism whose components are $-1$.

The\pageoriginale diagram
\[
\xymatrix{
G\circ F(K)[1]  \ar@{<->}[r]^{\rho[1]} \ar@{=}[d] & G(F(K))[1] \ar@{<->}[d]\\
G\circ F(K[1])  \ar@{<->}[r]^{\rho} & G(F(K[1]))
}
\]
is therefore commutative, which proves (ii).

\subsubsection{}\etiquette[1.1.6]{XVII.subsubsec1.1.6}
Let $s=+$ or $s=-$; let $I$ be a finite set, let
$\epsilon:I\rightarrow \{+,-\}$ be a function, let
$(\Aa_i)_{i\in I}$ and $\Aa$ be additive categories, and let $F$ be
a multiadditive multifunctor from the categories $\Aa_i$ to $\Aa$,
covariant in the arguments indexed by those $i$ for which
$\epsilon(i)=+$ and contravariant in the others. We shall define the
extension $K^{s}(F)$ (or simply $F$) of $F$ to complexes as an exact
multifunctor from the categories $K^{s\epsilon(i)}(\Aa_i)$
($s\epsilon(i)=\pm$, according to the sign rule) to $K^{s}(\Aa)$,
\footnote[1]{The construction also provides an extension of $F$ to a
  multifunctor from the categories $C^{s\epsilon(i)}(\Aa_i)$ to
  $C^{s}(\Aa)$.} having the same variance as $F$. We shall then explain
in what sense this construction is compatible with the composition of
functors. Other conventions for \emph{degrees} are possible and will
be used; that issue is independent of the questions of \emph{sign}
considered here.

\paragraph{}\etiquette[1.1.6.1]{XVII.par1.1.6.1}
If $F$ is contravariant in one variable, $K^{s}(F)$ is defined by the
formulas of \hyperref[XVII.subsubsec1.1.5]{1.1.5}.

\paragraph{}\etiquette[1.1.6.2]{XVII.par1.1.6.2}
If $F$ is covariant, the functor $K^{s}(F)$ is defined, by passage to
the quotient, from the composite of the evident functor from the
product of the categories $C^{s}(\Aa_i)$ to the category of naive
$I$-fold complexes, the functor
(\hyperref[XVII.subsubsec1.1.4]{1.1.4}) from the category of naive
$I$-fold complexes to that of\pageoriginale $I$-fold complexes
(\hyperref[XVII.par1.1.4.2.1]{1.1.4.2.1}), and the associated simple
complex functor \eqhyperref[XVII.eq1.1.4.0]{1.1.4.0}:
\[
K^{s}(F)=C^{s}(F)(\ast)_{s.a}{\quoi}.
\]

Let $K_j\in K^{s}(\Aa_j)$ be complexes, and let $i\in I$; set
$K_j[1_i]=K_j$ for $j\neq i$, and $K_i[1_i]=K_{i}[1]$. We then have
the evident isomorphisms
\[
C^{s}(F)(K_{j}[1_i])\simeq C^{s}(F)(K_j)[1_i]{\quoi},
\]
and
\[
C^{s}(F)(K_j)(\ast)[1_i]_{s.a}\simeq
C^{s}(F)(K_j)(\ast)_{s.a}[1]{\quoi},
\]
whence, via \eqhyperref[XVII.eq1.1.4.5]{1.1.4.5}, a
\emph{canonical isomorphism} $\sigma$ between
$K^{s}(F)(K_{j}[1_i])$ and $K^{s}(F)(K_{j})[1]$. It follows from
\hyperref[XVII.par1.1.4.7]{1.1.4.7} that these isomorphisms define a
grading of $K^{s}(F)$.

\paragraph{}\etiquette[1.1.6.3]{XVII.par1.1.6.3}
In the general case, $F$ is the composite of the canonical
contravariant functors $\Aa_i\rightarrow \Aa^{\circ}_i$ (for
$\epsilon(i)=-$) with a covariant functor $F^{+}$; define the functor
$K^{s}(F)$ by composition using
\hyperref[XVII.par1.1.6.1]{1.1.6.1} and
\hyperref[XVII.par1.1.6.2]{1.1.6.2}.

\begin{enonce*}{Proposition 1.1.7}\etiquette[1.1.7]{XVII.prop1.1.7}
Under the preceding hypotheses, the functor $K^{s}(F)$ is an exact
multiadditive functor.
\end{enonce*}

It suffices to verify \hyperref[XVII.prop1.1.7]{1.1.7} under the
hypotheses of \hyperref[XVII.par1.1.6.1]{1.1.6.1} (where the
proposition is \hyperref[XVII.lem1.1.5.4]{1.1.5.4}~(i)) or under those
of \hyperref[XVII.par1.1.6.2]{1.1.6.2}. To verify that the functors
deduced from $K^{s}(F)$ by fixing all variables except the $i$-th are
exact, one may describe $K^{s}(F)$ explicitly in terms of an order on
$I$ for which $i$ is the least element; a distinguished triangle of
the standard\pageoriginale form
\eqhyperref[XVII.eq1.1.1.1]{1.1.1.1} is then sent to a distinguished
triangle of the same form.

\setcounter{subsubsection}{7}
\subsubsection{}\etiquette[1.1.8]{XVII.subsubsec1.1.8}
Let $I$ and $J$ be two finite sets, let $\Psi:I\rightarrow J$ be a
surjective map, and let $\epsilon_{I}:I\rightarrow \{+,-\}$ and

% END INCLUDED SGA4_English_current_buildable_source_20260906/source/17/drafts/SGA4_XVII_LINES_261_600_English.texfrag

% BEGIN INCLUDED SGA4_English_current_buildable_source_20260906/source/17/drafts/SGA4_XVII_LINES_601_1000_English.texfrag
$\epsilon_{J}:J\rightarrow \{+,-\}$ be functions, and let us be given
multiadditive multifunctors, covariant in the arguments for which
$\epsilon=+$ and contravariant in the others:
\begin{align*}
G_j&: \prod_{\Psi(i)=j}\Aa_i\longrightarrow
\Bb_j\qquad (j\in J)\\
F&:  \prod_{j\in J} \Bb_j\longrightarrow \sisi{C}{\Cc.}
\end{align*}

Choose total orders on $I$ and $J$ such that $\Psi$ is an increasing
map. We then define the \emph{canonical isomorphism $\rho$ between the
  extension to complexes of $F\circ (G_j)$ and the composite of the
  extensions to complexes of $F$ and of the $G_j$} by the formula
$$\etiquette[1.1.8.1]{XVII.eq1.1.8.1}
\rho^{\sheafk }=\sisi{\prod_i \eta_i^{k_i}}{(-1)^{A(\sheafk )}}:\; \mbox{\rm automorphism of }
F\circ (G_j)(K_{i}^{k_i\epsilon\sisi{}{_I}(i)\epsilon\sisi{}{_J}(\Psi(i))}){\quoi}.\leqno(1.1.8.1)
$$%
where \sisi{$\eta_i=-1$ if $\epsilon_I(i)=\epsilon_{J}(\Psi(i))=-$,
$\eta_i=1$ otherwise.}{\[
A(\sheafk )=\sum_{\epsilon_I(i)=\epsilon_J(\Psi(i))=-} k_i + \sum_{\epsilon_J(j)=-} \sum_{\substack{\Psi(a)=\Psi(b)=j\\a<b}} k_a k_b.
\]}%See the Erratum in SGA 4 1/2, p. 312

This formula generalizes \eqhyperref[XVII.eq1.1.5.3]{1.1.5.3}; one verifies that $\rho$ does not depend
on the orders chosen on $I$ and $J$, and is an isomorphism of graded
functors (\cf \hyperref[XVII.lem1.1.5.4]{1.1.5.4} (ii)).

These isomorphisms $\rho$ satisfy the evident cocycle condition for a
triple composite.

\subsubsection{}\etiquette[1.1.9]{XVII.subsubsec1.1.9}
Let\pageoriginale $I$ be a finite set, let $i$ and $j$ be two distinct
elements of $I$, let $I'=I\backslash\{i,j\}$, let $(\Aa_i)_{i\in I}$ and
$\Aa$ be additive categories, and let $F:(\Aa_i)_{i\in
  I}\rightarrow \Aa$ and $G:(\Aa_i)_{i\in
  I'}\rightarrow \Aa$ be multiadditive multifunctors. Suppose that
$\Aa_i=\Aa_j$ and that $F$ is contravariant in the $i$-th variable and
covariant in the $j$-th. For $X$, $Y\in \Ob \Aa_i$, denote by
$F_{X,Y}$ the multifunctor from $(\Aa_i)_{i\in I'}$ to $\Aa$
obtained by fixing the $i$-th and $j$-th arguments:
\[
F_{X,Y}(\ldots)=F(\ldots\sisi{}{,} X,\ldots,Y,\ldots){\quoi}.
\]
A \emph{contraction morphism} $c:F\rightarrow G$ consists of, for
every $X\in \Aa_i$, a morphism of functors
\[
c_{X}:F_{X,X}\longrightarrow G{\quoi},
\]
such that, for every arrow $f:X\rightarrow Y$ of $\Aa_i$, the diagram
\[
\xymatrix{
F_{Y,X} \ar[d] \ar[r] & F_{X,X} \ar[d]^{c_X}\\
F_{Y,Y}  \ar[r]_{c_Y}     & G
}
\]
is commutative.

The conventions of \hyperref[XVII.subsubsec1.1.6]{1.1.6} are motivated by the

\begin{enonce*}{Lemma 1.1.9.1}\etiquette[1.1.9.1]{XVII.lem1.1.9.1}
Let $c:F\rightarrow G$ be a contraction morphism, and let
$(K_{\ell})_{\ell \in I}$ be complexes in the $\Aa_{\ell}$, with
$K_{i}=K_{j}$. Let $<$ be a total order on $I$ for which $i$ is the
predecessor of $j$. With the notation of
\hyperref[XVII.subsubsec1.1.6]{1.1.6}, for $k\in \ZZ$, let
\[
c^{k}:K^{s}(F)(K_{\ell})^{k}\longrightarrow K^{s}(G)(K_{\ell})^{k}
\]
be\pageoriginale the morphism whose restriction to $F(\ldots
K_{i}^{k_i}K_{j}^{k_j}\ldots)$ is $0$ if $k_i\neq k_j$, and is
$c_{K_i^{k_i}}$ if $k_{i}=k_{j}$. Then $(c^{k})_{k\in \ZZ}$ is a
morphism of complexes.
\end{enonce*}

We leave to the reader the verification of this lemma, as well as the
fact that the \emph{contraction morphism} $(c^{k})$ does not depend on
the order chosen on $I$.

The standard example of a contraction morphism is the composition law
\[
\Hom(Y,Z)\otimes \Hom(X,Y)\longrightarrow \Hom(X,Z)\sisi{}{.}
\]

\begin{enonce*}[remark]{Examples 1.1.10}\etiquette[1.1.10]{XVII.exem1.1.10}
\begin{enumerate}
\renewcommand{\theenumi}{\roman{enumi}}
\renewcommand{\labelenumi}{(\theenumi)}
\item Applying the definitions of \hyperref[XVII.subsubsec1.1.6]{1.1.6} to the
  \emph{tensor-product} functor gives the usual notion of the tensor
  product of two complexes, with differential
\[
d(x\otimes y)=dx\otimes y+(-1)^{\deg x}x\otimes dy{\quoi}.
\]

\item Apply the definitions of \hyperref[XVII.subsubsec1.1.6]{1.1.6} to the
  functor $\Hom(X,Y)$, \emph{regarding $Y$ as the first variable and
  $X$ as the second variable}, a convention that we shall always
  follow. This convention identifies, for two complexes $X$, $Y$, the
  $k$\sisi{-th}{-th} component of the complex $\Hom(X,Y)$ defined in
  \hyperref[XVII.subsubsec1.1.6]{1.1.6} with $\prod_n\Hom(X^{n},Y^{n+k})$, and gives the
  usual differential: for $f\in \Hom(X^{n},Y^{n+k})$, one has
\[
df=d\circ f-(-1)^{k}f\circ d{\quoi}.
\]

\item The preceding conventions have the virtue that the following
  systems of arrows are morphisms of complexes, without any signs
  having\pageoriginale to be inserted (\hyperref[XVII.subsubsec1.1.9]{1.1.9}):
\begin{align*}
& \Hom^{\bigdot}(L,M)\otimes \Hom^{\bigdot}(K,L)\longrightarrow
  \Hom^{\bigdot}(K,M):f\otimes g\mapsto f\circ g\\
& \Hom^{\bigdot}(L,M)\otimes L\longrightarrow M: f\otimes x\mapsto f(x){\quoi}.
\end{align*}
\end{enumerate}
\end{enonce*}

\subsubsection*{${}^{\ast}$1.1.10.1}\etiquette[*1.1.10.1]{XVII.par1.1.10.1}
The convention \eqhyperref[XVII.eq1.1.4.5]{1.1.4.5}--(\hyperref[XVII.par1.1.6.2]{1.1.6.2}) is justified as follows for
$F=\otimes$: an element of $((K\otimes L)[1])^{n}$ is written in the
form (\hyperref[XVII.subsubsec*1.1.2]{*1.1.2}) (recalling that $(K\otimes L)[1]=C(K\otimes
L\rightarrow 0)$ \eqhyperref[XVII.eq1.1.1.1]{1.1.1.1}) $\sum 1\otimes(x^{p}\otimes y^{q})$, an
element of $(K[1]\otimes L)^{n}$ in the form $\sum (1\otimes
x^{p})\otimes y^{q}$, and an element of $(K\otimes L[1])^{n}$ in the
form $\sum x^{p}\otimes (1\otimes y^{q})$; the sign appears when
$x^{p}$ and $1$ are interchanged. For $F=$ dual, the convention
\eqhyperref[XVII.eq1.1.5.2]{1.1.5.2} is justified as follows: an element of $K^{\ast}[1]$ is written
$1\otimes \omega$; its value on an element $1'\otimes x$ of $K[-1]$
($1'$ a cell of dimension $-1$) is $<1\otimes \omega, 1'\otimes
x>=(-1)^{\deg\omega}<1, 1'><\omega,x>=(-1)^{\deg\omega}<\omega,x>$;
this is the convention adopted.${}_{\ast}$

\begin{enonce*}[remark]{Remark 1.1.11}\etiquette[1.1.11]{XVII.rem1.1.11}
``Recall'' that a \emph{preadditive category} is a category whose sets
of arrows $\Hom(X,Y)$ are endowed with abelian-group structures, the
composition of morphisms being biadditive. An \emph{additive functor}
$F$ between preadditive categories is a functor satisfying the formula
\[
F(f+g)=F(f)+F(g)
\]
for $f$, $g\in \Hom(X,Y)$. The additive functors from a preadditive
category $\Aa$ to another $\Bb$ form a category denoted by
$\Hhom_{\text{\rm add}}(\Aa,\Bb)$.
\end{enonce*}

Let\pageoriginale $(D)$ be the unique preadditive category whose set of objects is
indexed by $\ZZ$, $(X\sisi{_n}{^n})_{n\in \ZZ}$, and for which
$$\etiquette[1.1.11.1]{XVII.eq1.1.11.1}
\begin{cases}
\Hom(X^{n},X^{n})=\ZZ,\;\mbox{\rm generated by the identity}\\
\Hom(X^{n},X^{n+1})=\ZZ,\;\mbox{\rm generated by an arrow
  denoted $d_n$}\\
\Hom(X^{i},X^{j})=0\;\mbox{if } j-i\neq 0\;\mbox{\rm or } 1{\quoi}.
\end{cases}\leqno(1.1.11.1)
$$

Let $D$ be the complex of $(D)$ defined by $D^{n}=X^{n}$ and
$d^{n}=d_n$.

For every preadditive category $\Aa$, the functor
$G\mapsto G(D)$ from $\Hhom_{\text{\rm add}}((D),\Aa)$ to
$C(\Aa)$ is an isomorphism.

The construction of \hyperref[XVII.subsubsec1.1.6]{1.1.6} extends to the case of a multiadditive functor
from preadditive categories $\Aa_i$ to an additive category $\Aa$.

Let
\[
F^{\bigdot}=(F^{n},d^{n})_{n\in \ZZ}
\]
be a complex of multiadditive functors from additive categories
$\Aa_i$ to an additive category $\Aa$. One may regard
$F^{\bigdot}$ as a functor from $(D)$ to a category of multifunctors,
or equivalently as a multiadditive functor
\[
F_1:\prod\Aa_i\times (D)\longrightarrow \Aa{\quoi}.
\]
The \emph{extension of $F\sisi{}{^{\bigdot}}$ to complexes} is defined by the formula
$$\etiquette[1.1.11.2]{XVII.eq1.1.11.2}
F^{\bigdot}(K_{1}\ldots K_{n})=F_{1}(K_{1}\ldots K_{n},D){\quoi}.\leqno(1.1.11.2)
$$

\setcounter{subsubsection}{11}
\subsubsection{}\etiquette[1.1.12]{XVII.subsubsec1.1.12}
Let $\Aa$ and $\Bb$ be two additive categories and let
\[
F^{\bigdot}=(F^{n},d^{n})_{n\in \ZZ}
\]
be a complex of additive functors from $\Aa$ to $\Bb$.
Suppose that the functors\pageoriginale $F^{i}$ admit right adjoints $F_i$; let
$\leftexp{t}{d}^{i}:F_{i+1}\rightarrow F_i$ be the transpose of
$d^{i}$, and let $F_{\bigdot}$ be the complex of additive functors from
$\Bb$ to $\Aa$ whose components are the $F_{i}$ and whose differential
operator is
\[
(d_\bigdot)^{i}=(-1)^{i+1}\,\leftexp{t}{d}^{-i-1}:F_{-i}\longrightarrow
F_{-i-1}{\quoi}.
\]
Let ${}^{t}$ be the \emph{contravariant} functor from $(D)$ to $(D)$
given by
\[
\leftexp{t}{X}^{n}=X^{-n}\quav \leftexp{t}{d}^{n}=d^{-n-1}{\quoi}.
\]
The following diagram is commutative up to canonical isomorphism:
$$\etiquette[1.1.12.1]{XVII.eq1.1.12.1}
\vcenter{
\xymatrix@C=6em{
\mathcal{A} \times (D) \times \mathcal{B}  \ar[d]^{(\Id,{}^t,\Id)}\ar[r]^(.55){(F^{n}(A), \Id)} &
\mathcal{B}\times \mathcal{B}  \ar[r]^{\Hom_{\mathcal{B}}}&
(\Ab)\ar@{=}[d] \\
\mathcal{A} \times (D) \times \mathcal{B}
\ar[r]^(.6){(\Id,(F_{\bigdot})^n(B))} & \mathcal{A} \times \mathcal{A}
\ar[r]^{\Hom_{\mathcal{A}}} & (\Ab) \rlap{.}
}}\leqno{(1.1.12.1)}
$$
This expresses the adjunction formula
\begin{equation*}
\Hom(F^{n}(A),B)\overset{\sim}{\longleftrightarrow}\Hom(A,F_{n}(B)){\quoi}.
\end{equation*}
The general conventions of \hyperref[XVII.def1.1.3]{1.1.3} therefore give an isomorphism
\begin{equation*}
\Hom^{\bigdot}(F^{\bigdot}(K),L)\overset{\sim}{\longleftrightarrow}\Hom\sisi{}{^\bigdot}(K,F_{\bigdot}(L)){\quoi}.
\end{equation*}

\begin{enonce*}{Definition 1.1.13}\etiquette[1.1.13]{XVII.def1.1.13}
If $K$ is a complex in an abelian category $\Aa$, the \emph{right
truncation (\resp left truncation)} of $K$ in dimension $n$, denoted
by $\tau_{\leq n}(K)$ \emph{(\resp $\tau_{\geq n}(K)$)}, means the
following complexes:
\begin{align*}
\tau_{\leq n}(K) &: \ldots K^{n-2}\longrightarrow
K^{n-1}\longrightarrow \Ker(d^n)\longrightarrow 0\\
\tau_{\geq n}(K) &: 0\longrightarrow \coker (d^{n-1})\longrightarrow
K^{n+1}\longrightarrow K^{n+2}\ldots.
\end{align*}
\end{enonce*}

\begin{enonce*}{Proposition 1.1.14}\etiquette[1.1.14]{XVII.prop1.1.14}
\begin{enumerate}
\renewcommand{\theenumi}{\roman{enumi}}
\renewcommand{\labelenumi}{(\theenumi)}
\item $H^{i}(\tau_{\leq n}(K))=0$\pageoriginale for $i>n$, and $H^{i}(\tau_{\geq
n}(K))=0$ for $i<n$.

\item The canonical morphism $\tau_{\leq n}(K)\rightarrow K$ (\resp
$K\rightarrow \tau_{\geq n}K$) induces isomorphisms on the $H^{i}$
for $i\leq n$ (\resp $i\geq n$).
\end{enumerate}
\end{enonce*}

The functor $\tau_{\leq n}$ (\resp $\tau_{\geq n}$) sends homotopic
morphisms to homotopic morphisms. It therefore follows from
\sisi{\hyperref[XVII.subsubsec1.1.4]{1.1.4}}{\hyperref[XVII.prop1.1.14]{1.1.14}} that the functor of \hyperref[XVII.def1.1.13]{1.1.13} remains
meaningful in the derived category.

\setcounter{subsubsection}{14}
\subsubsection{}\etiquette[1.1.15]{XVII.subsubsec1.1.15}
If $K$ is a double complex, with first (\resp second) differential
$d'$ (\resp $d''$), denote by $\tau''_{\leq n}(K)$ the double
subcomplex of $K$ such that
\begin{equation*}
\tau''_{\leq n}(K)^{pq}=
\begin{cases}
K^{pq}&\mbox{\rm if } q<n\\
\Ker(d'') &\mbox{\rm if } q=n\\
0 &\mbox{\rm if } q>n
\end{cases}
\end{equation*}
and define $\tau''_{\geq n}(K)$, $\tau'_{\leq n}(K)$, and
$\tau'_{\geq n}(K)$ similarly.

\subsubsection{}\etiquette[1.1.16]{XVII.subsubsec1.1.16}
We shall also have occasion to consider the ``stupid truncations'' of
a complex $K$, defined by
\begin{xalignat*}{2}
\sigma_{\leq n}(K)^{i}=
\begin{cases}
K^{i} &\mbox{\rm if } i\leq n\\
0 &\mbox{\rm if } i>n
\end{cases}
&&
\sigma_{\geq n}(K)^{i}=
\begin{cases}
K^{i} &\mbox{\rm if } i\geq n\\
0 &\mbox{\rm if } i<n{\quoi}.
\end{cases}
\end{xalignat*}

\subsection{Derived functors}\etiquette[1.2]{XVII.subsec1.2}\pageoriginale

Let $\Aa$ be a triangulated category and $S$ a saturated
multiplicative system (\cite{sga4.exp.XVII.bib.CD} \sisi{}{I }2.1.2). If $A\in \Ob \Aa$, denote
by $A^{+}$ and $A^{-}$ the \iind- and \ppro-objects
((\hyperref[XVII.subsec0.5]{0.5}) and (\hyperref[XVII.subsec0.6]{0.6}))
\begin{equation*}
A^{+}=\quotinjlim_{A\xrightarrow{s}A'} A' \qquad A^{-}=\quotprojlim_{A'\xrightarrow{s}A}A'{\quoi},
\end{equation*}
where $s$ ranges over the filtered category of arrows in $S$ with
source or target $A$, respectively.

The functors $A\mapsto A^{+}$ and $A\mapsto A^{-}$ make the elements
of $S$ invertible and define fully faithful functors from
$\Aa(S^{-1})$ to the categories of Ind- and \ppro-objects of $\Aa$,
respectively.

\begin{enonce*}{Definition 1.2.1}\etiquette[1.2.1]{XVII.def1.2.1}
Let $F$ be an exact covariant multifunctor from triangulated categories
$(\Aa_i)_{0<i\leq n}$ to the triangulated category $\Aa$, and
let $S_i$ and $S$ be saturated multiplicative systems in the
categories $\Aa_i$ and $\Aa$.
\begin{enumerate}
\renewcommand{\theenumi}{\roman{enumi}}
\renewcommand{\labelenumi}{(\theenumi)}
\item The right derived functor $RF$ of $F$ (relative to the $S_i$ and
$S$) is the functor from the categories $\Aa_i(S^{-1}_i)$ to the
category of \iInd-objects of $\Aa(S^{-1})$ that makes the following
diagram commutative.
\[
\xymatrix@C=5em{%csy
  (\Aa_i)  \ar[d]\ar[r]^{(A_i)\sisi{\simarrow}{\mapsto}  F(A^{+}_i)} & \Ind
  \Aa\ar[d] \\
\Aa_i(S^{-1}_i)\ar[r]^{RF} & \Ind\Aa(S^{-1}).
}
\]

\item We\pageoriginale say that $RF$ is defined at $(A_i)$ if the
\iInd-object $RF(A_i)$ of $\Aa(S^{-1})$ is ``essentially constant,''
\ie comes from an object of $\Aa(S^{-1})$; this object is called
\emph{the value of} $RF$ at $(A_i)$.

\item We say that $F$ is \emph{right derivable} if $RF$ is defined
everywhere; in that case we shall again denote by $RF$ the resulting
functor from the $\Aa_i(S^{-1}_i)$ to $\Aa(S^{-1})$.

\item A family $(A_i)$ of objects of the $\Aa_i$ is said to be
\emph{adapted to} $RF$ if the canonical morphism from $F(A_i)$ to
$RF(A_i)$ is an isomorphism.
\end{enumerate}
\end{enonce*}

We leave to the reader the task of defining by duality the left
derived functor $LF$ of $F$ (with values in $\Pro \Aa(S^{-1})$) and
of extending the preceding definitions to cases in which $F$ is
covariant in some arguments and contravariant in others.

When $RF$ is defined everywhere, the definition given here agrees
with Verdier's (\cite{sga4.exp.XVII.bib.CD} p.~39), as is readily verified.

\begin{enonce*}{Proposition 1.2.2}\etiquette[1.2.2]{XVII.prop1.2.2}
Let $(A'_1, A_1,A''_1)$ be a distinguished triangle in $\Aa_1$, and
let\break $A_i\in \Ob \Aa_i$ $(i\neq 1)$.
\begin{enumerate}
\renewcommand{\theenumi}{\roman{enumi}}
\renewcommand{\labelenumi}{(\theenumi)}
\item If $(A'_{1},(A_{i}))$ and $(A''_{1},(A_{i}))$ are adapted to
  $RF$, then the family $(A_{1},(A_{i}))$ is also adapted to $RF$.

\item If $RF$ is defined at $(A'_1, (A_i))$ and at $(A''_1, (A_i))$,
it is also defined at $(A_1, (A_i))$, and the triangle
$(RF(A'_1,(A_i)), RF(A_1,(A_i)),RF(A''_1,(A_i)))$ is distinguished in
$\Aa(S^{-1})$.
\end{enumerate}
\end{enonce*}

Let us prove\pageoriginale (ii) $\Rightarrow$ (i). Since $F$ is exact,
the arrow $$(F(A'_1,(A_i)), F(A_1,(A_i)), F(A''_1,(A_i)))\rightarrow
(RF(A'_1,(A_i)), RF(A_1,(A_i)), RF(A''_1,(A_i)))$$ is a morphism of
distinguished triangles. By hypothesis this arrow is an isomorphism
at two vertices; it is therefore an isomorphism at the third, as
desired.

\begin{enonce*}{Lemma 1.2.2.1}\etiquette[1.2.2.1]{XVII.lem1.2.2.1}
If $\Delta=(X,Y,Z)$ is a distinguished triangle in a triangulated
category $\Aa$ and $S$ is a saturated multiplicative system in $\Aa$,
then the triangle $\Delta^{+}=(X^{+},Y^{+},Z^{+})$ in $\Ind\Aa$ is a
limit of distinguished triangles in $\Aa$.
\end{enonce*}

Let $\mathcal{L}$ denote the category of morphisms of distinguished
triangles with source $\Delta$ whose arrows belong to $S$.
\begin{enumerate}
\renewcommand{\theenumi}{\alph{enumi}}
\renewcommand{\labelenumi}{(\theenumi)}
\item If $\Delta=(X,Y,Z)$ and $\Delta'=(X',Y',Z')$ are two
distinguished triangles in $\Aa$, and if $f:\Delta\rightarrow \Delta'$
is an $\Aa(S^{-1})$-morphism of triangles, there exist in $\Aa$
morphisms of triangles $s:\Delta'\rightarrow\Delta''$ and
$f_1:\Delta\rightarrow \Delta''$ such that the components of $s$
belong to $S$ and the diagram
\[
\vcenter{
\xymatrix{
\Delta                \ar[rd]^{f_1}  \ar[d]_{f}  &\\
\Delta'          \ar[r]^{s}    & \Delta''
}}\tag{1}
\]
is commutative in $\Aa(S^{-1})$.

% END INCLUDED SGA4_English_current_buildable_source_20260906/source/17/drafts/SGA4_XVII_LINES_601_1000_English.texfrag

% BEGIN INCLUDED SGA4_English_current_buildable_source_20260906/source/17/drafts/SGA4_XVII_LINES_1001_1400_English.texfrag

\sisi{\noindent}{}%
There is a commutative diagram
\[
\vcenter{
\xymatrix{
X \ar[rd]^{f_1} \ar[d]_{f}&  &\\
X' \ar[r]^{s\sisi{}{_X}} & X''\sisi{}{_0}&
}}(s\sisi{}{_X} \in S, f_1 \in \Fl \Aa)\tag{2}
\]
and\pageoriginale a morphism $\Delta'\xrightarrow{s\sisi{_1}{}}\Delta''\sisi{_1}{_0}$
in $S$, having $s\sisi{}{_X}$ as a component. Argue similarly at the
other vertices: one obtains a diagram of type (1), in which, however,
$f_1$ is not yet a morphism of triangles \emph{in}
$\Aa$, since the diagrams
\[
\vcenter{
\xymatrix{
Z  \ar[r]^{f\sisi{}{_1}} \ar[d]_{w}& Z''_{0} \ar[d]_{w\sisi{'}{''}}\\
X         \ar[r]^{f\sisi{}{_1}} & X''_{0}
}}\tag{3}
\]
may, for example, fail to be commutative \emph{in} $\Aa$,
being commutative only in $\Aa(S^{-1})$. There exists
$t\sisi{_0}{_X}:X''_0\rightarrow X''_1$ in $S$ such that $t\sisi{}{_X}w''f\sisi{}{_1}=t\sisi{}{_X}f\sisi{}{_1}w$
and $t:\Delta''_0\rightarrow \Delta''_1$ having $t\sisi{_0}{_X}$
as a component. Replace $\Delta''_0$ by $\Delta''_1$,
$s$ by $ts$, and $f_1$ by $tf_1$; in the resulting diagram of type (1),
(3) is now commutative; proceeding similarly at the other
vertices, one obtains (1). We leave it to the reader to verify
similarly:

\item If $f$, $g:\Delta\rightrightarrows \Delta'$ are two morphisms
of distinguished triangles in $\Aa$, equal in
$\Aa(S^{-1})$, then there exists a morphism of distinguished
triangles whose arrows are in $S$, $s:\Delta'\rightarrow \Delta''$, such
that $\sisi{S}{s}f=\sisi{S}{s}g$.

\item The category $\Ll$ is filtering on the right.

\item The proof of the lemma is completed by noting that
\[
\Delta^{+}=\quotinjlim_{s\in \Ll}{\mbox{\rm (target
of $s$).}}
\]
\end{enumerate}

Let us prove\pageoriginale \hyperref[XVII.prop1.2.2]{1.2.2} (ii). Let $X'$ and $X''$ be objects of
$\Aa(S^{-1})$ representing $RF(A'_1,(A_i))$ and $RF(A''_1,
(A_i))$, and complete the morphism of degree 1 from $X''$ to $X'$ to
a distinguished triangle $(X',X,X'')$. Replacing the $A_i$, $A'_1$,
$A_1$, and $A''_1$ by the target of an arrow in \sisi{$S$}{$S_i$ or $S_1$} with source $A_i$,
$A'_1$, $A_1$, or $A''_1$, we reduce to the case where there is a
commutative diagram
\[
\xymatrix{
F(A''_1, (A_i)) \ar[d] & X''  \ar[l]_-{a''} \ar[d]\\
F(A'_1, (A_i))    & X'  \ar[l]_-{a'}\rlap{,}
}
\]
where $a''$ and $a'$ are sections of arrows from $F(A''_1,(A_i))$ to
$X''$ and from $F(A'_1,(A_i))$ to $X'$. By axiom $TR$ 3 of
\cite{sga4.exp.XVII.bib.CD}, this diagram can be completed to a morphism of triangles
\begin{align*}
(X'XX'')& \longrightarrow
(F(A'_1,(A_i)),F(A_1,(A_i)),F(A''_1,(A_i)))\\
& \longrightarrow
(RF(A'_1,(A_i)),RF(A_1,(A_i)),RF(A''_1,(A_i))){\quoi}.
\end{align*}

The composite arrow is a morphism of triangles whose source is a
distinguished triangle and whose target is a limit (in
$\Ind\Aa(S^{-1})$) of distinguished triangles (\hyperref[XVII.lem1.2.2.1]{1.2.2.1}). For any
$Y$ in $\Aa(S^{-1})$, the following diagram is
commutative and its rows are exact:
\[
\text{\tiny
$\xymatrix@C=1em{
\Hom\sisi{'}{(}Y, X''[\sisi{}{-}1]) \ar[r]\ar[d]& \Hom(Y,X') \ar[r] \ar[d]& \Hom(Y,X)
\ar[r]\ar[d]& \Hom(Y,X'') \ar[r] \ar[d]&
\Hom(Y,X'[1]) \ar[d]\\
\Hom(Y,RF(A''\sisi{}{_1},A_i)[\sisi{}{-}1])\ar[r] & \Hom(Y,RF(A'_1,A_i)\sisi{}{)}\ar[r] &  \Hom(Y,RF(A_1,A_i))\ar[r] &
\Hom(Y,RF(A''_1,A_i)\sisi{}{)}\ar[r] & \Hom(Y,RF(A'_1,A_i)[1])
}$}
\]
By\pageoriginale assumption, all arrows other than the middle arrow
are isomorphisms. The \hyperref[XVII.sec5]{5}-lemma shows that the middle arrow
is also an isomorphism, for every $Y$, and hence that the distinguished
triangle $(X',X,X'')$ is isomorphic to the triangle ($RF(A'_1,(A_i))$,
$RF(A_1,(A_i))$, $RF(A''_1,(A_i))$). This completes the proof.

If $\Aa$ and $\Bb$ are two abelian categories and
$F$ is an additive covariant functor from $\Aa$ to $\Bb$,
$F$ defines exact functors from $K(\Aa)$ to
$K(\Bb)$, from $K^{+}(\Aa)$ to $K^{+}(\Bb)$, and
from $K^{-}(\Aa)$ to $K^{-}(\Bb)$ (\hyperref[XVII.prop1.1.7]{1.1.7}).

\begin{enonce*}{Definition 1.2.3}\etiquette[1.2.3]{XVII.def1.2.3}
\begin{enumerate}
\renewcommand{\theenumi}{\roman{enumi}}
\renewcommand{\labelenumi}{(\theenumi)}
\item The \emph{right derived functors of $F$} are the functors
from $D^{+}(\Aa)$ to $\Ind D^{+}\sisi{\Bb}{(\Bb)}$ and from
$D(\Aa)$ to $\Ind D(\Bb)$ right derived from the
extensions of $F$ to complexes $K^{+}(\Aa)\rightarrow
K^{+}(\Bb)$ and $K(\Aa)\rightarrow K(\Bb)$.

\item An object $A$ of $\Aa$ is said to be \emph{right acyclic}
for $F$, or, by abuse of language, \emph{acyclic for} $RF$, if the
complex concentrated at $A$ in degree $0$ is adapted to $RF$.
\end{enumerate}
\end{enonce*}

We shall be interested mainly in the case where $RF$ is defined everywhere; it
may then be regarded as a functor from $D^{+}(\Aa)$ to
$D^{+}(\Bb)$, or from $D(\Aa)$ to $D(\Bb)$, as the case may be.

\begin{enonce*}{Proposition 1.2.4}\etiquette[1.2.4]{XVII.prop1.2.4}
The following diagram is commutative.
\[
\xymatrix{
D^{+}(\mathcal{A})  \ar[r] \ar[d]^{RF}& D(\mathcal{A})\ar[d]^{RF}        \\
 \Ind \;D^{+}(\mathcal{B}) \ar[r] & \Ind \; D(\mathcal{B}).
}
\]
\end{enonce*}

If\pageoriginale $K$ is a complex that is zero in degrees $<n$ and
$s:K\rightarrow L$ is a quasi-isomorphism, the composite morphism
$K\xrightarrow{s}L\rightarrow \sigma_{\geq n}(L)$
(\hyperref[XVII.subsubsec1.1.6]{1.1.6}) is again a quasi-isomorphism. The
quasi-isomorphisms from $K$ to a complex bounded below therefore
form a cofinal (full) subcategory of the category of
all quasi-isomorphisms with source $K$. The proposition follows
formally.

In particular, if the functor $RF$\sisi{,}{} is defined everywhere on
$D(\Aa)$, it is defined everywhere on $D^{+}(\sisi{\Bb}{\Aa})$, and
the derived functor $RF:D^{+}(\Aa)\rightarrow
D^{+}(\Bb)$ is the restriction of the derived functor
$RF:D(\Aa)\rightarrow D(\Bb)$.

\setcounter{subsubsection}{4}
\subsubsection{}\etiquette[1.2.5]{XVII.subsubsec1.2.5}
We leave it to the reader to:
\begin{enumerate}
\renewcommand{\theenumi}{\roman{enumi}}
\renewcommand{\labelenumi}{(\theenumi)}
\item Generalize Definition \hyperref[XVII.def1.2.3]{1.2.3} to the case of a complex of
multifunctors, covariant in some variables and contravariant in
others;

\item Dualize Definition \hyperref[XVII.subsubsec1.2.5]{1.2.5} (i) to define the left derived
functors;

\item Extend Proposition \hyperref[XVII.prop1.2.4]{1.2.4} to the case of a bounded complex of
multifunctors;

\item Dualize Proposition \hyperref[XVII.subsubsec1.2.5]{1.2.5} (iii) by
considering left derived functors and the inclusion of
$D^{-}(\Bb)$ into $D(\Bb)$.
\end{enumerate}

\begin{enonce*}{Proposition 1.2.6}\etiquette[1.2.6]{XVII.prop1.2.6}
Let $\Aa$ and $\Bb$ be two abelian categories,
let $𝒫$ be a subset of $\Ob(\Aa)$, and let $F$ be an additive functor from
$\Aa$ to $\Bb$. Assume that
\begin{enumerate}
\renewcommand{\theenumi}{\alph{enumi}}
\renewcommand{\labelenumi}{(\theenumi)}
\item every object of $\Aa$ is a quotient of an element of $𝒫$;

\item for every short exact sequence
\[
0\longrightarrow P\longrightarrow Q\longrightarrow R\longrightarrow
0{\quoi},
\]
where $Q$ and $R$ belong to $𝒫$, $P$ also belongs to $𝒫$, and the
$0$-sequence $0\rightarrow FP\rightarrow FQ\rightarrow FR\rightarrow
0$\pageoriginale is exact. Then every object of $𝒫$ is
acyclic for $LF$.
\end{enumerate}
\end{enonce*}

If $A\in \Ob \Aa$, the left resolutions of $A$ by
objects of $𝒫$ form a cofinal system in the filtering
category of left resolutions of $A\sisi{_0}{}$. It therefore suffices to
prove that if a complex
\[
\ldots\longrightarrow A_2\longrightarrow A_1\longrightarrow A_0
\longrightarrow A\longrightarrow 0
\]
is acyclic and has its objects in $𝒫$, its image is acyclic. Splitting it
into short exact sequences
$$\etiquette[1.2.6.i]{XVII.eq1.2.6.i}
0\longrightarrow \Iim(d_{i-1})\longrightarrow A_i\longrightarrow
\Iim(d_i)\longrightarrow 0\qquad\sisi{(1.2.6. i)\quad}{}(i\geq 0)\sisi{}{\leqno(1.2.6.i)}%\tag{1.2.6.$i$}
$$
one sees by induction on $i$ that $\Iim(d_i)\in \PP$ and that the image
of \sisi{(\hyperref[XVII.prop1.2.6]{1.2.6} $i$)}{\eqhyperref[XVII.eq1.2.6.i]{1.2.6.i}} under $F$ is a short exact sequence. The image of the complex
is therefore acyclic.

\begin{enonce*}{Proposition 1.2.7}\etiquette[1.2.7]{XVII.prop1.2.7}
Let $F$ be an additive functor from an abelian category to an
abelian category $\Bb$. If every object of $\Aa$ is a
quotient (\resp subobject) of an object acyclic for $LF$ (\resp
$RF$) (\hyperref[XVII.def1.2.3]{1.2.3}), the derived functor $LF:D^{-}(\Aa)\rightarrow
D^{-}(\Bb)$ (\resp $RF:D^{+}(\Aa)\rightarrow
D^{+}(\Bb)$) is defined everywhere, and every complex bounded
above (\resp below) of objects acyclic for $LF$
(\resp $RF$) is adapted.
\end{enonce*}

It suffices to consider the left derived functor $LF$. If $K$ is a complex
bounded above, the quasi-isomorphisms $s:L\rightarrow K$ with
$L$ bounded above and consisting of acyclic objects form a
cofinal system in the category of quasi-isomorphisms with target $K$. It
suffices to show that, when $K$ consists of acyclic objects,
$F(s):F(L)\rightarrow F(K)$ is a quasi-isomorphism;\pageoriginale
then, in the general case, one will have $F(L)=LF(K)$. Let $M$ be the cone of
$s$. The complex consists of acyclic object\sisi{}{s} and $H^{\ast}(M)=0$; one must
prove that $H^{\ast}(F(M))=0$.

Let $i$ be an integer. The adapted complexes form a triangulated
subcategory, so that every bounded complex of acyclic objects,
in particular $\tau_{\geq i}(M)$, is adapted. Since $H^{j}(\tau_{\geq
i}(\sisi{L}{M}))=0$ for $j>i$, one has, for $j>i$,
$H^{j}(F(M))=H^{j}(F(\tau_{\geq i}(M)))=H^{j}LF(\tau_{\geq i}(M))=0$.

\begin{enonce*}{Proposition 1.2.8}\etiquette[1.2.8]{XVII.prop1.2.8}
Let $F^{\bigdot}$ be a complex bounded below of functors from
$\Aa$ to $\Bb$. If for every $K\in
K^{+}(\Aa)$ there exists a quasi-isomorphism $L\xrightarrow{s}K$,
where $L\in K^{+}(\Aa)$ is adapted to all the $RF^{i}$, then
the derived functor $RF^{\bigdot}:D^{+}(\Aa)\rightarrow
D^{+}(\Bb)$ exists, and every complex \emph{bounded
below} that is adapted to all the $RF^{i}$ is adapted to $RF^{\bigdot}$.
\end{enonce*}

\sisi{\noindent}{}%
The proof is left to the reader, as is that of the dual statement,
or of the variants obtained by taking $F$ contravariant, or by taking
a bounded complex of functors and working in $D(\Aa)$
and $D(\Bb)$.

\begin{enonce*}{Definition 1.2.9}\etiquette[1.2.9]{XVII.def1.2.9}
Let $F$ be a functor from $D^{b}(\Aa)$ to $D(\Bb)$, and let
$d$ be an interval in $\ZZ$. We say that $F$ has
\emph{cohomological amplitude $\subset d$} if, for every $A\in \Ob
\Aa$, $H^{i}F(A)=0$ for $i\not\in d$.
\end{enonce*}

Very often, $d$ will be an interval of the form $[0,x]$ or $[-x,0]$,
and we then say that $F$ has cohomological \resp
homological dimension $\leq x$. We say that $F$ has finite dimension if
there exists a finite interval $d$ in $\ZZ$ such that $F$ has
amplitude $\subset d$.

\begin{enonce*}{Proposition 1.2.10}\etiquette[1.2.10]{XVII.prop1.2.10}
Assume\pageoriginale that the hypotheses of \hyperref[XVII.prop1.2.7]{1.2.7} hold and that the
functor $RF:D^{+}(\Aa)\rightarrow D^{+}(\Bb)$ (\resp
$LF:D^{-}(\Aa)\rightarrow D^{-}(\Bb)$) has
finite dimension. Then every complex of objects acyclic for $RF$
(\resp $LF$) is adapted, and the functor $RF:D(\Aa)\rightarrow
D(\Bb)$ (\resp $LF:D(\Aa)\rightarrow D(\Bb)$) is defined everywhere.
\end{enonce*}

For the proof, we refer to Verdier \cite{sga4.exp.XVII.bib.CD}. There is an analogous statement
for bounded complexes of functors. The superscript $+$ (\resp $-$)
may also be replaced by $b$.

\section{Categories fibered in derived categories}\etiquette[2]{XVII.sec2}

\subsection{Introduction}\etiquette[2.1]{XVII.sec2.1}

The editor insists that the reader refrain from reading this \S. Here
we give the formalism of the trivial duality theorem and
resolve the perplexity raised by Artin in \hyperrefext[XII][XII.sec4]{4}.

For the notions of a category over another category, a fibered category, a
cofibered category, and a cleavage, we refer to \SGA 1 VI. Among other things,
that exposé establishes an equivalence between the notions of
normalized split fibered categories over a category $\sisi{C}{\Cc}$, and
pseudofunctors from $\sisi{C}{\Cc}^{\circ}$ to $(\Cat)$.

\begin{enonce*}{Definition 2.1.1}\etiquette[2.1.1]{XVII.def2.1.1}
If $X$ is a ringed site, $D(X)$ denotes the derived category
of the abelian category of sheaves of left modules on $X$.
\end{enonce*}

When confusion must be avoided, we instead write
$D(X,\Aa)$, where $\Aa$ denotes the sheaf of rings.

If\pageoriginale $f:X\rightarrow Y$ is a morphism of ringed sites,
$f$ defines derived functors
\begin{itemize}
\item[(1)] $Rf_{\ast}:D^{+}(X)\rightarrow D^{+}(Y)$ and

\item[(2)] $Lf^{\ast}:D^{-}(Y)\rightarrow D^{-}(X)$,
\end{itemize}
and, on passing to the opposite categor\sisi{y}{ies}, the functors
\begin{itemize}
\item[$\rm(1^{\circ})$] $Rf_{\ast}:D^{+}(X)^{\circ}\rightarrow
D^{+}(Y)^{\circ}$

\item[$\rm(2^{\circ})$] $Lf^{\ast}:D^{-}(Y)^{\circ}\rightarrow
D^{-}(X)^{\circ}$,
\end{itemize}
``functorial'' in $f$. More precisely, $(1^{\circ})$ and
$(2^{\circ})$ define two pseudofunctors from the category of
sites to that of triangulated categories, one covariant and the other
contravariant. One thereby obtains two categories over the category of sites,
one cofibered, with fibers $D^{+}(X)^{\circ}$ and direct-image functors
$Rf_{\ast}$, and the other fibered, with fibers
$D^{-}(X)^{\circ}$ and inverse-image functors
$\sisi{\LL}{L} f^{\ast}$. If the functor $f^{\ast}$ has finite
homological dimension (\ie has finite ``\tTor-dimension''), then in
(2) the superscript may be replaced by a $+$. Thus one obtains two
categories over the category $\Ss_{\tdf}$ o\sisi{f}{f} ringed sites and
morphisms of ringed sites of finite \tTor-dimension; the fibers of these
categories are the categories $D^{+}(X)^{\circ}$. The ``trivial
duality theorem'' of \textsc{Verdier} (\cite{sga4.exp.XVII.bib.CD} p.~48) states that these two
categories are canonically isomorphic; they are therefore fibered and
cofibered; the ``direct-image'' functors are the functors
$Rf_{\ast}$\sisi{}{ }(sic), and the ``inverse-image'' functors are the
functors $Lf^{\ast}$\sisi{}{ }(sic).

\setcounter{subsection}{1}
\setcounter{subsubsection}{1}
\subsubsection{}\etiquette[2.1.2]{XVII.subsubsec2.1.2}

The language of fibered categories readily resolves the
perplexity raised by \textsc{Artin} in \hyperrefext[XII][XII.sec4]{4}: it suffices to apply the
following proposition to the category fibered and cofibered over the category of
sites\pageoriginale whose fibers are the opposites of the categories of sheaves
on the corresponding sites (\cf \hyperlink{sga4.exp.XVII.subsection.0.2}{0.2}).

Let $\Ee$ be a category over $\Bb$, let $f:y\rightarrow x$
be an arrow of $\Bb$, and let $F\in \Ob(\Ee_x)$. Recall
that a pair $(G,\varphi)$ $(G\in \Ob \Ee_y, \varphi\in
\Hom_f(G,F))$ is a \emph{strict inverse image} of $F$
under $f$ if, for all $g:\sisi{Z}{z}\rightarrow y$ and $H\in \Ob
\Ee\sisi{_Z}{_z}$, one has
\[
\Hom_g(H,G)\xrightarrow[\sim]{\varphi}\Hom_{fg}(H,F){\quoi}.
\]
Similarly for direct images. If $\Ee$ is fibered or cofibered
over $\Bb$, this notion reduces to that of inverse image.

\begin{enonce*}{Proposition 2.1.3}\etiquette[2.1.3]{XVII.prop2.1.3}
Let $\epsilon$ be a category over $\Bb$, let $y\in \Ob
\Bb$, let $F\in \Ob \Ee_y$, and consider a commutative diagram in
$\Bb$:
$$\etiquette[2.1.3.1]{XVII.eq2.1.3.1}
\vcenter{
\xymatrix{
y'  \ar[r]^{g'} \ar[d]_{f'}& y\ar[d]^{f}         \\
x'          \ar[r]    & x.
}}\leqno{(2.1.3.1)}
$$
Assume that the direct and inverse images $f_{\ast}F$,
$g^{\ast}f_{\ast}F$, ${g'}^{\ast}F$, and $f'_{\ast}{g'}^{\ast}F$
exist in the strict sense. Then there exists a unique arrow
$\varphi \in \Fl \Ee_{x'}$, called the \emph{base-change
arrow}, making the diagram
$$\etiquette[2.1.3.2]{XVII.eq2.1.3.2}
\vcenter{
\xymatrix{
                             &  g'^{\ast}F \ar[ld]  \ar[r]    & F \ar[d]\\
f'_{\ast}g'^{\ast}F \ar[rd]_{\varphi}  &                                & f_{\ast}F\\
  & g^{\ast}f_{\ast}F \ar[ru]&.
}}\leqno{(2.1.3.2)}
$$
commutative. If\pageoriginale, in addition, the direct images $g'_{\ast}{g'}^{\ast
}F$ and $g'_{\ast}f'_{\ast}{g'}^{\ast}F$ exist in the strict sense, then the arrow $\varphi$ coincides with
the one defined in \emph{\hyperrefext[XII][XII.sec4]{4}}. If ${f'}^{\ast}f'_{\ast}{g'}^{\ast}F$
and ${f'}^{\ast}f'_{\ast}F$ exist in the strict sense, it also coincides
with the $2$\sisi{nd}{\textsuperscript{nd}} arrow defined in \emph{\hyperrefext[XII][XII.sec4]{4}}.

In particular, the two base-change arrows of
\emph{\hyperrefext[XII][XII.sec4]{4}} coincide.
\end{enonce*}

By hypothesis, the arrows in diagram \eqhyperref[XVII.eq2.1.3.2]{2.1.3.2} induce
bijections
\sisi{\begin{align*}
& \Hom_y(F,F)=\Hom_{g'}({g'}^{\ast}F,F)=\Hom_{fg'}({g'}^{\ast}F,f_{\ast}F)
\quad\mbox{\rm and}\\
&
\Hom_{x'}(f'_{\ast}{g'}^{\ast}F,g^{\ast}f_{\ast}F)=\Hom_{f'}({g'}^{\ast}F,
g^{\ast}f_{\ast}F)=\Hom_{gf'}({g'}^{\ast}F, f_{\ast}F){\quoi}.
\end{align*}}{\[\Hom_{x'}(f'_{\ast}{g'}^{\ast}F,g^{\ast}f_{\ast}F)=\Hom_{f'}({g'}^{\ast}F,
g^{\ast}f_{\ast}F)=\Hom_{gf'}({g'}^{\ast}F, f_{\ast}F).\]}
Since $fg'=gf'$, $\Hom_{fg'}({g'}^{\ast}F,f_{\ast}F)$ is identical
to $\Hom_{gf'}({g'}^{\ast}F,f_{\ast}F)$, and the arrows of diagram
\sisi{(2.1.2.2)}{\eqhyperref[XVII.eq2.1.3.2]{2.1.3.2}} induce \sisi{bijections}{maps}
\[
\Hom_y(F,F)\sisi{=}{\to}\Hom_{fg'}({g'}^{\ast}F,f_{\ast}F)=\Hom_{x'}(f'_{\ast}{g'}^{\ast}F,g^{\ast}f_{\ast}F){\quoi},
\]
and the image of $1_F$ is the unique arrow of $\Ee_{x'}$ making

% END INCLUDED SGA4_English_current_buildable_source_20260906/source/17/drafts/SGA4_XVII_LINES_1001_1400_English.texfrag

% BEGIN INCLUDED SGA4_English_current_buildable_source_20260906/source/17/drafts/SGA4_XVII_LINES_1401_1800_English.texfrag
\eqhyperref[XVII.eq2.1.3.2]{2.1.3.2} commutative.%See the Erratum in SGA 4 1/2, p. 312.

By definition, in \hyperrefext[XII][XII.sec4]{4}, the base change arrow $\psi$
was the one making the following diagram commutative
$$\etiquette[2.1.3.3]{XVII.eq2.1.3.3}
\vcenter{
\xymatrix{
f'_{\ast}g'{}^{\ast}F \ar[d]^{\psi}  & g'{}^{\ast}F \ar[d]^{u_4} \ar[l] \ar@{=}[r] &
g'{}^{\ast}F \ar[r] \ar[d]^{u_3}   & F \ar@{=}[r] \ar[d]^{u_2}&  F \ar[d]^{u_1}\\
g'_{\ast}f_{\ast}F   & {f'}^{\ast}g^{\ast}f_{\ast}F  \ar[l] \ar@{=}[r] &
g'f^{\ast}f_{\ast}F   \ar[r] & f^{\ast}f_{\ast}F \ar[r] &  f_{\ast}F.
}}\leqno{(2.1.3.3)}
$$
The horizontal arrows and $u_1$ are the canonical arrows; $u_2$
is therefore the adjunction arrow, $u_3={g\sisi{^{1}}{'}}^{\ast}(u_2)$, $u_4$ is
deduced from the isomorphism between ${f'}^{\ast}g^{\ast}$ and
${g'}^{\ast}f^{\ast}$, and the arrow $\psi$ is deduced from $u_4$ by
adjunction.

The composite\pageoriginale map with source ${g'}^{\ast}F$ and
target $f_{\ast}F$ is the one that appears in \eqhyperref[XVII.eq2.1.3.2]{2.1.3.2}, and the
same therefore holds in the following commutative diagram, whose left-hand
square is taken from \sisi{(2.1.2.3)}{\eqhyperref[XVII.eq2.1.3.3]{2.1.3.3}}
\[
\xymatrix{
g'{}^{\ast}F \ar[r]^{u_4} \ar[d]& f'{}^{\ast}g{}^{\ast}f_{\ast}F
\ar[r] \ar[d]&
f_{\ast}F \ar@{=}[d]\\
f'_{\ast}g'{}^{\ast}F  \ar[r]^{\psi} & g^{\ast}f_{\ast}F \ar[r] &
f_{\ast}F\rlap{.}
}
\]
This means that $\psi$ satisfies the defining property of
$\varphi$.

The definition of $\varphi$ is self-dual, so there is no need to prove
the second part of \hyperref[XVII.prop2.1.3]{2.1.3}.

\subsection{Categories fibered in triangulated categories}\etiquette[2.2]{XVII.subsec2.2}

\begin{enonce*}{Definition 2.2.1}\etiquette[2.2.1]{XVII.def2.2.1}
An “additive category over a category $\Bb$”
is a category $\Aa$ over $\Bb$ whose sets of
arrows $\Hom_f(X,Y)$ (for $f\in \Fl (\Bb)$, with $X$ (\resp
$Y$) over the source (\resp the target) of $f$) have been endowed with
abelian group structures, in such a way that:
\begin{enumerate}
\renewcommand{\theenumi}{\roman{enumi}}
\renewcommand{\labelenumi}{(\theenumi)}
\item For any $x\xrightarrow{f}y\xrightarrow{g}z$ in
$\Bb$ and $X$, $Y$, $Z$ in $\Aa$ lying over $\sisi{s}{x}$,
$y$, and $z$, respectively, the composition law
\[
\Hom_g(Y,Z)\times \Hom_f
(X,Y)\xrightarrow{\circ}\Hom_{gf}(X,Z)
\]
is biadditive.

\item The fiber categories are additive.
\end{enumerate}
\end{enonce*}

We shall not need here the fact that these abelian group structures, when
they exist, are uniquely determined by
$\Aa$, $\Bb$, and $\rho$.\nde{here $\rho$ denotes the structural functor $\Aa\to \Bb$.} A $\Bb$-functor
$F:\Aa_1\rightarrow \Aa_2$ between additive categories
over $\Bb$ is said to be \emph{additive} if it induces
homomorphisms
\[
F:\Hom_{\Aa_1,f}(X,Y)\longrightarrow
\Hom_{\Aa_{2},f}(X,Y){\quoi}.
\]

\begin{enonce*}{Definition 2.2.2}\etiquette[2.2.2]{XVII.def2.2.2}
An\pageoriginale “abelian category over a category
$\Bb$” is an additive category $\Aa$ over
$\Bb$ whose fibers are abelian, and such that for
every arrow $f:x\rightarrow y$ of $\Bb$ the bifunctor
\[
\Hom_f(X,Y):\Aa^{\circ}_x\times \Aa_y\longrightarrow
(\Ab)
\]
is left exact in $X$ and $Y$.
\end{enonce*}

We shall not need here the fact that if $\Aa$ is a category
fibered and cofibered over $\Bb$ (\SGA 1 VI 6), it is abelian
over $\Bb$ as soon as its fibers are abelian categories.

\begin{enonce*}{Definition 2.2.3}\etiquette[2.2.3]{XVII.def2.2.3}
A “triangulated category over a category $\Bb$” is an additive category $\Aa$ over $\Bb$, endowed
with a functor $T:\Aa\rightarrow \Aa$, and whose
fibers $\Aa_x$, for $x\in \Ob \Bb$, are endowed
with sets $\Delta_x$ of triangles, in such a way that:
\begin{enumerate}
\renewcommand{\theenumi}{\roman{enumi}}
\renewcommand{\labelenumi}{(\theenumi)}
\item $T$ is an additive $\Bb$-automorphism of $\Aa$.

\item Each fiber category $\Aa_x$, endowed with the functor $T_x$
induced by $T$ and with the set $\Delta_x$, is a triangulated
category \emph{(\cite{sga4.exp.XVII.bib.CD} \sisi{}{I 1.}1.1)}.

\item If the triangles $(X,Y,Z,u,v,w)$ and $(X',Y',Z',u',v',w')$
belong to $\Delta_x$ and $\Delta_{x'}$, respectively, every
diagram
\[
\xymatrix{
            & Z   \ar[ldd]_{w}&            &  &           & Z'   \ar[ldd]|\hole_(.3){w'} & \\
            &   & Y\ar[lu]_{v} \ar[rrrr]_(.4){\sisi{j}{g}}&  &     &       &Y'\ar[lu]_{v'} \\
X   \ar[rru]^{u} \ar[rrrr]_{f}&             &          &  &   X'   \ar[rru]^{u'} &              &
}
\]
in which $u'f=gu$ can be completed to a morphism of triangles
$(f,g,h):(X,Y,Z)\rightarrow (X',Y',Z')$.
\end{enumerate}
\end{enonce*}

$T$\pageoriginale will be called the translation functor; we shall
often write $A[1]$ rather than $T(A)$, with $A[n]$ then denoting $T^{n}(A)$
$(n\in\ZZ)$.

The triangles belonging to the sets $\Delta_x$ will be called
\emph{distinguished} or \emph{exact}.

\begin{enonce*}{Proposition 2.2.4}\etiquette[2.2.4]{XVII.prop2.2.4}
Let $\Aa$ be a triangulated category over $\Bb$, let $x$
be an object of $\Bb$, and let $(X,Y,Z,u,v,w)$ be a distinguished triangle of
$\Aa_x$. For every $f:y\rightarrow x$ and every object $A\in
\Ob \Aa_y$, the following infinite sequence is exact:
\begin{align*}
\ldots  \longrightarrow \Hom_f(A,X) \longrightarrow \Hom_f(A,Y)&\longrightarrow
\Hom_f(A,Z)\\
 & \longrightarrow \Hom_f(A,X[1])\longrightarrow
\Hom_f(A,Y[1])\longrightarrow\ldots
\end{align*}
\end{enonce*}

There is an analogous statement for $f:x\rightarrow y$ and $A\in \Ob \Aa_y$.

We leave it to the reader to prove this proposition by
paraphrasing \textsc{Verdier} \cite{sga4.exp.XVII.bib.CD} p.~4. We also leave it to the reader to give
a meaning to the following proposition and to verify it:

\begin{enonce*}{Proposition 2.2.5}\etiquette[2.2.5]{XVII.prop2.2.5}
Let $\Aa$ be an additive (\resp abelian, \resp
triangulated) category over $\Bb$. For every functor
$\Bb'\rightarrow \Bb$ with target $\Bb$, the
category $\Aa'=\fprod{\Aa}{\Bb'}{\Bb}$ is
additive (\resp abelian, \resp triangulated) over $\Bb$.
\end{enonce*}

\begin{enonce*}[remark]{Example 2.2.6}\etiquette[2.2.6]{XVII.exem2.2.6}
The category of sheaves of modules on varying ringed sites
is abelian over the category of ringed sites.
\end{enonce*}

\begin{enonce*}[remark]{Example 2.2.7}\etiquette[2.2.7]{XVII.exem2.2.7}
Let $\Aa$ be an additive category over $\Bb$. The
category $C(\Aa/\Bb)$ of complexes of $\Aa$
over $\Bb$ has as objects the complexes of the fiber
categories. An arrow $f$ from $(X^n,d)$ to $(Y^n,d)$ is a system
of\pageoriginale arrows $f^{n}:X^{n}\rightarrow Y^{n}$ such that
$\sisi{fd}{df}^{n}=f^{n+1} d$. We set $p(f)=p(f^{n})$ (this arrow does not depend
on $n$, since the $p(d)$ are
identities). $C(\Aa/\Bb)$ is an additive category
over $\Bb$. A “translation” functor is defined by setting
\[
(K[1])^{n}=K^{n+1}\quav (d[1])\sisi{_n}{^n}=-d\sisi{_{n+1}}{^{n+1}}{\quoi}.
\]
\end{enonce*}

\begin{enonce*}[remark]{Example 2.2.8}\etiquette[2.2.8]{XVII.exem2.2.8}
Let $\Aa$ be an additive category over $\Bb$. \emph{The
category $K(\Aa/\Bb)$ of complexes of $\Aa$
over $\Bb$ modulo homotopy} has the same objects as
$C(\Aa/\Bb)$. It is obtained by declaring the arrows homotopic
to zero to be zero in $C(\Aa/\Bb)$; in
other words, for $f:x\rightarrow y$ in $\Bb$,
$X^{\bigdot}\in \Ob C(\Aa/\Bb)_x$ and $Y^{\bigdot}\in
\Ob C(\Aa/\Bb)_y$, we set
\[
\Hom_f(X^{\bigdot},
Y^{\bigdot})=H^{0}\Hom^{\bigdot}_f(X^{\bigdot},Y^{\bigdot})
\]
\end{enonce*}

The translation functor descends to the quotient.

If $F:\Bb'\rightarrow \Bb$ is a functor with target
$\Bb$, and if
$\Aa'=\fprod{\Aa}{\Bb'}{\Bb}$, we
leave it to the reader to identify $C(\Aa'/\Bb')$
with $\fprod{C(\Aa/\Bb)}{\Bb'}{\Bb}$ and
$K(\Aa'/\Bb')$ with
$\fprod{K(\Aa/\Bb)}{\Bb'}{\Bb}$. In
particular, the fibers of the category $K(\Aa/\Bb)$
are the categories $K(\Aa_x)$ and, as such, are
triangulated.

\begin{enonce*}{Proposition 2.2.9}\etiquette[2.2.9]{XVII.prop2.2.9}
\begin{enumerate}
\renewcommand{\theenumi}{\roman{enumi}}
\renewcommand{\labelenumi}{(\theenumi)}
\item With the additional structures defined in \hyperref[XVII.exem2.2.8]{2.2.8}, the
category $K(\Aa/\Bb)$ is triangulated over $\Bb$.

\item The formation of the additive (\resp triangulated)
category $C(\Aa/\Bb)$ (\resp $K(\Aa/\Bb)$) is
compatible with every change of base category
$\Bb'\rightarrow \Bb$.
\end{enumerate}
\end{enonce*}

Conditions\pageoriginale (i) and (ii) of Definition \hyperref[XVII.def2.2.3]{2.2.3} are
verified fiber by fiber. Condition (iii) is verified by paraphrasing \textsc{Verdier}
\cite{sga4.exp.XVII.bib.CD} p.~4.

\begin{enonce*}[remark]{Remark 2.2.10}\etiquette[2.2.10]{XVII.rem2.2.10}
The preceding discussion also applies to the categor\sisi{y}{ies}
$K^{+}(\Aa/\Bb)$, $K^{-}(\Aa/\Bb)$, and
$K^{b}(\Aa/\Bb)$ obtained by restricting
respectively to complexes bounded below, bounded above, or bounded.
\end{enonce*}

\begin{enonce*}[remark]{Remark 2.2.11}\etiquette[2.2.11]{XVII.rem2.2.11}
Taking $\Bb$ to be the final category, one recovers the
usual categories of complexes.
\end{enonce*}

\setcounter{subsection}{2}
\setcounter{subsubsection}{11}
\subsubsection{}\etiquette[2.2.12]{XVII.subsubsec2.2.12}
Let $F$ be a pseudofunctor from a category $\Bb^{\circ}$ to
$(\Cat)$. It is known (\SGA 1 VI) that $F$ defines a category fibered over
$\Bb$, whose fibers are the categories $F(x)$. For any
$f:x\rightarrow y$ in $\Bb$ and $X\in \Ob F(\sisi{n}{x})$, $Y\in \Ob
F(y)$, we have by definition
$$\etiquette[2.2.12.1]{XVII.eq2.2.12.1}
\Hom_f(X,Y)=\Hom_{F(x)}(X,F(f)(Y)){\quoi}.\leqno(2.2.12.1)
$$
It is known (ibid.) that the preceding construction defines an
equivalence between the 2-category of pseudofunctors from
$\Bb^{\circ}$ to $(\Cat)$ and the 2-category of categories
fibered over $\Bb$.

Let us start with a pseudofunctor from $\Bb^{\circ}$ to the
2-category whose objects are additive (\resp
abelian, \resp triangulated) categories and whose 1-arrows are additive
(\resp left exact, \resp triangulated) functors. One immediately verifies that the
category defined by \sisi{(2.1.12.1)}{\eqhyperref[XVII.eq2.2.12.1]{2.2.12.1}} is additive (\resp abelian, \resp
triangulated). Conversely:

\begin{enonce*}{Proposition 2.2.13}\etiquette[2.2.13]{XVII.prop2.2.13}
The\pageoriginale preceding construction establishes an equivalence between
\begin{enumerate}
\renewcommand{\theenumi}{\alph{enumi}}
\renewcommand{\labelenumi}{(\theenumi)}
\item The $2$-category of pseudofunctors from $\Bb^{\circ}$
to the $2$-category whose objects are additive
(\resp abelian, \resp triangulated) categories and whose $1$-arrows are
additive (\resp left exact, \resp triangulated) functors.

\item The $2$-category of additive (\resp abelian,
\resp triangulated) categories over $\Bb$ that are fibered over $\Bb$.
\end{enumerate}
\end{enonce*}

The theory developed in \SGA 1 VI shows that it suffices to verify the
following point:

\setcounter{subsubsection}{13}
\setcounter{paragraph}{0}
\paragraph{}\etiquette[2.2.13.1]{XVII.par2.2.13.1}
Let $\Aa$ be an additive (\resp abelian, \resp
triangulated) category over $\Bb$, let $f:x\rightarrow y$ be an arrow of
$\Bb$, and suppose that an inverse image functor
$f^{\ast}:\Aa_y\rightarrow \Aa\sisi{_n}{_x}$ exists; then $f^{\ast}$
is additive (\resp left exact, \resp triangulated).

The functor $f^{\ast}$ gives rise to isomorphisms
\[
\Hom_f(X,Y)\simeq \Hom_x (X,f^{\ast}(Y)){\quoi}.
\]
In the triangulated case\sisi{}{}, it remains to prove that $f^{\ast}$ is triangulated
for the isomorphism $Tf^{\ast}=f^{\ast}T$ (where $T$ is the
translation functor) that makes the following diagrams commutative:
$$\etiquette[2.2.13.2]{XVII.eq2.2.13.2}
\vcenter{
\xymatrix{
\Hom_f(X,Y)  \ar@{=}[r] \ar@{=}[d]& \Hom_x(X,f^{\ast}Y)\ar@{=}[rd]\\
\Hom_f(TX,TY)\ar@{=}[rd]           &      & \Hom_x(TX,Tf^{\ast}Y) \\
             &      \Hom_x(TX,f^{\ast}TY)\ar@{=}[ru] &
}}\leqno{(2.2.13.2)}
$$

\begin{itemize}
\item[\emph{Additive case}:] Consider\pageoriginale a diagram
$$\etiquette[2.2.13.3]{XVII.eq2.2.13.3}
\vcenter{
\xymatrix@C=2.5cm@R=1.75cm{
f^{\ast} X \ar@<-3pt>[d]_{f^{*}(u)}\ar@<3pt>^{f^{\ast}(v)}[d]        \ar[r]^{p} & X
\ar@<-3pt>_{u}[d] \ar@<3pt>^{v}[d]\\
f^{\ast} X' \ar[r]^{p'}& X'\rlap{.}
}}\leqno{(2.2.13.3)}
$$
We have
$p'f^{\ast}(u+v)=(u+v)p=up+vp=p'f^{\ast}(u)+p'f^{\ast}(v)=p'(f^{\ast}(u)+f^{\ast}(v))$
and hence $f^{\ast}(u+v)=f^{\ast}(u)+f^{\ast}(v)$.

\item[\emph{Abelian case}:] by definition (\hyperref[XVII.subsec2.2]{2.2}), the functor
$\Hom(X,f^{\ast}Y)=\Hom_f(X,Y)$ is left exact in $Y$, for every
$X$, whence $f^{\ast}$ is left exact.

\item[\emph{Triangulated case}:] let $(Y',Y,Y'')$ be a distinguished triangle in
$\Aa_y$ and let \sisi{$(f^{\ast}Y',f^{\ast}Y)$}{$(f^{\ast}Y',f^{\ast}Y, X'')$ be a distinguished triangle in $\Aa_x$}. By \hyperref[XVII.def2.2.3]{2.2.3},
the canonical morphism between the bases of these triangles can be
completed to a morphism of triangles, giving rise to the following
diagram:
\[
\xymatrix@C=2em{
&  X''       \ar[rrr]^{\sisi{}{s}}\ar[ldd]   &     & &**[r]f^{\ast}Y''\ar[rrr]\ar[ldd]   &  & & Y''\ar[ldd] \\                      &     & f^{\ast}Y   \sisi{\ar[rrr]|\hole}{\ar@{=}[rrr]|\hole}\ar[lu]& &      &f^{\ast}Y \ar[rrr]|\hole\sisi{}{\ar[ul]}&&& Y\ar[lu]\\
  f^{\ast}Y' \sisi{\ar[rrr]}{\ar@{=}[rrr]}\ar[rru]         &&& f^{\ast}Y'\sisi{}{\ar[rru]}   \ar[rrr] && & Y'\ar[rru]&&.
}
\]

Let us prove that $s$ is an isomorphism, which will complete the
proof. For every $X$ in $\Ob \Aa_x$, we have the
commutative diagram:
\[
\text{\scriptsize
$\xymatrix@C=1em{
\Hom_x(X,f^{\ast}Y') \ar[r] \ar@{=}[d]& \Hom_x(X,f^{\ast}Y) \ar@{=}[d]\ar[r]& \Hom_x(X,X'')\ar[d]
\ar[r]& \Hom_x(X,f^{\ast}Y'[1])\ar[r]\ar@{=}[d]
& \Hom_x(X,f^{\ast}Y[1]) \ar@{=}[d]\\
\Hom_x(X,f^{\ast}Y') \ar[r]\ar@{=}[d]&  \Hom_x(X,f^{\ast}Y)\ar[r] \ar@{=}[d]&
\Hom_x(X,f^{\ast}Y'') \ar[r]\ar@{=}[d]&  \Hom_x(X,f^{\ast}Y'[1])\ar@{=}[d] \ar[r]&
\Hom_x(X,f^{\ast}Y[1]) \ar@{=}[d]\\
\Hom_f(X,Y)  \ar[r]& \Hom_f(X,Y)\ar[r]  &
\Hom_f(X,Y'') \ar[r]& \Hom_f(X,Y'[1])\ar[r] &
\Hom_x(X,Y[1]).
}$}
\]

The\pageoriginale first and last rows are exact by
\hyperref[XVII.prop2.2.4]{2.2.4}. By the five lemma, the map induced by $s$ from
$\Hom(X,X'')$ to $\Hom(X,f^{\ast}Y'')$ is always bijective, and $s$
is therefore an isomorphism.
\end{itemize}

There are, of course, dual results for cofibered categories.

\subsection{Trivial duality formula}\etiquette[2.3]{XVII.subsec2.3}

Let $\Aa$ be a category over a base category~$\Bb$,
and let $S$ be a set of arrows of $\Aa$, each of which
projects onto an identity of $\Bb$. The category
$\Aa(S^{-1})$ obtained from $\Aa$ by the calculus of fractions
is a category over $\Bb$ (by the universal property).

Let $\Bb'\rightarrow \Bb$ be a functor; denote by
$\Aa'$ the fiber product
$\fprod{\Aa}{\Bb'}{\Bb}$ and by $S'$ the inverse
image of $S$ in $\Aa'$. The composite functor
$\Aa'\rightarrow \Aa\rightarrow \Aa(S^{-1})$
makes the arrows of $S'$ invertible and hence factors through
$\Aa'({S'}^{-1})$. We thus obtain a canonical $\Bb'$-functor
as follows:
$$\etiquette[2.3.1]{XVII.eq2.3.1}
\Aa'(S^{-1})\longrightarrow
\fprod{\Bb'}{\Aa(\sisi{S}{{S'}}^{-1})}{\Bb}{\quoi}.\leqno(2.3.1)
$$

\begin{enonce*}{Lemma 2.3.2}\etiquette[2.3.2]{XVII.lem2.3.2}
If $\Bb'$ is a subcategory of $\Bb$, and if
whenever a composite arrow $h=f\circ g$ of $\Bb$
belongs to $\Bb'$, both $f$ and $g$ belong to
$\Bb'$, then the functor \eqhyperref[XVII.eq2.3.1]{2.3.1} is an isomorphism.
\end{enonce*}

Note\pageoriginale that \eqhyperref[XVII.eq2.3.1]{2.3.1} always induces a bijection on
objects, and that, in the present case,
$\fprod{\Bb'}{\Aa(S^{-1})}{\Bb}$ is a
subcategory of $\Aa(S^{-1})$.

Recall the “explicit” construction of $\Aa(S^{-1})$
\cite[\sisi{Chap.}{\chap}~I]{sga4.exp.XVII.bib.GZ}. Given two objects $X$ and $Y$ of $\Aa$, one
considers the “formal composites” $f_1\ldots f_n$, where each
$f_i$ is either an arrow of $\Aa$ or the inverse of an
arrow of $S$. The set $\Hom_{\Aa(S^{-1})}(X,Y)$ is the
quotient of the set of these “formal composites” by the
equivalence relation generated by the relations
\begin{alignat*}{2}
& f_1\ldots uv\ldots f_n\equiv f_1\ldots (u\circ v)\ldots f_n &\quad &
\mbox{\rm for } u,v\;\mbox{\rm in } \Fl(\Aa)\\
& f_1\ldots ss^{-1}\ldots f_n\equiv f_1\ldots f_n && \mbox{\rm for }
s\in S\\
& f_1\ldots s^{-1}s\ldots f_n\equiv f_1\ldots f_n && \mbox{\rm for }
s\in S{\quoi}.
\end{alignat*}

% END INCLUDED SGA4_English_current_buildable_source_20260906/source/17/drafts/SGA4_XVII_LINES_1401_1800_English.texfrag

% BEGIN INCLUDED SGA4_English_current_buildable_source_20260906/source/17/drafts/SGA4_XVII_LINES_1801_2200_English.texfrag

If these constructions are applied to $\Aa$ and to the
subcategory $\Aa'$, one checks that $\Aa'(S^{-1})$
is identical to the inverse-image subcategory of $\Aa(S^{-1})$
of $\Bb'$, since if a ``formal composite'' $f_1\ldots f_n$ has as
its composite in $\Bb$ an arrow of $\Bb'$, the image of each $f_i$
will also be an arrow of $\Bb'$.

Suppose now that $\Aa$ is a triangulated category over $\Bb$, and
that the sets $S_x=S\cap \Fl(\Aa_x)$, for $x\in \Ob(\Bb)$, are
saturated multiplicative systems (\cite{sga4.exp.XVII.bib.CD} \sisi{Ch.}{\chap}~I
\S~2\sisi{. 1}{ nᵒ 1}).

\begin{enonce*}{Definition 2.3.3}\etiquette[2.3.3]{XVII.def2.3.3}
Under the preceding hypotheses, the category $\Aa/\Bb$ is said to be
almost derivable (\resp derivable, \resp coderivable) relative to $S$
if, for every functor $F:\Bb'\rightarrow\Bb$, the functor
\eqhyperref[XVII.eq2.3.1]{2.3.1} is an isomorphism (\resp and if
$\Aa(S^{-1})$ is fibered, \resp cofibered over $\Bb$).
\end{enonce*}

We\pageoriginale denote by $I$ the category of the following diagram:
$$\etiquette[2.3.4]{XVII.eq2.3.4}
I:\underset{0}{\times} \xrightarrow{F} \underset{1}{\times}{\quoi}.\leqno{(2.3.4)}
$$

By \hyperref[XVII.prop2.2.13]{2.2.13}, to give a triangulated category
$\Aa$ over $I$, fibered (\resp cofibered) over $\sisi{\Aa}{I}$, is
equivalently to give its fibers $\Aa_0$ and $\Aa_1$ and the
triangulated inverse-image functor
$F^{\ast}:\Aa_1\rightarrow\Aa_0$ (\resp and the triangulated
direct-image functor $F_{\ast}:\Aa_0\rightarrow\Aa_1$). If
$A_0\in\Ob\Aa_0$ and $A_1\in\Ob\Aa_1$, then by
\eqhyperref[XVII.eq2.2.12.1]{2.2.12.1},
\[
\Hom_F(A_0,A_1)=\Hom(A_0,F^{\ast}A_1)
\]
(\resp $\Hom_F(A_0,A_1)=\Hom(F_{\ast}A_0,A_1)$).

\begin{enonce*}{Proposition 2.3.5}\etiquette[2.3.5]{XVII.prop2.3.5}
Let $\Aa$ be a triangulated category over $I$ and let $S$ be as in
(\hyperref[XVII.def2.3.3]{2.3.3}).
\begin{enumerate}
\renewcommand{\theenumi}{\roman{enumi}}
\renewcommand{\labelenumi}{(\theenumi)}
\item $\Aa$ is almost derivable relative to $S$, and
$\Aa(S^{-1})$ is triangulated over $I$.

\item If $\Aa$ is fibered (\resp cofibered), the functor $F^{\ast}$
(\resp $F_{\ast}$) is right derivable (\resp left derivable) relative
to $S_0$ and $S_1$ (\hyperref[XVII.def1.2.1]{1.2.1}) if and only if
$\Aa$ is derivable (\resp coderivable) over $I$.
\end{enumerate}
\end{enonce*}

By Lemma \hyperref[XVII.lem2.3.2]{2.3.2}, the fiber categories of
$\Aa(S^{-1})$ are $\Aa_0(S^{-1}_0)$ and $\Aa_1(S^{-1}_1)$, where
$S_i=S\cap\Fl(\Aa_i)$ $(i=0,1)$. With the notation of \S
\hyperref[XVII.sec1]{1} ${\rm n^{\circ}}$ \hyperref[XVII.sec2]{2}, the
functors $A_0\rightarrow A^{-}_0$ and $A_1\rightarrow A^{+}_1$
identify $\Aa_0(S^{-1}_0)$ and $\Aa_1(S^{-1}_1)$ with categories of
\ppro- and \iind-objects of $\Aa_0$ and $\Aa_1$. This makes it possible
to define a category over $I$, with fibers $\Aa_0(S^{-1}_0)$ and
$\Aa_1(S^{-1}_1)$, by setting, for $A_0\in\Aa_0$ and $A_1\in\Aa_1$,
\[
\Hom_F(A_0,A_1)=\Hom_F(A^{-}_0,A^{+}_1)=\varinjlim_{A'_0{\xrightarrow[s]{}}A_{\sisi{p}{0}}}
\varinjlim_{A_1{\xrightarrow[t]{}}{A'_1}}\Hom(A'_0,A'_1){\quoi},
\]
where\pageoriginale $s$ and $t$ range over the elements of $S_0$ and
$S_1$ with target (respectively source) $A_0$ (respectively $A_1$).
The category $\Aa$ maps to this category, and one checks immediately
that it satisfies the universal property of $\Aa(S^{-1})$, over $I$
and after every base change.

To verify (i), it suffices to prove that this category, namely
$\Aa(S^{-1})$, is triangulated, and it suffices to verify condition
\hyperref[XVII.def2.2.3]{2.2.3}~(iii). Let $(X_0,Y_0,Z_0)$ and
$(X_1,Y_1,Z_1)$ therefore be distinguished triangles of $\Aa_0$ and
$\Aa_1$, giving rise to the distinguished triangles
$(X^{-}_0,Y^{-}_0,Z^{-}_0)$ and $(X^{+}_1,Y^{+}_1,Z^{+}_1)$ of
$\Aa_0(S^{-1}_0)$ and $\Aa_1(S^{-1}_1)$.
$$\etiquette[2.3.5.1]{XVII.eq2.3.5.1}
\vcenter{
\xymatrix{
X^{-}_{0}  \ar[r]^{u} \ar[d]& X^{+}_1  \ar[d]\\
Y^{-}_{0}  \ar[r]^{v} & Y^{+}_1\rlap{.}\\
}}\leqno{(2.3.5.1)}
$$

By Lemma \hyperref[XVII.lem1.2.2.1]{1.2.2.1}, the triangles
$(X^{-}_0,Y^{-}_0,Z^{-}_0)$ and $(X^{+}_1,Y^{+}_1,Z^{+}_1)$ are limits
of distinguished triangles; there therefore exist morphisms of
distinguished triangles
\begin{equation*}
\begin{split}
(X^{-}_0,Y^{-}_0,Z^{-}_0) &\longrightarrow (X'_0, Y'_0,Z'_0)\\
(X'_1, Y'_1, Z'_1) &\longrightarrow (X^{+}_1, Y^{+}_1, Z^{+}_1){\quoi},
\end{split}\tag{2.3.5.2}
\end{equation*}\etiquette[2.3.5.2]{XVII.eq2.3.5.2}%Removing this comment would introduce unwanted white space.
and\pageoriginale a diagram
\[
\xymatrix{
X'_{0}  \ar[r]^{u'} \ar[d]& X'_1 \ar[d] \\
Y'_{0}  \ar[r]^{v'} & Y'_1\rlap{,}
}
\]
compatible with \eqhyperref[XVII.eq2.3.5.1]{2.3.5.1}. Since $\Aa$ is
triangulated over $I$, this diagram can be extended to a morphism of
triangles
\[
(X'_0,Y'_0,Z'_0)\longrightarrow (X'_1, Y'_1, Z'_1){\quoi}.
\]
Composing the latter with the morphisms
\eqhyperref[XVII.eq2.3.5.2]{2.3.5.2}, we see that diagram
\eqhyperref[XVII.eq2.3.5.1]{2.3.5.1} can be extended to a morphism of
triangles, which proves \hyperref[XVII.def2.2.3]{2.2.3}~(iii).

Let us verify (ii), assuming for example that $\Aa$ is fibered over
$I$. If $A_1\in\Ob(\Aa_1(S^{-1}_1))$, an inverse image of $A_1$ in
$\Aa(S^{-1})$ is an object $B$ of $\Aa_0(S^{-1}_0)$ such that, in
$\Aa_0(S^{-1}_0)$,
\begin{align*}
\Hom_{\Aa_0(S^{-1}_0)}(A_0,B) & =\Hom_F(A^{-}_0,
A^{+}_1)=\Hom(A^{-}_0,
F^{\ast}(A^{+}_1))\\
&=\Hom_{\Aa_0(S^{-1}_0)}
(A_0,F^{\ast}(A^{+}_1)){\quoi}.
\end{align*}
Such an object $B$ exists if and only if the \iInd-object
$F^{\ast}(A^{+}_1)$ of $\Aa_0(S^{-1}_0)$ is essentially constant
(\ie representable), and $B$ then represents this \iind-object, which
is precisely $RF(A_1)$ (\hyperref[XVII.def1.2.1]{1.2.1}); this proves
the assertion.

\begin{enonce*}{Corollary 2.3.6}\etiquette[2.3.6]{XVII.cor2.3.6}
With the notation of \hyperref[XVII.def2.3.3]{2.3.3}, if $\Aa$ is
almost derivable over $\Bb$, then $\Aa(S^{-1})$ is a triangulated
category over $\Bb$.
\end{enonce*}

First, the fiber category $\Aa(S^{-1})_x$ can be identified with
$\Aa_x(S^{-1}_x)$: apply the hypothesis to the base change
$\{x\}\rightarrow\Bb$. This\pageoriginale category is triangulated,
equipped with a set $\Delta_x$ of triangles, which gives a meaning to
the preceding statement. To verify condition \hyperref[XVII.def2.2.3]{2.2.3}~(iii), it suffices
to do so after every base change $I\rightarrow\Bb$, and the assertion
follows from \hyperref[XVII.prop2.3.5]{2.3.5}~(i).

It follows immediately from the proof of
\hyperref[XVII.prop2.3.5]{2.3.5}~(ii) that:

\begin{enonce*}{Theorem 2.3.7}\etiquette[2.3.7]{XVII.thm2.3.7}
\emph{(trivial duality formula).}

Let $\Aa$ be a triangulated category fibered and cofibered over $I$,
and denote by $F_{\ast}$ and $F^{\ast}$ the direct- and inverse-image
functors. Let $S=S_0\cup S_1$ be a multiplicative system as in
\hyperref[XVII.def2.3.3]{2.3.3}. Suppose that $LF_{\ast}$ is defined
at $A_0\in\Ob\Aa(S^{-1}_0)$ and that $RF^{\ast}$ is defined at
$A_1\in\Ob\Aa_1(S^{-1}_1)$. Then
\[
\Hom_{\Aa_0(S^{-1}_0)}(A_0,RF^{\ast}A_1)=\Hom_{\Aa\sisi{_{(S^{-1})}}{(S^{-1})}}(A_0,A_1)=\Hom_{\Aa_1(S^{-1}_1)}(LF_{\ast}A_0,A_1){\quoi}.
\]
\end{enonce*}

The positions of the symbols $R$, $L$, ${}^{\ast}$, and ${}_{\ast}$
in \hyperref[XVII.thm2.3.7]{2.3.7} are only apparently anomalous: in
fact, in the category \hyperlink{sga4.exp.XVII.subsection.0.2}{0.2} of sheaves on variable sites, fibered and
cofibered, the fiber categories are the opposites of the usual
categories of sheaves.

\subsection{Categories fibered in derived categories}\etiquette[2.4]{XVII.subsec2.4}

\setcounter{subsubsection}{-1}
\subsubsection{}\etiquette[2.4.0]{XVII.subsubsec2.4.0}
In this ${\rm n^{\circ}}$, $\Aa$ denotes a triangulated category
cofibered over $\Bb$, and $S$ a multiplicative system of morphisms of
$\Aa$, all lying over an identity, such that
$S_x=S\cap\Fl(\Aa_x)$ is a saturated multiplicative system of $\Aa_x$
\cite[\sisi{}{I }1.2.1]{sga4.exp.XVII.bib.CD} for every $x\in\Ob(\Bb)$.

We\pageoriginale seek a criterion for $\Aa$ to be coderivable relative
to $S$ (\hyperref[XVII.def2.3.3]{2.3.3}); this implies that
$\Aa(S^{-1})$ is triangulated over $\Bb$ when it is equipped with the
set of triangles isomorphic to the image of a distinguished triangle
of $\Aa$ (\hyperref[XVII.prop2.3.5]{2.3.5}).

\begin{enonce*}{Proposition 2.4.1}\etiquette[2.4.1]{XVII.prop2.4.1}
Let $\Aa$, $\Bb$, and $S$ be as above. For $\Aa$ to be coderivable
relative to $S$, it suffices that the following condition hold:
\begin{itemize}
\item[$(\ast)$] $\forall x\in\Ob\Bb$, $\forall A\in\Ob(\Aa_x)$,
$\exists s:A'\rightarrow A$ in $S$\sisi{}{,} $\forall g$,
$f:x\xrightarrow{f}y\xrightarrow{g}\sisi{Z}{z}$, $f_{\ast}A'$ is
left adapted to the functor $g_{\ast}$.
\end{itemize}
\end{enonce*}

Let $\Aa^{1}$ denote the full subcategory of $\Aa$ consisting of the
objects all of whose direct images are adapted to all direct-image
functors. Then \hyperref[XVII.prop1.2.2]{1.2.2} shows that $\Aa^{1}$ is
a triangulated subcategory of $\Aa$, evidently cofibered. Set
$S^{1}=S\cap\text{Fl}(\Aa^{1})$.

\begin{enonce*}{Lemma 2.4.1.1}\etiquette[2.4.1.1]{XVII.lem2.4.1.1}
Suppose that $S$ is stable under direct images. Then $\Aa/\Bb$ is
coderivable and $\Aa(S^{-1})$ can be computed by a calculus of right
fractions.
\end{enonce*}

Let us verify condition c) of \cite[I 2.2]{sga4.exp.XVII.bib.GZ}, possibly the least
trivial of the conditions. We are given $f$ and $s$ $(s\in S)$ and
must find
$$\etiquette[2.4.1.2]{XVII.eq2.4.1.2}
\vcenter{
\xymatrix{
A'        \ar@{-->}[r]      & B'          \\
A  \ar[u]^{s}   \sisi{\ar[r]}{\ar[r]_f}          & B\rlap{,}  \ar@{-->}[u]^{t}
}}\leqno{(2.4.1.2)}
$$
a\pageoriginale commutative diagram \eqhyperref[XVII.eq2.4.1.2]{2.4.1.2}
with $t\in S$. Let $F$ be the image of $f$ in $\Bb$. By hypothesis,
$F_{\ast}(s)$ still belongs to $S$, so by \textsc{Verdier} there is a
commutative diagram
\[
\xymatrix{
F_{\ast}(A')  \ar[r] & B'\\
F_{\ast}(A)    \ar[r]_{f}\sisi{\ar[u]_{s}}{\ar[u]^{F_\ast(s)}}   & B\ar[u]_{t}
}
\]
and hence also a diagram \eqhyperref[XVII.eq2.4.1.2]{2.4.1.2}.

Since $\Aa(S^{-1})$ is computed by a calculus of right fractions, it
is immediate that its formation commutes with base change
$c:\Bb'\rightarrow\Bb$.

This lemma applies in particular to $\Aa^{1}$ and $S^{1}$; the
following lemma therefore implies \hyperref[XVII.prop2.4.1]{2.4.1}:

\begin{enonce*}{Lemma 2.4.1.3}\etiquette[2.4.1.3]{XVII.lem2.4.1.3}
Over $\Bb$, and after every base change $c:\Bb'\rightarrow\Bb$, the
canonical functor
$\Phi:\Aa\sisi{_1}{^1}(\sisi{{S^1}}{(S^1)}^{-1})\rightarrow\Aa(S^{-1})$
is an equivalence of categories.
\end{enonce*}

Let us construct a quasi-inverse to $\Phi$. For each $A\in\Ob\Aa$,
choose $s_A\in S$ with target $A$ and source $A^1\in\Ob\Aa^{1}$,
with $s_A=\Id_A$ when $A\in\Ob\Aa^{1}$. Let $f$ be an arrow of $\Aa$
with target $A$, source $B$, and image $F$ in $\Bb$. Since
$s_B\in S$ and $F_{\ast}A^{1}\in\Ob\Aa^{1}$, one can find
$B\sisi{^{1}}{'}$, \sisi{$v$ and $s\in S^{1}$}{$s\in S^1$ and $u$}
making the following diagram commutative:
$$\etiquette[2.4.1.4]{XVII.eq2.4.1.4}
\vcenter{
\xymatrix{
A      \ar[r]^{f} &  B      \ar[r]_{s_B} & B\sisi{_1}{^1}       \\
A\sisi{_1}{^1} \ar[u]_{s_A}\ar[r]^{a} & F_{\ast}A\sisi{_1}{^1}  \ar[u]& B'\rlap{.} \ar[l]_{s} \ar[u]_u
}}\leqno{(2.4.1.4)}
$$

The\pageoriginale arrow $f^{1}=us^{-1}a$ of
$\Aa^{1}(\sisi{{S^1}}{(S^1)}^{-1})$ then depends only on $f$.

Let $f$ and $g$ be two composable arrows, with images $F$ and $G$ in
$\Bb$. To prove that $(g\circ f)^{1}=g^{1}\circ f^{1}$, consider the
following diagram, which contains diagrams
\eqhyperref[XVII.eq2.4.1.4]{2.4.1.4} for $f$ and $g$:
\[
\xymatrix@C=2em{
                      & A \ar[rr]^{f}&      & B\ar[rrr]^{g}     & &
                      & C  &       \\
 A^{1} \ar[ru] \ar[rr]^{a}&          &F_{\ast}A^{1}\ar[ru] \ar[rrr]^(.6){a'}& & & (GF)_{\ast}A^1\ar[ru] &
              & &C^{1}\ar[llu]\\
& & & B^1 \ar[rrr]^(.4){b}|(.58)\hole\ar[uu]|\hole & & & G_{\ast}B^1\ar[uu] &     \\
&  & B'\ar[uu]^{s}\ar[rrr]^{b'}\ar[ru]_{u}       &   & &G_{\ast}B'\ar[ru]\ar[uu]^(.3){G_{\ast}s}    &  &C''\ar[ll]^{r}\ar[r]_{w} & C'\ar[llu]_t\ar[uu]_{v}\rlap{.}
}
\]

Since $t\in S$, there exist $C''$, $r\in S^{1}$, and $w$ making the
diagram commutative. By definition,\break
$g^{1}\circ f^{1}=vt^{-1}b\,u\,s^{-1}a$ and
$(g\circ f)^{1}=vw(G_{\ast}(s)r)^{-1}a'a$, and the conclusion follows
from the commutativity of the diagram.

We have thus defined a functor from $\Aa$ to
$\Aa^{1}(\sisi{{S^1}}{(S^1)}^{-1})$; by the universal property of
$\Aa(S^{-1})$, it factors through
$\Psi:\Aa(S^{-1})\rightarrow\Aa^{1}(\sisi{{S^1}}{(S^1)}^{-1})$.
The composite $\Phi\Psi$ is isomorphic to the identity because its
composite with the functor $J:\Aa\rightarrow\Aa(S^{-1})$ is
isomorphic to $J$. One checks similarly that $\Psi\Phi$ is isomorphic
to the identity.

\begin{enonce*}{Corollary 2.4.2}\etiquette[2.4.2]{XVII.cor2.4.2}
With the notation of \hyperref[XVII.prop2.4.1]{2.4.1}, $\Aa$ is
coderivable relative to $S$ provided that, for every finite set $B$
of arrows of $\Bb$, condition $(\ast)$ holds when only direct images
under the arrows of $B$ are considered.
\end{enonce*}

For\pageoriginale every category $\Bb'$ with a finite set of arrows, Proposition
\hyperref[XVII.prop2.4.1]{2.4.1} is applicable after every base change
$\Bb'\rightarrow\Bb$. We shall take for $\Bb'$ the categories of the
following diagrams:
\[
  \begin{gathered}
    *\\
    (0)
  \end{gathered}\qquad
  \begin{gathered}
    *\to*\\
    (1)
  \end{gathered}\qquad
  \begin{gathered}
    *\to*\to*\\
    (2)
  \end{gathered}\qquad
  \begin{gathered}
    *\to*\to*\to*\\
    (3)
  \end{gathered}
\]
Consider the categories $\Aa_x(S^{-1}_x)$ $(x\in\Ob\Bb)$. Given
$f\in\Fl(\Bb)$ and $A_1,A_2$ in these categories over the source and
target of $f$, define $\Hom_f(A_1,A_2)$ to be the analogous set defined
after the base change of type (1) determined by $f$. Composition of
morphisms is defined using base changes of type (3). In this way one
constructs a category $\Aa(S^{-1})^{\ast}$ over $\Bb$, whose formation
is compatible with every base change; using base changes of types
(0), (1), and (2), one checks that it solves the universal problem
defining $\Aa(S^{-1})$.

The following result is dual to \hyperref[XVII.cor2.4.2]{2.4.2}.

\begin{enonce*}{Proposition 2.4.3}\etiquette[2.4.3]{XVII.prop2.4.3}
Let $\Aa$ be a triangulated category fibered over $\Bb$, and let $S$
be as in \hyperref[XVII.prop2.4.1]{2.4.1}. For $\Aa$ to be derivable
relative to $S$, it suffices that, for every finite set $B$ of arrows
of $\Bb$,
\begin{itemize}
\item[$(\ast')$] $\forall x\in\Ob(\Bb)$, $\forall A\in\Ob(\Aa_x)$,
$\exists s:A\rightarrow A'$ in $S$,
$\forall g,f:Z\xrightarrow{g}Y\xrightarrow{f}X$ in $B$, $f^{\ast}A'$
is right adapted to the functor $g^{\ast}$.
\end{itemize}
\end{enonce*}

\section{Gluing fibered or cofibered categories}\etiquette[3]{XVII.sec3}\pageoriginale

\subsection{Introduction}\etiquette[3.1]{XVII.subsec3.1}

Let $f:Y\rightarrow X$ be a separated morphism of finite type from a
scheme $Y$ to a scheme $X$, or a continuous map from a locally compact
space to a locally compact space. One can then define a ``direct image
with proper supports'' functor, denoted $f_{!}$, from the category of
abelian sheaves on $Y$ to the category of abelian sheaves on $X$ (\cf
\sisi{\hyperref[XVII.subsec5.3]{5.3}}{\hyperref[XVII.subsec6.1]{6.1}}
in the case of schemes, and \textsc{Verdier} \cite{sga4.exp.XVII.bib.V3} for separated
locally compact spaces). When $Y$ and $X$ are schemes of finite type
over $\CC$, these two definitions are compatible in the evident sense.
Unfortunately, for schemes, unlike for locally compact spaces, the
derived functors of $f_{!}$ are pathological; one can nevertheless
define a reasonable ``direct image with proper supports'' functor
directly from $D^{+}(Y)$ to $D^{+}(X)$. To define it, factor $f$ as an
open immersion followed by a proper morphism: $f=\bar f j$. The desired
functor $Rf_{!}$ is then defined as the composite of ``extension by
zero'' $j_{!}$ and the functor $R\bar f_{\ast}$, derived from direct
image under $\bar f$.

These ``direct image with proper supports'' functors $Rf_{!}$ are
usable only after one verifies for them a ``variance formalism'' that
includes, for a composite $f=gh$, a formula
\[
Rf_{!}=Rg_{!}Rh_{!}
\]
which is best expressed in terms of cofibered categories. Such a
formalism is available separately for ``extensions by zero'' and for
\pageoriginale direct images under proper morphisms, and the problem is to glue these
formalisms. To avoid later being bogged down by pointless computations,
we propose in this \S to list the elementary verifications required to
carry out such a gluing. We do so in a slightly more abstract setting
than is strictly necessary, in the hope that the result may be useful
again.

The construction of $Rf_{!}$ sketched above requires that the morphism
$f$ be factorable as an open immersion followed by a proper morphism.
Following \textsc{Hartshorne} \cite[III 8 p.~189]{sga4.exp.XVII.bib.RD}, in nᵒ
\hyperref[XVII.sec2]{2} we introduce the appropriate notion of a
compactifiable morphism.

\subsection{Compactifiable morphisms}\etiquette[3.2]{XVII.subsec3.2}

% END INCLUDED SGA4_English_current_buildable_source_20260906/source/17/drafts/SGA4_XVII_LINES_1801_2200_English.texfrag

% BEGIN INCLUDED SGA4_English_current_buildable_source_20260906/source/17/drafts/SGA4_XVII_LINES_2201_2600_English.texfrag
\begin{enonce*}{Definition 3.2.1}\etiquette[3.2.1]{XVII.def3.2.1}
A morphism $f$ from an $S$-scheme $X$ to a quasi-compact,
quasi-separated $S$-scheme $Y$ is said to be \emph{$S$-compactifiable} if there exists an
$S$-scheme $P$ proper over $S$ and a factorization of $f$ as a
separated quasi-finite morphism $j$ from $X$ to $\fprod{P}{Y}{S}$ followed by
the projection from $\fprod{P}{Y}{S}$ to $Y$
\[
\xymatrix{
                                & &  &\fprod{P}{Y}{S} \ar[rd]\ar[ld]& \\
X \ar[rr]^(.7){f} \ar[rrru]^{j}\ar[rrrd]& &Y \ar[rd]&                & P\ar[ld]\\
                                & &  &       S        &.
}
\]
\end{enonce*}

When $S=Y$, we shall simply speak of a compactifiable morphism.

\begin{enonce*}{Proposition 3.2.2}\etiquette[3.2.2]{XVII.prop3.2.2}
\begin{enumerate}
\renewcommand{\theenumi}{\roman{enumi}}
\renewcommand{\labelenumi}{(\theenumi)}
\item An\pageoriginale $S$-compactifiable morphism is separated and of finite
type.

\item If $Y$ is quasi-compact and quasi-separated, a separated quasi-finite
morphism with target $Y$ is \sisi{}{$S$-}compactifiable.

\item The composite of two $S$-compactifiable morphisms is again
$S$-compactifiable.

\item Let there be a diagram of quasi-compact, quasi-separated $S$-schemes
\[
\xymatrix{
\fprod{X}{Y'}{Y} \ar[d]_{f'} \ar[r] & X \ar[d]^{f}      \\
Y' \ar[r]& Y\rlap{.}
}
\]
If $f$ is $S$-compactifiable, then $f'$ is $S$-compactifiable.

\item Let $g$ be an $S$-compactifiable morphism;\nde{it suffices to assume that $g$ is a quasi-compact, quasi-separated morphism of $S$-schemes.} for a composite morphism
$g\circ f$ to be $S$-compactifiable, it is necessary that $\sisi{h}{f}$ be so.
\end{enumerate}
\end{enonce*}

Assertions (i), (ii), and (iv) are immediate. To prove (iii) and
(v), consider a composite $f=gh$, with $g$ \sisi{}{$S$-}compactifiable.

Suppose that $gh$ is \sisi{}{$S$-}compactifiable, whence a diagram
\[
\xymatrix{
                     & \fprod{P}{Y}{S} \ar[d] \ar[r]^{v} & \fprod{P}{Z}{S}\ar[d] \\  X  \ar[ru]^{u} \ar[r]^{h}   & Y                \ar[r]^{g} &    Z\rlap{.}
}
\]
The composite $vu$ is quasi-finite; hence $u$ is quasi-finite, and (v) is
proved.

If $h$ is compactifiable, we have a diagram
\[
\xymatrix{
  &&                & &\sisi{Q\times P \times Z}{\fprod{\fprod{Q}{P}{S}}{Z}{S}} \ar[d]&\\
  &&\fprod{Q}{Y}{S} \ar[d]\ar[rru]\ar[rdd]& & \fprod{P}{Z}{S} \ar[d]\ar[rdd] & \\
X \ar[rrrdd]\ar[rru] \ar[rr]^{h}&& Y \sisi{\ar[rru]}{\ar[rru]|(.2)\hole} \ar[rdd]\sisi{\ar[rr]^{g}}{\ar[rr]^g|(.25)\hole}            & & Z\ar[ldd]                & \\
  &&                &Q\ar[d]&                  &P\ar[lld] \\
  &&                &S&                   &,
}
\]
and\pageoriginale $f=fh$ factors through a separated quasi-finite morphism
to $\fprod{(\fprod{Q}{P}{S})}{Z}{S}$.

Nagata asserts in \sisi{[5]}{\cite{sga4.exp.XVII.bib.Na}} that every morphism of finite type between separated
integral noetherian schemes is compactifiable.\nde{more generally, every separated morphism of finite type whose target is quasi-compact and quasi-separated is compactifiable. See \nderef{XVII.ndeConrad}.} The editor admits to
not having understood the proof.

Now fix a scheme $S$ and denote by $(S)$ the
subcategory of the category of $S$-schemes whose objects are
the quasi-compact, quasi-separated $S$-schemes and whose arrows are
the $S$-compactifiable $S$-morphisms.\nde{Thanks to the compactification theorem, one may also work in the subcategory of the category of (absolute) schemes whose objects are the quasi-compact, quasi-separated schemes and whose arrows are the separated morphisms of finite type (\ie compactifiable morphisms). Note that the stability of compactifiable morphisms under composition is equivalent to the general form of the compactification theorem.}

\begin{enonce*}{Proposition 3.2.3}\etiquette[3.2.3]{XVII.prop3.2.3}
In the category $(S)$:
\begin{enumerate}
\renewcommand{\theenumi}{\roman{enumi}}
\renewcommand{\labelenumi}{(\theenumi)}
\item The open immersions (\resp the proper morphisms)
define a subcategory $(S,i)$ (\resp $(S,p)$) of $(S)$ having
the same objects as $(S)$.

\item Fiber products exist in $(S,p)$ and are fiber
products in $(S)$.

\item Every morphism $f$ is the composite $f=pj$ of a proper morphism
$p$ and an open immersion.
\end{enumerate}
\end{enonce*}

By\pageoriginale \hyperref[XVII.prop3.2.2]{3.2.2} (iii), (iv), and (v), fiber products
exist in $(S)$ and are fiber products in the category of
all schemes. Since the fiber product $X=\fprod{X_1}{X_2}{Y}$ of
two schemes proper over $Y$ is again proper over $Y$, and a
scheme is proper over $X$ if and only if it is proper over $X_1$
and $X_2$, condition (ii) is satisfied.

If $f$ is a morphism of $(S)$, by hypothesis it admits, in $(S)$,
a factorization $f=pj'$, where $p$ is proper and $j\sisi{}{'}$ is separated and quasi-finite.
By \EGA IV 18.12.13, $j'$ admits a factorization
$j'=qj$, where $q$ is finite and $j$ is an open immersion; moreover, by
\hyperref[XVII.prop3.2.2]{3.2.2} (ii), $q$ and $j$ belong to $(S)$. We then have $f=(pq)j$,
which proves (iii). Assertion (i) follows from \hyperref[XVII.prop3.2.2]{3.2.2} (iii).

\setcounter{subsubsection}{3}
\subsubsection{}\etiquette[3.2.4]{XVII.subsubsec3.2.4}
In the rest of this \S, suppose given a category $(S)$, equipped
with two subcategories $(S,i)$ and $(S,p)$, such that
\begin{enumerate}
\renewcommand{\theenumi}{\roman{enumi}}
\renewcommand{\labelenumi}{(\theenumi)}
\item $\Ob(S)=\Ob(S,i)=\Ob(S,p)$;

\item fiber products exist in $(S,p)$ and are fiber
products in $(S)$;

\item $\forall f\in \Fl(S)$, $\exists p\in \Fl(S,p)$, $\exists
j\in \Fl(S,i)$, $f=pj$.
\end{enumerate}

The arrows of $(S,i)$ (\resp of $(S,p)$) will be called open
immersions (\resp proper morphisms).

\begin{enonce*}{Definition 3.2.5}\etiquette[3.2.5]{XVII.def3.2.5}
\begin{enumerate}
\renewcommand{\theenumi}{\roman{enumi}}
\renewcommand{\labelenumi}{(\theenumi)}
\item A \emph{compactification} of a morphism $f$ of $(S)$ (\resp
of a pair $(f,g)$, \resp of a triple $(f,g,h)$ of composable
morphisms) is a commutative diagram \eqhyperref[XVII.eq3.2.5.1]{3.2.5.1} (\resp
\eqhyperref[XVII.eq3.2.5.2]{3.2.5.2}, \resp \eqhyperref[XVII.eq3.2.5.3]{3.2.5.3}) whose horizontal arrows are
open immersions and whose oblique arrows are proper.

\item \emph{A\pageoriginale morphism of compactifications} is a
  morphism of
diagrams \eqhyperref[XVII.eq3.2.5.1]{3.2.5.1} (\resp\ldots) that is the identity on the 1\sisi{st}{\textsuperscript{st}}
vertical column and is proper at every vertex
$$
\etiquette[3.2.5.1]{XVII.eq3.2.5.1}
\etiquette[3.2.5.2]{XVII.eq3.2.5.2}
\etiquette[3.2.5.3]{XVII.eq3.2.5.3}
\begin{gathered}[b]
\xymatrix{
Y        \ar[d]_{f} \ar@{^{(}->}[r]^{i}                & \bar{Y} \ar[ld]^{\bar{f}}\\
X      & &}\\
\hyperref[XVII.eq3.2.5.1]{(3.2.5.1)}
\end{gathered}\qquad
\begin{gathered}[b]
\xymatrix{
Z        \ar@{^{(}->}[r]^{j} \sisi{\ar[d]}{\ar[d]_g}& \bar{Z}\ar@{^{(}->}[r]^{j'}\sisi{\ar@{-}[ld]^{\bar{g}}}{\ar[ld]^{\bar{g}}}
&\bar{\bar{Z}} \ar[ld]^{\bar{\bar{g}}}\\
Y \ar[d]_{f}\ar@{^{(}->}[r]^{i} & {\bar{Y}} \ar[ld]^{\bar{f}}&\\
X      & &}\\
\hyperref[XVII.eq3.2.5.2]{(3.2.5.2)}
\end{gathered}\qquad
\begin{gathered}[b]
\xymatrix@=3em{
\bigdot \ar@{^{(}->}[r]_(.4){\sisi{h}{k}} \ar[d]_{h}& \bigdot \ar@{^{(}->}[r]_(.4){\sisi{h}{k}'}\ar[ld]^(.4){\bar{h}} & \bigdot
\ar@{^{(}->}[r]_(.4){\sisi{h}{k}''} \ar[ld]^(.4){\bar{\bar{h}}}&\bigdot \ar[ld]^(.4){\bar{\bar{\bar{h}}}}\\
\bigdot \ar@{^{(}->}[r]_(.4){j}\ar[d]_{g} & \bigdot \ar@{^{(}->}[r]_(.4){j'}
\ar[ld]^(.4){\bar{g}}& \bigdot \ar[ld]^(.4){\bar{\bar{g}}}\\
\bigdot
\ar@{^{(}->}[r]_(.4){i} \ar[d]_{f}&\bigdot \ar[ld]^(.4){\bar{f}}\\
\bigdot}\\
\hyperref[XVII.eq3.2.5.3]{(3.2.5.3)}
\end{gathered}
$$
\end{enumerate}
\end{enonce*}

\begin{enonce*}{Proposition 3.2.6}\etiquette[3.2.6]{XVII.prop3.2.6}
\begin{enumerate}
\renewcommand{\theenumi}{\roman{enumi}}
\renewcommand{\labelenumi}{(\theenumi)}
\item Every morphism (\resp every pair $(f,g)$, \resp every triple
$(f,g,h)$ of composable morphisms) admits a compactification.

\item The category of compactifications of $f$ is left
filtering.\nde{``left filtering'' means ``cofiltered'' (\hyperrefext[I][I.def2.7]{2.7}). For axioms $L_1$ and $L_2$ in the proof of \hyperref[XVII.prop3.2.6]{3.2.6}, \cf \hyperrefext[V][V.par8.1.1.1]{8.1.1.1}.}

\item If diagrams $\sisi{(3.2.4.2)}{\eqhyperref[XVII.eq3.2.5.2]{3.2.5.2}}_i$ $(i=1,2)$ are
compactifications of $(f,g)$, there exists a compactification
$\sisi{(3.2.4.2)}{\eqhyperref[XVII.eq3.2.5.2]{3.2.5.2}}_3$ of $(f,g)$ and morphisms of compactifications from
$\sisi{(3.2.4.2)}{\eqhyperref[XVII.eq3.2.5.2]{3.2.5.2}}_3$ to the compactifications $\sisi{(3.2.4.2)}{\eqhyperref[XVII.eq3.2.5.2]{3.2.5.2}}_i$.
\end{enumerate}
\end{enonce*}

\begin{enumerate}
\renewcommand{\theenumi}{\roman{enumi}}
\renewcommand{\labelenumi}{(\theenumi)}
\item The case of a morphism is precisely \hyperref[XVII.subsubsec3.2.4]{3.2.4}
(iii). Let us treat the case of a triple; with the notation of \eqhyperref[XVII.eq3.2.5.3]{3.2.5.3},
it suffices to construct successively compactifications of $f$,
$g$, and $h$, of $i\bar{g}$ and $j\bar{h}$, and finally of
$j'\bar{\bar{h}}$.

\item Let\pageoriginale $\eqhyperref[XVII.eq3.2.5.1]{3.2.5.1}_1$ and $\eqhyperref[XVII.eq3.2.5.1]{3.2.5.1}_2$ be two
compactifications of $f$. The projection $\sisi{\rho}{p}$ from
$\fprod{\bar{Y}_1}{\bar{Y}_2}{S}$ to $X$ is proper, and $Y$ maps
to it by $(i_1,i_2)$, which (\hyperref[XVII.subsubsec3.2.4]{3.2.4} (iii)) can be compactified as
$(i_1,i_2)=qj$. The compactification $f=(p\sisi{\cdot}{\circ} q)j$ dominates
$\sisi{(3.2.4.1)}{\eqhyperref[XVII.eq3.2.5.1]{3.2.5.1}}_1$ and $\sisi{(3.2.4.1)}{\eqhyperref[XVII.eq3.2.5.1]{3.2.5.1}}_2$.

If $r$ and $s$ are two morphisms from $\sisi{(3.2.4.1)}{\eqhyperref[XVII.eq3.2.5.1]{3.2.5.1}}_1$ to
$\sisi{(3.2.4.1)}{\eqhyperref[XVII.eq3.2.5.1]{3.2.5.1}}_2$, the kernel $K$ of the pair of arrows
$\bar{Y}_{1} \overset{r}{\underset{s}\rightrightarrows} \bar{Y}_{2}$ maps by
a proper morphism to $\bar{Y}_1$, and $i_1$ factors through $K$:
write $i_1=ki$. If $i=pi_3$ is a compactification of $i$,
the existence of the compactification $f=(\bar{f}_1kp)i_3$ shows that
axiom $L_2$ is satisfied.

\item Apply axiom $L_1$ for filtering categories to the
compactifications of $f$ and $g$ deduced from the $\sisi{(3.2.4.2)}{\eqhyperref[XVII.eq3.2.5.2]{3.2.5.2}}_i$
$(i=1,2)$. We thereby obtain $\bar{Y}_3$ and $\bar{Z}_3$ giving rise to
commutative diagrams $\eqhyperref[XVII.eq3.2.5.1]{3.2.5.1}_i$ $(i=1,2)$
\[
\xymatrix{
&         &\bar{Z}_i \ar[d] \ar@{^{(}->}[r]& \bar{\bar{Z}}_i\ar[d]\\
\bar{Z}_3 \ar[rru] \ar[d]^{\bar{g}_3} &&Y   \ar@{^{(}->}[r]      & \bar{Y}_i\\
Y \ar@{=}[rru]\ar[r]_{i_3} & \bar{Y}_3 \ar[rru]&&\sisi{(3.2.5.1)_i}{},
}
\]
and it remains to find a compactification of $i_3\bar{g}_3$ by
$\bar{\bar{Z}}_3$, such that $\bar{\bar{Z}}_3$ can map by a
proper morphism to $\bar{\bar{Z}}_i$, so as to give rise to
a commutative cube $(i=1,2)$. If the target of a morphism between
compactifications of $i_3\bar{g}_3$ has this property, its source has it
as well, and by (ii), it will suffice to consider separately the cases
$i=1$ and $i=2$. Let\pageoriginale $T$ be the fiber product of $\bar{Y}_3$
and $\bar{\bar{Z}}_i$ over $\bar{Y}_i$: then $T$ occurs in a
commutative cube analogous to the desired one, except that the arrow $k$ from
$\bar{Z}_3$ to $T$ might not be an open
immersion. Writing $k=pj$ ($p$ proper, $j$ an open immersion) and repla\c{}cing $T$
by the source $\bar{\bar{Z}}_3$ of $p$, we obtain the desired
cube.
\end{enumerate}

\begin{enonce*}{Remark 3.2.7}\etiquette[3.2.7]{XVII.rem3.2.7}
Proposition \hyperref[XVII.prop3.2.3]{3.2.3} would still have been valid if, in Definition
\hyperref[XVII.def3.2.1]{3.2.1} of compactifiable morphisms, we had omitted the requirement that $Y$
be quasi-compact and quasi-separated and required instead that $j$
be a quasi-compact open immersion and not merely a
quasi-finite morphism.
\end{enonce*}

\subsection{Gluing}\etiquette[3.3]{XVII.subsec3.3}

Recall that the hypotheses of \hyperref[XVII.subsubsec3.2.4]{3.2.4} are satisfied.

\subsubsection{}\etiquette[3.3.1]{XVII.subsubsec3.3.1}
Suppose given:
\begin{itemize}
\item[$\alpha$)] for each $X\in \Ob (S)$, a category $F(X)$,

\item[$\beta$)] an $(S,i)$-cofibered category $F_i$ and an $(S,p)$-
cofibered category $F_{p}$,

\item[$\gamma$)] for each $X\in \Ob(S)$, isomorphisms between
$F(X)$, $(F_{i})_X$, and $(F_{p})_X$,

\item[$\delta$)] a normalized cleavage of $F_{i}$ and a normalized cleavage
of $F_{p}$.
\end{itemize}

By (\SGA 1 VI), it ``amounts to the same thing'' to give instead $\alpha$) and:
\begin{itemize}
\item[(i)] For every proper morphism $p:X\rightarrow Y$, a functor
$p_{\ast}:F(X)\rightarrow F(Y)$,

\item[(ii)] for every composable pair $(p,q)$ of proper morphisms,
an\pageoriginale isomorphism of functors
$c_{p,q}:p_{\ast}q_{\ast}\leftrightarrow (pq)_{\ast}$,

\item[(i\sisi{}{'}) and (ii\sisi{)'}{')},] analogous data for $(S,i)$,
\end{itemize}
these data satisfying the following conditions:
\begin{itemize}
\item[(a)] if $p$ is an identity, then $p_{\ast}$, $c_{p,q}$, and
$c_{q,p}$ are identities,

\item[(b)] for every composable triple of proper morphisms, we have
\[
c_{p,qr}\circ (p_{\ast}\ast c_{q,r})=c_{pq,r}\circ (c_{p,q}\ast
r_{\ast}){\quoi},
\]

\item[(a\sisi{)'}{')} and (b\sisi{)'}{')},] analogous conditions for $(S,p)$.
\end{itemize}

Suppose \emph{given} in addition,
\begin{itemize}
\item[(iii)] for every commutative diagram $D$
\[
\vcenter{
\xymatrix{
X \ar[r]^{j'} \ar[d]^{p'}& Y\ar[d]^{p} \\
Z \ar[r]_{j}             &           T
}}\quad j,j' \in \Fl(S,i)\text{; }
  p,p' \in \Fl(S,p),
\]
an isomorphism of functors $d(D):p_{\ast}j'_{\ast}\leftrightarrow
j_{\ast}p'_{\ast}$.
\end{itemize}

\begin{enonce*}{Proposition 3.3.2}\etiquette[3.3.2]{XVII.prop3.3.2}
If data {\rm (i) (ii) ${\rm (i')}$ ${\rm (ii')}$ (iii)}
satisfy conditions {\rm (a) (b) ${\rm (a')}$ ${\rm (b')}$} as well
as {\rm(c)} and ${\rm (c')}$ stated below, then there exists one
and ``essentially'' only one category $F$ cofibered over $S$,
equipped with isomorphisms $\fprod{F}{(S,i)}{S}=F_i$ and
$\fprod{F}{(S,p)}{S}=F_p$, compatible with datum \hyperref[XVII.subsubsec3.3.1]{3.3.1} $\gamma$)
and such that the isomorphisms $d(D)$ of {\rm(iii)} are the
composite isomorphisms
$j_{\ast}p'_{\ast}=(jp')_{\ast}=(pj')_{\ast}=p_{\ast}j'_{\ast}$.
\end{enonce*}

\begin{itemize}
\item[(c)] for every commutative diagram of the following type
\[
\xymatrix{
X_1 \ar@{^{(}->}[r]^{i} \ar[d]^{p_1}& X_2 \ar@{^{(}->}[r]^{i'}
\ar[d]^{p_2}& X_3 \ar[d]^{p_3}\\
Y_1 \ar@{^{(}->}[r]^{j}& Y_2 \ar@{^{(}->}[r]^{j'}& Y_3,\\
}
\]
whose\pageoriginale horizontal arrows are open
immersions and whose vertical arrows are proper, the following
triangle of isomorphisms $d$ from (iii) is commutative:
\[
\xymatrix@R=1cm{
p_{3{\ast}} i'_{\ast} i_{\ast} \ar[rr] \ar@{-}[rd]&  & j'_{\ast}j_{\ast}p_{1{\ast}}\ar@{-}[ld]\\
                               &j'_{\ast}p_{2{\ast}}
                                i_{\ast}&.
}
\]

\item[(c')] For every commutative diagram of the following type
\[
\xymatrix{
X_1  \ar[r]^{p} \ar[d]^{i_1}& X_2  \ar[r]^{p'} \ar[d]^{i_2}& X_3 \ar[d]^{i_3}          \\
Y_1         \ar[r]^{q} & Y_2  \ar[r]^{q'} & Y_3\rlap{,}
}
\]
whose horizontal arrows are proper and whose vertical
arrows are open immersions, the following triangle
of isomorphisms $d$ from (iii) is commutative:
\[
\xymatrix{
i_{3{\ast}} p'_{\ast} p_{\ast} \ar[rr] \ar@{-}[rd]&    &q'_{\ast} q_{\ast}
i_{1{\ast}} \\
&q'_{\ast} i_{2{\ast}} p_{\ast}  \ar@{-}[ru]&.
}
\]
\end{itemize}

\setcounter{subsubsection}{2}
\paragraph{Definition of $f_{\ast}$}\etiquette[3.3.2.1]{XVII.par3.3.2.1}\pageoriginale
Let $f$ be a morphism from $Y$ to $X$. Each compactification of $f$,
with diagram \sisi{(\hyperref[XVII.def3.2.5]{3.2.5} (i))}{\eqhyperref[XVII.eq3.2.5.1]{3.2.5.1}}, defines a composite functor $\bar{f}_{\ast}i_{\ast}$ from $F(Y)$ to $F(X)$. A morphism of compactifications
defines a commutative diagram
$$\etiquette[3.3.2.2]{XVII.eq3.3.2.2}
\vcenter{
\xymatrix{
Y  \ar@{^{(}->}[r]^{i_1} \ar@{=}[d]& \bar{Y}_1  \sisi{\ar[d]}{\ar[d]^p}\\
Y  \ar@{^{(}->}[r]^{i_2} \ar[d]        & \bar{Y}_2 \ar[ld]^{\bar{f}_2}\\
X                         &.
}}\leqno{(3.3.2.2)}
$$
Datum (iii) therefore provides us with an isomorphism
\[
\bar{f}_{2{\ast}} i_{2{\ast}}=\bar{f}_{2\ast} \circ
(i_{2\ast}\circ 1_{Y{\ast}})=\bar{f}_{2{\ast}}\circ
(p_{\ast}\circ i_{1{\ast}})=\bar{f}_{1{\ast}}\circ
i_{1{\ast}}{\quoi}.
\]
Axiom ${\rm(c')}$ also ensures that the isomorphism associated with a
composite of morphisms of compactifications is the composite of the
isomorphisms associated with each of them. Since the category of morphisms of
compactifications is cofiltered (\hyperref[XVII.prop3.2.6]{3.2.6} (ii)), we thereby obtain a
transitive system of isomorphisms between the functors associated with the
various compactifications of $f$. Choosing one of them, we
define the functor $f_{\ast}$.

\setcounter{paragraph}{2}
\paragraph{Transitivity homomorphisms}\etiquette[3.3.2.3]{XVII.par3.3.2.3}

Every compactification of a pair $(f,g)$ of composable morphisms
(diagram \eqhyperref[XVII.eq3.2.5.2]{3.2.5.2}) defines a compactification of $f$, one of $g$,
one of $f\circ g=(\bar{f}\bar{\bar{g}})(j'j)$, and an isomorphism
$c_{f,g}:(fg)_{\ast}\isomto f_{\ast}g_{\ast}$, the composite of the
isomorphisms
\[
(fg)_{\ast}=(\bar{g}\,\bar{\bar{g}})_{\ast}(j'j)_{\ast}=\bar{f}_{\ast}(\bar{\bar{g}}_{\ast}j'_{\ast})j_{\ast}=\bar{f}_{\ast}(i_{\ast}\bar{g}_{\ast})j_{\ast}=(\bar{f}_{\ast}i_{\ast})(\bar{g}_{\ast}j_{\ast})=f_{\ast}g_{\ast}
\]
(where\pageoriginale the middle isomorphism comes from (iii)).

% END INCLUDED SGA4_English_current_buildable_source_20260906/source/17/drafts/SGA4_XVII_LINES_2201_2600_English.texfrag

% BEGIN INCLUDED SGA4_English_current_buildable_source_20260906/source/17/drafts/SGA4_XVII_LINES_2601_3000_English.texfrag

Let us prove that this isomorphism does not depend on the compactification
chosen for $(f,g)$. By \hyperref[XVII.prop3.2.6]{3.2.6} (iii), it will suffice to compare
the isomorphisms $c_{f,g}$ obtained from two compactifications
$\sisi{(3.2.4.2)}{\eqhyperref[XVII.eq3.2.5.2]{3.2.5.2}}_i$ $(i=1,2)$ such that there exists a morphism of
compactifications from the first to the second. The arrows in these
morphisms will be denoted by $p$.

Consider the following diagram of isomorphisms of functors. For
readability, the ${\ast}$ have been omitted.
\[
\xymatrix{
\bar{f}_2 \bar{\bar{g}}_2j'_2j_2\sisi{}{\ar@{}[rrd]|{(1)}} \ar@{-}[rr] \ar@{-}[ddd] \ar@{-}[rd]& &
\bar{f}_2\bar{\bar{g}}_2pj'_1j_1
\ar@{-}[rr] \ar@{-}[ld] \ar@{-}[rd]& &\bar{f}_1\bar{\bar{g}}_1j'_1j_1 \ar@{-}[ddd]\sisi{}{\ar@{}[lld]|{(2)}}\\
 \ar@{}[ddr]|{(3)} & \bar{f}_2\bar{\bar{g}}_2j'_2pj_1\ar@{-}[ru] \ar@{-}[dd]& &
\bar{f}_2p\bar{\bar{g}}_1j'_1j_1 \ar@{-}[dd]\ar@{-}[ru]& \ar@{}[ldd]|{(5)}\\
 &  &  (4)  & &\\
\bar{f}_2i_2\bar{g}_2j_2 \ar@{-}[r] & \bar{f}_2i_2\bar{g}_2pj_1 & &
\bar{f}_2pi_1\bar{g}_1j_1 \ar@{-}[r]&  \bar{f}_1i_1\bar{g}_1j_1\\
& &  \bar{f}_2i_21_Y\bar{g}_1j_1 \ar@{-}[ru]\ar@{-}[lu]& & .\\
}
\]

\begin{enumerate}
\renewcommand{\theenumi}{\alph{enumi}}
\renewcommand{\labelenumi}{\theenumi)}
\item The composite isomorphism\pageoriginale along the 1\sisi{st}{\textsuperscript{st}} row is the isomorphism
\hyperref[XVII.par3.3.2.1]{3.3.2.1} between two definitions of $(fg)_{\ast}$.

\item The composite isomorphism along the last row is the isomorphism
\hyperref[XVII.par3.3.2.1]{3.3.2.1} between two definitions of $f_{\ast}g_{\ast}$.

\item The two outer vertical isomorphisms are the isomorphisms
$c_{f,g}$ between $(fg)_{\ast}$ and $f_{\ast}g_{\ast}$ deduced from one
or the other compactification.
\end{enumerate}
The problem is therefore to prove that the
outer boundary of the diagram is commutative. A commutative diagram of
functors remains commutative when each of them is composed with the same
functor. Bearing this in mind, the commutativity of (\sisi{i}{1}) follows from (c)
applied to the diagram
\[
\xymatrix{
Z  \ar@{^{(}->}[r]^{j_1} \ar@{=}[d]& \bar{Z}_1  \ar@{^{(}->}[r]^{j'_1} \ar[d]& \bar{\bar{Z}}_1 \ar[d]\\
Z  \ar@{^{(}->}[r]^{j_2} & \bar{Z}_2  \ar@{^{(}->}[r]^{j'_2} & \bar{\bar{Z}}_2\rlap{.}
}
\]
The commutativity of (3) and (5) is trivial; that of (2), which
essentially concerns proper morphisms, is also trivial. It remains to
consider the hexagon (4).

Consider the following commutative cube
\[
\xymatrix{
                              &**[r]\bar{Z}_2  \ar@{^{(}->}[rrr]^{ j'_2}
                              \ar[dd]^(.35){\bar{g}_2}|\hole &    &   &
                              **[r]\bar{\bar{Z}}_2 \ar[dd]^{\bar{\bar{g}}_2}\\
  \bar{Z}_1\sisi{}{\ar@{-->}[rd]}  \ar@{^{(}->}[rrr]^(.65){ j'_1} \ar[ru] \ar[dd]^{\bar{g}_1}   &           &               &
                              \bar{\bar{Z}}_1 \ar[ru]\ar@{-->}[rd] \ar[dd]^(.35){\bar{\bar{g}}_1}|\hole&\\
& **[r] Y\ar[rrr]_{i_2}  \ar@{^{(}->}[rrr]  & && **[r]\bar{Y}_2 \\
Y  \ar@{^{(}->}[rrr]_{i_1}   \ar@{=}[ru]  &           &               &
                              \bar{Y}_1\ar[ru]  &.\\
}
\]

There\pageoriginale are six ways to follow the edges of this cube from
$\bar{Z}_1$ to $\bar{Y}_2$, each ``adjacent'' to two others;
there is an isomorphism between the composite functors associated with
two adjacent paths (\ie paths separated by a single face), and the
commutativity of (4) will follow from the commutativity of the resulting
hexagon.

The moral of this kind of diagram is that an edge represents a
functor, a face an isomorphism between composite functors, and a
volume a compatibility condition (between the faces forming its
boundary). We decompose the cube into two triangular prisms (along
the ``plane'' $\bar{Z}_1 \bar{\bar{Z}}_1 Y \bar{Y}_2$) in order
to reduce to the hypotheses ${\rm(c')}$.

More precisely, we complete the hexagon to form the following diagram
of isomorphisms (for readability, the ${\ast}$ have
been omitted):
\[
\xymatrix{
                   & & \bar{\bar{g}}_2 j'_2p      & &   \\
i_2\bar{g}_2 p  \ar@{-}[rru] \ar@{-}[rd]\ar@{-}[dd] &   &   &&   \bar{\bar{g}}_2 p j'_1 \ar@{-}[ld]\ar[dd]\ar@{-}[llu]\\
 & i_2(g_2p) \ar@{-}[rr]&      & (\bar{\bar{g}}_2p)j'_1 & \\
i_21_{Y}\bar{g}_1 \ar@{-}[ru] \ar@{-}[rrd]&      &&   &
  p\bar{\bar{g}}_1 j'_1 \ar@{-}[lu] \ar@{-}[lld]\\
&  &  pi_1 \bar{g}_1  & &
}
\]
The\pageoriginale triangles are trivially commutative, and the
pentagons are commutative by ${\rm(c')}$.

\paragraph{Cocycle condition}\etiquette[3.3.2.4]{XVII.par3.3.2.4}

Let us verify that the isomorphisms $c_{f,g}$ constructed in \hyperref[XVII.par3.3.2.3]{3.3.2.3}
satisfy a cocycle condition for a composable triple $(f,g,h)$ of
morphisms. Let us introduce a compactification of
$(f,g,h)$, with diagram \sisi{(3.2.4.3)}{\eqhyperref[XVII.eq3.2.5.3]{3.2.5.3}}, and consider the following diagram
of isomorphisms of functors:
\[
\sisi{\xymatrix{
\bar{f} i \bar{g} j \bar{h} k \ar@{-}[rr] \ar@{-}[dd]&
&\bar{f}\bar{\bar{g}} j' j\bar{h} k \ar@{-}[dd]\\
 & \bar{f}\bar{\bar{g}}j' \bar{\bar{h}}
k' k \ar@{-}[rd]&\\
\bar{f} i \bar{g}\bar{\bar{h}}k' k \ar@{-}[ru] \ar@{-}[rr]&    &
\bar{f}\bar{\bar{g}}\bar{\bar{h}}k'' k' k.
}}{\xymatrix{
\bar{f} i \bar{g} j \bar{h} k \ar@{-}[rr]^{c_{f,g}*h_*} \ar@{-}[dd]_{f_**c_{g,h}}&&
\bar{f}\bar{\bar{g}} j' j\bar{h} k \ar@{-}[dd]\ar@{-}[rddd]^{c_{fg,h}}\\
&\\
\bar{f} i \bar{g}\bar{\bar{h}}k' k \ar@{-}[rrrd]_{c_{f,gh}} \ar@{-}[rr] && \bar{f}\bar{\bar{g}}j' \bar{\bar{h}}
k' k \ar@{-}[rd]\\
& &   & \bar{f}\bar{\bar{g}}\bar{\bar{\bar{h}}}k'' k' k.
}}
\]
The commutativity of the square is trivial, and that of the triangles follows from
(c) and ${\rm(c')}$. Since the commutativity of the outer boundary is what
had to be proved, this completes the proof of \hyperref[XVII.prop3.3.2]{3.3.2}.

\subsubsection{}\etiquette[3.3.3]{XVII.subsubsec3.3.3}

Suppose that over $(S,i)$ and $(S,p)$ there are given, respectively, fibered
categories equipped with normalized cleavages $F_i$ and $F_p$, and suppose
that $F_i$ and $F_p$ have the same fiber $F(X)$ for every $X\in \Ob S$.

According to \SGA 1 VI, it ``amounts to the same thing'' to give instead:
\begin{itemize}
\item[(i bis)] for every proper morphism $p:X\rightarrow Y$, a
functor $p^{\ast}:F(Y)\rightarrow F(X)$,

\item[(ii bis)] for every composable pair $(p,q)$ of proper
morphisms, an isomorphism of functors
$c_{p,q}:q^{\ast}p^{\ast}\leftrightarrow (pq)^{\ast}$,

\item[(${\rm i'}$ bis)] and\pageoriginale (${\rm ii'}$ bis), analogous data for
$(S,i)$, these data satisfying the conditions (a bis) (b bis) (${\rm a'}$ bis)
(${\rm b'}$ bis) dual to \hyperref[XVII.subsubsec3.3.1]{3.3.1} (a) (b) ${\rm (a')}$ ${\rm (b')}$.

Suppose \emph{given} in addition,

\item[(iii bis)] for every commutative diagram $D$
\[
\vcenter{
\xymatrix{
X \ar[r]^{j'} \ar[d]^{p'}&  Y\ar[d]^{p} \\
Z \ar[r]^{j}& T
}}\quad j,j' \in \Fl(S,i)\text{; }p\sisi{}{,} p' \in \Fl(S,p)
\]
an isomorphism of functors $d(D):{j'}^{\ast}p^{\ast}\leftrightarrow
{p'}^{\ast}j^{\ast}$.
\end{itemize}

We leave to the reader the formulation of the conditions (c bis) and (${\rm c'}$
bis) dual to \hyperref[XVII.prop3.3.2]{3.3.2} (c) and ${\rm (c')}$.

Arguments parallel to the preceding ones then prove the following
proposition.

\begin{enonce*}{Proposition 3.3.4}\etiquette[3.3.4]{XVII.prop3.3.4}
If data {\rm (i bis) (ii bis) (${\rm i'}$ bis) (${\rm ii'}$
bis) (iii bis)} satisfy the conditions {\rm (a bis) (b bis) (${\rm
a'}$ bis) (${\rm b'}$ bis) (c bis) (${\rm c'}$ bis)}, there exists one and
essentially only one category $F$ cofibered over $S$ and equipped
with isomorphisms $\fprod{F}{(S,i)}{S}=F_i$, $\fprod{F}{(S,p)}{S}=F_p$,
these isomorphisms being compatible with the identifications
$F_{iX}=F_{pX}$ $(X\in \Ob S)$ and such that the isomorphisms $d(D)$ of
{\rm (iii bis)} become the composite isomorphisms
${p'}^{\ast}j^{\ast}=(jp')^{\ast}=(pj')^{\ast}={j'}^{\ast}p^{\ast}$.
\end{enonce*}

It should be noted, however, that \hyperref[XVII.prop3.3.4]{3.3.4} is not dual to \hyperref[XVII.prop3.3.2]{3.3.2},
since the hypotheses imposed on $(S,i)$ and $(S,p)$ are not self-dual.

\section{Resolutions. Application to the base change
morphism}\etiquette[4]{XVII.sec4}\pageoriginale

\subsection{Flat resolutions}\etiquette[4.1]{XVII.subsec4.1}

Recall that if $f:(X,\Aa)\rightarrow (Y,\Bb)$ is a
morphism of ringed sites (or ringed topoi), the inverse image under $f$
of a sheaf of left $\Bb$-modules $\Mm$ is the
sheaf of left $\Aa$-modules
$\Aa\otimes_{f^{\ast}\Bb} f^{\ast}\Mm$.

\begin{enonce*}{Proposition 4.1.1}\etiquette[4.1.1]{XVII.prop4.1.1}
Let $f:(X,\Aa)\rightarrow (Y,\Bb)$ be a morphism of
ringed sites, let $K$ be a sheaf of right $\Aa$-modules, and let
$0\rightarrow L\rightarrow M\rightarrow N\rightarrow 0$ be an exact sequence
of sheaves of left $\Bb$-modules. If $M$ and $N$
are flat, then
\begin{enumerate}
\renewcommand{\theenumi}{\roman{enumi}}
\renewcommand{\labelenumi}{(\theenumi)}
\item $L$ is flat,

\item the following sequence is exact:
\[
0\longrightarrow K\otimes_{\Aa}f^{\ast}L\longrightarrow
K\otimes_{\Aa}f^{\ast}M\longrightarrow
K\otimes_{\Aa}f^{\ast}N\longrightarrow 0{\quoi}.
\]
\end{enumerate}
\end{enonce*}

Factor $f$ into the arrows
$(X,\Aa)\xrightarrow{u}(X,f^{\sisi{\ast}{-1}}\Bb)\xrightarrow{v}(Y,\Bb)$.
The functor $v^{\ast}$ is exact and takes flat modules to
flat modules (\sisi{IV 15}{\hyperrefext[V][V.sec1]{1}}), so it suffices to prove \hyperref[XVII.prop4.1.1]{4.1.1} for
$u$. The isomorphism
$$\etiquette[4.1.1.1]{XVII.eq4.1.1.1}
K\otimes_{\Aa}(\Aa\sisi{_{f^{\ast}\Bb}\otimes}{\otimes_{f^{-1}\Bb}} P)=K\otimes_{f^{\sisi{\ast}{-1}}\Bb}P\leqno(4.1.1.1)
$$
then reduces us to the case where $f$ is the identity.

Take an exact sequence $0\rightarrow R\rightarrow P\rightarrow
K\rightarrow 0$ with $P$ flat\sisi{ (\hyperref[XVII.subsubsec4.1.2]{4.1.2})}{}. The rows and columns of the
following diagram are exact:
\[
\xymatrix{
 &      &                          &  0\ar[d]                        &  &\\
 &      & R\otimes L  \ar[r]\ar[d] & P\otimes L  \ar[r]\ar[d] & K\otimes L
 \ar[r]\ar[d]& 0\\
 & 0\ar[r]     & R\otimes M  \ar[r]\ar[d] & P\otimes M  \ar[r]\ar[d] & K\otimes M
 \ar[r]\ar[d]& 0\\
 & 0\ar[r]     & R\otimes N  \ar[r] \ar[d]& P\otimes N  \ar[r]\ar[d] & K\otimes M
 \ar[r] \ar[d]& 0\\
 &      & 0         &   0   &   0   &.
}
\]

Apply\pageoriginale the snake lemma to the first two
columns: we obtain
\[
0\xrightarrow{j}K\otimes L\rightarrow K\otimes M{\quoi},
\]
which proves (ii). Now drop the assumption that $P$ is flat. In the
preceding diagram the three columns and two rows are short exact
sequences; the first row is therefore exact, \ie $L$ is flat.

\setcounter{subsubsection}{1}
\subsubsection{}\etiquette[4.1.2]{XVII.subsubsec4.1.2}
Let $\Ss$ be a subcategory of the category of ringed
sites.\nde{\MakeLinkTarget*{sga4.exp.XVII.editor-note.4.1.2.ringed-site-category}\label{XVII.ndesous}the theory is valid more generally for categories $\Ss$ consisting of ringed sites.} The sheaves of modules on the ringed sites in
$\Ss$ form a category $\Mm$, abelian over
$\Ss$ (\hyperref[XVII.def2.2.2]{2.2.2}), fibered and cofibered, whose fibers are the
opposites of the usual categories of sheaves of modules on a
site. We denote by $K(\Ss)$ the category $K(\Mm)$
(\hyperref[XVII.exem2.2.8]{2.2.8}) and by $K^{+}(\Ss)$ (\resp $K^{-}(\Ss)$) the
full subcategory of $K(\Ss)$ whose fiber categories
are the opposites of the categories $K^{+}(S)$ (\resp $K^{-}(S)$) for
$S\in \Ob \Ss$. The categories $K(\Ss)$,
$K^{+}(\Ss)$, and $K^{-}(\Ss)$ are triangulated over
$\Ss$, fibered and cofibered.

Let\pageoriginale $f:S\rightarrow S'$ be an arrow of $\Ss$. Every
injective object is acyclic (\hyperref[XVII.def1.2.3]{1.2.3}) for the functor
$Rf_{\ast}$, and every flat sheaf of modules is acyclic for the
functor $Lf^{\ast}$ (\hyperref[XVII.prop4.1.1]{4.1.1} and \hyperref[XVII.prop1.2.6]{1.2.6}),
whence the derived functors (\hyperref[XVII.prop1.2.7]{1.2.7})
\begin{align*}
Rf_{\ast}:D^{+}(S) &\longrightarrow D^{+}(S')\tag{4.1.2.1}\\
Lf^{\ast}:D^{-}(S\sisi{}{'}) &\longrightarrow D^{-}(S\sisi{'}{}){\quoi}.\tag{4.1.2.2}
\end{align*}\etiquette[4.1.2.1]{XVII.eq4.1.2.1}
\etiquette[4.1.2.2]{XVII.eq4.1.2.2}%The vertical space following this paragraph is too large.

If $Rf_{\ast}$ (\resp $Lf^{\ast}$) has finite dimension, then
in \eqhyperref[XVII.eq4.1.2.1]{4.1.2.1} (\resp \eqhyperref[XVII.eq4.1.2.2]{4.1.2.2}) one may replace the exponent $+$ (\resp $-$) by
$b$, by $-$ (\resp $+$), or omit it (\hyperref[XVII.prop1.2.10]{1.2.10}).\nde{\MakeLinkTarget*{sga4.exp.XVII.editor-note.4.1.2.spaltenstein-resolutions}\label{XVII.ndeKS}using Spaltenstein resolutions, the exponents may be omitted unconditionally. In the commutative case, see M.~Kashiwara and P.~Schapira, \emph{Categories and Sheaves}, Springer, 2010, cited as [KS] below, 18.6.}

We denote by $D(\Ss)$ (\resp $D^{+}(\Ss)$, \resp
$D^{-}(\Ss)$) the category obtained from $K(\Ss)$ (\resp
$K^{+}(\Ss)$, \resp $K^{-}(\Ss)$) by inverting the
quasi-isomorphisms.

\begin{enonce*}{Scholium 4.1.3}\etiquette[4.1.3]{XVII.sch4.1.3}
(``Trivial duality theorem'').
\begin{enumerate}
\renewcommand{\theenumi}{\alph{enumi}}
\renewcommand{\labelenumi}{(\theenumi)}
\item Assume that all the functors $Rf_{\ast}$ (\resp
$Lf^{\ast}$) have finite dimension. Then:
\begin{itemize}
\item[(i)] The formation of $D(\Ss)$ commutes with every change
of base category $\Ss'\rightarrow \Ss$. The category
$D(\Ss)$ is triangulated over $\Ss$ \emph{(\hyperref[XVII.def2.2.3]{2.2.3})}, and its
fibers identify with the opposites of the categories $D(S)$ for $S\in \Ob
\Ss$.

\item[(ii)] The category $D(\Ss)$ is cofibered (\resp fibered)
over $\Ss$. Let $f:S\rightarrow S'$; if $Lf^{\ast}$ (\resp
$Rf_{\ast}$) has finite dimension, then every object of $D(S')$
(\resp $D(S)$) admits an inverse image (\resp direct image) under $f$ in the sense
of the category $D(\Ss)$ over $\Ss$. The
inverse image and direct image functors identify with $Lf^{\ast}$ and
$Rf_{\ast}$.

\item[(iii)] Every\pageoriginale object of $D^{-}(S')$ (\resp $D^{+}(S)$) admits an
inverse image (\resp direct image) in the sense of $D(\Ss)$, and the
inverse image (\resp direct image) functor identifies with $Lf^{\ast}$
(\resp $Rf_{\ast}$).\nde{the finite-dimension hypotheses in (a) are superfluous. \Cf the preceding N.D.E.}
\end{itemize}

\item \begin{itemize}
\item[(i)] The formation of $D^{-}(\Ss)$ (\resp
$D^{+}(\Ss)$) commutes with every change of base category
$\Ss'\rightarrow \Ss$. The category
$D^{-}(\Ss)$ (\resp $D^{+}(\Ss)$) is triangulated over
$\Ss$, and its fibers identify with the opposites of the categories
$D^{-}(S)$ (\resp $D^{+}(S)$).

\item[(ii)] The category $D^{-}(\Ss)$ (\resp
$D^{+}(\Ss)$) is fibered (\resp cofibered) over $\Ss$;
the ``inverse image'' (\resp ``direct image'')
functors identify with the derived functors $Lf^{\ast}$ (\resp $Rf_{\ast}$).

\item[(iii)] If $f:S\rightarrow S'$ is such that $Rf_{\ast}$ (\resp
$Lf^{\ast}$) has finite dimension, every object of $D^{-}(S)$ (\resp
$D^{+}(S')$) admits a direct image in the sense of $D^{-}(\Ss)$
(\resp an inverse image in the sense of $D^{+}(\Ss)$), and the
``direct image'' (\resp ``inverse image'')
functor identifies with $Rf_{\ast}$ (\resp $Lf^{\ast}$).
\end{itemize}
\end{enumerate}
\end{enonce*}

It suffices to verify that in each case the hypotheses of
\hyperref[XVII.prop2.4.1]{2.4.1} or
\hyperref[XVII.prop2.4.3]{2.4.3} are satisfied, the other assertions following, after a
base change $\{S\}\rightarrow \Ss$, or $I\rightarrow
\Ss$, from \hyperref[XVII.prop2.3.5]{2.3.5} and \hyperref[XVII.prop1.2.10]{1.2.10}.
In case (b), complexes of
flat (\resp flasque) sheaves satisfy the hypotheses of
\sisi{\hyperref[XVII.prop2.4.1]{2.4.1}}{\hyperref[XVII.prop2.4.3]{2.4.3} (\resp \hyperref[XVII.prop2.4.1]{2.4.1})}. In case (a), we use \sisi{\hyperref[XVII.prop2.4.3]{2.4.3}}{\hyperref[XVII.prop2.4.1]{2.4.1} (\resp \hyperref[XVII.prop2.4.3]{2.4.3})}
so that we need only
consider those $f$ for which the dimension of $Rf_{\ast}$ (\resp
$Lf^{\ast}$) is bounded by a fixed $N$. Let us prove that there are ``enough'' complexes $K$ whose components have their
direct (\resp inverse) images all acyclic for the
direct image (respectively, inverse image) functors under consideration. Let $K$ be a
complex\pageoriginale and let $K^{\ast}$ (\resp $K_{\ast}$) be a Cartan-Eilenberg resolution
of $K$ by complexes with flasque components
(\resp flat components).\nde{it is required only that the augmentation $K^p\to (K^*)^{p\bigdot}$ (\resp $(K_*)^{p\bigdot}\to K^p$) be a flasque (\resp flat) resolution for every $p$.} If $n>N$, we see that
$f_{\ast}\tau''_{<n}(K^{\ast})=\tau''_{<\sisi{}{n}}f_{\ast}(K^{\ast})$ (\resp
$f^{\ast}\tau''_{>-n}(K_{\ast})=\tau''_{>-n}f^{\ast}(K_{\ast})$), where
$\tau\sisi{^{\ast}}{''}$ denotes a truncation with respect to the second
differential. For $n>2N$, it follows that $\tau''_{<n}(K^{\ast})$
(\resp $\tau''_{>-n}(K_{\ast})$) satisfies the hypothesis of \sisi{\hyperref[XVII.prop2.4.3]{2.4.3}}{\hyperref[XVII.prop2.4.1]{2.4.1} (\resp \hyperref[XVII.prop2.4.3]{2.4.3})}, and every
complex $K$ is the source (\resp target) of a quasi-isomorphism
whose target (\resp source) is such a complex.

\setcounter{subsubsection}{3}
\subsubsection{}\etiquette[4.1.4]{XVII.subsubsec4.1.4}
Consider a commutative square of ringed sites \eqhyperref[XVII.eq4.1.4.1]{4.1.4.1}
(\cf~\hyperref[XVII.subsec0.11]{0.11})
$$\etiquette[4.1.4.1]{XVII.eq4.1.4.1}
\vcenter{
\xymatrix{
  X' \ar[r]^{g'}\ar[d]_{f'}& X\ar[d]^{f}\\
S' \ar[r]^{g} & S\rlap{.}
}}\leqno{(4.1.4.1)}
$$
We may consider the subcategory of the category of sites consisting
of this diagram;\nde{one must instead consider the category of ringed sites defined by this diagram, which is not strictly speaking a subcategory of the category of ringed sites in general. \Cf \nderef{XVII.ndesous}.} we then use \sisi{\hyperref[XVII.subsubsec2.1.2]{2.1.2}}{\hyperref[XVII.prop2.1.3]{2.1.3}} to
define a ``base change morphism'' from $Lg^{\ast}Rf_{\ast}$ to $Rf'_{\ast}L{g'}^{\ast}$ whenever \hyperref[XVII.sch4.1.3]{4.1.3} is
applicable:

\begin{enonce*}{Proposition 4.1.5}\etiquette[4.1.5]{XVII.prop4.1.5}
The base change morphism of \sisi{\hyperref[XVII.subsubsec2.1.2]{2.1.2}}{\hyperref[XVII.prop2.1.3]{2.1.3}}, from $Lg^{\ast}Rf_{\ast}$
to $Rf'_{\ast}L{g'}^{\ast}$, is defined in each of the
following cases:
\begin{enumerate}
\renewcommand{\theenumi}{\roman{enumi}}
\renewcommand{\labelenumi}{(\theenumi)}
\item We work in the categories $D^{+}$, and $Lg^{\ast}$ and
$L{g'}^{\ast}$ have finite dimension.

\item We work in the categories $D^{-}$, and $Rf_{\ast}$ and
$Rf'_{\ast}$ have finite dimension.

\item We work in the unbounded derived categories, and
$Rf_{\ast}$, $Rf'_{\ast}$, $Lg^{\ast}$, and $L{g'}^{\ast}$ have
finite dimension, as do $Lf^{\ast}$ and $L{f'}^{\ast}$ (or: as do
$Rg_{\ast}$ and $Rg'_{\ast}$).\nde{the finite-dimension hypotheses in (iii) are superfluous. \Cf \nderef{XVII.ndeKS}.}
\end{enumerate}
\end{enonce*}

It\pageoriginale should be noted that the last hypothesis in (iii) looks
rather like a practical joke. The editor does not know how to dispense with it except when
the sites under consideration have enough points (\hyperref[XVII.subsubsec4.2.12]{4.2.12}). The functorial
properties of this base change morphism follow
at once from the description in \sisi{\hyperref[XVII.subsubsec2.1.2]{2.1.2}}{\hyperref[XVII.prop2.1.3]{2.1.3}}.

\setcounter{subsubsection}{5}
\subsubsection{}\etiquette[4.1.6]{XVII.subsubsec4.1.6}
Let $(S,\Aa)$ be a ringed site, and consider the ``tensor product'' functor
\[

% END INCLUDED SGA4_English_current_buildable_source_20260906/source/17/drafts/SGA4_XVII_LINES_2601_3000_English.texfrag

% BEGIN INCLUDED SGA4_English_current_buildable_source_20260906/source/17/drafts/SGA4_XVII_LINES_3001_3400_English.texfrag
\otimes:\Mod\sisi{_d}{}(S,\Aa^{\circ})\times
\Mod\sisi{_g}{}(S,\Aa)\longrightarrow \Mod(S,C)
\]%For the first component of the source, take either right Aa-modules or left Aa-circ-modules.
where $C$ is the center of $\Aa$.

The recipe of \hyperref[XVII.subsubsec1.1.6]{1.1.6}, in applying which
the simple complex associated with a double complex is defined using a
sum, extends this functor to a functor still denoted $\otimes$:
\[
\otimes:K(S,\Aa^{\circ})\times K(S,\Aa)\longrightarrow
K(S,C){\quoi}.
\]

\begin{enonce*}{Proposition 4.1.7}\etiquette[4.1.7]{XVII.prop4.1.7}
With the preceding notation:
\begin{enumerate}
\renewcommand{\theenumi}{\roman{enumi}}
\renewcommand{\labelenumi}{(\theenumi)}
\item The tensor-product functor is left derivable
(\hyperref[XVII.def1.2.1]{1.2.1}) to a ``total tensor product'' functor
\[
{}\overset{L}{\otimes}{}:D^{-}(S,\Aa^{\circ})\times
D^{-}(S,\Aa)\longrightarrow D^{-}(S,C){\quoi}\text{;}
\]

\item a pair $(K,L)$ of bounded-above complexes is adapted to
$\overset{L}{\otimes}$ whenever either $K$ or $L$ has flat components.
\end{enumerate}
\end{enonce*}

The proof is standard.

\begin{enonce*}{Notation 4.1.8}\etiquette[4.1.8]{XVII.not4.1.8}
If\pageoriginale $K\in D^{-}(S,\Aa^{\circ})$ and
$L\in D^{-}(S,\Aa)$, denote by $K\overset{L}{\otimes}L$ the total
tensor product of $K$ and $L$. The sheaf
$H_{k}(K\overset{L}{\otimes}L)=H^{-k}(K\overset{L}{\otimes}L)$ is
denoted by $\Tor_k(K,L)$ and is called the $k$\textsuperscript{th}
local hypertor of $K$ and $L$.
\end{enonce*}

Let $K\in D^{-}(S,\Aa^{\circ})$ and $L\in D^{-}(S,\Aa)$. If
$H^{i}(K)=0$ for $i>k$ and $H^{i}(L)=0$ for $i>\ell$, one checks that
$H^{i}(K\overset{L}{\otimes}L)=0$ for $i>k+\ell$ and that
$H^{k+\ell}(K\overset{L}{\otimes}L)=H^{k}(K)\otimes H^{\ell}(L)$.
Fix $K$ and regard $K\overset{L}{\otimes}L$ as a functor of $L$. If
Definition \hyperref[XVII.def1.2.9]{1.2.9} of cohomological amplitude
as an interval in $\ZZ$ is applied, the upper endpoint of this interval
is trivial in nature: it is the largest integer $i$ such that
$H^{i}(K)\neq0$. The lower endpoint is more interesting:

\begin{enonce*}{Definition 4.1.9}\etiquette[4.1.9]{XVII.def4.1.9}
A complex $K\in\Ob D^{b}(S,\Aa^{\circ})$ is said to have
\emph{\tTor-dimension $\leq n$} if the following equivalent conditions
hold:
\begin{enumerate}
\renewcommand{\theenumi}{\roman{enumi}}
\renewcommand{\labelenumi}{(\theenumi)}
\item For every $\Aa$-module $L$, $\Tor_k(K,L)=0$ for $k>n$.

\item For every complex of $\Aa$-modules $L$ such that $H^i(L)=0$ for
$i<\ell$, one has $H^i(K\overset{L}{\otimes}L)=0$ for $i<\ell-n$.

\item There is a quasi-isomorphism with target $K$ and source a
\sisi{}{ bounded} complex $K'$ whose components are flat and vanish for
$i<-n$.
\end{enumerate}
\end{enonce*}

\setcounter{subsubsection}{9}
\subsubsection{}\etiquette[4.1.10]{XVII.subsubsec4.1.10}
Here are some possible variants:
\begin{enumerate}
\renewcommand{\theenumi}{\roman{enumi}}
\renewcommand{\labelenumi}{(\theenumi)}
\item If $K\in\Ob D^{b}(S,\Aa^{\circ})$ has finite \ttor-dimension
and $L\in\Ob D(S,\Aa)$, then the functor $\overset{L}{\otimes}$ is
defined (\hyperref[XVII.def1.2.1]{1.2.1}) at $(K,L)$.

\item If $K\in\Ob K^{b}(S,\Aa^{\circ})$ has components of finite
\tTor-dimension, the functor $L\mapsto K\otimes L$
$(L\in K(S,\Aa))$ is left derivable.

\item If every module has finite \tTor-dimension,\nde{this condition is
superfluous. In the case where $\Aa$ is commutative, see [KS], 18.6,
cited in \nderef{XVII.ndeKS}.} the functor $\otimes$ derives to a
functor
\[
{}\overset{L}{\otimes}{}:D(S,\Aa^{\circ})\times
D(S,\Aa)\longrightarrow D(S,\Cc){\quoi}.
\]
\end{enumerate}

\subsection{Godement flasque resolutions}\etiquette[4.2]{XVII.subsec4.2}\pageoriginale

\subsubsection{}\etiquette[4.2.1]{XVII.subsubsec4.2.1}
Let $S$ be a topos and $X$ a set. To give a family of point\sisi{}{s}
of $S$ indexed by $X$ is equivalently to give a morphism
$k:\Top(X)\rightarrow S$ from the topos defined by the discrete
topological space $X$ to $S$. For this family to be conservative
(IV \hyperref[IV.def6.4.1]{6.4.1}), it is necessary and sufficient that, for every sheaf $F$
on $S$, the adjunction arrow from $F$ to $k_{\ast}k^{\ast}F$ be a
monomorphism (\cf IV \hyperref[IV.subsubsec6.4.0]{6.4.0}).

\begin{enonce*}{Definition 4.2.2}\etiquette[4.2.2]{XVII.def4.2.2}
Let $k:\Top(X)\rightarrow S$ be a conservative family of points of a
topos $S$. The \emph{Godement resolution}, or \emph{canonical flasque
resolution} (relative to $X$), of an abelian sheaf $F$ on $S$, denoted
$\Cc_X^{\ast}(F)$ (or simply $\Cc^{\ast}(F)$), is the following right
resolution of $F$:
\begin{enumerate}
\renewcommand{\labelenumi}{\theenumi)}
\item $\Cc^{0}(F)=k_{\ast}k^{\ast}F$ and
$\epsilon:F\rightarrow\Cc^{0}(F)$ is the adjunction arrow.

\item Set $d^{-1}=\epsilon$; recursively, for $n\geq0$, define

$\Cc^{n+1}(F)=\Cc^{0}(\coker d^{n-1})$, and define $d^n$ as the
composite arrow
\[
d^{n}:\Cc^n(F)\rightarrow
\coker(d^{n-1})\rightarrow\Cc^{0}(\coker(d^{n-1})){\quoi}.
\]
\end{enumerate}
\end{enonce*}

\begin{enonce*}{Proposition 4.2.3}\etiquette[4.2.3]{XVII.prop4.2.3}
Under the hypotheses of \emph{\hyperref[XVII.def4.2.2]{4.2.2}},
\begin{enumerate}
\renewcommand{\theenumi}{\roman{enumi}}
\renewcommand{\labelenumi}{(\theenumi)}
\item $\Cc^n(F)$ is a flasque sheaf
\emph{\sisi{(V 3.7)}{\hyperrefext[V][V.def4.1]{4.1}}},

\item $\Cc^n(F)$ is an exact functor of $F$,

\item the stalk at a point $x\in X$ of the complex $\Cc^{\ast}(F)$ is
a canonically split resolution of $F_x$.
\end{enumerate}
\end{enonce*}

The\pageoriginale sheaves $\Cc^n(F)$ are flasque because they are
direct images of sheaves (automatically flasque) on the discrete
topological space $X$.

The functors $k^{\ast}$ and $k_{\ast}$ are exact. Hence the functor
$\Cc^0$ is exact. We prove by induction on $n\geq0$ that the functor
$\mathcal{Z}^{n}(F)=\coker(d^{n-1})$ is exact, and therefore so is
$\Cc^{n+1}=\Cc^0\mathcal{Z}^{n}$. Set $\Zz^{-1}(F)=F$. An exact
sequence of sheaves
\[
0\longrightarrow F'\longrightarrow F\longrightarrow F''\longrightarrow0
\]
defines, for $n\geq0$, a diagram
\[
\xymatrix{
        &   0   \ar[d]             &  0   \ar[d]            &  0 \ar[d]\\
0  \ar[r] & \Zz^{n-1}(F\sisi{^1}{'})  \ar[r] \ar[d]& \Zz^{n-1}(F)  \ar[r] \ar[d]& \Zz^{n-1}(F'')
\ar[r] \ar[d]& 0 \\
0  \ar[r] & \Cc^n(F')  \ar[r] \ar[d]& \Cc^n(F)  \ar[r] \ar[d]&
\Cc^n (F'')  \ar[r] \ar[d]& 0 \\
 0  \ar[r] & \Zz^{n}(F')  \ar[r]\ar[d] & \Zz^{n}(F)  \ar[r] \ar[d]& \Zz^{n}(F'')  \ar[r] \ar[d]&
0\\
        &  0               &  0             &  0      &,
}
\]
whose columns are exact, as are, by the induction hypothesis, the
first two rows. The 3\sisi{rd}{\textsuperscript{rd}} row is therefore
exact, proving (ii).

The short exact sequences
\[
0\longrightarrow\mathcal{Z}^{n+1}(F)\longrightarrow
\Cc^{n}(F)\longrightarrow\mathcal{Z}^{n}(F)\longrightarrow0
\]
are special cases of the short exact sequence
$$\etiquette[4.2.3.1]{XVII.eq4.2.3.1}
0\rightarrow F\xrightarrow{\epsilon}\Cc^{0}(F)\rightarrow
\mathcal{Z}^{0}(F)\rightarrow0\leqno(4.2.3.1)
$$
(make\pageoriginale the substitution $F\mapsto\mathcal{Z}^{n-1}(F)$).
To prove (iii), it remains to show that the inverse-image sequence of
\eqhyperref[XVII.eq4.2.3.1]{4.2.3.1} under $k$ is canonically split.
Indeed, the composite arrow
\[
k^{\ast}F\xrightarrow{k^{\ast}(\epsilon)}k^{\ast}\Cc^{0}F=k^{\ast}(k_{\ast}k^{\ast})F=(k^{\ast}k_{\ast})k^{\ast}F\longrightarrow
k^{\ast}F
\]
is the identity.

\begin{enonce*}[remark]{Remark 4.2.4}\etiquette[4.2.4]{XVII.rem4.2.4}
If $(x_i)_{i\in X}$ is a conservative family of points of a site, this
family corresponds to a conservative family of points of the
associated topos, which again makes it possible to define canonical
flasque resolutions.
\end{enonce*}

\begin{enonce*}[remark]{Remark 4.2.5}\etiquette[4.2.5]{XVII.rem4.2.5}
If $S$ is a ringed topos, the functors $\Cc^n$ carry modules to modules.
\end{enonce*}

\begin{enonce*}[remark]{Remark 4.2.6}\etiquette[4.2.6]{XVII.rem4.2.6}
Let $\Ss$ be a category consisting of sites. Suppose that, for each
$S\in\Ob\Ss$, a set $S^d$ and a conservative family
$k_S:\Top(S^d)\rightarrow S$ of points of $S$ indexed by $S^d$ are
given. Suppose that $S^d$ is functorial in $S$ and that the diagrams
\[
\xymatrix{
\Top(S^d)   \ar[r]^{f^d} \ar[d]^{k_S}& \Top(T^d)\ar[d]^{k_T}\\
S        \ar[r]^{f}  & T
}
\]
are commutative. The corresponding canonical flasque resolutions are
then ``functorial in $S$.''
\end{enonce*}

If\pageoriginale $\Ss$ is the category of topological spaces, one may
take $S^d$ to be the underlying set.

If $\Ss$ is the category of the étale sites of schemes $S$ locally of
finite type over a scheme $S_0$, and if, for each $s\in S_0$,
$\bar{k(s)}$ is an algebraic closure of $k(s)$, one may take for $S^d$
the set of points of $S$ with values in one of the $\bar{k(s)}$.

The condition ``of finite type'' may be replaced by a cardinality
condition if the algebraic closures $\bar{k(s)}$ are replaced by
sufficiently large algebraically closed extensions.

\setcounter{subsubsection}{6}
\subsubsection{}\etiquette[4.2.7]{XVII.subsubsec4.2.7}
One advantage of canonical flasque resolutions is that they make it
possible to construct complexes adapted simultaneously to functors of
``tensor-product'' type and to functors of ``direct-image'' type.
\sisi{It}{They} thereby give a perhaps more comprehensible
construction of the base-change morphisms
(\hyperref[XVII.prop4.1.5]{4.1.5}).

\begin{enonce*}{Lemma 4.2.8}\etiquette[4.2.8]{XVII.lem4.2.8}
Let $f:(S,\Aa)\rightarrow(S',\Bb)$ be a morphism of ringed sites, let
$K\in\Ob K(S,\Aa^{\circ})$ and $L\in\Ob K(S',\Bb)$ be complexes, and
let $k:\Top(X)\rightarrow S$ be a conservative set of points of $S$.
Suppose that one of the following conditions holds:
\begin{enumerate}
\renewcommand{\theenumi}{\alph{enumi}}
\renewcommand{\labelenumi}{(\theenumi)}
\item The cohomologies of $K$ and $L$ are bounded above;

\item $\Aa$ has finite \tTor-dimension over
$f^{\sisi{\ast}{-1}}\Bb$, and $K\in\Ob D^b(S,\Aa^{\circ})$ has finite
\tTor-dimension;

\item $L\in\Ob D^b(S',\Bb)$ has finite \tTor-dimension.
\end{enumerate}

Then\pageoriginale the functor $(K,L)\mapsto K\otimes f^{\ast}L$ is
derivable\sisi{}{ on the left} \emph{(\hyperref[XVII.def1.2.1]{1.2.1})} at
$(K,L)$; denote its derived functor by
${}\overset{\LL}{\otimes}\LL f^{\ast}$.
\begin{enumerate}
\renewcommand{\theenumi}{\roman{enumi}}
\renewcommand{\labelenumi}{(\theenumi)}
\item If $x$ is a point of $S$, then
$(K\overset{\LL}{\otimes}\LL f^{\ast}L)_x\simeq
K_x\overset{\LL}{\otimes}_{\Bb_{f(x)}}L_{f(x)}$.

\item The pair $(K,L)$ is left adapted to
${}\overset{\LL}{\otimes}f^{\ast}$
\eqhyperref[XVII.def1.2.1]{1.2.1} if and only if, for every $x\in X$,
the pair $(K_x,L_{f(x)})$ is left adapted to the functor
$\otimes_{\Bb_{f(x)}}$.

\item Under hypotheses $(a)$ or $(b)$, the functor
$K\overset{\LL}{\otimes}f^{\ast}L$ in $L$ has cohomological amplitude
$\subset d$ \emph{(\hyperref[XVII.def1.2.9]{1.2.9})} if and only if,
for every $x\in X$, the complex $K_x$ of $\Bb_{f(x)}$-modules has
\tTor-amplitude $\subset d$.\nde{hypotheses (a), (b), and (c) are all
superfluous. \Cf \nderef{XVII.ndeKS}.}
\end{enumerate}
\end{enonce*}

This lemma, whose proof is standard, expresses the ``pointwise nature''
of the functor under consideration.

\setcounter{subsubsection}{8}
\subsubsection{}\etiquette[4.2.9]{XVII.subsubsec4.2.9}

Let $L$ be a complex. The \emph{Godement resolution} (relative to $X$)
of $L$ is the simple complex (\emph{defined using sums}) associated
with the double complex $\sisi{C}{\Cc}^{\ast}L$
(\hyperref[XVII.subsubsec1.1.4]{1.1.4}). The \emph{Godement resolution
truncated at order} $n$, $\tau''_{\leq n}\sisi{C}{\Cc}^{\ast}L$, is
the simple complex associated with the double complex obtained from
$\sisi{C}{\Cc}^{\ast}L$ by truncation
(\hyperref[XVII.subsubsec1.1.15]{1.1.15}) in the direction of the
Godement differential. By \hyperref[XVII.prop4.2.3]{4.2.3}~(iii), $L$
is, point by point on $X$, homotopy equivalent to the complexes
$\tau''_{\leq n}\sisi{C}{\Cc}^{\ast}L$ and
$\sisi{C}{\Cc}^{\ast}L$; these complexes are therefore ``as good as
$L$'' with respect to functors of ``tensor-product'' type
(\hyperref[XVII.lem4.2.8]{4.2.8}).

\begin{enonce*}{Proposition 4.2.10}\etiquette[4.2.10]{XVII.prop4.2.10}
Let $(S,\Aa)$ be a ringed site, let $X$ be a conservative set of
points of $S$, and let $N,M\in\NN\cup\{\infty\}$, with $N$ or $M$
finite. Let $A$ be the subcategory of $K(S)$ consisting of complexes
$L$ such that:
\begin{enumerate}
\renewcommand{\theenumi}{\alph{enumi}}
\renewcommand{\labelenumi}{(\theenumi)}
\item If $N=\infty$, then $H^i(L)=0$ for $i$ sufficiently large; if
$M=\infty$, then $H^i(L)=0$ for $i$ sufficiently small.

\item For every\pageoriginale $x\in X$ and every $\Aa_x$-module $Q$
of \tTor-dimension $\leq N$, the complex $L_x$ is left adapted to the
functor $\sisi{M_x}{Q}\otimes_{\Aa_x}\ast$.

\item For every morphism of topoi
$f:\sisi{X^{\sim}}{S}\rightarrow E$ of cohomological dimension
$\leq M$, the complex $L$ is right adapted to the functor $f_{\ast}$.
\end{enumerate}

Then the inclusion of $A$ into $K(S)$ induces an equivalence between
the category obtained from $A$ by inverting quasi-isomorphisms and the
category $D^{\ast}(S)$, where $\ast=+$ if $N<M=\infty$, $\ast=-$ if
$M<N=\infty$, and $\ast$ is blank if $N,M<\infty$.
\end{enonce*}

Let $K$ be a complex satisfying (a). If $N=\infty$, then the canonical
arrow $a:K_1=\sisi{\sigma}{\tau}_{\leq i}K\rightarrow K$ is a
quasi-isomorphism for $i$ sufficiently large. If $N<\infty$, set
$K_1=K$. Let $K_1^{\ast}$ be a flat Cartan--Eilenberg resolution and
set $K_2=\tau''_{\geq\sisi{}{-}N}K_1$ (truncation in the direction of
the new differential). The composite arrow from $K_2$ to $K$ is a
quasi-isomorphism, and $K_2$ satisfies (b).

Let $K$ be a complex satisfying (a) and (b). If $M=\infty$, then
$N<\infty$; for $i$ sufficiently large,
$b:K\rightarrow\sisi{\sigma}{\tau}_{\geq i}K_1=K_3$ is a
quasi-isomorphism, \sisi{from}{and} $K_3$ still satisfies (a) and (b).
If $M<\infty$, set $K=K_3$. Let
$K_4=\sisi{\sigma}{\tau}''_{\leq M}\sisi{C}{\Cc}^{\ast}K_3$. Then,
by \hyperref[XVII.prop4.2.3]{4.2.3}~(iii), $K_4$ satisfies
\sisi{(a), (c), and (b)}{(a), (b), and (c)}.

The proposition therefore follows from two successive applications of
\cite[I 2.4.2]{sga4.exp.XVII.bib.CD}.

\begin{enonce*}[remark]{Variants 4.2.11}\etiquette[4.2.11]{XVII.var4.2.11}
If $N<\infty$ (\resp if $M<\infty$, \resp if $N,M<\infty$), one could
equally have worked with bounded-below (\resp bounded-above, \resp
bounded) complexes.
\end{enonce*}

\setcounter{subsubsection}{11}
\subsubsection{}\etiquette[4.2.12]{XVII.subsubsec4.2.12}
Here\pageoriginale is how these constructions may be used to define
the \emph{base-change morphism} of
\hyperref[XVII.subsubsec4.1.4]{4.1.4} under hypotheses more general
than those of \loccit, when the sites under consideration have enough
points.

Let there be a commutative diagram (\cf
(\hyperref[XVII.subsec0.11]{0.11})) of sites
$$\etiquette[4.2.12.1]{XVII.eq4.2.12.1}
\vcenter{
\xymatrix{
X'          \ar[r]^{g'} \ar[d]^{f'}& X\ar[d]^{f}\\
S'          \ar[r]^{g}  & S\rlap{,}
}}\leqno{(4.2.12.1)}
$$
let $\Aa$ be a sheaf of rings on $S$, $\Aa'$ a sheaf of rings on
$S'$, and $L$ a complex of $(\Aa',g^{\ast}\Aa)$-bimodules on $S'$.
Suppose that:
\begin{enumerate}
\renewcommand{\theenumi}{\roman{enumi}}
\renewcommand{\labelenumi}{(\theenumi)}
\item the sites $X$ and $X'$ admit conservative families of points,
$k:Q\rightarrow X$ and $k':Q'\rightarrow X'$, fitting into a
commutative diagram
$$\etiquette[4.2.12.2]{XVII.eq4.2.12.2}
\vcenter{
\sisi{\xymatrix{
X'         \ar[r]^{g'}  \ar[d]^{k'} & X     \ar[d]^{k}  & \\
Q'         \ar[r]^{g'_0} & Q &.
}}{\xymatrix{
Q'         \ar[r]^{g'_0}  \ar[d]^{k'} & Q     \ar[d]^{k}  & \\
X'         \ar[r]^{g'} & X \rlap{.}
}}}\leqno{(4.2.12.2)}
$$

\item One of the following hypotheses holds:
\begin{enumerate}
\renewcommand{\theenumi}{\alph{enumi}}
\renewcommand{\labelenumi}{(\theenumi)}
\item The functors $Rf_{\ast}$ and $Rf'_{\ast}$ have finite
cohomological dimension, and $L$ is bounded above. Set $s=-$.

\item $L$ is bounded below, with components of uniformly bounded
\tTor-dimension (as $g^{\ast}\Aa$-modules). Set
$s=+$.

% END INCLUDED SGA4_English_current_buildable_source_20260906/source/17/drafts/SGA4_XVII_LINES_3001_3400_English.texfrag

% BEGIN INCLUDED SGA4_English_current_buildable_source_20260906/source/17/drafts/SGA4_XVII_LINES_3401_3800_English.texfrag

\item $L$ is bounded, and the hypotheses of (a) and (b) hold. Set
$s=$ blank.
\end{enumerate}
\end{enumerate}

We\pageoriginale denote by
$L\overset{\LL}{\otimes}\LL g^{\ast}:D^{s}(S,\Aa)\rightarrow
D^{s}(S', \Aa')$ the left derived functor (\hyperref[XVII.def1.2.1]{1.2.1}) of the functor
$K\mapsto L\otimes_{g^{\ast}\Aa}g^{\ast}K$ from
$K^{s}(S,\Aa)$ to $K^{s}(S',\Aa')$.

We shall define a base-change morphism
$$\etiquette[4.2.12.3]{XVII.eq4.2.12.3}
\varphi:L\overset{\LL}{\otimes}\LL g^{\ast}Rf_{\ast}K\longrightarrow
Rf'_{\ast}({f'}^{\ast}L\overset{\LL}{\otimes}\LL{g'}^{\ast}K)\leqno(4.2.12.3)
$$
for $K$ in $D^{s}(X,\Aa)$.\nde{\MakeLinkTarget*{sga4.exp.XVII.editor-note.4.2.12.3.unbounded-adjunction}\label{XVII.ndecb}Adjunctions in unbounded derived categories make it possible to define such a morphism without hypotheses (i) and (ii). \Cf \nderef{XVII.ndeKS}.}

By \hyperref[XVII.prop4.2.10]{4.2.10}, it suffices to define it, functorially in $K$,
for the complexes belonging to the category $A$ of \hyperref[XVII.prop4.2.10]{4.2.10}, with
\begin{align*}
N & \geq\sup_i \text{ (tor\sisi{}{-}dimension of $L^{i}$ over
$g^{\ast}$)}\\
M & \geq\text{cohomological dimension of $Rf_{\ast}$}.
\end{align*}
Such a complex is adapted both to the functor $Rf_{\ast}$ and to
the functor ${f'}^{\ast}L\overset{\LL}{\otimes}_{\Aa'}\sisi{}{\LL}{g'}^{\ast}$
(\hyperref[XVII.lem4.2.8]{4.2.8}). The canonical morphisms
\begin{align*}
\alpha &: f_{\ast}K\longrightarrow Rf_{\ast}K\\
\beta &:
{f'}^{\ast}L\overset{\LL}{\otimes}\LL \sisi{g}{{g'}}^{\ast}K\longrightarrow
\sisi{}{{f'}^*}L\otimes \sisi{g}{{g'}}^{\ast}K
\end{align*}
are therefore quasi-isomorphisms. To compute the left-hand side
$I$ of \eqhyperref[XVII.eq4.2.12.3]{4.2.12.3}, it therefore suffices to take a left resolution
(\ie a quasi-isomorphism)
\[
\gamma:K_{1}\longrightarrow f_{\ast}K{\quoi},
\]
with $K_{1}$ adapted to
$L\overset{\LL}{\otimes}\LL g^{\ast}$; we
have
\[
I\simeq L\otimes \sisi{\LL}{} g^{\ast}K_{1}{\quoi}.
\]

Similarly,\pageoriginale to compute the right-hand side $II$ of
\eqhyperref[XVII.eq4.2.12.3]{4.2.12.3}, it suffices to take a right resolution
\[
\delta:{f'}^{\ast}L\otimes {g'}^{\ast}K\longrightarrow K_2{\quoi},
\]
with $K_{2}$ adapted to $Rf'_{\ast}$; then
\[
II\simeq f'_{\ast}K_{2}{\quoi}.
\]
We then define the base-change morphism as the composite
\[
I\simeq L\otimes g^{\ast}K_{1}\xrightarrow{\gamma}L\otimes
g^{\ast}f_{\ast}K\xrightarrow{c}f'_{\ast}({f'}^{\ast}L\otimes
{g'}^{\ast}K)\xrightarrow{\sisi{\alpha}{\delta}}f'_{\ast}K_{2}\simeq II{\quoi},
\]
where $c$ is the usual base-change morphism
at the level of actual complexes.

This base-change morphism has the following characteristic
property, which in principle makes it easy to verify its compatibility
properties.

\begin{enonce*}{Proposition 4.2.13}\etiquette[4.2.13]{XVII.prop4.2.13}
Under the preceding hypotheses, let $K\in \Ob
C^{s}(X\sisi{'}{},\Aa)$, $K_{1}\in \Ob C^{s}(\sisi{X}{S},\Aa)$ and
$K_{2}\in \Ob C^{s}(\sisi{S}{X'},\Aa)$. Moreover, let $u$ be an
$f$-morphism from $K$ to $K_{1}$, \ie $u\in
\Hom(K_{1},f_{\ast}K)\simeq \Hom(f^{\ast}K_{1},K)$, and let $v$ be a
morphism from ${f'}^{\ast}L\otimes {g'}^{\ast}K$ to $K_{2}$. From
$u$, $v$, and the usual base-change morphism, one obtains a
morphism of complexes $c_{u,v}:L\otimes g^{\ast}K_{1}\rightarrow
f'_{\ast}\sisi{k}{K}_2$. Let $u'$, $v'$, and $c'_{uv}$ denote the composite morphisms
\begin{align*}
u' &: K_{1}\xrightarrow{u} f_{\ast}K\rightarrow Rf_{\ast}K\\
v' &: L\overset{\LL}{\otimes}\LL{g'}^{\ast}K\rightarrow
L\otimes {g'}^{\ast}K\xrightarrow{v} K_{2}\\
c'_{u,v} &:
L\overset{\LL}{\otimes}\LL{g'}^{\ast}K_{1}\rightarrow
L\otimes {g'}^{\ast}K_{1}
\xrightarrow{c_{uv}}f'_{\ast}K_{2}\rightarrow Rf'_{\ast}K_{2}{\quoi}.
\end{align*}
The\pageoriginale following diagram is commutative:
$$\etiquette[4.2.13.1]{XVII.eq4.2.13.1}
\vcenter{
\xymatrix{
L \overset{\LL}{\otimes} \sisi{L}{\LL}g^{\ast}Rf_{\ast}K
\ar[r]^{\varphi} & Rf'_{\ast}(f'^{\ast}L
\overset{\LL}{\otimes}\sisi{L}{\LL} g'^{\ast}K) \ar[d]^{Rf'_{\ast}(v')}&\\
L \overset{\LL}{\otimes} \sisi{L}{\LL}g^{\ast}K_1 \ar[u]_{L \overset{\LL}{\otimes} \sisi{L}{\LL}g^{\ast}(u')}
\ar[r]^{c'_{uv}} & Rf'_{\ast}K_2.
}}\leqno{(4.2.13.1)}
$$
\end{enonce*}

Suppose first that $K_{1}=f_{\ast}K$ and that
$K_{2}={f'}^{\ast}L\otimes {g'}^{\ast}K$, and consider
the following commutative diagram of complexes:
\[
\xymatrix{
K^{a}  \ar[r] \ar[d]& K\ar[d]\\
K^{c}  \ar[r] & K^{b}
}
\]
such that $K^{a}$ is adapted to
$\sisi{L}{}{f'}^{\ast}L\overset{\LL}{\otimes}\sisi{L}{\LL}{g'}^{\ast}$,
$K^{b}$ is adapted to $Rf_{\ast}$, and $K^{c}$ is adapted
to both functors. To construct this diagram, we first
construct $K^{a}$, then apply a canonical flasque resolution
(truncated if necessary) to $K^{a}$ and $K$.

Apply the functor $f_{\ast}$ (\resp~$L\otimes {g'}^{\ast}$) to this
diagram, and resolve the complexes in the diagram
$(K_{1},K^{a}_{1},K^{b}_{1}, K^{c}_{1})$ (\resp $(K_{2}, K^{a}_{2},
K^{b}_{2}, K^{c}_{2})$) so as to obtain a commutative cube, in which
$\approx$ denotes a quasi-isomorphism:
\[
\xymatrix@C=2em{
          &         &K'_1 \ar[rdd] \ar[rrrrrr]^{\approx}&        && &         &
          & K_1\ar[rdd]\\  K_1^{a'}\ar[rru]\ar[rdd]\ar[rrrrrr]^(.75){\approx} &         &      &        && & K_1^{a}\ar[rru] \ar[rdd]&       &\\
  &         &      &K^{b'}_1 \ar[rrrrrr] ^(.25){\approx}&& &         &         &   & K^{b}_1\\
          &K_1^{c'} \ar[rru]^{\approx}\ar[rrrrrr] ^(.65){\approx}&      &        && &         &K_1^{c}\ar[rru]^{\approx} \\
}
\]
where\pageoriginale $K'_{1}$, $K^{a'}_{1}$, $K^{b'}_{1}$, and $K^{c'}_{1}$
are adapted to the functor
$L\overset{\LL}{\otimes}\sisi{L}{\LL}g^{\ast}$ (\resp
\[
\xymatrix@C=2em{
        &         &K_2 \ar[rdd] \ar[rrrrrr]^{\approx}&        && &         &
          & K'_2\ar[rdd]\\
K_2^{a}\ar[rru]\ar[rdd]^{\approx}\ar[rrrrrr]^(.75){\approx} &         &      &        && & K_2^{a'}\ar[rru] \ar[rdd]^(.65){\approx}&       &\\
  &         &      &K^{b}_2 \ar[rrrrrr] ^(.25){\approx}&& &         &         &   & K^{b'}_2\\
          &K_2^{c} \ar[rru]\ar[rrrrrr] ^(.65){\approx}&      &
        && &         &K_2^{c'}\ar[rru] & &, \\
}
\]
where $K'_{2}$, $K^{a'}_{2}$, $K^{b'}_{2}$, and $K^{c'}_{2}$ are adapted to
the functor $Rf'_{\ast}$). Using the usual
base-change morphisms, one obtains a commutative prism of complexes over
$S$.
%%%%%%
\newcommand{\circlenum}[1]{\bigcirc \kern -.255cm \lower
  .05em\hbox{${#1}$}}

\[
\def\objectstyle{\scriptstyle}
\xymatrix@C=-.27cm{
        &         & \scriptstyle{L\otimes g^{\ast}K'_1}\ar[rrrrrr]_*+[o][F-]{\scriptscriptstyle{4}} \ar[rdd]^(.35)*+[o][F-]{\scriptscriptstyle{2}}&        && &         &
          & \scriptstyle{L\otimes g^{\ast}K_1} \ar[rdd] \ar[rrrrr]_*+[o][F-]{\scriptscriptstyle{4}}& & & & & \scriptstyle{f'_{\ast}K_2} \ar[rrrrr]_*+[o][F-]{\scriptscriptstyle{4}}\ar[rdd] & & & && \scriptscriptstyle{f'_{\ast}K'_2} \ar[rdd]\\
\scriptstyle{L\otimes g^{\ast}K_1^{a'}} \ar[rru]\ar[rdd]\ar[rrrrrr] &         &
        &        && & \scriptstyle{L\otimes g^{\ast}K_1\sisi{}{^a}} \ar[rru] \ar[rdd]\ar[rrrrr]&       &
        &&& \scriptstyle{f'_{\ast}K^{a}_2}\ar[rdd]\ar[rru] \ar[rrrrr]&
        & & & &
        \scriptstyle{f'_{\ast}K^{a'}_2}\ar[rru]_*+[o][F-]{\scriptscriptstyle{3}}\ar[rdd]^(.7)*+[o][F-]{\scriptscriptstyle{3}}^{\quad \raisebox{1cm}{}{\scriptscriptstyle{-1}}}_(.7){\text{\rotatebox{110}{$\approx$}}}\\
  &         &      & \scriptstyle{L\otimes g^{\ast}K^{\sisi{c}{b}'}_1}  \ar[rrrrrr]
        && &         &         &   & \scriptstyle{L\otimes g^{\ast}K^{b}_2}
        \ar[rrrrr]& & & & & \scriptstyle{f'_{\ast}K^{b}_2} \ar[rrrrr]& & & & & \scriptstyle{f'_{\ast}K^{b'}_2}\\
          &\scriptstyle{L\otimes g^{\ast}K^{c'}_1}
 \ar[rrrrrr]_*+[o][F-]{\scriptscriptstyle{1}}\ar[rru]_{\text{\rotatebox{50}{$\approx$}}}^*+[o][F-]{\scriptscriptstyle{2}}^{\qquad\qquad \raisebox{-.3cm}{}{\scriptscriptstyle{-1}}}
&      &        && &         &\scriptstyle{L\otimes g^{\ast}K^{c}_2} \ar[rru] \ar[rrrrr]_*+[o][F-]{\scriptscriptstyle{1}}& & & & &\scriptstyle{f'_{\ast}K^{c}_2}\ar[rru]\ar[rrrrr]_*+[o][F-]{\scriptscriptstyle{1}}&
        & & & & \scriptstyle{f'_{\ast}K^{c'}_{2'}}\ar[rru]\\
}
\]
and\pageoriginale one concludes by noting that the morphisms in diagram
\eqhyperref[XVII.eq4.2.13.1]{4.2.13.1} are identified with the following composites:

for $\varphi$: the composite of the morphisms \ding{192};

for $L\overset{\LL}{\otimes}\sisi{L}{\LL}g^{\ast}(u')$: the
composite of the morphisms \ding{193}, one of which is the inverse of a
quasi-isomorphism;

for $Rf'_{\ast}(v')$: the composite of the morphisms \ding{194}, one of which is
the inverse of a quasi-isomorphism;

for $\sisi{C}{c}'_{uv}$: the composite of the morphisms \ding{195}.

In the general case, set $K'_{1}=f_{\ast}K$ and
$K'_{2}={f'}^{\ast}L\otimes g^{\ast}K$. We then have a
diagram
\[
\xymatrix{
&L \overset{\LL}{\otimes} \sisi{L}{\LL}g^{\ast}Rf_{\ast}K
\ar[r]^\varphi&
Rf'_{\ast}(f'^{\ast}L \overset{\LL}{\otimes}\sisi{L}{\LL}
g'^{\ast}K) \ar[d] \ar[rd]& \\
L \overset{\LL}{\otimes} \sisi{L}{\LL}g^{\ast}K_1 \ar[r]\ar[ru]& L
\overset{\LL}{\otimes} \sisi{L}{\LL}g^{\ast}K'_1 \ar[r] \ar[u]&
Rf'_{\ast}K'_2 \ar[r]&  Rf'_{\ast}K_2,
}
\]
whose triangles are commutative, as is the inner rectangle
by the preceding argument. The boundary is therefore commutative, and this proves
\hyperref[XVII.prop4.2.13]{4.2.13}.

\subsection{The base-change theorem}\etiquette[4.3]{XVII.subsec4.3}\pageoriginale

We shall reformulate in terms of derived categories the
base-change theorem for proper morphisms \hyperrefext[XII][XII.thm5.1]{5.1}
(iii).

\begin{enonce*}{Theorem 4.3.1}\etiquette[4.3.1]{XVII.thm4.3.1}
Let
\[
\xymatrix{
X' \ar[rr]^{g'} \ar[d]^{f'}&& X\ar[d]^{f}\\
S' \ar[rr]^{g}& & S
}
\]
be a Cartesian diagram of schemes, with $f$ proper. Let $\Aa$ and $\Aa'$ be sheaves
of torsion rings on $S$ and $S'$, and let $L$ be a bounded-above complex
of $(\Aa',g^{\ast}\Aa)$-bimodules. Assume that
the dimensions of the fibers of $f$ are bounded. Hypothesis
\hyperref[XVII.subsubsec4.2.12]{4.2.12} {\rm(a)} is then satisfied, and the base-change morphism
\eqhyperref[XVII.eq4.2.12.3]{4.2.12.3}, between functors from $D^{-}(X,f^{\ast}\Aa)$
to $D^{-}(S',\Aa')$,
\[
\varphi_K:L\overset{\LL}{\otimes}\sisi{L}{\LL}g^{\ast}\sisi{}{Rf_*}K\longrightarrow
Rf'_{\ast}({f'}^{\ast}L\overset{\LL}{\otimes}\sisi{L}{\LL}{g'}^{\ast}K){\quoi},
\]
is an isomorphism.
\end{enonce*}

The functors under consideration are triangulated and ``way-out'' \cite[I 7]{sga4.exp.XVII.bib.RD},
so, to verify that $\varphi_{K}$ is an isomorphism, it
suffices to verify that the $\varphi_{ℋ^{i}(K)}$ are
isomorphisms. We may therefore assume that $K$ is a sheaf
of torsion $f^{\ast}\Aa$-modules concentrated in degree $0$.

To\pageoriginale verify that $\varphi_{K}$ is an isomorphism,
it suffices to do so after applying the forgetful functor from the category of
$\Aa'$-modules to that of
$\ZZ/n\ZZ$-modules, where $n\Aa=n\Aa'=0$.

This reduces us to the case
$\Aa'=\ZZ/n\ZZ$, with $n\Aa=0$. We may
then replace $L$ by an isomorphic complex in the derived
category, bounded above and with components flat over $g^{\ast}\Aa$.
Since the functors under consideration are triangulated and ``way-
out'' in $L$, it suffices to verify that $\varphi_{K}$ is an
isomorphism after replacing $L$ by one of its components. We
may therefore assume that $L$ is a sheaf of flat
$g^{\ast}\Aa$-modules placed in degree $0$. We must
prove that the morphisms
\[
\varphi^{q}:L\otimes g^{\ast}R^{q}f_{\ast}K\longrightarrow
R^{q}f'_{\ast}({f'}^{\ast}L\otimes g^{\ast}K)
\]
are isomorphisms. Let $s$ be a geometric point of $S'$, let $X_s$
be the fiber of $X'$ at $s$, and let $K_s$ be the inverse image of $K$ on
$X_s$. By \hyperrefext[XII][XII.thm5.1]{5.1} (iii), the fiber at $s$ of $\varphi^{q}$ is
the morphism
\[
\varphi^{q}_s:L_{s}\otimes H^{q}(X_s,K)\longrightarrow H^{q}(X_s,
L_s\otimes K){\quoi}.
\]
By $D$. Lazard's theorem \cite{sga4.exp.XVII.bib.La}, the flat $\Aa_s$-module
$L_s$ is the filtered direct limit of free $\Aa_s$-modules
$L_{s,i}$; the functor $H^{q}(X_s,\ast)$ commutes with filtered
direct limits\sisi{}{ (\hyperrefext[VII][VII.prop3.3]{3.3})}. The morphism $\varphi^{q}$ is therefore the direct
limit of the morphisms
\[
\varphi^{q}_i:L_{s,i}\otimes H^{q}(X_s,K)\longrightarrow H^{q}(X_s,
L_{s,i}\otimes K){\quoi}.
\]
The latter are plainly isomorphisms; hence so is
$\varphi^{q}_s$, and $\varphi\sisi{}{^q}$ is an isomorphism.

\begin{enonce*}[remark]{Variants 4.3.2}\etiquette[4.3.2]{XVII.var4.3.2}
\begin{enumerate}
\renewcommand{\theenumi}{\roman{enumi}}
\renewcommand{\labelenumi}{(\theenumi)}
\item When\pageoriginale $L$ is a bounded complex whose components
have finite \tTor-dimension over $g^{\ast}\Aa$, hypotheses (c) of
\hyperref[XVII.subsubsec4.2.12]{4.2.12} are satisfied, and one may work in
the unbounded derived category.\nde{One may work in
the unbounded derived category without these hypotheses. \Cf \nderef{XVII.ndeKS}.}

\item If $L$ is bounded below, with components whose \tTor-dimension is \sisi{finite}{uniformly bounded}
over $g^{\ast}\Aa$, hypotheses (b) of \hyperref[XVII.subsubsec4.2.12]{4.2.12} are satisfied,
even if $\Aa$ and $\Aa'$ are not torsion; the
base-change morphism $\varphi_{K}$ $(K\in
D^{+}(X,f^{\ast}\Aa)$) will be an isomorphism if the sheaves
$ℋ^{i}(K)$ are torsion.
\end{enumerate}
\end{enonce*}

\section{The direct-image functors with proper supports}\etiquette[5]{XVII.sec5}

\subsection{The fundamental construction}\etiquette[5.1]{XVII.subsec5.1}

We shall follow the program set out in \hyperref[XVII.subsec3.1]{3.1}
and define the
functor $Rf_{!}$ for a compactifiable morphism $f$
(\hyperref[XVII.subsec3.2]{3.2}). In the appendix we shall
show how to treat the case of an arbitrary separated morphism of
finite type.\nde{\Cf \nderef{XVII.ndeConrad}.}

\subsubsection{}\etiquette[5.1.1]{XVII.subsubsec5.1.1}

Let $X$ be a topos, $U$ an open subtopos of $X$ (\hyperrefext[IV][IV.def8.3]{8.3}), and $j$ the
inclusion morphism from $U$ to $X$. In \hyperrefext[IV][IV.subsec5.2]{5.2} we defined an ``\emph{extension by the empty set}'' functor from the category of sheaves
of sets on $U$ to that of sheaves of sets on $X$. This
functor $j_{!}$ is left adjoint to the restriction functor
$j^{\ast}$.

We also denote by $j_{!}$ the functor from the category of
sheaves of pointed sets on $U$ (\resp sheaves of groups,
\resp sheaves of $j^{\ast}\Aa$-modules, for
$\Aa$ a sheaf of rings on $X$) to\pageoriginale the
analogous category on $X$, left adjoint to the restriction
functor $j^{\ast}$ (\sisi{IV\quad}{\hyperrefext[IV][IV.prop11.3.1]{11.3.1}}). These three functors $j_{!}$ are
compatible with the forgetful functors from the category of modules to
that of groups, and from the category of groups to that of
pointed sets \footnote[1]{But not to that of sets!
}. If $i:F\rightarrow X$ is the inclusion of the closed complement of
$U$, then $j^{\ast}j_{!}$ is (canonically isomorphic to) the identity
and $i^{\ast}j_{!}$ is the constant functor with value the initial object
(or the final object, which here amounts to the same thing). This also characterizes
$j_{!}$. It follows that the three functors $j_{!}$ under consideration are
exact and embed into the corresponding functors $j_{\ast}$. They are
called the \emph{extension by zero} functors. The formula
$(j_1j_2)^{\ast}=j^{\ast}_2 j^{\ast}_1$ transposes into a
transitivity isomorphism
\[
(j_1 j_2)_{!}=j_{1!}j_{2!}{\quoi}.
\]
The functors $j_{!}$ are compatible with base change:

\begin{enonce*}{Lemma 5.1.2}\etiquette[5.1.2]{XVII.lem5.1.2}
Let $f:X\rightarrow Y$ be a morphism of topoi, let $V$ be an open subtopos of $Y$,
and let $U$ be its inverse image in $X$:
\[
\xymatrix{
U \ar[rr]^{f'}  \ar[d]^{j'}& & V\ar[d]^{j}\\
X \ar[rr]^{f}   && Y
}
\]
There is then a unique base-change isomorphism
$f^{\ast}j_{!}\isomto j'_{!}{f'}^{\ast}$ making the
diagram
$$\etiquette[5.1.2.1]{XVII.eq5.1.2.1}
\vcenter{
\xymatrix{
f^{\ast}j_{!}    \ar[r]^{\sim} \ar@{^{(}->}[d] & j'_{!}f'^{\ast} \ar@{^{(}->}[d]     \\
f^{\ast}j_{\ast}  \ar[r]^{\alpha}&j'_{\ast}f'^{\ast},
}}\leqno{(5.1.2.1)}
$$
commutative, where $\alpha$ is the base-change morphism.
\end{enonce*}

\begin{enonce*}[remark]{Proof}
The\pageoriginale base-change morphism
$\varphi:j^{\ast}f_{\ast}\rightarrow f'_{\ast}{j'}^{\ast}$ is
trivially an isomorphism. It transposes into a base-change isomorphism
$\varphi':j'_{!}{f'}^{\ast}\isomto
f^{\ast}j_{!}$. To verify that diagram \eqhyperref[XVII.eq5.1.2.1]{5.1.2.1} is
commutative, it suffices to see this after restriction to $U$, where it
is trivial.
\end{enonce*}

\setcounter{subsubsection}{2}
\subsubsection{}\etiquette[5.1.3]{XVII.subsubsec5.1.3}
Under the hypotheses of \hyperref[XVII.subsubsec5.1.1]{5.1.1}, with $X$ ringed by a sheaf
of rings $\Aa$, the functor $j_{!}$ on modules is
exact, hence passes trivially to the derived categories and defines a
functor, still denoted $j_{!}$, from $D(U,j^{\ast}\Aa)$ to
$D(X,\Aa)$. We still have transitivity isomorphisms
$$\etiquette[5.1.3.1]{XVII.eq5.1.3.1}
(j_1 j_2)\sisi{}{!}=j_{1!}j_{2!}{\quoi}.\leqno(5.1.3.1)
$$

At the level of derived categories, we still have an
adjunction formula
$$\etiquette[5.1.3.2]{XVII.eq5.1.3.2}
\Hom_{D(X)}(j_{!}K,L)\isomto \Hom_{D(U)}
(K,j^{\ast}L){\quoi}.\leqno(5.1.3.2)
$$

\subsubsection{}\etiquette[5.1.4]{XVII.subsubsec5.1.4}
Let $f:X\rightarrow Y$ be a proper morphism, let $\Aa$ be a sheaf of
\emph{torsion} rings on $Y$, and endow $X$ with
the inverse-image sheaf of rings of $\Aa$. Assume that the
dimensions of the fibers of $f$ are bounded above by an integer $n$ (as is
the case if $Y$ is quasi-compact). We then know that, for every
sheaf of $f^{\ast}\Aa$-modules $F$ on $X$, the sheaves
$R^{q}f_{\ast}F$ vanish for $q>2n$ (\hyperrefext[XII][XII.corbis5.3]{5.3 bis}). The functor
$f_{\ast}$ therefore admits a derived functor (\hyperref[XVII.prop1.2.10]{1.2.10})
\[
Rf_{\ast}:D(X,f^{\ast}\Aa)\longrightarrow
D(Y,\Aa){\quoi}.
\]

% END INCLUDED SGA4_English_current_buildable_source_20260906/source/17/drafts/SGA4_XVII_LINES_3401_3800_English.texfrag

% BEGIN INCLUDED SGA4_English_current_buildable_source_20260906/source/17/drafts/SGA4_XVII_LINES_3801_4200_English.texfrag

If\pageoriginale $f$ is the composite $gh$ of two proper morphisms,
there is a transitivity formula
$Rf_{\ast}=Rg_{\ast}Rh_{\ast}$, which, as usual, can be expressed
in terms of cofibered categories (Scholium \sisi{\hyperref[XVII.subsubsec5.1.3]{5.1.3}}{\hyperref[XVII.sch4.1.3]{4.1.3}}).

\subsubsection{}\etiquette[5.1.5]{XVII.subsubsec5.1.5}
Consider a commutative diagram of schemes
$$\etiquette[5.1.5.1]{XVII.eq5.1.5.1}
\vcenter{
\xymatrix{
U        \ar@{^{(}->}[r]_{j_1} \ar[d]^{g}      & X  \ar[d]^{f}\\
V        \ar@{^{(}->}[r]_{j_2}       & Y,
}}\leqno{(5.1.5.1)}
$$
in which $j_1$ and $j_2$ are open immersions and $f$, $g$ are
proper morphisms whose fiber dimensions are bounded above by an
integer $n$. Let $\Aa$ be a sheaf of torsion rings on
$Y$, and equip $U$, $V$, and $X$ with the appropriate inverse images of
$\Aa$. By \eqhyperref[XVII.eq5.1.3.2]{5.1.3.2}, if $K\in D(U)$, defining a morphism
$$\etiquette[5.1.5.2]{XVII.eq5.1.5.2}
d:j_{2!}Rg_{\ast}K\longrightarrow Rf_{\ast}j_{1!}K\leqno{(5.1.5.2)}
$$
amounts to defining a morphism
\[
d':Rg_{\ast}K\longrightarrow j^{\ast}_2 Rf_{\ast}j_{1!}K{\quoi}.
\]

The fiber product $X'=V\times {}_YX$ fits into a commutative
diagram
$$\etiquette[5.1.5.3]{XVII.eq5.1.5.3}
\vcenter{
\xymatrix{
U \ar@{^{(}->}[r]^{k} \ar[d]^{g} & X'  \ar@{^{(}->}[r]_{j'_2} \ar[d]^{f'}& X \ar[d]^{f}          &\\
V         \ar@{=}[r] & V         \ar@{^{(}->}[r]_{j_2}  & Y   \rlap{,}
}}\leqno{(5.1.5.3)}
$$
and,\pageoriginale trivially, we have $j^{\ast}_2Rf_{\ast}j_{1!}K\simeq
Rf'_{\ast}(j'_2{}^\ast j_{1!}K)\simeq Rf'_{\ast}k_{!}K$.

The morphism $k$ is a proper open immersion; hence
$k_{!}=k_{\ast}$, giving an isomorphism $Rf'_{\ast}k_{!}K\simeq
Rf'_{\ast}k_{\ast}K\simeq Rg_{\ast}K$, which provides the desired isomorphism
$d'$ and the morphism $d$. Moreover, the restriction
of $d$ to $V$ is an isomorphism.

\begin{enonce*}{Lemma 5.1.6}\etiquette[5.1.6]{XVII.lem5.1.6}
The morphism \eqhyperref[XVII.eq5.1.5.2]{5.1.5.2} is an isomorphism.
\end{enonce*}

It remains to verify that the restriction of $d$ to $Y-V=F$ is an
isomorphism. Let $X''=\fprod{X}{F}{Y}$ and apply the proper
base-change theorem (\hyperref[XVII.thm4.3.1]{4.3.1}) to the Cartesian square
\[
\xymatrix{
X'' \ar[r]_{i'} \ar[d]^{f''}& X\ar[d]^{f\sisi{'}{}}\\
F   \ar[r]_{i}& Y\rlap{,}
}
\]
with $\Aa'=i^{\ast}\Aa$ and $L=\Aa'[0]$. We
obtain:
$i^{\ast}Rf_{\ast}j_{1!}K\simeq Rf''_{\ast}{i'}^{\ast}j_{1!}K\simeq
Rf''_{\ast}0\simeq 0$, so that both sides of \eqhyperref[XVII.eq5.1.5.2]{5.1.5.2} are
zero outside $V$ and $d$ is an isomorphism.

\setcounter{subsubsection}{6}
\subsubsection{}\etiquette[5.1.7]{XVII.subsubsec5.1.7}
Let $S$ be a scheme ringed by a sheaf of torsion rings
$\Aa$. For every $S$-scheme $\varphi:X\rightarrow S$, set
$\Aa_X=\varphi^{\ast}\Aa$. Work in the category $(S)$ of
quasi-compact, quasi-separated $S$-schemes and
$S$-compactifiable morphisms (\hyperref[XVII.subsec3.2]{3.2}). By
\hyperref[XVII.prop3.2.3]{3.2.3}, this category, the subcategory $(S,i)$ of open
immersions, and the subcategory $(S,p)$ of proper morphisms
satisfy the hypotheses of \hyperref[XVII.subsubsec3.2.4]{3.2.4}.

We\pageoriginale propose to apply the theory of \hyperref[XVII.prop3.3.2]{3.3.2} to the following data
(notation of \hyperref[XVII.subsubsec3.3.1]{3.3.1})

\begin{itemize}
\item[$\alpha$)] $F(X)=D(X,\Aa_X)$ for $\varphi:X\rightarrow
S$ in $S$:
\begin{itemize}
\item[(i)] For $p:X\rightarrow Y$ proper, take the functor
$p_{\ast}$ to be the functor $Rp_{\ast}$.

\item[(ii)] Use the usual transitivity isomorphisms
$R(pq)_{\ast}=Rp_{\ast}Rq_{\ast}$.

\item[${\rm(i')}$] For an open immersion $i:X\rightarrow Y$, take
the functor $i_{\ast}$ to be the functor $i_{!}$.

\item[${\rm(ii')}$] Use the transitivity isomorphisms
\eqhyperref[XVII.eq5.1.3.1]{5.1.3.1}.

\item[${\rm(iii)}$] For every commutative diagram \eqhyperref[XVII.eq5.1.5.1]{5.1.5.1}, take
the isomorphism \eqhyperref[XVII.eq5.1.5.2]{5.1.5.2} for $d$.
\end{itemize}
\end{itemize}

The conditions (a) (b) (c) ${\rm(a')}$ ${\rm(b')}$
${\rm(c')}$ of \hyperref[XVII.prop3.3.2]{3.3.2} are verified, proving part (a) of:

\begin{enonce*}{Theorem 5.1.8}\etiquette[5.1.8]{XVII.thm5.1.8}
\begin{enumerate}
\renewcommand{\theenumi}{\alph{enumi}}
\renewcommand{\labelenumi}{(\theenumi)}
\item In the category $(S)$ defined in \hyperref[XVII.subsubsec5.1.7]{5.1.7}, there exists one, and
essentially only one, way to define, for every arrow
$f:X\rightarrow Y$, a functor
\[
Rf_{!}:D(X,\Aa_X)\longrightarrow D(Y,\Aa_Y){\quoi},
\]
and, for every composite morphism
$f=gh:X\xrightarrow{h}Y\xrightarrow{g}Z$, transitivity
isomorphisms $c_{\sisi{f,g}{g,h}}$ between $Rf_{!}$ and $Rg_{!}$ $Rh_{!}$ such
that
\begin{enumerate}
\renewcommand{\theenumii}{\roman{enumii}}
\renewcommand{\labelenumii}{(\theenumii)}
\item The $c_{f,g}$ satisfy the ``cocycle condition''
\[
(c_{f,g}\ast Rh_{!})\circ c_{fg,h}=(Rf_{!}\ast c_{g,h})\circ
c_{f,gh}{\quoi}.
\]

\item When restricted to proper morphisms (\resp to
open immersions), one recovers the variance theory of
\hyperref[XVII.subsubsec5.1.4]{5.1.4} (\resp~\hyperref[XVII.subsubsec5.1.3]{5.1.3}).

\item For\pageoriginale every diagram \eqhyperref[XVII.eq5.1.5.1]{5.1.5.1}, the isomorphism
$d$ of \eqhyperref[XVII.eq5.1.5.2]{5.1.5.2} is the composite isomorphism $c_{f,j\sisi{'}{_1}} \circ c^{-1}_{j\sisi{,f'}{_2,g}}$.
\end{enumerate}

\item The functors $Rf_{!}$ are ``way out'' \emph{\cite[I 7]{sga4.exp.XVII.bib.RD}} and
triangulated.
\end{enumerate}
\end{enonce*}

Every functor $Rf_{!}$ is the composite of a functor $Rg_{\ast}$, for $g$
proper with quasi-compact target, hence with fibers of bounded dimension, and
of a functor $j_{!}$ for $j$ an open immersion. These functors are
way out and triangulated, and hence so is $Rf_{!}$.

\setcounter{subsubsection}{8}
\paragraph{}\etiquette[5.1.8.1]{XVII.par5.1.8.1}

The transitivity isomorphisms of \hyperref[XVII.thm5.1.8]{5.1.8} (a) are often used via
the ``spectral sequence of composite functors'' deduced from them,
valid for $K\in \Ob D(X,\Aa_X)$:
$$\etiquette[5.1.8.2]{XVII.eq5.1.8.2}
E^{pq}_2=R^{p}g_! R^{q}h_{!}(K)\Longrightarrow
R^{p+q}f_{!}(K){\quoi}.\leqno{(5.1.8.2)}
$$

To obtain it, note that by \hyperref[XVII.thm5.1.8]{5.1.8} b) and \textsc{Verdier}
\cite{sga4.exp.XVII.bib.V2}, for every
$L\in \Ob D(Y,\Aa_Y)$, there is a spectral sequence
$$\etiquette[5.1.8.3]{XVII.eq5.1.8.3}
E^{pq}_2=R^{p} g_!(ℋ^{q}(L))\Longrightarrow
R^{p+q}g_!(L){\quoi},\leqno{(5.1.8.3)}
$$
and \eqhyperref[XVII.eq5.1.8.2]{5.1.8.2} is the special case of \eqhyperref[XVII.eq5.1.8.3]{5.1.8.3} with $L=Rh_{!}K$.

\begin{enonce*}{Definition 5.1.9}\etiquette[5.1.9]{XVII.def5.1.9}
\begin{enumerate}
\renewcommand{\theenumi}{\roman{enumi}}
\renewcommand{\labelenumi}{(\theenumi)}
\item We call the ``\emph{
\sisi{total}{total} direct-image functors with proper supports}'' the functors constructed in \hyperref[XVII.thm5.1.8]{5.1.8}. We shall omit
the adjective ``total'' when there is no risk of
confusion.

\item The $q$\textsuperscript{th} \emph{direct-image functor with proper
support,} denoted $R^{q}f_{!}$, is the composite functor
$\Hh^{q}\circ Rf_{!}$.

\item If $X$ is a scheme over a separably closed field $k$ that is
embeddable in a proper scheme over $k$, identify sheaves of
modules on\pageoriginale $\Spec (k)$ with ordinary modules, and denote
the functor $R^{q}f_{!}$ by one of the notations $H^{q}_c(X,\quad)$,
$H^{q}_{!}(X,\quad)$, or simply $H^{q}_c(\quad)$ if there is no risk
of confusion. It is then called the ``$q$\textsuperscript{th} \emph{hypercohomology
group with proper supports},'' or the ``$q$\textsuperscript{th} \emph{cohomology group with proper supports}'' if it
is applied to a complex concentrated in degree $0$.
\end{enumerate}
\end{enonce*}

\setcounter{subsubsection}{9}
\setcounter{paragraph}{0}
\paragraph{}\etiquette[5.1.9.1]{XVII.par5.1.9.1}
The notation $Rf_{!}$ is abusive in that \emph{the functor $Rf_{!}$
is not the right derived functor of the functor $f_{!}$} that will be defined in
\hyperref[XVII.def6.1.2]{6.1.2}. Some prefer the notation $R_{!}f$.

\subsubsection{}\etiquette[5.1.10]{XVII.subsubsec5.1.10}
The ingredients in the proof of \hyperref[XVII.thm5.1.8]{5.1.8} are the
proper
base-change theorem and the identity between $j_{\ast}$ and $j_{!}$
when $j$ is a proper open immersion. This makes it possible to state
various variants of \hyperref[XVII.thm5.1.8]{5.1.8}; it is left to the reader
to give them a precise meaning and to verify them.

\begin{enonce*}{Variant 5.1.11}\etiquette[5.1.11]{XVII.var5.1.11}
In the category $(S)$, there exists one, and essentially only one,
way to define, for every arrow $f:X\rightarrow Y$, a functor
$f_{!}$ from the category of sheaves of pointed sets on $X$ to
the category of sheaves of pointed sets on $Y$, and, for every
composite morphism $f=gh$, transitivity isomorphisms satisfying
conditions {\rm (i)\sisi{,}{} (ii) (iii)} analogous to those of \sisi{\hyperref[XVII.subsubsec5.1.4]{5.1.4}}{\hyperref[XVII.thm5.1.8]{5.1.8}}.
\end{enonce*}

See \hyperref[XVII.def6.1.2]{6.1.2} for a direct description of this functor.

\begin{enonce*}{Variant 5.1.12}\etiquette[5.1.12]{XVII.var5.1.12}
For\pageoriginale every compactifiable morphism
$f:X\rightarrow Y$, one can define a functor $R^{1}f_{!}$ from the category
of sheaves of \iInd-finite groups on $X$ to the category of
sheaves of pointed sets on $Y$, together with, for every
compactification $f=\bar{f}j$ of $f$, an isomorphism
$R^{1}f_{!}\simeq R^{1}\bar{f}_{\ast}\circ j_{!}$.
\end{enonce*}

No transitivity formula is postulated here; the only problem is
to show that the functor $R^{1}f_{!}$ is independent of the chosen
compactification of $f$. Let $p:(\bar{f}_1, i_{1})\rightarrow
(\bar{f}_{2}, i_{2})$ be a morphism of compactifications, \ie a
diagram
$$\etiquette[5.1.12.1]{XVII.eq5.1.12.1}
\vcenter{
\xymatrix{
X   \ar[rr]^{i_1} \ar[rrd]^{i_2} \ar[dd]_{f}& & \bar{X}_1 \ar[d]^{p}\\
    & &  \bar{X}_2 \ar[lld]^{\bar{f}_2}\\
Y   & & .
}}\leqno{(5.1.12.1)}
$$
where $\bar{f}_{1}=\bar{f}_2\circ p$, with $p$ proper. If $G$ is a
sheaf \sisi{of}{of} groups on $X$, there is an exact sequence of pointed
sets (with $\alpha$ injective)
\[
0\rightarrow
R^{1}\sisi{f}{\bar{f}}_{2\ast}(p_{\ast}i_{1!}G)\xrightarrow{\alpha}R^{1}\sisi{f}{\bar{f}}_{1\ast}(i_{1!}G)\rightarrow
\bar{f}_{2\ast}R^{1}p_{\ast}(i_{1!}G){\quoi},
\]
and a morphism from $i_{2!}G$ to $p_{\ast}i_{1!}G$, whence, by
composition, a morphism $\beta$ from $R^{1}\sisi{f}{\bar{f}}_{2\ast}i_{2!}G$ to
$R^{1}\sisi{f}{\bar{f}}_{1\ast}i_{1!}G$. Reasoning as in \hyperref[XVII.par3.3.2.1]{3.3.2.1}, one sees that the
problem is to prove that this morphism is always an isomorphism
when $G$ is \iInd-finite.

Applying\pageoriginale the proper base-change theorem
to $p$ and the sheaf of sets $i_{1!}G$, one checks that
$p_{\ast}i_{1!}G=i_{2!}G$, yielding the exact sequence
\[
0\rightarrow R^{1}
\sisi{f}{\bar{f}}_{2\ast}(i_{2!}G)\xrightarrow{\beta}R^{1}\sisi{f}{\bar{f}}_{1\ast}(i_{1!}G)\rightarrow
\bar{f}_{2\ast}R^{1}p_{\ast}i_{1!}G{\quoi}.
\]
The morphism $\beta$ is injective (since $\alpha$ is). Applying
the proper base-change theorem to $p$ and the sheaf of
groups $G$, assumed \iInd-finite, one finds that
$R^{1}p_{\ast}i_{1!}G=0$. Thus $\beta$ is an epimorphism, which
completes the proof of \hyperref[XVII.var5.1.12]{5.1.12}.

\begin{enonce*}{Definition 5.1.13}\etiquette[5.1.13]{XVII.def5.1.13}
For every site $X$, denote by $D_{\tors}(X)$ the full
subcategory of the derived category of the category of abelian sheaves
on $X$ consisting of the complexes whose cohomology sheaves are
torsion.
\end{enonce*}

\begin{enonce*}{Variant 5.1.14}\etiquette[5.1.14]{XVII.var5.1.14}
In the category $(S)$, there exists one, and essentially only one,
way to define, for every arrow $f:X\rightarrow Y$, a
triangulated and ``way out'' functor \emph{\cite[I 7]{sga4.exp.XVII.bib.RD}}
\[
Rf_{!}:D^{+}_{\tors}(X)\longrightarrow D^{+}_{\tors}(Y)
\]
and, for every composite morphism $f=gh$, transitivity isomorphisms
satisfying conditions \mbox{\rm (i) (ii) (iii)}
analogous to those of \hyperref[XVII.subsubsec5.1.4]{5.1.4}.
\end{enonce*}

\begin{enonce*}{Variant 5.1.\sisi{}{1}5}\etiquette[5.1.15]{XVII.var5.1.15}
Work in the category of $S$-schemes ringed by
sheaves of torsion rings, the arrows being morphisms of ringed schemes
inducing (\hyperref[XVII.subsec0.8]{0.8}) an $S$-compactifiable morphism of
schemes. For every arrow
\[
f:(X,\Aa)\longrightarrow (Y,\Bb),
\]
let\pageoriginale $Rf_{!}$ be the composite of the forgetful functor from the category of
$\Aa$-modules to that of $f^{\sisi{\ast}{-1}}\Bb$-modules, followed
by the usual $Rf_{!}$. These functors have the variance properties of the
usual functors $Rf_{!}$.
\end{enonce*}

Note, however, that the notation $Rf_{!}$ is reasonable here only
if $\Aa$ is finite over $f^{\sisi{\ast}{-1}}\Bb$, as one
sees by taking $X=Y=\Spec(k)$, with $k$ algebraically closed.

\setcounter{subsubsection}{15}
\subsubsection{}\etiquette[5.1.16]{XVII.subsubsec5.1.16}
Let $f:X\rightarrow S$ be a compactifiable morphism, let $j:U\rightarrow
X$ be an open immersion into $X$, let $i:F\rightarrow X$ be the closed complement,
let $\Aa_S$ be a sheaf of torsion rings on $S$,
let $\Aa_X$ be its inverse image on $X$, and let $K\in \Ob
D(X,\Aa_X)$.

The exact sequence of complexes
\[
0\longrightarrow j_{!}j^{\ast}K\longrightarrow K\longrightarrow
i_{!}i^{\ast}K\longrightarrow 0
\]
gives rise to a distinguished triangle in $D(X,\Aa_X)$ by
\textsc{Verdier} \cite[\sisi{I 2.4}{II 1.1.5}]{sga4.exp.XVII.bib.CD}. Its image under $Rf_{!}$ is a distinguished triangle
(\hyperref[XVII.thm5.1.8]{5.1.8} (b))
$$\etiquette[5.1.16.1]{XVII.eq5.1.16.1}
R(fj)_{!}j^{\ast}K\longrightarrow Rf_{!}K\longrightarrow
R(fi)_{!}i^{\ast}K\longrightarrow R(fj)_{!}j^{\ast}K[1]\ldots\leqno{(5.1.16.1)}
$$
which defines a long exact cohomology sequence
$$\etiquette[5.1.16.2]{XVII.eq5.1.16.2}
\ldots \sisi{}{\to}R^{q}(fj)_{!}j^{\ast}K\rightarrow
R^{q}f_{!}K\rightarrow R^{q}(fi)_{!}i^{\ast}K \xrightarrow{j} R^{q+1}(fj)_{!}j^{\ast}K\rightarrow\ldots.\leqno{(5.1.16.2)}
$$

In the special case where $S$ is the spectrum of a separably
closed field and $K$ is a sheaf $𝒢$ concentrated in degree $0$, the
sequence \eqhyperref[XVII.eq5.1.16.2]{5.1.16.2} becomes
$$\etiquette[5.1.16.3]{XVII.eq5.1.16.3}
\ldots H^{q}_c(U,𝒢)\rightarrow H^{q}_c (X,𝒢)\rightarrow
H^{q}_c(F,𝒢)\xrightarrow{j} H^{q+1}_c(U,𝒢)\ldots\leqno{(5.1.16.3)}
$$

\subsubsection{}\etiquette[5.1.17]{XVII.subsubsec5.1.17}

Consider\pageoriginale a diagram of $S$-compactifiable schemes
\[
\xymatrix{
X \ar[rr]^{u} \ar[rd]_{f}& & Y\ar[ld]^{g}\\
  & S &,
}
\]
with $u$ \emph{proper}. Equip $X$ and $Y$ with the inverse images
of a sheaf of torsion rings $\Aa_S$ on $S$. For $K\in
D(Y,\Aa_Y)$, the evident (adjunction) morphism $K\rightarrow
Ru_{\ast}u^{\ast}K$ defines, upon applying the functor $Rg_{!}$, a
morphism
$$\etiquette[5.1.17.1]{XVII.eq5.1.17.1}
u^{\ast}:Rg_{!}K\longrightarrow
Rg_{!}Ru_{\ast}u^{\ast}K=Rg_{!}Ru_{!}u^{\ast}K\sisi{\eqsim}{\simeq} Rf_{!}u^{\ast}K\leqno(5.1.17.1)
$$
(\emph{contravariance of cohomology with proper supports with respect to
proper morphisms}).

When $S$ is the spectrum of an algebraically closed field and $K$ is
the sheaf $𝒢$ placed in degree $0$, this morphism becomes
$$\etiquette[5.1.17.2]{XVII.eq5.1.17.2}
u^{\ast}:H^{q}_c(Y,𝒢)\longrightarrow H^{q}_c(X,𝒢)\leqno(5.1.17.2)
$$

\subsection{The base-change theorem}\etiquette[5.2]{XVII.subsec5.2}

\subsubsection{}\etiquette[5.2.1]{XVII.subsubsec5.2.1}

Let $f:X\rightarrow Y$ be a morphism of topoi, let $V$ be an open of $Y$,
let $U=f^{-1}(V)$ be its inverse image in $X$, let $\Aa$ be a sheaf
of rings on $Y$, let $\Aa'$ be a sheaf of rings on $X$, and let $L$
be a bounded-above complex of $(\Aa',
f^{\ast}\Aa)$-bimodules:
$$\etiquette[5.2.1.1]{XVII.eq5.2.1.1}
\vcenter{
\xymatrix{
U \ar[rr]^{f'} \ar@{^{(}->}[d]^{j'}& & V \ar@{^{(}->}[d]^{j}\\
X \ar[rr]_{f} & & Y.
}}\leqno{(5.2.1.1)}
$$

There\pageoriginale is then a base-change isomorphism
\eqhyperref[XVII.eq5.1.2.1]{5.1.2.1} between functors from $D(V,\Aa)$ to
$D(X,f^{\ast}\Aa):f^{\ast}j_{!}\isomto
j'_{!}{f'}^{\ast}$. Moreover, there is an evident isomorphism
between functors from $D^{-}(U,{f'}^{\ast}\Aa)$ to
$D^{-}(X,\Aa'):L\overset{\LL}{\otimes}_{f^{\ast}\Aa}j'_{!}K\simeq
j'_{!}({j'}^{\ast}L\overset{\LL}{\otimes}_{f^{\ast}\Aa}K)$.
To construct it, it suffices to note that if a complex $K$ is
adapted to the functor
${j'}^{\ast}L\overset{\LL}{\otimes}_{f^{\ast}\Aa}$,
then the complex $j'_{!}K$ is adapted to the functor
$L\overset{\LL}{\otimes}_{f^{\ast}\Aa}$,
and that there is an isomorphism of the preceding type at the level of
complexes. Composing these two isomorphisms gives a
base-change isomorphism between functors from
$D^{-}(V,\Aa)$ to $D^{-}(X,\Aa')$:
$$\etiquette[5.2.1.2]{XVII.eq5.2.1.2}
L\overset{\LL}{\otimes}_{f^{\ast}\Aa}f^{\ast}j_{!}K\isomto j'_{!}({j'}^{\ast}L\overset{\LL}{\otimes}_{{f\sisi{'}{}}^{\ast}\Aa}{f'}^{\ast}K){\quoi}.\leqno(5.2.1.2)
$$

If $L$ is bounded with components of finite \ttor-dimension over

% END INCLUDED SGA4_English_current_buildable_source_20260906/source/17/drafts/SGA4_XVII_LINES_3801_4200_English.texfrag

% BEGIN INCLUDED SGA4_English_current_buildable_source_20260906/source/17/drafts/SGA4_XVII_LINES_4201_4600_English.texfrag
$f^{\ast}\Aa$,\nde{this condition is superfluous. \Cf
\nderef{XVII.ndeKS}.} this isomorphism is still defined for
$K\in\Ob D(V,\Aa)$.

\subsubsection{}\etiquette[5.2.2]{XVII.subsubsec5.2.2}

Let $f:X\rightarrow S$ and $g:S'\rightarrow S$ be morphisms of
schemes, and suppose that $f$ is the composite of an open immersion
$j$ and a proper morphism $p$ whose fibers have dimension bounded by
an integer $n$ (this is automatic when $S$ is quasi-compact). These
morphisms fit into a diagram with cartesian squares
\[
\xymatrix{
X' \ar[rrrr]_{g'} \ar@{^{(}->}[rd]^{j'} \ar[dd]_{f'}& &         &   &
X\ar@{^{(}->}[rd]^{j} \ar[dd]_(.35){f} & \\
             &\bar{X}'\ar[rrrr]|(.75)\hole_{\bar{g}}\ar[ld]^{p'}  &&   &   & \bar{X}\ar[ld]^{p}\\
S' \ar[rrrr]_{g}  & & &&  S &.
}
\]
Let\pageoriginale $\Aa$ be a sheaf of torsion rings on $S$, $\Aa'$ a
sheaf of torsion rings on $S'$, and $L$ a bounded-above complex of
$(\Aa',g^{\ast}\Aa)$-bimodules. We ring $X$ and $\bar X$
(\resp $X'$ and $\bar X'$) by the inverse images $\Aa_X$ and
$\Aa_{\bar X}$ (\resp $\Aa'_{X'}$ and $\Aa'_{\bar X'}$) of $\Aa$
(\resp $\Aa'$) on $X$ and $\bar X$ (\resp $X'$ and $\bar X'$). We
then have base-change isomorphisms
\eqhyperref[XVII.eq5.2.1.2]{5.2.1.2} between the functors
${p'}^{\ast}L\overset{\LL}{\otimes}\bar g^{\ast}j_{!}K$ and
$j'_{!}({f'}^{\ast}L\sisi{\otimes}{\overset{\LL}{\otimes}}
{g'}^{\ast}K)$ from $D^{-}(X,\Aa_X)$ to
$D^{-}(\bar X',\Aa'_{\bar X'})$, and base-change isomorphisms
(\hyperref[XVII.thm4.3.1]{4.3.1}) between the functors
$L\overset{\LL}{\otimes}g^{\ast}Rp_{\ast}K$ and
$Rp'_{\ast}({p'}^{\ast}L\overset{\LL}{\otimes}\bar g^{\ast}K)$ from
$D^{-}(\bar X,\Aa_{\bar X})$ to $D^{-}(S',\Aa')$. By composition,
one obtains a base-change isomorphism between functors from
$D^{-}(X,\Aa_X)$ to $D^{-}(S',\Aa')$:
$$\etiquette[5.2.2.1]{XVII.eq5.2.2.1}
L\overset{\LL}{\otimes}g^{\ast}Rp_{\ast}j_{!}K\isomto
Rp'_{\ast}j'_{!}({f'}^{\ast}L\overset{\LL}{\otimes}{g'}^{\ast}K){\quoi}.\leqno(5.2.2.1)
$$

When $L$ is bounded with components of finite \ttor-dimension, the
preceding arguments and \hyperref[XVII.var4.3.2]{4.3.2}~(i) define the
base-change isomorphism for $K\in\Ob D(X,\Aa_X)$.\nde{by combining the
preceding N.D.E. and \nderef{XVII.ndecb}, this isomorphism can be
defined without any condition on $L$ or the hypothesis that the
dimension of the fibers of $p$ is bounded.}

In order to refer to \hyperref[XVII.thm5.1.8]{5.1.8}, suppose that $S$
is quasi-compact and quasi-separated. The functors
$Rp_{\ast}j_{!}=Rf_{!}$ and $Rp'_{\ast}j'_{!}=Rf'_{!}$ then depend only
on $f$ and $f'$, giving a meaning to the following statement.

\begin{enonce*}{Lemma 5.2.3}\etiquette[5.2.3]{XVII.lem5.2.3}
The base-change isomorphism \eqhyperref[XVII.eq5.2.2.1]{5.2.2.1} does
not depend on the chosen compactification $f=pj$ of $X$.
\end{enonce*}

Let $q:\bar X_1\rightarrow\bar X_2$ be a morphism of compactifications
$$\etiquette[5.2.3.1]{XVII.eq5.2.3.1}
\vcenter{
\xymatrix{
X \ar[rr]^{j_1} \ar[rrd]|(.65)\hole^{j_2} \ar[dd]_{f}& & \bar{X}_1 \ar[d]_{q} \ar[lldd]^{p_1}\\
  & &  \bar{X}_2 \ar[lld]^{p_2}\\
S & & .
}}\leqno{(5.2.3.1)}
$$
The\pageoriginale base-change isomorphism $c_1$ relative to $\bar X_1$
is obtained by ``composing'' those relative to $j_1$ and $p_1$, or,
equivalently, by composing those relative to $j_1$, $q$, and $p_2$.
The isomorphism $c_2$ relative to $\bar X_2$ is obtained by composing
those relative to $j_2$ and $p_2$. To verify that $c_1=c_2$, it
suffices to verify that the isomorphism $c'_1$ relative to $j_2$ is the
composite $c'_2$ of those relative to $j_1$ and $q$. Use a prime to
denote the arrows obtained from the arrows
\eqhyperref[XVII.eq5.2.3.1]{5.2.3.1} by the base change $g$; $c'_1$ and
$c'_2$ are isomorphisms
\[
c'_{1}, c'_{2}:L\overset{\LL}{\otimes}\bar{g}_{2}^*\ast
j_{2!}K\isomto  j'_{2!}({j'}_{2}^*\ast
L\overset{\LL}{\otimes} {g'}^{\ast}K){\quoi}.
\]
By \eqhyperref[XVII.eq5.1.3.2]{5.1.3.2}, to verify that $c'_1=c'_2$ it
suffices to verify it after restricting to $\sisi{X}{\bar X}'_2$,
which is immediate.

The base-change isomorphisms relative to two compactifications between
which a morphism exists are therefore equal, and the conclusion
follows from \hyperref[XVII.prop3.2.6]{3.2.6}~(ii).

\begin{enonce*}{Lemma 5.2.4}\etiquette[5.2.4]{XVII.lem5.2.4}
Let $f=uv:X\rightarrow Y\rightarrow S$ be a composite of
$S$-compactifiable morphisms. The base-change isomorphisms
\eqhyperref[XVII.eq5.2.2.1]{5.2.2.1} relative to $u$, $v$, and $f$
give rise to a commutative diagram (notation as in
\eqhyperref[XVII.eq5.2.4.1]{5.2.4.1} below):
\[
\xymatrix{
L\overset{\LL}{\otimes} g^{\ast} R\sisi{uv}{(uv)}_{!} K
\ar@<-4pt>[rr]^{\sim} \ar@{=}[d]& & R\sisi{u'v'}{(u'v')}_{!}((u'v')^{\ast}L
\overset{\LL}{\otimes}g''^{\ast}K)\ar@{=}[d]\\
L\overset{\LL}{\otimes} g^{\ast}Ru_{!}Rv_{!}K \ar@<-4pt>[r]^(.45){\sim}& Ru'_{!}(u'^{\ast}L
\overset{\LL}{\otimes}g'^{\ast}Rv_{!}\sisi{}{K}) \ar@<-4pt>[r]^(.45){\sim}& Ru'_{!}Rv'_{!}(v'^{\ast}u'^{\ast}L\overset{\LL}{\otimes}g''^{\ast}K).
}
\]
\end{enonce*}

Choose\pageoriginale a compactification of $(u,v)$
(\hyperref[XVII.def3.2.5]{3.2.5}):
$$\etiquette[5.2.4.1]{XVII.eq5.2.4.1}
\vcenter{
\xymatrix{
X' \ar[rr]_{g''} \ar[d]_{\sisi{u}{v}'}& & X \ar@{^{(}->}[r]^{j_1} \ar[d]_{\sisi{u}{v}}& \bar{X}
\ar[ld]_{p_1}\ar@{^{(}->}[r]^{j_2} \ar[d]& \bar{\bar{X}} \ar[ld]^{p_3}& \\
Y' \ar[rr]_{g'} \ar[d]_{\sisi{v}{u}'}& & Y \ar@{^{(}->}[r]^{j_3}\ar[d]_{\sisi{v}{u}}& \bar{Y}\ar[ld]^{p_2} &  &\\
S' \ar[rr]_{g} &  & S &      &  .
}}\leqno{(5.2.4.1)}
$$

Lemma \hyperref[XVII.lem5.2.4]{5.2.4} says that the base-change
isomorphism $c_1$ relative to $f$ is the composite $c_2$ of those
relative to $u$ and $v$. The isomorphism $c_1$ is obtained by composing
the base-change isomorphisms relative to $j_1$, $j_2$, $p_3$, and
$p_2$. For $c_2$, one composes those relative to $j_1$, $p_1$, $j_3$,
and $p_2$. It therefore suffices to compare the isomorphisms $c'_1$,
the composite of those relative to $p_3$ and $j_2$, and $c'_2$, the
composite of those relative to $j_3$ and $p_1$.

Consider the cartesian diagram
\[
\xymatrix{
  \bar{X}' \ar[ddd]^{p'_1} \ar[rrrr]^{\bar{g}''} \ar@{^{(}->}[rrd]^{j'_2} &    &       &    & \bar{X}\ar@{^{(}->}[rrd]^{j_2} \ar[ddd]^{\sisi{}{p_1}}|(.38)\hole\\
                          & &\bar{\bar{X}}' \ar[rrrr]_(.3){\bar{\bar{g}}} \ar[ddd]^{p'_3}&   &&
& \bar{\bar{X}} \ar[ddd]^{\sisi{}{p_3}}\\
 & & &\\
  Y'  \ar@{^{(}->}[rrd]_{j'_3}\ar[rrrr]^(.7){g'}|(.52)\hole \ar[rdd]_{\sisi{v}{u}'} &    &&
& Y\ar[rdd]_(.65){\sisi{v}{u}} \ar@{^{(}->}[rrd]^{j_3}\\
  & &\bar{Y}' \ar[rrrr]_{\bar{g}'} \ar[ld]^{p'_2}&        &  & &\bar{Y}\ar[ld]^{p_2}\\
 &S'  \ar[rrrr]_{g} &&& & S &,
}
\]
$L'={p'_2}^{\ast}L$\sisi{}{, $K'=j_{1!}K$}\pageoriginale and
\[
c'_{1}, c'_{2}: L'\overset{\LL}{\otimes} {\bar{g}'}^{\ast}
R(j_{3}p_{1})_{!} K\sisi{}{'}\isomto  R(j'_{3} p'_{1})_{!} ((j'_{3}
p'_{1})^{\ast} L'\otimes {\bar{g}''}^{\ast}K\sisi{}{'}){\quoi}.
\]
By \eqhyperref[XVII.eq5.1.3.2]{5.1.3.2}, to verify that $c'_1=c'_2$ it
suffices to verify that ${j'_3}^{\ast}(c'_1)={j'_3}^{\ast}(c'_2)$,
which is immediate.

\begin{enonce*}{Lemma 5.2.5}\etiquette[5.2.5]{XVII.lem5.2.5}
Let there be morphisms of schemes ringed by sheaves of torsion rings
\[
(S'', \Aa'')\xrightarrow{g_{2}} (S',
\Aa')\xrightarrow{g_{1}} (S,\Aa){\quoi},
\]
and let $f:X\rightarrow S$ be a compactifiable morphism, giving a
diagram:
\[
\xymatrix{
X'' \ar[r]_{g'_2} \ar[d]^{f''}& X' \ar[r]_{g'_1}\ar[d]^{f'}& X \ar[d]^{f}\\
S'' \ar[r]_{g_2}& S' \ar[r]_{g_1}& S \rlap{.}
}
\]
Ring\pageoriginale $X$, $X'$, and $X''$ by $f^{\ast}\Aa$,
${f'}^{\ast}\Aa'$, and ${f''}^{\ast}\Aa''$, respectively. Then the
diagram
\[
\xymatrix{
L(g_1g_2)^{\ast}Rf_{!}  \ar[rr]^{\sim} \ar@{=}[d]&& Rf''_{!}L(g'_1g'_2)^{\ast}\ar@{=}[d]\\
Lg_{\sisi{1}{2}}^{\ast}Lg^{\ast}_{\sisi{2}{1}} R f_{!} \ar[r]^{\sim}& Lg^{\ast}_{\sisi{1}{2}} Rf'_{!}
L\sisi{g^\ast_2}{g'^\ast_1}  \ar[r]^{\sim} & Rf''_{!} Lg'^{\ast}_{\sisi{1}{2}}  L g'^{\ast}_{\sisi{2}{1}}
}
\]
is commutative.
\end{enonce*}

This lemma follows immediately from the analogous results for the
functors $j_{!}$ and $Rp_{\ast}$ with $p$ proper. We leave it to the
reader to formulate the precise version with modules.

The preceding statements are summarized in the following theorem.

\begin{enonce*}{Theorem 5.2.6}\etiquette[5.2.6]{XVII.thm5.2.6}
The functor $Rf_{!}$ commutes with base change.
\end{enonce*}

\begin{enonce*}[remark]{Variants 5.2.7}\etiquette[5.2.7]{XVII.var5.2.7}
\begin{enumerate}
\renewcommand{\theenumi}{\roman{enumi}}
\renewcommand{\labelenumi}{(\theenumi)}
\item The functors $f_{!}$ (for sheaves of pointed sets:
\hyperref[XVII.var5.1.11]{5.1.11}) and the functors $R^1f_{!}$ (for
sheaves of groups: \hyperref[XVII.var5.1.12]{5.1.12}) commute with
base change.

\item The functors $Rf_{!}$ (\hyperref[XVII.var5.1.14]{5.1.14})
commute with base change.

\item Under the hypotheses of (\hyperref[XVII.subsubsec5.2.2]{5.2.2}),
if $L$ is bounded and has components of finite \ttor-dimension,\nde{as
observed in \hyperref[XVII.subsubsec5.2.2]{5.2.2}, these hypotheses are
superfluous.} the functors $Rf_{!}:D(X)\rightarrow D(S)$ commute with
base change.
\end{enumerate}
\end{enonce*}

The\pageoriginale conjunction of the following two Propositions
\hyperref[XVII.prop5.2.8]{5.2.8} and
\hyperref[XVII.prop5.2.9]{5.2.9} is essentially equivalent to Theorem
\hyperref[XVII.thm5.2.6]{5.2.6}.

\begin{enonce*}{Proposition 5.2.8}\etiquette[5.2.8]{XVII.prop5.2.8}
Let $f:X\rightarrow S$ be a compactifiable morphism and
$\mathcal F$ a torsion sheaf on $X$. For every geometric point $s$ of
$S$, let $X_s$ be the geometric fiber of $X$ at $s$ and let
$\mathcal F_s$ be the induced sheaf. Then canonically
\[
(R^{q}f_{!}(\mathcal{F}))_s=H^{q}_c(X_s,\mathcal{F}_s){\quoi}.
\]
\end{enonce*}

\begin{enonce*}{Corollary 5.2.8.1}\etiquette[5.2.8.1]{XVII.cor5.2.8.1}
If the fibers of $f$ have dimension $\leq d$, the functor $Rf_{!}$ has
cohomological dimension $\leq2d$
(\hyperref[XVII.def1.2.9]{1.2.9}): for every torsion sheaf
$\sisi{F}{ℱ}$, one has $R^if_{!}\sisi{F}{ℱ}=0$ for $i>2d$. The functor
$R^{2d}f_{!}$ is therefore right exact on the category of torsion
sheaves.
\end{enonce*}

By base change (\hyperref[XVII.prop5.2.8]{5.2.8}), we reduce to the
case where $S$ is the spectrum of an algebraically closed field. If
$\bar X$ is then a proper scheme over $S$ containing $X$ as a dense
Zariski open subscheme, with $j:X\rightarrow\bar X$, then
$\dim(\bar X)=\dim(X)\leq\sisi{D}{d}$. Hence, for $ℱ$ a torsion sheaf
on $X$ and $i>2d$, by \hyperrefext[X][X.cor4.3]{4.3},
\refnoexpose{X.cor5.2}
\[
H^{i}_c(X,ℱ)=H^{i}(\bar{X},j_{!}\sisi{F}{ℱ})=0{\quoi}.
\]

\begin{enonce*}{Proposition 5.2.9}\etiquette[5.2.9]{XVII.prop5.2.9}
Let $f:X\rightarrow S$ be a compactifiable morphism, let $\Aa$ be a
sheaf of torsion rings on $S$, let $K$ be a bounded-above complex of
left $f^{\ast}\Aa$-modules, and let $L$ be a bounded-above complex of
right $\Aa$-modules. Then canonically
\[
L\overset{\LL}{\underset{\Aa}{\otimes}}
Rf_{!}(K)=Rf_{!}(f^{\ast}L\overset{\LL}{\underset{f^{\ast}\Aa}{\otimes}}K){\quoi}.
\]
\end{enonce*}

This\pageoriginale proposition formally implies the following result:

\begin{enonce*}{Theorem 5.2.10}\etiquette[5.2.10]{XVII.thm5.2.10}
Let $f:X\rightarrow S$ be a compactifiable morphism, let $\Aa$ be a
sheaf of torsion rings on $S$, let $K\in D^{-}(X,f^{\ast}\Aa)$, and
let $d\in\ZZ$. If $K$ has finite \tTor-dimension $\leq d$
(\hyperref[XVII.def4.1.9]{4.1.9}), then $Rf_{!}K$ has finite
\tTor-dimension $\leq d$.
\end{enonce*}

It suffices to verify that, for every sheaf $L$ of right $\Aa$-modules
on $S$,
\[
\sisi{H}{\Hh}^{-i}(L\overset{\LL}{\otimes}Rf_{!}K)=0\quad\mbox{\rm
for}\quad i>d{\quoi}.
\]
By \hyperref[XVII.prop5.2.9]{5.2.9}, we know that
\[
\sisi{H}{\Hh}^{-i}(L\overset{\sisi{L}{\LL}}{\otimes}Rf_{!}K)=R^{-i}f_{!}(f^{\ast}L\overset{\sisi{L}{\LL}}{\otimes}K){\quoi},
\]
and, by hypothesis,
\[
\sisi{H}{\Hh}^{-i}(f^{\ast}L\overset{\sisi{L}{\LL}}{\otimes}K)=0\quad\mbox{\rm for}\quad
i>d{\quoi}.
\]
The conclusion follows by observing that if a complex $M$ of sheaves
of modules on $X$ satisfies $\sisi{H}{\Hh}^{-i}(M)=0$ for $i>d$, then
one still has $\sisi{H}{\Hh}^{-i}(Rf_{!}M)=0$ for $i>d$.

For constant noetherian sheaves of rings, one has the following more
general result:

\begin{enonce*}{Theorem 5.2.11}\etiquette[5.2.11]{XVII.thm5.2.11}
Let $f:X\rightarrow S$ be a morphism of sites, let $A$ be a ring, let
$C$ be its center, and let $d\in\ZZ$. Suppose that
\begin{enumerate}
\renewcommand{\theenumi}{\roman{enumi}}
\renewcommand{\labelenumi}{(\theenumi)}
\item the functor $Rf_{\ast}:D^{+}(X,C)\rightarrow D^{+}(S,C)$ has
finite cohomological dimension;

\item $A$ is right noetherian;

\item $S$ has enough points.
\end{enumerate}
For\pageoriginale every complex $K\in\Ob D^{-}(X,A)$ of sheaves of
left $A$-modules, if $K$ has finite \tTor-dimension $\leq d$, then the
complex $Rf_{\ast}K\in\Ob D^{-}(S,A)$ has \tTor-dimension $\leq d$.
\end{enonce*}

\begin{enonce*}[remark]{Proof}
By (i), the functor $f_{\ast}$ is derivable to functors
\[
Rf_{\ast}:D(X,A)\longrightarrow D(S,A)
\]
and
\[
Rf_{\ast}:D(X,C)\longrightarrow D(S,C){\quoi}.
\]
\end{enonce*}

By \hyperref[XVII.lem4.2.8]{4.2.8}, it suffices to verify that, for
every point $s$ of $S$, the complex $(Rf_{\ast}K)_s$ of $A$-modules
has \tTor-dimension $\leq d$; for this it suffices that, for every
finitely generated right $A$-module $L$,
\[
H^{-i}(L\overset{L}{\otimes}Rf_{\ast}K)=0\quad\mbox{\rm for}\quad
i>d{\quoi}.
\]

Reasoning as in \hyperref[XVII.thm5.2.10]{5.2.10}, we see that it
suffices for the ``base-change arrow''
$$\etiquette[5.2.11.1]{XVII.eq5.2.11.1}
L\overset{L}{\otimes}Rf_{\ast}K\longrightarrow Rf_{\ast}(L\overset{L}{\otimes}K)\leqno(5.2.11.1)
$$
to be an isomorphism.

Let $L^{\ast}$ be a resolution of $L$ by finitely generated free
$A$-modules. Since the functors under consideration are way-out to the
right, it suffices to prove that
\eqhyperref[XVII.eq5.2.11.1]{5.2.11.1} is an isomorphism after replacing
$L$ by any component of $L\sisi{_{\ast}}{^\ast}$, which reduces us to
the immediate case where $L$ is finitely generated and free.

\begin{enonce*}[remark]{Remarks 5.2.12}\etiquette[5.2.12]{XVII.rems5.2.12}
\begin{enumerate}
\renewcommand{\theenumi}{\roman{enumi}}
\renewcommand{\labelenumi}{(\theenumi)}
\item The\pageoriginale hypothesis (iii) of
\hyperref[XVII.thm5.2.11]{5.2.11} is probably unnecessary.

\item If the noetherian hypothesis on $A$ is omitted,
\hyperref[XVII.thm5.2.11]{5.2.11} remains valid when the functors
$R^if_{\ast}$ commute with filtered direct limits.

\item The hypothesis that $S$ is ringed by a \emph{constant} sheaf of
rings is essential. A counterexample is obtained by taking for $S$ the
spectrum of a strictly local discrete valuation ring, for $f$ the
inclusion of the generic point $\eta$, for the sheaf of rings $\Aa$ the
sheaf equal to $\ZZ/n\ZZ$ at the generic point and to $\ZZ/n^2\ZZ$ at
the closed point $s$ ($n>1$ prime to the residue characteristics), and
for the sheaf the module $f^{\ast}\Aa$.
\end{enumerate}
\end{enonce*}

\subsection{Comparison theorem and finiteness theorem}\etiquette[5.3]{XVII.subsec5.3}

\subsubsection{}\etiquette[5.3.1]{XVII.subsubsec5.3.1}
To a scheme $X$ of finite type over $\CC$ are associated its étale
site $X_{\text{\rm et}}$, its transcendental site $X_{\cl}$, and a
morphism of sites (\hyperrefext[XI][XI.sec4]{4})
\[
\epsilon:X_{\cl}\longrightarrow X_{\et}{\quoi}.
\]

If $j:U\rightarrow X$ is an open immersion, the commutative diagram
$$
\xymatrix{
U_{\cl}  \ar[r]\ar[d]   & U_{\et} \ar[d]\\
X_{\cl}  \ar[r] & X_{\et}
}
$$
identifies $U^{\sim}_{\cl}$ with the open subtopos of
$X^{\sim}_{\cl}$ that is the inverse image of the open subtopos
$U_{\widetilde{\text{\rm et}}}$ of
$X_{\widetilde{\text{\rm et}}}$.

If $p:Y\rightarrow X$ is a proper morphism, the continuous map
$p^{\an}:Y^{\an}\rightarrow X^{\an}$ is proper and separated. To see
this, Chow's lemma reduces us to the case where $p$ is projective.

Let\pageoriginale $f:X\rightarrow S$ be a compactifiable morphism of
schemes of finite type over $\CC$, and let $f=pj$ be a compactification
of $f$.
$$\etiquette[5.3.1.1]{XVII.eq5.3.1.1}
\vcenter{
\xymatrix{
X_{\cl} \ar[rrr]_{\epsilon} \ar[rd]^{j_{\cl}} \ar[dd]_{f_{\cl}}&  & & X \ar[rd]^{j}\ar[dd]_(.6){f}|\hole & \\
            &\bar{X}_{\cl}\ar[ld]^{p_{\cl}}  \ar[rrr]_{\epsilon} &  & &\bar{X}\ar[ld]^{p} \\
S_{\cl}  \ar[rrr]_{\epsilon} &                 &   &  S & .
}}\leqno{(5.3.1.1)}
$$
The functor $p_{\cl\ast}j_{\cl !}$ is precisely the direct-image
functor with proper supports $f_{\cl !}$ (see \textsc{Verdier}
\cite{sga4.exp.XVII.bib.V3}). If $S_{\cl}$, and hence $X_{\cl}$, is separated, the
functor $j_{\cl !}$ carries soft sheaves on $X_{\cl}$ to soft sheaves
on $\bar X_{\cl}$ \cite[II 3.5.5]{sga4.exp.XVII.bib.TF}. Since the question is local on
$S$, it in any case carries flasque sheaves on $X_{\cl}$ to sheaves on
$\sisi{X}{\bar X}_{\cl}$ that are acyclic for the functor
$p_{\cl\ast}$. Hence (\cf \textsc{Verdier} \cite{sga4.exp.XVII.bib.V3})
$$
Rf_{\cl !}=Rp_{\cl\ast}j_{\cl !}{\quoi}.
$$

\subsubsection{}\etiquette[5.3.2]{XVII.subsubsec5.3.2}

% END INCLUDED SGA4_English_current_buildable_source_20260906/source/17/drafts/SGA4_XVII_LINES_4201_4600_English.texfrag

% BEGIN INCLUDED SGA4_English_current_buildable_source_20260906/source/17/drafts/SGA4_XVII_LINES_4601_5000_English.texfrag

By \hyperref[XVII.lem5.1.2]{5.1.2}, there is a base-change
isomorphism between functors from $D(X)$ to $D(\bar X_{\cl})$:
$$
\epsilon^{\ast}j_{!}\isomto j_{\cl !}\epsilon^{\ast}{\quoi}.
$$

By\pageoriginale \hyperref[XVII.subsubsec4.2.12]{4.2.12} or
\hyperref[XVII.prop4.1.5]{4.1.5}, there is a base-change morphism
between functors from $D^{+}(\bar X)$ to $D^{+}(S_{\cl})$:
$$
\epsilon^{\ast}Rp_{\ast}\longrightarrow Rp_{\cl
\ast}\epsilon^{\ast}{\quoi}.
$$
By composition one obtains a \emph{base-change morphism between
functors from} $D^{+}_{\tors}(X)$ to $D^{+}_{\tors}(S_{\cl})$
(\hyperref[XVII.def5.1.13]{5.1.13}):
$$\etiquette[5.3.2.1]{XVII.eq5.3.2.1}
\epsilon^{\ast}Rf_{!}\longrightarrow Rf_{\cl
!}\epsilon^{\ast}{\quoi}.\leqno(5.3.2.1)
$$

As in \hyperref[XVII.lem5.2.3]{5.2.3}, one checks that it does not
depend on the chosen compactification of $f$; it satisfies a
transitivity formula of the type in \hyperref[XVII.lem5.2.4]{5.2.4}.

\begin{enonce*}{Theorem 5.3.3}\etiquette[5.3.3]{XVII.thm5.3.3}
\emph{(Comparison theorem).} The morphism
\eqhyperref[XVII.eq5.3.2.1]{5.3.2.1} is an isomorphism.
\end{enonce*}

By construction, it suffices to treat the case where $f$ is proper.
The functors under consideration are triangulated and ``way-out''
\cite[I 7]{sga4.exp.XVII.bib.RD}. It therefore suffices to treat the case where $K$ is
a torsion sheaf placed in degree $0$, and the result follows from
\hyperrefext[XVI][XVI.thm4.1]{4.1}~(i).

\begin{enonce*}[remark]{Variants 5.3.4}\etiquette[5.3.4]{XVII.var5.3.4}
\begin{enumerate}
\renewcommand{\theenumi}{\roman{enumi}}
\renewcommand{\labelenumi}{(\theenumi)}
\item If $\Aa$ is a sheaf of torsion rings on $S_{\text{\rm et}}$,
one similarly obtains an isomorphism
\eqhyperref[XVII.eq5.3.2.1]{5.3.2.1} between functors from
$D(X,f^{\ast}\Aa)$ to $D(S,\epsilon^{\ast}\Aa)$.

\item In low degrees one obtains the usual variants for sheaves of
pointed sets and sheaves of \iInd-finite groups.
\end{enumerate}
\end{enonce*}

\begin{enonce*}{Corollary 5.3.5}\etiquette[5.3.5]{XVII.cor5.3.5}
Let\pageoriginale $X$ be a separated scheme of finite type over $\CC$
which admits a compactification in a proper scheme over $\CC$, and let
$F$ be a torsion sheaf on $X$. Then
$$
H^{q}_{!}(X,F)\isomto H^{q}_c(X_{\cl},\epsilon^{\ast}F){\quoi}.
$$
\end{enonce*}

\begin{enonce*}{Theorem 5.3.6}\etiquette[5.3.6]{XVII.thm5.3.6}
Let $f:X\rightarrow S$ be a compactifiable morphism of finite
presentation and let $A$ be a left noetherian torsion ring. The
functor $Rf_{!}:D(X,A)\rightarrow D(S,A)$ carries complexes with
constructible cohomology to complexes with constructible cohomology.
\end{enonce*}

The assertion is local on $S$, which may be assumed affine. By passage
to the limit, we reduce to the case where $S$ is noetherian. Complexes
with constructible cohomology form a triangulated subcategory of
$D(S,A)$. Since the functor $Rf_{!}$ is ``way-out,'' we reduce to
verifying that if $\mathcal F$ is a constructible sheaf of $A$-modules,
then the sheaves $R^qf_{!}\mathcal F$ are constructible. Write $f=pj$,
with $p$ proper and $j$ an open immersion. Then
$R^qf_{!}\mathcal F=R^qp_{\ast}(j_{!}\mathcal F)$. The sheaf
$j_{!}\mathcal F$ is constructible, and the conclusion follows from
\hyperrefext[XIV][XIV.thm1.1]{1.1}.

\setcounter{subsubsection}{6}
\subsubsection{}\etiquette[5.3.7]{XVII.subsubsec5.3.7}
In low degrees, there are analogous statements for direct images with
proper supports of a constructible sheaf of \iInd-finite sets or
groups.

\begin{enonce*}{Corollary 5.3.8}\etiquette[5.3.8]{XVII.cor5.3.8}
Let $X$ be a separated scheme of finite type over a separably closed
field $k$, and let $\Ff$ be a constructible torsion sheaf on $X$.
Suppose that $X$ can be compactified in a proper scheme over $k$.\nde{this
hypothesis is automatically satisfied. See \nderef{XVII.ndeConrad}.}
Then the groups $H^q_c(X,\mathcal F)$ are finite.
\end{enonce*}

\subsection{Künneth formula}\etiquette[5.4]{XVII.subsec5.4}\pageoriginale

\subsubsection{}\etiquette[5.4.1]{XVII.subsubsec5.4.1}
Let there be a commutative diagram (\hyperref[XVII.subsec0.11]{0.11})
of sites
$$\etiquette[5.4.1.1]{XVII.eq5.4.1.1}
\vcenter{
\xymatrix@R=.5cm@C=1cm{
             &  Z'\ar[ld]_{p'}\ar[rd]^{q'} \ar[dd]^{h}&    \\
X' \ar[dd]_{f}&                   &  Y' \ar[dd]^{g}\\
             &       Z\ar[ld]_{p} \ar[rd]^{q}   &    \\
X \ar[rd]           &                    & Y \ar[ld]\\
             &      S             & &.
}}\leqno{(5.4.1.1)}
$$
Suppose $S$ is ringed by a sheaf of rings $\Aa$ with center $\Cc$; we
again denote by $\Aa$ and $\Cc$ the inverse images of these sheaves on
the sites of \eqhyperref[XVII.eq5.4.1.1]{5.4.1.1}.

If $K$ is a sheaf of right $\Aa$-modules on $X$ and $L$ a sheaf of
left $\Aa$-modules on $Y$, set
\[
K\mathop{\boxtimes}_{\Aa}L=p^{\ast}\sisi{}{K}
\underset{\Aa}{\otimes}q^{\ast}L
\]
(external tensor product). Similarly for $K$ (\resp $L$) on $X'$
(\resp $Y'$). Now let $K$ be a sheaf of right $\Aa$-modules on $X'$
and $L$ a sheaf of left $\Aa$-modules on $Y'$. The base-change arrows
(\sisi{}{\hyperref[XVII.prop2.1.3]{2.1.3}, }
\sisi{\hyperrefext[XII][XII.eq4.1]{4.1}}{\neweqhyperrefext[XII][XII.eq4.1]{4.1}})
define
$$\etiquette[5.4.1.2]{XVII.eq5.4.1.2}
\begin{cases}
p^{\ast}f_{\ast}K\longrightarrow h_{\ast}{p'}^{\ast}K\\
q^{\ast}g_{\ast}L\longrightarrow h_{\ast}{q'}^{\ast}L{\quoi},
\end{cases}\leqno(5.4.1.2)
$$
and hence,\pageoriginale by taking their tensor product, an arrow
$$\etiquette[5.4.1.3]{XVII.eq5.4.1.3}
f_{\ast}K\boxtimes_{\Aa} g_{\ast}L\longrightarrow
h_{\ast}{p'}^{\ast}K\otimes_{\Aa} h_{\ast}
{q'}^{\ast}L\longrightarrow h_{\ast}(K\boxtimes_{\Aa}L){\quoi}.\leqno(5.4.1.3)
$$

Suppose that the functors
\begin{align*}
Rf_{\ast} &:D^{+}(X',\Aa^{\circ})\longrightarrow D^{+}(X,\Aa^{\circ})\\
Rg_{\ast} &:D^{+}(Y',\Aa)\longrightarrow D^{+}(Y,\Aa)\\
Rh_{\ast} &:D^{+}(Z',\Cc)\longrightarrow D^{+}(Z,\Cc)
\end{align*}
have finite cohomological dimension and therefore extend to the
unbounded derived categories. For $K\in\Ob D^{-}(X',\Aa^{\circ})$ and
$L\in\Ob D^{-}(Y',\Aa)$, one is then led to define a
\emph{Künneth arrow}
$$\etiquette[5.4.1.4]{XVII.eq5.4.1.4}
Rf_{\ast}K\overset{\LL}{\boxtimes}_{\Aa}
Rg_{\ast}L\longrightarrow Rh_{\ast}(K
\overset{\LL}{\boxtimes}_{\Aa}L){\quoi}.\leqno(5.4.1.4)
$$

We define it only when the sites under consideration have enough
points, so that \hyperref[XVII.prop4.2.10]{4.2.10} and
\hyperref[XVII.lem4.2.8]{4.2.8} may be used. By
\hyperref[XVII.prop4.2.10]{4.2.10}, it suffices to define it for a pair
$(K,L)$ left adapted to the functor
$(K,L)\mapsto K\boxtimes_{\Aa}L$, where $K$ and $L$ are bounded above
and right adapted to the functors $f_{\ast}$ and $g_{\ast}$.

Let $K_1$ be a flat resolution of $f_{\ast}K$, let $L_1$ be a flat
resolution of $g_{\ast}L$, and let
$u:K\overset{\LL}{\boxtimes}\sisi{}{_\Aa}L\rightarrow M$ be a
resolution of $K\overset{\LL}{\boxtimes}\sisi{}{_\Aa}L$ adapted to
the functor $Rh_{\ast}$. We have morphisms
\begin{gather*}
Rf_{\ast}K\overset{\LL}{\boxtimes}_{\Aa}Rg_{\ast}L\simeq
K_1\boxtimes_{\Aa}L_1\longrightarrow
f_{\ast}K\boxtimes_{\Aa}g_{\ast}L\longrightarrow\\
\longrightarrow h_{\ast}(K\boxtimes_{\Aa}L)\longrightarrow
h_{\ast}M\simeq Rh_{\ast}(K\overset{\LL}{\boxtimes}_{\Aa}L)
\end{gather*}
and define \eqhyperref[XVII.eq5.4.1.4]{5.4.1.4} to be their composite.

This\pageoriginale arrow has the following characteristic property,
whose proof is entirely analogous to that of
\hyperref[XVII.prop4.2.13]{4.2.13}:

\setcounter{paragraph}{4}
\paragraph{}\etiquette[5.4.1.5]{XVII.par5.4.1.5}
Let $K\in\Ob C^{-}(X',\Aa^{\circ})$ and
$L\in\Ob C^{-}(Y',\Aa)$. Set $K_1=f_{\ast}K$,
$L_1=g_{\ast}L$, and $M_2=K\boxtimes\sisi{}{_\Aa}L$. The following
diagram is then commutative:
$$
\xymatrix@R=1.25cm{
Rf_{\ast}K
\overset{\LL}{\boxtimes}_{\Aa}Rg_{\ast}
L \ar[rrr]& & & Rh_{\ast}(K
\overset{\LL}{\boxtimes}_{\Aa}L) \ar[d]\\
K_1\overset{\LL}{\boxtimes}_{\Aa}L_{\sisi{k}{1}}
\ar[r] \ar[u]&
K_1\boxtimes\sisi{}{_\Aa}L_1 \ar[r]& h_{\ast}M_2 \ar[r]& Rh_{\ast}M_2.
}
$$

There is an analogous commutativity (\cf
\hyperref[XVII.prop4.2.13]{4.2.13}) starting instead from morphisms
$K_1\rightarrow f_{\ast}K$, $L_1\rightarrow g_{\ast}L$, and
$K\boxtimes\sisi{}{_\Aa}L\rightarrow M_2$.

The commutativity in \hyperref[XVII.par5.4.1.5]{5.4.1.5} makes it, in
principle, easy to verify the compatibility properties of the Künneth
morphism.

\subsubsection{}\etiquette[5.4.2]{XVII.subsubsec5.4.2}
Let \eqhyperref[XVII.eq5.4.1.1]{5.4.1.1} be a commutative diagram of
schemes, and let $\Aa$ be a torsion sheaf on $S$. Suppose that $f$ and
$g$ are compactifiable (\hyperref[XVII.def3.2.1]{3.2.1}), that
$Z=\fprod{X}{Y}{S}$, and that $Z'=\fprod{X'}{Y'}{S}$. Factor $f$ and
$g$ as a proper morphism and an open immersion; this gives a diagram
$$
\xymatrix@C=1.75cm@R=1cm{
  X'      \sisi{\ar[r]^{p'}}{} \ar[d]^{j_1} \ar@/_1.75pc/[dd]_{f}  & Z'   \ar[d]^{j_3}\sisi{}{\ar[l]_{p'}\ar[r]^{q'}}
& \sisi{Y'\ar[l]_{q'}}{} \ar[d]^{j_2} \ar@/^1.75pc/[dd]^{g}\\
\bar{X}' \sisi{\ar[r]^{\bar{p}}}{} \ar[d]^{\bar{f}}  & \bar{Z}' \ar[d]^{\bar{h}}\sisi{}{\ar[l]_{\bar{p}}\ar[r]^{\bar{q}}} & \bar{Y}' \sisi{\ar[l]_{\bar{q}}}{}\ar[d]^{\bar{g}}\\
X            &    Z \ar[l]_{p} \ar[r]^{q}  & Y
}
$$
in\pageoriginale which $\bar Z'=\fprod{\bar X'}{\bar Y'}{S}$. For
complexes $K$ and $L$ on $X'$ and $Y'$, there is an isomorphism of an
immediate nature
$$\etiquette[5.4.2.1]{XVII.eq5.4.2.1}
j_{3!}(K\boxtimes_{\Aa}L)\isomto
j_{1!}K\boxtimes_{\Aa}j_{2!}L{\quoi}.\leqno(5.4.2.1)
$$

By ``composing'' the inverse of this isomorphism with the Künneth arrow
\eqhyperref[XVII.eq5.4.1.4]{5.4.1.4} relative to
$(\bar f,\bar g,\bar h,j_{1!}K,j_{2!}L)$, one obtains, \emph{for
$K\in\Ob D^{-}(X',\Aa)$ and $L\in\Ob D^{-}(Y',\Aa)$, a Künneth arrow}
$$\etiquette[5.4.2.2]{XVII.eq5.4.2.2}
Rf_{!}K\overset{\LL}{\boxtimes}\sisi{}{_\Aa}Rg_{!}L\longrightarrow
Rh_{!}(K\overset{\LL}{\boxtimes}\sisi{}{_\Aa}L){\quoi},\leqno(5.4.2.2)
$$
which one checks is independent of the chosen compactifications.

\begin{enonce*}{Theorem 5.4.3}\etiquette[5.4.3]{XVII.thm5.4.3}
Under the hypotheses of \hyperref[XVII.subsubsec5.4.2]{5.4.2}, the
Künneth arrow for cohomology with proper supports
\eqhyperref[XVII.eq5.4.2.2]{5.4.2.2} is an isomorphism.
\end{enonce*}

\begin{enonce*}[remark]{Proof}
\begin{enumerate}
\renewcommand{\theenumi}{\alph{enumi}}
\renewcommand{\labelenumi}{(\theenumi)}
\item By construction, it suffices to verify the result when $f$ and
$g$ (and hence also $h$) are proper.

\item Let there be a commutative diagram
$$\etiquette[5.4.3.1]{XVII.eq5.4.3.1}
\vcenter{
\xymatrix@C=1.75cm@R=1.25cm{
X_1 \ar[r]^{p_1} \ar[d]& S_1 \sisi{\ar[r]^{q_1}}{}\ar[d]& Y_1\ar[d]\sisi{}{\ar[l]_{q_1}}\\
X   \ar[r]^{p}& S  \sisi{\ar[r]^{q}}{} &  Y\sisi{}{\ar[l]_q}.
}}\leqno{(5.4.3.1)}
$$
and let $Z_1=\fprod{X_1}{Y_1}{S_1}$,
$X'_1=\fprod{X'}{X_1}{X}$, $Y'_1=\fprod{Y'}{Y_1}{Y}$, and
$Z'_1=\fprod{X'_1}{Y'_1}{S_1}=\fprod{Z'}{Z\sisi{_{1}}{}}{Z_1}$, so
that there is a morphism of diagrams
\sisi{(\hyperref[XVII.subsubsec5.4.1]{5.4.1})}{\eqhyperref[XVII.eq5.4.1.1]{5.4.1.1}},
\[
(S_1,X_1,Y_1,Z_1,X'_1,Y'_1,Z'_1)\longrightarrow
(S,X,Y,Z,X',Y',Z')
\]
all of whose arrows will be denoted uniformly by $C$.

We\pageoriginale shall admit that:

\begin{enonce*}{Lemma 5.4.3.2}\etiquette[5.4.3.2]{XVII.lem5.4.3.2}
The following diagram is commutative:
$$
\xymatrix{
C^{\ast}(Rf_{!} K \overset{L}{\boxtimes}_{\Aa}
Rg_{!} L) \ar@{=}[r] \ar[d]& C^{\ast}Rf_{!}K \overset{L}{\boxtimes}_{\Aa}
C^{\ast} Rg_{!} L \ar[r]^{\sim}&  Rf_{1!} C^{\ast}K
\overset{L}{\boxtimes}_{\Aa} Rg_{1!} C^{\ast}
L \ar[d]\\
C^{\ast}(Rh_{!}(K \overset{L}{\boxtimes}_{\Aa}
L)) \ar[r]^{\sim}& Rh_{\sisi{}{1}!}(C^{\ast}(K
 \overset{L}{\boxtimes}_{\Aa}  L))\ar@{=}[r] &Rh_{\sisi{}{1}!}(C^{\ast}K
 \overset{L}{\boxtimes}_{\Aa}  C^*L ).
}
$$
\end{enonce*}

In this diagram, the horizontal arrows are base-change arrows and the
vertical arrows are the Künneth arrows for the two diagrams
\eqhyperref[XVII.eq5.4.1.1]{5.4.1.1} under consideration.

To prove \hyperref[XVII.thm5.4.3]{5.4.3}, it suffices to prove that the
Künneth arrow is an isomorphism at every geometric point of $Z$. By
\hyperref[XVII.lem5.4.3.2]{5.4.3.2}, it suffices to verify this after
any base change \eqhyperref[XVII.eq5.4.3.1]{5.4.3.1} such that $S_1$
is the spectrum of an algebraically closed field and $p_1,q_1$ are
isomorphisms.

After such a base change, diagram
\eqhyperref[XVII.eq5.4.1.1]{5.4.1.1} reduces to a diagram
$$\vcenter{\etiquette[5.4.3.3]{XVII.eq5.4.3.3}
\xymatrix@R=1.5cm@C=2cm{
    & \fprod{X}{Y}{S} \ar[ld]_{p} \ar[rd]^{q} \ar[dd]^{f\times g}&     \\
X \ar[rd]_{f}  &                   &    Y \ar[ld]^{g}\\
    &      S             &\sisi{.}{}
}}\leqno{(5.4.3.3)}
$$
and the Künneth arrow reduces to an arrow
\eqhyperref[XVII.eq5.4.4.1]{5.4.4.1}
$$\etiquette[5.4.3.4]{XVII.eq5.4.3.4}
k:Rf_{!}K\overset{\LL}{\otimes}\sisi{}{_\Aa} Rg_{!} L\longrightarrow
R(f\times g)_{!}(K\overset{\LL}{\boxtimes}\sisi{}{_\Aa}L){\quoi}.\leqno(5.4.3.4)
$$

\item Let\pageoriginale \eqhyperref[XVII.eq5.4.3.3]{5.4.3.3} be a
diagram in which $f$ and $g$ are proper. We shall admit that:

\begin{enonce*}{Lemma 5.4.3.5}\etiquette[5.4.3.5]{XVII.lem5.4.3.5}
The following diagram is commutative:
$$
\xymatrix@C=1.75cm{
Rf_{!}K \overset{\sisi{L}{\LL}}{\underset{\Aa}{\otimes}} Rg_{!}L
\ar[r] \ar[d]& R(f\times g)_{!}(K
\overset{\sisi{L}{\LL}}{\underset{\Aa}{\otimes}} L) \ar@{=}[d]\\
Rf_{!}(K \overset{\sisi{L}{\LL}}{\sisi{\boxtimes}{\otimes}}_{\Aa}
f^{\ast}Rg_{!}L) \ar[r]& Rf_{!} Rp_{!}(p^{\ast}K \overset{\sisi{}{\LL}}{\otimes}\sisi{}{_\Aa} q^{\ast}L).
}
$$
\end{enonce*}

In this diagram, the horizontal arrows are \sisi{the arrows}{the arrow}
of Künneth and the base-change arrow. The vertical arrows are a
base-change arrow by $f$ and a transitivity isomorphism. The theorem
\sisi{\hyperref[XVII.subsubsec5.5.9]{5.5.9}}{\hyperref[XVII.thm5.4.3]{5.4.3}}
follows from this lemma and \hyperref[XVII.thm4.3.1]{4.3.1}.

The diligent reader may verify that this lemma remains valid if $f$
and $g$ are not necessarily proper.
\end{enumerate}
\end{enonce*}

\setcounter{subsubsection}{3}
\subsubsection{}\etiquette[5.4.4]{XVII.subsubsec5.4.4}
Let $S$ be a scheme, let $f_i:X'_i\rightarrow X_i$ $(i\in I)$ be a
finite family of compactifiable morphisms of $S$-schemes, let $X$ be
the product over $S$ of the $X_i$, let $X'$ be the product over $S$ of
the $X'_i$, and let $f:X'\rightarrow X$ be the product of the $f_i$.
If $S$ is ringed by $\Aa$ as in
\hyperref[XVII.subsubsec5.4.1]{5.4.1}, and the $X_i$ and $X'_i$ are
ringed by the inverse image of $\Aa$, there is more generally an
isomorphism
$$\etiquette[5.4.4.1]{XVII.eq5.4.4.1}
\mathop{\boxtimes}^{\LL}_{i\in I} Rf_{i!}(K_i)
\isomto Rf_{!}(\mathop{\boxtimes}^{\LL}_{i\in I}K_i){\quoi},\leqno(5.4.4.1)
$$
for $K_i\in\Ob D^{-}(X'_i,\Aa)$. Neither the definition of this
isomorphism nor that of the two sides of
\eqhyperref[XVII.eq5.4.4.1]{5.4.4.1} \emph{requires} the choice of a
total order on $I$.

\subsection{Cohomology of symmetric powers}\etiquette[5.5]{XVII.subsec5.5}
\pageoriginale

The main result of this ${\rm n^{\circ}}$ is Theorem
\hyperref[XVII.thm5.5.21]{5.5.21}, which in particular gives the
cohomology with proper supports of a symmetric power of a scheme $X$
over an algebraically closed field $k$, for suitable coefficients (for
example $\ZZ/n\ZZ$), in terms of the cohomology with proper supports
of $X$. The analogue of this theorem in semisimplicial topology is
well known, at least for cohomology with coefficients in a commutative
ring $A$; it can be verified at the level of semisimplicial cochains.
Indeed, if $X_{\ast}$ is a semisimplicial set, there is a canonical
continuous bijection from the symmetric $n\sisi{th}{-th}$ power
$(X_{\ast})^n/\Sf_n$ of $X_{\ast}$ to the
$n$\textsuperscript{th} symmetric power of the geometric realization
of $X_{\ast}$, and this bijection induces an isomorphism on cohomology.

\subsubsection{}\etiquette[5.5.1]{XVII.subsubsec5.5.1}
Recall some properties (due to \textsc{D. Lazard}) of flat modules,
some properties (due to \textsc{Roby}) of the functors $TS^n$ and
$T^n$, and some properties (due to \textsc{Dold} and \textsc{Puppe})
of the derived functors of nonadditive functors.

\paragraph{}\etiquette[5.5.1.1]{XVII.par5.5.1.1}
\emph{Every flat module over a ring $A$ is a filtered direct limit of
finitely generated free modules} (\textsc{Lazard} \cite[I 1.2]{sga4.exp.XVII.bib.La}).
This result implies the following statements:

\paragraph{}\etiquette[5.5.1.2]{XVII.par5.5.1.2}
The functor $\varinjlim$ is an equivalence from the category of
\iInd-objects (\hyperrefext[I][I.sec8.1]{8.1}) of the category of
finitely generated free (\resp finitely generated projective)
$A$-modules to the category of flat $A$-modules.

\paragraph{}\etiquette[5.5.1.3]{XVII.par5.5.1.3}
If\pageoriginale $A$ and $B$ are two rings, every functor
$T:\Mod(A)\rightarrow\Mod(B)$ that carries finitely generated free
modules to flat modules and commutes with filtered direct limits
carries flat modules to flat modules (\textsc{Lazard}
\cite[I 1.5]{sga4.exp.XVII.bib.La}).

\paragraph{}\etiquette[5.5.1.4]{XVII.par5.5.1.4}
Every exact sequence of $A$-modules, with $L''$ flat,

% END INCLUDED SGA4_English_current_buildable_source_20260906/source/17/drafts/SGA4_XVII_LINES_4601_5000_English.texfrag

% BEGIN INCLUDED SGA4_English_current_buildable_source_20260906/source/17/drafts/SGA4_XVII_LINES_5001_5400_English.texfrag
$$
0\longrightarrow L'\longrightarrow L\longrightarrow L''\longrightarrow
0
$$
is a filtered inductive limit of split exact sequences (\textsc{Lazard} \cite[I 2.3]{sga4.exp.XVII.bib.La}).

\subsubsection{}\etiquette[5.5.2]{XVII.subsubsec5.5.2}
Let $A$ be a commutative ring and $M$ an $A$-module.

\paragraph{}\etiquette[5.5.2.1]{XVII.par5.5.2.1}
For every integer $n$, denote by $TS^{n}(M)$ (or $TS^{n}_A(M)$) the
module of symmetric tensors of degree $n$ of $M$. Denoting by
$\Sf_n$ the symmetric group of degree $n$, by definition one has
$$
TS^{n}(M)=(\bigotimes^{n}M)^{\Sf_n}{\quoi}.
$$

If $B$ is an $A$-algebra, then $TS^{n}(B)$ is a subalgebra of
$\bigotimes\limits^{n}B$. If, moreover, $B$ is commutative,
then $\Spec (TS^{n}(B))$ is by definition the $n$\textsuperscript{th}
symmetric power of $\Spec(B)$ over $\Spec(A)$:
$$
\Spec(TS^{n}_{A}(B))=\Sym^{n}_{\Spec(A)}(\Spec(B)\underset{\dfn}{=}(\Spec(B)/\Spec(A))^{n}/\Sf_n{\quoi}.
$$

\paragraph{}\etiquette[5.5.2.2]{XVII.par5.5.2.2}
Recall (\textsc{Roby} \cite[I \sisi{1 pg.226}{2 p.~219}]{sga4.exp.XVII.bib.Ro}) that a \emph{polynomial law} $f$ from an
$A$-module $M$ to an $A$-module $N$ is the datum, for every
commutative $A$-algebra $A'$, of a map
$f_{A'}:M\otimes_{A}A'\rightarrow N\otimes_{A}A'$, these maps
being such that, for every homomorphism of $A$-algebras
$u:A'\rightarrow A''$, the diagram
\[
\xymatrix{
M \otimes_A A'  \ar[r]^{f_{A'}} \ar[d]^{u}& N \otimes_A A' \ar[d]^{u}\\
M \otimes_A A'' \ar[r]^{f_{A''}}& N \otimes_A A''
}
\]
is\pageoriginale commutative.

A polynomial law $f$ is said to be \emph{homogeneous of degree} $n$ (\textsc{Roby} \cite[I
\sisi{1}{8} p\sisi{g}{}.~226]{sga4.exp.XVII.bib.Ro}) if, for every $a'\in A'$, one has identically
\[
f_{A'}(a'x)={a'}^{n}f_{A'}(x){\quoi}.
\]
Such a polynomial law will also be called an \emph{$n$\textsuperscript{ic} map}
(quadratic, cubic, \ldots); their set will be denoted by
\[
\Hom n^{\text{\rm ique}} (M,N){\quoi}.
\]

More generally, if $(M_{i})_{i\in I}$ is a finite family of
$A$-modules and $\sheafn =(n_{i})_{i\in I}$ a family of integers, a
multihomogeneous polynomial law of multidegree $\sheafn $ from the modules $M_{i}$
to $N$ is a polynomial law from $\sisi{\sum}{\oplus} M_{i}$ to $N$ satisfying
identically, for $a_{i}\in A'$,
\[
f_{A'}((a_{i}x_{i})_{i\in I})=\prod_{i\in I} a^{n_{i}}_i\cdot
f_{A'}((x_{i})_{i\in I}){\quoi}.
\]
We shall denote by $\Hom \sheafn ^{\text{\rm ique}}((M_{i}),N)$ the set of these
polynomial laws.

\paragraph{}\etiquette[5.5.2.3]{XVII.par5.5.2.3}
For every module $M$, the covariant functor from $A$-modules to
sets
$$
N\longmapsto \Hom n^{\text{\rm ique}} (M,N)
$$
is represented by an $A$-module $\Gamma^{n}(M)$ (or
$\Gamma^{n}_{A}(M)$), equipped with a \emph{universal\pageoriginale
$n$\textsuperscript{ic} map} $\gamma^{n}:M\rightarrow
\Gamma^{n}(M)$ (\textsc{Roby} \cite[IV 1 \sisi{pg.226}{p.~265}]{sga4.exp.XVII.bib.Ro}). Given $(M_{i})_{i\in I}$
and $\sheafn $ as in \hyperref[XVII.par5.5.2.2]{5.5.2.2}, the covariant functor
$$
N\longmapsto \Hom \sheafn ^{\text{\rm ique}} ((M_{i}),N)
$$
is represented by the module $\bigotimes\limits_{i\in
I}\Gamma^{n_i}(M_{i})$, equipped with the \emph{universal multihomogeneous
polynomial law (of multidegree $\sheafn $)} $(m_{i})\mapsto
\bigotimes\limits_{i\in I}\gamma^{n_{i}}(m_{i})$.

\paragraph{}\etiquette[5.5.2.4]{XVII.par5.5.2.4}
The functors $TS^{n}$ and $\Gamma^{n}$ commute with filtered
inductive limits (for the latter, \textsc{Roby} \cite[IV \sisi{6 pg.275}{7 p.~277}]{sga4.exp.XVII.bib.Ro}). They
take free modules to free modules, and hence flat modules to
flat modules (\hyperref[XVII.par5.5.1.2]{5.5.1.2}).

\paragraph{}\etiquette[5.5.2.5]{XVII.par5.5.2.5}
The $n^{\text{\rm ique}}$ map from $M$ to $TS^{n}(M)$ given by
$m\mapsto \bigotimes\limits^{n}m$ determines a
map $\varphi$, called canonical, from $\Gamma^{n}(M)$ to
$TS^{n}(M)$, satisfying
$$
\varphi(\gamma^{n}(m))=\bigotimes^{n}m{\quoi}.
$$
This map is an isomorphism when $M$ is finite free
(\textsc{Roby} \cite[IV \sisi{pg.272}{5 p.~272}]{sga4.exp.XVII.bib.Ro}), and hence also when $M$ is flat (\hyperref[XVII.par5.5.2.4]{5.5.2.4}
and \hyperref[XVII.par5.5.1.1]{5.5.1.1}).

Let $M$ be an $A$-module and $P$ an exact sequence
$$
P:P_{1}\xrightarrow{d} P_{0} \xrightarrow{\epsilon} M\longrightarrow
0{\quoi}.
$$
Let $P'_{1}=P_{1}\times P_{0}$, and let $j_{0}$ and $j_{1}$ be the
morphisms from $P'_{1}$ to $P_{0}$ with coordinates $(d,\Id)$ and
$(0,\Id)$. The sequence
$$
\xymatrix@C=1.75cm{
\Gamma^{n}(P'_1) \ar@<2pt>[r]^{\Gamma^{n}(j_0)} \ar@<-2pt>[r]_{\Gamma^{n}(j_1)}& \Gamma^{n}(P_0) \ar[r]& \Gamma^{n}(M) \ar[r]& 0
}
$$
is then exact (\textsc{Roby} \cite[\sisi{pg.284}{IV 10 p.~284}]{sga4.exp.XVII.bib.Ro}). \emph{When $P$ is a flat presentation
of} $M$,\pageoriginale this yields a description of the functor
$\Gamma^{n}$ in terms of the functor $TS^{n}$
$$\etiquette[5.5.2.5.1]{XVII.eq5.5.2.5.1}
TS^{n}(P'_1)\rightrightarrows TS^{n}(P_{0})\rightarrow
\Gamma^{n}(M)\rightarrow 0{\quoi}.\leqno(5.5.2.5.1)
$$

\paragraph{}\etiquette[5.5.2.6]{XVII.par5.5.2.6}
If $M=M'\oplus M''$, the $n^{\text{\rm ique}}$ map
\[
m\longmapsto (\gamma^{n'}m\otimes \gamma^{n''}m)_{n'+n''=n}
\]
defines an isomorphism, called canonical (\textsc{Roby} \cite[III \sisi{4}{9 p.~262}]{sga4.exp.XVII.bib.Ro})
$$\etiquette[5.5.2.6.1]{XVII.eq5.5.2.6.1}
\Gamma^{n}(M)\isomto \sisi{\sum}{\bigoplus}_{n'+n''=n} \Gamma^{n'}(M)\otimes \Gamma^{n''}(M){\quoi}.\leqno(5.5.2.6.1)
$$

If $\Sigma$ is an exact sequence
\[
\Sigma:0\longrightarrow M'\longrightarrow M\longrightarrow
M''\longrightarrow 0{\quoi},
\]
the module $\Gamma^{n}(M)$ is equipped with a canonical
(decreasing) filtration and morphisms
$$\etiquette[5.5.2.6.2]{XVII.eq5.5.2.6.2}
\Gamma^{i}(M'')\otimes \Gamma^{n-i}(M')\longrightarrow \Gr^{i}(\Gamma^{n}(M)){\quoi}.\leqno(5.5.2.6.2)
$$
These data are characterized by being functorial
in $\Sigma$ and satisfying the evident compatibility with \eqhyperref[XVII.eq5.5.2.6.1]{5.5.2.6.1}
when $\Sigma$ is split. The arrows \eqhyperref[XVII.eq5.5.2.6.2]{5.5.2.6.2} are
isomorphisms when $\Sigma$ is split, and hence also when $M''$ is flat
(\hyperref[XVII.par5.5.1.4]{5.5.1.4}).

\paragraph{}\etiquette[5.5.2.7]{XVII.par5.5.2.7}
The functor $\Gamma^{n}$ commutes with extension of scalars (\textsc{Roby} \cite[III \sisi{3 pg.262}{8 p.~262}]{sga4.exp.XVII.bib.Ro}) (the analogous property for $TS^{n}$ is false in
general).

\paragraph{}\etiquette[5.5.2.8]{XVII.par5.5.2.8}
If $u:M_{1}\times M_{2}\rightarrow N$ is a bilinear map,
then $\gamma^{n}\circ u$ is a bihomogeneous polynomial law of bidegree
$(n,n)$ and therefore determines (\hyperref[XVII.par5.5.2.3]{5.5.2.3}) $u^{n}:\Gamma^{n}(M_{1})\otimes
\Gamma^{n}(M_{2})\rightarrow \Gamma^{n}(N)$ satisfying
$u^{n}(\gamma^{n}x\otimes
\gamma^{n}y)=\gamma^{n}u(x,y)$. In\pageoriginale particular, if $B$
is an $A$-algebra, multiplication $B\times B\rightarrow B$ induces a
product on $\Gamma^{n}(B)$. This product makes $\Gamma^{n}(B)$ an
$A$-algebra, and the canonical map from $\Gamma^{n}(B)$ to
$TS^{n}(B)$ is an algebra homomorphism.

\subsubsection{}\etiquette[5.5.3]{XVII.subsubsec5.5.3}
Let $\Aa$ be an abelian category. We shall denote by
$\text{\rm ssimpl}^{+}(\Aa)$ the category of
co-semisimplicial objects\nde{In modern terminology they are called \emph{cosimplicial} objects.} in $\Aa$, by $C^{\geq 0}(\Aa)$ the category of
differential complexes $K$ with $K^{i}=0$ for $i<0$, and by $K^{\geq
  0}(\Aa)$ the image of $C^{\geq 0}(\Aa)$ in $K(\Aa)$. If $K\in \Ob
\text{ssimpl}^{+}(\Aa)$, the associated normalized complex
$N(K)\in \Ob C^{\geq 0}(\Aa)$ is defined by
$$\etiquette[5.5.3.1]{XVII.eq5.5.3.1}
\begin{cases}
N(K)^{n} =\text{the intersection of the kernels of the degeneracy
operators}\\
\qquad\qquad\text{from $K^{n}$ to $K^{n-1}$}\\
d=\sum(-1)^{i}\partial_{i}{\quoi}.
\end{cases}\leqno(5.5.3.1)
$$

If $K$ also denotes the complex with components $K^{n}$ and
differential $d=\sum(-1)^{i}\partial_{i}$, then $N(K)$ is a
direct-summand subcomplex of $K$, and $K/N(K)$ is homotopic to
zero.

By \textsc{Dold-Puppe} \cite{sga4.exp.XVII.bib.DP}, the functor $N$ is an equivalence of
categories
$$
N:\;\text{\rm ssimpl}^{+}(\Aa)\longrightarrow C^{\geq
0}(\Aa){\quoi}.
$$
This equivalence and its inverse take homotopic morphisms to
homotopic morphisms. Moreover, for every $K\in \Ob\text{\rm
ssimpl}^{+}(\Aa)$, $K^{n}$ is a direct sum of copies of $N(K)^{i}$
$(i\leq n)$; if $\PP$ is a property of objects of
$\Aa$, stable under finite direct sums and direct summands, it
follows that the $K^{i}$ have property $\PP$ for $i\leq n$ if
and only if the $N(K)^{i}$ have this property for $i\leq n$.

The equivalence\pageoriginale $N$ is used by \textsc{Dold} and
\textsc{Puppe} to
extend every functor (not necessarily additive)
$T:\Aa\rightarrow \Bb$ from $\Aa$ to an
arbitrary category $\Bb$ to a functor $C^{\geq 0}
(T)$ (or simply $T$) from $C^{\geq 0}(\Aa)$ to $C^{\geq
0}(\Bb)$
$$\etiquette[5.5.3.2]{XVII.eq5.5.3.2}
C^{\geq 0}(T)=NTN^{-1}{\quoi}.\leqno(5.5.3.2)
$$

The functor $C^{\geq 0}(T)$ passes to the quotient to define
$$\etiquette[5.5.3.3]{XVII.eq5.5.3.3}
T\quad\mbox{\rm or}\quad K^{\geq 0}(T):K^{\geq
0}(\Aa)\longrightarrow K^{\geq 0}(\Bb){\quoi}.\leqno(5.5.3.3)
$$

\setcounter{paragraph}{3}
\paragraph{}\etiquette[5.5.3.4]{XVII.par5.5.3.4}
If $T$ is \emph{of degree} $\leq n$, \ie if the $(n+1)$\textsuperscript{th} “cross-effect” vanishes (\textsc{Dold-Puppe} \cite[4.18 \sisi{pg.232}{p.~232}]{sga4.exp.XVII.bib.DP})\sisi{}{,}
then, for every $K\in \Ob C^{\geq 0}(\Aa)$ such that $K^{i}=0$
for $i>m$, one has $[T(K)]^{i}=0$ for $i>nm$ (\textsc{Dold-Puppe} \cite[4.2\sisi{.}{}3
\sisi{pg.233}{p.~233}]{sga4.exp.XVII.bib.DP}).

\paragraph{}\etiquette[5.5.3.5]{XVII.par5.5.3.5}
In particular, if $T$ is a finite-degree functor from the category
of $A$-modules to that of $B$-modules, which takes flat objects
to flat objects, then the image of a bounded flat complex is a
bounded flat complex.

\begin{enonce*}{Lemma 5.5.4}\etiquette[5.5.4]{XVII.lem5.5.4}
Let $\Aa$ be an abelian category.
\begin{enumerate}
\renewcommand{\theenumi}{\roman{enumi}}
\renewcommand{\labelenumi}{(\theenumi)}
\item The quasi-isomorphisms in $K^{\geq 0}(\Aa)$
admit a calculus of left and right fractions.

\item The evident functor is an equivalence between the category
obtained from $K^{\geq 0}(\Aa)$ by inverting the
quasi-isomorphisms and the category $D^{\geq 0}(\Aa)$.
\end{enumerate}
\end{enonce*}

The\pageoriginale functor $\tau_{\geq 0}:C(\Aa)\rightarrow
C^{\geq 0}(\Aa)$ (\hyperref[XVII.def1.1.13]{1.1.13}) passes to the quotient to define a
functor $\tau_{\geq 0}:K(\Aa)\rightarrow K^{\geq
0}(\Aa)$, left adjoint to the inclusion of $K^{\geq
0}(\Aa)$ in $\sisi{\Aa}{K(\Aa)}$. The functor $\tau_{\geq 0}$
takes quasi-isomorphisms to quasi-isomorphisms, and the adjunction arrow
$X\rightarrow \tau_{\geq 0}(X)$ is a quasi-isomorphism
if and only if $H^{i}(X)=0$ for $i<0$. Lemma
\hyperref[XVII.lem5.5.4]{5.5.4} therefore follows
from \textsc{Verdier} \cite{sga4.exp.XVII.bib.CD} and parts (i) and (ii) of the following lemma, whose
proof is left to the reader:

\begin{enonce*}{Lemma 5.5.4.1}\etiquette[5.5.4.1]{XVII.lem5.5.4.1}
Let $K$ be a category, $K'$ a full subcategory, $S$ a
set of morphisms of $K$ equal to the inverse image of the set of
isomorphisms of $K(S^{-1})$ (\cf \cite[\sisi{chap 1}{\chap~1}]{sga4.exp.XVII.bib.GZ}), and $S'$ its
trace in $K'$. Suppose that the inclusion $i$ of $K'$ in $K$ has
a left adjoint $\tau$, and that $\tau(S)\subset S'$.
\begin{enumerate}
\renewcommand{\theenumi}{\roman{enumi}}
\renewcommand{\labelenumi}{(\theenumi)}
\item The evident functor from $K'({S'}^{-1})$ to $K(S^{-1})$ is
fully faithful, and has $\tau$ as a left adjoint.

\item If $K(S^{-1})$ admits a calculus of left fractions
(\resp right fractions), then $K'({S'}^{-1})$ has the same property.

If the inclusion $i$ does not necessarily have a left adjoint, and if
$\tau$ denotes the \sisi{partially}{not everywhere} defined functor left adjoint to $i$,
one has:

\item If $K(S^{-1})$ admits a calculus of left fractions
(\resp right fractions) and if, for every $u\in S$ with target (\resp source)
in $K'$, $\tau(u)$ is defined and belongs to $S'$, then $K'({S'}^{-1})$
admits a calculus of left fractions (\resp right fractions), and the functor
$K'({S'}^{-1})\rightarrow K(S^{-1})$ is fully faithful.
\end{enumerate}
\end{enonce*}

\setcounter{subsubsection}{4}
\subsubsection{}\etiquette[5.5.5]{XVII.subsubsec5.5.5}
Let\pageoriginale $T:\Aa\rightarrow \Bb$ be a functor
between abelian categories. As in \hyperref[XVII.def1.2.1]{1.2.1}, one defines the derived functors of
$T$
\begin{align*}
LT &: D^{\geq 0}(\Aa)\longrightarrow \Pro(D^{\geq
0}(\Bb))\\
RT &: D^{\geq 0}(\Aa)\longrightarrow \Ind(D^{\geq 0}(\Bb))
\end{align*}
by the formulas
\begin{align*}
LT(K) &= \quotprojlim_{s:K'\rightarrow K}T(K')\\
RT(K) &= \quotinjlim_{s:K\rightarrow K'}T(K'){\quoi},
\end{align*}
where $s$ ranges over the filtered category of quasi-isomorphisms with target
(\resp source) $K$.

We say that $T$ is \emph{left (\resp right) derivable at $K$ if
$LT(K)$} $(RT(K))$ is representable. We say that $K$ is \emph{left
(\resp right) adapted} to $T$ (or, by abuse of language,
adapted to $LT$ (\resp $RT$)) if the following canonical arrow is an
isomorphism:
\[
LT(K)\isomto  TK\text{ (\resp $TK \isomto  RT(K)$)}{\quoi}.
\]

\begin{enonce*}{Proposition 5.5.6}\etiquette[5.5.6]{XVII.prop5.5.6}
Let $F$ be a functor from the category of modules over a ring $A$
to an abelian category $\Bb$. Suppose that filtered
inductive limits in $\Bb$ are representable and
exact, and that $F$ commutes with filtered
inductive limits. Then every bounded complex $M^{\bigdot}\in \Ob K^{\geq
0}(A)$ of flat modules is adapted to $LF$.
\end{enonce*}

The\pageoriginale proposition follows immediately from the definitions, the
following lemma, which will be proved in \hyperref[XVII.par5.5.6.4]{5.5.6.4}, and the fact that the
quasi-isomorphisms $u:N^{\bigdot}\rightarrow M^{\bigdot}$, with
$N^{\bigdot}\in \sisi{B}{C}^{b,\geq 0}(A)$ having flat components, form a
cofinal set in the category of quasi-isomorphisms with target
$M^{\bigdot}$.

\begin{enonce*}{Lemma 5.5.6.1}\etiquette[5.5.6.1]{XVII.lem5.5.6.1}
Let $u:N^{\bigdot}\rightarrow M^{\bigdot}$ be a quasi-isomorphism
between flat complexes that are objects of $C^{b,\geq 0}(A)$. Then $F(u)$ is
a quasi-isomorphism.
\end{enonce*}

\begin{enonce*}{Lemma 5.5.6.2}\etiquette[5.5.6.2]{XVII.lem5.5.6.2}
Let $\Aa$ be an additive category and $d$ an interval in
$\ZZ$ that is bounded above. Denote by $C^{d}$ (\resp
$C^{d,b}$, \resp $K^{d}$, \resp $K^{d,b}$) the category of complexes
that vanish in degrees $i\not\in d$ (\resp and are bounded, \resp up to
homotopy). There is an equivalence
$$
\sisi{\sheafell }{\ell}:\Ind(C^{d,b}(\Aa)\sisi{}{)}\isomto
C^{d}(\Ind\Aa){\quoi};
$$
the inverse of this equivalence induces a functor
$$
\sisi{\sheafell }{\ell}':K^{d}(\Ind\Aa)\longrightarrow
\Ind(K^{d,b}(\Aa)){\quoi}.
$$
\end{enonce*}

\begin{enumerate}
\renewcommand{\theenumi}{\Alph{enumi}}
\renewcommand{\labelenumi}{\theenumi)}
\item Let us prove that $\sisi{\sheafell }{\ell}$ is fully faithful. Thus, let $K$ be a
bounded complex of objects of $\Aa$ and $L=\quotinjlim
L_{\alpha}$ an \iInd-complex. One has
\begin{align*}
\Hom(K,L)\overset{\dfn}{=} \varinjlim \Hom(K,L_{\alpha})  = & \varinjlim
Z^{0}(\Hom^{\bigdot} (K,L_{\alpha}))\simeq Z^{0}
\varinjlim_{\alpha} \Hom^{\bigdot} (K,L_{\alpha})\\
& \isomto
Z^{0} \Hom^{\bigdot} (K,\ell L)=\Hom(K,\ell L){\quoi},
\end{align*}
the penultimate isomorphism following from the fact that $K$ is finite, so
that $\varinjlim$ commutes with passage from the double complex to the total
complex. This proves the assertion.

\item We\pageoriginale leave it to the reader to verify (rather laboriously) that
$\sisi{\sheafell }{\ell}$ is essentially surjective by constructing,
from left to right and starting in an arbitrarily chosen degree, “enough” bounded complexes mapping to a
bounded-above complex of \iInd-objects $K$; and likewise to develop the
following point:

\item If $\ell f:\ell \varinjlim K_{\alpha}\rightarrow \ell \varinjlim
L_{\beta}$ is null-homotopic, $f=dH+Hd$, then, for every $\alpha$,
$H$ lifts to $H^{i}_{\alpha}:K^{i}_{\alpha}\rightarrow
L^{i\sisi{}{-1}}_{\beta(\alpha)}$ which, for $\beta$ sufficiently large, defines a
null-homotopy of $f_{\alpha}:K_{\alpha}\rightarrow L_{\beta(\alpha)}$.
\end{enumerate}

\setcounter{subsubsection}{6}
\setcounter{paragraph}{2}
\paragraph{}\etiquette[5.5.6.3]{XVII.par5.5.6.3}
Let $A$ be a ring, $d$ an interval in $\ZZ$ bounded above,
$P(A)$ the category of finite projective $A$-modules,
and $\Pl(A)$ the category of flat $A$-modules. Denote by
$D^{\Tor\subset d}(A)$ the image of $C^{d}(\Pl(A))$ in $D(A)$ and by
$D^{\parf\subset d}(A)$ the image of $C^{d,b}(P(A))$ in
$D(A)$. By \hyperref[XVII.par5.5.1.2]{5.5.1.2}, one has $\Pl(A)\simeq \Ind
P(A)$. On the
other hand, $K^{d,b}(P(A))\simeq D^{\parf\subset d}(A)$. It then follows
from \hyperref[XVII.lem5.5.6.2]{5.5.6.2}, and from the fact that $H^{i}$ commutes with filtered
inductive
limits, that there exists a commutative diagram of functors
$$
\xymatrix@C=2.25cm{
C^{d}(\text{Pl}(A)) \ar[d]& \Ind C^{d,b}(P(A)) \ar[d]\ar[l]_{\substack{\varinjlim\\\sim}}\\
K^{d}(\text{Pl}(A)) \ar[d]\ar[r]^{\ell'}& \Ind K^{d,b}(P(A)) \ar[d]^{\sisi{}{\simeq}}\\
D^{\Tor\subset d}(A) \ar[r]^{\ell''}& \Ind D^{\parf \subset d}(A).
}
$$

% END INCLUDED SGA4_English_current_buildable_source_20260906/source/17/drafts/SGA4_XVII_LINES_5001_5400_English.texfrag

% BEGIN INCLUDED SGA4_English_current_buildable_source_20260906/source/17/drafts/SGA4_XVII_LINES_5401_5800_English.texfrag

\setcounter{subsubsection}{6}
\paragraph{}\etiquette[5.5.6.4]{XVII.par5.5.6.4}
Let us prove\pageoriginale \hyperref[XVII.lem5.5.6.1]{5.5.6.1}. With
the notation of \hyperref[XVII.par5.5.6.3]{5.5.6.3}, if $d$ is an
interval $[0,k]$, then (because $F$ commutes with $\varinjlim$) there
are commutative diagrams
$$
\xymatrix@C=2.5cm{
K^{d}(\text{Pl}(A))          \ar[r]^{\ell'} \ar[d]_{K^{\geq 0}(F)} & \Ind K^d(P(A))  \ar[d]_{\Ind(K^{\geq 0}(F))}             \\
K^{+}(\Bb) \ar[d]_{H^i}                        & \Ind K^{+}(\Bb)     \ar[d]_{\Ind(H^i)}       \\
\Bb                  & \Ind \Bb  \ar[l]_{\varinjlim}.
}
$$
If $u$ is as in \hyperref[XVII.lem5.5.6.1]{5.5.6.1}, then by
\hyperref[XVII.par5.5.6.3]{5.5.6.3}, $\ell'(u)$ is an isomorphism, and
$H^iK^{\geq0}(F)(u)$ is therefore an isomorphism, as required.

\subsubsection{}\etiquette[5.5.7]{XVII.subsubsec5.5.7}
Let $X$ be a scheme over a scheme $S$, and let $k$ be an integer.
Suppose that the symmetric power $\Sym^k_S(X)$ is defined. A reader
who wishes not to make this hypothesis may work in the category of
algebraic spaces (\textsc{Artin}\sisi{ [1]}{\nde{the original text
cites an unspecified work here, probably one of the two works of Artin
cited in the notes on pages \pageref{XVII.noteArtin1} and
\pageref{XVII.noteArtin2}.}}). Ring $S$ by a sheaf $\Aa$ of
\emph{commutative} rings. Beginning with
\hyperref[XVII.subsubsec5.5.19]{5.5.19}, $\Aa$ will be assumed to be
\emph{torsion}. By abuse of notation, we shall \emph{still denote by}
$\Aa$ the inverse image of $\Aa$ on any $S$-scheme.

Let $ℱ$ be a sheaf of modules on $X$; the sheaf
$\mathop{\boxtimes}\limits^kℱ$, the external tensor product of $k$
copies of $ℱ$, is a sheaf on $X^k/S$, equivariant under the action of
the symmetric group $\Sf_k$ on $X^k/S$. Passing\pageoriginale to
invariants, it defines a sheaf $TS^k_{\ext}(ℱ)$ on $\Sym^k_S(X)$,
called the \emph{external symmetric tensor power of $ℱ$}\sisi{}{.}

To reason ``point by point,'' it is useful to work with the points of
$\Sym^k_S(X)$ described as follows.

\paragraph{}\etiquette[5.5.7.1]{XVII.par5.5.7.1}
Let $s$ be a geometric point of $S$. A point $x$ of $\Sym^k_S(X)$
with values in $k(s)$ defines a geometric point of $\Sym^k_S(X)$. Let
$(x_1\ldots x_k)$ denote a point of $(X/S)^k$ with values in $k(s)$
whose image is $x$, let $y_1\ldots y_\ell$ be the distinct elements of
the sequence $x_1\ldots x_k$, and let $n_i$ $(1\leq i\leq\ell)$ be the
multiplicity of $y_i$ in this sequence $(\sum n_i=k)$.

With the notation of \hyperref[XVII.par5.5.7.1]{5.5.7.1}, and writing
$\Sf$ for a symmetric group, one has
$$\etiquette[5.5.7.2]{XVII.eq5.5.7.2}
TS^{k}_{\ext}(F)_x\simeq (\bigotimes^{\ell}_{1}
(\bigotimes^{n_i} F_{y_i}))^{\prod^{\ell}_{1}\Sf_{n_i}}{\quoi},\leqno(5.5.7.2)
$$
where the tensor products are taken over $\Aa_S$. In particular,
\emph{for $F$ flat}, by \hyperref[XVII.par5.5.2.5]{5.5.2.5} one has
$$\etiquette[5.5.7.3]{XVII.eq5.5.7.3}
TS^{k}_{\ext}(F)_x\isomfrom \bigotimes^{\ell}_{1}
\Gamma^{n_i}F_{y_i}{\quoi},\leqno(5.5.7.3)
$$
using induction on $\ell$, the flatness of the
$\Gamma^{n_i}F_{y_i}$ (\hyperref[XVII.par5.5.2.4]{5.5.2.4}), and the
following fact: if $G$ (\resp $H$) is a finite group acting on an
$A$-module $M$\sisi{}{ }(\resp $N$), and if $M$ and $N^H$ are flat,
then $M^G\otimes_A N^H\isomto(M\otimes_A N)^{G\times H}$.

\subsubsection{}\etiquette[5.5.8]{XVII.subsubsec5.5.8}
Let $F$ be a module on $X$, and let $P$ be a flat presentation of $F$:
$$
P:P_1\longrightarrow P_0\longrightarrow F\longrightarrow0{\quoi}.
$$
Define\pageoriginale $P'_1$, $j_0$, and $j_1$ as in
\hyperref[XVII.par5.5.2.5]{5.5.2.5}, and define
$\Gamma^k_{\ext}(F)$ (the \emph{external $\Gamma$ functor}) to be the
cokernel of the pair of arrows
$(TS^k_{\ext}(j_0),TS^k_{\ext}(j_1))$. With the notation of
\hyperref[XVII.par5.5.7.1]{5.5.7.1}, one deduces from
\eqhyperref[XVII.eq5.5.7.3]{5.5.7.3} and
\hyperref[XVII.par5.5.2.5]{5.5.2.5} isomorphisms
$$\etiquette[5.5.8.1]{XVII.eq5.5.8.1}
\Gamma^k_{\ext}(F)_x\simeq\bigotimes^{\ell}_{1}
\Gamma^{n_i}(F_{y_i}){\quoi}.\leqno(5.5.8.1)
$$

This formula shows that $\Gamma^k_{\ext}(F)$, which a priori depends
functorially on $P$, depends, up to a unique isomorphism, only on $F$.

\subsubsection{}\etiquette[5.5.9]{XVII.subsubsec5.5.9}
There is a functorial morphism
$$
\alpha:\Gamma^k_{\ext}(F)\longrightarrow TS^k_{\ext}(F){\quoi}.
$$
The pair $(\Gamma,\alpha)$ is characterized by the following two
properties:
\begin{enumerate}
\renewcommand{\theenumi}{\roman{enumi}}
\renewcommand{\labelenumi}{(\theenumi)}
\item For $F$ flat, $\alpha$ is an isomorphism.

\item With the notation of \hyperref[XVII.par5.5.2.5]{5.5.2.5}, for
every presentation
$$
P_1\xrightarrow{d}P_0\xrightarrow{\epsilon}F\longrightarrow0
$$
of $F$, the sequence
$$
\Gamma^k_{\ext}(P'_1)\rightrightarrows
\Gamma^k_{\ext}(P_0)\rightarrow\Gamma^k_{\ext}(F)\rightarrow0
$$
is exact.

Moreover:

\item The functors $TS^k$ and $\Gamma^k$ commute with filtered direct
limits (by
\sisi{\hyperref[XVII.eq5.5.7.2]{5.5.7.2}}{\eqhyperref[XVII.eq5.5.7.2]{5.5.7.2}},
\sisi{\hyperref[XVII.eq5.5.8.1]{5.5.8.1}}{\eqhyperref[XVII.eq5.5.8.1]{5.5.8.1}},
\hyperref[XVII.par5.5.2.4]{5.5.2.4}).

\item The functor $\Gamma^k$ commutes with every extension
$\Aa\rightarrow\Aa'$ of the base ring.

\item The\pageoriginale functor $\Gamma^k$ commutes with every base
change $S'\rightarrow S$: for every cartesian square
$$
\xymatrix@C=2.5cm@R=1.5cm{
X' \ar[r]^{g'}\ar[d]^{f'}& X \ar[d]^{f}\\
S' \ar[r]^{g}& S,
}
$$
one has
$$
\Sym^k_g(g')^{\ast}\Gamma^k_{\ext}\simeq
\Gamma^k_{\ext}{g'}^{\ast}{\quoi}.
$$
(\sisi{\hyperref[XVII.eq5.5.8.1]{5.5.8.1}}{\eqhyperref[XVII.eq5.5.8.1]{5.5.8.1}}
and \hyperref[XVII.par5.5.2.7]{5.5.2.7}).

\item The functor $\Gamma^k$ commutes with inverse images under every
\emph{radicial} morphism $i:X'\rightarrow X$ (for example a closed
immersion): \sisi{O}{o}ne has
$$
\Sym^k_S(i)^{\ast}\Gamma^k_{\ext}\simeq
\Gamma^k_{\ext}i^{\ast}
$$
\eqhyperref[XVII.eq5.5.8.1]{5.5.8.1}.

\item For an \emph{immersion} $j:X'\rightarrow X$, one has
$$
\Sym^k_S(j)_{!}\Gamma^k_{\ext}\simeq\Gamma^k_{\ext}j_{!}
$$
\eqhyperref[XVII.eq5.5.8.1]{5.5.8.1}.
\end{enumerate}

\subsubsection{}\etiquette[5.5.10]{XVII.subsubsec5.5.10}

Let $k_1,k_2$, and $k=k_1+k_2$ be three integers $\geq0$. There is
then a finite morphism (\quad):\nde{I have not found references for the
finiteness of $\vee_{k_1k_2}$ in general.}
$$
\vee_{k_1k_2}:\Sym^{k_1}_S(X)\times_S\Sym^{k_2}_S(X)
\longrightarrow\Sym^k_S(X){\quoi}.
$$
If $F_i$ $(i=1,2)$ is a module on $\Sym_S^{k_i}(X)$, denote by
$F_1\vee F_2$ the sheaf on $\Sym_S^k(X)$
$$
F_1\vee F_2=\vee_{k_1k_2\ast}(F_1\boxtimes F_2){\quoi}.
$$

\subsubsection{}\etiquette[5.5.11]{XVII.subsubsec5.5.11}
Let\pageoriginale there be an exact sequence of modules on $X$
\[
\Sigma:0\longrightarrow F'\longrightarrow F\longrightarrow
F''\longrightarrow0{\quoi}.
\]
With the notation of \hyperref[XVII.par5.5.7.1]{5.5.7.1}, by
\eqhyperref[XVII.eq5.5.8.1]{5.5.8.1} and
\hyperref[XVII.par5.5.2.6]{5.5.2.6} there is a filtration on
$\Gamma^k_{\ext}(F)_x$, and there are morphisms
\[
\sum_{\sum^{\ell}_{1}i_\alpha=i}
\bigotimes^{\ell}_{1}\Gamma^{i_\alpha}(F'_{y_\alpha})\otimes
\bigotimes^{\ell}_{1}\Gamma^{n_i-i_\alpha}(F''_{y_\alpha})
\longrightarrow\Gr^i(\Gamma^k_{\ext}(F)_x){\quoi}.
\]
The first member is precisely
$(\Gamma^i_{\ext}(F')\vee\Gamma^{k-i}_{\ext}(F''))_x$. One readily
checks that these data, constructed point by point, arise from a
(unique) filtration with $k+1$ steps on $\Gamma^k_{\ext}(F)$ and from
(unique) morphisms
$$\etiquette[5.5.11.1]{XVII.eq5.5.11.1}
\Gamma^i_{\ext}(F')\vee\Gamma^{k-i}_{\ext}(F'')
\longrightarrow\Gr^i(\Gamma^k_{\ext}(F)){\quoi}.\leqno(5.5.11.1)
$$

If the exact sequence $\Sigma$ is split, one has more precisely
(verification point by point from
\sisi{\hyperref[XVII.eq5.5.2.6.2]{5.5.2.6.2}}{\eqhyperref[XVII.eq5.5.2.6.1]{5.5.2.6.1}}):
$$\etiquette[5.5.11.2]{XVII.eq5.5.11.2}
\Gamma^k_{\ext}(F'\oplus F'')\simeq
\sisi{\sum}{\bigoplus}_{k_1+k_2=k}\Gamma^{k_1}_{\ext}(F')\vee
\Gamma^{k_2}_{\ext}(F''){\quoi}.\leqno(5.5.11.2)
$$

If the exact sequence $\Sigma$ is split point by point, or if $F''$ is
flat, the morphisms \eqhyperref[XVII.eq5.5.11.1]{5.5.11.1} are
isomorphisms (\cf \hyperref[XVII.par5.5.1.4]{5.5.1.4}).

\subsubsection{}\etiquette[5.5.12]{XVII.subsubsec5.5.12}
Let $u:X\rightarrow Y$ be a morphism of $S$-schemes. Suppose that the
symmetric powers $\Sym^k_S(X)$ and $\Sym^k_S(Y)$ exist. Let $F$ be a
module on $X$, let $G$ be a module on $Y$, and let
$v\in\Hom(u^{\ast}G,F)\underset{\sisi{n}{}}{\simeq}\Hom(G,u_{\ast}F)$
be a $u$-homomorphism from $F$ to $G$ \hyperlink{sga4.exp.XVII.subsection.0.2}{(0.2)}. If $\prod^n u$ is the
$n$\textsuperscript{th} cartesian power of $u$, the homomorphism $v$
induces a $\prod^n u$-homomorphism $\mathop{\boxtimes}\limits^n v$
from $\mathop{\boxtimes}\limits^nF$ to
$\mathop{\boxtimes}\limits^nG$. Passing\pageoriginale to invariants,
the latter defines
\[
TS^n(v):TS^n_{\ext}(F)\xrightarrow[\Sym_S^n(u)]{}TS^n_{\ext}(G){\quoi}.
\]

It follows from Definition \hyperref[XVII.subsubsec5.5.8]{5.5.8} that
there exists a unique $\Sym^n(u)$-homomorphism
$\Gamma^n_{\ext}(v)$, functorial in $(F,G,v)$,
$$\etiquette[5.5.12.0]{XVII.eq5.5.12.0}
\Gamma^n_{\ext}(v):\Gamma^n_{\ext}(F)
\xrightarrow[\Sym^n_S(u)]{}\Gamma^n_{\ext}(G)\leqno(5.5.12.0)
$$
making the following diagram commutative (\cf \hyperlink{sga4.exp.XVII.subsection.0.2}{(0.2)} for the unusual
direction of the arrows):
$$
\xymatrix@C=3.5cm@R=2cm{
\Gamma^n_{\ext}(F) \ar[r]^{\Gamma^n_{\ext}(v)}_{\text{over $\Sym^n_S(u)$}}& \Gamma^n_{\ext}(G)\\
TS^n_{\ext}(F) \ar[r]^{TS^n_{\ext}(v)}_{\text{over $\Sym^n_S(u)$}}\ar[u]^{\alpha}_{\text{over $\Id$}} & TS^n_{\ext}(G)\ar[u]^{\alpha}_{\text{over $\Id$}}.
}
$$

We shall use this arrow mainly as a ``symmetric Künneth arrow'': for
$G=u_{\ast}F$ and $v$ the canonical arrow, the arrow
$\Gamma^n_{\ext}(v)$ identifies with a homomorphism\sisi{s}{} of
modules
$$\etiquette[5.5.12.1]{XVII.eq5.5.12.1}
k^n:\Gamma^n_{\ext}(u_{\ast}F)\longrightarrow
\Sym^n(u)_{\ast}\Gamma^n_{\ext}(F){\quoi}.\leqno(5.5.12.1)
$$

\subsubsection{}\etiquette[5.5.13]{XVII.subsubsec5.5.13}
Consider the condition:

\paragraph{}\etiquette[5.5.13.1]{XVII.par5.5.13.1}
$K\in\Ob K^{b,\geq0}(X,\Aa)$ and, for every geometric point $x$ of
$X$, the complex $K_x$ is homotopy equivalent to a complex of flat
modules $L_x\in\Ob C^{b,\geq0}(\Aa_x)$.

\begin{enonce*}{Lemma 5.5.13.2}\etiquette[5.5.13.2]{XVII.lem5.5.13.2}
The\pageoriginale subcategory $D^{b,\tor\leq0}(X,\Aa)$ of
$D^b(X,\Aa)$ consisting of complexes of \tTor-dimension $\leq0$
\emph{\hyperref[XVII.def4.1.9]{4.1.9}} is equivalent to the category
obtained from the category of complexes
$K\in\Ob K^{b,\geq0}(X,\Aa)$ satisfying
\emph{(\hyperref[XVII.par5.5.13.1]{5.5.13.1})} by inverting
quasi-isomorphisms. This construction may be performed indifferently
by a calculus of left or right fractions.
\end{enonce*}

Let $P$ be the subcategory of $K^{-}(X,\Aa)$ consisting of complexes
that are pointwise homotopy equivalent to a complex of flat modules.
It is known (\textsc{Verdier}
\sisi{\cite{sga4.exp.XVII.bib.CD}}{\cite[I 2.4.2]{sga4.exp.XVII.bib.CD}}) that $D^{-}(X,\Aa)$ is obtained
from $P$ by inverting quasi-isomorphisms, using a calculus of right or
left fractions. Apply \hyperref[XVII.lem5.5.4.1]{5.5.4.1} to the
inclusion of $K^{b,\geq0}(X,\Aa)\cap P$ into $P$. Indeed, if $K$ has
\tTor-dimension $\leq0$ and is pointwise homotopy equivalent to
complexes of flat modules $L_x$, then $\tau_{\geq0}(K)$ is pointwise
homotopy equivalent to the $\tau_{\geq0}(L_x)$, and the latter are
still complexes of flat modules, so $\tau_{\geq0}(K)$ satisfies
\hyperref[XVII.par5.5.13.1]{5.5.13.1}.

\subsubsection{}\etiquette[5.5.14]{XVII.subsubsec5.5.14}
By \hyperref[XVII.lem5.5.6.1]{5.5.6.1},
\sisi{\hyperref[XVII.eq5.5.7.2]{5.5.7.2}}{\eqhyperref[XVII.eq5.5.7.2]{5.5.7.2}}
or
\sisi{\hyperref[XVII.eq5.5.8.1]{5.5.8.1}}{\eqhyperref[XVII.eq5.5.8.1]{5.5.8.1}}
and \hyperref[XVII.par5.5.2.4]{5.5.2.4}, every quasi-isomorphism
$v:K'\rightarrow K$ between complexes satisfying
(\hyperref[XVII.par5.5.13.1]{5.5.13.1}) induces quasi-isomorphisms
\begin{align*}
TS^k(v)&:TS^k_{\ext}(K')\xrightarrow{\approx}TS^k_{\ext}(K)\\
\Gamma^k(v)&:\Gamma^k_{\ext}(K')\xrightarrow{\approx}\Gamma^k_{\ext}(K){\quoi}.
\end{align*}
It follows that the functors $\Gamma^k$ and $TS^k$ are derivable
(\hyperref[XVII.subsubsec5.5.5]{5.5.5}) on the subcategory
$D^{b,\tor\leq0}(X,\Aa)$ of $D^b(X,\Aa)$, and that complexes
satisfying \hyperref[XVII.par5.5.13.1]{5.5.13.1} are left adapted to
these functors. By \hyperref[XVII.par5.5.2.4]{5.5.2.4},\pageoriginale
$\Gamma^k$ and $TS^k$ induce functors
\begin{align*}
LTS^k &:D^{b,\tor\leq0}(X,\Aa)\longrightarrow
D^{b,\tor\leq0}(\Sym^k_S(X),\Aa)\tag{5.5.14.1}\\
L\Gamma^k &:D^{b,\tor\leq0}(X,\Aa)\longrightarrow
D^{b,\tor\leq0}(\Sym^k_S(X),\Aa){\quoi}.\tag{5.5.14.2}
\end{align*}\etiquette[5.5.14.1]{XVII.eq5.5.14.1}
\etiquette[5.5.14.2]{XVII.eq5.5.14.2}%Removing this comment would introduce unwanted white space.
Moreover, by \hyperref[XVII.subsubsec5.5.9]{5.5.9}~(i), the canonical
morphism $\alpha$ (\hyperref[XVII.subsubsec5.5.9]{5.5.9}) induces an
isomorphism between the functors \eqhyperref[XVII.eq5.5.14.1]{5.5.14.1}
and \eqhyperref[XVII.eq5.5.14.2]{5.5.14.2}:
$$\etiquette[5.5.14.3]{XVII.eq5.5.14.3}
\alpha:L\Gamma^k\isomto LTS^k{\quoi}.\leqno(5.5.14.3)
$$

\setcounter{paragraph}{3}
\paragraph{}\etiquette[5.5.14.4]{XVII.par5.5.14.4}
The isomorphisms of \hyperref[XVII.subsubsec5.5.9]{5.5.9}~(iv), (v),
(vi), and (vii) pass to the derived category and define isomorphisms:
\begin{itemize}
\item[${\rm(iv')}$] For an $\Aa$-algebra $\Aa'$, one has
$$
L\Gamma^k_{\ext_{\Aa}}(K)\overset{\LL}{\otimes}\Aa'\simeq
L\Gamma^k_{\ext_{\Aa'}}(K\overset{\LL}{\underset{\Aa}{\otimes}}\Aa'){\quoi}.
$$

\item[${\rm(v')}$] For every base change $g:S'\rightarrow S$, one has
$$
\Sym^k_g(g')^{\ast}L\Gamma^k_{\ext}K\simeq
L\Gamma^k_{\ext}({g'}^{\ast}K){\quoi}.
$$

\item[${\rm(vi')}$] For every \emph{radicial} morphism
$i:X'\rightarrow X$, one has
$$
\Sym^k_S(i)^{\ast}L\Gamma^k_{\ext}K\simeq
L\Gamma^k_{\ext}i^{\ast}K{\quoi}.
$$

\item[${\rm(vii')}$] For an \emph{immersion} $j:X'\rightarrow X$,
one has
$$
\Sym^k_S(j)_{!}L\Gamma^k_{\ext}\simeq
L\Gamma^k_{\ext}j_{!}{\quoi}.
$$
\end{itemize}

\subsubsection{}\etiquette[5.5.15]{XVII.subsubsec5.5.15}
Let there be an exact sequence of flat complexes in
$C^{\geq0}(X,\Aa)$,
$$
0\longrightarrow K'\longrightarrow K\longrightarrow K''\longrightarrow0{\quoi},
$$
and apply the functor $\Gamma^k_{\ext}$
(\hyperref[XVII.subsubsec5.5.3]{5.5.3}). One finds (as in
\hyperref[XVII.subsubsec5.5.11]{5.5.11}) that the semisimplicial module
$\Gamma^k_{\ext}(N^{-1}K)$ on $\Sym^k_S(X)$ has a filtration
whose\pageoriginale successive quotients are the semisimplicial
modules $\Gamma^{k'}_{\ext}(N^{-1}K')\vee
\Gamma^{k''}_{\ext}(N^{-1}K'')$ $(k'+k''=k)$.

Passing to the derived category, one finds that
$L\Gamma^k_{\ext}(K)$ is represented by a filtered complex whose
successive quotients represent the objects
$L\Gamma^i_{\ext}(K'\sisi{}{)}\overset{\LL}{\vee}
L\Gamma^{k-i}_{\ext}(K'')$ of the derived category (apply
\hyperref[XVII.par5.5.2.4]{5.5.2.4} and the
\textsc{Eilenberg--Zilber--Cartier} theorem (isomorphism in
$\mathcal K$ between a bisemisimplicial object and its diagonal:
\textsc{Dold\sisi{ }{-}Puppe} \cite[2.9
\sisi{pg.213}{p.~213}]{sga4.exp.XVII.bib.DP}).

In the derived category, this construction can already be made from
an exact sequence of complexes satisfying
\hyperref[XVII.par5.5.13.1]{5.5.13.1}.

${}^{\ast}$ In more intrinsic terms, let $\Delta$ be a ``\emph{true
triangle}'' of objects of $D^{b,\tor\leq0}(X,\Aa)$, \ie an object of
the ``filtered derived category'' with $\Gr^i(\Delta)=0$ for
$i\neq0,1$ and $\Gr^i(\Delta)\in\Ob D^{b,\tor\leq0}(X,\Aa)$
(\textsc{Deligne}, in the Appendix to \textsc{Verdier} \cite{sga4.exp.XVII.bib.V2}, or
\refvb[VI B nᵒ]{Vbis.subsubsec2.5.1},
\refnoexpose{Vbis.subsubsec2.5.2}).

The preceding construction then gives an object
$L\Gamma^k_{\ext}(\Delta)$ of the filtered derived category of
$\Sym^k_S(X)$, isomorphisms
\[
\Gr^i(L\Gamma^k_{\ext}(\Delta))\simeq
L\Gamma^i_{\ext}(\Gr^1(\Delta))\overset{\LL}{\vee}
L\Gamma^{k-i}_{\ext}(\Gr^0(\Delta)\sisi{}{)}
\]
and, denoting by $\lvert\quad\rvert$ the underlying complex of a
filtered complex, an isomorphism
\[
\lvert L\Gamma^k_{\ext}(\Delta)\rvert\simeq
L\Gamma^k_{\ext}(\lvert\Delta\rvert){\quoi}.{}_{\ast}
\]

\begin{enonce*}{Lemma 5.5.16}\etiquette[5.5.16]{XVII.lem5.5.16}
Let $N$ be an integer. The category $D^{b,\tor\leq0}(X,\Aa)$ is
obtained from the category of complexes
$K\in\Ob K^{b,\geq0}(X,\Aa)$ satisfying conditions {\rm(i)} and
{\rm(ii)} below by inverting quasi-isomorphisms:
\begin{enumerate}
\renewcommand{\theenumi}{\roman{enumi}}
\renewcommand{\labelenumi}{(\theenumi)}
\item condition\pageoriginale
(\hyperref[XVII.par5.5.13.1]{5.5.13.1}).

\item The complex $K$ is adapted to every direct-image functor of
cohomological dimension $\leq N$, relative to an isomorphism of topoi
$f:\widetilde X_{\text{\rm et}}\rightarrow T$.
\end{enumerate}
\end{enonce*}

The existence of a calculus of right fractions follows from
\hyperref[XVII.subsubsec5.5.13]{5.5.13} and from the fact that, by
truncated canonical flasque resolutions
(\hyperref[XVII.subsubsec4.2.9]{4.2.9}), every complex $K$ satisfying
(i) maps to a complex satisfying (i) and (ii).

\setcounter{subsubsection}{16}


% END INCLUDED SGA4_English_current_buildable_source_20260906/source/17/drafts/SGA4_XVII_LINES_5401_5800_English.texfrag

% BEGIN INCLUDED SGA4_English_current_buildable_source_20260906/source/17/drafts/SGA4_XVII_LINES_5801_6200_English.texfrag
\subsubsection{}\etiquette[5.5.17]{XVII.subsubsec5.5.17}
Let $u:X\rightarrow Y$ be a morphism of $S$-schemes. Suppose that
$\Sym^n_S(X)$ and $\Sym^n_S(Y)$ exist, that there is an integer $N$
such that the functor $u_{\ast}$ has dimension $\leq N$ on the category
of $\Aa$-modules, and that the functor $Ru_{\ast}$ carries bounded
complexes of \tTor-dimension $\leq0$ to complexes of \tTor-dimension
$\leq0$. The latter two conditions hold automatically if $u$ is
proper, $\Aa$ is torsion, and $Y$ is coherent (\ie quasi-compact and
quasi-separated) (\hyperref[XVII.cor5.2.8.1]{5.2.8.1},
\hyperref[XVII.thm5.2.10]{5.2.10}).\nde{the coherence of $Y$ may be
replaced by the weaker condition that the dimensions of the fibers of
$u$ are bounded.}

Let $K\in\Ob K^{b,\geq0}(X,\Aa)$ satisfy conditions (i) and (ii) of
\hyperref[XVII.lem5.5.16]{5.5.16}, with $N$ as above. By
\hyperref[XVII.subsubsec5.5.14]{5.5.14}, the arrow
$$
\beta:L\Gamma^n_{\ext}(K)\isomto\Gamma^n_{\ext}(K)
$$
is a quasi-isomorphism. Similarly, the arrow
$$
\gamma:u_{\ast}K\isomto Ru_{\ast}K
$$
is a quasi-isomorphism.

Define\pageoriginale the \emph{symmetric Künneth arrow} from
$L\Gamma^n_{\ext}Ru_{\ast}(K)$ to
$R\Sym^n(u)_{\ast}L\Gamma^n_{\ext}(K)$ to be the composite
\begin{multline*}
k^{n}: L\Gamma^{n}_{\ext}Ru_{\ast}K\xleftarrow[\gamma]{\sim}
L\Gamma^{n}_{\ext}u_{\ast}K\rightarrow
\Gamma^{n}_{\ext}u_{\ast}K\xrightarrow[\protect\hyperlink{sga4.semantic.XVII.eq5.5.12.1}{5.5.12.1}]{}\Sym^{n}_{S}(u)_{\ast}\Gamma^{n}_{\ext}(K)\\
\longrightarrow R
\Sym^{n}_{S}(u)_{\ast}\Gamma^{n}_{\ext}(K)\xleftarrow[\beta]{\sim} R \Sym^{n}_{S}(u)_{\ast}L\Gamma^{n}_{\ext}(K){\quoi}.\tag{5.5.17.1}
\end{multline*}\etiquette[5.5.17.1]{XVII.eq5.5.17.1}%The vertical space following this paragraph is too large.

It follows from \hyperref[XVII.lem5.5.16]{5.5.16} that the morphisms
\eqhyperref[XVII.eq5.5.17.1]{5.5.17.1} define a morphism between
functors from $D^{b,\tor\leq0}(X,\Aa)$ to
$D^{+}(\Sym^n_S(Y),\Aa)$:
$$\etiquette[5.5.17.2]{XVII.eq5.5.17.2}
k^n:L\Gamma^n_{\ext}Ru_{\ast}\longrightarrow
R\Sym^n_S(u)_{\ast}L\Gamma^n_{\ext}{\quoi},\leqno(5.5.17.2)
$$
also called the \emph{symmetric Künneth morphism}.

This morphism is characterized by Proposition
\hyperref[XVII.prop5.5.18]{5.5.18} below.

\setcounter{paragraph}{2}
\paragraph{}\etiquette[5.5.17.3]{XVII.par5.5.17.3}
Let $S$, $X$, $Y$, and $u$ be as above, let
$K\in\Ob C^{b,\geq0}(X,\Aa)$ be a complex of \tTor-dimension
$\leq0$, let $K_1\in\Ob C^{b,\geq0}(Y,\Aa)$, and let
$K_2\in\Ob C^b(\Sym^n_S(X),\Aa)$. Let $v$ be a $u$-homomorphism
\hyperlink{sga4.exp.XVII.subsection.0.2}{(0.2)} from $K$ to $K_1$, and let $t$ be a homomorphism (in the usual
sense, as are all those that follow) from $\Gamma^n_{\ext}(K)$ to
$K_2$. There are morphisms
\begin{enumerate}
\renewcommand{\theenumi}{\alph{enumi}}
\renewcommand{\labelenumi}{\theenumi)}
\itemsep=8pt
\item $\Gamma^n_{\ext}(v):\Gamma^n_{\ext}(K_1)\rightarrow
\Sym^n_S(u)_{\ast}\Gamma^n_{\ext}(K)$\qquad
\eqhyperref[XVII.eq5.5.12.0]{5.5.12.0}

\item
$\Sym^n_S(u)_{\ast}(t):\Sym^n_S(u)_{\ast}\Gamma^n_{\ext}(K)
\rightarrow\Sym^n_S(u)_{\ast}K_2$,
\end{enumerate}
whose composite
$$
\varphi':\Gamma^n_{\ext}(K_1)\longrightarrow
\Sym^n_S(u)_{\ast}K_2{\quoi},
$$
defines, by composition,
$$
\varphi:L\Gamma^n_{\ext}(K_1)\longrightarrow
R\Sym^n(u)_{\ast}K_2{\quoi}.
$$
Moreover,\pageoriginale $v$ and $t$ induce
$$
\begin{cases}
v':K_1\longrightarrow Ru_{\ast}K\\
t':L\Gamma^n_{\ext}K\longrightarrow K_2{\quoi}.
\end{cases}
$$

Then:

\begin{enonce*}{Proposition 5.5.18}\etiquette[5.5.18]{XVII.prop5.5.18}
Under the hypotheses of \hyperref[XVII.par5.5.17.3]{5.5.17.3}, the
following diagram is commutative in $D^{+}(\Sym^n_S(Y),\Aa)$:
\[
\xymatrix@C=4cm@R=1.75cm{
  L\Gamma^n_{\ext}Ru_{\ast}K \ar[r]^{k^n}_{\text{\textsuperscript{(symmetric Künneth)}}}& R\Sym^n_S(u)_{\ast}L
\Gamma^n_{\ext}K \ar[d]_{R
  \Sym^n_S(u)_{\ast}(t')}\\
L\Gamma^n_{\ext}K_1 \ar[r]^{\varphi}\ar[u]^{L\Gamma^n_{\ext}(v')} & R\Sym^n_S(u)_{\ast}K_2.
}
\]
\end{enonce*}

The proof is entirely similar to that of
\hyperref[XVII.par5.4.1.5]{5.4.1.5} and is left to the reader.

\setcounter{subsubsection}{18}
\subsubsection{Henceforth $\Aa$ is assumed to be
  torsion}\etiquette[5.5.19]{XVII.subsubsec5.5.19}

Let $u:X\rightarrow Y$ be a compactifiable morphism of $S$-schemes:
\[
\xymatrix@C=3cm@R=1.75cm{
   & \bar{X} \ar[rd]^{\bar{u}}&\\
X \ar@{^{(}->}[ru]^{j} \ar[rr]^{u} \ar[rd]&         & Y \ar[ld]\\
   &   S     & .
}
\]
Suppose\pageoriginale that $\Sym^n_S(X)$,
$\Sym^n_S(\bar X)$, and $\Sym^n_S(Y)$ exist. For every module $F$ on
$X$, one then has
$$\etiquette[5.5.19.1]{XVII.eq5.5.19.1}
\Gamma^n_{\ext}(j_{!}F)\simeq
\Sym^n_S(j)_{!}\Gamma^n_{\ext}(F){\quoi},\leqno(5.5.19.1)
$$
and this isomorphism passes immediately to the derived categories
\sisi{}{(\hyperref[XVII.par5.5.14.4]{5.5.14.4}~(vii'))}. Define the
\emph{symmetric Künneth morphism} between functors from
$D^{b,\tor\leq0}(X,\Aa)$ to
$D^{b,\tor\leq0}(\Sym^n_S(Y),\Aa)$\nde{How can one show that
$\Sym^n_S(u)$ is of finite type in general? This condition is
necessary for $R\Sym^n_S(u)_!$ to be defined.}
$$\etiquette[5.5.19.2]{XVII.eq5.5.19.2}
k^n_u\quad\mbox{\rm or}\quad
k^n:L\Gamma^n_{\ext}Ru_{!}\longrightarrow
R\Sym^n_S(u)_{!}L\Gamma^n_{\ext}\leqno(5.5.19.2)
$$
to be the composite
\begin{align*}
L\Gamma^n_{\ext}Ru_{!}\stackrel{\dfn}{\simeq}L\Gamma^n_{\ext}R\bar u_{\ast}j_{!}
&\xrightarrow{\eqhyperref[XVII.eq5.5.17.2]{5.5.17.2}}
R\Sym^n_S(\bar u)_{\ast}L\Gamma^n_{\ext}j_!\\
&\xrightarrow[\sisi{}{\text{(\hyperref[XVII.par5.5.14.4]{5.5.14.4} (vii'))}}]{\sisi{\substack{\eqhyperref[XVII.eq5.5.19.1]{5.5.19.1}\\\sim}}{\sim}}
R\Sym^n_S(\bar u)_*\Sym^n_S(j)_{!}L\Gamma^n_{\ext}\\
&\xrightarrow{\sisi{\text{(5.5.13.5 (vii'))}}{}}
R\Sym^n_S(u)_{!}L\Gamma^n_{\ext}.
\end{align*}

\setcounter{subsubsection}{19}
\subsubsection{}\etiquette[5.5.20]{XVII.subsubsec5.5.20}
We leave to the reader the following verifications, which are easy in
principle using \hyperref[XVII.prop5.5.18]{5.5.18}.

\paragraph{}\etiquette[5.5.20.1]{XVII.par5.5.20.1}
(Intrinsic character). The morphism
\eqhyperref[XVII.eq5.5.19.2]{5.5.19.2} is independent of the choice of
the compactification $\bar u$ of $u$.

\paragraph{}\etiquette[5.5.20.2]{XVII.par5.5.20.2}
(Transitivity). Let there be a diagram of $S$-schemes
$$
X\xrightarrow{u}Y\xrightarrow{v}Z{\quoi},
$$
with\pageoriginale $u$ and $v$ quasi-projective; suppose that
$\Sym^n_S(Z)$ exists (and hence also $\Sym^n_S(Y)$ and $\Sym^n_S(X)$
(\SGA 1 V 1.7, \EGA II 4.5.4)). The following diagram of functors is
then commutative:
$$
\xymatrix@C=1.5cm@R=1.75cm{
L\Gamma^n_{\ext}R(vu)_{!} \ar@{=}[d] \ar[rr]^{k^n_{vu}}& & R\Sym^n_S(vu)_{!}L\Gamma^n_{\ext}\ar@{=}[d]\\
L\Gamma^n_{\ext}Rv_{!}Ru_{!} \ar[r]^{k^n_v}& R\Sym^n_S(v)_!
L\Gamma^n_{\ext}Ru_{!} \ar[r]^{k^n_u}& R\Sym^n_S(v)_{!}R\Sym^n_S(u)_{!}L\Gamma^n_{\ext}.
}
$$

\paragraph{}\etiquette[5.5.20.3]{XVII.par5.5.20.3}
The symmetric Künneth morphism is compatible with changes of base or
space of the types \hyperref[XVII.subsubsec5.5.9]{5.5.9}~(iv), (v), and (vi). More precisely, one has
the following compatibilities with the morphisms
\sisi{5.5.13.5}{\hyperref[XVII.par5.5.14.4]{5.5.14.4}}
${\rm(iv')}$, ${\rm(v')}$, and ${\rm(vi')}$.
\begin{itemize}
\item[${\rm(iv'')}$] For every $\Aa$-algebra $\Aa'$ and every
quasi-projective $S$-morphism $u:X\rightarrow Y$, the diagram
\[\def\objectstyle{\scriptstyle}
\xymatrix@C=.8cm@R=1.75cm{
L\Gamma^n_{\ext}Ru_{!}(K)\overset{\LL}{\otimes}_{\Aa}\Aa'
\ar@{-}[r]^{\sim} \ar[d]^{k^n_u}&
L\Gamma^n_{\ext}(Ru_{!}(K)\overset{\LL}{\otimes}_{\Aa}\Aa')
\ar@{->}[r]^{\sim}&
L\Gamma^n_{\ext}Ru_{!}(K\overset{\LL}{\otimes}_{\Aa}\Aa') \ar[d]^{k^n_u}\\
R\Sym^n_S(u)_{!}L\Gamma^n_{\ext}(K)\overset{\LL}{\otimes}_{\Aa}\Aa'
\ar@{->}[r]^{\sim}&
R\Sym^n_S(u)_{!}(L\Gamma^n_{\ext}(K)\overset{\LL}{\otimes}_{\Aa}\Aa')
\ar@{-}[r]^{\sim}&
R\Sym^n_S(u)_{!}L\Gamma^n_{\ext}(K\overset{\LL}{\otimes}_{\Aa}\Aa')
}
\]
is commutative.

\item[${\rm(v'')}$] For\pageoriginale every base change
$g:S'\rightarrow S$, giving rise to a diagram
\sisi{cartesian}{with cartesian squares}
$$
\xymatrix{
X' \ar[rrr]^{g'}\ar[rd]^{u'}\ar[dd]&   & & X \ar[rd]^{u}\ar[dd]|\hole &   \\
   & Y' \ar[ld]\ar[rrr]^{g''}&  & & Y \ar[ld]\\
S' \ar[rrr]_{g}&  &  & S &   ,
}
$$
the diagram of functors, entirely analogous to ${\rm(iv'')}$, is
commutative:
$$\def\objectstyle{\scriptstyle}
\xymatrix@R=2.25cm{
\Sym^n_g(g'')^{\ast}L\Gamma^n_{\ext}Ru_{!} \ar@{-}[r]^{\sim} \ar[d]^{k^n_u}&
L\Gamma^n_{\ext}g''^{\ast}Ru_{!} \ar[r]^{\sim}& L\Gamma^n_{\ext}Ru'_{!}
g'^{\ast}\ar[d]^{k^n_{u'}} \\
\Sym^n_g(g'')^{\ast}R\Sym^n_S(u)_{!}L\Gamma^n_{\ext} \ar[r]^{\sim}& R
\Sym^n_S(u')_{!}\Sym^n_g(g')^{\ast}L\Gamma^n_{\ext} \ar@{-}[r]^\sim& R\Sym^n_S(u')_{!}L\Gamma^n_{\ext}g'^{\ast}.
}
$$

\item[${\rm(vi'')}$] For every \emph{radicial} base change
$g:Y'\rightarrow Y$, giving rise to a cartesian diagram of $S$-schemes
$$
\xymatrix@C=2.5cm@R=1.5cm{
X' \ar[r]^{g'}\ar[d]^{u'}& X \ar[d]^{u}\\
Y' \ar[r]^{g}& Y,
}
$$
the following diagram of functors is commutative:
\[
\xymatrix{
\Sym^n_S(g)^{\ast}L\Gamma^n_{\ext}Ru_{!} \ar@{-}[r]^{\sim} \ar[d]^{k^n_u}&
L\Gamma^n_{\ext}g^{\ast}Ru_{!} \ar[r]^{\sim}& L\Gamma^n_{\ext}
Ru'_{!}g'^{\ast} \ar[d]^{k^n_{u'}}\\
\Sym^n_S(g)^{\ast}R\Sym^n_S(u\sisi{_!)}{)_!L\Gamma^n_{\ext}} \ar[r]^{\sim}& R
\Sym^n_S(u')_{!}{g'}^{\ast}L\Gamma^n_{\ext} \ar@{-}[r]^{\sim}& R\Sym^n_S(u')_{!}L\Gamma^n_{\ext}g'^{\ast}
}
\]
\end{itemize}

\begin{enonce*}{Theorem 5.5.21}\etiquette[5.5.21]{XVII.thm5.5.21}
{\rm(Symmetric\pageoriginale Künneth formula).} Let $S$ be a coherent
scheme (= quasi-compact and quasi-separated) ringed by a sheaf $\Aa$
of commutative torsion rings whose stalks are noetherian rings. Ring
each $S$-scheme by the inverse image of $\Aa$. Let $u:X\rightarrow Y$
be a morphism of quasi-projective $S$-schemes. For every\sisi{}{ }
integer $n\geq0$ and every $K\in\Ob D^{b,\tor\leq0}(X,\Aa)$, the
symmetric Künneth arrow \eqhyperref[XVII.eq5.5.17.2]{5.5.17.2}
$$
k^n:L\Gamma^n_{\ext}Ru_{!}K\longrightarrow
R\Sym^n_S(u)_{!}L\Gamma^n_{\ext}K
$$
is an isomorphism in
$D^{b,\tor\leq0}(\Sym^n_S(Y),\Aa)$.
\end{enonce*}

The proof of this theorem will occupy the rest of \hyperref[XVII.subsec5.5]{5.5}.

\begin{enonce*}[remark]{Remark 5.5.21.1}\etiquette[5.5.21.1]{XVII.rem5.5.21.1}
The quasi-projectivity hypothesis serves only to ensure that $u$
admits a compactification $\bar u$
$$
X\xrightarrow{j}\bar X\xrightarrow{\bar u}Y
$$
such that the schemes $\Sym^n_S(X)$, $\Sym^n_S(\bar X)$, and
$\Sym^n_S(Y)$ exist for every $n$. The language of algebraic spaces
should make it possible to replace this hypothesis by the sole
hypothesis that $u$ is compactifiable.
\end{enonce*}

The\pageoriginale first part of the proof consists in reducing
\hyperref[XVII.thm5.5.21]{5.5.21}, by dévissage, to an
``irreducible'' case:

\begin{enonce*}{Lemma 5.5.22}\etiquette[5.5.22]{XVII.lem5.5.22}
Theorem \hyperref[XVII.thm5.5.21]{5.5.21} follows from the following
special case:
\begin{itemize}
\item[${(\ast)}$] For every algebraically closed field $k$, every
smooth projective curve $X$ over $k$, and every prime number $\ell$,
the symmetric Künneth arrows (obtained for
$Y=S=\Spec(k)$, $\Aa=\ZZ/\ell$, $K=\ZZ/\ell$; see the explanations
below)
$$\etiquette[5.5.22.1]{XVII.eq5.5.22.1}
L\Gamma^n(R\Gamma(X,\ZZ/\ell))\longrightarrow
R\Gamma(\Sym^n_k(X),\ZZ/\ell)\leqno(5.5.22.1)
$$
are isomorphisms $(n\geq0)$.
\end{itemize}
\end{enonce*}

In \sisi{\hyperref[XVII.eq5.5.22.1]{5.5.22.1}}{\eqhyperref[XVII.eq5.5.22.1]{5.5.22.1}},
the symbol $\Gamma$ has an unfortunate double use: in $R\Gamma$ it
denotes the ``global sections'' functor, whereas in $L\Gamma^n$ it
denotes the functor of \hyperref[XVII.par5.5.2.3]{5.5.2.3}. The
isomorphism $\Gamma^n_{\ext}(\ZZ/\ell)\simeq\ZZ/\ell$ is used in this
formula.

Let $S$, $X$, $Y$, $u$, and $\Aa$ be as in
\hyperref[XVII.thm5.5.21]{5.5.21}.

\begin{enonce*}{Lemma 5.5.22.2}\etiquette[5.5.22.2]{XVII.lem5.5.22.2}
Let there be an exact sequence of complexes in
$C^{b,\geq0}(X,\Aa)$, of \ttor-dimension $\leq0$:
$$
0\longrightarrow K'\longrightarrow K\longrightarrow K''\longrightarrow0{\quoi}.
$$
If all the Künneth arrows $k^n$ $(n\geq0)$ relative to two of the
complexes $K'$, $K$, and $K''$ are isomorphisms, then the Künneth
arrows relative to the third are also isomorphisms.
\end{enonce*}

We\pageoriginale reduce to the case where $u$ is projective and the
complexes $K'$, $K$, and $K''$ satisfy conditions (i) and (ii) of
\hyperref[XVII.lem5.5.16]{5.5.16}, for $N$ at least twice the fiber
dimension of $u$. The symmetric Künneth arrows are then given by
\eqhyperref[XVII.eq5.5.17.1]{5.5.17.1}.

We prove by induction on $n$ that if, for $i\leq n$ and two of the
complexes $K'$, $K$, and $K''$, the arrows $k^i$ are isomorphisms,
then the arrows $k^i$ $(i\leq n)$ are isomorphisms for the third.

To see this, note that by
\sisi{\hyperref[XVII.subsubsec5.5.14]{5.5.14}}{\hyperref[XVII.subsubsec5.5.15]{5.5.15}},
the morphism
\begin{equation*}
k^n:L\Gamma^n_{\ext}Ru_{\ast}K\longrightarrow
R\Sym^n(u)_{\ast}L\Gamma^n_{\ext}K\tag{a}
\end{equation*}
can be represented by a morphism of filtered complexes
$$
{k'}^n:K_0\longrightarrow K_1
$$
with finite filtration; on the successive quotients this morphism
induces arrows isomorphic to the composite arrows
\begin{equation*}
L\Gamma^i_{\ext}Ru_{\ast}K'\overset{\LL}{\vee}
L\Gamma^{n-i}_{\ext}Ru_{\ast}K''\longrightarrow
R\Sym^n(u)_{\ast}(L\Gamma^i_{\ext}K'\overset{\sisi{}{\LL}}{\vee}
L\Gamma^{n-i}_{\ext}K'')\tag{$\text{b}_i$}
\end{equation*}
as follows:
\begin{gather*}
L\Gamma^i_{\ext}Ru_{\ast}K'\overset{\LL}{\vee}
L\Gamma^{n-i}_{\ext}Ru_{\ast}K''\xrightarrow{k^i\overset{\LL}{\vee}k^{n-i}}
R\Sym^i(u)_{\ast}L\Gamma^i_{\ext}K'\overset{\LL}{\vee}
R\Sym^{n-i}(u)_{\ast}L\Gamma^{n-i}_{\ext}(K'')\\
\underset{\defn}{=}
\vee_{i,n-i{\ast}}(R\Sym^i(u)_{\ast}L\Gamma^i_{\ext}K'\overset{\LL}{\boxtimes}
R\Sym^{n-i}(u)_{\ast}L\Gamma^{n-i}_{\ext}(K''))\\
\xrightarrow{\substack{\hyperref[XVII.thm5.4.3]{5.4.3}\\\sim}}
\vee_{i,n-i{\ast}}\sisi{(R\Sym^i(u)_{\ast}\times R\Sym^{n-i}(u)_{\ast})}{(R\Sym^i(u)\times R\Sym^{n-i}(u))_\ast}
(L\Gamma^i_{\ext}K'\overset{\LL}{\boxtimes}
L\Gamma^{n-i}_{\ext}K'')\\
=R\Sym^n(u)_{\ast}\vee_{i,n-i{\ast}}(L\Gamma^i_{\ext}K'\overset{\LL}{\boxtimes}L\Gamma^{n-i}_{\ext}K'')
\overset{\dfn}{=}R\Sym^n(u)_{\ast}(L\Gamma^i_{\ext}K'\overset{\LL}{\vee}L\Gamma^{n-i}_{\ext}K''){\quoi}.
\end{gather*}

By\pageoriginale the induction hypothesis (and the Künneth formula),
these arrows are quasi-isomorphisms for $i\neq0,n$. Any two of the
following assertions therefore imply the third: (a) is a
quasi-isomorphism; $(\text{b}_0)$ is a quasi-isomorphism;
$(\text{b}_n)$ is a quasi-isomorphism. These are precisely the three
arrows in question, \cqfd

\setcounter{subsubsection}{22}
\setcounter{paragraph}{2}
\paragraph{}\etiquette[5.5.22.3]{XVII.par5.5.22.3}
Assume assertion $(\ast)$ of \hyperref[XVII.lem5.5.22]{5.5.22} and
prove \hyperref[XVII.thm5.5.21]{5.5.21}. Let $S$, $X$, $Y$, $\Aa$,
$K$, and $u$ therefore be as in \hyperref[XVII.thm5.5.21]{5.5.21}.
The problem is local on $S$, which may be assumed affine. Every complex
$K$ is a successive extension of complexes of the form $i_{!}i^{\ast}K$
for $i$ the inclusion in $X$ of a locally closed affine subset. By
dévissage (\hyperref[XVII.lem5.5.22.2]{5.5.22.2}), we may therefore
suppose that $K$ has the form $i_{!}K_1$, for an immersion
$i:X_1\rightarrow X$ with affine source. Compatibility
\hyperref[XVII.par5.5.20.2]{5.5.20.2} permits us to replace $X$ by
$X_1$ and $K$ by $K_1$, and hence to suppose that $u$ is affine.
Factor $u$ as a composite $u=u_1\ldots u_k$ of morphisms with fibers
of dimension $\leq1$. Compatibility
\hyperref[XVII.par5.5.20.2]{5.5.20.2} shows that it suffices to prove
that the arrows $k^n$ relative to the $u_i$ and to any complex $K$ are
isomorphisms. We are thus reduced to the case where $S$ is
\emph{affine and $u$ has relative dimension $\leq1$}.

\paragraph{}\etiquette[5.5.22.4]{XVII.par5.5.22.4}
For $k^n$ to be an isomorphism, it suffices that its restrictions to
the geometric fibers of $\Sym^n_S(Y)$ be isomorphisms; compatibility
\hyperref[XVII.par5.5.20.3]{5.5.20.3} ${\rm(v'')}$ therefore reduces
us to the case where $S$ is the \emph{spectrum of an algebraically
closed field $k$}.

\paragraph{}\etiquette[5.5.22.5]{XVII.par5.5.22.5}
For\pageoriginale every rational point $y$ of $\Sym^n_k(Y)$, there is
a finite reduced subscheme $Y'$ of $Y$ such that $\sisi{x}{y}$ belongs
to the image of $\Sym^n_k(Y')$. Compatibility
\hyperref[XVII.par5.5.20.3]{5.5.20.3} ${\rm(vi'')}$ therefore permits
us to suppose that $Y$ \emph{is a disjoint union of finitely many
copies of} $\Spec(k)$.

\paragraph{}\etiquette[5.5.22.6]{XVII.par5.5.22.6}
If $Y=\amalg Y_i$, then every complex $K$ on $X$ is a sum of complexes
supported on one of the schemes $X_i=u^{-1}(Y_i)$. By dévissage
(\hyperref[XVII.lem5.5.22.2]{5.5.22.2}) and compatibility
\hyperref[XVII.par5.5.20.3]{5.5.20.3} ${\rm(vi'')}$, it suffices to
prove \hyperref[XVII.thm5.5.21]{5.5.21} for each morphism
$u_i:X_i\rightarrow Y_i$. By \hyperref[XVII.par5.5.22.5]{5.5.22.5},
this reduces us to the case where
$$
\text{\rm $Y=S=\Spec(k)$, $k$ algebraically closed, $\dim X\leq1$.}
$$

The sheaf $\Aa$ now reduces to a noetherian ring $A$. The following
proposition will allow us to reduce to the case where $A$ is a field.

\begin{enonce*}{Proposition 5.5.23}\etiquette[5.5.23]{XVII.prop5.5.23}
Let $A$ be a (commutative) noetherian ring, let $K\in\Ob D^{-}(A)$ be
a bounded-above complex of $A$-modules, and let $n$ be an integer. The
following conditions are equivalent:

% END INCLUDED SGA4_English_current_buildable_source_20260906/source/17/drafts/SGA4_XVII_LINES_5801_6200_English.texfrag

% BEGIN INCLUDED SGA4_English_current_buildable_source_20260906/source/17/drafts/SGA4_XVII_LINES_6201_6600_English.texfrag
\begin{enumerate}
\renewcommand{\theenumi}{\roman{enumi}}
\renewcommand{\labelenumi}{(\theenumi)}
\item $K$ has \tTor-dimension $\leq n$

\item For every $x\in \Spec (A)$ and every $i<-n$,
$H^{i}(K\overset{\LL}{\underset{A}{\otimes}} k(x))=0$.
\end{enumerate}
\end{enonce*}

One has (i) $\Rightarrow$ (ii) trivially\sisi{,}{.} Suppose (ii). To prove
(i), we prove by noetherian induction on the closed subset $F$ of
$\Spec(A)$ that if a module $N$ satisfies $N_{x}=0$ for $x\not\in F$,
then $H^{i}(K\overset{\LL}{\otimes}_{A}N)=0$
for $i<-n$.

Let\pageoriginale $I$ be the unique ideal defining $F$ such that $A/I$
is reduced, and let $N_{k}$ be the submodule of $N$ annihilated by
$I^{k}$. One has $N=\varinjlim N_{k}$. Since the Tor functors
commute with filtered direct limits, we reduce to the case where there
exists $k$ such that $N=N_{k}$, so that $N$ is a successive extension
of modules annihilated by $I$; the long exact cohomology sequence
reduces us to the case where $N$ is already an $A/I$-module.

Let $\eta$ be a maximal point of $F$, and define $A/I$-modules
$M_{1}$ and $M_{2}$ by the exact sequence
$$\etiquette[5.5.23.1]{XVII.eq5.5.23.1}
0\longrightarrow M_{1}\longrightarrow N \xrightarrow{i} N\otimes
k(\eta)\longrightarrow M_{2} \longrightarrow 0{\quoi}.\leqno(5.5.23.1)
$$

One has $(M_{j})_{\eta}=0$ $(j=1,2)$. If, for every closed subset $G\subset F$
with $\eta\not\in G$, we denote by $M_{j,G}$ the submodule of
$M_{j}$ consisting of the sections of $M_{j}$ supported in $G$, then
$$
M_{j}=\varinjlim_{\eta\not\in G} M_{j,G}{\quoi}.
$$
By the induction hypothesis and passage to the limit, one therefore has
$$
H^{i}(K\overset{\LL}{\otimes}_{A} M_j)=0\quad\mbox{\rm
for}\quad i< -n{\quoi}.
$$

The module $N\otimes k(\eta)$ is a vector space over $k(\eta)$, hence a
direct sum of copies of $k(\eta)$. By hypothesis, one has
$$
H^{i}(K\overset{\LL}{\otimes} (N\otimes
k(\eta))\sisi{}{)}=0\quad\mbox{\rm for}\quad i<-n{\quoi}.
$$

If $I$ is the image of the arrow $i$ in \eqhyperref[XVII.eq5.5.23.1]{5.5.23.1}, there are
exact sequences
\begin{align*}
& H^{i-1}(K\overset{\LL}{\otimes}M_{2})\longrightarrow
H^{i}(K\overset{\LL}{\otimes}I)\longrightarrow
H^{i}(K\overset{\LL}{\otimes}(N\otimes k(\eta)))\\
& H^{i}(K\overset{\LL}{\otimes}M_{1})\longrightarrow
H^{i}(K\overset{\LL}{\otimes}N)\longrightarrow
H^{i}(K\overset{\LL}{\otimes}I){\quoi},
\end{align*}
from which\pageoriginale it follows that
$H^{i}(K\overset{\LL}{\otimes}N)=0$ for $i<-n$,
\cqfd

\begin{enonce*}{Corollary 5.5.24}\etiquette[5.5.24]{XVII.cor5.5.24}
For a morphism $u:K\rightarrow L$ in $D^{-}(A)$ to be an
isomorphism, it suffices that for every $x\in \Spec(A)$, the morphism
$$
u_{x}:K\overset{\LL}{\otimes}_{A}k(x)\longrightarrow
L\overset{\LL}{\otimes}_{A}k(x)
$$
be an isomorphism in $D^{-}(k(x))$.
\end{enonce*}

\sisi{\noindent}{}Let $M$ be the cone (mapping cylinder) of $u$. By hypothesis, one has
$M\overset{\LL}{\otimes}_{A}k(x)=0$ for every
$x\in\Spec(A)$. By \hyperref[XVII.prop5.5.23]{5.5.23}, one therefore has $M=0$, and $u$ is an
isomorphism.

\setcounter{subsubsection}{25}
\setcounter{paragraph}{0}
\paragraph{}\etiquette[5.5.25.1]{XVII.par5.5.25.1}
Let us continue the proof of \hyperref[XVII.lem5.5.22]{5.5.22}, interrupted
at \hyperref[XVII.prop5.5.23]{5.5.23}. By the compatibility
\hyperref[XVII.par5.5.20.3]{5.5.20.3} ${\rm(iv'')}$ and
\hyperref[XVII.cor5.5.24]{5.5.24}, the morphism $k^{n}$
associated with $K\in \Ob D^{b,\tor\geq 0}(X,A)$ is an isomorphism if and
only if the morphisms $k^{n}$ associated with
$K\overset{\LL}{\otimes}_{A}k(x)\in \Ob
D^{b,\tor\geq 0}(X,k(x))$ are isomorphisms, for $x\in
\Spec(A)$. This reduces us to the case where $A$ \emph{is a field of
characteristic} $\ell > 0$.

\paragraph{}\etiquette[5.5.25.2]{XVII.par5.5.25.2}
The exact sequence (where $C$ denotes a cone)
$$
0\longrightarrow K[-1]\longrightarrow CK[-1]\longrightarrow
K\longrightarrow 0
$$
shows that the morphisms $k^{n}$ are isomorphisms for $K$ if and
only if they are so for $K[-1]$ (\hyperref[XVII.lem5.5.22.2]{5.5.22.2}).

\paragraph{}\etiquette[5.5.25.3]{XVII.par5.5.25.3}
The\pageoriginale exact sequences
$$
0\longrightarrow \sigma_{\geq n} K\longrightarrow K\longrightarrow
\sigma_{\leq n} K\longrightarrow 0
$$
show that, for the Künneth morphisms to be isomorphisms when
$K$ is a flat complex, it suffices that they be so for the components
of $K$, placed in degree $0$ (\hyperref[XVII.lem5.5.22.2]{5.5.22.2}\sisi{)}{} and \sisi{(}{}\hyperref[XVII.par5.5.25.2]{5.5.25.2}).

\paragraph{}\etiquette[5.5.25.4]{XVII.par5.5.25.4}
Thus let $\sisi{F}{K}$ be a sheaf of $A$-vector spaces. This sheaf is a filtered
direct limit of constructible sheaves (\hyperrefext[IX][IX.cor2.7.2]{2.7.2}), automatically
flat since $A$ is a field, and, by passage to the limit, we
reduce to the case where $K$ is \emph{constructible} and concentrated in degree zero.

\paragraph{}\etiquette[5.5.25.5]{XVII.par5.5.25.5}
If a complex $K$ is a direct summand of a complex $L$, the morphism
$k^{n}$ associated with $K$ is a direct summand of the morphism $k^{n}$ associated
with $L$, hence is an isomorphism if the latter is one. The trace
method (\hyperrefext[IX][IX.rem5.8]{5.8}) therefore applies, and it suffices to treat the case where $F$
is of the form $j_{!}A_{X'}$, for $j:X'\rightarrow X$ quasi-finite. The
compatibility \hyperref[XVII.par5.5.20.2]{5.5.20.2} allows us to replace $X$ by $X'$, whence the
reduction
$$
\mbox{$X$ \emph{has dimension $\leq 1$ and} $K=A_{X}${\quoi}.}
$$

\paragraph{}\etiquette[5.5.25.6]{XVII.par5.5.25.6}
We may assume that $X$ is reduced. If $j:U\rightarrow X$ is the smooth
locus of $X$, with complement $i:F\rightarrow X$, the exact sequence
$$
0\longrightarrow j_{!}j^{\ast}A_{X}\longrightarrow
A_{X}\longrightarrow i_{!}i^{\ast}A_{X}\longrightarrow 0
$$
reduces us by dévissage to the cases where $X$ \emph{is either a smooth
curve or a finite set.} Moreover, the formula
$$
A_{X}=A\overset{\LL}{\otimes}_{\ZZ/\ell}(\ZZ/\ell)_{X}
$$
reduces us by \hyperref[XVII.par5.5.20.3]{5.5.20.3} ${\rm(iv'')}$ to the case where $A_{X}=\ZZ/\ell$.

If\pageoriginale $X$ is a curve, we embed $X$ in the corresponding
smooth projective curve, and a further dévissage reduces us to
one of the following two cases:
\begin{enumerate}
\renewcommand{\theenumi}{\roman{enumi}}
\renewcommand{\labelenumi}{(\theenumi)}
\item $X$ is a smooth projective curve, $Y=S=\Spec(k)$ with $k$
algebraically closed, $\Aa=\ZZ/\ell$, $K=\ZZ/\ell$.

\item $X$ is a finite disjoint union of copies of $\Spec(k)$,
$Y=S=\Spec(k)$ ($k$ algebraically closed),
$\Aa=\ZZ/\ell$, $K=\ZZ/\ell$.
\end{enumerate}

A final dévissage of $\sisi{}{(}\ZZ/\ell\sisi{}{)}_{X}$ into the sheaves
$\sisi{}{(}\ZZ/\ell\sisi{}{)}_{X\sisi{}{'}}$, for $X'$ a connected component of $X$, reduces us in
(i) and (ii) to assuming that $X$ is connected. Case (ii) is then
trivial. This proves \hyperref[XVII.lem5.5.22]{5.5.22}.

\begin{enonce*}{Lemma 5.5.26}\etiquette[5.5.26]{XVII.lem5.5.26}
Theorem \emph{\hyperref[XVII.thm5.5.21]{5.5.21}} follows from the following two special cases
of assertion $(\ast)$ of \emph{\hyperref[XVII.lem5.5.22]{5.5.22}}:
\begin{itemize}
\item[\emph{Case} 1.] Assume that $k=\CC$.

\item[\emph{Case} 2.] Assume that $k$ has characteristic $p>0$
and that $\ell=p$.
\end{itemize}
\end{enonce*}

With the notation of \hyperref[XVII.lem5.5.22]{5.5.22}, denote by $p$ the characteristic of
$k$. If $p=0$, the Lefschetz principle (relying here on
compatibility \hyperref[XVII.par5.5.20.3]{5.5.20.3} ${\rm(v'')}$ and on XII \hyperref[XII.cor5.4]{5.4}) allows us to
reduce to the case where $k=\CC$.

It remains to treat the case where $p>0$ and $\ell\neq p$. The curve $X$
lifts to a curve $X'$ over the ring $W(k)$ of Witt vectors
over $k$ (\SGA 1 III 7.4). Put $S=\Spec(W(k))$, and consider the diagram
$$
\xymatrix@C=2.25cm@R=1.5cm{
X'_{\eta} \ar[d]^{f_\eta}          \ar[r]              & X' \ar[d]^{f'}& X\ar[l]\ar[d]^{f}\\
\eta               \ar[r]         & S   & \Spec(k), \ar[l],
}
$$
and\pageoriginale let $f_{n}$, $f'_{n}$, and $f_{\eta n}$ be the
respective projections of $\Sym^{n}_{k}(X)$, $\Sym^{n}_{S}(X')$, and
$\Sym^{n}_{\eta}(X'_{\eta})$ onto $\Spec(k)$, $S$, and $\eta$. It is known
that $f'_{n}$ is a smooth projective morphism (projectivity is
clear; for smoothness, the simplest approach is to reduce to the case where
$X'=\PP^{1}_{S}$, in which case $\Sym^{n}_{S}(X')\simeq
\PP^{n}$, as follows from the theorem on elementary symmetric
functions). By the specialization theorem \hyperrefext[XVI][XVI.cor2.2]{2.2}, the
sheaves $R^{i}f'_{n{\ast}}$ $(\ZZ/\ell)$ are locally
constant on $S$. The same holds for the sheaves
$L^{i}\Gamma^{n}(Rf'_{\ast}\ZZ/\ell)$. To verify that the
morphism of \eqhyperref[XVII.eq5.5.22.1]{5.5.22.1} associated with $X$ over $k$ is an isomorphism, it
therefore suffices, by \hyperref[XVII.par5.5.20.3]{5.5.20.3} ${\rm(v'')}$, to prove the analogous statement
for the geometric generic fiber $X_{\bar{\eta}}$ of $X$ over
$S$. The field $k(\bar{\eta})$ has characteristic $0$, and, in view of the
first reduction, this proves \hyperref[XVII.lem5.5.26]{5.5.26}.

To prove \hyperref[XVII.thm5.5.21]{5.5.21} in the two special cases of
\hyperref[XVII.lem5.5.26]{5.5.26}, we
compare the symmetric Künneth morphism on the one hand with its
transcendental analogue (\hyperref[XVII.subsubsec5.5.27]{5.5.27} and
\hyperref[XVII.subsubsec5.5.28]{5.5.28}) and on the other hand with its coherent analogue
(\hyperref[XVII.subsubsec5.5.29]{5.5.29} through \hyperref[XVII.lem5.5.37]{5.5.37}).

\setcounter{subsubsection}{26}
\subsubsection{}\etiquette[5.5.27]{XVII.subsubsec5.5.27}
If $p:X\rightarrow S$ is a continuous map of topological
spaces, with $S$ ringed by a sheaf of commutative rings
$\Aa$, the constructions of \hyperref[XVII.subsubsec5.5.7]{5.5.7} and
\hyperref[XVII.subsubsec5.5.8]{5.5.8} apply unchanged
to define a functor $\Gamma^{n}_{\ext}$ from modules on
$X$ to modules on $\Sym^{n}_{S}(X)=X^{n}/\Sf_{n}$. The
formal arguments of \hyperref[XVII.subsubsec5.5.7]{5.5.7} through \hyperref[XVII.prop5.5.18]{5.5.18}
carry over. Likewise, if $S$, $X$, and $Y$
are locally compact Hausdorff topological spaces
$$
\xymatrix@C=1.5cm@R=1.25cm{
X \ar[rr]^{u}\ar[rd]&  &  Y\ar[ld]\\
            & S &,
}
$$
and\pageoriginale if the fibers of $u$ have bounded cohomological
dimension, the formalism of \hyperref[XVII.subsubsec5.5.19]{5.5.19} and
\hyperref[XVII.subsubsec5.5.20]{5.5.20} for the symmetric Künneth morphism in
cohomology with proper supports carries over.

If $p:X\rightarrow S$ is a morphism of schemes of finite type over
$\CC$, the functors $\Gamma^{n}$ and $L\Gamma^{n}$ commute with the
functors of ``passage to the transcendental setting''. Finally, if $X$, $Y$, $S$,
$\Aa$, and $K$ are as in \hyperref[XVII.thm5.5.21]{5.5.21}, with $S$ separated and of finite type
over $\CC$, and if we denote generically by $\epsilon$ the
morphisms of sites of type $\epsilon:X_{\cl}\rightarrow
X_{\et}$, the following diagram is commutative:
$$
\vcenter{\etiquette[5.5.27.1]{XVII.eq5.5.27.1}
\xymatrix@C=1.75cm@R=1.5cm{
\epsilon^{\ast} L\Gamma^{n}_{\ext} Ru_{!}K  \ar[r]^{k^{\sisi{k}{n}}} \ar[d]&
\epsilon^{\ast} R \Sym^{n}_{S}(u)_{!} L \Gamma^{n}_{\ext} K \ar[d]\\
L\Gamma^{n}_{\ext} R u^{\an}_{!} \epsilon^{\ast}K  \ar[r]^{k^{n \an}}
& R \Sym^{n}_{S}(u^{\an})_{!} L\Gamma^{n}_{\ext} \epsilon^{\ast}K.
}}\leqno{(5.5.27.1)}
$$

\subsubsection{}\etiquette[5.5.28]{XVII.subsubsec5.5.28}
By the compatibility \eqhyperref[XVII.eq5.5.27.1]{5.5.27.1} and an elementary
case of the comparison theorem
(\hyperrefext[XI][XI.thm4.4]{4.4} or \hyperrefext[XVI][XVI.thm4.1]{4.1}),
to treat Case 1 of \hyperref[XVII.lem5.5.26]{5.5.26},
it suffices to verify that if $X$ is a smooth complex projective
curve, then the topological analogue of \eqhyperref[XVII.eq5.5.22.1]{5.5.22.1} is an
isomorphism.

The topological space $X^{\an}$ is triangulable (this fact is
trivial here, since $X$ is a branched covering of the Riemann
sphere). The topological analogue of the dévissage lemma \hyperref[XVII.lem5.5.22.2]{5.5.22.2}
then reduces us to proving
\begin{itemize}
\item[$(\ast)$] If\pageoriginale $D$ is a closed
simplex (of dimension 0, 1, or 2), the topological symmetric Künneth
morphism
$$\etiquette[5.5.28.1]{XVII.eq5.5.28.1}
k^{n}:L\Gamma^{n} R\Gamma(D,\ZZ/\ell)\longrightarrow R\Gamma(\Sym^{n}(D\sisi{}{)},\ZZ/\ell)\leqno(5.5.28.1)
$$
is an isomorphism.
\end{itemize}

One has
$$
R\Gamma(D,\ZZ/\ell)=\ZZ/\ell\;\mbox{\rm placed in degree }
0{\quoi},
$$
and hence
$$
L\Gamma^{n}R\Gamma(D,\ZZ/\ell)=\ZZ/\ell\; \mbox{\rm
placed in degree $0$}{\quoi}.
$$

On the other hand, $\Sym^{n}(D)$ is contractible, and again one has
$$
R\Gamma(\Sym^{n}(D), \ZZ/\ell)=\ZZ/\ell\;\mbox{\rm placed
in degree $0$}{\quoi}.
$$

The $H^{i}$ of the two sides of \eqhyperref[XVII.eq5.5.28.1]{5.5.28.1} therefore vanish in nonzero degrees, and the morphism
induced on $H^0$ identifies with the identity of
$\ZZ/\ell$.

\subsubsection{}\etiquette[5.5.29]{XVII.subsubsec5.5.29}
For the application we have in mind, it is useful to regard a
quasi-coherent sheaf on a scheme $X$ as a
quasi-coherent sheaf of $\Oo$-modules on the étale site of $X$.

Let $p:X\rightarrow \sisi{X}{S}$ be a quasi-projective morphism of schemes (cf.
\hyperref[XVII.rem5.5.21.1]{5.5.21.1} \footnote[1]{Everything that follows, including
  \hyperref[XVII.prop5.5.34]{5.5.34}, remains valid
under the sole assumption that $X$ and $Y$ admit invertible sheaves
relatively ample over $S$ (without assuming that $X$ and $Y$ are of
finite type over $S$).}). If $ℱ$ is a sheaf of $\Oo_{X}$-modules
on $X_{\text{\rm et}}$, then
\begin{enumerate}
\renewcommand{\theenumi}{\alph{enumi}}
\renewcommand{\labelenumi}{\theenumi)}
\item We denote by $\bigotimes\limits^{n}ℱ$ the
sheaf of modules on $(X/S)^{n}_{\text{\rm et}}$ given by the tensor product
of the inverse-image $\Oo_{(X/S)^n}$-modules of $ℱ$ on
$(X/S)^{n}_{\et}$;

\item the sheaf $\bigotimes\limits^{n}ℱ$ is a
$\mathcal{G}_n$-equivariant sheaf. Passing to invariants, it
defines a sheaf $TS^{n}_{\text{\rm ext}}(ℱ)$ on $\Sym^{n}_{S}(X)$:
if $\sigma$ is the projection of $(X/S)^{n}$ onto $\Sym^{n}_{S}(X)$, one has
$$\etiquette[5.5.29.1]{XVII.eq5.5.29.1}
TS^{n}_{\ext}(ℱ)=(\sigma_{\ast}\bigotimes^{n}ℱ)^{\Sf^{n}}{\quoi}.\leqno(5.5.29.1)
$$
\end{enumerate}

\emph{Assume\pageoriginale that $p$ is flat}; as in
\hyperref[XVII.subsubsec5.5.8]{5.5.8}, one then defines the functor
$\Gamma^{m}_{\text{\rm ext}}$ to be the $0$\sisi{th}{\textsuperscript{th}} left derived
functor of $TS^{n}_{\ext}$. With the notation of \hyperref[XVII.par5.5.2.5]{5.5.2.5} and
\hyperref[XVII.subsubsec5.5.8]{5.5.8}, if
$$
P_{1}\longrightarrow P_{0}\longrightarrow ℱ\longrightarrow 0
$$
is a flat presentation of $ℱ$, then the sequence
$$
\xymatrix@C=1.75cm@M=.25cm{
TS^{n}_{\ext} (P'_1) \ar@<3pt>[r]^{j_0} \ar@<-3pt>[r]_{j_1}& TS^{n}_{\ext}(P_0)\ar[r] & \Gamma^{n}_{\ext}(\sisi{F}{ℱ})
}
$$
is exact. If $x$ is a geometric point of $\Sym^{n}_{S}(X)$ as
in \hyperref[XVII.par5.5.7.1]{5.5.7.1}, one has (\cf \eqhyperref[XVII.eq5.5.8.1]{5.5.8.1})
\begin{align*}
& \Gamma^{n}_{\text{\rm
ext}}(ℱ)_{x}=\Oo_{x}\otimes_{C}(\bigotimes^{\ell}_{1}
\Gamma_{\Oo_{S}}^{n_{i}}(\sisi{F}{ℱ}_{y_{i}})){\quoi},\tag{5.5.29.2}\\
& C=\bigotimes^{\ell}_{1} \Sym^{n_{i}}_{\Oo_{S}}(\Oo_{y_{i}}){\quoi}.\tag*{where}
\end{align*}\etiquette[5.5.29.2]{XVII.eq5.5.29.2}%Removing this comment would introduce unwanted white space.
Indeed, it suffices to verify this formula in the trivial case where
the fibers of $\sisi{F}{ℱ}$ are free. Note that $\Oo_{x}$ is flat over
$C$, being the strict henselization of $C$ at the closed point that is the image of $x$.

\subsubsection{}\etiquette[5.5.30]{XVII.subsubsec5.5.30}
Assume that $X$ and $S$ are affine, $p$ is flat, and $ℱ$ is quasi-coherent. Let
$S=\Spec(A)$, $X=\Spec(B)$, and let $F=H^0(X,ℱ)$, so that
$ℱ$ is the sheaf $F^{\sim}$ defined by the $B$-module $F$. It follows
from \eqhyperref[XVII.eq5.5.29.1]{5.5.29.1} that
$$
TS^{n}_{\ext}(ℱ)=TS^{n}_{A}(F)^{\sim}{\quoi},
$$
and, on deriving, one therefore has
$$\etiquette[5.5.30.1]{XVII.eq5.5.30.1}
\Gamma^{n}_{\text{\rm
ext}}(ℱ)=\Gamma^{n}_{A}(F)^{\sim}{\quoi}.\leqno(5.5.30.1)
$$
The\pageoriginale functor $\Gamma^{n}_{\text{\rm ext}}$ thus appears
as a functor $\Gamma^{n}$ ``relative to $A$'', as
\eqhyperref[XVII.eq5.5.29.2]{5.5.29.2} already expressed locally.

\subsubsection{}\etiquette[5.5.31]{XVII.subsubsec5.5.31}
We denote by $D^{b,\tor_{S}\leq 0}(X,\Oo_{X})$ the subcategory
of $D^{b}(X,\Oo_{X})$ formed by the complexes of $\Oo_{X}$-modules
which, as complexes of $p^{-1}(\Oo_{S})$-modules, have
\tTor-dimension $\leq 0$.

Consider the condition

\paragraph{}\etiquette[5.5.31.1]{XVII.par5.5.31.1}
$K\in \Ob C^{b,\geq 0}(X,\sisi{O}{\Oo}_{X})$, and, at every geometric point $x$
of $X$, $K_{x}$ is homotopic to a complex of $\Oo_{X,x}$-modules, in
nonnegative degrees, flat over $\Oo_{S,p(x)}$.

Since the morphism $p$ is flat, so that there are ``enough''
$\Oo_{X}$-modules flat over $S$, one checks as in \hyperref[XVII.lem5.5.13.2]{5.5.13.2} that
$D^{b,\tor_{S}\leq 0}(X,\Oo_{X})$ is obtained from the category of
complexes $K\in \Ob K^{b,\geq 0}(X,\Oo_{X})$ satisfying (\hyperref[XVII.par5.5.31.1]{5.5.31.1}) by
inverting quasi-isomorphisms.

It follows from \eqhyperref[XVII.eq5.5.29.2]{5.5.29.2} and \hyperref[XVII.lem5.5.6.1]{5.5.6.1} that if $u:K\rightarrow L$ is a
quasi-isomorphism between complexes satisfying (\hyperref[XVII.par5.5.31.1]{5.5.31.1}), then
$\Gamma^{n}_{\text{\rm ext}}(u)$ (in the sense of \textsc{Dold-Puppe}) is again a
quasi-isomorphism. The functor $\Gamma^{n}_{\text{\rm ext}}$ can
therefore be derived to give
$$\etiquette[5.5.31.2]{XVII.eq5.5.31.2}
L\Gamma^{n}_{\text{\rm ext}}:D^{b,\tor_{S}\leq
0}(X,\Oo_{X})\longrightarrow D^{b,\tor_{S}\leq 0}(\Sym^{n}_{S}(X), \Oo_{\Sym^{n}_{S}(X)})\leqno(5.5.31.2)
$$

\subsubsection{}\etiquette[5.5.32]{XVII.subsubsec5.5.32}
We denote by $D^{b,\tor_{S}\leq 0}_{\text{\rm qcoh}}(X,\Oo_{X})$
the subcategory of $D^{b,\tor_{S}\leq 0}(X,\Oo_{X})$ formed by
complexes with quasi-coherent cohomology.

% END INCLUDED SGA4_English_current_buildable_source_20260906/source/17/drafts/SGA4_XVII_LINES_6201_6600_English.texfrag

% BEGIN INCLUDED SGA4_English_current_buildable_source_20260906/source/17/drafts/SGA4_XVII_LINES_6601_7000_English.texfrag

Standard\pageoriginale arguments make it possible, as in
\hyperref[XVII.subsubsec5.5.17]{5.5.17}, to
define a symmetric Künneth arrow of the following type:

Let $u:X\rightarrow Y$ be a morphism of quasi-projective $S$-schemes
flat over $S$; for $K\in \Ob D^{b,\tor_{S}\leq 0}_{\text{\rm
qcoh}}(X,\Oo_{X})$, the symmetric Künneth arrow $k^{n}$ in
$D^{b,\tor_{S}\leq 0}_{\text{\rm qcoh}}(\Sym^{n}_{S}(Y),\Oo)$ is
an arrow
$$\etiquette[5.5.32.1]{XVII.eq5.5.32.1}
k^{n}:L\Gamma^{n}_{\text{\rm ext}}Ru_{\ast}K\longrightarrow
R\Sym^{n}_{S}u_{\ast}L\Gamma^{n}_{\text{\rm ext}} K{\quoi}.\leqno(5.5.32.1)
$$

We leave it to the reader to verify the formal properties of
\eqhyperref[XVII.eq5.5.32.1]{5.5.32.1} that will be used, notably the following one:

\begin{enonce*}{Lemma 5.5.33}\etiquette[5.5.33]{XVII.lem5.5.33}
If $S$ is of characteristic $p>0$, and if $u:X\rightarrow S$ is
proper and flat, the diagram of symmetric Künneth arrows
\eqhyperref[XVII.eq5.5.17.2]{5.5.17.2} and \eqhyperref[XVII.eq5.5.32.1]{5.5.32.1} is commutative ($Y=S, K=\sisi{\sheafZZ}{\ZZ}/p$ or
$\Oo_{X}$):
\[
\xymatrix{
L\Gamma^{n}_{\sisi{\sheafZZ}{\ZZ}/p}Ru_{\ast}(\sisi{\sheafZZ}{\ZZ}/p)
\ar[r]^{k^n} \ar[d]&R
\Sym^{n}_S(u)_{\ast} (\sisi{\sheafZZ}{\ZZ}/p) \ar[d]&\\
L\Gamma^{n}_{\Oo} Ru_{\ast}(\Oo) \ar[r]^{k^n}& R
\Sym^{n}_{S}(u)_{\ast}\Oo &.
}
\]
\end{enonce*}


\begin{enonce*}{Proposition 5.5.34}\etiquette[5.5.34]{XVII.prop5.5.34}
Under the hypotheses of \hyperref[XVII.subsubsec5.5.32]{5.5.32}, the symmetric Künneth arrow
\eqhyperref[XVII.eq5.5.32.1]{5.5.32.1} is an isomorphism.
\end{enonce*}

\begin{enumerate}
\renewcommand{\theenumi}{\alph{enumi}}
\renewcommand{\labelenumi}{\theenumi)}
\item The question is local on $S$, which may be assumed affine.

\item For every point $x$ of $\Sym^{n}_{S}(Y)$, there exists an affine
open $Y_{1}$ of $Y$ such that $x$ lies in $\Sym^{n}_{S}(Y_{1})$. By
localization, we may therefore assume that $Y$ is affine.

\item Every\pageoriginale complex $K$ of the type under consideration can be
represented by a complex $K_{0}$ of quasi-coherent modules on
$S$, zero in degrees $<0$. If $\mathcal{U}$ is an affine covering
of $X$, then $K$ admits a finite resolution by sums of
complexes of the form $j_{\ast}j^{\ast}K$, where $j$ is the inclusion of an
affine open (alternating Čech complex). Using
the quasi-coherent analogue of the dévissage lemma \hyperref[XVII.lem5.5.22.2]{5.5.22.2}, this
reduces us to the case where there exist an affine open $j:X_{1}\hookrightarrow X$
of $X$ and $K_{1}\in D^{b,\tor_{S}\leq 0}_{\text{\rm
qcoh}}(X_{1},\Oo)$ such that $K=Rj_{\ast}K_{1}$.

\item In view of the quasi-coherent analogue of transitivity \hyperref[XVII.par5.5.20.2]{5.5.20.2}, it
suffices to prove \hyperref[XVII.prop5.5.34]{5.5.34} for $uj$ and $j$. Reduction b), applied
to $j$, finally reduces us to the case where $X$, $Y$, and $S$ are affine. We
may then represent $K$ by a complex $K_{2}\in C^{b,\geq
0}(X,\Oo_{X})$ with \sisi{}{$S$-}flat quasi-coherent components (\sisi{\quad}{\hyperref[XVII.def4.1.9]{4.1.9}}). The
functors $Ru_{\ast}$ and $L\Gamma^{n}$ can then be computed without
having to resolve $K_{2}$ further, and it remains to note that, by \eqhyperref[XVII.eq5.5.30.1]{5.5.30.1},
for $ℱ$ quasi-coherent on $X$, we have
\[
\Gamma^{n}_{\text{\rm ext}} u_{\ast}ℱ\isomto
\Sym^{n}_{S}(u)_{\ast} \Gamma^{n}_{\text{\rm ext}}ℱ{\quoi},
\]
the two sides being identified with
\[
\Gamma^{n}_{A}(F)^{\sim} \quad\mbox{\rm for}\quad
ℱ=F^{\sim}{\quoi},
\]
\cqfd
\end{enumerate}

\setcounter{subsubsection}{34}
\subsubsection{}\etiquette[5.5.35]{XVII.subsubsec5.5.35}
Let us treat case $2$ of \hyperref[XVII.lem5.5.26]{5.5.26}.
With the notation of \hyperref[XVII.lem5.5.22]{5.5.22}, the
following diagram is commutative (\hyperref[XVII.lem5.5.33]{5.5.33}):
$$\vcenter{\etiquette[5.5.35.1]{XVII.eq5.5.35.1}
\xymatrix@C=2cm{
  L\Gamma^{n} (R\Gamma(X,\ZZ/p)) \ar[r]^{k^n}_*+[o][F-]{\scriptstyle{1}}
  \ar[d]& R\Gamma(\Sym^{n}_{k}(X),\ZZ/p)\ar[d] \\
L\Gamma^{n}(R\Gamma(X,\Oo)) \ar[r]^{k^n}_*+[o][F-]{\scriptstyle{2}} & R\Gamma(\Sym^{n}_{k}(X),\Oo),
}}\leqno{(5.5.35.1)}
$$
and\pageoriginale by \hyperref[XVII.prop5.5.34]{5.5.34}, arrow \ding{193} is an
isomorphism. The Frobenius endomorphism $F:x\rightarrow x^{p}$ of
$\Oo$ acts $p$-linearly on $R\Gamma(X,\Oo)$ and on
$R\Gamma(\Sym^{n}_{k}(X),\Oo)$. It therefore also acts $p$-linearly
on $L\Gamma^{n}(R\Gamma(X,\Oo))$. One checks that $k^{n}$ and $F$
commute.

Artin--Schreier theory provides exact sequences
$$
0\longrightarrow H^{i}(X,\ZZ/p)\longrightarrow
H^{i}(X,\Oo)\xrightarrow{F-I} H^{i}(X,\Oo)\longrightarrow
\sisi{\Oo}{0}{\quoi},
$$
and similarly for $\Sym^{n}_{S}(X)$. To deduce that \ding{192} is
an isomorphism from the fact that \ding{193} is one, it remains to
verify the following

\begin{enonce*}{Lemma 5.5.36}\etiquette[5.5.36]{XVII.lem5.5.36}
Let $k$ be a perfect field of characteristic $p>0$, let $K\in C^{\geq
0}(\ZZ/p)$ be a complex of finite-dimensional $\ZZ/p$-vector spaces,
let $L\in C^{\geq 0}(k)$ be a complex of finite-dimensional $k$-vector spaces,
let $F$ be a $p$-linear endomorphism of $L$, and let
$u:K\rightarrow L$ be a morphism. Suppose that
$H^{\ast}(Fu)=H^{\ast}(u)$, and that the sequences
$$
0\longrightarrow H^{i}(K)\longrightarrow
H^{i}(L)\xrightarrow{F-I}H^{i}(L)\longrightarrow 0
$$
are exact. Then\pageoriginale the sequences
$$
0\longrightarrow H^{i}(\Gamma^{n}K)\longrightarrow
H^{i}(\Gamma^{n}L)\xrightarrow{F-I}H^{i}(\Gamma^{n}L)\longrightarrow 0
$$
are exact.
\end{enonce*}

Since $H^{\ast}(Fu)=H^{\ast}(u)$ and $\ZZ/p$ is a field,
there exists a homotopy
$$
Fu-u=dH+Hd{\quoi}.
$$
Since the field $k$ is perfect, we can write
$$
FH_{0}-H_{0}=H{\quoi}.
$$
Replacing $u$ by $u-(dH_{0}+H_{0}d)$, we may assume that
$Fu=u$.

The functor $S$ that assigns to every $k$-vector space $V$ equipped with a
$p$-linear endomorphism $F$ the $\ZZ/p$-vector space $V^{S}$ of those $v\in
V$ such that $Fv=v$ is an exact functor. The hypothesis therefore means
that the morphism
$$
K\longrightarrow L^{S}
$$
is a quasi-isomorphism. It remains to prove that the functor $S$ commutes
with the functor $\Gamma^{n}$. This follows readily from the following well-known
lemma, whose proof is left to the reader:

\begin{enonce*}{Lemma 5.5.37}\etiquette[5.5.37]{XVII.lem5.5.37}
A finite-dimensional vector space $V$ over a perfect field $k$, equipped
with a $p$-linear endomorphism $F$, admits a unique
decomposition
$$
V=V_{0}+V_{n}
$$
with $V^{S}\otimes k\isomto  V_{0}$, $F(V_{n})\subset V_{n}$,
and $F|V_{n}$ nilpotent.
\end{enonce*}

This completes the proof of \hyperref[XVII.thm5.5.21]{5.5.21}.

\section[The functor $f_{!}$]{The functor $f_{!}$}\etiquette[6]{XVII.sec6}\pageoriginale

We propose to give a direct construction of the functor $f_{!}$
of \hyperref[XVII.var5.1.11]{5.1.11}. We will then use this construction to define the ``trace morphism'' (\hyperref[XVII.thm6.2.3]{6.2.3}) for quasi-finite flat morphisms of
finite presentation.

\subsection[New construction of the functor $f_{!}$]{New construction of the functor $f_{!}$}\etiquette[6.1]{XVII.subsec6.1}

Let $f:X\rightarrow S$ be a morphism locally of finite type. A
subset $P$ of $X$ is said to be \emph{proper} over $S$ if it is the image
of a subscheme of $X$ proper over $S$.

\begin{enonce*}{Proposition 6.1.1}\etiquette[6.1.1]{XVII.prop6.1.1}
Let $f:X\rightarrow S$ be a morphism locally of finite type and
quasi-separated.
\begin{enumerate}
\renewcommand{\theenumi}{\roman{enumi}}
\renewcommand{\labelenumi}{(\theenumi)}
\item A subset of $X$ proper over $S$ is locally closed, and is
closed if $f$ is separated.

\item If $f$ is separated, the union of two subsets of $X$ proper over
$S$ is proper over $S$.

\item For every $S$-scheme $U$, let $\Phi_{X}(U)$ denote the set
of subsets of $\fprod{U}{X}{S}$ proper over $U$. The functor
$\Phi_{X}$ is a sheaf for the \emph{fpqc} topology. If $f$ is
separated, $\Phi\sisi{}{_X}$ is a sheaf of filtered ordered sets for the
inclusion relation.
\end{enumerate}
\end{enonce*}

Assertions\pageoriginale (i) and (ii) are trivial, and the second
assertion of (iii) follows from (ii). The functor $U\mapsto
\mathcal{P}(\fprod{U}{X}{S})$ is a sheaf for every topology
whose coverings are surjective. It follows from \EGA
IV 2.3.14 and \EGA IV 2.7.1 (vii) that $\Phi_{X}$ is a
subsheaf thereof.

If $\Ff $ is an étale sheaf of pointed sets on a
scheme $X$, recall (IV \hyperref[IV.subsubsec8.5.2]{8.5.2}) that the support of a section $s$ of
$\Ff $ is the complement of the largest open set on which $s=0$,
where $0$ denotes the distinguished section.

\begin{enonce*}{Definition 6.1.2}\etiquette[6.1.2]{XVII.def6.1.2}
Let $f:X\rightarrow S$ be a separated morphism locally of finite type,
and let $\Ff $ be a sheaf of pointed sets on $X$. We
denote by $f_{!}\Ff $ the subsheaf of
$f_{\ast}\Ff $ whose sections over $U$ (étale over $S$) are
the sections of $\Ff $ over $\fprod{U}{X}{S}$ with support proper
over $U$.
\end{enonce*}

It follows from \hyperref[XVII.prop6.1.1]{6.1.1} that $f_{!}\Ff $ is indeed a
sheaf.

One readily verifies that:

\begin{enonce*}{Lemma 6.1.3}\etiquette[6.1.3]{XVII.lem6.1.3}
If $f$ is the composite $gh$ of two separated morphisms locally of
finite type, there exists a unique isomorphism of functors
$f_{!}\simeq g_{!}h_{!}$ making the following diagram commutative
$$
\xymatrix@C=1cm@R=1.25cm@M=.2cm{
f_{!}  \ar@{-}[r]^{\sim} \ar@{^{(}->}[d]& g_{!}h_{!} \ar@{^{(}->}[d]\\
f_{\ast} \ar@{=}[r]      & g_{\ast} h_{\ast}.
}
$$
\end{enonce*}

If\pageoriginale $f$ is an open immersion (\resp a proper morphism),
the functor $f_{!}$ is the extension-by-zero functor
(\resp the functor $f_{\ast}$). By \hyperref[XVII.lem6.1.3]{6.1.3}, definitions
\hyperref[XVII.def6.1.2]{6.1.2} and \hyperref[XVII.var5.1.11]{5.1.11} therefore agree on their
common domain of validity.

The properties of the functor $f_{!}$ verified in \hyperref[XVII.sec5]{\S 5} remain
valid in the more general case \hyperref[XVII.def6.1.2]{6.1.2}.

\begin{enonce*}{Proposition 6.1.4}\etiquette[6.1.4]{XVII.prop6.1.4}
Let $f:X\rightarrow S$ be a separated morphism locally of finite type.
\begin{enumerate}
\renewcommand{\theenumi}{\roman{enumi}}
\renewcommand{\labelenumi}{(\theenumi)}
\item The functor $f_{!}$ is left exact (\ie commutes with finite
projective limits) and commutes with filtered inductive limits.

\item If $f$ is quasi-finite, the functor $f_{!}$ is an exact functor
from the category of abelian sheaves on $X$ to that of
abelian sheaves on $S$.

\item For every Cartesian diagram
$$
\xymatrix@C=1.5cm@R=1.25cm{
X'          \ar[r]^{g'} \ar[d]^{f'}& X \ar[d]^{f}\\
S'          \ar[r]^{g}  & S
}
$$
there exists a base-change isomorphism making the
diagram
$$
\xymatrix@C=1.5cm@R=1.25cm{
g^{\ast}f_{!}  \sisi{\ar@{=}[r]}{\ar@{-}[r]^\sim} \ar@{^{(}->}[d]& f'_{!}g'^{\ast}\ar@{^{(}->}[d]\\
g^{\ast}f_{\ast}  \ar[r]    & f'_{\ast}g'^{\ast},
}
$$
commute, in which the second horizontal arrow is the base-change
arrow \emph{\neweqhyperrefext[XII][XII.eq4.1]{4.1}}.
\end{enumerate}
\end{enonce*}

\begin{enonce*}[remark]{Proof}
For\pageoriginale every $U\in \Ob S_{\et}$, we have
\[
f_{!}\Ff (U)=\varinjlim_{P\in
\Phi(U)}\Gamma_{\sisi{p}{P}}(\Ff ){\quoi},
\]
and for $U$ quasi-compact and quasi-separated, the functor
$\Ff \mapsto f_{!}\Ff (U)$ commutes with finite
projective limits and filtered inductive limits, as a filtered
inductive limit (\hyperref[XVII.prop6.1.1]{6.1.1}\sisi{)}{} (iii)\sisi{}{)} of functors having these properties.

This proves (i). Assertion (ii) follows at once from (iii) (take
the geometric points of $S$ as base changes).

Assertion (iii) is local on $S$, which may be assumed
affine. Let $I$ be the filtered ordered set of open subsets of $X$
of finite type over $S$. If
$U\overset{j}{\hookrightarrow}X$
is an open subset of $X$,
$(fj)_{!}(\Ff |U)=f_{!}(j_{!}(\Ff |U))$ maps to
$f_{!}(\Ff\sisi{|U}{})$. Moreover
$$
f_{!}(\Ff )=\varinjlim_{U\in I}
(fj\sisi{}{)}_{!}(\Ff |U){\quoi}.
$$
This allows us to assume that $X$ is of finite type over $S$ and to apply
Chow's lemma (\EGA II 2\sisi{nd}{\textsuperscript{nd}} ed.)\nde{see [C, 2.6], cited in \nderef{XVII.ndeConrad}.} to find a projective and
surjective morphism $p^{0}:X^{0}\rightarrow X$ such that $fp\sisi{_{0}}{^0}$
is compactifiable. Let $X^{1}=\fprod{X^{0}}{X^{0}}{X}$ and let $p^{1}$
be the natural arrow from $X^{1}$ to $X$. The sequence
$$
0\rightarrow \Ff \rightarrow
p^{0}_{\ast}p^{0{\ast}}\Ff \rightrightarrows
p^{1}_{\ast}p^{1{\ast}}\Ff
$$
is exact (\sisi{VII 9.2}{\hyperrefext[VIII][VIII.cor9.2]{9.2}}).

Taking into account the functoriality of the base-change arrow and
(i), it suffices to verify (iii) for the sheaves
$p^{0}_{\ast}p^{0^{\ast}}\Ff $ and
$p^{1}_{\ast}p^{1^{\ast}}\Ff $. By the change of notation
$(p^{1}\mapsto p^{0})$, we\pageoriginale are thus reduced to assuming
that $\Ff $ is of the form $p^{0}_{\ast}\mathcal{G}$. Assertion
(iii) for the pair $(f,p^{0}_{\ast}\mathcal{G})$ is equivalent to
assertion (iii) for the pairs $(p^{0}_{\ast}\sisi{}{,}\mathcal{G})$ and
$(fp^{0}_{\ast},\mathcal{G})$. We conclude by \hyperref[XVII.var5.2.7]{5.2.7} (i), which applies
because $p^{0}$ and $fp^{0}$ are compactifiable.
\end{enonce*}

\begin{enonce*}[remark]{Counterexample 6.1.6}\etiquette[6.1.6]{XVII.contrexem6.1.6}
The functor $Rf_{!}$ does not in general agree with the derived functor
of the functor $f_{!}$. Indeed, if $X$ is affine over an
algebraically closed field $k$, we have
$$
f_{!}\Ff =\sisi{\sum}{\bigoplus}_{x\in X(k)}H^{0}_{x}(\Ff ){\quoi},
$$
and the $i$\textsuperscript{th} derived functor of $f_{!}$ therefore takes the value
at $\Ff $
$$
\sisi{\sum}{\bigoplus}_{x\in X(k)}H^{i}_{x}(\Ff ){\quoi}.
$$
On the other hand, if $X$ is smooth, irreducible, and of dimension 1,
$$
H^{2}_{c}(X,\mu_{\ell^{n}})\simeq \ZZ/\ell^{n}
$$
for $\ell$ prime to the characteristic of $k$.
\end{enonce*}

\subsection{The trace morphism in relative dimension 0}\etiquette[6.2]{XVII.subsec6.2}

The main result of this ${\rm n^{\circ}}$ is Theorem
\hyperref[XVII.thm6.2.3]{6.2.3}. The reader is advised not to read the technical preliminaries
\hyperref[XVII.subsubsec6.2.1]{6.2.1} and
\hyperref[XVII.prop6.2.2]{6.2.2}. In the proof of
\hyperref[XVII.thm6.2.3]{6.2.3}, these allow one to consider only finite morphisms $f$, a reduction that,
via \EGA IV 18.12\sisi{.}{,} is in any event fairly evident.

\subsubsection{}\etiquette[6.2.1]{XVII.subsubsec6.2.1}
Let\pageoriginale $f:X\rightarrow S$ be a \emph{quasi-finite
separated} morphism. For every $U\in \Ob S^{\et}$, let
$\Psi_{f}(U)$ denote the category of pairs consisting of a finite pointed
set $I$, with distinguished point denoted $0$, and a decomposition of
$f^{\ast}U$ into open and closed subsets $f^{\ast}U=\prod_{i\in
I}V_{i}$, with $V_{i}$ finite over $U$ for $i\neq 0$. A morphism from
$\varphi=(V_{i})_{i\in I}$ to $\varphi'=(V'_{i})_{i\in I'}$ is a
map $\sigma$ from $I'$ to $I$ such that, for $i\in I$, one has
$V_{i}=\bigcup_{j\in \sigma^{-1}(i)}V'_{j}$. We say that $\varphi'$
\emph{refines} $\varphi$ if there exists a morphism from $\varphi$ to
$\varphi'$.

Let $\Ff $ be a sheaf of abelian groups on $S$, let $U\in
\Ob S_{\text{\rm et}}$, let $\varphi=(V_{i})_{i\in I}\in \Ob\Psi_{f}(U)$,
let $I^{\ast}=I\backslash \{0\}$, and let $f_{i}$ be the projection of $V_{i}$ onto
$U$. We have
$$
f_{!}f^{\ast}\Ff |U=f_{0!}f_{0}^{\ast}(\Ff |U)\oplus
\bigoplus_{i\neq 0} f_{i\ast}f_{i}^{\ast}(\Ff |U){\quoi},
$$
hence a composite arrow
$$\etiquette[6.2.1.1]{XVII.eq6.2.1.1}
\tau_{\varphi}:\Ff ^{I^{\ast}}|U\longrightarrow
\bigoplus_{i\neq 0}
f_{i\ast}f_{i}^{\ast}(\Ff |U)\longrightarrow f_{!}f^{\ast}\Ff |U{\quoi}.\leqno(6.2.1.1)
$$

If $\sigma:I\rightarrow J$ is a morphism of pointed sets, we
denote by $\leftexp{t}{\sigma}:\Ff ^{J^{\ast}}\rightarrow
\Ff ^{I^{\ast}}$ the arrow with coordinates $\leftexp{t}{\sigma}^i_{\phantom{i}j}=\Id$
if $\sigma(i)=j$, and $\leftexp{t}{\sigma}^i_{\phantom{i}j}=0$ if $\sigma(i)\neq
j$. If $\sigma:I'\rightarrow I$ is a morphism from $\varphi$ to
$\varphi'$, the diagram
$$
\vcenter{\etiquette[6.2.1.2]{XVII.eq6.2.1.2}
\xymatrix@R=1.25cm{
\Ff^{I^{\ast}}|U \ar[rr]^{\leftexp{t}{\sigma}}   \ar[rd] \ar[rd]_{\tau_{\varphi}}&&
\Ff^{I'^{\ast}}|U     \ar[ld]^{\tau_{\varphi'}}\\
      &          f_{!}f^{\ast}\Tt|U  &
}}\leqno{(6.2.1.2)}
$$
is\pageoriginale commutative. Moreover, ${}^{t}(\sigma\rho)={}^{t}\rho
{}^{t}\sigma$. The $\Psi_{f}(U)$ form a fibered category
$\Psi_{f}$ over $S_{\text{\rm et}}$; the preceding constructions

% END INCLUDED SGA4_English_current_buildable_source_20260906/source/17/drafts/SGA4_XVII_LINES_6601_7000_English.texfrag

% BEGIN INCLUDED SGA4_English_current_buildable_source_20260906/source/17/drafts/SGA4_XVII_LINES_7001_7400_English.texfrag
commute with localization and therefore define an $S_{\text{\rm
et}}$-functor $\varphi=(V_{i})_{i\in I} \in \Psi_{f}(U)\mapsto
\Ff ^{I^{\ast}}|U$ from $\Psi$ to the $S_{\text{\rm
et}}$-category of sheaves on étale open subschemes of $S$, and a
morphism of functors
$\tau_{\varphi}:\Ff ^{I^{\ast}}|U\rightarrow
f_{!}f^{\ast}\Ff |U$.

\begin{enonce*}{Proposition 6.2.2}\etiquette[6.2.2]{XVII.prop6.2.2}
\begin{enumerate}
\renewcommand{\theenumi}{\roman{enumi}}
\renewcommand{\labelenumi}{(\theenumi)}
\item The category $\Psi_{f}$ is locally filtering on the right
\emph{(\sisi{IV 14}{\hyperrefext[V][V.def8.1.1]{8.1.1}})} over $S_{\text{\rm et}}$.

\item The arrows \eqhyperref[XVII.eq6.2.1.1]{6.2.1.1} define an isomorphism\sisi{s}{}
\[
\tau:\varinjlim_{\varphi=(V_{i})_{i\in I}\in
\Psi_{f}(U)}\Ff ^{I^{\ast}}|U\isomto  f_{!}f^{\ast}\Ff\sisi{}{|U} {\quoi}.
\]
\end{enumerate}
\end{enonce*}

If $S$ is strictly local with closed point $s$, the $S$-scheme $X$
admits a decomposition $\varphi\in\Psi_{f}(S)$ into open and closed
parts $X=\coprod_{i\in I}X_{i}$ indexed by a pointed set $I$
such that the fiber of $X_{0}$ at $s$ is empty, the $X_{i}$ $(i\neq
0)$ are finite over $S$, and the fiber of $X_{i}$ $(i\neq 0)$ at $s$
consists of a single point (\EGA IV 18.12.1). Such a decomposition is
a final object of $\Psi\sisi{}{_f}(S)$. Moreover, by \hyperref[XVII.prop6.1.4]{6.1.4} (iii) (or
because this is evident), the fiber at $s$ of the arrow $\tau_{\varphi}$
\eqhyperref[XVII.eq6.2.1.1]{6.2.1.1} is an isomorphism.

In the general case, let $x$ be a geometric point of $S$ and let
$\Psi_{x}$ be the fiber of $\Psi_{f}$ at $x$,
\[
\Psi_{x}=\varinjlim_{U(x)} \Psi_{f}(U){\quoi},
\]
where $U$ ranges over the étale neighborhoods of $x$. Let $S_{x}$ be the strict localization of
$S$ at $x$, $X_{x}=\fprod{X}{S_{x}}{S}$, and $f_{x}:X_{x}\rightarrow
S_{x}$. The standard limit arguments show\pageoriginale that
the category $\Psi_{x}$ is equivalent to the category
$\Phi_{f_{x}}(S_{x})$, hence is filtering on the right. It follows that
$\Psi\sisi{}{_f}$ is locally filtering on the right. Moreover, the fiber at $x$ of
$\tau$ identifies with the arrow
$$
\tau_{x}:\varinjlim_{\Psi_{x}}\Ff _{x}^{I_{x}}\longrightarrow
(f_{!}f^{\ast}\Ff )_{x}
$$
and, by the preceding discussion, this arrow is an isomorphism.

\begin{enonce*}{Theorem 6.2.3}\etiquette[6.2.3]{XVII.thm6.2.3}
There is a unique way to define, for every morphism
$f:X\rightarrow S$ that is separated, flat, of finite presentation, and quasi-finite,
and every abelian sheaf $\Ff $ on $S$, a trace morphism
$$
\Tr_{f}:f_{!}f^{\ast}\Ff \longrightarrow \Ff
$$
such that the following conditions hold:
\begin{itemize}
\item[\textsuperscript{(Var 1)}] $\Tr_{f}$ is functorial in $\Ff $.

\item[\textsuperscript{(Var 2)}] $\Tr_{f}$ is compatible with every base change, \ie
for every cartesian diagram
$$
\xymatrix{
X' \ar[rr]^{g'} \ar[d]^{f'}& & X\ar[d]^{f}\\
S' \ar[rr]^{g} & &  S
}
$$
the following diagram of functors is commutative:
$$
\xymatrix{
f'_{!} f'^{\ast} g^{\ast} \ar@{=}[r] \ar[d]^{\Tr_{f'}\ast g^{\ast}}& f'_{!}g'^{\ast} f^{\ast}\ar@{=}[r] &
g^{\ast}f_{!}f^{\ast} \ar[d]^{g^{\ast}\ast
  \Tr_f}\\
g^{\ast} \ar@{=}[rr]&  & g^{\ast} &.
}
$$

\item[\textsuperscript{(Var 3)}] $\Tr_{f}$\pageoriginale is compatible with
composition, \ie if $f$ is the composite $f=gh$ of two morphisms
that are separated, flat, of finite presentation, and quasi-finite, the following diagram
is commutative:
$$
\xymatrix{
f_{!}f^{\ast}\ar[rrdd]_{\Tr_f} \ar@{=}[rr]& & g_{!}h_{!} h^{\ast}g_{\ast}\ar[d]^{g_{!}\ast\mathord{\Tr_{h}}\ast g^\ast}\\
              & &  g_{!}g^{\ast} \ar[d]^{\Tr_g}\\
              & & \Id .
}
$$

\item[\textsuperscript{(Var 4)}] (Normalization). If $f$ is finite of constant rank $n$,
the composite homomorphism
$$
\Ff \rightarrow
f_{\ast}f^{\ast}\Ff =f_{!}f^{\ast}\Ff \xrightarrow{\Tr_{f}}\Ff
$$
is multiplication by $n$.
\end{itemize}
\end{enonce*}

\sisi{\emph{First note that conditions \emph{(Var 2)} through \emph{(Var 4)}
imply that}}{First note that conditions (Var 2) through (Var 4)
imply that}

\begin{enonce*}{Lemma 6.2.3.1}\etiquette[6.2.3.1]{XVII.lem6.2.3.1}
If $X=X'\amalg X''$\sisi{}{,} then, setting $f'=f|X'$ and $f''=f|X''$, the
morphism $\Tr_{f}$:
$$
f_{!}f^{\ast}\Ff =f'_{!}{f'}^{\ast}\Ff \oplus
f''_{!}{f''}^{\ast}\Ff \longrightarrow \Ff
$$
is the sum of $\Tr_{f!}$ and $\Tr_{f''}$.
$$
\xymatrix@C=2.25cm@R=.75cm{
  X' \ar@{^{(}->}[rd]^{j'}&           &\\
   & X   \ar[r]^{f}  \ar[r]^{f}    & S\\
X'' \ar@{^{(}->}[ru]_{j''}& &.
}
$$
\end{enonce*}

If\pageoriginale one applies (Var 2) to $j'$, for base change by
$j'$, and (Var 4) to the identity map, one sees that
$$
\Tr_{j'}:j'_{!}{j'}^{\ast}f^{\ast}\Ff \longrightarrow
f^{\ast}\Ff
$$
is the canonical map defined by the decomposition of
$f^{\ast}\Ff $ as the direct sum
$$
f^{\ast}\Ff =j'_{!}{j'}^{\ast}f^{\ast}\Ff \oplus
j''_{!}{j''}^{\ast}f^{\ast}\Ff {\quoi}.
$$

If one applies (Var 3) to $f'=f\circ j'$, one sees that the following diagram
is commutative
\[
\xymatrix@R=1.25cm{
f'_{!} f'^{\ast}\Ff  \ar[rr] \ar[rd]_{\Tr_{f'}}& &f_{!}f^{\ast}\Ff \ar[ld]^{\Tr_f}\\
                                   & \Ff    &    ,
}
\]
which, together with the analogous statement for $f''$, is precisely \hyperref[XVII.lem6.2.3.1]{6.2.3.1}.

Let $f:X\rightarrow S$ be a separated, flat, quasi-finite morphism of
finite presentation. For every $U\in \Ob S_{\text{\rm et}}$, let
$\Psi'_{f}(U)$ be the full subcategory of $\Psi_{f}(U)$ (\sisi{\hyperref[XVII.prop6.1.1]{6.1.1}}{\hyperref[XVII.subsubsec6.2.1]{6.2.1}})
consisting of the decompositions $(V_{i})_{i\in I}$ such that, for $i\neq
0$, $V_{i}$ is locally free of constant rank over $U$. The
categories\pageoriginale $\Psi'_{f}(U)$ form a fibered subcategory
of $\Psi_{f}$, cofinal in $\Psi_{f}$.

Let $\Ff $ and $\mathcal{G}$ be abelian sheaves on
$S$. By \hyperref[XVII.prop6.2.2]{6.2.2} (ii), giving a morphism
$g:f_{!}f^{\ast}\Ff \rightarrow \mathcal{G}$ is equivalent to giving,
for each decomposition $\varphi=(V_{i})_{i\in I}\in\Psi'_{f}(U)$,
a morphism $g_{\varphi}:\Ff ^{I^{\ast}}|U\rightarrow
\mathcal{G}|U$, compatible with the morphisms ${}^{t}\sigma$ of
\eqhyperref[XVII.eq6.2.1.2]{6.2.1.2}. If $f_{i}$ is the projection from $V_{i}$ to $U$,
$g_{\varphi}$ is the arrow whose coordinate arrows are
the composites:
\[
g^{i}_{\varphi}:\Ff |U\rightarrow
f_{i\ast}f^{\ast}_{i}\Ff |U\rightarrow
f_{!}f^{\ast}\Ff |U\xrightarrow{g} \mathcal{G}|U{\quoi}.
\]

It follows that there exists a unique arrow $\Tr_{f}$, compatible
with étale localization and satisfying (Var 4) and \hyperref[XVII.lem6.2.3.1]{6.2.3.1}. It is clear
that this arrow satisfies (Var 1) and (Var~2).

To verify (Var 3), by base change one reduces to the
easy case in which $S$ is the spectrum of an algebraically closed field.

\setcounter{subsubsection}{3}
\subsubsection{}\etiquette[6.2.4]{XVII.subsubsec6.2.4}
A \emph{weighting} of a separated quasi-finite morphism
$f:X\rightarrow S$ is the datum, for each geometric point $x$ of
$X$, of an integer $n(x)$ depending only on the point of $X$ underlying $x$,
such that:
\begin{itemize}
\item[$(\ast)$] \emph{For every $U\in \Ob S_{\text{\rm et}}$ and
every open and closed subspace $X_{1}$ of $f^{-1}(U)$ that is finite over $U$, the
function assigning to each geometri\sisi{c}{cal} point $s$ of $U$ the number
$\sum_{x\in X_{1}(s)}n(x)$ is locally constant.}
\end{itemize}

If $f=gh$ is the composite of two weighted morphisms, we weight $f$
by setting $n_{f}(x)=n_{h}(x)\sisi{.}{\cdot}n_{g}(h(x))$.

If\pageoriginale the morphism $f'$ is obtained by a base change
$g$ from a weighted morphism $f$, we weight $f'$ by setting
$n_{f'}(x)=n_{f}(g'(x))$.

The preceding arguments in fact prove:

\begin{enonce*}{Proposition 6.2.5}\etiquette[6.2.5]{XVII.prop6.2.5}
There is a unique way to define, for every weighted, separated,
quasi-finite morphism $f:X\rightarrow S$ and every abelian sheaf
$\Ff $ on $S$, a trace morphism
$\Tr_{f}:f_{!}f^{\ast}\Ff \rightarrow \Ff $ that is
functorial in $\Ff $, compatible with base change and with
composition, and satisfies:
\begin{itemize}
\item[\textsuperscript{(Var $\sisi{\leftexp{\prime}{4}}{4'}$)}] If $f$ is finite and there exists an integer $n$ such
that, for every geometric point $s$ of $S$, one has $\sum_{x\in
X(s)}n(x)=n$, then the composite arrow
$$
\Ff \longrightarrow
f_{\ast}f^{\ast}\Ff =f_{!}f^{\ast}\Ff \xrightarrow{\Tr_{f}}\Ff
$$
is multiplication by $n$.
\end{itemize}
\end{enonce*}

\begin{enonce*}[remark]{Example 6.2.6}\etiquette[6.2.6]{XVII.exem6.2.6}
Let $S$ be a \sisi{}{reduced }noetherian geometrically unibranch scheme (for
example, a normal scheme), and let $f:X\rightarrow S$ be a separated
quasi-finite morphism. For every geometric point $x$ of $X$ with image $s$ in $S$,
let $\Oo_{x}$ (\resp $\Oo_{s}$) be the strictly local ring of $X$
at $x$ (\resp of $S$ at $s$). The scheme $\Spec(\Oo_{s})$ is reduced
and has a unique maximal point $\eta$. Moreover, $\Oo_{x}$ is finite over
$\Oo_{s}$. Set $n(x)=\rg_{\eta}(\Oo_{x})$. One verifies that
$n(x)$ is a weighting of $f$, whence a trace morphism $\Tr_{f}$.
\end{enonce*}

\setcounter{subsubsection}{6}
\subsubsection{}\etiquette[6.2.7]{XVII.subsubsec6.2.7}
Let\pageoriginale $u:X\rightarrow Y$ be a separated, quasi-finite,
flat morphism of finite presentation. Since the trace morphism (\hyperref[XVII.thm6.2.3]{6.2.3}) is
functorial, it defines, for every complex of abelian sheaves $K$
on $Y$, a morphism of complexes
$$\etiquette[6.2.7.1]{XVII.eq6.2.7.1}
\Tr_{u}:u_{!} u^{\ast}K\longrightarrow K{\quoi}.\leqno(6.2.7.1)
$$
We shall use the terms \emph{trace morphism, Gysin
morphism} and \emph{morphism} $u_{\ast}$ interchangeably for any morphism obtained from
\eqhyperref[XVII.eq6.2.7.1]{6.2.7.1} or \hyperref[XVII.thm6.2.3]{6.2.3} by applying a functor; for example, the
following ones.

Consider a commutative diagram
$$
\xymatrix@C=1.5cm@R=1.75cm{
X  \ar[rr]^{u} \ar[rd]_{f}&   & Y\ar[ld]^{g}\\
                          & S &
}
$$
with $g$ separated and locally of finite type. For every abelian sheaf
$\Ff $ on $Y$, the trace morphism defines a morphism
$$\etiquette[6.2.7.2]{XVII.eq6.2.7.2}
\Tr_{u}\quad\mbox{\rm or}\quad u_{\ast}:f_{!}u^{\ast}\Ff \simeq
g_{!}u_{!}u^{\ast}\Ff \xrightarrow{g_{!}(\Tr_{u})} g_{!}\Ff {\quoi}.\leqno(6.2.7.2)
$$

If $g$ is compactifiable, and if $K$ is a complex bounded
below with torsion cohomology, or a complex of
$g^{\ast}\Aa$-modules for $\Aa$ a sheaf of torsion rings on
$S$, the trace morphism \eqhyperref[XVII.eq6.2.7.1]{6.2.7.1} likewise induces
\begin{align*}
\Tr_{u}\quad\mbox{\rm or}\quad u_{\ast} &: Rf_{!}u^{\ast}K\simeq
Rg_{!}u_{!}u^{\ast}K\xrightarrow{Rg_{!}(\Tr_{u})}Rg_{!}K{\quoi},\tag{6.2.7.3}\\
\Tr_{u}\quad\mbox{\rm or}\quad u_{\ast} &: R^{q}f_{!}u^{\ast}K\simeq
R^{q}g_{!}u_{!}u^{\ast}K\xrightarrow{R^{q}g_{!}(\Tr_{u})}
R^{q}g_{!}K{\quoi}.\tag{6.2.7.4}
\end{align*}\etiquette[6.2.7.3]{XVII.eq6.2.7.3}
\etiquette[6.2.7.4]{XVII.eq6.2.7.4}%The vertical space following this paragraph is too large.

For\pageoriginale $S$ the spectrum of an algebraically closed field,
\eqhyperref[XVII.eq6.2.7.4]{6.2.7.4} is a morphism
$$\etiquette[6.2.7.5]{XVII.eq6.2.7.5}
u_{\ast}:\HH^{q}_{c}(X,u^{\ast}K)\longrightarrow \HH^{q}_{c}(Y,K){\quoi}.\leqno(6.2.7.5)
$$

Thus one sees that \emph{cohomology with proper supports is
covariant} with respect to quasi-finite morphisms.

\subsubsection{}\etiquette[6.2.8]{XVII.subsubsec6.2.8}
Let $u:X\rightarrow Y$ be as in \hyperref[XVII.subsubsec6.2.7]{6.2.7}, let $\Delta_{n}=[0,n]$
be the standard set with $\sisi{}{(}n+1\sisi{}{)}$ elements, and let $X_{n}$ be
the $(n+1)$\textsuperscript{fold} iterated fiber product of $X$ over $Y$. The $X_{n}$ form
a simplicial scheme augmented to $Y$ (\cf \sisi{VI${}_{\text{B}}$}{VBIS}). We
denote by $u_{n}:X_{n}\rightarrow Y$ the augmentation
morphism. If $\Ff $ is an abelian sheaf on $Y$, the
sheaves $\Ff _{n}=u_{n!}u_{n}^{\ast}\Ff $ form,
by \eqhyperref[XVII.eq6.2.7.2]{6.2.7.2} (applied after the change of notation
$S\sisi{\rightarrow}{\mapsto} Y$, $X\sisi{\rightarrow}{\mapsto} X_{n}$, $Y\sisi{\rightarrow}{\mapsto} Y$ or $X_{m}$),
a \emph{simplicial} sheaf augmented to $\Ff $. By definition, the
differential complex of ``\emph{nondegenerate chains}'' of this
simplicial sheaf has as its $n$\textsuperscript{th}
component the quotient of $\Ff _{n}$, namely the cokernel of the direct
sum of the degeneracy operators with target $\Ff _{n}$ and
of the morphisms $\sigma-\epsilon(\sigma)$ for $\sigma$ a
symmetry operator.

\begin{enonce*}{Proposition 6.2.9}\etiquette[6.2.9]{XVII.prop6.2.9}
Under the preceding hypotheses, if $u$ is separated, étale, of finite type,
and surjective, then for every abelian sheaf $\Ff $ on $Y$, one has
\begin{enumerate}
\renewcommand{\theenumi}{\roman{enumi}}
\renewcommand{\labelenumi}{(\theenumi)}
\item The sheaves $\Ff _{n}$ form a (left)
resolution of $\Ff $.

\item The differential complex of nondegenerate chains of the simplicial
sheaf $(\Ff _{n})$ is a (left) resolution of
$\Ff $. If all the geometric fibers of $u$ have at most
$d$ points, this resolution has length $\leq d$.
\end{enumerate}
\end{enonce*}

By\pageoriginale base change, it suffices to verify this statement
when $Y$ is the spectrum of an algebraically closed field, in which case
\hyperref[XVII.subsubsec6.2.8]{6.2.8} reduces to computing the simplicial homology
of a simplex.

Let $(U_{i})_{i\in I}$ be a finite open covering of the scheme $Y$,
let $\Ff $ be an abelian sheaf on $Y$, and, for each subset $P$
of $I$, let $j_{P}$ be the inclusion of $U_{P}=\cap_{i\in P}U_{i}$ into
$\sisi{X}{Y}$. Suppose that a total order on $I$ has been chosen. If $X=\coprod_{i\in
I}U_{i}$ and if $u$ is the projection from $X$ to $Y$, the resolution
\hyperref[XVII.prop6.2.9]{6.2.9} (ii) of $\Ff $ is given by
$$\etiquette[6.2.9.1]{XVII.eq6.2.9.1}
(\Ff _{n})_{n\,\text{deg}}=\bigoplus_{|P|=n+1}j_{P!}j^{\ast}_{P}\Ff {\quoi}.\leqno(6.2.9.1)
$$

This resolution still makes sense when $I$ is infinite.

\setcounter{subsubsection}{9}
\subsubsection{}\etiquette[6.2.10]{XVII.subsubsec6.2.10}
Under the hypotheses of \hyperref[XVII.prop6.2.9]{6.2.9}, let $\Aa$ be a sheaf
of torsion rings on $Y$ and let $K\in \Ob^{-}(Y,\Aa)$. We
regard $X$ as ringed by $u^{\ast}\Aa$. Since resolution \hyperref[XVII.prop6.2.9]{6.2.9} (i) (\resp
\hyperref[XVII.prop6.2.9]{6.2.9} (ii)) is functorial in $\Ff $, it defines a
resolution of $K$ by a double complex $K_{1}$. This resolution
is called the extraordinary \emph{Čech resolution} (\resp the
\emph{extraordinary alternating} Čech resolution) of $K$.

Consider a commutative diagram
$$
\xymatrix@C=1.5cm@R=1.75cm{
X  \ar[rr]^{u} \ar[rd]_{f}&   & Y\ar[ld]^{g}\\
                          & S & ,
}
$$
with $g$ compactifiable. With the notation of
\hyperref[XVII.subsubsec6.2.8]{6.2.8}, if $f_{n}$ is the projection from $X_{n}$ to
$S$, one has $Rf_{n!}=Rg_{!}\sisi{U}{u}_{n!}$. The complex $K_{1}$, filtered
by\pageoriginale the semisimplicial degree, therefore defines a spectral
sequence
$$\etiquette[6.2.10.1]{XVII.eq6.2.10.1}
E^{-p,q}_{1}=R^{q}f_{p!}(u^{\ast}_{p}K)\Longrightarrow R^{-p+q} g_{!}
K{\quoi}.\leqno(6.2.10.1)
$$
The differentials $d_{1}$ arise from the simplicial arrows, so
that
$$\etiquette[6.2.10.2]{XVII.eq6.2.10.2}
E^{-p,q}_{2}=\check{\Hh}^{p}(R^{q}f_{p!}(u^{\ast}_{p}K)\sisi{}{)}\Longrightarrow
R^{-p+q}g_{!}K{\quoi}.\leqno(6.2.10.2)
$$

Under the hypotheses of \eqhyperref[XVII.eq6.2.9.1]{6.2.9.1}, one similarly obtains a spectral
sequence
$$\etiquette[6.2.10.3]{XVII.eq6.2.10.3}
E^{-p,q}_{1}=\bigoplus_{|P|=p+1}R^{q}f_{p!}(j^{\ast}_{p}K)\Longrightarrow R^{-p+q}g_{!}K{\quoi}.\leqno(6.2.10.3)
$$

The spectral sequences \eqhyperref[XVII.eq6.2.10.1]{6.2.10.1} and \eqhyperref[XVII.eq6.2.10.3]{6.2.10.3} are called the
\emph{upper localization spectral sequences.}

\begin{enonce*}{Proposition 6.2.11}\etiquette[6.2.11]{XVII.prop6.2.11}
If $f$ is a separated étale morphism of finite type, the functor
$f_{!}$ is left adjoint to the functor $f^{\ast}$, and its
adjunction arrow is the trace morphism.
\end{enonce*}

It is known that if $X$ is an object of a site $S$, the functor $f^{\ast}$
restricting abelian sheaves to $X$ has a left adjoint
$f_{!}$. If $F$ is an abelian sheaf on $X$, $f_{!}F$ is the
sheaf generated by the presheaf
$$
V\longrightarrow \bigoplus_{\xi \in \Hom(V,X)}F(\xi){\quoi}.
$$
In the étale case, provisionally denote this functor by $f\sisi{''}{^0}_{!}$. Its
adjunction property immediately implies that its formation commutes with
every base change. The trace morphism defines a morphism $t$
from $f_{!}\Ff $ to $f^{0}_{!}\Ff $,
compatible\pageoriginale with every base change. It suffices to see
that $t$ is an isomorphism fiber by fiber, \ie when $S$ is the spectrum of a
separably closed field, where this is trivial.

\setcounter{subsubsection}{11}
\paragraph{}\etiquette[6.2.11.1]{XVII.par6.2.11.1}
One should note that this interpretation of $f^{0}_{!}$ is
specific to the case of the étale site and does not extend to the fpqf site of a
prescheme, for example.

\subsection{Variant of the trace for continuous coefficients}\etiquette[6.3]{XVII.subsec6.3}

In this no. we prove the assertions of \textsc{Grothendieck}
\cite[no. 6]{sga4.exp.XVII.bib.Gr}, and apply these results to the definition
of a ``trace morphism'' for suitable fppf sheaves.

% END INCLUDED SGA4_English_current_buildable_source_20260906/source/17/drafts/SGA4_XVII_LINES_7001_7400_English.texfrag

% BEGIN INCLUDED SGA4_English_current_buildable_source_20260906/source/17/drafts/SGA4_XVII_LINES_7401_7807_English.texfrag

\subsubsection{}\etiquette[6.3.1]{XVII.subsubsec6.3.1}
Let $A$ be a commutative ring. If $L$ and $M$ are two finitely
generated projective $A$-modules, the map $u\mapsto
\bigwedge\limits^{m}u$ from $\Hom(L,M)$ to
$\Hom(\bigwedge\limits^{m}L,\bigwedge\limits^{m}M)$ is
compatible with every base change and is $m$\sisi{-homogeneous}{\textsuperscript{homogeneous}} (\sisi{\hyperref[XVII.subsubsec5.5.2]{5.5.2}}{\hyperref[XVII.par5.5.2.2]{5.5.2.2}}). By
\sisi{\hyperref[XVII.subsubsec5.5.2]{5.5.2}}{\hyperref[XVII.par5.5.2.3]{5.5.2.3}}, it defines a morphism
$$\etiquette[6.3.1.1]{XVII.eq6.3.1.1}
\bigwedge^{m}:\Gamma^{m}\Hom(L,M)\longrightarrow
\Hom(\bigwedge^{m} L, \bigwedge^{m}M)\leqno(6.3.1.1)
$$
such that, after every base change, for $u\in \Hom(L,M)$ one has
$\bigwedge\limits^{m}(\gamma^{m}(u))=\bigwedge\limits^{m}u$.
Since $\Hom(L,M)$ is flat, this morphism is identified (\sisi{\hyperref[XVII.subsubsec5.5.2]{5.5.2}}{\hyperref[XVII.par5.5.2.5]{5.5.2.5}}) with a
morphism
$$\etiquette[6.3.1.2]{XVII.eq6.3.1.2}
\bigwedge^{m}:TS^{m}\Hom(L,M)\longrightarrow
\Hom(\bigwedge^{m}L, \bigwedge^{m}M){\quoi}.\leqno(6.3.1.2)
$$
N.B. Such a homomorphism can also be defined whenever $2$ is
invertible in $A$: a symmetric element of
$\bigotimes\limits^{m}\Hom(L,M)$ defines a symmetric
homomorphism from $\bigotimes\limits^{m}L$ to
$\bigotimes\limits^{m}M$, and this descends to the
quotient to define a homomorphism from
$\bigwedge\limits^{m}L$ to
$\bigwedge\limits^{m}M$.

Given\pageoriginale three finitely generated projective $A$-modules,
the diagram
$$
\vcenter{\etiquette[6.3.1.3]{XVII.eq6.3.1.3}
\xymatrix{
\Gamma^{m} \Hom(L,M) \otimes \Gamma^{m} \Hom(M,N)  \ar[r] \ar[d]^{\bigwedge\limits^{m}\otimes \bigwedge\limits^{m}}& \Gamma^{m}
\Hom(L,N) \ar[d]^{\bigwedge\limits^{m}}\\
\Hom(\bigwedge\limits^{m} L,
\bigwedge\limits^{m} M) \otimes
\Hom(\bigwedge\limits^{m} M,
\bigwedge\limits^{m} N)         \ar[r]     &
\Hom(\bigwedge\limits^{m} L,
\bigwedge\limits^{m} N),
}}\leqno{(6.3.1.3)}
$$
in which the horizontal arrows are deduced from the composition maps
(and, for the first, from (\hyperref[XVII.par5.5.2.8]{5.5.2.8})), is commutative. This
expresses the formula
$\bigwedge\limits^{m}(vu)=\bigwedge\limits^{m}v\circ
\bigwedge\limits^{m}u$.

If $L$ is a free $A$-module of rank $n$, \eqhyperref[XVII.eq6.3.1.2]{6.3.1.2} defines a
canonical morphism
$$\etiquette[6.3.1.4]{XVII.eq6.3.1.4}
\det:\Gamma^{n}\End_{A}(L)=TS^{n}\End_{A}(L)\longrightarrow
A{\quoi},\leqno(6.3.1.4)
$$
which is an algebra homomorphism by
\eqhyperref[XVII.eq6.3.1.3]{6.3.1.3}. By definition,
one has
$$\etiquette[6.3.1.5]{XVII.eq6.3.1.5}
\det (u^{\otimes n})=\det (u\text{; }L){\quoi}.\leqno(6.3.1.5)
$$

If, in addition, $L$ is a module over an $A$-algebra $B$, \eqhyperref[XVII.eq6.3.1.4]{6.3.1.4}
defines by composition an algebra homomorphism
$$\etiquette[6.3.1.6]{XVII.eq6.3.1.6}
{\det}_{L}:TS^{n}_{A}(B)\longrightarrow A{\quoi}.\leqno(6.3.1.6)
$$

\begin{enonce*}{Lemma 6.3.2}\etiquette[6.3.2]{XVII.lem6.3.2}
Suppose that $L$ occurs in an exact sequence
\[
\sisi{F}{E}:0\longrightarrow L'\longrightarrow L\longrightarrow
L''\longrightarrow 0
\]
with $L'$ (\resp $L''$) locally free of rank $n_{1}$ (\resp
$n_{2}$). Let $n=n_{1}+n_{2}$, let $\alpha$ be the
``associated graded'' morphism from the group $\End(E)$ of
endomorphisms\pageoriginale of the extension $E$ to
$\End(L')\oplus \End(L'')$, and, for two flat $A$-modules $X$ and $Y$,
let $\beta$ be the morphism from $TS^{n}(X\oplus Y)$ to
$TS^{n_{1}}(X)\otimes TS^{n_{2}}(Y)$ deduced from the isomorphism
(\hyperref[XVII.par5.5.2.6]{5.5.2.6})
\[
TS^{n}(X\oplus Y)\simeq \sisi{\sum}{\bigoplus}_{n'+n''=n} TS^{n'\sisi{_{1}}{}}(X)\otimes
TS^{n''}(Y){\quoi}.
\]
The following diagram is commutative:{\footnotesize
\[
\xymatrix@C=1cm@R=1.5cm{
  TS^{n}(\End\sisi{)}{} (E)) \ar[d] \ar[r]^(.4){TS^{n}(\alpha)} & TS^{n}(\End(L')\oplus
\Ens(L''))  \ar[r]^(.45){\beta} & TS^{n_1} (\End(L')) \otimes TS^{n_2}
(\End(L'')) \ar[d]^{\det \otimes \det}\\
TS^{n}(\sisi{L}{\End(L)})         \ar[rr]^{\det} &                 & A.
}
\]
}
\end{enonce*}

By the universal property \hyperref[XVII.subsubsec5.5.2]{5.5.2} of $TS^{n}$, it suffices to
check that the two images of $u^{\otimes n}$ coincide for $u\in
\End(E)$. If $\alpha{u}=(u',u'')$, these images are respectively
$\det(u)$ and $\det(u').\det(u'')$, which proves the assertion.

\begin{enonce*}{Lemma 6.3.3}\etiquette[6.3.3]{XVII.lem6.3.3}
Let $A$ be a\sisi{}{ commutative} ring, $B$ a\sisi{}{ commutative} $A$-algebra, $C$ a
not necessarily commutative $B$-algebra, $E$ a $B$-module that is free
of rank $n$ over $A$, and $F$ a $C$-module that is free of rank $m$
over $B$. Suppose that $C$ is flat over $A$. By the universal property
\hyperref[XVII.par5.5.2.2]{5.5.2.2} of $\Gamma^{nm}$ and \hyperref[XVII.par5.5.2.3]{5.5.2.3}, there is a unique arrow
$\sigma:TS^{nm}_{A}(C)\rightarrow TS^{n}_{A}TS^{\sisi{n}{m}}_{B}(C)$ such
that, after every ring extension $A\rightarrow A'$, one has
$\sigma(c^{\otimes nm})=(c^{\otimes m})^{\otimes n}$. The following
diagram is then commutative:
$$
\xymatrix@C=2cm@R=1.5cm{
TS^n_A (B) \ar[d]^{\det_E}  &  & TS^{n}_A TS^{m}_B (C) \ar[ll]_{TS^n_A
  (\det_F)}\\
A    & & TS^{nm}_A (C) \ar[ll]_{\det_{E\otimes_B F}} \ar[u]_{\sigma}.
}
$$
\end{enonce*}

\begin{enonce*}[remark]{Proof}
By\pageoriginale \hyperref[XVII.par5.5.2.2]{5.5.2.2} and \hyperref[XVII.par5.5.2.3]{5.5.2.3}, it suffices to show that the
two images of an element $u^{\otimes nm}$ of $TS^{nm}_{A}(C)$
coincide. These images are $\det_{E}(\det_{F}(u))$ and
$\det_{E\otimes_{B}F}(u)$; they are expressed in terms of $A$, $B$,
$E$, $F$, and $u\in \End_{B}(F)$, and, to check that they coincide,
one may first replace $B$ by its (commutative) image $B_{1}$ in
$\End(E)$ and $F$ by $F\otimes_{B}B_{1}$. The formula to be proved is
then the block-determinant formula (\textsc{Bourbaki},
\sisi{Alg.}{\emph{Alg.}} \sisi{Ch.}{\chap}~8 \S~12, Lemma 2, \sisi{p.140}{p.~140}).
\end{enonce*}

When $C$ is no longer necessarily flat over $A$, the same proof gives,
more generally, a commutative diagram of algebras
$$\vcenter{\etiquette[6.3.3.1]{XVII.eq6.3.3.1}
\xymatrix@C=3cm@R=1.5cm{
\Gamma^{n}_A (B)  \ar[d]^{\det_E}& \Gamma^{n}_A
\Gamma^n_{B}(C) \ar[l]_{\Gamma_A^{n}(\det_F)}\\
A            & \ar[l]_{\det_{E\otimes_B F}}  \Gamma^{nm}_A (C) \ar[u]_{\sigma}.
}}\leqno{(6.3.3.1)}
$$

\setcounter{subsubsection}{3}
\subsubsection{}\etiquette[6.3.4]{XVII.subsubsec6.3.4}
If $B$ is a commutative $A$-algebra, $\Spec(TS^{n}(B))$ is precisely
the $n$\textsuperscript{th} symmetric power of $\Spec(B)$ over
$\Spec(A)$, \ie the quotient of $(\Spec(B)/\Spec(A))^{n}$ by the
symmetric group (\SGA 1 \sisi{IV 3}{V 1}).

If\pageoriginale $X$ is a scheme separated\sisi{}{ over $S$}, a
quasi-coherent $\Oo_{X}$-module $\Ff $ on $X$ will be called
\emph{finite locally free} (\resp of rank $n$) \emph{over} $S$
if its support is finite over $S$ and its direct image is locally free
of finite rank (\resp of rank $n$) over $S$. By the
\emph{support} of $\Ff $ we shall again mean the \emph{subscheme} of
$X$ defined by the annihilator of $\Ff $ (which is a quasi-coherent
ideal, since $\Ff $ is quasi-coherent of finite type).

If $\Ff $ is finite locally free over $S$, of rank $n$ and with support
$E$, then $E$ is affine over $S$; by globalizing over $S$, $\Ff $
defines, by means of \eqhyperref[XVII.eq6.3.1.6]{6.3.1.6}, a section
$\sigma_{E}(\Ff )$ of $\Sym^{n}_{S}(E)$, hence, by functoriality, a
section $\sigma(\Ff )$ of the algebraic space
\footnote[1]{\MakeLinkTarget*{sga4.exp.XVII.footnote.6.3.4.artin-1}\label{XVII.noteArtin1}M. \textsc{Artin}, Algebraization of formal moduli
I\sisi{}{, in \emph{Global Analysis: Papers in Honor of K.~Kodaira}, Princeton University Press, 1969, 21--71}. We shall use the preceding construction only when
$\Sym^{n}_{S}(X)$ exists as a scheme, which is the case when $X$ is
quasi-projective over $S$.} $\Sym^{n}_{S}(X)$,
$$\etiquette[6.3.4.1]{XVII.eq6.3.4.1}
\sigma(\Ff )\in \Gamma(\Sym^{n}_{S}(X)/S){\quoi}.\leqno(6.3.4.1)
$$

In particular, if $f:S'\rightarrow S$ is finite locally free of rank
$n$, then, for $\Ff =\Oo_{S'}$, one obtains a canonical section
$$\etiquette[6.3.4.2]{XVII.eq6.3.4.2}
t\in\Gamma(\Sym^{n}_{S}(S')/S){\quoi}.\leqno(6.3.4.2)
$$

\begin{enonce*}{Proposition 6.3.5}\etiquette[6.3.5]{XVII.prop6.3.5}
Let $X$ be a scheme separated over $S$, and let $\Ff _{1}$ and
$\Ff _{2}$ be two $\Oo_{X}$-modules finite locally free of ranks
$n_{1}$ and $n_{2}$ over $S$. Put $n=n_{1}+n_{2}$. Let $d$ be the
canonical arrow $d:\Sym^{n_{1}}_{S}(X)\times
\Sym^{n_{2}}_{S}(X)\rightarrow \Sym^{n}_{S}(X)$. If $\Ff $ is an
extension of $\Ff _{1}$ by $\Ff _{2}$, then
$$
\sigma(\Ff )=d(\sigma(\Ff _{1}),\sigma(\Ff _{2})){\quoi}.
$$
\end{enonce*}

One\pageoriginale may replace $X$ by the support (finite over $S$) of
$\Ff _{1}\oplus \Ff _{2}$ and suppose that $S$ is affine. Translated
into ring-theoretic terms, the proposition then follows from
\hyperref[XVII.lem6.3.2]{6.3.2}.

\begin{enonce*}{Lemma 6.3.6}\etiquette[6.3.6]{XVII.lem6.3.6}
If $X$ is a scheme separated over $S$, $\Ff $ an $\Oo_{X}$-module
finite locally free of rank $n$ over $S$, and $\Ll$ an invertible
sheaf on $X$, then
$$
\sigma(\Ff )=\sigma(\Ff \otimes\Ll){\quoi}.
$$
\end{enonce*}

Replace $X$ by the support (\hyperref[XVII.subsubsec6.3.4]{6.3.4}) of
$\Ff $\sisi{ l}{. L}ocally over $S$, $\Ll$ is then isomorphic to
$\Oo_{X}$, and the assertion is trivial.

The reader may verify:

\begin{enonce*}{Proposition 6.3.7}\etiquette[6.3.7]{XVII.prop6.3.7}
For every diagram of $S$-schemes separated over $S$
\[
\xymatrix{
X        \ar[rr]^{u} \ar[rd]_{f} &        &  Y  \ar[ld]^{g}\\
                                & S      &,
}
\]
and every $\Oo_{X}$-module $\Ff $ finite locally free of rank $n$ over
$S$, the $\Oo_{Y}$-module $u_{\ast}\Ff $ is finite locally free of
rank $n$ over $S$, and
\[
\sisi{u}{(\Sym^n_S u)}(\sigma(\Ff ))=\sigma(u_{\ast}\Ff ){\quoi}.
\]
\end{enonce*}

\setcounter{subsubsection}{7}
\paragraph{}\etiquette[6.3.7.1]{XVII.par6.3.7.1}
Finally, note that the construction of the morphisms $\sigma$
\eqhyperref[XVII.eq6.3.4.1]{6.3.4.1} and $t$ \eqhyperref[XVII.eq6.3.4.2]{6.3.4.2} is \emph{compatible with every
base change}, which makes sense by \hyperref[XVII.par5.5.2.7]{5.5.2.7}.

\subsection*{Application 1. The case of smooth curves}\pageoriginale

\subsubsection{}\etiquette[6.3.8]{XVII.subsubsec6.3.8}
Let $f:X\rightarrow S$ be a smooth quasi-projective curve over $S$,
\ie a smooth quasi-projective morphism purely of relative dimension
$1$. The theory of Hilbert schemes \sisi{[}{}\textsc{Grothendieck} \cite{sga4.exp.XVII.bib.Gr}) shows that the
fpqc sheaf $\sheaf{\Div}^{n}_{X/S}$ which assigns to each $S$-scheme
$S'$ the set of relative divisors on $X'=\fprod{X}{S'}{S}$ that are
\emph{finite} over $S$ and of degree $n$ is representable.

For each relative divisor $E$ on $X$, finite of degree $n$ over $S$,
there is a canonical element \sisi{(6.3.4.3)}{\eqhyperref[XVII.eq6.3.4.2]{6.3.4.2}}
$t(E)\in \Sym^{n}_{S}(E)$, with image $\sigma(E)$ in
$\Sym^{n}_{S}(X)$. Since this construction is compatible with every
base change, it defines
$$\etiquette[6.3.8.1]{XVII.eq6.3.8.1}
\sigma:\sheaf{\Div}^{n}_{X/S}\longrightarrow
\Sym^{n}_{S}(X):E\longrightarrow \sigma(\Oo_{E}){\quoi}.\leqno(6.3.8.1)
$$

Each section $s$ of $f$ defines a relative divisor finite over $S$ of
degree $1$, namely the divisor $s(S)$. The morphism
\[
(X/S)^{n}\longrightarrow \sheaf{\Div}^{n}_{X/S}:(s_{i})_{0<i\leq
n}\longmapsto \sum s_{i}(S)
\]
is symmetric, and hence induces a morphism
$$\etiquette[6.3.8.2]{XVII.eq6.3.8.2}
\alpha:\Sym^{n}_{S}(X)\longrightarrow \sheaf{\Div}^{n}_{X/S}{\quoi}.\leqno(6.3.8.2)
$$

\begin{enonce*}{Proposition 6.3.9}\etiquette[6.3.9]{XVII.prop6.3.9}
The morphisms $\sigma$ and $\alpha$ of \hyperref[XVII.subsubsec6.3.8]{6.3.8} are inverse
isomorphisms.
\end{enonce*}

\begin{enonce*}[remark]{Proof}
\begin{enumerate}
\renewcommand{\theenumi}{\alph{enumi}}
\renewcommand{\labelenumi}{\theenumi)}
\item $\sigma\alpha=\Id$. It suffices to prove that the composite arrow
$$
(X/S)^{n}\xrightarrow{\pi}\Sym^{n}_{S}(X)\xrightarrow{\alpha}\sheaf{\Div}^{n}_{X/S}\xrightarrow{\sigma}\Sym^{n}_{S}(X)
$$
is\pageoriginale $\pi$. For $n=1$, the assertion is clear; the general
case follows from the

\begin{enonce*}{Lemma 6.3.9.1}\etiquette[6.3.9.1]{XVII.lem6.3.9.1}
With the notation of \hyperref[XVII.prop6.3.5]{6.3.5} and \eqhyperref[XVII.eq6.3.8.1]{6.3.8.1}, one has
\[
\sigma(E+F)=d(\sigma(E),\sigma(F)){\quoi}.
\]
\end{enonce*}

This lemma follows from \hyperref[XVII.prop6.3.5]{6.3.5}, \hyperref[XVII.lem6.3.6]{6.3.6}, and the exact sequence
\[
0\longrightarrow \Oo(-F)\otimes \Oo_{E}\longrightarrow
\Oo_{E+F}\longrightarrow \Oo_{F}\longrightarrow 0{\quoi}.
\]

\item It is known that $\sheaf{\Div}^{n}_{X/S}$ is smooth over $S$
(there is no obstruction to infinitesimally lifting divisors on a
curve (\cf \SGA 3 XI 1)). The morphism
$\alpha\pi:(X/S)^{n}\rightarrow \sheaf{\Div}^{n}_{X/S}$ is a
quasi-finite surjective morphism between smooth $S$-scheme\sisi{}{s} whose
source has equidimensional fibers. It is therefore flat (\EGA ${\rm
0_{IV}}$ 17.3.5), so $\alpha\pi$ is an epimorphism of fppf sheaves, and
hence so is $\alpha$. In view of a), this completes the proof of
\hyperref[XVII.prop6.3.9]{6.3.9}.
\end{enumerate}
\end{enonce*}

\begin{enonce*}{Lemma 6.3.10}\etiquette[6.3.10]{XVII.lem6.3.10}
Let $k$ be a commutative ring, $M$ a free $k$-module, and $u$ an
endomorphism of $M$. Denote by $N_{\det(u-T)=0}$ the norm from
$k[T]/\det(u-T)$ to $k$. Then, for $P\in k[T]$, one has
\[
N_{\det(u-T)=0} (P(T))=\det(P(u),M){\quoi}.
\]
\end{enonce*}

This amounts to checking an algebraic identity between the entries of
the matrix of $u$ and the coefficients of $P$. To prove
\hyperref[XVII.lem6.3.10]{6.3.10}, it therefore suffices (by familiar
arguments) to treat the case in which $k$ is algebraically closed and
$u$ is semisimple, with eigenvalues $a_{i}$.

One\pageoriginale then has
\begin{align*}
& \det(u-T)=\Pi(a_{i}-T)\\
& N_{\det(u-T)=0}(P(T))=\Pi P(a_{i})\quad\mbox{\rm and}\\
& \det(P(u),M)=\Pi P(a_{i}){\quoi}.
\end{align*}

\setcounter{subsubsection}{11}
\setcounter{paragraph}{-1}
\paragraph{}\etiquette[6.3.11.0]{XVII.par6.3.11.0}
Let $f:X\rightarrow S$ be a quasi-projective morphism of finite
presentation, flat with purely one-dimensional Cohen--Macaulay fibers;
these conditions are thus satisfied if $X$ is a smooth quasi-projective
curve over $S$. Let $\Ff $ be an $\Oo_{X}$-module finite locally free
\emph{over} $S$. Since $\Ff $ vanishes on a fiberwise dense open subset
of $X$, there is a canonical arrow $\psi:\Oo_{X}\rightarrow
\det^{\ast}(\Ff )$ (\SGA 6 XI or \textsc{Mumford} \sisi{[1 Ch.~5 \S~3]}{\cite[\chap~5 \S~3]{sga4.exp.XVII.bib.Mu}}). We
denote by $\delta(\Ff )$ the divisor defined by the annihilator ideal of
$\coker(\psi)$ (ibid.); this divisor is finite over $S$.

\begin{enonce*}{Proposition 6.3.11.1}\etiquette[6.3.11.1]{XVII.prop6.3.11.1}
Under the hypotheses of \hyperref[XVII.par6.3.11.0]{6.3.11.0} and with the notation of
\hyperref[XVII.subsubsec6.3.8]{6.3.8} and \eqhyperref[XVII.eq6.3.4.1]{6.3.4.1}, if $\Ff $ has rank $n$, then
$\delta(\Ff )$ has degree $n$ and
\[
\sigma(\Ff )=\sigma(\delta(\Ff )){\quoi}.
\]
\end{enonce*}

\begin{enonce*}[remark]{Proof}
\begin{enumerate}
\renewcommand{\theenumi}{\alph{enumi}}
\renewcommand{\labelenumi}{\theenumi)}
\item \emph{The case where $f$ is smooth.} By a standard reduction, we
may assume that $S$ is an Artinian local scheme with algebraically
closed residue field. The additivity (\hyperref[XVII.prop6.3.5]{6.3.5})
of $\sigma$ allows us to suppose that the support of $\Ff $ is
concentrated at a rational point $x$ of $X$. One checks at once that if
$T$ is a uniformizer at $x$, defining $p:X\rightarrow
\sisi{A}{\AAA}^{1}_{S}$, then $p_{\ast}(\Oo_{\sisi{S}{\delta}(\Ff )})\simeq
\Oo_{\delta(p_{\ast}\Ff )}$ (because the constructions in question
commute with passage to the completion). By this and
\hyperref[XVII.prop6.3.7]{6.3.7}, we reduce to the case
$X=\sisi{A'}{\AAA^1}_{S}$.

If\pageoriginale $X=\AAA\sisi{'}{^1}_{S}$ and $S$ is affine,
$S=\Spec(k)$, then $\Ff $ is identified with a locally free $k$-module
$M$ equipped with an endomorphism $u$ (the action of $T$). As a
$k[T]$-module, $M$ has the resolution
\[
0\rightarrow
k[T]\otimes_{k}M\xrightarrow{\varphi}k[T]\otimes_{k}M\xrightarrow{\epsilon}
M\rightarrow 0{\quoi},
\]
where $\epsilon(P\otimes m)=P(u)(m)$ and
$\varphi(1\otimes m)=1\otimes um-T\otimes m$. From the formula
$\varphi=u-T$, it follows that the divisor $\delta(\Ff )$ has equation
$\det(u-T,M)=0$.

To check \sisi{6.3.11}{\hyperref[XVII.prop6.3.11.1]{6.3.11.1}}, it suffices to show that, after every base
change and for every section $\varphi$ of $\Oo_{X}$, one has
\[
{\det}_{k}(\varphi,M)={\det}_{k}(\varphi,\Oo_{\delta(\Ff )})(=N_{\delta(\Ff )/S}(\varphi)){\quoi}.
\]
For $\sisi{f}{\varphi}=P(T)$, this is precisely the assertion of
\hyperref[XVII.lem6.3.10]{6.3.10}.

\item The general case. We first state more explicitly the assertion
to be proved, in the following more concrete form:

\begin{enonce*}{Corollary 6.3.11.2}\etiquette[6.3.11.2]{XVII.cor6.3.11.2}
Under the hypotheses of \hyperref[XVII.par6.3.11.0]{6.3.11.0} on
$f:X\rightarrow S$, let $M$ be a finite locally free module on $X$, and
let $u$ be an endomorphism of $M$ such that $\det u$ is a section of
$\Oo_{X}$ that is regular on every fiber $X_{s}$ (whence
$\Coker u=M/u(M)$ and $\Coker \det u=\Oo_{X}/\det u$ are both flat and
of finite presentation over $S$, with \sisi{quasi-finite fibers}{quasi-finite supports}). Suppose that the common support\nde{the support of $M/uM$ in the sense of \hyperref[XVII.subsubsec6.3.4]{6.3.4} is a subscheme of the support of $\Oo/\det u$ defined by a nilideal.} of these
sheaves is finite over $S$, so that they define finite locally free
sheaves over $S$. Then, for every section $\varphi$ of $\Oo_{X}$, the
determinants over $\Oo_{S}$ of $\varphi$ acting on $\Coker u$ and on
$\Coker \det u$ are equal:
$$\etiquette[6.3.11.3]{XVII.eq6.3.11.3}
{\det}_{S}(\varphi,M/u(M))={\det}_{S}(\varphi,\Oo_{X}/\det
u){\quoi}.\leqno(6.3.11.3)
$$
\end{enonce*}

% END INCLUDED SGA4_English_current_buildable_source_20260906/source/17/drafts/SGA4_XVII_LINES_7401_7807_English.texfrag

% BEGIN INCLUDED SGA4_English_current_buildable_source_20260906/source/17/drafts/SGA4_XVII_LINES_7808_8200_English.texfrag

We\pageoriginale note that this statement implies the equality of the ranks
of the two locally free modules under consideration over $S$, by
homogeneity (when $\varphi$ is replaced by $T\varphi$, with $T$ an
``indeterminate'' section of $\Oo_{S}$). Moreover, %\c
replacing $\varphi$ by $1+T\varphi$, where $\sisi{j}{T}$ is such an indeterminate (more
precisely, one makes a base change $S'=S[T]\rightarrow
S$), one deduces from \eqhyperref[XVII.eq6.3.11.3]{6.3.11.3} the apparently more general formula
$$\etiquette[6.3.11.4]{XVII.eq6.3.11.4}
\Pol_{S}(\varphi,M/uM)(T)=\Pol_{S}(\varphi,\Oo_{X}/\det u)(T)\leqno(6.3.11.4)
$$
for the characteristic polynomials.

To prove \eqhyperref[XVII.eq6.3.11.3]{6.3.11.3}, we shall reduce to the already proved smooth
case. By étale localization on $S$, we reduce to the case where $S$ is
strictly local with closed point $s$; then, by additivity, we reduce
to the case where the support of $M/u(M)$ has a single point on $X_{s}$, say
$x$. After making a flat extension of the base, we may assume
the residue field extension $k(x)/k(s)$ is trivial; and after replacing $X$
by a sufficiently small neighborhood of $x$, and adding to $\varphi$ a
function that induces the zero operation on $M/u(M)$ and on
$\Oo_{X}/ \det u$ (which does not change the formula to be proved), one
readily sees that we may assume that $\varphi$ induces a section of
$\Oo_{X_{s}}$ regular at $x$, and hence (after shrinking $X$)
defines a quasi-finite flat morphism
$$
p:X\longrightarrow X'=S[T]{\quoi}.
$$
By étale localization on $X$ and on $X'$, we may therefore assume
that we have a \emph{finite} and \emph{flat} morphism $p$ as above, with
$X'$ smooth over $S$, and $\varphi=p^{\ast}(\varphi')$, where $\varphi'$ is
a section of $\Oo_{X'}$. Set
\begin{align*}
& M'=p_{\ast}(M)\quav u'=p_{\ast}(u){\quoi},\\
& P=\Coker u\quav Q=\Coker \det u{\quoi},\\
& P'=\Coker u'\quav Q'=\Coker \det u'{\quoi}.
\end{align*}
We\pageoriginale know that
\[
\det u'=N_{\sisi{X'/X}{X/X'}}\det u,\quad\mbox{\rm hence}\quad Q'=\Coker
N_{\sisi{X'/X}{X/X'}}(\det u){\quoi}.
\]
We trivially have
\[
{\det}_{S}(\varphi,P)={\det}_{S}(\varphi',P\sisi{}{'}){\quoi},
\]
and by formula \eqhyperref[XVII.eq6.3.11.3]{6.3.11.3}, proved in the smooth case,
\[
{\det}_{S}(\varphi',P')={\det}_{S}(\varphi',Q'){\quoi},
\]
so it remains to prove the formula
\[
{\det}_{S}(\varphi',Q')={\det}_{S}(\varphi,Q){\quoi},
\]
which is \emph{equivalent} to formula \eqhyperref[XVII.eq6.3.11.3]{6.3.11.3} to be established. We thus see
that the validity of the latter depends only on the section $\det u$ of
$\Oo_{X}$, or equivalently that it is equivalent to the same formula in the case
where $M$ is replaced by $\Oo_{X}$, and $\varphi$ by $\det u$. But
in the case $M=\Oo_{X}$, formula \eqhyperref[XVII.eq6.3.11.3]{6.3.11.3} is tautological!
This completes the proof.
\end{enumerate}
\end{enonce*}

\begin{enonce*}[remark]{Remark 6.3.11.5}\etiquette[6.3.11.5]{XVII.rem6.3.11.5}
In\pageoriginale the case where $S$ is Artinian, the equality over $S$ of the ranks
of $\Coker u$ and $\Coker \det u$ (notation of \hyperref[XVII.cor6.3.11.2]{6.3.11.2}) is a
special case of \EGA IV 21.10.17.3, of which \eqhyperref[XVII.eq6.3.11.3]{6.3.11.3} may be
regarded as a ``relative'' version.
\end{enonce*}

\begin{enonce*}{Proposition 6.3.12}\etiquette[6.3.12]{XVII.prop6.3.12}
Let $X$ and $Y$ be two quasi-projective curves of finite
presentation, flat with Cohen-Macaulay fibers over $S$, and let $u:X\rightarrow Y$
be a quasi-finite flat morphism.

The following diagram
\[
\xymatrix@C=4em{
\sheaf{\Div}^{n}_{X/S}  \ar[r]^{u_{\ast}} \ar[d]_{\sigma_X} &
\sheaf{\Div}^{n}_{Y/S} \ar[d]_{\sigma_Y}\\
\Sym^{n}_S(X)     \ar[r]^{\Sym^{n}(u)} & \Sym^{n}_S(Y)
}
\]
is then commutative \footnote[1]{\MakeLinkTarget*{sga4.exp.XVII.footnote.6.3.12.1}\label{XVII.note192}The morphism $u_{\ast}$ is defined
because we restrict ourselves to divisors \emph{finite} over $S$. We take, for
example, $u_{\ast}(D)=\delta(u_{\ast}\Oo_{D})$ as the definition, a
formula that agrees with \EGA IV 21.5.5 when $u$ is finite.}.
\end{enonce*}

\sisi{\noindent}{}%
Since $u_{\ast}(D)=\delta(u_{\ast}\Oo_{D})$, \hyperref[XVII.prop6.3.12]{6.3.12} follows from
\hyperref[XVII.prop6.3.7]{6.3.7} and \hyperref[XVII.prop6.3.11.1]{6.3.11.1}.

\subsection*{Application 2. The trace morphism}

\setcounter{subsubsection}{12}
\subsubsection{}\etiquette[6.3.13]{XVII.subsubsec6.3.13}
Let $G$ be an abelian sheaf on the large fppf site of a scheme $S$
(\hyperref[XVII.subsec0.10]{0.10}). Suppose that $G$ satisfies the following condition.

\paragraph{}\etiquette[6.3.13.1]{XVII.par6.3.13.1}
\emph{For every affine flat morphism of finite presentation
$f:X\rightarrow Y$ of $S$-schemes, and for every integer $n\geq 1$, the
morphism}
\[
\Hom_{S}(\Sym^{n}_{Y}(X),G)\longrightarrow \Hom_{S}((X/Y)^{n},G)
\]
\emph{is\pageoriginale injective and has as its image the set of
symmetric morphisms from $(X/Y)^{n}$ to $G$.}

This condition is stable under base change.

Let $f:T\rightarrow S$ be a finite locally free morphism of rank
$n$. Let $\pr_{i}$ be the $i$\textsuperscript{th} projection from $(T/S)^{n}$
to $T$, and let $t$ be the section \sisi{(6.3.4.3)}{\eqhyperref[XVII.eq6.3.4.2]{6.3.4.2}} of $\Sym^{n}_{S}(T)$. If $u\in
G(T)$, then $\sum \pr^{\ast}_{i}(u)$ is a symmetric element of
$G\sisi{}{(}(T/S)^{n}$) and therefore defines
\[
u_{n}\in \sisi{}{G(}\Sym^{n}_{S}(T)\sisi{}{)}{\quoi}.
\]
We set
\[
\Tr_{f}(u)=t^{\ast}(u_{n}){\quoi}.
\]

When $f$ is finite locally free, we define $\Tr_{f}(u)$,
locally on $S$, by the preceding formula. This construction is
compatible with every base change and therefore defines a
\emph{trace morphism}
$$\etiquette[6.3.13.2]{XVII.eq6.3.13.2}
\Tr_{f}:f_{\ast}f^{\ast}G\longrightarrow G{\quoi}.\leqno(6.3.13.2)
$$

\subsubsection{}\etiquette[6.3.14]{XVII.subsubsec6.3.14}
Condition (\hyperref[XVII.par6.3.13.1]{6.3.13.1}) is satisfied in the following cases:
\begin{enumerate}
\renewcommand{\theenumi}{\alph{enumi}}
\renewcommand{\labelenumi}{\theenumi)}
\item If $G$ is representable, by the definition of symmetric
powers.

\item For $G$ the inverse image on $S_{\text{\rm fppf}}$ of a
sheaf on the small étale site of $S$ (\hyperref[XVII.subsec0.10]{0.10}).

\item For $G$ defined by a quasi-coherent $\Oo_{S}$-module $G_{0}$.
\end{enumerate}

Assertion c) reduces to (\hyperref[XVII.par5.5.2.7]{5.5.2.7}).

\begin{enonce*}{Proposition 6.3.15}\etiquette[6.3.15]{XVII.prop6.3.15}
For\pageoriginale $G$ satisfying (\hyperref[XVII.par6.3.13.1]{6.3.13.1}), we have:
\begin{enumerate}
\renewcommand{\theenumi}{\roman{enumi}}
\renewcommand{\labelenumi}{(\theenumi)}
\item The morphism $\Tr_{f}$ is functorial in $G$, and hence additive.

\item Its formation is compatible with every base change
$S'\rightarrow S$.

\item If $T=\coprod_{i}T_{i}$, so that $f_{\ast}G\simeq \sum
f_{i\ast}G$, then $\Tr_{f}=\sum \Tr_{f_{i}}$.

\item If $T$ has constant rank $n$ over $S$, the composite morphism
\[
G\longrightarrow f_{\ast}f^{\ast}G\xrightarrow{\Tr_{f}} G
\]
is multiplication by $n$.

\item The trace morphism is compatible with composition \emph{(\cf
\hyperref[XVII.thm6.2.3]{6.2.3} Var 3).}
\end{enumerate}
\end{enonce*}

Assertions (i) and (ii) are clear, and (iii) is an easy consequence
of \hyperref[XVII.prop6.3.5]{6.3.5}. With the notation
of \hyperref[XVII.subsubsec6.3.13]{6.3.13}, if $u\in G(S)$, then
$u_{n}$ is $n$ times the pullback of $u$, so that
$\Tr_{f}(f^{\ast}u)=t^{\ast}(u_{n})=nu$.

To prove (v), we reduce, by étale localization on $S$ and
using (iii), to proving that for
\[
f=gh:S''\xrightarrow{h} S'\xrightarrow{g}S{\quoi},
\]
with $\sisi{f}{g}$ and $\sisi{g}{h}$ finite locally free of \emph{constant ranks} $n$
and $m$, the diagram
$$\vcenter{\etiquette[6.3.15.1]{XVII.eq6.3.15.1}
\xymatrix{
G(S'')    \ar[rr]^{\Tr_h}\ar[rd]_{\Tr_f} &        &G(S') \ar[ld]^{\Tr_g}\\
                &  G(S)         &
}}\leqno{(6.3.15.1)}
$$
is commutative. For $u\in G(S'')$, consider the diagram
$$\vcenter{\etiquette[6.3.15.2]{XVII.eq6.3.15.2}
\xymatrix@C=2.5cm@R=1.75cm{
\Sym^{n}_S(S')  \ar[rr]^{\Sym^{n}(t)} \ar[rd]^{\Tr_{\sisi{g}{h}}(u)_n} & & \Sym^{n}_S
\Sym^{m}_{S'} (S'')  \ar[dd] \ar[ld]^{(u_m)_n}\\
                                          & G &\\
          S   \ar@<3pt>[ru]^{\Tr_{\sisi{f}{g}}\Tr_{\sisi{g}{h}}(u)} \ar@<-3pt>[ru]_{\Tr_{\sisi{fg}{gh}}(u)}  \ar[rr]^{t}  \ar[uu]_{t}   &        & \Sym^{nm}_S(S'') \ar[lu]_{u_{nm}}.
}}\leqno{(6.3.15.2)}
$$
The\pageoriginale right vertical arrow is defined because $S\sisi{'}{''}$ is
flat over $S$; every triangle is commutative, and (v) will therefore be a
consequence of the commutativity of the outer boundary.

Translating into ring-theoretic terms, this follows from the
special case $B=E$ and $C=F$ of \hyperref[XVII.lem6.3.3]{6.3.3}.

\setcounter{subsubsection}{15}
\subsubsection{}\etiquette[6.3.16]{XVII.subsubsec6.3.16}
Now let $f:T\rightarrow S$ be a \sisi{separated }{}quasi-finite flat morphism of
finite presentation. If $ℱ$ is an fppf sheaf on $T$, we define
its direct image with proper supports $f_{!}\sisi{T}{ℱ}$ on $S$ as the
presheaf that assigns to $S'/S$ the set of sections of $\sisi{F}{ℱ}$ on
$T'$ whose support is proper over $S'$. This presheaf is a
sheaf by \hyperref[XVII.prop6.1.1]{6.1.1}. With the notation of
\hyperref[XVII.prop6.1.1]{6.1.1}, it may also be
expressed as a local inductive limit (\sisi{IV 14.\quad}{\hyperrefext[V][V.subsec8.2]{8.2}}):
\[
f_{!}\sisi{F}{ℱ(U)}=\varinjlim_{V\in\Phi (U)}f_{\ast}(\sisi{F}{ℱ}|V){\quoi},
\]
a formula\pageoriginale that accounts for the fairly evident fact that the
properties of $f_{!}$ reduce to those of $f'_{\ast}$ when $f'$ is
finite. We thus obtain:

\begin{enonce*}{Proposition 6.3.17}\etiquette[6.3.17]{XVII.prop6.3.17}
There is a unique way to define, for every
$f:T\rightarrow S$ \sisi{}{separated }quasi-finite flat morphism of finite presentation and every
fppf sheaf $G$ on $S$ satisfying (\hyperref[XVII.par6.3.13.1]{6.3.13.1}), a trace morphism
$$
\Tr_{f}:f_{!}f^{\ast}G\longrightarrow G
$$
satisfying:
\begin{enumerate}
\renewcommand{\theenumi}{\roman{enumi}}
\renewcommand{\labelenumi}{(\theenumi)}
\item For finite $f$, one recovers the morphism \eqhyperref[XVII.eq6.3.13.2]{6.3.13.2}.

\item Conditions \emph{(Var 1)} through \emph{(Var 4)} of \hyperref[XVII.thm6.2.3]{6.2.3} are
satisfied, mutatis mutandis.
\end{enumerate}
\end{enonce*}

The quasi-finite morphism $f$ induces a commutative diagram (\cf
(\hyperref[XVII.subsec0.11]{0.11})) of sites (in which $\alpha_{\ast}$ is the ``restriction'' functor)
$$
\xymatrix{
T_{\text{\rm fppf}} \ar[rr]^{\alpha} \ar[d]^{f}& & T_{\text{et}} \ar[d]^{f_{\et}}\\
S_{\text{\rm fppf}} \ar[rr]^{\alpha}& & S_{\et}.
}
$$
For $G$ on $T_{\text{\rm fppf}}$, we have
$$
f_{!}\alpha_{\ast}\simeq \alpha_{\ast}f_{!}{\quoi}.
$$
We\pageoriginale also have a ``base-change''
morphism
$$
c:f^{\ast}_{\text{\rm et}}\alpha_{\ast}\longrightarrow
\alpha_{\ast}f^{\ast}{\quoi}.
$$
We leave it to the reader to verify the following compatibility
between the trace morphisms \hyperref[XVII.thm6.2.3]{6.2.3} and
\hyperref[XVII.prop6.3.17]{6.3.17}. For $G$ on $S_{\text{\rm
fppf}}$ satisfying (\hyperref[XVII.par6.3.13.1]{6.3.13.1}), the following diagram is commutative:
$$\MakeLinkTarget*{sga4.exp.XVII.equation.6.3.17.1}\label{XVII.eq6.3.17.1}
\vcenter{\xymatrix{
f_{\et!}f^{\ast}_{\et}\alpha_{\ast} G \ar[r]^{f_{!}c} \ar[d]_{\Tr_f}^{\hyperref[XVII.thm6.2.3]{(6.2.3)}}&
f_{\et!}\alpha_{\ast}f^{\ast}G  \ar@{<->}[r]^{\sim}&
\alpha_{\ast}f_!f^{\ast}G \ar[d]_{\alpha_{\ast}\Tr_f}^{\hyperref[XVII.prop6.3.17]{(6.3.17)}}\\
\alpha_{\ast}G                   \ar@{=}[rr] &        & \alpha_{\ast}G.
}}\leqno(6.3.17.1)
$$

In more concrete terms, this means that if $G$ is an fppf sheaf
on $S$ whose restriction to the étale site defines an étale sheaf
$\alpha_{\ast}G$, and if $u$ is a properly supported section of
the inverse image $f^{\ast}_{\text{\rm et}}(\alpha_{\ast}G)$, with image
$u_{1}$ in $f^{\ast}G$, then the elements $\Tr_{f}(u)$ (in the sense of \hyperref[XVII.thm6.2.3]{6.2.3}) and
$\Tr_{f}(u_{1})$ (in the sense of \hyperref[XVII.prop6.3.17]{6.3.17}) agree in $G(S)$.

\begin{enonce*}[remark]{Example 6.3.18}\etiquette[6.3.18]{XVII.exem6.3.18}
Let $f:S'\rightarrow S$ be a finite locally free morphism of rank
$n$ and let $u\in \Gamma(S',\Oo_{S'})=\Hom_{S}(S',\EE^{1}_{S})$.\nde{Here $\EE^1_S$ denotes the affine line over $S$.} The morphism $u$
defines a morphism
\[
\Sym(u):\Sym^{n}_{S}(S')\longrightarrow
\Sym^{n}_{S}(\EE^{1}_{S})=\Spec(S[\sigma_{1}\ldots
\sigma_{n}]\sisi{)}{\quoi},
\]
where the $\sigma_{i}$ are the elementary symmetric functions (\cf
\textsc{Bourbaki}, \emph{Alg.} 5, App. I no.~2\nde{reference to be checked. Probably IV \S~6 no.~1 in the new version.}). We then have (by an
easy calculation left to the reader):
$$\etiquette[6.3.18.1]{XVII.eq6.3.18.1}
[\Sym(u)\circ t]^{\ast}(\sigma_{i})=\sigma_{i}(u){\quoi},\leqno(6.3.18.1)
$$
where\pageoriginale $t\in \Gamma$ $(\Sym^{n}_{S}(S')/S)$ is the canonical
element \sisi{(6.3.4.3)}{\eqhyperref[XVII.eq6.3.4.2]{6.3.4.2}}, and where $\sigma_{i}(u)$ denotes the
$i$\textsuperscript{th} coefficient of the characteristic polynomial of
multiplication by $u$ on $f_{\ast}\Oo_{S'}$. It follows that the
``trace'' maps
\[
\begin{cases}%raghu
\GG_{a}(S')\longrightarrow \GG_{a}(S)\\
\GG_{m}(S')\longrightarrow \GG_{m}(S)
\end{cases}
\]
are respectively the usual \emph{trace} and \emph{norm}\sisi{}{}; indeed,
the morphism $\Sym^{n}_{S}(\EE^{1}_{S})\rightarrow
\EE^{1}_{S}$ induced by the additive law (\resp multiplicative law)
on $\EE^{1}_{S}$ is given by the function $\sigma_{1}$ (\resp
$\sigma_{n}$) on $\Sym^{n}_{S}(\EE^{1}_{S})\simeq \Spec
(S[\sigma_{1},\ldots,\sigma_{n}])$.

\sisi{\noindent}{}%
It then follows from \hyperref[XVII.eq6.3.17.1]{(6.3.17.1)} that, when $n$ is invertible on $S$, the
diagram of étale sheaves
$$\vcenter{\etiquette[6.3.18.2]{XVII.eq6.3.18.2}
\xymatrix{
0 \ar[r]& f_{\ast}\mu_n \ar[r] \ar[d]_{\Tr_f}^{\hyperref[XVII.thm6.2.3]{(6.2.3)}}& f_{\ast}\sisi{G}{\GG}_m
\ar[r] \ar[d]^{N_f}
& f_{\ast}\sisi{G}{\GG}_m \ar[d]^{N_f} \ar[r]& 0\\
0 \ar[r]& \mu_n \ar[r]        & \relax\sisi{G}{\GG}_m    \ar[r]     & \relax\sisi{G}{\GG}_m  \ar[r]       & 0
}}\leqno{(6.3.18.2)}
$$
is commutative.
\end{enonce*}

\begin{enonce*}{Proposition 6.3.19}\etiquette[6.3.19]{XVII.prop6.3.19}
Consider a commutative diagram
\[
\xymatrix{
X \ar[rr]^{u} \ar[rd]_{p}&   & Y\ar[ld]^{q}\\
  & S &
}
\]
in\pageoriginale which $X$ and $Y$ are smooth quasi-projective
curves over $S$ and $u$ is a quasi-finite morphism.
\begin{enumerate}
\renewcommand{\theenumi}{\roman{enumi}}
\renewcommand{\labelenumi}{(\theenumi)}
\item If $D$ is a divisor on $X$ finite over $S$ (see note
\emph{p.~\pageref{XVII.note192}}) and $g$ is a section of $G_{Y}$ on $Y$, then
\[
\Tr_{D/S}(u^{\ast}g)=\Tr_{(u_{\ast}D)/S}(g){\quoi}.
\]

\item If $E$ is a divisor on $Y$ finite over $S$, if $g$ is a section of
$G_{X}$ on $X$, and if $u$ is finite, then
\[
\Tr_{E/S}(\Tr_{X/Y}(g))=\Tr_{(u\sisi{_{\ast}}{^\ast}E)/S}(\sisi{G}{g}){\quoi}.
\]
\end{enumerate}
\end{enonce*}

\begin{enonce*}[remark]{Proof}
\begin{enumerate}
\renewcommand{\theenumi}{\roman{enumi}}
\renewcommand{\labelenumi}{(\theenumi)}
\item Let $g_{n}$ be the section of $G$ on $\Sym^{n}_{S}(Y)$ obtained from
$g$ (\hyperref[XVII.subsubsec6.3.13]{6.3.13}). By \hyperref[XVII.prop6.3.12]{6.3.12}, we have
\begin{align*}
\Tr_{D/S}(u^{\ast}g)\overset{\dfn}{=}
\sigma(D)^{\ast}(u^{\ast}g)_{n} & =\sigma(D)^{\ast}(\sisi{u}{\Sym^n_S(u)}^{\ast}g_{n})=\sisi{u}{(\Sym^n_S(u)}\circ
\sigma(D))^\ast
(g_{n})\\
& =\sigma(u_{\ast}D)^\ast(g_{n})=\Tr_{u_{\ast}D}(g_{n}){\quoi}.
\end{align*}

\item Formula (ii) follows from \hyperref[XVII.prop6.3.15]{6.3.15} (ii) and (v):
\[
\Tr_{E/S}(\Tr_{X/Y}(g))=\Tr_{E/S}(\Tr_{u^{\ast}E/E}(g))=\Tr_{u^{\ast}E/S}(g){\quoi}.
\]
\end{enumerate}
\end{enonce*}

\subsection*{Application 3. Trace of a torsor}

\setcounter{subsubsection}{19}
\subsubsection{}\etiquette[6.3.20]{XVII.subsubsec6.3.20}
Let $G$ be a sheaf on the large fppf site of $S$ that satisfies
(\hyperref[XVII.par6.3.13.1]{6.3.13.1}). Let $u:X\rightarrow Y$ be a finite locally free morphism
of $S$-schemes, and let $K$ be a $G_{X}$-torsor on $X$. Suppose that,
locally for the fppf topology on $Y$, the torsor $K$ on $X$ is trivial. This is
the case if $K$ is locally trivial for the étale topology.

\begin{enonce*}{Proposition 6.3.21}\etiquette[6.3.21]{XVII.prop6.3.21}
Under\pageoriginale the hypotheses of \hyperref[XVII.subsubsec6.3.20]{6.3.20}, there exists, up to unique
isomorphism, exactly one torsor $\Tr_{u}(K)$ under $G_{Y}$ on $Y$,
equipped with a morphism $\Tr_{u}$ or $\Tr^{K}_{u}:u_{\ast}K\rightarrow
\Tr_{u}(K)$, which, for $g$ a local section (in the fppf sense) of $G$ on
$Y$ and $k$ a local section of $u_{\ast}K$, satisfies
$$

% END INCLUDED SGA4_English_current_buildable_source_20260906/source/17/drafts/SGA4_XVII_LINES_7808_8200_English.texfrag

% BEGIN INCLUDED SGA4_English_current_buildable_source_20260906/source/17/drafts/SGA4_XVII_LINES_8201_8600_English.texfrag
\Tr^{K}_{u}(gk)=\Tr_{u}(g)\cdot\Tr_{u}\sisi{}{^K}(k){\quoi}.
$$
\end{enonce*}

The problem is local on $Y$; we may therefore assume that $K$ is trivial:
$K=G_{X}$. One solution to the problem is to take $K=G_{Y}$ and as
trace morphism the morphism \eqhyperref[XVII.eq6.3.13.2]{6.3.13.2}. One immediately verifies that it is
the only one.

The trace functor \sisi{\hyperref[XVII.subsubsec6.3.20]{6.3.20}}{\hyperref[XVII.prop6.3.21]{6.3.21}} is defined in particular
when $G$ satisfies the following condition

\setcounter{subsubsection}{21}
\paragraph{For every $S$-scheme $X$, every torsor under $G_{X}$ on $X$
is locally trivial for the étale topology}\etiquette[6.3.21.1]{XVII.par6.3.21.1}

This condition is satisfied when $G$ is a smooth group, when $G$ is the inverse
image of a sheaf on the small étale site of $S$, or when $G$ is
defined by a quasi-coherent sheaf on $S$.

By construction one has
$$\etiquette[6.3.21.2]{XVII.eq6.3.21.2}
\Tr_{u}(K)=u_{\ast}(K) \times_{u_{\ast}G_{X}}G_{Y}{\quoi}.\leqno(6.3.21.2)
$$

\subsubsection{}\etiquette[6.3.22]{XVII.subsubsec6.3.22}
To define a trace functor under more general hypotheses, the
key point is, for $f:X\rightarrow T$ a quasi-projective, flat
morphism of finite presentation of $S$-schemes, to define a
reasonable functor $K\mapsto K_{n}$ from the $G_{X}$-torsors on $X$ to the
$G_{Z}$-torsors on $Z=\Sym^{n}_{T}(X)$. We\pageoriginale first
have the following lemma, where $f_{n}$ denotes the projection from $\Sym^{n}_{T}(X)$ to $T$;
it is proved in the same way as \hyperref[XVII.prop6.3.21]{6.3.21}:

\hypertarget{sga4.exp.XVII.lemma.6.3.22.1}{}\begin{enonce*}{Lemma 6.3.22.1}\etiquette[6.3.22.1]{XVII.lem6.3.22.1}
With the preceding notation, and if, locally on $T$, the torsor
$K$ is trivial, then, up to unique isomorphism, there exists a unique
torsor $K_{n}$ on $\Sym^{n}_{T}(X)$, equipped with a morphism
$\alpha:f_{\ast}K\rightarrow f_{n{\ast}}(K_{n})$, such that for $g$ a
local \emph{(fppf)} section of $f_{\ast}G_{X}$ and $k$ a local section
of $\sisi{u}{f}_{\ast}K$, one has $\alpha(gk)=g_{n}\sisi{\circ}{\cdot} \alpha(k)$ (notation of
\hyperref[XVII.subsubsec6.3.13]{6.3.13} for $g_{n}$; the definition of $g_{n}$ here uses
the flatness hypothesis \etc).
\end{enonce*}

The general case will be handled by a somewhat contrived descent argument. First
of all:

\hypertarget{sga4.exp.XVII.proposition.6.3.23}{}\begin{enonce*}{Proposition 6.3.23}\etiquette[6.3.23]{XVII.prop6.3.23}
Consider a commutative diagram
$$\vcenter{\etiquette[6.3.23.1]{XVII.eq6.3.23.1}
\xymatrix{
X \ar[rr]^{u} \ar[rd]& & Y\ar[ld]\\
  & T &
}}\leqno{(6.3.23.1)}
$$
of flat, finitely presented, quasi-projective $T$-schemes, with
$u$ étale and surjective. Then the semisimplicial system of schemes
$\Sym^{n}_{T}((X/Y)^{k+1})$ is a hypercover for the étale
topology of the scheme $\Sym^{n}_{T}(Y)$.
\end{enonce*}

\begin{enonce*}[remark]{Proof}
\begin{enumerate}
\renewcommand{\theenumi}{\alph{enumi}}
\renewcommand{\labelenumi}{\theenumi)}
\item We shall use the following auxiliary result.

\setcounter{subsubsection}{23}\setcounter{paragraph}{1}
\paragraph{}\etiquette[6.3.23.2]{XVII.par6.3.23.2}
\emph{Let $(X,Y,T)$ be as above, let $t\in T$, and let $Z$ be a
$T$-scheme that is the disjoint union of finitely many $T$-schemes $\Spec(k_{i})$, where
$k_{i}$ is a finite extension of $k_{t}$,\nde{Here $k_t$ denotes the residue field of $T$ at $t$.} and\pageoriginale consider a
commutative diagram of $T$-schemes}
$$\vcenter{\etiquette[6.3.23.3]{XVII.eq6.3.23.3}
\xymatrix{
              & Z \ar@{^{(}->}[ld]_{x} \ar@{^{(}->}[rd]^{y}     \\
X_t \ar[rr]^{u}&    & Y_t,
}}\leqno{(6.3.23.3)}
$$
\emph{such that $x$ and $y$ are closed immersions. Let $z\in
\Sym^{n}_{T}(Z)$, and let $x(\sisi{Z}{z})$ and $y(\sisi{Z}{z})$ be its images in
$\Sym^{n}_{T}(X)$ and $\Sym^{n}_{T}(Y)$. Then $\Sym^{n}_{T}(u)$ is
étale at $x(\sisi{Z}{z})$.}

The formation of the scheme $\Sym^{n}_{T}(X)$ is compatible with every
base change $\sisi{S'}{T'}\rightarrow \sisi{S}{T}$ (\hyperref[XVII.par5.5.2.7]{5.5.2.7}). By reduction to the
noetherian case, it follows that $\Sym^{n}_{T}(X)$ is of finite
presentation over $\sisi{S}{T}$; moreover, it is flat over $\sisi{S}{T}$ (\hyperref[XVII.par5.5.2.4]{5.5.2.4}). By the
fiberwise flatness criterion, to prove
(\hyperref[XVII.par6.3.23.2]{6.3.23.2}) we reduce to the case where $\sisi{S}{T}$ is the spectrum of a field. The completion of
$\Sym^{n}_{T}(X)$ at $\sisi{Z}{z}$ then has as its affine algebra the inverse
limit of the affine algebras localized at $x(z)$ of the schemes
$\Sym^{n}_{T}(Z_{k})$, where $Z_{k}$ ranges over the infinitesimal neighborhoods of $Z$
in $X$. By hypothesis, $u$ induces isomorphisms between the
infinitesimal neighborhoods of $Z$ in $X$ and those in $Y$. The completion of
$\Sym^{n}_{T}(u)$ at $x(z)$ and $y(z)$ is therefore an isomorphism, and
this proves (\hyperref[XVII.par6.3.23.2]{6.3.23.2}).

It follows from \hyperref[XVII.par6.3.23.2]{6.3.23.2} that the scheme given by the disjoint union, over $z\in
\Sym^{n}_{T}(Z)$, of the henselizations at $y(z)$ of $\Sym^{n}_{T}(Y)$
depends only on $Z$ and on the $T$-scheme $X^{h}_{Z}$ that is the disjoint union of the henselizations of
$X$ at the various points of $Z$. We shall denote this scheme by
$\Sym^{h\,n}_{T}(X^{h}_{Z})$; it is functorial in
$X^{h}_{Z}$.

\item By\pageoriginale passing to the limit, to prove \hyperref[XVII.prop6.3.23]{6.3.23}, it
suffices to show that for $y\in \Sym^{n}_{T}(Y)$, closed in its
fiber, the pullback to the henselization of $\Sym^{n}_{T}(Y)$ at that point of the
simplicial scheme $\Sym^{n}_{T}((X/Y)^{k+1})$ is a hypercover
for the étale topology. This problem is local on $T$ for the
étale topology; we thereby reduce to the case where there exists a diagram
\eqhyperref[XVII.eq6.3.23.3]{6.3.23.3} with $y=x(z)$ (after killing the residue-field extensions by a
separable extension of $k_{t}$). The semisimplicial system under consideration
is then obtained from the system
$$
\Sym^{h\,n}_{T}((X^{h}_{u^{-1}(Z)}/Y^{h}_{Z})^{k+1}),\quad\mbox{\rm
augmented to}\quad \Sym^{h\,n}_{T}(Y^{h}_{Z})
$$
by base change, so it suffices to consider the
latter. The semisimplicial system of
$(X^{h}_{u^{-1}(Z)}/Y^{h}_{Z})^{k+1}$ is semisimplicially
homotopic to the constant system $Y^{h}_{Z}$, because
$X^{h}_{u^{-1}(Z)}/Y^{h}_{Z}$ admits a section. The system
$\Sym^{h\,n}_{T}((X^{h}_{u^{-1}(Z)}/Y^{h}_{Z})^{h+1})$ is therefore likewise
homotopic to $\Sym^{h\, n}_{T} (Y^{h}_{Z})$, and in particular is
a hypercover of it. This proves \sisi{6.3.2.3}{\hyperref[XVII.prop6.3.23]{6.3.23}}.
\end{enumerate}
\end{enonce*}

\subsubsection{}\etiquette[6.3.24]{XVII.subsubsec6.3.24}
Let $f:Y\rightarrow T$ be a flat, quasi-projective morphism of
finite presentation between $S$-schemes, and let $K$ be a $G_{Y}$-torsor on
$Y$. Consider the following problem:

\setcounter{paragraph}{-1}
\paragraph{}\etiquette[6.3.24.0]{XVII.par6.3.24.0}
Define, for every morphism $u:X\rightarrow Y$, with $X$
flat, quasi-projective, and of finite presentation over $T$, a torsor
$(u^{\ast}K)_{n}$ on $\Sym^{n}_{T}(X)$, and for every $v:X'\rightarrow
X$ an isomorphism $c_{v}:(v^{\ast}u^{\ast}K)_{n}\simeq
\Sym^{n}_{T}(v)^{\ast}\sisi{K}{(u^K)}_{n}$, such that
\begin{enumerate}
\renewcommand{\theenumi}{\roman{enumi}}
\renewcommand{\labelenumi}{(\theenumi)}
\item the\pageoriginale isomorphisms $c_{v}$ are compatible with
  composition;

\item when $u^{\ast}K$ is locally trivial on $T$, one recovers the
torsor of \hyperref[XVII.lem6.3.22.1]{6.3.22.1} and its functoriality.
\end{enumerate}

By \hyperref[XVII.prop6.3.23]{6.3.23}, the existence and uniqueness of a solution to
\hyperref[XVII.par6.3.24.0]{6.3.24.0}, for every torsor obtained from $K$ by an inverse image of the
type considered in \hyperref[XVII.par6.3.24.0]{6.3.24.0}, is a local question for the
étale topology on $Y$. In particular, if $G$ satisfies (\hyperref[XVII.par6.3.21.1]{6.3.21.1}), the problem
\hyperref[XVII.par6.3.24.0]{6.3.24.0} has a unique solution.

Now suppose that $G$ satisfies:

\paragraph*{(6.3.24.1)}\MakeLinkTarget*{sga4.exp.XVII.paragraph.6.3.24.1}\label{XVII.par6.3.24.1}
\emph{For every integer $n$ and every diagram of affine
$S$-schemes
\[
\xymatrix{
X \ar[rr]^{u} \ar[rd]_{f}& & Y\ar[ld]^{g}\\
  &T&,
}
\]
with $f$ and $g$ flat and of finite presentation and $u$ finite,
locally free, and surjective, the truncated semisimplicial scheme
$(\Sym^{n}_{T}((X/Y)^{k+1}))_{0\leq k\leq 2}$, augmented to
$\Sym^{n}_{T}(Y)$, is a descent diagram \emph{(\textsc{Giraud} \cite{sga4.exp.XVII.bib.Gi})} for the
category of $G$-torsors.}

\paragraph*{6.3.24.1 bis}\etiquette[6.3.24.1 bis]{XVII.par6.3.24.1bis}
It follows in the standard way from \hyperref[XVII.prop6.3.23]{6.3.23} that
condition (\ref{XVII.par6.3.24.1})
implies the analogous condition in which one assumes only that $f$ and $g$ are
flat, quasi-projective, and of finite presentation, and that $u$ is flat and
surjective. Indeed, locally for the étale topology, $u$ admits
finite locally free quasi-sections, and one concludes by \hyperref[XVII.prop6.3.23]{6.3.23} and
\textsc{Giraud} \cite{sga4.exp.XVII.bib.Gi}, applied to a diagram forming a square of descent data.

Given\pageoriginale a $G_{Y}$-torsor $K$ on a $T$-scheme $Y$ that is
flat, quasi-projective, and of finite presentation, there exists by definition
a faithfully flat morphism of finite presentation $u:X\rightarrow Y$
such that $u^{\ast}K$ is trivial.\nde{One may assume that $T$ is affine. Then there exists $u$ as stated, which may even be assumed affine so that the preceding paragraph can subsequently be applied.} If $u_{i}$ denotes the projection from
$(X/Y)^{i+1}$ to $Y$ $(i\geq 0)$, the torsors $(u^{\ast}_{i}K)_{n}$
(\hyperref[XVII.lem6.3.22.1]{6.3.22.1}) on the schemes $\Sym^{n}_{T}((X/Y)^{i+1})$ form a
descent datum $(i=0,1,2)$ on this semisimplicial scheme augmented
to $\Sym^{n}_{T}(Y)$. Suppose that $G$ satisfies the condition:

\paragraph*{(6.3.24.2)}\MakeLinkTarget*{sga4.exp.XVII.paragraph.6.3.24.2}\label{XVII.par6.3.24.2}
\emph{The descent data constructed above are all
effective} (\cf \ref{XVII.par6.3.24.1}).

Then we denote by $K_{n}$ the torsor on $\Sym^{n}_{T}(Y)$
(determined up to unique isomorphism) defined by this descent
datum.

\subsubsection{}\etiquette[6.3.25]{XVII.subsubsec6.3.25}
Suppose that $G$ satisfies (\hyperref[XVII.par6.3.21.1]{6.3.21.1}) or conditions
(\ref{XVII.par6.3.24.1}) and (\ref{XVII.par6.3.24.2}). We leave it to the
reader to make sense of the following formalism and to verify it:

\paragraph{}\etiquette[6.3.25.1]{XVII.par6.3.25.1}
\emph{Base change}. Let
\[
\xymatrix{
X' \ar[rr]^{g'} \ar[d]_{f'}& & X \ar[d]^{f}\\
T' \ar[rr]_{g}& & T.
}
\]
be a commutative square. Denote by $g'_{n}$ the arrow from $\Sym^{n}_{T'}(X')$ to
$\Sym^{n}_{T}(X)$. For $X$ a torsor under $G$ on $X$, one has
\[
{g'_n}^{\ast} K_{n}\simeq ({g'}^{\ast}K)_{n}{\quoi}.
\]

Now let\pageoriginale $f:X\rightarrow T$ be a morphism, and let $K$ be a
$G_{X}$-torsor on $X$.

\paragraph{}\etiquette[6.3.25.2]{XVII.par6.3.25.2}
\emph{Change of group.} Let $r:G_{1}\rightarrow G_{2}$ be a
morphism between groups both satisfying (\hyperref[XVII.par6.3.21.1]{6.3.21.1})
or (\ref{XVII.par6.3.24.1}), (\ref{XVII.par6.3.24.2}). One has
\[
r(K_{n})\simeq r(K)_{n}{\quoi}.
\]

\paragraph{}\etiquette[6.3.25.3]{XVII.par6.3.25.3}
\emph{Additivity in $K$}\sisi{:}{.} One has
\[
(K'+K'')_{n}\simeq K'_{n}+K''_{n}{\quoi}.
\]

\paragraph{}\etiquette[6.3.25.4]{XVII.par6.3.25.4}
\emph{Additivity in $n$}. Let $n_{1}$ and $n_{2}$ be two integers,
$n=n_{1}+n_{2}$, and let
\[
\alpha:\fprod{\Sym_{T}{}^{n_{1}}(X)}{\Sym_{T}{}^{n_{2}}(X)}{T}\longrightarrow
\Sym^{n}_{T}(X)
\]
be the canonical arrow; then
\[
\alpha^{\ast}K_{n}\simeq \pr^{\ast}_{1}
K_{n_{1}}+\pr^{\ast}_{2}K_{n_{2}}{\quoi}.
\]

\paragraph{}\etiquette[6.3.25.5]{XVII.par6.3.25.5}
\emph{Multiplicativity in $n$.} Let $n$ and $m$ be two integers, and let
$\alpha$ be the canonical arrow from $\Sym^{n}_{T}(\Sym^{m}_{T}(X))$ to
$\Sym^{nm}_{T}(X)$ (\cf \hyperref[XVII.lem6.3.3]{6.3.3}). One has
\[
\alpha^{\ast}(K_{nm})\simeq (K_{m})_{n}{\quoi}.
\]

\subsubsection{}\etiquette[6.3.26]{XVII.subsubsec6.3.26}
Suppose that $G$ satisfies (\hyperref[XVII.par6.3.21.1]{6.3.21.1}) or
(\ref{XVII.par6.3.24.1}),
(\ref{XVII.par6.3.24.2}). If
$f:X\rightarrow T$ is a finite locally free morphism of rank $n$
between $S$-schemes, if $t$ is the section \hyperref[XVII.eq6.3.4.2]{(6.3.4.3)} of
$\Sym^{n}_{T}(X)$, and if $K$ is a $G_{X}$-torsor on $X$, we
define a $G_{T}$-torsor $\Tr_{f}(K)$ on $T$ by the formula
$$\etiquette[6.3.26.1]{XVII.eq6.3.26.1}
\Tr_{f}(K)=t^{\ast}(K_{n}){\quoi}.\leqno(6.3.26.1)
$$

If\pageoriginale $f$ is locally free, we define $\Tr_{f}(K)$
locally on $T$ by the same formula. This definition, via a
canonical and functorial isomorphism, is equivalent to the definition
\hyperref[XVII.prop6.3.21]{6.3.21} on their common domain of validity
(the definition \hyperref[XVII.prop6.3.21]{6.3.21}
providing a solution to the descent problem posed). It follows
immediately from (\hyperref[XVII.par6.3.25.1]{6.3.25.1}), (\hyperref[XVII.par6.3.25.2]{6.3.25.2}),
and (\hyperref[XVII.par6.3.25.3]{6.3.25.3}) that the functor
$\Tr_{f}$ is \emph{compatible with base changes $T'\rightarrow
T$, compatible with changes of group $r:G_{1}\rightarrow G_{2}$,
and compatible with addition of torsors.}

The commutativity of the outer boundary of diagram \eqhyperref[XVII.eq6.3.15.2]{6.3.15.2}, together with
(\hyperref[XVII.par6.3.25.1]{6.3.25.1}) and (\hyperref[XVII.par6.3.25.4]{6.3.25.4}), yields
\emph{compatibility with composition of morphisms}:
$$
\Tr_{fg}(K)\simeq \Tr_{f}(\Tr_{g}(K)\sisi{}{)}{\quoi}.
$$

\subsubsection{}\etiquette[6.3.27]{XVII.subsubsec6.3.27}
In the particular case of relative curves, one also has:

\begin{enonce*}{Construction 6.3.27.1}\etiquette[6.3.27.1]{XVII.const6.3.27.1}
\emph{If} $f:X\rightarrow S$ is a smooth quasi-projective curve
over $S$, $D_{1}$ and $D_{2}$ are two relative divisors finite over $S$, and
$K$ is a $G$-torsor on $X$, then, setting
$\Tr_{D/S}(K)=\Tr_{D/S}(K|D)$, there is a canonical isomorphism
$$
\Tr_{D_{1}+D_{2}}(K)\simeq \Tr_{D_{1}}(K)+\Tr_{D_{2}}(K){\quoi}.
$$
\end{enonce*}

\sisi{\noindent}{}%
This isomorphism follows immediately from the isomorphism (\hyperref[XVII.par6.3.25.4]{6.3.25.4}) and
\hyperref[XVII.lem6.3.9.1]{6.3.9.1}. Finally, the arguments of \hyperref[XVII.prop6.3.19]{6.3.19} also give:

\begin{enonce*}{Construction 6.3.27.2}\etiquette[6.3.27.2]{XVII.const6.3.27.2}
Consider a commutative diagram
$$
\xymatrix{
X   \ar[rr]^{u} \ar[rd]_{p}&     &  Y \ar[ld]^{q}\\
     &   S &,
}
$$
in\pageoriginale which $X$ and $Y$ are smooth
quasi-projective curves over $S$ and $u$ is a quasi-finite morphism.
\begin{enumerate}
\renewcommand{\theenumi}{\roman{enumi}}
\renewcommand{\labelenumi}{(\theenumi)}
\item If $D$ is a divisor on $X$ finite over $S$ and $K$ is a torsor
under $G$ on $Y$, one has \emph{(\cf note p.~\pageref{XVII.note192})}:
$$
\Tr_{D/S}(u^{\ast}K)\simeq \Tr_{(u_{\ast}D)/S}(K){\quoi}.
$$

\item If $E$ is a divisor on $Y$ finite over $S$, $K$ is a torsor under
$G$ on $X$, and $u$ is finite, one has
$$
\Tr_{E/S}(\Tr_{X/Y}(K))\simeq \Tr_{u^{\ast}E/S}(K){\quoi}.
$$
\end{enumerate}
\end{enonce*}

\hypertarget{sga4.exp.XVII.proposition.6.3.28}{}\begin{enonce*}{Proposition 6.3.28}\etiquette[6.3.28]{XVII.prop6.3.28}
Let $A$ and $B$ be two (commutative) $R$-algebras and let $u:A\rightarrow
B$ be a homomorphism making $B$ a faithful, finite-type, locally free
$A$-module. Then the \sisi{}{co-}semisimplicial complex of algebras
$(TS^{n}_{R}(\bigotimes\limits^{k+1}_{A}B))_{k\geq 0}$
is, as a differential \sisi{}{co-}semisimplicial $TS^{n}_{R}(A)$-module,
a split resolution of $TS^{n}_{R}(A)$.
\end{enonce*}

\begin{enonce*}[remark]{Proof}
The complex of modules associated with the \sisi{}{co-}semisimplicial module
$TS^{n}_{R}(\bigotimes\limits^{k+1}_{A}B)$ depends only
on the $A$-module $B$, equipped with a module homomorphism $u:A\rightarrow
B$. As an $A$-module, $B$ can be written $B=A\oplus L$, and this
makes it possible to write down an explicit homotopy operator.
\end{enonce*}

Recall that:

\hypertarget{sga4.exp.XVII.proposition.6.3.29}{}\begin{enonce*}{Proposition 6.3.29}\etiquette[6.3.29]{XVII.prop6.3.29}
Let $A$ be a ring, let $B^{\ast}$ be a co-semisimplicial $A$-algebra,
and let $c:A\rightarrow B^{\ast}$ be the structure morphism. For the
inverse-image functor on modules ($=$ extension of scalars)
$\epsilon^{\ast}$ to be faithful (\resp fully faithful), it suffices that
the sequence of $A$-modules
$$
0\longrightarrow A\longrightarrow B^{0}
$$
(\resp\pageoriginale $0\rightarrow A\rightarrow
B^{0}\xrightarrow{j_{0}-j_{1}}B^{1}$)
be exact, and remain so after tensoring with any $A$-module.
\end{enonce*}

The proof is immediate and is left to the reader.

It follows in the standard way from (\hyperref[XVII.prop6.3.28]{6.3.28}) and (\hyperref[XVII.prop6.3.29]{6.3.29}) that

\hypertarget{sga4.exp.XVII.corollary.6.3.30}{}\begin{enonce*}{Corollary 6.3.30}\etiquette[6.3.30]{XVII.cor6.3.30}
Condition (\ref{XVII.par6.3.24.1}) is satisfied by every affine group $G$.
\end{enonce*}

(Use the fact that a torsor under $G$ is representable by a scheme
affine over $T$.)

It should not be difficult to extend \hyperref[XVII.cor6.3.30]{6.3.30} to every group $G$
representable by an algebraic space \footnote[1]{\MakeLinkTarget*{sga4.exp.XVII.footnote.6.3.30.artin-2}\label{XVII.noteArtin2}Algebraic Spaces
by M.~\textsc{Artin}. Yale University, Whittemore Lectures, May 1969.}. I
conjecture that condition (\ref{XVII.par6.3.24.2}) is satisfied by every group
representable by a flat algebraic space. We have the tongue-in-cheek
statement:

\hypertarget{sga4.exp.XVII.proposition.6.3.31}{}\begin{enonce*}{Proposition 6.3.31}\etiquette[6.3.31]{XVII.prop6.3.31}
Let the following sequence be exact on $S$:
$$
0\rightarrow G\xrightarrow{r}G'\xrightarrow{u}H\rightarrow 0{\quoi}.
$$ Suppose that $G'$ and $H$, and hence $G$, satisfy
(\hyperref[XVII.par6.3.13.1]{6.3.13.1}), that $G'$ satisfies (\hyperref[XVII.par6.3.21.1]{6.3.21.1}), and
that $G$ satisfies (\ref{XVII.par6.3.24.1}). Then $G$ satisfies
(\ref{XVII.par6.3.24.2}).
\end{enonce*}

This applies, for example, to an affine group that is the kernel of an epimorphism
of representable smooth groups.

% END INCLUDED SGA4_English_current_buildable_source_20260906/source/17/drafts/SGA4_XVII_LINES_8201_8600_English.texfrag

% BEGIN INCLUDED SGA4_English_current_buildable_source_20260906/source/17/drafts/SGA4_XVII_LINES_8601_9000_English.texfrag
\begin{enonce*}{Proof of 6.3.31}
\emph{Let\pageoriginale $K$ be a $G$-torsor on
$X$. By \hyperref[XVII.prop6.3.23]{6.3.23}, the question is local on $X$ for the
étale topology; it therefore suffices to treat the case where $r(K)$ is trivial. The
$G$-torsors $K$ such that $r(K)$ is trivial are identified with
trivial $G'$-torsors $K'$ equipped with an isomorphism
$k:H\isomto u(K')$. Put $Z=\Sym^{n}_{T}(X)$. The morphism
of trivial torsors $k$ induces $k_{n}:H_{Z}\isomto
u(K'_{n})$, and the corresponding $G$-torsor on $\Sym^{n}\sisi{}{_T}(X)$
represents the descent datum under consideration.}
\end{enonce*}

\section{Appendix}\etiquette[7]{XVII.sec7}\pageoriginale

\setcounter{subsection}{-1}
\subsection{Preliminaries}\etiquette[7.0]{XVII.subsec7.0}

In this appendix we shall use the results of \sisi{${\rm VI_{B}}$}{VBIS},
whose notation we retain.

\subsubsection{}\etiquette[7.0.1]{XVII.subsubsec7.0.1}
Let $\Sch_{!}$ be the category whose objects are quasi-compact,
quasi-separated schemes and whose morphisms are separated morphisms
of finite type. We denote by $\Ee$ the bifibred category over
$\Sch_{!}$, with fibre categories opposite to topoi, defined by étale
sheaves on schemes \sisi{${\rm VI_{B}}$ (4.3.0)}{(\hyperrefext[Vbis][Vbis.subsubsec4.3.0]{4.3.0})}. Fix once and for all an
object $R$ of $\Sch_{!}$, and denote by $\Ee_{R}$ the
bifibred category over $\Sch_{!R}$ deduced from
$\Ee$. Finally, let $\mathcal{A}$ be a ring in the topos
$\sheafGamma(\Ee^{\circ}_{R})$ satisfying the following conditions:
\begin{enumerate}
\renewcommand{\theenumi}{\roman{enumi}}
\renewcommand{\labelenumi}{(\theenumi)}
\item $\mathcal{A}$ is a \emph{cocartesian} section of
$\Ee^{\circ}_{R}$ over $\Sch^{\circ}_{!R}$.

\item For every scheme $X$ in $\Sch_{!R}$, $\mathcal{A}_{X}$ is a
\emph{torsion} sheaf.
\end{enumerate}

With this understood, we shall work with the category
$\sheaf{\Mod}(\sisi{}{\sheafGamma(}\Ee^{\circ}_{R}\sisi{}{)},\mathcal{A})$ (\refvb[VI B ]{Vbis.subsubsec1.3.5}).

Our aim is to define, for every morphism $f:X\rightarrow Y$ in
$\Sch_{!R}$, a functor $Rf_{!}:D(X,\mathcal{A}_X)\rightarrow
D(Y,\mathcal{A}_{Y})$, agreeing with the previously defined functor
when $f$ is compactifiable, and making the categories
$D(X,\mathcal{A}_{X})$, for varying $X$, into a cofibred category
over $\Sch_{!R}$.\nde{\MakeLinkTarget*{sga4.exp.XVII.editor-note.7.0.1.compactification}\label{XVII.ndeapp}Every morphism in $\Sch_{!}$ is compactifiable by Nagata's compactification theorem (see \nderef{XVII.ndeConrad}). This appendix shows how to avoid using the compactification theorem in the definition of $Rf_!$ by relying on Chow's lemma (\hyperref[XVII.lem7.0.6]{7.0.6}).}

\subsubsection{}\etiquette[7.0.2]{XVII.subsubsec7.0.2}
Although all the results of this appendix are stated in the setting
specified in \hyperref[XVII.subsubsec7.0.1]{7.0.1}, they could be transcribed, with obvious
modifications, into the following setting: take as a section of
$\sheafGamma(\Ee^{\circ}_{R})$ the constant section whose value is
the constant sheaf $\ZZ$, and, for every scheme $X$, restrict to the
derived category $D^{\ast}_{\tors}(X)$ defined by complexes of sheaves
of abelian groups whose cohomology sheaves are torsion sheaves
(\refvb[VI B ]{Vbis.subsubsec4.3.1}).

\begin{enonce*}{Lemma 7.0.3}\etiquette[7.0.3]{XVII.lem7.0.3}
Let\pageoriginale there be given a commutative diagram of schemes
$$\vcenter{\etiquette[7.0.3.1]{XVII.eq7.0.3.1}
\xymatrix{
X' \ar[rr]^{j'} \ar[d]_{f'}& & X \ar[d]^{f}\\
Y' \ar[rr]_{j}& & Y.
}}\leqno{(7.0.3.1)}
$$
where $f$ and $f'$ are proper and $j$ and $j'$ are immersions. Then,
for every complex $K^{\bigdot}$ in $D^{+}_{\tors}(X')$, the canonical
arrow\nde{The definition of this arrow is similar to that of \eqhyperref[XVII.eq5.1.5.2]{5.1.5.2}, up to an easy dévissage.}
$$
j_{!}\circ R^{+}f'\sisi{}{_*}(K^{\bigdot})\longrightarrow R^{+}f\sisi{}{_*}\circ
j_{!}(K^{\bigdot})
$$
is an isomorphism.
\end{enonce*}

By the ``lemma on way-out functors'' (\cite{sga4.exp.XVII.bib.RD} I \sisi{(}{}7.1), it
suffices to show that, for every integer $i$ and every torsion sheaf
$F$, the arrow
$$\etiquette[7.0.3.2]{XVII.eq7.0.3.2}
j_{!}\circ R^{i}f'\sisi{}{_*}(F)\longrightarrow R^{i}f\sisi{}{_*}\circ j_{!}(F)\leqno(7.0.3.2)
$$
is an isomorphism.

By replacing $j$ and $j'$ with their closed images, we are reduced to
the case where $j$ and $j'$ are open immersions with everywhere dense
images. The properness of $f$ and $f'$ then implies that diagram
\eqhyperref[XVII.eq7.0.3.1]{7.0.3.1} is cartesian. Under these conditions, the proper base-change
theorem, in the form \hyperrefext[XII][XII.cor5.2]{5.2}, shows that, for every geometric point
$\xi$ of $Y$, the arrow
$(j_{!}\circ R^{i}f'\sisi{}{_*}(F))_{\xi}\rightarrow (R^{i}f\sisi{}{_*}\circ
j_{!}(F))_{\xi}$ is an isomorphism.

\begin{enonce*}{Construction 7.0.4}\etiquette[7.0.4]{XVII.const7.0.4}
\emph{Let $D$ be a category, and let $X$ and $Y$ be two $D$-objects in
$\Sch_{!R}$ such that, for every arrow $\alpha\rightarrow \beta$ in
$D$, the morphisms $X_{\beta}\rightarrow X_{\alpha}$ and
$Y_{\beta}\rightarrow Y_{\alpha}$ are proper. Let $j:X\rightarrow
Y$ be a morphism such that, for every object $\alpha$ of $D$,
$j_{\alpha}:X_{\alpha}\rightarrow Y_{\alpha}$ is an immersion. The
isomorphism \eqhyperref[XVII.eq7.0.3.2]{7.0.3.2} for $i=0$ satisfies the usual compatibility
conditions and allows one to construct an \emph{exact} functor:}
\[
j_{!}:\Mod(\sheafGamma(\bar{X}),\mathcal{A})\longrightarrow
\Mod(\sheafGamma(\bar{Y}),\mathcal{A})
\]
\emph{(\cf \sisi{(}{}\refvb{Vbis.subsubsec1.2.5} for the notation)\sisi{.}{}
which\pageoriginale induces a triangulated functor, again denoted $j_{!}$:}
\[
j_{!}:D(\sheafGamma(\bar{X}),\mathcal{A})\longrightarrow
D(\sheafGamma(\bar{Y},\mathcal{A}))
\]

\emph{Then, for every object $\alpha$ of $D$, there is a commutative
diagram, where $j_{\alpha !}$ denotes the usual ``extension by $0$''
functor:}
\[
\xymatrix{
D(\sheafGamma(\bar{X}), \sisi{}{\Aa}) \ar[rr]^{j_!} \ar[d]_{R e^{\sisi{+}{*}}_{\alpha}}&
&D(\sheafGamma(\bar{Y}),\Aa) \ar[d]^{R e^{\sisi{+}{*}}_{\alpha}}\\
D(X_{\alpha}, \Aa_{X_{\alpha}})  \ar[rr]_{j_{\alpha !}} & &
  D(Y_{\alpha},\Aa_{Y_\alpha}).
}
\]
\end{enonce*}

\begin{enonce*}{Construction 7.0.5}\etiquette[7.0.5]{XVII.constru7.0.5}
Let there be given a commutative diagram in $\Sch_{!R}$
$$\vcenter{\etiquette[7.0.5.1]{XVII.eq7.0.5.1}
\xymatrix{
X'   \ar[rr]^{i'} \ar[d]_{f'}&    & X \ar[d]^{f}\\
Y'   \ar[rr]^{i} &    & Y
}}\leqno{(7.0.5.1)}
$$
where $i$ and $i'$ are immersions, $f$ is a proper morphism, and $f'$
is a proper \emph{surjective} morphism. It defines a commutative diagram
of simplicial objects in $\Sch_{!R}$
$$\vcenter{\etiquette[7.0.5.2]{XVII.eq7.0.5.2}
\xymatrix{
[X'|Y']      \ar[d]_{u_{f'}}  \ar[r]^{j} & [X|Y] \ar[d]^{u_f}\\
C^{\Delta}_{Y'}  \ar[r]_{i} & C^{\Delta}_{Y}
}}\leqno{(7.0.5.2)}
$$
where $j_{n}:[X'|Y']_{n}\rightarrow [X|Y]_{n}$ is an immersion
for every $n$. With this notation, there is a canonical isomorphism
\[
i_{!}\isomto  R^{+}\bar{u}_{f\sisi{+}{*}}\circ j_{!}\circ
L^{+}\bar{u}_{f'}\sisi{}{^*}{\quoi}.
\]
\end{enonce*}

By replacing $i$ and $i'$ with their closed images, we may suppose that
$i$ and $i'$ are open immersions with everywhere dense images. Diagram
\eqhyperref[XVII.eq7.0.5.1]{7.0.5.1} is then cartesian and $f$ is \emph{surjective}. Moreover,
for every integer $n$, $j_{n}$ is an\pageoriginale open immersion and
the diagram
\[
\xymatrix{
[X'|Y']_n        \ar[r]^{j_n} \ar[d]& [X|Y]_n \ar[d]\\
Y'  \ar[r]_{i} & Y
}
\]
is cartesian.

Consequently, $L^{+}\bar{u}_{f}^{\sisi{+}{*}}\circ i_{!}$ is isomorphic to $j_{!}\circ
L^{+}\bar{u}_{f'}^{\ast}$. Using the fact that $u_{f}$ is a
cohomological-descent augmentation (\sisi{${\rm VI_{B}}$ (\refnoexpose{Vbis.prop4.3.2})}{\refvb{Vbis.prop4.3.2}}), we
obtain a sequence of isomorphisms
\[
i_{!}\isomto R^{+}\bar{u}_{f\sisi{+}{*}}\circ
L^{+}(\bar{u}^{\sisi{+}{*}}_{f})\circ i_{!}\isomto
R^{+}\bar{u}_{f\ast}\circ j_{!}\circ L^{+}\bar{u}_{f'}^{\ast}
\]
which gives the desired isomorphism.\nde{More directly, without assuming $f'$ surjective, there is still a sequence of canonical arrows
\[i_!\to i_! R^+\bar{u}_{f'*}L^+\bar{u}^*_{f'}\isomto  R^+\bar{u}_{f*} j_! L^+u^*_{f'},\]
of which the second is an isomorphism (\hyperref[XVII.lem7.0.3]{7.0.3}). When $f'$ is surjective, the first is also an isomorphism by cohomological descent, and the composite agrees with the isomorphism constructed in the text.
}

Finally, we shall use the following variant of Chow's lemma, whose
proof will appear in \EGA II, 2\sisi{nd}{\textsuperscript{nd}} edition,\nde{See [C, 2.6], cited in \nderef{XVII.ndeConrad}.} and which is an
improvement of (\sisi{\hyperref[XVII.subsec7.1]{7.1}}{\hyperrefext[XII][XII.lem7.1]{7.1}}):

\begin{enonce*}{Lemma 7.0.6}\etiquette[7.0.6]{XVII.lem7.0.6}
Let $S$ be a quasi-compact and quasi-separated scheme and let
$f:X\rightarrow S$ be a separated morphism of finite type. Then there
exists a commutative diagram
\[
\xymatrix{
X' \ar[rr]^{p} \ar[d]_{f'}& & X \ar[lld]^{f}\\
S                         & &
}
\]
satisfying the following conditions:
\begin{enumerate}
\renewcommand{\theenumi}{\roman{enumi}}
\renewcommand{\labelenumi}{(\theenumi)}
\item $p$ is proper and \emph{surjective};

\item there is a dense open subset $U$ of $X$ such that $p$ induces an
isomorphism from $p^{-1}(U)$ onto $U$;

\item $f'$ is quasi-projective.
\end{enumerate}
\end{enonce*}


\subsection{Definition of $R^{+}f_{!}$}\etiquette[7.1]{XVII.subsec7.1}\pageoriginale

\begin{enonce*}{Definition 7.1.1}\etiquette[7.1.1]{XVII.def7.1.1}
Let $f:X\rightarrow S$ be a morphism of schemes. A
\emph{pseudo-compactification} of $f$ consists of a commutative diagram
$$\vcenter{
\xymatrix{
X_1 \ar[r]^{\sisi{}{j_1}} \ar[d]_{p_1}& T_1\ar[ldd]^{f_1}\\
X   \ar[d]_{f}& \\
S   &
}}\leqno{(7.1.1.0)}
$$\etiquette[7.1.1.0]{XVII.eq7.1.1.0}
where $p_{1}$ is proper and surjective, $j_{1}$ is an immersion, and
$f_{1}$ is proper.
\end{enonce*}

\setcounter{subsubsection}{1}
\subsubsection{}\etiquette[7.1.2]{XVII.subsubsec7.1.2}
Let $f:X\rightarrow S$ be a morphism in $\Sch_{!R}$. By Chow's lemma
(\hyperref[XVII.lem7.0.6]{7.0.6}), $f$ has a pseudo-compactification
$(p_{1},j_{1},f_{1})$. Diagram \eqhyperref[XVII.eq7.1.1.0]{7.1.1.0} defines a morphism
\[
[X_{1}|X]\xrightarrow{j_{1}}[T_{1}|S]
\]
again denoted by $j_{1}$, such that $j_{1n}$ is an immersion for every
$n$. We then put $R^{+(1)}f_{!}=\sisi{L^{+}\bar{u}_{p_{1}}\cdot
j_{1!}\circ R^{+}\bar{u}_{f_{1}}}{R^+\bar{u}_{f_1*}\circ j_{1!} \circ L^+\bar{u}_{p_1}^*}$, thereby obtaining a
\emph{triangulated} functor from $D^{+}(X,\mathcal{A}_{X})$ to
$D^{+}(S,\mathcal{A}_{S})$.

\begin{enonce*}{Construction 7.1.3}\etiquette[7.1.3]{XVII.constru7.1.3}
For every pair $((p_{1},j_{1},f_{1}),(p_{2},j_{2},f_{2}))$ of
pseudo-compactifications of $f$, there is an isomorphism
$\alpha_{12}:R^{+(1)}f_{!}\rightarrow R^{+(2)}f_{!}$ such that, for
every other pseudo-compactification $(p_{3}, j_{3},f_{3})$, one has
$\alpha_{23}\circ \alpha_{12}=\alpha_{13}$.
\end{enonce*}

\begin{enumerate}
\renewcommand{\theenumi}{\alph{enumi}}
\renewcommand{\labelenumi}{\theenumi)}
\item Define a bisimplicial system $[[X_{1},X_{2}|X]]$ by setting
\[
[[X_{1},X_{2}|X]]_{[n]\times [m]}=[X_{1}|X]_{n}\times
[X_{2}|X]_{n}
\]
Similarly, set
\[
[[X_{1}|X]]_{[n]\times[m]}=[X_{1}|S]_{n}{\quoi},
\]
all the morphisms $[[X_{1}|X]]_{[n]\times [m]}\rightarrow
[[X_{1}|X]]_{[n]\times [m']}$ induced by morphisms
$[m']\rightarrow [m]$ being equal to the identity. Define the
bisimplicial systems $[[T_{1},T_{2}|S]]$ and $[[T_{1}|S]]$ in the same
way, so that there is a commutative diagram
\[
\xymatrix@C=1em{
  & & [[T_1,T_2|S]] \ar[rrrr]\ar'[dr][rrrddd]_(.2){u_{f_1f_2}}|(.38)\hole& &   & &  [[T_1|S]]\ar[lddd]^{u_{f_1}}\\
  [[X_1,X_2|X]]
  \ar[rrrr]\ar[rru]^{j_{12}}\ar[rrrddd]_{u_{p_1 p_2}} & &  & & [[X_1|X]]\ar[rru]_{j_1}\ar[lddd]_(.75){u_{p_1}}  \\
  & & & & &   &\\
  & & & & & C^{\Delta\times\Delta}_S\\
  &&& C^{\Delta\times\Delta}_X \ar[rru]
}
\]
where\pageoriginale $j_{12}$ is an immersion degree by degree. The
computational procedure described in \sisi{VI b (\refnoexpose{Vbis.subsec2.3})}{\hyperrefext[Vbis][Vbis.subsec2.3]{2.3}} makes it possible to construct a
morphism
\[
R^{+(1)} f_{!}\xrightarrow{\alpha_{1}}R^{+}\bar{u_{f_{1}f_{2}}}\sisi{}{_*}\circ
j_{12!} \circ L^{+}\bar{u_{p_{1}p_{2}}}^{\sisi{+}{\ast}}=R^{{+}(1,2)}f_{!}
\]
which is obtained from a morphism of triple complexes of sheaves on
$S$.

\item $\alpha_{1}$ \emph{is an isomorphism}:

To prove this, one may fix the first index $n$ coming from
$\bigl[[T_{1},T_{2}]\bigr]_{[n]\times[m]}$ in the preceding morphism of
triple complexes and consider the resulting morphism of double
complexes. One checks that this morphism is precisely the isomorphism
defined in Lemma \hyperref[XVII.constru7.0.5]{7.0.5} for the diagram
\[
\xymatrix{
**[r][X_1|X]_n \times X_2  \ar[r] \ar[d] & **[r][T_1|S]_n \times
T_2 \ar[d]\\
**[r][X_1|X]_n             \ar[r] & **[r][T_1|S]_n
}
\]

\item In the same way, construct an isomorphism
\[
R^{+(2)}f_{!}\xrightarrow{\alpha_{2}}R^{+}\bar{u_{f_{1}f_{2}}}_{\sisi{+}{*}}\circ
j_{12!} \circ L^{+}\bar{u_{p_{1}p_{2}}}\sisi{}{^*}
\]
and set $\alpha_{12}=\alpha^{-1}_{2}\circ \alpha_{1}$.
\end{enumerate}

To verify the relation $\alpha_{23}\circ \alpha_{12}=\alpha_{13}$,
one uses a trisimplicial object. We leave it to the reader to check
that the resulting diagram is commutative:\pageoriginale
$$
\xymatrix{
  R^{+(2,3)}{f_!} \ar[rd]& R^{+(2)}f_! \ar[r] \ar[l]& R^{+(1,2)}f_! \ar[ld]\\
                &         R^{+(1,2,3)}f_!  &                  \\
R^{+(\sisi{2}{3})}f_! \ar[uu]\sisi{}{\ar[rd]}     &                         &  R^{+(1)}f_! \ar[uu]\sisi{}{\ar[ld]}\\
                &    R^{+(1,3)}f_! \ar[uu]        &
}
$$

\begin{enonce*}{Construction 7.1.4}\etiquette[7.1.4]{XVII.constru7.1.4}
Let $Z\xrightarrow{g}Y$ and $Y\xrightarrow{f}X$ be two morphisms in
$\Sch_{!R}$. For every triple $((p_{1},j_{1},f_{1}),
(p_{2},j_{2}g_{2}), (p_{3},j_{3},(fg)_{3}))$, where
$(p_{1},j_{1},f_{1})$ is a pseudo-compactification of $f$,
$(p_{2},j_{2},g_{2})$ a pseudo-compactification of $g$, and
$(p_{3},j_{3},(fg)_{3})$ a pseudo-compactification of $fg$, there is an
isomorphism $\sisi{R^{+(3)}fg_{!}\isomfrom  R^{+(2)}g_{!}\circ
R^{+(1)}f_{!}}{R^{+(1)}f_!\circ R^{+(2)}g_!\isomto R^{+(3)}(fg)_!}$ compatible with the previously defined
isomorphisms $\alpha_{ij}$.
\end{enonce*}

We proceed in several steps.

\setcounter{subsubsection}{4}
\paragraph{}\etiquette[7.1.4.1]{XVII.par7.1.4.1}
Let $(p_{1},j_{1},f_{1})$ (\resp $(p_{2},j_{2},f_{2})$) be a
pseudo-compactification of $f:Y\rightarrow X$ (\resp of
$g:Z\rightarrow Y$):
$$\vcenter{\etiquette[7.1.4.1.1]{XVII.eq7.1.4.1.1}
\xymatrix{
  &     & Z\ar[d]^{g}  & Z_2\ar[l]_{p_2}\ar@{^{(}->}[r]^{j_2}  & T_2 \ar[lld]^{g_2}\\
T_1  \ar[rrd]_{f_1}& Y_1 \ar[r]^{p_1} \ar@{_{(}->}[l]_{j_1}& Y \ar[d]^{f} &      &    \\
    &     & X        &      &
}}\leqno{(7.1.4.1.1)}
$$

The morphism $Z'=\fprod{Z}{Y_{1}}{Y}\rightarrow Y_{1}$ obtained from
$g$ by base change is equipped with a pseudo-compactification:
$$\vcenter{\etiquette[7.1.4.1.2]{XVII.eq7.1.4.1.2}
\xymatrix{
  T'_2 = \fprod{T_2}{Y_1}{Y}   \ar[rdd]_{g'_2}& \fprod{Z_2}{Y_1}{Y} =Z'_2 \ar[d]^{p'_2}
\ar@{_{(}->}[l]_{j'_2}\\
&  \fprod{Z}{Y_1}{Y} = Z' \ar[d]^{g'}\\
&    Y_1
}}\leqno{(7.1.4.1.2)}
$$


Now\pageoriginale choose a \sisi{}{pseudo-}compactification of
\sisi{$j_{1}\cdot g'_{2}$.
$$\etiquette[7.1.4.1.3]{XVII.eq7.1.4.1.3}
\vcenter{
\xymatrix{
T_{12}        \ar[d]_{g''_2} & \fprod{T_2}{Y_1}{Y}=T'_2
\ar@{_{(}->}[l]_(.6){j'_1} \ar[d]^{g'_2}\\
T_1            &   Y_1 \ar@{_{(}->}[l]
}}\leqno(7.1.4.1.3)
$$%
}{$j_1\circ g'_2$\nde{In the original text a compactification was chosen, which is possible by the compactification theorem. Since the purpose of the appendix is to avoid that theorem (\cf \nderef{XVII.ndeapp}), the editor has taken the liberty of replacing it with a pseudo-compactification and slightly modifying \hyperref[XVII.par7.1.4.2]{7.1.4.2}. Another solution would be to restrict to pseudo-compactifications $(p_1,j_1,f_1)$ for which $f_1$ is \emph{projective}.}
$$\etiquette[7.1.4.1.3]{XVII.eq7.1.4.1.3}
\vcenter{
\xymatrix{
T_{12}        \ar[dd]_{g''_2} &T_2''\ar[d]\ar[d]^p\ar@{_{(}->}[l]_{j'_1}\\
& \fprod{T_2}{Y_1}{Y}=T'_2\ar[d]^{g'_2}\\
T_1            &   Y_1 \ar@{_{(}->}[l]_{j_1}
}}\leqno(7.1.4.1.3)
$$}%
which defines a pseudo-compactification of the composite $fg$; denote
the corresponding functor by $R^{+(12)}\sisi{\int g}{(fg)}_{!}$.

\paragraph{}\etiquette[7.1.4.2]{XVII.par7.1.4.2}
\emph{The preceding data make it possible to define an isomorphism}
$R^{+(12)}\sisi{\int g}{(fg)}_{!}\isomfrom  R^{+(1)}f_{!}\sisi{\cdot}{\circ} R^{+(2)}g_{!}$:

There is a commutative diagram, where $Z_2''=\fprod{T''_2}{Z'_2}{T'_2}$ and all the arrows are
naturally deduced from \eqhyperref[XVII.eq7.1.4.1.1]{7.1.4.1.1}, \eqhyperref[XVII.eq7.1.4.1.2]{7.1.4.1.2}, and
\eqhyperref[XVII.eq7.1.4.1.3]{7.1.4.1.3}.
$$\etiquette[7.1.4.2.1]{XVII.eq7.1.4.2.1}
\vcenter{
\xymatrix{
[T_{12}|X] \ar[rd]_{g''_2} &   &[T\sisi{'}{''}_2|Y] \ar[rd]_{g'_2\sisi{}{p}} \ar[ll]_{j'_1}&
[Z\sisi{'}{''}_{2}|Z] \ar[r]^{r}\ar[l]_{j\sisi{'}{''}_2} \ar[d]_{l}&
C^{\Delta}_Z \ar[d]^{g}\\
            &    [T_1|X] \ar[rrrd]_{u_{f_1}} &  & [Y_1|Y]\ar[r]_{u_{p_1}}\ar[ll]_{j_1} & C_Y^{\Delta} \ar[d]^{f}\\
            &       &  &   &  C_X^{\Delta}
}}\leqno(7.1.4.2.1)
$$

We shall\pageoriginale define a natural arrow:

% END INCLUDED SGA4_English_current_buildable_source_20260906/source/17/drafts/SGA4_XVII_LINES_8601_9000_English.texfrag

% BEGIN INCLUDED SGA4_English_current_buildable_source_20260906/source/17/drafts/SGA4_XVII_LINES_9001_9400_English.texfrag
\[
\beta:R^{+(1)}f_{!}\circ R^{+(2)}g_{!}\longrightarrow R^{+(12)}\sisi{fg}{(fg)}_{!}
\]
which we shall show to be an isomorphism.

For this we shall use two lemmas:

\begin{enonce*}{Lemma 7.1.4.2.2}\etiquette[7.1.4.2.2]{XVII.lem7.1.4.2.2}
Let there be given a commutative diagram:
\[
\xymatrix{
X_1    \ar[d]_{p_1}     \ar[r]^{j_1} & T_1 \ar[d]^{h'}\\
X    \ar[d]_{f}       \ar[r]_{j}           & T \ar[ld]^{h}\\
    S       &
}
\]
where $h$ and $h'$ are proper, $p_{1}$ is proper and surjective, and
$j$ and $j_{1}$ are immersions. Then the canonical arrow
\[
R^{+}h\sisi{}{_*}\circ j_{!}\longrightarrow R^{+(1)}\sisi{(hh')}{f}_{!}
\]
is an isomorphism.
\end{enonce*}

Indeed, by (\hyperref[XVII.constru7.0.5]{7.0.5}), the arrow
\[
j_{!}\longrightarrow R^{+}\bar{u}_{h'\sisi{}{*}}\circ j_{1!}\circ
L^{+}\bar{u}_{p_{1}}\sisi{}{^*}
\]
is an isomorphism. The result follows by applying $R^{+}h\sisi{}{_*}$.

\begin{enonce*}{Lemma 7.1.4.2.3}\etiquette[7.1.4.2.3]{XVII.lem7.1.4.2.3}
With the notation of diagram \emph{\eqhyperref[XVII.eq7.1.4.2.1]{7.1.4.2.1}}, the canonical arrows
\begin{gather*}
R^{+(1)}f_{!}\circ R^{+}\bar{u}_{p_{1}\sisi{}{*}}\longrightarrow
R^{+}\bar{u}_{f_{1}\sisi{}{*}}\circ j_{1!}\\
R^{+(2)}g_{!}\circ R^{+}\bar{r}\sisi{}{_*}\longrightarrow
R^{+}\bar{u}_{p_{1}\sisi{}{*}}\circ R^{+}\sisi{g'_{2}}{(g_2'p)_*}\circ j\sisi{'}{''}_{2!}
\end{gather*}
are isomorphisms.
\end{enonce*}

We prove only the first isomorphism, leaving it to the reader to
deduce\pageoriginale the second.

With the evident notation\sisi{, put $C=\fprod{[Y_{1}|Y]}{X_{1}}{X}$
and $\bar{C}=\fprod{[T_{1}|S]}{T_{1}}{S}$, so that there is a
commutative diagram:
\[
\xymatrix{
                                & \bar{C} \ar[rdd]|\hole \ar[rr]  & & T_1 \ar[ldd]^{f_1}\\
C                 \ar[rr]  \ar[ru]^{i}  \ar[d]^{q}    &    &  Y_1  \ar[ru]_{j_1} \ar[d]^{p_1}&      \\
[Y_1|Y] \ar[rr]_{u_{p_1}}&  &      Y         &
}
\]
where $i$ is an immersion and $q$ is proper and surjective. One then
constructs a diagram
\[
\xymatrix{
\bigl[\bigl[C \underset{q}{\bigm|}[Y_1|Y]\bigr]\bigr] \ar[r]^{i} \ar[d]^{u_q}&
[[\bar{C}|C^{\Delta}_Y]] \ar[ldd]^{k}\\
|X_1 | X |                             \ar[d]         &\\
Y                       &                        &
}
\]
and it follows from (\hyperref[XVII.lem7.1.4.2.2]{7.1.4.2.2}) that the functor $R\bar{u}j_{1}\circ
j_{1}$ is identified with the functor $R^{+}\bar{k}\circ i_{!}\circ
L^{+}\bar{u}_{q}$ and, by (\hyperref[XVII.lem7.0.3]{7.0.3}) and the base-change theorem, the latter functor is identified with $R^{+(1)}f_{!}\circ
R^{+}\bar{u}_{p_{1}}$. }{ and the notation of \hyperref[XVII.constru7.1.3]{7.1.3} a), there is a commutative diagram
\[\xymatrix{&[Y_1|Y]\ar[rr]^{j_1}\ar[dd]^(.7){\bar{u}_{p_1}}|\hole &&[T_1|X]\ar[dd]^{\bar{u}_{f_1}}\\
[[Y_1,Y_1|Y]]\ar[ur]^{k}\ar[dd]_{q}\ar[rr]^(.7)i && [[T_1,T_1|X]]\ar[dd]\ar[ur]\\
&Y\ar[rr]^(.3){f}|\hole && X\\
[Y_1|Y]\ar[ur]^{\bar{u}_{p_1}}\ar[rr]^{j_1} &&[T_1|X]\ar[ur]_{\bar{u}_{f_1}} }\]
where $q$ and $k$ are the two projections from $[[Y_1,Y_1|Y]]$ to $[Y_1|Y]$. By cohomological descent, $R^+\bar{u}_{f_1 *} \circ j_{1!}$ is identified with $R^+\bar{u}_{f_1 *} \circ j_{1!} \circ R^+ q_* L^+ q^*$, which, by \hyperref[XVII.lem7.0.3]{7.0.3} applied to the front and top faces of the cube above, is identified with $R^+\bar{u}_{f_1*}\circ j_{1!}\circ R^+k_* \circ L^+q^*$. By the base-change theorem, this last functor is identified with $R^+\bar{u}_{f_1 *} \circ j_{1!} \circ L^+\bar{u}_{p_1*} \circ R^+\bar{u}_{p_1 *}$, which by definition is $R^{+(1)}f_{!}\circ
R^{+}\bar{u}_{p_{1}*}$.\nde{the proof in the original text is slightly different %and not very clear
and uses \hyperref[XVII.lem7.1.4.2.2]{7.1.4.2.2}.}

}%
Lemma (\hyperref[XVII.lem7.1.4.2.3]{7.1.4.2.3}) having been proved, there is a
sequence of isomorphisms
\begin{align*}
R^{+(\sisi{3}{12})} \sisi{fg}{(fg)}_{!} &= R^{+}\bar{u}_{f_{1}\sisi{}{*}}\circ R^{+}g''_{2}\circ
j'_{1!}\circ j\sisi{'}{''}_{2!}\circ L^{+}\bar{r}^{\sisi{+}{*}}\\
&\isomfrom  R^{+}\bar{u}_{f_{1}\sisi{+}{*}}\circ j_{1!}\circ R^{+}
\sisi{g'_{2}}{(g'_2p)_*}\circ j'_{2!}\circ L^{+}\bar{r}^{\sisi{+}{*}}\\
&\isomfrom  R^{+(1)}f_{!}\circ R^{+}\bar{u}_{p_{1}\sisi{+}{*}}\circ
R^{+}\sisi{g'_{2+}}{(g'_2p)_*}\circ j'_{2!}\circ L^{+}\bar{r}\sisi{}{^*}\\
&\isomfrom  R^{+(1)}f_{!}\circ R^{+(2)}g_{!}\circ
R^{+}\bar{r}\sisi{}{_*}\circ L^{+}\bar{r}^{\sisi{+}{*}}\\
&\isomfrom  R^{+(1)}f_{!}\circ R^{+(2)}g_{!}
\end{align*}
whose\pageoriginale composite gives $\beta$.

\paragraph{}\etiquette[7.1.4.3]{XVII.par7.1.4.3}
If $(p_{3},j_{3},(fg)_{3})$ is a pseudo-compactification of $fg$, an
isomorphism
\[
R^{+(1)}f_{!}\circ
R^{+(2)}g_{!}\xrightarrow{\gamma_{123}}R^{+(3)}fg_{!}
\]
is defined by composing $\alpha_{(12)3}$ with $\beta$.

\paragraph{}\etiquette[7.1.4.4]{XVII.par7.1.4.4}
Now let $(p_{1'},j_{1'},f_{1'})$ and $(p_{2'},j_{2'},g_{2'})$ be two
other pseudo-compactifications of $f$ and $g$, together with a diagram
\sisi{(\hyperref[XVII.par7.1.4.3]{7.1.4.3})}{\eqhyperref[XVII.eq7.1.4.1.3]{7.1.4.1.3}}; one checks that the following diagram is commutative:
\[
\xymatrix{
R^{+(12)}\sisi{fg}{(fg)}_1 \ar[d]_{\alpha_{(12)(1'2')}} & R^{+(1)}f_1 \circ
R^{+(2)}g_! \ar[d]^{\alpha_{11'}\ast \alpha_{22'}} \ar[l]_(.6){\beta}\\
R^{+(1'2')}\sisi{fg}{(fg)}_!  & R^{+(1')}f_! \circ R^{+(2')}g_! \ar[l]^(.6){\beta'}
}
\]
by considering a diagram of the type \sisi{(\hyperref[XVII.par7.1.4.4]{7.1.4.4})}{\eqhyperref[XVII.eq7.1.4.2.1]{7.1.4.2.1}} in which the
simplicial objects are replaced by bisimplicial objects.

It follows that the following diagram is commutative:
\[
\xymatrix{
R^{+(3)}\sisi{fg}{(fg)}_! \ar[d]_{\alpha_{33'}}& R^{+(1)}f_! \circ R^{+(2)}g_!
\ar[l]_(.6){\gamma_{123}} \ar[d]^{\alpha_{11'}\ast
  \alpha_{22'}}\\
R^{+(3')}\sisi{fg}{(fg)}_!  & R^{+(1')}f_! \circ R^{+(2')}g_! \ar[l]_(.6){\gamma_{1'2'3'}}
}
\]
for every pseudo-compactification $(p_{3'},j_{3'},(fg)_{3'})$ of
$fg$, which completes the proof of \hyperref[XVII.constru7.1.4]{7.1.4}.

\subsubsection{}\etiquette[7.1.5]{XVII.subsubsec7.1.5}
By choosing a pseudo-compactification for every morphism $f$ in
$\Sch_{!R}$, in such a way that the trivial pseudo-compactification is
associated with an identity morphism, one defines, for every morphism
$f:X\rightarrow Y$, a triangulated functor
\[
R^{+}f_{!}:D^{+}(X,\mathcal{A}_{X})\longrightarrow
D^{+}(Y,\mathcal{A}_{Y})
\]
and,\pageoriginale for every pair $(f,g)$ of morphisms, an isomorphism
\[
R^{+}fg_{!}\xrightarrow{\substack{c_{fg}\\\sim}}
R^{+}f_{!}\circ R^{+}g_{!}
\]
Moreover, if $f$ is an identity, $R^{+}f_{!}$ is the identity functor
and $c_{fg}$ is the identity.

\begin{enonce*}{Proposition 7.1.6}\etiquette[7.1.6]{XVII.prop7.1.6}
The correspondence $f\rightarrow R^{+}f_{!}$ defines a normalized
pseudofunctor from $\Sch_{!R}$ to the $2$-category of categories. In
other words, the categories $D^{+}(X,\mathcal{A}_{X})$, for $X$ an
object of $\Sch_{!R}$, can be organized into a cofibered category over
$\Sch_{!R}$ by setting, for every morphism $f:X\rightarrow Y$ in
$\Sch_{!R}$,
\[
\Hom_{f}(K,L)=\Hom(R^{+}f_{!}K,L)
\]
with $K\in D^{+}(X,\mathcal{A}_{X})$ and $L\in D^{+}(Y,\mathcal{A}_{Y})$.
\end{enonce*}

It remains to see that the identification
$R^{+}fg_{!}\xrightarrow{\substack{c_{fg}\\\sim}}
R^{+}f_{!}\circ R^{+}g_{!}$ is compatible with triple composites. Let
$h:T\rightarrow Z$, $g:Z\rightarrow Y$, and $f:Y\rightarrow X$ be
three morphisms in $\Sch_{!R}$; by what precedes, to verify this
compatibility it suffices to verify it after an arbitrary choice of
pseudo-compactifications for $h$, $g$, and $f$.

There is a diagram:
$$\vcenter{\etiquette[7.1.6.1]{XVII.eq7.1.6.1}
  \xymatrix{
 T_5 \ar[rd]& T_4 \ar[rd] \ar@{_{(}->}[l]& T_3 \ar[rd]\ar@{_{(}->}[l]
 & T'\ar@{_{(}->}[l] \ar[r]^{p_3} \ar[d]& T \ar[d]^{h} \\
        & T_3 \ar[rd]  & T_2 \ar@{_{(}->}[l] \ar[rd]  & Z'\ar[r]^{p_2}\ar@{_{(}->}[l] \ar[d]& Z
  \ar[d]^{g}\\
        &    & T_1 \ar[rrd]& Y_1\ar@{_{(}->}[l] \ar[r]^{p_1}& Y \ar[d]^{f}\\
        &    &     &     & X
}}\leqno{(7.1.6.1)}
$$
where all the oblique arrows are proper and $p_{1}$, $p_{2}$, and
$p_{3}$ are proper and surjective.

A diagram of the type (\hyperref[XVII.par7.1.4.4]{7.1.4.4}) constructed from
\eqhyperref[XVII.eq7.1.6.1]{7.1.6.1} then gives the conclusion.

\begin{enonce*}{Proposition 7.1.7}\etiquette[7.1.7]{XVII.prop7.1.7}
If\pageoriginale $f:X\rightarrow S$ is a compactifiable morphism, the
functor $Rf_{!}$ coincides with the one defined in Exposé \emph{XVII}.
\end{enonce*}

Indeed, there is a commutative diagram:
\[
\xymatrix{
X_1         \ar[r]^{j_1} \ar[d]_{p_1}& T_1 \ar[d]^{h'}\\
X           \ar[r]_{j}  \ar[d]_{f} & T \ar[ld]^{h}\\
S
}
\]
to which one applies (\hyperref[XVII.lem7.1.4.2.2]{7.1.4.2.2}).

\begin{enonce*}{Definition 7.1.8}\etiquette[7.1.8]{XVII.def7.1.8}
If $f:X\rightarrow Y$ is a morphism in $\Sch_{!}$, the functor
$R^{+}f_{!}$ just defined is called the \emph{direct-image functor with
proper supports} (in the sense of derived categories). The functor
$H^{q}\circ R^{+}f_{!}$, denoted $R^{q}f_{!}$, is called the
$q$\textsuperscript{th} \emph{direct-image functor with proper supports.}
\end{enonce*}

With this notation, we leave it to the reader to verify the

\begin{enonce*}{Construction 7.1.9}\etiquette[7.1.9]{XVII.constru7.1.9}
Let $(\pi_{1}, j_{1}, f_{1})$ be a pseudo-compactification of
$f:X\rightarrow Y$, and denote by $\pi^{n}_{1}$ the canonical morphism
$[X_{1}|X]_{n}\rightarrow X$. Then, for every complex $K^{\bigdot}$ in
$D^{+}(X,\mathcal{A}_{X})$, there are two biregular spectral sequences:
\begin{align*}
E^{pq}_{1} &= R^{q}f_{!}(K^{p})\Longrightarrow R^{p+q}
f_{!}(K^{\bigdot})\tag{7.1.9.1}\\
E^{pq}_{1} &= R^{q}(f\circ
\pi^{p}_{1})_{!}(\pi^{p\sisi{+}{*}}_{1}(K^{\bigdot}))\Longrightarrow
R^{p+q}f_{!}(K^{\bigdot})\tag{7.1.9.2}
\end{align*}\etiquette[7.1.9.1]{XVII.eq7.1.9.1}
\etiquette[7.1.9.2]{XVII.eq7.1.9.2}%The vertical space following this paragraph is too large.
\end{enonce*}

(Take, for example, the canonical flabby resolution of $j_{1!}\circ
L^{+}\bar{u}_{p_{1}}^{\sisi{+}{*}}(K^{\bigdot})$ (\cf XVII (\hyperref[XVII.subsec4.2]{4.2})), and then its
direct image under $u_{f_{1}}$.)

\subsection[Some properties of $R^{+}f_{!}$]{Some properties of
  \mathversion{bold}$R^{+}f_{!}$}\etiquette[7.2]{XVII.subsec7.2}
\pageoriginale

\begin{enonce*}{Proposition 7.2.1}\etiquette[7.2.1]{XVII.prop7.2.1}
Let $f:X\rightarrow Y$ be a morphism in $\Sch_{!R}$ of relative
dimension $\leq d$. Then the functor $R^{+}f_{!}$ has cohomological
dimension $\leq 2d$.
\end{enonce*}

It is enough to show that, for every sheaf $F$ of
$\mathcal{A}_{X}$-modules on $X$, one has $R^{i}f_{!}(F)=0$ for
$i>2d$. (It then follows by a standard argument that, for every complex
$K^{\bigdot}$ in $D^{+}(X,\mathcal{A}_{X})$ such that
$H^{i}(K^{\bigdot})=0$ for $i>k$, one has $R^{i}f_{!}(K^{\bigdot})=0$
for $i>k+2d$.)

Let $\mathcal{U}=(U_{i})_{i\in I}$ be a finite covering of $X$ by
affine open subsets. One defines (\hyperref[XVII.prop6.2.9]{6.2.9}) a left resolution
$K^{\bigdot}_{\mathcal{U}}\rightarrow F$ of $F$ by putting
\[
K^{n}_{\mathcal{U}}=\bigoplus_{|P|=\sisi{}{-}n+1} j_{P!}j^{*}_{P}(F)
\]
where, for every subset $P$ of $I$, $j_{P}$ is the inclusion of
$\cup_{i\in P} U_{i}$ into $X$.

From \eqhyperref[XVII.eq7.1.9.1]{7.1.9.1} one obtains a spectral sequence
\[
E^{pq}_{1}=\bigoplus_{\substack{|P|=\sisi{}{-}p+1\\ P\sisi{\in}{\subset} I\sisi{}{\text{ nonempty}}}} R^{p}(fj_{P})_{!}
j^{\ast}_{P}(F)\Longrightarrow R^{p+q}f_{!}(F)
\]
and the fact that $fj_{P}$ is compactifiable, a case covered by Exposé
XVII, shows that $E^{pq}_{1}=0$ for $p>0$ and $q>2d$; it follows that
$R^{n}f_{!}(F)=0$ for $n>2d$, as required.

\begin{enonce*}{Proposition 7.2.2}\etiquette[7.2.2]{XVII.prop7.2.2}
Let $f:X\rightarrow Y$ be a morphism in $\Sch_{!R}$ and let
$K^{\bigdot}$ be an object of $D^{b}(X,\mathcal{A}_{X})$ of finite
\tTor-dimension. Then $R^{+}f_{!}(K^{\bigdot})$ has finite
\tTor-dimension.
\end{enonce*}

As above, let $\mathcal{U}=(U_{i})_{i\in I}$ be a finite covering of
$X$ by affine open subsets. The functoriality of
$K^{n}_{\mathcal{U}}$ makes it possible to construct a bounded left
resolution of $K^{\bigdot}$ by complexes of finite \tTor-dimension,
since the latter property can be checked pointwise (\hyperref[XVII.def4.1.9]{4.1.9}). We
denote the double complex so obtained by
$K^{\bigdot\bigdot}_{\mathcal{U}}$, the \emph{second} index being the
one corresponding to $K^{\bigdot}$. With the conventions of
(\sisi{VI b.(\refnoexpose{Vbis.subsubsec2.3.2})}{\hyperrefext[Vbis][Vbis.subsubsec2.3.2]{2.3.2}}), for every integer $n$ there is an exact sequence, where the operation
$\sigma_{\leq n}$ is taken in
$C^{+}(C^{+}(\Mod(X,\mathcal{A}_{X}))$:
\[
\sisi{0\longrightarrow (\sigma_{\leq
n}(K^{\bigdot\bigdot}_{\mathcal{U}}))_{s}\longrightarrow (\sigma_{\leq
n+1}(K^{\bigdot\bigdot}_{\mathcal{U}}))_{s}\longrightarrow (-1)^{n+1}
K^{n+1,\bigdot}_{\mathcal{U}}\longrightarrow 0}{0\to (-1)^{n+1}K^{n+1,\bigdot}_{\mathcal{U}}\to(\sigma_{\leq
n+1}(K^{\bigdot\bigdot}_{\mathcal{U}}))_{s}\to(\sigma_{\leq
n}(K^{\bigdot\bigdot}_{\mathcal{U}}))_{s}\to 0}
\]
hence\pageoriginale a distinguished triangle in $D^{+}(X,\mathcal{A}_{X})$:
\[
\xymatrix@C=1em{
  & K^{n+1,\bigdot}_{\Uu} \sisi{\ar[ld]}{\ar[rd]}&    \\
  (\sigma_{\leq n}(K^{\bigdot\bigdot}_{\Uu}))_s \sisi{\ar[rr]}{\ar[ur]}&  &(\sigma_{\leq
  n+1}(K^{\bigdot\bigdot}_{\Uu}))_s \sisi{\ar[lu]}{\ar[ll]}\\
}
\]

Using the fact that $R^{+}f_{!}(K^{n+1,\bigdot}_{\mathcal{U}})$ has
finite \tTor-dimension for every $n$ (\hyperref[XVII.thm5.2.10]{5.2.10}), it follows that
$R^{+}f_{!}((\sigma_{\leq
0}(K^{\bigdot\bigdot}_{\mathcal{U}}))_{s})=R^{+}f_{!}(K^{\bigdot})$
has finite \tTor-dimension.

\begin{enonce*}[remark]{Remark 7.2.3}\etiquette[7.2.3]{XVII.rem7.2.3}
The preceding argument does not show that
$R^{+}f_{!}(K^{\bigdot})$ has \tTor-dimension $\leq k$ whenever
$K^{\bigdot}$ has \tTor-dimension $\leq k$ (\cf \sisi{(3.7)}{\hyperref[XVII.subsubsec7.3.7]{7.3.7}}).
\end{enonce*}

\begin{enonce*}{Proposition 7.2.4}\etiquette[7.2.4]{XVII.prop7.2.4}
Let there be given a Cartesian diagram in $\Sch_{!R}$:
\[
\xymatrix{
X'   \ar[r]^{g'}\ar[d]_{f'}& X \ar[d]^{f}\\
S'  \ar[r]_{g} & S
}
\]
Then, for every complex $K^{\bigdot}$ in $D^{+}(X,\mathcal{A}_{X})$,
the canonical morphism
\[
g^{\ast}\circ
R^{+}f_{!}(K^{\bigdot})\xrightarrow{\mathrm{ch}_{K^{\bigdot}}}R^{+}f_{!}\sisi{}{'}\circ
{g'}^{\ast}(K^{\bigdot})
\]
is an isomorphism.
\end{enonce*}

Let $(\pi_{1}, j_{1}, f_{1})$ be a pseudo-compactification of $f$; by
base change it gives a pseudo-compactification
$(\pi'_{1}, j'_{1}, f'_{1})$ of $f'$ and a commutative diagram:
\[
\xymatrix@R=2ex{
                 && T'_1 \ar'[ldddd]^(.8){f'_1}|(.55)\hole [lldddddddd]|(.25)\hole \ar[rrrr] &&&& T_1 \ar[lldddddddd]^{f_1}\\
                 &\\
X'_1 \ar[rruu]^{j'_1}\ar[ddd]_{\pi'_1} \ar[rrrr] &&&& X_1 \ar[rruu]_{j_1} \ar[ddd]^{\pi_1}&\\
                     & & &\\
                     & *{} & &\\
X'   \ar[rrrr]_{g'} \ar[ddd]_{f'} &&&& X  \ar[ddd]_{f}  &\\
                     & & &\\
                     & & &\\
S'   \ar[rrrr]_{g} &&&&  S         &
}
\]
which\pageoriginale makes it possible to define the arrow
$\mathrm{ch}_{K}$.

By \eqhyperref[XVII.eq7.1.9.2]{7.1.9.2}, there is then a morphism of
biregular spectral sequences:
\[
\xymatrix{
g^{\ast}\circ R^{q} (f \circ \pi^{p}_1)_! (\pi^{p\sisi{+}{\ast}}_1
(K^{\bigdot})) \ar[d]^{u^{pq}_1} \ar[r]&  g^{\ast} \circ R^{p+q}f_{!}(K^{\bigdot}) \ar[d]^{\mathrm{ch}_K}\\
R^{q}(f'\circ \pi'^{p}_1)_! (\pi'^{p\ast}_1 (g'^{\ast}(K^{\bigdot})))  \ar[r] &
  R^{+}f'_! \circ g'^{\ast}(K\sisi{}{^{\bigdot}})
}
\]
and $\mathrm{ch}_{K^{\bigdot}}$ is an isomorphism because the same is
true of the $u^{pq}_{1}$ by (\hyperref[XVII.thm5.2.6]{5.2.6}).

\subsection{Definition of $Rf_{!}$}\etiquette[7.3]{XVII.subsec7.3}

We now propose to extend the functor $R^{+}f_{!}$ ``in a reasonable
way'' to a functor
$Rf_{!}:D(X,\mathcal{A}_{X})\rightarrow D(Y,\mathcal{A}_{Y})$, for
every morphism $f$ in $\Sch_{!R}$.

\subsubsection{}\etiquette[7.3.1]{XVII.subsubsec7.3.1}
Let $(p_{1}, j_{1}, f_{1})$ be a pseudo-compactification of $f$, and
let $F$ be an $\mathcal{A}_{X}$-module. Taking the canonical flabby
resolution of $j_{1!}\circ L\bar{u}_{p_{1}}^{\ast}(F)$ (\cf
\hyperref[XVII.subsec4.2]{4.2}), then its direct image under $f_{1}$,
and then the simple complex associated with the double complex so
obtained gives a complex of $\mathcal{A}_{Y}$-modules:
$$\etiquette[7.3.1.1]{XVII.eq7.3.1.1}
0\longrightarrow
\Xx^{0}(F)\xrightarrow{S^{1}(F)}\Xx^{1}(F)\ldots
\Xx^{n}(F)\xrightarrow{S^{n}(F)}
\Xx^{n+1}(F)\ldots.\leqno(7.3.1.1)
$$
functorial\pageoriginale in $F$ and isomorphic in
$D^{+}(Y,\mathcal{A}_{Y})$ to $R^{+}f_{!}(F)$.

By \hyperref[XVII.prop7.2.1]{7.2.1}, $H^{i}(\Xx^{\bigdot}(F))=0$ for $i>2d$, where $d$
is an integer such that the relative dimension of $f$ is $\leq d$;
hence the complex
$$\etiquette[7.3.1.2]{XVII.eq7.3.1.2}
0\to \Xx^{0}(F)\to
\Xx^{1}(F)\to \ldots \sisi{}{\to}
\Xx^{2d-1}(F)\to \Ker S^{2d}(F)\to
0\leqno(7.3.1.2)
$$
is canonically isomorphic in $D^{+}(Y,\mathcal{A}_{Y})$ to the
complex \eqhyperref[XVII.eq7.3.1.1]{7.3.1.1}. Moreover:

\begin{enonce*}{Lemma 7.3.2}\etiquette[7.3.2]{XVII.lem7.3.2}
The functors $\Xx^{0}$,
$\Xx^{1},\ldots,\Xx^{2d-1}$, and $\Ker S^{2d}$ are exact
in $F$.
\end{enonce*}

% END INCLUDED SGA4_English_current_buildable_source_20260906/source/17/drafts/SGA4_XVII_LINES_9001_9400_English.texfrag

% BEGIN INCLUDED SGA4_English_current_buildable_source_20260906/source/17/drafts/SGA4_XVII_LINES_9401_9548_English.texfrag

The exactness of the functors $\Xx^{n}$ follows from the exactness
of the canonical flabby resolution and from the fact that the direct
image under a morphism of a short exact sequence of flabby sheaves is
again a short exact sequence.

To prove the exactness of $\Ker S^{2d}$, for a short exact sequence
$0\rightarrow F'\rightarrow F\rightarrow F''\rightarrow 0$ of
sheaves on $X$, consider the diagram
\[
\xymatrix{
\Ker S^{2d}(F')  \ar[r] \ar[d]& \Xx^{2d+1}(F') \ar[r] \ar[d]&
\Xx^{2d+2}(F') \ar[d] \ldots\\
\ker S^{2d}(F)   \ar[r] \ar[d]& \Xx^{2d+1}{(F)} \ar[r] \ar[d]&
\Xx^{2d+2}(F)\ar[d] \ldots \\
\ker S^{2d}(F'')   \ar[r] & \Xx^{2d+1}{(F'')} \ar[r] &
\Xx^{2d+2}{(F'')} \ldots \\
}
\]
and conclude by using the homology exact sequence associated with the
exact sequence
\[
0\longrightarrow \sigma_{\geq 2d+1}
(\Xx^{\bigdot}(F'))\longrightarrow \sigma_{\geq 2d+1}
(\Xx^{\bigdot}(F))\longrightarrow \sigma_{\geq
2d+1}(\Xx^{\bigdot}(F''))\to0.
\]

We now have the following general lemma:

\begin{enonce*}{Lemma 7.3.3}\etiquette[7.3.3]{XVII.lem7.3.3}
Let $A$ and $B$ be two abelian categories, and let
$\mathrm{Je}^{\bigdot}$ be an object of $C^{b}(\Ffex(A,B))$, where
$\Ffex(A,B)$ denotes the category of exact functors from $A$ to $B$.
Associating with every complex $K^{\bigdot}$ in $C(A)$ the simple
complex associated with the double complex
$\mathrm{Je}^{\bigdot}(K^{\bigdot})$ gives a triangulated functor
$K(A)\rightarrow K(B)$ that preserves quasi-isomorphisms.
\end{enonce*}

The functor in question may be written\pageoriginale as a composite
\[
C(A)\longrightarrow C^{b}(C(B))\xrightarrow{(\quad)_{s}}C(B),
\]
where $(\quad)_{s}$ is the associated-simple-complex functor, and one
uses (\sisi{VI b (\refnoexpose{Vbis.lem2.3.2.2})}{\hyperrefext[Vbis][Vbis.lem2.3.2.2]{2.3.2.2}}. Remark); the details are left to the
reader.

\begin{enonce*}[remark]{Exercise 7.3.4}\etiquette[7.3.4]{XVII.exer7.3.4}
Show that, for every complex $K^{\bigdot}$ in
$D^{+}(X,\mathcal{A}_{X})$, the canonical arrow
\[
\tau_{\leq 2d} \Xx^{\bigdot}(K^{\bigdot})\longrightarrow
\Xx^{\bigdot}(K^{\bigdot})
\]
induces a quasi-isomorphism on the associated simple complexes.
\end{enonce*}

For this, observe first that the canonical arrow
\[
\varinjlim_{k\geq 2d} (\tau_{\leq k}
\Xx^{\bigdot}(K^{\bigdot}))_{s}\longrightarrow
(\Xx^{\bigdot}(K^{\bigdot}))_{s}
\]
is an isomorphism, and then that the cohomology sheaves of a complex
commute with filtered direct limits.

\setcounter{subsubsection}{4}
\subsubsection{}\etiquette[7.3.5]{XVII.subsubsec7.3.5}
Let $R^{(1)}f_{!}:D(X,\mathcal{A}_{X})\rightarrow
D(Y,\mathcal{A}_{Y})$ be the functor defined by
\eqhyperref[XVII.eq7.3.1.2]{7.3.1.2} using
\sisi{\hyperref[XVII.rem7.2.3]{7.2.3}}{\hyperref[XVII.lem7.3.3]{7.3.3}}.
In view of
\sisi{\hyperref[XVII.prop7.2.4]{7.2.4}}{\hyperref[XVII.exer7.3.4]{7.3.4}},
it can be computed directly from
\eqhyperref[XVII.eq7.3.1.1]{7.3.1.1}, and one can check that it ``does
not depend'' on the chosen pseudo-compactification. One would then
have to verify statements \hyperref[XVII.constru7.1.3]{7.1.3},
\hyperref[XVII.constru7.1.4]{7.1.4}, and
\hyperref[XVII.prop7.1.6]{7.1.6} after deleting the superscripts $+$,
which the editor did not have the courage to do.

\begin{enonce*}{Definition 7.3.6}\etiquette[7.3.6]{XVII.def7.3.6}
The functor $Rf_{!}:D(X,\mathcal{A}_{X})\rightarrow
D(Y,\mathcal{A}_{Y})$ constructed in
\hyperref[XVII.subsubsec7.3.5]{7.3.5} is again called the
\emph{direct-image functor with proper supports.}
\end{enonce*}

\setcounter{subsubsection}{6}
\subsubsection{}\etiquette[7.3.7]{XVII.subsubsec7.3.7}
One can then show that the functor $Rf_{!}$ commutes with base change
and that statements (\hyperref[XVII.prop5.2.9]{5.2.9}) and
(\hyperref[XVII.thm5.2.10]{5.2.10}) remain valid for every morphism in
$\Sch_{!R}$, which justifies the constructions just made.

\begin{thebibliography}{\sisi{TF}{99}}
\sisi{\bibitem[1]{sga4.exp.XVII.bib.DP}}{\bibitem{sga4.exp.XVII.bib.DP}} A. Dold\pageoriginale and D. Puppe. Homologie nicht\sisi{ }{-}additiver
Funktoren. Anwendungen. Ann. Inst. Fourier \sisi{XI 1961}{\textbf{11} (1961), 201--312}.

\sisi{\bibitem[1]{sga4.exp.XVII.bib.GZ}}{\bibitem{sga4.exp.XVII.bib.GZ}} Gabriel-Zisman. Calculus of Fractions and Homotopy
Theory, Ergebnisse, Bd 35, Springer 1967.

\sisi{\bibitem[1]{sga4.exp.XVII.bib.Gi}}{\bibitem{sga4.exp.XVII.bib.Gi}} J. Giraud. Méthode de la descente, Mémoire nᵒ 2 de la Société Mathématique de France (1964).

\sisi{\bibitem[TF]{sga4.exp.XVII.bib.TF}}{\bibitem{sga4.exp.XVII.bib.TF}} R. Godement. Théorie des faisceaux, Hermann,
1958. Act. Scient. Ind. nᵒ 1252, (Paris)\sisi{ (cited [TF])}{}.

\sisi{\bibitem[1]{sga4.exp.XVII.bib.Gr}}{\bibitem{sga4.exp.XVII.bib.Gr}} A. Grothendieck. TDTE IV: les schémas de
Hilbert. Séminaire Bourbaki 221. May 1961. Reprinted in FGA.

\sisi{\bibitem[RD]{sga4.exp.XVII.bib.RD}}{\bibitem{sga4.exp.XVII.bib.RD}} R. Hartshorne. Residues and Duality. Lecture Notes
in Mathematics nᵒ 20. Springer 1966.

\sisi{\bibitem[1]{sga4.exp.XVII.bib.La}}{\bibitem{sga4.exp.XVII.bib.La}} D. Lazard. Thesis --- Paris 1968.

\sisi{}{\bibitem{sga4.exp.XVII.bib.Mu} D. Mumford, \emph{Geometric Invariant Theory}, Ergebnisse der Mathematik und ihrer Grenzgebiete 34, Springer-Verlag, 1965.}

\sisi{}{\bibitem{sga4.exp.XVII.bib.Na} M. Nagata, ``A generalization of the imbedding problem of an abstract variety in a complete variety'', \emph{J. Math. Kyoto Univ.} \textbf{3} (1963), 89--102}

\sisi{\bibitem[1]{sga4.exp.XVII.bib.Ro}}{\bibitem{sga4.exp.XVII.bib.Ro}} N. Roby. Lois polynômes et lois formelles en théorie
des modules. Ann. E.N.S. \textbf{80} \sisi{1963, p.}{(1963), }231--\sisi{2}{3}48.

\sisi{\bibitem[CD]{sga4.exp.XVII.bib.CD}}{\bibitem{sga4.exp.XVII.bib.CD}} J.L. Verdier. Catégories dérivées (état 0),
mimeographed by the I.H.E.S.

\sisi{\bibitem[2]{sga4.exp.XVII.bib.V2}}{\bibitem{sga4.exp.XVII.bib.V2}} J.L. Verdier. Catégories dérivées (to appear in
Grundlehren der Math. Wiss, Springer).

\sisi{\bibitem[3]{sga4.exp.XVII.bib.V3}}{\bibitem{sga4.exp.XVII.bib.V3}} J.L. Verdier. Dualité dans la cohomologie des
espaces localement compacts. Séminaire Bourbaki 300. Nov. 1965.
\end{thebibliography}


%0.2 no ref. 1957, 1329, 5575, 5594, 5821
%2.1.2.2 no ref. 1386
%2.1.2.3 no ref  1412
%2.1.12.1 no ref 1641
%3.2.4.3 no ref  2680
%5.1.8.2 no ref 3925
%XII 4.1 no ref 4690, 6830
%6.3.4.3 no ref 7561, 7876, 8060,8409
%6.3.11 no ref 7736, 7833
%6.2.17 no ref 8042
%6.3.17.1 no ref 8078
%6.3.2.3 no ref 8271
%4.3.0 no ref 8569
%3.7 no ref 9229


\ifx\danslelivre\undefined
\let\SGAXVIIenddocument\enddocument
\else
\let\SGAXVIIenddocument\relax
\fi
\SGAXVIIenddocument

% END INCLUDED SGA4_English_current_buildable_source_20260906/source/17/drafts/SGA4_XVII_LINES_9401_9548_English.texfrag


% END INCLUDED SGA4_English_current_buildable_source_20260906/source/17/17_English.tex


\makeatletter
\def\input@path{{../18/}}
\makeatother
% BEGIN INCLUDED SGA4_English_current_buildable_source_20260906/source/18/18_English.tex
\ifx\danslelivre\undefined
% Release engine: XeLaTeX (the LuaTeX compatibility layer breaks Xy-pic output).
\documentclass[10pt]{../fichiers-configuration-latex/smfart}
% BEGIN INCLUDED SGA4_English_current_buildable_source_20260906/source/fichiers-configuration-latex/paquets.tex
\usepackage{polyglossia}
%\usepackage[francais]{babel}
\usepackage{fontspec}
\usepackage[math-style=ISO,bold-style=upright,sans-style=upright]{unicode-math}
%\usepackage[mathletters]{ucs}
%\usepackage[utf8]{inputenc}
%\usepackage[T1]{fontenc}
\usepackage[unicode,hypertexnames=false,colorlinks=true,linkcolor=blue,citecolor=Sepia,urlcolor=magenta]{hyperref}
\setmainlanguage{french}
\setmainfont[Ligatures=TeX]{Linux Libertine O}
% variante possible : TeXGyrePagella ?
\setmathfont{XITS Math}

% maths
\usepackage{amsmath}
%\usepackage{amssymb}
%\usepackage{mathrsfs}
%\usepackage{stmaryrd}
%\usepackage{euscript}

% couleurs
\usepackage[usenames,dvipsnames]{xcolor}

% graphiques
\usepackage{graphicx}
\usepackage[all]{xy}
%\usepackage{tikz}
%\usetikzlibrary{matrix,arrows,calc}
%\tikzset{column sep=3em,row sep=2em,text height=1.5ex, text depth=0.25ex,>=angle 90}
%\tikzset{description/.style={fill=white,inner sep=1pt}}

% hyperliens, 'etiquettes
\usepackage{xr-hyper}
\usepackage{../fichiers-configuration-latex/smfthm}
%\ifx\cacheretiquettes\undefined
%\usepackage[inline]{showlabels} \else \fi % pour rendre visibles les 'etiquettes
%\usepackage{../fichiers-configuration-latex/makeidx} % pour faire un index
%\usepackage[pagebackref,unicode]{hyperref}  % pour voir o`u sont cit'ees les r'ef'erences
% hyperref is loaded once above with the cumulative-reader destination policy.
%\hypersetup{breaklinks=true} %permet le retour `a la ligne dans les liens trop longs. Incompatible avec pdflatex
%\renewcommand*{\backref}[1]{{\small [cit\'e page(s)~#1.]}}

\usepackage{xspace}
\usepackage{pifont} % pour expose XVII

% END INCLUDED SGA4_English_current_buildable_source_20260906/source/fichiers-configuration-latex/paquets.tex

% BEGIN INCLUDED SGA4_English_current_buildable_source_20260906/source/fichiers-configuration-latex/numerotation.tex
\renewcommand{\thefootnote}{\roman{footnote}} % pour mettre les notes en bas de page sous la forme (ii) et non (2), qui ressemble trop (?) a un exposant.

\newcounter{pageoriginal}
\newcommand{\pageoriginale}{\refstepcounter{pageoriginal}\marginpar{\footnotesize\textbf{\thepageoriginal}}\xspace}
\let\pageoriginaled\pageoriginale

\usepackage{ifthen}
\newboolean{original}
\def\sisi{\ifthenelse{\boolean{original}}}
\setboolean{original}{false}

\newcounter{cnote}
\def\nde#1{\ifthenelse{\boolean{original}}{}{\stepcounter{cnote}\def\thefootnote{\arabic{cnote}}\protect\footnote[\value{cnote}]{N.D.E.: #1}}}
\def\ndemark{\ifthenelse{\boolean{original}}{}{\stepcounter{cnote}\def\thefootnote{\arabic{cnote}}\protect\footnotemark[\value{cnote}]}}
\def\ndetext#1{\ifthenelse{\boolean{original}}{}{\def\thefootnote{\arabic{cnote}}\protect\footnotetext[\value{cnote}]{N.D.E.: #1}}}

% END INCLUDED SGA4_English_current_buildable_source_20260906/source/fichiers-configuration-latex/numerotation.tex

% BEGIN INCLUDED SGA4_English_current_buildable_source_20260906/source/fichiers-configuration-latex/macros.tex
\setmainlanguage{french}
\setotherlanguage{english}


%% Fixed, engine-independent release identifier.  The inherited dynamic Git
%% probe used Lua io.popen, exposed build-state-dependent output, and printed
%% a visible question mark outside a Git checkout.
\def\showgitstatus{English release r2, 2026-07-23}


% macros pour LaTeX
\renewcommand{\mathsf}{}

% Police différente pour U+2113 SCRIPT SMALL L
\setmathfont[range="2113]{latinmodern-math.otf}

% Give every inherited semantic label its own stable PDF destination instead
% of inheriting whichever automatic equation/page anchor happened to be live.
% The legacy label remains as a compatibility alias used by existing edges.
\makeatletter
\def\etiquette[#1]#2{%
  \Hy@raisedlink{\hyper@anchorstart{sga4.semantic.#2}\hyper@anchorend}%
  \def\@currentHref{sga4.semantic.#2}%
  \def\@currentlabel{#1}%
  \label{#2}%
}
\makeatother
\def\eqhyperref[#1]#2{(\hyperref[#1]{#2})}
\def\pagehyperref[#1]#2{(page \hyperref[#1]{#2})}
\def\hyperrefext[#1][#2]#3{\hyperref[#2]{#1 #3}}
\def\eqhyperrefext[#1][#2]#3{\hyperref[#2]{#1 #3}}     %cf. exposes 17
\def\neweqhyperrefext[#1][#2]#3{\hyperref[#2]{#1 #3}}  %cf. exposes 17

\def\shat#1{\protect\widehat{#1}}
\def\stilde#1{\protect\widetilde{#1}}
%%%%%%%%%%%%%%%%%%


\def\etcn{etc.\xspace}
\def\etco{etc...\xspace}
\def\etc{\sisi{\etco}{\etcn}}
\def\ldotsn{\sisi{...}{\etcn}}

\def\cf{{cf.}\kern.3em}
\def\Cf{{Cf.}\kern.3em}
\def\ie{{i.e.}\kern.3em}
\def\iev{{i.e.},\xspace}
\def\loccito{loc.\kern2pt cit.\xspace}
\def\loccitn{\textit{loc.\kern2pt cit.}\xspace}
\def\loccit{\sisi{\loccito}{\loccitn}}
\def\EGA{\textup{EGA}\kern.2em}
\def\SGA{\textup{SGA}\kern.3em}

\def\quoi{}

\def\quav{\quad,\quad}
\def\quaet{\quad \text{et}\quad}
\def\fibreproduct#1#2#3{#1 ×_{#3} #2}
\let\fprod\fibreproduct
\def\ccheck#1#2{\mbox{${\check{\check{\rm #1}}}^{\lower2pt\hbox{$\scriptstyle #2$}}$}}
\def\sheaf#1{\protect\mathcal{#1}}
\def\vieuxsheaf#1{{#1}} % cf. problème avec sheaf{a\Phi_i} par exemple...
\DeclareMathOperator{\Hom}{Hom}
\DeclareMathOperator{\can}{can}
\DeclareMathOperator{\Ind}{Ind}
\DeclareMathOperator{\ind}{ind}
%\DeclareMathOperator{\Hhomtop}{\mathcal{H}\!\mathit{omtop}}
\DeclareMathOperator{\Hhomtop}{\mathbf{Homtop}}
\DeclareMathOperator{\Point}{Point}
\DeclareMathOperator{\Pro}{Pro}
\DeclareMathOperator{\Top}{Top}
\DeclareMathOperator{\Mod}{Mod}
\DeclareMathOperator{\Cart}{Cart}
%\DeclareMathOperator{\Mmorsite}{\mathcal{M}\!\mathit{orsite}}
\DeclareMathOperator{\Mmorsite}{\mathbf{Morsite}}
\DeclareMathOperator{\loc}{loc}
\DeclareMathOperator{\Loc}{Loc}
\DeclareMathOperator{\Pic}{Pic}
\DeclareMathOperator{\Et}{Et}
\DeclareMathOperator{\card}{card}
\DeclareMathOperator{\Ens}{Ens}
\DeclareMathOperator{\h}{\underline{H}}
\DeclareMathOperator{\Fonct}{Fonct}
\DeclareMathOperator{\Ker}{Ker}
\DeclareMathOperator{\but}{but}
\DeclareMathOperator{\Ob}{Ob}
\DeclareMathOperator{\ob}{ob}
\DeclareMathOperator{\Lib}{Lib}
\DeclareMathOperator{\esj}{esj}
\DeclareMathOperator{\id}{id}
\DeclareMathOperator{\Id}{Id}
\DeclareMathOperator{\ad}{ad}
\DeclareMathOperator{\Iim}{Im}
\DeclareMathOperator{\pr}{pr}
\DeclareMathOperator{\Cat}{Cat}
\DeclareMathOperator{\car}{car}
\DeclareMathOperator{\SysInd}{SysInd}
\DeclareMathOperator{\Sup}{Sup}
\DeclareMathOperator{\Coker}{Coker}
\DeclareMathOperator{\Filt}{Filt}
\DeclareMathOperator{\acc}{acc}
\DeclareMathOperator{\sur}{sur}
\DeclareMathOperator{\cat}{cat}
\DeclareMathOperator{\Ord}{Ord}
\DeclareMathOperator{\Gr}{Gr}
\DeclareMathOperator{\Cov}{Cov}
\DeclareMathOperator{\Caf}{Caf}
\DeclareMathOperator{\coker}{coker}
\DeclareMathOperator{\coim}{coim}
\DeclareMathOperator{\im}{im}
\DeclareMathOperator{\Ouv}{Ouv}
\DeclareMathOperator{\diag}{diag}
\DeclareMathOperator{\Esp}{Esp}
\DeclareMathOperator{\TOP}{TOP}
\DeclareMathOperator{\Unif}{Unif}
\DeclareMathOperator{\Iind}{\mathbf{Ind}}
\DeclareMathOperator{\Hhomcart}{\mathbf{Homcart}}
\DeclareMathOperator{\Hhomtopess}{\mathbf{Homtopess}}
\DeclareMathOperator{\Ppoints}{\mathbf{Points}}
\DeclareMathOperator{\Ppoint}{\mathbf{Point}}
\DeclareMathOperator{\Ttors}{\mathbf{Tors}}
\DeclareMathOperator{\Ssex}{\mathbf{Sex}}
\DeclareMathOperator{\Ffib}{\mathbf{Fib}}
\DeclareMathOperator{\Vvois}{\mathbf{Vois}}
\DeclareMathOperator{\Ppointess}{\mathbf{Pointess}}
\DeclareMathOperator{\scrHom}{\mathbf{Hom}}
\DeclareMathOperator{\sheafPic}{\mathbf{Pic}}
\DeclareMathOperator{\sheafAut}{\mathbf{Aut}}
\DeclareMathOperator{\Hhomtoplocan}{\mathbf{Homtoplocan}}
\DeclareMathOperator{\Hhomtopan}{\mathbf{Homtopan}}
\DeclareMathOperator{\Pptgeom}{\mathbf{Ptgéom}}
\DeclareMathOperator{\Morsite}{Morsite}
\DeclareMathOperator{\Fib}{Fib}
\DeclareMathOperator{\Fais}{Fais}
\DeclareMathOperator{\kar}{kar}
\DeclareMathOperator{\Groupes}{Groupes}
\DeclareMathOperator{\Restr}{Restr}
\DeclareMathOperator{\Isom}{Isom}
\DeclareMathOperator{\Prol}{Prol}
\DeclareMathOperator{\Res}{Res}
\DeclareMathOperator{\Tors}{Tors}
\DeclareMathOperator{\esp}{esp}
\DeclareMathOperator{\es}{es}
\DeclareMathOperator{\od}{od}
\DeclareMathOperator{\gr}{gr}
\DeclareMathOperator{\sob}{sob}
\DeclareMathOperator{\dfn}{dfn}
\DeclareMathOperator{\Sch}{Sch}
\DeclareMathOperator{\ouv}{ouv}
\DeclareMathOperator{\cons}{cons}
\DeclareMathOperator{\Nerf}{Nerf}
\DeclareMathOperator{\Points}{Points}
\DeclareMathOperator{\Espsob}{Espsob}
\DeclareMathOperator{\Homtop}{Homtop}
\DeclareMathOperator{\Pt}{Pt}
\DeclareMathOperator{\Quot}{Quot}
\DeclareMathOperator{\glob}{glob}
\DeclareMathOperator{\Homfib}{Homfib}
\DeclareMathOperator{\Qouv}{Qouv}
\DeclareMathOperator{\Quotouv}{Quotouv}
\DeclareMathOperator{\ens}{ens}
\DeclareMathOperator{\ab}{ab}
\DeclareMathOperator{\an}{an}
\DeclareMathOperator{\fs}{fs}
\DeclareMathOperator{\red}{red}
\DeclareMathOperator{\Ab}{Ab}
\DeclareMathOperator{\cosq}{cosq}
\DeclareMathOperator{\Ann}{Ann}
\DeclareMathOperator{\Zar}{Zar}


%%%%%%%%%%%%
%%%%% SGA4-2
\DeclareMathOperator{\Sex}{Sex}
\DeclareMathOperator{\vois}{vois}
\DeclareMathOperator{\Vois}{Vois}
\DeclareMathOperator{\Ext}{Ext}
\DeclareMathOperator{\ext}{ext}
\DeclareMathOperator{\Fer}{Fer}
\DeclareMathOperator{\Add}{Add}
\DeclareMathOperator{\Cocart}{Cocart}
\DeclareMathOperator{\cosk}{cosk}
\DeclareMathOperator{\plat}{plat}

\DeclareMathOperator{\faisc}{faisc}
\DeclareMathOperator{\sk}{sk}
\DeclareMathOperator{\ch}{ch}

\def\HH{\mathbf{H}}
\def\qc{\mathrm{qc}}
\def\et{\mathrm{et}}
\def\coh{\mathrm{coh}}

\def\Hvr#1{%
\begin{array}[b]{l}
{\kern2pt\scriptstyle r}\\[-8pt]
{\kern1pt\scriptstyle \vee}\\[-5pt]
H^{#1}
\end{array}}

\def\Hvv#1{%
\begin{array}[b]{l}
\kern2pt\rotatebox{90}{${\scriptscriptstyle\ll}$}\\[-6pt]
H^{#1}
\end{array}}

\def\Hvi#1{%
\begin{array}[b]{l}
{\kern0.2pt\scriptscriptstyle \infty}\\[-8pt]
{\kern1pt\scriptstyle \vee}\\[-5pt]
H^{#1}
\end{array}}

\def\Hv#1{%
\begin{array}[b]{l}
{\kern1pt\scriptstyle \vee}\\[-5pt]
H^{#1}
\end{array}}

\DeclareMathOperator{\cart}{cart}
\DeclareMathOperator{\Ffaisc}{\mathbf{Faisc}}


\def\Ev#1{%
\begin{array}[b]{l}
{\kern1pt\scriptstyle \vee}\\[-5pt]
E'_{#1}
\end{array}}

\DeclareMathOperator{\Cont}{Cont}
\DeclareMathOperator{\Ccont}{\mathscr{C}\!\mathit{ont}}
\DeclareMathOperator{\Site}{Site}
\newcommand\MM{\mathbf{M}}
\def\limtop#1{{\displaystyle{\varprojlim_{#1}\!\mbox{\rm top }}}}
\def\Limtop{\mathop{\varprojlim\!\mbox{\rm top}}}
\DeclareMathOperator{\spec}{spec}
\def\ss{\scriptstyle{\mathscr{S}}}
\DeclareMathOperator{\Gal}{Gal}
\def\smallm{\scriptstyle{\mathscr{M}}}
\DeclareMathOperator{\End}{End}
\DeclareMathOperator{\ass}{ass}
\DeclareMathOperator{\zar}{zar}
\DeclareMathOperator{\Ogg}{Ogg}
\DeclareMathOperator{\res}{res}
\DeclareMathOperator{\tr}{tr}
\DeclareMathOperator{\Supp}{Supp}
\DeclareMathOperator{\codim}{codim}
\DeclareMathOperator{\rk}{rk}
\DeclareMathOperator{\cur}{cur}
\DeclareMathOperator{\cdqc}{cdqc}
\newcommand\DD{\mathbf{D}}
\DeclareMathOperator{\cl}{cl}
\def\JJ{\mathbf{J}}
\def\FF{\mathbf{F}}
\def\BB{\mathbf{B}}
\DeclareMathOperator{\Alget}{Alget}
\DeclareMathOperator{\Ass}{Ass}
\DeclareMathOperator{\supp}{supp}
\DeclareMathOperator{\rad}{rad}
\DeclareMathOperator{\prof}{prof}
\DeclareMathOperator{\cint}{int}
\DeclareMathOperator{\sic}{sic}
\DeclareMathOperator{\Tor}{Tor}
\DeclareMathOperator{\tors}{tors}
\DeclareMathOperator{\Sym}{Sym}
\DeclareMathOperator{\parf}{parf}
\DeclareMathOperator{\tor}{tor}
\DeclareMathOperator{\defn}{defn}
\DeclareMathOperator{\Div}{Div}
\DeclareMathOperator{\Pol}{Pol}
\DeclareMathOperator{\Fex}{Fex}
\DeclareMathOperator{\Lie}{Lie}
\DeclareMathOperator{\tot}{tot}
\DeclareMathOperator{\Aut}{Aut}
\DeclareMathOperator{\HOM}{HOM}
\DeclareMathOperator{\pch}{pch}
\DeclareMathOperator{\Ch}{Ch}
\DeclareMathOperator{\EXT}{EXT}
\DeclareMathOperator{\PIC}{PIC}
\DeclareMathOperator{\Rf}{Rf}
\DeclareMathOperator{\inj}{inj}
\newcommand\VV{\mathbf{V}}

\def\Xv{%
\begin{array}[b]{l}
{\kern1pt\scriptstyle \vee}\\[-5pt]
X
\end{array}}

%%%%%%%
\newcommand\TT{\mathbf{T}}
\newcommand\XX{\mathbf{X}}
\newcommand\EE{\mathbf{E}}
%%%%%%

\def\qedsymbol{\textup{C.Q.F.D.}}
\let\CQFD\qed
\let\cqfd\qed

%%%%%\bigdot
\font\bigsymb=cmsy10 at 4pt
\def\bigdot{{\kern1.2pt\raise 1.5pt\hbox{\bigsymb\char15}}}
%%%%%

\def\Lright{\mathop{\vcenter{\hbox{\hbox{$\longrightarrow$}\lower8pt \hbox{\kern-16pt\hbox{$\longrightarrow$}}}}}}

\def\Aa{{\mathcal A}}
\def\Bb{{\mathcal B}}
\def\Cc{{\mathcal C}}
\def\Dd{{\mathcal D}}
\def\Ee{{\mathcal E}}
\def\Ff{{\mathcal F}}
\def\Gg{{\mathcal G}}
\def\Hh{{\mathcal H}}
\def\Ii{{\mathcal I}}
\def\Jj{{\mathcal J}}
\def\Kk{{\mathcal K}}
\def\Ll{{\mathcal L}}
\def\Mm{{\mathcal M}}
\def\Nn{{\mathcal N}}
\def\Oo{{\mathcal O}}
\def\Pp{{\mathcal P}}
\def\Qq{{\mathcal Q}}
\def\Rr{{\mathcal R}}
\def\Ss{{\mathcal S}}
\def\Tt{{\mathcal T}}
\def\Uu{{\mathcal U}}
\def\Vv{{\mathcal V}}
\def\Ww{{\mathcal W}}
\def\Xx{{\mathcal X}}
\def\Yy{{\mathcal Y}}
\def\Zz{{\mathcal Z}}

\newcommand{\CC}{\mathbf{C}}
\newcommand{\GG}{\mathbf{G}}
\newcommand{\KK}{\mathbf{K}}
\newcommand{\LL}{\mathbf{L}}
\newcommand{\NN}{\mathbf{N}}
\newcommand{\PP}{\mathbf{P}}
\newcommand\QQ{\mathbf{Q}}
\newcommand\RR{\mathbf{R}}
\newcommand\ZZ{\mathbf{Z}}

\def\to{\mathchoice{\longrightarrow}{\rightarrow}{\rightarrow}{\rightarrow}}
\def\from{\mathchoice{\longleftarrow}{\leftarrow}{\leftarrow}{\leftarrow}}
\def\mto{\mathchoice{\longmapsto}{\mapsto}{\mapsto}{\mapsto}}
\def\mfrom{\mathrel{\from\joinrel\mapstochar\,}}
\def\hto{\mathrel{\lhook\joinrel\to}}
\def\hfrom{\mathrel{\from\joinrel\rhook}}

\def\implies{\mathchoice{\Longrightarrow}{\Rightarrow}{\Rightarrow}{\Rightarrow}}
\def\implied{\mathchoice{\Longleftarrow}{\Leftarrow}{\Leftarrow}{\Leftarrow}}
\def\iff{\mathchoice{\Longleftrightarrow}{\Leftrightarrow}{\Leftrightarrow}{\Leftrightarrow}}

\def\To#1{\mathchoice{\xrightarrow{\labelstyle~~#1~}}{\stackrel{#1}{\longrightarrow}}{}{}}
\def\Mto#1{\mathrel{\mapstochar\To{#1}}}
\def\From#1{\mathchoice{\xleftarrow{\labelstyle~~#1~}}{\stackrel{#1}{\longleftarrow}}{}{}}

\newcommand{\isomto}{\overset{\sim}{\longrightarrow}}
\newcommand{\isomfrom}{\overset{\sim}{\longleftarrow}}

\newcommand{\IFF}{\unskip~$\iff$~\ignorespaces}
\newcommand{\THEN}{\unskip~$\implies$~\ignorespaces}

\let\leq\leqslant
\let\geq\geqslant
\let\le\leqslant
\let\ge\geqslant

\def\cf{{cf.}\kern.3em}
\def\Cf{{Cf.}\kern.3em}
\def\ie{{i.e.}\kern.3em}
\def\iev{{i.e.},\xspace}
%\def\loccito{loc.\kern2pt cit.\xspace}
%\def\loccitn{\textit{loc.\kern2pt cit.}\xspace}
%\def\loccit{\sisi{\loccito}{\loccitn}}
\def\EGA{\textup{EGA}\kern.2em}
\def\SGA{\textup{SGA}\kern.3em}

\let\oldS\S
\def\S{\oldS\kern .15em}
\def\ptbl{.\kern .2em}
\def\resp{resp.\kern.3em}
\def\Exp{Exp.\kern.3em}

\let\bar\overline
\let\hat\widehat
\let\tilde\widetilde
\let\dpl\displaystyle
\let\ldots\dots
\let\emptyset\varnothing
\let\skpt\skippointrait
\let\epsilon\varepsilon

\def\sfrac#1#2{{#1}/{#2}}

\DeclareMathOperator{\cd}{cd}
%\DeclareMathOperator{\Hhom}{\mathcal{H}\!\mathit{om}}
\DeclareMathOperator{\Hhom}{\mathbf{Hom}}
%\DeclareMathOperator{\Eext}{\mathcal{E}\!\mathit{xt}}
\DeclareMathOperator{\Eext}{\mathbf{Ext}}
\DeclareMathOperator{\Spec}{Spec}
\DeclareMathOperator{\Tr}{Tr}

\let\emptyset=\varnothing

%%%%%%%%%%%%%%%
%%% expose XVII
%%%%%%%%%%%%%%%

\makeatletter
\let\old@optexpose\optexpose
\newcommand{\empty@optexpose}[1]{}
\newcommand{\paren@optexpose}[1]{\old@optexpose{#1}(}
\newcommand{\refnoexpose}[1]{\let\optexpose\empty@optexpose\ref{#1}\let\optexpose\old@optexpose}%\refnoexpose{X.cor5.2} engendre "5.2" au lieu de "X 5.2"
\newcommand{\neweqref}[1]{\let\optexpose\paren@optexpose\textup{\ref{#1})}\let\optexpose\old@optexpose}%\neweqref{XII.eq4.1} engendre "XII (4.1)" au lieu de "(XII 4.1)"
\makeatother
\newcommand{\refvb}[2][VI${}_{\text{B}}$ ]{\sisi{#1\refnoexpose{#2}}{\ref{#2}}}

\newcommand{\leftexp}[2]{{\vphantom{#2}}^{#1}{#2}}
\newcommand{\AAA}{\mathbf{A}}
\newcommand{\Sf}{\mathfrak{S}}
\DeclareMathOperator*{\quotprojlim}{\text{``$\varprojlim$''}}
\DeclareMathOperator*{\quotinjlim}{\text{``$\varinjlim$''}}
\DeclareMathOperator{\Fl}{F\sisi{\ell}{l}}
\DeclareMathOperator{\Pl}{P\sisi{\ell}{l}}
%\DeclareMathOperator{\cl}{c\sisi{\ell}{l}}
\DeclareMathOperator{\tdf}{tdf}
\DeclareMathOperator{\rg}{rg}
\newcommand{\iInd}{\sisi{I}{i}nd}
\newcommand{\iind}{ind}
\newcommand{\ppro}{pro}
\newcommand{\tTor}{\sisi{T}{t}or}
\newcommand{\ttor}{tor}
\newcommand{\chap}{Chap.}
\newcommand{\Ffex}{\Ff\mathit{ex}}
\newcommand{\sheafk}{\underline{k}}
\newcommand{\sheafn}{\underline{n}}
\newcommand{\sheafGamma}{\underline{\Gamma}}

% Repeated references to an already typeset footnote must not use
% \footnotemark without a matching \footnotetext: Hyperref would then create
% a live annotation to a destination that cannot exist.  This helper keeps the
% current printed footnote mark and links it to an explicit stable note target.
\makeatletter
\newcommand{\sgarepeatedfootnotemark}[1]{%
  \begingroup
  \edef\@thefnmark{\thefootnote}%
  \hyperlink{#1}{\@makefnmark}%
  \endgroup}
\makeatother

\sisi{
\DeclareMathOperator{\simarrow}{\rule{.15mm}{1.8mm}\kern -3.4pt\rightsquigarrow}
\newcommand{\sheafell}{\underline{\ell}}
\newcommand{\sheafZZ}{\underline{\ZZ}}
}{
\newcommand{\nderef}[1]{la N.D.E.~(\ref{#1}) page~\pageref{#1}}
}

%%%%%%%%%%%%%%%%%
%%%%% expose XVII
%%%%%%%%%%%%%%%%%

% END INCLUDED SGA4_English_current_buildable_source_20260906/source/fichiers-configuration-latex/macros.tex


\externaldocument{../01/01_English.tex}
\externaldocument{../02/02_English.tex}
\externaldocument{../03/03_English.tex}
\externaldocument{../04/04_English.tex}
\externaldocument{../05/05_English.tex}
\externaldocument{../05bis/05bis_English.tex}
\externaldocument{../06/06_English.tex}
\externaldocument{../07/07_English.tex}
\externaldocument{../08/08_English.tex}
\externaldocument{../09/09_English.tex}
%\externaldocument{../10/10_English.tex}
\externaldocument{../11/11_English.tex}
\externaldocument{../12/12_English.tex}
\externaldocument{../13/13_English.tex}
\externaldocument{../14/14_English.tex}
\externaldocument{../15/15_English.tex}
\externaldocument{../16/16_English.tex}
\externaldocument{../17/17_English.tex}

\setboolean{original}{false}

\begin{document}

\begin{center}
XVIII. The global duality formula\\
P. Deligne
\end{center}

\selectlanguage{english}version: \showgitstatus
\tableofcontents
\mainmatter

\else
\chapter{The global duality formula}
%\addtocontents{toc}{P. Deligne}
\begin{center}
P. Deligne
\end{center}
\fi
\phantomsection\MakeLinkTarget*{sga4.exp.XVIII}\label{sga4.exp.XVIII}

% The edition is fixed to the corrected (second) reading. This expandable
% selector is output-equivalent and avoids current Xy-pic conditional failures.
\renewcommand{\sisi}[2]{#2}
\let\Arrowvert\Vert

% BEGIN INCLUDED SGA4_English_current_buildable_source_20260906/source/18/drafts/SGA4_XVIII_LINES_44_450_English.texfrag
\setcounter{pageoriginal}{480}
\newcommand{\acorriger}[1]{\nde{#1}}

\setcounter{section}{-1}
\section{Introduction}\etiquette[0]{XVIII.sec0}\pageoriginale

\subsection{}\etiquette[0.1]{XVIII.subsec0.1}
The present exposé is devoted to Poincaré duality. This duality is
constructed on the model of \textsc{Verdier} [1], who treats the case of
locally compact Hausdorff topological spaces. We show a priori that the
functor $Rf_{!}$ of Exposé XVII admits a right adjoint $Rf^{!}$ and
establish many of its properties, before computing it more explicitly
in the case of smooth morphisms.

The possibility of such a construction seems to rest on the following
properties of the functor $Rf_{!}$:
\begin{enumerate}
\renewcommand{\theenumi}{\alph{enumi}}
\renewcommand{\labelenumi}{\theenumi)}
\item For the existence of an adjoint $Rf^{!}$ and its formal
  properties:
\begin{itemize}
\item[(I)] the functor $Rf_{!}$ commutes with base change;

\item[(II)] the functors $R^{i}f_{!}$ commute with filtered inductive
  limits; they vanish for sufficiently large $i$;

\item[(III)] the functor $Rf_{!}$ can be computed ``at the level of
  complexes'', \ie starting from a functor between categories of
  complexes.
\end{itemize}

\item For the computation of this adjoint when $f$ is smooth:
\begin{itemize}
\item[(IV)] the computation of $Rf_{!}((\ZZ/n\ZZ)_{U})$ for
  ``small'' $U$.
\end{itemize}
\end{enumerate}

\subsection{}\etiquette[0.2]{XVIII.subsec0.2}
In \hyperref[XVIII.sec3]{\S 3} nᵒ 1, which is independent of \S \hyperref[XVIII.sec1]{\S 1}, 2, we
construct the functor $Rf^{!}$; constructions specific to the situation
are used there, rather than only (I), (II), and (III). This makes this
nᵒ rather unpleasant; the author confesses that he did not
always fully understand what was happening.

\S \hyperref[XVIII.sec2]{2} is devoted to point (IV); it is purely formal
from\pageoriginale \hyperref[XVIII.sec1]{\S 1}, where the cohomology of curves is studied.

In \S \hyperref[XVIII.sec1]{1}, we prove substantially more than is needed for
\S \hyperref[XVIII.sec2]{2}, also studying the case of continuous
coefficients. This section was profoundly influenced by \textsc{Serre}
\cite{sga4.exp.XVIII.bib.6}. The points useful in what follows are:
\begin{enumerate}
\renewcommand{\theenumi}{\alph{enumi}}
\renewcommand{\labelenumi}{\theenumi)}
\item nᵒ \hyperref[XVIII.subsec1.1]{1.1}, devoted to the trace morphism in the
  case of curves, and generalized in \hyperref[XVIII.thm2.9]{2.9};

\item the vanishing theorem \hyperref[XVIII.lemfond1.6.9]{1.6.9}. A reader willing to accept a
  transcendental argument may read nᵒ \hyperref[XVIII.subsec1.6]{1.6}, beginning
  with \hyperref[XVIII.lem1.6.6]{1.6.6} (2\textsuperscript{nd} proof), independently of nᵒ \hyperref[XVIII.subsec1.2]{1.2}
  to \hyperref[XVIII.subsec1.5]{1.5}\nde{On a first reading, one cannot too strongly recommend
  following this route: after $\S$\hyperref[XVIII.subsec1.1]{1.1}, proceed
  directly to $\S$\hyperref[XVIII.sec2]{2} and $\S$\hyperref[XVIII.sec3]{3}.}.
  The essential ingredient in \hyperref[XVIII.lemfond1.6.9]{1.6.9} is either
  form of Poincaré duality for smooth curves over an algebraically
  closed field, together with the acyclicity lemma \hyperrefext[XV][XV.lem2.6]{2.6} (a preliminary
  to the smooth base-change theorem \hyperrefext[XVI][XVI.thm1.1]{1.1}). Its conclusion is
  generalized in \hyperref[XVIII.thm2.14]{2.14}.
\end{enumerate}

Number \hyperref[XVIII.subsec3.2]{3.2} collects the consequences of \S
\hyperref[XVIII.sec2]{2}.

\subsection{}\etiquette[0.3]{XVIII.subsec0.3}
\textsc{Verdier}'s method in Poincaré duality makes it possible to define the
functor $Rf^{!}$ for compactifiable $f$. The price of this generality is
that the compatibilities to be verified are unfamiliar and sometimes,
even for the most trivial ones, such as \hyperref[XVIII.lem3.2.3]{3.2.3}, require
outlandish proofs. The author admits that he has not proved as many
of them as he should have.

An interested reader will have no difficulty in extending, using the
appendix to Exposé XVII, the definition of $Rf^{!}$ to separated
morphisms of finite type whose target is a coherent scheme (\ie a
quasi-compact quasi-separated scheme).

\subsection{}\etiquette[0.4]{XVIII.subsec0.4}
The main result \hyperref[XVIII.thm3.2.5]{3.2.5} is stated in concrete terms that
would make it possible to return formally to the viewpoint taken in the
oral seminar, where the functor $Rf^{!}$ was defined only for\pageoriginale
smoothable morphisms, by means of a factorization through a smooth
morphism and an immersion. An outline of the proof from the oral
seminar can be found in \textsc{Verdier} \cite{sga4.exp.XVIII.bib.8}.

\section{Cohomology of curves}\etiquette[1]{XVIII.sec1}\pageoriginale

\subsection{The trace morphism}\etiquette[1.1]{XVIII.subsec1.1}

\subsubsection{}\etiquette[1.1.1]{XVIII.subsubsec1.1.1}
Let $n$ be an integer $\geq 1$ and let $X$ be a scheme on which $n$ is
invertible. The sheaf $𝛍_n$ (IX \hyperref[IX.eq3.1]{3.1}) on $X$ is then a sheaf of modules
locally free of rank 1 over the constant sheaf of rings $\ZZ/n$.
For $i\in \ZZ$, we denote by $\ZZ/n(i)$ its
$i^{\text{\rm th}}$ tensor power:
\begin{equation*}
\ZZ/n(i)=\Hhom(\ZZ/n,\GG_{m})^{\otimes
    i}{\quoi}.\tag{1.1.1.1}
\end{equation*}\etiquette[1.1.1.1]{XVIII.eq1.1.1.1}

If $n=dn'$, exponentiation by $d$ gives an isomorphism
$$
𝛍_n\otimes_{\ZZ/n}\ZZ/n'\xrightarrow{\sim}𝛍_{n'},
$$
and hence isomorphisms, called canonical:
\begin{equation*}
\ZZ/n(i)\otimes_{\ZZ/n}\ZZ/n'\xrightarrow{\sim} \ZZ/n'(i){\quoi}.\tag{1.1.1.2}
\end{equation*}\etiquette[1.1.1.2]{XVIII.eq1.1.1.2}

If $F$ is an abelian sheaf such that $nF=0$, set
\begin{equation*}
F(i)=F\otimes_{\ZZ/n}\ZZ/n(i){\quoi}.\tag{1.1.1.3}
\end{equation*}\etiquette[1.1.1.3]{XVIII.eq1.1.1.3}%
If already $n'F=0$, the isomorphism \eqhyperref[XVIII.eq1.1.1.2]{1.1.1.2} allows us to identify
$F\otimes_{\ZZ/n}\ZZ/n(i)$ with
${F\otimes_{\ZZ/n'}\ZZ/n'(i)}$, so that $F(i)$ depends only
on $F$ and not on the choice of $n$. More generally, if $F$ is a
torsion sheaf prime to the residual characteristics of X
(\hyperrefext[XVII][XVII.subsec0.13]{0.13}), and if ${}_{m}F$ is the kernel of multiplication
by $m$ on $F$, we define $F(i)$ to be the inductive limit, as $m$ tends
multiplicatively to infinity through the integers invertible on $X$:
\begin{equation*}
F(i)=\varinjlim_{m}{}_{m}F(i){\quoi}.\tag{1.1.1.4}
\end{equation*}\etiquette[1.1.1.4]{XVIII.eq1.1.1.4}

The functor $F\mapsto F(i)$ is exact. If $f:Y\rightarrow X$ is a
morphism of schemes (\resp a quasi-compact quasi-separated morphism
of schemes), and if $F$ is a torsion sheaf on $X$ (\resp on $Y$),
prime\pageoriginale to the residual characteristics of $X$, then
$f^{\ast}(F(i))\simeq(f^{\ast}F)(i)$ (\resp
$(f_{\ast}F)(i)\simeq f_{\ast}(F(i))$). If $\mathcal{A}$ is a sheaf of
rings on $X$, annihilated by an integer $n$ invertible on $X$, this
extends to an exact functor $K\mapsto K(i)$ from $D(X,\mathcal{A})$ to
$D(X,\mathcal{A})$. This functor commutes with inverse-image functors
and direct-image functors.

The functors $F\mapsto F(i)$ are sometimes called ``\emph{Tate twist}'' functors.

\begin{enonce*}{Definition 1.1.2}\etiquette[1.1.2]{XVIII.def1.1.2}
A \emph{flat curve} over a scheme $S$ is a flat, separated morphism
$f:X\rightarrow S$ of finite presentation whose fibers are of pure
dimension one.
\end{enonce*}

The expression ``of pure dimension one'' does not exclude the empty
scheme. Contrary to \EGA II 7.4.2, a curve over a field is allowed here
to be empty, but is required to be separated.

When $S$ is the spectrum of a field, we shall simply speak of a curve
over $S$; a curve over a field is quasi-projective\nde{In \EGA II
7.4.10, this statement is proved only under the additional assumption
that $X$ is normal, and it is announced (\EGA II 7.4.12) that a proof in
the general case will appear in \EGA V, which has not been published.
The reader may fill this gap as an exercise. In any case, this has no
bearing on the construction of the trace morphism, since it can be
carried out first locally (on quasi-projective open subschemes) and the
trace morphisms can then be glued as in the proof of
\hyperref[XVIII.prop1.1.6]{1.1.6}.}.

\setcounter{subsubsection}{2}
\subsubsection{}\etiquette[1.1.3]{XVIII.subsubsec1.1.3}
Let $X$ be a curve over an algebraically closed field $k$ of
characteristic exponent $p$ prime to an integer $n\geq 1$. First
suppose that $X$ is reduced, and let $\bar{X}$ be a complete curve over
$k$ in which $X$ is a dense \emph{open subscheme}.

The complement $Y=\bar{X}-X$ is of dimension $0$. The exact cohomology
sequence (\hyperrefext[XVII][XVII.eq5.1.16.3]{5.1.16.3}) therefore induces an isomorphism
\begin{equation*}
H^{2}_{c}(X,\ZZ/n(1))\xrightarrow{\sim} H^{2}(\bar{X},\ZZ/n(1)){\quoi}.\tag{1.1.3.1}
\end{equation*}\etiquette[1.1.3.1]{XVIII.eq1.1.3.1}%
We also have an isomorphism \hyperrefext[IX][IX.cor4.7]{4.7} (given by Kummer
theory)
\begin{equation*}
H^{2}(\bar{X},\ZZ/n(1))=\Pic(\bar{X})/n=(\ZZ/n)^{c}\tag{1.1.3.2}
\end{equation*}\etiquette[1.1.3.2]{XVIII.eq1.1.3.2}%
where $c$ denotes the set of irreducible components of $X$, or of
$\bar{X}$, which amounts to the same thing.

If\pageoriginale $X$ is no longer necessarily reduced, and if $c$ is
the set of its irreducible components, then for $i\in c$ let $n_{i}$
denote the multiplicity of $X$ at the generic point $\eta_{i}$ of the
$i^{\text{\rm th}}$ irreducible component of $X$:
$$
n_{i}=\lg(O_{X,\eta_{i}}){\quoi}.
$$

We know that the restriction map from $H^{2}_{c}(X,\ZZ/n(1))$ to
$H^{2}_{c}(X_{\red}, \ZZ/n(1))$ is an isomorphism. Thus,
via \eqhyperref[XVIII.eq1.1.3.1]{1.1.3.1} and \eqhyperref[XVIII.eq1.1.3.2]{1.1.3.2} applied to
$X_{\text{\rm red}}$, there is a canonical isomorphism between
$H^{2}_{c}(X,\ZZ/n(1))$ and $(\ZZ/n)^{c}$. We denote by $\Tr_{X}$,
or simply by $\Tr$, the composite map
\begin{equation*}
\Tr:H^{2}_{c}(X,\ZZ/n(1))\xrightarrow{\sim}(\ZZ/n)^{c}\xrightarrow{t}\ZZ/n{\quoi},\tag{1.1.3.3}
\end{equation*}\etiquette[1.1.3.3]{XVIII.eq1.1.3.3}%
where
$$
t((a_{i})_{i\in c})=\sum n_{i}a_{i}{\quoi}.
$$

If $F$ is an abelian group annihilated by $n$, there is an
isomorphism
$$
H^{2}_{c}(X,F(1))\mathop{\longleftrightarrow}^{\sim}F\otimes
H^{2}_{c}(X,\ZZ/n(1))
$$
and hence, via \eqhyperref[XVIII.eq1.1.3.3]{1.1.3.3}, a \emph{trace morphism}
\begin{equation*}
\Tr:H^{2}_{c}(X,F(1))\longrightarrow F{\quoi}.\tag{1.1.3.4}
\end{equation*}\etiquette[1.1.3.4]{XVIII.eq1.1.3.4}

If $n=n'd$, the diagram
$$\vcenter{\etiquette[1.1.3.5]{XVIII.eq1.1.3.5}
\xymatrix@R=1.25cm{
0  \ar[r] & 𝛍_n  \ar[r]\ar[d]_{x^d} & \GG_m \ar[d]_{x^d} \ar[r]^{x^n} & {\GG}_m 
\ar[r] \ar@{=}[d]& 0\\
0  \ar[r] & 𝛍_{n'}  \ar[r] & \GG_m  \ar[r]^{x^{n'}} & {\GG}_m 
\ar[r] & 0
}}\leqno{(1.1.3.5)}
$$
is commutative, and therefore so is the diagram
$$\vcenter{\etiquette[1.1.3.6]{XVIII.eq1.1.3.6}
  \xymatrix@C=2cm@R=1.25cm{
H^{2}_c(X,\ZZ/n(1)) \ar[d] \ar[r]^-{\sim} & (\ZZ/n)^{c} \ar[d]\\
H^{2}_c(X,\ZZ/n'(1)) \ar[r]^-{\sim}  & (\ZZ/n')^{c}
}}\leqno{(1.1.3.6)}
$$
by\pageoriginale the definition in \hyperrefext[IX][IX.cor4.7]{4.7} of the isomorphism
\eqhyperref[XVIII.eq1.1.3.2]{1.1.3.2}.

The map \eqhyperref[XVIII.eq1.1.3.4]{1.1.3.4} thus depends only on $F$, and not on the choice
of the integer $n$ prime to $p$ such that $nF=0$. By passage to the
limit, we define it for every torsion abelian group $F$ prime to $p$.

The following lemma follows immediately from the definitions.

\begin{enonce*}{Lemma 1.1.4}\etiquette[1.1.4]{XVIII.lem1.1.4}
Under the preceding hypotheses, if, for $i\in c$, $U_{i}$ is a nonempty
open subscheme of $X$ contained in its $i^{\text{\rm th}}$ irreducible
component, the diagram
$$
\xymatrix@R=1.5cm@C=1.25cm{
\displaystyle{\mathop{\oplus}_{i\in c}} H^2_c(U_i, F(1))
\ar[rr]^{\sim} \ar[rd]_{\displaystyle{\mathop{\oplus}_{i\in c}}
  Tr_{U_i}}& & H^2_c(X,F(1))\ar[ld]^{Tr_X}\\
 & F &
}
$$
is commutative.
\end{enonce*}

\begin{enonce*}{Lemma 1.1.5}\etiquette[1.1.5]{XVIII.lem1.1.5}
Let $X$ and $Y$ be two curves over an algebraically closed field $k$,
let $u:X\rightarrow Y$ be a quasi-finite flat morphism, and let $n$ be
an integer prime to the characteristic exponent $p$ of $k$.

The following diagram, in which $\Tr_{u}$ is the map
\emph{(\hyperrefext[XVII][XVII.thm6.2.3]{6.2.3}),} is then commutative:
$$
\xymatrix@R=1.5cm{
  H^2_c(X,\ZZ/n(1)) \ar[rr]^{Tr_u} \ar[rd]_{Tr_X}&              & H^2_c(Y,\ZZ/n(1)) \ar[ld]^{Tr_Y}\\
                         & \ZZ/n &. 
}
$$
\end{enonce*}

Lemma \hyperref[XVIII.lem1.1.4]{1.1.4} allows us to reduce to the case where $X$ and $Y$ are
irreducible, with $X_{\text{\rm red}}$ and $Y_{\text{\rm red}}$ smooth
as well. If $e$ (\resp $f$) is the multiplicity of $X$ (\resp $Y$) at
the generic point, and if $e=df$, the triangles marked $+$ in the
following diagram are commutative:
$$
\vcenter{\etiquette[1.1.5.1]{XVIII.eq1.1.5.1}
\xymatrix@R=1.5cm{
H^2_c(X,\ZZ/n(1))
\ar[rr]^{\sim}_(.5){{\displaystyle\mathop{\phantom{a}}_{\displaystyle\mathop{{\phantom{a}}}_{{\displaystyle\mathop{+}}}}}}  \ar[dd]_{\Tr_u} \ar[rd]_{\Tr_X}& &
H^2_c(X_{\red},\ZZ/n(1)) \ar[dd]^{d.\Tr_{u_{\red}}} \ar[ld]^{e.\Tr_{X_{\red}}}\\
                               & \ZZ/n &\\
H^2_c(Y,\ZZ/n(1)) \ar[rr]^{\displaystyle\mathop{\sim}^{\displaystyle\mathop{\phantom{a}}^{{\displaystyle\mathop{\phantom{a}}^{\textstyle{+}}}}}} \ar[ru]^{\Tr_Y}&  & H^2_c(Y_{\red},\ZZ/n(1))\ar[lu]_{f.\Tr_{Y_{\red}}}, 
}}\sisi{}{\leqno{(1.1.5.1)}}
$$

Its\pageoriginale outer boundary is commutative, as follows from the
commutativity of the diagram
$$
\xymatrix@C=2cm@R=1.5cm{
u_!\ZZ/n_X(1)           \ar[r]^{\sim} \ar[d]^{\Tr_u} &
u_{{\red}_!}\ZZ/n_{X_{\red}}(1) \ar[d]^{d.\Tr_{u_{\red}}}\\
\ZZ/n_Y(1)  \ar[r]^{\sim}& \ZZ/n_{Y_{\red}}(1).}
$$
It suffices to verify this last compatibility at the generic point of
$Y$, which is trivial.

\sisi{}{To establish the commutativity of the left-hand triangle in
\eqhyperref[XVIII.eq1.1.5.1]{1.1.5.1}, it therefore remains only to establish that of the
right-hand triangle. Hence we may now assume in addition that $X$ and $Y$
are reduced\nde{This paragraph was added in this edition, and the next
two were modified.}.}

Let $\bar{X}$ and $\bar{Y}$ be the complete \sisi{non-smooth}{smooth}
curves containing \sisi{$X_{\red}$}{$X$} and \sisi{$Y_{\red}$}{$Y$} as dense
open subschemes (\EGA II 7.4.11). The morphism \sisi{$u_{\red}$}{$u$}
extends to a flat morphism $\bar{u}:\bar{X}\rightarrow \bar{Y}$
(\EGA II 7.4.9), and by \hyperref[XVIII.lem1.1.4]{1.1.4} it suffices to verify
\hyperref[XVIII.lem1.1.5]{1.1.5} for $\bar{u}$.

By (XVII \hyperref[XVII.eq6.3.18.2]{6.3.18.2}), the diagram of sheaves on $\bar{Y}$ below, in
which $N$ denotes the norm,
$$
\xymatrix@R=1.25cm{
0  \ar[r] & \bar{u}_{\ast}𝛍_n  \ar[r] \ar[d]^{\Tr_{\bar{u}}}&
\bar{u}_{\ast}\GG_m  \ar[r]^{x^n} \ar[d]^{N}&
\bar{u}_{\ast} \GG_m  \ar[r] \ar[d]^{N}& 0 \\
0  \ar[r] & 𝛍_n  \ar[r] & \GG_m  \ar[r]^{x^n} & \GG_m 
\ar[r] & 0 ,
}
$$
is commutative, and hence so is the diagram
$$
\xymatrix@C=1.5cm@R=1.25cm{
\Pic(\bar{X})/n \ar[r]^(.45){\sim}\ar[d]^{N}& H^2_c(\bar{X},\ZZ/n(1))
\ar[d]^{\Tr_{\bar{u}}}\\ 
\Pic(\bar{Y})/n \ar[r]^(.45){\sim}& H^2_c(\bar{Y},\ZZ/n(1)),\\
}
$$
and\pageoriginale it remains to observe that if $\Ll$ is an invertible sheaf
of degree 1 on $\bar{X}$, then its norm again has degree $1$ (this fact
following from the compatibility \sisi{\EGA IV 21.10.7}{\EGA IV 21.10.17},
taking into account the definition \EGA IV 21.5.5 of the norm of a
divisor).

\begin{enonce*}{Proposition 1.1.6}\etiquette[1.1.6]{XVIII.prop1.1.6}
There is a unique way to define, for every compactifiable flat curve
$f:X\rightarrow S$ and every torsion sheaf $F$ on $S$ prime to the
residual characteristics of $S$, a trace morphism
$$
\Tr_{f}:R^{2}f_{!}(f^{\ast}F(1))\longrightarrow F{\quoi},
$$
functorial in $F$, whose formation is compatible with every base change
\emph{(\cf \hyperrefext[XVII][XVII.thm6.2.3]{6.2.3}, (VAR 2))}, and which, when $S$ is the spectrum
of an algebraically closed field, coincides with the trace morphism
\eqhyperref[XVIII.eq1.1.3.4]{1.1.3.4}.
\end{enonce*}

The morphism $\Tr_{f}$ is prescribed fiber by fiber; its uniqueness is
therefore clear.

Let $p:\PP^{1}_{S}\rightarrow S$ be the projective line over $S$, and let $n$ be
an integer invertible on $S$. We have
$\sheafPic_{S}(\PP^{1}_{S})=\ZZ_{S}$, whence Kummer theory
(\hyperrefext[IX][IX.Thr3.2]{3.2}) gives a morphism $\ZZ/n\rightarrow
R^{1}p_{\ast}\ZZ/n(1)$ which one checks fiber by fiber
(\hyperrefext[XII][XII.cor5.2]{5.2}) to be an isomorphism. Its inverse $\Tr_{p}$ induces the trace
morphism \eqhyperref[XVIII.eq1.1.3.4]{1.1.3.4} fiber by fiber.

Let $f:X\rightarrow S$ be a flat curve over $S$ such that there exists a
quasi-finite flat $S$-morphism $u$ from $X$ to $\PP^1_S$. By
\hyperref[XVIII.lem1.1.5]{1.1.5}, the composite morphism $\Tr_{p}\circ \Tr_{u}$
(\hyperrefext[XVII][XVII.thm6.2.3]{6.2.3}) has, at every geometric point of $S$, the trace morphism
\eqhyperref[XVIII.eq1.1.3.4]{1.1.3.4} as its fiber. In particular, it does not depend on the
choice of $u$.

Suppose\pageoriginale that $X$ is a union of open subschemes $U_{i}$ such
that there exists a quasi-finite flat morphism from $U_{i}$ to
$\PP^{1}_S$. If $\alpha_{i}$ (\resp $\alpha_{ij}$) is the inclusion of
$U_{i}$ (\resp $U_{ij}=U_{i}\cap U_{j}$) in $X$, the sequence
$$
\sum R^{2}(f\alpha_{ij})_{!}\ZZ/n(1)\rightrightarrows
\sum R^2(f\alpha_{i})_{!}\ZZ/n(1)\rightarrow
R^{2}f_{!}\ZZ/n(1)\rightarrow 0{\quoi}.
$$
is exact, by (\hyperrefext[XVII][XVII.prop6.2.9]{6.2.9}) and the right exactness of
$R^{2}f_{!}$ (since $R^{3}f_{!}=0$). One checks fiber by fiber, using
\hyperref[XVIII.lem1.1.5]{1.1.5}, that the sum of the trace morphisms
$\underset{i}\sum R^{2}(f\alpha_{i})_{!}\ZZ/n(1)\rightarrow
\ZZ/n$ factors through a trace morphism $\Tr_{f}$ from
$R^{2}f_{!}\ZZ/n(1)$ which induces \eqhyperref[XVIII.eq1.1.3.4]{1.1.3.4} fiber by fiber.

% END INCLUDED SGA4_English_current_buildable_source_20260906/source/18/drafts/SGA4_XVIII_LINES_44_450_English.texfrag

% BEGIN INCLUDED SGA4_English_current_buildable_source_20260906/source/18/drafts/SGA4_XVIII_LINES_451_855_English.texfrag

For a general $X$, let $j:X'\rightarrow X$ be the largest open subset
of $X$ such that $f|X'$ is a Cohen--Macaulay morphism (\cf \EGA IV
12.1.1 (vi)). Every point $x$ of $X'$ has a neighborhood $U_{x}$ in
$X'$ for which there exists a quasi-finite $S$-morphism $u_{x}$ from
$U_{x}$ to $\PP^{1}_{S}$. This morphism is automatically flat (\EGA IV
11.3.10 and 15.4.2), and $X'$ therefore satisfies the preceding
hypotheses. Moreover, $X'$ is fiberwise dense in $X$; the morphism from
$R^{2}(fj)_{!}\ZZ/n(1)$ to $R^{2}f_{!}\ZZ/n(1)$ is consequently an
isomorphism, as can be checked fiberwise, and the composite arrow
$$
\Tr_{f}:R^{2}f_{!}\ZZ/n(1)\xleftarrow{\sim}
R^{2}(fj)_{!}\ZZ/n(1)\xrightarrow{\Tr_{fj}}\ZZ/n
$$
induces \eqhyperref[XVIII.eq1.1.3.4]{1.1.3.4} on every fiber.

If $F$ satisfies $nF=0$, one has (\hyperrefext[XVII][XVII.thm5.2.6]{5.2.6})
\begin{equation*}
R^{2}f_{!}\ZZ/n(1)\otimes F\xrightarrow{\sim}
R^{2}f_{!}(f^{\ast}F(1))\tag{1.1.6.1}
\end{equation*}\etiquette[1.1.6.1]{XVIII.eq1.1.6.1}%
and $\Tr_{f}$ is defined to be the composite arrow
$$
\Tr_{f}:R^{2}f_{!}(f^{\ast}F(1))\xleftarrow{\sim}
R^{2}f_{!}\ZZ/n(1) \otimes F\rightarrow F{\quoi}.
$$
In the general case, $\Tr_{f}$ is defined as the filtered colimit of
the trace morphisms relative to the sheaves ${}_nF$, for $n$
invertible on $S$. These arrows induce \eqhyperref[XVIII.eq1.1.3.4]{1.1.3.4} on every fiber
and therefore commute with every base change, \cqfd

One\pageoriginale checks fiberwise, using \hyperref[XVIII.lem1.1.4]{1.1.4},
\hyperref[XVIII.lem1.1.5]{1.1.5}, and \eqhyperref[XVIII.eq1.1.6.1]{1.1.6.1}, the following results:

\begin{enonce*}{Lemma 1.1.7}\etiquette[1.1.7]{XVIII.lem1.1.7}
Let $f:Y\rightarrow S$ be a compactifiable flat curve over $S$, let
$u:X\rightarrow Y$ be a quasi-finite flat morphism of finite
presentation, and let $F$ be a torsion sheaf on $S$, of torsion prime
to the residue characteristics of $S$. The following diagram is
commutative:
$$
\xymatrix@C=1.5cm@R=1.25cm{
R^2(fu)_! (fu)^{\ast}F(1) \ar[d]^{\Tr_{fu}}\ar@{=}[r]& R^2f_!(u_!
u^{\ast})f^{\ast}F(1) \ar[r]^{\Tr_u}& R^2f_!f^{\ast}F(1)\ar[d]^{\Tr_f} \\
F     \ar@{=}[rr] &&  F,
}
$$
\end{enonce*}

\begin{enonce*}{Lemma 1.1.8}\etiquette[1.1.8]{XVIII.lem1.1.8}
Let $u:Y\rightarrow S$ be a quasi-finite flat morphism of finite
presentation, let $f:X\rightarrow Y$ be an $S$-compactifiable flat
curve \emph{(\hyperrefext[XVII][XVII.def3.2.1]{3.2.1})}, and let $F$ be as in \hyperref[XVIII.lem1.1.7]{1.1.7}. The
following diagram is commutative\nde{This compatibility merits a few
arguments. To verify it, one may first reduce by base change to the case
where $S$ is the spectrum of an algebraically closed field $k$. One may
then suppose that $Y$ is connected and nonempty, in which case $Y\to S$
has a (unique) section $S\to Y$ that is a closed immersion defined by a
nilpotent ideal. The arguments in the proof of \hyperref[XVIII.prop1.1.6]{1.1.6}
allow us to suppose that $X=\mathbf{P}^1_Y$. It remains only to observe
that the multiplicity of $Y$ equals the multiplicity of $X$ at its
generic point. See also the proof of the commutativity of
\eqhyperref[XVIII.eq2.9.3]{2.9.3}.}:
$$
\xymatrix@C=1.5cm@R=1.25cm{
R^{2}(uf)_! (uf)^{\ast}F(1)  \ar@{=}[r] \ar[d]^{\Tr_{uf}}&
u_!R^2f_!f^{\ast} u^{\ast}F(1) \ar[r]^{\Tr_f} &
u_!u^{\ast}F \ar[d]^{\Tr_u}\\
F        \ar@{=}[rr] & &  F
}
$$
\end{enonce*}

One similarly verifies from \eqhyperref[XVIII.eq1.1.3.2]{1.1.3.2}:

\begin{enonce*}{Lemma 1.1.9}\etiquette[1.1.9]{XVIII.lem1.1.9}
If $f:X\rightarrow S$ is a compactifiable smooth curve and the
geometric fibers of $f$ are irreducible, then the trace morphism
(\hyperref[XVIII.prop1.1.6]{1.1.6}) is an isomorphism.
\end{enonce*}

\subsection{$1$-acyclicity of projective
  space}\etiquette[1.2]{XVIII.subsec1.2}

In this section, the fppf topology is in principle always used
(\hyperrefext[XVII][XVII.subsec0.10]{0.10}).

\subsubsection{}\etiquette[1.2.1]{XVIII.subsubsec1.2.1}

Let $\epsilon$ be a locally free module on a scheme $S$, everywhere of
rank\pageoriginale $\geq 2$, let $p:\PP(\epsilon)\rightarrow S$ be the
corresponding projective bundle, and let $G$ be an abelian sheaf on the
big fppf site of $S$ (\hyperrefext[XVII][XVII.subsec0.10]{0.10}). Each pair $(K_{0},\varphi)$ consisting of a
torsor $K_{0}$ under $G$ and a homomorphism
$\varphi:\GG_{mS}\rightarrow G$ defines a $p^{\ast}G$-torsor
$K(K_{0},\varphi)$ on $\PP(\epsilon)$, the sum of $K_{0}$ and the image
under $\varphi$ of the $\GG_{m}$-torsor $\Oo(1)$\nde{In this section,
we implicitly identify line bundles and $\GG_{m}$-torsors by the
correspondence that associates with a line bundle the sheaf of its
invertible sections, on which $\GG_{m}$ acts by multiplication. A reader
concerned with details may wish to know that, following Grothendieck's
convention, the projective scheme $\PP(\epsilon)$ is the bundle of
hyperplanes in $\epsilon$, and likewise the punctured vector bundle
$\VV(\epsilon)^{\ast}\to \PP(\epsilon)$ is an open subset of
$\VV(\epsilon)$, which is the spectrum of the symmetric algebra of
$\epsilon$.}:
\begin{equation*}
K(K_{0},\varphi)=K_{0}+\varphi \Oo(1){\quoi}.\tag{1.2.1.1}
\end{equation*}\etiquette[1.2.1.1]{XVIII.eq1.2.1.1}%
The pairs $(K_{0},\varphi)$ are made into a category by setting
$$
\Hom((K_{0},\varphi),(K'_{0},\varphi'))=
\begin{cases}
∅\quad\mbox{\rm if}\quad \varphi\neq \varphi'\\
\Hom(K_{0},K'_{0})\quad\mbox{\rm if}\quad \varphi=\varphi'{\quoi}.
\end{cases}
$$
The torsor $K(K_{0},\varphi)$ then depends functorially on
$(K_{0},\varphi)$. It is compatible with localization and therefore
defines a functor from the stack over $S$ of pairs $(K_{0},\varphi)$ to
the stack over $S$ of $p^{\ast}G$-torsors on $\PP(\epsilon)$ (= the
stack whose sections over $U/S$ are the $p^{\ast}G$-torsors on
$\fprod{U}{\PP}{S}(\epsilon)$). These stacks are stacks in groupoids
(= every morphism over a $U/S$ is an isomorphism).

\begin{enonce*}{Theorem 1.2.2}\etiquette[1.2.2]{XVIII.thm1.2.2}
Under the hypotheses of \hyperref[XVIII.subsubsec1.2.1]{1.2.1}, if $G$ is a flat commutative
group scheme of finite presentation over $S$, then the functor $K$ of
\hyperref[XVIII.subsubsec1.2.1]{1.2.1} is an equivalence.
\end{enonce*}

\setcounter{subsubsection}{2}
\paragraph{}\etiquette[1.2.2.1]{XVIII.par1.2.2.1}

This assertion is equivalent to the conjunction of
\begin{gather*}
G\xrightarrow{\sim} p_{\ast}p^{\ast}G\quad\mbox{\rm and}\\
\Hhom(\GG_{m},G)\xrightarrow{\sim} R^{1}p_{\ast}p^{\ast}G{\quoi}.
\end{gather*}
Indeed, a morphism of stacks in groupoids is an equivalence if and
only if it induces an isomorphism on the sheaf generated by the
presheaf of isomorphism classes of objects (the 2\textsuperscript{nd}
formula) and an isomorphism on the sheaf of automorphisms of any local
object (the 1\textsuperscript{st} formula).

We shall give two proofs (\hyperref[XVIII.lem1.2.4]{1.2.4} and
\hyperref[XVIII.subsubsec1.2.5]{1.2.5}) of full faithfulness\nde{With regard to
essential surjectivity, the reader should be warned that the proof will
use the general proof technique explained in the introduction to the
\EGA.}.

\begin{enonce*}{Lemma 1.2.3}\etiquette[1.2.3]{XVIII.lem1.2.3}
Let\pageoriginale $f:X\rightarrow S$ be a proper flat morphism. Suppose
that $f_{\ast}\Oo_{X}=\Oo_{S}$ universally and
$R^{1}f_{\ast}\Oo_{X}=0$ universally. Every $S$-morphism from $X$ to a
finite-type $S$-group scheme $G$ then factors through $f$ and a section
of $G$; more precisely,
$G\xrightarrow{\sim} f_{\ast}f^{\ast}G$.
\end{enonce*}

It suffices to prove the 1\textsuperscript{st} assertion.

By the rigidity lemma (\textsc{Mumford} \cite[6.1. p.~115]{sga4.exp.XVIII.bib.4}; note
that \textsc{Mumford} does not use the Noetherian hypothesis imposed on
$S$), it suffices to treat the case where $S$ is the spectrum of an
algebraically closed field. Suppose first that $G$ is an abelian
scheme, and let $G^{\ast}$ be the dual abelian scheme. Since
$G=G^{{\ast}{\ast}}$, giving $\varphi:X\rightarrow G$ is equivalent to
giving an invertible sheaf $\Ll$ on $X\times G^{\ast}$, trivialized
along $X\times \{e\}$ and algebraically equivalent to zero on the
fibers of the projection $\pr_{1}$ from $X\times G^{\ast}$ to $X$. If
$x:S\to X$ is a rational point of $X$, then
$\Ll\sim \pr^{\ast}_{2}(x\times G^\ast)^{\ast}\Ll\otimes \Mm$, with
$\Mm$ trivialized along $x\times G^{\ast}$ and $X\times \{e\}$. Giving
$\Mm$ is equivalent to giving a morphism of schemes from $G^{\ast}$ to
$\sheafPic_{X/S}$ that takes $e$ to $e$. Since
$H^{1}(X,\Oo_{X})=0$, one has $\sheafPic^0_{X/S}=0$\nde{This follows
from the canonical isomorphism
$\Lie(\sheafPic_{X/S})=H^{1}(X,\Oo_{X})=0$, obtained by comparing both
groups with the kernel of $\Pic(X\times_k
\Spec (k[\epsilon]))\to \Pic(X)$ (where $S=\Spec(k)$).
For further details on the properties of Picard schemes, see
\cite{sga4.exp.XVIII.bib.2bis}.},
$\Mm\sim\Oo$, and $\Ll\sim \pr^{\ast}_{2}\Ll'$, with $\Ll'$
\sisi{equivalent}{algebraically equivalent} to zero on $G^{\ast}$,
hence defining $g\in G(S)$ such that $\varphi=g\circ f$.

In the general case, $G$ is an extension of an abelian scheme $A$ by
an affine group $G_0$. The preceding result reduces us to the case
$A=0$, and then
$$
\Hom_{S}(\Spec (H^0(X,\Oo_{X}),G)\xrightarrow{\sim}
\Hom_{S}(X,G){\quoi}.
$$

\begin{enonce*}{Lemma 1.2.4}\etiquette[1.2.4]{XVIII.lem1.2.4}
Under the hypotheses of \hyperref[XVIII.subsubsec1.2.1]{1.2.1}, if $G$ is a group scheme
of finite presentation over $S$ or is defined by a quasi-coherent sheaf
on $S$, then the functor \eqhyperref[XVIII.eq1.2.1.1]{1.2.1.1} is fully faithful.
\end{enonce*}

\begin{itemize}
\item[A.] We prove that, if $K(K_{0},\varphi)$ is isomorphic to
$K(K'_0,\varphi')$, then $\varphi=\varphi'$. For quasi-coherent $G$,
one has $\varphi=\varphi'=0$, and there is nothing to prove. For $G$ of
finite presentation, by (\SGA 3 IX 5.1), it suffices\pageoriginale to
prove the assertion when $S$ is the spectrum of an algebraically
closed field $k$. Let $T$ be the largest subtorus of $G$, and put
$H=G/T$. The cohomology exact sequence gives
$$
0\rightarrow H^{0}(\PP(\epsilon),T)\rightarrow
H^{0}(\PP(\epsilon),G) \xrightarrow{\alpha}
H^{0}(\PP(\epsilon),H)
 \xrightarrow{\delta}
H^{1}(\PP(\epsilon),T)\xrightarrow{\beta}H^{1}(\PP(\epsilon),G){\quoi}.
$$
By \hyperref[XVIII.lem1.2.3]{1.2.3}, $H^0(\PP(\epsilon),G)\sim G(k)$ and
$H^0(\PP(\epsilon),H)\sim H(k)$, so $\alpha$ is surjective and
$\delta$ is zero. If $Y(T)=\Hom(\GG_{m},T)$, then
$T\sim G_{m}\otimes Y(T)$ and $H^{1}(\PP(\epsilon),T)\sim Y(T)$.
Thus the injectivity of $\beta$ expresses the injectivity of the
canonical arrow defined by $\Oo(1)$,
$\Hom(\GG_{m},G)\simeq \Hom(\GG_{m},T)\rightarrow
H^{1}(\PP(\epsilon),G)$, proving the assertion.

\item[B.] Proving that the functor \eqhyperref[XVIII.eq1.2.1.1]{1.2.1.1} is fully faithful is
local on $S$. We may therefore suppose that $\PP(\epsilon)/S$ has a
section $s$\sisi{}{ and choose an isomorphism
$s^\ast\Oo(1)\sim \Oo$\nde{This additional necessary datum, which was
missing from the original edition, has been added in this edition.}}.
If $K(K_{0},\varphi)$ is isomorphic to $K(K'_{0},\varphi')$, then
$\varphi=\varphi'$ by A, and
$K_{0}\sim s^{\ast}K(K_{0},\varphi)\sim
s^{\ast}K(K'_{0},\varphi')\sim K'_{0}$. By the formula
\begin{equation*}
\Hom (K(K_{0},\varphi), K(K'_{0},\varphi'))\sim \Hom(K(K_{0}-K'_{0},\varphi-\varphi'),G_{\PP(\epsilon)}){\quoi},\tag{1.2.4.1}
\end{equation*}\etiquette[1.2.4.1]{XVIII.eq1.2.4.1}%
it remains only to show that
$$
\Hom_{\tors.}(G,G)\xrightarrow{\sim}
\Hom_{\tors.}(p^{\ast}G,p^{\ast}G){\quoi}.
$$
This follows from \hyperref[XVIII.lem1.2.3]{1.2.3}.
\end{itemize}

\setcounter{subsubsection}{4}
\subsubsection{}\etiquette[1.2.5]{XVIII.subsubsec1.2.5}

Here is another, more \sisi{Kunneth-like}{Künneth-like}, proof of full
faithfulness. The question is local on $S$; by \eqhyperref[XVIII.eq1.2.4.1]{1.2.4.1}, it therefore
suffices to check that
$$
\Hom_{\PP(\epsilon)}(G,K(G,\varphi))
\begin{cases}
=∅ &\text{\rm if}\quad \varphi\neq 0\\
\xleftarrow{\sim} G(S)\ &\text{\rm if}\quad \varphi= 0{\quoi}.
\end{cases}
$$
This is precisely what the following lemma says\nde{This translation
merits some explanation. The line bundle $\Oo(1)$ on
$\PP(\epsilon)$ becomes canonically trivialized after pullback by the
canonical projection $\VV(\epsilon)^{\ast}\to \PP(\epsilon)$. After
pullback to $\VV(\epsilon)^{\ast}$, the $G$-torsor $K(G,\varphi)$ is
therefore canonically trivialized as well. \emph{Via} this
identification, an isomorphism $G\to K(G,\varphi)$ thus induces an
automorphism of the trivial $G$-torsor over
$\VV(\epsilon)^{\ast}$, that is, an $S$-morphism
$\VV(\epsilon)^{\ast}\to G$. To descend from $\PP(\epsilon)$, this
morphism must satisfy a compatibility like the one stated, and the
converse follows by faithfully flat descent. Strictly speaking, it
seems to the reviewer that the correct condition should be
$\Psi(\lambda v)=\varphi(\lambda)^{-1}\Psi(v)$, but this has no effect,
since one can apply the lemma to $\Psi^{-1}$ instead of $\Psi$. The
condition $\Psi(\lambda v)=\varphi(\lambda)\Psi(v)$ in the original
edition also gives a correct statement (since $\varphi$ can be replaced
by $\varphi^{-1}$), but it has been changed so that the proof need not
be modified.}:

\begin{enonce*}{Lemma 1.2.6}\etiquette[1.2.6]{XVIII.lem1.2.6}
Let $\epsilon$ be a locally free module everywhere of rank $\geq 2$ on
$S$, let $G$ be a group scheme of finite presentation over $S$, let
$\varphi:\GG_{m}\rightarrow G$ be a group homomorphism, let
$\VV(\epsilon)$ be the vector bundle defined by $\epsilon$, let
$\VV(\epsilon)^{\ast}$\pageoriginale be the corresponding punctured
vector bundle, and let $\Psi:\VV(\epsilon)^{\ast}\rightarrow G$ be an
$S$-morphism satisfying the following identity between $S$-morphisms
from $\GG_{m}\times \VV(\epsilon)^{\ast}$ to $G$:
$$
\Psi(\lambda v)=\sisi{\varphi (\lambda)\Psi(v)}{\Psi(v)\varphi(\lambda)}{\quoi}.
$$
Then $\varphi=e$, and $\Psi$ factors through the projection $p$ from
$\PP(\epsilon)$ to $S$.
\end{enonce*}

It suffices to prove the second assertion. We reduce at once to the
case where $S$ is Noetherian and then, by (\SGA 3 IX 5.1) and the
rigidity lemma (\textsc{Mumford} \cite[6.1 p.~115]{sga4.exp.XVIII.bib.6}) applied to
$\PP(\epsilon)$, to the case where $S$ is the spectrum of a field. Let
$n$ be an integer, let $W_{n}$ and $G_{n}$ be the $n^{\text{\rm th}}$
infinitesimal neighborhoods of the zero section in $\VV(\epsilon)$ and
$G$, let $\mathfrak{g}_{n}$ be the affine algebra of $G_{n}$, and let
$P_{n}(\Psi):\VV(\epsilon)^{\ast}\rightarrow
\Hhom_{S}(W_{n},G)$ be the principal part of $n^{\text{\rm th}}$ order of
$\Psi$\nde{More concretely, $P_n(\Psi)$ is the morphism corresponding
by adjunction to the morphism $\VV(\epsilon)^{\ast}\times_S W_n\to G$
(which we also denote by $P_n$) obtained by composing $\Psi$ with the
morphism $\VV(\epsilon)^{\ast}\times_{S} W_n\to
\VV(\epsilon)^\ast$ induced on these subschemes by addition on
$\VV(\epsilon)$. In terms of the functors of points of these
$S$-schemes, this can be interpreted more concretely by saying that,
if $v$ is a section of $\VV(\epsilon)^{\ast}$ and $t$ a section of
$W_n$, then $P_n(\Psi)(v,t)=\Psi(v+t)$.}.
Every section of $\Hhom_{S}(W_{n},G)$ defines, by right translation by
the inverse of the image of the zero section, a morphism from $W_{n}$
to $G_{n}$ that takes $0$ to $e$; this defines
$$
\Psi_{1}=P_{n}(\Psi).\Psi^{-1}:\VV(\epsilon)^{\ast}\longrightarrow
\Hhom_{S}(W_{n},G_{n}){\quoi}.
$$
The hypothesis implies that $\Psi_{1}$ commutes with homotheties, which
act on the right-hand side by transport of structure from their action
on $W_{n}\subset \VV(\epsilon)$\nde{With the notation above,
$\Psi_1$ can be characterized by the relation
$\Psi_1(v,t).\Psi(v)=\Psi(v+t)$. It follows at once that, if $\lambda$
is a section of $\GG_{m}$, then
$\Psi_1(\lambda v,\lambda t)=\Psi_1(v,t)$, which is the asserted
commutation with homotheties provided that $\GG_{m}$ acts correctly on
$\Hhom_S(W_n,G_n)$, namely, that a section $\lambda$ of $\GG_{m}$ acts
by right composition with multiplication by the \emph{inverse} of
$\lambda$ on $W_n$.}. The affine algebra of $W_{n}$ is
$\sum^{n}_{i=0}\Sym^{i}(\epsilon)$. By homogeneity, $\Psi_{1}$ is
therefore defined by a section, on $\PP(\epsilon)$, of the
$\Oo_{\PP(\epsilon)}$-module
$$
\sum^{n}_{i=0}p^{\ast}\Hhom(\mathfrak{g}_{n},\Sym^{i}(\epsilon))(-i){\quoi}.\nde{To
verify this assertion in detail, one may begin by defining the pullback
of this section to $\VV(\epsilon)^{\ast}$; one then obtains
tautologically a section $s=(s_0,\dots,s_n)$ of the trivial bundle
obtained by tensoring the structure sheaf with the vector space
$\sum^{n}_{i=0}\Hhom(\mathfrak{g}_{n},\Sym^{i}(\epsilon))$. It satisfies
the relation $s_i(v)=\lambda^i s_i(\lambda v)$. This is indeed descent
data giving the desired section on $\PP(\epsilon)$. One can also argue
that the sections $s_i$ extend to all of $\VV(\epsilon)$ (since the
origin is a closed subset of codimension $\geq 2$ in
$\VV(\epsilon)$). By the density of $\GG_m$ in $\AAA^1$, the extended
sections $s_i$ satisfy the relation
$s_i(v)=\lambda^i s_i(\lambda v)$ for $\lambda\in\AAA^1$. Taking
$\lambda=0$ gives $s_i=0$ for $i>0$.}
$$
The components of index $i>0$ of this module have no section other
than the zero section, and it follows that
$$
\Psi_{1}=0.
$$
Since this holds for every $n$, it implies that $\Psi$ factors through
a section of $G$.

This second proof extends to prove the following analytic variant of
\hyperref[XVIII.lem1.2.4]{1.2.4}. When $G$ is affine, the proof of
\hyperref[XVIII.thm1.2.2]{1.2.2} can be given independently of this variant.

\setcounter{subsubsection}{6}
\subsubsection{}\etiquette[1.2.7]{XVIII.subsubsec1.2.7}

Let\pageoriginale $L$ be the fraction field of a complete discrete
valuation ring $\Oo$, let $U_{m}$ be the rigid analytic ``generic
fiber'' group of the formal completion of $\GG_{m}$, and let
$\Oo(1)^{\wedge}$ be the $U_{m}$-torsor on $\PP^{r}_{L}$ obtained by
formally completing $\Oo(1)$. If $G$ is a rigid analytic group over
$L$, the functor $(K_{0},\varphi)\mapsto
p^{\ast}K_{0}+\varphi\Oo(1)^{\wedge}$ (\cf \hyperref[XVIII.subsubsec1.2.1]{1.2.1}) is a fully
faithful functor from the category of pairs $(K_{0},\varphi)$
consisting of a $G$-torsor on $\Spec(L)$ and
$\varphi\in \Hom(U_m,G)$ to the category of $p^{\ast}G$-torsors on
$\PP^{r}_{L}$\acorriger{The reader may find an introduction to the
language of rigid analytic geometry in \cite{sga4.exp.XVIII.bib.5bis}. The nonexistence
of nontrivial sections of the sheaves $\Oo(n)$ for $n<0$ in this
context follows, for example, from a rigid version of GAGA; see
\cite{sga4.exp.XVIII.bib.3ter}.}.

\begin{enonce*}{Lemma 1.2.8}\etiquette[1.2.8]{XVIII.lem1.2.8}
Let $k$ be a field of characteristic exponent $p$, and let
$f:X\rightarrow \Spec(k)$ be a proper scheme over $k$ satisfying
$H^0(X,\Oo_{X})=k$ and $\sheafPic^{\tau}_{X}=0$\nde{A definition of
$\sheafPic^{\tau}$ in a more general setting will be found in
\hyperref[XVIII.def1.6.4]{1.6.4}.}. For every finite-type commutative
affine group scheme $G$ over $k$ containing no tori, one has
$R^{1}f_{\ast}(f^{\ast}G)=0$.
\end{enonce*}

% END INCLUDED SGA4_English_current_buildable_source_20260906/source/18/drafts/SGA4_XVIII_LINES_451_855_English.texfrag

% BEGIN INCLUDED SGA4_English_current_buildable_source_20260906/source/18/drafts/SGA4_XVIII_LINES_856_1250_English.texfrag

It suffices to show that after every base change $c:S\rightarrow
\Spec(k)$, the functor $f^{\ast}$ from the category of $G$-torsors on
$S$ to that of $G$-torsors on $X_{S}$ is an equivalence. This problem
is local on $\Spec (k)$, so that, in order to solve it, one may reduce
to the case where $X$ has a rational point, where $G$ is a successive
extension of the groups $\GG_{a}$, $\alpha_{p}$, $𝛍_{p}$,
$\ZZ/p$, and $\ZZ/\ell_{i}$, with $\ell_{i}$ prime to $p$,
and where $\ZZ/\ell_{i}\sim 𝛍_{\ell_i}$.

We know that $\Lie(\sheafPic_{X})=H^{1}(X,\Oo_{X})=0$. By
\hyperref[XVIII.lem1.2.3]{1.2.3}
and the existence of a section, the functors $f^{\ast}$ under
consideration are fully faithful; it therefore suffices to prove that
$R^{1}f_{\ast}(f^{\ast}G)=0$, or, by dévissage, that
$$
R^{1}f_{\ast}\GG_{a}=R^{1}f_{\ast}\alpha_{p}=R^{1}f_{\ast}𝛍_{p}=R^{1}f_{\ast}\ZZ/p=R^{1}f_{\ast}\ZZ/\ell_{i}=0{\quoi}.
$$

For $\GG_{a}$, this follows from $H^{1}(X,\Oo_{X})=0$. The case of
$\alpha_{p}$ (\resp $\ZZ/p$) follows from it via
\hyperref[XVIII.lem1.2.3]{1.2.3}, by the long exact cohomology sequence
defined by the exact sequence
\begin{align*}
& 0\longrightarrow \alpha_{p}\longrightarrow
  \GG_{a}\xrightarrow{x^{p}} \GG_{a}\longrightarrow 0\\
& 0\longrightarrow \ZZ/p\longrightarrow
  \GG_{a}\xrightarrow{x-x^{p}}\GG_{a}\longrightarrow
  0\qquad (p\neq 1)){\quoi}.\tag*{(\resp}
\end{align*}
The case of $\ZZ/\ell_{i}$ (\resp of $𝛍_{p}$) follows from the
isomorphism $\ZZ/\ell_{i}\sim 𝛍_{\ell_{i}}$,
from\pageoriginale \hyperref[XVIII.lem1.2.3]{1.2.3}, and from the long
exact sequence deduced from the Kummer exact sequence
$$
0\longrightarrow 𝛍_{n}\longrightarrow \GG_{m}\longrightarrow
\GG_{m}\longrightarrow 0{\quoi}.
$$

\begin{enonce*}{Lemma 1.2.9}\etiquette[1.2.9]{XVIII.lem1.2.9}
Let $S$ be a local Artinian scheme, let $\bar{s}$ be a geometric point
of $S$, let $f:X\rightarrow S$ be a proper flat morphism, and let $G$
be a flat commutative group scheme of finite type over $S$. Assume
that $f$ is smooth or that $G$ is affine. If
$H^0(X_{\bar{s}},\Oo)=k(\bar{s})$, if
$\sheafPic^{\tau}_{X_{\bar{s}}}=0$, and if $G_{\bar{s}}$ contains no
subgroup $\GG_{m}$, then the functor $f^{\ast}$ from the category of
$G$-torsors on $S$ to that of $f^{\ast}G$-torsors on $X$ is an
equivalence.
\end{enonce*}

\begin{itemize}
\item[a.] The problem in question is local on $S$ when the finite flat
  morphisms $u:S'\rightarrow S$ are taken as localization morphisms.
  This makes it possible to assume that $f$ admits a section
$x$\nde{The existence of a section after such a base change may be
deduced from \hyperref[XVIII.lem2.3]{2.3} and
\hyperref[XVIII.lem2.4]{2.4}.}. The functor $f^{\ast}$ is then fully
  faithful, by \hyperref[XVIII.lem1.2.3]{1.2.3} and the existence of
  $x$, and it remains to show that every torsor trivialized along $x$
  is trivial.

\item[b.] Suppose that $S$ is the spectrum of a field. If $G$ is
  affine, the assertion follows from
  \hyperref[XVIII.lem1.2.8]{1.2.8}. If $X$ is smooth and $G$
  is an extension of an abelian variety $A$ by an affine group
  $G_{0}$, denote by $G^{n}$ the inverse image of ${}_{n}A$ in $G$.
  The group $H^{1}(X,A)$ is torsion (\textsc{Raynaud}
  \cite[XIII 2.6]{sga4.exp.XVIII.bib.5}), and is therefore the union of the images of the
  $H^{1}(X,{}_{n}A)$. The group $H^{1}(X,G)$ is consequently the union
  of the images of the $H^{1}(X,G^{n})$. Since the groups $G^{n}$ are
  affine, the conclusion again follows from
  \hyperref[XVIII.lem1.2.8]{1.2.8}.

\item[c.] Let us prove the general case (assuming the existence of a
  section $x$) by induction on the length of $S$. Let
  $i:S_{0}\hookrightarrow S$ be a closed subscheme defined by an ideal
  $I$ of square zero and length 1.
\end{itemize}

By the induction hypothesis, the $G$-torsors $P$ on $X$ for which
$x^{\ast}P$ is trivial are deformations of the trivial $G$-torsor
$P_{0}$ on $X\times S_{0}$. If $G$ is smooth, and if $L$ is the $\Lie$
algebra of $G_{s}$, these deformations are classified by
$$
H^{1}(X_{s},f^{\ast}(L\otimes I))=0{\quoi},
$$
which\pageoriginale proves \hyperref[XVIII.lem1.2.9]{1.2.9} in this
case.

For $G$ flat over $S$, this argument generalizes in terms of the
\emph{Lie complex} $L$ of $G_{s}$.
% SGA-SRC-001128 — reconstructed by mathematical inference; scan inconclusive
If $\mathrm{L}\Lambda(G_{s}/k(s))$ denotes the relative cotangent complex
of $G_{s}$, this complex $L\in D^{b}(k(s))$ is the dual of
$\mathrm{L}e^{\ast}(\mathrm{L}\Lambda(G_{s}/k(s)))$. We have
$H^{i}(L)=0$ for $i\neq 0$, $1$\sisi{}{(\textsc{Illusie}
\cite[VII~2.4.1]{sga4.exp.XVIII.bib.3})}.

The deformations over $S$ (\resp over $X$) of the trivial $G$-torsor
over $S_{0}$ (\resp over $X_{0}$) are classified by $H^{1}(L\otimes I)$
(\resp by $H^{1}(X_{s},f^{\ast}(L\otimes I))$ (\textsc{Illusie}
\cite[VII~2.4.4~(ii)]{sga4.exp.XVIII.bib.3}).
By hypothesis, we have $H^{1}(X_{s},ℋ^0(f^{\ast}(L\otimes
I)))=0$ (because $H^{1}(X_{s},\Oo)=0$) and $H^{1}(L\otimes I)\sim
H^0(X_{s}, ℋ^{1}(f^{\ast}L\otimes I))$). We conclude that
$$
H^{1}(L\otimes I)\xrightarrow{\sim} H^{1}(X_{s},f^{\ast}(L\otimes
I)){\quoi},
$$
which means that every deformation over $X$ of the trivial $G$-torsor
over $X_{0}$ is the inverse image of one and only one deformation over
$S$ of the trivial $G$-torsor over $S_{0}$. If this deformation is
trivial along a section, it is therefore trivial, which proves
\hyperref[XVIII.lem1.2.9]{1.2.9}.

\setcounter{subsubsection}{9}
\subsubsection{}\etiquette[1.2.10]{XVIII.subsubsec1.2.10}
Let us prove \hyperref[XVIII.thm1.2.2]{1.2.2} when $S$ is Artinian. Let
$T$ be the largest subtorus of $G$; since full faithfulness is already
known (\hyperref[XVIII.lem1.2.4]{1.2.4}), the problem is local on $S$
for the étale topology: we may assume that $T$ is diagonalizable,
$T\sim \GG^{n}_{m}$. By \hyperref[XVIII.lem1.2.3]{1.2.3} and
\hyperref[XVIII.lem1.2.9]{1.2.9}, the exact cohomology sequence gives
$$
R^{1}f_{\ast}T\xrightarrow{\sim} R^{1}f_{\ast}G{\quoi},
$$
and the assertion follows, since it is trivial for $G=T$ and
$\Hom(\GG_{m},T)\xrightarrow{\sim}\Hom(\GG_{m},G)$ (\cf
\hyperref[XVIII.par1.2.2.1]{1.2.2.1}).

\subsubsection{}\etiquette[1.2.11]{XVIII.subsubsec1.2.11}
Let us prove \hyperref[XVIII.thm1.2.2]{1.2.2} when $S$ is complete
Noetherian local. Let $s$ be a section of~$p$\nde{As above, one must
also choose an isomorphism $s^\ast\Oo(1)\sim \Oo$.}.

\noindent
If $K$ is a $G$-torsor on $\PP(\epsilon)$, by
\hyperref[XVIII.subsubsec1.2.10]{1.2.10} there exists a unique pair
$(\widehat{\varphi},\widehat{\Psi})$ consisting of a formal
homomorphism $\widehat{\varphi}:\GG_{m}\rightarrow G$ and a formal
isomorphism\pageoriginale
$\widehat{\Psi}:\widehat{p^{\ast}s^{\ast}K}+
\widehat{\varphi}\widehat{\Oo(1)}\xrightarrow{\sim}
\widehat{K}$\nde{One must require this isomorphism to become the
obvious isomorphism after applying $s^\ast$.}\nde{What is meant here is
that, if $A$ is the ring of $S$, $I$ its maximal ideal, and
$S_n=\Spec(A/I^n)$, a formal morphism consists of a compatible system
of morphisms between objects over $S_n$ for every $n$. The question of
algebraicity is whether this compatible system is obtained by inverse
image from similar data over $S$.}. If $G$ is affine, then
$\widehat{\varphi}$ is algebraizable (\SGA 3 IX 7.1). Let us
temporarily admit that, as will be proved in
(\hyperref[XVIII.subsubsec1.2.13]{1.2.13}), $\widehat{\varphi}$ is
always algebraizable as $\varphi:\GG_{m}\rightarrow G$.

If $K_{0}=s^{\ast}K$, and if
$K_{1}=p^{\ast}K_{0}+\varphi\Oo(1)$, we have a formal isomorphism
$\widehat{\Psi}:\widehat{K_{1}}=\widehat{K(K_{0},\varphi)}\xrightarrow{\sim}
\widehat{K}$, \ie a formal trivialization of $K-K_{1}$. By formal
GAGA, this section is algebraic, which proves
\hyperref[XVIII.subsubsec1.2.11]{1.2.11}. More precisely, if $G$ is
quasi-affine, then $K-K_{1}$ is representable and one uses \EGA III
5.1.4. In the general case, one proceeds in the same way with the
algebraic space $K-K_{1}$.

\subsubsection{Algebraicity of $\widehat{\varphi}$: preliminary
  verification over a trait}\etiquette[1.2.12]{XVIII.subsubsec1.2.12}

Suppose that $S$ is a complete trait, $S=\Spec(V)$, and let $L$ be the
fraction field of~$V$.
\begin{itemize}
\item[a)] Over $L$, there exists a unique triple
  $(K_{0L},\varphi_{L},\Psi_{L})$:
$$
\Psi_{L}:K(K_{0L},\varphi_{L})\xrightarrow{\sim} K|\Spec (L){\quoi}.
$$

\item[b)] By (\hyperref[XVIII.subsubsec1.2.10]{1.2.10})
  (\hyperref[XVIII.subsubsec1.2.11]{1.2.11}), and with the notation of
  \hyperref[XVIII.subsubsec1.2.7]{1.2.7}, we have a rigid analytic map
  $\widehat{\varphi_{L}}:U_{m}\rightarrow \widehat{G}$, a
  $\widehat{G}$-torsor $\widehat{K_{0}}$, and a rigid analytic
  isomorphism
$$
\widehat{\Psi}:p^{\ast}\widehat{K_0}+\widehat{\varphi_{L}}
\widehat{\Oo(1)}\xrightarrow{\sim}
\widehat{K}{\quoi}.
$$
By the uniqueness assertion \hyperref[XVIII.subsubsec1.2.7]{1.2.7},
$\widehat{\varphi_{L}}$ is identified with the map deduced from
$\varphi_{L}$. Thus the formal morphism $\widehat{\Psi}$ of
\hyperref[XVIII.subsubsec1.2.11]{1.2.11} induces on $𝛍_{n}$ a map that
coincides with the map deduced from the algebraic map
$\varphi_{L}:\GG_{m}\rightarrow G$.
\end{itemize}

\subsubsection{Algebraicity of $\widehat{\varphi}$}\etiquette[1.2.13]{XVIII.subsubsec1.2.13}

The formal map $\widehat{\varphi}:\GG_{m}\rightarrow G$ of
\hyperref[XVIII.subsubsec1.2.11]{1.2.11} induces, for $\ell$ invertible
on $S$, maps
$$
\varphi_{\ell}:T_{\ell}(\widehat{\GG_{m}})=T_{\ell}(\GG_{m})\longrightarrow
T_{\ell}(G){\quoi}.
$$
By \hyperref[XVIII.subsubsec1.2.12]{(1.2.12)} and (\EGA II 7.1.9), for every point $t$ of $S$, there
exists an\pageoriginale algebraic morphism
$\varphi_{t}:\GG_{m_{t}}\rightarrow G_{t}$ such that
$T_{\ell}(\varphi_{t})=T_{\ell}(\widehat{\varphi})_{t}$. By (\SGA 3 XV 4.1
  bis), $\widehat{\varphi}$ is therefore algebraizable.

\subsubsection{}\etiquette[1.2.14]{XVIII.subsubsec1.2.14}
Let us prove \hyperref[XVIII.thm1.2.2]{1.2.2}. By passage to the limit,
we reduce to the case where $S$ is Noetherian, and then where $S$ is
Noetherian local with affine coordinate ring $A$. Let $s$ be a section
of $p$. Let $K$ be a $G$-torsor on $\PP(\epsilon)$. Over
$S'=\Spec(\widehat{A})$, there exists one and only one homomorphism
$\varphi':\GG_{mS'}\rightarrow G_{S'}$ such that the inverse image
$K'$ of $K$ on $\fprod{\PP(\epsilon)}{S'}{S}$ is of the form
$K(K'_{0},\varphi')$. This homomorphism $\varphi'$ descends to
$\varphi:\GG_{m}\rightarrow G$. To prove the essential surjectivity
of the functor \hyperref[XVIII.subsubsec1.2.1]{1.2.1}, we may replace
$K$ by $K-\varphi\Oo(1)$, \ie assume that $\varphi=0$. Over $S'$,
there then exists one and only one isomorphism
$\Psi:{p'}^{\ast}s^{\ast}K'\xrightarrow{\sim}K'$, inducing the
identity on the section $s$. This morphism descends to $S$ (fpqc
descent of morphisms of fppf torsors), and this proves
\hyperref[XVIII.thm1.2.2]{1.2.2}.

\begin{enonce*}{Variants 1.2.15}\etiquette[1.2.15]{XVIII.var1.2.15}
The conclusion of \hyperref[XVIII.thm1.2.2]{1.2.2} still holds in the
following cases.
\begin{enumerate}
\renewcommand{\theenumi}{\roman{enumi}}
\renewcommand{\labelenumi}{(\theenumi)}
\item $G$ is defined by a quasi-coherent module on $S$.

\item $G$ is the inverse image of a torsion sheaf on the small
  étale site of $S$.
\end{enumerate}
\end{enonce*}

In case (i), it is standard that the arrows
(\hyperref[XVIII.par1.2.2.1]{1.2.2.1}) are isomorphisms, taking into
account the fact that the fppf cohomology of $G$ coincides with its
cohomology for the Zariski topology.

In case (ii), the arrows (\hyperref[XVIII.par1.2.2.1]{1.2.2.1}) are
isomorphisms by the base change theorem for a proper morphism and the
simple connectedness of $\PP^{r}_{k}$ for $k$ algebraically closed.

\subsection{Integration of torsors}\etiquette[1.3]{XVIII.subsec1.3}

\subsubsection{}\etiquette[1.3.1]{XVIII.subsubsec1.3.1}
Let $S$ be a scheme and $G$ an abelian sheaf on the big fppf site of
$S$, satisfying the following three conditions.

\paragraph{}\etiquette[1.3.1.1]{XVIII.par1.3.1.1}
\begin{itemize}
\item[(i)] \emph{For\pageoriginale every projective space
$p:\PP^r_{S}\rightarrow S$, with $r\geq 1$, one has, on the big}
fppf \emph{site of} $S$ (\hyperrefext[XVII][XVII.subsec0.10]{0.10})
\begin{gather*}
G\xrightarrow{\sim} p_{\ast}p^{\ast}G\quad\mbox{and}\\
\Hhom(\GG_{m},G)\xrightarrow{\sim} R^{1}p_{\ast}p^{\ast}G
\end{gather*}

\item[(ii)] \emph{For every $n$ and every smooth affine $S$-morphism
  $f:X\rightarrow T$, if $K_{1}$ and $K_{2}$ are two $G$-torsors on
  $\Sym^{n}_{T}(X)$, the set $\Hom(K_{1},K_{2})$ is identified with
  the set of symmetric morphisms from the inverse image of $K_{1}$ on
  $(X/T)^{n}$ to the inverse image of $K_{2}$ on $(X/T)^{n}$.}

This condition is equivalent to the conjunction of:

\item[${\rm (ii')}$]
$\Hom_{S}(\Sym^{n}_{T}(X),G)\xrightarrow{\sim}\Hom_{S}((X/T)^{n},
G)^{\sigma_n}$, and

\item[${\rm (ii'')}$] If the inverse image on $(X/T)^{n}$ of a
torsor $K$ on $\Sym^{n}_{T}(X)$ admits a symmetric section, then $K$
is trivial.

\item[(iii)] $G$ \emph{satisfies either condition}
(\hyperrefext[XVII][XVII.par6.3.21.1]{6.3.21.1}), \emph{or conditions}
(\ref{XVII.par6.3.24.1}) \emph{and} (\ref{XVII.par6.3.24.2}).
\end{itemize}

The conclusion of \hyperref[XVIII.thm1.2.2]{1.2.2} therefore holds
(\cf \hyperref[XVIII.par1.2.2.1]{1.2.2.1}), as does the formula of
\hyperrefext[XVII][XVII.subsubsec6.3.26]{6.3.26}.

Condition (i) is discussed in \hyperref[XVIII.thm1.2.2]{1.2.2} and
\hyperref[XVIII.var1.2.15]{1.2.15}. Condition (ii) holds whenever $G$
is representable and of finite presentation, or is the inverse image
of a sheaf on the small étale site of $S$, or is flat
quasi-coherent. Condition \hyperref[XVIII.par1.3.1.1]{1.3.1.1}
therefore holds in each of the following cases.

\paragraph{}\etiquette[1.3.1.2]{XVIII.par1.3.1.2}
\begin{enumerate}
\renewcommand{\theenumi}{\alph{enumi}}
\renewcommand{\labelenumi}{(\theenumi)}
\item $G$ is smooth and of finite presentation
  (\hyperref[XVIII.thm1.2.2]{1.2.2} and
  \hyperrefext[XVII][XVII.par6.3.21.1]{6.3.21.1}).

\item $G$ is the inverse image of a torsion sheaf on the small
  étale site of $S$ (\hyperref[XVIII.var1.2.15]{1.2.15} (ii) and
  \hyperrefext[XVII][XVII.par6.3.21.1]{6.3.21.1}).

\item $G$ is affine, and the kernel of an epimorphism of smooth groups
  (\hyperref[XVIII.thm1.2.2]{1.2.2} and XVII \hyperref[XVII.eq6.3.3.1]{6.3.3.1}).

\item $G$ is defined by a flat quasi-coherent sheaf on $S$
  (\hyperref[XVIII.var1.2.15]{1.2.15} (i) and
  \hyperrefext[XVII][XVII.par6.3.21.1]{6.3.21.1}).
\end{enumerate}

\subsubsection{}\etiquette[1.3.2]{XVIII.subsubsec1.3.2}
Let\pageoriginale $f:X\rightarrow S$ be a smooth projective curve over
a scheme $S$, with nonempty geometrically connected fibers. Let $\Ll$
be an invertible sheaf on $X$ that is ``sufficiently'' relatively
ample in the following sense.

\paragraph{}\etiquette[1.3.2.1]{XVIII.par1.3.2.1}
For every geometric point $s$ of $S$, if $X_{s}$ has genus $g$, then
\begin{gather*}
\deg_{X_{s}}(\Ll)>2g-2,\\
\deg_{X_{s}}(\Ll)\geq 1\quad\mbox{\rm if}\quad g=0,\quaet
\deg_{X_{s}}(\Ll)\geq 2\quad\mbox{\rm if}\quad g=1
\end{gather*}

We then have $R^{1}f_{\ast}\Ll=0$, and $f_{\ast}\Ll$ is locally free,
its formation commutes with every base change, and it has rank
$\geq 2$ everywhere. If a section $s$ of $f_{\ast}\Ll$ vanishes at
no point, the scheme $D_{s}$ of zeros of $s$, regarded as a section of
$\Ll$ on $X$, is a relative divisor (\EGA IV 21.15.2), finite locally
free over $S$ since $X/S$ has relative dimension 1. We have canonically
$$
s:\Oo(D_{s})\xrightarrow{\sim}\Ll{\quoi}.
$$
The divisor $D_{s}$ depends on $s$ only up to multiplication by a
section of $\Oo^{\ast}_{S}$. The preceding construction therefore
gives us a relative divisor $Z_{\Ll}$ on the scheme obtained from $X$
by the base change
$$
p:\PP(f_{\ast}(\Ll)^{\vee})\longrightarrow S{\quoi}.
$$

Let $K$ be a torsor under $G_{X}$, and let $K|Z_{\Ll}$ be its inverse
image on $Z_{\Ll}$. The trace of this torsor
(\hyperrefext[XVII][XVII.subsubsec6.3.26]{6.3.26}), from $Z_{\Ll}$ to
$\PP(f_{\ast}(\Ll)^{\vee})$, is a torsor $K_{\Ll}$ on
$\PP(f_{\ast}(\Ll)^{\vee})$:
$$
K_{\Ll}=\Tr_{Z_{\Ll}/\PP}(K|Z_{\Ll}){\quoi}.
$$

By \hyperref[XVIII.thm1.2.2]{1.2.2}, there exists a triple, unique up
to unique isomorphism, consisting of a $G$-torsor
$\langle \Ll$, $K$] on $S$, a morphism
$\varphi_{\Ll,K}:\GG_{m}\rightarrow G$, and an isomorphism
$$
K_{\Ll}\simeq
p^{\ast}\langle \Ll,K]+\varphi_{\Ll,K}\Oo(1){\quoi}.
$$

\subsubsection{}\etiquette[1.3.3]{XVIII.subsubsec1.3.3}
If\pageoriginale the composition of torsors is written additively, we
have canonically
\begin{align*}
& \langle \Ll,K_{1}+K_{2}] \simeq
\langle \Ll,K_{1}]+\langle \Ll,K_{2}]\tag{1.3.3.1}\\
& \varphi_{\Ll, K_{1}+K_{2}}=\varphi_{\Ll,K_{1}}+\varphi_{\Ll,K_{2}}{\quoi}.\tag{1.3.3.2}
\end{align*}\etiquette[1.3.3.1]{XVIII.eq1.3.3.1}
\etiquette[1.3.3.2]{XVIII.eq1.3.3.2}

Let $\Ll_{1}$ and $\Ll_{2}$ be two invertible sheaves on $X$
satisfying (\hyperref[XVIII.par1.3.2.1]{1.3.2.1}), and let
$\Ll=\Ll_{1}\otimes \Ll_{2}$. The tensor product of sections defines
a morphism
$$
\xymatrix@R=1.5cm@C=.5cm{
 & **[r] \PP
  (f_{\ast}(\Ll_1)^{\vee}) \ar[ld]^{\pr_1}& \times_S \ar[dd]^{p_{12}}& **[l]\PP
  (f_{\ast}(\Ll_2)^{\vee})\ar[rr]^{\pi}\ar[rd]^{\pr_2} &     &
 \PP(f_{\ast}(\Ll)^{\vee}) \ar@{-}[dd]^{p}\\
**[r]\PP(f_*(\Ll_1)^{\vee})\ar[rrd]_{p_1}&          &                &                 &  **[l]\PP(f_{\ast}(\Ll_2)^{\vee}) \ar[lld]^{p_2}  &\\
 &          &  S     \ar@{=}[rrr]        &                 &      & S.
}
$$
We have $\pi^{\ast}Z_{\Ll}\simeq Z_{\Ll_{1}}+Z_{\Ll_{2}}$, whence a

% END INCLUDED SGA4_English_current_buildable_source_20260906/source/18/drafts/SGA4_XVIII_LINES_856_1250_English.texfrag

% BEGIN INCLUDED SGA4_English_current_buildable_source_20260906/source/18/drafts/SGA4_XVIII_LINES_1251_1650_English.texfrag
isomorphism (\hyperrefext[XVII][XVII.const6.3.27.1]{6.3.27.1})
\begin{equation*}
\pi^{\ast}K_{\Ll}=\pr^{\ast}_{1}K_{\Ll_{1}}+\pr^{\ast}_{2}K_{\Ll_{2}}{\quoi}.\tag{1.3.3.3}
\end{equation*}\etiquette[1.3.3.3]{XVIII.eq1.3.3.3}

There is a canonical isomorphism
$$
\pi^{\ast}\Oo(1)\simeq \pr^{\ast}_{1}\Oo(1)\otimes
\pr^{\ast}_{2}\Oo(1){\quoi},
$$
so that \eqhyperref[XVIII.eq1.3.3.3]{1.3.3.3} may be rewritten as
\begin{gather*}
\pi^{\ast}(p^{\ast}\langle \Ll,K]+\varphi_{\Ll,K}\Oo(1))=
p^{\ast}_{12}\langle \Ll,K]+\pr^{\ast}_{1}\varphi_{\Ll,K}\Oo(1)+\pr^{\ast}_{2}\varphi_{\Ll,K}\Oo(1) \\
= p^{\ast}_{12}\langle \Ll_{1},K]+p^{\ast}_{12}\langle \Ll_{2},K]+\pr^{\ast}_{1}\varphi_{\Ll_{1},K}\Oo(1)+\pr^{\ast}_{2}\varphi_{\Ll_{2},K}\Oo(1){\quoi}.
\end{gather*}

By applying \hyperref[XVIII.thm1.2.2]{1.2.2} twice, one obtains a canonical
isomorphism
\begin{gather*}
\langle \Ll_{1}\otimes \Ll_{2}, K]=\langle \Ll_{1},K]+\langle \Ll_{2},K]\tag{1.3.3.4}
\end{gather*}\etiquette[1.3.3.4]{XVIII.eq1.3.3.4}
and\pageoriginale an identity
$$
\varphi_{\Ll_{1}\otimes\Ll_{2},K}=\varphi_{\Ll_{1},K}=\varphi_{\Ll_{2},K}{\quoi}.
$$

Thus one sees that $\varphi_{\Ll,K}$ depends only on $K$, and one
sets
\begin{equation*}
\deg(K)=\varphi_{\Ll,K}{\quoi}.\tag{1.3.3.5}
\end{equation*}\etiquette[1.3.3.5]{XVIII.eq1.3.3.5}%
Consequently,
\begin{align*}
& K_{\Ll}\simeq
p^{\ast}\langle \Ll,K]+\deg(K)\Oo(1){\quoi}.\tag{1.3.3.6}\\
& \deg(K_{1}+K_{2})=\deg(K_{1})+\deg(K_{2}){\quoi}.\tag{1.3.3.7}
\end{align*}\etiquette[1.3.3.6]{XVIII.eq1.3.3.6}
\etiquette[1.3.3.7]{XVIII.eq1.3.3.7}


Let $\mathcal{U}$ be an invertible sheaf that is locally a direct
summand of $f_{\ast}\Ll$, defining a section $u$ of
$\PP((f_{\ast}\Ll)^{\vee})$. Let $D_{u}$ be the relative divisor
defined by the local invertible sections of $\mathcal{U}$.
$$
u^{\ast}\Oo(1)\simeq \mathcal{U}^{\vee}
$$
It therefore follows from \eqhyperref[XVIII.eq1.3.3.6]{1.3.3.6} that
\begin{equation*}
u^{\ast}K_{\Ll}=\Tr_{D_u/S}(K)\simeq
\langle \Ll,K]-\deg(K)\mathcal{U}{\quoi}\nde{The notation $\Tr_{D_u/S}(K)$
is introduced below in \eqhyperref[XVIII.eq1.3.5.0]{1.3.5.0}.}.\tag{1.3.3.8}
\end{equation*}\etiquette[1.3.3.8]{XVIII.eq1.3.3.8}

We leave it to the reader to verify the following compatibility:

\begin{enonce*}{Proposition 1.3.4}\etiquette[1.3.4]{XVIII.prop1.3.4}
The following diagram of the isomorphisms \eqhyperref[XVIII.eq1.3.3.1]{1.3.3.1} and \eqhyperref[XVIII.eq1.3.3.4]{1.3.3.4} is
commutative:
$$
\xymatrix@C=.1cm@R=.1cm{
\langle \Ll_{1}\otimes \Ll_{2}, K_{1}+K_{2}] & \simeq & \langle
\Ll_{1}\otimes \Ll_{2}, K_{1}]+\langle \Ll_{1} \otimes \Ll_{2}, K_{2}]\\
 & & \rotatebox{90}{$\simeq$} \\
\rotatebox{90}{$\simeq$} & & \langle \Ll_{1}, K_{1}]+\langle\Ll_{2}, K_{1}]+\langle
    \Ll_{1}, K_{2}]+\langle \Ll_{2}, K_{2}]\\
 & & \rotatebox{-50}{$\Arrowvert$}\\
\langle\Ll_{1}, K_{1}+K_{2}]+\langle\Ll_{2},K_{1}+K_{2}] & \simeq  &
    \langle\Ll_{1}, K_{1}]+\langle \Ll_{1},K_{2}]+\langle\Ll_{2},K_{1}]+\langle\Ll_{2},K_{2}]. 
}
$$%raghu

%$$
%\xymatrix{
%\langle \Ll_1\otimes \Ll_2, K_1 + K_2] \ar@{-}[dd]&\simeq& \langle \Ll_1\otimes \Ll_2, K_1] + \langle \Ll_! \otimes
%  \Ll_2, K_2]\ar@{-}[d]\\
%  & & \langle \Ll_1\otimes \Ll_2, K_1] + \langle \Ll_! \otimes
%  \Ll_2, K_2] \ar@{=}[d]\\
%\langle \Ll_1\otimes \Ll_2, K_1] + \langle \Ll_! \otimes
%  \Ll_2, K_2] & \simeq& \langle \Ll_1\otimes \Ll_2, K_1] + \langle \Ll_! \otimes
%  \Ll_2, K_2].
%}
%$$


\end{enonce*}

The isomorphisms \eqhyperref[XVIII.eq1.3.3.1]{1.3.3.1} and \eqhyperref[XVIII.eq1.3.3.4]{1.3.3.4} are moreover
compatible with the associativity and commutativity isomorphisms in a
sense whose precise formulation we leave to the reader. Formula
\eqhyperref[XVIII.eq1.3.3.4]{1.3.3.4} then makes it possible to extend by ``linearity'' the definition of
$\langle \Ll,K]$ when $\Ll$ no\pageoriginale longer necessarily
satisfies condition (\hyperref[XVIII.par1.3.2.1]{1.3.2.1}).

\setcounter{subsubsection}{4}
\subsubsection{}\etiquette[1.3.5]{XVIII.subsubsec1.3.5}
For every effective relative divisor $D$ on $X$, set
\begin{equation*}
\Tr_{D/S}(K)=\Tr_{D/S}(K|D){\quoi}.\tag{1.3.5.0}
\end{equation*}\etiquette[1.3.5.0]{XVIII.eq1.3.5.0}

If $\Oo(D)$ satisfies (\hyperref[XVIII.par1.3.2.1]{1.3.2.1}), the section
\hyperref[XVIII.sec1]{1} of $\Oo(D)$ is a
section of $f_{\ast}\Oo(D)$, and as such generates an invertible sheaf
that is locally a direct summand of $f_{\ast}\Oo(D)$, of which it
defines a trivialization. The isomorphism \eqhyperref[XVIII.eq1.3.3.8]{1.3.3.8} therefore
gives us an isomorphism
\begin{equation*}
\langle \Oo(D),K]\simeq \Tr_{D/S}(K){\quoi}.\tag{1.3.5.1}
\end{equation*}\etiquette[1.3.5.1]{XVIII.eq1.3.5.1}

One immediately verifies that the diagrams
$$\vcenter{
\xymatrix{
  \langle \Oo(D) , K_1 + K_2] \ar[r]^{\sim} \ar@{-}[d]_{\text{\rotatebox{90}{$\sim$}}}& \Tr_{D/S}(K_1 + K_2)\ar@{-}[d]^{\text{\rotatebox{90}{$\sim$}}}\\
\langle \Oo(D),K_1] + \langle \Oo(D) ,K_2] \ar[r]^{\sim}& \Tr_{D/S}(K_1) + \Tr_{D/S}(K_2)
}}\leqno{(1.3.5.2)}
$$\etiquette[1.3.5.2]{XVIII.eq1.3.5.2}%
and
$$\vcenter{
\xymatrix{
\langle \Oo(D_1 + D_2),K] \ar[r]^{\sim}
\ar@{-}[d]^{\text{\rotatebox{90}{$\sim$}}}& \Tr_{(D_1 +
  D_{2})/S}(K)\ar@{-}[d]^{\text{\rotatebox{90}{$\sim$}}\quad (XVII\
  \hyperref[XVII.const6.3.27.1]{6.3.27.1})}\\ 
\langle \Oo(D_1),K] + \langle \Oo(D_2),K] \ar[r]^{\sim}& \Tr_{D_1 /S}(K) +
    \Tr_{D_{2} /S} (K)
}}\leqno{(1.3.5.3)}
$$\etiquette[1.3.5.3]{XVIII.eq1.3.5.3}%
are commutative. By linearity, this makes it possible to define the
isomorphism \eqhyperref[XVIII.eq1.3.5.1]{1.3.5.1} for every relative divisor.

\subsubsection{}\etiquette[1.3.6]{XVIII.subsubsec1.3.6}
Let $D$ and $E$ be two effective relative Cartier divisors, and let
$f$ be a rational function such that
$$
\text{\rm div\,}(f)=D-E{\quoi},
$$
\ie an isomorphism between $\Oo(D)$ and $\Oo(E)$. The isomorphism $f$
between $\Oo(D)$ and $\Oo(E)$ defines, by \eqhyperref[XVIII.eq1.3.5.1]{1.3.5.1}, an isomorphism
\begin{equation*}
\langle f,K]:\Tr_{D/S}(K)\xrightarrow{\sim} \Tr_{E/S}(K){\quoi}.\tag{1.3.6.1}
\end{equation*}\etiquette[1.3.6.1]{XVIII.eq1.3.6.1}%
By\pageoriginale \eqhyperref[XVIII.eq1.3.5.2]{1.3.5.2} for the second formula, one has
\begin{align*}
& \langle fg,K]=\langle f,K].\langle g,K]\\
& \langle f,K_{1}+K_{2}]=\langle f,K_{1}]+\langle f,K_{2}]{\quoi}.
\end{align*}

If $\lambda$ is a section of $\Oo^{\ast}_{S}$, then, for every
invertible sheaf $\Ll$ on $X$, multiplication by
$\lambda$ is an automorphism of $\Ll$. By transport of
structure, this automorphism induces on the two sides of \eqhyperref[XVIII.eq1.3.3.8]{1.3.3.8}
automorphisms which correspond to one another under the isomorphism
\eqhyperref[XVIII.eq1.3.3.8]{1.3.3.8}. It follows that
$$
\langle \lambda, K]=\deg (K)(\lambda){\quoi}\;;
$$
if $\deg(K)\neq 0$, the isomorphism $\langle f,K]$ between $\Tr_{D/S}(K)$ and
$\Tr_{E/S}(K)$ thus depends on the choice of $f$.

Let $\lambda\in \Gamma(S,\Oo^{\ast}_{S})$, $g\in \Gamma(S,G)$,
let $\Ll$ be an invertible sheaf on $X$, and let $K$ be a $G$-torsor on
$X$. If $\langle \lambda,g]$ is the automorphism of $\langle \Ll,K]$ obtained
by transport of structure from the automorphisms $\lambda$ of
$\Ll$ and $g$ of $K$, then $\langle \lambda,g]$, identified with a
section of $G$, is given by
$\langle \lambda,g]=\langle \lambda,K]+\langle \Ll,g]$; since
$\langle \Oo(D),g]=\Tr_{D/S}(g)$, the preceding formula gives
\begin{equation*}
\langle \lambda,g]=\deg(K)(\lambda)+\deg(\Ll).g{\quoi}.\tag{1.3.6.2}
\end{equation*}\etiquette[1.3.6.2]{XVIII.eq1.3.6.2}

\subsubsection{}\etiquette[1.3.7]{XVIII.subsubsec1.3.7}
Let $f:X\rightarrow S$ be a smooth projective curve over $S$, and let
$G$ be a group over $S$ satisfying \hyperref[XVIII.par1.3.1.1]{1.3.1.1}. If
$X\xrightarrow{f_{0}}T\xrightarrow{\pi}S$ is the Stein factorization of
$f$, then $\pi$ is an étale covering of $S$, and the
sheaf $f_{\ast}G_{m}$ is the torus
$\prod_{T/S}\GG_{m}$. Locally on $S$ (for the étale topology),
$X$ is the disjoint union of a family of curves $(X_{i})_{i\in I}$ over
$S$ with nonempty geometrically connected fibers. If $K$ is a torsor
under $G$ and $\Ll$ an invertible sheaf on $X$, one then sets
$$
\langle \Ll,K]=\sum_{i}\langle \Ll|X_{i},K|X_{i}]
$$
and defines the degree of
$K:f_{\ast}\GG_{m}=\GG^{I}_{m}\rightarrow G$ to have
coordinates $\deg(K|X_{i}):\GG_{m}\rightarrow G$. The local
constructions globalize and give\pageoriginale
functors and morphisms
\begin{align*}
& \langle \Ll,K]_{X/S}=\Tr_{T/S}(\langle \Ll,K]_{X/T})\tag{1.3.7.1}\\
& \deg(K):f_{\ast}\GG_{m}=\prod_{T/S}\GG_{m}\longrightarrow G:\deg_{X/S}(K)=\prod_{T/S}\deg_{X/T}(K){\quoi}.\tag{1.3.7.2}
\end{align*}\etiquette[1.3.7.1]{XVIII.eq1.3.7.1}
\etiquette[1.3.7.2]{XVIII.eq1.3.7.2}

For an automorphism $\lambda$ of $\Ll$, defined by an
element $\lambda_{0}$ of $\Gamma(T,\Oo^{\ast}_{T})$, one has
\begin{equation*}
\langle \lambda,K]=\deg(K)(\lambda){\quoi}.\tag{1.3.7.3}
\end{equation*}\etiquette[1.3.7.3]{XVIII.eq1.3.7.3}

The total degree of $K$ is defined as the composite
$$
\deg \tot(K):\GG_{m}\longrightarrow
f_{\ast}\GG_{m}\xrightarrow{\deg(K)}G{\quoi}.
$$
The preceding results give:

\begin{enonce*}[remark]{Formulary 1.3.8}\etiquette[1.3.8]{XVIII.for1.3.8}
For every smooth projective curve $f:X\rightarrow S$, there is a
functor $\langle~,~]$, whose formation is compatible with every base
change, which associates to an invertible sheaf $\Ll$ on $X$ and a
torsor $K$ under $f^{\ast}G$ a torsor $\langle \Ll,K]$ under
$G$ on $S$. This functor is endowed with the following additional
structures and satisfies the following conditions:
\end{enonce*}

\setcounter{subsubsection}{8}
\paragraph{}\etiquette[1.3.8.1]{XVIII.par1.3.8.1}
$\langle \Ll,K]$ is biadditive in $\Ll$ and $K$; the
biadditivity isomorphisms are compatible with the associativity and
commutativity isomorphisms and with base changes, satisfy the
compatibility (\hyperref[XVIII.prop1.3.4]{1.3.4}), and are compatible with the
isomorphisms (\hyperref[XVIII.par1.3.8.2]{1.3.8.2}) below (via
\hyperrefext[XVII][XVII.par6.3.25.3]{6.3.25.3} and \hyperrefext[XVII][XVII.const6.3.27.1]{6.3.27.1}).

\paragraph{}\etiquette[1.3.8.2]{XVIII.par1.3.8.2}
For every effective relative divisor $E$ on $X/S$, there is a
canonical isomorphism
$$
\langle \Oo(E),K]\simeq \Tr_{E/S}(K|E){\quoi}.
$$

\paragraph{}\etiquette[1.3.8.3]{XVIII.par1.3.8.3}
The degree \eqhyperref[XVIII.eq1.3.7.2]{1.3.7.2} satisfies, for $\lambda\in f_{\ast}\GG_{m}$
an automorphism of $\Ll$,
$$
\langle \lambda,K]=\deg(K)(\lambda){\quoi}.
$$
For $\lambda\in \GG_{m}(S)$, the total degree satisfies
$$
\langle \lambda,K]=\deg\tot(K)(\lambda){\quoi}.
$$
Likewise,\pageoriginale for $g\in G(S)$ one has
$\langle \Ll,g]=\deg(\Ll).g$.

One deduces a canonical isomorphism:
\begin{equation*}
\langle f^{\ast}\Ll,K]\simeq \deg\tot(K)(\Ll){\quoi}.\tag{1.3.8.4}
\end{equation*}\etiquette[1.3.8.4]{XVIII.eq1.3.8.4}

It follows immediately from the definitions that, for every
homomorphism $\varphi:G\rightarrow G'$, there is a canonical
isomorphism, compatible with the preceding data:
\begin{equation*}
\langle \Ll,\varphi(K)]\simeq \varphi(\langle \Ll,K]){\quoi}.\tag{1.3.8.5}
\end{equation*}\etiquette[1.3.8.5]{XVIII.eq1.3.8.5}

\setcounter{paragraph}{5}
\paragraph{}\etiquette[1.3.8.6]{XVIII.par1.3.8.6}
Let $u:X\rightarrow Y$ be a flat morphism between smooth projective
curves over $S$. There are then isomorphisms compatible with
biadditivity and base changes:
\begin{itemize}
\item[--] For $\Ll$ an invertible sheaf on $Y$ and $K$ a
$G_{Y}$-torsor on $Y$, one has
$$
\langle \Ll, u^{\ast}K]\xrightarrow{\sim}\langle N_{X/Y}(\Ll),K]{\quoi}.
$$

\item[--] For $\Ll$ an invertible sheaf on $X$ and $K$ a
$G_{X}$-torsor on $X$, one has
$$
\langle u^{\ast}\Ll,K]\xrightarrow{\sim}
\langle \Ll,\Tr_{X/Y}K]{\quoi}.
$$
\end{itemize}
Moreover, these isomorphisms satisfy the evident compatibility for a
composite $uv$ of morphisms of curves; for $\Ll=\Oo(E)$, they
identify with the isomorphisms (\hyperrefext[XVII][XVII.const6.3.27.2]{6.3.27.2})
\begin{align*}
\Tr_{E/S}(u^{\ast}K) &\xrightarrow{\sim} \Tr_{u_{\ast}E/S}(K)\\
\Tr_{u^{\ast}E/S}(K) &\xrightarrow{\sim} \Tr_{E/S}\Tr_{X/Y}K{\quoi}.
\end{align*}
It remains to prove \hyperref[XVIII.par1.3.8.6]{1.3.8.6}. One reduces to the case in which $X$ and $Y$
have nonempty geometrically connected fibers, and notes that the
preceding explicit formulas give, for every section $s$ of $\Ll$
that does not vanish as a section of the direct image of $\Ll$, an
isomorphism $\Psi_{s}$ of the desired kind, homogeneous in $s$, and
therefore independent of $s$ when $\Ll$ satisfies
(\hyperref[XVIII.par1.3.2.1]{1.3.2.1}). The general case follows.


\subsubsection{}\etiquette[1.3.9]{XVIII.subsubsec1.3.9}
Under the hypotheses of \hyperref[XVIII.subsubsec1.3.7]{1.3.7},
we saw in (\hyperref[XVIII.for1.3.8]{1.3.8}) that every $G$-torsor $K$
on $X$ defines a functor $\langle \ast,K]$, denoted below by
$F(\ast)$,
endowed\pageoriginale with the following additional data and satisfying
the following conditions:
\begin{enumerate}
\renewcommand{\theenumi}{\roman{enumi}}
\renewcommand{\labelenumi}{(\theenumi)}
\item $F$ is a morphism of fppf stacks over $S$
  (\hyperrefext[XVII][XVII.subsec0.10]{0.10} and \textsc{Giraud}
  \cite{sga4.exp.XVIII.bib.1})
\begin{itemize}
\item[--] whose source is the stack whose $T$-objects, for an
$S$-scheme $T$, are the invertible sheaves on $X_{T}$;

\item[--] whose target is the stack of $G$-torsors.
\end{itemize}

\item The functor $F$ is endowed with additivity isomorphisms
$F(\Ll_{1}\otimes \Ll_{2})\xrightarrow{\sim}F(\Ll_{1})+F(\Ll_{2})$,
which are functorial, compatible with the associativity and
commutativity isomorphisms, and compatible with base changes.
\end{enumerate}

\begin{enonce*}{Theorem 1.3.10}\etiquette[1.3.10]{XVIII.thm1.3.10}
Under the hypotheses of \hyperref[XVIII.subsubsec1.3.7]{1.3.7}, the functor
$\Phi:K\mapsto \langle \ast,K]$ is
an equivalence between the category of $G$-torsors on $X$ and the
category of morphisms of stacks considered in
(\hyperref[XVIII.subsubsec1.3.9]{1.3.9}).
\end{enonce*}

One reduces easily to the case in which $X$ has nonempty connected fibers.

Let $F$ satisfy \hyperref[XVIII.subsubsec1.3.9]{1.3.9} (i) and (ii). For every
$S$-scheme $T$ and every section $t$ of $X_{T}/T$, denote also by $t$ the relative
divisor $t(T)$. Then $F(\Oo(t))$ is a $G_{T}$-torsor on $T$. We
denote by $\Psi(F)$ the torsor thus obtained on $X$, for $t$ the
``universal point'' of $X$ with values in $X$, defined by the
identity map from $X$ to $X$. For every section $t$ as above, one has
\begin{equation*}
F(\Oo(t))\simeq t^{\ast}\Psi(F){\quoi}.\tag{1.3.10.1}
\end{equation*}\etiquette[1.3.10.1]{XVIII.eq1.3.10.1}

The isomorphism \hyperref[XVIII.par1.3.8.2]{1.3.8.2} gives an isomorphism of functors
$\Psi\circ\Phi \simeq \Id$. We construct an isomorphism
$\Phi\circ\Psi\simeq \Id$.

For every section $t$ of $X_{T}$ over $T$, the isomorphism \eqhyperref[XVIII.eq1.3.10.1]{1.3.10.1}
is an isomorphism
$$
F(\Oo(t))\xrightarrow[\alpha]{\sim}\langle \Oo(t),\Psi(F)]{\quoi}.
$$

By additivity, every family $t_{1},\ldots, t_{n}$ of sections of
$X_{T}$ over $T$ defines an isomorphism
$$
F(\Oo(\sum t_{i}))\xrightarrow[\alpha]{\sim}\langle \Oo(\sum
t_{i}),\Psi(F)]{\quoi},
$$
and by \hyperref[XVIII.subsubsec1.3.9]{1.3.9} (ii), this isomorphism is independent
of the order of the $t_{i}$.

Take\pageoriginale $T=(X/S)^{n}$ and for $t_{1}\ldots t_{n}$ the
universal $n^{\text{\rm tuple}}$ of sections. If $T'=\Sym^{n}_{S}(X)$,
there is then on $X_{T'}$ a divisor $\sum t_{i}$ whose inverse image is
the divisor $\sum t_{i}$ on $X_{T}$. On $T$,
$\alpha$ is a symmetric isomorphism
$$
\alpha:F(\Oo(\sum t_{i}))\xrightarrow{\sim} \langle \Oo(\sum
t_{i}),\Psi(F)]
$$
between torsors that are inverse images of torsors on $T'$. By
hypothesis (\hyperref[XVIII.par1.3.1.1]{1.3.1.1}) (ii), this isomorphism comes from a unique
isomorphism on $T'$:
$$
\alpha:F(\Oo(\sum t_{i}))\xrightarrow{\sim} \langle \Oo (\sum
t_{i}),\Psi (F)]\qquad (\text{\rm on } T'){\quoi}.
$$

By \hyperrefext[XVII][XVII.prop6.3.9]{6.3.9}, one thus has, for every relative divisor $D$ on
$X$, an isomorphism
\begin{equation*}
\alpha_{D}:F(\Oo(D))\xrightarrow{\sim}
\langle \Oo(D),\Psi(F)]{\quoi},\tag{1.3.10.2}
\end{equation*}\etiquette[1.3.10.2]{XVIII.eq1.3.10.2}%
additive in $D$, whose formation is compatible with every base change.

Let $\Ll$ be a sufficiently ample invertible sheaf on $X$ in the
sense of \hyperref[XVIII.par1.3.2.1]{1.3.2.1}. Every nowhere-zero section $s$ of
$f_{\ast}\Ll$ defines, with the notation of \hyperref[XVIII.subsubsec1.3.2]{1.3.2}, an
isomorphism
$$
s:\Oo(D_{s})\xrightarrow{\sim}\Ll{\quoi},
$$
and hence, by \eqhyperref[XVIII.eq1.3.10.2]{1.3.10.2}, an isomorphism
\begin{equation*}
\alpha_{s}:F(\Ll)\xrightarrow{\sim}\langle \Ll,\Psi(F)]{\quoi}.\tag{1.3.10.3}
\end{equation*}\etiquette[1.3.10.3]{XVIII.eq1.3.10.3}

% END INCLUDED SGA4_English_current_buildable_source_20260906/source/18/drafts/SGA4_XVIII_LINES_1251_1650_English.texfrag

% BEGIN INCLUDED SGA4_English_current_buildable_source_20260906/source/18/drafts/SGA4_XVIII_LINES_1651_2050_English.texfrag

Let us prove that this isomorphism is independent of $s$. Indeed:
\begin{enumerate}
\renewcommand{\theenumi}{\roman{enumi}}
\renewcommand{\labelenumi}{(\theenumi)}
\item The construction of $\alpha_{s}$ commutes with every base change.

\item There exists a homomorphism $\varphi:\GG_{m}\rightarrow G$ such that
$$
\alpha_{\lambda s}=\varphi(\lambda).\alpha_{s}
$$
(by the functoriality of both sides of \eqhyperref[XVIII.eq1.3.10.3]{1.3.10.3}).
\end{enumerate}

One necessarily has $\varphi=0$, since otherwise the torsor on
$\PP((f_{\ast}\Ll)^{\vee})$ obtained by pulling back
$F(\Ll)-\langle \Ll,\Psi(F)]$ would have degree $\neq 0$. The
$\alpha_{s}$ therefore define a morphism from
$\PP((f_{\ast}\Ll)^{\vee})$ to $G$, and this morphism is
constant by hypothesis\pageoriginale (\hyperref[XVIII.par1.3.1.1]{1.3.1.1}) (i).
We may therefore set $\alpha_{\Ll}=\alpha_{s}$.

The isomorphisms $\alpha_{\Ll}$ thus form an isomorphism
of functors between the restrictions of the functors $F$ and $\langle *,\Psi(F)]$
to sufficiently ample sheaves. This isomorphism is compatible with the
additivity isomorphisms, which allows it to be extended by
linearity to all invertible sheaves and completes the construction
of
$$
\alpha:\Phi\circ\Psi\simeq \Id{\quoi}.
$$

\begin{enonce*}[remark]{Example 1.3.11}\etiquette[1.3.11]{XVIII.exem1.3.11}
$G=\GG_{m}$. Under the hypotheses of \hyperref[XVIII.subsubsec1.3.7]{1.3.7}, and in the
special case where $G=\GG_{m}$, we shall write $\langle~,~\rangle $ rather than
$\langle~,~]$; if $\Ll$ and $\Mm$ are two invertible sheaves
on $X$, then $\langle \Ll$, $\Mm\rangle $ is an invertible sheaf
on $S$. Moreover, there are isomorphisms (\hyperref[XVIII.par1.3.8.2]{1.3.8.2})
$$
\langle \Oo(D),\Mm\rangle \simeq N_{D/S}(\Mm){\quoi}.
$$

If $D$ and $E$ are disjoint effective relative divisors, then
$\Oo(E)|D$ is canonically isomorphic to $\Oo$, whence an isomorphism
\begin{equation}
\sigma_{D,E}:\langle \Oo(D),\Oo(E)\rangle \simeq N_{D/S}(\Oo(E))\simeq
N_{D/S}(\Oo)\simeq \Oo{\quoi}.\tag{1.3.11.1}
\end{equation}\etiquette[1.3.11.1]{XVIII.eq1.3.11.1}
\end{enonce*}

\begin{enonce*}{Proposition 1.3.12}\etiquette[1.3.12]{XVIII.prop1.3.12}
Let $(D_{i})_{i=1,2}$ and $E$ be effective relative divisors, with
$D_{i}$ disjoint from $E$ $(i=1,2)$. Let $f$ be a rational function
satisfying $\text{\rm div\,}(f)=D_{1}-D_{2}$, \ie an isomorphism
$f:\Oo(D_{1})\xrightarrow{\sim}\Oo(D_{2})$. The following diagrams
of isomorphisms are commutative:
\begin{equation*}
\vcenter{
\xymatrix@C=2.5cm@R=1.25cm{
\langle \Oo(D_1), \Oo(E)\rangle   \ar[d]^{\langle f,\Oo(E)\rangle } \ar[r]^(.65){\sigma_{D_{1},E}} & \Oo\ar[d]^{N_{E/S}(f)}\\
\langle \Oo(D_2),\Oo(E)\rangle    \ar[r]^(.65){\sigma_{D_{2},E}} & \Oo
}}\tag{1.3.12.1}
\end{equation*}\etiquette[1.3.12.1]{XVIII.eq1.3.12.1}
\begin{equation*}
\vcenter{
\xymatrix@C=2.5cm@R=1.25cm{
\langle \Oo(E), \Oo(D_1)\rangle   \ar[d]^{\langle \Oo(E),f\rangle } \ar[r]^(.65){\sigma_{E,D_{1}}} & \Oo\ar[d]^{N_{E/S}(f)}\\
\langle \Oo(E),\Oo(D_2)\rangle    \ar[r]^(.65){\sigma_{E,D_{2}}} & \Oo
}}\tag{1.3.12.2}
\end{equation*}\etiquette[1.3.12.2]{XVIII.eq1.3.12.2}
\end{enonce*}

The\pageoriginale formula \eqhyperref[XVIII.eq1.3.12.2]{1.3.12.2} simply expresses the
commutativity of the diagram
$$
\xymatrix@C=2.5cm@R=1.25cm{
N_{E/S}(\Oo(D_1))  \ar[r]^(.55){\sim} \ar[d]^{N_{E/S}(f)} & \Oo \ar[d]^{N_{E/S}(f)} \\         
N_{E/S}(\Oo(D_2))  \ar[r]^(.55){\sim} & \Oo.
}
$$

For every $S$-scheme $T$ and every invertible sheaf $\Ll$ on
$X_{T}$, set $F(\Ll)=N_{E_{T}/T}(\Ll)$. The
hypotheses of \hyperref[XVIII.thm1.3.10]{1.3.10} are satisfied by this functor, so
there exists a unique invertible sheaf $\Mm$ equipped with
an isomorphism of functors compatible with biadditivity:
$$
N_{E_{T}/T}(\Ll)\xrightarrow{\sim}\langle \Ll,\Mm\rangle {\quoi}.
$$

Moreover, if $d:T\rightarrow X_{T}$ is a section of $X$ over $T$, there
is a canonical isomorphism
$$
d^{\ast}\Mm\xrightarrow{\sim} N_{E_{T}/T}(\Oo(d)){\quoi}.
$$
Taking for $d$ the universal section of $X$, $\Delta:X\rightarrow
X_{X}$, we obtain
$$
\Mm\xrightarrow{\sim} N_{E_{X}/X}(\Oo(\Delta)){\quoi}.
$$

Locally on $S$, $E$ is a sum of sections $E=\sum t_{i}$
(\hyperrefext[XVII][XVII.subsec6.3]{6.3}); for such an $E$, one has
$$
N_{E_{X}/X}(\Oo(\Delta))=\otimes
t^{\ast}_{i}\Oo(\Delta)=\sum\Oo(t_{i})=\Oo(E){\quoi},
$$
and this isomorphism, being independent of the order of the $t_{i}$, defines
a canonical isomorphism\nde{The argument seems to assume that, locally,
the $t_i$ are well defined up to permutation, which is not true,
as can be seen in the case where $S=\Spec k[\epsilon]$ is the
spectrum of the algebra of dual numbers, $X=\PP^1_S$, and $E$ is the
closed subscheme of $\AAA^1_S=\Spec \Oo_S[T]\subset X$ defined by
the equation $(T-\varepsilon)(T+\varepsilon)=T^2$. It seems instead necessary
to apply the ``cancelled descent'' condition
\hyperref[XVIII.par1.3.1.1]{1.3.1.1}~(ii) to obtain the result in the situation of the
universal divisor on $X$ after base change
by $\Sym^n_S X\to S$, before deducing it for the divisor $E$.}
\begin{equation*}
N_{E_{X}/X}(\Oo(\Delta))\sim \Oo(E){\quoi},\tag{1.3.12.3}
\end{equation*}\etiquette[1.3.12.3]{XVIII.eq1.3.12.3}%
so that
\begin{equation*}
N_{E_{T}/T}(\Ll)\xrightarrow{\sim}\langle \Ll,\Oo(E)\rangle{\quoi}.%
\tag{1.3.12.4}
\end{equation*}\etiquette[1.3.12.4]{XVIII.eq1.3.12.4}

The isomorphism \eqhyperref[XVIII.eq1.3.12.3]{1.3.12.3} is compatible with biadditivity and
base changes and is characterized, in addition, by the fact that for
every\pageoriginale section $d$ of $X_{T}$ over $T$, the isomorphism
$$
N_{E/T}(\Oo(d))\xrightarrow{\sim}\langle \Oo(d),\Oo(E)\rangle
\xrightarrow{\sim} d^{\ast}\Oo(E)
$$
is the pullback of the universal isomorphism \eqhyperref[XVIII.eq1.3.12.3]{1.3.12.3}. In
particular, for $d$ disjoint from $E$, the diagram
$$
\xymatrix{
N_{E/T}(\Oo(d))  \ar[r]^{\sim} \ar@{=}[d]& \langle \Oo(d), \Oo(E)\rangle 
\ar[r]^{\sim} & d^\ast\Oo(E) \ar@{=}[d]\\
\Oo              \ar@{=}[rr] &          & \Oo  
}
$$
is commutative. By additivity\nde{The same objection raised above applies
here as well. One can argue in the same way by considering ``the universal
divisor disjoint from $E$,'' defined after base change by
$\Sym^n_S (X-E)\to S$.},
it follows that, for $D$ disjoint from
$E$, the diagram
\begin{equation*}
\vcenter{
\xymatrix{
N_{E/T}(\Oo(D))  \ar[r]^{\sim} \ar@{=}[d]& \langle \Oo(D), \Oo(E)\rangle 
\ar[r]^{\sim} & N_{D/T}(\Oo(E)) \ar@{=}[d]\\
\Oo              \ar@{=}[rr] &          & \Oo  
}}\tag{1.3.12.5}
\end{equation*}\etiquette[1.3.12.5]{XVIII.eq1.3.12.5}%
is commutative. The arrow $\sigma_{E,D}$ is therefore the composite of
the inverse of \hyperref[XVIII.eq1.3.12.4]{1.3.12.4} and an evident arrow
$$
\langle \Oo(D),\Oo(E)\rangle \xleftarrow{\sim} N_{E/S}(\Oo(D))\longrightarrow
\Oo
$$
and \eqhyperref[XVIII.eq1.3.12.1]{1.3.12.1} follows at once, since
\hyperref[XVIII.eq1.3.12.4]{1.3.12.4} is functorial.

We deduce from \hyperref[XVIII.prop1.3.12]{1.3.12} (\cf \textsc{Serre}
\cite[Ch.~III, prop.~7, p.~46]{sga4.exp.XVIII.bib.6})

\begin{enonce*}{Corollary 1.3.13}\etiquette[1.3.13]{XVIII.cor1.3.13}
Let $(D_{i})_{i=1,2}$ and $(E_{i})_{i=1,2}$ be effective relative
divisors, with $D_{i}$ disjoint from $E_{j}$, let $f$ be a rational
function satisfying $\text{\rm div\,}(f)=D_{1}-D_{2}$, and let $g$ be a
rational function satisfying $\text{\rm div\,}(g)=E_{1}-E_{2}$. Then
$$
N_{E_{1}/S}(f).N_{E_{2}/S}(f)^{-1}=N_{D_{1}/S}(g).N_{D_{2}/S}(g)^{-1}
$$
or, briefly,
\begin{equation*}
f(\text{\rm div\,}(g))=g(\text{\rm div\,}(f)){\quoi}.\tag{1.3.13.1}
\end{equation*}\etiquette[1.3.13.1]{XVIII.eq1.3.13.1}
\end{enonce*}

This formula, in the form
$$
N_{D_{2}/S}(g)N_{E_{1}/S}(f)=N_{E_{2}/S}(f)N_{D_{1}/S}(g){\quoi},
$$
expresses, via \eqhyperref[XVIII.eq1.3.12.1]{1.3.12.1} and \eqhyperref[XVIII.eq1.3.12.2]{1.3.12.2}, the commutativity of the
diagram
$$
\xymatrix@C=2.25cm@R=1.5cm{
\langle \Oo(D_1), \Oo(E_1)\rangle   \ar[r]^{\langle \Oo(D_1), g\rangle } \ar[d]^{\langle f,\Oo(E_1)\rangle }&
\langle \Oo(D_1),\Oo(E_2)\rangle  \ar[d]^{\langle f,\Oo(E_2)\rangle }\\
\langle \Oo(D_2), \Oo(E_1)\rangle   \ar[r]^{\langle \Oo(D_2),g\rangle } & \langle \Oo(D_2), \Oo(E_2)\rangle .
}
$$

\begin{enonce*}{Corollary 1.3.14}\etiquette[1.3.14]{XVIII.cor1.3.14}
For\pageoriginale $G=\GG_{m}$, the degree (\hyperref[XVIII.subsubsec1.3.7]{1.3.7}) coincides
with the usual degree of invertible sheaves.
\end{enonce*}

It suffices to compare formulas \eqhyperref[XVIII.eq1.3.7.3]{1.3.7.3} and \eqhyperref[XVIII.eq1.3.12.1]{1.3.12.1} for $f$
the pullback of $f_0\in \Gamma(S,\Oo^{\ast}_{S})$.

\setcounter{subsubsection}{14}
\subsubsection{}\etiquette[1.3.15]{XVIII.subsubsec1.3.15}
For any sufficiently ample $\Ll$ and $\Mm$,
there exist locally on $S$ disjoint divisors $D$ and $E$ and
isomorphisms
$$
\alpha:\Ll\longrightarrow \Oo(D)\quaet
\beta:\Mm\longrightarrow \Oo(E){\quoi}.
$$
Denote by $\tau$ the isomorphism making the diagram
$$
\xymatrix@C=1.5cm@R=1.25cm{
\langle \Ll,\Mm\rangle  \ar[r]^{\langle \alpha,\beta\rangle }
\ar[d]^{\tau}_{\text{\rotatebox{90}{$\backsim$}}}& \langle \Oo(D), \Oo(E)\rangle  \ar[r]^(.7){\sigma_{D,E}}& \Oo\ar@{=}[d]\\
\langle \Mm,\Ll\rangle  \ar[r]^{\langle \beta,\alpha\rangle }& \langle \Oo(E), \Oo(D)\rangle  \ar[r]^(.7){\sigma_{E,D}}& \Oo.
}
$$
It follows from \hyperref[XVIII.prop1.3.12]{1.3.12} that this isomorphism is independent
of the particular choice of $D$, $E$, $\alpha$, and $\beta$; it therefore globalizes to
a canonical isomorphism
\begin{equation*}
\tau:\langle \Ll,\Mm\rangle \xrightarrow{\sim}
\langle \Mm,\Ll\rangle {\quoi}.\tag{1.3.15.1}
\end{equation*}\etiquette[1.3.15.1]{XVIII.eq1.3.15.1}

This isomorphism satisfies $\tau^{2}=1$; it is compatible with
biadditivity, since \eqhyperref[XVIII.eq1.3.11.1]{1.3.11.1} is, and extends by linearity to
arbitrary invertible sheaves.

It follows from \eqhyperref[XVIII.eq1.3.12.5]{1.3.12.5} that the morphism
\eqhyperref[XVIII.eq1.3.12.4]{1.3.12.4} identifies with the composite
$$
N_{E_{T}/T}(\Ll)\xleftarrow{\sim}
\langle \Oo(E),\Ll\rangle \xrightarrow{\sim\tau}\langle \Ll,\Oo(E)\rangle 
$$
(verify this for $\Ll=\Oo(D)$ with $D$ disjoint from $E$).

\begin{enonce*}{Formulary 1.3.16}\etiquette[1.3.16]{XVIII.for1.3.16}
For\pageoriginale every smooth projective curve $f:X\rightarrow
S$, there is a functor $\langle~,~\rangle $, whose construction commutes with every
base change, that associates with two invertible sheaves $\Ll$ and
$\Mm$ on $X$ an invertible sheaf
$\langle \Ll,\Mm\rangle $ on $S$. This functor is endowed with the
following additional structures and satisfies the following
conditions:
\end{enonce*}

\setcounter{subsubsection}{16}
\paragraph{}\etiquette[1.3.16.1]{XVIII.par1.3.16.1}
With the notation of \hyperref[XVIII.for1.3.8]{1.3.8}, for $G=\GG_{m}$, one has
$$
\langle \Ll,\Mm\rangle =\langle \Ll,\Mm]{\quoi},
$$
in particular, $\langle \Ll,\Mm\rangle $ is biadditive in
$\Ll$, $\Mm$ in the sense specified in \hyperref[XVIII.par1.3.8.1]{1.3.8.1}, and this
biadditivity is compatible with that of the right-hand side of a canonical
isomorphism (functorial in $\Mm$)
$$
\langle \Oo(E),\Mm\rangle \xrightarrow{\sim} N_{E/S}(\Mm){\quoi}.
$$

\paragraph{}\etiquette[1.3.16.2]{XVIII.par1.3.16.2}
There is a symmetry isomorphism
$$
\tau:\langle \Ll,\Mm\rangle \xrightarrow{\sim}\langle \Mm,\Ll\rangle {\quoi},
$$
functorial in $\Ll$ and $\Mm$, compatible with
base changes and the biadditivity isomorphisms, and satisfying
$\tau^{2}=1$. It makes the following diagram commutative, for disjoint
divisors $D$ and $E$:
$$
\xymatrix{
N_{D/S}(\Oo(E))\ar@{-}[d]^{\text{\rotatebox{90}{$\backsim$}}}& \langle \Oo(D), \Oo(E)\rangle \ar[l]\ar[r]^{\tau} &
\langle \Oo(E), \Oo(D)\rangle  \ar[r]&
N_{E/S}(\Oo(D))\ar@{-}[d]^{\text{\rotatebox{90}{$\backsim$}}}\\
\Oo \ar@{=}[rrr]&    & & \Oo.
}
$$

These formulas imply formulas \eqhyperref[XVIII.eq1.3.12.1]{1.3.12.1}, \eqhyperref[XVIII.eq1.3.12.2]{1.3.12.2},
\eqhyperref[XVIII.eq1.3.13.1]{1.3.13.1}, and

\paragraph{}\etiquette[1.3.16.3]{XVIII.par1.3.16.3}
For $a$, $b\in \Gamma(S,\Oo^{\ast}_{S})$, denoting by
$\langle a,b\rangle $ the automorphism of $\langle \Ll,\Mm\rangle $, identified with
a section of $\GG_{m}$, obtained by functoriality from the
automorphisms $a$ and $b$ of $\Ll$ and $\Mm$, one has
$$
\langle a,b\rangle =a^{\deg(\Mm)} b^{\deg(\Ll)}{\quoi}.
$$

\paragraph{}\etiquette[1.3.16.4]{XVIII.par1.3.16.4}
If $\Ll_0$ (\resp $\Mm_0$) is an
invertible module on $S$, there is a canonical isomorphism
\begin{align*}
\langle f^{\ast}\Ll_0,\Mm\rangle  &\xrightarrow{\sim}
\Ll_0^{\otimes\deg(\Mm)}\\
\langle \Ll,f^{\ast}\Mm_0\rangle  &\xrightarrow{\sim} \Mm_0^{\otimes\deg(\Ll)}\qquad){\quoi}.\tag*{(\resp} 
\end{align*}

\paragraph{}\etiquette[1.3.16.5]{XVIII.par1.3.16.5}
If\pageoriginale $u:X\rightarrow Y$ is a finite flat morphism of
smooth projective curves over $S$, then, for $\Ll$ an
invertible sheaf on $X$ and $\Mm$ an invertible sheaf on
$Y$, the isomorphisms \hyperref[XVIII.par1.3.8.6]{1.3.8.6}
\begin{align*}
\langle \Ll,u^{\ast}\Mm\rangle  &\sim
\langle N_{X/Y}\Ll,\Mm\rangle \\
\langle u^{\ast}\Mm,\Ll\rangle  &\sim \langle \Mm,N_{X/Y}\Ll\rangle 
\end{align*}
are deduced from one another by the symmetry
\eqhyperref[XVIII.par1.3.16.2]{1.3.16.2} (verify this for
$\Ll=\Oo(E)$ and $\Mm=\Oo(F)$ with $E$ and $F$
disjoint).

\paragraph{}\etiquette[1.3.16.6]{XVIII.par1.3.16.6}
For an invertible sheaf $\Ll$ on $X$, the symmetry arrow
$$
\tau:\langle \Ll,\Ll\rangle \longrightarrow
\langle \Ll,\Ll\rangle 
$$
is multiplication by $(-1)^{\deg(\Ll)}$.

\subsubsection{}\etiquette[1.3.17]{XVIII.subsubsec1.3.17}
It remains to prove \hyperref[XVIII.par1.3.16.6]{1.3.16.6}. We shall give two proofs, the
first using the theory of the determinant (\SGA 6
XI)\nde{That exposé was not written, as explained in
\cite{sga4.exp.XVIII.bib.0ter},
which also contains a warning concerning
the multiplicativity property of the
determinant in distinguished triangles: it requires care, without
which counterexamples arise.
A reference on this subject is \cite{sga4.exp.XVIII.bib.3bis}.}.

\sisi{}{Before proceeding to each of the proofs, one may observe
that if $\Ll_1$ and $\Ll_2$ are two invertible sheaves on $X$ and
$\Ll=\Ll_1\otimes\Ll_2$, then the automorphism $\tau_{\Ll}\colon \langle
\Ll,\Ll\rangle \longrightarrow \langle \Ll,\Ll\rangle$ (identified with an
invertible function on $S$) is the product of the automorphisms
$\tau_{\Ll_1}$ and $\tau_{\Ll_2}$. In the situation where $S$ is the spectrum
of a field (and $X$ is geometrically connected),
one may therefore suppose that $\Ll=\Oo(D)$ with $D$ a
nonzero effective divisor. Since $\tau^2=\Id$, after replacing $\Ll$ by a
sufficiently large odd power (or writing it as a quotient of two ample
bundles), one may also make $\Ll$ as ample as desired,
if necessary.

The arguments below will also require reducing to the case where
$S$ is the spectrum of a field. For this, it suffices to reduce to the
case where $S$ is reduced, and this can be done by showing that fppf-locally
on $S$, the datum $(X,\Ll)$ comes from a similar datum
over a reduced base. Before making this reduction, one may reduce to the
case where $X\to S$ is equipped with a section $s\colon S\to
X$ and where the geometric fibers of $X\to S$ are connected and have some
genus $g$. Using \cite[\S 5.2]{sga4.exp.XVIII.bib.4}, one can obtain a smooth
projective curve $X_0\to S_0$ of genus $g$, with $S_0$ smooth of finite type
over $\Spec(\ZZ)$, such that \emph{locally}, every relative curve as
above is obtained from it by base change. Using the section
$s_0\colon S_0\to X_0$, one can rigidify the Picard scheme
$S_1=\sheafPic_{X_0/S_0}\to S_0$ and thus obtain a ``universal'' line bundle $\Ll_1$ on $X_1=X_0\times_{S_0}S_1$. Locally, every
relative curve as above equipped with a line bundle comes from
$(X_1,\Ll_1)$ over $S_1$, which is indeed reduced\nde{The preceding two
paragraphs were added in this edition.}.}

\emph{1\textsuperscript{st} proof}: Let $u:T\rightarrow S$ be a
finite locally free morphism. If $\Ll$ is an invertible sheaf on
$T$, there is an isomorphism
\begin{equation*}
N_{T/S}(\Ll)\simeq \det(u_{\ast}\Ll)\det(u_{\ast}\Oo_{T})^{-1}\tag{1.3.17.1}
\end{equation*}\etiquette[1.3.17.1]{XVIII.eq1.3.17.1}%
uniquely determined by the following three conditions:
\begin{enumerate}
\renewcommand{\theenumi}{\alph{enumi}}
\renewcommand{\labelenumi}{\theenumi)}
\item its construction commutes with base changes;

\item for $\Ll=\Oo_{T}$, it is the identity morphism of
$\Oo_{S}$;

\item it is functorial in $\Ll$.
\end{enumerate}

Let $D$ and $E$ be two effective relative divisors on $X$, let $i$
be the inclusion of $D$ in $X$, and let $u$ be the projection of $D$ onto $S$. The
short exact sequence
$$
0\longrightarrow \Oo (-E)\longrightarrow \Oo\longrightarrow
\Oo_{E}\longrightarrow 0
$$
allows the isomorphism
$$
\langle \Oo(D),\Oo(-E)\rangle \simeq N_{D/S}(\Oo(-E))\simeq
\det(u_{\ast}i^{\ast}\Oo(-E))\det (u_{\ast}i^{\ast}\Oo)^{-1}

% END INCLUDED SGA4_English_current_buildable_source_20260906/source/18/drafts/SGA4_XVIII_LINES_1651_2050_English.texfrag

% BEGIN INCLUDED SGA4_English_current_buildable_source_20260906/source/18/drafts/SGA4_XVIII_LINES_2051_2450_English.texfrag
$$
as\pageoriginale an isomorphism
$$
\langle \Oo(D),\Oo(-E)\rangle =\det Ru_{\ast}(Li^{\ast}\Oo_{E})^{-1}\simeq
\det Rf_{\ast}(\Oo_{D}\mathop{\otimes}^{\LL}
\Oo_{E})^{-1}{\quoi}.
$$
or
\begin{equation*}
\langle \Oo(D),\Oo(E)\rangle \simeq \det Rf_{\ast}(\Oo_{D}\mathop{\otimes}^{\LL}\Oo_{E}){\quoi}.\tag{1.3.17.2}
\end{equation*}\etiquette[1.3.17.2]{XVIII.eq1.3.17.2}

For disjoint $D$ and $E$, one has
$\Oo_{D}\displaystyle{\mathop{\otimes}^{\LL}}\Oo_{E}=0$, and
the resulting isomorphism of $\langle \Oo(D)$, $\Oo(E)\rangle $ with
$\Oo_{S}$ is the one already considered. It follows that, for disjoint
$D$ and $E$, the symmetry $\tau$ and the symmetry isomorphism
$\Oo_{D}\displaystyle{\mathop{\otimes}^{\LL}}\Oo_{E}\xrightarrow{\sim}\Oo_{E}\displaystyle{\mathop{\otimes}^{\LL}}\Oo_{D}$
are compatible via \eqhyperref[XVIII.eq1.3.17.2]{1.3.17.2}. By specialization\nde{It
seems that here one must first reduce to the case where $S$ is the
spectrum of a field, as explained above.
It is then possible to make $D$ and
$E$ occur in two families of divisors parametrized by an integral
scheme in such a way that the divisors in the two families are disjoint
over the generic point. This may be obtained, for example, by considering
the universal situation over which the condition that $D$ and $E$ be
disjoint holds on a dense open subset of $\Sym^n_S
X\times \Sym^m_S X$, which is integral.},
this remains true for arbitrary $D$ and $E$.

When $D=E$ is defined by an ideal $I$, one has
\begin{alignat*}{2}
&
ℋ^{0}(\Oo_{D}\mathop{\otimes}^{\LL}\Oo_{D})=\Oo_{D}
&\qquad& \text{\rm symmetry } = +\sisi{1}{\Id}\\
& ℋ^{-1}(\Oo_{D}\mathop{\otimes}^{\LL}\Oo_{D})=I/I^{2} &&
\text{\rm symmetry } =-\sisi{1}{\Id}\\
& ℋ^{i}(\Oo_{D}\mathop{\otimes}^{\LL}\Oo_{D})=0 &&
\text{\rm for } i\neq 0, -1{\quoi}.
\end{alignat*}
Consequently, the map
$$
\tau:\langle \Oo(D), \Oo(D)\rangle \xrightarrow{\sim}\langle \Oo(D), \Oo(D)\rangle 
$$
is $(-1)^{\rk I/I^{2}}=(-1)^{\deg(\Oo(D))}$, which proves \hyperref[XVIII.par1.3.16.6]{1.3.16.6}.

\smallskip
\noindent
\emph{2\textsuperscript{nd} proof}: By reduction to a “universal”
case, we may suppose that $S$ is reduced. This case reduces to that in which $S$
is the spectrum of an algebraically closed field $k$. In this case, let
$X_{1}$, $X_{2}$, $Y_{1}$, and $Y_{2}$ be four divisors on $X$, let $f$ and
$g$ be two rational functions, and suppose that
\begin{align*}
&\text{\rm div\,}(f)=X_{2}-X_{1}\\
&\text{\rm div\,}(g)=Y_{2}-Y_{1}\\
&(X_{1}\cup Y_{2})\cap (X_{2}\cup Y_{1})=∅{\quoi}.
\end{align*}
The maps $f$ and $g$ then define\nde{Compared with the original
edition, numerous “signs” have been modified here and there through
the end of this paragraph.}
$$
\langle f,g\rangle :\langle \Oo(X_{2}), \Oo(Y_{2})\rangle \longrightarrow
\langle \Oo(X_{1}),\Oo(Y_{1})\rangle {\quoi}.
$$
The\pageoriginale two sides are canonically isomorphic to $k$
\eqhyperref[XVIII.eq1.3.11.1]{1.3.11.1}, so that $\langle f,g\rangle $ is identified with an
element of $k^{\ast}$.

\hypertarget{sga4.exp.XVIII.lemma.1.3.18}{}\begin{enonce*}{Lemma 1.3.18}\etiquette[1.3.18]{XVIII.lem1.3.18}
One has
$$
\langle f,g\rangle =\prod_{x\in X_{1}\cup Y_{2}}(-1)^{v_{x}(f)v_{x}(g)}
(f^{v_{x}(g)}/g^{v_{x}(f)})(x){\quoi}.
$$
\end{enonce*}

Let us show that (\hyperref[XVIII.lem1.3.18]{1.3.18}) implies
(\hyperref[XVIII.par1.3.16.6]{1.3.16.6}). We may suppose that
$\Ll\sim \Oo(X)\sim \Oo(Y)$ with $X$ and $Y$ effective, and
that $X\cap Y=∅$\nde{We saw above that we could suppose that
$\Ll$ was very ample.}. Let $f$ be a rational function such that
$\text{\rm div\,}(f)=Y-X$\nde{Care should be taken not to confuse $X$ with the
ambient curve, which is also denoted by $X$.}. The commutative diagram
$$
\xymatrix@C=2.5cm@R=1.25cm{
\langle \Ll, \Ll\rangle  \ar[r]^{\tau} \ar@{=}[d]&
\langle \Ll,\Ll\rangle  \ar@{=}[d]\\
\langle \Oo(X), \Oo(Y)\rangle   \ar[r]^{\langle f^{-1},f\rangle } & \langle
\Oo(Y), \Oo(X)\rangle 
}
$$
shows that
$$
\tau=\langle f^{-1},f\rangle =\prod_{x\in
Y}(-1)^{-v_{x}(f)^{2}}=(-1)^{\deg(\Ll)}{\quoi}.
$$ Let us prove \hyperref[XVIII.lem1.3.18]{1.3.18}. We reduce to supposing that $X_{i}$ and
$Y_{i}$ are very ample. There are then factorizations $f=f_{1}f_{2}$;
$g=g_{1}g_{2}$ with
\begin{alignat*}{2}
&\text{\rm div\,}(f_{1}) = X'-X_{1} &\qquad& \text{\rm div\,}(g_{1})
= Y'-Y_{1}\\
& \text{\rm div\,}(f_{2}) = X_{2}-X' &\qquad& \text{\rm div\,}(g_{2})
= Y_{2}-Y'{\quoi},
\end{alignat*}
the divisors $X'$, $Y'$, and $X_{1}\cup X_{2}\cup Y_{1}\cup Y_{2}$
being pairwise disjoint.

Denote by $\langle~,~\rangle^{\ast}$ the right-hand side of
(\hyperref[XVIII.lem1.3.18]{1.3.18}), and let us prove that
$$ \langle f,g\rangle ^{\ast}=\langle f_{1},g_{1}\rangle ^{\ast}\cdot
\langle f_{2},g_{2}\rangle ^{\ast}{\quoi}.
$$
By \textsc{Serre} \cite[Ch.~III, prop.~6, p.~44]{sga4.exp.XVIII.bib.6}, one has
\begin{equation*}
\prod_{x}(-1)^{v_{x}(f_{1})v_{x}(g_{2})}
(f_{1}^{v_{x}(g_{2})}/g_{2}^{v_{x}(f_{1})}) (x)=1{\quoi}.\tag{1.3.18.1}
\end{equation*}\etiquette[1.3.18.1]{XVIII.eq1.3.18.1}%
It follows that
% SGA-SRC-001129 — exact continuation restored from the original scan
\begin{eqnarray*}
\langle f_{1},g_{1}\rangle ^{\ast}\cdot \langle f_{2},g_{2}\rangle
^{\ast} &= &g_{1}(X_{1})f_{1}(Y')g_{2}(X')f_{2}(Y_{2})\cdot\\
& & \prod_{X_{1}\cup X'\cup Y'\cup Y_{2}}(-1)^{v_{x}(f_{1})v_{x}(g_{2})}
(f_{1}^{v_{x}(g_{2})}\cdot g_{2}^{-v_{x}(f_{1})})(x)\\
&=&\prod_{x\in X'} g_{2}(x)^{-v_{x}(f_{2})}\cdot
g_{2}(x)^{-v_{x}(f_{1})}\cdot \prod_{x\in
Y'}f_{1}(x)^{v_{x}(g_{1})}\cdot f_{1}(x)^{v_{x}(g_{2})}\cdot\\
&&\prod_{x\in X_{1}\cup Y_{2}}(-1)^{v_{x}(f)v_{x}(g)}
((f_{1}f_{2})^{v_{x}(g_{2})}\cdot
(g_{1}g_{2})^{-v_{x}(f_{1})})(x)=\langle f_{1}f_{2},g_{1}g_{2}\rangle
^{\ast}{\quoi}. 
\end{eqnarray*}
One\pageoriginale has, by functoriality of $\langle~,~\rangle $,
$\langle f,g\rangle =\langle f_{1},g_{1}\rangle \langle
f_{2},g_{2}\rangle $, and it is readily checked that $\langle
f_{i},g_{i}\rangle =\langle f_{i},g_{i}\rangle ^{\ast}$\nde{This is the
special case in which one assumes in addition that $X_1\cap Y_2=∅$; it follows
from \hyperref[XVIII.eq1.3.12.1]{1.3.12.1} and \hyperref[XVIII.eq1.3.12.2]{1.3.12.2}.}. The lemma
follows.

\subsection{Strictly commutative Picard
  categories}\etiquette[1.4]{XVIII.subsec1.4}

Some results from “general topology” will be needed
to express the results of \hyperref[XVIII.subsec1.3]{1.3} by a “universal
coefficient formula.”

The results of this section are suggested in a letter from
A.~\textsc{Grothendieck} to J.L. \textsc{Verdier} (1966).

The notions of associativity and commutativity for bifunctors
(see below) were introduced by \textsc{Mac Lane}.

\subsubsection{}\etiquette[1.4.1]{XVIII.subsubsec1.4.1}
Let $C$ be a category, let $F:C\times C\rightarrow C$ be a functor, and let
$$
\sigma:F(F(X,Y),Z)\xrightarrow{\sim} F(X,F(Y,Z))
$$
be an isomorphism of trifunctors. The pair $(F,\sigma)$ is said to be
an \emph{associative functor}, or $\sigma$ is said to be an
\emph{associativity datum} on $F$, if the following condition holds:

\begin{itemize}
\item[(Ass)] For every family $(X_{i})_{i\in I}$ of objects of
$C$, denoting by $e:I\rightarrow M(I)$ the canonical map from
$I$ into the free monoid (without identity) generated by $I$, there exist a
map $ℱ:M(I)\rightarrow \Ob C$, isomorphisms
$a_{i}:ℱ(e_{i})\xrightarrow{\sim} X_{i}$, and isomorphisms
$a_{g,h}:ℱ(gh)\xrightarrow{\sim}F(ℱ(g),ℱ(h))$ such that the
diagrams
\begin{equation*}
\vcenter{
\xymatrix@C=1.25cm@R=1.25cm{
ℱ(f(gh))  \ar@{=}[d]\ar[r]^-{\sim}_-{a_{f,gh}} &
F(ℱ(f),ℱ(gh))
\ar[r]^-{\sim}_-{a_{g,h}}&
F(ℱ(f),F(ℱ(g),ℱ(h)))\\
ℱ((fg)h)  \ar[r]^-{\sim}_-{a_{fg,h}}& F(ℱ(fg),ℱ(h))
\ar[r]^-{\sim}_-{a_{f,g}}& F(F(ℱ(f),ℱ(g)),ℱ(h))\ar[u]_{\sigma}^{\text{\rotatebox{90}{$\backsim$}}}
}}\tag{1.4.1.1}
\end{equation*}\etiquette[1.4.1.1]{XVIII.eq1.4.1.1}%
are commutative.

We shall not need the fact that axiom (Ass) is equivalent to
the “pentagon axiom,” namely that the following diagram is
commutative:
$$
\xymatrix@C=-.5cm{
                                  & F(F(X,Y),F(Z,T))&  \\
F(X,F(Y,F(Z,T))) \ar@{<->}[ru]    &                 & F(F(F(X,Y),Z),T)\ar@{<->}[lu] \\
F(X,F(F(Y,Z),T)) \ar@{<->}[rr]\ar@{<->}[u]    &    & F(F(X,F(Y,Z)),T). \ar@{<->}[u] \\
}
$$

Now suppose\pageoriginale given functorial isomorphisms
\begin{align*}
\sigma_{X,Y,Z} &: F(F(X,Y),Z)\xrightarrow{\sim}F(X,F(Y,Z))\\
\tau_{X,Y} &: F(X,Y)\xrightarrow{\sim} F(Y,X){\quoi}.
\end{align*}
We shall say that $\sigma$ and $\tau$ make $F$ an
\emph{associative} and \emph{strictly commutative} functor if the
following condition holds.

\item[(Ass.s.Com)] For every family $(X_{i})_{i\in I}$
of objects of $C$, denoting by $e:I\rightarrow N(I)$ the canonical
map from $I$ into the free abelian monoid (without identity) generated
by $I$, there exist a map $ℱ:N(I)\rightarrow \Ob C$, isomorphisms
$a_{i}:ℱ(e_{i})\xrightarrow{\sim} X_{i}$, and isomorphisms
$a_{g,h}:ℱ(gh)\xrightarrow{\sim}
F(ℱ(g),ℱ(h))$
such that diagram \eqhyperref[XVIII.eq1.4.1.1]{1.4.1.1} is commutative, as is the
diagram
\begin{equation*}
\vcenter{
\xymatrix@C=2.25cm{
ℱ(gh) \ar@{=}[d] \ar[r]^(.4){a_{g,h}}& F(ℱ(g),ℱ(h))\ar[d]^{\tau}\\
ℱ(hg)    \ar[r]^(.4){a_{h,g}}& F(ℱ(h),ℱ(g)).        
}}\tag{1.4.1.2}
\end{equation*}\etiquette[1.4.1.2]{XVIII.eq1.4.1.2}
\end{itemize}

We shall not need the fact that this axiom is equivalent to the
conjunction of:
\begin{enumerate}
\renewcommand{\labelenumi}{\theenumi)}
\item the pentagon axiom;

\item $\tau_{X,X}:X+X\rightarrow X+X$ is the identity;

\item $\tau_{Y,X}\circ\tau_{X,Y}:X+Y\rightarrow Y+X\rightarrow X+Y$
is the identity;

\item the hexagon axiom: the diagram
$$
\xymatrix{
          &     X + (Y + Z)                 &\\
 (X+Y)+Z  \ar@{-}[ru]_{\sigma} \ar@{-}[d]^{\tau}& & X+(Z+Y)\ar@{-}[lu]^{\tau}\\
Z+(X+Y) \ar@{-}[rd]_{\sigma}  &  & (X+Z) + Y \ar@{-}[u]_{\sigma}\\
          &    (Z+X)+Y    \ar@{-}[ru]_{\tau}                  &
}
$$
is\pageoriginale commutative.
\end{enumerate}

Note that condition 2) is not very often satisfied
in practice (hence the terminology \emph{strictly} commutative).

\hypertarget{sga4.exp.XVIII.definition.1.4.2}{}\begin{enonce*}{Definition 1.4.2}\etiquette[1.4.2]{XVIII.def1.4.2}
A \emph{strictly commutative Picard category $\Pp$}
is a nonempty category all of whose morphisms are
isomorphisms, equipped with a functor $+:\Pp\times
\Pp\rightarrow \Pp:(X,Y)\mapsto X+Y$ and
functorial isomorphisms
\begin{align*}
\sigma &: (X+Y)+Z\xrightarrow{\sim} X+(Y+Z)\\
\tau &: X+Y\longrightarrow Y+X
\end{align*}
making $+$ an associative and strictly commutative functor, and
such that, moreover, for every $X\in \Ob \Pp$, the functor
$Y\mapsto X+Y$ is an equivalence of categories.
\end{enonce*}

The reader will verify:

\hypertarget{sga4.exp.XVIII.lemma.1.4.3}{}\begin{enonce*}{Lemma 1.4.3}\etiquette[1.4.3]{XVIII.lem1.4.3}
If $(X_{i})_{i\in I}$ is a family of objects of a strictly
commutative Picard category $\Pp$, there exist a map
$\Sigma:\ZZ^{(I)}\rightarrow \Ob \Pp$, isomorphisms
$a_{i}:\Sigma(e_{i})\xrightarrow{\sim} X_{i}$, and
$a_{\underline{n},\underline{m}}:\Sigma(\underline{n}+\underline{m})
\xrightarrow{\sim}\Sigma(\underline{n})+\Sigma(\underline{m})$
such that the diagrams of types \eqhyperref[XVIII.eq1.4.1.1]{1.4.1.1} and \eqhyperref[XVIII.eq1.4.1.2]{1.4.1.2}
are commutative. The system
$(\Sigma,(a_{i}),(a_{\underline{n},\underline{m}}))$ is unique up to unique
isomorphism, and is functorial in $(X_{i})_{i\in I}$.
\end{enonce*}

\setcounter{subsubsection}{3}
\subsubsection{}\etiquette[1.4.4]{XVIII.subsubsec1.4.4}
A strictly commutative Picard category admits, up to unique
isomorphism, exactly one \emph{neutral object}, which may here be defined
as a pair $(e,\varphi)$ consisting of an object $e$ and an
isomorphism $\varphi:e+e\xrightarrow{\sim}e$. If $(e,\varphi)$ is a
neutral object, there is a unique isomorphism of functors
$$
\alpha_{s}:e+X\xrightarrow{\sim} X
$$
making the diagram
$$
\xymatrix@C=2cm@R=1.25cm{
e + (e + X) \ar@{<->}[r] \ar[d]^{e+\alpha_s} & (e + e) + X \ar[d]^{\varphi}\\
e + X \ar@{=}[r]                       & e + X .
}
$$
commutative.

One\pageoriginale similarly defines $\alpha_{d}:X+e\xrightarrow{\sim}X$,
and $\varphi$ is a special case of both $\alpha_{s}$ and
$\alpha_{d}$; $\tau$ interchanges $\alpha_{s}$ and $\alpha_{d}$. The group
$\Aut(e)$ is abelian, and for every $X\in \Ob\Pp$, the
functors $X+$ and $+X$ induce the same isomorphism between $\Aut(e)$ and
$\Aut(X)$.

\hypertarget{sga4.exp.XVIII.definition.1.4.5}{}\begin{enonce*}{Definition 1.4.5}\etiquette[1.4.5]{XVIII.def1.4.5}
A \emph{strictly commutative Picard stack $\Pp$ on a}
site $\Ss$ is a stack in groupoids $\Pp$ on
$\Ss$ (\textsc{Giraud} \cite{sga4.exp.XVIII.bib.1}), equipped with a functor
$+:\fprod{\Pp}{\Pp}{\Ss}\rightarrow \Pp$ and
isomorphisms of functors
\begin{align*}
\tau_{x,y} &: x+y\longrightarrow y+x\\
\sigma_{x,y,z} &: (x+y)+z\longrightarrow x+(y+z)
\end{align*}
which, for every $U\in \Ob \Ss$, make $\Pp(U)$ a
strictly commutative Picard category.
\end{enonce*}

In what follows, we shall simply speak of a \emph{Picard stack},
without specifying “strictly commutative.” By
\hyperref[XVIII.subsubsec1.4.4]{1.4.4}, every Picard
stack $\Pp$ admits a global “neutral object”
$e$, and $\sheafAut(e)$ is an abelian sheaf.

\setcounter{subsubsection}{5}
\subsubsection{}\etiquette[1.4.6]{XVIII.subsubsec1.4.6}
An \emph{additive functor} $F:\Pp_{1}\rightarrow
\Pp_{2}$ between Picard stacks on $\Ss$ is an
$\Ss$-functor (necessarily Cartesian) equipped with an
isomorphism of functors
$$
F(x+y)\xrightarrow{\sim} F(x)+F(y)
$$
making the diagrams
$$
\xymatrix@C=2.25cm@R=1.5cm{
F(x+y)  \ar[r] \ar[d]^{F(\tau)}& F(x) + F(y) \ar[d]^{\tau}\\
F(y+x)         \ar[r] & F(y) + F(x)
}
$$
and
$$
\xymatrix@R=1.5cm{
  F((x+y)+z) \ar[r]\ar[d]^{F(\sigma)} & F(x+y)+F(z) \ar[r]&(F(x) + F(y)) + F(z)       \ar[d]^{\sigma}  &
\\
F(x + (y + z))    \ar[r]        & F(x) + F(y +z)  \ar[r] & F(x) + (F(y)
+ F(z)). 
}
$$
commutative.

An\pageoriginale \emph{morphism of additive functors}
$u:F\rightarrow G$ is a
morphism of $\Ss$-functors (automatically an isomorphism)
making the diagram
$$\vcenter{
\xymatrix@C=2cm{
F(x+y)       \ar[r]^{u_{x+y}} \ar[d]       & G(x+y) \ar[d]\\
F(x) + F(y)  \ar[r]^{u_x + u_y}     & G(x) + G(y). 
}}
$$
commutative.

\subsubsection{}\etiquette[1.4.7]{XVIII.subsubsec1.4.7}
If $\Pp_{1}$ and $\Pp_{2}$ are two Picard stacks
on $\Ss$, the Picard stack
$\HOM(\Pp_{1},\Pp_{2})$ is the following Picard stack:
\begin{enumerate}
\renewcommand{\theenumi}{\alph{enumi}}
\renewcommand{\labelenumi}{\theenumi)}
\item if $U\in \Ob \Ss$, its objects over $U$ are the additive
functors from $\Pp_{1}|U$ to $\Pp_{2}|U$; its
morphisms are morphisms of additive functors;

\item the sum of two additive functors $F_{1}$ and
$F_{2}$ is defined by the formula $(F_{1}+F_{2})(x)=F_{1}(x)+F_{2}(x)$;
the structural isomorphism is the one making the diagram
$$
\xymatrix@C=1cm{
\scriptstyle{(F_1 + F_2)(x+y)} \ar[r]^{\rm{\scriptstyle{structural\qquad}}}\ar@{=}[dd]& \scriptstyle{(F_1 + F_2)(x) + (F_1 + F_2)(y)}\ar@{=}[rd]& \\
                                 & &\scriptstyle{F_1(x) + F_2(x) +  F_1(y)+F_2(y)}\\
\scriptstyle{F_1(x+y)+F_2(x+y)}\ar[rd]                    & &\\
&\scriptstyle{F_1(x) + F_1(y) + F_2(x) + F_2(y)}\ar@{<->}[ruu]_{\tau}&
}
$$
commutative;

\item the associativity and commutativity isomorphisms are
defined by means of the analogous isomorphisms in $\Pp_{2}$.
\end{enumerate}

\subsubsection{}\etiquette[1.4.8]{XVIII.subsubsec1.4.8}
If $\Pp_{1}$, $\Pp_{2}$, and $\Pp$ are three
Picard stacks, a biadditive functor from $\Pp_{1}\times
\Pp_{2}$ to $\Pp$ is an $\Ss$-functor $F$
from $\fprod{\Pp_{1}}{\Pp_{2}}{\Ss}$ to $\Pp$,
equipped with isomorphisms of functors
\begin{align*}
F(x_{1}+y_{1},x_{2}) &\xrightarrow{\sim} F(x_{1},
x_{2})+F(y_{1},x_{2})\\
F(x_{1}, x_{2}+y_{2}) &\xrightarrow{\sim} F(x_{1}, x_{2})+F(x_{1},
y_{2}){\quoi},
\end{align*}

% END INCLUDED SGA4_English_current_buildable_source_20260906/source/18/drafts/SGA4_XVIII_LINES_2051_2450_English.texfrag

% BEGIN INCLUDED SGA4_English_current_buildable_source_20260906/source/18/drafts/SGA4_XVIII_LINES_2451_2850_English.texfrag
such that

\begin{enumerate}
\renewcommand{\theenumi}{\alph{enumi}}
\renewcommand{\labelenumi}{\theenumi)}
\item For fixed $x_{1}$ (\resp $x_{2}$), $F(x_{1},\ast)$
(\resp $F(\ast,x_{2})$) satisfies the compatibility conditions of
  \hyperref[XVIII.subsubsec1.4.6]{1.4.6}. 

\item For\pageoriginale $U\in \Ob \Ss$, $x_{1}$, $y_{1}\in
\Ob \Pp_{1}(U)$ and $x_{2}$, $y_{2}\in \Ob \Pp_{2}(U)$,
the diagram
$$
\xymatrix@C=0cm@R=.5cm{
\scriptscriptstyle{F(x_1+y_1, x_2+y_2)}\ar[dd] \ar[r] &\scriptscriptstyle{F(x_1+y_1, x_2) + F(x_1+y_1,y_2)}\ar[rd]\\
 & &\scriptscriptstyle{F(x_1,x_2) + F(y_1,x_2) +F(x_1,y_2)+F(y_1,y_2)}\\
\scriptscriptstyle{F(x_1,x_2+y_2) + F(y_1,x_2+y_2)}\ar[rd] \\
& \scriptscriptstyle{F(x_1,x_2)+F(x_1,y_2) + F(y_1,x_2)+F(y_1,y_2)}\ar@{<->}[ruu]_{\tau}
}
$$
is commutative.
\end{enumerate}

For example, the canonical functor
$$
\HOM(\Pp_{1},\Pp_{2})\times
\Pp_{1}\longrightarrow \Pp_{2}:(F,x)\longmapsto F(x)
$$
is biadditive.

As in \hyperref[XVIII.subsubsec1.4.7]{1.4.7}, one proves that the biadditive functors from
$\Pp_{1}\times \Pp_{2}$ to $\Pp$ form a
Picard stack $\HOM(\Pp_{1}, \Pp_{2};\Pp)$.

We shall see in (\hyperref[XVIII.const1.4.20]{1.4.20}) that there exists a ``$2$-universal'' biadditive functor: $\Pp_{1}\times \Pp_{2}\rightarrow
\Pp_{1}\otimes \Pp_{2}$; more precisely, there exist
a Picard stack $\Pp_{1}\otimes \Pp_{2}$ and a
biadditive functor $\otimes:\Pp_{1}\times
\Pp_{2}\rightarrow \Pp_{1}\otimes \Pp_{2}$ such
that, for every Picard stack $\Pp$, the functor defined by
$\otimes$:
\begin{equation*}
\HOM(\Pp_{1}\otimes \Pp_{2},\Pp)\longrightarrow \HOM (\Pp_{1},\Pp_{2};\Pp)\tag{1.4.8.1}
\end{equation*}\etiquette[1.4.8.1]{XVIII.eq1.4.8.1}%
is an equivalence. This stack is unique up to a unique equivalence
(itself unique up to a unique isomorphism).

\subsubsection{}\etiquette[1.4.9]{XVIII.subsubsec1.4.9}
Let $u:\Ss_{1}\rightarrow \Ss_{2}$ be a morphism of
sites. The direct image $u_{\ast}\Pp$ of a Picard stack
$\Pp$ on $\Ss_{1}$ is the Picard stack on
$\Ss_{2}$ defined by
$$
u_{\ast}\Pp(V)=\Pp(u^{\ast}V){\quoi}.
$$
Its objects over $V\in \Ob \Ss_{2}$ are the objects of
$\Pp$ over $u^{\ast}V$; its morphisms, its addition law, and
its compatibility isomorphisms are those of $\Pp$. The direct
image $u_{\ast}$ is a $2$-functor.

\subsubsection{}\etiquette[1.4.10]{XVIII.subsubsec1.4.10}
A\pageoriginale \emph{Picard prestack} $\Pp$ on a site
$\Ss$ is a prestack $\Pp$ on $\Ss$ (\textsc{Giraud}
\cite{sga4.exp.XVIII.bib.1}), equipped with a functor
$+:\fprod{\Pp}{\Pp}{\Ss}\rightarrow \Pp$ and
associativity and commutativity isomorphisms $\sigma$ and
$\tau$, such that, for every $U\in \Ob \Ss$, $\Pp(U)$
is a (strictly commutative) Picard category. If
$\Pp_{1}$ and $\Pp_{2}$ are two Picard prestacks,
one defines in the evident way (\cf \hyperref[XVIII.subsubsec1.4.6]{1.4.6},
\hyperref[XVIII.subsubsec1.4.7]{1.4.7}) the \emph{additive functors}
from $\Pp_{1}$ to $\Pp_{2}$ and the Picard prestack
$\HOM(\Pp_{1}, \Pp_{2})$ that they form. If
$\Pp$ is a Picard prestack, and $j:\Pp\rightarrow
a\Pp$ is the stack that it generates, there exists, up to an
essentially unique equivalence, one and only one pair consisting of a
Picard-stack structure on $a\Pp$ and an additive-functor structure
on $j$; equipped with these data, the pair $(j,a\Pp)$
is called \emph{the Picard stack generated by} $\Pp$. For
every Picard stack $\Pp_{1}$, one has an equivalence
\begin{equation*}
\HOM(a\Pp,\Pp_{1})\xrightarrow{\sim} \HOM(\Pp,\Pp_{1}){\quoi}.\tag{1.4.10.1}
\end{equation*}\etiquette[1.4.10.1]{XVIII.eq1.4.10.1}

\subsubsection{}\etiquette[1.4.11]{XVIII.subsubsec1.4.11}
We shall denote by $C^{[-1,0]}(\Ss)$ the category of complexes
of abelian sheaves $K$ on $\Ss$ such that $K^{i}=0$ for
$i\not\in [-1,0]$. To every complex $K\in \Ob C^{[-1,0]}(\Ss)$
$$
K:d:K^{-1}\longrightarrow K^{0}
$$
is associated the following Picard prestack $\pch(K)$:

\begin{enumerate}
\renewcommand{\theenumi}{\Roman{enumi}}
\renewcommand{\labelenumi}{(\theenumi)}
\item For $U\in \Ob \Ss$, one has $\Ob \pch(K)(U)=K^{0}(U)$.

\item If $x$, $y\in K^{0}(U)$, an arrow from $x$ to y is an element
$f$ of $K^{-1}(U)$ such that $df=y-x$.

\item The composition law for arrows is the addition law in
$K^{-1}(U)$.

\item The functor $+$ is given by the addition law in $K^{0}(U)$
and $K^{-1}(U)$.

\item The associativity and commutativity isomorphisms are
provided by the zero element of $K^{-1}(U)$.
\end{enumerate}

We shall denote by $\ch(K)$ the Picard stack generated
(\hyperref[XVIII.subsubsec1.4.10]{1.4.10}) by the Picard
prestack $\pch(K)$. The sheaf $ℋ^{0}(K)$ is interpreted
as the following associated sheaf
\begin{align*}
ℋ^{0}(K) & \simeq a\; \text{\rm ($U\longmapsto$ group of isomorphism
classes}\tag{1.4.11.1}\\
&\qquad\qquad \text{\rm of objects of $\ch(K)(U)$)}{\quoi}.
\end{align*}\etiquette[1.4.11.1]{XVIII.eq1.4.11.1}

On the\pageoriginale other hand,
\begin{equation*}
ℋ^{-1}(K)\simeq \sheafAut(e)\tag{1.4.11.2}
\end{equation*}\etiquette[1.4.11.2]{XVIII.eq1.4.11.2}

\subsubsection{}\etiquette[1.4.12]{XVIII.subsubsec1.4.12}
Every morphism of complexes $f:K\rightarrow L$ induces an additive
functor $\pch(f):\pch(K)\rightarrow \pch(L)$, hence an additive functor
$$
\ch(f):\ch(K)\longrightarrow \ch(L){\quoi}.
$$

It follows from \eqhyperref[XVIII.eq1.4.11.1]{1.4.11.1} and \eqhyperref[XVIII.eq1.4.11.2]{1.4.11.2} that $\ch(f)$ is an
equivalence if and only if $f$ is a quasi-isomorphism. From
\eqhyperref[XVIII.eq1.4.10.1]{1.4.10.1} one deduces that, for two morphisms $f$, $g:K\rightarrow L$,
one has
$$
\Hom(\ch(f),\ch(g))\xrightarrow{\sim} \Hom(\pch(f),\pch(g)){\quoi}.
$$

A morphism of \emph{functors} $h:\pch(f)\rightarrow \pch(g)$ is a
morphism of sheaves $h:K^{0}\rightarrow L^{-1}$ such that
\begin{equation*}
g(x)-f(x)=dh(x)\tag{1.4.12.1}
\end{equation*}\etiquette[1.4.12.1]{XVIII.eq1.4.12.1}%
and such that, for every triple $(x,y,u)$ with $y-x=du$, one
has
\begin{equation*}
h(y)+f(u)=g(u)+h(x){\quoi}\tag{1.4.12.2}
\end{equation*}\etiquette[1.4.12.2]{XVIII.eq1.4.12.2}

That $h$ is a morphism of additive functors means that
\begin{equation*}
h(x+y)=h(x)+h(y){\quoi}.\tag{1.4.12.3}
\end{equation*}\etiquette[1.4.12.3]{XVIII.eq1.4.12.3}%
Condition \eqhyperref[XVIII.eq1.4.12.3]{1.4.12.3} makes it possible to rewrite \eqhyperref[XVIII.eq1.4.12.2]{1.4.12.2} as
\begin{equation*}
g(u)-f(u)=h\sisi{du}{(du)} ;\tag{$1.4.12.2'$}
\NextLinkTarget{sga4.exp.XVIII.equation.1.4.12.2-prime}\label{1.4.12.2'}
\end{equation*}
conditions \eqhyperref[XVIII.eq1.4.12.1]{1.4.12.1} through \eqhyperref[XVIII.eq1.4.12.3]{1.4.12.3} therefore mean that $h$ is a
homotopy from $f$ to $g$, and
\begin{equation*}
\Hom(\ch(f),\ch(g))\simeq \{H:\;\text{\rm homotopy } K\longrightarrow
L| g-f=dH+Hd\}\tag{1.4.12.4}
\end{equation*}\etiquette[1.4.12.4]{XVIII.eq1.4.12.4}

\begin{enonce*}{Lemma 1.4.13}\etiquette[1.4.13]{XVIII.lem1.4.13}
\begin{enumerate}
\renewcommand{\theenumi}{\Roman{enumi}}
\renewcommand{\labelenumi}{(\theenumi)}
\item For every Picard stack $\Pp$, there exists a complex
$K$ such that $\Pp=\ch(K)$.

\item For every additive functor $F:\ch(K)\rightarrow \ch(L)$, there
exist a quasi-isomorphism $k:K'\rightarrow K$ and a morphism
$\ell:K'\rightarrow L$ such that $F\simeq \ch(\ell)\ch(k)^{-1}$.
\end{enumerate}
\end{enonce*}

Let us prove\pageoriginale (I). Let $(k_{i},U_{i})_{i\in I}$ be a family
such that
\begin{enumerate}
\renewcommand{\theenumi}{\alph{enumi}}
\renewcommand{\labelenumi}{(\theenumi)}
\item $k_{i}\in \Ob \Pp(U_{i})$;

\item every object of $\Pp$ is locally isomorphic to an inverse
image of one of the $k_{i}$.

Set $K^{0}=\displaystyle{\mathop{\oplus}_{i\in
I}}\ZZ_{U_{i}}$. It follows from \hyperref[XVIII.lem1.4.3]{1.4.3} that there exists an additive
functor $F$ from $\ch(0\rightarrow K^{0})$ to $\Pp$ taking
the basis $e_{i}$ of $\ZZ_{U_{i}}$ to $k_{i}$. Let $K^{-1}$ be the
sheaf of pairs consisting of a local section $x$ of $K^{0}$ and
an isomorphism $t:F(0)\xrightarrow{\sim} F(x)$. Addition on
$K^{-1}$ is defined by $(x_{1},
t_{1})+(x_{2},t_{2})=(x_{1}+x_{2},t)$, where $t$ makes the
diagram
$$
\xymatrix@C=1.5cm{
F(0) + F(0)    \ar[r]^{t_1+t_2} \ar@{<->}[d]   & F(x_1) + F(x_2)\ar@{<->}[d]\\
F(0+0)          \ar[r]^{t}         & F(x_1+x_2).
}
$$
commutative. Define $d:K^{-1}\rightarrow K^{0}$ by $d(x,t)=x$; $d$ is
additive.

If $y-x=dt$, there exists a unique morphism $F(t)$ making the
diagram
$$
\xymatrix@C=1.75cm{
F(x)           \ar[r]^{F(t)}   \ar@{<->}[d]& F(y) \ar@{<->}[d]\\
F(0)+F(x)      \ar[r]^{t+F(x)} & F(y-x)+F(x)
}
$$
commutative, and one verifies that this construction defines an equivalence
$$
F:\ch(K)\longrightarrow \Pp{\quoi}.
$$

Let us prove (II). Thus let $F:\ch(K)\rightarrow \ch(L)$ be an additive
functor. There then exists a family $(k_{i},\ell_{i}, \alpha_{i},
U_{i})_{i\in I}$ such that

\item $k_{i}\in K^{0}(U_{i})$, $\ell_{i}\in L^{0}(U_{i})$, and
$\alpha_{i}$ is a morphism from $F(k_{i})$ to $\ell_{i}$;

\item the homomorphism $k^{0}:\displaystyle{\mathop{\oplus}_{i\in
I}}\ZZ_{U_{i}}\rightarrow K^{0}$ with coordinates $k_{i}$ is an
epimorphism.

Set ${K'}^{0}=\displaystyle{\mathop{\oplus}_{i\in
I}}\ZZ_{U_{i}}$, and let $\ell^{0}:{K'}^{0}\rightarrow L^{0}$ be
the homomorphism with coordinates $\ell_{i}$. Since the functor $F$ is
additive, there exists a unique family of isomorphisms
$\alpha_{x}:Fk^{0}(x)\xrightarrow{\sim}\ell^{0}(x)$ ($x$ a local
section of ${K'}^{0}$) such that

\item for $e_{i}$ the \hyperref[XVIII.sec1]{section 1} of $\ZZ_{U_{i}}$, one has
$\alpha_{e_{i}}=\alpha_{i}$;

\item The\pageoriginale diagrams
$$
\xymatrix@C=2cm{
Fk^0(x+y)  \ar[r]^{\alpha_{x+y}} \ar@{<->}[d]& \ell^0(x+y)
\ar@{=}[d]\\
Fk^0(x)+ Fk^0(y)    \ar[r]^{\alpha_{x} + \alpha_{y}}     & \ell^0(x)+\ell^0(y)
}
$$
are commutative.

Let ${K'}^{-1}=\fprod{K^{-1}}{{K'}^{0}}{K^{0}}$; by (d), the
evident arrow is a quasi-isomorphism of complexes
$k:K'\rightarrow K$. If $x$, $y\in {K'}^{0}(U)$, and if $t\in
{K'}^{-1}(U)$ satisfies $y-x=dt$, then there exists a
unique section $\ell^{-1}(t;x,y)$ of $L^{-1}(U)$ defining the
morphism that makes the diagram

\item $\begin{array}[t]{c}
\xymatrix@C=2.75cm@R=1.75cm{
&F(k^0(x))  \ar[r]^{F(k^{-1}(t))} \ar[d]^{\alpha_x}&
  F(k^0(y))\ar[d]^{\alpha_y} \\ 
&\ell^0(x)     \ar[r]^{\ell^{-1}(t;x,y)}& \ell^0(y).
}
\end{array}$


commutative. The fact that $F$ is a functor gives

\item $\ell^{-1}(t;x,y)+\ell^{-1}(t';y,z)=\ell^{-1}(t+t';x,z)$.

The functoriality of the additivity isomorphism of $F$ gives

\item $\ell^{-1}(t_{1}+t_{2};x_{1}+x_{2},
y_{1}+y_{2})=\ell^{-1}(t_{1};x_{1},y_{1})+\ell^{-1}(t_{2};x_{2},y_{2})$.

From (h), one obtains
\begin{align*}
&\ell^{-1}(0;x,x)+\ell^{-1}(0;x,x)=\ell^{-1}(0;x,x)\qquad \text{\ie}\\
& \ell^{-1}(0;x,x)=0{\quoi}.
\end{align*}
It then follows from (i) that $\ell^{-1}(t;x,y)$ depends only on $t$:
$$
\ell^{-1}(t;x,y)=\ell^{-1}(t)
$$
and that $\ell^{-1}(t)$ is additive in $t$. One concludes that
$\ell=(\ell^{-1}, \ell^{0})$ is a morphism of complexes and,
by (f) and (g), the arrows $\alpha$ form an isomorphism of
additive functors $F\circ \ch(k)=\ch(\ell)$.
\end{enumerate}

\setcounter{subsubsection}{13}
\subsubsection{}\etiquette[1.4.14]{XVIII.subsubsec1.4.14}
Let $\Ch^{\flat}(\Ss)$ be the category whose objects
are the small Picard stacks on $\Ss$, and whose
arrows are the isomorphism classes of additive\pageoriginale
functors. The construction $\ch$ defines a functor from
$C^{[-1,0]}(\Ss)$ (\hyperref[XVIII.subsubsec1.4.11]{1.4.11}) to
$\Ch^{\flat}(\Ss)$. Let $D^{[-1,0]}(\Ss)$ be the
subcategory of the derived category of $\Ss$ formed by the
complexes $K$ such that $H^{i}(K)=0$ for $i\neq 0$ or $-1$. It follows
from \hyperref[XVIII.subsubsec1.4.12]{1.4.12} and \hyperref[XVIII.lem1.4.13]{1.4.13} that

\begin{enonce*}{Proposition 1.4.15}\etiquette[1.4.15]{XVIII.prop1.4.15}
The functor $\ch$ induces an equivalence of categories
$$
\ch:D^{[-1,0]}(\Ss)\xrightarrow{\sim}
\Ch^{\flat}(\Ss){\quoi}.
$$
\end{enonce*}

We shall denote by $\flat$ the inverse equivalence of $\ch$. For every
Picard stack $\Pp$ on $\Ss$,
$\Pp^{\flat}\in \Ob D^{[-1,0]}(\Ss)$ determines
$\Pp$ up to the isomorphism class of an equivalence.

\begin{enonce*}{Lemma 1.4.16}\etiquette[1.4.16]{XVIII.lem1.4.16}
Let $L\in \Ob C^{[-1,0]}(\Ss)$ be such that $L^{-1}$ is
injective.
\begin{enumerate}
\renewcommand{\theenumi}{\Roman{enumi}}
\renewcommand{\labelenumi}{(\theenumi)}
\item $\pch(L)$ is already a stack.

\item For $K\in \Ob C^{[-1,0]}(\Ss)$, every additive functor
$F:\ch(K)\rightarrow \ch(L)$ is isomorphic to a functor $\ch(f)$ for
$f$ a morphism from $K$ to $L$.
\end{enumerate}
\end{enonce*}

\begin{enumerate}
\renewcommand{\theenumi}{\Roman{enumi}}
\renewcommand{\labelenumi}{(\theenumi)}
\item For any $U\in \Ob \Ss$ and $x\in \Ob \ch(L)(U)$,
the sheaf of pairs consisting of a local section $\ell$ of
$L^{0}|U$ and an isomorphism between $x$ and $\ell$ is a principal
homogeneous space under $L^{-1}$; it therefore has a section, and
$j:\pch(L)\rightarrow \ch(L)$ is surjective on the isomorphism
classes of objects, hence is an equivalence.

\item Let $L'$ be a complex of injectives, equipped with a quasi-isomorphism
$\varphi:L{\xrightarrow[\approx]{}}L'$, such that $\varphi^{i}$ is an
isomorphism for $i\leq -1$. By \hyperref[XVIII.lem1.4.13]{1.4.13} and
\cite{sga4.exp.XVIII.bib.9}, there exists a diagram commutative up to homotopy
$$
\xymatrix@R=1cm@C=1.5cm{
L  \ar[r]^{\varphi}_{\approx} & L'\\
K' \ar[u]^{f_0}  \ar[r]^{\Psi}_{\approx}   & K\ar[u]_{f_1}
}
$$
such that $F\simeq \ch(f_{0})\ch(\Psi)^{-1}$. Moreover, $K=\tau_{\leq
0}K$ and $L=\tau_{\leq 0}L'$, so that $f_{1}$ factors through
$f:K\rightarrow L$, and $F\simeq \ch(f)$.
\end{enumerate}

The lemma refines \hyperref[XVIII.prop1.4.15]{1.4.15} by the following consequence
of \hyperref[XVIII.subsubsec1.4.12]{1.4.12}.

\begin{enonce*}{Corollary 1.4.17}\etiquette[1.4.17]{XVIII.cor1.4.17}
The\pageoriginale construction $\ch$ defines an equivalence of $2$-categories between
\begin{enumerate}
\renewcommand{\theenumi}{\alph{enumi}}
\renewcommand{\labelenumi}{\theenumi)}
\item the $2$-category of Picard stacks on $\Ss$.

\item the $2$-category whose objects and $1$-arrows are the
objects and arrows of the full subcategory of
$C^{[-1,0]}(\Ss)$ formed by the complexes $L$ with $L^{-1}$
injective, and whose $2$-arrows are homotopies between arrows:
$\Hom(f,g)=\{H|g-f=\text{\rm dH}+\text{\rm Hd}\}$.
\end{enumerate}
\end{enonce*}

\begin{enonce*}[remark]{Construction 1.4.18}\etiquette[1.4.18]{XVIII.const1.4.18}
One has
$$
\HOM(\Pp_{1}, \Pp_{2})^{\flat}\simeq \tau_{\leq
0}R\Hhom (\Pp^{\flat}_{1}, \Pp^{\flat}_{2}){\quoi}.
$$
\end{enonce*}

Let $K$, $L\in \Ob C^{[-1,0]}(\Ss)$, with $L^{-1}$ injective. By
\hyperref[XVIII.lem1.4.16]{1.4.16} (II) and \hyperref[XVIII.subsubsec1.4.12]{1.4.12}, one has
\begin{equation*}
\ch(\tau_{\leq
0}\Hhom(K,L))\xrightarrow{\sim}\HOM(\ch(K),\ch(L))\tag{1.4.18.1}
\end{equation*}\etiquette[1.4.18.1]{XVIII.eq1.4.18.1}%
and one deduces Construction \hyperref[XVIII.const1.4.18]{1.4.18} from this functorially.

\begin{enonce*}[remark]{Construction 1.4.19}\etiquette[1.4.19]{XVIII.const1.4.19}
For $f:\Ss_{1}\to\Ss_{2}$ a morphism of sites, one has
$$
(f_{\ast}\Pp)^{\flat}\simeq \tau_{\leq
0}\,Rf_{\ast}(\Pp^{\flat}){\quoi}.
$$
\end{enonce*}

Let $L\in \Ob C^{[-1,0]}$ with $L^{-1}$ injective. By
\hyperref[XVIII.lem1.4.16]{1.4.16} (I), one has

% END INCLUDED SGA4_English_current_buildable_source_20260906/source/18/drafts/SGA4_XVIII_LINES_2451_2850_English.texfrag

% BEGIN INCLUDED SGA4_English_current_buildable_source_20260906/source/18/drafts/SGA4_XVIII_LINES_2851_3250_English.texfrag
\begin{equation*}
\ch(f_{\ast}L)\xrightarrow{\sim} f_{\ast}\ch(L)\tag{1.4.19.1}
\end{equation*}\etiquette[1.4.19.1]{XVIII.eq1.4.19.1}%
and from this one deduces a functorial construction \hyperref[XVIII.const1.4.19]{1.4.19}.

We shall not need to use the following

\begin{enonce*}[remark]{Construction 1.4.20}\etiquette[1.4.20]{XVIII.const1.4.20}
The tensor product promised in \hyperref[XVIII.subsubsec1.4.8]{1.4.8} exists, and
$$
(\Pp_{1}\otimes \Pp_{2})^{\flat}\simeq \tau_{\geq
-1}\Pp^{\flat}_{1}\mathop{\otimes}^{\LL}\Pp^{\flat}_{2}{\quoi}.
$$
\end{enonce*}

One checks, as in \hyperref[XVIII.subsubsec1.4.12]{1.4.12} and
\hyperref[XVIII.lem1.4.13]{1.4.13} (II), that for all
$K_{1}$, $K_{2}$, and $L$ in $C^{[-1,0]}(\Ss)$, one has
\begin{enumerate}
\renewcommand{\theenumi}{\alph{enumi}}
\renewcommand{\labelenumi}{(\theenumi)}
\item Every morphism $f\in \Hom(K_{1}\otimes K_{2},L)\simeq
\Hom(\tau_{\geq-1}(K_{1}\otimes K_{2}),L)$ defines a biadditive
functor $\ch(f)\in \Hom(\ch(K_{1})$, $\ch(K_{2})$; $\ch(L)$); moreover,
morphisms of functors identify with homotopies;

\item For every biadditive functor $F:\ch(K_{1})$,
$\ch(K_{2})\rightarrow \ch(L)$, there exist quasi-isomorphisms
$k_{i}:K'_{i}\rightarrow K_{i}$ and $f:K'_{1}\otimes
K'_{2}\rightarrow L$ such that $F\simeq \ch(f)\circ
(\ch(k_{1})^{-1},\ch(k_{2})^{-1})$.
\end{enumerate}

If\pageoriginale now $K^0_{1}$ or $K^0_{2}$ is flat,
as is ${K'}^0_{1}$ or ${K'}^0_{2}$, then
$k_{1}\otimes k_{2}:\tau_{\geq -1}{K'}_{1}\otimes {K'}_{2}\rightarrow
\tau_{\geq -1}K_{1}\otimes K_{2}$ is a quasi-isomorphism, and it
follows that $\ch(\tau_{\geq -1}(K_{1}\otimes K_{2}))$ satisfies the
universal property \eqhyperref[XVIII.eq1.4.8.1]{1.4.8.1}, whence the existence of $\otimes$ and
of a functorial construction \hyperref[XVIII.const1.4.20]{1.4.20}.

\setcounter{subsubsection}{20}
\subsubsection{}\etiquette[1.4.21]{XVIII.subsubsec1.4.21}
Let $G$ be an abelian sheaf, and let $G[1]$ be the complex with value
$G$ placed in degree $-1$. The prestack $\pch(G[1])$ identifies
with the prestack of trivial torsors under $G$, and hence $\ch(G[1])$
``is none other than'' the stack of torsors under $G$, with its usual
addition law.

\subsubsection{}\etiquette[1.4.22]{XVIII.subsubsec1.4.22}
Let $K\in C^{[-1,0]}(\Ss)$, and let $G$ be an abelian sheaf. The
extensions $E$ of $K$ by $G$ form a Picard stack, $\EXT(K,G)$, under
% SGA-SRC-001136
Baer addition.
$$
E:0\rightarrow
G[0]\xrightarrow{\alpha}E\xrightarrow{\beta}K\rightarrow 0{\quoi}.
$$

To every extension $E$ one associates an additive functor from $\ch(K)$
to $\ch(G[1])$ which, when applied to a local section $x$ of $K^{0}$,
gives the torsor $\beta^{-1}(x)$. The reader will verify the following.

\begin{enonce*}{Proposition 1.4.23}\etiquette[1.4.23]{XVIII.prop1.4.23}
The construction outlined above is an equivalence of Picard stacks
$$
\EXT(K,G)\xrightarrow{\sim} \HOM(\ch(K),\ch(G[1])){\quoi}.
$$
\end{enonce*}

\setcounter{subsubsection}{23}
\subsubsection{}\etiquette[1.4.24]{XVIII.subsubsec1.4.24}
Let $F$ be an object of the site $\Ss$. We shall also write $F$ for
the sheaf represented by $F$, and write $\ZZ^{(F)}$ for the abelian
sheaf it generates. Let $C$ be a Picard stack on $\Ss$. One checks
easily (\cf \hyperref[XVIII.lem1.4.3]{1.4.3}) that giving any one of the following amounts to the
same thing:
\begin{enumerate}
\renewcommand{\theenumi}{\alph{enumi}}
\renewcommand{\labelenumi}{\theenumi)}
\item an additive functor $H:\ch(\ZZ^{(F)})\rightarrow C$;

\item a morphism of \emph{stacks}:

(stack of local sections of $F$, with only identity morphisms)
$\rightarrow C$;

\item an object of $C(F)$.
\end{enumerate}

Suppose (in order to be able to apply Definition
\hyperref[XVIII.subsubsec1.4.9]{1.4.9}) that fiber products exist in $\Ss$, and let $f$ be the
canonical morphism of sites $f:\Ss/F\rightarrow \Ss$. The preceding
construction gives an equivalence
\begin{equation*}
\HOM(\ch(\ZZ^{(F)}),C)\xrightarrow{\sim}
f_{\ast}f^{\ast}C{\quoi}.\tag{1.4.24.1}
\end{equation*}\etiquette[1.4.24.1]{XVIII.eq1.4.24.1}

The isomorphism\pageoriginale obtained from \eqhyperref[XVIII.eq1.4.24.1]{1.4.24.1} by the construction
$$
\tau_{\leq 0} R\Hhom(\ZZ^{(F)},C^{\flat})\xrightarrow{\sim}
\tau_{\leq 0} Rf_{\ast} f^{\ast}C^{\flat}{\quoi},
$$
may also be deduced from the more general isomorphism
\begin{equation*}
R\Hhom(\ZZ^{(F)},K)\xrightarrow{\sim}
Rf_{\ast}f^{\ast}K\quad (K\in \Ob D^{+}(\Ss)){\quoi}.\tag{1.4.24.2}
\end{equation*}\etiquette[1.4.24.2]{XVIII.eq1.4.24.2}

\subsection{The universal coefficient formula}\etiquette[1.5]{XVIII.subsec1.5}

\subsubsection{}\etiquette[1.5.1]{XVIII.subsubsec1.5.1}
Let $f:X\rightarrow S$ be a smooth projective curve over $S$. We
denote by $\PIC(X/S)$ the Picard stack on the big fppf site of $S$
(\hyperrefext[XVII][XVII.subsec0.10]{0.10}), namely the direct image (\hyperref[XVIII.subsubsec1.4.9]{1.4.9}) of the stack of
invertible sheaves on $X$. We have
\begin{align*}
\PIC(X/S)&= f_{\ast}\ch\GG_{m}[1]\tag{1.5.1.1}\\
\PIC(X/S)^{\flat} &=\tau_{\leq 0}
Rf_{\ast}(\GG_{m}[1]){\quoi}.\tag{1.5.1.2}
\end{align*}\etiquette[1.5.1.1]{XVIII.eq1.5.1.1}
\etiquette[1.5.1.2]{XVIII.eq1.5.1.2}

Each section $t$ of $X$ defines an invertible sheaf $\Oo(t)$ on $X$.
With the notation of \hyperref[XVIII.subsubsec1.4.24]{1.4.24} (\cf
\hyperref[XVIII.subsubsec1.4.24]{1.4.24} $a\leftrightarrow b$), this construction, being compatible with
every base change, defines a morphism of Picard stacks on $S$
\begin{equation*}
\ch \ZZ^{(X)}\longrightarrow \PIC(X/S){\quoi}.\tag{1.5.1.3}
\end{equation*}\etiquette[1.5.1.3]{XVIII.eq1.5.1.3}

If $C$ is any Picard stack on $S$, the morphisms
\eqhyperref[XVIII.eq1.5.1.3]{1.5.1.3} and \eqhyperref[XVIII.eq1.4.24.1]{1.4.24.1} define an additive functor
\begin{equation*}
\HOM(\PIC(X/S),C)\longrightarrow\HOM(\ch(\ZZ^{(X)}),
C)\xrightarrow{\sim} f_{\ast}f^{\ast}C{\quoi}.\tag{1.5.1.4}
\end{equation*}\etiquette[1.5.1.4]{XVIII.eq1.5.1.4}

Upon applying the construction $\flat$, the functor \eqhyperref[XVIII.eq1.5.1.3]{1.5.1.3}
defines a morphism\nde{Here the shift functor
$[1]$ is to be understood to have higher precedence than the truncation $\tau_{\leq 0}$.}
\begin{equation*}
\ZZ^{(X)}\longrightarrow \tau_{\leq 0}
Rf_{\ast}\GG_{m}[1]\longrightarrow
Rf_{\ast}\GG_{m}[1]{\quoi},\tag{1.5.1.5}
\end{equation*}\etiquette[1.5.1.5]{XVIII.eq1.5.1.5}%
and the analogue of the functor \eqhyperref[XVIII.eq1.5.1.4]{1.5.1.4} is, by \eqhyperref[XVIII.eq1.4.24.2]{1.4.24.2}, a morphism
\begin{equation*}
R\Hhom(\tau_{\leq 0} Rf_{\ast}(\GG_{m}[1]),K)\longrightarrow
Rf_{\ast} f^{\ast}K{\quoi}.\tag{1.5.1.6}
\end{equation*}\etiquette[1.5.1.6]{XVIII.eq1.5.1.6}

If $C$ is the stack $\ch(G[1])$ of torsors under a group $G$
satisfying (\hyperref[XVIII.par1.3.1.1]{1.3.1.1}), Theorem \hyperref[XVIII.thm1.3.10]{1.3.10}
asserts that the functor
\eqhyperref[XVIII.eq1.5.1.4]{1.5.1.4} is an equivalence,\pageoriginale the inverse equivalence being the one that
associates with a torsor $K$ under $G_{X}$ the additive functor
$\langle \ast,K]$. The assertion that \eqhyperref[XVIII.eq1.5.1.4]{1.5.1.4} is an equivalence amounts to
saying that the morphism deduced from \eqhyperref[XVIII.eq1.5.1.4]{1.5.1.4} or \eqhyperref[XVIII.eq1.5.1.6]{1.5.1.6}
\begin{equation*}
\tau_{\leq 1}R\Hhom(\tau_{\leq
0}Rf_{\ast}(\GG_{m}[1]),G)\longrightarrow \tau_{\leq 1}
Rf_{\ast}f^{\ast}G\tag{1.5.1.7}
\end{equation*}\etiquette[1.5.1.7]{XVIII.eq1.5.1.7}%
is an isomorphism.

\begin{enonce*}{Theorem 1.5.2}\etiquette[1.5.2]{XVIII.thm1.5.2}
\emph{(universal coefficient formula).}

Let $f:X\rightarrow S$ be a smooth projective curve over $S$, let $C$
be a Picard stack on $S_{\text{\rm fppf}}$, and let $K=C^{\flat}$
(\hyperref[XVIII.prop1.4.15]{1.4.15}). Suppose that $K$ is locally isomorphic (in
$D^{+}(S_{\text{\rm fppf}})$) to complexes of the form
$G_{-1}\rightarrow G_{0}$, where the sheaves $G_{i}$ satisfy
(\hyperref[XVIII.par1.3.1.1]{1.3.1.1}). Then:
\begin{enumerate}
\renewcommand{\theenumi}{\Roman{enumi}}
\renewcommand{\labelenumi}{(\theenumi)}
\item The functor \eqhyperref[XVIII.eq1.5.1.4]{1.5.1.4}
\begin{equation*}
\HOM(\PIC(X/S),C)\longrightarrow f_{\ast}f^{\ast}C\tag{1.5.2.1}
\end{equation*}\etiquette[1.5.2.1]{XVIII.eq1.5.2.1}%
is an equivalence of Picard stacks.

\item The morphism deduced from \eqhyperref[XVIII.eq1.5.1.6]{1.5.1.6} or \eqhyperref[XVIII.eq1.5.1.4]{1.5.1.4}
\begin{equation*}
\tau_{\leq 0}R\Hhom(\tau_{\leq 0}
Rf_{\ast}(\GG_{m}[1]),K)\longrightarrow \tau_{\leq 0}
Rf_{\ast}f^{\ast}K\tag{1.5.2.2}
\end{equation*}\etiquette[1.5.2.2]{XVIII.eq1.5.2.2}%
is an isomorphism.
\end{enumerate}
\end{enonce*}

This theorem makes $\tau_{\leq 0}Rf_{\ast}(\GG_{m}[1])$ play the
role played by homology in the classical universal coefficient
formula.

It is clear that, for every stack $C$, one has $(I)\leftrightarrow
(II)$, and both $(I)$ and $(II)$ are local in nature on $S$. If $C$
is of the form $\ch(G[1])$, then the theorem follows from
\hyperref[XVIII.thm1.3.10]{1.3.10}, as noted above. It remains to show that, if $K$ is of the form
$$
K:d:G_{-1}\longrightarrow G_{0}{\quoi},
$$
and if \eqhyperref[XVIII.eq1.5.2.2]{1.5.2.2} is an isomorphism for $G_{-1}[1]$ and $G_{0}[1]$,
then \eqhyperref[XVIII.eq1.5.2.2]{1.5.2.2} is an isomorphism for $K$. To verify this, it suffices to
compare the cohomology long exact sequence
$$
0\longrightarrow R^{-1}f_{\ast}K\longrightarrow
f_{\ast}G_{-1}\longrightarrow f_{\ast}G_{0}\longrightarrow
R^{0}f_{\ast}K\longrightarrow R^{1} f_{\ast}G_{-1}\longrightarrow
R^{1}f_{\ast}G_{0}
$$
with\pageoriginale the analogous long exact sequence for the left-hand
side of \eqhyperref[XVIII.eq1.5.2.2]{1.5.2.2}, and to apply the five lemma.

\setcounter{subsubsection}{2}
\subsubsection{}\etiquette[1.5.3]{XVIII.subsubsec1.5.3}
Let $u$ be a flat morphism of smooth projective curves over $S$
$$
\xymatrix@=1.25cm{
X              \ar[rr]^{u} \ar[rd]_{f_1}&       &Y \ar[ld]^{f_2}\\
                                       &  S    &           .
}
$$
The diagram of additive functors
$$\vcenter{\etiquette[1.5.3.1]{XVIII.eq1.5.3.1}
\xymatrix{
\ZZ^{(X)}  \ar[r]\ar[d]^{u} & \PIC(X/S) \ar[d]^{N_{X/Y}}\\
\ZZ^{(Y)}  \ar[r] & \PIC(Y/S)
}}\leqno{(1.5.3.1)}
$$
is then essentially commutative.

For every Picard stack $C$ on $S$, the diagram
$$\vcenter{\etiquette[1.5.3.2]{XVIII.eq1.5.3.2}
\xymatrix@R=1.75cm@C=2cm{
\HOM(\PIC(Y/S),C)         \ar[r] \ar[d]^{\HOM(N_{X/Y},C)}&
f_{2\ast}f^{\ast}_2 C \ar[d]^{u^{\ast}}\\
\HOM(\PIC(X/S), C)        \ar[r] & f_{1\ast}f^{\ast}_1 C
}}\leqno{(1.5.3.2)}
$$
is therefore essentially commutative, and for $K\in \Ob D^{+}(S)$,
the diagram
$$\vcenter{\etiquette[1.5.3.3]{XVIII.eq1.5.3.3}
\xymatrix@C=2cm@R=1.5cm{
R\Hhom(\tau_{\leq 0} Rf_{2\ast}\GG_m[1],K) \ar[r]\ar[d]^{R\Hhom(N_{X/Y},K)}&
Rf_{2\ast}f^{\ast}_2 K  \ar[d]^{u^{\ast}}\\
R\Hhom(\tau_{\leq 0} Rf_{1\ast}\GG_m[1],K) \ar[r]&
Rf_{1\ast}f^{\ast}_{1} K
}}\leqno{(1.5.3.3)}
$$
is commutative.

Now let $G$ be an abelian sheaf on $S$ satisfying
(\hyperref[XVIII.par1.3.1.1]{1.3.1.1}). The trace \hyperrefext[XVII][XVII.subsec6.3]{6.3} gives us an additive functor
\begin{equation*}
\Tr_{u}:f_{1\ast}f^{\ast}_{1}\ch(G[1])\longrightarrow f_{2\ast}f^{\ast}_{2}\ch(G[1])\tag{1.5.3.4}
\end{equation*}\etiquette[1.5.3.4]{XVIII.eq1.5.3.4}%
and\pageoriginale \hyperref[XVIII.par1.3.8.6]{1.3.8.6} gives us an isomorphism making the diagram
essentially commutative
$$\vcenter{\etiquette[1.5.3.5]{XVIII.eq1.5.3.5}
\xymatrix@R=1.75cm{
\HOM(\PIC(X/S), \ch(G[1])) \ar[d]^{\HOM(u^{\ast}, \ch(G[1]))}&  f_{1\ast}f^{\ast}_1
\ch(G[1]) \ar[l] \ar[d]^{\Tr_u}\\
\HOM(\PIC(Y/S), \ch(G[1])) & f_{2\ast}f^{\ast}_2 \ch(G[1]).\ar[l]
}}\leqno{(1.5.3.5)}
$$

The diagram
$$\vcenter{\etiquette[1.5.3.6]{XVIII.eq1.5.3.6}
\xymatrix@R=1.75cm{
\tau_{\leq 1} R\Hhom(\tau_{\leq 0} Rf_{1\ast} \GG_{m}[1],
G)  \ar[d]_{R\Hhom(u^{\ast},G)}  & \tau_{\leq 1} Rf_{1\ast}
f^{\ast}_1G  \ar[l]_-{\sim} \ar[d]^{Tr_u}\\
\tau_{\leq 1} R\Hhom(\tau_{\leq 0} Rf_{2\ast} \GG_m[1], G)&
  \tau_{\leq 1} Rf_{2\ast}f^{\ast}_2 G \ar[l]_-{\sim}
}}\leqno{(1.5.3.6)}
$$
is therefore commutative.

\subsubsection{}\etiquette[1.5.4]{XVIII.subsubsec1.5.4}
With the notation of \hyperref[XVIII.thm1.5.2]{1.5.2}, let $G$ be an abelian sheaf on
$S$ satisfying \hyperref[XVIII.par1.3.1.1]{1.3.1.1}. Locally on $S$ for the étale topology, one has
(noncanonically)
$$
\tau_{\leq 0}
Rf_{\ast}\GG_{m}[1]\simeq
f_{\ast}\GG_{m}[1]+R^{1}f_{\ast}\GG_{m}{\quoi}.
$$

At the level of Picard stacks, one indeed defines a section
$\ch(R^{1}f_{\ast}\GG_{m})\rightarrow
f_{\ast}\ch(\GG_{m}[1])$\nde{Recall that this stack identifies with
$\PIC(X/S)$, \cf \eqhyperref[XVIII.eq1.5.1.1]{1.5.1.1}.} by associating with each isomorphism class
of invertible sheaves the invertible sheaf belonging to that class,
rigidified along suitable sections of $X/S$.

The isomorphism \eqhyperref[XVIII.eq1.5.2.2]{1.5.2.2} therefore gives an isomorphism
\begin{equation*}
\Hhom(R^{1}f_{\ast}\GG_{m},G)\xrightarrow{\sim}\Hhom_{S}(X,G){\quoi}\tag{1.5.4.1}
\end{equation*}\etiquette[1.5.4.1]{XVIII.eq1.5.4.1}%
and an exact sequence
\begin{equation*}
0\longrightarrow
\Eext^{1}(R^{1}f_{\ast}\GG_{m},G)\longrightarrow
R^1f_{\ast}G\longrightarrow
\Hhom(f_{\ast}\GG_{m},G)\longrightarrow 0{\quoi}.\tag{1.5.4.2}
\end{equation*}\etiquette[1.5.4.2]{XVIII.eq1.5.4.2}

The morphisms in \eqhyperref[XVIII.eq1.5.4.1]{1.5.4.1} and \eqhyperref[XVIII.eq1.5.4.2]{1.5.4.2} are interpreted as
follows:
\begin{enumerate}
\renewcommand{\theenumi}{\alph{enumi}}
\renewcommand{\labelenumi}{\theenumi)}
\item The\pageoriginale morphism \eqhyperref[XVIII.eq1.5.4.1]{1.5.4.1} and the first
  morphism in \eqhyperref[XVIII.eq1.5.4.2]{1.5.4.2} are defined by the canonical map
  from $X$ to $\sheafPic_{X/S}$.

\item The second morphism in \eqhyperref[XVIII.eq1.5.4.2]{1.5.4.2} is the degree \eqhyperref[XVIII.eq1.3.7.2]{1.3.7.2}.
\end{enumerate}

Locally on $S$ for the étale topology, one has (noncanonically)
$\sheafPic_{X/S}\sim \sheafPic^{0}_{X/S}\times \ZZ^{k}$; hence
$$
\Eext^{1}(R^{1}f_{\ast}\GG_{m},G)\xrightarrow{\sim}\Eext^{1}(\sheafPic^{0}_{X/S}, G)
$$
and \eqhyperref[XVIII.eq1.5.4.2]{1.5.4.2} may also be written
\begin{equation*}
0 \longrightarrow \Eext^{1}(\sheafPic^{0}_{X/S},G)\longrightarrow
R^{1}f_{\ast}G\longrightarrow
\Hhom(f_{\ast}\GG_{m},G)\longrightarrow 0{\quoi}.\tag{1.5.4.3}
\end{equation*}\etiquette[1.5.4.3]{XVIII.eq1.5.4.3}

\subsubsection{}\etiquette[1.5.5]{XVIII.subsubsec1.5.5}
When, in \eqhyperref[XVIII.eq1.5.2.2]{1.5.2.2}, one takes $K=\GG_{m}[1]$, Theorem
\hyperref[XVIII.thm1.5.2]{1.5.2} appears as a self-duality theorem for $\tau_{\leq
0}(Rf_{\ast}\GG_{m}[1])$. Any complex representing
$\tau_{\leq 0} Rf_{\ast} \GG_{m}[1]$ admits a canonical three-step
filtration, whose successive quotients are, up to unique quasi-isomorphism,
$$
\sheafPic_{X/S}/\sheafPic^{0}_{X/S}\quav \sheafPic^{0}_{X/S}\quav
(f_{\ast}\GG_{m})[1]{\quoi}.
$$

The preceding self-duality corresponds to a duality (with values in
$\GG_{m}$) between $\sheafPic_{X/S}/\sheafPic^{0}_{X/S}$ and
$f_{\ast}\GG_{m}$, and to a self-duality (with values in
$\GG_{m}[1]$) on $\sheafPic^{0}_{X/S}$ (self-duality of the
Jacobian).

\subsubsection{}\etiquette[1.5.6]{XVIII.subsubsec1.5.6}
The situation is distinctly worse for smooth curves over
$S:f:X\rightarrow S$ that are not assumed proper over $S$. The
additive group seems to behave uncontrollably. The preceding arguments
nevertheless extend to the case of smooth curves obtained from a proper
curve over $S$ by removing a part finite over $S$, when $G$ is an
étale torsion sheaf whose torsion is prime to the residue
characteristics of $S$. One may hope that étale $p$-torsion sheaves
also behave reasonably; see \textsc{Serre}
\cite[Ch.~VI nᵒ 11 p.~126]{sga4.exp.XVIII.bib.6}.

The role of Theorem \hyperref[XVIII.thm1.2.2]{1.2.2} will here be played by the acyclicity
theorem\pageoriginale \hyperrefext[XV][XV.cor2.2]{2.2}, which has the following as a corollary.

\begin{enonce*}{Lemma 1.5.7}\etiquette[1.5.7]{XVIII.lem1.5.7}
Let $p:X\rightarrow S$ be the projection of a standard affine space
over $S$, and let $G$ be an abelian torsion sheaf on $S_{\text{\rm et}}$,
of torsion prime to the residue characteristics of $S$. We also write
$G$ for the inverse image of $G$ on the site $S_{\text{\rm
fppf}}$. Then the functor $p^{\ast}$ is an equivalence of categories
between the category of torsors under $G$ on $S_{\text{\rm
et}}$ (or on $S_{\text{\rm fppf}}$, which amounts to the same thing,
by {\textsc{Grothendieck} \cite[Corollaire~11.9]{sga4.exp.XVIII.bib.2}}), and the category of torsors under $G_{X}$
on $X$.
\end{enonce*}

\setcounter{subsubsection}{7}
\subsubsection{}\etiquette[1.5.8]{XVIII.subsubsec1.5.8}
Let $\bar{f}:\bar{X}\rightarrow S$ be a proper morphism of finite
presentation whose fibers are purely one-dimensional, let $Y$ be a
closed subscheme of $\bar{X}$ finite over $S$, defined by an ideal $\mathfrak{m}$,
and suppose that $X=\bar{X}-Y$ is a smooth curve over $S$. Let
furthermore $G$ be a torsion sheaf on $S_{\text{\rm et}}$, of torsion
prime to the residue characteristics of $S$. We shall also write $G$
for the sheaf on the big fppf site of $S$ (\hyperrefext[XVII][XVII.subsec0.10]{0.10}) obtained as
the inverse image of $G$.

Let $\GG_{m}(\mathfrak{m})$ be the subsheaf of $\GG_{m\bar{X}}$
consisting of sections equal to 1 on $Y$. A torsor under
$\GG_{m}(\mathfrak{m})$ identifies with an invertible sheaf on $\bar{X}$
trivialized on $Y$. Every relative divisor on $X$, finite over $S$,

% END INCLUDED SGA4_English_current_buildable_source_20260906/source/18/drafts/SGA4_XVIII_LINES_2851_3250_English.texfrag

% BEGIN INCLUDED SGA4_English_current_buildable_source_20260906/source/18/drafts/SGA4_XVIII_LINES_3251_3650_English.texfrag
defines such an invertible sheaf. Locally on $S$ for the
étale topology, there exist such divisors giving rise to relatively
ample invertible sheaves.

\begin{enonce*}{Lemma 1.5.9}\etiquette[1.5.9]{XVIII.lem1.5.9}
Under the hypotheses of \hyperref[XVIII.subsubsec1.5.8]{1.5.8}, with $S$ affine\nde{It is
in fact necessary to suppose that $S$ is local with closed point $s$. If one wishes to retain
the original hypothesis, then in the conclusion one should require that
$m$ be defined only over the inverse image of an open neighborhood of
$s$ in $S$, as occurs below in \hyperref[XVIII.par1.5.10.1]{1.5.10.1}.},
let $\Oo_{\bar{X}}(1)$ be a relatively ample invertible sheaf on
$\bar{X}$, trivialized along $Y$, and let $\Mm$ be an
invertible sheaf on $\bar{X}$, trivialized along $Y$. For
every point $s\in S$, there exists an integer $n_{0}$ such that, for every
integer $n\geq n_{0}$ and every section $m_{s}$ of
$\Mm(n)\otimes k(s)$ on the fiber $\bar{X}_{s}$ which is equal to $1$
on $Y_{s}$, there exists a section $m$ of $\Mm(n)$ which is equal to $1$
on $Y$ and lifts $m_{s}$.
\end{enonce*}

We reduce to the case where $S$ is Noetherian (since a family of
surjective morphisms remains a family of surjective morphisms after base
change\nde{Eliminating this Noetherian hypothesis is somewhat more
subtle than might appear at first sight,
because the question is affine rather than linear. If
$\bar{X}$, $Y$, and $\Mm$ are obtained by a base change $S\to S_0$
(which is assumed to be a local morphism between local schemes) from corresponding data
$\bar{X}_0$, $Y_0$, and $\Mm_0$ over a Noetherian $S_0$
(\cf \EGA IV 8), and if the desired assertion is known in this
Noetherian situation, then one can obtain the desired result over $S$ at
the cost of an unpleasant diagram chase, during which one may
have to consider a variant of the desired assertion in which the sections are required
to be equal to $0$ instead of $1$ on $Y_s$ or $Y$.}).
Let $\Nn$ be the coherent sheaf of local sections
of $\Mm$ vanishing on\pageoriginale $\bar{X}_{s}$ and $Y$. For
each $n$, the local solutions to the problem form a principal homogeneous
space under $\Nn(n)$, so that the obstruction to
lifting $m_{s}$ lies in
$H^{1}(X,\Nn(n))=H^{0}(S,R^{1}\bar{f}_{\ast}\Nn(n))$. But,
for $n$ sufficiently large, $R^{1}\bar{f}_{\ast}\Nn(n)=0$ (\EGA III 2.2.1).

\setcounter{subsubsection}{9}
\subsubsection{}\etiquette[1.5.10]{XVIII.subsubsec1.5.10}
Let $\Ll$ be an invertible sheaf on $\bar{X}$ trivialized
along $Y$, and let $s\in S$. Consider the conditions

\paragraph{}\etiquette[1.5.10.1]{XVIII.par1.5.10.1}
Every section of $\Ll\otimes k(s)$ on $X_{s}$ which is equal to 1 on
$Y_{s}$ lifts to a section of $\Ll$ on the inverse
image $U$ of a neighborhood of $s$, equal to 1 on $\fprod{Y}{U}{S}$.

\paragraph{}\etiquette[1.5.10.2]{XVIII.par1.5.10.2}
There exist four sections $m_{i}$ of $\Ll\otimes k(s)$ on $X_{s}$,
equal to 1 on $Y_s$, and linearly independent at the generic point of
every irreducible component of $X_{s}$.

Consider the following conditions, concerning respectively a
section $m$ of $\Ll$, a pair $(m_{1},m_{2})$, and a triple
$(m_{1},m_{2},m_{3})$ of sections of $\Ll$.

\paragraph{}\etiquette[1.5.10.3]{XVIII.par1.5.10.3}
$m$ is equal to 1 on $Y$, and the subscheme $Z(m)$ of $X$ defined by $m=0$
is a relative divisor over $S$ (which means that it is quasi-finite
over $S$, and hence finite over $S$, since $X=\bar{X}-Y$ is smooth over $S$).

\paragraph{}\etiquette[1.5.10.4]{XVIII.par1.5.10.4}
Let $\lambda$ be the canonical coordinate on the affine line
$g:\EE^{1}_{S}\rightarrow S$. The section $\lambda
m_{1}+(1-\lambda)m_{2}$ of ${g'}^{\ast}\Ll$ satisfies
(\hyperref[XVIII.par1.5.10.3]{1.5.10.3}).

\paragraph{}\etiquette[1.5.10.5]{XVIII.par1.5.10.5}
Similarly, $\lambda m_{1}+\mu m_{2}+(1-\lambda -\mu)m_{3}$ satisfies
(\hyperref[XVIII.par1.5.10.3]{1.5.10.3}).

We leave it to the reader to verify, by a general-position
argument\nde{It should be observed here that, in view of
(\hyperref[XVIII.par1.5.10.1]{1.5.10.1}) and the fact that the second part of condition
(\hyperref[XVIII.par1.5.10.3]{1.5.10.3}) can be checked in the local situation by
considering only the fiber over $s$, the lemma now involves only
sections of the restriction of $\Ll$ to the fiber
$X_s$. For example, for (I), condition (\hyperref[XVIII.par1.5.10.2]{1.5.10.2})
provides an \emph{affine} space $V$ of dimension $3$ of
sections of $\Ll\otimes k(s)$ on $X_s$ which are equal to $1$ on $Y_s$, and no
section belonging to this space $V$ vanishes at any generic point
of $X_s$. We seek to determine $k\in V$ such that, for
every $i$, $(m_i,k)$ satisfies (\hyperref[XVIII.par1.5.10.4]{1.5.10.4}), that is, such that the
affine line joining the points $m_i$ and $k$ does not pass through the origin
(in any of its incarnations associated with the different irreducible
components of $X_s$); in other words, it suffices that $k$ not belong
to the line joining $m_i$ and the origin. One thus obtains a conjunction of
conditions satisfied on nonempty open subsets of $V$. To be
completely rigorous and to account for the fact that the irreducible
components need not be geometrically irreducible, when one must take a point
of $V$ in an extension (if $k(s)$ is finite), it is in fact necessary to
introduce a finite Galois extension $s'\to s$ such that the irreducible
components of $X_{s'}$ are geometrically irreducible. The dense open subset of
$V_{s'}$ obtained by arguing on $X_{s'}$ is Galois-invariant and hence defines
the desired dense open subset of $V$.}, the following

\begin{enonce*}{Lemma 1.5.11}\etiquette[1.5.11]{XVIII.lem1.5.11}
Let $\Ll$ be an invertible sheaf on $X$, trivialized along
$Y$, and let $\bar{s}$ be a geometric point of $S$ with image $s$ in
$S$. If $\Ll$ satisfies (\hyperref[XVIII.par1.5.10.1]{1.5.10.1}) and
(\hyperref[XVIII.par1.5.10.2]{1.5.10.2}) at $s$, then:
\begin{enumerate}
\renewcommand{\theenumi}{\Roman{enumi}}
\renewcommand{\labelenumi}{(\theenumi)}
\item If $n$ sections $m_{i}$ of $\Ll$ satisfy (\hyperref[XVIII.par1.5.10.3]{1.5.10.3}),
there exist an étale neighborhood $U$ of $\bar{s}$ and a section $k$ of
$\Ll|\bar{f}^{-1}(U)$ such that the pairs $(m_{i},k)$ satisfy
(\hyperref[XVIII.par1.5.10.4]{1.5.10.4}).

\item Let\pageoriginale $m_{1}$, $m_{2}$, $k_{1}$, $k_{2}$ be four
sections of $\Ll$ such that the pairs $(m_{i},k_{j})$
satisfy (\hyperref[XVIII.par1.5.10.4]{1.5.10.4}). There then exist an étale neighborhood $U$ of
$\bar{s}$ and a section $k$ of $\Ll|\bar{f}^{-1}(U)$ such that the
triples $(m_{i}, k_{j}, k)$ satisfy (\hyperref[XVIII.par1.5.10.5]{1.5.10.5}). 
\end{enumerate}
\end{enonce*}

When $k(s)$ is infinite, one could simply take $U$ to be a
Zariski neighborhood.

\setcounter{subsubsection}{11}
\subsubsection{}\etiquette[1.5.12]{XVIII.subsubsec1.5.12}
Now let $\Ll$ be an invertible sheaf on $\bar{X}$
rigidified along $Y$, and let $K$ be a $G$-torsor on $X$, where $G$ is
as in \hyperref[XVIII.lem1.5.7]{1.5.7}. We propose to define a $G$-torsor
$\langle \Ll,K]$ on $S$, generalizing the one studied in \hyperref[XVIII.subsec1.3]{1.3}, and
additive in $\Ll$.

Let $(m_{1},m_{2})$ be a pair of sections of $\Ll$ that
satisfies (\hyperref[XVIII.par1.5.10.4]{1.5.10.4}). With the notation of
\hyperref[XVIII.par1.5.10.3]{1.5.10.3} and \hyperref[XVIII.par1.5.10.4]{1.5.10.4}, the
trace, from $\lambda m_{1}+(1-\lambda)m_{2}=0$ to $\EE^{1}_{S}$,
of the torsor ${g'}^{\ast}K$ is a torsor $K_{1,2}$ on
$\EE^{1}_{S}$. Its pullback by the section $\lambda =
0$ (\resp $\lambda = 1$) is the torsor $K_{2}=\Tr_{Z(m_{2})/S}(K)$
(\resp $K_{1}=\Tr_{Z(m_{1})/S}(K)$). By \hyperref[XVIII.lem1.5.7]{1.5.7}, the torsor
$K_{1,2}$ is canonically the pullback of a torsor under $G$
on $S$; hence there is a canonical isomorphism between $K_{1}$ and $K_{2}$.

If $m_{1}$, $m_{2}$, and $k$ are three sections of $\Ll$ such
that the pairs $(m_{i},k)$ satisfy (\hyperref[XVIII.par1.5.10.4]{1.5.10.4}), composition again
gives an isomorphism
\begin{equation*}
\Psi_{k}:K_{1}=\Tr_{Z(m_{1})/S}(K)\xrightarrow{\sim}
K_{2}=\Tr_{Z(m_{2})/S}(K){\quoi}.\tag{1.5.12.1}
\end{equation*}\etiquette[1.5.12.1]{XVIII.eq1.5.12.1}

If $E$ is a relative divisor on $X$, finite over $S$, the sections
$m'_{i}=m_{i}\otimes 1$ of $\Ll\otimes \Oo(E)$ still satisfy
(\hyperref[XVIII.par1.5.10.3]{1.5.10.3}), and we have (for $i=1,2$)
\begin{equation*}
K'_{i}=\Tr_{Z(m'_{i})/S}(K)\simeq K_{i}+\Tr_{E/S}(K)\tag{1.5.12.2}
\end{equation*}\etiquette[1.5.12.2]{XVIII.eq1.5.12.2}%
and the isomorphism $\Psi_{k\otimes 1}:K'_{1}\rightarrow K'_{2}$ is
deduced from $\Psi_{k}$:
\begin{equation*}
\Psi_{k\otimes 1}=\Psi_{k}+\id_{\Tr_{E/S}}(K){\quoi}.\tag{1.5.12.3}
\end{equation*}\etiquette[1.5.12.3]{XVIII.eq1.5.12.3}

To verify that $\Psi_{k}$ does not depend on $k$, it suffices to do so
in an étale neighborhood of every point $s$ of $S$, and, by \hyperref[XVIII.eq1.5.12.3]{1.5.12.3}
and \hyperref[XVIII.lem1.5.9]{1.5.9}, we\pageoriginale reduce to the case in which
conditions (\hyperref[XVIII.par1.5.10.1]{1.5.10.1}) and (\hyperref[XVIII.par1.5.10.2]{1.5.10.2}) hold. With
the notation of \hyperref[XVIII.lem1.5.11]{1.5.11} (II),
$\Psi_{\lambda k_{i}+(1-\lambda)k}-\Psi_{k_{i}}$ is a morphism from
$\EE^{1}_{S}$ to $G$, zero at $\lambda =1$, and hence zero; thus
$\Psi_{k_{i}}=\Psi_{k}$, and therefore $\Psi_{k_{1}}=\Psi_{k_{2}}$.

Let $m_{1}$ and $m_{2}$ be two sections of $\Ll$ satisfying
\hyperref[XVIII.par1.5.10.3]{1.5.10.3}, and let $K_{i}=\Tr_{Z(m_{i})/S}(K)$. To define
an isomorphism
$\Psi_{m_{2},m_{1}}:K_{1}\rightarrow K_{2}$, it suffices to do so in an
étale neighborhood of every point $s$ of $S$, compatibly with
base change. If $E$ is a relative divisor on $X$, finite over
$S$, such that $\Ll\otimes \Oo(E)$ satisfies (\hyperref[XVIII.par1.5.10.1]{1.5.10.1})
and (\hyperref[XVIII.par1.5.10.2]{1.5.10.2}) (see \hyperref[XVIII.lem1.5.9]{1.5.9}), and if the
pairs $(m_{i}\otimes 1,k)$
satisfy (\hyperref[XVIII.par1.5.10.4]{1.5.10.4}), we set (\cf \eqhyperref[XVIII.eq1.5.12.3]{1.5.12.3})
$$
\Psi_{m_{2},m_{1}}+\Tr_{E/S}(K)=\Psi_{k}{\quoi}.
$$
This defines $\Psi_{m_{2},m_{1}}$, and one verifies using
(\hyperref[XVIII.lem1.5.11]{1.5.11}. (I)) that $\Psi_{m_{3},
  m_{2}}\Psi_{m_{2},m_{1}}=\Psi_{m_{3},m_{1}}$. These
constructions are compatible with base changes. If there exists
a section $m$ of $\Ll$ satisfying (\hyperref[XVIII.par1.5.10.3]{1.5.10.3}), we set
$$
\langle \Ll,K]=\Tr_{Z(m)/S}(K)
$$
and, by what precedes, this torsor is independent, up to canonical
isomorphism, of the choice of $m$.

In the general case, note that to define
$\langle \Ll,K]$ it suffices to define it in an étale neighborhood of every point $s$
of $S$, compatibly with base change. If, as above,
$\Ll\otimes \Oo(E)$ satisfies (\hyperref[XVIII.par1.5.10.1]{1.5.10.1})
and (\hyperref[XVIII.par1.5.10.2]{1.5.10.2}),
then $\Ll\otimes \Oo(E)$ has sections satisfying
(\hyperref[XVIII.par1.5.10.3]{1.5.10.3}), and we set
$$
\langle \Ll,K]=\langle \Ll\otimes \Oo(E), K]-\Tr_{E/S}(K){\quoi}.
$$

\subsubsection{}\etiquette[1.5.13]{XVIII.subsubsec1.5.13}
Under the hypotheses of \hyperref[XVIII.subsubsec1.5.8]{1.5.8}, if $x$ and $y$ are
two sections of
$\bar{f}_{\ast}\GG_{m}(\mathfrak{m})$, then $\lambda x+(1-\lambda)y$ is
again a section for $\lambda$ in an open subset with nonempty fibers of
the affine line over $S$.

This can be expressed by saying that

\paragraph{}\etiquette[1.5.13.1]{XVIII.par1.5.13.1}
``$\bar{f}_{\ast}\GG_{m}(\mathfrak{m})$ is path-connected.''

Every\pageoriginale morphism from an open subset with nonempty fibers of the
affine line to $G$ is constant. There is therefore no
nontrivial homomorphism from $\bar{f}_{\ast}\GG_{m}(\mathfrak{m})$ to
$G$, and the automorphisms of $\Ll$ act trivially on
$\langle \Ll,K]$. The torsor $\langle \Ll,K]$ thus depends only on
$K$ and on the section of $R^{1}\bar{f}_{\ast}\GG_{m}(\mathfrak{m})$ defined
by $\Ll$. By globalization, this makes it possible to define a
symbol $\langle \lambda,K]$ for $\lambda$ a section of
$R^{1}\bar{f}_{\ast}\GG_{m}(\mathfrak{m})$. The torsor $\langle \lambda,K]$ is
additive in $\lambda$, and therefore corresponds to an extension $e(K)$ of
$R^{1}\bar{f}_{\ast}\GG_{m}(\mathfrak{m})$ by $G$ (see \hyperref[XVIII.prop1.4.23]{1.4.23}
or \SGA 7 VII 1.1.6 and 1.2).

The formation of $e(K)$ is compatible with every base change and with
addition of torsors, and hence gives an additive functor
$$
e:\bar{f}_{\ast}\ch(G[1])\longrightarrow
\EXT(R^{1}\bar{f}_{\ast}\GG_{m}(\mathfrak{m}),G){\quoi}.
$$

Let $j$ be the canonical map from $X$ into
$R^{1}\bar{f}_{\ast}\GG_{m}(\mathfrak{m})$, $j:t\rightarrow$ (class of
$\Oo(t)$). Pullback by $j$ is an additive functor
$$
j^{\ast}:\EXT(R^{1}\bar{f}_{\ast}\GG_{m}(\mathfrak{m}),G)\longrightarrow
\bar{f}_{\ast}\ch(G[1])
$$
and trivially $j^{\ast}e\sim \Id$. One verifies as in
\hyperref[XVIII.thm1.3.10]{1.3.10} and \hyperref[XVIII.thm1.5.2]{1.5.2} that, for the fppf topology,

\begin{enonce*}{Proposition 1.5.14}\etiquette[1.5.14]{XVIII.prop1.5.14}
Under the hypotheses of \hyperref[XVIII.subsubsec1.5.8]{1.5.8}, the functors $e$ and
$j$ of \hyperref[XVIII.subsubsec1.5.13]{1.5.13} are
quasi-inverse equivalences. It follows from
\hyperref[XVIII.const1.4.19]{1.4.19} and \hyperref[XVIII.prop1.4.23]{1.4.23} that there is an isomorphism
$$
\Ext^{1}(R^{1}\bar{f}_{\ast}\GG_{m}(\mathfrak{m}),G)\xrightarrow{\sim}
H^{1}(X,G){\quoi}.
$$
\end{enonce*}

\setcounter{subsubsection}{14}
\subsubsection{}\etiquette[1.5.15]{XVIII.subsubsec1.5.15}
Let $f_{i}:X_{i}\rightarrow S$ $(i=1,2)$ be two compactified curves
as in \hyperref[XVIII.subsubsec1.5.8]{1.5.8}.

Let $\bar{u}:\bar{X}_{1}\rightarrow \bar{X}_{2}$ be a morphism such that
$Y_{1}$ is a subscheme of the inverse image of $Y_{2}$ and which
induces a finite flat morphism from $X_{1}$ to $X_{2}$. Then,
for $\Ll$ an invertible sheaf on $\bar{X}_{2}$ rigidified along $Y_2$
and $K$ a torsor on $X_{1}$, we have
$$
\langle \Ll,\tr_{X_{1}/X_{2}}(K)]\xrightarrow{\sim}
\langle \bar{u}^{\ast}\Ll,K]{\quoi},
$$
and\pageoriginale we deduce the commutativity of the diagram
$$\vcenter{\etiquette[1.5.15.1]{XVIII.eq1.5.15.1}
\xymatrix@R=1.75cm{
\Ext^{1}(R^1\bar{f}_{1\ast} \GG_m (\mathfrak{m}_1), G) 
\ar[r]^-{\sim}\ar[d]^{\Ext^{1}(\bar{u}^{\ast},\Id)} &   H^{1}(X_1,G) \ar[d]^{\Tr_{X_1 /X_2}}\\
\Ext^1(R^1\bar{f}_{2\ast}\GG_m (\mathfrak{m}_2), G)  \ar[r]^-{\sim}&
H^1(X_2,G).
}}\leqno{(1.5.15.1)}
$$

If $\bar{u}:\bar{X}_{1}\rightarrow \bar{X}_{2}$ is finite and flat, and
satisfies $Y_{1}=\bar{u}^{-1}Y_{2}$ (as schemes), then, for
$\Ll$ an invertible sheaf on $X_{1}$ rigidified along $Y_1$
and $K$ a torsor on $X_{2}$, we have
$$
\langle \Ll,u^{\ast}K]\xrightarrow{\sim}
\langle N_{\bar{X}_{1}/\bar{X}_{2}}\Ll,K]
$$
and we deduce the commutativity of the diagram
\begin{equation*}
\vcenter{
\xymatrix@R=1.75cm@C=1.5cm{
\Ext^{1}(R^{1}\bar{f}_{2\ast} \GG_m (\mathfrak{m}_2) , G)
\ar[d]^{\Ext^{1}(N_{\bar{X}_1 /\bar{X}_2}, \Id)}\ar[r]^-{\sim} &
H^{1}(X_2,G) \ar[d]^{u^{\ast}}\\
\Ext^{1}(R^{1}\bar{f}_{1\ast} \GG_m (\mathfrak{m}_1), G)  \ar[r]^-{\sim} &
H^{1}(X_1, G).
}}\tag{1.5.15.2}
\end{equation*}\etiquette[1.5.15.2]{XVIII.eq1.5.15.2}

One could of course state this result more precisely in terms of Picard
stacks.

\subsection{A vanishing theorem}\etiquette[1.6]{XVIII.subsec1.6}

The present section, from \hyperref[XVIII.lem1.6.6]{1.6.6} (2\textsuperscript{nd}
proof) to the end,
can be read independently of \hyperref[XVIII.sec2]{sections 2} to 5.

\begin{enonce*}{Lemma 1.6.1}\etiquette[1.6.1]{XVIII.lem1.6.1}
Let $X$ be a scheme and $Y$ a closed subscheme of $X$ with complement
$U$, defined by an ideal $\mathfrak{m}$.
$$
U\mathop{\hookrightarrow}^{j} X\mathop{\hookleftarrow}^{i} Y{\quoi}.
$$
Let $\GG_{m}(\mathfrak{m})$ be the \emph{fpqc} sheaf whose sections
on an $X$-scheme $X_{1}$ are the sections of $\Oo^{\ast}_{X_{1}}$
equal to $1$ on the inverse image of $Y$. For $n\geq 1$ invertible
on $X$, the sequence
$$
0\rightarrow j_{!}𝛍_{n}\rightarrow
\GG_{m}(\mathfrak{m})\xrightarrow{n} \GG_{m}(\mathfrak{m})\rightarrow
0
$$
is an exact sequence of sheaves on the big étale site
\emph{(\hyperrefext[XVII][XVII.subsec0.10]{0.10})} of $X$.
\end{enonce*}

Let\pageoriginale $i$ be the morphism of big étale sites
$i:Y\rightarrow X$ induced by $i$:
$$
i^{\ast}(X_{1}/X)=\fprod{X_{1}}{Y}{X}{\quoi}.
$$
The functor $i_{\ast}$ is exact. In the diagram
$$\vcenter{\etiquette[1.6.1.1]{XVIII.eq1.6.1.1}
\xymatrix{
             &    0  \ar[d]           &  0  \ar[d]                   &    0\ar[d]    \\
0  \ar[r]    &j_!𝛍_n  \ar[r]\ar[d]  & \GG_m(\mathfrak{m})\ar[r]\ar[d]  &
             \GG_m(\mathfrak{m}) \ar[r] \ar[d] &0\\
0  \ar[r]     &𝛍_n  \ar[r]  \ar[d]  & \GG_m  \ar[r] \ar[d]  &
  \GG_m  \ar[r] \ar[d]  & 0\\
0  \ar[r]    & i_{\ast}𝛍_n  \ar[d] \ar[r] &
  i_{\ast}\GG_m  \ar[r]\ar[d]  & i_{\ast}\GG_m   \ar[r]\ar[d] & 0\\
             &    0             &  0                    &    0   \\
}}\leqno{(1.6.1.1)}
$$
the columns are exact, as are the last two rows
(\hyperrefext[IX][IX.Thr3.2]{3.2}). The first row is therefore exact.

\setcounter{subsubsection}{1}
\subsubsection{}\etiquette[1.6.2]{XVIII.subsubsec1.6.2}
Let $k$ be a field and
$$
G'_{1}\longrightarrow G'_{2}\longrightarrow G\longrightarrow
G''_{2}\longrightarrow G''_{1}
$$
an exact sequence of abelian sheaves on the big fppf site of
$\Spec (k)$. Recall that if $G'_{1}$ and $G'_{2}$ are
representable and of finite type, and $G''_{2}$ and $G''_{1}$ are representable
and locally of finite type, then $G$ is representable and locally of
finite type\acorriger{In the case where $G'_1=G'_2=0$, this is clear. In the
case where $G''_1=G''_2=0$, using \SGA 3 V 10.1.1 (new edition),
one can reduce to the case where $G'_1\to G'_2$ is a
monomorphism, which is treated in (\SGA 3 VI${}_{\text{A}}$ 3.2).
We are thus reduced to the case where $G'_1=G''_1=0$. One can then regard the
sheaf $G$ as a torsor under $G'_2$ over the scheme $G''_2$. If
$G'_2$ is assumed affine (the situation to which one is reduced,
by (\SGA 3 VI${}_{\text{B}}$ 11.17), if before making the reductions
$G'_1$ and $G'_2$ were assumed affine), then the representability of $G$
follows immediately from faithfully flat descent. This special case is
sufficient for the applications made in this paragraph. If $G'_2$ is
quasi-projective, the representability of $G$ can be obtained by using
(\SGA 1 VIII 7.7) and arguments similar to
those of \cite[Lemma~3.3.1]{sga4.exp.XVIII.bib.4bis}. In fact, every group scheme $G$
of finite type over a field is quasi-projective: one may suppose that $G$
is connected; this is then \cite[Corollary~1.2]{sga4.exp.XVIII.bib.0bis} if $G$ is smooth, and (M.
Raynaud) the general case reduces to it because, by (\SGA 3
VII${}_{\text{A}}$ 8.3), there exists
a finite morphism from $G$ to a smooth group scheme.}.

If $P$ is a connected commutative algebraic group of finite type over
$k$, and $n$ is an integer invertible in $k$, the sequence
\begin{equation*}
0\rightarrow {}_{n}P\rightarrow P\xrightarrow{n} P\rightarrow 0\tag{1.6.2.1}
\end{equation*}\etiquette[1.6.2.1]{XVIII.eq1.6.2.1}%
is exact. If $G$ is a finite étale commutative group over $k$, killed
by $n$, the first group and the last arrow of the exact sequence
$$
\Hom(P,G)\rightarrow \Hom({}_{n}P,G)\xrightarrow{j}
\Ext^{1}(P,G)\xrightarrow{n} \Ext^{1}(P,G)
$$
are zero, whence an isomorphism
\begin{equation*}
j:\Hom({}_{n}P,G)\xrightarrow{\sim} \Ext^{1}(P,G){\quoi}.\tag{1.6.2.2}

% END INCLUDED SGA4_English_current_buildable_source_20260906/source/18/drafts/SGA4_XVIII_LINES_3251_3650_English.texfrag

% BEGIN INCLUDED SGA4_English_current_buildable_source_20260906/source/18/drafts/SGA4_XVIII_LINES_3651_4050_English.texfrag
\end{equation*}\etiquette[1.6.2.2]{XVIII.eq1.6.2.2}

\subsubsection{}\etiquette[1.6.3]{XVIII.subsubsec1.6.3}
Under\pageoriginale the hypotheses of \hyperref[XVIII.lem1.6.1]{1.6.1}, suppose
that $X$ is proper over the spectrum of a field $k$
% SGA-SRC-001141
$$
p:X\longrightarrow \Spec(k){\quoi}.
$$
The exact sequence
$$
0\longrightarrow \GG_{m}(\mathfrak{m})\longrightarrow
\GG_{m}\longrightarrow i_{\ast}\GG_{m}\longrightarrow 0
$$
defines an exact sequence of fppf sheaf cohomology on $\Spec
(k)$, which is easily interpreted directly as
\begin{equation*}
 0\rightarrow p_{\ast}\GG_{m}(\mathfrak{m})\rightarrow
p_{\ast}\GG_{m_{X}}\rightarrow
(pi)_{\ast}\GG_{m_{Y}}\xrightarrow{j}\tag{1.6.3.1}
 R^{1}p_{\ast}\GG_{m}(\mathfrak{m})\rightarrow
\sheafPic_{X/k}\rightarrow \sheafPic_{Y/k}{\quoi}.
\end{equation*}\etiquette[1.6.3.1]{XVIII.eq1.6.3.1}

The sheaves $p_{\ast}\GG_{m_{X}}$ and
$(pi)_{\ast}\GG_{m_{Y}}$ are representable by group schemes of finite
type\sisi{}{ and affine}. The sheaves $\sheafPic_{X/k}$ and
$\sheafPic_{Y/k}$ are representable by group schemes locally of finite
type. By \hyperref[XVIII.subsubsec1.6.2]{1.6.2}, the sheaves
$p_{\ast}\GG_{m}(\mathfrak{m})$ and
$R^{1}p_{\ast}\GG_{m}(\mathfrak{m})$ are representable. By
\eqhyperref[XVIII.par1.5.13.1]{1.5.13.1}, the former is represented by a connected group scheme
(of finite type); the latter is locally of finite type.

\begin{enonce*}{Definition 1.6.4}\etiquette[1.6.4]{XVIII.def1.6.4}
Under the preceding hypotheses, we denote by
$\sheafPic_{\mathfrak{m},X/k}$ the group scheme representing
$R^{1}p_{\ast}\GG_{m}(\mathfrak{m})$. We denote by
$\sheafPic^{0}_{\mathfrak{m},X/k}$ its identity component and by
$\sheafPic^{\tau}_{\mathfrak{m},X/k}$ the inverse image in
$\sheafPic_{\mathfrak{m},X/k}$ of the torsion subgroup of
$\sheafPic_{\mathfrak{m},X/k}/\sheafPic^{0}_{\mathfrak{m},X/k}$.
\end{enonce*}

The scheme $\sheafPic_{\mathfrak{m},X/k}$ represents the fppf sheaf generated by
the presheaf which associates with every scheme $S$ over $k$ the set
of isomorphism classes of invertible modules on $\fprod{S}{X}{k}$,
trivialized along $\fprod{S}{Y}{k}$ (\hyperref[XVIII.subsubsec1.5.8]{1.5.8}).

The Kummer exact sequence \hyperref[XVIII.lem1.6.1]{1.6.1} gives, for $k$
algebraically closed
and $n$ invertible in $k$, an isomorphism
\begin{equation*}
H^{1}_{c}(U,𝛍_{n})\xrightarrow{\sim}
{}_{n}\sheafPic^{\tau}_{\mathfrak{m},X/k}(k){\quoi}.\tag{1.6.4.1}
\end{equation*}\etiquette[1.6.4.1]{XVIII.eq1.6.4.1}

\setcounter{subsubsection}{4}
\subsubsection{}\etiquette[1.6.5]{XVIII.subsubsec1.6.5}
Let\pageoriginale $X$ be a smooth curve over an algebraically closed
field $k$, let $\bar{X}$ be the smooth projective curve containing $X$ as a
dense open subscheme, let $Y$ be a closed subscheme of $\bar{X}$ such that $X=\bar{X}-Y$,
let $\mathfrak{m}$ be the ideal defining $Y$, let $n$ be an integer invertible in $k$, and
let $G$ be a finite abelian group killed by $n$.

The group $\sheafPic_{\mathfrak{m},X/k}$ is noncanonically the product of the
subgroup $\sheafPic^{0}_{\mathfrak{m},\bar{X}/k}$ and a free abelian
group. By \eqhyperref[XVIII.eq1.6.4.1]{1.6.4.1}, \eqhyperref[XVIII.eq1.6.2.2]{1.6.2.2}, and
\hyperref[XVIII.prop1.5.14]{1.5.14}, we therefore have a canonical composite isomorphism
\begin{gather*}
\Hom(H^{1}_{c}(X,𝛍_{n}),G)\xleftarrow{\sim}
\Hom({}_{n}\sheafPic^{\tau}_{\mathfrak{m},\bar{X}/k}(k),G)=\Hom({}_{n}\sheafPic^{0}_{\mathfrak{m},\bar{X}/k},G)\\
\xrightarrow{\sim}
\Ext^{1}(\sheafPic^{0}_{\mathfrak{m},\bar{X}/k},G)\xleftarrow{\sim}
\Ext^{1}(\sheafPic_{\mathfrak{m},\bar{X}/k},G)\xrightarrow{\sim} H^{1}(X,G)
\end{gather*}

If $u:X\rightarrow Y$ is a finite morphism between smooth curves over
$k$, we deduce from \hyperrefext[XVII][XVII.subsec6.3]{6.3}. and \eqhyperref[XVIII.eq1.5.15.2]{1.5.15.2} the commutativity of the
diagram
\begin{equation*}
\vcenter{
\xymatrix@C=2cm@R=1.5cm{
\Hom(H^{1}_c(Y,𝛍_n),G)  \ar[r]^(.55){\sim} \ar[d]^{\Hom(\Tr_u,G)} &
H^{1}(Y,G)               \ar[d]^{u^{\ast}}\\
\Hom(H^1_c(X,𝛍_n),G)    \ar[r]^(.55){\sim} & H^{1}(X,G).
}}\tag{1.6.5.1}
\end{equation*}\etiquette[1.6.5.1]{XVIII.eq1.6.5.1}

\begin{enonce*}{Lemma 1.6.6}\etiquette[1.6.6]{XVIII.lem1.6.6}
Let $X$ be a smooth curve over an algebraically closed field $k$, let $n$
be an integer invertible in $k$, and let $u:X'\rightarrow X$ be a tame
étale covering. The morphism
$$
\Tr_{u}:H^{1}_{c}(X',\ZZ/n)\longrightarrow
H^{1}_{c}(X,\ZZ/n)
$$
is then (noncanonically\nde{The isomorphism is canonical if an
isomorphism $𝛍_n\simeq \mathbf{Z}/n$ has been chosen.})
isomorphic to the transpose of the morphism
$$
u^{\ast}:H^{1}(X,\ZZ/n)\longrightarrow
H^{1}(X',\ZZ/n){\quoi}.
$$
\emph{1\textsuperscript{st} proof.} Since $𝛍_{n}$ is isomorphic to
$\ZZ/n$, \hyperref[XVIII.lem1.6.6]{1.6.6} is the special case of \eqhyperref[XVIII.eq1.6.5.1]{1.6.5.1} with
$G=\ZZ/n$.
\end{enonce*}

\noindent
\emph{2\textsuperscript{nd} proof.}
If $k$ is the field $\CC$ of complex numbers, then, by the
comparison theorems, the morphisms $\Tr_{u}$ and $u^{\ast}$
are identified with the analogous morphisms defined by the continuous map
of topological spaces $u:X'(\CC)\rightarrow
X(\CC)$:
\begin{align*}
\Tr_{u} &: H^{1}_{c}(X'(\CC),\ZZ/n) \longrightarrow
H^{1}_{c}(X(\CC), \ZZ/n)\quad \text{and}\\
u^{\ast} &: H^{1}(X(\CC),\ZZ/n)\longrightarrow H^{1}(X'(\CC),\ZZ/n){\quoi}.
\end{align*}
These\pageoriginale two homomorphisms are transposes of one another
under Poincaré duality.

By the Lefschetz principle, the assertion \hyperref[XVIII.lem1.6.6]{1.6.6} remains
true for $k$ of characteristic $0$ (the cohomology groups
under consideration are indeed invariant under extension of the algebraically closed
base field).

If $k$ has characteristic $p>0$, let $W(k)$ be the ring of Witt vectors
over $k$. Let $\bar{X}$ be the smooth projective completion
of $X$. By (\SGA 1 III 7.4), the curve $\bar{X}$ lifts to a
smooth projective curve $\bar{X}_{1}$ over $W(k)$. Since $\bar{X}$
is smooth, every point $s_{i}\in S=\bar{X}-X$ lifts to a
section $s'_{i}$ of $\bar{X}_{1}$ over $W(k)$; set
$S_{1}=⋃_{i}s'_{i}(\Spec(W(k)))$ and $X_{1}=\bar{X}_{1}-S_{1}$.

Recall (\SGA 1 XIII 2.4)
that the covering $X'$ of $X$ lifts to an
étale covering $X'_{1}$ of $X_{1}$, and that $X'_{1}$ is obtained from a
smooth projective curve $\bar{X}'_{1}$ over $W(k)$ by removing
the union $S'_{1}$ of finitely many disjoint sections.
$$
\xymatrix@=.5cm{
                  &     X_1 \ar[ddddd]^{f_1}       &                 &
                  X\ar[ll]\ar[ddddd]\\
 & & &\\
X'_1 \ar[ruu]_{u_1}\ar[rddd]_{f'_1}&                & X' \ar[ll]\ar[ruu]_{u} \ar[rddd] &   \\
& & &\\
& & &\\
                  &    \Spec(W(k)) &                 &\Spec(k).\ar[ll]
}
$$

The étale sheaves $R^{1}f_{1\ast}(\ZZ/n)$,
$R^{1}f'_{1\ast}(\ZZ/n)$, $R^{1}f_{1!}(\ZZ/n)$, and
$R^{1}f'_{1!}(\ZZ/n)$ are locally constant and their formation is
compatible with base change,\pageoriginale as follows from
the exact sequences relating the cohomology of $\bar{X}_{1}$, $X_{1}$,
and $S_{1}$, or of $\bar{X}'_{1}$, $X'_{1}$, and $S'_{1}$.

The special or generic geometric fibers of the morphisms
$$
\Tr_{u_{1}}:R^{1}f'_{1!}\ZZ/n\longrightarrow
R^{1}f_{1!}\ZZ/n
$$
or
$$
u^{\ast}_{1}:R^{1}f_{1\ast}\ZZ/n\longrightarrow
R^{1}f'_{1\ast}\ZZ/n
$$
are therefore isomorphic, and this reduces us to the case already treated in which $k$
has characteristic zero.

\begin{enonce*}{Lemma 1.6.7}\etiquette[1.6.7]{XVIII.lem1.6.7}
Let $X$ be a connected smooth curve over an algebraically closed field
$k$ of characteristic exponent $p$, and let $n$ be an integer prime to $p$. There
exists a principal covering \emph{(= finite étale Galois
(surjective))} $u:X'\rightarrow X$, of degree dividing a power of
$n$, such that the trace morphism
$$
\Tr_{u}:H^{1}_{c}(X', \ZZ/n)\longrightarrow
H^{1}_{c}(X,\ZZ/n)
$$
is zero.
\end{enonce*}

Recall that the group $H^{1}(X,\ZZ/n)$ is finite (\hyperrefext[IX][IX.thm4.6]{4.6} or
\hyperrefext[XVI][XVI.thm5.2]{5.2}). Let $u:X'\rightarrow X$ be the largest connected abelian principal
covering of $X$ whose Galois group is killed by $n$ (its Galois
group is denoted by $H_{1}(X,\ZZ/n)$). By construction, the
homomorphism
$$
u^{\ast}:H^{1}(X,\ZZ/n)\longrightarrow H^{1}(X', \ZZ/n)
$$
is zero. By \hyperref[XVIII.lem1.6.6]{1.6.6}, the homomorphism
$$
\Tr_{u}:H^{1}_{c}(X',\ZZ/n)\longrightarrow
H^{1}_{c}(X,\ZZ/n)
$$
is therefore zero.

Recall the following lemma (\EGA IV 15.6.5):

\begin{enonce*}{Lemma 1.6.8}\etiquette[1.6.8]{XVIII.lem1.6.8}
Let $f:X\rightarrow S$ be a smooth morphism and let $x$ be a section of
$f$. Then there exists a Zariski open subset $U$ of $X$, containing $x$,
such that the geometric fibers of $f|U$ are connected.
\end{enonce*}

\begin{enonce*}{Fundamental Lemma 1.6.9}\etiquette[1.6.9]{XVIII.lemfond1.6.9}
Let\pageoriginale $f:X\rightarrow S$ be a compactifiable smooth curve
(\hyperref[XVIII.def1.1.2]{1.1.2}) \emph{(\hyperrefext[XVII][XVII.subsec3.2]{3.2})},
let $x$ be a geometric point
of $X$, let $s$ be its image in $S$, and let $n\geq 1$ be an integer invertible on
$S$. There exist an étale neighborhood $V$ of $s$ in $S$ and an étale
neighborhood $U$ of $x$ in $f^{-1}(V)=\fprod{V}{X}{S}$
\begin{equation*}
\vcenter{
\xymatrix@R=1.5cm@C=1cm{
x \ar[r]\ar@{-}[d]& U\ar[rd]_{f'}\ar[r]^{u} & f^{-1}(V)\ar[r] \ar[d]^{f_v}& X\ar[d]^{f}\\
s \ar[rr]&   & V\ar[r]         & S,
}}\tag{1.6.9.1}
\end{equation*}\etiquette[1.6.9.1]{XVIII.eq1.6.9.1}%
such that $R^{0}f'_{!}\sheaf{\ZZ}/n=0$, the arrow
\emph{\hyperrefext[XVII][XVII.subsec6.2]{6.2}}
$$
\Tr_{u}:R^{1}f'_{!}\sheaf{\ZZ}/n\longrightarrow
R^{1}f_{v!}\sheaf{\ZZ}/n
$$
is zero, and the trace morphism
$$
\Tr_{f'}:R^{2}f'_{!}\sheaf{\ZZ}/n(1)\longrightarrow
\sheaf{\ZZ}/n
$$
is an isomorphism.
\end{enonce*}

For $R^{0}f'_{!}\sheaf{\ZZ}/n=0$, it suffices that $f'$ be
quasi-affine; for $\Tr_{f'}$ to be an isomorphism, it suffices that
the geometric fibers of $f'$ be connected
(\hyperref[XVIII.lem1.1.9]{1.1.9}); by
\hyperref[XVIII.lem1.6.8]{1.6.8}, these properties of $f'$ are easy to obtain
locally on
$S$ for the étale topology. It remains to prove, for $X$ with geometrically
connected fibers, the existence of a diagram \eqhyperref[XVIII.eq1.6.9.1]{1.6.9.1} for
which the arrow $\Tr_{u}$ under consideration is zero: it will suffice to
shrink $U$ and $V$ so that the other two required properties of
$f'$ hold.

The problem is local on $S$ in a neighborhood of $s$; by \hyperrefext[XVII][XVII.thm5.2.6]{5.2.6}
(base change), we immediately reduce to the case where $S$ is Noetherian. By
passage to the limit, and using \hyperrefext[XVII][XVII.thm5.3.6]{5.3.6} (constructibility), we
reduce to the case where $S$ is strictly local Noetherian, with closed point the image
of $s$.

Let $t$ be a geometric point of $S$. By \hyperref[XVIII.lem1.6.7]{1.6.7}, there
exists a principal covering of degree prime to the characteristic exponent
% SGA-SRC-001143
$p$ of $k(t)$, $u:X^{1}_{t}\rightarrow X_{t}$, such that the arrow
$$
\Tr_{u}:H^{1}_{c}(X^{1}_{t},\ZZ/n)\longrightarrow
H^{1}_{c}(X_{t},\ZZ/n)
$$
is\pageoriginale zero. By the acyclicity theorem \hyperrefext[XV][XV.thm2.1]{2.1} (or
\hyperrefext[XV][XV.lem2.6]{2.6}), there exists an étale neighborhood $v:U\rightarrow X$ of $x$ in
$X$ such that $X^{1}_{t}$ splits over $U_{t}$, so that
$v_{t}:U_{t}\rightarrow X_{t}$ factors through $X^{1}_{t}$. We therefore still have
$$
\Tr_{v_{t}}=0:H^{1}_{c}(U_{t},\ZZ/n)\longrightarrow
H^{1}_{c}(X_{t},\ZZ/n){\quoi}.
$$

For every geometric point $t$ of $S$, there therefore exists an étale
neighborhood of $x$ in $X$
$$
\xymatrix{
U \ar[r]^{v} \ar[rd]_{f'}& X \ar[d]^{f}\\
  & S
}
$$
such that the fiber at $t$ of
\begin{equation*}
\Tr_{v}:R^{1}f'_{!}\ZZ/n\longrightarrow
R^{1}f_{!}\ZZ/n\tag{1.6.9.2}
\end{equation*}\etiquette[1.6.9.2]{XVIII.eq1.6.9.2}%
is zero (\hyperrefext[XVII][XVII.thm5.2.6]{5.2.6}).

As $U$ varies, the images of the arrows \eqhyperref[XVIII.eq1.6.9.2]{1.6.9.2} form a
filtered decreasing system of subsheaves of the constructible sheaf
$R^{1}f_{!}\ZZ/n$, which becomes zero at any given geometric
point of $S$. By Noetherian induction, one verifies that such a
system is always stationary, with stationary value zero. Hence there
exists $U$ such that \eqhyperref[XVIII.eq1.6.9.2]{1.6.9.2} is zero, and this completes the
proof of \hyperref[XVIII.lemfond1.6.9]{1.6.9}.

\section{The trace morphism}\etiquette[2]{XVIII.sec2}\pageoriginale

This \S is devoted to the proof of Theorems
\hyperref[XVIII.thm2.9]{2.9} and \hyperref[XVIII.prop2.12]{2.12}. The proof of
\hyperref[XVIII.thm2.9]{2.9} uses only no.~1 of \S \hyperref[XVIII.sec1]{1}.

\begin{enonce*}{Lemma 2.1}\etiquette[2.1]{XVIII.lem2.1}
Let $f:X\rightarrow S$ be a compactifiable morphism
\emph{(\hyperrefext[XVII][XVII.def3.2.1]{3.2.1})} of relative dimension $\leq d$, and
let $j:U\hookrightarrow X$ be an open subscheme of finite type of $X$, whose
complement has fiberwise dimension $< d$ (for example, $U$ is
fiberwise dense in $X$). If $F$ is a torsion sheaf on
$X$, the canonical arrow
$$
R^{2d}(fj)_{!}j^{\ast}F\xrightarrow{\sim} R^{2d}f_{!}F
$$
is an isomorphism.
\end{enonce*}

Let $i$ be the inclusion of $Y=X-U$ into $X$. The exact sequence
(XVII \hyperref[XVII.eq5.1.16.2]{5.1.16.2}) gives us an exact sequence
$$
R^{2d-1}(fi)_{!}i^{\ast}F\longrightarrow
R^{2d}(fj)_{!}j^{\ast}F\longrightarrow R^{2d}f_{!}F\longrightarrow
R^{2d}(fi)_{!}i^{\ast}F
$$
whose outer terms are zero by (\hyperrefext[XVII][XVII.cor5.2.8.1]{5.2.8.1}).

\begin{enonce*}{Lemma 2.2}\etiquette[2.2]{XVIII.lem2.2}
Let $f:X\rightarrow S$ be a compactifiable morphism of relative
dimension $\leq d$, and let $u_{i}:X_{i}\rightarrow X$ be a family of
separated étale morphisms of finite type; let
$u_{ij}:X_{ij}=\fprod{X_{i}}{X_{j}}{X}\rightarrow X$. For every
torsion sheaf $F$ on $X$, the sequence
$$
\bigoplus_{i,j} R^{2d}(fu_{ij})_{!}u^{\ast}_{ij}F\rightrightarrows
\bigoplus_{i} R^{2d}(fu_{i})_{!}u^{\ast}_{i}F\rightarrow
R^{2d}f_{!}F\rightarrow 0
$$
of arrows \emph{(XVII \hyperref[XVII.eq6.2.7.2]{6.2.7.2})} is exact.
\end{enonce*}

This follows immediately from the right exactness of the functor $R^{2d}f_{!}$ and
from \hyperrefext[XVII][XVII.prop6.2.9]{6.2.9}.

\begin{enonce*}{Lemma 2.3}\etiquette[2.3]{XVIII.lem2.3}
Let $f:X\rightarrow S$ be a flat morphism of finite
presentation. The open subset of $X$ on which $f$ is Cohen--Macaulay \emph{(\EGA
IV 6.8.1 and 12.2.1 (vii)} is relatively of finite presentation and
fiberwise dense.
\end{enonce*}

The\pageoriginale question is local on $S$, which we reduce to assuming
Noetherian, whence the 1\textsuperscript{st} assertion. The second is checked
fiberwise, hence with $S$ the spectrum of a field. The local rings
of $X$ at the maximal points of $X$ are Artinian, hence
Cohen--Macaulay, which proves the assertion since ``Cohen--Macaulay'' is
an ``open'' property.

\begin{enonce*}{Lemma 2.4}\etiquette[2.4]{XVIII.lem2.4}
Let $f:X\rightarrow S$ be a flat morphism of finite presentation and
Cohen--Macaulay, with $S$ quasi-separated. Every point of $X$ has a
(Zariski) open neighborhood $U$ such that there exists a quasi-finite
flat $S$-morphism of finite presentation from $U$ to the standard vector bundle
$\EE^{d}_{S}$.
\end{enonce*}

This is a special case of \EGA IV 15.4.3.

\begin{enonce*}{Lemma 2.5}\etiquette[2.5]{XVIII.lem2.5}
Let $\underline{u}=(u_{1},\ldots,u_{d})$ and $\underline{v}=(v_{1},\ldots,v_{d})$
be two systems of parameters of a Cohen--Macaulay local ring $A$
of dimension $d$ and maximal ideal $\mathfrak{m}$. There exists a sequence
$\underline{w}_{0},\ldots,\underline{w}_{n}$ of systems of parameters of $A$ such that
$\underline{w}_{0}=\underline{u}$, $\underline{w}_{n}=\underline{v}$, and
each $\underline{w}_{i+1}$ is obtained from $\underline{w}_{i}$ by
changing only one of the parameters.
\end{enonce*}

We argue by induction on $d$, the assertion being vacuous for $d<2$;
we therefore suppose $d\geq 2$. For an element $x\in \mathfrak{m}$ to be
$A/u_{1}$-regular (\resp $A/v_{1}$-regular), it is necessary and
sufficient that $x$ belong to no member of $\Ass(A/u_{1})$ (\resp
$\Ass(A/v_{1})$); since no finite union of prime ideals $\neq \mathfrak{m}$
is equal to $\mathfrak{m}$, there exists $x\in \mathfrak{m}$ that is both $A/u_{1}$- and $A/v_{1}$-regular, and hence there also exist systems of parameters
$\underline{a}=(u_{1},x,\ldots)$ and $\underline{b}=(v_{1},x,\ldots)$. We conclude by
applying the induction hypothesis to $\underline{u}-\{u_{1}\}$ and
$\underline{a}-\{u_{1}\}$ in $A/u_{1}$, to $\underline{a}-\{x\}$ and
$\underline{b}-\{x\}$
in $A/x$, and to $\underline{b}-\{v_{1}\}$ and $\underline{v}-\{v_{1}\}$ in $A/v_{1}$.

\begin{enonce*}{Lemma 2.6}\etiquette[2.6]{XVIII.lem2.6}
Let $X$ be a Cohen--Macaulay scheme of finite type over an
algebraically closed field $k$, and let $u$, $v:X\rightarrow \EE^{d}_{k}$
be two quasi-finite flat morphisms from $X$ to a standard vector
bundle. Every closed point $x$ of $X$ has a (Zariski) open
neighborhood $U$ such that there exists a sequence\pageoriginale of morphisms
$w_{0},\ldots,w_{n}:U\rightarrow \EE^{d}_{k}$ such that $w_{0}=u|U$,
$w_{n}=v|U$, and $w_{i+1}$ differs from $w_{i}$ in only one
coordinate.
\end{enonce*}

Suppose first that $u(x)=v(x)=0$. The germs at $x$ of
quasi-finite flat morphisms from $X$ to $\AAA^{d}_{k}$,

% END INCLUDED SGA4_English_current_buildable_source_20260906/source/18/drafts/SGA4_XVIII_LINES_3651_4050_English.texfrag

% BEGIN INCLUDED SGA4_English_current_buildable_source_20260906/source/18/drafts/SGA4_XVIII_LINES_4051_4450_English.texfrag
sending $x$ to $0$, then identify with the systems of parameters of
the local ring $\Oo_{X,x}$, and \hyperref[XVIII.lem2.6]{2.6} follows from
\hyperref[XVIII.lem2.5]{2.5}. In the general
case, one first joins $u$ to $u-u(x)$ (\resp $v$ to $v-v(x)$)
by the chain of morphisms
$$
u-(u_{1}(x),\ldots,u_{i}(x),0,\ldots,0)(\text{\resp}
v-(v_{1}(x),\ldots,v_{i}(x),0,\ldots,0)){\quoi}.
$$

\setcounter{subsection}{6}
\subsection{}\etiquette[2.7]{XVIII.subsec2.7}
Let there be a pair of composable $S$-compactifiable morphisms
$X\xrightarrow{g} Y\xrightarrow{f}S$. Suppose that $f$ has relative
dimension $\leq d$ and $g$ has relative dimension $\leq
e$. For every torsion sheaf $F$ on $X$, the composition spectral
sequence
\begin{equation*}
R^{p}f_{!}R^{q}g_{!}F\Longrightarrow R^{p+q}(fg)_{!}F\tag{2.7.1}
\end{equation*}\etiquette[2.7.1]{XVIII.eq2.7.1}%
then gives us (by the maximum-cycle argument) an isomorphism
\begin{equation*}
R^{2d}f_{!}R^{2e}g_{!}F\simeq R^{2(d+e)}(fg)_{!}F{\quoi}.\tag{2.7.2}
\end{equation*}\etiquette[2.7.2]{XVIII.eq2.7.2}

\subsection{}\etiquette[2.8]{XVIII.subsec2.8}
Let
$$
a^{d}:\EE^{d}_{S}\longrightarrow S
$$
be the standard vector bundle over a coherent (\ie quasi-compact,
quasi-separated) scheme $S$. If $F$ is a torsion sheaf on $S$, prime
to the residual characteristics of $S$, we define as follows, by
induction on $d$, a trace morphism
\begin{equation*}
\Tr_{a^d}:R^{2d}a^{d}_{!}F(d)\xrightarrow{\sim} F{\quoi}.\tag{2.8.1}
\end{equation*}\etiquette[2.8.1]{XVIII.eq2.8.1}

For $d=0$, $\Tr_{a^0}$ is the identity. For $d=1$, $\Tr_{a^1}$
is the isomorphism (\hyperref[XVIII.prop1.1.6]{1.1.6}) (\cf
\hyperref[XVIII.lem1.1.9]{1.1.9}). Finally, there is a canonical isomorphism
$$
\EE^{d+1}_{S}=\fprod{\EE^{1}_{S}}{\EE^{d}_{S}}{S}=\EE^{1}_{\EE_{S}^d}
$$
and\pageoriginale we define $\Tr_{a^{d+1}}$ as the composite of the
isomorphisms $\Tr_{a^{d}_{S}}$ and
$Ra^{d}_{S!}(\Tr_{a^1_{\EE^d_S}})$
(\cf diagram Var 3 of \hyperref[XVIII.thm2.9]{2.9}).

The formation of the source (and the target) of \eqhyperref[XVIII.eq2.8.1]{2.8.1} is compatible with every
base change. Every algebraic group over $S$ acting on
$\EE^{d}_{S}$ therefore acts on these sheaves; a connected algebraic group
necessarily acts trivially. The group of coordinate
permutations, contained in the affine group, therefore acts
trivially on the source (and the target) of \eqhyperref[XVIII.eq2.8.1]{2.8.1}.

\setcounter{subsubsection}{1}
\subsubsection{}\etiquette[2.8.2]{XVIII.subsubsec2.8.2}
The isomorphism \eqhyperref[XVIII.eq2.8.1]{2.8.1} is invariant under permutation of the
coordinates.

This point could also be verified by interpreting
\eqhyperref[XVIII.eq2.8.1]{2.8.1} in terms
of cup products (XVII \hyperref[XVII.eq5.4.2.2]{5.4.2.2} and \hyperrefext[XVII][XVII.lem5.4.3.5]{5.4.3.5}).

\begin{enonce*}{Theorem 2.9}\etiquette[2.9]{XVIII.thm2.9}
Consider triples $(f,d,F)$ consisting of a compactifiable morphism
\emph{(\hyperrefext[XVII][XVII.subsec6.3]{6.3})} $f:X\rightarrow Y$, an integer $d$, and
a torsion sheaf
$F$ on $Y$ prime to the residual characteristics of
$Y$, where the morphism $f$ satisfies the condition $(\ast)_{d}$: there exists
an open subset $U$ of $X$ such that the restriction of $f$ to $U$ is a
flat morphism of finite presentation with fibers of dimension $\leq d$,
and such that the fibers of $X-U$ have dimension $< d$.

There is one and only one way to associate with every
triple $(f,d,F)$ as above a trace morphism
$$
\Tr_{f}:R^{2d}f_{!}f^{\ast}F(d)\longrightarrow F
$$
such that the following conditions are satisfied:
\begin{itemize}
\item[(Var 1)] (Functoriality). The trace morphism is functorial in $F$.

\item[(Var 2)] (Compatibility with base change). For every
Cartesian square
$$
\xymatrix{
X' \ar[r]_{g'}\ar[d]^{f'}& X\ar[d]^{f}\\
Y' \ar[r]_{g} & Y
}
$$
with\pageoriginale $Y'$ quasi-compact and quasi-separated, the diagram
$$
\xymatrix@R=1.25cm{
R^{2d}f'_!f'^{\ast}g^{\ast}F(d) \ar@{=}[r]\ar[d]^{Tr_{f'}}& R^{2d}f'_!
g'^{\ast} f^{\ast} F(d) & g^{\ast}
R^{2d}f_!f^{\ast} F(d) \ar[l]_{\sim} \ar[d]^{g^{\ast}(Tr_f)}\\
g^{\ast}F \ar@{=}[rr]&            & g^{\ast}F
}
$$
in which the horizontal isomorphism is the base-change arrow
\emph{\eqhyperrefext[XVII][XVII.eq5.2.2.1]{5.2.2.1}}, is commutative.

\item[(Var 3)] (Compatibility with composition). For every pair of
composable $Z$-compactifiable morphisms
$$
X\xrightarrow{g}Y\xrightarrow{f}Z{\quoi},
$$
integers $d$ and $e$ such that $f$ satisfies $(\ast)_{d}$ and $g$
satisfies $(\ast)_{e}$, and torsion sheaf $F$ on $Z$, the
composite morphism $fg$ satisfies $(\ast)_{d+e}$ and the diagram
$$
\xymatrix@C=2.25cm{
R^{2d}f_! R^{2e} g_! g^{\ast} f^{\ast}F(d)(e) \ar@{<->}[d]_{\text{\rotatebox{90}{$\backsim$}}} \ar[r]^(.6){R^{2d}f_!(\Tr_g)} &
R^{2d}f_! f^{\ast}F(d) \ar[d]^{\Tr_f}\\
R^{2(d+e)} (fg)_! (fg)^{\ast} F(d+e) \ar[r]^(.65){\Tr_{fg}} & F
}
$$
in which the vertical isomorphism is the arrow \eqhyperref[XVIII.eq2.7.2]{2.7.2}, is
commutative.

\item[(Var 4)] (Normalizations).
\begin{itemize}
\item[(I)] If $f$ is finite locally free of rank $n$ and $d=0$,
the composite morphism
$$
F\longrightarrow
f_{\ast}f^{\ast}F=f_{!}f^{\ast}F\xrightarrow{\tr_{f}}F
$$
is multiplication by $n$.

\item[(II)] If\pageoriginale $f$ is the standard affine line
$f:\EE^{1}_{Y}\rightarrow Y$, and if $d=1$, the trace morphism
$$
\Tr_{f}: R^{2}f_{!}f^{\ast}F(1)\longrightarrow F
$$
is the isomorphism \eqhyperref[XVIII.eq2.8.1]{2.8.1}.
\end{itemize}
\end{itemize}
\end{enonce*}

\begin{enumerate}
\renewcommand{\theenumi}{\alph{enumi}}
\renewcommand{\labelenumi}{\theenumi)}
\item For $d=0$, condition $(\ast)_{d}$ says that $f$ is
quasi-finite and flat of finite presentation. By (\hyperrefext[XVII][XVII.thm6.2.3]{6.2.3}), for
$d=0$, the trace morphism must coincide with the morphism constructed
in \loccit

When $f$ is the projection of standard affine space
$\AAA^{d}_{S}$ onto $S$, it follows from (Var~4)~(II) and (Var~3) that
$\Tr_{f}$ must be the isomorphism \eqhyperref[XVIII.eq2.8.1]{2.8.1}.

\item Let $f:X\rightarrow Y$ be a compactifiable morphism satisfying the
condition $(\ast)_{d}$: there exists a quasi-finite flat
$Y$-morphism of finite presentation $u$ from $X$ to the standard vector bundle
$a^{d}:\EE^{d}_{Y}\rightarrow Y$.

Let $t(f,u)$ be the ``composite'' morphism of the trace morphisms
\eqhyperref[XVIII.eq2.8.1]{2.8.1} and (\SGA \hyperrefext[XVII][XVII.thm6.2.3]{6.2.3}) (\cf (Var 3))
$$
t(f,u)=\Tr_{a^d}\circ R^{2d}a^{d}_{!}(\Tr_{u}){\quoi}.
$$

This morphism $t(f,u)$ satisfies (Var 1) and (Var 2). By base
change, to verify that $t(f,u)$ depends only on $f$ and $d$, and
not on $u$, it suffices to verify this when $Y$ is the spectrum of an
algebraically closed field $k$. By \hyperref[XVIII.lem2.2]{2.2}, the question is
local on $X$,
so by \hyperref[XVIII.lem2.6]{2.6}, it suffices to show that
$t(f,u)=t(f,v)$
when $u$ and $v$ differ in only one coordinate, which, by
\hyperref[XVIII.subsubsec2.8.2]{2.8.2}, may be assumed to be the first. Setting
$E=\EE^{d-1}_{Y}$, we thus have a commutative diagram
$$
\xymatrix{
                              &    &\EE^1_{E} \ar[rdd]_(.3){a} &  &  \\
X \ar[rrrr]_{v}\ar[rru]^{u} \ar[rdd]_{f} \ar[rrrd]_{x}&    &                 &
                              &\EE^{1}_E \ar[ld]_{a}\\
                              &    &                 &E \ar[lld]^w& \\
                              & Y  &                 &  &
}
$$
By\pageoriginale construction (\hyperref[XVIII.subsec2.8]{2.8}),
$$
\Tr_{wa}=\Tr_{w}\circ R^{2(d-1)}w_!(\Tr_{a}){\quoi}\;;
$$
we deduce that $t(f,u)$ is the ``composite'' of $\Tr_{w}$ and
$t(x,u)$:
\begin{align*}
t(f,u) &= \Tr_{wa}\circ R^{2d}(wa)_{!}(\Tr_{u})\\
       &= \Tr_{w}\circ R^{2(d-1)}w_{!}(\Tr_{a})\circ
       R^{2d}(wa)_{!}(\Tr_{u})\\
       &= \Tr_{w}\circ R^{2(d-1)}w_{!}(\Tr_{a}\circ
       R^{2}a_{!}(\Tr_{u})){\quoi}.
\end{align*}
Moreover (\hyperref[XVIII.lem1.1.7]{1.1.7}), $t(x,u)$ is precisely the morphism
$\Tr_{x}$ of \hyperref[XVIII.prop1.1.6]{1.1.6}; hence $t(x,u)=t(x,v)=\Tr_{x}$
and $t(f,u)=t(f,v)$.

We denote by $\Tr_{f}$ the morphism $t(f,u)$ constructed
above. By \hyperrefext[XVII][XVII.thm6.2.3]{6.2.3}, this trace morphism satisfies (Var 1), (Var 2),
(Var 3) for $e=0$, and (Var 4). Whenever it is defined, it necessarily
coincides with the trace morphism sought, if the latter exists.

\item Let $f:X\rightarrow Y$ be a compactifiable Cohen--Macaulay
morphism, flat of finite presentation and pure of relative dimension
$d$. Let $(U_{i})_{i\in I}$ be a cover of $X$ by quasi-compact
open subsets such that there exists a quasi-finite flat
$Y$-morphism from $U_{i}$ to $\EE^{d}_{Y}$ (\hyperref[XVIII.lem2.4]{2.4}). By
\hyperref[XVIII.lem2.2]{2.2} and b), there
is then a unique morphism $\Tr_{f}$ making the following diagram
commutative, where the notation of \hyperref[XVIII.lem2.2]{2.2} is used:
$$
\xymatrix{
\bigoplus_{i,j}
R^{2d}(fu_{ij})_!(fu_{ij})^{\ast} F(d) \ar@<2pt>[r]
\ar@<-2pt>[r]&\bigoplus_{i}
R^{2d}(fu_i)_!(fu_i)^{\ast}F(d)
\ar[rdd]_{\bigoplus_{i}\Tr_{fu_i}}\ar[r]
&
R^{2d}f_!f^{\ast}F(d) \ar[dd]^{\Tr_f}\\
& & \\
& & F
}
$$

This trace morphism does not depend on the chosen open cover
(compare the two with the sum cover); it satisfies the conditions
(Var 1), (Var 2), (Var 3) for $e=0$, and (Var 4), as follows
immediately from b).

\item Let $f:X\rightarrow Y$ be a compactifiable morphism satisfying
$(\ast)_{d}$, and let $j:U\hookrightarrow X$\pageoriginale be the open subset promised
by $(\ast)_{d}$. By \hyperref[XVIII.lem2.3]{2.3}, after replacing $U$
by a smaller open subset, we may suppose that $fj$ is Cohen--Macaulay and
pure of dimension $d$.

\noindent
By \hyperref[XVIII.lem2.1]{2.1}, the canonical morphism
$$
R^{2d}(fj)_{!}(fj)^{\ast}F(d)\xrightarrow{\sim}
R^{2d}f_{!}f^{\ast}F(d)
$$
is an isomorphism, and we define $\Tr_{f}$ to be the composite
$$
\Tr_{f}:R^{2d}f_{!}f^{\ast}F(d)\xleftarrow{\sim}
R^{2d}(fj)_{!}(fj)^{\ast}F(d)\xrightarrow{\Tr_{fj}}F{\quoi}.
$$
It is immediate, by (Var 3) for $e=0$ established in c), that
$\Tr_{f}$ does not depend on the choice of $U$. The trace morphism satisfies
(Var 1), (Var 2), (Var 3) for $e=0$, and (Var 4), and is the only one
that can satisfy the imposed conditions. It remains to prove that it
satisfies (Var 3), and for this purpose the preceding arguments reduce
us to considering only pairs $(f,d)$, $(g,e)$ satisfying
$(\ast)_{d}$ and $(\ast)_{e}$.

The reader who wishes to polish the arguments that follow will first verify the

\begin{enonce*}{Lemma 2.9.1}\etiquette[2.9.1]{XVIII.lem2.9.1}
Let $(f,g,h)$ be three composable $T$-compactifiable morphisms,
satisfying $(\ast)_{k}$, $(\ast)_{\ell}$, and $(\ast)_{m}$.
$$
X\xrightarrow{h}Y\xrightarrow{g}Z\xrightarrow{f}T{\quoi}.
$$
If the pairs $(f,g)$ and $(g,h)$ satisfy {\rm (Var 3)}, then
$(f,gh)$ satisfies {\rm (Var 3)} if and only if $(fg,h)$
satisfies {\rm (Var 3)}.
\end{enonce*}

The condition in the conclusion also says that
$\Tr_{fgh}$ is the ``composite'' of $\Tr_{f}$, $\Tr_{g}$, and $\Tr_{h}$.

\item Let there be a commutative diagram:
$$
\xymatrix@R=1.25cm@C=1.5cm{
X \ar[r]^{v} \ar[rd]_{g}& \EE^{e}_Y \ar[r]^{u'}\ar[d]^{c} & \EE_Z^{d+e}\ar[d]^{b} \\
            &     Y   \ar[r]^{u} \ar[rd]_{f}                       & \EE_Z^{d}\ar[d]^{a}  \\
            &                                  &     Z
}
$$
The\pageoriginale trace morphism relative to $fg$ is the
``composite'' of the trace morphisms relative to $ab$ and $u'v$; by
construction for $ab$, and by \hyperrefext[XVII][XVII.thm6.2.3]{6.2.3} for $u'v$, this is also the ``composite'' of the trace morphisms relative to $a$, $b$, $u'$, and
$v$. On the other hand, the ``composite'' of the trace morphisms relative to
$f$ and $g$ is the ``composite'' of those relative to $a$, $u$, $c$, and
$v$. We must therefore verify that, for every Cartesian diagram
\begin{equation*}
\vcenter{
\xymatrix@C=1.25cm@R=1cm{
\EE^{e}_S  \ar[d]^{b'} \ar[r]^{u'}   & \EE^{e}_T \ar[d]^{b}\\
    S      \ar[r]^{u}     &      T
}}\tag{2.9.2}
\end{equation*}\etiquette[2.9.2]{XVIII.eq2.9.2}%
with $u$ quasi-finite and flat of finite presentation, the diagram of
trace morphisms \eqhyperref[XVIII.eq2.8.1]{2.8.1} and \hyperrefext[XVII][XVII.subsec6.2]{6.2}.)
\begin{equation*}
\vcenter{
\xymatrix@C=2.25cm{
R^{2e}b_!u'_! u'^{\ast}b {}^{\ast} F(e) \ar[r]^(.55){R^{2e}b_!
  (\Tr_{u'})}\ar@{=}[d]& R^{2e}b_! b^{\ast} F(e)\ar[r]^{\Tr_b}
&  F \ar@{=}[d]\\
u_! R^{2e}b'_! b'^{\ast} u^{\ast} F(e)\ar[r]^(.55){u_!(\Tr_{b'})} & u_!
u^{\ast} F \ar[r]^{\Tr_u}& F
}}\tag{2.9.3}
\end{equation*}\etiquette[2.9.3]{XVIII.eq2.9.3}%
is commutative.
\end{enumerate}

By base change, it suffices to verify this when $T$ is the spectrum of an
algebraically closed field $k$. If $S$ is a sum of schemes $S_{i}$, we
are reduced to verifying the commutativity of \eqhyperref[XVIII.eq2.9.3]{2.9.3} only for the
diagrams \eqhyperref[XVIII.eq2.9.2]{2.9.2} with base $u_{i}:S_{i}\rightarrow T$. This makes it
possible to suppose that $S$ is the spectrum of a local Artinian
$k$-algebra $A$, of finite degree $n=[A:k]$. If we then identify sheaves on
$T$ (\resp on $\EE^{e}_{T}$) with sheaves on $S$ (\resp
on $\EE^{e}_{S}$) by the functor $u^{\ast}$ (\resp
${u'}^{\ast}$) (VIII \hyperref[VIII.thm1.1]{1.1}), the morphism $\Tr_{b}$ identifies with the
morphism $\Tr_{b'}$ (compatibility of \eqhyperref[XVIII.eq2.8.1]{2.8.1} with base
change), while $\Tr_{u}$ and $\Tr_{u'}$ identify with
multiplication by $n$, whence the assertion\nde{It was also possible to
reduce this compatibility to the case $e=1$, which is obtained by combining
\hyperref[XVIII.lem1.1.7]{1.1.7} and \hyperref[XVIII.lem1.1.8]{1.1.8}. The arguments given here
are moreover similar to those that were passed over in silence in
\hyperref[XVIII.lem1.1.8]{1.1.8}.}. This completes the proof
of \hyperref[XVIII.thm2.9]{2.9}.

We\pageoriginale have seen in the course of the proof that:

\begin{enonce*}{Proposition 2.10}\etiquette[2.10]{XVIII.prop2.10}
The trace morphisms \emph{\hyperrefext[XVII][XVII.thm6.2.3]{6.2.3},
  XVIII~\hyperref[XVIII.prop1.1.6]{1.1.6}}, and \eqhyperref[XVIII.eq2.8.1]{2.8.1} are special
  cases of the trace morphism \hyperref[XVIII.thm2.9]{2.9}.
\end{enonce*}

\begin{enonce*}[remark]{Remark 2.10.1}\etiquette[2.10.1]{XVIII.rem2.10.1}
It can be shown without difficulty that the trace morphism of
\hyperref[XVIII.thm2.9]{2.9} is an
isomorphism if and only if the geometric fibers of $f$ have
exactly one irreducible component of dimension $d$, and its ``multiplicity'' is prime to $n$. We shall not give here the direct proof of this fact, which will follow
immediately from the ``global duality theorem'' of \S \hyperref[XVIII.sec3]{3}.
\end{enonce*}

\setcounter{subsection}{10}
\subsection{}\etiquette[2.11]{XVIII.subsec2.11}
Let $f:X\rightarrow S$ and $g:Y\rightarrow S$ be two
$S$-compactifiable morphisms satisfying $(\ast)_{d}$ and
$(\ast)_{e}$, respectively.
$$
\xymatrix{
     & \fprod{X}{Y}{S} \ar[dd]^{f\times g}\ar[ld]_{\pr_1}\ar[rd]^{\pr_2} & \\
X\ar[rd]_{f} &                                        & Y\ar[ld]^{g} \\
             &    S                                   &
}
$$

The Künneth isomorphism (\hyperrefext[XVII][XVII.thm5.4.3]{5.4.3})
$$
Rf_{!}(\ZZ/n){\mathop{\otimes}^{\LL}}_{\ZZ/n}Rg_{!}(\ZZ/n)\xrightarrow{\sim}
R(f\times g)_{!}(\ZZ/n)
$$
induces an isomorphism
\begin{equation*}
R^{2d}f_{!}\ZZ/n\otimes
R^{2e}g_{!}\ZZ/n\xrightarrow{\sim} R^{2(d+e)}(f\times
g)_{!}(\ZZ/n)\tag{2.11.1}
\end{equation*}\etiquette[2.11.1]{XVIII.eq2.11.1}

It is clear that $f\times g$ satisfies $(\ast)_{d+e}$; we have the
compatibility:

\begin{enonce*}{Proposition 2.12}\etiquette[2.12]{XVIII.prop2.12}
The following diagram
$$\vcenter{\etiquette[2.12.1]{XVIII.eq2.12.1}
\xymatrix@R=1.25cm{
R^{2d}f_!\ZZ/n(d) \otimes R^{2e}g_!\ZZ/n (e)
\ar[r]^{\sim} \ar[d]^{\Tr_f\otimes \Tr_g} & R^{2(d+e)}(f\times g)_!
\ZZ/n(d+e) \ar[d]^{\Tr_{f\times g}}\\
\ZZ/n  \ar@{=}[r] & \ZZ/n,
}}\leqno{(2.12.1)}
$$
in which the first horizontal arrow is a twisted form of
\eqhyperref[XVIII.eq2.11.1]{2.11.1},\pageoriginale is commutative.
\end{enonce*}

To verify this, it suffices to use (Var 3) and the asymmetric
interpretation of the Künneth arrow used in the proof of
\hyperrefext[XVII][XVII.thm5.4.3]{5.4.3} (\hyperrefext[XVII][XVII.lem5.4.3.5]{5.4.3.5}).

\setcounter{subsection}{12}
\subsection{}\etiquette[2.13]{XVIII.subsec2.13}
Let $f:X\rightarrow S$ be an $S$-compactifiable morphism satisfying
$(\ast)_{d}$, let $\mathcal{A}$ be a sheaf of rings on $S$, and let $n$ be an
integer invertible on $S$ such that $n\mathcal{A}=0$.

For every $K\in D(S,\mathcal{A})$, we have (\hyperrefext[XVII][XVII.thm5.2.6]{5.2.6})
\begin{equation*}

% END INCLUDED SGA4_English_current_buildable_source_20260906/source/18/drafts/SGA4_XVIII_LINES_4051_4450_English.texfrag

% BEGIN INCLUDED SGA4_English_current_buildable_source_20260906/source/18/drafts/SGA4_XVIII_LINES_4451_4850_English.texfrag
K{\mathop{\otimes}^{\LL}}_{\ZZ/n}Rf_{!}\ZZ/n(d)\xrightarrow{\sim} Rf_{!}(f^{\ast}K(d)){\quoi}.\tag{2.13.1}
\end{equation*}\etiquette[2.13.1]{XVIII.eq2.13.1}
 
The trace morphism
$$
\Tr_{f}:R^{2d}f_{!}\ZZ/n(d)\longrightarrow \ZZ/n
$$
may also be interpreted as a morphism
$$
\Tr_{f}:Rf_{!}\ZZ/n(d)[2d]\longrightarrow \ZZ/n{\quoi}.
$$
Tensoring this morphism with $K$ defines, via
\eqhyperref[XVIII.eq2.13.1]{2.13.1}, a trace
morphism
\begin{equation*}
\Tr_{f}:Rf_{!}(f^{\ast}K(d)[2d])\longrightarrow K{\quoi}.\tag{2.13.2}
\end{equation*}\etiquette[2.13.2]{XVIII.eq2.13.2}

As we shall see in \hyperref[XVIII.thm3.2.5]{3.2.5}, the following theorem is
essentially equivalent to the global duality theorem in étale
cohomology.

\begin{enonce*}{Theorem 2.14}\etiquette[2.14]{XVIII.thm2.14}
Let $f:X\rightarrow S$ be a smooth compactifiable morphism
purely of relative dimension $d$, and let $n$ be an integer $\geq 1$
invertible on $S$. For every geometric point $x$ of $X$
with image $s$ in $S$, there exist an étale neighborhood $V$ of $s$ in $S$
and an étale neighborhood $U$ of $x$ in $f^{-1}(V)$
\begin{equation*}
\vcenter{
\xymatrix@C=1.5cm@R=1.25cm{
U \ar[rd]_{f'_V} \ar[r]^{j} & f^{-1}(V)   \ar[r] \ar[d]^{f_V}& X\ar[d]^{f}\\
                                             &     V           \ar[r]  & S
}}\tag{2.14.1}
\end{equation*}\etiquette[2.14.1]{XVIII.eq2.14.1}%
such\pageoriginale that the morphism
\sisi{\emph{(XVII 6.2.7.2)}}{\eqhyperrefext[XVII][XVII.eq6.2.7.4]{6.2.7.4}}
$$
R^{i}f'_{V!}\ZZ/n\longrightarrow R^{i}f_{V!}\ZZ/n
$$
is zero for $i<2d$, and the morphism
$$
\Tr_{f'_{V}}:R^{2d}f'_{V!}\ZZ/n(d)\longrightarrow \ZZ/n
$$
is an isomorphism.
\end{enonce*}

We first prove the following.

\begin{enonce*}{Lemma 2.14.2}\etiquette[2.14.2]{XVIII.lem2.14.2}
Let $\mathcal{A}$ be an abelian category, let $k$ be an integer $\geq 0$,
let $(K_{i})_{0\leq i\leq 2k}$ be objects of $D^{b}(\mathcal{A})$ such that
$H^{p}(K_{i})=0$ for $p\not\in [0,k]$, let $f_{i}:K_{i}\rightarrow
K_{i+1}$ $(0\leq i\leq 2k-1)$ be morphisms, and let $f$ be their composite. If
$H^{p}(f_{i})=0$ for $p<k$, then there exists in $D^{b}(\mathcal{A})$
a morphism $\varphi$ from the complex $H^{k}(K_{0})[-k]$, consisting of
$H^{k}(K_{0})$ placed in degree $k$, to $K_{2k}$, making
the diagram
$$
\xymatrix{
K_{0}  \ar[d]\ar[r]^{f} & K_{2k}\\
\sigma_{\geq k}(K_{0})  \ar@{=}[r] & H^{k}(K_{0})[-k] \ar[u]_{\varphi}.
}
$$
commute.
\end{enonce*}

The lemma is trivial for $k=0$. We prove it by induction on
$k$. Applying the induction hypothesis to the complexes $\sigma_{\geq
1}(K_{i})[1]$ gives a morphism
$\varphi':H^{k}(K_{0})[-k]\rightarrow \sigma_{\geq 1}K_{2k-2}$ making
the following diagram commute, where $f'=f_{2k-3}\ldots f_{0}$:
$$
\xymatrix{
K_{0}  \ar[d]\ar[r]^{f'} & \sigma_{\geq 1}K_{2k-2}            \\
\sigma_{\geq k}(K_{0})  \ar@{=}[r] & H^{k}(K_{0})[-k]. \ar[u]_{\varphi'}
}
$$

For every complex $L$, the distinguished triangle
$$
\xymatrix{
                            &\sigma_{\geq 1}(L)\ar[ld]_{+1}  &       \\
  \sigma_{\leq 0}(L) \ar[rr] &                    & L\ar[lu]
}
$$
gives,\pageoriginale for every complex $M$, an exact sequence
$$
\ldots\rightarrow \Hom(M,\sigma_{\leq 0}(L))\rightarrow \Hom
(M,L)\rightarrow \Hom (M,\sigma_{\geq
1}L)\xrightarrow{j}\Ext^{1}(M,\sigma_{\leq 0}L){\quoi}.
$$

In particular, we have a diagram
$${\def\objectstyle{\scriptscriptstyle}
\xymatrix@C=0.4cm{
\Hom(\sigma_{\geq k}K_{0}, \sigma_{\leq 0} K_{2k-2})\ar[r] \ar[d]&
\overset{\text{\ding{176}}}{\underset{\phantom{\text{\ding{176}}}}
{\Hom}}(\sigma_{\geq k}K_{0},
    K_{2k-2}) \ar[r]\ar[d] &
\overset{\text{\ding{174}}}{\underset{\phantom{\text{\ding{174}}}}
{\Hom}}(\sigma_{\geq
      k}K_{0}, \sigma_{\geq 1} K_{2k-2})\ar[r] \ar[d]&
\overset{\text{\ding{175}}}{\underset{\phantom{\text{\ding{175}}}}
{\Ext}^1}(\sigma_{\geq
      k}K_{0}, \sigma_{\leq 0} K_{2k-2})\ar[d]\\
\overset{\text{\ding{178}}}{\underset{\phantom{\text{\ding{178}}}}
{\Hom}}(K_0,\sigma_{\leq 0} K_{2k-2})\ar[r] &
\overset{\text{\ding{172}}}{\underset{\phantom{\text{\ding{172}}}}
{\Hom}}(K_0,K_{2k-2}) \ar[r] &
\overset{\text{\ding{173}\ding{177}}}{\underset{\phantom{\text{\ding{173}\ding{177}}}}
{\Hom}}(K_0, \sigma_{\geq
  1}K_{2k-2})\ar[r]& \Ext^{1}(K_0, \sigma_{\leq 0} K_{2k-2})
}}\leqno{(2.14.3)}\etiquette[2.14.3]{XVIII.eq2.14.3}
$$
which maps to the analogous diagrams for $K_{2k-1}$ and
$K_{2k}$. By hypothesis, the morphism $f'$ at \ding{172} has an image
at \ding{173} that lifts to \ding{174}. The obstruction
\ding{175} to the element \ding{174} lifting to \ding{176}
has zero image in diagram \eqhyperref[XVIII.eq2.14.3]{2.14.3} for
$K_{2k-1}$. In this new diagram, \ding{174} therefore lifts to
\ding{176}, and the difference between \ding{172} and the image of
\ding{176} has zero image at \ding{177}, and hence is the image of an
element \ding{178}. This element \ding{178} has zero image in
diagram \hyperref[XVIII.eq2.14.3]{2.14.3} for $K_{2k}$; in this new diagram,
$f$ is therefore the image of \ding{176}, which solves the problem.

We prove that, for a given value of $d$, assertion
\hyperref[XVIII.thm2.14]{2.14} implies the following.

\begin{enonce*}{Corollary 2.14.4}\etiquette[2.14.4]{XVIII.cor2.14.4}
Under the hypotheses of \hyperref[XVIII.thm2.14]{2.14}, there exists a diagram
\eqhyperref[XVIII.eq2.14.1]{2.14.1} such
that the morphism \emph{\hyperrefext[XVII][XVII.subsubsec6.2.4]{6.2.4}}, in
$D^{b}(V,\sheaf{\ZZ}/n)$, admits a factorization
$$
\xymatrix@C=.75cm@R=1cm{
Rf'_{V!}\sheaf{\ZZ}/n(d)  \ar[rr] \ar[rd]_{t}&  & Rf_{V!}\sheaf{\ZZ}/n(d)\\
                          &\sheaf{\ZZ}/n[-2d], \ar[ru]_{\varphi}
}
$$
where $t$ is the arrow
% SGA-SRC-001151
$$
t:Rf'_{V!}\ZZ/n(d)\rightarrow \sigma_{\geq
2d}Rf'_{V!}\ZZ/n(d)=H^{2d}(Rf'_{V!}\ZZ/n(d))[-2d]\xrightarrow{\Tr}\sheaf{\ZZ}/n[-2d]{\quoi}.
$$
\end{enonce*}

Set\pageoriginale $f'_{0}=f$ and construct, inductively, $4d$
diagrams \eqhyperref[XVIII.eq2.14.1]{2.14.1}
$$
\xymatrix@C=2cm@R=1.5cm{
U^{'}_{i+1} \ar[r]\ar[rd]_{f'_{i+1}}& f'^{-1}_{i}(V'_{i+1}) \ar[d]\ar[r]& U'_i \ar[d]^{f'_{i}}\\
            & V'_{i+1}  \ar[r]            & V'_i
}
$$
satisfying \hyperref[XVIII.thm2.14]{2.14}. Set $V=V'_{4d}$ and
$U_{i}=\fprod{U'_{i}}{V}{V'_{i}}$. We obtain a diagram
$$
\xymatrix@R=1.5cm@C=1cm{
U_{4d}\ar[r] \ar[rrrd]^{f_{4d}}& \cdots \ar[r]& U_1 \ar[r]\ar[rd]^{f_1}& U_0\ar[r]\ar[d]^{f^0} & X \ar[d]^{f}\\
            &              &           & V \ar@{-}[r]  & S ,
}
$$
with $U_{0}=f^{-1}(V)$, such that the arrows
$$
R^{i}f_{k+1!}\ZZ/n(d)\longrightarrow R^{i}f_{k!}\ZZ/n(d)
$$
are zero (for $i<2d$) and the arrow
$$
\Tr:R^{2d}f_{4d!}\sheaf{\ZZ}/n(d)\longrightarrow \sheaf{\ZZ}/n
$$
is an isomorphism. By \hyperref[XVIII.lem2.14.2]{2.14.2}, the corollary holds
for $U=U_{4d}$.

We are now \sisi{ready}{ready}
to prove \hyperref[XVIII.thm2.14]{2.14} by induction on
$d$. The case $d=0$ is trivial, and the case $d=1$ is precisely
\hyperref[XVIII.lemfond1.6.9]{1.6.9}. We may therefore suppose that $d\geq 2$.

Since the problem is local on $X$ and $S$, we may suppose
that $f$ admits a factorization $f=gh$, with $g$ and $h$ smooth
and compactifiable, purely of relative dimensions $d'$ and $d''$, and that $d'$,
$d''<d=d'+d''$. Apply \hyperref[XVIII.cor2.14.4]{2.14.4} to $g$ and to the
morphisms deduced from
$h$ by base change, so as to obtain a diagram
$$
\xymatrix@C=2.2cm@R=1.5cm{
  U  \ar[r]^{u} \ar[rd]^{h'_0} & f^{-1}W \ar[d]^{h_0}       \ar[r] & X \ar[d]^{h} \\
V  \ar[r]^{v} \ar[rd]^{g'_0} & g^{-1}W \ar[r] \ar[d]^{g_0}       & Y \ar[d]^{g}\\
               &    W   \ar[r]               & S,
}
$$
such\pageoriginale that the morphisms (XVII \hyperref[XVII.eq6.2.7.3]{6.2.7.3}) admit
factorizations
\begin{align*}
 Rh'_{0!}\ZZ/n(d'') & \xrightarrow{\Tr}
\ZZ/n[-2d'']\xrightarrow{\varphi} Rh_{0!}\ZZ/n(d'')\\
Rg'_{0!}\ZZ/n(d') & \xrightarrow{\Tr}\ZZ/n[-2d']\xrightarrow{\Psi} Rg_{0!}\ZZ/n(d')
\end{align*}
and such that $\Tr_{h'_0}$ and $\Tr_{g'_0}$ are isomorphisms.

Let $h_{1}$ and $h'_{1}$ be the morphisms deduced from $h_{0}$ and $h'_0$
by base change to $V$, and let $f'=g'_{0}h'_{1}$,
$f_{W}=g_{0}h_{0}$.

The following diagram is commutative:
$$
\xymatrix{
  Rf'_{!}\ZZ/n(d) \ar[rr]\ar@{=}[d] & & Rf_{W!}\ZZ/n(d) \ar@{=}[d]\\
Rg'_{0!} Rh'_1\ZZ/n(d) \ar[r]& Rg_{0!} Rh'_{0!}\ZZ/n(d) \ar[r]& Rg_{0!}
Rh_{0!}\ZZ/n(d),
}
$$
as is the diagram
$$
\xymatrix@C=2.1cm@R=1cm{
  Rg'_{0!} Rh'_{1!} \ZZ/n(d)   \ar@{=}[d] \ar@<-5ex>@/_2pc/[ddd]& &\\
Rg'_{0!} v^{\ast} Rh'_{0!} \ZZ/n(d)
 \ar[dd] \ar[r]_{Rg'_{0!}v^{\ast}(t)(d')}&
Rg'_{0!}\ZZ/n(d')[-2d''] \ar[d]^{t[-2d'']}& \\
   & \ZZ/n[-2d] \ar[d]^{\Psi[-2d'']}\\
Rg_{0!} Rh'_{0!} \ZZ/n(d) \ar[r]_{Rg_{0!}(t)(d')} \ar@/_2.2pc/[rr]& Rg_{0!}\ZZ/n(d')[-2d''] \ar[r]_{Rg_{0!}(\varphi)(d')}& Rg_{0!}Rh_{0!}\ZZ/n(d)
}
$$
For\pageoriginale the diagram
$$
\xymatrix@C=1.5cm@R=1.25cm{
U_1 \ar[r]\ar[rd]^{f'} & f^{-1}W \ar[r]\ar[d]^{f_W}& X\ar[d]^{f}\\
                       &  W \ar@{-}[r]     & S,
}
$$
the canonical arrow from $Rf'_{!}\ZZ/n(d)$ to
$Rf_{W!}\ZZ/n(d)$ therefore factors through $\ZZ/n[-2d]$; hence the
arrows from $R^{i}f'_{!}\ZZ/n(d)$ to
$R^{i}f_{W!}\ZZ/n(d)$ are zero for $i<2d$. Finally,
$\Tr_{f'}$ is an isomorphism, being the composite of the isomorphisms
$\Tr_{g'_{0}}$ and $R^{2d'}g'_{0!}\Tr_{h'_{1}}$.

This completes the proof of \hyperref[XVIII.thm2.14]{2.14}.

\section{The global duality theorem}\etiquette[3]{XVIII.sec3}\pageoriginale

\subsection{The functor $Rf^{!}$}\etiquette[3.1]{XVIII.sec3.1}

We shall need a variant of the canonical flasque resolutions.

\begin{enonce*}{Lemma 3.1.1}\etiquette[3.1.1]{XVIII.lem3.1.1}
Let $X$ be a coherent scheme (\ie quasi-compact and quasi-separated). The
functor $\varinjlim$ induces an equivalence between the category of
Ind-objects of the category of constructible sheaves of sets (\resp of
groups, \resp of abelian groups) and the category of sheaves
of sets (\resp of Ind-finite groups, \resp of torsion abelian groups).
\end{enonce*}

This follows from \hyperrefext[IX][IX.cor2.7.2]{2.7.2} and \hyperrefext[IX][IX.cor2.7.3]{2.7.3}.

\setcounter{subsubsection}{1}
\subsubsection{}\etiquette[3.1.2]{XVIII.subsubsec3.1.2}
Let $\Ff$ be a torsion abelian sheaf on a scheme $X$,
the filtered direct limit of constructible sheaves
$(\Ff_{i})_{i\in I}$, and let $X_{0}$ be a conservative set of
points of $X$. The \emph{modified canonical flasque
resolution} of $\mathcal{F}$ is defined by the formula
$$
C^{\ast}_{\ell}(\Ff)=\varinjlim_{i}C^{\ast}(\Ff_{i}){\quoi},
$$ where the $C^{\ast}(\Ff_{i})$ are the canonical flasque resolutions
of \hyperrefext[XVII][XVII.def4.2.2]{4.2.2}. By \hyperref[XVIII.lem3.1.1]{3.1.1}, for $X$
coherent, this definition is independent of the choice of the $\Ff_{i}$. It
also commutes with localization; this fact allows us to define
$C^{\ast}_{\ell}(\Ff)$ by globalization when $\Ff$ is a torsion
sheaf on an arbitrary scheme $X$, not necessarily coherent.

\emph{The complex $C^{\ast}_{\ell}(\Ff)$ is a resolution of
$\Ff$, functorial in $\Ff$, whose formation
commutes with localization and filtered direct limits; the
functors $C^{n}_{\ell}$ are exact.}

It follows that, for a sheaf of rings $\mathcal{A}$ on $X$,
the functors $C^{n}_{\ell}$ carry $\mathcal{A}$-Modules to
$\mathcal{A}$-Modules; they likewise carry torsion sheaves
to torsion sheaves. It is well known that

\begin{enonce*}{Lemma 3.1.3}\etiquette[3.1.3]{XVIII.lem3.1.3}
Let\pageoriginale $\mathcal{A}$ be an abelian $\sheaf{\mathcal{U}}$-category
having a small generating set and such that direct limits
indexed by a small filtered ordered set are representable in
$\mathcal{A}$, and are exact. A functor $F$ from $\mathcal{A}^{\circ}$
to $(\Ens)$ is representable if and only if it carries small direct
limits (arbitrary ones) to inverse limits.
\end{enonce*}

\Cf for example \SGA 4 I \hyperref[I.prop8.12.5]{8.12.5}, where $\mathcal{A}$ is not assumed
abelian, nor are filtered $\varinjlim$ assumed exact.

We shall apply this lemma to the category $\Mod(X,\mathcal{A})$ of
sheaves of modules on a ringed site $(X,\mathcal{A})$. In this special
case, if the functor $F$ is represented by the sheaf
$ℱ$, then
\begin{equation*}
ℱ(U)=F(\mathcal{A}_{U}){\quoi}.\tag{3.1.3.1}
\end{equation*}\etiquette[3.1.3.1]{XVIII.eq3.1.3.1}

This description of $ℱ$ gives a direct verification of
\hyperref[XVIII.lem3.1.3]{3.1.3} in the case at hand.

Recall that if $f:X\rightarrow S$ is a compactifiable morphism
(\hyperrefext[XVII][XVII.subsec3.2]{3.2}.), then $S$ is quasi-compact and quasi-separated (sic), and the
dimensions of the fibers of $f$ are bounded.

\begin{enonce*}{Theorem 3.1.4}\etiquette[3.1.4]{XVIII.thm3.1.4}
Let $f:X\rightarrow S$ be a compactifiable morphism and let $\mathcal{A}$
be a sheaf of torsion rings on $S$. Then the functor
\begin{equation*}
Rf_{!}:D(X,f^{\ast}\mathcal{A})\longrightarrow D(S,\mathcal{A})\tag{3.1.4.1}
\end{equation*}\etiquette[3.1.4.1]{XVIII.eq3.1.4.1}%
admits a partial right adjoint
\begin{equation*}
Rf^{!}:D^{+}(S,\mathcal{A})\longrightarrow D^{+}(X,f^{\ast}\mathcal{A}):\tag{3.1.4.2}
\end{equation*}\etiquette[3.1.4.2]{XVIII.eq3.1.4.2}%
for $K\in \Ob D(X,f^{\ast}\mathcal{A})$ and $L\in \Ob
D^{+}(S,\mathcal{A})$, there is a functorial isomorphism
\begin{equation*}
\Hom(Rf_{!}K,L)\simeq \Hom(K,Rf^{!}L){\quoi}.\tag{3.1.4.3}
\end{equation*}\etiquette[3.1.4.3]{XVIII.eq3.1.4.3}%
Endowed with the translation morphism making the diagram
\begin{equation*}
\vcenter{
\xymatrix{
  \Hom(Rf_!(K[1]),L)
  \ar@{<->}[d]_{\text{\rotatebox{90}{$\backsim$}}}\ar@{<->}[r]^-{\sim}&
  \Hom(K[1],Rf^{!}L)\ar@{-}[r]^-\sim&
\Hom(K,(Rf^{!}L)[-1])\ar@{<->}[d]^{\text{\rotatebox{90}{$\backsim$}}}\\
\Hom((Rf_!K)[1],L)\ar@{-}[r]^-\sim & \Hom(Rf_!K,L[-1])
\ar@{<->}[r]^-{\sim}_-{\txt{\rm adj.}} & \Hom(K,Rf^{!}(L[-1])),
}}\tag{3.1.4.4}
\end{equation*}\etiquette[3.1.4.4]{XVIII.eq3.1.4.4}%
commute,\pageoriginale this functor $Rf^{!}$ is a triangulated functor.
\end{enonce*}

N.B. The notation $Rf^{!}$ is abusive in that $Rf^{!}$ is not in
general the derived functor of a functor $f^{!}$.

\noindent
\emph{Proof} Let $d$ be an integer such that every fiber of $f$
has dimension $<d$, and choose a compactification
$$
\xymatrix@R=1.25cm@C=1.5cm{
X \ar@{^{(}->}[r]^{j} \ar[d]_{f}& \bar{X}\ar[ld]^{\bar{f}}\\
S.
}
$$

If $K$ is a complex of $f^{\ast}\mathcal{A}$-Modules on $X$, then
$$
Rf_{!}K\simeq R\bar{f}_{\ast}(j_{!}K){\quoi}.
$$

For every sheaf $F$ on $X$, the components $F^{i}$ of the complex
$\tau_{\leq 2d}C^\ast_{\ell}j_{!}F$ satisfy
\begin{equation*}
R^{k}\bar{f}_{\ast}F^{i}=0\quad \text{\rm for}\quad k>0{\quoi}.\tag{3.1.4.5}
\end{equation*}\etiquette[3.1.4.5]{XVIII.eq3.1.4.5}%
Indeed, for $i\neq 2d$, $F^{i}$ is a direct limit of flasque sheaves,
and for $i=2d$ and $k>0$, we have (\hyperrefext[XVII][XVII.cor5.2.8.1]{5.2.8.1})
$$
R^{k}\bar{f}_{\ast}F^{2d}=R^{k+2d}f_{!}F=0{\quoi}.
$$

The simple complex associated with the truncated modified canonical
flasque resolution double complex of $K$ is a resolution of $K$. By
\eqhyperref[XVIII.eq3.1.4.5]{3.1.4.5}, we therefore have (notation of \hyperrefext[XVII][XVII.subsubsec1.1.15]{1.1.15})
\begin{equation*}
R\bar{f}_{\ast}(j_{!}K)\xleftarrow{\sim} \bar{f}_{\ast}\tau''_{\leq
2d} C^{\ast}_{\ell}j_{!}K{\quoi}.\tag{3.1.4.6}
\end{equation*}\etiquette[3.1.4.6]{XVIII.eq3.1.4.6}

Denote by $f^{\bigdot}_{!}$ the functor from sheaves of modules
on $X$ to complexes\pageoriginale of sheaves on $S$ defined
by
\begin{equation*}
f^{\bigdot}_{!}(F)=\bar{f}_{\ast}\tau_{\leq
2d}C^{\ast}_{\ell}j_{!}F{\quoi}.\tag{3.1.4.7}
\end{equation*}\etiquette[3.1.4.7]{XVIII.eq3.1.4.7}

\begin{enonce*}{Lemma 3.1.4.8}\etiquette[3.1.4.8]{XVIII.lem3.1.4.8}
\begin{enumerate}
\renewcommand{\theenumi}{\roman{enumi}}
\renewcommand{\labelenumi}{(\theenumi)}
\item The functors $f^{i}_{!}$ are exact and commute with filtered
direct limits; we have $f^{i}_{!}=0$ for $i\not\in[0,2d]$.

\item Since the functor $f^{\bigdot}_{!}$ is bounded with exact
components, its derived functor \emph{(\hyperrefext[XVII][XVII.prop1.2.10]{1.2.10})} is obtained trivially and is denoted
$Rf^{\bigdot}_{!}$, or simply
$f^{\bigdot}_{!}:D(X,f^{\ast}\mathcal{A})\rightarrow
D(S,\mathcal{A})$. The arrow \eqhyperref[XVIII.eq3.1.4.6]{3.1.4.6} is an isomorphism of
functors
$$
Rf^{\bigdot}_{!}\xrightarrow{\sim} Rf_{!}{\quoi}.

% END INCLUDED SGA4_English_current_buildable_source_20260906/source/18/drafts/SGA4_XVIII_LINES_4451_4850_English.texfrag

% BEGIN INCLUDED SGA4_English_current_buildable_source_20260906/source/18/drafts/SGA4_XVIII_LINES_4851_5250_English.texfrag
$$
\end{enumerate}
\end{enonce*}

The functors $C^{i}_{\ell}$ are exact and commute with filtered
direct limits (\hyperref[XVIII.subsubsec3.1.2]{3.1.2}). They send every
sheaf to an acyclic sheaf. The functors $\bar{f}_{\ast}$
and $j_{!}$ commute with filtered direct limits. It follows that
\begin{enumerate}
\renewcommand{\theenumi}{\alph{enumi}}
\renewcommand{\labelenumi}{\theenumi)}
\item the functors $\bar{f}_{\ast}C^{i}_{\ell}j_{!}$, and hence also the
functors $f^{i}_{!}$, commute with filtered direct limits;

\item for $i<2d$, the functors
$f^{i}_{!}=\bar{f}_{\ast}C^{i}_{\ell}j_{!}$ are exact. For $i=2d$,
one verifies that $f^{i}_{!}$ is exact by \eqhyperref[XVIII.eq3.1.4.5]{3.1.4.5}.
\end{enumerate}

The remaining assertions of \hyperref[XVIII.lem3.1.4.8]{3.1.4.8} are trivial.

\hypertarget{sga4.exp.XVIII.remark.3.1.4.9}{}\begin{enonce*}[remark]{Remark 3.1.4.9}\etiquette[3.1.4.9]{XVIII.rem3.1.4.9}
The definition of $f^{\bigdot}_{!}$ is independent of the sheaf
of rings $\mathcal{A}$, and still makes sense for any abelian
sheaf. Assertion \hyperref[XVIII.lem3.1.4.8]{3.1.4.8} (i) remains valid in the
categories of torsion abelian sheaves on $X$ and $S$.
\end{enonce*}

By \hyperref[XVIII.lem3.1.4.8]{3.1.4.8} and \hyperref[XVIII.lem3.1.3]{3.1.3}, the functor
$$
f^{i}_{!}:\Mod(X,f^{\ast}\mathcal{A})\longrightarrow
\Mod(S,\mathcal{A})
$$
has a right adjoint $f^{!}_{i}$; since $f^{i}_{!}$ is exact,
$f^{!}_{i}$ sends injectives to injectives.

With the sign convention of (\hyperrefext[XVII][XVII.subsubsec1.1.12]{1.1.12}), the functors $f^{!}_{i}$
form a complex of left exact functors, zero in cohomological degree
$n\not\in [-2d,0]$:
\begin{equation*}
f^{!}_{\bigdot}:\Mod(S,\mathcal{A})\longrightarrow
C^{b}(X,f^{\ast}\mathcal{A}){\quoi}.\tag{3.1.4.10}
\end{equation*}\etiquette[3.1.4.10]{XVIII.eq3.1.4.10}

The\pageoriginale functor $f^{!}_{\bigdot}$ derives
(\hyperrefext[XVII][XVII.prop1.2.10]{1.2.10}) to a triangulated functor
\begin{equation*}
Rf^{!}:D^{+}(S,\mathcal{A})\longrightarrow
D^{+}(X,f^{\ast}\mathcal{A}){\quoi},\tag{3.1.4.11}
\end{equation*}\etiquette[3.1.4.11]{XVIII.eq3.1.4.11}%
right adjoint to the functor $Rf_{!}$
$(=Rf^{\bigdot}_{!}=Lf^{\bigdot}_{!})$ (``trivial duality
formula''). More precisely, if $I\in \Ob C^{+}(S,\mathcal{A})$
is a complex of sheaves of injective $\mathcal{A}$-modules, the
$\mathcal{A}$-modules $f^{!}_{i} I^{n}$ are injective, and there is
an isomorphism of triple complexes (\hyperrefext[XVII][XVII.subsubsec1.1.12]{1.1.12})
\begin{equation*}
\Hom^{\bigdot}(f^{\bigdot}_{!}K, I)\xrightarrow{\sim}
\Hom^{\bigdot}(K,f^{!}_{\bigdot}I){\quoi}.\tag{3.1.4.12}
\end{equation*}\etiquette[3.1.4.12]{XVIII.eq3.1.4.12}%
Passing to $H^{0}$ of the associated simple complex (computed using products)
gives the adjunction statement. We leave it to the reader to
verify the compatibility \eqhyperref[XVIII.eq3.1.4.4]{3.1.4.4}.

\hypertarget{sga4.exp.XVIII.remark.3.1.5}{}\begin{enonce*}[remark]{Remark 3.1.5}\etiquette[3.1.5]{XVIII.rem3.1.5}
Suppose that the functor $Rf^{!}$ has finite cohomological
amplitude. By \hyperrefext[XVII][XVII.prop1.2.10]{1.2.10}, the functor $f^{!}_{\bigdot}$ then admits
a derived functor
\begin{equation*}
Rf^{!}:D(S,\mathcal{A})\longrightarrow D(X,f^{\ast}\mathcal{A}){\quoi}.\tag{3.1.5.1}
\end{equation*}\etiquette[3.1.5.1]{XVIII.eq3.1.5.1}%
This derived functor is again right adjoint to the functor $Rf_{!}$
\eqhyperref[XVIII.eq3.1.4.1]{3.1.4.1}. Indeed, we have
\begin{align*}
& \Hom_{D(S)}(Rf_{!}K,L)\simeq
\varinjlim_{L\xrightarrow{\sim}L'}\Hom_{K(S)}(f^{\bigdot}_{!}K,L')\simeq
\varinjlim_{\substack{K'\xrightarrow{\sim}K\\ L\xrightarrow{\sim}
L'}}\Hom_{K(S)} (f^{\bigdot}_{!}K', L')\\
& \Hom_{D(X)}(K,Rf^{!}L)\simeq
\varinjlim_{K'\xrightarrow{\sim}K}\Hom_{K(X)}(K',Rf^{!}L)\simeq
\varinjlim_{\substack{K'\xrightarrow{\sim}K\\
L\xrightarrow{\sim}L'}}\Hom_{K(X)}(K', f^{!}_{\bigdot}L')
\end{align*}
and, by adjunction,
% SGA-SRC-001156
$$
\Hom_{K(S)}(f^{\bigdot}_{!}K',L')=H^{0}\Hom(f^{\bigdot}_{!}K',L')\simeq
H^{0}\Hom(K',f^{!}_{\bigdot}L')=\Hom_{K(X)}(K',f^{!}_{\bigdot}L'){\quoi}.
$$
\end{enonce*}

\hypertarget{sga4.exp.XVIII.definition.3.1.6}{}\begin{enonce*}{Definition 3.1.6}\etiquette[3.1.6]{XVIII.def3.1.6}
The triangulated functor
$$
Rf^{!}:D^{+}(S,\mathcal{A})\longrightarrow
D^{+}(X,f^{\ast}\mathcal{A})
$$
(or,\pageoriginale under the hypotheses of \hyperref[XVIII.rem3.1.5]{3.1.5},
$$
Rf^{!}:D(S,\mathcal{A})\longrightarrow
D(X,f^{\ast}\mathcal{A})){\quoi},
$$
right adjoint to the functor $Rf_{!}$ \eqhyperref[XVIII.eq3.1.4.1]{3.1.4.1}, is called the
\emph{exceptional inverse-image functor.}
\end{enonce*}

\hypertarget{sga4.exp.XVIII.proposition.3.1.7}{}\begin{enonce*}{Proposition 3.1.7}\etiquette[3.1.7]{XVIII.prop3.1.7}
Suppose that the compactifiable morphism $f:X\rightarrow S$ has relative
dimension $\leq d$.
\begin{enumerate}
\renewcommand{\theenumi}{\roman{enumi}}
\renewcommand{\labelenumi}{(\theenumi)}
\item If $L\in \Ob D^{+}(S,\mathcal{A})$ satisfies $H^{i}(L)=0$ for
$i\leq k$, then $H^{i}Rf^{!}L=0$ for $i\leq k-2d$.

\item For $L\in \Ob D^{b}(S,\mathcal{A})$, we have
$$
\dim \inj (Rf^{!}L)\leq \dim \inj (L){\quoi}.
$$
\end{enumerate}
\end{enonce*}

The functor $Rf^{!}$ is the right derived functor of
$f^{!}_{\bigdot}:\Mod(S,\mathcal{A})\rightarrow
C^{b}\Mod(X,f^{\ast}\mathcal{A})$. Assertion (i) follows from the fact that
$(f^{!}_{\bigdot})^{i}=0$ for $i<-2d$, and assertion (ii) follows from
the fact that $(f^{!}_{\bigdot})^{i}=0$ for $i>0$, and that the functors
$f^{!}_{i}$ send injectives to injectives.

\hypertarget{sga4.exp.XVIII.proposition.3.1.8}{}\begin{enonce*}{Proposition 3.1.8}\etiquette[3.1.8]{XVIII.prop3.1.8}
Let $f:X\rightarrow S$ be a quasi-finite compactifiable morphism.
\begin{enumerate}
\renewcommand{\theenumi}{\roman{enumi}}
\renewcommand{\labelenumi}{(\theenumi)}
\item The functor $f_{!}:\Mod(X,f^{\ast}\mathcal{A})\rightarrow
\Mod(S,\mathcal{A})$ admits a right adjoint
$f^{!}:\Mod(S,\mathcal{A})\rightarrow \Mod(X,f^{\ast}\mathcal{A})$, and
$Rf^{!}$ is the derived functor of $f^{!}$.

\item If $f$ is a closed immersion, then $f^{!}$ is the functor of
``sections with support in $X$.''

\item If $f$ is étale, then $f^{!}$ identifies with the functor
$f^{\ast}$, with
$$
\Tr_{f}:f_{!}f^{\ast}F\longrightarrow F
$$
as adjunction morphism.
\end{enumerate}
\end{enonce*}

Assertion (i) follows immediately from the special case $d=0$ of the
proof of \hyperref[XVIII.thm3.1.4]{3.1.4}: we have $f^{\bigdot}_{!}=f_{!}$ and
$f^{!}_{\bigdot}=f^{!}$. Equivalently, the existence of an
adjoint functor $f^{!}$ follows from \hyperref[XVIII.lem3.1.3]{3.1.3} and from the fact that $f_{!}$
is\pageoriginale exact and commutes with filtered direct limits;
the adjunction $(f_{!},Rf^{!})$ is a ``trivial duality
formula.''

Assertion (ii) is clear, and (iii) is \hyperrefext[XVII][XVII.prop6.2.11]{6.2.11}.

\setcounter{subsubsection}{8}
\subsubsection{}\etiquette[3.1.9]{XVIII.subsubsec3.1.9}
Let $f:X\rightarrow S$ be a compactifiable morphism and $\mathcal{A}$
a sheaf of torsion rings on $S$, and equip $X$ with the inverse-image
sheaf of rings of $\mathcal{A}$. Let $K\in \Ob
D^{-}(X,f^{\ast}\mathcal{A})$, $L\in \Ob D^{+}(X,f^{\ast}\mathcal{A})$,
and represent $L$ by a bounded-below complex of injective sheaves, so that
$$
R\Hhom(K,L)\simeq \Hhom(K,L){\quoi}.
$$
The complex $\Hhom(K,L)$ has flasque components, and therefore
\begin{equation*}
Rf_{\ast}R\Hhom(K,L)\simeq f_{\ast}\Hhom(K,L){\quoi}.\tag{3.1.9.1}
\end{equation*}\etiquette[3.1.9.1]{XVIII.eq3.1.9.1}%
Choose a compactification of $f$ and an integer $d$ as in the
proof of \hyperref[XVIII.thm3.1.4]{3.1.4}, giving a functor
$f^{\bigdot}_{!}$ \eqhyperref[XVIII.eq3.1.4.7]{3.1.4.7}.
\begin{equation*}
\vcenter{
\xymatrix@C=1.5cm@R=.5cm{
X_V \ar[rr]^{k_X} \ar@{^{(}->}[rd]_{j_V} \ar[ddd]^{f_V}&      & X\ar[ddd]^{f} \ar@{^{(}->}[rd]^{j}         \\
    & \bar{X}_V \ar[ldd]^{\bar{f}_V} \ar[rr]_(.3){k_{\bar{X}}}   &   & \bar{X} \ar[ldd]_{\bar{f}}\\
    &         &        &\\
V  \ar[rr]^{k} &      & S.
}}\tag{3.1.9.2}
\end{equation*}\etiquette[3.1.9.2]{XVIII.eq3.1.9.2}


For every $V\in \Ob S_{\text{\rm et}}$, there is an isomorphism\nde{This
morphism is given by the construction of \hyperrefext[XVII][XVII.rem4.2.6]{4.2.6}, but there is no
reason for it to be an isomorphism unless coherent choices of points have
first been made on the topoi under consideration. For this to be the case, the
square occurring in \hyperrefext[XVII][XVII.rem4.2.6]{4.2.6} must in some sense be
Cartesian. More precisely, one should fix here a family of points
$(p)_{p\in P}$
of the topos $S_{\text{\rm et}}$, and for every $V\in \Ob S_{\text{\rm et}}$ one
considers the modified canonical flasque resolution using the family of
points associated by the correspondence of \hyperrefext[IV][IV.subsec6.7]{6.7}
with the pairs $(p,\xi)$ where $p\in P$ and $\xi\in V_p$.}
\begin{equation*}
k^{\ast}f^{\bigdot}_{!}\xrightarrow{\sim}
f^{\bigdot}_{V!}k^{\ast}_{X}{\quoi},\tag{3.1.9.3}
\end{equation*}\etiquette[3.1.9.3]{XVIII.eq3.1.9.3}%
so that
\begin{equation*}
\Hhom^{\bigdot}(f^{\bigdot}_{!}K,f^{\bigdot}_{!}L)(V)\simeq
\Hhom^{\bigdot}(f^{\bigdot}_{V!}(k^{\ast}_{X}K),f^{\bigdot}_{V!}(k^{\ast}_{X}L)){\quoi}.
\end{equation*}%
The functoriality of $f^{\bigdot}_{V!}$ gives us
$$
\Hhom^{\bigdot}(K,L)(X_{V})=\Hom^{\bigdot}(k^{\ast}_{X}K,k^{\ast}_{X}L)\longrightarrow
\Hom^{\bigdot}(f^{\bigdot}_{V!}(k^{\ast}_{X}K),f^{\bigdot}_{V!}(k^{\ast}_{X}L)){\quoi},
$$
hence an arrow, at the level of complexes,
\begin{equation*}
f_{\ast}\Hhom^{\bigdot}(K,L)\longrightarrow
\Hhom^{\bigdot}(f^{\bigdot}_{!}K,f^{\bigdot}_{!}L){\quoi}.\tag{3.1.9.4}
\end{equation*}\etiquette[3.1.9.4]{XVIII.eq3.1.9.4}

By\pageoriginale \eqhyperref[XVIII.eq3.1.9.1]{3.1.9.1}, this arrow derives to
\begin{equation*}
Rf_{\ast}R\Hhom(K,L)\longrightarrow
R\Hhom(Rf_{!}K,Rf_{!}L){\quoi}.\tag{3.1.9.5}
\end{equation*}\etiquette[3.1.9.5]{XVIII.eq3.1.9.5}

Now let $K\in \Ob D^{-}(X,\mathcal{A})$ and $L\in \Ob
D^{+}(S,\mathcal{A})$. The adjunction arrow from $Rf_{!}Rf^{!}L$ to
$L$, together with \eqhyperref[XVIII.eq3.1.9.4]{3.1.9.4}, gives a composite arrow
$$
Rf_{\ast}R\Hhom(K,Rf^{!}L)\longrightarrow R\Hhom(Rf_{!}K,
Rf_{!}Rf^{!}L)\longrightarrow R\Hhom(Rf_{!}K,L){\quoi},
$$
the local form on $S$ of the adjunction isomorphism:
\begin{equation*}
Rf_{\ast}R\Hhom(K,Rf^{!}L)\longrightarrow
R\Hhom(Rf_{!}K,L){\quoi}.\tag{3.1.9.6}
\end{equation*}\etiquette[3.1.9.6]{XVIII.eq3.1.9.6}

\hypertarget{sga4.exp.XVIII.proposition.3.1.10}{}\begin{enonce*}{Proposition 3.1.10}\etiquette[3.1.10]{XVIII.prop3.1.10}
The arrow \hyperref[XVIII.eq3.1.9.6]{3.1.9.6} is an isomorphism.
\end{enonce*}

Let $V\in \Ob S_{\text{\rm et}}$ and retain the notation of
\eqhyperref[XVIII.eq3.1.9.2]{3.1.9.2}. The (``transitivity'') isomorphism
$$
k_{!}Rf_{V!}\simeq Rf_{!}k_{X!}{\quoi},
$$
whose existence is the basis of the theory of the functor $Rf_{!}$,
transposes to a \emph{localization isomorphism} (\hyperref[XVIII.prop3.1.8]{3.1.8}. (iii))
\begin{equation*}
k^{\ast}_{X}Rf^{!}\simeq Rf^{!}_{V}k^{\ast}{\quoi},\tag{3.1.10.1}
\end{equation*}\etiquette[3.1.10.1]{XVIII.eq3.1.10.1}%
represented at the level of complexes by a morphism
$$
k^{\ast}_{X}f^{!}_{\bigdot}\longrightarrow f^{!}_{V\bigdot}k^{\ast}{\quoi}.
$$

To verify that \eqhyperref[XVIII.eq3.1.9.6]{3.1.9.6} is an isomorphism, it suffices to verify
that for every $V\in \Ob S_{\text{\rm et}}$, \eqhyperref[XVIII.eq3.1.9.6]{3.1.9.6} induces an
isomorphism
\begin{align*}
\Hom &(k^{\ast}_{X}K,k^{\ast}_{X}Rf^{!}L)\simeq
 H^{0}(V,Rf_{\ast}R\Hhom (K,Rf^{!}L))\tag{3.1.10.2}\\
& \longrightarrow H^{0}(V,R\Hom(Rf_{!}K,L))\simeq
 \Hhom(k^{\ast}Rf_{!}K,k^{\ast}L)
\end{align*}\etiquette[3.1.10.2]{XVIII.eq3.1.10.2}

For $V=S$, this is precisely \eqhyperref[XVIII.eq3.1.4.3]{3.1.4.3}. The general case follows
from the following compatibility, \emph{which the editor has not verified}:

\hypertarget{sga4.exp.XVIII.lemma.3.1.10.3}{}\begin{enonce*}{Lemma 3.1.10.3}\etiquette[3.1.10.3]{XVIII.lem3.1.10.3}
The\pageoriginale diagram
% SGA-SRC-001158
$$
\xymatrix@C=2.5cm@R=1.5cm{
\Hom(k^{\ast}_XK, k^{\ast}_XRf^{!}L)
\ar[r]^{\scriptstyle{\hyperref[XVIII.eq3.1.10.2]{(3.1.10.2)}}} \ar[d]^{\scriptstyle{\hyperref[XVIII.eq3.1.10.1]{(3.1.10.1)}}}_{\text{\rotatebox{90}{$\backsim$}}}&
\Hom(k^{\ast}Rf_!K,k^{\ast}L) \ar[d]^{\txt{\rm (localization)}}_{\text{\rotatebox{90}{$\backsim$}}}\\
\Hom(k^{\ast}_XK,Rf^{!}_Vk^{\ast}L)
\ar@{<->}[r]^{\scriptstyle{\rm{\hyperref[XVIII.eq3.1.4.3]{(3.1.4.3\; on\; V)}}}}_{\sim}& \Hom(Rf_{V!}k^{\ast}_XK,k^{\ast}L)
}
$$
is commutative.
\end{enonce*}

\setcounter{subsubsection}{10}
\subsubsection{}\etiquette[3.1.11]{XVIII.subsubsec3.1.11}
Consider a Cartesian square
\begin{equation*}
\vcenter{
\xymatrix@C=2cm@R=1.5cm{
X' \ar[r]^{g'}\ar[d]^{f'}& X\ar[d]^{f}\\
S' \ar[r]^{g} & S,
}}\tag{3.1.11.1}
\end{equation*}\etiquette[3.1.11.1]{XVIII.eq3.1.11.1}%
with $f$ compactifiable, $\mathcal{A}$ a sheaf of torsion rings on
$S$, $\mathcal{B}$ a sheaf of torsion rings on
$S'$, and $M$ a bounded-above complex of
% SGA-SRC-001159
$(g^{\ast}\mathcal{A},\mathcal{B})$-bimodules. At the level of
complexes, there is then a morphism, functorial in $K$ $(K\in \Ob
C(X,f^{\ast}\mathcal{A}))$
\begin{equation*}
M\otimes_{\mathcal{A}}g^{\ast}f^{\bigdot}_{!}K\longrightarrow
f'_{!}{}^{\bigdot}
({f'}^{\ast}M\otimes_{\mathcal{A}}{g'}^{\ast}K){\quoi},\tag{3.1.11.2}
\end{equation*}\etiquette[3.1.11.2]{XVIII.eq3.1.11.2}%
which transposes to a morphism, functorial in $L$ $(L\in \Ob
C(S',\mathcal{B}))$
\begin{equation*}
g'_{\ast}\Hhom_{\mathcal{B}}({f'}^{\ast}M,f'_{\bigdot}{}^{!}L)\longrightarrow
f^{!}_{\bigdot}g_\ast\Hhom_{\mathcal{B}}(M,L){\quoi}.\tag{3.1.11.3}
\end{equation*}\etiquette[3.1.11.3]{XVIII.eq3.1.11.3}

Passing to derived functors gives a morphism, functorial in $L$, from $D^{+}$
$(S',\mathcal{B})$ to $D^{+}(X,f^{\ast}\mathcal{A})$:
\begin{equation*}
Rg'_{\ast}R\Hhom_{\mathcal{B}}({f'}^{\ast}M,Rf'{}^{!}L)\longrightarrow
Rf^{!}Rg_{\ast}R\Hom_{\mathcal{B}}(M,L){\quoi},\tag{3.1.11.4}
\end{equation*}\etiquette[3.1.11.4]{XVIII.eq3.1.11.4}%
which makes the following diagram commutative, where $K\in \Ob
D^{-}(X,f^{\ast}\mathcal{A})$ and $L\in \Ob D^{+}(S',\mathcal{B})$:
$$
\xymatrix{
\scriptscriptstyle{\Hom(K,Rg'_{\ast} R\Hhom_{\Bb}(f'^{\ast}M,Rf'^{!}L))}
\ar@{<->}[d] \ar[r]^{\text{\ding{172}}} &
\scriptscriptstyle{\Hom(K,Rf^{!}Rg_{\ast}R\Hhom_{\Bb}(M,L))}
\ar@{<->}[r]_{{\scriptscriptstyle{\rm adj}}}& \scriptscriptstyle{\Hom(Rf_!K, Rg_{\ast}
R\Hhom_{\Bb}(M,L))}\ar@{<->}[d]\\
\scriptscriptstyle{\Hom(g'^{\ast}K,R\Hhom_{\Bb}(f'^{\ast} M, Rf'^{!}L))}\ar@{<->}[d]
& &  \scriptscriptstyle{\Hom(g^{\ast}Rf_! K, R\Hhom_{\Bb} (M,L))} \ar@{<->}[d]\\
\scriptscriptstyle{\Hom(f'^{\ast}M\otimes_{\Aa}g'^{\ast}K,
Rf'^{!}L)} \ar@{<->}[r]_{\scriptscriptstyle{\rm adj}}& \scriptscriptstyle{\Hom(Rf'^{\cdot}_!(f'^{\ast}M\otimes^{L}_\Aa
g'^{\ast}K),L)}
\ar[r]^{\text{\ding{173}}} &
\scriptscriptstyle{\Hom_{\Bb}(M\mathop{\otimes}^{\LL}_\Aa g^{\ast}Rf_!K, L)}.
}
$$

In\pageoriginale this diagram, \ding{173} (deduced from the arrow
\hyperrefext[XVII][XVII.eq5.2.2.1]{5.2.2.1}) is an isomorphism (\hyperrefext[XVII][XVII.thm5.2.6]{5.2.6}). The
arrow \ding{172}
is therefore an isomorphism for every $K$, and this implies that
\eqhyperref[XVIII.eq3.1.11.4]{3.1.11.4} is an isomorphism.

The proper base change theorem thus transposes to the following

\hypertarget{sga4.exp.XVIII.proposition.3.1.12}{}\begin{enonce*}{Proposition 3.1.12}\etiquette[3.1.12]{XVIII.prop3.1.12}
The arrow \eqhyperref[XVIII.eq3.1.11.4]{3.1.11.4} is an isomorphism.
\end{enonce*}

This isomorphism is called the \emph{induction isomorphism}.

Here are three special cases.

\hypertarget{sga4.exp.XVIII.corollary.3.1.12.1}{}\begin{enonce*}{Corollary 3.1.12.1}\etiquette[3.1.12.1]{XVIII.cor3.1.12.1}
Let $f:X\rightarrow S$ be a compactifiable morphism, let $\mathcal{A}$
and $\mathcal{B}$ be two sheaves of torsion rings on $S$, and let
$\varphi:\mathcal{A}\rightarrow \mathcal{B}$ be a
homomorphism. Denote by $\rho$ the restriction-of-scalars functor from
$\mathcal{B}$ to $\mathcal{A}$. Then, for $K\in \Ob
D^{+}(S,\mathcal{B})$, we have
$$
\rho Rf^{!}K\xrightarrow{\sim} Rf^{!}(\rho K){\quoi}.
$$
\end{enonce*}

It suffices in \eqhyperref[XVIII.eq3.1.11.3]{3.1.11.3} to take $S=S'$, $X=X'$,
$\mathcal{A}=\mathcal{A}$, $\mathcal{B}=\mathcal{B}$, $M=\mathcal{B}$.

Thus the sheaf of rings plays no essential role in the
construction of $Rf^{!}$.

\hypertarget{sga4.exp.XVIII.corollary.3.1.12.2}{}\begin{enonce*}{Corollary 3.1.12.2}\etiquette[3.1.12.2]{XVIII.cor3.1.12.2}
Let $f:X\rightarrow S$ be a compactifiable morphism, let $\mathcal{A}$
be a sheaf of torsion rings on $S$, and let $K\in \Ob D^{-}$
$(S,\mathcal{A})$, $L\in \Ob D^{+}(S,\mathcal{A})$. We have the \emph{induction
formula}
\begin{gather*}
R\Hhom(f^{\ast}K,Rf^{!}L)\xrightarrow{\sim} Rf^{!}R\Hhom(K,L)\\
(\text{\rm take } S=S', X=X'){\quoi}.
\end{gather*}
\end{enonce*}

\hypertarget{sga4.exp.XVIII.corollary.3.1.12.3}{}\begin{enonce*}{Corollary 3.1.12.3}\etiquette[3.1.12.3]{XVIII.cor3.1.12.3}
% SGA-SRC-001161
For every square \hyperref[XVIII.subsubsec3.1.11]{3.1.11} and every
$L\in \Ob D^{+}(S',\mathcal{B})$, we have
\begin{gather*}
Rg'_{\ast}R{f'}^{!}L\xrightarrow{\sim} Rf^{!}Rg_{\ast}L\\
(\text{\rm take } \mathcal{B}=g^{\ast}\mathcal{A}=M){\quoi}.
\end{gather*}
\end{enonce*}

\setcounter{subsubsection}{12}
\subsubsection{}\etiquette[3.1.13]{XVIII.subsubsec3.1.13}
For\pageoriginale a composite $gh$ of $S$-compactifiable morphisms
$$
X\xrightarrow{h}Y\xrightarrow{g}Z
$$

% END INCLUDED SGA4_English_current_buildable_source_20260906/source/18/drafts/SGA4_XVIII_LINES_4851_5250_English.texfrag

% BEGIN INCLUDED SGA4_English_current_buildable_source_20260906/source/18/drafts/SGA4_XVIII_LINES_5251_5650_English.texfrag
the composition isomorphism
$$
c_{g,h}:Rg_{!}Rh_{!}\xrightarrow{\sim} R(gh)_{!}
$$
transposes to a composition morphism\nde{The transposed morphism actually goes
in the other direction, but this causes no problem since it is an
isomorphism.}
\begin{equation*}
c^{!}_{g,h}:Rh^{!}Rg^{!}\xrightarrow{\sim} R(gh)^{!}\tag{3.1.13.1}
\end{equation*}\etiquette[3.1.13.1]{XVIII.eq3.1.13.1}%
satisfying the usual cocycle condition for a triple composite. In
other words, the \emph{categories $D^{+}$
$(X,f^{\ast}\mathcal{A})$ form the fibers of a fibered and cofibered
category over the category whose objects are the $S$-schemes} $f:X\rightarrow S$, and whose arrows are
the $S$-compactifiable morphisms, \emph{the direct-image} (\resp
\emph{inverse-image}) \emph{functors being the functors} $Rf_{!}$
(\resp $Rf^{!}$). For every commutative square of compactifiable
morphisms,
$$
\xymatrix@=1.75cm{
X' \ar[r]_{g} \ar[d]^{f'}& X \ar[d]^{f}\\
S' \ar[r]_{g'}& S,
}
$$
we therefore have a cobase-change morphism (\hyperrefext[XVII][XVII.prop2.1.3]{2.1.3})
\begin{equation*}
ch^{!}:Rf'_{!}Rg^{!}\longrightarrow Rg^{!} Rf_{!}{\quoi}.\tag{3.1.13.2}
\end{equation*}\etiquette[3.1.13.2]{XVIII.eq3.1.13.2}

\subsubsection{}\etiquette[3.1.14]{XVIII.subsubsec3.1.14}
Consider a Cartesian square
$$
\xymatrix@=1.75cm{
X' \ar[r]^{g'} \ar[d]^{f'}& X \ar[d]^{f}\\
S' \ar[r]^{g}& S,
}
$$
with $f$ compactifiable, and let $\mathcal{A}$ be a sheaf of torsion
rings on $S$. Endow $S'$, $X$, and $X'$ with the inverse images of
$\mathcal{A}$ as their sheaves of rings, again denoted\pageoriginale by $\mathcal{A}$ by abuse
of notation. The base-change isomorphism
\begin{equation*}
g^{\ast}Rf_{!}\xrightarrow{\sim} Rf'_{!}{g'}^{\ast}\tag{3.1.14.1}
\end{equation*}\etiquette[3.1.14.1]{XVIII.eq3.1.14.1}%
allows us to regard the categories $D(S,\mathcal{A})$,
$D(X,\mathcal{A})$, $D(S',\mathcal{A})$, and $D(X',\mathcal{A})$ as
the fibers of a cofibered category over the category defined by the
commutative diagram
$$
\xymatrix@C=2cm@R=1cm{
X' \ar[d]& X\ar[l]\ar[d]\\
S' & S\ar[l],
}
$$
the direct-image functors being the functors $Rf_{!}$, $Rf'_{!}$,
$g^{\ast}$, and ${g'}^{\ast}$ (sic). When we restrict to the
categories $D^{+}$, we even obtain a fibered and cofibered category
over this diagram, the direct-image functors being the functors
$Rf^{!}$, $R{f'}^{!}$, $Rg_{\ast}$, and $Rg'_{\ast}$ (sic). We deduce
from this
\begin{enumerate}
\renewcommand{\theenumi}{\alph{enumi}}
\renewcommand{\labelenumi}{\theenumi)}
\item a ``transitivity isomorphism'' $(\sic)$ for the ``inverse-image''
functors
% SGA-SRC-001163
$$
Rg'_{\ast} R{f'}^{!}\xrightarrow{\sim} Rf^{!}Rg_{\ast}
$$
already obtained in \hyperref[XVIII.cor3.1.12.3]{3.1.12.3} (with the same definition by
transposition of \eqhyperref[XVIII.eq3.1.14.1]{3.1.14.1}).

\item a ``base-change morphism'' (\hyperrefext[XVII][XVII.prop2.1.3]{2.1.3})
\begin{equation*}
{g'}^{\ast}Rf^{!}\longrightarrow R{f'}^{!}g^{\ast}\tag{3.1.14.2}
\end{equation*}\etiquette[3.1.14.2]{XVIII.eq3.1.14.2}%
which, by \loccit, can be described in the following two ways:
\end{enumerate}

\begin{enumerate}
\renewcommand{\theenumi}{\alph{enumi}}
\renewcommand{\labelenumi}{\theenumi)}
\item by adjunction, giving a morphism \eqhyperref[XVIII.eq3.1.14.2]{3.1.14.2} amounts to
giving a morphism
$$
Rf'_{!}{g'}^{\ast}Rf^{!}\longrightarrow g^{\ast}{\quoi}.
$$
The base-change theorem provides such a morphism as the composite
$$
Rf'_{!}{g'}^{\ast}Rf^{!}\xleftarrow{\sim}
g^{\ast}Rf_{!}Rf^{!}\xrightarrow{g^{\ast}\ast \text{adj.}} g^{\ast}{\quoi}.
$$

\item by\pageoriginale adjunction, giving a morphism \eqhyperref[XVIII.eq3.1.14.2]{3.1.14.2}
amounts to giving a morphism
$$
Rf^{!}\longrightarrow Rg'_{\ast}R{f'}^{!}g^{\ast}{\quoi}.
$$
The isomorphism (\hyperref[XVIII.cor3.1.12.3]{3.1.12.3}) provides such a morphism as the composite
$$
Rf^{!}\xrightarrow[Rf^{!}\ast \text{adj.}]{} Rf^{!}Rg_{\ast}
g^{\ast}\xleftarrow{\sim} Rg'_{\ast}R{f'}^{!}g^{\ast}{\quoi}.
$$
\end{enumerate}

From their interpretations in terms of a fibered and cofibered
category, it is clear that \emph{the morphisms} (\hyperref[XVIII.cor3.1.12.3]{3.1.12.3}),
\eqhyperref[XVIII.eq3.1.13.2]{3.1.13.2}, \emph{and} \eqhyperref[XVIII.eq3.1.14.2]{3.1.14.2} \emph{each satisfy two
compatibilities of the type} \hyperrefext[XII][XII.prop4.4]{4.4} \emph{when $f$ or $g$ is a composite of two
morphisms\nde{These constructions can also be interpreted in the
context of exchange structures and cross functors, \cf \cite{sga4.exp.XVIII.bib.0}.}.}

The following results from general topology will be used in the
following subsection.

\subsubsection{}\etiquette[3.1.15]{XVIII.subsubsec3.1.15}
Let $f:X\rightarrow S$ be a morphism of sites. Denote by
$\Gamma(f)$ the following site:
\begin{enumerate}
\renewcommand{\theenumi}{\alph{enumi}}
\renewcommand{\labelenumi}{\theenumi)}
\item The category $\Gamma(f)$ is the category of triples
$(U,V,\varphi)$ with
$$
U\in \Ob (X), V\in \Ob S,\quad \varphi\in \Hom
(U,f^{\ast}V){\quoi}.
$$

\item A family of morphisms
$(u_{i},v_{i}):(U_{i},V_{i},\varphi_{i})\rightarrow (U,V,\varphi)$ is
covering if, in $X$, the morphisms $u_{i}:U_{i}\rightarrow U$ cover
$U$.
\end{enumerate}

We then have:

\paragraph{}\etiquette[3.1.15.1]{XVIII.par3.1.15.1}
The functor from $\Gamma(f)$ to $X$ given by
$$
(U,V,\varphi)\longmapsto U
$$
is an equivalence of sites $\gamma:X\rightarrow \Gamma(f)$.

\paragraph{}\etiquette[3.1.15.2]{XVIII.par3.1.15.2}
The functor $f^{\ast}_{\Gamma}$ from $S$ to $\Gamma(f)$
$$
V\longmapsto (f^{\ast}V,V,\Id_{f^{\ast}V})
$$
is a morphism of sites, denoted $f_{\Gamma}:\Gamma(f)\rightarrow S$, or
simply $f$ by abuse of notation. It satisfies
$f_{\Gamma}\gamma\simeq f$. The functor $f^{\ast}_{\Gamma}$ has a
left adjoint $f_{\Gamma\bigdot}$:
$$
f_{\Gamma\bigdot}:(U,V,\varphi)\longmapsto V{\quoi}.
$$
The\pageoriginale functor $f_{\Gamma\bigdot}$ is therefore a
comorphism of sites from $\Gamma(f)$ to $S$\nde{Since this terminology
is not familiar, it should be specified here that these are the notions
appearing in \hyperrefext[III][III.sec2]{2} and \hyperrefext[IV][IV.subsec4.7]{4.7} (one should bear in mind that
in the original edition, at the location of this second reference,
misprints made it impossible to understand in precisely which
direction the arrows went).},
and defines the same morphism of topoi as the morphism of sites
$f_{\Gamma}$.
\begin{equation*}
\vcenter{
\xymatrix@R=2.2cm@C=1.75cm{
X &         &\Gamma(f) \ar[ll]^{\approx}
\ar@<4pt>[ld]^{f_{\Gamma}}& \\
  &    S \ar[lu]^{f^{\ast}}  \ar@<4pt>[ru]^{f^{\ast}_{\Gamma}} & .
}}\tag{3.1.15.3}
\end{equation*}\etiquette[3.1.15.3]{XVIII.eq3.1.15.3}

We deduce the following rule for computing the inverse image of a
sheaf\nde{Here \hyperrefext[III][III.prop2.3]{2.3} or
\hyperrefext[IV][IV.eq4.7.3]{4.7.3} is being used.}

\setcounter{paragraph}{3}
\paragraph{}\etiquette[3.1.15.4]{XVIII.par3.1.15.4}
Let $\Ff$ be a sheaf on $S$, let
$f^{\bigdot}_{\Gamma}\Ff$ be the presheaf on $\Gamma(f)$ given
by
$$
f^{\bigdot}_{\Gamma}\Ff(U,V,\varphi)=\Ff(V){\quoi},
$$
and let $af^{\bigdot}_{\Gamma}\Ff$ be the associated sheaf. Then
$$
f^{\ast}\Ff(U)\simeq
af^{\bigdot}_{\Gamma}\Ff(U,V,\varphi){\quoi}.
$$

\subsubsection{}\etiquette[3.1.16]{XVIII.subsubsec3.1.16}
Let $f:X\rightarrow S$ be a morphism of sites, $\mathcal{A}$ a
sheaf of rings on $S$, $\mathcal{B}$ a sheaf of rings on $X$, and $d$
an integer. Let $K$ be an object of the following type:

\paragraph{}\etiquette[3.1.16.1]{XVIII.par3.1.16.1}
$K$ associates to every object $U$ of $X$ a complex $K(U)$ of
$(\mathcal{A},f_{|U\ast}(\mathcal{B}|U))$-modules on $S$, zero in
degrees $>d$, and $K(U)$ is a \emph{covariant} functor of $U$.

For every sheaf of $\mathcal{A}$-modules $ℱ$ on $S$, denote by
$K^{!}ℱ$ the complex of sheaves on $X$ generated by the complex of
presheaves of $\mathcal{B}$-modules
$$
U\longmapsto \Hom_{\mathcal{A}}(K(U),ℱ){\quoi}.
$$

Denote by $RK^{!}$ the derived functor of $K^{!}$
$$
RK^{!}:D^{+}(S,\mathcal{A})\longrightarrow
D^{+}(X,\mathcal{B}){\quoi}.
$$
If $U:K_{1}\rightarrow K_{2}$ is a morphism of objects \hyperref[XVIII.par3.1.16.1]{3.1.16.1},
then $U$ induces\pageoriginale morphisms of functors
$$
{}^{t}u:K^{!}_{2}\longrightarrow K^!_{1}\quaet RK^{!}_{2}\longrightarrow
RK^{!}_{1}{\quoi}.
$$

Under the preceding hypotheses, we have:

\begin{enonce*}{Proposition 3.1.17}\etiquette[3.1.17]{XVIII.prop3.1.17}
\begin{enumerate}
\renewcommand{\theenumi}{\roman{enumi}}
\renewcommand{\labelenumi}{(\theenumi)}
\item For every point $x$ of $X$ and every $L\in
\Ob D^{+}(S,\mathcal{A})$, denoting by $V(x)$ the category of
neighborhoods of $x$, we have
\begin{equation*}
ℋ^{n}(RK^{!}L)_{x}\simeq \varinjlim_{U\in \Ob V(x)}
\Ext^{n}(K(U),L){\quoi}.\tag{3.1.17.1}
\end{equation*}\etiquette[3.1.17.1]{XVIII.eq3.1.17.1}

\item We therefore have a spectral sequence
\begin{equation*}
% SGA-SRC-001166
E^{pq}_{2}=\varinjlim_{U\in \Ob V(x)}
\Ext^{p}(ℋ^{-q}(K(U)),L)\Longrightarrow
ℋ^{p+q}(RK^{!}L)_{x}{\quoi}.\tag{3.1.17.2}
\end{equation*}\etiquette[3.1.17.2]{XVIII.eq3.1.17.2}

\item If $X^{0}$ is a conservative set of points of $X$, and
$u:K_{1}\rightarrow K_{2}$ is a morphism of objects \eqhyperref[XVIII.eq3.1.11.1]{3.1.11.1}, then,
in order that
$$
{}^{t}u:RK^{!}_{2}\longrightarrow RK^{!}_{1}
$$
be an isomorphism, it is necessary and sufficient that, for every $n$
and every $x\in X^{0}$, the morphism of pro-objects of
$\Mod(S,\mathcal{A})$
\begin{equation*}
u_{x}:\mathop{``\varprojlim\mbox{''}}_{U\in \Ob V(x)}
ℋ^{n}K_{1}(U)\longrightarrow \mathop{``\varprojlim\mbox{''}}_{U\in
\Ob V(x)} ℋ^{n}K_{2}(U)\tag{3.1.17.3}
\end{equation*}\etiquette[3.1.17.3]{XVIII.eq3.1.17.3}%
be an isomorphism.
\end{enumerate}
\end{enonce*}

To prove (i), take $L$ to be a complex of injectives, so that
\eqhyperref[XVIII.eq3.1.17.1]{3.1.17.1} is immediate; it is clear that (i) $\Rightarrow$ (ii).

The sufficiency of condition \eqhyperref[XVIII.eq3.1.17.3]{3.1.17.3} in (iii) follows immediately
from \eqhyperref[XVIII.eq3.1.17.2]{3.1.17.2}. We shall not need the fact that it is necessary.
If ${}^{t}u$ is an isomorphism, then, by
(\hyperref[XVIII.par3.1.16.1]{3.1.16.1}), for every injective sheaf $I$
on $(S,\mathcal{A})$, we have
$$
\Hom(\mathop{``\varprojlim\mbox{''}}_{V(x)}ℋ^{n}K_{2}(U),I)\xrightarrow{\sim}\Hom(\mathop{``\varprojlim\mbox{''}}_{V(x)}
ℋ^{n}K_{1}(U),I){\quoi}, 
$$
and we conclude by noting that in every abelian category with enough
injectives (here, $\Mod(S,\mathcal{A})$), the functors $\Hom(\ast,I)$,
for $I$ an injective $\mathcal{A}$-module, form a conservative family
of functors on $\Pro(\mathcal{A})$.

\begin{enonce*}[remark]{Example 3.1.18}\etiquette[3.1.18]{XVIII.exem3.1.18}
Let\pageoriginale $f:X\rightarrow S$ be a compactifiable morphism of
schemes, and let $\mathcal{A}$ be a sheaf of torsion rings on $S$.
Endow $X$ with $f^{\ast}\mathcal{A}$ as its sheaf of rings. Choose an integer $d$ that bounds
the dimension of the fibers of $f$, as well as a compactification of
$f$, and set (\eqhyperref[XVIII.eq3.1.4.7]{3.1.4.7})
\begin{equation*}
K(U)=f^{\bigdot}_{!}((f^{\ast}\mathcal{A})_{U}){\quoi}.\tag{3.1.18.1}
\end{equation*}\etiquette[3.1.18.1]{XVIII.eq3.1.18.1}
\end{enonce*}

By formula \eqhyperref[XVIII.eq3.1.3.1]{3.1.3.1} for the adjoint of $f^{\bigdot}_{!}$, we have
an isomorphism
\begin{equation*}
f^{!}_{\bigdot}\simeq K^{!}:\Mod(S,\mathcal{A})\longrightarrow C^{b}\Mod(X,f^{\ast}\mathcal{A}){\quoi}.\tag{3.1.18.2}
\end{equation*}\etiquette[3.1.18.2]{XVIII.eq3.1.18.2}%
and hence
\begin{equation*}
Rf^{!}=RK^{!}:D^{+}(S,\mathcal{A})\longrightarrow
D^{+}(X,f^{\ast}\mathcal{A}){\quoi}.\tag{3.1.18.3}
\end{equation*}\etiquette[3.1.18.3]{XVIII.eq3.1.18.3}

\begin{enonce*}[remark]{Example 3.1.19}\etiquette[3.1.19]{XVIII.exem3.1.19}
Let $f:X\rightarrow S$, $\mathcal{A}$, $\mathcal{B}$, and $d$ be as
in \hyperref[XVIII.subsubsec3.1.16]{3.1.16}. Let $K$ be an object \sisi{of the}{of the}
following type.
\end{enonce*}

\setcounter{subsubsection}{19}
\setcounter{paragraph}{0}
\paragraph{}\etiquette[3.1.19.1]{XVIII.par3.1.19.1}
For every $W=(U,V,\varphi)\in \Ob \Gamma(f)$
(\hyperref[XVIII.subsubsec3.1.15]{3.1.15})\nde{By definition, $\varphi$ is an
$X$-morphism $U\to V\times_S X$. To understand this example properly, one must
adopt the equivalent viewpoint that interprets $\varphi$ as an
$S$-morphism $\varphi\colon U\to V$.}, $K(W)$ is a
complex of $(\mathcal{A}_{V},\varphi_{\ast}\mathcal{B}_{|U})$-modules on $V$,
functorial in $W$: for every morphism $(u,v):
W_{1}=(U_{1},V_{1},\varphi_{1})\to W_{2}=(U_{2},V_{2},\varphi_{2})$ it
is given by
\begin{align*}
K(u,v) &: K(W_{1})\longrightarrow v^{\ast}K(W_{2}),\quad\mbox{\ie}\\
K(u,v) &: v_{!}K(W_{1})\longrightarrow K(W_{2}){\quoi}.
\end{align*}

With every object $K$ of type (\hyperref[XVIII.par3.1.19.1]{3.1.19.1}) is associated an object $K'$
of type (\hyperref[XVIII.par3.1.16.1]{3.1.16.1})
relative to $f_{\Gamma}:\Gamma(f)\rightarrow S$, given by
\begin{equation*}
K'(U,V,\varphi)=j_{!}K(U,V,\varphi)\tag{3.1.19.2}
\end{equation*}\etiquette[3.1.19.2]{XVIII.eq3.1.19.2}%
where $j$ is the ``inclusion morphism'' of $V$ into $S$. We set
$$
K^{!}\mathop{=}_{\defn} {K'}^{!}\quaet RK^{!}\mathop{=}_{\defn}
R{K'}^{!}{\quoi}.
$$

\begin{enonce*}[remark]{Example 3.1.20}\etiquette[3.1.20]{XVIII.exem3.1.20}
With the notation of \hyperref[XVIII.exem3.1.19]{3.1.19}, take
$\mathcal{B}=f^{\ast}\mathcal{A}$ and set
\begin{equation*}
K(U,V,\varphi)=\mathcal{A}_{V}\quad\text{\rm(in degree 0)}\tag{3.1.20.1}
\end{equation*}\etiquette[3.1.20.1]{XVIII.eq3.1.20.1}
\end{enonce*}

By\pageoriginale \hyperref[XVIII.par3.1.15.4]{3.1.15.4}, we then have
isomorphisms of functors
\begin{gather*}
K^{!}\simeq f^{\ast}\tag{3.1.20.2}\\
RK^{!}\simeq f^{\ast}{\quoi}.\tag{3.1.20.3}
\end{gather*}\etiquette[3.1.20.2]{XVIII.eq3.1.20.2}
\etiquette[3.1.20.3]{XVIII.eq3.1.20.3}

\subsection{Poincaré duality}\etiquette[3.2]{XVIII.subsec3.2}

\subsubsection{}\etiquette[3.2.1]{XVIII.subsubsec3.2.1}
Let $f$ be a compactifiable flat morphism, purely of relative
dimension $d$ (or, more generally, satisfying condition
$(\ast)_{d}$ of \hyperref[XVIII.thm2.9]{2.9}):
$$
\xymatrix@C=2cm@R=1.5cm@M=.2cm{
X \ar@{^{(}->}[r]^{j} \ar[d]_{f}& \bar{X}\ar[ld]_{\bar{f}} &\\
S &.
}
$$
Let $n$ be an integer $\geq 1$ invertible on $S$, and suppose that $S$
is ringed by a sheaf of rings $\mathcal{A}$ such that
$n\mathcal{A}=0$. By \hyperref[XVIII.exem3.1.18]{3.1.18}, the functor
$Rf^{!}:D^{+}(S,\mathcal{A})\rightarrow D^{+}(X,f^{\ast}\mathcal{A})$ can be
computed by the procedure of \hyperref[XVIII.subsubsec3.1.16]{3.1.16}, for $K$ given
by \eqhyperref[XVIII.eq3.1.18.1]{3.1.18.1}.

Let $K'$ be the functor of type (\hyperref[XVIII.par3.1.19.1]{3.1.19.1}) that
associates to $(U,V,\varphi)\in \Ob \Gamma(f)$ the complex of sheaves
$\ZZ/n(-d)[-2d]$ on $V$, consisting of $\ZZ/n(-d)$ placed in degree
$2d$, and let $K''$ be the corresponding object of type
(\hyperref[XVIII.par3.1.16.1]{3.1.16.1}) under \eqhyperref[XVIII.eq3.1.19.2]{3.1.19.2}, relative to
$f:\Gamma(f)\rightarrow S$.

We may regard $K$ and $K''$ as objects of
\hyperref[XVIII.par3.1.16.1]{3.1.16.1} relative to
$\Gamma(f)\rightarrow S$. Moreover, for every $(U,V,\varphi)\in \Ob
\Gamma(f)$, the trace morphism \hyperref[XVIII.thm2.9]{2.9} defines a morphism
of sheaves on $V$
$$
\Tr:R^{2d}\varphi_{!}\ZZ/n\longrightarrow
\ZZ/n(-d){\quoi},
$$
whence, denoting by $j$ the morphism from $V$ to $S$, a morphism of
functors
\begin{gather*}
K(U)=\bar{f}_{\ast}\tau_{\leq
2d}C^{\ast}_{\ell}\ZZ/n_{U}\longrightarrow
ℋ^{2d}(K(U))[-2d]\simeq\tag{3.2.1.1}\\
R^{2d}f_{!}(\ZZ/n)_{U}\simeq j_{!}R^{2d}\varphi_{!} \ZZ/n
\xrightarrow{\Tr_{\varphi}} j_{!}
\ZZ/n(-d)=K''(U,V,\varphi){\quoi}.
\end{gather*}\etiquette[3.2.1.1]{XVIII.eq3.2.1.1}
 
A\pageoriginale twisted form of \hyperref[XVIII.exem3.1.20]{3.1.20} shows that
$$
{RK''}^{!}=Rf^{\ast}(d)[2d]{\quoi}.
$$
By \hyperref[XVIII.subsubsec3.1.16]{3.1.16}, we therefore have a composite morphism
\begin{equation*}

% END INCLUDED SGA4_English_current_buildable_source_20260906/source/18/drafts/SGA4_XVIII_LINES_5251_5650_English.texfrag

% BEGIN INCLUDED SGA4_English_current_buildable_source_20260906/source/18/drafts/SGA4_XVIII_LINES_5651_5928_English.texfrag
t_{f}:Rf^{\ast}(d)[2d]\simeq {RK''}^{!}\longrightarrow RK^{!}\simeq
Rf^{!}{\quoi}.\tag{3.2.1.2}
\end{equation*}\etiquette[3.2.1.2]{XVIII.eq3.2.1.2}

\begin{enonce*}[remark]{Remark 3.2.2}\etiquette[3.2.2]{XVIII.rem3.2.2}
If the functor $Rf^{!}$ has finite cohomological dimension, then
$t_f$ is still defined as a morphism between functors from
$D(S,\mathcal{A})$ to $D(X,f^{\ast}\mathcal{A})$ (definition in
\hyperref[XVIII.rem3.1.5]{3.1.5}).
\end{enonce*}

There is another method for constructing an arrow \eqhyperref[XVIII.eq3.2.1.2]{3.2.1.2},
by adjunction from \eqhyperref[XVIII.eq2.13.2]{2.13.2}
$$
\Tr:Rf_{!}(f^{\ast}L(d)[2d])\longrightarrow L{\quoi}.
$$
These two arrows coincide:

\begin{enonce*}{Lemma 3.2.3}\etiquette[3.2.3]{XVIII.lem3.2.3}
For $L\in \Ob D^{+}(S,\mathcal{A})$ (\resp $D(S,\mathcal{A})$ \emph{if
the functor $Rf^{!}$ has finite cohomological dimension}), the
diagram
$$
\xymatrix@C=2.5cm@R=1.75cm{
Rf_!(f^{\ast}L(d)[2d]) \ar[d]^{\Tr_f} \ar[r]^{Rf_!(t_f)}  & Rf_!
Rf^{!}L \ar[d]^{\txt{\rm adjunction}}\\
L                              \ar@{=}[r] & L
}
$$
is commutative.
\end{enonce*}

We first give a variant of the definition of $t_f$, in terms
of the functor $f^!_\bigdot$ \eqhyperref[XVIII.eq3.1.4.10]{3.1.4.10}.
Let $U\in \Ob X_{\text{\rm et}}$,
$V\in \Ob S_{\text{\rm et}}$, and $\varphi\in \Hom_{f}(U,V)$. By
definition,
$$
f^{!}_{\bigdot}(L)(U)=\Hom(f^{\bigdot}_{!}\mathcal{A}_{U},L){\quoi}\;;
$$
the trace morphism defines a morphism
$$
\Tr_{f}:f^{\bigdot}_{!}\mathcal{A}_{U}\longrightarrow
\mathcal{A}_{V}(-d)[-2d]
$$
whose transpose is
$$
t_{\varphi}:f^{!}_{\bigdot}(L)(U)\longleftarrow L(d)[2d](V){\quoi}.
$$
The\pageoriginale morphisms $t_{\varphi}$ define a morphism
$$
t_{f}:f^{\ast}L(d)[2d]\longrightarrow f^{!}_{\bigdot}L
$$
and we leave it to the reader to verify that, in the derived category,
it coincides with \eqhyperref[XVIII.eq3.2.1.2]{3.2.1.2}. This description shows that the following
diagrams are commutative at the level of complexes, for $L$ a
complex of modules:
$$
\xymatrix@R=1.25cm{
\Hom^{\bigdot}(\Aa_V,L) \ar[r] \ar[d]^{\Tr_f}& \Hom^{\bigdot}(f^{\ast}\Aa_V(d) [2d], f^{\ast}
L(d) [2d])\ar[d]^{t_f}\\
\Hom(f^{\bigdot}_!f^{\ast}\Aa_V(d)[2d],L) \ar@{-}[r]^{\sim}& \Hom^{\bigdot}(f^{\ast}\Aa_V(d)[2d], f^{!}_{\bigdot}L).
}
$$
By functoriality of the arrows in this diagram in $\mathcal{A}_{V}$,
it follows that for every sheaf, or every complex of sheaves $K$, the
following diagram is commutative:
$$
\xymatrix{
\Hom^{\bigdot}(K,L) \ar[r] \ar[d]^{\Tr_f}& \Hom^{\bigdot}(f^{\ast}K(d)[2d],
f^{\ast}L(d)[2d]) \ar[d]\\
\Hom^{\bigdot}(f^{\bigdot}_!f^{\ast}K(d)[2d],L) \ar@{<->}[r]^{\sim} &\Hom^{\bigdot}(f^{\ast}K(d)[2d], f^{!}_{\bigdot}L).
}
$$
Taking $K=L$ and $\Id:L\rightarrow L$, we obtain \hyperref[XVIII.lem3.2.3]{3.2.3}
at the level of complexes.

I would be grateful to anyone who has understood this proof if they
would explain it to me.

\setcounter{subsubsection}{3}
\subsubsection{}\etiquette[3.2.4]{XVIII.subsubsec3.2.4}
We deduce at once from this lemma and from the definition
(\hyperref[XVIII.subsubsec3.1.13]{3.1.13}) of the
composition isomorphisms by adjunction that if $f$ is the composite
of two flat, finitely presented, compactifiable morphisms that are
purely of relative dimensions $d_{1}$ and $d_{2}:f=gh$, and if
$d=d_{1}+d_{2}$, then the diagram
$$
\xymatrix{
h^{\ast}(g^{\ast}L(d_2)[2d_2])(d_1)[2d_1]
\ar@{-}[r]^(.6){\sim}
\ar[d]^{t_h\ast t_g}&(gh)^{\ast}L(d)[2d] \ar[d]^{t_{gh}}\\
Rh^{!} \; Rg^{!} \; L  \ar@{-}[r]^{\sim}&  R(gh)^{!}L
}
$$
is\pageoriginale commutative.

\begin{enonce*}{Theorem 3.2.5}[Poincaré duality]%
\etiquette[3.2.5]{XVIII.thm3.2.5}
Let $f:X\rightarrow S$ be a compactifiable smooth morphism
\emph{(thus $S$ is quasi-compact and quasi-separated (sic))}, let
$d$ be the locally constant function on $X$ ``relative dimension of
$f$'', let $n\geq 1$ be an integer invertible on $S$, and let
$\mathcal{A}$ be a sheaf of rings on $S$ satisfying
$n\mathcal{A}=0$. Then the functor
$$
Rf_{!}:D(X,f^{\ast}\mathcal{A})\longrightarrow D(S,\mathcal{A})
$$
admits as right adjoint the following functor from $D(S,\mathcal{A})$
to $D(X,f^{\ast}\mathcal{A})$:
$$
K\longmapsto f^{\ast}K(d)[2d]
$$
with the morphism \eqhyperref[XVIII.eq2.13.2]{2.13.2}
$$
\Tr_{f}:Rf_{!}f^{\ast}K(d)[2d]\longrightarrow K
$$
as adjunction arrow:
\hypertarget{sga4.exp.XVIII.equation.3.2.5.1}{}$$
\Hom(K,f^{\ast}L(d)[2d])\xrightarrow{\sim} \Hom(Rf_{!}K,L){\quoi}.
$$
\end{enonce*}

\noindent
\emph{Proof}. We reduce to the case where $d$ is constant. By the
description \hyperref[XVIII.subsubsec3.2.1]{3.2.1} of the arrow $t_{f}$
\eqhyperref[XVIII.eq3.2.1.2]{3.2.1.2}, Theorem \hyperref[XVIII.thm2.14]{2.14},
together with criterion \hyperref[XVIII.prop3.1.17]{3.1.17} (iii), implies that the arrow $t_{f}$ is an
isomorphism (for $L\in D^{+}(S,\mathcal{A})$)
$$
t_{f}:f^{\ast}L(d)[2d]\xrightarrow{\sim} Rf^{!}L{\quoi}.
$$

In particular, the functor $Rf^{!}$ has finite cohomological
dimension and is therefore defined on the unbounded derived category. The morphism
$t_{f}$ remains an isomorphism for $L\in D(S,\mathcal{A})$\nde{More
precisely, the fact
that $Rf^!$ has finite cohomological dimension makes it possible to show that one
can define the functor $Rf^!$ on the unbounded derived category
in such a way
that, for a complex $I^\bigdot$ of injective sheaves, $Rf^!I^\bigdot$
is the total complex associated with the double complex $f^!_\bigdot I^\bigdot$.
Then, for every $n\in\ZZ$, the morphism induced on the
$n$-th cohomology sheaf by $t_f$ evaluated at $I^\bigdot$
is identified with the same morphism for a stupid truncation $I^{\geq k}$ for
$k$ sufficiently small, which proves the assertion.}.

By\pageoriginale definition, $Rf^{!}$ is the right adjoint of the
functor $Rf_{!}$; by \hyperref[XVIII.lem3.2.3]{3.2.3}, the adjunction arrow
from $Rf_{!}Rf^{!}K$ to $K$ is identified,
via $t_{f}$, with $\Tr_{f}$, and this
completes the proof of \hyperref[XVIII.thm3.2.5]{3.2.5}.

Henceforth we identify, by means of $t_{f}$, the functors
$f^{\ast}L(d)[2d]$ and $Rf^{!}L$.

\setcounter{subsubsection}{5}
\subsubsection{}\etiquette[3.2.6]{XVIII.subsubsec3.2.6}
Let $S$ be the spectrum of an algebraically closed field $k$, let $X$
be a compactifiable scheme, smooth over $k$ and purely of relative
dimension $d$, let $n$ be an integer invertible in $k$, and let $F$ be a
locally constant constructible sheaf of $\ZZ/n$-modules, \ie a ``coefficient system'' killed by $n$. Denoting by
${}^\vee$ a $\ZZ/n$-dual, we have
$$
R\Hhom(F,\ZZ/n)=F^{\vee}\quad\text{\rm (in degree 0)}{\quoi}.
$$
This formula follows from the fact that $\ZZ/n$ is an
injective $\ZZ/n$-module.
% SGA-SRC-001173
The adjunction isomorphism of \hyperref[XVIII.thm3.2.5]{3.2.5}, or rather
the isomorphism deduced from it,
$$
R\Gamma\;R\Hhom(F,\ZZ/n(d)[2d])\xrightarrow{\sim}
R\Hom(R\Gamma_{c}(F),\ZZ/n){\quoi},
$$
is therefore written here as
\begin{equation*}
R\Gamma (F^{\vee}(d)[2d])\xrightarrow{\sim} R\Hom
(R\Gamma_{c}(F),\ZZ/n){\quoi}.\tag{3.2.6.1}
\end{equation*}\etiquette[3.2.6.1]{XVIII.eq3.2.6.1}%
Passing to cohomology groups, and using again the fact that
$\ZZ/n$ is an injective $\ZZ/n$-module, this gives
\begin{equation*}
H^{2d-i}(X,F^{\vee}(d))\xrightarrow{\sim}H^{i}_{c}(X,F)^{\vee}{\quoi},\tag{3.2.6.2}
\end{equation*}\etiquette[3.2.6.2]{XVIII.eq3.2.6.2}%
which is the usual form of Poincaré duality.

\begin{thebibliography}{99}\pageoriginale%

\sisi{}{%
\bibitem{sga4.exp.XVIII.bib.0} J. \textsc{Ayoub}, Les six opérations de Grothendieck et le
formalisme des cycles évanescents dans le monde motivique. I,
Astérisque 314 (2007).

\bibitem{sga4.exp.XVIII.bib.0bis} B. \textsc{Conrad}, A modern proof of Chevalley's theorem on
algebraic groups, Journal of the Ramanujan Mathematical Society 17 (2002),
p.~1--18.

\bibitem{sga4.exp.XVIII.bib.0ter} D. \textsc{Ferrand}, On the non additivity of the trace
in derived categories, \url{http://arxiv.org/abs/math/0506589}

}

\bibitem{sga4.exp.XVIII.bib.1} J. \textsc{Giraud}, Cohomologie non abélienne. Columbia
  University.

\bibitem{sga4.exp.XVIII.bib.2} A. \textsc{Grothendieck}, Le groupe de Brauer
  III. In: Dix exposés sur
la cohomologie des schémas. North Holl. Publ. Co.

\sisi{}%
{\bibitem{sga4.exp.XVIII.bib.2bis} A. \textsc{Grothendieck}, Technique de descente et théorème
d'existence en géométrie algébrique. VI. Les schémas de Picard : propriétés
générales, Séminaire Bourbaki 1961-1962, exposé nᵒ 236,
pages 221-243.}

\sisi{\bibitem{sga4.exp.XVIII.bib.3} L. \textsc{Illusie}, Thèse}%
{\bibitem{sga4.exp.XVIII.bib.3} L. \textsc{Illusie}, Complexe cotangent et déformations II.
Lecture Notes in Math. 283. Springer Verlag.}

\sisi{}{%
\bibitem{sga4.exp.XVIII.bib.3bis} F. \textsc{Knudsen}, D. \textsc{Mumford}, The
projectivity of the moduli space of stable curves. I : Preliminaries on
``det'' and ``Div'', Mathematica Scandinavica 39 (1976), p.~19--55.

\bibitem{sga4.exp.XVIII.bib.3ter} U. \textsc{Köpf},
Über eigentliche Familien algebraischer Varietäten
über affinoiden Räumen, Schriftenreihe des Mathematischen Instituts
der Universität Münster. 2. Serie, Heft 7 (1974).}

\bibitem{sga4.exp.XVIII.bib.4} D. \textsc{Mumford}, Geometric invariant
  theory. Springer-Verlag 1965.

\sisi{}{%
\bibitem{sga4.exp.XVIII.bib.4bis} F. \textsc{Orgogozo}, Isomotifs de dimension inférieure ou
égale à un, Manuscripta Mathematica 115 (2004), p.~339--360.}

\bibitem{sga4.exp.XVIII.bib.5} M. \textsc{Raynaud}, Faisceaux amples sur les schémas
  en groupes et les espaces homogènes. Lecture Notes in Math. 119. Springer Verlag.

\sisi{}{%
\bibitem{sga4.exp.XVIII.bib.5bis} M. \textsc{Raynaud}, Géométrie analytique rigide d'après
Tate, Kiehl. Mémoires de la Société Mathématique de France, tome 39--40
(1974), p.~319--327.}

\bibitem{sga4.exp.XVIII.bib.6} J.P. \textsc{Serre}, Groupes algébriques et corps de
classe. Publ. Inst. Math. Nancago. Hermann 1959.

\bibitem{sga4.exp.XVIII.bib.7} J.L. \textsc{Verdier}, Dualité dans la cohomologie des espaces
localement compacts. Sém. Bourbaki 300. Nov. 1965.

\bibitem{sga4.exp.XVIII.bib.8} J.L. \textsc{Verdier}, A duality theorem in the etale
  cohomology of schemes. Woods-hole seminar on alg. geometry. July 1964.

\bibitem{sga4.exp.XVIII.bib.9} Catégories dérivées, by J.L. \textsc{Verdier}. Notes
mimeographed by the IHÉS.

\end{thebibliography}


%2.17  - 106
%IX 3.1 - 138
%XVII  6.3.18.2 - 345
%1.2.0 - 855
%1.2.7.3 - 865
%1.2.2.11 - 892
%1.2.12 - 920
%6.3.3.1 - 1026
%1.3.12.4 - 1631, 1636, 1704
%1.3.2.7 - 1669
%1.2.9.7 - 1778
%1.2.16.2 - 1784
%2.12.3 - 4257
%3.1.9.10 -  5337
%3.2.5.1 - 5461

\ifx\danslelivre\undefined
\let\SGAXVIIIenddocument\enddocument
\else
\let\SGAXVIIIenddocument\relax
\fi
\SGAXVIIIenddocument

% END INCLUDED SGA4_English_current_buildable_source_20260906/source/18/drafts/SGA4_XVIII_LINES_5651_5928_English.texfrag


% END INCLUDED SGA4_English_current_buildable_source_20260906/source/18/18_English.tex


\makeatletter
\def\input@path{{../19/}}
\makeatother
% BEGIN INCLUDED SGA4_English_current_buildable_source_20260906/source/19/19_English.tex
\ifx\danslelivre\undefined
% Release engine: XeLaTeX (the LuaTeX compatibility layer breaks Xy-pic output).
\documentclass[10pt]{../fichiers-configuration-latex/smfart}
% BEGIN INCLUDED SGA4_English_current_buildable_source_20260906/source/fichiers-configuration-latex/paquets.tex
\usepackage{polyglossia}
%\usepackage[francais]{babel}
\usepackage{fontspec}
\usepackage[math-style=ISO,bold-style=upright,sans-style=upright]{unicode-math}
%\usepackage[mathletters]{ucs}
%\usepackage[utf8]{inputenc}
%\usepackage[T1]{fontenc}
\usepackage[unicode,hypertexnames=false,colorlinks=true,linkcolor=blue,citecolor=Sepia,urlcolor=magenta]{hyperref}
\setmainlanguage{french}
\setmainfont[Ligatures=TeX]{Linux Libertine O}
% variante possible : TeXGyrePagella ?
\setmathfont{XITS Math}

% maths
\usepackage{amsmath}
%\usepackage{amssymb}
%\usepackage{mathrsfs}
%\usepackage{stmaryrd}
%\usepackage{euscript}

% couleurs
\usepackage[usenames,dvipsnames]{xcolor}

% graphiques
\usepackage{graphicx}
\usepackage[all]{xy}
%\usepackage{tikz}
%\usetikzlibrary{matrix,arrows,calc}
%\tikzset{column sep=3em,row sep=2em,text height=1.5ex, text depth=0.25ex,>=angle 90}
%\tikzset{description/.style={fill=white,inner sep=1pt}}

% hyperliens, 'etiquettes
\usepackage{xr-hyper}
\usepackage{../fichiers-configuration-latex/smfthm}
%\ifx\cacheretiquettes\undefined
%\usepackage[inline]{showlabels} \else \fi % pour rendre visibles les 'etiquettes
%\usepackage{../fichiers-configuration-latex/makeidx} % pour faire un index
%\usepackage[pagebackref,unicode]{hyperref}  % pour voir o`u sont cit'ees les r'ef'erences
% hyperref is loaded once above with the cumulative-reader destination policy.
%\hypersetup{breaklinks=true} %permet le retour `a la ligne dans les liens trop longs. Incompatible avec pdflatex
%\renewcommand*{\backref}[1]{{\small [cit\'e page(s)~#1.]}}

\usepackage{xspace}
\usepackage{pifont} % pour expose XVII

% END INCLUDED SGA4_English_current_buildable_source_20260906/source/fichiers-configuration-latex/paquets.tex

% BEGIN INCLUDED SGA4_English_current_buildable_source_20260906/source/fichiers-configuration-latex/numerotation.tex
\renewcommand{\thefootnote}{\roman{footnote}} % pour mettre les notes en bas de page sous la forme (ii) et non (2), qui ressemble trop (?) a un exposant.

\newcounter{pageoriginal}
\newcommand{\pageoriginale}{\refstepcounter{pageoriginal}\marginpar{\footnotesize\textbf{\thepageoriginal}}\xspace}
\let\pageoriginaled\pageoriginale

\usepackage{ifthen}
\newboolean{original}
\def\sisi{\ifthenelse{\boolean{original}}}
\setboolean{original}{false}

\newcounter{cnote}
\def\nde#1{\ifthenelse{\boolean{original}}{}{\stepcounter{cnote}\def\thefootnote{\arabic{cnote}}\protect\footnote[\value{cnote}]{N.D.E.: #1}}}
\def\ndemark{\ifthenelse{\boolean{original}}{}{\stepcounter{cnote}\def\thefootnote{\arabic{cnote}}\protect\footnotemark[\value{cnote}]}}
\def\ndetext#1{\ifthenelse{\boolean{original}}{}{\def\thefootnote{\arabic{cnote}}\protect\footnotetext[\value{cnote}]{N.D.E.: #1}}}

% END INCLUDED SGA4_English_current_buildable_source_20260906/source/fichiers-configuration-latex/numerotation.tex

% BEGIN INCLUDED SGA4_English_current_buildable_source_20260906/source/fichiers-configuration-latex/macros.tex
\setmainlanguage{french}
\setotherlanguage{english}


%% Fixed, engine-independent release identifier.  The inherited dynamic Git
%% probe used Lua io.popen, exposed build-state-dependent output, and printed
%% a visible question mark outside a Git checkout.
\def\showgitstatus{English release r2, 2026-07-23}


% macros pour LaTeX
\renewcommand{\mathsf}{}

% Police différente pour U+2113 SCRIPT SMALL L
\setmathfont[range="2113]{latinmodern-math.otf}

% Give every inherited semantic label its own stable PDF destination instead
% of inheriting whichever automatic equation/page anchor happened to be live.
% The legacy label remains as a compatibility alias used by existing edges.
\makeatletter
\def\etiquette[#1]#2{%
  \Hy@raisedlink{\hyper@anchorstart{sga4.semantic.#2}\hyper@anchorend}%
  \def\@currentHref{sga4.semantic.#2}%
  \def\@currentlabel{#1}%
  \label{#2}%
}
\makeatother
\def\eqhyperref[#1]#2{(\hyperref[#1]{#2})}
\def\pagehyperref[#1]#2{(page \hyperref[#1]{#2})}
\def\hyperrefext[#1][#2]#3{\hyperref[#2]{#1 #3}}
\def\eqhyperrefext[#1][#2]#3{\hyperref[#2]{#1 #3}}     %cf. exposes 17
\def\neweqhyperrefext[#1][#2]#3{\hyperref[#2]{#1 #3}}  %cf. exposes 17

\def\shat#1{\protect\widehat{#1}}
\def\stilde#1{\protect\widetilde{#1}}
%%%%%%%%%%%%%%%%%%


\def\etcn{etc.\xspace}
\def\etco{etc...\xspace}
\def\etc{\sisi{\etco}{\etcn}}
\def\ldotsn{\sisi{...}{\etcn}}

\def\cf{{cf.}\kern.3em}
\def\Cf{{Cf.}\kern.3em}
\def\ie{{i.e.}\kern.3em}
\def\iev{{i.e.},\xspace}
\def\loccito{loc.\kern2pt cit.\xspace}
\def\loccitn{\textit{loc.\kern2pt cit.}\xspace}
\def\loccit{\sisi{\loccito}{\loccitn}}
\def\EGA{\textup{EGA}\kern.2em}
\def\SGA{\textup{SGA}\kern.3em}

\def\quoi{}

\def\quav{\quad,\quad}
\def\quaet{\quad \text{et}\quad}
\def\fibreproduct#1#2#3{#1 ×_{#3} #2}
\let\fprod\fibreproduct
\def\ccheck#1#2{\mbox{${\check{\check{\rm #1}}}^{\lower2pt\hbox{$\scriptstyle #2$}}$}}
\def\sheaf#1{\protect\mathcal{#1}}
\def\vieuxsheaf#1{{#1}} % cf. problème avec sheaf{a\Phi_i} par exemple...
\DeclareMathOperator{\Hom}{Hom}
\DeclareMathOperator{\can}{can}
\DeclareMathOperator{\Ind}{Ind}
\DeclareMathOperator{\ind}{ind}
%\DeclareMathOperator{\Hhomtop}{\mathcal{H}\!\mathit{omtop}}
\DeclareMathOperator{\Hhomtop}{\mathbf{Homtop}}
\DeclareMathOperator{\Point}{Point}
\DeclareMathOperator{\Pro}{Pro}
\DeclareMathOperator{\Top}{Top}
\DeclareMathOperator{\Mod}{Mod}
\DeclareMathOperator{\Cart}{Cart}
%\DeclareMathOperator{\Mmorsite}{\mathcal{M}\!\mathit{orsite}}
\DeclareMathOperator{\Mmorsite}{\mathbf{Morsite}}
\DeclareMathOperator{\loc}{loc}
\DeclareMathOperator{\Loc}{Loc}
\DeclareMathOperator{\Pic}{Pic}
\DeclareMathOperator{\Et}{Et}
\DeclareMathOperator{\card}{card}
\DeclareMathOperator{\Ens}{Ens}
\DeclareMathOperator{\h}{\underline{H}}
\DeclareMathOperator{\Fonct}{Fonct}
\DeclareMathOperator{\Ker}{Ker}
\DeclareMathOperator{\but}{but}
\DeclareMathOperator{\Ob}{Ob}
\DeclareMathOperator{\ob}{ob}
\DeclareMathOperator{\Lib}{Lib}
\DeclareMathOperator{\esj}{esj}
\DeclareMathOperator{\id}{id}
\DeclareMathOperator{\Id}{Id}
\DeclareMathOperator{\ad}{ad}
\DeclareMathOperator{\Iim}{Im}
\DeclareMathOperator{\pr}{pr}
\DeclareMathOperator{\Cat}{Cat}
\DeclareMathOperator{\car}{car}
\DeclareMathOperator{\SysInd}{SysInd}
\DeclareMathOperator{\Sup}{Sup}
\DeclareMathOperator{\Coker}{Coker}
\DeclareMathOperator{\Filt}{Filt}
\DeclareMathOperator{\acc}{acc}
\DeclareMathOperator{\sur}{sur}
\DeclareMathOperator{\cat}{cat}
\DeclareMathOperator{\Ord}{Ord}
\DeclareMathOperator{\Gr}{Gr}
\DeclareMathOperator{\Cov}{Cov}
\DeclareMathOperator{\Caf}{Caf}
\DeclareMathOperator{\coker}{coker}
\DeclareMathOperator{\coim}{coim}
\DeclareMathOperator{\im}{im}
\DeclareMathOperator{\Ouv}{Ouv}
\DeclareMathOperator{\diag}{diag}
\DeclareMathOperator{\Esp}{Esp}
\DeclareMathOperator{\TOP}{TOP}
\DeclareMathOperator{\Unif}{Unif}
\DeclareMathOperator{\Iind}{\mathbf{Ind}}
\DeclareMathOperator{\Hhomcart}{\mathbf{Homcart}}
\DeclareMathOperator{\Hhomtopess}{\mathbf{Homtopess}}
\DeclareMathOperator{\Ppoints}{\mathbf{Points}}
\DeclareMathOperator{\Ppoint}{\mathbf{Point}}
\DeclareMathOperator{\Ttors}{\mathbf{Tors}}
\DeclareMathOperator{\Ssex}{\mathbf{Sex}}
\DeclareMathOperator{\Ffib}{\mathbf{Fib}}
\DeclareMathOperator{\Vvois}{\mathbf{Vois}}
\DeclareMathOperator{\Ppointess}{\mathbf{Pointess}}
\DeclareMathOperator{\scrHom}{\mathbf{Hom}}
\DeclareMathOperator{\sheafPic}{\mathbf{Pic}}
\DeclareMathOperator{\sheafAut}{\mathbf{Aut}}
\DeclareMathOperator{\Hhomtoplocan}{\mathbf{Homtoplocan}}
\DeclareMathOperator{\Hhomtopan}{\mathbf{Homtopan}}
\DeclareMathOperator{\Pptgeom}{\mathbf{Ptgéom}}
\DeclareMathOperator{\Morsite}{Morsite}
\DeclareMathOperator{\Fib}{Fib}
\DeclareMathOperator{\Fais}{Fais}
\DeclareMathOperator{\kar}{kar}
\DeclareMathOperator{\Groupes}{Groupes}
\DeclareMathOperator{\Restr}{Restr}
\DeclareMathOperator{\Isom}{Isom}
\DeclareMathOperator{\Prol}{Prol}
\DeclareMathOperator{\Res}{Res}
\DeclareMathOperator{\Tors}{Tors}
\DeclareMathOperator{\esp}{esp}
\DeclareMathOperator{\es}{es}
\DeclareMathOperator{\od}{od}
\DeclareMathOperator{\gr}{gr}
\DeclareMathOperator{\sob}{sob}
\DeclareMathOperator{\dfn}{dfn}
\DeclareMathOperator{\Sch}{Sch}
\DeclareMathOperator{\ouv}{ouv}
\DeclareMathOperator{\cons}{cons}
\DeclareMathOperator{\Nerf}{Nerf}
\DeclareMathOperator{\Points}{Points}
\DeclareMathOperator{\Espsob}{Espsob}
\DeclareMathOperator{\Homtop}{Homtop}
\DeclareMathOperator{\Pt}{Pt}
\DeclareMathOperator{\Quot}{Quot}
\DeclareMathOperator{\glob}{glob}
\DeclareMathOperator{\Homfib}{Homfib}
\DeclareMathOperator{\Qouv}{Qouv}
\DeclareMathOperator{\Quotouv}{Quotouv}
\DeclareMathOperator{\ens}{ens}
\DeclareMathOperator{\ab}{ab}
\DeclareMathOperator{\an}{an}
\DeclareMathOperator{\fs}{fs}
\DeclareMathOperator{\red}{red}
\DeclareMathOperator{\Ab}{Ab}
\DeclareMathOperator{\cosq}{cosq}
\DeclareMathOperator{\Ann}{Ann}
\DeclareMathOperator{\Zar}{Zar}


%%%%%%%%%%%%
%%%%% SGA4-2
\DeclareMathOperator{\Sex}{Sex}
\DeclareMathOperator{\vois}{vois}
\DeclareMathOperator{\Vois}{Vois}
\DeclareMathOperator{\Ext}{Ext}
\DeclareMathOperator{\ext}{ext}
\DeclareMathOperator{\Fer}{Fer}
\DeclareMathOperator{\Add}{Add}
\DeclareMathOperator{\Cocart}{Cocart}
\DeclareMathOperator{\cosk}{cosk}
\DeclareMathOperator{\plat}{plat}

\DeclareMathOperator{\faisc}{faisc}
\DeclareMathOperator{\sk}{sk}
\DeclareMathOperator{\ch}{ch}

\def\HH{\mathbf{H}}
\def\qc{\mathrm{qc}}
\def\et{\mathrm{et}}
\def\coh{\mathrm{coh}}

\def\Hvr#1{%
\begin{array}[b]{l}
{\kern2pt\scriptstyle r}\\[-8pt]
{\kern1pt\scriptstyle \vee}\\[-5pt]
H^{#1}
\end{array}}

\def\Hvv#1{%
\begin{array}[b]{l}
\kern2pt\rotatebox{90}{${\scriptscriptstyle\ll}$}\\[-6pt]
H^{#1}
\end{array}}

\def\Hvi#1{%
\begin{array}[b]{l}
{\kern0.2pt\scriptscriptstyle \infty}\\[-8pt]
{\kern1pt\scriptstyle \vee}\\[-5pt]
H^{#1}
\end{array}}

\def\Hv#1{%
\begin{array}[b]{l}
{\kern1pt\scriptstyle \vee}\\[-5pt]
H^{#1}
\end{array}}

\DeclareMathOperator{\cart}{cart}
\DeclareMathOperator{\Ffaisc}{\mathbf{Faisc}}


\def\Ev#1{%
\begin{array}[b]{l}
{\kern1pt\scriptstyle \vee}\\[-5pt]
E'_{#1}
\end{array}}

\DeclareMathOperator{\Cont}{Cont}
\DeclareMathOperator{\Ccont}{\mathscr{C}\!\mathit{ont}}
\DeclareMathOperator{\Site}{Site}
\newcommand\MM{\mathbf{M}}
\def\limtop#1{{\displaystyle{\varprojlim_{#1}\!\mbox{\rm top }}}}
\def\Limtop{\mathop{\varprojlim\!\mbox{\rm top}}}
\DeclareMathOperator{\spec}{spec}
\def\ss{\scriptstyle{\mathscr{S}}}
\DeclareMathOperator{\Gal}{Gal}
\def\smallm{\scriptstyle{\mathscr{M}}}
\DeclareMathOperator{\End}{End}
\DeclareMathOperator{\ass}{ass}
\DeclareMathOperator{\zar}{zar}
\DeclareMathOperator{\Ogg}{Ogg}
\DeclareMathOperator{\res}{res}
\DeclareMathOperator{\tr}{tr}
\DeclareMathOperator{\Supp}{Supp}
\DeclareMathOperator{\codim}{codim}
\DeclareMathOperator{\rk}{rk}
\DeclareMathOperator{\cur}{cur}
\DeclareMathOperator{\cdqc}{cdqc}
\newcommand\DD{\mathbf{D}}
\DeclareMathOperator{\cl}{cl}
\def\JJ{\mathbf{J}}
\def\FF{\mathbf{F}}
\def\BB{\mathbf{B}}
\DeclareMathOperator{\Alget}{Alget}
\DeclareMathOperator{\Ass}{Ass}
\DeclareMathOperator{\supp}{supp}
\DeclareMathOperator{\rad}{rad}
\DeclareMathOperator{\prof}{prof}
\DeclareMathOperator{\cint}{int}
\DeclareMathOperator{\sic}{sic}
\DeclareMathOperator{\Tor}{Tor}
\DeclareMathOperator{\tors}{tors}
\DeclareMathOperator{\Sym}{Sym}
\DeclareMathOperator{\parf}{parf}
\DeclareMathOperator{\tor}{tor}
\DeclareMathOperator{\defn}{defn}
\DeclareMathOperator{\Div}{Div}
\DeclareMathOperator{\Pol}{Pol}
\DeclareMathOperator{\Fex}{Fex}
\DeclareMathOperator{\Lie}{Lie}
\DeclareMathOperator{\tot}{tot}
\DeclareMathOperator{\Aut}{Aut}
\DeclareMathOperator{\HOM}{HOM}
\DeclareMathOperator{\pch}{pch}
\DeclareMathOperator{\Ch}{Ch}
\DeclareMathOperator{\EXT}{EXT}
\DeclareMathOperator{\PIC}{PIC}
\DeclareMathOperator{\Rf}{Rf}
\DeclareMathOperator{\inj}{inj}
\newcommand\VV{\mathbf{V}}

\def\Xv{%
\begin{array}[b]{l}
{\kern1pt\scriptstyle \vee}\\[-5pt]
X
\end{array}}

%%%%%%%
\newcommand\TT{\mathbf{T}}
\newcommand\XX{\mathbf{X}}
\newcommand\EE{\mathbf{E}}
%%%%%%

\def\qedsymbol{\textup{C.Q.F.D.}}
\let\CQFD\qed
\let\cqfd\qed

%%%%%\bigdot
\font\bigsymb=cmsy10 at 4pt
\def\bigdot{{\kern1.2pt\raise 1.5pt\hbox{\bigsymb\char15}}}
%%%%%

\def\Lright{\mathop{\vcenter{\hbox{\hbox{$\longrightarrow$}\lower8pt \hbox{\kern-16pt\hbox{$\longrightarrow$}}}}}}

\def\Aa{{\mathcal A}}
\def\Bb{{\mathcal B}}
\def\Cc{{\mathcal C}}
\def\Dd{{\mathcal D}}
\def\Ee{{\mathcal E}}
\def\Ff{{\mathcal F}}
\def\Gg{{\mathcal G}}
\def\Hh{{\mathcal H}}
\def\Ii{{\mathcal I}}
\def\Jj{{\mathcal J}}
\def\Kk{{\mathcal K}}
\def\Ll{{\mathcal L}}
\def\Mm{{\mathcal M}}
\def\Nn{{\mathcal N}}
\def\Oo{{\mathcal O}}
\def\Pp{{\mathcal P}}
\def\Qq{{\mathcal Q}}
\def\Rr{{\mathcal R}}
\def\Ss{{\mathcal S}}
\def\Tt{{\mathcal T}}
\def\Uu{{\mathcal U}}
\def\Vv{{\mathcal V}}
\def\Ww{{\mathcal W}}
\def\Xx{{\mathcal X}}
\def\Yy{{\mathcal Y}}
\def\Zz{{\mathcal Z}}

\newcommand{\CC}{\mathbf{C}}
\newcommand{\GG}{\mathbf{G}}
\newcommand{\KK}{\mathbf{K}}
\newcommand{\LL}{\mathbf{L}}
\newcommand{\NN}{\mathbf{N}}
\newcommand{\PP}{\mathbf{P}}
\newcommand\QQ{\mathbf{Q}}
\newcommand\RR{\mathbf{R}}
\newcommand\ZZ{\mathbf{Z}}

\def\to{\mathchoice{\longrightarrow}{\rightarrow}{\rightarrow}{\rightarrow}}
\def\from{\mathchoice{\longleftarrow}{\leftarrow}{\leftarrow}{\leftarrow}}
\def\mto{\mathchoice{\longmapsto}{\mapsto}{\mapsto}{\mapsto}}
\def\mfrom{\mathrel{\from\joinrel\mapstochar\,}}
\def\hto{\mathrel{\lhook\joinrel\to}}
\def\hfrom{\mathrel{\from\joinrel\rhook}}

\def\implies{\mathchoice{\Longrightarrow}{\Rightarrow}{\Rightarrow}{\Rightarrow}}
\def\implied{\mathchoice{\Longleftarrow}{\Leftarrow}{\Leftarrow}{\Leftarrow}}
\def\iff{\mathchoice{\Longleftrightarrow}{\Leftrightarrow}{\Leftrightarrow}{\Leftrightarrow}}

\def\To#1{\mathchoice{\xrightarrow{\labelstyle~~#1~}}{\stackrel{#1}{\longrightarrow}}{}{}}
\def\Mto#1{\mathrel{\mapstochar\To{#1}}}
\def\From#1{\mathchoice{\xleftarrow{\labelstyle~~#1~}}{\stackrel{#1}{\longleftarrow}}{}{}}

\newcommand{\isomto}{\overset{\sim}{\longrightarrow}}
\newcommand{\isomfrom}{\overset{\sim}{\longleftarrow}}

\newcommand{\IFF}{\unskip~$\iff$~\ignorespaces}
\newcommand{\THEN}{\unskip~$\implies$~\ignorespaces}

\let\leq\leqslant
\let\geq\geqslant
\let\le\leqslant
\let\ge\geqslant

\def\cf{{cf.}\kern.3em}
\def\Cf{{Cf.}\kern.3em}
\def\ie{{i.e.}\kern.3em}
\def\iev{{i.e.},\xspace}
%\def\loccito{loc.\kern2pt cit.\xspace}
%\def\loccitn{\textit{loc.\kern2pt cit.}\xspace}
%\def\loccit{\sisi{\loccito}{\loccitn}}
\def\EGA{\textup{EGA}\kern.2em}
\def\SGA{\textup{SGA}\kern.3em}

\let\oldS\S
\def\S{\oldS\kern .15em}
\def\ptbl{.\kern .2em}
\def\resp{resp.\kern.3em}
\def\Exp{Exp.\kern.3em}

\let\bar\overline
\let\hat\widehat
\let\tilde\widetilde
\let\dpl\displaystyle
\let\ldots\dots
\let\emptyset\varnothing
\let\skpt\skippointrait
\let\epsilon\varepsilon

\def\sfrac#1#2{{#1}/{#2}}

\DeclareMathOperator{\cd}{cd}
%\DeclareMathOperator{\Hhom}{\mathcal{H}\!\mathit{om}}
\DeclareMathOperator{\Hhom}{\mathbf{Hom}}
%\DeclareMathOperator{\Eext}{\mathcal{E}\!\mathit{xt}}
\DeclareMathOperator{\Eext}{\mathbf{Ext}}
\DeclareMathOperator{\Spec}{Spec}
\DeclareMathOperator{\Tr}{Tr}

\let\emptyset=\varnothing

%%%%%%%%%%%%%%%
%%% expose XVII
%%%%%%%%%%%%%%%

\makeatletter
\let\old@optexpose\optexpose
\newcommand{\empty@optexpose}[1]{}
\newcommand{\paren@optexpose}[1]{\old@optexpose{#1}(}
\newcommand{\refnoexpose}[1]{\let\optexpose\empty@optexpose\ref{#1}\let\optexpose\old@optexpose}%\refnoexpose{X.cor5.2} engendre "5.2" au lieu de "X 5.2"
\newcommand{\neweqref}[1]{\let\optexpose\paren@optexpose\textup{\ref{#1})}\let\optexpose\old@optexpose}%\neweqref{XII.eq4.1} engendre "XII (4.1)" au lieu de "(XII 4.1)"
\makeatother
\newcommand{\refvb}[2][VI${}_{\text{B}}$ ]{\sisi{#1\refnoexpose{#2}}{\ref{#2}}}

\newcommand{\leftexp}[2]{{\vphantom{#2}}^{#1}{#2}}
\newcommand{\AAA}{\mathbf{A}}
\newcommand{\Sf}{\mathfrak{S}}
\DeclareMathOperator*{\quotprojlim}{\text{``$\varprojlim$''}}
\DeclareMathOperator*{\quotinjlim}{\text{``$\varinjlim$''}}
\DeclareMathOperator{\Fl}{F\sisi{\ell}{l}}
\DeclareMathOperator{\Pl}{P\sisi{\ell}{l}}
%\DeclareMathOperator{\cl}{c\sisi{\ell}{l}}
\DeclareMathOperator{\tdf}{tdf}
\DeclareMathOperator{\rg}{rg}
\newcommand{\iInd}{\sisi{I}{i}nd}
\newcommand{\iind}{ind}
\newcommand{\ppro}{pro}
\newcommand{\tTor}{\sisi{T}{t}or}
\newcommand{\ttor}{tor}
\newcommand{\chap}{Chap.}
\newcommand{\Ffex}{\Ff\mathit{ex}}
\newcommand{\sheafk}{\underline{k}}
\newcommand{\sheafn}{\underline{n}}
\newcommand{\sheafGamma}{\underline{\Gamma}}

% Repeated references to an already typeset footnote must not use
% \footnotemark without a matching \footnotetext: Hyperref would then create
% a live annotation to a destination that cannot exist.  This helper keeps the
% current printed footnote mark and links it to an explicit stable note target.
\makeatletter
\newcommand{\sgarepeatedfootnotemark}[1]{%
  \begingroup
  \edef\@thefnmark{\thefootnote}%
  \hyperlink{#1}{\@makefnmark}%
  \endgroup}
\makeatother

\sisi{
\DeclareMathOperator{\simarrow}{\rule{.15mm}{1.8mm}\kern -3.4pt\rightsquigarrow}
\newcommand{\sheafell}{\underline{\ell}}
\newcommand{\sheafZZ}{\underline{\ZZ}}
}{
\newcommand{\nderef}[1]{la N.D.E.~(\ref{#1}) page~\pageref{#1}}
}

%%%%%%%%%%%%%%%%%
%%%%% expose XVII
%%%%%%%%%%%%%%%%%

% END INCLUDED SGA4_English_current_buildable_source_20260906/source/fichiers-configuration-latex/macros.tex


\externaldocument{../01/01_English.tex}
\externaldocument{../02/02_English.tex}
\externaldocument{../03/03_English.tex}
\externaldocument{../04/04_English.tex}
\externaldocument{../05/05_English.tex}
\externaldocument{../05bis/05bis_English.tex}
\externaldocument{../06/06_English.tex}
\externaldocument{../07/07_English.tex}
\externaldocument{../08/08_English.tex}
\externaldocument{../09/09_English.tex}
%\externaldocument{../10/10_English.tex}
\externaldocument{../11/11_English.tex}
\externaldocument{../12/12_English.tex}
\externaldocument{../13/13_English.tex}
\externaldocument{../14/14_English.tex}
\externaldocument{../15/15_English.tex}
\externaldocument{../16/16_English.tex}
\externaldocument{../17/17_English.tex}
\externaldocument{../18/18_English.tex}

\setboolean{original}{false}

\begin{document}

\begin{center}
XIX. Cohomology of excellent preschemes of equal characteristic\\
M. Artin
\end{center}

\selectlanguage{english}version: \showgitstatus
\tableofcontents
\mainmatter

\else
\chapter{Cohomology of excellent preschemes of equal characteristic}
%\addtocontents{toc}{M. Artin}
\begin{center}
M. Artin
\end{center}
\fi
\phantomsection\MakeLinkTarget*{sga4.exp.XIX}\label{sga4.exp.XIX}

% The edition is fixed to the corrected (second) reading. This expandable
% selector is output-equivalent and avoids current Xy-pic conditional failures.
\renewcommand{\sisi}[2]{#2}

% BEGIN INCLUDED SGA4_English_current_buildable_source_20260906/source/19/drafts/SGA4_XIX_LINES_44_450_English.texfrag
\newcommand{\UU}{\mathbf{U}}

In\pageoriginale this exposé we develop the theory of étale
cohomology for excellent preschemes (\EGA IV 7.8.1) of equal
characteristic, \emph{assuming the theorem on resolution of
  singularities\pageoriginale for such preschemes}, in the following form:

Let $X=\Spec A$ be an excellent local scheme of equal characteristic,
and let $U\subset X$ be a regular open subset. There exist a proper
birational morphism $f:X'\rightarrow X$ and an open immersion
$U\rightarrow X'$ over $X$ such that $X'$ is regular and
$Y=X'-U$ is everywhere of codimension \hyperref[XIX.sec1]{1}, 
with normal crossings.

\emph{The proofs are therefore valid, “for the time being,” only in
  characteristic zero} \cite{sga4.exp.XIX.bib.4}. One could also deduce from them
results for preschemes of characteristic $p>0$ in low dimensions
\cite{sga4.exp.XIX.bib.1}.

Under the resolution hypothesis, we shall obtain more or less
satisfactory results for preschemes of equal characteristic and for
coefficients prime to the characteristic (\cf \hyperref[XIX.thm3.2]{3.2}, \hyperref[XIX.thm4.1]{4.1},
\hyperref[XIX.thm5.1]{5.1}, \hyperref[XIX.thm6.1]{6.1}). On the other hand, almost nothing is known about the cohomology
of general preschemes of unequal characteristic, even in dimension
\hyperref[XIX.sec2]{2}, where resolution is available \cite{sga4.exp.XIX.bib.2}. For example, the
purity theorem \hyperref[XIX.thm2.1]{2.1} is not known for any complete ring of unequal
characteristic of dimension $>1$. This result, for the formal power
series rings $k[[x_{1},\ldots, x_{n}]]$ (\hyperref[XIX.thm1.2]{1.2}), is one of the main
tools in the theory. Let us also note that the resolution theorem is
not used in the proof of \hyperref[XIX.thm1.2]{1.2}. \emph{This theorem is therefore
  proved in arbitrary characteristic.}

We shall often use without explicit mention the properties of
excellent\pageoriginale schemes developed in \EGA IV 7.8, such as
stability under extensions of finite type (\EGA IV 7.8.3(ii)) and
under strict localizations. The latter fact follows from resolution
and \EGA IV 7.9.5.

\section{Purity for the ring $k[[x_{1},\ldots,x_{n}]]$}\etiquette[1]{XIX.sec1}

Recall the following corollary of \hyperrefext[XVI][XVI.thm3.7]{3.7}:

\begin{enonce*}{Corollary 1.1}\etiquette[1.1]{XIX.cor1.1}
Let $A$ be a strictly local ring, and denote by $A\{x\}$ the strict
localization of $A[x]$ at the prime ideal generated by $\rad A$ and
$x$. Let $U=\Spec A\{x\}[1/x]$. Then, for $n$ prime to the residual
characteristic of $A$, we have
$$
H^{q}(U,𝛍_{n})\simeq
\begin{cases}
𝛍_{n}(A) & \text{\rm if}\quad q=0\\
\ZZ/n & \text{\rm if}\quad q=1\\
0            & \text{\rm if}\quad q>1{\quoi}.
\end{cases}
$$
\end{enonce*}

Indeed, let $X=\Spec A[x]$, and let $Y=V(x)$ be the zero locus of $x$
in $X$. Then $(Y,X)$ is a smooth pair (\hyperrefext[XVI][XVI.def3.1]{3.1}) over
$\Spec A$. By (V 4, VIII 5), $H^{q}(U,𝛍_{n})$
is identified with the fiber of the sheaf $ℋ^{q+1}_{Y}$
$(X,𝛍_{n})$ at a geometric point above the
point $(\rad A,x)$ of $X$, and the corollary therefore follows
immediately from (\hyperrefext[XVI][XVI.thm3.7]{3.7}). Note that $𝛍_{n}$ and $\ZZ/n$
are both cyclic groups of order $n$. We have used them here to give
the canonical\pageoriginale formulas (\cf \hyperref[XIX.thm3.4]{3.4}).

\begin{enonce*}{Theorem 1.2}\etiquette[1.2]{XIX.thm1.2}
Let $k$ be a separably closed field, and let $U=\Spec
k[[x_{1},\ldots,x_{r}]]$ $[1/x_{r}]$. Then, for $n$ prime to the
characteristic of $k$, we have
$$
H^{q}(U,𝛍_{n})=
\begin{cases}
𝛍_{n}(A) &\text{\rm if}\quad q=0\\
\ZZ/n &\text{\rm if}\quad q=1\\
0 &\text{\rm if}\quad q>1{\quoi}.
\end{cases}
$$
\end{enonce*}

\begin{enonce*}[remark]{Remark}
The proof applies equally to the ring of convergent power series when
$k=\CC$. Moreover, one easily avoids the hypothesis that $k$ be
separably closed by stating the result with a little more care. But
to establish the general purity theorem \hyperref[XIX.thm3.2]{3.2} (under the
resolution-of-singularities hypothesis, of course), we need the result
only when $k$ is separably closed, and these other assertions are
corollaries of that theorem \hyperref[XIX.thm3.2]{3.2}.
\end{enonce*}

In dimensions 0 and 1, the result can be proved for every regular
strictly local ring:

\begin{enonce*}{Lemma 1.3}\etiquette[1.3]{XIX.lem1.3}
Let $A$ be a regular strictly local ring and let $x\in \rad A$ be a
local parameter. Let $U=\Spec A[1/x]$. Then
$H^{0}(U,𝛍_{n})=𝛍_{n}(A)$, and
there are canonical isomorphisms
$$
\ZZ/n\simeq H^{0}(U,\GG_{m})/(H^{0}(U,\GG_{m}))^{n}\simeq
H^{1}(U,𝛍_{n}){\quoi},
$$
where\pageoriginale the generator of $\ZZ/n$ is the “residue” of $x\in H^{0}$ $(U,\GG_{m})$.
\end{enonce*}

\medskip
\noindent
\emph{Proof}. For $q=0$, the assertion is that $U$ is connected
and nonempty. This is true. For $q=1$, apply the Kummer exact sequence
(\hyperrefext[IX][IX.rems3.7]{3.7}). We have $\Pic U=0$ because $A$ is regular and
factorial, and $U$ is an open subset of $X=\Spec A$. It follows that
we have the exact sequence
$$
H^{0}(U,\GG_{m})\xrightarrow{n}H^{0}(U,\GG_{m})\rightarrow
H^{1}(U,𝛍_{n})\rightarrow 0{\quoi},
$$
and it remains to prove that
$H^{0}(U,\GG_{m})/(H^{0}(U,\GG_{m}))^{n}$ is isomorphic to
$\ZZ/n$ and generated by the residue of $x$. But there is an obvious
exact sequence
$$
\xymatrix@R=.1cm{
0 \ar[r]& H^0(X,\GG_m)\ar[r] & H^{0}(U,\GG_m)\ar[r] &
\ZZ \ar[r]& 0\\
& \text{\rotatebox{-90}{$\simeq$}} & \text{\rotatebox{-90}{$\simeq$}} &\\
& A^{\ast}  & A[1/x]^{\ast}   
}
$$
where the image of a section $a\in H^{0}$ $(U,\GG_{m})$ in
$\ZZ$ is the order of vanishing of $a$ along $\{x=0\}$. The
group $A^{\ast}$ is divisible by $n$, because $A$ is strictly local
and $n$ is prime to the residual characteristic. Indeed, if
$u\in H^{0}$ $(X,\GG_{m})$, the equation $Y^{n}-u=0$ defines a finite
étale extension of $A$, which is therefore completely split. It
follows that $H^{1}$ $(U,𝛍_{n})\simeq \ZZ/n$, proving the lemma.

\medskip
\noindent
\emph{Proof of} \hyperref[XIX.thm1.2]{1.2}. It remains to treat the case
$q>1$. Recall the terminology of (V, Appendix). We have
$$
H^{q}(U,𝛍_{n})=\varinjlim_{\UU}
H^{q}(\UU,𝛍_{n})
$$ 
where\pageoriginale $\xymatrix{\UU = \bigg\{U_0 & U_1
\ar@<3pt>[l] \ar@<-3pt>[l] & \ldots \ar@<5.5pt>[l]
\ar[l]\ar@<-5.5pt>[l]\bigg\}}$ ranges over the category (filtered up
to homotopy) of \emph{hypercoverings} of $U$. Let $\UU$ be such a
hypercovering, and take the $U_{i}$ to be separated and of finite type.
Since we wish to prove that $H^{q}$ $(U,𝛍_{n})=0$, it suffices to find
a hypercovering $\VV$ and a morphism $\VV\rightarrow \UU$ such that
the induced morphism
$$
H^{q}(\UU,𝛍_{n})\longrightarrow
H^{q}(\VV,𝛍_{n})
$$
is the zero morphism.

Let $\{W_{1},\ldots,W_{m}\}$ be the set of connected components of all
the $U_{i}$ for $i\leq q+1$, and let $B_{j}$ be the normalization of
$A=k[[x_{1},\ldots,x_{r}]]$ in the field $R(W_{j})$ of rational
functions on $W_{j}$. Then, since $W_{j}\rightarrow U$ is étale, the
scheme $W_{j}$ is an open subset of $\Spec B_{j}$, say the open
complement of $\Spec C_{j}$, where $C_{j}$ is the appropriate reduced
quotient ring of $B_{j}$.

Choose in addition an element $b_{j}$ of $B_{j}$ such that $b_{j}$ generates the
separable extension $R(W_{j})$ of
$k((x_{1},\ldots,x_{r}))$. Let $f_{j}(Y)\in A[Y]$ be the monic
irreducible polynomial of $b_{j}$ over $A$, so that $B_{j}$ is
isomorphic to the normalization of the ring $A[Y]/(f_{j})$. Finally,
let $d_{j}\in A$ be the discriminant of $f_{j}(Y)$, which is a nonzero
element of $A$.

Now we may replace the $x_{i}$ for $i<r$ by other elements of $A$, the
only condition being that $\{x_{1},\ldots,x_{r}\}$ form a system of
parameters of $A$.

It is easy to see that, after changing these elements if necessary,
we\pageoriginale may suppose that the following hypotheses are
satisfied:

\setcounter{subsection}{3}
\subsection{}\etiquette[1.4]{XIX.subsec1.4}
\begin{itemize}
\item [(i)] For every $j$ such that $d_{j}$ is not a unit,
$x_{1},\ldots,x_{r-1}$ and $d_{j}$ generate an ideal primary for
$\rad A$.

\item[(ii)] For every $j$, the images of $x_{1},\ldots,x_{r-1}$ in
$C_{j}$ generate an ideal primary for $\rad C_{j}$.
\end{itemize}
 We then have the following lemma:

\begin{enonce*}{Lemma 1.5}\etiquette[1.5]{XIX.lem1.5}
Let $A^{0}=k[[x_{1},\ldots,x_{r-1}]]$ $\{x_{r}\}$ be the henselization
of $k[[x_{1},\ldots,x_{r-1}]]$ $[x_{r}]$ at the point
$(x_{1},\ldots,x_{r})$, and let $X^{0}=\Spec A^{0}$ and $U^{0}=\Spec
A^{0}$ $[1/x_{r}]$. Under the conditions of $\hyperref[XIX.subsec1.4]{1.4}$ above, there exists
a hypercovering $\UU^{0}$ of $U^{0}$ (of order $q+1$) and an
isomorphism $\fprod{\UU^{0}}{X}{X^{0}}\simeq\UU$.
\end{enonce*}

Assume the lemma. The ring $A^{0}$ is excellent (\EGA IV 7.8.3
(ii), 7.9.5), and consequently the connected components of the
$U^{0}_{i}$ and of the $U_{i}$ correspond bijectively (\EGA IV
7.8.3 (vii)). The isomorphism $\fprod{\UU^{0}}{X}{X^{0}}\simeq
\UU$ therefore gives an isomorphism
$$
H^{q}(\UU^{0},𝛍_{n})\simeq
H^{q}(\UU,𝛍_{n}){\quoi}.
$$
Since $k$ is separably closed, $A^{0}$ is a strictly local ring, and
we may apply \hyperref[XIX.cor1.1]{1.1} to the ring $A^{0}$. It follows that
$H^{q}$ $(U^{0},𝛍_{n})=0$. But
$$
H^{q}(U^{0},𝛍_{n})\simeq
\varinjlim_{\VV^{0}}H^{q}(\VV^{0},𝛍_{n}){\quoi},
$$
where $\VV^{0}$ ranges over the category of hypercoverings of
$U^{0}$ of order $q+1$,\pageoriginale and $H^{q}$
$(\UU^{0},𝛍_{n})$ is a finite group. Consequently, there exists a
morphism of hypercoverings $\VV^{0}\rightarrow \UU^{0}$ such that
the induced morphism
$$
H^{q}(\UU^{0},𝛍_{n})\longrightarrow
H^{q}(\VV^{0},𝛍_{n})
$$
is the zero morphism. Put
$\VV=\fprod{\VV^{0}}{X}{X^{0}}$. Then $H^{q}$
$(\VV^{0},𝛍_{n})\simeq
H^{q}(\VV,𝛍_{n})$ (by the same argument as
for $\UU^{0}$), so
$$
H^{q}(\UU,𝛍_{n})\longrightarrow
H^{q}(\VV,𝛍_{n})
$$
is also the zero morphism, \cqfd

\medskip
\noindent
\emph{Proof of} \hyperref[XIX.lem1.5]{1.5}. It suffices to find finite
$A^{0}$-algebras $B^{0}_{j}$ and $C^{0}_{j}$ and $A$-isomorphisms
\begin{align*}
B^{0}_{j}\otimes_{A^{0}}A &\xrightarrow{\sim} B_{j}\\
C^{0}_{j}\otimes_{A^{0}}A &\xrightarrow{\sim} C_{j}{\quoi}.
\end{align*}
Indeed, by (\cite{sga4.exp.XIX.bib.3}~\hyperref[XIX.subsec1.4]{1.4}) the morphism
$B_{j}\rightarrow C_{j}$ is induced by a morphism
$B^{0}_{j}\rightarrow C^{0}_{j}$, so the open subset $W_{j}\subset
\Spec B_{j}$ is induced by an open subset $W^{0}_{j}\subset \Spec
B^{0}_{j}$. Moreover, the $U$-morphisms among the $W_{j}$ are induced
by morphisms among the $A$-algebras $B_{j}$, hence
(\cite{sga4.exp.XIX.bib.3}~\hyperref[XIX.subsec1.4]{1.4}) by morphisms among the $B^{0}_{j}$, and therefore by
$U^{0}$-morphisms among the $W^{0}_{j}$. It follows that the
$W^{0}_{j}$ form the set of connected components of a simplicial
object $\UU^{0}$ over $U^{0}$ that induces $\UU$, and one sees
immediately that it is a hypercovering of $U^{0}$.

Now \hyperref[XIX.subsec1.4]{1.4} (ii) does imply that $C_{j}$ is induced by an
$A^{0}$-algebra\pageoriginale $C^{0}_{j}$, as is easily verified
\footnote[1]{Indeed, \hyperref[XIX.subsec1.4]{1.4} (ii) implies that $C_{j}$ is finite over
$k[[X_{1},\ldots,X_{r-1}]]$, and a fortiori over $A^{0}$, and it
suffices to take $C^{0}_{j}=C_{j}$ with the $A^{0}$-module structure
obtained by restriction of scalars.}. For $B_{j}$, it suffices to find
a monic polynomial $f^{0}_{j}(Y)\in A^{0}[Y]$ such that the two
$A$-algebras
$$
A[Y]/(f_{j})\quav
A[Y]/(f^{0}_{j})=A\otimes_{A^{0}}A^{0}[Y]/(f^{0}_{j})
$$
are isomorphic. Indeed, $A^{0}[Y]/(f^{0}_{j})$ is then reduced, and
we take for $B^{0}_{j}$ the normalization (\EGA IV 7.8.3 (vi)) of
$A^{0}[Y]/(f^{0}_{j})$. The ring $B^{0}_{j}\otimes_{A^{0}}A$ is still
normal (\EGA IV 7.8.3 (vii)) and birational to $B_{j}$, hence equal to
that ring.

By \hyperref[XIX.subsec1.4]{1.4} (i), the ideal $I_{j}\subset A$ generated by $d_{j}$ is
induced by an ideal of $k[[x_{1},\ldots,x_{r-1}]]$ $[x_{r}]$
(Weierstrass preparation theorem \cite{sga4.exp.XIX.bib.5}\cite{sga4.exp.XIX.bib.6}), hence by an ideal
$I^{0}_{j}$ of $A^{0}$. Moreover, $A/I_{j}$ is a finite
$k[[x_{1},\ldots,x_{r-1}]]$-algebra, and is therefore isomorphic to
$A^{0}/I^{0}_{j}$. The latter ring is thus complete, and consequently
the $I^{0}_{j}$-adic completion of $A^{0}$ is complete and isomorphic
to $A$. We can therefore find a monic polynomial
$f^{0}_{j}\in A^{0}$ $[Y]$ such that in $A$ $[Y]$ we have
$$
f_{j}\equiv f^{0}_{j}(\text{\rm modulo } d_{j}{}^{2}\rad A){\quoi}.
$$
By (\cite{sga4.exp.XIX.bib.3}, \hyperref[XIX.lem1.3]{1.3}), $A[Y]/(f_{j})\simeq A[Y]/(f^{0}_{j})$, which
proves the lemma.

\section{The case of a strictly local ring}\etiquette[2]{XIX.sec2}

The theorem is as follows:

\begin{enonce*}{Theorem 2.1}\etiquette[2.1]{XIX.thm2.1}
\footnote[1]{Depends on resolution of singularities (\cf
Introduction).}.\pageoriginale Let $A$ be an excellent, regular, strictly
local ring of equal characteristic, and let $x$ be a “local
parameter”\footnote[2]{By which we mean here: an element belonging to a
\emph{regular} system of parameters.} of $A$. Let $U=\Spec
A[1/x]$. Then, for $n$ prime to the characteristic of $A$, we have
$$
H^{q}(U,𝛍_{n})=
\begin{cases}
𝛍_{n}(A) &\text{\rm if}\quad q=0\\
\ZZ/n &\text{\rm if}\quad q=1\\
0 &\text{\rm if}\quad q>1{\quoi}.
\end{cases}
$$
\end{enonce*}

The proof proceeds by induction on $q$. We denote by $P(N)$ the
assertion of \hyperref[XIX.thm2.1]{2.1} for the values $q\leq N$. 
By \hyperref[XIX.lem1.3]{1.3}, the assertion $P(1)$ is true.

\begin{enonce*}{Lemma 2.2}\etiquette[2.2]{XIX.lem2.2}
Suppose that $P(N)$ is true. Let $A\rightarrow A'$ be a local
morphism of excellent, regular, strictly local rings of equal
characteristic. Let $\{x_{1},\ldots,x_{r}\}$ be part of a regular
system of parameters of $A$ whose set of images in $A'$ is again part
of a regular system of parameters. Let
\begin{align*}
U &= \Spec A[1/x_{1},\ldots,1/x_{r}]\\
U'&= \Spec A'[1/x'_{1},\ldots,1/x'_{r}]{\quoi}.
\end{align*}
Then the canonical morphism
$$
H^{q}(U,𝛍_{n})\longrightarrow
H^{q}(U',𝛍_{n})
$$
is\pageoriginale bijective for every $q\leq N$.
\end{enonce*}

\medskip
\noindent
\emph{Proof}. \footnote[3]{A shorter proof than the one that follows,
but using the cup-product formalism, would consist in showing that,
under the hypothesis $P(N)$, the cohomology
$H^{\ast}(U,\ZZ/n\ZZ)$ is, in dimensions $\leq N$, the exterior
algebra (over $\ZZ/n\ZZ)$ on
$H^{1}(U,\ZZ/n\ZZ)\simeq (\ZZ/n\ZZ)^{r}$,
which reduces matters to establishing \hyperref[XIX.lem2.2]{2.2} for $q=1$, a well-known
case. The calculation indicated under the hypothesis $P(N)$ is made
explicit, for example, in \SGA 7.}
We argue by induction on $r$: the assertion is trivial for $r=0$.
Let $r>0$ and suppose the theorem has been proved for $r-1$. Let
$B=A/(x_{r})$, $V=\Spec A[1/x_{1},\ldots,1/x_{r-1}]$, and
$Y=\Spec B[1/x_{1},\ldots,1/x_{r-1}]$ (where the image of $x_{i}$ in
$B$ is denoted by the same symbol). There is an open immersion
$i:U\rightarrow V$, and $Y$ is the complementary closed subset,
defined by the element $x_{r}$.

Applying $P(N)$ to the strict localizations of $V$ at the geometric
points of $Y$, we find that
\begin{equation*}
R^{q}i_{\ast}(𝛍_{n})_{U}=
\begin{cases}
(𝛍_{n})_{V} &\text{\rm if}\quad q=0\\
(\ZZ/n)_{Y} &\text{\rm if}\quad q=1\\
0 &\text{\rm if}\quad 1<q\leq N{\quoi},
\end{cases}\tag{2.3}
\end{equation*}\etiquette[2.3]{XIX.eq2.3}
because these strict localizations are excellent regular rings of
equal characteristic, and the element $x_{r}$ is a local parameter at every
point of $Y$.

We use a “prime” to denote the analogous objects deduced from $A'$.
We obtain a morphism of spectral sequences
$$
\xymatrix@C=.4cm{
E^{pq}_2 \ar[d]^{\varphi^{pq}_2} & = \text{\phantom{ab}} H^{p}(V,R^{q}i_{\ast}
𝛍_n)\ar[d]^{\varphi} \ar@{=>}[r] & H^{p+q}(U,𝛍_n) \ar[d]^{\varphi}&\\
E'^{pq}_2 &= \text{\phantom{ab}} H^{p}(V',R^{q}i'_{\ast}𝛍_n) \ar@{=>}[r]& H^{p+q}(U',𝛍_n).
}
$$
Now\pageoriginale it is clear from \hyperref[XIX.lem1.3]{1.3} and \eqhyperref[XIX.eq2.3]{2.3}
that the base change $V\xleftarrow{g}V'$ induces an isomorphism
$$
g^{\ast}R^{q}i_{\ast}𝛍_{n}\xrightarrow{\sim}
R^{q}i'_{\ast}𝛍_{n}
$$
for $q\leq N$ (only for $q=0$, $1$ is there anything at all to
check). Consequently, the induction hypothesis, applied to the
morphism
$$
H^{q}(V,𝛍_{n})\longrightarrow
H^{q}(V',𝛍_{n}){\quoi},
$$
implies that $\varphi^{pq}_{2}$ is an isomorphism for $q=0$ and
$p\leq N$.

Similarly, we can apply the induction hypothesis to the morphism of
rings $B\rightarrow B'$; these are regular, excellent, strictly local rings of
equal characteristic, and the elements $x_{1},\ldots,x_{r-1}$ form part of a
system of parameters. Therefore
$$

% END INCLUDED SGA4_English_current_buildable_source_20260906/source/19/drafts/SGA4_XIX_LINES_44_450_English.texfrag

% BEGIN INCLUDED SGA4_English_current_buildable_source_20260906/source/19/drafts/SGA4_XIX_LINES_451_850_English.texfrag
H^{q}(Y,𝛍_{n})\xrightarrow{\sim} H^{q}(Y',
𝛍_{n})\quad\text{\rm if}\quad q\leq N{\quoi}.
$$
Since $(𝛍_{n})\simeq (\ZZ/n)_{Y}$, it follows that
$\varphi^{pq}_{2}$ is also bijective for $p\leq N$ and for $q=1$.

Since
$R^{q}i_{\ast}𝛍_{n}=R^{q}i'_{\ast}𝛍_{n}=0$
if $1<q\leq N$, we find $E^{pq}_{2}={E'}^{pq}_{2}=0$ if $1<q\leq N$,
and for $1\leq m\leq N$ the morphism of spectral sequences therefore
reduces to a morphism of exact sequences
$$
\xymatrix@C=.4cm{
H^{m-2}(Y,\ZZ/n) \ar[r] \ar[d]^{a}& H^{m}(V,𝛍_n) \ar[r]\ar[d]^{b}&
H^m(U,𝛍_n) \ar[r]^(.45){e} \ar[d]^{c}&
H^{m-1}(Y,\ZZ/n)  \ar[r] \ar[d]^{d}&
H^{m+1}(V,𝛍_n) \ar[d]\\
H^{m-2}(Y',\ZZ/n) \ar[r] & H^{m}(V',𝛍_n) \ar[r]&
H^m(U',𝛍_n) \ar[r]^(.45){e} & H^{m-1}(Y',\ZZ/n)  \ar[r]&
H^{m+1}(V',𝛍_n).
}
$$
Now\pageoriginale we have already seen that $a$, $b$, and $d$ are
bijective, and it follows that $c$ is \emph{injective}. To prove that
$c$ is \emph{surjective}, it suffices to prove that the morphism
$$
H^{m}(U,𝛍_{n})\xrightarrow{e}H^{m-1}(Y,\ZZ/n)
$$
is surjective for $m\leq N$.

% SGA-SRC-001196
Let $k\subset A$ be an arbitrary separably closed field, and let
$A^{0}=k\{x_{1},\ldots, x_{r}\}$ be the henselization of
$k[x_{1},\ldots,x_{r}]$ at the origin. Then the hypotheses of
\hyperref[XIX.thm2.1]{2.1} are satisfied for the morphism
$A^{0}\rightarrow A$ and the elements $x_{i}$. With the evident
notation, we have a commutative square
$$
\xymatrix@R=1.25cm{
H^{m}(U^{0},𝛍_n)  \ar[r]^(.45){e^0} \ar[d]^{c^0}& H^{m-1}(Y^0,
\ZZ/n) \ar[d]^{d^{0}}\\
H^{m}(U,𝛍_n)  \ar[r]^(.45){e} & H^{m-1}(Y,
\ZZ/n)
}
$$
and we already know that the arrow $d^{0}$ is bijective. It therefore
suffices to prove that the morphism $e^{0}$ is surjective. We are thus
reduced (dropping the superscript ${}^{0}$) to verifying surjectivity
in the case $A=k\{x_{1},\ldots,x_{r}\}$.

For this ring, which is a limit of schemes smooth over $k$, we proved
the purity theorem in \hyperrefext[XVI][XVI.thm3.7]{3.7}. Thus\pageoriginale
$R^{q}i_{\ast}𝛍_{n}=0$ for every $q>1$, and consequently there is an
exact sequence
$$
H^{m}(U,𝛍_{n})\xrightarrow{e}H^{m-1}(Y,\ZZ/n)\rightarrow
H^{m+1}(V,𝛍_{n})\xrightarrow{f}H^{m+1}(U,𝛍_{n})
$$
for every $m\geq 1$. The surjectivity of $e$ is therefore equivalent
to the injectivity of $f$.

Let $Z=\Spec k\{x_{1},\ldots,x_{r-1}\}$
$[1/x_{1},\ldots,1/x_{r-1}]$. The morphism $\Spec
k\{x_{1},\ldots,x_{r}\}\rightarrow \Spec k\{x_{1},\ldots,x_{r-1}\}$
given by the inclusion of rings is acyclic. This is the local
acyclicity of a smooth morphism (\hyperrefext[XVI][XVI.thm1.1]{1.1}).
Since the open subset V of $\Spec k\{x_{1},\ldots, x_{r}\}$ is the
inverse image of $Z$, the morphism
$$
H^{q}(Z,𝛍_{n})\longrightarrow
H^{q}(V,𝛍_{n})
$$
is bijective for every $q$ (\hyperrefext[XV][XV.def1.3]{1.3}). It
therefore suffices to prove that the morphism
$$
H^{q}(Z,𝛍_{n})\longrightarrow
H^{q}(U,𝛍_{n})
$$
is injective for every $q$.

Now $k\{x_{1},\ldots,x_{r}\}$ is the inductive limit of rings
$A_{\alpha}$ étale over $k\{x_{1},\ldots,x_{r-1}\}$
$[x_{r}]$, each equipped with a lifting to that ring of the point
$x_{1}=x_{2}=\ldots=x_{r}=0$. Consequently, $U$ is the projective
limit of the open subsets
$U_{\alpha}=\Spec A_{\alpha}[1/x_{1},\ldots,1/x_{r}]$, and we have
$$
H^{q}(U,𝛍_{n})=\varinjlim_{\alpha}H^{q}(U_{\alpha},
𝛍_{n}){\quoi}.
$$
% SGA-SRC-001198
But since $A_{\alpha}[1/x_{r}]$ is smooth over the henselian ring
$k\{x_{1},\ldots,x_{r-1}\}$,\pageoriginale and since the closed fiber
of $A_{\alpha}[1/x_{r}]$ is nonempty, we see immediately that the
$k\{x_{1},\ldots, x_{r-1}\}$-algebra $A_{\alpha}[1/x_{r}]$ admits
sections. Hence $U_{\alpha}$ admits sections over $Z$, and it follows
that
$$
H^{q}(Z,𝛍_{n})\longrightarrow
H^{q}(U_{\alpha},𝛍_{n})
$$
is injective for every $q$, and therefore that
$$
H^{q}(Z,𝛍_{n})\longrightarrow
H^{q}(U,𝛍_{n})
$$
is also injective, which proves the lemma.

\hypertarget{sga4.exp.XIX.lemma.2.4}{}\begin{enonce*}{Lemma 2.4}\etiquette[2.4]{XIX.lem2.4}
Suppose that $P(N)$ is true. Let $A\rightarrow A'$ be a regular
morphism \emph{(\EGA IV 6.8.1)} of excellent, strictly local rings of
equal characteristic (not necessarily regular), and let
$U\subset \Spec A$ be a regular open subset. Let
$U'=\fprod{U}{\Spec A'}{\Spec A}$. Then, for $n$ prime to the residue
characteristic of $A$, the morphism
$$
H^{q}(U,𝛍_{n})\longrightarrow
H^{q}(U',𝛍_{n})
$$
is bijective for $q\leq N$.
\end{enonce*}

\medskip
\noindent
\emph{Proof.} Let $f:X\rightarrow \Spec A$ be a resolution of
singularities of $\Spec A$ such that the open subset $U$ of $\Spec A$
lifts to an open immersion $i:U\rightarrow X$, and such that the closed
subset $Y=X-U$ is pure of codimension 1 and has normal crossings. We
denote by a prime the base change induced by $A\rightarrow A'$. The
morphism $g:X'\rightarrow X$ is regular (\EGA IV 6.8.3), so
$f':X'\rightarrow \Spec A'$ is a resolution of $\Spec A'$,
$i':U'\rightarrow X'$ is an open immersion, and\pageoriginale
$Y'=X'-U'$ is pure of codimension 1 and has normal crossings. All these
schemes are excellent.

Let $p'$ be a point of $X'$ and let $p$ be its image in $X$. Since $Y$
has normal crossings in $X$, there exist elements
$x_{1},\ldots,x_{r}\in O_{X,p}$ such that $\{x_{1},\ldots,x_{r}\}$ is
part of a system of parameters of $O_{X,p}$ and such that $\Spec
O_{X,p}\cap Y=V(x_{1},\ldots,x_{r})$. Moreover, it follows immediately
from the fact that $g:X'\rightarrow X$ is regular that the images
$x'_{i}$ of the $x_{i}$ in $O_{X',p'}$ form part of a system of
parameters of this ring, and that $\Spec O_{X',p'}\cap
Y'=V(x'_{1},\ldots,x'_{r})$.

Let $B'$ be the strict localization of $O_{X',p'}$ at a geometric point
$\bar{p}$ above $p'$, and let $B$ be the strict localization of
$O_{X,p}$ at the corresponding geometric point of $X$ (these are
excellent rings (\EGA IV 7.8.3 (ii), 7.9.5)). Let
$$
U_{\bar{p}}=\Spec B[1/x_{1},\ldots,1/x_{r}]\quav U'_{\bar{p}}=\Spec
B'[1/x'_{1},\ldots,1/x'_{r}]{\quoi}.
$$
Lemma \hyperref[XIX.lem2.2]{2.2} applies, and we deduce that the
morphisms
\begin{equation*}
H^{q}(U_{\bar{p}},𝛍_{n})\longrightarrow
H^{q}(U'_{\bar{p}},𝛍_{n})\tag{$\ast$}
\end{equation*}
are bijective for $q\leq N$.

Let us compute the sheaves $R^{q}i_{\ast}𝛍_{n}$ and
$R^{q}i'_{\ast}𝛍_{n}$ for the open immersions
$i:U\rightarrow X$ and $i':U'\rightarrow X'$: the stalk
$(R^{q}i'_{\ast}𝛍_{n})_{\bar{p}}$ at the geometric point $\bar{p}$ is
precisely $H^{q}(U'_{\bar{p}},𝛍_{n})$ (VIII 5). Likewise,
$(R^{q}i_{\ast}𝛍_{n})_{\bar{p}}=H^{q}(U_{\bar{p}},𝛍_{n})$. Thus
$(\ast)$ implies that the base-change morphisms
\begin{equation*}
g^{\ast}R^{q}i_{\ast}𝛍_{n}\longrightarrow R^{q}i'_{\ast}𝛍_{n}\tag{${\ast}{\ast}$}
\end{equation*}
are\pageoriginale bijective for $q\leq N$.

By the base-change theorem for a proper morphism
(\hyperrefext[XII][XII.cor5.5]{5.5}), we have (since $X$ and $X'$ are
proper over the spectra of strictly local rings, and the closed fibers
of $X$ and $X'$ are equal)
$$
H^{q}(X,F)\xrightarrow{\sim} H^{q}(X',g^{\ast}F)
$$
for every torsion sheaf $F$ and every $q$. Consequently,
$({\ast}{\ast})$ implies that in the morphism of spectral sequences
$$
\xymatrix@C=.5cm{
E^{pq}_2 \ar[d]& =& H^{p}(X,R^{q}i_{\ast}𝛍_n) \ar@{=>}[r]\ar[d]^{\varphi^{pq}_2}&
H^{p+q}(U,𝛍_n) \ar[d]^{\varphi}\\
E'^{pq}_2   &=&   H^{p}(X',R^{q}i'_{\ast}𝛍_n)\ar@{=>}[r]
& H^{p+q}(U',𝛍_n),
}
$$
the arrows $\varphi^{pq}_{2}$ are bijective for every $q\leq N$.
It is easily seen that this implies that $\varphi^{pq}_{\infty}$ is
bijective if $q\leq N$ and $p+q\leq N+1$. Consequently, the morphism
of the abutments
$$
H^{m}(U,𝛍_{n})\longrightarrow
H^{m}(U',𝛍_{n})
$$
is bijective for $m\leq N$, which proves the lemma.

The following lemma will complete the proof of \hyperref[XIX.thm2.1]{2.1}:

\hypertarget{sga4.exp.XIX.lemma.2.5}{}\begin{enonce*}{Lemma 2.5}\etiquette[2.5]{XIX.lem2.5}
$P(N)\Rightarrow P(N+1)$.
\end{enonce*}

\medskip
\noindent
\emph{Proof}. Let $A$ be strictly local, regular, excellent, and of
equal\pageoriginale characteristic, let $x\in \rad A$ be a local
parameter, and let $U=\Spec A[1/x]$. Let $\hat{A}$ be the completion of
$A$, let $\hat{x}$ be the image of $x$, let
$\hat{U}=\Spec \hat{A}[1/x]$, and let $g:\hat{U}\rightarrow U$ be the
canonical morphism. It suffices to prove that
\begin{equation*}
\begin{split}
g_{\ast}(𝛍_{n})_{\hat{U}}\simeq
(𝛍_{n})_{U}\quad\text{\rm and}\\
R^{q}g_{\ast}𝛍_{n}=0\quad\text{\rm if}\quad 1\leq
q\leq N{\quoi}.
\end{split}\tag{2.6}
\end{equation*}\etiquette[2.6]{XIX.eq2.6}
Indeed, the Leray spectral sequence
$$
E^{pq}_{2}=H^{p}(U,R^{q}g_{\ast}𝛍_{n})\Longrightarrow
H^{p+q}(\hat{U},𝛍_{n})
$$
will give isomorphisms
$$
H^{q}(U,𝛍_{n})\xrightarrow{\sim}
H^{q}(\hat{U},𝛍_{n})
$$
for $q\leq N$, and an exact sequence
$$
0\longrightarrow H^{N+1}(U,𝛍_{n})\longrightarrow
H^{N+1}(\hat{U},𝛍_{n})\longrightarrow
H^{0}(U,R^{N+1}g_{\ast}𝛍_{n}){\quoi}.
$$
But we already know that $P(1)$ is true (\hyperref[XIX.lem1.3]{1.3}),
and it is therefore enough to prove that $H^{q}(U,𝛍_{n})=0$ if
$1<q\leq N+1$. We may thus replace $A$ by $\hat{A}$, which is
isomorphic to a ring of formal power series
$\hat{A}\simeq k[[x_{1},\ldots,x_{r}]]$, with $k$ separably closed and
$\hat{x}=x_{r}$. This case was treated in \hyperref[XIX.thm1.2]{1.2}.

Let us therefore verify \eqhyperref[XIX.eq2.6]{2.6}: the fact that
$g_{\ast}(𝛍_{n})_{\hat{U}}\simeq
(𝛍_{n})_{U}$ is simply the assertion of \eqhyperref[XIX.eq2.3]{2.3}
for $q=0$. Recall that $R^{q}g_{\ast}𝛍_{n}$ is the sheaf associated
with the presheaf $\Rr^{q}$, where
$$
\Rr^{q}(V)=H^{q}(\hat{V},𝛍_{n})
$$
for\pageoriginale $V\rightarrow U$ étale. It suffices to take $V$
separated and of finite type; then $V$ is a disjoint union of integral
schemes, and it suffices to consider such $V$. Let $B$ be the
normalization of $A$ in the field of fractions of
$\Gamma(V,O_{V})$. Then $B$ is a finite $A$-algebra (\EGA IV 7.8.3
(vi)), and we see immediately that V is a regular open subset of
$\Spec B$ and that $\hat{V}$ is the corresponding open subset of
$\Spec \hat{B}$, where $\hat{B}=B\otimes_{A}\hat{A}$. It follows that
the hypotheses of \hyperref[XIX.lem2.4]{2.4} are satisfied for the open
subset $V\subset \Spec B$ and for the morphism
$B\rightarrow \hat{B}$. Hence
$$
\Rr^{q}(V)\simeq H^{q}(V, 𝛍_{n})\simeq
H^{q}(\hat{V}, 𝛍_{n})
$$
for $q\leq N$. But the sheaf associated with the presheaf
$V\mapsto H^{q}(V,𝛍_{n})$ is plainly zero, which proves the lemma.

\section{Purity}\etiquette[3]{XIX.sec3}

We shall adopt terminology analogous to that of
\hyperrefext[XVI][XVI.sec3]{3}:

\hypertarget{sga4.exp.XIX.definition.3.1}{}\begin{enonce*}{Definition 3.1}\etiquette[3.1]{XIX.def3.1}
A \emph{regular pair $(Y,X)$ of codimension $c$} is a closed immersion
$i:Y\rightarrow X$ of locally Noetherian regular schemes such that for
every $y\in Y$ one has $\codim_{y}(Y,X)=c$ (then there exists an open
neighborhood $X'$ of $y$ in $X$ and sections
$x_{1},\ldots,x_{c}\in \Gamma(X',O_{X})$ that generate the ideal of
$Y$, and which form part of a regular system of parameters of $X$ at
every point of $Y\cap X'$). We write $U=X-Y$, and
$j:U\rightarrow X$ for the corresponding open immersion.
\end{enonce*}

A\pageoriginale \emph{morphism} $(Y',X')\rightarrow (Y,X)$ of regular
pairs of codimension $c$ is a morphism $X'\rightarrow X$ such that
$Y'=\fprod{Y}{X'}{X}$. The notion of a regular triple, of a morphism of
regular triples, \etcn, is defined in the analogous evident way.

\hypertarget{sga4.exp.XIX.theorem.3.2}{}\begin{enonce*}{Theorem 3.2}\etiquette[3.2]{XIX.thm3.2}
\footnote[1]{Depends on resolution of singularities (\cf
Introduction).}. \emph{(Cohomological purity)} Let $(Y,X)$ be a regular
pair of codimension $c>0$, where $X$ is an excellent prescheme of equal
characteristic. Let $F$ be a locally constant sheaf on $X$ of cyclic
groups of order $n$ prime to the characteristic of $X$. Then, in the
notation of \emph{(V 4, VIII 5)},
\begin{align*}
ℋ^{q}_{Y} (X,F) &= (R^{q}i^{!})F=0\quad\text{\rm if}\quad q\neq
2c{\quoi},\\
(R^{q}j_{\ast})j^{\ast}F &= 0\quad\text{\rm if}\quad q\neq 0,
2c-1,\quad \text{and}\\
ℋ^{2c}_{Y}(X,F) &= i^{\ast}(R^{2c-1}j_{\ast})j^{\ast}F
\end{align*}
\noindent
is a locally constant sheaf of cyclic groups of order $n$ on $Y$.
\end{enonce*}

\medskip
\noindent
\emph{Proof}. (\cf \hyperrefext[XVI][XVI.thm3.7]{3.7}). We see
immediately that
$$
ℋ^{q}_{Y}(X,F)=0\quad\text{\rm if}\quad q=0,1{\quoi}.
$$
By (V \hyperref[V.cor4.5]{4.5}),
$$
ℋ^{q}_{Y}(X,F)=(R^{q}i^{!})F\simeq
i^{\ast}(R^{q-1}j_{\ast})j^{\ast}F\quad\text{\rm if}\quad q>1{\quoi},
$$
so the assertions for the three sheaves are equivalent. Since the
assertions are local on $X$ for the étale topology, we may
assume\pageoriginale that $F=𝛍_{n}$.

For the case $c=1$, the calculation of the stalks (VIII 5) of
$(R^{q}j_{\ast})j^{\ast}𝛍_{n}$ reduces us to the situation considered
in \hyperref[XIX.thm2.1]{2.1}, whence
$$
((R^{q}j_{\ast})j^{\ast}𝛍_{n})_{\bar{y}}=
\begin{cases}
\ZZ/n &\text{\rm if}\quad q=1\\
0 &\text{\rm if}\quad q>1
\end{cases}
$$
for every geometric point $\bar{y}$ of $Y$. Since these sheaves vanish
outside $Y$ in any case, this proves the assertion for $q>1$. For
$q=1$ $(=2c-1)$, we apply the criterion of
(\hyperrefext[IX][IX.prop2.13]{2.13} (i)) and the explicit calculation
of \hyperref[XIX.lem1.3]{1.3}.

The case $c>1$ is treated by induction (\cf
\hyperrefext[XVI][XVI.thm3.7]{3.7}): let $(Y,X)$ be a regular pair of
codimension $c>1$. By definition, it is clear that, locally on $X$, we
can find a regular triple $(Y,Z,X)$ in which $(Z,X)$ has codimension 1
and $(Y,Z)$ has codimension $c-1$. Let
$$
\xymatrix{
  Y \ar[rr]^{u} \ar[rd]_{i}   &   & Z\ar[ld]^{v}\\
      & X & 
}
$$
be the closed immersions. We have the spectral sequence
$$
E^{pq}_{2}=(R^{p}u^{!})(R^{q}v^{!})F\Longrightarrow
(R^{p+q}i^{!})F{\quoi}.
$$
By the induction hypothesis, $(R^{q}v^{!})F=0$ if $q\neq 2$, and is
locally constant \etc\ldots if $q=2$. Thus, again by the induction
hypothesis, applied to the morphism $u$,
$(R^{p}u^{!})(R^{2}v^{!})F=0$ if $(p,q)\neq (2c-2,2)$\pageoriginale
and is locally constant \etc\ldots if $(p,q)=(2c-2,2)$, whence the
result for the abutment $(R^{p+q}i^{!})F$, \cqfd

\setcounter{subsection}{2}
\subsection{}\etiquette[3.3]{XIX.subsec3.3}
It remains to determine the locally constant sheaves of
(\hyperref[XIX.thm3.2]{3.2}) up to \emph{canonical} isomorphism. For
each prescheme $S$ we introduce the locally constant sheaves
$(𝛍_{n}^{\otimes r})_{S}=$ the $r$-fold tensor product of the sheaf
$(𝛍_{n})_{S}$ over the constant sheaf of rings $\ZZ_{S}$. We also set
$(𝛍_{n}^{\otimes 0})_{S}=(\ZZ/n)_{S}$. It is immediately seen that if
$n$ is invertible on $S$, then $(𝛍_{n}^{\otimes r})_{S}$ is a locally
constant sheaf of cyclic groups of order $n$. We have formulas of the
form $(𝛍_{n}^{\otimes r})\otimes
(𝛍_{n}^{\otimes
s})=(𝛍_{n}^{\otimes r+s})$, and if $S'\rightarrow
S$, then $(𝛍_{n}^{\otimes r})_{S'}$ is canonically isomorphic to the
inverse image on $S'$ of $(𝛍_{n}^{\otimes r})_{S}$. We shall often
omit the symbol $S$ when no confusion can arise.

Let $(Y,X)$ be a regular pair of codimension 1, and suppose that $Y$ is
defined by an equation $x=0$. We have the exact sequence of sheaves
$$
0\longrightarrow (\GG_{m})_{X}\longrightarrow
j_{\ast}(\GG_{m})_{U}\longrightarrow \ZZ_{Y}\longrightarrow
0{\quoi},
$$
whence, by applying (\hyperref[XIX.lem1.3]{1.3}), canonical
isomorphisms
\begin{equation*}
(R^{1}j_{\ast})𝛍_{n}\simeq
j_{\ast}(\GG_{m})_{U}/(j_{\ast}(\GG_{m})_{U})^{n}\simeq
(\ZZ/n)_{Y}{\quoi},\tag{3.3.1}
\end{equation*}\etiquette[3.3.1]{XIX.eq3.3.1}
which plainly do not depend on the choice of $x$. Consequently, we have
a canonical isomorphism
\begin{equation*}
\varphi:\ZZ/n\xrightarrow{\sim}
(R^{2}i^{!})𝛍_{n}\simeq i^{\ast}(R^{1}j_{\ast})𝛍_{n}{\quoi}.\tag{3.3.2}
\end{equation*}\etiquette[3.3.2]{XIX.eq3.3.2}

\begin{enonce*}{Theorem 3.4}\etiquette[3.4]{XIX.thm3.4}

% END INCLUDED SGA4_English_current_buildable_source_20260906/source/19/drafts/SGA4_XIX_LINES_451_850_English.texfrag

% BEGIN INCLUDED SGA4_English_current_buildable_source_20260906/source/19/drafts/SGA4_XIX_LINES_851_1250_English.texfrag
\footnote[1]{Depends on resolution of singularities (\cf
Introduction). This theorem may also be regarded as the conjunction of
\hyperref[XIX.thm3.2]{3.2} and the theory of variance of \Exp XVIII by means of
``local fundamental classes,'' where homomorphisms
$\varphi(Y,X)$ (not necessarily bijective) are studied for $Y$ locally a
complete intersection in $X$.}.\pageoriginale Let $C$ be the family
of regular pairs of excellent schemes of equal characteristic. There exists a
unique function $\varphi$ assigning to every regular pair $(Y,X)$ of
codimension $c$ an isomorphism (called canonical) of sheaves
$$
\varphi(Y,X)=\varphi:\ZZ/n\xrightarrow{\sim}
ℋ^{2c}_{Y}(X,𝛍_{n}^{\otimes
c})=(R^{2c}i^{!})𝛍_{n}^{\otimes c}
$$
such that $\varphi$ satisfies the following conditions:

\begin{enumerate}
\renewcommand{\theenumi}{\roman{enumi}}
\renewcommand{\labelenumi}{(\theenumi)}
\item If $c=1$, and if $Y$ is defined in $X$ by an equation $x=0$,
% SGA-SRC-001210
then $\varphi$ is the morphism \eqhyperref[XIX.eq3.3.2]{3.3.2}.

\item $\varphi$ is compatible with smooth morphisms of pairs; that is, if
$(Y',X')\rightarrow (Y,X)$ is a morphism, let $g:Y'\rightarrow Y$ be the
induced morphism. The diagram of $Y'$-sheaves
$$
\xymatrix{
\ZZ/n  \ar[rr]^{\varphi(Y',X')} & &
(R^{2c}i^{!})𝛍_n^{{\boxtimes}c}\\
g^{\ast}(\ZZ/n) \ar[u]\ar[rr]^{g^{\ast}\varphi(Y,X)}& &
g^{\ast}(R^{2c}i^{!})𝛍_n^{{\boxtimes}c} \ar[u]
}
$$
is commutative (which implies in particular that the right-hand vertical
arrow is an isomorphism).

\item \emph{Transitivity:} let $(Z,Y,X)$ be a regular triple, that is, a
commutative diagram of closed immersions of regular preschemes \etcn,
$$
\xymatrix{
Z \ar[rr]^{v}\ar[rd]_{w}&    &Y\ar[ld]^{u}\\
            & X  & .
}
$$
Let\pageoriginale $a$, $b$, $c$ be the codimensions of $u$, $v$, $w$,
respectively \emph{(thus $c=a+b$)}. We have an isomorphism
\begin{align*}
\varphi(Y,X)\otimes 𝛍_{n} {}^{\otimes
b}:𝛍_{n} {}^{\otimes b} \longrightarrow &
 (R^{2a}u^{!})(𝛍_{n}^{\otimes a}) \otimes
𝛍_{n} {}^{\otimes b}\simeq\\
& \simeq (R^{2a}u^{!})𝛍_{n}{}^{\otimes c}{\quoi}.
\end{align*}
The spectral sequence $(R^{p}v^{!})(R^{q}u^{!})F\Rightarrow
(R^{p+q}w^{!})F$ gives an isomorphism
$$
\xi : (R^{2b}v^{!})(R^{2a}u^{!})𝛍_{n}{}^{\otimes
c}\xrightarrow{\sim} (R^{2c}w^{!})𝛍_{n}{}^{\otimes
c}{\quoi}.
$$
We thus obtain a diagram of isomorphisms
% SGA-SRC-001212 — tensor typography checked against the scan; missing \varphi reconstructed from transitivity
$$\vcenter{\etiquette[3.4.1]{XIX.eq3.4.1}
\xymatrix{
\ZZ/n \ar[r]^{\varphi(Z,Y)} \ar[d]^{\varphi(Z,X)}&
(R^{2b}v^{!})𝛍_{n}^{\otimes b} \ar[d]^{\eta =
  (R^{2b}v^{!})(\varphi(Y,X)\otimes 𝛍_{n}^{\otimes b})}\\
(R^{2c}w^{!})𝛍_{n}^{\otimes c} &
(R^{2b}v^{!})(R^{2a}u^{!})𝛍_{n}^{\otimes c}\ar[l]
}}\leqno{(3.4.1)}
$$
and the transitivity assertion is that this diagram is commutative.
\end{enumerate}
\end{enonce*}

\medskip
\noindent
\emph{Proof.} The proof of existence and uniqueness proceeds by induction on
the codimension $c$. For $c=1$, we define the morphism $\varphi$ locally by
\eqhyperref[XIX.eq3.3.1]{3.3.1}. Since this morphism does not depend on the
choice of the parameter $x$, gluing gives an isomorphism $\varphi$ for every
$(Y,X)$ of codimension 1; and it is clear from the definition of the morphism
and from \hyperref[XIX.lem1.3]{1.3} that the $\varphi$ so constructed satisfies
(ii). Transitivity does not enter when $c=1$.

Now\pageoriginale suppose that existence and uniqueness have already been
proved when the codimension is $<c$, and that $c>1$. Applying the induction
hypothesis, we immediately find that in the situation of
\hyperref[XIX.thm3.4]{3.4} (iii), the three arrows of
\eqhyperref[XIX.eq3.4.1]{3.4.1} that are already defined, \ie
$\varphi (Z,Y)$, $\xi$, $\eta$, are compatible with morphisms of regular
triples. In this situation set
$$
\Psi (Z,Y,X)=\xi \eta \varphi (Z,Y){\quoi},
$$
so $\Psi$ too is compatible with morphisms of triples. Now, if $(Z,X)$ is a
regular pair of codimension $c>1$, locally on $X$ one can find a triple
$(Z,Y,X)$ with $0<\codim(Y,X)<c$. It therefore remains only to prove that the
arrow $\Psi(Z,Y,X)$ is independent of $Y$. Indeed, this will prove the
uniqueness of $\varphi$ in codimension $c$, and existence will then follow by
gluing.

If $(Z,Y,X)$ is arbitrary with $\codim(Y,X)>1$, locally on $X$ one can find a
regular triple $(Y,W,X)$ in which $(W,X)$ has codimension 1. A transitivity
diagram chase, which we leave to the reader, therefore reduces us to checking
independence of $Y$ when $(Y,X)$ has codimension 1.

Let $(Z,Y_{0},X)$ and $(Z,Y_{1},X)$ be two such triples. It suffices to check
the assertion fiber by fiber, and we are thus reduced to the case in which
$X$ is strictly local and $Y_{i}$ is defined by an equation, say $f_{i}=0$
$(i=0,1)$. Let $\bar{X}=\sheaf{\Spec}O_{X}[t]$,
$\bar{Z}=\fprod{Z}{\bar{X}}{X}$, and let $\bar{Y}$ be the closed subscheme of
$\bar{X}$ defined by the equation
$$
(f_{1}-f_{0})t+f_{0}=0{\quoi}.
$$
Then\pageoriginale $\bar{Y}$ is plainly regular at the points lying over the
sections $\epsilon_{0}:\{t=0\}$ and $\epsilon_{1}:\{t=1\}$ of
$\bar{X}/X$. Let $C$ be the closed subset (\EGA IV 7.8.3 (iv)) of $\bar{Y}$
consisting of the points at which $\bar{Y}$ is not regular, and replace
$\bar{X}$, $\bar{Y}$, $\bar{Z}$ by $\bar{X}-C$, $\bar{Y}-C$,
$\bar{Z}-C\cap \bar{Z}$, respectively. Then $(\bar{Z},\bar{Y},\bar{X})$ is a
regular triple, and the sections $\epsilon_{i}$ give morphisms of triples
$$
(Z,Y_{i},X)\longrightarrow (\bar{Z},\bar{Y},\bar{X})\quad
(i=0,1){\quoi}.
$$

Let $F=(R^{2c}w^{!})𝛍_{n}{}^{\otimes c}$ be the sheaf on $Z$.
It is a locally constant sheaf isomorphic to $\ZZ/n$ (by either of the
isomorphisms $\Psi (Z,Y_{i},X)$), and one sees immediately that the sheaf
$\bar{F}=(R^{2c}\bar{w}^{!})𝛍_{n}{}^{\otimes c}$ on $\bar{Z}$ is the inverse
image on $\bar{Z}$ of $F$. Since the fibers of $\bar{Z}\rightarrow Z$ are
connected and nonempty and $\bar{Z}\rightarrow Z$ is smooth, hence open, it
follows that the induced morphism
$$
H^{0}(Z,F)\longrightarrow H^{0}(\bar{Z},\bar{F})
$$
is \emph{bijective}; hence each of the morphisms
$$\vcenter{
\xymatrix@R=.5cm{
                 &                                & H^{0}(Z,F)\\
H^{0}(Z,F) \ar[r]& H^{0}(\bar{Z},\bar{F})\ar[ru]^{\epsilon^{\ast}_0} \ar[rd]_{\epsilon^{\ast}_1} &           \\
                 &                                & H^{0}(Z,F)
}}\leqno{(*)}
$$
is bijective, and consequently $\epsilon^{\ast}_{0}=\epsilon^{\ast}_{1}$.

To give an isomorphism of a sheaf $F$ with $\ZZ/n$ is equivalent to giving a
global section (the image of 1) that is a generator. Let $\bar{\alpha}$,
$\alpha_{i}$ be the global sections of $\bar{F}$, $F$ given by the
isomorphisms\pageoriginale $\Psi(\bar{Z},\bar{Y},\bar{X})$,
$\Psi(Z,Y_{i},X)$, respectively. Since $\Psi(.,.,.)$ is compatible with
morphisms, the diagram $(\ast)$ implies that
$$
\alpha_{0}=\epsilon^{\ast}_{0}(\bar{\alpha})=\epsilon^{\ast}_{1}(\bar{\alpha})=\alpha_{1}{\quoi},
$$
whence $\Psi(Z,Y_{0},X)=\Psi(Z,Y_{1},X)$, \cqfd

\section{Local acyclicity of a regular morphism}\etiquette[4]{XIX.sec4}

\begin{enonce*}{Theorem 4.1}\etiquette[4.1]{XIX.thm4.1}
\footnote[1]{Depends on resolution of
singularities (\cf Introduction).} \emph{(local acyclicity).} Let
$g:X'\rightarrow X$ be a regular morphism \emph{(\EGA IV 6.8.1)} of
excellent schemes of equal characteristic. Then $g$ is universally locally
acyclic for $n$ prime to the characteristic.
\end{enonce*}

\medskip
\noindent
\emph{Proof}. Since the hypothesis is stable under base change of finite
type, it follows at once from \hyperrefext[XV][XV.prop1.13]{1.13} ii) that it
suffices to prove that $g$ is locally acyclic for $n$. Apply the criterion of
\hyperrefext[XV][XV.cor1.17]{1.17}: let $\bar{x}'$ be a geometric point of
$X'$ and $\bar{x}$ the corresponding geometric point of $X$. It follows
immediately from the definition of a regular morphism that the morphism of
strict localizations at these points induced by $g$ is again regular. We are
thus reduced to the case in which $X=\Spec A$ and $X'=\Spec A'$ are strictly
local, $g$ is a local morphism, and by \hyperrefext[XV][XV.cor1.17]{1.17} it
suffices to prove that the\pageoriginale geometric fibers of $g$ are acyclic
for $n$. Replacing $X$ by an integral closed subscheme and $X'$ by the
corresponding closed subscheme reduces us to the case of the geometric generic
fiber $X'_{\bar{x}}$. Take $\bar{x}=\Spec \bar{K}$, where $\bar{K}$ is the
separable closure of the field of rational functions on $X$. Then
$$
\bar{x}=\varprojlim_{\alpha}U_{\alpha}
$$
where $U_{\alpha}$ ranges over a filtered category of connected étale
$X$-schemes, which may be assumed \emph{regular} (\EGA IV 7.8.3 (v)). Let
$B_{\alpha}$ be the normalization of $A$ in the field of rational functions
$R(U_{\alpha})$. Set $U'_{\alpha}=\fprod{X'}{U_{\alpha}}{X}$ and
$B'_{\alpha}=A'\otimes_{A}B_{\alpha}$. Then $U_{\alpha}$ is a regular open
subscheme of $\Spec B_{\alpha}$ and $B_{\alpha}\rightarrow B'_{\alpha}$ is
% SGA-SRC-001217
regular, so by \hyperref[XIX.lem2.4]{Lemma 2.4}
$$
H^{q}(U_{\alpha},𝛍_{n})\xrightarrow{\sim}
H^{q}(U'_{\alpha},𝛍_{n})
$$
for every $q$. Now the fiber $X'_{\bar{x}}$ is the limit of the
$U'_{\alpha}$, hence
\begin{align*}
H^{q}(X'_{\bar{x}},𝛍_{n}) &= \varinjlim_{\alpha}
H^{q}(U'_{\alpha},𝛍_{n})\\
&= \varinjlim_{\alpha} H^{q}(U_{\alpha}, 𝛍_{n})\\
 &= H^{q}(\bar{x},𝛍_{n})=0\quad\mbox{\rm if}\quad q>0,\cqfd.
\end{align*}

\begin{enonce*}{Corollary 4.2}\etiquette[4.2]{XIX.cor4.2}
(base change by a regular morphism). Let
$$
\xymatrix@C=2.25cm@R=1cm{
X \ar[d]^{f}& X'\ar[l]_{g'} \ar[d]^{f'}\\
Y & Y'\ar[l]_{g}
}
$$
be a Cartesian diagram, where $g:Y'\rightarrow Y$ is a regular morphism of
excellent preschemes of equal characteristic, and $f$ is quasi-compact and
quasi-separated.\pageoriginale Then for every abelian torsion sheaf $F$ whose
torsion is prime to the characteristic, the base-change morphisms
\emph{(XII \hyperref[XII.eq4.2]{4.2})}
$$
g^{\ast}(R^{q}f_{\ast})F\longrightarrow (R^{q}f'_{\ast}){g'}^{\ast}F
$$
are bijective for every $q\geq 0$.
\end{enonce*}

This follows from \hyperref[XIX.thm4.1]{4.1} and
\hyperrefext[XVI][XVI.thm1.1]{1.1}.

From \hyperref[XIX.thm4.1]{4.1} and \hyperrefext[XV][XV.cor1.17]{1.17} we
immediately deduce

\begin{enonce*}{Corollary 4.3}\etiquette[4.3]{XIX.cor4.3}
Let $A$ be a strictly local excellent ring of equal characteristic, and let
$f:X\rightarrow \Spec A$ be a morphism. Let $A\rightarrow A'$ be a regular
morphism, where $A'$ too is strictly local, excellent, and of equal
characteristic. Denote by a $'$ the base change induced by $A\rightarrow A'$.
For every torsion sheaf $F$ on $X$, of torsion prime to the characteristic of
$A$, we have
$$
H^{q}(X,F)\xrightarrow{\sim} H^{q}(X',F')
$$
for every $q\geq 0$.
\end{enonce*}

This corollary applies in particular when $A'$ is the completion $\hat{A}$ of
$A$. (In this case, the hypothesis ``$A$ strictly local'' can moreover be
replaced by ``$A$ henselian local,'' as is seen immediately.)

\section{Finiteness theorem}\etiquette[5]{XIX.sec5}\pageoriginale

\begin{enonce*}{Theorem 5.1}\etiquette[5.1]{XIX.thm5.1}
\footnote[1]{\hypertarget{sga4.exp.XIX.note.resolution}{}Depends on resolution of singularities (\cf
Introduction).} \emph{(finiteness).} Let $f:X\rightarrow S$ be a morphism of
finite type between excellent schemes of equal characteristic, and let $F$ be
a constructible abelian torsion sheaf on $X$, of torsion prime to the
characteristic. Then $(R^{q}f_{\ast})F$ is likewise constructible for every
$q$.
\end{enonce*}

This result generalizes the theorem of \hyperrefext[XVI][XVI.thm5.1]{5.1}.
The proof is nearly the same: one applies the purity result
\hyperref[XIX.thm3.2]{3.2} instead of \hyperrefext[XVI][XVI.thm3.7]{3.7}.

\begin{enonce*}{Corollary 5.2}\etiquette[5.2]{XIX.cor5.2}
Let $A$ be an excellent strictly local ring of equal characteristic and let
$f:X\rightarrow \Spec A$ be a morphism of finite type. Let $F$ be a
constructible abelian sheaf on $X$, of torsion prime to the characteristic of
$A$. Then the groups $H^{q}(X,F)$ are finite for every $q$.
\end{enonce*}

\section{Cohomological dimension of affine morphisms}\etiquette[6]{XIX.sec6}

Let $f:X\rightarrow Y$ be a morphism of finite type between excellent
preschemes, and let $F$ be a sheaf on $X$. Recall that we defined
(\hyperrefext[XIV][XIV.def2.2]{2.2}) an integer $\delta(F)$ when $Y$ is local.
We may therefore regard\pageoriginale $\delta$ as a function $\delta(F,f)$
on an arbitrary $Y$, by setting
\begin{equation*}
\delta(F,f,y)=\delta(F'){\quoi},\tag{6.0}
\end{equation*}\etiquette[6.0]{XIX.eq6.0}
where $Y'=\Spec O_{Y,y}$, $X'\rightarrow Y'$ is the morphism obtained from
$f$ by the base change $Y'\rightarrow Y$, and $F'$ is the inverse-image sheaf
of $F$ on $X'$.

We have the following result, which generalizes
\hyperrefext[XVI][XVI.def3.1]{3.1}:

\begin{enonce*}{Theorem 6.1}\etiquette[6.1]{XIX.thm6.1}
\sgarepeatedfootnotemark{sga4.exp.XIX.note.resolution}. Let $f:X\rightarrow Y$ be an affine morphism of finite type,
where $Y$ is an excellent scheme of equal characteristic, and let $F$ be an
abelian torsion sheaf on $X$. Then $R^{q}f_{\ast}F$ vanishes at every point
$y\in Y$ such that $\delta(F,f,y)<q$.
\end{enonce*}

The local variant, plainly equivalent to \hyperref[XIX.thm6.1]{6.1}, is the
following:

\begin{enonce*}{Theorem 6.1 bis}\etiquette[6.1 bis]{XIX.thmbis6.1}
\sgarepeatedfootnotemark{sga4.exp.XIX.note.resolution}. Let $A$ be a strictly local excellent ring of equal
characteristic, and let $f:X\rightarrow Y=\Spec A$ be an affine morphism of
finite type. Let $F$ be a torsion sheaf on $X$. Then
$$
H^{q}(X,F)=0\quad\text{\rm if}\quad q>\delta(F){\quoi}.
$$
\end{enonce*}

Recall (\hyperrefext[X][X.thm5.1]{5.1}) that assertions of this kind are of no
interest for $F$ of $p$-torsion, where $p$ is the characteristic.\pageoriginale
We may therefore suppose that $F$ is of torsion prime to the characteristic.

We have the following corollary:

\begin{enonce*}{Corollary 6.2}\etiquette[6.2]{XIX.cor6.2}
Let $A$ be a henselian, excellent ring of equal characteristic, with residue
field $k$, and let $U\subset \Spec A$ be an affine open subscheme. Let
$\ell \in \PP$. We have
$$
\cd_{\ell}U\leq \cd_{\ell}k+\dim A{\quoi}.
$$
\end{enonce*}

Indeed, let $\widetilde{A}$ be the strict localization of $A$, which is an
infinite Galois étale covering with group $G=G(\bar{k}/k)$, and let
$\widetilde{U}$ be the affine open subscheme of $\Spec \widetilde{A}$ induced
by $U$. Applying the Hochschild--Serre spectral sequence (VIII \hyperref[VIII.prop8.4]{8.4})
$$
H^{p}(G,H^{q}(\widetilde{U},.))\Longrightarrow H^{p+q}(U,.)
$$ 
reduces us to treating the case $A=\widetilde{A}$, and this is then a special
case of \hyperref[XIX.thm6.1]{6.1} bis.

With the notation above, let $R(A)$ be the ring of rational functions on
$A$. We have $\Spec R(A)=\varprojlim_{\alpha}U_{\alpha}$, where $U_{\alpha}$
is an affine open subscheme of $\Spec A$. The theory of passage to the limit
therefore gives

\begin{enonce*}{Corollary 6.3}\etiquette[6.3]{XIX.cor6.3}
Let $A$ be henselian, excellent, and of equal characteristic, with residue
field $k$. Let $R(A)$ be the field of rational functions on $\Spec A$. Then,
for $\ell \in \PP$, we have
$$
\cd_{\ell}R(A)\leq \cd_{\ell}k+\dim A
$$
with\pageoriginale equality if $\ell$ is invertible in $k$, and if moreover
$\ell\neq 2$ or $k$ is not orderable.
\end{enonce*}

Equality in the latter case follows from \hyperrefext[X][X.cor2.4]{2.4}.

Finally, applying \hyperrefext[X][X.cor4.2]{4.2},
\hyperrefext[X][X.cor4.4]{4.4}, we find

\begin{enonce*}{Corollary 6.4}\etiquette[6.4]{XIX.cor6.4}
Let $X$ be a noetherian excellent scheme of equal characteristic, and let
$R(X)$ be the ring of rational functions on $X$. Let $\ell\in \PP$ be
invertible on $X$, and suppose that $\ell\neq 2$ or that no residue field of
$X$ is orderable. Then
$$
\cd_{\ell}(X)\leq \cd_{\ell}R(X)+\dim X\leq 2\cd_{\ell}R(X){\quoi}.
$$
\end{enonce*}

\medskip
\noindent
\emph{Proof of} \hyperref[XIX.thm6.1]{6.1}. We begin with a reduction
analogous to that of \hyperrefext[XIV][XIV.sec3]{3}. Induction on $\delta$:
Clearly, \hyperref[XIX.thm6.1]{6.1} for $\delta\leq d$ is equivalent to
(\hyperref[XIX.thm6.1]{6.1} bis) for $\delta\leq d$. For $\delta=0$ this is
trivial. We treat the case $\delta\leq 1$ in the form
\hyperref[XIX.thm6.1]{6.1} bis. Since $F$ is the filtered direct limit of its constructible
subsheaves (\hyperrefext[IX][IX.prop2.9]{2.9} (iii)), we may suppose that $F$
is constructible; then the support of $F$ is contained in a closed subset
$Z\subset X$ with $\delta(z)\leq 1$ for every $z\in Z$
(\hyperrefext[XIV][XIV.def2.2]{2.2}). Each reduced irreducible component
$Z_{i}$ of $Z$ is then either an affine scheme of finite type and dimension
$\leq 1$ over the closed point $y$ of $Y$, or a scheme quasi-finite over a
reduced closed subscheme $C$ of $Y$ of dimension 1. In the latter case,
$C=\Spec B$, where $B$ is a strictly local integral domain of dimension 1.
Such a scheme $Z_{i}$ is the spectrum of a strictly local integral domain of
dimension $\leq 1$, or of a field finite over the fraction field of $B$.
Thus in all cases the cohomological dimension of $Z_{i}$ is at
most\pageoriginale 1, by \hyperrefext[X][X.thm2.3]{2.3} and
\hyperrefext[IX][IX.cor5.7]{5.7}. We have an exact sequence
$$
0\longrightarrow F\longrightarrow \prod_{i} F_{i}\longrightarrow

% END INCLUDED SGA4_English_current_buildable_source_20260906/source/19/drafts/SGA4_XIX_LINES_851_1250_English.texfrag

% BEGIN INCLUDED SGA4_English_current_buildable_source_20260906/source/19/drafts/SGA4_XIX_LINES_1251_1650_English.texfrag
\epsilon \longrightarrow 0
$$
where $F_{i}$ is the restriction of $F$ to $Z_{i}$, from which the
result for $Z$ follows.

Suppose now that the theorem has been proved for the values
$\delta<d$, and let us prove it for $\delta(F)\leq d$, assuming $d\geq 2$.

\hypertarget{sga4.exp.XIX.lemma.6.5}{}\begin{enonce*}{Lemma 6.5}\etiquette[6.5]{XIX.lem6.5}
It suffices to treat the case where
$X=\EE^{1}_{Y}=\sheaf{\Spec}O_{Y}[t]$ is affine space of
dimension $1$, and where $Y$ is strictly local.
\end{enonce*}

The proof is that of \hyperrefext[XIV][XIV.lem4.2]{4.2}: It is clear that one may
suppose $Y$ strictly local, and that $\delta(F)\leq d$. We must
prove that $H^{q}(X,F)=0$ if $q>d$. We may embed $X$ in
$\EE^{N}_{Y}$, so it suffices to treat the case
$X=\EE^{N}_{Y}$. Induction on $N$: Consider the diagram
$$
\xymatrix@C=1.2cm@R=1cm{
\EE^{N}_Y \ar[rr]^{g} \ar[rd]_{f}&   & \EE^{N-1}_Y \ar[ld]^{h}\\
                     & Y   &  
}
$$
so that $\EE^{N}_{Y}$ is affine space of dimension 1
over $\EE^{N-1}_{Y}$. Applying the induction hypothesis and
the Leray spectral sequence
$$
H^{p}(\EE^{N-1},R^{q}g_{\ast}F)\Longrightarrow
H^{p+q}(\EE^{N},F)
$$
shows that it suffices to prove that
$$
\delta(R^{q}g_{\ast}F)\leq d-q{\quoi}.
$$
This\pageoriginale means that if $z\in \EE^{N-1}$ is a
point such that $\delta(z)>d-q$, then $R^{q}g_{\ast}F=0$ at the
point $z$. But the induction hypothesis on $N$ applies to the morphism
$g$, and hence $R^{q}g_{\ast}F=0$ at the point $z$ if $\delta(F,g,z)<
q$. Thus it suffices to prove the inequality
$$
\delta(F,g,z)+\delta(z)\leq \delta(F){\quoi}.
$$
We leave this pleasure to the reader.

\hypertarget{sga4.exp.XIX.lemma.6.6}{}\begin{enonce*}{Lemma 6.6}\etiquette[6.6]{XIX.lem6.6}
To prove $\hyperref[XIX.thm6.1]{6.1}$ for the values $\delta\leq d$, it suffices to
prove the following: Let $B$ be a complete normal local ring of equal
characteristic and dimension $d$. Let $U\subset \Spec B$ be an
affine open subset. Suppose that there exists an element $x\in \rad B$ which is
invertible on $U$. Then
$$
H^{q}(U,\ZZ/\ell)=0\quad\text{\rm if}\quad q>d{\quoi}.
$$
\end{enonce*}

This lemma is moreover a special case of Corollary \hyperref[XIX.cor6.2]{6.2}.

\medskip
\noindent
\emph{Proof} (\cf \hyperrefext[XIV][XIV.lem4.4]{4.4}). By \hyperref[XIX.lem6.5]{6.5},
it suffices to treat the case where $Y$ is strictly local and
$X=\EE^{1}_{Y}$. Consider the diagram
$$
\xymatrix@C=1cm@R=.8cm{
\EE^{1}_Y \ar[rr]^{i} \ar[rd]_{f}&   & \PP^{1}_Y \ar[ld]^{g}\\
          & Y &,
}
$$
where $\PP^{1}_{Y}$ is the projective line and $\EE^{1}_{Y}$
is obtained from $\PP^{1}$ by removing the section $Y_{\infty}$ at
infinity. Let $F$ be a torsion sheaf on $\EE^{1}$, with
$\delta(F)\leq d$. We have the spectral sequence 
$$
H^{p}(\PP^{1},R^{q}i_{\ast}F)\Longrightarrow
H^{p+q}(\EE^{1},F){\quoi}.
$$
Now\pageoriginale the $R^{q}i_{\ast}F$ for $q>0$ are concentrated on
the section $Y_{\infty}$, which is isomorphic to $Y$. Since $Y$ is
strictly local, it follows that $H^{p}(\PP^{1},
R^{q}i_{\ast}F)=0$ if $p>0$ and $q>0$. Furthermore, by \hyperrefext[XII][XII.cor5.3]{5.3}, we have
$H^{p}(\PP^{1}, i_{\ast}F)=0$ if $p>2$. Since we wish to prove
that $H^{n}(\EE^{1},F)=0$ if $n>d$, and $d\geq 2$, it suffices
to prove that
$$
H^{0}(\PP^{1},
R^{q}i_{\ast}F)=H^{0}(Y_{\infty},R^{q}i_{\ast}F)=0\quad\text{\rm
if}\quad q>d{\quoi}.
$$
This group is isomorphic to the stalk of $R^{q}i_{\ast}F$ at the
point $y_{\infty}$ at infinity of $\PP^{1}$ in the closed fiber.

Now it suffices (VII \hyperref[VII.prop3.3]{3.3} and \hyperrefext[IX][IX.prop2.9]{2.9} (iii)) to treat the case where $F$
is constructible. Then, since $\delta(F)\leq d$, there exists a
closed subscheme $X$ of $\PP^{1}$ with $\dim X\leq d$ such that
$X$ contains the support of $F$. Let $\Spec A=\widetilde{X}$ be the
strict localization of $X$ at the point $y_{\infty}$,
$\widetilde{U}=\widetilde{X}-\fprod{\widetilde{X}}{Y_{\infty}}{\PP^{1}}$,
and let $\widetilde{F}$ be the sheaf induced by $F$ on $\widetilde{U}$. Note that $\widetilde{U}$
is obtained from $\widetilde{X}$ by localizing with respect to an element
$x\in \Gamma(X,O_{X})$. We have
$$
(R^{q}i_{\ast}F)_{y_{\infty}}\simeq
H^{q}(\widetilde{U},\widetilde{F}){\quoi}.
$$
Thus it suffices to prove that $H^{q}(\widetilde{U},\widetilde{F})=0$, if $q>d$,
for every torsion sheaf $\widetilde{F}$ on $\widetilde{U}$.

We may now apply \hyperrefext[IX][IX.prop5.6]{5.6}, with $\varphi=\delta$, and we
find that it suffices to prove that
$$
H^{q}(V,\ZZ/\ell)=0\quad\text{\rm if}\quad q>\delta(V)
$$
for\pageoriginale every $V$ finite and integral over
$\widetilde{U}$ and for invertible $\ell$. Let $x:\bar{V}\rightarrow
V$ be the normalization of $V$. We have the exact sequence
$$
0\longrightarrow (\ZZ/\ell)_{V}\longrightarrow
x_{\ast}(\ZZ/\ell)_{\bar{V}}\longrightarrow C\longrightarrow 0
$$
with $\delta(C)<\delta(V)$, and it follows from this (by induction on
$\delta$) that it suffices to take $V$ \emph{normal}.

Let $B$ be the normalization of $A$ in the field of rational functions
on $V$. Then $B$ is an excellent strictly local ring of equal
characteristic, and $V$ is an affine open subset of $\Spec B$, obtained by
localizing with respect to $x$. Let $\hat{B}$ be the completion of $B$, and
$\hat{V}$ the induced open subset (which is obtained from $\Spec \hat{B}$ by
localizing with respect to the image of $x$ in $\hat{B}$). By \hyperref[XIX.cor4.3]{4.3},
we have
$$
H^{q}(V,\ZZ/\ell)\simeq H^{q}(\hat{V},\ZZ/\ell){\quoi}.
$$
We are therefore reduced to proving that
$$
H^{q}(\hat{V},\ZZ/\ell)=0\quad\text{\rm if}\quad q>\delta
(V)=\dim B=\dim \hat{B}{\quoi}.
$$
By the induction hypothesis, the theorem is true if $\delta(V)<d$,
so Theorem \hyperref[XIX.thm6.1]{6.1} bis applies in this situation. Consequently,
it suffices to treat the case $\dim \hat{B}=d$, and since
$\hat{B}$ is normal, we are in the situation of the lemma, whence the
result.

\setcounter{subsection}{7}
\setcounter{subsubsection}{-1}
\subsubsection{}\etiquette[6.7.0]{XIX.subsubsec6.7.0}
We shall first prove \hyperref[XIX.lem6.6]{6.6} under additional
hypotheses, when 
$B$ has characteristic $p>0$. Let $A$ be an
equidimensional, complete ring of equal characteristic, and let $k\subset A$ be a
field such that $A/\rad A$ is a finite extension of $k$. We say that
the\pageoriginale $k$-algebra $A$ is \emph{formally separable}
if there exists a $k$-subalgebra of $A$ of the form
$k[[x_{1},\ldots,x_{n}]]$ such that the morphism $\Spec A\rightarrow
\Spec k[[x]]$ is finite and étale at the generic point of $\Spec k[[x]]$
$(x=(x_{1},\ldots,x_{n}))$.

\hypertarget{sga4.exp.XIX.lemma.6.7}{}\begin{enonce*}{Lemma 6.7}\etiquette[6.7]{XIX.lem6.7}
Let $k\subset B$ be a field over which $B/\rad B$ is
finite. Under the induction hypothesis on $d$ above, the conclusion
of \hyperref[XIX.lem6.6]{6.6} holds if the hypotheses of $\hyperref[XIX.lem6.6]{6.6}$ are satisfied and if,
in addition, either of the following conditions is satisfied:
\begin{enumerate}
\renewcommand{\theenumi}{\roman{enumi}}
\renewcommand{\labelenumi}{(\theenumi)}
\item $k$ has characteristic $0$.

\item There exists an element $x\in \rad B$ which is invertible on $U$ and
such that $\bar{B}=B/(x)$ is a formally separable $k$-algebra.
\end{enumerate}
\end{enonce*}

\medskip
\noindent
\emph{Proof}. Choose elements $x_{1}=x$,
$x_{2},\ldots, x_{d}$ of $\rad B$ such that $\{x_{1},\ldots,x_{d}\}$
generates a primary ideal for $\rad B$. Then $B$ is a finite $A$
-algebra, where $A=k[[x_{1},\ldots,x_{d}]]$, and $\bar{B}$ is the
$k[[x_{2},\ldots, x_{d}]]=\bar{A}$-algebra $B/x_{1}B$. Since in
case (ii) $\bar{B}$ is formally separable, we may in that case
choose the elements $x_{2},\ldots,x_{d}$ in such a way that
$\bar{B}$ becomes a generically étale $k[[x_{2},\ldots,x_{d}]]$-algebra.
It then follows immediately that $B$ is
also a generically étale $A$-algebra, which is of course also true
in case (i).

Let $b\in B$ be an element which generates the field extension $R(B)$ over
$k((x))$. Let $\epsilon \in k[[x]]$ be the discriminant of the
monic irreducible equation of $b$ over $k[[x]]$. In case (ii),
we may choose $b$ so that $x_{1}$ does not divide
$\epsilon$. Write moreover\pageoriginale $\Spec B-U=V(x_{1})\cup Y$,
where $Y$ contains no irreducible component of $V(x_{1})$. Let
$C$ be the ring $\Gamma(Y,O_{Y})$, where $Y$ is given its induced
reduced structure. Changing the $x_{2},\ldots,x_{d}$ if necessary, we
may suppose that $V(x_{d})$ is not contained in the image of $Y$
in $\Spec A$, and that, moreover, in case (ii), $V(x_{d})$ is not
contained in $V(\epsilon)$.

\hypertarget{sga4.exp.XIX.lemma.6.8}{}\begin{enonce*}{Lemma 6.8}\etiquette[6.8]{XIX.lem6.8}
Let $A_{0}=k[[x_{1},\ldots,x_{d-1}]]$ $\{x_{d}\}$ be the strict localization
of $k[[x_{1},\ldots,x_{d-1}]]$ $[x_{d}]$ at the point
$x_{1}=x_{2}=\ldots=x_{d}=0$. Under the conditions above, there exist
an $A^{0}$-algebra $B^{0}$, an affine open subset $U^{0}\subset \Spec
B^{0}$, and an isomorphism $A\otimes_{A^{0}}B^{0}\simeq B$ such that
the open subset of $\Spec B$ induced by $U^{0}$ is $U$.
\end{enonce*}

The proof of this lemma is analogous to that of \hyperref[XIX.lem1.5]{1.5}:
it suffices to find $A^{0}$-algebras $B^{0}$, $C^{0}$ and
isomorphisms $A\otimes_{A^{0}}B^{0}\simeq B$,
$A\otimes_{A^{0}}C^{0}\simeq C$, because then there will exist
(\cite{sga4.exp.XIX.bib.3}~\hyperref[XIX.subsec1.4]{1.4}) a unique (surjective) morphism
$B^{0}\rightarrow C^{0}$ inducing the morphism $B\rightarrow C$. We
therefore take $U^{0}=\Spec B^{0}-\{(\Spec C^{0})\cup V(x_{1})\}$, and
immediately verify that $U$ is induced by $U^{0}$. The existence of
$C^{0}$, and of $B^{0}$ in case (ii), is proved as in
\hyperref[XIX.lem1.5]{1.5}. The existence of $B^{0}$ in case (i) is a
consequence of (\cite{sga4.exp.XIX.bib.3} 
\hyperref[XIX.thm5.1]{5.1}). (It may also be proved directly by
applying Abhyankar's lemma to the ramification along $\{x_{1}=0\}$,
and descent.)

We can now complete the proof of
\hyperref[XIX.lem6.7]{6.7}. Applying 
\hyperref[XIX.cor4.3]{4.3} once more, we find that
$$
H^{q}(U,\ZZ/\ell)\simeq H^{q}(U^{0},\ZZ/\ell){\quoi},
$$
from which\pageoriginale it follows that it suffices to prove that the latter
group is zero for $q>d$. Now $\Spec A^{0}$ is a limit of schemes
affine and étale over $k[[x_{1},\ldots,x_{d-1}]]$ $[x_{d}]$,
so $\Spec B^{0}$ is a limit of schemes of finite type over
$\Spec k[[x_{1},\ldots,x_{d-1}]]$, each fiber of which has dimension
$\leq 1$. Since $U^{0}\subset \Spec B^{0}[1/x_{1}]$, we have
$$
U^{0}=\varprojlim U^{0}_{\alpha}
$$
where $U^{0}_{\alpha}$ is affine of finite type over the scheme
$V^{0}=\Spec k[[x_{1},\ldots, x_{d-1}]]$ $[1/x_{1}]$ and where each
fiber of $U^{0}_{\alpha}/V^{0}$ has dimension $\leq 1$.

Let $f^{0}:U^{0}\rightarrow V^{0}$ be the structural morphism. Since
the fibers of $f^{0}_{\alpha}$ all have dimension $\leq 1$, for
$v\in V^{0}$ we have
$$
\delta(\ZZ/\ell, f^{0}_{\alpha},v)\leq \dim O_{V^{0},v}+1\leq
d-1
$$
(because $\dim O_{V^{0},v}\leq d-2$). The induction hypothesis on
% SGA-SRC-001224
$d$ is therefore applicable to the morphism $f^{0}_{\alpha}$, and, by
applying \hyperref[XIX.thm6.1]{6.1}, we find
$$
R^{q}f^{0}_{\alpha\ast}(\ZZ/\ell)=0\quad\text{\rm at the
point}\quad v\quad\text{\rm if}\quad q>\dim O_{V^{0},v}+1{\quoi},
$$
\ie if $q>d-\dim\bar{v}$. Thus $R^{q}f^{0}_{\alpha
\ast}(\ZZ/\ell)\neq 0$ at the point $v$ implies that $\dim
\bar{v}\geq d-q$. Therefore, for the inclusion $i:V^{0}\rightarrow \Spec
k[[x_{1},\ldots, x_{d-1}]]$ (\cf \eqhyperref[XIX.eq6.0]{6.0}), we have
$$
\delta (R^{q}f^{0}_{\alpha\ast}(\ZZ/\ell))\leq d-q{\quoi},
$$
and the induction hypothesis, in the form of \hyperref[XIX.thm6.1]{6.1} bis, applied to the
morphism\pageoriginale $i$, implies that
$$
H^{p}(V^{0},R^{q}f^{0}_{\alpha\ast}(\ZZ/\ell))=0\quad\text{\rm
if}\quad p+q>d{\quoi}.
$$
The Leray spectral sequence
$$
E^{pq}_{2}=H^{p}(V^{0},R^{q}f^{0}_{\alpha\ast}(\ZZ/\ell))\Longrightarrow
H^{p+q}(U^{0},\ZZ/\ell)
$$
therefore implies that $H^{n}(U^{0}_{\alpha},\ZZ/\ell)=0$ if $n>d$,
and since $U^{0}$ is the limit of the $U^{0}_{\alpha}$, we likewise find
$$
H^{n}(U^{0},\ZZ/\ell)=0\quad\text{\rm if}\quad n>d{\quoi}.
$$

This completes the proof of \hyperref[XIX.lem6.7]{6.7}, and hence of Theorem \hyperref[XIX.thm6.1]{6.1} in
characteristic 0.

\section{Affine morphisms --- end of the proof}\etiquette[7]{XIX.sec7}

Recall that to complete the proof of \hyperref[XIX.thm6.1]{6.1} in characteristic
$p>0$ (still under the resolution hypothesis), we had reduced by
\hyperref[XIX.lem6.6]{6.6} to proving the following assertion for $d\geq 2$:

\hypertarget{sga4.exp.XIX.lemma.7.1}{}\begin{enonce*}{Lemma 7.1}\etiquette[7.1]{XIX.lem7.1}
Let $B$ be a complete normal local ring of equal characteristic, with
separably closed residue field, and of dimension $d$. Let $U\subset
\Spec B$ be an affine open subset such that there exists $x\in \rad B$ which is
invertible on $U$. Then 
$$
H^{q}(U,\ZZ/\ell)=0\quad\text{\rm if}\quad q>d{\quoi}.
$$

Moreover,\pageoriginale by the induction hypothesis, we may suppose 
that $\hyperref[XIX.thm6.1]{6.1}$ is true for the values $\delta<d$, and that $\hyperref[XIX.lem7.1]{7.1}$ has already been proved
in the particular case $\hyperref[XIX.lem6.7]{6.7}$ {\rm (ii)}.
\end{enonce*}

First note that the condition that $B$ be normal is not
important. Indeed, let $B$ be arbitrary, let $\pi:\bar{B}\rightarrow
B$ be the normalization of $B$, and let $\bar{U}\subset\Spec \bar{B}$ be the inverse
image of $U$. We have an exact sequence
$$
0\longrightarrow (\ZZ/\ell)_{U}\longrightarrow
\pi_{\ast}(\ZZ/\ell)_{\bar{U}}\longrightarrow C\longrightarrow
0{\quoi},
$$
where the dimension of the support of $C$ is $<d$, so $H^{q}(U,C)=0$ if
$q>d-1$. Since $H^{q}(U,\pi_{\ast}\bigdot)\simeq
H^{q}(\bar{U},\bigdot)$ (VIII \hyperref[VIII.cor5.6]{5.6}), we see that
\begin{equation*}
\begin{split}
& H^{q}(U,\ZZ/\ell)=0\quad \text{\rm if}\quad q>d\quad
\text{\rm is equivalent to}\\
& H^{q}(\bar{U},\ZZ/\ell)=0\quad\text{\rm if}\quad q>d{\quoi}.
\end{split}\tag{7.2}
\end{equation*}\etiquette[7.2]{XIX.eq7.2}

\setcounter{subsection}{3}
\setcounter{subsubsection}{-1}
\subsubsection{}\etiquette[7.3.0]{XIX.subsubsec7.3.0}
We shall briefly recall, without proof, the properties of the notion of
formal separability that we need. The reader may consult
(\EGA IV 18.11.10, 18.11.11): Let $A$ be a complete noetherian local
$k$-algebra whose residue field is a finite extension
of $k$, and let $A'$ denote the completion of the local ring
$A\otimes_{k}k'$ (where $k'=k^{1/p}$). Then the extension $A\rightarrow
A'$ is \emph{radicial} (possibly infinite), as is easily seen
--- this is the well-known example of Nagata. It is therefore
harmless, for the calculation of étale cohomology, to replace $A$ by
$A'$ (VIII 1).

Now the condition that an equidimensional algebra $A$ of
dimension\pageoriginale $d$ be formally \emph{separable},
that is, be a finite generically étale extension of
$k[[x_{1},\ldots, x_{d}]]$, can be expressed in terms of the module
$\widehat{\Omega}^{1}_{A/k}$ (\EGA IV 18.11.10), or of the Jacobian
criterion (\EGA $0_{IV}$ 22.6, IV 7.1)\footnote[1]{The co-editor of this
seminar has not succeeded in elucidating the meaning of this
reference to a situation which seems quite different to him, nor the
justification of the assertion that follows. A. G.}, and one sees that if $A$
is not formally separable, then the $k'$-algebra $A'$ is not
reduced. For the proof, one may for example reduce to the
case where $A$ is given by an equation $f=0$, $f\in k[[x_{1},\ldots,
x_{d+1}]]$, and apply the Jacobian criterion and the following lemma,
whose pleasant proof we leave to the reader:

\hypertarget{sga4.exp.XIX.lemma.7.3}{}\begin{enonce*}{Lemma 7.3}\etiquette[7.3]{XIX.lem7.3}
Let $f(x_{1},\ldots,x_{n})\in k[[x_{1},\ldots,x_{n}]]$ be a formal
power series. Suppose that for each $i=1,\ldots,n$ there exists an
invertible element $u_{i}\in k[[x_{1},\ldots,x_{n}]]$ such that $u_{i}f$
is a series in $x_{i}^{p}$ and in the
$x_{1},\ldots,x_{i-1},x_{i+1},\ldots,x_{n}$. Then there exists an
invertible element $v\in k[[x_{1},\ldots,x_{n}]]$ such that
$$
vf\in k[[x_{1}^{p},\ldots,x_{n}^{p}]]{\quoi}.
$$
\end{enonce*}

It follows easily that, by applying a finite number of times the
% SGA-SRC-001226
procedure of replacing $(A,k)$ by $(\bar{A}',k')$, where $\bar{A}'$ is the
normalization of $A'$, one is reduced to the situation where $A$ is indeed
formally separable. Similarly, one can prove that for a given $x\in
\rad A$, the ring $(A/(x))_{\text{\rm réd}}$ becomes a formally
separable algebra by the same procedure, applied to
the algebra $A$.

Applying VIII \hyperref[VIII.thm1.1]{1.1} and \hyperref[VIII.def7.2]{7.2}, we find

\hypertarget{sga4.exp.XIX.corollary.7.4}{}\begin{enonce*}{Corollary 7.4}\etiquette[7.4]{XIX.cor7.4}
It suffices to prove $\hyperref[XIX.lem7.1]{(7.1)}$ under the additional condition that
$B$ contain a field $k$ over which $B/\rad B$ is finite, and
such that the $k$-algebra $(B/(x))_{\text{\rm réd}}$ is formally
separable.\pageoriginale
\end{enonce*}
 
The proof of \hyperref[XIX.lem7.1]{7.1} in this case consists in a
reduction to the particular
case treated in \hyperref[XIX.lem6.7]{6.7} (ii):

With the notation of \hyperref[XIX.lem7.1]{7.1}, \hyperref[XIX.cor7.4]{7.4}, let
$$
D(x)=\sum_{i}r_{i}X_{i}
$$
be the divisor of $x$, so that each $X_{i}$ is a closed irreducible
subscheme of $\Spec B$ of codimension 1 which is formally separable over
$k$ (that is, such that
$B/\mathfrak{J}(X_{i})$ is a formally separable $k$-algebra,
which is moreover equivalent to the assertion that $(B/(x))_{\text{\rm
réd}}$ is formally separable). Choose $y\in\rad B$ such that
we have
$$
D(y)=\sum_{i}(r_{i}-1)X_{i}+\Gamma

% END INCLUDED SGA4_English_current_buildable_source_20260906/source/19/drafts/SGA4_XIX_LINES_1251_1650_English.texfrag

% BEGIN INCLUDED SGA4_English_current_buildable_source_20260906/source/19/drafts/SGA4_XIX_LINES_1651_1920_English.texfrag
$$
where $\Gamma$ contains none of the $X_{i}$, and consider the blow-up
$$
f:Z\longrightarrow \Spec B
$$
of the ideal $(x,y)$. The spectra of the rings $B[y/x]$ and $B[x/y]$
form an affine open covering of $Z$.

Let $B_{1}$ be the strict localization of $B[x/y]$ at the maximal ideal
generated by $(\rad B,x/y)$. For each closed point $Q$ of the closed
fiber of $\Spec B[y/x]\rightarrow \Spec B$, let $A_{Q}$ denote the
strict localization of $B[y/x]$ at $Q$, and let $U_{Q}$ be the affine
open subset of $\Spec A_{Q}$ that is the inverse image of $U$.

\hypertarget{sga4.exp.XIX.lemma.7.5}{}\begin{enonce*}{Lemma 7.5}\etiquette[7.5]{XIX.lem7.5}
To\pageoriginale prove $\hyperref[XIX.lem7.1]{7.1}$ for the ring $B$, it suffices to prove
that for $q>d(\geq 2)$ we have
\begin{equation*}
H^{q}(\Spec B_{1}[y/x],\ZZ/\ell)=0\tag{i}
\end{equation*}
and
\begin{equation*}
H^{q}(U_{Q},\ZZ/\ell)=0\quad\text{for each $Q$}{\quoi}.\tag{ii}
\end{equation*}
\end{enonce*}

\medskip
\noindent
\emph{Proof.} We have open immersions
$$
U\hookrightarrow \Spec B[1/x]\hookrightarrow \Spec B[y/x]{\quoi},
$$
and hence a diagram
$$
\xymatrix{
U \ar[rr]^{i} \ar[rd]_{k}&         & \Spec B[y/x] \ar[ld]^j\\
                        &   Z     &      . 
}
$$
If $Q$ is a closed point of the closed fiber of $\Spec
B[y/x]\rightarrow \Spec B$, we have
$(R^{q}i_{\ast}\ZZ/\ell)_{Q}=0$ if $q>d$, by hypothesis. For
each point $P$ of $\Spec B[y/x]$ which is not closed in the closed
fiber of $\Spec B[y/x]\rightarrow \Spec B$, we have $\delta(F,i,P)<d$,
where $F=(\ZZ/\ell)_{U}$ (notation of \eqhyperref[XIX.eq6.0]{6.0}). Thus, applying
the induction hypothesis and \hyperref[XIX.thm6.1]{6.1} to the morphism $i$, we find that for
the morphism $\Spec B[y/x]\rightarrow \Spec B$ we have
$$
\delta(R^{q}i_{\ast}\ZZ/\ell)\leq d-q{\quoi}.
$$
Consequently, the induction hypothesis in the form \hyperref[XIX.thm6.1]{6.1} bis,
applied to the latter morphism, implies that
$$
H^{p}(\Spec B[y/x],R^{q}i_{\ast}\ZZ/\ell)=0
$$
if\pageoriginale $q>0$ and $p>d-q$. With the Leray spectral sequence
$$
E^{pq}_{2}=H^{p}(\Spec
B[y/x],R^{q}i_{\ast}\ZZ/\ell)\Longrightarrow
H^{p+q}(U,\ZZ / \ell){\quoi}.
$$
we are reduced, in proving \hyperref[XIX.lem7.1]{7.1}, to proving that
$$
H^{p}(\Spec B[y/x],i_{\ast}\ZZ/\ell)=0\quad\text{\rm if}\quad
p>d{\quoi}.
$$

We have an exact sequence
\begin{equation*}
0\longrightarrow \ZZ/\ell\longrightarrow
i_{\ast}(\ZZ/\ell)_{U}\longrightarrow C\longrightarrow 0\tag{7.6}
\end{equation*}\etiquette[7.6]{XIX.eq7.6}
where $\delta(C)<d$. It therefore suffices (again by induction) to prove
that
\begin{equation*}
H^{p}(\Spec B[y/x],\ZZ/\ell)=0\quad\text{\rm if}\quad
p>d{\quoi}.\tag{7.7}
\end{equation*}\etiquette[7.7]{XIX.eq7.7}

Now the open immersion $\Spec B[y/x]\xrightarrow{j}Z$ is obtained by
removing from $Z$ the closed subset $C$ ``at infinity,'' and in the
other open subset $\Spec B[x/y]$ of $Z$ it is identified with the open
subset $C=V(x/y)$. The $R^{q}j_{\ast}\ZZ/\ell$ are concentrated on
$C$, and the morphism $C\rightarrow \Spec B$ is a closed immersion.
Consequently, we have $H^{p}(Z,R^{q}j_{\ast}\ZZ/\ell)=0$
if $p>0$ and $q>0$. Since the closed fiber of $Z\rightarrow \Spec B$ has
dimension 1 and this morphism is proper, \hyperrefext[XII][XII.cor5.3]{5.3} bis implies
that $H^{p}(Z,j_{\ast}\ZZ/\ell)=0$ if $p>2$. The spectral sequence
$$
E^{pq}_{2}=H^{p}(Z,R^{q}j_{\ast}\ZZ/\ell)\Longrightarrow
H^{p+q}(\Spec B[y/x],\ZZ/\ell)
$$
therefore gives isomorphisms
$$
H^{q}(\Spec B[y/x],\ZZ/\ell)\simeq
H^{0}(Z,R^{q}j_{\ast}\ZZ/\ell)\simeq
(R^{q}j_{\ast}\ZZ/\ell)_{P}
$$
if\pageoriginale $q>2$, where $P$ is the point ``at infinity'' of
$Z$. These stalks are none other (VIII 5) than
% SGA-SRC-001231
$$
(R^{q}j_{\ast}\ZZ/\ell)_{P}\simeq H^{q}(\Spec
B_{1}[y/x],\ZZ/\ell){\quoi}.
$$

By \eqhyperref[XIX.eq7.7]{7.7}, it therefore suffices to prove that these latter groups
vanish if $q>d$ $(\geq 2)$. We have an exact sequence analogous to \eqhyperref[XIX.eq7.6]{7.6}
which, together with the induction hypothesis, shows that it suffices that
$$
H^{q}(\Spec B_{1}[y/x],\ZZ/\ell)=0\quad\mbox{\rm if}\quad
q>d{\quoi},
$$
whence the lemma.

\medskip
\noindent
\emph{Proof of} \hyperref[XIX.lem7.5]{7.5} (i).

Consider the closed subset $C=V(x/y)$ of $\Spec B[x/y]$, with its
induced reduced structure. It is immediate that the closed immersion
of $C$ into $\Spec B$ induced by the morphism $Z\rightarrow
\Spec B$ identifies $C$ with the closed subscheme $X=\bigcup X_{i}$ of
$\Spec B$. Since
\begin{align*}
D(x) &= \sum r_{i}X_{i}\\
D(y) &= \sum (r_{i}-1)X_{i}+\Gamma{\quoi},
\end{align*}
the element $x/y$ generates the ideal $\mathfrak{J}(X_{i})$ locally at
the generic point of $X_{i}$, where the local ring is a discrete
valuation ring because $B$ is normal. Consequently, the
morphism $\Spec B[x/y]\rightarrow \Spec B$ is an isomorphism in a
neighborhood of such a point, and the element $x/y$ vanishes to order 1
on each of the irreducible\pageoriginale components $C_{i}$
(corresponding to $X_{i}$) of $C$.

Since the closed subset $C$ lifts to $\Spec B_{1}$, this remains
true if the ring $B[x/y]$ is replaced by $B_{1}$, and hence
$B_{1}/(x/y)$ is a \emph{complete} ring which is a formally separable
$k$-algebra.

To prove \hyperref[XIX.lem7.5]{7.5} (i), we may first apply \hyperref[XIX.cor4.3]{4.3},
which allows us to
replace $B_{1}$ by its completion $\hat{B}_{1}$, without affecting the
ring $B_{1}/(x/y)$. By \eqhyperref[XIX.eq7.2]{7.2}, we may moreover replace
$\hat{B}_{1}$ by its normalization $\bar{B}_{1}$. Since $B_{1}$ (and hence
$\hat{B}_{1}$) is already regular at the generic points of the $C_{i}$, the
morphism $\Spec \bar{B}_{1}\rightarrow \Spec \hat{B}_{1}$ is an
isomorphism at these points, and so $\bar{B}_{1}/(x/y)$ is still a
formally separable $k$-algebra. We are thus reduced to the special case
\hyperref[XIX.lem6.7]{6.7} (ii), which proves the result.

\medskip
\noindent
\emph{Proof of} \hyperref[XIX.lem7.5]{7.5} (ii).

Choose the point $Q$ of the closed fiber of $\Spec
B[y/x]\rightarrow \Spec B$. The strict localization $A_{Q}$ may be written as
a limit
$$
A_{Q}=\varinjlim A
$$
where $A$ ranges over a filtered category of rings that are étale and
of finite type over $B[y/x]$. Let $U_{A}$ be the affine open subset of
$\Spec A$ that is the inverse image of $U$. Since $x$ is invertible on
$U_{A}$, the morphism $U_{A}\rightarrow \Spec B$ is étale, and $U_{A}$
is normal. Consequently, if $C$ denotes the normalization of $B$
in the field $R(A)$ of rational functions on $A$, we have an open
immersion
$$
U_{A}\longrightarrow \Spec C{\quoi}.
$$
It\pageoriginale suffices (VII \hyperref[VII.thm5.7]{5.7}) to prove that for each $A$,
the image of $H^{q}(U_{A},\ZZ/\ell)$ in
$H^{q}(U_{Q},\ZZ/\ell)$ is zero if $q>d$. Since $x$ is
invertible on $U_{Q}$, the morphism $U_{Q}\rightarrow \Spec A$
factors through $\Spec A[1/f]$, if $f\in B$ is any element
such that $f$ is divisible by $x$ in $B[y/x]$ and
$f/x$ does not vanish at the point $Q$. It thus suffices to prove that
for a suitable such $f$, we have
$$
H^{q}(U_{A,f},\ZZ/\ell)=0\quad\text{\rm if}\quad q>d{\quoi},
$$
where $U_{A,f}=U_{A}-V(f)$ is the affine open subset of $\Spec A[1/f]$, which is
also an affine open subset of $\Spec C$.

Let $c\in k$ be an element such that $y/x+c$ does not vanish at $Q$, and
let $J$ be the ideal in $B$ of the closed subset $\Gamma\cap X$. Then if
\begin{equation*}
f\equiv y+cx(\text{\rm mod }J^{N})\quad N>0{\quoi},\tag{7.8}
\end{equation*}\etiquette[7.8]{XIX.eq7.8}
the element $f$ satisfies the condition above: Indeed, it is
immediate that in $\Spec B[y/x]$ we have $V(J)=V(x)$. Thus \eqhyperref[XIX.eq7.8]{7.8}
does imply that $f$ is divisible by $x$, whence
$$
f/x\equiv y/x+c(\text{\rm mod }J){\quoi}.
$$
Since $J$ vanishes at $Q$, and $y/x+c$ is nonzero at $Q$, this
proves our assertion.

Now $\Gamma\cap X$ is a closed subset of $\Spec B$ of codimension
2. Consequently, we can choose an $f$ satisfying \eqhyperref[XIX.eq7.8]{7.8},
call it $z$, such that $V(z)\cap V(y)$ is likewise of
codimension\pageoriginale 2. Also choose elements
$x_{3},\ldots,x_{d}$ such that $y$, $z$, $x_{3},\ldots,x_{d}$ form a
system of parameters of the ring $B$. Then $B$, and hence $C$ as well, is
a finite $k[[y,z,x_{3},\ldots,x_{d}]]$-algebra.

Let $C^{0}$ be the normalization of $k[[y,z,x_{3},\ldots,x_{d}]]$ in the
separable closure of $k$ $((y,z,x))$ in the field of rational
functions $R(C)$, so that $C$ is a radicial extension of
$C^{0}$, and $C^{0}$ is a generically étale extension of $k$
$((y,z,x))$ (where $x=x_{3},\ldots,x_{d}$).

Choose a linear combination
$$
f=az+by\quad (a,b\in k){\quoi},
$$
such that the morphism
$$
\Spec C^{0}\longrightarrow \Spec k[[z,y,x]]
$$
is étale over the generic point of $V(f)$, and such that $f/x$ does not
vanish at $Q$ in $\Spec B[y/x]$. It suffices to prove that
$$
H^{q}(U_{A,f},\ZZ/\ell)=0\quad\text{\rm if}\quad q>d{\quoi},
$$
and hence (VIII \hyperref[VIII.thm1.1]{1.1}) to prove that
$$
H^{q}(U^{0}_{A,f},\ZZ/\ell)=0\quad\text{\rm if}\quad
q>d{\quoi},
$$
where $U^{0}_{A,f}$ is the open subset of $\Spec C^{0}$ corresponding to
the open subset $U_{A,f}$ of $\Spec C$, which is affine, as follows by
descent (\EGA IV 2.7.1 (xiii)). But by construction,
the element $f$ of $C^{0}$ is invertible on $U^{0}_{A,f}$, and
$C^{0}/(f)$ is a formally separable $k$-algebra. We are therefore
in the situation of \hyperref[XIX.lem6.7]{6.7} (ii), which completes
the proof of \hyperref[XIX.lem7.5]{7.5} (ii), and hence that of \hyperref[XIX.thm6.1]{6.1}.

\begin{thebibliography}{99}
\bibitem{sga4.exp.XIX.bib.1} \textsc{Abhyankar}, S.\pageoriginale Local uniformization on
algebraic surfaces over ground fields of characteristic $p\neq
0$. Annals of Math. vol. 63 (1963), and Resolution of singularities of
embedded algebraic surfaces, Academic Press (1966).

\bibitem{sga4.exp.XIX.bib.2} \textsc{Abhyankar}, S. Resolution of singularities of
arithmetical surfaces. Arithmetical Algebraic Geometry, New York
(1966).

\bibitem{sga4.exp.XIX.bib.3} \textsc{Artin}, M. Etale coverings of schemes over hensel
rings, Amer, Journal Math. vol 88 ${\rm n^{\circ}}$ 4 p 915-934
(1966).

\bibitem{sga4.exp.XIX.bib.4} \textsc{Hironaka}, H. Resolution of singularities of an
algebraic variety over a field of characteristic zero I, II. Annals of
Math. vol. 79, ${\rm n^{\circ}}$ 1, 2 (1964).

\bibitem{sga4.exp.XIX.bib.5} \textsc{Weierstrass}, K. Einige auf die analytischen
Functionen mehrerer Veränderlichen sich beziehende
Sätze. Mathematische Werke von K. Weierstrass, Bd 2, Berlin (1895).

\bibitem{sga4.exp.XIX.bib.6} \textsc{Bourbaki}, N. Commutative Algebra, Chap. 3.
\end{thebibliography}


%XII 4.2 no ref. 1072

\ifx\danslelivre\undefined
\let\SGAXIXenddocument\enddocument
\else
\let\SGAXIXenddocument\relax
\fi
\SGAXIXenddocument

% END INCLUDED SGA4_English_current_buildable_source_20260906/source/19/drafts/SGA4_XIX_LINES_1651_1920_English.texfrag


% END INCLUDED SGA4_English_current_buildable_source_20260906/source/19/19_English.tex


\backmatter

\end{document}

% SGA-CUMULATIVE-END 04_SGA4_EN_COMPLETE_READER.tex

% SGA-CUMULATIVE-BEGIN 05_SGA4HALF_EN_COMPLETE_READER.tex
% Cumulative English reader for Exposes I--VIII.
% The ordered component list is recorded in reader-components-1-to-8.txt.
\documentclass[oneside]{book}
\usepackage{sga-style}
\geometry{text={6.3in,8in},centering}
\setlength{\emergencystretch}{1.5em}
\title{Séminaire de Géométrie Algébrique du Bois-Marie\\[20pt]
Étale Cohomology\\[20pt]
(SGA $4\frac 1 2$)}
\author{by P. Deligne\\[10pt]
with the collaboration of J.~F. Boutot,\\
A. Grothendieck, L. Illusie, and J.~L. Verdier}
\date{1977}

\hypersetup{
  pdftitle={SGA 4 1/2: Étale Cohomology — English Translation},
  pdfauthor={Pierre Deligne; J. F. Boutot; Alexander Grothendieck; Luc Illusie; Jean-Louis Verdier},
  pdfsubject={Algebraic geometry; étale cohomology},
  pdfkeywords={SGA 4 1/2, étale cohomology, algebraic geometry, English translation},
  bookmarksopen=true,
  bookmarksnumbered=true
}

\begin{document}

\pagenumbering{alph}
\maketitle

\pagenumbering{roman}
\tableofcontents
\clearpage

\pagenumbering{arabic}
% BEGIN INCLUDED components/sga4half_1_en.tex
% !TEX root = ../SGA4half_English_Reader_1_to_8.tex

\chapter{Étale cohomology: the starting points}\label{I}
\blfootnote{by P.\ Deligne, written up by J.F.\ Boutot}

\section{Grothendieck topologies}\label{I:1}

Grothendieck topologies first appeared as the foundation of Grothendieck's
theory of descent (cf. SGA 1 VI, VIII); their use in the corresponding
cohomology theories came later. We follow the same path here: by formalizing
the classical notions of localization, local property, and patching
(\ref{I:1-1}, \ref{I:1-2}, \ref{I:1-3}), we obtain the general concept of a
Grothendieck topology (\ref{I:1-6}); to justify its introduction into
algebraic geometry, we prove a faithfully flat descent theorem
(\ref{I:1-4}), generalizing the classical theorem 90 of Hilbert
(\ref{I:1-5}).

The reader will find a more complete, but concise, exposition of the
formalism in Giraud \cite{gi64}. M. Artin's notes, “Grothendieck
topologies” \cite{ar62} (Chapters I to III), also remain useful. The
866 pages of Exposés I to V of SGA 4 are valuable when one considers exotic
topologies, such as the one giving rise to crystalline cohomology; in order
to use the étale topology, which is so close to classical intuition, it is
not essential to read them.

\subsection{Sieves}\label{I:1-1}

Let $X$ be a topological space and let $f:X\to \dR$ be a real-valued
function on $X$. The continuity of $f$ is a local property; in other words,
if $f$ is continuous on every sufficiently small open subset of $X$, then
$f$ is continuous on all of $X$. To formalize the notion of a “local
property,” we introduce some definitions.

A set $\cU$ of open subsets of $X$ is called a \emph{sieve} if, for every
$U\in\cU$ and $V\subset U$, one has $V\in\cU$. A sieve is called
\emph{covering} if the union of all the open subsets belonging to it is
equal to $X$.

Given a family $\{U_i\}$ of open subsets of $X$, the sieve generated by
$\{U_i\}$ is, by definition, the set of open subsets $U$ of $X$ such that
$U$ is contained in one of the $U_i$.

A property $P(U)$, defined for every open subset $U$ of $X$, is called
\emph{local} if, for every covering sieve $\cU$ of every open subset $U$
of $X$, $P(U)$ is true if and only if $P(V)$ is true for every
$V\in \cU$. For example, given $f:X\to \dR$, the property “$f$ is
continuous on $U$” is local.

\subsection{Sheaves}\label{I:1-2}

Let us make precise the notion of a function given locally on $X$.

\subsubsection{The sieve point of view}\label{I:1-2-1}

Let $\cU$ be a sieve of open subsets of $X$. A function given
$\cU$-locally on $X$ consists of a function $f_U$ on $U$ for every
$U\in \cU$, such that if $V\subset U$ then $f_V=f_U|V$.

\subsubsection{The Čech point of view}\label{I:1-2-2}

If the sieve $\cU$ is generated by a family of open subsets $U_i$ of $X$,
then giving a function $\cU$-locally amounts to giving a function $f_i$
on each $U_i$ such that
$f_i|{U_i\cap U_j} = f_j|{U_i\cap U_j}$.

In other words, if $Z=\coprod U_i$, then giving a function
$\cU$-locally amounts to giving a function on $Z$ that is constant on the
fibres of the natural projection $Z\to X$.

\subsubsection{}\label{I:1-2-3}

The continuous functions form a sheaf. This means that for every covering
sieve $\cU$ of an open subset $V$ of $X$, and every function
$\{f_U\}$ given $\cU$-locally such that each $f_U$ is continuous on
$U$, there exists a unique continuous function $f$ on $V$ such that
$f|U=f_U$ for every $U\in \cU$.

\subsection{Stacks}\label{I:1-3}

Let us now make precise the notion of a vector bundle given locally on $X$.

\subsubsection{The sieve point of view}\label{I:1-3-1}

Let $\cU$ be a sieve of open subsets of $X$. A vector bundle given
$\cU$-locally on $X$ consists of
\begin{enumerate}[\indent a)]
  \item a vector bundle $E_U$ on each $U\in \cU$,
  \item if $V\subset U$, an isomorphism
    $\rho_{U,V}:E_V\iso E_U|V$, satisfying
  \item if $W\subset V\subset U$, the diagram
    \[\xymatrix{
      E_W \ar[r]^-{\rho_{U,W}} \ar[dr]_-{\rho_{V,W}}
        & E_U|W \\
      & E_V|W \ar[u]_-{\rho_{U,V}|W}
    }\]
    commutes, that is,
    $\rho_{U,W} = (\mbox{$\rho_{U,V}$ restricted to $W$}) \circ \rho_{V,W}$.
\end{enumerate}

\subsubsection{The Čech point of view}\label{I:1-3-2}

If the sieve $\cU$ is generated by a family of open subsets $U_i$ of
$X$, then giving a vector bundle $\cU$-locally amounts to giving:
\begin{enumerate}[\indent a)]
  \item a vector bundle $E_i$ on each $U_i$,
  \item if $U_{ij}=U_i\cap U_j=U_i\times_X U_j$, an isomorphism
    $\rho_{ji}:E_i|U_{ij}\iso E_j|U_{ij}$, such that
  \item if $U_{ijk}=U_i\times_X U_j\times_X U_k$, the diagram
    \[\xymatrix{
      E_i |U_{ijk} \ar[r]^-{\rho_{ki}|U_{ijk}} \ar[dr]_-{\rho_{ji}|U_{ijk}}
        & E_k|U_{ijk} \\
      & E_j|U_{ijk} \ar[u]_-{\rho_{kj}|U_{ijk}}
    }\]
    commutes, that is, $\rho_{ki}=\rho_{kj}\circ\rho_{ji}$ on
    $U_{ijk}$.
\end{enumerate}

In other words, if $Z=\coprod U_i$ and $\pi:Z\to X$ is the natural
projection, then giving a vector bundle $\cU$-locally amounts to giving:
\begin{enumerate}[\indent a)]
  \item a vector bundle $E$ on $Z$,
  \item if $x$ and $y$ are two points of $Z$ such that
    $\pi(x)=\pi(y)$, an isomorphism $\rho_{yx}:E_x\iso E_y$ between the
    fibres of $E$ at $x$ and $y$, depending continuously on $(x,y)$,
    and such that
  \item if $x$, $y$, and $z$ are three points of $Z$ such that
    $\pi(x)=\pi(y)=\pi(z)$, then
    $\rho_{zx}=\rho_{zy}\circ\rho_{yx}$.
\end{enumerate}

\subsubsection{}\label{I:1-3-3}

A vector bundle $E$ on $X$ defines a vector bundle $E_\cU$ given
$\cU$-locally, namely the system of restrictions $E_U$ of $E$ to the
objects of $\cU$. The fact that the notion of a vector bundle is local can
be expressed as follows: for every covering sieve $\cU$ of $X$, the
functor $E\mapsto E_\cU$ from vector bundles on $X$ to vector bundles
given $\cU$-locally is an equivalence of categories.

\subsubsection{}\label{I:1-3-4}

If in \ref{I:1} one replaces “open subset of $X$” by “subset of $X$,”
one obtains the notion of a sieve of subspaces of $X$. In this setting
there are also patching theorems. For example, let $X$ be a space and
let $\cC$ be a sieve of subspaces of $X$ generated by a locally finite
closed covering of $X$. Then the functor $E\mapsto E_\cC$ from vector
bundles on $X$ to vector bundles given $\cC$-locally is an equivalence
of categories.

In algebraic geometry, it is also useful to consider “sieves of spaces
over $X$”; this is what we shall see in the next paragraph.

\subsection{Faithfully flat descent}\label{I:1-4}

\subsubsection{}\label{I:1-4-1}

In the setting of schemes, the Zariski topology is not fine enough for the
study of nonlinear problems, and one is led to replace the open immersions
in the preceding definitions by more general morphisms. From this point of
view, descent techniques appear as localization techniques. Thus the
following descent statement can be expressed by saying that the properties
under consideration are local for the faithfully flat topology (a morphism
of schemes is called faithfully flat if it is flat and surjective).

\begin{proposition}\label{I:1-4-2}
Let $A$ be a ring and $B$ a faithfully flat $A$-algebra. Then:
\begin{enumerate}[(i)]
  \item A sequence $\Sigma=(M'\to M\to M'')$ of $A$-modules is exact if
    the sequence $\Sigma_{(B)}$ deduced from it by extension of scalars to
    $B$ is exact.
  \item An $A$-module $M$ is of finite type (resp. of finite presentation,
    flat, locally free of finite rank, invertible (i.e. locally free of rank
    one)) if the $B$-module $M_{(B)}$ has the corresponding property.
\end{enumerate}
\end{proposition}
\begin{proof}
(i) Since the functor $M\mapsto M_{(B)}$ is exact (by flatness of $B$),
it suffices to show that if an $A$-module $N$ is nonzero, then
$N_{(B)}$ is nonzero. If $N$ is nonzero, $N$ contains a nonzero cyclic
submodule $A/\fa$; hence $N_{(B)}$ contains the cyclic submodule
$(A/\fa)_{(B)}=B/\fa B$, which is nonzero by surjectivity of the
structural morphism $\varphi:\spec(B)\to\spec(A)$ (if $V(\fa)$ is
nonempty, then $\varphi^{-1}(V(\fa))=V(\fa B)$ is nonempty).

(ii) For every family $(x_i)$ of elements of $M_{(B)}$, there exists a
finitely generated submodule $M'$ of $M$ such that $M'_{(B)}$ contains
the $x_i$. If $M_{(B)}$ is of finite type and the $x_i$ generate
$M_{(B)}$, then $M'_{(B)}=M_{(B)}$, and hence $M'=M$ and $M$ is of
finite type.
\end{proof}

If $M_{(B)}$ is of finite presentation, the preceding argument yields a
surjection $A^n\to M$. If $N$ is the kernel of this surjection, then the
$B$-module $N_{(B)}$ is of finite type, so $N$ is of finite type and
$M$ is of finite presentation. The assertion for “flat” follows at once
from (i); “locally free of finite rank” means “flat and of finite
presentation,” and the rank can be tested after extension of scalars to
fields.

\subsubsection{}\label{I:1-4-3}

Let $X$ be a scheme and let $\cS$ be a class of $X$-schemes stable under
fibre products over $X$. A class $\cU\subset\cS$ is a \emph{sieve} on
$X$ (relative to $\cS$) if, for every morphism $\varphi:V\to U$ of
$X$-schemes with $U,V\in\cS$ and $U\in\cU$, one has $V\in\cU$.
The sieve \emph{generated} by a family $\{U_i\}$ of $X$-schemes in
$\cS$ is the class of $V\in\cS$ for which there exists a morphism of
$X$-schemes from $V$ to one of the $U_i$.

\subsubsection{}\label{I:1-4-4}

Let $\cU$ be a sieve on $X$. A quasi-coherent module given
$\cU$-locally on $X$ consists of
\begin{enumerate}[\indent a)]
  \item a quasi-coherent module $E_U$ on every $U\in\cU$,
  \item for every $U\in\cU$ and every morphism $\varphi:V\to U$ of
    $X$-schemes in $\cS$, an isomorphism
    $\rho_\varphi:E_V\iso\varphi^*E_U$, such that
  \item if $\psi:W\to V$ is a morphism of $X$-schemes in $\cS$, the
    diagram
    \[\xymatrix{
      E_W \ar[r]^-{\rho_{\varphi\circ\psi}} \ar[dr]_-{\rho_\psi}
        & \psi^*\varphi^* E_U \\
      & \psi^* E_V \ar[u]_-{\psi^*\rho_\varphi}
    }\]
    commutes, that is,
    $\rho_{\varphi\circ\psi}=(\psi^*\rho_\varphi)\circ\rho_\psi$.
\end{enumerate}

If $E$ is a quasi-coherent module on $X$, we write $E_\cU$ for the
module given $\cU$-locally whose value is $\varphi_U^*E$ on
$\varphi:U\to X$ and such that, for every morphism $\psi:V\to U$, the
restriction isomorphism $\rho_\psi$ is the canonical isomorphism
$E_V=(\varphi_U\circ\psi)^*E\iso\psi^*\varphi_U^*E=\psi^*E_U$.

\begin{theorem}\label{I:1-4-5}
Let $\{U_i\}\in\cS$ be a finite family of $X$-schemes flat over $X$
such that $X$ is the union of the images of the $U_i$, and let $\cU$
be the sieve generated by $\{U_i\}$. Then the functor
$E\mapsto E_\cU$ is an equivalence from the category of quasi-coherent
modules on $X$ to the category of quasi-coherent modules given
$\cU$-locally.
\end{theorem}
\begin{proof}
We treat only the case in which $X$ is affine and $\cU$ is generated by
an affine $X$-scheme $U$ faithfully flat over $X$. The reduction to this
case is formal. Put $X=\spec(A)$ and $U=\spec(B)$.

If $U\to X$ admits a section, then $X$ belongs to the sieve $\cU$ and
the assertion is clear. We reduce to this case.

A quasi-coherent module given $\cU$-locally defines modules $M'$,
$M''$, and $M'''$ on $U$, $U\times_XU$, and
$U\times_XU\times_XU$, and isomorphisms
$\rho:p^*M^\bullet\simeq M^\bullet$ for every projection morphism $p$
between these spaces; this is a \emph{cartesian diagram}
\[\xymatrix{
  M^* : M' \ar@<2pt>[r] \ar@<-2pt>[r]
    & M'' \ar[r] \ar@<4pt>[r] \ar@<-4pt>[r]
    & M'''
}\]
over
\[\xymatrix{
  U_* : U
    & U\times_X U \ar@<2pt>[l] \ar@<-2pt>[l]
    & U\times_X U\times_X U \ar[l] \ar@<4pt>[l] \ar@<-4pt>[l] \mbox{.}
}\]
Conversely, $M^*$ determines the module given $\cU$-locally: for
$V\in\cU$ there exists $\varphi:V\to U$, and we put
$M_V=\varphi^*M'$; for $\varphi_1,\varphi_2:V\to U$ there is a natural
identification
$\varphi_1^*M'\simeq(\varphi_1\times\varphi_2)^*M''\simeq\varphi_2^*M'$,
and by using $M'''$ one sees that these identifications are compatible,
so the definition is legitimate. In short, giving a module
$\cU$-locally is equivalent to giving a cartesian diagram $M^*$ over
$U_*$.

In algebraic terms, giving $M^*$ amounts to giving a cartesian diagram of
modules
\[\xymatrix{
  M' \ar@<3pt>[r]|-{\partial_0} \ar@<-3pt>[r]|-{\partial_1}
    & M'' \ar@<6pt>[r]|-{\partial_0} \ar[r]|-{\partial_1} \ar@<-6pt>[r]|-{\partial_2}
    & M'''
}\]
over the diagram of rings
\[\xymatrix{
  B \ar@<3pt>[r]|-{\partial_0} \ar@<-3pt>[r]|-{\partial_1}
    & B\otimes_A B \ar@<6pt>[r]|-{\partial_0} \ar[r]|-{\partial_1} \ar@<-6pt>[r]|-{\partial_2}
    & B\otimes_A B\otimes_A B
}\]
(more precisely, $\partial_i(bm)=\partial_i(b)\cdot\partial_i(m)$, the
usual identities such as $\partial_0\partial_1=\partial_0\partial_0$
hold, and “cartesian” means that the morphisms
$\partial_i:M'\otimes_{B,\partial_i}(B\otimes_A B)\to M''$ and
$M''\otimes_{B\otimes_A B,\partial_i}(B\otimes_A B\otimes_A B)\to M'''$
are isomorphisms).

The functor $E\mapsto E_\cU$ becomes the functor that associates to an
$A$-module $M$ the diagram
\[\xymatrix{
  M^* = (M\otimes_A B \ar@<2pt>[r] \ar@<-2pt>[r]
    & M\otimes_A B\otimes_A B \ar@<-4pt>[r] \ar[r] \ar@<4pt>[r]
    & M\otimes_A B\otimes_A B\otimes_A B)\text{.}
}\]
Its right adjoint is the functor
\[\xymatrix{
  (M' \ar@<2pt>[r] \ar@<-2pt>[r]
    & M'' \ar@<-4pt>[r] \ar[r] \ar@<4pt>[r]
    & M''') \ar@{|->}[r]
    & \ker(M' \ar@<2pt>[r] \ar@<-2pt>[r]
    & M''')\text{.}
}\]
We must prove that the adjunction maps
\[
  M \to\ker(M\otimes_A B\rightrightarrows M\otimes_A B\otimes_A B)
\]
and
\[
  \ker(M'\rightrightarrows M'')\otimes_A B \to M'
\]
are isomorphisms. By (\ref{I:1-4-2}.i), it suffices to prove this after
a faithfully flat base change $A\to A'$ (with $B$ becoming
$B'=B\otimes_AA'$). Taking $A'=B$ reduces us to the case in which
$U\to X$ admits a section.
\end{proof}

\subsection{A special case: Hilbert's theorem 90}\label{I:1-5}

\subsubsection{}\label{I:1-5-1}

Let $k$ be a field, let $k'$ be a Galois extension of $k$, and let
$G=\gal(k'/k)$. Then the homomorphism
\begin{align*}
  k'\otimes_k k' &\to \oplus_{\sigma\in G} k' \\
  x\otimes y     &\mapsto \{x\cdot \sigma(y)\}_{\sigma\in G}
\end{align*}
is bijective.

It follows that giving a module locally for the sieve on $\spec(k)$
generated by $\spec(k')$ is equivalent to giving a $k'$-vector space
equipped with a semilinear action of $G$, that is:
\begin{enumerate}[\indent a)]
  \item a $k'$-vector space $V'$,
  \item for every $\sigma\in G$, an endomorphism $\varphi_\sigma$ of the
    underlying group of $V'$ such that
    $\varphi_\sigma(\lambda v)=\sigma(\lambda)\varphi_\sigma(v)$ for all
    $\lambda\in k'$ and $v\in V'$, satisfying the condition
  \item for all $\sigma,\tau\in G$,
    $\varphi_{\tau\sigma}=\varphi_\tau\circ\varphi_\sigma$.
\end{enumerate}

Let $V={V'}^G$ be the group of invariants under this action of $G$. It
is a $k$-vector space, and by Theorem \ref{I:1-4-5} one has:

\begin{proposition}\label{I:1-5-2}
The inclusion of $V$ in $V'$ defines an isomorphism
$V\otimes_k k'\iso V'$.
\end{proposition}

In particular, if $V'$ has dimension $1$ and $v'\in V'$ is nonzero, then
$\varphi_\sigma$ is determined by the constant
$c(\sigma)\in {k'}^\times$ such that
$\varphi_\sigma(v')=c(\sigma)v'$, and condition c) becomes
\[
  c(\tau\sigma) = c(\tau)\cdot \tau(c(\sigma))\text{.}
\]
By the proposition, there exists a nonzero invariant vector
$v=\mu v'$, $\mu\in {k'}^\times$. Hence, for every $\sigma\in G$,
\[
  c(\sigma) = \mu\cdot \sigma(\mu^{-1})\text{.}
\]
In other words, every $1$-cocycle of $G$ with values in
${k'}^\times$ is a coboundary:

\begin{corollary}\label{I:1-5-3}
One has $\h^1(G,{k'}^\times)=0$.
\end{corollary}

\subsection{Grothendieck topologies}\label{I:1-6}

We now transpose the definitions of the preceding paragraphs to an
abstract setting that encompasses both topological spaces and schemes.

\subsubsection{}\label{I:1-6-1}

Let $\cS$ be a category and let $U$ be an object of $\cS$. A
\emph{sieve} on $U$ is a subset $\cU$ of $\ob(\cS/U)$ such that if
$\varphi:V\to U$ belongs to $\cU$ and $\psi:W\to V$ is a morphism in
$\cS$, then $\varphi\circ\psi:W\to U$ belongs to $\cU$.

If $\{\varphi_i:U_i\to U\}$ is a family of morphisms, the sieve generated
by the $U_i$ is, by definition, the set of morphisms $\varphi:V\to U$
that factor through one of the $\varphi_i$.

If $\cU$ is a sieve on $U$ and $\varphi:V\to U$ is a morphism, the
restriction $\cU_V$ of $\cU$ to $V$ is, by definition, the sieve on
$V$ consisting of the morphisms $\psi:W\to V$ such that
$\varphi\circ\psi:W\to U$ belongs to $\cU$.

\subsubsection{}\label{I:1-6-2}

A \emph{Grothendieck topology} on $\cS$ consists of the choice, for every
object $U$ of $\cS$, of a set $C(U)$ of sieves on $U$, called
covering sieves, such that the following axioms hold:
\begin{enumerate}[\indent a)]
  \item The sieve generated by the identity of $U$ is covering.
  \item If $\cU$ is a covering sieve on $U$ and $V\to U$ is a
    morphism, then the sieve $\cU_V$ is covering.
  \item A locally covering sieve is covering. In other words, if $\cU$
    is a covering sieve on $U$ and $\cU'$ is a sieve on $U$ such that,
    for every $V\to U$ belonging to $\cU$, the sieve $\cU_V'$ is
    covering, then $\cU'$ is covering.
\end{enumerate}

A category equipped with a Grothendieck topology is called a \emph{site}.

\subsubsection{}\label{I:1-6-3}

Given a site $\cS$, a presheaf on $\cS$ is a contravariant functor
$\sF$ from $\cS$ to the category of sets. For every object $U$ of
$\cS$, the elements of $\sF(U)$ are called sections of $\sF$ over
$U$. For every morphism $V\to U$ and every $s\in\sF(U)$, we denote by
$s|V$ ($s$ restricted to $V$) the image of $s$ in $\sF(V)$.

If $\cU$ is a sieve on $U$, a section given $\cU$-locally consists,
for every $V\to U$ belonging to $\cU$, of a section
$s_V\in\sF(V)$ such that, for every morphism $W\to V$,
$s_V|W=s_W$. The presheaf $\sF$ is called a \emph{sheaf} if, for every
object $U$ of $\cS$, every covering sieve $\cU$ on $U$, and every
section $\{s_V\}$ given $\cU$-locally, there exists a unique section
$s\in\sF(U)$ such that $s|V=s_V$ for every $V\to U$ belonging to
$\cU$.

\emph{Abelian sheaves} are defined similarly by replacing the category of
sets with the category of abelian groups. One shows that the category of
abelian sheaves on $\cS$ is an abelian category with enough injectives. A
sequence $\sF\xrightarrow f\sG\xrightarrow g\sH$ of sheaves is exact if,
for every object $U$ of $\cS$ and every $s\in\sG(U)$ such that
$g(s)=0$, there exists locally a $t$ such that $f(t)=s$; i.e. there
exists a covering sieve $\cU$ on $U$ and, for every $V\in\cU$, a
section $t_V$ of $\sF$ on $V$ such that $f(t_V)=s|V$.

\subsubsection{Examples}\label{I:1-6-4}

We have seen two examples above.

\begin{enumerate}[\indent a)]
  \item Let $X$ be a topological space and let $\cS$ be the category
    whose objects are the open subsets of $X$ and whose morphisms are the
    natural inclusions. The Grothendieck topology on $\cS$ corresponding
    to the usual topology of $X$ is the one for which a sieve $\cU$ on
    an open subset $U$ of $X$ is covering if the union of the open
    subsets belonging to the sieve equals $U$. It is clear that the
    category of sheaves on $\cS$ is equivalent to the category of sheaves
    on $X$ in the usual sense.
  \item Let $X$ be a scheme and let $\cS$ be the category of schemes over
    $X$. The fpqc (faithfully flat quasi-compact) topology on $\cS$ is
    the Grothendieck topology for which a sieve on an $X$-scheme $U$ is
    covering if it is generated by a finite family of flat morphisms whose
    images cover $U$.
\end{enumerate}

\subsubsection{Cohomology}\label{I:1-6-5}

We shall always assume that the category $\cS$ has a final object $X$.
The global sections of an abelian sheaf $\sF$, denoted by $\Gamma\sF$ or
$\h^0(X,\sF)$, form the group $\sF(X)$. The functor
$\sF\mapsto\Gamma\sF$ is a left exact functor from the category of
abelian sheaves on $\cS$ to the category of abelian groups; its derived
functors (or satellites) are denoted by $\h^i(X,-)$. These cohomology
groups represent the obstructions to passing from the local to the global.
By definition, if $0\to\sF\to\sG\to\sH\to0$ is an exact sequence of
abelian sheaves, then there is a long exact cohomology sequence:
\[\xymatrix{
  0 \ar[r]
    & \h^0(X,\sF) \ar[r]
    & \h^0(X,\sG) \ar[r]
    & \h^0(X,\sH) \ar[r]
    & \h^1(X,\sF) \ar[r]
    & \cdots \\
  \ldots \ar[r]
    & \h^n(X,\sF) \ar[r]
    & \h^n(X,\sG) \ar[r]
    & \h^n(X,\sH) \ar[r]
    & \h^{n+1}(X,\sF) \ar[r]
    & \cdots
}\]

\subsubsection{}\label{I:1-6-6}

Given an abelian sheaf $\sF$ on $\cS$, an \emph{$\sF$-torsor} is a
sheaf $\sG$ equipped with an action $\sF\times\sG\to\sG$ of $\sF$
such that locally (after restriction to all the objects of a covering
sieve of the final object $X$), $\sG$ with its $\sF$-action is
isomorphic to $\sF$ with the canonical translation action
$\sF\times\sF\to\sF$.

One can show that $\h^1(X,\sF)$ may be interpreted as the set of
isomorphism classes of $\sF$-torsors.

\section{The étale topology}\label{I:2}

We specialize the definitions of the preceding chapter to the étale
topology of a scheme $X$ (\ref{I:2-1}, \ref{I:2-2}, \ref{I:2-3}). When
$X$ is the spectrum of a field $K$, the corresponding cohomology agrees
with the Galois cohomology of $K$ (\ref{I:2-4}).

\subsection{The étale topology}\label{I:2-1}

We begin by recalling some facts about étale morphisms.

\begin{definition}\label{I:2-1-1}
Let $A$ be a (commutative) ring. An $A$-algebra $B$ is called étale if
$B$ is an $A$-algebra of finite presentation and the following equivalent
conditions hold:
\begin{enumerate}[\indent a)]
  \item For every $A$-algebra $C$ and every square-zero ideal $J$ of $C$,
    the canonical map
    \[
      \hom_{A\textnormal{-alg}}(B,C) \to \hom_{A\textnormal{-alg}}(B,C/J)
    \]
    is a bijection.
  \item $B$ is a flat $A$-module and $\Omega_{B/A}=0$ (where
    $\Omega_{B/A}$ denotes the module of relative differentials).
  \item Let $B=A[X_1,\dotsc,X_n]/I$ be a presentation of $B$. Then, for
    every prime ideal $\fp$ of $A[X_1,\dotsc,X_n]$ containing $I$, there
    exist polynomials $P_1,\dotsc,P_n\in I$ such that $I_\fp$ is generated
    by the images of $P_1,\dotsc,P_n$ and
    $\det(\partial P_i/\partial X_j)\notin\fp$.
\end{enumerate}
\end{definition}
(cf. \cite[I]{sga1} or \cite[V]{ra70})
% [cf. SGA I, Exposé I, or M. {\sc Raynoud}, \emph{Henselian Local Rings}, Chapter V].

A morphism of schemes $f:X\to S$ is called \emph{étale} if, for every
$x\in X$, there exist an affine open neighbourhood $U=\spec(A)$ of
$f(x)$ and an affine open neighbourhood $V=\spec(B)$ of $x$ in
$X\times_SU$ such that $B$ is an étale $A$-algebra.

\subsubsection{Examples}\label{I:2-1-2}
\begin{enumerate}[\indent a)]
  \item If $A$ is a field, an $A$-algebra $B$ is étale if and only if it
    is a finite product of separable extensions of $A$.
  \item If $X$ and $S$ are schemes of finite type over $\dC$, a morphism
    $f:X\to S$ is étale if and only if its analytification
    $f^{\text{an}}:X^{\text{an}}\to S^{\text{an}}$ is a local
    isomorphism.
\end{enumerate}

\subsubsection{Consequences}\label{I:2-1-3}
\begin{enumerate}[\indent a)]
  \item (base change) If $f:X\to S$ is an étale morphism, then so is
    $f_{S'}:X\times_SS'\to S'$ for every morphism $S'\to S$.
  \item If $f:X\to S$ and $g:Y\to S$ are two étale morphisms, then every
    $S$-morphism from $X$ to $Y$ is étale.
  \item (descent) Let $f:X\to S$ be a morphism. If there exists a
    faithfully flat morphism $S'\to S$ such that
    $f_{S'}:X\times_SS'\to S'$ is étale, then $f$ is étale.
\end{enumerate}

\subsubsection{}\label{I:2-1-4}

Let $X$ be a scheme, and let $\cS$ be the category of étale
$X$-schemes. By (\ref{I:2-1-3}.c), every morphism in $\cS$ is étale.
The \emph{étale topology} on $\cS$ is the topology for which a sieve on
$U$ is covering if it is generated by a finite family of morphisms
$\varphi_i:U_i\to U$ such that the union of the images of the
$\varphi_i$ covers $U$. The \emph{étale site} of $X$, denoted by
$\et X$, is the site defined by $\cS$ with the étale topology.

\subsection{Examples of sheaves}\label{I:2-2}

\subsubsection{Constant sheaf}\label{I:2-2-1}

Let $C$ be an abelian group and, for simplicity, suppose that $X$ is
noetherian. We denote by $\const C_X$ (or simply $C$ if there is no
ambiguity) the sheaf defined by $U\mapsto C^{\pi_0(U)}$, where
$\pi_0(U)$ is the (finite) set of connected components of $U$. The most
important case will be $C=\dZ/n$. Thus, by definition,
\[
  \h^0(X,\dZ/n)=\left(\dZ/n\right)^{\pi_0(X)}\text{.}
\]
Moreover, $\h^1(X,\dZ/n)$ is the set of isomorphism classes of
$\dZ/n$-torsors (\ref{I:1-6-6}), in other words, of finite Galois étale
coverings of $X$ with group $\dZ/n$. In particular, if $X$ is connected
and $\pi_1(X)$ is its fundamental group for a chosen base point, then
\[
  \h^1(X,\dZ/n) = \hom(\pi_1(X),\dZ/n)\text{.}
\]

\subsubsection{Multiplicative group}\label{I:2-2-2}

We denote by $\dG_{m,X}$ (or $\dG_m$ if there is no ambiguity) the sheaf
defined by $U\mapsto\Gamma(U,\sO_U^\times)$; it is indeed a sheaf by the
faithfully flat descent theorem (\ref{I:1-4-5}). By definition,
\[
  \h^0(X,\dG_m) = \h^0(X,\sO_X)^\times\text{;}
\]
in particular, if $X$ is reduced, connected, and proper over an
algebraically closed field $k$, then
\[
  \h^0(X,\dG_m) = k^\times\text{.}
\]

\begin{proposition}\label{I:2-2-3}
There is an isomorphism
\[
  \h^1(X,\dG_m) = \pic(X)\text{,}
\]
where $\pic(X)$ is the group of isomorphism classes of invertible sheaves
on $X$.
\end{proposition}
\begin{proof}
Let $*$ be the functor that associates to an invertible sheaf $\sL$ on
$X$ the following presheaf $\sL^*$ on $\et X$: for
$\varphi:U\to X$ étale,
\[
  \sL^*(U)=\operatorname{Isom}_U(\sO_U,\varphi^*\sL)\text{.}
\]
% The original references here are just to (4.2) and (4.5), but the context
% makes it seem that they refer to the results in Chapter 1. Likewise a little
% later.
By (\ref{I:1-4-2}.i) and (\ref{I:1-4-5}) (full faithfulness), this
presheaf is a sheaf; indeed, it is a $\dG_m$-torsor. One checks
immediately that
\begin{enumerate}[\indent a)]
  \item the functor $*$ is compatible with (étale) localization;
  \item it induces an equivalence from the category of trivial invertible
    sheaves (i.e. isomorphic to $\sO_X$) to the category of trivial
    $\dG_m$-torsors: $\sL$ is trivial if and only if $\sL^*$ is trivial.
\end{enumerate}

Moreover, by (\ref{I:1-4-2}.ii) and (\ref{I:1-4-5}),
\begin{enumerate}[\indent a)]
\setcounter{enumi}{2}
  \item the notion of an invertible sheaf is local for the étale topology.
\end{enumerate}

It follows formally from a), b), and c) that $*$ is an equivalence between
the category of invertible sheaves on $X$ and that of $\dG_m$-torsors on
$\et X$; it induces the desired isomorphism. The inverse equivalence is
constructed as follows. If $\sT$ is a $\dG_m$-torsor, there exists a
finite étale covering $\{U_i\}$ of $X$ such that the torsors
$\sT/U_i$ are trivial; $\sT$ is then trivial on every $V$ étale over
$X$ belonging to the sieve $\cU\subset\et X$ generated by $\{U_i\}$.
On every $V\in\cU$, $\sT|V$ corresponds to an invertible sheaf $\sL_V$
(by b)), and the $\sL_V$ form an invertible sheaf given $\cU$-locally,
$\sL_\cU$ (by a)). By c), the latter comes from an invertible sheaf
$\sL(\sT)$ on $X$, and $\sT\mapsto\sL(\sT)$ is the desired inverse
of $*$.
\end{proof}

\subsubsection{Roots of unity}\label{I:2-2-4}

For every integer $n>0$, the sheaf of $n$th roots of unity, denoted by
$\dmu_n$, is the kernel of the $n$th-power map on $\dG_m$.
If $X$ is a scheme over a separably closed field $k$ and $n$ is
invertible in $k$, the choice of a primitive $n$th root of unity
$\zeta\in k$ defines an isomorphism $i\mapsto\zeta^i$ from $\dZ/n$ to
$\dmu_n$.

The relation between cohomology with coefficients in $\dmu_m$ and
cohomology with coefficients in $\dG_m$ is given by the cohomology exact
sequence deduced from the following theorem.

\begin{theorem}[Kummer theory]\label{I:2-2-5}
If $n$ is invertible on $X$, the $n$th-power map on $\dG_m$ is an
epimorphism of sheaves. Hence there is an exact sequence
\[
  0 \to \dmu_n \to \dG_m \to \dG_m \to 0\text{.}
\]
\end{theorem}
\begin{proof}
Let $U\to X$ be an étale morphism and let
$a\in\dG_m(U)=\Gamma(U,\sO_U^\times)$. Since $n$ is invertible on $U$,
the equation $T^n-a=0$ is separable; in other words,
$U'=\spec\left(\sO_U[T]/(T^n-a)\right)$ is étale over $U$. Moreover,
$U'\to U$ is surjective and this element has an $n$th root on $U'$, which proves
the result.
\end{proof}

\subsection{Fibres, direct images}\label{I:2-3}

\subsubsection{}\label{I:2-3-1}

A \emph{geometric point} of $X$ is a morphism $\bar x\to X$, where
$\bar x$ is the spectrum of a separably closed field $k(\bar x)$. By
abuse of notation, we denote it by $\bar x$, leaving the morphism
$\bar x\to X$ implicit. If $x$ is the image of $\bar x$ in $X$, we say
that $\bar x$ is centred at $x$. If $k(\bar x)$ is an algebraic extension
of the residue field $k(x)$, then $\bar x$ is called an
\emph{algebraic} geometric point of $X$.

An \emph{étale neighbourhood} of $\bar x$ is a commutative diagram
\[\xymatrix{
  & U \ar[d] \\
  \bar x \ar[ur] \ar[r]
  & X\text{,}
}\]
where $U\to X$ is étale.

The \emph{strict localization} of $X$ at $\bar x$ is the ring
$\sO_{X,\bar x}=\varinjlim\Gamma(U,\sO_U)$, where the direct limit is
taken over the étale neighbourhoods of $\bar x$. It is a strictly henselian
local ring whose residue field is the separable closure in $k(\bar x)$
of the residue field $k(x)$ of $X$ at $x$. It plays the role of the local
ring for the étale topology.

\subsubsection{}\label{I:2-3-2}

Given a sheaf $\sF$ on $\et X$, the \emph{fibre} of $\sF$ at $\bar x$
is the set (resp. group, \dots)
$\sF_{\bar x}=\varinjlim\sF(U)$, where the direct limit is again taken
over the étale neighbourhoods of $\bar x$.

A homomorphism of sheaves $\sF\to\sG$ is a mono-/epi-/isomorphism if and
only if the induced morphisms $\sF_{\bar x}\to\sG_{\bar x}$ on fibres
have the corresponding property at every geometric point of $X$. If $X$
is of finite type over an algebraically closed field, it suffices that
this hold at the rational points of $X$.

\subsubsection{}\label{I:2-3-3}

If $f:X\to Y$ is a morphism of schemes and $\sF$ is a sheaf on
$\et X$, the \emph{direct image} $f_*\sF$ of $\sF$ under $f$ is the
sheaf on $\et Y$ defined by
$f_*\sF(V)=\sF(X\times_YV)$ for every $V$ étale over $Y$. The functor
$f_*:\sh(\et X)\to\sh(\et Y)$ is left exact. Its right derived functors
$\eR^qf_*$ are called the higher direct images. If $\bar y$ is a
geometric point of $Y$, then
\[
  (\eR^q f_* \sF)_{\bar y} = \varinjlim \h^q(V\times_Y X,\sF)\text{,}
\]
where the direct limit is taken over the étale neighbourhoods $V$ of
$\bar y$.

Let $\sO_{Y,\bar y}$ be the strict localization of $Y$ at $\bar y$,
let $\widetilde Y=\spec(\sO_{Y,\bar y})$, and let
$\widetilde X=X\times_Y\widetilde Y$. The sheaf $\sF$ may be extended
to $\et{\widetilde X}$ (this is a special case of the general notion of
inverse image) as follows. Let $\widetilde U$ be a scheme étale over
$\widetilde X$. Then there exist an étale neighbourhood $V$ of $\bar y$
and a scheme $U$ étale over $X\times_YV$ such that
$\widetilde U=U\times_V\widetilde Y$; set
\[
  \sF(\widetilde U)=\varinjlim \sF\left(U\times_V V'\right)\text{,}
\]
where the direct limit is taken over the étale neighbourhoods $V'$ of
$\bar y$ that dominate $V$. With this definition,
\[
  \left(\eR^q f_* \sF\right)_{\bar y} = \h^q(\widetilde X,\sF)\text{.}
\]

The functor $f_*$ has a left adjoint $f^*$, the “inverse image” functor.
If $\bar x$ is a geometric point of $X$ and $f(\bar x)$ its image in
$Y$, then $(f^*\sF)_{\bar x}=\sF_{f(\bar x)}$. This formula shows that
$f^*$ is exact. Consequently, $f_*$ sends injective sheaves to injective
sheaves, and the spectral sequence for the composite functor
$\Gamma\circ f_*$ (resp. $g_*f_*$) gives the following theorem.

\begin{theorem}[Leray spectral sequence]\label{I:2-3-4}
Let $\sF$ be an abelian sheaf on $\et X$ and let $f:X\to Y$ be a
morphism of schemes (resp. let
$X\xrightarrow fY\xrightarrow gZ$ be morphisms of schemes). There is a
spectral sequence
\begin{align*}
  E_2^{pq} &= \h^p(Y,\eR^q f_* \sF) \Rightarrow \h^{p+q}(X,\sF) \\
  \text{(resp.}\qquad E_2^{pq} &= \eR^p g_* \eR^q f_* \sF \Rightarrow \eR^{p+q}(gf)_* \sF\text{).}
\end{align*}
\end{theorem}

\begin{corollary}\label{I:2-3-5}
If $\eR^qf_*\sF=0$ for every $q>0$, then
$\h^p(Y,f_*\sF)=\h^p(X,\sF)$ (resp.
$\eR^pg_*(f_*\sF)=\eR^p(gf)_*\sF$) for every $p\geqslant0$.
\end{corollary}

This applies in particular in the following case:

\begin{proposition}\label{I:2-3-6}
Let $f:X\to Y$ be a finite morphism (or even, by passage to the limit,
an integral morphism), and let $\sF$ be an abelian sheaf on $X$. Then
$\eR^qf_*\sF=0$ for every $q>0$.
\end{proposition}

Indeed, let $\bar y$ be a geometric point of $Y$, let $\widetilde Y$ be
the spectrum of the strict localization of $Y$ at $\bar y$, and let
$\widetilde X=X\times_Y\widetilde Y$. By the preceding discussion, it
suffices to prove that $\h^q(\widetilde X,\sF)=0$ for every $q>0$. But
$\widetilde X$ is the spectrum of a product of strictly henselian local
rings (cf. \cite[I]{ra70}), and the functor
$\Gamma(\widetilde X,-)$ is exact because every surjective étale
$\widetilde X$-scheme admits a section. This proves the assertion.

\subsection{Galois cohomology}\label{I:2-4}

For $X=\spec(K)$, the spectrum of a field, we shall see that étale
cohomology can be identified with Galois cohomology.

\subsubsection{}\label{I:2-4-1}

We begin with a topological analogy. If $K$ is the function field of an
integral affine algebraic variety $Y=\spec(A)$ over $\dC$, then
$K=\varinjlim_{f\in A}A[1/f]$.

In other words, $X=\varprojlim U$, where $U$ runs over the open subsets
of $Y$. It is known that there exist arbitrarily small Zariski open
subsets that are $K(\pi,1)$ spaces for the classical topology. It is
therefore natural to regard $\spec(K)$ itself as a $K(\pi,1)$, where
$\pi$ is the fundamental group (in the algebraic sense) of $X$, namely
the Galois group of $\bar K/K$, where $\bar K$ is a separable closure of
$K$.

\subsubsection{}\label{I:2-4-2}

More precisely, let $K$ be a field, let $\bar K$ be a separable closure
of $K$, and let $G=\gal(\bar K/K)$ be the topological Galois group. To
every finite étale $K$-algebra $A$ (a finite product of separable
extensions of $K$), associate the finite set $\hom_K(A,\bar K)$. The
Galois group $G$ acts on this set through a discrete (hence finite)
quotient. If $A=K[T]/(F)$, this set is identified with the set of roots
of $F$ in $\bar K$. Galois theory, in the form given to it by
Grothendieck, says:

\begin{proposition}\label{I:2-4-3}
The functor
\[
  \left\{\begin{array}{c}
          \textnormal{$K$-algebras} \\
          \textnormal{finite étale}
        \end{array}\right\}
  \to
  \left\{\begin{array}{c}
          \textnormal{finite sets on which} \\
          \textnormal{$G$ acts continuously}
        \end{array}\right\}
\]
that associates $\hom_K(A,\bar K)$ to an étale algebra $A$ is an
anti-equivalence of categories.
\end{proposition}

This yields an analogous description of sheaves for the étale topology on
$\spec(K)$:

\begin{proposition}\label{I:2-4-4}
The functor
\[
  \left\{\begin{array}{c}
          \textnormal{étale sheaves} \\
          \textnormal{on $\spec(K)$}
        \end{array}\right\}
  \to
  \left\{\begin{array}{c}
          \textnormal{sets on which} \\
          \textnormal{$G$ acts continuously}
        \end{array}\right\}
\]
that associates to a sheaf $\sF$ its fibre $\sF_{\bar K}$ at the
geometric point $\spec(\bar K)$ is an equivalence of categories.
\end{proposition}

The group $G$ is said to act continuously on a set $E$ if the stabilizer
of every element of $E$ is an open subgroup of $G$. The inverse functor
is described in the evident way: let $A$ be a finite étale
$K$-algebra, let $U=\spec(A)$, and let
$U(\bar K)=\hom_K(A,\bar K)$ be the corresponding $G$-set; then
$\sF(U)=\hom_{G\text{-set}}(U(\bar K),\sF_{\bar K})$.

In particular, if $X=\spec(K)$, then $\sF(X)=\sF_{\bar K}^G$. Restricting
to abelian sheaves and passing to derived functors gives canonical
isomorphisms
\[
  \h^q(\et X,\sF) = \h^q(G,\sF_{\bar K})
\]

\subsubsection{Examples}\label{I:2-4-5}

\begin{enumerate}[\indent a)]
  \item The constant sheaf $\dZ/n$ corresponds to $\dZ/n$ with the
    trivial action of $G$.
  \item The sheaf $\dmu_n$ of $n$th roots of unity corresponds to the
    group $\dmu_n(\bar K)$ of $n$th roots of unity in $\bar K$, with
    the natural action of $G$.
  \item The sheaf $\dG_m$ corresponds to the group $\bar K^\times$ with
    the natural action of $G$.
\end{enumerate}

\section{The cohomology of curves}\label{I:3}

For topological spaces, unwindings using the K\"unneth formula and
simplicial decompositions reduce the calculation of cohomology to the
interval $I=[0,1]$, for which $\h^0(I,\dZ)=\dZ$ and
$\h^q(I,\dZ)=0$ for $q>0$.

In our case, the unwindings lead to more complicated objects, namely curves
over an algebraically closed field; in this chapter we calculate their
cohomology. The situation is more complicated than in the topological case,
because the cohomology groups vanish only for $q>2$. The essential
ingredient in the calculations is the vanishing of the Brauer group of the
function field of such a curve (Tsen's theorem, \ref{I:3-2}).

\subsection{The Brauer group}\label{I:3-1}

We first recall its classical definition.

\begin{definition}\label{I:3-1-1}
Let $K$ be a field and let $A$ be a finite-dimensional $K$-algebra. The
algebra $A$ is called a central simple algebra over $K$ if the following
equivalent conditions hold:
\begin{enumerate}[\indent a)]
  \item $A$ has no nontrivial two-sided ideal, and its centre is $K$.
  \item There exists a finite Galois extension $K'/K$ such that
    $A_{K'}=A\otimes_KK'$ is isomorphic to an algebra of square matrices
    over $K'$.
  \item $A$ is $K$-isomorphic to an algebra of square matrices over a
    division algebra with centre $K$.
\end{enumerate}
\end{definition}

Two such algebras are called equivalent if the division algebras associated
with them by c) are $K$-isomorphic. If the algebras have the same
dimension, this is equivalent to their being $K$-isomorphic. Tensor product
induces, on passing to the quotient, an abelian group structure on the set
of equivalence classes. This group is classically called the
\emph{Brauer group} of $K$ and is denoted by $\br(K)$.

\subsubsection{}\label{I:3-1-2}

We denote by $\br(n,K)$ the set of $K$-isomorphism classes of
$K$-algebras $A$ such that there exists a finite Galois extension $K'$
of $K$ for which $A_{K'}$ is isomorphic to the algebra $M_n(K')$ of
$n\times n$ square matrices over $K'$. By definition, $\br(K)$ is the
union of the subsets $\br(n,K)$ for $n\in\dN$. Let $\bar K$ be an
algebraic closure of $K$ and let $G=\gal(\bar K/K)$. The set $\br(n,K)$
is the set of “forms” of $M_n(\bar K)$, and hence is canonically
isomorphic to $\h^1\left(G,\aut(M_n(\bar K))\right)$.

Every automorphism of $M_n(\bar K)$ is inner. Consequently,
$\aut(M_n(\bar K))$ is identified with the projective linear group
$\pgl(n,\bar K)$, and there is a canonical bijection
\[
  \theta_n : \br(n,K) \iso \h^1\left(G,\pgl(n,\bar K)\right)\text{.}
\]

On the other hand, the exact sequence
\begin{equation*}\tag{3.1.2.1}\label{I:eq:3-1-2-1}
  1 \to \bar K^\times \to \gl(n,\bar K) \to \pgl(n,\bar K) \to 1\text{,}
\end{equation*}
defines a coboundary operator
\[
  \Delta_n : \h^1\left(G,\pgl(n,\bar K)\right) \to \h^2(G,\bar K^\times)\text{.}
\]

Composing $\theta_n$ and $\Delta_n$ gives a map
\[
  \delta_n : \br(n,K) \to \h^2(G,\bar K^\times)\text{.}
\]
It is readily checked that the maps $\delta_n$ are compatible and define
a group homomorphism
\[
  \delta : \br(K) \to \h^2(G,\bar K^\times)\text{.}
\]

\begin{proposition}\label{I:3-1-3}
The homomorphism $\delta:\br(K)\to\h^2(G,\bar K^\times)$ is bijective.
\end{proposition}

This follows from the next two lemmas.

\begin{lemma}\label{I:3-1-4}
The map
$\Delta_n:\h^1\left(G,\pgl(n,\bar K)\right)\to
\h^2(G,\bar K^\times)$ is injective.
\end{lemma}

By \cite{se94}, Corollary to Proposition I-44, it suffices to check that
whenever the exact sequence (\ref{I:eq:3-1-2-1}) is twisted by an element
of $\h^1\left(G,\pgl(n,\bar K)\right)$, the $\h^1$ of the middle group
is trivial. This middle group is the group of $\bar K$-points of the
multiplicative group of a central simple algebra $A$ of rank $n^2$ over
$K$. To prove that $\h^1(G,A_{\bar K}^\times)=0$, interpret $A^\times$
as the automorphism group of the free $A$-module $L$ of rank $1$, and
$\h^1$ as the set of “forms” of $L$---$A$-modules of rank $n^2$ over
$K$, which are automatically free.

\begin{lemma}\label{I:3-1-5}
Let $\alpha\in\h^2(G,\bar K^\times)$, let $K'$ be a finite extension of
$K$ contained in $\bar K$, let $n=[K':K]$, and let
$G'=\gal(\bar K/K')$. If the image of $\alpha$ in
$\h^2(G',\bar K^\times)$ is zero, then $\alpha$ belongs to the image of
$\Delta_n$.
\end{lemma}

First note that
\[
  \h^2(G',\bar K^\times) \simeq \h^2\left(G,(\bar K\otimes_K K')^\times\right)\text{.}
\]
(Geometrically, if $x=\spec(K)$, $x'=\spec(K')$, and
$\pi:x'\to x$ is the canonical morphism, then
$\eR^q\pi_*(\dG_{m,x'})=0$ for $q>0$ and consequently
$\h^q(x',\dG_{m,x'})\simeq\h^q(x,\pi_*\dG_{m,x'})$ for
$q\geqslant0$.)

Moreover, the choice of a basis of $K'$ as a vector space over $K$
defines a homomorphism
\[
  (\bar K\otimes_K K')^\times \to \gl(n,\bar K)
\]
that sends an element $x$ to the endomorphism of $\bar K\otimes_KK'$
given by multiplication by $x$. There is then a commutative diagram with
exact rows:
\[\xymatrix{
  1 \ar[r]
    & \bar K^\times \ar[r] \ar@{=}[d]
    & (\bar K\otimes_K K')^\times \ar[r] \ar[d]
    & (\bar K\otimes_K K')^\times / \bar K^\times \ar[r] \ar[d]
    & 1 \\
  1 \ar[r]
    & \bar K^\times \ar[r]
    & \gl(n,\bar K) \ar[r]
    & \pgl(n,\bar K) \ar[r]
    & 1
}\]
The lemma follows from the commutative diagram obtained by passing to
cohomology:
\[\xymatrix{
  \h^1\left(G,(\bar K\otimes_K K')^\times/\bar K^\times\right) \ar[r] \ar[d]
    & \h^2(G,\bar K^\times) \ar[r] \ar@{=}[d]
    & \h^2\left(G,(\bar K\otimes_K K')^\times\right) \\
  \h^1\left(G,\pgl(n,\bar K)\right) \ar[r]^-{\Delta_n}
    & \h^2(G,\bar K^\times) \text{.}
}\]

\begin{proposition}\label{I:3-1-6}
Let $K$ be a field, let $\bar K$ be an algebraic closure of $K$, and let
$G=\gal(\bar K/K)$. Suppose that $\br(K')=0$ for every finite extension
$K'$ of $K$. Then:
\begin{enumerate}[\indent i)]
  \item $\h^q(G,\bar K^\times)=0$ for every $q>0$.
  \item $\h^q(G,\sF)=0$ for every torsion $G$-module $\sF$ and every
    $q\geqslant2$.
\end{enumerate}
\end{proposition}

(For the proof, cf. \cite{se94}.)

\subsection{Tsen's theorem}\label{I:3-2}

\begin{definition}\label{I:3-2-1}
A field $K$ is called $C_1$ if every nonconstant homogeneous polynomial
$f(x_1,\dotsc,x_n)$ of degree $d<n$ has a nontrivial zero.
\end{definition}

\begin{proposition}\label{I:3-2-2}
If a field $K$ is $C_1$, then $\br(K)=0$.
\end{proposition}

We must show that every division algebra $D$ with centre $K$ and finite
over $K$ is equal to $K$. Let $r^2$ be the degree of $D$ over $K$ and
let $\operatorname{Nrd}:D\to K$ be the reduced norm.

(Locally for the étale topology on $K$, $D$ is isomorphic---not
canonically---to a matrix algebra $M_r$, and the reduced norm agrees with
the determinant map. This is well defined independently of the chosen
isomorphism between $D$ and $M_r$, because every automorphism of $M_r$
is inner and similar matrices have the same determinant. This map, defined
locally for the étale topology, descends, by its local uniqueness, to a map
$\operatorname{Nrd}:D\to K$.)

The only zero of $\operatorname{Nrd}$ is the zero element of $D$, since
for $x\ne0$ one has
$\operatorname{Nrd}(x)\cdot\operatorname{Nrd}(x^{-1})=1$. On the other
hand, if $\{e_1,\dotsc,e_{r^2}\}$ is a basis of $D$ over $K$ and
$x=\sum x_ie_i$, then $\operatorname{Nrd}(x)$ is a homogeneous
polynomial $\operatorname{Nrd}(x_1,\dotsc,x_{r^2})$ of degree $r$
(this is clear locally for the étale topology). Since $K$ is $C_1$, one
has $r^2\leqslant r$, that is, $r=1$ and $D=K$.

\begin{theorem}[Tsen]\label{I:3-2-3}
Let $k$ be an algebraically closed field and let $K$ be an extension of
$k$ of transcendence degree $1$. Then $K$ is $C_1$.
\end{theorem}

First suppose that $K=k(X)$. Let
\[
  f(T) = \sum a_{i_1,\dotsc i_n} T_1^{i_1} \dotsm T_n^{i_n}
\]
be a homogeneous polynomial of degree $d<n$ with coefficients in
$k(X)$. After multiplying the coefficients by a common denominator, we
may suppose that they lie in $k[X]$. Let
$\delta=\sup\deg(a_{i_1\dotsc i_n})$. We seek a nontrivial zero in
$k[X]$ by the method of undetermined coefficients, writing each $T_i$
($i=1,\dotsc,n$) as a polynomial of degree $N$ in $X$. The equation
$f(T)=0$ then becomes a system of homogeneous equations in the
$n\times(N+1)$ coefficients of the polynomials $T_i(X)$, expressing the
vanishing of the coefficients of the polynomial in $X$ obtained by
replacing $T_i$ with $T_i(X)$. This polynomial has degree at most
$\delta+Nd$, so there are $\delta+Nd+1$ equations in
$n\times(N+1)$ variables. Since $k$ is algebraically closed, this system
has a nontrivial solution if $n(N+1)>Nd+\delta+1$, which holds for
sufficiently large $N$ when $d<n$.

To prove the theorem in general, it plainly suffices to prove it when
$K$ is a finite extension of the purely transcendental extension
$k(X)$ of $k$. Let $f(T)=f(T_1,\dotsc,T_n)$ be a homogeneous
polynomial of degree $d<n$ with coefficients in $K$. Let
$s=[K:k(X)]$ and let $e_1,\dotsc,e_s$ be a basis of $K$ over $k(X)$.
Introduce $sn$ new variables $U_{ij}$ such that
$T_i=\sum U_{ij}e_j$. For $f(T)$ to have a nontrivial zero in $K$, it
suffices that the polynomial $g(U_{ij})=N_{K/k(X)}(f(T))$ have a nontrivial
zero in $k(X)$. But $g$ is a polynomial of degree $sd$ in $sn$
variables, which proves the result.

\begin{corollary}\label{I:3-2-4}
Let $k$ be an algebraically closed field and let $K$ be an extension of
$k$ of transcendence degree $1$. Then the étale cohomology groups
$\h^q(\spec(K),\dG_m)$ vanish for every $q>0$.
\end{corollary}

\subsection{Cohomology of smooth curves}

From now on, unless expressly stated otherwise, all cohomology groups are
étale cohomology groups.

\begin{proposition}\label{I:3-3-1}
Let $k$ be an algebraically closed field and let $X$ be a connected
nonsingular projective curve over $k$. Then
\begin{align*}
  \h^0(X,\dG_m) &= k^\times \text{,} \\
  \h^1(X,\dG_m) &= \pic(X) \text{,} \\
  \h^q(X,\dG_m) &= 0 \text{ for $q\geqslant 2$.}
\end{align*}
\end{proposition}

Let $\eta$ be the generic point of $X$, let $j:\eta\to X$ be the
canonical morphism, and let $\dG_{m,\eta}$ be the multiplicative group of
the function field $K(X)$. For every closed point $x$ of $X$, let
$i_x:x\to X$ be the canonical immersion and let $\dZ_x$ be the constant
sheaf with value $\dZ$ on $x$. Thus $j_*\dG_{m,\eta}$ is the sheaf of
nonzero meromorphic functions on $X$ and
$\oplus_{x\in X}i_{x*}\dZ_x$ is the sheaf of divisors; hence there is an
exact sequence of sheaves:
\begin{equation*}\tag{3.3.1.1}\label{I:eq:3-3-1-1}
\xymatrix{
  0 \ar[r]
    & \dG_m \ar[r]
    & j_* \dG_{m,\eta} \ar[r]^-{\dv}
    & \displaystyle\bigoplus_{x\in X} i_{x*} \dZ_x \ar[r]
    & 0\text{.}
}
\end{equation*}

\begin{lemma}\label{I:3-3-2}
One has $\eR^qj_*\dG_{m,\eta}=0$ for every $q>0$.
\end{lemma}

It suffices to show that the fibre of this sheaf at every closed point
$x$ of $X$ is zero. If $\tilde\sO_{X,x}$ is the henselization of $X$ at
$x$ and $K$ is the fraction field of $\tilde\sO_{X,x}$, then
\[
  \spec(K) = \eta\times_X \spec(\tilde\sO_{X,x})\text{,}
\]
and hence $(\eR^qj_*\dG_{m,\eta})_x=\h^q(\spec(K),\dG_m)$.

Now $K$ is an algebraic extension of $k(X)$, hence an extension of $k$
of transcendence degree $1$; the lemma follows from (\ref{I:3-2-4}).

\begin{lemma}\label{I:3-3-3}
One has $\h^q(X,j_*\dG_{m,\eta})=0$ for every $q>0$.
\end{lemma}

Indeed, (\ref{I:3-3-2}) and the Leray spectral sequence for $j$ give
\[
  \h^q(X, j_* \dG_{m,\eta}) = \h^q(\eta,\dG_{m,\eta})
\]
for every $q\geqslant0$, and the right-hand side vanishes for $q>0$ by
(\ref{I:3-2-4}).

\begin{lemma}\label{I:3-3-4}
One has
$\h^q\left(X,\bigoplus_{x\in X}i_{x*}\dZ_x\right)=0$ for every $q>0$.
\end{lemma}

Indeed, for every closed point $x$ of $X$ one has
$\eR^qi_{x*}\dZ_x=0$ for $q>0$, because $i_x$ is a finite morphism
(\ref{I:2-3-6}), and
\[
  \h^q(X,i_{x*}\dZ_x) = \h^q(x,\dZ_x)\text{.}
\]
The right-hand side vanishes for every $q>0$ because $x$ is the spectrum
of an algebraically closed field. (This also shows that the lemma holds
more generally for every skyscraper sheaf on $X$.)

The preceding lemmas and the exact sequence (\ref{I:eq:3-3-1-1}) give
\[
  \h^q(X,\dG_m) = 0 \text{ for $q\geqslant 2$,}
\]
and the low-degree exact cohomology sequence
\[
  1 \to \h^0(X,\dG_m) \to \h^0(X, j_*\dG_{m,\eta}) \to \h^0\left(X,\bigoplus_{x\in X} i_{x*}\dZ_x\right) \to \h^1(X,\dG_m) \to 1
\]
which is precisely the exact sequence
\[
  1 \to k^\times \to k(X)^\times\to \Div(X) \to \pic(X) \to 1\text{.}
\]
Proposition \ref{I:3-3-1} implies that the cohomology groups of $X$ with
values in $\dZ/n$, for $n$ prime to the characteristic of $k$, have
reasonable values:

\begin{corollary}\label{I:3-3-5}
If $X$ has genus $g$ and $n$ is invertible in $k$, then
$\h^q(X,\dZ/n)$ vanishes for $q>2$ and is free over $\dZ/n$ of ranks
$1,2g,1$ for $q=0,1,2$. Replacing $\dZ/n$ by the isomorphic group
$\dmu_n$, there are canonical isomorphisms
\begin{align*}
  \h^0(X,\dmu_n) &= \dmu_n \\
  \h^1(X,\dmu_n) &= \pic^0(X)_n \\
  \h^2(X,\dmu_n) &= \dZ/n \text{.}
\end{align*}
\end{corollary}

Since $k$ is algebraically closed, $\dZ/n$ is isomorphic
(noncanonically) to $\dmu_n$. The Kummer exact sequence
\[
  0 \to \dmu_n \to \dG_m \to \dG_m \to 0\text{,}
\]
and Proposition \ref{I:3-3-1} yield the equalities
\[
  \h^q(X,\dZ/n) = 0 \text{ for $q>2$,}
\]
and, in low degrees, the exact sequences
\[\xymatrix{
  0 \ar[r]
    & \h^0(X,\dmu_n) \ar[r]
    & k^\times \ar[r]^-n
    & k^\times \ar[r]
    & 0
}\]
\[\xymatrix{
  0 \ar[r]
    & \h^1(X,\dmu_n) \ar[r]
    & \pic(X) \ar[r]^-n
    & \h^2(X,\dmu_n) \ar[r]
    & 0\text{.}
}\]
Moreover, there is an exact sequence
\[\xymatrix{
  0 \ar[r]
    & \pic^0(X) \ar[r]
    & \pic(X) \ar[r]^-{\deg}
    & \dZ \ar[r]
    & 0 \text{,}
}\]
and $\pic^0(X)$ is identified with the group of $k$-rational points of
an abelian variety of dimension $g$, the Jacobian of $X$. In such a
group, multiplication by $n$ is surjective and its kernel is a free
$\dZ/n\dZ$-module of rank $2g$ (because $n$ is invertible in $k$);
this proves the corollary.

A clever unwinding using the “trace method” yields the following
proposition as a corollary.

\begin{proposition}[{\cite[IX 5.7]{sga4}}]\label{I:3-3-6}
Let $k$ be an algebraically closed field, let $X$ be an algebraic curve
over $k$, and let $\sF$ be a torsion sheaf on $X$. Then:
\begin{enumerate}[\indent i)]
  \item $\h^q(X,\sF)=0$ for $q>2$.
  \item If $X$ is affine, then even $\h^q(X,\sF)=0$ for $q>1$.
\end{enumerate}
\end{proposition}

For the proof, as well as for an exposition of the “trace method,” we
refer to \cite[IX 5]{sga4}.

\subsection{Unwindings}\label{I:3-4}

To calculate the cohomology of varieties of dimension $>1$, one uses
fibrations by curves, thereby reducing to the study of morphisms whose
fibres have dimension $\leqslant1$. This principle has several variants;
we mention some of them.

\subsubsection{}\label{I:3-4-1}

Let $A$ be a $k$-algebra of finite type, and let $a_1,\dotsc,a_n$ be
generators of $A$. If we put $X_0=\spec(k)$,
$X_i=\spec(k[a_1,\dotsc,a_i])$, and $X_n=\spec(A)$, the canonical
inclusions
$k[a_1,\dotsc,a_i]\to k[a_1,\dotsc,a_i,a_{i+1}]$ define morphisms
$X_n\to X_{n-1}\to\cdots\to X_1\to X_0$ whose fibres have dimension
$\leqslant1$.

\subsubsection{}\label{I:3-4-2}

For a smooth morphism, one can be more precise. An
\emph{elementary fibration} is a morphism of schemes $f:X\to S$ that
can be embedded in a commutative diagram
\[\xymatrix{
  X \ar@{^{(}->}[r]^-j \ar[dr]_-f
    & \bar X \ar[d]^-{\bar f}
    & \ar@{_{(}->}[l]_-i \ar[dl]^-g Y \\
  & S
}\]
satisfying the following conditions:
\begin{enumerate}[\indent i)]
  \item $j$ is an open immersion, dense in every fibre, and
    $X=\bar X\setminus Y$.
  \item $\bar f$ is smooth and projective, with irreducible geometric
    fibres of dimension $1$.
  \item $g$ is an étale covering and no fibre of $g$ is empty.
\end{enumerate}

A \emph{good neighbourhood} relative to $S$ is an $S$-scheme $X$ for
which there exist $S$-schemes $X=X_n,\dotsc,X_0=S$ and elementary
fibrations $f_i:X_i\to X_{i-1}$, $i=1,\dotsc,n$. One can show
\cite[XI 3.3]{sga4} that if $X$ is a smooth scheme over an algebraically
closed field $k$, then every rational point of $X$ has an open
neighbourhood that is a good neighbourhood (relative to $\spec(k)$).

\subsubsection{}\label{I:3-4-3}

A proper morphism $f:X\to S$ can be unwound as follows. By Chow's lemma,
there is a commutative diagram
\[\xymatrix{
  X \ar[dr]_-f
    & \ar[l]_-\pi \ar[d]^-{\bar f} \bar X \\
  & S
}\]
in which $\pi$ and $\bar f$ are projective morphisms, with $\pi$ in
addition an isomorphism over a dense open subset of $X$. Locally on
$S$, $\bar X$ is a closed subscheme of a standard projective space
$\dP_S^n$.

The latter is unwound by considering the projection
$\varphi:\dP_S^n\to\dP_S^1$ that sends the point with homogeneous
coordinates $(x_0,x_1,\dotsc,x_n)$ to $(x_0,x_1)$. This is a rational
map defined outside the closed subset $Y\simeq\dP_S^{n-2}$ of
$\dP_S^n$ given by the homogeneous equations $x_0=x_1=0$. Let
$u:P\to\dP_S^n$ be the blow-up with centre $Y$; the fibres of $u$ have
dimension $\leqslant1$. Moreover, there is a natural morphism
$v:P\to\dP_S^1$ extending the rational map $\varphi$, and $v$ makes
$P$ a $\dP_S^1$-scheme locally isomorphic to the standard projective
space $\dP^{n-1}$, which may in turn be projected onto a $\dP^1$, and
so on.

\subsubsection{}\label{I:3-4-4}

A smooth projective variety $X$ can be swept out by a
\emph{Lefschetz pencil}. The blow-up $\widetilde X$ of the intersection
of the axis of the pencil with $X$ maps to $\dP^1$, and the fibres of
this map are the hyperplane sections of $X$ by the hyperplanes of the
pencil.

\section{The proper base change theorem}\label{I:4}

\subsection{Introduction}\label{I:4-1}

This chapter is devoted to the proof and applications of the following
theorem.

\begin{theorem}\label{I:4-1-1}
Let $f:X\to S$ be a proper morphism of schemes and let $\sF$ be a
torsion abelian sheaf on $X$. Then, for every $q\geqslant0$, the fibre
of $\eR^qf_*\sF$ at a geometric point $s$ of $S$ is isomorphic to the
cohomology $\h^q(X_s,\sF)$ of the fibre
$X_s=X\otimes_S\spec k(s)$ of $f$ at $s$.
\end{theorem}

For a proper separated continuous map $f:X\to S$ between topological
spaces (separated means that the diagonal of $X\times_SX$ is closed) and
an abelian sheaf $\sF$ on $X$, the analogous result is well known and
elementary. Since $f$ is closed, the $f^{-1}(V)$, with $V$ ranging over
the neighbourhoods of $s$, form a fundamental system of neighbourhoods
of $X_s$, and one checks that
$\h^\bullet(X_s,\sF)=\varinjlim_U\h^\bullet(U,\sF)$, where $U$ ranges
over the neighbourhoods of $X_s$. In practice, $X_s$ even has a
fundamental system $\cU$ of neighbourhoods $U$ of which it is a
deformation retract, and for $\sF$ constant one therefore has
$\h^\bullet(X_s,\sF)=\h^\bullet(U,\sF)$. Figuratively speaking, the
special fibre swallows the general fibre.

For schemes the proof is more delicate, and it is essential to suppose
that $\sF$ is torsion (\cite[XII.2]{sga4}). In view of the description of
the fibres of $\eR^qf_*\sF$ (\ref{I:2-3-3}), Theorem
(\ref{I:4-1-1}) is essentially equivalent to the following theorem.

\begin{theorem}\label{I:4-1-2}
Let $A$ be a strictly henselian local ring and let $S=\spec(A)$. Let
$f:X\to S$ be proper, and let $X_0$ be the closed fibre of $f$. Then, for every
torsion abelian sheaf $\sF$ on $X$ and every $q\geqslant0$, one has
$\h^q(X,\sF)\iso\h^q(X_0,\sF)$.
\end{theorem}

By passage to the limit, it suffices to prove the theorem when $A$ is the
strict henselization of a finitely generated $\dZ$-algebra at a prime
ideal. We first treat the case $q=0$ or $1$ and $\sF=\dZ/n$
(\ref{I:4-2}). An argument based on constructible sheaves (\ref{I:4-3})
shows moreover that it suffices to consider constant $\sF$. On the other
hand, the unwinding in (\ref{I:3-4-3}) allows us to suppose that $X_0$
is a curve; in that case it remains only to prove the theorem for $q=2$
(\ref{I:4-4}).

Among other applications (\ref{I:4-6}), the theorem allows one to define
cohomology with proper support (\ref{I:4-5}).

\subsection{Proof for \texorpdfstring{$q=0$}{q=0} or \texorpdfstring{$1$}{1} and \texorpdfstring{$\sF=\dZ/n$}{\sF=Z/n}}\label{I:4-2}

The result for $q=0$ and $\sF$ constant is equivalent to the following
proposition (Zariski's connectedness theorem).

\begin{proposition}\label{I:4-2-1}
Let $A$ be a noetherian henselian local ring and let $S=\spec(A)$. Let
$f:X\to S$ be a proper morphism, and let $X_0$ be the closed fibre of
$f$. Then the sets of connected components $\pi_0(X)$ and
$\pi_0(X_0)$ are in bijection.
\end{proposition}

Equivalently, we must show that the sets $\operatorname{Of}(X)$ and
$\operatorname{Of}(X_0)$ of open-and-closed subsets are in bijection.
The set $\operatorname{Of}(X)$ is known to correspond bijectively to the
set of idempotents of $\Gamma(X,\sO_X)$, and similarly
$\operatorname{Of}(X_0)$ corresponds bijectively to the set of
idempotents of $\Gamma(X_0,\sO_{X_0})$. Thus we must show that the
canonical map
\[
  \operatorname{Idem} \Gamma(X,\sO_X) \to \operatorname{Idem}\Gamma(X_0,\sO_{X_0})
\]
is bijective.

Let $\fm$ denote the maximal ideal of $A$, let
$\Gamma(X,\sO_X)^\wedge$ denote the completion of $\Gamma(X,\sO_X)$
for the $\fm$-adic topology, and, for every integer $n\geqslant0$, let
$X_n=X\otimes_AA/\fm^{n+1}$. By the finiteness theorem for proper
morphisms \cite[III.3.2]{ega}, $\Gamma(X,\sO_X)$ is a finite
$A$-algebra; since $A$ is henselian, the canonical map
\[
  \operatorname{Idem} \Gamma(X,\sO_X) \to \operatorname{Idem}\Gamma(X,\sO_X)^\wedge
\]
is bijective. In particular, the canonical map
\[
  \operatorname{Idem} \Gamma(X,\sO_X)^\wedge \to \varprojlim \operatorname{Idem}\Gamma(X_n,\sO_{X_n})
\]
is bijective. But since $X_n$ and $X_0$ have the same underlying
topological space, the canonical map
\[
  \operatorname{Idem}\Gamma(X_n,\sO_{X_n}) \to \operatorname{Idem} \Gamma(X_0,\sO_{X_0})
\]
is bijective for every $n$, which completes the proof.

Since $\h^1(X,\dZ/n)$ is in bijection with the set of isomorphism classes
of Galois étale coverings of $X$ with group $\dZ/n$, the theorem for
$q=1$ and $\sF=\dZ/n$ follows from the next proposition.

\begin{proposition}\label{I:4-2-2}
Let $A$ be a noetherian henselian local ring and let $S=\spec(A)$. Let
$f:X\to S$ be a proper morphism, and let $X_0$ be the closed fibre of $f$. Then
the restriction functor
\[
  \operatorname{Rev.et}(X) \to \operatorname{Rev.et}(X_0)
\]
is an equivalence of categories.
\end{proposition}

(If $X_0$ is connected and a geometric point of $X_0$ has been chosen
as base point, this says that the canonical map
$\pi_1(X_0)\to\pi_1(X)$ on (profinite) fundamental groups is bijective.)

Proposition (\ref{I:4-2-1}) shows that this functor is fully faithful.
Indeed, if $X'$ and $X''$ are two étale coverings of $X$, an
$X$-morphism from $X'$ to $X''$ is determined by its graph, which is an
open-and-closed subset of $X'\times_XX''$.

It remains to show that every étale covering $X_0'$ of $X_0$ extends to
an étale covering of $X$. Étale coverings are insensitive to nilpotent
elements \cite[ch.1]{sga1}; consequently, $X_0'$ lifts uniquely to an
étale covering $X_n'$ of $X_n$ for every $n\geqslant0$, in other words,
to an étale covering $\fX'$ of the formal scheme $\fX$ obtained by
completing $X$ along $X_0$. By Grothendieck's algebraization theorem for
formal coherent sheaves (the existence theorem, \cite[III.5]{ega}),
$\fX'$ is the formal completion of an étale covering $\bar X'$ of
$\bar X=X\otimes_A\hat A$.

By passage to the limit, it suffices to prove the proposition when $A$ is
the henselization of a finitely generated $\dZ$-algebra. We can then
apply Artin's approximation theorem to the functor
$\sF:(\textnormal{$A$-algebras})\to(\textnormal{sets})$ that associates
to an $A$-algebra $B$ the set of isomorphism classes of étale coverings
of $X\otimes_AB$. Indeed, this functor is locally of finite presentation:
if $B_i$ is a filtered direct system of $A$-algebras and
$B=\varinjlim B_i$, then $\sF(B)=\varinjlim\sF(B_i)$. By Artin's
theorem, given an element $\bar\xi\in\sF(\hat A)$, here the isomorphism class
of $\bar X'$, there exists $\xi\in\sF(A)$ having the same image as
$\bar\xi$ in $\sF(A/\fm)$. In other words, there exists an étale
covering $X'$ of $X$ whose restriction to $X_0$ is isomorphic to
$X_0'$.

\subsection{Constructible sheaves}\label{I:4-3}

In this paragraph, $X$ is a \emph{noetherian} scheme, and by a sheaf on
$X$ we mean an \emph{abelian} sheaf on $\et X$.

\begin{definition}\label{I:4-3-1}
A sheaf $\sF$ on $X$ is called \emph{locally constant constructible}
(abbreviated l.c.c.) if it is represented by an étale covering of $X$.
\end{definition}

\begin{definition}\label{I:4-3-2}
A sheaf $\sF$ on $X$ is called \emph{constructible} if it satisfies the
following equivalent conditions:
\begin{enumerate}[\indent (i)]
  \item There exists a finite surjective family of subschemes $X_i$ of
    $X$ such that the restriction of $\sF$ to $X_i$ is l.c.c.
  \item There exist a finite family of finite morphisms $p_i:X_i'\to X$,
    for each $i$ a constant constructible sheaf $C_i$ (i.e. one defined
    by a finite abelian group) on $X_i'$, and a monomorphism
    $\sF\to\prod p_{i*}C_i$.
\end{enumerate}
\end{definition}

It is readily checked that the category of constructible sheaves on $X$
is an abelian category. Moreover, if $u:\sF\to\sG$ is a homomorphism of
sheaves and $\sF$ is constructible, then $\im(u)$ is constructible.

\begin{lemma}\label{I:4-3-3}
Every torsion sheaf $\sF$ is a filtered direct limit of constructible
sheaves.
\end{lemma}

Indeed, if $j:U\to X$ is an étale scheme of finite type over $X$, an
element $\xi\in\sF(U)$ such that $n\xi=0$ defines a homomorphism of
sheaves $j_!\const{\dZ/n}\to\sF$ whose image (the smallest subsheaf of
$\sF$ for which $\xi$ is a local section) is a constructible subsheaf
of $\sF$. Clearly $\sF$ is the direct limit of such subsheaves.

\begin{definition}\label{I:4-3-4}
Let $\cC$ be an abelian category and let $T$ be a functor from $\cC$ to
the category of abelian groups. The functor $T$ is called
\emph{effaceable} in $\cC$ if, for every object $A$ of $\cC$ and every
$\alpha\in T(A)$, there exists a monomorphism $u:A\to M$ in $\cC$ such
that $T(u)\alpha=0$.
\end{definition}

\begin{lemma}\label{I:4-3-5}
The functors $\h^q(X,-)$ for $q>0$ are effaceable in the category of
constructible sheaves on $X$.
\end{lemma}

It suffices to observe that, if $\sF$ is a constructible sheaf, there is
necessarily an integer $n>0$ such that $\sF$ is a sheaf of
$\dZ/n$-modules. There is then a monomorphism $\sF\hookrightarrow\sG$,
where $\sG$ is a sheaf of $\dZ/n$-modules and
$\h^q(X,\sG)=0$ for every $q>0$. For example, one can take for $\sG$
the Godement resolution $\prod_{x\in X}i_{x*}\sF_{\bar x}$, where $x$
runs over the points of $X$ and $i_x:\bar x\to X$ is a geometric point
centred at $x$. By (\ref{I:4-3-3}), $\sG$ is a direct limit of
constructible sheaves, whence the lemma, because the functors
$\h^q(X,-)$ commute with direct limits.

\begin{lemma}\label{I:4-3-6}
Let $\varphi^\bullet:T^\bullet\to{T'}^\bullet$ be a morphism of
cohomological functors defined on an abelian category $\cC$ with values
in the category of abelian groups. Suppose that $T^q$ is effaceable for
$q>0$, and let $\cE$ be a subset of the objects of $\cC$ such that
every object of $\cC$ is contained in an object belonging to $\cE$.
Then the following conditions are equivalent:
\begin{enumerate}[\indent (i)]
  \item $\varphi^q(A)$ is bijective for every $q\geqslant0$ and every
    $A\in\ob\cC$.
  \item $\varphi^0(M)$ is bijective and $\varphi^q(M)$ is surjective for
    every $q>0$ and every $M\in\cE$.
  \item $\varphi^0(A)$ is bijective for every $A\in\ob\cC$, and
    ${T'}^q$ is effaceable for every $q>0$.
\end{enumerate}
\end{lemma}

The proof is by induction on $q$ and presents no difficulty.

\begin{proposition}\label{I:4-3-7}
Let $X_0$ be a subscheme of $X$. Suppose that, for every
$n>0$ and every scheme $X'$ finite over $X$, the canonical map
\[
  \h^q(X',\dZ/n) \to \h^q(X_0',\dZ/n)\text{,}
\]
where $X_0'=X'\times_XX_0$, is bijective for $q=0$ and surjective for
$q>0$. Then, for every torsion sheaf $\sF$ on $X$ and every
$q\geqslant0$, the canonical map
\[
  \h^q(X,\sF) \to \h^q(X_0,\sF)
\]
is bijective.
\end{proposition}

By passage to the limit, it suffices to prove the assertion for
constructible $\sF$. Apply Lemma (\ref{I:4-3-6}) with $\cC$ the category
of constructible sheaves on $X$, $T^q=\h^q(X,-)$,
${T'}^q=\h^q(X_0,-)$, and $\cE$ the set of constructible sheaves of the
form $\prod p_{i*}C_i$, where $p_i:X_i'\to X$ is finite and $C_i$ is a
finite constant sheaf on $X_i'$.

\subsection{End of the proof}\label{I:4-4}

By the method of fibration by curves (\ref{I:3-4-3}), we reduce to proving
the theorem in relative dimension $\leqslant1$. By the preceding
paragraph, it suffices to show the following. If $S$ is the spectrum of a
strictly henselian noetherian local ring, $f:X\to S$ is a proper
morphism whose closed fibre $X_0$ has dimension $\leqslant1$, and $n$
is an integer $>0$, then the canonical homomorphism
\[
  \h^q(X,\dZ/n) \to \h^q(X_0,\dZ/n)
\]
is bijective for $q=0$ and surjective for $q>0$.

The cases $q=0$ and $1$ were treated above, and
$\h^q(X_0,\dZ/n)=0$ for $q\geqslant3$; it therefore suffices to treat
$q=2$. We may clearly suppose that $n$ is a power of a prime. If
$n=p^r$, where $p$ is the characteristic of the residue field of $S$,
Artin--Schreier theory shows that $\h^2(X_0,\dZ/p^r)=0$. If
$n=\ell^r$, $\ell\ne p$, Kummer theory gives a commutative diagram
\[\xymatrix{
  \pic(X) \ar[d] \ar[r]^-\alpha
    & \h^2(X,\dZ/\ell^r) \ar[d] \\
  \pic(X_0) \ar[r]^-\beta
    & \h^2(X_0,\dZ/\ell^r)
}\]
in which $\beta$ is surjective (this was seen in Chapter \ref{I:3} for a
smooth curve over an algebraically closed field, but similar arguments
apply to any curve over a separably closed field).

Thus it suffices to prove:

\begin{proposition}\label{I:4-4-1}
Let $S$ be the spectrum of a henselian noetherian local ring, and let
$f:X\to S$ be a proper morphism whose closed fibre $X_0$ has dimension
$\leqslant1$. Then the canonical restriction map
\[
  \pic(X)\to \pic(X_0)
\]
is surjective. (In fact, it is enough that $f$ be a \emph{separated}
morphism of finite type.)
\end{proposition}

To simplify the proof, we suppose that $X$ is integral, although this is
not necessary. Every invertible sheaf on $X_0$ is associated with a
Cartier divisor (because $X_0$ is a curve, hence quasi-projective), so it
suffices to show that the canonical map $\Div(X)\to\Div(X_0)$ is
surjective.

Every divisor on $X_0$ is a linear combination of divisors whose support
is concentrated at a single nonisolated closed point of $X_0$. Let $x$
be such a point, let $t_0\in\sO_{X_0,x}$ be a regular nonunit in
$\sO_{X_0,x}$, and let $D_0$ be the divisor concentrated at $x$ with
local equation $t_0$. Let $U$ be an open neighbourhood of $x$ in $X$
on which there is a section $t\in\Gamma(U,\sO_U)$ lifting $t_0$. Let
$Y$ be the closed subset of $U$ defined by $t=0$; after making $U$
sufficiently small, we may suppose that $x$ is the only point of
$Y\cap X_0$. Then $Y$ is quasi-finite over $S$ at $x$. Since $S$ is
the spectrum of a henselian local ring, it follows that
$Y=Y_1\amalg Y_2$, where $Y_1$ is finite over $S$ and $Y_2$ does not
meet $X_0$. Moreover, since $X$ is separated over $S$, $Y_1$ is closed
in $X$.

After replacing $U$ by a smaller open neighbourhood of $x$, we may
suppose that $Y=Y_1$, in other words, that $Y$ is closed in $X$. We then
define a divisor $D$ on $X$ lifting $D_0$ by setting
$D|X\setminus Y=0$ and $D|U=\dv(t)$; this makes sense because $t$ is
invertible on $U\setminus Y$.

\subsubsection{Remark}\label{I:4-4-2}

When $f$ is proper, one could also give a proof similar to that of
Proposition (\ref{I:4-2-2}). Indeed, since $X_0$ is a curve, there is no
obstruction to lifting an invertible sheaf on $X_0$ to the infinitesimal
neighbourhoods $X_n$ of $X_0$, and hence to the formal completion $\fX$
of $X$ along $X_0$. One then concludes by applying successively
Grothendieck's existence theorem and Artin's approximation theorem.

\subsection{Cohomology with proper support}\label{I:4-5}

\begin{definition}\label{I:4-5-1}
Let $X$ be a separated scheme of finite type over a field $k$. By a
theorem of Nagata, there exist a scheme $\bar X$ proper over $k$ and an
open immersion $j:X\to\bar X$. For every torsion sheaf $\sF$ on $X$,
let $j_!\sF$ denote the extension by $0$ of $\sF$ to $\bar X$, and
define the groups of \emph{cohomology with proper support} by
\[
  \h_c^q(X,\sF)=\h^q(\bar X, j_! \sF) \text{.}
\]
\end{definition}

Let us show that this definition is independent of the chosen
compactification $j:X\to\bar X$. Let $j_1:X\to\bar X_1$ and
$j_2:X\to\bar X_2$ be two compactifications. Then $X$ maps into
$\bar X_1\times\bar X_2$ by $x\mapsto(j_1(x),j_2(x))$, and
the closed image $\bar X_3$ of $X$ under this map is a compactification
of $X$. We thus have a commutative diagram
\[\xymatrix{
  & \bar X_1 \\
  X \ar[ur]^-{j_1} \ar[rr]^-{(j_1,j_2)} \ar[dr]_-{j_2}
    & & \bar X_3 \ar[ul]_-{p_1} \ar[dl]^-{p_2} \\
  & \bar X_2
}\]
where $p_1$ and $p_2$, the restrictions to $\bar X_3$ of the natural
projections, are proper morphisms. It therefore suffices to treat the case
of a commutative diagram
\[\xymatrix{
  X \ar[r]^-{j_2} \ar[dr]_-{j_1}
    & \bar X_2 \ar[d]^-p \\
  & \bar X_1
}\]
with $p$ proper.

\begin{lemma}\label{I:4-5-2}
One has $p_*(j_{2!}\sF)=j_{1!}\sF$ and
$\eR^qp_*(j_{2!}\sF)=0$ for $q>0$.
\end{lemma}

Let us note at once that the lemma suffices. Using the Leray spectral
sequence for $p$, it gives, for every $q\geqslant0$,
\[
  \h^q(\bar X_2, j_{2!} \sF) = \h^q(\bar X_1, j_{1!} \sF) \text{.}
\]

To prove the lemma, argue fibre by fibre using the base change theorem
(\ref{I:4-1-1}) for $p$. The result is immediate because above a point
of $X$, $p$ is an isomorphism, whereas above a point of
$\bar X_1\setminus X$, $j_{2!}\sF$ vanishes on the fibre of $p$.

\subsubsection{}\label{I:4-5-3}

Similarly, if $f:X\to S$ is a separated morphism of finite type between
noetherian schemes, there exist a proper morphism $\bar f:\bar X\to S$
and an open immersion $j:X\to\bar X$. The higher direct images with
proper support $\eR^qf_!$ are defined, for every torsion sheaf $\sF$ on
$X$, by
\[
  \eR^q f_! \sF = \eR^q \bar f_*(j_! \sF) \text{.}
\]
As before, one checks that this definition is independent of the chosen
compactification.

\begin{theorem}\label{I:4-5-4}
Let $f:X\to S$ be a separated morphism of finite type between noetherian
schemes, and let $\sF$ be a torsion sheaf on $X$. Then the fibre of
$\eR^qf_!\sF$ at a geometric point $s$ of $S$ is isomorphic to the
cohomology with proper support $\h_c^q(X_s,\sF)$ of the fibre $X_s$ of
$f$ at $s$.
\end{theorem}

This is a simple variant of the proper base change theorem
(\ref{I:4-1-1}). More generally, if
\[\xymatrix{
  X \ar[d]^-f
    & X' \ar[l]_-{g'} \ar[d]^-{f'} \\
  S
    & S' \ar[l]^-g
}\]
is a cartesian diagram, there is a canonical isomorphism
\begin{equation*}\tag{4.5.4.1}\label{I:eq:4-5-4-1}
  g^* (\eR^q f_! \sF) \simeq \eR^q f'_!({g'}^* \sF) \text{.}
\end{equation*}

\subsection{Applications}\label{I:4-6}

\begin{theorem}[Vanishing theorem]\label{I:4-6-1}
Let $f:X\to S$ be a separated morphism of finite type whose fibres have
dimension $\leqslant n$, and let $\sF$ be a torsion sheaf on $X$. Then
$\eR^qf_!\sF=0$ for $q>2n$.
\end{theorem}

By the base change theorem, we may suppose that $S$ is the spectrum of a
separably closed field. If $\dim X=n$, there exists an affine open
$U$ of $X$ such that $\dim(X\setminus U)<n$. We then have an exact
sequence $0\to\sF_U\to\sF\to\sF_{X\setminus U}\to0$, and induction on
$n$ reduces us to proving the theorem for $X=U$ affine. The method of
fibration by curves (\ref{I:3-4-1}) and the base change theorem then
reduce us to a curve over a separably closed field, for which the desired
result follows from Tsen's theorem (\ref{I:3-3-6}).

\begin{theorem}[Finiteness theorem]\label{I:4-6-2}
Let $f:X\to S$ be a separated morphism of finite type and let $\sF$ be
a constructible sheaf on $X$. Then the sheaves $\eR^qf_!\sF$ are
constructible.
\end{theorem}

We consider only the case in which $\sF$ is annihilated by an integer
invertible on $X$.

To prove the theorem, one reduces to the case in which $\sF$ is the
constant sheaf $\const{\dZ/n}$ and $f:X\to S$ is proper and smooth with
geometrically connected curve fibres of genus $g$. If $n$ is invertible
on $X$, the sheaves $\eR^qf_*\sF$ are then locally free of finite rank
and vanish for $q>2$ (\ref{I:4-6-1}). Replacing $\dZ/n$ by the sheaf
$\dmu_n$, which is locally isomorphic to it (over $S$), there are
canonical isomorphisms
\begin{equation*}\tag{4.6.2.1}\label{I:eq:4-6-2-1}
\begin{aligned}
  \eR^0 f_* \dmu_n &= \dmu_n \\
  \eR^1 f_* \dmu_n &= \const\pic(X/S)_n \\
  \eR^2 f_* \dmu_n &= \dZ/n \text{.}
\end{aligned}
\end{equation*}

\begin{theorem}[Comparison with classical cohomology]\label{I:4-6-3}
Let $f:X\to S$ be a separated morphism of schemes of finite type over
$\dC$, and let $\sF$ be a torsion sheaf on $X$. Denote by a superscript
$\an{(-)}$ the functor passing to the usual topological spaces, and by
$\eR^q\an f_!$ the derived functors of direct image with proper support
under $\an f$. Then
\[
  \an{(\eR^q f_! \sF)} \simeq \eR^q \an f_! \an \sF \text{.}
\]
In particular, when $S=$ a point and $\sF$ is the constant sheaf
$\dZ/n$,
\[
  \h_c^q(X,\dZ/n) \simeq \h_c^q(\an X,\dZ/n)\text{.}
\]
\end{theorem}

Unwindings using the base change theorem reduce us to the case in which
$X$ is a smooth proper curve, $S=$ a point, and $\sF=\dZ/n$. The
cohomology groups in question then vanish for $q\ne0,1,2$, and we invoke
\cite{se55}: indeed, if $X$ is proper over $\dC$, then
$\pi_0(X)=\pi_0(\an X)$ and $\pi_1(X)=$ the profinite completion of
$\pi_1(\an X)$, which proves the assertion for $q=0,1$. For $q=2$,
one uses the Kummer exact sequence and the fact that, again by
\cite{se55}, $\pic(X)=\pic(\an X)$.

\begin{theorem}[Cohomological dimension of affine schemes]\label{I:4-6-4}
Let $X$ be an affine scheme of finite type over a separably closed field
and let $\sF$ be a torsion sheaf on $X$. Then
$\h^q(X,\sF)=0$ for $q>\dim(X)$.
\end{theorem}

For the very elegant proof, we refer to \cite{sga4}, XIV \S 2 and 3.

\subsubsection{Remark}\label{I:4-6-5}

This theorem is in a sense a substitute for Morse theory. Indeed, consider
the classical case in which $X$ is smooth and affine over $\dC$, embedded
in a standard affine space $\dC^N$. Then, for almost every point
$p\in\dC^N$, the function “distance from $p$” on $X$ is a Morse
function, and the indices of its critical points are smaller than
$\dim(X)$. Thus $X$ is obtained by attaching handles of index less than
$\dim(X)$, which gives the classical analogue of (\ref{I:4-6-4}).

\section{Local acyclicity of smooth morphisms}\label{I:5}

Let $X$ be a complex analytic variety and let $f:X\to D$ be a morphism
from $X$ to the disk. Write $[0,t]$ for the closed line segment in $D$
with endpoints $0$ and $t$, and $(0,t]$ for the half-open segment. If
$f$ is \emph{smooth}, the inclusion
\[
  j : f^{-1}\left((0,t]\right) \hookrightarrow f^{-1}\left([0,t]\right)
\]
is a homotopy equivalence; the special fibre
$X_0=f^{-1}(0)$ can be pushed into $f^{-1}([0,t])$.

In practice, for sufficiently small $t$, $f^{-1}((0,t])$ will be a
bundle over $(0,t]$, so that the inclusion
\[
  X_t = f^{-1}(t) \hookrightarrow f^{-1}\left((0,t]\right)
\]
will also be a homotopy equivalence. The \emph{cospecialization morphism}
is then the homotopy class of maps
\[
  \cosp : X_0 \hookrightarrow f^{-1}\left([0,t]\right) \xleftarrow\sim f^{-1}\left((0,t]\right) \xleftarrow\sim X_t \text{.}
\]
Figuratively speaking,
for a smooth morphism the general fibre swallows the special fibre.

No longer suppose that $f$ is necessarily smooth (but suppose that
$f^{-1}((0,t])$ is a bundle over $(0,t]$). A morphism
$\cosp^\bullet$ can still be defined in cohomology whenever
$j_*\dZ=\dZ$ and $\eR^qj_*\dZ=0$ for $q>0$. Under these hypotheses,
the Leray spectral sequence for $j$ shows that
\[
  \h^\bullet\left(f^{-1}\left([0,t]\right),\dZ\right) \iso \h^\bullet\left(f^{-1}\left((0,t]\right),\dZ\right)
\]
and $\cosp^\bullet$ is the composite morphism
\[
  \cosp^\bullet : \h^\bullet(X_t,\dZ)
    \xleftarrow\sim \h^\bullet\left(f^{-1}\left((0,t]\right),\dZ\right)
    \xleftarrow\sim \h^\bullet\left(f^{-1}\left([0,t]\right),\dZ\right)
    \to \h^\bullet(X_0,\dZ)
\]
The fibre of $\eR^qj_*\dZ$ at a point $x\in X_0$ is calculated as
follows. In an ambient space, take a ball $B_\varepsilon$ centred at $x$
with sufficiently small radius $\varepsilon$, and for sufficiently small
$\eta$ put $E=X\cap B_\varepsilon\cap f^{-1}(\eta t)$; this is the
\emph{vanishing-cycle variety} at $x$. One has
\[
  (\eR^q j_* \dZ)_x \xleftarrow\sim \h^q\left(X\cap B_\varepsilon\cap f^{-1}\left((0,\eta t]\right),\dZ\right) \xrightarrow\sim \h^q(E,\dZ)
\]
and the cospecialization morphism is defined in cohomology whenever the
vanishing-cycle varieties are acyclic
($\h^0(E,\dZ)=\dZ$ and $\h^q(E,\dZ)=0$ for $q>0$), which is expressed
by saying that $f$ is \emph{locally acyclic}.

This chapter is devoted to the analogue of this situation for a smooth
morphism of schemes in étale cohomology. In this setting, however, it is
essential to restrict to torsion coefficients of order \emph{prime to the
residue characteristics}. Paragraph \ref{I:5-1} gives generalities on
locally acyclic morphisms and cospecialization maps. In Paragraph
\ref{I:5-2}, we prove that a smooth morphism is locally acyclic. In
Paragraph \ref{I:5-3}, this result is combined with those of the preceding
chapter to obtain two applications: a specialization theorem for
cohomology groups (the cohomology of the geometric fibres of a smooth
proper morphism is locally constant), and a base change theorem for a
smooth morphism.

Throughout what follows, we fix an integer $n$, and “scheme” means
“scheme on which $n$ is invertible.” A “geometric point” will always mean
an “algebraic geometric point” (\ref{I:2-3-1})
$x:\spec(k)\to X$, with $k$ algebraically closed.

\subsection{Locally acyclic morphisms}\label{I:5-1}

\subsubsection{Notation}\label{I:5-1-1}

Given a scheme $S$ and a geometric point $s$ of $S$, we denote by
$\widetilde S^s$ the spectrum of the strict localization of $S$ at $s$.

\begin{definition}\label{I:5-1-2}
A geometric point $t$ of $S$ is called a \emph{generization} of $s$ if
it is defined by an algebraic closure of the residue field of a point of
$\widetilde S^s$. We also say that $s$ is a \emph{specialization} of
$t$, and call the \emph{specialization map} the $S$-morphism
$t\to\widetilde S^s$.
\end{definition}

\begin{definition}\label{I:5-1-3}
Let $f:X\to S$ be a morphism of schemes. Let $s$ be a geometric point
of $S$, $t$ a generization of $s$, $x$ a geometric point of $X$ over
$s$, and
$\widetilde X_t^x=\widetilde X^x\times_{\widetilde S^s}t$. Then
$\widetilde X_t^x$ is called a \emph{vanishing-cycle variety} of $f$ at
$x$.

The morphism $f$ is called \emph{locally acyclic} if the reduced
cohomology of every vanishing-cycle variety $\widetilde X_t^x$ vanishes:
\begin{equation*}\tag{5.1.3.1}\label{I:eq:5-1-3-1}
  \hat\h^\bullet(\widetilde X_t^x,\dZ/n) = 0\text{,}
\end{equation*}
i.e. $\h^0\left(\widetilde X_t^x,\dZ/n\right)=\dZ/n$ and
$\h^q(\widetilde X_t^x,\dZ/n)=0$ for $q>0$.
\end{definition}

\begin{lemma}\label{I:5-1-4}
Let $f:X\to S$ be a locally acyclic morphism and let $g:S'\to S$ be a
quasi-finite morphism (or a projective limit of quasi-finite morphisms).
Then the morphism $f':X'\to S'$ obtained from $f$ by base change is
locally acyclic.
\end{lemma}

Indeed, every vanishing-cycle variety of $f'$ is a vanishing-cycle
variety of $f$.

\begin{lemma}\label{I:5-1-5}
Let $f:X\to S$ be a locally acyclic morphism. For every geometric point
$t$ of $S$ giving rise to a cartesian diagram
\[\xymatrix{
  X_t \ar[r]^-{\varepsilon'} \ar[d]
    & X \ar[d]^-f \\
  t \ar[r]^-\varepsilon
    & S
}\]
one has $\varepsilon_*'\dZ/n=f^*\varepsilon_*\dZ/n$ and
$\eR^q\varepsilon_*'\dZ/n=0$ for $q>0$.
\end{lemma}

Let $\bar S$ be the closure of $\varepsilon(t)$, let $S'$ be the
normalization of $\bar S$ in $k(t)$, and consider the cartesian diagram
\[\xymatrix{
  X_t \ar[r]^-{i'} \ar[d]
    & X' \ar[r]^-{\alpha'} \ar[d]^-{f'}
    & X \ar[d]^-f \\
  t \ar[r]^-i
    & S' \ar[r]^-\alpha
    & S
}\]
The local rings of $S'$ are normal with separably closed fraction fields.
They are therefore strictly henselian, and the local acyclicity of $f'$
(\ref{I:5-1-4}) gives
$i'_*\dZ/n={f'}^*i_*\dZ/n$ and $\eR^q i'_*\dZ/n=0$ for $q>0$. Since $\alpha$
is integral, it follows that
\[
  \eR^q\varepsilon_*'\dZ/n
    = \alpha_*' \eR^q i_*' \dZ/n
    = \alpha_*' {f'}^* \eR^q i_* \dZ/n
    = f^* \alpha_* \eR^q i_* \dZ/n
    = f^* \eR^q \varepsilon_* \dZ/n \text{,}
\]
which proves the lemma.

\subsubsection{}\label{I:5-1-6}

Given a locally acyclic morphism $f:X\to S$ and a specialization map
$t\to\widetilde S^s$, we shall define canonical homomorphisms, called
cospecialization maps,
\[
  \cosp^\bullet : \h^\bullet(X_t,\dZ/n) \to \h^\bullet(X_s,\dZ/n)\text{,}
\]
relating the cohomology of the general fibre $X_t=X\times_St$ to that of
the special fibre $X_s=X\times_Ss$.

Consider the cartesian diagram
\[\xymatrix{
  X_t \ar[d] \ar[r]^-{\varepsilon'}
    & \widetilde X \ar[d]^-{f'}
    & \ar[l] X_s \ar[d] \\
  t \ar[r]^-\varepsilon
    & \widetilde S^s
    & \ar[l] s
}\]
obtained from $f$ by base change. By (\ref{I:5-1-4}), $f'$ is again
locally acyclic. The definition of local acyclicity immediately shows that
the restrictions to $X_s$ of the sheaves $\eR^q\varepsilon_*'\dZ/n$ are
$\dZ/n$ for $q=0$ and $0$ for $q>0$. By (\ref{I:5-1-5}), we even know
that $\eR^q\varepsilon_*'\dZ/n=0$ for $q>0$. We define
$\cosp^\bullet$ as the composite map
\begin{equation*}\tag{5.1.6.1}\label{I:eq:5-1-6-1}
  \h^\bullet(X_t,\dZ/n) \simeq \h^\bullet(\widetilde X,\varepsilon_*' \dZ/n) \to \h^\bullet(X_s,\dZ/n) \text{.}
\end{equation*}

\paragraph{Variant}
Let $\bar S$ be the closure of $\varepsilon(t)$ in
$\widetilde S^s$, let $S'$ be the normalization of $\bar S$ in $k(t)$,
and let $X'/S'$ be obtained from $X/S$ by base change. Diagram
\eqref{I:eq:5-1-6-1} may also be written
\[
  \h^\bullet(X_t,\dZ/n) \simeq \h^\bullet(X',\dZ/n) \to \h^\bullet(X_s,\dZ/n) \text{.}
\]

\begin{theorem}\label{I:5-1-7}
Let $S$ be a locally noetherian scheme, $s$ a geometric point of $S$,
and $f:X\to S$ a morphism. Suppose that
\begin{enumerate}[\indent a)]
  \item $f$ is locally acyclic,
  \item for every specialization map $t\to\widetilde S^s$ and every
    $q\geqslant0$, the cospecialization maps
    $\h^q(X_t,\dZ/n)\to\h^q(X_s,\dZ/n)$ are bijective.
\end{enumerate}
Then the canonical homomorphism
$(\eR^qf_*\dZ/n)_s\to\h^q(X_s,\dZ/n)$ is bijective for every
$q\geqslant0$.
\end{theorem}

To prove the theorem, we may clearly suppose that
$S=\widetilde S^s$. We shall in fact show that, for every sheaf of
$\dZ/n\dZ$-modules $\sF$ on $S$, the canonical homomorphism
$\varphi^q(\sF):(\eR^qf_*f^*\sF)_s\to\h^q(X_s,f^*\sF)$ is bijective.

Every sheaf of $\dZ/n\dZ$-modules is a filtered direct limit of
constructible sheaves of $\dZ/n\dZ$-modules (\ref{I:4-3-3}). Moreover,
every constructible sheaf of $\dZ/n\dZ$-modules embeds in a sheaf of the
form $\prod i_{\lambda *}C_\lambda$, where
$i_\lambda:t_\lambda\to S$ is a finite family of generizations of $s$
and $C_\lambda$ is a free $\dZ/n\dZ$-module of finite rank on
$t_\lambda$. By the definition of the cospecialization maps, condition
b) says that the homomorphisms $\varphi^q(\sF)$ are bijective when
$\sF$ has this form.

We conclude by a variant of Lemma (\ref{I:4-3-6}):

\begin{lemma}\label{I:5-1-8}
Let $\cC$ be an abelian category in which filtered direct limits exist.
Let $\varphi^\bullet:T^\bullet\to{T'}^\bullet$ be a morphism of
cohomological functors commuting with filtered direct limits, defined on
$\cC$ with values in the category of abelian groups. Suppose that there
are two subsets $\cD$ and $\cE$ of the objects of $\cC$ such that:
\begin{enumerate}[\indent a)]
  \item every object of $\cC$ is a filtered direct limit of objects
    belonging to $\cD$,
  \item every object belonging to $\cD$ is contained in an object
    belonging to $\cE$.
\end{enumerate}
Then the following conditions are equivalent:
\begin{enumerate}[\indent (i)]
  \item $\varphi^q(A)$ is bijective for every $q\geqslant0$ and every
    $A\in\ob\cC$.
  \item $\varphi^q(M)$ is bijective for every $q\geqslant0$ and every
    $M\in\cE$.
\end{enumerate}
\end{lemma}

The proof of the lemma proceeds by passage to the direct limit, induction
on $q$, and repeated application of the five lemma to the diagram of
exact cohomology sequences deduced from an exact sequence
$0\to A\to M\to A'\to0$, with $A\in\cD$, $M\in\cE$, and
$A'\in\ob\cC$.

\begin{corollary}\label{I:5-1-9}
Let $S$ be the spectrum of a strictly henselian noetherian local ring and
let $f:X\to S$ be a locally acyclic morphism. Suppose that for every
geometric point $t$ of $S$ one has $\h^0(X_t,\dZ/n)=\dZ/n$ and
$\h^q(X_t,\dZ/n)=0$ for $q>0$ (in other words, the geometric fibres of
$f$ are acyclic). Then $f_*\dZ/n=\dZ/n$ and
$\eR^qf_*\dZ/n=0$ for $q>0$.
\end{corollary}

\begin{corollary}\label{I:5-1-10}
Let $f:X\to Y$ and $g:Y\to Z$ be morphisms of locally noetherian
schemes. If $f$ and $g$ are locally acyclic, then so is $g\circ f$.
\end{corollary}

We may suppose that $X$, $Y$, and $Z$ are strictly local and that $f$
and $g$ are local morphisms. It is then enough to show that, if $z$ is
an algebraic geometric point of $Z$, then
$\h^0(X_z,\dZ/n)=\dZ/n$ and $\h^q(X_z,\dZ/n)=0$ for $q>0$.

Since $g$ is locally acyclic, one has $\h^0(Y_z,\dZ/n)=\dZ/n$ and
$\h^q(Y_z,\dZ/n)=0$ for $q>0$. Moreover, the morphism
$f_z:X_z\to Y_z$ is locally acyclic (\ref{I:5-1-4}), and its geometric
fibres are acyclic because they are vanishing-cycle varieties of $f$.
By (\ref{I:5-1-9}), $\eR^q{f_z}_*\dZ/n=0$ for $q>0$. In addition,
${f_z}_*\dZ/n$ is constant with fibre $\dZ/n$ on $Y_z$. The result
follows from the Leray spectral sequence for $f_z$.

\subsection{Local acyclicity of a smooth morphism}\label{I:5-2}

\begin{theorem}\label{I:5-2-1}
A smooth morphism is locally acyclic.
\end{theorem}

Let $f:X\to S$ be a smooth morphism. The assertion is local for the étale
topology on $X$ and $S$, so we may suppose that $X$ is standard affine
space of dimension $d$ over $S$. By passage to the limit, we may suppose
that $S$ is noetherian, and the transitivity of local acyclicity
(\ref{I:5-1-10}) shows that it suffices to treat $d=1$.

Let $s$ be a geometric point of $S$ and $x$ a geometric point of $X$
centred at a closed point of $X_s$. We must show that the geometric fibres
of $\widetilde X^x\to\widetilde S^s$ are acyclic. Henceforth put
$S=\widetilde S^s=\spec(A)$ and $X=\widetilde X^x$. Then
$X\simeq\spec A\{T\}$, where $A\{T\}$ is the henselization of $A[T]$
at the point $T=0$ above $s$.

If $t$ is a geometric point of $S$, the fibre $X_t$ is a projective
limit of smooth affine curves over $t$. Thus $\h^q(X_t,\dZ/n)=0$ for
$q\geqslant2$, and it suffices to show that
$\h^0(X_t,\dZ/n)=\dZ/n$ and $\h^1(X_t,\dZ/n)=0$ for $n$ prime to the
residue characteristic of $S$. This follows from the next two
propositions.

\begin{proposition}\label{I:5-2-2}
Let $A$ be a strictly henselian local ring, $S=\spec(A)$, and
$X=\spec A\{T\}$. Then the geometric fibres of $X\to S$ are connected.
\end{proposition}

By passage to the limit, we may reduce to the case in which $A$ is a
strict henselization of a finitely generated $\dZ$-algebra.

Let $\bar t$ be a geometric point of $S$ lying over $t$, and let $k'$
be a finite separable extension of $k(t)$ in $k(\bar t)$. Put
$t'=\spec(k')$ and $X_{t'}=X\times_S\spec(k')$. We must check that, for
all $\bar t$ and $t'$, $X_{t'}$ is connected (by which we mean connected
and nonempty). Let $A'$ be the normalization of $A$ in $k'$, i.e. the
ring of elements of $k'$ integral over the image of $A$ in $k(t)$. One
has $A\{T\}\otimes_AA'\iso A'\{T\}$: indeed, the left-hand side is
henselian local (because $A'$ is finite over $A$ and local) and is a
limit of local étale algebras over
$A'\{T\}=A\{T\}\otimes_AA'$. Thus $X_{t'}$ is still the fibre at $t'$
of $X=\spec\left(A'\{T\}\right)$ over $S'=\spec(A')$. The local scheme
$X'$ is normal, hence integral; its localization $X_{t'}$ is still
integral, and a fortiori connected.

\begin{proposition}\label{I:5-2-3}
Let $A$ be a strictly henselian local ring, $S=\spec(A)$, and
$X=\spec\left(A\{T\}\right)$. Let $\bar t$ be a geometric point of $S$
and $X_{\bar t}$ the corresponding geometric fibre. Then every Galois
étale covering of $X_{\bar t}$ of order prime to the characteristic of
the residue field of $A$ is trivial.
\end{proposition}

\begin{lemma}[Zariski--Nagata purity theorem in dimension $2$] \label{I:5-2-4} % originally 5.2.3.1
Let $C$ be a regular local ring of dimension $2$ and let $C'$ be a
finite normal $C$-algebra étale over the open complement of the closed
point of $\spec(C)$. Then $C'$ is étale over $C$.
\end{lemma}

Indeed, $C'$ is normal of dimension $2$, and hence
$\operatorname{prof}(C')=2$. Since
$\operatorname{prof}(C')+\dim\operatorname{proj}(C')=\dim(C)=2$, it
follows that $C'$ is free over $C$. The set of points of $C$ where $C'$
is ramified is then defined by an equation, the discriminant; since it
contains no point of height $1$, it is empty.

\begin{lemma}[Special case of Abhyankar's lemma]\label{I:5-2-5} % originally 5.2.3.2
Let $S=\spec(V)$ be a trait, let $\pi$ be a uniformizer, let $\eta$ be
the generic point of $S$, let $X$ be smooth over $S$, irreducible, and
of relative dimension $1$, let $\widetilde X_\eta$ be a Galois étale
covering of $X_\eta$ of degree $n$ invertible on $S$, and let
$S_1=\spec\left(V[\pi^{1/n}]\right)$. Use a subscript $(-)_1$ to denote
base change from $S$ to $S_1$. Then $\widetilde X_{1\eta}$ extends to an
étale covering of $X_1$.
\end{lemma}

Let $\widetilde X_1$ be the normalization of $X_1$ in
$\widetilde X_{1\eta}$. By the structure of the tame inertia groups of
the discrete valuation rings obtained by localizing $X$ at the generic
points of the special fibre $X_s$, $\widetilde X_1$ is étale over $X_1$
on the general fibre and at the generic points of the special fibre. By
(\ref{I:5-2-4}), it is étale everywhere.

\subsubsection{}\label{I:5-2-6} % originally 5.2.3.3

Let $t$ denote the point over which $\bar t$ lies. We are free to replace
$A$ by the normalization of $A$ in a finite separable extension $k(t')$ of
$k(t)$ in $k(\bar t)$ (cf. \ref{I:5-2-2}). This and a preliminary
passage to the limit allow us to suppose that
\begin{enumerate}[\indent a)]
  \item $A$ is normal and noetherian, and $t$ is the generic point of $S$.
  \item The étale covering under consideration of $X_{\bar t}$ comes from
    an étale covering of $X_t$.
  \item It comes from an étale covering $\widetilde X_U$ of the inverse
    image $X_U$ of a nonempty open subset $U$ of $S$ (this follows from
    b): $t$ is the limit of the $U$).
  \item The complement of $U$ has codimension $\geqslant2$ (at the price
    of enlarging $k(t)$, by applying (\ref{I:5-2-5}) to the discrete
    valuation rings obtained by localizing $S$ at the points of
    $S\setminus U$ of codimension $1$ in $S$; there are only finitely many
    such points).
  \item The covering $\widetilde X_U$ is trivial over the subscheme $T=0$.
    (This may require enlarging $k(t)$ once more.)
\end{enumerate}

Under these conditions, we shall show that the covering
$\widetilde X_U$ is trivial.

\begin{lemma}\label{I:5-2-7} % originally 5.2.3.4
Let $A$ be a normal noetherian strictly henselian local ring, let $U$
be an open subset of $\spec(A)$ whose complement has codimension
$\geqslant2$, let $V$ be its inverse image in
$X=\spec\left(A\{T\}\right)$, and let $V'$ be an étale covering of
$V$. If $V'$ is trivial over $T=0$, then $V'$ is trivial.
\end{lemma}

Let $B=\Gamma(V',\sO)$. Since $V'$ is the inverse image of $V$ in
$\spec(B)$, it suffices to show that $B$ is finite étale over $A\{T\}$
(and hence split, since $A\{T\}$ is strictly henselian). Let
$\widehat X=\spec A\llbracket T\rrbracket$, and denote by $(-)^\wedge$ base
change from $X$ to $\widehat X$. The scheme $\widehat X$ is faithfully
flat over $X$. Hence
$\Gamma(\widehat V',\sO)=B\otimes_{A\{T\}}A\llbracket T\rrbracket$,
and it suffices to show that this ring $\widehat B$ is finite étale over
$A\llbracket T\rrbracket$.

Let $V_m$ (resp. $V_m'$) be the subscheme of $\widehat V$ (resp.
$\widehat V'$) defined by $T^{m+1}=0$. By hypothesis, $V_0'$ is a
trivial covering of $V_0$: a sum of $n$ copies of $V_0$. The same is true
of $V_m'/V_m$, because étale coverings are insensitive to nilpotents.
We obtain
\[
  \varphi : \Gamma(\widehat V',\sO) \to \varprojlim_m \Gamma(V_m',\sO) = \left(\varprojlim_m \Gamma(V_m,\sO)\right)^n \text{.}
\]
By hypothesis, the complement of $U$ has depth $\geqslant2$: one has
$\Gamma(V_m,\sO)=A[T]/(T^{m+1})$, and $\varphi$ is a homomorphism from
$\widehat B$ to $A\llbracket T\rrbracket^n$. Over $U$, it provides $n$
distinct sections of $\widehat V'/\widehat V$; hence $\widehat V'$ is
trivial, a sum of $n$ copies of $\widehat V$. Since the complement of
$\widehat V$ in $\widehat X$ still has codimension $\geqslant2$ (and
therefore depth $\geqslant2$), it follows that
$\widehat B=A\llbracket T\rrbracket^n$, proving the lemma.

\subsection{Applications}\label{I:5-3}

\begin{theorem}[Specialization of cohomology groups]\label{I:5-3-1}
Let $f:X\to S$ be a proper locally acyclic morphism, for example a smooth
proper morphism. Then the sheaves $\eR^qf_*\dZ/n$ are locally constant
constructible, and for every specialization map $t\to\widetilde S^s$,
the cospecialization maps
$\h^q(X_t,\dZ/n)\to\h^q(X_s,\dZ/n)$ are bijective.
\end{theorem}

This follows immediately from the definition of the cospecialization maps
and the finiteness and proper base change theorems.

\begin{theorem}[Base change by a smooth morphism]\label{I:5-3-2}
Let
\[\xymatrix{
  X' \ar[r]^-{g'} \ar[d]^-{f'}
    & X \ar[d]^-f \\
  S' \ar[r]^-{g}
    & S
}\]
be a cartesian diagram with $g$ smooth. For every sheaf $\sF$ on $X$
that is torsion of order prime to the residue characteristics of $S$,
one has
\[
  g^* \eR^q f_* \sF \iso \eR^q f'_*({g'}^* \sF) \text{.}
\]
\end{theorem}

After passing to an open covering of $X$, we reduce to the case where
$X$ is affine, and then, by passage to the limit, to the case where $X$
is of finite type over $S$. The morphism $f$ then factors as an open
immersion $j:X\to\bar X$ followed by a proper morphism
$\bar f:\bar X\to S$. The Leray spectral sequence for $\bar f\circ j$
and the proper base change theorem show that it suffices to prove the
theorem when $X\to S$ is an open immersion.

In this case, if $\sF$ has the form $\varepsilon_*C$, where
$\varepsilon:t\to X$ is a geometric point of $X$, the theorem is a
corollary of (\ref{I:5-1-5}). The general case follows from Lemma
(\ref{I:5-1-8}).

\begin{corollary}\label{I:5-3-3}
Let $K/k$ be an extension of separably closed fields, let $X$ be a
$k$-scheme, and let $n$ be an integer prime to the characteristic of
$k$. Then the canonical map
$\h^q(X,\dZ/n)\to\h^q(X_K,\dZ/n)$ is bijective for every
$q\geqslant0$.
\end{corollary}

It suffices to observe that $\bar K$ is a direct limit of smooth
$\bar k$-algebras.

\begin{theorem}[Relative purity]\label{I:5-3-4}
Let
\begin{equation*}\tag{5.3.4.1}\label{I:eq:5-3-4-1}\xymatrix{
  U \ar@{^{(}->}[r]^-j \ar[dr]
    & X \ar[d]^-f
    & \ar@{_{(}->}[l]_-i Y \ar[dl]^-h \\
  & S
}\end{equation*}
be a commutative diagram with $f$ smooth of pure relative dimension $N$,
$h$ smooth of pure relative dimension $N-1$, $i$ a closed immersion,
and $U=X\setminus Y$. If $n$ is prime to the residue characteristics of
$S$, then
\begin{equation*}
\begin{aligned}
  j_* \dZ/n     &= \dZ/n \\
  \eR^1 j_*\dZ/n &= \dZ/n(-1)_Y \\
  \eR^q j_*\dZ/n &= 0 \qquad \text{for $q\geqslant 2$}
\end{aligned}
\end{equation*}
\end{theorem}

In these formulas, $\dZ/n(-1)$ denotes the $\dZ/n$-dual of $\dmu_n$.
If $t$ is a local equation for $Y$, the isomorphism
$\eR^1j_*\dZ/n\simeq\dZ/n(-1)_Y$ is defined by the map
$a:\dZ/n\to\eR^1j_*\dmu_n$ sending $1$ to the class of the
$\dmu_n$-torsor of $n$th roots of $t$.

The question is local. This allows us to replace $(X,Y)$ by a locally
isomorphic pair, for example
\[\xymatrix{
  \dA_T^1 \ar@{^{(}->}[r]^-j \ar[dr]_-g
    & \dP_T^1 \ar[d]^-f
    & \ar@{_{(}->}[l]_-i T \ar[dl] \\
  & T
}\]
with $T=\dA_S^{N-1}$ and $i=$ the section at infinity. Corollary
(\ref{I:5-1-9}) applies to $g$ and gives
$\eR^qg_*\dZ/n=\dZ/n$ for $q=0$ and $0$ for $q>0$. For $f$, on the
other hand, \eqref{I:eq:4-6-2-1} gives
\[
  \eR^q f_* \dZ/n =
  \begin{cases}
    \dZ/n & q=0,\\
    0 & q=1,\\
    \dZ/n(-1) & q=2,\\
    0 & q>2.
  \end{cases}
\]
It is readily checked that $j_*\dZ/n=\dZ/n$, and that the
$\eR^qj_*\dZ/n=0$ are concentrated on $i(T)$ for $q>0$. The Leray
spectral sequence
$E_2^{pq}=\eR^pf_*\eR^qj_*\dZ/n\Rightarrow\eR^{p+q}g_*\dZ/n$ therefore
reduces to
\[\xymatrix{
  i^* \eR^q j_* \dZ/n & 0 & \cdots \\
  i^* \eR^1 j_* \dZ/n \ar[drr]^-{d_2} & 0 & \cdots \\
  \dZ/n              & 0 & \dZ/n(-1) & 0 & \cdots
}\]
$\eR^qj_*\dZ/n=0$ for $q\geqslant2$, and $\eR^1j_*\dZ/n$ is the
extension by zero of a locally free sheaf of rank one on $T$ (isomorphic,
via $d_2$, to $\dZ/n(-1)$). Since the map $a$ defined above is injective
(as may be checked fibre by fibre), it is an isomorphism, which proves
(\ref{I:5-3-4}).

\subsubsection{}\label{I:5-3-5}

We refer to \cite[XVI.4,5]{sga4} for proofs of the following applications
of the acyclicity theorem (\ref{I:5-2-1}).

\begin{theorem}\label{I:5-3-6} % originally 5.3.5.1
Let $f:X\to S$ be a morphism of schemes of finite type over $\dC$ and
let $\sF$ be a constructible sheaf on $X$. Then
\[
  \an{(\eR^q f_*\sF)} \sim \eR^q \an f_* (\an\sF)
\]
(cf. \ref{I:4-6-3}: for ordinary cohomology it is necessary to assume
that $\sF$ is constructible, not merely torsion).
\end{theorem}

\begin{theorem}\label{I:5-3-7} % originally 5.3.5.2
Let $f:X\to S$ be a morphism of schemes of finite type over a field $k$
of characteristic $0$, and let $\sF$ be a constructible sheaf on $X$.
Then the $\eR^qf_*\sF$ are also constructible.
\end{theorem}

The proof uses resolution of singularities and (\ref{I:5-3-4}). It is
generalized to a morphism of finite type between excellent schemes of
characteristic $0$ in \cite[XIX.5]{sga4}. Another proof, independent of
resolution, is given in this volume (\nameref{VII}, \ref{VII:1-1})
% NOTE: originally (Th. finitude, 1.1),
and applies to a morphism of schemes of finite type over a field or over
a Dedekind ring.

\section{Poincaré duality}

\subsection{Introduction}\label{I:6-1}

Let $X$ be an oriented topological manifold of pure dimension $N$, and
suppose that $X$ admits a finite open covering
$\cU=(U_i)_{1\leqslant i\leqslant K}$ such that every nonempty
intersection of the open subsets $U_i$ is homeomorphic to a ball. For
such a manifold, the Poincaré duality theorem can be presented as follows:
\begin{enumerate}[\indent A]
  \item The cohomology of $X$ is the \v Cech cohomology associated with
    the covering $\cU$. It is the cohomology of the complex
    \begin{equation}\label{I:eq:6-1}
      0 \to \dZ^{A_0} \to \dZ^{A_1} \to \cdots
    \end{equation}
    where
    $A_k=\{(i_0,\dotsc,i_k):i_0<\cdots<i_k\text{ and }U_{i_0}\cap\cdots\cap U_{i_k}\ne\varnothing\}$.
  \item For $a\in A_k$, $a=(i_0,\dotsc,i_k)$, let
    $U_a=U_{i_0}\cap\cdots\cap U_{i_k}$ and let $j_a$ be the inclusion
    of $U_a$ into $X$. The constant sheaf $\dZ$ on $X$ admits the
    (left) resolution
    \begin{equation}\label{I:eq:6-2}
      \xymatrix{
        \cdots \ar[r]
          & \displaystyle\bigoplus_{a\in A_1} j_{a!}\dZ \ar[r]
          & \displaystyle\bigoplus_{a\in A_0} j_{a!} \dZ \ar[r] \ar[d]
          & 0 \\
        & & \dZ \text{.}
      }
    \end{equation}
    The compactly supported cohomology $\h_c^\bullet(X,j_{a!}\dZ)$ is
    precisely the cohomology with proper support of the (oriented) ball
    $U_a$:
      \[
        \h_c^i(X,j_{a!}\dZ) = \begin{cases}
                                0   & \text{if $i\ne N$} \\
                                \dZ & \text{if $i = N$}
                              \end{cases}
      \]
\end{enumerate}

The hypercohomology spectral sequence for the complex
\eqref{I:eq:6-2} and cohomology with proper support therefore says that
$\h_c^i(X)$ is the $(i-N)$th cohomology group of the complex
\begin{equation}\label{I:eq:6-3}
  \cdots \to \dZ^{A_1} \to \dZ^{A_0} \to 0 \text{.}
\end{equation}
This complex is dual to the complex \eqref{I:eq:6-1}, whence Poincaré
duality.

The essential points of this construction are
\begin{enumerate}[\indent a)]
  \item the existence of a theory of cohomology with proper support;
  \item the fact that every point $x$ of a manifold $X$ of pure dimension
    $N$ has a fundamental system of open neighbourhoods $U$ for which
    \begin{equation}\label{I:eq:6-4}
      \h_c^i(U) = \begin{cases}
                    0   & \text{for $i\ne N$,} \\
                    \dZ & \text{for $i=N$.}
                  \end{cases}
    \end{equation}
\end{enumerate}

Poincaré duality in étale cohomology can be constructed on this model.
For $X$ smooth of pure dimension $N$ over an algebraically closed field
$k$, $n$ invertible on $X$, and $x$ a closed point of $X$, the key
point is to calculate the projective limit over the étale neighbourhoods
$U$ of $x$:
\begin{equation}\label{I:eq:6-5}
  \varprojlim \h_c^i(U,\dZ/n)
    = \begin{cases}
        0     & \text{if $i\ne 2N$} \\
        \dZ/n & \text{if $i=2N$.}
      \end{cases}
\end{equation}

Just as for topological manifolds one must first treat directly the case
of an open ball (or even simply the interval $(0,1)$), here one must
first treat directly the case of curves (\ref{I:6-2}). The local
acyclicity theorem for smooth morphisms then reduces
\eqref{I:eq:6-5} to this special case (\ref{I:6-3}).

The isomorphisms \eqref{I:eq:6-4} and \eqref{I:eq:6-5} are not
canonical: they depend on the choice of an \emph{orientation} of $X$.
For $n$ invertible on a scheme $X$, $\dmu_n$ is a sheaf of free
$\dZ/n$-modules of rank one. We denote its $N$th tensor power by
$\dZ/n(N)$ ($N\in\dZ$). The intrinsic form of the second line of
\eqref{I:eq:6-5} is
\begin{equation}\label{I:eq:6-5'}
  \varprojlim \h_c^{2N} \left(U,\dZ/n(N)\right) = \dZ/n \text{,}
\end{equation}
and $\dZ/n(N)$ is called the \emph{orientation sheaf} of $X$. Since the
sheaf $\dZ/n(N)$ is constant and isomorphic to $\dZ/n$, it may be taken
outside the limit sign (with $N$ fixed), and one may instead write
\begin{equation}\label{I:eq:6-5''}
  \varprojlim \h_c^{2N} (U,\dZ/n) = \dZ/n(-N) \text{,}
\end{equation}
and Poincaré duality will take the form of a perfect duality, with values
in $\dZ/n(-N)$, between $\h^i(X,\dZ/n)$ and
$\h_c^{2N-i}(X,\dZ/n)$.

\subsection{The case of curves}\label{I:6-2}

\subsubsection{}\label{I:6-2-1}

Let $\bar X$ be a smooth projective curve over an algebraically closed
field $k$, and let $n$ be invertible on $\bar X$. The proof of
(\ref{I:3-3-5}) gives, for connected $\bar X$, a canonical isomorphism
\[\xymatrix{
  \h^2(\bar X,\dmu_n) = \pic(\bar X)/n\pic(\bar X) \ar[r]^-{\deg}_-{\sim}
    & \dZ/n
}\]
Let $D$ be a reduced divisor of $\bar X$ and let
$X=\bar X\setminus D$:
\[\xymatrix{
  X \ar@{^{(}->}[r]^-j
    & \bar X
    & \ar@{_{(}->}[l]_-i D \text{.}
}\]
The exact sequence
$0\to j_!\dmu_n\to\dmu_n\to i_*\dmu_n\to0$ and the fact that
$\h^i(\bar X,i_*\dmu_n)=\h^i(D,\dmu_n)=0$ for $i>0$ give an
isomorphism
\[\xymatrix{
  \h_c^2(X,\dmu_n) = \h^2(\bar X,j_! \dmu_n) \ar[r]^-\sim
    & \h^2(\bar X,\dmu_n) \ar[r]^-\sim
    & \dZ/n \text{.}
}\]
For disconnected $\bar X$, similarly
\[
  \h_c^2(X,\dmu_n) = (\dZ/n)^{\pi_0(X)}
\]
and the \emph{trace morphism} is defined to be the sum
\[\xymatrix{
  \tr : \h_c^2(X,\dmu_n) \simeq (\dZ/n)^{\pi_0(X)} \ar[r]^-\Sigma
    & \dZ/n\text{.}
}\]

\begin{theorem}\label{I:6-2-2}
The pairing $\tr(a\smallsmile b)$ identifies each of the two groups
$\h^i(X,\dZ/n)$ and $\h_c^{2-i}(X,\dmu_n)$ with the dual (with values
in $\dZ/n$) of the other.
\end{theorem}

\emph{Transcendental proof.} Suppose that $\bar X$ is a smooth projective
curve over a trait $S$ and that $j:X\hookrightarrow\bar X$ is the
inclusion of the complement of a divisor $D$ étale over $S$. The
cohomology (resp. cohomology with proper support) of the special and
generic geometric fibres of $X/S$ are “the same,” i.e. the fibres of
locally constant sheaves on $S$. This follows from the analogous facts
for $\bar X$ and $D$, by means of the exact sequence
$0\to j_!\dZ/n\to\dZ/n\to{\dZ/n}_D\to0$ (for cohomology with proper
support) and the formulas $j_*\dZ/n=\dZ/n$,
$\eR^1j_*\dZ/n\simeq{\dZ/n}_D(-1)$, and
$\eR^ij_*\dZ/n=0$ ($i\geqslant2$) (for ordinary cohomology)
(\ref{I:5-3-4}).

This specialization principle reduces (\ref{I:6-2-2}) to the case in
which $k$ has characteristic $0$. By (\ref{I:5-3-3}), this case reduces
to $k=\dC$. Finally, when $k=\dC$, the groups
$\h^\bullet(X,\dZ/n)$ and $\h_c^\bullet(X,\dmu_n)$ agree with the
groups of the same name calculated for the classical topological space
$X_{\textnormal{cl}}$, and, through the isomorphism
$\dZ/n\to\dmu_n:x\mapsto\exp\left(\frac{2\pi ix}{n}\right)$, the trace
morphism is identified with “integration over the fundamental class.”
Thus (\ref{I:6-2-2}) follows from Poincaré duality for
$X_{\textnormal{cl}}$.

\subsubsection{Algebraic proof}\label{I:6-2-3}

For a very economical proof, see (\nameref{V}, \S\ref{V:2}).
% originally (Duality \S 2)
% NOTE: should be hyperlinked V:2
Here is another, related to the autoduality of the Jacobian.

Return to the notation of (\ref{I:6-2-1}). We may suppose---and do
suppose---that $X$ is connected. The cases $i=0$ and $i=2$ being
trivial, we also suppose that $i=1$. Define $_D\dG_m$ by the exact
sequence $0\to_D\dG_m\to\dG_m\to i_*\dG_m\to0$ (sections of $\dG_m$
congruent to $1\mod D$). The group $\h^1(\bar X,_D\dG_m)$ classifies
invertible sheaves on $\bar X$ trivialized on $D$. It is the group of
points of $\pic_D(\bar X)$, an extension of $\dZ$ (the degree) by the
group of points of Rosenlicht's generalized Jacobian
$\pic_D^0(\bar X)$ (corresponding to conductor $1$ at each point of
$D$), itself an extension of the abelian variety $\pic^0(\bar X)$ by
the torus $\dG_m^D/(\text{diagonal $\dG_m$})$.

\begin{enumerate}[\indent a)]
  \item The exact sequence
    $0\to j_!\dmu_n\to_D\dG_m\xrightarrow{n}_D\dG_m\to0$ gives an
    isomorphism
    \begin{equation*}\tag{6.2.3.1}\label{I:eq:6-2-3-1}
      \h_c^1(X,\dmu_n) = \pic_D^0(\bar X)_n \text{.}
    \end{equation*}
  \item The map sending $x\in X(k)$ to the class of the invertible sheaf
    $\sO(x)$ on $\bar X$, trivialized by $1$ on $D$, comes from a
    morphism
    \[
      f:X \to \pic_D(\bar X) \text{.}
    \]
\end{enumerate}
For the remainder, fix a base point $0$ and put
$f_0(x)=f(x)-f(0)$. For every homomorphism
$v:\pic_D^0(\bar X)_n\to\dZ/n$, let
$\bar v\in\h^1\left(\pic_D^0(\bar X),\dZ/n\right)$ be the image under
$v$ of the class in
$\h^1\left(\pic_D^0(\bar X),\pic_D^0(\bar X)_n\right)$ of the torsor
defined by the extension
\[\xymatrix{
  0 \ar[r]
    & \pic_D^0(\bar X)_n \ar[r]
    & \pic_D^0(\bar X) \ar[r]^-n
    & \pic_D^0(\bar X) \ar[r]
    & 0
}\]
Geometric class field theory (as presented in Serre \cite{se59}) shows
that the map $v\mapsto f_0^*(\bar v)$
\begin{equation*}\tag{6.2.3.2}\label{I:eq:6-2-3-2}
  \hom\left(\pic_D^0(\bar X)_n,\dZ/n\right) \to \h^1(X,\dZ/n)
\end{equation*}
is an isomorphism. To deduce (\ref{I:6-2-2}) from
\eqref{I:eq:6-2-3-1} and \eqref{I:eq:6-2-3-2}, it remains to know that
\begin{equation*}\tag{2.3.3}\label{I:eq:6-2-3-3}
  \tr\left(u \smallsmile f_0^*(\bar v)\right) = -v(u) \text{.}
\end{equation*}

% originally (Duality, 3.2.4).

\subsection{The general case}\label{I:6-3}

Let $X$ be a smooth algebraic variety of pure dimension $N$ over an
algebraically closed field $k$. To state the Poincaré duality theorem,
one must first define the trace morphism
\[
  \tr : \h_c^{2N}\left(X,\dZ/n(N)\right) \to \dZ/n\text{.}
\]
The definition is a painful unwinding starting from the case of curves
\cite[XVIII \S 2]{sga4}. One then has:

\begin{theorem}\label{I:6-3-1}
The pairing $\tr(a\smallsmile b)$ identifies each of the groups
$\h_c^i\left(X,\dZ/n(N)\right)$ and $\h^{2N-i}(X,\dZ/n)$ with the dual
of the other.
\end{theorem}

Let $x\in X$ be a closed point and let $X_x$ be the strict localization
of $X$ at $x$. For $U$ ranging over the étale neighbourhoods of $x$,
put
\begin{equation*}\tag{6.3.1.1}\label{I:eq:6-3-1-1}
  \h_c^\bullet(X_x,\dZ/n) = \varprojlim \h_c^\bullet(U,\dZ/n) \text{.}
\end{equation*}
It would be preferable to consider instead the pro-object
$\text{“}\varprojlim\text{”}\h_c^\bullet(U,\dZ/n)$, but since the
groups involved are finite, the difference is inessential. As we tried
to explain in the introduction, (\ref{I:6-3-1}) follows from the facts
that
\begin{equation*}\tag{6.3.1.2}\label{I:eq:6-3-1-2}
\begin{aligned}
  \h_c^i(X_x,\dZ/n) &= 0 \text{ for $i\ne 2N$ and that } \\
  \tr : \h_c^{2 N}\left(X_x,  \dZ/n(N)\right) &\iso \dZ/n \text{ is an isomorphism.}
\end{aligned}
\end{equation*}

The case $N=0$ is trivial. If $N>0$, let $Y_y$ be the strict
localization at a closed point of a smooth scheme $Y$ of pure dimension
$N-1$, and let $f:X_x\to Y_y$ be an essentially smooth morphism (of
relative dimension one). The proof uses the Leray spectral sequence for
cohomology with proper support for $f$ to reduce to the case of curves.
Because the “cohomology with proper support” under consideration is
defined as limits \eqref{I:eq:6-3-1-1}, the existence of such a spectral
sequence raises various passage-to-the-limit problems, treated in too much
detail in \cite[XVIII]{sga4}. For every geometric point $z$ of $Y_y$,
one has
\[
  (\eR^if_! \dZ/n)_z = \h_c^i(f^{-1}(z),\dZ/n) \text{.}
\]
The geometric fibre $f^{-1}(z)$ is a projective limit of smooth curves
over an algebraically closed field. It satisfies Poincaré duality. Its
ordinary cohomology is given by the local acyclicity theorem for smooth
morphisms:
\[
  \h^i(f^{-1}(z),\dZ/n)
    = \begin{cases}
        \dZ/n & \text{for $i=0$} \\
        0     & \text{for $i>0$.}
      \end{cases}
\]
By duality,
\[
  \h_c^i(f^{-1}(z),\dZ/n)
    = \begin{cases}
        \dZ/n(-1) & \text{for $i=2$} \\
        0         & \text{for $i\ne 2$.}
      \end{cases}
\]
and the Leray spectral sequence becomes
\[
  \h_c^i\left(X_x,\dZ/n(N)\right) = \h_c^{i-2}\left(Y_y,\dZ/n(N-1)\right) \text{.}
\]
The result follows by induction on $N$.

\subsection{Variants and applications}\label{I:6-4}

One can construct in étale cohomology a “duality formalism”
($=$ functors $\eR f_*$, $\eR f_!$, $f^*$, and $\eR f^!$ satisfying various
compatibilities and adjunction formulas) parallel to the one that exists
in coherent cohomology. In this language, the results of the preceding
paragraph can be rewritten as follows: if $f:X\to S$ is smooth of pure
relative dimension $N$ and $S=\spec(k)$ with $k$ algebraically closed,
then
\[
  \eR f^! \dZ/n = \dZ/n [2N](N) \text{.}
\]
This result holds without any hypothesis on $S$. It has the following
corollary.

\begin{corollary}\label{I:6-4-1}
If $f$ is smooth of pure relative dimension $N$ and the sheaves
$\eR^if_!\dZ/n$ are locally constant, then the sheaves
$\eR^if_*\dZ/n$ are also locally constant, and
\[
  \eR^i f_* \dZ/n = \const\hom\left(\eR^{2 N-i} f_! \dZ/n(N),\dZ/n\right)\text{.}
\]
\end{corollary}

In particular, under the hypotheses of the corollary, the sheaves
$\eR^if_*\dZ/n$ are constructible. From this one can show that, if $S$
is of finite type over the spectrum of a field or of a Dedekind ring,
then, for every morphism of finite type $f:X\to S$ and every constructible
sheaf $\sF$ on $X$, the sheaves $\eR^if_*\sF$ are constructible
(\nameref{VII}, \ref{VII:1-1}).
% originally (Th. finitude, 1.1).

% END INCLUDED components/sga4half_1_en.tex

% BEGIN INCLUDED components/sga4half_2_en.tex
% !TEX root = ../SGA4half_English_Reader_1_to_8.tex

\chapter{Report on the trace formula}\label{II}

In this text, I have attempted to present Grothendieck's cohomological theory
of $L$-functions as directly as possible. I follow very closely some of the
lectures given by Grothendieck at the IHES in the spring of 1966. In the
spirit of this volume, no appeal will be made to SGA 5---apart from two
references to passages in Exposé XI, which is otherwise independent of that
seminar.

In \S 10 (pp. 81--90) of \cite{de73}, the reader will find a summary of the
theory in the crucial case of curves.










\section{Some reminders}\label{II:1}





\subsection{Notation}\label{II:1-1}

$p,q,\dF,\dF_q$: $p$ is a prime number, $q=p^f$ a power of $p$, and $\dF$
an algebraic closure of the prime field $\dF_p$. For every power $p^n$ of
$p$, denote by $\dF_{p^n}$ the subfield with $p^n$ elements of $\dF$.

$X_0,X$: An object over $\dF_q$ will often be denoted by a symbol carrying a
subscript $0$. The same symbol without the $0$ then denotes the object
obtained from it by extension of scalars to $\dF$. For example, if $\sF_0$ is
a sheaf on a scheme $X_0$ over $\dF_q$, denote by $\sF$ its inverse image on
$X=X_0\otimes_{\dF_q} \dF$.

$F$: If $X_0$ is a scheme over $\dF_q$, the endomorphism of $X_0$ that
``sends the point with coordinates $x$ to the point with coordinates $x^q$''
is called the \emph{Frobenius morphism} and denoted by $F$: it is the identity
on the underlying topological space and, for $x$ a local section of the
structure sheaf, $F^* x = x^q$.

The endomorphism of $X$ obtained from it by extension of scalars is also
denoted by $F$. On the set $X(\dF) = X_0(\dF)$ of rational points of $X$,
$F$ acts as the Frobenius substitution
$\varphi\in\gal(\dF/\dF_q)$ defined by $\varphi(x)=x^q$. In particular,
$X^F=X_0(\dF_q)$.

In case of ambiguity, we shall write $F_{(q)}$ rather than $F$.










\subsection{Frobenius correspondence}\label{II:1-2}

Let $X_0$ be a scheme over $\dF_q$, and let $\sF_0$ be a sheaf on $X_0$. I
like to view the \emph{Frobenius correspondence} (an $F$-endomorphism of
$\sF_0$) as follows:
\begin{equation*}\tag{1.2.1}\label{II:eq:1-2-1}
  F^* : F^*\sF_0 \iso \sF_0
\end{equation*}

Regard $\sF_0$ as the sheaf of local sections of a space $[\sF_0]$ étale over
$X_0$ ($[\sF_0]$ is an algebraic space in the sense of M. Artin). The
functoriality of $F$ gives a commutative diagram
\[\xymatrix{
  [\sF_0] \ar[r]^-F \ar[d]
    & [\sF_0] \ar[d] \\
  X_0 \ar[r]^F
    & X_0
}\]
Because $[\sF_0]$ is étale over $X_0$, this diagram is Cartesian; it gives an
isomorphism $[\sF_0]\iso F^*[\sF_0]$ whose inverse is
\eqref{II:eq:1-2-1}. We also denote by $F^*$ the correspondences or morphisms
deduced from $F$ by extension of scalars or by functoriality, such as
$F^*:\h_c^\bullet(X,\sF)\to \h_c^\bullet(X,\sF)$.

For a formal definition, I refer to \cite[XI.1,2]{sga5}. Only the case
$q=p$ is considered there, but the general case is treated in the same way.






\subsection{}\label{II:1-3}

The Frobenius correspondence is functorial in $(X_0,\sF_0)$. If
$f_0:X_0\to Y_0$ is a morphism of $\dF_q$-schemes, $\sF_0$ a sheaf on $X_0$,
$\sG_0$ a sheaf on $Y_0$, and
$u\in\hom(f_0^*\sG_0,\sF_0)$ an $f_0$-morphism from $\sG_0$ to $\sF_0$, there
is a commutative diagram of spaces $f_0 F = F f_0$, and above it a
commutative diagram of sheaves
\[\xymatrix{
  X_0 \ar[r]^-{f_0} \ar[d]^-F
    & Y_0 \ar[d]^-F \\
  X_0 \ar[r]^-{f_0}
    & Y_0
}\qquad\qquad
\xymatrix{
  \sF_0
    & \ar[l]_-u \sG_0 \\
  \sF_0 \ar[u]_-{F^*}
    & \ar[l] \sG_0 \ar[u]^-{F^*}
}\]
In particular, the correspondence on $f_{0*}\sF_0$ deduced from the
Frobenius correspondence of $(X_0,\sF_0)$ is the Frobenius correspondence of
$(Y_0,f_{0*}\sF_0)$.






\subsection{}\label{II:1-4}

The natural morphism $Y\to Y_0$ is a limit of étale morphisms. It follows
that for every abelian sheaf $\sF_0$ on $X_0$, the inverse images on $Y$ of
the $\eR^if_0\sF_0$ are the $\eR^if_*\sF$.

On $\eR^if_*\sF$, one has
\begin{enumerate}[\indent a)]
  \item the correspondence
    $F^*\eR^i f_*\sF \iso \eR^i f_*\sF \xrightarrow{F^*} \eR^if_*\sF$
    deduced by functoriality from the Frobenius correspondence;
  \item the correspondence $F^*\eR^if_*\sF\to\eR^if_*\sF$ that is the
    inverse image on $Y$ of the Frobenius correspondence on
    $\eR^i f_{0*}\sF_0$.
\end{enumerate}

It follows from (\ref{II:1-3}) that these correspondences coincide. They
induce the initial term of a Frobenius action on the entire spectral sequence
\[
  \h^p(Y,\eR^q f_*\sF) \rightarrow \h^{p+q}(X,\sF) \text{.}
\]
Likewise for cohomology with proper support, and for the spectral sequence
\[
  \h_c^p(Y,\eR^q f_* \sF) \rightarrow \h_c^{p+q}(X,\sF)
\]





\subsection{Change of finite field}\label{II:1-5}

Let $(X_0,\sF_0)$ be over $\dF_q$, let $n$ be an integer, and let
$(X_1,\sF_1)$ over $\dF_{q^n}$ be obtained from $(X_0,\sF_0)$ by extension of
scalars. On the one hand there is the Frobenius correspondence
$F_{(q)}:X_0\to X_0)$, $F_{(q)}^*:F_{(q)}^*\sF_0\to \sF_0$, and on the other
there is $F_{(q^n)}:X_1\to X_1$,
$F_{(q^n)}^*:F_{(q^n)}^*\sF_1\to\sF_1$. It is easily verified that
$F_{(q^n)}$ is obtained from the $n$-th power of $F_{(q)}$ by extension of
scalars from $\dF_q$ to $\dF_{q^n}$; after extension of scalars to $\dF$, one
still has the same identity $F_{(q^n)}=F_{(q)}^n$ between endomorphisms of
$X$ and correspondences on $X$.

If $x$ is a closed point of $X_0$, with $[k(x),\dF_q]=n$, and
$\bar x\in X_0(\dF) = X(\dF)$ lies over $x$, then $\bar x$ is a fixed point
of $F_{(q^n)}=F_{(q)}^n$. Denote by $F_x^*$ the endomorphism of
$\sF_{\bar x}$ induced by $F_{(q^n)}^*$. Up to isomorphism,
$(F_{\bar x}^*,\sF_{\bar x})$ does not depend on the choice of $\bar x$.
The trace, etc. of $F_x^*$ is denoted by notation such as
$\tr(F_x^*,\sF)$.






\subsection{\texorpdfstring{$L$}{L}-function}\label{II:1-6}

Let $X_0$ be a scheme of finite type over $\dF_q$. Denote by $|X_0|$ the set
of its closed points, and for $x\in |X_0|$ set
$\deg(x)=[k(x),\dF_p]$. If $\sF_0$ is a $\dQ_\ell$-sheaf on $X_0$ (for this
notion, see below)---or a constructible flat sheaf of modules over a
commutative Noetherian ring $A$---denote by $L(X_0,\sF_0)$ the formal series
with constant term $1$, in $\dQ_\ell\llbracket t\rrbracket$ or in
$A\llbracket t\rrbracket$,
\[
  L(X_0,\sF_0) = \prod_{x\in|X_0|} \det\left(1-F_x^* t^{\deg(x)},\sF\right)^{-1} \text{.}
\]






\subsection{Remark}\label{II:1-7}

In definition (\ref{II:1-6}), one has $\deg(x)=[k(x),\dF_p]$ (absolute
degree), and not $\deg(x)=[k(x),\dF_q]$ (degree relative to $\dF_q$). The
effect is that $L(X_0,\sF_0)$ depends only on the scheme $X_0$ and the sheaf
$\sF_0$, not on $q$ and the structure morphism
$X_0\to\spec(\dF_q)$. The indeterminate $t$ is an avatar of $p^{-s}$.






\subsection{The Galois point of view}\label{II:1-8}

This subsection (n$^\circ$) will not be used in the remainder of the report. For $\sF_0$
an abelian sheaf on $X_0$, the Galois group $\gal(\dF/\dF_q)$ acts on
$\dF$, and by transport of structure on $\h_c^i(X,\sF)$. In particular, the
Frobenius substitution $\varphi$ (\ref{II:1-1}) induces an automorphism, also
denoted by $\varphi$, of $\h_c^i(X,\sF)$. The Frobenius correspondence, on the
other hand, defines an endomorphism $F^*$ of $\h_c^i(X,\sF)$. One has
\begin{equation*}\tag{1.8.1}\label{II:eq:1-8-1}
  {F^*}^{-1} = \varphi \qquad \text{(on $\h_c^i(X,\sF)$).}
\end{equation*}

To see this, apply (\ref{II:1-4}) to the case where $Y_0$ is a point:
$Y_0=\spec(\dF_q)$. One finds that $\h_c^i(X,\sF)$ is the fiber, at the
geometric point $\spec(\dF)$ of $Y_0$, of $\eR^i f_! \sF_0$, and that $F^*$
is the fiber of the Frobenius correspondence of $\eR^i f_!\sF_0$. Let
$[\eR^i f_! \sF_0]$ be as in (\ref{II:1-2}). One has
$\h_c^i(X,\sF) = [\eR^i f_!\sF_0](\dF)$ and
\begin{enumerate}[\indent a)]
  \item $\varphi$ acts by transport of structure, via its action on $\dF$;
  \item $F^*$ is the inverse of
    $F:[\eR^i f_!\sF_0](\dF) \to [\eR^i f_!\sF_0](\dF)$.
\end{enumerate}

The conclusion follows from the identity $F=\varphi$ of (\ref{II:1-1}).

One advantage of the Galois point of view is that it permits reasoning by
transport of structure.










\section{\texorpdfstring{$\dQ_\ell$}{Ql}-sheaves}\label{II:2}

The notion of $\dQ_\ell$-sheaves is developed in detail in
\cite[V,VI]{sga5}. We shall use only the definitions and results below.






\begin{definition_}\label{II:2-1}
Let $X$ be a Noetherian scheme. A $\dZ_\ell$-sheaf $\sF$ on $X$ is an
inverse system of sheaves $\sF_n$ ($n\geqslant 0$), with $\sF_n$ a
constructible sheaf of $\dZ/\ell^{n+1}$-modules \cite[IX.2]{sga4}, such that
the transition morphism $\sF_n\to\sF_{n-1}$ factors through an isomorphism
\[
  \sF_n\otimes_{\dZ/\ell^{n+1}} \dZ/\ell^n \iso \sF_{n-1}\text{.}
\]
The sheaf $\sF$ is said to be \emph{lisse} if the $\sF_n$ are locally
constant.
\end{definition_}

The terminology of \cite{sga5} is different: ``constructible
$\dZ_\ell$-sheaf'' is used there for ``$\dZ_\ell$-sheaf,'' and
``constructibly twisted constant'' for ``lisse.''

Likewise below, for $\dQ_\ell$-sheaves.






\subsection{}\label{II:2-2}

If $k$ is a separably closed field, a sheaf of $\dZ/\ell^n$-modules on
$\spec(k)$ is identified with a $\dZ/\ell^n$-module, and the functor
$\sF\mapsto \varprojlim \sF_n$ identifies the $\dZ_\ell$-sheaves on
$\spec(k)$ with the $\dZ_\ell$-modules of finite type. This justifies the
following definition: the \emph{fiber} of the $\dZ_\ell$-sheaf $\sF$ at the
geometric point $x$ of $X$ is the $\dZ_\ell$-module
$\varprojlim (\sF_n)_x$.






\subsection{}\label{II:2-3}

Suppose $X$ is connected, and let $x$ be a geometric point of $X$. It is
known that the ``fiber at $x$'' functor is an equivalence between the
category of constructible locally constant sheaves of
$\dZ/\ell^n$-modules and the category of finite-type
$\dZ/\ell^n$-modules equipped with a continuous action of $\pi_1(X,x)$.
Passing to the limit gives the following proposition.






\begin{proposition_}\label{II:2-4}
Under the hypotheses above, the ``fiber at $x$'' functor is an equivalence
between the category of lisse $\dZ_\ell$-sheaves on $X$ and that of
finite-type $\dZ_\ell$-modules equipped with a continuous action of
$\pi_1(X,x)$.
\end{proposition_}






\begin{proposition_}\label{II:2-5}
Let $\sF$ be a $\dZ_\ell$-sheaf on a Noetherian scheme $X$. There is a finite
partition of $X$ into locally closed parts $X_i$ such that the $\sF|X_i$ are
lisse.
\end{proposition_}

Set $\gr_\ell^n\sF=\ell^n\sF_m/\ell^{m+1}\sF_m$ for $m\geqslant n$, and
$\gr_\ell\sF=\bigoplus \gr_\ell^n\sF$. It is a sheaf of modules over the
Noetherian ring $\gr_\ell\dZ=\dZ/\ell[t]$, the graded ring of $\dZ$ for the
$\ell$-adic filtration. It is a quotient of
$\gr_\ell^0\sF \otimes_{\dZ/\ell} \dZ/\ell[t]$, hence a constructible sheaf
of $\gr_\ell\dZ$-modules \cite[IX.2]{sga4}; and there is a finite partition
of $X$ into locally closed parts $X_i$ such that
$\gr_\ell \sF|X_i$ is locally constant. Each of the sheaves
$\gr_\ell^n\sF|X_i$ is then locally constant, as are the $\sF_m|X_i$, which
are successive extensions of them.






\subsection{}\label{II:2-6}

If $\sF$ and $\sG$ are two $\dZ_\ell$-sheaves on $X$, set
$\hom(\sF,\sG) = \varprojlim(\sF_n,\sG_n)$; this is a $\dZ_\ell$-module.
Since $\hom(\sF_n,\sG_n) = \hom(\sF_m,\sG_n)$ for $m\geqslant n$, one also
has
\[
  \hom(\sF,\sG) = \varprojlim_n \varinjlim_m \hom(\sF_m,\sG_n) \text{,}
\]
so that the functor $\sF\mapsto \text{``}\varprojlim\text{''}\sF_n$ is fully
faithful from the category of $\dZ_\ell$-sheaves to that of pro-sheaves on
$X$ ($=$ pro-objects of the category of sheaves on $X$). We shall use it to
identify $\dZ_\ell$-sheaves with certain pro-sheaves (those such that every
$\sF_n=\sF/\ell^{n+1}\sF$ is a sheaf, and
$\sF\iso \text{``}\varprojlim\text{''}\sF_n$). It is easily verified that if
$\sF+\text{``}\varprojlim\text{''}\sF_n$, with $\sF_n$ a constructible sheaf
of $\dZ/\ell^n$-modules, then $\sF$ is a $\dZ_\ell$-sheaf if and only if the
inverse systems (in $n$) $\sF_n/\ell^{k+1}\sF_n$ are essentially constant:
the inverse system of
$\sF_k'=\varprojlim_n\sF_n/\ell^{k+1}\sF_n$ then satisfies the conditions of
(\ref{II:2-1}), and
$\sF=\text{``}\varprojlim\text{''}\sF_k'$.

In what follows, we abandon the notation $\sF_n$ of (\ref{II:2-1}); if
$\sF$ is a $\dZ_\ell$-sheaf, the sheaf $\sF_n$ of (\ref{II:2-1}) will be
denoted by $\sF\otimes\dZ/\ell^{n+1}$.






\subsection{}\label{II:2-7}

In the abelian category of pro-sheaves, the subcategory of
$\dZ_\ell$-sheaves is stable under kernels, cokernels, and extensions. This is
clear for cokernels. For kernels, if $f:\sF\to\sG$ is a morphism of
$\dZ_\ell$-sheaves, it follows readily from
(\ref{II:2-4}, \ref{II:2-5}) and Artin--Rees that the inverse system of
$\ker(f:\sF\otimes\dZ/\ell^n\to\sG\otimes\dZ/\ell^n)$ satisfies the
criterion given in (\ref{II:2-6}) above. There is even an integer $r$ such
that for every $n$
\[
  \ker(f)\otimes\dZ/\ell^n = \ker(f:\sF\otimes\dZ/\ell^{n+r}\to\sG\otimes\dZ/\ell^{n+r})\otimes\dZ/\ell^n \text{.}
\]
The case of extensions is left to the reader.






\subsection{}\label{II:2-8}

A $\dZ_\ell$-sheaf $\sF$ is \emph{torsion-free} if
$\ker(\ell:\sF\to\sF) = 0$. Each $\sF\otimes\dZ/\ell^n$ is then flat over
$\dZ/\ell^n$ (i.e. its fibers are finite-rank free
$\dZ/\ell^n$-modules). It follows from (\ref{II:2-4}, \ref{II:2-5}) that for
every $\dZ_\ell$-sheaf $\sF$ there is an integer $n$ such that
$\sF/\ker(\ell^n)$ is torsion-free.






\subsection{}\label{II:2-9}

The abelian category of \emph{$\dQ_\ell$-sheaves} on $X$ is the quotient of
that of $\dZ_\ell$-sheaves by the subcategory of torsion
$\dZ_\ell$-sheaves. Equivalently, its objects are the $\dZ_\ell$-sheaves,
and, denoting by $\sF\otimes\dQ_\ell$ the $\dZ_\ell$-sheaf $\sF$, viewed as
an object of the category of $\dQ_\ell$-sheaves, one has
\[
  \hom(\sF\otimes\dQ_\ell,\sG\otimes\dQ_\ell) = \hom(\sF,\sG)\otimes_{\dZ_\ell}\dQ_\ell \text{.}
\]

The \emph{fiber} of a $\dQ_\ell$-sheaf $\sF\otimes\dQ_\ell$ at a geometric
point $x$ of $X$ is the finite-dimensional $\dQ_\ell$-vector space
$\sF_x\otimes_{\dZ_\ell}\dQ_\ell$.






\subsection{}\label{II:2-10}

Let $X$ be a separated scheme of finite type over an algebraically closed
field $k$, and let $\sF$ be a constructible $\dQ_\ell$-sheaf on $X$. Let
$\sF'$ be a constructible $\dZ_\ell$-sheaf such that
$\sF\sim \sF'\otimes_{\dZ_\ell}\dQ_\ell$. We shall see in
(\ref{II:4-11}) that
\[
  \left(\varprojlim \h_c^q(X,\sF'\otimes \dZ/\ell^n)\right)\otimes_{\dZ_\ell}\dQ_\ell
\]
is a finite-dimensional $\dQ_\ell$-vector space. It depends only on $\sF'$,
and is denoted by $\h_c^q(X,\sF)$. The proof does not use the fact that
$\ell$ is prime to the characteristic, but only that case concerns us.






\subsection{}\label{II:2-11}

In this exposé, all important calculations will be carried out at the level
of $\dZ/\ell^n$-sheaves, with passage to $\dQ_\ell$-sheaves taking place only
at the very end. In practice, however, it is often convenient to work
systematically with $\dQ_\ell$-sheaves. In Exposés V and VI of
\cite{sga5}, Jouanolou gives the formalism that makes this possible. One of
the essential results (\hyperref[VI:2-2]{VI 2.2}) is that if $f:X\to Y$ is a separated morphism
of finite type of Noetherian schemes, and $\sF$ a $\dZ_\ell$-sheaf on $Y$,
then, for every $q$, the pro-sheaf
$\text{``}\varprojlim\text{''}\eR^q f_! (\sF\otimes\dZ/\ell^n)$ is a
$\dZ_\ell$-sheaf on $Y$. He even proves a stronger regularity property (of
Artin--Rees type) for the inverse system of the
$\eR^q f_!(\sF\otimes\dZ/\ell^n)$.

Jouanolou even works in a more general setting, encompassing the very useful
case of $E_\lambda$-sheaves, for $E_\lambda$ a finite extension of
$\dQ_\ell$. They are defined as the $\dQ_\ell$-sheaves were, with
$\dQ_\ell$ replaced by $E_\lambda$, $\dZ_\ell$ by the ring $\sO_\lambda$ of
the valuation $\lambda$, and $\dZ/\ell^n$ by $\sO_\lambda/\pi^n$, where
$\pi$ is a uniformizer. It would amount to the same thing to define an
$E_\lambda$-sheaf as a $\dQ_\ell$-sheaf $\sF$ equipped with a homomorphism
$E_\lambda\to\operatorname{End}(\sF)$.









\section{Trace formula and \texorpdfstring{$L$}{L}-functions}\label{II:3}

\emph{We use the notation of \S\ref{II:1}, and $\ell$ is a prime number
$\ne p$.}

Let $X_0$ be a separated scheme of finite type over $\dF_q$ ($q=p^f$), and
let $\sF_0$ be a constructible $\dQ_\ell$-sheaf on $X_0$. Grothendieck's
cohomological interpretation of $L$-functions is the following theorem:






\begin{theorem_}\label{II:3-1}
$L(X_0,\sF_0) = \prod_i \det\left(1-F^* t^f, \h_c^i(X,\sF)\right)^{(-1)^{i+1}}$.
\end{theorem_}

It is obtained as a consequence of the trace formula:






\begin{theorem_}\label{II:3-2}
For every integer $n$,
$\sum_{x\in X^{F^n}} \tr\left({F^n}^*,\sF_x\right) = \sum_i (-1)^i \tr\left({F^n}^*,\h_c^i(X,\sF)\right)$.
\end{theorem_}






Let us briefly recall the classical method for deducing (\ref{II:3-1}) from
(\ref{II:3-2}). Set $T=t^f$; both sides are formal series in $T$ with
constant term $1$. Since $\dQ_\ell$ has characteristic $0$, it suffices to
prove that the logarithmic derivatives $T\frac{d}{dT}\log$ of the two sides
are equal.

To compute these logarithmic derivatives, we shall use the formula
\[
  T \frac{d}{dT}\log\det(1-f T)^{-1} = \sum_{n>0} \tr(f^n) T^n \text{,}
\]
% NOTE: the following lines are not typesetting correctly
whose proof (in a slightly more general setting) is recalled below. On the
left-hand side, one finds
\begin{align*}
  T\frac{d}{dT}\log L(X_0,\sF_0)
    &= \sum_{x\in |X_0|} \sum_n [k(x):\dF_q] \tr\left(F_x^n,\sF\right) T^{n\cdot[k(x):\dF_q]} \\
    &= \sum_{n>0} T^n \sum_{x\in X^{F^n}} \tr\left({F^n}^*,\sF_x\right)
\end{align*}
and on the right-hand side
\[
  \sum_{n>0} T^n \sum (-1)^i \tr\left({F^n}^*,\h_c^i(X,\sF)\right) \text{ ; }
\]
the terms are compared one by one.






\begin{proposition_}\label{II:3-3}
Let $A$ be a commutative ring, $M$ a projective $A$-module of finite type,
and $f$ an endomorphism of $M$. For every invertible formal series $Q$, set
$\frac{d}{dT}\log Q = Q^{-1}\frac{d}{dT}Q$. One has
\begin{equation*}\tag{3.3.1}\label{II:eq:3-3-1}
  T\frac{d}{dT}\log\det(1-f T)^{-1} = \sum_{n>0} \tr(f^n) T^n \text{.}
\end{equation*}
\end{proposition_}

Let $M'$ be an $A$-module such that $M\oplus M'$ is free of finite rank.
Replacing $(M,f)$ by $(M\oplus M',f\oplus 0)$ changes neither side of
(\ref{II:3-3}), and reduces us to the case where $M$ is free. For $M=A^d$,
\eqref{II:eq:3-3-1} is an algebraic identity in the coefficients $f_i^j$ of
$f$. By the principle of extension of algebraic identities, it suffices to
verify it for $A$ an algebraically closed field of characteristic $0$. Since
both sides are additive in $M$ in an extension
$0\to M'\to M\to M''\to 0$, it even suffices to treat only the case where
$M$ has rank $1$. If $f$ is multiplication by $a$,
\eqref{II:eq:3-3-1} then reduces to the formula
\[
  T\frac{d}{dT}\log\left((1-a T)^{-1}\right) = \sum_{n>0} a^n T^n \text{.}
\]






\begin{corollary_}\label{II:3-4}
Let $n$ be an integer, and let $f^{(n)}$ be the endomorphism
$(x_1,\dotsc,x_n)\mapsto \left(f(x_n),x_1,\dotsc,x_{n-1}\right)$ of $M^n$.
One has
\begin{equation*}\tag{3.4.1}\label{II:eq:3-4-1}
  \det(1-f T^n) = \det(1-f^{(n)} T) \text{.}
\end{equation*}
\end{corollary_}

Proceeding as in (\ref{II:3-3}), we reduce to assuming that $A$ has
characteristic $0$. It then suffices to verify that the logarithmic
derivatives $-T\frac{d}{dT}\log$ of the two sides are equal.
\begin{align*}
  -T\frac{d}{dT}\log\det(1-f T^n) &= - n T^n \frac{d}{dT^n} \log\det(1-f T^n) = n \sum_{m>0} \tr(f^m) T^{n m} \text{,} \\
  -T \frac{d}{dT}\log\det\left(1-f^{(n)} T\right) &= \sum_{m>0} \tr\left(f^{(n) m}\right) T^m \text{,}
\end{align*}
and one observes that
\begin{align*}
  \tr(f^{(n) m}) &= 0 \qquad \text{if $n\nmid m$} \\
  \tr(f^{(n) n m}) &= n \tr(f^m) \text{.}
\end{align*}






\subsection{Remark}\label{II:3-5}

For $A=\dQ_\ell$, \eqref{II:eq:3-4-1} is the special case
$X_0=\spec(\dF_{q^n})$ of (\ref{II:3-1}). Using this identity, it is easy to
verify directly that the right-hand side of (\ref{II:3-1}) is independent of
$q$ and of the structure morphism $X_0\to\spec(\dF_q)$.






\subsection{Remark}\label{II:3-6}

If $j:X_0\hookrightarrow \bar X_0$ is a compactification of $X_0$, one has
$L(X_0,\sF_0) = L(\bar X_0, j_!\sF_0)$ and
$\h_c^i(X,\sF) = \h^i(\bar X,j_!\sF)$. This reduces (\ref{II:3-1}) to the
proper case, and explains the use of cohomology with proper support. For $X$
proper, (\ref{II:3-2}) has the form of a Lefschetz trace formula. Note that
the local terms $\tr(F^{n *},\sF_x)$ ($x\in X^{F^n}$) depend only on the
fiber of $\sF$ at $x$ and carry no multiplicity. For $X$ and $\sF$ smooth,
this corresponds to the fact that the graph of $F$ is transverse to the
diagonal (since $d F = 0$).






\subsection{Remark}\label{II:3-7}

If the formalism of $\dQ_\ell$-sheaves is admitted (cf. \ref{II:2-11}; one
needs the Leray spectral sequence and the base-change theorem for the higher
direct images $\eR^q f_!$), it is easy to reduce the proof of
(\ref{II:3-1}, \ref{II:3-2}) to the case where $X_0$ is a smooth curve and
$\sF_0$ is lisse. This is clearly explained in \cite[\S 5]{gr95} (for
\ref{II:3-1}; \ref{II:3-2} is treated in the same way).






\subsection{}\label{II:3-8}

There are two methods for proving (\ref{II:3-2}).



\paragraph{Lefschetz--Verdier}
If $X_0$ is proper (cf. \ref{II:3-6}), the general Lefschetz--Verdier trace
formula expresses the right-hand side of (\ref{II:3-2}) as a sum of local
terms, one for each point of $X^{F^n}$. In the original version of
\cite{sga5}, this formula was proved only modulo resolution of singularities.
The reader will find an unconditional proof in the definitive version. In
the case of curves, to which one can reduce (\ref{II:3-7}), all the
ingredients of the proof were in fact already available.

To deduce (\ref{II:3-2}) from the Lefschetz--Verdier formula, one must be
able to compute its local terms. For a curve and the Frobenius endomorphism,
this had been done by Artin and Verdier \cite{ve67}.



\paragraph{Nielson--Wecken}
A method inspired by the work of Nielsen and Wecken makes it possible to
reduce (\ref{II:3-2}) to a special case proved by Weil; this is what will be
explained in the following sections.











\section{Reduction to theorems in \texorpdfstring{$\dZ/\ell^n$}{Z/l n}-cohomology}\label{II:4}

We shall prove (\ref{II:3-2}) by passage to the limit from an analogous
theorem in $\dZ/\ell^n$-cohomology. Its statement presents the following
difficulty. Let $X_0$ be a separated scheme of finite type over $\dF_q$, and
let $\sF_0$ be a flat constructible sheaf of $\dZ/\ell^n$-modules. Its fibers
are finite-rank free $\dZ/\ell^n$-modules, and the left-hand side of
(\ref{II:3-2}) therefore makes sense (it is an element of $\dZ/\ell^n$).
The $\h_c^i(X,\sF)$, on the other hand, will not in general be free, so the
right-hand side is not defined. Overcoming this difficulty is one of
Grothendieck's essential uses of the theory of derived categories.






\subsection{}\label{II:4-1}

For the definition of the derived category $\eD(\cA)$ of an abelian category
$\cA$, and that of the subcategories $\eD^+(\cA)$, $\eD^-(\cA)$,
$\eD^b(\cA)$, I refer to Verdier's text in this volume. For $\Lambda$ a ring,
write $\eD(\Lambda)$ instead of
$\eD(\text{category of left $\Lambda$-modules})$; for $X$ a site, write
$\eD(X,\Lambda)$ instead of
$\eD($category of sheaves of left $\Lambda$-modules$)$. Likewise for
$\eD^-$, etc\ldots

The sign $\eR$ (resp. $\mathsf{L}$) indicates a right-derived (resp.
left-derived) functor. For example, $\lotimes$ denotes the derived functor of
the tensor product: for $K$ in $\eD^-(\Lambda^\circ)$ and $L$ in
$\eD^-(\Lambda)$ ($\Lambda^\circ$ the opposite ring of $\Lambda$),
$K\lotimes L$ is computed as follows: take quasi-isomorphisms $K'\to K$,
$L'\to L$, with $K'$ and $L$ bounded above and one of them having flat
components, and take the simple complex $s(K'\otimes_\Lambda L)$ associated
with the double complex $K'\otimes_\Lambda L'$.






\subsection{}\label{II:4-2}

Even to treat only the case of $\dZ/\ell^n$-cohomology, we shall need the
theory of traces of endomorphisms of modules over rings that are not
necessarily commutative (finite-group algebras $\dZ/\ell^n[G]$). These
traces were introduced by Stallings and Hattori \cite{ba76}.

Let $\Lambda$ be a ring (with identity, but not necessarily commutative).
Denote by $\Lambda^\natural$ the quotient of the additive group of $\Lambda$
by the subgroup generated by the commutators $a b-b a$. If $f$ is an
endomorphism of the free left $\Lambda$-module $\Lambda^n$, with matrix
$f_i^j$, denote by $\tr(f)$ the image of $\sum_i f_i^i$ in
$\Lambda^\natural$. If $f$ and $g$ are two homomorphisms
$f:\Lambda^n \leftrightarrows \Lambda^m : g$, one immediately verifies that
$\tr(f g) = \tr(g f)$.

Let $f$ be an endomorphism of a projective $\Lambda$-module of finite type
$P$. Write $P$ as a direct factor of a free module $\Lambda^n$: this means
choosing a module $P'$ and an isomorphism
$\alpha:P\oplus P'\iso \Lambda^n$, or equivalently choosing a homomorphism
$a:P\to \Lambda^n$ and a retraction $b:\Lambda^n\to P$ (take
$P'=\ker(b)$). Let $f'=\alpha(f\oplus 0)\alpha^{-1} =af b$. Set
$\tr(f) = \tr(f')$. This element of $\Lambda^\natural$ depends only on $f$:
if another choice $c:P\leftrightarrows \Lambda^m:d$ is made, one has
$a=a(dc)(ba)$, whence
\[
 \tr(a f b) = \tr(adcbafb) = \tr(cbafbad) = \tr(cfd) \text{.}
\]

If $f$ is an endomorphism of a $\dZ/2$-graded projective
$\Lambda$-module of finite type, with components $f_i^j:P^j\to P^i$, set
\[
  \tr(f) = \tr(f_0^0) - \tr(f_1^1) \text{.}
\]
If there is a risk of ambiguity, we shall instead write $\tr(f;P)$. If the
homomorphisms $f:P\leftrightarrows Q:g$ are homogeneous, reduction to the
free case immediately verifies that
\begin{equation*}\tag{4.2.1}\label{II:eq:4-2-1}
  \tr(f g) = \tr\left((-1)^{\deg f \deg g} g f\right) \text{.}
\end{equation*}

If $f$ is an endomorphism of a filtered $\dZ/2$-graded $\Lambda$-module
$(P,F)$, with finite filtration compatible with the grading, and the
$\gr_F^i(P)$ are projective of finite type, then $P$ is projective of finite
type, and
\begin{equation*}\tag{4.2.2}\label{II:eq:4-2-2}
  \tr(f;P) = \sum \tr\left(f; \gr_F^i(P) \right)
\end{equation*}
(split the filtration to reduce to the case of a direct sum).






\subsection{}\label{II:4-3}

A $\dZ$-graded module will be regarded as $\dZ/2$-graded by the parity of
the degree. In particular, if $f$ is an endomorphism of a bounded complex of
projective $\Lambda$-modules of finite type, set
\begin{equation*}\tag{4.3.1}\label{II:eq:4-3-1}
  \tr(f) = \sum (-1)^i \tr(f^i) \text{.}
\end{equation*}

By \eqref{II:eq:4-2-1}, if $f$ is homotopic to zero, $\tr(f) = 0$:
\begin{equation*}\tag{4.3.2}\label{II:eq:4-3-2}
  f = d H + H d \rightarrow \tr(f) = 0 \text{.}
\end{equation*}

Let $\mathsf{K}_\text{parf}(\Lambda)$ be the category whose objects are
bounded complexes of projective $\Lambda$-modules of finite type and whose
arrows are morphisms of complexes taken up to homotopy. The functor
$\mathsf{K}_\text{parf}(\Lambda)\to\eD(\Lambda)$ is fully faithful. Denote its
essential image by $\mathsf{D}_\text{parf}(\Lambda)$. Formulas
\eqref{II:eq:4-3-1} and \eqref{II:eq:4-3-2} make it possible to define
$\tr(f)$ for $f$ an endomorphism of
$K\in\ob \mathsf{D}_\text{parf}(\Lambda)$ (and this trace depends only on the
isomorphism class of $(K,f)$).






\subsection{}\label{II:4-4}

Recall that the filtered derived category $\mathsf{DF}(\cA)$ of an abelian
category $\cA$ is the category obtained from that of complexes of objects of
$\cA$ with finite filtration by inverting the filtered quasi-isomorphisms
($=$ morphisms of filtered complexes inducing a quasi-isomorphism on the
associated graded). There are functors ``underlying complex'':
$\mathsf{DF}(\cA)\to\eD(\cA): (K,F)\mapsto K$, and ``$p$-th component of the
associated graded'':
$\mathsf{DF}(\cA)\to\eD(\cA): (K,F)\mapsto \gr_F^p(K)$.

Let $\mathsf{KF}_\text{parf}(\Lambda)$ be the category whose objects are
bounded filtered complexes $(K,F)$ with finite filtration, with the
$\gr_F^p(K^q)$ projective of finite type, and whose arrows are
filtration-preserving morphisms of complexes taken up to a
filtration-preserving homotopy. The functor
$\mathsf{KF}_\text{parf}(\Lambda)\to\mathsf{DF}(\Lambda)$ is fully faithful;
its essential image $\mathsf{DF}_\text{parf}(\Lambda)$ consists of the $(K,F)$ in
$\mathsf{DF}(\Lambda)$ such that the $\gr_F^p(K)$ belong to
$\mathsf{D}_\text{parf}(\Lambda)$ and are almost all zero; if $(K,F)$ belongs to
$\mathsf{DF}_\text{parf}(\Lambda)$, then $K$ belongs to
$\mathsf{D}_\text{parf}(\Lambda)$.

By \eqref{II:eq:4-2-2}, if $f$ is an endomorphism of
$(K,F)\in\ob\mathsf{DF}_\text{parf}(\Lambda)$, one has
\begin{equation*}\tag{4.4.1}\label{II:eq:4-4-1}
 \tr(f,K) = \sum_p \tr(f,\gr_F^p(K)) \text{.}
\end{equation*}






\subsection{}\label{II:4-5}

Let $\Lambda$ be a ring, and let $K\in\ob\eD^-(\Lambda)$. We say that $K$
has tor-dimension $\leqslant r$ if, for every right $\Lambda$-module $N$, the
hypertors $\h^i(N\lotimes_\Lambda K)$ vanish for $i<-r$. It is said to have
finite tor-dimension if this condition holds for some $r$.

The same terminology is used for a complex of sheaves
$K\in\ob\eD^-(X,\Lambda)$.






\begin{lemma}\label{II:4-5-1}
Let $\Lambda$ be a left Noetherian ring, and let
$K\in\ob\eD^-(\Lambda)$. Then $K$ belongs to
$\mathsf{D}_\text{parf}(\Lambda)$ if and only if it has finite tor-dimension
and the $\h^i(K)$ are of finite type.
\end{lemma}

After replacing $K$ by an isomorphic complex (in $\eD^-(\Lambda)$), we may
suppose that $K$ is bounded above and has finite-rank free components (cf.
\ref{II:4-7} below). If $K$ has tor-dimension $\leqslant r$, the
$\h^i(K)$ vanish for $i<-r$, $K^{-r}/\im(d)$ admits the free resolution
\[
  (\cdots \to K^{-r-1} \to K^{-r}) \to K^{-r}/\im(d) \text{,}
\]
and for every right $\Lambda$-module $N$,
$\h^{-r-k}(N\lotimes_\Lambda K) = \h^{-r-k}(N\otimes_\Lambda K)
= \tor_k(N,K^{-r}/\im d)$ for $k\geqslant 1$.

By hypothesis, these groups vanish; the module $K^{-r}/\im d$ is therefore
flat and finitely presented, i.e. projective of finite type. The quotient
\[
  0 \to K^{-r}/\im(d) \to K^{-r+1} \to K^{-r+2}\to \cdots
\]
of $K$ is quasi-isomorphic to $K$, and this shows that
$K\in\ob\mathsf{D}_\text{parf}(\Lambda)$.






\begin{prop-def_}\label{II:4-6}
Let $X$ be a Noetherian scheme, and let $\Lambda$ be a left Noetherian ring.
Denote by $\eD_\textnormal{ctf}^b(X,\Lambda)$ the subcategory of
$\eD^-(X,\Lambda)$ formed by the complexes $\sK$ satisfying the following
equivalent conditions:
\begin{enumerate}[\indent i)]
  \item $\sK$ is isomorphic (in $\eD^-(X,\Lambda)$) to a bounded complex of
    flat constructible sheaves of $\Lambda$-modules.
  \item $\sK$ has finite tor-dimension and the sheaves $\h^i(K)$ are
    constructible.
\end{enumerate}
\end{prop-def_}

The same argument as in (\ref{II:4-5}), starting from the following lemma.






\begin{lemma_}\label{II:4-7}
If a complex $\sK$ of sheaves of $\Lambda$-modules is such that the
$\h^i(\sK)$ are constructible and vanish for large $i$, there is a
quasi-isomorphism $\sK'\to \sK$, with $\sK'$ bounded above and having flat
constructible components.
\end{lemma_}

Construct $\sK'$ by proceeding from right to left. For an $m$ such that the
$\h^i(\sK)$ ($i\geqslant m$) vanish, take ${\sK'}^i = 0$ for
$i\geqslant m$. Next suppose that the ${\sK'}^i$ have already been
constructed for $i>n$:
\[\xymatrix{
  \cdots \ar[r]
    & \sK^{n-1} \ar[r]
    & \sK^n \ar[r]
    & \sK^{n+1} \ar[r]
    & \cdots \\
  & & & {\sK'}^{n+1} \ar[u] \ar[r]
    & \cdots \text{,}
}\]
with $\h^i(\sK')\iso \h^i(\sK)$ for $i>n+1$, and
$\ker(d)\to\h^{n+1}(\sK)$ surjective in degree $n+1$. Define sheaves $\sA$
and $\sB$ by the Cartesian diagram
\[\xymatrix{
  \sK^n \ar[r]
    & \sK^n/\im(d) \ar[r]
    & \ker(d) \\
  \sA \ar[u] \ar[r]^u
    & \sB \ar[u]\ar[r]
    & \ker(d) \ar[u] \text{.}
}\]

Then $\sB$ is constructible, $u$ is surjective, and to advance one step it
suffices to construct ${\sK'}^n$ flat and constructible and
$v:{\sK'}^n\to \sA$ such that $u v$ is surjective. The possibility of doing so
is guaranteed by the following lemma.






\begin{lemma_}\label{II:4-8}
Let $v:\sA\to \sB$ be an epimorphism of sheaves of $\Lambda$-modules, with
$\sB$ constructible. There is then $u:\sK\to \sA$, with $\sK$ flat and
constructible, such that $v u$ is an epimorphism.
\end{lemma_}

The sheaf $\sA$ is a quotient of a sum of sheaves of the form
$\varphi_! \Lambda$, for $\varphi:U\to X$ étale. Since $\sB$ is Noetherian
(\cite[IX.2.10]{sga4}), a finite sum of these already maps onto $\sB$, which
proves the lemma.






\begin{theorem_}\label{II:4-9}
Let $f:X\to Y$ be a separated morphism of finite type of Noetherian schemes.
If $\sK\in\ob\eD_\textnormal{ctf}^b(X,\Lambda)$, then
$\eR f_!\sK\in \ob\eD_\textnormal{ctf}^b(Y,\Lambda)$.
\end{theorem_}

Recall that if $j:X\hookrightarrow \bar X \xrightarrow{\bar f} Y$ is a
relative compactification of $X/Y$, then by definition
$\eR f_! = \eR \bar f_* \circ j_!$, where $\eR \bar f_*$ is the derived
functor of $\bar f_*$. This formula reduces the proof of (\ref{II:4-9}) to
the case where $f$ is proper. For $f$ proper, $f_*$ has finite cohomological
dimension; this makes it possible to define $\eR f_*$ on the entire derived
category, and also as a functor
\[
  \eR f_* : \eD^-(X,\Lambda) \to \eD^-(Y,\Lambda) \text{.}
\]
By composition, one likewise defines in general
\[
  \eR f_! : \eD^-(X,\Lambda) \to \eD^-(Y,\Lambda) \text{.}
\]

For every right $\Lambda$-module $N$, one has (proof given below)
\begin{equation*}\tag{4.9.1}\label{II:eq:4-9-1}
  N\lotimes \eR f_! \sK \iso \eR f_! (N\lotimes \sK) \text{.}
\end{equation*}

Let us deduce (\ref{II:4-9}) from \eqref{II:eq:4-9-1}.
\begin{enumerate}[\indent a)]
  \item The spectral sequence
    \[
      E_2^{p q} = \eR^p f_! \h^q(\sK) \rightarrow \h^{p+q} \eR f_! \sK \text{.}
    \]
    and the finiteness theorem for $\eR f_!$ show that
    $\eR f_!\sK\in \ob \eD^b(Y,\Lambda)$ and has constructible cohomology.
    It remains to prove finite tor-dimension (\ref{II:4-6}).
  \item If $\sK$ has tor-dimension $\leqslant -r$, the same is true of
    $\eR f_! \sK$: one has
    $\h^i(N\lotimes \eR f_! \sK) = \h^i(\eR f_!(N\lotimes \sK))$, and in the
    spectral sequence above for $N\lotimes \sK$, the $E_2^{p q}$ vanish for
    $q<r$, hence a fortiori for $p+q<r$.
\end{enumerate}

Let us prove \eqref{II:eq:4-9-1}, remaining at the level of complexes.
\begin{enumerate}[\indent a)]
  \item Using the defining formula $\eR f_! = \eR f_*\circ j_!$, we reduce to
    assuming that $f$ is proper.
  \item Represent $\sK$ by a complex bounded above whose components are
    acyclic for $f_*:\eR^p f_* \sK^q = 0$ for $p>0$. One then has
    $\eR f_* K \sim f_* \sK$. It is possible to work in this way with
    complexes bounded above only because $f_*$ has finite cohomological
    dimension.
  \item If $L$ is a free right $\Lambda$-module, one has
    $\eR^p f_*(L\otimes \sK^q) = L\otimes \eR^p f_* \sK^q$: this is trivial
    when $\Lambda$ is free of finite type, and the general case follows by
    passage to the limit. The sheaf $L\otimes \sK^q$ is acyclic for $f_*$,
    and $f_*(L\otimes \sK^q) = L\otimes f_* \sK^q$.
  \item Let $N_\bullet$ be a free resolution of $N$. One has
    $N\lotimes \sK\sim s(N_\bullet\otimes \sK)$ (the simple complex
    associated with the double complex $N_\bullet\otimes \sK$). By c),
    \[
      \eR f_*(N\lotimes \sK) \sim \eR f_* (s(N_\bullet\otimes \sK)) \sim f_*(s(N_\bullet\otimes \sK)) \sim s(N_\bullet\otimes f_* \sK) \sim N\lotimes \eR f_* \sK \text{,}
    \]
    which is the asserted formula.
\end{enumerate}

When $Y$ is the spectrum of a separably closed field, the functor $\eR f_!$
is identified with a functor denoted by
$\eR \Gamma_c : \eD_\text{ctf}^b(X,\Lambda)\to\eD_\text{parf}(\Lambda)$.






\begin{theorem_}\label{II:4-10}
With the notation of \S\ref{II:1}, let $X_0$ be a separated scheme of finite
type over $\dF_q$, and let $\Lambda$ be a Noetherian torsion ring annihilated
by an integer prime to $p$. Let
$\sK_0\in\ob\eD_\textnormal{ctf}^b(X_0,\Lambda)$. With the definition of
traces in (\ref{II:4-3}), one has for every $N$
\begin{equation*}\label{II:eq:4-10}
  \sum_{x\in X^{F^n}} \tr\left({F^n}^*,\sK_x\right) = \tr\left({F^n}^*,\eR\Gamma_c(X,\sK)\right) \text{.}
\end{equation*}
\end{theorem_}

In the remainder of this \S, we shall deduce \ref{II:3-2} from
\ref{II:4-10}, and prove \ref{II:2-10}.






\subsection{Proof of \texorpdfstring{\ref{II:2-10}}{2.10}}\label{II:4-11}

Let $X$ be a separated scheme of finite type over an algebraically closed
field $k$, and let $\sG$ be a $\dQ_\ell$-sheaf on $X$. By \ref{II:2-8}, it
can be written in the form $\sG=\sF\otimes\dQ_\ell$, with $\sF$ a
torsion-free $\dZ_\ell$-sheaf. Set
\[
  K_n = \eR \Gamma_c(X,\sF\otimes\dZ/\ell^{n+1})\in\eD(\dZ/\ell^{n+1}) \text{.}
\]
By \ref{II:4-9}, one has
$K_n\in\ob\eD_\text{parf}(\dZ/\ell^{n+1})$, and it follows from
\eqref{II:eq:4-9-1} that the natural arrow
\begin{equation*}\tag{4.11.1}\label{II:eq:4-11-1}
\xymatrix{
  K_n\lotimes_{\dZ/\ell^{n+1}} \dZ/\ell^n \ar[r] & K_{n-1} \text{,}
}
\end{equation*}
is an isomorphism. The problem of defining this arrow is discussed in
\ref{II:4-12}.

We may now invoke \cite[XV.3.3]{sga5} (p. 31, lines 1--2 from the end; this
passage is independent of the rest of \cite{sga5}) to realize each $K_n$ by
a bounded complex of finite-rank free $\dZ/\ell^{n+1}$-modules, and the
arrows \eqref{II:eq:4-11-1} by \emph{isomorphisms of complexes}
\[\xymatrix{
  K_n\otimes \dZ/\ell^n \ar[r]^-\sim
    & K_{n-1} \text{.}
}\]
Let $K$ be the complex $\varprojlim K_n$. It is a bounded complex of free
$\dZ_\ell$-modules, and
$K_n\simeq K\otimes_{\dZ_\ell}\dZ/\ell^{n+1}$. One has
\[
  \h^i(K) = \varprojlim \h^i(K_n) \text{,}
\]
(indeed, every inverse system of finite groups satisfies the Mittag--Leffler
condition), and $\h^i(K)$ is a $\dZ_\ell$-module of finite type. One has
\[
  \h_c^i(X,\sG) = \left(\varprojlim \h^i(K_n) \right)\otimes_{\dZ_\ell} \dQ_\ell \text{:}
\]
this is a finite-dimensional $\dQ_\ell$-vector space.






\subsection{The arrow \texorpdfstring{\eqref{II:eq:4-11-1}}{(4.11.1)}}\label{II:4-12}

To \emph{define} the arrow \eqref{II:eq:4-11-1}, one cannot proceed as in
\ref{II:4-9}, i.e. resolve $\dZ/\ell^n$ by a complex of free
$\dZ/\ell^{n+1}$-modules, since one wishes to construct an isomorphism in
$\eD_\text{parf}(\dZ/\ell^n)$. The desired arrow is a special case of the
following scalar-extension arrow.

(*) Let $\Lambda\to\Lambda'$ be a homomorphism of Noetherian torsion rings.
For $X$ a separated scheme of finite type over an algebraically closed $k$,
and $\sK\in\ob\eD_\text{ctf}(X,\Lambda)$, one has an isomorphism
\[\xymatrix{
  \eR\Gamma_c(X,\sK)\lotimes_\Lambda \Lambda' \ar[r]^-\sim & \eR \Gamma_c(X,\sK\lotimes_\Lambda \Lambda') \text{.}
}\]

More general arrows are constructed in \cite[XVII 4.2.12]{sga4}. The method
is as follows:
\begin{enumerate}[\indent a)]
  \item The definition of $\eR\Gamma_c$ and the formula
    $j_!(\sK\otimes_\Lambda\Lambda') = (j_! \sK)\otimes_\Lambda\Lambda'$ for
    an open immersion reduce us to the case where $X$ is proper.
  \item For $X$ proper, one must replace $\sK$ by a bounded complex whose
    components are acyclic for the functor $\Gamma$, and whose fiber at every
    geometric point (or merely every closed point) is homotopic to a complex
    of flat $\Lambda$-modules. For such a complex, one has
    $\Gamma(X,\sK)\sim \eR\Gamma(X,\sK)$ and
    $K\otimes_\Lambda\Lambda'\sim \sK\lotimes_\Lambda\Lambda'$; the arrow
    \eqref{II:eq:4-11-1} is defined as the composite
    \[
      \eR\Gamma(X,\sK)\lotimes_\Lambda \Lambda' \to \Gamma(X,\sK)\otimes_\Lambda\Lambda' \to \Gamma(X,\sK\otimes_\Lambda\Lambda') \xleftarrow\sim \eR\Gamma(X,\sK\lotimes_\Lambda\Lambda') \text{.}
    \]
  \item It remains to show that $\sK$ can be replaced by a complex of the
    desired kind \cite[XVII.4.2.10]{sga4}. First replace $\sK$ by a bounded
    complex with flat components (\ref{II:4-6}), then take a truncated
    canonical flasque resolution of it (this does not modify the homotopy
    type of the fiber complexes, and one truncates far enough for the
    acyclicity condition for $\Gamma$ to hold (i.e. beyond the cohomological
    dimension of $\Gamma$)).
\end{enumerate}






\subsection{Proof of \texorpdfstring{(\ref{II:3-2})}{(3.2)}}\label{II:4-13}

To simplify notation, we shall treat only the case $n=1$ of
(\ref{II:3-2}). The general case follows by replacing $F$ below by $F^n$.

Thus let $X_0$ be separated and of finite type over $\dF_q$, and let $\sG_0$
be a $\dQ_\ell$-sheaf on $X_0$. One has
$\sG_0=\sF_0\otimes\dQ_\ell$, for a suitable torsion-free
$\dZ_\ell$-sheaf $\sF_0$. Let
$K_n=\eR\Gamma_c(X,\sF\otimes\dZ/\ell^n)$. The Frobenius morphisms induce
morphisms
\[
  F^* : K_n\to K_n \qquad \text{(in $\eD_\text{parf}(\dZ/\ell^{n+1})$),}
\]
which are deduced from one another via the isomorphisms
\eqref{II:eq:4-11-1}.

As in (\ref{II:4-11}), realize the $K_n$ and the isomorphisms
\eqref{II:eq:4-11-1} at the level of complexes. The already cited passage
\cite[XV.3.3]{sga5} makes it possible to realize the morphisms $F^*$ by
endomorphisms of complexes, again denoted by $F^*$, which reduce to one
another. Again denote by $F^*$ their inverse limit and the endomorphisms
deduced from it. Since
$\h_c^i(X,\sG) = \h^i(K)\otimes\dQ_\ell = \h^i(K\otimes\dQ_\ell)$, one has
(with the sign convention of \ref{II:eq:4-3-1})
\[
  \tr\left(F^*, \h_c^\bullet(X,\sG)\right) = \tr\left(F^*,K^\bullet\otimes\dQ_\ell\right) = \varprojlim_n \tr\left(F^*,K_n^\bullet\right) = \varprojlim_n \tr\left(F^*,\eR\Gamma_c(X,\sF\otimes\dZ/\ell^{n+1})\right) \text{.}
\]
Apply (\ref{II:4-10}) to $\sF_0\otimes\dZ/\ell^{n+1}$ (viewed as a complex
concentrated in degree $0$), for $\Lambda=\dZ/\ell^{n+1}$. One finds that
\begin{align*}
  \tr\left(F^*,\eR\Gamma_c(X,\sF\otimes\dZ/\ell^{n+1})\right)
    &= \sum_{x\in X^F} \tr\left(F^*,\sF_x\otimes\dZ/\ell^{n+1}\right) \\
    &= \sum_{x\in X^F} \tr\left(F^*,\sG_x\right) \mod \ell^{n+1} \text{,}
\end{align*}
and (\ref{II:3-2}) follows by passage to the limit.











\section{The Nielsen--Wecken method}\label{II:5}

In this chapter, we prove two special cases of (\ref{II:4-10}). The general
case will be considered in the next chapter.






\begin{lemma_}\label{II:5-1}
Theorem (\ref{II:4-10}) is true for $X_0$ of dimension $0$.
\end{lemma_}

In dimension $0$, formula \eqref{II:eq:4-10} holds for every correspondence,
and not only for the iterates of a Frobenius. It is a special case of the
following statement (applied to $X(\dF)$, $K$, and the ${F^*}^n$).

(*) Let $X$ be a finite set, $F:X\to X$,
$\sK\in\ob\eD_\text{parf}(X,\Lambda)$, and $F^*:F^* \sK\to \sK$. One has
\[
  \sum_{X^F} \tr\left(F^*,\sK_x\right) = \tr\left(F^*,\Gamma(X,K)\right) \text{.}
\]

One has $\Gamma(X,\sK)=\bigoplus_{x\in X} \sK_x$, and the formula follows
from the fact that for every morphism
$u:\bigoplus_{x\in X} \sK_x\to \bigoplus_{x\in X} \sK_x$, with matrix
$u_x^y$, one has $\tr(u)=\sum_{x\in X} \tr\left(u_x^x\right)$.






\begin{lemma_}\label{II:5-2}
Let $\Lambda$ be a Noetherian torsion ring annihilated by a power of the prime
number $\ell\ne p$, let $X_0$ be a smooth projective connected curve over
$\dF_q$, let $j:U_0\hookrightarrow X_0$ be a dense open subscheme of $X_0$,
and let $\sG_0$ be a locally constant sheaf of projective
$\Lambda$-modules on $U_0$. One has
\[
  \sum_{x\in U^F} \tr\left(F^*,\sG_x\right) = \tr\left(F^*,\eR\Gamma_c(U,\sG)\right) \text{.}
\]
\end{lemma_}

More precisely, we shall deduce (\ref{II:5-2}) from the following theorem.






\begin{theorem_}\label{II:5-3}
Let $X$ be a nonsingular projective curve over an algebraically closed field
$k$, let $f$ be an endomorphism of $X$ with isolated fixed points, and let
$v(f)$ be the number of fixed points of $f$, each counted with its
multiplicity: $v(f)$ is the intersection number of the graph of $f$ with the
diagonal of $X\times X$. Then, for $\ell$ prime to the characteristic
exponent of $k$, one has
\[
  v(f) = \sum (-1)^i \tr\left(f^*,\h^i(X,\dQ_\ell)\right) \text{.}
\]
\end{theorem_}

This theorem is due to Weil (\cite{we48}, n$^\circ$ 43, 66, and 68), and Weil's
proof is reproduced in Lang \cite{la83} (VI \S3 combined with VII \S2,
Th. 3). It can also be deduced from the formalism of the cohomology class
associated with a cycle (\nameref{IV}, \ref{IV:3-7}).






\begin{corollary_}\label{II:5-4}
Let $j:U\hookrightarrow X$ be a dense open subscheme of $X$, such that
$f^{-1}(U)=U$. If the fixed points of $f$ in $X\setminus U$ have multiplicity
one, then the number $v_U(f)$ of fixed points of $f$ in $U$ (each counted
with its multiplicity) is
\[
  v_U(f) = \sum (-1)^i \tr\left(f^*,\h_c^i(U,\dQ_\ell)\right) \text{.}
\]
\end{corollary_}

Let $F$ be the finite set $X\setminus U$. The cohomology exact sequence
defined by the short exact sequence
$0\to j_! \dQ_\ell\to\dQ_\ell\to{\dQ_\ell}_F\to 0$,
\[\xymatrix{
  \ar[r]^-\partial
    & \h_c^i(U,\dQ_\ell) \ar[r]
    & \h^i(X,\dQ_\ell) \ar[r]
    & \h^i(F,\dQ_\ell) \ar[r]^-\partial
    &
}\]
gives, with the sign convention of (\ref{II:3-1}),
\[
  \tr\left(f^*,\h_c^\bullet(U,\dQ_\ell)\right)
    = \tr\left(f^*,\h^\bullet(X,\dQ_\ell)\right) - \tr\left(f^*,\h^\bullet(F,\dQ_\ell)\right)
    = v(f) - \tr(f^*,\dQ_\ell^F) \text{.}
\]
The trace of $f^*$ on $\dQ_\ell^F$ is the number $v_F(f)$ of fixed points of
$f$ in $F$; the multiplicity-one hypothesis therefore ensures that
\[
  \tr\left(f^*,\h_c^\bullet(U,\dQ_\ell)\right)
    = v(f) - v_F(f) = v_U(f) \text{,}
\]
as promised.

The proof of (\ref{II:5-2}) uses some lemmas concerning a special case of
noncommutative traces (\ref{II:4-2}), which we now develop. \emph{For
brevity, we shall say projective for projective of finite type}.






\subsection{}\label{II:5-5}

Let $\Lambda$ be a ring, and let $H$ be a finite group. The map
$\sum a_h \cdot a\mapsto $ class of $a_e$, from the group algebra
$\Lambda[H]$ to $\Lambda^\natural$, vanishes on the commutators $ab-ba$ of
$\Lambda[H]$. It therefore factors through
\[
  \varepsilon : \Lambda[H]^\natural \to \Lambda^\natural \text{.}
\]
If $F$ is an endomorphism of a projective $\Lambda[H]$-module $P$, set
\[
  \tr_\Lambda^H(F) = \varepsilon \tr_{\Lambda[H]} (F)
\]
where the trace of (\ref{II:4-2}) appears on the right-hand side. To avoid
ambiguity, we shall sometimes write this trace as $\tr_\Lambda^H(F,P)$.






\begin{proposition_}\label{II:5-6}
$\tr_\Lambda(F) = |H|\tr_\Lambda^H(F)$.
\end{proposition_}

The formal properties of traces reduce us to treating only the case
$P=\Lambda[H]$. The endomorphism $F$ is then right multiplication by an
element $\sum a_h\cdot h$ of $\Lambda[H]$, and
$\tr_\Lambda(F) = |H| a_e$, $\tr_\Lambda^H(F) = a_e$.






\subsection{}\label{II:5-7}

Let $A$ be a commutative ring, and let $\Lambda$ be an $A$-algebra.
Multiplication by an element of $A$ passes to the quotient to define an
$A$-module structure on $A^\natural$. If $H$ is a finite group, $P$ a
projective $A[H]$-module, and $M$ a $\Lambda[H]$-module projective as a
$\Lambda$-module, it is known that the $\Lambda$-module $P\otimes_A M$,
equipped with the diagonal action of $H$, is a projective
$\Lambda[H]$-module. Indeed:
\begin{enumerate}[\indent a)]
  \item It suffices to verify this for $P=A[H]$.
  \item The $\Lambda$-module $P\otimes_A M$, equipped with the action of $H$
    defined by its action on the first factor, is clearly a projective
    $\Lambda[H]$-module; denote it by $(P\otimes_A M)'$.
  \item The map $h\otimes m\mapsto h\otimes h^{-1} m$ induces an isomorphism
    \begin{equation*}\tag{5.7.1}\label{II:eq:5-7-1}
      A[H]\otimes_A M \iso (A[H]\otimes_A M)' \text{.}
    \end{equation*}
\end{enumerate}






\begin{proposition_}\label{II:5-8}
With the preceding notation, if $u$ is an endomorphism of $P$ and $v$ an
endomorphism of $M$, then
\[
  \tr_\Lambda^H(u\otimes v) = \tr_A^H(u)\tr_\Lambda(v) \text{.}
\]
\end{proposition_}

We reduce to assuming that $P=A[H]$; the endomorphism $u$ is then right
multiplication $m_x$ by an element $x=\sum a_h h$ of $A[H]$. Isomorphism
\eqref{II:eq:5-7-1} takes $u\otimes v$ to
$\sum a_h m_h\otimes h^{-1} v$, whose trace is
$a_e\cdot\tr_\Lambda(v)$.






\subsection{}\label{II:5-9}

Let $G$ be a group that is an extension of $\dZ$ by a finite group $H$
\[
  1 \to H \to G \to \dZ \to 1 \text{,}
\]
and let $G^+$ be the inverse image of $\dN$. Let $\Lambda$ be a ring, and let
$P$ be a $\Lambda[G^+]$-module projective as a $\Lambda[H]$-module. If
$g\in G^+$, denote by $Z_g$ the centralizer of $g$ in $H$; multiplication by
$g$ is an endomorphism of $P$, viewed as a $\Lambda[Z_g]$-module. Since $P$
is $\Lambda[Z_g]$-projective, a trace $\tr_\Lambda^{Z_g}(g)$ is defined. By
\ref{II:5-6}, one has
\begin{equation*}\tag{5.9.1}\label{II:eq:5-9-1}
  \tr_\Lambda(g) = |Z_g| \cdot \tr_\Lambda^{Z_g}(g) \text{.}
\end{equation*}

Suppose that $\Lambda$ is an algebra over a commutative ring $A$. For $P$ an
$A[G^+]$-module projective as an $A[H]$-module, and $M$ a
$\Lambda[G^+]$-module projective as a $\Lambda$-module, the
$\Lambda[G^+]$-module $P\otimes_A M$ (diagonal action of $G^+$) is projective
as a $\Lambda[H]$-module and, by \ref{II:5-8}, for every $g\in G^+$ one has
\begin{equation*}\tag{5.9.2}\label{II:eq:5-9-2}
  \tr_\Lambda^{Z_g}(g,P\otimes_A M) = \tr_A^{Z_g}(g,P) \cdot \tr_\Lambda(g,M) \text{.}
\end{equation*}






\subsection{}\label{II:5-10}

Let $P$ be a $\Lambda[G^+]$-module projective as a $\Lambda[H]$-module, and
let $P_H$ be the coinvariants of $H$ in $P$. This is a projective
$\Lambda$-module. The action of $g\in G^+$ passes to the quotient and defines
an endomorphism of $P_H$ that depends only on the image of $g$ in $\dN$. We
denote by $F$ the action of the elements $g\in G^+$ with image $1$.






\begin{proposition_}\label{II:5-11}
$\tr_\Lambda(F,P_H) = \sum_{g\mapsto 1}' \tr_\Lambda^{Z_g}(g,P)$, where
$\sum_{g\mapsto 1}'$ denotes a sum over the $H$-conjugacy classes of elements
of $G^+$ with image $1$ in $\dN$.
\end{proposition_}

After multiplication by $|H|$, this formula can be verified as follows:
\[
  |H|\cdot \tr_\Lambda(F,P_H)
    = \tr_\Lambda\left(\sum_{g\mapsto 1} g,P\right)
    = \sum_{g\mapsto 1}' \frac{|H|}{|Z_g|}\tr_\Lambda(g,P)
    = \sum_{g\mapsto 1}' |H|\cdot \tr_\Lambda^{Z_g}(g,P) \text{.}
\]

Choose an element $g\in G^+$ with image $1$ in $\dN$, and let $\sigma$ be
the automorphism of $\Lambda[H]$ induced by the automorphism
$\operatorname{ad} g$ of $H$. Giving a $\Lambda[G^+]$-module amounts to
giving a $\Lambda[H]$--module equipped with a $\sigma$-linear endomorphism
$\gamma$: let $g^n h$ act by $\gamma^n h$. In this language, formula
\ref{II:5-11} takes the form: for $P$ a projective $\Lambda[H]$-module and
$\gamma$ a $\sigma$-linear endomorphism of $P$, one has
\[
  \tr_\Lambda(\gamma,P_H) = \sum_h' \tr_\Lambda^{Z_g}(h\gamma,P)\text{,}
\]
where $\sum'$ is a sum over the $\operatorname{ad}g$-conjugacy classes in
$H$ (orbits of $H$ acting on itself by the maps
$x\mapsto h x \operatorname{ad} g(h)^{-1}$).

In this form, the properties of traces immediately reduce us to the case
$P=\Lambda[H]$. The endomorphism $\gamma$ then has the form
$x\mapsto \sigma(x) a$, with $a=\sum a_h h$ in $\Lambda[H]$. It suffices to
treat the universal case where the $a_h$ are indeterminates. In this case,
$\Lambda=\dZ\langle a_h\rangle_{h\in H}$ (noncommutative polynomials), and
$\Lambda^\natural$ is torsion-free: one obtains \ref{II:5-11} by the
argument at the beginning and division by $|H|$. One can also calculate
directly.






\subsubsection{Remark}\label{II:5-11-1}

Everything said above applies equally well to $\dZ/2$-graded modules, or to
complexes of modules (cf. \ref{II:4-2}, \ref{II:4-3}).






\subsection{}\label{II:5-12}

Retain the notation of \ref{II:5-2}, and let $A=\dZ/\ell^n$, with $n$ large
enough that $\Lambda$ is an $A$-algebra. Let $f:V_0\to U_0$ be a finite
\'etale connected covering of $U_0$, Galois with Galois group $H$, and such
that the sheaf $f^*\sG_0$ is constant. Denote its constant value
$\h^0(V_0,f^*\sG_0)$ by $M$; it is a $\Lambda[H]$-module projective as a
$\Lambda$-module.

The sheaf $f_* A$ on $U_0$, equipped with the natural action of $H$, is a
locally free sheaf (of rank one) of $A[H]$-modules. The complex
$\eR\Gamma_c(U,f_* A)$ is therefore an object of
$\eD_\text{parf}\left(A[H]\right)$. It is equipped with a Frobenius
endomorphism $F^*$.

For the natural action of $H$ on $f_* f^*\sG_0$, the trace morphism
$f_* f^*\sG_0\to\sG_0$ factors through an isomorphism
\[
  \left(f_* f^*\sG_0\right)_H \to \sG_0 \text{.}
\]
Moreover, one has $f_* f^*\sG_0=f_* A\otimes_A M$ (with diagonal action of
$H$). The isomorphism
\[\xymatrix{
  \sG_0
    & \ar[l]_-\sim (f_* A\otimes_A M)_H = (f_* A\otimes_A M)\otimes_{\Lambda[H]} \Lambda
}\]
derives (\ref{II:4-12}) to an isomorphism
\[
  \eR\Gamma_c(U,\sG) \sim \left(\eR\Gamma_c(U,f_* A)\otimes_A M\right)\otimes_{\Lambda[H]} \Lambda \text{.}
\]

The endomorphism $F^*$ of the left-hand side is deduced from the
endomorphism $F^*$ of $\eR\Gamma_c(U,f_* A)$. The complex of
$\Lambda[H]$-modules $\eR\Gamma_c(U,f_* A)$, equipped with $F^*$, can be
viewed as a complex of $\Lambda[G^+]$-modules, for $G^+=H\times \dN$.
Formulas \ref{II:5-10} and \eqref{II:eq:5-9-2}, amplified by
\ref{II:5-11-1}, therefore give
\[
  \tr\left(F^*,\eR\Gamma_c(U,\sG)\right) = \sum_h' \tr_A^{Z_h}\left(h F^*,\eR\Gamma_c(U,f_* A)\right)\cdot \tr_\Lambda(h,M)
\]
where $\sum'$ is a sum over the conjugacy classes in $H$. Moreover,
\[
  |Z_h|\cdot \tr_A^{Z_h}\left(h F^*,\eR\Gamma_c(U,f_* A)\right) = \tr_A\left((F h^{-1})^*,\eR\Gamma_c(V,A)\right) \text{.}
\]
This formula remains true for $A=\dZ/\ell^m$, with $m$ increasingly large.
Passing to the limit, one finds that
$\tr_A^{Z_h}\left(h F^*,\eR\Gamma_c(U,f_* A)\right)$ is equal to
\[
  \frac{1}{|Z_h|} \sum (-1)^i \tr\left((F h^{-1})^*,\h_c^i(V,\dQ_\ell)\right)
\]
(an element of $\dZ_\ell$), and
\begin{equation*}\tag{5.12.1}\label{II:eq:5-12-1}
\tr\left(F^*,\eR\Gamma_c(U,\sG)\right) = \sum_h' \frac{1}{|Z_h|} \tr\left((F h^{-1})^*,\h_c^\bullet(V,\dQ_\ell)\right)\cdot \tr_\Lambda(h,M) \text{.}
\end{equation*}






\subsection{}\label{II:5-13}

It remains to apply \ref{II:5-4} (note that if $\bar V$ is the nonsingular
projective completion of $V$, all the fixed points of $F h^{-1}$ have
multiplicity one). This can be done without any calculation:
\begin{enumerate}[\indent A.]
  \item If $U^F=\varnothing$,
    $\tr\left(F^*,\eR\Gamma_c(U,\sG)\right) = 0$. \end{enumerate}

Indeed, $F h^{-1}$ has no fixed point on $V$; \ref{II:5-4} shows that
$\sum (-1)^i \tr\left((F h^{-1})^*,\h_c^i(V,\dQ_\ell)\right)=0$, and one
applies \eqref{II:eq:5-12-1}.
\begin{enumerate}[\indent A.]
\setcounter{enumi}{1}
  \item In the general case, let $j:U'=U\setminus U^F\hookrightarrow U$.
    The filtration $j_! j^* \sG_0\subset \sG_0$ of $\sG_0$ induces a
    filtration of $\eR\Gamma_c(U,\sG)$ with successive quotients
    $\eR\Gamma_c(U',\sG)$ and $\eR\Gamma_c(U^F,\sG)$. By
    \eqref{II:eq:4-4-1} and \ref{II:5-1}, one has
    \[
      \tr\left(F^*,\eR\Gamma_c(U,\sG)\right) = \tr\left(F^*,\eR\Gamma_c(U',\sG)\right) + \sum_{x\in U^F} \tr(F^*,\sG_x) \text{,}
    \]
    and, since ${U'}^F=0$, the first term on the right-hand side is $0$.
\end{enumerate}











\section{End of the proof of \ref{II:4-10}}\label{II:6}

Let us prove \ref{II:4-10} under the additional hypotheses that the ring
$\Lambda$ is annihilated by a power of a prime number $\ell\ne p$, and that
$n=1$. The general case will follow: decompose $\Lambda$ into its
$\ell$-primary components (with $\ell$ running through the prime numbers),
and replace $\dF_q$ by $\dF_{q^n}$ and $X_0$ by
$X_0\otimes_{\dF_q}\dF_{q^n}$.

For $X_0$ separated and of finite type over $\dF_q$, and
$\sK_0\in\ob\eD_\text{ctf}(X_0,\Lambda)$, set
\[
  T'(X_0,\sK_0) = \sum_{x\in X^F} \tr(F^*,\sK_x)
\]
and
\[
  T''(X_0,\sK_0) = \tr(F^*,\eR\Gamma_c(X,\sK)) \text{.}
\]
We must prove the identity
\[
  (1;X_0,\sK_0) \qquad T'(X_0,\sK_0) = T''(X_0,\sK_0) \text{.}
\]
We begin by observing that $T'$ and $T''$ (denoted generically by $T$) obey
the same formalism:

(A) If $j:U_0\hookrightarrow X_0$ is an open subscheme of $X_0$, with
closed complement $Y_0$, one has
\[
  T(X_0,\sK_0) = T(U_0,\sK_0|U_0) + T(Y_0,\sK_0|Y_0) \text{.}
\]
This is clear for $T'$. For $T''$, observe that the filtration
$0\subset j_! j^* \sK_0\subset \sK_0$ of $\sK_0$ induces a filtration of
$\eR\Gamma_c(X,\sK)$ with successive quotients $\eR\Gamma_c(U,\sK)$ and
$\eR\Gamma_c(Y,\sK)$, and invoke \eqref{II:eq:4-4-1} (more precisely stated:
$\sK_0$, equipped with the said filtration, defines an object of
$\mathsf{DF}(X,\Lambda)$; applying $\eR\Gamma_c$ gives an object of
$\mathsf{DF}_\text{parf}(\Lambda)$ whose underlying object is
$\eR\Gamma_c(X,\sK)$ and whose associated graded is as stated).

(B) If $\sK_0$ is a bounded complex with constructible flat components,
\[
  T(X_0,\sK_0) = \sum (-1)^i T(X_0,\sK^i) \text{.}
\]
This is clear for $T'$; for $T''$, apply the argument above to the stupid
filtration of $\sK_0$ to obtain
\[
  T''(X_0,\sK_0) = \sum T''(X_0,\sK_0^i[-i]) = \sum (-1)^i T''(X_0,\sK_0^i) \text{.}
\]

(C) For a morphism $f:X_0\to Y_0$, one has
\[
  T''(X_0,\sK_0) = T''(Y_0,\eR f_! \sK_0)
\]
and the same identity holds for $T'$ if the fibers of $f$ at the rational
points of $Y_0$ satisfy
$\left(1;f^{-1}(y),\sK_0|f^{-1}(y)\right)$: use the fact that
\[
  T'(X_0,\sK_0) = \sum_{y\in Y^F} T'\left(f^{-1}(y), \sK_0|f^{-1}(y)\right)
\]
and
\[
  \eR\Gamma_c(Y, \eR f_! \sK) = \eR\Gamma_c(X,\sK) \text{.}
\]

Using this formalism and \ref{II:4-7}, the proof of \ref{II:4-10} is reduced
to the two special cases \ref{II:5-1} and \ref{II:5-2}.

% END INCLUDED components/sga4half_2_en.tex

% BEGIN INCLUDED components/sga4half_3_en.tex
% !TEX root = ../SGA4half_English_Reader_1_to_8.tex

\chapter{\texorpdfstring{$L$}{L}-functions modulo \texorpdfstring{$\ell^n$}{l n} and modulo \texorpdfstring{$p$}{p}}\label{III}

In this exposé, I prove a formula analogous to the main theorem of the report on
the trace formula (this volume; cited as ``Report'' below) for a coefficient
ring $A$ that is torsion prime to $p$, or a field of characteristic $p$.
Essential use will be made of \cite[XVII 5.5]{sga4}.










\section{Symmetric tensors}\label{III:1}





\subsection{}\label{III:1-1}

Let $A$ be a commutative ring. For a list of properties of the functors
$\Gamma^n:(\text{$A$-modules})\to (\text{$A$-modules})$, I refer to
\cite[XVII 5.5.1 and 2]{sga4}. We recall only that for $M$ flat (in fact, only
the case where $M$ is projective of finite type concerns us), one has
\[
  \Gamma^n M = \left(M^{\otimes n}\right)^{S_n} \qquad \text{(symmetric tensors of degree $n$).}
\]






\subsection{}\label{III:1-2}

Let $X$ be an $S$-scheme, let $X^n/S$ be the $n$-th fiber power of $X$, and
let $\pi$ be the projection from $X^n/S$ onto
$(X^n/S)/S_n=\operatorname{Sym}_S^n(X)$ (assumed to exist). If $\sG$ is a
sheaf of $A$-modules on $X$, the external tensor product of $n$ copies of
$\sG$, namely $\sG^{\boxtimes n}$, is an $S_n$-equivariant sheaf on $X^n/S$.
We shall have to use the \emph{external} functor $\Gamma^n$ of
\cite[XVII 5.5.7 to 9]{sga4}. We recall only that for $\sG$ a sheaf of flat
$A$-modules (in fact, only the constructible flat case concerns us), one has
\[
  \Gamma_\text{ext}^n(\sG) = \left(\pi_*\sG^{\boxtimes n}\right)^{S_n} \text{.}
\]






\begin{lemma_}\label{III:1-3}
For $S=\spec(k)$, with $k$ algebraically closed, $X$ a finite disjoint union of
copies of $S$, and $G$ a flat sheaf of $A$-modules on $X$, one has
\[
  \Gamma\left(\operatorname{Sym}_S^n(X),\Gamma_\text{ext}(\sG)\right) = \Gamma^n \Gamma(X,\sG)
\]
\end{lemma_}

Indeed,
\begin{align*}
  \Gamma\left(\operatorname{Sym}_S^n(X),\Gamma_\text{ext}^n(\sG)\right)
    &= \Gamma\left(\operatorname{Sym}_S^n(X),\pi_*\sG^{\boxtimes n}\right)^{S_n} \\
    &= \Gamma\left(X^n/S,\sG^{\boxtimes n}\right)^{S_n}
    = \left(\Gamma(X,\sG)^{\otimes n}\right)^{S_n}
    = \Gamma^n \Gamma(X,\sG) \text{.}
\end{align*}






\subsection{}\label{III:1-4}

We shall have to use the derived functors of the nonadditive functors
$\Gamma^n$ and $\Gamma_\text{ext}^n$.

The theory of such derived functors is due to Dold and Puppe
\cite{do61}. For more complete reminders than those given below, I refer to
\cite[XVII 5.5.3]{sga4}.

Let $\cA$ be an abelian category (or even an additive category in which every
idempotent endomorphism is the projection onto a direct factor). Denote by $N$
the normalization functor: $N:($cosimplicial complexes of objects of
$\cA)\to($differential complexes of objects of $\cA$, with $K^i$ for
$i<0)$. It is an equivalence of categories. Like its inverse, it takes
homotopic morphisms to homotopic morphisms. If $\Gamma$ is a functor
$\cA\to\cB$, and if $K$ is a complex of objects of $\cA$, in nonnegative
degrees ($K^i=0$ for $i<0$), we set $T K=N T N^{-1} K$.





\subsection{Example}\label{III:1-5}

If $K$ is concentrated in $M$ in degree $0$, $T K$ is concentrated in $T M$
in degree $0$.






\subsection{}\label{III:1-6}

If $K$ is a bounded complex, in nonnegative degrees, of projective
$A$-modules of finite type, the complex of $A$-modules $\Gamma^n K$ has the
same properties.

Denote by $\eD_\text{parf}^{\geqslant 0}(A)$ the subcategory of
$\eD_\text{parf}(A)$ (Report, \ref{II:4-3}) consisting of complexes of
tor-dimension $\leqslant 0$. One sees that $\Gamma^n$ induces a functor
\[
  \eL\Gamma^n : \eD_\text{parf}^{\geqslant 0}(A) \to \eD_\text{parf}^{\geqslant 0}(A) \text{.}
\]

For the analogous, but more complicated, definition of
$\eL\Gamma_\text{ext}^n$, I refer to \cite[XVII 5.5.14]{sga4}.






\subsection{}\label{III:1-7}

Let $K$ be a bounded complex of projective $A$-modules of finite type, and let
$u$ be an endomorphism of $K$. We denote by $\det(1-u t,K)$ the formal series
with constant term $1$
\[
  \det(1-u t,K) = \prod_i \det(1-u t,K^i)^{(-1)^i} \text{.}
\]

The verification of the following properties is left to the reader.


\subsubsection{}\label{III:1-7-1}

Let $F$ be a finite filtration stable under $u$, such that the
$\gr_F^p(K^q)$ are projective of finite type. Then
\[
  \det(1-u t,K) = \prod_p \det(1-u t,\gr_F^p(K)) \text{.}
\]


\subsubsection{}\label{III:1-7-2}

If $K$ is shifted $q$ places to the left, one has
\[
  \det(1-u t,K[q]) = \det(1-u t,K)^{(-1)^q} \text{.}
\]






\begin{proposition_}\label{III:1-8}
Let $u$ be an endomorphism of a bounded complex, in nonnegative degrees, of
projective $A$-modules of finite type. One has \
\begin{equation*}\tag{1.8.1}\label{III:eq:1-8-1}
  \det(1-u t,K)^{-1} = \sum_n \tr(u,\Gamma^n K) t^n \text{.}
\end{equation*}
\end{proposition_}

Denote the right-hand side of \eqref{III:eq:1-8-1} by $D(u,K)$.






\begin{lemma_}\label{III:1-9}
Let $u$ be an endomorphism of a short exact sequence of bounded complexes in
degrees $\geqslant 0$ of projective $A$-modules of finite type
\[
  0 \to K' \to K \to K'' \to 0 \text{.}
\]
One has $D(u,K)=D(u,K')\cdot D(u,K'')$.
\end{lemma_}

If $0\to M'\to M\to M''\to 0$ is a split short exact sequence of
$A$-modules, $\Gamma^n M$ has a natural filtration whose successive quotients
are the $\Gamma^i M'\otimes \Gamma^j M''$ ($i+j=n$). Consequently,
$\Gamma^n K$ has a natural filtration whose successive quotients are the
complexes $N(N^{-1} \Gamma^i K'\otimes N^{-1}\Gamma^j K'')$ ($i+j=n$),
canonically homotopic to the simple complexes
$s(\Gamma^i K'\otimes \Gamma^i K'')$ obtained from the double complexes
$\Gamma^i K'\otimes \Gamma^i K''$, and
\[
  \tr(u,\Gamma^n K) = \sum_{i+j=n} \tr(u,\Gamma^i K')\cdot \tr(u,\Gamma^j K'') \text{,}
\]
whence the lemma. From \ref{III:1-9}, one obtains by induction.






\begin{lemma_}\label{III:1-10}
Let $F$ be a finite filtration of $K$ stable under $u$, such that the
$\gr_F^p(K^q)$ are projective of finite type. Then
\[
  D(u,K) = \prod_p D\left(u,\gr_F^p(K)\right) \text{.}
\]
\end{lemma_}






\begin{lemma_}\label{III:1-11}
Let $u$ be an endomorphism of a projective $A$-module of finite type $M$, and
let $M[-q]$ be the complex concentrated in $M$ in degree $q$. One has
\[
  D\left(u,M[-q]\right) = \left(\sum \tr(u,\Gamma^n M) t^n\right)^{(-1)^q} \qquad \text{($q\geqslant 0$.)}
\]
\end{lemma_}






By \ref{III:1-5}, this is true for $q=0$. Proceeding by induction, it remains
to show that
\[
  D\left(u,M[-q]\right) \cdot D\left(u,M[-(q+1)]\right) = 1 \text{.}
\]
This formula follows from \ref{III:1-9}, applied to the complex $K$
concentrated in $M$ in degrees $q$ and $q+1$ (with
$d=\operatorname{id}$), and to its stupid filtration:
\[
  (\cdots 0 \to M \cdots) \to (\cdots M \to M) \to (\cdots M \to 0 \to 0 \cdots) \text{.}
\]
Indeed, $K$ is homotopic to $0$. In the derived category, one therefore has
$\Gamma^n(K)=\Gamma^n(0)$: $0$ for $n>0$ and $A$ for $n=0$, and $D(u,K)=1$.





\subsection{}\label{III:1-12}

Let us prove \ref{III:1-8}. By \ref{III:1-10}, applied to the ``stupid''
filtration of $K$ by the subcomplexes
\[
  0 \to \cdots \to 0 \to K^q \to K^{q+1} \to \cdots \text{,}
\]
one has
\[
  D(u,K) = \prod_q D\left(u,K^q[-q]\right) =_{\text{(\ref{III:1-10})}} \prod_q \left(\sum \tr\left(u,\Gamma^n(K^q)\right) t^n \right)^{(-1)^q} \text{.}
\]
It remains to prove that for $u$ an endomorphism of a projective module of
finite type $M$, one has
\begin{equation*}\tag{1.12.1}\label{III:eq:1-12-1}
  \det\left(1-u t,M\right)^{-1} = \sum \tr(u,\Gamma^n M) t^n \text{.}
\end{equation*}
If $u=0$, both sides are equal to $1$. By adjoining to $M$ a module $N$ on
which $u=0$, and applying \ref{III:1-9}, we reduce to the assumption that $M$
is free. Choose a basis for it. It is then a matter of proving algebraic
identities among the coordinates of $u$. The principle of extension of
algebraic identities reduces us to assuming that $A$ is an algebraically
closed field. The module $M$ then admits a filtration stable under $u$ with
successive quotients of rank $1$, and \ref{III:1-10} reduces us to the case
where $M$ has rank $1$. Formula \eqref{III:eq:1-12-1} then reduces to the
identity
\[
  \frac{1}{1-a t} = \sum a^n t^n \text{.}
\]





\begin{corollary_}\label{III:1-13}
Let $u$ and $u'$ be endomorphisms of bounded complexes $K$ and $K'$ of
projective $A$-modules of finite type. If, in the derived category, $K$,
equipped with $u$, is isomorphic to $K'$, equipped with $u'$, then
\[
  \det(1-u t,K) = \det(1-u' t,K') \text{.}
\]
\end{corollary_}

Applying \ref{III:1-7-2}, we reduce to assuming that $K$ and $K'$ are in
nonnegative degrees. One then observes that the right-hand side of
\ref{III:1-8} has the required invariance property.






\subsection{}\label{III:1-14}

This corollary makes it possible to define $\det(1-u t,K)$ for $u$ an
endomorphism of $K\in \ob\eD_\text{parf}(A)$. This construction is ad hoc; in
fact, for $v$ an automorphism of $L\in\ob\eD_\text{parf}(\Lambda)$, one can
define $\det(v)\in\Lambda^\times$, and $\det(1-u t,K)$ is obtained by taking
$v=1-u t$, $\Lambda=A\llbracket t\rrbracket$, and
$L=K\otimes_A\Lambda$.









\section{The theorem}\label{III:2}






\subsection{}\label{III:2-1}

We retain the notation of (Report, \S\ref{II:1}). For $A$ a commutative
Noetherian torsion ring, $X_0$ a separated scheme of finite type over
$\dF_q$, and $\sK_0\in\ob\eD_\text{ctf}(X_0;A)$, set
\[
  L(X_0,\sK_0) = \prod_{x\in |X_0|} \det\left(1-F_x^* t^{\deg(x)}, \sK\right)^{-1}
\]
(\ref{III:1-7}, \ref{III:1-13}, and Report, \ref{II:1-5}, \ref{II:1-6}).





\begin{theorem_}\label{III:2-2}
In each of the following two cases
\begin{enumerate}[\indent a)]
  \item $A$ is torsion prime to $p$;
  \item $A$ is reduced and of characteristic $p$
\end{enumerate}
one has
\begin{equation*}\tag{2.2.1}\label{III:eq:2-2-1}
  L(X_0,\sK_0) = \det\left(1-F^* t^f, \eR\Gamma_c(X,\sK)\right)^{-1} \text{.}
\end{equation*}
\end{theorem_}






\subsection{Remark}\label{III:2-3}

In case a), the method of Report, \S\ref{II:3} makes it possible to deduce
from the trace formula (Report \ref{II:4-10}) that the logarithmic derivatives
of the two sides of \eqref{III:eq:2-2-1} are equal. Since the ring $A$ is a
torsion ring, this does not suffice to prove \ref{III:2-2}.






\subsection{Remark}\label{III:2-4}

If $A$ is a field, one has
\[
  L(X_0,\sK_0) = \prod_i L\left(X_0,\underline\h^i(\sK_0)\right)^{(-1)^i}
\]
and the spectral sequence
$E_2^{p q} = \h_c^p(X,\underline\h^q(\sK))\Rightarrow \h_c^{p+q}(X,\sK)$ shows
that
\begin{align*}
  \det\left(1-F^* t, \eR\Gamma_c(X,\sK)\right)
    &= \prod_n \det\left(1-F^* t^f,\h_c^n(X,\sK)\right)^{(-1)^n} \\
    &= \prod_{p,q} \det\left(1-F^* t,\h_c^p(X,\underline\h^q(\sK))\right)^{(-1)^{p+1}} \text{.}
\end{align*}
This reduces \ref{III:2-2} (for $A$ a field) to the following special case.


\subsubsection{}\label{III:2-4-1}

If $\sG_0$ is a constructible sheaf of $A$-vector spaces, one has
\[
  L(X_0,\sG_0) = \prod_n \det\left(1-F^* t^f,\h_c^n(X,\sG)\right)^{(-1)^{n+1}} \text{.}
\]






\subsection{Remark}\label{III:2-5}

It is easy to reduce case b) to the case where $A$ is a field (or even a
finite field) of characteristic $p$; cf. \ref{III:4-2}.






\subsection{Remark}\label{III:2-6}

The special case of b) where $A=\dZ/p$, and where $X$ is reduced to the
constant sheaf $\dZ/p$ in degree $0$,
\[
  Z(X_0,t) = \prod_{x\in |X_0|} \left(1-t^{\deg(x)}\right)^{-1}
      \equiv \prod_n \det\left(1-F^* t^f,\h_c^n(X,\dZ/p)\right)^{(-1)^{n+1}} \pmod p
\]
was proved by N. Katz \cite[XXII 3.1]{sga7}. His method, entirely different
from the one followed here, starts from Dwork's $p$-adic theory of
$Z(X_0,t)$.






\subsection{Remark}\label{III:2-7}


Contrary to what I had rashly asserted to friends, in b) one cannot omit the
hypothesis ``$A$ reduced.'' For a counterexample, cf. \ref{III:4-5}.






\subsection{Outline of the proof of \ref{III:2-2}}\label{III:2-8}

A shift in degrees \ref{III:1-7-2} reduces us to assuming that $\sK_0$ is a
bounded complex, in nonnegative degrees, of constructible flat sheaves of
$A$-modules. One then has
$\eR\Gamma_c(X,\sK)\in \ob\eD_\text{parf}^{\geqslant 0}(A)$ (Report, proof of
\ref{II:4-9}, b). We shall use \ref{III:1-8} to expand both sides of
\eqref{III:eq:2-2-1} as power series in $T=t^f$. The equality of the
coefficients of $T^n$ will follow from \cite[XVII 5.5.21]{sga4} and a trace
formula on $\operatorname{Sym}^n(X)$: in case a), the formula (Report,
\ref{II:4-10}), and in case b), a formula to be proved in the next chapter.
It is first necessary to reduce to the case where the symmetric powers of
$X_0$ exist (for example, the case where $X$ is affine), using the
multiplicativity in $X_0$ of both sides of \eqref{III:eq:2-2-1}: for $U_0$ an
open subscheme of $X_0$, with closed complement $Y_0$, one has
$L(X_0,\sK_0)=L(U_0,\sK_0)\cdot L(Y_0,\sK_0)$, and the analogous formula for
the right-hand side follows from \ref{III:1-7-1} by the method of (Report
\S\ref{II:6}A).






\subsection{The case where $X_0$ is finite}\label{III:2-9}

The multiplicativity in $X_0$ of both sides of \eqref{III:eq:2-2-1} reduces
this case to the case where $X_0=\spec(\dF_{q^k})$. Returning to the
definitions, one sees that \eqref{III:eq:2-2-1} is then equivalent to Report
\ref{II:3-4}.






\subsection{The left-hand side}\label{III:2-10}

By \ref{III:1-8}, one has
\[
  L(X_0,\sK_0) = \prod_{x\in |X_0|} \sum_n \tr\left(F_x^*,\Gamma^n\sK\right) t^{m\deg(x)} \text{.}
\]
The coefficient of $T^n$ is therefore the sum, over those formal linear
combinations $\sum_{x\in |X_0|} n_x\cdot x$ ($n_x\geqslant 0$, with almost
all the $n_x$ zero) of elements of $|X_0|$ such that
$\sum n_x [k(x):\dF_q] = n$, of the products
\begin{equation*}\tag{2.10.1}\label{III:eq:2-10-1}
  \prod_x \tr\left(F_x^*,\Gamma^{n_x} \sK\right) \text{.}
\end{equation*}

For $x\in |X_0|$, denote by $(x)$ the $0$-cycle consisting of the
$[k(x):\dF_q]$ points of $X$ above $x$ (an orbit of $F$). The map
$\sum n_x\mapsto \sum n_x\cdot (x)$ identifies the formal linear combinations
above, with $\sum n_x [k(x):\dF_q]=n$, with the effective $0$-cycles of degree
$n$ on $X$ fixed under $F$, or equivalently with the fixed points of $F$ in
$\operatorname{Sym}^n(X)$. The coefficient of $T^n$ in $L(X_0,\sK_0)$ thus
appears as a sum of products \eqref{III:eq:2-10-1}, indexed by
$\operatorname{Sym}^n(X)^F$. More precisely,





\begin{lemma_}\label{III:2-11}
 $L(X_0,\sK_0) = \sum_n \sum_{x\in\operatorname{Sym}^n(X)^F} \tr\left(F_x^*,\eL\Gamma_\textnormal{ext}^n \sK\right) T^n$.
\end{lemma_}

For $Y_0$ an increasingly large finite subscheme of $X_0$, one has
\[
  L(X_0,\sK_0) = \lim L(Y_0,\sK_0) \text{,}
\]
i.e., for every $m$, there is a finite subscheme $Z_0\subset X_0$ such that
if $Y_0\supset Z_0$, then $L(X_0,\sK_0)$ and $L(Y_0,\sK_0|Y_0)$ are congruent
$\mod T^m$.
It suffices for $Z_0$ to contain the closed points of degree $\leqslant m$.
Since the same fact holds for the right-hand side, it suffices to prove
\ref{III:2-11} for $X_0$ finite. One then has
\begin{align*}
  L(X_0,\sK_0)
    &=_\text{\ref{III:2-9}} \det\left(1-F^* T,\Gamma(X,\sK)\right)^{-1} \\
    &=_\text{\ref{III:1-8}} \sum_n \tr\left(F^*,\Gamma^n \Gamma(X,\sK)\right) T^n \\
    &=_\text{\ref{III:1-3}} \sum_n \tr\left(F^*,\Gamma(\operatorname{Sym}^n(X),\Gamma_\text{ext}^n(\sK)\right) T^n \\
    &= \sum_n \sum_{x\in \operatorname{Sym}^n(X)^F} \tr\left(F_x^*,\Gamma_\text{ext}^n(\sK)\right) T^n
\end{align*}
(by the trace formula in the trivial case of a finite set).






\subsection{The right-hand side}\label{III:2-12}

By \ref{III:1-8} and \cite[XVII 5.5.12]{sga4}, one has
\begin{align*}
  \det\left(1-F^* T,\eR\Gamma_c(X,\sK)\right)^{-1}
    &= \sum_n \tr\left(F^*,\eL\Gamma^n\eR\Gamma_c(X,\sK)\right) T^n \\
    &= \sum_n \tr\left(F^*,\eR\Gamma_c(\operatorname{Sym}^n(X),\Gamma_\text{ext}^n\sK)\right) T^n \text{.}
\end{align*}






\subsection{}\label{III:2-13}

The theorem is therefore equivalent to the trace formulas
\[
  \tr\left(F^*,\eR\Gamma_c(\operatorname{Sym}^n(X,\Gamma_\text{ext}^n \sK)\right) = \sum_{x\in\operatorname{Sym}^n(X)^F} \tr\left(F_x^*,\Gamma_\text{ext}^n\sK\right) \text{.}
\]
In case a), they follow from Report \ref{II:4-10}. The analogous formula in
case b) will be proved in \S\ref{III:4}.










\section{Artin--Schreier theory}\label{III:3}






\subsection{}\label{III:3-1}

In what follows, a coherent sheaf on a scheme $S$ will always be regarded as
a sheaf on $S_\text{ét}$. Let $S$ be a scheme of characteristic $p>0$. If
$G$ is a locally constant sheaf of finite-rank $\dZ/p$-vector spaces,
$\sG=G\otimes_{\dZ/p}\sO$ is a locally free coherent sheaf on $S$. We shall
denote by $\Phi$ both the endomorphism $f\mapsto f^p$ of $\sO$ and its tensor
product with the identity of $G$: $\phi:\sG\to\sG$. Artin--Schreier theory
\cite[IX 3.5]{sga4} asserts that the sequence
\[\xymatrix{
  0 \ar[r]
    & \dZ/p \ar[r]
    & \sO \ar[r]^-{\Phi-1}
    & \sO \ar[r]
    & 0
}\]
is exact. Tensoring with $G$ gives an exact sequence
\[\xymatrix{
  0 \ar[r]
    & G \ar[r]
    & \sG \ar[r]^-{\Phi-1}
    & \sG \ar[r]
    & 0
}\]






\subsection{}\label{III:3-2}

Suppose that $S$ is an open subscheme of a Noetherian scheme $\bar S$. Let
$j:S\hookrightarrow \bar S$ be the inclusion morphism, and let $I$ be a
sheaf of ideals defining the closed complement $\bar S\setminus S$. We again
denote by $\Phi$ the extension of $\Phi$ to $j_*\sG$ by functoriality.






\begin{lemma_}\label{III:3-3}
There exist coherent extensions $\bar\sG\subset j_*\sG$ of $\sG$ such that
\begin{equation*}\tag{3.3.1}\label{III:eq:3-3-1}
  \Phi(\bar\sG)\subset I\cdot \bar\sG\text{.}
\end{equation*}
If $\bar\sG$ is such an extension, the sequence
\begin{equation*}\tag{3.3.2}\label{III:eq:3-3-2}
\xymatrix{
  0 \ar[r]
    & j_! G \ar[r]
    & \bar\sG \ar[r]^-{\Phi-1}
    & \bar\sG \ar[r]
    & 0
}
\end{equation*}
is exact.
\end{lemma_}
\begin{enumerate}[a)]
  \item \emph{Existence}: Let $\sG'\subset j_*\sG$ be a coherent extension
    of $\sG$ to $\bar S$. For every $n$, one has
    $\Phi(I^n \sG')\subset I^{p^n}\Phi(\sG')$. Let $n$ be large enough that
    $I^{n(p-1)-1}\Phi(\sG')\subset \sG'$, and set $\bar\sG=I^n\sG'$.
    One has $\Phi(\bar\sG)\subset I^{p n}\Phi(\sG')
    =I\cdot I^n\cdot I^{n(p-1)-1}\Phi(\sG')\subset I\cdot I^n\sG'
    = I\bar\sG$.
  \item \emph{Kernel of $\Phi-1$}: On $S$, the assertion follows from
    \ref{III:3-1}; it remains to verify that a section $x$ of
    $\ker(\Phi-1)$, on $U$ étale over $\bar S$, vanishes in a neighborhood of
    the inverse image of $\bar S\setminus S$. Indeed, since $x=\Phi(x)$, one
    has by \eqref{III:eq:3-3-1} that $x\in I\bar\sG$; moreover,
    $x\in I^n\bar\sG\Rightarrow x=\Phi(x)\in\Phi(I^n\bar\sG)\subset I^{n p}\bar\sG$, and $x$ lies in every $I^n\bar\sG$, hence is zero near
    $\bar S\setminus S$.
  \item \emph{Surjectivity of $\Phi-1$}: We may suppose $\bar S$ affine, and
    write $\sG$ as a quotient of $\sO^n$; $\Phi$ then lifts to a $p$-linear
    map $\widetilde\Phi:\sO^n\to\sO^n$, and it suffices to verify the
    surjectivity of $\widetilde\Phi-1$:
    \[\xymatrix{
      \sO^n \ar[r]^-{\widetilde \Phi-1} \ar[d]
        & \sO^n \ar[d] \\
      \sG \ar[r]^-{\Phi-1}
        & \sG
    }\]
    The sheaf $\sO^n$ is the sheaf of local sections of the group scheme
    $\dG_a^n$, and $\widetilde\Phi-1$ is induced by a homomorphism from
    $\dG_a^n$ to $\dG_a^n$. This homomorphism is étale. To verify that it is
    surjective, it suffices to check this after every base change
    $\spec(k)\to\bar S$ ($k$ algebraically closed), and over an algebraically
    closed $k$, every étale homomorphism between connected groups of the same
    dimension is surjective. Finally, an étale surjective morphism of
    $\bar S$-schemes induces an epimorphism between the corresponding sheaves
    on $\et{\bar S}$.
\end{enumerate}






\subsection{Frobenius and Frobenius}\label{III:3-4}

The absolute Frobenius morphism $F_\text{abs}:S\to S$ is the identity on the
underlying topological space of $S$, and $f\mapsto f^p$ on the structure
sheaf. For $U/S$ étale, one has a canonical isomorphism
$F_\text{abs}^* U=U$: the diagram
\[\xymatrix{
  U \ar[r]^-{F_\text{abs}} \ar[d]
    & S \ar[d] \\
  S \ar[r]^-{F_\text{abs}}
    & S
}\]
is Cartesian. The endomorphism of the ringed site
$(\et{\bar S},\sO)\to(\et{\bar S},\sO)$ defined by $F_\text{abs}$ can
therefore be described as follows:
\begin{center}
  $F_\text{abs}^*$ is the identity on $\et S$; the homomorphism
  $F_\text{abs}^*\sO=\sO\to\sO$ is $\Phi$.
\end{center}
With this description, for every étale sheaf $\sG$ on $S$, the Frobenius
correspondence $F_\text{abs}^*\sG\to\sG$ is the identity.

Giving a $p$-linear morphism of quasi-coherent sheaves
$u:\sM\to\sN$ on $S$ is equivalent to giving a linear morphism
$u':F_\text{abs}^*\sM\to\sN$. In particular, with the notation of
\ref{III:3-5}, the homomorphism $\Phi:\sG\to\sG$ corresponds to
$\Phi':F_\text{abs}^*\sG\to\sG$.






\begin{lemma_}\label{III:3-5}
With the notation of \ref{III:3-4}, the diagram
\[\xymatrix{
  F_\textnormal{abs}^* G \ar[r] \ar[d]
    & G \ar[d] \\
  F_\textnormal{abs}^*\sG \ar[r]^-{\Phi'}
    & \sG
}\]
is commutative.
\end{lemma_}

It suffices to verify this locally. We may therefore suppose $G$ constant, in
which case the lemma is evident.

The same theory applies to the powers of $F_\text{abs}$.






\subsection{}\label{III:3-6}

Suppose now that $\bar S$ is a scheme over $\dF_q$ ($q=p^f$), and in
accordance with the conventions of Report \S\ref{II:1}, write
$S_0,\bar S_0,G_0,\sG_0,\bar\sG_0$ in place of $S,\dots$. Lemma
\ref{III:3-5}, for $F_\text{abs}^F:S_0\to S_0$, gives a commutative diagram
\[\xymatrix{
  F^* G_0 \ar[r] \ar[d]
    & G_0 \ar[d] \\
  F^*\sG_0 \ar[r]^-{(\Phi^f)'} \ar[r]
    & \sG_0
}\qquad\text{, and then}\qquad
\xymatrix{
  F^*(j_! G_0) \ar[r] \ar[d]
    & j_!G_0 \ar[d] \\
  F^*\bar\sG_0 \ar[r]^-{(\Phi^f)'}
    & \sG_0
}\text{,}
\]
and then, by extension of scalars from $\dF_q$ to $\dF$ and passage to
cohomology,
\[\xymatrix{
  F^*(j_! G) \ar[r]^-{F_*} \ar[d]
    & j_!G \ar[d] \\
  F^*\bar\sG \ar[r]
    & \bar\sG
}\qquad\text{and}\qquad
\xymatrix{
  \h^\bullet(\bar S,j_! G) \ar[r]^-{F^*} \ar[d]
    & \h^\bullet(\bar S,j_!G) \ar[d] \\
  \h^\bullet(\bar S,\bar\sG) \ar[r]
    & \h^\bullet(\bar S,\bar\sG)
}\text{,}
\]
where the second row is obtained by extension of scalars from $\dF_q$ to
$\dF$ of the $\dF_q$-linear endomorphism $\Phi^f$ of
$\h^\bullet(\bar S_0,\bar G_0)$.

Artin--Schreier theory provides a long exact sequence
\[\xymatrix{
  \cdots \ar[r]
    & \h^i(\bar S,j_! G) \ar[r]
    & \h^i(\bar S,\bar\sG) \ar[r]^-{\Phi-1}
    & \h^i(\bar S,\bar\sG) \ar[r]
    & \cdots
}\]
where $\Phi$ is the $p$-linear extension to
$\h^\bullet(\bar S,\bar \sG)= \h^\bullet(\bar S_0,\bar\sG_0)\otimes_{\dF_q}\dF$
of the $p$-linear endomorphism $\Phi$ of
$\h^\bullet(\bar S_0,\bar\sG_0)$.






\subsection{}\label{III:3-7}

If $\bar S_0$ is proper over $\dF_q$, the $\h^i(\bar S,\bar\sG)$ are
finite-dimensional and $\Phi-1$ is surjective (cf. \ref{III:3-6}b). Thus
\[
  \h_c^i(S,G) = \ker\left(\Phi-1 : \h^i(\bar S,\bar G) \to \h^i(\bar S,\bar\sG)\right) \text{,}
\]
and $F^*$ is induced by the $\dF$-linear extension to
$\h^i(\bar S,\bar\sG) = \h^i(\bar S_0,\bar\sG_0)\otimes_{\dF_q}\dF$ of the
$\dF_q$-linear endomorphism $\Phi^f$ of $\h^i(\bar S_0m\bar\sG_0)$.

Applying Lemma \ref{III:3-8} below, one deduces the identity
\begin{equation*}\tag{3.7.1}\label{III:eq:3-7-1}
  \tr\left(F^*,\h_c^i(S,G)\right) = \tr\left(\Phi^f,\h^i(\bar S_0,\bar\sG_0)\right) \text{.}
\end{equation*}






\begin{lemma_}\label{III:3-8}
Let $\Phi$ be a $p$-linear endomorphism of a vector space $V$ over
$\dF_q$: $\phi(\lambda x) = \lambda^p\phi(x)$. Set $F=\Phi^f$ ($F$ is
linear), and again denote by $\Phi$ and $F$ the respective $p$-linear
extensions of $\Phi$ and $F$ to $V\otimes_{\dF_q}\dF$. Then $F$ stabilizes
$\ker(\Phi-1)\subset V\otimes_{\dF_q}\dF$ and
\[
  \tr\left(F,\ker(\Phi-1)\right) = \tr(F,V) = \tr\left(F,V\otimes_{\dF_q}\dF\right) \text{.}
\]
\end{lemma_}

On $V\otimes_{\dF_q}\dF$, one has $F\Phi = \Phi F$: the $p$-linear
endomorphism $F\Phi-\Phi F$ vanishes on $V$, hence everywhere. Write
$V=V'+V''$, with $F$ invertible on $V'$ and nilpotent on $V''$. The theory
of $p$-linear endomorphisms ensures that
\[
  \ker(\Phi-1)\otimes_{\dF_p}\dF \iso V'\otimes_{\dF_q}\dF \text{.}
\]
Since $\tr(F,V'') = 0$, one therefore has
\begin{align*}
  \tr\left(F,\ker(\Phi-1)\right)
    &= \tr\left(F,\ker(\Phi-1)\otimes_{\dF_q}\dF\right) = \tr\left(F,V'\otimes_{\dF_q}\dF\right) \\
    &= \tr\left(F,V\otimes_{\dF_q}\dF\right) = \tr(F,V) \text{.}
\end{align*}











\section{Trace formula modulo \texorpdfstring{$p$}{p}}\label{III:4}






\begin{theorem_}\label{III:4-1}
Let $X_0$ be a separated scheme of finite type over $\dF_q$, let $A$ be a
commutative reduced Noetherian ring of characteristic $p$, and let
$\sK_0\in\ob\eD_\textnormal{ctf}(X_0,A)$. One has
\begin{equation*}\tag{4.1.1}\label{III:eq:4-1-1}
  \sum_{x\in X^F} \tr\left(F_x^*,\sK_x\right) = \tr\left(F^*,\eR\Gamma_c(X,\sK)\right) \text{.}
\end{equation*}
\end{theorem_}






\subsection{Reduction to the case where \texorpdfstring{$A$}{A} is a finite field}\label{III:4-2}

%NOTE: in this paragraph and others, (3.1.1) seems to be referring to (4.1.1)
A standard passage-to-the-limit argument reduces us to the case where $A$ is
of finite type over $\dZ/p$. For every maximal ideal $\fm$ of $A$, the image
of each side of \eqref{III:eq:4-1-1} in $A/\fm$ is given by the same formula,
with $\sK_0$ replaced by
$\sK_0\lotimes_A A/\fm\in\ob\eD_\text{ctf}(X_0,A/\fm)$ (for the right-hand
side, cf. Report \ref{II:4-12}). Since the $A/\fm$ are finite, and the
intersection of the maximal ideals of $A$ is reduced to $0$, one obtains the
stated reduction.






\subsection{Reduction to the case where \texorpdfstring{$A=\dZ/p$}{A=Z/p}}\label{III:4-3}

Suppose that $A$ is a finite field, and denote by $T_A'(\sK_0)$ and
$T_A''(\sK_0)$ the two sides of \eqref{III:eq:4-1-1}. If, by restriction of
scalars, $\sK_0$ is regarded as a complex of sheaves of $\dZ/p$-modules, one
has
\[
  T_{\dZ/p}'(\sK_0) = \tr_{A/(\dZ/p)} T_A'(\sK_0)
\]
and likewise for $T''$. If \eqref{III:eq:4-1-1} holds for $A=\dZ/p$, one
therefore has
\begin{equation*}\tag{4.3.1}\label{III:eq:4-3-1}
  \tr_{A/(\dZ/p)} T_A'(\sK_0) = \tr_{A/(\dZ/p)} T_A''(\sK_0) \text{.}
\end{equation*}
I claim that, for every $\lambda\in A$, one even has
\begin{equation*}\tag{4.3.2}\label{III:eq:4-3-2}
  \tr_{A/(\dZ/p)}\left(\lambda\cdot T_A'(\sK_0)\right) = \tr_{A/(\dZ/p)}\left(\lambda\cdot T_A''(\sK_0)\right) \text{,}
\end{equation*}
and hence $T_A'(\sK_0) = T_A''(\sK_0)$. This is clear for $\lambda=0$. For
$\lambda$ invertible, denote by $A(\lambda)$ the inverse image on $X$ of the
free rank-one sheaf of $A$-modules on $\spec(\dF_q)$ for which
$F^*=\lambda$. Formula \eqref{III:eq:4-3-2} is precisely
\eqref{III:eq:4-3-1} for $A(\lambda)\otimes_A \sK_0$.






% originally theorem 4.1bis (i.e. theorem 4.1 again)
\begin{theorem_}\label{III:4-1bis}
Let $\bar S_0$ be a smooth projective curve, let
$j:S_0\hookrightarrow \bar S_0$ be a dense open subscheme of $\bar S_0$,
and let $G_0$ be a locally constant sheaf of finite-rank $\dZ/p$-vector
spaces on $S_0$. One has
\begin{equation*}\tag{4.4.1}\label{III:eq:4-4-1}
  \sum_{x\in S^F} \tr\left(F_x^*,G_x\right) = \sum_i (-1)^i \tr\left(F^*,\h_c^i(S,G)\right) \text{.}
\end{equation*}
\end{theorem_}
\begin{proof}
Apply \ref{III:3-3} to $(S_0,\bar S_0,G_0)$. The coherent extension
$\bar\sG_0$ of $\sG_0$ whose existence is guaranteed by \ref{III:3-3} is
locally free because it is torsion-free (we have
$\bar\sG_0\subset j_*\sG_0$) and $\bar S_0$ is regular of dimension one. The
$q$-linear endomorphism $\Phi^f$ of $\bar\sG_0$ can be interpreted as a
morphism
\[
  u : F^*\bar\sG_0 \to \bar\sG_0 \text{.}
\]
It factors through $I\bar G$, where $I$ defines the closed subset
$\bar S\setminus S$. Let $v:j_!G_0 \to \bar\sG_0$ be the unique extension of
the natural map $G_0 \to \sG_0 = G_0\otimes_{\dZ/p}\sO$. The diagram
\begin{equation*}\tag{4.4.2}\label{III:eq:4-4-2}
\xymatrix{
  F^* j_! G_0 \ar[r]^-{F^*} \ar[d]^-{F^* v}
    & j_! G_0 \ar[d]^-v \\
  F^*\bar \sG_0 \ar[r]^-u
    & \bar\sG_0
}
\end{equation*}
Finally, denote by $u F^*$ the composite map
\[\xymatrix{
  \h^\bullet(\bar S,\bar\sG) \ar[r]^-{F^*}
    & \h^\bullet(\bar S,F^*\bar\sG) \ar[r]^-u
    & \h^\bullet(\bar S,\bar\sG) \text{,}
}\]
one has by \eqref{III:eq:3-7-1}
\begin{equation*}\tag{4.4.3}\label{III:eq:4-4-3}
  \tr\left(F^*,\h_c^i(S,G)\right) = \tr\left(u F^*,\h^i(\bar S,\bar\sG)\right) \text{.}
\end{equation*}

%NOTE: trace down reference for Woods Hole seminar
For every closed point $i_x:x\hookrightarrow \bar S$ of $\bar S$, set
$\bar\sG_x = i_x^*\bar\sG$. The Lefschetz trace formula in coherent
cohomology (``Woods Hole fixed point formula''), applied to $F$ and $u$, says
that
\begin{equation*}\tag{4.4.4}\label{III:eq:4-4-4}
  \sum_i (-1)^i \tr\left(uF^*, \h^i(\bar S,\bar\sG)\right) = \sum_{x\in \bar S^F} \frac{\tr(u_x,\sG_x)}{\det(1-d F_x)} \text{.}
\end{equation*}
On the right-hand side,
\begin{enumerate}[\indent a)]
  \item since $d F=0$, the denominators are equal to 1;
  \item for $x\in \bar S^F\setminus S^F$, the endomorphism $u_x$ of
    $\bar\sG_x$ is zero, and therefore so is $\tr(u_x,\bar\sG_x)$;
  \item for $x\in S^F$, \eqref{III:eq:4-4-4} gives a contribution
    \[
      \tr(u_x,\bar\sG_x) = \tr(F_x^*,G_x) \text{.}
    \]
    Identity \eqref{III:eq:4-4-1} therefore follows from
    \eqref{III:eq:4-4-3} and \eqref{III:eq:4-4-4}.
\end{enumerate}
\end{proof}






\subsection{A counterexample}\label{III:4-5}

An example shows that the hypothesis that $A$ be reduced cannot be omitted
from \ref{III:4-1}.

Let $Y_0$ be the elliptic curve over $\dF_2$ for which the eigenvalues of
Frobenius are $\pm \sqrt{-2}$. It has $3$ rational points and is
supersingular. The involution $\sigma:x\mapsto -x$ therefore has only one
fixed point, the point $0$. The quotient of $Y_0$ by $\sigma$ is a projective
line; the quotient of $Y_0\setminus 0$ by $\sigma$ is therefore an affine
line $\dF_0^1$, and $Y_0\setminus 0$ is an unramified double covering of
$\dF_0^1$:
\[
  \pi:Y_0\setminus 0 \to \dA_0^1 \qquad \text{, $\pi=\pi\sigma$.}
\]

The sheaf $\pi_*\dZ/2$ is a free module of rank one over the algebra $A$ of
the two-element group $\{1,\sigma\}$ over $\dZ/2$. This algebra is isomorphic
to the algebra of numbers over $\dZ/2$.

The sheaf of $A$-modules $\pi_*\dZ/2$ is our example:
\begin{enumerate}[\indent a)]
  \item One has
    $\h_c^\bullet(\dA^1,\pi_*\dZ/2) = \h_c^\bullet(Y\setminus 0,\dZ/2)$.
    Since $Y_0$ is supersingular, one has $\h^i(Y,\dZ/2)=0$ for $i>0$, and
    the long exact sequence
    \[
      \cdots \to \h_c^i(Y\setminus 0,\dZ/2) \to \h^i(Y,\dZ/2) \to \h^i(\{0\},\dZ/2) \to \cdots
    \]
    shows that $\h_c^\bullet(Y\setminus 0,\dZ/2)=0$. Thus
    $\h_c^\bullet(\dA^1,\pi_*\dZ/2)=0$ and
    \[
      \tr_A\left(F^*,\eR\Gamma_c(\dA^1,\pi_*\dZ/2)\right) = 0 \text{.}
    \]
  \item Since $Y_0\setminus 0$ has two rational points, at one of the
    rational points of $\dA_1^0$, the Frobenius substitution is the identity,
    and at the other, it is $\sigma$. Thus
    \[
      \sum_{x\in {\dA^1}^F} \tr\left(F^*,\pi_*\dZ/2\right) = 1+\sigma \text{,}
    \]
    and $1+\sigma\ne 0$.
\end{enumerate}

The equation of the covering $\pi$ is $y^2-y=x^3$.

% END INCLUDED components/sga4half_3_en.tex

% BEGIN INCLUDED components/sga4half_4_en.tex
% !TEX root = ../SGA4half_English_Reader_1_to_8.tex

\chapter{The cohomology class associated with a cycle}\label{IV}
\blfootnote{by A.\ Grothendieck, written up by P.\ Deligne}





This exposé is inspired by notes of Grothendieck that formed a preliminary
version of \cite[IV]{sga5}. It defines the cohomology class of a cycle $X$ in
a separated smooth scheme of finite type over a field, and proves that
intersection corresponds to the cup product.

Chapter \ref{IV:1} contains some general formalities. In Chapter \ref{IV:2},
the class of a cycle is defined in several situations more general than the
one just mentioned. The principal compatibility not considered is that
between the direct image of a cycle and the trace morphism in cohomology. In
Chapter \ref{IV:3}, this formalism is used to deduce the Lefschetz trace
formula for an endomorphism with isolated fixed points of a proper smooth
scheme over an algebraically closed $k$---and for a Frobenius endomorphism of
a curve.

We adopt the following conventions:
\begin{enumerate}[\indent 1)]
  \item ``scheme'' means a separated Noetherian scheme (this is largely a
    hypothesis of convenience).
  \item In Chapter \ref{IV:2}, an integer $n$ is fixed, and $n$ is invertible
    on all schemes under consideration.
  \item In Chapter \ref{IV:3}, a prime number $\ell$ is fixed, and $\ell$ is
    invertible on all schemes under consideration. The cohomology used is
    always $\ell$-adic cohomology:
    \[
      \h^\bullet(X) = \dQ_\ell\otimes_{\dZ_\ell} \varprojlim \h^\bullet(X,\dZ/\ell^n) \text{.}
    \]
\end{enumerate}
Cohomology classes of cycles, trace morphisms, \ldots{} are defined by passage
to the limit from the case of finite coefficients $\dZ/\ell^n$ (cf.
\cite[VI]{sga5}).




















\section{Cohomology with support and cup products}\label{IV:1}

This section contains reminders from general topology, which the reader is
invited to consult only as the need arises.










\subsection{\texorpdfstring{$H^1$}{H1} and torsors}\label{IV:1-1}






\subsubsection{}\label{IV:1-1-1}

Let $\sF$ be an abelian sheaf on a site $X$. It is known that
$\h^1(X,\sF)$ classifies $\sF$-torsors on $X$. We normalize ($=$ choose the
sign of) the isomorphism $($set of isomorphism classes of
$\sF$-torsors$)\to \h^1(X,\sF)$ in such a way that for every exact sequence
$0\to \sF\xrightarrow\alpha\sG\xrightarrow\beta\sH\to 0$ and every
$h\in\h^0(X,\sH)$, the $\sF$-torsor $\beta^{-1}(h)\subset \sG$, on which
$\sF$ acts by $(f,x)\mapsto \alpha(f)+x$, has class $\partial h$.






\subsubsection{}\label{IV:1-1-2}

Let $\sP$ be an $\sF$-torsor. If $(U_i)$ is an open covering of $X$, and
$p_i$ a section of $\sP$ on $U_i$, one associates to $\sP$ the \v{C}ech
cocycle
\[
  p_{i j} = p_j - p_i \qquad \text{($p_{i j}\in \h^0(U_i\times U_j,\sF)$).}
\]
If, according to the usual rule, the map
$\operatorname{\check H}^\bullet(X,\sF)\to \h^\bullet(X,\sF)$ is defined so
as to be a morphism of $\delta$-functors, the image of
$(p_{i j})\in \operatorname{\check H}^1(X,\sF)$ in $\h^1(X,\sF)$ is the class
of $\sP$, normalized as in \ref{IV:1-1-1}.






\subsubsection{}\label{IV:1-1-3}

The definition of $\h^\bullet(X,\sF)$ is as follows: for $\sF^\bullet$ a
resolution of $\sF$ with acyclic components,
$\h^i(X,\sF) = \h^i\Gamma(X,\sF^\bullet)$. The $\delta$-functor structure is
obtained by associating to a short exact sequence of sheaves a short exact
sequence of resolutions that remains exact after application of the functor
$\Gamma$. If $\sF^\bullet$ is a resolution of $\sF$, the connecting
homomorphism $\partial$ associated with
\[\xymatrix{
  0 \ar[r]
    & \sF \ar[r]
    & \sF^0 \ar[r]^-d
    & \ker(d) \ar[r]
    & 0
}\]
induces the negative of the defining isomorphism
$\h^1(X,\sF) = \Gamma\left(X,\ker(d)\right)/d \Gamma(X,\sF^0)$.






\subsubsection{}\label{IV:1-1-4}

Let $U$ be an open part of $X$ (a subsheaf of the final sheaf), and let $D$
be the ``closed complement.'' It is known that $\h_D^1(X,\sF)$ classifies
the $\sF$-torsors on $X$ trivialized on $U$. For every short exact sequence
$0\to \sF\xrightarrow\alpha \sG\xrightarrow\beta\sH\to 0$, and every
section with support in $D$, $h\in\h_D^0(X,\sH)$, the torsor
$\beta^{-1}(h)$, trivialized on $U$ by the section $0$, has class
$\partial h$.

The long exact sequence of cohomology with support
\[\xymatrix{
  \cdots \ar[r]^-\partial
    & \h_D^i(X,\sF) \ar[r]
    & \h^i(X,\sF) \ar[r]
    & \h^i(U,\sF) \ar[r]^-\partial
    & \cdots
}\]
is defined from the sequence of functors
\[\xymatrix{
  0 \ar[r]
    & \Gamma_D \ar[r]
    & \Gamma \ar[r]
    & \Gamma(U,-) \ar[r]
    & 0
}\]
(exact on injective sheaves). For every section $f\in \h^0(U,\sF)$,
$\partial f\in \h_D^1(X,\sF)$ is the class of the trivial torsor $\sF$,
trivialized on $U$ by the section $f$.






\subsubsection{}\label{IV:1-1-5}

Let $j:U\hookrightarrow X$. If $\sF$ injects into $j_* j^* \sF$, the exact
sequence
\[\xymatrix{
  0 \ar[r]
  & \sF \ar[r]
  & j_* j^* \sF \ar[r]
  & j_* j^* \sF / \sF \ar[r]
  & 0
}\]
gives
$\partial:\h^0(j_* j^*\sF/\sF) = \h_D^0(j_* j^* \sF/\sF) \to \h_D^1(X,\sF)$.
The composite map
\[\xymatrix{
  \h^0(U,\sF) = \h^0(X,j_* j^*\sF) \ar[r]
    & \h^0(X,j_* j^*\sF / \sF) \ar[r]^-\partial
    & \h_D^1(X,\sF)
}\]
is the negative of the map considered in \ref{IV:1-1-4}.






\subsubsection{}\label{IV:1-1-6}

Recall that if $\sL$ is an invertible sheaf on a scheme $X$, the
$\dG_m$-torsor corresponding to it is the sheaf
$\operatorname{Isom}(\sO,\sL)$, on which $\dG_m$ acts by
$(\lambda,f)\mapsto f\circ (\lambda\cdot) = \lambda f$. Recall also that if
$D$ is a Cartier divisor on $X$, and $j:U\hookrightarrow X$ is the inclusion
of the complementary open subscheme, the invertible sheaf $\sO(D)$ is the
subsheaf of $j_*\sO_U$ consisting of local sections $s$ such that $s f$ lies
in $\sO_X$, where $f$ is a local equation for $D$.











\subsection{Cup products}\label{IV:1-2}

In this subsection, we develop some observations about cup products in
cohomology with support that will be used again in [\nameref{V}] to relate
Poincaré duality for curves and self-duality of the Jacobian.

\subsubsection{}\label{IV:1-2-1}

Let $X$ be a site, $Y$ a closed part of $X$, and $\sF,\sG$ two abelian
sheaves on $X$. For example: $X$ a scheme, $Y$ a closed subscheme, and
$\sF,\sG$ sheaves on $\et X$. Define a product
\begin{equation*}\tag{1.2.1.1}\label{IV:eq:1-2-1-1}
  \Gamma_Y(X,\sF)\otimes \Gamma(Y,\sG) \to \Gamma_Y(X,\sF\otimes\sG) \text{.}
\end{equation*}
The product of a section $s$ of $\sF$ with support in $Y$ by a section $t$ of
$\sG$ on $Y$ is obtained as follows: locally on $X$, $t$ is the restriction
to $Y$ of a section $t'$ of $\sG$ on $X$, and one forms the product
$s\otimes t'$. It has support in $Y$ and does not depend on the choice of
$t'$, which justifies the definition and permits it to be globalized.

Let $i:Y\hookrightarrow X$ be the inclusion morphism. The local analogue of
\eqref{IV:eq:1-2-1-1} is the product
\begin{equation*}\tag{1.2.1.2}\label{IV:eq:1-2-1-2}
  i^!\sF\otimes i^*\sG \to i^!(\sF\otimes\sG) \text{,}
\end{equation*}
from which \eqref{IV:eq:1-2-1-1} follows by applying $\Gamma(Y,-)$.






\subsubsection{}\label{IV:1-2-2}

Let us derive these arrows. Let $\sK$, $\sL$, and $\sM$ be in the derived
category, with a bilinear map $\sK\lotimes\sL\to\sM$. Deriving
\eqref{IV:eq:1-2-1-1} gives
\begin{equation*}\tag{1.2.2.1}\label{IV:eq:1-2-2-1}
  \eR\Gamma_Y(X,\sK) \lotimes \eR\Gamma(Y,\sL) \to \eR\Gamma_Y(X,\sM)
\end{equation*}
inducing
\begin{equation*}\tag{1.2.2.2}\label{IV:eq:1-2-2-2}
  \h_Y^i(X,\sK)\otimes \h^j(Y,\sL) \to \h_Y^{i+j}(X,\sM)
\end{equation*}
(the cup product). The local arrow \eqref{IV:eq:1-2-1-2} gives
\begin{equation*}\tag{1.2.2.3}\label{IV:eq:1-2-2-3}
  \eR i^! \sK\lotimes i^*\sL \to \eR i^!\sM \text{,}
\end{equation*}
from which \eqref{IV:eq:1-2-2-1} follows by applying $\eR\Gamma(Y,-)$.





\subsubsection{}\label{IV:1-2-3}

Above, I passed over in silence the conflicts caused by using right-derived
functors ($\eR\Gamma$) and left-derived functors ($\lotimes$) in the same
formula.
\begin{enumerate}[\indent a)]
  \item In the next subsection, all the right-derived functors under
    consideration will have finite cohomological dimension. This makes it
    possible to work systematically in the derived categories $\eD^-$. To
    have complexes that are both flat and flasque, one uses canonical
    flasque resolutions as in \cite[XVII]{sga4}.
  \item For a more general theory, one can no longer interpret bilinear maps
    from $\sK$ and $\sL$ to $\sM$ as morphisms from
    $\sK\lotimes\sL$ to $\sM$. For example,
    $\eR\Gamma_Y(X,\sK)\lotimes\eR\Gamma(Y,\sL)$ is not defined if both
    factors lie in $\eD^+$. One solution is to work in $\eD^+$, define
    \[
      \operatorname{Bil}(\sK,\sL;\sM) = \varinjlim \hom(\sK'\otimes\sL',\sM') \text{,}
    \]
    where the limit is taken over the quasi-isomorphisms $\sK'\iso \sK$,
    $\sL'\iso \sL$, $\sM\iso \sM'$ ($\sK'$, $\sL'$, $\sM'$ bounded
    below), where $\hom$ means ``morphism of complexes, up to homotopy,'' and
    systematically use such bilinear maps without ever mentioning
    $\otimes$.
\end{enumerate}





\subsubsection{Second theme}\label{IV:1-2-4}

Let $U$ be an open subscheme of $X$, and let $j:U\hookrightarrow X$ be the
inclusion morphism. For $\sK$ in the derived category on $U$, set
$\eR\Gamma_!(U,\sK)=\eR\Gamma(X,j_!\sK)$. For $\sK$, $\sL$, and $\sM$ on
$U$, with a bilinear map $\sK\lotimes\sL\to\sM$, we wish to define
\begin{equation*}\tag{1.2.4.1}\label{IV:eq:1-2-4-1}
  \eR\Gamma(U,\sK)\lotimes\eR\Gamma_!(U,\sL)\to\eR\Gamma_!(U,\sM) \text{.}
\end{equation*}
At the level of sheaves and their global sections, such a product follows
from the isomorphism
$j_*\sF\otimes j_!\sG\xleftarrow\sim j_!(\sF\otimes \sG)$, but care must be
taken because $\eR\Gamma_!$ is not in general the derived functor of
$\Gamma(X,j_!-)$.

We begin by defining
\[\xymatrix{
  \eR j_* \sK\lotimes j_! \sL
    & \ar[l]_-\sim j_!(\sK\lotimes \sL) \ar[r]
    & j_!\sM \text{.}
}\]
Applying $\eR\Gamma(X,-)$ gives
\[\xymatrix{
  \eR\Gamma(U,\sK)\lotimes\eR\Gamma_!(U,\sL) = \eR\Gamma(X,\eR j_*\sK)\lotimes \eR\Gamma(X,j_!\sL) \ar[r]
    & \eR\Gamma(X,j_! \sM) \text{.}
}\]

Here again, the conflict between left and right will be resolved in the next
subsection by working in $\eD^-$.






\subsubsection{Coda}\label{IV:1-2-5}

Let $j:U\hookrightarrow X$ be an open part, and let $i:Y\hookrightarrow U$ be
a closed part of $U$. Let $\bar Y$ be a closed subset of $X$ such that
$\bar Y\cap U=Y$ (for example, the complement of $U\setminus Y$). Let $\sK$
be on $U$. Set $\eR\Gamma_{Y!}(U,\sK)=\eR\Gamma_!(Y,\eR i^!\sK)$, where
$\eR\Gamma_!$ is relative to the inclusion of $Y$ into $\bar Y$. For every
open subset $V$ of $U$ containing $Y$, $\eR\Gamma_{Y!}(U,\sK)$ maps to
$\eR\Gamma_!(V,\sK)$: if $i$ also denotes the inclusion of $\bar Y$ into
$X$, $j$ that of $Y$ into $\bar Y$, and $k$ that of $V$ into $X$, one has
$\eR i^!\sK = j^*\eR i^! k_!(k^*\sK)$, whence a morphism
$j_!\eR i^!\sK \to \eR i^! k_!(k^*\sK)$. Applying
$\eR\Gamma(\bar Y,-)$ to it gives
$\eR\Gamma_{Y!}(U,\sK)\to \eR\Gamma_{\bar Y}(k_! k^* \sK) \to \eR\Gamma(k_! k^*\sK) = \eR\Gamma_!(V,\sK)$.

For $\sK,\sL,\sM$ on $U$, with a bilinear map
$\sK\lotimes\sL\to\sM$, we wish to define
\begin{equation*}\tag{1.2.5.1}\label{IV:eq:1-2-5-1}
  \h^n(U,\sK)\otimes \h_!^m(Y,\sL) \to \h_{Y!}^{n+m}(U,\sM)
\end{equation*}
(the last group itself mapping to $\h_!^{n+m}(V,\sM)$). In the derived
category, this means defining
\begin{equation*}\tag{1.2.5.2}\label{IV:eq:1-2-5-2}
  \eR\Gamma_Y(U,\sK) \lotimes \eR\Gamma_!(Y,\sL) \to \eR\Gamma_{Y!}(U,\sM) \text{.}
\end{equation*}
We identify $\eR\Gamma_Y$ with $\eR\Gamma(Y,\eR i^!\sK)$. The desired
product is then of type \eqref{IV:eq:1-2-4-1}, relative to the local product
\eqref{IV:eq:1-2-2-3} on $Y$:
$\eR i^! \sK\lotimes\sL \to \eR i^! \sM$.






\subsubsection{}\label{IV:1-2-6}

Above, we developed the absolute formalism. There is a parallel relative
formalism, with $\Gamma$ replaced by $f_*$ for $f$ a morphism $X\to S$.










\subsection{The Koszul rule}\label{IV:1-3}

Let $A$ be a commutative ring, and let $(V_i)_{i\in I}$ be a finite family of
graded (or $\dZ/2$-graded) $A$-modules. Recall the definition of the graded
tensor product $\bigotimes_{i\in I} V_i$, in the sense of the Koszul rule
(cf. \cite[XVII 1.1]{sga4}). For each total ordering $a$ on $I$, we shall
define a module $V(a)$. We shall also define a transitive system of
isomorphisms $\varphi_{a b}:V(b)\iso V(a)$, and the graded tensor product of
the $V_i$ will be the ``common value'' $\varprojlim V(a)$ of this system of
modules. We take
\begin{enumerate}[\indent a)]
  \item $V(a) = \bigotimes_{i\in I} |V_i|$ (the ordinary tensor product of
    the underlying ungraded modules of the $V_i$).
  \item if the $x_i\in V_i$ are homogeneous, we take
    $\varphi_{a b}(\otimes x_i) = (-1)^N \otimes x_i$, where $N$ is the sum
    of the $\deg(x_i)\deg(x_j)$ over the pairs $(i,j)$ such that
    $i<_a j$ and $i>_b j$.
\end{enumerate}






\subsubsection*{\hypertarget{sga4half.e04.example.1-3-1}{Example 1}}

Take $I=\{1,2\}$, and let $a$ be the ordering in which $1<2$, and $b$ the
ordering in which $1>2$. Let $v_i\in V_i$ be homogeneous, and denote by
$v_1\otimes v_2$ (resp. $v_2\otimes v_1$) the image in the graded tensor
product of the product of the $v_i$ in $V(a)$ (resp. $V(b)$). One has
\[
  v_1\otimes v_2 = (-1)^{\deg(v_1)\deg(v_2)} v_2\otimes v_1
\]
(the Koszul rule).






\subsubsection*{\hypertarget{sga4half.e04.example.1-3-2}{Example 2}}

If all the $V_i$ have degree $1$, there is a canonical isomorphism relating
the module underlying the graded tensor product and the ordinary tensor
product of the modules $|V_i|$ underlying the $V_i$:
\[
  \left|\bigotimes V_i\right| \simeq \bigotimes |V_i|\otimes_{\dZ}\bigwedge^{|I|}\dZ^I \text{.}
\]






\subsubsection*{\hypertarget{sga4half.e04.example.1-3-3}{Example 3}}

For $A$ a field, and $X_i$ a finite family of spaces, the K\"unneth formula
is written
$\h^\bullet\left(\prod X_i,A\right) = \bigotimes \h^\bullet(X_i,A)$.

Let $V^\vee$ be the graded dual of $V$. The canonical map
\begin{equation*}\tag{1.3.1}\label{IV:eq:1-3-1}
  V^\vee\otimes V = V\otimes V^\vee \to A
\end{equation*}
is $v'\otimes v\to v'(v)$.

Suppose now that $A$ is a field, and consider only finite-dimensional vector
spaces. The canonical isomorphism
\begin{equation*}\tag{1.3.2}\label{IV:eq:1-3-2}
  W\otimes V^\vee \to \hom(V,W)
\end{equation*}
is $w\otimes v'\mapsto \left(v\mapsto w\cdot v'(v)\right)$. Via this
isomorphism, the composition
$\hom(Y,Z)\otimes \hom(X,Y)\to \hom(X,Z)$ is identified with the morphism
induced by \eqref{IV:eq:1-3-1}:
$Z\otimes Y^\vee\otimes Y\otimes X^\vee \to Z\otimes X^\vee$.

The trace of $f:V\to V$ (zero for homogeneous $f$ of degree $\ne 0$) is
the image of $f$ under
\begin{equation*}\tag{1.3.4}\label{IV:eq:1-3-4}
  \tr:\hom(V,V) {\xleftarrow\sim}_\text{\eqref{IV:eq:1-3-2}} V\otimes V^\vee = V^\vee \otimes V \to_\text{\eqref{IV:eq:1-3-1}} A .
\end{equation*}

It is easily verified that, for $f$ of degree $0$,
\begin{equation*}\tag{1.3.5}\label{IV:eq:1-3-5}
  \tr(f,V) = \sum (-1)^i \tr(f,V^i) \text{.}
\end{equation*}

Expressing the commutativity of the two morphisms obtained by composing the
morphisms \eqref{IV:eq:1-3-1}
$V^\vee\otimes V\otimes W^\vee\otimes W\to k$, one finds that for homogeneous
$f:V\to W$ and $g:W\to V$, one has
\begin{equation*}\tag{1.3.6}\label{IV:eq:1-3-6}
  \tr(f g) = (-1)^{\deg(f)\deg(g)} \tr(g f) \text{.}
\end{equation*}




















\section{The cohomology class associated with a cycle}\label{IV:2}










\subsection{The class of a divisor}\label{IV:2-1}






\subsubsection{}\label{IV:2-1-1}

Let $D$ be a Cartier divisor in a scheme $X$. Away from $D$, the invertible
sheaf $\sO(D)$ is trivialized by the section $1$. The class $\cl(D)$ of $D$
in $\h_D^1(X,\dG_m)$ is the class of the corresponding $\dG_m$-torsor
trivialized on $X\setminus D$ (\ref{IV:1-1-6} and \ref{IV:1-1-4}).

Let $\partial:\h^i(X\setminus D,\dG_m)\to \h_D^i(X,\dG_m)$ be the morphism
of \ref{IV:1-1-4}. If $D$ admits a global equation $f$, multiplication by
$f$ is an isomorphism from $\sO(D)$, trivialized by $1$ on $X\setminus D$,
to $\sO$, trivialized by $f$ on $X\setminus D$. By \ref{IV:1-1-4}, one
therefore has
\begin{equation*}\tag{2.1.1}\label{IV:eq:2-1-1}
  \cl(D) = \partial f \text{.}
\end{equation*}

For every morphism $u:X'\to X$ such that $u^* D$ is again a Cartier divisor
(i.e., $u^{-1} D$ is disjoint from $\operatorname{Ass}(X')$), one has
$\cl(u^* D) = u^*\cl(D)$. If such functoriality were desired for every
morphism $u$, one would have to consider not Cartier divisors, but more
generally invertible sheaves equipped with a section.

Recall that the integer $n$ is henceforth assumed invertible on the schemes
under consideration. Let
$\partial:\h_D^i(X,\dG_m) \to \h_D^{i+1}(X,\dmu_n)$ be the boundary map for
the Kummer exact sequence $0 \to \dmu_n \to \dG_m \to \dG_m \to 0$.






\begin{definition}\label{IV:2-1-2}
The \emph{class $\cl_n(D)$} of $D$ in $\h_D^2(X,\dmu_n)$ is
$\partial \cl(D)$.
\end{definition}

When there is no risk of confusion, the subscript $n$ will be omitted.






\subsubsection{}\label{IV:2-1-3}

The diagram
\[\xymatrix{
  \h^0(X\setminus D,\dG_m) \ar[r]^-\partial \ar[d]^-\partial
    & \h_D^1(X,\dG_m) \ar[d]^-\partial \\
  \h^1(X\setminus D,\dmu_n) \ar[r]^-\partial
    & \h_D^2(X,\dmu_n)
}\]
is anticommutative. If $D$ admits a global equation $f$, $\cl_n(D)$ is
therefore the negative of the image under $\partial$ of the class in
$\h^1(X\setminus D,\dmu_n)$ of the $\dmu_n$-torsor of $n$-th roots of $f$.






\begin{proposition}\label{IV:2-1-4}
Let $i$ be the inclusion of $D$ into $X$. If $D$ and $X$ are regular, the
cohomology sheaves with support $\eR^p i^! \dmu_n$ are zero for $p=0,1$, and
$\eR^2 i^!\dmu_n = \underline{\dZ/n}$, generated by $\cl_n(D)$.
\end{proposition}

It suffices to prove that, for $X$ strictly local and $D$ defined by a regular
parameter, one has $\h_D^p(X,\dmu_n) = 0$ for $p=0,1$ and
$\h_D^2(X,\dmu_n) = \dZ/n$, generated by $\cl(D)$. Denoting reduced
cohomology by $\sim$, one has
$\widetilde\h^{p-1}(X\setminus D,\dmu_n)\iso \h_D^p(X,\dmu_n)$. The assertion
for $p=0,1$ says that $D$ does not disconnect $X$, and for $p=2$ it follows,
via \ref{IV:2-1-3}, from Abhyankar's lemma.

This is a partial analogue of the relative theorem
(\hyperref[I]{Étale cohomology}, \ref{I:5-3-4}). Grothendieck conjectures
that the $\eR^p i^! \dmu_n$ vanish for $p\ne 2$, at least for $X$ excellent
(purity conjecture), but this is known only in characteristic $0$
(\cite[XIX]{sga4}).






\begin{theorem}[Fundamental compatibility]\label{IV:2-1-5}
Let $X$ be a smooth curve over an algebraically closed field $k$, let $P$ be a
closed point of $X$, and let $\tr$ be the composite
$\h_P^2(X,\dmu_n) \to \h_c^2(X,\dmu_n) \xrightarrow{\tr} \dZ/n$. One has
\[
  \tr \cl(P) = 1 \text{.}
\]
\end{theorem}

Let $\bar X$ be the smooth projective curve completing $X$. The formula says
that the invertible sheaf $\sO(P)$ on $\bar X$ has degree $1$.











\subsection{Cohomological method}\label{IV:2-2}






\subsubsection{}\label{IV:2-2-1}

Let $X$ be a (Noetherian) scheme. Recall that a subscheme $Y$ of $X$ is said
to be a local complete intersection of codimension $c$ if, locally (on $Y$),
it is defined by a regular sequence of $c$ equations in $X$. For $X$ the
spectrum of a local ring $A$ with maximal ideal $\fm$, and $Y$ with ideal
$\fa$, this means that $\Ext^i(A/\fa,A) = 0$ for
$i<\dim(\fa/\fm\fa)=c$, and every sequence of elements of $\fa$ whose image
in $\fa/\fm\fa$ is a basis of $\fa/\fm\fa$ is a regular sequence of
equations for $Y$.






\subsubsection{}\label{IV:2-2-2}

Let $i:Y\hookrightarrow X$ be a local complete intersection of codimension
$c$. We propose to define a local fundamental class $\cl(Y)$ that is a
global section of the cohomology sheaf with support
$\eR^{2c}i^! \dZ/n(c)$ (recall that $\dZ/n(c)=\dmu_n^{\otimes c}$).

Locally, $Y$ is the intersection of a sequence of $c$ divisors $D_i$, and we
define $\cl(Y)$ as the cup product of the $\cl(D_i)$. Each $\cl(D_i)$ has
support in $D_i$; their product has support in $Y$. The fact that this product
is locally independent of the choice of the $D_i$ follows from
\ref{IV:2-2-3} below and the following invariance properties:
\begin{enumerate}[\indent a)]
  \item compatibility with localization;
  \item independence of the ordering of the $D_i$ (the $\cl(D_i)$ have
    degree $2$, which is even, so the cup product is commutative);
  \item the product depends only on the ``flag''
    $D_1\supset D_1\cap D_2\supset \cdots\supset Y$.
\end{enumerate}
To prove c), note the existence of a product (a variant of
\ref{IV:1-2-1})
\[
  \h_{D_1}^\bullet(X)\otimes \h_{D_1\cap D_2}^\bullet(D_1)\otimes \cdots \otimes \h_Y^\bullet(D_1\cap \cdots \cap D_{c-1}) \to \h_Y^\bullet(X) \text{;}
\]
the product of the $\cl(D_i)$ is also the product of the
($\cl(D_i)$ restricted to
$D_1\cap \cdots \cap D_{i-1})=(\cl(D_1\cap \cdots \cap D_i)$ in
$D_1\cap \cdots \cap D_{i-1})$.






\begin{lemma}\label{IV:2-2-3}
Let $A$ be a local ring with maximal ideal $\fm$, and let
$u=(u_1,\dots,u_c)$, $v=(v_1,\dots,v_c)$ be two regular sequences generating
the same ideal $\fa$. There is then a sequence $w_i$
($1\leqslant i\leqslant N$) of such sequences joining them, such that
$w_{i+1}$ is obtained from $w_i$ by a permutation or by changing only the
last element.
\end{lemma}

Since permutations are available, it would amount to the same thing to allow
one element, rather than the last element, to be changed. By
\ref{IV:2-2-1}, we are thus reduced to verifying that, in the vector space
$\fa/\fm\fa$, one can pass from one basis to another by a sequence of
permutations and changes of basis modifying only one vector. Indeed, the
linear group is generated by diagonal, elementary, and permutation matrices.






\subsubsection{}\label{IV:2-2-4}

The same methods make it possible to define a local fundamental class for
every $Y\subset X$ locally definable by $c$ equations (where fewer equations
suffice, the class is zero).






\subsubsection{}\label{IV:2-2-5}

One can pass from such a local fundamental class in
$\h^0\left(Y,\eR^{2 c} i^!\dZ/n(c)\right)$ to a global fundamental class in
$\h_Y^{2 c}\left(X,\dZ/n(c)\right)$ when semipurity results are available:






\begin{proposition}\label{IV:2-2-6}
\begin{enumerate}[(i)]
  \item If $\eR^p i^!\dZ/n=0$ for $p<2 c$, then
    \[
      \h_Y^{2 c}\left(X,\dZ/n(c)\right) \iso \h^0\left(Y,\eR^{2 c} i^! \dZ/n(c)\right) \text{.}
    \]
  \item Let $Z$ be a closed part of $Y$, with complement $V$ in $Y$, and let
    $k$ be the inclusion of $Z$ into $X$. If $\eR^p i^! \dZ/n = 0$ for
    $p\leqslant 2 c$, then
    \[
      \h_Y^{2 c}\left(X,\dZ/n(c)\right) \hookrightarrow \h_V^{2 c}\left(X,\dZ/n(c)\right) \text{.}
    \]
    If $\eR^p i^!\dZ/n=0$ for $p\leqslant 2 c+1$, this arrow is an
    isomorphism.
\end{enumerate}
\end{proposition}

The hypothesis $\eR^p i^!\dZ/n=0$ is equivalent to
$\eR^p i^!\dZ/n(c)=0$. With this understood, (i) follows from the spectral
sequence $\h^p(Y,\eR^q i^!)\Rightarrow \h_Y^{p+q}(X,-)$. If $k_1$ is the
inclusion of $Z$ into $Y$, one has $k=i k_1$, whence
$\eR k^! \dZ/n=\eR k_1^! \eR i^!$, and the long exact cohomology sequence for
$Z\subset Y$ gives a long exact sequence
\[\xymatrix{
  \cdots \ar[r]
    & \h^i\left(Z,\eR k^! \dZ/n(c)\right) \ar[r]
    & \h_Y^i\left(X,\dZ/n(c)\right) \ar[r]
    & \h_V^i\left(X,\dZ/n(c)\right) \ar[r]
    & \cdots
}\]
from which (ii) follows.






\subsubsection{Amplification}\label{IV:2-2-7}

Proposition \ref{IV:2-2-6} remains valid for $i$ any separated morphism of
finite type, and $2 c$ any integer (positive or negative), provided that
$\h_Y^{2 c}\left(X,\dZ/n(c)\right)$ is replaced there by
$\h^{2 c}\left(Y,\eR i^!\dZ/n(c)\right)$, and likewise for $\h_V$.

The following semipurity results follow from \cite[1.8, 1.10, 1.15]{sga2}.
We recall their proof.






\begin{theorem}\label{IV:2-2-8}
Let there be a morphism of $S$-schemes of finite type
\[\xymatrix{
  Y \ar[r]^-i \ar[dr]_-f
    & X \ar[d]^-g \\
  & S
}\]
Suppose that $X$ is smooth, purely of relative dimension $N$, and that $Y$
has fiberwise dimension $\leqslant d$. Let $c=N-d$.
\begin{enumerate}[\indent (i)]
  \item One has $\eR^p i^!\dZ/n=0$ for $p<2 c$; likewise with $\dZ/n$
    replaced by a sheaf $g^* \sF$.
  \item If, over a dense open subset $U$ of $S$, $Y$ has fiberwise dimension
    $<d$, one has $\eR^p i^!\dZ/n=0$ for $p\leqslant 2 c$. If in addition
    the complement $Z$ of $U$ does not locally disconnect $S$, one has
    $\eR^pi^!\dZ/n=0$ for $p\leqslant 2 c+1$.
\end{enumerate}
\end{theorem}

Since $\eR g^! \sF=g^*\sF(N)[2N]$, the transitivity formula
$\eR i^!\eR g^!=\eR f^!$ shows that
$\eR^{2 c+q} i^!(g^*\sF)=0\Leftrightarrow \eR^{-2 d+q}f^!(\sF)$. This reduces
us to studying $f$, i.e. to assuming that $X=S$. Assertion (i) is then
\cite[XVIII 3.17]{sga4}.

Let $Y'$ be the inverse image of $Z$ in $Y$:
\[\xymatrix{
  Y' \ar[r]^-v \ar[d]^-f
    & Y \ar[d]^-f \\
  Z \ar[r]^-u
    & S
}\]
By (i), the $\eR^p f^! \dZ/n$ have support in $Y'$ for
$p\leqslant -2 d+1$; and the spectral sequence
$\eR^a v^!\eR^b f^! \Rightarrow \eR^{a+b} (f v)^!$ shows that they coincide
with the $\eR^p(f v)^!\dZ/n=\eR^p(u f)^!\dZ/n$. Applying (i) to $Y'/Z$ and
to the spectral sequence
$\eR^a f^!\eR^b u^! \Rightarrow \eR^{a+b} (u f)^!$, one finds that (ii)
follows from the vanishing of $\eR^b u^! \dZ/n$ for $b=0$, or for $b=0$ and
$1$, as the case may be.






\subsubsection{}\label{IV:2-2-9}

Grothendieck conjectures the following absolute analogue of
\ref{IV:2-2-8} (semipurity conjecture---a consequence of the purity
conjecture): for $Y$ of codimension $\geqslant c$ in a regular $X$, one has
$\h_Y^i(X,\dZ/n)=0$ for $i<2 c$, at least if $X$ is excellent.






\subsubsection{}\label{IV:2-2-10}

We are now in a position to define the class of a cycle $Y$ of codimension
$c$ in $X$ smooth over a field $k$. Write $Y=\sum d_i Y_i$, where the $Y_i$
are reduced and irreducible. An open subset $U_i$ of $Y_i$, whose complement
has codimension $>c$, is then a local complete intersection in $X$. This
makes it possible to define the local fundamental class of $U_i$. By
\ref{IV:2-2-6} and \ref{IV:2-2-8}, it comes from a unique fundamental class
$\cl(Y_i)\in \h_{Y_i}^{2 c}\left(X,\dZ/n(c)\right)$, and we set
\[
  \cl(Y) = \sum d_i \cl(Y_i) \in \h_{|Y|}^{2 c}\left(X,\dZ/n(c)\right) \text{.}
\]






\subsubsection{}\label{IV:2-2-11}

In complex analytic geometry, a construction of Baum, Fulton, and
Mac Pherson \cite{bfm75} makes it possible to define without restriction the
cohomology class of a local complete intersection $Y\subset X$. Suppose $Y$
is purely of codimension $c$, and let $\sN$ be the vector bundle on $Y$ that
is the normal bundle of $Y$ in $X$. Its local sections are those of
$(\sI/\sI^2)^\vee$, where $\sI$ is the sheaf of ideals of $Y$. Recall that
the sheaf of $C^\infty$ functions on $X$ is defined locally, in terms of local
embeddings of $X$ into $\dC^n$, as the restriction to $X$ of the quotient of
the sheaf of $C^\infty$ functions on $\dC^n$ by the ideal generated by the
real and imaginary parts of the equations defining $X$. The bundle $\sN$
extends to a complex $C^\infty N$ vector bundle on a neighborhood $U$ of
$Y$ in $X$, and for $U$ sufficiently small, there exist sections $f$ of $N$
whose zero locus is $Y$ and such that, on $Y$, $d f:\sN\to N$ is the
identity. Two choices of $N$ and $f$ are homotopic on a sufficiently small
neighborhood of $Y$.

The cohomology class $\cl(U)$ of the section $0$ of $N$ (denoted
$z:U\to N$) is defined: $U$ is a local complete intersection in $N$, and
$\eR^i z^!\dZ = 0$ for $i\ne 2 c$, $\eR^{2 c} z^!\dZ = \dZ$. We have
$f^*:\h_U^\bullet(N,\dZ) \to \h_Y^\bullet(X,\dZ)$ and set
\[
  \cl(Y) = f^*\cl(U) \text{.}
\]

The same construction works whenever an analytic subspace $Y$ of $X$ has
the following normal structure: a locally free sheaf $\sC$ of rank $c$ on
$Y_\text{red}$, and an epimorphism
$\sC\to \sI/\sI\cdot \sI_\text{red}$.










\subsection{Homological method}\label{IV:2-3}






\subsubsection{}\label{IV:2-3-1}

Let $f:Y\to S$ be a flat morphism of finite type with fibers of dimension
$\leqslant d$. In \cite[XVIII 2.9]{sga4}, we defined a trace morphism
\begin{equation*}\tag{2.3.1.1}\label{IV:eq:2-3-1-1}
  \tr_f:\eR^{2 d} f_! \dZ/n(d) \to \dZ/n \text{.}
\end{equation*}
One has $\eR^i f_!\dZ/n(d)=0$ for $i>2 d$, and
$\eR^i f^!\dZ/n=0$ for $i<-2 d$. This, together with the adjunction between
$\eR f_!$ and $\eR f^!$, gives isomorphisms
\begin{equation*}\tag{2.3.1.2}\label{IV:eq:2-3-1-2}
\begin{aligned}
  \hom\left(\eR^{2 d}f_!\dZ/n(d),\dZ/n\right)
    &= \hom\left(\eR f_!\dZ/n(d),\dZ/n[-2 d]\right) \\
    &= \hom\left(\dZ/n,\eR f^!\dZ/n(-d)[-2 d]\right) \\
    &= \h^0\left(Y,\eR^{-2 d}f^!\dZ/n\right) \text{.}
\end{aligned}
\end{equation*}
At least when $S$ is a point, the image of $\tr_f$ in the last two groups
deserves the name fundamental class of $Y$ (in homology).

We again denote by $\tr_f$ the image of $\tr_f$ in the second group, and the
morphisms deduced from it by functoriality. For example, for $Y/S$ proper,
the morphism
\begin{equation*}\tag{2.3.1.3}\label{IV:eq:2-3-1-3}
\xymatrix{
  \h^i\left(Y,\dZ/n(d)\right) = \h^i\left(S,\eR f_!\dZ/n(d)\right) \ar[r]
    & \h^{i-2 d}(S,\dZ/n) \text{.}
}
\end{equation*}

Suppose $Y$ is contained in $X$, smooth over $S$ and purely of relative
dimension $N$, and set $c=N-d$.
\[\xymatrix{
  Y \ar@{^{(}->}[r]^-i \ar[dr]_-f
    & X \ar[d] \\
  & S
}\]
One has $\eR f^! = \eR i^!\eR g^!$, and
$\eR g^! \dZ/n=\dZ/n(-N)[-2 N]$. The fundamental class of $Y$ is identified
with an element of
$\h^0\left(Y,\eR i^!\dZ/n(c)[2 c]\right)=\h_Y^{2 c}\left(X,\dZ/n(c)\right)$,
the class $\cl(Y)$ of $Y$ in $X$. We shall see below that it depends only on
$Y\subset X$, not on the projection of $X$ onto $S$, and that for $Y$ a local
complete intersection, it induces the local class of the preceding
subsection.

Let us make explicit the isomorphism
\[
  \h_Y^{2 c}\left(X,\dZ/n(c)\right) \iso \hom\left(\eR^{2 d}f_!\dZ/n(d),\dZ/n\right)
\]
that takes $\cl(Y)$ to $\tr_f$: via the isomorphisms
\begin{align*}
  \h_Y^{2 c}\left(X,\dZ/n(c)\right)
    &= \h_Y^{2 c}\left(\eR i^!\dZ/n(c)\right) \\
    &= \hom_Y\left(\dZ/n(d)[2 d],\eR i^!\dZ/n(N)[2 N]\right) \\
    &= \hom_X\left(\dZ/n(d)[2 d]_Y,\dZ/n(N)[2 N]\right) \text{,}
\end{align*}
it is the top row of
\[\xymatrix{
  \h_Y^{2 c} \ar[r] \ar[dr]_-{(3)}
    & \hom\left(\eR f_!\dZ/n(d)[2 d],\eR f_! \eR i^!\dZ/n(N)[2 N]\right) \ar[r]^-{(1)} \ar[d]^-{\tr_i} \ar@{}[dr]|-{(2)}
    & \hom\left(\eR f_!\dZ/n(d)[2d],\dZ/n\right) \ar@{=}[d] \\
  & \hom\left(\eR g_!\dZ/n(d)[2 d]_Y,\eR g_!\dZ/n(N)[2 N]\right) \ar[r]^-{\tr_g}
    & \hom\left(\eR g_!\dZ/n(d)[2 d]_Y,\dZ/n\right) \text{,}
}\]
where (1) is the adjunction arrow $\eR f_!\eR f^!\to \text{id}$. The
commutativity of (2) expresses the fact that the isomorphism
$\eR f^!=\eR i^!\eR g^!$ is defined by adjunction. Finally, the arrow
deduced from (3)
\[
  \h_Y^{2c}\left(X,\dZ/n(c)\right) \otimes \eR^{2 d} f_!\dZ/n(d) \to \eR^{2 N} g_!\dZ/n(N)
\]
is interpreted as a cup product (cf. \ref{IV:1-2}).






\begin{definition}\label{IV:2-3-2}
The class $\cl(Y)\in\h_Y^{2 c}\left(X,\dZ/n(c)\right)$ is characterized by
the property that, for every local section $u$ of
$\eR^{2 d}f_! \dZ/n(d)$, one has
\[
  \tr_f(u) = \tr_g\left(\cl(Y)\smallsmile u\right) \text{.}
\]
\end{definition}

Because $\cl(Y)$ already has an image $\tr_f$ in
$\hom\left(\eR f_!\dZ/n(d)[2 d],\dZ/n\right)$, formula
\ref{IV:2-3-2} holds for the arrows deduced from $\tr_f$ by functoriality.
For example, for $X$ proper over $S$ and
$u\in \h^\bullet(Y,\dZ/n)$, the formula holds in
$\h^\bullet\left(S,\dZ/n(-d)\right)$. For
$v\in \h^\bullet(X,\dZ/n)$, it gives
\[
  \tr_f(i^* v) = \tr_g\left(\cl(Y)\smallsmile v\right) \text{,}
\]
where the cup product can be computed in $\h^\bullet(X,\dZ/n)$.






\subsubsection{}\label{IV:2-3-3}

We shall define the trace morphisms, and hence the class $\cl(Y)$, under more
general hypotheses. Consider a diagram
\[\xymatrix{
  Y \ar@{^{(}->}[r] \ar[dr]_-f
    & X \ar[d]^-g \\
  & S
}\]
with $g$ smooth and purely of relative dimension $N$, and $Y$ closed in $X$
with fiberwise dimension $\leqslant d$. Equip $Y$ in addition with a
``weight'' $\sK$ of the following kind:
$\sK\in\eD_{\text{parf},X}^b$ is a bounded complex of sheaves of
$\sO$-modules on $X$, of finite $\tor$-dimension over $S$ (or over $X$,
which amounts to the same thing), whose cohomology sheaves are coherent and
have support in $Y$. We propose to define a morphism
$\tr_{f,\sK}:\eR^{2 d}f_! \dZ/n(d) \to \dZ/n$. For $Y$ flat over $S$ and
$\sK=\sO_Y$, this will be the preceding trace morphism. In general, it
depends only on the lengths of $\sK$ at the generic points $y$ of the
irreducible components of
$Y$ ($\operatorname{lg}_y(\sK) = \sum (-1)^i \operatorname{lg} \h^i(\sK)_y$).
% NOTE: notation is not clear, is it a calligraphic H?
Finally, $\cl(Y,\sK)$ will denote the class with support in $Y$ such that
\begin{equation*}\tag{2.3.3.1}\label{IV:eq:2-3-3-1}
  \tr_g\left(\cl(Y,\sK)\smallsmile u\right) = \tr_{f,\sK}(u) \text{.}
\end{equation*}






\subsubsection{}\label{IV:2-3-4}

The construction of $\tr_{f,\sK}$ is parallel to that of
\cite[XVIII.2]{sga4}; we shall indicate only its main outlines.

\paragraph{A}
$d=0$ ($f$ quasi-finite). Let $x$ be a geometric point of $Y$, $s$ its image
in $S$, $\sK_{(x)}$ the inverse image of $\sK$ on the strict localization of
$X$ at $x$, $\sK_{(x)s} = \sK_{(x)}\lotimes_{\sO_{S,s}^h} k(s)$ its inverse
image on the geometric fiber at $s$, and
$n(x)=\sum (-1)^i \dim_{k(s)} \h^i(\sK_{(x)s})$. The function
$x\mapsto n(x)$ is a weighting of $f$, and we take the corresponding trace
morphism \cite[XVII.6.2.5]{sga4}. Formation of the weighting $n(x)$ and of
$\tr_{f,\sK}$ is compatible with every base change.

\paragraph{B}
General case. If $u:X\to \dA_S^d$ is such that $u i$ is quasi-finite, set
$\tr_{f,\sK} = \tr_{\dA_S^d}\circ \tr_{u i,\sK}$. This trace morphism is
clearly compatible with every base change $S'/S$. To prove that it does not
depend on $u$, we may therefore suppose that $S$ is the spectrum of an
algebraically closed field $k$. In that case, let $Y_i$ be the irreducible
components of $Y$, and let $\operatorname{lg}_i$ be the length of $\sK$ at
the generic point of $Y_i$. If $\alpha_i$ is the inclusion of
$(Y_i)_\text{red}$ into $Y$, one has
\begin{equation*}\tag{2.3.4}\label{IV:eq:2-3-4}
  \tr_{f,\sK} = \sum \operatorname{lg}_i \tr_{Y_{i,\text{red}/S}} \alpha_i^\ast
\end{equation*}
(this formula follows from the analogous and easy formula for $\tr_{u i}$).

This defines $\tr_{f,\sK}$ locally on $Y$; one then proceeds as in
\cite[XVIII.2.9]{sga4}.






\begin{lemma}\label{IV:2-3-5}
If $g=g'g'':X\xrightarrow{g''} S'\xrightarrow{g'} S$, with $g'$ and $g''$
smooth and purely of relative dimensions $N'$ and $N''$, and if $Y$ still
has fiberwise codimension $\geqslant c$ over $S'$, then the class
$\operatorname{cl}(Y,\sK)$ is the same whether computed in terms of $g$ or
of $g''$.
\end{lemma}

Let $f'=g'' i$ and $d'=d-N'$. The formula
$\tr_g=\tr_{g'}\tr_{g''}$ (\cite[XVIII.2.9]{sga4} Var 3) ensures the
commutativity of
\[\xymatrix{
  \h_Y^{2 c}\left(X,\dZ/n(c)\right) \ar[d] \\
  \hom_{S'}\left(\eR f_!\dZ/n(d')[2 d'],\eR g_!''\dZ/n(N'')[2 N'']\right) \ar[r] \ar[d]^-{\eR g_!'(N')}
    & \hom_{S'}\left(\eR^{2 d'} f_!' \dZ/n,\dZ/n\right) \ar[d]^-{\tr_{g'}\circ \eR^{2 N'} g_!'} \\
  \hom_S\left(\eR f_!\dZ/n(d)[2d],\eR g_!'\dZ/n(N)[2 N]\right) \ar[r]
    & \hom_S\left(\eR^{2 d} f_!\dZ/n,\dZ/n\right)
}\]
while the analogous and easily verified formula
$\tr_{f,\sK}=\tr_{g'}\tr_{f',\sK}$ ensures that the image of
$\tr_{f',\sK}$ is $\tr_{f,\sK}$. The assertion follows.






\begin{lemma}\label{IV:2-3-6}
If $Y$ is a divisor in $X$, the class of \ref{IV:2-3-2} coincides with the
class of \ref{IV:2-1-2}.
\end{lemma}

By semipurity (\ref{IV:2-2-6}(i) and \ref{IV:2-2-8}(i)), the problem is local
on $Y$; using \ref{IV:2-3-5} to replace $S$ by a suitable $S'$, this allows
us to suppose that $N=1$. Using semipurity once more
(\ref{IV:2-2-6}(ii) and \ref{IV:2-2-8}(ii)), it suffices to prove
\ref{IV:2-3-6} over the generic points of $S$. An étale localization then
reduces us to assuming that $S$ is the spectrum of an algebraically closed
field $k$. Since the classes of \ref{IV:2-3-2} and \ref{IV:2-1-2} are each
additive in $Y$, this reduces us to the case where $Y$ is a closed point on a
smooth curve over $k$. The fact that the characteristic property
\ref{IV:2-3-2} is satisfied is then the fundamental compatibility
\ref{IV:2-1-5}.






\begin{lemma}\label{IV:2-3-7}
The formation of $\cl(Y,\sK)$ is compatible with every base change $S'/S$.
\end{lemma}

This follows from the same assertion for the trace morphisms.






\begin{theorem}\label{IV:2-3-8}
(i) $\cl(Y,\sK)$ depends only on $Y\subset X$ and on the lengths of $\sK$ at
the generic points of $Y$. This class is additive in $\sK$. For
$\sK=\sO_Y$ and $Y$ a local complete intersection, it induces the local
class of \ref{IV:2-2-2}.

(ii) Let there be a commutative diagram
\[\xymatrix{
  X' \ar[r]^-u \ar[d]
    & X \ar[d] \\
  S' \ar[r]
    & S
}\]
with $X'$ smooth over $S'$. If $u^{-1}(Y)$ still has fiberwise codimension
$\geqslant c$, then
\[
  \cl\left(u^{-1}(Y),\eL u^\ast\sK\right) = u^\ast \cl(Y,\sK) \text{.}
\]

(iii) Let $Y'\subset X$ have fiberwise codimension $\geqslant c'$, let
$\sK'$ be a weight on $Y'$, and suppose that $Y\cap Y'$ has fiberwise
codimension $\geqslant c+c'$. Then
\[
\cl\left(Y\cap Y',\sK\lotimes \sK'\right) = \cl(Y,\sK) \smallsmile \cl(Y',\sK') \text{.}
\]
\end{theorem}



% The following paragraphs were originally numbered

% (2.3.8.1)
\paragraph{(A)}
Given $Y\subset X/S$ as above, it follows from semipurity
(\ref{IV:2-2-6}, \ref{IV:2-2-8}) that to verify that two classes in
$\h_Y^{2 c}\left(X,\dZ/n(c)\right)$ coincide, it suffices to verify this
locally at the generic points of $Y$, or even after every base change
$s\to S$ ($s$ a geometric generic point of $S$), at the generic points of
$Y$.

% (2.3.8.2)
\paragraph{(B) Proof of (ii) for $S=S'$ and $u$ a closed immersion}
By localization (A), we may suppose that $X'$ is the inverse image of the
section $0$ under a smooth morphism $v:X\to \dA_S^{N'}$
\[\xymatrix{
  X' \ar[r]^-u \ar[d]
    & X \ar[d]^-v \\
  S \ar[r]^-0
    & \dA_S^{N'} \ar[r]^-\pi
    & S
}\]
Apply \ref{IV:2-3-5} to $g=\pi v$ and \ref{IV:2-3-7} to the base change
$S\to \dA_S^{N'}$.

% (2.3.8.3)
\paragraph{(C) In a product situation:}
$X=X'\times_S X''$, $Y=Y'\times_S Y''$, $\sK=\sK'\lotimes_S \sK''$, one has
$\cl(Y,\sK) = \cl(Y',\sK')\smallsmile \cl(Y'',\sK'')$ (external cup product,
i.e. $\operatorname{pr}_1^\ast \cl(Y',\sK')\smallsmile \operatorname{pr}_2^\ast \cl(Y'',\sK'')$).

The functorial properties of $\tr_g$ \cite[XVIII.2.12]{sga4} imply the
commutativity of
\small
\[\xymatrix@=0.5cm{
  \h_{X'}^{2c'}\otimes \h_{X''}^{2 c''}\ar[r] \ar[d]
    & \hom_S\left(\eR f_!' \dZ/n(d')[2 d'],\eR g_!'\dZ/n(N')[2 N']\right)\otimes \cdots \ar[r] \ar[d]
    & \hom_S\left(\eR^{2 d'} f_!' \dZ/n,\dZ/n\right) \otimes \cdots \ar[d] \\
  \h_X^{2 c} \ar[r]
    & \hom_S\left(\eR f_!\dZ/n(d)[2 d],\eR g_!\dZ/n(N)[2 N]\right) \ar[r]
    & \hom_S\left(\eR^{2 d}f_!\dZ/n,\dZ/n\right)
}\]
\normalsize
and one verifies that $\tr_f$ is the image of
$\tr_{f'}\otimes \tr_{f''}$, whence the assertion.

\paragraph{Proof of (ii)}
(\ref{IV:2-3-7}) and a preliminary base change reduce us to assuming that
$S=S'$. Factor $u$ through its graph:
\[\xymatrix{
  X' \ar[r]^-{(\text{id},u)}
    & X'\times X \ar[r]^-{\operatorname{pr}_2}
    & X \text{.}
}\]
This leads us to treat $\operatorname{pr}_2$ and $(\text{id},u)$ separately:
the former is covered by (C) with $Y'=X'$, $\sK'=\sO_{X'}$ (with
$\cl(Y')=1\in \h^0(X',\dZ/n)$), and the latter by (B).

\paragraph{Proof of (iii)}
Let $\delta:X\to X\times_S X$ be the diagonal inclusion. We use the formula
$\delta^\ast(Y_1\times_S Y_2) = Y_1\cap Y_2$ to reduce to (B) and (C).

\paragraph{Proof of (i)}
The base change $S_\text{red}\to S$ replaces $X$ by $X_\text{red}$; this
reduces us to assuming $X$ and $S$ reduced. Localization near the generic
points of $Y$ then reduces us to the case where $Y_\text{red}$ is irreducible
and the complete intersection of $c$ divisors $D_i$. The same localization,
together with the definition of $\tr$ in the quasi-finite case, then shows
that
$\cl(Y,\sK)=\operatorname{lg}\cl(Y_\text{red},\sO_{Y_\text{red}})$, where
$\operatorname{lg}$ is the length of $\sK$ at the generic point of $Y$.
By (iii), one therefore has
\[
  \cl(Y,\sK) = \operatorname{lg} \prod_i \cl(D_i,\sO_{D_i})
\]
and the conclusion follows from \ref{IV:2-3-6}.






\subsubsection{Remark}\label{IV:2-3-9}

For $S$ the spectrum of a field, with the class of a cycle defined as in
\ref{IV:2-2-10}, assertion (iii) says that if two cycles $Y'$ and $Y''$
intersect with the correct dimension, then
\[
  \cl(Y'\cap Y'') = \cl(Y')\smallsmile \cl(Y'')\text{,}
\]
provided that the multiplicities of the components of $Y'\cap Y''$ are
computed as alternating sums of $\tor$.






\subsubsection{Remark}\label{IV:2-3-10}

Let $X/\spec(k)$, with $k$ algebraically closed. If two cycles of codimension
$c$ in $X$ are algebraically equivalent, the existence of the relative
theory above implies that the images in $\h^{2 c}(X,\dZ/n(c))$ of their
classes are the fibers, at two points of a suitable connected $S$, of a
section over $S$ of the constant sheaf $\h^{2 c}(X,\dZ/n(c))$ (the sheaf
$\eR^{2 c} f_\ast \dZ/n(c)$ for
$f=\operatorname{pr}_2:X\times S\to S$): two algebraically equivalent cycles
have the same class in $\h^{2 c}(X,\dZ/n(c))$.





















\section{Application: the Lefschetz trace formula in the proper smooth case}\label{IV:3}






\subsection{}\label{IV:3-1}

Let $X$ and $Y$ be proper smooth algebraic varieties over an algebraically
closed field $k$, purely of dimensions $N$ and $M$. Via the K\"unneth
isomorphism $\h^\bullet(X\times Y)=\h^\bullet(X)\otimes \h^\bullet(Y)$, the
trace morphism $\h^\bullet(X\times Y)(M)\to \h^\bullet(X)$ is simply
$\h^\bullet(X)\otimes ($the trace morphism
$\h^\bullet(Y)(M)\to \dQ_\ell)$. Since the morphism is homogeneous of even
degree, there is no sign problem. We denote it by $\int_Y$.






\subsection{}\label{IV:3-2}

A class $\eta\in \h^{2 q}(X\times Y)(q)$ defines a morphism of degree
$2(q-M)$, $\eta_{XY}:\h^\bullet(Y)(M-q)\to \h^\bullet(X)$, by the formula
\[
  \beta \mapsto \int_Y \eta\cdot \operatorname{pr}_2^\ast \beta \text{.}
\]






\begin{proposition_}\label{IV:3-3}
Let $p$ and $q$ be two integers, with $p+q=N+M$, and let
$\varepsilon\in \h^{2 p}(X\times Y)(p)$,
$\eta\in \h^{2 q}(X\times Y)(q)$. With the trace of
$\eta_{X Y}\varepsilon_{Y X}:\h^\bullet(X)\to \h^\bullet(X)$ understood in
the sense of \ref{IV:1-3}, one has
\[
  \tr_{X\times Y}(\eta\cdot \varepsilon) = \tr\left(\eta_{X Y}\varepsilon_{Y X},\h^\bullet(X)\right) \text{.}
\]
\end{proposition_}

To exorcise the signs, it is helpful to retain from the grading of
$\h^\bullet$ only the underlying $\dZ/2$-grading. To avoid carrying Tate
twists along, also fix an isomorphism $\dQ_\ell(1)=\dQ_\ell$. Let $\alpha_Y$
be the homomorphism $\h^\bullet(Y)\to \h^\bullet(Y)^\vee$ (an isomorphism)
that makes the diagram
\[\xymatrix{
  \h^\bullet(Y)\otimes \h^\bullet(Y) \ar[r] \ar[d]^-{\alpha_Y\otimes 1}
    & \h^\bullet(Y) \ar[r]^-{\tr}
    & \dQ_\ell \ar@{=}[d] \\
  \h^\bullet(Y)^\vee \otimes \h^\bullet(Y) \ar[rr]^-{\text{(\ref{IV:eq:1-3-1})}}
    & & \dQ_\ell \text{.}
}\]
commute. The map $\eta_{X Y}$ is the image of $\eta$ under
\[\xymatrix{
  \h^\bullet(X\times Y) = \h^\bullet(X)\otimes \h^\bullet(Y) \ar[r]^-{\alpha_Y}
    & \h^\bullet(X) \otimes \h^\bullet(Y)^\vee \ar@{=}[r]^-{\text{(\ref{IV:eq:1-3-2})}}
    & \hom\left(\h^\bullet(Y),\h^\bullet(X)\right) \text{,}
}\]
and likewise for $\varepsilon$. Definition \ref{IV:eq:1-3-4} of the trace
then reduces \ref{IV:3-3} to the formula
$\tr_{X\times Y}=\tr_X\otimes \tr_Y$.






\subsection{Remark}\label{IV:3-4}

The proof of \ref{IV:3-3} does not use the fact that $\alpha_X$ and
$\alpha_Y$ are isomorphisms. It is valid without assuming $X$ and $Y$
smooth. Outside the smooth case, however, it is difficult to find
\emph{cohomology} classes to which one would want to apply \ref{IV:3-3}.






\subsection{Remark}\label{IV:3-5}

If $\varepsilon,\eta$ are the cohomology classes of algebraic cycles $A$ and
$B$ on $X\times Y$, of codimensions $p$ and $q$, and $A\cap B$ has dimension
$0$, then $\tr_{X\times Y}(\eta\varepsilon)$ is the intersection
multiplicity $A\cdot B$, computed by an alternating sum of $\tor$
(\ref{IV:2-3-9}).






\begin{proposition_}\label{IV:3-6}
For $q=M$ and $\eta=\cl(B)$ the class of a subscheme
$B\subset X\times Y$, finite over $X$, the map $\eta_{X Y}$ is the composite
\[\xymatrix{
  \eta_{X Y} : \h^\bullet(Y) \ar[r]^-{\operatorname{pr}_2^\ast}
    & \h^\bullet(B) \ar[r]^-{\tr_{B/X}}
    & \h^\bullet(X) \text{.}
}\]
\end{proposition_}

Indeed, by \ref{IV:3-2} and \eqref{IV:eq:2-3-3-1} for $\sK=\sO_B$, one has
\[
  \eta_{X Y}(\beta) = \tr_{X\times Y/X}\left(\cl(B)\cdot \operatorname{pr}_2^\ast\beta\right) = \tr_{B/X}\left(\operatorname{pr}_2^\ast \beta\right) \text{.}
\]

The most important case is that in which $B$ is the graph of a map
$f:X\to Y$. In this case, \ref{IV:3-6} shows that $\eta_{X Y}=f^\ast$.






\begin{corollary_}\label{IV:3-7}
Let $f$ be an endomorphism of $X$ proper and smooth over an algebraically
closed field $k$. Suppose that the fixed points of $f$ are isolated. Then the
trace $\sum (-1)^i \tr(f^\ast,\h^i(X))$ is the number of fixed points of $f$,
each counted with its multiplicity.
\end{corollary_}

In \ref{IV:3-5}, take $A=$ the graph of the identity and $B=$ the graph of
$f$, and apply \ref{IV:3-3}, \ref{IV:3-6}.






\begin{corollary_}\label{IV:3-8}
Let $X$ be a curve over an algebraically closed $k$, obtained by extension of
scalars from $X_0$ over $\dF_q$, and let $f$ be the Frobenius morphism. Then
$\sum (-1)^i \tr\left(f^\ast,\h_c^i(X)\right)$ is the number of fixed
points of $f$.
\end{corollary_}

One can replace $X$ by $X_\text{red}$, and hence suppose that $X$ is reduced.
Let $U$ be the open subset of $X$ on which $X$ is smooth of dimension $1$,
and let $\bar U$ be the smooth projective completion of $U$. The long exact
cohomology sequences for $(X,U)$ and $(\bar U,U)$ give
\[
  \tr\left(f^\ast,\h_c^\bullet(X)\right) = \tr\left(f^\ast,\h_c^\bullet(U)\right) + \tr\left(f^\ast,\h_c^\bullet(X\setminus U)\right)
\]
and
\[
  \tr\left(f^\ast,\h_c^\bullet(\bar U)\right) = \tr\left(f^\ast,\h_c^\bullet(U)\right) + \tr\left(f^\ast,\h_c^\bullet(\bar U\setminus U)\right) \text{.}
\]
The same formulas hold for the numbers of fixed points. Since formula
\ref{IV:3-8} is clear for a scheme of dimension $0$, this reduces us to
proving \ref{IV:3-8} for $\bar U$. This case follows from \ref{IV:3-7}. The
fact that $d f=0$ guarantees that all fixed points have multiplicity one.






\subsection{Remark}\label{IV:3-9}

For $X$ of dimension $\leqslant 1$ over an algebraically closed $k$, and
$f:X\to X$ such that $d f=0$, it is not difficult to verify that
\ref{IV:3-8} is still valid.

% END INCLUDED components/sga4half_4_en.tex

% BEGIN INCLUDED components/sga4half_5_en.tex
% !TEX root = ../SGA4half_English_Reader_1_to_8.tex

\chapter{Duality}\label{V}

This exposé contains several theorems and compatibility results, all related
to Poincaré duality. In Section \ref{V:1}, we give the local biduality theorem
in dimension $1$ \cite[I.5.1]{sga5}, together with some computations of
duals. In Section \ref{V:2}, we give a very economical proof of Poincaré
duality for curves, taught to me by M.\ Artin. In Section \ref{V:3}, we prove
a compatibility linking two definitions of the pairing giving rise to
Poincaré duality for curves: one by cup product, the other by autoduality of
the Jacobian. Finally, in Section \ref{V:4}, we prove the compatibility
announced in the title.

Throughout the exposé, all schemes are noetherian and separated, and $n$ is
an integer invertible on every scheme under consideration.


\section{Local biduality in dimension \texorpdfstring{$1$}{1}}\label{V:1}


\subsection{}\label{V:1-1}

Let $S$ be a regular scheme, purely of dimension $1$. We shall show that the
complex with value $\dZ/n$ concentrated in degree $0$ is dualizing; that is,
for $\sK\in \ob\eD_c^b(S,\dZ/n)$ (where $(-)_c$ means constructible), if we
put $D\sK=\rHom(\sK,\dZ/n)$, then $D\sK$ is again
constructible with bounded cohomology, and the canonical morphism $\alpha$
from $\sK$ to $D D\sK$ is an isomorphism.

For $\sK$ in $\eD^-$ and arbitrary $\sL$, we have
\[
  \hom(\sK\lotimes \sL,\dZ/n) = \hom(\sL,\Hom(\sK,\dZ/n)) \text{.}
\]
The morphism from $\sK$ to $DD\sK$ is defined when $D\sK$ lies in $\eD^-$:
if $\beta:D\sK\lotimes\sK\to \dZ/n$ is the canonical pairing, it is the
morphism defined by the pairing
$\sK\lotimes D\sK=D\sK\lotimes \sK \xrightarrow\beta \dZ/n$.


\subsection{}\label{V:1-2}

Let $f:X\to S$ be a separated morphism of finite type. Put
$\sK_X=\eR f^!\dZ/n$. For $\sK\in \ob\eD^-(X,\dZ/n)$, put
$D\sK = \rHom(\sK,\sK_X)$. The adjunction between $\eR f_!$ and $\eR f^!$
gives
\[\xymatrix{
  \eR f_\ast \sK\lotimes \eR f_\ast D\sK \ar[r]
    & \eR f_! (\sK\lotimes D\sK) \ar[r]
    & \eR f_! \eR f^! \dZ/n \ar[r]
    & \dZ/n \text{.}
}\]
The symmetry of this description shows that, when $f$ is proper, the diagram
\begin{equation*}\tag{1.2.1}\label{V:eq:1-2-1}
\xymatrix{
  \eR f_\ast \sK \ar[r] \ar[d]
    & D D \eR f_\ast \sK \ar[d]^-\wr \\
  \eR f_\ast D D \sK \ar[r]^\sim
    & D\eR f_\ast D\sK
}
\end{equation*}
is commutative.

If $f$ is the inclusion of a closed point, we know that
$\eR f^!\dZ/n=\dZ/n(-1)[-2]$---that is, up to a twist and a shift, $\dZ/n$.
Over the spectrum of a field this complex is dualizing (Pontryagin duality
for $\dZ/n$-modules). By \eqref{V:eq:1-2-1}, we therefore have
$\sK\iso D D\sK$ for $\sK$ of the form $\eR f_!\sL$, and hence whenever the
support of $\underline\h^\bullet(\sK)$ is finite.


\begin{theorem_}\label{V:1-3}
Let $j:U\hookrightarrow S$ be a dense open subscheme of $S$, and let $\sF$
be a locally constant constructible sheaf of $\dZ/n$-modules on $U$. Then
$\eD j_\ast \sF = j_\ast D \sF$; that is,
$\Hom(j_\ast \sF,\dZ/n) = j_\ast\Hom(\sF,\dZ/n)$ and
$\Ext^i(j_\ast\sF,\dZ/n) = 0$ for $i>0$.
\end{theorem_}

On $U$, $\Ext^i(j_\ast\sF,\dZ/n)=0$ for $i>0$, since $\sF$ is locally
constant (and $\dZ/n$ is an injective $\dZ/n$-module), while for $i=0$ it is
the dual $\sF^\vee$ of $\sF$. One checks that
$\hom(j_\ast\sF,\dZ/n)=j_\ast\sF^\vee$, and it remains to verify the
vanishing of the $\Ext^i$ for $i>0$ at the points of $S\setminus U$. The
problem is local at these points. We are thus reduced to assuming that $S$
is a strictly local trait and that $U$ consists of its generic point $\eta$.
Let $I=\gal(\bar\eta/\eta)$. The sheaf $\sF$ is identified with the Galois
module $\sF_{\bar\eta}$, and the special fiber of $j_\ast \sF$ with
$\sF_{\bar\eta}^I$.

Let $i$ be the inclusion of the closed point $s$, and apply $D$ to the exact
sequences
\[\xymatrix{
  0 \ar[r]
    & j_! \sF \ar[r]
    & j_\ast \sF \ar[r]
    & i_\ast \sF_{\bar\eta}^I \ar[r]
    & 0 \\
  0 \ar[r]
    & j_! \sF_{\bar\eta}^I \ar[r] \ar[u]
    & \sF_{\bar\eta}^I \ar[r] \ar[u]
    & i_\ast \sF_{\bar\eta}^I \ar[r] \ar[u]
    & 0 \text{.}
}\]

We obtain a morphism of triangles
\[\xymatrix{
  i_\ast\left(\sF_{\bar\eta}^I\right)^\vee (-1)[-2] \ar[r] \ar@{=}[d]
    & D j_\ast\sF \ar[r] \ar[d]
    & \eR j_\ast \sF^\vee \ar[d] \\
  i_\ast\left(\sF_{\bar\eta}^I\right)^\vee(-1)[-2] \ar[r]
    & \left(\sF_{\bar\eta}^I\right)^\vee \ar[r]
    & \eR j_\ast \sF_{\bar\eta}^I \text{.}
}\]
The long exact sequence arising from the first row gives the vanishing of the
$\Ext^i$ for $i>2$, and, on taking the fiber at $s$, we find
\[\xymatrix@=0.7cm{
  0 \ar[r]
    & \left(\underline\h^1(D j_\ast\sF)\right)_s \ar[r]
    & \h^1\left(I,\sF_{\bar\eta}^\vee\right) \ar[r]^-\partial \ar[d]
    & \left(\sF_{\bar\eta}^I\right)^\vee(-1) \ar[r] \ar@{=}[d]
    & \left(\underline\h^2 D j_\ast\sF\right)_s \ar[r]
    & 0 \\
  & 0 \ar[r]
    & \h^1\left(I,(\sF_{\bar\eta}^I)^\vee\right) \ar[r]^-\partial
    & \left(\sF_{\bar\eta}^I\right)^\vee(-1) \ar[r]
    & 0 \text{.}
}\]
Since
$\left(\sF_{\bar\eta}^I\right)^\vee = \left(\sF_{\bar\eta}^\vee\right)_I$,
putting $M=\sF_{\bar\eta}^\vee$, it remains only to check that
\[\xymatrix{
  \h^1(I,M) \ar[r]^-\sim
    & \h^1(I,M_I) \text{.}
}\]
If $p$ is the residual characteristic exponent, then $I$ is an extension of
a group isomorphic to $\widehat\dZ_{p'} = \varprojlim_{(m,p)=1} \dZ/m$ by a
$p$-group $P$. Since the order of $P$ is prime to the order of $M$, we have
$\h^i(P,M)=0$ for $i>0$ and $(M_I)^P=(M^P)_I$. This allows us to replace $M$
by $M^P$ and $I$ by $I/P$. Finally, there is a functorial isomorphism
$\h^1(\widehat\dZ_{p'},M)\sim ($coinvariants of $\widehat\dZ_{p'}$ in $M)$,
which proves the theorem.


\begin{theorem_}\label{V:1-4}
For $\sK\in \ob\eD_c^b(S,\dZ/n)$, one has $\sK\iso D D \sK$.
\end{theorem_}

By dévissage, we reduce to the case in which $\sK$ is a constructible sheaf
$\sF$ concentrated in degree $0$, and $\sF$ either has finite support
(\ref{V:1-2}) or is of the form $j_\ast\sF_1$ as in \ref{V:1-3}. In the
second case, \ref{V:1-3} reduces us to local biduality for $\sF$ locally
constant on $U$.


\section{Poincaré duality for curves, after M.\ Artin}\label{V:2}


\subsection{}\label{V:2-1}

Let $X$ be a smooth projective curve over an algebraically closed field $k$.
Put $\sK_X=\dZ/n(1)[2]$ and, for $\sK\in\ob \eD_c^b(X,\dZ/n)$,
$D\sK=\rHom(\sK,\sK_X)$. For $M$ a $\dZ/n$-module, also put
$D M = \hom(M,\dZ/n)$, and similarly for complexes of modules. The trace
morphism $\h^0(X,\sK_X) \to \dZ/n$, or
$\eR\Gamma(X,\sK_X)\to \dZ/n$, defines a pairing
$\eR\Gamma(X,\sK)\lotimes\eR\Gamma(X,D\sK) \to \dZ/n$. Poincaré duality
between the cohomology and compactly supported cohomology of an open
subscheme $j:U\hookrightarrow X$ of the curve $X$ says that, for
$\sK=j_!\dZ/n$, this pairing identifies each factor with the dual of the
other.

I give below a proof, communicated to me by M.\ Artin, that for
$\sK\in\ob\eD_c^b(X,\dZ/n)$ this pairing is always perfect; that is, it
defines an isomorphism
\begin{equation*}\tag{2.1.1}\label{V:eq:2-1-1}
\xymatrix{
  \eR\Gamma(X,\sK) \ar[r]^-\sim
    & D \eR\Gamma(X,D\sK) \text{.}
}
\end{equation*}

For every constructible sheaf $\sF$, put
\[
  '\h^i(X,\sF) = \text{the $\dZ/n$-dual of $\h^{-i}(X,D\sF)$.}
\]
The assertion that \eqref{V:eq:2-1-1} is an isomorphism is equivalent to the
following theorem.


\begin{theorem_}\label{V:2-2}
For a constructible sheaf $\sF$ of $\dZ/n$-modules, one has
\begin{equation*}\tag{2.2.1}\label{V:eq:2-2-1}
\xymatrix{
  \h^i(X,\sF) \ar[r]^-\sim
    & '\h^i(X,\sF) \text{.}
}
\end{equation*}
\end{theorem_}


\begin{lemma_}\label{V:2-3}
Let $f:X\to Y$ be a generically étale morphism between smooth curves over
$k$. Then $\sK_X=f^\ast\sK_Y=\eR f^!\sK_Y$, with $\tr_f$ as the adjunction
morphism $\eR f_!\eR f^!\sK_Y\to \sK_Y$.
\end{lemma_}

% NOTE: \rhom or \rHom ? originally \rhom
We reduce to checking that, for $f$ finite,
$f_\ast \eR f^!\sK_Y = \rHom(f_\ast\dZ/n,\sK_Y)$ is $f_\ast\dZ/n$, with
$\tr_f$ as the adjunction morphism. The vanishing of the $\Ext^i$ for $i>0$
follows from \ref{V:1-3}, and the assertion follows.

More explicitly, $\eR f^!$ is the derived functor of the functor
$f^!:f^!\sF(U)=\hom(f_!\dZ_U,\sF)$, and \ref{V:1-3} implies the vanishing of
$\eR^i f^!\dZ/n$ for $i>0$.


\begin{corollary_}\label{V:2-4}
Let $f:X'\to X$ be a generically étale morphism of smooth projective curves
over $k$. Then $'\h^i(X,f_\ast \sF)='\h^i(X',\sF)$; this isomorphism is
functorial, and the diagram
\[\xymatrix{
  \h^i(X,f_\ast\sF) \ar[r] \ar@{=}[d]
    & '\h^i(X,f_\ast\sF) \ar@{=}[d] \\
  \h^i(X',\sF) \ar[r]
    & '\h^i(X',\sF)
}\]
is commutative.
\end{corollary_}

The isomorphism is given by \ref{V:1-2}: $D f_\ast\sF=f_\ast D\sF$.


\subsection{}\label{V:2-5}

We prove \ref{V:2-2}. If $\sF$ is $\dZ/n$ at one point, extended by $0$,
then $\h^\bullet(D\sF) = \ext^\bullet(\sF,\sK_X)=\dZ/n$ in degree $0$, and
in the pairing $\h^\bullet(\sF)\otimes\h^\bullet(D\sF) \to \dZ/n$, one has
$1\otimes 1\mapsto 1$ (\hyperref[IV]{Cycle}, \ref{IV:2-1-5}): the duality is
perfect. The case in which $\sF$ has finite support is treated in the same
way.

For arbitrary constructible $\sF$, let $\sA$ be the subsheaf of its sections
with finite support, and let $\sG=\sF/\sA$. The exact sequence
\begin{equation*}\label{V:eq:2-5-1}\tag{2.5.1}
\xymatrix{
  0 \ar[r]
    & \sA \ar[r]
    & \sF \ar[r]
    & \sG \ar[r]
    & 0
}
\end{equation*}
gives a diagram
\begin{equation*}\tag{2.5.2}\label{V:eq:2-5-2}
\xymatrix@=.5cm{
  & & 0 \ar[r]
    & \h^0(X,\sA) \ar[r] \ar[d]^-\wr
    & \h^0(X,\sF) \ar[r] \ar[d]
    & \h^0(X,\sG) \ar[r] \ar[d]
    & 0 \\
  0 \ar[r]
  & '\h^{-1}(X,\sF) \ar[r]
  & '\h^{-1}(X,\sG) \ar[r]
  & '\h^0(X,\sA) \ar[r]
  & '\h^0(X,\sF) \ar[r]
  & '\h^0(X,\sG) \ar[r]
  & 0
}
\end{equation*}
and compatible isomorphisms $\h^i(X,\sF)\iso \h^i(X,\sG)$,
$'\h^i(X,\sF) \iso '\h^i(X,\sG)$ for $i\ne 0,-1$. If $j:U\to X$ is an
open subscheme on which $\sF$ is locally constant, there is an exact sequence
\begin{equation*}\tag{2.5.3}\label{V:eq:2-5-3}
\xymatrix{
  0 \ar[r]
    & \sG \ar[r]
    & j_\ast j^\ast \sF \ar[r]
    & \sB \ar[r]
    & 0
}
\end{equation*}
with $\sB$ of finite support. Since
$'\h^i(X,j_\ast j^\ast\sF) = D \h^{2-i} (X,j_\ast(j^\ast\sF)^\vee(1))$
by \ref{V:1-3}, one has $'\h^i=0$ for $i<0$; then
\eqref{V:eq:2-5-2} and the long exact sequence defined by
\eqref{V:eq:2-5-3} show that $'\h^i$ always vanishes for $i<0$.

If $\sF=\dZ/n$, the theorem is true for $i=0$ and $2$. This expresses the
fact that, for connected $X$, the trace morphism
$\h^2(X,\dmu_n) \iso \dZ/n$ is an isomorphism. By \ref{V:2-4}, the theorem
remains true for $i=0$ and $\sF=f_\ast\dZ/n$.

For arbitrary $\sF$ again, with $\sG$ as above, there exists
$f:X'\to X$ as in \ref{V:2-4} such that $\sG$ embeds into
$f_\ast\dZ/n$:
\[\xymatrix{
  0 \ar[r]
    & \sG \ar[r]
    & f_\ast\dZ/n \ar[r]
    & \sQ \ar[r]
    & 0 \text{.}
}\]
After replacing $X'$ by an $X''/X'$, we may assume that
$\h^i(X,\sF) \to \h^i(X,f_\ast\dZ/n) = \h^i(X',\dZ/n)$ is zero for $i>0$.
Consider then
\small
\[\xymatrix@=.3cm{
  0 \ar[r]
    & \h^0(X,\sG) \ar[r] \ar[d]
    & \h^0(X',\dZ/n) \ar[r] \ar[d]^-\wr
    & \h^0(X,\sQ) \ar[r] \ar[d]
    & \h^1(X,\sG) \ar[r]^-0 \ar[d]
    & \h^1(X',\dZ/n) \ar[r] \ar[d]
    & \h^1(X,\sQ) \ar[r] \ar[d]
    & \cdots \\
  0 \ar[r]
    & '\h^0(X,\sG) \ar[r]
    & '\h^0(X',\dZ/n) \ar[r]
    & '\h^0(X,\sQ) \ar[r]
    & '\h^1(X,\sG) \ar[r]
    & '\h^1(X',\dZ/n) \ar[r]
    & '\h^1(X,\sQ) \ar[r]
    & \cdots
}\]
\normalsize

The isomorphism at $\h^0(X')$ gives
$\h^0(X,\sG)\hookrightarrow '\h^0(X,\sG)$, and the same injectivity for
$\sF$ by \eqref{V:eq:2-5-2}. Applied to $\sQ$, this gives
$\h^1(X,\sG)\hookrightarrow '\h^1(X,\sG)$. We apply this to
$\h^1(X',\dZ/n)$. This group has the same order as
$'\h^1(X',\dZ/n)=D(\h^1(X',\dZ/n)(1))$, so we obtain an isomorphism, and
$\h^1(X,\sG) \iso '\h^1(X,\sG)$. The same holds for $\sF$. For $i=2$, we
already know that $\h^2(X',\dZ/n) \iso '\h^2(X',\dZ/n)$, and proceed in the
same way, then stop for lack of combatants.


\section{Poincaré duality for curves}\label{V:3}

In this section we prove the compatibility
(\hyperref[I]{Arcata} \ref{I:eq:6-2-3-3}).


\subsection{}\label{V:3-1}

We use the notation of (\hyperref[I]{Arcata} \ref{I:6-2-3}): $\bar X$ is a
smooth, connected, projective curve over an algebraically closed field $k$,
$D$ is a reduced divisor on $\bar X$, and $0$ is a point of
$X=\bar X\setminus D$,
\[\xymatrix{
  X \ar@{^{(}->}[r]^-j
    & \bar X
    & \ar[l]_-i D \text{,}
}\]
and $\pic_D(\bar X)=\h^1(\bar X,_D\dG_m)$, where $_D\dG_m$ is the sheaf of
sections of $\dG_m$ congruent to $1\mod D$. Recall that
$\h_c^1(X,\dmu_n) = \pic_D^0(\bar X)_n$, and that $x\to \sO(x)$ defines a
canonical map $f:X\to \pic_D(\bar X)$. Put
$f_0(x)=f(x)-f(0):X\to \pic_D^0(\bar X)$.

Again write $j$ for the inclusion of $X\times X$ in $X\times \bar X$. The
diagonal $\Delta$ of $X\times X$ is closed in $X\times \bar X$; it defines a
class in
$\h_\Delta^2(X\times X,\dmu_n)=\h_\Delta^2(X\times \bar X,j_!\dmu_n)$. Let
$c$ denote its image in $\h^2(X\times \bar X,j_!\dmu_n)$.

The K\"unneth formula gives
\[
  \h^\bullet(X\times \bar X,j_!\dmu_n) = \h^\bullet(X,\dZ/n)\otimes \h^\bullet(\bar X,j_!\dmu_n) = \h^\bullet(X,\dZ/n)\otimes \h_c^\bullet(X,\dmu_n) \text{.}
\]
We shall compute the $(1,1)$-component of $c$,
\[
  c^{1,1} \in \h^1(X,\dZ/n)\otimes \h_c^1(X,\dmu_n) = \h^1(X,\h_c^1(X,\dmu_n)) = \h^1\left(X,\pic_D^0(\bar X)_n\right) \text{.}
\]

Consider the exact sequence
\[\xymatrix{
  0 \ar[r]
    & \pic_D^0(\bar X)_n \ar[r]
    & \pic_D^0(\bar X) \ar[r]^n
    & \pic_D^0(\bar X) \ar[r]
    & 0
}\]
This makes $\pic_D^0(\bar X)$ a torsor under $\pic_D^0(\bar X)_n$ over
$\pic_D^0(\bar X)$. Let $u$ be the class, in
$\h^1(X,\h_c^1(X,\dmu_n))=\h^1(X,\pic_D^0(\bar X)_n)$, of its inverse image
under $f_0$.


\begin{proposition_}\label{V:3-2}
With the preceding notation, $c^{1,1}=-u$.
\end{proposition_}

Let $e$ be the difference between the images in
$\h^2(X\times \bar X,j_!\dmu_n)$ of the cohomology classes of $\Delta$ and
$X\times\{0\}$. Since the latter is of K\"unneth type $(0,2)$, one has
$c^{1,1}=e^{1,1}$. Consider the Leray spectral sequence for
$\pr_1:X\times \bar X\to X$ and the sheaf $j_!\dmu_n$:
\[
  E_2^{pq} = \h^p(X,\eR^q {\pr_1}_\ast j_!\dmu_n) = \h^p(X,\h_c^q(X,\dmu_n)) \Rightarrow \h^{p+q}(X\times \bar X,j_!\dmu_n) \text{.}
\]
It degenerates: it is the tensor product of $\h^\bullet(X)$ with the
(trivial) Leray spectral sequence for $\bar X\to \spec(k)$ and the sheaf
$j_!\dmu_n$. Fiber by fiber, the divisor $\Delta-X\times\{0\}$ has degree
$0$, so the image of $e$ in $E^{02}$ is zero. We may therefore speak of its
image $\bar e$ in $E^{1,1}$, and we must show that $\bar e = -u$.

Let $\sG$ on $X\times\bar X$ be the subsheaf of $\dG_m$ consisting of local
sections whose restriction to the closed subscheme $X\times D$ is $1$.
There is again an exact sequence
\begin{equation*}\tag{3.2.1}\label{V:eq:3-2-1}
\xymatrix{
  0 \ar[r]
    & j_!\dmu_n \ar[r]
    & \sG \ar[r]
    & \sG \ar[r]
    & 0
}
\end{equation*}
and $e$ is the image under $\partial$ of the class
$e_1\in \h^1(X\times\bar X,\sG)$ of the invertible sheaf
$\sO(\Delta-X\times\{0\})$, trivialized by $1$ on $X\times D$.

Applying $\eR\pr_{1\ast}$ to the distinguished triangle defined by
\eqref{V:eq:3-2-1} gives a distinguished triangle
\begin{equation*}\tag{3.2.2}\label{V:eq:3-2-2}
\xymatrix{
  & \ar[r]^-\partial
    & \eR \pr_{1\ast} j_!\dmu_n \ar[r]
    & \eR \pr_{1\ast}\sG \ar[r]
    & \eR \pr_{1\ast}\sG \ar[r]^-\partial
    &
}
\end{equation*}
and the diagram
\[\xymatrix{
  \h^1(X\times \bar X,\sG) \ar[r]^-\partial \ar@{=}[d]^-{(1)}
    & \h^2(X\times \bar X,j_!\dmu_n) \ar@{=}[d] \\
  \h^1(X,\eR\pr_{1\ast} \sG) \ar[r]^-\partial
    & \h^2(X,\eR\pr_{1\ast} j_!\dmu_n)
}\]
is commutative. Our problem is thus to identify the image in
$E^{11}=\h^1(X,\h_c^1(X,\dmu_n))$ (the image in $E^{02}$ being zero) of the
image under
$\partial_{\text{\eqref{V:eq:3-2-2}}}:\h^1(X,\eR\pr_{1\ast}\sG) \to \h^2(X,\eR\pr_{1\ast} j_!\dmu_n)$
of the class, still denoted by $e_1$, corresponding to $e_1$ under (1).

The sheaf $\eR^1\pr_{1\ast}\sG$ is defined by the group scheme over $X$
obtained by pulling back the algebraic group $\pic_D(\bar X)$ over
$\spec(k)$, and the image of $e_1$ in
$\h^0(X,\eR^1\pr_{1\ast}\sG) = \hom(X,\pic_D(\bar X))$ is $f_0$.

Represent the triangle \eqref{V:eq:3-2-2} by a short exact sequence of
complexes of sheaves:
\[\xymatrix{
  0 \ar[r]
    & \sA \ar[r]
    & \sB \ar[r]
    & \sC \ar[r]
    & 0 \text{.}
}\]
Let $\sB'$ be the subsheaf of
\[
  \tau_{\leqslant 1} \sB : \cdots \sB^0 \to \ker(d) \to 0 \cdots
\]
obtained by replacing the degree-one term $\ker(d)$ by the inverse image, under
$\ker(d)\to \underline\h^1(\sB) = \pic_D(\bar X)$ over $X$, of
$\pic_D^0(\bar X)$ over $X$. Let $\sA'=\sA\cap \sB'$ and let $\sC'$ be the
image of $\sB'$. We have $\sA'=\tau_{\leqslant 1}(\sA)$, and $\sC'$ is
obtained from $\sC$ in the same way as $\sB'$ from $\sB$. Hence there is a
commutative diagram of short exact sequences
\begin{equation*}\tag{3.2.3}\label{V:eq:3-2-3}
\xymatrix{
  0 \ar[r]
    & \sA \ar[r]
    & \sB \ar[r]
    & \sC \ar[r]
    & 0 \\
  0 \ar[r]
    & \tau_{\leqslant 1} \sA \ar[r] \ar[u] \ar[d]
    & \sB' \ar[r] \ar[u] \ar[d]
    & \sC' \ar[r] \ar[u] \ar[d]
    & 0 \\
  0 \ar[r]
    & \h_c^1(X,\dmu_n)[-1] \ar[r]
    & \pic_D^0(\bar X)[-1] \ar[r]
    & \pic_D^0(\bar X)[-1] \ar[r]
    & 0
}
\end{equation*}
where the last row is to be viewed as an exact sequence of complexes of
sheaves concentrated in degree $1$ on $X$.
\begin{enumerate}[a)]
  \item $e\in \hh^2(X,\sA)$ comes from
    $\widetilde e\in \hh^2(X,\tau_{\leqslant 1}\sA)$ (because its image in
    $E^{02}$ is zero). We seek the image $\bar e$ of $\widetilde e$ in
    $\hh^2(X,\h_c^1(X,\dmu_n)[-1]) = \h^1(X,\h_c^1(X,\dmu_n))$.
  \item $e_1\in \hh^1(X,\sC)$ comes from
    $\widetilde e_1\in \hh^1(X,\sC')$, because its image $f_0$ in
    $\h^0(X,\pic_D(\bar X))$ lies in $\h^0(X,\pic_D^0(\bar X))$. We may
    take $\widetilde e=\partial \widetilde e_1$.
  \item Thus $\bar e = \partial f_0$, where $\partial$ is defined by
    \eqref{V:eq:3-2-2}, and where we use the usual isomorphisms
    $\hh^2(X,\sF[-1]) = \h^1(X,\sF)$. These isomorphisms anticommute with
    $\partial$, so that $\bar e=-\partial f_0$, for $\partial$ defined by
    the exact sequence
    \[\xymatrix{
      0 \ar[r]
        & \h_c^1(X,\dmu_n) \ar[r]
        & \pic_D^0(\bar X) \ar[r]
        & \pic_D^0(\bar X) \ar[r]
        & 0
    }\]
    The conclusion follows from (\hyperref[IV]{Cycles}, \ref{IV:1-1-1}).
\end{enumerate}


\subsection{}\label{V:3-3}

For every homomorphism $\varphi:\h_c^1(X,\dmu_n)\to \dZ/n$, let
$\varphi(u)\in \h^1(X,\dZ/n)$ be the image of $u$ under
$\varphi:\h^1(X,\h_c^1(X,\dmu_n)) \to \h^1(X,\dZ/n)$. If
$x\in \h_c^1(X,\dmu_n)$, the cup product $\varphi(u)\smallsmile x$ lies in
$\h_c^2(X,\dmu_n)$. We shall prove the following.


\begin{proposition_}\label{V:3-4}
$\tr(\varphi(u)\smallsmile x) = \varphi(x)$.
\end{proposition_}

This identity follows, after applying $\varphi$, from
\begin{equation*}\tag{3.4.1}\label{V:eq:3-4-1}
  \tr(u\smallsmile x) = x
\end{equation*}
where $\tr$ now denotes the map
\[\xymatrix{
  \tr:\h_c^2(X,\h_c^1(X,\dmu_n)\otimes\dmu_n) \ar[r]
    & \h_c^1(X,\dmu_n) \text{,}
}\]
so that $x\mapsto \tr(u\smallsmile x)$ is the composite
\small
\begin{equation*}\tag{3.4.2}\label{V:eq:3-4-2}
\begin{aligned}
&\xymatrix{
  \h_c^1(X,\dmu_n) \ar[r]^-{u\otimes x}
    & \h^1(X,\h_c^1(X,\dmu_n)) \otimes \h_c^1(X,\dmu_n)
    & \ar[l]_-\sim \h^1(X,\dZ/n)\otimes \h_c^1(X,\dmu_n) \otimes \h_c^1(X,\dmu_n)} \\ &\xymatrix{
    \ar[r]^-{(1)}
      & \h_c^2(X,\dmu_n) \otimes \h_c^1(X,\dmu_n) \ar[r]^-{\tr\otimes\operatorname{id}}
      & \h_c^1(X,\dmu_n)
}
\end{aligned}
\end{equation*}
\normalsize
where (1) is given by
$a\otimes b\otimes c\mapsto (a\smallsmile c)\otimes b$.

We express this morphism in terms of the trace morphism
$\tr:\h_c^3(X\times X,\dmu_n^{\otimes 2}) \to \h_c^2(X,\dmu_n)$ relative to
the second projection. For $x\in \h_c^3$, $\tr(x)$ is computed as follows:
among its K\"unneth components, retain only the component of type $(2,1)$,
and apply $\tr:\h_c^2(X,\dmu_n) \to \dZ/n$.

Let $u'$ be the image of $u$ under the map
\begin{align*}
&\xymatrix{
  \h^1(X,\h_c^1(X,\dmu_n))
    & \ar[l]_-\sim \h^1(X,\dZ/n) \otimes \h_c^1(X,\dmu_n) = \h^1(X,\dZ/n)\otimes \h^1(\bar X,j_!\dmu_n)
} \\
&\xymatrix{
  \ar[r]
    & \h^2(X\times \bar X,\dZ/n\boxtimes j_!\dmu_n) \text{.}
}
\end{align*}
In view of the relation between the K\"unneth formula and cup product, one
has
\[
  \tr(u\smallsmile x) = -\tr(u'\smallsmile \pr_1^\ast x) \text{.}
\]
The second cup product is the one in
(\hyperref[IV]{Cycle} \ref{IV:1-2-4}):
\[\xymatrix{
  \h^2(X\times \bar X,\dZ/n\boxtimes j_!\dmu_n) \otimes \h^1(\bar X\times \bar X,j_!\dmu_n\boxtimes \dZ/n) \ar[r]
    & \h^3(\bar X\times \bar X,j_!\dmu_n \boxtimes j_!\dmu_n) \text{.}
}\]
The minus sign comes from permuting two symbols of degree $1$ in (1). Since
$c\smallsmile\pr_1^\ast x$ and $c^{1,1}\smallsmile \pr_1^\ast x$ have the
same $(2,1)$-component in the K\"unneth decomposition, and since
$u'=-c^{1,1}$, one has
\[
  \tr(u\smallsmile x) = \tr(c\smallsmile \pr_1^\ast x) \text{.}
\]
The cup product on the right can now be interpreted as the cup product
(\hyperref[IV]{Cycle} \ref{IV:1-2-5}) of
$\cl(\Delta)\in \h_\Delta^2(X\times X,\dmu_n)$ with
$\pr_1^\ast x\in \h_c^1(\Delta,\dmu_n)$. On $\Delta$, one has
$\pr_1=\pr_2$, whence
\[
  \tr(c\smallsmile \pr_1^\ast x) = \tr(\cl(\Delta)\smallsmile \pr_1^\ast x) = \tr(\cl(\Delta)\smallsmile \pr_2^\ast x) = \tr(\cl(\Delta))\smallsmile x = x \text{.}
\]


\section{Compatibility with SGA 4 XVIII 3.1.10.3}\label{V:4}

In \cite[XVIII]{sga4}, overcome by discouragement, I did not verify the
compatibility announced in the title. Although it turns out to be no more
than smoke, I feel bound to remedy that omission here. For $f:X\to S$
separated and of finite type, the point was to verify that the adjunction
morphism $\eR f_!\eR f^!\sL\to \sL$ is compatible with every étale
localization $k:V\to S$. Using $k$ and $f$ for several arrows, as below, we
must verify commutativity:
\[\xymatrix{
  X_V \ar[r]^-k \ar[d]^-f
    & X \ar[d]^-f \\
  V \ar[r]^-k
    & S
} \qquad
\xymatrix{
  k^\ast \eR f_! \eR f^!\sL \ar[r]^-{k^\ast(\text{adj})} \ar[d]^-\sim
    & k^\ast \sL \ar@{=}[d] \\
  \eR f_! \eR f^! k^\ast \sL \ar[r]^-{\text{adj}}
    & k^\ast \sL \text{.}
}\]
At the level of complexes, the required commutativity is written (loc. cit.)
\[\xymatrix{
  k^\ast f_!^\bullet f^!_\bullet \sL \ar[r] \ar[d]^-{(1)}
    & k^\ast \sL \ar@{=}[d] \\
  f_!^\bullet f^!_\bullet k^\ast \sL \ar[r]
    & k^\ast \sL \text{.}
}\]
It is checked component by component, reducing the problem to a
compatibility for $\sL$ a sheaf and for each pair of adjoint functors
$(f_i^!,f_!^i)$, denoted simply by $f^!$ and $f_!$. The arrow (1) is defined
using the isomorphism $(2):k^\ast f_!=f_! k^\ast$, and a morphism
$(3):k^\ast f^! \to f^! k^\ast$ deduced from it. In loc. cit., one first
introduces a morphism $(4):k_! f_!\to f_! k_!$, deduced from $(2)$ by the
adjunction of $k^\ast=k^!$ and
$k_!:k_! f_!\to k_! f_! k^\ast k_!\to k_! k^\ast f_! k_! \to f_! k_!$, and
then deduces $(3)$ from $(4)$ by adjunction. Equivalently, one may deduce
$(3)$ from $(2)$ by the adjunction of $f_!$ and
$f^!:k^\ast \to f^! f_! k^\ast f^!\to f^! k^\ast f_! f^! \to f^! k^\ast$.

Writing $k$ for $k^\ast$ and omitting $\sL$, the required commutativity is
that of the outer boundary of
\[\xymatrix@=.5cm{
  k f_! f^! \ar[r] \ar@{=}[d]
    & f_! k f^! \ar[r]
    & f_! f^! f_! k f^! \ar[r]
    & f_! f^! k f_! f^! \ar@{=}[d] \\
  k f_! f^! \ar[d]
    & & & \ar[lll] f_! f^! k f_! f^! \ar[d] \\
  k
    & & & \ar[lll] f_! f^! k \text{.}
}\]
In the first row, we find a homomorphism
$k f_! f^! \to f_! f^! k f_! f^!$, defined because $k f_! f^!$ is of the
form $f_! X$ (with $X=k f^!$). It admits the adjunction morphism in the
second row as a retraction, and the commutativity of the lower square is the
functoriality of this adjunction morphism.

% END INCLUDED components/sga4half_5_en.tex

% BEGIN INCLUDED components/sga4half_6_en.tex
% !TEX root = ../SGA4half_English_Reader_1_to_8.tex

\chapter{Applications of the trace formula to trigonometric sums}\label{VI}
\chaptermark{Applications to trigonometric sums}










In this exposé, I explain how the trace formula can be used to calculate 
or study various trigonometric sums and how, together with the 
Weil conjecture, it can be used to bound them.

The first two sections provide a ``user's guide'' to these tools. 
\hyperref[VI:3]{Section 3} presents, in cohomological language, Weil's results on sums in 
$1$ variable. Sections $4$ through $6$ give a 
detailed study of Gauss and Jacobi sums -- including 
Weil's old and recent results on Hecke characters 
defined by Jacobi sums. In Section $7$, we study a 
generalization of Kloosterman sums to several variables. Finally, 
Section $8$ gives some indications of other uses that have been, or can be, 
made of these methods. 










\section*{Notations}\label{VI:0}





\subsection{}\label{VI:0-1}

We use the notation of \hyperref[II]{Report}, Section \ref{II:1}. We 
will often have to consider a finite extension $\dF_{q^n}\subset \dF$ of 
$\dF_q$. A subscript $0$ will denote an object over $\dF_q$, and a subscript $1$ 
an object over $\dF_{q^n}$. Replacing a subscript $0$ by a subscript $1$ (resp. 
omitting the subscript) means extending scalars to $\dF_{q^n}$ (resp. 
to $\dF$).





\subsection{}\label{VI:0-2}

We denote by $\ell$ a prime number $\ne p$. We freely use the 
language of $\dQ_\ell$-sheaves, as well as that of $E_\lambda$-sheaves, 
where $E_\lambda$ is a finite extension of $\dQ_\ell$ (cf. \hyperref[II]{Report}, 
Section \ref{II:2}, especially \ref{II:2-11}). 





\subsection{}\label{VI:0-3}

Let $H^\bullet$ be a graded vector space. If $T$ is an endomorphism of 
$H^\bullet$, we put (cf. \hyperref[IV]{Cycle}, \ref{IV:eq:1-3-5}) 
\[
  \tr(T,H^\bullet) = \sum (-1)^i \tr(T,H^i) \text{.}
\]










\section{Principles}\label{VI:1}





\subsection{}\label{VI:1-1}

Let $X_0$ be a separated scheme of finite type over $\dF_q$, let $E_\lambda$ be a 
finite extension of $\dQ_\ell$, and let $\sF_0$ be an $E_\lambda$-sheaf on $X_0$. The 
trace formula says that 
\begin{equation*}\tag{1.1.1}\label{VI:eq:1-1-1}
  \sum_{x\in X^F} \tr(F_x^\ast,\sF) = \sum (-1)^i \tr\left(F^\ast,\h_c^i(X,\sF)\right) \text{.}
\end{equation*}

With the notation \ref{VI:0-3}, the right-hand side is simply written 
$\tr(F^\ast,\h_c^\bullet(X,\sF))$. 

This trace formula for $E_\lambda$-sheaves can either be 
deduced from \hyperref[II]{Report} \ref{II:4-10} by passage to the limit (cf. 
\hyperref[II]{Report} \ref{II:4-11} through \ref{II:4-13}), or be deduced 
from the trace formula for $\dQ_\ell$-sheaves (\hyperref[II]{Report} 
\ref{II:3-2}) by the method of (\hyperref[III]{Functions $L$ mod. $\ell^n$}, 
\ref{III:4-3}). 

We shall interpret various trigonometric sums as the left-hand side of 
\eqref{VI:eq:1-1-1}, for suitable $X_0$ and $\sF_0$. 





% NOTE: in this section, following the source (but not convention in the 
% LaTeX version, torsors are denoted in roman (not script) letters
\subsection{}\label{VI:1-2}

Let $A$ be a finite group. An \emph{$A$-torsor} on a scheme $X$ is a 
sheaf $T$ on $X$, equipped with a right action of $A$, which, locally 
(for the étale topology) on $X$, is isomorphic to the constant sheaf $A$ on 
which $A$ acts by right translation. 

If $\tau:A\to B$ is a homomorphism and $T$ is an $A$-torsor, there exists, 
up to unique isomorphism, a unique $B$-torsor $\tau(T)$, equipped with a 
morphism of sheaves $\tau:T\to \tau(T)$ such that 
\begin{equation*}\tag{1.2.1}\label{VI:eq:1-2-1}
  \tau(t a) = \tau(t) \tau(a)\text{.}
\end{equation*}

If $T$ is an $A$-torsor on $X$ and $\rho$ is a linear representation 
of $A$, $\rho:A\to \gl(V)$ (where $V$ is a finite-dimensional vector space over 
$E_\lambda$), there exists, up to unique isomorphism, a unique 
$E_\lambda$-sheaf $\sF$, smooth of rank $\dim(V)$, equipped with a morphism of 
sheaves 
\begin{equation*}\tag{1.2.2}\label{VI:eq:1-2-2}
  \rho:T\to \underline{\operatorname{Isom}}(V,\sF) \text{.}
\end{equation*}
such that $\rho(t a) = \rho(t) \rho(a)$. We denote it by $\rho(T)$. For 
$\tau:A\to B$ and $\rho$ a representation of $B$, we have a canonical 
isomorphism $\rho(\tau(T)) = (\rho\tau)(T)$. 

In this paragraph (except the appendix), we only consider the case where $A$ 
is commutative and where $V=E_\lambda$: $\rho$ is a character 
$A\to E_\lambda^\times$, and $\rho(T)$ is a smooth $E_\lambda$-sheaf of rank 
one. The morphism \eqref{VI:eq:1-2-2} is interpreted as an everywhere 
nonzero morphism 
\begin{equation*}\tag{1.2.3}\label{VI:eq:1-2-3}
  \rho:T\to \rho(T) 
\end{equation*}
such that $\rho(t a) = \rho(a)\rho(t)$. 





\subsection{}\label{VI:1-3}

Let $S$ be a scheme, let $G$ be a commutative group scheme over $S$, and let $G'$ 
be an extension of commutative group schemes over $S$
\[\xymatrix{
  0 \ar[r] 
    & A \ar[r] 
    & G' \ar[r]^-\pi 
    & G \ar[r] 
    & 0 \text{.}
}\]
The sheaf $T$ of local sections of $\pi$ is then an $A$-torsor on 
$G$. 

If $X$ is a scheme over $S$ and $f,g$ are two $S$-morphisms from $X$ to $G$, 
the fact that $T$ is defined by an extension gives a canonical isomorphism 
\begin{equation*}\tag{1.3.1}\label{VI:eq:1-3-1}
  (f+g)^\ast T = f^\ast T + g^\ast T \text{.}
\end{equation*}
On the left, $f+g$ is the sum of $f$ and $g$ with respect to the group law on $G$; 
on the right, $+$ denotes a sum of torsors. 

If $f:X\to G$ factors through $G'$, $f^\ast T$ is trivial (and a 
factorization gives a trivialization); up to (non-unique) isomorphism, 
the $A$-torsor $f^\ast T$ depends only on the image of $f$ in 
$\hom_S(X,G)/\pi\hom_S(X,G')$: it is given by the morphism $\partial$ in 
the exact sequence 
\[\xymatrix{
  \hom(X,G') \ar[r] 
    & \hom(X,G) \ar[r]^-\partial 
    & \h^1(X,A) \text{.}
}\]

We will apply these constructions to the exact Kummer sequences 
$0 \to \dmu_n \to \dG_m \to \dG_m \to 0$, and to Lang torsors. 





\subsection{The Lang torsor}\label{VI:1-4}

Let $G_0$ be a connected commutative algebraic group over $\dF_q$ (for the 
noncommutative case, see the appendix). The group law is written 
\emph{multiplicatively}. Lang's isogeny 
\[
\fL:G_0 \to G_0: x\mapsto F x\cdot x^{-1}
\]
is étale; its image, an open subgroup of $G_0$, can only be 
$G_0$ itself; its kernel is the finite group $G_0(\dF_q)$. The \emph{Lang 
torsor} $L$ is the $G_0(\dF_q)$-torsor on $G_0$ defined by the exact sequence 
\begin{equation*}\tag{1.4.1}\label{VI:eq:1-4-1}
\xymatrix{
  0 \ar[r] 
    & G_0(\dF_q) \ar[r] 
    & G_0 \ar[r]^-\fL 
    & G_0 \ar[r] 
    & 0 \text{.}
}
\end{equation*}

\paragraph{Notation}
We also write $L$ and $\fL$ (or $L_{(q)}$ and $\fL_{(q)}$) for the objects obtained 
from $L$ and $\fL$ by extending the base field. When necessary, we 
specify them with a subscript $0$ or $1$.


\paragraph{Examples}
For $G_0=\dG_a$ or $\dG_m$, the sequences \eqref{VI:eq:1-4-1} are written 
\begin{align*}\tag{1.4.2}\label{VI:eq:1-4-2}
&\xymatrix{
  0 \ar[r] 
    & \dF_q \ar[r] 
    & \dG_a \ar[r]^-{x^q-x} 
    & \dG_a \ar[r] 
    & 0 
} \\
\tag{1.4.3}\label{VI:eq:1-4-3}
&\xymatrix{
  0 \ar[r] 
    & \dmu_{q-1} \ar[r] 
    & \dG_m \ar[r]^-{x^{q-1}} 
    & \dG_m \ar[r] 
    & 0
}
\end{align*}





\subsection{}\label{VI:1-5}

Let us compute the endomorphism $F^\ast$ of the fiber of the Lang torsor at 
$\gamma\in G^F=G_0(\dF_q)$. If $g\in \fL^{-1}(\gamma)\subset G$, we have 
$F g = F g\cdot g^{-1}\cdot g = \gamma g$. Therefore (\hyperref[II]{Report}, 
\ref{II:1-2}) 
\begin{equation*}\tag{1.5.1}\label{VI:eq:1-5-1}
  \text{on $L_0(G_0)_\gamma\simeq \fL^{-1}(\gamma)$, $F^\ast$ is $g\mapsto g\gamma^{-1}$.}
\end{equation*}





\subsection{}\label{VI:1-6}

On $\dF_{q^n}$, the identity $F_{(q^n)}=F_{(q)}^n$ between endomorphisms of 
$G_1$ implies that $\fL_{(q^n)}=\fL_{(q)}\circ \prod_{i=0}^{n-1} F_{(q)}^i$. 
On $G_0(\dF_{q^n})$, $F_{(q)}^i$ acts as the element 
$x\mapsto x^{q^i}$ of $\gal(\dF_{q^n}/\dF_q)$. On $G_0(\dF_{q^n})$, 
$\prod F_{(q)}^i$ is therefore the composite of the norm 
$G_0(\dF_{q^n})\to G_0(\dF_q)$ and the inclusion 
$G_0(\dF_q)\subset G_0(\dF_{q^n})$: the diagram 
\begin{equation*}\tag{1.6.1}\label{VI:eq:1-6-1}
\xymatrix{
  0 \ar[r] 
    & G_0(\dF_{q^n}) \ar[r] \ar[d]^-N 
    & G_1 \ar[r]^-{\fL_{(q^n)}} \ar[d]^-{\prod_{i=0}^{n-1} F_{(q)}^i} 
    & G_1 \ar[r] \ar@{=}[d] 
    & 0 \\
  0 \ar[r] 
    & G_0(\dF_q) \ar[r] 
    & G_1 \ar[r]^-{\fL_{(q)}} 
    & G_1 \ar[r] 
    & 0
}
\end{equation*}
is commutative. In the language of torsors, this gives a canonical isomorphism 
between $G_0(\dF_q)$-torsors on $G_1$ 
\begin{equation*}\tag{1.6.2}\label{VI:eq:1-6-2}
  N L_{(q^n)} = L_{(q)} \text{.}
\end{equation*}





\begin{definition_}\label{VI:1-7}
Let $f_0:X_0\to G_0$ be a morphism and 
$\chi:G_0(\dF_q)\to E_\lambda^\times$ a character. We set 
$\sF(\chi,f_0) = \chi^{-1}(f_0^\ast (L_0(G_0))) = f_0^\ast\chi^{-1}(L_0(G_0))$. 
\end{definition_}

We have the following functoriality properties:
\begin{enumerate}[a)]
  \item $\sF(\chi,f_0)$ is bimultiplicative in $\chi$ and $f_0$: 
    \begin{align*}\tag{1.7.1}\label{VI:eq:1-7-1}
      \sF(\chi,f_0'\cdot f_0'') &= \sF(\chi,f_0')\otimes \sF(\chi,f_0'') \text{,} \\
      \tag{1.7.2}\label{VI:eq:1-7-2}
      \sF(\chi'\chi'',f_0) &= \sF(\chi',f_0)\otimes \sF(\chi'',f_0) \text{.}
    \end{align*}
    This follows from the definition of $\chi(T)$, for $T$ a torsor, combined 
    with \eqref{VI:eq:1-3-1} for \eqref{VI:eq:1-7-1}. 
  \item For $g_0:Y_0\to X_0$ a morphism, we have 
    \begin{equation*}\tag{1.7.3}\label{VI:eq:1-7-3}
      \sF(\chi,f_0\circ g_0) = g_0^\ast \sF(\chi,f_0) \text{.}
    \end{equation*}
    In particular, $\sF(\chi,f_0) = f_0^\ast \sF(\chi)$, where $\sF(\chi)$ 
    denotes $\sF(\chi,\operatorname{id}_{G_0})$. 
  \item For $u_0:G_0\to H_0$ a morphism, and 
    $\chi:H_0(\dF_q) \to E_\lambda^\times$, we have 
    \begin{equation*}\tag{1.7.4}\label{VI:eq:1-7-4}
      \sF(\chi,u_0 f_0) = \sF(\chi u_0,f_0) \text{.}
    \end{equation*}
  \item For $G_0=\prod_{i\in I} G_0^i$, with $\chi$ having coordinates $\chi_i$ 
    ($i\in I$) and $f_0$ having coordinates $f_0^i$, it follows from 
    \eqref{VI:eq:1-7-1} through \eqref{VI:eq:1-7-4} that 
    \begin{equation*}\tag{1.7.5}\label{VI:eq:1-7-5}
      \sF(\chi,f_0) = \bigotimes_{i\in I} \sF(\chi_i,f_0^i) \text{.}
    \end{equation*}    
\end{enumerate}

Note the $\chi^{-1}$ in Definition \ref{VI:1-7}. It ensures that, for 
$x\in X^F$, on $\sF(\chi,f_0)_x$ one has 
\begin{equation*}\tag{1.7.6}\label{VI:eq:1-7-6}
  F_x^\ast = \chi f_0(x)
\end{equation*}
(use the fact that the fiber at $x$ of the morphism \eqref{VI:eq:1-2-2}, for 
$\rho=\chi^{-1}$, commutes with $F_x^\ast$, together with \eqref{VI:eq:1-5-1}). 

If $\sF(\chi,f_0)_1$ is the sheaf obtained from $\sF(\chi,f_0)$ by extension 
from the base field to $\dF_{q^n}$, we deduce from \eqref{VI:eq:1-6-2} that 
\begin{equation*}\tag{1.7.7}\label{VI:eq:1-7-7}
  \sF(\chi,f_0)_1 = \sF(\chi\circ N,f_1) \text{.}
\end{equation*}





\subsection{Abuse of notations}\label{VI:1-8}

(i) If $\Xi$ denotes the composite map 
$\chi f_0:X_0(\dF_q) \to E_\lambda^\times$, we will sometimes write $\sF(\Xi)$ 
instead of $\sF(\chi,f_0)$. By \eqref{VI:eq:1-7-1} through \eqref{VI:eq:1-7-5}, 
there is little risk of ambiguity. For example:
\begin{enumerate}[a)]
  \item we write $\sF(\chi)$ for $\sF(\chi,\operatorname{id}_{G_0})$ 
    (notation already used in \ref{VI:1-7}b)); 
  \item we write $\sF(\chi f_0)$ for $\sF(\chi,f_0)$;
  \item with the notation of \ref{VI:1-7}d), we write 
    $\sF(\prod \chi_i f_0^i)$ for $\sF(\chi,f_0)$. 
\end{enumerate}

With this notation, \eqref{VI:eq:1-7-4} shows that the notation 
$\sF(\chi u_0 f_0)$ is unambiguous, \eqref{VI:eq:1-7-1} \eqref{VI:eq:1-7-2} 
\eqref{VI:eq:1-7-5} express the multiplicativity of $\sF(\Xi)$ with respect to $\Xi$, 
and \eqref{VI:eq:1-7-6} becomes $F_x^\ast = \Xi(x)$ on $\sF(\Xi)_x$.

(ii) If $X$ is a scheme over an extension $k$ of $\dF_q$, and $f$ is a 
morphism of $X$ into $G_0\otimes_{\dF_q} k$, we shall again denote by $\sF(\chi,f)$, 
$\sF(\chi f)$, or $\sF(\chi)$ (for $f$ an inclusion) the inverse image under 
$f$ of the $E_\lambda$-sheaf obtained from $\chi^{-1}(L_0(G))$ by extension of 
scalars from $\dF_q$ to $k$. If $f_0$ is the composite 
$X\to G_0\otimes_{\dF_q} k \to G_0$, this is again the $\sF(\chi,f_0)$ of 
\ref{VI:1-7}. If $k$ is a finite extension of $\dF_q$, we have 
$\sF(\chi f) = \sF(\chi N_{k/\dF_q} f)$. 

Applying \eqref{VI:eq:1-1-1} to \eqref{VI:eq:1-7-6} and \eqref{VI:eq:1-7-7}, 
we find: 





% NOTE: originally a 'scholium'
\begin{theorem_}\label{VI:1-9}
Let $S_0$ be a separated scheme of finite type over $\dF_q$, let $G_0$ be a 
connected commutative algebraic group over $\dF_q$, let $f_0:S_0\to G_0$ be a morphism, and let 
$\chi:G_0(\dF_q) \to E_\lambda^\times$. We have 
\begin{equation*}\tag{1.9.1}\label{VI:eq:1-9-1}
  \sum_{s\in S_0(\dF_q)} \chi f_0(s) = \tr(F^\ast,\h_c^\bullet(S,\sF(\chi,f_0))) 
\end{equation*}
and, for any integer $n\geqslant 1$,
\begin{equation*}\tag{1.9.2}\label{VI:eq:1-9-2}
  \sum_{s\in S_0(\dF_{q^n})} \chi\circ N_{\dF_{q^n}/\dF_q} f_0(s) = \tr({F^\ast}^n, \h_c^\bullet\left(S,\sF(\chi,f_0))\right) \text{.}
\end{equation*}
\end{theorem_}





\addtocounter{subsubsection}{2}
\subsubsection{Remark}\label{VI:1-9-3}

Take $G_0$ to be a product. With the notation of \ref{VI:1-7}d) and 
\ref{VI:1-8}c), the formula \eqref{VI:eq:1-9-2} becomes 
\[
  \sum_{s\in S_0(\dF_{q^n})} \prod_i \chi_i\circ N_{\dF_{q^n}/\dF_q} (f_0^i(s)) = \tr\Big({F^\ast}^n, \h_c^\bullet\Big(S,\sF\Big(\prod_i \chi_i f_0^i\Big)\Big)\Big) \text{.}
\]





\subsubsection{Remark}\label{VI:1-9-4}

Take $G_0=\{e\}$. Formula \eqref{VI:eq:1-9-1} becomes 
\[
  \left|S_0(\dF_q)\right| = \sum_{s\in S_0(\dF_q)} 1 = \tr(F^\ast,\h_c^\bullet(S,\dQ_\ell)) \text{.}
\]

It is rare to be able to compute the right-hand side of 
\eqref{VI:eq:1-9-1} explicitly. Here is an amusing example, with $G_0=\{e\}$. 





\subsection{Example}\label{VI:1-10}

A nonsingular projective quadric of odd dimension over \texorpdfstring{$\dF_q$}{Fq} has the same number of rational points as projective space of the same dimension over \texorpdfstring{$\dF_q$}{Fq}.

If, over $\dC$, $X$ is a nonsingular quadric hypersurface in 
the projective space $\dP^{2 N}$, and $Y$ a hyperplane, we know that for 
$i\leqslant 2\dim(X)=2\dim(Y)$, the inclusions 
$X\hookrightarrow \dP^{2 N}\hookleftarrow Y$ induce isomorphisms (in 
ordinary cohomology) 
\[\xymatrix{
  \h^i(X,\dQ) 
    & \ar[l]_-\sim \h^i(\dP^{2 N},\dQ) \ar[r]^-\sim 
    & \h^i(Y,\dQ) \text{.}
}\]

By specialization, it follows that if $X_0'$ is a nonsingular quadric 
hypersurface in projective space $\dP_0^{2 N}$ over $\dF_q$, 
and $Y_0'$ a hyperplane, the inclusions 
$X_0'\hookrightarrow \dP_0^{2 N}\hookleftarrow Y_0'$ induce 
isomorphisms 
\[\xymatrix{
  \h^i(X',\dQ_\ell) 
    & \ar[l]_-\sim \h^i(\dP^{2 N},\dQ_\ell) \ar[r]^-\sim 
    & \h^i(Y',\dQ_\ell) \text{.}
}\]
These isomorphisms commute with $F^\ast$, and the trace formula applies. 





\subsection{}\label{VI:1-11}

The trace formula can also be used to understand the classical formulas 
giving the number of rational points of linear groups over 
finite fields, and of certain homogeneous spaces (see \hyperref[VI:8]{Section 8}). 





\subsection{}\label{VI:1-12}

There is a dictionary that translates various kinds of classical 
manipulations of trigonometric sums into cohomological terms. This dictionary 
will be given in \hyperref[VI:2]{Section 2}. Since the cohomological statement is more 
``geometric,'' it can sometimes be applied 
to situations where the classical argument only applies after a 
extension of the finite base field. Here is an example (a special case of a 
theorem of Kazhdan).

\begin{theorem_}[Kazhdan]\label{VI:1-13}
Let $X_0$ be a scheme over $\dF_q$. Suppose there is an action $\rho$ of 
$\dG_a$ on $X$, and a morphism $f:X\to Y$ of schemes over $\dF$, making 
$X$ a $\dG_a$-torsor over $Y$. Then the number of rational points of $X_0$ 
is divisible by $q$.
\end{theorem_}

First suppose that $\rho$, $Y$, and $f$ are defined over $\dF_q$. 


\paragraph{Classical argument}
The fibers of $f:X_0(\dF_q)\to Y_0(\dF_q)$ all have $q$ elements: 
$|X_0(\dF_q)| = q\cdot |Y_0(\dF_q)|$, hence the divisibility. 


\paragraph{Cohomological translation}
The higher direct image sheaves $\eR^i f_!\dQ_\ell$ are 
\[
  \eR^i f_!\dQ_\ell = \begin{cases}
                       \dQ_\ell(-1) & \text{if $i = 2$} \\
                       0            & \text{for $i\ne 2$}
                     \end{cases}
\]
The Leray spectral sequence for $f$ therefore degenerates, giving an isomorphism 
\[
  \h_c^i(X,\dQ_\ell) = \h_c^{i-2}(Y,\dQ_\ell)(-1) \text{.}
\]
To understand the effect of a Tate twist on the eigenvalues of 
Frobenius, it is most convenient to use the Galois point of view 
(\hyperref[II]{Report} \ref{II:1-8}) and note that the Frobenius substitution 
$\varphi$ acts on $\dQ_\ell(1)$ by multiplication by $q$. The 
eigenvalues of $F^\ast$ on $\h_c^p(X,\dQ_\ell)$ are therefore obtained by 
multiplying by $q$ the eigenvalues of $F^\ast$ acting on $\h_c^{i-2}(Y,\dQ_\ell)$. We 
know that the latter are algebraic integers \cite[XXI 5.2.2]{sga7}; hence 

\begin{equation*}\tag{$*$}\label{VI:eq:*}
  \substack{\displaystyle\text{the eigenvalues of $F^\ast$ acting on $\h_c^i(X,\dQ_\ell)$ are} \\
  \displaystyle\text{algebraic integers divisible by $q$.}}
\end{equation*}


\paragraph{Descent}
Return to the hypotheses of the theorem. For a suitable $n$, we may 
suppose that $\rho$, $Y$, and $f$ are defined over $\dF_{q^n}$. Since the 
Frobenius morphism relative to $\dF_{q^n}$ is the $n$-th power of the one 
relative to $\dF_q$, statement \eqref{VI:eq:*} for the scheme over 
$\dF_{q^n}$ obtained from $X_0$ by extension of scalars gives: 
\begin{equation*}\tag{$**$}\label{VI:eq:**}
  \substack{\displaystyle\text{the $n$-th powers of the eigenvalues of $F^\ast$ acting on} \\ \displaystyle\text{$\h_c^i(X,\dQ_\ell)$ are algebraic integers divisible by $q^n$.}}
\end{equation*}

This assertion implies \eqref{VI:eq:*}. The number of rational points of 
$X_0$, namely $\sum (-1)^i \tr(F^\ast,\h^i(X,\dQ_\ell))$, is therefore an 
algebraic integer divisible by $q$. It follows that, as a rational integer, 
it is divisible by $q$. 





\subsection{}\label{VI:1-14}

Formula \ref{VI:1-9-3} controls the dependence on $n$ of the trigonometric 
sum on the left-hand side. It implies identities between a 
trigonometric sum and those which are deduced from it by ``extension of 
scalars.'' The most famous of these identities is that of 
Hasse--Davenport. Let $\chi:\dF_q^\times \to \dC^\times$ and 
$\psi:\dF_q\to \dC^\times$ be characters, with $\psi$ nontrivial, and 
define the Gauss sum $\tau(\chi,\psi)$ by 
\begin{equation*}\tag{1.14.1}\label{VI:eq:1-14-1}
  \tau(\chi,\psi) = - \sum_{x\in \dF_q^\times} \psi(x) \chi^{-1}(x) \text{.}
\end{equation*}





\begin{theorem_}[Hasse-Davenport]\label{VI:1-15}
We have 
\begin{equation*}\tag{1.15.1}\label{VI:eq:1-15-1}
  \tau\left(\chi\circ N_{\dF_{q^n}/\dF_q},\psi\circ \tr_{\dF_{q^n}/\dF_q}\right) = \tau(\chi,\psi)^n \text{.}
\end{equation*}
\end{theorem_}

Let $E\subset \dC$ be a number field containing the values of $\chi$ and 
$\psi$, and let $E_\lambda$ be the completion of $E$ at a place of residue 
characteristic $\ell\ne p$. Identity \eqref{VI:eq:1-15-1} is an identity in 
$E$. Proving it in $\dC$ is equivalent to proving it in $E_\lambda$. We 
shall prove it in $E_\lambda$, regarding $\chi$ and $\psi$ as taking values 
in $E_\lambda^\times$.

Apply \ref{VI:1-9-3} with $X_0=\dG_m$ over $\dF_q$ and 
$G_0=\dG_m\times \dG_a$. The groups $\h_c^i(\dG_m,\sF(\chi^{-1}\psi))$ vanish for 
$i\ne 1$, and $\h_c^1$ has dimension $1$ (\ref{VI:4-2}). By 
\ref{VI:1-9-3}, 
$\tau(\chi\circ N_{\dF_{q^n}/\dF_q},\psi\circ \tr_{\dF_{q^n}/\dF_q})$ is 
the unique eigenvalue of $(F^\ast)^n$ on $\h_c^1$, and \eqref{VI:eq:1-15-1} 
follows. 

Similar examples are given in Weil \cite[App V]{we74}. 





\subsection{}\label{VI:1-16}

An $E_\lambda$-sheaf $\sF_0$ on $X_0$ is said to be \emph{pointwise pure of 
weight $n$} if for any closed point $x\in |X_0|$, the eigenvalues of 
$F_x^\ast$ (\hyperref[II]{Report} \ref{II:1-2}) are algebraic numbers 
all of whose complex conjugates have absolute value 
$q_x^{n/2}$, where $q_x$ is the number of elements of $k(x)$. The 
following theorem will be proved in \cite{de80}. 





\begin{theorem_}\label{VI:1-17}
If $\sF_0$ is pointwise pure of weight $n$, then for any eigenvalue $\alpha$ 
of the endomorphism $F^\ast$ of $\h_c^i(X,\sF)$, there exists an integer 
$m\leqslant n+i$ such that the complex conjugates of $\alpha$ are all of 
absolute value $q^{m/2}$.
\end{theorem_}

The sheaves considered in \ref{VI:1-9} have weight $0$. The theorem 
therefore gives, for the sum \ref{VI:eq:1-9-1} (or rather, for all its 
complex conjugates), the bound 
\begin{equation*}\tag{1.17.1}\label{VI:eq:1-17-1}
  \left|\sum_{s\in S_0(\dF_q)} \prod_i \chi_i(f^i(s)) \right| \leqslant \sum_i \dim \h_c^i\left(S,\sF\left(\prod \chi_i f^i\right)\right)\cdot q^{i/2} \text{.}
\end{equation*}





\subsection{Remarks}\label{VI:1-18}

\begin{enumerate}[a)]
  \item If $\dim S=n$, for any $E_\lambda$-sheaf $\sF$ on $S$, we have 
    \[
      \h_c^i(S,\sF) = 0 \qquad \text{for $i\notin [0,2n]$.}
    \]
  \item The group $\h_c^0(X,\sF)$ is the group of global sections of $\sF$ 
    on $S$ whose support is proper. If $S$ is connected, not proper and 
    $\sF$ is smooth, this group is zero; the same holds if $S$ is connected, 
    $\sF$ is smooth of rank one, and is nonconstant.
  \item If $S$ is smooth and purely of dimension $n$, and $\sF$ is smooth, 
    Poincaré duality says that the vector spaces 
    $\h_c^i(X,\sF)$ and $\h^{2n-i}(X,\sF^\vee)(n)$ are dual to each other 
    (for the effect of a Tate twist on the eigenvalues of Frobenius, 
    cf. the $2$nd part of the proof of \ref{VI:1-13}).
  \item If $\dim S=n$, Poincaré duality allows us to compute 
    $\h_c^{2n}(S,\sF)$ as follows: take an open subset $U$ of $S_\text{red}$ 
    whose complement has dimension $<n$; require this open subset to be smooth and 
    purely of dimension $n$, with $\sF$ smooth on it. We then have 
    $\h_c^{2n}(U,\sF) \iso \h_c^{2n}(S,\sF)$, because in the long exact sequence 
    of cohomology, $\h_c^i(S\setminus U,\sF) = 0$ for $i=2n,2n-1$. 
    Poincaré duality therefore gives 
    \[
      \h_c^{2n}(S,\sF) \iso \left(\h^0(U,\sF^\vee)(n)\right)^\vee \text{.}
    \]
    
    Let $U$ be connected, and let $u$ be a geometric point of $U$. The 
    ``local system'' $\sF$ corresponds to a representation of 
    $\pi_1(U,u)$ on $\sF_u$, and $\h^0(U,\sF^\vee)$ is 
    $(\sF_u^\vee)^{\pi_1(U,u)}=((\sF_u)_{\pi_1(U,u)})^\vee$ (dual of 
    coinvariants). We therefore have 
    \[
      \h_c^{2n}(S,\sF) \simeq (\sF_u)_{\pi_1(U,u)}(-n) \text{.}
    \]
  \item These remarks often allow the calculation of $b_0$ and $b_{2n}$. 
    Computing the other $b_i$ can be difficult. If $b_{2n}=0$ and 
    $b_{2n-1}=0$, the bound \eqref{VI:eq:1-17-1} is a Lang--Weil bound, 
    with, for large $q$, a gain of $\sqrt q$ over the 
    trivial $O(q^n)$ bound. One should bear in mind, however, that 
    \eqref{VI:eq:1-17-1} does not formally imply the trivial bound 
    \[
      \left|\sum_{s\in S_0(\dF_q)} \prod_i \chi_i(f^i(s))\right| \leqslant \left| S_0(\dF_q)\right| \text{,}
    \]
    so it may be useful to take this into account (cf. the proof of \ref{VI:3-8}). 
\end{enumerate}

Below I explain a method for proving that $b_i=0$ for $i\ne n$. 
When it applies, it gives an estimate of order $O(q^{n/2})$. The constant 
implicit in $O$ is $b_n=(-1)^n \chi(S,\sF(\prod \chi_i f^i))$, which may 
be difficult to compute. 





\begin{proposition_}\label{VI:1-19}
Let $\bar X$ be a proper scheme over an algebraically closed field $k$, 
let $j:X\hookrightarrow \bar X$ be an open immersion with domain $X$, and let $\sF$ be an $E_\lambda$-sheaf 
on $X$. Assume that 
\begin{enumerate}[\indent a)]
  \item $(j_\ast\sF)_x = 0$ for $x\in \bar X\setminus X$, i.e. 
    $j_!\sF\iso j_\ast\sF$; 
  \item $\eR^i j_\ast \sF = 0$ for $i>0$.
\end{enumerate}
Then the maps 
\[
  \h_c^i(X,\sF) \to \h^i(X,\sF)
\]
are isomorphisms. 
\end{proposition_}

Assumptions a) and b) mean that $j_! \sF \iso \eR j_\ast \sF$, hence 
\[
  \h_c^i(X,\sF) = \h^i(\bar X,j_!\sF) \iso \hh^i(\bar X,\eR j_\ast \sF) = \h^i(X,\sF) \text{.}
\]
In other words, the Leray spectral sequence for $j$ 
\[
  E_2^{pq} = \h^p(\bar X,\eR^q j_\ast \sF) \Rightarrow \h^{p+q}(X,\sF) 
\]
degenerates, by b) into isomorphisms 
\[
  \h^i(\bar X,j_\ast\sF) \iso \h^i(X,\sF) \text{,}
\]
while, by a) 
\[
  \h_c^i(X,\sF) = \h^i(\bar X,j_!\sF) \iso \h^i(\bar X,j_\ast\sF) \text{.}
\]





\subsubsection{Example}\label{VI:1-19-1}

If $X$ is a curve, or if $\bar X$ is smooth, $X$ is the complement of a 
normal-crossings divisor, and $\sF$ is tamely ramified, 
hypothesis a) of \ref{VI:1-19} (no invariants under local monodromy) 
implies hypothesis b). 





\begin{proposition_}\label{VI:1-20}
Let $X_0$ be a smooth affine scheme, purely of dimension $n$, over $\dF_q$, and let 
$\sF_0$ be a smooth $E_\lambda$-sheaf, pointwise pure of weight $m$, on $X_0$. We 
assume that the maps $\h_c^i(X,\sF) \to \h^i(X,\sF)$ are 
isomorphisms (cf. \ref{VI:1-17}). Then, $\h_c^i(X,\sF) = 0$ for $i\ne n$, and 
the complex conjugates of the eigenvalues of $F^\ast$ on $\h_c^n(X,\sF)$ 
are of absolute value $q^{(m+n)/2}$. 
\end{proposition_}

Poincaré duality pairs $\h_c^i(X,\sF)$ (resp. 
$\h_c^i(X,\sF^\vee)$) with $\h^{2n-i}(X,\sF^\vee(n))$ (resp. 
$\h^{2n-i}(X,\sF(n))$). According to \cite[XIV 3.2]{sga4}, the ordinary $\h^i$ groups 
vanish for $i>n$. By duality, the $\h_c^i$ are zero for $i<n$; since 
$\h_c^i(X,\sF) \iso \h^i(X,\sF)$, these groups are zero for $i\ne n$. 
Apply \ref{VI:1-17} to $\sF_0$ and to $\sF_0^\vee(n)$ (of weight 
$-m-2n$). We find that the complex conjugates $\alpha$ of the eigenvalues 
of $F^\ast$ on $\h_c^n(X,\sF) = \h_c^n(X,\sF^\vee(n))^\vee$ satisfy 
\begin{align*}
  |\alpha| &\leqslant q^{(m+n)/2} && \text{and} \\
  |\alpha^{-1}| &\leqslant q^{((-m-2n)+n)/2} = q^{-(m+n)/2} \text{,}
\end{align*}
hence the assertion. 





\subsection{Remark}\label{VI:1-21}

The last argument shows that if $X_0$ is a smooth separated scheme over 
$\dF_q$, if $\sF_0$ is a smooth $E_\lambda$-sheaf pointwise pure of weight 
$m$ on $X_0$, and if $\h_c^i(X,\sF) \hookrightarrow \h^i(X,\sF)$, then the 
complex conjugates of the eigenvalues of $F^\ast$ on $\h_c^i(X,\sF)$ are 
of absolute value $q^{(m+i)/2}$. 





\subsection*{Appendix. The Lang isogeny in the noncommutative case}




\subsection{}\label{VI:1-22}

Let $G_0$ be a connected algebraic group over $\dF_q$ and let $\fL$ be the morphism 
$x\mapsto F x\cdot x^{-1}$ from $G_0$ to itself. It is a morphism of 
$G_0$-homogeneous spaces if, on the source, $G_0$ acts by left translation 
$x\mapsto (y\mapsto x y)$, and on the target by 
$x\mapsto (y\mapsto F x\cdot y\cdot x^{-1})$). We have $\fL(xy) = \fL(x)$ for 
$y\in G_0(\dF_q)$, and $\fL$ induces an isomorphism $G_0/G_0(\dF_q)\iso G_0$. 

As in the commutative case, $\fL$ therefore makes $G_0$ a $G_0(\dF_q)$-torsor 
on $G_0$, the \emph{Lang torsor} $L$. If 
$\rho:G_0(\dF_q) \to \gl(n,E_\lambda)$ is a linear representation of 
$G_0(\dF_q)$, we denote by $\sF(\rho)$ the $E_\lambda$-sheaf $\rho(L)$ 
(\ref{VI:1-2}). 





\subsection{}\label{VI:1-23}

Let $\gamma\in G^F=G_0(\dF_q)$, and compute the endomorphism $F^\ast$ of 
$\sF(\rho)_\gamma$. Let $g\in G$ be such that $\fL(g)=\gamma$, and choose the frame 
$\rho(g)\in \operatorname{Isom}(E_\lambda^n,\sF(\rho)_\gamma)$. For 
$e\in E_\lambda^n$, we have 
\[
  F(\rho(g)(e)) = \rho(F g)(e) = \rho(\gamma g)(e) = \rho(g\cdot g^{-1}\gamma g)(e) \text{.}
\]
We will see that $g^{-1} \gamma g\in G_0(\dF_q)$, hence 
\[
  F(\rho(g)(e)) = \rho(g) \rho(g^{-1}\gamma g)(e)
\]
and 
\begin{equation*}\tag{1.23.1}\label{VI:eq:1-23-1}
  \rho(g)^{-1} (F_\gamma^\ast)^{-1} \rho(g) = \rho(g^{-1} \gamma g) \text{:}
\end{equation*}
$\rho(g)$ identifies the inverse of $F^\ast$ at $\gamma$ with the automorphism 
$\rho(g^{-1} \gamma g)$ of $E_\lambda^n$. It remains to understand 
$g^{-1} \gamma g$.

\begin{lemma_}\label{VI:1-24}
\begin{enumerate}[(i)]
  \item If $F g\cdot g^{-1} \in G_0(\dF_q)$, we have 
    $g^{-1} (F g\cdot g^{-1}) g = g^{-1} F g\in G_0(\dF_q)$.
  \item Let $\gamma\in G_0(\dF_q)$, and $g$ such that 
    $F g\cdot g^{-1} = \gamma$. The conjugation class of $\gamma'=g^{-1} Fg$ 
    in $G_0(\dF_q)$ only depends on that of $\gamma$. 
  \item The map induced by $\gamma\mapsto \gamma'$ on the set 
    $G_0(\dF_q)^\natural$ of conjugacy classes of $G_0(\dF_q)$ 
    is bijective. Its inverse is the corresponding map for the opposite 
    group. 
\end{enumerate}
\end{lemma_}

(i) If $\gamma= F g\cdot g^{-1}\in G_0(\dF_q)$, the identities $F g=\gamma g$ 
and 
$F\gamma=\gamma$ give $F(g^{-1} F g) = (F g)^{-1} F F g = g^{-1} \gamma^{-1} F(\gamma g) = g^{-1} \gamma^{-1} F \gamma F g = g^{-1} F g$, 
so that $g^{-1} F g \in G_0(\dF_q)$. 

(ii) For $\alpha,\gamma\in G_0(\dF_q)$, the $g$ such that 
$F g\cdot g^{-1} = \gamma$ form a right coset under $G_0(\dF_q)$, and 
if $F g\cdot g^{-1} = \gamma$, we have 
$F(\alpha g \alpha^{-1})(\alpha g \alpha^{-1})^{-1} = \alpha \gamma\alpha^{-1}$. 
As $\gamma$ ranges over a conjugacy class in $G_0(\dF_q)$, the $g$ 
such that $F g\cdot g^{-1} = \gamma$ therefore range over a double coset under 
$G_0(\dF_q)$, and $\gamma'= g^{-1} F g$ ranges over a conjugacy class. 

(iii) Finally, the inverse of $F g\cdot g^{-1}\mapsto g^{-1} F g$ is 
$g^{-1} F g\mapsto F g\cdot g^{-1}$. 





\subsection{}\label{VI:1-25}

Formula \eqref{VI:eq:1-23-1} can be read succinctly: $F^\ast$ at 
$\gamma$ is $\rho(\gamma')^{-1}$. If $\gamma$ belongs to the identity component 
of its centralizer, we can take $g$ in $Z(\gamma)$. We then have 
$\gamma=\gamma'$. If this condition is not met, $\gamma$ and $\gamma'$ 
need not be conjugate. 





\subsection{An example}\label{VI:1-26}

For $G_0=\operatorname{SL}(2)$, $q$ odd, 
$\alpha\in \dF_q^\times\setminus {\dF_q^\times}^2$, and $\bar\gamma$ is the 
conjugacy class of $-1(1+N)$, with $N^2=0$, then $\bar\gamma'$ is the 
conjugacy class of $-1\cdot (1+\alpha N)$, so that $\bar\gamma'\ne \bar\gamma$ 
if $N\ne 0$. 










\section{Dictionary}\label{VI:2}

Elementary manipulations of trigonometric sums can often 
be seen, via the trace formula, as counterparts of statements 
in cohomology. In this section, I list such statements and their 
counterparts; they carry the same number, with a $\ast$ appended to 
the cohomological statement. 





\subsection{}\label{VI:2-1}

``A sum of algebraic integers is an algebraic integer'' has as its 
cohomological analogue the following theorem, proved in 
\cite[XXI 5.2.2]{sga7} (see the use of \ref{VI:2-1}$^\ast$ in 
\ref{VI:1-13}). 





\begin{theorem*}[2.1*]\label{VI:2-1*}
Let $X_0$ be a separated scheme of finite type over $\dF_q$ and let $\sF_0$ be an 
$E_\lambda$-sheaf on $X_0$. If, for all $x\in |X_0|$, the 
eigenvalues of $F_x^\ast$ on $\sF_0$ are algebraic integers, then the 
eigenvalues of $F^\ast$ on $\h_c^i(X,\sF)$ are algebraic integers. 
\end{theorem*}





\subsection{}\label{VI:2-2}

Other integrality results are proved in 
\cite[XXI, 5.2.2 et 5.4]{sga7}: if $\dim(X_0)\leqslant n$, and $\alpha$ is 
an eigenvalue of $F^\ast$ on $\h_c^i(X,\sF)$, we have: 
\begin{enumerate}[a)]
  \item under the assumptions of \ref{VI:2-1}$^\ast$, and if $i>n$, then 
    $q^{i-n}$ divides $\alpha$;
  \item if the inverses of the eigenvalues of $F_x^\ast$ are algebraic 
    integers, then $\alpha^{-1}$ is integral away from $p$. More 
    precisely, $q^{\inf(n,i)}\alpha^{-1}$ is an algebraic integer. 
\end{enumerate}

For the sheaves considered in \ref{VI:1-9}, the eigenvalues of 
$F_x^\ast$ are roots of unity, so the assumptions of 
the theorem and of b) above are satisfied. 





\subsection{}\label{VI:2-3}

Let $f:A\to B$ be a map from one finite set to another, and 
$\varepsilon$ a function on $A$. We have 
\begin{equation*}\tag{2.3.1}\label{VI:eq:2-3-1}
  \sum_{a\in A} \varepsilon(a) = \sum_{b\in B} \sum_{f(a)=b} \varepsilon(a) \text{.}
\end{equation*}





\subsection*{2.3*}\label{VI:2-3*}

Let $f:X\to Y$ be a morphism of separated schemes of finite type over an 
algebraically closed field $k$, and let $\sF$ be an $E_\lambda$-sheaf on $X$. The 
Leray spectral sequence for $f$ in cohomology with compact support is 
\begin{align*}\tag{2.3.1*}\label{VI:eq:2-3-1*}
  E_2^{p q} = \h_c^p(Y,\eR^q f_! \sF) \Rightarrow \h_c^{p+q}(X,\sF) \text{.}
\end{align*}

Let $f_0:X_0 \to Y_0$ be a morphism of separated schemes of finite type over 
$\dF_q$ and let $\sF_0$ be an $E_\lambda$-sheaf on $X_0$. Between 
\eqref{VI:eq:2-3-1} and \eqref{VI:eq:2-3-1*} there is the following compatibility. 

\begin{enumerate}[a)]
  \item The sheaf $\eR^q f_! \sF$ is obtained by extending scalars from 
    $\dF_q$ to $\dF$ in the sheaf $\eR^q f_{0!} \sF_0$ on $Y_0$. As 
    such, it is equipped with a Frobenius correspondence 
    $F^\ast:F^\ast\eR^q f_! \sF \to \eR^q f_! \sF$. It is obtained by applying 
    the functoriality of $\eR f_!$ to the Frobenius correspondence of $\sF$: 
    namely the composite 
    \[\xymatrix{
      F^\ast \eR^q f_! \sF \ar[r]^-\sim 
        & \eR^q f_! F^\ast \sF \ar[r]^-\sim 
        & \eR^q f_! \sF \text{.}
    }\]
  \item The trace formula for $\sF_0$ reads 
    \begin{align*}\tag{1}\label{VI:eq:1}
      \sum_{x\in X_0(\dF_q)} \tr(F_x^\ast,\sF_0) = \tr(F^\ast,\h_c^\bullet(X,\sF)) \text{;}
    \end{align*}
    while that for $\eR^q f_! \sF$ reads 
    \begin{align*}\tag{2$_q$}\label{VI:eq:2_q} 
      \sum_{y\in Y_0(\dF_q)} \tr(F_y^\ast,\eR^q f_! \sF) = \tr(F^\ast, \h_c^\bullet(Y,\eR^q f_! \sF)) \text{.}
    \end{align*}
\end{enumerate}

We have, for $y\in Y(\dF)$, $(\eR^q f_! \sF)_y = \h_c^q(f^{-1}(y),\sF)$, and if 
$y\in Y_0(\dF_q)$, 
\[
  \tr(F_y^\ast,\eR^q f_! \sF) = \tr(F^\ast,\h_c^q(f^{-1}(y),\sF)) \text{,}
\]
hence by the trace formula on the fiber 
\[
  \sum_{\substack{x\in X_0(\dF_q) \\ f(x) = y}} \tr(F_x^\ast,\sF_0) = \sum (-1)^q \tr(F_y^\ast,\eR^q f_! \sF) \text{.}
\]
Taking the alternating sum of the identities \eqref{VI:eq:2_q}, we therefore find 
\begin{align*}\tag{2}\label{VI:eq:2}
  \sum_{y\in Y_0(\dF_q)} \sum_{\substack{x \in X_0(\dF_q) \\ f(x) = y}} \tr(F_x^\ast,\sF_0) = \tr(F^\ast,\h_c^\bullet(Y,\eR^\bullet f_! \sF)) \text{.}
\end{align*}

Equality of the left-hand sides of \eqref{VI:eq:1} and \eqref{VI:eq:2} 
follows from \eqref{VI:eq:2-3-1}; equality of the right-hand sides follows from 
\eqref{VI:eq:2-3-1*} (taking into account that $F^\ast$ is an endomorphism of 
the entire spectral sequence). 





\subsection{}\label{VI:2-4}

Let $(A_i)_{i\in I}$ be a finite family of finite sets, 
$A=\prod_{i\in I} A_i$, $\varepsilon_i$ a function on $A_i$, and 
$\varepsilon$ the function $a\mapsto \prod \varepsilon_i(a_i)$ on $A$. We have 
\begin{align*}\tag{2.4.1}\label{VI:eq:2-4-1}
  \sum_{a\in A} \varepsilon(a) = \prod_{i\in I} \sum_{a_i\in A_i} \varepsilon_i(a_i) \text{.}
\end{align*}





\subsection*{2.4*}\label{VI:2-4*}

Let $(X_i)_{i\in I}$ be a finite family of separated schemes of finite type 
over an algebraically closed field $k$, let $X=\prod_{i\in I} X_i$, let $\sF_i$ be an 
$E_\lambda$-sheaf on $X_i$ and $\sF$ the external tensor product 
$\bigboxtimes_{i\in I}\sF_i = \bigotimes_{i\in I} \pr_i^\ast(\sF_i)$ 
of the $\sF_i$. 
The K\"unneth formula gives 
\begin{align*}\tag{2.4.1*}\label{VI:eq:2-4-1*}
  \h_c^\bullet(X,\sF) = \bigotimes_{i\in I} \h_c^\bullet(X_i,\sF_i) \text{.}
\end{align*}

For the isomorphism \eqref{VI:eq:2-4-1*} to be canonical and independent of 
the choice of a total ordering of $I$, the tensor product on the right must be 
understood as a graded tensor product, with the Koszul sign rule 
(\hyperref[IV]{Cycle}, \ref{IV:1-3}). 





\subsection{}\label{VI:2-5}

Let $A$ be a finite set, let $B$ be a subset of $A$, and let $\varepsilon$ be a 
function on $A$. We have 
\begin{align*}\tag{2.5.1}\label{VI:2-5-1}
  \sum_{a\in A} \varepsilon(a) = \sum_{a\in A\setminus B}\varepsilon(a) + \sum_{a\in B} \varepsilon(a) \text{.}
\end{align*}





\subsection*{2.5*}\label{VI:2-5*}

Let $X$ be a separated scheme of finite type over an algebraically closed field $k$, 
let $Y$ be a closed subscheme, and let $\sF$ be an $E_\lambda$-sheaf on $X$. We have 
a long exact cohomology sequence 
\begin{align*}\tag{2.5.1*}\label{VI:eq:2-5-1*}
\xymatrix{
  \cdots \ar[r]^-\partial 
    & \h_c^i(X\setminus Y,\sF) \ar[r] 
    & \h_c^i(X,\sF) \ar[r] 
    & \h_c^i(Y,\sF) \ar[r]^-\partial 
    & \cdots 
}
\end{align*}

More generally, given a finite filtration of $X$ by closed subsets 
$X_p$ ($p\in \dZ$; assume that $X_p\supset X_{p+1}$, that $X_p=X$ 
for $p$ sufficiently small, and that $X_p=\varnothing$ for $p$ sufficiently large), there is a 
spectral sequence 
\begin{align*}\tag{2.5.2*}\label{VI:eq:2-5-2*}
  E_1^{p q} = \h_c^{p+q}(X_p\setminus X_{p+1},\sF) \Rightarrow \h_c^{p+q}(X,\sF) \text{.}
\end{align*}





\subsection{}\label{VI:2-6}

Let $A$ be a finite set, $(A_i)_{i\in I}$ a finite covering of $A$ and 
$\varepsilon$ a function on $A$. For $J\subset I$, denote by $A_J$ 
the intersection of $A_j$ ($j\in J$). We have 
\begin{align*}\tag{2.6.1}\label{VI:eq:2-6-1}
  \sum_{J\subset I} (-1)^{|J|} \sum_{a\in A_J} \varepsilon(a) = 0 \text{.}
\end{align*}





\subsection*{2.6*}\label{VI:2-6*}

Let $X$ be a separated scheme of finite type over an algebraically closed field $k$, 
$(X_i)_{i\in I}$ a closed (resp. open) finite covering of $X$ by 
subschemes, and $\sF$ an $E_\lambda$-sheaf on $X$. For $J\subset I$, 
let $X_J$ be the intersection of the $X_j$ ($j\in J$). We have the respective 
Leray spectral sequences 
\begin{align*}\tag{2.6.1*}\label{VI:eq:2-6-1*}
  E_1^{p q} &= \bigoplus_{|J|=p+1>0} \h_c^q(X_J,\sF) \otimes_\dZ \textstyle\bigwedge^{|J|} \dZ^J \Rightarrow \h_c^{p+q}(X,\sF) \\ 
  \tag{2.6.2*}\label{VI:eq:2-6-2*}
  E_1^{p q} &= \bigoplus_{|J|=1-p>0} \h_c^q(X_J,\sF) \otimes_\dZ \textstyle\bigwedge^{|J|} \dZ^J \Rightarrow \h_c^{p+q}(X,\sF)
\end{align*}
If $J=\varnothing$ were not excluded, there would also be spectral 
sequences converging to $0$. The factors $\bigwedge^{|J|} \dZ^J$ are included 
to avoid choosing a total order on $I$. 





\subsection{}\label{VI:2-7}

Let $A$ be a finite commutative group, and $\chi:A \to E_\lambda^\times$ a 
nontrivial character. We have 
\begin{align*}\tag{2.7.1}\label{VI:eq:2-7-1}
  \sum_{a\in A} \chi(a) = 0 \text{.}
\end{align*}
We will prove the cohomological analogue of \eqref{VI:eq:2-7-1} by transposing 
the following proof: if $x\in A$, $a\mapsto x a$ is a permutation of $A$, and 
\begin{align*}
  \sum_{a\in A}\chi(a) = \sum_{a\in A} \chi(x a) &= \chi(x) \sum_{a\in A} \chi(a) && \text{, whence} \\
  (\chi(x)-1)\sum_{a\in A}\chi(a) &= 0 \text{,}
\end{align*}
and by hypothesis there exists $x$ such that $\chi(x)-1\ne 0$. 





\begin{theorem*}[2.7*]\label{VI:2-7*}
Let $G_0$ be a connected commutative algebraic group over $\dF_q$ and 
$\chi:G_0(\dF_q) \to E_\lambda^\times$ a nontrivial character. We have 
\begin{align*}\tag{2.7.1*}\label{VI:eq:2-7-1*}
  \h_c^\bullet(G,\sF(\chi)) = 0 \text{.}
\end{align*}
\end{theorem*}

For $x$ a rational point of $G$, let $t_x$ be the translation $t_x(g) = x g$ 
of $G$. The formula $\fL t_x = t_{\fL(x)} \fL$ expresses that $(t_x,t_{\fL(x)})$ 
is an automorphism of the diagram $G\xrightarrow\fL G$ ($G$ equipped with the torsor of 
Lang). Let $\rho(g)$ be the automorphism of $(G,\sF(\chi))$ induced by it. 
For $g\in G_0(\dF_q)$, i.e. for $\fL(g) = e$, it is the identity on $G$, and 
multiplication by $\chi(g)^{-1}$ on $\sF(\chi)$. 

Let $\rho_H(g)$ denote the automorphism of $\h_c^\bullet(G,\sF(\chi))$ induced by 
$\rho(g)$. A homotopy argument (Lemma \ref{VI:2-8} below) shows that 
$\rho_H(g) = \rho_H(e)$, and hence is the identity. On the other hand, for 
$g\in G_0(\dF_q)$, $\rho_H(g)$ is multiplication by 
$\chi^{-1}(g)$: on $\h^\bullet(G,\sF(\chi))$, multiplication by 
$(\chi^{-1}(g)-1)$ is zero. Taking $g$ such that $\chi(g)\ne 1$, we obtain 
\ref{VI:2-7*}. 





\begin{lemma_}\label{VI:2-8}
Let $X$ and $Y$ be two schemes over an algebraically closed field $k$, with $X$ 
separated and of finite type and $Y$ connected. Let $\sF$ be a sheaf on $X$ and 
let $(\rho,\varepsilon)$ be a family of endomorphisms of $(X,\sF)$ parametrized 
by $Y$: 
\begin{align*} 
  \rho:Y\times_k X &\to Y\times_k X && \text{is a $Y$-morphism, and} \\
  \varepsilon:\rho^\ast \pr_2^\ast \sF &\to \pr_2^\ast\sF && \text{a morphism of sheaves.} 
\end{align*}
Assume that $\rho$ is proper. For $y\in Y(k)$, let $\rho_H(y)^\ast$ denote 
the endomorphism of $\h_c^\bullet(X,\sF)$ induced by $\rho_y:X\to X$ and 
$\varepsilon_y:\rho_y^\ast\sF\to \sF$. Then, $\rho_H(y)^\ast$ is 
independent of $y$. 
\end{lemma_}

Indeed, $\eR^p \pr_{1!} \pr_2^\ast\sF$ is the 
constant sheaf on $Y$ with fiber $\h_c^p(X,\sF)$, and $\rho_H(y)^\ast$ is the 
fiber at $y$ of the endomorphism 
\[\xymatrix{
  \eR^p \pr_{1!} \pr_2^\ast \sF \ar[r]^-{\rho^\ast} 
    & \eR^p \pr_{1!} \rho^\ast \pr_2^\ast \sF \ar[r]^-\varepsilon 
    & \eR^p \pr_{1!} \pr_2^\ast \sF 
}\]
of this sheaf. 





\subsection{Remark}\label{VI:2-9}

Let $G_0$ be a connected algebraic group over $\dF_q$ and 
$\rho:G_0(\dF_q) \to \operatorname{GL}(V)$ a linear representation 
such that $V^{G_0(\dF_q)}=0$. Then again $\h_c^\bullet(G,\rho(L))=0$. 










\section{Sums in one variable}\label{VI:3}





\subsection{}\label{VI:3-1}

Weil was the first to apply ``cohomological'' methods to 
the study of trigonometric sums in one variable; since he had 
proved the analogue of the Riemann hypothesis for Artin's $L$ functions 
over function fields, these methods gave him excellent 
estimates (of order $O(\text{square root of the number of terms})$) and an 
understanding of how these sums behave under ``extension of the base field.'' 

The substance of this section is a presentation, in cohomological language, of 
his results. 





\subsection{}\label{VI:3-2}

Cohomological methods here lead us to compute the groups 
$\h_c^i(X,\sF)$, where $\sF_0$ is an $E_\lambda$-sheaf on a curve $X_0$ over 
$\dF_q$. These groups vanish for $i\ne 0,1,2$, and for 
$i=0,2$, they have a simple interpretation (\ref{VI:1-8} a,b,c). Furthermore, the 
Euler-Poincaré characteristic 
$\chi_c(\sF) = \sum (-1)^i \dim \h_c^i(X,\sF)$ (which is also equal to 
$\chi(\sF) = \sum (-1)^i \dim \h^i(X,\sF)$) can be computed in terms of 
$X$, the rank of $\sF$ at the generic points of $X$, and the 
ramification of $\sF$. The essential result is the following. 

Let $\bar X$ be a smooth connected projective curve of genus $g$ over an 
algebraically closed field $k$, let $X$ be the dense open subset of $\bar X$ 
complementary to a finite set $S$ of points, and let $\sF$ be a 
smooth $E_\lambda$-sheaf on $X$. Let $\chi(X)=2-2 g-\# S$ be the 
Euler--Poincaré characteristic of $X$, and let $\operatorname{rg}(\sF)$ be the 
rank of $\sF$. For each point $s\in S$, we define an integer 
$\swan_s(\sF)$, the Swan conductor, measuring the wild ramification of 
$\sF$, and 
\begin{align*}\tag{3.2.1}\label{VI:eq:3-2-1}
  \chi_c(\sF) = \operatorname{rg}(\sF) \cdot \chi(X) - \sum_{s\in S} \swan_s(\sF) 
\end{align*}
\cite{ra65}. 





\subsection{}\label{VI:3-3}

Let $\bar X$, $j:X\hookrightarrow\bar X$ and $\sF$ as in \ref{VI:3-2}. The 
Poincaré duality theorem takes the particularly convenient form: 
$\h^i(\bar X,j_\ast \sF)$ and $\h^{2-i}(\bar X,j_\ast(\sF^\vee(1)))$ are 
dual to each other (\hyperref[V]{Duality}, \ref{V:1-3}). 





\subsection{}\label{VI:3-4}

Take $k=\dF$, and suppose that $X$, $\bar X$, and $\sF$ arise from 
$X_0$, $\bar X_0$, and $\sF_0$ over $\dF_q$. If $\sF_0$ becomes trivial on a 
finite covering of $X_0$, Weil proved that the complex conjugates of 
the eigenvalues of $F^\ast$ on $\h^i(\bar X,j_\ast\sF)$ have absolute value 
$q^{i/2}$. Such sheaves $\sF_0$ correspond to Artin $L$-functions 
over the function field of $\bar X_0$. 





\subsection*{3.5}\label{VI:3-5_} % originally one of two 3.5s

Under the same hypotheses, or more generally if $E_\lambda$ is the 
completion at a place $\lambda$ of a number field $E$ and if $\sF_0$ 
belongs to an infinite compatible system of $v$-adic representations 
($v$ a place of $E$), we can compute 
\[
  \prod_i \det(-F^\ast,\h^i(X,\sF))^{(-1)^i} 
\]
from \emph{local} information on $\sF_0$. This is explained in 
\cite{de73}. For trivial $\sF_0$ on a finite covering of $X_0$, it is 
the expression of the constant in the functional equation of an Artin 
$L$-function as a product of local constants. 





\subsection{Example}\label{VI:3-5}

Let $\psi$ be the character 
$\exp\left(\frac{2\pi i}{p} \tr_{\dF_q/\dF_p}\right)$ of $\dF_q$; we have 
$\psi(a^p-a)=1$. Let $X_0$ be a smooth projective geometrically 
irreducible curve of genus $g$ over $\dF_q$, and let $f$ be a rational function on 
$X_0$, i.e. a morphism $f:X_0 \to \dP^1$, not identically equal to 
$\infty$. We are interested in the sum 
\[
  S_f' = \sum_{\substack{x\in X_0(\dF_q) \\ f(x)\ne \infty}} \psi(f(x)) 
  \text{.}
\]
This sum is zero for $f$ of the form $g^p-g$. This suggests modifying 
the sum $S_f'$ as follows: 
\begin{enumerate}[\indent a)]
  \item for any closed point $x$ of $X_0$, set $v_x(f) = $ the order of the 
    pole of $f$ at $x$ if $f(x)=\infty$, and $v_x(f)=0$ otherwise; also set 
    $v_x^\ast(f) = \inf v_x(f+g^p-g)$ (the infimum being over $g$). 
  \item If $v_x^\ast(f)=0$, $x\in X_0(\dF_q)$, and $f+g^p-g$ is 
    regular at $x$, put $\psi(f(x)) = \psi((f+g^p-g)(x))$. Finally set 
    \begin{equation*}\tag{3.5.1}\label{VI:eq:3-5-1}
      S_f = \sum'_{x\in X_0(\dF_q)} \psi(f(x)) \text{,}
    \end{equation*}
    where $(-)'$ indicates that the sum runs over those $x$ for which 
    $v_x^\ast(f)=0$. 
\end{enumerate}

We have $S_f=S_{f+g^p-g}$. We shall verify that if $f$ is not of the 
form $g^p-g+c^{t e}$, then 
\begin{equation*}\tag{3.5.2}\label{VI:eq:3-5-2}
  |S_f|\leqslant \left(2 g-2+\sum_{v_x^\ast(f)\ne 0} [k(x):\dF_q] (1+v_x^\ast(f))\right) q^{1/2} \text{.}
\end{equation*}

We begin by passing from complex coefficients to $\ell$-adic coefficients, replacing $\psi$ by 
a character of the form $\psi_0\circ \tr_{\dF_q/\dF_p}$, with 
$\psi_0:\dF_p\hookrightarrow E_\lambda^\times$. For $j:U_0\hookrightarrow X_0$ 
the inclusion of the open where $f\ne\infty$, we then prove that 
\begin{equation*}\tag{3.5.3}\label{VI:eq:3-5-3}
  S_f = \sum (-1)^i \tr(F^\ast,\h^i(X,j_\ast \sF(\psi f))) \text{.} 
\end{equation*}

The trace formula reduces this statement to the following (for 
$x\in X_0(\dF_q)$): 
\begin{enumerate}[\indent a)]
  \item If $v_x^\ast(f)=0$, then $F_x^\ast$ on $j_\ast \sF(\psi f)$ is 
    $\psi(f(x))$; 
  \item If $v_x^\ast(f)\ne 0$, then $j_\ast \sF(\psi f)$ is ramified at $x$: 
    $(j_\ast \sF(\psi f))_{\bar x}=0$. 
\end{enumerate}

Formula a) follows from \eqref{VI:eq:1-7-6} if $v_x(f)=0$, and the general 
case reduces to this one by replacing $f$ by $f+g^p-g$: up to isomorphism, this 
does not change $\sF(\psi f)_0$ on the open where $f$ and $g$ are regular 
(\ref{VI:1-3}). 

Formula b) follows from 
\begin{equation*}\tag{3.5.4}\label{VI:eq:3-5-4}
  \swan_x \sF(f,\psi) = v_x^\ast(f) \text{.} 
\end{equation*}
After reduction to the case $v_x(f)=v_x^\ast(f)$ (hence $p\nmid v_x(f)$) 
and extension of scalars to $\dF$, this formula can be found in 
\cite[4.4]{se61}.

Finally, we check that the sheaf $\sF(\psi f)$ on $U$ is nonconstant if 
$f$ is not of the form $g^p-g+c^{t e}$: since $\psi_0$ is injective, the 
triviality of $\sF(\psi f)=\sF(\psi_0 f)$ (\ref{VI:1-8}.ii) is equivalent to 
that of the $\dF_p$-torsor on $U$ defined by $T^p-T-f=0$. If it were trivial, 
the torsor on $U_0$ defined by the same equation would be the inverse image of an 
$\dF_q$-torsor on $\spec(\dF_q)$, with equation $T^p-T-\lambda=0$, and the 
$\dF_p$-torsor equation $T^p-T-(f-\lambda)=0$ on $U_0$ would be trivial 
and $f$ would therefore be of the form $g^p-g+\lambda$. The 
$\h^i(\bar X,j_\ast \sF(\psi f))$ therefore vanish for $i\ne 1$, and formula 
\eqref{VI:eq:3-5-3} reduces to 
\begin{equation*}\tag{3.5.5}\label{VI:eq:3-5-5}
  S_f = -\tr(F^\ast,\h^1(\bar X,j_\ast \sF(\psi f))) 
\end{equation*}
and, by \eqref{VI:eq:3-2-1} and \eqref{VI:eq:3-5-4}, the dimension of $\h^1$ is 
$2 g-2+\sum_{v_x^\ast(f)>0} [k(x):\dF_q](1+v_x^\ast(f))$; moreover, $-S_f$ 
is the sum of this many eigenvalues of $F^\ast$, each having complex 
absolute value $q^{1/2}$. 





\subsection{}\label{VI:3-6}

Suppose that, for an automorphism $\sigma$ of $X_0$, we have 
\[
  f(\sigma x) = - f(x) \text{.} 
\]
The sum $S_f$ is then real. If $\sigma$ is an involution, Poincaré 
duality allows us to say a little more: set 
$\sF_0=j_\ast \sF(\psi f)$, $\sG_0 = j_\ast \sF(\psi(-f))$, we have 
\begin{enumerate}[\indent a)]
  \item $\sF_0$ and $\sG_0$ are in duality, 
  \item $\sigma^\ast \sF_0 \iso \sG_0$ and $\sigma^\ast \sG_0 \iso \sF_0$, for 
    natural isomorphisms such that the composite 
    $\sF_0=(\sigma^2)^\ast \sF_0 \iso \sigma^\ast \sG_0 \iso \sF_0$ 
    is the identity. 
  \item the pairing $\sF_0\otimes \sG_0 \to E_\lambda$ satisfies 
    $\sigma^\ast(f\cdot g) = \sigma^\ast(f) \cdot \sigma^\ast (g)$. 
\end{enumerate} 

Passing to cohomology, we find that $\h^1(X,\sF)$ and $\h^1(X,\sG)$ are in 
perfect duality with values in $E_\lambda(1)$, and that the bilinear form 
$\alpha\cdot \sigma^\ast\beta$ on $\h^1(X,\sF)$ is \emph{alternating}: 
$\sigma^\ast$ acts trivially on $E_\lambda(1)$, and 
\[
  \alpha\cdot \sigma^\ast \beta = \sigma^\ast(\alpha\cdot \sigma^\ast \beta) = \sigma^\ast \alpha\cdot \beta = -\beta\cdot \sigma^\ast \alpha \text{.} 
\]
We conclude that $\h^1(X,\sF)$ is of even dimension (we also check 
    easily that each nonzero $v_x^\ast(f)$ is odd) and that the 
eigenvalues of $F^\ast$ on $\h^1(X,\sF)$ occur in pairs $\alpha$ and 
$q/\alpha$. 





\subsection{Example}\label{VI:3-7}

This applies to the Kloosterman sums 
$\sum_{x\in \dF_q^\times} \psi\left(x+\frac a x\right)$ (take 
$X_0=\dP^1$, $f=x+\frac a x$, $\sigma x=-x$; the poles of $f$ are $0$ and 
$\infty$, and at each of them $v_x^\ast(f)=1$; the dimension of $\h^1$ is 
$-2+2+2=2$). We therefore have (for $a\ne 0$) 
\[
  \sum_{\substack{x y=a \\ x,y\in \dF_{q^n}}} \exp\left(\frac{2\pi i}{p} \tr_{\dF_{q^n}/\dF_p}(x+y)\right) = -(\alpha^n+\bar\alpha^n) \qquad \text{, }\alpha\bar \alpha = q \text{.} 
\]
This result is due to L.\ Carlitz \cite{ca69}. 

We now give an application of a Lang--Weil method. 





\begin{proposition_}\label{VI:3-8}
Let $P\in \dF_q[X_1,\dots,X_n]$ be a polynomial in $n$ variables, of degree 
$d$, which is not of the form $Q^p-Q+C^{te}$. Again setting 
$\psi(x) = \exp\left(\frac{2\pi i}{p} \tr_{\dF_q/\dF_p}(x)\right)$, we have 
\[
  \left|\sum_{x_i\in \dF_q} \psi P(x_1,\dots,x_n)\right| \leqslant (d+1) q^{n-\frac 1 2} \text{.}
\]
\end{proposition_}

For the left-hand side we have the trivial estimate $q^n$. It is therefore enough to 
prove \ref{VI:3-8} when $d-1<q^{1/2}$; assume therefore that 
$d<q+1$, and let $P_d$ be the homogeneous part of degree $d$ of $P$. Recall: 





\begin{lemma_}\label{VI:3-9}
A hypersurface of degree $d$ in $\dP^r$ ($r\geqslant 1$) cannot contain 
all rational points over $\dF_q$ unless $d\geqslant q+1$. 
\end{lemma_}

We argue by induction: if there exists a rational hyperplane not 
entirely contained in the hypersurface (this is not the case for $r=1$), we 
apply the induction hypothesis to the intersection of the hypersurface with this 
hyperplane. Otherwise, the degree $d$ is $\geqslant$ the number of rational 
hyperplanes, $\geqslant q+1$. 

Let us now distinguish two cases. 


\paragraph{Case 1: $p\nmid d$.}
Applying the lemma to $P_d$, we see that, after making a linear change 
of variables if necessary, we may assume that $P_d(1,0,\ldots,0)\ne 0$, i.e. 
that the coefficient of $X_1^d$ in $P$ is nonzero. A polynomial in $1$ 
variable 
\[
  S(X) = \sum_{i=0}^d a_i x^i \text{,} 
\]
with $a_d\ne 0$ ($p\nmid d$) is never of the form $Q^p-Q+C^{t e}$, and 
the \eqref{VI:eq:3-5-2} estimate is reduced to 
\[
  \left| \sum \psi(S(x)) \right| \leqslant (d-1) q^{1/2} \text{.} 
\]
Applying this estimate to the partial sums obtained by varying 
only $x_1$, we find 
\[
  \left| \sum \psi P(x_1,\dots,x_n) \right| = \left| \sum_{x_2,\dots,x_n} \sum_{x_1} \psi P(x_1,\dots,x_n)\right| \leqslant q^{n-1} (d-1) q^{1/2} 
\]
as promised. 


\paragraph{Case 2: $p\mid d$ ($d>0$).}
If $P_d$ is a $p$-th power, replacing $P$ by $P-(P_d-P_d^{1/p})$ 
does not change the sum under consideration and lowers the degree: this case 
is disposed of by induction on $d$. Otherwise, the differential of 
$P_d$ is not identically zero; applying \ref{VI:3-9}, one can assume, 
after a linear change of variables if necessary, that it is nonzero 
at the point $(1,0,\dots,0)$. If we write 
\[
  P_d = \sum_{i=0}^d X_1^{d-i} S_i(X_2,\dots,X_n) \text{,} 
\]
this means that the linear form $S_1$ is not identically zero. 

The form $S_0$ is constant, and by replacing $P$ with 
$P-(S_0 X_1^d -S_0^{1/p} X_1^{d/p})$, we can assume that it is zero. 
Finally, denote by $-\lambda$ the coefficient of $X_1^{d-1}$ in $P$. For fixed 
$x_2,\dots,x_n$, we have 
\[
  P(X_1,x_2,\dots,x_n) = (S_1(x_2,\dots,x_n)-\lambda) X_1^{d-1} + \text{terms of lower degree in $X_1$} 
\]
and if $S_1(x_2,\dots,x_n)\ne \lambda$, we therefore have 
\[
  \left| \sum_{x_1} \psi P(x_1,\dots,x_n)\right| \leqslant (d-2) q^{1/2} \text{.} 
\]
In total, 
\begin{align*}
  \left|\sum \psi P(x_1,\dots,x_n)\right| 
    &\leqslant \left|\sum_{S_1=\lambda} \psi P(x_1,\dots,x_n)\right| + \sum_{S_1(X_2,\dots,X_n)\ne\lambda} \left| \sum_{x_1} \psi P(x_1,\dots,x_n)\right| \\
    &\leqslant q^{n-1} + (q^{n-1}-q^{n-2})(d-2) q^{1/2} \\
    &< (d-2) q^{n-1/2} + q^{n-1} \\
    &< (d-1) q^{n-1/2} \text{.}
\end{align*}

Another result of this nature is given by R.A.\ Smith 
\cite{sm70}. 










\section{Gauss and Jacobi sums}\label{VI:4}





\subsection{}\label{VI:4-1}

Let $k$ be a finite field of characteristic $p$, let $\chi$ be a character of 
$k^\times$, and let $\psi$ be a nontrivial character of the additive group of $k$. 
We take the following as the definition of Gauss sums: 
\begin{equation*}\tag{4.1.1}\label{VI:eq:4-1-1}
  \tau(\chi,\psi) = -\sum_{x\in k^\times} \psi(x) \chi^{-1}(x) 
\end{equation*}
(note the sign). Classically, $\chi$ and $\psi$ are complex-valued. 
We take them to have values in $E_\lambda^\times$ (see the proof of 
\ref{VI:1-15}). Regard $-\tau(\chi,\psi)$ as a sum over the rational 
points of the scheme $\dG_m$ over $k$. By Section \ref{VI:1}, 
we have $-\tau(\chi,\psi) = \tr(F^\ast,\h_c^\ast(\dG_m,\sF(\psi \chi^{-1}))$. Moreover,

\begin{proposition_}\label{VI:4-2}
The cohomology of $\dG_m$ with coefficients in $\sF(\psi\chi^{-1})$ satisfies 
\begin{enumerate}[(i)]
  \item $\h_c^i=0$ for $i\ne 1$, and $\dim \h_c^1=1$. 
  \item $F^\ast$, acting on $\h_c^1$, is multiplication by 
    $\tau(\chi,\psi)$. 
  \item If $\chi$ is nontrivial, we have $\h_c^\bullet \iso \h^\bullet$. 
\end{enumerate}
\end{proposition_}
\begin{proof}[Proof]
For any $n$ prime to $p$, denote by $K_n$ the $\dmu_n$-torsor on $\dG_m$ 
defined by the exact K\"ummer sequence 
$0 \to \dmu_n \to \dG_m \xrightarrow{x^n} \dG_m \to 0$. If $k$ has $q$ 
elements, then over $k$ one has $\dmu_{q-1}=k^\times$, and the Lang torsor on 
$\dG_m/k$ is $K_{q-1}$. Hence $\sF(\chi^{-1})=\chi(K_{q-1})$. 
Assertion \ref{VI:4-2}(ii) follows from (i,iii), which in turn follow from 
the following geometric statement, where $\sF(\psi)$ denotes the sheaf 
on $\dG_m/\dF$ obtained by extending scalars from $k$ to $\dF$ in 
$\sF(\psi)$ on $\dG_m/k$ (\ref{VI:1-8}.b). 
\end{proof}





\begin{proposition_}\label{VI:4-3}
Let $\chi:\dmu_n \to E_\lambda^\times$. The cohomology of $\dG_m$ with 
coefficients in $\sF(\psi)\otimes \chi(K_n)$ satisfies 
\begin{enumerate}[\indent (i)]
  \item $\h_c^i=0$ for $i\ne 1$, and $\dim \h_c^1{} = 1$. 
  \item If $\chi$ is nontrivial, we have $\h_c^\bullet{}\iso \h^\bullet{}$. 
\end{enumerate}
\end{proposition_}
\begin{proof}[Proof]
The sheaf $\sF(\psi)$ is the restriction to $\dG_m$ of a sheaf locally 
constant on $\dG_a$, wildly ramified at infinity, with Swan conductor 
$1$. The sheaf $\chi(K_n)$ is constant if $\chi=1$; if $\chi\ne 1$, it 
is tamely ramified at $0$ and $\infty$. 

The sheaf $\sF(\psi)\otimes\chi(K_n)$ is therefore ramified at $\infty$; 
applying \ref{VI:1-18}(b,c), we find that 
$\h_c^i(\dG_m,\sF(\psi)\otimes \chi(K_n))=0$ for $i\ne 1$. If $\chi\ne 1$, it 
is ramified at $0$ and $\infty$, and (ii) follows from \ref{VI:1-19}.a. 

The Swan conductors are $0$ at $0$ and $1$ at $\infty$. By 
\eqref{VI:eq:3-2-1}, the Euler--Poincaré characteristic is therefore $-1$, 
which completes the proof. 
\end{proof}





\subsection{Remark}\label{VI:4-4}

If $\chi$ is nontrivial, \ref{VI:4-2}(iii) and Poincaré duality 
show that $\h_c^1(\dG_m,\sF(\psi\chi^{-1}))$ and 
$\h_c^1(\dG_m,\sF(\psi^{-1}\chi))$ are in duality (duality with values 
in $E_\lambda(-1)$). We therefore have 
$\tau(\chi,\psi)\cdot \tau(\chi^{-1},\psi^{-1}) = q$, i.e. 
$|\tau(\chi,\psi)|=q^{1/2}$. 





\subsection{}\label{VI:4-5}

If $k$ is an extension of degree $N$ of $\dF_q$, we can also regard 
\eqref{VI:eq:4-1-1} as a sum in $N$ variables over $\dF_q$. More 
generally, let $k$ be an étale algebra over $\dF_q$, of degree $N$ over 
$\dF_q$. It is a product of fields $k_i$ of degrees $N_i$, and we set 
\begin{equation*}\tag{4.5.1}\label{VI:eq:4-5-1}
  \varepsilon(k) = (-1)^{\sum (N_i+1)} \text{.} 
\end{equation*}
This is the sign of the permutation of $S=\hom_{\dF_q}(k,\dF)$ induced by the 
Frobenius substitution $\varphi\in \gal(\dF/\dF_q)$. 

Let $\psi:\dF_q\to E_\lambda^\times$ be a nontrivial character, and let 
$\chi:k^\times \to E_\lambda^\times$. We set 
\begin{equation*}\tag{4.5.2}\label{VI:eq:4-5-2}
  \tau_{\dF_q}(\chi,\psi) = (-1)^N \sum_{x\in k^\times} \psi \tr_{k/\dF_q}(x) \cdot \chi^{-1}(x) \text{.} 
\end{equation*}

If $\chi$ has coordinates $\chi_i:k_i^\times \to E_\lambda^\times$, we 
have the elementary identity 
\begin{equation*}\tag{4.5.3}\label{VI:eq:4-5-3}
  \tau_{\dF_q}(\chi,\psi) = \varepsilon(k) \prod \tau(\chi_i,\psi\circ \tr_{k_i/\dF_q}) \text{.}
\end{equation*}

After some preliminaries, we give in \ref{VI:4-10} a 
cohomological interpretation of the sums \eqref{VI:eq:4-5-2}. In 
\ref{VI:4-12}, we will interpret the Hasse--Davenport identity as a 
twisted form (cf. \ref{VI:1-12}) of the following special case of 
\eqref{VI:eq:4-5-3}: for $k=\dF_q^N$, and $\chi$ a character of 
$\dF_q^\times$, we have 
$\tau_{\dF_q}(\chi\circ N_{k/\dF_q},\psi)=\tau(\chi,\psi)^N$. 





\subsection{}\label{VI:4-6}

Recall that for $M$ a finitely generated projective module over a ring $A$, the 
functor $\spec(B)\mapsto M\otimes_A B$ from affine schemes over $\spec(A)$ 
to $\mathsf{Ens}$ is represented by the affine scheme $\dV(M^\vee)$ 
(EGA notation), the spectrum of $\operatorname{Sym}_A^\bullet(M^\vee)$. If 
$M=A^n$, this is the standard affine space of dimension $n$. 

Suppose that $M$ is a unital $A$-algebra, and let 
$V=\dV(M^\vee)$. By definition, for any $B$ equipped with an $A$-algebra structure, the set 
$V(B)$ of points of $V$ with coordinates in $B$ is $M\otimes_A B$. Since the 
functor $B\mapsto M\otimes_A B$ takes values in unital rings, 
$V$ is a unital ring scheme over $\spec(A)$. The 
norm and trace morphisms $M\otimes_A B\to B$ are functorial in $B$; they 
therefore correspond to morphisms of schemes $N$ and $T$ from $V$ to 
$\dG_a$. We shall consider the following schemes obtained from $V$: 
\begin{enumerate}[a)]
  \item $V^\ast$ is the open subscheme of invertible elements of $V$ (a group 
    scheme under $\cdot$);
  \item $W$ is the hyperplane of equation $T=0$ and $W^\ast=W\cap V^\ast$; 
  \item $P$ is the hyperplane at infinity $V\setminus \{0\}/\dG_m$ of the affine 
    space $V$ over $\spec(A)$. If $Q$ is the hyperplane at infinity of $W$, 
    $W^\ast/\dG_m$ and $V^\ast/\dG_m$ are open subsets of the projective spaces 
    $Q$ and $P$ over $\spec(A)$. 
    \begin{equation*}\tag{4.6.1}\label{VI:eq:4-6-1}
    \xymatrix{
      W^\ast \ar@{^{(}->}[r] \ar[d]^-\pi 
        & V^\ast \ar[d]^-\pi \\
      W^\ast/\dG_m \ar@{^{(}->}[r] 
        & V^\ast/\dG_m 
    }\qquad\subset\qquad
    \xymatrix{
      W\setminus \{0\} \ar@{^{(}->}[r] \ar[d]^-\pi 
        & V\setminus \{0\} \ar[d]^-\pi \\
      Q \ar@{^{(}->}[r] 
        & P \text{.}
    }
    \end{equation*}
\end{enumerate}

If $M=A^I$, $I$ a finite set, then $V\simeq \dG_a^I$, $V^\ast=\dG_m^I$, 
$V^\ast/\dG_m$ is the torus $\dG_m^I/(\dG_m\text{ diagonal})$, $N$ and $T$ 
are given by $\prod x_i$ and $\sum x_i$, and $Q$ is therefore the projective 
hyperplane with equation $\sum x_i=0$ in $P$. 

The case of interest is that in which $M$ is finite étale over $A$. The 
previous description then applies locally to $\spec(A)$ (for the 
étale topology) and we can write $V^\ast=\dG_m^I$, for $I$ a locally 
constant sheaf of finite sets on $\spec(A)$. If, for example, $A$ is a 
field with algebraic closure $\bar A$, and $M$ is a separable $A$-algebra, 
we put $I=\hom_A(M,\bar A)$, and over $\bar A$ we have 
$V\sim \dG_a^I$. Via this isomorphism, $N$ and $T$ are given by $\prod x_i$ and 
$\sum x_i$. 





\subsection{The Kummer torsor}\label{VI:4-7}

Let $S$ be a scheme, let $n$ be an integer invertible on $S$, and let $G$ be a 
commutative group scheme with connected fibers (written multiplicatively) over $S$. 
The sequence 
\[\xymatrix{
  0 \ar[r] 
    & G_n \ar[r] 
    & G \ar[r]^-{x^n} 
    & G \ar[r] 
    & 0 
}\]
is exact. It defines a $G_n$-torsor $K_n(G)$ (or simply $K_n$) on 
$G$. For $\chi:G_n \to E_\lambda^\times$ a homomorphism from the étale 
$S$-sheaf $G_n$ to the constant sheaf $E_\lambda^\times$, we denote by 
$\sK_n(\chi)$ the $E_\lambda$-sheaf $\chi^{-1}(K_n)$. 

The $K_n$ torsors form an inverse system of torsors under the 
projective system of groups $G_n$ (transition morphisms 
$x\mapsto x^d:G_{n d}\to G_n$). This expresses the commutativity of the diagrams 
\[\xymatrix{
  0 \ar[r] 
    & G_{n d} \ar[r] \ar[d]^-{x^d} 
    & G \ar[r]^-{x^{n d}} \ar[d]^-{x^d} 
    & G \ar[r] \ar@{=}[d] 
    & 0 \\
  0 \ar[r] 
    & G_n \ar[r] 
    & G \ar[r]^-{x^n} 
    & G \ar[r] 
    & 0 \text{.} 
}\]
We therefore have $\sK_n(\chi) = \sK_{n d}(\chi\circ x^d)$. 





\subsection{}\label{VI:4-8}

Apply the construction of \ref{VI:4-6} with $A=\dF_q$ and with $M=k$ an 
étale algebra over $\dF_q$. In accordance with our general conventions, we denote 
by a subscript $0$ the $\dF_q$-schemes denoted $V,V^\ast,\ldots$ in 
\ref{VI:4-6}. The same letters without a subscript denote the schemes over 
$\dF$ obtained by extension of scalars. 

We put $I=\hom(k,\dF)$; then $V\sim \dG_a^I$ and $T$ is given by 
$(x_i)\mapsto \sum x_i$. 
\begin{equation*}\tag{4.8.1}\label{VI:eq:4-8-1}
  \sF(\psi\circ \tr_{k/\dF_q}) = T^\ast \sF(\psi) = \bigotimes_i \pr_i^\ast \sF(\psi) \text{.} 
\end{equation*}
Since $V^\ast\sim \dG_m^I$, a character $\chi$ of $V_n^\ast\sim \dmu_n^I$, 
with values in $E_\lambda^\times$, amounts to a family 
$(\chi_i)_{i\in I}$ of characters of $\dmu_n$ indexed by $I$. It is 
defined over $\dF_q$ if, for all $\sigma\in \gal(\dF/\dF_q)$, we have 
$\chi_{\sigma_i} = \chi_i\circ \sigma^{-1}$. It then defines a 
$E_\lambda$-sheaf $\sK_n((\chi_i)_{i\in I})$ on $V_0^\ast$. If the product 
of the $\chi_i$ is trivial, $\chi$ factors through a character of 
$(V_0^\ast/\dG_m)_n$ and $\sK_n((\chi_i)_{i\in I})$ is the inverse image 
of a sheaf, denoted in the same way, on $V_0^\ast/\dG_m$. This construction 
is related to the Lang torsor as follows. 





\begin{lemma_}\label{VI:4-9}
\begin{enumerate}[(i)]
  \item For $n$ sufficiently divisible, we have a commutative diagram 
    \[\xymatrix{
      0 \ar[r] 
        & (V_0^\ast)_n \ar[r] \ar[d]^-\tau 
        & V_0^\ast \ar[r]^-{x^n} \ar[d]^-\tau 
        & V_0^\ast \ar[r] \ar@{=}[d] 
        & 0 \\
      0 \ar[r] 
        & k^\times \ar[r] 
        & V_0^\ast \ar[r]^-\fL 
        & V_0^\ast \ar[r] 
        & 0 \text{;} 
    }\]
    if $\chi$ is a character $\chi:k^\times \to E_\lambda^\times$, we have 
    $\sF(\chi) = \sK_n(\chi\circ\tau)$. 
  \item For sufficiently divisible $n$, $\tau$ identifies $k^\times$ with the coinvariants 
    of $\gal(\dF/\dF_q)$ acting on $V_n^\ast\sim \dmu_n^I$. 
  \item For $k$ a field, $N=[k:\dF_q]$, $\omega\in I$ an embedding of $k$ 
    in $\dF$, and $n=q^N-1$, the index component $\omega$ of 
    $\chi\circ\tau$ is $\chi\circ \omega^{-1}$. 
  \item If $\chi$ is nontrivial on each factor of $k$, the 
    $(\chi\circ \tau)_i$ are all nontrivial. 
\end{enumerate}
\end{lemma_}

It is enough to prove the lemma when $k$ is a field of degree $N$ over 
$\dF_q$. In this case, $n$ is ``sufficiently divisible'' if $q^N-1\mid n$. 
Choose an embedding $\omega$ of $k$ into $\dF$; then $I$ is identified 
with $\dZ/N$: the element $i\in \dZ/N$ corresponds to $\omega_i=\omega^{q^i}$. Via 
the isomorphism $V_0^\ast(\dF) = {\dF^\times}^I$, we have 
\begin{enumerate}[\indent a)]
  \item $x\in k^\times$ corresponds to $(\omega_i(x))\in {\dF^\times}^I$; 
  \item $F((x_i)_{i\in \dZ/N})=(x_{i-1}^q)_{i\in \dZ/N}$. 
\end{enumerate}

A character $\chi=(\chi_i)$ of $V_n^\ast\sim \dmu_n^I$ is defined over 
$\dF_q$ if $\chi_{-i}(x) =\chi_0(x^{q^i})$ ($i\in \dZ$). If $q^N-1\mid n$, there are 
$q^N-1$ such characters: $\chi_0$ factors through $\dmu_{q^N-1}$ and 
determines the $\chi_i$. To prove (ii), it is therefore sufficient to verify (i) 
and (iii) (or its corollary (iv)) which ensures that $\chi\mapsto \chi\circ \tau$ 
is injective. 

Let us prove (i) and (iii) for $n=q^N-1$. Write the group of 
endomorphisms of $V^\ast\sim \dG_m^I$ additively; in particular, denote by $n$ 
the $x\mapsto x^n$ operator. If $\alpha$ is the cyclic permutation operator 
$(x_i)\mapsto (x_{i-1})$, we have 
$q^N-1=(q\alpha)^N-1 = (q\alpha-1)((q\alpha)^{N-1}+\cdots + 1)$. This 
determines $\tau$: we have $\tau((x_i)) = \prod_{0\leqslant j<N} x_{i-j}^{q^j}$, 
and the induced map from $(V_0^\ast)_n$ to $k^\times$ is [as in a) above] 
$(x_i)\mapsto \prod x_{-i}^{q^i}$; hence (iii) follows. 





\begin{proposition_}\label{VI:4-10}
The cohomology of $V^\ast$ with coefficients in 
$\sF(\chi^{-1}\cdot \psi\tr_{k/\dF_q})$ satisfies 
\begin{enumerate}[\indent (i)]
  \item $\h_c^i = 0$ for $i\ne N$, and $\dim \h_c^N=1$. 
  \item On $\h_c^N$, $F^\ast$ is multiplication by 
    $\tau_{\dF_q}(\chi,\psi)$. 
  \item If $\chi$ is nontrivial on each factor of $k$, we have 
    $\h_c^\bullet \iso \h^\bullet$. 
\end{enumerate}
\end{proposition_}

Applying \ref{VI:4-9}(i), over $\dF$, if $\chi\circ \tau=(\chi_i)_{i\in I}$, 
we have $\sF(\chi) = \sK_n((\chi_i)_{i\in I})$. Points (i) and (iii) 
therefore follow from the following more geometric statement (for (iii), 
apply \ref{VI:4-9}(iv)), and (ii) follows from the trace formula: 





\begin{proposition_}\label{VI:4-11}
Let $(\chi_i)_{i\in I}$ be a family of characters of $\dmu_n(\dF)$. The 
cohomology of $V^\ast\simeq \dG_m^I$ with coefficients in 
$\sK_n((\chi_i)_{i\in I})\otimes \sF(\psi\tr_{k/\dF_q})$ satisfies 
\begin{enumerate}[\indent (i)]
  \item $\h_c^i{}=0$ for $i\ne N$, and $\dim \h_c^N{} = 1$. 
  \item If the $\chi_i$ are all nontrivial, we have 
    $\h_c^\bullet{}\iso \h^\bullet$. 
\end{enumerate}
\end{proposition_}

We have $V^\ast\sim \dG_m^I$, 
$\sK((\chi_i)) = \bigotimes_i \pr_i^\ast \chi_i(K_n(\dG_m))$ and 
$\sF(\psi\circ \tr_{k/\dF_q})=T^\ast \sF(\psi) = \bigotimes_i \pr_i^\ast \sF(\psi)$. 
The K\"unneth formula then applies, and \ref{VI:4-11} follows from 
\ref{VI:4-3}. 





\subsection{}\label{VI:4-12}

From \ref{VI:2-4*}* and (\hyperref[IV]{Cycle}, \ref{IV:1-3} Example 2), it also follows 
that the group of permutations $\sigma$ of $I$ such that 
$\chi_i = \chi_{\sigma i}$ ($i\in I$) acts on $\h_c^N$ by multiplication by 
the sign $\varepsilon(\sigma)$. We deduce a second proof 
of the identity of Hasse-Davenport. 

If $\chi$ is a character of $\dF_q^\times$, the fact that, over $\dF$, $N$ 
is given by $(x_i)\mapsto \prod x_i$ gives 
$\sF(\chi\circ N)=N^\ast\sF(\chi) = \bigotimes_i \pr_i^\ast\sF(\chi)$ 
and the K\"unneth formula gives 
\[
  \h_c^\bullet(V^\ast,\sF(\chi^{-1}\circ N,\psi\circ \tr)) \sim \h_c^\bullet(\dG_m,\sF(\chi^{-1}\psi))^{\otimes I} \text{,} 
\]
where the tensor product on the right-hand side is taken in the sense of 
\ref{VI:2-4*}*. Since $\h_c^1(\dG_m,\sF(\psi\chi^{-1}))$ is of dimension 
$1$, its tensor power $\otimes I$, in the ordinary sense, depends only on 
the cardinality $n$ of $I$. Applying (\hyperref[IV]{Cycle}, \ref{IV:1-3} Example 
2), we obtain a canonical isomorphism 
\begin{equation*}\tag{4.12.1}\label{VI:eq:4-12-1}
  \h_c^N\left(V^\ast,\sF(\chi^{-1}\circ N)\otimes T^\ast \sF(\psi)\right) \sim \h_c^1(\dG_m,\sF(\psi\chi^{-1})) \otimes_\dZ \textstyle \bigwedge^N \dZ^I \text{.}
\end{equation*}
To compute the action of $F^\ast$, it is most convenient to adopt the Galois 
point of view and observe that, since the isomorphism \eqref{VI:eq:4-12-1} is 
canonical, it is compatible with the transport-of-structure action of 
$\gal(\dF/\dF_q)$. If $\varepsilon(k)$ is the sign of the permutation 
$\varphi$ of $I$, we find that 
\begin{align*} 
  \tau_{\dF_q}(\chi\circ N_{k/\dF_q},\psi) 
    &= \tr\left(\varphi^{-1},\h_c^N\left(V^\ast,N^\ast\sF(\chi^{-1})\otimes T^\ast \sF(\psi)\right)\right) \\ 
    &= \varepsilon(k) \tr\left(\varphi^{-1},\h_c^1(\dG_m,\sF(\chi^{-1}\psi))\right)^N \\
    &= \varepsilon(k) \tau(\chi,\psi)^N \text{,} 
\end{align*}
that is, 
\begin{equation*}\tag{4.12.2}\label{VI:eq:4-12-2}
  \tau_{\dF_q}\left(\chi\circ N_{k/\dF_q},\psi\right) = \varepsilon(k) \tau(\chi,\psi)^N \text{.} 
\end{equation*}

If $k$ is a field, $\varphi$ is a cyclic permutation of $I$, 
$\varepsilon(k) = (-1)^{N+1}$, and we obtain the Hasse--Davenport identity: 
\[
  \tau\left(\chi\circ N_{k/\dF_q},\psi\circ \tr_{k/\dF_q}\right) = (-1)^{N+1} \tau_{\dF_q}\left(\chi\circ N_{k/\dF_q},\psi\right) = \tau(\chi,\psi)^N \text{.} 
\]





\begin{lemma_}\label{VI:4-13}
For $\chi\circ \tau=(\chi_i)_{i\in I}$ as in \ref{VI:4-9}(i), the
following conditions are equivalent 
\begin{enumerate}[\indent (i)]
  \item $\chi| \dF_q^\times$ is trivial. 
  \item The product of $\chi_i$ is trivial. 
  \item $\sF(\chi) = \sK_n((\chi_i))$ is the inverse image of a unique 
    sheaf on $V^\ast/\dG_m$. 
  \item The inverse image of $\sF(\chi)$ on $\dG_m$ (embedded in 
    $V^\ast$ by the structural morphism $\dF_q \to k$) is trivial. 
\end{enumerate}
\end{lemma_}

(i)$\Rightarrow$(ii). We regard $\chi$ as a character of 
$V^\ast/\dG_m(\dF_q)=k^\times/\dF_q^\times$, and we take $\sF(\chi)$ on 
$V^\ast/\dG_m$. 

(ii)$\Rightarrow$(iii). Likewise, we regard $(\chi_i)$ as a character 
of $(V^\ast/\dG_m)_n$. The uniqueness in (iii) follows from the fact that $V^\ast$ is 
a fibration with connected fibers over $V^\ast/\dG_m$, and (iii)$\Rightarrow$(iv)  
is trivial. 

(iv)$\Rightarrow$(i), (ii). This inverse image is 
$\sF(\chi|\dF_q^\times)$ and $\sK_n(\prod \chi_i)$. 





\subsection{}\label{VI:4-14}

Let $k$ be an étale algebra over $\dF_q$, of 
dimension $N+1$, and let $\chi$ be a nontrivial character of $k^\times$, trivial on 
$\dF_q^\times$. The Jacobi sum $J(\chi)$ is defined by 
\begin{equation*}\tag{4.14.1}\label{VI:eq:4-14-1}
  J(\chi) = (-1)^{N-1} \sum_{\substack{x\in k^\times/\dF_q^\times \\ \tr(x)=0}} \chi^{-1}(x) \text{.} 
\end{equation*}

We have the following identity between Gauss sums and Jacobi sums. 





\begin{proposition_}\label{VI:4-15}
For $\chi$ as above and $\psi$ a nontrivial additive character of 
$\dF_q$, we have 
\[
  q J(\chi) = \tau_{\dF_q}(\chi,\psi) \text{.} 
\]
\end{proposition_}

In the special case $k=\dF_q^n$, formula \eqref{VI:eq:4-5-3} rewrites this as 
\begin{equation*}\tag{4.15.1}\label{VI:eq:4-15-1}
  q J(\chi) = \prod_i \tau(\chi_i,\psi) \text{,} 
\end{equation*}
or equivalently as 
\begin{align*}
  \chi_0(-1) J(\chi) &= (-1)^{N-1} \sum_{\substack{x_1,\dots,x_N\in \dF_q^\times \\ \sum x_i = 1}} \prod_{i=1}^N \chi_i^{-1}(x_i) = \tau(\chi_0^{-1},\psi)^{-1} \prod_{i=1}^N \tau(\chi_i,\psi) \\
  \chi_0^{-1} &= \prod_{i=1}^N \chi_i 
\end{align*}





\begin{proof}[Proof of \ref{VI:4-15}]
\[
  \tau_{\dF_q}(\chi,\psi) = (-1)^{N+1} \sum_{x\in k^\times} \chi(x)^{-1} \psi \tr(x) = (-1)^{N+1} \sum_{x\in k^\times/\dF_q^\times} \chi(x)^{-1} \sum_{\lambda\in \dF_q^\times} \psi(\lambda \tr x) \text{.}
\]
The sum $\sum_{\lambda\in \dF_q^\times} \psi(\lambda \tr x)$ equals $q-1$ if 
$\tr(x)=0$, and $-1$ if $\tr(x)\ne 0$. Therefore, 
\[
  \tau_{\dF_q}(\chi,\psi) = (-1)^{N+1} \left(q\sum_{\substack{x\in k^\times/\dF_q^\times \\ \tr(x)=0}} \chi(x)^{-1} - \sum_{x\in k^\times/\dF_q^\times} \chi(x)^{-1}\right) \text{.}
\]
The second term on the right-hand side is zero, since it is the sum of the values 
of a character, and \ref{VI:4-15} follows. 
\end{proof}

Proposition \ref{VI:4-15} therefore has the following cohomological counterpart: 





\begin{proposition_}\label{VI:4-16}
The cohomology of $W^\ast/\dG_m$ (\ref{VI:4-6}, \ref{VI:4-8}) with coefficients 
in $\sF(\chi^{-1})$ satisfies 
\begin{enumerate}[\indent (i)]
  \item $\h_c^i=0$ for $i\ne N-1$, and $\dim \h_c^{N-1}{}=1$. 
  \item On $\h_c^{N-1}$, $F^\ast$ is multiplication by $J(\chi)$. 
  \item $\h_c^{N-1}$ is canonically isomorphic to 
    $\h_c^{N+1}\left(V^\ast,\sF(\chi^{-1},\psi\circ \tr_{k/\dF_q})\right)(1)$. 
\end{enumerate}
\end{proposition_}

Assertion (ii) follows from (i) and the trace formula; (i) and (iii) 
follow from \ref{VI:4-11} and the following more geometric statement. 





\begin{proposition_}\label{VI:4-17}
Let $(\chi_i)_{i\in I}$ be a family of characters of 
$\dmu_n(\dF)$ that are not all trivial and whose product is $1$. Then 
$\h_c^{i-1}\left(W^\ast/\dG_m,\sK_n((\chi_i)_{i\in I})\right)$ is 
canonically isomorphic to 
\newline
$\h_c^{i+1}\left(V^\ast,\sK_n((\chi_i)_{i\in I})\otimes \sF(\psi\tr_{k/\dF_q})\right)(1)$. 
\end{proposition_}

The first line of \ref{VI:4-15}, namely the projection formula, gives ($\pi$ as in 
\eqref{VI:eq:4-6-1}). 
\begin{equation*}\tag{4.17.1}\label{VI:eq:4-17-1}
  \eR^i \pi_!\left(\sK_n((\chi_i)_{i\in I})\otimes \sF(\psi \tr_{k/\dF_q})\right) = \sK_n((\chi_i)_{i\in I}) \otimes \eR^i \pi_! \sF(\psi\tr_{k/\dF_q}) 
\end{equation*}
(on $V^\ast/\dG_m$). Let us compute the sheaves 
$\eR^i\pi_! \sF(\psi\tr_{k/\dF_q}) = \eR^i \pi_! T^\ast \sF(\psi)$. Let 
$v_0:\widetilde V_0 \to V_0$ be the blow-up of $V_0$ at $\{0\}$. In the 
diagram 
\begin{equation*}\tag{4.17.2}\label{VI:eq:4-17-2}
\xymatrix{
  V_0\setminus \{0\} \ar@{^{(}->}[r] \ar[dr]_-\pi 
    & \widetilde V_0 \ar[r]^-{v_0} \ar[d]^-{\bar\pi} 
    & V_0 \\
  & P_0 
}
\end{equation*}
$\pi$ is a fibration whose fibers are punctured affine lines, $\widetilde V_0$ 
is the corresponding line bundle, and $Z_0=v_0^{-1}(0)$ is its section $0$. 
The $T^\ast\sF(\psi)$ sheaf on $V_0\setminus \{0\}$ extends to 
$(T v_0)^\ast \sF(\psi)$ on $\widetilde V_0$, and the restriction of this sheaf 
to $Z_0$ is the constant sheaf $E_\lambda$, because $T_0 v_0$ vanishes on 
$Z_0$. The long exact sequence for cohomology with compact supports is therefore 
\begin{equation*}\tag{4.17.3}\label{VI:eq:4-17-3}
\xymatrix{
  \ar[r] 
    & \eR^i \pi_! T^\ast \sF(\psi) \ar[r] 
    & \eR^i \bar\pi_! (T v_0)^\ast \sF(\psi) \ar[r] 
    & (E_\lambda\text{ for $i=0$, $0$ otherwise}) \ar[r] 
    & 
}
\end{equation*}

The inverse image $\bar\pi^{-1}(Q_0)$ is the blow-up $\widetilde W_0$ of 
$W_0$ at $0$; $T_0 v_0$ vanishes on $\widetilde W_0$; hence 
$(T v_0)^\ast \sF(\psi)=E_\lambda$ on $\bar\pi^{-1}(Q_0)$ and, since $\bar\pi$ is an 
affine-line fibration, 
\begin{equation*}\tag{4.17.4}\label{VI:eq:4-17-4}
  \eR^i \bar\pi_! (T v_0)^\ast \sF(\psi)|Q = 
    \begin{cases}
      0 & \text{for $i\ne 2$} \\
      E_\lambda(-1) & \text{for $i=2$.} 
    \end{cases} 
\end{equation*}

If $x\in P$, $x\notin Q$, the line $D=\pi^{-1}(x)$ is mapped 
isomorphically onto $\dG_a$ by $T v_0$; therefore (by 
\ref{VI:2-7}*) 
\begin{align} \notag 
  \h_c^\bullet(D,(T v_0)^\ast \sF(\psi)) &= \h_c^\bullet(\dG_a,\sF(\psi)) = 0 && \text{and} \\
  \tag{4.17.5}\label{VI:eq:4-17-5}
  \eR^i \bar\pi_! (T v_0)^\ast \sF(\psi) 
    &= \begin{cases} 
         0 & \text{for $i\ne 2$} \\
         E_\lambda(-1)_Q & \text{for $i=2$} 
       \end{cases} 
\end{align}

Combining \eqref{VI:eq:4-17-3} and \eqref{VI:eq:4-17-5}, we finally find 
\begin{equation*}\tag{4.17.6}\label{VI:eq:4-17-6}
  \eR^i \pi_! T^\ast\sF(\psi) = 
    \begin{cases}
      0 & \text{for $i\ne 1,2$} \\
      E_\lambda & \text{for $i=1$} \\
      E_\lambda(-1)_Q & \text{for $i=2$.} 
    \end{cases}
\end{equation*}





\subsection{}\label{VI:4-18}

Let us compute the compactly supported cohomology of the sheaf 
$\sK_n((\chi_i))\otimes T^\ast \sF(\psi)$ on $V^\ast$ using the 
Leray spectral sequence of $\pi:V^\ast\to V^\ast/\dG_m$. Applying 
\eqref{VI:eq:4-17-1} and \eqref{VI:eq:4-17-6}, we obtain as initial terms 
$E_2^{p,1} = \h_c^p(V^\ast/\dG_m,\sK_n(\chi_i)) = 0$ (because $(\chi_i)\ne 1$) 
and $E_2^{p,2} = \h_c^p(W^\ast/\dG_m,\sK(\chi_i))(-1)$. 

The spectral sequence reduces to an isomorphism, and \ref{VI:4-17} follows. 





\subsection{}\label{VI:4-19}

The sheaves $\sK_n((\chi_i)_{i\in I})$ make sense over any field, 
or even over any base scheme. This allows us to generalize 
\ref{VI:4-16}(i). With the notation of \ref{VI:4-6}, suppose that $M$ is finite 
étale of constant rank $N$ over $A$, and let $n$ be invertible in $A$. 
Let $a$ be the projection of $W^\ast/\dG_m$ to $\spec(A)$. Let also $\chi$ be 
a character (morphism of sheaves) 
$\chi:(V^\ast/\dG_m)_n \to E_\lambda^\times$, and $\sK_n(\chi)$ the sheaf 
corresponding. Locally for the étale topology, we can regard $\chi$ 
as a family of characters $(\chi_i)_{i\in I}$ of $\dmu_n$, whose product is 
trivial.

\begin{proposition_}\label{VI:4-20}
\begin{enumerate}[(i)]
  \item If $\chi$ is everywhere nontrivial, we have 
    $\eR^i a_! (\sK_n(\chi)) = 0$ for $i\ne N-2$, and 
    $\eR^{N-2} a_!(\sK_n(\chi))$ is smooth of rank $1$. 
  \item An automorphism $\sigma$ of $M$ which preserves $\chi$ acts on this 
    $\eR^{N-2} a_!$ by multiplication by the sign $\varepsilon(\sigma)$ 
    of $\sigma$, seen as a permutation of $I$. 
  \item If the $\chi_i$ are all nontrivial, we have $\eR a_! \iso \eR a_\ast$. 
\end{enumerate}
\end{proposition_}
\begin{proof}[Proof]
The scheme $W^\ast/\dG_m$ is the complement of a relative normal-crossings 
divisor in the scheme $Q$, which is proper and smooth over $\spec(A)$, and 
$\sK_n(\chi)$ is locally constant on $W^\ast/\dG_m$ and is tamely ramified 
along the boundary. Consequently, the $\eR^i a_!$ and 
$\eR^i a_\ast$ are smooth sheaves, compatible with every change of 
base. A standard argument then reduces us to the case in which $A$ is a 
finite field, and (i) follows from \ref{VI:4-17} and \ref{VI:4-11}(i). 
To prove (ii), use the compatibility of the action of $\sigma$ with 
the isomorphism \ref{VI:4-17}, together with \ref{VI:4-12}. For (iii), note that if the 
$\chi_i$ are all nontrivial, then $\sK_n(\chi)$ on 
$W^\ast/\dG_m\subset Q$ is ramified along every divisor at infinity, 
and apply \ref{VI:1-19-1}. 
\end{proof}










\section{Hecke characters}\label{VI:5}





\subsection{}\label{VI:5-1}

Let $F$ be a number field (of finite degree over $\dQ$) and let $k$ be a field of 
characteristic $0$. Here are several equivalent descriptions of an 
\emph{algebraic} homomorphism $\alpha:F^\times \to k^\times$. 
\begin{enumerate}[(i)]
  \item Let $\{e_i\}$ be a basis of $F$ over $\dQ$. The homomorphism $\alpha$ is 
    algebraic if it is given by a formula 
    $\alpha(\sum x^i e_i) = A(x^i)$, $A\in k(X^i)$. 
\end{enumerate}

This means that $\alpha$ agrees on $F^\times$ with a rational map defined over $k$ 
from $\res_{F/\dQ}(\dG_m)$ to $\dG_m$. By the Zariski density of 
$F^\times$, and because a rational homomorphism is defined everywhere, such 
a map is a homomorphism of group schemes: 
\begin{enumerate}[(i)]
\setcounter{enumi}{1}
  \item $\alpha$ is induced by a homomorphism of group schemes over $k$ 
    \[
      \res_{F/\dQ}(\dG_m)\otimes_\dQ k \to \dG_m \text{.} 
    \] 
\end{enumerate}

If $k$ is a number field, the adjunction property of restriction of 
scalars $R$ shows that this is equivalent to 
\begin{enumerate}[(i)]
\setcounter{enumi}{2}
  \item ($k$ a number field) $\alpha$ is induced by a homomorphism of 
    group schemes over $\dQ$: $\res_{F/\dQ}(\dG_m) \to \res_{k/\dQ}(\dG_m)$. 
\end{enumerate}

Let $\bar k$ be an algebraic closure of $k$, and let $I=\hom(F,\bar k)$. Over 
$\bar k$, the character group of $\res_{F/\dQ}(\dG_m)$ is $\dZ^I$, with basis 
the embeddings of $F$ into $\bar k$. The characters defined over $k$ 
are those invariant under $\gal(\bar k/k)$; they can be described either as 
those sending $F^\times$ into $k^\times\subset \bar k$, or in terms of 
orbits of $\gal(\bar k/k)$ on $I$, which in turn correspond to the factors of 
$F\otimes k$. 
\begin{enumerate}[(i)]
\setcounter{enumi}{3}
  \item $\alpha$ is of the form 
    $\alpha=\prod_{\omega\in I} \omega^{n_\omega}$. The allowed families of exponents 
    $(n_\omega)$ are those such that $n_\omega=n_{\omega'}$ for 
    $\omega$ and $\omega'$ in the same orbit of $\gal(\bar k/k)$. Equivalently, these 
    are those such that $\prod_\omega \omega(x)^{n_\omega}\in k$ for every 
    $x\in F$. 
  \item Write $F\otimes k=\prod_{j\in J} F_j$, where the $F_j$ are fields. 
    Then $\alpha$ can be written $\alpha=\prod N_{F_j/k}^{m_j}$. 
\end{enumerate}

In the special case where $k$ contains a normal closure of $F$, these 
expressions become: $\alpha=\prod \omega^{n_\omega}$, where $\omega$ 
runs over the $[F:\dQ]$ embeddings of $F$ into $k$. 





\subsection{}\label{VI:5-2}

Let $k$ be a number field, and let $\alpha$ be an algebraic homomorphism 
from $F^\times$ to $k^\times$. There is then a unique 
homomorphism, also denoted $\alpha$, of the group $I(F)$ 
of fractional ideals of $F$ to that of $k$, such that $\alpha((x))=(\alpha(x))$. 
Uniqueness follows from the fact that every ideal $a$ has a power that is 
principal, and that $I(k)$ is torsion-free. To prove 
existence, we use for example \ref{VI:5-1}(v): we have 
$\alpha(a) = \prod N_{F_j/k}^{m_j}((a))$. 





\subsection{}\label{VI:5-3}

Let us recall the definition of algebraic Hecke characters, also called 
Weil Hecke characters (or gr\"ossencharaktere) of type $A_0$. Let $F$ be 
a number field, let $\fm$ be an ideal of $F$ (i.e., of the ring of integers of 
$F$), let $I_\fm$ be the group of fractional ideals of $F$ prime to $\fm$, 
and let $k$ be a field of characteristic $0$. A homomorphism 
$\chi:I_\fm \to k^\times$ is an \emph{algebraic Hecke character} (of 
conductor $\leqslant \fm$) if there exists an algebraic homomorphism 
$\chi_\text{alg}:F^\times \to k^\times$ satisfying 
\begin{equation*}\tag{$*$}\label{VI:eq:5-3-1}
\text{For $x\in F^\times$, prime to $\fm$, totally positive, and 
$\equiv 1\pmod\fm$, one has $\chi((x))=\chi_\text{alg}(x)$.}
\end{equation*}

By the density of the set of $x$ in \eqref{VI:eq:5-3-1}, $\chi_\text{alg}$ 
is entirely determined by $\chi$. This is the \emph{algebraic part} 
of $\chi$. If $\chi((x))=\chi_\text{alg}(x)$ for $x$ totally positive and 
$\equiv 1\pmod{\fm'}$, $\chi$ extends to a Hecke character of 
conductor $\leqslant \inf(\fm,\fm')$: $I_{\fm+\fm'}\to k^\times$. We 
identify Hecke characters that coincide on their common domain of 
definition, and we call \emph{conductor} of $\chi$ the smallest $\fm'$ 
such that $\chi$ is of conductor $\leqslant \fm'$. 





\subsection{Remark}\label{VI:5-4}

If an algebraic Hecke character $\chi$ takes its values in a 
subfield $k'$ of $k$, it is already a Hecke character with values in 
$k'$: it suffices to observe that $\chi_\text{alg}:\res_{F/\dQ}(\dG_m) \to \dG_m$ 
is already defined over $k'$, and this follows from the Zariski density 
in $\res_{F/\dQ}(\dG_m)$ of the set of totally positive $x$ 
$\equiv 1\pmod\fm$. We can always take for $k'$ a subfield of $k$ of 
finite degree over $\dQ$. 





\subsection{}\label{VI:5-5}

If $\varepsilon$ is a totally positive unit $\equiv 1\pmod\fm$, we have 
$\chi_\text{alg}(\varepsilon) = \chi((\varepsilon)) = 1$. The homomorphism 
$\chi_\text{alg}$ therefore factors through the quotient $T_\fm$ of 
$\res_{F/\dQ}(\dG_m)$ by the Zariski closure of the group 
$E_\fm\subset F^\times$ of totally positive units $\equiv 1\pmod\fm$. 

According to Serre \cite[II.3]{se68}, if $k$ is an algebraic closure of 
$\dQ$, and $\fm$ is large enough, the characters $\prod \omega^{n_\omega}$ 
of $\res_{F/\dQ}(\dG_m)$ of the form $\chi_\text{alg}$ are characterized 
as follows: there must exist an integer $N$, the \emph{weight} of $\chi$ (or 
$\chi_\text{alg}$), such that for any element $\sigma$ of $\gal(k/\dQ)$ 
conjugate to complex conjugation, we have $n_\omega+n_{\sigma \omega}=N$. If 
$F_1$ is the largest subfield of $F$ which is a quadratic extension 
that is totally imaginary over a totally real field $F_1'$, this amounts to 
  saying that $\chi_\text{alg}$ is of the form $\chi_1\circ N_{F/F_1}$, and that 
  $\chi_1|F_1'=(N_{F_1'/\dQ})^N$. 

If $\chi$ is of weight $N$, for any ideal $\fa$ of $F$, $\chi(\fa)$ is a 
algebraic number all of whose conjugates have absolute value 
$N(\fa)^{N/2}$ \cite[II.3 prop.2]{se68}. 





\subsection{}\label{VI:5-6}

For a fixed number field $k$, additional conditions are 
imposed on $\chi_\text{alg}$. For any ideal $\fa$ of $F$ prime to the 
conductor, we have in fact 
\begin{equation*}\tag{5.6.1}\label{VI:eq:5-6-1}
  \chi_\text{alg}(\fa) = (\chi(\fa)) \text{,} 
\end{equation*}
so that $\chi_\text{alg}(\fa)$ is principal. Since the group of 
  ideal classes of $k$ is torsion, it is enough to prove the $n$-th power 
of \eqref{VI:eq:5-6-1}, for a suitable $n\ne 0$; this allows us 
to replace $\fa$ by $\fa^n$, therefore to assume $\fa=(x)$, with $x$ totally 
positive $\equiv 1\pmod\fm$. All that remains is to use the definitions. 

Together with \ref{VI:5-5}, this shows that $\chi_\text{alg}$ determines the 
valuations of $\chi(\fa)$ at all places of $k$. 





\subsection{}\label{VI:5-7}

Serre's group $S_\fm$ can be characterized as the 
group of multiplicative type whose group of characters (on any 
field) is the group of algebraic Hecke characters of conductor 
$\leqslant \fm$ (see \cite[II 2.1 et 2.2]{se68}). Its relation with 
$\ell$-adic representations is explained in \cite[II 2.3]{se68}.





\begin{theorem_}[{\cite{se68}}]\label{VI:5-8}
Let $\chi$ be an algebraic Hecke character of $F$ with values in $E_\lambda$, for 
$E_\lambda$ a finite extension of $\dQ_\ell$. There is then a unique 
  homomorphism $\chi_\lambda:\gal(\bar F/F)^\textnormal{ab} \to E_\lambda^\times$, such 
that 
\begin{enumerate}[\indent (i)]
  \item $\chi_\lambda$ is unramified outside the conductor $\ff$ of 
    $\chi$ and outside $\ell$. 
  \item For $\fp$ a prime ideal of $F$ prime to $\ff$ and to $\ell$, 
    and $F_\fp\in \gal(\bar F/F)^\textnormal{ab}$ the geometric Frobenius 
    at $\fp$, we have 
    \[
      \chi_\lambda(F_\fp) = \chi(\fp) \text{.}
    \] 
\end{enumerate}
\end{theorem_}





\subsection{}\label{VI:5-9}

At the end of this section, we give a criterion for a 
$\lambda$-adic representation to arise from a Hecke character, 
and in the next section we apply it to Jacobi sums. We 
thereby recover, in a somewhat generalized form, Weil's results 
\cite{we52,we74-2}, with the second part of Weil's proof 
replaced by a $\ell$-adic cohomology argument. 





\begin{theorem_}\label{VI:5-10}
Let $F$ and $k$ be two number fields, let $\lambda$ be a place of $k$ of 
residue characteristic $\ell$, and let 
$\chi_\lambda:\gal(\bar F/F)^\textnormal{ab}\to k_\lambda^\times$ be a 
homomorphism, unramified outside $\ell$ and a finite set $S$ of 
places. We assume that there exists a set $T$ of places (containing $S$ and the 
places above $\ell$), of density $0$, and an algebraic homomorphism 
$\chi_0:F^\times \to k^\times$ such that 
\begin{enumerate}[(i)]
  \item For $\bar k$ an algebraic closure of $k$, 
    $\chi_0:F^\times \to k^\times \hookrightarrow \bar k^\times$ is of the type 
    considered in \ref{VI:5-5}, of weight $N$; 
  \item For $\fp$ prime $\notin T$, $\chi_\lambda(F_\fp)$ lies in 
    $k^\times$, $(\chi_\lambda(F_\fp))=\chi_0(\fp)$, and all complex conjugates 
    of $\chi_\lambda(F_\fp)$ have absolute value $(N \fp)^{N/2}$. 
\end{enumerate}

Then $\chi_\lambda$ is defined (\ref{VI:5-8}) by an algebraic Hecke character 
of $F$ with values in $k$, whose algebraic part is $\chi_0$. 
\end{theorem_}

Let $\chi'$ be an algebraic Hecke character with values in a 
finite Galois extension $k'$ of $k$, of algebraic part $\chi_0$ 
(\ref{VI:5-5}). After enlarging $T$ if necessary, we may suppose that the conductor of 
$\chi'$ is supported in $T$. Let $\lambda'$ be a place of $k'$ above 
$\lambda$, and let $\chi_{\lambda'}$ be the composite 
$\gal(\bar F/F)^\text{ab}\to k_\lambda^\times \to {k'}_{\lambda'}^\times$, and 
$\chi_{\lambda'}':\gal(\bar F/F)^\text{ab}\to {k'}_{\lambda'}^\times$ defined 
by $\chi'$ (\ref{VI:5-8}). Set 
$\varepsilon = \chi_{\lambda'} \cdot {\chi'}_{\lambda'}^{-1}$.

Hypothesis (ii), together with \ref{VI:5-5}, \ref{VI:5-6}, and \ref{VI:5-8}, ensures that 
for $\fp\notin T$, we have $\varepsilon(F_\fp)\in {k'}^\times$, and that in every 
completion of $k'$ this number has absolute value $1$. It follows that 
$\varepsilon(F_\fp)$ belongs to the finite group $\dmu$ of roots of unity 
in $k'$. By the Chebotarev density theorem and the density-zero hypothesis 
on $T$, the $F_\fp$ ($\fp\notin T$) are dense in 
  $\gal(\bar F/F)^\text{ab}$, so that $\varepsilon(\sigma)\in \dmu$ for all 
$\sigma\in \gal(\bar F/F)^\text{ab}$: the character $\varepsilon$ is a 
  finite order character $\gal(\bar F/F)^\text{ab}\to \dmu$. By 
class field theory, it corresponds to a finite-order character 
of the idèle class group of $F$, i.e. to an algebraic Hecke character 
with trivial algebraic part. After correcting $\chi'$ by this 
character, we are reduced to the case in which 
$\chi_{\lambda'}=\chi_{\lambda'}'$, and it remains to show that $\chi'$, 
a priori taking values in $k'$, in fact takes values in $k$. If 
$\sigma\in \gal(k'/k)$, $\chi'$ and ${\chi'}^\sigma$ coincide on every prime 
ideal $\fp\notin T$, since 
$\chi'(\fp) = \chi'_{\lambda'}(F_\fp) = \chi_{\lambda'}(F_\fp) \in k^\times$. 
The characters $\chi_{\lambda'}'$, and ${\chi_{\lambda'}'}^\sigma$, coincide 
on a dense subset of $\gal(\bar F/F)^\text{ab}$ and hence are equal. Thus 
$\chi'={\chi'}^\sigma$, so $\chi'$ takes values in $k$. 










\section{Hecke characters defined by Jacobi sums}\label{VI:6}

The idea of this section is as follows. Consider a Gauss sum 
$\tau_{\dF_q}(\chi,\psi)$. Both this sum and the corresponding sheaf 
$\sF(\chi^{-1}\psi)$ are tied in two ways to the characteristic $p$ and to the 
base field $\dF_q$: 
\begin{enumerate}[\indent a)]
  \item because the additive character $\psi$ and the sheaf 
    $\sf(\psi)$ occur in it; 
  \item because $\sF(\chi)$ is defined from the exact Lang sequence 
    \[\xymatrix{
      0 \ar[r] 
        & k^\times \ar[r] 
        & V^\ast \ar[r] 
        & V^\ast \ar[r] 
        & 0 \text{.} 
    }\]
\end{enumerate}
These present two difficulties if we wish to lift the corresponding 
cohomological constructions to a number field $F$ and find 
$\ell$-adic representations of $\gal(\bar F/F)$ whose Frobenius eigenvalues 
are Gauss sums. 

If $\chi$ is trivial on $\dF_q^\times$, the sum $\tau_{\dF_q}(\chi,\psi)$ 
is independent of $\psi$ and is $q$ times a Jacobi sum. The latter is 
related to the cohomology of $W^\ast/\dG_m$, with coefficients in a sheaf 
obtained from $\chi$: thus $\psi$, and with it difficulty a), disappear. To 
resolve b), it suffices to describe the sheaves in question by means of 
exact K\"ummer sequences (which make sense in every characteristic). This is 
possible by \ref{VI:4-9}. 

Once the $\ell$-adic representations have been found, we use 
\ref{VI:5-10} and Stickelberger's theorem to show that they 
come from algebraic Hecke characters. 





\subsection{}\label{VI:6-1}

If $E$ is a field with algebraic closure $\bar E$, and $\lambda$ is 
a character, of order dividing $n$, of $\mu_d(\bar E)$, we denote by $_n\lambda$ the 
character of $\mu_n(\bar E)$ such that $\lambda$ and $_n \lambda$ induce the 
same character of $\widehat\dZ(1)_{\bar E} = \varinjlim \mu_n(\bar E)$. The 
same notation is used when $\lambda$ is a character of $\widehat\dZ(1)_{\bar E}$.

Let $F$ and $k$ be two number fields, $\bar F$ an algebraic closure of 
$F$, $I$ a finite set equipped with an action of $\gal(\bar F/F)$ and 
$\lambda=(\lambda_i)_{i\in I}$ a family of characters 
$\lambda_i:\widehat\dZ(1)_{\bar F} \to k^\times$. It is assumed that 
$\lambda_i\ne 1$ ($i\in I$), that the product of $\lambda_i$ is trivial and that 
the family $\lambda$ is defined over $F$, i.e. that for every 
$\sigma\in \gal(\bar F/F)$, we have 
$\lambda_{\sigma(i)}=\lambda_i\circ \sigma^{-1}$. 

The Galois set $I$ is the set of homomorphisms into $\bar F$ of an 
algebra $E$ separable over $F$. If $C$ is the set of orbits of 
$\gal(\bar F/F)$ in $I$, we have a decomposition of $E$ as a product of fields 
$E=\prod_{\alpha\in C} E_\alpha$, where $\alpha$ is identified with the set of 
embeddings of $E_\alpha$ in $\bar F$. If the characters $\lambda_i$, for 
$i\in \alpha$, are of order $d_\alpha$, the field $E_\alpha$ contains the 
$d_\alpha$-th roots of $1$ and there exists 
$\lambda_\alpha:\mu_{d_\alpha}(E_\alpha) \hookrightarrow k^\times$ such that, if 
$i\in \alpha$ corresponds to the embedding $\sigma_i$ of $E_\alpha$ in $\bar F$, 
then $_{d_\alpha} \lambda_i = \lambda_\alpha \sigma_i^{-1}$. 





\subsection{}\label{VI:6-2}

Let $n>0$ be a multiple of the $d_\alpha$, let $\Sigma$ be the set of places of $F$ 
at which $E/F$ ramifies or which divide $n$, let $\mu$ be a place of $k$ of 
residue characteristic $\ell$, and let $A$ be the ring of elements of $F$ 
integral outside $\Sigma$ and $\ell$. The integral closure $M$ of $A$ 
is finite étale of rank $N=|I|$ over $A$. With the notation of \ref{VI:4-6}, 
\ref{VI:4-8}, $_n \lambda_i$ define 
\[
  _n\lambda:(V^\ast/\dG_m)_n \to k^\times \subset k_\mu^\times \text{.} 
\]

Let $a$ be the projection of $W^\ast/\dG_m$ onto $\spec(A)$. By 
\ref{VI:4-20}, the sheaf $\eR^{N-2} a_! \sK_n(_n \lambda)$ is smooth of rank 
one; it defines an $\ell$-adic representation 
\[
  j_\mu[\lambda]:\gal(\bar F/F)^\text{ab} \to k_\mu^\times 
\]
  (or simply $j_\mu$) unramified outside $\Sigma$ and $\ell$. 





\subsection{}\label{VI:6-3}

Let us compute $j_\mu(F_\fp)$. By \ref{VI:4-20}, it suffices to do so 
after reduction modulo $\fp$. Put $f=A/\fp$, $e=M/\fp  M$, and let 
$\bar f$ be the algebraic closure of $f$ determined by a place of $\bar F$ 
above $\fp$. We still have $I=\hom_f(e,\bar f)$. After reduction, 
$_n\lambda$ still has a description as in \ref{VI:6-1}: 
\begin{enumerate}[a)]
  \item Let $D$ be the set of orbits of $\gal(\bar f/f)$ in $I$. The 
    decomposition of $e$ as a product of fields has the form 
    $e=\prod_{\beta\in D} e_\beta$, where $\sigma_i$ ($i\in \beta$) 
    are identified with the embeddings of $e_\beta$ into $\bar f$. For $\beta\in D$, 
    denote by $\alpha(\beta)$ the element of $C$ containing $\beta$, and put 
    $d_\beta=d_{\alpha(\beta)}$. Let 
    \[\xymatrix{
      \lambda_\beta : \mu_{d_\beta}(e_\beta) 
        & \ar[l]_-\sim \mu_{d_\beta}(E_{\alpha(\beta)}) \ar[r]^-{\lambda_{\alpha(\beta)}} 
        & k^\times \text{.} 
    }\]
  \item If $i\in \beta$ corresponds to the embedding $\sigma_i$ of $e_\beta$ into 
    $\bar f$, the reduction 
    $_{d_\beta} \bar\lambda_i:\mu_{d_\beta}(\bar f) \to k^\times$ 
    of $_{d_\beta} \lambda_i$ is 
    $\lambda_\beta \circ \sigma_i^{-1}$. 
\end{enumerate}

Let $q_\beta$ be the number of elements of $e_\beta$, 
  $\chi_\beta={}_{q_\beta-1}\lambda_\beta$, and 
$\chi:e^\times \to k^\times$ the character with coordinates $\chi_\beta$. On the 
affine space over $A/\fp$ defined by $e$ (\ref{VI:4-6}), or rather on the torus 
$e_\beta^\times$, we have isomorphisms 
  $\sK_n((_n\bar\lambda_i)_{i\in\beta}) = \sK_{d_\beta}((\lambda_\beta \circ \sigma_i^{-1})_{i\in\beta}) =\sK_{q_\beta-1}((\chi_\beta\circ \sigma_i^{-1})_{i\in\beta}) = \sF(\chi_\beta)$ (\ref{VI:4-9}). 
  On $e^\times$, we therefore have $\sK_n((_n\bar\lambda_i)_{i\in I}) = \sF(\chi)$, 
and this isomorphism descends to the reduction modulo $\fp$ of $V^\ast/\dG_m$. 
It follows that





\begin{proposition_}\label{VI:6-4}
With the notations of \ref{VI:6-2} and \ref{VI:6-3}, $j_\mu(F_\fp)$ is the 
Jacobi sum $J(\chi)$. 
\end{proposition_}

Let us rewrite the formulas linking $(\lambda_i)_{i\in I}$ to 
  $(\chi_\beta)_{\beta\in D}$: for $i\in \beta\subset \alpha$, 
\begin{align*}
  _{d_\alpha} \lambda_i &= \lambda_\alpha\circ \sigma_i^{-1} && \lambda_\alpha:\mu_{d_\alpha}(E_\alpha) \to k^\times \\
  \lambda_\beta &= \text{reduction of $\lambda_\alpha$} && :\mu_{d_\alpha}(e_\beta) \to k^\times \\
  \chi_\beta &= {}_{q_\beta-1}\lambda_\beta: x\mapsto \lambda_\beta(x^{(q_\beta-1)/d_\alpha}) && :e_\beta^\times \to k^\times 
\end{align*}





\begin{theorem_}\label{VI:6-5}
The representation $j_\mu$ is defined by an algebraic Hecke character 
of $F$ with values in $k^\times$. 
\end{theorem_}

We shall verify this using \ref{VI:5-10}. Here we follow Weil's 
proof \cite{we52,we74-2}, using Stickelberger's theorem, 
and I shall give only a few indications. 





\subsection{}\label{VI:6-6}

Given $j:\gal(\bar F/F) \to k_\lambda^\times$ and a finite extension 
$F'/F$, one verifies using \ref{VI:5-10} that $j$ is defined by a 
Hecke character if and only if $j|\gal(\bar F/F')$ is. This 
remark allows us, in \ref{VI:6-5}, to reduce to the case where Galois acts 
trivially on $I$ and where $F$ contains the $n$-th roots of 
unity. We may then reduce to the case $F=\dQ(\sqrt[n]{1})$ and 
take $k=F$. In this case, each $_n \lambda_i$ is defined by 
$a_i\in \dZ/n$: it is $x\mapsto x^{a_i}:\mu_n \to k^\times$. We have 
$\sum a_i=0$, and $a_i\ne 0$. 





\subsection{}\label{VI:6-7}

In a number field, the places of degree greater than $1$ always form a 
set of density $0$. By applying \ref{VI:5-10}, we may therefore disregard 
them and use Stickelberger only for a prime field. Let 
$p$ be prime, let $0<a<p-1$, and let $\widetilde x:\dF_p^\times \to \dQ_p^\times$ be the 
``multiplicative lift'' character; put 
  $g=-\sum_{x\in \dF_p^\times} \widetilde x^{-a} \zeta^x\in \dQ_p(\zeta)$, for 
$\zeta$ a $p$-th root of $1$. Again denote by $x$ the integer in $[0,p-1]$ 
  lifting $x$, and expand $\zeta^x=(1+\pi)^x$ by the binomial formula. 
  We find $g=-\sum_{i=0}^{p-1} A_i\pi^i$ with $A_i\in \dZ_p$, whose reduction 
$\mod p$ is the sum $\sum_{x\in \dF_p^\times} x^{-a} \binom x i$. For $b$ not 
divisible by $p-1$, we have $\sum_{x\in \dF_p^\times} x^b=0$. If $i<a$, 
$\binom x i$ is a polynomial in $x$ of degree $i<a$, whence 
$A_i\equiv 0\pmod p$. Likewise, $A_a\equiv \frac{p-1}{a!} \pmod p$, and 
$g\sim \frac{\pi^a}{a!}$: the valuation of $g$ is given by 
$v(g)/v(p) = \frac{a}{p-1}$. 





\subsection{}\label{VI:6-8}

Let $\sigma_i$ ($i\in (\dZ/n)^\times$) be the element of 
$\gal(\dQ(\sqrt[n]{1})/\dQ)$ such that $\sigma_i(\zeta)=\zeta^i$ for $\zeta^n=1$, 
and write the group of algebraic homomorphisms 
$\dQ(\sqrt[n]{1})^\times \to \dQ(\sqrt[n]{1})^\times$ additively. For $\lambda$ as in 
\ref{VI:6-6}, defined by a family 
$(a_i)_{i\in I}$ of nonzero integers $\mod n$ whose sum is zero, 
\ref{VI:5-10}, \eqref{VI:eq:4-15-1} and \ref{VI:6-7} 
show that $j_\mu[\lambda]$ is defined by an algebraic Hecke character 
$j[a]$ and that, if $N$ denotes the norm, its algebraic part is 
given by 
\begin{equation*}\tag{6.8.1}\label{VI:eq:6-8-1}
  (N j[a])_\text{alg} = \sum_{i\in (\dZ/n)^\times} \sum_{j\in I} \left\langle \frac{i a_j}{n} \right\rangle \sigma_i^{-1} 
\end{equation*}
  where $\langle \cdot \rangle$ is the fractional part. 





\subsection{}\label{VI:6-9}

The most interesting case is that in which $F\subset \dQ(\sqrt[n]{1})$ and 
the family of $_n \lambda_i$ is constructed as follows: 
\begin{enumerate}[a)]
  \item Take a finite family $(A_j)_{j\in J}$ of orbits of 
    $\gal(\dQ(\sqrt[n]{1})/F)=H\subset (\dZ/n)^\times$ in $\dZ/n$. Assume that 
    $A_j\ne\{0\}$ and that $\sum_{j\in J} \sum_{a\in A_j} a=0$. 
  \item Take $I=\sqcup_{j\in J} A_j$; Galois acts on each $A_j$. 
  \item Take $\bar F\supset \dQ(\sqrt[n]{1})$, $k=\dQ(\sqrt[n]{1})$ and, if 
    $i\in I$ corresponds to $a\in A_j\subset (\dZ/n)^\times$, 
    $_n\lambda_i(x)=x^a$. 
\end{enumerate}

Using \ref{VI:4-20}, one verifies that a general $j_\mu$ (relative to 
$F_1$) is obtained from a $j_\mu$ as above (for 
$F\subset \dQ(\sqrt[n]{1})$, with $F_1$ an extension of $F$) by restriction to 
$\gal(\bar F/F_1)$ and multiplication by a character of order two. 

For $a\in \gal(k/\dQ)=(\dZ/n)^\times$, $a$ transforms the family of $A_j$ 
into the family of $a A_j$. If $a\in H$, the family of $A_j$ is therefore preserved, 
and the $j_\mu(F_\fp)$ are invariant under $H$: 





\begin{proposition_}\label{VI:6-10}
Under the assumptions of \ref{VI:6-9}, $j_\mu$ is defined by a character 
that is an algebraic Hecke character of $F$ with values in $F$. 
\end{proposition_}

Let us again take $F=\dQ(\sqrt[n]{1})$ and consider the Hecke characters 
$j[a]$ defined by a family $(a_i)_{i\in I}$, $i\in \dZ/n$, 
$a_i\ne 0$, $\sum a_i=0$. The following theorem was obtained 
independently by Vishik. 





\begin{theorem_}\label{VI:6-11}
Every algebraic Hecke character of $\dQ(\sqrt[n]{1})$ has a power 
in the group generated by the $j[a]$ and the norm. 
\end{theorem_}

Since we work ``up to torsion,'' only the algebraic part of the 
Hecke characters under consideration matters. Applying \eqref{VI:eq:6-8-1} to 
$j[a]$, with $a=$ an element of $\dZ/n$ repeated $n$ times, we 
are reduced to the following lemma.





\begin{lemma_}\label{VI:6-12}
Write $\sum x_j \sigma_j$ for the elements of the group algebra 
  $\dQ[(\dZ/n)^\times]$, and let $X$ be the subgroup consisting of those for which 
$x_j+x_{-j}=C^{\underline{t e}}$. Then $X$ is generated over $\dQ$ by 
$N=\sum \sigma_j$ and by the elements 
  $g_a=\sum \left\langle \frac{i a}{n}\right\rangle \sigma_i^{-1}$ 
($a\in \dZ/n$, $a\ne 0$). 
\end{lemma_}

Both $X$ and $Y=\langle N,(g_a)\rangle$ are ideals of the group 
algebra, and it suffices to check that they are annihilated by the same 
complex characters $\chi$ of $(\dZ/n)^\times$. This amounts to saying that a 
odd character $(\chi(j)=-\chi(-j)$) never annihilates all the $g_a$. If 
$\chi$ is primitive, we have 
\[
  \chi(g_1) = \sum \frac i n \chi^{-1}(i) = -L(\chi,0)\ne 0 \text{.} 
\]
In the general case, if $\chi$ comes from a primitive character of 
  $(\dZ/d)^\times$, $\chi(g_{n/d})$ reduces to an analogous sum over 
  $(\dZ/d)^\times$, hence the lemma. 










\section{Generalized Kloosterman sums}\label{VI:7}





\subsection{}\label{VI:7-1}

Let $\dF_q$ be a finite field with $q$ elements, and let $\psi$ be a nontrivial 
additive character. In this section, we undertake a cohomological study of 
generalized Kloosterman sums 
\begin{equation*}\tag{7.1.1}\label{VI:eq:7-1-1}
  \kloos_{n,a} = \sum_{x_1\dotsm x_n = a} \psi(x_1+\cdots + x_n) \qquad (n\geqslant 1) \text{.} 
\end{equation*}
For $a=0$, this sum is elementary (see \ref{VI:7-7}): 
\begin{equation*}\tag{7.1.2}\label{VI:eq:7-1-2}
  \kloos_{n,0} = (-1)^{n-1} \text{.} 
\end{equation*}
The interesting case is $a\ne 0$. The $x_i$ then belong to 
$\dF_q^\times$. In this case we shall prove the bound 
\begin{equation*}\tag{7.1.3}\label{VI:eq:7-1-3}
  \left|\kloos_{n,a}\right| \leqslant n q^{\frac{n-1}{2}} \text{.} 
\end{equation*}

Our essential tools will be the cohomological counterparts of the 
following identities: 
\begin{align*}\tag{7.1.4}\label{VI:eq:7-1-4}
  \kloos_{n,a} &= \sum_{x_1} \psi(x_1) \sum_{x_2\dotsm x_n = a/x_1} \psi(x_2+\cdots + x_n) \\ \notag
    &= \sum_{x\in \dF_q^\times} \psi(x) \kloos_{n-1,a/x} \qquad (a\ne 0,n\geqslant 2) \\ \tag{7.1.5}\label{VI:eq:7-1-5}
  \sum_a \kloos_{n,a} &= \sum_a \sum_{x_1 \dotsm x_n = a} \psi\left(\sum x_i\right) = \sum_{x_1,\dots,x_n} \psi\left(\sum x_i\right) = 0 
\end{align*}
For $\chi$ a nontrivial character of $\dF_q^\times$, 
\begin{equation*}\tag{7.1.6}\label{VI:eq:7-1-6}
  \sum \chi(a) \kloos_{n,a} = \sum_{x_1,\dots,x_n} \chi\left(\prod x_i\right) \cdot \psi\left(\sum x_i\right) = (-\tau(\chi^{-1},\psi))^n 
\end{equation*}
(Gauss sum). 





\subsection{}\label{VI:7-2}

Let $k$ be an étale algebra of degree $n$ over $\dF_q$, and let $\varepsilon(k)$ 
be as in \eqref{VI:eq:4-5-1}. Put 
\[
  \kloos_{k,a} = \sum_{N_{k/\dF_q}(x)=a} \psi \tr(x) \text{.} 
\]
We likewise have 
\begin{align*}\tag{7.2.1}\label{VI:eq:7-2-1}
  \kloos_{k,0} &= (-1)^{n-1} \cdot \varepsilon(k) \\ \tag{7.2.2}\label{VI:eq:7-2-2}
  \sum_a \kloos_{k,a} &= 0 \\ \tag{7.2.3}\label{VI:eq:7-2-3}
  \sum_a \chi(a) \kloos_{k,a} &= (-1)^n \tau_{\dF_q}\left(\chi^{-1}\circ N_{k/\dF_q},\psi\right) \text{.} 
\end{align*}

By Fourier inversion on $\dF_q^\times$, we deduce an expression of 
$\kloos_{k,a}$ in terms of Gauss sums: 
\begin{align*} 
  \kloos_{k,a} &= \frac{1}{q-1} \sum_\chi \chi(a) \sum_{b\in \dF_q^\times} \chi^{-1}(b) \kloos_{k,b} \\ \tag{7.2.4}\label{VI:eq:7-2-4}
  (-1)^n \kloos_{k,a} &= \frac{1}{q-1} \left(\varepsilon(k) + \sum_{\chi\ne 1} \chi(a) \tau_{\dF_q}\left(\chi\circ N_{k/\dF_q},\psi\right) \right) \text{.} 
\end{align*}

The Hasse--Davenport identity now makes it possible to compare $\kloos_{k,a}$ 
with $\kloos_{n,a} = \kloos_{\dF_{q^n},a}$: 
\begin{equation*}\tag{7.2.5}\label{VI:eq:7-2-5}
  \kloos_{k,a} = \varepsilon(k) \kloos_{n,a} \text{.} 
\end{equation*}

I know of no cohomological interpretation of \eqref{VI:eq:7-2-4}, but 
I shall give one for \eqref{VI:eq:7-2-5}. Let us return to the sums $\kloos_{n,a}$. 





\subsection{}\label{VI:7-3}

Let $\dF$ be an algebraic closure of $\dF_q$ and, for $a\in \dF$, let 
$V_a^{n-1}$, or simply $V_a$, be the hypersurface in $\dA^n$ with equation 
$x_1 \dotsm x_n=a$. As usual, we shall regard $\psi$ as taking 
$\ell$-adic rather than complex values. With the notation of 
\ref{VI:1-8}(ii), \eqref{VI:eq:7-1-3} follows from 





\begin{theorem_}\label{VI:7-4}
The cohomology of $V_a^{n-1}$ with coefficients in $\sF(\psi(\sum x_i))$ satisfies 
\begin{enumerate}[\indent (i)]
  \item $\h_c^i=0$ for $i\ne n-1$; 
  \item $\h_c^\bullet{} \iso \h^\bullet{}$; 
  \item for $a\ne 0$, $\dim \h_c^{n-1}{}=n$; 
  \item for $a=0$, $\h_c^{n-1}{}$ is canonically isomorphic to $E_\lambda$. 
\end{enumerate}
\end{theorem_}





\subsection{}\label{VI:7-5}

As usual, this theorem also controls the dependence of 
$\kloos_{n,a}$ on $q$: for each $a\in \dF_q^\times$, there are $n$ Frobenius 
eigenvalues $\alpha_1,\dots,\alpha_n$, of complex absolute value 
$q^{\frac{n-1}{2}}$, such that 
\[
  \sum_{\substack{x_1\dotsm x_n = a \\ x_i\in \dF_{q^m}}} \psi\circ \tr\left(\sum x_i\right) = (-1)^{n-1} \left(\alpha_1^m + \cdots + \alpha_n^m \right) \text{.} 
\]

For even $n$, $x\mapsto -x$ is an involution of $V_a$, which transforms 
$\sF(\psi)$ into its dual. Arguing as in \ref{VI:3-6}, we deduce a 
canonical nondegenerate alternating form on $\h_c^{n-1}(V_a,\sF(\psi))$, 
with values in $\dQ_\ell(-(n-1))$, and hence the following. 





\begin{corollary_}\label{VI:7-6}
With the notations of \ref{VI:7-5}, if $n$ is even, the $\alpha_i$ group together 
into $\frac n 2$ pairs of roots $\alpha,\bar\alpha$ whose product is 
$q^{n-1}$. 
\end{corollary_}





\subsection{}\label{VI:7-7}

Let us verify the case $a=0$ of \ref{VI:7-5}. The hypersurface $V_0$ is the union 
of the coordinate hyperplanes $x_i=0$. Let us compute the Leray spectral sequence 
\eqref{VI:eq:2-6-1*} for this covering of $V_0$, for cohomology with or 
without supports and with coefficients in $\sF(\psi)$. By \ref{VI:2-7*}*, all 
the initial terms vanish except the cohomology of the intersection $\{0\}$ of 
all the coordinate hyperplanes. What remains is an isomorphism 
$\h_c^0{}\iso \h^0{}=E_\lambda$ which justifies \ref{VI:7-4}. 

We shall prove \ref{VI:7-4} and \ref{VI:7-8} below by simultaneous induction 
on $n$, using the spectral sequence of which 
\eqref{VI:eq:7-1-4} is the counterpart. This requires control, as $a$ varies, 
of $\kloos_{n,a}$. 

Let $\pi:\dA^n \to \dA$ be the map ``product of the coordinates,'' and let 
$\sigma$ be the sum of the coordinates. The hypersurfaces $V_a$ are the fibers of 
$\pi$.





\begin{theorem_}\label{VI:7-8}
\begin{enumerate}[(i)]
  \item The sheaf $\eR^{n-1}\pi_! \sF(\psi\sigma)$ is smooth of rank $n$ on 
    $\dA^1\setminus \{0\}$. 
  \item Its extension by $0$ to $\dP^1$ is the direct image of its 
    restriction to $\dA^1\setminus \{0\}$. 
  \item At $0$, the monodromy is unipotent, with a single Jordan block.
  \item At $\infty$, wild inertia acts without fixed points $\ne 0$, and the 
    Swan conductor is $1$. 
  \item We have $\eR^i \pi_! \sF(\psi\sigma) \iso \eR^i \pi_\ast \sF(\psi\sigma)$ 
    (zero for $i\ne n-1$).
\end{enumerate}
\end{theorem_}
\begin{proof}[Proof]
that \ref{VI:7-4} (for a given value of $n$) implies \ref{VI:7-8} 
(for the same value of $n$). 

The fiber of $\eR^i \pi_! \sF(\psi\sigma)$ at $a$ is 
$\h_c^i(V_a,\sF(\psi\sigma))$, and on the dense open subset of $\dA^1$ where the 
  $\eR^i \pi_!\sF(\psi\sigma)$ are locally constant, the fiber of 
$\eR^i \pi_\ast \sF(\psi\sigma)$ is $\h^i(V_a,\sF(\psi\sigma))$ 
(\hyperref[VII]{Finiteness\ Theorem} \ref{VII:2-1}). The hypothesis \ref{VI:7-4}($n$) 
therefore gives us the following. 





\begin{lemma_}\label{VI:7-9}
\begin{enumerate}[(i)]
  \item $\eR^i \pi_! \sF(\psi\sigma)=0$ for $i\ne n-1$. 
  \item For $i=n-1$, the fiber of this sheaf at any point $a\ne 0$ has 
    constant rank $n$. At $0$, it has rank $1$. 
  \item On a dense open $U$, we have 
    $\eR^i \pi_! \sF(\psi\sigma) \iso \eR^i \pi_\ast \sF(\psi\sigma)$. 
\end{enumerate}
\end{lemma_}





\subsection{}\label{VI:7-10}

We shall now use the Leray spectral sequence of $\pi$ for 
cohomology with or without compact supports, with coefficients in $\sF(\psi\sigma)$. 
The abutment is known: 
$\h_c^\bullet(\dA^n,\sF(\psi\sigma))=\h^\bullet(\dA^n,\sF(\psi\sigma)) = 0$ 
(\ref{VI:2-7*}*). For cohomology with compact supports, \ref{VI:7-9}(i) ensures that 
$E_2^{pq}=0$ for $q\ne n-1$. The terms $E_2$ are therefore all zero: 
\begin{equation*}\tag{7.10.1}\label{VI:eq:7-10-1}
  \h_c^p(\dA^1,\eR^{n-1} \pi_!\sF(\psi\sigma)) = 0 \text{.} 
\end{equation*}

  For $p=0$, this means that $\eR^{n-1}\pi_! \sF(\psi\sigma)$ has no sections 
with point support. In view of the constancy of rank in \ref{VI:7-9}(ii), we obtain 
\begin{equation*}\tag{7.10.2}\label{VI:eq:7-10-2}
  \begin{array}{c}
    \text{The sheaf $\eR^{n-1} \pi_! \sF(\psi\sigma)$ on $\dA^1$ is smooth on $\dA^1\setminus \{0\}$, and is } \\ 
    \text{a subsheaf of the direct image $\sG$ of its restriction to $\dA^1\setminus \{0\}$.}
  \end{array}
\end{equation*}

The long exact cohomology sequence of 
\[\xymatrix{
  0 \ar[r] 
    & \eR^{n-1} \pi_! \sF(\psi\sigma) \ar[r] 
    & \sG \ar[r] 
    & \sQ \ar[r] 
    & 0 \text{.} 
}\]
where the support of $\sQ$ is concentrated at $0$, gives 
\[\xymatrix{
  \h_c^0(\dA^1,\sG) \ar[r] 
    & \sQ_0 \ar[r] 
    & \h_c^1(\dA^1,\eR^{n-1} \pi_! \sF(\psi\sigma)) \text{.} 
}\]
Since the outer terms vanish, $\sQ=0$ and 
\begin{equation*}\tag{7.10.3}\label{VI:eq:7-10-3}
  \text{$\eR^{n-1} \pi_! \sF(\psi\sigma)$ is the direct image of its restriction to $\dA^1\setminus \{0\}$.}
\end{equation*}

Let $j$ be the inclusion of $\dA^1$ into $\dP^1$. The ``derived-category 
version'' of the Leray spectral sequences is 
\begin{align*}
  \hh^\bullet(\dP^1,j_! \eR\pi_! \sF(\psi\sigma)) &= \hh_c^\bullet(\dA^1,\eR \pi_! \sF(\psi\sigma)) = \h_c^\bullet(\dA^n,\sF(\psi\sigma)) = 0 \\
  \hh^\bullet(\dP^1,\eR j_\ast \eR \pi_\ast \sF(\psi\sigma)) &= \hh^\bullet(\dA^1,\eR \pi_\ast \sF(\psi\sigma)) = \h^\bullet(\dA^n,\sF(\psi\sigma)) = 0 \text{.}
\end{align*}
Let $\Delta$ be the mapping cylinder of 
$j_! \eR\pi_! \sF(\psi\sigma) \to \eR j_\ast \eR \pi_\ast \sF(\psi\sigma)$. 
According to \ref{VI:7-9}(iii), the cohomology sheaves of this complex of 
sheaves have finite support. On the other hand, the long exact sequence in 
cohomology of the triangle 
$(j_!\eR \pi_! \sF(\psi\sigma),\eR j_\ast \eR \pi_\ast \sF(\psi\sigma),\Delta)$ 
shows that $\hh^\bullet(\dP^1,\Delta)=0$. We have 
$\hh^\bullet(\dP^1,\Delta) = \h^0(\dP^1,\sH^\bullet(\Delta))$, and finally 
$\Delta=0$: 
\[\xymatrix{
  j_! \eR\pi_! \sF(\psi\sigma) \ar[r]^-\sim 
    & \eR j_\ast \eR \pi_\ast \sF(\psi\sigma) \text{.} 
}\]

In particular, $\eR\pi_! \sF(\psi\sigma)\iso \eR\pi_\ast \sF(\psi\sigma)$ and, 
since these complexes of sheaves have only one nonzero cohomology sheaf, 
\begin{equation*}\tag{7.10.4}\label{VI:eq:7-10-4}
\xymatrix{
  j_! \eR \pi_! \sF(\psi\sigma) \ar[r]^-\sim 
    & j_\ast \eR \pi_! \sF(\psi\sigma) \text{.} 
}
\end{equation*}
This completes the proof of (i)(ii)(v). 





\subsection{}\label{VI:7-11}

Let $\swan_0$ and $\swan_\infty$ be the Swan conductors at $0$ and at 
$\infty$ of $\eR^{n-1} \pi_! \sF(\psi\sigma)$. We have 
$\chi_c(\eR^{n-1} \pi_! \sF(\psi\sigma))=0$ \eqref{VI:eq:7-10-1} and the formula 
of Euler--Poincaré \eqref{VI:eq:3-2-1} reduces to 
\begin{equation*}\tag{7.11.1}\label{VI:eq:7-11-1}
  \swan_0 + \swan_\infty =1 \text{:} 
\end{equation*}
one of these conductors is $0$ (tame ramification), and the other is $1$. 





\subsection{}\label{VI:7-12}

Let $\sG$ be a sheaf of rank $1$ on $\dG_m$, of Kummer type 
\ref{VI:4-7} ($\sG=\sK_n(\chi)$) and nonconstant ($\chi\ne 1$). The sheaf 
$\sF(\psi\sigma)\otimes \pi^\ast \sG$ on 
$V^\ast=\pi^{-1}(\dG_m) = \dA^n\setminus \pi^{-1}(0)$ is then of the type 
encountered in the study of Gauss sums, and \ref{VI:4-11}(ii) gives us 
\begin{equation*}\tag{7.12.1}\label{VI:eq:7-12-1}
\xymatrix{
  \h_c^\bullet(V^\ast,\pi^\ast\sG\otimes \sF(\psi\sigma)) \ar[r]^-\sim 
    & \h^\bullet(V^\ast,\pi^\ast \sG\otimes \sF(\psi\sigma)) \text{.} 
}
\end{equation*}
The Leray spectral sequence of $\pi$ transforms this statement into 
\begin{equation*}\tag{7.12.2}\label{VI:eq:7-12-2}
\xymatrix{
  \h_c^\bullet(\dG_m,\sG\otimes \eR^{n-1} \pi_! \sF(\psi\sigma)) \ar[r]^-\sim 
    & \h^\bullet(\dG_m,\sG\otimes \eR^{n-1} \pi_\ast \sF(\psi\sigma)) \text{.} 
}
\end{equation*}

Let $i$ be the inclusion of $\dG_m$ into $\dP^1$, and let $\Delta$ be the mapping cylinder 
of 
  $\delta:i_!(\sG\otimes \eR^{n-1}\pi_! \sF(\psi\sigma))\to \eR i_\ast (\sG\otimes \eR^{n-1} \pi_! \sF(\psi\sigma))$
(where $\eR^{n-1} \pi_! = \eR^{n-1}\pi_\ast$). Reasoning as in 
\ref{VI:7-10}, we deduce from \eqref{VI:eq:7-12-2} that $\Delta=0$. In 
particular 
\begin{equation*}\tag{7.12.3}\label{VI:eq:7-12-3}
  \text{At $0$ and at $\infty$, inertia acts without fixed points $\ne 0$ on $\sG\otimes \eR^{n-1}\pi_! \sF(\psi\sigma)$.}
\end{equation*}

At $\infty$, this holds even if $\sG$ is trivial \eqref{VI:eq:7-10-4}, 
so wild inertia acts without a fixed point. In particular, the 
ramification is wild: by \ref{VI:7-11}, we have $\swan_\infty=1$, and the 
ramification at $0$ is tame. 

At $0$, \eqref{VI:eq:7-12-3} forces the ramification -- which is tame -- of 
$\eR^{n-1} \pi_! \sF(\psi\sigma)$ to be unipotent. By 
\eqref{VI:eq:7-10-3} and \ref{VI:7-9}(ii), there is a single Jordan block. This 
completes the proof of \ref{VI:7-8}($n$). 





\subsection{}\label{VI:7-13}

For $n=1$, \ref{VI:7-4} is trivial. To complete the proof of \ref{VI:7-4} and 
\ref{VI:7-8}, we still have to prove that \ref{VI:7-8}, for a given value 
of $n$, implies \ref{VI:7-4} for $n+1$. The case $a=0$ having been 
  treated (\ref{VI:7-7}), we assume that $a\ne 0$. We shall write $x_0,\dots,x_n$ for 
the coordinates of $\dA^{n+1}$. Let $g=x_0:V_a\to \dG_m$, $\tau$ 
the involution $x\mapsto a x^{-1}$ of $\dG_m$, and let $\pi$, as above, be the 
product of the coordinates: $\dA^{n+1} \to \dA^1$, with $\sigma$ their sum. We 
shall denote by $\sF(\psi\sigma)$ the sheaf $\sF(\psi(\sum x_i))$ both on $\dA^{n+1}$ 
and on $\dA^n$. 





\subsection{}\label{VI:7-14}

Let us write down the Leray spectral sequence for $g$. Since the sheaf $\sF(\psi\sigma)$ on 
$\dA^{n+1}$ is the external tensor product of the sheaves $\sF(\psi)$ on 
$\dA^1$ and $\sF(\psi\sigma)$ on $\dA^n$, we find (cf. \eqref{VI:eq:7-1-4}). 
\begin{align*}
  'E_2^{p q} &= \h_c^p(\dG_m,\sF(\psi)\otimes \tau^\ast \eR^q \pi_! \sF(\psi\sigma)) \Rightarrow \h_c^{p+q}(V_a,\sF(\psi\sigma)) \\
  ''E_2^{p q} &= \h^p(\dG_m,\sF(\psi)\otimes \tau^\ast \eR^q \pi_\ast \sF(\psi\sigma)) \Rightarrow \h^{p+q}(V_a,\sF(\psi\sigma)) \text{.} 
\end{align*}
By hypothesis, 
$\eR^q \pi_! \sF(\psi\sigma) \iso \eR^q \pi_\ast \sF(\psi\sigma)$, and this 
sheaf is zero for $q\ne n-1$. The sheaf $\sF(\psi)$ on $\dG_m$ is 
wildly ramified at $\infty$ and unramified at $0$, while 
$\tau^\ast \eR^q \pi_! \sF(\psi\sigma)$ is wildly ramified at $0$ (without 
invariants under wild inertia, and with Swan conductor $1$) and tamely 
ramified at $\infty$. Their tensor product is therefore 
\begin{enumerate}[\indent a)]
  \item wildly ramified at $0$ and $\infty$, without invariants under wild 
    inertia; 
  \item of Swan conductor $1$ at $0$ and $n$ (the rank of 
    $\eR^q \pi_! \sF(\psi\sigma))$ at $\infty$. 
\end{enumerate}
We have $'E_2^{p q} \iso ''E_2^{p q}$, and $'E_2^{p q}=0$ except for $p=1$ and 
$q=n-1$. The Euler-Poincaré formula finally gives 
\[
  \dim{'E_2^{1,n-1}} = 1+n \text{;} 
\]
this completes the proof. 
\end{proof}





\subsection{Remark}\label{VI:7-15}

The isomorphism obtained in \ref{VI:7-14} 
\begin{equation*}\tag{7.15.1}\label{VI:eq:7-15-1}
  \h_c^n(V_a^n,\sF(\psi\sigma)) = \h_c^1(\dG_m,\sF(\psi)\otimes \tau^\ast \eR^{n-1} \pi_! \sF(\psi\sigma)) 
\end{equation*}
allows us, with the notation of \ref{VI:7-5}, to calculate the product $\det(F)$ 
of the Frobenius eigenvalues $\alpha_i$. The formalism of \cite{de73} 
applies, since $\sF(\psi)$ and $\eR^{n-1} \pi_! \sF(\psi\sigma)$ 
belong to infinite compatible systems of $\ell$-adic 
representations. This reduces us to local problems, which can be 
simplified by noting that (cf. \cite[9.5]{de73}) 
\begin{enumerate}[\indent a)]
  \item the global constants for 
    $\sF(\psi)$ and $\tau^\ast \eR^{n-1} \pi_! \sF(\psi\sigma)$ are easy to 
    calculated easily, since the cohomology of these sheaves is elementary. 
  \item at every point, for either of these sheaves, the 
    corresponding $\ell$-adic representation of the decomposition group 
    has unramified semisimplification. 
\end{enumerate}

The result, for the sum with $n$ variables, is that 
\begin{equation*}\tag{7.15.2}\label{VI:eq:7-15-2}
  \det\left(F,\h_c^{n-1}(V_a^{n-1},\sF(\psi))\right) = q^{\frac{n(n-1)}{2}} \text{.} 
\end{equation*}





\subsection{Remark}\label{VI:7-16}

For $p=2$, the sheaf $\sf(\psi)$ is orthogonal. For $a\ne 0$ and $n$ 
odd, we deduce from \ref{VI:7-4}(ii) and Poincaré duality a 
symmetric bilinear form on $\h_c^{n-1}(V_a^{n-1},\sF(\psi))$ with 
values in $\bar\dQ_\ell(1-n)$. With respect to this form, $F$ is an 
orthogonal similitude of multiplier $q^{n-1}$, and $q^{\frac{1-n}{2}} F$ 
belongs to the special orthogonal group, by \eqref{VI:eq:7-15-2}. One 
of the eigenvalues of $F$ is therefore $q^{\frac{n-1}{2}}$, and the others 
occur in pairs of roots $\alpha$, $\bar\alpha$ whose product is $q^{n-1}$. 

For $n=3$, L.\ Carlitz \cite{ca69-2} obtained a more precise result, 
equivalent to the following proposition. 





\begin{proposition_}\label{VI:7-17}
For $p=2$, and $\psi(x)=(-1)^{\tr_{\dF_q/\dF_2}(x)}$, we have an isomorphism 
\[
  \h_c^2(V_a^2,\sF(\psi\sigma)) = \operatorname{Sym}^2 \h_c^1(V_a^1,\sF(\psi\sigma)) \text{.} 
\]
\end{proposition_}

The second tensor power of $\h_c^1(V_a^1,\sF(\psi\sigma))$ is 
$\h_c^2(V_a^1\times V_a^1,\sF(\psi\sigma)\boxtimes \sF(\psi\sigma))$ 
(K\"unneth), with the symmetry $x\otimes y\mapsto y\otimes x$ represented 
geometrically by $-\tau^\ast$, where $\tau$ is the automorphism 
$(x,y)\mapsto (y,x)$ of $V_a\times V_a$. If $X'$ is the quotient of 
$V_a^1\times V_a^1$ by $\{\operatorname{id}{},\tau\}$ and $\pi$ is the 
projection of $V_a^1\times V_a^1$ onto $X'$, the group 
$\operatorname{Sym}^2 \h_c^1(V_a^1,\sF(\psi\sigma))$ is therefore identified with the 
cohomology of $X'$ with coefficients in the part anti-invariant under $\tau$ of 
$\pi_\ast(\sF(\psi\sigma)\boxtimes \sF(\psi\sigma))$. Denote by $s$ the function 
on $X'$ such that $s\pi =\sigma\pr_1+\sigma\pr_2$. We have 
$\sF(\psi\sigma)\boxtimes \sF(\psi\sigma) = \sF(\psi(\sigma\pr_1+\sigma\pr_2)) = \sF(\psi s \pi) = \pi^\ast \sF(\psi s)$ 
(an isomorphism compatible with $\tau$). On the image $\pi(\Delta)$ of the 
diagonal, the desired sheaf of anti-invariants is therefore zero, while on 
the complement $X=X'\setminus \pi(\Delta)$ it is the tensor product of 
$\sF(\psi s)$ by $\varepsilon(P)$, where $\varepsilon$ is the unique character 
of order $2$ of $\{1,\tau\}$, and $P$ the $\{1,\tau\}$-torsor on $X$ given by the 
double covering induced by $V_a^1\times V_a^1$. We have 
\begin{equation*}\tag{7.17.1}\label{VI:eq:7-17-1}
  \operatorname{Sym}^2 \h_c^1(V_a^1,\sF(\psi\sigma)) = \h_c^2(X,\sF(\psi s)\otimes \varepsilon(P)) \text{.}
\end{equation*}

Let us compute $X$, $s$ and $\varepsilon(P)$. On $V_a^1\times V_a^1$, put 
$x_1=x_1\circ\pr_1$, $x_2=x_2\circ\pr_2$, $x_1'=x_1\circ \pr_2$, 
$x_2'=x_2\circ \pr_1$. We have $x_1 x_2=x_1'x_2'=a$. Identify functions on 
$X'$ with functions on $X$ invariant under $\tau$. Let $s_1=x_1+x_1'$, 
$s_2=x_2+x_2'$, $p_1=x_1 x_1'$, $p_2=x_2 x_2'$. We have $s_1 p_2 = a s_2$, 
$s_2 p_1 = a s_1$ and $X$ is defined in $X'$ by $s_1\ne 0$, $s_2\ne 0$. On 
$X$, we therefore have $p_1 = a s_1 s_2^{-1}$, and $(s_1,s_2)$ identifies $X$ with 
$\dG_m\times \dG_m$. On $X$, the covering $P$ has equation 
$T^2-s_1 T+a s_1 s_2^{-1}$ (take $T$ to be the coordinate $x_1$). If 
$T=s_1 t$, this is rewritten $t^2-t+a s_1^{-1} s_2^{-1}$. We therefore have 
$\varepsilon(P) = \sF(\psi(a s_1^{-1} s_2^{-1}))$. Since $s=s_1+s_2$, we have 
$\sF(\psi s)\otimes \varepsilon(P) = \sF(\psi(s_1+s_2+ a s_1^{-1} s_2^{-1}))$ 
and the right-hand side of \eqref{VI:eq:7-17-1} is identified with the left-hand side of 
\ref{VI:7-17}. 


\paragraph{Remark}

$\h_c^1(V_a^1,\sF(\psi\sigma))$ is also the first cohomology group of 
the elliptic curve obtained by completing the affine curve with equations 
\begin{align*}
  y^2-y &= x_1+x_2 \\
  x_1 x_2 &= a \text{;} 
\end{align*}
its modular invariant is $j=a^{-2}$. 





\subsection{Remark}\label{VI:7-18}

For $\chi$ a multiplicative character of $\dF_q^\times$, and $a\ne 0$, the 
method of \ref{VI:7-14} applies to the study of the sum 
\[
  S = \sum_{x_0\dotsm x_n = a} \chi(x_0) \psi(x_0+\cdots + x_n) 
\]
(replace the sheaf $\sF(\psi)$ on $\dG_m$ by $\sF(\chi\cdot \psi)$). We 
again find a cohomology group of dimension $n+1$, and 
$|S|\leqslant (n+1)q^{n/2}$. Such sums were considered by 
Salie and Mordell \cite{mo73}.

One may hope that our methods make it possible to study the sums 
\[
  S_{k,\chi,a} = \sum_{N(x)=a} \chi(x) \psi \tr(x) 
\]
($k$ an étale algebra over $\dF_q$, $a\in \dF_q^\times$, $\chi$ a character of 
$k^\times$). The analog of \ref{VI:7-4} (for $a\ne 0$) should be true. 





\subsection{}\label{VI:7-19}

The symmetric group $S_n$ acts on $\dA^n$, preserving $V_a$ and 
the map $\sigma$, and hence also the sheaf $\sf(\psi\sigma)$, the inverse 
image under $\sigma$ of the sheaf $\sf(\psi)$ on $\dG_a$. It therefore acts 
on $\h_c^{n-1}(V_a,\sF(\psi\sigma))$. Formula \eqref{VI:eq:7-2-5} admits 
the following cohomological interpretation.





\begin{proposition_}\label{VI:7-20}
The action of $S_n$ on $\h_c^{n-1}(V_a,\sF(\psi\sigma))$ is multiplication 
by the sign character. 
\end{proposition_}

It suffices to show that a transposition $\tau$ acts by multiplication 
by $-1$. Since $\tau^2=1$, the only possible eigenvalues of $\tau$ are 
$\pm 1$. Let $V_a/\tau$ be the quotient of $V_a$ for 
$\{\operatorname{id},\tau\}$. The map $\sigma$, and hence the sheaf 
$\sf(\psi\sigma)$, descend to the quotient, and the subspace of 
  $\h_c^{n-1}(V_a,\sF(\psi\sigma))$ fixed by 
$\tau$ is $\h_c^{n-1}(V_a/\tau,\sF(\psi\sigma))$. We must prove that 
this cohomology is zero. 

Suppose $a\ne 0$. For $\tau=(1,2)$ the quotient $V_a/\tau$ is identified 
with $\dG_a\times \dG_m^{n-2}$, the quotient map being 
$(x_1,x_2,x_3,\dots,x_n)\mapsto (x_1+x_2,x_3,\dots,x_n)$. Indeed, $x_1$ and 
$x_2$ are determined up to order by $x_1+x_2$ and 
$x_1x_2 = a/x_3\dotsm x_n$. In the coordinates $(t,x_3,\dots,x_n)$, the 
sheaf takes the form 
$\sF(\psi(t+x_3+\cdots + x_n)) = \sF(\psi(t))\otimes \sF(\psi(x_3+\cdots + x_n))$. 
  The K\"unneth formula and $\h_c^\bullet(\dG_a,\sF(\psi)) = 0$ therefore give 
the required vanishing. 

For $a=0$, we can deduce \ref{VI:7-20} from calculation \ref{VI:7-7}. 





\subsection{}\label{VI:7-21}

The sums $\kloos_{n,a}$ in this section were studied by 
S.\ Sperber, in his thesis, and in \cite{sp77} by methods 
$p$-adic methods inspired by Dwork and by Bombieri's article \cite{bo66}. 
For $p\ne 2$, these methods determine the dependence on $q$ in \ref{VI:7-5}, 
the functional equation $\alpha_i = q^{n-1} \beta_i^{-1}$ between the 
Frobenius eigenvalues for the sums $\kloos_{n,a}$ and 
$\kloos_{n,b}$ if $b=(-1)^n a$ (see \ref{VI:7-6}), and the value \ref{VI:7-15} 
of the product of the roots. For $p\geqslant n+3$, they also determine the 
$p$-adic valuations of the roots. The result is very beautiful: if we arrange 
the roots $\alpha_i$ in a suitable order, with 
  $0\leqslant i<n$, $\alpha_i q^{-i}$ are units. The 
restrictions on $p$ are probably not essential. 










\section{Other applications}\label{VI:8}





\subsection{}\label{VI:8-1}

The identity \ref{VI:1-9-4} can sometimes be used to calculate 
the number of rational points of an algebraic variety over $\dF_q$. In 
most cases where an exact result has been obtained, cohomology 
is expressed in terms of algebraic cycles; this is the case for rational 
surfaces, quadrics, smooth intersections of two quadrics of 
odd dimension, flag varieties\ldots. This is not, however, 
always the case; in Lusztig's beautiful article \cite{lu76}, algebraic cycles 
do not appear. 





\subsection{}\label{VI:8-2}

For reductive groups, cohomology can be calculated by lifting to 
characteristic $0$: if $B$ is a Borel subgroup of $G$, $G$ is a 
$B$-torsor over $G/B$; the cohomology of the smooth projective variety 
$G/B$ is invariant under specialization, and we use the Leray spectral sequence 
of $G\to G/B$ to prove the same result for $G$. It can be expressed 
conveniently in terms of \emph{the} maximal torus $T$ of $G$ and the action 
\emph{of the} Weyl group $W$ on $T$ (for the definition \emph{of the} maximal 
torus, cf. \cite[VIII \S 2 Rem.2]{bo05} or \cite[p.105]{dl76}). It is 
the exterior algebra on its primitive part, which agrees up to a shift 
with the quotient $I(\h^\bullet(B G,\dQ_\ell))$ of the cohomology of 
$B G$ formed by the ``indecomposable elements,'' and 
$\h^\bullet(B G,\dQ_\ell)=\h^\bullet(B T,\dQ_\ell)^W$. If $X$ is the group of 
characters of $T$, we have altogether 
\[
  \h^\bullet(G,\dQ_\ell) = \bigwedge \left(I\left(\operatorname{Sym}^\bullet(X\otimes \dQ_\ell(-1))^W\right)[-1]\right) \text{.}
\]
If $G$ is defined over a finite field $\dF_q$, $\gal(\dF/\dF_q)$ acts on $X$ 
and this isomorphism is Galois-compatible. Passing to 
cohomology with compact supports, we recover the classical formulas for the 
number of rational points of $G$. If $F^\ast:X\to X$ is defined by the 
Frobenius morphism $F:T\to T$, and $\deg_G(F)$ is the degree 
$q^{\dim G}$ of $F:G\to G$, we have 
\begin{equation*}\tag{8.2.1}\label{VI:eq:8-2-1}
  \left|G_0(\dF_q)\right| = \deg_G(F) \cdot \det\left(1-{F^\ast}^{-1},I(\operatorname{Sym}^\bullet(X\otimes \dQ)^W)\right) \text{.}
\end{equation*}

The same formula applies to the Ree and Suzuki groups; these groups are 
of the form $G^F$ (reductive $G$, $F^2$ a Frobenius), and it suffices to 
verify that the Lefschetz trace formula holds for $F$ acting 
on $G$. If $B$ is a Borel subgroup stable by $F$, we verify it 
directly for $F$ acting on $B$; for the action on $G/B$, which is proper and 
smooth, we can invoke the general theorems. 





\subsection{}\label{VI:8-3}

The trace formula was used by Deligne--Lusztig \cite{dl76}, 
Kazhdan \cite{ka77} and Springer \cite{sp76} in the study of 
complex representations of finite groups $G_0(\dF_q)$, for $G_0$ 
reductive over $\dF_q$. The work of Kazhdan and Springer contains 
admirable examples of how the trace formula permits ``analytic 
continuation'' from a split situation to a nonsplit situation 
(see \ref{VI:1-13}).

% END INCLUDED components/sga4half_6_en.tex

% BEGIN INCLUDED components/sga4half_7_en.tex
% !TEX root = ../SGA4half_English_Reader_1_to_8.tex


\chapter{Finiteness theorems in \texorpdfstring{$\ell$}{l}-adic cohomology}\label{VII}











\section{Statement of the theorems}\label{VII:1}

Throughout this exposé, $S$ will be a Noetherian scheme and $A$ a left
Noetherian torsion ring annihilated by an integer invertible on $S$. Our main
result is the following.






\begin{theorem_}\label{VII:1-1}
Suppose that $S$ is regular of dimension $0$ or $1$. Let $f:X\to Y$ be a
morphism of $S$-schemes of finite type, and let $\sF$ be a constructible sheaf
of left $A$-modules on $X$. Then the sheaves $\eR^i f_\ast \sF$ are
constructible.
\end{theorem_}

The proof will be given in Section \ref{VII:2} and in \ref{VII:3-10}.






\subsection{Remark}\label{VII:1-2}

In characteristic $0$, this result is less general than
\cite[XIX paragraphe 5]{sga4}, which proves the conclusion of the theorem for
every morphism of finite type of excellent schemes of characteristic $0$.
The proof of \cite[XIX]{sga4} uses, on the one hand, resolution of
singularities and, on the other, the assumption that the schemes have equal
characteristic (in order to deduce the ``purity theorem'' from resolution).






\subsection{Remark}\label{VII:1-3}

We shall say that a complex $K\in\ob\eD(X,A)$ is \emph{constructible} if its
cohomology sheaves are constructible, and shall use a subscript $c$ for the
subcategory of $\eD(X,A)$ (or $\eD^+$, $\eD^-$, $\eD^b$) consisting of
constructible complexes. In the language of derived categories, which we
shall use freely, \ref{VII:1-1} says that
$\eR f_\ast:\eD^+(X,A) \to \eD^+(Y,A)$ sends $\eD_c^+$ into $\eD_c^+$.
More explicitly:
\begin{enumerate}[\indent a)]
  \item For $K$ concentrated in a sheaf $\sF$ in degree $0$, one has
    $\sH^i\eR f_\ast K = \eR^i f_\ast \sF$; the derived statement therefore
    implies the theorem.
  \item Conversely, one invokes the spectral sequence
    \[
      E_1^{p q} = \eR^q f_\ast \sH^p(K) \Rightarrow \sH^{p+q} \eR f_\ast K \text{.}
    \]
\end{enumerate}

Speaking the language of derived categories has the advantage of replacing,
by the simple transitivity formula $\eR(f g)_\ast = \eR f_\ast \eR g_\ast$,
the Leray spectral sequence
$E_2^{p q} = \eR^p f_\ast \eR^q g_\ast \Rightarrow \eR^{p+q} (f g)_\ast$.

Under the hypotheses of \ref{VII:1-1}, it follows from \cite{sga4} that
$\eR f_\ast$ has finite cohomological dimension. This makes it possible above
to replace $\eD^+$ by $\eD$, $\eD^-$, or $\eD^b$.






\subsection{Ideas of the proof}\label{VII:1-4}
\begin{enumerate}[\indent a)]
  \item Poincaré duality treats the case where $X$ is smooth, $\sF$ is
    locally constant, and $Y=S$. The dimension hypothesis is seen na\"ively
    when calculating the $\rHom$ on $S$.
  \item Factor $f$ as $g j$: $g$ proper and $j$ an open immersion. One has
    $\eR f_\ast = \eR g_\ast \eR j_\ast$, and the finiteness theorem for
    proper morphisms controls $\eR g_\ast$. By dévissage, one reduces to
    assuming that $X$ is smooth, $\sF$ is locally constant, $f$ is an open
    immersion, and $Y$ is proper over $S$.
  \item Naively, an induction on $\dim X$ (with a change of $S$) makes it
    possible to suppose that the theorem is true outside a part of $Y$ finite
    over $S$. Denoting the structure morphisms by $a:X\to S$ and $b:Y\to S$,
    one then shows that if the theorem were false for $\sF$ and
    $\eR f_\ast$, it would also be false for $\sF$ and
    $\eR b_\ast \eR f_\ast = \eR a_\ast$: a contradiction.
\end{enumerate}






\begin{corollary_}\label{VII:1-5}
Under the hypotheses of the theorem, the categories $\eD_c(X,A)$ and
$\eD_c(Y,A)$ are taken into one another by the $4$ operations
$\eR f_\ast$, $\eR f_!$, $f^\ast$, $\eR f^!$.
\end{corollary_}

The case of $\eR f_\ast$ is treated above, that of $\eR f_!$ in
\cite[XVII 5.3]{sga4}, and that of $f^\ast$ is clear. It remains to consider
$\eR f^!$. The problem is local, which makes it possible to suppose that $f$
factors as $X\xrightarrow i Z \xrightarrow g Y$ ($i$ a closed immersion and
$g$ smooth, purely of relative dimension $n$). Poincaré duality
$\eR g^! K \iso g^\ast K(n)[2 n]$ and the transitivity
$\eR f^! = \eR i^! \eR g^!$ reduce us to proving \ref{VII:1-5} for $i$.
For every $L\in \eD(Z,A)$, if $j$ is the inclusion into $Z$ of
$U=Z\setminus X$, then $i_\ast \eR i^! L$ is the mapping cylinder of
$L\to \eR j_\ast j^\ast L$, and \ref{VII:1-5} for $i$ follows from
\ref{VII:1-1} for $j$.






\begin{corollary_}\label{VII:1-6}
Let $X$ be as in the theorem, and suppose that $A$ is commutative. Then, if
$\sF$ and $\sG$ are constructible sheaves of $A$-modules on $X$, the
$\Ext_A^i(\sF,\sG)$ are constructible: $\rHom$ sends
$\eD_c^-(X,A)\times \eD_c^+(X,A)$ into $\eD_c^+(X,A)$.
\end{corollary_}

By dévissage of $\sF$, we reduce to assuming that $\sF$ has the form
$j_!\sF_1$ ($j:Y\hookrightarrow X$ a locally closed immersion and $\sF_1$
locally constant on $Y$). One then has
\[
  \rHom(j_! \sF_1,\sG) = \eR j_\ast \Hom(\sF_1,\eR j^! \sG) \text{:}
\]
by \ref{VII:1-3}, we are reduced to the case where $\sF$ is locally
constant---or even constant, since the problem is local. For $\sF$ locally
constant and constructible, one has
$\Hom(\sF,\sG)_x=\hom(\sF_x,\sG_x)$, and likewise for $\rHom$. For $\sF$
constant, the $\Ext^i$ can therefore be calculated using a projective
resolution of finite type of its constant value, and the $\Ext^i$ are
constructible if $\sG$ is---whence the corollary.






\subsection{Remark}\label{VII:1-7}

Since $\eR f_\ast$ has finite cohomological dimension, it takes complexes of
finite tor-dimension (resp. $\leqslant d$) to complexes of finite
tor-dimension (resp. $\leqslant d$) \cite[XVII 5.2.11]{sga4}.
\cite[XVII 5.2.10]{sga4} and the proof of \ref{VII:1-5} then show that finite
tor-dimension is stable under the $4$ operations $\eR f_\ast$,
$\eR f_!$, $f^\ast$, $\eR f^!$. A variant of the argument of
\ref{VII:1-6} (apply dévissage to $K$ along a partition of $X$ so that the
cohomology sheaves become locally constant, then localize so as to replace $K$
by a finite complex of locally free sheaves of finite type) then shows that,
for $A$ commutative, $\rHom$ induces
\[
  \rHom:\eD_\text{ctf}^b(X,A)\times \eD_\text{tf}^b(X,A) \to \eD_\text{tf}^+(X,A)
\]
(t.f.: finite tor-dimension).






\subsection{}\label{VII:1-8}

Under the hypotheses of the theorem, the same method allows us to prove a
local biduality theorem $K\iso D D K$ (Section \ref{VII:4}). We also prove a
finiteness theorem for vanishing-cycle sheaves. For arbitrary $S$, one still
obtains a ``generic'' theorem:






\begin{theorem_}\label{VII:1-9}
Let $f:X\to Y$ be a morphism of $S$-schemes of finite type, and let $\sF$ be
a constructible sheaf of $A$-modules on $X$. There is a dense open subset
$U$ of $S$ such that
\begin{enumerate}[\indent (i)]
  \item Above $U$, the $\eR^i f_\ast \sF$ are constructible, and all but
    finitely many of them vanish.
  \item Formation of the $\eR^i f_\ast \sF$ is compatible with every base
    change $S' \to U\subset S$.
\end{enumerate}
\end{theorem_}

For $S$ the spectrum of a field, one has $U=S$.

If $S$ is the spectrum of a field, a passage-to-the-limit argument gives
compatibility with $S$-base changes for the $\eR^i f_\ast$ for every
quasi-compact quasi-separated morphism of $S$-schemes and every sheaf $\sF$.






\begin{corollary_}\label{VII:1-10}
Let $X$ be a scheme of finite type over a separably closed $k$, and let
$\sF$ be a constructible sheaf of $A$-modules. Then the $\h^i(X,\sF)$ are of
finite type.
\end{corollary_}

This is the special case $S=\spec(k)$, $Y=S$.

\begin{corollary_}\label{VII:1-11}
Let $X$ and $Y$ be two schemes of finite type over a separably closed $k$,
and let $K\in \ob\eD^-(X,A^\circ)$,
$L\in \ob\eD^-(Y,A)$ be complexes of sheaves of right and left modules. The
K\"unneth arrow
\[
  \eR \Gamma(X,K) \lotimes \eR \Gamma(Y,L) \to \eR\Gamma(X\times Y,\pr_1^\ast K \lotimes \pr_2^\ast L)
\]
is an isomorphism.
\end{corollary_}

Proceed as in \cite[XVII 5.4.3]{sga4}:
\[\xymatrix{
  X\times Y \ar[r]^-{\text{pr}_1} \ar[d]^-{\text{pr}_2}
    & X \ar[d]^-a \\
  Y \ar[r]^-b
    & \spec(k)
}\]
base change by $b$ shows that $\eR\pr_{2,\ast}{} \pr_1^\ast K$ is
$b^\ast \eR\Gamma(X,K)$. One then has (cf.
\cite[XVII 5.2.11]{sga4} and the proof)

\begin{align*}
  \eR\Gamma(X\times Y,\pr_1^\ast K\lotimes \pr_2^\ast L)
    &= \eR\Gamma(Y,\eR \pr_{2,\ast}(\pr_1^\ast K\lotimes \pr_2^\ast L)) \\
    &= \eR\Gamma(Y,(\eR\pr_{2,\ast}\pr_1^\ast K)\lotimes L) \\
    &= \eR\Gamma(Y,b^\ast \eR\Gamma(X,K)\lotimes L) \\
    &= \eR\Gamma(X,K)\lotimes \eR\Gamma(Y,L)
\end{align*}






\subsection{}\label{VII:1-12}

Finally, as a local counterpart of \ref{VII:1-9}(ii), we shall prove a
generic local acyclicity theorem (\ref{VII:2-13}).











\section{Generic theorems}\label{VII:2}

In this section, we prove \ref{VII:1-9} and the generic local acyclicity theorem
\ref{VII:2-13}. To prove \ref{VII:1-9}, we immediately reduce to the case where
$S$ is integral. Let $\eta$ be its generic point.





\subsection{}\label{VII:2-1}

We begin by proving \ref{VII:1-9} under the following additional assumptions:
$X$ is smooth over $S$, purely of relative dimension $n$, $A=\dZ/m$, $\sF$ is
locally constant, and $Y=S$. Let $\sF'=\Hom(\sF,\dZ/m)$. We shall show that if
the $\eR^i f_! \sF'$ are locally constant, the conclusions of \ref{VII:1-9}
hold with $U=S$. Recall that if $\sL$ is a constructible locally constant
sheaf, the $\Ext^i(\sL,\sG)$ may be computed fiber by fiber, and are
constructible if $\sG$ is. Recall also that $\dZ/m$ is an injective
$\dZ/m$-module and is the dualizing module, so that
$\sF=\rHom(\sF',\dZ/m)$. Poincaré duality
\[
  \eR f_\ast \rHom(K,\eR f^! L) = \rHom(\eR f_! K,L) \text{,} 
\]
for $K=\sF'$, $L=\dZ/m$, $\eR f^! L = \dZ/m[2n](n)$, therefore gives
\[
  \eR^{2n-i} f_\ast \sF = \Hom(\eR^i f_! \sF',\dZ/m)(-n) \text{.} 
\]

The sheaf on the right is locally constant and its formation commutes with
every base change, which proves the assertion.





\subsection{}\label{VII:2-2}

Let us prove \ref{VII:1-9} under the additional assumptions that $X$ is smooth
over $S$, $\sF$ is locally constant, and $Y=S$.
\begin{enumerate}[a)]
  \item By decomposing $X$ into connected components, we reduce to the case
    where it is purely of some relative dimension $n$ over $S$.
  \item By decomposing $A$ as a product, we reduce to the case where
    $\ell^m A=0$, with $\ell$ a prime invertible on $S$; we then filter $\sF$
    by the $\ell^k \sF$ and use the corresponding spectral sequence to reduce
    to the case $\ell\sF=0$. We may then replace $A$ by $A/\ell A$ and assume
    that $\ell A=0$.
  \item After replacing $S$ by a suitable $U$, we may assume that there exists
    a finite étale Galois covering $X_1/X$, with Galois group $G$, such that
    the inverse image $\sF_1$ of $\sF$ on $X_1$ is a constant sheaf, with
    constant value $F$. Writing $f_1$ for the projection of $X_1$ onto $S$,
    we then have
    \[
      \eR^i {f_1}_\ast \sF_1 = (\eR^i {f_1}_\ast \dZ/\ell)\otimes_{\dZ/\ell} F \text{.} 
    \]
    We also have the Hochschild--Serre spectral sequence (or the Leray
    spectral sequence for the covering $X_1/X$)
    \[
      \underline\h^p(G,\eR^q {f_1}_\ast \sF_1) \Rightarrow \eR^{p+q} f_\ast \sF \text{.} 
    \]
    By \ref{VII:2-1}, after shrinking $U$ if necessary, we may assume that
    the $\eR^i {f_1}_\ast \dZ/\ell$ are locally constant and that their
    formation commutes with every base change. The same property then holds
    for the $\eR^i f_\ast \sF$. Finally, if $\bar\eta$ is a geometric point
    localized at the generic point $\eta$ of $S$, almost all the
    $\eR^i f_\ast\sF$ vanish, because the
    $(\eR^i f_\ast \sF)_{\bar\eta} = \h^i(X_{\bar \eta},\sF)$ do.
\end{enumerate}





\subsection{}\label{VII:2-3}

We prove by induction on $n$ that
\begin{equation*}\tag{$\ast_n$}\label{VII:eq:2-3-1}
  \begin{array}{c}
    \text{The conclusions of \ref{VII:1-9} hold when $\dim{X_\eta}\leqslant n$} \\
    \text{and $f$ is an open immersion with dense image.}
  \end{array}
\end{equation*}

For $n=0$, after shrinking $S$ if necessary, we have $X=Y$: $(\ast_0)$ is
clear. Assume $(\ast_{n-1})$, and prove \eqref{VII:eq:2-3-1}. In
($\ast_{n-1}$), ``open immersion'' may be replaced by ``immersion,'' as is
seen by factoring it as an open immersion followed by a closed immersion.





\begin{lemma_}\label{VII:2-4}
After shrinking $S$ if necessary, the conclusions of \ref{VII:1-9} hold over
$Y'\subset Y$, where the complement $Y_1$ of $Y'$ is finite over $S$.
\end{lemma_}

The assertion is local on $Y$, which we may assume affine:
$Y\subset \dA_S^n$. The induction hypothesis ($\ast_{n-1}$) applies to
\[\xymatrix{
  X \ar[r]^-f \ar[dr] 
    & Y \ar[d]^-{\pr_i} \\
  & \dA_S^1 
}\]
Thus, for each $i$, there exists a dense open subset $U_i$ of $\dA_S^1$ such
that the conclusions of \ref{VII:1-9} hold over $\pr_i^{-1}(U_i)$; they hold
over the union of the $\pr_i^{-1}(U_i)$, and \ref{VII:2-4} follows.





\addtocounter{subsection}{1} % a section is skipped in the original
\subsection{}\label{VII:2-6}

Let us prove \eqref{VII:eq:2-3-1} when $X$ is smooth over $S$ and $\sF$ is
locally constant on $X$. Since the problem is local on $Y$, we may first
assume $Y$ affine and then projective (replacing $Y$ by its closure in a
projective space). Let $i:Y_1\hookrightarrow Y$ and $j:Y' \to Y$ be as in
\ref{VII:2-4}.
\[\xymatrix{
  X \ar@{^{(}->}[r] \ar[dr]_-a 
    & Y \ar[d]^-b 
    & \ar[l]_-i \ar[dl]^-{b_1} Y_1 \\
  & S 
}\] 
After shrinking $S$ if necessary, we know that $j^\ast \eR f_\ast \sF$ is
constructible and that its formation commutes with every base change in $S$;
we also know that $\eR a_\ast \sF = \eR b_\ast \eR f_\ast \sF$ is
constructible and that its formation commutes with every base change.

Apply $\eR b_\ast$ to the triangle defined by the exact sequence
\begin{equation*}\tag{1}\label{VII:eq:2-6-1}
\xymatrix{
  0 \ar[r] 
    & j_! j^\ast \eR f_\ast \sF \ar[r] 
    & \eR f_\ast \sF \ar[r] 
    & i_! i^\ast \eR f_\ast \sF \ar[r] 
    & 0 \text{:} 
}
\end{equation*}
we obtain a triangle
\begin{equation*}\tag{2}\label{VII:eq:2-6-2}
\xymatrix{
  \ar[r] 
    & \eR b_\ast j_! j^\ast \eR f_\ast \sF \ar[r] 
    & \eR a_\ast \sF \ar[r] 
    & {b_1}_\ast i^\ast \eR f_\ast \sF \ar[r] 
    & 
}
\end{equation*}
in which the first two terms are constructible and their formation commutes
with every base change in $S$ (for the first, by the finiteness theorem for
the proper morphism $b$). The same therefore holds for the third. Since $b_1$
is finite, it follows that $i^\ast \eR f_\ast \sF$ is constructible and its
formation commutes with every base change in $S$; the same holds for
$\eR f_\ast \sF$ by \eqref{VII:eq:2-6-1}.





\subsection{}\label{VII:2-7}

Let us prove \eqref{VII:eq:2-3-1} in general. We first reduce to the case where
$X$ contains a dense open subset $V$ smooth over $S$. When $S$ is the spectrum
of a perfect field, it is enough to replace $X$ by $X_\text{red}$ (and $Y$ by
$Y_\text{red}$). In general, we must shrink $S$, make a finite, radicial,
surjective base change $S' \to S$, and replace $X$ and $Y$ by $X_\text{red}$
and $Y_\text{red}$. Since the étale topology is insensitive to finite,
radicial, surjective morphisms, this is harmless. After shrinking $V$, we may
assume that $\sF$ is locally constant on $V$.
\[\xymatrix{
  V \ar@{^{(}->}[r]^j 
    & X \ar@{^{(}->}[r]^-f 
    & Y \text{.} 
}\]
Define $\Delta$ by the triangle
\begin{equation*}\tag{1}\label{VII:eq:2-7-1}
\xymatrix{
  \ar[r] 
    & \sF \ar[r] 
    & \eR j_\ast j^\ast \sF \ar[r] 
    & \Delta \ar[r] 
    & \text{.} 
}
\end{equation*}
The cohomology sheaves of $\Delta$ are supported on $X\setminus V$, and
$\dim(X\setminus V)_\eta < n$. The induction hypothesis therefore allows us
to assume that $\eR f_\ast \Delta$ is constructible. Applying $\eR f_\ast$ to
the triangle \eqref{VII:eq:2-7-1}, we obtain a triangle
\[\xymatrix{
  \ar[r] 
    & \eR f_\ast \sF \ar[r] 
    & \eR(f j)_\ast j^\ast \sF \ar[r] 
    & \eR f_\ast \Delta \ar[r] 
    & 
}\]
in which two of the vertices are constructible and their formation commutes
with every base change in $S$. The same therefore holds for the third.





\subsection{}\label{VII:2-8}

Let us prove \ref{VII:1-9}. The problem is local on $Y$, which we may assume
affine. Taking an affine covering of $X$ and invoking the Leray spectral
sequence, we reduce to the case where $X$ is affine as well. This ensures that
we can factor $f$ as an open immersion followed by a proper morphism:
$f=g j$, hence $\eR f_\ast = \eR g_\ast \eR j_\ast$. The open immersion is
covered by \eqref{VII:eq:2-3-1}, and the proper morphism by the finiteness
theorem.

The following corollaries are proved as in \hyperref[VII:1]{Section 1}.





\begin{corollary_}\label{VII:2-9}
Under the hypotheses of the theorem, for $K$ in $\eD_c^b(X,A)$ or
$\eD_c^b(Y,A)$ respectively, there exists a nonempty open subset $U$ of $S$
over which $\eR f_\ast K$, $\eR f_! K$, $f^\ast K$, and $\eR f^! K$ belong to
$\eD_c^b(Y,A)$ or $\eD_c^b(X,A)$, and their formation commutes with every base
change $S' \to U\subset S$.
\end{corollary_}





\begin{corollary_}\label{VII:2-10}
Let $X$ be as in the theorem, and suppose that $A$ is commutative. Then, if
$\sF$ and $\sG$ are constructible sheaves of $A$-modules on $X$, over a dense
open subset $U$ of $S$, the $\Ext_A^i(\sF,\sG)$ are again constructible and
their formation commutes with every base change $S' \to U\subset S$.
\end{corollary_}





\subsection{}\label{VII:2-11}

If $x$ is a geometric point of a scheme $X$, we denote by $X_x$ the strict
henselization of $X$ at $x$. For $f:X\to S$ and $t$ a geometric point of $S$,
we denote by $X_t$ the geometric fiber of $X$ at $t$. Finally, for $x$ a
geometric point of $X$ and $t$ a geometric point of $S_{f(x)}$, $(X_x)_t$ is
the fiber at $t$ of $X_x \to S_{f(x)}$.





\begin{definition_}\label{VII:2-12}
Let $f:X\to S$ and $K\in \ob\eD^+(X,A)$. We say that $f$ is \emph{locally
acyclic at $x$ relative to $K$} if, for every geometric point $t$ of
$S_{f(x)}$, we have $K_x=\eR\Gamma(X_x,K)\iso \eR\Gamma(X_{x,t},K)$. We say
that $f$ is \emph{locally acyclic relative to $K$} if this holds for every
$x$, and \emph{universally locally acyclic relative to $K$} if it remains true
after every base change $S' \to S$.
\end{definition_}





\begin{theorem_}\label{VII:2-13}
Let $f:X\to S$ be a morphism of finite type and let $\sF$ be constructible on
$X$. There exists a dense open subset $U$ of $S$ over which $f$ is universally
locally acyclic relative to $\sF$.
\end{theorem_}

We shall use the following result, proved in the appendix.





\begin{lemma_}\label{VII:2-14}
Let $X\xrightarrow f S_1 \xrightarrow g S_2$. If $g$ is smooth and $f$ is
universally locally acyclic relative to $K$, then so is $g f$.
\end{lemma_}

We reduce to the case where $S$ is integral, with generic point $\eta$, and
argue by induction on $\dim{X_\eta}$. We first deduce from the induction
hypothesis that
\begin{equation*}\tag{A}\label{VII:eq:2-14-1}
  \begin{array}{c}
    \text{After shrinking $S$, there exists $T\subset X$, finite over $S$, such that $f$} \\
    \text{is universally locally acyclic outside $T$.} 
  \end{array}
\end{equation*}

Since this question is local, we localize on $X$ and factor $f$ as
$X\xrightarrow u \dA_S^1 \to S$. The induction hypothesis applies
to $u$, and we conclude by \ref{VII:2-14} and the usual arguments.

To prove the theorem, we may assume that $X$ is proper, since the problem is
local. We shrink $S$ so that the $\eR^i f_\ast \sF$ are locally constant and
\eqref{VII:eq:2-14-1} applies, and use the following lemma.





\begin{lemma_}\label{VII:2-15}
Let $f:X\to S$ be proper and let $\sF$ be constructible. Suppose that the
$\eR^i f_\ast \sF$ are locally constant and that, outside $T\subset X$ finite
over $S$, $f$ is locally acyclic relative to $\sF$. Then $f$ is locally
acyclic relative to $\sF$.
\end{lemma_}

Let $s$ be a geometric point of $S$, let $S_s$ be the strict localization of
$S$ at $s$, and let $t$ be a geometric point of $S_s$. Let $X_{(s)}$ be the
inverse image of $X$ on $S_s$, let $\bar j:X_t \to X_{(s)}$, and let
$i:X_s \hookrightarrow X_{(s)}$. Again denoting by $\sF$ the inverse image of
$\sF$ on $X_{(s)}$ or $X_t$, we must prove that
$i^\ast \sF \iso i^\ast \eR \bar j_\ast \sF$. Let $\Delta$ be the mapping
cylinder of this map. Its cohomology sheaves are supported on $T_s$. To prove
that $\Delta=0$, it therefore suffices to prove that
$\eR\Gamma(X_s,\Delta)=0$. This is the mapping cylinder of
$(\eR f_\ast \sF)_s = \eR\Gamma(X_s,\sF) \to \eR\Gamma(X_s,i^\ast \eR\bar j_\ast \sF) = \eR\Gamma(X_{(s)},\eR\bar j_\ast \sF) = \eR\Gamma(X_t,\sF) = (\eR f_\ast \sF)_t$.
By hypothesis this specialization morphism is an isomorphism, and
$\Delta=0$.





\begin{corollary_}\label{VII:2-16}
When $S$ is the spectrum of a field, every $S$-scheme $X$ is universally
locally acyclic relative to $K$, for every $K\in\ob\eD^+(X,A)$.
\end{corollary_}

This follows from \ref{VII:2-13} by passage to the limit (which is possible
because the $U$ of \ref{VII:2-13} is always equal to $S$.








\section{Proof of \texorpdfstring{\ref{VII:1-1}}{1.1} and constructibility of the vanishing-cycle sheaves}\label{VII:3}






\subsection{}\label{VII:3-1}

Let $S$ be a strictly local trait (the spectrum of a henselian discrete
valuation ring with separably closed residue field). Denote its closed and
generic points by $s$ and $\eta$, and let $\bar\eta$ be a geometric point
localized at $\eta$ (a separable closure of $k(\eta)$). Let $X$ be over $S$
and let $\sF$ be a sheaf on $X_\eta$. Recall the definition of the
vanishing-cycle sheaves $\eR^i \Psi_\eta(\sF)$ (sheaves on $X_s$, endowed with
an action of $\gal(\bar\eta/\eta)$). Let $i:X_s\hookrightarrow X$,
$j:X_\eta \hookrightarrow X$, and let $\bar j$ be the composite
$X_{\bar\eta}\to X_\eta \hookrightarrow X$; then
\[
  \eR^i \Psi_\eta(\sF) = i^\ast \eR^i \bar j_\ast (\bar j^\ast \sF) \text{.} 
\]

Here is the main theorem of this subsection. It improves
\cite[XIII 2.3.1, 2.4.2]{sga7}.






\begin{theorem_}\label{VII:3-2}
Suppose that $X$ is of finite type over $S$ and that $\sF$ is constructible.
Then the vanishing-cycle sheaves $\eR^i \Psi_\eta(\sF)$ are constructible.
\end{theorem_}

We argue by induction on the dimension $n$ of $X_\eta$. We may moreover
assume---and do assume---that $X_\eta$ is dense in $X$ (replacing $X$ by
$\bar X_\eta$ does not change the $\eR^i \Psi$, which are supported on
$\bar X_\eta$).






\begin{lemma_}\label{VII:3-3}
If \ref{VII:3-2} holds when $X_\eta<n$, and if
$\dim X_\eta=n$, then there exist constructible subsheaves $\sG_i$ of the
$\eR^i \Psi_\eta(\sF)$ such that the supports of the local sections of
$\eR^i\Psi_\eta(\sF)/\sG_i$ are finite.
\end{lemma_}

Let $s'$ be a geometric generic point of the affine line over $s$, and let
$S'$ be the strict localization at $s'$ of the affine line $\dA_S^1$ over
$S$. It is again a strictly local trait, and uniformizers for $S$ are
uniformizers on $S'$.
\[\xymatrix{
  & \spec(k') \ar[rr] \ar[dd] \ar[dr] 
    & & (S',\eta',s') \ar[dr] \ar[dd] \\
  \bar\eta' \ar[ur] \ar[dr] 
    & & \dA_{\bar S}^1 \ar[d]  
    & & \dA_S^1 \ar[dl] \\
  & \bar\eta\ar@{^{(}->}[r] 
    & \bar S \ar[r] 
    & (S,\eta,s) 
}\]
Let $k'=k(\bar\eta)\otimes_{k(\eta)}k(\eta')$. If $\bar S$ is the
normalization of $S$ in $\bar\eta$, then $k'$ is the fraction field of the
strict localization of $\dA_{\bar S}^1$ at $s'$; it is therefore a field, and
$\gal(\bar\eta/\eta)=\gal(k'/\eta')$. Let $\bar\eta'$ be the spectrum of an
algebraic closure of $k'$. The Galois group $P=\gal(\bar\eta'/k')$ is a
pro-$p$ group, where $p$ is the characteristic exponent of $k(s)$.






\begin{lemma_}\label{VII:3-4}
For every scheme $X'$ over $S'$, and every sheaf $\sF$ on
$X_\eta'=X_{\eta'}'$, the vanishing-cycle sheaves for $X'/S'$ and for $X'/S$
are related by
$\eR^i \Psi_\eta(\sF) =\eR^i\Psi_{\eta'}(\sF)^P$.
\end{lemma_}

The following diagram compares the morphisms used in the definition of
$\eR\Psi$ for $X'/S'$ and $X'/S$:
\[\xymatrix{
  X_{\bar\eta'} \ar[rr] \ar[d] 
    & & X_{\eta'}' \ar@{=}[dd] \ar@{^{(}->}[dr] 
    & & S' \\ 
  X_{\eta'}\times_{\eta'}\spec(k') \ar@{=}[d] \ar[urr] 
    & & & X \ar[ur] \ar[dr] \\
  X_\eta' \ar[rr] 
    & & X_\eta' \ar@{^{(}->}[ur] 
    & & S
}\]
and the lemma follows from the Hochschild--Serre spectral sequence, since $P$
is a pro-$p$ group and $p$ is invertible in $A$.

Let us prove \ref{VII:3-3}. The question is local; this allows us to assume
that $X$ is affine, $X\subset \dA_S^n$. Let $f$ be one of the projections
$X\subset \dA_S^n\to \dA_S^1$. Let $X'$ be the ``localization''
$X_{\dA_S^1} S'$ of $X$, and let $\sF'$ be the inverse image of $\sF$ on $X'$.
\[\xymatrix{
  \sF'\text{ on }X' \ar[r] \ar[d]^-\lambda 
    & S' \ar[d]^-\lambda \\
  \sF \text{ on }X \ar[r]^-f 
    & \dA_S^1 \ar[r] 
    & S \text{.} 
}\]
On $X_s'=X_{s'}'$, we have
\[
  \lambda^\ast \eR^i\Psi_\eta(\sF) = \eR^i \Psi_\eta(\sF') = \eR^i\Psi_{\eta'}(\sF')^P \text{:} 
\]
this sheaf is constructible because the induction hypothesis applies to
$X'/S'$, and we use the following lemma, applied to
$X_s\subset \dA_S^n$ and to the vanishing-cycle sheaves.






\begin{lemma_}\label{VII:3-5}
Let $\dA^n$ be affine $n$-space over a field $k$, let $X\subset \dA^n$, let
$\sF$ be a sheaf of $A$-modules on $X$, and let $\bar\eta$ be a geometric
generic point of $\dA^1$. Let $X_{\bar\eta,i}$ be the geometric generic fiber
of $X\subset \dA^n\xrightarrow{\pr_i} \dA^1$, and let
$\sF_{\bar\eta,i}$ be the restriction of $\sF$ to $X_{\bar\eta,i}$. If the
$\sF_{\bar\eta,i}$ are constructible, there exists a constructible
$\sF'\subset \sF$ such that the local sections of $\sF/\sF'$ have finite
support.
\end{lemma_}

By passage to the limit, for each $i$
\begin{enumerate}[a)]
  \item There exists an étale neighborhood $U$ of $\bar\eta$ and a
    constructible sheaf $\sH$ on
    $X_{U,i} = X\times_{\dA^1,\pr_i} U$ whose restriction
    $\sF_{\bar\eta,i}$ to $X_{\bar\eta,i}$ is the given one.
  \item For a suitable $U$, the isomorphism
    $\sH_{\bar\eta,i} \iso \sF_{\bar\eta,i}$ comes from
    $u:\sH\to \sF_{U,i}$, where $\sF_{U,i}$ is the inverse image of $\sF$ on
    $X_{U,i}$. Let $\varphi:X_{U,i} \to X$. The morphism $u$ defines
    $v_i:\varphi_! \sH \to \sF$, with image $\sF_i'$. We have
    $(\sF/\sF_i')_{\bar\eta,i} = 0$, and for $\sF'$ we take the sum of the
    $\sF_i'$.
\end{enumerate}






\subsection{}\label{VII:3-6}

Let us prove \ref{VII:3-2} in dimension $n$. We may assume that $X$ is affine
(because the problem is local on $X_s$), and then projective over $S$ (by
taking the closure of $X$ in a projective embedding). We then have a spectral
sequence \cite[I 2.2.3]{sga7} (the Leray spectral sequence for $\bar j$, taking
into account the base-change theorem for the proper morphism $X\to S$).
\begin{equation*}\tag{1}\label{VII:eq:3-6-1}
  E_2^{p q} = \h^p(X_s,\eR^q\Psi_\eta(\sF)) \Rightarrow \h^{p+q}(X_{\bar\eta},\sF) \text{.} 
\end{equation*}
Let $\sG^q\subset \eR^q \Psi_\eta(\sF)$ be as in \ref{VII:3-3}, and let
$\sH^q$ be the quotient sheaf. We have $\h^i(X_s,\sH^q)=0$ for $i\ne 0$, and
for $\sH^q$, hence $\eR^q\Psi_\eta(\sF)$, to be constructible, it suffices
that $\h^0(X_s,\sH^q)$ be of finite type.

Compute \eqref{VII:eq:3-6-1} modulo modules of finite type (i.e. in the
quotient of the category of $A$-modules by the thick subcategory of
$A$-modules of finite type). The long cohomology exact sequence deduced from
$0 \to \sG^q \to \eR^q\Psi_\eta(\sF) \to \sH^q \to 0$, together with
\ref{VII:1-10} applied to $\sG^q$, gives
$E_2^{p q} \sim \h^p(X_s,\sH^q)$. In particular,
$E_2^{p q}\sim 0$ for $p\ne 0$,
$E_2^{0 q}\sim \h^q(X_{\bar\eta},\sF) \sim 0$, and this completes the proof.

Parallel arguments give the following result, which improves
\cite[XIII 2.1.12, 2.4.2]{sga7}.






\begin{proposition_}\label{VII:3-7}
The formation of the vanishing-cycle sheaves commutes with changes of traits.
\end{proposition_}

Let $g:(S',\eta',s') \to (S,\eta,s)$ be a (surjective) morphism of strictly
local traits, let $X/S$, let $\sF$ be on $X_\eta$, of torsion prime to the
residual characteristic, and let $(X',\sF')$ be their inverse images on $S'$.
Again denoting by $g$ the morphisms parallel to $g$, such as
$X_{s'}' \to X_s$, we must prove that
\begin{equation*}\tag{3.7.1}\label{VII:eq:3-7-1}
\xymatrix{
  g^\ast \eR\Psi_\eta(\sF) \ar[r]^-\sim 
    & \eR\Psi_{\eta'}(\sF') \text{.} 
}
\end{equation*}
This morphism is defined by the commutative diagram
\[\xymatrix{
  X_{s'}' \ar@{^{(}->}[r]^-{i'} \ar[d]^-g 
    & X' \ar[d]^-g 
    & X_{\eta'}' \ar[l]_-{\bar j'} \ar[d]^-g \\
  X_s \ar@{^{(}->}[r]^-i 
    & X 
    & \ar[l]_-{\bar j} X_{\bar\eta} 
}\]

A passage to the limit reduces us to the case where $X$ is of finite type
over $S$, and we argue by induction on $\dim{X_\eta}$. Since the question is
local, we may assume that $X$ is affine, and then projective. Assume
\eqref{VII:eq:3-7-1} when $\dim{X_\eta}<n$. When $\dim{X_\eta}=n$, it then
follows from the following two assertions, where $\Delta$ is the mapping
cylinder of \eqref{VII:eq:3-7-1}.

\begin{enumerate}[(A)]
  \item The support of the local sections of the cohomology sheaves of
    $\Delta$ is finite.
    \phantomsection\hypertarget{VII:3-A}{}
  \item $\eR\Gamma(X_s,\Delta)=0$, hence $\Delta=0$, by
    (\hyperlink{VII:3-A}{A}).
    \phantomsection\hypertarget{VII:3-B}{}
\end{enumerate}





\subsection{}\label{VII:3-8}

The proof of (\hyperlink{VII:3-A}{A}) parallels that of \ref{VII:3-3}. We
localize on $X$ and consider maps $X\to \dA_S^1$; denoting by $S(x)$ the
strict henselization of $\dA_S^1$ at the generic point of the special fiber,
we apply the induction hypothesis to $X_1/S(x)$
($X_1=X\times_{\dA^1} S(x)$) and to $S'(x)$. Taking invariants under a pro-$p$
group, we find that $\Delta$ vanishes on the geometric generic fiber of
$X_{s'} \to \dA_s^1$. Since this holds locally on $X$ and for every
projection, we obtain (\hyperlink{VII:3-A}{A}).






\subsection{}\label{VII:3-9}

For (\hyperlink{VII:3-B}{B}), the properness of $X$ over $S$ ensures that
\[
  \eR\Gamma(X_{s'}',\eR\Psi_{\eta'}(\sF')) = \eR\Gamma(X_{\bar\eta}',\sF') 
\]
and
\[
  \eR\Gamma(X_s,\eR\Psi_\eta(\sF)) = \eR\Gamma(X_{\bar\eta},\sF) \text{.} 
\]
We have
$\eR\Gamma(X_s,\eR\Psi_\eta(\sF))\iso \eR\Gamma(X_{s'}',g^\ast\eR\Psi_\eta(\sF))$
(invariance under change of the separably closed base field); thus
$\eR\Gamma(X_{s'}', \Delta)$ is the mapping cylinder of
$\eR\Gamma(X_{\bar\eta},\sF) \to \eR\Gamma(X_{\bar\eta}',\sF')$, an
isomorphism by the same invariance theorem.






\subsection{}\label{VII:3-10}

Let us prove \ref{VII:1-1}. We reduce to the case where $S$ is connected,
hence integral. If $\dim S=0$ ($S$ is the spectrum of a field), apply
\ref{VII:1-9}, \ref{VII:1-10}. If $S$ has dimension $1$, Theorem
\ref{VII:1-9} ensures that the $\eR^i f_\ast \sF$ are constructible over the
complement of a finite set $T$ of points of $S$. It remains to see that their
restrictions to the $Y_t$ ($t\in T$) are constructible. It suffices to see
this after localization at $t$: we may assume that $S$ is a strictly local
trait. The essential case is the following.






\begin{lemma_}\label{VII:3-11}
Let $X$ be of finite type over $S$, let $j:X_\eta\hookrightarrow X$,
$i:X_s\hookrightarrow X$, and let $\sF$ be constructible on $X_\eta$. Then
$i^\ast\eR j_\ast \sF$ is constructible.
\end{lemma_}

Let $I=\gal(\bar\eta/\eta)$ be the inertia group. It is known to be an
extension of $\widehat\dZ_{p'}(1)$ by a pro-$p$ group $P$ ($p$ being the
characteristic exponent of the residue field). The Hochschild--Serre spectral
sequence gives
\[
  E_2^{p q} = \sH^p(I,\eR^q\Psi_\eta(\sF))\Rightarrow i^\ast \eR^{p+q}j_\ast\sF \text{.} 
\]
Let $R_t^q = \eR^q\Psi_\eta(\sF)^P$ ($t$ for tame). If $\sigma$ is a
generator of $\widehat\dZ_{p'}(1)$, then
$E_2^{0 q} = \ker(\sigma-1,R_t^q)$,
$E_2^{1 q} = \operatorname{coker}(\sigma-1,R_t^q)$, and $E_2^{p q} = 0$ for
$p\ne 0,1$. These sheaves are constructible because the vanishing-cycle
sheaves are, and \ref{VII:3-11} follows.






\subsection{}\label{VII:3-12}

Let us treat the general case, with $S$ still a strictly henselian trait. If
$X=X_\eta$, then $f$ is the composite $X\to Y_\eta \to Y$, and we apply
\ref{VII:1-9} and \ref{VII:3-11}. In the general case, let
$j:X_\eta\hookrightarrow X$ and let $\Delta$ be the mapping cylinder of
$\sF \to \eR j_\ast j^\ast \sF$. Its cohomology sheaves are supported on
$X_s$; by \ref{VII:1-9} applied to $X_s \to Y_s$,
$\eR f_\ast \Delta$ is therefore constructible. The complex
$\eR f_\ast \eR j_\ast j^\ast \sF = \eR(f j)_\ast j^\ast \sF$ is also
constructible, and we conclude from the triangle
\[\xymatrix{
  \ar[r] 
    & \eR f_\ast \sF \ar[r] 
    & \eR f_\ast \eR j_\ast j^\ast \sF \ar[r] 
    & \eR f_\ast \Delta \ar[r] 
    & \text{.} 
}\]










\section{Local biduality}\label{VII:4}






\subsection{}\label{VII:4-1}

Let $S$ be regular of dimension $0$ or $1$, and put $K_S=A$ (the constant
sheaf with value $A$). For $K\in\ob\eD_\text{ctf}^b(S,A)$, put
$D K=\rHom(K,K_S)\in\ob\eD^+(S,A^\circ)$ (a complex of sheaves of right
$A$-modules). We still have $D K\in\ob\eD_\text{ctf}^b(S,A^\circ)$, and an explicit
calculation, possible because $S$ has dimension $1$, shows that
$K\iso D D K$ (for $A=\dZ/n$, \hyperref[V]{Duality} \ref{V:1-4}).





\subsection{}\label{VII:4-2}

We assume that $A$ is commutative, an assumption that is probably
unnecessary. For $a:X\to S$ of finite type over $S$, put
$K_X=\eR a^! K_S$. For $K\in\ob\eD_\text{ctf}^b(X,A)$, put
$D K=\rHom(K,K_X)\in\ob\eD_\text{ctf}^b(X,A)$. Our main result is the
following.





\begin{theorem_}\label{VII:4-3}
We have $K\iso D D K$ ($K\in \ob\eD_\textnormal{ctf}^b(X,A)$).
\end{theorem_}





\begin{lemma_}\label{VII:4-4}
If $a$ is proper, then $\eR a_\ast K \iso \eR a_\ast D D K$.
\end{lemma_}

Poincaré duality
\[
  \eR a_! \rHom(K,\eR a^! K_S) = \rHom(\eR a_! K,K_S) 
\]
gives us an isomorphism $\eR a_\ast D = D \eR a_! = D \eR a_\ast$. The lemma
follows from the commutativity of the diagram (\hyperref[V]{Duality},
\ref{V:1-2})
\[\xymatrix{
  \eR a_\ast K \ar[r] \ar[d]^-\wr 
    & \eR a_\ast D D K \ar[d]^-\wr \\
  D D \eR a_\ast K \ar[r]^-\sim 
    & D \eR a_\ast D K \text{.} 
}\]





\subsection{}\label{VII:4-5}

By localization, it is enough to treat the cases where $S$ is the spectrum
of a field or a trait. When $S$ is the spectrum of a field, we argue by
induction on $\dim X$. Since the problem is local, we may assume that $X$ is
proper. Moreover, the induction hypothesis, by the usual arguments, ensures
that the mapping cylinder $\Delta$ of $K\to D D K$ has skyscraper cohomology
sheaves. The lemma then ensures that $\eR a_\ast \Delta=0$, hence that
$\Delta=0$.





\subsection{}\label{VII:4-6}

When $S$ is a trait, we already know that \ref{VII:4-3} holds on the generic
fiber (\ref{VII:4-5}). We may also assume that $X$ is proper over $S$, and we
argue by induction on $\dim{X_s}$, with $X$ proper over $S$. The mapping
cylinder $\Delta$ again has skyscraper cohomology sheaves on $X_s$, and we
conclude as above.





\subsection{}\label{VII:4-7}

Suppose that $A=\dZ/m$. We then have a variant of the preceding theory in
which $K$ is arbitrary in $\eD_c^b(X,A)$. The complex $\eR a^! \dZ/m$ has
finite injective dimension \cite[XVIII 3.1.7]{sga4}. This ensures that $D K$
still belongs to $\eD_c^b(X,A)$. To prove local biduality over $S$, we use the
fact that $\dZ/m$ is the dualizing $\dZ/m$-module. We then proceed as above.










\section{Appendix, by L.\ Illusie}\label{VII:5}

In this appendix, which takes up certain parts of the late
\cite[II]{sga5}, and a letter from P.\ Deligne to the author, we prove
\ref{VII:2-14}, as well as various generalizations and variants of the
specialization theorems \cite[XVI 2.1]{sga4} and
(\hyperref[I]{Étale cohomology} \ref{I:5-1-7}).

The schemes and morphisms considered will be assumed quasi-compact and
quasi-separated. If $X$ is a scheme and $x$ a geometric point of $X$, we
denote by $X(x)$ the strict localization of $X$ at $x$.

Fix a base scheme $S$ and a sheaf of rings $A$ on $S$ (not necessarily
commutative or noetherian). If $X$ is an $S$-scheme, we write $\eD(X)$ for
$\eD(X,A_X)$.





\subsection{Cohomological properness}\label{VII:5-1}






\begin{proposition}\label{VII:5-1-1}
Let $f:X\to Y$ be an $S$-morphism and let $E\in\ob \eD^+(X)$. The following
conditions are equivalent:
\begin{enumerate}[\indent (i)]
  \item The formation of $\eR f_\ast E$ commutes with every finite base change
    $S'\to S$.
  \item The formation of $\eR f_\ast E$ commutes with every quasi-finite base
    change (or inverse limit of quasi-finite morphisms) $S'\to S$.
  \item For every specialization arrow $i:t\to S(s)$
    (\hyperref[I]{Étale cohomology}, \ref{I:5-1-2}), if $\bar S$ denotes
    the normalization in $k(\bar t)$ of the integral scheme that is the
    closure of $i(t)$ in $S(s)$ (by
    \cite[IV, 6.15.5, 18.6.11, 18.8.16]{ega}, $\bar S$ is therefore local and
    integral, with closed point radicial over $s$), if
    $\bar f:\bar X\to \bar Y$ is the morphism deduced from $f$ by the base
    change $\bar S\to S$, and if $f_s:X_s\to Y_s$ is the fiber of $f$ at $s$,
    then the base-change arrow
    \[
      (\eR \bar f_\ast (E|\bar X))|Y_s \to \eR {f_s}_\ast(E|X_s) 
    \]
    is an isomorphism.
\end{enumerate}
\end{proposition}

The equivalence of (i) and (ii) follows from Zariski's Main Theorem, and it is
clear that (ii) implies (iii). Let us prove that (iii) implies (ii). For every
$S'\to S$ that is an inverse limit of quasi-finite morphisms, and every
geometric point $s$ of $S'$, denote by $C(S',s)$ the cone of the base-change
arrow
\[
  (\eR f'_\ast(E|X'))|Y_s \to \eR{f_s}_\ast (E|X_s) \text{,} 
\]
where $f':X'\to Y'$ is deduced from $f$ by the base change $S'\to S$, and
$f_s:X_s\to Y_s$ is the fiber of $f$ at $s$. We must prove that, for every
$(S',s)$ as above, $C(S',s)$ is acyclic. We shall prove by induction on $N$ that
$\h^i C(S',s) = 0$ for $i\leqslant N$. By the Main Theorem, passage to the
limit, and elimination of radicial morphisms, we may assume that $S$ and $S'$
are strictly local, that $S'$ is integral with closed point $s$, and that
$S'\to S$ is finite. Let $t$ be a geometric generic point of $S'$, and let
$\bar S$ be the normalization of $S'$ in $k(t)$. Let $\bar S_\bullet$ be the
simplicial scheme that is the coskeleton of $\bar S\to S'$
($\bar S_0=\bar S$,\ldots $\bar S_n=(\bar S/S')^{n+1}$,\ldots), and let
$\bar X_\bullet=X'\times_{S'} \bar S_\bullet$. Let
$g_\bullet:\bar X_\bullet \to X'$ be the canonical projection. Put
$E|X'=E'$. The bicomplex ${g_\bullet}_\ast g_\bullet^\ast E'$, whose columns
are the ${g_n}_\ast g_n^\ast E'$, is a resolution of $E'$
\cite[VIII 8]{sga4}. We may therefore represent $C(S',s)$ by a bicomplex whose
columns identify with the $C(\bar S_n,s)$. By (iii),
$C(\bar S_0,s)=C(\bar S,s)$ is acyclic. On the other hand, for $n>0$, we have
$\h^i C(\bar S_n,s)=0$ for $i\leqslant N$ by the induction hypothesis. It
follows that $\h^i C(S',s)=0$ for $i\leqslant N+1$, which completes the proof.





\subsubsection{Remark}\label{VII:5-1-2}

Let $(f,E)$ satisfy condition (i) of \ref{VII:5-1-1}.
\begin{enumerate}[a)]
  \item If $E\in \ob \eD^b(X)$ and $f_\ast$ has finite cohomological
    dimension, then for every $F\in \ob\eD^-(S,A^\circ)$, the canonical arrow
    \begin{equation*}\tag{$*$}\label{VII:eq:5-1-1-1}
      \eR f_\ast(E)\lotimes_A q^\ast F \to \eR f_\ast(E\lotimes_A p^\ast F) \text{,} 
    \end{equation*}
    where $p:X\to S$ and $q:Y\to S$ are the projections, is an isomorphism.
    The same conclusion holds with $E\in\ob\eD^-(X)$.
\end{enumerate} 

Indeed, by dévissage, it is enough to verify the assertion for $F$ of the
form $i_! A_{S'}$, with $i:S'\to S$ quasi-finite. By the Main Theorem, we
reduce to the case where $i$ is finite, and the conclusion is then equivalent
to the commutation of $\eR f_\ast(E)$ with base change by $i$.
\begin{enumerate}[a)]
\setcounter{enumi}{1}
  \item Suppose that $A$ is noetherian and constructible. Then, for every
    constructible $F\in\ob\eD^b(S,A^\circ)$ of finite tor-dimension, the arrow
    \eqref{VII:eq:5-1-1-1} above is an isomorphism.
\end{enumerate}

Indeed, by dévissage, we reduce to the case where $F$ is of the form $i_! M$,
with $i:S'\to S$ locally closed, $A$ locally constant on $S'$, and $M$
constructible, of finite tor-dimension, and with locally constant cohomology
on $S'$. After base change from $S$ to $S'$, we are thus reduced to the case
where $F$ is perfect \cite[I]{sga6}, hence finally, by localization and
dévissage, to the case $F=A$, for which the conclusion is immediate.





\subsubsection{}\label{VII:5-1-3}

Let $f:X\to Y$ be an $S$-morphism and let $E\in\ob\eD^+(X)$. We say that
$(f,E)$ is \emph{cohomologically proper} rel.\ to $S$ if the formation of
$\eR f_\ast E$ commutes with every base change $S'\to S$. Equivalently, after
every base change $S'\to S$, the formation of $\eR f_\ast' E'$ (where
$f'=f\times_S S'$ and $E'=$ the inverse image of $E$ on $X\times_S S'$)
commutes with every finite base change $S'' \to S'$ (cf.
\cite[XII 6.1]{sga4}). Here are some examples of cohomological properness
(for a study of the analogous notion for sheaves of sets or of
noncommutative groups, the reader may consult \cite[XIII]{sga1}).
\begin{enumerate} % originally 1.3.1, 1.3.2, ...
  \item Suppose that $A$ is torsion and $f$ is proper. Then, for every
    $E\in\ob\eD^+(X)$, $(f,E)$ is cohomologically proper rel.\ to $S$
    (proper base-change theorem, \cite[XII 5.1]{sga4}).
  \item Suppose that $S$ is finite, that $X$ and $Y$ are of finite type over
    $S$, and that $A$ is annihilated by an integer invertible on $S$. Then,
    for every $E\in\ob\eD^+(X)$, $(f,E)$ is cohomologically proper rel.\ to
    $S$ (\ref{VII:1-9}).
  \item Suppose that $A$ is constant, noetherian, and annihilated by an
    integer $n$ invertible on $S$, and that
    $f:X=Y\setminus D\to Y$ is the inclusion of the complement of a divisor
    $D$ on $Y$ with normal crossings relative to $S$
    \cite[XIII 2.1]{sga1}. Let $E$ be a sheaf of $A$-modules on $X$ satisfying
    one of the following conditions:
    \begin{enumerate}[(i)]
      \item $E$ is constructible locally constant and tamely ramified along
        $D$;
      \item $E$ is, locally for the étale topology on $Y$ and on $S$, the
        inverse image of a sheaf of $A$-modules on $S$.
    \end{enumerate}
    Then $(f,E)$ is cohomologically proper rel.\ to $S$, and $\eR f_\ast E$
    is constructible (compare \cite[XIII 2.4]{sga1}).
\end{enumerate} 

Let us sketch the proof in case (i). The question is local near a point $y$
of $Y$. We may assume that $D$ is a sum of smooth divisors,
$D=\sum_{i=1}^m D_i$. By the relative Abhyankar lemma
\cite[XIII 5.5]{sga1}, near $y$ we may trivialize $E$ by a finite covering
$\widetilde Y\to Y$ of the form
$\widetilde Y=Y[T_1,\dots,T_r]/(T_1^{n_1}-t_1,\dots,T_r^{n_r}-t_r)$, where the
$t_i$ are local equations of the smooth divisors passing through $y$, and the
$n_i$ are integers prime to the char.\ of $k(y)$. The inverse image
$\widetilde D$ of $D$ in $\widetilde Y$ still has normal crossings rel.\ to $S$, and
if $g:\widetilde X=\widetilde Y\setminus \widetilde D\to X$ is the projection,
$g^\ast E$ is constant. Since $E$ injects into $g_\ast g^\ast E$ and the
quotient is tamely ramified, an easy dévissage reduces us to the case where
$E$ is constant. Denote the projections by $p:X\to S$, $q:Y\to S$, and, for
$1\leqslant i\leqslant m$, let $f_i:Y\setminus D_i \to Y$ be the canonical
inclusion. From the relative purity theorem \cite[XVI 3.7]{sga4}, one easily
deduces by induction on $m$ that, for every locally constant $A$-module $M$ on
$S$, the canonical arrow
\begin{equation*}\tag{5.1.3.1}\label{VII:eq:5-1-3-3-1} % originally 1.3.3.1
  \eR f_\ast (\dZ/n)\lotimes q^\ast M\to \eR f_\ast (p^\ast M) 
\end{equation*}
is an isomorphism, and that on the other hand we have:
\begin{align*}\tag{5.1.3.2}\label{VII:eq:5-1-3-3-2}
  {f_i}_\ast(\dZ/n) &= \dZ/n \text{,} \\
  \eR^1 {f_i}_\ast (\dZ/n) &= (\dZ/n)(-1)_{D_i} \text{,} 
\end{align*}
where a canonical isomorphism is given by the fundamental class $\cl(D_i)$
(\hyperref[IV]{Cycle} \ref{IV:2-1-4}),
\begin{align*}
  \eR^q {f_i}_\ast(\dZ/n) &= 0\text{ for $q>1$,} \\
  \eR f_\ast (\dZ/n) &= \bigotimes_i^\eL \eR {f_i}_\ast (\dZ/n) \text{ , i.e.:} \\
  \eR^1 f_\ast(\dZ/n) &= \bigoplus_{i=1}^m (\dZ/n)(-1)_{D_i} \text{,} \\
  \textstyle\bigwedge^\bullet \eR^1 f_\ast(\dZ/n) &\iso \eR^\bullet f_\ast (\dZ/n) \text{.} 
\end{align*}
The conclusion of 3, in case (i), therefore follows from
\eqref{VII:eq:5-1-3-3-1} and \eqref{VII:eq:5-1-3-3-2}. In case (ii), by
passage to the limit we reduce to the case where $E=p^\ast M$, with $M$
constructible, and then, by dévissage and base change (using (i)), to the
already treated case where $M$ is constant.

It follows from the preceding proof that \eqref{VII:eq:5-1-3-3-1} is an
isomorphism for every $M\in\ob\eD^+(S)$.

One may compare the formulas \eqref{VII:eq:5-1-3-3-2} with the analogous
formulas in integral cohomology when $S=\spec(\dC)$ (see, for example,
\cite[3.1]{de71-2}), which follow from the fact that, locally for the classical
topology, the complement of $D$ has the homotopy type of a torus.

By passage to the limit, one deduces from \eqref{VII:eq:5-1-3-3-2} similar
formulas in $\ell$-adic cohomology, i.e. with $\dZ/n$ replaced by $\dZ_\ell$
($\ell$ a prime invertible on $S$). Suppose that $q:Y\to S$ is proper and
smooth. It then follows from the Weil conjectures \cite{de74,de80} that the
Leray spectral sequence
\[
  E_2^{i j} = \eR^i q_\ast \eR^j f_\ast \dZ_\ell \Rightarrow \eR^\bullet p_\ast \dZ_\ell 
\]
degenerates at $E_3$ modulo torsion \cite[\S 6]{de71-1}. The same phenomenon
occurs in integral cohomology when $S=\spec(\dC)$; see (loc.\ cit.) and
\cite[3.2.13]{de71-2}.





\subsection{Specialization and cospecialization}\label{VII:5-2}

Let $f:X\to S$ and $E\in\ob\eD^+(X)$.






\subsubsection{}\label{VII:5-2-1}

Suppose that the formation of $\eR f_\ast E$ commutes with every finite base
change. For every specialization $u:t\to S(s)$ of geometric points of $S$,
we then define an arrow, called the \emph{specialization arrow},
\begin{equation*}\tag{5.2.1.1}\label{VII:eq:5-2-1-1}
  \spe(u):\eR\Gamma(X_s,E|X_s) \to \eR\Gamma(X_t,E|X_t) \text{,}
\end{equation*}
as follows. Let, as in \ref{VII:5-1-1}(iii), $\bar S$ be the normalization in
$k(t)$ of the closure of $u(t)$ in $S(s)$, and let
$\bar f:\bar X\to \bar S$ be the morphism deduced from $f$ by the base change
$\bar S\to S$. By the immediate implication (i)$\Rightarrow$(iii) of
\ref{VII:5-1-1}, the base-change arrow
\begin{equation*}\tag{5.2.1.2}\label{VII:eq:5-2-1-2}
  \eR\Gamma(\bar X,E|\bar X)=\eR\bar f_\ast (E|\bar X)_s \to \eR\Gamma(X_s,E|X_s) 
\end{equation*}
is an isomorphism. We define \eqref{VII:eq:5-2-1-1} as the composite of the
inverse of \eqref{VII:eq:5-2-1-2} with the restriction arrow
\begin{equation*}\tag{5.2.1.3}\label{VII:eq:5-2-1-3}
  \eR\Gamma(\bar X,E|\bar X) \to \eR\Gamma(X_t,E|X_t) \text{.}
\end{equation*}






\subsubsection{}\label{VII:5-2-2}

Suppose that $f$ is locally acyclic rel.\ to $E$ (\ref{VII:2-12}). For
every specialization $u:t\to S(s)$ of geometric points of $S$, we then define
an arrow, called the \emph{cospecialization arrow},
\begin{equation*}\tag{5.2.2.1}\label{VII:eq:5-2-2-1}
  \cosp(u):\eR\Gamma(X_t,E|X_t) \to \eR\Gamma(X_s,E|X_s) \text{,} 
\end{equation*}
as follows. Let $g:X_t\to X$ and $g':\bar X\to X$ be the canonical arrows.
Viewed vis-\`a-vis these two maps, the canonical arrow
\begin{equation*}\tag{5.2.2.2}\label{VII:eq:5-2-2-2}
  \eR g'_\ast(E|\bar X)\to \eR g_\ast(E|X_t) 
\end{equation*}
is an isomorphism by the local-acyclicity hypothesis. Indeed, we may assume
that $S$ is strictly local and integral, with closed point $s$, and that $t$
is a geometric generic point. Compute the fiber of
\eqref{VII:eq:5-2-2-2} at a geometric point $x$ of $X$, with image $y$ in
$S$. After replacing $S$ by the strict localization of $S$ at $y$, we may
assume that $y=s$. Since $\bar S\to S$ is an inverse limit of finite
morphisms and is radicial at $s$, we find that
$\eR g'_\ast(E|\bar X)_x = E_x$ and that the fiber of
\eqref{VII:eq:5-2-2-2} at $x$ identifies with the restriction
$E_x=\eR\Gamma(X(x),E)\to \eR\Gamma(X(x)_t,E|X(x)_t)$, which is an
isomorphism by hypothesis. Applying $\eR\Gamma(X,-)$ to
\eqref{VII:eq:5-2-2-2}, we deduce that the restriction arrow
\eqref{VII:eq:5-2-1-3} is an isomorphism. We define
\eqref{VII:eq:5-2-2-1} as the composite of the inverse of
\eqref{VII:eq:5-2-1-3} with the base-change arrow
\eqref{VII:eq:5-2-1-2}.






\subsubsection{}\label{VII:5-2-3}

If $s'' \xrightarrow{u'} s' \xrightarrow u s$ are specialization arrows of
geometric points of $S$, one verifies that
\begin{align*}
  \spe(u')\spe(u) &= \spe(u u') \text{,} \\
  \cosp(u) \cosp(u') &= \cosp(u u') \text{,} 
\end{align*}
whenever both sides are defined.






\subsubsection{}\label{VII:5-2-4}

Suppose that the formation of $\eR f_\ast E$ commutes with every finite base
change and that $f$ is locally acyclic rel.\ to $E$. It follows from the
definitions that, for every specialization $u:t\to S(s)$ of geometric points
of $S$, the arrows $\spe(u)$ and $\cosp(u)$ are isomorphisms, inverse to one
another.

The preceding hypotheses hold in particular when $A$ is noetherian and
annihilated by an integer $n$ invertible on $S$, when $f$ is proper and smooth,
and when $E$ has locally constant cohomology (cohomological properness of
proper morphisms (\ref{VII:5-1-3}), and local acyclicity of smooth morphisms
\cite[XV]{sga4}). This recovers the ``specialization theorem''
\cite[XVI 2.1]{sga4}.

On the other hand, \ref{VII:5-1-1} has the following consequence:






\begin{proposition}\label{VII:5-2-5}
Suppose that $f$ is locally acyclic relative to $E$ and that, for every
specialization $u:t\to S(s)$, the cospecialization arrow $\cosp(u)$ is an
isomorphism. Then the formation of $\eR f_\ast E$ commutes with every finite
base change.
\end{proposition}

In view of \ref{VII:5-2-4}, \ref{VII:5-2-5} generalizes
(\hyperref[I]{Étale cohomology} \ref{I:5-1-7}).






\begin{corollary}\label{VII:5-2-6}
Suppose that $f$ is locally acyclic rel.\ to $E$. Let $x$ be a geometric
point of $X$, with image $s$ in $S$, and denote by
$f_{(x)}:X(x)\to S(s)$ the morphism induced by $f$. Then the formation of
$\eR {f_{(x)}}_\ast(E|X(x))$ commutes with every finite base change
$S'\to S(s)$.
\end{corollary}

By \ref{VII:5-2-5}, it is enough to verify that the cospecialization arrows
for $(f_{(x)},E|X(x))$ are isomorphisms. By the transitivity of the
cospecialization arrows (\ref{VII:5-2-3}), it is enough to show that if
$u:t\to S(s)$ is a specialization, then
$\cosp(u):\eR\Gamma(X(x)_t,E|X(x)_t)\to \eR\Gamma(X(x)_s,E|X(x)_s)$ is an
isomorphism. But since $s$ is closed in $S(s)$, $X(x)_s$ is strictly local
with closed point $x$, $\eR\Gamma(X(x)_s,E|X(x)_s)=E_x$, and $\cosp(u)$ is an
isomorphism, inverse to the restriction isomorphism
$E_x\to \eR\Gamma(X(x)_t,E|X(x)_t)$.






\begin{corollary}\label{VII:5-2-7}
Let $f:X\to Y$ be an $S$-morphism, let $p:X\to S$, $q:Y\to S$ be the
projections, and let $E\in\ob\eD^+(X)$. Suppose that $A$ is locally constant,
that $f$ is locally acyclic rel.\ to $E$, and that $q$ is locally acyclic at
every geometric point $y$ of $Y$ rel.\ to every sheaf locally constant
in a neighborhood of $y$. Then $p$ is locally acyclic relative to $E$.
\end{corollary}

Let $x$ be a geometric point of $X$, with images $y$ and $s$ in $Y$ and $S$,
let $t\to S(s)$ be a specialization, and let
$f_{(x)}:X(x)\to Y(y)$ be the morphism induced by $f$. We must show that the
restriction
\begin{equation*}\tag{$*$}\label{VII:eq:5-2-7-1}
  E_x \to \eR\Gamma(X(x)_t,E|X(x)_t) 
\end{equation*}
is an isomorphism. We may assume that $A$ is constant. By \ref{VII:5-2-4} and
\ref{VII:5-2-5}, for every specialization $z\to Y(y)$, the specialization
arrow
\[
  E_x = \eR {f_{(x)}}_\ast (E|X(x))_y \to (\eR {f_{(x)}}_\ast (E|X(x)))_z = \eR\Gamma(X(x)_z,E|X(x)_z) 
\]
is an isomorphism; in other words, $\eR{f_{(x)}}_\ast(E|X(x))$ is constant
with value $E_x$. The acyclicity hypothesis on $q$ therefore plays the
decisive r\^ole and ensures that the restriction
\[
  E_x \to \eR\Gamma(Y(y)_t,(E_x)|Y(y)_t) = \eR\Gamma\left(Y(y)_t,\eR {f_{(x)}}_\ast (E|X(x))|Y(y)_t\right)
\]
is an isomorphism. But this arrow identifies with
\eqref{VII:eq:5-2-7-1}, which proves the assertion.

The hypothesis of \ref{VII:5-2-7} on $q$ holds in particular when $q$ is
smooth and $A$ is annihilated by an integer $n$ invertible on $S$
\cite[XV]{sga4}. On the other hand, if the hypotheses of \ref{VII:5-2-7} hold
after every base change $S'\to S$, then it follows that $p$ is universally
locally acyclic rel.\ to $E$. These remarks imply, in particular, the
statement of \ref{VII:2-14}.






\begin{corollary}\label{VII:5-2-8}
Let $f:X\to S$ and $E\in \ob\eD^+(X)$. Suppose that $f$ is universally locally
acyclic rel.\ to $E$, that the formation of $\eR f_\ast E$ commutes with
every finite base change, and that, for every morphism $s'\to s$, where $s$ is
a geometric point of $S$ and $s'$ is the spectrum of a separably closed
field, the base-change arrow
$\eR\Gamma(X_s,E|X_s)\to \eR\Gamma(X_{s'},E|X_{s'})$ is an isomorphism. Then
$(f,E)$ is cohomologically proper (\ref{VII:5-1-3}).
\end{corollary}

Let $f':X'\to S'$ be the morphism deduced from $f$ by a base change
$S'\to S$; we must show that the formation of $\eR f'_\ast(E|X')$ commutes
with every finite base change. By \ref{VII:5-2-5}, it is enough to show that,
for every specialization $u':t'\to S'(s')$, the cospecialization arrow
$\cosp(u')$:
$\eR\Gamma(X_{t'}',E|X_{t'}') \to \eR\Gamma(X_{s'}',E|X_{s'}')$ is an
isomorphism. Let $s$ (resp.\ $t$) be the geometric point that is the image of
$s'$ (resp.\ $t'$) in $S$, and let $u:t\to S(s)$ be the
specialization defined by $u'$. One readily sees that there is a commutative
square
\[\xymatrix{
  \eR\Gamma(X_t,E|X_t) \ar[r]^-{\cosp(u)} \ar[d] 
    & \eR\Gamma(X_s,E|X_s) \ar[d] \\
  \eR\Gamma(X_{t'}',E|X_{t'}') \ar[r]^-{\cosp(u')} 
    & \eR\Gamma(X_{s'}',E|X_{s'}') 
}\]
in which the vertical arrows are the base-change arrows. By hypothesis, these
are isomorphisms. On the other hand, (\ref{VII:5-2-4}) $\cosp(u)$ is an
isomorphism, hence $\cosp(u')$ is an isomorphism, which completes the proof.

Here is an example of an application of \ref{VII:5-2-8}. Suppose that $A$ is
annihilated by an integer $n$ invertible on $S$, and that $f$ is locally of
finite type and universally locally acyclic. Let $x$ be a geometric point of
$X$, with image $s$ in $S$, and let $f_{(x)}:X(x)\to S(s)$ be the morphism
induced by $f$. It follows from \ref{VII:5-2-6} and \ref{VII:1-9} (by passage
to the limit) that the hypotheses of \ref{VII:5-2-8} are satisfied by the
pair $(f_{(x)},E|X(x))$, which is therefore cohomologically proper.






\subsubsection{}\label{VII:5-2-9}

Suppose that $A$ is constant. Let $f:X\to S$, and let $E\in\ob\eD^+(X)$ be a
complex of finite tor-dimension. We say that $f$ is \emph{strongly locally
acyclic} rel.\ to $E$ if, for every geometric point $x$ of $X$, with image
$s$ in $S$, every specialization $t\to S(s)$, and every $A^\circ$-module $M$,
the restriction arrow
\[
  E_x\lotimes_A M \to \eR\Gamma(X(x)_t,(E|X(x)_t)\lotimes_A M) 
\]
is an isomorphism. I do not know whether ``locally acyclic'' implies
``strongly locally acyclic.''






\begin{proposition}\label{VII:5-2-10}
Under the hypotheses of \ref{VII:5-2-9}, suppose that $A$ is noetherian and
torsion, that $S$ is noetherian, and that $f$ is strongly locally acyclic
rel.\ to $E$. Then, for every diagram with cartesian squares
\[\xymatrix{
  X \ar[d]^-f 
    & \ar[l]_-j X' \ar[d]^-{f'} 
    & \ar[l]_-{j'} X'' \ar[d]^-{f''} \\
  S 
    & \ar[l]_-i S' 
    & \ar[l]_-{i'} S'' \text{,} 
}\]
with $i$ quasi-finite and $i'$ an open immersion, and every
$F\in\ob\eD^+(S'',A^\circ)$, the canonical arrow
\begin{equation*}\tag{$*$}\label{VII:eq:5-2-10-1}
  j^\ast E \lotimes {f'}^\ast \eR i'_\ast F \to \eR j'_\ast ((j j')^\ast E\lotimes {f''}^\ast F) 
\end{equation*}
is an isomorphism.
\end{proposition}

Compare this statement with \cite[XV 1.17]{sga4}.

Let us prove \ref{VII:5-2-10}. We readily reduce to the case where $S=S'$ is
strictly local, and to $F$ constructible and concentrated in degree $0$. We then
resolve $F$ on the right by finite sums of sheaves of the form $u_\ast M$,
where $u:t\to S$ is a geometric point and $M$ is an $A^\circ$-module. The
hypothesis, vis-\`a-vis this resolution, implies that
\eqref{VII:eq:5-2-10-1} is an isomorphism for $F=u_\ast M$, and the result
follows.

% END INCLUDED components/sga4half_7_en.tex

% BEGIN INCLUDED components/sga4half_8_en.tex
% !TEX root = ../SGA4half_English_Reader_1_to_8.tex

\chapter{Derived categories}\label{VIII}
\blfootnote{by Jean-Louis Verdier}

\section{Triangulated categories: definitions and examples}\label{VIII:1}

\subsection{Definition of triangulated categories}\label{VIII:1-1}

\addtocounter{subsubsection}{-1}
\subsubsection{}\label{VIII:1-1-0}

Let $\cA$ be an additive category and let $T$ be an additive automorphism of
$\cA$. For two objects $X$ and $Y$ of $\cA$, we set
\[
  \hom^i(X,Y) = \hom_\cA(X,T^i(Y)) \qquad i\in \dZ
\]
Let $X$, $Y$, and $Z$ be three objects of $\cA$, let
$\alpha\in \hom^i(X,Y)$, and let $\beta\in \hom^j(Y,Z)$. We define the
composite $\beta\circ\alpha$ belonging to $\hom^{i+j}(X,Z)$ by
\[
  \beta\circ \alpha = T^i(\beta)\circ\alpha \text{.}
\]
One checks that this gives a new category whose groups of morphisms between
objects are graded. By an abuse of language, this new category will be called
the category $\cA$ graded by the translation functor $T$.

Let $\cA$ be a category graded by a translation functor $T$. Whenever a
morphism is named or denoted without its degree being specified, it will
always mean a morphism of degree zero.

Let $\cA$ and $\cA'$ be two additive categories graded by the functors $T$
and $T'$, and let $F$ be an additive functor from $\cA$ to $\cA'$. We say
that $F$ \emph{is graded} if an isomorphism of
functors\footnote{Text amended in 1976: the original text required an
equality $F\circ T=T'\circ F$. The case of functors in several variables is
discussed in \cite[XVIII 0.2]{sga4}. It raises a problem of signs.}
has been specified:
\[
  F\circ T = T'\circ F
\]

The definition of morphisms of graded functors is given in
\ref{VIII:4-1-1}.

Let $\cA$ be an additive category graded by the functor $T$. A triangle is
a collection of three objects $X$, $Y$, $Z$ and three morphisms
$u:X\to Y$, $v:Y\to Z$, $w:Z\to T(X)$ ($\deg w=1$). We use either the
notation $(X,Y,Z,u,v,w)$ or the following diagram for a triangle:
\[\xymatrix{
  & Z \ar[dl]_-w \\
  X \ar[rr]^-u
    & & Y \ar[lu]_-v
}\qquad \deg(w)=1
\]
Let $(X,Y,Z,u,v,w)$ and $(X',Y',Z',u',v',w')$ be two triangles of $\cA$.
A morphism from the first triangle to the second is a collection of three
morphisms $f:X\to X'$, $g:Y\to Y'$, $h:Z\to Z'$ such that the following
diagram is commutative:
\[\xymatrix{
  X \ar[r]^-u \ar[d]^-f
    & Y \ar[r]^-v \ar[d]^-g
    & Z \ar[r]^-w \ar[d]^-h
    & T(X) \ar[d]^-{T(f)} \\
  X' \ar[r]^-{u'}
    & Y' \ar[r]^-{v'}
    & Z' \ar[r]^-{w'}
    & T(X')
}\]

\subsubsection{}\label{VIII:1-1-1}

A \emph{triangulated category} is an additive category graded by a
translation functor and equipped with a family of triangles, called the
family of distinguished triangles. This family must also satisfy the
following axioms.

\paragraph{TR1}
\phantomsection\hypertarget{VIII:TR1}{}

Every triangle isomorphic to a distinguished triangle is distinguished.
Every morphism $u:X\to Y$ is contained in a distinguished triangle
$(X,Y,Z,u,v,w)$. The triangle
$(X,X,0,\operatorname{id}_X,0,0)$ is distinguished.

\paragraph{TR2}
\phantomsection\hypertarget{VIII:TR2}{}

A triangle $(X,Y,Z,u,v,w)$ is distinguished if and only if
$(Y,Z,T(X),v,w,-T(u))$ is distinguished.

\paragraph{TR3}
\phantomsection\hypertarget{VIII:TR3}{}

If $(X,Y,Z,u,v,w)$ and $(X',Y',Z',u',v',w')$ are distinguished, then for
every morphism $(f,g):u\to u'$, there exists a morphism $h:Z\to Z'$ such
that $(f,g,h)$ is a morphism of triangles.

\paragraph{TR4}
\phantomsection\hypertarget{VIII:TR4}{}

\[\xymatrix{
  & Z' \ar[dl]|{\deg=1} \\
  X \ar[rr]^-u
    & & Y \ar[ul]_-i
}\quad
\xymatrix{
  & X' \ar[dl]|{\deg(j)=1} \\
  Y \ar[rr]^-v
    & & Z \ar[ul]
}\quad
\xymatrix{
  & Y' \ar[dl]|{\deg=1} \\
  X \ar[rr]^-w
    & & Z \ar[ul]
}\]
If these are three distinguished triangles such that $w=v\circ u$, then
there exist two morphisms $f:Z'\to Y'$, $g:Y' \to X'$ such that
\begin{enumerate}
  \item $(\operatorname{id}_X, v, f)$ is a morphism of triangles;
  \item $(u,\operatorname{id}_Z,g)$ is a morphism of triangles;
  \item $(Z',Y',X',f,g,T(i) j)$ is a distinguished triangle.
\end{enumerate}

If $\cC$ and $\cC'$ are two triangulated categories, an \emph{exact
functor} from $\cC$ to $\cC'$ is an additive graded functor that takes
distinguished triangles to distinguished triangles. Morphisms of exact
functors are morphisms of graded functors.

The following property follows immediately from axioms
\hyperlink{VIII:TR1}{TR1}, \hyperlink{VIII:TR2}{TR2}, and
\hyperlink{VIII:TR3}{TR3}.

\begin{proposition}\label{VIII:1-1-2}
Let $(X,Y,Z,u,v,w)$ be a distinguished triangle and let $M$ be an object
of a triangulated category. Then the following sequences are exact:
\[\xymatrix{
  \cdots \ar[r]
    & \hom^i(M,X) \ar[r]^-{\hom(M,u)}
    & \hom^i(M,Y) \ar[r]^-{\hom(M,v)}
    & \hom^i(M,Z) \ar[r]^-{\hom(M,w)}
    & \hom^{i+1}(M,X) \ar[r]
    & \cdots
}\]
\[\xymatrix{
  \cdots \ar[r]
    & \hom^i(Z,M) \ar[r]^-{\hom(v,M)}
    & \hom^i(Y,M) \ar[r]^-{\hom(u,M)}
    & \hom^i(X,M) \ar[r]^-{\hom(w,M)}
    & \hom^{i+1}(Z,M) \ar[r]
    & \cdots
}\]
\end{proposition}

From this proposition one obtains statements of the following kind:
\begin{itemize}
  \item If $(f,g,h)$ is a morphism of distinguished triangles and $f,g$
    are isomorphisms, then $h$ is an isomorphism.
  \item The distinguished triangles constructed on a morphism
    (\hyperlink{VIII:TR1}{TR1}) are all isomorphic (\emph{but the
    isomorphisms are not, in general, uniquely determined}).
  \item When one of the morphisms in a distinguished triangle has a kernel
    or a cokernel, the latter splits; that is, the triangle is of the form
    \[\xymatrix{
      & Z+T(X) \ar[dl] \\
      X+Y \ar[rr]
        & & Y+Z \ar[ul]
    }\]
  \item The sum of two distinguished triangles is distinguished. (One uses
    axiom \hyperlink{VIII:TR4}{TR4}.)
\end{itemize}

\subsection{Examples}\label{VIII:1-2}

\subsubsection{}\label{VIII:1-2-1}

We first fix some notation. Let $\cA$ be an additive category.
The following category will be denoted by $\eC(\cA)$:
\begin{itemize}
  \item The objects of $\eC(\cA)$ are the complexes of $\cA$, with no
    restriction on degree and with differential of degree $+1$.
  \item The morphisms of $\eC(\cA)$ are the degree-preserving morphisms of
    complexes (those commuting with the differential).
\end{itemize}

Let $T$ be the following functor. For every object $X^\bullet$ of
$\eC(\cA)$,
\begin{align*}
  T(X^\bullet)_i &= (X^\bullet)_{i+1} \\
  d_i(T(X^\bullet)) &= -d_{i+1}(X^\bullet)  \text{.}
\end{align*}
On morphisms the functor $T$ acts as follows: for every morphism
$f:X^\bullet \to Y^\bullet$, the morphism
$T(f):T(X^\bullet) \to T(Y^\bullet)$ is defined by
\[
  T(f)_i = f_{i+1} \text{.}
\]
It is immediate that $T$ is an additive automorphism of $\eC(\cA)$.
The morphisms of degree $n$ defined using this translation functor are
then the morphisms that raise degree by $n$ and commute or anticommute
with the differential according to the parity of $n$. Thus one recovers
the usual definition.

Let $\sC$ denote the following family of triangles. A triangle
$(X^\bullet,Y^\bullet,Z^\bullet,u,v,w)$ belongs to the family
if\footnote{Text amended in 1976 to ensure the validity of
\ref{VIII:1-2-4}.}
\begin{itemize}
  \item $Z^\bullet$ is the simple complex associated with the double
    complex
    \[\xymatrix{
      \cdots \ar[r]
        & 0 \ar[r]
        & X^\bullet \ar[r]^-u
        & Y^\bullet \ar[r]
        & 0 \ar[r]
        & \cdots
    }\]
    in which the objects of $Y^\bullet$ have first degree zero and the
    objects of $X^\bullet$ have first degree $-1$.
  \item $v$ is the image, under the functor: double complex $\mapsto$
    simple complex, of the following morphism of double complexes:
    \[\xymatrix{
      \cdots \ar[r]
        & 0 \ar[r] \ar[d]
        & 0^\bullet \ar[r] \ar[d]
        & Y^\bullet \ar[r] \ar[d]
        & 0 \ar[r] \ar[d]
        & \cdots \\
      \cdots \ar[r]
        & 0 \ar[r]
        & X^\bullet \ar[r]
        & Y^\bullet \ar[r]
        & 0 \ar[r]
        & \cdots
    }\]
  \item $w$ is the negative of the image, under the same functor, of the
    following morphism of double complexes:
    \[\xymatrix{
      \cdots \ar[r]
        & 0 \ar[r]
        & X^\bullet \ar[r]
        & 0 \ar[r]
        & 0 \ar[r]
        & \cdots \\
      \cdots  \ar[r]
        & 0 \ar[r] \ar[u]
        & X^\bullet \ar[r] \ar[u]
        & Y^\bullet \ar[r] \ar[u]
        & 0 \ar[r] \ar[u]
        & \cdots
    }\]
\end{itemize}

\subsubsection{}\label{VIII:1-2-2}

The set of null-homotopic morphisms is a two-sided ideal (if $f$ and $g$
are composable and either $f$ or $g$ is null-homotopic, then the composite
is null-homotopic; if $f$ and $g$ are null-homotopic, then
$f\oplus g$ is null-homotopic). For every pair of objects $X^\bullet$,
$Y^\bullet$, let $\operatorname{Htp}(X^\bullet,Y^\bullet)$ denote the
subgroup of $\hom_{\eC(\cA)}(X^\bullet,Y^\bullet)$ consisting of the
null-homotopic morphisms. The objects of $\eK(\cA)$ are then the objects
of $\eC(\cA)$, and its group of morphisms with source $X^\bullet$ and
target $Y^\bullet$ is
\[
  \hom_{\eC(\cA)}(X^\bullet,Y^\bullet) / \operatorname{Htp}(X^\bullet,Y^\bullet) \text{.}
\]
The composition law on the morphisms of $\eK(\cA)$ is defined by passing
the composition law on the morphisms of $\eC(\cA)$ to the quotient.
The category $\eK(\cA)$ is additive. The translation functor $T$ on
$\eC(\cA)$ passes to the quotient (the translate of a null-homotopic
morphism is null-homotopic). The resulting functor is an additive
automorphism of $\eK(\cA)$, again denoted by $T$. By an abuse of notation,
$\eK(\cA)$ will henceforth denote the category $\eK(\cA)$ graded by the
functor $T$.

A distinguished triangle in $\eK(\cA)$ is any triangle \emph{isomorphic}
to a triangle arising from $\sC$. The family of distinguished triangles
satisfies axioms \hyperlink{VIII:TR1}{TR1},
\hyperlink{VIII:TR2}{TR2}, \hyperlink{VIII:TR3}{TR3}, and
\hyperlink{VIII:TR4}{TR4}. By an abuse of notation, $\eK(\cA)$ will denote
the category $\eK(\cA)$ triangulated by the family of triangles defined
above.

\subsubsection{}\label{VIII:1-2-3}

There is also a whole arsenal of subcategories:
\begin{itemize}
  \item A complex is said to be bounded below if all but finitely many of
    its objects of negative degree are zero. Bounded-above and bounded
    complexes are defined similarly. Let $\eK^+(\cA)$, $\eK^-(\cA)$,
    and $\eK^b(\cA)$ denote the full subcategories of $\eK(\cA)$ generated
    by these families of objects. These subcategories are
    “\emph{triangulated subcategories}” of $\eK(\cA)$: every distinguished
    triangle two of whose objects belong to the subcategory is isomorphic
    to a triangle all three of whose objects belong to the subcategory.
  \item Suppose that $\cA$ is an abelian category. A complex is said to be
    cohomologically bounded below if all but finitely many of its
    cohomology objects of negative degree are zero. Cohomologically
    bounded-above complexes, cohomologically bounded complexes, and
    acyclic complexes are defined similarly. Let
    $\eK^{\infty,+}(\cA)$, $\eK^{\infty,-}(\cA)$,
    $\eK^{\infty,b}(\cA)$, and $\eK^{\infty,\varnothing}(\cA)$ denote the
    full subcategories generated by these families of objects. They are
    triangulated subcategories.

    We shall also use the triangulated subcategories $\eK^{+,b}(\cA)$,
    $\eK^{+,\varnothing}(\cA)$, \ldots etc \ldots (the first superscript
    gives information about the objects of the complexes, and the second
    about the cohomology objects).
  \item Without assuming that $\cA$ is abelian, let $O$ be a set of objects
    of $\cA$ stable under isomorphisms and direct sums. Let $O(\cA)$ denote
    the full subcategory of $\eC(\cA)$ generated by the complexes whose
    objects in every degree belong to $O$. Let $\underline O(\cA)$ denote
    the corresponding triangulated subcategory of $\eK(\cA)$. When $\cA$
    is abelian, one may take for $O$ either the set $I$ of injectives or
    the set $P$ of projectives.
\end{itemize}

\subsubsection{}\label{VIII:1-2-4}

The category $\eK(\cA)$ is uniquely determined by the data of $\eC(\cA)$
and the family $\sC$. More precisely, every graded additive functor from
$\eC(\cA)$ to a triangulated category that takes every triangle of $\sC$
to a distinguished triangle factors uniquely through $\eK(\cA)$, and the
resulting functor is exact.

For every complex $X^\bullet$ of $\eC(\cA)$, let $\bar X^\bullet$ denote
the complex obtained from $X^\bullet$ by setting the differential equal
to zero. A three-term sequence
\[\xymatrix{
  0 \ar[r]
    & X^\bullet \ar[r]
    & Y^\bullet \ar[r]
    & Z^\bullet \ar[r]
    & 0
}\]
is called a semi-split sequence if the sequence
\[\xymatrix{
  0 \ar[r]
    & \bar X^\bullet \ar[r]
    & \bar Y^\bullet \ar[r]
    & \bar Z^\bullet \ar[r]
    & 0
}\]
is split; that is, if it is isomorphic to the sequence
\[\xymatrix{
  0 \ar[r]
    & \bar X^\bullet \ar[r]
    & \bar X^\bullet + \bar Z^\bullet \ar[r]
    & \bar Z^\bullet \ar[r]
    & 0 \text{.}
}\]

Let $0 \to X^\bullet \to Y^\bullet \to Z^\bullet \to 0$ be a semi-split
sequence. Choose a splitting, that is, for every $n$ an isomorphism
\[
  (Y^\bullet)^n \simeq (X^\bullet)^n + (Z^\bullet)^n \text{.}
\]
The matrix of the differential of $Y^\bullet$ in the above direct-sum
decomposition has the form
\[
  \begin{pmatrix}
    d_{X^\bullet} & \varphi_n \\
    0 & d_{Z^\bullet}
  \end{pmatrix}
\]
where $\varphi_n$ is a morphism from $(Z^\bullet)^n$ to
$(X^\bullet)^{n+1}$. The $\varphi_n$ determine a morphism of complexes
\[
  \varphi:Z^\bullet \to T(X^\bullet) \text{.}
\]
When the splitting is changed, the homotopy class of $\varphi$ does not
change. Moreover, if $\underline u$, $\underline v$, and
$\underline\varphi$ denote the images in $\eK(\cA)$ of the morphisms $u$,
$v$, and $\varphi$, then the triangle
$(X^\bullet,Y^\bullet,Z^\bullet,\underline u,\underline v,\underline\varphi)$
is distinguished. Let $\mathsf{SSS}(\cA)$ denote the category of
semi-split sequences of $\eC(\cA)$ and let $\mathsf{TrK}(\cA)$ denote the
category of distinguished triangles of $\eK(\cA)$. The preceding
construction defines a functor
\[
  \rho:\mathsf{SSS}(\cA) \to \mathsf{TrK}(\cA) \text{.}
\]
This functor is essentially surjective.

\subsection{Examples of cohomological and exact functors}\label{VIII:1-3}

\begin{definition}
Let $\cA$ be a triangulated category and let $\cB$ be an abelian category.
An additive functor $F:\cA\to \cB$ is called a \emph{cohomological
functor} if, for every distinguished triangle $(X,Y,Z,u,v,w)$, the
sequence
\[\xymatrix{
  F(X) \ar[r]^-{F(u)}
    & F(Y) \ar[r]^-{F(v)}
    & F(Z)
}\]
is exact.
\end{definition}

The functor $F\circ T^i$ will often be denoted by $F^i$. By axiom
\hyperlink{VIII:TR2}{TR2} for triangulated categories, there is an
unbounded exact sequence
\[\xymatrix{
  \cdots \ar[r]
    & F^i(X) \ar[r]
    & F^i(Y) \ar[r]
    & F^i(Z) \ar[r]
    & F^{i+1}(X) \ar[r]
    & \cdots
}\]

\subsubsection{Examples}\label{VIII:1-3-2}

\begin{enumerate}
  \item Let $\cA$ be a triangulated category. For every object $X$ of
    $\cA$, the functor $\hom_\cA(X,-)$ is cohomological. The functor
    $\hom_\cA(-,X)$ is a cohomological functor
    $\cA^\circ\to\mathsf{Ab}$.
  \item Let $\cA$ be an abelian category and let $\eC(\cA)$ be the category
    of complexes of $\cA$. Define the functor
    $\h^0:\eC(\cA)\to\cA$ as follows. For an object $X^\bullet$ of
    $\eC(\cA)$, set
    \[
      \h^0(X^\bullet) = \ker d^0 / \im{d^{-1}} \text{.}
    \]
    The functor $\h^0$ annihilates null-homotopic morphisms and therefore
    factors uniquely through $\eK(\cA)$. The resulting functor will again
    be denoted by $\h^0:\eK(\cA) \to \cA$.
\end{enumerate}

\subsubsection{Example of an exact bifunctor}\label{VIII:1-3-3}

Let $\cA$, $\cA'$, and $\cA''$ be three additive categories, and let
\[
  P:\cA\times \cA'\to \cA''
\]
be a bilinear functor (that is, additive in each argument separately).
It induces the bilinear functor
\[
  P^\ast:\eC(\cA)\times \eC(\cA') \to \eC(\cA'')
\]
as follows.

Let $X^\bullet$ be an object of $\eC(\cA)$ and let $Y^\bullet$ be an
object of $\eC(\cA')$. Then $P(X^\bullet,Y^\bullet)$ is a double complex
of $\cA''$. Set $P^\ast(X^\bullet,Y^\bullet) = $ the simple complex
associated with $P(X^\bullet,Y^\bullet)$.

Let $f$ be a null-homotopic morphism of $\eC(\cA)$ (respectively
$\eC(\cA')$), and let $Z^\bullet$ be an object of $\eC(\cA')$
(respectively $\eC(\cA)$). The morphism $P^\ast(f,Z^\bullet)$
(respectively $P^\ast(Z^\bullet,f)$) is then null-homotopic. It follows
that $P^\ast$ uniquely defines a functor
\[
  P^\bullet:\eK(\cA) \times \eK(\cA') \to \eK(\cA'') \text{.}
\]
The bifunctor $P^\bullet$ is exact.

In particular, let $\cA$ be an additive category. For $P$ one may take
the functor
\begin{align*}
  \cA^\circ\times \cA &\to \mathsf{Ab} \\
  (X,Y) &\mapsto \hom(X,Y)
\end{align*}

The preceding construction then gives a functor
\[
  \hom^\bullet:\eK(\cA)^\circ \times \eK(\cA) \to \eK(\mathsf{Ab})
\]
which, when composed with the functor
$\h^0:\eK(\mathsf{Ab})\to \mathsf{Ab}$, plainly gives back the functor
$\hom_{\eK(\cA)}$.

\section{Quotient categories}\label{VIII:2}

\subsection{Thick subcategories. Multiplicative systems of morphisms}\label{VIII:2-1}

We begin with two definitions.

\begin{definition}\label{VIII:2-1-1}
A subcategory $\cB$ of a triangulated category $\cA$ is said to be
\emph{thick} if $\cB$ is a full triangulated subcategory of $\cA$ and,
in addition, $\cB$ has the following property:

For every morphism $f:X\to Y$ that factors through an object of $\cB$ and
is contained in a distinguished triangle $(X,Y,Z,f,g,h)$ in which $Z$ is
an object of $\cB$, the source of $f$ and the target of $f$ are objects of
$\cB$.
\end{definition}

The set of thick subcategories of $\cA$ will be denoted by $\sN$.

The intersection of an arbitrary family of thick subcategories is a thick
subcategory.

\begin{definition}
Let $\cA$ be a triangulated category. A set $S$ of morphisms of $\cA$ is
called a \emph{multiplicative system of morphisms} if it has the following
five properties.

\paragraph{FR1}
\phantomsection\hypertarget{VIII:FR1}{}
If $f,g\in S$ and $f$ and $g$ are composable, then $f\circ g\in S$. For
every object $X$ of $\cA$, the identity morphism of $X$ belongs to $S$.

\paragraph{FR2}
\phantomsection\hypertarget{VIII:FR2}{}
In the category $\cA$, every diagram
\[\xymatrix{
  & Y \ar[d]^-{s\in S} \\
  Z \ar[r]^-f
    & X
}\]
can be completed to a commutative diagram
\[\xymatrix{
  P \ar[r]^-g \ar[d]^-{t\in S}
    & Y \ar[d]^-{s\in S} \\
  Z \ar[r]^-f
    & X
}\]
Moreover, the symmetric property holds.

\paragraph{FR3}
\phantomsection\hypertarget{VIII:FR3}{}
If $f$ and $g$ are morphisms, the following properties are equivalent:
\begin{enumerate}[\indent i)]
  \item there exists $s\in S$ such that $s\circ f=s\circ g$;
  \item there exists $t\in S$ such that $f\circ t=g\circ t$.
\end{enumerate}

\paragraph{FR4}
\phantomsection\hypertarget{VIII:FR4}{}
Let $T$ be the translation functor of $\cA$. For every element $s$ of $S$,
one has $T(s)\in S$.

\paragraph{FR5}
\phantomsection\hypertarget{VIII:FR5}{}
Let $(X,Y,Z,u,v,w)$ and $(X',Y',Z',u',v',w')$ be two distinguished
triangles, and let $f$ and $g$ be two elements of $S$ such that the pair
$(f,g)$ is a morphism from $u$ to $u'$. There exists a morphism $h\in S$
such that $(f,g,h)$ is a morphism of triangles.
\end{definition}

A multiplicative system of morphisms is said to be \emph{saturated} if it
has the following property:

A morphism $f$ belongs to $S$ if and only if there exist two morphisms $g$
and $g'$ such that $g\circ f\in S$ and $f\circ g'\in S$.

The intersection of an arbitrary family of saturated multiplicative
systems is a saturated multiplicative system.

The set of saturated multiplicative systems will be denoted by $\sS$.

Let $\cB$ be a thick subcategory. Denote by $\varphi(\cB)$ the set of
morphisms $f$ contained in a distinguished triangle $(X,Y,Z,f,g,h)$ in
which $Z$ is an object of $\cB$. Then $\varphi(\cB)$ is a saturated
multiplicative system.

Let $S$ be a saturated multiplicative system. Denote by $\psi(S)$ the
full subcategory generated by the objects $Z$ contained in a distinguished
triangle $(X,Y,Z,f,g,h)$ in which $f$ is an element of $S$. The
subcategory $\psi(S)$ is thick.

The main result of this subsection is the following: $\varphi$ is an
order-preserving isomorphism from $\sN$ onto $\sS$, where the order
relation is defined by inclusion. Its inverse is precisely $\psi$.

\subsection{Universal problems}\label{VIII:2-2}

Let $\cA$ and $\cA'$ be two triangulated categories and let $F$ be an exact
functor from $\cA$ to $\cA'$. Let $S(F)$ be the set of morphisms of $\cA$
that $F$ takes to isomorphisms, and let $\cB(F)$ be the full subcategory
generated by the objects of $\cA$ that $F$ takes to zero objects of
$\cA'$. Then $S(F)$ is a saturated multiplicative system and $\cB(F)$ is
a thick subcategory. Moreover, $S(F) = \varphi(\cB(F))$. Now let $\cA$ be
a triangulated category, let $\cB$ be a thick subcategory, and let
$S=\varphi(\cB)$ be the corresponding saturated multiplicative system.
The following two universal problems are equivalent.

\paragraph{Problem 1}
\phantomsection\hypertarget{VIII:prob1}{}
Find a triangulated category $\cA/\cB$ and an exact functor
$Q:\cA\to \cA/\cB$ such that every exact functor from $\cA$ to a
triangulated category $\cA'$ that takes the objects of $\cB$ to zero
objects of $\cA'$ admits a unique factorization of the form $G\circ Q$,
where $G$ is exact.

\paragraph{Problem 2}
\phantomsection\hypertarget{VIII:prob2}{}
Find a triangulated category $\cA_S$ and an exact functor
$Q:\cA\to \cA_S$ such that every exact functor from $\cA$ to a
triangulated category $\cA'$ that takes the morphisms of $S$ to
isomorphisms of $\cA'$ admits a unique factorization of the form
$G\circ Q$, where $G$ is exact.

(The author asks the reader to excuse this use of obsolete and reactionary
language.) In the next subsection we shall show that
\hyperlink{VIII:prob2}{Problem 2} has a solution. Consequently, the
corresponding \hyperlink{VIII:prob1}{Problem 1} has the same solution.

\subsection{Calculus of fractions}\label{VIII:2-3}

\subsubsection{}\label{VIII:2-3-1}

Let $\cA$ be a triangulated category and let $S$ be a saturated
multiplicative system. Denote by $\cA[S^{-1}]$ the category of fractions
of $\cA$ with respect to the set of morphisms $S$, and by $Q$ the
canonical functor from $\cA$ to $\cA[S^{-1}]$. The pair
$(\cA[S^{-1}],Q)$ solves the following problem: every functor from $\cA$
to an arbitrary category (not necessarily additive) that takes morphisms
of $S$ to isomorphisms factors uniquely through $(\cA[S^{-1}],Q)$. Such
a pair always exists, without any hypothesis on the set of morphisms $S$
\cite[1.1]{gz67}\footnote{The original text references [C.G.G] here. Most
likely, this is the lost work \emph{Catégories et foncteurs} by Chevalley,
Gabriel and Grothendieck, referenced in\cite{sc72}. }. However, because the
set of morphisms $S$ is a saturated multiplicative system, it has
properties \hyperlink{VIII:FR1}{FR1}, \hyperlink{VIII:FR2}{FR2}, and
\hyperlink{VIII:FR3}{FR3}. Consequently, by \cite{gz67}, the category
$\cA[S^{-1}]$ can be obtained by a right or left calculus of fractions.
We recall how $\cA[S^{-1}]$ is obtained by a right calculus of fractions.
(The left calculus of fractions is obtained by reversing the arrows.)

\subsubsection{}\label{VIII:2-3-2}

The objects of $\cA[S^{-1}]$ are the objects of $\cA$.
\begin{itemize}
  \item For every object $X$ of $\cA$, let $S_X$ denote the category of
    arrows belonging to $S$ and having target $X$. To every object $Y$ of
    $\cA$ one associates, functorially in $Y$, the following set-valued
    functor on $S_X^\circ$ ($S_X^\circ$ denotes the opposite category):
    \[
      H_Y:s\mapsto \hom(\operatorname{source}(s),Y) \text{.}
    \]
    Set
    \[
      \hom_{\cA[S^{-1}]}(X,Y) = \varinjlim_{S_X^\circ} H_Y \text{.}
    \]
    For all matters concerning direct limits, we refer to \cite{ar62}.

    One then observes with pleasure, using \hyperlink{VIII:FR1}{FR1},
    \hyperlink{VIII:FR2}{FR2}, and \hyperlink{VIII:FR3}{FR3}, that the
    category $S_X^\circ$ has properties L1, L2, and L3. Direct limits
    therefore have all the good properties of direct limits over filtered
    ordered sets.
  \item Let $X$, $Y$, and $Z$ be three objects of $\cA$, let $a$ be an
    element of $\hom_{\cA[S^{-1}]}(X,Y)$, and let $b$ be an element of
    $\hom_{\cA[S^{-1}]}(Y,Z)$. Let $s$ be an object of $S_X$, let
    $\underline a$ be an element of
    $\hom(\operatorname{source}(s),Y)$ whose image is $a$, and let an object
    of $S_Y$, $\underline b$ be an element of
    $\hom(\operatorname{source}(t),Z)$ whose image is $b$. There is then a
    diagram
    \begin{equation*}\tag{1}\label{VIII:eq:2-3-2-1}
    \xymatrix{
      & \bullet \ar[dl]_-s \ar[dr]^-{\underline a}
        &
        & \bullet \ar[dl]_-t \ar[dr]^-{\underline b} \\
      X
        &
        & Y
        &
        & Z
    }
    \end{equation*}
    which, by \hyperlink{VIII:FR2}{FR2}, can be completed to a diagram
    \[\xymatrix{
      & & \bullet \ar[dl]_-{t'} \ar[dr]^-{\underline c} \\
      & \bullet \ar[dl]_-s \ar[dr]^-{\underline a}
        & & \bullet \ar[dl]_-t \ar[dr]^-{\underline b} \\
      X
        & & Y
        & & Z
    }\]
    Let $b\circ a$ denote the image in $\hom_{\cA[S^{-1}]}(X,Z)$ of the
    morphism $\underline{b\circ c}$. Using properties
    \hyperlink{VIII:FR1}{FR1}, \hyperlink{VIII:FR2}{FR2}, and
    \hyperlink{VIII:FR3}{FR3}, one checks that $b\circ a$ is independent
    of the chosen representatives $\underline a$ and $\underline b$, and
    also of the way in which diagram \eqref{VIII:eq:2-3-2-1} is completed.
    One further checks that this defines a category $\cA[S^{-1}]$.
\end{itemize}

We denote by $Q$ the evident functor $\cA\to \cA[S^{-1}]$.

Since $\cA$ is additive, $\cA[S^{-1}]$ is an additive category. Moreover,
using \hyperlink{VIII:FR4}{FR4}, one proves that there exists a unique
translation functor on $\cA[S^{-1}]$, denoted by $T$, satisfying
\[
  Q\circ T = T\circ Q \text{.}
\]
Finally, using \hyperlink{VIII:FR5}{FR5}, one proves that there exists a
unique triangulated structure on $\cA[S^{-1}]$ for which $Q$ is exact.
The distinguished triangles of $\cA[S^{-1}]$ are precisely the triangles
isomorphic to the images under $Q$ of the distinguished triangles of
$\cA$. Denoting by $\cA_S$ the category $\cA[S^{-1}]$ equipped with this
triangulated structure, one proves without difficulty that the pair
$(\cA_S,Q)$ is a solution to \hyperlink{VIII:prob2}{Problem 2}.

\begin{definition}\label{VIII:2-3-3}
Let $\cA$ be a triangulated category, let $\cB$ be a thick subcategory of
$\cA$, and let $(Q,\cA/\cB)$ be the solution to
\hyperlink{VIII:prob1}{Problem 1} (\S \ref{VIII:2-2}). The category
$\cA/\cB$ is called the \emph{quotient category} of $\cA$ by $\cB$, and
$Q$ is called the canonical quotient functor. More generally, let $\cN$
be a triangulated subcategory of $\cA$, and let $\underline\cN$ be the
smallest thick subcategory containing $\cN$. By definition, the quotient
category of $\cA$ by $\cN$, denoted $\cA/\cN$, is
$\cA/\underline\cN$.
\end{definition}

\subsection{Properties of quotient categories}\label{VIII:2-4}

\subsubsection{}\label{VIII:2-4-1}

Let $\cA$ be a triangulated category and let $H$ be a cohomological functor
with values in an abelian category $\cG$. Let $\cB_H$ be the full
subcategory generated by the objects all of whose translates are taken by
$H$ to zero objects of $\cG$, and let $S_H$ be the set of morphisms all of
whose translates are taken by $H$ to isomorphisms of $\cG$. Then $\cB_H$
is a thick subcategory of $\cA$, and $S_H$ is a saturated multiplicative
system. Moreover, $S_H=\varphi(\cB_H)$ (\S \ref{VIII:2-1}). The functor
$H$ factors uniquely through $(Q,\cA/\cB_H)$.

\begin{theorem}\label{VIII:2-4-2}
Let $\cA$ be a triangulated category, let $\cB$ be a triangulated
subcategory, let $\cN$ be a thick subcategory of $\cA$, and let
$S=\varphi(\cN)$ be the corresponding saturated multiplicative system.
\begin{enumerate}[(a)]
  \item The category $\cN\cap\cB$ is a thick subcategory of $\cB$. The
    corresponding multiplicative system is $S\cap \cB$.
  \item The following two properties are equivalent:
    \begin{enumerate}[(i)]
      \item For every object $X$ of $\cB$ and every morphism $s:R\to X$,
        where $s$ belongs to $S$, there exists a morphism $s':R'\to R$ such
        that $s\circ s'$ belongs to $S\cap \cB$; and
        $R'\in\operatorname{Ob}\cB$.
      \item Every morphism from an object $X$ of $\cB$ to an object $Y$ of
        $\cN$ factors through an object of $\cN\cap \cB$.
    \end{enumerate}
  \item The properties (i)' and (ii)' obtained from (i) and (ii) by
    reversing the arrows are equivalent.
  \item If properties (i) and (ii) hold, or if properties (i)' and (ii)'
    hold, the canonical functor
    \[\xymatrix{
      \cB/\cN\cap \cB \ar[r] & \cA/\cN
    }\]
    is faithful.

    Since this functor is injective on objects, it realizes
    $\cB/\cN\cap\cB$ as a subcategory of $\cA/\cN$.
  \item If, in addition, $\cB$ is a full subcategory of $\cA$, then
    $\cB/\cN\cap \cB$ is a full subcategory of $\cA/\cN$.
\end{enumerate}
\end{theorem}

\begin{corollary}\label{VIII:2-4-3}
Let $\cA\subset \cB\subset \cC$ be three triangulated categories, with
$\cA$ a thick subcategory of $\cB$ and $\cB$ a thick subcategory of
$\cC$. Then $\cA$ is a thick subcategory of $\cC$, and $\cB/\cA$ is a
thick subcategory of $\cC/\cA$. The canonical functor
$\cC/\cB \to (\cC/\cA)/(\cB/\cA)$ is an isomorphism.
\end{corollary}

\subsection{Properties of the quotient functor}\label{VIII:2-5}

Let $\cA$ be a triangulated category, let $\cB$ be a thick subcategory,
let $S$ be the corresponding multiplicative system, and let
$Q:\cA\to \cA/\cB$ be the canonical quotient functor.

\begin{proposition}\label{VIII:2-5-1}
For objects $X$ and $Y$ of $\cA$, the following assertions are equivalent:
\begin{enumerate}[(a)]
  \item Every morphism $f:X\to Y$ for which there exists a morphism
    $s:Z\to X$ in $S$ satisfying $f\circ s=0$ is zero.
  \item Every morphism $f:X\to Y$ for which there exists a morphism
    $s:Y\to Z$, $s\in S$, satisfying $s\circ f=0$ is zero.
  \item Every morphism $f:X\to Y$ that factors through an object of $\cB$
    is zero.
  \item The canonical map
    \[
      \hom_\cA(X,Y) \to \hom_{\cA/\cB}(Q(X),Q(Y))
    \]
    is injective.
\end{enumerate}
\end{proposition}

\begin{proposition}\label{VIII:2-5-2}
Let $X$ and $Y$ be two objects of $\cA$. The following assertions are
equivalent:
\begin{enumerate}[(a)]
  \item Every diagram ($s\in S$)
    \[\xymatrix{
      & Z \ar[dl]_-s \ar[dr]^-f \\
      X
        & & Y
    }\]
    can be completed to a diagram: $(s,t\in S$):
    \[\xymatrix{
      & Z' \ar[d]_-t \\
      & Z \ar[dl]_-s \ar[dr]^-f \\
      X
        & & Y
    }\]
    in which $f\circ t$ factors through $(X,s\circ t)$.
  \item The assertion obtained by reversing the direction of the arrows in
    (a) and interchanging $X$ and $Y$.
  \item Every diagram
    \[\xymatrix{
      X \ar[r]^-f
        & N \ar[r]^-g
        & Y
    }\]
    in which $g\circ f=0$ and $N$ is an object of $\cB$ can be completed
    to a diagram
    \[\xymatrix{
      & N' \ar[d]^-i \\
      X \ar[ur]^-h \ar[r]^-f
        & N \ar[r]^-g
        & Y
    }\]
    in which $f=i\circ h$ and $g\circ i=0$.
  \item The assertion obtained by reversing the arrows in (c) and
    interchanging $X$ and $Y$.
  \item The canonical map
    \[\xymatrix{
      \hom_\cA(X,Y) \ar[r]^-Q
        & \hom_{\cA/\cB}(Q(X),Q(Y))
    }\]
    is surjective.
\end{enumerate}
\end{proposition}

\begin{proposition}\label{VIII:2-5-3}
Let $X$ be an object of $\cA$. The following assertions are equivalent:
\begin{enumerate}[(a)]
  \item For every object $Y$ of $\cA$, the canonical map
    \[\xymatrix{
      \hom_\cA(X,Y) \ar[r]^-Q
        & \hom_{\cA/\cB}(Q(X),Q(Y))
    }\]
    is an isomorphism.
  \item Every morphism from $X$ to an object of $\cB$ is zero.
\end{enumerate}
Similarly, the following two assertions are equivalent:
\begin{enumerate}[(a)']
  \item For every object $Y$ of $\cA$, the canonical map
    \[\xymatrix{
      \hom_\cA(Y,X) \ar[r]^-Q
        & \hom_{\cA/\cB}(Q(Y),Q(X))
    }\]
    is an isomorphism.
  \item Every morphism from an object of $\cB$ to $X$ is zero.
\end{enumerate}
\end{proposition}

\begin{definition}\label{VIII:2-5-4}
Every object having the equivalent properties (a) and (b) of the preceding
proposition is called an \emph{object free on the left over $\cA/\cB$},
or a \emph{left $Q$-free object}. Similarly, every object having the
equivalent properties (a)' and (b)' is called an \emph{object free on the
right over $\cA/\cB$}, or a \emph{right $Q$-free object}. It is determined
up to isomorphism by the knowledge of
$Q(X)\in\operatorname{Ob}{\cA/\cB}$.
\end{definition}

\subsection{Functors adjoint to the quotient functor}\label{VIII:2-6}

\begin{definition}\label{VIII:2-6-1}
Two triangulated subcategories $\cN$ and $\cN'$ of a triangulated category
$\cA$ are said to be \emph{orthogonal} if, for every object $X$ of $\cN$
and every object $Y$ of $\cN'$, one has
\[
  \hom_\cA(X,Y) = 0
\]
The category $\cN'$ is then said to be right orthogonal to $\cN$, and
$\cN$ left orthogonal to $\cN'$.
\end{definition}

\begin{proposition}\label{VIII:2-6-2}
Let $\cN$ be a triangulated subcategory of a triangulated category $\cA$.
The full category $\cN^\bot$ (respectively $^\bot\cN$) generated by the
objects $X$ of $\cA$ such that, for every object $Y$ of $\cN$,
$\hom_\cA(Y,X)=0$ (respectively $\hom_\cA(X,Y)=0$), is a thick subcategory
of $\cA$. By an abuse of language, the category $\cN^\bot$ is called the
right orthogonal of $\cN$.
\end{proposition}

This proposition allows us to restate Proposition \ref{VIII:2-5-3}.

\begin{proposition}\label{VIII:2-6-3}
Let $\cB\subset\cA$ be a thick subcategory of a triangulated category
$\cA$. The full category generated by the objects free on the right over
$\cA/\cB$ (respectively on the left) is precisely the right
(respectively left) orthogonal of $\cB$.
\end{proposition}

\begin{proposition}\label{VIII:2-6-4}
Let $\cA$ be a triangulated category, let $\cB$ be a thick subcategory of
$\cA$, let $\cB^\bot$ be the right orthogonal of $\cB$, let $S(\cB)$ and
$S(\cB^\bot)$ be the corresponding multiplicative systems, and let $Q$
and $Q^\bot$ be the canonical quotient functors. For an object $X$ of
$\cA$, consider the following properties:
\begin{enumerate}[(i)]
  \item $X$ maps by a morphism of $S(\cB)$ to an object of $\cB^\bot$.
  \item $X$ receives, by a morphism of $S(\cB^\bot)$, an object of
    $\cB$.
  \item The category $S_X(\cB)$ of arrows of $S(\cB)$ with source $X$ has
    a final object.
  \item The category $S^X(\cB^\bot)$ of arrows of $S(\cB^\bot)$ with
    target $X$ has an initial object.
  \item The category $\cB/X$ of objects of $\cB$ over $X$ has a final
    object.
  \item The category $\cB^\bot/X$ of objects of $\cB^\bot$ under $X$ has
    an initial object.
\end{enumerate}
Then
\[\xymatrix@=0.5cm{
  & & & & (iv) \\
  (i) \ar@{<=>}[r]
    & (ii) \ar@{<=>}[r]
    & (iii) \ar@{<=>}[r]
    & (v) \ar@{=>}[ur] \ar@{=>}[dr] \\
  & & & & (vi)
}\]

If, in addition, $^\bot(\cB^\bot)=\cB$, then all the properties are
equivalent.
\end{proposition}

Let $\cB$ be a thick subcategory of a triangulated category $\cA$, let
$i$ be the inclusion functor from $\cB$ into $\cA$, let $Q$ be the
quotient functor from $\cA$ to $\cA/\cB$, let $\cB^\bot$ be the right
orthogonal of $\cB$, and let $i^\bot$ and $Q^\bot$ be the corresponding
functors. Proposition \ref{VIII:2-6-4} immediately gives the following
propositions.

\begin{proposition}\label{VIII:2-6-5}
The following two properties are equivalent:
\begin{enumerate}[(i)]
  \item The functor $i$ has a right adjoint.
  \item The functor $Q$ has a right adjoint.
\end{enumerate}
\end{proposition}

The proposition remains true when “right” is replaced by “left.”

\begin{proposition}\label{VIII:2-6-6}
The following two properties are equivalent:
\begin{enumerate}[(i)]
  \item The functor $i$ has a right adjoint.
  \item The functor $i^\bot$ has a left adjoint. The left orthogonal of the
    category $\cB^\bot$ is equal to $\cB$.
\end{enumerate}
\end{proposition}

\begin{proposition}\label{VIII:2-6-7}
Suppose the properties of Propositions \ref{VIII:2-6-5} and
\ref{VIII:2-6-6} hold. Let $i^\ast$ and $Q^\ast$ be the right adjoints of
$i$ and $Q$. Similarly, let $^\ast i^\bot$ and $^\ast Q^\bot$ be the left
adjoints of $i^\bot$ and $Q^\bot$. All these functors are exact. Moreover:

There is a functorial isomorphism
$Q^\ast\circ Q\iso i^\bot\circ ^\ast i^\bot$ such that the following
diagram is commutative:
\[\xymatrix{
  & \operatorname{id} \ar[dl] \ar[dr] \\
  Q^\ast \circ Q \ar[rr]^\sim
    & & i^\bot\circ ^\ast i^\bot
}\]
(the oblique arrows are defined by the adjunction properties).

Similarly, there is an isomorphism
$^\ast Q^\bot \circ Q^\bot \iso i\circ i^\ast$ such that the following
diagram is commutative:
\[\xymatrix{
  ^\ast Q^\bot\circ Q^\bot \ar[rr]^-\sim \ar[dr]
    & & i\circ i^\ast \ar[dl] \\
  & \operatorname{id}{}
}\]

Finally, there is a functorial morphism of degree $1$,
$\partial:i^\bot\circ ^\ast i^\bot \to i\circ i^\ast$, such that the
following triangle is distinguished $(\deg\partial=1$):
\[\xymatrix{
  i^\bot \circ ^\ast i^\bot \ar[rr]^-\partial
    & & i\circ i^\ast \ar[dl] \\
  & \operatorname{id}{} \ar[ul]
}\]

Such a morphism $\partial$ is unique and satisfies
$\partial(TX)=-T\partial(X)$ ($X\in\operatorname{Ob}\cA$, $T$ the
translation).
\end{proposition}

\section{Derived categories of an abelian category}\label{VIII:3}

\subsection{Derived categories}\label{VIII:3-1}

We use the notation of \ref{VIII:1-2-3}. We shall also use the following
notation.

\subsubsection{Notation}\label{VIII:3-1-1}

Let $\cA$ be an abelian category. Set
\begin{align*}
  \eD(\cA) &= \eK^{\infty,\infty}(\cA)/\eK^{\infty,\varnothing}(\cA) \\
  \eD^b(\cA) &= \eK^{b,b}(\cA)/\eK^{b,\varnothing}(\cA) \\
  \eD^+(\cA) &= \eK^{+,+}(\cA) / \eK^{+,\varnothing}(\cA) \\
  \eD^-(\cA) &= \eK^{-,-}(\cA) / \eK^{-,\varnothing}(\cA) \text{.}
\end{align*}

\begin{proposition}\label{VIII:3-1-2}
The natural functors appearing in the following diagram
\[\xymatrix@=0.5cm{
  & \eD(\cA) \\
  \eD^+(\cA) \ar[ur]
    & & \eD^-(\cA) \ar[ul] \\
  & \eD^b(\cA) \ar[ul] \ar[ur]
}\]
are fully faithful.

The natural functors
\begin{align*}
  \eD^b(\cA) &= \eK^{b,b}(\cA)/\eK^{b,\varnothing}(\cA) \to \eK^{+,b}(\cA)/\eK^{+,\varnothing}(\cA) \\
  \eD^b(\cA) &= \eK^{b,b}(\cA)/\eK^{b,\varnothing}(\cA) \to \eK^{-,b}(\cA) / \eK^{-,\varnothing}(\cA)
\end{align*}
are equivalences of categories.

Since the functors in the diagram
\[\xymatrix@=0.5cm{
  & \eD(\cA) \\
  \eD^+(\cA) \ar[ur]
    & & \eD^-(\cA) \ar[ul] \\
  & \eD^b(\cA) \ar[ul] \ar[ur]
}\]
are injective on objects, they realize the categories $\eD^b(\cA)$,
$\eD^+(\cA)$, and $\eD^-(\cA)$ as full subcategories of $\eD(\cA)$.
\end{proposition}

The proof of this proposition is contained in Theorem \ref{VIII:2-4-2}.

\begin{definition}\label{VIII:3-1-3}
The category $\eD(\cA)$ is called the \emph{derived category} of $\cA$.
The objects of $\eD(\cA)$ isomorphic to objects of $\eD^b(\cA)$ are
called the \emph{bounded} objects of $\eD(\cA)$. The objects of
$\eD(\cA)$ isomorphic to objects of $\eD^+(\cA)$ are called the
\emph{left-bounded} objects, and the objects of $\eD(\cA)$ isomorphic to
objects of $\eD^-(\cA)$ are called the \emph{right-bounded} objects.
\end{definition}

Let $D$ denote the canonical functor
\[
  D : \eC(\cA) \to \eD(\cA) \text{.}
\]
The “restriction” of $D$ to the full subcategory $\cA$ of $\eC(\cA)$
will again be denoted by $D$. The functor $D$ restricted to $\cA$ is fully
faithful. The functor $D$ factors uniquely through $\eK(\cA)$. Let $X$
and $Y$ be two objects of $\cA$.

One readily checks that, for every $m>0$, the groups
\[
  \hom_{\eD(\cA)}(D(X),T^{-m} D(Y)) \qquad \text{($T$ is the translation functor)}
\]
are zero. The cohomological dimension of $\cA$ is the least integer $n$
such that, for every $m>n$,
\[
  \hom_{\eD(\cA)}(D(X),T^m D(Y)) = 0 \qquad\text{for every pair $X,Y$ of objects of $\cA$.}
\]
If there is no such integer, the cohomological dimension of $\cA$ is said
to be infinite.

Let $\eI^+(\cA)$ (respectively $\eI(\cA)$) denote the full triangulated
subcategory of $\eK^+(\cA)$ (respectively of $\eK(\cA)$) defined by the
complexes whose objects are injective in every degree.

\begin{proposition}\label{VIII:3-1-4}
Suppose that the category $\cA$ has enough injectives. The natural functor
\[
  Q^+:\eK^+(\cA) \to \eD^+(\cA)
\]
induces an equivalence of categories
\[
  \eI^+(\cA) \to \eD^+(\cA)  \text{.}
\]
The quasi-inverse functor is a right adjoint of $Q^+$.
\end{proposition}

If, in addition, $\cA$ has finite cohomological dimension, the preceding
assertion remains true after the superscripts $+$ are omitted.

There is an analogous statement concerning projectives.

The proof of Proposition \ref{VIII:3-1-4} relies on
\ref{VIII:2-6-4}.

The triangulated structure of $\eD(\cA)$ is made more explicit by the
following proposition.

Let $\eE(\cA)$ be the category of three-term exact sequences in
$\eC(\cA)$. The category $\mathsf{SSS}(\cA)$ of semi-split exact
sequences (\S\ref{VIII:1-2-4}) is a full subcategory of $\eE(\cA)$.
We defined (loc. cit.) a functor
\[
  \rho:\mathsf{SSS}(\cA) \to \mathsf{TrK}(\cA)
\]
where $\mathsf{TrK}(\cA)$ denotes the category of distinguished triangles
of $\eK(\cA)$. This gives a functor
\[
  \sigma : \mathsf{SSS}(\cA) \to \mathsf{TrD}(\cA)
\]
where $\mathsf{TrD}(\cA)$ denotes the category of distinguished triangles
of $\eD(\cA)$.

\begin{proposition}\label{VIII:3-1-5}
There exists a unique functor
\[
  \Pi:\eE(\cA) \to \mathsf{TrD}(\cA)
\]
of the following form ($\deg \delta(S)=1$):
\[
(S=0 \to X^\bullet \to Y^\bullet \to Z^\bullet \to 0)
\mapsto
\xymatrix{
  & D(Z^\bullet) \ar[dl]_-{\delta(S)} \\
  D(X^\bullet) \ar[rr]^-{D(u)}
    & & D(Y^\bullet) \ar[ul]_-{D(v)}
}
\]
whose restriction to $\mathsf{SSS}(\cA)$ is $\sigma$. \emph{This functor
is essentially surjective.}
\end{proposition}

\begin{proposition}\label{VIII:3-1-6}
The canonical cohomological functor
\[
  \h^0:\eK(\cA) \to \cA
\]
factors uniquely through a functor that we again denote by $\h^0$,
\[
  \h^0:\eD(\cA) \to \cA \text{.}
\]
On $\eD(\cA)$, the functor $\h^0$ and its translates $\h^i$ form a
conservative system; that is, a morphism $f:X^\bullet \to Y^\bullet$ is
an isomorphism if and only if, for every integer $i$,
\[
  \h^i(f):\h^i(X) \to \h^i(Y)
\]
is an isomorphism.
\end{proposition}

\begin{definition}\label{VIII:3-1-7}
A morphism $f:X^\bullet \to Y^\bullet$ of $\eC(\cA)$ (respectively of
$\eK(\cA)$) is called a \emph{quasi-isomorphism} if, for every integer
$i$,
\[
  \h^i(f):\h^i(X) \to \h^i(Y)
\]
is an isomorphism.
\end{definition}

By the preceding proposition, the quasi-isomorphisms are the morphisms
that become isomorphisms in $\eD(\cA)$.

\begin{definition}\label{VIII:3-1-8}
Let $\sigma$ be a set of objects of $\cA$ and let $X^\bullet$ be an
object of $\eC(\cA)$ (respectively $\eK(\cA)$). A right (respectively
left) resolution of type $\sigma$ of $X^\bullet$ is a quasi-isomorphism
$X^\bullet \to V^\bullet$ (respectively $V^\bullet \to X^\bullet$), where
$V^\bullet$ is a complex all of whose objects belong to $\sigma$. We say
\emph{injective resolution} (respectively projective resolution) instead
of right (respectively left) resolution of injective (respectively
projective) type.
\end{definition}

In general, we shall omit the words “right” or “left” when no ambiguity
can result.

\begin{proposition}\label{VIII:3-1-9}
\begin{enumerate}
  \item Suppose that every object of $\cA$ embeds in an object of $\sigma$
    (respectively is a quotient of an object of $\sigma$). Then every
    object of $\eC^+(\cA)$ (respectively $\eC^-(\cA)$) has a right
    (respectively left) resolution of type $\sigma$ in $\eC^+(\cA)$
    (respectively in $\eC^-(\cA)$).
  \item Suppose, in addition, that there exists an integer $n$ such that
    for every exact sequence
    \begin{align*}
      Y^0 &\to Y^1 \to \cdots \to Y^{n-1} \to Y^n \to 0 \\
      \text{(respectively } 0 &\to Y^n \to Y^{n-1} \to \cdots \to Y^1 \to Y^0 \text{)}
    \end{align*}
    of objects of $\cA$ in which, for every $i$, $0\leqslant i\leqslant n-1$,
    $Y^i$ is an object of $\sigma$, the object $Y^n$ is also an object of
    $\sigma$. Then every object of $\eC(\cA)$ has a right (respectively
    left) resolution of type $\sigma$.
\end{enumerate}
\end{proposition}

\subsection{Study of $\ext$}\label{VIII:3-2}

Let $X^\bullet$ and $Y^\bullet$ be two objects of $\eC(\cA)$
(respectively $\eK(\cA)$), and let
$D:\eC(\cA) \to \eD(\cA)$ (respectively
$Q:\eK(\cA) \to \eD(\cA)$) be the canonical functor.

\begin{definition}\label{VIII:3-2-1}
We call the \emph{$i$th hyper-Ext}, and denote by
\[
  (X^\bullet,Y^\bullet)\mapsto \ext^i(X^\bullet,Y^\bullet)
\]
the functor $\hom_{\eD(\cA)}(D(X^\bullet),T^i D(Y^\bullet))$
(respectively
$\hom_{\eD(\cA)}(Q(X^\bullet),T^i Q(Y^\bullet))$).
\end{definition}

\begin{proposition}\label{VIII:3-2-2}
Let $0 \to X^\bullet \to Y^\bullet \to Z^\bullet \to 0$ be an exact
sequence in $\eC(\cA)$, and let $V^\bullet$ be an object of $\eC(\cA)$.
There are unbounded exact sequences
\[\xymatrix@=0.5cm{
  \cdots \ar[r]
    & \ext^i(V^\bullet,X^\bullet) \ar[r]
    & \ext^i(V^\bullet,Y^\bullet) \ar[r]
    & \ext^i(V^\bullet,Z^\bullet) \ar[r]^-\delta
    & \ext^{i+1}(V^\bullet,X^\bullet) \ar[r]
    & \cdots \\
  \cdots \ar[r]
    & \ext^i(Z^\bullet,V^\bullet) \ar[r]
    & \ext^i(Y^\bullet,V^\bullet) \ar[r]
    & \ext^i(X^\bullet,V^\bullet) \ar[r]^-\delta
    & \ext^{i+1}(Z^\bullet,V^\bullet) \ar[r]
    & \cdots
}\]
\end{proposition}

Let $X^\bullet$ and $Y^\bullet$ be two objects of $\eK(\cA)$. Denote by
$\mathsf{Qis}/X^\bullet$ (respectively $\mathsf{Qis}^+/X^\bullet$,
$\mathsf{Qis}^-/X^\bullet$, $\mathsf{Qis}^b/X^\bullet$) the category of
quasi-isomorphisms with target $X^\bullet$ and source in $\eK(\cA)$
(respectively $\eK^+(\cA)$, $\eK^-(\cA)$, $\eK^b(\cA)$). Similarly,
define the categories $Y^\bullet/\mathsf{Qis}$,
$Y^\bullet/\mathsf{Qis}^+$,\ldots (quasi-isomorphisms with source
$Y^\bullet)$.

The following proposition summarizes the various equivalent definitions
of $\ext^i$.

\begin{proposition}\label{VIII:3-2-3}
1. There are isomorphisms of functors
\[\xymatrix{
  \ext^0(X^\bullet,Y^\bullet) = \hom_{\eD(\cA)}(Q(X^\bullet),Q(Y^\bullet)) \ar[r]^-\sim
    & \varinjlim_{(\mathsf{Qis}/X^\bullet)^\circ} \hom_{\eK(\cA)}(-,Y^\bullet) \ar[r]^-\sim
    & \\
  \varinjlim_{Y^\bullet/\mathsf{Qis}} \hom_{\eK(\cA)}(X^\bullet,-) \ar[r]^-\sim
    & \varinjlim_{(\mathsf{Qis}/X^\bullet)^\circ \times Y^\bullet/\mathsf{Qis}} \hom_{\eK(\cA)}(-,-) \text{.}
}\]
2. If $X^\bullet$ is an object of $\eK^-(\cA)$,
\[
  \ext^0(X^\bullet,Y^\bullet) \iso \varinjlim_{(\mathsf{Qis}^-/X^\bullet)^\circ} \hom_{\eK(\cA)}(-,Y^\bullet)
\]
If $Y^\bullet$ is an object of $\eK^+(\cA)$,
\[
  \ext^0(X^\bullet,Y^\bullet) \iso \varinjlim_{Y^\bullet/\mathsf{Qis}} \hom_{\eK(\cA)}(X^\bullet,-)
\]
3. If $X^\bullet$ is an object of $\eK^-(\cA)$ and $Y^\bullet$ is an object
of $\eK^b(\cA)$,
\[
  \ext^0(X^\bullet,Y^\bullet) \iso \varinjlim_{Y^\bullet/\mathsf{Qis}^b} \hom_{\eK(\cA)}(X^\bullet,-)
\]
If $X^\bullet$ is an object of $\eK^b(\cA)$ and $Y^\bullet$ is an object
of $\eK^+(\cA)$,
\[
  \ext^0(X^\bullet,Y^\bullet) \iso \varinjlim_{(\mathsf{Qis}^b/X^\bullet)^\circ} \hom_{\eK(\cA)}(-,Y^\bullet)
\]
4. If the category $\cA$ has enough injectives and $Y^\bullet$ is an object
of $\eK^+(\cA)$, then $Y^\bullet$ has an injective resolution, unique up
to isomorphism in $\eK^+(\cA)$,
\[
  Y^\bullet \to I(Y^\bullet)
\]
and
\[
  \ext^0(X^\bullet,Y^\bullet) \iso \hom_{\eK(\cA)}(X^\bullet,I(Y^\bullet)) \text{.}
\]
There is an analogous statement for projectives.
5. If the category $\cA$
has enough injectives and if $\cA$ has \emph{finite cohomological dimension}, then
every object $Y^\bullet$ of $\eK(\cA)$ has a unique injective resolution
\[
  Y^\bullet \to I(Y^\bullet)
\]
and
\[
  \ext^0(X^\bullet,Y^\bullet) \iso \hom_{\eK(\cA)}(X^\bullet,I(Y^\bullet))
\]
There is an analogous statement for projectives.
\end{proposition}

\paragraph{Remark:}
On making explicit the isomorphism of Proposition
\ref{VIII:3-2-3}(3), one recovers the Yoneda definition.

\section{Derived functors}\label{VIII:4}

\subsection{Definition of derived functors}\label{VIII:4-1}

\begin{definition}\label{VIII:4-1-1}
Let $\cC$ and $\cC'$ be two graded categories (we denote by $T$ the
translation functor of $\cC$ and $\cC'$), and let $F$ and $G$ be two
graded functors from $\cC$ to $\cC'$. A \emph{morphism of graded functors}
is a morphism of functors
\[
  u:F\to G
\]
having the following property:

For every object $X$ of $\cC$, the following diagram is commutative:
\[\xymatrix{
  F(T X) \ar[r]^-{u(T X)}
    & G(T X) \\
  T F(X) \ar[r]^-{T u(X)} \ar[u]_-\wr
    & T G(X) \ar[u]_-\wr
}\]
\end{definition}

Let $\cC$ and $\cC'$ be two triangulated categories. Denote by
$\mathsf{Fex}(\cC,\cC')$ the category of exact functors from $\cC$ to
$\cC'$, the morphisms between two functors being morphisms of graded
functors.

Let $\cA$ and $\cB$ be two abelian categories and let
$\Phi:\eK^\ast(\cA)\to \eK^{\ast'}(\cB)$ be an exact functor ($\ast$ and
$\ast'$ denote one of the signs $+$, $-$, $b$, or $v$, “empty”). By
composition, the canonical functor
$Q:\eK^\ast(\cA)\to \eD^\ast(\cA)$ gives a functor
\[\xymatrix{
  \mathsf{Fex}(\eD^\ast(\cA),\eD^{\ast'}(\cB)) \ar[r]
    & \mathsf{Fex}(\eK^\ast(\cA),\eD^{\ast'}(\cB))
}\]
and hence (also denoting by $Q'$ the canonical functor
$\eK^{\ast'}(\cB) \to \eD^{\ast'}(\cB)$) a functor $\%$
(respectively $\%'$)
$\mathsf{Fex}(\eD^\ast(\cA),\eD^{\ast'}(\cB)) \to \mathsf{Ab}$:
\begin{align*}
  \Psi  &\mapsto \hom(Q'\circ \Phi,\Psi\circ Q) \\
  (\text{respectively }\Psi &\mapsto \hom(\Psi\circ Q,Q'\circ \Phi)\text{ ).}
\end{align*}

\begin{definition}\label{VIII:4-1-2}
We say that $\Phi$ admits a \emph{total right derived functor}
(respectively total left derived functor) if the functor $\%$
(respectively $\%'$) is representable. An object representing $\%$
(respectively $\%'$) is called the total right (respectively left) derived
functor of $\Phi$ and is denoted by $\eR \Phi$ (respectively
$\eL \Phi$).\footnote{A more manageable definition (and one equivalent in
practice) is given in \cite[XVII 1.2]{sga4}.}
\end{definition}

\subsubsection{Notation}\label{VIII:4-1-3}

Let $\cA$ and $\cB$ be two abelian categories and let $f:\cA\to \cB$ be
an additive functor. The \emph{total right derived functor} of
$\eK^+(f):\eK^+(\cA) \to \eK^+(\cB)$ is denoted by $\eR^+ f$. The
functors $\eR^- f$, $\eR f$, $\eL^+ f$, $\eL^- f$, and $\eL f$ are
defined similarly.

Likewise, let $F:\eK^\ast(\cA) \to \cB$ be a cohomological functor.
By composition, the canonical quotient functor
$Q:\eK^\ast(\cA) \to \eD^\ast(\cA)$ defines a functor
\[\xymatrix{
  \mathsf{Foco}(\eD^\ast(\cA),\cB) \ar[r]
    & \mathsf{Foco}(\eK^\ast(\cA),\cB)
}\]
(where $\mathsf{Foco}(-,-)$ denotes the category of cohomological
functors), and hence a functor $\%$ (respectively $\%'$)
$\mathsf{Foco}(\eD^\ast(\cA),\cB) \to \mathsf{Ab}$ given by
$G\mapsto \hom(F,G\circ Q)$ (respectively
$G\mapsto \hom(G\circ Q,F)$).

\begin{definition}\label{VIII:4-1-4}
We say that the functor $F$ admits a \emph{right cohomological derived
functor} (respectively left cohomological derived functor) if the functor
$\%$ (respectively $\%'$) is representable. An object representing $\%$
(respectively $\%'$) is called the right (respectively left)
cohomological derived functor and is denoted by $\eR F$ (respectively
$\eL F$).
\end{definition}

\subsubsection{Notation}\label{VIII:4-1-5}

Let $\cA$ and $\cB$ be two abelian categories and let $f:\cA\to \cB$ be
an additive functor. The right cohomological derived functor of
\[
  \h^0{}\circ \eK^+ f:\eK^+(\cA) \to \cB
\]
is denoted by $\eR^+ f$. When no confusion can result, the functor
$\eR^+ f\circ T^i$ ($T$ is the translation functor) is denoted by
$\eR^i f$. The right cohomological derived functor of
$\h{}\circ \eK f$ is denoted by $\eR f$. When no confusion can result,
the functor $\eR f\circ T^i$ is denoted by $\eR^i f$.

The functors $\eL^- f$, $\eL f$, and $\eL^i f$ are defined similarly.

\subsubsection{Remarks}\label{VIII:4-1-6}

As we shall see in the next subsection, these definitions include the
classical definition of derived functors when the categories have enough
injectives or projectives and the complexes are suitably bounded. They are
nevertheless more general, but without further hypotheses they are not
very manageable. For example, let $f:\cA\to \cB$ be an additive functor
between two abelian categories. Suppose that the functors $\eR^+ f$ and
$\eR f$ exist. I do not know how to prove that $\eR f$ induces on
$\eD^+(\cA)$ the functor $\eR^+ f$. Suppose, in addition, that the functor
$\eR^+ f$ (cohomological derived) exists. I do not know how to prove that
the functor $\h^0{}\circ \eR^+ f$ is isomorphic to $\eR^+ f$. In practice,
however, these properties will hold.

\subsection{Existence of derived functors}\label{VIII:4-2}

\begin{proposition}\label{VIII:4-2-1}
The hypotheses are those of Definition \ref{VIII:4-1-4}. Suppose, in
addition, that $\cA$ is a $\sU$-category, where $\sU$ is a universe such
that the set $\dZ$ of rational integers belongs to $\sU$, and such that
the set of objects of $\cA$ belongs to $\sU$. Suppose also that $\cB$ is
the category of abelian $\sU$-groups, or more generally the category of
sheaves of abelian $\sU$-groups on an arbitrary $\sU$-site. (A $\sU$-site
is a $\sU$-category whose set of objects belongs to $\sU$, equipped with
a topology.) The functor $F$ admits a right cohomological derived functor.
This functor is calculated as follows:

Let $X$ be an object of $\eD^\ast(\cA)$, and let $\underline X$ be the
object of $\eK^\ast(\cA)$ above $X$. There is a functorial isomorphism
\[
  \eR F(X) = \varinjlim_{\underline X/\mathsf{Qis}^\ast} F(-) \text{.}
\]
\end{proposition}

In particular, let $f:\cA\to \cB$ be an additive functor. The functors
$\eR f$ and $\eR^+ f$ exist and are calculated as above. The functor
$\eR f$ induces on $\eD^+(\cA)$ the functor $\eR^+ f$. Since no confusion
can then result, we shall use the notation $\eR^i f$.

\subsubsection{Remarks}\label{VIII:4-2-2}

Proposition \ref{VIII:4-2-1} can be expressed in a more general form in
the setting of triangulated categories. The hypotheses on the category
$\cB$ needed to ensure the validity of Proposition \ref{VIII:4-2-1} are
hypotheses asserting the existence and exactness of direct limits over
pseudo-filtered $\sU$-categories whose sets of objects belong to $\sU$.
Proposition \ref{VIII:4-2-1} introduces an asymmetry between right and
left derived functors, at least in practice.

\begin{theorem}\label{VIII:4-2-3} % labelled 4.2.2 in original
The data are those of Definition \ref{VIII:4-1-2}. Suppose, in addition,
that a full triangulated subcategory of $\eK^\ast(\cA)$, denoted $\cC$, is given,
and denote by $\cN$ the full triangulated subcategory of the acyclic
objects of $\cC$. Suppose that
\begin{enumerate}[(1)]
  \item Every object of $\cN$ is taken by $\Phi$ to an acyclic object of
    $\eK^{\ast'}(\cB)$.
  \item Every object $X$ of $\eK^\ast(\cA)$ is the source (respectively
    target) of a quasi-isomorphism whose target (respectively source) is
    an object of $\cC$.
\end{enumerate}
\begin{enumerate}[a)]
  \item Then $\Phi$ admits a total right (respectively left) derived
    functor.
  \item The functor $\eR \Phi$ (respectively $\eL \Phi$) is obtained as
    follows:

    The restriction of the functor
    $Q'\circ \Phi:\eK^\ast(\cA) \to \eD^{\ast'}(\cB)$ to $\cC$ vanishes on the
    objects of $\cN$ and therefore factors through $\cC/\cN$. This gives
    a functor $\Phi':\cC/\cN\to \eD^{\ast'}(\cB)$. But the natural
    functor $\cC/\cN\to \eD^\ast(\cA)$ is an equivalence of categories
    (Theorem \ref{VIII:2-4-2}). Composing with its quasi-inverse gives the
    functor $\eR\Phi$ (respectively $\eL \Phi$).
  \item Let $X$ be an object of $\eK^\ast(\cA)$, let $Y$ be an object of
    $\cC$, and let $X\to Y$ (respectively $Y\to X$) be a
    quasi-isomorphism. The object $\eR\Phi\circ Q(X)$ (respectively
    $\eL \Phi\circ Q(X)$) is isomorphic to $Q'\circ \Phi(Y)$.
  \item The functor $X\mapsto \eR\Phi\circ Q(X)$ (respectively
    $X\mapsto \eL \Phi\circ Q(X)$) is isomorphic to the functor
    $X\mapsto \varinjlim_{X/\mathsf{Qis}^\ast} Q'\circ\Phi(-)$
    (respectively
    $X\mapsto \varinjlim_{\mathsf{Qis}^\ast/X} Q'\circ \Phi(-)$).
  \item The functor $\h^0{}\circ \Phi:\eK^\ast(\cA)\to \cB$ admits a
    right (respectively left) cohomological derived functor, which is
    precisely $\h^0{}\circ\eR\Phi$ (respectively
    $\h^0{}\circ \eL\Phi$).
\end{enumerate}
\end{theorem}

\begin{corollary}\label{VIII:4-2-3-1} % originally corollary 1
Suppose that $\ast=+$ (respectively $\ast=-$) in the hypotheses of the
preceding theorem and that $\cA$ has enough injectives (respectively
projectives). Then Theorem \ref{VIII:4-2-3} applies with $\cC$ the
category of complexes of injectives (respectively projectives). Suppose
in addition that $\cA$ has finite cohomological dimension. Then the
theorem applies without restriction on $\ast$, with $\cC$ the category
of complexes of injectives (respectively projectives).
\end{corollary}

\begin{definition}\label{VIII:4-2-5} % originally 2.4
Let $\cA$ and $\cB$ be two abelian categories and let $f:\cA\to \cB$ be
an additive functor.

We say that $f$ has \emph{finite cohomological dimension} on the right
(respectively left) if $\eR^+ f$ (respectively $\eL^- f$)
exists and there is an integer $m\geqslant 0$ such that, for every integer
$n>m$, $\eR^n f$ (respectively $\eL^{-n} f$) is zero on the objects
arising from $\cA$ under the functor $D$. The
\emph{cohomological dimension} on the right (respectively left) of $f$ is then the
least integer $m$ having this property. If $\eR^+ f$ (respectively
$\eL^- f$) exists and $f$ does not have finite cohomological dimension,
we say that $f$ has infinite cohomological dimension.

An object $X$ of $\cA$ is said to be \emph{$f$-acyclic} on the right
(respectively on the left) if, for every integer $n>0$,
$\eR^n f(D(X)) = 0$ (respectively $\eL^{-n} f(D(X))=0$).
\end{definition}

\begin{corollary}\label{VIII:4-2-3-2} % originally corollary 2
The data are those of Definition \ref{VIII:4-2-5}. Suppose there exists a
set $\sigma$ of objects of $\cA$, stable under direct sums, having the
following properties:
\begin{enumerate}
  \item Every acyclic complex
    $\in \operatorname{Ob}{\eK^+(\cA)}$ (respectively
    $\in\operatorname{Ob}{\eK^-(\cA)}$) of objects of $\sigma$ is taken by
    $f$ to an acyclic complex.
  \item Every object $X$ of $\cA$ embeds in (respectively is a quotient of)
    an object of $\sigma$. The hypotheses of Theorem \ref{VIII:4-2-3}
    hold with $\ast=+$ (respectively $\ast=-$) and with $\cC$ the category
    of complexes whose objects in every degree belong to $\sigma$. The
    elements of $\sigma$ are right (respectively left) $f$-acyclic
    objects.
\end{enumerate}

If, in addition, $f$ has finite right (respectively left) cohomological
dimension, the set of right (respectively left) $f$-acyclic objects has,
in addition to properties (1) and (2) above, property (2) of Proposition
\ref{VIII:3-1-9}. Theorem \ref{VIII:4-2-3} then applies to the functor
$\eK f$, with $\cC$ the category of complexes whose objects are right
(respectively left) $f$-acyclic in every degree. The functor $\eR f$
(respectively $\eL f$) exists and induces on $\eD^+(\cA)$ the functor
$\eR^+ f$ (respectively $\eL^+ f$), and on $\eD^-(\cA)$ the functor
$\eR^- f$ (respectively $\eL^- f$).
\end{corollary}

\subsection{Product. Composition}\label{VIII:4-3}

Let $\cA$ be an abelian category, let
$F:\eD(\cA) \to \mathsf{Ab}$ be a cohomological functor, and let
$X^\bullet$ and $Y^\bullet$ be two objects of $\eC(\cA)$. The functor $F$
defines a map
\[
  \ext^0(X^\bullet,Y^\bullet) \to \hom_{\mathsf{Ab}}(F(D(X^\bullet)),F(D(Y^\bullet)))
\]
and hence a pairing
$F(D(X^\bullet))\times \ext^0(X^\bullet,Y^\bullet) \to F(D(Y^\bullet))$.

Applying translations to this pairing gives pairings
\[
  F^i(D(X^\bullet))\times \ext^j(X^\bullet,Y^\bullet) \to F^{i+j}(D(Y^\bullet)) \text{.}
\]
By definition, when these pairings are applied to cohomological derived
functors, they are precisely the Yoneda products. The properties of these
products follow immediately from this definition. We shall not develop
them here.

Let $\cA$, $\cA'$, and $\cA''$ be three abelian categories, and let
$f:\cA\to \cA'$ and $g:\cA'\to \cA''$ be two additive functors. Suppose
that the functors $\eR^+ f$ and $\eR^+ g$ exist. I do not know how to
deduce that the functor $\eR^+ g\circ f$ exists, and even if the latter
functor exists, the formula
\[\xymatrix{
  \eR^+ g\circ f\ar[r]^-\sim
    & \eR^+ g \circ \eR^+ f
}\]
is probably not true in general. Nevertheless, one has the following
proposition.

\begin{proposition}\label{VIII:4-3-1}
Suppose that the category $\cA'$ has enough right $g$-acyclic objects and
that $\cA$ has enough right $f$-acyclic objects that $f$ takes to right
$g$-acyclic objects. Then the functor $\eR^+ g\circ f$ exists and there
is an isomorphism
\[\xymatrix{
  \eR^+ g\circ f \ar[r]^-\sim
    & \eR^+ g \circ \eR^+ f \text{.}
}\]
\end{proposition}

There are analogous statements for left derived functors.

\section{Examples}\label{VIII:6} % originally chapter 2, section 3

\subsection{The functor \texorpdfstring{$\mathsf{Rhom}^\bullet$}{Hom*}}\label{VIII:6-1}

Let $\cA$ be an abelian category with enough injectives.

The functor $Y\mapsto \hom^\bullet(X^\bullet,Y^\bullet)$
(\ref{VIII:1-3}), where $X^\bullet$ is an object of $\eK(\cA)$ and
$Y^\bullet$ an object of $\eK^+(\cA)$, has a total right derived functor
(Corollary \ref{VIII:4-2-3-1}), denoted by
$\rhom^\bullet(X^\bullet,-)$. Let $Y$ be an object of $\eD^+(\cA)$.
The functor
\begin{align*}
  \eK(\cA) &\to \eD(\mathsf{Ab}) \\
  X^\bullet &\mapsto \rhom^\bullet(X^\bullet,Y)
\end{align*}
is exact and vanishes on acyclic complexes, hence gives an exact
bifunctor
\[\xymatrix{
  \eD(\cA)^\circ \times \eD^+(\cA) \ar[r]
    & \eD(\mathsf{Ab})
}\]
which we denote by $\mathsf{Rhom}^\bullet$.

When $\cA$ has finite cohomological dimension, the functor
$\mathsf{Rhom}^\bullet$ extends to
$\eD(\cA)^\circ\times \eD(\cA)$ (loc. cit.).

If, in addition, $\cA$ has enough projectives, the functor $\hom$ can be
derived with respect to the first argument. The two definitions coincide
on their common domain of validity.

\subsection{\texorpdfstring{$\Tor$}{Tor} of sheaves}\label{VIII:6-2}

Let $E$ be a ringed site (in particular, a ringed topological space), let
$\sA$ be its sheaf of rings, and let $\mathsf{Mod}(E)$ be the category of
$\sA$-modules. An object $X$ of $\mathsf{Mod}(E)$ is said to be
\emph{$\sA$-flat} if, for every exact sequence
\[\xymatrix{
  0 \ar[r]
    & Y' \ar[r]
    & Y \ar[r]
    & Y'' \ar[r]
    & 0
}\]
the sequence
$0 \to Y'\otimes_\sA X \to Y\otimes_\sA X \to Y''\otimes_\sA X\to 0$
is exact. There are enough $\sA$-flat objects: the direct sums of objects
$\sA_U$, where $U$ is an object of $E$ and $\sA_U$ is the sheaf $\sA$
“extended by zero outside $U$.” Thus Theorem \ref{VIII:4-2-3} applies,
and one can define the functor
\begin{align*}
  \eD^-(E,\sA) \times \eD^-(E,\sA) &\to \eD^-(E,\sA) \\
  (X,Y) &\mapsto X\lotimes Y
\end{align*}
the total tensor product, which is the left derived functor of the tensor
product. Passing to cohomology defines the $\Tor_i^\sA(X,Y)$, the local
hyper-Tor groups.

Since the category $\mathsf{Mod}(E)$ has enough injectives, one can derive
the functor $\hom^\bullet(X,Y)$, the complex of local homomorphisms, and
thus obtain $\mathsf{Rhom}(X,Y)$ ($X\in \eD$, $Y\in \eD^+$).

Let $X$ and $Y$ be two objects of $\eD^-(E,\sA)$ and let $Z$ be an object
of $\eD^+(E,\sA)$. There is a functorial isomorphism
\[\xymatrix{
  \hom_{\eD(E,\sA)}(X\lotimes Y,Z) \ar[r]^-\sim
    &\hom_{\eD(E,\sA)}(X,\mathsf{Rhom}(Y,Z)) \text{.}
}\]

\subsection{Left derived functor of the inverse image}\label{VIII:6-3}

Let $f:E\to E'$ be a map of ringed topological spaces. The inverse-image
functor
\[
  f^\ast:\mathsf{Mod}(E') \to \mathsf{Mod}(E)
\]
is not necessarily exact. But one knows how to define $f^\ast$-flat
objects, and there are enough such objects. Thus one can define the
functor $\eL^- f^\ast$. There is likewise a duality formula:

Let $X$ be an object of $\eD^-(E',\sA')$ and let $Y$ be an object of
$\eD^+(E,\sA)$. There is a functorial isomorphism
\[\xymatrix{
  \hom_{\eD(E,\sA)}(\eL^- f^\ast(X),Y) \ar[r]^-\sim
    & \hom_{\eD(E',\sA')}(X,\eR^+ f_\ast(Y)) \text{.}
}\]

% END INCLUDED components/sga4half_8_en.tex


\bibliographystyle{amsplain}
\bibliography{sga-sources}

\end{document}

% SGA-CUMULATIVE-END 05_SGA4HALF_EN_COMPLETE_READER.tex

% SGA-CUMULATIVE-BEGIN 06_SGA5_EN_COMPLETE_PUBLISHED_CONTENT_READER.tex
\documentclass[11pt]{article}
\usepackage[T1]{fontenc}
\usepackage[utf8]{inputenc}
\usepackage[english]{babel}
\usepackage{lmodern}
\usepackage{geometry}
\geometry{margin=1in}
\usepackage{microtype}
\usepackage[normalem]{ulem}
\usepackage{amsmath,amssymb,mathtools,mathrsfs}
\usepackage{enumitem}
\usepackage{amsthm}
\usepackage{tikz-cd}
\DeclareMathOperator{\Reg}{Reg}
\DeclareMathOperator{\Supp}{Supp}
\DeclareMathOperator{\Div}{Div}
\DeclareMathOperator{\codim}{codim}
\providecommand{\cO}{\mathcal{O}}
\providecommand{\fm}{\mathfrak{m}}
\providecommand{\wt}[1]{\widetilde{#1}}
\providecommand{\Cons}{\operatorname{Cons}}
\providecommand{\Hh}{\mathrm{H}}
\providecommand{\uH}{\underline{\mathrm H}}
\providecommand{\uExt}{\underline{\Ext}}
\providecommand{\uHom}{\underline{\Hom}}
\providecommand{\uRHom}{\underline{\RHom}}
\providecommand{\qedtext}{\hfill$\square$}
\providecommand{\ZZ}{\mathbb{Z}}
\providecommand{\QQ}{\mathbb{Q}}
\providecommand{\RR}{\mathbb{R}}
\providecommand{\FF}{\mathbb{F}}
\providecommand{\NN}{\mathbb{N}}

\newenvironment{statement}[1]{\par\medskip\noindent\textbf{#1.}\ }{\par\medskip}
\newenvironment{prooftext}[1][Proof]{\par\noindent\emph{#1.}\ }{\par\medskip}

\usepackage{hyperref}
\hypersetup{
  colorlinks=true,
  linkcolor=blue,
  citecolor=blue,
  urlcolor=blue,
  hypertexnames=false,
  pdftitle={SGA 5: Cohomologie l-adique et fonctions L - English Translation},
  pdfauthor={SGA 5 authors and redactors},
  pdfsubject={English translation of SGA 5},
  pdfkeywords={SGA 5, l-adic cohomology, L-functions, algebraic geometry, English translation}
}
\setlength{\parindent}{0pt}
\setlength{\parskip}{0.55em}
\let\accentH\H
\renewcommand{\H}{\mathrm{H}}
\DeclareMathOperator{\Hom}{Hom}
\DeclareMathOperator{\Ext}{Ext}
\DeclareMathOperator{\Tor}{Tor}
\DeclareMathOperator{\Spec}{Spec}
\DeclareMathOperator{\cd}{cd}
\newcommand{\R}{\mathrm{R}}
\newcommand{\D}{\mathrm{D}}
\newcommand{\Z}{\mathbb{Z}}
\newcommand{\F}{\mathbb{F}}
\newcommand{\Q}{\mathbb{Q}}
\newcommand{\E}{\mathrm{E}}
\newcommand{\ob}{\operatorname{ob}}
\newcommand{\charac}{\operatorname{char}}
\newcommand{\ol}[1]{\overline{#1}}
\newcommand{\SGAA}{SGA~A}

\providecommand{\Ker}{\operatorname{Ker}}
\providecommand{\Coker}{\operatorname{Coker}}
\providecommand{\Gal}{\operatorname{Gal}}
\providecommand{\car}{\operatorname{car}}


\providecommand{\RHom}{\mathrm{RHom}}
\providecommand{\Dctf}{D_{\mathrm{ctf}}}
\providecommand{\Id}{\mathrm{Id}}
\providecommand{\pr}{\operatorname{pr}}
\providecommand{\Gr}{\Gamma}
\providecommand{\Fix}{\operatorname{Fix}}
\providecommand{\Loc}{\operatorname{Loc}}
\providecommand{\lto}{\longrightarrow}
\providecommand{\xto}[1]{\xrightarrow{#1}}
\providecommand{\lisom}{\xrightarrow{\sim}}
\providecommand{\ctf}{\mathrm{ctf}}
\providecommand{\et}{\mathrm{\acute{e}t}}
\providecommand{\Sfrac}{SGA~$4^{1/2}$}


\providecommand{\cl}{\operatorname{cl}}
\providecommand{\Tr}{\operatorname{Tr}}
\providecommand{\EP}{\operatorname{EP}}
\providecommand{\gr}{\operatorname{gr}}
\providecommand{\DF}{\operatorname{DF}}

\providecommand{\OO}{\mathcal{O}}
\providecommand{\LL}{\mathbf L}
\providecommand{\parf}{\mathrm{parf}}
\providecommand{\coh}{\mathrm{coh}}
\providecommand{\frob}{\operatorname{Frob}}

\providecommand{\RGamma}{\mathrm R\Gamma}
\providecommand{\Rj}{\mathrm Rj}
\providecommand{\Tr}{\operatorname{Tr}}
\providecommand{\id}{\operatorname{id}}
\providecommand{\cone}{\operatorname{cone}}
\providecommand{\lten}{\overset{\mathrm L}{\otimes}}
\providecommand{\Dctf}{D_{\mathrm{ctf}}}
\providecommand{\Ob}{\operatorname{ob}}

\providecommand{\Tr}{\operatorname{Tr}}
\providecommand{\RHom}{\mathrm{RHom}}
\providecommand{\Hom}{\operatorname{Hom}}
\providecommand{\ob}{\operatorname{ob}}
\providecommand{\lten}{\overset{\mathrm L}{\otimes}}
\providecommand{\nat}{\natural}
\providecommand{\dfn}{\mathrm{dfn}}

\providecommand{\Fl}{\operatorname{Fl}}
\providecommand{\MLAR}{\mathrm{MLAR}}
\providecommand{\AR}{\mathrm{AR}}
\providecommand{\NN}{\mathbb N}
\providecommand{\Ens}{\mathrm{Ens}}
\providecommand{\cC}{\mathcal C}
\providecommand{\cE}{\mathcal E}
\providecommand{\cF}{\mathcal F}
\providecommand{\cG}{\mathfrak G}
\providecommand{\source}{\operatorname{source}}
\providecommand{\op}{\mathrm{op}}


\providecommand{\End}{\operatorname{End}}
\providecommand{\id}{\operatorname{id}}
\providecommand{\Im}{\operatorname{Im}}
\providecommand{\Ker}{\operatorname{Ker}}
\providecommand{\Coker}{\operatorname{Coker}}
\providecommand{\gr}{\operatorname{gr}}

\renewcommand{\Im}{\operatorname{Im}}

\providecommand{\Rlim}{\varprojlim}
\providecommand{\adn}{\operatorname{adn}}
\providecommand{\adf}{\operatorname{adf}}
\providecommand{\AR}{\mathrm{AR}}
\providecommand{\MLAR}{\mathrm{MLAR}}
\providecommand{\Ecat}{\mathcal E}
\providecommand{\Dcat}{\mathcal D}
\providecommand{\Ccat}{\mathcal C}
\providecommand{\Bcat}{\mathcal B}
\providecommand{\Pcat}{\mathcal P}
\providecommand{\J}{J}
\providecommand{\gr}{\operatorname{gr}}
\providecommand{\grs}{\operatorname{grs}}
\providecommand{\Homcat}{\operatorname{Hom}}
\providecommand{\Ker}{\operatorname{Ker}}
\providecommand{\Coker}{\operatorname{Coker}}
\providecommand{\Im}{\operatorname{Im}}
\providecommand{\can}{\operatorname{can}}
\providecommand{\epi}{\operatorname{epi}}
\providecommand{\id}{\operatorname{id}}


\providecommand{\RGamma}{\mathrm R\Gamma}
\providecommand{\RHom}{\mathrm R\mathcal{H}om}
\providecommand{\Zell}{\mathbb Z_\ell}
\providecommand{\Qell}{\mathbb Q_\ell}
\providecommand{\Ql}{\mathbb Q_\ell}
\providecommand{\lten}{\overset{\mathrm L}{\otimes}}
\providecommand{\Ob}{\operatorname{ob}}
\providecommand{\Ens}{\operatorname{Ens}}
\providecommand{\Ensf}{\operatorname{Ensf}}
\providecommand{\Aut}{\operatorname{Aut}}
\providecommand{\End}{\operatorname{End}}
\providecommand{\Sp}{\operatorname{Sp}}
\providecommand{\Ab}{\operatorname{Ab}}
\providecommand{\Abc}{\operatorname{Abc}}
\providecommand{\Abf}{\operatorname{Abf}}
\providecommand{\fct}{\operatorname{fct}}
\providecommand{\Lc}{\operatorname{Lc}}
\providecommand{\GL}{\operatorname{GL}}
\providecommand{\car}{\operatorname{car}}
\providecommand{\cHom}{\mathcal{H}om}
\providecommand{\fc}{\operatorname{fc}}
\providecommand{\adc}{\operatorname{adc}}
\providecommand{\ad}{\operatorname{ad}}
\providecommand{\adn}{\operatorname{adn}}
\providecommand{\modn}{\operatorname{mod}}
\providecommand{\AR}{\operatorname{AR}}


\providecommand{\R}{\mathrm R}
\providecommand{\V}{\mathrm V}
\providecommand{\A}{\mathbb A}
\providecommand{\P}{\mathbb P}
\providecommand{\Z}{\mathbb Z}
\providecommand{\muS}{\mu_S}
\providecommand{\cupdot}{\smile}
\providecommand{\id}{\operatorname{id}}
\providecommand{\Tr}{\operatorname{Tr}}
\providecommand{\Spec}{\operatorname{Spec}}
\providecommand{\Supp}{\operatorname{Supp}}
\providecommand{\coker}{\operatorname{coker}}
\providecommand{\lten}{\overset{\mathrm L}{\otimes}}
\newenvironment{proposition}[1][]{\par\medskip\noindent\textbf{Proposition #1.}\ }{\par\medskip}
\newenvironment{corollaire}[1][]{\par\medskip\noindent\textbf{Corollaire #1.}\ }{\par\medskip}
\newenvironment{theoreme}[1][]{\par\medskip\noindent\textbf{Théorème #1.}\ }{\par\medskip}
\newenvironment{lemme}[1][]{\par\medskip\noindent\textbf{Lemme #1.}\ }{\par\medskip}


\providecommand{\Ahat}{\widehat A}
\providecommand{\Gm}{\mathbb G_m}
\providecommand{\Et}{\mathrm{et}}
\providecommand{\Pic}{\operatorname{Pic}}
\providecommand{\cl}{\operatorname{cl}}
\providecommand{\EP}{\operatorname{EP}}
\providecommand{\Ch}{\operatorname{Ch}}
\providecommand{\lten}{\overset{\mathrm L}{\otimes}}
\providecommand{\rank}{\operatorname{rank}}
\providecommand{\Sym}{\operatorname{Sym}}
\providecommand{\Spec}{\operatorname{Spec}}
\providecommand{\pr}{\operatorname{pr}}
\providecommand{\id}{\operatorname{id}}
\providecommand{\im}{\operatorname{Im}}
\providecommand{\coker}{\operatorname{Coker}}


\providecommand{\cC}{\mathcal C}
\providecommand{\cD}{\mathcal D}
\providecommand{\cA}{\mathcal A}
\providecommand{\cB}{\mathcal B}
\providecommand{\Ab}{\operatorname{Ab}}
\providecommand{\Ob}{\operatorname{Ob}}
\providecommand{\parf}{\mathrm{parf}}
\providecommand{\coh}{\mathrm{coh}}
\providecommand{\lten}{\overset{\mathrm L}{\otimes}}


\providecommand{\Mod}{\operatorname{Mod}}
\providecommand{\tr}{\operatorname{tr}}
\providecommand{\Sw}{\operatorname{Sw}}
\providecommand{\Art}{\operatorname{Art}}
\providecommand{\etG}{\mathrm{\acute et},G}
\providecommand{\cT}{\mathcal T}


\providecommand{\clK}[1]{\operatorname{cl}_{K(#1)}}
\providecommand{\clKstar}[1]{\operatorname{cl}_{K^{\bullet}(#1)}}
\providecommand{\epsD}{\varepsilon^\Delta}
\providecommand{\alphaD}{\alpha^\Delta}
\providecommand{\bbone}{\mathbf 1}
\providecommand{\sw}{\operatorname{sw}}
\providecommand{\art}{\operatorname{art}}
\providecommand{\RG}{\mathrm R\Gamma}
\providecommand{\RGc}{\mathrm R\Gamma_c}
\providecommand{\RHom}{\mathrm{RHom}}
\providecommand{\Rpi}{\mathrm R\pi}
\providecommand{\rG}{\mathrm R\Gamma^G}
\providecommand{\wh}[1]{\widehat{#1}}
\providecommand{\checkSw}{\check{\operatorname{Sw}}}


\providecommand{\La}{\Lambda}
\providecommand{\bZ}{\mathbb Z}
\providecommand{\bQ}{\mathbb Q}
\providecommand{\bZell}{\mathbb Z_\ell}
\providecommand{\fm}{\mathfrak m}
\providecommand{\card}{\operatorname{card}}
\providecommand{\im}{\operatorname{im}}
\providecommand{\Aut}{\operatorname{Aut}}
\providecommand{\Hom}{\operatorname{Hom}}
\providecommand{\Spec}{\operatorname{Spec}}
\providecommand{\cl}{\operatorname{cl}}
\providecommand{\NW}{\operatorname{NW}}
\providecommand{\Tr}{\operatorname{Tr}}
\providecommand{\Ob}{\operatorname{ob}}
\providecommand{\car}{\operatorname{car}}
\providecommand{\qedtext}{\hfill$\square$}


\providecommand{\fr}{\operatorname{fr}}
\providecommand{\Fr}{\operatorname{Fr}}
\providecommand{\Sch}{\operatorname{Sch}}
\providecommand{\Proj}{\operatorname{Proj}}
\providecommand{\Sym}{\operatorname{Sym}}
\providecommand{\PP}{\mathbb P}
\providecommand{\JJ}{\mathcal J}
\providecommand{\XX}{\mathcal X}
\providecommand{\cL}{\mathcal L}
\providecommand{\cM}{\mathcal M}
\providecommand{\cN}{\mathcal N}
\providecommand{\cS}{\mathcal S}
\providecommand{\cX}{\mathcal X}
\providecommand{\cY}{\mathcal Y}
\providecommand{\cZ}{\mathcal Z}
\providecommand{\bF}{\mathbb F}
\providecommand{\bbZ}{\mathbb Z}
\providecommand{\simeqto}{\xrightarrow{\sim}}


\providecommand{\Fr}{\operatorname{Fr}}
\providecommand{\fr}{\operatorname{fr}}
\providecommand{\Sch}{\operatorname{Sch}}
\providecommand{\Spec}{\operatorname{Spec}}
\providecommand{\id}{\operatorname{id}}
\providecommand{\RGamma}{\mathrm R\Gamma}
\providecommand{\Ob}{\operatorname{Ob}}
\providecommand{\Gal}{\operatorname{Gal}}
\providecommand{\et}{\acute{e}t}
\providecommand{\lten}{\overset{\mathrm L}{\otimes}}
\providecommand{\boxt}{\boxtimes}
\begin{document}

\begin{titlepage}
\centering
\vspace*{2.1in}
{\LARGE\textbf{SGA 5}}\\[0.6em]
{\Large $\ell$-adic Cohomology and $L$-Functions}\\[1.5em]
{\large Complete English reader}\\[1.2em]
{\normalsize All ten published exposés}
\vfill
\end{titlepage}
\tableofcontents
\clearpage

\phantomsection\addcontentsline{toc}{section}{Exposé I --- Dualizing Complexes}
\begin{center}
{\Large \textbf{SGA 5 --- Exposé I}}\\[0.4em]
{\large \textbf{Dualizing Complexes}}\\[0.3em]
by A.~Grothendieck\\
(written up by L.~Illusie)
\end{center}

\section*{Introduction}

This exposé is devoted to the biduality theorem in étale cohomology. The theory developed here and in SGA~A~XVIII (SGA~A = SGA~4), for the étale topology and constructible torsion sheaves on preschemes, is formally very analogous to the one developed in HARTSHORNE~[H] for the Zariski topology and coherent sheaves: the latter served, to a very large extent, as its model. But, whereas in the case of ``continuous'' coefficients the existence of dualizing complexes presents no difficulties, in the case of ``discrete'' coefficients, on the other hand, it is much more delicate to prove, and it seems to require resolution of singularities in an essential way. See however Deligne's biduality theorem (SGA~$4^{1/2}$, Finiteness Theorem~4.3.).

Paragraph~1 contains the definition of dualizing complexes and the formal properties that follow from it, notably the interpretation of dualizing functors as ``exchangers of functors.''

In paragraph~2, one shows that, on a connected locally noetherian prescheme, a dualizing complex is determined uniquely, up to a translation of degrees and tensoring by an invertible sheaf.

In paragraph~3, one comes to the central theorem of the exposé (3.4.1.), which asserts that, on a ``good'' regular prescheme $X$, for example on an excellent noetherian regular prescheme of characteristic zero\footnote{This last restriction will become unnecessary once one has resolution of singularities à la Hironaka for these preschemes, and the ``purity theorem.''}, the complex $(\Z/n\Z)_X$ is dualizing.

Paragraph~4 gives the applications of the theory to the local duality theorem. The latter expresses that, on a ``good'' prescheme $X$ endowed with a closed point $x$ whose residue field is separably closed, and for a constructible sheaf $F$, one has a perfect pairing between $\underline{\H}^i(F')_x$ and $\underline{\H}^{-i}_x(F)$, where $F'$ is the dual of $F$, i.e.\ $F' = \R\,\underline{\Hom}(F,K_X)$, where $K_X$ is a dualizing complex.

Finally, in paragraph~5, one proves directly that on a regular prescheme of dimension~1 the complex $(\Z/n\Z)_X$ is dualizing ($n$ prime to the residual characteristics).

\section{Definition and formal properties of dualizing complexes}

We fix a coefficient ring $A = \Z/n\Z$. If $X$ is a prescheme, we denote by $A_X$ the constant sheaf on $X_{\mathrm{\acute{e}t}}$ with value $A$, and by $\D(X)$ the derived category of the category of $A_X$-modules. We denote, on the other hand, by $\D_c(X)$ the full triangulated subcategory of $\D(X)$ formed by the complexes $F$ such that $\H^i(F)$ is a constructible $A_X$-module for every $i$. We put $\D_c^{*}(X) = \D^{*}(X) \cap \D_c(X)$, where $* = \varnothing$, $+$, $-$, or $b$.

\subsection*{Definition 1.1}
An $A_X$-module $I$ is said to be \emph{quasi-injective} if one has
\[
\underline{\Ext}^i(F,I) = 0
\]
for every \emph{constructible} $A_X$-module $F$ and every $i > 0$.

\subsection*{Remarks 1.2}
\begin{enumerate}[label=\alph*)]
\item Contrary to what happens in the case of ``continuous'' coefficients (cf.~[H]~II.7.17), a quasi-injective sheaf is not injective in general. For example, let $X = \Spec(\R)$ and $A = \Z/2\Z$: the sheaf $A_X$ is quasi-injective, but is not injective, nor even of finite injective dimension, since $\H^{*}(X; A_X) = \H^{*}(\Z/2\Z; \Z/2\Z)$ is a polynomial algebra on one generator of dimension~1.

\item Let $I$ be a quasi-injective $A_X$-module. It is false in general that the relation $\underline{\Ext}^i(F,I) = 0$ for $i > 0$, valid for constructible $F$, is valid for every $F$. Indeed, take $A = \F_{\ell}$. Let $k$ be a field, $\ol{k}$ a separable closure of $k$, and $f : \ol{x} = \Spec(\ol{k}) \to X = \Spec(k)$ the canonical morphism, and suppose that, for every connected étale covering $X' \to X$, there exists a connected étale covering $X'' \to X'$ such that $[X'':X']$ is divisible by $\ell$ (take, for example, $k = \Q$). Consider the exact sequence
\[
(*) \qquad 0 \longrightarrow A_X \xrightarrow{\;\varepsilon\;} f_{*}f^{*}A_X \longrightarrow F \longrightarrow 0
\]
defined by the adjunction arrow $\varepsilon$. In view of the hypothesis, the section of $\underline{\Ext}^1(F,A_X)$ defined by $(*)$ does not vanish above any $U$ étale over $X$; hence $\underline{\Ext}^1(F,A_X)_{\ol{x}} \neq 0$. But one checks trivially that $A_X$ is quasi-injective.
\end{enumerate}

\subsection*{Proposition 1.3 {\normalfont (cf.~[H]~1.7.6)}}
Let $X$ be a prescheme and $N \in \Z$. The following conditions on $K \in \ob \D^{+}(X)$ are equivalent:
\begin{enumerate}[label=(\roman*)]
\item $K$ is isomorphic, in $\D(X)$, to a bounded complex $I$ such that $I^i$ is quasi-injective for every $i$ and zero for $i > N$;
\item for every $F \in \ob \D_c^{b}(X)$ such that $\H^i(F) = 0$ for $i < 0$, one has $\underline{\Ext}^i(F,K) = 0$ for $i > N$;
\item for every constructible $A_X$-module $F$, one has $\underline{\Ext}^i(F,K) = 0$ for $i > N$.
\end{enumerate}

\emph{Proof.} $(i) \Rightarrow (ii)$. Let $F \in \ob \D_c^{b}(X)$ be such that $\H^i(F) = 0$ for $i < 0$, and let $I \in \ob \D^{+}(X)$ be bounded and such that $I^i$ is quasi-injective for every $i$ and zero for $i > N$; we show that $\underline{\Ext}^i(F,I) = 0$ for $i > N$. Arguing by induction on the number of indices $i$ such that $I^i \neq 0$, one is reduced, by dévissage, to the case where $N = 0$ and $I$ is concentrated in degree $0$. One then has a biregular spectral sequence
\[
\E_2^{pq} = \underline{\Ext}^p(\H^{-q}(F),I) \Rightarrow \underline{\Ext}^{*}(F,I).
\]
Since $I$ is quasi-injective, $\E_2^{pq} = 0$ for $p > 0$ and every $q$, and hence one has $\E_2^{0i} = \underline{\Hom}(\H^{-i}(F),I) \xrightarrow{\sim} \underline{\Ext}^i(F,I)$, which gives the conclusion.

$(ii) \Rightarrow (iii)$ is trivial.

$(iii) \Rightarrow (i)$ We may suppose that $K$ is bounded below and has injective components. The hypothesis, applied to $F = A_X$, implies that $\H^i(K) = 0$ for $i > N$. We may therefore replace $K$ by an isomorphic complex $K'$, bounded, and such that $K'^i$ is injective for $i < N$ and zero for $i > N$. It is then observed that $K'^N$ is quasi-injective, which completes the proof.

\subsection*{Remark 1.3.1}
The writer does not know whether (1.3) is still true when, in condition (ii), the hypothesis ``$F \in \ob \D_c^{b}(X)$'' is replaced by ``$F \in \ob \D_c^{+}(X)$''.

\subsection*{Definition 1.4}
The \emph{quasi-injective dimension} of $K$, denoted $\dim.\,q.\,inj(K)$, is the lower bound (in $\overline{\R}$) of the integers $N$ satisfying the equivalent conditions of (1.3).

If $K$ is of finite quasi-injective dimension, the functor $\R\,\underline{\Hom}(\quad,K) : \D(X) \to \D(X)$ induces a functor $\D_c^{b}(X) \to \D^{b}(X)$.

We shall see below that if $X$ is smooth of dimension $d$ over a field $k$ and $n$ is prime to the characteristic of $k$, then $\dim.\,q.\,inj(A_X) = 2d$.

The example of (1.2~a)) shows that a complex of finite quasi-injective dimension is not necessarily of finite injective dimension. Nevertheless one has:

\subsection*{Proposition 1.5}
Let $X$ be a locally noetherian prescheme and $K \in \ob \D^{+}(X)$. Suppose that there exists an integer $N$ such that $\cd_{\mathbf{n}}(X) \leqslant N$, where $\mathbf{n}$ is the set of prime divisors of $n$ (notation of SGAA~X~1). Then, for $N' \in \Z$, one has
\[
\dim.\,q.\,inj(K) \leqslant N' \quad \Longrightarrow \quad \dim.\,inj(K) \leqslant N + N'.
\]

We shall use, in the proof, the following

\subsection*{Lemma 1.5.1}
Let $C$ be an abelian category satisfying (AB5) and possessing a family of generators $(U_i)_{i \in I}$. Let $K \in \ob \D^{+}(C)$ and $N \in \Z$. The following conditions are equivalent:
\begin{enumerate}[label=(\roman*)]
\item $\dim.\,inj(K) \leqslant N$;
\item for every $i \in I$ and every quotient $F$ of $U_i$, one has
\[
\Ext^n(F,K) = 0 \qquad \text{for } n > N.
\]
\end{enumerate}

\emph{Proof.} It suffices to prove $(ii) \Rightarrow (i)$. By a well-known argument (cf.\ for example [H]~1.7.6), one is reduced to showing that if $K' \in \ob(C)$ is such that $\Ext^1(F,K') = 0$ every time $F$ is a quotient of a $U_i$, then $K'$ is injective. In other words, one has to prove that, if every morphism from a subobject $V$ of a $U_i$ into $K'$ extends to $U_i$, then $K'$ is injective. But this follows immediately from the proof of Lemma~1, p.~136 of (Tohoku).

\emph{Proof of (1.5).} Suppose $\dim.\,q.\,inj(K) \leqslant N'$ and prove $\dim.\,inj(K) \leqslant N + N'$. Since the category of $A_X$-modules admits as generators the $A_{U,X}$ where $U$ is étale of finite type over $X$, it suffices, by (1.5.1), to prove that $\Ext^i(X; F,K) = 0$ for $i > N + N'$ when $F$ is a quotient of such an $A_{U,X}$, hence constructible (SGAA~IX~2.9). Consider the spectral sequence
\[
\E_2^{pq} = \H^p(X; \underline{\Ext}^q(F,K)) \Rightarrow \Ext^{*}(X; F,K).
\]
The hypothesis $\cd_{\mathbf{n}}(X) \leqslant N$ (resp. $\dim.\,q.\,inj(K) \leqslant N'$) implies $\E_2^{pq} = 0$ for $p > N$ (resp. $q > N'$). Hence $\E_2^{pq} = 0$ for $p+q > N+N'$, whence the conclusion.

\subsection*{Remark 1.5.2}
The hypothesis of (1.5) is verified (with $N = 2d$) if $X$ is a prescheme of finite type and dimension $d$ over a separably closed field $k$, with $n$ prime to $\charac(k)$ (SGAA~X~4).

\subsection*{Proposition 1.6}
Let $f : X \to Y$ be a separated quasi-finite morphism of noetherian preschemes, and let $K \in \ob \D^{+}(Y)$. Then, for $N \in \Z$, one has
\[
\dim.\,q.\,inj(K) \leqslant N \quad \Longrightarrow \quad \dim.\,q.\,inj(\R^{!}f(K)) \leqslant N.
\]

\emph{Proof.} According to the ``Main theorem'' (EGA~IV~8.12.6), there exists a factorization of $f$ as $f = uf'$, where $f'$ is an open immersion and $u$ is a finite morphism. We have $\R^{!}f = \R^{!}f' \, \R^{!}u$, which reduces us to examining separately the case of an open immersion and that of a finite morphism. If $f$ is an open immersion, the assertion is trivial (since every constructible sheaf on $X$ is then induced by a constructible sheaf on $Y$). Suppose therefore that $f$ is a finite morphism, and let $F$ be a constructible $A_X$-module. We then have the duality isomorphisms (SGAA~XVIII~3.1.9.6)
\[
f_{*}\underline{\Ext}^i(F, \R^{!}f(K)) \xrightarrow{\sim} \underline{\Ext}^i(f_{*}F, K).
\]
Since $f_{*}(F)$ is constructible (SGAA~IX~2.14), it follows that, if $\dim.\,q.\,inj(K) \leqslant N$, then $f_{*}\underline{\Ext}^i(F, \R^{!}f(K)) = 0$ for $i > N$, hence $\underline{\Ext}^i(F, \R^{!}f(K)) = 0$ for $i > N$ (SGAA~VIII~5.5), as required.

Let now $X$ be a prescheme, and let $K \in \ob \D^{+}(X)$. Denote by $\underline{D}_K$ the functor $\R\,\underline{\Hom}(\quad, K)$. We shall define, for $F \in \ob \D(X)$, a functorial homomorphism $F \to \underline{D}_K \underline{D}_K F$ (cf.~[H]~V.1.2). The construction which follows would be valid more generally on a ringed topos. We may first suppose that $K$ is formed of injectives, which allows us to ``remove the $\R$''. The identity morphism $\underline{\Hom}(F,K) \to \underline{\Hom}(F,K)$ defines, by the formula dear to Cartan, a homomorphism $F \otimes \underline{\Hom}(F,K) \to K$, whence, by a second application of the said formula, a homomorphism $F \to \underline{\Hom}(\underline{\Hom}(F,K), K)$, which is the announced homomorphism, and which we baptize the canonical homomorphism. We say that the pair $(F,K)$ is \emph{bidualizing}, or \emph{satisfies biduality}, or again that $F$ is \emph{reflexive} relative to $K$, if the canonical homomorphism $F \to \underline{D}_K \underline{D}_K F$ is an isomorphism.

\textbf{\uline{From now on, unless expressly mentioned otherwise, all preschemes considered will be assumed locally noetherian.}}

\subsection*{Definition 1.7}
Let $X$ be a prescheme. A complex $K \in \ob \D^{+}(X)$ is said to be \emph{dualizing} if the following conditions are satisfied:
\begin{enumerate}[label=(\roman*)]
\item $K$ is of finite quasi-injective dimension;
\item for every $F \in \ob \D_c^{-}(X)$, one has $\underline{D}_K(F) \in \ob \D_c^{+}(X)$;
\item every $F \in \ob \D_c^{b}(X)$ is reflexive relative to $K$.
\end{enumerate}

\subsection*{Remarks 1.8}
\begin{enumerate}[label=\alph*)]
\item By a standard dévissage (truncating $F$ and using the fact that $\D_c(X)$ is triangulated), one sees that condition (ii) is equivalent to

(ii~bis): for every constructible $A_X$-module $F$ and every $i \in \Z$, $\underline{\Ext}^i(F,K)$ is constructible.

\item Condition (ii~bis) implies $K \in \ob \D_c^{+}(X)$. Contrary to the example of coherent sheaves (EGA~$0_{\mathrm{III}}$~12.3.3), the converse is perhaps not automatically verified. It is however verified (see no.~3) in the ``good'' cases, for example when one has on $X$ resolution of singularities and purity (in particular if $X$ is an excellent prescheme of characteristic zero).

\item By dévissage on $F$, one sees that condition (iii) is equivalent to (iii~bis): every constructible $A_X$-module $F$ is reflexive relative to $K$.

\item If $K \in \ob \D^{+}(X)$ satisfies (i) and (ii) (resp. is dualizing), the functor $\underline{D}_K$ induces a functor (resp. an equivalence of categories) $(\D_c^{b}(X))^{\circ} \to \D_c^{b}(X)$.

\item If $K \in \ob \D^{+}(X)$ is of finite \emph{injective} dimension (cf.~1.5.2) and satisfies (ii), then, for every $F \in \ob \D_c(X)$, one has $\underline{D}_K(F) \in \D_c(X)$, and $\underline{D}_K$ induces a functor $(\D_c^{+}(X))^{\circ} \to \D_c^{-}(X)$.

If moreover $K$ is dualizing, every $F \in \ob \D_c(X)$ is reflexive relative to $K$, and $\underline{D}_K$ induces equivalences of categories
\[
(\D_c(X))^{\circ} \xrightarrow{\approx} \D_c(X), \quad (\D_c^{-}(X))^{\circ} \xrightarrow{\approx} \D_c^{+}(X), \quad (\D_c^{+}(X))^{\circ} \xrightarrow{\approx} \D_c^{-}(X).
\]
It is not known whether this conclusion remains valid without the restriction that $K$ be of finite injective dimension.
\end{enumerate}

We now turn to the ``exchange formulas.'' We shall not return to the notion of the tor-dimension of a complex, see SGA A, Exposé XVII, 4.1.9.\footnote{See also SGA 6, I, 5.} Retain only the following point: if $X$ is a prescheme and $G\in\ob\D^-(X)$, the condition $\underline{\Tor}_i(F,G)=0$ for $i\geqslant N$ and all $A_X$-modules $F$ is equivalent to the same condition with $F$ constructible, by SGA A, Exposé IX, 2.9 and the commutation of $\otimes^{\mathrm L}$ with filtered inductive limits. Thus there is no need to introduce a notion of ``quasi-tor-dimension.''

\begin{statement}{Lemma 1.9, cf. [H] II.4.3}
\begin{enumerate}[label=\alph*)]
\item Let $X$ be a prescheme, let $F\in\ob\D_c(X)$, and let $G\in\ob\D_c^-(X)$. One has $F\otimes^{\mathrm L}G\in\ob\D_c(X)$ if $F\in\ob\D_c^-(X)$, or if $G$ has finite tor-dimension.
\item Let $f:X\to Y$ be a morphism of preschemes and let $G\in\ob\D_c(Y)$. Then $f^*(G)\in\ob\D_c(X)$.
\end{enumerate}
\end{statement}

\begin{prooftext}
Part (b) follows trivially from SGA A, Exposé IX, 2.4. We prove (a). Either hypothesis on $(F,G)$ implies that the spectral sequence
\[
   E_2^{pq}=\sum_{q_1+q_2=q}\underline{\Tor}_{-p}\bigl(\Hh^{q_1}(F),\Hh^{q_2}(G)\bigr)
       \Rightarrow \Hh^*(F\otimes^{\mathrm L}G)
\]
is biregular. Hence one is reduced to the case where $F$ and $G$ are concentrated in degree $0$, that is, are constructible $A_X$-modules. The question is local, so one may suppose $X$ noetherian. Then $F$ admits, by SGA A, Exposé IX, 2.7, a left resolution $F'$ in which the $F'^i$ are of the form $A_{U,X}$ with $U$ étale and of finite type over $X$. We have
\[
   F\otimes^{\mathrm L}G\xleftarrow{\sim}F'\otimes G,
\]
and the result follows since the $A_{U,X}\otimes G$ are constructible. \qedtext
\end{prooftext}

\begin{statement}{Notation 1.10}
If $X$ is a prescheme, we denote by $\D^b(X)_{\mathrm{torf}}$, respectively $\D^b(X)_{\mathrm{qinjf}}$, the full subcategory of $\D(X)$ consisting of objects of finite tor-dimension, respectively of finite quasi-injective dimension.
\end{statement}

\begin{statement}{Proposition 1.11, cf. [H] V.2.6}
Let $X$ be a prescheme and $K\in\ob\D^+(X)$.
\begin{enumerate}[label=\alph*)]
\item For $F,G\in\ob\D^-(X)$, respectively for $F\in\ob\D(X)$ and $G\in\ob\D^b(X)_{\mathrm{torf}}$, there is a functorial isomorphism
\[
   \underline D_K(F\otimes^{\mathrm L}G)\xrightarrow{\sim}
   \R\,\underline{\Hom}(F,\underline D_K(G)).
\]

\item If $K$ satisfies conditions (i) and (ii) of Definition 1.7, then
\[
   G\in\ob\D_c^b(X)_{\mathrm{torf}}
   \Longrightarrow
   \underline D_K(G)\in\ob\D_c^b(X)_{\mathrm{qinjf}}
\]
and $\underline D_K(G)$ satisfies condition 1.7(ii).

\item If $K$ is dualizing, then for $F\in\ob\D_c^-(X)$ and $G\in\ob\D_c^b(X)$, respectively for $F\in\ob\D(X)$ and $G\in\ob\D_c^b(X)$ with $\underline D_K(G)\in\ob\D_c^b(X)_{\mathrm{torf}}$, there is a functorial isomorphism
\[
   \underline D_K\bigl(F\otimes^{\mathrm L}\underline D_K(G)\bigr)
      \xrightarrow{\sim}\R\,\underline{\Hom}(F,G).
\]
Moreover,
\[
   \underline D_K(G)\in\ob\D_c^b(X)_{\mathrm{torf}}
   \Longrightarrow
   G\in\ob\D_c^b(X)_{\mathrm{qinjf}},
\]
and $G$ satisfies condition 1.7(ii).

\item If $K$ is dualizing and has finite injective dimension, then the implications in (b) and (c) are equivalences.
\end{enumerate}
\end{statement}

\begin{prooftext}
Part (a) is precisely Cartan's adjunction formula, SGA 6, I, 7.4.

For (b), let $G\in\ob\D_c^b(X)_{\mathrm{torf}}$. If $F$ is a variable constructible $A_X$-module, then the cohomology objects of $F\otimes^{\mathrm L}G$ are constructible, by Lemma 1.9, and vanish outside two fixed bounds independent of $F$. The same therefore holds for the cohomology objects of $\underline D_K(F\otimes^{\mathrm L}G)$, and (a) gives the assertion.

For (c), the first assertion follows from (a) applied to $\underline D_K(G)$ in place of $G$, together with the reflexivity of $G$. The second follows by a similar argument to (b).

For (d), prove for instance the converse implication of (b). Suppose that $\underline D_K(G)\in\ob\D_c^b(X)_{\mathrm{qinjf}}$ and satisfies condition 1.7(ii). Then, as $F$ varies through constructible modules, the cohomology objects of $\underline D_K(F\otimes^{\mathrm L}G)$ are constructible and vanish outside two fixed bounds. The same is consequently true for the cohomology objects of $\underline D_K\underline D_K(F\otimes^{\mathrm L}G)$; but
\[
   F\otimes^{\mathrm L}G\xrightarrow{\sim}\underline D_K\underline D_K(F\otimes^{\mathrm L}G)
\]
because $K$ is dualizing and of finite injective dimension, cf. 1.8(d). Hence $G$ has finite tor-dimension. The converse implication in (c) is proved analogously. \qedtext
\end{prooftext}

\begin{statement}{Remarks 1.11.1}
\begin{enumerate}[label=\alph*)]
\item The editor does not know whether assertion (d) of Proposition 1.11 is valid under the sole hypothesis that $K$ be dualizing.
\item One expresses the preceding proposition pictorially by saying that the dualizing functor $\underline D_K$ exchanges the functors $\otimes^{\mathrm L}$ and $\R\,\underline{\Hom}$, as well as the notions of finite tor-dimension and finite quasi-injective dimension. The exchange is, however, fully satisfactory only when $K$ has finite injective dimension.
\end{enumerate}
\end{statement}

\begin{statement}{Proposition 1.12}
Let $Y$ be a noetherian prescheme, let $f:X\to Y$ be a morphism of finite type, and let $K\in\ob\D^+(Y)$. Suppose that $n$ is prime to the residual characteristics of $Y$, and put
\[
   K_X=\R^!f(K_Y),\qquad \underline D_X=\underline D_{K_X},\qquad \underline D_Y=\underline D_{K_Y},
\]
in the notation of SGA A, Exposé XVIII, 3.1.
\begin{enumerate}[label=\alph*)]
\item For $F\in\ob\D^-(X)$ there is a functorial isomorphism
\begin{equation}
\R f_*\underline D_X(F)\xrightarrow{\sim}\underline D_Y\R_!f(F).\tag{i}
\end{equation}
If, moreover, $K_X$ and $K_Y$ are dualizing,\footnote{We shall see below that under fairly general conditions $K_Y$ dualizing implies $K_X$ dualizing.} then for $F\in\ob\D_c^b(X)$ there is a functorial isomorphism
\begin{equation}
\R_!f\,\underline D_X(F)\xrightarrow{\sim}\underline D_Y\R f_*(F).\tag{ii}
\end{equation}

\item For $F\in\D_c^-(Y)$ there is a functorial isomorphism\footnote{If the functor $\R^!f$ has finite cohomological dimension, as for example when $Y$ is of finite type over a field (SGA A, Exposé X, 4.3 and Exposé XVIII, 3.1.7), the isomorphism (b)(i) is valid for $F\in\ob\D_c(Y)$.}
\begin{equation}
\R^!f\,\underline D_Y(F)\xrightarrow{\sim}\underline D_X f^*(F).\tag{i}
\end{equation}
If, moreover, $K_X$ and $K_Y$ are dualizing, then for $F\in\D_c^b(Y)$ there is a functorial isomorphism
\begin{equation}
f^*\underline D_YF\xrightarrow{\sim}\underline D_X\R^!f\,F.\tag{ii}
\end{equation}
\end{enumerate}
\end{statement}

\begin{prooftext}
The isomorphism (a)(i) is just the duality isomorphism of SGA A, Exposé XVIII, 3.1.9.6.

The isomorphism (a)(ii) is obtained by applying $\underline D_Y$ to both sides of (a)(i), written with argument $\underline D_XF$, and using the finiteness theorem SGA A, Exposé XVII, 5.3.6, which asserts that $\R_!f(\underline D_X(F))\in\ob\D_c(Y)$.

The isomorphism (b)(i) is the induction formula of SGA A, Exposé XVIII, 3.1.12.2.

The isomorphism (b)(ii) is obtained by applying $\underline D_X$ to both sides of (b)(i), written with argument $\underline D_YF$, and using 1.9(b). \qedtext
\end{prooftext}

\begin{statement}{Scholium 1.12.1}
In figurative language, the dualizing functors $\underline D_X$ and $\underline D_Y$ exchange the usual and unusual direct images, $\R f_*$ and $\R_!f$, and also the usual and unusual inverse images, $f^*$ and $\R^!f$.
\end{statement}

\begin{statement}{Corollary 1.13}
Under the hypotheses of Proposition 1.12, suppose in addition that $f$ is proper, and let $F\in\ob\D^-(X)$. If $(F,K_X)$ is bidualizing, then $(\R f_*F,K_Y)$ is bidualizing. Conversely, if $f$ is finite and $(\R f_*F,K_Y)$ is bidualizing, then $(F,K_X)$ is bidualizing.
\end{statement}

\begin{prooftext}
Suppose $(F,K_X)$ is bidualizing. Applying $\R f_*$ to the canonical isomorphism $F\xrightarrow{\sim}\underline D_X^2F$, one obtains\footnote{\emph{Editorial note.} The printed source gives the intermediate term $\underline D_Y\R f_*\underline D_XF$ and cites Proposition~1.12(a)(i) for both successive arrows. This edition reproduces that chain. The mathematically coherent emendation is $\underline D_Y\R_!f\,\underline D_XF$ in the intermediate line, with Proposition~1.12(a)(ii) cited for the final arrow.}
\begin{align*}
\R f_*F &\xrightarrow{\sim}\R f_*\underline D_X\underline D_XF \\
&\xrightarrow{\sim}\underline D_Y\R f_*\underline D_XF \qquad \text{by (a)(i)}\\
&\xrightarrow{\sim}\underline D_Y\underline D_Y\R f_*F \qquad \text{by (a)(i)}.
\end{align*}
The courageous reader will convince himself that the isomorphism obtained is the canonical one; this will prove the first assertion.

Now suppose $f$ finite and $(f_*F,K_Y)$ bidualizing, since $\R f_*=f_*$ for $f$ finite. We prove that the canonical morphism $F\to\underline D_X^2F$ is an isomorphism. By Lemma 1.14(a) below, it suffices to show that the induced morphism
\[
   f_*F\longrightarrow f_*\underline D_X^2F
\]
is an isomorphism. But by applying (a)(i) twice, one convinces oneself that this morphism is precisely the canonical isomorphism
\[
   f_*F\longrightarrow \underline D_Y^2f_*F,
\]
which completes the proof. \qedtext
\end{prooftext}

\begin{statement}{Lemma 1.14}
Let $f:X\to Y$ be a finite morphism of preschemes.
\begin{enumerate}[label=\alph*)]
\item For a morphism $u:F\to G$ in $\D(X)$, $u$ is an isomorphism if and only if $f_*u$ is an isomorphism.
\item For $F\in\ob\D(X)$, one has $F\in\ob\D_c(X)$ if and only if $f_*F\in\ob\D_c(Y)$.
\end{enumerate}
\end{statement}

\begin{prooftext}
For (a), since $f_*$ is exact, one may restrict to the case where $F$ and $G$ are concentrated in degree zero. The assertion then follows immediately from the computation of the fibres of $f_*F$ and $f_*G$, SGA A, Exposé VIII, 5.5.

For (b), again one may restrict to the case where $F$ is concentrated in degree zero. Necessity is known, SGA A, Exposé IX, 2.14. We prove sufficiency: $f_*F$ constructible implies $F$ constructible. The question is local upstairs, hence a fortiori downstairs, so we may suppose $Y$ and hence $X$ affine. By EGA IV, 17.16.6, there exists a finite family $(Y_i)$ of affine subpreschemes of $Y$, pairwise disjoint and with union $Y$, such that, if $X_i=X\times_YY_i$, the induced morphism $f_i:X_i\to Y_i$ factors as
\[
   X_i\xrightarrow{u_i}X'_i\xrightarrow{v_i}Y_i,
\]
where $u_i$ is finite radicial and surjective and $v_i$ is finite étale. Put $F|_{X_i}=F_i$. By the base-change theorem for finite morphisms, SGA A, Exposé VIII, 5.5, one has
\[
   f_{i*}(F_i)\xrightarrow{\sim}(f_*F)|_{Y_i},
\]
and by SGA A, Exposé IX, 2.8, one is reduced to proving the assertion for $(f_i,F_i)$. Put $F'_i=u_{i*}(F_i)$, so that $f_{i*}(F_i)=v_{i*}(F'_i)$. It is clear from the definition of constructibility, respectively follows from SGA A, Exposé VIII, 1.1, that $F_i$, respectively $F'_i$, is constructible if and only if $F'_i$, respectively $v_{i*}(F'_i)$, is constructible. This proves the claim. \qedtext
\end{prooftext}

\begin{statement}{Corollary 1.15}
Let $f:X\to Y$ be a separated quasi-finite morphism of noetherian preschemes. Let $K_Y\in\ob\D^+(Y)$, and put, as in Proposition 1.12,
\[
   K_X=\R^!f(K_Y),\qquad \underline D_X=\underline D_{K_X},\qquad \underline D_Y=\underline D_{K_Y}.
\]
If $K_Y$ is dualizing, then $K_X$ is dualizing.
\end{statement}

\begin{prooftext}
We already know, by Proposition 1.6, that if $K_Y$ has finite quasi-injective dimension, then the same is true of $K_X$. We show that if $K_Y$ satisfies condition (ii), respectively condition (iii), of Definition 1.7, then $K_X$ satisfies the same condition.

By the Main Theorem, EGA IV, 8.12.6, $f$ factors as $f=f'i$, where $i$ is an open immersion and $f'$ is finite. Since
\[
   \R^!f\simeq\R^!i\,\R^!f',
\]
one must treat separately the case of an open immersion and the case of a finite morphism. If $f$ is an open immersion, the assertion is trivial, since a constructible sheaf on $X$ is induced by a constructible sheaf on $Y$, and formation of $\R\,\underline{\Hom}$ commutes with restriction to an open.

Suppose then that $f$ is finite. It follows from Corollary 1.13 that if $K_Y$ satisfies (iii), then $K_X$ satisfies (iii), taking account of the finiteness theorem SGA A, Exposé IX, 2.14. It remains to prove that if $K_Y$ satisfies (ii), then $K_X$ satisfies (ii). Let $F\in\ob\D_c^-(X)$. We must show that $\underline D_X(F)\in\ob\D_c^+(X)$. By the finiteness theorem, $f_*F\in\ob\D_c^-(Y)$. Hence, by the hypothesis on $K_Y$,
\[
   \underline D_Y f_*F\in\ob\D_c^+(Y).
\]
But
\[
   \underline D_Y f_*F\xrightarrow{\sim}f_*\underline D_X(F)
\]
by Proposition 1.12(a)(i), and Lemma 1.14(b) completes the proof. \qedtext
\end{prooftext}



\section*{2. Uniqueness of the dualizing complex}

In this section we suppose that $A$ is local, that is,
\[
   A=\ZZ/\ell^\nu\ZZ,
\]
where $\ell$ is prime.

\begin{statement}{Theorem 2.1, cf. [H] V.3.1}
Let $X$ be a connected locally noetherian prescheme, and let $K$ be a dualizing complex on $X$ in the sense of Definition 1.7. Let also $K'\in\ob\D_c^+(X)$. Then $K'$ is dualizing if and only if there exist an invertible $A_X$-module $L$ and an integer $r$, necessarily determined up to unique isomorphism, such that
\[
   K'\xrightarrow{\sim}K\otimes^{\mathrm L}L[r].
\]
\end{statement}

\begin{prooftext}
Let $L$ be an invertible $A_X$-module and $r$ an integer, and put $L^\vee=\underline{\Hom}(L,A_X)$. For $F\in\ob\D(X)$ and $G\in\ob\D^+(X)$ there are canonical isomorphisms
\begin{equation}
\R\,\underline{\Hom}\bigl(F\otimes^{\mathrm L}L^\vee[-r],G\bigr)
   \xrightarrow{\sim}
\R\,\underline{\Hom}\bigl(F,G\otimes^{\mathrm L}L[r]\bigr)
   \xrightarrow{\sim}
\R\,\underline{\Hom}(F,G)\otimes^{\mathrm L}L[r]. \tag{*}
\end{equation}
It follows immediately that, if $K'\simeq K\otimes^{\mathrm L}L[r]$, then $K'$ is dualizing.

Moreover, if we write $\underline D_X=\underline D$ and $\underline D_{K'}=\underline D'$, then for $F\in\ob\D(X)$ the isomorphism $(*)$ gives
\[
   \underline D'F\xrightarrow{\sim}\underline DF\otimes^{\mathrm L}L[r].
\]
Since $\underline DK\simeq A_X$, one obtains
\[
   \underline D'K\xrightarrow{\sim}L[r],
\]
which shows that $L[r]$, hence $L$ and $r$, are determined essentially uniquely by $K'$.

Conversely, suppose that $K'$ is dualizing, and put
\[
   P=\underline D'K=\underline D'\underline DA_X=\R\,\underline{\Hom}(K,K').
\]
We shall prove:
\begin{align*}
\text{(a)}\qquad K'&\xrightarrow{\sim}K\otimes^{\mathrm L}P,\\
\text{(b)}\qquad P&\xrightarrow{\sim}L[r]\quad\text{for an invertible }A_X\text{-module }L.
\end{align*}
This will complete the proof of the theorem.

First observe that, by Proposition 1.9(b)(ii), for $F\in\ob\D_c(X)$ there is an isomorphism
\begin{equation}
F\otimes^{\mathrm L}P\xrightarrow{\sim}\underline D'\underline DF. \tag{**}
\end{equation}
Applying $(**)$ to $F=K$ gives
\[
   K\otimes^{\mathrm L}P\xrightarrow{\sim}\underline D'\underline DK\xrightarrow{\sim}\underline D'A_X=K',
\]
which proves (a).

Applying $(**)$ to
\[
   F=P'=\underline D\underline D'A_X
\]
gives
\[
   P'\otimes^{\mathrm L}P
      =\underline D'\underline D\underline D\underline D'A_X
      \xrightarrow{\sim}\underline D'\underline D'A_X
      \xrightarrow{\sim}A_X.
\]
Thus (b) follows from the following lemma. \qedtext
\end{prooftext}

\begin{statement}{Lemma 2.2, cf. [H] V.3.3}
Let $X$ be a connected locally noetherian prescheme, and let $P,P'\in\ob\D_c^-(X)$ be such that
\[
   P\otimes^{\mathrm L}P'\xrightarrow{\sim}A_X.
\]
Then there exist an invertible $A_X$-module $L$ and an integer $r$ such that
\[
   P\xrightarrow{\sim}L[r],\qquad P'\xrightarrow{\sim}L^\vee[-r].
\]
\end{statement}

\begin{prooftext}
The proof has two steps.

First suppose that $X=\Spec(k)$, with $k$ separably closed. One merely repeats the easy proof given in [H]. The essential point is this: by EGA $0_{\mathrm I}$, 5.4.3, if $M$ and $M'$ are finite-type $A$-modules such that
\[
   M\otimes M'\xrightarrow{\sim}A,
\]
then $M$ and $M'$ are invertible and $M'\simeq M^\vee$. This is where the hypothesis that $A$ is local is used.

We now pass to the general case. By the special case just treated, for every $x\in X$ there exists an integer $r(x)$ such that
\[
   P_{\ol{x}}\simeq A_{\ol{x}}[r(x)],\qquad
   P'_{\ol{x}}\simeq A_{\ol{x}}[-r(x)].
\]
Equivalently, $\uH^i(P_{\ol{x}})=0$, respectively $\uH^i(P'_{\ol{x}})=0$, except for $i=-r(x)$, respectively $i=r(x)$, and the corresponding nonzero cohomology modules are isomorphic to $A_{\ol{x}}$. It remains to show that the function $x\mapsto r(x)$ is locally constant.

Assume this for the moment. Since $X$ is connected, $r$ is constant, say equal to $r_0$. Then
\[
   \uH^i(P)=0\ (i\ne-r_0),\qquad \uH^i(P')=0\ (i\ne r_0),
\]
and
\[
   \uH^{-r_0}(P)\otimes\uH^{r_0}(P')\xrightarrow{\sim}A.
\]
Put $L=\uH^{-r_0}(P)$ and $L'=\uH^{r_0}(P')$. Let $u:\ol{x}\to\ol{y}$ be a specialization morphism, and let
\[
   u^*:L_{\ol{y}}\to L_{\ol{x}},\qquad u'^*:L'_{\ol{y}}\to L'_{\ol{x}}
\]
be the corresponding specialization morphisms. There is a commutative square
\begin{equation}
\begin{tikzcd}
L_{\ol{y}}\otimes L'_{\ol{y}} \arrow[r,"\sim"] \arrow[d,"u^*\otimes u'^*"'] & A_{\ol{y}} \arrow[d] \\
L_{\ol{x}}\otimes L'_{\ol{x}} \arrow[r,"\sim"] & A_{\ol{x}}
\end{tikzcd}
\tag{*}
\end{equation}
\par\medskip
Choose a basis $e$, respectively $f$, of $L_{\ol{x}}$, respectively $L_{\ol{y}}$, and let $e'$, respectively $f'$, be the dual basis of $L'_{\ol{x}}$, respectively $L'_{\ol{y}}$, meaning that $e\otimes e'$ and $f\otimes f'$ map to $1$. If $u^*(f)=\lambda e$ and $u'^*(f')=\lambda'e'$, the commutativity of $(*)$ gives $\lambda\lambda'=1$. Hence $u^*$ is an isomorphism, and $u'^*$ is the contragredient isomorphism. By SGA A, Exposé IX, 2.13, it follows that $L$ and $L'$ are locally constant, hence of finite presentation as $A_X$-modules. Consequently, for every $x\in X$ the canonical map
\[
   \underline{\Hom}(L,F)_{\ol{x}}\longrightarrow\Hom(L_{\ol{x}},F_{\ol{x}})
\]
is an isomorphism for every $A_X$-module $F$. Since the geometric fibres of $L$ are invertible, $L$ itself is invertible. Similarly $L'$ is invertible, and $L'\simeq L^\vee$. Hence
\[
   P\simeq L[-r_0],\qquad P'\simeq L^\vee[r_0].
\]

It remains only to prove that $r$ is locally constant. For $m\in\ZZ$, the condition $r(x)=m$ is equivalent to
\[
   \uH^m(P')_{\ol{x}}\ne0.
\]
Thus, by SGA A, Exposé IX, 2.4, the function $r$ is locally constructible. Therefore, to verify that it is locally constant, EGA IV, 1.10.1, allows one to suppose that $X$ is local. Let $x$ be its closed point and $U$ the open complement.

We prove that $r$ is constant. By induction on the dimension of $X$, we may already suppose that $r(y)$ is constant and equal to $r_0$ for $y\in U$; the point is to prove $r(x)=r_0$.

Assume the contrary, and derive a contradiction. By hypothesis
\[
   P|_U\xrightarrow{\sim}L[r_0]
\]
with $L$ an invertible $A_U$-module. After translating degrees we may suppose $r_0=0$. On the other hand $r(x)\ne0$; write $r(x)=a>0$. Let
\[
   i:U\to X,
   \qquad
   j:x\to X
\]
be the canonical inclusions. Consider the exact sequence
\begin{equation}
0\longrightarrow i_!(P|_U)\longrightarrow P\longrightarrow j_*j^*P\longrightarrow0.\tag{**}
\end{equation}
By hypothesis the only nonzero cohomology objects of $P$ are $\uH^0(P)$ and $\uH^{-a}(P)$, and
\[
   \uH^0(P)\xrightarrow{\sim}i_!L.
\]
It follows that $(**)$ splits naturally, via the map $P\to\uH^0(P)$, hence
\[
   P\xrightarrow{\sim} i_!(P|_U)\oplus j_*j^*P.
\]
Therefore
\begin{align*}
P\otimes^{\mathrm L}P'
&\xrightarrow{\sim} i_!(P|_U)\otimes^{\mathrm L}P'
      \oplus j_*j^*P\otimes^{\mathrm L}P' \\
&\xrightarrow{\sim} i_!\bigl(P|_U\otimes^{\mathrm L}P'|_U\bigr)
      \oplus j_*\bigl(j^*P\otimes^{\mathrm L}j^*P'\bigr)\\
&\xrightarrow{\sim} i_!A_U\oplus j_*A_x,
\end{align*}
which is impossible, since $i_!A_U\oplus j_*A_x$ is plainly not isomorphic to $A_X$; already the sheaves $\uH^0$ are not isomorphic. This proves Lemma 2.2, and hence Theorem 2.1. \qedtext
\end{prooftext}

\section*{3. Existence of dualizing complexes}

\subsection*{3.1. Preliminaries}

As was said in the introduction, the object of this section is to prove that, on a ``good'' regular prescheme $X$, the sheaf $A_X$ is dualizing, where $n$ is assumed prime to the residual characteristics.

In fact, one would like this to hold on excellent regular preschemes of equal characteristic; but there is no doubt that the proof given below will extend to them once resolution of singularities and the purity theorem are available in that case.

We first present the various ingredients which will be used.

\begin{statement}{3.1.1. Singular locus}
Let $X$ be a prescheme, locally noetherian. We shall say that $X$ satisfies condition \emph{(reg)} if the following equivalent properties hold (EGA IV, 6.12.3):
\begin{enumerate}[label=(\roman*)]
\item for every sub-prescheme $Y$ of $X$, the set $\Reg(Y)$ of regular points of $Y$ is open in $Y$;
\item for every closed integral sub-prescheme $Y$ of $X$, the set $\Reg(Y)$ contains a non-empty open subset of $Y$.
\end{enumerate}
The condition (reg) is satisfied if $X$ is locally of finite type over a complete noetherian local ring, for example over a field (EGA IV, 6.12.8), or over a Dedekind ring whose field of fractions has characteristic zero, for example over $\ZZ$ (EGA IV, 6.12.6); more generally, it is satisfied if $X$ is excellent (EGA IV, 7.8.6).
\end{statement}

\begin{statement}{3.1.2. Dévissage of constructible sheaves}
Let $X$ be a noetherian prescheme, and let $C$ be a strictly full subcategory of the category $\operatorname{Cons}(X)$ of constructible $A_X$-modules. Suppose that $C$ is stable under direct factors and extensions. The following conditions are equivalent:
\begin{enumerate}[label=(\roman*)]
\item $C=\operatorname{Cons}(X)$;
\item $C$ contains sheaves of the form $i_!(F)$, where $i:Y\to X$ is the immersion of an integral sub-prescheme and $F$ is locally constant on $Y$;
\item $C$ contains sheaves of the form $f_*i_!(A'_U)$, where $i:U\to Y$ is an open subset of an integral prescheme, $f:Y\to X$ is a finite morphism such that the generic point of $Y$ is separable over its image, and $A'=\ZZ/\ell\ZZ$, for $\ell$ a prime divisor of $n$.
\end{enumerate}
If $X$ satisfies (reg) (3.1.1), these conditions are again equivalent to the following:
\begin{enumerate}[label=(\roman* bis)]
\item the analogue of (ii), with $Y$ moreover regular;
\item the analogue of (iii), with $U$ regular.
\end{enumerate}
If, in addition, $X$ is affine, then in (ii), respectively in (ii bis), one may suppose $Y$ closed in an open subset of the form $D(f)$, with $f\in \Gamma(X,\cO_X)$.
\end{statement}

\begin{prooftext}
The equivalence of (i)--(iii) was proved in SGA A, Exposé IX, 2.5 and 5.8. The remaining assertions follow by noetherian induction. \qedtext
\end{prooftext}

\begin{statement}{3.1.3. Cohomological dimension}
Let $X$ be a locally noetherian prescheme. We denote by $(\mathrm{cdloc})_n$ the following condition. For every sub-prescheme $Y$ of $X$, every geometric point $\ol y$ of $Y$, and every
\[
   f\in\Gamma\bigl(\wt Y(\ol y),\cO_{Y,\ol y}\bigr),
\]
where $\wt Y(\ol y)$ denotes the strict localization of $Y$ at $\ol y$ (SGA A, Exposé VIII, 7.1), one has
\[
   \cd_{\mathbf n}\bigl(\wt Y(\ol y)-V(f)\bigr)
       \le \dim \cO_{Y,\ol y}
       \quad (=\dim \cO_{Y,y}\text{ by EGA IV, 6.1.3}),
\]
where $\mathbf n$ is the set of prime divisors of $n$.

The condition $(\mathrm{cdloc})_n$ is satisfied if $X$ is locally of finite type over a field (SGA A, Exposé XIV, 3.5), or is an excellent prescheme of characteristic zero (SGA A, Exposé XIX, 6.2).\footnote{Added in 1977: or, by a recent result of G. Ofer, if $X$ is excellent, regular, and of dimension $2$.}
\end{statement}

\begin{statement}{3.1.4. Absolute cohomological purity}
Let $X$ be a regular prescheme, and let $Y$ be a regular closed sub-prescheme of $X$, of codimension $d$ at every point. One says that the pair $(X,Y)$ satisfies the absolute cohomological purity theorem, relative to $n$, if
\[
   \uH^i_Y(A_X)=0\qquad (i\ne 2d),
\]
and if the ``local fundamental class'' homomorphism (SGA $4^{1/2}$, Cycle 2.2)
\[
   A_Y\longrightarrow \uH^{2d}_Y\bigl((\mu_n)^{\otimes d}_X\bigr)
\]
is an isomorphism.

One says that the absolute cohomological purity theorem is true on $X$ if every pair $(U,Y)$ as above satisfies it, where $U$ is an open subset of $X$.

Absolute purity is true on $X$ if $X$ is smooth over a perfect field of characteristic prime to $n$ (SGA A, Exposé XVI, 3.9), or if $X$ is an excellent prescheme of characteristic zero (SGA A, Exposé XIX, 3.2), or if $X$ is regular of dimension $1$ (cf. 5.1).\footnote{Added in 1977: or, by a recent result of G. Ofer, if $X$ is excellent, regular, and of dimension $2$.}
\end{statement}

\begin{statement}{3.1.5. Hironaka-style resolution of singularities}
\begin{enumerate}[label=\alph*)]
\item Let $X$ be a regular prescheme, and let $D\ge0$ be a divisor on $X$. One says that $D$ has \emph{strict normal crossings} if there exists a finite family $(f_i)_{i\in I}$ of elements of $\Gamma(X,\cO_X)$ such that
\[
   D=\sum_{i\in I}\Div(f_i)
\]
(EGA IV, 21), and such that, for every $x\in\Supp(D)$, the germs $(f_i)_x$ which lie in $\fm_x$ form part of a regular system of parameters. One says that $D$ has \emph{normal crossings} if, locally for the étale topology, $D$ has strict normal crossings.

The standard example of a normal crossings divisor is furnished by a locally finite union of coordinate hyperplanes in affine space.

\item Let $X$ be a locally noetherian prescheme. We shall say that $X$ is \emph{strongly desingularizable} if the following condition is satisfied: for every finite morphism $f:Y\to X$, where $Y$ is integral and its generic point is separable over its image, and every regular open subset $U$ of $Y$, there exists a proper birational morphism $h:Y'\to Y$, with $Y'$ regular, such that $h$ induces an isomorphism from $U'=h^{-1}(U)$ onto $U$, and $Y'-U'$ is a normal crossings divisor.
\end{enumerate}
Here are cases where $X$ is strongly desingularizable:
\begin{enumerate}[label=\alph*)]
\item $X$ is excellent of characteristic zero;
\item $X$ is locally noetherian, regular, and of dimension $\le1$;
\item $X$ is of finite type over $\ZZ$ or over a perfect field, and of dimension $\le2$.\footnote{Added in 1977: more generally, if $X$ is excellent noetherian of dimension $\le2$, by H. Hironaka, \emph{Desingularization of excellent surfaces}, notes by B. Bennett, Bowdoin College, 1967, to be supplemented if necessary by H. Hironaka, \emph{Certain numerical characters of singularities}, J. Math. Kyoto Univ. 10 (1970), 151--187.}
\end{enumerate}
\end{statement}

\begin{statement}{3.1.6. Cohomological finiteness}
Let $f:X\to Y$ be a morphism of preschemes, and let $F\in\ob\D_c^+(X)$. One has $\R f_*(F)\in\ob\D_c^+(Y)$ in the following cases:
\begin{enumerate}[label=\alph*)]
\item $f$ is proper and of finite presentation (SGA A, Exposé XIV, 1.1);
\item ``when resolution of singularities and the purity theorem are available'', in particular:
\begin{enumerate}[label=(\roman*)]
\item if $X$ and $Y$ are locally of finite type over a field $k$ of characteristic zero and $f$ is a $k$-morphism of finite type (SGA A, Exposé XVI, 5.1);
\item if $Y$ is excellent of characteristic zero and $f$ is of finite type (SGA A, Exposé XIX, 5.1);
\item if $X$ is locally of finite type over a field and $\dim(X)\le2$ (SGA A, Exposé XIX, 5.1).
\end{enumerate}
\end{enumerate}
\end{statement}

\subsection*{3.2. The quasi-injective dimension of $A_X$}

\begin{statement}{Proposition 3.2.1}
Let $X$ be a regular prescheme whose residual characteristics are prime to $n$. Suppose that $X$ satisfies (reg) (3.1.1) and $(\mathrm{cdloc})_n$ (3.1.3), and that absolute purity is true on $X$ (3.1.4). Then, for every constructible $A_X$-module $F$ and every $x\in X$, one has
\begin{equation}
\uExt^i(F,A_X)_x=0\qquad\text{for } i>2\dim\cO_{X,x}.\tag{*}
\end{equation}
\end{statement}

\begin{prooftext}
The question being local, we may suppose $X$ affine. By dévissage (3.1.2), it is enough to verify $(*)$ for $F=i_!(G)$, where $i:Y\to X$ is the immersion of a regular integral sub-prescheme, closed in an open subset $D(f)$, $f\in A(X)$, and $G$ is a locally constant constructible sheaf on $Y$. We must compute
\[
   \uExt^q(F,A_X)=\Hh^q\bigl(\R\,\underline{\Hom}(F,A_X)\bigr).
\]
The projection formula for $i$ (SGA A, Exposé XVIII) gives
\[
 \R\,\underline{\Hom}(i_!G,A_X)
   \xrightarrow{\sim}
 \R i_*\R\,\underline{\Hom}(G,\R^!i(A_X)).
\]
By purity,
\[
   \R^!i(A_X)\xrightarrow{\sim}(\mu_n)^{\otimes -d}_Y[-2d],
\]
where $d=\codim(Y,X)$. Put $T=(\mu_n)^{\otimes -d}_Y$. Since $G$ is locally constant constructible, for every geometric point $\ol y$ of $Y$ and every $p$ one has
\[
   \uExt^p(G,T)_{\ol y}\xrightarrow{\sim}\Ext^p(G_{\ol y},T_{\ol y}).
\]
But $T_{\ol y}$ is injective, being isomorphic to $A=\ZZ/n\ZZ$. Hence $\uExt^p(G,T)=0$ for $p\ne0$, and
\[
   \R\,\underline{\Hom}(G,T)\xrightarrow{\sim}\underline{\Hom}(G,T).
\]
Thus
\[
   \R\,\underline{\Hom}(F,A_X)
      \xrightarrow{\sim} \R i_*(\mathcal H)[-2d],
   \qquad \mathcal H=\underline{\Hom}(G,T).
\]
Let $Z$ be the schematic closure of $Y$ in $X$ (EGA I, 9.5.10). Then $Y=Z\cap D(f)$ (EGA $0_{\mathrm I}$, 2.1.6). Denote by $j:Y\to Z$ and $k:Z\to X$ the canonical immersions. We have $\R i_*=k_*\R j_*$, so we are reduced to computing $\R j_*(\mathcal H)$. If $\ol x$ is a geometric point of $Z$, let $\wt Z=\Spec(\cO_{Z,\ol x})$ be the strict localization, let $\wt f$ be the image of $f$ in $\cO_{Z,\ol x}$, let $\wt Y=\wt Z\times_XY=\wt Z-V(\wt f)$, and let $\wt{\mathcal H}$ be the inverse image of $\mathcal H$ on $\wt Y$. By SGA A, Exposé VIII, 5.2,
\[
   \R^qj_*(\mathcal H)_{\ol x}
      \xrightarrow{\sim}\Hh^q(\wt Y,\wt{\mathcal H}),
\]
hence
\[
   \uExt^q(F,A_X)_{\ol x}
      \xrightarrow{\sim}\Hh^{q-2d}(\wt Y,\wt{\mathcal H}).
\]
The condition $(\mathrm{cdloc})_n$ implies that this group is zero for
\[
   q-2d>\dim\cO_{Z,x}.
\]
Using
\[
   d=\codim_x(Z,X)\le\dim\cO_{X,x}
\]
(EGA IV, 5.1.3) and
\[
   \dim\cO_{Z,x}
     =\dim\cO_{X,x}-\codim_x(Z,X)
\]
(EGA $0_{\mathrm{IV}}$, 16.5.12 and EGA IV, 5.1.9), one obtains
\[
   \uExt^q(F,A_X)_x=0\qquad(q>2\dim\cO_{X,x}),
\]
as required. \qedtext
\end{prooftext}

\begin{statement}{Remark 3.2.2}
The bound obtained in 3.2.1 is the best possible. Indeed, let $x$ be a point of $X$, and let $Y$ be the closed integral sub-prescheme of $X$ whose generic point is $x$. Since $X$ satisfies (reg), there exists a regular open subset $U$ of $Y$ containing $x$. By purity,
\[
   \uExt^{2d}(A_U,A_X)|U=(\mu_n)^{\otimes d}_U,
\]
where $d=\dim\cO_{X,x}=\codim(Y,X)$.
\end{statement}

\begin{statement}{Corollary 3.2.3}
Let $X$ be a regular prescheme of dimension $d$. If $X$ is smooth over a field $k$ of characteristic prime to $n$, or if $X$ is excellent of characteristic zero, then $X$ satisfies the conclusion of 3.2.1; consequently
\[
   \dim.\,q.\,inj(A_X)\le 2d.
\]
\end{statement}

\begin{prooftext}
If $X$ is excellent of characteristic zero, or if $k$ is perfect of characteristic prime to $n$ and $X$ is smooth over $k$, then the hypotheses of 3.2.1 are satisfied by 3.1.1, 3.1.3, and 3.1.4. The conclusion is still valid when $X$ is smooth over an arbitrary field $k$, as follows immediately after the base change $\Spec(k')\to\Spec(k)$, where $k'$ is a perfect closure of $k$ (SGA A, Exposé VIII, 1.1). \qedtext
\end{prooftext}

\begin{statement}{Lemma 3.2.4}
Let $X$ be a prescheme satisfying the conclusion of 3.2.1 and of dimension $\le d$. If there exists an integer $N$ such that $X$ has cohomological dimension $\le N$, then
\[
   \dim.\,inj(A_X)\le 2d+N.
\]
If, moreover, for every closed subset $Y$ of $X$, the implication
\[
   \codim(Y,X)\ge p \quad\Rightarrow\quad \cd_{\mathbf n}(Y)\le N-2p
\]
holds, then
\[
   \dim.\,inj(A_X)\le N.
\]
\end{statement}

\begin{prooftext}
The first assertion follows immediately from 1.5. For the second, one argues as in 1.5. We are reduced to proving, for $F$ constructible on $X$, the relation
\[
   \Ext^i(X;F,A_X)=0\qquad(i>N).
\]
For this, inspect the spectral sequence
\[
   E_2^{pq}=\Hh^p\bigl(X;\uExt^q(F,A_X)\bigr)
      \Rightarrow \Ext^*(X;F,A_X).
\]
By 3.2.1, the sheaves $\uExt^{2i}(F,A_X)$ and $\uExt^{2i-1}(F,A_X)$ are concentrated in codimension $\ge i$. Therefore $E_2^{p,2i}=E_2^{p,2i-1}=0$ for $p>N-2i$, by the hypothesis. Hence $E_2^{pq}=0$ for $p+q>N$, and the assertion follows. \qedtext
\end{prooftext}

\begin{statement}{Corollary 3.2.5}
If $X$ is a smooth prescheme of dimension $d$ over a separably closed field $k$ of characteristic prime to $n$, then
\[
   \dim.\,inj(A_X)=2d.
\]
\end{statement}

\begin{prooftext}
By 3.2.3 and SGA A, Exposé X, 4.2, we are in the situation of 3.2.4. The hypotheses of SGA A, Exposé X, 4.2 are satisfied by SGA A, Exposé X, 2.1 and 3.2; hence $\dim.\,inj(A_X)\le2d$. The equality follows from purity: if $x$ is a closed point of codimension $d$ in $X$, then
\[
   \Ext^{2d}(X,A_x,A_X)=\Hh^{2d}_x(A_X)
      \xrightarrow{\sim}(\mu_n)^{\otimes d}(k)\ne0.
\]
\qedtext
\end{prooftext}

\subsection*{3.3. Constructibility of the sheaves $\uExt^i(F,A_X)$}

\begin{statement}{Proposition 3.3.1}
Let $X$ be a prescheme.
\begin{enumerate}[label=\alph*)]
\item The following conditions are equivalent:
\begin{enumerate}[label=(\roman*)]
\item for $F\in\ob\D_c^-(X)$ and $G\in\ob\D_c^+(X)$, one has
\[
   \R\,\underline{\Hom}(F,G)\in\ob\D_c^+(X);
\]
\item[(i bis)] for every pair $(F,G)$ of constructible $A_X$-modules, one has $\R\,\underline{\Hom}(F,G)\in\ob\D_c^+(X)$, i.e. the sheaves $\uExt^i(F,G)$ are constructible for all $i$;
\item for every open immersion $i:U\to X$ and every constructible $A_U$-module $F$, the sheaves $\R^pi_*(F)$ are constructible for all $p$;
\item[(ii bis)] for every immersion $i:Y\to X$ and every constructible $A_Y$-module $F$, the sheaves $\R^pi_*(F)$ are constructible for all $p$;
\item[(ii ter)] for every immersion $i:Y\to X$ and every locally constant constructible $A_Y$-module $F$, the sheaves $\R^pi_*(F)$ are constructible for all $p$;
\item for every finite morphism, respectively every closed immersion, $f:Y\to X$, and every $F\in\ob\D_c^+(X)$, one has $\R^!f(F)\in\ob\D_c^+(Y)$.
\end{enumerate}
\item These conditions are satisfied in each of the following two cases:
\begin{enumerate}[label=(\roman*)]
\item $\dim(X)\le1$;
\item resolution of singularities and purity are available on $X$ in the sense of 3.1.6(b).
\end{enumerate}
\end{enumerate}
\end{statement}

\begin{prooftext}
The equivalence of (i) and (i bis) follows from the standard spectral sequence passing from sheaf-Ext to hyper-Ext.

(i bis) $\Rightarrow$ (ii). Let $\ol F$ be a constructible $A_X$-module extending $F$, for example $\ol F=i_!F$. By duality (SGA A, Exposé XVIII, 3.1.10),
\[
   \R i_*F
    =\R i_*\R\,\underline{\Hom}(A_U,i^!\ol F)
    \xrightarrow{\sim}
      \R\,\underline{\Hom}(i_!A_U,\ol F),
\]
and this proves the claim.

(ii) $\Rightarrow$ (ii bis). Factor $i$ as a closed immersion followed by an open immersion.

(ii bis) $\Rightarrow$ (i bis). The question is local, so we may suppose $X$ affine. By the usual dévissage (SGA A, Exposé IX, 2.5), it is enough to verify (ii) for $F=i_!M$, where $i:Y\to X$ is the immersion of a sub-prescheme and $M$ is a locally constant constructible $A_Y$-module. By the projection formula (SGA A, Exposé XVIII),
\[
   \R\,\underline{\Hom}(i_!M,G)
      \xrightarrow{\sim}
   \R i_*\R\,\underline{\Hom}(M,\R i^!G),
\]
and by a spectral sequence we are reduced to proving that
\[
   \R\,\underline{\Hom}(M,\R i^!G)\in\ob\D_c^+(Y).
\]
First prove $\R i^!G\in\ob\D_c^+(Y)$. After factoring $i$, we may suppose that $Y$ is closed. Let $j:U=X-Y$ be the complementary open immersion. Applying $\R\,\underline{\Hom}(\ ,G)$ to
\[
   0\longrightarrow A_{U,X}\longrightarrow A_X\longrightarrow A_Y\longrightarrow0
\]
gives a distinguished triangle.
\[
\begin{tikzcd}[column sep=large,row sep=large]
& \R\,\underline{\Hom}(A_{U,X},G) \arrow[dl] & \\
\R\,\underline{\Hom}(A_Y,G) \arrow[rr] & & \R\,\underline{\Hom}(A_X,G) \arrow[ul]
\end{tikzcd}
\]
Since $\R\,\underline{\Hom}(A_X,G)=G$, and
\[
   \R\,\underline{\Hom}(A_{U,X},G)\xrightarrow{\sim}\R j_*j^!G
\]
belongs to $\D_c^+(X)$ by hypothesis, it follows that
\[
   \R\,\underline{\Hom}(A_Y,G)\xrightarrow{\sim}i_*\R i^!G
\]
belongs to $\D_c^+(X)$. Resolving $M$ locally on $Y$, for the étale topology, by free $A_Y$-modules, another dévissage gives
\[
   \R\,\underline{\Hom}(M,\R i^!G)\in\ob\D_c^+(Y).
\]
This proves the implication.

(ii bis) $\Rightarrow$ (ii ter) is trivial.

(ii ter) $\Rightarrow$ (ii). Let $i:U\to X$ be an open immersion and let $F$ be a constructible $A_U$-module. We must prove $\R i_*F\in\ob\D_c^+(X)$. The question being local, suppose $X$ noetherian and argue by noetherian induction on $X$. We may suppose $U\ne\varnothing$. Then there exists a non-empty open subset $j:V\to U$ such that $j^*F$ is locally constant (SGA A, Exposé IX, 2.4). Let $M$ be the mapping cylinder of
\[
   F\longrightarrow \R j_*j^*F.
\]
Applying $\R i_*$ gives a distinguished triangle
\[
\begin{tikzcd}[column sep=large,row sep=large]
& \R i_*(M) \arrow[dl,"{(+1)}"'] & \\
\R i_*(F) \arrow[rr] & & \R i_*\R j_*j^*(F) \arrow[ul]
\end{tikzcd}
\]
Condition (ii ter) implies
\[
   \R i_*\R j_*j^*F\xrightarrow{\sim}\R(ij)_*j^*F\in\ob\D_c^+(X).
\]
Consequently,
\[
   \R j_*j^*F\xrightarrow{\sim}i^*\R(ij)_*j^*F\in\ob\D_c^+(U).
\]
Thus $M\in\ob\D_c^+(U)$, since $F\in\ob\D_c^+(U)$. Moreover, the cohomology of $M$ is supported on a closed subset $Z$ of $U$ distinct from $U$. Applying the induction hypothesis to $U\cap\ol Z\to\ol Z$, we deduce $\R i_*M\in\ob\D_c^+(X)$, and hence $\R i_*F\in\ob\D_c^+(X)$.

The equivalence of (i)--(ii ter) with condition (iii) is left to the reader; use 1.14.

For b)(i), normalization reduces one to the case where $X$ is regular, and one uses purity, which will be proved in 5.1. For b)(ii), see 3.1.6(b). \qedtext
\end{prooftext}

\subsection*{3.4. The biduality theorem on regular preschemes}

\begin{statement}{Theorem 3.4.1, cf. [H], V, 2.2}
Let $X$ be a regular locally noetherian prescheme of finite dimension. Suppose that $X$ is strongly desingularizable (3.1.5), satisfies (reg) (3.1.1) and $(\mathrm{cdloc})_n$ (3.1.3), and that absolute purity (3.1.4) is true on every prescheme smooth over $X$. Then $A_X$ is dualizing.
\end{statement}

\begin{prooftext}
We must prove that $A_X$ satisfies conditions (i)--(iii) of 1.6. Condition (i) follows from 3.2.1. Condition (ii) follows from 3.3.1(b). It remains to prove condition (iii). Thus, by 1.7, for every constructible $A_X$-module $F$ we must prove that the canonical homomorphism
\[
   F\longrightarrow \underline D_X\underline D_XF,
   \qquad \underline D_X=\R\,\underline{\Hom}(\ ,A_X),
\]
is an isomorphism. The question being local, suppose $X$ noetherian. By dévissage (3.1.2, iii bis), it is enough to treat sheaves
\[
   F=f_*i_!(A'_U),
\]
where $i:U\to Y$ is the inclusion of a regular open subset of an integral prescheme, $f:Y\to X$ is finite and the generic point of $Y$ is separable over its image, and $A'=\ZZ/\ell\ZZ$ for a prime divisor $\ell$ of $n$.

Because $X$ is strongly desingularizable, there is a proper birational morphism $h:Y'\to Y$, with $Y'$ regular, inducing an isomorphism from $U'=h^{-1}U$ to $U$, and such that $Y'-U'$ is a normal crossings divisor. By base change (SGA A, Exposé XVII),
\[
   i_!(A'_U)\xrightarrow{\sim}\R h_*i'_!(A'_{U'}),
\]
where $i':U'\to Y'$ is the inclusion. Hence, with $g=fh$,
\[
   f_*i_!(A'_U)
     \xrightarrow{\sim} f_*\R h_*i'_!(A'_{U'})
     \xrightarrow{\sim} \R g_*i'_!(A'_{U'}).
\]
By 1.11, it suffices to prove that the pair
\[
   \bigl(i'_!(A'_{U'}),\R^!g(A_X)\bigr)
\]
is bidualizing. Locally immersing $Y'$ in a prescheme $X'$ smooth over $X$, and applying purity on $X'$, one sees that $\R^!g(A_X)$ is locally isomorphic to $A_{Y'}$, up to a translation of degrees. We are therefore reduced to the following special case.
\end{prooftext}

\begin{statement}{Lemma 3.4.2}
Let $X$ be a locally noetherian regular prescheme on which absolute purity is true. Let $Y=Y_1\cup\cdots\cup Y_p$ be a strict normal crossings divisor on $X$, and put $U=X-Y$. Then the pair $(A'_{U,X},A_X)$ is bidualizing.
\end{statement}

\begin{prooftext}
By the exact sequence
\[
   0\longrightarrow A'_{U,X}\longrightarrow A'_X\longrightarrow A'_Y\longrightarrow0,
\]
it is equivalent to verify that $(A'_Y,A_X)$ is bidualizing. We argue by induction on $p$.

If $p=1$, let $i:Y\to X$ be the immersion. By 1.11, it is enough to verify that
\[
   (A'_Y,\R^!i(A_X))
\]
is bidualizing. By purity,
\[
   \R^!i(A_X)\xrightarrow{\sim}(\mu_n)^{\otimes -1}_Y[-2].
\]
Thus the canonical homomorphism $A'_Y\to\underline D_Y\underline D_Y(A'_Y)$ is an isomorphism, as is checked trivially fibre by fibre, using Pontryagin duality for $A$-modules.

Suppose $p\ge2$, and put $Y'=Y_1\cup\cdots\cup Y_{p-1}$. We have an exact sequence
\[
   0\longrightarrow A'_{(Y-Y'),X}\longrightarrow A'_Y\longrightarrow A'_{Y'}\longrightarrow0.
\]
By the induction hypothesis, $(A'_{Y'},A_X)$ is bidualizing. By the five lemma, we are reduced to biduality for $(A'_{(Y-Y'),X},A_X)$. But again there is an exact sequence
\[
   0\longrightarrow A'_{(Y-Y'),X}\longrightarrow A'_{Y_p}\longrightarrow A'_{Y_p\cap Y'}\longrightarrow0.
\]
Since $(A'_{Y_p},A_X)$ is bidualizing, it remains to prove the same for $(A'_{Y_p\cap Y'},A_X)$. Applying 1.11 again and using purity, we are finally reduced to biduality for $(A'_{Y_p\cap Y'},A_{Y_p})$ on $Y_p$, where the induction hypothesis applies. \qedtext
\end{prooftext}

This completes the proof of Theorem 3.4.1.

\begin{statement}{Corollary 3.4.3}
Let $S$ be an excellent noetherian regular prescheme of characteristic zero. Then $A_S$ is dualizing. Moreover, if $f:X\to S$ is a morphism of finite type, then $\R^!f(A_S)$ is dualizing.
\end{statement}

\begin{prooftext}
The first assertion follows because $S$ satisfies the hypotheses of 3.4.1, as recalled in 3.1. For the second, one must first verify that $\R^!f(A_S)$ has finite quasi-injective dimension. Cover $S$ by finitely many affine open subsets $S_i$ such that $X_i=f^{-1}(S_i)$ is a finite union of affine open subsets $X_{ij}$ immersed as closed subschemes in preschemes $X'_{ij}$ smooth of finite dimension over $S_i$. From 3.2.3 and 1.6 it follows that
\[
   \dim.\,q.\,inj\bigl(\R^!f(A_S)\bigr)
       \le 2\sup \dim X'_{ij}.
\]
It remains to verify conditions (ii) and (iii) of 1.7. Since these are local on $X$ for the étale topology, we may suppose that $f$ factors as $f=f'i$, where $f':X'\to S$ is smooth of relative dimension $d$, and $i:X\to X'$ is a closed immersion. Then
\[
   \R^!f(A_S)\simeq \R^!i\,\R^!f'(A_S).
\]
But $\R^!f'(A_S)$ is locally isomorphic to $A_{X'}$, up to the translation by $2d$, hence is dualizing by the first assertion. By 1.15, $\R^!f(A_S)$ is dualizing, and in particular satisfies conditions (ii) and (iii) of 1.7. \qedtext
\end{prooftext}

\begin{statement}{Corollary 3.4.4}
Let $k$ be a field, and let $f:X\to S=\Spec(k)$ be a morphism locally of finite type, with $\dim(X)\le2$. Then $\R^!f(A_S)$ is dualizing.
\end{statement}

\begin{prooftext}
Replacing $k$ by its perfect closure, which is harmless for the étale topology (SGA A, Exposé VIII, 1.1), we may suppose $k$ perfect. Cover $X$ by affine open subsets $X_i$ of finite type over $k$. Since $\dim(X_i)\le2$, the normalization lemma makes $X_i$ finite over a scheme $X'_i$ smooth over $S$ and of dimension $\le2$. By 1.15 this reduces us to the case where $f$ is smooth, hence locally
\[
   f^!(A_S)\simeq A_X[2d].
\]
Thus it remains to show that $A_X$ is dualizing in this case; this is exactly 3.4.1, since the hypotheses are satisfied (see 3.1). \qedtext
\end{prooftext}

\begin{statement}{Remark 3.4.5}
We shall see below that if $X$ is regular of dimension $\le1$, then $A_X$ is dualizing.
\end{statement}

\section*{4. Local duality}

The object of this section is to present the notion of ``biduality'' developed in the preceding sections in a form closer to geometric intuition.

\subsection*{4.1}
Let $X$ be a prescheme, let $i:x\to X$ be a closed point whose residue field is separably closed, and let $K_X$ be a dualizing complex on $X$ (1.7). We know (1.15) that $K_{\{x\}}=\R^{!}i(K_X)$ is a dualizing complex on $x$, hence, by (2.1) and (3.4.4), isomorphic to $A[r]$ for some $r\in\Z$ (if $A$ is local). We shall say that $K_X$ is \emph{normalized at $x$} if $r=0$ and if an isomorphism $K_{\{x\}}\xrightarrow{\sim}A$ has been given.

For example, if $X$ is smooth of relative dimension $d$ over the spectrum of a separably closed field, the complex $(\mu_n)_X^{\otimes d}[2d]$ is, by the purity theorem (SGA~XVI~3, SGA~$4^{1/2}$ Cycle~2.3), canonically normalized at every closed point of $X$.

\subsection*{4.2}
Let $X$ be a prescheme, $i:Y\to X$ a closed sub-prescheme, and $K_X$ a dualizing complex on $X$. Then (1.15) $K_Y=\R^{!}i(K_X)$ is a dualizing complex on $Y$, and one has the \emph{complementary induction formula} (1.12 b) (ii)), for $F\in \ob\D_c^b(X)$,
\[
(4.2.1)\qquad i^*\underline{D}_X(F)\xrightarrow{\sim}\underline{D}_Y\R^{!}i(F),
\]
where $\underline{D}_X=\R\,\underline{\Hom}(\ ,K_X)$ and $\underline{D}_Y=\R\,\underline{\Hom}(\ ,K_Y)$.

Apply (4.2.1) in particular in the case where $Y=\{x\}$ is a closed point of $X$ with separably closed residue field, $K_X$ being normalized at $x$ (4.1). Taking into account that $A_x$ is injective, (4.2.1) is then written
\[
(4.2.2)\qquad \underline{D}_X(F)_x\xrightarrow{\sim}\Hom^*(\R^{!}i(F),A).
\]
Noting that $\R^{!}i(F)$ is also written, by definition, $\R\Gamma_x(F)$, one therefore has a perfect pairing in $\D(A)$:
\[
(4.2.2)'\qquad \underline{D}_X(F)_x\times \R\Gamma_x(F)\longrightarrow A,
\]
giving rise to perfect pairings between finite groups:
\[
(4.2.2)''\qquad \uH^i(\underline{D}_X(F))_x\times \uH_x^{-i}(F)\longrightarrow A.
\]
For example, let $X$ be a prescheme smooth of dimension $d$ over a separably closed field, and let $x\in X$ be a closed point. Then, for $F\in\ob\D_c^b(X)$, one has a canonical perfect duality between finite groups:
\[
(4.2.2)'''\qquad \Ext^{2d-i}(F,(\mu_n)_X^{\otimes d})_x\times \uH_x^i(F)\longrightarrow A.
\]

\subsection*{Exercise 4.2.3}
Let $X$ be a strictly local scheme (SGAA~VIII~4.2), with closed point $x$, and let $i:x\to X$ be the canonical morphism. Let $K$ be a dualizing complex on $X$, normalized at $x$. Denote by
\[
\mathbf L\, i_? : \D^b(A)\longrightarrow \D^-(X)
\]
the functor defined by
\[
\mathbf L\, i_?(M)=M\overset{\mathbf L}{\otimes}_A K.
\]
(If $K$ has finite tor-dimension over $A$, $\mathbf L\, i_?$ extends to a functor from $\D^+(A)$ to $\D(X)$.)

Show that there exists, for $L\in\D^-(X)$ and $M\in\D^b(A)$, a canonical functorial isomorphism
\[
(*)\qquad \R\Hom(L,\mathbf L\, i_?(M))\xrightarrow{\sim}\R\Hom(\R^{!}i(L),M),
\]
(assuming of course that $\R^{!}i:\D^-(X)\to\D^-(A)$ is defined \dots).

\emph{Hint:} First define a ``trace'' morphism $\R^{!}i\,\mathbf L\, i_?(M)\to M$; then, to check that $(*)$ is an isomorphism, reduce by dévissage to the case where $L\in\ob\D_c^b(X)$ and $M=A$, and apply (4.2.2). Note that, for $M=A$, $(*)$ reduces to the perfect pairing (``local duality'', cf.~[H]~V~6.2)
\[
\R\Gamma_x(L)\times \R\Hom(L,K)\longrightarrow A,
\]
or
\[
\H_x^i(L)\times \Ext^{-i}(L,K)\longrightarrow A.
\]

\subsection*{4.3}
We shall now generalize (4.2.2) to the case of a geometric point $\ol{x}$ of $X$ localized at a point $x$ which is not necessarily closed. For this we need a few preliminaries.

\begin{statement}{Lemma 4.3}
Let $X$ be a prescheme, $X'$ the strict localization of $X$ at a geometric point $\ol{x}$, and $f:X'\to X$ the canonical map. Then, for $F\in\ob\D_c^-(X)$ and $G\in\ob\D^+(X)$, the canonical functorial homomorphism (SGAA~VI~I~3.7)
\[
f^*\R\,\underline{\Hom}(F,G)\longrightarrow \R\,\underline{\Hom}(f^*F,f^*G)
\]
is an isomorphism.
\end{statement}

\begin{prooftext}[Proof]
As usual, one reduces to $F$ and $G$ concentrated in degree $0$. Since the question is local in a neighborhood of $x$, one may suppose $X$ affine. Then $F$ admits (SGAA~IX~2.7) a left resolution by sheaves of the form $i_!(A_U)$, where $i:U\to X$ is étale of finite presentation; hence, by dévissage, one is reduced to $F=i_!(A_U)$. Form the cartesian square
\[
\begin{tikzcd}
U \arrow[d,"i"'] & U' \arrow[l,"g"'] \arrow[d,"i'"] \\
X & X' \arrow[l,"f"]
\end{tikzcd}
\]
One has canonical isomorphisms, trivial special cases of duality (SGAA~XVIII~3.1.10):
\[
\R\,\underline{\Hom}(F,G)\xleftarrow{\sim}\R i_* i^*G,
\qquad
\R\,\underline{\Hom}(f^*F,f^*G)\xleftarrow{\sim}\R i'_* g^*i^*G.
\]
Under these isomorphisms, the canonical arrow
\[
f^*\R\,\underline{\Hom}(F,G)\longrightarrow \R\,\underline{\Hom}(f^*F,f^*G)
\]
is identified with the base-change arrow (SGAA~XVII)
\[
f^*\R i_*i^*(G)\longrightarrow \R i'_*g^*i^*(G).
\]
But this is an isomorphism, as one sees at once by passing to the limit (SGAA~VII~5.11).
\end{prooftext}

\begin{statement}{Lemma 4.4}
Let $X$ be a prescheme, and let $K\in\ob\D_c^+(X)$ be a complex satisfying conditions (i) and (ii) of (1.7). In order that $K$ be dualizing, it is necessary and sufficient that, for every strict localization $f:\ol{X}(\ol{x})\to X$, $f^*K$ be dualizing.
\end{statement}

\begin{prooftext}[Proof]
Use the fact that every constructible sheaf on a strict localization $\ol{X}(\ol{x})$ is induced by a constructible sheaf on an étale $\ol{x}$-pointed scheme over $X$ (SGAA~IX~2.7.4), and apply (4.3).
\end{prooftext}

\subsection*{4.5}
Let $X$ be a prescheme, $\ol{x}$ a geometric point of $X$, and $f:X'=\ol{X}(\ol{x})\to X$ the corresponding strict localization.

Let $Y$ be the integral closed sub-prescheme which is the closure of $x$ in $X$; denote by $g:\ol{x}\to Y$ and $i:Y\to X$ the canonical maps, and form the cartesian square
\[
\begin{tikzcd}
Y \arrow[r,"i"] & X \\
\ol{x} \arrow[u,"g"] \arrow[r,"i'"'] & X' \arrow[u,"f"']
\end{tikzcd}
\]
($\ol{x}$ is both the closed point of $X'$ and the strict localization of $Y$ at $\ol{x}$).

We shall denote by
\[
\R\Gamma_{\ol{x}}:\D^+(X)\longrightarrow \D^+(\ol{x})
\]
the functor defined by
\[
(4.5.1)\qquad \R\Gamma_{\ol{x}}=g^*\R^{!}i.
\]
Using the well-known relation between the functors $\R^{!}i$ and $\R j_*j^*$, where $j:U\to X$ is the open complement of $Y$, and applying passage to the limit (SGAA~VII~5.11), one obtains a canonical isomorphism
\[
(4.5.2)\qquad \R\Gamma_{\ol{x}}\xrightarrow{\sim}\R^{!}i' f^*.
\]
When $x$ is a closed point with separably closed residue field, the functor $\R\Gamma_{\ol{x}}$ coincides with the usual functor $\R\Gamma_x$ (SGAA~V~4.3), which justifies the notation.

Let $K_X\in\ob\D^+(X)$, and put $K_Y=\R^{!}i(K_X)$, $K_{X'}=f^*K_X$, and $K_{\ol{x}}=\R\Gamma_{\ol{x}}(K_X)$ (warning: $K_{\ol{x}}$ is not the fiber of $K_X$ at $\ol{x}$!). Suppose that $K_X$ is dualizing. Then, by (1.15) and (4.4), the same is true of $K_{X'}$, $K_Y$, and $K_{\ol{x}}$. Hence (2.1) one has $K_{\ol{x}}\simeq A[d]$ for some $d\in\Z$.

We shall say that $K_X$ is \emph{normalized at $\ol{x}$} if $d=0$ and if an isomorphism $K_{\ol{x}}\simeq A$ has been chosen. In the case where $x$ is closed with separably closed residue field, this recovers the notion introduced in (4.1).

Let $K_X$ be a dualizing complex, and put $\underline{D}_X=\R\underline{\Hom}(\ ,K_X)$, $\underline{D}_Y=\R\underline{\Hom}(\ ,K_Y)$, etc. Let $F\in\ob\D_c^b(X)$; let us compute $(\underline{D}_XF)_{\ol{x}}$. One has:
\[
\begin{aligned}
(\underline{D}_XF)_{\ol{x}}&\xrightarrow{\sim}g^*i^*\underline{D}_X F\\
&\xrightarrow{\sim}g^*\underline{D}_Y\R^{!}iF\qquad\text{(by the complementary induction formula (4.2.1))}\\
&\xrightarrow{\sim}\underline{D}_{\ol{x}}g^*\R^{!}i(F)\qquad\text{(by (4.3)), that is,}
\end{aligned}
\]
\[
(4.5.3)\qquad (\underline{D}_X F)_{\ol{x}}\xrightarrow{\sim}\underline{D}_{\ol{x}}\R\Gamma_{\ol{x}}(F).
\]
Of course, one could have made an analogous computation by passing through $X'$ instead of $Y$; one would then have obtained a canonical isomorphism
\[
(\underline{D}_XF)_{\ol{x}}\xrightarrow{\sim}\underline{D}_{\ol{x}}\R^{!}i'f^*(F),
\]
compatible with (4.5.3) by means of the identification (4.5.2).

If moreover $K_X$ is normalized at $\ol{x}$, one obtains perfect pairings generalizing (4.2.2)' and (4.2.2)'':
\[
(4.5.3)'\qquad \underline{D}_X(F)_{\ol{x}}\times \R\Gamma_{\ol{x}}(F)\longrightarrow A,
\]
\[
(4.5.3)''\qquad \uH^i(\underline{D}_X(F)_{\ol{x}})\times \uH_{\ol{x}}^{-i}(F)\longrightarrow A,
\]
where $\uH_{\ol{x}}^i=\uH^i\R\Gamma_{\ol{x}}$.

Conversely, let $K_X\in\ob\D^+(X)$ satisfy conditions (i) and (ii) of (1.7), and suppose that $K_{\ol{x}}=\R\Gamma_{\ol{x}}(K_X)$ is dualizing for every geometric point $\ol{x}$ of $X$. Then, for every $\ol{x}$, one can define (cf.~1.12 b) (ii)) a canonical arrow
\[
(4.5.3)^{\mathrm{bis}}\qquad \underline{D}_X(F)_{\ol{x}}\longrightarrow \underline{D}_{\ol{x}}\R\Gamma_{\ol{x}}(F)\qquad (F\in\ob\D_c^b(X)).
\]
If this is an isomorphism for every $\ol{x}$ and every $F$, then $K_X$ is dualizing. Indeed, it remains to verify that the canonical arrow
\[
F\longrightarrow \underline{D}_X\underline{D}_X(F)
\]
is an isomorphism for every $F\in\ob\D_c^b(X)$. At an arbitrary geometric point $\ol{x}$, one has
\[
\begin{aligned}
F_{\ol{x}}&\xrightarrow{\sim}\underline{D}_{\ol{x}}\underline{D}_{\ol{x}}(F_{\ol{x}})\qquad\text{(since $K_{\ol{x}}$ is dualizing)}\\
&\xrightarrow{\sim}\underline{D}_{\ol{x}}\R\Gamma_{\ol{x}}\underline{D}_X(F)\qquad\text{(by (4.3) and the induction formula 1.12 b) (i))}\\
&\xrightarrow{\sim}\underline{D}_X\underline{D}_X(F)_{\ol{x}}\qquad\text{(by hypothesis)},
\end{aligned}
\]
and therefore, modulo a compatibility check, $(F,K_X)$ is bidualizing, which proves the assertion. In summary:

\begin{statement}{Proposition 4.5.4}
Let $X$ be a prescheme, and let $K_X\in\ob\D^+(X)$ be a complex satisfying conditions (i) and (ii) of (1.7). In order that $K_X$ be dualizing, it is necessary and sufficient that, for every geometric point $\ol{x}$ of $X$, the complex $K_{\ol{x}}$ be dualizing and the canonical arrow (4.5.3)$^{\mathrm{bis}}$ be an isomorphism for every $F\in\ob\D_c^b(X)$.
\end{statement}

\begin{statement}{Remark 4.5.5}
The existence of dualizing complexes should be regarded, in a sense, as an important complement to the global duality theorem. Indeed, the latter expresses that, for a separated morphism of finite type $f:X\to Y$ ($Y$ quasi-compact), one has a canonical isomorphism
\[
\underline{D}_Y\R_!f(F)\xrightarrow{\sim}\R f_*\underline{D}_X(F)
\]
(with the notation of 1.12). Its use for computing $\R_!f(F)$ (when $F\in\ob\D_c^b(X)$) requires information about $\underline{D}_X(F)$. When $K_X$ is dualizing, this information is supplied by the isomorphism (4.5.3), which makes possible, at least in theory, a computation of the fibers of $\underline{D}_X(F)$ as invariants of ``purely spatial'' cohomology.
\end{statement}

\begin{statement}{Exercise 4.5.6}
Let $X$ be a prescheme endowed with a dualizing complex (1.7). The family of functors (cf.~(4.5.3)'')
\[
\underline{\H}_{\ol{x}}^i:\D_c^b(X)\longrightarrow (A\text{-modules})
\]
($i$ running through $\Z$ and $\ol{x}$ through the set of geometric points of $X$) is ``conservative'', i.e. a morphism $F\to G$ in $\D_c^b(X)$ such that the maps $\H_{\ol{x}}^i(F)\to\H_{\ol{x}}^i(G)$ are isomorphisms, is an isomorphism.
\end{statement}

\subsection*{4.6}
Let $X$ be a prescheme, $j:Y\to X$ a closed sub-prescheme, and $i:\ol{x}\to Y$ a geometric point of $Y$. Let, on the other hand, $K_X$ be a dualizing complex on $X$, normalized at $\ol{x}$ (cf.~(4.5)). We saw in (4.5) that, if $Y$ is the closure of $x$, one has, for $F\in\ob\D_c^b(X)$, a perfect pairing
\[
\R^{!}j(F)_{\ol{x}}\times \underline{D}_X(F)_{\ol{x}}\longrightarrow A.
\]
It is not difficult to define an analogous pairing in the general case. Introduce the factorization of $i$ as
\[
\ol{x}\xrightarrow{i'}\ol{\{x\}}\xrightarrow{i''}Y.
\]
For $F\in\ob\D_c^b(X)$, one has canonical isomorphisms:\footnote{\emph{Editorial note.} The ordinary-induction line is not unambiguously legible in the printed source. The current source-checked French TeX reads $i'^*\underline D_{\ol{x}}i''^*\R^!j(F)$; type checking through $i'':\ol{\{x\}}\to Y$ instead suggests the subscript $\ol{\{x\}}$, while the source glyph can resemble $X$ and its first pullback mark is also unclear. This edition retains the French-control reading without claiming that the scan adjudicates it.}
\[
\begin{aligned}
\R\Gamma_{\ol{x}}j^*\underline{D}_X(F)&\xrightarrow{\sim}\R\Gamma_{\ol{x}}\underline{D}_Y\R^{!}j(F)\qquad\text{(complementary induction formula (4.2.1))}\\
&\xrightarrow{\sim}i'^*\R^{!}i''\underline{D}_Y\R^{!}j(F)\qquad\text{(definition of $\R\Gamma_{\ol{x}}$ (4.5.1))}\\
&\xrightarrow{\sim}i'^*\underline{D}_{\ol{x}}i''^*\R^{!}j(F)\qquad\text{(ordinary induction (1.12 b) (i))}\\
&\xrightarrow{\sim}\underline{D}_{\ol{x}}i^*\R^{!}j(F)\qquad\text{(by 4.3)), i.e. finally a}
\end{aligned}
\]
canonical isomorphism
\[
(4.6.1)\qquad \R^{!}j(F)_{\ol{x}}\xrightarrow{\sim}\underline{D}_{\ol{x}}\R\Gamma_{\ol{x}}j^*\underline{D}_X(F),
\]
and hence perfect pairings
\[
(4.6.1)'\qquad \R^{!}j(F)_{\ol{x}}\times \R\Gamma_{\ol{x}}j^*\underline{D}_X(F)\longrightarrow A,
\]
\[
(4.6.1)''\qquad \uH_Y^i(F)_{\ol{x}}\times \uH_{\ol{x}}^{-i}(\underline{D}_X(F)|Y)\longrightarrow A.
\]

\begin{statement}{Example 4.6.2}
Let $X$ be a regular prescheme, let $\ol{x}$ be a geometric point of $X$, and put $d=\dim(\cO_{X,\ol{x}})$. Set $K_X=(\mu_n)_X^{\otimes d}[2d]$. If $X$ satisfies condition (reg) (3.1.1) and the purity theorem (3.1.4), then $K_X$ is canonically normalized at $\ol{x}$. Thus, if $K_X$ is dualizing, one has perfect pairings
\[
\uH_Y^i(F)_{\ol{x}}\times \uH_{\ol{x}}^{2d-i}(F'|Y)\longrightarrow A,
\]
for every closed subset $Y$ containing $x$ and every $F\in\ob\D_c^b(X)$, with $F'=\R\Hom(F,(\mu_n)_X^{\otimes d})$.
\end{statement}

To finish, let us give a strictly local variant of (4.6.1).

\subsection*{4.7}
Let $X$ be a noetherian strictly local scheme (SGAA~VIII~4.2), with closed point $x$. Let $Y$ be a non-empty closed subset of $X$, and consider the canonical immersions
\[
U=X-Y\xrightarrow{i}V=X-\{x\}\xrightarrow{j}X.
\]
Finally, let $K_X$ be a dualizing complex normalized at $x$ (4.1), giving dualizing complexes ``corresponding'' on $U,V,Y,x$ (and by hypothesis $\underline{D}_x=\R\Hom(\ ,A)$):
\[
K_V=j^*K_X,
\qquad K_U=i^*K_V,
\qquad K_Y=\R\Gamma_Y(K_X).
\]
Let $G\in\ob\D_c^b(U)$, and put $G'=\underline{D}_U(G)$. We propose to define a natural pairing
\[
(4.7.1)\qquad \R\Gamma(U,G)\times \R\Gamma(V,i_!G')[-1]\longrightarrow A,
\]
i.e. a functorial homomorphism
\[
(4.7.1)'\qquad \R\Gamma(U,G)\overset{\mathbf L}{\otimes}_A\R\Gamma(V,i_!G')[-1]\longrightarrow A.
\]
(Since the ring $A=\Z/\ell\Z$ has infinite cohomological dimension for $\ell\geqslant2$, one knows how to define only the tensor products
\[
\overset{\mathbf L}{\otimes}:\D^-(A)\times\D^-(A)\longrightarrow\D^-(A)
\]
and
\[
\overset{\mathbf L}{\otimes}:\D(A)\times\D^b(A)_{\mathrm{torf}}\longrightarrow\D(A)\quad(*).
\]
We therefore leave it to the reader to make suitable hypotheses ensuring that the two factors of $\overset{\mathbf L}{\otimes}$ appearing in (4.7.1)' are in $\D^b(A)$, or that one of them has finite tor-dimension.)

First note that, by the complementary duality formula (1.12 a) (ii), one has a canonical isomorphism
\[
(4.7.2)\qquad i_!(G')\xrightarrow{\sim}\underline{D}_V(\R i_*G).
\]
On the other hand, one has the ordinary duality isomorphism
\[
(4.7.3)\qquad \R\Hom(A_U,G)\xrightarrow{\sim}\R\Hom(A_V,\R i_*G).
\]
By the definition of $\underline{D}_V$ (and independently of the fact that $K_V$ is dualizing), one has a canonical arrow
\[
(4.7.4)\qquad \R i_*G\overset{\mathbf L}{\otimes}\underline{D}_V(\R i_*G)\longrightarrow K_V.
\]
From (4.7.2), (4.7.3), and (4.7.4) one obtains a (canonical) arrow
\[
(4.7.5)\qquad \R\Gamma(U,G)\overset{\mathbf L}{\otimes}_A\R\Gamma(V,i_!G')\longrightarrow \R\Hom(A_V,K_V).
\]
But $K_V=j^!K_X$, and one has the duality isomorphism (SGAA~3.1.10):
\[
(4.7.6)\qquad \R\Hom(A_V,K_V)\xrightarrow{\sim}\R\Hom(A_{V,X},K_X).
\]
Finally, the exact sequence
\[
(4.7.7)\qquad 0\longrightarrow A_{V,X}\longrightarrow A_X\longrightarrow A_x\longrightarrow 0
\]
gives a homomorphism of degree $+1$:
\[
(4.7.8)\qquad \R\Hom(A_{V,X},K_X)\longrightarrow \R\Hom(A_x,K_X)\xrightarrow{\sim}A.
\]
Composing (4.7.5), (4.7.6), and (4.7.8), one obtains the announced arrow (4.7.1)'.

\emph{(*) (added in 1977) Inexact: in fact it is easy to extend $\overset{\mathbf L}{\otimes}:\D^-(A)\times\D^-(A)\to\D^-(A)$ to $\overset{\mathbf L}{\otimes}:\D^-(A)\times\D(A)\to\D(A)$.}

We shall now see that the pairing (4.7.1) is ``perfect'', i.e. that the arrow (4.7.1)' identifies $\R\Gamma(V,i_!G')[-1]$ with $\underline{D}_x\R\Gamma(U,G)$. Indeed suppose, as is allowed, that
\[
G=(ji)^*F,
\]
with $F\in\ob\D_c^b(X)$. Put $F'=\underline{D}_X(F)$ and $F'_{U,X}=(ji)_!G'$. One has the exact sequence
\[
(4.7.9)\qquad 0\longrightarrow F'_{U,X}\longrightarrow F'\longrightarrow F'|Y\longrightarrow0,
\]
which gives rise to a distinguished triangle
\[
\begin{tikzcd}
& \R\Gamma_x(F'_{U,X}) \arrow[dr] \\
\R\Gamma_x(F'|Y) \arrow[ur,"{+1}"] && \R\Gamma_x(F') \arrow[ll]
\end{tikzcd}
\qquad (b)
\]
By definition of $\R\Gamma_x$, one has
\[
\R\Gamma_x(F'_{U,X})=\R\Hom(A_x,F'_{U,X});
\]
hence, using (4.7.7) and the fact that $(F'_{U,X})_x=\R\Gamma(X,F'_{U,X})=0$, one has a canonical isomorphism
\[
\R\Hom(A_{V,X},F'_{U,X})[1]\xrightarrow{\sim}\R\Gamma_x(F'_{U,X}).
\]
But one has
\[
\begin{aligned}
\R\Hom(A_{V,X},F'_{U,X})&=\R\Hom(j_!A_V,j_!i_!G')\\
&\xrightarrow{\sim}\R\Hom(A_V,i_!G')\qquad\text{by duality (SGAA~3.1.10)},
\end{aligned}
\]
and hence finally a canonical isomorphism
\[
(4.7.10)\qquad \R\Gamma(V,i_!G')[-1]\xrightarrow{\sim}\R\Gamma_x(F'_{U,X}).
\]
On the other hand one has the exact sequence
\[
(4.7.11)\qquad 0\longrightarrow A_{U,X}\longrightarrow A_X\longrightarrow A_Y\longrightarrow0,
\]
which gives rise to a distinguished triangle
\[
\begin{tikzcd}
& \R\Gamma(U,F|U) \arrow[dl,"{+1}"'] \\
\R\Gamma_Y(F) \arrow[rr] && \R\Gamma_X(F) \arrow[ul]
\end{tikzcd}
\qquad (a)
\]
By virtue of (4.7.1) and (4.7.10), one has a natural pairing
\[
(4.7.12)\qquad \R\Gamma(U,F|U)\times\R\Gamma_x(F'_{U,X})\longrightarrow A.
\]
By virtue of (4.6.1), one has a perfect pairing
\[
(4.7.13)\qquad \R\Gamma_Y(F)\times\R\Gamma_x(F'|Y)\longrightarrow A.
\]
Finally, by the ordinary induction formula (1.12 b) (i)), one has a perfect pairing
\[
(4.7.14)\qquad \R\Gamma(X,F)\times\R\Gamma_x(F')\longrightarrow A.
\]
The reader who has the patience to orient himself in this maze of ``canonical arrows'' will check that the pairings (4.7.12), (4.7.13), and (4.7.14) are ``compatible'' with the arrows of the triangles (a) and (b), and consequently that the pairings (4.7.12) and (4.7.1) are perfect.

The pairing of the triangles (a) and (b) gives rise to two exact sequences ``transpose to one another'':
\[
(4.7.15)\qquad \cdots\longrightarrow \H_Y^i(X,F)\longrightarrow\H^i(X,F)\longrightarrow\H^i(U,F)\longrightarrow\cdots
\]
\[
\cdots\longleftarrow \H_x^{-i}(Y,F'|Y)\longleftarrow\H_x^{-i}(X,F')\longleftarrow\H^{-i-1}(V,F'_{U,V})\longleftarrow\cdots.
\]
In summary, we may state:

\begin{statement}{Proposition 4.7.16}
Let $X$ be a noetherian strictly local scheme, with closed point $x$. Let $Y$ be a non-empty closed subset of $X$, let $U=X-Y$, $V=X-x$, and let $i:U\to V$ and $j:V\to X$ be the canonical inclusions. Let $K_X$ be a dualizing complex on $X$ normalized at $x$ (hence giving associated dualizing complexes on $V,U,Y$). For $G\in\ob\D_c^b(U)$, one has a perfect pairing:
\[
(*)\qquad \R\Gamma(U,G)\times\R\Gamma(V,i_!G')[-1]\longrightarrow A
\]
(where $G'=\underline{D}_U(G)$). If $G=F|U$, where $F\in\ob\D_c^b(X)$, one has a canonical isomorphism
\[
(**)\qquad \R\Gamma(V,i_!G')[-1]\xrightarrow{\sim}\R\Gamma_x(F'_{U,X})
\]
(where $F'=\underline{D}_X(F)$ and $F'_{U,X}=(ji)_!(F'|U)=(ji)_!(G')$), and a perfect pairing between the distinguished triangles
\begin{center}
\begin{tikzpicture}[>=latex,scale=0.92,baseline=(current bounding box.center)]
\node (aT) at (0,1.2) {$\R\Gamma(U,F|U)$};
\node (aL) at (-1.75,0) {$\R\Gamma_Y(F)$};
\node (aR) at (1.75,0) {$\R\Gamma(X,F)$};
\draw[->] (aL) -- (aR);
\draw[->] (aR) -- (aT);
\draw[->] (aT) -- node[above left] {$+1$} (aL);
\node at (3.2,0.55) {and};
\node (bT) at (5.9,1.2) {$\R\Gamma_x(F'_{U,X})$};
\node (bL) at (4.25,0) {$\R\Gamma_x(F'|Y)$};
\node (bR) at (7.55,0) {$\R\Gamma_x(F')$};
\draw[->] (bT) -- (bR);
\draw[->] (bR) -- (bL);
\draw[->] (bL) -- node[above left] {$+1$} (bT);
\end{tikzpicture}
\end{center}
compatible with the pairing $(*)$ and the isomorphism $(**)$, and giving rise to two exact sequences transpose to one another (4.7.15).
\end{statement}

\begin{statement}{Remark 4.7.17}
In the particular case where $Y=\{x\}$, the pairing (4.7.1) is defined under the sole hypothesis that $K_X$ is normalized at $x$ (and not necessarily dualizing), i.e. such that $\R\Gamma_x(K_X)\simeq A$ (the isomorphism being given): the isomorphism (4.7.2) (which used biduality) then reduces to the identity. Since the pairing (4.7.14) is in any case perfect, one sees therefore that local duality on $X$ (in the form (4.2.2)') is equivalent to the ``global'' duality on $U=X-x$, i.e. to the assertion that the pairing (4.7.1)
\[
\R\Gamma(U,G)\times\R\Gamma(U,G')[-1]\longrightarrow A
\]
is perfect, or equivalently that the pairings
\[
\H^i(U,G)\times\H^{-i-1}(U,G')\longrightarrow A
\]
are perfect.

Suppose that one has on $X$ a dualizing complex $K_X$ normalized at $x$. If $X$ is excellent and $U$ is regular of dimension $d$ ($x$ may therefore be an isolated singularity), one should be able to check that $K_U=K_X|U$ is canonically isomorphic to $(\mu_n)_U^{\otimes d}[2d]$ (this is easily checked, at any rate, if $X$ is the strict localization of a prescheme locally of finite type over a field of characteristic~$0$). Then local duality on $X$ implies that, for every locally constant constructible $A_U$-module $G$, one has perfect pairings:
\[
\H^i(U,G)\times\H^{2d-1-i}(U,G^\circ)\longrightarrow A,
\]
where $G^\circ=\uHom(G,(\mu_n)_U^{\otimes d})$, corresponding formally to Poincaré duality on a scheme of dimension $2d-1$.
\end{statement}

\section*{5. Local duality on curves}

\begin{statement}{Theorem 5.1}
Let $X$ be a noetherian regular prescheme of dimension $1$. Suppose that $n$ is prime to the residual characteristics of $X$. Then the purity theorem is true on $X$ (3.1.4), and the complex $A_X$ is dualizing (1.7).
\end{statement}

\begin{prooftext}[Proof]
Suppose purity has been proved. Since $X$ plainly satisfies condition (reg) (3.1.1), and since condition $(\mathrm{cdloc})_n$ is verified by (SGAA~X~2.2), it follows from (3.2.1) that $\dim.q.inj(A_X)=2$. On the other hand, condition (ii) of (1.7) is verified by (3.3.1 b) (i)) (which was proved assuming purity). Hence, by (4.4), we are reduced to proving that $A_X$ is dualizing when $X$ is strictly local. But, to prove purity, which is a local question in a neighborhood of a closed point, one may equally well suppose $X$ strictly local. Finally, we are reduced to proving the theorem when $X$ is strictly local.

Thus suppose that $X$ is strictly local. Let $x$ be its closed point, let $(\pi=\car(k(x)))$, let $\eta$ be its generic point, and let $U=X-x=\Spec(k(\eta))$. Let $G=\Gal(k(\ol{\eta})/k(\eta))$ be the fundamental group of $U$. It is known [CL, Chap.~IV \S~2 Exer.~2] that one has a canonical exact sequence
\[
(5.1.1)\qquad 1\longrightarrow P\longrightarrow G\longrightarrow H\longrightarrow1,
\]
where $H=\prod_{\ell\ne\pi}\Z_\ell(1)$, $\Z_\ell(1)=\varprojlim \mu_{\ell^\nu}$, and $P$ is a pro-$\pi$-group. (Of course, $H$ is isomorphic, but not canonically, to $\prod_{\ell\ne\pi}\Z_\ell$.)

It is also known (SGAA~VIII~2) that the étale cohomology of $U$ is interpreted as the Galois cohomology of $G$, sheaves (respectively constructible sheaves) of $A_U$-modules corresponding to $A$-modules (respectively finite type $A$-modules) on which $G$ acts continuously.

This being so, the $\uH_x^i(A_X)$ are determined by the exact sequence
\[
(5.1.2)\qquad 0\longrightarrow \uH_x^0(A_X)\longrightarrow \H^0(X,A_X)\longrightarrow \H^0(U,A_U)\longrightarrow \uH_x^1(A_X)\longrightarrow0,
\]
and the isomorphisms
\[
(5.1.3)\qquad \uH_x^i(A_X)\xrightarrow{\sim}\H^{i-1}(U,A_U)\qquad\text{for }i\geqslant2.
\]
But trivially $\H^0(X,A_X)=\H^0(U,A_U)=A$, so $\uH_x^0(A_X)=\uH_x^1(A_X)=0$.

On the other hand, for a Galois $G$-$A$-module $M$, one has the Hochschild--Serre spectral sequence
\[
E_2^{pq}=\H^p(H,\H^q(P,M))\Longrightarrow \H^{*}(G,M).
\]
Since $n$ is prime to $\pi$ and $P$ is a pro-$\pi$-group, one has
\[
\H^q(P,M)=0\quad\text{for }q\ne0,
\]
hence $E_2^{pq}=0$ for $q\ne0$, giving a canonical isomorphism
\[
(5.1.4)\qquad \H^p(G,M)\xrightarrow{\sim}\H^p(H,M^P).
\]
By an analogous argument, one sees moreover that $\H^p(H,N)$ is identified with $\H^p(\widehat{\Z}(1),N)$ when the factor $\widehat{\Z}(1)$ is made to act trivially on $N$. Thus, finally, with the preceding convention, one has a canonical isomorphism
\[
(5.1.5)\qquad \H^p(G,M)\xrightarrow{\sim}\H^p(\widehat{\Z}(1),M^P).
\]
Now the Galois cohomology of $\widehat{\Z}$ is well known (CG chap.~I p.~31): for a $G$-$A$-module $N$, one has
\[
(5.1.6)\qquad
\H^i(\widehat{\Z},N)=0\quad\text{for }i\geqslant2,
\qquad
\H^0(\widehat{\Z},N)=N^{\widehat{\Z}},
\qquad
\H^1(\widehat{\Z},N)=N_{\widehat{\Z}}.
\]
From (5.1.3), (5.1.5), and (5.1.6) one immediately deduces that
\[
\H^i(U,A_U)=0\quad\text{for }i\geqslant2,
\]
and
\[
\H^1(U,A_U)\xrightarrow{\sim}A.
\]
This last isomorphism is not canonical, since it depends on the choice of an identification of $\widehat{\Z}(1)$ with $\widehat{\Z}$. But, using the Kummer exact sequence, one finds (SGAA~XIX~1.3) a canonical isomorphism (given by the ``local fundamental class'')
\[
(5.1.7)\qquad \H^1(U,\mu_n)\xrightarrow{\sim}A.
\]
Purity is therefore proved.

Let us now prove that $A_X$ is dualizing. By (4.7.17), it is a matter of showing that, for every locally constant constructible $A_U$-module $M$, the pairing
\[
(5.1.8)\qquad \H^i(U,M)\times\H^{1-i}(U,\uHom(M,A_U))\longrightarrow \H^1(U,A_U)\xrightarrow{\sim}(\mu_n)_X^{\otimes -1}
\]
is a perfect duality. We shall interpret this pairing in terms of Galois cohomology. We shall use the following elementary lemma.
\end{prooftext}

\begin{statement}{Lemma 5.1.9}
Let $D=\Hom(\ ,A)$ denote the Pontrjagin dualizing functor on the category of finite type $A$-modules. Let $M$ be a finite type $A$-module on which a group $G$ acts $A$-linearly.
\begin{enumerate}[label=(\roman*)]
\item One has a canonical isomorphism
\[
D(M^G)\xrightarrow{\sim}(DM)_G.
\]
\item If $G$ is a profinite group acting continuously on $M$, and if $G$ has pro-order prime to $n$, then the canonical arrow
\[
M^G\longrightarrow M_G
\]
is an isomorphism.
\end{enumerate}
\end{statement}

\begin{prooftext}[Proof]
By definition one has
\[
M^G=\varprojlim_{s\in G}\Ker(M\xrightarrow{s-1}M),
\]
hence
\[
D(M^G)\xrightarrow{\sim}\varinjlim_{s\in G}\Coker(DM\xrightarrow{s-1}DM)\xrightarrow{\sim}(DM)_G,
\]
which proves (i).

Let us prove (ii). Let $M'$ be the $A$-module defined by the exact sequence
\[
0\longrightarrow M'\longrightarrow M\longrightarrow M_G\longrightarrow0.
\]
One has $\H^1(G,M')=0$ because $G$ has pro-order prime to $n$. The exact cohomology sequence therefore implies that the arrow $M^G\to M_G$ is an epimorphism. But, applying the functor $D$ and taking (i) into account, one deduces that it is a monomorphism, which completes the proof.
\end{prooftext}

Applying (5.1.9) to $M$ and $P$, one finds that $\uHom(M,A)^P$ is canonically identified with $\uHom(M^P,A)$. If one chooses an isomorphism from $\widehat{\Z}(1)$ to $\widehat{\Z}$, the pairing (5.1.8) is written, taking (5.1.5) into account, as
\[
(5.1.8)'\qquad \H^i(\widehat{\Z},M^P)\times\H^{1-i}(\widehat{\Z},D(M^P))\longrightarrow\H^1(\widehat{\Z},A)\xrightarrow{\sim}A.
\]
It follows immediately from (5.1.6) and (5.1.9(i)) that (5.1.8)' is a perfect pairing. This completes the proof of (5.1).

\section*{Bibliography}

\begin{description}[leftmargin=2.5cm,style=nextline]
\item[{[2]}] Hironaka, H. \emph{Resolution of singularities of an algebraic variety over a field of characteristic zero}, Ann. of Math., vol. 79 (1964), p.~109.
\item[{[1]}] Abhyankar, S. \emph{Local uniformization on algebraic surfaces over ground fields of characteristic $p\ne0$}, Ann. of Math., vol. 63 (1956), p.~491--526, and \emph{Resolution of singularities of arithmetical surfaces}, Arithmetical Algebraic Geometry, Harper \& Row, New York (1965).
\item[{[H]}] Hartshorne, R. \emph{Residues and duality}, Lecture Notes no.~20, Springer (1966).
\item[{[SGAA]}] Séminaire de Géométrie Algébrique de l'IHES, \emph{Cohomologie étale des schémas}, by M.~Artin and A.~Grothendieck, 1963--64.
\item[{[CL]}] Serre, J.-P. \emph{Corps locaux}, Hermann (1962).
\item[{[CG]}] Serre, J.-P. \emph{Cohomologie galoisienne}, Lecture Notes no.~5, Springer (1964).
\end{description}

\section*{Appendix}
\begin{center}
by L.~Illusie
\end{center}

The preceding theory extends easily to base Rings more general than the constant Ring $\Z/n\Z$. We shall briefly indicate how the theory is written for locally constant base Rings, killed by $n>0$, and whose geometric fibres are noetherian rings. This is, moreover, only a provisional framework\footnote{According to Grothendieck!}: sooner or later one will have to globalize, by taking as ``coefficients'' on a prescheme $X$ relative preschemes over the site $(X_{\mathrm{\acute et}},\Z/n\Z)$.

\subsection*{1. Constructible sheaves}

\textbf{1.1.} Let $X$ be a prescheme, and let $A_X$ be a sheaf of rings on $X$ (for the étale topology), locally constant, whose geometric fibres are left noetherian rings. We denote by ${}_A X$ the category of left $A_X$-Modules, and by $\D({}_A X)$ (or $\D(X)$ if there is no risk of confusion) the corresponding derived category. The notion of a constructible left $A_X$-Module is defined as in (SGAA IX 2.3), by replacing the expression ``locally constant with value of finite presentation'' by ``of finite presentation as an $A_{U_i}$-Module''.

We denote by $\D_c(X)$ the full subcategory of $\D(X)$ consisting of complexes $E$ such that $\uH^i(E)$ is a constructible $A_X$-Module for every $i$. It is a triangulated subcategory of $\D(X)$.

\textbf{1.2.} Let $E\in\ob\D(X)$. We say (cf. Exp.~II 3.11) that $E$ is pseudo-coherent if, for every $i$, $\uH^i(E)$ is an $A_X$-Module of finite presentation (or locally constant constructible, if one prefers). We denote by $\D(X)_{\mathrm{coh}}$ (or $\D_{\mathrm{coh}}(X)$) the full subcategory of $\D(X)$ consisting of pseudo-coherent objects. It follows from (SGAA IX 2.1) that this is a triangulated subcategory of $\D_c(X)$.

\textbf{Proposition 1.3} (dévissage lemma). Let $X$ be a quasi-compact and quasi-separated prescheme (respectively a noetherian prescheme), and let $C$ be a strictly full subcategory of the category $\Cons(X)$ of constructible left $A_X$-Modules. Suppose that $C$ is stable under extensions. The following conditions are equivalent:
\begin{enumerate}[label=(\roman*)]
\item $C=\Cons(X)$;
\item $C$ contains the sheaves of the form $i_!(F)$, where $i:Y\to X$ is a sub-prescheme of finite presentation (respectively an integral sub-prescheme), and $F$ is locally constant constructible on $Y$.
\end{enumerate}
If $A_X$ is constant with value $A$, a left noetherian ring killed by an integer $n>0$, the preceding conditions are equivalent to:
\begin{enumerate}[label=(\roman*),start=3]
\item $C$ is stable under direct factors and contains the sheaves of the form $f_*i_!(M_U)$, where $i:U\to Y$ is an open immersion (respectively an open subset of an integral prescheme), $f:Y\to X$ is a finite morphism (respectively finite and such that the generic point of $Y$ is separable over its image), and $M_U$ is a left $A_U$-Module constant with value a monogenic $A$-module killed by a prime divisor of $n$.
\end{enumerate}
Finally, if $X$ is noetherian and satisfies condition (reg) (I 3.1.1), condition (ii) (respectively (iii)) is equivalent to the analogous condition in which one further assumes $Y$ (respectively $U$) regular.

\emph{Proof.} The equivalence of conditions (i) and (ii) is proved as in the case where $A_X$ is constant; see (SGAA IX 2.4 and 2.5).

Suppose that $A_X$ is constant with value $A$ and killed by $n$, and prove that (iii) implies (ii). Let $E\in\ob\Cons(X)$; we show that $E\in\ob C$. By the equivalence (i)$\Leftrightarrow$(ii), we may suppose that $E=j_!F$, where $j:Z\to X$ is the immersion of a sub-prescheme of finite presentation (respectively an integral one), and $F$ is a constructible locally constant $A_Z$-Module, which we may even suppose, after a little more dévissage, to be trivialized by an étale covering $g:Z'\to Z$ with group $G$. The sheaf $F$ is therefore defined by a finite type $A$-module $M$ on which $G$ acts. If $\ell$ is a prime divisor of $n$, the subgroup of $\ell$-torsion in $M$ is a sub-$G$-$A$-module $M_\ell$ of $M$, and
\[
        M=\bigoplus_\ell M_\ell,
\]
where $\ell$ runs over the set of prime divisors of $n$. Since $C$ is stable under direct factors, we are reduced to the case where $M$ is $\ell$-torsion. Let $H$ be an $\ell$-Sylow subgroup of $G$, and let $h:U=Z'/H\to Z$ be the corresponding intermediate covering. Since the degree of $h$ is prime to $\ell$, the ``trace method'' (SGAA IX 5.1) shows that the adjunction map
\[
        F\longrightarrow h_*h^*(F)
\]
is an isomorphism onto a direct factor; it is therefore enough to prove that $j_!h_*h^*(F)\in\ob C$. By the ``Main theorem'' (EGA IV 8.12.6), there exists a commutative diagram
\[
\begin{tikzcd}
U \arrow[r,"i"] \arrow[d,"h"'] & Y \arrow[d,"f"]\\
Z \arrow[r,"j"'] & X,
\end{tikzcd}
\]
where $i$ is an open immersion and $f$ is finite. We have
\[
        j_!h_*h^*(F) \simeq f_* i_! h^*(F).
\]
But $h^*(F)$ is defined by $M$ considered as an $H$-$A$-module. Thus (iii)$\Rightarrow$(ii) follows from the next lemma.

\textbf{Lemma 1.3.1.} Let $A$ be a left noetherian ring, $\ell$ a prime number, $H$ a finite $\ell$-group, and $M$ a left $H$-$A$-module of finite type killed by $\ell^\nu$ ($\nu\geq1$). There exists a finite filtration of $M$ by sub-$H$-$A$-modules such that the successive quotients are monogenic $A$-modules killed by $\ell$, and on which $H$ acts trivially.

\emph{Proof.} The exact sequence of $H$-$A$-modules
\[
        0\longrightarrow \ell M\longrightarrow M\longrightarrow M/\ell M\longrightarrow 0
\]
reduces, by induction on $\nu$, to the case where $M$ is killed by $\ell$. We show that then $M$ has a finite filtration by sub-$H$-$A$-modules on whose successive quotients $H$ acts trivially. Arguing by noetherian induction on the set of sub-$H$-$A$-modules of $M$ having the stated property, we are reduced to proving that $M\neq0$ implies $M^H\neq0$. But the kernel of the augmentation $\F_\ell[H]\to\F_\ell$ is plainly nilpotent, and Nakayama's lemma implies that every nonzero $H$-$\F_\ell$-module has a nonzero invariant vector. This proves the assertion. Thus $M$ has a finite filtration whose successive quotients are finite type $A$-modules killed by $\ell$ and on which $H$ acts trivially. By noetherian induction one may even arrange for the successive quotients to be monogenic $A$-modules, completing the proof of (1.3.1).

The last assertion of (1.3) is immediate and left to the reader.

\subsection*{2. Quasi-injective dimension}

With the notation of (1.1), one defines, as in (I 1.1 and 1.4), the notions of a quasi-injective $A_X$-Module and of quasi-injective dimension of an object of $\D(X)$. Propositions (I 1.3, 1.5 and 1.6) extend without change.

It is sometimes convenient to introduce a notion of pointwise quasi-injective dimension. Let $x\in X$ and let $K\in\ob\D^+(X)$. We say that $K$ has quasi-injective dimension $\leq N$ at $\bar{x}$ if, for every constructible $A_X$-Module $F$, one has
\[
        \underline{\Ext}^{i}(F,K)_{\bar{x}}=0\qquad\text{for }i>N.
\]
The quasi-injective dimension of $K$ at $\bar{x}$, denoted $\dim.\mathrm{q.inj}_{\bar{x}}(K)$, is the greatest lower bound of the $N$ satisfying the preceding condition. Of course,
\[
        \dim.\mathrm{q.inj}(K)=\sup_{\bar{x}}\dim.\mathrm{q.inj}_{\bar{x}}(K),
\]
where $\bar{x}$ runs through the set of geometric points of $X$.

It is interesting to compare this with the notion of strict topological dimension introduced by Verdier in [1]. We say that $K\in\ob\D^+(X)$ has strict topological dimension $\leq N$ at a geometric point $\bar{x}$ of $X$ if, for every sub-prescheme $Y$ of $X$,
\[
        \underline{\Ext}^{i}(A_Y,K)_{\bar{x}}=\uH^i_Y(K)_{\bar{x}}=0\qquad\text{for }i>N.
\]
The strict topological dimension of $K$ at $\bar{x}$, denoted $\dim.\mathrm{top}_{\bar{x}}(K)$, is the greatest lower bound of the $N$ satisfying the preceding condition. Whereas quasi-injective dimension depends a priori on the base Ring, strict topological dimension depends only on the underlying complex of abelian sheaves.

One plainly has
\[
        \dim.\mathrm{top}_{\bar{x}}(K)\leq \dim.\mathrm{q.inj}_{\bar{x}}(K).
\]
We shall see below examples where equality holds.

\subsection*{3. Dualizing complexes: definition and formal properties}

Keep the preceding notation, but suppose $A_X$ is commutative. The notions of bidualizing pair and dualizing complex are defined as in (I 1.7). The statements (I 1.9, 1.11, 1.14) remain valid without modification. The statements (I 1.12, 1.13, 1.15) are also valid, but on condition that $A$ be killed by an integer $n>0$ prime to the residual characteristics of $Y$. Finally, the uniqueness theorem (I 2.1) is valid provided one assumes $\Spec(A_{\bar{x}})$ connected for every geometric point $\bar{x}$ of $X$.

\subsection*{4. Absolute purity}

Return to the situation of absolute cohomological purity (I 3.1.4). Let $X$ be a locally noetherian regular prescheme, let $n>0$ be an integer prime to the residual characteristics of $X$, and let $i:Y\to X$ be a regular sub-prescheme of codimension $d>0$ at every point. To say that the pair $(Y,X)$ satisfies absolute cohomological purity means, recall, that the ``local fundamental class'' homomorphism (SGA $4^{1/2}$ Cycle 2.2)
\begin{equation}\tag{4.1}
        (\mu_n^{\otimes -d})_Y[-2d]\longrightarrow \R i^!(\Z/n\Z)_X
\end{equation}
is an isomorphism.

Now put on $X$ a sheaf of rings $A_X$ killed by $n$. For the moment set $G=(\Z/n\Z)_X$. Then, for every $F\in\ob\D_A^+(X)$, one can define a canonical arrow in $\D_A^+(X)$:
\begin{equation}\tag{4.2}
        i^*(F)\otimes^L_{\Z/n\Z} \R i^!(G)\longrightarrow \R i^!\bigl(F\otimes^L_{\Z/n\Z}G\bigr).
\end{equation}
Indeed, by duality (SGAA XVIII 3.1.6), this is equivalent to defining an arrow
\[
        i_!\bigl(i^*(F)\otimes^L_{\Z/n\Z}\R i^!(G)\bigr)\longrightarrow F\otimes^L_{\Z/n\Z}G;
\]
we take this to be the composite of the Künneth arrow (SGAA XVII 5.4.2.2)
\[
        i_!\bigl(i^*(F)\otimes^L_{\Z/n\Z}\R i^!(G)\bigr)\longrightarrow F\otimes^L_{\Z/n\Z}i_!\R i^!(G),
\]
and the arrow deduced from the trace homomorphism (SGAA XVIII 3.1.6)
\[
        F\otimes^L_{\Z/n\Z}i_!\R i^!(G)\longrightarrow F\otimes^L_{\Z/n\Z}G.
\]
Notice that in the definition of (4.2) only the fact that $G$ and $\R i^!(G)$ have finite Tor-dimension over $\Z/n\Z$ was used. Suppose on the other hand that the functor $i^!$ has finite cohomological dimension on the category of $(\Z/n\Z)_X$-Modules: this is the case, for example, if $X$ is of finite type over a field (SGAA XIV 3.1), or more generally if $X$ is excellent of equal characteristic (SGAA XIX 6.1) when resolution of singularities is available, or again if $X$ is of finite type over $\Spec(\Z)$ (SGAA X 6). Then the arrow (4.2) is defined for all $F\in\ob\D({}_A X)$.

Composing (4.1) and (4.2), one obtains, for $F\in\ob\D_A^+(X)$, a canonical arrow
\begin{equation}\tag{4.3}
        i^*(F)\otimes^L_{\Z/n\Z}(\mu_n^{\otimes -d})_Y[-2d]\longrightarrow \R i^!(F).
\end{equation}
Assume that $A_X$ is locally constant, and let $F\in\ob\D_A^+(X)$ be such that every $\uH^i(F)$ is a locally constant $A_X$-Module. Then, if $i^!$ has finite cohomological dimension, the arrow (4.2), and hence the arrow (4.3) if $(Y,X)$ satisfies purity, is an isomorphism.

Indeed, by dévissage, one first reduces to the case where $F$ is concentrated in degree zero. Since the question is local, one may suppose $X$ affine and $A_X$ and $F$ constant. Resolving $F$ on the left by free $A_X$-Modules and applying the hypothesis that $i^!$ has finite cohomological dimension, one reduces to the case where $F$ is free. By an immediate passage to the limit, one reduces to $F$ free of finite type, and finally to $F=A_X$. It is therefore enough to prove the assertion when $A_X=(\Z/n\Z)_X$. But the preceding dévissage reduces to $F=(\Z/n\Z)_X$, and in that case (4.2) is the identity.

\textbf{Definition 4.4.} Given the Ring $A_X$ (locally constant and killed by $n$), we shall say that the pair $(Y,X)$ satisfies cohomological purity in the strong sense if, for every $F\in\ob\D_A^+(X)$ such that the $\uH^i(F)$ are locally constant, the canonical arrow (4.3) is an isomorphism.

We shall say that the purity theorem is true in the strong sense on $X$ if, locally for the étale topology, every pair $(Y,X)$ as above satisfies purity in the strong sense. From what has been shown, for $(Y,X)$ to satisfy purity in the strong sense it is enough that $(Y,X)$ satisfy ordinary purity (I 3.1.) and that the functor $i^!$ have finite cohomological dimension.

\subsection*{5. Bound for quasi-injective dimension}

\textbf{Proposition 5.1.} Let $X$ be a locally noetherian regular prescheme, equipped with a locally constant sheaf of rings $A_X$, killed by an integer $n>0$ prime to the residual characteristics of $X$, whose geometric fibres are left noetherian rings. Suppose that $X$ satisfies conditions (reg) (I 3.1.1) and $(\mathrm{cdloc})_n$ (I 3.1.3), and that the purity theorem in the strong sense is true on $X$ (4.4) (conditions all satisfied if $X$ is smooth over a perfect field, or if $X$ is excellent of characteristic zero).

Let $F\in\ob\D_A^+(X)$ be such that the $\uH^i(F)$ are locally constant, and let $x$ be a point of $X$ such that $F_{\bar{x}}\neq0$ and $\dim.\mathrm{inj}_{A_{\bar{x}}}(F_{\bar{x}})<+\infty$. Then
\[
        \dim.\mathrm{q.inj}_{\bar{x}}(F)=2\dim(\mathcal{O}_{X,x})+\dim.\mathrm{inj}(F_{\bar{x}}),
\]
with the notation of no.~2.

\emph{Proof.} One merely repeats step by step the proof of (I 3.2.1). First show the inequality
\[
        \dim.\mathrm{q.inj}_{\bar{x}}(F)\leq 2\dim(\mathcal{O}_{X,x})+\dim.\mathrm{inj}(F_{\bar{x}})=N,
\]
that is, the relation
\begin{equation}\tag{*}
        \underline{\Ext}^{i}(E,F)_{\bar{x}}=0\quad\text{for } i>N\text{ and every constructible }A_X\text{-Module }E.
\end{equation}
Since the question is local, by dévissage (1.3) one may restrict to checking $(*)$ for $E=i_!(M)$, where $i:Y\to X$ is an integral regular sub-prescheme closed in an open $D(f)$ (with $f\in\Gamma(X,\mathcal{O}_X)$), of codimension $d$ at every point of $D(f)$, and $M$ is a locally constant constructible $A_Y$-Module ($A_Y=i^*A_X$).

Now there is the duality isomorphism (SGAA XVIII)
\[
        \R\underline{\Hom}(i_!(M),F)\simeq \R i_*\R\underline{\Hom}(M,\R i^!F).
\]
By the purity theorem, there is also the isomorphism (4.3)
\[
        i^*(F)\otimes^L_{\Z/n\Z}(\mu_n^{\otimes -d})_Y[-2d]\xrightarrow{\sim}\R i^!(F).
\]
Let $\bar{y}$ be a geometric point of $Y$. Since $M$ is locally constant constructible, one has canonically
\[
        \R\underline{\Hom}(M,\R i^!F)_{\bar{y}}\simeq \R\underline{\Hom}(M_{\bar{y}},(\R i^!F)_{\bar{y}}),
\]
hence, by (4.3),
\[
        \R\underline{\Hom}(M,\R i^!F)_{\bar{y}}\simeq \R\underline{\Hom}(M_{\bar{y}},F_{\bar{y}})[-2d].
\]
After replacing $X$ by a connected neighbourhood of $x$, one has $F_{\bar{x}}\simeq F_{\bar{y}}$ and $A_{\bar{x}}\simeq A_{\bar{y}}$, since $A_X$ and the $\uH^i(F)$ are locally constant. Thus
\[
        \underline{\Ext}^i(M_{\bar{y}},F_{\bar{y}})=0\quad\text{for }i>\dim.\mathrm{inj}(F_{\bar{x}}),
\]
and hence
\[
        \underline{\Ext}^{q}(M,\R i^!F)=0\quad\text{for }q>2d+\dim.\mathrm{inj}(F_{\bar{x}}).
\]
Set $G=\R\underline{\Hom}(M,\R i^!F)$ and compute $\R i_*(G)_{\bar{x}}$. Introduce, as in (I 3.2.1), the schematic closure $Z$ of $Y$ in $X$, giving a cartesian diagram of immersions
\[
\begin{tikzcd}
Y \arrow[r] \arrow[d,"j"'] & D(f) \arrow[d]\\
Z \arrow[r,"k"'] & X.
\end{tikzcd}
\]
Then $\R i_*(G)\simeq k_*\R j_*(G)$. Assume of course that $x\in Z$, since otherwise $\R i_*(G)=0$. Let $\bar{Z}=\Spec(\mathcal{O}_{Z,\bar{x}})$ be the strict localization of $Z$ at $\bar{x}$, let $\bar{f}$ be the image of $f$ in $\mathcal{O}_{Z,\bar{x}}$, let $\bar{Y}=\bar{Z}\times_ZY=\bar{Z}\cap D(\bar{f})$, and let $\bar{G}$ be the inverse image of $G$ on $\bar{Y}$. By (SGAA VIII 5.2),
\[
        \R i_*(G)_{\bar{x}}=\R j_*(G)_{\bar{x}}\simeq \R\Gamma(\bar{Y},\bar{G}).
\]
Condition $(\mathrm{cdloc})_n$ therefore implies, by a spectral sequence argument, that
\[
        R^p i_*(G)_{\bar{x}}=0\quad\text{for }p>2d+\dim.\mathrm{inj}(F_{\bar{x}})+\dim(\mathcal{O}_{Z,x}).
\]
But (cf. the proof of I 3.2.1)
\[
        \dim(\mathcal{O}_{Z,x})=\dim(\mathcal{O}_{X,x})-d\geq0,
\]
so
\[
        R^p i_*(G)_{\bar{x}}=0\quad\text{for }p>d+\dim(\mathcal{O}_{X,x})+\dim.\mathrm{inj}(F_{\bar{x}}),
\]
and a fortiori for $p>2\dim(\mathcal{O}_{X,x})+\dim.\mathrm{inj}(F_{\bar{x}})$. This proves $(*)$.

It remains to prove that the bound $(*)$ is the best possible. Put $f=\dim.\mathrm{inj}(F_{\bar{x}})$, and choose a finite type $A_{\bar{x}}$-Module $M'$ such that
\[
        \underline{\Ext}^f(M',F_{\bar{x}})\neq0.
\]
Let $Z$ be the integral closed sub-prescheme of $X$ which is the closure of $x$, and let $Y=U\cap Z$ be a regular open subset of $Z$ containing $x$ and of codimension $d=\dim(\mathcal{O}_{X,x})$ at every point. After possibly shrinking $Y$, there exists a locally constant constructible $A_Y$-Module $M$ such that $M_{\bar{x}}=M'$. Applying the purity theorem, one finds
\[
        \underline{\Ext}^{2d+f}(M,F|U)_{\bar{x}}\simeq \underline{\Ext}^f(M',F_{\bar{x}})\neq0,
\]
which completes the proof.

If $F$ is an $A_X$-Module such that $F_{\bar{x}}$ is injective, then in the last part of the proof one may take $M=A_Y$; hence:

\textbf{Corollary 5.1.1.} Under the hypotheses of (5.1), suppose that $F$ is a locally constant $A_X$-Module such that $F_{\bar{x}}$ is injective. Then
\[
        \dim.\mathrm{q.inj}_{\bar{x}}(F)=\dim.\mathrm{top}_{\bar{x}}(F)=2\dim(\mathcal{O}_{X,x}),
\]
with the notation of no.~2.

\textbf{Remarks 5.1.2.} a) With the notation of (5.1), if $F_{\bar{x}}=0$, then $F=0$ in a neighbourhood of $x$, hence $\R\underline{\Hom}(E,F)_{\bar{x}}=0$ for all $E\in\ob\D^{-}(X)$.

b) If $\dim.\mathrm{inj}(F_{\bar{x}})=+\infty$, then $\dim.\mathrm{q.inj}_{\bar{x}}(F)=+\infty$; indeed, for every $p$ there exists a finite type $A_{\bar{x}}$-module $M'$ such that $\underline{\Ext}^p(M',F_{\bar{x}})\neq0$, hence a constructible $A_X$-Module $M$, locally constant in a neighbourhood of $x$, such that $\underline{\Ext}^p(M,F)_{\bar{x}}\neq0$.

\subsection*{6. Constructibility of $\R\underline{\Hom}$}

The equivalence of conditions a) of (I 3.3.1) remains valid when the Ring $(\Z/n\Z)_X$ is replaced by a commutative Ring, locally constant, whose geometric fibres are noetherian rings: nothing has to be changed in the proof.

The preceding conditions are verified when one has resolution of singularities and purity in the strong sense (4.4), in particular if $X$ has dimension $\leq1$, or if $X$ is excellent of characteristic zero, or locally of finite type over a field and of dimension $\leq2$.

\subsection*{7. Existence of dualizing complexes}

In this number, $X$ denotes a prescheme equipped with a commutative sheaf of rings $A_X$, locally constant, whose geometric fibres are noetherian rings.

\textbf{Lemma 7.1.} If $F\in\ob\D^{-}_{\mathrm{coh}}(X)$ and $G\in\ob\D^{+}_{\mathrm{coh}}(X)$, then $\R\underline{\Hom}(F,G)\in\ob\D^{+}_{\mathrm{coh}}(X)$ (notation of (1.2)).

\emph{Proof.} This is standard. One reduces, by a spectral sequence, to $F$ and $G$ concentrated in degree zero, and then argues as in (EGA $0_{\mathrm{III}}$ 12.3.3).

\textbf{Lemma 7.2.} Let $g:X'\to X$ be a morphism of preschemes, and endow $X'$ with the sheaf of rings $A_{X'}=g^*(A_X)$. Then, for $F\in\ob\D^{-}_{\mathrm{coh}}(X)$ and $G\in\ob\D^{+}(X)$, the canonical map
\[
        g^*\R\underline{\Hom}(F,G)\longrightarrow \R\underline{\Hom}(g^*F,g^*G)
\]
is an isomorphism.

\emph{Proof.} This too is standard: one reduces by dévissage to $F=A_X$ (cf. EGA $0_{\mathrm{III}}$ 12.3.5).

\textbf{Remark 7.2.1.} The preceding lemmas are special cases of more general statements appearing in (SGA 6 I).

\textbf{Lemma 7.3.} Suppose $X$ is quasi-compact, and let $K\in\ob\D^{+}_{\mathrm{coh}}(X)$ be such that, for every geometric point $\bar{x}$ of $X$, $K_{\bar{x}}$ is a dualizing complex in the sense of [H] V 2.1. Then, for every $F\in\ob\D^{b}_{\mathrm{coh}}(X)$, one has $\R\underline{\Hom}(F,K)\in\ob\D^{b}_{\mathrm{coh}}(X)$, and $F$ is reflexive relative to $K$.

\emph{Proof.} It suffices, of course, to check this for $F$ concentrated in degree zero. By (7.1), $\R\underline{\Hom}(F,K)\in\ob\D^{+}_{\mathrm{coh}}(X)$. To prove boundedness, let $\bar{x}$ be a geometric point of $X$. By (7.2),
\[
        \R\underline{\Hom}(F,K)_{\bar{x}}\simeq \R\underline{\Hom}(F_{\bar{x}},K_{\bar{x}}).
\]
Putting $\dim.\mathrm{inj}(K_{\bar{x}})=N(\bar{x})$, one obtains
\[
        \underline{\Ext}^i(F,K)_{\bar{x}}=0\quad\text{for }i>N(\bar{x}).
\]
Since the sheaves $\underline{\Ext}^{i}(F,K)$ are locally constant, there is an open neighbourhood $U$ of $x$ such that $\underline{\Ext}^{i}(F,K)|U=0$ for $i>N(\bar{x})$, and quasi-compactness of $X$ gives the assertion. It remains to check that the canonical arrow
\[
        F\longrightarrow \R\underline{\Hom}(\R\underline{\Hom}(F,K),K)
\]
is an isomorphism. It is enough to evaluate at a geometric point $\bar{x}$ of $X$. In the commutative diagram
\[
\begin{tikzcd}
F_{\bar{x}} \arrow[r] \arrow[d,"\mathrm{Id}"'] & \R\underline{\Hom}(\R\underline{\Hom}(F,K),K)_{\bar{x}} \arrow[d]\\
F_{\bar{x}} \arrow[r] & \R\underline{\Hom}(\R\underline{\Hom}(F_{\bar{x}},K_{\bar{x}}),K_{\bar{x}}),
\end{tikzcd}
\]
the right vertical arrow is an isomorphism by (7.2), and the bottom horizontal arrow is an isomorphism because $K_{\bar{x}}$ is dualizing; hence the top horizontal arrow is an isomorphism, q.e.d.

\textbf{Lemma 7.4.} Suppose $X$ is a noetherian regular prescheme, $A_X$ is killed by an integer $n>0$ prime to the residual characteristics of $X$, and the purity theorem in the strong sense (4.4) is true on $X$.

Let $D$ be a divisor with normal crossings on $X$ (I 3.1.5 a)), and denote by $j:Y=V(D)\to X$ and $i:X-Y=U\to X$ the canonical inclusions. Let $K\in\ob\D^{+}_{\mathrm{coh}}(X)$ be such that, for every geometric point $\bar{x}$ of $X$, $K_{\bar{x}}$ is a dualizing complex for the ring $A_{\bar{x}}$ (in the sense of [H] V 2.1). Then, for every $E\in\ob\D^{b}_{\mathrm{coh}}(X)$, one has $\R\underline{\Hom}(i_!i^*(E),K)\in\ob\D^{b}_{c}(X)$, and the pair $(i_!i^*(E),K)$ is bidualizing.

\emph{Proof.} We may suppose $E$ is concentrated in degree zero, i.e. a locally constant constructible $A_X$-Module. Since the question is local on $X$ for the étale topology, we may suppose that (a) $A_X$ is constant with value $A$, and $E$ is constant with value $E$; and (b) $D$ has strict normal crossings, say $D=\sum_{i\in I}\mathrm{div}(s_i)$ with $s_i\in\Gamma(X,\mathcal{O}_X)$.

Argue by induction on $p=\operatorname{card}(I)$. For $p=0$, the assertions follow from (7.3). Suppose $I=\{1,\ldots,p\}$, put $Y_i=V(s_i)$, $Y'=V(\sum_{1\leq i\leq p-1}\mathrm{div}(s_i))$, and $U'=X-Y'$. We get a cartesian diagram of immersions
\[
\begin{tikzcd}
\varnothing \arrow[r] \arrow[d] & U'\cap Y_p \arrow[r,"i'_p"] \arrow[d,"j'_p"'] & Y_p \arrow[d,"j_p"]\\
U \arrow[r] & U' \arrow[r,"i'"'] & X.
\end{tikzcd}
\]
Consider the exact sequence
\[
        0\longrightarrow E_{U,X}\longrightarrow E_{U',X}\longrightarrow E_{U'\cap Y_p,X}\longrightarrow0.
\]
By the induction hypothesis the theorem is true for $(E_{U',X},K)$. For fixed $K$, the $F$ such that $\R\underline{\Hom}(F,K)\in\ob\D_c^b(X)$ and $(F,K)$ is bidualizing form a triangulated subcategory of $\D_c(X)$; hence we are reduced to the case $F=E_{U'\cap Y_p,X}$. But
\[
        E_{U'\cap Y_p,X}=j_{p*}i'_{p!}(E_{U'\cap Y_p}),
\]
and, by (I 1.13), it is equivalent to prove the theorem for $(i'_{p!}(E_{U'\cap Y_p}),\R j_p^!(K))$ on $Y_p$. By purity, however,
\[
        \R j_p^!(K)\simeq j_p^*(K)\otimes^L_{\Z/n\Z}(\mu_n)_{Y_p}^{\otimes -1}[-2].
\]
Thus $\R j_p^!(K)\in\ob\D^{+}_{\mathrm{coh}}(Y_p)$ and, for every geometric point $\bar{y}$ of $Y_p$, $(\R j_p^!(K))_{\bar{y}}$ is dualizing. The induction hypothesis then applies.

\textbf{Theorem 7.5.} Assume the following: $X$ is a noetherian regular prescheme; $A_X$ is killed by an integer $n>0$ prime to the residual characteristics of $X$; $X$ is strongly desingularizable (I 3.1.5), satisfies conditions (reg) (I 3.1.1) and $(\mathrm{cdloc})_n$ (I 3.1.3); and the purity theorem in the strong sense (4.4) is true on every prescheme smooth over $X$ (conditions satisfied if $X$ is excellent of characteristic zero, or of dimension $\leq2$ and smooth over a perfect field).

Then, if $K\in\ob\D^{+}_{\mathrm{coh}}(X)$ is such that $K_{\bar{x}}$ is dualizing (in the sense of [H] V 2.1) for every geometric point $\bar{x}$ of $X$, the complex $K$ is dualizing.

\emph{Proof.} We must show that $K$ satisfies conditions (i) to (iii) of (I 1.7). For every geometric point $\bar{x}$ of $X$, $K_{\bar{x}}$ has finite injective dimension, and since $K\in\ob\D^{+}_{\mathrm{coh}}(X)$, the function $x\mapsto\dim.\mathrm{inj}(K_{\bar{x}})$ is locally constant; it is therefore bounded because $X$ is noetherian. It follows from (5.1) that $K$ has finite quasi-injective dimension.

By what was seen in no.~6, condition (ii) of loc. cit. is also satisfied.

It remains to prove condition (iii), i.e. that every $F\in\ob\D_c^b(X)$ is reflexive relative to $K$. By dévissage we may suppose $F$ is concentrated in degree zero. Since the question is local on $X$ for the étale topology, we may suppose $A_X$ constant with value $A$. Dévissage (1.3) then reduces us to $F=f_*i_!(E_U)$, where $f:Y\to X$ is finite, $Y$ is integral with generic point separable over its image, $i:U\to Y$ is a nonempty regular open, and $E_U$ is an $A_U$-Module constant with value a finite type $A$-module $E$. Since $X$ is strongly desingularizable, there exists (I 3.1.5) a projective birational morphism $h:Y'\to Y$ such that $Y'$ is regular, $h$ induces an isomorphism $h^{-1}(U)\to U$, and $h^{-1}(U)$ is the complement of a normal crossings divisor in $Y'$. Identifying $h^{-1}(U)$ with $U$ via $h$, we have a commutative diagram
\[
\begin{tikzcd}
& Y' \arrow[d,"h"] \arrow[dr,"g"] & \\
U \arrow[ur,"i'"] \arrow[r,"i"'] & Y \arrow[r,"f"'] & X.
\end{tikzcd}
\]
Consequently,
\[
        f_*i_!(E_U)\simeq f_*\R h_* i'_!(E_U)\simeq \R g_* i'_!(E_U)\simeq \R g_* i'_! i'^*(E_{Y'}).
\]
By (I 1.13), it is enough to prove biduality for the pair $(i'_!i'^*(E_{Y'}),\R g^!(K))$. But, locally embedding $Y'$ into a prescheme $X'$ smooth over $X$ and applying strong purity on $X'$ (4.4), one finds that $\R g^!(K)$ is locally isomorphic, up to degree shift, to $g^*(K)$. Hence $\R g^!(K)\in\ob\D^{+}_{\mathrm{coh}}(Y')$ and $\R g^!(K)_{\bar{x}}$ is dualizing for every geometric point $\bar{x}$ of $Y'$. It remains only to apply (7.4).

The preceding proof also shows that $\R\underline{\Hom}(F,K)\in\ob\D_c^b(X)$ (using the finiteness theorem for a proper morphism (SGAA XIV)), which will reassure the reader who did not trust the spiel of no.~6.

This completes the proof.

\begin{center}
$-:\!-:\!-$
\end{center}

\section*{Bibliography}

\begin{enumerate}[label={[\arabic*]}]
\item Verdier, J.-L., \emph{Dimension des espaces localement compacts}, Note aux C.R. Ac. Sc. Paris, 261, pp.~5293--5296 (1965). See also: Verdier, J.-L., \emph{Dualité dans la cohomologie des espaces localement compacts}, Séminaire Bourbaki, no.~300, 1965--66.
\end{enumerate}

SGA A = SGA 4.

\begin{center}
$-:\!-:\!-$
\end{center}

\clearpage
\phantomsection\addcontentsline{toc}{section}{Exposé III --- Lefschetz Formula}
\begin{center}
{\Large \textbf{SGA 5 --- Exposé III}}\\[0.4em]
{\large \textbf{LEFSCHETZ FORMULA}}\\[0.3em]
by A.~Grothendieck, written up by L.~Illusie
\end{center}

In this exposé we prove the Lefschetz formula in étale cohomology, for correspondences between complexes of sheaves, in Verdier's formulation (cf.~[18]). In an earlier version, distributed by the IHES about ten years ago, the formula was stated, under resolution hypotheses, but the required compatibilities were not proved. Deligne's finiteness theorems (SGA~$4^{1/2}$, Finiteness) made it possible to dispense with these hypotheses, provided one works with schemes separated and of finite type over a field. As for the compatibilities, the editor has verified them (painfully).

Let us say at once that the Lefschetz formula in question here, just like the trace formula for constant coefficients (SGA~$4^{1/2}$, Cycle), which it generalizes, is essentially trivial: it only translates certain compatibilities in the formalism of duality. Roughly speaking, it says that, under certain conditions, the trace of an endomorphism of a hypercohomology complex $R\Gamma(X,L)$, where $L$ is a complex of sheaves on $X$, can be calculated as an ``integral'' of certain cohomology classes on $X\times X$. In practice, the endomorphisms one considers are associated with correspondences $C\subset X\times X$, and the classes recovered on the product have support in $C\cap\Delta$, where $\Delta$ is the diagonal; these classes are ``local in nature'' in a neighbourhood of the fixed points, which in principle facilitates their calculation. See 4.2 and 4.7 for precise statements. In fact, the calculation of these ``local'' classes in terms of more comprehensible invariants, without which the present formula would have no interest, has so far been undertaken only in very special situations (see 4.10, 4.11, and the following exposé, where in particular the case of the Frobenius correspondence on a curve is treated).

For lack of a good derived category of $\ell$-adic sheaves (Jouanolou's thesis unfortunately not having been published), we work systematically with torsion coefficients (prime to the residual characteristics). Despite minor technical difficulties in the definition of traces, this point of view in fact leads to finer formulae, also more manageable insofar as, for the calculation of local terms, one wants to trivialize the sheaves involved by coverings (see notably (SGA~$4^{1/2}$, Report), III~B, and XII).

In no.~1, we fix the notation and recall various Künneth formulae from (SGA~$4^{1/2}$, Finiteness), which play an essential role in what follows. We also introduce a certain number of fibred and cofibred categories, whose consideration later facilitates certain verifications of compatibilities. No.~2, which is very technical, may be skipped on first reading. Its main result (2.5) expresses a compatibility between external tensor product of $\RHom$, direct image, and base change. In no.~3, we define cohomological correspondences between complexes, and show how they act on hypercohomology; the result quoted above makes it possible to describe this action in several equivalent ways. The core of the exposé is no.~4, where Verdier's pairing between cohomological correspondences is defined and the Lefschetz formula is proved, in a relative situation appreciably more general than the one usually considered, with a view to later applications to computations of local terms. We give the most evident corollaries for correspondences with constant coefficients. In no.~5, we indicate rapidly how to generalize to correspondences between complexes certain usual constructions on correspondences with constant coefficients (cf.~[12]), such as passage to the dual, composition, and external tensor product. Most of the verifications, analogous to those in the preceding numbers, have been omitted. Finally, in the appendix, we paraphrase the preceding theory in the context of coherent sheaves, using the formalism of duality of [8]. Here one may work over any noetherian base $S$, provided that, when necessary, one imposes hypotheses of properness, finite tor-dimension, and tor-independence on the $S$-schemes considered. We show that the very general Lefschetz formula thereby obtained contains as a special case the Woods Hole formula [1], and we point out some applications.

\subsection*{Contents}
\begin{enumerate}[label=\arabic*.]
\item Notation and recalls of Künneth formulae.
\item Functoriality of $\RHom$ and external tensor products.
\item Cohomological correspondences.
\item Pairings of correspondences. Lefschetz formula.
\item Complements.
\item Appendix. Lefschetz formula for coherent sheaves.
\end{enumerate}

\section{Notation and recalls of Künneth formulae}

\subsection*{1.1}
We denote by $S$ the spectrum of a field. Unless otherwise stated, the $S$-schemes considered in what follows will be $S$-schemes separated and of finite type.

We denote by $A$ a commutative noetherian ring, annihilated by an integer invertible on $S$. If $X$ is an $S$-scheme, we denote by $D(X)=D(X,A)$ the derived category of complexes of $A$-modules on $X$ (for the étale topology), and by $D_{\ctf}(X)$ the full subcategory consisting of complexes of finite tor-dimension, with constructible cohomology. The category $D_{\ctf}(X)$ is a triangulated subcategory of $D(X)$. If $E,F\in\ob D_{\ctf}(X)$, it follows readily from the definitions that $E\otimes^L F\in\ob D_{\ctf}(X)$. According to (SGA~$4^{1/2}$, Finiteness Theorem 1.6, 1.7), one also has $\RHom(E,F)\in\ob D_{\ctf}(X)$.

\subsection*{1.2}
Let $f:X\to Y$ be an $S$-morphism. One has trivially
\[
f^{*}D_{\ctf}(Y) \subset D_{\ctf}(X),
\]
and, by (loc. cit.),
\[
Rf_{*}D_{\ctf}(X) \subset D_{\ctf}(Y),\qquad
Rf_{!}D_{\ctf}(X) \subset D_{\ctf}(Y),\qquad
Rf^{!}D_{\ctf}(Y) \subset D_{\ctf}(X),
\]
(the inclusion relative to $Rf_{!}$ uses only the finiteness theorem for proper morphisms (SGA~4~XIV) and the sorites (SGA~4~XVII~5.2.11), and hence holds more generally for every compactifiable morphism whose fibres have bounded dimension). In what follows, we shall sometimes write $f_{*}$, $f_{!}$, $f^{!}$ instead of $Rf_{*}$, $Rf_{!}$, $Rf^{!}$; the hypotheses of 1.1 imply that these functors are defined on the whole derived category.

\subsection*{1.3}
For $f:X\to S$, we shall put
\[
K_X=f^{!}A,
\]
and shall denote by $D_X$ (or $D$) the functor $\RHom(-,K_X)$. According to (SGA~$4^{1/2}$, Finiteness Theorem 4.3), the complex $K_X$ is dualizing: for $E\in\ob D_{\ctf}(X)$, one has canonically
\[
E \xrightarrow{\sim} DDE.
\]

\subsection*{1.4}
Let $f:X\to Y$ be an $S$-morphism, $L\in\ob D(X)$, $M\in\ob D(Y)$. We put
\[
\Hom_f^{(1)}(L,M)=\Hom(f^{*}M,L)\quad (\simeq \Hom(M,f_{*}L)),
\]
\[
\Hom_f^{(2)}(L,M)=\Hom(L,f^{*}M),
\]
\[
\Hom_f^{(3)}(L,M)=\Hom(f_{!}L,M)\quad (\overset{\text{(SGA~4~XVIII~3.1.5)}}{\simeq} \Hom(L,f^{!}M)),
\]
\[
\Hom_f^{(4)}(L,M)=\Hom(f^{!}M,L).
\]
The elements of $\Hom_f^{(i)}(L,M)$ will be called \emph{arrows of type $i$ from $L$ to $M$ over $f$} (or $f$-arrows of type $i$ from $L$ to $M$). The transitivity isomorphisms for the functors $f^{*}$, $f_{*}$, $f_{!}$, $f^{!}$ allow one to compose arrows of type $i$ for fixed $i$, and one thus obtains certain fibred or cofibred categories over the category of $S$-schemes: for example, for $i=1$ (resp. $i=3$) one obtains the fibred and cofibred category considered in (SGA~4~XVII~4.1.3) (resp. (SGA~4~XVIII~3.1.13)), whose fibre over $X$ is the category opposite to $D(X)$ (resp. $D(X)$).

\subsection*{1.5}
Let, for $i=1,2$, $X_i$ be an $S$-scheme, $X=X_1\times_SX_2$, and $L_i\in\ob D^{-}(X_i)$. We put
\[
L_1\otimes^L_SL_2=\pr_1^{*}L_1\otimes^L\pr_2^{*}L_2,
\]
where $\pr_i:X\to X_i$ is the canonical projection.

\subsection*{1.6}
Let, for $i=1,2$, $f_i:X_i\to Y_i$ be an $S$-morphism,
\[
f=f_1\times_Sf_2:X=X_1\times_SX_2\to Y=Y_1\times_SY_2,
\]
$L_i\in\ob D^{-}(X_i)$, $M_i\in\ob D^{-}(Y_i)$. The evident canonical isomorphism
\begin{equation}\tag{1.6.1}
f_1^{*}M_1\otimes^L_Sf_2^{*}M_2\xrightarrow{\sim}f^{*}(M_1\otimes^L_SM_2)
\end{equation}
allows one to define, for $u_i\in\Hom_{f_i}^{(1)}(L_i,M_i)$ (resp. $\Hom_{f_i}^{(2)}(L_i,M_i)$),
\begin{equation}\tag{1.6.2}
u_1\otimes^L_Su_2\in\Hom_f^{(1)}(L,M)\quad(\text{resp. }\Hom_f^{(2)}(L,M)),
\end{equation}
where $L=L_1\otimes^L_SL_2$, $M=M_1\otimes^L_SM_2$. If $u_i$, $v_i$ are composable, one has
\begin{equation}\tag{1.6.3}
(v_1u_1)\otimes^L_S(v_2u_2)=(v_1\otimes^L_Sv_2)(u_1\otimes^L_Su_2).
\end{equation}
Taking $M_i=f_{i*}L_i$, and $u_i$ of type~1, given by the identity of $f_{i*}L_i$, 1.6.2 gives a \emph{Künneth arrow} (cf.~SGA~4~XVII~5.4.1.4)
\begin{equation}\tag{1.6.4}
f_{1*}L_1\otimes^L_Sf_{2*}L_2\longrightarrow f_{*}(L_1\otimes^L_SL_2).
\end{equation}
According to (SGA~$4^{1/2}$, Finiteness Theorem 1.9), 1.6.4 is an isomorphism: by 1.6.3 one is indeed reduced to the case where $X_2=Y_2$, $f_2=1$, and one applies (loc. cit. App.). When $f_1$ and $f_2$ are proper, it is not necessary to invoke (loc. cit.): 1.6.4 is a special case of the Künneth isomorphism (SGA~4~XVII~5.4.3).

\subsection*{1.7}
Choose a compactification $f_i=p_i j_i$, where $j_i$ is an open immersion and $p_i$ is a proper morphism. Composing the evident canonical isomorphism (where $j=j_1\times_Sj_2$)
\[
j_{1!}L_1\otimes^L_Sj_{2!}L_2\xrightarrow{\sim}j_{!}(L_1\otimes^L_SL_2)
\]
with the isomorphism 1.6.4 relative to $(p_i,j_{i!}L_i)$, one obtains a Künneth isomorphism (SGA~4~XVII~5.4.3)
\begin{equation}\tag{1.7.1}
f_{1!}L_1\otimes^L_Sf_{2!}L_2\xrightarrow{\sim}f_{!}(L_1\otimes^L_SL_2).
\end{equation}
One checks easily (loc. cit.) that 1.7.1 does not depend on the compactifications chosen, and is compatible with the transitivity isomorphisms for composites $g_if_i$ (cf.~(loc. cit. 5.2.4)). Thanks to 1.7.1, one can define, for $u_i\in\Hom_{f_i}^{(3)}(L_i,M_i)$,
\begin{equation}\tag{1.7.2}
u_1\otimes^L_Su_2\in\Hom_f^{(3)}(L,M)
\end{equation}
with $L$ and $M$ as in 1.6.2, and one has a transitivity formula analogous to 1.6.3.

For $L_i=f_i^{!}M_i$, and $u_i$ given by the identity of $L_i$, 1.7.2 gives a \emph{Künneth arrow}
\begin{equation}\tag{1.7.3}
f_1^{!}M_1\otimes^L_Sf_2^{!}M_2\longrightarrow f^{!}(M_1\otimes^L_SM_2).
\end{equation}

\subsection*{Proposition 1.7.4}
The arrow 1.7.3 is an isomorphism.

By the transitivity formula, it suffices to treat the case where $X_2=Y_2$, $f_2=1$. By dévissage, one reduces to $M_2=g_{*}j_{!}A_V$, where $j:V\to U$ is an open immersion and $g:U\to Y_2$ a finite morphism, and then, by (SGA~4~XVIII~3.1.12.3), to the case where $M_2=A_{Y_2}$. By localization on $X_2$, one may suppose moreover that $f_1=g_1j_1$, where $j_1$ is a closed immersion and $g_1$ a smooth morphism purely of relative dimension $n_1$. By transitivity, it suffices to treat the cases
\[
\text{a) } f_1=g_1,\qquad \text{b) } f_1=j_1.
\]
In case a), the trace morphism supplies by adjunction (SGA~4~XVIII~3.2.5) an isomorphism
\[
f_1^{*}M_1(n_1)[2n_1]\xrightarrow{\sim}f_1^{!}M_1;
\]
this is compatible with the base change $Y_2\to S$, because the trace morphism is so (SGA~4~XVIII~2.9), whence 1.7.4 in this case. In case b), if $i_1$ denotes the inclusion of the complementary open, one has an exact sequence
\[
0\longrightarrow j_{1*}Rj_1^{!}M_1\longrightarrow M_1\longrightarrow Ri_{1*}i_1^{*}M_1\longrightarrow0,
\]
and the compatibility of $Ri_{1*}$ with base change, a special case of 1.6.4, gives the conclusion.

Thanks to 1.7.3, one can define, for $u_i\in\Hom_f^{(4)}(L_i,M_i)$,
\begin{equation}\tag{1.7.5}
u_1\otimes^L_Su_2\in\Hom_f^{(4)}(L,M)
\end{equation}
($L$ and $M$ as in 1.6.2), with a transitivity formula analogous to 1.6.3.

On the other hand, with the notation of 1.3, 1.7.3 gives an isomorphism
\begin{equation}\tag{1.7.6}
K_{X_1}\otimes^L_SK_{X_2}\xrightarrow{\sim}K_X.
\end{equation}

\subsection*{1.8}
Let $f:X\to Y$ be an $S$-morphism, $L\in\ob D(X)$, $M\in\ob D(Y)$, $S'$ an $S$-scheme, $P\in\ob D(S')$, and $f':X'\to Y'$ the morphism deduced from $f$ by the base change $S'\to S$. The identity arrow of $P$ may be regarded indifferently as an arrow of type $i$ ($1\leq i\leq4$) over the identity arrow of $S'$; for $u\in\Hom_f^{(i)}(L,M)$ ($1\leq i\leq4$), one may therefore, by 1.6.2, 1.7.2, 1.7.5, consider
\begin{equation}\tag{1.8.1}
u\otimes^L_S\Id_P\in\Hom_{f'}^{(i)}(L\otimes^L_SP,M\otimes^L_SP).
\end{equation}
We shall sometimes write $u\otimes^L_SP$ instead of $u\otimes^L_S\Id_P$.

\section{Functoriality of \texorpdfstring{$\RHom$}{RHom} and external tensor products}

\subsection*{2.1}
Let $f:X\to Y$ be an $S$-morphism, $P\in\ob D(Y)$, $Q\in\ob D^{+}(Y)$. The evident canonical arrow at the level of complexes
\[
f^{*}\mathcal Hom^{\bullet}(P,Q)\longrightarrow\mathcal Hom^{\bullet}(f^{*}P,f^{*}Q)
\]
gives by derivation a canonical arrow
\begin{equation}\tag{2.1.1}
f^{*}\RHom(P,Q)\longrightarrow\RHom(f^{*}P,f^{*}Q).
\end{equation}
Let $L\in\ob D(X)$, $M\in\ob D^{+}(X)$. Put
\begin{equation}\tag{2.1.2}
\Hom_f^{(i,j)}((L,M),(P,Q))=\Hom_f^{(i)}(L,P)\times\Hom_f^{(j)}(M,Q),
\end{equation}
with the notation of 1.4. For $(u,v)\in\Hom_f^{(2,1)}((L,M),(P,Q))$ one defines
\begin{equation}\tag{2.1.3}
\RHom(u,v)\in\Hom_f^{(1)}(\RHom(L,M),\RHom(P,Q))
\end{equation}
by composing 2.1.1 with the arrow $\RHom(f^{*}P,f^{*}Q)\to\RHom(L,M)$ defined by $u$ and $v$. For $L=f^{*}P$, $Q=f_{*}M$, $u$ (resp. $v$) given by the identity of $L$ (resp. $Q$), $\RHom(u,v)$ is the ``trivial duality'' isomorphism
\begin{equation}\tag{2.1.4}
\RHom(P,f_{*}M)\xrightarrow{\sim}f_{*}\RHom(f^{*}P,M).
\end{equation}
For $(u,v)\in\Hom_f^{(3,4)}((L,M),(P,Q))$, one defines
\begin{equation}\tag{2.1.5}
\RHom(u,v)\in\Hom_f^{(1)}(\RHom(L,M),\RHom(P,Q))
\end{equation}
as the composite of the arrows
\[
\RHom(P,Q)\xrightarrow{a}\RHom(f_{!}L,Q)\xrightarrow{b}f_{*}\RHom(L,f^{!}Q)\xrightarrow{c}f_{*}\RHom(L,M),
\]
where $a$ (resp. $c$) is defined by $u$ (resp. $v$), and $b$ is the inverse of the duality isomorphism (SGA~4~XVIII~3.1.9.6).

It follows easily from the definitions that the formation of $\RHom(u,v)$ is compatible with composition: for composable $(u,v)$, $(u',v')$ of type $(2,1)$ (resp. $(3,4)$), one has
\begin{equation}\tag{2.1.6}
\RHom(u'u,v'v)=\RHom(u',v')\,\RHom(u,v).
\end{equation}

\subsection*{2.2}
Let $X$ be an $S$-scheme. Recall that for $L\in\ob D^{-}(X)$, $M\in\ob D(X)$, $N\in\ob D^{+}(X)$, one has the adjunction isomorphism (``dear to Cartan,'' as Grothendieck says)
\begin{equation}\tag{2.2.1}
\Hom(L\otimes^{L}M,N)=\Hom(L,\RHom(M,N)).
\end{equation}
For $E\in\ob D^{-}(X)$, $F\in\ob D^{+}(X)$, the identity arrow of $\RHom(E,F)$ defines via 2.2.1 an arrow
\begin{equation}\tag{2.2.2}
E\otimes^{L}\RHom(E,F)\longrightarrow F,
\end{equation}
called the \emph{evaluation arrow}: it is derived from the arrow at the level of complexes
\[
\mathcal Hom^{\bullet}(E,F)\otimes E\longrightarrow F,
\qquad f\otimes x\longmapsto f(x).
\]
Let, for $i=1,2$, $E_i,F_i\in\ob D_{\ctf}(X)$. By adjunction (2.2.1), the arrow
\[
E_1\otimes^{L}E_2\otimes^{L}\RHom(E_1,F_1)\otimes^{L}\RHom(E_2,F_2)\longrightarrow F_1\otimes^{L}F_2,
\]
composed of a symmetry isomorphism and of the tensor product of the evaluation arrows relative to $(E_1,F_1)$ and $(E_2,F_2)$, supplies an arrow
\begin{equation}\tag{2.2.3}
\RHom(E_1,F_1)\otimes^{L}\RHom(E_2,F_2)\longrightarrow\RHom(E_1\otimes^{L}E_2,F_1\otimes^{L}F_2).
\end{equation}
Let, for $i=1,2$, $X_i$ be an $S$-scheme, $X=X_1\times_SX_2$, $p_i:X\to X_i$ the canonical projection, $E_i,F_i\in\ob D_{\ctf}(X_i)$, $E=E_1\otimes^{L}_{S}E_2$, $F=F_1\otimes^{L}_{S}F_2$. One has a canonical arrow
\begin{equation}\tag{2.2.4}
\RHom(E_1,F_1)\otimes^{L}_{S}\RHom(E_2,F_2)\longrightarrow\RHom(E,F),
\end{equation}
obtained by composing the tensor product of the arrows $p_i^{*}\RHom(E_i,F_i)\to\RHom(p_i^{*}E_i,p_i^{*}F_i)$ and 2.2.3 for $(p_i^{*}E_i,p_i^{*}F_i)$.

\subsection*{Proposition 2.3}
The arrow 2.2.4 is an isomorphism.

By localization and dévissage, one reduces to $E_i=f_{i!}A_{U_i}$ for $f_i:U_i\to X_i$ étale. One has canonical isomorphisms (duality)
\begin{equation}\tag{1}
\RHom(f_{i!}A_{U_i},F_i)\xrightarrow{\sim}f_{i*}f_i^{*}F_i,
\qquad
\RHom(E,F)\xrightarrow{\sim}f_{*}f^{*}F\quad(f=f_1\times_S f_2).
\end{equation}
Taking into account the Künneth formula 1.6.4, it suffices, to conclude, to verify that these isomorphisms make the square
\begin{equation}\tag{2}
\begin{tikzcd}[column sep=large,row sep=large]
\RHom(f_{1!}A_{U_1},F_1)\otimes^{L}_{S}\RHom(f_{2!}A_{U_2},F_2) \arrow[r,"2.2.3"] \arrow[d,"\sim"'] & \RHom(f_{!}A_U,F) \arrow[d,"\sim"]\\
f_{1*}f_1^{*}F_1\otimes^{L}_{S}f_{2*}f_2^{*}F_2 \arrow[r,"1.6.4"'] & f_{*}f^{*}F
\end{tikzcd}
\end{equation}
commute, where $U=U_1\times_SU_2$. Since $f_i$ is étale, the isomorphism (1) relative to $f_i$ (resp. $f$) is adjoint to the composite arrow
\begin{equation}\tag{3}
f_i^{*}\RHom(f_{i!}A_{U_i},F_i)\xrightarrow{\sim}\RHom(f_i^{*}f_{i!}A_{U_i},f_i^{*}F_i)\longrightarrow f_i^{*}F_i,
\end{equation}
where the second arrow is deduced from the adjunction $A_{U_i}\to f_i^{*}f_{i!}A_{U_i}$ (resp. the analogous composite arrow with $f$). The commutativity of (2) is equivalent to that of the square
\begin{equation}\tag{4}
\begin{tikzcd}[column sep=large,row sep=large]
f_1^{*}\RHom(f_{1!}A_{U_1},F_1)\otimes^{L}_{S}f_2^{*}\RHom(f_{2!}A_{U_2},F_2) \arrow[r] \arrow[d] & f^{*}\RHom(f_{!}A_U,F) \arrow[d]\\
f_1^{*}F_1\otimes^{L}_{S}f_2^{*}F_2 \arrow[r,"\sim"'] & f^{*}F,
\end{tikzcd}
\end{equation}
where the vertical arrows are given by (3). But the commutativity of (4) follows from the fact that 2.2.3 is functorial in $E_i$ and reduces to the identity for $E_i=A$, whence the proposition.

\subsection*{2.4}
Let, for $i=1,2$, $f_i:X_i\to Y_i$ be an $S$-morphism,
\[
f=f_1\times_S f_2:X\to Y,
\]
$E_i,F_i\in\ob D_{\ctf}(X_i)$, $P_i,Q_i\in\ob D_{\ctf}(Y_i)$,
\[
E=E_1\otimes^{L}_{S}E_2,
\quad F=F_1\otimes^{L}_{S}F_2,
\quad P=P_1\otimes^{L}_{S}P_2,
\quad Q=Q_1\otimes^{L}_{S}Q_2.
\]
Consider the diagram
\begin{equation}\tag{2.4.0}
\resizebox{0.96\textwidth}{!}{%
\begin{tikzpicture}[>=latex,baseline=(current bounding box.center)]
\node (EF1) at (-5.2,2.0) {$(E_1,F_1)$};
\node (X1)  at (-3.55,2.0) {$X_1$};
\node (X1Y2) at (-1.0,2.0) {$X_1\times_S Y_2$};
\node (X) at (1.9,2.0) {$X$};
\node (EF) at (3.15,2.35) {$(E,F)$};
\node (PQ1) at (-5.2,0.45) {$(P_1,Q_1)$};
\node (Y1) at (-3.55,0.45) {$Y_1$};
\node (Y) at (-1.0,0.45) {$Y$};
\node (Y1X2) at (1.9,0.45) {$Y_1\times_S X_2$};
\node (S) at (-3.55,-1.25) {$S$};
\node (Y2) at (-1.0,-1.25) {$Y_2$};
\node (X2) at (1.9,-1.25) {$X_2$};
\node (PQ) at (-1.0,-0.25) {$(P,Q)$};
\node (PQ2) at (-1.0,-1.9) {$(P_2,Q_2)$};
\node (EF2) at (2.95,-1.9) {$(E_2,F_2)$};
\draw[->] (X1Y2) -- (X1);
\draw[->] (X) -- (X1Y2);
\draw[->] (X1) -- node[left] {$f_1$} (Y1);
\draw[->] (X1Y2) -- (Y);
\draw[->] (X) -- (Y1X2);
\draw[->] (X) -- node[below,sloped,pos=.58] {$f$} (Y);
\draw[->] (Y) -- (Y1);
\draw[->] (Y1X2) -- (Y);
\draw[->] (Y1) -- (S);
\draw[->] (Y) -- (Y2);
\draw[->] (Y1X2) -- (X2);
\draw[->] (Y2) -- (S);
\draw[->] (X2) -- node[below] {$f_2$} (Y2);
\end{tikzpicture}}
\end{equation}
Give ourselves $(u_1,v_1)\in\Hom_{f_1}^{(3,4)}((E_1,F_1),(P_1,Q_1))$, $(u_2,v_2)\in\Hom_{f_2}^{(2,1)}((E_2,F_2),(P_2,Q_2))$ (notation of 2.1.2 and 1.4). By 1.8, one deduces from $(u_i,v_i)$ a square
\begin{equation}\tag{2.4.1}
\begin{tikzcd}[column sep=huge,row sep=large]
(E,F) \arrow[d,"{(u_1\otimes^L_S E_2,\,v_1\otimes^L_S F_2)}"] & (E_1\otimes^L_S P_2,F_1\otimes^L_S Q_2) \arrow[l,"{(E_1\otimes^L_S u_2,\,F_1\otimes^L_S v_2)}"'] \arrow[d,"{(u_1\otimes^L_S P_2,\,v_1\otimes^L_S Q_2)}"]\\
(P_1\otimes^L_S E_2,Q_1\otimes^L_S F_2) \arrow[r,"{(P_1\otimes^L_S u_2,\,Q_1\otimes^L_S v_2)}"'] & (P,Q),
\end{tikzcd}
\end{equation}
where the vertical arrows (resp. horizontal arrows) are of type $(3,4)$ (resp. $(2,1)$).

We may now state the main compatibility of this number, which will play an essential role in the verification of the Lefschetz--Verdier formula.

\subsection*{Proposition 2.5}
The diagram 2.5.1 below, where the square on the right is deduced from 2.4.1 by applying $\RHom$ in the sense of 2.1.3 and 2.1.5, is commutative:
\begin{equation}\tag{2.5.1}
\resizebox{0.98\textwidth}{!}{%
\begin{tikzpicture}[>=latex,baseline=(current bounding box.center)]
\footnotesize
\node (A) at (0,2.6) {$\RHom(E_1,F_1)\otimes^L_S\RHom(E_2,F_2)$};
\node (B) at (0,1.3) {$\RHom(u_1,v_1)\otimes^L_S\RHom(u_2,v_2)$};
\node (C) at (0,0) {$\RHom(P_1,Q_1)\otimes^L_S\RHom(P_2,Q_2)$};
\node (D) at (8.8,2.6) {$\RHom(E,F)$};
\node (E) at (6.2,1.3) {$\RHom(E_1\otimes^L_SP_2,F_1\otimes^L_SQ_2)$};
\node (F) at (12.1,1.3) {$\RHom(P_1\otimes^L_SE_2,Q_1\otimes^L_SF_2)$};
\node (G) at (8.8,0) {$\RHom(P,Q)$};
\draw[->] (A) -- node[above] {$2.2.4$} (D);
\draw[->] (A) -- (B);
\draw[->] (B) -- (C);
\draw[->] (C) -- node[below] {$2.2.4$} (G);
\draw[->] (D) -- (E);
\draw[->] (D) -- (F);
\draw[->] (E) -- (G);
\draw[->] (F) -- (G);
\end{tikzpicture}}
\end{equation}

\subsection*{2.6}
The proof of 2.5 will occupy the rest of no.~2. It follows from 1.6.3 that, if $w_i:L_i\to M_i$ is an arrow of type~1 over $f_i$, the arrow $w=w_1\otimes^{L}_{S}w_2$ makes the diagram
\[
\begin{tikzcd}[column sep=large,row sep=large]
& L_1\otimes^L_S L_2 \arrow[dl,"{L_1\otimes^L_S w_2}"'] \arrow[dr,"{w_1\otimes^L_S L_2}"] & \\
L_1\otimes^L_S M_2 \arrow[dr,"{w_1\otimes^L_S M_2}"'] & & M_1\otimes^L_S L_2 \arrow[dl,"{M_1\otimes^L_S w_2}"]\\
& M_1\otimes^L_S M_2 &
\end{tikzcd}
\]
commute. Applying this remark to $w_i=\RHom(u_i,v_i)$, one sees that it suffices to prove the following special cases of 2.5.

\subsection*{Corollary 2.6.1}
Let $f:X\to Y$ be an $S$-morphism, $S'$ an $S$-scheme, $f':X'\to Y'$ the morphism deduced from $f$ by the base change $S'\to S$, $E,F\in\ob D_{\ctf}(X)$, $P,Q\in\ob D_{\ctf}(Y)$, $P',Q'\in\ob D_{\ctf}(S')$, and $(u,v)\in\Hom_f^{(2,1)}((E,F),(P,Q))$ (resp. $\Hom_f^{(3,4)}((E,F),(P,Q))$). Then one has a commutative square, whose vertical arrows are of type~1 over $f'$:
\begin{equation}\tag{2.6.1.1}
\begin{tikzcd}[column sep=large,row sep=large]
\RHom(E,F)\otimes^L_S\RHom(P',Q') \arrow[r,"2.2.4"] \arrow[d,"{\RHom(u,v)\otimes^L_S\RHom(P',Q')}"'] & \RHom(E\otimes^L_SP',F\otimes^L_SQ') \arrow[d,"{\RHom(u\otimes^L_SP',v\otimes^L_SQ')}"]\\
\RHom(P,Q)\otimes^L_S\RHom(P',Q') \arrow[r,"2.2.4"'] & \RHom(P\otimes^L_SP',Q\otimes^L_SQ').
\end{tikzcd}
\end{equation}
First examine the case where $(u,v)$ is of type $(2,1)$. By functoriality, one may suppose that $E=f^{*}P$, $F=f^{*}Q$, $(u,v)$ being given by the identity of $(f^{*}P,f^{*}Q)$. The commutativity of 2.6.1.1 in the case $P'=Q'=A_{S'}$ is evident; to establish it in the general case, it therefore suffices to verify that, for $L,M\in\ob D_{\ctf}(Y)$, the square
\[
\begin{tikzcd}[column sep=large,row sep=large]
f^{*}\RHom(E,F)\otimes^L f^{*}\RHom(L,M) \arrow[r,"{f^*(2.2.3)}"] \arrow[d,"{2.1.1\otimes 2.1.1}"'] & f^{*}\RHom(E\otimes^LL,F\otimes^LM) \arrow[d,"2.1.1"]\\
\RHom(f^{*}E,f^{*}F)\otimes^L\RHom(f^{*}L,f^{*}M) \arrow[r,"2.2.3"'] & \RHom(f^{*}(E\otimes^LL),f^{*}(F\otimes^LM))
\end{tikzcd}
\]
is commutative. But one sees immediately that the two composites of this square are adjoint to the tensor product of the evaluation arrows
\[
f^{*}E\otimes^L f^{*}\RHom(E,F)\longrightarrow f^{*}F
\quad\text{and}\quad
f^{*}L\otimes^L f^{*}\RHom(L,M)\longrightarrow f^{*}M,
\]
which proves the assertion.

Now pass to the case where $(u,v)$ is of type $(3,4)$. By functoriality, one may suppose that $P=f_{!}E$, $F=f^{!}Q$, $(u,v)$ being defined by $(\Id_P,\Id_F)$. One must prove the commutativity of the square
\[
\begin{tikzcd}[column sep=large,row sep=large]
f_*\RHom(E,f^!Q)\otimes^L_S\RHom(P',Q') \arrow[r,"b"] \arrow[d,"{a\otimes^L_S\RHom(P',Q')}"'] & f'_*\RHom(E\otimes^L_SP',f^!Q\otimes^L_SQ') \arrow[d,"a'"]\\
\RHom(f_!E,Q)\otimes^L_S\RHom(P',Q') \arrow[r,"2.2.4"'] & \RHom(f_!E\otimes^L_SP',Q\otimes^L_SQ'),
\end{tikzcd}
\]
where $b$ is the composite of $f'_*(2.2.4)$ and a base-change isomorphism, $a$ is the duality isomorphism (SGA~4~XVIII~3.1.9.6), and $a'$ is the composite of the duality isomorphism (relative to $E\otimes^L_SP'$ and $Q\otimes^L_SQ'$) and the base-change isomorphisms
\[
f_!E\otimes^L_SP'\xrightarrow[\sim]{c}f'_!(E\otimes^L_SP'),
\qquad
f^!Q\otimes^L_SQ'\xrightarrow[\sim]{d}f'^!(Q\otimes^L_SQ').
\]
Recall that $a$ is the composite of the canonical arrow (SGA~4~XVIII~3.1.9.5)
\[
f_*\RHom(E,f^!Q)\longrightarrow\RHom(f_!E,f_!f^!Q)
\]
and of the arrow defined by the adjunction arrow $f_!f^!Q\to Q$. There is an analogous description of $a'$, as the composite of the quoted arrow
\begin{equation}\tag{*}
f'_*\RHom(E\otimes^L_SP',f^!Q\otimes^L_SQ')\longrightarrow
\RHom(f'_!(E\otimes^L_SP'),f'_!(f^!Q\otimes^L_SQ')),
\end{equation}
of the arrows defined by $c$ and $d$, and of the arrow defined by the adjunction arrow
\[
f'_!f'^!(Q\otimes^L_SQ')\longrightarrow Q\otimes^L_SQ'.
\]
By definition, $d$ is adjoint to the arrow
\[
f'_!(f^!Q\otimes^L_SQ')\longrightarrow Q\otimes^L_SQ'
\]
defined by the base-change isomorphism $f'_!(f^!Q\otimes^L_SQ')\simeq f_!f^!Q\otimes^L_SQ'$ and the adjunction arrow $f_!f^!Q\to Q$; in other words, the square
\[
\begin{tikzcd}[column sep=large,row sep=large]
f'_!(f^!Q\otimes^L_SQ') \arrow[r,"{f'_!d}"] \arrow[d,"{\text{base ch.}\ \sim}"'] & f'_!f'^!(Q\otimes^L_SQ') \arrow[d,"\mathrm{adj}"]\\
f_!f^!Q\otimes^L_SQ' \arrow[r,"{\mathrm{adj}\otimes 1}"'] & Q\otimes^L_SQ'
\end{tikzcd}
\]
is commutative. It follows that one may also describe $a'$ as the composite of (*), of the arrow defined by $c$, and of those defined by the base-change isomorphism
\[
f'_!(f^!Q\otimes^L_SQ')\xrightarrow{\sim}f_!f^!Q\otimes^L_SQ'
\]
and the adjunction arrow $f_!f^!Q\to Q$. This decomposition of $a'$ and that of $a$ given above, together with the commutativity of the square
\[
\begin{tikzcd}[column sep=large,row sep=large]
\RHom(f_!E,f_!f^!Q)\otimes^L_S\RHom(P',Q') \arrow[r] \arrow[d] & \RHom(f_!E\otimes^L_SP',f_!f^!Q\otimes^L_SQ') \arrow[d]\\
\RHom(f_!E,Q)\otimes^L_S\RHom(P',Q') \arrow[r] & \RHom(f_!E\otimes^L_SP',Q\otimes^L_SQ'),
\end{tikzcd}
\]
defined by 2.2.4 and the adjunction arrow $f_!f^!Q\to Q$, show that, to establish the commutativity of 2.6.2, it suffices to prove the following compatibility.

\subsection*{Lemma 2.6.3}
Let $f$, $f'$, $E$, $F$, $P'$, $Q'$ be as in 2.6.1. Then one has a commutative square
\[
\begin{tikzcd}[column sep=large,row sep=large]
f_*\RHom(E,F)\otimes^L_S\RHom(P',Q') \arrow[r,"(1)"] \arrow[d,"(2)"'] & f'_*\RHom(E\otimes^L_SP',F\otimes^L_SQ') \arrow[d,"(4)"]\\
\RHom(f_!E,f_!F)\otimes^L_S\RHom(P',Q') \arrow[r,"(3)"'] & \RHom(f_!E\otimes^L_SP',f_!F\otimes^L_SQ'),
\end{tikzcd}
\]
where (1) is the composite of the base-change isomorphism
\[
f_*\RHom(E,F)\otimes^L_S -\xrightarrow{\sim}f'_*\bigl(\RHom(E,F)\otimes^L_S-\bigr)
\]
and $f'_*(2.2.4)$, (2) is deduced from the canonical arrow
\[
f_*\RHom(E,F)\longrightarrow\RHom(f_!E,f_!F)
\]
(SGA~4~XVIII~3.1.9.5) by applying $\otimes^L_S\RHom(P',Q')$, (3) is given by 2.2.4, and finally (4) is the composite of the canonical arrow
\[
f'_*\RHom(E\otimes^L_SP',F\otimes^L_SQ')
\longrightarrow
\RHom(f'_!(E\otimes^L_SP'),f'_!(F\otimes^L_SQ'))
\]
(loc. cit.) and the arrows defined by the base-change isomorphisms
\[
f'_!(E\otimes^L_SP')\xrightarrow{\sim}f_!E\otimes^L_SP',
\qquad
f'_!(F\otimes^L_SQ')\xrightarrow{\sim}f_!F\otimes^L_SQ'.
\]
Returning to the definition of the arrow (SGA~4~XVIII~3.1.9.5), choose a compactification of $f$\footnote{The arrow quoted is defined with the aid of such a compactification, but it does not depend on it: this point, tacitly admitted in (loc. cit.), is easily verified.}; one is then reduced to proving the lemma separately in the case where $f$ is an open immersion and in the case where $f$ is a proper morphism.

\emph{a) Case where $f$ is an open immersion.} By adjunction (2.2.1), it suffices to see that, if one tensors the square 2.6.3 by $f_!E\otimes^L_SP'$ and composes with the evaluation arrow 2.2.2
\[
(f_!E\otimes^L_SP')\otimes^L\RHom(f_!E\otimes^L_SP',f_!F\otimes^L_SQ')\longrightarrow f_!F\otimes^L_SQ',
\]
the square obtained is commutative. But the sums of this square are zero outside $X'$, so that it suffices to verify that the square induced on $X'$ commutes, which is trivial.

\emph{b) Case where $f$ is a proper morphism.} Then $f_!=f_*$, $f'_!=f'_*$. One sees easily that the composite of
\[
f_*\RHom(E,F)\longrightarrow\RHom(f_*E,f_*F)
\]
with the trivial-duality isomorphism
\[
\RHom(f_*E,f_*F)\xrightarrow{\sim}f_*\RHom(f^{*}f_*E,F)
\]
(2.1.4) is none other than the arrow deduced from the adjunction arrow $f^{*}f_*E\to E$ by applying $f_*\RHom(-,F)$. Similarly, the composite of (4) with the trivial-duality isomorphism
\[
\RHom(f_*E\otimes^L_SP',f_*F\otimes^L_SQ')\xrightarrow{\sim}f'_*\RHom(f^{*}f_*E\otimes^L_SP',F\otimes^L_SQ')
\]
($f_*F\otimes^L_SQ'$ being identified with $f'_*(F\otimes^L_SQ')$ by the canonical isomorphism) is none other than the arrow deduced from the adjunction arrow $f^{*}f_*E\to E$ by applying $f'_*\RHom(-\otimes^L_SP',F\otimes^L_SQ')$. Thus one is reduced to verifying the commutativity of the square
\[
\begin{tikzcd}[column sep=large,row sep=large]
\RHom(f_*E,f_*F)\otimes^L_S\RHom(P',Q') \arrow[r,"(3)"] \arrow[d] & \RHom(f_*E\otimes^L_SP',f_*F\otimes^L_SQ') \arrow[d]\\
f_*\RHom(f^{*}f_*E,F)\otimes^L_S\RHom(P',Q') \arrow[r,"(1')"'] & f'_*\RHom(f^{*}f_*E\otimes^L_SP',F\otimes^L_SQ'),
\end{tikzcd}
\]
where the vertical arrows are the trivial-duality isomorphisms and $(1')$ is (1) with $E$ replaced by $f^{*}f_*E$. But this is a special case of 2.6.1.1 in the case, already treated, where $(u,v)$ is of type $(2,1)$ (take $u:f^{*}f_*E\to f_*E$ given by the identity of $f^{*}f_*E$, and $v:F\to f_*F$ given by the identity of $f_*F$). We have therefore proved 2.6.3. As we have seen, the commutativity of 2.6.2 follows, which completes the proof of 2.6.1 and, finally, of 2.5.

\clearpage
\section{Cohomological Correspondences}

\subsection*{3.1}
Let, for $i=1,2$, $X_i$ be an $S$-scheme and $L_i\in\ob\Dctf(X_i)$, and let
\[
X=X_1\times_S X_2,
\qquad
p_i:X\longrightarrow X_i,
\qquad
q_i:X_i\longrightarrow S
\]
be the canonical projections:
\[
\begin{tikzcd}[column sep=1.8em,row sep=1.5em]
& X \arrow[dl,"P_1"'] \arrow[dr,"P_2"] & \\
L_1\quad X_1 \arrow[dr,"q_1"'] & & X_2\quad L_2 \arrow[dl,"q_2"] \\
& S &
\end{tikzcd}
\]
With the notation of 1.3, one has a canonical isomorphism (2.2.4)
\[
DL_1\overset{\mathbf L}{\otimes}_S L_2
\xrightarrow{\sim}
\RHom\bigl(L_1\overset{\mathbf L}{\otimes}_S A_{X_2},
q_1^!A_S\overset{\mathbf L}{\otimes}_S L_2\bigr).
\]
Composing with the Künneth isomorphism
\[
q_1^!A_S\overset{\mathbf L}{\otimes}_S L_2
\xrightarrow{\sim}p_2^!L_2 \qquad (1.7.3),
\]
one obtains an isomorphism
\begin{equation}
DL_1\overset{\mathbf L}{\otimes}_S L_2
\xrightarrow{\sim}
\RHom(p_1^*L_1,p_2^!L_2).
\tag{3.1.1}
\end{equation}
This isomorphism will play a fundamental role throughout. For $X_i=S$, 3.1.1 reduces to
\[
L_1^\vee\overset{\mathbf L}{\otimes}L_2\xrightarrow{\sim}\RHom(L_1,L_2),
\]
a special case of an isomorphism valid for perfect complexes on an arbitrary ringed topos (SGA 6 I 7.7). We set
\begin{equation}
\RHom(p_1^*L_1,p_2^!L_2)=\RHom_S(L_1,L_2),
\qquad
\Hom(p_1^*L_1,p_2^!L_2)=\Hom_S(L_1,L_2).
\tag{3.1.2}
\end{equation}
The elements of $\Hom_S(L_1,L_2)$ will be called \emph{cohomological correspondences from $L_1$ to $L_2$}. Instead of ``cohomological correspondence from $A_{X_1}$ to $A_{X_2}$'', we shall simply say \emph{cohomological correspondence from $X_1$ to $X_2$} (understood: with coefficients in $A$).

If $X_1$ is \emph{smooth}, purely of relative dimension $r_1$, then the same is true of $p_2$, and by the duality theorem (SGA 4 XVIII 3.2.5) one has
\[
p_2^!L_2\xrightarrow{\sim}p_2^*L_2(r_1)[2r_1],
\]
therefore
\begin{equation}
\RHom_S(L_1,L_2)=\RHom(p_1^*L_1,p_2^*L_2)(r_1)[2r_1],
\qquad
\Hom_S(L_1,L_2)=\H^{2r_1}\bigl(X,\RHom(p_1^*L_1,p_2^*L_2)(r_1)\bigr).
\tag{3.1.3}
\end{equation}
In particular, the group of cohomological correspondences from $X_1$ to $X_2$ is identified in this case with $\H^{2r_1}(X,A(r_1))$.

\subsection*{3.2}
The cohomological correspondences met in practice are associated with $S$-morphisms $X_2\to X_1$, or more generally with cycles on $X$, or with schemes over $X$. Let
\[
c:C\longrightarrow X
\]
be a scheme over $X$ (one sometimes says that $c$ is a correspondence between $X_1$ and $X_2$), and put $c_i=p_ic:C\to X_i$ for the projections. We shall call \emph{cohomological correspondences from $L_1$ to $L_2$ with support in $c$} (or in $C$, if there is no risk of confusion) the elements of the group
\[
\H^0\bigl(X,c_!\RHom(c_1^*L_1,c_2^!L_2)\bigr)
\quad(=\Hom(c_1^*L_1,c_2^!L_2)\text{ if }c\text{ is proper}).
\]
By the induction formula (SGA 4 XVIII 3.1.12.2), one has a canonical isomorphism
\begin{equation}
c^!\RHom(p_1^*L_1,p_2^!L_2)
\xrightarrow{\sim}
\RHom(c_1^*L_1,c_2^!L_2),
\tag{3.2.1}
\end{equation}
whence, applying $c_!$ and composing with the adjunction arrow $c_!c^!\to \Id$, canonical arrows
\begin{equation}
c_!\RHom(c_1^*L_1,c_2^!L_2)
\longrightarrow
\RHom(p_1^*L_1,p_2^!L_2),
\tag{3.2.2}
\end{equation}
\begin{equation}
\H^0\bigl(X,c_!\RHom(c_1^*L_1,c_2^!L_2)\bigr)
\longrightarrow
\Hom_S(L_1,L_2).
\tag{3.2.3}
\end{equation}
Thus every correspondence from $L_1$ to $L_2$ with support in $C$ canonically defines, by 3.2.3, a correspondence (without support) from $L_1$ to $L_2$.

\emph{Examples.} a) Suppose that $c$ is proper and that $c_2$ has finite tor-dimension and relative dimension $\leq N$. The trace morphism (SGA $4^{1/2}$ Cycle 2.3.3)
\[
\mathrm{Tr}_{c_2}:c_{2!}A_C\longrightarrow A_{X_2}(-N)[-2N]
\]
defines by adjunction a cohomological correspondence from $A_{X_1}$ to $A_{X_2}(-N)[2N]$ with support in $C$, which will be called the \emph{class of $(c,N)$}, and will be denoted
\begin{equation}
\operatorname{cl}(c,N)\in \H^{-2N}\bigl(C,c_2^!A_{X_2}(-N)\bigr).
\tag{3.2.4}
\end{equation}
The most important case is that in which $N=0$, i.e. $c_2$ is finite (and of finite tor-dimension); in that case $\operatorname{cl}(c)$ is a cohomological correspondence from $X_1$ to $X_2$ with support in $C$. If $X_1$ is smooth, purely of relative dimension $r_1$, then, by 3.1.3 and 3.2.1,
\begin{equation}
\H^{-2N}\bigl(C,c_2^!A_{X_2}(-N)\bigr)
=
\H^{2(r_1-N)}\bigl(C,c^!A_X(r_1-N)\bigr).
\tag{3.2.5}
\end{equation}
If moreover $c$ is a closed immersion, the second member rewrites as
\[
\H^{2(r_1-N)}_C\bigl(X,A_X(r_1-N)\bigr),
\]
and the class of $(c,N)$, considered as an element of this group, is precisely the cohomology class associated with $c$ in the sense of (SGA $4^{1/2}$ Cycle 2.3.1), as follows immediately from the definitions.

b) Let $F:X_2\to X_1$ be an $S$-morphism, and denote by $c(F):C(F)\to X$ its graph, i.e. the image of the section of $p_2$ defined by $F$ (``the set of points $(Fx,x)$''): $c_2$ is an isomorphism and $c_1c_2^{-1}=F$. One therefore has canonical isomorphisms
\begin{equation}
\RHom\bigl(c_1(F)^*L_1,c_2(F)^!L_2\bigr)=\RHom(F^*L_1,L_2),
\qquad
\Hom\bigl(c_1(F)^*L_1,c_2(F)^!L_2\bigr)=\Hom(F^*L_1,L_2).
\tag{3.2.6}
\end{equation}
In other words, the cohomological correspondences from $L_1$ to $L_2$ with support in the graph of $F$ are the ``lifts'' of $F$ to a homomorphism $F^*L_1\to L_2$. Here are two important examples of such correspondences:

(i) $X_1=X_2$, $L_1=L_2$, $F=1$ is the identity of $X_1$, whose graph $\Delta$ is the diagonal. The correspondence with support in $\Delta$ defined by the identity of $L_1$ is called the \emph{diagonal correspondence}. It is the element
\[
1\in \Hom(L_1,L_1)=\H^0\bigl(X,\Delta_*\Delta^!\RHom_S(L_1,L_1)\bigr)
\]
(a group that is identified with $\H^{2r}_\Delta(X,\RHom_S(L_1,L_1)(r))$ when $X_1$ is smooth, purely of relative dimension $r$).

(ii) Suppose that $S$ is the spectrum of a finite field $\mathbb F_q$, that $X_1=X_2$, $L_1=L_2$, and take for $F:X_1\to X_1$ the relative Frobenius endomorphism over $\mathbb F_q$, the identity on the underlying space and raising to the $q$-th power on $\mathcal O_{X_1}$ (cf. XIV), and take as lift of $F$ to $L_1$ the canonical isomorphism
\[
F^*L_1\xrightarrow{\sim}L_1
\]
(loc. cit. §2, Prop. 1) (more precisely, the $a$-th iterate of the isomorphism of (loc. cit.) if $q=p^a$). The cohomological correspondence from $L_1$ to $L_1$ with support in the graph of $F$ so defined is called the \emph{Frobenius correspondence} (relative to $\mathbb F_q$).

c) Certain cohomological correspondences between sheaves on curves, related to Hecke operators, have been studied by Langlands in ([13] §7).

d) Cohomological correspondences between constant sheaves (the most frequent ones in nature) furnish in a natural way, by ``integration'', cohomological correspondences between complexes, as we shall now see.

\subsection*{3.3}
Let, for $i=1,2$, $f_i:X_i\to Y_i$ be an $S$-morphism, and put
\[
f=f_1\times_S f_2:X\to Y,
\qquad
f'_1=f_1\times_S Y_2,
\qquad
f'_2=Y_1\times_S f_2,
\]
\[
p'_1=X_1\times_S f_2,
\qquad
p'_2=f_1\times_S X_2.
\]
\begin{equation}
\begin{tikzcd}[column sep=3.8em,row sep=2.4em]
X_1\times_S Y_2 \arrow[d,"f'_1"'] & X \arrow[l,"p'_1"'] \arrow[dl,"f" description] \arrow[d,"p'_2"] \\
Y & Y_1\times_S X_2 \arrow[l,"f'_2"]
\end{tikzcd}
\tag{3.3.0}
\end{equation}
Let $L_i\in\ob\Dctf(X_i)$. There exists a unique isomorphism
\begin{equation}
f_*\RHom_S(L_1,L_2)
\xrightarrow{\sim}
\RHom_S(f_{1!}L_1,f_{2*}L_2)
\tag{3.3.1}
\end{equation}
whose inverse makes the square
\[
\begin{tikzcd}[column sep=4.5em,row sep=2.6em]
D(f_{1!}L_1)\overset{\mathbf L}{\otimes}_S f_{2*}L_2
  \arrow[r,"3.1.1"] \arrow[d,"(1)"']
& \RHom_S(f_{1!}L_1,f_{2*}L_2) \arrow[d] \\
f_{1*}DL_1\overset{\mathbf L}{\otimes}_S f_{2*}L_2
  \arrow[r,"(2)"']
& f_*\RHom_S(L_1,L_2),
\end{tikzcd}
\]
commutative, where (1) is deduced from the inverse of the duality isomorphism (SGA 4 XVIII 3.1.9.6)
\[
Df_{1!}L_1\xrightarrow{\sim}f_{1*}DL_1
\]
by applying $-\overset{\mathbf L}{\otimes}_S f_{2*}L_2$, and (2) is the composite of the Künneth isomorphism
\[
f_{1*}DL_1\overset{\mathbf L}{\otimes}_S f_{2*}L_2
\xrightarrow{\sim}
f_*(DL_1\overset{\mathbf L}{\otimes}_S L_2)
\]
and $f_*(3.1.1)$. Applying $R\Gamma(Y,-)$ to 3.3.1, one obtains an isomorphism
\begin{equation}
f_*:\Hom_S(L_1,L_2)
\xrightarrow{\sim}
\Hom_S(f_{1!}L_1,f_{2*}L_2).
\tag{3.3.2}
\end{equation}
For $u\in\Hom_S(L_1,L_2)$, the correspondence $f_*u$ will be called the \emph{direct image} of $u$ by $f$. When $Y_1=Y_2=S$, we shall write $u_*$ rather than $f_*u$. If $Y_1=Y_2=S$ is the spectrum of a separably closed field, 3.3.2 reads
\begin{equation}
\Hom_S(L_1,L_2)
\xrightarrow{\sim}
\Hom\bigl(R\Gamma_c(X_1,L_1),R\Gamma(X_2,L_2)\bigr).
\tag{3.3.3}
\end{equation}
Thus cohomological correspondences ``operate'' on cohomology. The following statement will allow us to interpret this operation in several ways, and in particular to relate it to the usual definitions in situations such as a) or b) above.

\begin{statement}{Proposition 3.4}
With the notation of 3.3.0, one has a commutative diagram of isomorphisms in $\Dctf(Y)$:
\begin{equation}
\begin{tikzcd}[column sep=4.0em,row sep=2.8em]
& \RHom_S(f_{1!}L_1,f_{2*}L_2)
  \arrow[dl,"(1)"'] \arrow[dr,"(2)"] \arrow[dd,"(5)" description] & \\
f'_{1*}\RHom_S(L_1,f_{2*}L_2) \arrow[dr,"f'_{1*}(3)"'] &&
f'_{2*}\RHom_S(f_{1!}L_1,L_2) \arrow[dl,"f'_{2*}(4)"] \\
& f_*\RHom_S(L_1,L_2) &
\end{tikzcd}
\tag{3.4.1}
\end{equation}
where (5) is the inverse of 3.3.1 and the arrows (1) to (4) are defined by the commutative diagrams below; in these diagrams ``base ch.'' denotes an arrow deduced from a base-change or Künneth isomorphism of one of the types considered in no.~1, ``dual.'' an arrow deduced from a duality isomorphism of type 2.1.4 or (SGA 4 XVIII 3.1.9.6), and ``triv.'' a trivial isomorphism:
\[
\begin{tikzcd}[column sep=3.7em,row sep=1.7em]
\RHom(f_{1!}L_1\overset{\mathbf L}{\otimes}_S A_{Y_2},K_{Y_1}\overset{\mathbf L}{\otimes}_S f_{2*}L_2)
  \arrow[r,"\mathrm{base\ ch.}"] \arrow[d,"\mathrm{base\ ch.}"']
& \RHom_S(f_{1!}L_1,f_{2*}L_2) \arrow[ddd,"(1)"] \\
\RHom(f'_{1!}(L_1\overset{\mathbf L}{\otimes}_S A_{Y_2}),K_{Y_1}\overset{\mathbf L}{\otimes}_S f_{2*}L_2) \arrow[d,"\mathrm{dual.}"']
& \\
f'_{1*}\RHom(L_1\overset{\mathbf L}{\otimes}_S A_{Y_2},f_1'^!(K_{Y_1}\overset{\mathbf L}{\otimes}_S f_{2*}L_2)) \arrow[d,"\mathrm{base\ ch.}"']
& \\
f'_{1*}\RHom(L_1\overset{\mathbf L}{\otimes}_S A_{Y_2},K_{X_1}\overset{\mathbf L}{\otimes}_S f_{2*}L_2) \arrow[r,"\mathrm{base\ ch.}"']
& f'_{1*}\RHom_S(L_1,f_{2*}L_2),
\end{tikzcd}
\]
\[
\begin{tikzcd}[column sep=3.7em,row sep=1.7em]
\RHom(f_{1!}L_1\overset{\mathbf L}{\otimes}_S A_{Y_2},K_{Y_1}\overset{\mathbf L}{\otimes}_S f_{2*}L_2)
  \arrow[r,"\mathrm{base\ ch.}"] \arrow[d,"\mathrm{base\ ch.}"']
& \RHom_S(f_{1!}L_1,f_{2*}L_2) \arrow[ddd,"(2)"] \\
\RHom(f_{1!}L_1\overset{\mathbf L}{\otimes}_S A_{Y_2},f'_{2*}(K_{Y_1}\overset{\mathbf L}{\otimes}_S L_2)) \arrow[d,"\mathrm{dual.}"']
& \\
f'_{2*}\RHom(f_2'^*(f_{1!}L_1\overset{\mathbf L}{\otimes}_S A_{Y_2}),K_{Y_1}\overset{\mathbf L}{\otimes}_S L_2) \arrow[d,"\mathrm{triv.}"']
& \\
f'_{2*}\RHom(f_{1!}L_1\overset{\mathbf L}{\otimes}_S A_{X_2},K_{Y_1}\overset{\mathbf L}{\otimes}_S L_2) \arrow[r,"\mathrm{base\ ch.}"']
& f'_{2*}\RHom_S(f_{1!}L_1,L_2),
\end{tikzcd}
\]
\[
\begin{tikzcd}[column sep=3.7em,row sep=1.7em]
\RHom(L_1\overset{\mathbf L}{\otimes}_S A_{Y_2},K_{X_1}\overset{\mathbf L}{\otimes}_S f_{2*}L_2)
  \arrow[r,"\mathrm{base\ ch.}"] \arrow[d,"\mathrm{base\ ch.}"']
& \RHom_S(L_1,f_{2*}L_2) \arrow[ddd,"(3)"] \\
\RHom(L_1\overset{\mathbf L}{\otimes}_S A_{Y_2},p'_{1*}(K_{X_1}\overset{\mathbf L}{\otimes}_S L_2)) \arrow[d,"\mathrm{dual.}"']
& \\
p'_{1*}\RHom(p_1'^*(L_1\overset{\mathbf L}{\otimes}_S A_{Y_2}),K_{X_1}\overset{\mathbf L}{\otimes}_S L_2) \arrow[d,"\mathrm{triv.}"']
& \\
p'_{1*}\RHom(L_1\overset{\mathbf L}{\otimes}_S A_{X_2},K_{X_1}\overset{\mathbf L}{\otimes}_S L_2) \arrow[r,"\mathrm{base\ ch.}"']
& p'_{1*}\RHom_S(L_1,L_2),
\end{tikzcd}
\]
\[
\begin{tikzcd}[column sep=3.7em,row sep=1.7em]
\RHom(f_{1!}L_1\overset{\mathbf L}{\otimes}_S A_{X_2},K_{Y_1}\overset{\mathbf L}{\otimes}_S L_2)
  \arrow[r,"\mathrm{base\ ch.}"] \arrow[d,"\mathrm{base\ ch.}"']
& \RHom_S(f_{1!}L_1,L_2) \arrow[ddd,"(4)"] \\
\RHom(p'_{2!}(L_1\overset{\mathbf L}{\otimes}_S A_{X_2}),K_{Y_1}\overset{\mathbf L}{\otimes}_S L_2) \arrow[d,"\mathrm{dual.}"']
& \\
p'_{2*}\RHom(L_1\overset{\mathbf L}{\otimes}_S A_{X_2},p_2'^!(K_{Y_1}\overset{\mathbf L}{\otimes}_S L_2)) \arrow[d,"\mathrm{base\ ch.}"']
& \\
p'_{2*}\RHom(L_1\overset{\mathbf L}{\otimes}_S A_{X_2},K_{X_1}\overset{\mathbf L}{\otimes}_S L_2) \arrow[r,"\mathrm{base\ ch.}"']
& p'_{2*}\RHom_S(L_1,L_2).
\end{tikzcd}
\]
\end{statement}
This is only an explication of 2.5 in the particular case where $(u_1,v_1):(L_1,K_{X_1})\to(f_{1!}L_1,K_{Y_1})$ is given by the identity of $f_{1!}L_1$ and the canonical isomorphism $f_1^!K_{Y_1}\xrightarrow{\sim}K_{X_1}$, and where $(u_2,v_2):(A_{X_2},L_2)\to(A_{Y_2},f_{2*}L_2)$ is given by the canonical isomorphism $A_{X_2}\xrightarrow{\sim}f_2^!A_{Y_2}$ and the identity of $f_{2*}L_2$.

\begin{statement}{Corollary 3.5}
Assume that $Y_1=Y_2=S$. Let $u\in\Hom_S(L_1,L_2)$. The direct image $f_*u\in\Hom(f_{1!}L_1,f_{2*}L_2)$ (3.3.2) can be obtained in the following two equivalent ways:
\end{statement}
\begin{enumerate}[label=\alph*)]
\item Identify $u$, by trivial duality, with an arrow
\[
L_1\longrightarrow p_{1*}p_2^!L_2;
\]
compose with the base-change isomorphism
\[
p_{1*}p_2^!L_2\xrightarrow{\sim}f_1^!f_{2*}L_2;
\]
$f_*u$ is the adjoint of this composite arrow.
\item Identify $u$, by adjunction, with an arrow
\[
p_{2!}p_1^*L_1\longrightarrow L_2;
\]
compose with the base-change isomorphism
\[
f_2^*f_{1!}L_1\xrightarrow{\sim}p_{2!}p_1^*L_1;
\]
$f_*u$ is the adjoint of this composite arrow.
\end{enumerate}
Description a) (respectively b)) indeed corresponds to the path $(f_{1*}(3)\circ(1))^{-1}$ (respectively $f_{2*}(4)\circ(2)$) in 3.4.1.

\subsection*{3.6}
In practice, it is mainly description 3.5 b) that is used. Let us spell it out at the level of cohomology classes, assuming, in order to lighten the notation, that $S$ is the spectrum of a separably closed field. Let
\[
u\in\Hom_S(L_1,L_2)=\Hom(p_1^*L_1,p_2^!L_2).
\]
For $x\in\H_c^i(X_1,L_1)$, $u_*(x)\in\H^i(X_2,L_2)$ is defined as follows: consider the element $p_1^*x\in\H^i(X_2,p_{2!}p_1^*L_1)$ deduced from $x$ by inverse image, apply $u$ to it, i.e. form the product
\[
u\cdot p_1^*x\in \H^i(X_2,p_{2!}p_2^!L_2),
\]
then apply the arrow
\[
a:\H^i(X_2,p_{2!}p_2^!L_2)\longrightarrow \H^i(X_2,L_2)
\]
defined by the adjunction arrow $p_{2!}p_2^!\to\Id$; in short:
\begin{equation}
u_*(x)=a(u\cdot p_1^*x).
\tag{3.6.1}
\end{equation}
Let $c:C\to X$ be a proper morphism, and let $u\in\Hom(c_1^*L_1,c_2^!L_2)$. Denote by
\[
u_*\in\Hom\bigl(R\Gamma_c(X_1,L_1),R\Gamma(X_2,L_2)\bigr)
\]
the image of $u$ under the composite of 3.2.3 and 3.3.3. It follows from 3.6.1 that, for $x\in\H_c^i(X_1,L_1)$, one has
\begin{equation}
u_*(x)=a(u\cdot c_1^*x),
\tag{3.6.2}
\end{equation}
where $c_1^*x$ is the inverse image of $x$ in $\H^i(X_2,c_{2!}c_1^*L_1)$ and
\[
a:\H^i(X_2,c_{2!}c_2^!L_2)\longrightarrow \H^i(X_2,L_2)
\]
is the arrow deduced from the adjunction $c_{2!}c_2^!L_2\to\Id$. In the situation of Example 3.2 a), one thus finds
\begin{equation}
\operatorname{cl}(c,N)_*(x)=\mathrm{Tr}_{c_2}(c_1^*x),
\tag{3.6.3}
\end{equation}
where
\[
\mathrm{Tr}_{c_2}:\H^i(X_2,Rc_{2!}A_C)\longrightarrow\H^{i-2N}(X_2,A(-N))
\]
is the arrow induced by the trace morphism relative to $c_2$. For $N=0$, this formula agrees with (SGA $4^{1/2}$ Cycle 3.6): $\operatorname{cl}(c)_*$ is the homomorphism denoted $\eta_{X_2,X_1}$ there. If $X_1$ is smooth, purely of relative dimension $r_1$, and $c$ is a closed immersion, it follows from 3.6.1 and from the remark made at the end of 3.2 a) that one also has
\begin{equation}
\operatorname{cl}(c,N)_*(x)=\mathrm{Tr}_{p_2}(\operatorname{cl}(c)\,p_1^*x),
\tag{3.6.4}
\end{equation}
where $\operatorname{cl}(C)\in \H_C^{2(r_1-N)}(X,A(r_1-N))$ is the cohomology class of $C$ and
\[
\mathrm{Tr}_{p_2}:\H^*(X_2,Rp_{2!}A_X)\longrightarrow \H^{*-2r_1}(X_2,A(-r_1))
\]
is the arrow defined by the trace morphism relative to $p_2$. This recovers the usual definition of the homomorphism induced on cohomology by a correspondence ([12] 1.3), (SGA $4^{1/2}$ Cycle 3.2).

\subsection*{3.6.5}
In the situation of 3.2 b), it follows from 3.6.2 that, for $u\in\Hom(F^*L_1,L_2)$,
\[
u_*:\H_c^i(X_1,L_1)\longrightarrow \H^i(X_2,L_2)
\]
is the composite of the ``inverse image'' homomorphism
\[
\H_c^i(X_1,L_1)\longrightarrow \H^i(X_2,F^*L_1)
\]
and the homomorphism induced by $u$ (more generally, $f_*u:Rf_{1!}L_1\to Rf_{2*}L_2$ is the composite of the canonical homomorphism
\[
Rf_{1!}L_1\longrightarrow Rf_{2*}F^*L_1
\]
and of $Rf_{2*}u$). In particular, the diagonal correspondence induces the canonical arrow
\[
Rf_{1!}L_1\longrightarrow Rf_{1*}L_1
\]
(the identity if $X_1$ is proper over $S$).

\subsection*{3.7}
Return to the situation of the beginning of 3.3, assuming in addition that a commutative square
\begin{equation}
\begin{tikzcd}[column sep=3em,row sep=2em]
X \arrow[d,"f"'] & C \arrow[l,"c"'] \arrow[d,"g"] \\
Y & D \arrow[l,"d"]
\end{tikzcd}
\tag{3.7.1}
\end{equation}
is given, with $c$ proper. Put $c_i=p_ic$, $d_i=q_id$, where $p_i:X\to X_i$, $q_i:Y\to Y_i$ are the projections. From 3.7.1 one deduces a commutative diagram whose inner square is cartesian:
\begin{equation}
\begin{tikzpicture}[>=latex,baseline=(current bounding box.center),scale=1]
\node (X) at (0,0) {$X$};
\node (Y) at (0,-1.45) {$Y$};
\node (Cp) at (1.65,0) {$C'$};
\node (D) at (1.65,-1.45) {$D$};
\node (C) at (2.55,1.1) {$C$};
\draw[->] (C) -- node[above,sloped,pos=.55] {$c$} (X);
\draw[->] (C) -- node[above,sloped,pos=.50] {$i$} (Cp);
\draw[->] (C) -- node[right,pos=.55] {$g$} (D);
\draw[->] (Cp) -- node[above] {$c'$} (X);
\draw[->] (Cp) -- node[right] {$g'$} (D);
\draw[->] (X) -- node[left] {$f$} (Y);
\draw[->] (D) -- node[below] {$d$} (Y);
\end{tikzpicture}
\tag{3.7.2}
\end{equation}
For $K\in\ob\Dctf(X)$, there is a canonical arrow
\begin{equation}
f_*c_!c^!K\longrightarrow d_*d^!f_*K,
\tag{3.7.3}
\end{equation}
which is the composite of the arrow
\[
f_*c_!c^!K=f_*c'_*(i_*i^!)c'^!K\longrightarrow f_*c'_*c'^!K=d_*g'_*c'^!K
\]
defined by the adjunction arrow $i_*i^!\to\Id$, and the isomorphism
\[
d_*g'_*c'^!K\xrightarrow{\sim}d_*d^!f_*K
\]
defined by the canonical isomorphism $g'_*d'^!\xrightarrow{\sim}d^!f_*$ (SGA 4 XVIII 3.1.12.3). Taking $K=\RHom_S(L_1,L_2)$, and composing 3.7.3 with the isomorphism deduced from 3.3.1, one obtains a canonical arrow
\begin{equation}
f_*c_*c^!\RHom_S(L_1,L_2)
\longrightarrow
 d_*d^!\RHom_S(f_{1!}L_1,f_{2*}L_2),
\tag{3.7.4}
\end{equation}
which, by 3.2.1, may be rewritten as
\begin{equation}
f_*c_!\RHom(c_1^*L_1,c_2^!L_2)
\longrightarrow
 d_*\RHom(d_1^*(f_{1!}L_1),d_2^!(f_{2*}L_2)).
\tag{3.7.5}
\end{equation}
One deduces
\begin{equation}
f_*:\Hom(c_1^*L_1,c_2^!L_2)
\longrightarrow
\Hom(d_1^*(f_{1!}L_1),d_2^!(f_{2*}L_2)).
\tag{3.7.6}
\end{equation}
The arrow 3.7.5 (respectively 3.7.6) generalizes 3.3.1 (respectively 3.3.2), with which it coincides when $c$ and $d$ are the identity. Also note that, for $c$ proper, 3.2.2 (respectively 3.2.3) is nothing other than 3.7.5 (respectively 3.7.6) when $f_1$, $f_2$, and $d$ are the identity arrows. When $Y_1=Y_2=S=D$, we shall abbreviate $f_*u$ as $u_*$, as above.

\section{Pairings of Correspondences. Lefschetz Formula}

\subsection*{4.1}
Let, for $i=1,2$, $X_i$ be an $S$-scheme, $X=X_1\times_SX_2$, $p_i:X\to X_i$ the projections, $L_i\in\ob\Dctf(X_i)$, and $K_{X_i}$ the dualizing complex (1.3). The canonical arrows
\begin{equation}
p_i^*\RHom(L_i,K_{X_i})\overset{\mathbf L}{\otimes}p_i^*L_i
\longrightarrow p_i^*K_{X_i},
\tag{4.1.1}
\end{equation}
composed from 2.1.1 and the evaluation arrow 2.2.2, give, by tensor product on $X$ and composition with the Künneth isomorphism 1.7.6, a pairing
\begin{equation}
(DL_1\overset{\mathbf L}{\otimes}_S L_2)\overset{\mathbf L}{\otimes}(L_1\overset{\mathbf L}{\otimes}_S DL_2)
\longrightarrow K_X,
\tag{4.1.2}
\end{equation}
which may be rewritten, by 3.1.1 and 3.1.2, as
\begin{equation}
\RHom_S(L_1,L_2)\overset{\mathbf L}{\otimes}\RHom_S(L_2,L_1)
\longrightarrow K_X.
\tag{4.1.3}
\end{equation}
One deduces a pairing, called cup product,
\begin{equation}
\langle\ ,\ \rangle:\Hom_S(L_1,L_2)\otimes\Hom_S(L_2,L_1)
\longrightarrow \H^0(X,K_X).
\tag{4.1.4}
\end{equation}
One may also regard 4.1.2 as the external tensor product of the evaluation arrows
\[
DL_i\overset{\mathbf L}{\otimes}L_i\longrightarrow K_{X_i}.
\]
In view of the Künneth isomorphism 2.2.4, one thus sees that 4.1.2 is a ``perfect duality'', i.e. identifies by the adjunction isomorphism 2.2.1 each of the complexes $DL_1\overset{\mathbf L}{\otimes}_S L_2$ and $L_1\overset{\mathbf L}{\otimes}_S DL_2$ with the dual of the other (with values in $K_X$).

When $X_1=X_2=S$, 4.1.3 is a special case of a pairing valid for perfect complexes on an arbitrary ringed topos (SGA 6 I 7.8), and, with the notation of (SGA 6 I 8.3), one has
\begin{equation}
\langle u,v\rangle=\mathrm{Tr}(vu)=\mathrm{Tr}(uv).
\tag{4.1.5}
\end{equation}

\subsection*{4.2}
Let $c:C\to X$, $d:D\to X$, and $e=c\times_Xd:E=C\times_XD\to X$. For $P,Q\in\ob\Dctf(X)$, one has a canonical arrow (which is not in general an isomorphism)
\begin{equation}
c^!P\overset{\mathbf L}{\otimes}_X d^!Q
\longrightarrow e^!(P\overset{\mathbf L}{\otimes}Q),
\tag{4.2.1}
\end{equation}
defined, in the same way as 1.7.3, as the adjoint of the composite arrow
\[
e_!(c^!P\overset{\mathbf L}{\otimes}_X d^!Q)
\xrightarrow[(1)]{\sim}
 c_!c^!P\overset{\mathbf L}{\otimes}d_!d^!Q
\xrightarrow{(2)} P\overset{\mathbf L}{\otimes}Q,
\]
where (1) is the Künneth isomorphism (SGA 4 XVII 5.4.2.2) and (2) is the tensor product of the adjunction arrows. Applying $e_!$ to 4.2.1 and composing with the Künneth isomorphism just cited, one obtains a canonical arrow
\begin{equation}
c_!c^!P\overset{\mathbf L}{\otimes}d_!d^!Q
\longrightarrow e_!e^!(P\overset{\mathbf L}{\otimes}Q).
\tag{4.2.2}
\end{equation}
Similarly, composing the Künneth arrow
\[
c_*c^!P\overset{\mathbf L}{\otimes}d_*d^!Q
\longrightarrow e_*(c^!P\overset{\mathbf L}{\otimes}_X d^!Q)
\]
with the arrow deduced from 4.2.1 by applying $e_*$, one obtains a canonical arrow
\begin{equation}
c_*c^!P\overset{\mathbf L}{\otimes}d_*d^!Q
\longrightarrow e_*e^!(P\overset{\mathbf L}{\otimes}Q).
\tag{4.2.3}
\end{equation}
For $P=DL_1\overset{\mathbf L}{\otimes}_S L_2$, $Q=L_1\overset{\mathbf L}{\otimes}_S DL_2$, one deduces from 4.2.2 and 4.2.3, by composition with 4.1.2 and the isomorphism $e^!K_X=K_E$, and with the identifications 3.1.1 and 3.2.1, pairings
\begin{equation}
c_!\RHom(c_1^*L_1,c_2^!L_2)\overset{\mathbf L}{\otimes}
 d_!\RHom(d_2^*L_2,d_1^!L_1)
\longrightarrow e_!K_E,
\tag{4.2.4}
\end{equation}
\begin{equation}
c_*\RHom(c_1^*L_1,c_2^!L_2)\overset{\mathbf L}{\otimes}
 d_*\RHom(d_2^*L_2,d_1^!L_1)
\longrightarrow e_*K_E,
\tag{4.2.4'}
\end{equation}
\begin{equation}
\langle\ ,\ \rangle:
\Hom(c_1^*L_1,c_2^!L_2)\otimes\Hom(d_2^*L_2,d_1^!L_1)
\longrightarrow \H^0(E,K_E).
\tag{4.2.5}
\end{equation}
These pairings are compatible with étale localization: given a commutative square (not necessarily cartesian)
\[
\begin{tikzcd}[column sep=2.5em,row sep=1.5em]
& E' \arrow[dl] \arrow[dr] & && & E \arrow[dl] \arrow[dr] & \\
C' \arrow[dr] && D' \arrow[dl] & \text{ étale over } & C \arrow[dr] && D \arrow[dl] \\
& X &&&& X &
\end{tikzcd}
\]
if $c'$, $d'$, $e'$ denote the composites $C'\to C\to X$, $D'\to D\to X$, $E'\to E\to X$, one has a commutative square.


\begin{equation}
\tag{4.2.6}
\begin{tikzcd}[column sep=large,row sep=large]
\Hom(c_1^*L_1,c_2^!L_2)\otimes\Hom(d_2^*L_2,d_1^!L_1)
  \arrow[r,"4.2.5"] \arrow[d] &
\H^0(E,K_E)\arrow[d] \\
\Hom(c_1'{}^*L_1,c_2'{}^!L_2)\otimes\Hom(d_2'{}^*L_2,d_1'{}^!L_1)
  \arrow[r] &
\H^0(E',K_{E'})
\end{tikzcd}
\end{equation}
where the vertical arrows are the restrictions, and the lower horizontal arrow is the composite of the arrow analogous to 4.2.5 and the restriction
\[
\H^0(E'',K_{E''})\longrightarrow \H^0(E',K_{E'}),
\]
with $E''=C'\times_{X'}D'$ (there are analogous commutative squares with arrows of type 4.2.4, 4.2.4', which we leave to the reader to write down). If $Z$ is étale over $E$, we shall denote, for
\[
u\in\Hom(c_1^*L_1,c_2^!L_2),\qquad
v\in\Hom(d_2^*L_2,d_1^!L_1),
\]
\begin{equation}
\tag{4.2.7}
\langle u,v\rangle_Z\in \H^0(Z,K_Z)
\end{equation}
the image of $\langle u,v\rangle$ in $\H^0(Z,K_Z)$ by restriction (in practice, $Z$ will be a connected component of $E$). If $z$ is an isolated closed point of $E$, and $\bar z$ a geometric point above $z$, it follows from 4.2.6 that
\[
\langle u,v\rangle_z\in \H^0(z,K_z)=A
\]
depends only on the situation deduced by strict localization at the images of $\bar z$ in $C,D,X$.

\begin{statement}{Example 4.3}
Assume that $X_i$ is smooth, purely of relative dimension $r_i$, and that $c$ and $d$ are closed immersions, with $c_2$ (respectively $d_1$) finite. For $L_1=A_{X_1}$, $L_2=A_{X_2}$, we have, by 3.1.3,
\[
DL_1\otimes_S^{\mathbf L}L_2=A_X(r_1)[2r_1],\qquad
L_1\otimes_S^{\mathbf L}DL_2=A_X(r_2)[2r_2],
\]
and 4.1.2 is identified with the product
\[
A_X(r_1)[2r_1]\otimes^{\mathbf L}A_X(r_2)[2r_2]\longrightarrow A_X(r)[2r],
\]
where $r=r_1+r_2$, and hence, taking 3.2.1 into account, 4.2.5 is identified with the usual cup product
\[
\H_C^{2r_1}(X,A(r_1))\otimes \H_D^{2r_2}(X,A(r_2))
\longrightarrow
\H_{C\cap D}^{2r}(X,A(r)).
\]
Take as cohomological correspondence from $X_1$ to $X_2$ (respectively from $X_2$ to $X_1$) the class of $c$ (respectively $d$) (3.2 a)). It follows from (SGA $4^{1/2}$ Cycle 2.3.8) that, at an isolated point $z$ of $C\cap D$, $\langle \cl(c),\cl(d)\rangle_z$ is nothing but the image in $A$ of the intersection multiplicity of $C$ and $D$ at $z$.
\end{statement}

We now give the main result of the exposé.

\begin{statement}{Theorem 4.4}
Let, for $i=1,2$, $f_i:X_i\longrightarrow Y_i$ be a proper $S$-morphism,
\[
f=f_1\times_S f_2:X\longrightarrow Y,\qquad
p_i:X\longrightarrow X_i,\qquad q_i:Y\longrightarrow Y_i
\]
the projections. Let moreover
\begin{equation}
\tag{4.4.0}
\begin{tikzpicture}[>=latex,baseline=(current bounding box.center),scale=0.98]
\node (Cp) at (0,2.1) {$C'$};
\node (C) at (3.0,2.1) {$C$};
\node (X) at (-1.25,1.0) {$X$};
\node (Cpp) at (1.45,1.0) {$C''$};
\node (Dp) at (0,0) {$D'$};
\node (D) at (3.0,0) {$D$};
\node (Y) at (-1.25,-1.15) {$Y$};
\node (Dpp) at (1.45,-1.15) {$D''$};
\draw[->] (C) -- (Cp);
\draw[->] (Cp) -- node[left,pos=.55] {$c'$} (X);
\draw[->] (C) -- node[above,sloped,pos=.45] {$c$} (X);
\draw[->] (C) -- (Cpp);
\draw[->] (Cpp) -- node[below,pos=.45] {$c''$} (X);
\draw[->] (Cp) -- node[right] {$f'$} (Dp);
\draw[->] (Cpp) -- node[right] {$f''$} (Dpp);
\draw[->] (C) -- node[right] {$g$} (D);
\draw[->] (X) -- node[left] {$f$} (Y);
\draw[->] (D) -- (Dp);
\draw[->] (Dp) -- node[left,pos=.55] {$d'$} (Y);
\draw[->] (D) -- node[above,sloped,pos=.45] {$d$} (Y);
\draw[->] (D) -- (Dpp);
\draw[->] (Dpp) -- node[below] {$d''$} (Y);
\end{tikzpicture}
\end{equation}
be a commutative cube with cartesian horizontal faces, with $c'$, $c''$, $d'$, $d''$ proper. Put
\[
c_i'=p_i c',\qquad c_i''=p_i c'',\qquad d_i'=q_i d',\qquad d_i''=q_i d''.
\]
Finally, let $L_i\in\ob \Dctf(X_i)$. Then there is a commutative square
\begin{equation}
\tag{4.4.1}
\begin{tikzcd}[column sep=huge,row sep=huge]
f_*c'_*\RHom(c_1'{}^*L_1,c_2'{}^!L_2)\otimes^{\mathbf L}
f_*c''_*\RHom(c_2''{}^*L_2,c_1''{}^!L_1)
 \arrow[r,"(1)"] \arrow[d,"(2)"'] &
f_*c_*K_C \arrow[d,"(4)"] \\
d'_*\RHom(d_1'{}^*M_1,d_2'{}^!M_2)\otimes^{\mathbf L}
d''_*\RHom(d_2''{}^*M_2,d_1''{}^!M_1)
 \arrow[r,"(3)"'] &
d_*K_D ,
\end{tikzcd}
\end{equation}
where $M_i=f_{i*}L_i$, $(3)$ is given by 4.2.4, $(1)$ is deduced from 4.2.4 by applying $f_*$ and composing with the canonical arrow
\[
f_*c'_*(-)\otimes^{\mathbf L} f_*c''_*(-)\longrightarrow f_*(c'_*(-)\otimes^{\mathbf L}c''_*(-)),
\]
$(2)$ is the tensor product of the arrows 3.7.5 relative to the faces containing $f$, and $(4)$ is defined by the adjunction arrow $g_*K_C=g_!g^!K_D\longrightarrow K_D$.
\end{statement}

\begin{prooftext}[Proof]
Put
\[
DL_1\otimes_S^{\mathbf L}L_2=P,\qquad
L_1\otimes_S^{\mathbf L}DL_2=Q,\qquad
DM_1\otimes_S^{\mathbf L}M_2=E,\qquad
M_1\otimes_S^{\mathbf L}DM_2=F.
\]
Consider the diagram
\begin{equation}
\tag{4.4.2}
{\scriptsize
\begin{tikzcd}[column sep=2.6em,row sep=2.6em]
f_*c'_*c'{}^!P\otimes^{\mathbf L} f_*c''_*c''{}^!Q
  \arrow[rr,"(1)"] \arrow[d,"(8)"'] &&
f_*c_*c^!\bigl(P\otimes^{\mathbf L}Q\bigr)
  \arrow[r,"(2)"] \arrow[d,"(9)"'] &
f_*c_*c^!K_X \arrow[d,"(10)"]\\
d'_*d'{}^!f_*P\otimes^{\mathbf L}d''_*d''{}^!f_*Q
  \arrow[r,"(3)"'] \arrow[d,"(11)"'] &
d_*d^!\bigl(f_*P\otimes^{\mathbf L}f_*Q\bigr)
  \arrow[r,"(4)"'] \arrow[d,"(12)"'] &
d_*d^!f_*\bigl(P\otimes^{\mathbf L}Q\bigr)
  \arrow[r,"(5)"'] &
d_*d^!f_*K_X \arrow[d,"(13)"]\\
d'_*d'{}^!E\otimes^{\mathbf L}d''_*d''{}^!F
  \arrow[r,"(6)"'] &
d_*d^!\bigl(E\otimes^{\mathbf L}F\bigr)
  \arrow[rr,"(7)"'] &&
d_*d^!K_Y,
\arrow[from=1-1,to=2-2, phantom, "\scriptstyle (A)"]
\arrow[from=1-3,to=2-4, phantom, "\scriptstyle (B)"]
\arrow[from=2-1,to=3-2, phantom, "\scriptstyle (C)"]
\arrow[from=2-2,to=3-4, phantom, "\scriptstyle (D)"]
\end{tikzcd}
}
\end{equation}
The arrows are as follows: $(1)$ is deduced from 4.2.2 by applying $f_*$; $(2)$ is defined by 4.1.2; $(3)$ is given by 4.2.2; $(4)$ is given by the canonical arrow $f_*P\otimes^{\mathbf L}f_*Q\longrightarrow f_*(P\otimes^{\mathbf L}Q)$; $(5)$ by 4.1.2; $(6)$ by 4.2.2; $(7)$ by 4.1.2; $(8)$ is the tensor product of arrows 3.7.3; $(9)$ and $(10)$ are defined by 3.7.3; $(11)$ and $(12)$ by the isomorphisms $f_*P\longrightarrow E$, $f_*Q\longrightarrow F$ (3.3.1); and $(13)$ by the adjunction arrow $f_*K_X\longrightarrow K_Y$. By definition, 4.4.1 is identified, via 3.1.1 and 3.2.1, with the outer boundary of 4.4.2. It remains only to prove that the squares $(A)$, $(B)$, $(C)$, $(D)$ are commutative. The commutativity of $(B)$ and $(C)$ is clear by functoriality. It remains to prove that of $(A)$ and $(D)$.

a) Commutativity of $(D)$. It is enough to prove that the square
\begin{equation}
\tag{4.4.3}
\begin{tikzcd}[column sep=large,row sep=large]
f_*P\otimes^{\mathbf L}f_*Q \arrow[r] \arrow[d] & f_*K_X \arrow[d]\\
E\otimes^{\mathbf L}F \arrow[r] & K_Y
\end{tikzcd}
\end{equation}
(corresponding to $(D)$ for $d=\Id$) commutes. Recall that the pairing
\[
P\otimes^{\mathbf L}Q\longrightarrow K_X
\quad(\text{respectively }E\otimes^{\mathbf L}F\longrightarrow K_Y)
\]
is the external tensor product of the evaluation arrows
\[
DL_i\otimes^{\mathbf L}L_i\longrightarrow K_{X_i}
\quad(\text{respectively }DM_i\otimes^{\mathbf L}M_i\longrightarrow K_{Y_i}).
\]
By the Künneth isomorphisms
\[
f_{1*}DL_1\otimes_S^{\mathbf L} f_{2*}L_2 \lisom f_*P,\qquad
f_{1*}L_1\otimes_S^{\mathbf L} f_{2*}DL_2 \lisom f_*Q,
\]
4.4.3 is therefore identified with the external tensor product of the squares
\begin{equation}
\tag{4.4.4}
\begin{tikzcd}[column sep=large,row sep=large]
f_{i*}DL_i\otimes^{\mathbf L} f_{i*}L_i
 \arrow[r,"(1)"] \arrow[d,"(3)"'] &
f_{i*}K_{X_i} \arrow[d,"(4)"]\\
DM_i\otimes^{\mathbf L}M_i \arrow[r,"(2)"'] & K_{Y_i},
\end{tikzcd}
\end{equation}
where $(2)$ is the evaluation arrow, $(1)$ the composite of the canonical arrow
\[
f_{i*}DL_i\otimes^{\mathbf L} f_{i*}L_i\longrightarrow f_{i*}(DL_i\otimes^{\mathbf L}L_i)
\]
and the arrow defined by the evaluation arrow, $(3)$ is the tensor product of the duality isomorphism $f_{i*}DL_i\longrightarrow DM_i$ and the identity arrow of $M_i$, and $(4)$ is the adjunction arrow. It is therefore enough to verify the commutativity of the squares 4.4.4. Since the above duality isomorphism is the composite of the canonical arrow (SGA 4 XVIII 3.1.9.5)
\[
f_{i*}\RHom(L_i,K_{X_i})\longrightarrow \RHom(f_{i*}L_i,f_{i*}K_{X_i})
\]
and the arrow deduced from the adjunction arrow
\[
f_{i*}K_{X_i}\longrightarrow K_{Y_i},
\]
it is enough to see that the following square commutes:
\begin{equation}
\tag{4.4.5}
\begin{tikzcd}[column sep=large,row sep=large]
f_{i*}\RHom(L_i,K_{X_i})\otimes^{\mathbf L} f_{i*}L_i
 \arrow[r,"(1)"] \arrow[d,"(3)"'] &
f_{i*}K_{X_i} \arrow[d,"\Id"]\\
\RHom(f_{i*}L_i,f_{i*}K_{X_i})\otimes^{\mathbf L} f_{i*}L_i
 \arrow[r,"(2)"'] &
f_{i*}K_{X_i},
\end{tikzcd}
\end{equation}
where $(1)$ is the arrow $(1)$ of 4.4.4, $(2)$ the evaluation arrow, and $(3)$ the tensor product by $f_{i*}L_i$ of the canonical arrow just mentioned. More generally, for any arrow $N_i\longrightarrow f_{i*}L_i$ in $\Dctf(X_i)$, i.e. any arrow $L_i\xrightarrow{u}N_i$ of type $1$ over $f_i$, the triangle
\[
\begin{tikzcd}[column sep=large,row sep=large]
f_{i*}\RHom(L_i,K_{X_i})\otimes^{\mathbf L}N_i
 \arrow[dr] \arrow[d] & \\
\RHom(N_i,f_{i*}K_{X_i})\otimes^{\mathbf L}N_i \arrow[r] &
f_{i*}K_{X_i}
\end{tikzcd}
\]
deduced from 4.4.5 commutes: indeed, by functoriality, one may assume that $L_i=f_i^*N_i$, with $u$ given by the identity of $L_i$; the arrow
\[
f_{i*}\RHom(L_i,K_{X_i})\longrightarrow\RHom(N_i,f_{i*}K_{X_i})
\]
is then the inverse of the trivial duality isomorphism, and the commutativity of the triangle is verified immediately at the level of complexes, resolving $N_i$ flatly and $K_{X_i}$ injectively.

b) Commutativity of $(A)$. Referring to the definition of 3.7.3, one sees that it is enough to verify this commutativity in each of the following cases:
\[
\text{(i) } X_i=Y_i,\quad f_i=\text{the identity arrow of }X_i;\qquad
\text{(ii) the lateral faces of 4.4.0 are cartesian.}
\]
Place ourselves in case (i). The square $(A)$ is the outer boundary of the diagram
\[
\begin{tikzcd}[column sep=large,row sep=large]
c'_*c'{}^!P\otimes^{\mathbf L}c''_*c''{}^!Q
 \arrow[r,"\sim"] \arrow[d,"(8)"'] &
c_*\bigl(c'{}^!P\otimes_X^{\mathbf L}c''{}^!Q\bigr)
 \arrow[r] \arrow[d] &
c_*c^!(P\otimes^{\mathbf L}Q)\arrow[d,"(9)"]\\
d'_*d'{}^!P\otimes^{\mathbf L}d''_*d''{}^!Q
 \arrow[r,"\sim"'] &
d_*\bigl(d'{}^!P\otimes_X^{\mathbf L}d''{}^!Q\bigr)
 \arrow[r] &
d_*d^!(P\otimes Q),
\arrow[from=1-1,to=2-2,phantom,"\scriptstyle A_1"]
\arrow[from=1-2,to=2-3,phantom,"\scriptstyle A_2"]
\end{tikzcd}
\]
in which the horizontal isomorphisms of $A_1$ are Künneth isomorphisms, the horizontal arrows of $A_2$ are given by 4.2.1, and the middle vertical arrow is deduced by applying $d_*$ to the arrow
\[
g_!(c'{}^!P\otimes_X^{\mathbf L}c''{}^!Q)\longrightarrow d'{}^!P\otimes_X^{\mathbf L}d''{}^!Q
\]
defined, via Künneth, by the external tensor product of the adjunction arrows
\[
f'_!c'{}^!P\longrightarrow d'{}^!P,\qquad
f''_!c''{}^!Q\longrightarrow d''{}^!Q.
\]
These adjunction arrows and the identity arrows of $d'_!d'{}^!P$, $d''_!d''{}^!Q$ define commutative triangles of arrows of type $3$ (1.4)
\[
\begin{tikzcd}[column sep=large,row sep=large]
& c'{}^!P \arrow[d] \arrow[dl] & & c''{}^!Q \arrow[d]\arrow[dl]\\
d'_!d'{}^!P & d'{}^!P\arrow[l] & d''_!d''{}^!Q & d''{}^!Q\arrow[l]
\end{tikzcd}
\]
over
\[
\begin{tikzcd}[column sep=large,row sep=large]
& C' \arrow[d,"f'"] \arrow[dl,"c'"']\\
X & D'\arrow[l,"d'"]
\end{tikzcd}
\qquad
\begin{tikzcd}[column sep=large,row sep=large]
& C'' \arrow[d,"f''"] \arrow[dl,"c''"']\\
X & D''\arrow[l,"d''"]
\end{tikzcd}.
\]
By external tensor product (over $X$), these triangles give a commutative triangle
\[
\begin{tikzcd}[column sep=huge,row sep=large]
& c'{}^!P\otimes_X^{\mathbf L}c''{}^!Q \arrow[d] \arrow[dl]\\
d'_!d'{}^!P\otimes d''_!d''{}^!Q &
d'{}^!P\otimes_X^{\mathbf L}d''{}^!Q\arrow[l]
\end{tikzcd}
\qquad\text{over}\qquad
\begin{tikzcd}[column sep=large,row sep=large]
& C\arrow[d,"g"]\arrow[dl,"c"']\\
X & D\arrow[l,"d"]
\end{tikzcd}.
\]
The commutativity of this triangle is equivalent to that of the square deduced from $A_1$ by inverting the horizontal arrows. On the other hand, there are commutative triangles of arrows of type $3$
\[
\begin{tikzcd}[column sep=large,row sep=large]
& c'{}^!P \arrow[d]\arrow[dl]\\
P & d'{}^!P\arrow[l]
\end{tikzcd}
\qquad
\begin{tikzcd}[column sep=large,row sep=large]
& c''{}^!Q \arrow[d]\arrow[dl]\\
Q & d''{}^!Q\arrow[l]
\end{tikzcd}
\]
defined by the identity arrows of $d'{}^!P$, $c'{}^!P$, $d''{}^!Q$, $c''{}^!Q$. By external tensor product, one obtains a commutative triangle
\[
\begin{tikzcd}[column sep=huge,row sep=large]
& c'{}^!P\otimes_X^{\mathbf L}c''{}^!Q \arrow[d]\arrow[dl]\\
P\otimes^{\mathbf L}Q & d'{}^!P\otimes_X^{\mathbf L}d''{}^!Q\arrow[l]
\end{tikzcd}
\]
which can be interpreted as a morphism of arrows of type $3$ over $g:C\longrightarrow D$:
\[
\begin{tikzcd}[column sep=large,row sep=large]
c'{}^!P\otimes_X^{\mathbf L}c''{}^!Q \arrow[r]\arrow[d] &
c'{}^!(P\otimes^{\mathbf L}Q)\arrow[d]\\
d'{}^!P\otimes_X^{\mathbf L}d''{}^!Q \arrow[r] &
d'{}^!(P\otimes Q),
\end{tikzcd}
\]
or again as a commutative square over $D$:
\[
\begin{tikzcd}[column sep=large,row sep=large]
g_*(c'{}^!P\otimes_X^{\mathbf L}c''{}^!Q) \arrow[r]\arrow[d] &
g_*c'{}^!(P\otimes^{\mathbf L}Q)\arrow[d]\\
d'{}^!P\otimes_X^{\mathbf L}d''{}^!Q \arrow[r] &
d^!(P\otimes^{\mathbf L}Q).
\end{tikzcd}
\]
But the square deduced from the latter by applying $d_*$ is none other than $A_2$, whence the commutativity of $(A)$ in case (i).

It remains to treat case (ii). Since the faces of 4.4.0 containing $f$ are cartesian, the arrows
\[
f'_*c'{}^!P\longrightarrow d'{}^!f_*P,\qquad
f''_*c''{}^!Q\longrightarrow d''{}^!f_*Q
\]
(3.7.3) are isomorphisms. Their inverses define arrows of type $1$
\[
c'{}^!P\longrightarrow d'{}^!f_*P,\qquad
c''{}^!Q\longrightarrow d''{}^!f_*Q
\]
over $f'$, $f''$, whence, by direct image, commutative squares of arrows of type $1$
\[
\begin{tikzcd}[column sep=large,row sep=large]
c'_*c'{}^!P \arrow[d] & c'{}^!P\arrow[l]\arrow[d]\\
d'_*d'{}^!P & d'{}^!f_*P\arrow[l]
\end{tikzcd}
\qquad
\begin{tikzcd}[column sep=large,row sep=large]
c''_*c''{}^!Q \arrow[d] & c''{}^!Q\arrow[l]\arrow[d]\\
d''_*d''{}^!Q & d''{}^!f_*Q\arrow[l]
\end{tikzcd}
\]
over the faces of 4.4.0 containing $f$. Taking the tensor product, one obtains a commutative square in $\Dctf(Y)$
\begin{equation}
\tag{4.4.6}
{\scriptsize
\begin{tikzcd}[column sep=large,row sep=large]
d'_*d'{}^!f_*P\otimes^{\mathbf L}d''_*d''{}^!f_*Q
 \arrow[r]\arrow[d] &
d_*\bigl(d'{}^!f_*P\otimes_Y^{\mathbf L}d''{}^!f_*Q\bigr)\arrow[d]\\
f_*c'_*c'{}^!P\otimes^{\mathbf L}f_*c''_*c''{}^!Q
 \arrow[r] &
d_*g_*\bigl(c'{}^!P\otimes_X^{\mathbf L}c''{}^!Q\bigr),
\end{tikzcd}
}
\end{equation}
where the left vertical arrow is the inverse of $(8)$, the horizontal arrows are given by Künneth, and the right vertical arrow by the external tensor product over $Y$ of the arrows
\[
f'_*c'{}^!P\longrightarrow d'{}^!f_*P,\qquad
f''_*c''{}^!Q\longrightarrow d''{}^!f_*Q.
\]
The same tensor product defines the left vertical arrow of the square
\begin{equation}
\tag{4.4.7}
\begin{tikzcd}[column sep=large,row sep=large]
d'{}^!f_*P\otimes_Y^{\mathbf L}d''{}^!f_*Q
 \arrow[r]\arrow[d] &
d^!f_*(P\otimes^{\mathbf L}Q)\arrow[d]\\
g_*(c'{}^!P\otimes_X^{\mathbf L}c''{}^!Q)
 \arrow[r] &
g_*c'{}^!(P\otimes^{\mathbf L}Q),
\end{tikzcd}
\end{equation}
where the right vertical arrow is 3.7.3, and the horizontal arrows are defined by 4.2.1. The outer boundary of the diagram composed of 4.4.6 and $d_*(4.4.7)$ is the square deduced from $(A)$ by inverting the vertical arrows. It is therefore enough to prove that 4.4.7 commutes, in other words that the arrows 4.2.1
\[
c'{}^!P\otimes_X^{\mathbf L}c''{}^!Q\longrightarrow c'{}^!(P\otimes^{\mathbf L}Q),\qquad
d'{}^!f_*P\otimes_Y^{\mathbf L}d''{}^!f_*Q\longrightarrow d^!f_*(P\otimes^{\mathbf L}Q)
\]
define a morphism from
\[
c'{}^!P\otimes_X^{\mathbf L}c''{}^!Q\longrightarrow d'{}^!f_*P\otimes_Y^{\mathbf L}d''{}^!f_*Q
\]
to
\[
c'{}^!(P\otimes^{\mathbf L}Q)\longrightarrow d^!f_*(P\otimes^{\mathbf L}Q),
\]
viewed as arrows of type $1$ over $g$. By adjunction, we are reduced to verifying that the tensor products of adjunction arrows
\[
c'_*c'{}^!P\otimes_X^{\mathbf L}c''_*c''{}^!Q\longrightarrow P\otimes^{\mathbf L}Q,
\qquad
d'_*d'{}^!f_*P\otimes_Y^{\mathbf L}d''_*d''{}^!f_*Q
\longrightarrow f_*P\otimes^{\mathbf L}f_*Q
\]
define a morphism, between arrows of type $1$ over $f$, from the tensor product of the arrows
\[
c'_*c'{}^!P\longrightarrow d'_*d'{}^!f_*P,\qquad
c''_*c''{}^!Q\longrightarrow d''_*d''{}^!f_*Q
\]
to the tensor product of the arrows
\[
P\longrightarrow f_*P,\qquad Q\longrightarrow f_*Q.
\]
For this it is enough to verify that the adjunction arrows
\[
c'_*c'{}^!P\longrightarrow P,\qquad d'_*d'{}^!f_*P\longrightarrow f_*P
\]
define a morphism from
\[
c'_*c'{}^!P\longrightarrow d'_*d'{}^!f_*P
\]
to
\[
P\longrightarrow f_*P,
\]
and likewise for the pair $(c'',d'')$. It is enough to consider the case of $(c',d')$, the statement being the same in the two cases. We have to see that the triangle
\[
\begin{tikzcd}[column sep=large,row sep=large]
& f_*c'_*c'{}^!P \arrow[d,"\sim"]\arrow[dl]\\
f_*P & d'_*d'{}^!f_*P\arrow[l]
\end{tikzcd}
\]
where the horizontal and oblique arrows are given by adjunction and the vertical arrow by the isomorphism 3.7.3, is commutative. Taking into account the definition of 3.7.3 (SGA 4 XVIII 3.1.12.3) as the transpose of a base-change isomorphism, we are reduced, by applying the functor $\Hom(R,-)$ ($R\in\ob\Dctf(Y)$), to verifying the commutativity of the square
\[
\begin{tikzcd}[column sep=large,row sep=large]
c'_*c'{}^*f^*R \arrow[r,"\Id"] &
c'_*f'{}^*d'{}^*R\\
f^*R \arrow[u] \arrow[r] &
f^*d'_*d'{}^*R\arrow[u,"\sim"']
\end{tikzcd}
\]
where the left vertical arrow (respectively the lower horizontal arrow) is defined by the adjunction arrow $1\longrightarrow c'_*c'{}^*$ (respectively $1\longrightarrow d'_*d'{}^*$), and the right vertical arrow by the base-change isomorphism
\[
f^*d'_*\lisom c'_*f'{}^*.
\]
But this square commutes by the very definition of the base-change arrow. The commutativity of $(A)$ is therefore established, and this completes the proof of the theorem.
\end{prooftext}

\begin{statement}{Corollary 4.5}
Under the hypotheses of 4.4, there is a commutative square
\[
\begin{tikzcd}[column sep=large,row sep=large]
\Hom(c_1'{}^*L_1,c_2'{}^!L_2)\otimes
\Hom(c_2''{}^*L_2,c_1''{}^!L_1)
 \arrow[r,"4.2.5"] \arrow[d,"3.7.6\otimes 3.7.6"'] &
\H^0(C,K_C)\arrow[d,"g_*"]\\
\Hom(d_1'{}^*f_{1*}L_1,d_2'{}^!f_{2*}L_2)\otimes
\Hom(d_2''{}^*f_{2*}L_2,d_1''{}^!f_{1*}L_1)
 \arrow[r,"4.2.5"'] &
\H^0(D,K_D),
\end{tikzcd}
\]
where $g_*$ is defined by the adjunction arrow
\[
g_*K_C=g_*g^!K_D\longrightarrow K_D.
\]
In other words, for
\[
u\in\Hom(c_1'{}^*L_1,c_2'{}^!L_2),\qquad
v\in\Hom(c_2''{}^*L_2,c_1''{}^!L_1),
\]
one has
\begin{equation}
\tag{4.5.1}
\langle f_*u,f_*v\rangle=g_*\langle u,v\rangle.
\end{equation}
\end{statement}

For every proper $S$-scheme $t:T\longrightarrow S$, denote by
\begin{equation}
\tag{4.6}
\int_T:\H^0(T,K_T)\longrightarrow \H^0(S,K_S)=\H^0(S,A)
\end{equation}
the arrow defined by the adjunction arrow
\[
t_!K_T\longrightarrow A.
\]

\begin{statement}{Corollary 4.7 -- Lefschetz--Verdier formula}
Under the hypotheses of 4.4, suppose that $Y_1=Y_2=S=D'=D''$. For
\[
u\in\Hom(c_1'{}^*L_1,c_2'{}^!L_2),\qquad
v\in\Hom(c_2''{}^*L_2,c_1''{}^!L_1),
\]
one has
\begin{equation}
\tag{4.7.1}
\langle u_*,v_*\rangle=\int_C\langle u,v\rangle,
\end{equation}
where
\[
u_*\in\Hom(f_{1*}L_1,f_{2*}L_2),\qquad
v_*\in\Hom(f_{2*}L_2,f_{1*}L_1)
\]
are the homomorphisms deduced from $u$, $v$ by direct image over $S$ (3.7.6), and
\[
\langle u_*,v_*\rangle=\Tr(u_*v_*)=\Tr(v_*u_*)
\]
is the cup product, in the sense of (SGA 6 I 8.3), of the homomorphisms of perfect complexes $u_*$, $v_*$ (4.1.5).
\end{statement}

Taking 3.6.5 into account, in particular one has:

\begin{statement}{Corollary 4.8}
Let $t:T\longrightarrow S$ be a proper $S$-scheme, $\Delta\subset T\times_S T$ the diagonal, $F:T\longrightarrow T$ an $S$-endomorphism of $T$, $T^F=(\text{graph of }F)\times_{(T\times_S T)}\Delta$ the scheme of fixed points of $F$, $L\in\ob\Dctf(T)$, $u\in\Hom(F^*L,L)$ a correspondence from $L$ to $L$ with support in the graph of $F$, and $1\in\Hom(L,L)$ the identity correspondence, with support in $\Delta$ (3.2 b)). Then
\[
\Tr(u_*)=\int_{T^F}\langle u,1\rangle,
\]
where $u_*$ is the endomorphism of $t_*L$ obtained by composing $t_*L\longrightarrow t_*F^*L$ (inverse image) with $t_*u:t_*F^*L\longrightarrow t_*L$.
\end{statement}

From the remarks of 3.2 a) and 3.6 one also deduces:

\begin{statement}{Corollary 4.9}
Under the hypotheses of 4.7, assume that $X_i$ is smooth over $S$, purely of relative dimension $r_i$, and that $C'$ (respectively $C''$) is a closed subscheme of $X$ finite and of finite tor-dimension over $X_2$ (respectively $X_1$). Let
\[
\cl(C')\in \H_{C'}^{2r_1}(X,A(r_1)),\qquad
\cl(C'')\in \H_{C''}^{2r_2}(X,A(r_2))
\]
be the cohomology classes of $C'$ and $C''$, and let
\[
\cl(C')_*\in\Hom(\R f_{1*}A,\R f_{2*}A),\qquad
\cl(C'')_*\in\Hom(\R f_{2*}A,\R f_{1*}A)
\]
be the corresponding homomorphisms (3.2 a), 3.7.6). Then
\[
\langle \cl(C')_*,\cl(C'')_*\rangle
=
\int_{C'\cap C''}\cl(C)\cl(D),
\]
where
\[
\cl(C)\cl(D)\in
\H_{C'\cap C''}^{2(r_1+r_2)}(X,A(r_1+r_2))
=
\H^0(C'\cap C'',K_{C'\cap C''})
\]
is the usual cup product.
\end{statement}

Taking 4.3 into account, in particular one has:

\begin{statement}{Corollary 4.10}
Under the hypotheses of 4.9, suppose that $C'\cap C''$ is finite. For $z\in C'\cap C''$, denote by $(C'\!\cdot C'')_z$ the image in $A$ of the intersection multiplicity of $C'$ and $C''$ at $z$. Then
\[
\langle \cl(C')_*,\cl(C'')_*\rangle
=
\sum_{z\in C'\cap C''}(C'\!\cdot C'')_z.
\]
\end{statement}

\begin{statement}{4.11}
Under the hypotheses of 4.9, one has in general
\begin{equation}
\tag{4.11.1}
\langle \cl(C')_*,\cl(C'')_*\rangle
=
\sum_{Z\in\pi_0(C'\cap C'')}(\cl(C')\cl(C''))_Z,
\end{equation}
where $\pi_0(C'\cap C'')$ denotes the set of connected components of $C'\cap C''$, and $(\cl(C')\cl(C''))_Z$ the restriction of $\cl(C')\cl(C'')$ to $\H^0(Z,K_Z)$. When $Z$ is a component of dimension $\geq 1$, the calculation of $(\cl(C')\cl(C''))_Z$ presents problems, because one is in an ``excess'' situation. If $Z$ is smooth, purely of dimension $d$, if $C'$ and $C''$ are smooth in a neighbourhood of $Z$, and if $X$ is projective, one can show that
\begin{equation}
\tag{4.11.2}
\int_Z(\cl(C')\cl(C''))_Z=\int_Z c_d(N_Z),
\end{equation}
where $N_Z$ is the locally free sheaf of rank $d$ defined by the exact sequence of normal bundles
\[
0\longrightarrow N_{Z/C'}\oplus N_{Z/C''}\longrightarrow N_{Z/X}\longrightarrow N_Z\longrightarrow 0,
\]
and $c_d(N_Z)\in \H^{2d}(Z,A(d))$ denotes its $d$-th Chern class (SGA 5 VII) (use (SGA 6 XIV 6.3) to reduce to a calculation of intersection number in K-theory, and apply (SGA 6 VII 2.5), cf. ([10] 4.1); the projectivity hypothesis is no doubt unnecessary).

Let $T$ be a smooth projective $S$-scheme; take as correspondence $C'$ (respectively $C''$) the graph of an $S$-endomorphism $F$ (respectively the diagonal), and suppose that the scheme of fixed points $T^F=C'\cap C''$ is smooth. The sheaf $N_{(T^F)}$ above is nothing other than the tangent sheaf to $T^F$, so that 4.11.2 implies, by the Gauss--Bonnet formula (SGA 5 VII 4.9), the analogue of a formula well known to topologists:
\begin{equation}
\tag{4.11.3}
\Tr(F,\R\Gamma(T,A))=\EP(T^F),
\end{equation}
where $\EP(T^F)$ denotes the Euler--Poincaré characteristic of $T_s^F$, for $s$ a geometric point above $S$ (loc. cit.).
\end{statement}

The calculation of $(\cl(C')\cl(C''))_Z$ for a non-smooth component $Z$ of dimension $\geq 1$ has not been treated in general, at least to the knowledge of the editor.

\begin{statement}{4.12}
In the situation of 4.7, let $z$ be an isolated point of $C$, with projection $z_i$ on $X_i$, $a'$ (respectively $a''$) on $C'$ (respectively $C''$), and let $\bar z$ be a geometric point above $z$, with image $\bar z_i$ above $z_i$. Assume that $c_2'$ (respectively $c_1''$) is quasi-finite and flat at $a'$ (respectively $a''$), that $L_i$ is perfect in a neighbourhood of $\bar z_i$, and that $u$ (respectively $v$) is given in a neighbourhood of $a'$ (respectively $a''$) by
\[
\Tr_{c_2'}:c_{2!}'c_2'{}^*L_2\longrightarrow L_2
\quad\text{and}\quad
s\in\Hom(c_1'{}^*L_1,c_2'{}^*L_2)
\]
(respectively
\[
\Tr_{c_1''}:c_{1!}''c_1''{}^*L_1\longrightarrow L_1
\quad\text{and}\quad
t\in\Hom(c_2''{}^*L_2,c_1''{}^*L_1)).
\]
Let $s_{\bar z}\in\Hom(L_{1\bar z_1},L_{2\bar z_2})$ (respectively $t_{\bar z}\in\Hom(L_{2\bar z_2},L_{1\bar z_1})$) be the homomorphism deduced from $s$ (respectively $t$) by passage to the fibres. Let again, as in 3.2 a), $\cl(c')$ (respectively $\cl(c'')$) denote the correspondence from $X_1$ to $X_2$ (respectively from $X_2$ to $X_1$) defined, in a neighbourhood of $a'$ (respectively $a''$), by $c_2'$ (respectively $c_1''$), and let $(C'\!\cdot C'')_z$ be the corresponding local term $\langle\cl(c'),\cl(c'')\rangle_z$ (equal to the image in $A$ of the intersection multiplicity of $C'$ and $C''$ at $z$ when $C'$ and $C''$ are closed subschemes of $X$, with $X_1$ and $X_2$ smooth). Then:
\begin{equation}
\tag{4.12.1}
\langle u,v\rangle_z=(C'\!\cdot C'')_z\langle s_{\bar z},t_{\bar z}\rangle.
\end{equation}
This follows easily from 4.2.6.
\end{statement}

\begin{statement}{4.13}
If $T$ is a scheme over $S$, denote by $DF(T)$ the filtered derived category of complexes of $A_T$-modules with finite filtration ([9] V 1), and by $DF_{\ctf}(T)$ the full subcategory formed by the $L$ such that $\gr(L)\in\ob\Dctf(T)$. The kind of compatibility of filtered derived functors with passage to associated graded objects implies that $DF_{\ctf}$ is stable under the six operations $f^*$, $f_*$, $f_!$, $f^!$, $\otimes^{\mathbf L}$, $\RHom$. The definitions and results of the preceding numbers are valid, mutatis mutandis, in $DF_{\ctf}$. The pairing 4.1.2 is compatible with passage to associated graded objects ($K_X$ being regarded as endowed with the trivial filtration), and it follows that, with the notation of 4.2.5, for $L_i\in\ob DF_{\ctf}(X_i)$,
\[
u\in\Hom_{DF}(c_1^*L_1,c_2^!L_2),\qquad
v\in\Hom_{DF}(d_2^*L_2,d_1^!L_1),
\]
one has the following equality in $\H^0(E,K_E)$:
\begin{equation}
\tag{4.13.1}
\langle u,v\rangle=\sum_i\langle \gr^i u,\gr^i v\rangle.
\end{equation}
As Deligne observes, this remark applies in particular to the situation considered by Langlands in ([13] §7). The correspondences that he considers in (loc. cit. 7.11) indeed preserve, near the fixed points, finite filtrations on the sheaves with associated graded objects inducing, on the strict localization at the image of a fixed point, either a constant sheaf or a constant sheaf on the complement of the closed point extended by zero. The calculation of the local terms is easily reduced, by 4.13, to 4.12.1 and to the case of sheaves concentrated on the closed points that are the images of the fixed points, a case treated directly. Applying 4.7, one obtains the Lefschetz formula required in ([13] 7.11). We point out that the general statement (loc. cit. 7.12) is true; the proof, more delicate, will be given in III B.
\end{statement}



\section*{5. Complements}

\subsection*{5.1. Dual of a correspondence}

Let, for $i=1,2$, $X_i$ be an $S$-scheme, let $X=X_1\times_S X_2$, and let $p_i:X\to X_i$ be the projections. Let $L_i\in\operatorname{ob}D_{\mathrm{ctf}}(X_i)$. Every cohomological correspondence $u$ from $L_1$ to $L_2$ gives, by applying $D_X$ and identifying $Dp_2^!$ with $p_2^*D$ and $Dp_1^*$ with $p_1^!D$, a correspondence $Du$ from $DL_2$ to $DL_1$, called the \emph{dual of $u$}. In other words, $D$ defines a homomorphism
\begin{equation}
\tag{5.1.1}
D:\Hom(p_1^*L_1,p_2^!L_2)\longrightarrow \Hom(p_2^*DL_2,p_1^!DL_1).
\end{equation}
One checks that 5.1.1 is obtained, by applying $H^0(X,-)$, from the isomorphism
\begin{equation}
\tag{5.1.2}
D:\RHom(p_1^*L_1,p_2^!L_2)\xrightarrow{\sim}\RHom(p_2^*DL_2,p_1^!DL_1)
\end{equation}
which makes the square
\begin{equation}
\tag{5.1.3}
\begin{tikzcd}[column sep=large,row sep=large]
p_1^*DL_1\otimes^{\mathbf L}p_2^*L_2 \arrow[r,"3.1.1"] \arrow[d]
& \RHom(p_1^*L_1,p_2^!L_2) \arrow[d] \\
p_2^*DDL_2\otimes^{\mathbf L}p_1^*DL_1 \arrow[r,"3.1.1"]
& \RHom(p_2^*DL_2,p_1^!DL_1),
\end{tikzcd}
\end{equation}
commutative. Here the left vertical arrow is the composite of the symmetry isomorphism and the one defined by the biduality isomorphism $L_2\to DDL_2$. It follows easily from the definitions that the following square is commutative:
\begin{equation}
\tag{5.1.4}
\begin{tikzcd}[column sep=large,row sep=large]
\RHom_S(L_1,L_2)\otimes^{\mathbf L}\RHom_S(L_2,L_1) \arrow[r,"4.1.3"] \arrow[d,"D\otimes^{\mathbf L} D"']
& K_X \arrow[d,"\operatorname{Id}"] \\
\RHom_S(DL_2,DL_1)\otimes^{\mathbf L}\RHom_S(DL_1,DL_2) \arrow[r,"4.1.3"]
& K_X.
\end{tikzcd}
\end{equation}
Let $c:C\to X$ be an $S$-morphism and put $c_i=p_ic$. Applying $c_*c^!$ to 5.1.2, and using (3.2.1), one deduces an isomorphism
\begin{equation}
\tag{5.1.5}
D:\Hom(c_1^*L_1,c_2^!L_2)\xrightarrow{\sim}\Hom(c_2^*DL_2,c_1^!DL_1),
\end{equation}
which is checked to identify with the one obtained by applying the functor $D_X$ and the identifications of $Dc_1^*$ with $c_1^!D$ and $Dc_2^!$ with $c_2^*D$. If $c':C'\to X$ and $c'':C''\to X$ are $S$-morphisms, it follows from 5.1.4 that, for
\[
u\in\Hom(c_1'{}^*L_1,c_2'{}^!L_2),\qquad v\in\Hom(c_2''{}^*L_2,c_1''{}^!L_1),\]
one has
\begin{equation}
\tag{5.1.6}
\langle u,v\rangle=\langle Du,Dv\rangle,
\end{equation}
with the notation of 4.2.5.

On the other hand, one checks that, in the situation of 3.3, the following square is commutative:
\begin{equation}
\tag{5.1.7}
\begin{tikzcd}[column sep=large,row sep=large]
f_*\RHom_S(L_1,L_2) \arrow[r,"3.3.1"] \arrow[d,"f_*D"']
& \RHom_S(f_{1!}L_1,f_{2*}L_2) \arrow[d,"D"] \\
f_*\RHom_S(DL_2,DL_1) \arrow[r,"3.3.1"]
& \RHom_S(f_{2!}DL_2,f_{1*}DL_1),
\end{tikzcd}
\end{equation}
and consequently, in the situation of 3.7.6, one has
\begin{equation}
\tag{5.1.8}
f_*Du=Df_*u.
\end{equation}

\subsection*{5.2. Composition of correspondences}

Let, for $i=1,2,3$, $X_i$ be an $S$-scheme, let $X_{ij}=X_i\times_SX_j$, $X=X_1\times_SX_2\times_SX_3$, and let $p_i:X\to X_i$, $p_{ij}:X\to X_{ij}$ be the projections. Let $L_i\in\operatorname{ob}D_{\mathrm{ctf}}(X_i)$. One has $X=X_{12}\times_{X_2}X_{23}$, and
\[
(DL_1\otimes_S^{\mathbf L}L_2)\otimes_{X_2}^{\mathbf L}(DL_2\otimes_S^{\mathbf L}L_3)
=p_1^*DL_1\otimes^{\mathbf L}p_2^*L_2\otimes^{\mathbf L}p_2^*DL_2\otimes^{\mathbf L}p_3^*L_3.
\]
The evaluation arrow $L_2\otimes^{\mathbf L}DL_2\to K_{X_2}$ therefore defines a pairing
\begin{equation}
\tag{5.2.1}
(DL_1\otimes_S^{\mathbf L}L_2)\otimes_{X_2}^{\mathbf L}(DL_2\otimes_S^{\mathbf L}L_3)\longrightarrow p_{13}^!(DL_1\otimes_S^{\mathbf L}L_3),
\end{equation}
the second member being identified with $(DL_1\otimes_S^{\mathbf L}L_3)\otimes_S^{\mathbf L}K_{X_2}$ by Künneth (1.7.3). Taking 3.1.1 and 3.1.2 into account, one may rewrite 5.2.1 in the form
\begin{equation}
\tag{5.2.2}
\RHom_S(L_1,L_2)\otimes_{X_2}^{\mathbf L}\RHom_S(L_2,L_3)\longrightarrow p_{13}^!\RHom_S(L_1,L_3)=\RHom(p_1^*L_1,p_3^!L_3),
\end{equation}
the identification with the second member being given by the induction isomorphism (SGA~4~XVIII 3.1.12.2). Hence one obtains a pairing, called \emph{composition of correspondences},
\begin{equation}
\tag{5.2.3}
\Hom_S(L_1,L_2)\otimes\Hom_S(L_2,L_3)\longrightarrow\Hom(p_1^*L_1,p_3^!L_3),
\end{equation}
which will be denoted $(u,v)\mapsto vu$. One checks that, if $p_{ij,i}:X_{ij}\to X_i$ and $p_{ij,j}:X_{ij}\to X_j$ denote the projections, 5.2.2 can also be obtained by composing the arrow
\[
\RHom_S(L_1,L_2)\otimes_{X_2}^{\mathbf L}\RHom_S(L_2,L_3)
\longrightarrow
\RHom(p_1^*L_1,p_{12}^*p_{12,2}^!L_2)
\otimes
\RHom(p_{23}^!p_{23,2}^*L_2,p_3^!L_3)
\]
(the tensor product of the canonical arrows $p_{12}^*\RHom_S(L_1,L_2)\to\RHom(p_1^*L_1,p_{12}^*p_{12,2}^!L_2)$ and $p_{23}^*\RHom_S(L_2,L_3)\to\RHom(p_{23}^!p_{23,2}^*L_2,p_3^!L_3)$), the arrow defined by the base-change arrow $p_{12}^!p_{12,2}^*L_2\to p_{23}^*p_{23,2}^*L_2$, and a usual composition arrow on the $\RHom$'s.

Let $c':C'\to X_{12}$ and $c'':C''\to X_{23}$ be correspondences, and let $c:C=C'\times_{X_2}C''\to X$ be the ``composite'' correspondence. From 5.2.2 one deduces, or else defines directly by a composition analogous to the one just described, a pairing
\begin{equation}
\tag{5.2.4}
\RHom(c_1'{}^*L_1,c_2'{}^!L_2)\otimes_{X_2}^{\mathbf L}\RHom(c_2''{}^*L_2,c_3''{}^!L_3)
\longrightarrow\RHom(c_1^*L_1,c_3^!L_3),
\end{equation}
hence a ``composition'' pairing $(u,v)\mapsto vu$:
\begin{equation}
\tag{5.2.5}
\Hom(c_1'{}^*L_1,c_2'{}^!L_2)\otimes\Hom(c_2''{}^*L_2,c_3''{}^!L_3)
\longrightarrow\Hom(c_1^*L_1,c_3^!L_3).
\end{equation}

One checks that composition of correspondences is compatible with direct images in the following sense. Let, for $1\leq i\leq 3$, $f_i:X_i\to Y_i$ be an $S$-morphism, with $f_2$ proper, and let $c'$, $c''$, $c$ be morphisms as above, with $c'$ and $c''$ proper. Let $d':D'\to Y_{12}$ and $d'':D''\to Y_{23}$ be proper morphisms, $d:D=D'\times_{Y_2}D''\to Y$, and suppose given the commutative squares
\[
\begin{tikzcd}[column sep=large,row sep=large]
X_{12} \arrow[d,"f_{12}"'] & C' \arrow[l] \arrow[d] \\
Y_{12} & D' \arrow[l]
\end{tikzcd}
\qquad
\begin{tikzcd}[column sep=large,row sep=large]
X_{23} \arrow[d,"f_{23}"'] & C'' \arrow[l] \arrow[d] \\
Y_{23} & D'' \arrow[l]
\end{tikzcd}
\]
Then one has a commutative square
\begin{equation}
\tag{5.2.6}
\begin{tikzcd}[column sep=large,row sep=large]
\Hom(c_1'{}^*L_1,c_2'{}^!L_2)\otimes\Hom(c_2''{}^*L_2,c_3''{}^!L_3) \arrow[r,"5.2.5"] \arrow[d,"f_{12*}\otimes f_{23*}"']
& \Hom(c_1^*L_1,c_3^!L_3) \arrow[d,"f_{13*}"] \\
\Hom(d_1'{}^*M_1,d_2'{}^!M_2)\otimes\Hom(d_2''{}^*M_2,d_3''{}^!M_3) \arrow[r,"5.2.5"]
& \Hom(d_1^*M_1,d_3^!M_3),
\end{tikzcd}
\end{equation}
where $M_1=f_{1!}L_1$, $M_2=f_{2*}L_2$, and $M_3=f_{3*}L_3$.

Finally, suppose that $X_1=X_3$ and $L_1=L_3$. The fundamental pairing 4.1.2 may be interpreted as the composite of the pairing
\begin{equation}
\tag{5.2.7}
q_1^*DL_1\otimes^{\mathbf L}q_2^*L_2\otimes^{\mathbf L}q_2^*DL_2\otimes^{\mathbf L}q_1^*L_1
\longrightarrow
q_1^*DL_1\otimes^{\mathbf L}q_2^*K_{X_2}\otimes^{\mathbf L}q_1^*L_1
\end{equation}
defined by $L_2\otimes^{\mathbf L}DL_2\to K_{X_2}$, or, equivalently, induced by 5.2.1 on $X_{12}$ via the diagonal $X_{12}\to X_1\times_SX_2\times_SX_1$ ($q_i:X_{12}\to X_i$ denoting the projection), and the pairing
\begin{equation}
\tag{5.2.8}
q_1^*DL_1\otimes^{\mathbf L}q_2^*K_{X_2}\otimes^{\mathbf L}q_1^*L_1\longrightarrow K_{X_{12}}
\end{equation}
defined by $DL_1\otimes^{\mathbf L}L_1\to K_{X_1}$. It follows that, for $u\in\Hom_S(L_1,L_2)$ and $v\in\Hom_S(L_2,L_1)$, one has, in $H^0(X_1\times_SX_2,K_{X_1\times_SX_2})$,
\begin{equation}
\tag{5.2.9}
\langle u,v\rangle=\langle vu,1\rangle,
\end{equation}
where $1$ denotes the identity correspondence from $L_1$ to $L_1$ on $X_1\times_SX_1$, supported on the diagonal (3.2 b)), and the cup product in the second member is taken in the sense of 4.2.5, taking account of the fact that the fibred product $X\times_{(X_1\times_SX_1)}X_1$ (via $p_{13}:X\to X_{13}=X_1\times_SX_1$ and the diagonal $X_1\to X_1\times_SX_1$) identifies with $X_1\times_SX_2$. Likewise, more generally, if $c':C'\to X_{12}$ and $c'':C''\to X_{21}$ are correspondences, if $c:C=C'\times_{X_2}C''\to X_1\times_SX_2\times_SX_1$ is the composite correspondence, and if $d:C'\times_{X_{12}}C''=D\to X_{12}$ is the intersection correspondence, then, for $u\in\Hom(c_1'{}^*L_1,c_2'{}^!L_2)$ and $v\in\Hom(c_2''{}^*L_2,c_1''{}^!L_1)$, one has the following equality in $H^0(D,K_D)$:
\begin{equation}
\tag{5.2.10}
\langle u,v\rangle=\langle vu,1\rangle,
\end{equation}
where $1$ denotes, as above, the diagonal correspondence, and $D$ is identified with the fibred product of $C\to X_1\times_SX_1$ and the diagonal $X_1\to X_1\times_SX_1$. Formulas 5.2.9 and 5.2.10 generalize a well-known formula in the case of constant coefficients ([12] 1.3.6 (ii) c)).

\subsection*{5.3. External tensor product of correspondences}

Let, for $1\leq i\leq 4$, $X_i$ be an $S$-scheme, let $X=X_1\times_SX_2\times_SX_3\times_SX_4$, $X_{ij}=X_i\times_SX_j$, and let $L_i\in\operatorname{ob}D_{\mathrm{ctf}}(X_i)$. The Künneth isomorphism (1.7.6, 2.2.4)
\[
DL_1\otimes_S^{\mathbf L}DL_3\xrightarrow{\sim}D(L_1\otimes_S^{\mathbf L}L_3)
\]
defines an isomorphism on $X$
\begin{equation}
\tag{5.3.1}
(DL_1\otimes_S^{\mathbf L}L_2)\otimes_S^{\mathbf L}(DL_3\otimes_S^{\mathbf L}L_4)
\xrightarrow{\sim}D(L_1\otimes_S^{\mathbf L}L_3)\otimes_S^{\mathbf L}(L_2\otimes_S^{\mathbf L}L_4),
\end{equation}
which, by 3.1.1, may be rewritten in the form
\begin{equation}
\tag{5.3.2}
\RHom_S(L_1,L_2)\otimes_S^{\mathbf L}\RHom_S(L_3,L_4)
\xrightarrow{\sim}\RHom_S(L_1\otimes_S^{\mathbf L}L_3,L_2\otimes_S^{\mathbf L}L_4).
\end{equation}
Hence a pairing, called the \emph{external tensor product},
\begin{equation}
\tag{5.3.3}
\Hom_S(L_1,L_2)\otimes\Hom_S(L_3,L_4)
\longrightarrow\Hom_S(L_1\otimes_S^{\mathbf L}L_3,L_2\otimes_S^{\mathbf L}L_4),
\end{equation}
which will be denoted $(u,v)\mapsto u\otimes_S^{\mathbf L}v$. The pairing 5.3.2 may also be interpreted as the composite of the external tensor product
\[
\RHom(E_1,E_2)\otimes_S^{\mathbf L}\RHom(E_3,E_4)
\xrightarrow{\sim}\RHom(E_1\otimes_S^{\mathbf L}E_3,E_2\otimes_S^{\mathbf L}E_4)
\]
(where $E_1=p_{12,1}^*L_1$, etc., with the evident notation) and of the isomorphism defined by Künneth
\[
E_2\otimes_S^{\mathbf L}E_4\xrightarrow{\sim}p_{24}^!(L_2\otimes_S^{\mathbf L}L_4)\quad (1.7.3).
\]

Let $c':C'\to X_{12}$ and $c'':C''\to X_{34}$ be correspondences, and let $c:C=C'\times_SC''\to X=X_{13}\times_SX_{24}$ be the product correspondence. From 5.3.2 one deduces, or else defines analogously, a pairing
\begin{equation}
\tag{5.3.4}
\RHom(c_1'{}^*L_1,c_2'{}^!L_2)\otimes_S^{\mathbf L}\RHom(c_3''{}^*L_3,c_4''{}^!L_4)
\longrightarrow\RHom(c_{13}^*(L_1\otimes_S^{\mathbf L}L_3),c_{24}^!(L_2\otimes_S^{\mathbf L}L_4)),
\end{equation}
hence an external tensor product pairing, $(u,v)\mapsto u\otimes_S^{\mathbf L}v$,
\begin{equation}
\tag{5.3.5}
\Hom(c_1'{}^*L_1,c_2'{}^!L_2)\otimes\Hom(c_3''{}^*L_3,c_4''{}^!L_4)
\longrightarrow\Hom(c_{13}^*(L_1\otimes_S^{\mathbf L}L_3),c_{24}^!(L_2\otimes_S^{\mathbf L}L_4)).
\end{equation}

One checks that these pairings are compatible with direct images, in the following sense. Let, for $1\leq i\leq4$, $f_i:X_i\to Y_i$ be an $S$-morphism, let $c'$, $c''$, $c$ be as above, with $c'$ and $c''$ proper, let $d':D'\to Y_{12}$ and $d'':D''\to Y_{34}$ be proper morphisms, let $d:D=D'\times_SD''\to Y=Y_{12}\times_SY_{34}$, and suppose given the commutative squares
\[
\begin{tikzcd}[column sep=large,row sep=large]
X_{12} \arrow[d,"f_{12}"'] & C' \arrow[l] \arrow[d] \\
Y_{12} & D' \arrow[l]
\end{tikzcd}
\qquad
\begin{tikzcd}[column sep=large,row sep=large]
X_{34} \arrow[d,"f_{34}"'] & C'' \arrow[l] \arrow[d] \\
Y_{34} & D'' \arrow[l]
\end{tikzcd}
\]
Then one has a commutative square
\begin{equation}
\tag{5.3.6}
\begin{tikzcd}[column sep=large,row sep=large]
\Hom(c_1'{}^*L_1,c_2'{}^!L_2)\otimes\Hom(c_3''{}^*L_3,c_4''{}^!L_4) \arrow[r,"5.3.5"] \arrow[d,"f_{12}\otimes f_{34}"']
& \Hom(c_{13}^*L_{13},c_{24}^!L_{24}) \arrow[d,"f_{13,24*}"] \\
\Hom(d_1'{}^*M_1,d_2'{}^!M_2)\otimes\Hom(d_3''{}^*M_3,d_4''{}^!M_4) \arrow[r,"5.3.5"]
& \Hom(d_{13}^*M_{13},d_{24}^!M_{24}),
\end{tikzcd}
\end{equation}
where $M_i=f_{i!}L_i$ if $i=1,3$ and $f_{i*}L_i$ otherwise, $L_{ij}=L_i\otimes_S^{\mathbf L}L_j$, and $M_{ij}=M_i\otimes_S^{\mathbf L}M_j$ (there is also an analogous square with the $\RHom$'s).

Suppose that $X_3=X_2$, $X_4=X_1$, $L_3=L_2$, and $L_4=L_1$. A reasoning analogous to that of the last paragraph of 5.2 shows that, for $u\in\Hom_S(L_1,L_2)$ and $v\in\Hom_S(L_2,L_1)$, one has
\begin{equation}
\tag{5.3.7}
\langle u,v\rangle=\langle u\otimes_S^{\mathbf L}v,1\rangle,
\end{equation}
where $1$ denotes the diagonal correspondence from $L_{12}$ to $L_{12}$ on $X_{12}$, and the cup product in the second member is taken in the sense of 4.2.5. More generally, if $c':C'\to X_{12}$ and $c'':C''\to X_{21}$ are correspondences, their ``intersection'' $C=C'\times_{X_{12}}C''$ identifies with the fibred product of $C'\times_SC''\to X=X_{12}\times_SX_{21}$ and the diagonal $X_{12}\to X$, and, for $u\in\Hom(c_1'{}^*L_1,c_2'{}^!L_2)$ and $v\in\Hom(c_2''{}^*L_2,c_1''{}^!L_1)$, one has the following equality in $H^0(C,K_C)$:
\begin{equation}
\tag{5.3.8}
\langle u,v\rangle=\langle u\otimes_S^{\mathbf L}v,1\rangle,
\end{equation}
where $1$ denotes, as above, the diagonal correspondence. Thus there are two ``diagonal'' procedures (5.2.10, 5.3.8) for the calculation of the ``Verdier local term'' $\langle u,v\rangle$.

\section*{6. Appendix. Lefschetz formula for coherent sheaves}

In this number, $S$ denotes a noetherian scheme. The $S$-schemes considered will be assumed separated and of finite type. Recall that, by Nagata, every $S$-morphism between such schemes is compactifiable, i.e. factors as an open immersion followed by a proper morphism. For every $S$-scheme $X$, we shall write $D(X)$ for the derived category of the category of $\cO_X$-modules.

\subsection*{6.1}
Recall ([8] App.) that for every $S$-morphism $f:X\to Y$ one has a functor
\[
f^!:D^+(Y)\longrightarrow D^+(X)
\]
(and transitivity isomorphisms $(gf)^! = f^!g^!$ for a composite $gf$, with a cocycle condition). If $f$ is proper, $f^!$ is right ``adjoint'' to the functor $f_*:D^+(X)_{\mathrm{coh}}\to D^+(Y)_{\mathrm{coh}}$ (where the subscript ``coh'' denotes the full subcategory formed by complexes with coherent cohomology), this adjunction being made precise by an isomorphism
\begin{equation}
\tag{6.1.1}
f_*\RHom(E,f^!F)\xrightarrow{\sim}\RHom(f_*E,F)
\end{equation}
for $E\in\operatorname{ob}D(X)_{\mathrm{coh}}$, $F\in\operatorname{ob}D^+(Y)$, defined analogously to the isomorphism (SGA~4 XVIII 3.1.9.6). The considerations of ([8] App.5), together with the duality theorem for projective space, show that, if $f$ is smooth, purely of dimension $d$, there is a canonical isomorphism
\begin{equation}
\tag{6.1.2}
f^!L\xleftarrow{\sim} f^*L\otimes\Omega^d_{X/Y}[d];
\end{equation}
and, if moreover $f$ is proper, the arrow
\begin{equation}
\tag{6.1.3}
\operatorname{Tr}_f:f_*(f^*L\otimes\Omega^d_{X/Y}[d])\longrightarrow L
\end{equation}
deduced, by 6.1.2, from the adjunction arrow $f_*f^!\to I$, coincides with the trace morphism ([8] VI).

\subsection*{6.2}
Contrary to what happens in the case of discrete coefficients, there is no reasonable subcategory of $D$ here that is stable under the ``six operations'' $(f^*,f_*,f^!,f_!,\otimes^{\mathbf L},\RHom)$. The most natural finiteness notion seems to be relative perfection (SGA~6 III 4). If $X$ is an $S$-scheme and $L\in\operatorname{ob}D(X)$, recall that $L$ is said to be perfect relative to $S$ if $L$ has finite tor-dimension relative to $S$ and coherent cohomology. We shall denote by $D(X)_{\mathrm{parf}/S}$ the full subcategory of $D(X)$ formed by the complexes perfect relative to $S$. It is not stable in general either under $\otimes^{\mathbf L}$ or under $\RHom$. However, if $E$ is perfect in the absolute sense and $F$ is perfect relative to $S$, then $\RHom(E,F)$ is perfect relative to $S$ (SGA~6 III 4.9). On the other hand, let, for $i=1,2$, $X_i$ be an $S$-scheme and $L_i\in\operatorname{ob}D^{-}(X_i)$. We shall put, as in 1.5,
\[
L_1\otimes_S^{\mathbf L}L_2=Lp_1^*L_1\otimes^{\mathbf L}Lp_2^*L_2,
\]
where $p_i:X_1\times_SX_2\to X_i$ is the projection. Considering this object is reasonable only if $X_1$ and $X_2$ are tor-independent over $S$, i.e. (SGA~6 III 1.5) if $\Tor_i^S(\cO_X,\cO_Y)=0$ for $i\ne0$. Suppose this is the case: then, if $L_1$ and $L_2$ are perfect relative to $S$, so is $L_1\otimes_S^{\mathbf L}L_2$ (SGA~6 III 4.7).

Let $f:X\to Y$ be an $S$-morphism. If $f$ is proper, then $f_*D(X)_{\mathrm{parf}/S}\subset D(Y)_{\mathrm{parf}/S}$ (SGA~6 III 4.8). If $f$ has finite tor-dimension, then $f^*D(Y)_{\mathrm{parf}/S}\subset D(X)_{\mathrm{parf}/S}$ (SGA~6 III 4.5), and $f^!D(Y)_{\mathrm{parf}/S}\subset D(X)_{\mathrm{parf}/S}$ (SGA~6 III 4.9).

For every $S$-scheme $a:X\to S$, we put
\begin{equation}
\tag{6.2.1}
K_X=a^!\cO_S,
\end{equation}
and write $D_X$ (or simply $D$) for the functor $\RHom(-,K_X)$. If $X$ has finite tor-dimension over $S$, then $K_X$ is perfect relative to $S$ and, for every $L\in\operatorname{ob}D(X)_{\mathrm{parf}/S}$, one has $DL\in\operatorname{ob}D(X)_{\mathrm{parf}/S}$, and the canonical arrow $L\to DDL$ is an isomorphism (SGA~6 III 4.9).

Let $f:X\to Y$ be an $S$-morphism, $L\in\operatorname{ob}D(X)$, and $M\in\operatorname{ob}D(Y)$. If $f^*M$ is defined (for example if $M\in\operatorname{ob}D^-(Y)$, or if $f$ has finite tor-dimension), put
\[
\Hom_f^{(1)}(L,M)=\Hom(f^*M,L),\qquad
\Hom_f^{(2)}(L,M)=\Hom(L,f^*M).
\]
If $f^!M$ is defined (for example if $M\in\operatorname{ob}D^+(Y)$, or if $f$ is smooth, or finite of finite tor-dimension, or a composite of such morphisms, since in these cases $f^!$ extends, by 6.1.1 and 6.1.2, to a functor defined on the whole category $D(Y)$), put
\[
\Hom_f^{(3)}(L,M)=\Hom(L,f^!M),\qquad
\Hom_f^{(4)}(L,M)=\Hom(f^!M,L).
\]
The elements of $\Hom_f^{(i)}(L,M)$ will be called, as in 1.4, \emph{arrows of type $i$ from $L$ to $M$ over $f$} (or $f$-arrows of type $i$ from $L$ to $M$). We leave to the reader the task of making explicit the composition of arrows of type $i$, for fixed $i$, and the corresponding fibred and cofibred categories. We recall only the partial adjunction formulas
\[
\Hom_f^{(1)}(L,M)=\Hom(M,f_*L)\quad(M\in\operatorname{ob}D^{-}(Y),\ L\in\operatorname{ob}D^+(X)),
\]
\[
\Hom_f^{(3)}(L,M)=\Hom(f_*L,M)\quad(f\text{ proper},\ L\in\operatorname{ob}D(X)_{\mathrm{coh}},\ M\in\operatorname{ob}D^+(Y)).
\]

\subsection*{6.3}
In the situation of 1.6, suppose that $X_1$ and $X_2$ are tor-independent over $S$, and likewise $Y_1$ and $Y_2$. Then, for $L_i\in\operatorname{ob}D(X_i)_{\mathrm{parf}/S}$, the Künneth arrow 1.6.4 is defined (use the fact that finite relative tor-dimension is preserved by direct image (SGA~6 III 3.7.1)) and is an isomorphism (to verify this, one may assume $S$ and the $Y_i$ affine, and, for a Čech calculation, reduce to also assuming $X_1$ and $X_2$ affine; the conclusion is then a trivial consequence of the tor-independence hypotheses). On the other hand, for $M_i\in\operatorname{ob}D(Y_i)_{\mathrm{parf}/S}$, with $f_1^!M_1$ or $f_2^!M_2$ in $D^b$, the Künneth arrow 1.7.3 is defined\footnote{In the proper or lissifiable case, the definition of the arrow is immediate; the general case can be reduced to the lissifiable case by cohomological descent.}, and is an isomorphism when $f_1$ and $f_2$ have finite tor-dimension, as is easily checked with the help of the Künneth isomorphism of type 1.6.4 just indicated, conjugated with 6.1.2, and with 6.1.1 in the case where $f_i$ is a closed immersion. In particular, if $X_1$ and $X_2$ have finite tor-dimension over $S$ and are tor-independent, 1.7.6 is an isomorphism.

One defines $\RHom(u,v)$ as in 2.1 for $(u,v)$ of type $(2,1)$ (respectively $(3,4)$, with $f$ proper in order to avoid the pro-objects caused by an $f_!$). Under suitable degree hypotheses (see for example ([9] V 2.3)), the evaluation arrows (2.2.2) and tensor-product arrows (2.2.3, 2.2.4) are defined, and the arrows 2.2.3 and 2.2.4 (provided they are defined) are isomorphisms as soon as $E_1$ and $E_2$ are perfect in the absolute sense.

In the situation of 2.4, suppose that the pairs $(X_1,X_2)$, $(X_1,Y_2)$, $(Y_1,X_2)$, $(Y_1,Y_2)$ are tor-independent over $S$, that $f_1$ is proper, that $E_i,F_i$, $P_i,Q_i$, $\RHom(E_i,F_i)$, $\RHom(P_i,Q_i)$ are perfect relative to $S$, and that the Künneth arrows
\[
f_1^!Q_1\otimes_S^{\mathbf L}Q_2\to(f_1\times_SY_2)^!(Q_1\otimes_S^{\mathbf L}Q_2),\qquad
f_1^!Q_1\otimes_S^{\mathbf L}F_2\to(f_1\times_SY_2)^!(Q_1\otimes_S^{\mathbf L}F_2)
\]
are isomorphisms. Then 2.4.1 and the diagram of 2.5 are defined, and the latter is commutative (however, the horizontal arrows are not necessarily isomorphisms). The verification is the same as for 2.5.

\subsection*{6.4}
Let, for $i=1,2$, $X_i$ be an $S$-scheme, let $X=X_1\times_SX_2$, and let $p_i:X\to X_i$, $q_i:X_i\to S$ be the canonical projections. Suppose that $X_1$ and $X_2$ are tor-independent and that $q_1$ has finite tor-dimension. Let $L_i\in\operatorname{ob}D(X_i)_{\mathrm{parf}/S}$. By the analogue of 2.2.4, one has a canonical arrow
\begin{equation}
\tag{6.4.1}
DL_1\otimes_S^{\mathbf L}L_2\longrightarrow \RHom(L_1\otimes_S^{\mathbf L}\cO_{X_2},q_1^!\cO_S\otimes_S^{\mathbf L}L_2),
\end{equation}
and hence, by composing with the Künneth isomorphism (6.3) $q_1^!\cO_S\otimes_S^{\mathbf L}L_2\xrightarrow{\sim}p_2^!L_2$, an arrow
\begin{equation}
\tag{6.4.2}
DL_1\otimes_S^{\mathbf L}L_2\longrightarrow\RHom(p_1^*L_1,p_2^!L_2).
\end{equation}
The arrow 6.4.1, and therefore also 6.4.2, is an isomorphism, as one sees by reducing to the case where $q_1$ is a closed immersion.

As in 3.1.2, we shall put
\begin{equation}
\tag{6.4.3}
\RHom(p_1^*L_1,p_2^!L_2)=\RHom_S(L_1,L_2),\qquad
\Hom(p_1^*L_1,p_2^!L_2)=\Hom_S(L_1,L_2),
\end{equation}
and call the elements of $\Hom_S(L_1,L_2)$ \emph{cohomological correspondences from $L_1$ to $L_2$}. If $c:C\to X$ is a correspondence (proper in practice), then by the induction formula ([8] III 8.8)\footnote{The ``lissifiability'' hypothesis of (loc. cit.) is unnecessary.}
\begin{equation}
\tag{6.4.4}
c^!\RHom(p_1^*L_1,p_2^!L_2)=\RHom(c_1^*L_1,c_2^!L_2),
\end{equation}
and the elements of
\[
\Hom(c_1^*L_1,c_2^!L_2)=H^0(X,c_*c^!\RHom(p_1^*L_1,p_2^!L_2))
\]
will be called \emph{cohomological correspondences from $L_1$ to $L_2$ with support in $C$}. We shall see below examples analogous to those of 3.2.

\subsection*{6.5}
Let, for $i=1,2$, $f_i:X_i\to Y_i$ be a proper $S$-morphism, and let $f=f_1\times_Sf_2:X\to Y$. Suppose that the $X_i$ and $Y_i$ have finite tor-dimension over $S$, and that the squares of 2.4.0 are tor-independent. Let $L_i\in\operatorname{ob}D(X_i)_{\mathrm{parf}/S}$. As in 3.3.1, one defines an isomorphism
\begin{equation}
\tag{6.5.1}
f_*\RHom_S(L_1,L_2)\xrightarrow{\sim}\RHom_S(f_{1*}L_1,f_{2*}L_2),
\end{equation}
hence, by applying $R\Gamma(Y,-)$, an isomorphism
\begin{equation}
\tag{6.5.2}
f_*:\Hom_S(L_1,L_2)\xrightarrow{\sim}\Hom_S(f_{1*}L_1,f_{2*}L_2),
\end{equation}
which we shall simply denote by $u\mapsto u_*$ when $Y_1=Y_2=S$. There is a commutative diagram analogous to 3.4.1, as follows from the analogue of 2.5 indicated in 6.3. One deduces from it a corollary analogous to 6.5 and, in the case where $Y_1=Y_2=S$ is the spectrum of a field, a description, analogous to that of 3.6, of the operations $u_*$ on cohomology.

Given a commutative square 3.7.1, with $c$ proper, one has, for $K\in\operatorname{ob}D^+(C)$, a canonical arrow
\begin{equation}
\tag{6.5.3}
f_*c_*c^!K\longrightarrow d_*d^!f_*K,
\end{equation}
defined in the same way as 3.7.3 (the isomorphism $g_*'c'^!\xrightarrow{\sim}d^!f_*$ is also valid in the present context: definition of the arrow by transposition of a base-change arrow, and verification that it is an isomorphism by reduction to the cases where $d$ is smooth and where $d$ is a closed immersion). Taking $K=\RHom_S(L_1,L_2)$, one deduces from 6.5.3 and 6.5.1 canonical arrows
\begin{equation}
\tag{6.5.4}
f_*c_*\RHom(c_1^*L_1,c_2^!L_2)\longrightarrow d_*\RHom(d_1^*(f_{1*}L_1),d_2^!(f_{2*}L_2)),
\end{equation}
\begin{equation}
\tag{6.5.5}
f_*:\Hom(c_1^*L_1,c_2^!L_2)\longrightarrow\Hom(d_1^*(f_{1*}L_1),d_2^!(f_{2*}L_2)),
\end{equation}
analogous to 3.7.5 and 3.7.6.

\subsection*{6.6}
Let, for $i=1,2$, $X_i$ be $S$-schemes of finite tor-dimension and tor-independent over $S$, and let $L_i\in\operatorname{ob}D(X_i)_{\mathrm{parf}/S}$. The external tensor product of the evaluation arrows $DL_i\otimes^{\mathbf L}L_i\to K_{X_i}$ defines, as in 4.1, ``cup-product'' pairings analogous to 4.1.2, 4.1.3, and 4.1.4.

Let $c:C\to X$ and $d:D\to X$ be proper morphisms, let $e=c\times_Xd:E\to X$, and let $P,Q\in\operatorname{ob}D(X)_{\mathrm{parf}/S}$. By 6.1.1, one has
\[
c_*c^!P=\RHom(Rc_*\cO_C,P)
\]
(and an analogous isomorphism with $d$). Assuming $c$ and $d$ of finite tor-dimension at the points above $e(E)$, one deduces, by tensor product, a canonical arrow
\begin{equation}
\tag{6.6.1}
c_*c^!P\otimes^{\mathbf L}d_*d^!Q\longrightarrow
\RHom(Rc_*\cO_C\otimes^{\mathbf L}Rd_*\cO_D,P\otimes^{\mathbf L}Q).
\end{equation}
When $c$ and $d$ are tor-independent, the second member is rewritten, by Künneth and 6.1.1, as $e_*e^!(P\otimes^{\mathbf L}Q)$, and 6.6.1 takes the form 4.2.3. For $P=DL_1\otimes_S^{\mathbf L}L_2$, $Q=L_1\otimes_S^{\mathbf L}DL_2$, one deduces from 6.6.1 and the pairings indicated above pairings analogous to 4.2.4 and 4.2.5:
\begin{equation}
\tag{6.6.2}
c_*\RHom(c_1^*L_1,c_2^!L_2)\otimes^{\mathbf L}d_*\RHom(d_2^*L_2,d_1^!L_1)
\longrightarrow\RHom(Rc_*\cO_C\otimes^{\mathbf L}Rd_*\cO_D,K_X),
\end{equation}
\begin{equation}
\tag{6.6.3}
\langle\ ,\ \rangle:\Hom(c_1^*L_1,c_2^!L_2)\otimes\Hom(d_2^*L_2,d_1^!L_1)
\longrightarrow\Hom(Rc_*\cO_C\otimes^{\mathbf L}Rd_*\cO_D,K_X).
\end{equation}
Just like the pairings of 4.2, these are compatible with étale localization. When $c$ and $d$ are tor-independent, the second member of 6.6.2 (respectively 6.6.3) identifies canonically with $e_*K_E$ (respectively $H^0(E,K_E)$).

Theorem 4.4 has the following analogue:

\begin{statement}{Theorem 6.7}
Let, for $i=1,2$, $f_i:X_i\to Y_i$ be a proper $S$-morphism, let $f=f_1\times_Sf_2:X\to Y$, and let $p_i:X\to X_i$, $q_i:Y\to Y_i$ be the projections. Suppose given, on the other hand, a cube 4.4.0 with cartesian horizontal faces, with $c'$, $c''$, $d'$, $d''$ proper, and with $c'$, $c''$ (respectively $d'$, $d''$) of finite tor-dimension at the points above $c(C)$ (respectively $d(D)$). Put $c_i'=p_ic'$, $c_i''=p_ic''$, $d_i'=q_id'$, and $d_i''=q_id''$. Finally, let $L_i\in\operatorname{ob}D(X_i)_{\mathrm{parf}/S}$. Assume that the pairs $(X_1,X_2)$, $(X_1,Y_2)$, $(Y_1,X_2)$, $(Y_1,Y_2)$ are tor-independent over $S$, and that $X_i$ and $Y_i$ have finite tor-dimension over $S$. Then one has a commutative square
\begin{equation}
\tag{6.7.1}
\begin{tikzcd}[column sep=small,row sep=large]
f_*c'_*\RHom(c_1'{}^*L_1,c_2'{}^!L_2)\otimes^{\mathbf L} f_*c''_*\RHom(c_2''{}^*L_2,c_1''{}^!L_1) \arrow[r,"(1)"] \arrow[d,"(2)"']
& f_*\RHom(Rc'_*\cO_{C'}\otimes^{\mathbf L}Rc''_*\cO_{C''},K_X) \arrow[d,"(4)"] \\
d'_*\RHom(d_1'{}^*M_1,d_2'{}^!M_2)\otimes^{\mathbf L} d''_*\RHom(d_2''{}^*M_2,d_1''{}^!M_1) \arrow[r,"(3)"]
& \RHom(Rd'_*\cO_{D'}\otimes^{\mathbf L}Rd''_*\cO_{D''},K_Y),
\end{tikzcd}
\end{equation}
where $M_i=f_{i*}L_i$, $(3)$ is given by 6.6.2, $(1)$ is deduced from 6.6.2 by application of $f_*$ and composition with the canonical arrow $f_*c'_*\otimes^{\mathbf L} f_*c''_*\to f_*(c'_*\otimes^{\mathbf L} c''_*)$, $(2)$ is the tensor product of the arrows 6.5.4 relative to the faces containing $f$, and $(4)$ is the composite of the duality isomorphism 6.1.1 and of the arrow defined by the canonical arrow
\[
Rd'_*\cO_{D'}\otimes^{\mathbf L}Rd''_*\cO_{D''}\longrightarrow f_*(Rc'_*\cO_{C'}\otimes^{\mathbf L}Rc''_*\cO_{C''})
\]
which is adjoint to the tensor product of the base-change arrows $f^*Rd'_*\cO_{D'}\to Rc'_*\cO_{C'}$, $f^*Rd''_*\cO_{D''}\to Rc''_*\cO_{C''}$.
\end{statement}

The proof is essentially the same as that of 4.4; we leave to the reader the small modifications to be made.

When $c'$ and $c''$ are tor-independent, and likewise $d'$ and $d''$, the arrow (4) of 6.7.1 is rewritten in the form $f_*c_*K_C\to d_*K_D$; it is simply defined by the adjunction arrow $g_*K_C=g_*g^!K_D\to K_D$ (as is arrow (4) of 4.4.1). However, it would not be natural to consider only this case, since that would exclude correspondences with excess intersection.

From 6.7 one deduces formulas analogous to 4.5.1 and 4.7.1: under the hypotheses of 6.7, one has, for $u\in\Hom(c_1'{}^*L_1,c_2'{}^!L_2)$ and $v\in\Hom(c_2''{}^*L_2,c_1''{}^!L_1)$,
\begin{equation}
\tag{6.7.2}
\langle f_*u,f_*v\rangle=f_*\langle u,v\rangle,
\end{equation}
where $f_*u$ (respectively $f_*v$) denotes the direct image correspondence of $u$ (respectively $v$) in the sense of 6.5.5, and $f_*$ in the second member is the homomorphism
\begin{equation}
\tag{6.7.3}
\Hom(Rc'_*\cO_{C'}\otimes^{\mathbf L}Rc''_*\cO_{C''},K_X)
\longrightarrow
\Hom(Rd'_*\cO_{D'}\otimes^{\mathbf L}Rd''_*\cO_{D''},K_Y)
\end{equation}
which is the composite of the duality isomorphism 6.1.1 and the arrow induced by the canonical arrow
\[
Rd'_*\cO_{D'}\otimes^{\mathbf L}Rd''_*\cO_{D''}\longrightarrow f_*(Rc'_*\cO_{C'}\otimes^{\mathbf L}Rc''_*\cO_{C''})
\]
cited in 6.7. When $Y_1=Y_2=S=D'=D''$, 6.7.2 is the \emph{Lefschetz-Verdier formula}.

As an illustration, we shall show that this formula implies the Woods Hole formula [1].

\subsection*{6.8}
For this purpose we shall need a few reminders on the ``Hodge'' cohomology class associated with a regularly embedded subscheme of a smooth scheme and on Grothendieck residues ([8] III 9). Let $f:X\to S$ be a smooth morphism purely of relative dimension $N$, and let $i:Y\hookrightarrow X$ be a regular immersion of codimension $d$, i.e. ((SGA~IV 16.9.2), (SGA~6 VII 1)) a closed subscheme defined locally by a regular sequence of $d$ equations, or, in equivalent terminology ((8 III 1), (SGA~6 VIII 1)), a closed immersion, locally a complete intersection, of codimension $d$. Put $g=fi$. Then one associates with $Y$ an element
\begin{equation}
\tag{6.8.1}
\cl^d_{X/S}(Y)\in\Ext^d_{\cO_X}(\cO_Y,\Omega^d_{X/S}),
\end{equation}
called the \emph{cohomology class associated with $Y$}, defined as follows. By the ``purity'' theorem ([8] III 7.2), one has
\begin{equation}
\tag{6.8.2}
\Ext^i_{\cO_X}(\cO_Y,\Omega^d_{X/S})=
\begin{cases}
0,& i\ne d,\\
\Hom(\bigwedge^d N_{Y/X},\Omega^d_{X/S}\otimes\cO_Y),& i=d,
\end{cases}
\end{equation}
where $N_{Y/X}=I/I^2$ is the conormal sheaf ($I=\Ker\,\cO_X\to\cO_Y$). Therefore
\begin{equation}
\tag{6.8.3}
\Ext^d_{\cO_X}(\cO_Y,\Omega^d_{X/S})=H^0(Y,\underline{\Hom}(\bigwedge^dN_{Y/X},\Omega^d_{X/S}\otimes\cO_Y)).
\end{equation}

The differential $d_{X/S}:\OO_X\to\Omega^1_{X/S}$ induces an $\OO_Y$-linear map
\[
d_{X/S}\otimes 1:N_{Y/X}\longrightarrow \Omega^1_{X/S}\otimes\OO_Y.
\]
The $d$-th exterior power of this map is a global section, on $Y$, of $\Lambda^d N_{Y/X}^{\vee}\otimes\Omega^d_{X/S}$; its image in $\Ext^d_{\OO_X}(\OO_Y,\Omega^d_{X/S})$ under the isomorphism 6.8.3 is, by definition, the class 6.8.1.

Assume $g$ finite and flat (hence $d=N$). Let $\omega\in H^0(Y,\Lambda^d N_{Y/X}^{\vee}\otimes\Omega^d_{X/S})$. By 6.8.2 and 6.1, one may view $\omega$ as a homomorphism
\[
\OO_Y\longrightarrow g^!\OO_S(=\Lambda^d N_{Y/X}^{\vee}\otimes\Omega^d_{X/S}),
\]
and therefore, by adjunction, as a homomorphism $g_*\OO_Y\to\OO_S$. The residue (relative to $g$) of $\omega$ is defined as the section of $\OO_S$ which is the image of $1$ under this homomorphism, and is denoted $\operatorname{Res}_{Y/S}(\omega)$ (cf. ([8] III 9)). Let, on the other hand, $i_*\omega\in H^N(X,\Omega^N_{X/S})$ denote the image of $\omega$ by the composite of the isomorphism 6.8.3 and the canonical arrow
\[
\Ext^N_{\OO_X}(\OO_Y,\Omega^N_{X/S})\longrightarrow H^N(X,\Omega^N_{X/S}).
\]
When $f$ is proper, it follows from transitivity for trace morphisms (or adjunction arrows) $i_*i^!\to 1$, $f_*f^!\to 1$, that
\begin{equation}
\tag{6.8.4}
\operatorname{Res}_{Y/S}(\omega)=\operatorname{Tr}_f(i_*\omega),
\end{equation}
where $\operatorname{Tr}_f:H^N(X,\Omega^N_{X/S})\to H^0(S,\OO_S)$ is defined by the adjunction arrow $f_*f^!\OO_S=Rf_*\Omega^N_{X/S}[N]\to\OO_S$ (6.1). On the other hand, without any properness hypothesis, one has the important compatibility ([8] III 9 R6)
\begin{equation}
\tag{6.8.5}
\operatorname{Res}_{Y/S}(y\cl_{X/S}(Y))=\operatorname{Tr}_{Y/S}(y),
\end{equation}
for $y\in H^0(Y,\OO_Y)$, where $\operatorname{Tr}_{Y/S}:g_*\OO_Y\to\OO_S$ denotes the trace in the classical sense, i.e. $\operatorname{Tr}_{Y/S}(y)$ is the trace of the endomorphism of $g_*\OO_Y$ defined by multiplication by $y$ (this compatibility is not proved in loc. cit.; the case $N=1$ is treated in [Raynaud, Anneaux locaux henséliens, VII 1, Lecture Notes in Math. n\textsuperscript{o} 169, Springer Verlag, 1970]). In other words, $\operatorname{Tr}_{Y/S}$ corresponds, by adjunction, to the homomorphism $\OO_Y\to g^!\OO_S$ defined by multiplication by $\cl_{X/S}(Y)$. It follows that, for $F\in\ob D(S)_{\coh}$, the arrow
\begin{equation}
\tag{6.8.6}
g^*F\longrightarrow g^!F=g^*F\otimes^{\LL}g^!\OO_S
\end{equation}
given by multiplication by $\cl_{X/S}(Y)$ corresponds by adjunction to the arrow
\begin{equation}
\tag{6.8.7}
g_*g^*F=F\otimes^{\LL}g_*\OO_Y\longrightarrow F
\end{equation}
defined by $\operatorname{Tr}_{Y/S}:g_*\OO_Y\to\OO_S$.


\subsection*{6.9}

Let $X$ (respectively $Y$) be a proper smooth $S$-scheme of relative dimension $m$ (respectively $n$), $Z=X\times_S Y$, and let $a:A\hookrightarrow Z$, $b:B\hookrightarrow Z$ be regular immersions such that $a_2:A\to Y$, $b_1:B\to X$ are finite and flat (write $p_1:Z\to X$, $p_2:Z\to Y$ for the projections, and $a_i=p_i a$, $b_i=p_i b$). Let $c:C=A\times_Z B\to Z$, $L\in\ob D(X)_{\parf}$, $M\in\ob D(Y)_{\parf}$ (the subscript parf means ``perfect in the absolute sense'', which is also equivalent to being perfect rel. to $S$, since $X$ and $Y$ are smooth). Assume that $C$ is finite and flat over $S$. The hypotheses made therefore imply that $c$ is a regular immersion and that $a$ and $b$ are tor-independent.

One has $p_1^*\Omega^1_{X/S}=\Omega^1_{Z/Y}$, $p_2^*\Omega^1_{Y/S}=\Omega^1_{Z/X}$, and a decomposition
\[
\Omega^1_{Z/S}=\Omega^1_{Z/Y}\oplus\Omega^1_{Z/X},
\]
whence, for every integer $r$, a decomposition
\begin{equation}
\tag{6.9.0}
\Omega^r_{Z/S}=\bigoplus_{i+j=r}\Omega^i_{Z/Y}\otimes\Omega^j_{Z/X}.
\end{equation}
On the other hand, by 6.1.2, one has $p_2^!M=p_2^*M\otimes\Omega^m_{Z/Y}[m]$, hence
\[
H^0(Z,a_*a^!\RHom(p_1^*L,p_2^!M))=\Ext^m_{\OO_Z}(a_1^*L,a_2^*M\otimes\Omega^m_{Z/Y}),
\]
whence, by 6.4.4 and 6.8.2, an isomorphism
\begin{equation}
\tag{6.9.1}
\Hom(a_1^*L,a_2^!M)=H^0(A,\RHom(a_1^*L,a_2^*M)\otimes\Lambda^m N_{A/Z}^{\vee}\otimes\Omega^m_{Z/Y}).
\end{equation}
Similarly one has an isomorphism
\begin{equation}
\tag{6.9.2}
\Hom(b_2^*M,b_1^!L)=H^0(B,\RHom(b_2^*M,b_1^*L)\otimes\Lambda^n N_{B/Z}^{\vee}\otimes\Omega^n_{Z/X}).
\end{equation}
Let $u\in\Hom(a_1^*L,a_2^*M)$, $v\in\Hom(b_2^*M,b_1^*L)$, and write
\begin{equation}
\tag{6.9.3}
\langle u,v\rangle\in H^0(C,\OO_C)
\end{equation}
for the section $\operatorname{Tr}(u_Cv_C)=\operatorname{Tr}(v_Cu_C)$, where $u_C\in\Hom(c_1^*L,c_2^*M)$ and $v_C\in\Hom(c_2^*M,c_1^*L)$ are induced by $u$ and $v$. Consider the cohomological correspondence $u^!\in\Hom(a_1^*L,a_2^!M)$ (respectively $v^!\in\Hom(b_2^*M,b_1^!L)$) deduced from $u$ (respectively $v$) by composition with the trace homomorphism 6.8.7 $a_{2*}a_2^*M\to M$ (respectively $b_{1*}b_1^*L\to L$), or equivalently (6.8.6) by product with $\cl_{Z/Y}(A)$ (respectively $\cl_{Z/X}(B)$). Note in passing that $\cl_{Z/Y}(A)$ (respectively $\cl_{Z/X}(B)$) is the component in $H^0(A,\Lambda^mN_{A/Z}^{\vee}\otimes\Omega^m_{Z/Y})$ (respectively $H^0(B,\Lambda^nN_{B/Z}^{\vee}\otimes\Omega^n_{Z/X})$) of $\cl_{Z/S}(A)$ (respectively $\cl_{Z/S}(B)$) defined by the projection $\Omega^m_{Z/S}\to\Omega^m_{Z/Y}$ (respectively $\Omega^n_{Z/S}\to\Omega^n_{Z/X}$) of 6.9.0. One checks easily that the cup-product
\[
\langle u^!,v^!\rangle\in\Ext^{m+n}_{\OO_Z}(\OO_C,\Omega^{m+n}_{Z/S})=H^0(C,\Lambda^{m+n}N_{C/Z}^{\vee}\otimes\Omega^{m+n}_{Z/S})
\]
defined in 6.6.3 is given by
\begin{equation}
\tag{6.9.4}
\langle u^!,v^!\rangle=\langle u,v\rangle\cl_{Z/Y}(A)\cl_{Z/X}(B).
\end{equation}
On the other hand one checks that the homomorphism 6.7.3
\[
\Ext^{m+n}_{\OO_Z}(\OO_C,\Omega^{m+n}_{Z/S})\longrightarrow H^0(S,\OO_S)
\]
is the composite of the canonical arrow $\Ext^{m+n}_{\OO_Z}(\OO_C,\Omega^{m+n}_{Z/S})\to H^{m+n}(Z,\Omega^{m+n}_{Z/S})$ and the trace morphism $\operatorname{Tr}_{Z/S}:H^{m+n}(Z,\Omega^{m+n}_{Z/S})\to H^0(S,\OO_S)$. The Lefschetz-Verdier formula 6.7.2 therefore gives the relation
\begin{equation}
\tag{6.9.5}
\langle (u^!)_*,(v^!)_*\rangle=\operatorname{Tr}_{Z/S}(\langle u,v\rangle\cl_{Z/Y}(A)\cl_{Z/X}(B)).
\end{equation}
Thanks to 6.8.4 one may rewrite the second member as a residue; thus, in summary, one has obtained:


\begin{statement}{Theorem 6.10}
Under the hypotheses of the first paragraph of 6.9, let $u\in\Hom(a_1^*L,a_2^*M)$, $v\in\Hom(b_2^*M,b_1^*L)$, and consider the cohomological correspondence $u^!$ (respectively $v^!$) from $L$ to $M$ (respectively from $M$ to $L$), with support in $A$ (respectively $B$), deduced from $u$ (respectively $v$) by composition with the trace $a_{2*}a_2^*M\to M$ (respectively $b_{1*}b_1^*L\to L$) (6.8.7), and the homomorphism $(u^!)_*:f_*L\to g_*M$ (respectively $(v^!)_*:g_*M\to f_*L$) deduced from $u^!$ (respectively $v^!$) by direct image (6.5.5) ($f:X\to S$, $g:Y\to S$ denoting the projections). Then, with the notation of 6.8, 6.9.3, one has
\begin{equation}
\tag{6.10.1}
\langle (u^!)_*,(v^!)_*\rangle=\operatorname{Res}_{C/S}(\langle u,v\rangle\cl_{Z/Y}(A)\cl_{Z/X}(B)).
\end{equation}
\end{statement}

\begin{statement}{Remarks 6.11}
a) As in 3.6, using the description 3.5 b) of the direct image of a correspondence, one checks that, if $h:Z\to S$ is the projection, $(u^!)_*$ (respectively $(v^!)_*$) is the composite of the arrows
\[
f_*L\xrightarrow{(1)}(ha)_*a_1^*L\xrightarrow{(ha)_*u}(ha)_*a_2^*M=g_*a_{2*}a_2^*M\xrightarrow{(2)}g_*M,
\]
\[
\text{(respectively }g_*M\xrightarrow{(1)}(hb)_*b_2^*M\xrightarrow{(hb)_*v}(hb)_*b_1^*L=f_*b_{1*}b_1^*L\xrightarrow{(2)}f_*L\text{)},
\]
where (1) is the usual functoriality arrow (defined by the adjunction arrow $1\to a_{1*}a_1^*$ (respectively $1\to b_{2*}b_2^*$)) and (2) is given by the trace $a_{2*}a_2^*M\to M$ (respectively $b_{1*}b_1^*L\to L$) (6.8.7).

b) Suppose that $S$ is the spectrum of a field. The intersection $C=A\times_{(X\times_S Y)}B$ is then formed of isolated points. At each point $z$ of $C$, choose, near the images of $z$, regular sequences $(s_1,\ldots,s_m)$, $(t_1,\ldots,t_n)$ of equations of $A,B$. Then 6.10.1 is written, with the notation of the residue symbols of ([8] III 9), as
\begin{equation}
\tag{6.11.1}
\langle (u^!)_*,(v^!)_*\rangle=
\sum_{z\in C}\operatorname{Res}_z\left(\langle u,v\rangle
\prod_{1\le i\le m}\frac{d_{Z/Y}s_i}{s_i}
\prod_{1\le j\le n}\frac{d_{Z/X}t_j}{t_j}\right),
\end{equation}
where $d_{Z/Y}s_i$ (respectively $d_{Z/X}t_j$) denotes the component of type $(1,0)$ (respectively $(0,1)$) of $d_{Z/S}s_i$ (respectively $d_{Z/S}t_j$) in the decomposition 6.9.0.
\end{statement}

\begin{statement}{Corollary 6.12}
(Woods-Hole formula [1]). Assume that $S$ is the spectrum of a field. Let $X$ be a proper smooth $S$-scheme, $F$ an $S$-endomorphism of $X$, $L$ an object of $D(X)_{\parf}$ (i.e. a perfect complex on $X$), and $u\in\Hom(F^*L,L)$. Assume that the scheme of fixed points $X^F$ is finite and consists of transverse points (i.e. points where the graph of $F$ meets the diagonal transversely). Then
\[
\sum_i(-1)^i\operatorname{Tr}((F,u),H^i(X,L))=
\sum_{x\in X^F}\frac{\operatorname{Tr}(u_x)}{\det(1-dF_x)},
\]
where $(F,u)$ is the endomorphism of $H^i(X,L)$ which is the composite of $H^i(X,L)\to H^i(X,F^*L)$ (inverse image) and $H^i(X,u)$, $\operatorname{Tr}(u_x)$ denotes the value at $x$ of the endomorphism induced by $u$ on the fibre of $L$ at $x$, and $\det(1-dF_x)$ is the value at $x$ of the determinant of the map $(1-\text{map induced by }F\text{ on }\Omega^1_{X,x})$.

\emph{Proof.} Apply 6.10, taking account of 6.11 a), b). Put $Z=X\times_S X$. Let $x\in X^F$. Near $x$, choose a system of local coordinates $(x_1,\ldots,x_n)$ (defining an étale morphism into $\mathbb A_S^n$), put $F_i=x_iF$, and write $(x_1,\ldots,x_n,y_1,\ldots,y_n)$ for the corresponding local coordinates on $Z$ near $z=(x,x)$ ($x_i=x_i\otimes1$, $y_i=1\otimes x_i$). Then, near $z$, $(x_1-y_1,\ldots,x_n-y_n)$ (respectively $(y_1-F_1,\ldots,y_n-F_n)$) is a regular system of equations for the diagonal $A$ (respectively for the graph $B$ of $F$). The contribution at $z$ of the second member of 6.11.1 is
\[
\langle u,1\rangle_z=
\operatorname{Res}_z\left(\operatorname{Tr}(u)\frac{dx_1\cdots dx_n\,dy_1\cdots dy_n}{(x_1-y_1)\cdots(x_n-y_n)(y_1-F_1)\cdots(y_n-F_n)}\right).
\]
Observe that in the numerator one may replace $dy_i$ by
\[
dy_i-dF_i=dy_i-\sum_{1\le j\le n}(dF_i/dx_j)dx_j.
\]
Applying then formula ([8] III 9 R3) to the situation $(B\hookrightarrow Z\to S)$, $(s_1,\ldots,s_p)=(y_1-F_1,\ldots,y_n-F_n)$, $(t_1,\ldots,t_n)=(x_1-y_1,\ldots,x_n-y_n)$, $\omega=(\operatorname{Tr}(u)dx_1\cdots dx_n)$, one obtains
\begin{equation}
\tag{*}
\langle u,1\rangle_z=
\operatorname{Res}_x\left(\operatorname{Tr}(u)\frac{dx_1\cdots dx_n}{(x_1-F_1)\cdots(x_n-F_n)}\right).
\end{equation}
Since $x$ is transverse, the $x_i-F_i$ form a regular system of parameters of $\OO_{X,x}$, and hence one may write, near $x$,
\[
x_i=\sum_j(x_j-F_j)c_{ij},
\]
where $(c_{ij})$ is a matrix of sections of $\OO_X$ invertible at $x$, whose determinant has value at $x$ equal to $\det(1-dF_x)^{-1}$. Thanks to ([8] III 9 R1), one therefore deduces from (*) that
\[
\langle u,1\rangle_z=\operatorname{Res}_x\left(\operatorname{Tr}(u)\det(c_{ij})\frac{dx_1\cdots dx_n}{x_1\cdots x_n}\right),
\]
whence, by ([8] III 9 R6),
\[
\langle u,1\rangle_z=\operatorname{Tr}(u(x))\det(c_{ij}(x))=\operatorname{Tr}(u(x))\det(1-dF_x)^{-1}.
\]
The corollary therefore follows from 6.11.1.
\end{statement}

\begin{statement}{Remark 6.12.1}
If $X^F$ is assumed only finite, the preceding calculation shows that one has
\[
\sum_i(-1)^i\operatorname{Tr}((F,u),H^i(X,L))=
\sum_{x\in X^F}\operatorname{Res}_x\left(\operatorname{Tr}(u)\frac{dx_1\cdots dx_n}{(x_1-F_1)\cdots(x_n-F_n)}\right),
\]
where $(x_i)$ is a system of local coordinates at $x$, and the $F_i$ are the coordinates of $F$ (near $x$) in this system.

The reader will convince himself, moreover, that the preceding formula is still valid if $X$ is proper and is assumed only to be smooth at the points of $X^F$.
\end{statement}

\subsection*{6.13}
Assume that $S$ is the spectrum of an algebraic closure of the finite field $\mathbb F_q$ and that $X/S$ is obtained by extension of scalars from a proper smooth scheme $X_0/\mathbb F_q$. Take for $F$ the Frobenius endomorphism of $X/S$ ``defined by raising coordinates to the $q$-th power''. One has $X^F=X_0(\mathbb F_q)$, and $X^F$ consists of transverse points since $dF=0$. If $L$ is a locally free sheaf of finite type on $X$ and $u\in\Hom(F^*L,L)$, the Woods-Hole formula is therefore written as
\begin{equation}
\tag{6.13.1}
\sum_i(-1)^i\operatorname{Tr}((F,u),H^i(X,L))=
\sum_{x\in X_0(\mathbb F_q)}\operatorname{Tr}(u_x,L_x).
\end{equation}
In particular, for $L=\OO_X$, and $u$ the canonical isomorphism $F^*\OO_X\xrightarrow{\sim}\OO_X$, one finds:
\begin{equation}
\tag{6.13.2}
\sum_i(-1)^i\operatorname{Tr}(F,H^i(X,\OO_X))=\operatorname{Card}(X_0(\mathbb F_q))\quad(\operatorname{mod} p).
\end{equation}
Artin-Schreier theory makes it possible to rewrite the first member in the form $\sum_i(-1)^i\operatorname{Tr}(F,H^i(X,\mathbb F_p))$, where $H^*(X,\mathbb F_p)$ denotes the étale cohomology of $X$ with values in the constant sheaf $\mathbb F_p$. By a refinement of this argument, Deligne, in (SGA $4^{1/2}$ Fonctions L), deduced from 6.13.1 the following generalization of 6.13.2:


\begin{statement}{Theorem 6.13.3 (Deligne)}
Let $X_0$ be a separated scheme of finite type over $\mathbb F_q$, let $X$ be the scheme deduced from $X_0$ by extension of scalars to an algebraic closure $k$ of $\mathbb F_q$, and let $F$ be the Frobenius endomorphism of $X/k$. Then, for every $L_0\in\ob D_{\ctf}(X_0,\mathbb F_p)$, if $L$ denotes the inverse image of $L_0$ on $X$, one has
\[
\sum_i(-1)^i\operatorname{Tr}(F,H_c^i(X,L))=
\sum_{x\in X^F}\operatorname{Tr}(F,L_x).
\]
\end{statement}

Let us point out that the particular case of 6.13.3 corresponding to $L_0=\mathbb F_p$ was established by Katz in (SGA 7 XXII) by a completely different method, based on Dwork's calculation of the zeta function of a hypersurface. Moreover, 6.13.2 in the proper smooth case is also an immediate consequence of the Lefschetz formula in crystalline cohomology ([5] VII 3.1.9). In fact, crystalline cohomology gives the more precise information that each factor $\det(1-Ft,H^i(X_0/W))\in W[t]$ of the zeta function is congruent mod $p$ to $\det(1-Ft,H^i(X_0,\OO_{X_0}))$ (at least if $H^*(X_0/W)$ is torsion-free). Katz showed [11] that, using estimates of B. Mazur [14], one could obtain higher congruences (involving the $H^i(X_0,\Omega^j_{X_0})$ and the Cartier operator) for smooth hypersurfaces (or complete intersections) for which the Hodge numbers (in the interesting dimension) are assumed to vanish up to a certain level.

On the other hand, if $X_0$ is a proper scheme over $\mathbb F_q$, geometrically integral, and such that $H^i(X_0,\OO_{X_0})=0$ for $i>0$, the Lefschetz-Verdier formula 6.7.2, applied to $X=X_0\otimes k$ and the Frobenius endomorphism of $X$, immediately implies that $X_0$ has at least one rational point over $\mathbb F_q$ (since $\operatorname{Tr}(F,R\Gamma(X,\OO_X))=1$). Serre asks whether this conclusion is still true when one replaces $\mathbb F_q$ by a field satisfying $(C_1)$ ([17] II 3), for example $\mathbb C(t)$ (Tsen), or $\mathbb C((t))$ (Lang) (see loc. cit. for other examples); indeed, he observes that if $X_0$ is a complete intersection of multidegree $(m_1,\ldots,m_r)$ in $\mathbb P^n$, the condition of triviality of cohomology is equivalent ([16] no. 78) to the condition $m_1+\cdots+m_r<n+1$, which is also the condition that occurs when one writes condition $(C_1)$.

\subsection*{6.14}
The Woods Hole formula has many other applications, for which we refer to [2] and to Beauville's exposé [4].


\subsection*{6.15}

The editor has not examined the case of components of fixed points of dimension $\geq1$, but it is not impossible that one could deduce from 6.7.2 (probably conjugated with Riemann-Roch) Lefschetz-Riemann-Roch formulae in the style of those of Donovan [7], Iversen [10], Nielsen [15] (at least with values in the base field). Nor is it impossible, on the other hand, that the formula 6.7.2, applied over the dual numbers, might provide residue formulae for vector fields analogous to those of Bott [6], [3].


\section*{Bibliography}
\begin{enumerate}[label={[\arabic*]},leftmargin=3.5em]
\item Atiyah, M. and Bott, R. On the Woods Hole Fixed point theorem, Proceedings of the Woods Hole Conference on Algebraic Geometry, 1964.
\item Atiyah, M. and Bott, R., Notes on the Lefschetz fixed point theorem for elliptic complexes, Harvard University, 1964.
\item Baum, P. and Bott, R., On the zeroes of meromorphic vector-fields, in Essays on topology and related topics, memoirs dedicated to Georges De Rham, Springer Verlag, 1970.
\item Beauville, A., Fixed point formulae in coherent cohomology, in Séminaire de géométrie algébrique, Université de Paris-Sud, Orsay, 1970.
\item Berthelot, P., Crystalline cohomology of schemes of characteristic $p>0$, Lecture Notes in Mathematics no. 407, Springer-Verlag, 1974.
\item Bott, R., A residue formula for holomorphic vector-fields, J. of Diff. Geometry Vol 1, p. 311--330, 1967.
\item Donovan, P., The Lefschetz-Riemann-Roch formula, Bull. S.M.F., 97, p. 257--273, 1969.
\item Hartshorne, R., Residues and duality, Lecture Notes in Mathematics no. 20, Springer-Verlag, 1966.
\item Illusie, L., Cotangent complex and deformations I, Lecture Notes in Math. no. 239, Springer-Verlag, 1971.
\item Iversen, B., A fixed point formula for action of tori on algebraic varieties, Inv. Math., 16, p. 229--236, 1972.
\item Katz, N., Letter to B. Mazur, 1971.
\item Kleiman, S., Algebraic cycles and the Weil conjectures, in Dix exposés sur la cohomologie des schémas, North Holland Pub. Company, 1968.
\item Langlands, R.P., Modular forms and $\ell$-adic representations, in Modular functions of one variable II, Lecture Notes in Math. no. 349, Springer-Verlag, 1973.
\item Mazur, B., Frobenius and the Hodge filtration (estimates), Ann. of Math., 98, p. 58--95, 1973.
\item Nielsen, H.A., Diagonalizably linearized coherent sheaves, Bull. SMF, 102, p. 85--97, 1974.
\item Serre, J.-P., Coherent algebraic sheaves, Ann. of Math., 61, p. 197--278, 1955.
\item Serre, J.P., Galois cohomology, Lecture Notes in Math. no. 5, Springer Verlag, 1964.
\item Verdier, J.-L., The Lefschetz fixed point formula in étale cohomology, Proceedings of a conference on local fields, Springer-Verlag, 1967.
\end{enumerate}

SGA $4^{1/2}$ Étale cohomology, by P. Deligne, Séminaire de Géométrie Algébrique du Bois-Marie, Lecture Notes in Mathematics no. 569, Springer-Verlag, 1977.

In SGA $4^{1/2}$: Finitude: Finiteness theorem in $\ell$-adic cohomology.

Cycle: The cohomology class associated with a cycle, by A. Grothendieck, edited by P. Deligne (replaces the absent exposé SGA 5 IV).

Rapport: Report on the trace formula.

Fonctions L: L-functions mod $\ell^n$ and mod $p$.

\clearpage
\phantomsection\addcontentsline{toc}{section}{Exposé III B --- Calculations of Local Terms}
\section*{Exposé III B\\Calculations of Local Terms}
\begin{center}
by Luc Illusie
\end{center}

\section*{Introduction}

This exposé, written in January 1977, does not correspond to any oral exposé of the seminar. It comprises two independent parts. In the first, we compute the Lefschetz--Verdier local terms for cohomological correspondences between smooth curves over an algebraically closed field, supported on correspondences satisfying certain general-position hypotheses. The formula that we obtain gives as corollaries Verdier's formula~[6] for transversal endomorphisms (which applies in particular to the case of the Frobenius endomorphism), and the formula stated by Langlands in ([4]~7.12). The method of proof is parallel to that followed by Artin--Verdier in [6]. One reduces to proving the vanishing of the local term when the given complexes of sheaves have zero fibre at the images of the fixed point under consideration, and in fact one proves a stronger property, namely that a certain restriction map is zero (2.5). To establish this property, one first reduces to the tamely ramified case by covers of the same degree of the two given curves (so as not to destroy the general-position hypotheses), and then treats this case directly by means of Abhyankar's lemma and the relative purity theorem. The second part of this exposé, of a much more technical nature, is inspired by the method used by Grothendieck to establish the Lefschetz formula for certain cohomological correspondences on curves (see XII and (SGA~$4^{1/2}$ Rapport)). In no.~5, for complexes of modules on a topos, we develop a theory of non-commutative traces which generalizes the commutative theory of (SGA~6~I) and, in the case of a punctual topos, recovers the non-commutative theory developed by Grothendieck in the oral seminar (this was written up in Exposé XI, unfortunately lost). We apply the techniques of no.~5 to define, in no.~6, Lefschetz--Verdier local terms for cohomological correspondences between complexes of modules over rings that are not necessarily commutative. These local terms behave like non-commutative traces and make it possible, in particular when one takes as base ring the algebra of a group, to divide canonically the ordinary local terms attached to equivariant correspondences. As an application, we prove Grothendieck's divisibility conjecture of (XII~4.5), in a more general form. We also show that the local terms defined by Grothendieck in the Lefschetz formula of XII are indeed the Lefschetz--Verdier local terms, and we generalize the formula of (loc. cit.).

I thank P.~Deligne for the useful suggestions he made to me during many discussions; I am particularly grateful to him for having pointed out, and helped me correct, an error in an early proof of Theorem~1.2.

\section*{Summary}
\begin{tabular}{@{}p{0.84\textwidth}r@{}}
I. Correspondences in general position between curves & 3\\
1. Statement of the theorem and corollaries & 3\\
2. Reduction to a vanishing theorem & 9\\
3. Reduction to the tame case & 15\\
4. End of the proof of 1.2 & 20\\
II. Equivariant correspondences & 25\\
5. Non-commutative traces & 25\\
6. Equivariant correspondences and divisibility of local terms & 46
\end{tabular}

\section*{I. Correspondences in General Position between Curves}

In this part, $S$ denotes the spectrum of an algebraically closed field $k$ of characteristic $p$, and $\Lambda$ a noetherian commutative ring killed by an integer prime to $p$. For every scheme $X$ over $S$ we denote by $D(X)$ the derived category of sheaves of $\Lambda$-modules on $X$, and by $D_{\ctf}(X)$ the full subcategory of $D(X)$ formed by complexes of finite tor-dimension and constructible cohomology.

\subsection*{1. Statement of the theorem and corollaries}

\subsubsection*{1.1}
Let $X$ (resp. $Y$) be the strict localization, at a closed point, of a smooth separated $S$-curve, let $A$ (resp. $B$) be a strictly local $S$-scheme, and let
\[
a:A\longrightarrow X\times_S Y,
\qquad
b:B\longrightarrow X\times_S Y
\]
be $S$-morphisms. We denote by $p_1:X\times_S Y\to X$, $p_2:X\times_S Y\to Y$ the projections, put $a_i=p_i a$, $b_i=p_i b$, and denote by $x$ (resp. $y$) the closed point of $X$ (resp. $Y$). We assume that $a_2$ and $b_1$ are finite and surjective, and that $A\times_{(X\times_S Y)}B$ is reduced to the closed point $z$ whose image in $X\times_S Y$ is $(x,y)$.

Let $L\in \ob D_{\ctf}(X)$, $M\in \ob D_{\ctf}(Y)$, $u\in\Hom(a_1^*L,a_2^!M)$, and $v\in\Hom(b_2^*M,b_1^!L)$. Denote by
\[
u_z\in\Hom(L_x,M_y)\qquad (\text{resp. }v_z\in\Hom(M_y,L_x))
\]
the composite homomorphism
\[
L_x=R\Gamma(X,L)\xrightarrow{a_1^*}R\Gamma(A,a_1^*L)
\xrightarrow{R\Gamma(A,u)}R\Gamma(A,a_2^!M)
\xrightarrow{a_{2*}}R\Gamma(Y,M)=M_y,
\]
where $a_{2*}$ is defined by the adjunction map $a_{2*}a_2^!M\to M$ (resp. the analogous composite homomorphism defined by means of $(b,v)$). Since $L_x$ and $M_y$ are perfect complexes of $\Lambda$-modules, one can form the cup-product (SGA~6~I~8.3)
\begin{equation}\tag{1.1.1}
\langle u_z,v_z\rangle=\Tr(u_zv_z)=\Tr(v_zu_z)\in\Lambda.
\end{equation}
Put $\Lambda_X(1)[2]=K_X$, $\Lambda_Y(1)[2]=K_Y$. By (I~5) (or (SGA~$4^{1/2}$ Duality~1, or Finitude~4.1)) the functor $D_X=\RHom(-,K_X)$ (resp. $D_Y=\RHom(-,K_Y)$) on $D_{\ctf}(X)$ (resp. $D_{\ctf}(Y)$) is dualizing, i.e. for every $E\in\ob D_{\ctf}(X)$ (resp. $E\in\ob D_{\ctf}(Y)$), one has $E\xrightarrow{\sim}D_XD_XE$ (resp. $E\xrightarrow{\sim}D_YD_YE$). We shall sometimes write $D$ instead of $D_X$, $D_Y$.

On the other hand, denote by $p_i^!$ the functor $p_i^*(-)(1)[2]$. We have trivially
\[
K_X\otimes_S^{\LL}M\xrightarrow{\sim}p_2^!M,
\qquad
L\otimes_S^{\LL}K_Y\xrightarrow{\sim}p_1^!L,
\]
and it follows from (III~2.3), by passage to the limit, that the Künneth maps analogous to (III~3.1.1)
\[
DL\otimes_S^{\LL}M\longrightarrow \RHom(p_1^*L,p_2^!M),
\qquad
L\otimes_S^{\LL}DM\longrightarrow \RHom(p_2^*M,p_1^!L)
\]
are isomorphisms. Since $a_2$ is finite, $a$ is finite, and hence the functor $a^!$ is defined; one verifies easily that one has the transitivity isomorphism $a^!p_2^!=a_2^!$. By the induction formula for $a$ (SGA~4~XVIII~3.1.12.2), one therefore has a canonical isomorphism analogous to (III~3.2.1)
\[
a^!\RHom(p_1^*L,p_2^!M)\xrightarrow{\sim}\RHom(a_1^*L,a_2^!M),
\]
whence an isomorphism
\begin{equation}\tag{*}
\Hom(a_1^*L,a_2^!M)\xrightarrow{\sim}H^0(X\times_S Y,a_*a^!(DL\otimes_S^{\LL}M)).
\end{equation}
Likewise one has an isomorphism
\begin{equation}\tag{**}
\Hom(b_2^*M,b_1^!L)\xrightarrow{\sim}H^0(X\times_S Y,b_*b^!(L\otimes_S^{\LL}DM)).
\end{equation}
As in (III~4.1.2), one defines a natural pairing
\[
(DL\otimes_S^{\LL}M)\otimes^{\LL}(L\otimes_S^{\LL}DM)
\longrightarrow K_X\otimes_S^{\LL}K_Y=\Lambda_{X\times_S Y}(2)[4]\stackrel{\mathrm{dfn}}{=}K_{X\times_S Y}.
\]
Composing with the pairing analogous to (III~4.2.3)
\[
a_*a^!P\otimes^{\LL} b_*b^!Q\longrightarrow c_*c^!(P\otimes^{\LL}Q),
\]
where $P=DL\otimes_S^{\LL}M$, $Q=L\otimes_S^{\LL}DM$, and $c:Z=A\times_{(X\times_S Y)}B\to X\times_S Y$ is the canonical projection, and taking into account the isomorphisms (*), (**), one deduces a pairing analogous to (III~4.2.5)
\[
\langle\ ,\ \rangle:\Hom(a_1^*L,a_2^!M)\otimes\Hom(b_2^*M,b_1^!L)
\longrightarrow H^0(X\times_S Y,c_*c^!K_{X\times_S Y}).
\]
One verifies easily, by passage to the limit, that $c^!K_{X\times_S Y}=\Lambda_z$, so that $H^0(X\times_S Y,c_*c^!K_{X\times_S Y})=\Lambda$. Thus one has a ``cup-product''
\begin{equation}\tag{1.1.2}
\langle u,v\rangle\in\Lambda.
\end{equation}

\subsubsection*{1.1.3}
Let $X'$ (resp. $Y'$) be a smooth $S$-curve, $A'$ (resp. $B'$) a separated $S$-curve,
\[
a':A'\longrightarrow X'\times_S Y',
\qquad
b':B'\longrightarrow X'\times_S Y'
\]
finite $S$-morphisms, $x'$ (resp. $y'$) a closed point of $X'$ (resp. $Y'$), $z'$ an isolated closed point of $A'\times_{(X'\times_S Y')}B'$, with image $s'$ (resp. $t'$) in $A'$ (resp. $B'$), and such that $a'_2(=p_2a')$ (resp. $b'_1(=p_1b')$) is quasi-finite at $s'$ (resp. $t'$),
\[
L'\in\ob D_{\ctf}(X'),\quad M'\in\ob D_{\ctf}(Y'),\quad
u'\in\Hom(a_1'^*L',a_2'^!M'),\quad v'\in\Hom(b_2'^*M',b_1'^!L').
\]
One can then consider the local term (III~4.2.7) $\langle u',v'\rangle_{z'}\in\Lambda$. On the other hand, if $X,Y,A,B$ are the strict localizations of $X',Y',A',B'$ at $x',y',s',t'$, if $a:A\to X\times_S Y$, $b:B\to X\times_SY$ are the morphisms deduced from $a'$, $b'$, if $L=L'|X$, $M=M'|Y$, and if $u\in\Hom(a_1^*L,a_2^!M)$, $v\in\Hom(b_2^*M,b_1^!L)$ are the maps deduced from $u'$, $v'$, then the data $(X,Y,a,b,u,v)$ satisfy the hypotheses at the beginning of 1.1, and hence a cup-product 1.1.2 $\langle u,v\rangle\in\Lambda$ is defined. It follows from the compatibility of the formation of local terms (III~4.2.7) with localization (III~4.2.6) that $\langle u',v'\rangle_{z'}=\langle u,v\rangle$. It is also easy to see that any data $(X,Y,a,b,u,v)$ satisfying the hypotheses at the beginning of 1.1 are induced by data $(X',Y',a',b',x',y',z',s',t',u',v')$ as above, where one may also assume, if desired, that $A'\times_{(X'\times_S Y')}B'$ is reduced to $z'$.

In general the cup-products 1.1.1 and 1.1.2 are unequal. Nevertheless one has the following theorem, which is the main result of this first part:

\begin{statement}{Theorem 1.2}
Under the hypotheses at the beginning of 1.1, assume moreover that $A$ and $B$ are in general position, i.e. that the intersection of the tangent cones at $(x,y)$ to the images of $A$ and $B$ in $X\times_S Y$ is reduced to $(x,y)$, and that likewise $A$ and $X\times\{y\}$ are in general position, as are $\{x\}\times Y$ and $B$. Then
\begin{equation}\tag{1.2.1}
\langle u_z,v_z\rangle=\langle u,v\rangle.
\end{equation}
\end{statement}

By the Lefschetz--Verdier formula (III~4.7), taking 1.1.3 into account, one deduces:

\begin{statement}{Corollary 1.3}
Let $F$ be an endomorphism of a proper smooth curve $X/S$ such that the fixed-point scheme $X^F$ is finite and formed of points where the graph of $F$ is transverse to the diagonal. Let $L\in\ob D_{\ctf}(X)$, and $u\in\Hom(F^*L,L)$. Then
\[
\Tr((F,u),R\Gamma(X,L))=\sum_{x\in X^F}\Tr(u_x,L_x),
\]
where the trace on both sides is taken in the sense of (SGA~6~I~8).
\end{statement}

For $L$ concentrated in degree $0$, 1.3 is the result of Artin--Verdier ([6]~4.1). By a passage to the limit explained in XV, one deduces from 1.3 an analogous result in $\ell$-adic cohomology ([6]~1.1), (XV~\S2 no.~3 prop.~2):

\begin{statement}{Corollary 1.3.1}
With $X$ and $F$ as in 1.3, let $L$ be a constructible $\mathbb Q_\ell$-sheaf on $X$ ($\ell$ prime to $p$), and let $u\in\Hom(F^*L,L)$. Then:
\[
\sum_i(-1)^i\Tr((F,u),H^i(X,L))=\sum_{x\in X^F}\Tr(u_x,L_x).
\]
\end{statement}

The hypotheses of 1.3 and 1.3.1 apply in particular to the case where $(F,u)$ is a Frobenius correspondence (XV~2), whence Grothendieck's trace formula ([2]~5.1), (SGA~$4^{1/2}$ Rapport~3.2).

\subsubsection*{1.4}
With the notation of 1.1, assume that $b_2$ and $a_1$ are finite and surjective (instead of $a_2$ and $b_1$), and consider the dual correspondences (III~5.1)
\[
Du\in\Hom(a_2^*DM,a_1^!DL),\qquad Dv\in\Hom(b_1^*DL,b_2^!DM),
\]
and the homomorphisms $(Du)_z\in\Hom((DM)_y,(DM)_x)$, $(Dv)_z\in\Hom((DL)_x,(DM)_y)$. Assume that $(A,\{x\}\times Y)$ is in general position, as are $(X\times\{y\},B)$ and $(A,B)$. Applying 1.2 to $Du,Dv$, one obtains
\[
\langle Du,Dv\rangle=\langle (Du)_z,(Dv)_z\rangle,
\]
whence, by (III~5.1.6) and (III~4.2.6),
\begin{equation}\tag{1.4.1}
\langle u,v\rangle=\langle (Du)_z,(Dv)_z\rangle.
\end{equation}

From this remark and from 1.2 one deduces, by the Lefschetz--Verdier formula (III~4.7):

\begin{statement}{Corollary 1.5}
Let $X,Y$ be proper smooth curves over $S$,
\[
a:A\longrightarrow X\times_S Y,
\qquad
b:B\longrightarrow X\times_S Y
\]
correspondences where $A$ and $B$ are proper curves over $S$, let $L\in\ob D_{\ctf}(X)$, $M\in\ob D_{\ctf}(Y)$, $u\in\Hom(a_1^*L,a_2^!M)$, and $v\in\Hom(b_2^*M,b_1^!L)$. Assume that $C=A\times_{(X\times_S Y)}B$ is finite over $S$ and of the form $C=C_1\amalg C_2$, with the following properties: for $z\in C_1$, the morphisms deduced from $a_2,b_1$ by strict localization at the images of $z$ are finite, and, if $(x,y)$ denotes the image of $z$ in $X\times_S Y$, each of the pairs $(A,B)$, $(A,X\times\{y\})$, $(\{y\}\times Y,B)$ is in general position at $(x,y)$; for $z\in C_2$, the morphisms deduced from $a_1,b_2$ by strict localization at the images of $z$ are finite, and each of the pairs $(A,B)$, $(A,\{x\}\times Y)$, $(B,X\times\{y\})$ is in general position at $(x,y)$ (``general position'' meaning that the intersection of the tangent cones at $(x,y)$ to the images of the curves of the pair under consideration is reduced to $(x,y)$). Then, if $u_*\in\Hom(R\Gamma(X,L),R\Gamma(Y,M))$, $v_*\in\Hom(R\Gamma(Y,M),R\Gamma(X,L))$ are the morphisms defined by $u,v$ (III~3.7.6), one has, with the notation of 1.1.1, (III~5.1.5),
\begin{equation}\tag{1.5.1}
\langle u_*,v_*\rangle=\sum_{z\in C_1}\langle u_z,v_z\rangle+\sum_{z\in C_2}\langle (Du)_z,(Dv)_z\rangle.
\end{equation}
\end{statement}

\subsubsection*{1.6}
Assuming the formalism of $\ell$-adic cohomology, one deduces from 1.5, by the passage-to-the-limit technique of XIV, that, under the hypotheses of 1.5 on $X,Y,a,b$, for $L\in\ob D^b(X,\mathbb Q_\ell)$ ($\ell$ prime $\ne p$), $M\in\ob D^b(Y,\mathbb Q_\ell)$, $u\in\Hom(a_1^*L,a_2^!M)$, $v\in\Hom(b_2^*M,b_1^!L)$, one has the analogue of 1.5.1 in $\mathbb Q_\ell$. This formula contains, as a special case, Proposition 7.12 of Langlands [4]: apply the formula with $X=Y$, $M$ the sheaf $F$ considered by Langlands; $L=\alpha M$ ($\alpha\in\mathbb Q_\ell^*$), $A$ the diagonal of $X\times X$, $b$ the correspondence $\phi$ of (loc. cit.), $u$ the correspondence supported on $A$ given by $\alpha^{-1}:L\to M$, $v$ the correspondence $\Phi$ of (loc. cit.); with the notation of (loc. cit.), $C_1$ (resp. $C_2$) is the set of fixed points where $a<d$ (resp. $a>d$).

\subsubsection*{1.7}
Under the hypotheses of 1.2, take for $u$ (resp. $v$) the correspondence $\cl(A)$ (resp. $\cl(B)$) defined by $\Tr_{a_2}:a_{2*}\Lambda\to\Lambda$ (resp. $\Tr_{b_1}:b_{1*}\Lambda\to\Lambda$). Formula 1.2.1 becomes
\begin{equation}\tag{1.7.1}
\deg(a_2)\deg(b_1)=\langle \cl(A),\cl(B)\rangle,
\end{equation}
where $\deg(a_2)$ (resp. $\deg(b_1)$) denotes the generic degree of $a_2$ (resp. $b_1$). Taking into account the hypotheses of general position, the left hand side of 1.7.1 is none other than the intersection multiplicity at $(x,y)$ of the cycles $a_*A$, $b_*B$, the direct images of $A$ and $B$. It is not difficult to deduce 1.7.1 directly from the compatibility of the formation of the cohomology class associated with a cycle with intersections (SGA~$4^{1/2}$ Cycle~2.3.8).


\subsection*{2. Reduction to a vanishing theorem}

\subsubsection*{2.1}
Under the hypotheses of 1.1, let $s$ (resp. $t$) denote the image of $z$ in $A$ (resp. $B$), and let $i:U=X-\{x\}\hookrightarrow X$ and $j:V=Y-\{y\}\hookrightarrow Y$ be the canonical inclusions. Since $a_2$ and $b_1$ are finite, one has $a_2^{-1}(y)=s$, $b_1^{-1}(x)=t$, hence
\[
a_1^{-1}(U)\subset a_2^{-1}(V),\qquad b_2^{-1}(V)\subset b_1^{-1}(U),
\]
and there are commutative diagrams
\[
\begin{tikzcd}[column sep=large,row sep=large]
U \arrow[d,"i"] \arrow[rd,phantom,"(1)"{description}] & a_1^{-1}(U) \arrow[l,"a'_1"'] \arrow[r,"a'_2"] \arrow[d] \arrow[rd,phantom,"(2)"{description}] & V \arrow[d,"j"]\\
X & A \arrow[l,"a_1"] \arrow[r,"a_2"'] & Y
\end{tikzcd}
\qquad
\begin{tikzcd}[column sep=large,row sep=large]
U \arrow[d,"i"] \arrow[rd,phantom,"(3)"{description}] & b_2^{-1}(V) \arrow[l,"b'_1"'] \arrow[r,"b'_2"] \arrow[d] \arrow[rd,phantom,"(4)"{description}] & V \arrow[d,"j"]\\
X & B \arrow[l,"b_1"] \arrow[r,"b_2"'] & Y ,
\end{tikzcd}
\]
where squares (1) and (4) are cartesian.

\begin{statement}{Lemma 2.2}
To prove 1.2, it is enough to consider the following two cases:
\[
\text{a) } L_x=0,\quad M_y=0;\qquad
\text{b) } L|U=0,\quad M|V=0 .
\]
\end{statement}

Consider the decreasing filtrations on $L$ and $M$ defined by the subcomplexes
\[
F^nL=
\begin{cases}
L,& n\le 0,\\
i_!i^*L,& n=1,\\
0,& n>1,
\end{cases}
\qquad
F^nM=
\begin{cases}
M,& n\le 0,\\
j_!j^*M,& n=1,\\
0,& n>1.
\end{cases}
\]
The complex $a_{2*}a_1^*L$ (resp. $b_{1*}b_2^*M$) is then filtered by
\[
F^n a_{2*}a_1^*L=a_{2*}a_1^*F^nL.
\]
We have $F^n a_{2*}a_1^*L=a_{2*}a_1^*L$ for $n\le0$, and $0$ for $n>1$; moreover, since $a_2$ is finite and square (1) is cartesian,
\[
F^1a_{2*}a_1^*L=j_!(a'_2)_!(a'_1)^*i^*L.
\]
Similarly,
\[
F^n b_{1*}b_2^*M=
\begin{cases}
b_{1*}b_2^*M,& n\le 0,\\
i_!(b'_1)_!(b'_2)^*j^*M,& n=1,\\
0,& n>1.
\end{cases}
\]

We may suppose that $u$ (resp. $v$) is defined by a homomorphism of complexes
\[
u:a_{2*}a_1^*L\longrightarrow M
\quad
(\text{resp. }v:b_{1*}b_2^*M\longrightarrow L).
\]
The preceding formulas show that the composite
\[
F^1a_{2*}a_1^*L\longrightarrow a_{2*}a_1^*L\xrightarrow{u}M
\]
(resp. $F^1b_{1*}b_2^*M\to b_{1*}b_2^*M\xrightarrow{v}L$) factors uniquely through $F^1M$ (resp. $F^1L$); in other words, $u$ (resp. $v$) is a homomorphism of filtered complexes. Thus $u$ (resp. $v$) underlies a filtered correspondence, in the sense of (III 4.13), which we again denote by
\[
u\in\Hom_{DF}(a_1^*L,a_2^!M),\qquad
v\in\Hom_{DF}(b_2^*M,b_1^!L).
\]
By additivity of cup-products ([3] V 3.8.7), (III 4.13.1), one has
\[
\langle u_z,v_z\rangle
 =\langle \operatorname{gr}^0u_z,\operatorname{gr}^0v_z\rangle
 +\langle \operatorname{gr}^1u_z,\operatorname{gr}^1v_z\rangle,
\]
and
\[
\langle u,v\rangle
 =\langle \operatorname{gr}^0u,\operatorname{gr}^0v\rangle
 +\langle \operatorname{gr}^1u,\operatorname{gr}^1v\rangle.
\]
The lemma follows, since $\operatorname{gr}^1L=i_!i^*L$ (resp. $\operatorname{gr}^1M=j_!j^*M$) has zero fibre at $x$ (resp. $y$), whereas $\operatorname{gr}^0L$ (resp. $\operatorname{gr}^0M$) is concentrated at $x$ (resp. $y$).

\begin{statement}{Lemma 2.3}
Under the hypotheses of 1.1, suppose that $L|U=0$ and $M|V=0$. Then
\[
\langle u_z,v_z\rangle=\langle u,v\rangle .
\]
\end{statement}

The idea of the proof is that then $u$ and $v$ are direct images of correspondences between complexes concentrated at the closed points, and that the formula to be proved merely expresses the Lefschetz formula for the projections $x\to S$, $y\to S$. In view of 1.1.3, the data of 1.1 may be replaced by the following data: $X$ (resp. $Y$) is a smooth $S$-curve endowed with a closed point $x$ (resp. $y$), $A$ (resp. $B$) is a separated $S$-curve, $a:A\to X\times_S Y$, $b:B\to X\times_S Y$ are finite $S$-morphisms such that $A\times_{X\times_SY}B$ is reduced to a closed point $z$ above $(x,y)$,
\[
u\in\Hom(a_1^*L,a_2^!M),\qquad v\in\Hom(b_2^*M,b_1^!L),
\]
with
\[
L=h_{1*}L',\quad M=h_{2*}M',\quad
L'\in\Ob D_{\mathrm{ctf}}(\{x\}),\quad M'\in\Ob D_{\mathrm{ctf}}(\{y\}),
\]
where $h_1:\{x\}\to X$ and $h_2:\{y\}\to Y$ are the canonical inclusions. We must show that
\[
\langle u_z,v_z\rangle=\langle u,v\rangle_z,
\]
where
\[
u_z\in\Hom(\mathrm R\Gamma(X,L),\mathrm R\Gamma(Y,M)),\qquad
v_z\in\Hom(\mathrm R\Gamma(Y,M),\mathrm R\Gamma(X,L))
\]
are defined as in 1.1, and $\langle u,v\rangle_z$ is the value at $z$ of the cup-product (III 4.2.7). We have $\mathrm R\Gamma(X,L)=\mathrm R\Gamma_c(X,L)=L_x$, $\mathrm R\Gamma(Y,M)=\mathrm R\Gamma_c(Y,M)=M_y$, and $u_z$ (resp. $v_z$) is nothing other than the direct image $u_*$ (resp. $v_*$) of $u$ (resp. $v$) on $S$ in the sense of (III 3.7.6), as is shown by the calculation procedure (III 3.5 b)).

On the other hand, if $h=h_1\times_Sh_2:(x,y)\to X\times_SY$ is the inclusion, and if $a':A'\to(x,y)$, $b':B'\to(x,y)$ denote the fibres of $a$, $b$ at $(x,y)$, each reduced to a point, then there are canonical isomorphisms
\[
h_*a'_*a'{}^!(DL_x\lten_{\Lambda}M_y)
 \xrightarrow{\sim} a_*a^!(DL\lten_{\Lambda}M),
\]
\[
h_*b'_*b'{}^!(L_x\lten_{\Lambda}DM_y)
 \xrightarrow{\sim} b_*b^!(L\lten_{\Lambda}DM),
\]
and the arrows
\[
h_*:\Hom(a_1'{}^*L_x,a_2'{}^!M_y)\to \Hom(a_1^*L,a_2^!M),
\]
\[
h_*:\Hom(b_2'{}^*M_y,b_1'{}^!L_x)\to \Hom(b_2^*M,b_1^!L)
\]
(III 3.7.6) are isomorphisms. The correspondences $u$ and $v$ are therefore written uniquely as
\begin{equation}\tag{1}
u=h_*u',\qquad v=h_*v',
\end{equation}
with
\[
u'\in\Hom(a_1'{}^*L_x,a_2'{}^!M_y),\qquad
v'\in\Hom(b_2'{}^*M_y,b_1'{}^!L_x).
\]
By the Lefschetz formula (III 4.5.1) for $(h_1,h_2,A'\to A,B'\to B)$, one has
\begin{equation}\tag{2}
\langle u,v\rangle_z=\langle u',v'\rangle_z .
\end{equation}
On the other hand, by (1), $u_*=u'_*$ and $v_*=v'_*$; hence, by the Lefschetz formula for $\{x\}\to S$, $\{y\}\to S$, one has
\begin{equation}\tag{3}
\langle u_*,v_*\rangle=\langle u',v'\rangle_z .
\end{equation}
In view of the preceding remark, the conclusion follows by comparing (2) and (3).

\begin{statement}{Lemma 2.4}
To prove 1.2, one may suppose that $a$ and $b$ are closed immersions, and it is enough to prove that, if $L_x=0$ and $M_y=0$, the restriction map
\begin{equation}
H_A^0(X\times_SY,DL\lten_SM)
 \longrightarrow H_z^0(B,b^*(DL\lten_SM))
\tag{2.4.1}
\end{equation}
is zero.
\end{statement}

Put $X\times_SY=Z$. The closed subschemes $a':A'\to Z$, $b':B'\to Z$, images of $a$ and $b$, satisfy the same hypotheses as $a$ and $b$. Let $u'$ (resp. $v'$) be the correspondence supported on $A'$ (resp. $B'$) which is the direct image of $u$ (resp. $v$) by the arrow deduced from the adjunction arrow $f_*f^!\to1$ (resp. $g_*g^!\to1$), where $f:A\to A'$ (resp. $g:B\to B'$) is the projection. It is immediate that $u_z=u'_z$ and $v_z=v'_z$. On the other hand, it follows easily from the definition of the cup-product 1.1.2 that
\[
\langle u,v\rangle=\langle u',v'\rangle,
\]
hence the first assertion.

By 2.2 and 2.3, one may suppose that $L_x=0$ and $M_y=0$, and the conclusion of 1.2 then becomes $\langle u,v\rangle=0$. Let $P,Q\in\Ob D_{\mathrm{ctf}}(Z)$. The cup-product
\[
H_A^0(Z,P)\otimes H_B^0(Z,Q)\longrightarrow H_z^0(Z,P\lten Q)
\]
decomposes as
\[
H_A^0(Z,P)\otimes H_B^0(Z,Q)
 \xrightarrow{\,b^*\otimes 1\,}
H_z^0(B,b^*P)\otimes H_B^0(Z,Q)
 \longrightarrow H_z^0(Z,P\lten Q),
\]
where $b^*$ is the restriction and the second arrow, a variant of the product considered in (SGA $4^{1/2}$ Cycle 1.2.1), is derived from the evident product at the level of sections of sheaves. This assertion is easily checked, for example by means of the flasque resolutions of (SGA 4 XVII 4.2). If $a':\{z\}\to B$, $b':\{z\}\to A$, $c:\{z\}\to Z$ are the inclusions, one may also invoke the standard decomposition of the Künneth arrow (III 4.2.1)
\[
a'{}^!P\lten_Z b'{}^!Q\longrightarrow c^!(P\lten Q)
\]
into a ``base-change arrow''
\[
a'{}^!P\lten_Z b'{}^!Q\longrightarrow a'{}^!b^*P\lten a^*b'{}^!Q,
\]
followed by ``projection arrows''
\[
a'{}^!b^*P\lten a^*b'{}^!Q
 \longrightarrow a'{}^!(b^*P\lten b'{}^!Q)
 \longrightarrow a'{}^!b'{}^!(P\lten Q)).
\]
Taking $P=DL\lten_SM$, $Q=L\lten_SDM$, and returning to the definition of $\langle u,v\rangle$ (III 4.2.5), one sees that it is enough to prove that $b^*u=0$, whence the conclusion.

In view of 2.4, 1.2 will therefore follow from the following more precise result.

\begin{statement}{Proposition 2.5}
Under the hypotheses of 1.2, suppose that $a$ and $b$ are closed immersions and that $L_x=0$, $M_y=0$. Let $T$ be the strict localization of $X\times_SY$ at $z=(x,y)$, and identify $A$, $B$ with closed subschemes of $T$. Then the restriction arrows
\[
b^*:\mathrm R\Gamma_A(T,(DL\lten_SM)|T)
 \longrightarrow \mathrm R\Gamma_z(B,b^*(DL\lten_SM)),
\]
\[
a^*:\mathrm R\Gamma_B(T,(L\lten_SDM)|T)
 \longrightarrow \mathrm R\Gamma_z(A,a^*(L\lten_SDM))
\]
are zero.
\end{statement}

\subsubsection*{2.6}
The proof of 2.5 will occupy the rest of number 2 and numbers 3 and 4. By symmetry of the statement, it is enough to prove the assertion concerning $b^*$. We shall rewrite it in a form which will be more convenient. First note that
\[
\mathrm R\Gamma(T,(DL\lten_SM)|T)
 =(DL\lten_SM)_z
 =(DL)_x\lten M_y=0,
\]
because $M_y=0$, and similarly
\[
\mathrm R\Gamma(B,b^*(DL\lten_SM))
 =(DL\lten_SM)_z=0.
\]
Consequently, the relative cohomology triangles give isomorphisms
\[
\mathrm R\Gamma(T-A,(DL\lten_SM)|T-A)
 \xrightarrow{\sim}
\mathrm R\Gamma_A(T,(DL\lten_SM)|T)[1],
\]
\[
\mathrm R\Gamma(B-z,b^*(DL\lten_{\Lambda}M)|B-z)
 \xrightarrow{\sim}
\mathrm R\Gamma_z(B,b^*(DL\lten_{\Lambda}M))[1],
\]
under which the arrow $b^*$ of 2.5, shifted by $1$, is identified with the restriction
\begin{equation}
b^*:\mathrm R\Gamma(T-A,(DL\lten_SM)|T-A)
 \longrightarrow \mathrm R\Gamma(B-z,(DL\lten_SM)|B-z).
\tag{2.6.1}
\end{equation}
Since $L_x=0$, $M_y=0$, one may write
\begin{equation}
L=i_!DP,\qquad M=j_!Q,
\tag{2.6.2}
\end{equation}
where $i:X-x\to X$, $j:Y-y\to Y$ are the inclusions, and $P\in\Ob D_{\mathrm{ctf}}(X-x)$, $Q\in\Ob D_{\mathrm{ctf}}(Y-y)$. By the duality isomorphism (SGA 4 XVIII 3.19.6), together with passage to the limit, one has
\[
DL=Ri_*DDP=Ri_*P,
\]
hence
\[
DL\lten_SM=Ri_*P\lten_Sj_!Q.
\]
Put $X\times_SY=Z$, and identify $X\times\{y\}$ and $\{x\}\times Y$ with closed subschemes of $Z$ (resp. $T$), which we shall simply denote by $X$ and $Y$. Let
\[
f':Z-Y\hookrightarrow Z,\qquad
g':Z-(X\cup Y)\hookrightarrow Z-Y,
\]
\[
f:T-Z\hookrightarrow T,\qquad
g:T-(X\cup Y)\hookrightarrow T-Y
\]
be the inclusions. The Künneth formula (III 1.6.4) gives, by passage to the limit,
\[
DL\lten_SM=Rf'_*g'_!(P\lten_SQ),
\]
and hence
\begin{equation}
(DL\lten_SM)|T=Rf_*g_!((P\lten_SQ)|T-(X\cup Y)).
\tag{2.6.3}
\end{equation}
With this identification, 2.6.1 becomes
\begin{equation}
b^*:\mathrm R\Gamma(T-(Y\cup A),g_!(E)|T-(Y\cup A))
 \longrightarrow \mathrm R\Gamma(B-z,g_!(E)|B-z),
\tag{2.6.4}
\end{equation}
where
\[
E=(P\lten_SQ)|T-(X\cup Y).
\]
In the next two numbers, we shall show that the arrow 2.6.4 is zero: first we shall reduce to the case where the cohomology sheaves of $E$ have tame ramification, more precisely where the inertia action occurs through an $\ell$-group, for a prime number $\ell\ne p$, and then we shall treat that case directly by means of a purity theorem.

\subsection*{3. Reduction to the tame case}

\subsubsection*{3.1}
Recall that the base ring $\Lambda$ is annihilated by an integer prime to $p$, hence is a direct product of noetherian rings $\Lambda_i$, each $\Lambda_i$ being annihilated by a power of a prime number $\ell_i\ne p$. For every scheme $S'$ over $S$, the derived category $D(S',\Lambda)$ is the product of the categories $D(S',\Lambda_i)$, so that, in order to prove the vanishing of 2.6.4, we may suppose that $\Lambda$ is annihilated by a power of a prime number $\ell\ne p$. We may also suppose that $P$ (resp. $Q$) is bounded and has locally free components of finite type.

\subsubsection*{3.2}
Put $X-x=U$, $Y-y=V$. Let $U''/U$ be a finite Galois covering with group $G$ which trivializes the components of $P$, let $H$ be an $\ell$-Sylow subgroup of $G$, let $U'=U''/H$ be the quotient covering, and let $X'$ be the normalized strictly local trait of $X$ in $U'$. Then $X'/X$ is a finite connected covering of degree $d_1$ prime to $\ell$, and the representation of $\pi_1(U')$ defined by the components of $P'=P|U'$ factors through a finite $\ell$-group. Proceeding similarly with $Q$, one finds a finite connected covering $Y'/Y$ of degree $d_2$ prime to $\ell$, inducing on $V$ an étale covering $V'$ such that the representation of $\pi_1(V')$ defined by the components of $Q'=Q|V'$ factors through a finite $\ell$-group. Let
\[
d=d_1e_1=d_2e_2=\operatorname{lcm}(d_1,d_2).
\]
After adjoining an $e_1$-th root (resp. an $e_2$-th root) of a uniformizer of $X'$ (resp. $Y'$), we may suppose that $X'$ (resp. $Y'$) has the same properties as above, and moreover that $d_1=d_2=d$.

Let $x'$ (resp. $y'$) be the closed point of $X'$ (resp. $Y'$), $T'$ the strict localization of $X'\times_S Y'$ at $(x',y')$, $z'$ the closed point of $T'$, $A'$ (resp. $B'$) the inverse image of $A$ (resp. $B$) in $T'$, and
\[
E'=(P\lten_{\Lambda}Q)|T'-(X'\cup Y')
   =(P'\lten_{\Lambda}Q')|T'-(X'\cup Y').
\]
We then have a commutative square of restriction arrows
\[
\begin{tikzcd}[column sep=huge,row sep=large]
\mathrm R\Gamma(T-(Y\cup A),g_!(E)|T-(Y\cup A))
 \arrow[r,"b^*"] \arrow[d]
&
\mathrm R\Gamma(B-z,g_!(E)|B-z) \arrow[d,"(1)"]\\
\mathrm R\Gamma(T'-(Y'\cup A'),g'_!(E')|T'-(Y'\cup A'))
 \arrow[r,"b'{}^*"]
&
\mathrm R\Gamma(B'-z',g'_!(E')|B'-z') ,
\end{tikzcd}
\]
where $b':B'\to T'$, $g':T'-(X'\cup Y')\to T'-Y'$ are the inclusions. The covering $T'/T$ is finite, of degree $d^2$, and étale outside $X\cup Y$, hence induces an étale covering of degree $d^2$ from $B'-z'$ onto $B-z$. The composite of (1) and the trace morphism
\[
\mathrm R\Gamma(B'-z',g'_!(E')|B'-z')
 \longrightarrow
\mathrm R\Gamma(B-z,g_!(E)|B-z)
\]
is therefore multiplication by $d^2$, which is an isomorphism, since $\ell$ does not divide $d$. To prove that the arrow 2.6.4 is zero, it is therefore enough to prove that $b'{}^*=0$. We are thus reduced to the case where the representation of $\pi_1(U)$ (resp. $\pi_1(V)$) defined by the components of $P$ (resp. $Q$) factors through a finite $\ell$-group, provided one ensures that $A'$, $B'$ satisfy the general-position hypotheses of 1.2. This will follow from the following lemma.

\begin{statement}{Lemma 3.3}
Under the hypotheses of 1.2, $a$ and $b$ being closed immersions, let $X'/X$, $Y'/Y$ be finite coverings of strictly local $S$-traits, of the same degree $d$. Let $A'$ (resp. $B'$) denote the inverse image of $A$ (resp. $B$) in $X'\times_S Y'$. Then the pairs $(A',B')$, $(A',X')$, $(Y',B')$ are each in general position in the sense of 1.2.\footnote{With the usual abuse of notation $X'=X'\times\{y'\}$, etc. (where $x'$ (resp.\ $y'$) is the closed point of $X'$ (resp.\ $Y'$).}
\end{statement}

The proof of 3.3 will be given after a few preliminaries.

\begin{statement}{Lemma 3.4}
Let $X\xleftarrow{s_1}W\xrightarrow{s_2}Y$ be finite morphisms of $S$-traits, let $s=(s_1,s_2):W\to X\times_SY$, and let $u,v,w$ be uniformizers of $X,Y,W$ respectively, so that
\[
s_1^*u=aw^m\pmod {w^{m+1}},\qquad
s_2^*v=bw^n\pmod {w^{n+1}},
\]
with $a,b\in k^*$. Then: (i) if $m<n$, $s(W)$ is tangent to $X$ at $z=(x,y)$; (ii) if $m=n$, $s(W)$ is tangent to the subscheme of $X\times_SY$ with equation
\[
ap_2^*(v)-bp_1^*(u)=0.
\]
\end{statement}

Let $Z=X\times_SY$, and let $W'$ be the reduced closed subscheme which is the image of $s$. In the tangent space to $Z$ at $z$, the slope, with respect to the tangent to $X$, of the tangent to $W'$ is the value at the origin of the function induced on $W'$ by $p_2^*(v)/p_1^*(u)$; this is also the value at the origin of the function induced by $p_2^*(v)/p_1^*(u)$ on $W$, and is therefore $0$ (resp. $b/a$) in case (i) (resp. (ii)).

In the situation of 3.4, we shall define the slope of $W$ with respect to $((X,u),(Y,v))$ as the element of $k\cup\{\infty\}$ equal to $0$ if $m<n$, to $\infty$ if $m>n$, and to $b/a$ if $m=n$. This element is independent of the choice of $w$.

\begin{statement}{Lemma 3.5}
Let $f:X'\to X$, $g:Y'\to Y$ be finite morphisms of $S$-traits with the same ramification index $d$, and let $u,v,u',v'$ be uniformizers of $X,Y,X',Y'$ respectively. Then there exists $C\in k^*$ such that, for every commutative diagram of finite morphisms of $S$-traits
\[
\begin{tikzcd}[column sep=large,row sep=large]
X' & W' \arrow[l,"s'_1"'] \arrow[r,"s'_2"] \arrow[d,"h"] & Y' \arrow[d,"g"]\\
X & W \arrow[l,"s_1"] \arrow[r,"s_2"'] & Y ,
\end{tikzcd}
\]
one has $\lambda=C{\lambda'}^d$, where $\lambda$ (resp. $\lambda'$) is the slope of $W$ (resp. $W'$) with respect to $((X,u),(Y,v))$ (resp. $((X',u'),(Y',v'))$).
\end{statement}

By hypothesis,
\[
f^*u=\alpha u'{}^d\pmod {u'{}^{d+1}},\qquad
g^*v=\beta v'{}^d\pmod {v'{}^{d+1}},
\]
with $\alpha,\beta\in k^*$. We show that $C=\beta/\alpha$ has the required properties. Let $w$ (resp. $w'$) be a uniformizer of $W$ (resp. $W'$). We have
\[
s_1^*u=aw^m\pmod {w^{m+1}},\qquad
s_2^*v=bw^n\pmod {w^{n+1}},
\]
\[
s'_1{}^*u'=a'w'{}^{m'}\pmod {w'{}^{m'+1}},\qquad
s'_2{}^*v'=b'w'{}^{n'}\pmod {w'{}^{n'+1}},
\]
\[
h^*w=\gamma w'{}^r\pmod {w'{}^{r+1}},
\]
with $a,b,a',b',\gamma\in k^*$. Commutativity of the squares in the diagram implies
\[
m'd=rm,\qquad n'd=rn,
\]
hence $m'/n'=m/n$. Suppose first that $m=n$. Then $\lambda=b/a$, $\lambda'=b'/a'$. If $N=m'd=rm=n'd=rn$, the preceding congruences give
\[
(fs'_1)^*u=\alpha a'{}^d w'{}^N\pmod {w'{}^{N+1}},\qquad
(s_1h)^*u=a\gamma^m w'{}^N\pmod {w'{}^{N+1}},
\]
whence $\alpha a'{}^d=a\gamma^m$. Similarly $\beta b'{}^d=b\gamma^m$, whence $\lambda=C{\lambda'}^d$, for $\lambda\ne0,\infty$. But, since $m'/n'=m/n$, $\lambda=0$ (resp. $\infty$) is equivalent to $\lambda'=0$ (resp. $\infty$), which completes the proof.

We now prove 3.3. Let $Z=X\times_SY$, $Z'=X'\times_SY'$, let $\widetilde A$ (resp. $\widetilde B$) be the normalization of $A$ (resp. $B$), and let $\widetilde A'$ (resp. $\widetilde B'$) be the normalization of $A'$ (resp. $B'$). Since $\widetilde A'$ (resp. $\widetilde B'$) is also the normalization of $\widetilde A\times_ZZ'$ (resp. $\widetilde B\times_ZZ'$), we have commutative squares
\begin{equation}\tag{*}
\begin{tikzcd}[column sep=large,row sep=large]
\widetilde A' \arrow[r] \arrow[d] & Z' \arrow[d] & \widetilde B' \arrow[l] \arrow[d]\\
\widetilde A \arrow[r] & Z & \widetilde B . \arrow[l]
\end{tikzcd}
\end{equation}
Let $(x_i)_{1\le i\le m}$ (resp. $(y_j)_{1\le j\le n}$) be the points of $\widetilde A$ (resp. $\widetilde B$) above $x$ (resp. $y$), and let $(x'_i)_{1\le i\le m'}$ (resp. $(y'_j)_{1\le j\le n'}$) be the points of $\widetilde A'$ (resp. $\widetilde B'$) above $x'$ (resp. $y'$). Since $A$ and $B$ are strictly local, we have decompositions
\[
\widetilde A=\coprod_{1\le i\le m}A_i,\qquad
\widetilde B=\coprod_{1\le j\le n}B_j,\qquad
\widetilde A'=\coprod_{1\le i\le m'}A'_i,\qquad
\widetilde B'=\coprod_{1\le j\le n'}B'_j,
\]
where $A_i,B_j,A'_i,B'_j$ are strictly local traits with respective closed points $x_i,y_j,x'_i,y'_j$. Choose uniformizers $u,v,u',v'$ of $X,Y,X',Y'$ respectively; let $\lambda_i$ (resp. $\mu_j$) be the slope of $A_i$ (resp. $B_j$) with respect to $((X,u),(Y,v))$, and let $\lambda'_i$ (resp. $\mu'_j$) be that of $A'_i$ (resp. $B'_j$) with respect to $((X',u'),(Y',v'))$. The general-position hypotheses on $A$ and $B$ are expressed by
\begin{equation}
\lambda_i\ne0\text{ for all }i,\qquad
\mu_j\ne\infty\text{ for all }j,\qquad
\lambda_i\ne\mu_j\text{ for all }i,j,
\tag{1}
\end{equation}
and we have to show that
\begin{equation}
\lambda'_i\ne0\text{ for all }i,\qquad
\mu'_j\ne\infty\text{ for all }j,\qquad
\lambda'_i\ne\mu'_j\text{ for all }i,j.
\tag{2}
\end{equation}
Let $i'\in[1,m']$, $j'\in[1,n']$, and let $x_i$ (resp. $y_j$) be the image of $x'_{i'}$ (resp. $y'_{j'}$). The preceding diagram provides commutative squares of finite morphisms of $S$-traits
\[
\begin{tikzcd}[column sep=large,row sep=large]
X' & A'_{i'} \arrow[l] \arrow[r] \arrow[d] & Y' \arrow[d]\\
X & A_i \arrow[l] \arrow[r] & Y
\end{tikzcd}
\qquad
\begin{tikzcd}[column sep=large,row sep=large]
X' & B'_{j'} \arrow[l] \arrow[r] \arrow[d] & Y' \arrow[d]\\
X & B_j \arrow[l] \arrow[r] & Y .
\end{tikzcd}
\]
By 3.5, there exists $C\in k^*$ such that
\[
\lambda_i=C{\lambda'_{i'}}^d,\qquad
\mu_j=C{\mu'_{j'}}^d.
\]
Thus (1) implies (2), which completes the proof of 3.3.

\subsection*{4. End of the proof of 1.2}

\subsubsection*{4.1}
To complete the proof of 2.5, hence of 1.2, it remains to establish the vanishing of 2.6.4 under the following additional hypotheses: $\Lambda$ is annihilated by a power of a prime number $\ell\ne p$, and $P$ and $Q$ are bounded, with locally free components of finite type, such that the corresponding representations of the fundamental groups factor through finite $\ell$-groups. We have
\[
g_!(E)|T-(Y\cup A)=g_{1!}(E_1),
\]
where $g_1:T-(X\cup Y\cup A)\hookrightarrow T-(Y\cup A)$ is the inclusion, and
\[
E_1=(P\lten_{\Lambda}Q)|T-(X\cup Y\cup A).
\]
Now $E_1$ is a complex with bounded, constructible, locally constant cohomology such that the corresponding representation of $\pi_1(T-(X\cup Y\cup A))$ factors through an $\ell$-group. The assertion to be proved will therefore follow from the following more general statement.

\begin{statement}{Proposition 4.2}
Let $T$ be the strict localization at a closed point of a smooth $S$-scheme of dimension $2$, let $t$ be the closed point of $T$, and let $A$, $B$, $X$ be closed subschemes of dimension $1$ of $T$. Let $i:T-(A\cup X)\hookrightarrow T-A$ be the inclusion. Suppose that $X$ is regular, and that $A$ is in general position with respect to $B$ and $X$; that is, the intersection of the tangent cones at $t$ to $A$ and $B\cup X$ is reduced to $t$. Let $E\in\Ob D^b(T-(A\cup X))$ be such that $H^*(E)$ is constructible and locally constant, and that the corresponding representation of $\pi_1(T-(A\cup X))$ factors through an $\ell$-group. Then the restriction arrow
\[
\mathrm R\Gamma(T-A,i_!(E))
\longrightarrow
\mathrm R\Gamma(B-t,i_!(E)|B-t)
\]
is zero.
\end{statement}

The proof will rely on the following lemma.

\begin{statement}{Lemma 4.3}
Let $Z$ be a connected regular noetherian scheme of dimension $2$, with residual characteristics prime to $\ell$, and let $X$, $Y$ be regular divisors meeting transversely at a point $z$, giving a diagram of inclusions
\[
\begin{tikzcd}[column sep=large,row sep=large]
Z-(X\cup Y) \arrow[r,"i'"] \arrow[d,"j"'] & Z-Y \arrow[d,"j"]\\
Z-X \arrow[r,"i"'] \arrow[d,"m"'] & Z \arrow[d,"m"]\\
Y-z \arrow[r,"i_1"'] & Y .
\end{tikzcd}
\]
Assume the purity hypothesis
\[
(P)\qquad R^qj_*(\mathbb Z/\ell)=0\quad\text{for }q>1.
\]
Let $L\in\Ob D^b(Z-(X\cup Y),\Lambda)$ have constructible locally constant cohomology, with the corresponding representation of $\pi_1(Z-(X\cup Y))$ factoring through an $\ell$-group. Then:
\begin{enumerate}[label=\alph*)]
\item There is a canonical isomorphism
\[
i_{1!}m'^*Rj_*i'_!L\xrightarrow{\sim} m^*Rj_*(i'_!L),
\]
and the complex $m'^*Rj_*i'_!L$ has constructible, locally constant cohomology.
\item Suppose that $Y$ is connected, and let $\bar s$ be a geometric point of $Y$ localized at a closed point $s\ne z$. Then the restriction arrow
\begin{equation}
\mathrm R\Gamma(Y,m^*Rj_*(i'_!L))\longrightarrow (Rj_*(i'_!L))_{\bar s}
\tag{4.3.1}
\end{equation}
is zero.
\item Under the hypotheses of b), let $C$ be a closed subscheme of dimension $1$ of $Z$ such that $Y\cap C=s$, and let $\widetilde C$ be the strict localization of $C$ at $\bar s$. Then the restriction arrow
\begin{equation}
\mathrm R\Gamma(Z-Y,i'_!L)
\longrightarrow
\mathrm R\Gamma(\widetilde C-\bar s,i'_!(L)|\widetilde C-\bar s)
\tag{4.3.2}
\end{equation}
is zero.
\end{enumerate}
\end{statement}

The trivial isomorphism
\[
i^*Rj_*i'_!L\xrightarrow{\sim}Rj_*(i'^*i'_!L)=Rj_*L
\]
gives an isomorphism
\[
m'^*Rj_*i'_!L\xrightarrow{\sim} i_1^*m^*Rj_*i'_!L,
\]
hence, by adjunction, an arrow
\[
i_{1!}m'^*Rj_*i'_!L\longrightarrow m^*Rj_*(i'_!L),
\]
which induces an isomorphism outside $z$. We show that it is an isomorphism, that is, that if $\bar z$ is a geometric point above $z$, then
\[
(Rj_*(i'_!L))_{\bar z}=0.
\]
For this, one may replace $\Lambda$ by $\mathbb Z/\ell^\nu$, where $\ell^\nu$ annihilates $\Lambda$, and by dévissage and passage to the limit reduce to the case where $L$ is a locally constant $\mathbb Z/\ell^\nu$-module of finite type corresponding to a representation of $\pi_1$ through an $\ell$-group. Since the only finite abelian group of $\ell$-torsion which is simple under the action of an $\ell$-group is $\mathbb Z/\ell$ with trivial action, a further dévissage reduces to the case where $L$ is the constant sheaf $\mathbb Z/\ell$.

Let $\widetilde Z$ be the strict localization of $Z$ at $\bar z$, $\widetilde X=X\times_Z\widetilde Z$, $\widetilde Y=Y\times_Z\widetilde Z$, and let $i':\widetilde Z-(\widetilde X\cup\widetilde Y)\hookrightarrow \widetilde Z-\widetilde Y$ be the inclusion. One has
\[
(Rj_*(i'_!(\mathbb Z/\ell)))_{\bar z}
=\mathrm R\Gamma(\widetilde Z-\widetilde Y,i'_!(\mathbb Z/\ell)),
\]
so it remains to prove that the restriction arrow
\[
H^n(\widetilde Z-\widetilde Y,\mathbb Z/\ell)
\longrightarrow
H^n(\widetilde X-\bar z,\mathbb Z/\ell)
\]
is an isomorphism for all $n$. For $n\le1$ this is a consequence of Abhyankar's lemma, and for $n>1$ both sides vanish by the purity hypothesis. The first assertion of a) is therefore proved. We prove the second. By dévissage, we may suppose $L$ concentrated in degree $0$. The hypothesis (P) implies, by dévissage, that $R^qj_*L=0$ for $q>2$. It remains to prove constructibility and local constancy of $m'^*R^qj_*i'_!L$ for $q\le1$, which is standard. For constructibility, resolving $L$ on the right by sheaves of the form $p_*p^*L$, where $p:Z'\to Z$ is a finite covering tamely ramified along $Y$, reduces to the case where $L$ is constant, of value $L_0$: then
\[
m'^*j_*i'_!L=(L_0)_{Y-z},
\qquad
m'^*R^1j_*i'_!L=(L_0)_{Y-z}(-1)
\]
by Abhyankar. For local constancy, one has to see that the specialization arrows are isomorphisms. One may, for this, replace $\Lambda$ by $\mathbb Z/\ell^\nu$ as above, and then reduce by dévissage to the case where $L$ is the constant sheaf $\mathbb Z/\ell$; the conclusion follows from Abhyankar's lemma.

We have therefore proved a). We prove b). Let $\bar\eta$ be a geometric generic point of $Y$. Choose specialization arrows
\[
a:\bar\eta\to Y(\bar z),\qquad b:\bar\eta\to Y(\bar s),
\]
where $Y(\bar z)$ (resp. $Y(\bar s)$) is the strict localization of $Y$ at $\bar z$ (resp. $\bar s$). Put
\[
M=m^*Rj_*(i'_!L).
\]
The commutative square
\[
\begin{tikzcd}
Y \arrow[r] & Y(\bar s)\\
Y(\bar z) \arrow[u] \arrow[r,"a"'] & \bar\eta \arrow[u,"b"']
\end{tikzcd}
\]
induces a commutative square
\[
\begin{tikzcd}[column sep=large,row sep=large]
\mathrm R\Gamma(Y,M) \arrow[r,"4.3.1"] \arrow[d] & M_{\bar s} \arrow[d,"b^*"]\\
M_{\bar z} \arrow[r,"a^*"'] & M_{\bar\eta}.
\end{tikzcd}
\]
Since $M|Y-z$ has locally constant cohomology, $b^*$ is an isomorphism. On the other hand, since $M=i_{1!}i_1^*M$, one has $M_{\bar z}=0$, hence the arrow 4.3.1 is zero. It remains to prove c). We have $\widetilde C=C\times_Z\widetilde Z$, where $\widetilde Z$ denotes, as above, the strict localization of $Z$ at $\bar z$. Put $\widetilde Y=Y\times_Z\widetilde Z$. Then there is a commutative diagram of restriction arrows
\[
\begin{tikzcd}[column sep=large,row sep=large]
\mathrm R\Gamma(Z-Y,i'_!L)=\mathrm R\Gamma(Z,Rj_*(i'_!L))
 \arrow[r,"4.3.2"] \arrow[d]
&
\mathrm R\Gamma(\widetilde C-\bar s,i'_!(L)|\widetilde C-\bar s)\\
\mathrm R\Gamma(Y,m^*Rj_*(i'_!L)) \arrow[r,"4.3.1"']
&
(Rj_*(i'_!L))_{\bar s}
  =\mathrm R\Gamma(\widetilde Z-\widetilde Y,(i'_!L)|\widetilde Z-\widetilde Y) \arrow[u]
\end{tikzcd}
\]
The vanishing of 4.3.2 therefore follows from b), completing the proof of 4.3.

We prove 4.2. Let $f:T'\to T$ be the blowing-up of the maximal ideal of $T$, let $D=f^{-1}(z)(\simeq\mathbb P^1_S)$ be the exceptional divisor, and let $A'$, $B'$, $X'$ be the respective strict transforms of $A$, $B$, $X$. Then $T'$ is regular, connected, of dimension $2$, $X'$ is a regular divisor meeting $D$ transversely at one point $z'$, and the general-position hypotheses on $A$, $B$, $X$ mean that $A'$ is disjoint from $B'\cup X'$. Let $E'$ be the inverse image of $E$ on $T'-(A'\cup D\cup X')$, and let
\[
i':T'-(A'\cup D\cup X')\hookrightarrow T'-(A'\cup D)
\]
be the inclusion. We must prove that the restriction arrow
\begin{equation}
\mathrm R\Gamma(T'-(A'\cup D),i'_!(E'))
\longrightarrow
\mathrm R\Gamma(B'\cup D,i'_!(E')|B'-(B'\cup D))
\tag{*}
\end{equation}
is zero. But $B'$ is a disjoint union of strictly local schemes $B'_j$ $(1\le j\le n)$, of dimension $1$, such that for every $j$, $B'_j\cap D$ is reduced to a point $b_j$. Applying 4.3 c) to $(Z=T'-A', X=X', Y=D-(D\cap A'), L=E', C=B'_j)$, the hypothesis (P) being verified by virtue of the relative purity theorem (SGA 4 XVI 3), since $(Z,Y)$ is a filtered projective limit of smooth pairs of codimension $1$ over $S$, with affine étale transition morphisms, one obtains that the restriction arrow
\[
\mathrm R\Gamma(T'-(A'\cup D),i'_!(E'))
\longrightarrow
\mathrm R\Gamma(B'_j-b_j,i'_!(E')|B'_j-b_j)
\]
is zero. Therefore (*) is zero, which completes the proof of 4.2, hence of 2.5 and 1.2.

\section*{II. Equivariant correspondences}

\subsection*{5. Non-commutative traces}

Throughout this number, fix a topos $T$ and a commutative unital ring $K$ of $T$. The rings under consideration will be assumed associative and unital, but not necessarily commutative. By a $K$-algebra we shall mean any (associative, unital) ring $A$ of $T$ endowed with a (unital) homomorphism $K\to A$ whose image is contained in the center.

\subsubsection*{5.0. Introduction}

Let $E$ be a $K$-module locally a direct factor of a free module of finite type. There is then a trace homomorphism, local in nature,
\begin{equation}\tag{5.0.1}
\Tr_K:\uHom_K(E,E)\longrightarrow K,
\end{equation}
defined by composing the evaluation arrow
\begin{equation}\tag{5.0.2}
\uHom_K(E,K)\otimes_KE\longrightarrow K,\qquad f\otimes x\longmapsto f(x),
\end{equation}
with the inverse of the canonical tensor-product isomorphism
\begin{equation}\tag{5.0.3}
\uHom_K(E,K)\otimes_KE\longrightarrow\uHom_K(E,E),\qquad f\otimes y\longmapsto (x\longmapsto f(x)y).
\end{equation}
The arrows 5.0.2 and 5.0.3 are defined more generally for every $K$-module $E$, and from them one obtains, for every complex of $K$-modules $E$, homomorphisms of complexes
\begin{align}
\uHom_K^\bullet(E,K)\otimes_KE&\longrightarrow K,\qquad f\otimes x\longmapsto f(x),\tag{5.0.4}\\
\uHom_K^\bullet(E,K)\otimes_KE&\longrightarrow \uHom_K^\bullet(E,E),\qquad f\otimes y\longmapsto (x\longmapsto f(x)y).\tag{5.0.5}
\end{align}
Upon deriving, 5.0.4 and 5.0.5 give, for $E\in\ob D^-(K)$, arrows in $D(K)$
\begin{align}
\uRHom_K(E,K)\lten_KE&\longrightarrow K,\tag{5.0.6}\\
\uRHom_K(E,K)\lten_KE&\longrightarrow\uRHom_K(E,E).\tag{5.0.7}
\end{align}
If $E$ is perfect (SGA 6 I), 5.0.7 is an isomorphism, as is readily checked by reduction to the case $E=K$, and, by composing the inverse of 5.0.7 with 5.0.6, one obtains a trace homomorphism, global in nature,
\begin{equation}\tag{5.0.8}
\Tr_K:\uRHom_K(E,E)\longrightarrow K,
\end{equation}
which generalizes 5.0.1 (for more details, see SGA 6 I 7, 8).

We propose here to examine to what extent, in the preceding constructions, $K$ can be replaced by a $K$-algebra $A$ (we shall be interested chiefly in the case where $A$ is the $K$-algebra of an ordinary finite group), and in particular to define a trace arrow analogous to 5.0.8, generalizing the notion of non-commutative trace introduced by Stallings in [5] and developed by Grothendieck in Expos\'e XI of this seminar. As was said in the introduction, that expos\'e, written by I. Bucur, was unfortunately lost; the reader may nevertheless refer to Deligne's expos\'e (SGA $4^{1/2}$, Rapport), where the essential points are indicated. Although replacing $K$ by $A$ in 5.0.3 presents no difficulty, the same is not true in 5.0.2: if $E$ is a left $A$-module, the evaluation arrow
\[
\uHom_A(E,A)\otimes_KE\longrightarrow A,
\qquad f\otimes x\longmapsto f(x),
\]
does not in general factor through $\uHom_A(E,A)\otimes_AE$, but, upon passage to the quotient, gives an arrow
\begin{equation}\tag{5.0.9}
\uHom_A(E,A)\otimes_AE\longrightarrow A_\natural,
\end{equation}
where $A_\natural$ is the quotient $K$-module of $A$ by the sub-$K$-module locally generated by local sections of the form $ab-ba$. If $E$ is locally a direct factor of a free left $A$-module of finite type, the tensor-product arrow, analogous to 5.0.3,
\begin{equation}\tag{5.0.10}
\uHom_A(E,A)\otimes_AE\longrightarrow\uHom_A(E,E)
\end{equation}
is an isomorphism, and, by composing the inverse of 5.0.10 with 5.0.9, one obtains a trace homomorphism
\begin{equation}\tag{5.0.11}
\Tr_A:\uHom_A(E,E)\longrightarrow A_\natural,
\end{equation}
which reduces to 5.0.1 when $A=K$. One also checks readily (see 5.7) that, when $T$ is punctual, 5.0.11 coincides with the trace homomorphism of Stallings and Grothendieck. In this number we shall examine the extension to complexes and the derivation of the arrows 5.0.9, 5.0.10, 5.0.11, and the behavior of the resulting arrows with respect to extension and restriction of scalars and to certain external tensor products. In fact, we shall derive arrows slightly more general than 5.0.9 and 5.0.10, so that the constructions may be applied, in no.~6, to the definition of ``non-commutative'' Lefschetz--Verdier local terms.

\subsubsection*{5.1. Notation}

If $A$ is a $K$-algebra, we denote by $A^\circ$ the algebra opposite to $A$, and by $A^e=A\otimes_KA^\circ$ the enveloping algebra of $A$. Giving an $A$-bimodule, a left $A^e$-module, or a right $A^e$-module amounts to the same thing. In particular, $A$ is an $A$-bimodule, and hence both a left and a right $A^e$-module. If $E$ is a bimodule, recall ([1] IX 4) that
\[
H_0(A,E)=A\otimes_{A^e}E=E\otimes_{A^e}A
\]
is the quotient $K$-module of $E$ by the sub-$K$-module locally generated by local sections of the form $ax-xa$, for $x\in E$, $a\in A$\footnote{If $M$ is a sheaf, the notation $x\in M$ means that $x$ is a local section of $M$.}. We shall sometimes write $E_\natural$ instead of $H_0(A,E)$. In particular, $A_\natural$ is the $K$-module defined in 5.0.9. If $A=K[G]$, where $G$ is an ordinary finite group, $A_\natural$ is simply the free $K$-module on the set $G_\natural$ of conjugacy classes of $G$.

If $A$, $B$ are $K$-algebras, we shall use the Cartan--Eilenberg notation ${}_AE$, $E_B$, ${}_AE_B$ to denote respectively a left $A$-module, a right $B$-module, and a module that is left over $A$ and right over $B$. Finally, we shall denote by $D(A)$ (resp. $D(A,B)$) the derived category of the category of left $A$-modules (resp. of $A$-$B$-bimodules, left over $A$ and right over $B$).

\subsubsection*{5.2}

Let $A$, $B$ be $K$-algebras. Recall that, in the situation $({}_AE,{}_BF_A,{}_BG)$, there is a canonical isomorphism
\begin{equation}\tag{5.2.1}
\uHom_A(E,\uHom_B(F,G))\xrightarrow{\sim}\uHom_B(F\otimes_AE,G),\qquad
f\longmapsto (y\otimes x\longmapsto f(x)(y)).
\end{equation}
This derives to an isomorphism
\begin{equation}\tag{5.2.2}
\uRHom_A(E,\uRHom_B(F,G))\xrightarrow{\sim}\uRHom_B(F\lten_AE,G)
\end{equation}
for $E\in\ob D^-(A)$, $F\in\ob D(B,A)$, $G\in\ob D^+(B)$. Indeed, one may suppose that $E$ is bounded above with flat terms and that $G$ is bounded below with injective terms; one then has $\uRHom_B(F\lten_AE,G)=\Hom_B^\bullet(F\otimes_AE,G)$, $\uRHom_B(F,G)=\Hom_B^\bullet(F,G)$, and $\uRHom_A(E,\Hom_B^\bullet(F,G))=\Hom_A^\bullet(E,\Hom_B^\bullet(F,G))$, since, if $E$ is acyclic, $F\otimes_AE$ is acyclic, and hence $\Hom_A^\bullet(E,\Hom_B^\bullet(F,G))\simeq\Hom_B^\bullet(F\otimes_AE,G)$ is acyclic. For $E$ and $G$ as above, 5.2.2 is given by the isomorphism of complexes
\[
\Hom_A^\bullet(E,\Hom_B^\bullet(F,G))\xrightarrow{\sim}\Hom_B^\bullet(F\otimes_AE,G)
\]
defined by 5.2.1. Note that, if $A$ and $B$ are flat over $K$, 5.2.2 still holds under the conditions $E\in\ob D(A)$, $F\in\ob D^-(B,A)$, $G\in\ob D^+(B)$, as is seen by resolving $F$ by flat $(B\otimes_KA^\circ)$-modules. From 5.2.2 one deduces an isomorphism
\begin{equation}\tag{5.2.3}
\uHom_A(E,\uRHom_B(F,G))\xrightarrow{\sim}\uHom_B(F\lten_AE,G).
\end{equation}

\subsubsection*{5.3}

Let $A$ be a $K$-algebra. In the situation $(E_A,{}_AF)$, there is a canonical isomorphism
\begin{equation}\tag{5.3.1}
E\otimes_AF\xrightarrow{\sim}(E\otimes_KF)\otimes_{A^e}A,
\qquad x\otimes y\longmapsto (x\otimes y)\otimes 1.
\end{equation}
On the other hand, in the situation $({}_AE,{}_AF_A)$, the evaluation arrow
\begin{equation}\tag{5.3.2}
\uHom_A(E,F)\otimes_KE\longrightarrow F,
\qquad f\otimes x\longmapsto f(x)
\end{equation}
is $A^e$-linear, and therefore induces, via 5.3.1, upon applying $-\otimes_{A^e}A$, an arrow
\begin{equation}\tag{5.3.3}
\uHom_A(E,F)\otimes_AE\longrightarrow F\otimes_{A^e}A\;(=F_\natural).
\end{equation}
In other words, the composite of 5.3.2 and the canonical arrow $F\to F_\natural$, $x\mapsto x\otimes 1$, factors uniquely through 5.3.3.

Suppose that $A$ is flat over $K$. Then 5.3.1 derives to an isomorphism
\begin{equation}\tag{5.3.4}
E\lten_AF\xrightarrow{\sim}(E\lten_KF)\lten_{A^e}A
\end{equation}
for $E\in\ob D^-(A^\circ)$, $F\in\ob D^-(A)$ (resolve $E$ and $F$ by flat modules). On the other hand, 5.3.3 derives to an arrow of $D(K)$, called the evaluation arrow,
\begin{equation}\tag{5.3.5}
\uRHom_A(E,F)\lten_AE\longrightarrow F\lten_{A^e}A,
\end{equation}
for $E\in\ob D^-(A)$, $F\in\ob D^b(A^e)$ such that $\uRHom_A(E,F)\in\ob D^b(A^\circ)$. One first derives 5.3.2 to an arrow of $D^-((A^e)^\circ)$,
\begin{equation}\tag{5.3.6}
\uRHom_A(E,F)\lten_KE\longrightarrow F,
\end{equation}
by resolving $E$ by flat modules and $F$ by injective $A^e$-modules, and then one applies the functor $-\lten_{A^e}A$ to 5.3.6, taking 5.3.4 into account.

\subsubsection*{5.4}

Let $A$ be a $K$-algebra. In the situation $({}_AE,{}_AF_A,{}_AG)$, there is a canonical arrow
\begin{equation}\tag{5.4.1}
\uHom_A(E,F)\otimes_AG\longrightarrow\uHom_A(E,F\otimes_AG),
\qquad u\otimes y\longmapsto(x\longmapsto u(x)\otimes y),
\end{equation}
which corresponds by adjunction (5.2.1) to the $A$-linear arrow
\[
E\otimes_K\uHom_A(E,F)\otimes_AG\longrightarrow F\otimes_AG
\]
defined by the evaluation arrow $E\otimes_K\uHom_A(E,F)\to F$. The arrow 5.4.1 derives to an arrow of $D(K)$, called the tensor-product arrow,
\begin{equation}\tag{5.4.2}
\uRHom_A(E,F)\lten_AG\longrightarrow\uRHom_A(E,F\lten_AG),
\end{equation}
for $E\in\ob D^-(A)$, $F\in\ob D^+(A^e)$, $G\in\ob D^-(A)$, with $F\lten_AG\in\ob D^+(A)$. If $E$ and $G$ are bounded above with flat terms, and $F$ is bounded below with injective terms, the first member of 5.4.2 identifies with $\Hom_A^\bullet(E,F)\otimes_AG$, and it is sent to the second member by the composite of 5.4.1 and the natural arrow
\[
\Hom_A^\bullet(E,F\otimes_AG)\longrightarrow \uRHom_A(E,F\otimes_AG)=\uRHom_A(E,F\lten_AG).
\]
The arrow 5.4.2 is not an isomorphism in general. It is, however, an isomorphism in certain important cases, for example if $E$ is perfect (as is checked by reducing to $E=A$). We shall see another interesting example later (6.2.3).

The arrows 5.3.5 and 5.4.2 will play an essential role throughout the rest of the expos\'e. The examples of objects $F$ that we have in mind are: a) $F=A$, or $F=A$ regarded as a right $A$-module in the natural way and as a left $A$-module through a homomorphism of $K$-algebras $\varphi:A\to A$ (cf. 5.13); b) $F=M\lten_KA$, with $M\in\ob D^+(K)$; in no.~6, $M$ will be a dualizing complex, see especially 6.3.2, 6.4.1, 6.5.6.

\subsubsection*{5.5}

Denote by $DF(A)$ the derived category of complexes of left $A$-modules endowed with a finite filtration ([3] V). Then, in the situation of 5.2, there is a canonical isomorphism in $DF(K)$, analogous to 5.2.2, for $E\in\ob D^-F(A)$, $F\in\ob D^-F(B,A)$ ($=D^-F(B\otimes_KA^\circ)$), $G\in\ob D^+F(B)$. Similarly, there are ``filtered derived'' arrows analogous to 5.3.4, 5.3.5, 5.3.6, 5.4.2. These arrows are compatible with passage to associated graded objects, in the sense of ([3] V 1,2).

\subsubsection*{5.6}

Assume that $A$ is flat over $K$. Denote by $D(A)_{\mathrm{parf}}$ the full subcategory of $D(A)$ formed by complexes of finite perfect amplitude, that is, perfect complexes of finite tor-dimension (SGA 6 I). We define a canonical arrow, called the trace,
\begin{equation}\tag{5.6.1}
\Tr_A=\uRHom_A(E,F\lten_AE)\longrightarrow F\lten_{A^e}A
\end{equation}
for $E\in\ob D(A)_{\mathrm{parf}}$, $F\in\ob D^b(A^e)$, by composing the inverse isomorphism of 5.4.2 for $G=E$ with the evaluation arrow 5.3.5. From this one deduces a trace arrow
\begin{equation}\tag{5.6.2}
\Tr_A:\uHom_A(E,F\lten_AE)\longrightarrow H^0(T,F\lten_{A^e}A).
\end{equation}
For $F=A$, by composing with the canonical arrow $A\lten_{A^e}A\to A\otimes_{A^e}A=A_\natural$ (5.1), one obtains arrows, again called traces,
\begin{align}
\Tr_A:\uRHom_A(E,E)&\longrightarrow A_\natural,\tag{5.6.3}\\
\Tr_A:\uHom_A(E,E)&\longrightarrow H^0(T,A_\natural).\tag{5.6.4}
\end{align}
Notice that 5.6.4 is a rather innocuous arrow, local in nature, whereas 5.6.3 is a ``serious'' arrow, global in nature.

Denote by $DF(A)_{\mathrm{parf}}$ the full subcategory of $DF(A)$ formed by filtered complexes whose associated graded object has finite perfect amplitude. Then there is an arrow in $DF(K)$ analogous to 5.6.1 for $E\in\ob DF(A)_{\mathrm{parf}}$, $F\in\ob D^bF(A^e)$, whence arrows analogous to 5.6.2, 5.6.3, 5.6.4. It follows from the remarks in 5.5 that the filtered arrow 5.6.1 is compatible with passage to the associated graded object (cf. [3] V 3.7.2), and one deduces, as in loc. cit., that, if the filtration on $F$ is concentrated in degree $0$, then, for $u\in\Hom_{DF(A)}(E,F\lten_AE)$, one has
\begin{equation}\tag{5.6.5}
\Tr_A(u)=\sum_i\Tr_A(\operatorname{gr}^i u).
\end{equation}
This formula implies in particular that, for $E\in\ob D(A)_{\mathrm{parf}}$, $F\in\ob D^b(A^e)$, $u\in\uHom_A(E,F\lten_AE)$, one has
\begin{equation}\tag{5.6.6}
\Tr(u[1])=-\Tr(u),
\end{equation}
where $u[1]:E[1]\to(F\lten_AE)[1]$ is the translate of $u$.

\subsubsection*{5.7}

The traces defined in 5.6 simultaneously generalize the usual commutative traces ((SGA 6 I), ([3] V 3)) and the non-commutative traces of Stallings [5] and Grothendieck.

(a) Suppose that $A=K$. Then $A^e=A_\natural=K$, the arrow 5.6.1 is written
\[
\Tr_K:\uRHom_K(E,F\lten_KE)\longrightarrow F
\]
and coincides with the trace arrow defined in ([3] V 3.7.1), in the special case where the complex $K$ of loc. cit. is reduced to $A$. In particular,
\[
\Tr_K:\uRHom_K(E,E)\longrightarrow K
\]
coincides with the trace defined in (SGA 6 I 8).

(b) Assume again that $A$ is a flat $K$-algebra, not necessarily commutative. We propose to make 5.6.4 explicit for $E$ strictly perfect, that is, bounded with terms locally direct factors of free modules of finite type. First, take $E=A$; $\uRHom_A(A,A)=\uHom_A(A,A)$ identifies with $A$ by $f\mapsto f(1)$, $a\mapsto(x\mapsto xa)$, and 5.6.3 identifies with the canonical projection $A\to A_\natural$, which we shall denote by $a\mapsto a_\natural$; thus, if $d_a$ denotes right multiplication by $a$,
\begin{equation}\tag{5.7.1}
\Tr_A(d_a)=a_\natural.
\end{equation}
For $1\le i\le n$, let $E_i$ be a copy of $A$, and let $E=\bigoplus_{1\le i\le n}E_i$. The isomorphism 5.4.2, for $F=A$, $G=E$, is written
\[
\bigoplus_{1\le i,j\le n}\uHom_A(E_i,A)\otimes_AE_j\longrightarrow\uHom_A(E,E),
\]
and, for $u=\sum u_{ij}\otimes 1\in\uHom_A(E,E)$ $(u_{ij}\otimes 1\in\uHom_A(E_i,A)\otimes_A E_j)$, one has
\begin{equation}\tag{5.7.2}
\Tr_A(u)=\sum_i (u_{ii})_\natural,
\end{equation}
by definition of the evaluation arrow 5.3.5. Now take for $E$ a direct factor of $A^n$, the image of a projector $p$, and denote by $i:E\to A^n$ the inclusion; then, for $u\in\uHom_A(E,E)$, one has
\begin{equation}\tag{5.7.3}
\Tr_A(u)=\Tr_A(iup),
\end{equation}
as is seen either with the aid of 5.6.5 or directly from the definition. Finally, suppose that $E$ is strictly perfect; if $u\in\Hom_A(E,E)$ is a homomorphism of complexes, one has, by 5.6.5 and 5.6.6,
\begin{equation}\tag{5.7.4}
\Tr_A(u)=\sum_i(-1)^i\Tr_A(u^i).
\end{equation}
The fact that 5.6.4 is obtained, by applying $H^0$, from a homomorphism of complexes $\uHom_A^\bullet(E,E)\to A$ (5.6.3) implies that, if $u$, $v\in\Hom_A(E,E)$ are homotopic, then
\begin{equation}\tag{5.7.5}
\Tr_A(u)=\Tr_A(v).
\end{equation}
Since every endomorphism of $E$ in $D(A)$ is locally defined by an endomorphism of complexes, locally unique up to homotopy (SGA 6 I 4.10), the formulas 5.7.2 to 5.7.4 therefore permit the calculation of $\Tr_A(u)$ for every endomorphism $u$ of $E$ in $D(A)$ and, combined with 5.7.5, provide a direct definition of $\Tr_A(u)$. When the topos $T$ is punctual, this recovers the definition of (SGA $4^{1/2}$, Rapport 4.3).

\subsubsection*{5.8}

Let $P,Q$ be $A$-bimodules. Then $P\otimes_AQ$ and $Q\otimes_AP$ are $A$-bimodules, and there are canonical isomorphisms
\begin{equation}\tag{5.8.1}
(P\otimes_AQ)\otimes_{A^e}A\xrightarrow{\sim}P\otimes_{A^e}Q\xrightarrow{\sim}(Q\otimes_AP)\otimes_{A^e}A\xrightarrow{\sim}A\otimes_{A^e}(P\otimes_KQ)\otimes_{A^e}A,
\end{equation}
each of the objects above identifying canonically with the quotient $K$-module of $P\otimes_KQ$ by the submodule generated by sections of the form $axb\otimes y-x\otimes bya$, for $x\in P$, $y\in Q$, $a,b\in A$. For $A$ flat over $K$, these isomorphisms derive to isomorphisms
\begin{equation}\tag{5.8.2}
(P\lten_AQ)\lten_{A^e}A\simeq P\lten_{A^e}Q\simeq(Q\lten_AP)\lten_{A^e}A\simeq A\lten_{A^e}(P\lten_KQ)\lten_{A^e}A
\end{equation}
for $P,Q\in\ob D^-(A^e)$. On the other hand, the composition arrow at the level of modules
\[
\uHom_A(E,P\otimes_AF)\otimes_K\uHom_A(F,Q\otimes_AE)
\longrightarrow\uHom_A(E,P\otimes_AQ\otimes_AE),
\]
\[
f\otimes g\longmapsto (P\otimes_Ag)f,
\]
derives to
\begin{equation}\tag{5.8.3}
\uRHom_A(E,P\lten_AF)\lten_K\uRHom_A(F,Q\lten_AE)
\longrightarrow\uRHom_A(E,P\lten_AQ\lten_AE)
\end{equation}
for $P,Q\in\ob D^b(A^e)$, $E,F\in\ob D^b(A)$ such that the tensor products and $\uRHom$'s appearing above belong to $D^b$. When $E$ and $F$ are perfect, this arrow identifies, via 5.4.2, with the arrow deduced from the evaluation arrow $F\lten_K\uRHom_A(F,Q)\to Q$; from this one readily deduces that there is a commutative square
\begin{equation}\tag{5.8.4}
\begin{tikzcd}[column sep=large,row sep=large]
\uRHom_A(E,P\lten_AF)\lten_K\uRHom_A(F,Q\lten_AE) \arrow[r,"(1)"] \arrow[d,"(2)"'] & \uRHom_A(E,P\lten_AQ\lten_AE) \arrow[d,"\Tr_A"]\\
\uRHom_A(F,Q\lten_AP\lten_AF) \arrow[r,"\Tr_A"'] & P\lten_{A^e}Q,
\end{tikzcd}
\end{equation}
where (1) and (2) are given by 5.8.3. The point is that the two composite arrows in 5.8.4 are defined by the tensor product of the evaluation arrows
\[
F\lten_K\uRHom_A(F,Q)\longrightarrow Q,\qquad
E\lten_K\uRHom_A(E,P)\longrightarrow P.
\]
It follows that, for $u\in\uHom_A(E,P\lten_AF)$, $v\in\uHom_A(F,Q\lten_AE)$, one has
\begin{equation}\tag{5.8.5}
\Tr_A(vu)=\Tr_A(uv).
\end{equation}
We shall sometimes denote the common value in 5.8.5 by $\langle u,v\rangle_A$.

\subsubsection*{5.9. Traces and extension of scalars}

Let $A\to B$ be a morphism of flat $K$-algebras. It induces a morphism of $K$-algebras $A^e\to B^e$, hence an extension-of-scalars functor $-\lten_{A^e}B^e:D(A^e)\longrightarrow D(B^e)$.
For $E,F$ as in 5.3.5, if $F\lten_{A^e}B^e\in\ob D^b(B^e)$ and $\RHom_B(B\lten_AE,F\lten_{A^e}B^e)\in\ob D^b((B^e)^\circ)$, there is a commutative square in $D(K)$
\begin{equation}\tag{5.9.1}
\begin{tikzcd}[column sep=large,row sep=large]
\RHom_A(E,F)\lten_{A^e}E \arrow[r,"(1)"] \arrow[d,"(3)"'] & F\lten_{A^e}A \arrow[d,"(4)"]\\
\RHom_B(B\lten_AE,F\lten_{A^e}B^e)\lten_{B^e}(B\lten_AE) \arrow[r,"(2)"'] & (F\lten_{A^e}B^e)\lten_{B^e}B,
\end{tikzcd}
\end{equation}
where the arrows (1) and (2) are defined by 5.3.6, (4) by $A\to B$, and (3) by the canonical arrow $F\to F\lten_{A^e}B^e$. The verification is immediate. Similarly, for $E,F,G$ as in 5.4.2, if $F\lten_{A^e}B^e\in\ob D^+(B^e)$ and $(F\lten_{A^e}B^e)\lten_B(B\lten_AG)\in\ob D^+(B)$, one has a commutative square in $D(K)$
\begin{equation}\tag{5.9.2}
\begin{tikzcd}[column sep=large,row sep=large]
\RHom_A(E,F)\lten_AG \arrow[r,"(1)"] \arrow[d,"(3)"'] & \RHom_A(E,F\lten_AG) \arrow[d,"(4)"]\\
\RHom_B(B\lten_AE,F\lten_{A^e}B^e)\lten_B(B\lten_AG) \arrow[r,"(2)"'] & \RHom_B(B\lten_AE,(F\lten_{A^e}B^e)\lten_B(B\lten_AG)),
\end{tikzcd}
\end{equation}
where (1) and (2) are defined by 5.4.2, and (3) and (4) by $F\to F\lten_{A^e}B^e$. From this one deduces, for $E,F$ as in 5.6.1, with $F\lten_{A^e}B^e\in\ob D^b(B^e)$, a commutative square in $D(K)$, analogous to the preceding ones, in which the vertical arrows are defined by 5.6.1:
\begin{equation}\tag{5.9.3}
\begin{tikzcd}[column sep=large,row sep=large]
\RHom_A(E,F\lten_AE) \arrow[r,"\Tr_A"] \arrow[d] & F\lten_{A^e}A \arrow[d]\\
\RHom_B(B\lten_AE,(F\lten_{A^e}B^e)\lten_B(B\lten_AE)) \arrow[r,"\Tr_B"'] & (F\lten_{A^e}B^e)\lten_{B^e}B.
\end{tikzcd}
\end{equation}
Finally, if one is given $F'\in\ob D^b(B^e)$ such that $\RHom_A(E,F')\in\ob D^b((B^\circ))$, and $f:F\lten_{A^e}B^e\to F'$, then there is a compatibility slightly more general than 5.9.1, expressed by an analogous commutative square, with $F'$ in place of $F\lten_{A^e}B^e$ in the lower row and with the vertical arrows defined by $f$. One has similarly commutative squares analogous to 5.9.2 and 5.9.3 with $F'$ in place of $F\lten_{A^e}B^e$. Taking in particular $F=A$ and, for $f:A\lten_{A^e}B^e\to B$, the arrow defined by $A\to B$, one obtains a commutative square
\begin{equation}\tag{5.9.4}
\begin{tikzcd}[column sep=large,row sep=large]
\RHom_A(E,E) \arrow[r,"\Tr_A"] \arrow[d] & A\lten_{A^e}A \arrow[d]\\
\RHom_B(B\lten_AE,B\lten_AE) \arrow[r,"\Tr_B"'] & B\lten_{B^e}B,
\end{tikzcd}
\end{equation}
where the left vertical arrow, respectively the right vertical arrow, is given by extension of scalars, respectively is the canonical arrow deduced from $A\to B$.

\subsubsection*{5.10. Traces and restriction of scalars}

Let $A\to B$ be a morphism of flat $K$-algebras such that $B$, as a left $A$-module, is flat and of finite presentation, i.e. locally a direct factor of a free $A$-module of finite type. This is the case for example if $A\to B$ is the morphism $K[H]\to K[G]$ defined by an injective morphism of discrete groups $H\hookrightarrow G$ of finite index. The scalar-restriction homomorphism
\[
\Hom_B(B,B)\longrightarrow\Hom_A(B,B),
\]
where $B$ is considered as a left module, identifies, by the canonical isomorphisms
\[
B\xrightarrow{\sim}\Hom_B(B,B),\qquad \Hom_A(B,A)\otimes_AB\xrightarrow{\sim}\Hom_A(B,B),
\]
with an arrow of $B$-bimodules
\begin{equation}\tag{5.10.1}
B\longrightarrow\Hom_A(B,A)\otimes_AB.
\end{equation}
The composite of 5.10.1 with the evaluation arrow $\Hom_A(B,A)\otimes_AB\to A_\natural$ (5.3.3) is the homomorphism $f\mapsto\Tr_A(f)$ ($=\Tr_A(d_f)$, where $d_f$ is right multiplication by $f$); hence, by the formula $\Tr_A(f_g)=\Tr_A(gf)$ (5.8.5), it factors through $B_\natural$ as an arrow again denoted $\Tr_{B/A}:B_\natural\to A_\natural$. Thus there is a commutative square
\begin{equation}\tag{5.10.2}
\begin{tikzcd}[column sep=large,row sep=large]
B \arrow[r,"\Tr_B"] \arrow[d,"5.10.1"'] & B_\natural \arrow[d,"\Tr_{B/A}"]\\
\Hom_A(B,A)\otimes_AB \arrow[r,"\mathrm{ev}"'] & A_\natural .
\end{tikzcd}
\end{equation}
In the situation $({}_BE,{}^BP_A)$, one has a homomorphism of right $B$-modules
\begin{equation}\tag{5.10.3}
\Hom_B(E,P)\otimes_AB\longrightarrow\Hom_B(E,P\otimes_AB),\qquad
u\otimes b\longmapsto(x\longmapsto u(x)\otimes b),
\end{equation}
which is an isomorphism, as is checked immediately: it is enough to see that the underlying $K$-linear homomorphism is an isomorphism, and this follows from the fact that $B$ is locally a direct factor of a free $A$-module of finite type. This isomorphism derives to an isomorphism of $D(B)$
\begin{equation}\tag{5.10.4}
\RHom_B(E,P)\lten_AB\xrightarrow{\sim}\RHom_B(E,P\lten_AB),
\end{equation}
for $E\in\ob D^-(B)$, $P\in\ob D^+(B\otimes_KA^\circ)$: resolve $E$ by modules of the form $B\otimes_AL$, where $L$ is flat, and $P$ by injective $(B\otimes_KA^\circ)$-modules, which are necessarily injective as $B$-modules, since $A$ is flat over $K$; this allows one to remove the $R$ on the left, and the resolution chosen for $E$ also allows one to remove it on the right.

For $E\in\ob D^-(B)$, $M\in\ob D^+(K)$, $F\in\ob D^-(B)$, define an arrow in $D(K)$\footnote{In no.~6 we shall apply the constructions that follow to the case where $M$ is a dualizing complex (6.7).}
\begin{equation}\tag{5.10.5}
\RHom_B(E,B\lten_KM)\lten_BF\longrightarrow\RHom_A(E,A\lten_KM)\lten_AF
\end{equation}
as the composite of the arrows
\[
\begin{tikzcd}[column sep=small]
\RHom_B(E,B\lten_KM)\lten_BF \arrow[r,"(1)"] & \RHom_B(E,\Hom_A(B,A)\lten_A(B\lten_KM))\lten_BF \\
\arrow[r,"(2)"] & \RHom_B(E,\Hom_A(B,A)\lten_KM)\lten_A B\lten_BF \\
\arrow[r,"(3)"] & \RHom_B(E,\Hom_A(B,A\lten_KM))\lten_AF \\
\arrow[r,"(4)"] & \RHom_A(E,A\lten_KM)\lten_AF,
\end{tikzcd}
\]
where (1) is defined by 5.10.1, (2) by the inverse isomorphism of 5.10.4 with $P=\Hom_A(B,A)\lten_KM$, (3) by the canonical isomorphism
\[
\Hom_A(B,A)\lten_KM\xrightarrow{\sim}\Hom_A(B,A\lten_KM)
\]
(5.4.2), and (4) by the adjunction isomorphism 5.2.2.

\begin{statement}{Proposition 5.10.6}
Suppose that $M\lten_KF\in\ob D^+(B)$. Then there is a commutative square in $D(K)$
\begin{equation}\tag{5.10.7}
\begin{tikzcd}[column sep=large,row sep=large]
\RHom_B(E,B\lten_KM)\lten_BF \arrow[r,"5.4.2"] \arrow[d,"5.10.5"'] & \RHom_B(E,M\lten_KF) \arrow[d,"\mathrm{res}"]\\
\RHom_A(E,A\lten_KM)\lten_AF \arrow[r,"5.4.2"'] & \RHom_A(E,M\lten_KF),
\end{tikzcd}
\end{equation}
where the right vertical arrow is restriction of scalars.
\end{statement}

One has to verify that there is a commutative square
\[
\begin{tikzcd}[column sep=large,row sep=large]
E\lten_K\RHom_B(E,B\lten_KM)\lten_BF \arrow[r] \arrow[d] & M\lten_KF \arrow[d,"\mathrm{res}"]\\
E\lten_K\RHom_A(E,A\lten_KM)\lten_AF \arrow[r] & M\lten_KF,
\end{tikzcd}
\]
where the horizontal arrows are given by the evaluations (5.3.6). This square is deduced, by applying $-\lten_BF$, from the analogous square with $F=B$, so that one may suppose $F=B$. The commutativity of 5.10.7 in this case is equivalent to that of the triangle
\[
\begin{tikzcd}[column sep=huge,row sep=large]
\RHom_B(E,B\lten_KM) \arrow[d,"(1)"'] \arrow[dr,"(2)"] & \\
\RHom_B(E,\Hom_A(B,B\lten_KM)) \arrow[r,"(3)"'] & \RHom_A(E,B\lten_KM),
\end{tikzcd}
\]
where (1) is defined by 5.10.1, (2) is restriction of scalars, and (3) is the adjunction isomorphism. But the arrow $B\lten_KM\to\Hom_A(B,B\lten_KM)$ given by 5.10.1 is nothing other than the adjunction arrow corresponding to the functors restriction of scalars and $\Hom_A(B,-)$; the commutativity follows at once.

Let $E\in\ob D^-(B)$, $M\in\ob D^b(K)$ be such that $\RHom_B(E,B\lten_KM)\in\ob D^b(B^\circ)$. Then, by 5.3.5, there is an evaluation arrow
\[
\RHom_B(E,B\lten_KM)\lten_{B^e}E\longrightarrow M\lten_K(B\lten_{B^e}B),
\]
which, by composition with the canonical arrow $B\lten_{B^e}B\to B_\natural$, gives an arrow
\begin{equation}\tag{5.10.8}
\RHom_B(E,B\lten_KM)\lten_BE\longrightarrow M\lten_KB_\natural.
\end{equation}

\begin{statement}{Proposition 5.10.9}
Suppose that $A$ is of finite presentation and flat as a $K$-module, and that $A_\natural$, $B_\natural$ are flat over $K$. Let $E\in\ob D^-(B)$, $M\in\ob D^b(K)$ be such that $\RHom_B(E,B\lten_KM)\in\ob D^b(B^\circ)$. Then $\RHom_A(E,A\lten_KM)\in\ob D^b(A^\circ)$, and there is a commutative square
\begin{equation}\tag{5.10.10}
\begin{tikzcd}[column sep=large,row sep=large]
\RHom_B(E,B\lten_KM)\lten_BE \arrow[r,"5.10.8"] \arrow[d,"5.10.5"'] & M\lten_KB_\natural \arrow[d,"M\otimes_K\Tr_{B/A}"]\\
\RHom_A(E,A\lten_KM)\lten_AE \arrow[r,"5.10.8"'] & M\lten_KA_\natural .
\end{tikzcd}
\end{equation}
\end{statement}

The additional hypotheses on $A$ and $B$ (verified, for example, if $A$ and $B$ are $K$-algebras of finite groups) are no doubt useless, and one should have a commutative square analogous to 5.10.10 with $A_\natural$ (resp. $B_\natural$) replaced by $A\lten_{A^e}A$ (resp. $B\lten_{B^e}B$).

The first assertion of 5.10.9 is trivial. We prove the second. We may suppose that $M$ has injective components and that $E$ has components of the form $B\otimes_KL$, with $L$ flat. Then
\[
\RHom_B(E,B\lten_KM)\lten_BE=\Hom_B^\bullet(E,B\otimes_KM)\otimes_BE,
\]
\[
\RHom_A(E,A\lten_KM)\lten_AE=\Hom_A^\bullet(E,A\otimes_KM)\otimes_AE,
\]
and the arrows in 5.10.10 are defined at the level of complexes. Thus we are reduced to verifying that, if $L,M$ are $K$-modules and $E$ a left $B$-module, there is a commutative square
\[
\begin{tikzcd}[column sep=large,row sep=large]
|[alias=star]| {(*)} & \Hom_B(B\otimes_KL,B\otimes_KM)\otimes_B(B\otimes_KL) \arrow[r,"5.3.3"] \arrow[d,"(1)"'] & M\otimes_KB_\natural \arrow[d,"M\otimes_K\Tr_{B/A}"]\\
& \Hom_A(B\otimes_KL,A\otimes_KM)\otimes_A(B\otimes_KL) \arrow[r,"5.3.3"'] & M\otimes_KA_\natural,
\end{tikzcd}
\]
where the arrow (1) is defined by sending $B\otimes_KM$ into $\Hom_A(B,A)\otimes_AB\otimes_KM$ by 5.10.1 and composing with an isomorphism of the type 5.10.3 and an adjunction isomorphism. Since $A$ is flat and of finite presentation as a $K$-module, the ``tensor-product'' arrows
\[
\text{(2)}\quad B\otimes_K\Hom_K(L,M)=\Hom_B(B,B)\otimes_K\Hom_K(L,M)\to\Hom_B(B\otimes_KL,B\otimes_KM),
\]
\[
\text{(3)}\quad \Hom_A(B,A)\otimes_K\Hom_K(L,M)\to\Hom_A(B\otimes_KL,A\otimes_KM)
\]
are isomorphisms, by which (1) identifies with the tensor product (over $K$) of 5.10.1 by $\Hom_K(L,M)\otimes_KL$. On the other hand, the upper (resp. lower) horizontal arrow of (*) identifies by (2) (resp. (3)) with the tensor product (over $K$) of $\Tr_B:B\to B_\natural$ (resp. $\mathrm{ev}:\Hom_A(B,A)\otimes_AB\to A_\natural$) by the evaluation arrow $\Hom_K(L,M)\otimes_KL\to M$. The commutativity therefore follows from that of 5.10.2, completing the proof of 5.10.9.

\begin{statement}{Corollary 5.10.11}
Under the hypotheses of 5.10.9, for $E\in\ob D(B)_{\mathrm{parf}}$, $M\in\ob D^b(K)$, there is a commutative square in $D(K)$
\begin{equation}\tag{5.10.12}
\begin{tikzcd}[column sep=large,row sep=large]
\RHom_B(E,M\lten_KE) \arrow[r,"\Tr_B"] \arrow[d] & M\lten_KB_\natural \arrow[d,"M\lten\Tr_{B/A}"]\\
\RHom_A(E,M\lten_KE) \arrow[r,"\Tr_A"'] & M\lten_KA_\natural,
\end{tikzcd}
\end{equation}
where the left vertical arrow is restriction of scalars.
\end{statement}

It is enough indeed to conjugate 5.10.6 and 5.10.9.

\subsubsection*{5.11. Application to finite groups: divisibility properties}

Let $G$ be a finite group, with neutral element denoted $e$, and put $K[G]=B$. The $K$-module $B_\natural$ is free with basis the set $G_\natural$ of conjugacy classes of $G$; every section $s$ of $B_\natural$ is written uniquely
\[
s=\sum_{x\in G_\natural}s(x)x.
\]

\begin{statement}{Proposition 5.11.1}
Cf. SGA $4^{1/2}$, Rapport 4.5. Let $E\in\ob D(B)_{\mathrm{parf}}$, $M\in\ob D^b(K)$. For $u\in\Hom_B(E,M\lten_KE)$, one has, in $H^0(T,M)$,
\begin{equation}\tag{*}
\Tr_K(u)=|G|\footnote{If $X$ is a set, $|X|$ denotes the cardinality of $X$.}\,\Tr_B(u)(e).
\end{equation}
\end{statement}
It is enough to apply 5.10.11 with $A=K$, noting that $\Tr_{B/K}:B_\natural\to K$ is given by
\begin{equation}\tag{5.11.2}
\Tr_{B/K}(g)=
\begin{cases}
|G|,& g=e,\\
0,& g\ne e.
\end{cases}
\end{equation}

\begin{statement}{Proposition 5.11.3}
Under the hypotheses of 5.11.1, let $g\in G$, let $Z_g$ be the centralizer of $g$ in $G$, and put $A=K[Z_g]$. Let $u\in\Hom_B(E,M\lten_KE)$, and let $gu\in\Hom_A(E,M\lten_KE)$ be the composite of the image of $u$ in $\Hom_A(E,M\lten_KE)$ with left multiplication by $g$. Then, in $H^0(T,M)$,
\[
\Tr_A(gu)(e)=\Tr_B(u)(g^{-1}).
\]
\end{statement}

We shall show more generally that there is a commutative square
\begin{equation}\tag{5.11.4}
\begin{tikzcd}[column sep=3.0em,row sep=2.2em]
\RHom_B(E,M\lten_KE) \arrow[r,"\Tr_B"] \arrow[d] &
M\lten_KB_\natural \arrow[r,"{x\mapsto x(g^{-1})}"] &
M \arrow[d,equal] \\
\RHom_A(E,M\lten_KE) \arrow[r,"\Tr_A"'] &
M\lten_KA_\natural \arrow[r,"{x\mapsto x(e)}"'] &
M,
\end{tikzcd}
\end{equation}
where the left vertical arrow is the composite of restriction of scalars and left multiplication by $g$. This square is the boundary of the diagram
\[
\begin{tikzcd}[column sep=2.8em,row sep=2.0em]
\RHom_B(E,M\lten_KE) \arrow[r,"\Tr_B"] \arrow[d,"\mathrm{rest}"'] &
M\lten_KB_\natural \arrow[r,"{x\mapsto x(g^{-1})}"] \arrow[d,"{\underset{(1)}{\Tr_{B/A}}}" description] &
M \arrow[d,equal] \\
\RHom_A(E,M\lten_KE) \arrow[r,"\Tr_A"'] \arrow[d,"{\RHom_A(E,s_g)}"'] &
M\lten_KA_\natural \arrow[d,"{\underset{(2)}{s_g}}" description] \arrow[dr,phantom,"\scriptstyle(3)" description] &
M \\
\RHom_A(E,M\lten_KE) \arrow[r,"\Tr_A"'] &
M\lten_KA_\natural \arrow[r,"{x\mapsto x(e)}"'] &
M,
\end{tikzcd}
\]
where $\mathrm{rest}$ denotes restriction of scalars and $s_g$ left multiplication by $g$. Square (1) commutes by 5.10.11, and (2) by the definition of $\Tr_A$. It therefore remains to prove the commutativity of (3), for which it is enough to show that, for $h\in G$,
\[
\Tr_A(x\mapsto gxh)(e)=\text{the value at }g^{-1}\text{ of the image of }h\text{ in }B,
\]
i.e.
\begin{equation}\tag{*}
\Tr_A(x\mapsto gxh)(e)=
\begin{cases}
1,& g^{-1}\text{ is conjugate to }h,\\
0,& \text{otherwise}.
\end{cases}
\end{equation}
For this one may replace $T$ by the punctual topos and $K$ by $\mathbb Z$, and it is then enough to prove (*) after multiplication by $|Z_g|$. In view of 5.11.1, we are therefore reduced to showing that
\[
\Tr_{\mathbb Z}(x\mapsto gxh)=
\begin{cases}
|Z_g|,& g^{-1}\text{ is conjugate to }h,\\
0,& \text{otherwise}.
\end{cases}
\]
But this follows from the fact that the left hand side is the number of fixed points of the map $x\mapsto gxh$ from $G$ to $G$, whence the commutativity of 5.11.4 and the conclusion of 5.11.3.

\begin{statement}{Corollary 5.11.5}
Under the hypotheses of 5.11.3, one has
\begin{equation}\tag{5.11.5}
\Tr_K(gu)=|Z_g|\,\Tr_B(u)(g^{-1})=|Z_g|\,\Tr_A(gu)(e).
\end{equation}
\end{statement}

\subsubsection*{5.12. Traces and external tensor products}

Let, for $i=1,2$, $A_i$ be a flat $K$-algebra, and put $A=A_1\otimes_KA_2$. Let $E_i\in\ob D^-(A_i)$, $F_i\in\ob D^b(A_i^e)$ be such that $\RHom_{A_i}(E_i,F_i)\in\ob D^b(A_i^\circ)$, let $F=F_1\lten_KF_2\in\ob D^b(A^e)$, and suppose $\RHom_A(E,F)\in\ob D^b(A^\circ)$, where $E=E_1\lten_KE_2$. Then there is a commutative square
\begin{equation}\tag{5.12.1}
\begin{tikzcd}[column sep=large,row sep=large]
(\RHom_{A_1}(E_1,F_1)\lten_{A_1}E_1)\lten_K(\RHom_{A_2}(E_2,F_2)\lten_{A_2}E_2) \arrow[r,"(1)"] \arrow[d,"(2)"'] & (F_1\lten_{A_1^e}A_1)\lten_K(F_2\lten_{A_2^e}A_2) \arrow[d,"(3)"]\\
\RHom_A(E,F)\lten_AE \arrow[r] & F\lten_{A^e}A,
\end{tikzcd}
\end{equation}
where (1) is the tensor product of the arrows 5.3.5, (2) is defined by the tensor-product arrow
\begin{equation}\tag{5.12.2}
\RHom_{A_1}(E_1,F_1)\lten_K\RHom_{A_2}(E_2,F_2)\longrightarrow\RHom_A(E,F),
\end{equation}
and (3) is the evident canonical isomorphism. The verification is essentially trivial. One checks in the same way that, for $E_i\in\ob D^-(A_i)$, $F_i\in\ob D^+(A_i^e)$, $G_i\in\ob D^-(A_i)$, with $F_i\lten_{A_i}G_i\in\ob D^b(A_i)$, $\RHom_{A_i}(E_i,F_i)\in\ob D^b(A_i^\circ)$, $\RHom_{A_i}(E_i,F_i\lten_{A_i}G_i)\in\ob D^b(K)$, $F\lten_AG\in\ob D^b(A)$, where $G=G_1\lten_KG_2$, there is a commutative square
\begin{equation}\tag{5.12.3}
{\scriptsize
\begin{tikzcd}[column sep=2.8em,row sep=large]
\begin{gathered}
(\RHom_{A_1}(E_1,F_1)\lten_{A_1}G_1)\lten_K{}\\[-.15em]
(\RHom_{A_2}(E_2,F_2)\lten_{A_2}G_2)
\end{gathered}
\arrow[r,"(1)"] \arrow[d,"(2)"']
&
\begin{gathered}
\RHom_{A_1}(E_1,F_1\lten_{A_1}G_1)\lten_K{}\\[-.15em]
\RHom_{A_2}(E_2,F_2\lten_{A_2}G_2)
\end{gathered}
\arrow[d,"(3)"]\\
\RHom_A(E,F)\lten_AG \arrow[r,"5.4.2"'] & \RHom_A(E,F\lten_AG),
\end{tikzcd}
}
\end{equation}
where (2) and (3) are tensor-product arrows of the type 5.12.2, and (1) is the tensor product of the arrows 5.4.2. From the commutativity of the squares 5.12.1 and 5.12.3 it follows that, for $E_i\in\ob D(A_i)_{\mathrm{parf}}$, $F_i\in\ob D^b(A_i^e)$ such that $F=F_1\lten_KF_2\in\ob D^b(A^e)$, there is a commutative square
\begin{equation}\tag{5.12.4}
\begin{tikzcd}[column sep=large,row sep=large]
\RHom_{A_1}(E_1,F_1\lten_{A_1}E_1)\lten_K\RHom_{A_2}(E_2,F_2\lten_{A_2}E_2) \arrow[r,"(1)"] \arrow[d,"(2)"'] & (F_1\lten_{A_1^e}A_1)\lten_K(F_2\lten_{A_2^e}A_2) \arrow[d,"(3)"]\\
\RHom_A(E,F\lten_AE) \arrow[r,"5.6.1"'] & F\lten_{A^e}A,
\end{tikzcd}
\end{equation}
where (1) is the tensor product of the arrows $\Tr_{A_i}$ (5.6.1), (2) is defined by 5.12.2, and (3) is the canonical isomorphism appearing in 5.12.1. Consequently, for $u_i\in\Hom_{A_i}(E_i,F_i\lten_{A_i}E_i)$, one has
\begin{equation}\tag{5.12.5}
\Tr_A(u_1\lten_Ku_2)=\Tr_{A_1}(u_1)\Tr_{A_2}(u_2),
\end{equation}
the product on the right being defined by the canonical arrow
\[
H^0(T,F_1\lten_{A_1^e}A_1)\otimes H^0(T,F_2\lten_{A_2^e}A_2)
\longrightarrow H^0(T,F\lten_{A^e}A).
\]

\subsubsection*{5.13. Traces of semilinear endomorphisms}

Let $A$ be a $K$-algebra and $\varphi$ an endomorphism of $A$. If $E$, $F$ are left $A$-modules, recall that an additive homomorphism $u:E\to F$ is called $\varphi$-$A$-linear if $u(ax)=\varphi(a)u(x)$ for all $a\in A$, $x\in E$, i.e. if $u$ is an $A$-linear homomorphism from $E$ to the $A$-module ${}_{(\varphi)}F$ deduced from $F$ by restriction of scalars along $\varphi$. If $F$ is an $A$-bimodule and $G$ a left $A$-module, there is an evident isomorphism
\[
{}_{(\varphi)}F\otimes_AG\xrightarrow{\sim}{}_{(\varphi)}(F\otimes_AG),
\]
which derives trivially to
\[
{}_{(\varphi)}F\lten_AG\xrightarrow{\sim}{}_{(\varphi)}(F\lten_AG).
\]
Assuming $A$ flat over $K$, one deduces from 5.4.2 the arrow
\begin{equation}\tag{5.13.1}
\RHom_A(E,{}_{(\varphi)}F)\lten_AG\longrightarrow\RHom_A(E,{}_{(\varphi)}(F\lten_AG))
\end{equation}
for $E\in\ob D^-(A)$, $F\in\ob D^+(A^e)$, $G\in\ob D^-(A)$, with $F\lten_AG\in\ob D^+(A)$. For $E\in\ob D(A)_{\mathrm{parf}}$, 5.13.1 is an isomorphism, which allows one to define, by 5.3.5, a trace arrow
\begin{equation}\tag{5.13.2}
\Tr_{A,\varphi}:\RHom_A(E,{}_{(\varphi)}(F\lten_AE))\longrightarrow {}_{(\varphi)}F\lten_{A^e}A.
\end{equation}
In particular, for $F=A$, one obtains
\begin{equation}\tag{5.13.3}
\Tr_{A,\varphi}:\Hom_A(E,{}_{(\varphi)}E)\longrightarrow H^0(T,A_{\natural,\varphi}),
\end{equation}
where $A_{\natural,\varphi}={}_{(\varphi)}A\otimes_{A^e}A$ is the quotient $K$-module of $A$ by the submodule locally generated by local sections of the form $ab-\varphi(b)a$. In the case where $A=K[G]$, $G$ is a discrete group, and $\varphi$ is defined by an endomorphism $\varphi$ of $G$, $A_{\natural,\varphi}$ is the free $K$-module with basis $G_{\natural,\varphi}$, where $G_{\natural,\varphi}$ is the quotient of $G$ by the action $(g,x)\mapsto\varphi(g)xg^{-1}$; the orbits are sometimes called $\varphi$-conjugacy classes of $G$. The traces 5.13.2 coincide, for $T$ punctual, with those considered by Grothendieck and discussed in Exposé XII.

In the situation of 5.10, if $\varphi$, $\psi$ are respective endomorphisms of $A$, $B$, compatible with $A\to B$, there is a unique arrow
\[
\Tr_A:B_{\natural,\psi}\longrightarrow A_{\natural,\varphi}
\]
making the square, analogous to 5.10.2,
\begin{equation}\tag{5.13.4}
\begin{tikzcd}[column sep=large,row sep=large]
{}_{(\psi)}B=\Hom_B(B,{}_{(\psi)}B) \arrow[r,"\Tr_{B,\psi}"] \arrow[d] & B_{\natural,\psi} \arrow[d,"\Tr_A"]\\
\Hom_A(B,{}_{(\varphi)}A)\otimes_AB \arrow[r,"\mathrm{ev}"'] & A_{\natural,\varphi},
\end{tikzcd}
\end{equation}
commute, where $\mathrm{ev}$ is the evaluation arrow, and the left vertical arrow is restriction of scalars. One defines an arrow analogous to 5.10.5, with $A\otimes_KM$ (resp. $B\otimes_KM$) replaced by ${}_{(\varphi)}A\otimes_KM$ (resp. ${}_{(\psi)}B\otimes_KM$), and there are compatibilities analogous to 5.10.6, 5.10.9, 5.10.11, which we leave to the reader to state.

Let $G$ be a finite group and $\varphi$ an endomorphism of $G$. For $g\in G$, let $Z_{g,\varphi}$ be the $\varphi$-centralizer of $g$ in $G$, i.e. the set of $x\in G$ such that $\varphi(x)gx^{-1}=g$. Put $A=K[G]$, and let $E\in\ob D(A)_{\mathrm{parf}}$, $M\in\ob D^b(K)$. The preceding compatibilities easily imply that, for $u\in\Hom_A(E,{}_{(\varphi)}(M\lten_KE))$ and $g\in G$,
\begin{equation}\tag{5.13.5}
\Tr_K(gu)=|Z_{g^{-1},\varphi}|\,\Tr_{A,\varphi}(u)(g^{-1}),
\end{equation}
the details of the verification being left to the reader.

\section*{6. Equivariant Correspondences and Divisibility of Local Terms}

\subsection*{6.1. Notation}

In this number, $S$ denotes the spectrum of a field of characteristic $p$, and $\Lambda$ a noetherian commutative ring annihilated by an integer not divisible by $p$. The $S$-schemes considered will be assumed separated and of finite type. The $\Lambda$-algebras considered will be associative, unital, and of finite type, but not necessarily commutative. Most often, and in particular for everything concerning “non-commutative Lefschetz--Verdier local terms,” they will be assumed flat; in applications, they will be algebras of finite groups.

If $X$ is an $S$-scheme and $A$ a $\Lambda$-algebra, we write $D({}_AX)$ (resp. $D(X_A)$) for the derived category of the category of sheaves of left (resp. right) $A$-modules on $X$. For $A$ commutative, we shall also write $D(X,A)$ for $D({}_AX)$, and $D(X)=D(X,\Lambda)$. We denote by $D_c(-)$ (resp. $D_{\mathrm{ctf}}(-)$) the full subcategory formed by complexes with constructible cohomology (resp. finite tor-dimension and constructible cohomology). If $A$, $B$ are $\Lambda$-algebras, we write $D({}_AX_B)$ for the derived category of sheaves of $(A,B)$-bimodules, left over $A$ and right over $B$. Finally, we use the notations $A^o$, $A^e$, $A_{\natural}$ of 5.1.

\subsection*{6.2}
The formalism of (III 1, 2, 3) extends without difficulty to the “non-commutative” case. We briefly indicate the main points of this generalization.

Let $X$ be an $S$-scheme and let $A$, $B$ be $\Lambda$-algebras. If $L\in\Ob D_{\mathrm{ctf}}(X_A)$ and $M\in\Ob D({}_AX_B)$, with $M$ of finite tor-dimension and constructible cohomology as an object of $D(X_B)$, then
\[
L\lten_A M\in\Ob D(X_B)_{\mathrm{ctf}}
\]
as is immediate. On the other hand, it follows from (SGA $4^{1/2}$ Finiteness 1.6) that, for $M$ as above and $L\in\Ob D_{\mathrm{ctf}}({}_AX)$, one has
\[
\RHom_A(L,M)\in\Ob D_{\mathrm{ctf}}(X_B).
\]
It also follows from loc. cit. that, if $f$ is an $S$-morphism, $D_{\mathrm{ctf}}({}_A-)$ is stable under the operations $(f^*,f_*,f_!,f^!)$, the case of $f^*$ being trivial.

Let, for $i=1,2$, $f_i:X_i\to Y_i$ be an $S$-morphism,
\[
f=f_1\times_S f_2:X=X_1\times_SX_2\longrightarrow Y=Y_1\times_SY_2,
\]
and let $A$ be a $\Lambda$-algebra. For $L_1\in\Ob D^-(X_{1A})$, $L_2\in\Ob D^-({}_AX_2)$, there is a Künneth arrow in $D^-(Y)$
\begin{equation}
f_{1!}L_1\lten_{S,A}f_{2!}L_2
\longrightarrow
f_*(L_1\lten_{S,A}L_2),
\tag{6.2.1}
\end{equation}
where
\[
L_1\lten_{S,A}L_2
\stackrel{\mathrm{dfn}}{=}p_1^*L_1\lten_Ap_2^*L_2.
\]
As in (III 1.6), one deduces that it is an isomorphism. If $B$, $C$ are $\Lambda$-algebras, with $B$ or $C$ assumed flat over $\Lambda$ in order to define
\[
\lten_A:D^-({}_B{-}_A)\times D^-({}_A{-}_C)\longrightarrow D^-({}_B{-}_C),
\]
there is an isomorphism analogous to 6.2.1 in $D^-({}_BY_C)$, for $L_1\in\Ob D^-({}_BX_{1A})$, $L_2\in\Ob D^-({}_AX_{2B})$. One checks similarly that, for $M_1\in\Ob D^-(Y_{1A})$, $M_2\in\Ob D^-({}_AY_2)$, there is a Künneth isomorphism in $D^-(X)$
\begin{equation}
f^!M_1\lten_{S,A}f^!M_2
\longrightarrow
f^!(M_1\lten_{S,A}M_2),
\tag{6.2.2}
\end{equation}
and an analogous isomorphism in $D^-({}_B X_C)$ for $M_1\in\Ob D^-({}_BY_{1A})$, $M_2\in\Ob D^-({}_AY_{2C})$.

Finally, let, for $i=1,2$, $A_i$ be a $\Lambda$-algebra, $A_2$ being flat, $E_1\in\Ob D_{\mathrm{ctf}}({}_{A_1}X_1)$, $F_2\in\Ob D^+({}_{A_2}X_2)$, $F_1\in\Ob D({}_{A_1}X_{1A_2})$, with $F_1$ of finite tor-dimension and constructible cohomology as an object of $D(X_{1A_2})$. Then
\[
\RHom_{A_1}(E_1,F_1)\in\Ob D_{\mathrm{ctf}}(X_{1A_2}),
\]
as seen above, and there is a canonical arrow in $D^+(X)$ ($D^+(X,A_2)$ when $A_2$ is commutative)
\begin{equation}
\RHom_{A_1}(E_1,F_1)\lten_{S,A_2}F_2
\longrightarrow
\RHom_{A_1}(p_1^*E_1,F_1\lten_{S,A_2}F_2),
\tag{6.2.3}
\end{equation}
defined by composing the “inverse-image” arrow (III 2.1.1)
\[
p_1^*\RHom_{A_1}(E_1,F_1)\lten_{A_2}p_2^*F_2
\longrightarrow
\RHom_{A_1}(p_1^*E_1,p_1^*F_1)\lten_{A_2}p_2^*F_2
\]
with the arrow analogous to 5.4.2
\[
\RHom_{A_1}(p_1^*E_1,p_1^*F_1)\lten_{A_2}p_2^*F_2
\longrightarrow
\RHom_{A_1}(p_1^*E_1,p_1^*F_1\lten_{A_2}p_2^*F_2),
\]
adjoint, by 5.2.3, to the arrow
\[
p_1^*E_1\lten\RHom_{A_1}(p_1^*E_1,p_1^*F_1)\lten_{A_2}p_2^*F_2
\longrightarrow
p_1^*F_1\lten_{A_2}p_2^*F_2
\]
defined by the evaluation arrow 5.3.6. One deduces from 6.2.1, as in (III 2.3), that 6.2.3 is an isomorphism. In particular, for $A_2=\Lambda$, $F_2=\Lambda_{X_2}$, 6.2.3 gives an isomorphism
\begin{equation}
p_1^*\RHom_{A_1}(E_1,F_1)
\xrightarrow{\sim}
\RHom_{A_1}(p_1^*E_1,p_1^*F_1).
\tag{6.2.4}
\end{equation}

\subsection*{6.3}
If $t:T\to S$ is an $S$-scheme, and $A$ is a $\Lambda$-algebra, we put
\begin{equation}
t_A^! = KA_T \qquad(\in\Ob D^+({}_AT_A)).
\tag{6.3.1}
\end{equation}
For $A=\Lambda$, we write $K_T$ instead of $K\Lambda_T$. It follows from 6.2.2 that there is a canonical isomorphism in $D^+({}_AT_A)$
\begin{equation}
K_T\lten A\xrightarrow{\sim}KA_T.
\tag{6.3.2}
\end{equation}
We denote by $D_A$ the functor
\[
\RHom_A(-,KA_T).
\]
It sends $D_{\mathrm{ctf}}({}_AT)$ into $D_{\mathrm{ctf}}(T_A)$ and conversely, and the proof of (SGA $4^{1/2}$ Finiteness Theorem 4) shows that, for $L\in\Ob D_{\mathrm{ctf}}({}_AT)$, the canonical arrow
\[
L\longrightarrow D_AD_AL
\]
is an isomorphism. For $A=\Lambda$, we write $D$ instead of $D_A$.

\subsection*{6.4}
Let $A$ be a flat $\Lambda$-algebra, and let, for $i=1,2$, $X_i$ be an $S$-scheme, $X=X_1\times_SX_2$, $p_i:X\to X_i$ the projections, and $L_i\in\Ob D_{\mathrm{ctf}}({}_AX_i)$. By 6.2.3 applied to $A_1=A_2=A$, $E_1=L_1$, $F_2=L_2$, $E_2=KA_{X_1}$, there is an isomorphism
\[
D_A L_1\lten_{S,A}L_2
\xrightarrow{\sim}
\RHom_A(p_1^*L_1,KA_{X_1}\lten_{S,A}L_2),
\]
whence, by composing with the Künneth isomorphism 6.2.2
\[
KA_{X_1}\lten_{S,A}L_2\xrightarrow{\sim}p_2^!L_2,
\]
an isomorphism
\begin{equation}
D_AL_1\lten_{S,A}L_2
\xrightarrow{\sim}
\RHom_A(p_1^*L_1,p_2^!L_2),
\tag{6.4.1}
\end{equation}
which generalizes (III 3.1.1).

Let $c:C\to X$ be an $S$-morphism, and put $c_i=p_ic$. By the induction formula (SGA 4 XVIII 3.1.12.2), there is a canonical isomorphism, analogous to (III 3.2.1),
\begin{equation}
c^!\RHom_A(p_1^*L_1,p_2^!L_2)
\xrightarrow{\sim}
\RHom_A(c_1^*L_1,c_2^!L_2),
\tag{6.4.2}
\end{equation}
and hence, by composing with 6.4.1, an isomorphism
\begin{equation}
H^0(X,c_*c^!(D_AL_1\lten_{S,A}L_2))
\xrightarrow{\sim}
\Hom_A(c_1^*L_1,c_2^!L_2).
\tag{6.4.3}
\end{equation}
The elements of $\Hom_A(c_1^*L_1,c_2^!L_2)$ will be called cohomological correspondences from $L_1$ to $L_2$ with support in $c$. The terminology differs slightly from that of loc. cit.; the two terminologies coincide when $c$ is proper, which is the case in practice.

Given $S$-morphisms $f_i:X_i\to Y_i$ $(i=1,2)$ and a commutative square (III 3.7.1) with $c$ proper, one defines, as in loc. cit., direct-image arrows analogous to (III 3.7.5), (III 3.7.6), for $L_i\in\Ob D_{\mathrm{ctf}}({}_AX_i)$.

\subsection*{6.5}
Let $A$ be a flat $\Lambda$-algebra, $X$ an $S$-scheme, $X^2=X\times_SX$, $p_1,p_2:X^2\to X$ the canonical projections, $\delta:X\to X^2$ the diagonal, and
\[
c:C\longrightarrow X^2
\]
an $S$-morphism. Put $c_i=p_ic$, and let $X^c=C\times_{X^2}X$ be the scheme of fixed points of $c$, giving a cartesian square
\begin{equation}
\begin{tikzcd}[column sep=large,row sep=large]
X^c \arrow[r,"i^c"] \arrow[d,"\delta^c"'] & X \arrow[d,"\delta"]\\
C \arrow[r,"c"'] & X^2 .
\end{tikzcd}
\tag{6.5.1}
\end{equation}
Put
\begin{equation}
KA_{X,L_{\natural}}=K_X\lten_{\Lambda}(A\lten_{A^e}A),
\qquad
KA_{X,\natural}=K_X\lten_{\Lambda}A_{\natural}.
\tag{6.5.2}
\end{equation}
Let $L\in\Ob D_{\mathrm{ctf}}({}_AX)$. We define a “local trace” homomorphism
\begin{equation}
\Tr_A:\delta^*\RHom_A(p_1^*L,p_2^!L)
\longrightarrow
KA_{X,L_{\natural}},
\tag{6.5.3}
\end{equation}
by composing the isomorphism
\begin{equation}
\delta^*\RHom_A(p_1^*L,p_2^!L)
\xrightarrow{\sim}
\delta^*(D_AL\lten_{S,A}L)
\tag{6.5.4}
\end{equation}
deduced from the inverse of 6.4.1 by applying $\delta^*$, the trivial isomorphism
\begin{equation}
\delta^*(D_AL\lten_{S,A}L)=D_AL\lten_AL,
\tag{6.5.5}
\end{equation}
and the evaluation arrow 5.3.5
\begin{equation}
D_AL\lten_AL\longrightarrow KA_{X,L_{\natural}}.
\tag{6.5.6}
\end{equation}
(Here $KA_X\lten_{A^e}A$ is identified with $KA_{X,L_{\natural}}$ by means of 6.3.2, $K_X\otimes A=KA_X$.)
It follows at once from the definition that, for $X=S$, 6.5.3 coincides with the arrow
\[
\Tr_A:\RHom_A(L,L)\longrightarrow A\lten_{A^e}A
\]
of 5.6.1. Composing the arrow deduced from 6.5.3 by applying $i^cc^!$ with the arrow
\begin{equation}
\delta^*c_*^c\RHom_A(c_1^*L,c_2^!L)
\longrightarrow
i_*^ci^{c!}\delta^*\RHom_A(p_1^*L,p_2^!L)
\tag{6.5.7}
\end{equation}
deduced from 6.4.2 and the base-change arrow
\[
\delta^*c_*c^!\longrightarrow i_*^ci^{c!}\delta^*,
\]
one obtains an arrow
\begin{equation}
\Tr_A:\delta^*c_*\RHom_A(c_1^*L,c_2^!L)
\longrightarrow
i_*^ci^{c!}KA_{X,L_{\natural}}
=i_*^cKA_{X^c,L_{\natural}},
\tag{6.5.8}
\end{equation}
and hence the arrow
\begin{equation}
\Tr_A:\Hom_A(c_1^*L,c_2^!L)
\longrightarrow
H^0(X^c,KA_{X^c,L_{\natural}}).
\tag{6.5.9}
\end{equation}
We again write $\Tr_A$ for the composites of 6.5.3, 6.5.8, 6.5.9 with the arrows deduced from the canonical arrow
\[
KA_{X,L_{\natural}}\longrightarrow KA_{X,\natural}.
\]
One checks easily that, for $A=\Lambda$, one has, for $u\in\Hom(c_1^*L,c_2^!L)$,
\begin{equation}
\Tr(u)=\langle u,1\rangle,
\tag{6.5.10}
\end{equation}
with the notation of (III 4.2.5), where $1$ denotes the diagonal correspondence. More precisely, the arrow
\[
c_*\RHom(c_1^*L,c_2^!L)\longrightarrow \delta_*i_*^cK_{X^c}
\]
adjoint to 6.5.8 is the one defined by (III 4.2.4') with $X_1=X_2$, $L_1=L_2$, $d=\delta$, and the diagonal correspondence $1\in\Hom(L,L)$.

\subsection*{6.6}
Let $A$ be a flat $\Lambda$-algebra, and let, for $i=1,2$, $X_i$ be an $S$-scheme, $X=X_1\times_SX_2$, and $L_i\in\Ob D_{\mathrm{ctf}}({}_AX_i)$. We define a pairing
\begin{equation}
\RHom_A(p_1^*L_1,p_2^!L_2)
\lten_{\Lambda}
\RHom_A(p_2^*L_2,p_1^!L_1)
\longrightarrow
KA_{X,L_{\natural}}
\tag{6.6.1}
\end{equation}
as follows. One identifies the left hand side with
\[
\RHom_A(p_1^*L_1,p_1^*KA_{X_1})
\lten_A p_2^*L_2
\lten_{\Lambda}
\RHom_A(p_2^*L_2,p_2^*KA_{X_2})
\lten_A p_1^*L_1
\]
by the isomorphisms 6.4.1, then with
\[
A\lten_{A^e}\RHom_A(\ )\lten_{\Lambda}p_2^*L_2
 \lten_{\Lambda}\RHom_A(\ )\lten_{\Lambda}p_1^*L_1\lten_{A^e}A
\]
by 5.3.1. This object is sent to
\[
A\lten_{A^e}(KA_{X_1}\lten_{\Lambda}KA_{X_2})\lten_{A^e}A=R
\]
by the tensor product of the evaluation arrows 5.3.6. Finally $R$ is identified with
\[
A\lten_{A^e}(KA_{X_1}\lten_AKA_{X_2})
\]
by 5.8.2, and then with $KA_{X,L_{\natural}}$ by 6.5.2 and (III 1.7.6).

If $c:C\to X$, $d:D\to X$ are correspondences, and
\[
e=c\times_Xd:E\to X,
\]
one deduces from 6.6.1, as in (III 4.2), a pairing
\begin{equation}
c_*\RHom_A(c_1^*L_1,c_2^!L_2)
\lten_{\Lambda}
d_*\RHom_A(d_2^*L_2,d_1^!L_1)
\longrightarrow
e_*KA_{E,L_{\natural}},
\tag{6.6.2}
\end{equation}
and hence a cup-product
\begin{equation}
\langle\ ,\ \rangle_A:
\Hom_A(c_1^*L_1,c_2^!L_2)\otimes
\Hom_A(d_2^*L_2,d_1^!L_1)
\longrightarrow
H^0(E,KA_{E,L_{\natural}}).
\tag{6.6.3}
\end{equation}
It is not difficult to extend to the non-commutative case the notion of composition of correspondences studied in (III 5.2), and one checks that, for
\[
u\in\Hom_A(c_1^*L_1,c_2^!L_2),\qquad
v\in\Hom_A(d_2^*L_2,d_1^!L_1),
\]
one has the formula analogous to (III 5.2.9)
\begin{equation}
\langle u,v\rangle_A=\Tr_A(\bar vu)=\Tr_A(u\bar v).
\tag{6.6.4}
\end{equation}
The Lefschetz formula (III 4.4) generalizes without difficulty in this context. Given proper $S$-morphisms $f_i:X_i\to Y_i$ $(i=1,2)$ and a commutative cube (III 4.4.0), for $L_i\in\Ob D_{\mathrm{ctf}}({}_AX_i)$ one has a commutative square analogous to (III 4.4.1), with $K_C$ (resp. $K_D$) replaced by $KA_{C,L_{\natural}}$ (resp. $KA_{D,L_{\natural}}$), and a formula analogous to (III 4.5.1), in $H^0(D,KA_{D,L_{\natural}})$:
\begin{equation}
\langle f_*u,f_*v\rangle_A=g_*\langle u,v\rangle_A.
\tag{6.6.5}
\end{equation}
The arrow
\[
f_*c_*KA_{C,L_{\natural}}\longrightarrow d_*KA_{D,L_{\natural}}
\]
is deduced from arrow (4) of (III 4.4.1) by tensoring with
\[
A\lten_{A^e}A .
\]

\subsection*{6.7}
Let $X$ and $c$ be as in 6.5, and let $A\to B$ be a morphism of flat $\Lambda$-algebras. With respect to extension and restriction of scalars, the arrows $\Tr_A$ (6.5.3, 6.5.8, 6.5.9) behave like the arrows $\Tr_A$ of 5.6.1, 5.6.2. The point is that, if $f$ is an $S$-morphism, the functors $(f^*,f_*,f_!,f^!)$ commute with extension and restriction of scalars; consequently, the study of the behavior of $\Tr_A$ (6.5.3) with respect to extension and restriction of scalars reduces to the behavior of the “tensor-product” arrow
\[
\RHom_A(p_1^*L,p_1^*KA_X)\lten_Ap_2^*L
\longrightarrow
\RHom_A(p_1^*L,p_1^*KA_X\lten_Ap_2^*L)
\]
and of the “evaluation” arrow
\[
\RHom_A(L,KA_X)\lten_A L\longrightarrow KA_{X,L_{\natural}}.
\]
Thus the commutativity of the squares 5.9.1, 5.9.2 implies the commutativity of the square
\begin{equation}
\begin{tikzcd}[column sep=huge,row sep=large]
\delta^*c_*\RHom_A(c_1^*L,c_2^!L) \arrow[r,"\Tr_A"] \arrow[d]
& i_*^cKA_{X^c,L_{\natural}} \arrow[d]\\
\delta^*c_*\RHom_B(c_1^*(B\lten_A L),c_2^!(B\lten_A L)) \arrow[r,"\Tr_B"']
& i_*^cKB_{X^c,L_{\natural}},
\end{tikzcd}
\tag{6.7.1}
\end{equation}
where the left vertical arrow is extension of scalars and the right vertical arrow is defined by the canonical arrow
\[
A\lten_{A^e}A\longrightarrow B\lten_{B^e}B.
\]
It follows that, for $u\in\Hom_A(c_1^*L,c_2^!L)$,
\begin{equation}
\Tr_B(B\lten_Au)=\varphi\Tr_A(u),
\tag{6.7.2}
\end{equation}
where
\[
\varphi:H^0(X^c,KA_{X^c,L_{\natural}})\longrightarrow H^0(X^c,KB_{X^c,L_{\natural}})
\]
is the canonical map.

Assume $A$ is finite type and flat as a $\Lambda$-module, $B$ finite type and flat as a left $A$-module, and $A_{\natural}$, $B_{\natural}$ flat over $\Lambda$. There is a map $\Tr_{B/A}:B_{\natural}\to A_{\natural}$; we shall again denote by
\begin{equation}
\Tr_{B/A}:KB_{X,\natural}\longrightarrow KA_{X,\natural}
\tag{6.7.3}
\end{equation}
the arrow deduced from it by applying $K_X\lten -$ (6.5.2). It follows from 5.10.6 and 5.10.9 that, for $L\in\Ob D_{\mathrm{ctf}}({}_BX)$, there is a commutative square
\begin{equation}
\begin{tikzcd}[column sep=huge,row sep=large]
\delta^*c_*\RHom_B(c_1^*L,c_2^!L) \arrow[r,"\Tr_B"] \arrow[d]
& i_*^cKB_{X^c,\natural} \arrow[d,"\Tr_{B/A}"]\\
\delta^*c_*\RHom_A(c_1^*L,c_2^!L) \arrow[r,"\Tr_A"']
& i_*^cKA_{X^c,\natural},
\end{tikzcd}
\tag{6.7.4}
\end{equation}
where the left vertical arrow is restriction of scalars. Consequently, for $u\in\Hom_B(c_1^*L,c_2^!L)$, one has, in $H^0(X^c,KA_{X^c,\natural})$,
\begin{equation}
\Tr_A(u)=\Tr_{B/A}(\Tr_B(u)).
\tag{6.7.5}
\end{equation}

\subsection*{6.8}
Let $A$, $B$ be flat $\Lambda$-algebras, $X$, $Y$ $S$-schemes, $Z=X\times_SY$, and let $c:C\to X^2$, $d:D\to Y^2$ be correspondences. Let
\[
e=c\times_Sd:E=C\times_SD\longrightarrow Z^2
\]
be their product. Then
\[
Z^e=X^c\times_SY^d.
\]
Put $R=A\otimes_{\Lambda}B$. As in (III 5.3), one defines the external tensor product
\begin{equation}
\Hom_A(c_1^*L,c_2^!L)\otimes \Hom_B(d_1^*M,d_2^!M)
\xrightarrow{\ \lten_S\ }
\Hom_R(e_1^*(L\lten_SM),e_2^!(L\lten_SM)),
\tag{6.8.1}
\end{equation}
and one deduces easily from the commutativity of the squares 5.12.1, 5.12.3 that, for
\[
u\in\Hom_A(c_1^*L,c_2^!L),\qquad v\in\Hom_B(d_1^*M,d_2^!M),
\]
one has, in $H^0(Z^e,KR_{Z^e,L_{\natural}})$,
\begin{equation}
\Tr_R(u\lten_Sv)=\Tr_A(u)\Tr_B(v),
\tag{6.8.2}
\end{equation}
the product on the right being defined by the canonical arrow
\begin{equation}
H^0(X^c,KA_{X^c,L_{\natural}})\otimes
H^0(Y^d,KB_{Y^d,L_{\natural}})
\longrightarrow
H^0(Z^e,KR_{Z^e,L_{\natural}}),
\tag{6.8.3}
\end{equation}
coming from the isomorphism
\[
KA_{X^c,L_{\natural}}\lten_S KB_{Y^d,L_{\natural}}\longrightarrow K(A\otimes B)_{Z^e,L_{\natural}}
\]
given by (III 1.7.6) and arrow (4) of 5.9.1.

\subsection*{6.9}
Let $G$ be a finite group. If $X$ is a $G$-$S$-scheme, that is, an $S$-scheme endowed with an action of $G$ by $S$-automorphisms,\footnote{We agree that $G$ acts on spaces on the right.} and $A$ is a $\Lambda$-algebra,\footnote{We are chiefly interested in the case of commutative $\Lambda$-algebras.} we denote by $D(G,{}_AX)$ the derived category of the category of sheaves of left $G$-$A$-modules on $X$. For the notions of $G$-sheaf of sets, rings, modules, and so on, on $X$ for the étale topology, we refer the reader to (X 1). We denote by the subscript $c$ (resp. ctf) the full subcategory formed by complexes of $G$-$A$-modules which, as complexes of $A$-modules, have constructible cohomology (resp. finite tor-dimension and constructible cohomology).

Similarly, if $B$ is a $\Lambda$-algebra, we denote by $D(G,{}_AX_B)$ the derived category of sheaves of $G$-$(A,B)$-bimodules on $X$. If $G$ acts trivially on $X$, a left $G$-$A$-module on $X$ is just a left $A[G]$-module on $X$ without operators, hence
\[
D(G,{}_AX)=D({}_{A[G]}X).
\]
Since $G$ is finite, one also has
\[
D_c({}_AX)=D_c({}_{A[G]}X),
\]
but only an inclusion
\[
D_{\mathrm{ctf}}({}_{A[G]}X)\subset D_{\mathrm{ctf}}(G,{}_AX).
\]

The duality formalism of (SGA 4 XVII, XVIII) extends without difficulty to $G$-$S$-schemes. Since the functor forgetting the action of $G$ is faithful, it is enough to extend naturally the definition of the “six operations” $(\lten,\RHom,f^*,f_*,f_!,f^!)$ and the canonical arrows considered in loc. cit. (base change, Künneth, duality, etc.). This extension presents no difficulty: for the definition of the functors $f_!$, $f^!$, it is enough to use equivariant compactifications. We leave to the reader the task of paraphrasing the definitions and results 6.2 to 6.8 in the setting of $G$-$S$-schemes.

If, for $i=1,2$, $X_i$ is a $G$-$S$-scheme, $c:C\to X=X_1\times_SX_2$ is a $G$-$S$-morphism, and $L_i\in\Ob D_{\mathrm{ctf}}(G,{}_AX_i)$, the elements of
\[
\Hom_A(c_1^*L_1,c_2^!L_2),
\]
where $\Hom_A$ is taken in the category $D(G,{}_A-)$, will be called equivariant cohomological correspondences from $L_1$ to $L_2$ with support in $c$. Note that, when $G$ acts trivially on $X_i$ $(i=1,2)$ and $L_i\in\Ob D_{\mathrm{ctf}}({}_{A[G]}X_i)$, the isomorphism 6.4.1 in $D(G,X)$ refines to an isomorphism
\[
D_{A[G]}L_1\lten_{S,A[G]}L_2
\xrightarrow{\sim}
\RHom_{A[G]}(p_1^*L_1,p_2^!L_2).
\]
It follows that, in a $G$-equivariant situation of type 6.5, with trivial action of $G$ on the spaces, and with $L\in\Ob D_{\mathrm{ctf}}({}_{A[G]}X)\subset D_{\mathrm{ctf}}(G,{}_AX)$, one may, for
\[
u\in\Hom_{A[G]}(c_1^*L,c_2^!L),
\]
define a local term
\[
\Tr_{A[G]}(u)\in H^0(X^c,KA[G]_{X^c,L_{\natural}}),
\]
finer, in the evident sense, than the local term
\[
\Tr_A(u)\in H^0(G,X^c,KA_{X^c,L_{\natural}})
\]
obtained by using only finite tor-dimension of $L$ with respect to $A$.

\subsection*{6.10}
Let $X$, $c$ be as in 6.5, and let $L\in\Ob D_{\mathrm{ctf}}({}_{\Lambda[G]}X)$. For
\[
u\in\Hom_{\Lambda[G]}(c_1^*L,c_2^!L),
\]
6.5.9 gives a local trace
\begin{equation}
\Tr_{\Lambda[G]}(u)\in
H^0(X^c,K\Lambda[G]_{X^c,\natural})
=H^0(X^c,K_{X^c}\otimes\Lambda[G_{\natural}]),
\tag{6.10.1}
\end{equation}
where $G_{\natural}$ is the set of conjugacy classes of $G$. For every $x\in G_{\natural}$, one may consider the “value of $\Tr_{\Lambda[G]}(u)$ at $x$,”
\[
\Tr_{\Lambda[G]}(u)(x)\in H^0(X^c,K_{X^c}),
\]
the image of $\Tr_{\Lambda[G]}(u)$ under the projection
\[
H^0(X^c,K_{X^c}\otimes\Lambda[G_{\natural}])
\longrightarrow H^0(X^c,K_{X^c})
\]
defined by $x$. It follows from 6.7.5 and 5.11.2 that
\begin{equation}
\Tr_{\Lambda}(u)=|G|\,\Tr_{\Lambda[G]}(u)(e),
\tag{6.10.2}
\end{equation}
where $e$ denotes the identity element of $G$.

More generally, let $g\in G$, let $Z_g$ be the centralizer of $g$, and let
\[
gu\in\Hom_{\Lambda[Z_g]}(c_1^*L,c_2^!L)
\]
be the correspondence obtained by composing $u$ with left multiplication by $g:c_1^*L\to c_1^*L$. We have $L\in\Ob D_{\mathrm{ctf}}({}_{\Lambda[Z_g]}X)$, and hence 6.10.2 implies
\begin{equation}
\Tr_{\Lambda}(gu)=|Z_g|\,\Tr_{\Lambda[Z_g]}(gu)(e).
\tag{6.10.3}
\end{equation}
By the same argument as in 5.11.3, one deduces from 6.7.4 the formula
\begin{equation}
\Tr_{\Lambda[Z_g]}(gu)(e)=\Tr_{\Lambda[G]}(u)(g^{-1}),
\tag{6.10.4}
\end{equation}
which allows 6.10.3 to be rewritten, if desired, in the form
\begin{equation}
\Tr_{\Lambda}(gu)=|Z_g|\,\Tr_{\Lambda[G]}(u)(g^{-1}).
\tag{6.10.5}
\end{equation}

\subsection*{6.11}
Let $f:X\to Y$ be a proper $G$-$S$-morphism, and let
\begin{equation}
\begin{tikzcd}[column sep=large,row sep=large]
C \arrow[r,"c"] \arrow[d] & X\times_SX \arrow[d,"f\times_Sf"]\\
D \arrow[r,"d"'] & Y\times_SY
\end{tikzcd}
\tag{6.11.1}
\end{equation}
be a commutative $G$-equivariant square of $G$-$S$-schemes, with $c$ and $d$ proper. Assume that $G$ acts trivially on $Y$ and $D$. We denote by
\begin{equation}
f^{c/d}\quad(\text{or } f^c):X^c\longrightarrow Y^d
\tag{6.11.2}
\end{equation}
the morphism induced by 6.11.1 on the schemes of fixed points of $c$ and $d$, with the notation of 6.5.1. Let $L\in\Ob D_{\mathrm{ctf}}(G,{}_AX)$, and let
\[
u\in\Hom(c_1^*L,c_2^!L)
\]
be a $G$-equivariant correspondence with support in $c$. By direct image, it gives a $G$-equivariant correspondence
\[
f_*u\in\Hom_{\Lambda[G]}(d_1^*f_*L,d_2^!f_*L).
\]
For $g\in G$, let $Z_g$ denote the centralizer of $g$, and let
\begin{equation}
gu\in\Hom((gc)_1^*L,c_2^!L)
\tag{6.11.3}
\end{equation}
be the $Z_g$-equivariant correspondence with support in
\begin{equation}
gc\stackrel{\mathrm{dfn}}{=}(gc_1=c_1g,c_2):C\to X\times_SX,
\tag{6.11.4}
\end{equation}
defined as the composite of the arrows
\[
g^*(c_1^*L)\xrightarrow{\alpha_g}c_1^*L\xrightarrow{u}c_2^!L,
\]
where $\alpha_g$ is the arrow giving the action of $G$ on $c_1^*L$. It is clear that
\begin{equation}
f_*(gu)=g(f_*u).
\tag{6.11.5}
\end{equation}

\begin{statement}{Proposition 6.12}
Under the hypotheses of 6.11, assume that
\[
f_*L\in\Ob D_{\mathrm{ctf}}({}_{\Lambda[G]}Y).
\]
Then, for every $g\in G$,
\begin{equation}
f_*^{gc}\Tr_{\Lambda}(gu)
=
|Z_g|\,\Tr_{\Lambda[G]}(f_*u)(g^{-1}),
\tag{6.12.1}
\end{equation}
where
\[
f_*^{gc}:H^0(X^{gc},K_{X^{gc}})\longrightarrow H^0(Y^d,K_{Y^d})
\]
is the ``Gysin'' morphism defined by the adjunction arrow
\[
f_*^{gc}K_{X^{gc}}=f_*^{gc}f^{gc!}K_{Y^d}\longrightarrow K_{Y^d}.
\]
\end{statement}

Indeed, by 6.10.5 applied to $f_*u$, and by 6.11.5, one has
\[
\Tr_{\Lambda}(f_*(gu))=|Z_g|\,\Tr_{\Lambda[G]}(f_*u)(g^{-1}),
\]
and 6.12.1 therefore follows from the Lefschetz formula (III 4.5.1).

\subsection*{6.13}
Let $\ell$ be a prime number distinct from $p$. For every integer $n\ge1$, put
\[
\Lambda_n=\mathbb Z/\ell^n.
\]
Under the hypotheses of 6.11, let $L=(L_n)$ be a projective system of objects of
\[
D_{\mathrm{ctf}}(G,X,\Lambda_n)
\]
such that
\[
L_m\lten_{\Lambda_m}\Lambda_n\xrightarrow{\sim}L_n\quad (m\ge n),
\]
and let
\[
u\in\Hom(c_1^*L,c_2^!L)
\]
be a projective system of equivariant correspondences $u_n\in\Hom(c_1^*L_n,c_2^!L_n)$ such that, for $m\ge n$,
\[
u_m\lten_{\Lambda_m}\Lambda_n=u_n
\]
modulo the isomorphism $L_m\lten_{\Lambda_m}\Lambda_n=L_n$. By compatibility of local traces with extension of scalars, the local terms, for $g\in G$,
\[
\Tr_{\Lambda_n}(gu_n)\in H^0(X^{gc},K\Lambda_{n,X^{gc}})
\]
define an element
\[
\Tr_{\mathbb Z_{\ell}}(gu)\in
H^0(X^{gc},K_{\ell,X^{gc}})
\stackrel{\mathrm{dfn}}{=}
\varprojlim_n H^0(X^{gc},K\Lambda_{n,X^{gc}}),
\]
and the terms
\[
f_*^{gc}\Tr_{\Lambda_n}(gu_n)=\Tr_{\Lambda_n}(gf_*u_n)
\]
an element
\[
f_*\Tr_{\mathbb Z_{\ell}}(gu)\in
H^0(Y^d,K_{\ell,Y^d})
\stackrel{\mathrm{dfn}}{=}
\varprojlim_n H^0(Y^d,K\Lambda_{n,Y^d}).
\]
Assume moreover that
\[
f_*M_n\in\Ob D_{\mathrm{ctf}}({}_{\Lambda_n[G]}Y)
\]
for all $n$. Then the local traces
\[
\Tr_{\Lambda_n[G]}(f_*u_n)
\]
similarly define an element
\[
\Tr_{\mathbb Z_{\ell}[G]}(f_*u)\in
H^0(Y^d,K\mathbb Z_{\ell}[G]_{Y^d,\natural})
\stackrel{\mathrm{dfn}}{=}
\varprojlim_n H^0(Y^d,K\Lambda_n[G]_{Y^d,\natural}).
\]
It follows from 6.12 that
\begin{equation}
f_*^{gc}\Tr_{\mathbb Z_{\ell}}(gu)
=
|Z_g|\,\Tr_{\mathbb Z_{\ell}[G]}(f_*u)(g^{-1}).
\tag{6.13.1}
\end{equation}

\subsection*{6.14}
Let $X$ be a $G$-$S$-scheme, let $j:U\to X$ be an equivariant open immersion, and let
\[
c:C\to X\times_SX
\]
be a morphism of $G$-$S$-schemes such that $c_2(=p_2c):C\to X$ is quasi-finite and of finite tor-dimension. Suppose that
\[
c_1^{-1}(U)\subset c_2^{-1}(U).
\]
Then there are commutative squares
\begin{equation}
\begin{tikzpicture}[baseline=(current bounding box.center),>=stealth]
\node (U1) at (0,3.0) {$U$};
\node (C1U) at (2.8,3.0) {$c_1^{-1}(U)$};
\node (C2U) at (2.8,2.15) {$c_2^{-1}(U)$};
\node (U2) at (5.2,2.15) {$U$};
\node (X1) at (0,1.15) {$X$};
\node (C) at (2.8,1.15) {$C$};
\node (X2) at (5.2,1.15) {$X$};
\node (XU1) at (0,0) {$X-U$};
\node (C2XU) at (2.8,0) {$c_2^{-1}(X-U)$};
\node (XU2) at (5.2,0) {$X-U$};
\draw[->] (C1U) -- (U1);
\draw[->] (C1U) -- (C2U);
\draw[->] (C2U) -- (U2);
\draw[->] (U1) -- node[left] {$i$} (X1);
\draw[->] (U2) -- (X2);
\draw[->] (C2U) -- (C);
\draw[->] (C2XU) -- (C);
\draw[->] (C) -- node[above] {$c_1$} (X1);
\draw[->] (C) -- node[above] {$c_2$} (X2);
\draw[->] (XU1) -- node[left] {$j$} (X1);
\draw[->] (XU2) -- (X2);
\draw[->] (C2XU) -- (XU1);
\draw[->] (C2XU) -- (XU2);
\end{tikzpicture}
\tag{6.14.1}
\end{equation}
The trace morphism
\[
\Tr_{c_2}:c_{2!}\Lambda_C\longrightarrow \Lambda_X
\]
then defines a filtered $G$-equivariant correspondence
\begin{equation}
\underline c\in\Hom(c_1^*\Lambda_X,c_2^!\Lambda_X),
\tag{6.14.2}
\end{equation}
where $\Lambda_X$ is endowed with the filtration defined by the submodule $i_!\Lambda_U$, that is,
\[
F^n\Lambda_X=\Lambda_X\quad(n\le0),\qquad
F^1\Lambda_X=i_!\Lambda_U,\qquad
F^n\Lambda_X=0\quad(n\ge1).
\]
The associated graded consists of two $G$-equivariant correspondences
\[
\operatorname{gr}^0\underline c=\underline c_{X-U,X}
\in\Hom(c_1^*\Lambda_{X-U,X},c_2^!\Lambda_{X-U,X}),
\]
\begin{equation}
\operatorname{gr}^1\underline c=\underline c_{U,X}
\in\Hom(c_1^*\Lambda_{U,X},c_2^!\Lambda_{X-U,X}),
\tag{6.14.3}
\end{equation}
where we have put
\[
\Lambda_{X-U,X}=j_*\Lambda_{X-U},\qquad
\Lambda_{U,X}=i_!\Lambda_U.
\]
We have
\begin{equation}
\underline c_{X-U,U}=j_*\underline c_{X-U},
\tag{6.14.4}
\end{equation}
where
\[
\underline c_{X-U}\in\Hom(c_1^*\Lambda_{X-U},c_2^!\Lambda_{X-U})
\]
is the correspondence defined by $\Tr_{c_2,X-U}$. We write $\underline c$, $\underline c_{U,X}$, and so on, when we want to specify the base ring.

Forgetting the action of $G$, there are local terms
\begin{equation}
\Tr_{\Lambda}(\underline c),\quad
\Tr_{\Lambda}(\underline c_{U,X}),\quad
\Tr_{\Lambda}(\underline c_{X-U,X})
\in H^0(X^c,K_{X^c}),
\tag{6.14.5}
\end{equation}
which, by additivity of traces (III 4.13.1), satisfy the relation
\begin{equation}
\Tr_{\Lambda}(\underline c)
=
\Tr_{\Lambda}(\underline c_{U,X})
+
\Tr_{\Lambda}(\underline c_{X-U,X}).
\tag{6.14.6}
\end{equation}
Moreover, by 6.14.4,
\begin{equation}
\Tr_{\Lambda}(\underline c_{X-U,X})
=
j_*^c\Tr_{\Lambda}(\underline c_{X-U}),
\tag{6.14.7}
\end{equation}
where
\[
j_*^c:H^0((X-U)^c,K_{(X-U)^c})\longrightarrow H^0(X^c,K_{X^c})
\]
is the direct-image morphism defined by $j^c:(X-U)^c\to X^c$. Recall (III 4.3) that, if $X$ is smooth, $c$ a closed immersion, and $X^c$ finite, $\Tr_{\Lambda}(\underline c)$ is the image in $\Lambda$ of the intersection multiplicity $(C\cdot X)$ of $C$ with the diagonal. If, moreover, $X-U$ is smooth, $\Tr_{\Lambda}(\underline c_{X-U})$ is the image in $\Lambda$ of $(c_2^{-1}(X-U)\cdot(X-U))$. Note also that, if $X-U$ is proper over $S$ and $c|X-U$ is the diagonal, then the Lefschetz formula (III 4.7) implies that
\[
\int_{X-U}\Tr(\underline c_{X-U})\in\Lambda
\]
is the Euler--Poincaré characteristic of $X-U$ with values in $\Lambda$.

\subsection*{6.15}
Let there be a $G$-equivariant diagram 6.11.1, with $f$, $c$, $d$ proper, and an equivariant open immersion $U\to X$. Suppose that $c_2$ is quasi-finite and of finite tor-dimension, that
\[
c_1^{-1}(U)\subset c_2^{-1}(U),
\]
that $G$ acts trivially on $C$ and $D$, and that
\[
f|U:U\longrightarrow f(U)=V
\]
is an étale Galois covering with group $G$. This last hypothesis implies that
\[
Rf_*(\Lambda_{U,X})\in\Ob D_{\mathrm{ctf}}({}_{\Lambda[G]}Y):
\]
indeed $Rf_*(\Lambda_{U,X})$ is concentrated in degree $0$ and is extension by zero of a sheaf on $V$ locally isomorphic to $\Lambda[G]$. Hence one may apply 6.12 to $L=\Lambda_{U,X}$, $u=\underline c_{U,X}$, and, in view of 6.14.6 and 6.14.7, one obtains:

\begin{statement}{Proposition 6.16}
Under the hypotheses of 6.15, for every $g\in G$ one has
\begin{equation}
f_*^{gc}\Tr_{\Lambda}(g\underline c)
-(fj)_*^{gc}\Tr_{\Lambda}(g\underline c_{X-U})
=
|Z_g|\,\Tr_{\Lambda[G]}(f_*\underline c_{U,X})(g^{-1}),
\tag{6.16.1}
\end{equation}
where $j:X-U\to X$ is the inclusion, and $Z_g$ denotes the centralizer of $g$.
\end{statement}

Applying 6.16 to the projective system of correspondences $\underline c_{\ell}=(\underline c_{\Lambda_n})$, where $\Lambda_n=\mathbb Z/\ell^n$ with $\ell$ prime distinct from $p$, and taking into account that 6.14.6 and 6.14.7 give, after passage to the limit and application of $f_*^{gc}$,
\[
f_*^{gc}\Tr_{\mathbb Z_{\ell}}(g\underline c_{\ell})
=
f_*^{gc}\Tr_{\mathbb Z_{\ell}}(g\underline c_{\ell,U,X})
+(fj)_*^{gc}\Tr_{\mathbb Z_{\ell}}(g\underline c_{\ell,X-U}),
\]
one deduces, by 6.13:

\begin{statement}{Corollary 6.17}
Under the hypotheses of 6.15, for every $g\in G$ one has
\begin{equation}
f_*^{gc}\Tr_{\mathbb Z_{\ell}}(g\underline c_{\ell})
-(fj)_*^{gc}\Tr_{\mathbb Z_{\ell}}(g\underline c_{\ell,X-U})
=
|Z_g|\,\Tr_{\mathbb Z_{\ell}[G]}(f_*\underline c_{\ell,U,X})(g^{-1}),
\tag{6.17.1}
\end{equation}
with the notation of 6.16.
\end{statement}

In view of the compatibility of local traces with localization and of the remarks made in 6.14, one sees that 6.17 gives an affirmative solution, in a more general formulation, to Grothendieck's divisibility conjecture (XII 4.5), in the case where the endomorphism $\varphi$ of $G$ considered in loc. cit. is the identity.

In fact, by following 5.13, it is not difficult to give divisibility statements analogous to 6.12, 6.13, 6.16, 6.17, in the case where $c$ is not $G$-equivariant, but where $c_2$ is and $c_1$ satisfies
\[
c_1(xg)=c_1(x)\varphi(g)
\]
for every point $x$ of $G$, $\varphi$ being a given endomorphism of $G$. The divisibility formulas are then written as above, with $Z_g$ replaced by the $\varphi$-centralizer of $g^{-1}$ (5.13.5).

\subsection*{6.18}
Put
\begin{equation}
N_G=\sum_{g\in G} g\in\Lambda[G].
\tag{6.18.1}
\end{equation}
Let $X$ be an $S$-scheme. For $L\in\Ob D({}_{\Lambda[G]}X)$, left multiplication by $N_G$ induces an arrow in $D(X)$
\begin{equation}
N_G:\Lambda\lten_{\Lambda[G]}L
\longrightarrow
\RHom_{\Lambda[G]}(\Lambda,L),
\tag{6.18.2}
\end{equation}
where $\Lambda$ is regarded as a $\Lambda[G]$-algebra by the natural augmentation sending every $g\in G$ to $1$.

\begin{statement}{Lemma 6.18.3}
If $L\in\Ob D_{\mathrm{ctf}}({}_{\Lambda[G]}X)$, the arrow 6.18.2 is an isomorphism.
\end{statement}

By dévissage, one reduces to the case where $L$ is perfect, then to $L=\Lambda[G]$, and the assertion follows from the fact that $N_G$ (6.18.1) is a basis of the module of $G$-invariants of the left $\Lambda[G]$-module $\Lambda[G]$.

\begin{statement}{Proposition 6.19}
In the situation of 6.11, let $M\in\Ob D_{\mathrm{ctf}}(Y)$,
\[
v\in\Hom(d_1^*M,d_2^!M),
\]
and let $a:M\to f_*L$ be an arrow in $D({}_{\Lambda[G]}Y)$, where $M$ is regarded as an object of $D({}_{\Lambda[G]}Y)$ by the augmentation $\Lambda[G]\to\Lambda$, such that the square
\[
\begin{tikzcd}[column sep=large,row sep=large]
d_1^*M \arrow[r,"v"] \arrow[d] & d_2^!M \arrow[d] & {}\\
d_1^*f_*L \arrow[r,"f_*u"'] & d_2^!f_*L
\end{tikzcd}
\]
is commutative. Suppose that $f_*M\in\Ob D_{\mathrm{ctf}}({}_{\Lambda[G]}Y)$ and that $a$ induces an isomorphism
\[
M\xrightarrow{\sim}\RHom_{\Lambda[G]}(\Lambda,f_*L).
\]
Then
\begin{equation}
\Tr_{\Lambda}(v)=
\sum_{g\in G_{\natural}}\Tr_{\Lambda[G]}(f_*u)(g^{-1}).
\tag{6.19.1}
\end{equation}
\end{statement}

By 6.18.3, $a$, composed with $N_G^{-1}$, gives an isomorphism
\[
M\xrightarrow{\sim}\Lambda\lten_{\Lambda[G]}f_*L,
\]
under which $v$ is identified with
\[
\Lambda\lten_{\Lambda[G]}f_*u.
\]
By 6.7.2, $\Tr_{\Lambda}(v)$ is therefore the image of $\Tr_{\Lambda[G]}(f_*u)$ under the arrow induced by the augmentation
\[
\Lambda[G]\longrightarrow \Lambda.
\]
This augmentation, by definition, sends every conjugacy class to $1$, whence the conclusion.

\subsection*{6.20}
Let, for $i=1,2$, $X_i$ be a $G$-$S$-scheme, and let
\[
X=X_1\times_SX_2,\qquad c:C\to X
\]
be a $G$-$S$-morphism such that $c_2$ is quasi-finite and of finite tor-dimension. Let $L_i\in\Ob D^+(G,{}_AX_i)$. The trace morphism
\[
\Tr_{c_2}:c_{2!}c_2^*L_2\longrightarrow L_2
\]
defines by adjunction an arrow in $D(G,X)$, again denoted
\begin{equation}
\Tr_{c_2}:c_2^*L_2\longrightarrow c_2^!L_2,
\tag{6.20.1}
\end{equation}
which associates to every $G$-equivariant morphism
\[
u\in\Hom(c_1^*L_1,c_2^*L_2)
\]
a $G$-equivariant correspondence
\begin{equation}
u^! = \Tr_{c_2}\circ u\in\Hom(c_1^*L_1,c_2^!L_2).
\tag{6.20.2}
\end{equation}

\begin{statement}{Proposition 6.21}
Let $f:X\to Y$ be a proper $G$-morphism, and let 6.11.1 be an equivariant commutative square, with trivial action of $G$ on $Y$ and $D$. Suppose $c$, $d$ are proper, $c_2$, $d_2$ quasi-finite and of finite tor-dimension, and the square
\[
\begin{tikzcd}[column sep=large,row sep=large]
C \arrow[r,"c_2"] \arrow[d,"h"'] & X \arrow[d,"f"]\\
D \arrow[r,"d_2"'] & Y
\end{tikzcd}
\]
is cartesian and tor-independent (SGA 6 III 1.5). Let $M\in\Ob D_{\mathrm{ctf}}(Y)$,
\[
v\in\Hom(d_1^*M,d_2^*M),
\]
hence a $G$-equivariant homomorphism
\[
h^*v\in\Hom(c_1^*f^*M,c_2^*f^*M),
\]
and $G$-equivariant correspondences
\[
(h^*v)^!\in\Hom(c_1^*f^*M,c_2^!f^*M),
\qquad
f_*((h^*v)^!)\in\Hom(d_1^*(f_*f^*M),d_2^!(f_*f^*M)).
\]
Then there is a commutative square of arrows in $D({}_{\Lambda[G]}Y)$, where $M$ is regarded as an object of $D({}_{\Lambda[G]}Y)$ by the trivial action,
\[
\begin{tikzcd}[column sep=huge,row sep=large]
d_1^*M \arrow[r,"v^!"] \arrow[d] & d_2^!M \arrow[d] & {}\\
d_1^*f_*f^*M \arrow[r,"f_*((h^*v)^!)"'] & d_2^!f_*f^*M,
\end{tikzcd}
\]
where the vertical arrows are given by the adjunction arrow $M\to f_*f^*M$.
\end{statement}

Put $h^*v=u$, $f^*M=L$, and consider the diagram
\[
\begin{tikzcd}[column sep=large,row sep=large]
d_1^*M \arrow[r,"v"] \arrow[d] &
d_2^*M \arrow[r,"\Tr_{d_2}"] \arrow[d] &
d_2^!M \arrow[d] & {}\\
d_1^*f_*L \arrow[r] &
h_*c_1^*L \arrow[r,"h_*\Tr_{c_2}"] &
h_*c_2^!L \arrow[r] &
d_2^!f_*L .
\end{tikzcd}
\]
The arrows (1) and (3) are given by the adjunction arrow $M\to f_*L$, (2) by the adjunction arrow $d_2^*M\to h_*h^*(d_2^*M)=h_*c_2^*L$, (4) by the base-change arrow, and (5) is the canonical isomorphism (SGA 4 XVIII 3.1.12.3). By definition, $v^!$ is the composite of the upper row. On the other hand, it follows from the interpretation of the direct image of a correspondence (III 3.5 b), 3.7) that $f_*u^!$ is the composite of the lower row. It remains only to prove commutativity of the squares (A) and (B). The commutativity of (A) follows immediately from the definition of the base-change arrow. As for (B), it follows from the fact that, in view of the tor-independence hypothesis, $\Tr_{d_2}$ is compatible with base change by $f$ (SGA $4^{1/2}$ Cycle 2.3.3).

Applying 6.19, one obtains:

\begin{statement}{Corollary 6.22}
Under the hypotheses of 6.21, assume that
\[
f_*f^*M\in\Ob D_{\mathrm{ctf}}({}_{\Lambda[G]}Y),
\]
and that the adjunction arrow
\[
M\to f_*f^*M
\]
induces an isomorphism
\[
M\xrightarrow{\sim}\RHom_{\Lambda[G]}(\Lambda,f_*f^*M).
\]
Then
\begin{equation}
\Tr_{\Lambda}(v^!)=
\sum_{g\in G_{\natural}}\Tr_{\Lambda[G]}(f_*((f^*v)^!))(g^{-1}).
\tag{6.22.1}
\end{equation}
\end{statement}

\begin{statement}{Proposition 6.23}
Under the hypotheses of 6.21, suppose moreover that there is an open immersion
\[
i:V\hookrightarrow Y
\]
such that $(fc_1)^{-1}(V)$ is connected, that $d_1^{-1}(V)\subset d_2^{-1}(V)$, and that $f$ induces a principal Galois covering with group $G$,
\[
f_V:U=f^{-1}(V)\to V.
\]
Suppose also that $M$ is of the form $M=i_!E$, where $E\in\Ob D(V)$, and that there is an isomorphism
\[
\alpha:f_V^*E\xrightarrow{\sim}\Lambda_X\lten_{\Lambda}E_0
\]
in $D(G,X)$, with $E_0\in\Ob D(\Lambda[G])$, perfect over $\Lambda$. Let
\[
v\in\Hom(d_1^*M,d_2^*M)
\]
be as in 6.21, and let
\[
v_0\in\Hom_{\Lambda[G]}(E_0,E_0)
\]
be the homomorphism defined, thanks to $\alpha$, by $f^*v|c_1^{-1}(U)$. Then, if $\underline c_{U,X}$ denotes the correspondence defined in 6.14.3, with the notation 6.20.2 one has
\begin{equation}
\Tr_{\Lambda}(v^!)
=
\sum_{g\in G_{\natural}}
\Tr_{\Lambda[G]}(f_*\underline c_{U,X})(g^{-1})\Tr_{\Lambda}(gv_0).
\tag{6.23.1}
\end{equation}
\end{statement}

Note that, by the central character of $\Tr_{\Lambda}$, the right hand side does not depend on the choice of $\alpha$.

By the projection formula
\[
f_*f^*\Lambda_Y\lten M\xrightarrow{\sim}f_*f^*M
\]
and 6.15, $M$ satisfies the hypotheses of 6.22. On the other hand, $\alpha$ defines an isomorphism
\[
f^*M\xrightarrow{\sim}\Lambda_{U,X}\otimes E_0
\]
under which $(f^*v)^!$ is identified with the external tensor product $\underline c_{U,X}\otimes v_0$. By the projection formula, $f_*f^*M$ is therefore identified with $f_*(\Lambda_{U,X})\otimes E_0$, and $f_*((f^*v)^!)$ with the external tensor product $f_*\underline c_{U,X}\otimes v_0$. For $g\in G$, the $Z_g$-equivariant correspondence $gf_*((f^*v)^!)$ is therefore identified with
\[
gf_*\underline c_{U,X}\otimes gv_0.
\]
Since, by 6.10.4,
\[
\Tr_{\Lambda[G]}(f_*((f^*v)^!))(g^{-1})
=
\Tr_{\Lambda[Z_g]}(gf_*((f^*v)^!))(e),
\]
it is enough, in view of 6.22.1, to prove the formula
\[
\Tr_{\Lambda[Z_g]}(gf_*\underline c_{U,X}\otimes gv_0)(e)
=
\Tr_{\Lambda[G]}(f_*\underline c_{U,X})(g^{-1})\Tr_{\Lambda}(gv_0),
\]
or, equivalently by 6.10.4,
\begin{equation}
\Tr_{\Lambda[Z_g]}(gf_*\underline c_{U,X}\otimes gv_0)(e)
=
\Tr_{\Lambda[Z_g]}(gf_*\underline c_{U,X})(e)\Tr_{\Lambda}(gv_0).
\tag{6.23.2}
\end{equation}
Replacing $G$ by $Z_g$ if necessary, we may suppose that $g=e$. Returning to the definition of the traces, one sees that it is enough to establish the following lemma.

\begin{statement}{Lemma 6.23.3}
(cf. SGA $4^{1/2}$ Report 4.7). Let $P$ be a $G_V$-module locally a direct factor of a free module of finite type, let $P'=i_!P$, and let $Q$ be a complex of $\Lambda[G]$-modules, perfect over $\Lambda$. Then $P\lten_{\Lambda}Q$, considered as an object of $D({}_{\Lambda[G]}V)$ by the diagonal action of $G$, is perfect over $\Lambda[G]$, and there is a commutative square
\begin{equation}\tag{*}
\begin{tikzcd}[column sep=huge,row sep=large]
\delta^*\RHom_{\Lambda[G]}(p_1^*P',p_2^!P')\lten_{\Lambda}
\RHom_{\Lambda[G]}(Q,Q) \arrow[r,"(1)"] \arrow[d,"(2)"']
& K_Y \arrow[d,equal]\\
\delta^*\RHom_{\Lambda[G]}(p_1^*P'\lten_{\Lambda}Q,p_2^!P'\lten_{\Lambda}Q) \arrow[r,"(3)"']
& K_Y ,
\end{tikzcd}
\end{equation}
where $\delta:Y\to Y\times_SY$ is the diagonal, the upper horizontal arrow is the tensor product of the arrows
\[
\delta^*\RHom_{\Lambda[G]}(p_1^*P',p_2^!P')\xrightarrow{6.5.3}
K\Lambda[G]_{Y,\natural}\xrightarrow{x\mapsto x(e)}K_Y
\]
and
\[
\RHom_{\Lambda[G]}(Q,Q)\xrightarrow{\Tr_{\Lambda}}\Lambda,
\]
the vertical arrow is deduced from the external tensor product by restriction along the diagonal
\[
\Lambda[G]\to\Lambda[G]\otimes\Lambda[G],
\]
and the lower horizontal arrow is the composite of 6.5.3 (defined because $P'\lten_{\Lambda}Q\in\ob D_{\mathrm{ctf}}({}_{\Lambda[G]}Y)$) and of $x\mapsto x(e)$.\footnote{One should be able to replace $P$ by a perfect complex (over $\Lambda[G]$) (or even $P'$ by an object of $D_{\mathrm{ctf}}({}_{\Lambda[G]}Y)$).}
\end{statement}

The first assertion is evident, for one reduces to $P=\Lambda[G]_V$, and $P\lten_{\Lambda}Q$ is then isomorphic to the tensor product of $P$ by $Q$ considered with the trivial action of $G$. For the second, first note that, by 6.5.4 and 6.5.5, one has
\[
\delta^*\RHom_{\Lambda[G]}(p_1^*P',p_2^!P')
\xrightarrow{\sim}
D_{\Lambda[G]}(i_!P)\lten_{\Lambda[G]} i_!P
\xrightarrow{\sim}
i_!(D_{\Lambda[G]}P\lten_{\Lambda[G]}P),
\]
so that we are reduced to checking that the square induced by (*) on $V$ commutes. We may therefore suppose that $V=Y$. The square (*) then reads
\begin{equation}\tag{**}
\begin{tikzcd}[column sep=huge,row sep=large]
\RHom_{\Lambda[G]}(P,K_Y\lten P)\lten\RHom_{\Lambda[G]}(Q,Q) \arrow[r,"(1)"] \arrow[d,"(2)"']
& K_Y \arrow[d,equal]\\
\RHom_{\Lambda[G]}(P\lten Q,K_Y\lten P\lten Q) \arrow[r,"(3)"']
& K_Y ,
\end{tikzcd}
\end{equation}
where (2) is the external tensor product, (1) is given by
\[
\Tr_{\Lambda[G]}(-)(e)\otimes\Tr_{\Lambda},
\]
(traces in the sense of no.~5), and (3) by $\Tr_{\Lambda[G]}(-)(e)$. By descent, one reduces to $P=\Lambda[G]_Y$, and, by the initial remark, one may then suppose that $G$ acts trivially on $Q$. The assertion is then an immediate consequence of 5.12.5.

\subsection*{6.24}
If $\Lambda$ is annihilated by a power of a prime number $\ell\ne p$, the right hand side of 6.23.1 may be rewritten in the form
\[
\sum_{g\in G_{\natural}}
\Tr_{\mathbb Z_{\ell}[G]}(f_*\underline c_{\ell,U,X})(g^{-1})
\Tr_{\Lambda}(gv_0),
\]
with the notation of 6.17. Formula 6.17.1 shows that the local terms defined by Grothendieck in (XII 6.2) coincide with Verdier's local terms, at least when the endomorphism $\varphi$ of $G$ considered in loc. cit. is the identity; but, as already observed, this restriction is not serious. If $Y$ is proper over $S$, then by conjugating 6.23.1 with the Lefschetz formula for $Y\to S$ (III 4.7), one obtains a formula generalizing (XII 6.3), and a fortiori the Lefschetz formula for Frobenius proved in (SGA $4^{1/2}$ Report) and Part I of the present exposé.

\section*{Bibliography}

\begin{enumerate}[label={[\arabic*]}]
\item Cartan, H. and Eilenberg, S. -- \textit{Homological Algebra}, Princeton University Press, 1956.
\item Grothendieck, A. -- Formule de Lefschetz et rationalité des fonctions $L$, Séminaire Bourbaki 64/65, no. 279, in \textit{Dix exposés sur la cohomologie des schémas}, North Holland, Pub. Co., 1968.
\item Illusie, L. -- Complexe cotangent et déformations I, II, Lecture Notes in Mathematics 239, 283, Springer-Verlag, 1971, 1972.
\item Langlands, R. P. -- Modular forms and $\ell$-adic representations, in \textit{Modular Functions of One Variable II}, Lecture Notes in Mathematics 349, Springer-Verlag, 1973.
\item Stallings, J. -- Centerless groups -- an algebraic formulation of Gottlieb's theorem, \textit{Topology} 4, p. 129--134, 1965.
\item Verdier, J.-L. -- The Lefschetz fixed point formula in étale cohomology, in \textit{Proceedings of a conference on local fields}, Springer-Verlag, 1967.
\end{enumerate}

SGA $4^{1/2}$: \textit{Cohomologie étale}, by P. Deligne, Lecture Notes in Mathematics 569, Springer-Verlag, 1977.

In SGA $4^{1/2}$:
\[
\text{Finiteness Theorem: finiteness theorems in }\ell\text{-adic cohomology;}
\]
\[
\text{Cycle: the cohomology class associated with a cycle;}
\]
\[
\text{Report: report on the trace formula.}
\]


\phantomsection\addcontentsline{toc}{section}{Exposé V --- $J$-adic Projective Systems}
\section*{Exposé V. $J$-adic Projective Systems}
\begin{center}by J.-P. Jouanolou\end{center}

Throughout this exposé, the letter $A$ denotes a ring with unit. The present exposé is intended to develop a formalism which will be used in the following exposé for the study of $\ell$-adic cohomology.

\subsection*{1. Generalities on abelian $A$-categories}

\begin{statement}{Definition 1.1}
An $A$-category (resp. an abelian $A$-category) is a pair $(\cC,\rho)$ formed by an additive (resp. abelian) category $\cC$ and a homomorphism of rings
\[
\rho:A\longrightarrow \End(\id_{\cC})
\]
from $A$ to the ring of endomorphisms of the identity functor from $\cC$ to $\cC$.
\end{statement}

Given two $A$-categories (resp. abelian $A$-categories) $(\cC,\rho)$ and $(\cC',\rho')$, a morphism of $A$-categories (resp. of abelian $A$-categories) from $(\cC,\rho)$ to $(\cC',\rho')$ is an additive (resp. exact) functor $F:\cC\to\cC'$ such that, for every object $M$ of $\cC$ and every element $a$ of $A$, the equality below is satisfied:
\[
F(\rho(a)_M)=\rho'(a)_{F(M)}.
\]

\subsubsection*{Example 1.1.1}
Let $\cC$ be an additive (resp. abelian) category. Recall that the category of left $(A,\cC)$-modules, denoted $(A,\cC)$-mod, is the category described as follows:
\begin{enumerate}[label=(\roman*)]
\item Its objects are pairs $(X,\alpha)$ consisting of an object $X$ of $\cC$ and a homomorphism of rings
\[
\alpha:A\longrightarrow \Hom_{\cC}(X,X).
\]
\item A morphism $(X,\alpha)\to(Y,\beta)$ is a $\cC$-morphism $u:X\to Y$ such that, for every element $a$ of $A$, the following diagram is commutative:
\[
\begin{tikzcd}
X \arrow[r,"u"] \arrow[d,"\alpha(a)"'] & Y \arrow[d,"\beta(a)"]\\
X \arrow[r,"u"'] & Y .
\end{tikzcd}
\]
\end{enumerate}
The category $(A,\cC)$-mod is evidently an $A^{\natural}$-category (resp. an abelian $A^{\natural}$-category), where $A^{\natural}$ is the centre of $A$. Moreover, it has the following properties:
\begin{enumerate}[label=(\alph*)]
\item The functor $(A,\cC)$-mod $\to\cC$ obtained by forgetting the $A$-module structure commutes with inductive and projective limits.
\item Let $I$ be a category. In order that the $I$-indexed inductive (resp. projective) limits of $(A,\cC)$-mod be representable, it is enough that this be so for the $I$-indexed inductive (resp. projective) limits of $\cC$.
\end{enumerate}

\subsubsection*{1.2. External tensor product}
Let $(\cC,\rho)$ be an abelian $A$-category. If $M$ is a right $A$-module and $X$ an object of $\cC$, the external tensor product of $X$ by $M$, denoted $M\otimes_A X$, is the covariant functor from $\cC$ to the category of sets, defined on objects $Y$ of $\cC$ by the equality below and in the evident way on arrows:
\[
(M\otimes_A X)(Y)=\Hom_A(M,\Hom_{\cC}(X,Y)),
\]
where $\Hom_{\cC}(X,Y)$ is regarded as a right $A$-module via the action of $A$ on $X$. Interpreting $M\otimes_AX$ and $Y$ as objects of the category $\widehat{\cC}=\Hom(\cC,\Ens)^{\circ}$, the preceding equality takes the more familiar form
\[
\Hom_{\widehat{\cC}}(M\otimes_A X,Y)=\uHom_A(M,\Hom_{\cC}(X,Y)).
\]
If the functor $M\otimes_A X$ is representable, a representative will be denoted and named in the same way.

If $u:M\to M'$ and $v:X\to X'$ are two morphisms, respectively in $(A\text{-mod})$ and in $\cC$, one defines in the usual way a $\cC$-morphism
\[
u\otimes_A v:M\otimes_A X\longrightarrow M'\otimes_A X',
\]
which allows us to speak of the external tensor-product bifunctor, covariant in both variables. Its properties are summarized in the following proposition.

\begin{statement}{Proposition 1.2.1}
The tensor-product bifunctor is additive and commutes with inductive limits.
\end{statement}
The proof, classical in the case of modules over a ring, is immediate.

\begin{statement}{Proposition 1.2.2 (Representability criteria)}
\begin{enumerate}[label=(\alph*)]
\item For every object $X$ of $\cC$, $A\otimes_AX$ is represented by $X$.
\item For every finitely presented $A$-module $M$ and every object $X$ of $\cC$, the functor $M\otimes_AX$ is representable.
\item If the category $\cC$ has arbitrary direct sums, the functor $M\otimes_A X$ is representable for every pair $(M,X)$.
\end{enumerate}
\end{statement}

\begin{prooftext}[Proof]
Assertion (a) is clear. The other two assertions then follow from the lemma below, whose proof is immediate.
\end{prooftext}

\begin{statement}{Lemma 1.2.3}
Let $M\to N\to P\to0$ be an exact sequence of $A$-modules and let $X$ be an object of $\cC$. If the functors $M\otimes_A X$ and $N\otimes_A X$ are represented by objects $Y$ and $Z$ of $\cC$, then the functor $P\otimes_A X$ is represented by the cokernel of the morphism $Y\to Z$ corresponding to $\alpha\otimes_A\id_X$.
\end{statement}

\subsubsection*{Remark 1.2.4}
Let $\cG$ be an ideal of finite type of $A$. For every object $X$ of $\cC$, the functor $(A/\cG)\otimes_A X$ is representable, and the arrow
\[
p_X:X\longrightarrow (A/\cG)\otimes_A X
\]
coming from the canonical epimorphism from $A$ onto $A/\cG$ and from part (a) of 1.2.2 is an epimorphism. The same statement remains valid when the ideal $\cG$ is no longer assumed to be of finite type, but $\cC$ has arbitrary direct sums. Put $\cG X=\Ker(p_X)$. When the category has filtered inductive limits, it is clear that
\[
\cG X=\sum_{a\in\cG}aX.
\]
If $u:X\to Y$ is an epimorphism, $u$ maps $\cG X$ epimorphically onto $\cG Y$. Finally, for an object $X$ to be annihilated by every element of $\cG$, it is necessary and sufficient that $\cG X=0$. Then every morphism from an object $Z$ to $X$ factors, by means of $p_Z$, through a morphism from $(A/\cG)\otimes_A Z$ to $X$.

\subsection*{2. The Mittag-Leffler-Artin-Rees condition}
Throughout this paragraph, the letter $\cC$ denotes an abelian category and we put
\[
P=\Hom(\NN^\circ,\cC),
\]
the category of projective systems indexed by the ordered set $\NN$ of positive integers and with values in $\cC$.

\subsubsection*{2.1. The Mittag-Leffler-Artin-Rees condition}
Let $X=(X_n,u_n)_{n\ge0}$ be an object of the category $P$. If $r$ denotes an integer $\ge0$, we agree to write $X[r]$ for the projective system $(X_{n+r},u_{n+r})_{n\ge0}$. A morphism $h:X\to Y$ of $P$ defines in the evident way a morphism denoted $h[r]$ from $X[r]$ to $Y[r]$, functorially. If $r$ and $s$ are two integers satisfying $s\ge r\ge0$, the transition morphisms $X_{n+s}\to X_{n+r}$ $(n\ge0)$ define a morphism
\[
w_{sr}:X[s]\longrightarrow X[r],
\]
and it is clear that, if $t$ is a third integer with $t\ge s\ge r$, the equality below is satisfied:
\[
w_{ts}\circ w_{sr}=w_{tr}.
\]
In other words, we have thus associated to $X$ an object $\overline X$ of the category $\Hom(\NN^\circ,P)$. Moreover, if $h:X\to Y$ is a morphism of $P$, there corresponds to it in the evident way a morphism $\overline h$ from $\overline X$ to $\overline Y$, functorially. Finally, if
\[
0\longrightarrow X\stackrel{u}{\longrightarrow}Y\stackrel{v}{\longrightarrow}Z\longrightarrow0
\]
is an exact sequence of $P$, the corresponding sequence below in $\overline P=\Hom(\NN^\circ,P)$ is exact:
\[
0\longrightarrow\overline X\stackrel{\overline u}{\longrightarrow}\overline Y\stackrel{\overline v}{\longrightarrow}\overline Z\longrightarrow0.
\]

\begin{statement}{Definition 2.1.1}
Let an object of the category $P$ be given. It is said to satisfy the Mittag-Leffler-Artin-Rees condition (abbreviated MLAR) if there exists an integer $r$ such that, for every integer $s\ge r$, the equality below is satisfied:
\[
\Im(X[s]\xrightarrow{w_{s0}}X)=\Im(X[r]\xrightarrow{w_{r0}}X).
\]
\end{statement}
Equivalently, for every integer $n\ge0$, the images of $X_{n+r}$ and $X_{n+s}$ in $X_n$ by the ``transition morphisms'' are the same. Finally, a third interpretation is that the associated projective system $\overline X$ satisfies the Mittag-Leffler condition (abbreviated ML) (EGA $0_{\mathrm{III}}$ 13.1.2). The following proposition is deduced from this by (EGA $0_{\mathrm{III}}$ 13.2.1).

\begin{statement}{Proposition 2.1.2}
Let
\[
0\longrightarrow X\longrightarrow Y\longrightarrow Z\longrightarrow0
\]
be an exact sequence of $P$.
\begin{enumerate}[label=(\roman*)]
\item If $Y$ satisfies (MLAR), then $Z$ satisfies (MLAR).
\item If $X$ and $Z$ satisfy (MLAR), then $Y$ satisfies (MLAR).
\end{enumerate}
\end{statement}

\subsubsection*{2.2. AR-null projective systems}
\begin{statement}{Definition 2.2.1}
Let $X$ be an object of the category $P$. It is called AR-null if there exists an integer $r\ge0$ such that the canonical morphism $X[r]\to X$ is zero.
\end{statement}
It is clear in particular that an AR-null projective system satisfies the condition (MLAR).

\begin{statement}{Proposition 2.2.2}
Let
\[
0\longrightarrow X\longrightarrow Y\longrightarrow Z\longrightarrow0
\]
be an exact sequence of $P$. The following assertions are equivalent:
\begin{enumerate}[label=(\roman*)]
\item $Y$ is AR-null;
\item $X$ and $Z$ are AR-null.
\end{enumerate}
In other words, the full subcategory $P_0$ of $P$, whose objects are the AR-null projective systems, is thick in $P$.
\end{statement}

\begin{prooftext}[Proof]
It is clear that (i) implies (ii). Conversely, let $r$ and $t$ be integers satisfying the condition of Definition 2.2.1 for $X$ and $Z$ respectively. One sees immediately that the integer $r+t$ satisfies the conditions of Definition 2.2.1 for $Y$.
\end{prooftext}

\subsubsection*{2.3. Reminders on the calculus of fractions and applications}
Let $\cE$ be a category and $S$ a subset of $\Fl(\cE)$. A category of fractions of $\cE$ with respect to $S$ is a pair $(p,\cE(S^{-1}))$ consisting of a category $\cE(S^{-1})$ and a functor $p:\cE\to\cE(S^{-1})$ satisfying the following universal property: for a functor $u$ from $\cE$ to a category $\cF$ to factor through $p$, it is necessary and sufficient that it make all arrows belonging to $S$ invertible; the factorization is then unique. The uniqueness of such a pair up to isomorphism is clear. We refer the reader to Exposé I of the Séminaire Homotopique de Strasbourg for the proof of existence, and also for most of the results stated without proof in what follows in this paragraph. See also P. Gabriel and M. Zisman, \emph{Calculus of fractions and homotopy theory}, Erg. der Math., 35, Springer-Verlag, 1967.

Certain properties of the multiplicative part $S$ ensure a simple description of the category $\cE(S^{-1})$.

\begin{statement}{Definition 2.3.1}
A subset $S$ of $\Fl(\cE)$ is said to admit a right calculus of fractions if the following conditions hold:
\begin{enumerate}[label=(\alph*)]
\item The identity arrows of $\cE$ belong to $S$.
\item If $u:X\to Y$ and $v:Y\to Z$ belong to $S$, the composite $v\circ u$ belongs to $S$.
\item Every diagram of type $(D_1)$ below, in which $t$ belongs to $S$, can be completed to a diagram of type $(D_2)$, in which $s$ belongs to $S$:
\[
\begin{tikzcd}[column sep=large]
(D_1)& X \arrow[r,"u"] & Y & Y' \arrow[l,"t"']
\end{tikzcd}
\qquad
\begin{tikzcd}[column sep=large,row sep=large]
(D_2)& X' \arrow[r,"u'"] \arrow[d,"s"'] & Y' \arrow[d,"t"]\\
& X \arrow[r,"u"'] & Y .
\end{tikzcd}
\]
\item For every commutative diagram
\[
\begin{tikzcd}[column sep=large]
X \arrow[r,shift left,"f"] \arrow[r,shift right,"g"'] & Y \arrow[r,"s"] & Z,
\end{tikzcd}
\]
in which $s$ belongs to $S$, there exists a commutative diagram
\[
\begin{tikzcd}[column sep=large]
T \arrow[r,"t"'] & X \arrow[r,shift left,"f"] \arrow[r,shift right,"g"'] & Y,
\end{tikzcd}
\]
in which $t$ belongs to $S$.
\end{enumerate}
\end{statement}

Dually one defines the following notion.

\begin{statement}{Definition 2.3.2}
A subset $S$ of $\Fl(\cE)$ is said to admit a left calculus of fractions if the following conditions hold:
\begin{enumerate}[label=(\alph*)]
\item[(a')] The identity arrows of $\cE$ belong to $S$.
\item[(b')] The composite of two arrows belonging to $S$ belongs to $S$.
\item[(c')] Every diagram of type $(D_1)'$ below, in which $s$ belongs to $S$, can be completed to a diagram of type $(D_2)'$, in which $t$ belongs to $S$:
\[
\begin{tikzcd}[column sep=large,row sep=large]
(D_1)'& X' \arrow[r,"u'"] \arrow[d,"s"'] & Y'\\
& X
\end{tikzcd}
\qquad
\begin{tikzcd}[column sep=large,row sep=large]
(D_2)'& X' \arrow[r,"u'"] \arrow[d,"s"'] & Y' \arrow[d,"t"]\\
& X \arrow[r,"u"'] & Y .
\end{tikzcd}
\]
\item[(d')] For every commutative diagram $\Delta$ of the form below, in which $s$ belongs to $S$, there exists a morphism $t$ with source $Y$ and belonging to $S$ such that $t\circ f=t\circ g$:
\[
(\Delta)\quad
\begin{tikzcd}[column sep=large]
X' \arrow[r,"s"] & X \arrow[r,shift left,"f"] \arrow[r,shift right,"g"'] & Y .
\end{tikzcd}
\]
\end{enumerate}
\end{statement}

\subsubsection*{2.3.3}
Suppose that $S$ is a subset of $\Fl(\cE)$ admitting a right calculus of fractions. It follows from part (d) of 2.3.1 that, for every object $X$ of $\cE$, the category $(S/X)^\circ$, opposite to the category of arrows over $X$ belonging to $S$, is filtering (SGA 4 I 2.7). Consider then the category $\cE_S$ described as follows:
\begin{itemize}
\item the objects of $\cE_S$ are the objects of $\cE$;
\item if $X$ and $Y$ are two objects of $\cE$,
\[
\Hom_{\cE_S}(X,Y)=\varinjlim_{(S/X)^\circ}\Hom_{\cE}(\source(\cdot),Y).
\]
\end{itemize}
If $P_S$ denotes the evident functor from $\cE$ to $\cE_S$, the pair $(P_S,\cE_S)$ is a solution of the universal problem described at the beginning of the paragraph. In this case we agree to write $\cE(S^{-1})$ for this particular solution.

\subsubsection*{2.3.4}
Let now $S$ be a subset of $\Fl(\cE)$ admitting a left calculus of fractions. Then the functor $P_S:\cE\to\cE(S^{-1})$ commutes with finite inductive limits. In particular, if $\cE$ is additive, $\cE(S^{-1})$ is additive and $P_S$ is an additive functor.

\subsubsection*{2.3.5}
If $\cE$ is an abelian category and if $S$ is a subset of $\Fl(\cE)$ admitting both a right calculus of fractions and a left calculus of fractions, the category $\cE(S^{-1})$ is abelian and the functor $P_S$ is exact. Moreover, the pair $(P_S,\cE(S^{-1}))$ is a quotient of $\cE$ by the thick abelian subcategory $P_S^{-1}(0)$.

\subsubsection*{2.4. The category $P_{AR}=\Hom_{AR}(\NN^\circ,\cC)$}
Returning to the notation of 2.1, let $AR$ be the set of canonical arrows (cf. 2.1)
\[
v_{rX}:X[r]\longrightarrow X,
\]
where $X$ runs through the set of objects of $P$ and $r$ through the set of integers $\ge0$.

\begin{statement}{Proposition 2.4.1}
The subset $AR$ of $\Fl(P)$ admits both a right calculus of fractions and a left calculus of fractions.
\end{statement}

\begin{prooftext}[Proof]
The parts (a) $\Rightarrow$ (a') and (b) $\Rightarrow$ (b') of 2.3.1 and 2.3.2 are clear.

We prove (c). Let $(D_1)$ be the commutative diagram below, in which $v_{rY}$ denotes the canonical arrow from $Y[r]$ to $Y$:
\[
\begin{tikzcd}[column sep=large,row sep=large]
& Y[r] \arrow[d,"{v_{rY}}"]\\
X \arrow[r] & Y .
\end{tikzcd}
\]
The diagram $(D_2)$ below is commutative, and $v_{rX}$, the canonical morphism from $X[r]$ to $X$, belongs to $AR$, whence the assertion:
\[
\begin{tikzcd}[column sep=large,row sep=large]
X[r] \arrow[r,"{u[r]}"] \arrow[d,"{v_{rX}}"'] & Y[r] \arrow[d,"{v_{rY}}"]\\
X \arrow[r,"u"'] & Y .
\end{tikzcd}
\]
We prove (d). Let a commutative diagram be given:
\[
\begin{tikzcd}[column sep=large]
X \arrow[r,shift left,"f"] \arrow[r,shift right,"g"'] & Z[r] \arrow[r,"{v_{rZ}}"] & Z .
\end{tikzcd}
\]
Applying translation by $r$, one obtains a commutative diagram
\[
\begin{tikzcd}[column sep=large]
X[r] \arrow[r,shift left,"{f[r]}"] \arrow[r,shift right,"{g[r]}"'] & Z[2r] \arrow[r,"{v_{rZ}[r]}"] & Z[r].
\end{tikzcd}
\]
It is clear that the following diagram is commutative, so that $f\circ v_{rX}=g\circ v_{rX}$:
\[
\begin{tikzcd}[column sep=large,row sep=large]
X[r] \arrow[r,shift left,"{f[r]}"] \arrow[r,shift right,"{g[r]}"'] \arrow[d,"{v_{rX}}"'] & Z[2r] \arrow[d,"{v_{rZ}[r]}"]\\
X \arrow[r,shift left,"f"] \arrow[r,shift right,"g"'] & Z[r] .
\end{tikzcd}
\]

We prove (c'). First observe that, if $Y$ denotes an object of $P$ and $r$ a positive integer, there exists an object $T$ of $P$ such that $Y=T[r]$ and such that, for every object $X$ of $P$, the homomorphism given by translation by $r$
\[
\Hom_P(X,T)\longrightarrow\Hom_P(X[r],Y)
\]
is a bijection. The projective system $(T_n,w_n)$ defined as follows plainly works. If $Y=(Y_n,u_n)_{n\ge0}$, then $T_n$ is the zero object of $\cC$ for $n\ge r-1$ and is $Y_{n-r}$ otherwise, and
\[
w_n=0\quad\text{for }n\ge r-1,\qquad w_n=u_{n-r}.
\]

This said, suppose a diagram of the following type is given, and show that it can be completed as indicated in the statement of (c'):
\[
\begin{tikzcd}[column sep=large,row sep=large]
X[r] \arrow[r,"u'"] \arrow[d,"{v_{rX}}"'] & Y'\\
X .
\end{tikzcd}
\]
For this, let $Y$ be a projective system such that $Y'=Y[r]$ and such that the canonical map
\[
\Hom_P(X,Y)\longrightarrow\Hom_P(X[r],Y')
\]
is a bijection. Then it is clear that the following diagram is commutative, where $u$ denotes the unique morphism from $X$ to $Y$ such that $u'=u[r]$:
\[
\begin{tikzcd}[column sep=large,row sep=large]
X[r] \arrow[r,"u'"] \arrow[d,"{v_{rX}}"'] & Y'=Y[r] \arrow[d,"{v_{rY}}"]\\
X \arrow[r,"u"'] & Y .
\end{tikzcd}
\]

We prove (d'). Consider the commutative diagram below, and show that there exists a morphism $w$ from $Y$ such that $w\circ f=w\circ g$:
\[
\begin{tikzcd}[column sep=large]
X[r] \arrow[r,"{v_{rX}}"] & X \arrow[r,shift left,"f"] \arrow[r,shift right,"g"'] & Y .
\end{tikzcd}
\]
It is first of all clear that the following diagram is commutative:
\[
\begin{tikzcd}[column sep=large,row sep=large]
X[r] \arrow[r,shift left,"{f[r]}"] \arrow[r,shift right,"{g[r]}"'] \arrow[d,"{v_{rX}}"'] & Y[r] \arrow[d,"{v_{rY}}"]\\
X \arrow[r,shift left,"f"] \arrow[r,shift right,"g"'] & Y .
\end{tikzcd}
\]
In particular $v_{rY}\circ f[r]=v_{rY}\circ g[r]$. Let now $T$ be the system associated to $Y$ and $r$ as at the beginning of the proof of (c'), so that in particular $Y=T[r]$. The translates by $r$ of the morphisms $v_{rT}\circ f$ and $v_{rT}\circ g$ are respectively equal to $v_{rY}\circ f$ and $v_{rY}\circ g[r]$, hence are equal. It follows from the fact that translation by $r$ is injective on morphisms from $X$ to $T$ that the equality below holds, whence the assertion:
\[
v_{rT}\circ f=v_{rT}\circ g.
\]
\end{prooftext}

\begin{statement}{Definition 2.4.2}
Let $\cC$ be an abelian category. The category of projective systems of $\cC$ up to translation, denoted $\Hom_{AR}(\NN^\circ,\cC)$, is the category described as follows:
\begin{enumerate}[label=(\roman*)]
\item The objects of $\Hom_{AR}(\NN^\circ,\cC)$ are the objects of $\Hom(\NN^\circ,\cC)$.
\item If $X$ and $Y$ are two objects of $\Hom(\NN^\circ,\cC)$, a morphism from $X$ to $Y$ in the category $\Hom_{AR}(\NN^\circ,\cC)$ is an element of the set
\[
\Hom_{AR}(X,Y)=\varinjlim_r\Hom(X[r],Y),
\]
where, for every pair $(r,s)$ of integers with $s\ge r$, the map
\[
\Hom(X[r],Y)\longrightarrow\Hom(X[s],Y)
\]
comes from the canonical morphism
\[
X[s]\longrightarrow X[r].
\]
\end{enumerate}
The composition of morphisms is defined in the evident way.
\end{statement}
In what follows, the category $\cC$ being fixed, we put, to simplify notation,
\[
P=\Hom(\NN^\circ,\cC),\qquad P_{AR}=\Hom_{AR}(\NN^\circ,\cC).
\]

\subsubsection*{2.4.3}
We have seen (2.4.1) that the subset $AR$ of $\Fl(P)$ admits a right calculus of fractions. This allows us (2.3.3) to consider the category
\[
P'_{AR}=P(AR^{-1}),
\]
which, recall, has the same objects as $P$. Let $X$ and $Y$ be two objects of $P$, and denote by $\Hom'_{AR}(X,Y)$ the set of isomorphisms from $X$ to $Y$ in $P'_{AR}$. An element of $\Hom_{AR}(X,Y)$ defines in the evident way an element of $\Hom'_{AR}(X,Y)$. More precisely, one thus defines a functor
\[
m:P_{AR}\longrightarrow P'_{AR}.
\]

\begin{statement}{Proposition 2.4.3}
The functor $m$ is an isomorphism of categories.
\end{statement}

\begin{prooftext}[Proof]
We must see that, for every pair $(X,Y)$ of objects of $P$, the corresponding map
\[
m_{XY}:\Hom_{AR}(X,Y)\longrightarrow\Hom'_{AR}(X,Y)
\]
is a bijection. It is plainly a surjection. For injectivity, one is immediately reduced to seeing that, for every commutative diagram of the type below, there exists an integer $t$ greater than or equal to $r$ and $s$ such that the composite of $\rho$ and the canonical morphism from $X[t]$ to $X[s]$ is the canonical morphism from $X[t]$ to $X[r]$:
\[
\begin{tikzcd}[column sep=large,row sep=large]
X[r] \arrow[d,"{v_{rX}}"'] & X[s] \arrow[l,"\rho"'] \arrow[dl,"{v_{sX}}"]\\
X & {}
\end{tikzcd}
\]
But this follows from commutativity of the diagram below, in which the letter $c$ denotes all canonical morphisms involved:
\[
\begin{tikzcd}[column sep=large,row sep=large]
X[r+s] \arrow[d,"c"'] & X[2s] \arrow[l,"\rho(s)"] \arrow[d,"c"] \arrow[dl,"c" description]\\
X[r] & X[s] \arrow[l,"\rho"]
\end{tikzcd}
\]
\end{prooftext}

\begin{statement}{Proposition 2.4.4}
\begin{enumerate}[label=(\roman*)]
\item The category $P_{AR}$ is abelian.
\item The functor $p_{AR}:P\to P_{AR}$ is exact.
\item For an object $X$ of $P$ to be AR-null, it is necessary and sufficient that $p_{AR}(X)=0$.
\item The pair $(p_{AR},P_{AR})$ is a quotient of the abelian category $P$ by the thick abelian subcategory of AR-null objects.
\end{enumerate}
\end{statement}

\begin{prooftext}[Proof]
If (iii) is admitted, the other assertions are consequences of (2.3.5). It follows from (2.3.4) that, for an object $X$ of $P$ to be AR-null, it is necessary and sufficient that $p_{AR}(\id_X)=0$, whence (iii).
\end{prooftext}

\begin{statement}{Corollary 2.4.5}
Let $u:X\to Y$ be a morphism in the category $P$. The following assertions are equivalent:
\begin{enumerate}[label=(\roman*)]
\item $p_{AR}(u)$ is an isomorphism.
\item The kernel and cokernel of $u$ are AR-null.
\end{enumerate}
A morphism $u:X\to Y$ in the category $P$ of projective systems, satisfying the equivalent conditions (i) and (ii), will be called an AR-isomorphism.
\end{statement}

\begin{statement}{Definition 2.4.6}
Two objects $X$ and $Y$ of the category $P$ are called AR-isomorphic if $p_{AR}(X)$ and $p_{AR}(Y)$ are isomorphic in the category $P_{AR}$. If $X$ and $Y$ are two AR-isomorphic objects of the category $P$, there need not exist in general an AR-isomorphism $u$ from $X$ to $Y$.
\end{statement}

\subsubsection*{2.5. Characterization of projective systems satisfying the Mittag-Leffler-Artin-Rees condition}
\begin{statement}{Proposition 2.5.1}
The following assertions for an object $X$ of $P$ are equivalent:
\begin{enumerate}[label=(\roman*)]
\item $X$ satisfies (MLAR) (2.1.1).
\item There exists a strict projective system $Y$ and an AR-isomorphism $u:Y\to X$ (2.4.5).
\item There exists a strict projective system $Y$ such that $X$ and $Y$ are isomorphic in the category $P_{AR}$.
\item $X$ satisfies the Mittag-Leffler condition and, denoting by $X'$ the projective system of universal images of $X$, the projective system $X/X'$ is AR-null.
\end{enumerate}
\end{statement}

\begin{prooftext}[Proof]
We shall prove the chain of implications
\[
(i)\Rightarrow(iv)\Rightarrow(iii)\Rightarrow(ii)\Rightarrow(i).
\]
$(i)\Rightarrow(iv)$. There exists an integer $m$ such that the canonical morphism $w_{mX}$ from $X[m]$ to $X$ factors through the canonical inclusion $i$ of $X'$ in $X$. It follows that $p_{AR}(i)$ is a direct epimorphism and therefore, since $p_{AR}$ is exact, that $p_{AR}(X/X')=0$. We conclude by (2.4.4) (iii).

$(iv)\Rightarrow(iii)$. Condition (iv) means, with the preceding notation, that $p_{AR}(i)$ is an epimorphism. But because $p_{AR}$ is exact, it is also a monomorphism, hence an isomorphism. For $Y$ one takes the projective system $X'$.

$(iii)\Rightarrow(ii)$. The $P_{AR}$-isomorphism from $Y$ to $X$ whose existence is assumed is the image under $p_{AR}$ of a $P$-morphism $Y[r]\to X$, where $r$ denotes a positive or zero integer. The conclusion follows from the fact that $Y[r]$ is strict.

$(ii)\Rightarrow(i)$. We may suppose $Y$ contained in $X$ and $u$ the canonical injection. Let $T$ be the cokernel of $u$. The hypothesis means that $T$ is AR-null, hence there exists an integer $r$ such that, for every $s\ge r$, the canonical morphism from $T[s]$ to $T$ is zero. It follows that, for $s\ge r$, the image of $X[s]$ in $X$ is contained in $Y$; since it also plainly contains the image, equal to $Y$, of $Y[s]$ in $Y$, it is equal to $Y$. Thus $Y$ is the image of $X[s]$ for $s$ sufficiently large, whence $X$ satisfies (MLAR), and $Y$ is the projective system of universal images of $X$.
\end{prooftext}

\begin{statement}{Corollary 2.5.2}
Let $u:Y\to X$ be an AR-isomorphism, with $Y$ strict. Under these conditions, $X$ satisfies (MLAR) and $u(Y)$ is the projective system of universal images of $X$.
\end{statement}

\subsection*{3. $J$-adic and AR-$J$-adic projective systems}
Let $A$ be a commutative ring with unit, let $J$ be an ideal of $A$, and let $(\cC,\rho)$ be an abelian $A$-category. We suppose either that $J$ is of finite type, or that $\cC$ has infinite direct sums. The notation is the same as in the preceding paragraph. In particular, the symbols $P$ and $P_{AR}$ denote respectively the categories $\Hom(\NN^\circ,\cC)$ and $\Hom_{AR}(\NN^\circ,\cC)$.

\subsubsection*{3.1. $J$-adic projective systems}
\begin{statement}{Definition 3.1.1}
Let $X$ be an object of $P$. It is called $J$-adic if it satisfies the following conditions:
\begin{enumerate}[label=(\roman*)]
\item for every integer $n\ge0$, $J^{n+1}X_n=0$;
\item for every pair $(m,n)$ of integers with $m\ge n\ge0$, the morphism
\[
A/J^{n+1}\otimes_A X_m\longrightarrow X_n
\]
deduced from the transition morphism $X_m\to X_n$ is an isomorphism.
\end{enumerate}
\end{statement}
We shall write $J$-ad$(\cC)$ for the full subcategory of $P$ generated by the $J$-adic projective systems.

\subsubsection*{Remark 3.1.2}
Let $I$ and $J$ be two ideals of $A$ and suppose either that $I$ and $J$ are of finite type, or that $\cC$ has infinite direct sums. Then one can show that, if $I$ and $J$ define the same topology on $A$, the categories $I$-ad$(\cC)$ and $J$-ad$(\cC)$ are equivalent.

\begin{statement}{Proposition 3.1.3 (Stability under exact sequences)}
Let
\[
0\longrightarrow X\stackrel{u}{\longrightarrow}Y\stackrel{v}{\longrightarrow}Z\longrightarrow0
\]
be an exact sequence of $P$.
\begin{enumerate}[label=(\roman*)]
\item If $X$ is strict and $Y$ is $J$-adic, then $Z$ is $J$-adic.
\item If $Y$ is strict and $Z$ is $J$-adic, then $X$ is strict.
\item If $X$ and $Z$ are $J$-adic, then $Y$ is $J$-adic.
\end{enumerate}
\end{statement}

\begin{prooftext}[Proof]
If $m\ge n$, it is clear that the diagram $(T)$ below, in which the vertical arrows denote the transition morphisms, is commutative and exact:
\[
\begin{tikzcd}[column sep=large,row sep=large]
{} & A/J^{n+1}\otimes_A X_m \arrow[r,"\id\otimes u_m"] \arrow[d,"f_{nm}"'] & A/J^{n+1}\otimes_A Y_m \arrow[r,"\id\otimes v_m"] \arrow[d,"g_{nm}"] & A/J^{n+1}\otimes_A Z_m \arrow[r] \arrow[d,"h_{nm}"] & 0\\
0 \arrow[r] & X_n \arrow[r] & Y_n \arrow[r] & Z_n \arrow[r] & 0 .
\end{tikzcd}
\]
In case (i), $f_{nm}$ is an epimorphism and $g_{nm}$ an isomorphism. In case (ii), $g_{nm}$ is an epimorphism and $h_{nm}$ an isomorphism. In case (iii), $f_{nm}$ and $h_{nm}$ are isomorphisms. In all three cases, the snake lemma applied to the diagram $(T)$ gives the conclusion.
\end{prooftext}

\subsubsection*{3.2. AR-$J$-adic projective systems}
\begin{statement}{Proposition 3.2.1}
The following assertions, for an object $X$ of $P$ such that $J^{n+1}X_n=0$ for every $n\ge0$, are equivalent:
\begin{enumerate}[label=(\roman*)]
\item There exists a $J$-adic projective system $Y$ and an AR-isomorphism $u:Y\to X$ (2.4.5).
\item There exists a $J$-adic projective system $Y$, isomorphic to $X$ in the category $P_{AR}$.
\end{enumerate}
\end{statement}

\begin{prooftext}[Proof]
We have to show that (ii) implies (i). A $P_{AR}$-isomorphism from $Y$ to $X$ is represented, for some integer $r$, by an AR-isomorphism $v:Y[r]\to X$. Since $J^{n+1}X_n=0$ for every $n\ge0$ and $Y$ is $J$-adic, $v$ is the composite of a $P$-morphism $u:Y\to X$ and the canonical morphism from $Y[r]$ to $Y$. It is then clear that $u$ is an AR-isomorphism, whence the assertion.
\end{prooftext}

\begin{statement}{Definition 3.2.2}
A projective system indexed by $\NN$ with values in $\cC$ is called AR-$J$-adic if:
\begin{enumerate}[label=(\alph*)]
\item $J^{n+1}X_n=0$ for every $n\ge0$;
\item it satisfies the equivalent conditions (i) and (ii) of Proposition 3.2.1.
\end{enumerate}
\end{statement}

\begin{statement}{Proposition 3.2.3}
Let $X$ be a projective system indexed by $\NN$ with values in $\cC$. Suppose that $J^{n+1}X_n=0$ for every $n\ge0$. Then, in order that $X$ be AR-$J$-adic, it is necessary and sufficient that it satisfy (MLAR) and that, denoting by $X'$ its projective system of universal images, there exist an integer $r\ge0$ such that, for every pair $(m,n)$ of integers with $m\ge n+r$, the ``transition morphism'' below is an isomorphism:
\[
X'_m/J^{n+1}X'_m\longrightarrow X'_{n+r}/J^{n+1}X'_{n+r}.
\]
\end{statement}

\begin{prooftext}[Proof]
\begin{enumerate}[label=(\alph*)]
\item The conditions of the statement are sufficient. If $X$ satisfies (MLAR), it is AR-isomorphic to its projective system of universal images $X'$. Moreover, the hypothesis made on $X'$ means that the projective system
\[
\bigl(X'_{n+r}/J^{n+1}X'_{n+r}\bigr)_{n\ge0},
\]
with the evident transition morphisms, is $J$-adic. It is clear that this projective system is AR-isomorphic (2.4.6) to $X'$.
\item The conditions of the statement are necessary. If $X$ is AR-$J$-adic, it is AR-isomorphic to a strict projective system, hence satisfies (MLAR). The projective system $X'$ of universal images of $X$ is AR-isomorphic to $X$, hence is AR-$J$-adic. To prove the second point, it remains to see that, if $X$ is a strict AR-$J$-adic projective system, there exists an integer $r$ such that the projective system $\ell_r(X)$ defined by
\[
\ell_r(X)_n=X_{n+r}/J^{n+1}X_{n+r}
\]
is $J$-adic. For this, let $u:Y\to X$ be an AR-isomorphism, with $Y$ $J$-adic. Since the kernel $T$ of $u$ is AR-null, let $r$ be an integer such that the canonical morphism $T[r]\to T$ is zero. It is immediate that the canonical morphism
\[
\ell_r(u):\ell_r(Y)\longrightarrow\ell_r(X)
\]
is an isomorphism, whence the assertion.
\end{enumerate}
\end{prooftext}

\begin{statement}{Proposition 3.2.4 (Stability under exact sequences)}
Let
\[
0\longrightarrow X\stackrel{u}{\longrightarrow}Y\stackrel{v}{\longrightarrow}Z\longrightarrow0
\]
be an exact sequence of $P$, the projective systems $X$, $Y$, and $Z$ being such that
\[
J^{n+1}X_n=J^{n+1}Y_n=J^{n+1}Z_n=0\qquad(n\ge0).
\]
\begin{enumerate}[label=(\roman*)]
\item If $X$ satisfies (MLAR) and $Y$ is AR-$J$-adic, then $Z$ is AR-$J$-adic.
\item If $Y$ satisfies (MLAR) and $Z$ is AR-$J$-adic, then $X$ satisfies (MLAR).
\end{enumerate}
\end{statement}

\begin{prooftext}[Proof]
We use freely, for every projective system $X$ and every integer $r$, the notation $\ell_r(X)$ already made explicit.

Proof of (i). If $X'$ and $Y'$ denote the projective systems of universal images of $X$ and $Y$, then $u$ maps $X'$ into $Y'$, whence the evident exact commutative diagram below, in which the arrows are evident:
\[
\begin{tikzcd}
0 \arrow[r] & X' \arrow[r] \arrow[d] & Y' \arrow[r] \arrow[d] & \widetilde{Z'} \arrow[r] \arrow[d] & 0\\
0 \arrow[r] & X \arrow[r] & Y \arrow[r] & Z \arrow[r] & 0 .
\end{tikzcd}
\]
Since the category $P_{AR}$ is abelian and the functor $P\to P_{AR}$ is exact, it follows from this diagram, and from the fact that the inclusions of $X'$ and $Y'$ into $X$ and $Y$ respectively are AR-isomorphisms (2.5.1), that the morphism $\widetilde Z\to Z$ is an AR-isomorphism.

This allows us to reduce to the case where $X$ and $Y$, and therefore $Z$, are strict. Then there exists an integer $r$ such that the projective system $\ell_r(Y)$ is $J$-adic. The functor $\ell_r$ is plainly right exact, whence an exact sequence
\[
\ell_r(X)\stackrel{\ell_r(u)}{\longrightarrow}\ell_r(Y)\stackrel{\ell_r(v)}{\longrightarrow}\ell_r(Z)\longrightarrow0.
\]
After replacing $\ell_r(X)$ by its image under $\ell_r(u)$, we are under the hypotheses for applying (3.1.3 (i)), which shows that $\ell_r(Z)$ is $J$-adic, whence the assertion.

(ii) Let $Y'$ and $Z'$ be the projective systems of universal images of $Y$ and $Z$. Consider the evident exact commutative diagram below, in which the two vertical arrows on the right are the canonical inclusions:
\[
\begin{tikzcd}
0 \arrow[r] & \widetilde X \arrow[r] \arrow[d] & Y' \arrow[r] \arrow[d] & Z' \arrow[r] \arrow[d] & 0\\
0 \arrow[r] & X \arrow[r,"u"'] & Y \arrow[r,"v"'] & Z \arrow[r] & 0 .
\end{tikzcd}
\]
It shows that the morphism $\widetilde X\to X$ is an AR-isomorphism, which allows us to reduce to the case where $Y$ and $Z$ are strict. In this case, since $Z$ is strict and AR-$J$-adic, it follows from 3.2.3 that there exists an integer $r$ such that $\ell_r(Z)$ is $J$-adic. Considering the exact commutative diagram
\[
\begin{tikzcd}
0 \arrow[r] & \widetilde{\widetilde X} \arrow[r] \arrow[d] & \ell_r(Y) \arrow[r] \arrow[d] & \ell_r(Z) \arrow[r] \arrow[d] & 0\\
0 \arrow[r] & X \arrow[r,"u"'] & Y \arrow[r,"v"'] & Z \arrow[r] & 0,
\end{tikzcd}
\]
in which, as before, $\widetilde{\widetilde X}\to X$ is an AR-isomorphism, allows us to reduce to the case where $Y$ is strict and $Z$ is $J$-adic. Then (3.1.3 (ii)) gives the conclusion.
\end{prooftext}

\subsection*{4. Filtrations and graduations}
\subsubsection*{4.1. Generalities}
\subsubsection*{4.1.1}
Let $\cC$ be an abelian category and let $X=(X_n)_{n\in\ZZ}$ be a projective system indexed by $\ZZ$ with values in $\cC$. Suppose that
\[
X_i=0\quad(i<0).
\]
For every element $p$ of $\ZZ$, define a subobject $F^pX$ of $X$ in the category $P_1=\Hom(\ZZ^\circ,\cC)$ by setting
\[
(F^pX)_n=F^pX_n=
\begin{cases}
\Ker(X_n\to X_{p-1}),& n\ge p-1,\\
0,& n<p-1,
\end{cases}
\]
with the transition morphisms defined in the evident way. Thus one obtains a decreasing filtration $(F^pX)_{p\in\ZZ}$ of $X$ satisfying
\begin{equation}\tag{a}
F^pX=X\quad\text{if }p\le0,
\end{equation}
\begin{equation}\tag{b}
(F^pX)_{p-1}=0\quad\text{for all }p.
\end{equation}
Moreover, the snake lemma applied to the exact commutative diagram $(D)$ below proves that, as soon as $n\ge p$, the canonical morphism
\[
F^pX_n/F^{p+1}X_n\longrightarrow \Im(X_n\to X_p)\cap\Ker(X_p\to X_{p-1})
\]
is an isomorphism:
\[
\begin{tikzcd}[column sep=large,row sep=large]
0 \arrow[r] & \Ker(X_n\to X_p) \arrow[r] \arrow[d] & X_n \arrow[r] \arrow[d] & \Im(X_n\to X_p) \arrow[r] \arrow[d] & 0\\
0 \arrow[r] & 0 \arrow[r] & X_{p-1} \arrow[r,"\sim"'] & X_{p-1} \arrow[r] & 0 .
\end{tikzcd}
\]
If now $X$ is a projective system indexed by $\NN$ with values in $\cC$, we extend it to a projective system indexed by $\ZZ$ by putting $X_n=0$ for $n<0$, with the evident transition morphisms, which allows the preceding constructions to be applied to it. This allows us also to consider the category
\[
P=\Hom(\NN^\circ,\cC)
\]
as a full subcategory of $P_1$.

\begin{statement}{Definition 4.1.2}
Let $X$ be a projective system indexed by $\ZZ$. We say that $X$ is essentially constant if there exists an integer $p$ such that, for every pair $(r,s)$ of integers greater than or equal to $p$ and satisfying $s\ge r$, the transition morphism
\[
X_s\longrightarrow X_r
\]
is an isomorphism. It is clear that it is enough to require this when $r=p$.
\end{statement}

\begin{statement}{Proposition 4.1.3 (The property (MLAR) and associated gradeds)}
The following assertions for an object $X$ of the category $P$ are equivalent:
\begin{enumerate}[label=(\roman*)]
\item $X$ satisfies (MLAR).
\item The projective system with values in the category of objects of $\cC$ graded by $\NN$,
\[
r\longmapsto (\gr^\alpha(X_{\alpha+r}))_{\alpha\in\NN},
\]
with the evident transition morphisms, is essentially constant.
\item The projective system with values in the category of objects of $\cC$ graded by $\NN$,
\[
n\longmapsto (\gr^p(X_n))_{p\in\NN},
\]
with the evident transition morphisms, satisfies (MLAR).
\end{enumerate}
\end{statement}

\begin{prooftext}[Proof]
$(i)\Rightarrow(ii)$ follows from the expression of $F^pX_n/F^{p+1}X_n$ given in 4.1.1.

$(ii)\Rightarrow(iii)$. Let $r$ be an integer such that, for every $s>r$ and every $p$, the canonical morphism
\[
\gr^p(X_{p+s})\longrightarrow \gr^p(X_{p+r})
\]
is an isomorphism. We shall see that, for every pair $(m,n)$ of integers with $m\ge n+r$, the equality below is satisfied for every integer $p$:
\[
\Im(\gr^p(X_m)\to\gr^p(X_n))=
\Im(\gr^p(X_{n+r})\to\gr^p(X_n)).
\]
Indeed:

(a) If $p\ge n+1$, then $\gr^p(X_n)=0$,

(b) if $p\le n$, the chain of inequalities $m\ge n+r\ge p+r$ shows that the canonical morphism
\[
\gr^p(X_m)\longrightarrow\gr^p(X_{n+r})
\]
is an isomorphism, whence the assertion.

$(iii)\Rightarrow(i)$. Denote by $Y$ the projective system of (iii), and let $r$ be an integer such that
\[
\Im(Y[s]\to Y)=\Im(Y[r]\to Y)\qquad(s\ge r).
\]
If $n$ and $\ell$ are two integers, with $\ell\ge n$, the filtration induced on $\Im(X_\ell\to X_n)$ by that of $X_n$ is also the ``quotient'' filtration of that of $X_\ell$, and consequently, endowing $\Im(X_\ell\to X_n)$ with this filtration, the equality below is satisfied:
\[
\gr^p(\Im(X_\ell\to X_n))=
\Im(\gr^p(X_\ell)\to\gr^p(X_n)).
\]
In particular, condition (iii) is expressed by saying that, as soon as $\ell\ge n+r$, the canonical monomorphism
\[
j:\Im(X_\ell\to X_n)\longrightarrow\Im(X_{n+r}\to X_n)
\]
induces an isomorphism on associated gradeds. Since the gradings are finite, it follows that $j$ itself is an isomorphism, whence (i).
\end{prooftext}


\subsubsection*{4.1.4}
If \(X\) is an object of \(\mathcal P\) satisfying (ML), it follows from the expression of the associated graded recalled in (4.1.1) that, for every integer \(p\), the projective system
\[
\gr^p(X)=\bigl(\gr^p(X_n)\bigr)_{n\in\mathbb N}
\]
is essentially constant.

Its projective limit, which exists without any hypothesis on the abelian category \(\mathcal C\), will be called the \(p\)-th component of the strict associated graded of \(X\), and will be denoted
\[
\grs^p(X).
\]

\subsection*{4.2. Projective systems adapted to \(J\) and the Artin-Rees lemma}

We suppose given a commutative ring \(A\), an ideal \(J\) of \(A\), and an abelian \(A\)-category \(\mathcal C\).

\begin{statement}{Definition 4.2.1}
An object \(X\) of the category \(\mathcal P=\Hom(\mathbb N^\circ,\mathcal C)\) is said to be adapted to \(J\) if its filtration \((F^pX)_{p\in\mathbb N}\) as a projective system (4.1.1) is compatible with the \(J\)-adic filtration of \(X\), i.e. if for every pair \((p,q)\) of non-negative integers the inclusion below holds:
\[
J^pF^q(X)\subset F^{p+q}(X).
\]
\end{statement}

\begin{statement}{Lemma 4.2.2}
Let \(X\) be an object of the category \(\mathcal P\). The following assertions are equivalent:
\begin{enumerate}[label=\textup{(\roman*)}]
\item \(X\) is adapted to \(J\).
\item For every pair \((n,p)\) of integers with \(p\le n+1\), one has
\[
J^pF^{\,n-p+1}(X_n)=0 .
\]
\end{enumerate}
\end{statement}

\begin{prooftext}
\((i)\Rightarrow(ii)\): then
\[
J^pF^{\,n-p+1}(X_n)\subset F^{n+1}(X_n)=0.
\]

\((ii)\Rightarrow(i)\): if \(n<p+q\), \(F^{p+q}(X_n)=0\) by definition. Moreover
\[
J^pF^q(X_n)
\]
is contained in \(J^{\,n+1-q}F^q(X_n)\), and the latter is zero by hypothesis. If \(n>p+q\), the image of \(J^pF^q(X_n)\) in \(X_{p+q-1}\) under the transition morphism is contained in \(J^pF^q(X_{p+q-1})\), which is zero by hypothesis. The inclusion
\[
J^pF^q(X_n)\subset F^{p+q}(X_n)
\]
therefore follows by definition of the filtration of a projective system.

In particular, it follows from (4.2.2) that every projective subsystem of a projective system adapted to \(J\) is adapted to \(J\).
\end{prooftext}

\paragraph{Examples.}
a) A \(J\)-adic projective system is adapted to \(J\). It follows that every projective subsystem of a \(J\)-adic projective system is adapted to \(J\).

b) If a projective system \(X\) satisfies condition (ML) and is adapted to \(J\), then its projective system of universal images is also adapted to \(J\).

\subsubsection*{4.2.3. The associated graded of a projective system adapted to \(J\)}
Let \(X\) be a projective system adapted to \(J\). Since the filtration of \(X\) is compatible with \(J\), it is clear that the graded object
\[
\bigl(\gr^p(X)\bigr)_{p\in\mathbb N}
\]
is canonically endowed with the structure of a graded \(\gr_J(A)\)-module. In particular, by passage to the limit, one deduces that when \(X\) satisfies condition (ML), the graded object
\[
\grs(X)=\bigl(\grs^p(X)\bigr)_{p\in\mathbb N}
\]
is canonically endowed with the structure of a graded \(\gr_J(A)\)-module.

Let us briefly spell out how, when moreover \(X\) is strict, one defines the corresponding morphisms:
\[
\omega_{pq}: J^p/J^{p+1}\otimes_A F^q(X_q)\longrightarrow F^{p+q}(X_{p+q}).
\]
Since the projective system \(X\) is strict, we know (4.1.1 and 4.1.2) that the following canonical morphism is an isomorphism:
\[
\varphi:F^q(X_{p+q})/F^{q+1}(X_{p+q})\longrightarrow F^q(X_q).
\]
Moreover, the fact that the filtration of \(X\) is adapted to \(J\) gives a morphism
\[
J^p\otimes_A F^q(X_{p+q})\longrightarrow F^{p+q}(X_{p+q})
\]
which factors, since
\[
J^{p+1}F^q(X_{p+q})\subset F^{p+q+1}(X_{p+q})=0,\qquad
J^pF^{q+1}(X_{p+q})\subset F^{p+q+1}(X_{p+q})=0,
\]
through
\[
J^p/J^{p+1}\otimes_A
\bigl(F^q(X_{p+q})/F^{q+1}(X_{p+q})\bigr)
\]
as a morphism
\[
\beta:(J^p/J^{p+1})\otimes_A
\bigl(F^q(X_{p+q})/F^{q+1}(X_{p+q})\bigr)
\longrightarrow F^{p+q}(X_{p+q}).
\]
The morphism \(\omega_{pq}\) is then given by the following expression:
\[
\omega_{pq}=\beta\circ(\id_{J^p/J^{p+1}}\otimes_A\varphi)^{-1}.
\]

\begin{statement}{Lemma 4.2.4}
Let \(X\) be a \(J\)-adic projective system. The canonical morphism
\[
\gr_J(A)\otimes_A X_0\longrightarrow \grs(X),
\]
coming from the graded \(\gr_J(A)\)-module structure of \(\grs(X)\), is an epimorphism.
\end{statement}

\begin{prooftext}
This follows immediately from the description of the morphisms \(\omega_{pq}\) given in (4.2.3).
\end{prooftext}

\begin{statement}{Proposition 4.2.5}
Characterization of strict AR-\(J\)-adic projective systems by means of their associated gradeds. Let \(X\) be a strict projective system indexed by \(\mathbb N\), with values in \(\mathcal C\). Suppose that for every integer \(n\), one has
\[
J^{n+1}X_n=0.
\]
The following assertions are equivalent:
\begin{enumerate}[label=\textup{(\roman*)}]
\item \(X\) is AR-\(J\)-adic.
\item There exists an integer \(r>0\) such that, for every integer \(n\), the following inclusion holds:
\[
F^n(X_n)\subset J^{\,n-r}X_n.
\]
\item There exists an integer \(r\) such that, for every integer \(m\), the following canonical morphism is an isomorphism whenever \(n\ge m+r\):
\[
\gr^n(J^mX)\longrightarrow \gr^n(X).
\]
\end{enumerate}
In the statement of (ii), we agree to set \(J^{\,n-r}X_n=X_n\) for \(n<r\).
\end{statement}

\begin{prooftext}
The equivalence of (ii) and (iii) is clear, since \(X\) and \(J^mX\) are strict. Let us prove the equivalence of (i) and (ii). We know (3.2.3) that (i) is equivalent to the existence of an integer \(r\) such that, for every pair \((m,n)\) of integers with \(m\ge n+r\), the transition morphism below is an isomorphism:
\[
X_{m+1}/J^{n+1}X_{m+1}\longrightarrow X_m/J^{n+1}X_m.
\]
Consider the following exact commutative diagram:
\[
\begin{array}{ccccccccccc}
&&0&&0\\
&&\downarrow&&\downarrow\\
&&J^{n+1}X_{m+1}&\longrightarrow&J^{n+1}X_m&\longrightarrow&0\\
&&\downarrow&&\downarrow\\
0&\longrightarrow&F^{m+1}(X_{m+1})&\longrightarrow&X_{m+1}&\longrightarrow&X_m&\longrightarrow&0\\
&&\downarrow&&\downarrow\\
&&X_{m+1}/J^{n+1}X_{m+1}&\longrightarrow&X_m/J^{n+1}X_m&\longrightarrow&0\\
&&\downarrow&&\downarrow\\
&&0&&0 .
\end{array}
\]
It shows the existence of a canonical isomorphism between
\[
\Ker\bigl(X_{m+1}/J^{n+1}X_{m+1}
       \longrightarrow X_m/J^{n+1}X_m\bigr)
\]
and
\[
F^{m+1}(X_{m+1})/
\bigl(F^{m+1}(X_{m+1})\cap J^{n+1}X_{m+1}\bigr).
\]
In particular, for the transition morphism after reduction modulo \(J^{n+1}\) to be an isomorphism, it is necessary and sufficient that \(F^{m+1}(X_{m+1})\) be contained in \(J^{n+1}X_{m+1}\), which is precisely assertion (ii).
\end{prooftext}

\begin{statement}{Theorem 4.2.6}
Artin-Rees lemma. Let \(X\) be a projective system indexed by \(\mathbb N\), with values in \(\mathcal C\). Suppose that:
\begin{enumerate}[label=\textup{(\roman*)}]
\item \(X\) satisfies (MLAR);
\item \(X\) is adapted to \(J\);
\item \(\grs(X)\) is a graded \((\gr_J(A),\mathcal C)\)-module of finite type, i.e. for every increasing sequence \((M_n)_{n\in\mathbb N}\) of graded sub-\(\gr_J(A)\)-modules of \(\grs(X)\) whose union is \(\grs(X)\), there exists an integer \(p\) such that
\[
\grs(X)=M_p.
\]
\end{enumerate}
Then \(X\) is AR-\(J\)-adic.
\end{statement}

\begin{prooftext}
Recall that if \(X'\) is the projective system of universal images of \(X\), we have set
\[
\grs(X)=\gr(X').
\]
Moreover, \(X'\) is adapted to \(J\), so it remains to show that \(X'\) is AR-\(J\)-adic. Suppose that the following lemma has been proved.
\end{prooftext}

\begin{statement}{Lemma 4.2.7}
Let \(X\) be a strict projective system adapted to \(J\). For every integer \(s\), the graded subobject \(\bigl((M_s)_n\bigr)_{n\in\mathbb N}\) of \(\grs(X)\) defined by the equalities
\[
(M_s)_n=J^{\,n-s}X_n\cap F^n(X_n)\qquad(J^i=A\text{ if }i<0),
\]
is a graded \(\gr_J(A)\)-submodule of \(\grs(X)\).
\end{statement}

To prove the theorem, we may suppose \(X\) strict. Then \(\grs(X)\) is the union of the increasing sequence of the \(M_s\), hence is equal to one of them. This means that there exists an integer \(s_0\) such that
\[
F^n(X_n)\subset J^{\,n-s_0}X_n
\]
for every integer \(n\). We conclude by (4.2.5).

\begin{prooftext}[Proof of 4.2.7]
Return to the notation of (4.2.3). We have to see that the morphism
\[
\omega_{pq}:J^p/J^{p+1}\otimes_A F^q(X_q)\longrightarrow F^{p+q}(X_{p+q})
\]
sends
\[
J^p/J^{p+1}\otimes_A
\bigl(J^{q-s}X_q\cap F^q(X_q)\bigr)
\]
into
\[
J^{p+q-s}X_{p+q}\cap F^{p+q}(X_{p+q}).
\]
Now the canonical morphism
\[
\rho:X_{p+q}/F^{q+1}(X_{p+q})\longrightarrow X_q
\]
is an isomorphism and sends
\[
\bigl(J^{q-s}X_{p+q}+F^{q+1}(X_{p+q})\bigr)/F^{q+1}(X_{p+q})
\]
onto \(J^{q-s}X_q\). It follows that the inverse image under
\(\id(J^p/J^{p+1})\otimes_A\varphi\) (4.2.3) of
\[
J^p/J^{p+1}\otimes_A
\bigl(J^{q-s}X_q\cap F^q(X_q)\bigr)
\]
is
\[
J^p/J^{p+1}\otimes_A
\frac{(J^{q-s}X_{p+q}\cap F^q(X_{p+q}))+F^{q+1}(X_{p+q})}
     {F^{q+1}(X_{p+q})}.
\]
It is now clear that the image under \(\beta\) (4.2.3) of this last object is contained in
\[
J^{p+q-s}X_{p+q}\cap J^pF^q(X_{p+q}),
\]
hence a fortiori in
\[
J^{p+q-s}X_{p+q}\cap F^{p+q}(X_{p+q}).
\]
This proves the lemma.
\end{prooftext}

\subsection*{5. Noetherian \(J\)-adic and AR-\(J\)-adic projective systems}

Throughout this paragraph, a noetherian commutative ring \(A\), an ideal \(J\) of \(A\), and an abelian \(A\)-category \(\mathcal C\) are fixed.

\subsubsection*{5.1. Generalities}

\begin{statement}{Definition 5.1.1}
An object \(X\) of the category \(\Hom(\mathbb N^\circ,\mathcal C)\) is said to be noetherian \(J\)-adic (resp. artinian) if it is \(J\)-adic and if its components are noetherian (resp. artinian) objects of the category \(\mathcal C\).
\end{statement}

We shall denote by \(J\)-adn\((\mathcal C)\) (resp. \(J\)-adf\((\mathcal C)\)) the full subcategory of
\(\mathcal P=\Hom(\mathbb N^\circ,\mathcal C)\) whose objects are noetherian (resp. artinian) \(J\)-adic projective systems. It is clear (3.1.3) that an extension of noetherian (resp. artinian) \(J\)-adic projective systems is of the same type.

\begin{statement}{Proposition 5.1.2}
Let \(X\) be a \(J\)-adic object of \(\mathcal P\). In order for \(X\) to be noetherian (resp. artinian) \(J\)-adic, it is necessary and sufficient that it be \(J\)-adic and that \(X_0\) be noetherian (resp. artinian).
\end{statement}

\begin{prooftext}
Indeed, let \(X\) be a \(J\)-adic projective system such that \(X_0\) is noetherian (resp. artinian). Consider the canonical epimorphism (4.2.4):
\[
\gr_J(A)\otimes_A X_0\longrightarrow \grs(X).
\]
Since \(X_0\) is noetherian (resp. artinian) and \(J^n\) is of finite type for every \(n\), the existence of this epimorphism shows that for every integer \(n\), the kernel of the canonical morphism
\[
X_n\longrightarrow X_{n-1}
\]
is noetherian (resp. artinian). It follows at once by induction on \(n\) that the \(X_n\) are noetherian (resp. artinian).
\end{prooftext}

\begin{statement}{Definition 5.1.3}
Let \(X\) be a projective system indexed by \(\mathbb N\), with values in \(\mathcal C\). It is said to be noetherian (resp. artinian) AR-\(J\)-adic if it is AR-\(J\)-adic and if its components are noetherian (resp. artinian) objects.
\end{statement}

We shall denote by AR-\(J\)-adn\((\mathcal C)\) (resp. AR-\(J\)-adf\((\mathcal C)\)) the full subcategory of \(\Hom_{\mathrm{AR}}(\mathbb N^\circ,\mathcal C)\) (2.4) generated by the images of noetherian (resp. artinian) AR-\(J\)-adic projective systems.

It is then clear that the canonical functors
\[
P^n_{\mathrm{AR}}:J\text{-adn}(\mathcal C)\longrightarrow \mathrm{AR}\text{-}J\text{-adn}(\mathcal C),
\]
\[
P^f_{\mathrm{AR}}:J\text{-adf}(\mathcal C)\longrightarrow \mathrm{AR}\text{-}J\text{-adf}(\mathcal C),
\]
obtained by restriction of \(P_{\mathrm{AR}}\) (2.4), are equivalences of categories.

Let us briefly recall the statement and proof of Hilbert's theorem.

\begin{statement}{Theorem 5.1.4}
Hilbert's theorem. Let \(A\) be a commutative ring, \(\mathcal C\) an abelian \(A\)-category, and \(M\) an object of finite type (resp. noetherian) of \(\mathcal C\). For every \(A\)-algebra of finite type \(B\), graded by \(\mathbb N\), the graded \((B,\mathcal C)\)-module
\[
B\otimes_A M
\]
is of finite type (resp. noetherian).

N.B. If \(B=\bigoplus_{n\ge0}B_n\), then \(B\otimes_A M\) denotes the graded object \(\bigl(B_n\otimes_A M\bigr)_{n\in\mathbb N}\).
\end{statement}

\begin{prooftext}
One reduces immediately to the case where \(B\) is the polynomial ring \(A[T]\), endowed with its usual grading. The \((B,\mathcal C)\)-module \(A[T]\otimes_A M=(M_n)_{n\in\mathbb N}\) is then obtained by putting \(M_n=M\) for every \(n\), with \(T\) sending \(M_n\) to \(M_{n+1}\) by the identity of \(M\), for every integer \(n\).

Let \(N\) be a graded sub-\((A[T],\mathcal C)\)-module of \(A[T]\otimes_A M\). The components \((N_n)_{n\in\mathbb N}\) of \(N\) define an increasing sequence
\[
N_0\subset N_1\subset\cdots\subset N_n\subset\cdots
\]
of subobjects of \(M\). Suppose now that \(M\) is noetherian (resp. of finite type), and let
\[
N^1\subset N^2\subset\cdots\subset N^p\subset\cdots
\]
be an increasing sequence of graded sub-\((A[T],\mathcal C)\)-modules of \(A[T]\otimes_A M\), whose union, in the first hypothesis, when \(M\) is of finite type, is \(A[T]\otimes_A M\).

Consider the following double-entry table:
\[
\begin{array}{cccccc}
N_0^1&\subset&N_0^2&\subset\cdots\subset&N_0^p&\subset\cdots\\
\cap&&\cap&&\cap\\
N_1^1&&N_1^2&&N_1^p\\
\vdots&&\vdots&&\vdots\\
\cap&&\cap&&\cap\\
N_n^1&\subset&N_n^2&\subset\cdots\subset&N_n^p&\subset\cdots .
\end{array}
\]
It is clear by hypothesis on \(M\) that there exists a pair \((p_0,n_0)\) such that \(N^{p_0}_{n_0}\) is maximal (resp. equal to \(M\)). If \(p_1\) is an integer \(\ge p_0\) such that the \(N_n^{p_1}\) (\(n\le n_0\)) are the largest elements of the sequences \((N_n^p)_{p\in\mathbb N}\), it is clear that \(N^p=N^{p_1}\) as soon as \(p\ge p_1\), which allows one to conclude in both cases.
\end{prooftext}

\begin{statement}{Corollary 5.1.5}
Suppose moreover that the category \(\mathcal C\) has denumerable direct sums. Then for every object \(M\) of finite type (resp. noetherian) of \(\mathcal C\), and every \(A\)-algebra of finite type \(B\), the \((B,\mathcal C)\)-module \(B\otimes_A M\) is of finite type (resp. noetherian).
\end{statement}

\begin{prooftext}
One reduces to the case where \(B=A[T]\). Let \(N\) be a sub-\((A[T],\mathcal C)\)-module of \(A[T]\otimes_A M\), considered as an object of \(\mathcal C\). We shall denote by \(\gr(N)\) the associated graded of \(N\) for the filtration induced by the natural filtration of the graded object \(A[T]\otimes_A M\): \(\gr(N)\) is a graded sub-\((A[T],\mathcal C)\)-module of the graded \((A[T],\mathcal C)\)-module \(A[T]\otimes_A M\).

Now let
\[
N^0\subset N^1\subset\cdots\subset N^p\subset\cdots
\]
be an increasing sequence of sub-\((A[T],\mathcal C)\)-modules of \(A[T]\otimes_A M\), whose union is \(A[T]\otimes_A M\) in the first hypothesis. In the first hypothesis (\(M\) of finite type), there exists an integer \(i\) such that \(\gr(N^i)=A[T]\otimes_A M\), whence
\[
N^i=A[T]\otimes_A M.
\]
In the second hypothesis (\(M\) noetherian), the sequence of the \(\gr(N^i)\) is stationary, hence the sequence of the \(N^i\) is stationary as well.

In what follows, Hilbert's theorem will be useful to us especially through the following proposition.
\end{prooftext}

\begin{statement}{Proposition 5.1.6}
Let \(A\) be a noetherian commutative ring, \(J\) an ideal of \(A\), and \(\mathcal C\) an abelian \(A\)-category. For every noetherian \(J\)-adic projective system \(X\) (5.1.1), the graded object \(\grs(X)\) is a noetherian \((\gr_J(A),\mathcal C)\)-module.
\end{statement}

\begin{prooftext}
Since \(\gr_J(A)\) is an \(A\)-algebra of finite type, Hilbert's theorem (5.1.4) implies that
\[
\gr_J(A)\otimes_A X_0
\]
is a noetherian graded \((\gr_J(A),\mathcal C)\)-module. The assertion then follows from the existence of the epimorphism (4.2.4):
\[
\gr_J(A)\otimes_A X_0\longrightarrow \grs(X).
\]
\end{prooftext}

\subsubsection*{5.2. Structure of \(J\)-adn\((\mathcal C)\) and AR-\(J\)-adn\((\mathcal C)\)}

In this section, we shall denote by \(\mathcal E_{\mathcal C}\), or by \(\mathcal E\) if there is no possible confusion, the full subcategory of \(\Hom(\mathbb N^\circ,\mathcal C)\) whose objects are noetherian AR-\(J\)-adic projective systems.

\begin{statement}{Proposition 5.2.1}
The category \(\mathcal E_{\mathcal C}\) is stable under kernels and cokernels in \(\Hom(\mathbb N^\circ,\mathcal C)\). In other words, \(\mathcal E_{\mathcal C}\) is an abelian category and the inclusion functor
\[
\mathcal E_{\mathcal C}\longrightarrow \Hom(\mathbb N^\circ,\mathcal C)
\]
is exact.
\end{statement}

\begin{prooftext}
Let \(u:X\to Y\) be a morphism of \(\mathcal E_{\mathcal C}\). We shall successively see that the cokernel, the image, and the kernel of \(u\) are noetherian AR-\(J\)-adic. Of course, the proof that the image of \(u\) is noetherian AR-\(J\)-adic is only a technical intermediate step.

a) \(\Coker(u)\) is noetherian AR-\(J\)-adic. Since \(X\) satisfies (MLAR) and \(Y\) is AR-\(J\)-adic, \(\Coker(u)\) is AR-\(J\)-adic (3.2.4 (i)).

b) \(\Im(u)\) is noetherian AR-\(J\)-adic. From the exact sequence
\[
0\longrightarrow \Im(u)\longrightarrow Y\longrightarrow \Coker(u)\longrightarrow 0,
\]
the assertion will follow from Lemma 5.2.1.1 below.
\end{prooftext}

\begin{statement}{Lemma 5.2.1.1}
Let
\[
0\longrightarrow L\longrightarrow M\longrightarrow N\longrightarrow0
\]
be an exact sequence in the category \(\Hom(\mathbb N^\circ,\mathcal C)\). If \(M\) and \(N\) are noetherian AR-\(J\)-adic, then \(L\) is noetherian AR-\(J\)-adic.
\end{statement}

\begin{prooftext}
It is plainly enough to see that it is AR-\(J\)-adic. One reduces, in the same way as in the proof of (3.2.4), to the case where \(M\) and \(N\) are strict, and then even noetherian \(J\)-adic. In this case, it follows from 3.1.3 (ii) that \(L\) is strict. The projective system \(L\) is strict, adapted to \(J\), and its associated graded, as a graded subobject of \(\grs(M)\), is noetherian: \(\grs(M)\), a quotient of \(\gr_J(A)\otimes_A M_0\), is noetherian (5.1.6). The assertion then follows from the Artin-Rees lemma (4.2.6).
\end{prooftext}

c) \(\Ker(u)\) is noetherian AR-\(J\)-adic. It is enough to apply (5.2.1.1) to the evident exact sequence below:
\[
0\longrightarrow \Ker(u)\longrightarrow X\longrightarrow \Im(u)\longrightarrow0.
\]

\begin{statement}{Proposition 5.2.2}
The objects of the abelian category AR-\(J\)-adn\((\mathcal C)\) are noetherian.
\end{statement}

\begin{prooftext}
One reduces immediately to proving the following. Suppose we have a noetherian \(J\)-adic object \(X\), a sequence of noetherian \(J\)-adic objects \((Y^n)_{n\in\mathbb N}\), and morphisms in \(\Hom(\mathbb N^\circ,\mathcal C)\)
\[
v_n:Y^n\longrightarrow X
\]
such that
\begin{enumerate}[label=\textup{(\roman*)}]
\item the \(v_n\) become monomorphisms in AR-\(J\)-adn\((\mathcal C)\);
\item for every integer \(n\), \(v_n\) factors through \(v_{n+1}\) in AR-\(J\)-adn\((\mathcal C)\),
\end{enumerate}
then the \(v_n\) remain stationary, up to isomorphism, for sufficiently large \(n\), in AR-\(J\)-adn\((\mathcal C)\).

Now, to say that \(v_n\) factors through \(v_{n+1}\) in AR-\(J\)-adn\((\mathcal C)\) means that there exists an integer \(r\) and a morphism
\[
f:Y^n[r]\longrightarrow Y^{n+1}
\]
such that, with \(a\) denoting the canonical morphism
\[
Y^n[r]\longrightarrow Y^n,
\]
the following equality holds:
\[
v_n\circ a=v_{n+1}\circ f .
\]
Since the morphism \(a\) is strict, it follows in particular that the image of \(v_n\) is contained in that of \(v_{n+1}\). Replacing each \(Y^n\) by \(\Im(v_n)\), we are reduced to seeing that every increasing sequence of strict projective subsystems of \(X\) is stationary. But if \((Z^n)_{n\in\mathbb N}\) is such a sequence, then, since \(\grs(X)\) is noetherian, the increasing sequence of the \(\grs(Z^n)\) is stationary, hence so is the sequence of the \(Z^n\).
\end{prooftext}

Summarizing:

\begin{statement}{Theorem 5.2.3}
The categories \(J\)-adn\((\mathcal C)\) and AR-\(J\)-adn\((\mathcal C)\) are abelian and noetherian.
\end{statement}

As for AR-\(J\)-adn\((\mathcal C)\), this is what we have just proved. As for \(J\)-adn\((\mathcal C)\), it is equivalent to it.

\begin{statement}{Proposition 5.2.4}
Stability under extension. Let
\[
0\longrightarrow X\longrightarrow Y\longrightarrow Z\longrightarrow0
\]
be an exact sequence in \(\Hom(\mathbb N^\circ,\mathcal C)\), with \(J^{n+1}Y_n=0\) for every \(n\). Suppose that:
\begin{enumerate}[label=\textup{(\roman*)}]
\item \(Z\) is AR-\(J\)-adic;
\item \(X\) is artinian AR-\(J\)-adic.
\end{enumerate}
Then \(Y\) is AR-\(J\)-adic.
\end{statement}

\begin{prooftext}
As in the proof of (3.2.4), one reduces to the case where \(Y\) and \(Z\) are strict and \(Z\) is \(J\)-adic. Then \(X\) is strict by (3.1.3 (ii)). There then exists an integer \(r\) such that the projective system \((X_{n+r}/J^{n+1}X_{n+r})_{n\in\mathbb N}\) is artinian \(J\)-adic. Since \(X\) is a quotient of it and is AR-isomorphic to it, we are reduced to proving the following lemma.
\end{prooftext}

\begin{statement}{Lemma 5.2.4.1}
Let
\[
0\longrightarrow T\longrightarrow X\longrightarrow Y\longrightarrow Z\longrightarrow0
\]
be an exact sequence in \(\Hom(\mathbb N^\circ,\mathcal C)\), with \(J^{n+1}Y_n=0\) for every \(n\). If \(X\) is artinian \(J\)-adic, \(Z\) is \(J\)-adic, and \(T\) is AR-zero, then \(Y\) is AR-\(J\)-adic.
\end{statement}

\begin{prooftext}
To prove the lemma, we first show that, for every integer \(n\), the evident projective system
\[
\bigl(Y_{n+k}/J^{n+1}Y_{n+k}\bigr)_{k\in\mathbb N}
\]
is essentially constant. Let \(S_n^k\) denote the kernel of the evident morphism
\[
X_{n+k}/J^{n+1}X_{n+k}\longrightarrow
Y_{n+k}/J^{n+1}Y_{n+k}.
\]
Consider the following evident exact commutative diagram:
\[
\begin{tikzcd}[column sep=1.35em,row sep=large]
0 \arrow[r] & S_n^k \arrow[r] \arrow[d] & X_{n+k}/J^{n+1}X_{n+k} \arrow[r] \arrow[d,"\sim"'] & Y_{n+k}/J^{n+1}Y_{n+k} \arrow[r] \arrow[d] & Z_{n+k}/J^{n+1}Z_{n+k} \arrow[r] \arrow[d,"\sim"'] & 0\\
0 \arrow[r] & T_n \arrow[r] & X_n \arrow[r] & Y_n \arrow[r] & Z_n \arrow[r] & 0 .
\end{tikzcd}
\]
With \(n\) fixed, the objects \(S_n^k\) are identified with subobjects of \(T_n\) and form a decreasing sequence. Since \(T_n\) is artinian, there exists an integer \(k_0\) such that \(S_n^k=S_n^{k_0}\) for every integer \(k\ge k_0\); it follows at once that the canonical morphisms
\[
Y_{n+k}/J^{n+1}Y_{n+k}
\longrightarrow
Y_{n+k_0}/J^{n+1}Y_{n+k_0}
\]
are isomorphisms.

This being said, denote by \(Y'\) the projective system defined by
\[
Y'_n=\varprojlim_r Y_{n+r}/J^{n+1}Y_{n+r}
\]
on objects, and in the evident way on morphisms. The projective system \(Y'\) is \(J\)-adic. Moreover, by passage to the limit one obtains an exact commutative diagram of the following type:
\[
\begin{array}{ccccc}
0\to S\to X\to Y'\to Z\to0\\
\downarrow\quad\downarrow\id\quad\downarrow\quad\downarrow\id\\
0\to T\to X\to Y\to Z\to0 .
\end{array}
\]
Since \(T\) is AR-zero, the snake lemma applied to this diagram shows that the evident morphism \(Y'\to Y\) is an AR-isomorphism, whence \(Y\) is AR-\(J\)-adic.
\end{prooftext}

\begin{statement}{Corollary 5.2.5}
Let
\[
0\longrightarrow X\longrightarrow Y\longrightarrow Z\longrightarrow0
\]
be an exact sequence in \(\Hom(\mathbb N^\circ,\mathcal C)\), with \(J^{n+1}Y_n=0\) for every \(n\). Suppose that:
\begin{enumerate}[label=\textup{(\roman*)}]
\item \(X\) is artinian AR-\(J\)-adic (resp. artinian and noetherian);
\item \(Z\) is artinian AR-\(J\)-adic (resp. noetherian).
\end{enumerate}
Then \(Y\) is artinian AR-\(J\)-adic (resp. noetherian).
\end{statement}

\subsubsection*{5.3. Noetherian AR-\(J\)-adic projective systems and \(\delta\)-functors}

Let \(\mathcal C\) and \(\mathcal D\) be two abelian \(A\)-categories, and let \(F:\mathcal C\to\mathcal D\) be an additive functor. The functor \(F\) extends in the clear way to a functor, again denoted \(F\),
\[
F:\Hom(\mathbb N^\circ,\mathcal C)\longrightarrow\Hom(\mathbb N^\circ,\mathcal D),
\]
which is additive and “commutes with translations”.

\begin{statement}{Proposition 5.3.1}
Let \(\mathcal C\) and \(\mathcal D\) be two abelian \(A\)-categories and
\[
T=(T^i)_{i\in\mathbb N}
\]
an exact \(A\)-linear \(\delta\)-functor from \(\mathcal C\) to \(\mathcal D\). Suppose that:
\begin{enumerate}[label=\textup{(\roman*)}]
\item For every noetherian object \(M\) of \(\mathcal C\), annihilated by a power of \(J\), and every integer \(n\), \(T^n(M)\) is noetherian.
\item For every graded \(A\)-algebra \(B\) of finite type, annihilated by \(J\), and every noetherian graded \((B,\mathcal C)\)-module \(N\), the graded \((B,\mathcal D)\)-modules \(T^n(N)\), \(n\in\mathbb N\), are noetherian.
\end{enumerate}
Then, for every noetherian AR-\(J\)-adic object \(X\) of \(\Hom(\mathbb N^\circ,\mathcal C)\), the \(T^n(X)\) are noetherian AR-\(J\)-adic.
\end{statement}

\begin{prooftext}
Since the functors \(T^n\) commute with translations, it is enough to prove the assertion when \(X\) is noetherian \(J\)-adic. In that case, since \(\grs(X)\) is a noetherian graded \(\gr_J(A)\)-module (5.1.6), it follows from hypothesis (ii) and from a variant of a corollary of Shih's theorem (cf. Appendix A.3) that for every integer \(n\) the following assertions hold:
\[
\text{(a) }\quad T^n(X)\text{ satisfies (MLAR),}
\]
\[
\text{(b) }\quad \grs\,T^n(X)\text{ is a noetherian graded }(\gr_J(A),\mathcal D)\text{-module.}
\]
The Artin-Rees lemma (4.2.6) allows us to conclude.
\end{prooftext}

\begin{statement}{Corollary 5.3.2}
The \(\delta\)-functor
\[
T=(T^i)_{i\in\mathbb N}:\Hom(\mathbb N^\circ,\mathcal C)\longrightarrow\Hom(\mathbb N^\circ,\mathcal D)
\]
induces an exact \(A\)-linear \(\delta\)-functor
\[
\mathcal E_{\mathcal C}\longrightarrow\mathcal E_{\mathcal D},
\]
commuting with translations, and hence, by passage to the quotient, a new exact \(A\)-linear \(\delta\)-functor, again denoted \((T^i)_{i\in\mathbb N}\),
\[
\mathrm{AR}\text{-}J\text{-adn}(\mathcal C)\longrightarrow
\mathrm{AR}\text{-}J\text{-adn}(\mathcal D).
\]
\end{statement}

\section*{Appendix: Shih's theorem}

The purpose of the present appendix is to give a proof of Shih's theorem in the form used in the proof of (5.3.1). This form being more general than the one appearing in EGA \(0_{\mathrm{III}}\), a direct proof, adapted from that of loc. cit., has proved necessary.

\paragraph{Terminology.}
Let \(r_0\) be an integer \(\ge1\). The datum, for every integer \(r\ge r_0\), of a bigraded object
\[
E_r=\bigl(E_r^{pq}\bigr)_{p\in\mathbb Z,\ q\in\mathbb Z},
\]
endowed with a differential \(d_r\) of bidegree \((r,1-r)\), and, for every \(r\), of an isomorphism
\[
H(E_r)\xrightarrow{\sim}E_{r+1},
\]
will be called below a spectral filter. The notion of a morphism of spectral filters is defined in the evident way.

\subsubsection*{A.1. Recollections on the spectral sequences associated with \(\delta\)-functors}

Let \(\mathcal C\) be an abelian category and \(X\) an object of \(\mathcal C\), endowed with a decreasing filtration \((F^pX)_{p\in\mathbb Z}\). For every exact \(\delta\)-functor \(T\) from \(\mathcal C\) to an abelian category \(\mathcal D\), the method of “exact couples” (cf., e.g., Maclane, \emph{Homology}, p.~336) makes it possible to describe a spectral sequence
\[
E(X):\quad E_1^{pq}=T^{p+q}(F^pX/F^{p+1}X)\Rightarrow E^n=T^n(X),
\]
for which the filtration of \(E^n\) is defined as follows:
\[
F^p(E^n)=\Im\bigl(T^n(F^pX)\longrightarrow T^n(X)\bigr)
=\Ker\bigl(T^n(X)\longrightarrow T^n(X/F^pX)\bigr).
\]
Moreover, if the filtration of \(X\) is finite, the spectral sequence so obtained is weakly convergent (EGA \(0_{\mathrm{III}}\) 11.1.3).

\paragraph{Functoriality.}
If \(f:X\to Y\) is a morphism of filtered objects of \(\mathcal C\), one constructs functorially a morphism of spectral sequences
\[
E(f):E(X)\longrightarrow E(Y),
\]
coinciding with \(T^{p+q}(\gr^p f)\) (resp. \(T^n(f)\)) on \(E_1^{pq}(X)\) (resp. \(E^n(X)\)). In particular, if \(u:X\to X\) is an endomorphism of \(X\) such that
\[
u(F^pX)\subset F^{p+h}(X)
\]
for an integer \(h\) independent of \(p\), then \(E(u)\) is an endomorphism of bidegree \((h,-h)\) of the spectral filter associated with \(E(X)\), commuting with the differentials.

\subsubsection*{A.2. The case of projective systems}
Let now \(X=(X_n)_{n\in\mathbb Z}\) be a projective system of objects of an abelian category \(\mathcal C\), such that
\[
X_i=0\qquad(i<0).
\]
If \(T\) denotes an exact \(\delta\)-functor from \(\mathcal C\) to an abelian category \(\mathcal D\), \(T\) extends naturally to an exact \(\delta\)-functor, again denoted \(T\),
\[
\Hom(\mathbb Z^\circ,\mathcal C)\longrightarrow \Hom(\mathbb Z^\circ,\mathcal D).
\]
Thus, endowing \(X\) with its natural filtration as a projective system (EGA \(0_{\mathrm{III}}\) 13.4.1), one obtains a spectral sequence
\[
E(X)=E(X_n)_{n\in\mathbb Z},
\]
where the \(E(X_n)\) denote the spectral sequences associated with \(T\) and the \(X_n\) endowed with the induced filtration.

\paragraph{Properties of \(E(X)\).}
a) For every integer \(n\), the filtration of the abutment \(E^n(X)\) in the spectral sequence \(E(X)\) is identical with its natural filtration as a projective system when \(X\) is strict.

Indeed, denoting by the letter \(F\) (as in A.1) the first filtration and by the letter \(\Phi\) the second, one has for every integer \(\alpha\):
\[
F^pE^n(X)_\alpha
=\Ker\bigl(T^n(X_\alpha)\longrightarrow T^n(X_\alpha/F^pX_\alpha)\bigr)
\]
\[
=
\begin{cases}
0,& \alpha\le p-1,\\
\Ker\bigl(T^n(X_\alpha)\longrightarrow T^n(X_{p-1})\bigr),& \alpha\ge p-1.
\end{cases}
\]
In both cases one recognizes the expression of \(\Phi^pE^n(X)\).

This shows in particular that, for \(\beta\ge\alpha\ge p\), the transition morphism
\[
E_\infty^{pq}(X_\beta)\longrightarrow E_\infty^{pq}(X_\alpha)
\]
which is identified with
\[
\gr_T^pT^{p+q}(X_\beta)\longrightarrow
\gr_T^pT^{p+q}(X_\alpha),
\]
is a monomorphism, by EGA \(0_{\mathrm{III}}\) 13.4.1.3.

b) Fixing \(\alpha\in\mathbb Z\), it is clear (cf. the adjoining diagram) that, since \(F^i(X_\alpha)=0\) for \(i\ge\alpha+1\), one has:
\[
\begin{tikzpicture}[baseline=(current bounding box.center),x=1.0cm,y=1.0cm,>=latex]
\draw[->] (-2.0,0) -- (2.2,0);
\draw[->] (0,-1.0) -- (0,1.6);
\node[below left] at (0,0) {$O$};
\node[below] at (0,0) {$\alpha$};
\draw[->] (0,0) -- (0,1.1);
\draw[->] (0,1.1) -- (1.35,0);
\draw[->] (0,0) -- (1.35,0);
\node[above] at (.22,1.1) {$(p,q)$};
\node[left] at (0,0.55) {$1-r$};
\node[above] at (.72,0.08) {$r$};
\end{tikzpicture}
\]
\[
B_r^{pq}(X_\alpha)=B_{r+1}^{pq}(X_\alpha)\quad\text{for }r>p,
\]
\[
Z_r^{pq}(X_\alpha)=Z_{r+1}^{pq}(X_\alpha)\quad\text{for }r>\alpha-p.
\]

c) When \(X\) satisfies the Mittag-Leffler condition, the various projective systems \(\gr^p(X)\) are essentially constant (EGA \(0_{\mathrm{III}}\) 13.4.3). It follows at once that, for \(r<+\infty\), the projective systems \(E_r^{pq}(X)\) are essentially constant. We shall denote by
\[
\mathcal E_r^{pq}(X)
\]
their projective limits. Since the differential operators pass to the limit, one thus obtains a spectral filter \(\mathcal E_r^{pq}(X)\), \(p,q\in\mathbb Z,\ r\ge1\). The objects
\[
\mathcal Z_r^{pq}(X),\quad \mathcal B_r^{pq}(X)\quad(r<+\infty)
\]
associated with this spectral filter are precisely the limits of the essentially constant projective systems \(Z_r^{pq}(X)\) and \(B_r^{pq}(X)\).

It also follows from part b) that
\[
B_\infty^{pq}(X)=B_{p+1}^{pq}(X),
\]
which allows one to put
\[
\mathcal B_\infty^{pq}(X)=\mathcal B_{p+1}^{pq}(X).
\]
Still supposing that \(X\) satisfies (ML), since, for every \(\alpha\),
\[
B_r^{pq}(X_\alpha)=B_{r+1}^{pq}(X_\alpha)\qquad(r>p),
\]
one has a sequence of canonical monomorphisms
\begin{equation}
\cdots\subset\mathcal E_r^{pq}(X)\subset\cdots\subset\mathcal E_{p+2}^{pq}(X)\subset\mathcal E_{p+1}^{pq}(X).
\tag{S}
\end{equation}
Suppose from now on that \(X\) is strict. The following expressions for \(Z_r^{pq}(X_n)\) and \(B_r^{pq}(X_n)\):
\[
Z_r^{pq}(X_n)=
\Im\!\left[
T^{p+q-1}\bigl(F^{p-r+1}X_n/F^pX_n\bigr)
\xrightarrow{\partial}
T^{p+q}\bigl(F^pX_n/F^{p+1}X_n\bigr)
\right],
\]
\[
B_r^{pq}(X_n)=
\Im\!\left[
T^{p+q}\bigl(F^pX_n/F^{p+r}X_n\bigr)
\longrightarrow
T^{p+q}\bigl(F^pX_n/F^{p+1}X_n\bigr)
\right],
\]
show that
\begin{equation}
\left\{
\begin{array}{ll}
Z_r^{pq}(X_m)\xrightarrow{\sim}Z_r^{pq}(X_n)&\text{for }m\ge n\ge p,\\[3pt]
B_r^{pq}(X_m)\xrightarrow{\sim}B_r^{pq}(X_n)&\text{for }m\ge n\ge p+r-1,
\end{array}
\right.
\tag{\(d_1\)}
\end{equation}
where the canonical morphisms are isomorphisms. In particular, the canonical morphism
\[
\mathcal E_r^{pq}(X)\longrightarrow E_r^{pq}(X_{p+r-1})
\]
is an isomorphism. Moreover, it follows from b) that
\begin{equation}
\left\{
\begin{array}{ll}
Z_\infty^{pq}(X_{p+r-1})=Z_r^{pq}(X_{p+r-1})&(r\ge1),\\[3pt]
B_\infty^{pq}(X_{p+r-1})=B_r^{pq}(X_{p+r-1})&(r>p).
\end{array}
\right.
\tag{\(d_2\)}
\end{equation}
This makes it possible to define, for every \(r\), a canonical epimorphism
\[
\theta_r^{pq}:\mathcal E_r^{pq}(X)\longrightarrow E_\infty^{pq}(X_{p+r-1})
\]
which is an isomorphism when \(r>p\). In addition, if \(a\) is an integer such that
\[
\mathcal B_r^{pq}(X)=\mathcal B_a^{pq}(X)\qquad(r\ge a),
\]
for example \(a=p+1\) suffices, so that
\begin{equation}
\cdots\hookrightarrow\mathcal E_r^{pq}(X)\hookrightarrow\cdots\hookrightarrow
\mathcal E_{a+1}^{pq}(X)\hookrightarrow\mathcal E_a^{pq}(X),
\tag{S bis}
\end{equation}
then, for \(s\ge r\ge a\), the diagrams
\begin{equation}
\begin{array}{ccc}
\mathcal E_s^{pq}(X)&\xrightarrow[\epi]{\theta_s^{pq}}&E_\infty^{pq}(X_{p+s-1})\\
\downarrow&&\downarrow\can\\
\mathcal E_r^{pq}(X)&\xrightarrow[\epi]{\theta_r^{pq}}&E_\infty^{pq}(X_{p+r-1})
\end{array}
\tag{\(d_3\)}
\end{equation}
are commutative.

\subsubsection*{A.3. A variant of Shih's theorem}
Suppose given a ring \(S\) filtered by ideals \(F^i(S)\), \(i\ge0\), acting on \(X\) compatibly with its filtration as a projective system.

It follows from the generalities of A.1 that, for every integer \(i\), \(\gr^i(S)\) acts on the bigraded objects \(E_r(X)\) by morphisms of bidegree \((i,-i)\). In particular, for every integer \(n\),
\[
E_r^{(n)}(X)=\bigl(E_r^{p,q}(X)\bigr)_{p+q=n}
\]
is canonically endowed with the structure of a graded \((\gr^\bullet(S),\mathcal D)\)-module. The same holds by passage to the limit for
\[
\mathcal E_r^{(n)}(X)=\bigl(\mathcal E_r^{p,q}(X)\bigr)_{p+q=n}.
\]
The ring \(S\), acting on \(X\) compatibly with its filtration as a projective system, likewise acts on \(T^n(X)\) for every integer \(n\). Consequently, the strict associated graded \(\grs(X)\) is canonically endowed with the structure of a graded \((\gr^\bullet(S),\mathcal C)\)-module, and similarly, when \(T^n(X)\) satisfies (ML), the strict associated graded \(\grs\,T^n(X)\) is canonically endowed with the structure of a graded \((\gr^\bullet(S),\mathcal D)\)-module.

\begin{statement}{Theorem A.3}
Suppose that, for an integer \(n\), the graded \((\gr^\bullet(S),\mathcal D)\)-modules
\[
\mathcal E_1^{(n)}(X)=T^n(\grs^\bullet X)
\quad\text{and}\quad
\mathcal E_1^{(n-1)}(X)=T^{n-1}(\grs^\bullet X)
\]
are noetherian. Then:
\begin{enumerate}[label=\textup{(\roman*)}]
\item \(T^n(X)\) satisfies (MLAR);
\item \(\grs\,T^n(X)\) is a noetherian \((\gr^\bullet(S),\mathcal D)\)-module.
\end{enumerate}
\end{statement}

\begin{prooftext}
i) For every pair \((m,r)\) of integers,
\[
\mathcal B_r^{(m)}(X)=\bigl(\mathcal B_r^{p,m-p}(X)\bigr)_{p\in\mathbb N}
\]
is a graded sub-\((\gr^\bullet(S),\mathcal D)\)-module of \(\mathcal E_1^{(m)}(X)\). In particular, for every integer \(m\), the \(\mathcal B_r^{(m)}(X)\) form an increasing sequence, indexed by \(r\), of graded sub-\((\gr^\bullet(S),\mathcal D)\)-modules of \(\mathcal E_1^{(m)}(X)\). Using the fact that \(\mathcal E_1^{(n)}(X)\) and \(\mathcal E_1^{(n-1)}(X)\) are noetherian graded \((\gr^\bullet(S),\mathcal D)\)-modules, one deduces that the sequences
\[
\bigl[\mathcal B_r^{(n)}(X)\bigr]_{r\in\mathbb N}
\qquad\text{and}\qquad
\bigl[\mathcal B_r^{(n-1)}(X)\bigr]_{r\in\mathbb N}
\]
are stationary.

Thus there exists an integer \(a\ge1\) such that, for every \((p,q)\), with \(p+q=n\) or \(n-1\),
\begin{equation}
\mathcal B_r^{pq}(X)=\mathcal B_a^{pq}(X)\qquad(r\ge a).
\tag{A.3.1}
\end{equation}
Consequently, for \(r\ge a+1\), the arrows \(d_r\) with source and target \(\mathcal E_r^{pq}(X)\), \(p+q=n\), are zero. In other words, the monomorphisms
\[
\mathcal E_s^{pq}(X)\hookrightarrow\mathcal E_r^{pq}(X)\qquad(s\ge r\ge a+1)
\]
deduced from (A.3.1) with \(p+q=n\) are in fact isomorphisms. For \(s\ge r\ge a+1\), the commutativity of diagram \((d_3)\), with \(p+q=n\),
\[
\begin{array}{ccc}
\mathcal E_s^{pq}(X)&\xrightarrow[\epi]{\theta_s^{pq}}&E_\infty^{pq}(X_{p+s-1})\\
\downarrow\sim&&\downarrow\\
\mathcal E_r^{pq}(X)&\xrightarrow[\epi]{\theta_r^{pq}}&E_\infty^{pq}(X_{p+r-1})
\end{array}
\]
shows that the transition morphism
\[
E_\infty^{pq}(X_{p+s-1})\longrightarrow E_\infty^{pq}(X_{p+r-1}),
\]
which is in any case a monomorphism (cf. a)), is an epimorphism, hence an isomorphism. In other words, the projective system
\[
r\longmapsto\left[\bigoplus_p\gr^p T^n(X_{p+r})\right]_{r\ge1}
\]
is constant from \(r=a\) on, and therefore \(T^n(X)\) satisfies (MLAR), by (4.1.3).

ii) Denoting by
\[
\mathcal E_\infty^{(n)}(X),\quad
\mathcal Z_\infty^{(n)}(X),\quad
\mathcal B_\infty^{(n)}(X)
\]
the “limits” of the various projective systems indexed by \(r\) considered above, we see that
\[
\mathcal E_\infty^{(n)}(X)=
\mathcal Z_\infty^{(n)}(X)/\mathcal B_\infty^{(n)}(X)
\]
is a noetherian \((\gr^\bullet(S),\mathcal D)\)-module, since \(\mathcal Z_\infty^{(n)}(X)\) is contained in \(\mathcal E_1^{(n)}(X)\). Moreover, it follows from A.2 a) that
\[
\mathcal E_\infty^{(n)}(X)=\grs\,T^n(X).
\]
\end{prooftext}

\phantomsection\addcontentsline{toc}{section}{Exposé VI --- $\ell$-adic Cohomology}
\section*{Exposé VI}

\begin{center}
\textbf{\(\ell\)-ADIC COHOMOLOGY}\\[0.5em]
by J.-P. Jouanolou
\end{center}

Throughout this exposé, a prime number \(\ell\) is fixed.

As usual, \(X_{\mathrm{\acute et}}\) denotes the étale site of a prescheme \(X\).

\subsection*{1. Constructible \(\ell\)-adic sheaves}

Throughout this paragraph, fix once and for all a locally noetherian prescheme \(X\).

\subsubsection*{1.1. Definition and generalities}

\begin{statement}{Definition 1.1.1}
A constructible \(\ell\)-adic sheaf on \(X\) is a projective \((\ell\mathbb Z)\)-adic system of constructible abelian sheaves on \(X_{\mathrm{\acute et}}\) (V 3.2.1).
\end{statement}

For a projective \((\ell\mathbb Z)\)-adic system
\[
F=(F_n)_{n\in\mathbb N}
\]
of abelian sheaves on \(X_{\mathrm{\acute et}}\) to be a constructible \(\ell\)-adic sheaf, it is enough that \(F_0\) be constructible. Indeed, the assertion is local and follows, in the noetherian case, from (V 5.1.2).

Let \(\Abc(X)\) denote the category of constructible abelian sheaves on \(X\). We shall denote by
\[
\ell\text{-}\adc(X)
\]
the full subcategory of \(\Hom(\mathbb N^\circ,\Abc(X))\) generated by the constructible \(\ell\)-adic sheaves. With the notation of Exposé V,
\[
\ell\text{-}\adc(X)=(\ell\mathbb Z)\text{-}\ad(\Abc(X)).
\]
Letting \(\Ab(X)\) denote the category of abelian sheaves on \(X_{\mathrm{\acute et}}\), one knows (SGA 4 IX 2.9) that, when \(X\) is noetherian, the constructible abelian sheaves on \(X\) are the noetherian abelian sheaves. In this case, the category \(\ell\text{-}\adc(X)\) is clearly identified with the category denoted
\[
(\ell\mathbb Z)\text{-}\adn(\Ab(X))
\]
in Exposé V.

\begin{statement}{Proposition 1.1.2}
Let \(F=(F_n)_{n\in\mathbb N}\) be an object of the category \(\Hom(\mathbb N^\circ,\Ab(X))\). For \(F\) to be constructible \(\ell\)-adic, it is necessary and sufficient that it be so locally for the étale topology of \(X\).
\end{statement}

\begin{prooftext}[Proof]
It is clear that the property of \(F\) of being \(\ell\)-adic may be read above a covering of the final object of the site. Moreover, it follows from (SGA 4 IX 2.8) that the \(F_n\) are constructible if they are so locally.
\end{prooftext}

\begin{statement}{Proposition 1.1.3}
The category \(\ell\text{-}\adc(X)\) of constructible \(\ell\)-adic sheaves on \(X\) is abelian. It is noetherian when \(X\) is noetherian.
\end{statement}

\begin{prooftext}[Proof]
When \(X\) is noetherian, the two assertions are consequences of (V 5.2.3). In the general case, one must see that every morphism
\[
u:F\longrightarrow G
\]
of constructible \(\ell\)-adic sheaves on \(X\) has a kernel and a cokernel.

As for the cokernel, this follows from (V 3.1.3). Let us prove the existence of the kernel of \(u\). For every noetherian open \(V\) of \(X\), one knows that the restriction of \(u\) to \(V\),
\[
u_V:F_V\longrightarrow G_V,
\]
has a kernel in the category \(\ell\text{-}\adc(V)\), admitting the following simple description. If \(K\) is the kernel of \(u_V\) in the category \(\Hom(\mathbb N^\circ,\Ab(V))\), \(K\) satisfies the Mittag-Leffler-Artin-Rees condition and, denoting by \(K'\) its projective system of universal images, there exists an integer \(r\) such that the projective system
\[
m_r(K')=\bigl(K'_{n+r}/\ell^{\,n+1}K'_{n+r}\bigr)_{n\in\mathbb N}
\]
(with the evident transition morphisms) is constructible \(\ell\)-adic. The evident composite below is a kernel of \(u_V\) in the category \(\ell\text{-}\adc(V)\):
\[
m_r(K')\longrightarrow K'\longrightarrow K.
\]
The preceding construction shows in particular that, if \(W\) is an open contained in \(V\), the restriction to \(W\) of a kernel for \(u_V\) is a kernel for \(u_W\). Let \((V_i)_{i\in I}\) be a covering of \(X\) by noetherian opens. If \(i\) and \(j\) are two elements of \(I\), the restrictions to \(V_i\cap V_j\) of the kernels of \(u_{V_i}\) and \(u_{V_j}\) are canonically isomorphic, which allows them to be glued to construct a kernel of \(u\).
\end{prooftext}

\subsubsection*{1.1.4}
Let \(F=(F_n)_{n\in\mathbb N}\) be a constructible \(\ell\)-adic sheaf. Since the ring \(\mathbb Z_\ell\) acts evidently on each sheaf \(F_n\), it also acts on \(F\). Moreover, if \(F\) and \(G\) are two constructible \(\ell\)-adic sheaves and
\[
u:F\longrightarrow G
\]
is a morphism, \(u\) is compatible with the actions of \(\mathbb Z_\ell\) on \(F\) and \(G\), respectively. This shows that \(\ell\text{-}\adc(X)\) is canonically endowed with the structure of an abelian \(\mathbb Z_\ell\)-category.

This justifies the following terminology.

\begin{statement}{Terminology}
In what follows, constructible \(\ell\)-adic sheaves will be called constructible \(\mathbb Z_\ell\)-sheaves.
\end{statement}

One should take care not to confuse this notion with that of a \((\mathbb Z_\ell)_X\)-module, which will not be used below.

The category \(\ell\text{-}\adc(X)\) will also be denoted
\[
\mathbb Z_\ell\text{-}\fc(X).
\]

\subsubsection*{1.2. Constructible twisted constant \(\mathbb Z_\ell\)-sheaves}

\begin{statement}{Definition 1.2.1}
A constructible \(\mathbb Z_\ell\)-sheaf
\[
F=(F_n)_{n\in\mathbb N}
\]
is said to be twisted constant if and only if the sheaves \(F_n\) are locally constant on \(X_{\mathrm{\acute et}}\).
\end{statement}

The category of constructible twisted constant \(\mathbb Z_\ell\)-sheaves is the full subcategory of \(\mathbb Z_\ell\text{-}\fc(X)\) generated by the constructible twisted constant \(\mathbb Z_\ell\)-sheaves.

\begin{statement}{Example 1.2.2}
The \(\mathbb Z_\ell\)-sheaf \(\mathbb Z_\ell(1)\).
\end{statement}

Assume that the integer \(\ell\) is prime to the residual characteristics of \(X\). For every integer \(n\), raising to the \(\ell\)-th power defines a morphism of locally constant sheaves:
\[
\mu_{\ell^{n+1}}\longrightarrow \mu_{\ell^n}.
\]
This evidently defines a projective system:
\[
\mathbb Z_\ell(1)=(\mu_{\ell^{n+1}})_{n\in\mathbb N}.
\]
Reducing, for each integer \(n\), to the case where \(\mu_{\ell^n}\) and \(\mu_{\ell^{n+1}}\) are constant, one sees that the projective system \(\mathbb Z_\ell(1)\) is \(\ell\)-adic.

When \(X=\Spec(K)\), where \(K\) is a field of finite type over the prime field, the \(\mathbb Z_\ell\)-sheaf \(\mathbb Z_\ell(1)\) is an example of a constructible twisted constant \(\ell\)-adic sheaf (1.2.1) for which there is no finite étale extension \(K'\) of \(K\) such that the inverse images of the \(F_n\) on \(X'=\Spec(K')\) are constant. Indeed, \(K'\) would contain, for every integer \(n\), a primitive \(\ell^n\)-th root of unity, which is not possible for a finite-type extension of the prime field.

\begin{statement}{Reminders 1.2.3}
(SGA 1 V 7). One knows that the category of constructible locally constant abelian torsion sheaves on a prescheme \(T\) is equivalent to the category of étale coverings of \(T\) endowed with the structure of an abelian \(X\)-group.
\end{statement}

Now let \(X\) be a connected locally noetherian prescheme. Let \(R_X\) denote the category of étale coverings of \(X\), and let \(\Ensf\) denote the category of finite sets. If \(a\) is a geometric point of \(X\), let \(F_a\) be the functor
\[
F_a:R_X\longrightarrow \Ensf,
\]
which associates to every étale covering of \(X\) its geometric fibre at \(a\). The fundamental group
\[
\pi_1(X)=\pi_1(X,a)
\]
of \(X\) at \(a\), being by definition the group of automorphisms of \(F_a\), suitably topologized, allows the functor \(F_a\) to be interpreted as a functor
\[
R_X\longrightarrow \Ensf(\pi_1)
\]
from the category \(R_X\) to the category of finite sets endowed with a continuous action of \(\pi_1=\pi_1(X,a)\).

Recall the following proposition.

\begin{statement}{Proposition 1.2.3.1}
The functor
\[
F_a:R_X\longrightarrow \Ensf(\pi_1)
\]
is an equivalence of categories. It commutes with inductive limits and finite projective limits.
\end{statement}

Let \(Lc(X)\) and \(\Abf(\pi_1)\) denote respectively the category of constructible locally constant abelian torsion sheaves and the category of finite abelian groups endowed with a continuous operation of the fundamental group. The functor \(F_a\) induces an exact functor, still denoted \(F_a\),
\[
F_a:Lc(X)\longrightarrow \Abf(\pi_1),
\]
which is an equivalence of categories.

\subsubsection*{1.2.4. Twisted constant \(\mathbb Z_\ell\)-sheaves on a connected prescheme}

The reminders in (1.2.3) make the lemma below immediate.

\begin{statement}{Lemma 1.2.4.1}
The extension \(\widehat F_a\) of \(F_a\) to the category of projective systems indexed by \(\mathbb N\) induces an equivalence from the category of constructible twisted constant \(\ell\)-adic sheaves to the category
\[
(\ell\mathbb Z)\text{-}\ad(\Abf(\pi_1)).
\]
\end{statement}

Now let \(F=(F_n)_{n\in\mathbb N}\) be an object of the category
\[
(\ell\mathbb Z)\text{-}\ad(\Abf(\pi_1)).
\]
Its projective limit
\[
\varprojlim_{\mathbb N}F_n
\]
is naturally endowed with a structure of \(\mathbb Z_\ell\)-module of finite type on which \(\pi_1\) acts continuously for the \(\ell\)-adic topology. Recall that if \(G\) and \(M\) denote respectively a profinite group and a \(\mathbb Z_\ell\)-module of finite type, a continuous operation of \(G\) on \(M\) means a continuous map
\[
G\times M\longrightarrow M,\qquad (g,m)\longmapsto g\cdot m,
\]
such that
\[
\text{(i)}\quad g\cdot(m+n)=g\cdot m+g\cdot n\quad (g\in G,\ m,n\in M),
\]
and
\[
\text{(ii)}\quad (gh)\cdot m=g\cdot(h\cdot m)\quad (g,h\in G,\ m\in M).
\]
It is then clear that, for every \(\ell\)-adic integer \(a\),
\[
g\cdot(am)=a(g\cdot m)\qquad (g\in G,\ m\in M).
\]
Moreover, the group of invertible \(\mathbb Z_\ell\)-linear maps from \(M\) to \(M\), denoted \(\GL(M)\), is identified with the projective limit of the groups \(\GL(M/\ell^{n+1}M)\). Since these groups are finite, the group \(\GL(M)\) is thus naturally endowed with the structure of a profinite group, and in particular with the structure of a compact topological group. One checks easily that to give a continuous operation of the profinite group \(G\) on \(M\) is the same as to give a morphism of topological groups
\[
G\longrightarrow \GL(M),
\]
where the latter group is endowed with the topology just described. This said, one evidently defines a functor
\[
\varprojlim_{\mathbb N}:(\ell\mathbb Z)\text{-}\ad(\Abf(\pi_1))\longrightarrow
\mathbb Z_\ell\text{-}\modn(\pi_1),
\]
where \(\mathbb Z_\ell\text{-}\modn(\pi_1)\) denotes the category of finite-type \(\mathbb Z_\ell\)-modules endowed with a continuous operation of \(\pi_1\).

This allows us to state the following lemma.

\begin{statement}{Lemma 1.2.4.2}
The preceding functor
\[
\varprojlim_{\mathbb N}:(\ell\mathbb Z)\text{-}\ad(\Abf(\pi_1))
\longrightarrow \mathbb Z_\ell\text{-}\modn(\pi_1)
\]
is an equivalence of categories.
\end{statement}

\begin{prooftext}[Proof]
It is clear that it is essentially surjective: a finite-type \(\mathbb Z_\ell\)-module \(E\) on which \(\pi_1\) acts continuously is the image of the projective system
\[
(E/\ell^{n+1}E)_{n\in\mathbb N},
\]
with the evident operations of \(\pi_1\). Now let
\[
F=(F_n)_{n\in\mathbb N},\qquad G=(G_n)_{n\in\mathbb N}
\]
be two objects of the category \((\ell\mathbb Z)\text{-}\ad(\Abf(\pi_1))\), and, for simplicity, write \(\Hom_{\pi_1}(F,G)\) for the set of morphisms from \(F\) to \(G\) in this category. Forgetting the action of \(\pi_1\), one may regard \(F\) and \(G\) as objects of the category of \(\ell\)-adic projective systems of finite abelian groups, and write \(\Hom_0(F,G)\) for the set of morphisms from \(F\) to \(G\) in this category. The group \(\pi_1\) acts on \(\Hom_0(F,G)\) as follows: an element \(s\) of \(\pi_1\) clearly defines two automorphisms, denoted in the same way, of \(F\) and \(G\) respectively in the category of \(\ell\)-adic projective systems of finite abelian groups, and if \(u\) belongs to \(\Hom_0(F,G)\), one puts
\[
s\cdot u=sus^{-1}.
\]
Under these conditions, it is clear that \(\Hom_{\pi_1}(F,G)\) is canonically identified with the set of \(\pi_1\)-invariants in \(\Hom_0(F,G)\), and similarly that \(\Hom_{\pi_1}(F_n,G_n)\) is identified for every \(n\) with
\[
\Hom_{\mathbb Z}(F_n,G_n)^{\pi_1},
\]
the operation of \(\pi_1\) being defined analogously. It follows, for example, from (Serre, Groupes proalgébriques, Prop. 12) that, denoting by \(M\) and \(N\) the projective limits of \(F\) and \(G\) respectively, one has
\[
\Hom_{\mathbb Z_\ell}(M,N)=\varprojlim_n\Hom_{\mathbb Z}(F_n,G_n).
\]
Taking the invariants of both sides under \(\pi_1\) and using the fact that this operation, being associated with a projective limit, commutes with projective limits, one obtains the assertion.
\end{prooftext}

Lemmas 1.2.4.1 and 1.2.4.2 allow one to state the following less precise proposition.

\begin{statement}{Proposition 1.2.5}
The category of constructible twisted constant \(\mathbb Z_\ell\)-sheaves on a connected locally noetherian prescheme \(X\) is equivalent to the category of finite-type \(\mathbb Z_\ell\)-modules on which \(\pi_1(X)\) acts continuously for the \(\ell\)-adic topology.
\end{statement}

\begin{statement}{Remark}
Let \(F\) be a constructible twisted constant \(\mathbb Z_\ell\)-sheaf, and let \(M\) be the corresponding finite-type \(\mathbb Z_\ell\)-module endowed with a continuous operation of \(\pi_1\). For \(F\) to be locally constant, it is clear that it is enough for \(\pi_1\) to act on \(M\) through a finite quotient; this is not always the case (1.2.2). It is even necessary when \(X\) is normal. After decomposing \(X\) into connected components, one may, in order to see this, reduce to the case where \(X\) is integral. I claim that then, if \(K\) denotes its field of rational functions, and if there exists a finite separable extension \(K'\) of \(K\) such that the inverse image \(F_{K'}\) of \(F\) above \(\Spec(K')\) is constant, then the group \(\pi_1(X)\) acts on \(M\) through a finite group. Indeed (SGA 1 V 8.2), \(\pi_1(K)\to\pi_1(X)\) is surjective, and the hypothesis means that the kernel of the composite
\[
\pi_1(K)\longrightarrow \pi_1(X)\longrightarrow \Aut(M)
\]
has finite index. The same is therefore true for the kernel of \(\pi_1(X)\to\Aut(M)\).
\end{statement}

\begin{statement}{Proposition 1.2.6}
Let \(X\) be a locally noetherian prescheme and let \(F\) be a projective system, indexed by \(\mathbb N\), of abelian sheaves on \(X_{\mathrm{\acute et}}\). The following assertions are equivalent:
\begin{enumerate}[label=(\roman*)]
\item \(F\) is a constructible \(\mathbb Z_\ell\)-sheaf;
\item every noetherian open \(U\) of \(X\) is a finite union of locally closed subsets \((Z_i)_{1\le i\le q}\) above which the inverse image of \(F\) is a constructible twisted constant \(\mathbb Z_\ell\)-sheaf;
\item \(X\) admits a covering by opens, each a finite union of locally closed subsets, above which the inverse image of \(F\) is a constructible twisted constant \(\mathbb Z_\ell\)-sheaf.
\end{enumerate}
\end{statement}

\begin{prooftext}[Proof]
It is clear that (ii) implies (iii).

(iii) \(\Longrightarrow\) (i): If \(F\) satisfies (iii), its geometric fibres are \(\ell\)-adic, hence \(F\) is \(\ell\)-adic by (SGA 4 VIII 3.6). Moreover, its components are constructible (SGA 4 IX 2.8).

To see that (i) implies (ii), one immediately reduces to the case where \(X\) is noetherian. If \(F\) is constructible \(\ell\)-adic, its strict graded object \(\operatorname{gr}^s(F)\) is a noetherian \((\operatorname{gr}_\ell\mathbb Z,\Ab(X))\)-module (V 5.1.6). It follows (SGA 4 IX 2.9) that \(X\) is a finite union of locally closed subsets \((Z_i)_{1\le i\le q}\) above which the inverse image of \(\operatorname{gr}^s(F)\) is a locally constant sheaf. Since, for a sheaf with finite fibres, the property of being locally constant is stable under extensions, one deduces that, for every integer \(n\), the inverse images on each \(Z_i\) of the component \(F_n\) of \(F\) are locally constant sheaves.
\end{prooftext}

\subsubsection*{1.3. Operations \(\otimes\) and \(\Hom\) on \(\mathbb Z_\ell\)-sheaves}

\paragraph{1.3.1. Tensor product.}
Let \(F\) and \(G\) be two objects of the category \(\Hom(\mathbb N^\circ,\Ab(X))\). If
\[
F=(F_n)_{n\in\mathbb N}\quad\text{and}\quad G=(G_n)_{n\in\mathbb N},
\]
the tensor product \(F\otimes G\) of \(F\) and \(G\) is defined by the following formula on objects, and in the evident way on arrows:
\[
(F\otimes G)_n=F_n\otimes G_n.
\]
More precisely one clearly defines a right exact additive bifunctor, called the tensor product:
\[
\otimes:\Hom(\mathbb N^\circ,\Ab(X))\times \Hom(\mathbb N^\circ,\Ab(X))
\longrightarrow \Hom(\mathbb N^\circ,\Ab(X)).
\]

\begin{statement}{Proposition 1.3.2}
If \(F\) and \(G\) are two constructible \(\mathbb Z_\ell\)-sheaves, then the tensor product \(F\otimes G\) is a constructible \(\mathbb Z_\ell\)-sheaf.
\end{statement}

\begin{prooftext}[Proof]
The \(\ell\)-adic property is clear. Moreover, it follows immediately from the definition of constructible sheaves (SGA 4 IX 2.3) that a tensor product of two constructible sheaves is constructible.
\end{prooftext}

By restriction one therefore defines a right exact additive bifunctor, again called tensor product:
\[
\otimes:\mathbb Z_\ell\text{-}\fc(X)\times \mathbb Z_\ell\text{-}\fc(X)
\longrightarrow \mathbb Z_\ell\text{-}\fc(X).
\]

\subsubsection*{1.3.3. Sheaf of homomorphisms}

Let \(F=(F_n)_{n\in\mathbb N}\) and \(G=(G_n)_{n\in\mathbb N}\) be two constructible \(\mathbb Z_\ell\)-sheaves. Suppose that \(F\) is locally free, meaning that, for every integer \(n\), the sheaf \(F_n\) is locally free over the constant sheaf
\[
(\mathbb Z/\ell^{\,n+1}\mathbb Z)_X.
\]
For every integer \(n\), the structural exact sequence
\[
G_{n+1}\xrightarrow{\ell^{\,n+1}}G_{n+1}\longrightarrow G_n\longrightarrow 0
\]
gives, since \(F_{n+1}\) is locally free over \((\mathbb Z/\ell^{\,n+1}\mathbb Z)_X\), an exact sequence
\[
\cHom(F_{n+1},G_{n+1})
 \xrightarrow{\ell^{\,n+1}}
\cHom(F_{n+1},G_{n+1})
 \longrightarrow
\cHom(F_{n+1},G_n)
 \longrightarrow 0.
\tag{1}
\]
Since \(\ell^{\,n+1}G_n=0\), the structural exact sequence
\[
F_{n+1}\xrightarrow{\ell^{\,n+1}}F_{n+1}\longrightarrow F_n\longrightarrow 0
\]
clearly defines an isomorphism
\[
\cHom(F_{n+1},G_n)\xleftarrow{\sim}\cHom(F_n,G_n).
\tag{2}
\]
Combining (1) and (2), one obtains, for every integer \(n\), a morphism
\[
\cHom(F_{n+1},G_{n+1})\longrightarrow \cHom(F_n,G_n).
\]
Hence a projective system
\[
\cHom(F,G)=\bigl(\cHom(F_n,G_n)\bigr)_{n\in\mathbb N},
\]
which is clearly a \(\mathbb Z_\ell\)-sheaf.

\subsubsection*{1.3.4. Generalities on invertible \(\ell\)-adic sheaves. Application to the sheaves \(\mathbb Z_\ell(r)\) (\(r\in\mathbb Z\))}

Let \(L\) be a constructible \(\mathbb Z_\ell\)-sheaf. We shall say that it is invertible if, for every integer \(n\), its \(n\)-th component \(L_n\) is an invertible \((\mathbb Z/\ell^{\,n+1}\mathbb Z)_X\)-module, i.e. locally free of rank one. This is equivalent to saying that \(L\) is locally free in the sense of the preceding paragraph and that the geometric fibres of the \(L_n\) are nonzero cyclic abelian groups.

If \(L\) is an invertible \(\mathbb Z_\ell\)-sheaf, we agree to put
\[
L^{-1}=\cHom(L,\mathbb Z_\ell).
\]
It is clear that there is an isomorphism
\[
L\otimes L^{-1}\xrightarrow{\sim}\mathbb Z_\ell,
\]
deduced componentwise from the analogous isomorphism (EGA \(0_I\) 5.4.3.1). Following (EGA \(0_I\) 5.4.4), we agree to put \(L^{\otimes0}=\mathbb Z_\ell\) and, for \(n\ge1\),
\[
L^{\otimes -n}=(L^{-1})^{\otimes n}.
\]
With this notation, there is then a canonical functorial isomorphism
\[
L^{\otimes m}\otimes L^{\otimes n}\xrightarrow{\sim}L^{\otimes(m+n)}
\]
for all rational integers \(m\) and \(n\). Moreover, this isomorphism has the usual associativity properties.

In particular, the graded objects
\[
(L^{\otimes i})_{i\in\mathbb N}\quad\text{and}\quad(L^{\otimes i})_{i\in\mathbb Z}
\]
are \(\mathbb Z_\ell\)-algebras, graded respectively by \(\mathbb N\) and \(\mathbb Z\).

\begin{statement}{Example}
The \(\mathbb Z_\ell\)-sheaf \(\mathbb Z_\ell(1)\) defined in (1.2.2) is invertible, and in what follows we shall agree to put
\[
\mathbb Z_\ell(r)=\bigl(\mathbb Z_\ell(1)\bigr)^{\otimes r}\qquad (r\in\mathbb Z).
\]
\end{statement}

\subsubsection*{1.4. Various categories constructed from the category of constructible \(\mathbb Z_\ell\)-sheaves}

\paragraph{1.4.1. Constructible \(A\)-sheaves (\(A\) a finite \(\mathbb Z_\ell\)-algebra).}
Let \(A\) be a finite \(\mathbb Z_\ell\)-algebra. Define a category denoted
\[
A\text{-}\fc(X),
\]
whose objects are called constructible \(A\)-sheaves, as follows.

\begin{enumerate}[label=(\roman*)]
\item An object of \(A\text{-}\fc(X)\) is a pair \((F,u)\) consisting of a constructible \(\mathbb Z_\ell\)-sheaf \(F\) and a morphism of \(\mathbb Z_\ell\)-algebras
\[
u:A\longrightarrow \End(F).
\]
\item If \((F,u)\) and \((F',u')\) are two constructible \(A\)-sheaves, an arrow \((F,u)\to(F',u')\) is a morphism of \(\mathbb Z_\ell\)-sheaves
\[
f:F\longrightarrow F'
\]
such that, for every element \(a\) of \(A\), the diagram below is commutative:
\[
\begin{tikzcd}
F \arrow[r,"f"] \arrow[d,"u(a)"'] & F' \arrow[d,"u'(a)"]\\
F \arrow[r,"f"'] & F'.
\end{tikzcd}
\]
\end{enumerate}

A constructible twisted constant \(A\)-sheaf is a constructible \(A\)-sheaf whose underlying constructible \(\mathbb Z_\ell\)-sheaf is twisted constant. The category of constructible twisted constant \(A\)-sheaves is the full subcategory of the category of constructible \(A\)-sheaves generated by the constructible twisted constant \(A\)-sheaves.

It follows immediately from (1.2.4) that, if \(X\) is connected, the category of constructible twisted constant \(A\)-sheaves is equivalent to the category of finite-type \(A\)-modules endowed with a continuous \(A\)-linear operation of \(\pi_1(X)\) on the underlying finite-type \(\mathbb Z_\ell\)-module (1.2.4).

\paragraph{1.4.2. Constructible \(\mathbb Q_\ell\)-sheaves.}
It is clear that the full subcategory of \(\mathbb Z_\ell\text{-}\fc(X)\) generated by the constructible \(\mathbb Z_\ell\)-sheaves of torsion, i.e. killed by a power of \(\ell\), is thick.

\begin{statement}{Definition 1.4.3}
The category of constructible \(\mathbb Q_\ell\)-sheaves on \(X\), denoted
\[
\mathbb Q_\ell\text{-}\fc(X),
\]
is the quotient abelian category, in the sense of Gabriel's thesis, of \(\mathbb Z_\ell\text{-}\fc(X)\) by the thick abelian subcategory of \(\mathbb Z_\ell\)-sheaves of torsion. An object of \(\mathbb Q_\ell\text{-}\fc(X)\) is called a constructible \(\mathbb Q_\ell\)-sheaf.
\end{statement}

The category \(\mathbb Q_\ell\text{-}\fc(X)\) admits the following simple description. It has the same set of objects as \(\mathbb Z_\ell\text{-}\fc(X)\), and, if \(F\) and \(G\) are two such objects,
\[
\Hom_{\mathbb Q_\ell}(F,G)
=\mathbb Q_\ell\otimes_{\mathbb Z_\ell}\Hom_{\mathbb Z_\ell}(F,G),
\]
where \(\Hom_{\mathbb Z_\ell}(F,G)\) and \(\Hom_{\mathbb Q_\ell}(F,G)\) denote respectively the morphisms in \(\mathbb Z_\ell\text{-}\fc(X)\) and \(\mathbb Q_\ell\text{-}\fc(X)\). Indeed, if
\[
u:\mathbb Z_\ell\text{-}\fc(X)\longrightarrow D
\]
denotes an exact functor from \(\mathbb Z_\ell\text{-}\fc(X)\) to an abelian category \(D\), then to say that \(u\) makes multiplication by the elements of \(\mathbb Z_\ell\) invertible is the same as to say that \(u\) vanishes on the torsion objects; this immediately shows that the abelian category described above by objects and morphisms solves the universal problem of passage to the quotient.

A constructible \(\mathbb Q_\ell\)-sheaf is said to be twisted constant if it is the image of a constructible twisted constant \(\mathbb Z_\ell\)-sheaf. The category of constructible twisted constant \(\mathbb Q_\ell\)-sheaves is the full subcategory of
\[
\mathbb Q_\ell\text{-}\fc(X)
\]
generated by the constructible twisted constant \(\mathbb Q_\ell\)-sheaves.

It follows from (1.2.4) that, if \(X\) is connected, the category of constructible twisted constant \(\mathbb Q_\ell\)-sheaves is equivalent to the category of finite-dimensional \(\mathbb Q_\ell\)-vector spaces \(V\) endowed with a continuous operation of \(\pi_1(X)\), i.e. a continuous map
\[
u:\pi_1(X)\times V\longrightarrow V
\]
satisfying the properties made explicit in (1.2.4), which in particular implies that \(u(g,v)\) is \(\mathbb Q_\ell\)-linear in \(v\) for every \(g\in\pi_1(X)\).

It is immediate to see that this is equivalent to giving a morphism of topological groups
\[
\pi_1(X)\longrightarrow \GL(V),
\]
the latter being endowed with its natural topology as an \(\ell\)-adic analytic group.

\paragraph{1.4.3. Constructible \(L\)-sheaves (\(L\) a finite \(\mathbb Q_\ell\)-algebra).}
Let \(L\) be a finite \(\mathbb Q_\ell\)-algebra. In a manner analogous to what was done in (1.4.1) starting from constructible \(\mathbb Z_\ell\)-sheaves, one defines, starting from the category of constructible \(\mathbb Q_\ell\)-sheaves, the category of constructible \(L\)-sheaves, and then the category of constructible twisted constant \(L\)-sheaves. If \(X\) is connected, the category of constructible twisted constant \(L\)-sheaves is equivalent to the category of finite-dimensional \(L\)-vector spaces endowed with a continuous operation of \(\pi_1(X)\), i.e. a continuous operation of \(\pi_1(X)\) on the underlying \(\mathbb Q_\ell\)-vector space, commuting with the operation of \(L\).

\subsubsection*{1.4.4}
If \(A\) is a finite commutative \(\mathbb Z_\ell\)-algebra and \(F\), \(G\) are two constructible \(A\)-sheaves, the projective system
\[
F\otimes_A G=(F_n\otimes_{A_X} G_n)_{n\in\mathbb N},
\]
with the evident transition morphisms, is a \(\mathbb Z_\ell\)-sheaf. Moreover \(A\) acts evidently on each of its components, which allows it to be endowed naturally with the structure of an \(A\)-sheaf. Thus one defines a right exact bifunctor, called tensor product,
\[
A\text{-}\fc(X)\times A\text{-}\fc(X)\longrightarrow A\text{-}\fc(X).
\]

Now let \(F\) and \(G\) be two \(A\)-sheaves and suppose that \(F\) is locally free, i.e. that, for every integer \(n\), \(F_n\) is a locally free \((A/\ell^{\,n+1}A)_X\)-module. Then, as in (1.3.2), one sees that the projective system
\[
\cHom_A(F,G)=\bigl(\cHom_{A_n}(F_n,G_n)\bigr)_{n\in\mathbb N},
\]
with the evident transition morphisms, is a \(\mathbb Z_\ell\)-sheaf. Moreover, since \(A\) acts naturally on its components, it is canonically endowed with the structure of an \(A\)-sheaf.

As in (1.3.3), one defines the notion of an invertible \(A\)-sheaf and, if \(L\) is such an \(A\)-sheaf, gives a meaning to the expression \(L^{\otimes r}\) for every natural integer \(r\). Paragraph (1.3.3) adapts without difficulty to this situation.

\subsubsection*{1.4.5}
Let \(F\) be a constructible \(\mathbb Z_\ell\)-sheaf, and let \(p\) be the canonical functor
\[
\mathbb Z_\ell\text{-}\fc(X)\longrightarrow \mathbb Q_\ell\text{-}\fc(X).
\]
The composite of \(p\) with the functor \(G\mapsto F\otimes G\) makes multiplication by the elements of \(\mathbb Z_\ell\) invertible and hence defines a functor
\[
\mathbb Q_\ell\text{-}\fc(X)\longrightarrow \mathbb Q_\ell\text{-}\fc(X),
\]
which is right exact. Applying the same argument to the first variable, one thus defines a right exact bifunctor, called tensor product,
\[
\mathbb Q_\ell\text{-}\fc(X)\times \mathbb Q_\ell\text{-}\fc(X)
\longrightarrow \mathbb Q_\ell\text{-}\fc(X).
\]
If \(X\) is connected and \(F\) and \(G\) are two constructible twisted constant \(\mathbb Q_\ell\)-sheaves corresponding to representations of \(\pi_1(X)\) on vector spaces \(V\) and \(W\), respectively, then the tensor product \(F\otimes G\) is again locally constant and corresponds to the diagonal representation of \(\pi_1(X)\) on \(V\otimes W\).

\subsubsection*{1.4.6}
Let \(F\) be a locally free constructible \(\mathbb Z_\ell\)-sheaf (1.3.3). The composite of \(p\) with the functor \(G\mapsto\cHom(F,G)\) (1.3.2) makes multiplication by elements of \(\mathbb Z_\ell\) invertible and hence defines a functor
\[
\cHom(F,\cdot):\mathbb Q_\ell\text{-}\fc(X)\longrightarrow \mathbb Q_\ell\text{-}\fc(X).
\]
If \(F\) and \(F'\) are two locally free constructible \(\mathbb Z_\ell\)-sheaves defining isomorphic \(\mathbb Q_\ell\)-sheaves, then there exists a morphism of \(\mathbb Z_\ell\)-sheaves
\[
u:F\longrightarrow F'
\]
whose kernel and cokernel are torsion. This immediately shows that the functors \(\cHom(F,\cdot)\) and \(\cHom(F',\cdot)\) from \(\mathbb Q_\ell\text{-}\fc(X)\) to itself are isomorphic.

Now let \(F\) be a constructible twisted constant \(\mathbb Q_\ell\)-sheaf. If we again denote by \(F\) the corresponding \(\mathbb Z_\ell\)-sheaf, it follows from the interpretation of the restrictions of \(F\) to the connected components \(Y\) of \(X\) as continuous \(\pi_1(Y)\)-modules that \(F\) may be written uniquely, up to isomorphism, in the form
\[
F=P\oplus T,
\]
where \(P\) is locally free and \(T\) is torsion. This allows one to define, for every \(\mathbb Q_\ell\)-sheaf \(G\), a \(\mathbb Q_\ell\)-sheaf
\[
\cHom(F,G)=\cHom(P,G).
\]
An isomorphism \(u:F\to F'\) of \(\mathbb Q_\ell\)-sheaves defines, for every \(G\), an isomorphism
\[
\cHom(F',G)\xrightarrow{\sim}\cHom(F,G).
\]
More precisely, one thereby defines a bifunctor
\[
\cHom:(\mathbb Q_\ell\text{-}\fct(X))^\circ\times \mathbb Q_\ell\text{-}\fc(X)
\longrightarrow \mathbb Q_\ell\text{-}\fc(X),
\]
where \(\mathbb Q_\ell\text{-}\fct(X)\) denotes the category of constructible twisted constant \(\mathbb Q_\ell\)-sheaves. This bifunctor has the usual exactness properties.

\subsubsection*{1.4.7}
Let \(L\) now be a finite-type extension of \(\mathbb Q_\ell\), and let \(A\) denote the integral closure of \(\mathbb Z_\ell\) in \(L\), so that \(A\) is a free \(\mathbb Z_\ell\)-algebra of finite type and a Dedekind ring. The category of constructible \(L\)-sheaves is isomorphic to the quotient category of the category of constructible \(A\)-sheaves by the category of constructible torsion \(A\)-sheaves. This allows one, as in (1.4.6), to define from the tensor product of (1.4.4) a tensor product:
\[
L\text{-}\fc(X)\times L\text{-}\fc(X)\longrightarrow L\text{-}\fc(X).
\]
Since \(A\) is a Dedekind ring, the argument used in (1.4.6) shows that every twisted constant \(A\)-sheaf is the direct sum of a locally free \(A\)-sheaf and a torsion \(A\)-sheaf. Denoting by \(L\text{-}\fct(X)\) the category of constructible twisted constant \(L\)-sheaves, one defines, again as in (1.4.6), a bifunctor
\[
\cHom:(L\text{-}\fct(X))^\circ\times L\text{-}\fc(X)
\longrightarrow L\text{-}\fc(X),
\]
with the usual exactness properties.

\subsubsection*{1.5. Constructible \((AR,\mathbb Z_\ell)\)-sheaves}

\begin{statement}{Definition 1.5.1}
Let \(X\) be a locally noetherian prescheme. An object \(F\) of \(\Hom(\mathbb N^\circ,\Ab(X))\) is called a constructible \((AR,\mathbb Z_\ell)\)-sheaf if \(X\) admits a covering
\[
(U_i)_{i\in I}
\]
by noetherian opens such that the restriction of \(F\) above each \(U_i\) is a \((AR,\ell\mathbb Z)\)-noetherian object (V 1.3).
\end{statement}

It follows immediately from (V 5.2.1) that the full subcategory of \(\Hom(\mathbb N^\circ,\Ab(X))\) generated by the constructible \((AR,\mathbb Z_\ell)\)-sheaves is stable under kernels and cokernels. In particular, it is an abelian category.

\subsubsection*{1.5.2. Constructible \((AR,\mathbb Z_\ell)\)-sheaves which are locally \(AR\)-null}

It is clear that an \(AR\)-null object (V 2.2.1) of \(\Hom(\mathbb N^\circ,\Ab(X))\) is a constructible \((AR,\mathbb Z_\ell)\)-sheaf if and only if its components are constructible. It follows in particular from (V 2.2.2) and (SGA 4 IX 2.6) that the category of \(AR\)-null constructible \((AR,\mathbb Z_\ell)\)-sheaves is a thick abelian subcategory of the full category generated by the constructible \((AR,\mathbb Z_\ell)\)-sheaves in \(\Hom(\mathbb N^\circ,\Ab(X))\).

Agreeing to say that an object \(P\) of \(\Hom(\mathbb N^\circ,\Ab(X))\) is locally \(AR\)-null if there exists an open covering of \(X\) for the étale topology (or, equivalently, for the Zariski topology) above each open of which \(P\) is \(AR\)-null, one sees, by reducing through localization to the preceding case, that the category of locally \(AR\)-null constructible \((AR,\mathbb Z_\ell)\)-sheaves is a thick abelian subcategory of that of constructible \((AR,\mathbb Z_\ell)\)-sheaves.

Observe that, if \(X\) is noetherian, the constructible \((AR,\mathbb Z_\ell)\)-sheaves are the noetherian \((AR,\ell\mathbb Z)\)-adic objects of \(\Hom(\mathbb N^\circ,\Ab(X))\), and locally \(AR\)-null is equivalent to \(AR\)-null.

\begin{statement}{Definition 1.5.3}
Let \(X\) be a locally noetherian prescheme. The category of constructible \((AR,\mathbb Z_\ell)\)-sheaves on \(X\), denoted
\[
(AR,\mathbb Z_\ell)\text{-}\fc(X),
\]
is the quotient abelian category of the full subcategory generated by the constructible \((AR,\mathbb Z_\ell)\)-sheaves in
\[
\Hom(\mathbb N^\circ,\Ab(X))
\]
by the full thick subcategory generated by the locally \(AR\)-null constructible \((AR,\mathbb Z_\ell)\)-sheaves.
\end{statement}

\begin{statement}{Notation 1.5.4}
Whenever useful, \(U(X)\) will denote the full subcategory of \(\Hom(\mathbb N^\circ,\Ab(X))\) generated by the constructible \((AR,\mathbb Z_\ell)\)-sheaves.
\end{statement}

Since every constructible \(\mathbb Z_\ell\)-sheaf is in particular a constructible \((AR,\mathbb Z_\ell)\)-sheaf, one defines in the evident way a functor
\[
u:\mathbb Z_\ell\text{-}\fc(X)\longrightarrow (AR,\mathbb Z_\ell)\text{-}\fc(X).
\]

\begin{statement}{Proposition 1.5.5}
The above functor \(u\) is an equivalence of categories.
\end{statement}

\begin{prooftext}[Proof]
To prove that \(u\) is fully faithful and essentially surjective, one reduces ``by gluing'' as in the proof of 1.1.3 to the analogous statement (V 5.1.3).
\end{prooftext}

\subsection*{2. Formalism of \(\ell\)-adic cohomology}

\subsubsection*{2.1. Inverse image}

Let
\[
f:X\longrightarrow Y
\]
be a morphism of locally noetherian preschemes. The inverse-image functor
\[
f^*:\Ab(Y)\longrightarrow \Ab(X)
\]
is exact and therefore extends in the evident way to an exact functor, denoted in the same way,
\[
f^*:\Hom(\mathbb N^\circ,\Ab(Y))\longrightarrow \Hom(\mathbb N^\circ,\Ab(X)).
\]
This latter functor evidently carries constructible \(\mathbb Z_\ell\)-sheaves to constructible \(\mathbb Z_\ell\)-sheaves, and locally \(AR\)-null projective systems to locally \(AR\)-null projective systems. With the notation of (1.5.4), it therefore induces two ``inverse image'' functors
\[
f^*:\mathbb Z_\ell\text{-}\fc(Y)\longrightarrow \mathbb Z_\ell\text{-}\fc(X),
\]
\[
f^*:U(Y)\longrightarrow U(X),
\]
the second of which is exact.

The equivalence of categories (1.5.5), together with the fact that \(f^*\) carries locally \(AR\)-null projective systems to locally \(AR\)-null projective systems, then shows that the inverse-image functor
\[
f^*:\mathbb Z_\ell\text{-}\fc(Y)\longrightarrow \mathbb Z_\ell\text{-}\fc(X)
\]
is also exact.

If \(g:Y\to Z\) is another morphism of locally noetherian preschemes, the functors \(f^*\) and \(g^*\) are evidently related by an isomorphism
\[
(g\circ f)^*\xrightarrow{\sim}f^*\circ g^*,
\]
with the usual associativity property (SGA 1 VI 7 ff.).

\subsubsection*{2.2. Higher direct images}

Let
\[
f:X\longrightarrow Y
\]
be a separated morphism of finite type between locally noetherian preschemes. Recall the following result of étale cohomology (SGA 4 XVII; for the proper case, cf. SGA 4 XIV 1.1).

\begin{statement}{Lemma 2.2.1}
Let \(A\) be a commutative ring annihilated by a nonzero integer \(n\), and let \(M\) be a constructible \((A,\Ab(X))\)-module. For every integer \(i\), the sheaf
\[
R^i_!f(M)
\]
is a constructible \((A,\Ab(Y))\)-module.
\end{statement}

The following lemma follows from this by applying (V 5.3.1) after restriction over the noetherian opens of \(Y\).

\begin{statement}{Lemma 2.2.2}
For every integer \(i\) and every constructible \((AR,\mathbb Z_\ell)\)-sheaf \(F\) on \(X\), the projective system
\[
R^i_!f(F)
\]
is a constructible \((AR,\mathbb Z_\ell)\)-sheaf.
\end{statement}

This said, the exact \(\delta\)-functor
\[
(R^i_!f)_{i\in\mathbb Z}:\Ab(X)\longrightarrow \Ab(Y)
\]
extends in the evident way to an exact \(\delta\)-functor, denoted in the same way,
\[
\Hom(\mathbb N^\circ,\Ab(X))\longrightarrow \Hom(\mathbb N^\circ,\Ab(Y)),
\]
which by restriction gives an exact \(\delta\)-functor
\[
U(X)\longrightarrow U(Y).
\tag{1}
\]
Moreover, the functors \(R^i_!f\) carry locally \(AR\)-null projective systems to locally \(AR\)-null projective systems; hence, after passage to the quotient, the \(\delta\)-functor (1) gives an exact \(\delta\)-functor
\[
(R^i_!f)_{i\in\mathbb Z}:(AR,\mathbb Z_\ell)\text{-}\fc(X)
\longrightarrow (AR,\mathbb Z_\ell)\text{-}\fc(Y).
\]
Using now the equivalence of categories (1.5.5), this latter \(\delta\)-functor is interpreted as an exact \(\delta\)-functor
\[
\mathbb Z_\ell\text{-}\fc(X)\longrightarrow \mathbb Z_\ell\text{-}\fc(Y).
\]

\begin{statement}{Definition 2.2.3}
The functor
\[
R^i_!f:\mathbb Z_\ell\text{-}\fc(X)\longrightarrow \mathbb Z_\ell\text{-}\fc(Y)
\]
is called the \(i\)-th direct image functor with proper supports of \(f\).
\end{statement}

\begin{statement}{Remark}
Assuming the analogue of Lemma (2.2.1) proved for the ordinary higher direct images when \(X\) and \(Y\) are excellent, which is done in characteristic zero (SGA 4 XIX 5.1), one defines in the same way higher direct images in \(\ell\)-adic cohomology in that case. One could also develop the result of the following paragraph in this framework.
\end{statement}

\begin{statement}{Proposition 2.2.4 (Leray spectral sequence)}
Let \(f:X\to Y\) and \(g:Y\to Z\) be two separated morphisms of finite type between locally noetherian preschemes, and let \(F\) be a constructible \(\mathbb Z_\ell\)-sheaf, respectively a constructible \((AR,\mathbb Z_\ell)\)-sheaf. There exists in \(\mathbb Z_\ell\text{-}\fc(Z)\) a biregular spectral sequence (EGA \(0_{\mathrm{III}}\) 11.1.3), functorial in \(F\):
\[
R^i_!g\bigl(R^j_!f(F)\bigr)\Longrightarrow R^{i+j}_!(g\circ f)(F).
\]
\end{statement}

\begin{prooftext}[Proof]
By the equivalence of categories (1.5.5) and the way in which higher direct images are defined for constructible \(\mathbb Z_\ell\)-sheaves, it is enough to show the analogous statement for constructible \((AR,\mathbb Z_\ell)\)-sheaves. For any abelian category \(C\), let us agree to denote by \(\Sp(C)\) the category of spectral sequences beginning at degree \(2\) with values in \(C\). Associating to every object \(F\) of \(\Hom(\mathbb N^\circ,\Ab(X))\) the ordinary Leray spectral sequence, obtained by extension to projective systems,
\[
R^i_! g(R^j_! f(F))\Longrightarrow R^{i+j}_!(g\circ f)(F),
\tag{S}
\]
it is clear that one obtains a functor
\[
\Hom(\mathbb N^\circ,\Ab(X))\longrightarrow \Sp(\Hom(\mathbb N^\circ,\Ab(Z))),
\]
which by restriction gives, by (2.2.1) and because \(U(Z)\) is thick in \(\Hom(\mathbb N^\circ,\Ab(Z))\), a functor
\[
U(X)\longrightarrow \Sp(U(Z)).
\tag{1}
\]
Since the canonical functor
\[
U(Z)\longrightarrow (AR,\mathbb Z_\ell)\text{-}\fc(Z)
\]
is exact, (1) yields an additive functor
\[
E:U(X)\longrightarrow \Sp((AR,\mathbb Z_\ell)\text{-}\fc(Z)).
\tag{2}
\]
For every object \(F\) of \(\Hom(\mathbb N^\circ,\Ab(X))\), the associated spectral sequence \((S)\) is biregular: this follows from the analogous assertion for the components of the projective system \(F\) and from the fact that the spectral sequence under consideration is supported in the first quadrant. This shows that, for every object of \(U(X)\), the corresponding spectral sequence is biregular.

Moreover, if \(u:F\to G\) is a morphism of \(U(X)\) whose kernel and cokernel are locally \(AR\)-null, then, because the higher direct images carry locally \(AR\)-null projective systems to locally \(AR\)-null projective systems, \(E(u)\) is an isomorphism on the terms \(E_2^{pq}\); it is therefore an isomorphism, since the spectral sequences considered are biregular. It follows that (2) defines, by passage to the quotient, a functor
\[
(AR,\mathbb Z_\ell)\text{-}\fc(X)
\longrightarrow \Sp((AR,\mathbb Z_\ell)\text{-}\fc(Z)),
\tag{3}
\]
which associates in a precise way to every constructible \((AR,\mathbb Z_\ell)\)-sheaf \(F\) a biregular spectral sequence
\[
R^i_!g(R^j_!f(F))\Longrightarrow R^{i+j}_!(g\circ f)(F).
\]
\end{prooftext}

\subsubsection*{2.2.3. Other properties}

\paragraph{A) Cohomological dimension.}
Let \(f:X\to Y\) be a separated morphism of finite type between locally noetherian preschemes. Suppose that the relative dimension of \(f\) is at most an integer \(d\). Then, for every constructible \((AR,\mathbb Z_\ell)\)-sheaf \(F\),
\[
R^i_!f(F)=0\qquad\text{for }i>2d.
\]
This follows from the analogous assertion for the components of \(F\) (SGA 4 X 4.1 and 5.2).

\paragraph{B) Proper base change.}
Consider a cartesian square
\[
\begin{tikzcd}
X' \arrow[r,"g'"] \arrow[d,"f'"'] & X \arrow[d,"f"]\\
Y' \arrow[r,"g"'] & Y .
\end{tikzcd}
\]
Let \(f:X\to Y\) and \(g:Y'\to Y\) be morphisms of locally noetherian preschemes, with \(f\) separated of finite type. Let \(X'=X\times_Y Y'\), which is locally noetherian since \(f\) is of finite type, and let
\[
f':X'\to Y',\qquad g':X'\to X
\]
be the canonical morphisms. For every integer \(i\) and every constructible \(\mathbb Z_\ell\)-sheaf \(F\), there exists an isomorphism in \(\mathbb Z_\ell\text{-}\fc(Y')\), functorial in \(F\):
\[
w_i:g^*(R^i_!f(F))\xrightarrow{\sim} R^i_!f'(g'^*F).
\]
Moreover, if
\[
0\longrightarrow F'\longrightarrow F\longrightarrow F''\longrightarrow 0
\]
is an exact sequence of \(\mathbb Z_\ell\text{-}\fc(X)\) (resp. of \((AR,\mathbb Z_\ell)\text{-}\fc(X)\)), then the \((w_i)_{i\in\mathbb N}\) commute with boundary operators; in other words, the diagrams below are commutative:
\[
\begin{tikzcd}[column sep=huge]
g^*(R^i_!f(F'')) \arrow[r,"w_i"] \arrow[d,"g^*(\delta)"']
&
R^i_!f'(g'^*F'') \arrow[d,"\delta"]\\
g^*(R^{i+1}_!f(F')) \arrow[r,"w_{i+1}"']
&
R^{i+1}_!f'(g'^*F').
\end{tikzcd}
\]

\begin{prooftext}[Proof]
It is enough to prove the analogous assertion for constructible \((AR,\mathbb Z_\ell)\)-sheaves. But it is deduced by passage to the quotient from the analogous assertion for projective systems of torsion sheaves, which is an immediate consequence of the proper base-change theorem (cf. SGA 4 XII 5.1 in the proper case and SGA 4 XVII in the general case).
\end{prooftext}

\paragraph{C) Exact sequence relative to an open and its complement.}
Let \(f:X\to Y\) be a compactifiable morphism of locally noetherian preschemes, let \(U\) be an open of \(X\), and let \(Z\) be the closed complement. For every constructible \(\mathbb Z_\ell\)-sheaf \(F\), there exists an exact sequence in \(\mathbb Z_\ell\text{-}\fc(Y)\), functorial in \(F\):
\[
\cdots\longrightarrow R^i_!f_U(F/U)
\longrightarrow R^i_!f(F)
\longrightarrow R^i_!f_Z(F/Z)
\longrightarrow R^{i+1}_!f_U(F/U)
\longrightarrow\cdots .
\]

\begin{prooftext}[Proof]
Analogous to that of B).
\end{prooftext}

\subsubsection*{2.2.5}
The formalism developed in this paragraph for \(\mathbb Z_\ell\)-sheaves extends without difficulty to \(A\)-sheaves, where \(A\) is a finite \(\mathbb Z_\ell\)-algebra, and to \(L\)-sheaves, where \(L\) is a finite extension of \(\mathbb Q_\ell\). For example, in the case of \(\mathbb Q_\ell\)-sheaves, the fundamental point in the proofs is that all the functors considered for \(\mathbb Z_\ell\)-sheaves carry \(\mathbb Z_\ell\)-sheaves of torsion to \(\mathbb Z_\ell\)-sheaves of torsion.

\subsection*{3. \(\ell\)-adic cohomology class associated with a cycle}

\subsubsection*{3.1}
If \(X\) is a proper scheme, or a smooth finite-type scheme, over an algebraically closed field \(k\), one knows that, for every constructible abelian torsion sheaf \(M\) prime to the characteristic of \(k\), the cohomology groups
\[
H^i(X,M)
\]
are abelian groups of finite length. This permits (2.2.1, Remark) the definition, for every \(\mathbb Z_\ell\)-sheaf \(F\), with \(\ell\) prime to the characteristic of \(k\), of a noetherian \((AR,\ell)\)-adic projective system of abelian groups
\[
H^i(X,F)
\]
with the usual properties. In particular, the method used in number 2 makes it possible to define the cup-product, the inverse-image homomorphism for a morphism of schemes over \(k\), and the direct-image homomorphism for a proper morphism of smooth schemes over \(k\) (cf. Exposé IV for the definitions in the case of torsion sheaves).

Moreover, one knows (1.2.4) that the category of \(\mathbb Z_\ell\)-sheaves on an algebraically closed field is equivalent to the category of finite-type \(\mathbb Z_\ell\)-modules, the equivalence being supplied by the projective-limit functor. In practice it will be convenient to identify the finite-type \(\mathbb Z_\ell\)-module \(H^i(X,F)\) with the corresponding \(\mathbb Z_\ell\)-sheaf.

\subsubsection*{3.2}
Let \(\ell\) be a prime number \(\ne\car k\), and let \(X\) be a smooth scheme of finite type over \(k\). If \(Z\) is an algebraic cycle of codimension \(d\) on \(X\), then in Exposé IV\footnote{See (SGA \(4^{1/2}\) Cycle).} one associated to it, for every integer \(n\), a cohomology class
\[
\gamma_X^{(n)}(Z)\in H^{2d}\bigl(X,(\mu_X^{(n)})^{\otimes d}\bigr),
\]
where \(\mu_X^{(n)}\) denotes the locally free \((\mathbb Z/\ell^{\,n+1}\mathbb Z)_X\)-module of the \(\ell^{\,n+1}\)-st roots of unity. Moreover, one saw that, if \(n'\ge n\), \(\gamma_X^{(n)}(Z)\) is the image of \(\gamma_X^{(n')}(Z)\) in
\[
H^{2d}\bigl(X,(\mu_X^{(n)})^{\otimes d}\bigr)
\]
by the morphism
\[
H^{2d}\bigl(X,(\mu_X^{(n')})^{\otimes d}\bigr)
\longrightarrow
H^{2d}\bigl(X,(\mu_X^{(n)})^{\otimes d}\bigr)
\]
deduced from raising to the power \(\ell^{\,n'-n}\):
\[
\mu_X^{(n')}\longrightarrow \mu_X^{(n)}.
\]
Thus one associates to \(Z\) an element, denoted \(\gamma_X(Z)\):
\[
\gamma_X(Z)\in H^{2d}(X,\mathbb Z_\ell(d)),
\]
where \(\mathbb Z_\ell(d)\) is defined in (1.3.4). More generally, if \(Z\) is a cycle on \(X\), one associates to it by linearity an \(\ell\)-adic cohomology class in
\[
\bigoplus_d H^{2d}(X,\mathbb Z_\ell(d)),
\]
which is in fact the projective limit of the cohomology classes associated with \(Z\) in
\[
\bigoplus_d H^{2d}\bigl(X,(\mu_X^{(n)})^{\otimes d}\bigr).
\]

\subsubsection*{3.3}
By passage to the limit, it is evident that the \(\ell\)-adic cohomology classes of the preceding paragraph satisfy the same formal properties as the cohomology classes studied in Exposé IV, namely:

\begin{enumerate}[label=(\roman*)]
\item If \(f:X\to Y\) is a proper morphism of smooth preschemes over \(k\) and \(Z\) is a cycle on \(X\), then
\[
f_*\bigl(\gamma_X(Z)\bigr)=\gamma_Y(f_*(Z)).
\tag{*}
\]

\item If \(f:X\to Y\) is a morphism of smooth preschemes over \(k\) and \(Z\) is a cycle on \(Y\), then
\[
\gamma_X(f^*Z)=f^*\bigl(\gamma_Y(Z)\bigr).
\tag{**}
\]

\item If \(X\) is a smooth prescheme over \(k\) and \(Z_1\) and \(Z_2\) are two cycles on \(X\), then
\[
\gamma_X(Z_1\cdot Z_2)=\gamma_X(Z_1)\cdot \gamma_X(Z_2),
\tag{**}
\]
where the dot on the left-hand side denotes the intersection product of cycles, and the one on the right-hand side the cup-product in cohomology.

\item Two algebraically equivalent cycles on a smooth \(k\)-prescheme have the same \(\ell\)-adic cohomology class.
\end{enumerate}

\begin{flushleft}
\(\quad (*)\) This compatibility is not proved in (SGA \(4^{1/2}\) Cycle), but is not difficult to deduce from the results of loc. cit.

\(\quad (**)\) In the non-excessive case, see (SGA \(4^{1/2}\) Cycle 2.3.8) for a precise statement.
\end{flushleft}


\phantomsection\addcontentsline{toc}{section}{Exposé VII --- Cohomology and Chern Classes}
\section*{Exposé VII}

\begin{center}
\textbf{COHOMOLOGY OF SOME CLASSICAL SCHEMES}\\[0.5em]
\textbf{AND}\\[0.5em]
\textbf{THE COHOMOLOGICAL THEORY OF CHERN CLASSES}\\[1em]
by\\[0.5em]
J.-P. Jouanolou
\end{center}

Throughout the text, an integer $\nu>0$ is fixed and we put
\[
A=\mathbb Z/\nu\mathbb Z;
\]
the residual characteristics of all schemes considered are assumed to be prime to the integer $\nu$. If $X$ is such a scheme, $\mu_X$ denotes the sheaf of $\nu$-th roots of unity and $A^{\bullet}(X)$ denotes the cohomology ring
\[
\bigoplus_i H^{2i}(X,\mu_X^{\otimes i}),
\]
with multiplication given by cup-product.

\subsection*{1. Vector bundles}

Let $S$ be a scheme and let $E$ be a locally free $\mathcal O_S$-module, understood to be of finite type. Put $X=\V(E)$ and $U=\V(E)^*$, the complement of the zero section, and denote by
\[
p:X\to S,
\qquad
q:U\to S
\]
the canonical projections.

\begin{statement}{Proposition 1.1}
Let $L$ be an $A$-module.
\begin{enumerate}[label=(\roman*)]
\item
\begin{enumerate}[label=\alph*)]
\item For every integer $j>0$, one has
\[
R^j p_*(p^*L)=0.
\]
\item The adjunction morphism
\[
L\longrightarrow p_*p^*(L)
\]
is an isomorphism.
\end{enumerate}
\item Suppose that $E$ has constant rank $r$.
\begin{enumerate}[label=\alph*)]
\item
\[
R_!^i p(p^*L)=0\quad\text{if }i\ne 2r.
\]
\item The trace morphism (SGA 4 XVIII)
\[
R_!^{2r}p(p^*L)\longrightarrow L\otimes\mu_S^{\otimes -r}
\]
is an isomorphism.
\end{enumerate}
\end{enumerate}
\end{statement}

\emph{Proof.} Part (i) follows immediately by localization for the Zariski topology from (SGA 4 XV 2.2). Assertion (ii) is essentially contained in (SGA 4 XVIII).

\begin{statement}{Corollary 1.2}
Let $L$ be an $A_S$-module. For every integer $i$, the inverse-image morphism
\[
H^i(S,L)\longrightarrow H^i(X,p^*L)
\]
is an isomorphism.
\end{statement}

\emph{Proof.} By (1.1), the adjunction morphism $L\to Rp_*(p^*L)$ is an isomorphism. Since $H^*(X,p^*L)\simeq H^*(S,Rp_*p^*(L))$, the assertion follows.

\begin{statement}{Proposition 1.3}
Suppose that $E$ has constant rank $r$. For every $A_S$-module $L$, the following properties hold.
\begin{enumerate}[label=(\roman*)]
\item
\begin{enumerate}[label=\alph*)]
\item
\[
R^i q_*(q^*L)=0\quad\text{if }i\ne 0,2r-1.
\]
\item The adjunction morphism
\[
L\longrightarrow q_*q^*(L)
\]
is an isomorphism.
\item If $Y$ denotes the complement of $U$ in $X$, then the morphism
\[
R^{2r-1}q_*(q^*L)\longrightarrow L\otimes\mu_S^{\otimes -r}
\]
obtained as the composite of the boundary morphism
\[
R^{2r-1}q_*(q^*L)\longrightarrow R_Y^{2r}p_*(p^*L)
\tag{1}
\]
and the purity isomorphism (SGA 4 XVIII)
\[
R_Y^{2r}p_*(p^*L)\longrightarrow L\otimes\mu_S^{\otimes -r}
\]
is an isomorphism.
\end{enumerate}
\item
\begin{enumerate}[label=\alph*)]
\item
\[
R_!^i q(q^*L)=0\quad\text{if }i\ne 1,2r.
\]
\item The boundary morphism
\[
L\longrightarrow R_!^1 q(q^*L)
\]
coming from the unlimited exact sequence, where $f:Y\to S$ is the canonical projection,
\[
\cdots\longrightarrow R_!^i q(q^*L)\longrightarrow R_!^i p(p^*L)
\longrightarrow R^i f_*(f^*L)\longrightarrow R_!^{i+1}q(q^*L)\longrightarrow\cdots
\]
is an isomorphism.
\item The trace morphism
\[
R_!^{2r}q(q^*L)\longrightarrow L\otimes\mu_S^{\otimes -r}
\]
is an isomorphism.
\end{enumerate}
\end{enumerate}
\end{statement}

\emph{Proof.} Let $\alpha:U\to X$ and $\beta:Y\to X$ be the canonical immersions. One knows (SGA 4 XVI 3.7) that, putting $F=p^*(L)$,
\[
F\xrightarrow{\sim}\alpha_*\alpha^*(F)
\quad\text{and}\quad
(R^j\alpha_*)\alpha^*(F)=0\quad\text{if }j\ne0,2r-1.
\]
Moreover, the sheaves $(R^j\alpha_*)\alpha^*F$ for $j\ge1$ have support in $Y$, and hence are acyclic for the functor $p_*$. From these remarks one obtains assertions (i) a) and b) by using the Leray spectral sequence
\[
R^i p_*\bigl(R^j\alpha_*\bigr)(q^*L)\Longrightarrow R^n q_*(q^*L).
\]
Assertion (i) c) follows from (1.1 (i)) and from the unlimited exact sequence
\[
\cdots\longrightarrow R_Y^i p_*(p^*L)\longrightarrow R^i p_*(p^*L)\longrightarrow R^i q_*(q^*L)
\longrightarrow R_Y^{i+1}p_*(p^*L)\longrightarrow\cdots.
\]
The sheaf $R_Y^{2r}p_*(p^*L)$ is associated with the presheaf
\[
U\longmapsto H^{2r}_{Y\times_S U}(X\times_S U,p^*L).
\]

We prove (ii). Assertion c) comes from the fact that $q$ is smooth and has connected geometric fibres of dimension $r$. The exact sequence mentioned in b) is obtained by applying the functor $R_!p$ to the canonical exact sequence
\[
0\longrightarrow \alpha_!(q^*L)\longrightarrow p^*L\longrightarrow \beta_!(\beta^*p^*L)\longrightarrow 0.
\]
Using (1.1 (i)), one obtains assertions (ii) a) and b) without difficulty from this unlimited exact sequence.

\begin{statement}{Remark 1.4}
For $S=\operatorname{Spec}(\mathbb C)$, (1.3) says that the complement of the origin in $\mathbb C^r$ has the same cohomology as the real sphere of dimension $2r-1$, a deformation retract of $\mathbb C^r-0$.
\end{statement}

\begin{statement}{Corollary 1.5 (Gysin exact sequence)}
The Leray spectral sequence
\[
H^i(S,R^j q_*(q^*L))\Rightarrow H^{i+j}(U,q^*L)
\]
defines an unlimited exact sequence, called the Gysin exact sequence,
\[
\cdots\longrightarrow H^i(S,L)\longrightarrow H^i(U,q^*L)\longrightarrow
H^{i-2r+1}(S,L\otimes\mu_S^{\otimes -r})\longrightarrow H^{i+1}(S,L)\longrightarrow\cdots.
\]
\end{statement}

\emph{Proof.} Lemma of two lines.

\subsection*{2. Projective schemes}

In what follows, the mode of exposition adopted is based on a suggestion of L. Illusie, and the method followed is close to a method of P. Deligne for proving the degeneration of certain spectral sequences [6].

\subsubsection*{2.1}
Let $X$ and $Y$ be two schemes and let $f:X\to Y$ be a morphism. We shall define an arrow
\[
Rf_*(A_X)\otimes^L Rf_*(A_X)\longrightarrow Rf_*(A_X)
\tag{2.1.1}
\]
in $D(Y)$, having the usual associativity and commutativity properties, and therefore endowing $Rf_*(A_X)$ with a structure of generalized algebra. For this, write
\[
C^*(X,A_X)
\]
for the Godement resolution of $A_X$ (SGA 4 XVII). Proceeding as in (TF 6.6), one defines a morphism of resolutions
\[
C^*(X,A_X)\otimes C^*(X,A_X)\longrightarrow C^*(X,A_X\otimes A_X).
\]
Composing it with the morphism
\[
C^*(X,A_X\otimes A_X)\longrightarrow C^*(X,A_X)
\]
deduced from the multiplication of $A_X$, one finally obtains a morphism
\[
C^*(X,A_X)\otimes C^*(X,A_X)\longrightarrow C^*(X,A_X).
\tag{2.1.2}
\]
Applying the functor $f_*$, one obtains without difficulty a morphism of complexes of $A_Y$-modules
\[
f_*C^*(X,A_X)\otimes f_*C^*(X,A_X)\longrightarrow f_*C^*(X,A_X),
\]
and hence (2.1.1), using flat resolutions of the components of the expression on the left.

\subsubsection*{2.2}
Let $S$ be a scheme and let $E$ be a locally free $\mathcal O_S$-module of rank $r+1$. We denote by
\[
P=P_S(E)\xrightarrow{p}S
\]
the corresponding projective bundle. The canonical $\mathcal O_P$-module $\mathcal O_P(1)$ defines, as is known, by means of the Kummer exact sequence, an element
\[
\xi\in H^2(P,\mu),
\]
and we shall denote by $\eta$ the image of $\xi$ under the canonical morphism
\[
H^2(P,\mu)\longrightarrow H^0(S,R^2p_*(\mu))
\]
coming from the Leray spectral sequence
\[
H^i(S,R^jp_*(\mu))\Rightarrow H^{i+j}(P,\mu).
\]
The class $\xi$ is identified with an arrow $A_P\to\mu_P[2]$ in $D(P)$, which by transposition defines $\mu_P^{\otimes -1}[-2]\to A_P$, and hence by the usual adjunction an arrow
\[
c:\mu_S^{\otimes -1}[-2]\longrightarrow Rp_*(A_P),
\]
coinciding with $\eta$ on cohomology in dimension $2$.

By tensor product, the arrow $c$ defines, by means of (2.2.1), arrows
\[
c_i:\mu_S^{\otimes -i}[-2i]\longrightarrow Rp_*(A_P),
\]
for every integer $i\ge1$, and we agree to denote by $c_0$ the adjunction arrow
\[
A_S\longrightarrow Rp_*(A_P).
\]
Thus one defines an arrow
\[
\gamma=\bigoplus_{i=0}^{r}c_i:
\bigoplus_{i=0}^{r}\mu_S^{\otimes -i}[-2i]\longrightarrow Rp_*(A_P)
\]
of $D(S)$.

\begin{statement}{Theorem 2.2.1}
The arrow $\gamma$ is an isomorphism.
\end{statement}

Once the arrow $\gamma$ has been defined, (2.2.1) is equivalent to the following corollary.

\begin{statement}{Proposition 2.2.2}
Under the preceding hypotheses, one has:
\begin{enumerate}[label=\alph*)]
\item
\[
R^{2i+1}p_*(A_P)=0\quad\text{for every }i\in\mathbb N.
\]
\item When $\bigoplus_i R^{2i}p_*(\mu_P^{\otimes i})$ is endowed with its $A_S$-algebra structure deduced from cup-product, the morphism of $A_S$-algebras
\[
A_S[T]\longrightarrow \bigoplus_i R^{2i}p_*(\mu_P^{\otimes i})
\]
defined by $\eta$ has kernel the ideal $(T^{r+1})$ and hence defines an isomorphism of $A_S$-algebras
\[
\theta:A_S[T]/(T^{r+1})\xrightarrow{\sim}\bigoplus_i R^{2i}p_*(\mu_P^{\otimes i}).
\]
\item The morphism
\[
A_S\longrightarrow R^{2r}p_*(\mu^{\otimes r})
\]
obtained by composing $\theta$ with the restriction to $A_S$ of multiplication by $T^r$ on $A_S[T]/(T^{r+1})$ is the inverse of the trace morphism (SGA 4 XVIII)
\[
R^{2r}p_*(\mu^{\otimes r})\longrightarrow A_S;
\]
in other words,
\[
\operatorname{Tr}(\eta^r)=1.
\]
\end{enumerate}
\end{statement}

\emph{Proof.} Since the data are stable under base change, the proper base-change theorem (SGA 4 XII 5.1) and compatibility of the trace morphism with base change (SGA 4 XVIII) allow one to reduce to checking the analogous assertion on the fibres of $P$. In other words, one may suppose that $S$ is the spectrum of a separably closed field $k$ and that $p$ is the canonical morphism $\mathbb P^r_k\to k$. In this case, the proof uses the following lemma.

\begin{statement}{Lemma 2.2.3}
Let $k$ be a separably closed field and let $p:\mathbb P_k^r\to k$ be the canonical projection. Then:
\begin{enumerate}[label=\alph*)]
\item
\[
R^{2i+1}p_*(A_P)=0\quad\text{for every }i.
\]
\item
\[
R^{2i}p_*(A_P)\simeq A_k\quad\text{for }i\in[0,r],
\qquad
R^{2i}p_*(A_P)=0\quad\text{if }i>r.
\]
\end{enumerate}
\end{statement}

The lemma, manifestly true for $r=0$, is proved by induction on $r$. Suppose therefore that it is true for $r-1$ and prove it for $r$. The projective space $\mathbb P_k^{r-1}$ embeds as a closed subspace in $\mathbb P_k^r$ with complement the open $\mathbb E_k^r$, whence an exact sequence of cohomology with proper supports
\[
\cdots\longrightarrow H_!^i(\mathbb E_k^r)\longrightarrow H^i(\mathbb P_k^r)
\longrightarrow H^i(\mathbb P_k^{r-1})\longrightarrow H_!^{i+1}(\mathbb E_k^r)\longrightarrow\cdots.
\]
The computation of cohomology with proper supports of $\mathbb E_k^r$ (1.1) shows that, if $i\ne(2r,2r-1)$, the inverse-image map
\[
H^i(\mathbb P_k^r)\longrightarrow H^i(\mathbb P_k^{r-1})
\]
is an isomorphism, whence the assertion in this case. If $i=2r-1$, one similarly sees that $H^{2r-1}(\mathbb P_k^r)$ is contained in $H^{2r-1}(\mathbb P_k^{r-1})$, hence is zero by the induction hypothesis. Finally, if $i=2r$, the relations
\[
H^{2r-1}(\mathbb P_k^{r-1})=H^{2r}(\mathbb P_k^{r-1})=0
\]
show that
\[
H_!^{2r}(\mathbb E_k^r)\xrightarrow{\sim}H^{2r}(\mathbb P_k^r),
\]
whence the assertion by (1.1).

The lemma implies part a) of 2.2.1. We now prove that $\theta$ is an isomorphism. The morphism $\theta$ may be interpreted as a morphism of rings
\[
A[T]/(T^{r+1})\longrightarrow\bigoplus_i H^{2i}(P,\mu^{\otimes i}).
\]


Thus, by (2.2.3), the assertion will be proved if one shows that, for every $i\in[0,r]$, $\xi^i$ is a free element of the $A$-module $H^{2i}(P,\mu^{\otimes i})$. It is enough to check this for $\xi^r$, which will be an immediate consequence of 2.2.2 c).

We now prove part c). In practice, this is the definition of the trace morphism in this situation; in any case, it expresses the compatibility of the trace morphism and the fundamental class in the diagram
\[
\begin{tikzcd}[column sep=large,row sep=large]
x_0 \arrow[r,hook,"\mathrm{inj}"] \arrow[dr] & \mathbb P^r_k \arrow[d,"p"]\\
& k ,
\end{tikzcd}
\]
where $\mathrm{inj}:x_0\to \mathbb P^r_k$ is the injection of a closed point of $\mathbb P^r_k$.

\begin{statement}{Corollary 2.2.4}
Let $L$ be a complex of $A_S$-modules. The arrow
\[
\mathrm R\Gamma(S,\gamma\lten_A\mathrm{id}(L)):
\bigoplus_{i=0}^r \mathrm R\Gamma(S,L\otimes\mu_S^{\otimes -i})[-2i]
\longrightarrow
\mathrm R\Gamma(P,p^*(L))
\]
is an isomorphism.
\end{statement}

Indeed, it is enough to apply the functor $\mathrm R\Gamma_S$ to the isomorphism $\gamma\lten_A\mathrm{id}_L$.

\subsubsection*{2.2.5}
To state the next corollary, for every scheme $X$ we shall put
\[
A^*(X)=\bigoplus_i H^{2i}(X,\mu_X^{\otimes i})
\]
and shall use this notation systematically in what follows.

\begin{statement}{Corollary 2.2.6}
The inverse-image homomorphism
\[
p^*:A^*(S)=\bigoplus_i H^{2i}(S,\mu_S^{\otimes i})
\longrightarrow
\bigoplus_i H^{2i}(P,\mu_P^{\otimes i})=A^*(P)
\]
is an injective morphism of unital rings, and the elements $1,\xi,\ldots,\xi^r$ form a basis of the $A^*(S)$-module $A^*(P)$. In particular, $A^*(P)$ is a free $A^*(S)$-module of the same rank as $E$.
\end{statement}

\section*{3. Chern Classes}

Let $S$ be a scheme and let $E$ be a locally free $\mathcal O_S$-module of constant rank $r$. Denoting by $\check E$ the dual of $E$, put $P=P_S(\check E)$. We have seen (2.2.6) that, if $\xi$ denotes the class of $\mathcal O_P(1)$ in $H^2(P,\mu)$, then the $A^*(S)$-module $A^*(P)$ is free with basis $1,\xi,\ldots,\xi^{r-1}$. In particular, it is immediate for reasons of homogeneity that there exists a unique relation of the form
\[
\xi^r+c_1\xi^{r-1}+\cdots+c_i\xi^{r-i}+\cdots+c_r=0,
\tag{3.1}
\]
with $c_i\in H^{2i}(S,\mu^{\otimes i})$ for every $i\in[0,r]$.

Now let $E$ be any locally free $\mathcal O_S$-module. For every integer $r$, the set $S_r$ of points of $S$ at which $E$ has rank $r$ is open. Thus one obtains a partition
\[
S=\coprod_r S_r
\]
of $S$ into disjoint opens; in particular, for every $A_S$-module $M$ and every integer $i$, the product of the inverse-image morphisms
\[
H^i(S,M)\longrightarrow \prod_r H^i(S_r,M/S_r)
\]
is a bijection.

The preceding considerations justify the following definition.

\begin{statement}{Definition 3.2}
Let $S$ be a scheme and $E$ a locally free $\mathcal O_S$-module. For every integer $i\ge0$, the $i$-th Chern class of $E$, denoted $c_i(E)$, is the element of $H^{2i}(S,\mu^{\otimes i})$ whose restriction, for every integer $r\ge0$, to the open subset $S_r$ of $S$ consisting of the points where $E$ has rank $r$, is the element $c_i$ defined by relation 3.1. If the rank of $E$ is bounded, the total Chern class of $E$ is the element
\[
c(E)=\sum_i c_i(E)
\]
of $A^*(S)$; in the general case, this denotes the corresponding element of
\[
\widehat A^*(S)=\prod_i A^i(S).
\]
\end{statement}

\subsubsection*{3.3}
It is evident that, for all $S$ and $E$, $c_0(E)=1$, and that, if the rank of $E$ is bounded by $d$, then $c_i(E)=0$ for $i>d$.

\begin{statement}{Proposition 3.4}
The Chern classes just defined satisfy the following properties.
\begin{enumerate}[label=(\roman*)]
\item \emph{Normalization.} Let $L$ be an invertible $\mathcal O_S$-module. If $[L]$ denotes the image of the class of $L$ under the boundary morphism
\[
H^1(S,\mathbb G_m)\longrightarrow H^2(S,\mu),
\]
then
\[
c(L)=1+[L].
\]
\item \emph{Stability under base change.} Let $f:S'\to S$ be a morphism of schemes and let $E$ be a locally free $\mathcal O_S$-module. For every integer $i$,
\[
f^*(c_i(E))=c_i(f^*E).
\]
\item \emph{Additivity.} Let
\[
0\to E'\to E\to E''\to0
\]
be an exact sequence of locally free $\mathcal O_S$-modules. For every integer $n$, one has
\[
c_n(E)=\sum_{i+j=n}c_i(E')c_j(E'').
\]
\end{enumerate}
\end{statement}

Before beginning the proof, we extract the following corollary.

\begin{statement}{Corollary 3.5}
For every exact sequence $0\to E'\to E\to E''\to0$ of locally free $\mathcal O_S$-modules, one has
\[
c(E)=c(E')c(E'')
\]
in $\widehat A^*(S)$.
\end{statement}

\emph{Proof of 3.4.} Following A. Grothendieck: \emph{La théorie des classes de Chern}, Bull. S.M.F. (1958), pp. 137--154.

(i) Let $E$ be a locally free $\mathcal O_S$-module; put $P=P_S(\check E)$ and let $p:P\to S$ be the canonical projection. Recall that there is a canonical surjective homomorphism
\[
p^*(\check E)\longrightarrow \mathcal O_P(1).
\]
In particular, if $E$ is invertible and we put $E=L$, in agreement with the notation of (i), then $P=S$ and there is an isomorphism
\[
L\xrightarrow{\sim}\mathcal O_P(1).
\]
Therefore $\xi+[L]=0$, and the assertion follows from the definition of Chern classes.

(ii) Immediate verification, left to the reader.

(iii) By using a partition of $S$ into opens on which the modules in question have constant rank, we may suppose that $E,E',E''$ have constant ranks $r,r',r''$, respectively. Suppose that the following lemma has been proved; we show how it implies the assertion.

\begin{statement}{Lemma 3.6}
Let $E$ be a locally free $\mathcal O_S$-module endowed with a finite filtration
\[
0=E_0\subset E_1\subset\cdots\subset E_i\subset\cdots\subset E_r=E
\]
whose consecutive quotients $L_i=E_i/E_{i-1}$ $(1\le i\le r)$ are invertible. Then
\[
c(E)=(1+[L_1])(1+[L_2])\cdots(1+[L_r]).
\]
\end{statement}

Let $h:Z\to S$ be the fibre product of the flag schemes of $E'$ and $E''$. By the definition of flag schemes, it is clear that $h^*(E')$ and $h^*(E'')$ admit finite filtrations with invertible quotients $(L'_i)$ $(1\le i\le r')$ and $(L''_j)$ $(1\le j\le r'')$, respectively. It follows that $h^*(E)$ admits a filtration with consecutive quotients the $L'_i$ and the $L''_j$. Applying Lemma 3.6 then shows that
\[
c(h^*(E))=
\prod_{1\le i\le r'}c(L'_i)\prod_{1\le j\le r''}c(L''_j)
=c(h^*(E'))c(h^*(E'')),
\]
hence, by (ii),
\[
h^*(c(E))=h^*(c(E')c(E'')).
\tag{3.5.1}
\]
But $h$ is a composite of morphisms of the form
\[
P_T(F)\longrightarrow T
\qquad (F\text{ locally free }\mathcal O_T\text{-module}),
\]
so by (2.2.5) the inverse-image homomorphism $h^*:A^*(S)\to A^*(Z)$ is injective. This gives the desired assertion by (3.5.1).

We prove Lemma 3.6. By the preceding argument, it is enough to prove the analogous relation
\[
c(p^*E)=\prod_i c(p^*(L_i)),
\]
where, recall, $p:P=P_S(\check E)\to S$ is the canonical projection. For every $i\in[1,r]$, the monomorphism $E_i\to E$ defines by duality an epimorphism $\check E\to\check E_i$; putting $Y_i=P_S(\check E_i)$, one obtains a sequence of closed subschemes of $P$:
\[
S\simeq Y_1\subset Y_2\subset\cdots\subset Y_i\subset\cdots\subset Y_r=P.
\]
For every pair $(i,j)$ of integers with $1\le i\le j\le r$, let
\[
\alpha_{ji}:Y_i\to Y_j
\]
be the canonical closed immersion, and put $\alpha_i=\alpha_{r,i}$. For every integer $i\in[1,r-1]$, it follows from the evident lemma (3.6.1) below that $Y_i$ is defined in $Y_{i+1}$ by an invertible ideal isomorphic to
\[
\alpha_{i+1}^*(p^*(\check L_{i+1})\otimes\mathcal O_P(-1)).
\]

\begin{statement}{Lemma 3.6.1}
Let $0\to L\to E\to F\to0$ be an exact sequence of locally free $\mathcal O_S$-modules. Then the sequence
\[
0\to L\otimes_S S(E)\xrightarrow{u}S(E)\xrightarrow{v}S(F)\to0,
\]
where $v$ denotes the evident morphism and $u$ is the degree-one morphism corresponding to the inclusion of $L$ in $E$, is an exact sequence of graded $\mathcal O_S$-modules.
\end{statement}

Moreover, $Y_1$ is identified with $S$ and we have an isomorphism
\[
\alpha_1^*(p^*(\check L_1))\simeq\mathcal O_{Y_1}(1),
\]
i.e.
\[
\alpha_1^*(p^*(\check L_1)\otimes\mathcal O_P(-1))\simeq \mathcal O_{Y_1}.
\tag{3.6.2}
\]
We shall be able to conclude by using the following two properties of the Gysin homomorphism (SGA 4 XVIII) associated with a closed immersion $u:Z\to X$ of smooth $S$-schemes defined by an ideal $J$ of $\mathcal O_X$. It satisfies the projection formula
\[
u_*(x\,u^*(y))=u_*(x)y,
\tag{3.6.3}
\]
and, if the ideal $J$ is invertible, the relation
\[
u_*(1_Z)=-[J]=[\check J].
\tag{3.6.4}
\]
Indeed, from this we shall prove, by induction increasing on the integer $i$, that
\[
(\alpha_i)_*(\alpha_i)^*\Bigl(\prod_{j\le i}[p^*(\check L_j)\otimes\mathcal O_P(-1)]\Bigr)=0.
\tag{3.6.5}
\]
For $i=r$, assertion (3.6.5) says
\[
\prod_i[p^*(\check L_i)\otimes\mathcal O_P(-1)]=0,
\]
or equivalently
\[
\prod_i(\xi+c_1(p^*(\check L_i)))=0,
\]
which immediately gives (3.6). Assertion (3.6.5) for $i=1$ follows from (3.6.2). Suppose it true for $i$ and prove it for $i+1$. For simplicity put
\[
N=\prod_{j\le i}[p^*(\check L_j)\otimes\mathcal O_P(-1)].
\]
The equality (3.6.5) for $i$ can also be written
\[
(\alpha_{i+1})_*(\alpha_{i+1,i})_*(\alpha_{i+1,i})^*(\alpha_{i+1})^*(N)=0,
\]
hence, in view of the projection formula for $\alpha_{i+1,i}$,
\[
(\alpha_{i+1})_*\bigl[(\alpha_{i+1})^*(N)(\alpha_{i+1})^*(p^*(\check L_{i+1})\otimes\mathcal O_P(-1))\bigr]=0,
\]
which is exactly the desired result.

\begin{statement}{Corollary 3.7}
\begin{enumerate}[label=(\roman*)]
\item For every locally free $\mathcal O_S$-module $E$ and every integer $i$,
\[
c_i(\check E)=(-1)^i c_i(E).
\]
\item If $E$ is a locally free $\mathcal O_S$-module of rank $r$, then
\[
c_1(E)=c_1(\Lambda^rE).
\]
\item More generally, the Chern classes of a tensor product or of an exterior power are expressed by means of the universal polynomials made explicit in (SGA 6 V 6.1).
\end{enumerate}
\end{statement}

\emph{Proof.} We prove, for example, (ii); the other assertions are seen similarly. Using the flag scheme of $E$, one reduces to the case where $E$ admits a filtration with invertible quotients $L_i$. The assertion then follows from the well-known isomorphism
\[
\Lambda^r(E)\simeq\bigotimes_i L_i.
\]

\subsubsection*{3.8}
Given a scheme $X$ locally of finite type over $\operatorname{Spec}(\mathbb C)$, the set $X(\mathbb C)$ of its rational points is canonically endowed with the structure of a locally compact topological space. Thus, by algebraic topology, one has a theory of Chern classes [3], assigning to every continuous complex vector bundle $V$ on $X(\mathbb C)$ classes $c_i(V)\in H^{2i}(X(\mathbb C),\mathbb Z)$. Now let $\mathcal O_\mathbb C$ be the sheaf of continuous complex-valued functions on $X(\mathbb C)$, and let $\mathcal E$ be a locally free $\mathcal O_\mathbb C$-module of finite type. One defines Chern classes
\[
c_i(\mathcal E)\in H^{2i}(X(\mathbb C),\mathbb Z),
\]
namely those of the complex vector bundle whose sheaf of continuous sections is $\mathcal E$. Finally, using the evident map $\mathbb Z\to\mathbb Z/\nu\mathbb Z$, one obtains Chern classes, denoted by the same symbol when no confusion is possible,
\[
c_i(\mathcal E)\in H^{2i}(X(\mathbb C),\mathbb Z/\nu\mathbb Z).
\tag{3.8.1}
\]
On the other hand, when $X_{\et}$ is equipped with its structural sheaf $\mathcal O_a$ and $X(\mathbb C)$ with the ring $\mathcal O_\mathbb C$, there is a morphism of ringed topoi (SGA 4 XI)
\[
\varepsilon:X(\mathbb C)\to X_{\et},
\]
the structural morphism of rings
\[
\varepsilon^*(\mathcal O_a)\to\mathcal O_\mathbb C
\tag{3.8.2}
\]
coming from the fact that every algebraic function is continuous. The morphism $\varepsilon$ induces a morphism of abelian sheaves
\[
\varepsilon^*(\mu_{X_{\et}})\to \mu_{X(\mathbb C)},
\tag{3.8.3}
\]
where the letter $\mu$ denotes, in each of the two contexts, the sheaf of $\nu$-th roots of unity. Finally, we shall canonically identify $(\mathbb Z/\nu\mathbb Z)_{X(\mathbb C)}$ and $\mu_{X(\mathbb C)}$ by means of the isomorphism
\[
(\mathbb Z/\nu\mathbb Z)_{X(\mathbb C)}
 \xrightarrow{\sim}\mu_{X(\mathbb C)}
\tag{3.8.4}
\]
sending the section $1$ to the section $e^{2i\pi/\nu}$.

\begin{statement}{Proposition 3.8.5}
Let $E$ be a locally free $\mathcal O_X$-module. For every integer $i\ge0$, let
\[
\varepsilon^*:H^{2i}(X_{\et},\mu_X^{\otimes i})
\to H^{2i}(X(\mathbb C),\mathbb Z/\nu\mathbb Z)
\]
be the inverse-image morphism deduced from 3.8.3 and 3.8.4. Then
\[
c_i(\varepsilon^*E)=\varepsilon^*c_i(E).
\]
\end{statement}

\emph{Proof.} It is enough to check this when $E$ is invertible. Then, with the usual identifications, $\varepsilon^*(E)$ is the image of $E$ under the canonical map
\[
\varepsilon^*:H^1(X_{\et},\mathbb G_m)\to H^1(X(\mathbb C),\mathcal O_\mathbb C^*),
\tag{3.8.6}
\]
and one has to see that the diagram
\[
\begin{tikzcd}[column sep=large,row sep=large]
H^1(X,\mathbb G_m) \arrow[r,"\varepsilon^*"] \arrow[d,"\partial"'] &
H^1(X(\mathbb C),\mathcal O_\mathbb C^*) \arrow[d,"c_1"]\\
H^2(X,\mu_X) \arrow[r,"\varepsilon^*"'] &
H^2(X(\mathbb C),\mathbb Z/\nu\mathbb Z)
\end{tikzcd}
\tag{3.8.7}
\]
where $\partial$ is the boundary morphism coming from the canonical exact sequence
\[
0\to \mu_X\to\mathbb G_m\xrightarrow{\nu}\mathbb G_m\to0
\tag{3.8.8}
\]
and $c_1$ is the first Chern class (3.8.1), is commutative. In fact, using that the morphism "first Chern class"
\[
H^1(X(\mathbb C),\mathcal O_\mathbb C^*)\to H^2(X,\mathbb Z)
\]
is the boundary morphism of the canonical exact sequence
\[
0\to\mathbb Z\to\mathcal O_\mathbb C\xrightarrow{\exp}\mathcal O_\mathbb C^*\to0,
\qquad
t\mapsto 2i\pi t,
\]
one easily sees that the morphism $c_1$ of (3.8.7) is nothing other than the boundary morphism associated with the exact sequence
\[
0\to\mathbb Z/\nu\mathbb Z\to\mathcal O_\mathbb C^*\xrightarrow{\nu}\mathcal O_\mathbb C^*\to0,
\qquad
t\mapsto e^{2i\pi t/\nu},
\]
which, in view of (3.8.4), is the analogue of (3.8.8). The assertion follows at once.

\subsubsection*{3.9}
Let $X$ be a scheme and $E$ a locally free $\mathcal O_X$-module. Immediate passage to the limit allows one to define, for every prime number $\ell$ prime to the residual characteristics of $X$, $\ell$-adic Chern classes
\[
c_i(E)\in H^{2i}(X,\mathbb Z_\ell(i))
\stackrel{\mathrm{def}}{=}\varprojlim_\nu H^{2i}(X,\mu_{\ell^\nu}^{\otimes i}),
\]
which have the usual properties made explicit in 3.4 and 3.7 and which, in the context of 3.8, are manifestly compatible, in an evident sense, with the Chern classes defined topologically.

Tensoring by $\mathbb Q_\ell$, one obtains $\ell$-adic Chern classes modulo torsion
\[
c_i(E)\in H^{2i}(X,\mathbb Q_\ell(i))
=H^{2i}(X,\mathbb Z_\ell(i))\otimes_{\mathbb Z_\ell}\mathbb Q_\ell,
\]
which manifestly satisfy the same properties.

Now suppose that $X$ is a quasi-projective smooth scheme over a field $k$, and denote here by $C^*(X)$ its Chow ring (Sém. Chevalley: Anneaux de Chow et applications III 4). Recall that in IV we defined a morphism of graded rings
\[
\varphi_X:C^*(X)\to A^*(X)
\]
compatible with direct-image and inverse-image morphisms. Moreover, the canonical isomorphism
\[
\varepsilon_X:\Pic(X)\xrightarrow{\sim} C^1(X)
\]
which sends every invertible $\mathcal O_X$-module $L$ to the class of the divisor of a rational section of $L$ not vanishing on any connected component of $X$, gives, on the model of the aforementioned article of Grothendieck, where $k$ was assumed algebraically closed, a theory of Chern classes with values in $C^*(X)$. The composite
\[
\varphi_X\circ\varepsilon_X:\Pic(X)\longrightarrow H^2(X,\mathbb Z_\ell(1))
\]
is just the first $\ell$-adic Chern class. Indeed, it is enough to check this on elements of $\Pic(X)$ of the form $\check J$, with $J$ an invertible ideal of $\mathcal O_X$, and then this follows from (3.6.4). By the uniqueness properties of theories of Chern classes, it follows that the morphism $\varphi_X$ sends Chern classes to Chern classes.

\section*{4. Self-intersection Formula and Applications}

Let $S$ be a scheme, $X$ an $S$-scheme, and $Y$ a closed subscheme of $X$, smooth over $S$ and everywhere of codimension $d$. Finally suppose that $X$ is smooth over $S$ in a neighbourhood of $Y$. Let $i:Y\to X$ be the canonical closed immersion and $J$ the ideal of $\mathcal O_X$ defining $Y$. Recall (SGA $4^{1/2}$ Cycle) that in this case one has Gysin isomorphisms
\[
i_*:H^{2p}(Y,\mu_Y^{\otimes p})
\to H_Y^{2p+2d}(X,\mu_X^{\otimes(p+d)}),
\]
namely cup-products with the fundamental class. If no confusion is possible, we shall denote by the same symbol the Gysin morphisms
\[
i_*:H^{2p}(Y,\mu_Y^{\otimes p})
\to H^{2p+2d}(X,\mu_X^{\otimes(p+d)})
\]
obtained from the preceding ones by forgetting supports.

\begin{statement}{Theorem 4.1}
In the preceding situation, if
\[
N=J/J^2
\]
denotes the normal sheaf of $i$, then for every $y\in A^*(Y)$ one has
\[
i^*i_*(y)=y\,c_d(\check N).
\]
\end{statement}

Before beginning the proof of this theorem, we indicate a few variants, proved in exactly the same way.

a) If $\ell$ is a prime number prime to the residual characteristics of $S$, the analogous formula is true in $H^*(Y,\mathbb Z_\ell(*))$.

b) One may put in supports. If, for example, $T$ is a closed subset of $Y$ and $y\in H_T^{2p}(Y,\mu_Y^{\otimes p})$, formula 4.1 is still true in $H_T^{2p+2d}(Y,\mu_Y^{\otimes(p+d)})$. Of course, $i_*$ then denotes the Gysin morphism
\[
H_T^{2p}(Y,\mu_Y^{\otimes p})
\to H_T^{2p+2d}(X,\mu_X^{\otimes(p+d)}),
\]
which is again an isomorphism (SGA $4^{1/2}$ Cycle).

c) Finally, the two cases may be combined.

Since the fundamental tool in the proof of 4.1 is the theory of blown-up schemes, we now recall a few facts about it. Let
\[
f:X'\to X
\]
be the $X$-scheme obtained by blowing up $Y$, and let
\[
\begin{tikzcd}[column sep=large,row sep=large]
Y' \arrow[r,"j"] \arrow[d,"g"'] & X' \arrow[d,"f"]\\
Y \arrow[r,"i"'] & X
\end{tikzcd}
\]
be the cartesian diagram constructed from $f$ and $i$. Recall (cf. SGA 6 VII 3.1, whose notation we adopt) that $j$ is a regular immersion of codimension $1$ defined by an ideal
\[
J'\simeq\mathcal O_{X'}(1)
\]
and that $Y'\simeq P_Y(N)$. Moreover, one has an exact sequence of $\mathcal O_{Y'}$-modules
\[
0\to F\to g^*(N)\to \mathcal O_{Y'}(1)\to0.
\]

\begin{statement}{Lemma 4.2}
For every $A_X$-module $M$, the homomorphisms
\[
f^*:H_Y^*(X,M)\to H_{Y'}^*(X',f^*(M))
\]
and
\[
f_*:H_{Y'}^*(X',f^*(M))\to H_Y^*(X,M)
\]
are related by $f_*f^*=\mathrm{id}$. Moreover, $f_*$ has degree $0$.
\end{statement}

\emph{Proof.} Since $f$ is a complete-intersection morphism of relative dimension zero (SGA 6 VII 1.8), the last assertion follows from the definition of the Gysin homomorphism (SGA $4^{1/2}$ Cycle). For the first, the projection formula for $f$,
\[
f_*(x\,f^*(x'))=f_*(x)x',
\]
where $x\in H^*(X,A)$ and $x'\in H_Y^*(X,M)$, reduces the claim to proving that $f_*(1)=1$. It is enough to check the analogous assertion over $U=X-Y$, which is evident because $f$ induces an isomorphism from $U'=X'-Y'$ onto $U$.

\begin{statement}{Remark 4.3}
Of course, the analogue of 4.2 obtained by eliminating supports is also true and is proved in the same way.
\end{statement}

We now come to the proof of 4.1. First, by a general result of (SGA $4^{1/2}$ Cycle), which is proved in an essentially formal way, one has
\[
i^*i_*(y)=y\,i^*i_*(1),
\]
so it is enough to prove
\[
i^*i_*(1_Y)=c_d(\check N),
\]
where $i_*$ here denotes the Gysin isomorphism
\[
H^0(Y,A_Y)\to H_Y^{2d}(X,\mu_X^{\otimes d}).
\]
Since $f^*i_*1_Y\in H_{Y'}^{2d}(X',\mu_{X'}^{\otimes d})$, and since the Gysin morphism
\[
j_*:H^{2d-2}(Y',\mu_{Y'}^{\otimes(d-1)})
\to H_{Y'}^{2d}(X',\mu_{X'}^{\otimes d})
\]
is an isomorphism, there exists
\[
u\in H^{2d-2}(Y',\mu_{Y'}^{\otimes(d-1)})
\]
such that
\[
f^*i_*(1_Y)=j_*(u).
\tag{4.1.1}
\]
Applying $j^*$ to both sides of 4.1.1, one obtains
\[
j^*f^*i_*(1_Y)=j^*j_*(u).
\]
Putting $\xi=c_1(\mathcal O_{Y'}(1))$, it follows from 3.6.4 that
\[
g^*i^*i_*(1_Y)=-u\,\xi.
\tag{4.1.2}
\]
The element $u$ admits a unique decomposition of the form
\[
u=a_0+a_1\xi+\cdots+a_{d-1}\xi^{d-1},
\qquad
a_i\in H^{2d-2-2i}(Y,\mu_Y^{\otimes(d-1-i)}).
\tag{4.1.3}
\]
In view of the defining relation for the Chern classes of $\check N$,
\[
\xi^d+c_1(\check N)\xi^{d-1}+\cdots+c_{d-1}(\check N)\xi+c_d(\check N)=0,
\tag{4.1.4}
\]
we deduce the equality
\[
u\xi=-a_{d-1}c_d(\check N)+(a_0-a_{d-1}c_{d-1}(\check N))\xi
+\cdots+(a_{d-2}-a_{d-1}c_1(\check N))\xi^{d-1}.
\]
Comparing this last expression with (4.1.2), and using the injectivity of $g^*$ (2.2.6), one obtains
\[
i^*i_*(1_Y)=a_{d-1}c_d(\check N).
\tag{4.1.5}
\]
It remains to determine $a_{d-1}$. First, the homomorphism $g_*$ lowers degrees by $2d-2$, hence
\[
g_*(\xi^i)=0\quad(0\le i\le d-2),
\]
and moreover (2.2.2)
\[
g_*(\xi^{d-1})=1.
\]
Applying $g_*$ to both sides of 4.1.3, we obtain
\[
g_*(u)=a_{d-1}.
\tag{4.1.6}
\]
On the other hand, applying $f_*$ to both sides of 4.1.1, we obtain, in view of 4.2,
\[
i_*(1_Y)=i_*g_*(u).
\tag{4.1.7}
\]
But the homomorphism $i_*$ in question is injective; hence
\[
a_{d-1}=g_*(u)=1_Y,
\]
which completes the proof.

\begin{statement}{Remarks 4.4}
a) The essential ingredient in what precedes is the purity theorem. Once it is available in full generality, statement 4.1 will be valid with $Y$ regular and purely of codimension $d$ in $X$, and $X$ excellent and regular in a neighbourhood of $Y$. In this form, the result is in any case established in characteristic $0$, by (SGA 4 XIX).

b) When $S$ is the spectrum of a separably closed field and $X$ and $Y$ are quasi-projective and smooth over $S$, the analogous statement in the Chow ring of $Y$ was established by Mumford (cf. 9.2).
\end{statement}

\begin{statement}{Corollary 4.5}
Under the hypotheses of 4.1, let $x$ and $y$ belong to $A^*(Y)$. Then
\[
i_*(x)i_*(y)=i_*(xy\,c_d(\check N)).
\]
\end{statement}

\emph{Proof.} The projection formula for the morphism $i$ shows that
\[
i_*(x)i_*(y)=i_*(x\,i^*i_*(y)),
\]
and the conclusion follows from 4.1.

\begin{statement}{Corollary 4.6}
Let $S$ be a scheme, let $E$ be a locally free $\mathcal O_S$-module of constant rank $r$, and let $g:X=V(\check E)\to S$ be the vector bundle whose sheaf of sections is $E$. For every section $f$ of $X$ over $S$, one has
\[
f_*(1_S)=g^*(c_r(E))
\]
in $H^{2r}(X,\mu_X^{\otimes r})$.
\end{statement}

\emph{Proof.} Since $g^*$ is an isomorphism (1.2) and $g\circ f=\mathrm{id}$, $f^*$ is bijective. It is therefore enough to check that
\[
f^*f_*(1_S)=f^*g^*(c_r(E)).
\]
If $N$ denotes the normal sheaf of $f$, this last equality is equivalent, by 4.1, to
\[
c_r(\check N)=c_r(E).
\]
In fact one even has an isomorphism $\check N\simeq E$, since the exact sequence of (EGA IV 17.2.5), applied to the situation
\[
S\xrightarrow{f} X\xrightarrow{g} S,
\]
gives an isomorphism $N\simeq f^*\Omega^1_{X/S}$, and moreover $\Omega^1_{X/S}\simeq g^*\check E$ (EGA IV 16.4.8).

\begin{statement}{Corollary 4.7}
Under the hypotheses of 4.6, assume in addition that $S$ is smooth and quasi-projective over a separably closed field. Let $\sigma_X$ denote the zero-section cycle of $X$, and let $Y=f^{-1}(\sigma_X)$ be the cycle of zeros of $f$. Then
\[
\cl_S(Y)=c_r(E)
\]
in $H^{2r}(S,\mu_S^{\otimes r})$.
\end{statement}

\emph{Proof.} By IV,
\[
\cl_S(Y)=f^*\cl_X(\sigma_X).
\]
Let $s_0$ be the zero-section morphism of $X$. Then
\[
\cl_X(\sigma_X)\stackrel{\mathrm{dfn}}{=}(s_0)_*(1_S)\stackrel{(4.6)}{=}g^*(c_r(E)),
\]
and the assertion follows by applying $f^*$ to the extreme terms.

\begin{statement}{Remark 4.8}
Let $\ell$ be a prime number distinct from the residual characteristics of $S$. By simple passage to the projective limit, one deduces from 4.6 and 4.7 the corresponding statements in $\ell$-adic cohomology.
\end{statement}

Now let $\ell$ be a prime number and let $X$ be a proper scheme over a separably closed field $k$ of characteristic prime to $\ell$. Recall (IV) that the $\ell$-adic Euler-Poincaré characteristic of $X$ is the integer
\[
\operatorname{EP}_\ell(X)=\sum_i(-1)^i\dim H^i(X,\mathbb Q_\ell).
\]

\begin{statement}{Corollary 4.9}
Let $X$ be a smooth, proper, connected scheme of dimension $n$ over a separably closed field $k$, with $\operatorname{char}(k)\ne\ell$. If $p:X\to k$ denotes the canonical projection, then
\[
\operatorname{EP}_\ell(X)=p_*c_n(\check\Omega^1_X)
\]
in $\mathbb Q_\ell=H^0(k,\mathbb Q_\ell)$.
\end{statement}

\emph{Proof.} Let $j:X\to X\times X$ be the diagonal immersion of $X$, and let $q:X\times X\to k$ be the canonical projection. Applying 4.1 to the morphism $j$ gives
\[
c_n(\check\Omega^1_X)=j^*j_*(1_X),
\]
hence
\[
p_*(c_n(\check\Omega^1_X))=p_*j^*j_*(1_X)=q_*j_*j^*j_*(1_X).
\]
Using the projection formula for $j$, one deduces
\[
p_*(c_n(\check\Omega^1_X))=q_*(j_*(1_X)\cup j_*(1_X)),
\]
where, by IV, the second member is equal to the $\ell$-adic Euler-Poincaré characteristic of $X$.

\subsubsection*{4.10}
When $X$ is projective over $k$, the result recalled from IV shows that $\operatorname{EP}_\ell(X)$ is, in the Chow ring of $X\times X$, the self-intersection of the diagonal cycle, and consequently is independent of the prime number $\ell\ne\operatorname{char}(k)$. We shall therefore be able to speak of the Euler-Poincaré characteristic $\operatorname{EP}(X)$ of $X$, without reference to the integer $\ell$.

\begin{statement}{Proposition 4.11}
Let $X$ be a smooth, projective, connected scheme of dimension $n$ over a separably closed field $k$. Then
\[
\operatorname{EP}(X)=\sum_{p,q}(-1)^{p+q}\dim H^p(X,\Omega_X^q).
\]
\end{statement}

\emph{Proof.} We use the notation of the proof of 4.9, and let $J$ be the ideal of the diagonal immersion of $X$. As observed in 4.10, $\operatorname{EP}(X)$ is, in the Chow ring of $X\times X$, the intersection product $\Delta\cdot\Delta$ of the diagonal cycle $\Delta$ with itself. To compute it, one may neglect torsion and therefore work in $\operatorname{Gr}_{\mathrm{top}}^\bullet(X)$. This shows that $\operatorname{EP}(X)$ is the image under the direct-image morphism
\[
q_*:K(X\times X)\to K(k)=\mathbb Z
\]
of the square of the class of $\mathcal O_{X\times X}/J$. But, by (SGA 6 VII 2.5), this square is equal to
\[
\sum_i(-1)^i\lambda^i([\Omega_X^1]),
\]
and the desired formula follows immediately, by the definition of $q_*$.

\section*{5. Flag Schemes}

Let $S$ be a scheme, let $E$ be a locally free $\mathcal O_S$-module of rank $r$, and let $(p_i)_{1\le i\le m}$ be $m$ integers $>0$ satisfying
\[
p_1+p_2+\cdots+p_m=r.
\]
Recall that a flag of type $(p_1,\ldots,p_m)$ of $E$ is a finite filtration
\[
0=E_m\subset E_{m-1}\subset\cdots\subset E_0=E
\]
of $E$ such that, for every integer $i\in[1,m]$, the quotient $E_{i-1}/E_i$ is locally free of rank $p_i$.

One defines a functor $(\operatorname{Sch}/S)^\circ\to(\operatorname{Ens})$ by assigning to an $S$-scheme $u:T\to S$ the set of flags of type $(p_1,\ldots,p_m)$ of $u^*(E)$, and to a morphism of $S$-schemes the inverse image, in the evident sense, on flags. It is proved (EGA I, 2nd edition) that this functor is representable, and any representative is called an $S$-scheme of flags of type $(p_1,\ldots,p_m)$ of $E$.

In what follows, $f:X\to S$ will denote the $S$-scheme of flags of type $(p_1,\ldots,p_m)$ of $E$. Moreover, the scheme of flags of type $(1,1,\ldots,1)$ of $E$ will be denoted $D$ and called the flag scheme of $E$, without further specification, if no confusion is possible.

It is clear from the definition that the $\mathcal O_X$-module $f^*(E)$ admits a canonical finite filtration
\[
0=F_m\subset F_{m-1}\subset\cdots\subset F_0=f^*(E),
\]
for which the quotient $E_i=F_{i-1}/F_i$ $(1\le i\le m)$ is locally free of rank $p_i$. For every integer $i\in[1,m]$, let
\[
g_i:D_i\to X
\]
be the $X$-scheme of flags of $E_i$, and let $g:D_1\times_X\cdots\times_X D_m\to X$ be the fibre product of the $X$-schemes $D_i$. It is clear that the composite
\[
h=f\circ g:\prod_iD_i\to S
\tag{5.1}
\]
is $S$-isomorphic to the $S$-scheme of flags of $E$; up to $S$-isomorphism, one may therefore write
\[
D=\prod_iD_i.
\]

\begin{statement}{Proposition 5.2}
In $D^b(S,A)$ there is an isomorphism
\[
\mathrm Rf_*(A)\simeq\bigoplus_i R^{2i}f_*(A)[-2i].
\]
\end{statement}

\emph{Proof.} Since $h$ is a composite of finitely many morphisms of the form
\[
P_T(F)\to T\qquad(F\text{ a locally free }\mathcal O_T\text{-module}),
\tag{Q}
\]
it follows from (2.2.1) that
\[
\mathrm Rh_*(A)\simeq\bigoplus_j R^jh_*(A)[-j].
\]
For the same reason, one also has
\[
\mathrm Rg_*(A)\simeq\bigoplus_j R^jg_*(A)[-j].
\]
Using ([6] 2.16), one obtains an isomorphism
\[
\mathrm Rf_*(g_*A)\simeq\bigoplus_j R^j f_*(g_*A)[-j].
\]
Since $g$ is smooth with geometrically connected fibres, $A_X\simeq g_*(A_D)$; hence it remains to show that
\[
R^{2i+1}f_*(A)=0\qquad (i\in\mathbb N).
\tag{5.2.1}
\]
For this, one immediately reduces, using conservativity of the system of fibre functors (SGA 4 VIII 3.5), to the case where $S$ is the spectrum of a separably closed field $k$. Then, since $g$ is a composite of projective bundles, (2.2.4) shows that the inverse-image map
\[
H^{2i+1}(X,A_X)\to H^{2i+1}(D,A_D)
\]
is a monomorphism, so one may suppose that $X=D$. In this case, since the projection $D\to k$ is itself a composite of projective bundles, the assertion follows by induction from (2.2.4).

\begin{statement}{Corollary 5.3}
For every morphism of schemes $u:S\to T$, the Leray spectral sequence
\[
E_2^{pq}=R^p u_*\,R^q f_*(A)\Rightarrow R^{p+q}(u\circ f)_*(A)
\]
is degenerate.
\end{statement}

\emph{Proof.} Cf. [6] 1.2.

We shall now specify the multiplicative structure of the cohomology of $X$, first examining a particular case. With the preceding notation, one obtains a morphism of rings
\[
\alpha:A^*(S)[T_{ij_i}]_{1\le i\le m,\,1\le j_i\le p_i}\to A^*(X)
\]
by sending $A^*(S)$ into $A^*(X)$ by inverse image $f^*$ and, for every pair $(i,j_i)$, as above, sending the element $T_{ij_i}$ to the $j_i$-th Chern class $c_{j_i}(E_i)$. For every integer $p$, one has, by (3.4 (iii)),
\[
c_p(f^*E)=\sum_{i_1+i_2+\cdots+i_m=p}c_{i_1}(E_1)\cdots c_{i_m}(E_m),
\]
so that, by functoriality of Chern classes, the morphism vanishes on the ideal $J_X$ generated by the elements
\[
c_p(E)-\sum_{i_1+\cdots+i_m=p}T_{1,i_1}\cdots T_{m,i_m}
\]
for varying $p$. Hence there is a morphism of rings
\[
\varphi_X:\bigl(A^*(S)[T_{ij_i}]_{(i,j_i)}\bigr)/J_X\to A^*(X).
\]

\begin{statement}{Proposition 5.4}
The above morphism of rings $\varphi_X$ is an isomorphism. The $A^*(S)$-module $A^*(X)$ is free of rank
\[
\frac{r!}{p_1!p_2!\cdots p_m!}.
\]
\end{statement}

\emph{Proof.} We follow Grothendieck's exposé (Séminaire Chevalley 1958, Anneaux de Chow et applications, Exp. IV). We shall need the following lemma (cf. Lemma 1, p. 19 of loc. cit.).

\begin{statement}{Lemma 5.4.1}
a) Let $R$ be a commutative ring, and let $(a_i)_{1\le i\le r}$ be $r$ elements of $R$. Let $\sigma_i$ be the $i$-th elementary symmetric function of the elements $(U_i)_{1\le i\le r}$, and let $I$ be the ideal of the polynomial ring $R[U_1,U_2,\ldots,U_r]$ generated by the elements $\sigma_i-a_i$. Then the $R$-algebra
\[
R[U_1,U_2,\ldots,U_r]/I
\]
is a free $R$-module of rank $r!$.

b) Let $r,p_1,p_2,\ldots,p_m$ be integers $>0$ satisfying $p_1+\cdots+p_m=r$. Then the morphism of polynomial rings
\[
V=R[T_{i,j_i}]_{1\le i\le m,\,1\le j_i\le p_i}
\longrightarrow R[U_i]_{1\le i\le r}=W
\]
obtained by sending $T_{i,j_i}$ to the $j_i$-th elementary symmetric function of the elements
\[
(U_{p_1+\cdots+p_{i-1}+s})_{1\le s\le p_i}
\]
is injective and makes $W$ a free $V$-module of rank $p_1!\cdots p_m!$.

In particular, the morphism of $R$-algebras
\[
R[T_1,\ldots,T_r]\to R[U_1,\ldots,U_r]
\]
obtained by sending $T_i$ to the $i$-th elementary symmetric function of the $U_j$ is injective and makes the second member a free module of rank $r!$ over the first.
\end{statement}

We show how (5.4.1) implies (5.4). First suppose that $X$ is the scheme $D$ of flags of $E$, and for simplicity put $T_{i1}=U_i$. Since the canonical projection $D\to S$ is a composite of $r-1$ morphisms of the form
\[
P_{Y_i}(M_i)\to Y_i\quad(1\le i\le r-1),
\]
where $M_i$ is a locally free $\mathcal O_{Y_i}$-module of rank $r-i+1$, it follows from (2.2.6) that $A^*(D)$ is a free $A^*(S)$-module of rank $r!$, and that the morphism $\varphi_D$ is surjective. But Lemma 5.4.1 a) shows that the source of $\varphi_D$ is also a free $A^*(S)$-module of rank $r!$; hence the assertion in the case $X=D$.

In the general case, it is clear that the diagram
\[
\begin{tikzcd}[column sep=large,row sep=large]
A^*(X) \arrow[r,"g^*"] & A^*(D)\\
C=A^*(S)[T_{i,j_i}]/J_X \arrow[u,"\varphi_X"] \arrow[r,"\gamma"'] &
A^*(S)[U_k]/J_D \arrow[u,"\varphi_D"']
\end{tikzcd}
\]
commutes; here the map $\gamma$ is obtained by passage to the quotient from the morphism of $A^*(S)$-algebras
\[
A^*(S)[T_{i,j_i}]\to A^*(S)[U_k]
\]
defined by sending $T_{i,j_i}$ to the $j_i$-th elementary symmetric function of the elements
\[
(U_{p_1+\cdots+p_{i-1}+s})_{1\le s\le p_i}.
\]
Since $\varphi_D$ is an isomorphism, 5.4.1 b) implies that $A^*(D)$ is a free $C$-module of rank $p_1!\cdots p_m!$. On the other hand, expressing $g$ as a fibre product of $X$-schemes of flags shows that $A^*(D)$ is also a free $A^*(X)$-module of rank $p_1!\cdots p_m!$. It follows at once that $\varphi_X$ is an isomorphism.

In the statement of the next proposition, we agree to put, for every scheme $Y$ and every $A_Y$-module $M$,
\[
Q(Y,M)=\bigoplus_{n,i}H^n(Y,M\otimes\mu_Y^{\otimes i}).
\]
The $A$-module $Q(Y,M)$ is canonically endowed with a structure of $A^*(Y)$-module, coming from cup-product. If $u:Z\to Y$ is a morphism of schemes, the inverse-image morphism
\[
u^*:Q(Y,M)\to Q(Z,u^*M)
\]
is compatible with the $A^*(Y)$-module structure on $Q(Y,M)$ and the $A^*(Z)$-module structure on $Q(Z,u^*M)$. One deduces a morphism of $A^*(Z)$-modules, called canonical,
\[
A^*(Z)\otimes_{A^*(Y)}Q(Y,M)\to Q(Z,u^*M).
\]

\begin{statement}{Proposition 5.5}
Let $S$ be a scheme, and let $E$ be a locally free $\mathcal O_S$-module of rank $r$. Let $f:X\to S$ be the $S$-scheme of flags of type $(p_1,\ldots,p_m)$ of $E$. Then, for every $A_S$-module $L$, the canonical morphism
\[
\alpha:A^*(X)\otimes_{A^*(S)}Q(S,L)\to Q(X,f^*L)
\]
is an isomorphism.
\end{statement}

\emph{Proof.} When $X=P_S(E)$, this is an immediate consequence of (2.2.4) and (2.2.6). Hence it follows that it is true for $X=D$. In the general case, since the canonical morphisms $g:D\to X$ and $h:D\to S$ are composites of projective bundles, the canonical morphisms
\[
A^*(D)\otimes_{A^*(S)}Q(S,L)\to Q(D,h^*L)
\]
and
\[
A^*(D)\otimes_{A^*(X)}Q(X,g^*L)\to Q(D,g^*L)
\]
are isomorphisms. Consequently $\mathrm{id}_{A^*(D)}\otimes\alpha$ is an isomorphism, whence the result, since $A^*(D)$ is free over $A^*(X)$.

We shall now give a relative variant of (5.4) and (5.5). With the preceding notation, let $u:S\to T$ be a morphism of schemes. Put
\[
A^*(X/T)=\bigoplus_i R^{2i}(u\circ f)_*(A_X)
\]
and similarly
\[
A^*(S/T)=\bigoplus_i R^{2i}u_*(A_S).
\]
Finally, for every $A_S$-module $L$, put
\[
Q(S/T,L)=\bigoplus_{n,i}R^n u_*(L\otimes\mu_S^{\otimes i})
\]
and
\[
Q(X/T,f^*L)=\bigoplus_{n,i}R^n(u\circ f)_*(f^*L\otimes\mu_X^{\otimes i}).
\]

For every pair $(i,j)$ of integers, with $1\le i\le m$ and $1\le j\le p_i$, the image of the Chern class $c_j(E_i)$ by the canonical morphism
\[
H^{2j}(X,\mu_X^{\otimes j})
\to H^0(T,R^{2j}(u\circ f)_*(\mu_X^{\otimes j}))
\]
defines an $A_T$-morphism
\[
c_{ij}:A_T\to R^{2j}(u\circ f)_*(\mu_X^{\otimes j}).
\]
The data of the $c_{ij}$ and the inverse image $f^*$ allow one to define a morphism of $A^*(S/T)$-algebras
\[
\theta_1:A^*(S/T)[T_{i,j_i}]_{(i,j_i)}\to A^*(X/T).
\]
Denoting by $c_i$ the sections of $A^*(S/T)$ defined by the Chern classes $c_i(E)$, one sees as in (5.4) that $\theta_1$ vanishes on the ideal $J$ locally generated by the relations
\[
\sum_{i_1+\cdots+i_m=p}T_{1,i_1}\cdots T_{m,i_m}=c_p\qquad(p\in\mathbb N).
\]
Hence there is a morphism of $A^*(S/T)$-algebras
\[
\theta:A^*(S/T)[T_{i,j_i}]_{(i,j_i)}/J\to A^*(X/T).
\]

\begin{statement}{Proposition 5.6}
With the preceding notation:
\begin{enumerate}[label=\alph*)]
\item The morphism of $A^*(S/T)$-algebras
\[
\theta:A^*(S/T)[T_{i,j_i}]_{(i,j_i)}/J\to A^*(X/T)
\]
is an isomorphism.
\item For every $A_S$-module $L$, the canonical morphism of $A^*(X/T)$-modules
\[
A^*(X/T)\otimes_{A^*(S/T)}Q(S/T,L)\to Q(X/T,f^*L)
\]
is an isomorphism.
\end{enumerate}
\end{statement}

\emph{Proof.} Simple sheafification of (5.4) and (5.5), respectively.

\section*{6. Group Schemes}

We begin with a general property of objects twisted by torsors.

\begin{statement}{Proposition 6.1}
Let $S$ be a scheme, $G$ a quasi-compact smooth $S$-group scheme with connected fibres, and $X$ an $S$-scheme endowed with a left action of $G$. If $f:E\to S$ is a right $S$-torsor under $G$, denote by $X_E$ the $S$-scheme obtained from $X$ by twisting by means of $E$ and the action of $G$ on $X$. Let $p:X\to S$ and $q:X_E\to S$ be the canonical projections. Under these conditions, for every $A_S$-module $M$ and every integer $i$, the $A_S$-module $R^iq_*(q^*M)$ is canonically isomorphic to $R^ip_*(p^*M)$.
\end{statement}

\emph{Proof.}
\[
\begin{tikzcd}[column sep=large,row sep=large]
E\times_S X \arrow[r,"\varphi"] \arrow[d,"\pr_1"'] & X_E \arrow[d,"q"]\\
E \arrow[r,"f"'] & S
\end{tikzcd}
\qquad
\begin{tikzcd}[column sep=large,row sep=large]
E\times_S X \arrow[r,"\pr_2"] \arrow[d,"\pr_1"'] & X \arrow[d,"p"]\\
E \arrow[r,"f"'] & S .
\end{tikzcd}
\]
It is known that diagram $(D_1)$, where $\varphi$ denotes the canonical passage-to-the-quotient map, is cartesian. Since $f$ is smooth, the smooth base-change theorem gives an isomorphism
\[
u:f^*R^iq_*(q^*M)\xrightarrow{\sim}R^i(\pr_1)_*(\varphi^*q^*M).
\]
Similarly, the cartesian diagram $(D_2)$ gives an isomorphism
\[
v:f^*R^ip_*(p^*M)\xrightarrow{\sim}R^i(\pr_1)_*((\pr_2)^*p^*M).
\]
But $p\circ\pr_2=q\circ\varphi$, so after the evident identification of $\varphi^*\circ q^*$ with $(\pr_2)^*\circ p^*$ one obtains an isomorphism
\[
v^{-1}\circ u:f^*R^iq_*(q^*M)\xrightarrow{\sim} f^*R^ip_*(p^*M).
\]
Since the geometric fibres of $f$ are non-empty and connected, it follows from (SGA 4 XII 6.5 (i)) that the adjunction morphism
\[
\operatorname{adj}:\mathrm{id}\to f_*f^*
\]
is an isomorphism, and therefore
\[
(\operatorname{adj})^{-1}\circ f_*(v^{-1}\circ u)\circ(\operatorname{adj})
\]
is the desired isomorphism.

\begin{statement}{Proposition 6.2}
(Mme Raynaud [5].) Let $S$ be a scheme and let $p:G\to S$ be a reductive $S$-group scheme (SGA 3 XIX). For every locally constant constructible $A_S$-module $M$ and every integer $i$, the $A_S$-module $R^ip_*(p^*M)$ is locally constant and constructible. Moreover, for every morphism $u:S'\to S$ defining a cartesian diagram
\[
\begin{tikzcd}[column sep=large,row sep=large]
G' \arrow[r,"u'"] \arrow[d,"p'"'] & G \arrow[d,"p"]\\
S' \arrow[r,"u"'] & S,
\end{tikzcd}
\]
the base-change morphism
\[
u^*R^ip_*(p^*M)\to R^ip'_*(p\circ u')^*(M)
\]
is an isomorphism.
\end{statement}

\emph{Proof.} The two assertions of the statement are local for the étale topology of $S$. By (SGA 3 XII 2.3), we may therefore suppose that $G$ is split. Let $B$ be a Borel subgroup of $G$ of the type indicated in (SGA 3 XXII 5.1.1), so that the quotient $G/B$ exists by (SGA 3 XXII 5.8.5), and let $q:G\to G/B$ and $r:G/B\to S$ be the canonical morphisms. Since $r$ is proper and smooth, it satisfies the proper base-change theorem and the functors $R^ir_*$ send locally constant constructible abelian sheaves to sheaves of the same type (SGA 4 XVI 2.2). Moreover, by (SGA 4 XXII 5.5.1), the $S$-scheme $B$ is isomorphic to a finite product of vector bundles and punctured vector bundles, and consequently it follows immediately from (1.1) and (1.2) that, if $t:B\to S$ denotes the canonical morphism, the sheaves $R^it_*(t^*M)$, with $M$ a locally constant constructible $A_S$-module, are locally constant constructible and stable under base change. Now let $M$ be a locally constant constructible $A_S$-module. The preceding considerations and (6.1) imply that the sheaves $R^iq_*(p^*M)$ are locally constant constructible; using the Leray spectral sequence
\[
R^ir_*(R^jq_*)(p^*M)\Rightarrow R^{i+j}p_*(p^*M),
\]
one deduces the same for the sheaves $R^ip_*(p^*M)$. Finally, by (6.1), the sheaves $R^iq_*(p^*M)$ are stable under base change; using the proper base-change theorem for $r$, it follows that the sheaves $R^ip_*(p^*M)$ are also stable under base change.

\begin{statement}{Remark 6.3}
Since every reductive $S$-group is locally, for the étale topology, an inverse image by the canonical morphism $S\to\operatorname{Spec}(\mathbb Z)$ of a reductive $\operatorname{Spec}(\mathbb Z)$-scheme (SGA 3 XXIII 5.7 and XXV 1.3), (6.2) shows that the computation of the cohomology, with coefficients of torsion prime to the characteristic, of reductive group schemes over an algebraically closed field reduces to the transcendental problem of computing the cohomology of reductive group schemes of the same type over the field of complex numbers; this is known by the methods of algebraic topology, in view of the comparison theorem (SGA 4 XI 4.4).
\end{statement}

\begin{statement}{Corollary 6.4}
Let $S$ be a scheme and let $p:G\to S$ be an $S$-group scheme which is an extension of a finite étale $S$-group scheme $F$ by a reductive $S$-group scheme $H$. Then the conclusions of (6.2) are valid for every locally constant constructible $A_S$-module $M$.
\end{statement}

\emph{Proof.} After the base change $F\to S$, we may suppose that $F$ is of the form $\Gamma_S$, where $\Gamma$ is an ordinary finite group. Then there is an isomorphism of $S$-schemes
\[
G\simeq\coprod_{\gamma\in\Gamma}H,
\]
and the conclusion follows from (6.2).

\begin{statement}{Corollary 6.5}
Let $S$ be a scheme, $G$ a reductive $S$-group scheme, and $H$ a reductive sub-$S$-group scheme of $G$, such that the quotient sheaf of $G$ by $H$ for the faithfully flat quasi-compact topology is representable. Then, if
\[
p:G/H\to S
\]
is the canonical morphism, the conclusions of (6.2) are valid for every locally constant constructible $A_S$-module $M$.
\end{statement}

\emph{Proof.} Let $f:G\to G/H$ be the canonical morphism; it is smooth, hence of descent for étale sheaves (SGA 4 VIII 9.1). For every integer $r\ge0$, put
\[
G_r=G\times_{G/H}G\times_{G/H}\cdots\times_{G/H}G
\quad(r+1\text{ times}),
\]
which, by SGA 3 V 1.5, is canonically identified with
\[
G\times_S H\times_S\cdots\times_S H,
\]
where $H$ occurs $r$ times, and let $f_r:G_r\to G/H$ and $g_r:G_r\to S$ be the canonical morphisms, so that in particular $f_0=f$. Finally, for every integer $q\ge0$, let $H^q(M)$ denote the complex
\[
R^q(g_0)_*(f_0^*M)\to R^q(g_1)_*(f_1^*M)\rightrightarrows\cdots\to R^q(g_r)_*(f_r^*M)\cdots .
\]
% SGA-SRC-000365: the scan leaves the citation blank; reconstructed from the mathematics.
By (SGA 4, Exposé V bis, 2.5.5 and 2.5.6), one has a spectral sequence of descent, or Čech spectral sequence,
\[
E(M):E_2^{pq}=H^p(H^q(M))\Rightarrow E^n=R^np_*(M).
\]
Since the schemes $G_r$ are reductive group schemes, it follows from 6.2 that the terms $E_2^{pq}$ are locally constant constructible and compatible with every base change. It follows at once that the same is true for the $E^n$.

\section*{7. Complete Intersections}

In this paragraph a separably closed field $k$ is given. If $\ell$ is a prime number $\ne\operatorname{char}(k)$, we denote, for every proper $k$-scheme $X$ and every integer $i\ge0$, by
\[
b_\ell^i(X)\quad\text{or}\quad b^i(X)
\]
when no confusion is possible, the $i$-th Betti number of $X$ for coefficients $\mathbb Z_\ell$. If $X$ is projective over $k$, a hyperplane section of $X$ means the zero locus of a section transverse to the zero section of a vector bundle $V(\check{\mathcal L})$, with $\mathcal L$ an ample invertible $\mathcal O_X$-module.

\begin{statement}{Theorem 7.1}
(Lefschetz theorem.) Let $X$ be a smooth, connected, projective $k$-scheme of dimension $n$, let $Y$ be a subscheme of $X$ which is the intersection of $d$ hyperplane sections of $X$, and let $L$ be a locally constant constructible $A_X$-module. Then:
\begin{enumerate}[label=(\roman*)]
\item The inverse-image homomorphism
\[
H^i(X,L)\to H^i(Y,L|Y)
\]
is an isomorphism for $i\le n-d-1$ and a monomorphism for $i=n-d$.
\item Suppose $Y$ is smooth over $k$, of dimension $n-d$. Then the Gysin homomorphism
\[
H^i(Y,(L|Y)\otimes\mu_Y^{\otimes -d})\to H^{i+2d}(X,L)
\]
is an isomorphism for $i\ge n-d+1$ and an epimorphism for $i=n-d$.
\end{enumerate}
\end{statement}

\emph{Proof.} It is well known that the complement of a hyperplane section is an affine open; consequently $U=X-Y$ is a union of $d$ affine opens. Assertion (i) then follows immediately from the exact cohomology sequence
\[
\cdots\to H_!^i(U,L|U)\to H^i(X,L)\to H^i(Y,L|Y)
\xrightarrow{\partial} H_!^{i+1}(U,L|U)\to\cdots
\]
and from the following lemma.

\begin{statement}{Lemma 7.1.1}
Let $U$ be a $k$-scheme of finite type of dimension $n$, which is a union of $d$ affine open subsets, and let $F$ be an $A_U$-module. Then:
\begin{enumerate}[label=\alph*)]
\item if $F$ is constructible,
\[
H^i(U,F)=0\qquad\text{for }i\ge n+d;
\]
\item if $F$ is locally constant and $U$ is smooth over $k$,
\[
H^i_c(U,F)=0\qquad\text{for }i\le n-d.
\]
\end{enumerate}
\end{statement}

For $d=1$, assertion a) is just (SGA 4 XIV 3.2). In the general case, writing $\mathfrak U$ for a covering of $U$ by $d$ affine open subsets, one deduces a) from the Čech spectral sequence
\[
H^p(\mathfrak U,\mathcal H^q(F))\Rightarrow H^{p+q}(U,F).
\]
We prove b). Since the functors $H_c^i(U,\cdot)$ commute with filtered inductive limits (SGA 4 VII 3.3), one may suppose $F$ constructible; the assertion then follows from a), by Poincaré duality (SGA 4 XVIII).

Finally, in view of the definition of the Gysin morphism through the canonical isomorphism
\[
H^{i-2d}(Y,(L|Y)\otimes\mu_Y^{\otimes -d})\simeq H^i_Y(X,L),
\tag{7.1.2}
\]
assertion (ii) follows from 7.1.1 a) and from the exact cohomology sequence associated with an open subset and its complement
\[
\cdots\to H^{i-1}(U,L|U)\xrightarrow{\partial}H^i_Y(X,L)\to H^i(X,L)\to H^i(U,L|U)\to\cdots .
\]

\begin{statement}{Remark 7.1.3}
In fact, it follows from excision that the canonical isomorphism (7.1.2) generalizes to the case where the smoothness hypothesis on $X$, and the hypothesis that $L$ is constructible locally constant, are satisfied only in a neighbourhood of $Y$; hence 7.1 (ii) generalizes in this setting.

For substantially more general formulations, see (SGA 2 XIV 4.6), which, however, is proved only by means of resolution of singularities, and hence provisionally in characteristic $0$.
\end{statement}

\begin{statement}{Corollary 7.2}
Let $X$ be a smooth connected projective $k$-scheme of dimension $n$, let $Y$ be a subscheme of $X$ which is the intersection of $d$ hyperplane sections of $X$, and let $L$ be a constructible twisted constant $\mathbb Z_\ell$-sheaf $(\ell\ne\operatorname{car}(k))$. Then:
\begin{enumerate}[label=(\roman*)]
\item the inverse-image homomorphism
\[
\alpha_i:H^i(X,L)\to H^i(Y,L|Y)
\]
is an isomorphism for $i\le n-d-1$, and a monomorphism for $i=n-d$. If moreover $L$ is constructible locally free, the cokernel of $\alpha_{n-d}$ is torsion-free;
\item suppose that $Y$ is smooth over $k$, of dimension $n-d$. Then the Gysin homomorphism
\[
\beta_i:H^i(Y,(L|Y)(-d))\to H^{i+2d}(X,L)
\]
is an isomorphism for $i\ge n-d+1$ and an epimorphism for $i=n-d$. If moreover $L$ is constructible locally free, the epimorphism $\beta_{n-d}$ is split and its kernel is torsion-free.
\end{enumerate}
\end{statement}

\emph{Proof.} We prove (i), assertion (ii) following from it by Poincaré duality. In the general case, a simple passage to the limit starting from 7.1 (i) gives the conclusion. When $L$ is a locally free $\mathbb Z_\ell$-sheaf, apply the following lemma to the inverse-image morphism
\[
R\Gamma(X,L)\to R\Gamma(Y,L|Y)
\]
in $D(\mathbb Z_\ell)$; hypothesis (iii) is satisfied by 7.1 (i).

\begin{statement}{Lemma 7.2.1}
Let $V$ be a discrete valuation ring, let $k_V$ be its residue field, and let $u:K\to L$ be a morphism in $D(V)$. Suppose that $K$ and $L$ have cohomology of finite type, and denote by $M$ the mapping-cylinder of $u$. Then for every $m\in\mathbb Z$ the following conditions are equivalent:
\begin{enumerate}[label=(\roman*)]
\item $H^p(M)=0$ for $p\le m-1$, and $H^m(M)$ is torsion-free;
\item for every $V$-module $N$, the morphism
\[
H^p(u\otimes^L\operatorname{id}_N):H^p(K\otimes^L N)\to H^p(L\otimes^L N)
\]
is an isomorphism for $p\le m-1$ and a monomorphism for $p=m$;
\item[(ii bis)] for every $V$-module $N$, one has
\[
H^p(M\otimes^L N)=0\qquad(p\le m-1);
\]
\item the morphism
\[
H^p(u\otimes^L\operatorname{id}):H^p(K\otimes^L k_V)\to H^p(L\otimes^L k_V)
\]
is an isomorphism for $p\le m-1$ and a monomorphism for $p=m$;
\item[(iii bis)] one has
\[
H^p(M\otimes^L k_V)=0\qquad(p\le m-1).
\]
\end{enumerate}
Moreover these equivalent conditions imply:
\begin{enumerate}[label=(\roman* bis)]
\item the morphism
\[
H^p(u):H^p(K)\to H^p(L)
\]
is an isomorphism for $p\le m-1$, and a monomorphism with free cokernel for $p=m$.
\end{enumerate}
\end{statement}

\emph{Proof.} The following diagram of implications is evident:
\[
\begin{array}{ccc}
\text{(ii)}&\Leftrightarrow&\text{(ii bis)}\\
\Downarrow&&\\
\text{(iii)}&\Leftrightarrow&\text{(iii bis)}
\end{array}
\qquad\qquad
\text{(i)}\Rightarrow\text{(i bis).}
\]
We shall prove separately that (i)$\Rightarrow$(ii bis) and that (iii bis)$\Rightarrow$(i), which will suffice. The implication (i)$\Rightarrow$(ii bis) is immediate from the Künneth exact sequence
\[
0\to H^p(M)\otimes_V N\to H^p(M\otimes^L N)\to \operatorname{Tor}_1^V(H^{p+1}(M),N)\to0.
\]
The same exact sequence, for $N=k_V$, shows that if (iii bis) holds, then
\[
H^p(M)\otimes k_V=0\quad(p\le m-1),\qquad
\operatorname{Tor}_1^V(H^m(M),k_V)=0.
\]
One deduces (i) by Nakayama's lemma and the local criterion of flatness (EGA $0_{\mathrm{III}}$ 10.1.3).

Under the preceding hypotheses, one sees that knowledge of the cohomology groups of $X$ with coefficients in $\mathbb Q_\ell$ implies knowledge of those of $Y$, except in dimension $n-1$. The calculation of that remaining group is immediately reduced, under these conditions, to the calculation of the Euler-Poincaré characteristic of $Y$. This is the calculation to which we now turn, in a slightly more general setting.

\begin{statement}{Proposition 7.3}
Let $X$ be a proper smooth connected $k$-scheme of dimension $n$, let $E$ be a locally free $\mathcal O_X$-module of constant rank $r$, and let $Y$ be the zero-locus of a section of $V(\check E)$ transverse to the zero-section. Let $p:X\to k$ be the canonical morphism and $p_*$ the corresponding Gysin morphism. Then
\[
\operatorname{EP}(Y)=p_*\bigl(c(\check\Omega_X^1)c(E)^{-1}c_r(E)\bigr)
\]
in $H^0(k,\mathbb Q_\ell)=\mathbb Q_\ell$.
\end{statement}

\emph{Proof.} By (EGA IV 17.13.2), $Y$ is smooth and purely of dimension $n-r$. If $i:Y\to X$ denotes the canonical immersion, the conormal sheaf of $i$ is the inverse image by $i$ of the conormal sheaf $\check E$ to the zero-section of $V(\check E)$. Thus one has an exact sequence
\[
0\to i^*(\check E)\to i^*(\Omega_X^1)\to\Omega_Y^1\to0,
\tag{7.3.1}
\]
whence immediately
\[
c(\check\Omega_Y^1)=i^*\bigl(c(\check\Omega_X^1)c(E)^{-1}\bigr).
\tag{7.3.2}
\]
By (3.10) we therefore get
\[
\operatorname{EP}(Y)=p_*i_*i^*\bigl(c(\check\Omega_X^1)c(E)^{-1}\bigr),
\]
or, using the projection formula for $i$,
\[
\operatorname{EP}(Y)=p_*\bigl(c(\check\Omega_X^1)c(E)^{-1}i_*(1_Y)\bigr).
\]
The stated formula follows immediately from (4.7).

\begin{statement}{Corollary 7.4}
Let $X$ be a proper smooth connected $k$-scheme of dimension $n$, and let $Y$ be a smooth closed subscheme of dimension $n-1$ of $X$, hence defined by a section of an invertible $\mathcal O_X$-module $L$. Let $p$ be the canonical morphism from $X$ to $k$ and $p_*$ the corresponding Gysin morphism. Then
\[
\operatorname{EP}(X)-\operatorname{EP}(Y)=p_*\bigl(c(\check\Omega_X^1)c(L)^{-1}\bigr)
\tag{7.4.1}
\]
in $H^0(k,\mathbb Q_\ell)=\mathbb Q_\ell$. Moreover,
\[
b^{n-1}(Y)=2b^{\,n-1}(X)-b^n(X)+(-1)^n p_*\bigl(c(\check\Omega_X^1)c(L)^{-1}\bigr).
\tag{7.4.2}
\]
\end{statement}

\emph{Proof.} Formula (7.4.1) follows from 7.3 and 4.2, taking into account the evident identity
\[
1-c(L)^{-1}c_1(L)=c(L)^{-1}.
\]
Using (7.2), one sees that
\[
b^{n-1}(Y)=b^{n-1}(X)-b^n(X)+b^{n+1}(X)+(-1)^n(\operatorname{EP}(X)-\operatorname{EP}(Y)).
\]
But $b^{n-1}(X)=b^{n+1}(X)$ by Poincaré duality, and (7.4.2) follows.

\begin{statement}{Corollary 7.5}
Let $X$ be a smooth connected subscheme of dimension $n$ of $\mathbb P_k^{\,n+r}$, a global complete intersection of $r$ divisors $X_1,X_2,\ldots,X_r$ of $\mathbb P_k^{\,n+r}$. For each $i\in\{1,\ldots,r\}$, let $d_i$ be the degree of $X_i$. Then:
\begin{enumerate}[label=(\roman*)]
\item the cohomology groups $H^p(X,\mathbb Z_\ell)$ are free $\mathbb Z_\ell$-modules of finite type;
\item if $p$ is an integer $\le 2n$ and $\ne n$, then
\[
H^p(X,\mathbb Z_\ell)=
\begin{cases}
0,&\text{if }p\text{ is odd},\\
\mathbb Z_\ell,&\text{if }p\text{ is even};
\end{cases}
\]
\item the Euler-Poincaré characteristic of $X$ is given by
\[
\operatorname{EP}(X)
=d_1\cdots d_r
\sum_{\substack{0\le i\le n\\ j_1+\cdots+j_r=n-i}}
(-1)^{\,n-i}\binom{n+r+1}{i}
d_1^{j_1}\cdots d_r^{j_r}.
\]
\end{enumerate}
\end{statement}

\emph{Proof.} Assertion (i) follows immediately from (7.2), and so does (ii), in view of the cohomology of projective spaces (no. 2). We prove (iii). Let $P$ be the ambient projective space, let $p:P\to k$ be the canonical projection, let $p_*$ be the corresponding Gysin morphism, and let $\xi$ be the class of $\mathcal O_P(1)$. For every $i\in\{1,\ldots,r\}$, the scheme $X_i$ is defined by a section of an invertible $\mathcal O_P$-module $L_i$, so that $X$ is the zero-locus of a section of $V(\check L_1\oplus\cdots\oplus\check L_r)$ transverse to the zero-section. On the other hand, one has an exact sequence
\[
0\to\Omega_P^1\otimes\mathcal O_P(1)\to\mathcal O_P^{\,n+r+1}\to\mathcal O_P(1)\to0,
\]
hence
\[
c(\check\Omega_P^1)=(1+\xi)^{\,n+r+1}.
\]
Formula (7.3) gives
\[
\operatorname{EP}(X)
=p_*\bigl(c(\check\Omega_P^1)c(\check L_1\oplus\cdots\oplus\check L_r)^{-1}
c_1(\check L_1)\cdots c_1(\check L_r)\bigr)
\]
\[
=p_*\bigl(d_1\cdots d_r\,\xi^r(1+\xi)^{\,n+r+1}
(1+d_1\xi)^{-1}\cdots(1+d_r\xi)^{-1}\bigr).
\]
Using the fact that $p_*(\xi^{n+r})=1$ (2.2.2 c)), the formula of the statement follows without difficulty.

\begin{statement}{Remark 7.6}
One may also obtain the expression for the Euler-Poincaré characteristic by using (4.11) and by transcribing Hirzebruch's calculations ([3], pp. 160--161), which use the Riemann-Roch theorem. For a more complete calculation and the determination of the Hodge cohomology of complete intersections, the reader is referred to (SGA 7 XI).
\end{statement}

\subsection*{7.7}

By analogy with what happens in the transcendental case, where for constant coefficients one has another well-known theorem of Lefschetz ([4]), one is led to conjecture the following statement. If $X$ is a smooth connected projective $k$-scheme of dimension $n$, and if $\xi$ is the class in $H^2(X,\mathbb Z_\ell(1))$ of a very ample invertible $\mathcal O_X$-module, then for every twisted constant $\mathbb Q_\ell$-sheaf $L$ and every integer $p\le n$, multiplication by $\xi^{n-p}$:
\[
H^p(X,L)\to H^{2n-p}(X,L(n-p))
\]
is an isomorphism. As Deligne observed, the statements (7.7) are false without an additional hypothesis on $L$. However, Deligne proved them in \emph{La conjecture de Weil II} when, in addition, $L$ is pure.

This conjectural statement has the following consequences, likewise not proved in $\operatorname{car}\ne0$.

Let $X$ be a smooth connected projective $k$-scheme of dimension $n$, let $Y$ be a smooth hyperplane section of $X$, and let $L$ be a twisted constant $\mathbb Q_\ell$-sheaf on $X$.
\begin{enumerate}[label=(\roman*)]
\item The inverse-image morphism
\[
H^p(X,L)\to H^p(Y,L|Y)
\]
is an epimorphism for $p\ge n$.
\item The Gysin morphism
\[
H^p(Y,(L|Y)(-1))\to H^{p+2}(X,L)
\]
is a monomorphism for $p\le n-2$.
\end{enumerate}

\emph{Proof.} Let $j:X\hookrightarrow\mathbb P^r$ be a closed immersion and let $\xi_X$ (resp. $\xi_Y$) be the class in $H^2(X,\mathbb Q_\ell(1))$ (resp. $H^2(Y,\mathbb Q_\ell(1))$) of $\mathcal O_X(1)$ (resp. $\mathcal O_Y(1)$). Let $\alpha:Y\hookrightarrow X$ denote the canonical injection. For $q\le n-1$, the diagram
\[
\begin{array}{ccc}
H^q(X,L)&\xrightarrow{\alpha^*}&H^q(Y,L|Y)\\
\cup\,\xi_X^{\,n-1-q}\downarrow&&\downarrow\cup\,\xi_Y^{\,n-1-q}\\
H^{2n-2-q}(X,L(n-1-q))&\xrightarrow{\alpha^*}&
H^{2n-2-q}(Y,(L|Y)(n-1-q))
\end{array}
\]
is commutative and the right column is conjecturally an isomorphism. Put $p=2n-2-q$. For $q\le n-2$, hence $p\ge n$, the top row is an isomorphism by (7.2), giving (i). Assertion (ii) is proved similarly by considering, this time, the following diagram, with $p\le n-1$:
\[
\begin{array}{ccc}
H^p(Y,(L|Y)(-1))&\xrightarrow{\mathrm{Gysin}}&H^{p+2}(X,L)\\
\cup\,\xi_Y^{\,n-1-p}\downarrow&&\downarrow\cup\,\xi_X^{\,n-1-p}\\
H^{2n-2-p}(Y,(L|Y)(n-2-p))&\xrightarrow{\mathrm{Gysin}}&
H^{2n-p}(X,L(n-1-p)).
\end{array}
\]
In this diagram, the bottom row is an isomorphism for $2n-2-p\ge n$, i.e. for $p\le n-2$ (7.2 (ii)).

For a discussion of these conjectures and their consequences, the reader is referred to Kleiman's exposé, \emph{Algebraic cycles and the Weil conjectures}, in \emph{Dix exposés sur la cohomologie des schémas}.

\section*{8. Blown-up varieties}

In this section we give a cohomological variant of the calculations, for the group $K$ and for the ``group of cycle classes'' of a blown-up scheme, given in SGA 6 VII 3.7 and 4.8. The hypotheses and notations are those of the beginning of paragraph 4 (4.1 to 4.3). Under these conditions, the ``fundamental class'' morphism
\[
\gamma:j_*(\mu_{Y'}^{\otimes -1})[-2]\to A_{X'}
\]
gives, after applying the functor $Rf_*$, a morphism
\[
Rf_*(\gamma):i_*Rg_*(\mu_{Y'}^{\otimes -1})[-2]\to Rf_*(A_{X'}).
\]

\begin{statement}{Theorem 8.1}
The morphism in $D^+(X,A)$
\[
A_X\oplus i_*Rg_*(\mu_{Y'}^{\otimes -1})[-2]\to
Rf_*(A_{X'})\oplus i_*(\mu_Y^{\otimes -d})[-2d]
\tag{8.1.1}
\]
defined by the matrix
\[
\begin{pmatrix}
\operatorname{adj}&Rf_*(\gamma)\\
0&i_*(\operatorname{trace})
\end{pmatrix}
\]
is an isomorphism.
\end{statement}

\emph{Proof.} By conservativity of the system of fibre functors, it is enough to verify the assertion above each point $x$ of $X$. It is evident above $U=X-Y$. For $x\in Y$, proper base change reduces us to proving the following lemma, with $r=d-1$.

\begin{statement}{Lemma 8.2}
Let $k$ be a separably closed field, let $r\ge1$ be an integer, and let $P=\mathbb P_k^r$. Then:
\begin{enumerate}[label=\alph*)]
\item the canonical map $A\to H^0(P,A)$ is a bijection;
\item $H^1(P,A)=H^{2r-1}(P,A)=0$;
\item for $0\le p\le2r-2$, multiplication by $\xi=c_1(\mathcal O_P(1))\in H^2(P,\mu)$
\[
H^p(P,\mu_P^{\otimes -1})\to H^{p+2}(P,A)
\]
is an isomorphism;
\item the trace morphism
\[
H^{2r}(P,\mu_P^{\otimes -1})\to H^0(k,\mu_k^{\otimes -(r+1)})
\]
is a bijection.
\end{enumerate}
\end{statement}

The lemma is an immediate consequence of 2.2.2.

For the statement of the following corollary, we agree that, for every $X$-scheme $h:T\to X$, every closed subset $Z$ of $T$, and every complex $M$ of $A_X$-modules, we write
\[
H^*(T,M)=H^*(T,h^*(M)),\qquad
H^*_Z(T,M)=H^*_Z(T,h^*(M)).
\]

\begin{statement}{Corollary 8.3}
For every complex $M$ of $A_X$-modules and every $p\in\mathbb Z$, the maps
\[
H^p(X,M)\oplus H^{p-2}(Y',M\otimes\mu^{\otimes -1})
\to
H^p(X',M)\oplus H^{p-2d}(Y,M\otimes\mu^{\otimes -d})
\tag{8.3.1}
\]
and
\[
H_Y^p(X,M)\oplus H^{p-2}(Y',M\otimes\mu^{\otimes -1})
\to
H_{Y'}^p(X',M)\oplus H^{p-2d}(Y,M\otimes\mu^{\otimes -d})
\tag{8.3.2}
\]
defined by the matrix
\[
\begin{pmatrix}
f^*&j_*\\
0&g_*
\end{pmatrix}
\]
are bijections.
\end{statement}

\emph{Proof.} Temporarily denote by $u$ the morphism (8.1.1). By the projection formula, the arrows (8.3.1) and (8.3.2) are those induced on the $p$-th cohomology objects by the morphisms $R\Gamma(X,u\otimes\operatorname{id}_M)$ and $R\Gamma_Y(X,u\otimes\operatorname{id}_M)$ respectively. The assertion follows.

\subsection*{8.4}

We now make explicit the inverse isomorphism of (8.3.2), and for this we need a few lemmas.

\begin{statement}{Lemma 8.4.1}
Let $S$ be a scheme, let $E$ be a locally free $\mathcal O_S$-module of constant rank $d$, and let $p:P=\mathbb P_S(E)\to S$ be the associated projective bundle. Let $F$ be the $\mathcal O_P$-module defined by the canonical exact sequence
\[
0\to F\to p^*(E)\to\mathcal O_P(1)\to0.
\]
Then
\[
p_*(c_{d-1}(\check F))=1.
\]
\end{statement}

\emph{Proof.} Let $c_i$ denote the Chern classes of $E$, and put $\xi=c_1(\mathcal O_P(1))$. From
\[
p^*(c(E))=c(F)(1+\xi)
\]
it follows that
\[
c(F)=p^*(c(E))(1+\xi)^{-1},
\]
hence
\[
c_{d-1}(F)=\sum_{0\le i\le d-1}(-1)^{d-1-i}c_i\,\xi^{d-1-i}.
\tag{8.4.2}
\]
Now $p_*$ lowers degrees by $2d-2$, and $p_*(\xi^{d-1})=1$ (2.2.2). Applying $p_*$ to both sides of (8.4.2) and using the projection formula gives the lemma.

\begin{statement}{Lemma 8.4.3}
Let $F$ denote the locally free $\mathcal O_{Y'}$-module defined by the canonical exact sequence
\[
0\to F\to g^*(N)\to\mathcal O_{Y'}(1)\to0.
\]
Let $M$ be an $A_X$-module which is locally constant constructible in a neighbourhood of $Y$, and let $p\ge0$ be an integer. For every $y\in H^p(Y,M)$ one has
\[
f^*(i_*(y))=j_*(g^*(y)c_{d-1}(\check F))
\]
in $H_{Y'}^{p+2d}(X',M\otimes\mu^{\otimes d})$.
\end{statement}

\emph{Proof.} Put
\[
z=f^*(i_*(y))-j_*(g^*(y)c_{d-1}(\check F)).
\]
Then
\[
j^*(z)=g^*i^*i_*(y)-j^*j_*(g^*(y)c_{d-1}(\check F)),
\]
so, by (4.1) and the linearity lemma of (SGA 4 XVIII),
\[
j^*(z)=g^*(y)\bigl(c_d(g^*\check N)-c_{d-1}(\check F)c_1(\mathcal O_{Y'}(-1))\bigr)=0.
\tag{8.4.4}
\]
Writing $U=X-Y$ and $U'=X'-Y'$, one has an evident exact commutative diagram
\[
\begin{array}{ccccc}
H^{p+2d-1}(U')&\xrightarrow{\partial}&H^{p+2d}_{Y'}(X')&\xrightarrow{j^*}&H^{p+2d}(Y')\\
\uparrow\wr&&\uparrow f^*&&\uparrow g^*\\
H^{p+2d-1}(U)&\xrightarrow{\partial}&H_Y^{p+2d}(X)&\xrightarrow{i^*}&H^{p+2d}(Y),
\end{array}
\tag{8.4.5}
\]
with evident simplifications of notation. Considering (8.4.5) and (8.4.4), one sees that $z$ is of the form
\[
z=f^*(t),
\]
with $t\in H_Y^{p+2d}(X,M\otimes\mu^{\otimes d})$. By (4.2), $t=f_*(z)$, so it is enough to see that $f_*(z)=0$. But
\[
f_*(z)=f_*f^*i_*(y)-f_*j_*(g^*(y)c_{d-1}(\check F)),
\]
and using (4.2) gives
\[
f_*(z)=i_*(y)-i_*g_*(g^*(y)c_{d-1}(\check F)).
\]
The conclusion follows from the projection formula for $g$ and from (8.4.2).

\begin{statement}{Proposition 8.5}
Let $M$ be an $A_X$-module which is locally constant constructible in a neighbourhood of $Y$.
\begin{enumerate}[label=\alph*)]
\item Let
\[
\gamma:H_{Y'}^*(X',M\otimes\mu^{\otimes d})\to
H^*(Y',M\otimes\mu^{\otimes(d-1)})(-2)
\]
be the unique map defined by the relation
\[
f^*\circ f_* = j_*\circ\gamma+\operatorname{id}.
\tag{8.5.1}
\]
Then
\[
g_*\circ\gamma=0,
\tag{8.5.2}
\]
and the inverse matrix of (8.3.2) is
\[
\begin{pmatrix}
f_*&-i_*\\
-\gamma&g^*(\cdot)c_{d-1}(\check F)
\end{pmatrix}.
\tag{8.5.3}
\]
\item For every integer $p\ge0$, the sequence
\[
0\to H^{p-2d}(Y,M(-d))
\xrightarrow{\lambda}
H^{p-2}(Y',M(-1))\oplus H^p(X,M)
\xrightarrow{\mu}
H^p(X',M)\to0,
\]
where
\[
\lambda(x)=\bigl(g^*(x)c_{d-1}(\check F),-i_*(x)\bigr),
\qquad
\mu(x,y)=j_*(x)+f^*(y),
\]
is exact. Moreover, the map $\lambda$ has the left inverse
\[
\lambda'(x,y)=g_*(x).
\]
\end{enumerate}
\end{statement}

\emph{Proof.} We prove a). In view of 4.2 ($f_*f^*=\operatorname{id}$), in order to see that
\[
\begin{pmatrix}
f_*&-i_*\\
-\gamma&g^*(\cdot)c_{d-1}(\check F)
\end{pmatrix}
\circ
\begin{pmatrix}
f^*&j_*\\
0&g_*
\end{pmatrix}
=\operatorname{id},
\]
it is enough to prove
\[
\gamma j_*+\operatorname{id}=g^*(g_*(\cdot)c_{d-1}(\check F)),
\tag{8.5.4}
\]
and
\[
\gamma f^*=0.
\tag{8.5.5}
\]
From (8.5.1) one deduces
\[
f^*f_*f^*=j_*\gamma f^*+f^*,
\]
and hence, by (4.2), $j_*\gamma f^*=0$; this gives (8.5.5), since $j_*$ is an isomorphism. We prove (8.5.4). Again (8.5.1) gives
\[
f^*f_*j_*=j_*\gamma j_*+j_*,
\]
so, since $f_*j_*=i_*g_*$ and using (8.4.3),
\[
j_*g^*(g_*(\cdot)c_{d-1}(\check F))=j_*\gamma j_*+j_*.
\]
Multiplying on the left by the inverse of $j_*$ gives (8.5.4). Finally, (8.5.2) is obtained by writing that the product of the two matrices in the other order is the identity.

We prove b). From 8.4.3 and 8.4.1 respectively one gets $\mu\circ\lambda=0$ and $\lambda'\circ\lambda=\operatorname{id}$. Moreover, $\mu$ is surjective by (8.3). It remains to prove that $\Ker(\mu)\subset\operatorname{Im}(\lambda)$. But $\mu\circ(\lambda\circ\lambda')=0$ and $\lambda'\circ(\lambda\circ\lambda')=\operatorname{id}$, so $\mu(x)=0$ implies
\[
(\mu\oplus\lambda')(\lambda\circ\lambda'(x))=(\mu\oplus\lambda')(x),
\]
and hence $x=\lambda\circ\lambda'(x)$ by 8.3.

\subsection*{8.6}

We shall now try to determine the multiplicative structure of the cohomology of $X'$. For every $p\in\mathbb Z$ and every $A_X$-module $M$ which is locally constant constructible in a neighbourhood of $Y$, put
\[
\nabla^p(X',M)=H^{p-2}(Y',M(-1))\oplus H^p(X,M).
\]
For $p,q\in\mathbb Z$ and two such modules $M,N$, define a composition law
\[
\cup:\nabla^p(X',M)\otimes_A\nabla^q(X',N)\to
\nabla^{p+q}(X',M\otimes_A N)
\tag{8.6.1}
\]
by
\[
(y_1',x_1)\cup(y_2',x_2)=
\bigl(-y_1'y_2'\xi+y_1'g^*i^*(x_2)+(-1)^{pq}y_2'g^*i^*(x_1),\,x_1x_2\bigr).
\tag{8.6.2}
\]
This composition law is clearly anticommutative, and an easy but tedious calculation shows that it is associative. In particular,
\[
\nabla^*(X',A)=\bigoplus_{p\in\mathbb Z}\nabla^p(X',A)
\]
is also endowed with a structure of graded-commutative $A$-algebra, in the sense of algebras graded by $\mathbb Z$, and
\[
\nabla^*(X',M)=\bigoplus_{p\in\mathbb Z}\nabla^p(X',M)
\]
is endowed with a graded module structure over the graded ring $\nabla^*(X',A)$.

\begin{statement}{Proposition 8.6.3}
\begin{enumerate}[label=\alph*)]
\item For every pair $(M,N)$ of $A_X$-modules locally constant constructible in a neighbourhood of $Y$, the diagram
\[
\begin{tikzcd}
\nabla^*(X',M)\otimes\nabla^*(X',N)\arrow[r,"\mu\otimes\mu"]\arrow[d,"{(8.6.1)}"']&
H^*(X',M)\otimes H^*(X',N)\arrow[d,"\cup"]\\
\nabla^*(X',M\otimes N)\arrow[r,"\mu"]&
H^*(X',M\otimes N)
\end{tikzcd}
\]
is commutative. In particular, the surjective map
\[
\mu_A=\mu:\nabla^*(X',A)\to H^*(X',A)
\]
is a ring morphism, and
\[
\mu_M=\mu:\nabla^*(X',M)\to H^*(X',M)
\]
is a morphism of modules over variable rings.
\item If $\nabla^*(X',A)$ is made to act on $H^*(Y,M(-d))(-2d)$ by the map
\[
\nabla^p(X',A)\otimes H^{q-2d}(Y,M(-d))
\to H^{p+q-2d}(Y,M(-d)),
\qquad
(y',x)\otimes y\mapsto i^*(x)y,
\]
then the map $\lambda$ of 8.5 b) becomes a morphism of $\nabla^*(X',A)$-modules.
\end{enumerate}
\end{statement}

\emph{Proof.} We prove a). With the notation of 8.6.2, one must prove the equality
\[
j_*\bigl(-y_1'y_2'\xi+y_1'g^*i^*(x_2)+(-1)^{pq}y_2'g^*i^*(x_1)\bigr)+f^*(x_1x_2)
\]
\[
=\bigl(j_*(y_1')+f^*(x_1)\bigr)\bigl(j_*(y_2')+f^*(x_2)\bigr).
\]
It follows at once from
\[
\begin{cases}
f^*(x_1)f^*(x_2)=f^*(x_1x_2),&\text{multiplicativity of }f^*,\\
j_*(-y_1'y_2'\xi)=j_*(y_1')j_*(y_2'),&\text{(4.5)},\\
j_*(y'g^*i^*(x))=j_*(y'j^*f^*(x))=j_*(y')f^*(x),&\text{projection formula}.
\end{cases}
\]
We prove b). The equality to be proved,
\[
\lambda(i^*(x)\cdot y)=(y',x)\cdot\lambda(y),
\]
breaks up into the two following equalities:
\[
g^*(i^*(x)\cdot y)c_{d-1}(\check F)
=
-y'g^*(y)c_{d-1}(\check F)\xi
+y'g^*i^*(-i_*(y))
+g^*(i^*x)g^*(y)c_{d-1}(\check F),
\]
and
\[
-i_*(i^*(x)y)=-x\,i_*(y).
\]
The second is just the projection formula for $i$. The first follows immediately from
\[
-\xi\,c_{d-1}(\check F)=g^*[c_d(\check N)]
\quad\text{(exact sequence of 8.4.3)}
\]
and
\[
i^*i_*(y)=y\,c_d(\check N)\qquad\text{(4.1)}.
\]

\section*{9. Chow ring of a blown-up variety and the self-intersection formula in the Chow ring}

In this paragraph we propose to give Mumford's proof of the analogue of 4.1 in the Chow ring. We shall then show how one can deduce from it the ``key formula'' (8.4.3), and hence determine the structure of the Chow ring of a blown-up variety.

In accordance with the conventions of 3.9, $C^*(X)$ denotes the Chow ring of a quasi-projective smooth scheme $X$ over a field $k$. The theory of Chern classes used is that of [1]. We shall use the following lemma several times.

\begin{statement}{Lemma 9.1}
Let $S$ be a quasi-projective smooth scheme over a separably closed field $k$, let $E$ be a locally free $\mathcal O_S$-module of constant rank $r+1$, and let $p:P=\mathbb P_S(E)\to S$ be the associated projective bundle. Let $F$ be the $\mathcal O_P$-module defined by the canonical exact sequence
\[
0\to F\to p^*(E)\to\mathcal O_P(1)\to0,
\tag{9.1.0}
\]
and put $\xi=c_1(\mathcal O_P(1))$. Then:
\begin{enumerate}[label=(\roman*)]
\item for every $y\in C^*(S)$,
\[
y=p_*\bigl(p^*(y)c_r(\check F)\bigr);
\tag{9.1.1}
\]
\item suppose that $E=N\oplus\mathcal O_S$, with $N$ a locally free $\mathcal O_S$-module of rank $r$, so that $P$ identifies with the projective completion $\widehat N$ of $V(N)$, and let
\[
s:S\to\mathbb P_S(E)
\]
be the section obtained by composing the zero-section of $V(N)$ with the canonical inclusion $V(N)\hookrightarrow\widehat N$. Then:
\begin{enumerate}[label=(\alph*)]
\item for every $y\in C^*(S)$ one has
\[
s_*(y)=p^*(y)c_r(\check F)=p^*(y)\left(\sum_{i=0}^r p^*(c_i(\check N))\xi^{r-i}\right)
\quad\text{in }C^*(\widehat N),
\tag{9.1.2}
\]
\[
s^*s_*(y)=y\,c_r(\check N)\quad\text{in }C^*(S);
\tag{9.1.3}
\]
\item for every $z\in C^*(\widehat N)$,
\[
s^*(z)=p_*(z\,c_r(\check F))\quad\text{in }C^*(S).
\tag{9.1.4}
\]
\end{enumerate}
\end{enumerate}
\end{statement}

\emph{Proof.} We prove (9.1.1). By the projection formula for $p$, it is enough to show that $p_*(c_r(\check F))=1$. For this it suffices to copy the proof of 8.4.1, using the relations ([CH] 4 (17)). The equality of the first and third members of (9.1.2) is just Lemma 3 of [1]. Paraphrasing the proof of (8.4.1), one sees that
\[
c_r(\check F)=\sum_{i=0}^r p^*(c_i(\check N))\xi^{r-i},
\tag{9.1.2 bis}
\]
which proves (9.1.2). Applying $s^*$ to the two extreme members of (9.1.2), and using $p\circ s=\operatorname{id}$, one sees that (9.1.3) will follow from
\[
s^*\left(\sum_{i=0}^{r-1}p^*(c_i(\check N))\xi^{r-i}\right)=0,
\]
and a fortiori from
\[
s^*(\xi^i)=0\qquad(i>0).
\tag{9.1.5}
\]
We prove (9.1.5). If $\alpha:V(N)\hookrightarrow\widehat N$ denotes the canonical inclusion, then evidently $\alpha^*(\mathcal O_P(1))\simeq\mathcal O_{V(N)}$, whence
\[
\alpha^*(\xi)=\alpha^*(c_1(\mathcal O_P(1)))=c_1(\mathcal O_{V(N)})=0
\]
by functoriality of Chern classes. The assertion follows immediately from the definition of $s$. We prove (9.1.4). We shall use the equality
\[
\xi\,c_r(\check F)=c_{r+1}(p^*(\check N)\oplus\mathcal O'_P)=0,
\tag{9.1.6}
\]
which follows at once from the exact sequence (9.1.0). The element $z$ has a unique decomposition
\[
z=\sum_{i=0}^r p^*(a_i)\xi^i,\qquad a_i\in C^*(S),
\]
([CH] 4 Th. 1). It follows from (9.1.5) that $s^*(z)=a_0$. On the other hand, (9.1.6) shows that $z\,c_r(\check F)=p^*(a_0)c_r(\check F)$, and the conclusion follows from (9.1.1).

We now state the self-intersection formula in the Chow ring.

\begin{statement}{Theorem 9.2 (Mumford)}
Let $k$ be a separably closed field and let $i:Y\hookrightarrow X$ be a closed immersion, purely of codimension $d$, of quasi-projective smooth $k$-schemes. If $N$ denotes the normal sheaf of $i$, then for every $y\in C^*(Y)$ one has
\[
i^*i_*(y)=y\,c_d(\check N)
\]
in $C^*(Y)$.
\end{statement}

For the proof we shall need a number of notations and auxiliary schemes. Embed $X$ in $Z=X\times_k\mathbb P^1_k$ by the morphism $u:t\mapsto(t,0)$, so that there is a cartesian diagram
\[
\begin{tikzcd}
\widehat N \arrow[r,"j_1"] \arrow[d,"g_1"'] & Z' \arrow[d,"f_1"]\\
Y \arrow[r,"i_1"'] & Z ,
\end{tikzcd}
\]
where $Z'$ denotes the blowing-up of $Z$ along $Y$, and $\widehat N=\mathbb P_Y(N\oplus\mathcal O_Y)$ is the projective completion of the vector bundle $V(N)$. As in the statement of 9.1, let $s:Y\hookrightarrow\widehat N$ be the morphism obtained by composing the zero-section of $V(N)$ with the canonical inclusion $V(N)\subset\widehat N$.

The inclusion $Y\subset X$ identifies $Y\times\mathbb P^1$ with a closed subscheme of $X\times\mathbb P^1$. Let $W'$ be the strict transform of $Y\times\mathbb P^1$ by $f_1$, i.e. the subscheme of $Z'$ obtained by blowing up $Y\times\mathbb P^1$ along $Y=Y\times0$. It is clear that $W'\cap\widehat N$ identifies with the zero-section of $\widehat N$ and that, since $Y$ has codimension $1$ in $Y\times\mathbb P^1$, the canonical morphism $f_2:W'\to Y\times\mathbb P^1$ is an isomorphism. The remaining notation relative to $W'$ is fixed by the canonical diagram $(D_1)$.
\[
\begin{tikzcd}[column sep=3.2em,row sep=2.4em,cells={nodes={inner sep=1pt}}]
\text{(D}_1\text{)}\arrow[from=1-3,to=5-2,"g_1"']
& & \widehat N \arrow[rr,"j_1"] & & Z' \arrow[ddd,"f_1"] \\
& Y \arrow[ur,"s"] \arrow[rr,"\delta"'] \arrow[ddd,"\mathrm{id}"'] & &
W' \arrow[ur,"\beta"'] \arrow[ddd,"f_2"'] & \\
& & & & \\
& & & & X\times\mathbb P^1=Z \\
& Y \arrow[rr,"\gamma"'] & & Y\times\mathbb P^1 \arrow[ur,"\alpha"'] &
\end{tikzcd}
\]

Similarly, if $X'$ denotes the closed subscheme of $Z'$ which is the blowing-up of $X=X\times0$ along $Y$, it is clear that $\mathbb P(N)=X'\cap\widehat N$ identifies with the hyperplane at infinity $H$ of $\widehat N$. The remaining notation relative to $X'$ is fixed by the canonical diagram $(D_2)$.
\[
\begin{tikzcd}[column sep=3.2em,row sep=2.4em,cells={nodes={inner sep=1pt}}]
\text{(D}_2\text{)}\arrow[from=1-3,to=5-2,"g_1"']
& & \widehat N \arrow[rr,"j_1"] & & Z' \arrow[ddd,"f_1"] \\
& H \arrow[ur,"k"] \arrow[rr,"j"'] \arrow[ddd,"g"'] & &
X' \arrow[ur,"v"'] \arrow[ddd,"f"'] & \\
& & & & \\
& & & & X\times\mathbb P^1=Z \\
& Y \arrow[rr,"i"'] & & X \arrow[ur,"u"'] &
\end{tikzcd}
\]

We now give, in lemma form, a number of relations useful below. Denote by $E$ the locally free $\mathcal O_{\widehat N}$-module defined by the canonical exact sequence
\[
0\to E\to g_1^*(N\oplus\mathcal O_Y)\to\mathcal O_{\widehat N}(1)\to0,
\tag{9.3.0}
\]
and put
\[
\bar\xi=c_1(\mathcal O_{\widehat N}(1)).
\tag{9.3.1}
\]

\begin{statement}{Lemma 9.3}
With the preceding notation, one has
\[
j_1^*(\beta_*(1))=c_d(\check E)\quad\text{in }C^*(\widehat N),
\tag{9.3.2}
\]
\[
j_1^*(v_*(1))=\bar\xi\quad\text{in }C^*(\widehat N),
\tag{9.3.3}
\]
\[
f_1^*u_*(1)=(j_1)_*(1)+v_*(1)\quad\text{in }C^*(Z').
\tag{9.3.4}
\]
\end{statement}

\emph{Proof.} All the $k$-schemes under consideration are smooth, so it is immediate for dimension reasons that the morphisms $j_1$ and $\beta$ (resp. $j_1$ and $v$) are transverse. From ([CH] 4 Cor. 1 to Prop. 2) one obtains
\[
\begin{cases}
j_1^*\beta_*(1)=s_*(1),\\
j_1^*v_*(1)=k_*(1).
\end{cases}
\]
Equations (9.3.2) and (9.3.3) then come from (9.1.2) and from the normalization of Chern classes ([CH] 4--18), respectively. We prove (9.3.4). Since the cycle $X'$ of $Z'$ is the strict transform of $X\times0$ by $f_1$, one has a priori an equality of the form
\[
f_1^*u_*(1)=(j_1)_*g_1^*(\mu)+v_*(1),
\tag{9.3.5}
\]
with $\mu\in\Gamma(Y,\mathbb N)$; it remains to show that $\mu=1$. Applying $j_1^*$ to both sides of (9.3.5), one obtains, cf. $(D_2)$,
\[
g_1^*i^*u^*u_*(1)=j_1^*(j_1)_*g_1^*(\mu)+j_1^*v_*(1).
\tag{9.3.6}
\]
A very special case of (9.1.3), or a direct verification, shows that $u^*u_*(1)=0$. Using (9.3.3), it follows from (9.3.6) that
\[
0=j_1^*(j_1)_*g^*(\mu)+\bar\xi.
\tag{9.3.7}
\]
After restricting to a neighbourhood of a connected component of $Y$, one may suppose $Y$ connected, so that (9.3.7) takes the form
\[
0=\mu\,j_1^*(j_1)_*(1)+\bar\xi,\qquad \mu\in\mathbb N.
\tag{9.3.7 bis}
\]
Since $j_1^*(j_1)_*(1)$ is homogeneous of degree $1$, it is of the form
\[
j_1^*(j_1)_*(1)=g_1^*(\rho)+\nu\bar\xi,
\]
with $\rho\in C^1(Y)$ and $\nu\in\mathbb Z$. Substituting this expression into (9.3.7 bis) gives at once $\rho=0$ and $\mu\nu=-1$, with $\mu\in\mathbb N$; hence $\mu=1$ and $\nu=-1$, completing the proof of (9.3.4).

The key to the proof of 9.2 lies in the proof of the following lemma, which would be evident if one knew, as in cohomology, that
\[
j_1^*(j_1)_*(y)=-y\bar\xi
\]
for every $y\in C^*(\widehat N)$.

\begin{statement}{Lemma 9.4}
For every $y\in C^*(\widehat N)$, one has
\[
c_d(\check E)\,j_1^*(j_1)_*(y)=0.
\]
\end{statement}

\emph{Proof.} The element $y$ has a unique decomposition in $C^*(\widehat N)$ of the form
\[
y=\sum_{0\le p\le d}g_1^*(a_p)\bar\xi^p,
\]
or, by (9.3.3),
\[
y=\sum_{0\le p\le d}g_1^*(a_p)\bigl(j_1^*v_*(1)\bigr)^p.
\tag{9.4.1}
\]
Using the projection formula ([CH] 4 I 9) for $j_1$ and (9.1.6), it follows that
\[
c_d(\check E)\,j_1^*(j_1)_*(y)=c_d(E)\,j_1^*(j_1)_*g_1^*(a_0);
\]
in other words, we are reduced to proving 9.4 when $y$ is of the form $g_1^*(a)$, with $a\in C^*(Y)$. Then $j_1^*(j_1)_*g_1^*(a)$ has in turn a unique decomposition in $C^*(\widehat N)$ of the form
\[
j_1^*(j_1)_*g_1^*(a)=\sum_{0\le p\le d}g_1^*(b_p)\bar\xi^p,
\tag{9.4.2}
\]
and, using (9.1.6) once again, it remains to prove that $g_1^*(b_0)c_d(\check E)=0$. In fact, we shall prove the stronger statement $b_0=0$. In view of (9.3.3), we may rewrite (9.4.2) as
\[
j_1^*(j_1)_*g_1^*(a)=g_1^*(b_0)+z\,j_1^*v_*(1),
\tag{9.4.3}
\]
with $z\in C^*(\widehat N)$. Applying $(j_1)_*$ to both sides of (9.4.3) and using the projection formula for $j_1$, one gets
\[
(j_1)_*g_1^*(b_0)=\bigl[(j_1)_*g_1^*(a)\bigr](j_1)_*(1)-v_*(1)(j_1)_*(z).
\tag{9.4.4}
\]
Substituting in (9.4.4) the expression for $(j_1)_*(1)$ deduced from (9.3.4) gives
\[
(j_1)_*g_1^*(b_0)=
(j_1)_*g_1^*(a)f_1^*u_*(1)-v_*(1)\bigl[(j_1)_*g_1^*(a)+(j_1)_*(z)\bigr].
\tag{9.4.5}
\]
Since $X\times0$ is linearly equivalent to $X\times1$ in $X\times\mathbb P^1$, one has
\[
(j_1)_*g_1^*(a)f_1^*u_*(1)=0.
\tag{9.4.6}
\]
Comparing with (9.4.5), we get
\[
(j_1)_*g_1^*(b_0)=-v_*(1)\bigl[(j_1)_*g_1^*(a)+(j_1)_*(z)\bigr].
\tag{9.4.7}
\]
Moreover, it is clear that $W'\cap X'=\emptyset$: they could meet at most above $Y$, where one gives the zero-section and the other the hyperplane at infinity of $\widehat N$. Hence $\beta^*v_*(1)=0$, and from (9.4.7)
\[
\beta^*(j_1)_*g_1^*(b_0)=0.
\tag{9.4.8}
\]
But, as already observed in the proof of 9.3.2, $\beta$ and $j_1$ are transverse; thus (9.4.8) gives
\[
0=\delta_*s^*g_1^*(b_0)=\delta_*(b_0).
\tag{9.4.9}
\]
Finally, applying $(f_2)_*$ to (9.4.9), one obtains $\gamma_*(b_0)=0$, whence $b_0=0$, since $\gamma_*$ is injective ($\gamma$ admits an evident retraction).

\subsection*{9.5}

We finally show how the preceding results imply the theorem. Let $p:X\times\mathbb P^1_k\to X$ be the first projection, so that $i=p\circ i_1$. If $y\in C^*(Y)$, then evidently
\[
i^*i_*(y)=(i_1)^*(\bar x),\qquad \bar x=p^*i_*(y).
\tag{9.5.1}
\]
By (9.1.1) it follows that
\[
i^*i_*(y)=(g_1)_*\bigl[g_1^*i_1^*(\bar x)c_d(\check E)\bigr],
\]
hence
\[
i^*i_*(y)=(g_1)_*\bigl[j_1^*f_1^*(\bar x)\cdot c_d(\check E)\bigr].
\tag{9.5.2}
\]
To prove the theorem, by linearity we may suppose that $y$ is the irreducible cycle associated with a subvariety $T$ of $Y$, so that $\bar x$ corresponds to the cycle $T\times\mathbb P^1$ of $Z=X\times\mathbb P^1$. Let $\theta$ denote the strict transform of $T\times\mathbb P^1$ by $f_1$, i.e. the closure of the inverse image of $T\times\mathbb P^1$ over $Z\setminus Y$. Then
\[
f_1^*(\bar x)=\theta+(j_1)_*(w),\qquad w\in C^*(\widehat N).
\tag{9.5.3}
\]
Substituting this expression into (9.5.2), 9.4 gives
\[
i^*i_*(y)=(g_1)_*\bigl[j_1^*(\theta)c_d(\check E)\bigr],
\tag{9.5.4}
\]
and, by (9.1.4),
\[
i^*i_*(y)=s^*j_1^*(\theta).
\tag{9.5.5}
\]
But the reduced scheme underlying $\theta$ is the blowing-up of $T\times\mathbb P^1$ along $T=(T\times\mathbb P^1)\cap Y$; it meets the zero-section of $\widehat N$ transversely along the image of $T$ by that zero-section. Hence
\[
j_1^*(\theta)=s_*(y),
\]
which, together with (9.5.5) and (9.1.3), completes the proof of the theorem.

We now turn to the ``key formula'' mentioned in the introduction. The adapted notation, recalled for the most part, is that of 4. Observe that it is compatible with the notation introduced in the present paragraph, except that the scheme denoted $Y'$ in 4 is here denoted $H$ (for ``hyperplane''). We shall use both notations concurrently.

\begin{statement}{Proposition 9.6}
Let $k$ be a separably closed field and let $i:Y\hookrightarrow X$ be a closed immersion, purely of codimension $d$, of quasi-projective smooth $k$-schemes. Let $f:X'\to X$ be the $X$-scheme obtained by blowing up $X$ along $Y$, and let
\[
\begin{tikzcd}
Y'\arrow[r,"j"]\arrow[d,"g"']&X'\arrow[d,"f"]\\
Y\arrow[r,"i"']&X
\end{tikzcd}
\]
be the cartesian diagram obtained from $f$ and $i$. If $F$ denotes the locally free $\mathcal O_{Y'}$-module defined by the canonical exact sequence
\[
0\to F\to g^*(N)\to\mathcal O_{Y'}(1)\to0,
\tag{9.6.1}
\]
then
\[
f^*i_*(y)=j_*(g^*(y)c_{d-1}(\check F))
\tag{9.6.2}
\]
in $C^*(X')$, for every $y\in C^*(Y)$.
\end{statement}

Before proving 9.6, we isolate in a lemma some properties of the diagrams $(D_1)$ and $(D_2)$ which have not yet been established.

\begin{statement}{Lemma 9.7}
One has
\[
c_d(\check E)=k_*(c_{d-1}(\check F))+(g_1)^*(c_d(\check N)),
\tag{9.7 i}
\]
\[
\beta_*f_2^*\gamma_*(y)=f_1^*\alpha_*\gamma_*(y)\qquad\text{for every }y\in C^*(Y),
\tag{9.7 ii}
\]
\[
f_1^*u_*(t)=(j_1)_*((i\circ g_1)^*(t))+v_*f^*(t)\qquad\text{for every }t\in C^*(X).
\tag{9.7 iii}
\]
\end{statement}

\emph{Proof.} We prove (i). Put $\xi=c_1(\mathcal O_{Y'}(1))=k^*(\bar\xi)$. Paraphrasing the proof of (8.4.1), one obtains
\[
c_d(\check E)=\sum_{p=0}^d g_1^*(c_p(\check N))\bar\xi^{\,d-p},
\tag{9.7.1}
\]
\[
c_{d-1}(\check F)=\sum_{q=0}^{d-1}g^*(c_q(\check N))\xi^{d-1-q}.
\tag{9.7.2}
\]
From (9.7.2) one gets
\[
k_*[c_{d-1}(\check F)]
=\sum_{q=0}^{d-1}k_*k^*[g_1^*(c_q(\check N))\bar\xi^{\,d-1-q}],
\]
and hence, by the projection formula for $k$ and the fact that $k_*(1)=\bar\xi$,
\[
k_*[c_{d-1}(\check F)]=\sum_{q=0}^{d-1}g_1^*(c_q(\check N))\bar\xi^{\,d-q}.
\tag{9.7.3}
\]
Assertion (i) follows at once by comparing (9.7.1) and (9.7.3). We prove (ii). We may suppose that $y$ is an irreducible cycle of $Y$, corresponding to a subvariety denoted by the same letter. Then the cycle $\gamma_*(y)=y\times0$ is rationally equivalent to $y\times1$, so it is enough to verify the equality of cycles
\[
\beta_*f_2^*(y\times1)=f_1^*\alpha_*(y\times1).
\tag{9.7.4}
\]
Since $(y\times1)\cap(y\times0)=\emptyset$, $f_1$ is an isomorphism above an open neighbourhood of $y\times1$, which gives (9.7.4). We prove (iii). If $p:X\times\mathbb P^1\to X$ is the first projection, then (9.1.2) gives
\[
u_*(t)=p^*(t)u_*(1),
\tag{9.7.5}
\]
and comparing with (9.3.4) yields
\[
f_1^*u_*(t)=(p\circ f_1)^*(t)\bigl[(j_1)_*(1)+v_*(1)\bigr].
\tag{9.7.6}
\]
The projection formulas for $j_1$ and $v$ then give
\[
f_1^*u_*(t)=(j_1)_*\bigl[(p\circ f_1\circ j_1)^*(t)\bigr]+v_*\bigl[(p\circ f_1\circ v)^*(t)\bigr],
\tag{9.7.7}
\]
which is (iii), since $p\circ f_1\circ j_1=i\circ g_1$ and $p\circ f_1\circ v=f$.

\subsection*{9.8}

We now prove 9.6. We shall first prove
\[
v_*f^*i_*(y)=v_*j_*(g^*(y)c_{d-1}(\check F)),
\tag{9.8.1}
\]
and then prove that $v_*$ is injective. From (9.7)(iii) one gets
\[
v_*f^*i_*(y)=f_1^*u_*i_*(y)-(j_1)_*[g_1^*i^*i_*(y)],
\]
and therefore, by (9.2),
\[
v_*f^*i_*(y)=f_1^*(i_1)_*(y)-(j_1)_*[g_1^*(y\,c_d(\check N))].
\tag{9.8.2}
\]
On the other hand,
\[
v_*j_*(g^*(y)c_{d-1}(\check F))
=(j_1)_*\bigl[k_*k^*g_1^*(y)c_{d-1}(\check F)\bigr],
\]
or, by the projection formula for $k$,
\[
v_*j_*(g^*(y)c_{d-1}(\check F))
=(j_1)_*\bigl[g_1^*(y)k_*(c_{d-1}(\check F))\bigr].
\]
Using (9.7)(i), this gives
\[
v_*j_*(g^*(y)c_{d-1}(\check F))
=(j_1)_*[g_1^*(y)c_d(\check E)]-(j_1)_*[g_1^*(y\,c_d(\check N))].
\tag{9.8.3}
\]
Comparing (9.8.2) and (9.8.3), one sees that (9.8.1) is equivalent to
\[
f_1^*(i_1)_*(y)=(j_1)_*[g_1^*(y)c_d(\check E)],
\tag{9.8.4}
\]
i.e. to the equality corresponding to (9.6.2) when $X$ is replaced by $X\times\mathbb P^1$. By (9.1.2),
\[
(j_1)_*[g_1^*(y)c_d(\check E)]=(j_1)_*s_*(y)=\beta_*\delta_*(y).
\]
But it is immediate (cf. $(D_1)$) that $\delta_*=f_2^*\gamma_*$, whence:
\[
(j_1)_*[g_1^*(y)c_d(\check E)]=\beta_*f_2^*\gamma_*(y).
\tag{9.8.5}
\]
Also, evidently,
\[
f_1^*(i_1)_*(y)=f_1^*\alpha_*\gamma_*(y),
\tag{9.8.6}
\]
and (9.8.4) follows by comparing (9.8.5), (9.8.6), and (9.7)(ii).

\begin{statement}{Lemma 9.8.7}
The map $v_*$ is injective.
\end{statement}

\emph{Proof.} Let $Z''$ be the blowing-up of $Z'$ along $W'$, or equivalently along $\widehat N\cup W'=f_1^{-1}(Y\times\mathbb P^1)$. The strict transform $X''$ of $X$ in $Z''$ is also the strict transform of $X'$ in $Z''$. The notation is determined by the canonical diagram
\[
\begin{tikzcd}
X''\arrow[r,"v''"]\arrow[d,"f'"',"\wr"]&Z''\arrow[d,"f'_1"]\\
X'\arrow[r,"v"']&Z'.
\end{tikzcd}
\]
Since $X'\cap W'=\emptyset$, the morphism $f'_1$ is an isomorphism above an open neighbourhood of $X'$, and in particular $f'$ is an isomorphism. Moreover, since $f'_1$ is transverse to $v$, one has
\[
(f'_1)^*\circ v_*=v''_*\circ(f')^*,
\tag{9.8.8}
\]
so it is enough to show that $v''_*$ is injective. But, by the commutativity property of blow-ups (see, for example, Hironaka: \emph{Smoothing of algebraic cycles}, Amer. Journal of Math. (1968), Lemma 4.1), $Z''$ can also be obtained by first blowing up $X\times\mathbb P^1$ along $Y\times\mathbb P^1$, yielding a scheme $Z'_1$, and then blowing up $Z'_1$ along the inverse image of $Y$. Clearly $Z'_1=X'\times\mathbb P^1$, and the projection $Z'_1\to X\times\mathbb P^1$ is the product $f\times\operatorname{id}_{\mathbb P^1}$; hence the evident cartesian diagram
\[
\begin{tikzcd}[column sep=huge,row sep=large]
H\arrow[r,hook,"{t\mapsto(t,0)}"]\arrow[d,"g"']&H\times\mathbb P^1\arrow[r,hook,"{j\times\operatorname{id}}"]\arrow[d,"{g\times\operatorname{id}}"']&X'\times\mathbb P^1=Z'_1\arrow[d,"{f\times\operatorname{id}}"]\\
Y\arrow[r,hook,"\gamma"]&Y\times\mathbb P^1\arrow[r,hook,"{i\times\operatorname{id}}"]&X\times\mathbb P^1
\end{tikzcd}
\]
The strict transform of $X\times0$ by $f\times\operatorname{id}$ identifies with $X'\times0$, which contains $H=H\times0=(f\times\operatorname{id})^{-1}(Y)$. If we now blow up $Z'_1$ along $H$, then, since $H$ is purely of codimension $1$ in $X'$, the strict transform of $X'$ is canonically isomorphic to $X'$ (as was already known); hence the canonical commutative diagram
\[
\begin{tikzcd}[column sep=large,row sep=large]
&X'\arrow[r,hook,"v''"]\arrow[d,"f''"']&Z''\arrow[d,"f''_1"]\\
H\arrow[r,hook]&X'\arrow[r,hook,"v'"]&X'\times\mathbb P^1\\
&t\arrow[r,mapsto]&{(t,0)}
\end{tikzcd}
\]
Recall that we wish to prove that $v''_*$ is injective. Comparing with $(D_2)$, we see that this is equivalent to proving that $v_*$ is injective when, in addition, $Y$ has codimension $1$ in $X$; this is what we now do. But then the identity $(f_1)_*\circ v_*=u_*\circ f_*$ gives the conclusion, since $f_*$ is an isomorphism and $u_*$ is a direct monomorphism ($u$ admits a retraction).

We now state the theorem corresponding to (8.5 b).

\begin{statement}{Theorem 9.9}
For every integer $p\ge0$, the sequence
\[
0\to C^{p-d}(Y)\xrightarrow{\lambda} C^{p-1}(Y')\oplus C^p(X)\xrightarrow{\mu} C^p(X')\to0
\]
in which $\lambda$ and $\mu$ are the maps
\[
\lambda:x\mapsto\bigl(g^*(x)c_{d-1}(\check F),-i_*(x)\bigr),
\qquad
\mu:(x,y)\mapsto j_*(x)+f^*(y)
\]
is exact. Moreover, the map $\lambda$ has the left inverse
\[
\lambda':(x,y)\mapsto g_*(x).
\]
\end{statement}

\emph{Proof.} The last assertion is just (9.1.1). Likewise, (9.6.2) says that $\mu\circ\lambda=0$. We prove that $\mu$ is surjective. We shall need
\[
f_*f^*(z)=z\qquad(z\in C^*(X)),
\tag{9.9.1}
\]
which follows immediately from the projection formula for $f$ and the evident fact that $f_*(1)=1$. Let $t\in C^p(X')$. By (9.9.1),
\[
f_*(t-f^*f_*(t))=0.
\]
Since $f$ induces an isomorphism from $U'=X'-Y'$ to $U=X-Y$, the restriction of $t-f^*f_*(t)$ to $U'$ is zero, and hence ([CH] 4 th. 2) there is an equality
\[
t-f^*f_*(t)=j_*(a)
\qquad\text{with }a\in C^*(Y').
\]
We now prove exactness in the middle. Let $a\in C^{p-1}(Y')$ and $b\in C^p(X)$ satisfy
\[
j_*(a)+f^*(b)=0.
\tag{9.9.2}
\]
Applying $f_*$ to both sides gives
\[
f_*j_*(a)+f_*f^*(b)=0,
\]
and in view of (9.9.1),
\[
b=-i_*g_*(a).
\tag{9.9.3}
\]
Putting $x=g_*(a)$, it remains only to prove that
\[
a=g^*(x)c_{d-1}(\check F).
\tag{9.9.4}
\]
Put
\[
z=a-g^*(x)c_{d-1}(\check F)\qquad(x=g_*(a)),
\]
and prove $z=0$. We shall first show that $j_*(z)=0$ and $g_*(z)=0$, and then that an element $z$ satisfying these two relations is necessarily zero. By definition,
\[
j_*(z)=j_*(a)-j_*(g^*(x)\cdot c_{d-1}(\check F)),
\]
so, by (9.6.2),
\[
j_*(z)=j_*(a)-f^*i_*(x).
\]
Replacing $i_*(x)$ by $-b$ using (9.9.3), one deduces from (9.9.2) that $j_*(z)=0$. By the definition of $x$,
\[
g_*(z)=x-g_*(g^*(x)c_{d-1}(\check F))=0
\]
by the last assertion of the theorem. It remains to explain why the vanishing of $z$ follows from $j_*(z)=0$ and $g_*(z)=0$. Since $Y'=\mathbb P_Y(N)$, $z$ has a unique expression
\[
z=g^*(a_0)+g^*(a_1)\xi+\cdots+g^*(a_{d-1})\xi^{d-1},
\]
where $a_i\in C^{p-1-i}(Y)$. Since $g_*(z)=0$, the projection formula for $g$, together with the fact that $g_*$ lowers degrees by $d-1$, gives $a_{d-1}=0$. On the other hand, by (9.2),
\[
z\,\xi=j^*j_*(z)=0,
\]
hence
\[
g^*(a_0)\xi+\cdots+g^*(a_{d-1})\xi^{d-1}=0.
\]
Thus all $a_i$ are zero and $z=0$.

\subsection*{9.10}

With the notation of (9.9), the map
\[
\omega=\lambda\circ\lambda'
\]
is a projector satisfying $\operatorname{Im}(\omega)=\operatorname{Im}(\lambda)$. Therefore $\mu$ induces an isomorphism
\[
\bar\mu:\Ker(\omega)\xrightarrow{\sim}C^p(X').
\]
But
\[
\Ker(\omega)=\Ker(g_*)\oplus C^p(X)\hookrightarrow C^{p-1}(Y')\oplus C^p(X).
\tag{9.10.1}
\]
Indeed, an element $(x,y)$ of $C^{p-1}(Y')\oplus C^p(X)$ belongs to $\Ker(\omega)$ if and only if
\[
g^*g_*(x)\cdot c_{d-1}(\check F)=-i_*g_*(x)=0,
\]
that is, in view of (9.1.1), if and only if $g_*(x)=0$. On the other hand, the structure of the Chow ring of a projective bundle, and in particular the relations
\[
g_*(\xi^i)=0\quad\text{for }i<d-1,\qquad g_*(\xi^{d-1})=1,
\]
show that
\[
\Ker(g_*)=g^*C^{p-1}(Y)\oplus g^*C^{p-2}(Y)\xi\oplus\cdots\oplus g^*C^{p-d+1}(Y)\xi^{d-2}.
\]
In other words, $\bar\mu$ may be interpreted as an isomorphism
\[
C^{p-1}(Y)\oplus C^{p-2}(Y)\oplus\cdots\oplus C^{p-d+1}(Y)\oplus C^p(X)
\xrightarrow{\sim}C^p(X'),
\tag{9.10.2}
\]
whose inverse is easy to exhibit.

As Murre pointed out to me, the statement (9.10.2) was proved for $d=2$ by Samuel in \emph{Relations d'équivalence}, Proceedings of the International Congress, Edinburgh 1958.

Of course, in the framework of étale or $\ell$-adic cohomology one has an analogue of (9.10.2), deduced this time from (8.5).

\subsection*{Bibliography}

\begin{enumerate}[label={[\arabic*]}]
\item Grothendieck, \emph{La théorie des classes de Chern}, Bull. Soc. Math. de France \textbf{86} (1958), pp. 137--154.
\item Hodge, \emph{The theory and applications of harmonic integrals}, Cambridge University Press, 1952.
\item Hirzebruch, \emph{Topological Methods in Algebraic Geometry}, Springer Verlag.
\item Lefschetz, \emph{Algebraic Geometry}, Princeton University Press, 1953.
\item Mme Raynaud, \emph{Modules projectifs universels}, Inventiones Mathematicae vol. 6. fasc. 1 (1968), pp. 1--26.
\item[{[TF]}] Godement, \emph{Théorie des faisceaux}.
\item[{[CH]}] \emph{Anneaux de Chow et applications}, Séminaire Chevalley 1958.
\item[{[6]}] \emph{Théorème de Lefschetz et critères de dégénérescence}, Pub. I.H.E.S. n$^\circ$ 35.
\end{enumerate}


\newpage
\phantomsection\addcontentsline{toc}{section}{Exposé VIII --- Perfect Complexes}
\section*{Exposé VIII}

\begin{center}
\textbf{GROUPS OF CLASSES OF ABELIAN}\\
\textbf{AND TRIANGULATED CATEGORIES. PERFECT COMPLEXES.}\\[1ex]
by A. Grothendieck, written up by I. Bucur
\end{center}

In the present exposé, we quickly outline the notions and results which will be indispensable to us in the following exposés, concerning in particular pseudo-coherent and perfect complexes. A more complete exposé on the latter, with more detailed proofs, will appear in the first exposés of the seminar on the Riemann-Roch theorem for schemes, which will follow the present seminar.\footnote{SGA 6 I}

\subsection*{1. The case of abelian categories}

Let $\cC$ be an abelian category and let $\cD$ be a full subcategory of $\cC$. A function $f$ from $\cD$ to an abelian group $G$ is said to be additive if, for every exact sequence in $\cC$
\[
0\longrightarrow E'\longrightarrow E\longrightarrow E''\longrightarrow 0,
\]
such that $E'$, $E''$, and $E$ are objects of $\cD$, one has
\[
f(E)=f(E')+f(E'').
\]
Let $f$ be an additive function from $\cD$ to an abelian group $G$. The following properties are easy consequences of the definition:
\[
\alpha)\quad f(0)=0,\qquad
\beta)\quad f(X)=f(Y)\quad\text{if }X\text{ and }Y\text{ are isomorphic,}
\]
\[
\gamma)\quad f(X\oplus Y)=f(X)+f(Y).
\]

\[
\delta)\quad
\text{Let }E=(E^i,d^i)\text{ be a finite complex in }\cD
\]
such that the objects $B^{i+1}=\operatorname{Im}(d^i)$, $Z^i=\operatorname{Ker}(d^i)$ and $H^i(E)=Z^i/B^i$ are in $\cD$. Then
\[
\sum_i (-1)^i f(E^i)=\sum_i (-1)^i f(H^i(E)).
\]

\[
\varepsilon)\quad
\text{If }F\in\Ob \cD\text{ has a finite filtration}
\]
\[
0=F_0\subset F_1\subset\cdots\subset F_n=F,
\]
and if one assumes that, for $0<i\le n$, $F_i$ and $F_i/F_{i-1}$ are objects of $\cD$, then
\[
f(F)=\sum_{i=1}^n f(F_i/F_{i-1}).
\]

\[
\varphi)\quad
\text{If }\cD\text{ is stable under infinite direct sums, then }f=0.
\]

To prove $\delta)$, one uses the exact sequences
\[
0\longrightarrow Z^i\longrightarrow E^i\longrightarrow B^{i+1}\longrightarrow 0,
\]
\[
0\longrightarrow B^i\longrightarrow Z^i\longrightarrow H^i(E)\longrightarrow 0.
\]
The proof of $\varepsilon)$ is by induction on $n$, using the exact sequence
\[
0\longrightarrow F_{n-1}\longrightarrow F\longrightarrow F/F_{n-1}\longrightarrow 0.
\]
To prove $\varphi)$, let $E\in\Ob\cD$ and
\[
F=\bigoplus_{n\in\mathbb Z^+}E_n,\qquad E_n=E,\quad n\in\mathbb Z^+.
\]
Then $E\oplus F\simeq F$, and applying $\gamma)$ gives
\[
f(F)=f(F)+f(E),
\]
hence $f(E)=0$.

If, for every abelian group $G$, one considers the abelian group $A_{\cD}(G)$ of additive functions from $\cD$ to $G$, one obtains in an evident way a covariant functor
\[
A_{\cD}:\Ab\longrightarrow\Ab.
\]
The functor $A_{\cD}$ is representable. A representing pair $(K_{\cD}(\cC),\varphi)$ may be obtained as follows. Let $L$ be the free abelian group generated by the objects of $\cD$, and let $R$ be the subgroup of $L$ generated by the elements of the form $E-E'-E''$ for each exact sequence
\[
0\longrightarrow E'\longrightarrow E\longrightarrow E''\longrightarrow 0
\]
of $\cC$. Put
\[
K_{\cD}(\cC)=L/R.
\]
The map $\varphi$ is obtained by composing the two canonical maps
\[
\Ob\cD\longrightarrow L\longrightarrow L/R.
\]
When there is no danger of confusion, one may abbreviate the notation:
\[
K_{\cD}(\cC)=K(\cD).
\]

\subsection*{2. The case of triangulated categories}

Let $\cC$ be a triangulated category and let $G$ be an abelian group. A function $f$ from $\cC$ to $G$ is said to be additive if, for every distinguished triangle
\[
\begin{tikzcd}[column sep=small,row sep=small]
& Z \arrow[dl,"w"']\\
X \arrow[rr,"u"'] && Y \arrow[ul,"v"']
\end{tikzcd}
\]
one has
\[
f(Y)=f(X)+f(Z).
\]

Given an additive function $f$, the following relations are immediate consequences of the definition:
\[
\text{a) } f(0)=0,\qquad
\text{b) } f(X[n])=(-1)^n f(X),
\]
\[
\text{c) } f(X)=f(Y)\text{ if }X\text{ and }Y\text{ are isomorphic,}
\qquad
\text{d) } f(X\oplus Y)=f(X)+f(Y).
\]
Indeed the following triangles are distinguished:
\[
\begin{tikzcd}[column sep=small,row sep=small]
&0 \arrow[dl]\\ X \arrow[rr,"1_X"']&&X\arrow[ul]
\end{tikzcd}
\qquad
\begin{tikzcd}[column sep=small,row sep=small]
&X[1]\arrow[dl]\\ X\arrow[rr]&&0\arrow[ul]
\end{tikzcd}
\]
\[
\begin{tikzcd}[column sep=small,row sep=small]
&0\arrow[dl]\\ X\arrow[rr,"u"']&&Y\arrow[ul]
\end{tikzcd}
\qquad
\begin{tikzcd}[column sep=small,row sep=small]
&Y[1]\arrow[dl]\\ X\oplus Y\arrow[rr]&&X\arrow[ul]
\end{tikzcd}
\]
($u$ an isomorphism). To form the last triangle, we used the fact that the sum of two distinguished triangles is distinguished.

As in the case of abelian categories, one obtains in an evident way a functor
\[
A:\Ab\longrightarrow\Ab
\]
which associates to each abelian group $G$ the abelian group $A(G)$ of additive functions on $\cC$ with values in $G$. This functor is representable, and it is easy to construct a representing pair $(K(\cC),\operatorname{cl}_{\cC})$. Namely, one first considers the free abelian group $\mathbb Z^{(\Ob\cC)}$ generated by the objects of $\cC$, and then the subgroup $R$ generated by the elements of the form $X-Y-Z$, where
\[
\begin{tikzcd}[column sep=small,row sep=small]
&Z\arrow[dl]\\ Y\arrow[rr]&&X\arrow[ul]
\end{tikzcd}
\]
is a distinguished triangle. The quotient group $\mathbb Z^{(\Ob\cC)}/R$ is $K(\cC)$, and $\operatorname{cl}_{\cC}$ is obtained by composing the two canonical maps
\[
\Ob\cC\longrightarrow \mathbb Z^{(\Ob\cC)}\longrightarrow K(\cC).
\]

\subsection*{3. Functorial character}

Let $\cC$ and $\cC'$ be two triangulated categories, and let $T:\cC\to\cC'$ be an exact functor from $\cC$ to $\cC'$, that is, an additive ``graded'' functor carrying distinguished triangles to distinguished triangles. Denote, as above, by
\[
A,A':\Ab\longrightarrow\Ab
\]
the functors associated with the categories $\cC$ and $\cC'$. There results a functorial morphism $A'\to A$ given by composition with $T$, hence a homomorphism
\[
K(T):K(\cC)\longrightarrow K(\cC')
\]
uniquely determined by the commutativity of the diagram
\[
\begin{tikzcd}[column sep=large,row sep=large]
\Ob\cC \arrow[r,"T"] \arrow[d,"\operatorname{cl}_{\cC}"']&
\Ob\cC' \arrow[d,"\operatorname{cl}_{\cC'}"]\\
K(\cC) \arrow[r,"K(T)"']& K(\cC').
\end{tikzcd}
\]
It follows that, if the functors $T,T_1:\cC\to\cC'$ are isomorphic, then $K(T)=K(T_1)$. In particular, if $T$ is an equivalence of categories, then $K(T)$ is an isomorphism.

If $\cC''$ is a third triangulated category and
\[
F:\cC\times\cC'\longrightarrow \cC''
\]
is a bifunctor exact in each variable, then it is clear that there results a bilinear map
\[
K(F):K(\cC)\times K(\cC')\longrightarrow K(\cC'')
\]
such that the following relation holds:
\[
K(F)(\operatorname{cl}_{\cC}(X),\operatorname{cl}_{\cC'}(X'))
=
\operatorname{cl}_{\cC''}(F(X,X')).
\]

\begin{statement}{Proposition 3.1}
Let $\cA$ be a triangulated category, $\cB$ a thick triangulated subcategory of $\cA$, $\cA/\cB$ the quotient category of $\cA$ by $\cB$, $Q:\cA\to\cA/\cB$ the canonical quotient functor, and $i:\cB\to\cA$ the inclusion functor. Then the following sequence is exact:
\[
K(\cB)\xrightarrow{K(i)}K(\cA)\xrightarrow{K(Q)}K(\cA/\cB)\longrightarrow 0.
\]
\end{statement}

\emph{Proof.} It is clear that there is a natural homomorphism
\[
\operatorname{Coker}(K(i))\xrightarrow{\alpha}K(\cA/\cB).
\]
On the other hand, there is an evident map
\[
\Ob(\cA/\cB)\xrightarrow{\gamma}\operatorname{Coker}(K(i)),
\]
since $\Ob(\cA/\cB)=\Ob(\cA)$. Let $X,Y\in\Ob(\cA/\cB)$ and suppose that they are isomorphic. We want to show that $\gamma(X)=\gamma(Y)$. The category $\cA/\cB$ is obtained by a calculus of fractions on the right or on the left, the multiplicative system $S$ being formed by the morphisms $f$ which are contained in a distinguished triangle $(X',Y',Z',f,g,h)$ in which $Z'$ is an object of $\cB$. ($S$ is a saturated multiplicative system.) It follows that there exist $X_1$ and morphisms
\[
X_1\xrightarrow{s}X,\qquad X_1\xrightarrow{t}Y
\]
such that $s,t\in S$. It is therefore enough to prove that $\gamma(X)=\gamma(X_1)$, and this follows from the fact that there is a distinguished triangle of the form
\[
(X_1,X,Z,s,\ldots),
\]
with $Z\in\Ob\cB$.

The map $\gamma$ is an additive function. Indeed, if $(X,Y,Z,f,g,h)$ is a distinguished triangle of $\cA/\cB$, one must prove that
\[
\gamma(X)+\gamma(Z)=\gamma(Y).
\]
But the relation now follows because the triangle $(X,Y,Z,f,g,h)$ is isomorphic to a distinguished triangle of $\cA$. Consequently one obtains
\[
\beta:K(\cA/\cB)\longrightarrow \operatorname{Coker}(K(i)),
\]
and it is immediate to check that $\beta$ is the inverse of $\alpha$.

\subsection*{4. Comparison with the case of abelian categories}

Let $\cC$ be an abelian category and $D(\cC)$ the derived category of $\cC$. As usual, denote by $D^b(\cC)$ the full subcategory of $D(\cC)$ generated by the bounded complexes. By associating with every object $X$ of $\cC$ the complex whose components are all zero except the component in dimension zero, which is equal to $X$, one obtains a homomorphism
\[
K(\cC)\xrightarrow{\mu}K(D^b(\cC)).
\]
This homomorphism is in fact an isomorphism. Indeed, we may use Proposition 3.1 to construct an inverse to $\mu$, taking into account that
\[
D^b(\cC)\simeq K^{b,b}(\cC)/K^{b,\phi}(\cC),
\]
where $K^{b,\phi}(\cC)$ is the thick subcategory of $K^{b,b}(\cC)$ consisting of the acyclic complexes. Let $X^\bullet=(X_i)\in\Ob(K^{b,b}(\cC))$. If one sets
\[
f(X^\bullet)=\sum_i(-1)^i\operatorname{cl}_{\cC}(X_i),
\]
one obtains an additive function from $K^{b,b}(\cC)$ with values in $K(\cC)$. Indeed, one first observes that, if $X^\bullet$ and $Y^\bullet$ are homotopic, then $f(X^\bullet)=f(Y^\bullet)$, since $f(X^\bullet)=f(H(X^\bullet))$. It follows that there is a homomorphism
\[
K(D^b(\cC))\xrightarrow{\nu}K(\cC),
\]
and it is immediate to check that it is the inverse of $\mu$.

Since the invariant $K(\cC)$ is often of no interest when $\cC$ is too large, it is useful to relativize it with respect to a subcategory $\cC_o$ of $\cC$. Thus let $\cC_o$ be a full subcategory of $\cC$. We may define, on the one hand, the group
\[
K_o(\cC)=K_{\cC_o}(\cC).
\]
Suppose, on the other hand, that the following condition is satisfied:
\[
(A)\qquad \cC_o\text{ is stable under finite direct sums.}
\]
It is then easy to define the triangulated category $D_o^b(\cC)$ by using complexes in $\cC$ with bounded cohomology and with components in the subcategory $\cC_o$, and passing to the quotient by quasi-isomorphisms of such complexes. As above, there results a homomorphism
\[
(*)\qquad K_o(\cC)\longrightarrow K(D_o^b(\cC)).
\]
Suppose that the subcategory $\cC_o$ also satisfies the following condition (B):
\[
(B)\qquad
\text{If }X,Y\in\Ob\cC_o\text{ and }X\xrightarrow{u}Y
\text{ is an epimorphism (in }\cC\text{, of course)},
\]
then $\operatorname{Ker}u\in\Ob\cC_o$. Reasoning as above, one can show that the homomorphism $(*)$ is an isomorphism.

The canonical inclusion functor
\[
D_o^b(\cC)\longrightarrow D^b(\cC)
\]
is not fully faithful in general. However:

\begin{statement}{Proposition 4.1}
Suppose that the subcategory $\cC_o$ satisfies conditions (A), (B), and the following condition (C):
\[
(C)\qquad
\text{Every diagram }
\begin{tikzcd}[column sep=large,row sep=large]
A \arrow[r,"u"]&B\\
&L \arrow[u,"v"']
\end{tikzcd}
\]
where $u$ is an epimorphism and $L\in\Ob\cC_o$, can be completed to a commutative diagram
\[
\begin{tikzcd}[column sep=large,row sep=large]
A \arrow[r,"u"]&B\\
L_1 \arrow[u] \arrow[r,"w"']&L \arrow[u,"v"']
\end{tikzcd}
\]
such that $L_1\in\Ob\cC_o$ and $w$ is an epimorphism. Under these conditions, the canonical functors
\[
D_o^b(\cC)\longrightarrow D^b(\cC),\qquad
D_o^-(\cC)\longrightarrow D^-(\cC)
\]
are fully faithful. Moreover, the essential image of $D_o^b(\cC)$ (resp. $D_o^-(\cC)$) is stable under direct factors.
\end{statement}

The proof of this proposition essentially uses a criterion of Verdier ([2], 4--2 Theorem) and will be given in the first exposé of the seminar on Riemann-Roch (SGA 6 I 2.7).

\begin{statement}{Corollary}
Suppose, in addition, that the subcategory $\cC_o$ is stable under $\Ker$ and $\Coker$, hence is an abelian category, and denote by $D^b(\cC)_o$ the full triangulated subcategory of $D^b(\cC)$ generated by the complexes of $D^b(\cC)$ whose cohomology objects are isomorphic to objects of $\cC_o$. Then the inclusion functor
\[
D^b(\cC_o)\longrightarrow D^b(\cC)_o
\]
is an equivalence between the two categories, and consequently one obtains the sequence of isomorphisms
\[
K(\cC_o)\simeq K_o(\cC)\simeq K(D^b(\cC_o))\simeq K(D^b(\cC)_o).
\tag{*}
\]
\end{statement}

\emph{Observation.} Suppose that the subcategory $\cC_o$ satisfies conditions (A) and (B), and that every object of $\cC$ admits a finite resolution on the left by elements of $\cC_o$. In this case, $\cC_o$ also satisfies condition (C), and the canonical functor
\[
D_o^b(\cC)\longrightarrow D^b(\cC)
\]
is in fact an equivalence of categories. Consequently, one has the isomorphism
\[
K_{\cC_o}(\cC)\simeq K(\cC).
\]
\footnote{cf. (SGA 6 IV 1.6).}

\subsection*{5. Pseudo-coherent complexes}

Let $\cC$ be an abelian category and $\cC_o$ a full subcategory of $\cC$ satisfying conditions (A), (B), (C). We shall denote by $D_o^-(\cC)_{\coh}$ the essential image of $D_o^-(\cC)$ under the canonical functor
\[
D_o^-(\cC)\longrightarrow D(\cC).
\]
The objects of $D_o^-(\cC)_{\coh}$ will be called $\cC_o$-pseudo-coherent complexes. Thus a complex of $\cC$ is $\cC_o$-pseudo-coherent if it is isomorphic in $D(\cC)$ to a bounded-above complex whose components are objects of $\cC_o$.

The category $D_o^-(\cC)_{\coh}$ is a triangulated subcategory of $D(\cC)$, in the sense that every distinguished triangle of $D(\cC)$ with two vertices in the subcategory $D_o^-(\cC)_{\coh}$ is isomorphic to a triangle whose three objects are in the subcategory $D_o^-(\cC)_{\coh}$.

First, it is clear that the subcategory $D_o^-(\cC)_{\coh}$ is stable under the translation automorphism $T:D(\cC)\to D(\cC)$. Therefore, to prove that $D_o^-(\cC)_{\coh}$ is a triangulated subcategory of $D(\cC)$, it is enough to check that if $U^\bullet$, $V^\bullet$ are objects of $D_o^-(\cC)_{\coh}$ and if the triangle
\[
\begin{tikzcd}[column sep=small,row sep=small]
& W^\bullet \arrow[dl,"w"']\\
U^\bullet \arrow[rr,"u"']&& V^\bullet \arrow[ul,"v"']
\end{tikzcd}
\]
is distinguished for the triangulated category structure of $D(\cC)$, then $W^\bullet$ is also an object of $D_o^-(\cC)_{\coh}$. By hypothesis, there exist complexes $X^\bullet$, $Y^\bullet$, isomorphic in $D(\cC)$ to $U^\bullet$ and $V^\bullet$ respectively, whose components are objects of $\cC_o$. We obtain a commutative square in $D(\cC)$:
\[
\begin{tikzcd}[column sep=large,row sep=large]
U^\bullet \arrow[r,"u"] \arrow[d,"\sim"']&V^\bullet \arrow[d,"\sim"]\\
X^\bullet \arrow[r,"f"']&Y^\bullet .
\end{tikzcd}
\]
But, after replacing $X^\bullet$ by another isomorphic complex, we may suppose, using the fact that the inclusion functor
\[
D_o^-(\cC)\longrightarrow D^-(\cC)
\]
is fully faithful (Prop. 4.1), that $f$ is an actual morphism of complexes. Let $Z^\bullet$ be the mapping cylinder of $f$, and let $g:Y^\bullet\to Z^\bullet$ be the canonical morphism. By the construction itself, the components of $Z^\bullet$ are objects of $\cC_o$ (condition (A)). On the other hand,
\[
\begin{tikzcd}[column sep=small,row sep=small]
&Z^\bullet \arrow[dl,"h"']\\
X^\bullet \arrow[rr,"f"']&&Y^\bullet \arrow[ul,"g"']
\end{tikzcd}
\]
being a distinguished triangle in $D(\cC)$, it follows that $Z^\bullet$ is isomorphic to $W^\bullet$.

The category $D_o^-(\cC)_{\coh}$ is therefore endowed with a structure of triangulated category: a triangle in $D_o^-(\cC)_{\coh}$ is distinguished if and only if it is distinguished in the triangulated category $D(\cC)$.

The intersection of $D_o^-(\cC)_{\coh}$ with $D^b(\cC)_{\coh}$ will be denoted by $D_o^b(\cC)_{\coh}$.

An object $M$ of $\cC$ will be called pseudo-coherent if it is so as a complex reduced to degree zero; one proves that it is necessary and sufficient for this that it admit a resolution on the left whose components are objects of $\cC_o$.

\begin{statement}{Proposition 5.1}
A complex $X$ is $\cC_o$-pseudo-coherent if and only if it is isomorphic to a complex which is bounded on the right and whose components are $\cC_o$-pseudo-coherent.
\end{statement}

The proof of this proposition will appear in the first exposé of the seminar on Riemann-Roch (SGA 6 I 2.7).

\subsection*{6. Perfect complexes}

An object of the category $D(\cC)$ will be called a $\cC_o$-perfect complex if it is isomorphic to an object of $D_o^b(\cC)$, that is, to a bounded complex all of whose components are in $\cC_o$. We shall denote by $D_o(\cC)_{\parf}$ the full subcategory of $D(\cC)$ generated by the perfect complexes.

Thus the category $D_o(\cC)_{\parf}$ is nothing other than the essential image of $D_o^b(\cC)$ under the inclusion functor
\[
D_o^b(\cC)\longrightarrow D(\cC).
\]
It is easy to see that $D_o(\cC)_{\parf}$ is a triangulated subcategory of $D(\cC)$, and hence is endowed with a structure of triangulated category.

\subsection*{7. An important special case}

Let $A$ be a ring and let $\cC$ be the category of $A$-modules, on the left for example. Denote by $\cC_o$ the full subcategory of $\cC$ generated by the projective $A$-modules of finite type. The subcategory $\cC_o$ satisfies conditions (A), (B), (C). In this case we use the following notation:
\[
D^-(A)_{\coh}=D_o^-(\cC)_{\coh},\qquad
D^b(A)_{\coh}=D_o^b(\cC)_{\coh},
\]
\[
D(A)_{\parf}=D_o(\cC)_{\parf}.
\]
The complexes of $D(\cC)$ which are $\cC_o$-pseudo-coherent (resp. $\cC_o$-perfect) will simply be called pseudo-coherent (resp. perfect) complexes.

It will be useful to consider the following groups:
\[
K_\bullet(A)=K(D^b(A)_{\coh}),
\]
\[
K^\bullet(A)=K(D(A)_{\parf}).
\]
As for the group $K^\bullet(A)$, we have, by Proposition 4.1 and its corollary,
\[
K^\bullet(A)=K(D(A)_{\parf})\simeq K(D_o^b(\cC))\simeq K_{\cC_o}(\cC),
\]
which shows that it is also the classical group of classes of projective modules of finite type over $A$.

In the special case where $A$ is noetherian, the group $K_\bullet(A)$ also admits a classical description. For this, let $\cC_1$ be the abelian category of modules of finite type. Then
\[
K_\bullet(A)\simeq K(\cC_1).
\]
Indeed, using Proposition 4.1, one has
\[
K(\cC_1)\simeq K(D^b(\cC_1))\simeq K(D^b_{\cC_1}(\cC)),
\]
because $D^b(\cC_1)\simeq D^b_{\cC_1}(\cC)$ immediately.

In fact, if $A$ is noetherian and $L^\bullet\in D^-(A)$, then it follows from Proposition 5.1 that $L^\bullet$ is pseudo-coherent if and only if the $H^i(L^\bullet)$ are of finite type.

For an arbitrary ring $A$, one has a homomorphism
\[
K^\bullet(A)\xrightarrow{\widetilde\gamma_A}K_\bullet(A),
\]
which comes from the inclusion functor
\[
(*)\qquad D(A)_{\parf}\longrightarrow D^b(A)_{\coh}.
\]
In general, the homomorphism $\widetilde\gamma_A$ is neither injective nor surjective. In the special case of a regular ring $A$, every $A$-module $M$ of finite type has finite projective dimension and it follows, thanks to 5.1, that the inclusion functor $(*)$ is in fact an equivalence between the two categories. In that case, $\widetilde\gamma_A$ is therefore an isomorphism.

\subsection*{8. Tor of complexes}

Let $A$ be a not necessarily commutative ring, and let ${}_A\cC$ (resp. $\cC_A$) be the abelian category of left (resp. right) $A$-modules. By $K^-({}_A\cC)$ we shall denote the triangulated category of complexes of ${}_A\cC$ bounded above, and by ${}_A\cD$ the full triangulated subcategory of $K^-({}_A\cC)$ consisting of complexes whose components are flat left $A$-modules.

Let $X^\bullet$ be an object of $K^-(\cC_A)$, and consider the functor
\[
\Phi_{X^\bullet}=X^\bullet\otimes_A\,\cdot\;:\;K^-({}_A\cC)\longrightarrow K^-(\Ab).
\]
Under these conditions, the following propositions hold.

1) Every acyclic object of ${}_A\cD$ is transformed by $\Phi_{X^\bullet}$ into an acyclic object of $K^-(\Ab)$. Indeed, if $N^\bullet$ is an object of ${}_A\cD$ such that $Z^i(N^\bullet)$, $B^i(N^\bullet)$, and $H^i(N^\bullet)$ are flat, one has the classical isomorphism
\[
\bigoplus_{p+q=n}H_p(X^\bullet)\otimes H_q(N^\bullet)
\longrightarrow H_n(X^\bullet\otimes_A N^\bullet).
\]

2) Every object $Y^\bullet$ of $K^-({}_A\cC)$ is the target of a quasi-isomorphism whose source is an object of ${}_A\cD$. It is enough for this to take a Cartan-Eilenberg resolution. We can therefore apply Theorem 2 of §2, no. 2 of [2] to conclude that $\Phi_{X^\bullet}$ admits a left derived functor:
\[
\mathbf L\Phi_{X^\bullet}:D^-({}_A\cC)\longrightarrow D^-(\Ab).
\]
The object $\mathbf L\Phi_{X^\bullet}(Y^\bullet)$ may be obtained, up to isomorphism, as follows. Let $P^\bullet$ be an object of ${}_A\cD$ and let $P^\bullet\to Y^\bullet$ be a quasi-isomorphism. Then $\mathbf L\Phi_{X^\bullet}(Y^\bullet)$ is isomorphic to
\[
Q(X^\bullet\otimes_A P^\bullet),
\]
where, as usual, $Q$ denotes the canonical functor
\[
K^-(\Ab)\longrightarrow D^-(\Ab).
\]
It follows that for each object $Y^\bullet$ of $D^-({}_A\cC)$, the functor
\[
K^-(\cC_A)\longrightarrow D^-(\Ab),\qquad
X^\bullet\longmapsto \mathbf L\Phi_{X^\bullet}(Y^\bullet)
\]
is exact and vanishes on acyclic complexes, hence gives an exact bifunctor
\[
D^-(\cC_A)\times D^-({}_A\cC)\longrightarrow D^-(\Ab),
\]
\[
(X^\bullet,Y^\bullet)\longmapsto X^\bullet\overset{\mathbf L}{\otimes}_A Y^\bullet.
\]
Passing to homology, one defines the $\operatorname{Tor}_i^A(X^\bullet,Y^\bullet)$ (cf. also [1], 6.3). In the particular case of a commutative ring $A$, one obtains in fact an exact bifunctor
\[
D^-(\cC)\times D^-(\cC)\longrightarrow D^-(\cC),
\]
the tensor product in the sense of derived categories ($\cC={}_A\cC=\cC_A$).

\begin{statement}{Definition}
We shall say that the complex $X$ of $D(\cC_A)$ has finite Tor-dimension if it lies in $D^-$ and if there exists $n_o\in\mathbb Z^+$ such that
\[
\operatorname{Tor}_n^A(X,Y)=0
\]
for every $Y\in\Ob({}_A\cC)$ and every $n>n_o$.
\end{statement}

We shall use the notion of Tor-dimension to obtain a characterization of perfect complexes.

\begin{statement}{Proposition 8.1}
Let $A$ be a not necessarily commutative ring and let $L^\bullet$ be a complex of $A$-modules. The following conditions are equivalent:
\[
\text{a) } L^\bullet\text{ is perfect;}
\]
\[
\text{b) } L^\bullet\text{ is pseudo-coherent, has bounded cohomology, and has finite Tor-dimension.}\footnote{cf. SGA 6 I 5.8.1.}
\]
\end{statement}

For the proof, one uses the following lemma.

\begin{statement}{Lemma}
(Bourbaki, Alg. Comm. I, §2, ex. 15). Every $A$-module $E$ which is flat and of finite presentation is projective.
\end{statement}

\emph{Proof of the lemma.} Let $u:F\to F'$ be an epimorphism. We must prove that
\[
\Hom_A(E,F)\longrightarrow\Hom_A(E,F')
\]
is an epimorphism. It is enough to prove that, for every divisible group $G$,
\[
\Hom_{\mathbb Z}(\Hom_A(E,F'),G)\longrightarrow
\Hom_{\mathbb Z}(\Hom_A(E,F),G)
\]
is a monomorphism. But the diagram
\[
\begin{tikzcd}[column sep=large,row sep=large]
\Hom_{\mathbb Z}(\Hom_A(E,F'),G)
 \arrow[r]\arrow[d,"\sim"']&
\Hom_{\mathbb Z}(\Hom_A(E,F),G)\arrow[d,"\sim"]\\
\Hom_{\mathbb Z}(F',G)\otimes_A E
 \arrow[r]&
\Hom_{\mathbb Z}(F,G)\otimes_A E
\end{tikzcd}
\]
is commutative, and the vertical arrows are isomorphisms, by the hypothesis on $E$.

\emph{Proof of Proposition 5.2.} The implication a) $\Rightarrow$ b) is immediate. To prove b) $\Rightarrow$ a), let
\[
\cdots\longrightarrow P_n\xrightarrow{d_n}\cdots\longrightarrow P_i\longrightarrow 0\longrightarrow\cdots
\]
be a complex isomorphic to $L^\bullet$ in $D(\cC)$, where the $P_j$ are projective $A$-modules of finite type. Since $L^\bullet$ has finite Tor-dimension and bounded cohomology, it follows that there exists $n_o\in\mathbb Z$ such that:

I) The sequence
\[
\cdots\longrightarrow P_{n_o-p}\xrightarrow{d_{n_o-p}}\cdots\longrightarrow
P_{n_o}\longrightarrow \Ker d_{n_o+1}\longrightarrow0
\]
is exact.

II) $\Ker d_{n_o+1}$ is flat.

III) The complex
\[
0\longrightarrow \Ker d_{n_o+1}\longrightarrow P_{n_o+1}\longrightarrow\cdots\longrightarrow P_i\longrightarrow0\longrightarrow\cdots
\]
is isomorphic to $L^\bullet$ in $D(\cC)$.

But $\Ker d_{n_o+1}$ is of finite presentation, hence projective of finite type by the lemma, which proves 8.1.

If $A$ is a commutative ring, and if $X^\bullet$, $Y^\bullet$ are two perfect complexes, then their tensor product
\[
X^\bullet\overset{\mathbf L}{\otimes}_A Y^\bullet
\]
in the derived category is again a perfect complex. We obtain a bifunctor
\[
D(A)_{\parf}\times D(A)_{\parf}\longrightarrow D(A)_{\parf}
\]
which is exact in each variable. Consequently, one defines a product on $K^\bullet(A)$, and it is easy to see that $K^\bullet(A)$ is thereby endowed with the structure of a commutative ring with unit. In fact, $K^\bullet(A)$ can be endowed with the structure of a $\lambda$-ring; for details, see the seminar on Riemann-Roch following this one.\footnote{SGA 6 V 2.}

The tensor product
\[
X^\bullet\overset{\mathbf L}{\otimes}_A Y^\bullet
\]
of a perfect complex $X^\bullet$ and a bounded pseudo-coherent complex $Y^\bullet$ is a bounded pseudo-coherent complex ($A$ still a commutative ring). Consequently, one can introduce on $K_\bullet(A)$ the structure of a $K^\bullet(A)$-module.

\subsection*{9. Functorial properties}

a) Let $A$ and $B$ be two not necessarily commutative rings, and let $u:A\to B$ be a homomorphism. Every left (resp. right) $B$-module $M$ can be endowed with the structure of a left (resp. right) $A$-module, denoted $M_{[u]}$. In particular, $B_{[u]}$ is a right $A$-module.

The functor
\[
E\longmapsto B_{[u]}\otimes_A E
\]
allows one to define the exact functor
\[
D(A)_{\parf}\longrightarrow D(B)_{\parf}
\]
and consequently the homomorphism
\[
K^\bullet(A)\xrightarrow{u^*}K^\bullet(B).
\]
Thus one obtains a covariant functor from the category of rings (resp. commutative rings) to the category of abelian groups (resp. commutative rings).

b) Suppose now that the left $A$-module $B_{[u]}$ is pseudo-coherent. It follows that the functor
\[
F\longmapsto F_{[u]}
\]
transforms a free $B$-module of finite type into a pseudo-coherent $A$-module. By using Proposition 5.1, one obtains an exact functor
\[
D^b(B)_{\coh}\longrightarrow D^b(A)_{\coh}
\]
and consequently a canonical homomorphism
\[
K_\bullet(B)\xrightarrow{u_*}K_\bullet(A).
\]
The homomorphisms $u_*$ and $u^*$ satisfy the following identity (``the projection formula''):
\[
u_*(u^*(\alpha)\cdot \beta)=\alpha\,(u_*(\beta))
\quad\text{for }\alpha\in K^\bullet(A),\ \beta\in K_\bullet(B).
\]
This easy formula will be proved in a more general form in the exposés of the seminar on Riemann-Roch (SGA 6 IV 2.11.1.2).

\subsection*{10. Relative construction\footnote{cf. (SGA 6 III, IV).}}

Let $u:A\to B$ be a homomorphism of rings, and let
\[
D(B)\longrightarrow D(A)
\]
be the associated exact functor. In $D(B)$, we can consider the complexes that are pseudo-coherent relative to $A$, those of finite Tor-dimension relative to $A$, and also the complexes perfect relative to $A$. In particular, it is natural to consider the triangulated subcategories
\[
D(B)_{A-\parf},\qquad D^b(B)_{A-\coh}
\]
of $D(B)$ generated by the complexes that are perfect, respectively pseudo-coherent and bounded, relative to $A$. These two triangulated categories allow one to define the following groups:
\[
K^\bullet(B\operatorname{\,rel\,}A)=K(D(B)_{A-\parf}),
\]
\[
K_\bullet(B\operatorname{\,rel\,}A)=K(D^b(B)_{A-\coh}).
\]

\emph{Examples.} Let $G$ be a group and let $\Lambda$ be a commutative ring. The homomorphism
\[
\Lambda\xrightarrow{\varepsilon}\Lambda[G]
\]
therefore allows one to define the groups
\[
K^\bullet(\Lambda[G]\operatorname{\,rel\,}\Lambda),\qquad
K_\bullet(\Lambda[G]\operatorname{\,rel\,}\Lambda),
\]
which will be denoted by
\[
R^\bullet_\Lambda(G),\qquad R^\Lambda_\bullet(G).
\]
In fact, $R^\bullet_\Lambda(G)$ is a commutative ring and $R^\Lambda_\bullet(G)$ is a module over it.

\emph{Remarks.}

a) If the ring $\Lambda$ is not necessarily commutative, one may set, by definition,
\[
\Lambda[G]=\Lambda\otimes_{\mathbb Z}\mathbb ZG
\]
and consequently obtains a homomorphism
\[
\Lambda\xrightarrow{\varepsilon}\Lambda[G].
\]
In this way, one may give a meaning to the groups $R^\bullet_\Lambda(G)$ and $R^\Lambda_\bullet(G)$ even when $\Lambda$ is non-commutative.

b) The preceding construction can be generalized to the case of a group prescheme, or of a pro-prescheme in groups, $G$ over $\Spec\Lambda$, by starting from the abelian category of $\Lambda$-modules on which $G$ operates; this is no longer made explicit in terms of modules over a ring. One then uses the natural inclusion functor into the category of $\Lambda$-modules. If $\Lambda$ is a field, it can be made explicit in terms of modules over a pro-ring, and in the general case, in terms of comodules, at least if $G$ is affine over $\Lambda$; cf. SGA 3 I 4.7.2).

\subsection*{Bibliography}

\begin{enumerate}
\item[{[1]}] A. Grothendieck and J. Dieudonné, \emph{Elements de Géométrie Algébrique}, Publ. Math. I.H.E.S. IV (Second Part).
\item[{[2]}] J. L. Verdier, \emph{Catégories Dérivées}, I.H.E.S., 1963.\footnote{Published in SGA $4^{1/2}$.}
\end{enumerate}

\phantomsection\addcontentsline{toc}{section}{Exposé X --- Euler--Poincaré Formula in Étale Cohomology}
\section*{Exposé X --- Euler--Poincaré Formula in Étale Cohomology}

\begin{center}
\textbf{EULER--POINCARÉ FORMULA IN ÉTALE COHOMOLOGY}

by A. Grothendieck

(written up by I. Bucur)
\end{center}

The purpose of the present exposé is to state and prove formula (7.2), of Euler--Poincaré type. This formula generalizes a formula obtained independently by Ogg and Chafarevitch, whose results are presented (from the viewpoint of étale cohomology) in M. Raynaud's exposé [1].

The Euler--Poincaré formula uses in an essential way the formalism of derived categories. As in loc. cit. the proof is made through a variant (5.1) of a well-known formula of Weil.

\subsection*{1. Sheaves on a scheme with operators}

Let $X$ be a scheme endowed with a group $G$ acting by automorphisms. Proceeding as in ([T] chap. V), one may define the notion of an étale sheaf on $X$ with group of operators $G$, or simply a “sheaf on $(X_{\et},G)$”.

By definition, to give an étale sheaf on $X$ with group of operators $G$ is to give a sheaf $F\in \Ob(\widetilde X_{\et})$ and a morphism
\[
S_g:F\longrightarrow g^*(F)
\]
for each $g\in G$, in such a way that the following conditions are satisfied:
\[
\text{a) }\quad S_e=1_F,
\]
and, for $g,h\in G$, the diagram
\[
\begin{tikzcd}[column sep=large,row sep=large]
F \arrow[r,"S_{hg}"] \arrow[d,"S_h"'] & (hg)^*(F) \arrow[d,"C_{h,g}(F)"]\\
h^*(F) \arrow[r,"h^*(S_g)"'] & h^*(g^*(F))
\end{tikzcd}
\]
is commutative, where
\[
C_{h,g}(F):(hg)^*(F)\longrightarrow h^*(g^*(F))
\]
is the canonical isomorphism.

Morphisms between two sheaves $(F,S_g)$, $(F',S'_g)$ on $X$ with group of operators $G$ are defined in the evident way, and one obtains a category $\widetilde X_{\et,G}$ which is a topos, as follows readily from Giraud's criterion (SGA 4 II 4.12).

Given a ring $A$ of $\widetilde X_{\et,G}$, one may consider the abelian category of $A$-modules, which will be denoted $\Mod(A)$.

For example, given a ring $\Lambda$, one may equip the constant sheaf $\Lambda_X$ associated with $\Lambda$ with the trivial structure of a sheaf with group of operators $G$, using the canonical transitivity isomorphisms
\[
S_g:\Lambda_X\xrightarrow{\sim}g^*(\Lambda_X).
\]
All the results proved in [1], chap. V, can be extended to this setting, and we leave the details to the reader.

\subsection*{2. Where one proves that a certain complex of $\Lambda[G]$-modules is perfect}

Recall the following result of SGA 4 XVII.

\begin{statement}{Lemma 2.1}
Let $X$ and $Y$ be two schemes and let $f:X\to Y$ be a proper morphism, with $Y$ quasi-compact. Let $A$ be a torsion ring, and let $F^\bullet$ be a complex of sheaves of $A$-modules on $X$ of finite Tor-dimension (SGA 4 XVII). Then $Rf_*(F^\bullet)$ is of finite Tor-dimension.
\end{statement}

\begin{statement}{Proposition 2.2}
Let $X$ be a scheme proper over a separably closed field $k$, let $G$ be a finite group acting on $X$, and let $V$ be an open subset of $X$ stable under $G$ and such that $G$ acts admissibly (SGA 1 V 1.7) and freely on $V$. Let $\Lambda=\mathbb Z/\ell^n\mathbb Z$, with $\ell$ a prime number different from the characteristic of $k$, and let $\Lambda_V$ be the constant sheaf on $V$ associated with $\Lambda$, endowed with its trivial structure of a sheaf with group of operators $G$. Let
\[
i:V\longrightarrow X
\]
be the immersion of $V$ into $X$, and put $\Lambda_{V,X}=i_!(\Lambda_V)$. Under these conditions the complex of $\Lambda[G]$-modules
\[
R\Gamma_X(\Lambda_{V,X})
\]
is a perfect complex (VIII 7).
\end{statement}

\begin{prooftext}
If $p:X\to Y=X/G$ is the canonical projection, it is well known that there is a canonical isomorphism
\[
R\Gamma_X(\Lambda_{V,X})\simeq R\Gamma_Y(Rp_*(\Lambda_{V,X}))
\]
in the category $D(\Lambda[G])$. On the other hand, the complex $Rp_*(\Lambda_{V,X})$ of $\Lambda[G]$-modules on $Y$ is of finite Tor-dimension. Indeed, since $p:X\to Y$ is a finite morphism of schemes, one has
\[
R^q p_*(\Lambda_{V,X})=0\qquad(q\ne0)
\]
(SGA 4 VIII 5.5). Hence
\[
Rp_*(\Lambda_{V,X})=p_*(\Lambda_{V,X}).
\]
SGA 4 VIII 5.5 also gives the possibility of determining the fibres of the sheaf $p_*(\Lambda_{V,X})$, and, taking account of the fact that $G$ acts freely on $V$, one finds
\[
\bigl(p_*(\Lambda_{V,X})\bigr)_{\bar y}\simeq
\begin{cases}
\displaystyle \prod_{x\in p^{-1}(y)}\Lambda_{V,X,\bar x}\simeq \Lambda[G],& y\in p(V),\\[1em]
0,& y\notin p(V),
\end{cases}
\]
where $\bar y$ is the geometric point of $Y$ localized at $y$.

Consequently $p_*(\Lambda_{V,X})$ is flat over $\Lambda[G]$, fibre by fibre (SGA 4 XVII), and therefore is of finite Tor-dimension. Applying 2.1 to the canonical morphism $Y\to\Spec(k)$, we deduce that
\[
R\Gamma_X(\Lambda_{V,X})=R\Gamma_Y(Rp_*(\Lambda_{V,X}))
\]
is of finite Tor-dimension.

To conclude that $R\Gamma_X(\Lambda_{V,X})$ is perfect, it remains to show that it is pseudo-coherent and has bounded cohomology (VIII 8.1). Since $\Lambda[G]$ is noetherian, to verify that $R\Gamma_X(\Lambda_{V,X})$ is pseudo-coherent it is enough to show that the groups $H^q(X,\Lambda_{V,X})$ are of finite type over $\Lambda[G]$ (VIII 7.2), or equivalently over $\Lambda$; this follows from (SGA 4 XIV 1.1), applied to the $\Lambda$-module $\Lambda_{V,X}$.
\end{prooftext}

The Weil formula 6.1 will give an expression for the class
$\operatorname{cl}_{K^\bullet(\Lambda[G])}(R\Gamma_X(\Lambda_{V,X}))$,
when $X$ is a curve, in terms of known global invariants and local invariants
that we shall define in no.~5.

\subsection*{3. Reminders on linear representations of finite groups}

In this number, $\Lambda$ denotes a commutative ring, and $G$ a group.

\subsubsection*{3.1. Characters}

Consider the full subcategory $\mathcal E$ of $\Mod(\Lambda[G])$ consisting of $\Lambda[G]$-modules $M$ whose underlying $\Lambda$-module is projective of finite type. For such a module $M$ and for $g\in G$, consider the endomorphism, for the $\Lambda$-module structure,
\[
g_M:M\longrightarrow M,
\qquad g_M(x)=gx.
\]
One obtains a function
\[
\chi_M:G\longrightarrow \Lambda,
\qquad \chi_M(g)=\Tr^*_\Lambda(g),
\]
which is a central function on $G$ and is called the character of the $\Lambda[G]$-module $M$.

Associating to each $M\in\Ob(\mathcal E)$ its character $\chi_M$ gives an additive function (VIII 1). Using the formalism of traces for endomorphisms of perfect complexes, one even obtains an additive function $\Tr$ on the category of complexes of $\Lambda[G]$-modules which are perfect as complexes of $\Lambda$-modules. Hence an homomorphism
\[
\Tr:R^{\bullet}(G)\longrightarrow \mathcal F(G,\Lambda),
\]
where $R^{\bullet}_{\Lambda}(G)$ is the ring of representations of $G$ over $\Lambda$ (VIII 10), and $\mathcal F(G,\Lambda)$ is the module of central functions $G\to\Lambda$.

When $G$ is finite, the composite of the canonical homomorphism
\[
K^{\bullet}(\Lambda[G])\longrightarrow R^{\bullet}_{\Lambda}(G)
\]
with $\Tr$ will be denoted by $\tr$.

Recall the following result (IX th. 1, cor. 3).

\begin{statement}{Proposition 3.2}
Let $\mathbb Z_\ell$ be the ring of $\ell$-adic integers and let $G$ be a finite group. Then the homomorphism
\[
\tr:K^{\bullet}(\mathbb Z_\ell[G])\longrightarrow \mathcal F(G,\mathbb Z_\ell)
\]
is a monomorphism. Two projective $\mathbb Z_\ell[G]$-modules having the same trace function in $\mathcal F(G,\mathbb Z_\ell)$ are isomorphic.
\end{statement}

\subsubsection*{3.3. Restriction and induction}

Let $\alpha:G\to G'$ be a homomorphism from a group $G$ to a group $G'$. By restriction of operators one obtains an exact functor
\[
J:\Mod(\Lambda[G'])\longrightarrow\Mod(\Lambda[G]).
\]
This functor has a left adjoint
\[
I:\Mod(\Lambda[G])\longrightarrow\Mod(\Lambda[G']),
\]
namely
\[
I(F)=\Lambda[G']\otimes_{\Lambda[G]}F.
\]
If $F$ is projective of finite type, respectively of finite type, then so is $I(F)$.

If $\varphi\in\mathcal F(G',\Lambda)$ is a central function on $G'$, the function $\varphi\circ\alpha$ is a central function on $G$, and therefore one obtains an application
\[
\alpha^*: \mathcal F(G',\Lambda)\longrightarrow\mathcal F(G,\Lambda).
\]
If $\varphi$ is the character of a $\Lambda[G']$-module $M$ in $\mathcal C(G')$, then $\alpha^*(\varphi)$ is the character of the $\Lambda[G]$-module $J(M)$.

Now suppose that $\alpha:G\to G'$ is injective and that $G$ has finite index in $G'$. If $F$ is a $\Lambda[G]$-module in $\mathcal C(G)$ whose character is $\psi:G\to\Lambda$, then $I(F)$ lies in $\mathcal C(G')$ and the character $\psi^*:G'\to\Lambda$ of $I(F)$ is the character induced by $\psi$, i.e. the function on $G'$ deduced from $\psi$ by the formula
\[
\psi^*(g)=\sum_{t\in G'/G}\psi(t^{-1}gt),
\tag{3.3.1}
\]
with the convention that $\psi(t^{-1}gt)=0$ if $t^{-1}gt\notin G$.

\subsubsection*{3.4. Additional structures on complexes}

Let $X^\bullet$ and $Y^\bullet$ be two complexes of $A$-modules. As in the case of modules, if one of the complexes $X^\bullet$, $Y^\bullet$ has an additional structure of an $A$-$B$-bimodule, it is possible to introduce on $\Hom_A^\bullet(X^\bullet,Y^\bullet)$ a structure of complex of $B$-modules, and consequently a corresponding structure on $\RHom_A(X^\bullet,Y^\bullet)$.

For example, suppose that $X^\bullet$ is an object of $D^-(A)$ and that $Y^\bullet$ is a complex of $A$-$B$-bimodules; to be precise, suppose that we are in the case described by the symbol ${}_{A-B}Y^\bullet$. In particular, $Y^\bullet$ is an object of $D(A)$. Let $P^\bullet\to X^\bullet$ be a projective resolution of $X^\bullet$. Then $\Hom_A^\bullet(P^\bullet,Y^\bullet)$ is evidently a complex of left $B$-modules, i.e. an object of $D(B)$. If another projective resolution $P_1^\bullet\to X^\bullet$ is used, the complex of $B$-modules $\Hom_A^\bullet(P_1^\bullet,Y^\bullet)$ is canonically isomorphic in $D(B)$ to $\Hom_A^\bullet(P^\bullet,Y^\bullet)$. It follows that one may define $\RHom_A(X^\bullet,Y^\bullet)$ as an object of $D(B)$, determined up to unique isomorphism, and functorial in $X^\bullet$ in $D^-(A)$ and in $Y^\bullet$ in $D(A\otimes_{\mathbb Z}B)$.

Analogous considerations apply to the complex $X^\bullet\otimes_A^{\mathbf L}Y^\bullet$. For example, if $Y^\bullet$ is a complex of $A$-$B$-bimodules as above and $X^\bullet$ is a right-bounded complex of right $A$-modules, one may define $X^\bullet\otimes_A^{\mathbf L}Y^\bullet$ as an object of $D(B)$ determined up to unique isomorphism.

\subsubsection*{3.5}

Let $\Lambda$ be a commutative ring, let $A$ and $B$ be two algebras over $\Lambda$, and put $C=A\otimes_\Lambda B$. Thus $C$-modules are $A$-$B$-bimodules.

If $\Delta$ is a full triangulated subcategory of $D(A)$, denote by $D(C)_{\Delta}$ the full triangulated subcategory of $D(C)$ generated by the complexes which, when regarded as objects of $D(A)$, belong to $\Delta$.

\begin{statement}{Proposition 3.6}
With the notation of 3.4 and 3.5, let $\Delta$ be a full triangulated subcategory of $D(A)$ stable under direct factors. If $P^\bullet$ is an object of $D(B)_{\parf}$ (VIII 7) and $M^\bullet$ is an object of $D(C)_{\Delta}$, then the complex of $A$-modules $\RHom_B(P^\bullet,M^\bullet)\in D(A)$ defined in 4.1 is an object of $\Delta$. In particular, there is a pairing of triangulated categories
\[
\RHom:D(B)_{\parf}\times D(C)_{\Delta}\longrightarrow \Delta,
\tag{3.6.1}
\]
that is, a bifunctor exact in each variable, and consequently (VIII 3) a bilinear map
\[
K^{\bullet}(B)\times K^{\bullet}(D(C)_{\Delta})\xrightarrow{\operatorname{hom}} K(\Delta)
\tag{3.6.2}
\]
such that
\[
K(a)\bigl(\cl_{K^{\bullet}(B)}(P^\bullet),\cl_{K(D(C)_{\Delta})}(M^\bullet)\bigr)
=\cl_{K(\Delta)}\bigl(\RHom_B(P^\bullet,M^\bullet)\bigr).
\tag{3.6.3}
\]
One obtains an analogous statement on replacing $B$ by its opposite $B^\circ$ and $\RHom_B(P^\bullet,M^\bullet)$ by $P^\bullet\otimes_B^{\mathbf L}M^\bullet$.
\end{statement}

\begin{prooftext}[Proof]
First one proves that $\RHom_B(P^\bullet,M^\bullet)$ is an object of $\Delta$ in the special case where $P^\bullet$ has a single nonzero component, and then argues by induction on the number of nonzero components of $P^\bullet$.
\end{prooftext}

\subsubsection*{3.7}

Let $P$ be a $\Lambda[G]$-module. It is in particular a $\Lambda$-module. The $\Lambda$-module $\check P=\Hom_\Lambda(P,\Lambda)$ can be endowed with a left $\Lambda[G]$-module structure by putting, for every $u\in\Hom_\Lambda(P,\Lambda)$ and $g\in G$, $(gu)(x)=u(g^{-1}x)$, i.e.
\[
g_{\check P}={}^{t}(g^{-1})_P.
\tag{3.7.1}
\]
The $\Lambda[G]$-module $\check P$ thus obtained is called the dual of $P$.

One may therefore define the dual $\check P^\bullet=\Hom_\Lambda^\bullet(P^\bullet,\Lambda)$ of a complex $P^\bullet$ of $\Lambda[G]$-modules as a complex of $\Lambda[G]$-modules.

If $\chi_P$ is the character of $P$, the character $\chi_{\check P}$ of its dual is given by
\[
\chi_{\check P}(g)=\chi_P(g^{-1}),
\tag{3.7.2}
\]
which follows immediately from (4.3.1), using the formula
\[
\Tr_\Lambda^*(v)=\Tr_\Lambda^*({}^{t}v).
\]
If $P$ is a projective $\Lambda[G]$-module of finite type, it is clear that $\check P$ is also a projective $\Lambda[G]$-module of finite type. Moreover, in this case the canonical homomorphism of $\Lambda$-modules
\[
P\longrightarrow\check{\check P}
\]
is in fact an isomorphism of $\Lambda[G]$-modules. In this way one defines on $K^{\bullet}(\Lambda[G])$ an involution
\[
x\longmapsto \check x,
\]
which associates with the class of $P$ the class of $\check P$. It is immediate that if $\mu:\Lambda\to\Gamma$ is a homomorphism of commutative rings and $u:\Lambda[G]\to\Gamma[G]$ the associated homomorphism, then the homomorphism
\[
u^*:K^{\bullet}(\Lambda[G])\longrightarrow K^{\bullet}(\Gamma[G])
\]
commutes with the involutions $x\longmapsto\check x$ on $K^{\bullet}(\Lambda[G])$ and $K^{\bullet}(\Gamma[G])$.

Let $A$ be an algebra over $\Lambda$, and put
\[
A[G]=\Lambda[G]\otimes_\Lambda A.
\]
Given a complex $X^\bullet$ of $A[G]$-modules, denote by $X^{\bullet G}$ the complex of $A$-modules which is the image of $X^\bullet$ under the functor that associates to each $A[G]$-module $X$ the $A$-module formed by the elements of $X$ invariant under $G$.

\begin{statement}{Proposition 3.8}
Let $G$ be a finite group, let $P^\bullet$ be a bounded complex of $\Lambda[G]$-modules all of whose components are projective of finite type, and let $M^\bullet$ be a complex of $A[G]$-modules. Under these conditions there are canonical isomorphisms of complexes of $A$-modules
\[
\Hom_{\Lambda[G]}^\bullet(P^\bullet,M^\bullet)
\simeq \Hom_\Lambda^\bullet(P^\bullet,M^\bullet)^G
\simeq (\check P^\bullet\otimes_\Lambda M^\bullet)^G
\simeq \check P^\bullet\otimes_{\Lambda[G]}M^\bullet.
\tag{3.8.1}
\]
\end{statement}

The proof uses the well-known isomorphisms in the case of modules and is left to the reader. Note that, in the second and third terms, the operations of $G$ on $\Hom_\Lambda^\bullet(P^\bullet,M^\bullet)$, respectively on $\check P^\bullet\otimes_\Lambda M^\bullet$, are defined by transport of structure from its operations on $P^\bullet$ and $M^\bullet$, respectively $\check P^\bullet$ and $\check M^\bullet$.

\begin{statement}{Corollary 3.9}
If $P^\bullet\in\Ob(D(\Lambda[G])_{\parf})$ and $M^\bullet\in D^+(A[G])$, then in the derived category $D(A)$ one has the following isomorphisms:
\[
\RHom_{\Lambda[G]}(P^\bullet,M^\bullet)
=R\Gamma^G(\check P^\bullet\otimes_\Lambda^{\mathbf L}M^\bullet)
\simeq \check P^\bullet\otimes_{\Lambda[G]}M^\bullet.
\tag{3.9.1}
\]
\end{statement}

Indeed, since $P^\bullet\in D(\Lambda[G])_{\parf}$, we may suppose that $P^\bullet$ is bounded and that all its components are projective of finite type. Hence
\[
\RHom_{\Lambda[G]}(P^\bullet,M^\bullet)=\Hom_{\Lambda[G]}^\bullet(P^\bullet,M^\bullet)
\]
and one can use the isomorphisms (4.3.4). But $\check P^\bullet$ is, like $P^\bullet$, bounded with projective finite-type components, hence
\[
\check P^\bullet\otimes_{\Lambda[G]}M^\bullet\simeq
\check P^\bullet\otimes_{\Lambda[G]}^{\mathbf L}M^\bullet.
\]
Similarly, $\check P^\bullet\otimes_\Lambda M^\bullet=\check P^\bullet\otimes_\Lambda^{\mathbf L}M^\bullet$. Moreover the components of the complex $\check P^\bullet\otimes_\Lambda M^\bullet$ are cohomologically trivial ([3] IX 3), and therefore
\[
R\Gamma^G(\check P^\bullet\otimes_\Lambda^{\mathbf L}M^\bullet)
\simeq (\check P^\bullet\otimes_\Lambda M^\bullet)^G.
\]

\subsection*{4. The Swan representation}

Let $E$ be a smooth connected algebraic curve over the algebraically closed field $k$. If $x\in E$ is a closed point of $E$, its local ring $\mathcal O_{E,x}$, or simply $\mathcal O_x$, is a discrete valuation ring whose field of fractions is isomorphic to the field of rational functions on $E$ and whose residue field is $k$. The corresponding valuation will be denoted $v_x$.

\begin{statement}{Lemma 4.1}
Let $E$ be a smooth algebraic curve over the algebraically closed field $k$, and let $g:E\to E$ be an endomorphism of $E$ different from the identity. Suppose that $x\in E$ is an isolated fixed point of $g$, and let $\pi$ be a uniformizer of $\mathcal O_{E,x}$. If $v_x$ is the corresponding valuation of $\mathcal O_{E,x}$, then the multiplicity $v_x(g)$ of the fixed point $x$ of the map $g$ (IV) is equal to $v_x(g(\pi)-\pi)$.
\end{statement}


In the product $E\times E$, the diagonal and the graph of the endomorphism $g$ will be denoted respectively by $\Delta$ and $\Gamma_g$. Let
\[
\gamma:E\longrightarrow E\times E,\qquad
\delta:E\longrightarrow E\times E
\]
be the corresponding immersions. On the completed local rings at $x$ and at $(x,x)$, they are given by homomorphisms
\[
\gamma^*,\delta^*:A\widehat\otimes A\longrightarrow A
\qquad (A=\widehat{\mathcal O}_x),
\]
where
\[
\delta^*(a\otimes b)=ab,\qquad
\gamma^*(a\otimes b)=a\,g^*(b).
\]
We have to evaluate the multiplicity $v_x(g)$ of the point $\delta(x)$ in the intersection $\Delta\Gamma_g$ on $E\times E$. The divisor $\Delta$ corresponds to the ideal generated by the element $1\otimes\pi-\pi\otimes1$ of $A\widehat\otimes A$. Hence the multiplicity of $\delta(x)$ in the intersection $\Delta\Gamma_g$ is
\[
v_x(g)=v_x\bigl(\gamma^*(1\otimes\pi-\pi\otimes1)\bigr)
      =v_x(g(\pi)-\pi),
\]
as required.

\subsubsection*{4.2}

Let $C$ be a smooth, proper, connected algebraic curve over the algebraically closed field $k$. If $\eta$ is its generic point, let $K=k(\eta)$ be its field of rational functions.

Let
\[
K'\supset K,\qquad [K':K]=n
\]
be a finite Galois extension of degree $n$ with Galois group $G$. Let $C'$ be the normalization of $C$ in $K'$, and let
\[
p:C'\longrightarrow C
\]
be the canonical projection. The Galois group $G$ acts naturally on $C'$. The stabilizer of $y'\in C'$ in $G$ will be denoted by $G_{y'}$.

If $\pi$ is a uniformizer of $\mathcal O_{C',y'}$ and if $s\in G_{y'}$, $s\ne e$, put
\[
v_{y'}(s(\pi)-\pi)=v_{y'}(s)
\tag{4.2.1}
\]
(the multiplicity of the fixed point $y'$ for the map $s:C'\to C'$ (Lemma 5.1)).

The Swan representation associated with $y'$ is defined by its character $\sigma_{y'}$, given by
\[
\sigma_{y'}(s)=
\begin{cases}
1-v_{y'}(s),& s\ne e,\\[0.4em]
\displaystyle\sum_{s\ne e}\bigl(v_{y'}(s)-1\bigr),& s=e.
\end{cases}
\tag{4.2.2}
\]
The results of Artin--Serre--Swan (IX Th. 4) show that, for a prime number $\ell$ distinct from the characteristic of $k$, there exists a projective $\mathbb Z_\ell[G_{y'}]$-module of finite type $Sw_{y'}$, determined up to isomorphism, whose character is $\sigma_{y'}$.

It is also useful to consider the $\mathbb Z_\ell[G_{y'}]$-module
\[
Sw'_{y'}=Sw_{y'}\oplus\mathbb Z_\ell[G_{y'}].
\tag{4.2.3}
\]
The character $\sigma'_{y'}$ of $Sw'_{y'}$ is given by
\[
\sigma'_{y'}(s)=
\begin{cases}
1-v_{y'}(s),& s\ne e,\\[0.4em]
\displaystyle 1+\sum_{s\ne e}v_{y'}(s)
=1+v_{y'}(\mathfrak D_{C'/C}),& s=e,
\end{cases}
\tag{4.2.4}
\]
where $\mathfrak D_{C'/C}$ is the different of the covering $C'/C$; the last formula is just [3, Chap. VI 2, prop. 4] with $H=\{1\}$.

Finally, the function
\[
a_{y'}=\sigma_{y'}+u_{G_{y'}}=\sigma'_{y'}-1_{G_{y'}},
\tag{4.2.5}
\]
where $u_{G_{y'}}$ denotes the character of the augmentation ideal of $\mathbb Z_\ell[G_{y'}]$, is also the character of a $\mathbb Q_\ell[G_{y'}]$-module, the Artin representation of $G_{y'}$, which will be denoted by $Art_{y'}$.

Thus
\[
a_{y'}(s)=
\begin{cases}
-v_{y'}(s),& s\ne e,\\[0.4em]
\displaystyle\sum_{s\ne e}v_{y'}(s)=v_{y'}(\mathfrak D_{C'/C}),& s=e.
\end{cases}
\tag{4.2.6}
\]

Now consider the inclusion $G_{y'}\subset G$. We put
\[
Sw_y=Sw_{y'}\otimes_{\mathbb Z_\ell[G_{y'}]}\mathbb Z_\ell[G],
\qquad
Sw'_y=Sw'_{y'}\otimes_{\mathbb Z_\ell[G_{y'}]}\mathbb Z_\ell[G],
\tag{4.2.7}
\]
\[
Art_y=Art_{y'}\otimes_{\mathbb Q_\ell[G_{y'}]}\mathbb Q_\ell[G].
\tag{4.2.8}
\]
The characters of the preceding representations are denoted by $\sigma_y$, $\sigma'_y$, and $a_y$.

\subsubsection*{4.3}

The modules so obtained depend, up to isomorphism, only on the point $y\in C$, and not on the choice of the point $y'$ and $C'$ above $y$. Let $y''$ be another such point. There exists $g\in G$ such that $y''=g y'$. The situation $(C',y'',G)$ is deduced from the situation $(C,y',G)$ by the isomorphism $g$ on $C'$ and by $\operatorname{int}(g):h\mapsto ghg^{-1}$ on $G$. It follows by transport of structure that the modules $Sw_y$, etc., associated with $y''$ are deduced from those associated with $y'$ by extension of the operators by means of $\operatorname{int}(g):G\to G$. This extension is known to transform a module into an isomorphic module.

\begin{statement}{Proposition 4.4}
With the notation of 3.7, with $(\Lambda=\mathbb Z_\ell)$ and with $M^\bullet$ an object of $D^+(A[G])$, one has isomorphisms
\[
(*)\qquad Sw_y\simeq \check{Sw}_y,
\]
\[
(**)\qquad
\Hom^\bullet_{\mathbb Z_\ell[G]}(Sw_y,M^\bullet)
\simeq Sw_y\otimes_{\mathbb Z_\ell[G]}M^\bullet.
\]
\end{statement}

Proof. For two projective $\mathbb Z_\ell[G]$-modules of finite type to be isomorphic, it is enough that they have the same character (3.2). If $\check\sigma_y$ denotes the character of $\check{Sw}_y$, then by (3.7.2)
\[
\check\sigma_y(s)=\sigma_y(s^{-1}).
\]
On the other hand
\[
v_{y'}(s^{-1}(\pi)-\pi)
=v_{y'}(s(s^{-1}(\pi)-\pi))
=v_{y'}(\pi-s(\pi))
=v_{y'}(s(\pi)-\pi),
\]
so $Sw_{y'}\simeq\check{Sw}_{y'}$, whence formula $(*)$ by definition (4.2.7). Formula $(**)$ follows from 3.8.

\subsubsection*{4.5}

We shall need below an explicit formula for the character of the Swan representation. In fact one has the formulas
\[
\sigma^{*}_{y'}(s)=\sigma^{*}_{y''}(s)=\sigma_y(s)=
\begin{cases}
\displaystyle\sum_{\substack{p(z')=y\\s(z')=z'}}\bigl(1-v_{z'}(s)\bigr),& s\ne e,\\[0.8em]
\displaystyle\sum_{p(z')=y}\bigl(1+v_{z'}(\mathfrak D_{C'/C})\bigr)-n,& s=e,
\end{cases}
\tag{4.5.1}
\]
\[
\sigma^{\prime*}_{y'}(s)=\sigma^{\prime*}_{y''}(s)=\sigma'_y(s)=
\begin{cases}
\displaystyle\sum_{\substack{p(z')=y\\s(z')=z'}}\bigl(1-v_{z'}(s)\bigr),& s\ne e,\\[0.8em]
\displaystyle\sum_{p(z')=y}\bigl(1+v_{z'}(\mathfrak D_{C'/C})\bigr),& s=e,
\end{cases}
\tag{4.5.2}
\]
\[
a^{*}_{y'}(s)=a^{*}_{y''}(s)=a_y(s)=
\begin{cases}
\displaystyle-\sum_{\substack{p(z')=y\\s(z')=z'}}v_{z'}(s),& s\ne e,\\[0.8em]
\displaystyle\sum_{p(z')=y}v_{z'}(\mathfrak D_{C'/C}),& s=e.
\end{cases}
\tag{4.5.3}
\]
To deduce these formulas one uses (3.3.1) and takes into account the following fact: if one fixes a point $y'_0\in C'$, then for every $y'\in C'$ such that $p(y')=p(y'_0)$ one can find an element $g'\in G$ such that $g'(y')=y'_0$. The elements $g'$ form a system of representatives for the right cosets of $G_{y'_0}$ in $G$, and moreover
\[
G_{y'_0}=g'G_{y'}g'^{-1}.
\]
This system of representatives can be used in the second member of formula (3.3.1), and one obtains the formulas displayed above.

\subsubsection*{4.6}

The invariants $Sw_y$, $Sw'_y$, $Art_y$ can be defined under more general conditions. Let $C'$ be a smooth algebraic curve over $k$, not necessarily connected. Suppose that the finite group $G$ acts on $C'$ and that there exists an open subset $U$ of $C=C'/G$, dense in $C$, such that $C'|U$ is a principal covering of $U$ with group $G$. For a point $y\in C$, one can define the modules $Sw_y$, $Sw'_y$, $Art_y$ by proceeding as above. One can reduce to the case where $C'$ is connected in the following way. Let $y'\in C'$ be such that $p(y')=y$, let $C'_1$ be the irreducible component of $C'$ containing $y'$, and let $C_1$ be the irreducible component of $C$ containing $y$. If $G_{C'_1}$ is the stabilizer of the component $C'_1$, and if $\eta'_1$ (resp. $\eta_1$) is the generic point of $C'_1$ (resp. $C_1$), then the extension
\[
K'_1=k(\eta'_1)\supset k(\eta_1)=K_1
\]
is Galois with Galois group $G_{C'_1}$ and $C'_1$ is the normalization of $C_1$ in $K'_1$ (SGA 1 I 10). One can therefore apply the preceding considerations to the curves $C'_1$, $C_1$ and the group $G_{C'_1}$. In this way one obtains projective $\mathbb Z_\ell[G_{C'_1}]$-modules of finite type
\[
Sw_{G_{C'_1},y},\qquad Sw'_{G_{C'_1},y},\qquad Art_{G_{C'_1},y},
\]
and $Sw_y$, $Sw'_y$, $Art_y$ are the induced modules associated with the inclusion $G_{C'_1}\subset G$.

Remark 4.6. The modules $Sw_{y'}$, $Sw'_{y'}$, $Art_{y'}$ are purely local invariants: they are known from the group $G_{y'}$ acting on the completion of the local ring $\mathcal O_{C',y'}$. Moreover $Sw_{y'}=0$ if and only if $p:C'\to C$ is tamely ramified at $y'$, or equivalently, if $G_{y'}$ has order prime to the characteristic of $k$ [3 VI].

\subsection*{5. The Weil formula (compare [3 VI §4])}

\subsubsection*{5.0}

All schemes considered in this number will be defined over an algebraically closed base field $k$.

We shall use the following notation:
\[
\Lambda_\nu=\mathbb Z/\ell^\nu\mathbb Z,
\]
where $\ell$ is a fixed prime number distinct from the characteristic of $k$.

Let $\mathbb Z_\ell\to\Lambda_\nu$ be the canonical homomorphism. It induces a homomorphism
\[
\lambda_\nu:\mathbb Z_\ell[G]\longrightarrow \Lambda_\nu[G]
\]
and consequently [VIII 9] a reduction homomorphism
\[
\lambda_\nu^*:K^{\bullet}(\mathbb Z_\ell[G])\longrightarrow K^{\bullet}(\Lambda_\nu[G]).
\]
The images by $\lambda_\nu^*$ of the elements associated with the modules $Sw_y$, $Sw'_y$ will be denoted by
\[
Sw_y^{\Lambda_\nu},\qquad Sw_y'{}^{\Lambda_\nu}.
\]
The reduction homomorphism $\lambda_\nu^*$ is in fact an isomorphism [IX 1.5]. The inverse of $\lambda_\nu^*$ will be denoted by $\mu_\nu$.

The constant sheaf on the scheme $X$ associated with the ring $\Lambda_\nu$ will be denoted by $\Lambda_{\nu,X}$.

Let $V$ be an open subset of $X$, and let $j:V\hookrightarrow X$ be the canonical immersion. We put
\[
j_!(\Lambda_{\nu,V})=\Lambda_{\nu,V,X}.
\]

\begin{statement}{Proposition 5.1}
Let $C'$ be a smooth proper algebraic curve over the algebraically closed field $k$, let $G$ be a finite group acting on $C'$, let $C=C'/G$, let $p:C'\to C$ be the canonical projection, let $U$ be a dense open subset of $C$ such that $U'=p^{-1}(U)$ is a principal covering of $U$ with group $G$, let $Y=C-U$, let $Y'=p^{-1}(Y)=C'-U'$, and let $\EP(C)$ be the Euler--Poincaré characteristic of $C$ (equal to $2-2g$ if $C$ is connected of genus $g$).
Under these conditions one has the Weil formula:
\[
\clKstar{\Lambda_\nu[G]}\bigl(\RG_{C'}(\Lambda_{\nu,U',C'})\bigr)
=
\EP(C)\clKstar{\Lambda_\nu[G]}(\Lambda_\nu[G])
-\sum_{y\in Y}Sw_y'{}^{\Lambda_\nu},
\tag{5.2}
\]
where the first member of (5.2) is defined using 2.1.
\end{statement}

Proof. Denote by $S_1$ (resp. $S_2$) the first (resp. second) member of the formula to be proved. Since the reduction homomorphism $\lambda_\nu^*$ is in fact an isomorphism, it is enough to prove that
\[
\mu_\nu(S_1)=\mu_\nu(S_2).
\]
On the other hand, by (3.3.1), the trace homomorphism
\[
K^{\bullet}(\mathbb Z_\ell[G])\xrightarrow{\ tr\ }\mathcal F(G,\mathbb Z_\ell)
\]
is a monomorphism. Consequently, it is enough to prove that
\[
\operatorname{tr}(\mu_\nu(S_1))=\operatorname{tr}(\mu_\nu(S_2)).
\]
The character $\sigma_2=\operatorname{tr}(\mu_\nu(S_2))$ is plainly given by
\[
\sigma_2=\EP(C)r_G-\sum_{y\in Y}\sigma'_y,
\tag{5.2.1}
\]
where $\sigma'_y$ is the character of $Sw'_y$ ((4.2.7) and (4.2.3)), and $r_G$ is the regular representation.

Now consider the reduction homomorphisms for traces:
\[
\mathcal F(G,\mathbb Z_\ell)\xrightarrow{\lambda_n^*}\mathcal F(G,\mathbb Z/\ell^n\mathbb Z).
\]
One plainly has
\[
\mathcal F(G,\mathbb Z_\ell)=\varprojlim_n\mathcal F(G,\mathbb Z/\ell^n\mathbb Z).
\]
It is therefore enough to prove, for every $n$, that
\[
\lambda_n^*(\operatorname{tr}(\mu_\nu(S_1)))=\lambda_n^*(\sigma_2).
\]
The first member is equal to $\operatorname{tr}(\lambda_n^*(\mu_\nu(S_1)))$, while
\[
\lambda_n^*(\mu_\nu(S_1))
=
\clKstar{\Lambda_n[G]}\bigl(\RG_{C'}(\Lambda_{n,U',C'})\bigr),
\]
by the Künneth formula (SGA 4 XVIII)
\[
\RG_{C'}(\Lambda_{n,U',C'})\otimes_{\Lambda_n}^{L}\Lambda_\nu
\simeq \RG_{C'}(\Lambda_{\nu,U',C'})
\quad\text{for } n\ge\nu.
\]
Thus we are reduced to proving the equality
\[
\operatorname{tr}\bigl(\clKstar{\Lambda_n[G]}(\RG_{C'}(\Lambda_{n,U',C'}))\bigr)
=
\lambda_n^*(\sigma_2).
\]
Let $\tau$ denote the first member of this equality, and let $g\in G$. The element $g$ plainly induces an endomorphism of the distinguished triangle
\[
\begin{tikzcd}[column sep=huge,row sep=large]
& \RG_{Y'}(\Lambda_{n,Y'}) \arrow[dl] \arrow[dr] & \\
\RG_{C'}(\Lambda_{n,U',C'}) \arrow[rr] && \RG_{C'}(\Lambda_{n,C'})
\end{tikzcd}
\]
Therefore
\[
\tau(g)=\Tr^*g_{\RG_{C'}(\Lambda_{n,U',C'})}
=\Tr^*g_{\RG_{C'}(\Lambda_{n,C'})}
-\Tr^*g_{\RG_{Y'}(\Lambda_{n,Y'})}.
\]
For brevity, $\Tr^*$ has been written instead of $\Tr^*_{\Lambda_n}$.

Applying a trivial Lefschetz formula to the discrete space $Y'$ on which $g$ acts gives
\[
\begin{aligned}
\Tr^*g_{\RG_{Y'}(\Lambda_{n,Y'})}
&=\text{the number of fixed points of the restriction $g_{Y'}$ of $g$ to $Y'$}\\
&=\sum_{\substack{y'\in Y'\\gy'=y'}}1.
\end{aligned}
\]
On the other hand, the Lefschetz formula (III 4.7) for $g$ acting on $C'$ gives, if $g\ne e$,
\[
\Tr^*g_{\RG_{C'}(\Lambda_{n,C'})}
=
\sum_{y'\in Y'} v_{y'}(g),
\]
the sum of the multiplicities of the fixed points of the map $g_{C'}:C'\to C'$. Also
\[
\Tr^*g_{\RG_{C'}(\Lambda_{n,C'})}=\EP(C')\cdot1_\Lambda
\quad\text{for }g=e,
\]
by the explicit computation of $H^1(C',\Lambda_{n,C'})$, via (SGA 4 IX 47) and Künneth (SGA 4 XVII). We therefore obtain the following formula for $\tau$:
\[
\tau(g)=
\begin{cases}
\displaystyle\sum_{\substack{y'\in Y'\\g(y')=y'}}\bigl(v_{y'}(g)-1\bigr),& g\ne e,\\[0.8em]
\EP(C')-\operatorname{card}(Y'),& g=e.
\end{cases}
\tag{5.2.2}
\]
Comparison with (5.2.1) gives easily $\tau=\lambda_n^*(\sigma_2)$. Indeed, if $g\ne e$, the equality between $\tau(g)$ and $\lambda_n^*(\sigma_2)$ is immediate from (4.5.2). If $g=e$, the formula to be checked is
\[
\EP(C')-\operatorname{card}(Y')
=
N\EP(C)-\sum_{y\in Y}\sigma'_y(e),
\qquad N=\operatorname{card}(G),
\]
that is, by (4.5.2),
\[
\EP(C')-N\EP(C)=-\deg(\mathfrak D_{C'/C}),
\]
which is nothing other than the classical Hurwitz formula.

\subsection*{6. Definition of the local terms $\varepsilon_x^\Delta(F)$}

\subsubsection*{6.1}

Let $C$ be a connected smooth algebraic curve over the algebraically closed field $k$, let $A$ be a $\Lambda$-algebra ($\Lambda=\mathbb Z/\ell^\nu\mathbb Z$, with $\ell$ distinct from the characteristic of $k$), and let $\Delta\subset D^+(A)$ be a triangulated subcategory of $D^+(A)$ stable under direct factors. Let $F$ be a complex of sheaves of $A$-modules on $C$ satisfying the following conditions:

i) For every geometric point $\bar x\in C(k)$, one has $F_{\bar x}\in\Ob(\Delta)$.

ii) There exists a non-empty open subset of $C$ on which $F$ is locally constant; equivalently, there exist a non-empty open subset $U$ of $C$ and a connected Galois covering $C'\xrightarrow{p}C$ of $C$, with group $G$, such that, if $F'=p^*(F)$ is the inverse image of $F$ by $p$, then $F'|U'$ is constant $(U'=p^{-1}(U))$.

The sheaf $F'|U'$ is therefore defined by a complex $M$ of $A$-modules on which $G$ acts. If $\eta$ (resp. $\eta'$) denotes the generic point of $C$ (resp. $C'$), and if $\bar\eta$ (resp. $\bar\eta'$) is a geometric point above $\eta$ (resp. $\eta'$), then
\[
p(\eta')=\eta,\qquad
F_{\bar\eta}\simeq F'_{\bar\eta'}=M.
\tag{6.1.1}
\]
As a complex of $A$-modules, $M$ is isomorphic to $F_{\bar\eta}$ and, by condition i), is an object of $\Delta$. Hence, by (3.6), for every $x\in C(k)$,
\[
\Hom^\bullet_{\Lambda[G]}(Sw_x^\Lambda,M)
\]
is an object of $\Delta$, and so defines an element of $K(\Delta)$.

In fact the object of $\Delta$ defined by $\Hom^\bullet_{\Lambda[G]}(Sw_x^\Lambda,M)$ is independent, up to isomorphism, of the Galois covering $C'\to C$ which trivializes $F|U$. To show this it is enough to compare the elements associated with two coverings
\[
C'\xrightarrow{p}C,\qquad C''\xrightarrow{p'}C'\xrightarrow{p}C.
\]
In this case, if $G'$ is the Galois group of the covering $C''\xrightarrow{p''}C$, then $G$ is a quotient group of $G'$. Let
\[
\varphi:\mathbb Z_\ell[G']\longrightarrow \mathbb Z_\ell[G]
\]
be the associated canonical epimorphism. The Swan modules are then related by
\[
Sw_{G,x}\simeq Sw_{G',x}\otimes_{\mathbb Z_\ell[G']}\mathbb Z_\ell[G],
\]
which follows immediately from the fact that the character $a_{G,x}$ is the character induced from $G'$ to $G$ ([3] VI Prop. 3).

On the other hand, there is an isomorphism of complexes of $A[G']$-modules
\[
F_{\bar\eta''}\simeq [F_{\bar\eta'}]_{(\varphi)}.
\]
It follows that there is a canonical isomorphism of complexes of $A$-modules between
\[
\Hom^\bullet_{\Lambda[G]}(Sw_{G,x}^\Lambda,F_{\bar\eta})
\quad\text{and}\quad
\Hom^\bullet_{\Lambda[G']}(Sw_{G',x}^\Lambda,F_{\bar\eta''})
\]
(the adjunction formula).

\subsubsection*{6.2}

We shall use the following notation:
\[
\alpha_x^\Delta(F)
=
\clK{\Delta}\bigl(\Hom^\bullet_{\Lambda[G]}(Sw_x^\Lambda,F_{\bar\eta})\bigr),
\tag{6.2.1}
\]
\[
\varepsilon_x^\Delta(F)
=
\alpha_x^\Delta(F)+
\bigl(\clK{\Delta}(F_{\bar\eta})-\clK{\Delta}(F_x)\bigr).
\tag{6.2.2}
\]
Of course, in formula (6.2.1), $F_{\bar\eta}$ is endowed with its $A[G]$-module structure by the isomorphism (6.1.1). Note that if $x\in C(k)$ is such that $C'$ is tamely ramified at $x$, then by (4.6) $Sw_x=0$, hence $\alpha_x^\Delta(F)=0$, and in this case, in conformity with Ogg-Shafarevich, one has
\[
\varepsilon_x^\Delta(F)=\clK{\Delta}(F_{\bar\eta})-\clK{\Delta}(F_x).
\tag{6.2.3}
\]
It is clear that $\varepsilon_x^\Delta(F)=0$ if $x\notin Y=C-U$, and therefore only finitely many $x$ satisfy $\varepsilon_x^\Delta(F)\ne0$.

We now give examples of categories $\Delta$ and complexes of sheaves $F$ for which the conditions above are satisfied.

a) Let $\Delta=D^b(A)_{\coh}$ (VIII 5), and let $F$ be a constructible complex of sheaves of $A$-modules (SGA 4 XVII) with bounded cohomology. In this case $\varepsilon_x^\Delta(F)\in K_\bullet(A)$.

b) Take $\Delta=D(A)_{\parf}$ (VIII 6), and let $F$ be a constructible complex of sheaves of $A$-modules of finite tor-dimension. In this case $\varepsilon_x^\Delta(F)\in K^{\bullet}(A)$.

c) Let $A$, $B$ be two $\Lambda$-algebras and let $\varphi:A\to B$ be a homomorphism of $\Lambda$-algebras. Take $\Delta=D^b(B\operatorname{ rel }A)_{\coh}$ (VIII 10), and let $F$ be a complex of sheaves of $B$-modules, constructible as a sheaf of complexes of $A$-modules and with bounded cohomology. One obtains $\varepsilon_x^\Delta(F)\in K_\bullet(B\operatorname{ rel }A)$. Of course, one can vary this as above to obtain, for example, elements $\varepsilon_x^\Delta(F)\in K^{\bullet}(B\operatorname{ rel }A)$.

\subsection*{7. Euler--Poincaré formula}

\begin{statement}{Theorem 7.1}
Let $C$ be a connected, smooth, proper algebraic curve over the algebraically closed field $k$, let $\eta$ be the generic point of $C$, let $\Delta\subset D^+(A)$ be a triangulated subcategory of $D^+(A)$ stable under direct factors, and let $F$ be a complex of sheaves of $A$-modules satisfying conditions i), ii) of 7.1. Under these conditions one has $\RG_C(F)\in\Ob(\Delta)$, and the following formula (Euler--Poincaré formula) is true:
\[
\clK{\Delta}(\RG_C(F))
=
\EP(C)\clK{\Delta}(F_{\bar\eta})
-\sum_{x\in C(k)}\varepsilon_x^\Delta(F),
\tag{7.2}
\]
where the $\varepsilon_x^\Delta(F)$ are defined in 6.2 (formula 6.2.2).
\end{statement}

The proof of this theorem will be made in a few steps.

a) Let $\Psi$ be a bounded below complex of sheaves of $\Lambda[G]$-modules on $C$. There is an isomorphism
\[
R\Gamma^G(\RG_C(\Psi))\simeq\RG_C(R\underline{\Gamma}^G(\Psi)),
\tag{7.3}
\]
where $\underline{\Gamma}^G$ is the functor ``subsheaf of $G$-invariant sections.'' Indeed, $\Gamma_C$ sends an injective sheaf of $\Lambda[G]$-modules to an injective $\Lambda[G]$-module, and $\underline{\Gamma}^G$ sends an injective sheaf of $\Lambda[G]$-modules to an injective sheaf of $\Lambda$-modules (using, for example, the evident adjoints of $\Gamma_C$ and $\underline{\Gamma}^G$). Moreover
\[
\Gamma^G\circ\Gamma_C=\Gamma_C\circ\underline{\Gamma}^G.
\]
Taking derived functors on both sides and applying [5] prop. 3.1) gives (7.3).

b) Let $\Phi$ be a bounded below complex of sheaves of $\Lambda$-modules on $C$, zero outside $U$. There is an isomorphism
\[
\Phi\simeq R\underline{\Gamma}^G(p_*(\Phi')),
\qquad \Phi'=p^*(\Phi).
\tag{7.4}
\]
First note that the functors $p^*$ and $\underline{\Gamma}^G\circ p_*$ establish quasi-inverse equivalences between the category of sheaves of $\Lambda$-modules on $C'$ on which $G$ acts and which are zero outside $U'=p^{-1}(U)$, and the category of sheaves of $\Lambda$-modules on $C$ which are zero outside $U$. Hence one has the isomorphism
\[
\Phi\simeq R\bigl((\underline{\Gamma}^G\circ p_*)\circ p^*\bigr)(\Phi).
\]
On the other hand, one can again apply [5] Prop. 3.1), which gives, since $Rp^*=p^*$ and $Rp_*=p_*$:
\[
\Phi\simeq R(\underline{\Gamma}^G\circ p_*)(p^*(\Phi))
\simeq (R\underline{\Gamma}^G)(p_*p^*(\Phi)),
\]
which is the desired isomorphism.

c) Let $\Phi$ be a complex of sheaves as in b). One has the isomorphism
\[
\RG_C(\Phi)=(R\Gamma^G)(\RG_{C'}(\Phi')),
\qquad \Phi'=p^*(\Phi).
\tag{7.5}
\]
Indeed,
\[
(R\Gamma^G)(\RG_{C'}(\Phi'))
\simeq
(R\Gamma^G)(\RG_C(p_*(\Phi')))
\simeq
\RG_C(R\underline{\Gamma}^G(p_*(\Phi')))
\]
by (7.3), and (7.5) follows from the isomorphism (7.4).

d) Let $i':U'\to C'$, $i:U\to C$ be the canonical inclusions and put
\[
\Phi=F_{U,C}=i_!(F|U).
\]
To compute $\RG_C(\Phi)$, one may apply formula (7.5). On the other hand, with the notation of 7.1,
\[
\Phi'\simeq \Lambda_{U',C'}\otimes_\Lambda^L M,
\]
and the Künneth formula (SGA 4 XVII) gives
\[
\RG_{C'}(\Phi')
\simeq
\RG_{C'}(\Lambda_{U',C'}\otimes_\Lambda^L M)
\xleftarrow{\sim}
\RG_{C'}(\Lambda_{U',C'})\otimes_\Lambda^L M.
\tag{7.6}
\]
This last isomorphism is in fact an isomorphism between the objects of $D(A[G])$ defined by the two sides. Combining (7.6), (7.5), and formula (3.9.1), which is applicable by 2.1, one obtains
\[
\RG_C(F_{U,C})
\simeq
R\Gamma^G\bigl(\RG_{C'}(\Lambda_{U',C'})\otimes_\Lambda^L M\bigr)
\simeq
\RHom_{\Lambda[G]}\bigl(\RG_{C'}(\Lambda_{U',C'})^\vee,M\bigr).
\tag{7.7}
\]
Thus, since $M\simeq F_{\bar\eta}\in\Ob(\Delta)$ (condition (i)) and by (3.6), we find
\[
\RG_C(F_{U,C})\in\Ob(\Delta).
\]
On the other hand, there is the distinguished triangle (where $Y=C-U$)
\[
\begin{tikzcd}[column sep=huge,row sep=large]
& \RG_Y(F|Y) \arrow[dl] & \\
\RG_C(F_{U,C}) \arrow[rr] && \RG_C(F) \arrow[ul]
\end{tikzcd}
\tag{7.8}
\]
where obviously
\[
\RG_Y(F)=\coprod_{y\in Y}F_y\in\Ob(\Delta),
\]
because $F_y\in\Ob(\Delta)$ (condition (i)). It follows that $\RG_C(F)\in\Ob(\Delta)$, which is the first assertion of the theorem. The distinguished triangle (7.8) furthermore gives the relation
\[
\clK{\Delta}(\RG_C(F))
=
\clK{\Delta}(\RG_C(F_{U,C}))
+\sum_{y\in Y}\clK{\Delta}(F_y),
\]
and, using the isomorphism 7.7,
\[
\clK{\Delta}(\RG_C(F))
=
\clK{\Delta}\bigl(\RHom_{\Lambda[G]}(\RG_{C'}(\Lambda_{U',C'})^\vee,F_{\bar\eta})\bigr)
+\sum_{y\in Y}\clK{\Delta}(F_y).
\]
To obtain the Euler--Poincaré formula it remains only to use the Weil formula (5.2). Indeed, first 4.4, (5.2), and 3.6 (in the particular case $B=\Lambda[G]$) give
\[
\clKstar{\Lambda[G]}\bigl(\RG_{C'}(\Lambda_{U',C'})\bigr)
=
\clKstar{\Lambda[G]}\bigl(\RG_{C'}(\Lambda_{U',C'})^\vee\bigr).
\]
\[
\begin{aligned}
\clK{\Delta}(\RG_C(F))
&=\operatorname{hom}\Bigl(
\clKstar{\Lambda[G]}\bigl(\RG_{C'}(\Lambda_{U',C'})\bigr),
\cl_{K(D(A[G])_\Delta)}(F_{\bar\eta})\Bigr)
 +\sum_{y\in Y}\clK{\Delta}(F_y)\\
&=\sum_{y\in Y}\clK{\Delta}(F_y)\\
&\quad +\operatorname{hom}\Bigl(
\EP(C)\clKstar{\Lambda[G]}(\Lambda[G])
 -\sum_{y\in Y}\bigl(Sw_y^{\Lambda}+\clKstar{\Lambda[G]}(\Lambda[G])\bigr),\\
&\hspace{10em}\cl_{K(D(A[G])_\Delta)}(F_{\bar\eta})\Bigr)\\
&=\sum_{y\in Y}\clK{\Delta}(F_y)
 +\EP(C)\clK{\Delta}(F_{\bar\eta})
 -\sum_{y\in Y}\bigl(\alpha_y^\Delta(F)+\clK{\Delta}(F_{\bar\eta})\bigr)\\
&=\EP(C)\clK{\Delta}(F_{\bar\eta})
 -\sum_{y\in Y}\varepsilon_y^\Delta(F).
\end{aligned}
\]
Since $\varepsilon_x^\Delta(F)=0$ for $x\notin Y$, this is exactly formula (7.2). For a complex of sheaves $F$ satisfying the conditions of 7.1, we introduce the notation
\[
\EP_\Delta(C,F)=\clK{\Delta}(\RG_C(F)).
\]

\begin{statement}{Corollary 7.9}
If $F$, $F'$ are complexes of sheaves satisfying the conditions of 7.1 and if they are locally isomorphic (for the étale topology) as complexes of $A$-modules, then
\[
\EP_\Delta(C,F)=\EP_\Delta(C,F').
\tag{7.10}
\]
In particular, if $F$ is locally constant of value $M^\bullet$, then
\[
\EP_\Delta(C,F)=\EP(C)\clK{\Delta}(M^\bullet).
\tag{7.11}
\]
\end{statement}

\begin{statement}{Corollary 7.12}
Let $C$ be a connected smooth proper algebraic curve over $k$, let $V$ be a non-empty open subset of $C$, let $\Delta$ be a subcategory of $D^+(A)$ as in 7.1, and let $F$ be a complex of sheaves of $A$-modules on $V$ such that its restriction to a suitable non-empty open subset of $V$ is trivialized by a Galois covering of $V$. Suppose moreover that $F_{\bar x}\in\Ob(\Delta)$ for every $x\in V$. Under these conditions one has
\[
(*)\qquad \RG_!V(F),\ \RG_V(F)\in\Ob(\Delta),
\]
\[
(**)\qquad
\clK{\Delta}(\RG_!V(F))=\clK{\Delta}(\RG_V(F)).
\]
Put $\EP_\Delta(V,F)=\clK{\Delta}(\RG_V(F))$ and
\[
\EP(V)=\EP_{D(F_\ell)_{\coh}}(V,F_\ell)=\EP(C)-\operatorname{card}(Z),
\]
where $Z=C-V$ and where $F_\ell$ is the sheaf on $C$ defined in dimension zero by the constant sheaf associated with $F_\ell=\mathbb Z/\ell\mathbb Z$. Then one has the formula
\[
\EP_\Delta(V,F)
=
\EP(V)\clK{\Delta}(F_{\bar\eta})
-\sum_{x\in V(k)}\varepsilon_x^\Delta(F)
-\sum_{x\in Z}\alpha_x^\Delta(F),
\tag{7.13}
\]
where $\varepsilon_x^\Delta$ and $\alpha_x^\Delta$ are defined by (6.2.1), (6.2.2).
\end{statement}

We shall use the following lemma.

\begin{statement}{Lemma 7.14}
Let $K$ be the field of fractions of a strictly local discrete valuation ring $R$, let $W=\Spec K$, let $A$ be a noetherian ring annihilated by an integer $n$ prime to the residual characteristic $p$ of $R$, and let $F$ be a complex of sheaves of $A$-modules on $W$ trivialized by a finite extension of $K$. If $\Delta$ is a triangulated subcategory of $D^+(A)$ stable under direct factors and if $F_{\bar\eta}\in\Ob(\Delta)$ ($\eta$ the generic point of $S$), then $\RG_W(F)\in\Ob(\Delta)$ and
\[
\clK{\Delta}(\RG_W(F))=0.
\]
\end{statement}

Proof of 7.14. Let $\pi$ be the fundamental group of $W$. Then $F$ corresponds to a complex of $A$-modules $M^\bullet$ with $\pi$ as group of operators. The underlying complex of $A$-modules of $M^\bullet$ is isomorphic to $F_{\bar\eta}$ and consequently belongs to $\Ob(\Delta)$. Introduce the quotient group $\bar\pi$ of $\pi$ by its unique Sylow $p$-subgroup $P$ [4, II 33], that is, the projective limit of discrete quotient groups of $\pi$ which are of order prime to $p$, where $p$ is the characteristic exponent of the residue field of $R$ (if $p=1$, put $\bar\pi=\pi$). Thus there is an exact sequence
\[
1\longrightarrow P\longrightarrow \pi\longrightarrow \bar\pi\longrightarrow0,
\]
and one obtains
\[
\RG_W(F)\simeq R\Gamma^\pi(M^\bullet)
\simeq R\Gamma^{\bar\pi}(R\Gamma^P(M^\bullet))
\]
(Hochschild--Serre isomorphism). Now $P$ acts on $M^\bullet$ through a finite quotient group, which is therefore a $p$-group, hence of order prime to $n$. Thus
\[
R\Gamma^P(M^\bullet)\simeq \Gamma^P(M^\bullet),
\]
and $\Gamma^P(M^\bullet)$ is a direct factor of $M^\bullet$, hence belongs to $\Ob(\Delta)$. On the other hand
\[
\bar\pi=\prod_{\ell\ne p}\mathbb Z_\ell
\]
([4], II 33). We are therefore reduced to proving that if $M^\bullet\in\Ob(\Delta)$, then $R\Gamma^{\bar\pi}(M^\bullet)\in\Ob(\Delta)$ and
\[
\clK{\Delta}(R\Gamma^{\bar\pi}(M^\bullet))=0.
\]
For this, let $f$ be a topological generator of $\bar\pi$. We may suppose that the components $M^i$ of $M$ are $\bar\pi$-modules such that
\[
H^j(\bar\pi,M^i)=0\quad\text{if }j\ge1.
\]
There is then a distinguished triangle
\[
\begin{tikzcd}[column sep=huge,row sep=large]
& M^\bullet \arrow[dl] &\\
R\Gamma^{\bar\pi}(M^\bullet) \arrow[rr] && M^\bullet \arrow[ul,"1-f"]
\end{tikzcd}
\tag{7.15}
\]
Indeed, on the one hand it is clear that $R\Gamma^{\bar\pi}=\Ker(1-f)$. On the other hand $1-f:M^\bullet\to M^\bullet$ is surjective. Indeed
\[
H^1(\bar\pi,M)=M^\bullet/(1-f)M^\bullet
\]
(see for this calculation [4] pp. 197, in a substantially analogous case), and therefore $M^\bullet=(1-f)M^\bullet$. The triangle (7.15) gives at the same time $R\Gamma^{\bar\pi}(M^\bullet)\in\Ob(\Delta)$ and
\[
\clK{\Delta}(R\Gamma^{\bar\pi}(M^\bullet))=0,
\]
which proves 7.14.

Proof of 7.12. The fact that $\RG_!V(F)\in\Ob(\Delta)$ follows immediately from 7.1 applied to the complex $i_!(F)$, where $i:V\to C$; indeed
\[
\RG_!V(F)=\RG_C(i_!F).
\]
To prove that $\RG_V(F)\in\Ob(\Delta)$, consider the distinguished triangle
\[
\begin{tikzcd}[column sep=huge,row sep=large]
& Q^\bullet \arrow[dl] &\\
Ri_!(F) \arrow[rr] && Ri_*(F) \arrow[ul],
\end{tikzcd}
\]
where $Q^\bullet$ is concentrated on $Z=C-V$. Applying $\RG_C$ gives the distinguished triangle
\[
\begin{tikzcd}[column sep=huge,row sep=large]
& \RG_Z(Q^\bullet|Z) \arrow[dl] &\\
\RG_!V(F) \arrow[rr] && \RG_V(F) \arrow[ul].
\end{tikzcd}
\tag{7.16}
\]
Since $\RG_!V(F)\in\Ob(\Delta)$, it follows that $\RG_V(F)\in\Ob(\Delta)$ if and only if $\RG_Z(Q^\bullet|Z)\in\Ob(\Delta)$, that is, if and only if $(Ri_*(F))_{\bar x}\in\Ob(\Delta)$ for every $x\in Z$. Recall that to calculate $(Ri_*(F))_{\bar x}$, one considers the cartesian square
\[
\begin{tikzcd}
\bar V \arrow[r,hook,"j"] \arrow[d,"u"'] & C_x \arrow[d]\\
V \arrow[r,hook] & C
\end{tikzcd}
\]
where $C_x$ is the strictly localized scheme corresponding to the geometric point $x$. One has
\[
Ri_*(F)_{\bar x}\simeq \RG_{\bar V}(u^*(F)),
\]
by virtue of 7.14, and therefore $(Ri_*(F))_{\bar x}\in\Ob(\Delta)$ and
\[
\clK{\Delta}((Ri_*(F))_{\bar x})=0.
\]
It follows that
\[
\clK{\Delta}(\RG_Z(Q^\bullet|Z))=0,
\]
and the distinguished triangle (7.16) gives
\[
\clK{\Delta}(\RG_!V(F))=\clK{\Delta}(\RG_V(F)).
\]
Finally, to establish (7.13), it is enough to use the Euler--Poincaré formula (7.2) for the complex of sheaves $i_!(F)$:
\[
\EP_\Delta(V,F)
=
\clK{\Delta}(\RG_!V(F))
=
\clK{\Delta}(\RG_C(i_!(F)))
=
\EP(C)\clK{\Delta}(F_{\bar\eta})
-\sum_{x\in C(k)}\varepsilon_x^\Delta(i_!(F)).
\]
But
\[
\sum_{x\in C(k)}\varepsilon_x^\Delta(i_!(F))
=
\sum_{x\in V(k)}\varepsilon_x^\Delta(F)
+\sum_{x\in Z}\varepsilon_x^\Delta(F).
\]
For $x\in Z$ one has $(i_!(F))_{\bar x}=0$, and consequently, for such $x$,
\[
\varepsilon_x^\Delta(i_!(F))
=
\alpha_x^\Delta(F)+\clK{\Delta}(F_{\bar\eta}),
\]
which gives formula (7.13).

\section*{References}

\begin{enumerate}
\item[{[1]}] M. Raynaud, Caractéristique d'Euler-Poincaré d'un faisceau et cohomologie des variétés abéliennes. Séminaire Bourbaki 1965, n\textsuperscript{o} 286.
\item[{[2]}] A. Grothendieck, Sur quelques points d'algèbre homologique, Tohoku Math. Journal, vol. 9 (1957), p. 119--221, (cited [T]).
\item[{[3]}] Serre, J.-P., Corps locaux, Act. Sci. Ind., 1296 (1962), Hermann, Paris.
\item[{[4]}] Serre, J.-P., Cohomologie Galoisienne, Lecture Notes in Mathematics 5, 1964 (Springer).
\item[{[5]}] Verdier, J.-L., Catégories dérivées, I H E S 1963 (published in SGA $4^{1/2}$).
\end{enumerate}


\phantomsection\addcontentsline{toc}{section}{Exposé XII --- Nielsen--Wecken and Lefschetz Formulas}
\section*{Exposé XII}

\begin{center}
\Large\bfseries NIELSEN--WECKEN AND LEFSCHETZ FORMULAS IN ALGEBRAIC GEOMETRY
\end{center}

\begin{center}
by A. Grothendieck

written up by I. BUCUR
\end{center}

\bigskip

\section*{1.}
Let $X'$ be a scheme proper over the algebraically closed field $k$, and let $G$ be a finite group acting on $X'$.

For each element $g\in G$, denote by
\[
g_{X'}:X'\to X'
\]
the automorphism of $X'$ induced by $g$.

Suppose that there exists an open subset $U'$ of $X'$, stable under $G$, such that $G$ acts freely on $U'$, and let
\[
\La_n=\bZ/\ell^n\bZ,
\]
where $\ell$ is a prime number different from the characteristic of $k$.

Let $\La_{n,U',X'}$ denote the sheaf $i_!(\La_{n,U'})$ on $X'$, where $\La_{n,U'}$ is the constant sheaf on $U'$ associated with the ring $\La_n$ and where $i:U'\hookrightarrow X'$ is the canonical immersion of $U'$ into $X'$. The sheaf $\La_{n,U',X'}$ is endowed in a natural way with the structure of a sheaf with group $G$ of operators. Consequently the complex
\[
H_{n,U',X'}=R\Gamma_{X'}(\La_{n,U',X'})
\]
is a complex of $\La_n[G]$-modules, and more precisely a perfect complex (VIII, 7) of $\La_n[G]$-modules (IX, 2).

We shall see below that it is more convenient for our purpose to consider, instead of the complex $R\Gamma_{X'}(\La_{n,U',X'})$, its dual
\[
H_{n,U',X'}^\vee=R\Gamma_{X'}(\La_{n,U',X'})^\vee
\]
with respect to the ring $\La_n$, in the sense of derived categories:
\[
H_{n,U',X'}^\vee
=R\Gamma_{X'}(\La_{n,U',X'})^\vee
=R\Hom_{\La_n}(R\Gamma_{X'}(\La_{n,U',X'}),\La_n).
\]
The complex $R\Gamma_{X'}(\La_{n,U',X'})^\vee$ is a perfect complex of $\La_n[G]$-modules.

Suppose now that an endomorphism $f':X'\to X'$ of the scheme $X'$ and an endomorphism $\varphi:G\to G$ of the group $G$ are given, satisfying the following two conditions:
\begin{equation}\tag{1.1}
f'g_{X'}=\varphi(g)_{X'}f'\qquad \forall g\in G .
\end{equation}
\begin{equation}\tag{1.2}
\text{The complement }Y'=X'-U'\text{ is stable under }f'.
\end{equation}
Under the hypotheses already made, $Y'$ is stable under $G$. In general one does not require the open subset $U'$ to be stable under $f'$.

We therefore obtain homomorphisms of sheaves
\begin{align*}
f'_{\La_{n,X'}} &: \La_{n,X'}\longrightarrow \La_{n,X'},\qquad
f'_{\La_{n,Y'}} : \La_{n,Y'}\longrightarrow \La_{n,Y'},\\
f'_{\La_{n,U',X'}} &: \La_{n,U',X'}\longrightarrow \La_{n,U',X'},
\qquad
g_{\La_{n,U',X'}} : \La_{n,U',X'}\longrightarrow \La_{n,U',X'},
\end{align*}
such that the following diagram, whose rows are exact, is commutative:
\[
\begin{tikzcd}
0 \arrow[r] & \La_{n,U',X'} \arrow[r] \arrow[d,"f'_{\La_{n,U',X'}}"'] & \La_{n,X'} \arrow[r] \arrow[d,"f'_{\La_{n,X'}}"] & \La_{n,Y'} \arrow[r] \arrow[d,"f'_{\La_{n,Y'}}"] & 0 \\
0 \arrow[r] & \La_{n,U',X'} \arrow[r] & \La_{n,X'} \arrow[r] & \La_{n,Y'} \arrow[r] & 0 .
\end{tikzcd}
\tag{1.3}
\]
We deduce in particular homomorphisms of complexes of $\La_n$-modules
\begin{align*}
f'_{H_{n,U',X'}}&=R\Gamma_{X'}(f'_{\La_{n,U',X'}}):H_{n,U',X'}\to H_{n,U',X'},\\
f'^\vee_{H_{n,U',X'}}&=(f'_{H_{n,U',X'}})^\vee:H_{n,U',X'}^\vee\to H_{n,U',X'}^\vee,\\
g_{H_{n,U',X'}}&=R\Gamma_{X'}(g_{\La_{n,U',X'}}):H_{n,U',X'}\to H_{n,U',X'},\\
g^\vee_{H_{n,U',X'}}&=(g_{H_{n,U',X'}})^\vee:H_{n,U',X'}^\vee\to H_{n,U',X'}^\vee,
\end{align*}
such that the following relations hold:
\begin{align}
\tag{1.4}g_{H_{n,U',X'}}\circ f'_{H_{n,U',X'}}&=f'_{H_{n,U',X'}}\circ \varphi(g)_{H_{n,U',X'}},\\
\tag{1.5}f'^\vee_{H_{n,U',X'}}\circ g^\vee_{H_{n,U',X'}}&=\varphi(g)^\vee_{H_{n,U',X'}}\circ f'^\vee_{H_{n,U',X'}} .
\end{align}
Relation (1.5) shows (cf. Exposé XI on the general theory of traces, 5.2.1, 7.1) that the trace
\[
\Tr^*_{\La_n[G],\varphi}(f'^\vee_{H_{n,U',X'}})
\in \La_n[G]_{\natural,\varphi}\simeq \mathcal F_!(G_{\natural,\varphi},\La_n)
\qquad (*)\footnote{Cf. III B 5.12.}
\]
is defined.

Our first aim will be to give, under some additional hypotheses, an explicit formula for $\Tr^*_{\La_n[G],\varphi}(f'^\vee_{H_{n,U',X'}})$ in terms of local invariants associated with $f'$, terms which we shall define below.

\section*{2.}
In fact, one can do better. For this, let $n,m$ be two positive integers such that $m\ge n$. The ring homomorphisms
\[
\alpha_{m,n}:\La_m\longrightarrow \La_n
\]
give rise to homomorphisms of sheaves of rings on $X'$:
\begin{align*}
\alpha_{m,n,X'}&:\La_{m,X'}\longrightarrow \La_{n,X'},\\
\alpha_{m,n,U',X'}&:\La_{m,U',X'}\longrightarrow \La_{n,U',X'} .
\end{align*}
Application of the Künneth formula gives
\begin{align}
\tag{2.1}R\Gamma_{X'}(\La_{m,X'})\otimes_{\La_m}^{\mathbb L}\La_n&\simeq R\Gamma_{X'}(\La_{n,X'}),\\
\tag{2.2}R\Gamma_{X'}(\La_{m,U',X'})\otimes_{\La_m}^{\mathbb L}\La_n&\simeq R\Gamma_{X'}(\La_{n,U',X'}),\\
\tag{2.3}R\Gamma_{X'}(f'_{\La_{m,X'}})\otimes_{\La_m}^{\mathbb L}\La_n&\simeq R\Gamma_{X'}(f'_{\La_{n,X'}}),\\
\tag{2.4}R\Gamma_{X'}(f'_{\La_{m,U',X'}})\otimes_{\La_m}^{\mathbb L}\La_n&\simeq R\Gamma_{X'}(f'_{\La_{n,U',X'}}).
\end{align}
One deduces easily that
\begin{align}
\tag{2.5}R\Hom_{\La_m}(R\Gamma_{X'}(\La_{m,U',X'}),\La_m)\otimes_{\La_m}^{\mathbb L}\La_n&\simeq R\Hom_{\La_n}(R\Gamma_{X'}(\La_{n,U',X'}),\La_n),\\
\tag{2.6}f'^\vee_{H_{m,U',X'}}\otimes_{\La_m}^{\mathbb L}\La_n&\simeq f'^\vee_{H_{n,U',X'}}.
\end{align}
The morphism of perfect complexes $R\Gamma_{X'}(f'_{\La_{n,X'}})$ (resp. $R\Gamma_{X'}(f'_{\La_{n,U',X'}})$, $f'^\vee_{H_{n,U',X'}}$) is therefore obtained by extension of scalars (XI, 6.3) from the morphisms $R\Gamma_{X'}(f'_{\La_{m,X'}})$ (resp. $R\Gamma_{X'}(f'_{\La_{m,U',X'}})$, $f'^\vee_{H_{m,U',X'}}$) by means of the morphisms $\alpha_{m,n}$ (resp. $\alpha_{m,n}$, $\alpha_{m,n,G}:\La_m[G]\to\La_n[G]$).

Using (XI, 6.3, 6.3.2, 6.3.3), 2.3, and 2.6, one obtains the relations
\begin{align}
\tag{2.7}\Tr^*_{\La_n}(R\Gamma_{X'}(f'_{\La_{n,X'}}))&=\alpha_{m,n}(\Tr^*_{\La_m}(R\Gamma_{X'}(f'_{\La_{m,X'}}))),\\
\tag{2.8}\Tr^*_{\La_n[G],\varphi}(f'^\vee_{H_{n,U',X'}})&=\alpha_{m,n,G}(\Tr^*_{\La_m[G],\varphi}(f'^\vee_{H_{m,U',X'}})).
\end{align}
The system of functions
\[
(\Tr_{\La_n[G],\varphi}(f'^\vee_{H_{n,U',X'}}))_{n>0}
\]
defines a function
\[
\Tr_{\bZ_\ell[G],\varphi}(f'^\vee_{H_{U',X'}(\ell)})
\in \mathcal F_!(G_{\natural\varphi},\bZ_\ell),
\]
which is $\varphi$-central on $G$ and has values in the ring $\bZ_\ell$ of $\ell$-adic integers.

Similarly, the elements $\Tr^*_{\La_n}(R\Gamma_{X'}(f'_{\La_{n,X'}}))$, $n\ge1$, define an $\ell$-adic integer
\[
\Tr^*_{\bZ_\ell}(f'_{R\Gamma_{X'}(\bZ_\ell,X')}).
\]
Our aim will be to find a formula for
\[
\Tr^*_{\bZ_\ell[G],\varphi}(f'^\vee_{H_{U',X'}(\ell)}).
\]

\section*{3. The local Nielsen--Wecken invariant (cf. [1])}
Suppose, in addition to the hypotheses described in number 1, that every orbit of $G$ is contained in an affine open subset. Then a quotient $X=X'/G$ exists and, by formula 1.1, the endomorphism $f'$ of $X'$ descends to the quotient:
\[
f:X\longrightarrow X .
\]
We shall use the following notation:
\begin{align*}
p &: X'\longrightarrow X=X'/G\quad \text{for the canonical projection,}\\
Y&=p(Y')=Y'/G,\\
U&=X-Y.
\end{align*}
The open $U'$ is a principal covering with structural group $G$. Let $x\in U$ be a fixed point of $f$. We shall associate with it a $\varphi$-conjugacy class (XI, 7.1), denoted $\mathrm{NW}_{f}^{G,\varphi}(x)$, as follows.

Let $x'\in U'$ be such that $p(x')=x$. Since $U'$ is a principal covering with structural group $G$ and since $x$ is fixed by $f$, there exists a unique element $g\in G$ such that $f'(x')=gx'$. The $\varphi$-conjugacy class of $g$ depends only on $x$. Indeed, if $p(x'')=x$, there exists $h\in G$ such that $hx''=x'$. Consequently, if $f'(x'')=g'x''$, then
\begin{align*}
gx'&=f'(x')=f'(hx'')=\varphi(h)f'(x'')=(\varphi(h)g'h^{-1})x',\\
g&=\varphi(h)g'h^{-1},
\end{align*}
that is, $g$ and $g'$ are $\varphi$-conjugate. Thus one obtains a well-defined $\varphi$-conjugacy class, the local Nielsen--Wecken invariant associated with the fixed point $x\in U$ of $f$, which will be denoted by $\mathrm{NW}_{f}^{G,\varphi}(x)$.

The $\varphi$-conjugacy class $\mathrm{NW}_{f}^{G,\varphi}(x)$ defines a $\varphi$-central function (cf. XI, 7.1)
\[
\mathrm{NW}_{f}^{G,\varphi}(x):G\longrightarrow \bZ
\]
with integer values, characterized by the following relations:
\begin{equation}\tag{3.1}
\mathrm{NW}_{f}^{G,\varphi}(x)(g)=
\begin{cases}
1&\text{if there exists }x'\in X'\text{ such that }(g^{-1}f')(x')=x'\\
 &\qquad\text{and }p(x')=x,\\
0&\text{otherwise.}
\end{cases}
\end{equation}

\section*{4. Generalities on local terms}
Suppose that $X'$ is smooth over $k$ and, in addition, that the number of fixed points of $f$ is finite. It follows in particular that the number of fixed points of $g^{-1}f'$ is finite for every $g\in G$.

For each fixed point $x\in X$ of $f$, introduce the $\varphi$-central function
\[
S_x(f',f,G,\varphi):G\longrightarrow \bQ,
\]
with rational values, defined by the formula
\begin{equation}\tag{4.1}
S_x(f',f,G,\varphi)(g)=
\frac{1}{c^\varphi(g)}
\sum_{x'\in {X'_x}^{g^{-1}f'}}(\nu_{x'}(g^{-1}f')-1).
\end{equation}
We have used the following notation:
\begin{equation}\tag{4.2}
c^\varphi(g)=\card G_g^\varphi,
\qquad
G_g^\varphi=\{x\in G\mid \varphi(x)gx^{-1}=g\},
\end{equation}
\begin{equation}\tag{4.3}
{X'_x}^{g^{-1}f'}=\{x'\in X'\mid p(x')=x,\\ (g^{-1}f')(x')=x'\},
\end{equation}
\begin{equation}\tag{4.4}
\nu_{x'}(g^{-1}f')=\text{the multiplicity of the fixed point }x'\text{ of }g^{-1}f'.
\end{equation}
By composition with the canonical inclusions
\[
\bQ\hookrightarrow \bQ_\ell,
\]
the function $S_x(f',f,G,\varphi)$ gives $\varphi$-central functions $S_{x,\ell}(f',f,G,\varphi)$ on $G$ with values in the field $\bQ_\ell$ of $\ell$-adic numbers.

\medskip
\noindent
\textbf{Conjecture 4.5.} \emph{For each $\ell\ne\car(k)$, the values of the function $S_{x,\ell}(f',f,G,\varphi)$ are $\ell$-adic integers.}

\medskip
We shall prove below that Conjecture 4.5 is true when the dimension of $X$ is equal to $1$.

\medskip
\noindent
\textbf{Proposition 4.6.} \emph{Let $R$ be a discrete valuation ring, $\fm$ its maximal ideal, $\pi$ a uniformizer, $k=R/\fm$ its residue field, assumed algebraically closed, $v$ its valuation, $f'_R:R\to R$ an endomorphism of $R$, and $G$ a finite group of automorphisms of $R$. Suppose that the endomorphism $f'_R$ and the automorphisms $g_R$ induce the identity on $k$, and suppose in addition that an endomorphism of $G$ is given,
\[
\varphi:G\longrightarrow G,
\]
satisfying
\[
f'_R\circ g_R=\varphi(g)_R\circ f'_R.
\]
If, for every $g\in G$, one writes
\[
\nu(g_R^{-1}f'_R)=v(g_R^{-1}f'_R(\pi)-\pi),
\]
and if $c^\varphi(g)$ (cf. 4.2) is decomposed in the form
\[
c^\varphi(g)=p^sq,
\qquad (p,q)=1,
\qquad p=\car(k),
\]
then $q\mid(\nu(g_R^{-1}f'_R)-1)$.}

\medskip
To prove Proposition 4.6 one reduces to the case $g=e$ and $\varphi=\id_G$. Indeed one first reduces to the case $g=e$ by associating with the situation $(f'_R,G,g,\varphi)$ the situation
\[
(f''_R=g_R^{-1}f'_R,
\; G,
\; e,
\; \varphi'=\mathrm{int}(g)\circ\varphi),
\]
where $\varphi'(a)=g^{-1}\varphi(a)g$, and by observing easily that
\[
\nu(g_R^{-1}f')-1=\nu(f''_R)-1
\]
and
\[
c^\varphi(g)=c^{\varphi'}(e)\qquad (G_g^\varphi=G_e^{\varphi'})\qquad (\text{cf. 4.2}).
\]
To reduce simultaneously to the case $g=e$, $\varphi=\id_G$, note that one can associate with the situation $(f''_R,G,e,\varphi')$ the situation
\[
(f''_R,G_e^{\varphi'},e,\id_{G_e^{\varphi'}}),
\]
and that
\[
G_e^{\varphi'}=(G_e^{\varphi'})_e^{\id_{G_e^{\varphi'}}}.
\]
The proof of Proposition 4.6 essentially uses the following lemma.

\medskip
\noindent
\textbf{Lemma 4.7.} \emph{With the hypotheses and notation of Proposition 4.6, suppose moreover that $\varphi=\id_G$, that is, that the operations of $G$ on $R$ commute with $f'_R$, and that $r=\nu(f'_R)\ge2$. Let
\[
c=\card(\im(G\longrightarrow \Aut_k(\fm/\fm^2))).
\]
Then the following divisibility relation holds:
\[
c\mid (r-1).
\]}

\medskip
\noindent
\textbf{Proof.} It is clear that the integer $r$ satisfies the following conditions:
\begin{enumerate}
\item[a)] the homomorphism
\[
f'_r:A/\fm^r\longrightarrow A/\fm^r
\]
induced by $f'_R$ is the identity;
\item[b)] the homomorphism
\[
f'_{r+1}:A/\fm^{r+1}\longrightarrow A/\fm^{r+1}
\]
is different from the identity.
\end{enumerate}
On the other hand, it is evident that $G$ acts on all the vector spaces $\fm^i/\fm^{i+1}$ over $k$.

We first show that under the hypotheses of the lemma there exists a unique homomorphism of $k[G]$-modules
\[
u(f'_R):\fm/\fm^2\longrightarrow \fm^r/\fm^{r+1}
\]
satisfying the two following conditions:
\begin{enumerate}
\item[1)] $u(f'_R)\ne0$;
\item[2)] $u(f'_R)(x\bmod \fm^2)=(f'_R(x)-x)\bmod \fm^{r+1}$, for every $x\in\fm$.
\end{enumerate}
It is clear that condition a) implies $f'_R(x)-x\in\fm^r$ for every $x\in R$, which gives meaning to condition 2). This observation also gives the possibility of defining a homomorphism of abelian groups
\[
\fm\longrightarrow \fm^r/\fm^{r+1},
\qquad
x\longmapsto (f'_R(x)-x)\bmod \fm^{r+1}.
\]
This homomorphism vanishes on $\fm^2$. Indeed, if $x_1,x_2\in\fm$, then $f'_R(x_1x_2)-x_1x_2\in\fm^{r+1}$ because, by condition a),
\[
f'_R(x_1)=x_1+\alpha_1,
\quad \alpha_1\in\fm^r,
\]
\[
f'_R(x_2)=x_2+\alpha_2,
\quad \alpha_2\in\fm^r,
\]
and hence
\[
f'_R(x_1x_2)=f'_R(x_1)f'_R(x_2)=(x_1+\alpha_1)(x_2+\alpha_2)
=x_1x_2+x_1\alpha_2+x_2\alpha_1+\alpha_1\alpha_2,
\]
and, plainly, $x_1\alpha_2+x_2\alpha_1+\alpha_1\alpha_2\in\fm^{r+1}$.

It follows by passage to the quotient that we get a homomorphism of abelian groups
\[
u(f'_R):\fm/\fm^2\longrightarrow \fm^r/\fm^{r+1}.
\]
The homomorphism $u(f'_R)$ is in fact a homomorphism of vector spaces over $k$. To see this, one must show that for all $x\in\fm$, $\alpha\in R$,
\[
f'_R(\alpha x)-\alpha x\equiv \alpha(f'_R(x)-x)\pmod{\fm^{r+1}},
\]
that is,
\begin{equation}
\tag{4.7.1}
f'_R(\alpha x)-\alpha f'_R(x)\in\fm^{r+1}.
\end{equation}
But
\[
f'_R(\alpha x)-\alpha f'_R(x)=(f'_R(\alpha)-\alpha)f'_R(x),
\]
which gives (4.7.1), since
\[
f'_R(\alpha)-\alpha\in\fm^r\quad\text{(condition a))},
\qquad
f'_R(x)\in\fm\quad(x\in\fm).
\]

The homomorphism $u(f'_R)$ is a homomorphism of $k[G]$-modules. Indeed,
for $x\in\fm$ one has
\[
\begin{aligned}
u(f'_R)\bigl(g(x\bmod\fm^2)\bigr)
&=(f'_R(gx)-gx)\bmod\fm^{r+1}\\
&=g(f'_R(x)-x)\bmod\fm^{r+1}
=g\,u(f'_R)(x\bmod\fm^2).
\end{aligned}
\]

Finally, $u(f'_R)\ne0$. Indeed, by condition b) there exists an element
$x_0\in R$ such that $f'_R(x_0)-x_0\notin\fm^{r+1}$. But we need such an
element in $\fm$. To obtain one, write $x_0$ in the form
\[
x_0=\frac{x}{x'},\qquad x,x'\in\fm,\quad x'\notin\fm^{r+1}.
\]
At least one of the elements $x,x'$ satisfies the required condition.

Suppose, for contradiction, that simultaneously
\[
f'_R(x)-x\in\fm^{r+1},
\qquad
f'_R(x')-x'\in\fm^{r+1}.
\]
Then, in the field of fractions $K$ of $A$, one has
\[
\frac{f'_R(x)}{x}\in1+\fm^r,
\qquad
\frac{f'_R(x')}{x'}\in1+\fm^r.
\]
We may suppose that $f'_R(x')\ne0$, since otherwise $x'$ is the desired
element. Since $1+\fm^r$ is a multiplicative group, it follows that
\[
\frac{x'}{f'_R(x')}\in1+\fm^r,
\]
and consequently
\[
\frac{f'_R(x)}{x}\cdot\frac{x'}{f'_R(x')}\in1+\fm^r.
\]
This last relation implies
\[
\frac{f'_R(x_0)}{x_0}\in1+\fm^r.
\]
Indeed,
\[
\begin{aligned}
\frac{f'_R(x_0)}{x_0}
-\frac{f'_R(x)}{x}\cdot\frac{x'}{f'_R(x')}
&=\frac{f'_R(\frac{x}{x'})}{\frac{x}{x'}}
-\frac{f'_R(x)}{x}\cdot\frac{x'}{f'_R(x')}\\
&=\frac{x'f'_R(\frac{x}{x'})}{x}
-\frac{f'_R(x)}{x}\cdot\frac{x'}{f'_R(x')}\\
&=\frac1x\left(
\frac{x'f'_R(x')f'_R(\frac{x}{x'})-x'f'_R(x)}
{f'_R(x')}\right)=0.
\end{aligned}
\]
Thus $f'_R(x_0)-x_0\in\fm^{r+1}$, a contradiction.

If the vector space
\[
\Hom_k(\fm/\fm^2,\fm^r/\fm^{r+1})
\]
is endowed with its usual $G$-module structure, its element $u(f'_R)$
belongs to the subspace
\[
\Hom_k(\fm/\fm^2,\fm^r/\fm^{r+1})^G
\]
of $G$-invariants.

It is classical that the vector space $\fm^i/\fm^{i+1}$ identifies
canonically with the $i$-th tensor power
$\bigotimes_k^i(\fm/\fm^2)$. Consequently
$\Hom_k(\fm/\fm^2,\fm^r/\fm^{r+1})$ identifies by the homomorphism of
contraction with the $(r-1)$-st tensor power
$\bigotimes_k^{r-1}(\fm/\fm^2)$:
\[
\gamma:\Hom_k(\fm/\fm^2,\fm^r/\fm^{r+1})
\longrightarrow\bigotimes_k^{r-1}(\fm/\fm^2).
\]
Through $\gamma$, the $G$-module structure on the Hom-space is
transported to the $G$-module structure on
$\bigotimes_k^{r-1}(\fm/\fm^2)$ induced by that on $\fm/\fm^2$:
\[
g(\xi_1\otimes\xi_2\otimes\cdots\otimes\xi_{r-1})
=g\xi_1\otimes\cdots\otimes g\xi_{r-1}.
\]
Denote by $\mu(f'_R)$ the image of $u(f'_R)$ under $\gamma$. Thus:
\begin{enumerate}
\item[I)] $0\ne\mu(f'_R)\in\bigotimes_k^{r-1}(\fm/\fm^2)$;
\item[II)] $\mu(f'_R)$ is invariant under $G$.
\end{enumerate}

We can now conclude the proof of Lemma 4.7. Since the vector space
$\fm/\fm^2$ has dimension one,
\[
\Aut_k(\fm/\fm^2)=k^*.
\]
Therefore the image of $G\to\Aut_k(\fm/\fm^2)$ is a finite cyclic
subgroup of $k^*$, generated by a primitive $c$-th root of unity
$\varepsilon$. It is generated by the automorphism
\[
\chi:\fm/\fm^2\longrightarrow\fm/\fm^2,
\qquad \xi\longmapsto\varepsilon\xi.
\]
Since $\mu(f'_R)$ is invariant under $G$,
\[
\begin{aligned}
\mu(f'_R)
&=\chi\mu(f'_R)\\
&=\chi\sum_i\xi_{1i}\otimes\xi_{2i}\otimes\cdots\otimes\xi_{r-1,i}\\
&=\sum_i\varepsilon\xi_{1i}\otimes\varepsilon\xi_{2i}
  \otimes\cdots\otimes\varepsilon\xi_{r-1,i}\\
&=\varepsilon^{r-1}\mu(f'_R),
\end{aligned}
\]
that is,
\[
(1-\varepsilon^{r-1})\mu(f'_R)=0.
\]
Since $\mu(f'_R)\ne0$, it follows that
$1-\varepsilon^{r-1}=0$, that is, $c\mid(r-1)$.

\medskip
\noindent
\textbf{Proof of Proposition 4.6.} We have already noted that one may reduce to the case $g=e$, $\varphi=\id_G$. If $G_1$ denotes the kernel of the homomorphism
\[
G\longrightarrow\Aut_k(\fm/\fm^2),
\]
then ramification theory gives
\[
\card(G_1)=p^t.
\]
Consequently, with the notation of Lemma 4.7,
\[
\card(G)=p^tc.
\]
Since, on the other hand, $c^{\id_G}(e)=\card(G)$, the assertion of Proposition 4.6 in the case $g=e$, $\varphi=\id_G$, and hence in general, follows from Lemma 4.7.

\medskip
\noindent
\textbf{Corollary 4.8.} \emph{For each $\ell\ne\car(k)$, the image of the rational number
\[
\frac{\nu(g_R^{-1}f'_R)-1}{c^\varphi(g)}
\]
under the canonical inclusion
\[
\bQ\hookrightarrow \bQ_\ell
\]
is an $\ell$-adic integer.}

\medskip
\noindent
\textbf{Proposition 4.9.} \emph{Conjecture 4.5 is true if the dimension of the scheme $X$ is equal to $1$.}

\medskip
\noindent
\textbf{Proof.} We introduce the endomorphism
\[
\psi:G\longrightarrow G
\]
defined by
\[
\psi=\mathrm{int}_g\circ\varphi,
\qquad \mathrm{int}_g(\xi)=g^{-1}\xi g.
\]
It is easy to see that:
\begin{enumerate}
\item[i)] $\psi(G_{x'})\subset G_{g^{-1}f'x'}$, where $G_{x'}$ is the stabilizer of the point $x'$;
\item[ii)] $c^\varphi(g)=\card(G^\psi)$, where $G^\psi=\{\xi\mid\psi(\xi)=\xi\}$;
\item[iii)] $G^\psi$ acts on ${X'_x}^{g^{-1}f'}$;
\item[iv)] $\nu_{x'}(g^{-1}f')=\nu_{hx'}(g^{-1}f')$ for every $h\in G^\psi$;
\item[v)] $(g^{-1}f')(g_1x)=\psi(g_1)(g^{-1}f')(x)$.
\end{enumerate}
Consequently one can decompose ${X'_x}^{g^{-1}f'}$ as the union of its orbits and obtain
\begin{equation}\tag{4.9.1}
\sum_{x'\in {X'_x}^{g^{-1}f'}}(\nu_{x'}(g^{-1}f')-1)
=
\sum_{x'\in {X'_x}^{g^{-1}f'}/G^\psi}
\frac{\nu_{x'}(g^{-1}f')-1}{\card(G^\psi_{x'})}\,\card(G^\psi),
\end{equation}
where
\[
\nu_{x'}(g^{-1}f')=\nu\bigl((g^{-1}f')_{\mathcal O_{x'}}\bigr),
\]
that is,
\begin{equation}\tag{4.9.2}
S_x(f',f,G,\varphi)(g)=\sum_{x'\in {X'}^{g^{-1}f'}/G^\psi}
\frac{\nu_{x'}(g^{-1}f')-1}{\card(G^\psi_{x'})}.
\end{equation}
Proposition 4.9 follows from Proposition 4.6.

\section*{5. Generalization of a Nielsen--Wecken type formula {\normalfont (cf. [1])}}

\begin{statement}{Proposition 5.1}
With the hypotheses and notation of 1, 2, 3, and 4, suppose in addition that $Y'$ is a finite subset of $X'$.

If $\ell$ is a prime number different from the characteristic of $k$, then:
\begin{enumerate}[label=\alph*)]
\item the values of the $\varphi$-central function
\[
\sum_{x\in Y^f}S_{x,\ell}(f,f',G,\varphi)
\]
are $\ell$-adic integers;
\item the following formula holds:
\begin{equation}\tag{5.1.1}
\Tr_{\bZ_\ell[G],\varphi}\bigl(f'^\vee_{H_{U',X'}(\ell)}\bigr)
=\sum_{x\in U^f}\nu_x(f)\mathrm{NW}^{G,\varphi}_f(x)
+\sum_{x\in Y^f}S_{x,\ell}(f,f',G,\varphi).
\end{equation}
\end{enumerate}
\end{statement}

\noindent
\textbf{Proof.} It is clear that if equality b) is proved with everything regarded as functions with values in $\bQ_\ell$, assertion a) follows at once, since $\mathrm{NW}^{G,\varphi}_{f}(x)$ and $\Tr_{\bZ_\ell[G],\varphi}(f'^\vee_{H_{U',X'}(\ell)})$ are functions with values in $\bZ_\ell$.

First, reducing modulo $\ell^n$ for every $n\ge1$ and using (XI, 7.2.3.6), one sees that for every $g\in G$,
\begin{equation}\tag{5.2}
\begin{aligned}
c^\varphi(g)\,\Tr_{\bZ_\ell[G],\varphi}(f'^\vee_{H_{U',X'}(\ell)})(g)
&=\Tr_{\bZ_\ell}(g_{X'}^{-1}f')^\vee_{H_{U',X'}(\ell)}\\
&=\Tr_{\bZ_\ell}(g_{X'}^{-1}f')_{H^\vee_{U',X'}(\ell)}.
\end{aligned}
\end{equation}
The commutative diagram (1.3), whose rows are exact, gives (XI, \S4)
\begin{equation}\tag{5.3}
\Tr_{\bZ_\ell}(g_{X'}^{-1}f')_{H^\vee_{U',X'}(\ell)}
=\Tr_{\bZ_\ell}(g_{X'}^{-1}f')_{R\Gamma_{X'}(\bZ_{\ell,X'})}
-\Tr^*_{\bZ_\ell}(g_{Y'}^{-1}f'_{Y'})_{R\Gamma_{Y'}(\bZ_{\ell,Y'})}.
\end{equation}
% SGA-SRC-000386: the scan leaves the citation blank; reconstructed from the mathematics.
The two terms on the right-hand side of (5.3) can be evaluated using the classical Lefschetz formula (cf. III, 4.8 and 4.10). (For the last term this is in fact a trivial Lefschetz formula):
\begin{equation}\tag{5.4}
\Tr_{\bZ_\ell}(g_{X'}^{-1}f')_{R\Gamma_{X'}(\bZ_{\ell,X'})}
=\sum_{x\in X^f}\sum_{x'\in {X'_x}^{g_{X'}^{-1}f'}}\nu_{x'}(g_{X'}^{-1}f'),
\end{equation}
\begin{equation}\tag{5.5}
\Tr^*_{\bZ_\ell}(g_{Y'}^{-1}f'_{Y'})_{R\Gamma_{Y'}(\bZ_{\ell,Y'})}
=\text{the number of fixed points of }g_{Y'}^{-1}f'_{Y'}.
\end{equation}
Consequently,
\begin{equation}\tag{5.6}
\Tr_{\bZ_\ell}(g_{X'}^{-1}f')_{H^\vee_{U',X'}(\ell)}
=\sum_{x\in U^f}\sum_{x'\in {X'_x}^{g_{X'}^{-1}f'}}\nu_{x'}(g_{X'}^{-1}f')
-\sum_{x\in Y^f}\sum_{x'\in {X'_x}^{g_{X'}^{-1}f'}}(1-\nu_{x'}(g_{X'}^{-1}f')).
\end{equation}
But for a given $x\in U^f$, the set ${X'_x}^{g_{X'}^{-1}f'}$ is nonempty if and only if $\mathrm{NW}_{f}^{G,\varphi}(x)=1$. Moreover,
\begin{equation}\tag{5.7}
\card {X'_x}^{g_{X'}^{-1}f'}=c^\varphi(g),
\qquad x\in U,
\end{equation}
\begin{equation}\tag{5.8}
\nu_{x'}(g_{X'}^{-1}f')=\nu_x(f),
\qquad x\in U.
\end{equation}
Relations (5.2) and (5.6) therefore give
\begin{equation}\tag{5.9}
\Tr_{\bZ_\ell[G],\varphi}(f'^\vee_{H_{U',X'}(\ell)})(g)
=\sum_{x\in U^f}\nu_x(f)\mathrm{NW}_f^{G,\varphi}(x)(g)
-\sum_{x\in Y^f}\sum_{x'\in {X'_x}^{g_{X'}^{-1}f'}}
\frac{1-\nu_{x'}(g_{X'}^{-1}f')}{c^\varphi(g)},
\end{equation}
which is the formula to be proved.

\subsection*{5.10. Special cases}
\subsubsection*{a) The unramified case: $Y'=\emptyset$, that is, $Y=\emptyset$.}
In this case formula (5.1.1) becomes
\begin{equation}\tag{5.11}
\Tr_{\bZ_\ell[G],\varphi}(f'^\vee_{H_{U',X'}(\ell)})
=
\sum_{x\in X^f}\nu_x(f)\mathrm{NW}_f^{G,\varphi}(x).
\end{equation}

\subsubsection*{b) The case of multiplicities $1$: $\nu_{x'}(g_{X'}^{-1}f')=1$ for every $x'\in {Y'}^{g_{Y'}^{-1}f'_{Y'}}$ and every $g\in G$.}
In this case formula 5.1.1 reduces to formula 5.11. This case occurs in particular if $\dim X=1$ and $\nu_x(f)=1$ for every $y\in Y$.

\section*{6. Application to a Lefschetz formula}

\subsection*{6.1}
We keep the notation and the conditions of number 1\footnote{$X^f$ denotes the set of closed fixed points of $X$.}, except where the contrary is expressly stated. Let $A$ be a commutative noetherian $\La_n$-algebra, let $F$ be a constructible sheaf, for the étale topology, of perfect $A$-modules on $X$, and let
\[
u:f^*(F)\longrightarrow F
\]
be a homomorphism of sheaves of $A$-modules.

In particular, one obtains an endomorphism of the pair $(X,F)$:
\[
(u,f):(X,F)\longrightarrow (X,F),
\]
and consequently an endomorphism of complexes of $A$-modules:
\[
R\Gamma_X(u,f):R\Gamma_X(F)\longrightarrow R\Gamma_X(F).
\]
Since the scheme $X$ is proper over $k$ by hypothesis and the sheaf $F$ is constructible and torsion, the finiteness theorem (SGA 4, XVII 5) ensures that the complex $R\Gamma_X(F)$ is a perfect complex of $A$-modules. The trace
\[
\Tr_A(R\Gamma_X(u,f))
\]
is therefore defined, with
\[
\Tr_A(R\Gamma_X(u,f))\in A.
\]
Our aim will be to give, under some additional hypotheses, an explicit formula for $\Tr_A(R\Gamma_X(u,f))$ in terms of local invariants associated with the fixed points of $f$.

\subsection*{6.2. Definition of the local terms}
Suppose that $X'$ is irreducible of dimension one and that the restriction to $U'$ of the sheaf $F'=p^*(F)$ is constant.

If the number of fixed points of $f'$ is finite, as will be assumed from now on, one may always suppose that $X^f\subset Y^f$.

The sheaf $F'|_{U'}$ is therefore defined by an $A$-module $M$ on which $G$ acts:
\[
F'|_{U'}\simeq M_{U'}.
\]
Indeed, if $\eta$ (resp. $\eta'$) denotes the generic point of $X$ (resp. $X'$), and $\bar\eta$ (resp. $\bar\eta'$) a geometric point over $\eta$ (resp. $\eta'$), then
\[
p(\eta')=\eta,\qquad F_{\bar\eta}\simeq F'_{\bar\eta'}.
\]
The complex $F'$ is isomorphic to $M$, and consequently $M$ is a perfect $A$-module.

On the other hand, if the morphism $f'$ is not constant, it fixes the generic point of $X'$. Consequently the morphism $(p^*(u),f')$ of the pair $(X',F')$ induces a morphism on the fibre $F'_{\bar\eta'}$,
\[
u'_{\bar\eta'}:F'_{\bar\eta'}\longrightarrow F'_{\bar\eta'},
\]
and hence, evidently, a morphism of $A$-modules
\[
u_M:M\longrightarrow M.
\]
One easily sees that the following relation is satisfied:
\begin{equation}\tag{6.2.1}
u_M(\varphi(g)m)=g u_M(m) \qquad \forall m\in M,\quad \forall g\in G.
\end{equation}
The dual $M^\vee$ of $M$ with respect to $A$ is therefore endowed with an $A[G]$-module structure and an endomorphism
\[
u_M^\vee:M^\vee\longrightarrow M^\vee
\]
such that
\[
u_M^\vee(gx)=\varphi(g)u_M^\vee(x) \qquad \forall g\in G,\quad \forall x\in M^\vee.
\]
Moreover, $M^\vee$ is a perfect $A$-module, and consequently the trace $\Tr_A(g^{-1}u_M^\vee)$ is defined for every $g\in G$.

Let $x\in X^f$ be a closed fixed point of $f$. The morphism $u$ induces an endomorphism of the module $F_x$, the fibre of the sheaf $F$ at the point $x$ (the closed points of $X$ are identified with the points of $X(x)$):
\[
u_x:F_x\longrightarrow F_x.
\]
The trace $\Tr_A(u_x)$ gives a first local invariant associated with the fixed point $x$.

We introduce another local invariant $\alpha_x(f,u,f',G,\varphi)$, an element of the ring $A$, by the formula
\begin{equation}\tag{6.2.2}
\alpha_x(f,u,f',G,\varphi)=
\sum_{g\in G_{\natural,\varphi}}
S_{x,\ell^n}(f,f',G,\varphi)(g)\Tr_A(g^{-1}u_M^\vee),
\end{equation}
where the function
\[
S_{x,\ell^n}(f,f',G,\varphi):G\longrightarrow\La_n
\]
is obtained from the function $i_\ell\circ S_x(f,F',G,\varphi)$ (cf. 4) by reduction modulo $\ell^n\bZ$, i.e.
\[
S_{x,\ell^n}(f,f',G,\varphi)(g)=
(i_\ell\circ S_x(f,f',G,\varphi))(g)\bmod \ell^n\bZ.
\]

\subsection*{6.3. A Lefschetz formula}
Under the hypotheses and with the notation of 5.1, 6.1, and 6.2, one has the formula
\begin{equation}
\tag{6.3.1}
\Tr_A(\RG_X(u,f))=\sum_{x\in X^f}\tau_x(f,u,f',G,\varphi),
\end{equation}
where
\begin{equation}
\tag{6.3.2}
\tau_x(f,u,f',G,\varphi)=\Tr_A(u_x)-\alpha_x(f,u,f',G,\varphi).
\end{equation}

\emph{Proof.} Since the subscheme $Y$ is stable under $f$, the morphism
\[
(u,f):(X,F)\longrightarrow(X,F)
\]
induces an endomorphism of distinguished triangles
\begin{equation}
\tag{6.3.3}
\begin{tikzpicture}[baseline=(current bounding box.center),>=stealth]
\node (Y) at (0,1.75) {$\RG_Y(F|Y)$};
\node (U) at (-2.45,0) {$\RG_X(F_{U,X})$};
\node (X) at (2.35,0) {$\RG_X(F)$};
\node at (2.75,2.25) {$\scriptstyle\RG_Y(u|Y,f|Y)$};
\node at (-5.05,0.02) {$\scriptstyle\RG_X((u,f)_{U,X})$};
\node at (5.05,0.02) {$\scriptstyle\RG_X(u,f)$};
\draw[->] (Y) to[out=145,in=35,looseness=5.8] (Y);
\draw[->] (U) to[out=210,in=-30,looseness=5.5] (U);
\draw[->] (X) to[out=25,in=-25,looseness=7] (X);
\draw[->] (Y) -- (U);
\draw[->] (U) -- (X);
\draw[->] (X) -- (Y);
\end{tikzpicture}
\end{equation}
and hence the relation
\begin{equation}
\tag{6.3.4}
\Tr_A(\RG_X(u,f))=
\Tr_A(\RG_X((u,f)_{U,X}))+
\Tr_A(\RG_Y(u|Y,f|Y)).
\end{equation}
The last term in (6.3.4) can be made explicit by the elementary Lefschetz formula and the fact that one may always suppose $X^f=Y^f$:
\begin{equation}
\tag{6.3.5}
\Tr_A(\RG_Y(u|_Y,f|_Y))=\sum_{x\in X^f}\Tr_A(u_x).
\end{equation}
On the other hand,
\begin{equation}
\tag{6.3.6}
\RG_X((u,f)_{U,X})\simeq \R\Gamma^G\RG_{X'}((u',f')_{U',X'}),\quad (u'=p^*(u)),
\end{equation}
\begin{equation}
\tag{6.3.7}
F_{U',X'}\simeq \La_{n,U',X'}\otimes_{\La_n}^{\mathbf L}M,
\qquad
(u',f')_{U',X'}=f'_{\La_{n,U',X'}}\otimes_{\La_n}M,
\end{equation}
whence, by the Künneth formula,
\begin{equation}
\tag{6.3.8}
\RG_{X'}(u',f')_{U',X'}\simeq
\RG_{X'}(f'_{\La_{n,U',X'}})\otimes_{\La_n}^{\mathbf L}u_M.
\end{equation}
It follows from (6.3.8) that
\begin{equation}
\tag{6.3.9}
\RG_{X'}((u',f')_{U',X'})^\vee\simeq
(\RG_{X'}(f'_{\La_{n,U',X'}}))^\vee\otimes_{\La_n}^{\mathbf L}A\otimes_A^{\mathbf L}u_M^\vee.
\end{equation}
Using formula 6.3.6, one obtains
\begin{equation}
\tag{6.3.10}
\RG_X((u,f)_{U,X})=\RG_{X'}(f'_{\La_{n,U',X'}})\otimes_{A[G]}^{\mathbf L}u_M^\vee,
\end{equation}
and hence, taking account of the general formula (cf. Exposé XI, 7.5.9),
\begin{equation}
\tag{6.3.11}
\Tr_A((u,f)_{U,X})=
\sum_{g\in G_{\natural,\varphi}}\Tr_{\La_n[G],\varphi}(f'{}^\vee_{H_{n,U',X'}})(g)\Tr_A(g^{-1}u_M^\vee).
\end{equation}
But the trace $\Tr_{\La_n[G],\varphi}(f'{}^\vee_{H_{n,U',X'}})$ is calculated using the Nielsen--Wecken formula (cf. Proposition 5.1 b)), evidently after reduction modulo $\ell^n\mathbb Z$.

Taking account of the relation $X^f=Y^f$, one obtains
\begin{equation}\tag{6.3.12}
\Tr_{\La_n[G],\varphi}(f'{}^\vee_{H_{n,U',X'}})(g)=\sum_{x\in X^f}S_{x,\ell^n}(f,f',G,\varphi)(g).
\end{equation}

Combining relations 6.3.12, 6.3.11, 6.3.5, and 6.3.4 gives the formula in the statement.

\subsection*{6.4. Remarks on the local terms}
\subsubsection*{a)} It is easy to see that, if $g\in G$, then one has the relation
\begin{equation}\tag{6.4.1}
\alpha_x(f,u,g^{-1}f',G,\varphi_g)
=
\alpha_x(f,u,f',G,\varphi),
\end{equation}
where
\begin{equation}\tag{6.4.2}
\varphi_g(\xi)=(\mathrm{int}_g\circ\varphi)(\xi)=g^{-1}\varphi(\xi)g,
\qquad \forall\xi\in G.
\end{equation}
Relation 6.4.1 allows us always to reduce, after replacing $f'$ by $g^{-1}f'$ and $\varphi$ by $\varphi_g$, to the case where there exists at least one point $x'\in {X'_x}^{f'}$.

\subsubsection*{b)} Let $x'\in X_x'{}^{f'}$ and let $G_{x'}$ be the stabilizer of the point $x'$. Evidently $G_{x'}$ is stable under $\varphi$, and hence one obtains an endomorphism
\[
\varphi_{x'}=\varphi|_{G_{x'}}:G_{x'}\longrightarrow G_{x'}.
\]
Let $X'/G_{x'}$ denote the quotient scheme obtained from $X'$ by the equivalence relation defined by $G_{x'}$, let $F'/G_{x'}$ be the endomorphism of $X'/G_{x'}$ obtained from the endomorphism $f'$ of $X'$ by the same equivalence relation, and let $x'$ denote the class of the point $x'$ modulo this equivalence relation.

The canonical projection
\[
p:X'\longrightarrow X'/G=X
\]
factors as
\[
p:X'\xrightarrow{\;p_1\;}X'/G_{x'}\xrightarrow{\;p_2\;}X=X'/G.
\]
Taking the inverse image by $p_2$ of the homomorphism $u$, one obtains a homomorphism
\[
v=p_2(u):(f'/G_{x'})^*(p_2^*(F))\longrightarrow p_2^*(F).
\]
It is easy to see that
\begin{equation}\tag{6.4.3}
\alpha_{x'}(f'/G_{x'},v,f',G_{x'},\varphi_{x'})
=
\alpha_x(f,u,f',G,\varphi)
\qquad (x=p(x')).
\end{equation}
Thus one may always reduce to the ``totally ramified'' case, that is, to the case where the fibre of the point $x$ is reduced to one point.

\subsubsection*{c) The ``tamely ramified'' case.} Suppose that there exists $x'\in p^{-1}(x)$ such that
\[
f'(x')=x',\qquad p(x')=x,
\]
a case to which one can always reduce, cf. case a), and suppose that the order of the group $G_{x'}$ is prime to the characteristic of the field $k$. Then one has
\begin{equation}\tag{6.4.4}
\alpha_x(f,u,f',G,\varphi)=S_{x,\ell^n}(f,f',G,\varphi)(e)\,\Tr_A(u_M^\vee).
\end{equation}
Taking case b) into account, one may suppose that $G=G_{x'}$. Formula (6.4.4) then follows from the following lemma.

\subsection*{Lemma 6.4.5}
\emph{If the following relations hold:}
\begin{align*}
\nu_{x'}(f')&>1,\\
\nu_{x'}(g^{-1}f')&>1,
\end{align*}
\emph{then $g$ belongs to the first ramification group $G_1$ (cf. 4.7):}
\[
g\in G_1,
\]
\emph{that is, $g$ induces the identity on the vector space $\fm_{x'}/\fm_{x'}^2$.}

\subsection*{Proof}
Let $\widehat{\mathcal O}_{x'}=k[[T]]$. The automorphism of $\fm/\fm^2$ induced by $g^{-1}$ is determined by a root of unity $\varepsilon(g^{-1})\in k$ by the condition
\[
g^{-1}T=\varepsilon(g^{-1})T\mod \fm^2,
\]
where $\fm$ is the maximal ideal of the ring $k[[T]]$.

On the other hand, the hypothesis $\nu_{x'}(f')>1$ gives
\[
f'(T)=T+a_\nu T^\nu+a_{\nu+1}T^{\nu+1}+\cdots,
\qquad \nu\ge2.
\]
Consequently
\[
g^{-1}f'(T)=\varepsilon(g^{-1})T+b_\nu T^\nu+b_{\nu+1}T^{\nu+1}+\cdots,
\]
and the hypothesis $\nu(g^{-1}f')>1$ implies $\varepsilon(g^{-1})=1$, as required.

\subsubsection*{d) The case of a sheaf $F$ which is not ramified at the fixed points of the endomorphism $f$.}
In this case the local term $\tau_x(f,u,f',G,\varphi)$ is given by the formula
\begin{equation}\tag{6.4.6}
\tau_x(f,u,f',G,\varphi)=\nu_x(f)\,\Tr_A(u_x),
\end{equation}
and consequently depends only on the pair $(f,u)$.

Indeed, in this case one has
\begin{equation}\tag{6.4.7}
S'_{x,\ell^n}(f,f',G,\varphi)(g)=
\begin{cases}
0&\text{if there exists }x'\in X'\text{ such that }(g^{-1}f')(x')=x'\text{ or}\\
&\text{equivalently if the $\varphi$-conjugacy class of $g$ coincides with}\\
&\text{the class }\mathrm{NW}_{f}^{G,\varphi}(x)\text{ (cf. 3); }\\
\nu_x(f)-1&\text{otherwise.}
\end{cases}
\end{equation}
It follows in particular that
\begin{equation}\tag{6.4.8}
S_{x,\ell^n}(f,f';G,\varphi)(g)=0
\qquad \text{if } g\notin \mathrm{NW}_{f}^{G,\varphi}(x).
\end{equation}
We therefore find
\begin{equation}\tag{6.4.9}
\begin{aligned}
\tau_x(f,u,f',G,\varphi)
&=\nu_x(f)\Tr_A(g^{-1}u_M^\vee)
=\nu_x(f)\Tr_A(g^{-1}u')_{x'}\\
&=\nu_x(f)\Tr_A(u_x).
\end{aligned}
\end{equation}

\section*{7. Comments on the conditions of validity of the Lefschetz formula}

\subsection*{7.1}
Of course, in practice one starts with a scheme $X$ proper over $k$ and an endomorphism
\[
(u,f):(X,F)\longrightarrow (X,F),
\]
such that:
\begin{itemize}
\item $F$ is a constructible sheaf of $A$-modules,
\item the set $X^f$ of fixed points of $f$ is finite.
\end{itemize}
The problem is to define, under these conditions, local terms $\tau_x(u,f)$, $x\in X^f$, such that
\[
\text{(*)}\qquad \Tr_A(R\Gamma_X(u,f))=\sum_{x\in X^f}\tau_x(u,f).
\]
Using the two duality theorems, global and local, Verdier succeeds in defining terms $\tau_x(u,f)$ such that formula (*) is true (cf. III). Unfortunately, the explicit determination of the $\tau_x(u,f)$, as defined by Verdier, in terms of known local invariants seems difficult; hence the interest of the definitions proposed in the preceding chapter.\footnote{See, however, III B for the comparison of the local terms of 6.3.1 with Verdier's local terms.}

By contrast, these latter definitions require additional hypotheses on $F$, $f$, and $u$, namely the existence of $U,U',G,f',\varphi$ having the properties required in 5.1, 6.1, 6.2.

\begin{statement}{Proposition 7.2}
Let $X$ be a smooth scheme of dimension $1$, let $F$ be constructible, and let
\[
(u,f):(X,F)\longrightarrow (X,F)
\]
be an endomorphism such that the morphism of sheaves
\[
u:f^*(F)\longrightarrow F
\]
is an isomorphism. Under these conditions there exist $U,U',G,f',\varphi$ satisfying the conditions of 5.1, 6.1, 6.2.
\end{statement}

\begin{prooftext}[Proof]
Let $U$ be the largest open subset of $X$ on which $F$ is locally constant, and let $Y$ denote its complement. The set $Y$ is stable under $f$:
\[
f(Y)\subset Y.
\]
Indeed, suppose by contradiction that there is a $y\in Y$ such that $f(y)\in U$. The sheaf $f^*(F)$ is then locally constant in a neighbourhood of $y$, and using the isomorphism $u$ it follows that $F$ is locally constant in a neighbourhood of $y$, which is absurd.

If $f$ is constant, the proof of the proposition is trivial. We shall therefore suppose that $f$ is not constant. If $\eta$ is the generic point of $X$, then $f(\eta)=\eta$.

Let $i_\eta:\Spec(k(\eta))\hookrightarrow X$ be the canonical inclusion, put $F_\eta=i_\eta^*(F)$, and let
\[
u_\eta:f_\eta^*(F_\eta)\longrightarrow F_\eta
\]
be the morphism, in fact the isomorphism, of sheaves induced by $u$, where $f_\eta$ is of course the restriction of $f$ to $\Spec(k(\eta))$.

Translating everything into terms of $H$-modules,
\[
H=\Gal(k(\eta)_s/k(\eta)),\qquad k(\eta)_s\text{ a separable closure of }k(\eta),
\]
one obtains the isomorphisms
\[
\begin{tikzcd}[column sep=large]
H \arrow[r,"\Psi"] & H
\end{tikzcd}
\qquad
\begin{tikzcd}[column sep=large]
F_{\bar\eta} \arrow[r,"u_{\bar\eta}"] & F_{\bar\eta},
\end{tikzcd}
\]
where $F_{\bar\eta}$ has an $H$-module structure and the following relation holds:
\[
u_{\bar\eta}(gx)=\Psi(g)\,u_{\bar\eta}x,
\qquad g\in H,
\quad x\in F_{\bar\eta}.
\]
If $H_1$ is the largest distinguished subgroup of $H$ of finite index such that $F_{\bar\eta}^{H_1}=F_{\bar\eta}$, then $\Psi(H_1)=H_1$. Such a subgroup exists because the sheaf $F_\eta$ is constructible.

Let $K\supset k(\eta)$ be the Galois extension of $k(\eta)$ corresponding to the subgroup $H_1$, and let $X'$ be the normalization of $X$ in $K$. It is clear that one can define an endomorphism
\[
f':X'\longrightarrow X'
\]
and an endomorphism
\[
H/H_1=G\xrightarrow{\varphi}G
\]
such that the system $U,U',f',G,\varphi$ satisfies the required conditions.
\end{prooftext}

\begin{statement}{Example 7.3}
Let $X$ be a prescheme defined over the finite field $\mathbb F_p$, where $p$ is a prime number, let $\fr_X:X\to X$ be the Frobenius endomorphism (cf. XIV), and let $F$ be a sheaf of $\Lambda$-modules on $X$.

If $Y,S$ are schemes over $\mathbb F_p$, and if $Y$ is a scheme over $S$, we denote by
\[
\Fr_{Y/S}:Y\longrightarrow Y^{(p/S)}
\]
the relative Frobenius morphism (cf. XIV 2).

We obtain a morphism of sheaves on $X_{\et}$
\[
(\fr_X)_*(F)\longrightarrow F
\]
by setting, for every étale $X$-prescheme $U$,
\[
(\fr_X)_*(F)(U)=F(\fr_X^{-1}(U))=F(U^{(p/X)})
 \xrightarrow{\,F(\Fr_{U/X})\,}F(U).
\]
This morphism is in fact an isomorphism. Its inverse
\[
F(\Fr_{/X})^{-1}:F\longrightarrow (\fr_X)_*(F)
\]
defines by adjunction a morphism
\[
\Fr_{F/X}^*:(\fr_X)^*(F)\longrightarrow F
\]
which is an isomorphism.

If in particular $X$ is smooth of dimension $1$ and $F$ is constructible, the pair
\[
(\Fr_{F/X}^*,\fr_X):(X,F)\longrightarrow (X,F)
\]
therefore falls under Proposition 7.2.
\end{statement}

\noindent[1] Wecken, F., Fixpunktklassen I, Math. Ann. 117 (1941), p. 659--671; Fixpunktklassen II, III, Math. Ann. 118 (1942--1943), p. 216--234 and p. 544--577.

\phantomsection\addcontentsline{toc}{section}{Exposé XV --- Frobenius and Rationality of $L$-Functions}
\section*{Exposé XV}
\begin{center}
\Large\bfseries FROBENIUS MORPHISM AND RATIONALITY OF $L$-FUNCTIONS

\medskip
by C. Houzel
\end{center}

\subsection*{§ 1. Frobenius morphism}

\subsubsection*{No. 1. Frobenius endomorphism}

\begin{statement}{Definition 1}
Let $p$ be a prime number. A prescheme $X$ is said to be of characteristic $p$ if $p\cdot 1_{\OO_X}=0$.
\end{statement}

This condition is equivalent to the following one: the unique morphism $X\to\Spec(\bbZ)$ factors through $\Spec(\bF_p)\to\Spec(\bbZ)$; moreover, such a factorization is unique.

We shall denote by $(\Sch)_p$ the full subcategory of the category of preschemes whose objects are the preschemes of characteristic $p$. By the preceding remark, the category $(\Sch)_p$ is naturally isomorphic to the category of $\bF_p$-preschemes.

If $X$ is a prescheme of characteristic $p$, the pair
\[
\fr_X=(\id_{|X|},\psi),
\qquad \psi:\OO_X\longrightarrow\OO_X,
\qquad \psi(f)=f^p
\]
for every section $f$ of $\OO_X$ over an open set, is an endomorphism of $X$, because $\psi$ is visibly an endomorphism of the sheaf of rings $\OO_X$.

\begin{statement}{Definition 2}
Let $X$ be a prescheme of characteristic $p$. The endomorphism $\fr_X$ thus defined is called the Frobenius endomorphism of $X$.
\end{statement}

The Frobenius endomorphism $\fr_X$ depends functorially on the prescheme $X$ of characteristic $p$; that is, for every morphism $g:X\to Y$ of preschemes of characteristic $p$, the following diagram is commutative:
\[
\begin{tikzcd}[column sep=large,row sep=large]
X \arrow[r,"\fr_X"] \arrow[d,"g"'] & X \arrow[d,"g"]\\
Y \arrow[r,"\fr_Y"'] & Y.
\end{tikzcd}
\]
In other words, we have defined an endomorphism $\fr$ of the identity functor $\id_{(\Sch)_p}$.

\subsubsection*{No. 2. Relative Frobenius morphism}

Let $S$ be a prescheme of characteristic $p$; all $S$-preschemes are then also of characteristic $p$. If $X$ is an $S$-prescheme, one sees that $\fr_X$ factors through the projection
\[
\pi_{X/S}:X\times_S(S,\fr_S)\longrightarrow X,
\]
where $(S,\fr_S)$ denotes $S$ regarded as an $S$-prescheme by means of the morphism $\fr_S$. We shall denote the fibre product $X\times_S(S,\fr_S)$ by $X^{(p/S)}$, or even by $X^{(p)}$ if no confusion can result. The unique $S$-morphism from $X$ to $X^{(p/S)}$ whose composite with $\pi_{X/S}$ is $\fr_X$ is denoted by $\Fr_{X/S}$; thus one obtains a commutative diagram in which the square is cartesian.
\[
\begin{tikzpicture}[baseline=(current bounding box.center),>=stealth]
\node (XL) at (0,2.0) {$X$};
\node (XM) at (3.0,2.0) {$X^{(p/S)}$};
\node (XR) at (6.0,2.0) {$X$};
\node (SL) at (0,0) {$S$};
\node (SM) at (3.0,0) {$S$};
\draw[->] (XM) -- node[below] {$\pi_{X/S}$} (XL);
\draw[->] (XR) -- node[below] {$\Fr_{X/S}$} (XM);
\draw[->] (XR) .. controls (5.1,3.2) and (0.9,3.2) .. node[above] {$\fr_X$} (XL);
\draw[->] (XL) -- node[left] {$g$} (SL);
\draw[->] (XM) -- node[right] {$g^{(p)}$} (SM);
\draw[->] (XR) -- node[pos=.55,above right] {$g$} (SM);
\draw[->] (SM) -- node[below] {$\fr_S$} (SL);
\end{tikzpicture}
\]

It is clear that $X^{(p/S)}$ depends functorially on the $S$-prescheme $X$; the functor
\[
p_S:X\longmapsto X^{(p/S)}
\]
is simply base change by the morphism $\fr_S:S\to S$. Like every base change, it commutes with projective limits and with sums, and preserves properties such as being of finite type, separated, affine, projective, flat, étale, and so on. Commutation with projective limits shows that if $X$ is endowed with an algebraic structure, for example a group or algebra structure, then $X^{(p/S)}$ is naturally endowed with an algebraic structure of the same kind.

Using the functoriality of $\fr_X$ with respect to $X$, one sees that the morphism $\Fr_{X/S}$ depends functorially on the $S$-prescheme $X$, so that a functorial morphism
\[
\Fr_{X/S}:\id_{(\Sch)/S}\longrightarrow p_S
\]
has been defined.

\begin{statement}{Definition 3}
Let $S$ be a prescheme of characteristic $p$ and $X$ an $S$-prescheme. The $S$-morphism
\[
\Fr_{X/S}:X\longrightarrow X^{(p/S)}
\]
is called the Frobenius morphism of $X$ relative to $S$.
\end{statement}

The first formal properties of the Frobenius morphism are stated in the following proposition.

\begin{statement}{Proposition 1}
\begin{enumerate}[label=\alph*)]
\item \emph{(Transitivity.)} Let $S$ be a prescheme of characteristic $p$, and let $g:X\to Y$ be an $S$-morphism of $S$-preschemes. The natural diagram relating $X$, $X^{(p/S)}$, $X^{(p/Y)}$, $Y$, $Y^{(p/S)}$, and $S$ is commutative and its three squares are cartesian.
\[
\begin{tikzcd}[column sep=huge,row sep=large]
X \arrow[d,"g"'] & X^{(p/S)} \arrow[l,"\pi_{X/S}"'] \arrow[d,"g^{(p/S)}"] & X^{(p/Y)} \arrow[l] \arrow[d,"g^{(p/Y)}"] & X \arrow[l,"\Fr_{X/Y}"'] \arrow[ll,bend left=35,"\Fr_{X/S}"{description,pos=0.62}] \arrow[dl,"g"] \\
Y \arrow[d] & Y^{(p/S)} \arrow[l,"\pi_{Y/S}"] \arrow[d] & Y \arrow[l,"\Fr_{Y/S}"] \arrow[dl] & {} \\
S & S \arrow[l,"\fr_S"] & {} & {}
\end{tikzcd}
\]

\item \emph{(Base change.)} Let $g:S'\to S$ be a morphism of preschemes of characteristic $p$, and let $X$ be an $S$-prescheme. There is a canonical isomorphism
\[
X^{(p/S)}\times_S S'\simeq (X\times_S S')^{(p/S')}
\]
functorial in $X$, that is, $(X^{(p/S)})'\simeq X'{}^{(p/S')}$ when the result of base change by $g$ is marked with a prime; moreover
\[
\Fr_{X'/S'}=(\Fr_{X/S})'
\]
modulo the identification allowed by the preceding isomorphism.
\end{enumerate}
\end{statement}

The proof is immediate.

\begin{statement}{Proposition 2}
Let $S$ be a prescheme of characteristic $p$ and $X$ an $S$-prescheme.
\begin{enumerate}[label=\alph*)]
\item The relative Frobenius morphism
\[
\Fr_{X/S}:X\longrightarrow X^{(p/S)}
\]
is integral, surjective, and radicial; hence it is a universal homeomorphism.

\item If $X\to S$ is an open immersion, then $X^{(p)}=X$ and $\Fr_{X/S}=\id_X$. If $X\to S$ is a closed immersion, identifying $X$ with a closed ringed subspace $(|X|,\OO_S/\JJ)$ of $S$, then
\[
X^{(p)}=(|X|,\OO_S/\JJ^p)
\]
and $\Fr_{X/S}=(\id_{|X|},\varphi)$, where $\varphi:\OO_S/\JJ^p\to\OO_S/\JJ$ is the evident augmentation.

\item Suppose that $X$ is locally of finite presentation over $S$.
\begin{enumerate}[label=\arabic*)]
\item The following three conditions are equivalent:
\[
\begin{array}{ll}
\text{(i)}& X\text{ is unramified over }S,\\
\text{(ii)}& \Fr_{X/S}\text{ is unramified},\\
\text{(iii)}& \Fr_{X/S}\text{ is a monomorphism}.
\end{array}
\]
\item Likewise, the following three conditions are equivalent:
\[
\begin{array}{ll}
\text{(i')}& X\text{ is étale over }S,\\
\text{(ii')}& \Fr_{X/S}\text{ is étale},\\
\text{(iii')}& \Fr_{X/S}\text{ is an isomorphism}.
\end{array}
\]
\end{enumerate}
\end{enumerate}
\end{statement}

We prove a). The composite $\pi_{X/S}\circ\Fr_{X/S}=\fr_X$ is visibly integral, surjective, and radicial; it follows that $\Fr_{X/S}$ is radicial (EGA I 3.5.6 (ii)). Moreover $\pi_{X/S}$ is separated and radicial, because it is deduced from $\fr_S$ by the base change $X\to S$; it follows that $\Fr_{X/S}$ is integral (EGA II 6.1.5 (v)) and surjective.

The proof of b) is evident, observing that
\[
\JJ\OO_{(S,\fr_S)}=\JJ^p
\]
for every ideal $\JJ$ of $\OO_S$.

Proof of c). It is known that for a morphism to be a monomorphism, respectively an isomorphism, it is necessary and sufficient that it be radicial and unramified, respectively radicial, étale, and surjective. This gives the equivalence of (ii) and (iii), respectively of (ii') and (iii'). It remains to prove that (i) is equivalent to (ii), and (i') to (ii').

For this, one uses the differential criterion of unramifiedness (SGA I 3) and the isomorphism
\[
\Omega^1_{X/S}\simeq \Omega^1_{X/X^{(p)}}
\]
given by the canonical exact sequence
\[
\Fr_{X/S}^*(\Omega^1_{X^{(p)}/S})\longrightarrow \Omega^1_{X/S}
\longrightarrow \Omega^1_{X/X^{(p)}}\longrightarrow 0
\]
(SGA II 4), in which the first arrow is visibly zero, because $\Omega^1_{X^{(p)}/S}$ is generated by the differentials of sections of $\OO_{X^{(p)}}$, and the differential $d_{X/S}$ vanishes on $(\OO_X)^p$ and on $\OO_S$. Thus (i) is equivalent to (ii).

Now suppose $X$ is étale over $S$; the same is true of $X^{(p)}$ by base change, and consequently the morphism $\Fr_{X/S}:X\to X^{(p)}$ is étale. Conversely, suppose that $\Fr_{X/S}$ is étale; it is unramified, hence $X$ is unramified over $S$. In other words, the structural morphism $g:X\to S$ is locally on $X$, for the étale topology, a closed immersion (SGA I, cor. 7.8). We may therefore suppose that $g$ is a closed immersion of finite presentation, and it remains to see that, if $\Fr_{X/S}$ is an isomorphism, then $g$ is an open immersion. Now $X$ is defined by an ideal $\JJ$ of finite type,
\[
X=(|X|,\OO_S/\JJ),
\]
and, by b),
\[
X^{(p)}=(|X|,\OO_S/\JJ^p),
\qquad
\Fr_{X/S}=(\id_{|X|},\OO_S/\JJ^p\to\OO_S/\JJ).
\]
To say that $\Fr_{X/S}$ is an isomorphism is therefore to say that $\JJ=\JJ^p$; since $\JJ$ is of finite type, Nakayama's lemma implies that $\JJ_s=0$, or $\JJ_s=\OO_{S,s}$, for every point $s\in S$, whence the conclusion.

\subsubsection*{No. 3. Effect on a module. Examples}

Let $S$ be a prescheme of characteristic $p$ and $X$ an $S$-prescheme. For every $\OO_X$-Module $\cE$, a sheaf for the Zariski topology, set
\[
\cE^{(p/S)}=\pi_{X/S}^*(\cE)=\cE\otimes_{\OO_X}\OO_{X^{(p)}}
 =\cE\otimes_{\OO_S}\OO_{(S,\fr_S)}.
\]
This is an $\OO_{X^{(p)}}$-Module, and one has
\[
\cE^{(p/X)}\simeq \Fr_{X/S}^*(\cE^{(p/S)}).
\]

The most interesting case is the case $X=S$; one may then write $\cE^{(p)}$ instead of $\cE^{(p/S)}$ without risk of confusion, and $\cE^{(p)}=\fr_S^*(\cE)$. The adjunction morphism
\[
\cE\longrightarrow \fr_{S*}\fr_S^*(\cE)
\]
is a morphism of abelian sheaves $\cE\to\cE^{(p)}$ which is not $\OO_S$-linear, but is ``$p$-linear'' in the evident sense; it is visibly a universal $p$-linear morphism from $\cE$ to an $\OO_S$-Module.

It is clear that $\cE^{(p)}$ depends functorially on $\cE$, and that the functor $\cE\mapsto\cE^{(p)}$ is right exact, but not exact in general; it commutes with tensor products:
\[
(\cE\otimes_{\OO_X}\cF)^{(p)}\simeq
\cE^{(p)}\otimes_{\OO_{X^{(p)}}}\cF^{(p)},
\]
and preserves properties such as being quasi-coherent, of finite type, of finite presentation, flat, locally free, and so on. In particular, if $\cL$ is an invertible sheaf on $S$, one immediately sees that $\cL^{(p)}$ is canonically isomorphic to $\cL^{\otimes p}$, because the morphism $\cL\to\cL^{\otimes p}$ defined by raising sections to their $p$-th tensor power is a universal $p$-linear morphism.

If $\mathcal{A}$ is an $\OO_S$-Algebra, then $\mathcal{A}^{(p)}$ has an $\OO_S$-Algebra structure deduced from that of $\mathcal{A}$. Suppose $\mathcal{A}$ is quasi-coherent, respectively graded and quasi-coherent; the same is true of $\mathcal{A}^{(p)}$, and
\[
\Spec(\mathcal{A}^{(p)})=\Spec(\mathcal{A}\otimes_{\OO_S}\OO_{(S,\fr_S)})
 \simeq \Spec(\mathcal{A})\times_S(S,\fr_S)=\Spec(\mathcal{A})^{(p)},
\]
respectively
\[
\Proj(\mathcal{A}^{(p)})=\Proj(\mathcal{A}\otimes_{\OO_S}\OO_{(S,\fr_S)})
 \simeq \Proj(\mathcal{A})\times_S(S,\fr_S)=\Proj(\mathcal{A})^{(p)}.
\]
The projection $\pi_{\Spec(\mathcal{A})/S}:\Spec(\mathcal{A})^{(p)}\to\Spec(\mathcal{A})$, respectively $\pi_{\Proj(\mathcal{A})/S}$, is identified with the morphism deduced from $\mathcal{A}\to\mathcal{A}^{(p)}$. On the other hand, $\Fr_{\Spec(\mathcal{A})/S}$, respectively $\Fr_{\Proj(\mathcal{A})/S}$, comes from a homomorphism of $\OO_S$-Algebras
\[
\mathcal{A}^{(p)}=\mathcal{A}\otimes_{\OO_S}\OO_{(S,\fr_S)}\longrightarrow \mathcal{A}
\]
which sends $a\otimes f$ to $a^p f$, for every section $a$ of $\mathcal{A}$ and every section $f$ of $\OO_S$. To see this, observe that the composite of this homomorphism with $\mathcal{A}\to\mathcal{A}^{(p)}$, on the right, is raising to the $p$-th power in $\mathcal{A}$.

Consider the particular case where $\mathcal{A}$ is the symmetric algebra $\Sym(\cE)$ of a quasi-coherent $\OO_S$-Module $\cE$; it is clear that $\mathcal{A}^{(p)}\simeq\Sym(\cE^{(p)})$, hence
\[
\mathbb V(\cE^{(p)})\simeq\Spec(\mathcal{A}^{(p)})\simeq\Spec(\mathcal{A})^{(p)}=\mathbb V(\cE)^{(p)},
\]
\[
\mathbb P(\cE^{(p)})\simeq\Proj(\mathcal{A}^{(p)})\simeq\Proj(\mathcal{A})^{(p)}=\mathbb P(\cE)^{(p)}.
\]
As for $\Fr_{\mathbb V(\cE)/S}$, respectively $\Fr_{\mathbb P(\cE)/S}$, it factors in the following way:
\[
\mathbb V(\cE)\longrightarrow \mathbb V(\Sym^p(\cE))\longrightarrow \mathbb V(\cE^{(p)}),
\]
respectively
\[
\mathbb P(\cE)\longrightarrow \mathbb P(\Sym^p(\cE))\longrightarrow \mathbb P(\cE^{(p)}),
\]
where the first arrow is defined by the canonical morphism
\[
\Sym(\Sym^p(\cE))\longrightarrow \Sym(\cE)
\]
which comes from $\Sym^p(\cE)\to\Sym(\cE)$, while the second is defined by the morphism of $\OO_S$-Modules
\[
\cE^{(p)}\longrightarrow \Sym^p(\cE)
\]
coming from the $p$-linear morphism $\cE\to\Sym^p(\cE)$ given by raising sections to their $p$-th symmetric power. In the case $X=\mathbb P(\cE)$, let us compute the fundamental sheaf $\OO_{X^{(p)}}(1)$:
\[
\OO_{X^{(p)}}(1)=(\Sym(\cE)^{(p)}(1))^{\sim}
=(\Sym(\cE)(1)^{(p)})^{\sim}
\simeq \OO_X(1)^{(p/S)}.
\]
When $\cE=\OO_S^n$, $\Sym(\cE)=\OO_S[T_1,\ldots,T_n]$ and $X=\mathbb V(\cE)=S[T_1,\ldots,T_n]$, respectively $X=\mathbb P(\cE)=\mathbb P_S^{n-1}$, it is clear that $\cE^{(p)}\simeq\OO_S^n=\cE$, so that $X^{(p)}=X$. The Frobenius morphism $\Fr_{X/S}$ comes from the endomorphism of the algebra $\OO_S[T_1,\ldots,T_n]$ which sends $T_i$ to $T_i^p$ $(i=1,\ldots,n)$.

Now consider a prescheme $X$, affine of finite type, respectively projective, over $S$, defined by an ideal, respectively a homogeneous ideal,
\[
\JJ\subset \OO_S[T_1,\ldots,T_n].
\]
Thus $X=\Spec(\mathcal{A})$, respectively $X=\Proj(\mathcal{A})$, where
\[
\mathcal{A}=\OO_S[T_1,\ldots,T_n]/\JJ.
\]
One sees that
\[
\mathcal{A}^{(p)}
\simeq \OO_S[T_1,\ldots,T_n]/\JJ',
\]
where $\JJ'$ is the image in $\OO_S[T_1,\ldots,T_n]$ of
$\JJ^{(p)}=\JJ\otimes_{\OO_S}\OO_{(S,\fr_S)}$; the ideal $\JJ'$ is
deduced from $\JJ$ by raising the coefficients of the polynomials to
the $p$-th power. One has $X^{(p)}=\Spec(\mathcal{A}^{(p)})$, respectively $\Proj(\mathcal{A}^{(p)})$, and the Frobenius morphism
\[
\Fr_{X/S}:X\longrightarrow X^{(p)}
\]
is defined by raising the coordinates $T_i$ to the $p$-th power; this determines a morphism $\mathcal{A}^{(p)}\to\mathcal{A}$.


\subsection*{4. Iterated Frobenius morphism; perfect preschemes}

For every integer $\nu\geq0$, the iterated endomorphism $\fr_X^\nu$ depends functorially on the prescheme $X$ of characteristic $p$. Thus, if $S$ is a prescheme of characteristic $p$ and $X$ is an $S$-prescheme, $\fr_X^\nu$ factors, in a unique way, as an $S$-morphism
\[
\Fr_{X/S}^\nu:X\longrightarrow X^{(p^\nu/S)},
\]
where
\[
X^{(p^\nu/S)}=X\times_S(S,\fr_S^\nu),
\]
followed by the projection $X^{(p^\nu/S)}\to X$. The morphism $\Fr_{X/S}^\nu$ depends functorially on the $S$-prescheme $X$ and is compatible with change of the base $S$ (cf. no. 2, prop. 1b)). Moreover it is clear that
\[
X^{(p^\nu/S)}=X^{(p^\nu)}\simeq\bigl(X^{(p^{\nu-1})}\bigr)^{(p)}
\]
and that
\[
\Fr_{X/S}^\nu=(\Fr_{X/S}^{\nu-1})^{(p)}\circ\Fr_{X/S},
\]
modulo the identification allowed by the preceding canonical isomorphism.

\begin{statement}{Definition 4}
A prescheme $S$ of characteristic $p$ is said to be perfect if $\fr_S$ is an automorphism of $S$.
\end{statement}

If $S$ is a perfect prescheme of characteristic $p$ and $X$ is an $S$-prescheme, one sees immediately that $X^{(p/S)}$ is identified with the prescheme $X$ considered as an $S$-prescheme by means of the morphism $\fr_S^{-1}\circ g:X\to S$, where $g$ is the structural morphism of $X$ into $S$; modulo this identification, $\Fr_{X/S}=\fr_X$.

As an example of a perfect prescheme, we shall meet $S=\Spec(\mathbb F_q)$, where $q=p^\nu$ is a power of the characteristic $p$; indeed, in this case $\fr_S^\nu=\id_S$. If $X$ is an $S$-prescheme, one sees that
\[
X^{(p^\nu/S)}=X,
\qquad
\Fr_{X/S}^\nu=\fr_X^\nu.
\]
Note that for $\nu=1$ one obtains $S=\Spec(\mathbb F_p)$, for which $\fr_S=\id_S$; the $S$-preschemes are the preschemes of characteristic $p$; over this base $S$, $X^{(p/S)}=X$ and $\Fr_{X/S}=\fr_X$.

\section*{§ 2. Frobenius correspondence}

\subsection*{1. Definition}

Let $X$ be a prescheme of characteristic $p$, and let $F$ be a sheaf of sets on $X_{\et}$. For every étale $X$-prescheme $U$ one has
\[
(\fr_X)_*(F)(U)=F(\fr_X^{-1}(U))=F(U^{(p/X)}),
\]
and the Frobenius morphism $\Fr_{U/X}:U\to U^{(p/X)}$ defines a map
\[
F(\Fr_{U/X}):\fr_{X*}(F)(U)\longrightarrow F(U),
\]
which is visibly functorial in $U$ (since $\Fr_{U/X}$ is functorial in $U$). Thus one obtains a morphism of sheaves on $X_{\et}$:
\[
(\fr_X)_*(F)\longrightarrow F.
\]
Now, when $U$ is étale over $X$, $\Fr_{U/X}$ is an isomorphism (§ 1 no. 2, prop. 2c)), hence the same is true of $F(\Fr_{U/X})$; it follows that the preceding morphism is an isomorphism. In the case where $F$ is endowed with an algebraic structure, for example if it is a sheaf of $\Lambda$-modules, where $\Lambda$ is a given base ring, this isomorphism is compatible with the algebraic structure. The inverse
\[
F(\Fr_{/X})^{-1}:F\longrightarrow \fr_{X*}(F)
\]
of this isomorphism defines, since $(\fr_X)^*$ is left adjoint to the functor $(\fr_X)_*$, a morphism
\[
\Fr^*_{F/X}:(\fr_X)^*(F)\longrightarrow F.
\]
One can even say more: since $\fr_X$ is a universal homeomorphism (§ 1 no. 2, prop. 2a)), the functors $(\fr_X)^*$ and $(\fr_X)_*$ are quasi-inverse to each other; consequently $\Fr^*_{F/X}$ is an isomorphism, just as is the morphism $F(\Fr_{/X})^{-1}$ from which it comes.

\begin{statement}{Definition 1}
Given a prescheme $X$ of characteristic $p$ and a sheaf of sets $F$ on $X_{\et}$, the Frobenius correspondence on $(X,F)$ is the pair $(\fr_X,\Fr^*_{F/X})$ formed by the Frobenius endomorphism $\fr_X$ of $X$ and the isomorphism
\[
\Fr^*_{F/X}=(F(\Fr_{/X})^{-1})^\sharp:(\fr_X)^*(F)\longrightarrow F.
\]
\end{statement}

Of course one defines in the same way a class of iterated Frobenius correspondence $(\fr_X^\nu,(\Fr^*_{F/X})^\nu)$ for every integer $\nu\geq0$; the isomorphism
\[
(\Fr^*_{F/X})^\nu:(\fr_X^\nu)^*(F)\longrightarrow F
\]
is equal to
\[
(\Fr^*_{F/X})^{\nu-1}\circ (\fr_X^{\nu-1})^*(\Fr^*_{F/X})
\]
or also to
\[
\Fr^*_{F/X}\circ\fr_X^*((\Fr^*_{F/X})^{\nu-1})
\]
for $\nu\geq1$.

Let us examine the particular case where $F$ is representable by an étale prescheme $V$ over $X$; then $(\fr_X)^*(F)=V^{(p/X)}$, $F(V)=\Hom_X(V,V)$,
\[
(\fr_X)^*(F)(V)=F(V^{(p/X)})=\Hom_X(V^{(p/X)},V),
\]
and $F(\Fr_{V/X})$ is the map from $\Hom_X(V^{(p/X)},V)$ to $\Hom_X(V,V)$ which sends a morphism $\psi$ to $\psi\circ\Fr_{V/X}$. Thus $\Fr^*_{F/X}$ is identified with a morphism from $V^{(p/X)}$ to $V$, which is the image of $\id_V$ under the inverse map $F(\Fr_{V/X})^{-1}$ of the preceding one. This shows that for $F=V$ étale over $X$ one has
\[
\Fr^*_{V/X}=(\Fr_{V/X})^{-1}.
\]

Note that if the sheaf $F$ is constant, then it is clear that
\[
(\fr_X)^*(F)=F,
\qquad
\Fr^*_{F/X}=\id_F.
\]

\begin{statement}{Proposition 1}
\begin{enumerate}[label=\alph*)]
\item The isomorphism $\Fr^*_{F/X}$ depends functorially on the sheaf $F$; in other words there is a functorial isomorphism
\[
\Fr^*_{/X}:\fr_X^*\longrightarrow \id_{\widetilde X_{\et}},
\]
where $\widetilde X_{\et}$ is the category of sheaves of sets on $X_{\et}$.
\item The Frobenius correspondence $(\fr_X,\Fr^*_{F/X})$ on $(X,F)$ is compatible with change of the base $X$. That is, if $\mathcal F$ denotes the fibred category over $(\Sch)_p$ of sheaves, for the étale topology, on variable preschemes of characteristic $p$, the functors $\fr_X^*$, $X\in\Ob(\Sch)_p$, define a cartesian $(\Sch)_p$-functor
\[
\fr^*: \mathcal F\longrightarrow\mathcal F,
\]
and the morphisms $\Fr^*_{F/X}$, $X\in\Ob(\Sch)_p$, $F\in\Ob\widetilde X_{\et}$, define a functorial $(\Sch)_p$-morphism
\[
\Fr^*_{/}:\fr^*\longrightarrow\id_{\mathcal F}
\]
which is an isomorphism.
\item Conversely, these properties characterize $\Fr^*_{/}$; that is, there exists a unique functorial $(\Sch)_p$-morphism from $\fr^*$ to $\id_{\mathcal F}$.
\end{enumerate}
\end{statement}

Notice that statement b) implies statement a); the latter is entirely evident.

Let us prove b). Let $g:X'\to X$ be a morphism of preschemes of characteristic $p$, and let $F$ be a sheaf on $X_{\et}$. Put $F'=g^*(F)$ and, more generally, mark by a prime the inverse image by $g$ of any object given over $X$. Since
\[
g\circ\fr_{X'}=\fr_X\circ g,
\]
one has
\[
\fr_{X'}^*(F')=\fr_{X'}^*(g^*F)=g^*(\fr_X^*F)=(\fr_X^*F)',
\]
whence the fact that the $\fr_X^*$ define a cartesian functor $\fr^*$ over $(\Sch)_p$. It remains to prove that
\[
\Fr^*_{F'/X'}=(\Fr^*_{F/X})'.
\]
Using the properties of adjoint functors, one sees that this amounts to showing that the image of $F'(\Fr_{/X'})^{-1}$ under the bijection
\[
\Hom_{X'}(F',\fr_{X'*}(F'))
 =\Hom_{X'}(g^*F,\fr_{X'*}g^*F)
 \simeq \Hom_X(F,g_*\fr_{X'*}g^*F)
\]
(with $g^*$ left adjoint to $g_*$) coincides with the image of $F(\Fr_{/X})^{-1}$ under the map
\[
\Hom_X(F,\fr_{X*}F)\longrightarrow
\Hom_X(F,\fr_{X*}g_*g^*F)=
\Hom_X(F,g_*\fr_{X'*}g^*F)
\]
defined by the adjunction morphism $F\xrightarrow{\alpha}g_*g^*F$; in other words, one must show that
\[
\alpha\circ F(\Fr_{/X})=g_*(F'(\Fr_{/X'}))\circ\fr_{X*}(\alpha),
\]
This says that for every prescheme $U$ étale over $X$ there is a commutative diagram
\[
\begin{tikzcd}[column sep=large,row sep=large]
F(U^{(p/X)}) \arrow[r,"F(\Fr_{U/X})"] \arrow[d]
& F(U) \arrow[d]\\
F'(U'{}^{(p/X')}) \arrow[r,"F'(\Fr_{U'/X'})"']
& F'(U') ,
\end{tikzcd}
\]
where the vertical arrows are the canonical maps; but this follows immediately from
\[
\Fr_{U'/X'}=(\Fr_{U/X})'
\]
(§ 1, no. 2, prop. 1b)).

To prove c), after composing with $(\Fr^*_{/})^{-1}$, we are reduced to proving that $\id_{\mathcal F}$ has only one $(\Sch)_p$-endomorphism, namely the identity. If $\psi$ is a $(\Sch)_p$-endomorphism of $\id_{\mathcal F}$, let $F$ be a sheaf on a prescheme $X$ of characteristic $p$, and let $U$ be an étale prescheme over $X$. We know that $F(U)$ is identified with $\Hom_U(U,F|U)$, and $\psi_F(U):F(U)\to F(U)$ with the map
\[
s\longmapsto \psi_{F|U}\circ s=s\circ\psi_U,
\]
where $s\in\Hom_U(U,F|U)$; since the final sheaf $U$ of $\widetilde U_{\et}$ has no endomorphism other than the identity, one has $\psi_U=\id_U$. It follows that $\psi_F(U)$ is also the identity of $F(U)$, and hence that $\psi=\id_{\mathcal F}$.

Of course, if the sheaf $F$ is endowed with an algebraic structure, the isomorphism $\Fr^*_{F/X}$ is compatible with this structure; for example, if $F$ is a sheaf of groups on the prescheme $X$ of characteristic $p$, $\Fr^*_{F/X}$ is an isomorphism of sheaves of groups. Let $\Lambda$ be a ring and let $F$ be a sheaf of $\Lambda$-modules on $X$; then $\Fr^*_{F/X}$ is an isomorphism of sheaves of $\Lambda$-modules. More generally, if $F^\bullet$ is a complex of $\Lambda$-modules, one immediately defines an isomorphism
\[
\Fr^*_{F^\bullet/X}:\fr_X^*(F^\bullet)\longrightarrow F^\bullet
\]
which is equal to $\Fr^*_{F^i/X}$ on the component of degree $i$, thanks to proposition 1a). In this case the pair $(\fr_X,\Fr^*_{F^\bullet/X})$ is indeed a correspondence, in the usual sense, on $(X,F^\bullet)$; this correspondence is functorial in $F^\bullet$ and compatible with change of the base $X$ (Proposition 1a), b)).

\subsection*{2. Effect on cohomology}

\begin{statement}{Proposition 2}
Let $X$ be a prescheme of characteristic $p$.
\begin{enumerate}[label=\alph*)]
\item For every sheaf of sets $F$ on $X_{\et}$, the endomorphism of $H^0(X,F)$ defined by $(\fr_X,\Fr^*_{F/X})$, that is, the composite
\[
H^0(X,F)\longrightarrow H^0(X,\fr_X^*(F))\longrightarrow H^0(X,F),
\]
where the first arrow is the canonical map arising from the contravariant character of $H^0(X,F)$ with respect to $X$, and the second is $H^0(X,\Fr^*_{F/X})$, is the identity of $H^0(X,F)$.
\item For every sheaf of groups $F$ on $X_{\et}$, the endomorphism of $H^1(X,F)$ defined by $(\fr_X,\Fr^*_{F/X})$, that is,
\[
H^1(X,F)\longrightarrow H^1(X,\fr_X^*F)\longrightarrow H^1(X,F),
\]
is the identity.
\item For every left bounded complex $F$ of $\Lambda$-modules on $X_{\et}$, the endomorphism of $\RGamma_X(F)$ defined by $(\fr_X,\Fr^*_{F/X})$, that is, the composite
\[
\RGamma_X(F)\longrightarrow \RGamma_X(\fr_X^*F)\longrightarrow \RGamma_X(F),
\]
is the identity.
\end{enumerate}
\end{statement}

We first establish a); b) and c) will be deduced from it. Let $\psi_F$ be the considered endomorphism of $H^0(X,F)$. It is clear that $\psi_F$ depends functorially on $F$ and that $\psi_X=\id_{H^0(X,X)}$. Every element $s$ of $H^0(X,F)$ is identified with a morphism of sheaves $s:X\to F$ such that $s$ is the image under $H^0(X,s)$ of the unique element of $H^0(X,X)$; since
\[
\psi_F\circ H^0(X,s)=H^0(X,s)\circ\psi_X=H^0(X,s),
\]
by the preceding remarks, one sees that $\psi_F(s)=s$ for every $s$, hence $\psi_F=\id_{H^0(X,F)}$.

Note that $\psi_F$ is also the composite
\[
H^0(X,F)\longrightarrow H^0(X,\fr_{X*}F)
\longrightarrow H^0(X,F),
\]
where the first arrow is $H^0(X,F(\Fr_{/X})^{-1})$ and the second is the canonical map, which is nothing but the identity. Similarly, the endomorphisms considered in b) and c) factor as
\[
H^1(X,F)\longrightarrow H^1(X,\fr_{X*}F)
\longrightarrow H^1(X,F),
\]
and
\[
\RGamma_X(F)\longrightarrow \RGamma_X(\fr_{X*}F)
\longrightarrow \RGamma_X(F).
\]

To establish b), we shall use a monomorphism $F\to G$ of $F$ into a sheaf of groups $G$ such that $H^1(X,G)=0$. Then $\fr_{X*}F\to\fr_{X*}G$ is a monomorphism and $H^1(X,\fr_{X*}G)=0$. Introduce the quotient homogeneous sheaf of sets $C=G/F$. Since $\fr_X$ is integral, surjective, and radicial, the functors $\fr_X^*$ and $\fr_{X*}$ are quasi-inverse equivalences of categories (SGA VIII theorem 1.1). It follows that
\[
\fr_{X*}(C)\simeq\fr_{X*}(G)/\fr_{X*}(F),
\]
and the exact cohomology sequence (SGA XII proposition 3.1) shows that the considered endomorphism of $H^1(X,F)$ arises, by passage to the quotient, from
\[
\psi_C:H^0(X,C)\longrightarrow H^0(X,C),
\]
which is the identity by a). This endomorphism is therefore $\id_{H^1(X,F)}$.

Assertion c) is proved analogously, using a complex $I$ that is injective in each degree and isomorphic to $F$ in the derived category $D^+$, built from the $\Lambda$-modules. The complex $\fr_{X*}(I)$ is injective in each degree and isomorphic to $\fr_{X*}(F)$ in $D^+$, since $\fr_{X*}$ is an equivalence of categories. One may then identify $\RGamma_X(F)$ with $H^0(X,I)$ and $\RGamma_X(\fr_{X*}F)$ with $H^0(X,\fr_{X*}I)$; one sees that the considered endomorphism of $\RGamma_X(F)$ is identified with $\psi_I$, which is the identity. This endomorphism is therefore $\id_{\RGamma_X(F)}$.

\subsection*{3. Behavior of products}

Let $X$ and $X'$ be preschemes of characteristic $p$. Their product $X\times X'$ is identical with the fibred product $X\times_eX'$, where $e=\Spec(\mathbb F_p)$. It is clear that
\[
\fr_{X\times X'}=\fr_X\times\fr_{X'}=(\fr_X\times\id_{X'})\circ(\id_X\times\fr_{X'})
=(\id_X\times\fr_{X'})\circ(\fr_X\times\id_{X'}).
\]
It follows immediately that, for every $X$-prescheme $Y$ and every $X'$-prescheme $Y'$, there is a canonical isomorphism
\[
(Y\times Y')^{(p/X\times X')}
\simeq Y^{(p/X)}\times Y'{}^{(p/X')};
\]
modulo this isomorphism one can write
\[
\Fr_{Y\times Y'/X\times X'}=\Fr_{Y/X}\times\Fr_{Y'/X'}.
\]

\begin{statement}{Proposition 3}
Let $X$ and $X'$ be preschemes of characteristic $p$, let $\Lambda$ be a commutative ring, let $F$ be a sheaf of $\Lambda$-modules on $X_{\et}$, and let $F'$ be a sheaf of $\Lambda$-modules on $X'_{\et}$. One has
\[
\begin{aligned}
\Fr^*_{F\boxt_\Lambda F'/X\times X'}
&=\Fr^*_{F/X}\otimes\Fr^*_{F'/X'}\\
&=(\Fr^*_{F/X}\otimes\id_{\fr_{X'}^*F'})
  \circ(\id_F\otimes\Fr^*_{F'/X'})\\
&=(\id_{\fr_X^*F}\otimes\Fr^*_{F'/X'})
  \circ(\Fr^*_{F/X}\otimes\id_{F'}).
\end{aligned}
\]
\end{statement}

To write these equalities, one has identified
\[
\fr_{X\times X'}^*(F\boxt_\Lambda F')
\]
with the tensor product $\fr_X^*(F)\otimes_\Lambda\fr_{X'}^*(F')$ by means of the canonical isomorphism.

The two sides of the first equality depend functorially on $F$ and $F'$, and are compatible with change of the base $X$ or the base $X'$. Moreover they are compatible with passage to the inductive limit in $F$ or in $F'$, since the tensor product commutes with inductive limits and it is clear that
\[
\Fr^*_{(\varinjlim F_i)/X}=\varinjlim\Fr^*_{F_i/X};
\]
this permits reduction to the case where $F=\Lambda_X$ and $F'=\Lambda_{X'}$, in which case it is evident.

\begin{statement}{Corollary}
With the data of proposition 3, the endomorphisms of $\RGamma_{X\times X'}(F\boxt F')$ defined respectively by
\[
(\id_X\times\fr_{X'},\,\id_F\otimes\Fr^*_{F'/X'})
\]
and by
\[
(\fr_X\times\id_{X'},\,\Fr^*_{F/X}\otimes\id_{F'})
\]
are inverse isomorphisms of each other.
\end{statement}

Denoting these endomorphisms by $\psi'$ and $\psi$, let us show for example that $\psi\circ\psi'=\id$. This follows from the commutative diagram
\[
\resizebox{\textwidth}{!}{%
\begin{tikzpicture}[baseline=(current bounding box.center),>=stealth,font=\scriptsize]
\node (A) at (0,3.6) {$\RGamma_{X\times X'}(F\boxt F')$};
\node (B) at (4.7,3.6) {$\RGamma_{X\times X'}(F\boxt\fr_{X'}^*(F'))$};
\node (C) at (9.4,3.6) {$\RGamma_{X\times X'}(F\boxt F')$};
\node (D) at (3.3,1.9) {$\RGamma_{X\times X'}(\fr^*_{X\times X'}(F\boxt F'))$};
\node (E) at (9.4,1.9) {$\RGamma_{X\times X'}(\fr_X^*(F)\boxt F')$};
\node (F) at (4.4,0.25) {$\RGamma(\Fr^*_{F\boxt F'/X\times X'})$};
\node (G) at (9.4,0.25) {$\RGamma_{X\times X'}(F\boxt F')$};
\draw[->] (A) -- node[above] {$\id_X\times\fr_{X'}$} (B);
\draw[->] (B) -- node[above] {$\RGamma(\id_F\otimes\Fr^*_{F'/X'})$} (C);
\draw[->] (A) -- node[below left] {$\fr_{X\times X'}$} (D);
\draw[->] (B) -- node[left] {$\fr_X\times\id_{X'}$} (D);
\draw[->] (C) -- node[right] {$\fr_X\times\id_{X'}$} (E);
\draw[->] (D) -- node[above] {$\RGamma(\id_{\fr_X^*(F)}\otimes\Fr^*_{F'/X'})$} (E);
\draw[->] (E) -- node[right] {$\RGamma(\Fr^*_{F/X}\otimes\id_{F'})$} (G);
\draw[->] (D) -- (G);
\end{tikzpicture}%
}
\]
which is commutative because
\[
\fr_{X\times X'}=(\fr_X\times\id_{X'})\circ(\id_X\times\fr_{X'})
\]
for the upper triangle and because
\[
\Fr^*_{F\boxt F'/X\times X'}=(\Fr^*_{F/X}\otimes\id_{F'})\circ(\id_{\fr_X^*F}\otimes\Fr^*_{F'/X'})
\]
for the lower triangle, and in which the composite of the two oblique arrows is the identity by proposition 2c). With an analogous diagram exchanging the roles of $(X,F)$ and $(X',F')$, one obtains $\psi'\circ\psi=\id$.

We shall apply this result to the particular case where
$X'=\bar e=\Spec(\overline{\mathbb F}_p)$ is the spectrum of an
algebraic closure $\overline{\mathbb F}_p$ of $\mathbb F_p$ and where
$F'=\Lambda_{\bar e}$ is the constant sheaf with fibre $\Lambda$.
In this case put
\[
\overline X=X\times_e\bar e
=X\otimes_{\mathbb F_p}\overline{\mathbb F}_p,
\qquad
\overline F=F\otimes_\Lambda\Lambda_{\bar e},
\]
the inverse image of $F$ by the projection $\overline X\to X$.
The Frobenius endomorphism $\fr_{\bar e}$ is identified with the
canonical generator
\[
f:\lambda\longmapsto\lambda^p
\qquad(\lambda\in\overline{\mathbb F}_p)
\]
of the Galois group
\[
\pi_e=\Gal(\overline{\mathbb F}_p/\mathbb F_p),
\]
where $f$ determines an isomorphism
$\widehat{\mathbb Z}\simeq\pi_e$, and
\[
\Fr^*_{\Lambda_{\bar e}/\bar e}=\id_{\Lambda_{\bar e}}
\]
since $\Lambda_{\bar e}$ is a constant sheaf. Put
\[
\overline{\fr}_X=\fr_X\times\id_{\bar e},
\]
the inverse image of $\fr_X$ by $\overline X\to X$,
\[
f_X=\id_X\times\fr_{\bar e},
\]
the automorphism of $\overline X$ defined by $f$ by transport of
structure, and
\[
\overline{\Fr}^*_{F/X}
=\Fr^*_{F/X}\otimes\id_{\Lambda_{\bar e}},
\]
the inverse image of $\Fr^*_{F/X}$ by $\overline X\to X$. We may write
\[
\fr_{\overline X}
=\overline{\fr}_X\circ f_X
=f_X\circ\overline{\fr}_X,
\qquad
\Fr^*_{\overline F/\overline X}
=\overline{\Fr}^*_{F/X}.
\]
The corollary of Proposition 3 shows that the endomorphism of
$\RGamma_{\overline X}(\overline F)$ defined by the pair
$(\overline{\fr}_X,\overline{\Fr}^*_{F/X})$ is the inverse of the
automorphism defined by $f_X$. We shall say that
$(\overline{\fr}_X,\overline{\Fr}^*_{F/X})$ is the geometric
Frobenius correspondence; it is the inverse image by
$\overline X\to X$ of the Frobenius correspondence on $(X,F)$. Denote by
\[
\fr_{\RGamma_{\overline X}(\overline F)}
\]
the automorphism of $\RGamma_{\overline X}(\overline F)$ which it
defines. The arithmetic Frobenius operation, coming by transport of
structure from $f\in\pi_e$, is its inverse.

\section*{§ 3. The $L$-function}

\subsection*{1. Definition}

Let $p$ be a prime number; consider a prime number $\ell\ne p$ and a finite extension $\Omega$ of $\mathbb Q_\ell$. To every pair $(X,F)$ consisting of a scheme $X$ of finite type over $e=\Spec(\mathbb F_p)$ and a constructible $\Omega$-sheaf $F$ on $X$, we shall associate a formal series
\[
L_F\in\Omega[[t]],
\]
with constant term $1$.


First consider the special case where $X=e$; let $\bar{\mathbb F}_p$ be an algebraic closure of $\mathbb F_p$ and let $\bar e=\Spec(\bar{\mathbb F}_p)$. The sheaf $F$ is entirely determined by its fibre $F_{\bar e}$, regarded as a finite-dimensional $\Omega$-vector space on which the Galois group
\[
\pi_e=\Gal(\bar{\mathbb F}_p/\mathbb F_p)
\]
acts continuously. Since $\pi_e$ is topologically generated by the Frobenius element
\[
f:\lambda\longmapsto\lambda^p\qquad(\lambda\in\bar{\mathbb F}_p),
\]
the data of the sheaf $F$ is equivalent to that of the vector space $F_{\bar e}$ endowed with an automorphism $f_{F_{\bar e}}$ which is ``topologically unipotent'', that is, whose $k$-th power tends to the identity in the analytic group $\Aut(F_{\bar e})$ as $k$ tends ``multiplicatively'' to infinity. We set
\[
L_F(t)=1/\det(1-f_{F_{\bar e}}^{-1}t)\in\Omega[[t]].
\]
Now suppose that $X=\Spec(\mathbb F_q)$, with $q=p^d$ ($d$ an integer $\ge1$). Let $\bar x$ be a geometric point of $X$, defined by an embedding of $\mathbb F_q$ into $\bar{\mathbb F}_p$; the Galois group
\[
\pi_x=\Gal(\bar{\mathbb F}_p/\mathbb F_q)
\]
is identified with the subgroup of $\pi_e$ topologically generated by $f^d$. Thus the data of the sheaf $F$ is equivalent to that of its fibre $F_{\bar x}$, a finite-dimensional vector space over $\Omega$, endowed with a topologically unipotent automorphism $(f^d)_{F_{\bar x}}$, which we write as $f^d_{F_{\bar x}}$ for simplicity. Let $g$ be the unique morphism from $X$ to $e$; the direct image $g_*(F)$ is likewise defined by its fibre $(g_*F)_{\bar e}$ with the operations of $\pi_e$, which is computed as an ``induced module''
\[
(g_*F)_{\bar e}\simeq F_{\bar x}\otimes_{\pi_x}\pi_e.
\]
We set
\[
L_F(t)=L_{g_*F}(t)=1/\det(1-f^{-1}_{(g_*F)_{\bar e}}t);
\]
an easy determinant calculation shows that this expression is equal to
\[
L_F(t)=1/\det(1-f^{-d}_{F_{\bar x}}t^d).
\]

In the general case of a scheme $X$ of finite type over $e$, endowed with a constructible $\Omega$-sheaf $F$, define
\[
L_F=\prod_{x\in X^0}L_{F_x},
\]
where $X^0$ denotes the set of closed points of $X$; we shall see that this product is multipliable in $\Omega[[t]]$. Indeed, let us make it explicit by introducing, for each $x\in X^0$, a geometric point $\bar x$ of $X$ localized at $x$, defined by an embedding of the residue field $k(x)$ into $\bar{\mathbb F}_p$, and by setting
\[
d(x)=[k(x):\mathbb F_p]
\]
(the degree of $x$). With this notation
\[
L_{F_x}(t)=1/\det(1-f_{F_{\bar x}}^{-d(x)}t^{d(x)})
 =1+\Tr(f_{F_{\bar x}}^{-d(x)})t^{d(x)}+\cdots .
\]
Since for every integer $n$ the set of points $x\in X^0$ such that $d(x)\le n$ is finite, it follows that the family $(L_{F_x})_{x\in X^0}$ is multipliable in $\Omega[[t]]$. Thus the $L$-function of $(X,F)$ is
\[
L_F(t)=\prod_{x\in X^0}1/\det(1-f_{F_{\bar x}}^{-d(x)}t^{d(x)}).
\]

The first formal properties of the $L$-function are collected in the following statement.

\begin{statement}{Proposition 1}
\begin{enumerate}[label=\alph*)]
\item Multiplicativity in $F$. Let
\[
0\longrightarrow F'\longrightarrow F\longrightarrow F''\longrightarrow0
\]
be an exact sequence of constructible $\Omega$-sheaves on $X$; then
\[
L_F=L_{F'}L_{F''}.
\]
\item Let $Y$ be a closed subscheme of $X$ and let $U=X-Y$ be its open complement. Then
\[
L_F=L_{F|U}\,L_{F|Y}
\]
for every constructible $\Omega$-sheaf $F$.
\item Let $X\to S$ be a morphism of schemes of finite type over $\mathbb F_p$, and let $F$ be a constructible $\Omega$-sheaf on $X$; then
\[
L_F=\prod_{s\in S^0}L_{F|X_s}.
\]
\end{enumerate}
\end{statement}

Assertion a) is clear; b) follows from $X^0=U^0\cup Y^0$ as a disjoint union, and c) follows from
\[
X^0=\coprod_{s\in S^0}X_s^0
\]
using associativity of multipliable infinite products. These formulas on closed points are valid because the image of a closed point under a morphism of preschemes of finite type over $\mathbb F_p$ is again a closed point.

\subsection*{2. Statement of the rationality theorem and reduction to the case of a curve}

\begin{statement}{Theorem}
Let $X$ be a scheme of finite type over $\mathbb F_p$, of dimension $n$, let $g:X\to e$ be its structural morphism, and let $F$ be a constructible $\Omega$-sheaf on $X$. Then
\[
L_F=\prod_{i=0}^{2n}\bigl(L_{R^i_!g(F)}\bigr)^{(-1)^i}.
\tag{1}
\]
\end{statement}

In this statement the $R^i_!g(F)$ are the higher direct images with proper support, in the sense of $\ell$-adic cohomology. We put
\[
H^i_!(\bar X,\bar F)=R^i_!\bar g(\bar F)=R^i_!g(F)_{\bar e}
\]
(by the base-change theorem), where $\bar X$, $\bar F$, and $\bar g$ are obtained from $X$, $F$, and $g$ by the base change $\bar e\to e$, with $\bar e$ the spectrum of an algebraic closure of $\mathbb F_p$. With this notation, by no.~1,
\[
L_{R^i_!g(F)}(t)=1/\det(1-f^{-1}_{H^i_!(\bar X,\bar F)}t)
 =L_{f^{-1}_{H^i_!(\bar X,\bar F)}}(t),
\]
where $L_u(t)=1/\det(1-ut)$ for every endomorphism $u$ of a finite-dimensional vector space, according to the formalism developed in an earlier exposé.

More generally, consider a morphism $u:X\to Y$ of schemes of finite type over $\mathbb F_p$ and a constructible $\Omega$-sheaf $F$ on $X$. Let $(P_{F,u})$ denote the following property:
\[
(P_{F,u}):\qquad
L_F=\prod_i\bigl(L_{R^i_!u(F)}\bigr)^{(-1)^i}.
\]
The theorem for the pair $(X,F)$ is precisely $(P_{F,g})$, where $g:X\to e$ is the structural morphism of $X$. We shall reduce it to the case where $X$ is the projective line by means of a few lemmas.

\begin{statement}{Lemma 1}
Let $u:X\to Y$ be a morphism of schemes of finite type over $\mathbb F_p$, and let $F$ be a constructible $\Omega$-sheaf. If $(P_{F|X_y,u_y})$ is true for every $y\in Y^0$, then $(P_{F,u})$ is true.
\end{statement}

Indeed,
\[
L_F=\prod_{y\in Y^0}L_{F|X_y}
\]
(no.~1, Proposition 1c), and hence
\[
=\prod_{y\in Y^0}\prod_i\bigl(L_{R^i_!u_y(F|X_y)}\bigr)^{(-1)^i}
\]
using $(P_{F|X_y,u_y})$; but
\[
R^i_!u_y(F|X_y)\simeq R^i_!u(F)_y
\]
by the base-change theorem, which gives
\[
L_F=\prod_i\left(\prod_{y\in Y^0}L_{(R^i_!u(F))_y}\right)^{(-1)^i}
=\prod_i\bigl(L_{R^i_!u(F)}\bigr)^{(-1)^i}
\]
by the definition of the $L$-function on $Y$.

\begin{statement}{Corollary}
Let $u:X\to Y$ be a finite morphism of schemes of finite type over $\mathbb F_p$, and let $F$ be a constructible $\Omega$-sheaf on $X$. The property $(P_{F,u})$ is true.
\end{statement}

The lemma allows us to reduce to the case where $Y$ is punctual: $Y=\Spec(\mathbb F_q)$. Then $X$ is finite and
\[
u_*(F)\simeq \bigoplus_{x\in X}u_*(F_x).
\]
To establish the formula $(P_{F,u})$, namely $L_F=L_{u_*F}$, one is therefore reduced to the case where $X$ itself is punctual: $X=\Spec(\mathbb F_{q'})$. Let $g:X\to e$ and $h:Y\to e$ be the structural morphisms. By definition (no.~1),
\[
L_F=L_{g_*F}=L_{h_*u_*F}=L_{u_*F},
\]
since $g=h\circ u$.

\begin{statement}{Lemma 2}
Let $u:X\to Y$ and $v:Y\to Z$ be morphisms of schemes of finite type over $\mathbb F_p$, and let $F$ be a constructible $\Omega$-sheaf on $X$. If the properties $(P_{F,u})$ and $(P_{R^i_!u(F),v})$ for all $i$ are true, then $(P_{F,v\circ u})$ is true.
\end{statement}

Indeed,
\[
L_F=\prod_i\bigl(L_{R^i_!u(F)}\bigr)^{(-1)^i}
 =\prod_i\prod_j\bigl(L_{R^j_!v(R^i_!u(F))}\bigr)^{(-1)^i(-1)^j}
\]
\[
=\prod_{i,j}\bigl(L_{R^j_!v(R^i_!u(F))}\bigr)^{(-1)^{i+j}}
\]
by $(P_{F,u})$ and $(P_{R^i_!u(F),v})$. The Leray spectral sequence
\[
R^j_!v(R^i_!u(F))\Longrightarrow R^k_!(v\circ u)(F)
\]
and the multiplicativity of $L_F$ with respect to $F$ (no.~1, Proposition 1a) finally give
\[
L_F=\prod_k\bigl(L_{R^k_!(v\circ u)(F)}\bigr)^{(-1)^k},
\]
as required.

\begin{statement}{Lemma 3}
Let $X$ be a scheme of finite type over $\mathbb F_p$, and let $g:X\to e$ be its structural morphism.
\begin{enumerate}[label=\alph*)]
\item Let $F$ be a constructible $\Omega$-sheaf on $X$, let $Y$ be a closed subscheme of $X$, and let $U=X-Y$ be the open complement. If the theorem is true for two of the three pairs $(X,F)$, $(Y,F|Y)$, $(U,F|U)$, then it is true for the third.
\item Let
\[
0\longrightarrow F'\longrightarrow F\longrightarrow F''\longrightarrow0
\]
be an exact sequence of constructible $\Omega$-sheaves on $X$. If the theorem is true for two of the three pairs $(X,F')$, $(X,F)$, $(X,F'')$, then it is true for the third.
\end{enumerate}
\end{statement}

We shall use only a), but we also give b) for symmetry. The proofs of a) and b) are analogous; let us prove a), for example. The exact sequence
\[
\cdots\longrightarrow R^i_!g_U(F|U)\longrightarrow R^i_!g(F)
\longrightarrow R^i_!g_Y(F|Y)\longrightarrow R^{i+1}_!g_U(F|U)\longrightarrow\cdots
\]
and Proposition 1a) (no.~1) give
\[
\prod_i\bigl(L_{R^i_!g(F)}\bigr)^{(-1)^i}
=\prod_i\bigl(L_{R^i_!g_U(F|U)}\,L_{R^i_!g_Y(F|Y)}\bigr)^{(-1)^i},
\]
and Proposition 1b) (no.~1) then allows one to conclude.

The arguments of Lemmas 1, 2, and 3 may be summarized by saying that, by Proposition 1 of no.~1, the function $L_F$ has formal properties analogous to those of cohomology with proper support.

We now use the preceding lemmas to reduce to the case where $X$ is the projective line.

First reduction: by noetherian induction one may construct a partition $(U_j)_{1\le j\le r}$ of $X$ by affine schemes such that $U_{j+1}$ is open in $X-\bigcup_{i=1}^jU_i$ for $j=1,\ldots,r-1$. Suppose that the theorem has been established for every pair $(U,G)$, where $U$ is an affine scheme of finite type over $\mathbb F_p$ and $G$ is a constructible $\Omega$-sheaf on $U$. It is then true for $(U_j,F|U_j)$, $j=1,\ldots,r$; by Lemma 3a), descending induction on $j$ shows that it is true for $(\bigcup_{i\ge j}U_i,F|\bigcup_{i\ge j}U_i)$, and hence for $(X,F)$, since $X=\bigcup_{i\ge1}U_i$. It is therefore enough to prove the theorem in the case where $X$ is affine.

Second reduction: suppose $X$ affine of dimension $n$. There exists a finite morphism $u:X\to\mathbb E^n$ from $X$ into numerical space of dimension $n$. We reason by induction on $n$. For $n=0$, $X$ is finite over $e$ and the theorem is true by the corollary to Lemma 1. If $n\ge1$, suppose the theorem established for affine schemes of dimension $n-1$; let $u:X\to\mathbb E^n$ be a finite morphism and let $v:\mathbb E^n\to\mathbb E^1$ be a projection morphism, defined from an isomorphism $\mathbb E^n\simeq\mathbb E^1\times\mathbb E^{n-1}$. Composing, one obtains a morphism $w=v\circ u:X\to\mathbb E^1$ whose fibres are affine schemes of dimension $n-1$. Lemma 1 then shows that the property $(P_{F,w})$ is true by the induction hypothesis. By Lemma 2, it is enough to establish the theorem for the pairs $(\mathbb E^1,R^i_!w(F))$. Thus one is reduced to the case where $X=\mathbb E^1$ is the affine line.

Third reduction: suppose $X=\mathbb E^1$. Let $i:\mathbb E^1\to\mathbb P^1$ be the immersion of the affine line into the projective line $\mathbb P^1$, and let $Z=\mathbb P^1-i(\mathbb E^1)$ be the reduced closed subscheme at the point at infinity of $\mathbb P^1$. Since the theorem is true for $(Z,i_!(F)|Z)$, it is equivalent, by Lemma 3a), to establish it for $(\mathbb E^1,F)$ or for $(\mathbb P^1,i_!(F))$; since $i$ is an open immersion one has $F=i_!(F)|\mathbb E^1$.

\subsection*{3. Proof of the theorem}

After these reductions, it remains to establish the theorem in the case where $X$ is the projective line. We shall first transform formula (1) of the statement into a Lefschetz-type formula involving traces; this latter formula will be proved in the particular case where $X$ is a projective smooth curve.

Return to formula (1):
\[
L_F=\prod_i\bigl(L_{R^i_!g(F)}\bigr)^{(-1)^i}
=\frac{P_F^1\cdots P_F^{2n-1}}{P_F^0\cdots P_F^{2n}},
\quad\text{where}\quad
P_F^i(t)=\det(1-f^{-1}_{H^i_!(\bar X,\bar F)}t).
\]
The logarithm of the left-hand side is
\[
\sum_{x\in X^0}\sum_{n=1}^{\infty}\Tr(f_{F_{\bar x}}^{-n d(x)})\,t^{n d(x)}/n
=\sum_{n=1}^{\infty}\left(\sum_{d(x)|n} d(x)\Tr(f_{F_{\bar x}}^{-n})\right)t^n/n,
\]
while that of the right-hand side is
\[
\sum_{n\ge1}\left(\sum_i(-1)^i\Tr(f^{-n}_{H^i_!(\bar X,\bar F)})\right)t^n/n.
\]
Comparing the coefficients of $t^n/n$, one sees that it is necessary to establish
\[
\sum_{d(x)|n}d(x)\Tr(f_{F_{\bar x}}^{-n})
=\sum_i(-1)^i\Tr(f^{-n}_{H^i_!(\bar X,\bar F)})
\tag{2}
\]
for all $n\ge1$; these calculations are legitimate because the field $\Omega$ has characteristic zero.

Transform formula (2) further by supposing that $X$ is proper, which
allows us to write ordinary cohomology instead of cohomology with
proper support. Recall that the category of $\mathbb Q_\ell$-sheaves
on $X$ is a category of fractions of the category of $\ell$-adic
sheaves. By abuse of notation, we shall henceforth denote by
$F=(F_\nu)_{\nu\in\mathbb N}$ the $\ell$-adic
sheaf\footnote{To simplify the notation, for the remainder of the
proof we shall suppose that $\Omega=\mathbb Q_\ell$. In the general
case, one would have to replace $\mathbb Z_\ell$ by the normal closure
of $\mathbb Z_\ell$ in $\Omega$. (Details left to the reader.)}
on $X$ corresponding to the $\mathbb Q_\ell$-sheaf underlying $F$.
Then $\bar F=(\bar F_\nu)$ and the Frobenius morphisms
$\overline{\Fr}^*_{F_\nu/X}$ define a morphism of $\ell$-adic sheaves
\[
\overline{\Fr}^*_{F/X}:
\fr_X^*(\bar F)\longrightarrow\bar F,
\]
hence an $\ell$-adic correspondence on $(\bar X,\bar F)$. For every
integer $\nu$, the endomorphism
$\fr_{R\Gamma_{\bar X}(\bar F_\nu)}$ defined by
$(\bar{\fr}_X,\overline{\Fr}^*_{F_\nu/X})$ is the inverse of
$f_{R\Gamma_{\bar X}(\bar F)}$ (§2, no.~3, Corollary to Proposition 3).
Using the definition of $\ell$-adic cohomology, it follows that the
endomorphism $\fr_{H^i(\bar X,\bar F)}$ of $H^i(\bar X,\bar F)$
defined by $(\bar{\fr}_X,\overline{\Fr}^*_{F/X})$ is the inverse of
$f_{H^i(\bar X,\bar F)}$; thus
\[
f^{-n}_{H^i(\bar X,\bar F)}
=\fr^n_{H^i(\bar X,\bar F)}
\]
for every $n\ge1$ and every $i$.

Apply this result in the particular case where $X=x=\Spec(\mathbb F_q)$ ($q=p^d$) is reduced to a point whose degree $d$ divides $n$. We obtain
\[
f^{-n}_{F_{\bar x}}=\fr^n_{F_{\bar x}}=(\Fr^*_{F_{\bar x}/\bar x})^n
\]
($\bar x$ is the geometric point defined by an embedding of $\mathbb F_q$ into $\bar{\mathbb F}_p$).

On the other hand, $\bar{\fr}_X$ and $f_X$ act in inverse ways on the points of $\bar X$. It follows that the projection $\bar X\to X$ establishes a bijective correspondence between the set of closed points of degree $d$ of $X$ and the set of $d$-element orbits of $\bar{\fr}_X$ acting on $\bar X$. Noting that the set $\bar X^{\bar{\fr}_X^n}$ of points of $\bar X$ fixed by $\bar{\fr}_X^n$ is the union of the orbits of $\bar{\fr}_X$ whose cardinality divides $n$, one can rewrite (2) in the ``geometric'' form
\[
\sum_{\bar x\in\bar X^{\bar{\fr}_X^n}}\Tr(\fr^n_{F_{\bar x}})
=\sum_i(-1)^i\Tr(\fr^n_{H^i(\bar X,\bar F)}).
\tag{2'}
\]
This is a Lefschetz-type formula on the scheme $\bar X$ endowed with the sheaf $\bar F$ and the correspondence $(\bar{\fr}_X,\overline{\Fr}^*_{F/X})$.

To lighten notation, we shall henceforth write $X$, $F$, $h$, and $h_F$ instead of $\bar X$, $\bar F$, $\bar{\fr}_X^n$, and $(\overline{\Fr}^*_{F/X})^n$; since $X$ and $F$ will no longer appear in their former sense, no confusion can arise. We shall suppose that $X$ is a smooth projective curve over $k=\bar{\mathbb F}_p$. The fixed points of $h$ on $X$ all have multiplicity $1$: indeed, $h$ acts by $T\mapsto T^{p^n}$ on the completed local ring of a fixed point, where $T$ is a uniformizer. The theorem will follow from the following result.

\begin{statement}{Proposition 2}
Let $X$ be a smooth projective curve over an algebraically closed field $k$ of characteristic $p$, let $F$ be a constructible $\ell$-adic sheaf on $X_{\mathrm{et}}$ (with $\ell$ a prime number distinct from $p$), let $h$ be an endomorphism of $X$ whose fixed points all have multiplicity $1$, and let $h_F:h^*(F)\to F$ be a morphism of sheaves. The following formula, where $h_{H^i(X,F)}$ denotes the endomorphism of $H^i(X,F)$ defined by $(h,h_F)$, is true:
\[
\sum_{x\in X^h}\Tr((h_F)_x)=\sum_i(-1)^i\Tr(h_{H^i(X,F)}).
\tag{2''}
\]
\end{statement}

In this seminar, an analogous Lefschetz formula has been established for a constructible torsion sheaf $F_\nu$, annihilated by $\ell^{\nu+1}$ on the curve $X$, endowed with a morphism $h_{F_\nu}:h^*(F_\nu)\to F_\nu$. It is written
\[
\sum_{x\in X^h}\Tr((h_{F_\nu})_x)=\Tr^*h_{R\Gamma_X(F_\nu)},
\tag{3}
\]
where both sides belong to $\mathbb Z/\ell^{\nu+1}\mathbb Z$.

Suppose that the $\ell$-adic sheaf in the statement is written as $F=(F_\nu)_{\nu\in\mathbb N}$, with $F$ annihilated by $\ell^{\nu+1}$. For each $\nu$ we may write formula $(3)_\nu$, and one sees that the left-hand sides of these formulas, as $\nu$ varies, are the components of an $\ell$-adic integer which is precisely the left-hand side of formula (2'') to be proved. It remains to establish that the right-hand sides of the formulas $(3)_\nu$ define, as $\nu$ varies, an $\ell$-adic integer equal to
\[
\sum_i(-1)^i\Tr(h_{H^i(X,F)}),
\quad\text{where}\quad
H^i(X,F)=\varprojlim_\nu H^i(R\Gamma_X(F_\nu)).
\]
This follows from purely algebraic lemmas that we shall now prove to finish.

Consider a projective system $(K_\nu)_{\nu\in\mathbb N}$ of perfect complexes, with $K_\nu\in D_{\mathrm{parf}}(\mathbb Z/\ell^{\nu+1}\mathbb Z)$, such that for each $\nu$ the transition morphism $K_{\nu+1}\to K_\nu$ is isomorphic, in the sense of derived categories, to the canonical morphism
\[
K_{\nu+1}\to K_{\nu+1}/\ell^{\nu+1}K_{\nu+1}.
\]
Suppose each complex $K_\nu$ is endowed with an endomorphism $h_\nu\in\mathrm{Fl}(D_{\mathrm{parf}}(\mathbb Z/\ell^{\nu+1}\mathbb Z))$, so that $(h_\nu)_{\nu\in\mathbb N}$ is compatible with the transition morphisms, thereby defining an endomorphism of the projective system $(K_\nu)$. In the application to the proof of Proposition 2 one takes
\[
K_\nu=R\Gamma_X(F_\nu),\qquad h_\nu=h_{R\Gamma_X(F_\nu)};
\]
the condition on transition morphisms follows by the ``Künneth formula'' from the analogous condition verified by the transition morphisms of the $\ell$-adic projective system $(F_\nu)$.

We propose to find a projective system $(K'_\nu)_{\nu\in\mathbb N}$, where $K'_\nu$ is a complex of free $(\mathbb Z/\ell^{\nu+1}\mathbb Z)$-modules of finite type, zero outside an interval of degrees $[a,b]$ independent of $\nu$, and endowed with an endomorphism of complexes $h'_\nu:K'_\nu\to K'_\nu$, such that the following conditions are satisfied:
\begin{enumerate}[label=\arabic*)]
\item the transition morphism $K'_{\nu+1}\to K'_\nu$ is isomorphic, as a morphism of complexes, to
\[
K'_{\nu+1}\to K'_{\nu+1}/\ell^{\nu+1}K'_{\nu+1}
\]
for every $\nu\in\mathbb N$;
\item $(h'_\nu)_{\nu\in\mathbb N}$ is compatible with the transition morphisms;
\item $(K'_\nu,h'_\nu)$ is isomorphic to $(K_\nu,h_\nu)$ in $D_{\mathrm{parf}}(\mathbb Z/\ell^{\nu+1}\mathbb Z)$ for every $\nu$, these isomorphisms being choosable so as to be compatible with the transition morphisms of $(K'_\nu)$ and $(K_\nu)$.
\end{enumerate}
Assuming the construction of the systems $(K'_\nu)$ and $(h'_\nu)$, we may write
\[
\Tr^*h_\nu=\Tr^*h'_\nu=\sum_i(-1)^i\Tr h_\nu'^i
\]
for every $\nu\in\mathbb N$. Let
\[
K'=\varprojlim_\nu K'_\nu,
\qquad h'=\varprojlim_\nu h'_\nu
\]
be the induced endomorphism of $K'$. Then $(\Tr^*h_\nu)_\nu$ is an $\ell$-adic integer equal to $\sum_i(-1)^i\Tr h'^i$. Moreover, the projective system $(K'_\nu)$ satisfies the Mittag-Leffler condition, and so does the system $(H^{i-1}(K'_\nu))_\nu$ for every $i$; it follows that
\[
H^i(K')\simeq\varprojlim_\nu H^i(K'_\nu)
\]
for every $i$ (EGA $0_{\mathrm{III}}$ 13.2.3). In the case of interest to us, where $K_\nu=R\Gamma_X(F_\nu)$ and $h_\nu=h_{R\Gamma_X(F_\nu)}$, $H^i(K')$ is identified with $H^i(X,F)$, and the endomorphism of $H^i(K')$ defined by $h'$ is identified with $h_{H^i(X,F)}$. Consequently
\[
\sum_i(-1)^i\Tr h'^i=
\sum_i(-1)^i\Tr h_{H^i(X,F)},
\]
and formula (2'') in the statement is thereby proved.

The existence of the systems $(K'_\nu)$ and $(h'_\nu)$ rests on two lemmas, in which we consider a commutative ring $A$, an ideal $I$ of $A$, and the functor
\[
T:L\longmapsto L\otimes_A A_0
\]
where $L$ is a complex of $A$-modules bounded on the right, with $A_0=A/I$.

\begin{statement}{Lemma 1}
Suppose that the ideal $I$ is nilpotent. Consider a complex $L\in D^-(A)$ flat in every degree, a complex $M\in D^-(A_0)$ projective in every degree, and an isomorphism
\[
u:T(L)\xrightarrow{\sim}M
\]
in $D^-(A_0)$. There exist a complex $L'$ of $A$-modules, bounded on the right and projective in every degree, an isomorphism
\[
v:L\xrightarrow{\sim}L'
\]
in $D^-(A)$, and an isomorphism of complexes
\[
w:T(L')\xrightarrow{\sim}M,
\]
such that $u=w\circ T(v)$ in $D^-(A_0)$.
\end{statement}

\begin{statement}{Remark}
If $M$ is assumed to be free, respectively of finite type, in every degree, then the complex $L'$ is also free, respectively of finite type, in every degree. Indeed, if a flat $A$-module $N$ is such that $T(N)$ is free, respectively of finite type, it is easy to show that $N$ is free, respectively of finite type (Bourbaki, Commutative Algebra, Chap.~II, §3, Prop.~5). Note the corollary: let $L\in\Ob D^-(A)$ be such that $L\lten_A A_0$ is a pseudo-coherent, respectively perfect, complex of $A_0$-modules, that is, isomorphic in $D^-(A)$ to a complex bounded above, respectively bounded, and projective of finite type in every degree. Then $L$ is a pseudo-coherent, respectively perfect, complex of $A$-modules.
\end{statement}

\begin{statement}{Lemma 2}
Consider a complex $L$ of $A$-modules bounded on the right, projective in every degree, endowed with an endomorphism $\varphi$, and an endomorphism $\varphi_0$ of $M=T(L)$ coinciding with $T(\varphi)$ in the sense of the derived category $D^-(A_0)$. There exists an endomorphism $\varphi'$ of $L$ which coincides with $\varphi$ in the sense of $D^-(A)$ and satisfies $T(\varphi')=\varphi_0$.
\end{statement}

Before proving these lemmas, let us indicate how they are applied to prove the existence of the systems $(K'_\nu)$ and $(h'_\nu)$. We argue by induction on $\nu$, supposing that $K'_0,h'_0,\ldots,K'_{\nu-1},h'_{\nu-1}$ have already been constructed. Apply the lemmas with
\[
A=\mathbb Z/\ell^{\nu+1}\mathbb Z,
\qquad
I=\ell^\nu\mathbb Z/\ell^{\nu+1}\mathbb Z
\]
so that $I^2=0$, and
\[
A_0=\mathbb Z/\ell^\nu\mathbb Z.
\]
First take $L=K_\nu$ and $M=K'_{\nu-1}$ in the first lemma. In view of the remark, this gives the existence of $K'_\nu=L'$ with transition morphism $K'_\nu\to K'_{\nu-1}$ given by $w$, so that condition 1) above is satisfied. The isomorphism $v$ in $D^-(A)$ allows $h_\nu$ to be transported to an endomorphism of $K'$ in the sense of $D^-(A)$. We may suppose that the latter comes from an endomorphism of complexes $h''_\nu:K'_\nu\to K'_\nu$. Applying the second lemma with $L=K'_\nu$, $\varphi=h''_\nu$, and $\varphi_0=h'_{\nu-1}$, we obtain the existence of $h'_\nu:K'_\nu\to K'_\nu$ satisfying condition 2); condition 3) is also satisfied. Note that, with the notation of Lemma 1, if $M^i=0$ for an index $i$, then $L'^i=0$ follows from Nakayama's lemma; the complexes $K'_\nu$ are therefore indeed all zero outside one and the same interval of degrees.

To prove Lemmas 1 and 2 we shall need the following result.

\begin{statement}{Lemma 3}
\begin{enumerate}[label=(\roman*)]
\item Let $P$ and $Q$ be $A$-modules. If $P$ is projective, then for every $A_0$-linear map $f_0:T(P)\to T(Q)$ there exists an $A$-linear map $f:P\to Q$ such that $T(f)=f_0$.
\item Suppose that the ideal $I$ is nilpotent. For every projective $A_0$-module $P_0$ there exists a projective $A$-module $P$ such that $T(P)\simeq P_0$. If moreover $P'$ is a flat $A$-module such that $T(P')\simeq P_0$, then $P'\simeq P$.
\end{enumerate}
\end{statement}

Assertion (i) is clear: if $u_P$ (resp. $u_Q$) denotes the canonical map from $P$ onto $T(P)$ (resp. from $Q$ onto $T(Q)$), then the composite $f_0\circ u_P$ is an $A$-linear map from $P$ to $T(Q)$, and $u_Q$ is surjective. Since $P$ is projective, there is a factorization of $f_0\circ u_P$ through $u_Q$, namely $f_0\circ u_P=u_Q\circ f$, where $f$ is an $A$-linear map from $P$ to $Q$; then $T(f)=f_0$.

To prove (ii), first observe that the case where $P_0$ is free is evident, with no hypothesis on the ideal $I$. The general case, where $P_0$ is projective, is reduced to this by introducing a free $A_0$-module $L_0$ of which $P_0$ is a direct factor, and a free $A$-module $L$ such that $T(L)\simeq L_0$. The module $P_0$ is identified with the image of a projector $e_0:L_0\to L_0$, and it remains to see that there exists a projector $e:L\to L$ lifting $e_0$. This follows from Bourbaki, Commutative Algebra, Chap.~III, §4, no.~6, Lemma 2, applied to the sub-$A$-algebra $B$ of $\End(L)$ generated by an endomorphism of $L$ lifting $e_0$ (whose existence is ensured by (i)) and to the nilpotent ideal which is the trace on $B$ of the kernel of the canonical homomorphism $\End(L)\to\End(L_0)$. This proves the existence of $P$ in (ii). If moreover $P'$ is as in (ii), then the composite isomorphism $T(P)\simeq P_0\simeq T(P')$ comes, by (i), from a homomorphism $f:P\to P'$. Since $P$ and $P'$ are flat over $A$ and $T(f)$ is an isomorphism, it follows that $\operatorname{Gr}_I(f)$ is an isomorphism; hence $f$ is an isomorphism, since $I$ is nilpotent. This proves (ii).

\begin{statement}{Corollary}
Suppose the ideal $I$ is nilpotent. For every acyclic complex $N_0$ of projective $A_0$-modules, bounded on the right, there exists an acyclic complex $N$ of projective $A$-modules, bounded on the right, such that $T(N)\simeq N_0$.
\end{statement}

The complex $N_0$, being acyclic, decomposes into short exact sequences
\[
0\longrightarrow Z_0^i\longrightarrow N_0^i\longrightarrow Z_0^{i+1}\longrightarrow0,
\]
where $Z_0^i$ denotes, for each $i$, the submodule of cycles of $N_0^i$; moreover $N_0^i=0$ for all sufficiently large $i$. By descending induction on $i$, one sees that the preceding exact sequences are split and that $Z_0^i$ is projective for all $i$. One lifts $N_0$ by lifting all these exact sequences step by step, by descending induction on $i$. It is therefore enough to show that if
\[
S_0:0\longrightarrow P_0\longrightarrow Q_0\longrightarrow R_0\longrightarrow0
\]
is an exact sequence of $A_0$-modules with $Q_0$ projective, and if $R$ is an $A$-module such that $T(R)\simeq R_0$, then there exists an exact sequence
\[
S:0\longrightarrow P\longrightarrow Q\longrightarrow R\longrightarrow0
\]
of $A$-modules, with $Q$ projective, such that $T(S)\simeq S_0$. Now $Q_0$ can be lifted to a projective $A$-module $Q$ by Lemma 3(ii), and $Q_0\to R_0$ can be lifted to an $A$-linear map $Q\to R$ by Lemma 3(i). This map is surjective by Nakayama's lemma; one sets $P=\ker(Q\to R)$ to complete the exact sequence $S$.

Proof of Lemma 2: since $M=T(L)$ is projective in every degree, the hypothesis of the lemma says that $T(\varphi)$ is homotopic to $\varphi_0$. Thus there exists a homotopy operator $k_0=(k_0^i)_i$, with $k_0^i:M^i\to M^{i-1}$, such that
\[
T(\varphi)-\varphi_0=d_0k_0+k_0d_0,
\]
where $d_0$ is the differential of $M$. By Lemma 3(i), $k_0^i$ lifts to a linear map $k^i:L^i\to L^{i-1}$. This gives a homotopy operator $k=(k^i)_i$ in $L$. It is enough to set
\[
\varphi'=\varphi-kd-dk
\]
where $d$ is the differential of $L$, to verify Lemma 2.

Proof of Lemma 1: after replacing $L$ by a quasi-isomorphic complex if necessary, we may suppose that the complex $L$, and hence also $T(L)$, is projective in every degree; we shall assume this henceforth. Thus we may suppose that the isomorphism $u$ of $D^-(A_0)$ comes from a quasi-isomorphism $T(L)\to M$, again denoted by $u$. We shall first reduce to the case where $u$ is surjective. The complex $N_0$ defined by
\[
N_0^i=M^i\oplus M^{i-1},
\quad
 d^i(x,y)=(d^i x,x-d^{i-1}y)
\]
($x\in M^i$, $y\in M^{i-1}$) is visibly acyclic and projective in every degree; an epimorphism of complexes $N_0\to M$ is defined by the projection $N_0^i\to M^i$ in degree $i$. The corollary to Lemma 3 allows us to lift $N_0$ to an acyclic complex $N$, bounded on the right and projective in every degree. We may then replace $L$ by $L\oplus N$, which is quasi-isomorphic to it, since $N$ is acyclic, and projective in every degree. At the same time, replace $u$ by the quasi-isomorphism
\[
T(L\oplus N)=T(L)\oplus N_0\longrightarrow M
\]
defined by $u$ on $T(L)$ and by the epimorphism $N_0\to M$ on $N_0$; it is clear that this last quasi-isomorphism is surjective.

Suppose therefore that $u:T(L)\to M$ is a surjective quasi-isomorphism, and let $R_0$ be its kernel. The complex $R_0$ is acyclic because $u$ is a quasi-isomorphism, and it is projective in every degree because the exact sequence
\[
0\longrightarrow R_0\longrightarrow T(L)\longrightarrow M\longrightarrow0
\]
is split, since $M$ is projective in every degree. Hence, by the corollary to Lemma 3, $R_0$ can be lifted to a complex $R$ of $A$-modules, bounded on the right, acyclic, and projective in every degree. We now show that the morphism $f_0:R_0\to T(L)$ lifts to a morphism of complexes $f:R\to L$. For this, observe that $R_0$ is homotopically trivial, so that $f_0$ is homotopic to $0$ and can be written
\[
f_0=h_0d_{R_0}+d_{T(L)}h_0,
\]
where $h_0=(h_0^i)_i$, with $h_0^i:R_0^i\to T(L)^{i-1}$. Lemma 3(i) allows $h_0$ to be lifted to $h=(h^i)_i$, with $h^i:R^i\to L^{i-1}$, and one sees that
\[
f=hd_R+d_Lh
\]
is a morphism of complexes lifting $f_0$. A standard lemma (compare SGA 60/61, IV 5.7) then implies that $f$ is injective and that $\coker f=L'$ is flat in every degree. Since
\[
T(L')\simeq\coker T(f)=\coker f_0\simeq M
\]
is projective in every degree, it follows, by Lemma 3(ii), that the same is true of $L'$. On the other hand, the canonical homomorphism $L\to L'$ is a quasi-isomorphism, being surjective with acyclic kernel $R$; thus $L'$ satisfies the desired conclusion of Lemma 1, q.e.d.

\section*{Terminological Index}
\newcommand{\sgaindexentry}[2]{\noindent\begin{minipage}[t]{0.62\linewidth}#1\end{minipage}\hfill\begin{minipage}[t]{0.32\linewidth}#2\end{minipage}\par\smallskip}
\sgaindexentry{biduality (theorem of)}{I 3.4.1}
\sgaindexentry{$A$-category, abelian $A$-category}{V 1.1}
\sgaindexentry{class of $(c,N)$, class of $c$}{III 3.2.4}
\sgaindexentry{cohomology class associated with $Y$}{III 6.8}
\sgaindexentry{Chern class}{VII 3.2, 3.9}
\sgaindexentry{composition (of correspondences)}{III 5.2}
\sgaindexentry{condition (reg)}{I 3.1.1}
\sgaindexentry{Mittag-Leffler-Artin-Rees condition}{V 2.1.1}
\sgaindexentry{Grothendieck divisibility conjecture}{III B 6, XII 4.5}
\sgaindexentry{twisted constant ($\ZZ_\ell$-sheaf)}{VI 1.2.1}
\sgaindexentry{cohomological correspondence}{III 3.1, 3.2, 6.4}
\sgaindexentry{equivariant correspondence}{III B 6}
\sgaindexentry{diagonal correspondence}{III 3.2 b)}
\sgaindexentry{Frobenius correspondence}{III 3.2 b), XIV 2.1}
\sgaindexentry{cup-product of correspondences}{III 4.1.4}
\sgaindexentry{quasi-injective dimension}{I 1.4}
\sgaindexentry{normal-crossings divisor}{I 3.1.5}
\sgaindexentry{dual (correspondence)}{III 5.1}
\sgaindexentry{dualizing (complex)}{I 1.7}
\sgaindexentry{local duality}{I 4.2}
\sgaindexentry{Euler-Poincaré (characteristic)}{VII 4.11}
\sgaindexentry{Euler-Poincaré (formula)}{X 7.1}
\sgaindexentry{constructible $\ell$-adic sheaf}{VI 1.1.1}
\sgaindexentry{$\ZZ_\ell$-sheaf}{VI 1.1.4}
\sgaindexentry{constructible $A$-sheaf}{VI 1.4.1, 1.4.3}
\sgaindexentry{constructible $\QQ_\ell$-sheaf, twisted constant}{VI 1.4.2}
\sgaindexentry{constructible $AR$-$\ZZ_\ell$-sheaf}{VI 1.5.1}
\sgaindexentry{arrow of type $i$ over $f$}{III 1.4}
\sgaindexentry{Künneth arrow}{III 1.6, 1.7}
\sgaindexentry{evaluation arrow}{III 2.2.2, III B 5.3.5}
\sgaindexentry{tensor-product arrow}{III B 5.4.2}
\sgaindexentry{strongly desingularizable}{I 3.1.5}
\sgaindexentry{Frobenius (endomorphism)}{XIV 1}
\sgaindexentry{Frobenius (correspondence)}{III 3.2 b), XIV 2.1}
\sgaindexentry{relative Frobenius (morphism)}{XIV 2 def 3}
\sgaindexentry{arithmetic, geometric Frobenius}{XIV p. 20}
\sgaindexentry{Gauss-Bonnet (formula)}{VII 4.9}
\sgaindexentry{Gysin (exact sequence)}{VII 1.5}
\sgaindexentry{direct image (of a cohomological correspondence)}{III 3.3}
\sgaindexentry{Lefschetz-Verdier (formula)}{III 4.5.1, 4.7, 6.7.2}
\sgaindexentry{Nielsen-Wecken (local invariant)}{XII 3.1}
\sgaindexentry{Nielsen-Wecken (formula)}{XII 5.1}
\sgaindexentry{normalized at $x$ (dualizing complex)}{I 4.1, 4.5}
\sgaindexentry{external tensor product (of correspondences)}{III 5.3}
\sgaindexentry{absolute cohomological purity}{I 3.1.4, I App 4.4}
\sgaindexentry{quasi-injective}{I 1.1}
\sgaindexentry{rationality of $L$-functions (theorem)}{XIV 2 no. 2}
\sgaindexentry{Artin representation}{X 4.2.5}
\sgaindexentry{Swan representation}{X 4.2.2}
\sgaindexentry{self-intersection (formula)}{VII 4.1, 9.2}
\sgaindexentry{Shih (theorem)}{V A 3.1}
\sgaindexentry{AR-null projective system}{V 2.2.1}
\sgaindexentry{$J$-adic, AR-$J$-adic projective system}{V 3.1.1, 3.2.2, 5.1.1, 5.1.3}
\sgaindexentry{local Lefschetz-Verdier term}{III 4.2.7, 6.6.3, III B 6}
\sgaindexentry{trace}{III B 5.6.1, 6}
\sgaindexentry{Woods-Hole (formula)}{III 6.12}

\section*{Index of Notation}
\sgaindexentry{$\dim\mathrm{q.inj}$}{I 1.4}
\sgaindexentry{$\underline D_K$}{I 1.6, 1.7}
\sgaindexentry{$D^b(X)_{\mathrm{torf}}$, $D^b(X)_{\mathrm{qinjf}}$}{I 1.10}
\sgaindexentry{$K_X$, $\underline D_X$, $D$}{I 1.12, III 1.3}
\sgaindexentry{$(\mathrm{cdloc})_n$}{I 3.1.3}
\sgaindexentry{$D_{\mathrm{ctf}}(X)$}{III 1.1}
\sgaindexentry{$f_*, f_!, f^!$}{III 1.2}
\sgaindexentry{$\Hom_f^{(i)}(L,M)$}{III 1.4}
\sgaindexentry{$L_1\lten_S L_2$}{III 1.5}
\sgaindexentry{$u_1\lten_S L_2$}{III 1.6.2, 1.7.2, 1.7.5}
\sgaindexentry{$u\lten_S P$}{III 1.8}
\sgaindexentry{$R\underline{\Hom}(u,v)$}{III 2.1.3, 2.1.5}
\sgaindexentry{$R\underline{\Hom}_S(L_1,L_2)$, $\Hom_S(L_1,L_2)$}{III 3.1.2, 6.4.3}
\sgaindexentry{$\cl(c,N)$, $\cl(c)$}{III 3.2.4}
\sgaindexentry{$\cl_{X/S}(Y)$}{III 6.8.1}
\sgaindexentry{$f_*u$, $u_*$}{III 3.3, 6.5}
\sgaindexentry{$\langle u,v\rangle$}{III 4.1.4, III B 6.6.3}
\sgaindexentry{$\langle u,v\rangle_Z$}{III 4.2.7, III B 6.6.3}
\sgaindexentry{$\operatorname{Res}_{Y/S}(\omega)$}{III 6.8}
\sgaindexentry{$A^\circ$, $A^e$, $A_\natural$, $G_\natural$}{III B 5.1}
\sgaindexentry{$\Tr_A$}{III B 5.6, 6.5.3}
\sgaindexentry{$\Tr_{B/A}$}{III B 5.10.2}
\sgaindexentry{$A_{\natural,\varphi}$, $G_{\natural,\varphi}$}{III B 5.13}
\sgaindexentry{$K\!A_{X,L_\natural}$, $K\!A_{X,\natural}$}{III B 6.5.2}
\sgaindexentry{MLAR}{V 2.1.1}
\sgaindexentry{$\Hom_{AR}(N^\circ,C)$}{V 2.4.2}
\sgaindexentry{$J$-ad$(C)$}{V 3.1.1}
\sgaindexentry{$AR$-$J$-ad$(C)$}{V 3.2.1}
\sgaindexentry{$AR$-$J$-adn$(C)$}{V 5.1.3}
\sgaindexentry{$\ell$-adc$(X)$}{VI 1.1.1}
\sgaindexentry{$A$-fc$(X)$}{VI 1.4.1}
\sgaindexentry{$\QQ_\ell$-fc$(X)$}{VI 1.4.3}
\sgaindexentry{$AR$-$\ZZ_\ell$-fc$(X)$}{VI 1.5.3}
\sgaindexentry{$A^\bullet(X)$}{VII 2.2.5}
\sgaindexentry{$c_i(E)$, $c(E)$}{VII 3.2, 3.9}
\sgaindexentry{$EP(X)$}{VII 4.11}
\sgaindexentry{$K_D(C)$, $K(C)$}{VII 1.2}
\sgaindexentry{$\mathcal F(G,\Lambda)$}{X 3.1}
\sgaindexentry{$\Sw_y$}{X 4.2.2, 4.3}
\sgaindexentry{$\Art_y$}{X 4.2.5}
\sgaindexentry{$\Tr^*_{\Lambda_n[G],\varphi}$}{XII 1}
\sgaindexentry{$\Tr^*_{\ZZ_\ell[G],\varphi}$}{XII 2}
\sgaindexentry{$NW_f^{G,\varphi}(x)$}{XII 3.1}
\sgaindexentry{$s_x(f',f,G,\varphi)$}{XII 4.1, 4.4}
\sgaindexentry{$c^\varphi(g)$}{XII 4.2}
\sgaindexentry{$\alpha_x(f,u,f',G,\varphi)$}{XII 6.2.2}
\sgaindexentry{$\fr_X$}{XIV 1}
\sgaindexentry{$\Fr_{X/S}$}{XIV 2}
\sgaindexentry{$X^{(p/S)}$}{XIV 2}
\sgaindexentry{$\Fr^*_{F/X}$}{XIV 2 def 1}
\sgaindexentry{$L_F(t)$}{XIV p. 21}
\sgaindexentry{$|X|$}{III B 5.11.1}

\end{document}

% SGA-CUMULATIVE-END 06_SGA5_EN_COMPLETE_PUBLISHED_CONTENT_READER.tex

% SGA-CUMULATIVE-BEGIN 07_SGA6_EN_CURRENT_LAYERED_READER.tex
\documentclass[11pt]{article}
\usepackage[utf8]{inputenc}
\usepackage[T1]{fontenc}
\usepackage[a4paper,margin=1in]{geometry}
\usepackage{lmodern}
\usepackage{microtype}
\usepackage{amsmath,amssymb,mathtools}
\usepackage{mathrsfs}
\usepackage{tikz-cd}
\usepackage{enumitem}
\usepackage{longtable,booktabs,array}
\usepackage[hidelinks]{hyperref}
\usepackage{hyperxmp}
% Deterministic, zero-content semantic destinations.
\newcommand{\sgaStableTarget}[1]{\hypertarget{#1}{}}
\newcommand{\sgaStableTargetInline}[1]{\pdfdest name {#1} xyz\relax}
\setlength{\parindent}{0pt}
\setlength{\parskip}{0.55em}
\emergencystretch=3em
\hbadness=10000
\hfuzz=20pt
\newcommand{\K}{K}
\newcommand{\Gr}{\operatorname{Gr}}
\newcommand{\Fil}{\operatorname{Fil}}
\newcommand{\ch}{\operatorname{ch}}
\newcommand{\Todd}{\operatorname{Todd}}
\newcommand{\cl}{\operatorname{cl}}
\newcommand{\Parf}{\operatorname{Parf}}
\newcommand{\Tor}{\operatorname{Tor}}
\newcommand{\Spec}{\operatorname{Spec}}
\newcommand{\RRR}{\mathrm{RRR}}
\providecommand{\codim}{\operatorname{codim}}
\providecommand{\Inv}{\operatorname{Inv}}
\providecommand{\Pic}{\operatorname{Pic}}
\providecommand{\Z}{\mathbb Z}
\providecommand{\Q}{\mathbb Q}
\providecommand{\R}{\mathbb R}
\providecommand{\C}{\mathbb C}
\providecommand{\Ahat}{\widehat A}
\providecommand{\Ahp}{\widehat A^{+}}
\providecommand{\That}{\widehat T}
\providecommand{\rad}{\operatorname{rad}}
\providecommand{\Gl}{\operatorname{Gl}}
\providecommand{\GL}{\operatorname{GL}}
\providecommand{\ch}{\operatorname{ch}}
\providecommand{\Todd}{\operatorname{Todd}}
\providecommand{\Gr}{\operatorname{Gr}}

\providecommand{\Atil}{\widetilde A}
\providecommand{\Ctil}{\widetilde C}
\providecommand{\K}{\mathrm K}
\providecommand{\cC}{\mathcal C}
\providecommand{\cF}{\mathscr F}
\providecommand{\eps}{\varepsilon}
\providecommand{\ellc}{\ell}
\providecommand{\Tor}{\operatorname{Tor}}
\providecommand{\Coker}{\operatorname{Coker}}
\providecommand{\Ker}{\operatorname{Ker}}
\providecommand{\Supp}{\operatorname{Supp}}
\providecommand{\codim}{\operatorname{codim}}
\providecommand{\Gr}{\operatorname{G}}
\long\def\SGAArchiveOnly#1{}
\providecommand{\rank}{\operatorname{rank}}
\providecommand{\rang}{\operatorname{rang}}
\providecommand{\cO}{\mathcal O}

\title{SGA 6: Intersection Theory and the Riemann--Roch Theorem}
\author{P. Berthelot, A. Grothendieck, L. Illusie}
\date{22 July 2026}

\hypersetup{
  hypertexnames=false,
  pdftitle={SGA 6: Intersection Theory and the Riemann--Roch Theorem -- English Translation},
  pdfauthor={P. Berthelot, A. Grothendieck, L. Illusie},
  pdfsubject={English translation of SGA 6},
  pdfkeywords={SGA 6, intersection theory, Riemann--Roch, English translation}
}
\providecommand{\N}{\mathbb N}
\providecommand{\Zell}{\mathbb Z_\ell}
\providecommand{\Qell}{\mathbb Q_\ell}
\providecommand{\C}{\mathbb C}
\providecommand{\Ch}{\operatorname{Ch}}
\providecommand{\Sym}{\operatorname{Sym}}
\providecommand{\Td}{\operatorname{Td}}

\begin{document}

% Inherited, partially source-synchronized English control through sourcePDF525.
% BEGIN INCLUDED SGA6_English_current_buildable_source_20260823/SGA6_sourcePDF001_525_English_Inherited_PartiallySourceSynchronized_fragment.tex
\begin{center}
{\Large Seminar on Algebraic Geometry of Bois Marie}\par
{\large 1966--67}\par
\vspace{1.5em}
{\Large Intersection Theory and the Riemann--Roch Theorem}\par
{\large (SGA 6)}\par
\vspace{1.5em}
a seminar directed by\par
P. Berthelot, A. Grothendieck, L. Illusie\par
with the collaboration of\par
D. Ferrand, J.-P. Jouanolou, O. Jussila, S. Kleiman, M. Raynaud, J.-P. Serre
\end{center}

\sgaStableTarget{sga6.front.preface}\section*{Preface}
The present edition of \emph{Séminaire de Géométrie Algébrique 6} reproduces, up to minor changes, the first edition, distributed by the I.H.E.S. in the form of mimeographed exposés. The corrections made concern essentially printing errors, except in \hyperlink{sga6.e01.expose}{Exposés I} and III, where slight modifications were made by L. Illusie to a few statements which were incorrect in the primitive version. Footnotes have also been added, especially in \hyperlink{sga6.e14.expose}{Exposé XIV} by A. Grothendieck, in order to point out recent progress on questions that were still open at the time of the seminar.

Finally, note that Exposé XI, by D. Ferrand, on determinant theory for perfect complexes, was not written up and therefore will not be published; to avoid errors in references, we have left unchanged the numbering of the following exposés.

\begin{flushright}June 1971\end{flushright}

\sgaStableTarget{sga6.front.introduction}\section*{Introduction}
The purpose of the present seminar is to develop a global theory of intersections, and a Riemann--Roch formula, for arbitrary schemes. The reader will find a description of the programme of the seminar in \hyperlink{sga6.e00.expose}{Exposé 0}. The possibility in principle of proving a Riemann--Roch formula for general regular schemes, along the lines of the well-known Borel--Serre report of 1958, was already clear at that time, at least for a projective morphism. The extension to the general case is less evident; the programme in which it belongs, and which is partly carried out in this seminar, goes back to 1960. As was also the case for the duality theory of coherent sheaves, an essential tool for formulating a satisfactory theory is Verdier's theory of derived categories; familiarity with it is indispensable for understanding the seminar.

Apart from this theory, the study of the seminar requires little more than a general knowledge of the foundations of algebraic geometry as set out in EGA, chapters I, II, III. We shall only occasionally refer to EGA IV, for certain facts concerning flat or smooth morphisms, which for the most part also appear in the first exposés of SGA 1. Finally, to develop some results with the full generality desirable for applications, we sometimes use the notion of a ringed site and of a ringed topos, for which we refer to SGA 4, Exposés I to IV. A reader unfamiliar with the language of sites and topoi may everywhere replace these objects by ordinary topological spaces, the objects of the topos then being replaced by open subsets of these spaces; nevertheless, we recommend that the reader assimilate the language of topoi, which provides an extremely convenient unifying principle.

The fundamental notion for the theory presented here is that of a perfect complex of modules, together with its various variants, developed in \hyperlink{sga6.e01.expose}{Exposés I}, II, III. It seems clear that these notions, also important in other contexts in algebraic geometry, notably for formulas of Lefschetz--Verdier type in cohomology with discrete coefficients (SGA 5) or coherent coefficients, will also have a role to play outside algebraic geometry, especially in the formulation of a complex analytic variant of the Riemann--Roch theorem presented here, or of suitable variants of the Atiyah--Singer index theorem. Some indications in this direction will be given in \hyperlink{sga6.e02.expose}{Exposé II}. Unfortunately, at present there is still no statement, even heuristic, which would encompass these two theorems, the first of which remains conjectural for the moment. A new idea seems to be missing, just as one is also missing in algebraic geometry for proving Riemann--Roch beyond projective hypotheses (cf. \hyperlink{sga6.e14.sec.2}{Exposé XIV, no. 2)}.

We shall make no use in this seminar of the local theory of intersections, set out in J.-P. Serre, \emph{Algèbre locale. Multiplicités} (Lecture Notes in Mathematics no. 11, Springer), and merely point out here that Serre's course had an evident influence on the genesis of the theory in 1957. Nor shall we use the theory of the Chow ring developed in the Chevalley Seminar of 1958, Exposés 2 and 3 (École Normale Supérieure). That theory is essentially tied to nonsingularity hypotheses, whereas the aim of our seminar is, on the contrary, to develop a theory of intersections on general schemes, and even on general locally ringed topoi. On this matter see the comments in \hyperlink{sga6.e14.expose}{Exposé XIV}, nos. 4 and 8, which give the relations between our theory and that of the Chow ring and raise the question of a generalization of the latter. Thus one may say that the seminar presents a coherent and self-contained theory of intersections, but one which is by no means exhaustive of the different points of view useful, indeed indispensable, in intersection theory; consequently it should be complemented by the cited exposés of Serre and Chevalley.

We have appended to the seminar, as an appendix to \hyperlink{sga6.e00.expose}{Exposé 0}, Grothendieck's 1957 report, ``Classes of sheaves and the Riemann--Roch theorem'', which served as the basis for the Borel--Serre seminar at Princeton in the same year and for their already cited article. That report sketches two proofs of the Riemann--Roch theorem for nonsingular quasi-projective varieties. The first, valid for the time being only in characteristic zero but giving, in return, a slightly more precise result in the case of an immersion, has not yet appeared in the literature, apart from the work of Borel--Serre and the present seminar. The report presupposes, of course, none of the other exposés of the seminar, and can even serve, like \hyperlink{sga6.e00.expose}{Exposé 0}, as an introduction to the study of the latter, especially for a reader not yet familiar with Borel--Serre. The proof just alluded to applies essentially unchanged to the case of a projective morphism of complex analytic spaces, and will doubtless also apply in arbitrary characteristic once the problem of the operations ``exterior powers'' in the derived category is solved (cf. \hyperlink{sga6.e14.expose}{Exposé XIV}, nos. 1, 2, 3). The absence of a study of this question is no doubt the most serious gap in this seminar, from the point of view we have adopted.

One will notice the absence, in the table of contents of the present seminar, of two of the participants mentioned on the cover page. One of them, J.-P. Jouanolou, took an active part in the technical elaboration of the first part of the seminar, but was prevented from taking part in the oral exposés; in particular, he is responsible for the linear-independence assertion contained in the important statement \hyperlink{sga6.e06.thm.1-1}{VI 1.1}, giving the structure of the ring $K^\circ$ of classes of locally free modules on a projective bundle, which had previously been proved only under the assumption that the base scheme has an ample invertible module. The other, J.-P. Serre, gave two oral exposés which were logically independent of the rest of the seminar, and therefore preferred to publish them in the form of a separate article (cf. J.-P. Serre, ``Grothendieck groups of split reductive group schemes'', to appear in \emph{Publications Mathématiques}, no. 34).

As in the other parts of the \emph{Séminaire de Géométrie Algébrique du Bois Marie}, the sigla SGA 1 to SGA 6 refer to the different parts of that seminar; the Roman numeral following the siglum indicates the number of the exposé, and the Arabic numerals following it correspond to the internal numbering of the exposé in question. For internal references within the present seminar the siglum SGA 6 is omitted, and references of the form I 4.9 are used, meaning: Exposé I, statement 4.9. The siglum EGA refers to the \emph{Éléments de Géométrie Algébrique} of J. Dieudonné and A. Grothendieck.

\sgaStableTarget{sga6.front.table-of-contents}\section*{Table of contents}
\begin{longtable}{@{}p{0.16\linewidth}p{0.64\linewidth}r@{}}
\toprule
Exposé & Title and author & Page \\
\midrule
0 & Outline of a programme for a theory of intersections on general schemes, by A. Grothendieck & 1 \\
Appendix & Classes of sheaves and the Riemann--Roch theorem, by A. Grothendieck & 20 \\
I & Generalities on finiteness conditions in derived categories, by L. Illusie & 56 \\
II & Existence of global resolutions, by L. Illusie & 160 \\
III & Relative finiteness conditions, by L. Illusie & 190 \\
IV & Grothendieck groups of ringed topoi, by L. Illusie & 222 \\
V & Generalities on $\lambda$-rings, by P. Berthelot & 297 \\
VI & The $K^\circ$ of a projective bundle: computations and consequences, by P. Berthelot & 365 \\
VII & Regular immersions and the computation of $K^\circ$ of a blow-up, by P. Berthelot & 416 \\
VIII & The Riemann--Roch theorem, by P. Berthelot & 466 \\
IX & Some computations of $K_\bullet$-groups, by P. Berthelot & 498 \\
X & Formalism of intersections on proper algebraic schemes, by O. Jussila, with an appendix by A. Grothendieck, ``Specialization in intersection theory'' & 519 \\
XI & Not written up & -- \\
XII & A relative representability theorem for the Picard functor, by M. Raynaud, written by S. Kleiman & 595 \\
XIII & Finiteness theorems for the Picard functor, by S. Kleiman & 616 \\
XIV & Open problems in intersection theory, by A. Grothendieck & 667 \\
\bottomrule
\end{longtable}

\sgaStableTarget{sga6.e00.expose}\section*{Exposé 0}
\begin{center}
{\large Outline of a programme for a theory of intersections on general schemes}\par
by A. Grothendieck
\end{center}

The present exposé is introductory in nature, and its reading is not logically indispensable for the study of the seminar. It is addressed more particularly to readers familiar with the provisional version of the Riemann--Roch theorem, as it is presented in the Borel--Serre report or in Grothendieck's report already mentioned in the introduction, cited below as [RRR], and reproduced as an appendix at the end of the present exposé.

\sgaStableTarget{sga6.e00.sec.1}\subsection*{1. Reminder on the Riemann--Roch formula}
Recall the Riemann--Roch formula for a proper morphism
\[
f:X\longrightarrow Y
\]
of quasi-projective smooth schemes over a field $k$, and a coherent sheaf $F$ on $X$, whose class in the group $K(X)$ of classes of coherent sheaves on $X$ is denoted by $\cl(F)$:
\sgaStableTarget{sga6.e00.eq.1-1}\begin{equation}\tag{1.1}
\Todd(T_Y)\,\ch_Y\bigl(f_!(\cl(F))\bigr)=f_*\bigl(\Todd(T_X)\,\ch_X(F)\bigr).
\end{equation}
This formula is valid in $A(Y)\otimes\mathbb Q$, where $A(Y)$ is the Chow ring of $Y$. The $f_*$ on the right hand side is obtained by tensoring with $\mathbb Q$ the homomorphism ``direct image of cycles''
\[
f_*:A(X)\longrightarrow A(Y),
\]
while the term on the left is the Euler--Poincaré characteristic of $F$ relative to $f$:
\[
f_!(\cl(F))=\sum_i(-1)^i\cl\bigl(R^if_*(F)\bigr).
\]
The symbols $\ch_X$, $\ch_Y$ denote the Chern characters on $X$ and on $Y$, and $T_X$, $T_Y$ the tangent bundles of $X$ and $Y$. As is known, $\Todd(-)$ and $\ch(-)$ are universal polynomials with coefficients in $\mathbb Q$ in the Chern classes of the argument. Since the constant term of $\Todd(-)$ is $1$, it is an invertible element for every value of the argument. Thus, after multiplying by $\Todd(T_Y)^{-1}$, formula \hyperlink{sga6.e00.eq.1-1}{(1.1)} takes the equivalent form, more convenient for our purposes,
\sgaStableTarget{sga6.e00.eq.1-2}\begin{equation}\tag{1.2}
\ch_Y\bigl(f_!(\cl(F))\bigr)=f_*\bigl(\Todd(T_f)\,\ch_X(F)\bigr),
\end{equation}
where
\sgaStableTarget{sga6.e00.eq.1-3}\begin{equation}\tag{1.3}
T_f=T_X-f^*T_Y\in K(X).
\end{equation}
Thus $T_f$ plays the role of a virtual relative tangent bundle of $X$ over $Y$. When $f$ is smooth, one simply has $T_f=T_{X/Y}$, whereas when $f:X\to Y$ is an immersion, one obtains
\[
T_f=-\check N_{X/Y},
\]
where $N_{X/Y}$ denotes the normal sheaf of $X$ in $Y$.

One of the main aims of the seminar is to generalize \hyperlink{sga6.e00.eq.1-2}{(1.2)} simultaneously in two directions:
\begin{enumerate}[label=\alph*)]
\item to get rid of the hypothesis that a base field $k$ exists;
\item to replace the regularity hypotheses on $Y,X$ by a hypothesis of ``local regularity'' on $f$.
\end{enumerate}
Finally, along the way, we shall also deal with the following problem:
\begin{enumerate}[label=\alph*),start=3]
\item to eliminate the quasi-projectivity hypotheses which, in the absence of a base field, are expressed by the existence of ample invertible modules on $X$ and $Y$.
\end{enumerate}

\sgaStableTarget{sga6.e00.sec.2}\subsection*{2. Removing the base field}
Let us first examine generalization a), while retaining the hypotheses of regularity and of the existence of ample invertible modules on $X$ and on $Y$.

The definitions of $K(X)$, $K(Y)$ and of the homomorphism $f_!:K(X)\to K(Y)$ then present no new problem, thanks to the fact that $X$ and $Y$ are regular. The most natural way to give a meaning to \hyperlink{sga6.e00.eq.1-2}{(1.2)} would seem to consist in defining Chow rings $A(X),A(Y)$ and a group homomorphism
\[
f_*:A(X)\longrightarrow A(Y),
\]
in establishing a theory of Chern classes, giving maps
\[
c_i:K(X)\longrightarrow A(X),
\]
and similarly for $Y$, and finally in describing an element, the virtual relative tangent bundle,
\[
T_f\in K(X).
\]

\sgaStableTarget{sga6.e00.sec.2.1}\subsubsection*{2.1}
For the latter, a definition presents itself rather evidently. One uses the fact that the morphism $f$, thanks to the hypothesis that ample invertible modules exist on $X$ and on $Y$, may be factored as
\[
X\xrightarrow{i}X'\xrightarrow{f'}Y,
\]
where $i$ is a closed immersion and $f'$ is smooth; for example, one may take $X'=\mathbb P^r_Y$ for suitable $r$. One then sets
\sgaStableTarget{sga6.e00.eq.2-1}\begin{equation}\tag{2.1}
T_f=i^*T_{X'/Y}-\check N_{X/X'}\in K(X),
\end{equation}
where $\Omega^1_{X'/Y}$ is the locally free module of relative $1$-differentials of $X'$ over $Y$, or the relative cotangent module of $X'$ over $Y$, and $N_{X/X'}$ is the conormal module of $X$ in $X'$, by definition equal to $J/J^2$, where $J$ is the ideal on $X'$ defining $X$. The regularity hypotheses on $X$ and $X'$ also imply that $N_{X/X'}$ is locally free. Finally, $\check N$ denotes the dual of $N$. One verifies without difficulty, in \hyperlink{sga6.e08.expose}{Exposé VIII}, that the element $T_f$ defined by \hyperlink{sga6.e00.eq.2-1}{(2.1)} is independent of the chosen factorization of $f$.

\sgaStableTarget{sga6.e00.sec.2.2}\subsubsection*{2.2}
As for a definition, under the conditions considered, of a Chow ring and a corresponding theory of Chern classes, although it is plausible that such a theory should be developable on the model of smooth varieties over a field, it seems that this work has not been done at the present time. We shall therefore adopt another viewpoint, by directly defining a ring which replaces the Chow ring, and whose tensor product with $\mathbb Q$ is indeed isomorphic to $A(X)\otimes\mathbb Q$ in the classical case.

A first natural filtration of $K(X)$, the one considered in [RRR], is the topological filtration by the codimension of the support of coherent sheaves: $\Fil^iK(X)$ is generated by the $\cl(F)$ with $F$ coherent and $\operatorname{codim}(\operatorname{Supp}F,X)\ge i$. A first question is the compatibility of this filtration with the multiplicative structure of $K(X)$. In [RRR], this compatibility is obtained, for $X$ smooth over a field, by means of Chow's moving lemma. Assuming this assertion, one obtains a graded ring $\Gr_{\mathrm{top}}(X)$ and a canonical surjective homomorphism
\[
\varphi:A(X)\longrightarrow \Gr_{\mathrm{top}}(X),
\]
sending the class of a prime cycle $Z$ to $\cl(\mathcal O_Z)$; it then turns out, \hyperlink{sga6.e14.para.4-2}{Exposé XIV 4.2}, that the kernel is a torsion subgroup. Moreover, [RRR] shows how the images in $\Gr_{\mathrm{top}}(X)$ of the Chern classes $c_i(F)$ can be computed directly in terms of $\cl(F)\in K(X)$ and the exterior-power operations $\lambda^i$ in $K(X)$.

This shows that, for a noetherian regular scheme $X$ endowed with an ample invertible module, the ring $\Gr_{\mathrm{top}}(X)$ is a convenient substitute for the hypothetical Chow ring $A(X)$. It also has the advantage over the latter that certain computations of [RRR], leading to important formulas in this ring, are performed directly in it--formulas which, even for $X$ projective and smooth over a field, have not at present been established in $A(X)$ itself (\hyperlink{sga6.e14.para.4-3}{Exp. XIV, 4.3} and \hyperlink{sga6.e14.para.6-4}{6.4}). To justify this viewpoint completely, however, one would have to solve the problem mentioned above of the compatibility of the topological filtration of $K(X)$ with its ring structure, a problem which seems essentially connected with Chow's moving lemma, itself the essential technical ingredient of the theory of the Chow ring which we have claimed to be able to avoid.

\sgaStableTarget{sga6.e00.sec.2.3}\subsubsection*{2.3}
To avoid this problem, one endows $K(X)$ with a second filtration, finer than the topological filtration, and describable in purely algebraic terms from the structure of $K(X)$ as an augmented $\lambda$-ring. The definition of this filtration is suggested by the expression given in [RRR] for the Chern classes of an element $x\in K(X)$ in $\Gr_{\mathrm{top}}(X)$:
\sgaStableTarget{sga6.e00.eq.2-3}\begin{equation}\tag{2.3}
c_i(x)=\gamma^i(x-\epsilon(x))\mod \Fil^{i+1}K(X),
\end{equation}
where $\gamma^i$ denotes a certain linear combination of the $\lambda^j$ ($j\le i$), made explicit in loc. cit., and where $\epsilon$ is the augmentation. One shows that $\gamma^i(x-\epsilon(x))$ is indeed of topological filtration at least $i$, which gives meaning to formula \hyperlink{sga6.e00.eq.2-3}{(2.3)}. One then defines the $\lambda$-filtration of $K(X)$ by the ideals $\Fil^i_\lambda K(X)$, where $\Fil^i_\lambda K(X)$ consists of finite sums of expressions of the form
\[
\gamma^{i_1}(x_1-\epsilon(x_1))\cdots \gamma^{i_r}(x_r-\epsilon(x_r)),\qquad i_1+\cdots+i_r\ge i.
\]
The associated graded ring for this filtration will be denoted $\Gr^\bullet(X)$; it is this ring which will replace the Chow ring of $X$ in the seminar. In the case under consideration here, one can moreover show that the canonical homomorphism
\sgaStableTarget{sga6.e00.eq.2-4}\begin{equation}\tag{2.4}
\Gr^\bullet(X)\longrightarrow \Gr_{\mathrm{top}}(X)
\end{equation}
is an isomorphism modulo torsion, justifying the viewpoint that, after tensoring with $\mathbb Q$, this ring can play exactly the same role as the Chow ring tensored with $\mathbb Q$. On the other hand, exactly what was needed has been done so that formula \hyperlink{sga6.e00.eq.2-3}{(2.3)}, with $\Fil_{\mathrm{top}}$ replaced by $\Fil_\lambda$, defines a theory of Chern classes with values in $\Gr^\bullet(X)$; the Chern classes in $\Gr_{\mathrm{top}}(X)$ are then deduced by means of the homomorphism \hyperlink{sga6.e00.eq.2-4}{(2.4)}. It is shown in \hyperlink{sga6.e05.expose}{Exposé V}, using the fact proved in \hyperlink{sga6.e06.expose}{Exposé VI} that $K(X)$ is a special $\lambda$-ring in the sense of [RRR], that is, a $\lambda$-ring in the sense of \hyperlink{sga6.e05.expose}{Exposé V}, that this notion of Chern class satisfies all the usual formal properties reviewed in [RRR].

\sgaStableTarget{sga6.e00.sec.2.4}\subsubsection*{2.4}
One should note a very appreciable advantage of the definition of $\Gr^\bullet(X)$ over that of $A(X)$ and $\Gr_{\mathrm{top}}(X)$: it remains meaningful for an absolutely arbitrary scheme $X$, and more generally for an arbitrary locally ringed topos $X$. One then takes for $K(X)$ the group of classes constructed from locally free modules on $X$, which is again an augmented $\lambda$-ring, hence equipped with a corresponding filtration and an associated graded $\Gr^\bullet(X)$. The price of this generality, however, is that when one is not willing to neglect torsion phenomena, i.e. to tensor with $\mathbb Q$, the properties of $\Gr^\bullet(X)$ are in several respects less satisfactory than those of its two competitors (\hyperlink{sga6.e14.sec.4}{Exp. XIV no. 4)}. Thus, assuming again that $f:X\to Y$ is a proper morphism of noetherian regular schemes endowed with ample invertible modules, it is not true in general that the homomorphism
\[
f_*:K(X)\longrightarrow K(Y)
\]
maps $\Fil^i_\lambda K(X)$ into $\Fil^{i-d}_\lambda K(Y)$, where $d$ is the relative dimension of $X$ over $Y$, and hence defines, by passage to associated gradeds, a degree $-d$ homomorphism from $\Gr^\bullet(X)$ to $\Gr^\bullet(Y)$. However, \hyperlink{sga6.e07.expose}{Exposé VII} proves that this is true after tensoring with $\mathbb Q$, so that one obtains a degree $-d$ homomorphism
\sgaStableTarget{sga6.e00.eq.2-5}\begin{equation}\tag{2.5}
f_*:\Gr^\bullet(X)\otimes\mathbb Q\longrightarrow \Gr^\bullet(Y)\otimes\mathbb Q,
\end{equation}
which plays the role of the homomorphism ``direct image of cycles'' $A(X)\to A(Y)$.

% SGA-SRC-001469
Having said this, we see that we possess all the structural elements needed to give meaning to formula \hyperlink{sga6.e00.eq.1-2}{(1.2)} when the existence of a base field is not postulated, with $f:X\to Y$ a proper morphism of noetherian regular schemes admitting ample invertible modules.

\sgaStableTarget{sga6.e00.sec.3}\subsection*{3. The case of general schemes}
We now show how one can still give meaning to formula \hyperlink{sga6.e00.eq.1-2}{(1.2)} when the regularity hypothesis on $X,Y$ is abandoned and replaced by a local hypothesis on the morphism $f$, to be specified below.

\sgaStableTarget{sga6.e00.sec.3.1}\subsubsection*{3.1}
A first difficulty comes from the fact that the observation which served as the starting point of the theory [RRR], namely that the group $K(X)$ of classes of sheaves on $X$ is the same whether it is defined in terms of coherent sheaves or in terms of locally free sheaves, is essentially tied to the regularity hypothesis on $X$. For a general scheme $X$, say noetherian for simplicity, one must distinguish between these two groups; we shall denote them respectively by $K_\bullet(X)$ and $K^\bullet(X)$. The first is covariant in $X$ for proper morphisms, the second contravariant in $X$ for arbitrary morphisms. Notice that $K^\bullet(X)$ is naturally endowed with a ring structure, indeed with an augmented $\lambda$-ring structure, whereas $K_\bullet(X)$ no longer has, in general, a ring structure, but only a module structure over $K^\bullet(X)$. The homomorphism $x\mapsto x\cdot\cl(\mathcal O_X)$ from $K^\bullet(X)$ to $K_\bullet(X)$ is no longer an isomorphism in general. For our formulation of the generalized Riemann--Roch theorem, we shall work exclusively with $K^\bullet(X)$ and $K^\bullet(Y)$, to the exclusion of $K_\bullet(X)$ and $K_\bullet(Y)$. As noted in \hyperlink{sga6.e00.sec.2.3}{2.3}, these augmented $\lambda$-rings, equipped with their natural $\lambda$-filtration, give rise to graded rings $\Gr^\bullet(X)$, $\Gr^\bullet(Y)$, and to Chern-class homomorphisms
\[
c_i:K^\bullet(X)\longrightarrow \Gr^\bullet(X),
\]
and similarly for $Y$. The difficulty, however, lies in defining a direct-image homomorphism
\sgaStableTarget{sga6.e00.eq.3-1}\begin{equation}\tag{3.1}
f_*:K^\bullet(X)\longrightarrow K^\bullet(Y),
\end{equation}
which can no longer be defined by the naive formula
% SGA-SRC-001470
\[
f_*(\cl(F))=\sum_i(-1)^i\cl\bigl(R^if_*(F)\bigr),
\]
because even if $F$ is locally free, the modules $R^if_*(F)$ are in general no longer locally free, nor even of finite Tor-dimension, and the preceding formula loses its meaning. A solution to this difficulty is furnished by the viewpoint of derived categories. Indeed, whereas the $R^if_*(F)$ are in general ``bad'' sheaves, say of infinite Tor-dimension, to which one would not know how to associate a class in $K^\bullet(Y)$, the complex $Rf_*(F)$ on $Y$, the direct image of $F$ in the sense of derived categories, whose cohomology sheaves are the $R^if_*(F)$, tends to be excellent. More precisely, suppose that $F$ has finite Tor-dimension over $\mathcal O_Y$; this is the case if $\mathcal O_X$ has finite Tor-dimension over $\mathcal O_Y$ and, in addition, $F$ is locally free. Then one proves (\hyperlink{sga6.e03.expose}{Exp. III}) that the complex $Rf_*(F)$ is ``perfect'', by which we mean here that it has coherent cohomology sheaves (as was already known in any case) and globally finite Tor-dimension; equivalently, on every affine open $U$ of $Y$, this complex is isomorphic, in the derived category $D(U)$, to a bounded complex of locally free modules of finite type. From the viewpoint of homological algebra, perfect complexes are ``as good as'' locally free modules, of which they are the natural generalization. They moreover form a ``triangulated subcategory'' of the derived category $D(Y)$, i.e. a subcategory stable under the cones (or ``mapping cylinders'') of morphisms. On the other hand, to every triangulated category $C$ one associates a group $K(C)$ of classes, a natural variant of the construction of the group of classes of an abelian category, by taking the free group generated by the objects of $C$, modulo the relations
\[
\cl(K)-\cl(K')-\cl(K'')
\]
for every distinguished triangle
\[
\begin{tikzcd}[column sep=large]
 & K'' \arrow[dl] & \\
K' \arrow[rr] && K \arrow[ul]
\end{tikzcd}
\]
What then replaces the fundamental observation of [RRR] in the present case is the fact that, when the scheme $Y$ admits an ample invertible module, every perfect complex $L_\bullet$ on $Y$ is globally (and not merely locally) isomorphic in $D(Y)$ to a bounded complex $L'$ of locally free modules of finite type (\hyperlink{sga6.e02.expose}{Exp. II}). It follows easily (\hyperlink{sga6.e04.expose}{Exp. IV}) that the canonical homomorphism
\[
% SGA-SRC-001471
K^\bullet(Y)\longrightarrow K(\Parf(Y))
\]
is an isomorphism, where $\Parf(Y)$ is the triangulated category of perfect complexes on $Y$. Thus to every perfect complex $L_\bullet$ on $Y$ one can associate a class
\[
\cl(L_\bullet)\in K^\bullet(Y),
\]
which is none other than
\[
\cl(L_\bullet)=\sum_i(-1)^i\cl(L'_i),
\]
where $L'$ is as above, so that the $L'_i$ are indeed locally free modules on $Y$ whose classes may be taken in $K^\bullet(Y)$.

\sgaStableTarget{sga6.e00.sec.3.2}\subsubsection*{3.2}
To give a meaning to \hyperlink{sga6.e00.eq.1-2}{(1.2)}, it remains to impose on $f$ a hypothesis making it possible to give a meaning to the homomorphism \hyperlink{sga6.e00.eq.2-5}{(2.5)}, i.e. ensuring that the homomorphism \hyperlink{sga6.e00.eq.3-1}{(3.1)} just defined is, up to a shift and modulo torsion, compatible with the filtrations of these rings; and finally one must again define an element, the ``virtual relative tangent bundle'',
\[
T_f\in K^\bullet(X).
\]
For the latter, returning to the definition proposed in \hyperlink{sga6.e00.sec.2.1}{2.1}, one sees that a hypothesis is needed ensuring that, with the notation of loc. cit., the conormal module $N_{X/X'}$ is locally free. The most natural hypothesis to this end is that $i$ be a ``regular immersion'', i.e. identify $X$ with a subscheme of $X'$ defined by an ideal which locally can be defined by an $\mathcal O_{X'}$-regular sequence of sections of $\mathcal O_{X'}$. This hypothesis, purely local in nature, in fact does not depend on the choice of the factorization of $f$ into an immersion followed by a smooth morphism (\hyperlink{sga6.e08.expose}{Exp. VIII}), and, as above, the same is true for the class $T_f$ defined by formula \hyperlink{sga6.e00.eq.2-1}{(2.1)}. When the hypothesis just considered is satisfied, the morphism $f$ is called a ``locally complete intersection'' morphism. Such a morphism is also locally of finite Tor-dimension.

We shall also allow ourselves to omit the word ``locally'' in this terminology, which cannot cause any confusion.

If $f$ is a locally complete intersection morphism, one introduces a natural notion of ``virtual relative dimension'' $d(f)$ by the formula
\[
d(f)(x)=\dim_x\bigl(f^{-1}(f(x))\bigr)-\codim_x(X,X').
\]
It is an integer-valued function, locally constant at the point $x\in X$, hence constant if $X$ is connected; moreover it does not depend on the choice of the factorization \hyperlink{sga6.e00.eq.2-1}{(2.1)}. When this virtual relative dimension $d$ is constant, when $f$ is proper, and when $X$ and $Y$ admit ample invertible modules, one can prove (\hyperlink{sga6.e08.expose}{Expos\'e VIII}), although this is far from a trivial result, that
\begin{equation}\tag{\sgaStableTargetInline{sga6.e00.eq.3-4}3.4}
f_*\bigl(\Fil^i K^\bullet(X)_{\mathbb Q}\bigr)\subset \Fil^{i-d}K^\bullet(Y)_{\mathbb Q},
\end{equation}
whence a homomorphism \hyperlink{sga6.e00.sec.3.3}{(3.3)} of degree $-d$.

\sgaStableTarget{sga6.e00.sec.3.3}\subsubsection*{3.3}
We have therefore assembled all the structural elements needed to give a meaning to formula \hyperlink{sga6.e00.eq.1-2}{(1.2)} in the case of a proper locally complete intersection morphism of noetherian schemes $X,Y$, both admitting ample invertible modules. With a little additional care, this formulation in fact retains a meaning independently of any noetherian hypothesis.

In this general form, this Riemann--Roch formula will be proved in \hyperlink{sga6.e08.expose}{Expos\'e VIII}. The reader will observe that the proof given in this seminar combines elements of the two proofs which appear in [RRR] and which were discussed in the Introduction. The need to resort to the methods of the first of these proofs, involving more $K$-formalism than the second, followed in Borel--Serre, is due to the need for a proof of the inclusion \hyperlink{sga6.e00.eq.3-4}{(3.4)}, a question which did not arise in the more restricted framework of [RRR] and Borel--Serre [2], where one could work with the Chow ring or the ring $\Gr^\bullet_{\mathrm{top}}$. From the second proof, we retain the introduction of the scheme obtained by blowing up $Y$ along $X$ in the case where $f$ is an immersion. This technical artifice will no doubt disappear, and the proof in the case of an immersion will extend to the case where one no longer postulates the existence of an ample invertible module on $Y$, once the question raised in \hyperlink{sga6.e14.sec.1}{Expos\'e XIV, no. 1}, concerning the introduction of exterior-power operations in the derived category $D^-(Y)$ has been resolved.

\sgaStableTarget{sga6.e00.sec.4}\subsection*{4. Removing the ample invertible module hypothesis}
We must finally give a few indications of how one can still give a meaning to \hyperlink{sga6.e00.eq.1-2}{(1.2)} in the absence of a hypothesis asserting the existence of ample invertible modules on $X$ and $Y$.

\sgaStableTarget{sga6.e00.sec.4.1}\subsubsection*{4.1}
First note that for an arbitrary scheme $X$, or more generally an arbitrary locally ringed topos, it seems beyond doubt that the ``right'' invariant which should replace the ring $K(X)$ of [RRR] is not the ring $K^\bullet(X)$ considered above, defined in terms of locally free modules on $X$, but rather the group $K(\Parf(X))$ considered in \hyperlink{sga6.e00.sec.3.1}{3.1}, defined in terms of perfect complexes on $X$. We shall denote this group here by $K^\bullet(X)$, and distinguish it from the ``naive'' group of \hyperlink{sga6.e00.sec.3.1}{3.1}, denoted $K^\bullet_{\mathrm{naive}}(X)$. This point of view is forced by the need to define a direct-image homomorphism \hyperlink{sga6.e00.eq.3-1}{(3.1)} for a proper morphism
\[
f:X\longrightarrow Y
\]
of finite Tor-dimension, say between noetherian schemes; the non-noetherian case requires somewhat stronger local finiteness properties on $f$, cf. \hyperlink{sga6.e03.expose}{Expos\'e III}. With the definition of $K^\bullet(X)$ now adopted, this homomorphism $f_*$ is indeed defined in a trivial way by formula \hyperlink{sga6.e00.sec.3.2}{(3.2)}, where now $F$ is an arbitrary perfect complex on $X$, which implies that $Rf_*(F)$ is also perfect.

\sgaStableTarget{sga6.e00.sec.4.2}\subsubsection*{4.2}
This new definition of $K^\bullet(X)$ raises, however, a new problem in homological algebra: that of defining a $\lambda$-structure on $K^\bullet(X)$ compatible with its natural ring structure coming from tensor product in the derived category. If $L$ and $L'$ are two perfect complexes, then so is their total tensor product $L\otimes^{\mathbf L}L'$. More precisely, one would need to define ``exterior-power'' operations in the derived category $D^-(X)$, transforming perfect complexes into perfect complexes. This question has not yet been clarified at present, and, as was said in the Introduction, it constitutes the most serious gap in the seminar. Once this question is solved, and proceeding as in \hyperlink{sga6.e00.sec.2.3}{2.3}, one obtains a natural filtration on $K^\bullet(X)$, allowing one to form the associated graded ring $\Gr^\bullet(X)$, which will play the role of the Chow ring, and to define a theory of Chern classes
\[
c_i:K^\bullet(X)\longrightarrow \Gr^i(X).
\]

\sgaStableTarget{sga6.e00.sec.4.3}\subsubsection*{4.3}
To give a meaning to formula \hyperlink{sga6.e00.eq.1-2}{(1.2)} for a proper locally complete intersection morphism
\[
f:X\longrightarrow Y
\]
of noetherian schemes, one must still include in the statement of this formula the nontrivial inclusions \hyperlink{sga6.e00.eq.3-4}{(3.4)}, where $d$ is the virtual relative dimension, in order to be able to define \hyperlink{sga6.e00.sec.3.3}{(3.3)}; and finally one must define the virtual relative tangent class
\[
T_f\in K^\bullet(X).
\]
When $f$ admits a factorization \hyperlink{sga6.e00.eq.2-1}{(2.1)} through an immersion into a relative smooth scheme, the definition of $T_f$ presents no new difficulty. It can moreover be reformulated in that case, in a way more consonant with the spirit of derived categories, by considering the well-known canonical homomorphism
\sgaStableTarget{sga6.e00.eq.4-1}\begin{equation}\tag{4.1}
N_{X/X'}\longrightarrow \Omega^1_{X'/Y}\otimes_{\mathcal O_{X'}}\mathcal O_X
\end{equation}
induced by the exterior differential, and the chain complex $L_{X/Y}$ defined from \hyperlink{sga6.e00.eq.4-1}{(4.1)} by ``extending by zeros'', with the complex equal in degree $1$, respectively $0$, to the source, respectively the target, of the homomorphism \hyperlink{sga6.e00.eq.4-1}{(4.1)}. Formula \hyperlink{sga6.e00.eq.2-2}{(2.2)} then takes the more satisfactory form
\sgaStableTarget{sga6.e00.eq.4-2}\begin{equation}\tag{4.2}
T_f=\cl\bigl(L_{X/Y}^{\vee}\bigr),
\end{equation}
a formula which makes sense in $K^\bullet(X)$ when $f$ is a locally complete intersection morphism, since in that case the complex $L_{X/Y}$ is manifestly perfect. The assertion that $T_f$ is invariant relative to the chosen factorization of $f$ can then be made precise as an invariance assertion for the complex $L_{X/Y}$: up to canonical isomorphism in the derived category $D(X)$, this complex is likewise independent of the chosen factorization (\hyperlink{sga6.e08.expose}{Expos\'e VIII}).

\sgaStableTarget{sga6.e00.sec.4.4}\subsubsection*{4.4}
When $f:X\to Y$ is an arbitrary morphism of schemes, one can still construct a canonical object $L_{X/Y}$ of the derived category $D(X)$, defined up to unique isomorphism, which on each affine open of $X$ lying above an affine open of $Y$ is isomorphic to the complex constructed according to the principle explained in the preceding paragraph. In the general case, however, ``smooth'' should be replaced by ``formally smooth'', so that every algebra $B$ over a ring $A$ occurs as a quotient of a formally smooth algebra, for example a polynomial algebra. For such a definition it is not possible to proceed by simple gluing from the affine constructions, since the objects of the derived category $D(X)$ are, essentially, not glueable.

Here is a general construction of $L_{X/Y}$, valid more generally for every morphism
\[
f:X\longrightarrow Y
\]
of commutatively ringed topoi. Let $B=\mathcal O_X$ and $A=f^{-1}\mathcal O_Y$, so that $A$ is a ring on the space $X$ and $B$ is an $A$-algebra. Let $T\to B$ be a homomorphism of sheaves of sets which generates $B$ as an $A$-algebra, i.e. such that the corresponding homomorphism of $A$-algebras
\[
C=A[T]\longrightarrow B
\]
is an epimorphism of sheaves of sets; here $A[T]$ denotes the commutative $A$-algebra freely generated by $T$, that is, the sheaf associated with the presheaf
\[
U\longmapsto A(U)[T(U)].
\]
Let $J$ be the kernel of $C\to B$, and consider, as in \hyperlink{sga6.e00.sec.4.3}{4.3}, the chain complex of length $1$ defined by the natural homomorphism
\[
J/J^2\longrightarrow \Omega^1_{C/A}\otimes_C B.
\]
It is not difficult to see that, up to canonical isomorphism in the category $D(X)$, the complex so constructed is independent of the choice of $T$ and of $T\to B$, so that it deserves the designation of the relative cotangent complex $L_{X/Y}$ of $X$ over $Y$. One also verifies that, on every affine open of $X$ over an affine open of $Y$, it is given, up to canonical isomorphism, by the procedure of 4.3.

In the case where $f$ is locally complete intersection, one sees moreover that the result is a perfect complex, which makes it possible again to define an element $T_f$ of $K^\bullet(X)$ by formula \hyperlink{sga6.e00.eq.4-2}{(4.2)}. Here again, in general, it would not have been possible to define $T_f$ as an element of $K^\bullet_{\mathrm{naive}}(X)$, and nothing permits one to suppose that $T_f$ is globally isomorphic in $D(X)$ to a bounded complex with locally free components of finite type. This is therefore an additional justification for the point of view advanced in \hyperlink{sga6.e00.sec.4.1}{4.1} in favour of the ``sophisticated'' $K^\bullet(X)$, in preference to $K^\bullet_{\mathrm{naive}}(X)$.

\sgaStableTarget{sga6.e00.sec.4.5}\subsubsection*{4.5}
Subject to the question raised in \hyperlink{sga6.e00.sec.4.2}{4.2}, we have therefore again assembled all the structural elements needed to give a meaning to \hyperlink{sga6.e00.eq.1-2}{(1.2)} for a proper locally complete intersection morphism of arbitrary schemes; the noetherian hypothesis implicitly made in the preceding sections can in fact be eliminated without any essential difficulty. Unfortunately, this does not mean that, even subject to this reservation, the formula in question is established, even in the case where $Y$ is the spectrum of an algebraically closed field of characteristic $p>0$ and $X$ is, say, a proper smooth algebraic scheme of dimension $3$. See the comments in \hyperlink{sga6.e14.sec.2}{Expos\'e XIV, no. 2}.

\sgaStableTarget{sga6.e00.sec.5}\subsection*{5. Analytic variant}
In the form presented in no. 4, the Riemann--Roch formula \hyperlink{sga6.e00.eq.1-2}{(1.2)} also has a meaning for a proper locally complete intersection morphism of complex analytic spaces, or of rigid analytic spaces in the sense of Tate. But of course the proof given in [RRR] applies to this case only when one assumes in addition that the morphism $f$ is projective; and, as in the algebraic case, it seems that an essentially new idea is needed to treat the general case. The formula we have in view here is a formula with values in an ``analytic Chow ring'' $\Gr^\bullet(X)$, or rather $\Gr^\bullet(X)_{\mathbb Q}$, constructed according to the principle explained in \hyperlink{sga6.e00.sec.4.2}{4.2}, the difficulty encountered there disappearing, moreover, because one is in characteristic zero, as is noted in \hyperlink{sga6.e14.sec.1}{Expos\'e XIV, no. 1}.

When $X$ and $Y$ are nonsingular, it is possible to give another version of the Riemann--Roch formula by proceeding as follows. First let $X$ be an arbitrary complex analytic space, and let $X^{\mathrm{top}}$ be the same space equipped with the sheaf of rings of continuous complex-valued functions. We therefore have a canonical morphism of ringed spaces
\[
X^{\mathrm{top}}\longrightarrow X,
\]
which, by the evident contravariant nature of the functor $K^\bullet(X)$ in the variable ringed space $X$, defines a homomorphism of rings
\sgaStableTarget{sga6.e00.eq.5-1}\begin{equation}\tag{5.1}
K^\bullet(X)\longrightarrow K^\bullet(X^{\mathrm{top}}).
\end{equation}
For simplicity suppose that $X^{\mathrm{top}}$ is compact; then, as in the case of a scheme admitting an ample invertible module, one proves (\hyperlink{sga6.e02.expose}{Expos\'e II}) that every perfect complex on $X^{\mathrm{top}}$ is globally isomorphic to a bounded complex, locally free in each degree, i.e. to a complex of complex vector bundles on $X^{\mathrm{top}}$, and hence that $K^\bullet(X^{\mathrm{top}})$ is canonically isomorphic to the well-known group $K(X^{\mathrm{top}})$ of topologists [1], defined in terms of complex vector bundles on the compact space $X^{\mathrm{top}}$. Thus one may interpret \hyperlink{sga6.e00.eq.5-1}{(5.1)} as a homomorphism from the $K^\bullet(X)$ defined in terms of the analytic structure to the well-known ring $K(X^{\mathrm{top}})$ of topologists. If now
\[
f:X\longrightarrow Y
\]
is a morphism of compact nonsingular analytic spaces, one knows how to associate to it, thanks to Atiyah--Hirzebruch [1], a group homomorphism
\sgaStableTarget{sga6.e00.eq.5-2}\begin{equation}\tag{5.2}
f^{\mathrm{top}}_*:K^\bullet(X^{\mathrm{top}})\longrightarrow K^\bullet(Y^{\mathrm{top}}).
\end{equation}
Using likewise the homomorphism $f_*$ of \hyperlink{sga6.e00.eq.3-1}{(3.1)}, defined by formula \hyperlink{sga6.e00.sec.3.2}{(3.2)}, which here implicitly uses Grauert's finiteness theorem [4] for a proper morphism of analytic spaces, one obtains a square of natural homomorphisms involving only groups of type $K^\bullet$, to the exclusion of rings of Chow-ring or cohomology type:
\sgaStableTarget{sga6.e00.eq.5-3}\begin{equation}\tag{5.3}
\begin{array}{ccc}
K^\bullet(X) & \longrightarrow & K^\bullet(X^{\mathrm{top}}) \\
\downarrow f_* && \downarrow f^{\mathrm{top}}_* \\
K^\bullet(Y) & \longrightarrow & K^\bullet(Y^{\mathrm{top}}).
\end{array}
\end{equation}
The Riemann--Roch type formula to which we are alluding consists in asserting the commutativity of this square. Note the difference in nature between this assertion, which does not neglect torsion phenomena and is situated in a group $K^\bullet(Y^{\mathrm{top}})$ of an essentially discrete nature, and the Riemann--Roch formula of the type considered in no. 4, which is situated in a group of ``Chow'' type no longer of a discrete nature, but which, on the other hand, must be tensored with $\mathbb Q$. Let us note that formula \hyperlink{sga6.e00.eq.5-3}{(5.3)} is proved at any rate when $X$ and $Y$ are nonsingular projective varieties [7], as is the formula of the type considered in no. 4, proved in this case in [RRR]. What is involved here is a plausible variant of the Atiyah--Singer index theorem, and of its generalizations due to Shih [7] and Atiyah, which obviously calls for a common generalization. Once the commutativity of \hyperlink{sga6.e00.eq.5-3}{(5.3)} is proved, one would immediately deduce from it, using the ``differentiable Riemann--Roch theorem'' of Atiyah--Hirzebruch [1], a cohomological variant of the complex analytic Riemann--Roch formula, namely the commutativity of the diagram
\sgaStableTarget{sga6.e00.eq.5-4}\begin{equation}\tag{5.4}
\begin{array}{ccc}
K^\bullet(X) & \xrightarrow{\;x\mapsto \ch(x)\Todd(T_X)\;} & H^{2\bullet}(X,\mathbb Q) \\
\downarrow f_* && \downarrow f_* \\
K^\bullet(Y) & \xrightarrow{\;y\mapsto \ch(y)\Todd(T_Y)\;} & H^{2\bullet}(Y,\mathbb Q),
\end{array}
\end{equation}
where the horizontal arrows are deduced from cohomological Chern classes, themselves deduced from the homomorphism \hyperlink{sga6.e00.eq.5-1}{(5.1)} and from the ordinary Chern-class homomorphism
\[
K^\bullet(X^{\mathrm{top}})\longrightarrow H^{2\bullet}(X,\mathbb Z),
\]
the first vertical arrow is the same as in \hyperlink{sga6.e00.eq.5-3}{(5.3)}, and the second is the Gysin homomorphism in rational cohomology.

\sgaStableTarget{sga6.e00.sec.6}\subsection*{6. The covariant invariant $K_\bullet(X)$}
In all that precedes, scarcely anything has been said except about the contravariant invariant $K^\bullet(X)$, to the exclusion of the covariant invariant $K_\bullet(X)$ whose existence was indicated in 3.1. If one regards $K^\bullet(X)$ as giving an analogue of the integral cohomology ring of a topological space, one should regard $K_\bullet(X)$ as giving an analogue of integral homology. The module structure of $K_\bullet(X)$ over $K^\bullet(X)$ then corresponds to the cap-product operation. The fact that, when $X$ is regular, the natural homomorphism $K^\bullet(X)\to K_\bullet(X)$ is an isomorphism should be regarded as the analogue of the Poincar\'e duality theorem for oriented topological varieties, asserting that cap product with the fundamental homology class defines an isomorphism from cohomology to homology. These analogies moreover suggest the usefulness of also defining, for a finite polyhedron $X$, say, a covariant invariant $K_\bullet(X)$ whose relations to integral homology -- spectral sequence, Chern homomorphism, and so on -- would be analogous to those existing between $K^\bullet(X)$ and integral cohomology. Nor does it seem excluded that there might exist, in terms of these groups $K_\bullet(X)$, variants of the differentiable Riemann--Roch theorem valid for spaces more general than compact differentiable manifolds.

Returning to the case where $X$ is an arbitrary noetherian scheme, one should endow $K_\bullet(X)$ with an increasing filtration, $\Fil_iK_\bullet(X)$ being generated by classes $\cl_\bullet(F)$, with $F$ a coherent sheaf on $X$ such that $\dim \operatorname{Supp}F\le i$. One will prove (\hyperlink{sga6.e10.expose}{Expos\'e X}) the compatibility of this filtration with the module structure and with the filtration of $K^\bullet(X)$, i.e. the formula
\sgaStableTarget{sga6.e00.eq.6-1}\begin{equation}\tag{6.1}
\Fil^iK^\bullet(X)\cdot \Fil_jK_\bullet(X)\subset \Fil_{j-i}K_\bullet(X),
\end{equation}
which allows the graded group associated with the filtered module $K_\bullet(X)$, denoted $\Gr_\bullet(X)$, to be endowed with a module structure over $\Gr^\bullet(X)$. Its elements may in fact be interpreted as classes of algebraic cycles on $X$, for an equivalence relation less fine than rational equivalence when $X$ is of finite type over a base field $k$. If
\[
f:X\longrightarrow Y
\]
is a proper morphism of noetherian schemes, then
\sgaStableTarget{sga6.e00.eq.6-2}\begin{equation}\tag{6.2}
f_*:K_\bullet(X)\longrightarrow K_\bullet(Y)
\end{equation}
is compatible with the filtrations and defines a homomorphism of graded groups
\sgaStableTarget{sga6.e00.eq.6-3}\begin{equation}\tag{6.3}
f_*:\Gr_\bullet(X)\longrightarrow \Gr_\bullet(Y),
\end{equation}
giving rise to a ``projection formula'', i.e. a formula which says that it is $\Gr^\bullet(Y)$-linear; this formula is obtained by passage to associated gradeds from the analogous projection formula for the homomorphism \hyperlink{sga6.e00.eq.6-2}{(6.2)}, which is $K^\bullet(Y)$-linear.

From the point of view proposed in the seminar, a manageable theory of intersections on noetherian schemes $X$, not necessarily regular, should consist in the combined study of the contravariant invariants $K^\bullet(X),\Gr^\bullet(X)$ and the covariant invariants $K_\bullet(X),\Gr_\bullet(X)$. These play roles of essentially different nature, and their properties complement each other. Thus the covariant groups $K_\bullet(X),\Gr_\bullet(X)$ are easier to compute for certain nonproper fibrations, thanks to the ``homotopy exact sequence'' which one can establish for them ([RRR] and \hyperlink{sga6.e09.expose}{Expos\'e IX}). For example, one proves (\hyperlink{sga6.e09.expose}{Expos\'e IX}) canonical isomorphisms such as
\[
K_\bullet(X[t])\simeq K_\bullet(X),\qquad \Gr_\bullet(X[t])\simeq \Gr_\bullet(X),
\]
which fail in general for $K^\bullet$ when $X$ is a nonregular scheme. On the other hand, $K^\bullet(X)$ and $\Gr^\bullet(X)$ lend themselves well to computations, thanks to their ring structures.

When one considers schemes $X$ proper over a field, the projection
\[
\pi_X:X\longrightarrow (\operatorname{point})
\]
defines a homomorphism, the ``Euler--Poincar\'e characteristic'',
\sgaStableTarget{sga6.e00.eq.6-4}\begin{equation}\tag{6.4}
\chi_X=\pi_{X*}:K_\bullet(X)\longrightarrow K_\bullet(\operatorname{point})=\mathbb Z,
\end{equation}
which induces on the subgroup $\Gr_0(X)=\Fil_0K_\bullet(X)$ the homomorphism ``degree of zero-cycles''
\begin{equation}\tag{\sgaStableTargetInline{sga6.e00.eq.6-5}6.5}
\deg_X:\Gr_0(X)\longrightarrow \mathbb Z.
\end{equation}
The multiplication homomorphism
\[
\Gr^i(X)\times \Gr_i(X)\longrightarrow \Gr_0(X),
\]
composed with the degree homomorphism \hyperlink{sga6.e00.eq.6-5}{(6.5)}, then furnishes a pairing
\sgaStableTarget{sga6.e00.eq.6-6}\begin{equation}\tag{6.6}
\Gr^i(X)\times \Gr_i(X)\longrightarrow \mathbb Z,
\end{equation}
which explains that, in certain respects, $\Gr^\bullet(X)$ and $\Gr_\bullet(X)$ play dual roles, just as cohomology and homology do. Thus, if $f:X\to Y$ is a morphism of proper algebraic schemes, the two homomorphisms
\[
f^*:\Gr^\bullet(Y)\longrightarrow \Gr^\bullet(X),\qquad f_*:\Gr_\bullet(X)\longrightarrow \Gr_\bullet(Y)
\]
are formally transposes of one another, in the sense of the pairings \hyperlink{sga6.e00.eq.6-6}{(6.6)}.

The homomorphisms \hyperlink{sga6.e00.eq.6-4}{(6.4)} and \hyperlink{sga6.e00.eq.6-5}{(6.5)}, and the pairing \hyperlink{sga6.e00.eq.6-6}{(6.6)} deduced from them, may be regarded as the source of numerical relations in intersection theory on proper algebraic schemes over a given field $k$. This point of view will be developed in some detail in \hyperlink{sga6.e10.expose}{Expos\'e X}, which is essentially independent of \hyperlink{sga6.e06.expose}{Expos\'es VI}, VII, VIII, centred on the proof of the Riemann--Roch theorem, and of \hyperlink{sga6.e09.expose}{Expos\'e IX}. Some notions concerning the behaviour of $K^\bullet(X)$ and $\Gr^\bullet(X)$ under specialization will also be given there.

\sgaStableTarget{sga6.e00.sec.7}\subsection*{7. Picard groups and the determinant functor}
As in every intersection theory, it is important to specify the place occupied in the theory by the Picard group, or group of classes of divisors in the classical terminology. Subject to a slight restriction, automatically satisfied when $X$ is for instance quasi-compact, one establishes in \hyperlink{sga6.e10.expose}{Expos\'e X} a canonical isomorphism
\sgaStableTarget{sga6.e00.eq.7-1}\begin{equation}\tag{7.1}
c_1:\Pic(X)\xrightarrow{\sim}\Gr^1(X),
\end{equation}
which shows the exact role of the Picard group in the theory we are proposing. This role justifies a more detailed study of the Picard group in \hyperlink{sga6.e12.expose}{Expos\'es XII} and XIII, which are logically independent of the text of the seminar. Kleiman, in \hyperlink{sga6.e13.expose}{Expos\'e XIII}, studies numerical equivalence and certain finiteness theorems for the Picard group, announced without proof in [6], using purely numerical techniques, not relying on the theory of the Picard functor and its representability, due to himself and Mumford, and markedly simpler than Grothendieck's initial proofs. The only result refractory to Kleiman's methods, and which he is obliged to use, is [6, C-07 (i)]; its proof and certain refinements, including certain representability theorems for the Picard functor, occupy \hyperlink{sga6.e12.expose}{Expos\'e XII}. The latter is moreover of a spirit and technique entirely foreign to the rest of the seminar, and is connected with [5]; this expos\'e is logically independent of all the others.

Finally, in Expos\'e XI, of a nature appreciably more general than XII and XIII and closer to the general spirit of the seminar, one studies, on an arbitrary locally ringed topos, a remarkable functor called the ``determinant functor'':
\sgaStableTarget{sga6.e00.eq.7-2}\begin{equation}\tag{7.2}
\det^\bullet:\Parf_{\mathrm{is}}(X)\longrightarrow \Inv(X)
\end{equation}
from the category of perfect complexes on $X$, whose morphisms are only the isomorphisms in the sense of the derived category $D(X)$, to the category of invertible modules on $X$. This functor is described, up to very little, by the requirement that it be of ``local nature'' and that it be given, for a bounded complex $L$ with locally free components, by the formula
\sgaStableTarget{sga6.e00.eq.7-3}\begin{equation}\tag{7.3}
\det^\bullet(L)=\bigotimes_i \det(L_i)^{(-1)^i},
\end{equation}
where $\det(L_i)$ denotes the exterior power of maximal degree of the locally free module $L_i$. Although this expos\'e is also logically independent of the rest of the seminar, except for \hyperlink{sga6.e01.expose}{Expos\'e I}, it nevertheless constitutes an important element of the general ``yoga'' proposed in the seminar and developed essentially in \hyperlink{sga6.e01.expose}{Expos\'es I} to XI.

\sgaStableTarget{sga6.e00.bibliography}\section*{Bibliography}
\begin{enumerate}[label={[\arabic*]},leftmargin=2.5em]
\item M. Atiyah and F. Hirzebruch, \emph{Vector bundles and homogeneous spaces}, Symposia in Pure Mathematics, vol. 3, Differential Geometry, pp. 7--38, 1961.
\item A. Borel and J.-P. Serre, \emph{Le th\'eor\`eme de Riemann--Roch}, Bull. Soc. math. France, t. 86, pp. 97--136 (1958).
\item H. Grauert, \emph{Ein Theorem der analytischen Garbentheorie}, Pub. Math. no. 5, pp. 5--64 (1960).
\item A. Grothendieck, \emph{Classes de faisceaux et th\'eor\`eme de Riemann--Roch}, mimeographed notes, Princeton 1957, cited as [RRR] and reproduced as an appendix to Expos\'e 0 of the seminar.
\item A. Grothendieck, \emph{Technique de descente et th\'eor\`emes d'existence en G\'eom\'etrie Alg\'ebrique}, in \emph{Fondements de la G\'eom\'etrie Alg\'ebrique}, extracts from the Bourbaki Seminar 1957--1962, Secr\'etariat Math\'ematique, Paris.
\item W. Shih, S\'eminaire de Topologie \`a l'I.H.E.S., 1966/67.
\item M. Atiyah and F. Hirzebruch, \emph{The Riemann--Roch theorem for analytic embeddings}, Topology, vol. 1, pp. 151--166 (1962).
\end{enumerate}



\clearpage



\sgaStableTarget{sga6.e00.appendix-to-expose-0-rrr}\section*{Appendix to Exposé 0: RRR}
\begin{center}
{\Large Classes of Sheaves and the Riemann--Roch Theorem}\par
\vspace{0.5em}
by A. Grothendieck\footnote{This is the literal reproduction of the report cited in the Introduction to the present seminar. The footnotes indicated by the signs (*), (**) were added in October 1967.}
\end{center}

\begin{center}
Demonstration of the Riemann--Roch theorem (1 November 1957).
\end{center}

\sgaStableTarget{sga6.e00.chapter-i-rings-formal-preliminaries}\subsection*{Chapter I. $\lambda$-rings (formal preliminaries)}
\begin{longtable}{@{}p{0.74\linewidth}r@{}}
1. Definitions & 1\\
2. Examples & 4\\
3. The $\lambda$-ring defined by a graded ring & 8\\
4. The operations $\lambda^p(N,x)$ & 13\\
\end{longtable}

\sgaStableTarget{sga6.e00.chapter-ii-classes-of-coherent-algebraic-sheaves}\subsection*{Chapter II. Classes of coherent algebraic sheaves and Chern classes}
\begin{longtable}{@{}p{0.74\linewidth}r@{}}
1. The theory of Chow & 18\\
2. Definition of the Chern classes of coherent algebraic sheaves & 22\\
3. Functorial generalities on $K(X)$ & 27\\
4. Some technical results & 34\\
5. Sheaf-theoretic definition of Chern classes. Application to the study of immersion morphisms & 43\\
6. The Riemann--Roch theorem & 49
\end{longtable}

\sgaStableTarget{sga6.e00.chapter-i-rings-formal-preliminaries.occ2}\section*{Chapter I. $\lambda$-rings (formal preliminaries)}

\sgaStableTarget{sga6.e00.sec.1.occ2}\subsection*{1. Definitions}

A \emph{$\lambda$-ring}\footnote{In this seminar we rather say \emph{pre-$\lambda$-ring} (\hyperlink{sga6.e05.def.2-1}{V 2.1}), reserving the name $\lambda$-ring for those satisfying the supplementary condition of page 3 (cf. \hyperlink{sga6.e05.def.2-4}{V 2.4}).} means a commutative ring $K$ with unit, provided with a family of maps
\sgaStableTarget{sga6.e00.eq.1-1.occ2}\begin{equation}\tag{1.1}
\lambda^i:K\longrightarrow K\qquad (i\text{ an integer }\geq 0)
\end{equation}
satisfying the following conditions:
\sgaStableTarget{sga6.e00.eq.1-2.occ2}\begin{equation}\tag{1.2}
\lambda_x^0=1,\qquad \lambda_x^1=x,\qquad
\lambda^n(x+y)=\sum_{i=0}^n \lambda^i_x\lambda^{n-i}_y .
\end{equation}
For every $x\in K$ put
\sgaStableTarget{sga6.e00.eq.1-3.occ2}\begin{equation}\tag{1.3}
\lambda_t(x)=\sum_{n\geq 0}\lambda^n(x)t^n\in K[[t]].
\end{equation}
The preceding relations then say that $x\mapsto \lambda_t(x)$ is a homomorphism from the additive group $K$ to the multiplicative group $1+K[[t]]^+$ of formal series over $K$ with augmentation $1$, lifting the canonical homomorphism
\[
1+\sum_{i>0}x_it^i\longmapsto x_1.
\]

Thus on the ring $\Z$ there exists a unique $\lambda$-structure compatible with its ring structure and such that
\sgaStableTarget{sga6.e00.eq.1-4}\begin{equation}\tag{1.4}
\lambda_t(1)=1+t .
\end{equation}
One then has
\sgaStableTarget{sga6.e00.eq.1-5}\begin{equation}\tag{1.5}
\lambda_t(n)=(1+t)^n,
\end{equation}
whence
\[
\lambda^i(n)=\binom{n}{i}.
\]

The notion of a $\lambda$-homomorphism of $\lambda$-rings is clear; an \emph{augmented $\lambda$-ring} is a $\lambda$-ring provided with a $\lambda$-homomorphism to $\Z$.

To define the notion of a \emph{special} $\lambda$-ring, the following considerations are useful. Let $k$ be a commutative ring with unit and let
\[
K=1+k[[t]]^+
\]
be the group of formal series with coefficients in $k$ and augmentation $1$. We shall introduce on it a $\lambda$-ring law whose additive structure is the multiplication of formal series. The multiplication of this ring will be denoted $f\circ g$. It is entirely determined by the condition that the coefficients $c_n=c_n(f\circ g)$ of $f\circ g$ are given by universal polynomials with integral coefficients in the coefficients $a_i=c_i(f)$ and $b_i=c_i(g)$ of $f$ and $g$, respectively:
\sgaStableTarget{sga6.e00.eq.1-6}\begin{equation}\tag{1.6}
c_n=P_n(a_1,\ldots,a_n;b_1,\ldots,b_n),
\end{equation}
that it be distributive with respect to the ``addition'' $fg$, and be given for linear factors by
\sgaStableTarget{sga6.e00.eq.1-7}\begin{equation}\tag{1.7}
(1+at)\circ(1+bt)=1+abt.
\end{equation}
The $P_n$ are computed by calculations with symmetric functions; $P_n$ is isobaric of weight $n$ in the $a_i$ and isobaric of weight $n$ in the $b_i$ (where $a_i$ and $b_i$ have weight $i$), and moreover is symmetric in $a=(a_i)$ and $b=(b_i)$.

Similarly, the $\lambda^i(f)$ are determined without ambiguity by the condition that the coefficients
\[
c_n\bigl(\lambda^i(f)\bigr)=c_n\bigl(\lambda^i(f)\bigr)
\]
of $\lambda^i(f)$ be calculated by universal polynomials with integral coefficients in the coefficients $c_n(f)=a_n$ of $f$:
\sgaStableTarget{sga6.e00.eq.1-8}\begin{equation}\tag{1.8}
c_n=P_{i,n}(a_1,\ldots,a_{in}),
\end{equation}
that the relations \hyperlink{sga6.e00.eq.1-2.occ2}{(1.2)} be satisfied, and that, for $i>1$, one have
\sgaStableTarget{sga6.e00.eq.1-9}\begin{equation}\tag{1.9}
\lambda^i(1+at)=1
\end{equation}
($1$ is the zero element of the ring $K=1+k[[t]]^+$). The $P_{i,n}$ are again computed by calculations with symmetric functions, and one observes that $P_{i,n}$ is an isobaric polynomial of weight $in$ in the $a_i$. Provided with the preceding operations, $K$ becomes a $\lambda$-ring; its zero element is $1$ and its unit element is $1+t$. Also note the relation
\sgaStableTarget{sga6.e00.eq.1-7bis}\begin{equation}\tag{1.7 bis}
(1+at)\circ f=f(at).
\end{equation}

Let now $K$ be an arbitrary $\lambda$-ring. We say that $K$ is \emph{special} if the additive homomorphism $\lambda_t$ from $K$ to $1+K[[t]]^+$ is a $\lambda$-homomorphism. Thus this means that one has the formulas
\sgaStableTarget{sga6.e00.eq.1-10}\begin{equation}\tag{1.10}
\lambda_t(1)=1+t,\quad
\lambda_t(xy)=\lambda_t(x)\circ\lambda_t(y),\quad
\lambda_t(\lambda^i(x))=\lambda^i(\lambda_t(x)),
\end{equation}
or, more explicitly,
\sgaStableTarget{sga6.e00.eq.1-11}\begin{equation}\tag{1.11}
\begin{cases}
\lambda^i(1)=0 & \text{if }i>1,\\
\lambda^n(xy)=P_n(\lambda^1x,\ldots,\lambda^n x;\lambda^1y,\ldots,\lambda^n y),\\
\lambda^n(\lambda^i(x))=P_{i,n}(\lambda^1x,\ldots,\lambda^{in}x),
\end{cases}
\end{equation}
where the $P_n$ and $P_{i,n}$ are the universal polynomials considered above. In particular the formulas \hyperlink{sga6.e00.eq.1-5}{(1.5)} follow. One may show that the $\lambda$-ring $1+k[[t]]^+$ constructed from a commutative ring $k$ is special, but we shall not need this fact.

\sgaStableTarget{sga6.e00.sec.2.occ2}\subsection*{2. Examples}

\sgaStableTarget{sga6.e00.ex.2-1}1) Let $G$ be a group and $k$ a commutative ring. Consider the group $K(G)$, the quotient of the free group generated by the classes of representations of $G$ in finite-type $k$-modules by the subgroup generated by elements of the form
\[
(\rho\oplus\rho')-(\rho)-(\rho')
\]
(where $\oplus$ denotes the direct sum). Tensor product and exterior powers define on $K(G)$ a $\lambda$-ring structure. One may also pass to the quotient by elements of the form
\[
(\rho)-(\rho')-(\rho''),
\]
where $\rho$ is a representation which is an extension of $\rho'$ by $\rho''$, trivial as an extension of $k$-modules. One obtains the group $K_r(G)$ of reduced classes of representations of $G$, and one checks that the $\lambda$-ring law passes to the quotient.

When $k$ is an algebraically closed field of characteristic $0$,\footnote{It is unnecessary for $k$ to be algebraically closed; cf. \hyperlink{sga6.e06.thm.3-3}{VI 3.3} for a still more general result, encompassing all those that will be used in the present report.} one can show that $K(G)$, and a fortiori $K_r(G)$, is a special $\lambda$-ring. This is no longer true if the characteristic is not $0$; however, if $k$ is algebraically closed, $K_r(G)$ is still a special $\lambda$-ring. To prove this one is reduced to the case where $G$ is a product $\Gl(m,k)\times\Gl(n,k)$ and where only rational representations of $G$ are considered; our assertion follows from the explicit determination of $K_r(G)$ to be made below (\hyperlink{sga6.e00.ex.2-3}{Example 3)}. Finally, when $k$ has characteristic $0$, complete reducibility of rational representations of $G$ shows that $K(G)=K_r(G)$.

Variants of the preceding example: $G$ is an algebraic group defined over $k$ and one considers only rational representations of $G$ defined over $k$; one may consider the case of a topological group, or of a Lie group, and continuous representations, etc.\footnote{One could build other examples of $\lambda$-rings; for example, the set of equivalence classes of quadratic forms, the set of representation classes of a Lie algebra, etc. We promise that we shall not need all these examples in order to treat the Riemann--Roch theorem.}

\sgaStableTarget{sga6.e00.ex.2-2}2) Let $X$ be an algebraic variety, irreducible or not, defined over the field $k$. We may consider the quotient group of the free group generated by the classes of vector bundles on $X$ (defined over $k$) by the subgroup generated by the elements of the form
\[
(E\oplus E')-(E)-(E'),
\]
on which the $\lambda$-ring operations are defined by tensor product and exterior powers. If $X$ is complete, this is a free group generated by the indecomposable vector bundles on $X$ (thanks to Krull--Remak--Schmidt); knowledge of this ring and of its basis is equivalent to the complete classification of vector bundles on $X$ and to knowledge of the elementary tensor operations on these bundles. If $k$ is algebraically closed of characteristic $0$, all other tensor operations on classes of vector bundles are then known by means of what follows in 3). One may make a more drastic passage to the quotient by the group generated by the elements
\[
(E)-(E')-(E''),
\]
where $E$ is a vector bundle extension of $E'$ by $E''$. One then obtains a $\lambda$-ring, denoted $K(\xi(X))$, where $\xi(X)$ denotes the additive category of vector bundles on $X$ defined over $k$. Using the results of \hyperlink{sga6.e00.ex.2-1}{Example 1)}, one shows that, if $k$ is algebraically closed,\footnote{Unnecessary condition; cf. the preceding footnote.} the $\lambda$-ring $K(\xi(X))$ is special, and that the same is true of the first $\lambda$-ring defined here if the characteristic of $k$ is $0$, but not in the general case.

\sgaStableTarget{sga6.e00.ex.2-3}3) Determination of $K(G)$ and of $K_r(G)$ in some important cases. Suppose the field $k$ is algebraically closed and that the group $G$ is affine algebraic, connected, and ``globally reductive'' (the radical of $G$ is isomorphic to $(k^*)^s$), and consider only rational linear representations. In characteristic $0$, complete reducibility of rational linear representations of $G$ shows that $K(G)=K_r(G)$ is the free group generated by the classes of irreducible rational representations. If $T$ is a maximal torus of $G$, then
\[
K(T)=K_r(T)=\Z(\That),
\]
the algebra of the group $\That$ of rational characters of $T$; one has $\That\simeq \Z^r$ if $r$ is the rank of $G$ (this is independent of the characteristic). The Weyl group $W$ operates on $T$, hence on $K(T)$, in a way which depends only on the operations of $W$ on the free group $\That$. There is an evident homomorphism of $\lambda$-rings, deduced from the injection of $T$ into $G$:
\sgaStableTarget{sga6.e00.eq.1-13}\begin{equation}\tag{1.13}
K_r(G)\longrightarrow K_r(T)^W=K(T)^W,
\end{equation}
where $K(T)^W$ denotes the set of fixed points of $W$ in $K(T)$. The theory of linear representations of $G$\footnote{Cf. Chevalley Seminar 1956/58, \emph{Groupes de Lie algébriques} (École Normale Supérieure).} then proves the following theorem.

\sgaStableTarget{sga6.e00.thm.1-1}\textbf{Theorem 1.1.} The homomorphism \hyperlink{sga6.e00.eq.1-13}{(1.13)} is an isomorphism.

Thus one sees that $K_r(G)$, once the ``type'' of $G$ is known (characterized by $W$ operating on a group $\Z^r$), does not depend on the characteristic, and that it is a special $\lambda$-ring. There is a canonical basis of $K_r(G)$ corresponding to the orbits of $W$ in $\That$; another basis is obtained as follows. Let $(\rho_i)_{1\leq i\leq r'}$ be the classes of representations of $G$ corresponding to the fundamental representations of $G/\rad G$, and let $(\sigma_j)_{1\leq j\leq r''}$ be a basis of the group of rational homomorphisms from $G$ to $k^*$. Thus $r'+r''=r$ (the rank of $G$).

\sgaStableTarget{sga6.e00.cor.unnamed}\textbf{Corollary.} The ring $K_r(G)$ is identified with the subring
\[
\Z[\rho_1,\ldots,\rho_{r'};\sigma_1,\ldots,\sigma_{r''},\sigma_1^{-1},\ldots,\sigma_{r''}^{-1}]
\]
of the field of rational functions in the $\rho_i$ and the $\sigma_j$ with coefficients in $\Q$.\footnote{The derived group $G'$ of $G$ must be supposed simply connected for this statement to be correct.}

In particular, $K_r(G)$ has no zero divisors. If $G=\prod_i\Gl(n_i,k)$, the ring $K_r(G)$ is made explicit as follows. Let $\rho_i$ be the representation of $G$ deduced from the identical representation of $\Gl(n_i,k)$. If one then considers the field of rational functions over $\Q$ in the indeterminates $\lambda^j(\rho_i)$, $1\leq j\leq n_i$, then $K_r(G)$ is identified with the subring generated by these indeterminates and the inverses of the $\lambda^{n_i}(\rho_i)$. This justifies our assertion that, at least in characteristic $0$, the operations $\lambda^i$ make it possible to reconstruct all tensor operations on vector bundles; this remains true over an algebraically closed field of arbitrary characteristic, provided one takes reduced classes.

\sgaStableTarget{sga6.e00.rem.unnamed}\textbf{Remarks.} a) In \hyperlink{sga6.e00.ex.2-1}{Example 1)}, the ring $K(G)$ is augmented, the augmentation being defined by the degree of representations; by passage to the quotient, this augmentation defines an augmentation on $K_r(G)$. Likewise, the $\lambda$-rings considered in \hyperlink{sga6.e00.ex.2-2}{Example 2)} are augmented when $X$ is $k$-irreducible; in $K(\xi(X))$, the augmentation is defined by the rank of a vector bundle, that is, the dimension of its fibres.

b) When the field $k$ is not algebraically closed, it is plausible that $K(G)$ is a special $\lambda$-ring in characteristic $0$, and likewise for $K_r(G)$ in all characteristics.

c) In non-zero characteristic, \hyperlink{sga6.e00.thm.1-1}{Theorem 1.1} was graciously supplied by Chevalley.

\sgaStableTarget{sga6.e00.sec.3.occ2}\subsection*{3. The $\lambda$-ring defined by a graded ring}

Let $A$ be a graded commutative ring with unit. To simplify, assume that the degrees of $A$ are positive and that $A^0=\Z$. Denote by $\Ahat$ the product of the $A^i$ (this is a commutative ring with unit), and by $\Ahp$ the kernel of the evident augmentation $\Ahat\to A^0=\Z$. Then $1+\Ahp$ is a subgroup of the multiplicative group of invertible elements of $\Ahat$. Consider the abelian group
\sgaStableTarget{sga6.e00.eq.1-14}\begin{equation}\tag{1.14}
\widetilde A=\Z\times (1+\Ahp),
\end{equation}
whose elements will be denoted $[n,x]$, with $n\in\Z$ and
\[
x=1+\sum_{i\geq 1}x^i\in 1+\Ahp\qquad (x^i\in A^i).
\]
We write this group additively, thus:
\sgaStableTarget{sga6.e00.eq.1-15}\begin{equation}\tag{1.15}
[n,x]+[n',x']=[n+n',xx'].
\end{equation}
The zero element of this group is $[0,1]$. We shall introduce on it an augmented $\lambda$-ring structure, compatible with its additive structure, the augmentation being $\varepsilon:[n,x]\mapsto n$. The product $[m,x][n,y]$ is completely characterized by the fact that it is written
\[
[mn,x^n y^m(x*y)],
\]
where $x*y$ is an element of $1+\Ahp$ whose components are expressed as universal polynomials with integral coefficients in the $x^i$ and the $y^i$:
\sgaStableTarget{sga6.e00.eq.1-16}\begin{equation}\tag{1.16}
(x*y)^i=Q_i(x^1,\ldots,x^i;y^1,\ldots,y^i)\qquad (i\geq 1),
\end{equation}
the polynomial $Q_i$ being isobaric of weight $i$. Moreover, $x*y$ must be distributive with respect to the ``addition'' $xy$ and, for ``linear'' factors, must reduce to
\sgaStableTarget{sga6.e00.eq.1-17}\begin{equation}\tag{1.17}
[1,1+x^1][1,1+y^1]=[1,1+x^1+y^1].
\end{equation}
The polynomials $Q_i$ are evaluated by calculations with elementary symmetric functions. This is the calculation of the Chern classes of a tensor product of vector bundles by means of the Chern classes $x^i$ and $y^i$ of the factors and of their respective dimensions $m$ and $n$. From \hyperlink{sga6.e00.eq.1-17}{(1.17)} one deduces
\sgaStableTarget{sga6.e00.eq.1-17bis}\begin{equation}\tag{1.17 bis}
(1+x^1)*(1+y^1)=\frac{1+x^1+y^1}{(1+x^1)(1+y^1)},
\end{equation}
and an arbitrary product $x*y$ is computed term by term by decomposing formally
\[
x=\prod_{i=1}^n(1+\alpha_i),\qquad y=\prod_{i=1}^n(1+\beta_i),
\]
where the $\alpha_i$ and $\beta_i$ have degree $1$.

In this way $\widetilde A$ becomes a commutative augmented ring with unit $[1,1]$; it may be interpreted as the ring obtained by adjoining a unit to the ring $1+\Ahp$. Moreover, on $\widetilde A$ one defines a filtration by putting $\widetilde A^0=\widetilde A$ and denoting by $\widetilde A^n$ the set of the $[0,x]$ with $x^i=0$ for $1\leq i<n$. This filtration is compatible with the product, for
\sgaStableTarget{sga6.e00.eq.1-18}\begin{equation}\tag{1.18}
\left(1+\sum_{i\geq m}x^i\right)*\left(1+\sum_{j\geq n}y^j\right)
=1-\frac{(m+n-1)!}{(m-1)!(n-1)!}x^m y^n+\cdots,
\end{equation}
the omitted terms having degree $>m+n$. This formula shows that the graded ring $G(\widetilde A)$ associated with $\widetilde A$ is identified with $A$ from the additive point of view, but not from the multiplicative point of view. One can only say that the map $\eta'$ from $A$ to $G(\widetilde A)$ defined by
\sgaStableTarget{sga6.e00.eq.1-19}\begin{equation}\tag{1.19}
\eta'(x^i)=(-1)^{i-1}(i-1)!\,x^i\qquad (x^i\in A^i)
\end{equation}
is a ring homomorphism.

On $\widetilde A$ introduce operations $\lambda^i$ satisfying the axioms of $\lambda$-rings (Section 1, formulas \hyperlink{sga6.e00.prop.1-2}{(1.2)}), well determined by the condition that the components of $\lambda^i[0,x]$ be deduced from the components $x^i$ of $x$ by universal polynomials with integral coefficients:
\sgaStableTarget{sga6.e00.eq.1-20}\begin{equation}\tag{1.20}
\lambda^i[0,x]=[0,\lambda^i x],\qquad (\lambda^i x)^{(n)}=Q_{i,n}(x^1,\ldots,x^n),
\end{equation}
where $Q_{i,n}$ is isobaric of weight $n$, and such that
\sgaStableTarget{sga6.e00.eq.1-21}\begin{equation}\tag{1.21}
\lambda^i[1,1+x^1]=0\qquad\text{for }i>1.
\end{equation}
The polynomials $Q_{i,n}$ are evaluated by a calculation with elementary symmetric functions: it is the calculation of the Chern classes of the exterior powers of a vector bundle by means of the Chern classes and the rank of that bundle. One then observes that $\widetilde A$ is a special augmented $\lambda$-ring and that the $\widetilde A^i$ are stable under the operations $\lambda^j$.

The augmented $\lambda$-ring $\widetilde A$ may be considered as a functor in $A$. Consequently, the principal automorphism
\[
x=\sum_{i\geq 0}x^i\longmapsto \check x=\sum_{i\geq 0}(-1)^i x^i
\]
of $A$ defines an automorphism of $\widetilde A$, denoted by the same symbol $\check{\ }$.

\sgaStableTarget{sga6.e00.prop.1-2}\textbf{Proposition 1.2.} Let $x\in\widetilde A$ be of the form
\[
x=[n,1+x^1+\cdots+x^n]
\]
(so $x^i=0$ for $i>n$). Then $\lambda^i(x)=0$ for $i>n$, one has $\lambda^n(x)=[1,1+x^1]$, and finally
\sgaStableTarget{sga6.e00.eq.1-22}\begin{equation}\tag{1.22}
\sum_{i=0}^{n}(-1)^i\lambda^i(x)=[0,1-(n-1)!x^n+\cdots].
\end{equation}
From this proposition, and from the fact that the $\widetilde A^i$ are stable under the operations $\lambda^j$, it follows that formula \hyperlink{sga6.e00.eq.1-22}{(1.22)} is valid whenever $\varepsilon(x)=n$, whether or not $x^i=0$ for $i>n$.

Formula \hyperlink{sga6.e00.eq.1-22}{(1.22)} admits the following generalizations. Let $K$ be any $\lambda$-ring, and let $n$ be arbitrary. For $x\in K$, put
\sgaStableTarget{sga6.e00.eq.1-23}\begin{equation}\tag{1.23}
\gamma^n(x)=(-1)^n\sum_{i=0}^{n}(-1)^i\lambda^i(x+n)=\lambda^n(x+n-1)
\end{equation}
(the identity of the last two members follows from an immediate induction). In particular $\gamma^0(x)=1$ and $\gamma^1(x)=x$. Next put
\sgaStableTarget{sga6.e00.eq.1-24}\begin{equation}\tag{1.24}
\gamma_t(x)=\sum_{n\geq 0}\gamma^n(x)t^n\in K[[t]].
\end{equation}
One deduces the formula
\sgaStableTarget{sga6.e00.eq.1-25}\begin{equation}\tag{1.25}
\gamma_t(x)=\lambda_{t/(1-t)}(x),
\end{equation}
which is equivalent to
\sgaStableTarget{sga6.e00.eq.1-25bis}\begin{equation}\tag{1.25 bis}
\lambda_s(x)=\gamma_{s/(1+s)}(x).
\end{equation}
This proves that the $\lambda^i$ are expressed in terms of the $\gamma^i$, and that
\[
\gamma_t(x+y)=\gamma_t(x)\gamma_t(y).
\]
Consequently the maps $\gamma^i$ define a new $\lambda$-ring structure on $K$. Formula \hyperlink{sga6.e00.eq.1-22}{(1.22)} then shows that, for $K=\widetilde A$ and every $x\in\widetilde A$, one has
\sgaStableTarget{sga6.e00.eq.1-22bis}\begin{equation}\tag{1.22 bis}
\gamma^n(x-\varepsilon(x))=[0,1+(-1)^{n-1}(n-1)!x^n+\cdots],
\end{equation}
and hence the component $x^n$ is obtained, up to the factor $(-1)^{n-1}(n-1)!$, by reducing modulo $\widetilde A^{n+1}$ the element $\gamma^n(x-\varepsilon(x))$ of $\widetilde A^n$.

\sgaStableTarget{sga6.e00.the-chern-homomorphism}\paragraph{The Chern homomorphism.}
For every graded ring $A$ satisfying the preceding conditions, we shall define a homomorphism of augmented rings
\sgaStableTarget{sga6.e00.eq.1-26}\begin{equation}\tag{1.26}
\ch:\widetilde A\longrightarrow \Ahat\otimes\Q.
\end{equation}
This homomorphism is defined without ambiguity by the condition that it be an additive homomorphism, send $1$ to $1$, and be such that the components of $\ch(x)$ in degree $>0$ are given by universal isobaric polynomials with rational coefficients in the components of $x$, and that
\sgaStableTarget{sga6.e00.eq.1-27}\begin{equation}\tag{1.27}
\ch([1,1+x^1])=\exp x^1=\sum_{n\geq 0}(x^1)^n/n!.
\end{equation}
One immediately obtains the explicit expression of $\ch$. Let $\eta$ be the additive endomorphism of $\Ahat\otimes\Q$ defined, for $x^i$ of degree $i$, by
\sgaStableTarget{sga6.e00.eq.1-28}\begin{equation}\tag{1.28}
\eta(x^i)=\frac{1}{(-1)^{i-1}(i-1)!}x^i .
\end{equation}
Then
\sgaStableTarget{sga6.e00.eq.1-29}\begin{equation}\tag{1.29}
\ch\left([n,1+\sum_{i\geq 1}x^i]\right)
=n+\eta\left(\log\left(1+\sum_{i\geq 1}x^i\right)\right).
\end{equation}
Since this formula can be inverted, at least if $A$ has only finitely many homogeneous components, one sees that in this case one has an isomorphism $\ch$ from the ring $\widetilde A\otimes\Q$ onto the ring $\Ahat\otimes\Q$, which completely elucidates the structure of $\widetilde A\otimes\Q$.

Recall that, according to Hirzebruch [1], a formal series $f\in\Q[[t]]$ defines by universal polynomial formulas an additive homomorphism
\[
1+\Ahp\longrightarrow 1+(\Ahat\otimes\Q)^+,
\]
denoted $\mathfrak C_f$. When $f(t)=t/(1-\exp(-t))$, we write simply $\mathfrak C_f=\mathfrak C$ (the initial of Todd). We extend the definition of $\mathfrak C(x)$ to $x\in\widetilde A$. This said, we have the following result.

\sgaStableTarget{sga6.e00.prop.1-3}\textbf{Proposition 1.3.} Let $N\in\widetilde A$ be of the form
\[
N=[q,1+N^1+\cdots+N^q]
\]
(so $N^i=0$ for $i>q$). If one puts
\[
\lambda_{-1}(N)=\sum_{i\geq 0}(-1)^i\lambda^i(N)=\sum_{i=0}^{n}(-1)^i\lambda^i(N),
\]
then
\sgaStableTarget{sga6.e00.eq.1-30}\begin{equation}\tag{1.30}
\ch(\lambda_{-1}(N))=(-1)^qN^q\mathfrak C(\check N)^{-1},
\end{equation}
which is equivalent to
\sgaStableTarget{sga6.e00.eq.1-30bis}\begin{equation}\tag{1.30 bis}
\ch(\lambda_{-1}(\check N))=N^q\mathfrak C(N)^{-1}.
\end{equation}

\sgaStableTarget{sga6.e00.sec.4.occ2}\subsection*{4. The operations $\lambda^p(N,x)$}

Let $A$ be the ring of formal series, with coefficients in $\Z$, in two infinite sequences of indeterminates denoted $\lambda^iN$ and $\lambda^i x$ ($i\geq 1$). Put $\lambda^1N=N$ and $\lambda^1x=x$, and introduce on $A$ the unique special $\lambda$-ring structure for which
\[
\lambda^i(N)=\lambda^iN,
\qquad
\lambda^i(x)=\lambda^i x,
\]
and for which the operations $\lambda^i$ are continuous. If one restricted oneself to polynomials in the indeterminates $\lambda^iN$ and $\lambda^i x$, one would have the $\lambda$-ring freely generated by $N$ and $x$; our ring $A$ is its completion. One may then form
\sgaStableTarget{sga6.e00.eq.1-31}\begin{equation}\tag{1.31}
\lambda_{-1}(N)=\sum_{i\geq 0}(-1)^i\lambda^iN,
\end{equation}
which is a formal series with constant term $1$, hence invertible. Thus an element $\lambda^p(N,x)$ of $A$ can be defined without ambiguity by putting
\sgaStableTarget{sga6.e00.eq.1-32}\begin{equation}\tag{1.32}
\lambda^p(N,x)\lambda_{-1}(N)=\lambda^p(x\lambda_{-1}(N)).
\end{equation}

\sgaStableTarget{sga6.e00.thm.1-4}\textbf{Theorem 1.4.} Let $J_q$ be the closed ideal of $A$ generated by the $\lambda^iN$ with $i>q$. Identify $A/J_q$ with the ring of formal series in the $\lambda^i x$ and the $\lambda^jN$ ($1\leq j\leq q$). Then $\lambda^p(N,x)$ reduced modulo $J_q$ is a polynomial, isobaric of weight $p$ with respect to the variables $\lambda^i x$.

Let $N$ and $F$ be two vector spaces of finite dimensions $q,q'$ over an algebraically closed field $k$ of characteristic $0$. Let $N^{(p)}$ be the kernel of the map
\[
(a_1,\ldots,a_p)\longmapsto a_1+\cdots+a_p
\]
from $N^p$ to $N$. We make the symmetric group $\mathfrak S_p$ operate on $N^p$, hence on $N^{(p)}$, hence on $\Lambda(N^{(p)})$, and then on
\[
\Lambda(N^{(p)})\otimes\bigotimes^p F.
\]
Consider this vector space as a graded module over $\mathfrak S_p\times\Gl(N)\times\Gl(F)$, and take its ``alternating part'', which is a graded module over $\Gl(N)\times\Gl(F)=G$. Then
\sgaStableTarget{sga6.e00.eq.1-33}\begin{equation}\tag{1.33}
\lambda^p(N,F)=E_G\left(\left(\Lambda(N^{(p)})\otimes\bigotimes^p F\right)^{\mathrm{alt}}\right).
\end{equation}

N.B. If $M$ is a complex of $G$-modules, $E_G(M)$ denotes the alternating sum in $K(G)$ of the classes of the components of $M$, that is, the Euler--Poincaré characteristic of $M$. In formula \hyperlink{sga6.e00.eq.1-33}{(1.33)}, the first member $\lambda^p(N,F)$ denotes the element of the ring $K(G)$ obtained by substituting, for the variables $\lambda^i_N$ and $\lambda^i_F$, the classes in $K(G)$ of the $G$-modules $\lambda^i(N)$ and $\lambda^i(F)$, respectively.

\emph{Proof.} Expanding the second member of \hyperlink{sga6.e00.eq.1-32}{(1.32)} by the second formula \hyperlink{sga6.e00.eq.1-11}{(1.11)}, one sees that the first assertion is proved if one proves the analogous assertion obtained by setting $x=1$, i.e. if one proves that in the $\lambda$-ring $\Z[\lambda^1N,\ldots,\lambda^qN]$, where $\lambda^iN=0$ for $i>q$, the element $\lambda^p(\lambda_{-1}(N))$ is divisible by $\lambda_{-1}(N)$. A priori, the quotient
\[
\lambda^p(N,1)=\lambda^p(\lambda_{-1}(N))/\lambda_{-1}(N)
\]
is a formal series contained in the quotient field $K$ of the ring of formal series in the $\lambda^iN$ for $1\leq i\leq q$. On the other hand, also denote by $N$ a vector space of dimension $q$ over an algebraically closed field $k$ of characteristic $0$, and consider the element
\[
\mathscr L^p(N)=E_G((\Lambda(N^{(p)}))^{\mathrm{alt}})\in K(G),
\]
where now $G=\Gl(N)$. We know (Section 2, \hyperlink{sga6.e00.thm.1-1}{Theorem 1.3)} that $K(G)$ is identified with the subring of $K$ generated by the $\lambda^i_N$ for $1\leq i\leq q$ and by the inverse of $\lambda_N^q$. Moreover, in this ring one verifies easily that
\[
\mathscr L^p(N)\lambda_{-1}(N)=\lambda^p(\lambda_{-1}(N)).
\]
Since we are in a large field $K$ and $\lambda_{-1}(N)\neq 0$, it follows that $\mathscr L^p(N)=\lambda^p(N,1)$. Both members lie in the intersection of $K(G)$ with the ring of formal series in the $\lambda_N^i$, hence are polynomials in the $\lambda_N^i$.

To prove \hyperlink{sga6.e00.eq.1-33}{(1.33)} in the general case, one verifies that the second member satisfies in $K(G)$ the formula analogous to \hyperlink{sga6.e00.eq.1-32}{(1.32)}, that is, that its product by $\lambda_{-1}(N)$ is $\lambda^p(F\lambda_{-1}(N))$. Likewise, the product of $\lambda^p(N,F)\in K(G)$ by $\lambda_{-1}(N)$ is $\lambda^p(F\lambda_{-1}(N))$, as follows from the fact that $K(G)$ is special. Since $K(G)$ is integral (cf. Section 2), and $\lambda_{-1}(N)\neq 0$ (same reference), formula \hyperlink{sga6.e00.eq.1-33}{(1.33)} is proved. QED.

\sgaStableTarget{sga6.e00.rem.unnamed.occ2}\textbf{Remarks.} a) When the dimension $q'$ of $F$ is at least $p$, formula \hyperlink{sga6.e00.eq.1-33}{(1.33)} characterizes the polynomial $\lambda^p(N,x)$ in the $\lambda_N^i$ ($1\leq i\leq q$) and the $\lambda_x^j$ ($1\leq j\leq p$).

b) It is plausible that \hyperlink{sga6.e00.eq.1-33}{(1.33)} remains valid if the field $k$ is not algebraically closed.\footnote{This is in fact proved by the same method.} One would have to resolve this question if one wanted to prove the Riemann--Roch theorem over a non-algebraically-closed field of characteristic $0$ by the method to be explained. The passage from an $\mathfrak S_p$-module to its alternating part has reasonable meaning only when the base field has characteristic $0$.

Let now $K$ be any special $\lambda$-ring and let $N$ be an element of $K$ such that $\lambda^iN=0$ for all sufficiently large $i$. Then for every $x\in K$ and every integer $p\geq 1$ one can consider the elements $\lambda^p(N,x)$, which are expressed by universal polynomials with integral coefficients in the $\lambda_N^i$ and the $\lambda_x^j$. One has $\lambda^1(N,x)=x$ and formula \hyperlink{sga6.e00.eq.1-32}{(1.32)} in $K$. To express the properties of the operations $x\mapsto\lambda^p(N,x)$, it is convenient to introduce a new $\lambda$-ring, denoted $K_N$. As a commutative ring, $K_N$ is obtained by adjoining a unit to the ring whose underlying additive group is $K$ and whose multiplication is given by
\sgaStableTarget{sga6.e00.eq.1-34}\begin{equation}\tag{1.34}
x_Ny=xy\mu\qquad(\mu=\lambda_{-1}(N)).
\end{equation}
The map $\lambda_t$ from $K_N$ to $K_N[[t]]$ is defined by
\sgaStableTarget{sga6.e00.eq.1-35}\begin{equation}\tag{1.35}
\lambda_t(n,x)=(1+t)^n\sum_{p\geq 0}\lambda^p(N,x)t^p,
\end{equation}
where $\lambda^0(N,x)=(1,0)$ (the unit element of $K_N$). I say that, with these operations, $K_N$ is a special $\lambda$-ring. Thus one has to verify the formulas
\[
\lambda_t(X+Y)=\lambda_t(X)\lambda_t(Y),\qquad
\lambda_t(XY)=\lambda_t(X)\lambda_t(Y),\qquad
\lambda_t(\lambda^i(X))=\lambda^i(\lambda_t(X))
\]
for $X,Y\in K_N$. One may restrict oneself to verifying these formulas when $X$ and $Y$ lie in $K$, and one checks that it suffices to verify them when $K$ is the free special $\lambda$-ring generated by $X,Y$ and $N$ subject to the relations $\lambda^iN=0$ for $i>q$. But the group homomorphism
\sgaStableTarget{sga6.e00.eq.1-36}\begin{equation}\tag{1.36}
K_N\longrightarrow K,
\end{equation}
which sends unit to unit and on $K$ reduces to the homomorphism $x\mapsto x\mu$, is a homomorphism compatible with the ring structures and the maps $\lambda^i$, as one verifies immediately. In the present case, its restriction to $K$ is injective. Since $K$ is a special $\lambda$-ring, it follows easily that the above relations are indeed verified for $X,Y\in K$, completing the proof of our assertion.

According to formula \hyperlink{sga6.e00.eq.1-23}{(1.23)}, it is natural, for $x\in K$ and $n\geq 1$, to put
\sgaStableTarget{sga6.e00.eq.1-37}\begin{equation}\tag{1.37}
\gamma^n(N,x)=\sum_{i=0}^{n-1}\binom{n-1}{i}\lambda^i(N,x),
\end{equation}
in other words, $\gamma^n(N,x)=\gamma^n(x_N)$, where $x_N$ denotes $x$ considered as an element of the $\lambda$-ring $K_N$.

\sgaStableTarget{sga6.e00.prop.1-5}\textbf{Proposition 1.5.} Let $A$ be a graded commutative ring with unit, let $N$ be an element of $\widetilde A$ of the form
\[
N=[q,1+N^1+\cdots+N^q],
\]
and let $x$ be an arbitrary element of $\widetilde A$:
\[
x=[\nu,1+\sum_{i\geq 1}x^i].
\]
Then, for every integer $j\geq 0$, one has $\gamma^{q+j}(N,x)\in\widetilde A^j$, and the component of degree $j$ of $\gamma^{q+j}(N,x)$ is of the form
\sgaStableTarget{sga6.e00.eq.1-38}\begin{equation}\tag{1.38}
\bigl(\gamma^{q+j}(N,x)\bigr)^{(j)}=(-1)^{j-1}(j-1)!\,G_{q,j}(\nu,x^1,\ldots,x^j;N^1,\ldots,N^j),
\end{equation}
where $G_{q,j}$ is a universal polynomial with integral coefficients characterized by the condition
\sgaStableTarget{sga6.e00.eq.1-39}\begin{equation}\tag{1.39}
G_{q,j}(\nu,x^1,\ldots,N^j)\,(N^q)=(-1)^q(x*\lambda_{-1}(N))^{(q+j)}.
\end{equation}

\emph{Proof.} We may suppose that $A$ is the ring of polynomials with integral coefficients in indeterminates $x^i$ and $N^j$ ($1\leq j\leq q$). Since, combining \hyperlink{sga6.e00.eq.1-32}{(1.32)} and \hyperlink{sga6.e00.eq.1-37}{(1.37)}, one has
\sgaStableTarget{sga6.e00.eq.1-40}\begin{equation}\tag{1.40}
\gamma^n(N,x)\lambda_{-1}(N)=\gamma^n(x\lambda_{-1}(N)),
\end{equation}
and, for $n=q+j$, the second member has filtration $\geq q+j$ (use formula (\hyperlink{sga6.e00.eq.1-22}{1.22} bis), taking account of the fact that $x\lambda_{-1}(N)$ has zero augmentation), while $\lambda_{-1}(N)$ has filtration $q$ (formula \hyperlink{sga6.e00.eq.1-22}{(1.22)}), it follows that $\gamma^n(N,x)$ has filtration $j$, since by \hyperlink{sga6.e00.eq.1-18}{(1.18)} the associated graded of $A$ is integral here. Using the preceding formulas, one immediately obtains the indicated form for $(\gamma^{q+j}(N,x))^{(j)}$. QED.

\sgaStableTarget{sga6.e00.chapter-ii-classes-of-coherent-algebraic-sheaves.occ2}\section*{Chapter II. Classes of coherent algebraic sheaves and Chern classes}

\sgaStableTarget{sga6.e00.sec.1.occ3}\subsection*{1. The theory of Chow}

Throughout this chapter, a base field $k$ is fixed; for simplicity it will be assumed algebraically closed. A large part of what follows, perhaps all of it, is nevertheless valid without this restriction. The algebraic spaces and the regular maps considered will be defined over $k$.

An algebraic space is called \emph{quasi-projective} if it is isomorphic to a locally closed subset of a projective space.

Let $X$ be a nonsingular connected quasi-projective variety of dimension $n$. Denote by $A(X)$, and call the \emph{Chow ring} of $X$, the graded ring of classes of cycles on $X$ modulo rational equivalence. Thus $A^i(X)$ is formed from the classes of cycles of dimension $n-i$. The fact that $A(X)$ is indeed a graded ring follows from the following theorem of Chow.

\sgaStableTarget{sga6.e00.thm.2-1}\textbf{Theorem 2.1.}\footnote{As a result of a misunderstanding on the author's part, the stated theorem is more than is presently known. In fact, one uses only the ``moving lemma'' of Chow, saying that two cycles can be moved in their classes so that their intersection is not excessive; cf. Chevalley Seminar 1958, \emph{Anneau de Chow et applications}.} Every cycle on $X$ is rationally equivalent to a nonsingular cycle. Given elements of $A^i(X)$ and $A^j(X)$, one can find nonsingular representatives of these elements, say $Y^{n-i}$ and $Z^{n-j}$, which meet everywhere transversely.

A cycle is called nonsingular if its components are nonsingular subvarieties. Two cycles are said to meet everywhere transversely if every component of one cuts every component of the other transversely, i.e. if the two secant varieties are nonsingular at their intersection points and their tangent varieties there are in general position.

Let $f$ be a morphism, i.e. a regular map, from $X$ to $Y$ (where $X$ and $Y$ are quasi-projective and nonsingular). Then $f$ defines, using the preceding theorem, a homomorphism of graded rings
\sgaStableTarget{sga6.e00.eq.2-1.occ2}\begin{equation}\tag{2.1}
f^*:A(Y)\longrightarrow A(X),
\end{equation}
namely inverse image of cycle classes. Thus $A(X)$ becomes a contravariant functor in $X$.

Let $f$ be a proper morphism from $X^n$ to $Y^m$ (i.e. in Weil's terminology, the graph of $f$ is complete over $Y$, cf. [3]). Then $f$ also defines an additive homomorphism
\sgaStableTarget{sga6.e00.eq.2-2}\begin{equation}\tag{2.2}
f_*:A(X^n)\longrightarrow A(Y^m),
\end{equation}
preserving the dimension of cycles, hence increasing the degree by $m-n$.




Thus, relative to proper morphisms, and ignoring multiplicative structures and gradings, $A(X)$ is a covariant functor in $X$. Let us record the classical formula
\sgaStableTarget{sga6.e00.eq.2-3.occ2}\begin{equation}\tag{2.3}
 f_*(x\cdot f^*(y))=f_*(x)\cdot y .
\end{equation}

We shall also use Chow's theory of Chern classes. It consists in attaching to every vector bundle $E$ on $X$ (still assumed quasi-projective and nonsingular) Chern classes
\[
 c^i(E)\in A^i(X)\qquad (i\geq 1),
\]
in such a way as to satisfy the following conditions. Put $c^0(E)=1$ and
\sgaStableTarget{sga6.e00.eq.2-4.occ2}\begin{equation}\tag{2.4}
 \Ctil(E)=\left[\rank E,\sum_{i\geq 0}c^i(E)\right]\in \Atil(X)
\end{equation}
(see Chapter I, \S 3 for the definition of $\Atil(X)$). The conditions to be verified by a theory of Chern classes are then expressed as follows:
\sgaStableTarget{sga6.e00.eq.2-5.occ2}\begin{equation}\tag{2.5}
 \Ctil(E)=\Ctil(E')+\Ctil(E'')
\end{equation}
if the vector bundle $E$ is an extension of $E'$ by $E''$. If one puts
\[
 C(E)=\sum_i c^i(E),
\]
this condition is expressed by $C(E)=C(E')C(E'')$. Formula \hyperlink{sga6.e00.eq.2-5}{(2.5)} permits one to extend $\Ctil$ to an additive homomorphism from $K(\xi(X))$ (cf. Chapter I, \S 2, \hyperlink{sga6.e00.ex.2-2}{Example 2)} to $\Atil(X)$:
\sgaStableTarget{sga6.e00.eq.2-6}\begin{equation}\tag{2.6}
 \Ctil:K(\xi(X))\longrightarrow \Atil(X),
\end{equation}
and one wants $\Ctil$ to be a $\lambda$-homomorphism of $\lambda$-rings. This last condition may also be written
\sgaStableTarget{sga6.e00.eq.2-7}\begin{equation}\tag{2.7}
 \Ctil(E\otimes F)=\Ctil(E)\Ctil(F),\qquad
 \Ctil(\lambda^iE)=\lambda^i\Ctil(E),
\end{equation}
for two vector bundles $E$ and $F$; of course $\Ctil$ is automatically compatible with augmentations. Let $D$ be a divisor on $X$ and let $L(D)$ be the vector bundle it defines. One wants
\sgaStableTarget{sga6.e00.eq.2-8}\begin{equation}\tag{2.8}
 c^1(L(D))=\ellc(D),\qquad c^i(L(D))=0\quad (i>1),
\end{equation}
where, for any cycle $Z$, $\ellc(Z)$ denotes its class in $A(X)$. Finally, one imposes that the $c^i$ be functorial, i.e. that if $f$ is a morphism from $X$ to $Y$ and $E$ is a vector bundle on $Y$, then
\sgaStableTarget{sga6.e00.eq.2-9}\begin{equation}\tag{2.9}
 c^i(f^!(E))=f^*(c^i(E)),
\end{equation}
where $f^!(E)$ is the inverse image of $E$ on $X$. One may also write \hyperlink{sga6.e00.eq.2-9}{(2.9)} in the following form: the notion of inverse image of vector bundles defines a $\lambda$-homomorphism of augmented $\lambda$-rings
\sgaStableTarget{sga6.e00.eq.2-10}\begin{equation}\tag{2.10}
 f^!:K(Y)\longrightarrow K(X),\footnote{The notation $f^!$ used here is in conflict with that used in duality theory (cf., for example, R. Hartshorne, \emph{Residues and Duality}, Lecture Notes in Mathematics, no. 20, Springer, 1966). This is why in this seminar we replace it by the notation $f^*$.}
\end{equation}
where here one writes $K(X)$ instead of $K(\xi(X))$, and \hyperlink{sga6.e00.eq.2-9}{(2.9)} means that there is commutativity in the diagram
\[
\begin{tikzcd}[column sep=large,row sep=large]
K(Y) \arrow[r,"f^!"] \arrow[d,"\Ctil_Y"'] & K(X) \arrow[d,"\Ctil_X"]\\
A(Y) \arrow[r,"f^*"'] & A(X).
\end{tikzcd}
\]
Here the space on which $\Ctil$ is considered has been put as a subscript.

\sgaStableTarget{sga6.e00.thm.2-2}\textbf{Theorem 2.2.} There exists a theory of Chern classes satisfying conditions \hyperlink{sga6.e00.eq.2-5.occ2}{(2.5)}, \hyperlink{sga6.e00.eq.2-7}{(2.7)}, \hyperlink{sga6.e00.eq.2-8}{(2.8)}, and \hyperlink{sga6.e00.eq.2-9}{(2.9)}.

We shall also admit the following proposition.

\sgaStableTarget{sga6.e00.prop.2-3}\textbf{Proposition 2.3.} Let $Y$ be a closed subset of $X$ and let $U$ be its complement. If $x$ is an element of $A(X)$ inducing $0$ on $U$, then there exists a representative of $x$ which is a cycle carried by $Y$.

\sgaStableTarget{sga6.e00.cor.unnamed.occ2}\textbf{Corollary.} Let $p=\dim X-\dim Y$. Then, if $i<p$, the homomorphism
\[
 A^i(X)\longrightarrow A^i(U)
\]
is bijective. Every element of the kernel of $A^p(X)\to A^p(U)$ is a linear combination of components of dimension $n-p$ of $Y$; hence it is proportional to $\ellc(Y)$ if $Y$ is irreducible.

\sgaStableTarget{sga6.e00.sec.2.occ3}\subsection*{2. Definition of the Chern classes of coherent algebraic sheaves}

Let $\cC$ be an abelian category, and let $\cC_0$ be a subclass of $\cC$. Denote by $K(\cC_0)$ the quotient group of the free group generated by the isomorphism classes of elements of $\cC_0$ by the subgroup generated by the elements of the form
\[
(E)-(E')-(E''),
\]
where $E$ is an extension of $E'$ by $E''$ and $E,E',E''\in\cC_0$. We must suppose that the classes of elements of $\cC_0$, for the isomorphism relation, form a set. Denote by $\gamma_{\cC_0}(E)$ the class of an object $E\in\cC_0$ in $K(\cC_0)$, and omit the subscript $\cC_0$ when there is no danger of confusion. The group $K(\cC_0)$ is the solution of a universal problem relative to functions
\[
E\longmapsto f(E)
\]
defined on $\cC_0$, with values in an arbitrary abelian group $A$, which are additive in the sense that $f(E)=f(E')+f(E'')$ whenever $E$ is an extension of $E'$ by $E''$.

Let $F$ be an additive functor from $\cC$ to $\cC'$ and let $\cC'_0$ be a subclass of $\cC'$ satisfying the same condition as $\cC_0$, such that for every $A\in\cC_0$ the $R^qF(A)$ are in $\cC'_0$ and are zero for $q$ sufficiently large. We suppose that injective resolutions exist in $\cC$, so that the derived functors $R^qF$ of $F$ exist. Then $F$ defines a homomorphism from $K(\cC_0)$ to $K(\cC'_0)$, still denoted $F$, by the formula
\sgaStableTarget{sga6.e00.eq.2-11}\begin{equation}\tag{2.11}
 F(\gamma_{\cC_0}(E))=\sum_{i\geq 0}(-1)^i\gamma_{\cC'_0}(R^iF(E)).
\end{equation}
More generally, if one has a cohomological functor $(F^i)$ from $\cC$ to $\cC'$ such that for every $E\in\cC_0$ the $F^i(E)$ are in $\cC'_0$ and are zero except for a finite number of $i$, one obtains a homomorphism from $K(\cC_0)$ to $K(\cC'_0)$. These remarks extend to multifunctors; moreover, the classical spectral sequences give transitivity properties which we shall not make explicit.

In particular, if $F$ is an exact functor, formula \hyperlink{sga6.e00.eq.2-11}{(2.11)} simplifies to
\sgaStableTarget{sga6.e00.eq.2-11bis}\begin{equation}\tag{2.11 bis}
 F(\gamma_{\cC_0}(E))=\gamma_{\cC'_0}(F(E)).
\end{equation}
This applies in particular if $\cC=\cC'$, if $F$ is the identity functor, and if $\cC_0\subset \cC'_0$; one has a canonical homomorphism from $K(\cC_0)$ to $K(\cC'_0)$.

\sgaStableTarget{sga6.e00.thm.2-3}\textbf{Theorem 2.3.} Let $\cC$ be an abelian category and let $\cC_0\subset \cC_1$ be two subclasses such that the isomorphism classes of objects of $\cC_1$ form a set. Suppose that the following conditions are satisfied:

\begin{enumerate}[label=(\roman*)]
\item For every exact sequence $0\to A'\to A\to A''\to 0$ in $\cC$ such that $A$ and $A''$ are in $\cC_0$ (respectively $\cC_1$), one has $A'\in\cC_0$ (respectively $A'\in\cC_1$).
\item Let $u:A\to B$ and $v:C\to B$, with $A,B,C\in\cC_1$ and $u$ surjective. Then there exist $L\in\cC_0$ and homomorphisms $u':L\to C$ and $v':L\to A$ such that $u'$ is surjective and $vu'=uv'$.
\item For every $A\in\cC_1$ there exists an integer $d=d(A)$ such that for every left resolution $L$ of $A$ by objects of $\cC_0$ one has $Z_d(L)\in\cC_0$, where $Z_d$ denotes the cycles in dimension $d$.
\end{enumerate}
Under these conditions, the canonical homomorphism from $K(\cC_0)$ to $K(\cC_1)$ is an isomorphism.

\emph{Principle of the proof.} Let $A\in\cC_1$. By (ii) and (iii) there exists a finite resolution $L$ of $A$ by objects of $\cC_0$; put
\[
 \ell(A)=\sum_{i\geq 0}(-1)^i\gamma_{\cC_0}(L_i).
\]
One shows that the element of $K(\cC_0)$ thus constructed is independent of the resolution $L$ chosen. If one has a second resolution $L'$, one ``caps'' $L$ and $L'$ by a third resolution $L''$, with surjective homomorphisms from $L''$ to $L$ and $L'$; this is possible by (i) and (ii). One is then reduced to proving that $E_{\cC_0}(L)=E_{\cC_0}(L'')$. But the difference of the two is $E_{\cC_0}(N)$, where $N$ is the kernel of the surjective homomorphism from $L''$ to $L$; this kernel is in $\cC_0$ by (i). Since $N$ is an acyclic complex in all dimensions, $E_{\cC_0}(N)=0$, whence the conclusion. Here $E_{\cC_0}(L)$, for a finite complex with values in $\cC_0$, denotes the Euler--Poincar\'e characteristic
\[
 \sum_i (-1)^i\gamma_{\cC_0}(L_i).
\]
One proves analogously that the function $\ell:\cC_1\to K(\cC_0)$ thus constructed is additive; hence it defines a homomorphism from $K(\cC_1)$ to $K(\cC_0)$, and it is immediate that this homomorphism is inverse to the canonical homomorphism from $K(\cC_0)$ to $K(\cC_1)$. QED.

\sgaStableTarget{sga6.e00.rem.unnamed.occ3}\textbf{Remark.} To verify (ii), it is enough to verify that $\cC_0$ is contained in a subclass $\cC_2$ of $\cC$ satisfying the two conditions
\begin{enumerate}[label=(ii \alph*)]
\item every $A\in\cC_1$ is isomorphic to a quotient of an $L\in\cC_2$;
\item if $u:A\to B$, with $A,B\in\cC_2$, then $\Ker u\in\cC_1$.
\end{enumerate}

\sgaStableTarget{sga6.e00.cor.unnamed.occ3}\textbf{Corollary.} Let $X$ be a quasi-projective algebraic space. Let $\cF(X)$ be the category of coherent algebraic sheaves on $X$, let $\xi(X)$ be the category of coherent locally free algebraic sheaves, equivalently vector bundles on $X$, and let $\cF_0(X)$ be the category of coherent algebraic sheaves $F$ such that, for every $x\in X$, the stalk $F_x$ is an $\cO_x$-module of finite cohomological dimension. Thus $\xi(X)\subset\cF_0(X)\subset\cF(X)$. Then the canonical homomorphism
\sgaStableTarget{sga6.e00.eq.2-12}\begin{equation}\tag{2.12}
 K(\xi(X))\longrightarrow K(\cF_0(X))
\end{equation}
is an isomorphism.

We verify the conditions of \hyperlink{sga6.e00.thm.2-3}{Theorem 2.3}. Conditions (i) and (iii) follow from the definition of cohomological dimension and from the fact that a finite-type projective module over a noetherian local ring $\cO_x$ is free. To verify (ii), we shall prove that $\cC=\cF(X)$ satisfies conditions (ii a) and (ii b) of the remark. The second is trivial; therefore it remains to prove the following result.

\sgaStableTarget{sga6.e00.prop.2-4}\textbf{Proposition 2.4 (Serre).} Let $X$ be a quasi-projective algebraic space. Then every coherent algebraic sheaf on $X$ is isomorphic to the quotient of a locally free algebraic sheaf.

This proposition is true if $X$ is a closed subset of projective space by [2], since, for $n$ large enough, $F(n)$ is generated by its sections, which means that $F$ is a quotient of $\cO(-n)^k$. In the general case, $X$ is an open subset of a closed subset of a projective space, and it suffices to use the following result, due to Cartier and Serre in particular cases.

\sgaStableTarget{sga6.e00.prop.2-5}\textbf{Proposition 2.5.} Let $F$ be a coherent algebraic sheaf on an open subset $U$ of an algebraic space $X$. Then $F$ is isomorphic to the restriction of a coherent algebraic sheaf on $X$.

By applying Zorn's lemma to the opens containing $U$ on which one can extend $F$, one is reduced to the case where $X$ is affine. Let $\overline F$ be the algebraic sheaf on $X$ whose sections on an open $V$ are the sections of $F$ on $U\cap V$. One checks easily that $\overline F$ is quasi-coherent, i.e. defined on every affine open $V$---hence here on all of $X$---by a module, not necessarily of finite type, over the coordinate ring of that open. This is immediate if $U$ itself is affine, since $\overline F$ is then defined on the open $V$ by $\Gamma(U\cap V,F)$ considered as a module over $\Gamma(V,\cO_X)$. In the general case, $U$ is a finite union of affine opens $U_i$, hence $\overline F$ is the kernel of the evident homomorphism of quasi-coherent sheaves
\[
 \prod_i \overline F\big|_{U_i}\longrightarrow \prod_{i,j}\overline F\big|_{U_i\cap U_j},
\]
and is therefore itself quasi-coherent. Since $X$ is affine, $\overline F$ is generated by its sections. Since $\overline F|_U=F$ is coherent, there exists a finite number of sections of $\overline F$ which generate $F$ on $U$. These sections generate a coherent subsheaf of $\overline F$ whose restriction to $U$ is isomorphic to $F$. QED.

Now suppose that $X$ is nonsingular and quasi-projective. By the theorem on syzygies, one has, with the notation of the corollary to \hyperlink{sga6.e00.thm.2-3}{Theorem 2.3},
\[
 \cF_0(X)=\cF(X).
\]
In this case, therefore,
\sgaStableTarget{sga6.e00.eq.2-13}\begin{equation}\tag{2.13}
 K(\xi(X))=K(\cF(X)).
\end{equation}
One may express this isomorphism as follows.

\sgaStableTarget{sga6.e00.cor.2}\textbf{Corollary 2 to Theorem 2.3.} Let $X$ be a nonsingular quasi-projective variety. For every abelian group $A$, there is a one-to-one correspondence between additive functions $f$ of vector bundles on $X$ with values in $A$, and additive functions $g$ of coherent algebraic sheaves on $X$ with values in $A$. To every $f$ there corresponds the unique $g$ such that
\sgaStableTarget{sga6.e00.eq.2-14}\begin{equation}\tag{2.14}
 f(E)=g(\cO_X(E))
\end{equation}
for every vector bundle $E$ on $X$, where $\cO_X(E)$ denotes the sheaf of germs of regular sections of $E$ on $X$.

To compute $g(F)$ for an arbitrary coherent algebraic sheaf, one considers a finite resolution of $F$ by sheaves of the form $\cO_X(E_i)$, and one then has
\sgaStableTarget{sga6.e00.eq.2-14bis}\begin{equation}\tag{2.14 bis}
 g(F)=\sum_i (-1)^i f(E_i).
\end{equation}

In particular, if $X$ is connected, in \S 1 we considered an additive function $\Ctil(E)$ of the vector bundle $E$, with values in the group $A=\Atil(X)$. We deduce from it an additive function of coherent sheaves, say $\Ctil(F)$, and, by passing to components, functions
\[
 F\longmapsto c^i(F)\in A^i(X).
\]
The $c^i(F)$ are called the \emph{Chern classes of the coherent sheaf} $F$. Theoretically they are computed from the Chern classes of vector bundles by (\hyperlink{sga6.e00.eq.2-14}{2.14} bis), with the understanding that in the second member one must pass to multiplicative notation. As for the component of $\Ctil(F)$ in $\Z$, it is the rank of the sheaf $F$, corresponding to the rank of a module over an integral domain: if $X$ is affine, to which one can always reduce, $F$ is defined by a finite-type module over the coordinate ring of $X$, and one takes the rank of this module.

\sgaStableTarget{sga6.e00.sec.3.occ3}\subsection*{3. Functorial generalities on $K(X)$}

If $X$ is an algebraic space, put, for brevity,
\[
 K(X)=K(\cF(X)).
\]
When $X$ is nonsingular and quasi-projective, identify this group with $K(\xi(X))$ by the \hyperlink{sga6.e00.cor.2}{corollary 2} to \hyperlink{sga6.e00.thm.2-3}{Theorem 2.3}. Since $K(\xi(X))$ is not only an abelian group but a $\lambda$-ring, it follows by transport of structure that $K(X)$ is also a $\lambda$-ring, and is therefore provided with a multiplicative structure and operations $\lambda^i$. It is important to define these structures directly, without passing through vector bundles; this will be done below.

If $E$ is a vector bundle on $X$, denote by $\gamma(E)$ its class in $K(X)$. If $F$ is a coherent algebraic sheaf on $X$, $\gamma(F)$ denotes in the same way its class in $K(X)$. If $Y$ is an irreducible subvariety of $X$, put $\gamma(Y)=\gamma(\cO_Y)$, where $\cO_Y$ denotes the sheaf of local rings of $Y$, considered as a coherent algebraic sheaf on $X$. By linearity, one defines the symbol $\gamma(Y)$ when $Y$ is any cycle on $X$. When there is risk of confusion, one writes $\gamma_X$ for $\gamma$.

Let $X$ be an arbitrary algebraic space. We introduce on $K(X)$ a filtration, denoting by $K(X)^i$ the subgroup of $K(X)$ generated by the $\gamma(F)$ with
\[
 \dim \Supp F\leq n-i,
\]
where $n$ is the dimension of $X$. From Serre's d\'evissage lemma, in the form given in [3], follows the following result.

\sgaStableTarget{sga6.e00.prop.2-6}\textbf{Proposition 2.6.} The group $K(X)^i$ is generated by the $\gamma(Y)$, where $Y$ runs through the subvarieties of dimension $n-i$ of $X$. If $X$ is irreducible, then $\Gr^0(K(X))=\Z$.

The latter isomorphism is given by the homomorphism from $K(X)$ to $\Z$ defined with the aid of the notion of rank of a coherent algebraic sheaf.

Now suppose that $X$ is nonsingular. Following the developments at the beginning of \S 2, define a multiplication in $K(X)$ by putting
\sgaStableTarget{sga6.e00.eq.2-15}\begin{equation}\tag{2.15}
 \gamma(F)\gamma(G)=\sum_{i\geq 0}(-1)^i\gamma(\Tor_i(F,G)),
\end{equation}
the $\Tor_i$ being taken, of course, with respect to the local rings of $X$. The exact sequence of the $\Tor_i$ shows that formula \hyperlink{sga6.e00.eq.2-15}{(2.15)} indeed defines a biadditive map from $K(X)\times K(X)$ to $K(X)$; associativity follows from the spectral sequence of multiple Tors. One also verifies in this way that one can compute the product of several factors by the formula
\sgaStableTarget{sga6.e00.eq.2-16}\begin{equation}\tag{2.16}
 \gamma(F_1)\cdots\gamma(F_n)=\sum_{i\geq 0}(-1)^i\gamma(\Tor_i(F_1,\ldots,F_n)),
\end{equation}
where, in the second member, the $\Tor_i$ are simultaneous Tors, computed by simultaneous projective resolution of the arguments $F_i$.

If in \hyperlink{sga6.e00.eq.2-15}{(2.15)} one supposes that $G$ is locally free, one finds
\sgaStableTarget{sga6.e00.eq.2-15bis}\begin{equation}\tag{2.15 bis}
 \gamma(F)\gamma(G)=\gamma(F\otimes G),
\end{equation}
which shows in particular that $\gamma(\cO_X)$ is the unit element of $K(X)$, and that the canonical homomorphism from $K(\xi(X))$ to $K(X)$ is a homomorphism of rings.

Formula \hyperlink{sga6.e00.eq.2-15}{(2.15)} is evidently suggested by Serre's ``Tor formula'' in intersection theory. This formula proves the second part of the following proposition, the first part following immediately from \hyperlink{sga6.e00.eq.2-15}{(2.15)} and from a local Tor calculation.

\sgaStableTarget{sga6.e00.prop.2-7}\textbf{Proposition 2.7.} Let $X$ be a nonsingular variety, and let $Y$ and $Z$ be two cycles in $X$. If $Y$ and $Z$ meet everywhere transversely, then
\sgaStableTarget{sga6.e00.eq.2-17}\begin{equation}\tag{2.17}
 \gamma(Y)\gamma(Z)=\gamma(Y\cdot Z).
\end{equation}
If $Y$ and $Z$ are homogeneous and of respective dimensions $n-i$ and $n-j$, and if they meet without an excessive component, then
\sgaStableTarget{sga6.e00.eq.2-18}\begin{equation}\tag{2.18}
 \gamma_X(Y)\gamma_X(Z)=\gamma_X(Y\cdot Z)\mod K(X)^{i+j+1}.
\end{equation}
One must be careful: equality does not in general hold in \hyperlink{sga6.e00.eq.2-18}{(2.18)}. We shall see in \S 4 that, if $X$ is nonsingular and quasi-projective, the ring structure of $K(X)$ is compatible with its filtration, and that the associated graded ring $\Gr(K(X))$ is identified with a quotient of $A(X)$ by a subgroup of torsion, perhaps always zero.\footnote{No: for a counterexample see this seminar, \hyperlink{sga6.e14.para.4-7}{XIV 4.7}.}

I do not know how to define directly the $\lambda^i(F)$ in terms of sheaves except when the field $k$ has characteristic $0$; this is the reason that has prevented us from proving the Riemann--Roch theorem when $k$ is not of characteristic $0$. Let $F$ be a coherent algebraic sheaf on $X$ on which the group $\mathfrak S_p$ of permutations of $p$ letters operates. Denote by $F^{\mathrm{alt}}$ the set of $x\in F$ such that
\[
 s\cdot x=\varepsilon(s)x
\]
for every $s\in\mathfrak S_p$, where $\varepsilon(s)$ is the signature of $s$. It is evidently a coherent algebraic subsheaf of $F$. This said, one has:

\sgaStableTarget{sga6.e00.thm.2-8}\textbf{Theorem 2.8.} Let $X$ be a nonsingular quasi-projective variety over a field $k$ of characteristic $0$. Let $F$ be a coherent algebraic sheaf on $X$. Then
\sgaStableTarget{sga6.e00.eq.2-19}\begin{equation}\tag{2.19}
 \lambda^p(\gamma(F))=\sum_{i\geq 0}(-1)^i\gamma\left((\Tor_i(F,F,\ldots,F))^{\mathrm{alt}}\right).
\end{equation}
Here it is clear that the group $\mathfrak S_p$ operates on the sheaf $\Tor_i(F,\ldots,F)$.

\emph{Principle of the proof.} One considers a finite resolution $L$ of $F$ by locally free sheaves $L_i$, and uses the following general lemma.

\sgaStableTarget{sga6.e00.lem.2-9}\textbf{Lemma 2.9.} Let $L=(L_i)$ be a graded locally free coherent algebraic sheaf with only finitely many nonzero components. In $K(X)$ one has the formula
\sgaStableTarget{sga6.e00.eq.2-20}\begin{equation}\tag{2.20}
 E\left((\otimes^p L)^{\mathrm{alt}}\right)=\lambda^p(E(L)).
\end{equation}
In this formula, $\otimes^pL$ is considered as a sheaf on which the group $\mathfrak S_p$ acts, taking the gradings into account; the gradings therefore introduce signs in the usual way, cf. Cartan--Eilenberg. Moreover, if $M$ is a graded coherent algebraic sheaf on $X$, put
\[
 E(M)=\sum_i(-1)^i\gamma(M_i).
\]

To prove formula \hyperlink{sga6.e00.eq.2-20}{(2.20)}, one is reduced to proving the analogous formula in $K(G)$, where $G$ is an algebraic group, for instance a product of groups $\operatorname{Gl}(n,k)$, and where $L$ is a finite-dimensional $G$-module over $k$ (cf. Chapter I, \S 2, \hyperlink{sga6.e00.ex.2-2}{Example 2}, for the definition of $K(G)$). In that case, this formula has already been used in \hyperlink{sga6.e00.thm.1-4}{Theorem 1.4} to show that the second member of \hyperlink{sga6.e00.eq.1-33}{(1.33)} satisfies the same functional equation as the first member; the verification is elementary. QED.

Let $f$ be a morphism from $X$ to $Y$, where $Y$ is an algebraic space. The notion of inverse image of vector bundle defines a homomorphism
\[
 f^!:K(\xi(Y))\longrightarrow K(\xi(X))
\]
of $\lambda$-rings. Identifying vector bundles with locally free coherent algebraic sheaves, the notion of inverse image of vector bundle corresponds to the notion of inverse image of algebraic sheaf, already used in [3, no. 8]. When $Y$ is nonsingular, one defines more generally a homomorphism of groups
\sgaStableTarget{sga6.e00.eq.2-21}\begin{equation}\tag{2.21}
 f^!:K(Y)\longrightarrow K(X)
\end{equation}
by a formula involving an alternating sum of Tors, inspired by the developments at the beginning of \S 2, and which the reader will spell out. One then has commutativity in the following diagram:
\[
\begin{tikzcd}[column sep=large,row sep=large]
K(\xi(Y)) \arrow[r,"f^!"] \arrow[d] & K(\xi(X)) \arrow[d]\\
K(Y) \arrow[r,"f^!"'] & K(X),
\end{tikzcd}
\]
which shows in particular, when $X$ and $Y$ are both nonsingular and quasi-projective, that the two definitions of $f^!$ coincide under the identifications of $K(X)$ with $K(\xi(X))$ and of $K(Y)$ with $K(\xi(Y))$.

Another case where one can define a homomorphism \hyperlink{sga6.e00.eq.2-21}{(2.21)}, without $Y$ necessarily nonsingular, is the case where the functor $F\mapsto f^!(F)$ from $\cF(Y)$ to $\cF(X)$ is exact, i.e. where the local rings of $X$ are flat modules over the local rings of $Y$. It is enough then to put
\[
 f^!(\gamma_Y(F))=\gamma_X(f^!(F)).
\]
This case occurs in particular when $X$ is a locally trivial fibre space over $Y$; to see this, one is reduced to the case of a product $X=Y\times T$. Note that one also verifies that in this case one has $f^!(\cO_Z)=\cO_{f^{-1}(Z)}$ if $Z$ is an irreducible subvariety of $Y$, whence
\sgaStableTarget{sga6.e00.eq.2-22}\begin{equation}\tag{2.22}
 f^!(\gamma_Y(Z))=\gamma_X(f^{-1}(Z))
\end{equation}
if $X$ is flat over $Y$. Without this latter condition, formula \hyperlink{sga6.e00.eq.2-22}{(2.22)} becomes false in general, but one has the following analogue of \hyperlink{sga6.e00.prop.2-7}{Proposition 2.7}.

\sgaStableTarget{sga6.e00.prop.2-8}\textbf{Proposition 2.8.} Let $f:X\to Y$ be a morphism of nonsingular varieties, and let $Z$ be a cycle on $Y$. If $Z$ is everywhere transverse to $f$, formula \hyperlink{sga6.e00.eq.2-22}{(2.22)} is exact. If $Z$ is homogeneous of codimension $i$ and if $f^{-1}(Z)$ has no excessive component, then
\sgaStableTarget{sga6.e00.eq.2-23}\begin{equation}\tag{2.23}
 f^!(\gamma_Y(Z))=\gamma_X(f^{-1}(Z))\mod K(X)^{i+1}.
\end{equation}
The proof is analogous to that of \hyperlink{sga6.e00.prop.2-7}{Proposition 2.7}.

Finally, let $X$ and $Y$ again be arbitrary algebraic spaces, and let $f$ be a proper morphism from $X$ to $Y$. For every coherent algebraic sheaf $F$ on $X$, the direct image sheaf $f_*(F)$ is coherent algebraic, and the same holds more generally for the sheaves $R^if_*(F)$ (cf. [3]). Recall that, by definition, $R^if_*(F)$ is the sheaf on $Y$ associated with the presheaf which sends an open $U$ to the group $H^i(f^{-1}(U),F)$; in particular
\sgaStableTarget{sga6.e00.eq.2-24}\begin{equation}\tag{2.24}
 R^0f_*(F)=f_*(F),\qquad R^qf_*(F)=0\quad(q>\dim X),
\end{equation}
and moreover
\sgaStableTarget{sga6.e00.eq.2-25}\begin{equation}\tag{2.25}
 \Gamma(U,R^qf_*(F))=H^q(f^{-1}(U),F)
\end{equation}
if the open $U$ is affine. According to \S 2, one defines a homomorphism from $K(X)$ to $K(Y)$, still denoted $f_*$, by the formula
\sgaStableTarget{sga6.e00.eq.2-26}\begin{equation}\tag{2.26}
 f_*(\gamma(F))=\sum_{q=0}^{\infty}(-1)^q\gamma(R^qf_*(F)).
\end{equation}
Thus $K(X)$ becomes a covariant functor in $X$ relative to proper morphisms. The relation $(fg)_*=f_*g_*$, just like the relation $(fg)^!=g^!f^!$ which we forgot to mention in its place, is a particular case of the transitivity relations alluded to in \S 2 and which follow from the spectral sequences of composed functors.

Let us note the following formula, valid for $Y$ nonsingular and particularly easy to verify when $Y$ is also quasi-projective by taking for $y$ the class of a vector bundle on $Y$:
\sgaStableTarget{sga6.e00.eq.2-27}\begin{equation}\tag{2.27}
 f_*(x\cdot f^!(y))=f_*(x)\cdot y\qquad (x\in K(X),\ y\in K(Y)).
\end{equation}
In the particular case where the inverse images of the points of $Y$ by $f$ are finite and contained in an affine open,\footnote{This latter condition is evidently a consequence of the first, something the author still did not know in 1957!} one checks that the functor $F\mapsto f_*(F)$ is exact and that $R^qf_*(F)=0$ for $q>0$. In this case, therefore, $f_*(\gamma_X(F))=\gamma_Y(f_*(F))$. In particular, if $X$ is an algebraic subspace of $Y$ and if $i$ is the canonical injection of $X$ into $Y$, then
\sgaStableTarget{sga6.e00.eq.2-28}\begin{equation}\tag{2.28}
 i_*(\gamma_X(F))=\gamma_Y(F),
\end{equation}
where, in the second member, $F$ is considered as a sheaf on $Y$ which is zero outside $X$.

\sgaStableTarget{sga6.e00.sec.4.occ3}\subsection*{4. Some technical results}

Let $X$ be an algebraic space. We have seen in \S 3 that the projection $p$ from $X\times k$ to $X$ defines a homomorphism
\sgaStableTarget{sga6.e00.eq.2-29}\begin{equation}\tag{2.29}
 p^!:K(X)\longrightarrow K(X\times k).
\end{equation}

\sgaStableTarget{sga6.e00.prop.2-9}\textbf{Proposition 2.9.} The homomorphism \hyperlink{sga6.e00.eq.2-29}{(2.29)} is surjective.

We argue by induction on $n=\dim X$. The assertion is trivial if $n<0$; suppose $n\geq 0$. By \hyperlink{sga6.e00.prop.2-6}{Proposition 2.6}, it is enough to verify that for every irreducible closed subset $Y$ of $X\times k$, $\gamma(Y)$ is in the image of $p^!$. It is enough to continue the reasoning by supposing $X$ irreducible. If $p(Y)$ is not dense in $X$, our assertion follows from the induction hypothesis. Otherwise, either $Y=X\times k$ and the assertion is trivial, or $Y$ is a hypersurface in $X\times k$. Let $U$ be an affine open of $X$ formed of nonsingular points. If $D$ is the divisor on $U\times k$ induced by $Y$, one knows by commutative algebra that it is linearly equivalent to a divisor of the form $D'\times k$, where $D'$ is a divisor on $U$. Then
\[
 \gamma_{U\times k}(D)-\gamma_{U\times k}(D'\times k)
\]
belongs to $K(U\times k)^2$ by the corollary to \hyperlink{sga6.e00.prop.2-10}{Proposition 2.10}, and by what has already been proved this element is of the form $p^!(x')$ with $x'\in K(U)$. Thus
\[
 \gamma_{U\times k}(D)=p^!(x'+\gamma_U(D')),
\]
but $x'$ is the restriction of an element $y$ of $K(X)$ by \hyperlink{sga6.e00.prop.2-10}{Proposition 2.10}, and consequently the restriction of $\gamma_{X\times k}(D)-p^!(y)$ to $U\times k$ is zero. By \hyperlink{sga6.e00.prop.2-10}{Proposition 2.10}, $\gamma_{X\times k}(D)-p^!(y)$ comes from an element of
\[
 K(X\times k-U\times k)=K(X'\times k),\qquad X'=X-U.
\]
The proof is finished by applying the induction hypothesis again. QED.

\sgaStableTarget{sga6.e00.cor.unnamed.occ4}\textbf{Corollary.} Let $X$ be an algebraic space and let $p$ be the projection from $X\times k^n$ to $X$. Then the homomorphism
\[
 p^!:K(X)\longrightarrow K(X\times k^n)
\]
is surjective; it is even bijective if $X$ is nonsingular.

The first assertion is by induction on $n$. For the second, put $f(x)=(x,0)$, whence $pf=1$ and $f^!p^!=1$, which proves that $p^!$ is injective.

\sgaStableTarget{sga6.e00.prop.2-10}\textbf{Proposition 2.10.} Let $X$ be an algebraic space, let $U$ be an open subset of $X$, and let $Y$ be its complement in $X$. Then the sequence of natural homomorphisms
\sgaStableTarget{sga6.e00.eq.2-30}\begin{equation}\tag{2.30}
 K(Y)\longrightarrow K(X)\longrightarrow K(U)\longrightarrow 0
\end{equation}
is exact.

N.B. Compare with \hyperlink{sga6.e00.prop.2-3}{Proposition 2.3}.

The fact that $K(X)\to K(U)$ is surjective follows at once from \hyperlink{sga6.e00.prop.2-5}{Proposition 2.5}, or also from \hyperlink{sga6.e00.prop.2-6}{Proposition 2.6}. It remains to prove that if $x\in K(X)$ has zero restriction in $K(U)$, then it is in the image of $K(Y)$. For this, return to the explicit definition of $K(X)$ and $K(U)$. There are two coherent algebraic sheaves $F$ and $G$ on $X$ with
\[
 x=\gamma_X(F)-\gamma_X(G)
\]
and
\[
 \gamma_U(F)-\gamma_U(G)=0.
\]
This latter relation means that one can find coherent sheaves $P'$ and $Q'$ on $U$, provided with filtrations whose factors are pairwise isomorphic on $U$, and such that one has an isomorphism on $U$
\[
 P=F+Q'\longrightarrow Q=G+P'.
\]
By \hyperlink{sga6.e00.prop.2-5}{Proposition 2.5}, the sheaves $P'$ and $Q'$ are restrictions to $U$ of coherent sheaves on $X$, denoted by the same letters; in the same way, their filtrations come from filtrations on the corresponding sheaves on $X$, by \hyperlink{sga6.e00.lem.2-11}{Lemma 2.11} below. One then has
\[
 \gamma_X(F)-\gamma_X(G)=(\gamma_X(P)-\gamma_X(Q))+\bigl(\gamma_X(P')-\gamma_X(Q')\bigr).
\]
Introducing the factors $P_i$ and $Q_i$ of the filtrations of $P'$ and $Q'$ on $X$, the last term is the sum of terms $\gamma_X(P_i)-\gamma_X(Q_i)$, where $P_i|_U\simeq Q_i|_U$. We are thus reduced to proving the following particular case: if $R$ and $S$ are coherent sheaves on $X$ such that $R|_U\simeq S|_U$, then $\gamma_X(R)-\gamma_X(S)$ is in the image of $K(Y)$. Let $T$ be the graph of the given isomorphism from $R|_U$ to $S|_U$; it is a coherent subsheaf of $(R\times S)|_U$, hence by \hyperlink{sga6.e00.lem.2-11}{Lemma 2.11} it is the restriction to $U$ of a coherent subsheaf, again denoted $T$, of $R\times S$. Then
\[
 \gamma_X(R)-\gamma_X(S)=(\gamma_X(R)-\gamma_X(T))-(\gamma_X(S)-\gamma_X(T)),
\]
and it is enough to prove, for instance, that the first term of the second member is in the image of $K(Y)$. But the projection $u:T\to R$ is an isomorphism over $U$, hence
\[
 \gamma_X(R)-\gamma_X(T)=\gamma_X(\Coker u)-\gamma_X(\Ker u),
\]
and $\Ker u$ and $\Coker u$ have support in $Y$.

It remains, then, to prove the following lemma.

\sgaStableTarget{sga6.e00.lem.2-11}\textbf{Lemma 2.11.} Let $X$ be an algebraic space, let $F$ be a coherent algebraic sheaf on $X$, let $U$ be an open subset of $X$, and let $G$ be a coherent subsheaf of $F|_U$. Then $G$ is the restriction to $U$ of a coherent subsheaf of $F$.

Let $\overline G$ be the subsheaf of $F$ whose sections on an open $V$ of $X$ are the sections of $F$ on $V$ whose restriction to $U\cap V$ is a section of $G$. Everything amounts to proving that this sheaf $\overline G$ is coherent. For this one may suppose $X$ and $U$ affine, because $U$ is a union of affine opens $U_i$ and $\overline G$ is then the intersection of the analogous sheaves $\overline{G|U_i}$. Then the coordinate ring of $U$ is a ring of fractions $A_S$ of the coordinate ring $A$ of $X$, and one immediately verifies that if $F$ is defined by the $A$-module $M$, then $F|_U$ is defined by the $A_S$-module
\[
 M_S=A_S\otimes_A M.
\]
The subsheaf $G$ of $F|_U$ is then defined by a submodule $N'$ of $M_S$, and one sees that $\overline G$ is precisely defined by the submodule $N$ of $M$, inverse image of $N'$ under the canonical homomorphism from $M$ to $M_S$. QED.

\sgaStableTarget{sga6.e00.cor.1}\textbf{Corollary 1 to Proposition 2.10.} Let $x\in K(X)$. Then $x\in K(X)^i$ if and only if there exists a closed subset $Y$ of $X$, of dimension $\leq n-i$, such that the restriction of $x$ to $X-Y$ is zero.

\sgaStableTarget{sga6.e00.cor.2.occ2}\textbf{Corollary 2.} Suppose $X$ is connected of dimension $n$, with no singularities in dimension $n-1$. If $D$ and $D'$ are two linearly equivalent divisors on $X$, then
\[
 \gamma(D)=\gamma(D')\mod K(X)^2.
\]
By \hyperlink{sga6.e00.cor.1}{Corollary 1}, one may suppose $X$ nonsingular. The corollary then follows from the general formula
\sgaStableTarget{sga6.e00.eq.2-31}\begin{equation}\tag{2.31}
 \gamma(D)=1-\gamma(L(-D))\mod K(X)^2,
\end{equation}
where $L(-D)$ denotes the rank-one vector bundle defined by the divisor $-D$. To verify \hyperlink{sga6.e00.eq.2-31}{(2.31)}, one writes it in the form
\sgaStableTarget{sga6.e00.eq.2-31bis}\begin{equation}\tag{2.31 bis}
 \gamma(L(-D))=1-\gamma(D)\mod K(X)^2,
\end{equation}
and notes that the two members are homomorphisms from the group of divisors to the multiplicative group of invertible elements of the ring $K(X)/K(X)^2$. Here $K(X)^1/K(X)^2$ has square zero, which is very easily verified on the explicit formula \hyperlink{sga6.e00.eq.2-15}{(2.15)}, using \hyperlink{sga6.e00.cor.1}{Corollary 1}. Thus one is reduced to the case where $D$ is an irreducible hypersurface. More generally, if $D$ is a hypersurface whose components are disjoint, one has the equality
\sgaStableTarget{sga6.e00.eq.2-32}\begin{equation}\tag{2.32}
 \gamma(D)=1-\gamma(L(-D)),
\end{equation}
as follows from the exact sequence of sheaves
\[
0\longrightarrow \cO_X(L(-D))\longrightarrow \cO_X\longrightarrow \cO_D\longrightarrow 0.
\]

\sgaStableTarget{sga6.e00.thm.2-12}\textbf{Theorem 2.12.} Let $X$ be a nonsingular quasi-projective variety, and let $Z$ and $Z'$ be two homogeneous cycles of codimension $p$ on $X$ which are linearly equivalent. Then
\sgaStableTarget{sga6.e00.eq.2-33}\begin{equation}\tag{2.33}
 \gamma_X(Z)=\gamma_X(Z')\mod K(X)^{p+1}.
\end{equation}
By hypothesis, one can find a cycle $Y$ of dimension $n-p+1$ on $X\times k$ ($n$ being the dimension of $X$), and two points $a,a'\in k$, such that
\[
 Z=Y\cdot (X\times a),\qquad Z'=Y\cdot (X\times a').
\]
By formula \hyperlink{sga6.e00.eq.2-23}{(2.23)}, one then has
\[
 \gamma_X(Z)=i_a^!(\gamma_{X\times k}(Y))\mod K(X)^{p+1},
\]
and an analogous formula for $\gamma_X(Z')$. It is therefore enough to prove that $i_a^!=i_{a'}^!$, which follows at once from \hyperlink{sga6.e00.prop.2-9}{Proposition 2.9}. QED.

\sgaStableTarget{sga6.e00.cor.1.occ2}\textbf{Corollary 1.} The filtration of $K(X)$ is compatible with its ring structure, i.e. one has
\[
 K(X)^iK(X)^j\subset K(X)^{i+j}
\]
for all integers $i,j\geq 0$. Moreover there is a natural surjective homomorphism of graded rings
\sgaStableTarget{sga6.e00.eq.2-34}\begin{equation}\tag{2.34}
 \varphi:A(X)\longrightarrow \Gr(K(X)),
\end{equation}
where the second member denotes the graded ring associated with $K(X)$, provided with the filtration $(K(X)^i)_{i\geq 0}$, defined by the condition
\[
 \varphi(\ellc(Z^{n-p}))=\gamma(Z^{n-p})\mod K(X)^{p+1}.
\]
This corollary follows immediately from \hyperlink{sga6.e00.thm.2-12}{Theorem 2.12}, Propositions \hyperlink{sga6.e00.prop.2-6}{2.6} and \hyperlink{sga6.e00.prop.2-7}{2.7}, and Chow's \hyperlink{sga6.e00.thm.2-1}{theorem 2.1}, by induction on the integer $i+j$.

\sgaStableTarget{sga6.e00.cor.2.occ3}\textbf{Corollary 2.} Let $X$ and $Y$ be two nonsingular quasi-projective varieties and let $f$ be a morphism from $X$ to $Y$. Then the homomorphism $f^!$ from $K(Y)$ to $K(X)$ is compatible with the filtrations, and one has a commutative diagram
\[
\begin{tikzcd}[column sep=large,row sep=large]
A(Y) \arrow[r,"\varphi_Y"] \arrow[d,"f^*"'] & \Gr(K(Y)) \arrow[d,"\Gr(f^!)"]\\
A(X) \arrow[r,"\varphi_X"'] & \Gr(K(X)).
\end{tikzcd}
\]
Using the graph of $f$, decompose $f$ as the product of an immersion of $X$ into $X\times Y$ and of a projection of $X\times Y$ onto $Y$. For the latter the corollary is trivial, and for an immersion one proceeds as for \hyperlink{sga6.e00.cor.1}{Corollary 1}, using \hyperlink{sga6.e00.prop.2-8}{Proposition 2.8} in place of \hyperlink{sga6.e00.prop.2-7}{Proposition 2.7}.

\sgaStableTarget{sga6.e00.prop.2-13}\textbf{Proposition 2.13 (``cellular decomposition'').} Let $P$ be an algebraic space, irreducible or not, provided with an increasing family
\[
P^0\subset P^1\subset\cdots\subset P^n=P
\]
of closed subsets, such that for every $i$ with $0\leq i\leq n$, the connected components of $P^i-P^{i-1}$ are algebraic spaces $V_{i,\alpha}$ isomorphic to $k^i$; put $P^{-1}=\varnothing$. Then:
\begin{enumerate}[label=\alph*)]
\item For every algebraic space $X$, the natural map from $K(X)\otimes K(P)$ to $K(X\times P)$, deduced from the tensor product of sheaves on $X$ and $P$, is surjective.
\item The $\gamma_P(V_{i,\alpha})$ form a system of generators of the group $K(P)$.
\end{enumerate}
The proof is by induction on $n$, using the corollary to \hyperlink{sga6.e00.prop.2-9}{Proposition 2.9} and \hyperlink{sga6.e00.prop.2-10}{Proposition 2.10}. N.B. Unfortunately, I do not know whether the $\gamma_P(V_{i,\alpha})$ form a basis of $K(P)$ or whether the homomorphism considered in a) is bijective. We shall not need to know this below.

\sgaStableTarget{sga6.e00.thm.2-14}\textbf{Theorem 2.14.} Let $x_i$ ($1\leq i\leq q$) and $E_j$ ($1\leq j\leq r$) be letters, and let $n_j$ ($1\leq j\leq r$) be positive integers. Consider the symbols $\lambda^k(x_i)$ with $k\geq 1$ and $\lambda^k(E_j)$ with $1\leq k\leq n_j$ as indeterminates, and let $R$ be a polynomial with integral coefficients in these indeterminates. Suppose that, for every graded ring $A$, when one substitutes in $R$ elements $\xi_i$ for the $x_i$ and elements
\[
 \eta_j=[n_j,1+\eta_j^{(1)}+\cdots+\eta_j^{(n_j)}]
\]
of $\Atil$ for the $E_j$, subject to the restrictions $\varepsilon(\xi_i)=0$ and $\eta_j^{(k)}=0$ for $k>n_j$, one has
\[
 R\bigl((\lambda^k(\xi_i)),(\lambda^k(\eta_j))\bigr)\in \Atil^p.
\]
Then, if $X$ is an irreducible nonsingular quasi-projective variety, for arbitrary elements $x'_i$ of $K(X)^1$ and vector bundles $E'_j$ of dimension $n_j$ on $X$ ($1\leq i\leq q$, $1\leq j\leq r$), one has
\sgaStableTarget{sga6.e00.eq.unnamed}\begin{equation}\tag{*}
 R\bigl((\lambda^k(x'_i)),(\lambda^k(\gamma(E'_j)))\bigr)\in K(X)^p.
\end{equation}

\emph{Proof.} By \S 2, the $x'_i$ are themselves differences of elements of the form $\gamma(E)$, with $E$ a vector bundle. On the other hand, assuming $X$ locally closed in a projective space $P$, it follows from \S 2 and \hyperlink{sga6.e00.prop.2-5}{Proposition 2.5} that, for every $F\in\cF(X)$, the sheaf $F(n)$ is generated by its sections for $n$ sufficiently large. If $F$ is defined by a vector bundle $E$, this proves that $E(n)$ is the inverse image of the canonical vector bundle on a Grassmannian variety $G$ by a morphism from $X$ to $G$. Considering the corresponding morphism from $X$ to $P\times G$, one sees that $F$ is the inverse image of a vector bundle on $P\times G$. Applying this to all vector bundles which occur, and using \hyperlink{sga6.e00.cor.2}{Corollary 2} to \hyperlink{sga6.e00.thm.2-12}{Theorem 2.12}, we are reduced to the case where $X$ is a product of Grassmannians. In this case the conclusion will follow from the following lemma.

\sgaStableTarget{sga6.e00.lem.2-15}\textbf{Lemma 2.15.} Suppose that $X$ is a product of Grassmannians. Then one can find a $\lambda$-homomorphism from $K(X)$ to a torsion-free $\lambda$-ring of the form $\Atil$, where $A$ is a graded ring, compatible with filtrations, and such that the corresponding homomorphism $\Gr(K(X))\to\Gr(A)$ is injective.

To prove the lemma, we shall use Chow's theory of Chern classes. If one could prove the lemma, or \hyperlink{sga6.e00.thm.2-14}{Theorem 2.14} directly, one would deduce a purely ``sheaf-theoretic'' theory of Chern classes; cf. \S 5.

We shall consider the homomorphism $\Ctil$ from $K(X)$ to $\Atil(X)$ and put $A=A(X)$. It remains to prove that the homomorphism $\psi=\Gr(\Ctil)$ from $\Gr(K(X))$ to $\Gr(A(X))\simeq A(X)$ is injective. The fact that $\Ctil$ is compatible with filtrations is true for every nonsingular quasi-projective variety and follows easily from the corollary to \hyperlink{sga6.e00.prop.2-3}{Proposition 2.3}. On the other hand, since $X$ is a product of Grassmannians $G_j$, there are vector bundles $E_j$ on $X$, inverse images of the canonical bundles on the factors. We shall use the following known facts:
\begin{enumerate}[label=\alph*)]
\item On $X$, linear equivalence of cycles is equivalent to numerical equivalence, and $A(X)$ is a group of finite rank.
\item The $c^i(E_j)$ generate the ring $A(X)$.
\end{enumerate}
From a) it follows that
\[
 \varphi:A(X)\longrightarrow \Gr(K(X))
\]
is injective, and hence bijective, by the following lemma.

\sgaStableTarget{sga6.e00.lem.2-15.occ2}\textbf{Lemma 2.15.} Let $X^n$ be a nonsingular projective variety and let $Z^{n-p}$ be a cycle on $X$ such that $\varphi(Z)\in\Gr^p(K(X))$ is zero. Then $Z^{n-p}$ is numerically equivalent to zero.

To prove this lemma, it suffices to use the additive function
\[
 \chi(F)=\sum_{i\geq 0}(-1)^i\dim H^i(X,F)
\]
and the fact that this function is equal to $1$ for $F=\cO_{(x)}$ for every $x\in X$.

Thus, in the case under consideration, $A(X)$ and $\Gr(K(X))$ have the same finite rank in every degree, and $\Gr(K(X))$ is torsion-free, since this is true of $A(X)$. From assertion b) above it follows, in a formal way, that the map $\Ctil$ from $K(X)$ to $\Atil(X)$ is surjective modulo finite groups. Since $\Gr(K(X))$ and $A(X)$ have the same finite rank in every degree, it then follows easily that the homomorphism $\Gr(\Ctil)$ from $\Gr(K(X))$ to $\Gr(\Atil(X))$ is bijective modulo finite groups in every degree; one uses induction on the degree. Since $\Gr(K(X))$ is torsion-free, this homomorphism $\Gr(\Ctil)$ is therefore injective. QED.

N.B. The homomorphism $\Ctil$ is not surjective in general, even if $X$ is a projective space of dimension $2$.

\sgaStableTarget{sga6.e00.sec.5.occ2}\subsection*{5. Sheaf-theoretic definition of Chern classes. Application to the study of injection morphisms}

\sgaStableTarget{sga6.e00.thm.2-16}\textbf{Theorem 2.16.} Let $X$ be a nonsingular quasi-projective variety. Then, for every $x\in K(X)$ and every integer $i\geq 0$, the element $\gamma^i(x-\varepsilon(x))$ is of filtration $\geq i$. Let $c^i(x)\in\Gr^i(K(X))$ be the class of this element modulo $K(X)^{i+1}$. If one puts
\[
 \widetilde c(x)=\left[\varepsilon(x),1+\sum_{i\geq 0}c^i(x)\right]\in \widetilde{\Gr(K(X))},
\]
then $\widetilde c$ satisfies conditions analogous to those of \hyperlink{sga6.e00.thm.2-2}{Theorem 2.2}.

The first assertion follows from \hyperlink{sga6.e00.thm.2-14}{Theorem 2.14}, as does the fact that $x\mapsto\widetilde c(x)$ is a homomorphism of $\lambda$-rings; formula \hyperlink{sga6.e00.eq.1-25}{(1.25)} implies $\gamma_t(x+y)=\gamma_t(x)\gamma_t(y)$, which already shows that $\widetilde c$ is an additive homomorphism. The functoriality of $\widetilde c$ is verified immediately, and \hyperlink{sga6.e00.prop.2-8}{(2.8)} follows from \hyperlink{sga6.e00.eq.2-31}{(2.31)}. QED.

\sgaStableTarget{sga6.e00.cor.1.occ3}\textbf{Corollary 1.} One has
\[
 c^i(x)=\varphi(c^i(x)).
\]
This follows from a uniqueness theorem for Chern classes, proved as in Hirzebruch by passage to flag varieties. One only has to show that, if $Y$ is a flag variety associated with a vector bundle on $X$, then $\Gr(K(X))\to\Gr(K(Y))$ is injective; this is verified by the explicit determination of $K(Y)$, with all its structures, from $K(X)$. I shall not give the details here.

\sgaStableTarget{sga6.e00.cor.2.occ4}\textbf{Corollary 2.} Suppose that for every element $Z\in A(X)$ there exists a morphism $f$ from $X$ to a product $X'$ of Grassmannians such that $Z$ is in the image of $f^*$. Then the canonical homomorphisms
\[
 \varphi:A(X)\longrightarrow \Gr(K(X)),\qquad
 \psi:\Gr(K(X))\longrightarrow A(X)
\]
satisfy, in dimension $i$, the relations
\sgaStableTarget{sga6.e00.eq.2-35}\begin{equation}\tag{2.35}
 \psi\varphi=(-1)^{i-1}(i-1)!\,1,
 \qquad
 \varphi\psi=(-1)^{i-1}(i-1)!\,1.
\end{equation}
Since $\varphi$ is surjective, it is clearly enough to prove the first formula, and for this one may suppose that $X$ is a product of Grassmannians. Since then $\varphi$ is injective, it is enough to prove the second formula \hyperlink{sga6.e00.eq.2-35}{(2.35)}, which by \hyperlink{sga6.e00.cor.1}{Corollary 1} is written
\[
 \gamma^i(x-\varepsilon(x))=(-1)^{i-1}(i-1)!\,x\mod K(X)^{i+1}
\]
for every $x\in K(X)^i$. Since $\psi$ is injective, it is enough to prove the analogous relation for $\Ctil(x)\in A(X)^i$, and this is then a consequence of (\hyperlink{sga6.e00.eq.1-22}{1.22} bis).

\sgaStableTarget{sga6.e00.rem.unnamed.occ4}\textbf{Remark.} \hyperlink{sga6.e00.cor.2}{Corollary 2} makes plausible the validity of \hyperlink{sga6.e00.eq.2-35}{(2.35)} in all cases; there should exist a direct proof of it. In every case one proves easily the second formula \hyperlink{sga6.e00.eq.2-35}{(2.35)} modulo the torsion subgroup of $\Gr(K(X))$, and in characteristic $0$ this same formula is a consequence of the Riemann--Roch theorem. The interest of \hyperlink{sga6.e00.eq.2-35}{(2.35)}, when true, lies in the fact that it shows that $\varphi$ is injective modulo torsion groups, perhaps always zero. Thus one does not lose much by defining Chern classes as elements of $\Gr(K(X))$. If one tensors with $\Q$, as will be done later for the Riemann--Roch theorem, the two definitions of Chern classes would then be identical.

Let $X$ still be nonsingular and quasi-projective. Let in addition $Y$ be a nonsingular subvariety of $X$, and let $i$ be the injection of $Y$ into $X$. For $i_*$ and $i^!$ one then has the following formulas:
\begin{enumerate}[label=(\roman*)]
\item $i^!$ is a homomorphism of $\lambda$-rings.
\item $i_*(x\cdot i^!(y))=i_*(x)\cdot y$, whence in particular
\[
 i_*i^!(y)=\xi y,
\qquad \text{with }\xi=i_*i^!(1)=\gamma_X(N).
\]
\item $i^!i_*(x)=x\cdot\lambda_{-1}(N)$, whence in particular
\[
 i^!(\xi)=i^!i_*(1)=\lambda_{-1}(N).
\]
\item $i_*(x)i_*(x')=i_*(xx'\lambda_{-1}(N))$.
\end{enumerate}
Here, by definition, $N$ is the class in $K(Y)$ of the conormal bundle of $Y$ in $X$. Formulas (i) and (ii) are true for arbitrary morphisms, not necessarily injections, and (iii) and (iv) are verified from the \hyperlink{sga6.e00.eq.2-15}{definition (2.15)} and the analogous formula for inverse images, taking $x$ and $x'$ to be classes of vector bundles and using the fact that
\sgaStableTarget{sga6.e00.eq.2-36}\begin{equation}\tag{2.36}
 \Tor(\cO_Y,\cO_Y)=\Lambda(N).
\end{equation}
Formula (iv) expresses that $i_*$ is a homomorphism of rings if one introduces on $K(Y)$ a new multiplicative structure, by the definition
\[
 x\cdot_N x'=xx'\lambda_{-1}(N)
\]
(cf. Chapter I, \S 4). It is completed by the formula
\begin{enumerate}[label=(\roman*)]
\setcounter{enumi}{4}
\item $i_*(\lambda^p(N,x))=\lambda^p(i_*(x))$ for $p\geq 1$,
\end{enumerate}
which means that $i_*$ defines a $\lambda$-homomorphism from $K(Y)_N$ into $K(X)$. This formula (v), certainly true in all cases, is for the moment proved only in characteristic $0$, because it follows from the conjunction of Theorems \hyperlink{sga6.e00.thm.1-4}{1.4} and 2.8.

There are analogous formulas for $i_*$ and $i^*$:
\begin{enumerate}[label=(\roman* bis)]
\item $i^*$ is a homomorphism of graded rings.
\item $i_*(x\cdot i^*(y))=i_*(x)y$, whence in particular
\[
 i_*i^*(y)=\xi' y,
\qquad \text{with }\xi'=i_*(1).
\]
\item $i^*i_*(x)=x\cdot (-1)^qN^q$, whence in particular
\[
 i^*(\xi')=(-1)^qN^q.
\]
\item $i_*(x)i_*(x')=i_*(xx'(-1)^qN^q)$.
\item $i_*$ increases degrees by $q$ units.
\end{enumerate}
In these formulas, $q$ is the codimension of $Y$ in $X$, and one puts $N^i=c^i(N)$. Only formulas (iii bis) and (iv bis) require proof. I do not know whether they are well known, or even whether they are true, and I confine myself to noting that they follow by reduction from formulas (ii) and (iv), provided one works in $\Gr(K(Y))$ and $\Gr(K(X))$ and not in $A(Y)$ and $A(X)$. One then uses the fact that $i_*$ and $i^*$ come by passage to the quotient from $i_*$ and $i^!$, noting that $i_*$ shifts the filtration by $q$ units.

All universal formulas relating $i_*$ and $i^!$, respectively $i_*$ and $i^*$, are formal consequences of the preceding formulas, as one can convince oneself without great difficulty.

Formula (v) gives
\[
 i_*(\gamma^p(N,x))=\gamma^p(i_*(x)).
\]
The conjunction of \hyperlink{sga6.e00.thm.2-14}{Theorem 2.14} and \hyperlink{sga6.e00.prop.1-5}{Proposition 1.5} shows that $\gamma^p(N,x)$, for $p=q+j$, is of filtration $\geq j$. If one denotes by $G_{q,j}(N,x)$ its class in $\Gr^j(K(Y))$, \hyperlink{sga6.e00.thm.2-16}{Theorem 2.16} shows that
\[
 i_*(G_{q,j}(N,x))=c^p(i_*(x)).
\]
Finally, \hyperlink{sga6.e00.prop.1-5}{Proposition 1.5} joined with \hyperlink{sga6.e00.thm.2-14}{Theorem 2.14} permits one to express $G_{q,j}(N,x)$ as a polynomial with integral coefficients in the $c^i(N)$, the $c^i(x)$, and $\varepsilon(x)$. One deduces the following theorem.

\sgaStableTarget{sga6.e00.thm.2-17}\textbf{Theorem 2.17.} Suppose the field $k$ has characteristic $0$, and let $i:Y\to X$ be an injection morphism of nonsingular quasi-projective varieties. Then one has formulas (i) to (v), and moreover, for every integer $j\geq 0$, the formula
\sgaStableTarget{sga6.e00.eq.2-37}\begin{equation}\tag{2.37}
 c^{q+j}(i_*(x))=
 i_*\left(\left(\frac{\widetilde c(x)*\lambda_{-1}(\widetilde c(N))}{(-1)^qN^q}\right)^{(j)}\right).
\end{equation}
The second member of formula \hyperlink{sga6.e00.eq.2-37}{(2.37)} is an abbreviated way of denoting the polynomial in the components of $\widetilde c(x)$ and $\widetilde c(N)$ described in \hyperlink{sga6.e00.prop.1-5}{Proposition 1.5}.

\sgaStableTarget{sga6.e00.cor.unnamed.occ5}\textbf{Corollary.} Under the preceding conditions one has
\sgaStableTarget{sga6.e00.eq.2-38}\begin{equation}\tag{2.38}
 \ch(\widetilde c(i_*(x)))=i_*\left(\ch(\widetilde c(x))\,\mathfrak C(\check N)^{-1}\right),
\end{equation}
where $\ch$ is the Chern homomorphism defined by formula \hyperlink{sga6.e00.eq.1-26}{(1.26)}.

It is enough to use \hyperlink{sga6.e00.eq.2-37}{(2.37)} written in the form
\sgaStableTarget{sga6.e00.eq.2-37bis}\begin{equation}\tag{2.37 bis}
 \widetilde c(i_*(x))=1+i_*\left(\frac{\widetilde c(x)*\lambda_{-1}(\widetilde c(N))-1}{(-1)^qN^q}\right),
\end{equation}
the explicit formula \hyperlink{sga6.e00.eq.1-29}{(1.29)}, formula (iv bis), and finally \hyperlink{sga6.e00.eq.1-30}{(1.30)}.

One may write \hyperlink{sga6.e00.eq.2-38}{(2.38)} in the form
\sgaStableTarget{sga6.e00.eq.2-38bis}\begin{equation}\tag{2.38 bis}
 \ch_X(F)=i_*\left(\ch_Y(F)\,\mathfrak C(T_{X/Y})^{-1}\right),
\end{equation}
where $F$ is a coherent algebraic sheaf on $Y$ (considered in the first member as a sheaf on $X$ which is zero outside $Y$), where $\ch_X$ and $\ch_Y$ are the Chern homomorphisms taken respectively on $X$ and $Y$, and where $T_{X/Y}$ is the normal bundle of $Y$ in $X$.

\sgaStableTarget{sga6.e00.rem.unnamed.occ5}\textbf{Remarks.} Contrary to \hyperlink{sga6.e00.eq.2-38}{(2.38)}, formula \hyperlink{sga6.e00.eq.2-37}{(2.37)} is a ``formula without denominators''. It has been proved, even in characteristic $0$, only by placing oneself in $\Gr(K(X))$ and not in $A(X)$, which may not be the same thing. It is certainly true in every characteristic and in $A(X)$. Let us note that one verifies at once that the images of the two members by $i^*$ are the same; hence, applying $i_*$, the product of their difference by $N^q$ is zero. This suggests a completely different method of proof, by trying to obtain $Y$ as a transverse inverse image of a subvariety $Y'$ in a variety $X'$ such that, in dimension $\leq n=\dim X$, the class of $Y'$ in $A(X')$ is not a divisor of $0$. Unfortunately one cannot hope for $X'$ to be a Grassmannian or something close to it.

Note that, if torsion elements are neglected, formula \hyperlink{sga6.e00.eq.2-38}{(2.38)} is equivalent to \hyperlink{sga6.e00.eq.2-37}{(2.37)}.




% Additional macro definitions
\providecommand{\cS}{\mathcal S}
\providecommand{\cD}{\mathcal D}
\providecommand{\cE}{\mathscr E}
\providecommand{\cG}{\mathscr G}
\providecommand{\cL}{\mathscr L}
\providecommand{\cM}{\mathscr M}
\providecommand{\cP}{\mathscr P}
\providecommand{\cK}{\mathcal K}
\providecommand{\Db}{\mathrm D}
\providecommand{\Kb}{\mathrm K}
\providecommand{\Trc}{\operatorname{Tr}}
\providecommand{\ob}{\operatorname{ob}}
\providecommand{\Coh}{\operatorname{Coh}}
\providecommand{\Cone}{\operatorname{Cone}}
\providecommand{\im}{\operatorname{Im}}
\providecommand{\Id}{\operatorname{Id}}
\providecommand{\RR}{\mathrm{R.R.}}
\providecommand{\RRR}{\mathrm{R.R.R.}}
\providecommand{\GrR}{\mathrm{G}}
\providecommand{\pr}{\operatorname{pr}}


\sgaStableTarget{sga6.e00.sec.6.occ2}\section*{6. The Riemann--Roch theorem}

Its statement is as follows: let \(f\) be a proper morphism from a non-singular quasi-projective variety \(Y\) into a variety \(X\) of the same kind. Then one has
\sgaStableTarget{sga6.e00.eq.2-39}\begin{equation}\tag{2.39}
 f_*\bigl(\ch_Y(x)\,\mathfrak C(Y)\bigr)
 =
 \ch_X\bigl(f_!(x)\bigr)\,\mathfrak C(X).
\end{equation}
Both sides are in \(A(X)\otimes \mathbb Q\), and \(\mathfrak C(X)=\mathfrak C(T_X)\).

This formula makes it possible to compute, up to torsion, the Chern classes of \(f_!(x)\) from the Chern classes of \(x\), those of \(T_X\) and \(T_Y\), and the homomorphism \(f_*\). Or, equivalently, it makes it possible to determine the homomorphism
\[
  f_!:K(Y)\otimes \mathbb Q \longrightarrow K(X)\otimes \mathbb Q
\]
when these groups are identified, by means of the homomorphisms \(\ch_Y\) and \(\ch_X\), respectively with \(A(Y)\otimes \mathbb Q\) and \(A(X)\otimes \mathbb Q\) (provided that formula \hyperlink{sga6.e00.eq.2-35}{(2.35)} is true), from \(f_*\) and the Chern classes of \(T_X\) and \(T_Y\).

I can prove formula \hyperlink{sga6.e00.eq.2-39}{(2.39)} only in characteristic \(0\), and as an equality in \(\GrR(K(X))\otimes\mathbb Q\). For this, since \(Y\) is locally closed in a projective space \(P\), \(f\) is decomposed as the product of an injection of \(Y\) into \(P\times X\) and a projection morphism from \(P\times X\) to \(X\), and \hyperlink{sga6.e00.eq.2-39}{(2.39)} is proved for each of these two morphisms.

For an injection morphism, one sees immediately that formulas \hyperlink{sga6.e00.eq.2-38}{(2.38)} and \hyperlink{sga6.e00.eq.2-39}{(2.39)} are equivalent. The characteristic \(0\) hypothesis enters only through \hyperlink{sga6.e00.thm.2-8}{Theorem 2.8}, which permits one to express the classes \(\lambda^i(\gamma(F))\), for a coherent algebraic sheaf \(F\), by local cohomological invariants of \(F\).

For a projection morphism \(f:P\times X\to X\), \hyperlink{sga6.e00.prop.2-13}{Proposition 2.13} permits us to suppose that \(x\) is of the form \(y\otimes z\), with \(y\in K(P)\) and \(z\in K(X)\). Then the Künneth formula gives
\[
  f_!(x)=\chi(y)\,z,
\]
where, for a coherent algebraic sheaf \(F\) on \(P\), one puts
\[
  \chi(F)=\sum_{i\geq 0}(-1)^i\dim H^i(P,F).
\]
One is therefore at once reduced to proving the formula
\sgaStableTarget{sga6.e00.eq.2-40}\begin{equation}\tag{2.40}
  \chi(F)=\pi_r\bigl(\ch_p(F)\,\mathfrak C(P)\bigr),
\end{equation}
where \(\pi_r\) denotes the component of degree \(r=\dim P\). By \hyperlink{sga6.e00.prop.2-13}{Proposition 2.13}, one may suppose that \(F=\mathcal O_Q\), where \(Q\) is a linear subvariety of \(P\). Then \(\chi(F)=1\), since this is the arithmetic genus of \(Q\); it remains only to see that the second member of \hyperlink{sga6.e00.eq.2-40}{(2.40)} is \(1\).

Indeed, if \(V\) is a hyperplane of \(P\) and if one puts
\[
  \xi=\gamma_P(\mathcal O_V),\qquad L=\ell(V),
\]
then
\[
  1-\xi=\gamma_P(E(-L)),
\]
where \(E(-L)\) is the vector bundle defined by \(-L\). Hence
\[
  C_p(1-\xi)=[1,1-L],
\]
and consequently
\[
  \ch_p(1-\xi)=e^{-L},\qquad \ch_p(\xi)=1-e^{-L},
\]
and therefore
\[
  \ch_p(\xi^p)=(1-e^{-L})^p.
\]
If \(p\) is the dimension of \(Q\) in \(P\), one has \(\xi^p=\gamma_P(Q)\); hence
\[
  \ch_p(F)\,\mathfrak C(P)
  =
  \bigl(L/(1-e^{-L})\bigr)^{r+1}(1-e^{-L})^p.
\]
It remains to verify that the coefficient of \(L^r\) in this formal series in \(L\) is \(1\). This follows from the characteristic property of the formal series \(L/(1-e^{-L})\), namely that the coefficient of \(L^q\) in \(\bigl(L/(1-e^{-L})\bigr)^{q+1}\) is \(1\) (put \(q=r+1-p\)).

\sgaStableTarget{sga6.e00.bibliographical-references}\subsection*{Bibliographical references}
\begin{enumerate}[label={[\arabic*]}]
\item Hirzebruch: \emph{Neue topologische Methoden in der algebraischen Geometrie}.
\item Serre: Faisceaux algébriques cohérents, \emph{Ann. of Math.}
\item Grothendieck: Sur les faisceaux algébriques cohérents, in Séminaire Cartan 1956/57.
\end{enumerate}

\sgaStableTarget{sga6.e00.proof-of-the-riemann-roch-theorem}\section*{Proof of the Riemann--Roch theorem}
\begin{center}
(Grothendieck, 1 November 1957)
\end{center}

\sgaStableTarget{sga6.e00.reminder-of-notation}\subsection*{Reminder of notation}
These are the notations of the ``Report on Riemann--Roch,'' cited as \(\RRR\) below.

All varieties considered are non-singular and quasi-projective, that is, isomorphic to a locally closed subvariety of a projective space. The base field is algebraically closed, of arbitrary characteristic.

If \(X\) is such a variety, \(K(X)\) denotes the abelian group generated by the classes of coherent sheaves on \(X\), modulo the identification coming from an extension into a sum; the alternating sum of Tor makes \(K(X)\) a ring. If \(F\) is a sheaf, respectively \(E\) a vector bundle, respectively \(Y\) a subvariety of \(X\), one denotes by \(\gamma_X(F)\), respectively \(\gamma_X(E)\), respectively \(\gamma_X(Y)\), the element of \(K(X)\) corresponding to \(F\), respectively to \(E\), respectively to \(\mathcal O_Y\). In fact, one often writes simply \(F,E\) instead of \(\gamma_X(F)\).

Let \(A(X)\) be the graded ring of classes of cycles on \(X\) for rational equivalence (cf. Chow); \(A^i(X)\) is formed by the cycles of codimension \(i\). If \(E\) is a vector bundle, \(C(E)=\sum C^i(E)\), with \(C^i(E)\in A^i(X)\), is its Chern class (defined as inverse image of Schubert cycles). By linearity one deduces the classes \(C^i(y)\), for all \(y\in K(X)\). If
\[
  C(y)=\prod (1+a_i),
\]
one sets, following Hirzebruch,
\[
  \ch(y)=\sum e^{a_i},\qquad
  T(y)=\prod\frac{a_i}{1-e^{-a_i}}.
\]
These are elements of \(A(X)\otimes\mathbb Q\). When \(y\) is the tangent bundle of \(X\), one writes \(T(X)\) instead of \(T(y)\).

If \(G\) is a vector bundle, \(\lambda^pG\) denotes its \(p\)-th exterior power and
\[
  \lambda_{-1}(G)=\sum(-1)^p\lambda^p(G);
\]
this is an element of \(K(X)\).

If \(f\) is a proper regular map from a variety \(Y\) into a variety \(X\), one defines homomorphisms
\[
  f^*:A(X)\longrightarrow A(Y),\qquad
  f_*:A(Y)\longrightarrow A(X).
\]
The first is homogeneous of degree \(0\) and multiplicative. The second is homogeneous of degree \(\dim X-\dim Y\). One also defines homomorphisms
\[
  f^*:K(X)\longrightarrow K(Y),\qquad
  f_!:K(Y)\longrightarrow K(X).
\]
If \(F\) is a coherent sheaf on \(X\), \(f^*(F)\) is, by definition, the alternating sum of the \(\Tor_p^{\mathcal O_X}(\mathcal O_Y,F)\). For vector bundles, this is simply the inverse-image operation.

If \(G\) is a coherent sheaf on \(Y\), \(f_*(G)\) is coherent on \(X\). If \(R^pf_*\) denotes the \(p\)-th derived functor of the functor \(f_*\) so defined, one has, by definition,
\[
  f_!(G)=\sum(-1)^p R^pf_*(G).
\]

\sgaStableTarget{sga6.e00.riemann-roch-theorem}\subsection*{Riemann--Roch theorem}
Let \(f:Y\to X\) be a proper morphism (regular map) from a non-singular quasi-projective variety \(Y\) into another such variety \(X\), and let \(y\in K(Y)\). Then
\sgaStableTarget{sga6.e00.eq.1}\begin{equation}\tag{1}
  f_*\bigl(\ch(y)T(Y)\bigr)=\ch(f_!(y))T(X)
  \quad\text{in } A(X)\otimes\mathbb Q.
\end{equation}

\emph{Proof} (abridged).

\sgaStableTarget{sga6.e00.lem.1}\textbf{Lemma 1.} Consider proper morphisms of varieties
\[
  Z \xrightarrow{g} Y \xrightarrow{f} X.
\]
If the theorem \(\RR\) is true for \(f\) and for \(g\), it is true for \(fg\). If it is true for \(fg\) and for \(g\), it is true for \(f\). If \(f_*\) is injective, and if the theorem \(\RR\) is true for \(fg\) and for \(f\), it is true for \(g\).

This is trivial. In fact, when one supposes \(\RR\) true for \(fg\) and \(g\), one must still suppose that \(g_!\) is surjective in order to be able to assert that \(\RR\) is true for \(f\).

Since a proper morphism is composed of an injection and a projection \(X\times P\to X\), where \(P\) is a projective space, one is reduced to treating these two cases separately. The second, as was seen above, is reduced to the standard calculation, by means of \hyperlink{sga6.e00.prop.2-13}{Proposition 2.13} (``cellular decomposition'') of \(\RRR\). Thus suppose \(Y\subset X\), with injection morphism \(i\). Formula (1) becomes
\sgaStableTarget{sga6.e00.eq.2}\begin{equation}\tag{2}
  \ch(i_!(y))=i_*\bigl(\ch(y)T(E)^{-1}\bigr),
\end{equation}
where \(E=T_{X/Y}\) is the normal bundle of \(Y\) in \(X\).

\sgaStableTarget{sga6.e00.lem.2}\textbf{Lemma 2.} Suppose that there exists on \(X\) a vector bundle \(G\) of rank \(p=\codim_XY\), such that
\begin{equation}\tag{\sgaStableTargetInline{sga6.e00.eq.3}3}
  \gamma_X(Y)=i_*(1)=\lambda_{-1}(G),\qquad
  i^*(G)=E,\qquad
  (-1)^pC^p(G)=Y.
\end{equation}
If then
\sgaStableTarget{sga6.e00.eq.4}\begin{equation}\tag{4}
  y=i^*(x),\qquad x\in K(X),
\end{equation}
formula (2) is valid.

The proof is immediate from the fact that \(\ch\) is a homomorphism of rings, and that
\sgaStableTarget{sga6.e00.eq.5}\begin{equation}\tag{5}
  \ch(\lambda_{-1}(G))=(-1)^pC^p(G)\,T(G)^{-1}.
\end{equation}

\sgaStableTarget{sga6.e00.cor.unnamed.occ6}\textbf{Corollary.} \(\RR\) is true if \(Y\) is a hypersurface and if \(y\) is of the form \(i^*(x)\), with \(x\in K(X)\).

\sgaStableTarget{sga6.e00.lem.3}\textbf{Lemma 3.} Suppose that \(X=Y\times P\), where \(P\) is a projective space, and that \(i\) is the injection \(u\mapsto (u,v_0)\), \(v_0\in P\). Then \(i_*\) is injective and \(\RR\) is true.

This is pure standard calculation.

Suppose again that \(X,Y,y\) are arbitrary, always with \(Y\subset X\). Let \(X'\) be the variety obtained by blowing up \(X\) along \(Y\), and let \(Y'\) be the hypersurface inverse image of \(Y\). One has a diagram of maps
\[
\begin{tikzcd}
Y' \arrow[r,"j"] \arrow[d,"g"'] & X' \arrow[d,"f"]\\
Y \arrow[r,"i"'] & X .
\end{tikzcd}
\]
The inverse image of \(E\) in \(Y'\) admits a subbundle of rank \(1\), whose dual is denoted \(L\). Moreover \(L^{-1}\) is also the restriction to \(Y'\) of the line bundle \(L(Y')\) defined on \(X'\) by the divisor \(Y'\). Put \(F=E/L^{-1}\). Elementary computations give the following lemma.

\sgaStableTarget{sga6.e00.lem.4}\textbf{Lemma 4.} One has the formulas, where \(p=\codim_XY\):
\sgaStableTarget{sga6.e00.eq.6}\sgaStableTarget{sga6.e00.eq.7}\sgaStableTarget{sga6.e00.eq.8}\sgaStableTarget{sga6.e00.eq.9}\begin{align}
  f^*i_!(y)&=j_*\bigl(\lambda_{-1}(\check F)y\bigr), \tag{6}\\
  \lambda_{-1}(\check F)&\equiv \gamma^{p-1}(E)\lambda^p(\check E)
    \pmod{1-L}, \tag{7}\\
  f_*(1)&=1,\quad\text{hence }f_*f^*=1, \tag{8}\\
  f^*i_*(\alpha)&\equiv j_*\bigl(\alpha C^{p-1}(F)\bigr)
    \pmod{\Ker f_*}\quad\text{for all }\alpha\in A(Y). \tag{9}
\end{align}
To simplify notation, \(K(Y)\) has been made to operate on \(K(Y')\), and \(A(Y)\) on \(A(Y')\), by the homomorphisms \(g^*\). In formula (7), it seems that \(\gamma^{p-1}(E)\) denotes \(\lambda^{p-1}(E-2)\), that is,
\[
  \sum_{i+j=p-1}(\,\cdots\,)\lambda^i(E);
\]
the congruence takes place in \(K(Y')\).

Formula (6) follows from the equality
\sgaStableTarget{sga6.e00.eq.6bis}\begin{equation}\tag{6 bis}
  \Tor_i^{\mathcal O_X}\bigl(\mathcal O_Y(v),\mathcal O_Y(v')\bigr)
  =
  \lambda^i\check F(v')
  \quad (v'\in Y',\ v=g(v')).
\end{equation}
Formula (7) follows from the formulas
\sgaStableTarget{sga6.e00.eq.10}\begin{equation}\tag{10}
  g_*(L^{-i})=0\quad\text{for }0<i\leq p-1,\qquad g_*(1)=1,
\end{equation}
themselves consequences of Künneth and of the cohomology of projective spaces.

Formula (8) is trivial. Finally, (9) follows from \(g_*(C^{p-1}(F))=1\), which is deduced, for example, from the immediate formulas
\sgaStableTarget{sga6.e00.eq.10bis}\begin{equation}\tag{10 bis}
  g_*(\xi^i)=0\quad\text{if }0\leq i<p-1,\qquad
  g_*(\xi^{p-1})=1,
\end{equation}
with \(\xi=C^1(L)\in A^1(Y')\). In fact, (9) is very likely an identity, and not merely a congruence. The unnecessarily complicated formula (7) may be replaced, for example, by \hyperlink{sga6.e07.lem.4-1}{VII 4.1}.

From formula (7) it follows at once that \(\lambda_{-1}(\check F)\in j^*(K(X'))\), provided that
\sgaStableTarget{sga6.e00.eq.11}\begin{equation}\tag{11}
  y\,\gamma^{p-1}(E)\lambda^p(\check E)\in i^*(K(X)).
\end{equation}
Suppose this condition (11) satisfied. By the corollary to \hyperlink{sga6.e00.lem.2}{Lemma 2}, \(\RR\) is true for \((Y',X',y\lambda_{-1}(\check F))\). On the other hand, to prove (2), formula (8) reduces the matter to proving that
\sgaStableTarget{sga6.e00.eq.12}\begin{equation}\tag{12}
  f^*\ch(i_!(y))
  \equiv
  f^*\bigl(i_*(\ch(y)T(E)^{-1})\bigr)
  \pmod{\Ker f_*}.
\end{equation}
This last formula is verified immediately, using (6), (9), \(\RR\) for \((Y',X',y\lambda_{-1}(\check F))\), and the formula
\sgaStableTarget{sga6.e00.eq.13}\begin{equation}\tag{13}
  \ch(\lambda_{-1}(\check F))=C^{p-1}(F)T(F)^{-1}
\end{equation}
(cf. formula (5)).

Thus \(\RR\) is proved every time \(y\,\gamma^{p-1}(E)=0\), in particular every time \(p-1>\dim Y\), since \(\gamma^{p-1}(E)\) is of filtration \(\geq p-1\) in \(K(Y)\), filtered by the codimension of supports; that is, every time
\[
  \dim Y < \frac{\dim X-1}{2}.
\]
One is reduced to this case by combining \hyperlink{sga6.e00.lem.1}{Lemma 1} and \hyperlink{sga6.e00.lem.3}{Lemma 3}, taking for \(P\) a projective space of sufficiently large dimension, for example \(r>\dim X+1\). This proves the assertion.

\sgaStableTarget{sga6.e00.rem.unnamed.occ6}\textbf{Remark.} We have used the fact that \(\gamma^{p-1}(E)\) was of filtration \(\geq p-1\). One can prove this fact directly, without passing through Grassmannians as in \(\RRR\), 2.16, by studying the projective bundle associated to \(E\). This method more generally gives \hyperlink{sga6.e00.thm.2-14}{Theorem 2.14} of \(\RRR\).

\sgaStableTarget{sga6.e00.extract-from-a-letter-of-grothendieck}\subsection*{Extract from a letter of Grothendieck}
``Morally,'' the new method rests on the determination of \(K(X)\) and \(A(X)\) when \(X\) is either the projective bundle associated to a vector bundle, or a variety obtained by blowing up a non-singular subvariety. In this second case I have not made the complete determination; a small equality is still missing. With the notation of the short note, one must prove formula (9), but as an equality:
\sgaStableTarget{sga6.e00.eq.9.occ2}\begin{equation}\tag{9}
  f^*(i_*(\alpha))=j_*(\alpha C^{p-1}(F)),\qquad \alpha\in A(Y).
\end{equation}
One observes that the two members have the same image by \(f_*\), and the same image by \(j^*\), so that the question is equivalent to
\sgaStableTarget{sga6.e00.eq.9bis}\begin{equation}\tag{9 bis}
  \Ker f_*\cap\Ker j^*=0.
\end{equation}
In any case, (9) is true after reducing to \(G(K(X'))\) (which destroys a little torsion, at most), for it then follows from (6); there is therefore hardly any doubt that it is true.

I also point out that equality (9) follows very easily from (6), without passing through the troublesome (6 bis).

\begin{flushleft}
P.C.C.J-P.S.
\end{flushleft}

\newpage
\sgaStableTarget{sga6.e01.expose}\section*{Exposé I. Generalities on finiteness conditions in derived categories}
\begin{center}
by L. Illusie
\end{center}

\sgaStableTarget{sga6.e01.sec.0}\subsection*{0. Introduction}
The object of \hyperlink{sga6.e01.expose}{Exposés I} to IV is to develop, with the degree of generality appropriate in this Seminar, the formalism of finiteness conditions in derived categories of ringed topoi.

As was said in \hyperlink{sga6.e00.expose}{Exposé 0}, the need to define, on arbitrary schemes, ``Grothendieck groups'' possessing good variance properties forced one to generalize and to soften the finiteness notions used up to now.

A classical notion such as that of a coherent sheaf on a ringed space \((X,\mathcal O_X)\) makes sense only when \(\mathcal O_X\) is coherent. One might think of replacing it by the notion of finite presentation. But this has the inconvenience that the kernel of an epimorphism of modules of finite presentation is no longer, in general, of finite presentation. A satisfactory notion is reached by observing that, when \(\mathcal O_X\) is coherent, a coherent sheaf \(F\) is not only of finite presentation but of \(n\)-finite presentation for every \(n\in\mathbb N\), where \(n\)-finite presentation means that locally there exists an exact sequence
\[
  L_n\longrightarrow L_{n-1}\longrightarrow \cdots \longrightarrow L_0\longrightarrow F\longrightarrow 0,
\]
with the \(L_i\) free of finite type. When \(\mathcal O_X\) is no longer assumed coherent, say that an \(\mathcal O_X\)-module \(F\) is \emph{pseudo-coherent} if it satisfies the preceding condition, that is, if it is of \(n\)-finite presentation for every \(n\). Then the notion of pseudo-coherence has, with respect to short exact sequences, the same stability property as coherence: if two terms of a short exact sequence of \(\mathcal O_X\)-modules are pseudo-coherent, then so is the third.

The preceding notion is easily generalized to complexes of sheaves. A complex \(F\) of \(\mathcal O_X\)-modules is said to be \emph{pseudo-coherent} if it is \(n\)-pseudo-coherent for every \(n\in\mathbb Z\), where \(n\)-pseudo-coherent means that there locally exists a quasi-isomorphism \(L\to F\), with \(L\) a complex bounded above and with components free of finite type in degrees \(>n\). For a complex \(F\) reduced to degree \(0\), this is the same as saying that \(F\) is pseudo-coherent as a complex or as a module.

The notion of pseudo-coherent complex has excellent stability properties, described in \hyperlink{sga6.e01.expose}{Exposé I}, §§1 and 2. First of all it is stable under isomorphism in the derived category \(\Db(X)\) of the category of \(\mathcal O_X\)-modules; better still, if two objects of a distinguished triangle of \(\Db(X)\) are pseudo-coherent, then so is the third. In other words, the full subcategory \(\Db(X)_{\mathrm{coh}}\) of \(\Db(X)\) formed by pseudo-coherent complexes is a triangulated subcategory. Moreover, pseudo-coherence is preserved under inverse image and derived tensor product. In the case where \(\mathcal O_X\) is coherent, to say that a complex \(F\) is pseudo-coherent simply means that \(F\) has locally bounded-above cohomology and that all the \(H^i(F)\) are coherent.

If \(\mathcal O_X\) is a sheaf of regular local rings, every coherent \(\mathcal O_X\)-module \(F\) admits locally a finite resolution to the left by finite-type free modules, that is, there exists locally an exact sequence
\[
  0\longrightarrow L_n\longrightarrow \cdots \longrightarrow L_0\longrightarrow F\longrightarrow 0,
\]
where the \(L_i\) are free of finite type. It is therefore natural, when \(\mathcal O_X\) is an arbitrary sheaf of local rings, to study the modules which have the preceding property. Such modules are called \emph{perfect}. More generally, a complex \(F\) of \(\mathcal O_X\)-modules is said to be \emph{perfect} if there exists locally a quasi-isomorphism \(L\to F\), where \(L\) is a bounded complex with components free of finite type. It is shown in \hyperlink{sga6.e01.expose}{Exposé I}, §3, that the full subcategory \(\Db(X)_{\mathrm{parf}}\) of \(\Db(X)\) formed by perfect complexes is, like \(\Db(X)_{\mathrm{coh}}\), a triangulated subcategory, stable under inverse image and derived tensor product. One obviously has
\[
  \Db(X)_{\mathrm{parf}}\subset \Db(X)_{\mathrm{coh}},
\]
but the inclusion is strict in general. Pseudo-coherence can be redefined ``by passage to the limit'' from perfection: a complex \(F\) is pseudo-coherent if and only if, for every \(n\in\mathbb Z\), \(F\) can be locally approximated to order \(n\) by a perfect complex, meaning that there exists locally a distinguished triangle
\[
\begin{tikzcd}
L \arrow[rr] && F \arrow[dl]\\
& R \arrow[ul] &
\end{tikzcd}
\]
with \(L\) perfect and \(R\) acyclic in degrees \(>n\). Conversely, perfection is recovered from pseudo-coherence by means of an additional regularity condition: a complex is perfect if and only if it is pseudo-coherent and locally of finite Tor-dimension (\hyperlink{sga6.e01.expose}{Exposé I}, §5). Thus one recovers the fact that, if the local rings are regular, every coherent sheaf is perfect, and more generally every pseudo-coherent complex with locally bounded cohomology is perfect.

Pseudo-coherence and perfection are the two fundamental finiteness notions with which we shall work in this Seminar. Sections 1 to 5 of \hyperlink{sga6.e01.expose}{Exposé I} are devoted to their definition and to establishing their elementary stability properties. For this purpose, only two or three local properties of the category of finite-type free modules, relative to the category of all modules, are used. It is therefore convenient and useful, notably in \hyperlink{sga6.e02.expose}{Exposé II}, to axiomatize the situation by introducing a notion of pseudo-coherence (respectively of perfection) in a category fibered over a site, relative to a fibered subcategory satisfying suitable conditions.

Sections 6 to 8 of \hyperlink{sga6.e01.expose}{Exposé I} generalize to perfect complexes a certain number of notions well known for finite-type locally free sheaves: rank, duality, and the trace of an endomorphism.

In \hyperlink{sga6.e02.expose}{Exposé II}, one examines the problem of the existence of ``global resolutions'': under what conditions, on a ringed topos \((X,\mathcal O_X)\), is a pseudo-coherent (respectively perfect) complex globally isomorphic in \(\Db(X)\) to a complex of locally free modules of finite type? The importance of this question lies essentially in the fact that at present one does not know how to generalize to perfect complexes certain usual constructions on locally free modules, such as exterior and symmetric powers or Chern classes. That is why it is useful to have workable sufficient criteria for the existence of global resolutions, which reduce certain questions about pseudo-coherent (respectively perfect) complexes to analogous questions about locally free sheaves of finite type. Such criteria are given in \hyperlink{sga6.e02.expose}{Exposé II}.

In particular, one shows that global resolutions exist in the following cases: \(X\) is the Zariski topos of an affine scheme, or more generally of a quasi-compact scheme possessing an ample invertible sheaf, or even an ample family of invertible sheaves in the sense of Kleiman, for example a regular scheme; or again \(X\) is the topos of sheaves of sets on a compact topological space and \(\mathcal O_X\) is the sheaf of continuous functions with complex values. By way of illustration, and to show the flexibility of the notions introduced, the appendix gives a ``purely sheaf-theoretic'' definition of the index of a family of elliptic operators.

\hyperlink{sga6.e03.expose}{Exposé III} studies the stability of finiteness conditions under derived direct image. To obtain reasonable statements, one must relativize the notions of pseudo-coherence and perfection. We place ourselves here in the framework of ordinary schemes, which is sufficient for the Seminar, but there is no doubt that one will sooner or later need to develop an analogous theory for relative analytic schemes or spaces. Let \(p:X\to S\) be an \(S\)-scheme locally of finite type, and let \(F\) be a complex of \(\mathcal O_X\)-modules. One says that \(F\) is pseudo-coherent (respectively perfect) relative to \(p\) if one can locally extend \(X\), by a closed immersion, into a smooth \(S\)-scheme \(X'\), so that \(F\), extended by zero to \(X'\), is pseudo-coherent (respectively perfect). The notion of pseudo-coherence relative to \(p\) is especially interesting when \(S\) is not locally noetherian; in the opposite case it coincides with the ordinary notion of pseudo-coherence. Moreover, as expected, to say that \(F\) is perfect relative to \(p\) is equivalent to saying that \(F\) is pseudo-coherent relative to \(p\) and locally of finite Tor-dimension relative to \(p\) (i.e. relative to the sheaf of rings \(p^{-1}(\mathcal O_S)\)).

The central theorem of \hyperlink{sga6.e03.expose}{Exposé III} is the following finiteness theorem, stated below in a slightly more precise form: if \(f:X\to Y\) is a proper morphism of \(S\)-schemes locally of finite type, the functor \(Rf_*\) transforms complexes pseudo-coherent relative to \(S\) into complexes pseudo-coherent relative to \(S\). This theorem is in fact a conjecture, but it is nevertheless proved in the two particular cases \(S\) locally noetherian and \(f\) projective. It seems likely, unfortunately, that the extension to the general case is of the same order of difficulty as the analogous theorem in analytic geometry, Grauert's theorem. Combining the finiteness theorem with an essentially trivial formula called the projection formula, one obtains workable criteria for the stability of relative perfection under direct image. As a corollary one recovers Grauert's continuity and semicontinuity theorems (EGA III 7.6).

\hyperlink{sga6.e04.expose}{Exposé IV} translates the results of \hyperlink{sga6.e01.expose}{Exposés I} to III into the language of ``Grothendieck groups.'' On a ringed topos \((X,\mathcal O_X)\), write \(K^*(X)\) (respectively \(K_\bullet(X)\)) for the Grothendieck group of the category of perfect complexes and complexes of finite Tor-dimension (respectively pseudo-coherent complexes with bounded homology). The group \(K^*(X)\) is a ring, and \(K_\bullet(X)\) is a module over \(K^*(X)\). The functor \(K^*\) is contravariant, as suggested above, while \(K_\bullet\) is covariant for proper pseudo-coherent morphisms of schemes, or of analytic spaces, subject to the availability of the finiteness theorem. There are also unusual variants: covariance of \(K^*\), contravariance of \(K_\bullet\), for morphisms satisfying suitable regularity hypotheses. Finally, the globalization criteria of \hyperlink{sga6.e02.expose}{Exposé II} make it possible to link these groups with the ``Grothendieck groups'' defined naïvely from coherent sheaves or finite-type locally free sheaves.

\sgaStableTarget{sga6.e01.sec.1}\section*{1. Preliminary definitions}

\sgaStableTarget{sga6.e01.sec.1.1}\subsection*{1.1. Fibered categories with additive fibers (respectively abelian, respectively triangulated fibers)}

\sgaStableTarget{sga6.e01.para.1-1-1}\textbf{1.1.1.} Let \(S\) be a category, and let \(\mathcal C\) be an \(S\)-category fibered in categories with additive (respectively triangulated) fibers (SGA 1 VI 6.1). We shall say that \(\mathcal C\) is an \emph{additive} (respectively \emph{triangulated}) \(S\)-category, or that \(\mathcal C\) is additive (respectively triangulated) over \(S\), if, for every arrow \(f:X\to Y\) of \(S\), the inverse-image functor
\[
  f^*:\mathcal C_Y\longrightarrow \mathcal C_X
\]
(determined up to unique isomorphism) is additive (respectively exact). We shall say that \(\mathcal C\) is an \emph{abelian} \(S\)-category if \(\mathcal C\) is an additive \(S\)-category with abelian fibers, and that \(\mathcal C\) is a \emph{flat abelian} \(S\)-category if, moreover, the inverse-image functors are exact.

\sgaStableTarget{sga6.e01.para.1-1-2}\textbf{1.1.2.} Let \(\mathcal C\) be a triangulated \(S\)-category. One defines, for example with the aid of a cleavage of \(\mathcal C\) (SGA 1 VI 7.1), a fibered \(S\)-category \(\Trc(\mathcal C)\) whose fiber over each \(X\in \ob(S)\) is the category \(\Trc(\mathcal C_X)\) of distinguished triangles of \(\mathcal C_X\).

\sgaStableTarget{sga6.e01.para.1-1-3}\textbf{1.1.3.} Let \(\mathcal C\) and \(\mathcal C'\) be additive (respectively triangulated) \(S\)-categories, and let \(F:\mathcal C\to\mathcal C'\) be a functor. We say that \(F\) is an additive (respectively exact) \(S\)-functor if \(F\) is a Cartesian \(S\)-functor that is additive (respectively exact) on each fiber.



% Additional macro definitions
\providecommand{\cC}{\mathcal C}
\providecommand{\cD}{\mathcal D}
\providecommand{\cK}{\mathcal K}
\providecommand{\cO}{\mathcal O}
\providecommand{\cS}{\mathcal S}
\providecommand{\Cpx}{\mathrm C}
\providecommand{\Kpx}{\mathrm K}
\providecommand{\Dpx}{\mathrm D}
\providecommand{\D}{\mathrm D}
\providecommand{\Ho}{\mathrm H}
\providecommand{\RHom}{\mathrm R\!\operatorname{Hom}}
\providecommand{\Ob}{\operatorname{ob}}
\providecommand{\Coker}{\operatorname{Coker}}
\providecommand{\Ker}{\operatorname{Ker}}
\providecommand{\Im}{\operatorname{Im}}
\providecommand{\Hom}{\operatorname{Hom}}
\providecommand{\Cone}{\operatorname{Cone}}
\providecommand{\id}{\operatorname{Id}}
\providecommand{\qis}{\xrightarrow{\sim}}


\sgaStableTarget{sga6.e01.para.1-1-3.occ2}\textbf{1.1.3 (continued).} If \(\mathcal C\) and \(\mathcal C'\) are triangulated \(S\)-categories and if \(F\) is an exact \(S\)-functor, then \(F\) defines a Cartesian \(S\)-functor
\[
  \Trc(\mathcal C)\longrightarrow \Trc(\mathcal C').
\]

\sgaStableTarget{sga6.e01.para.1-1-4}\textbf{1.1.4.} Let \(\mathcal C\) be a triangulated \(S\)-category, and let \(\mathcal C'\) be a full fibered sub-\(S\)-category of \(\mathcal C\). We say that \(\mathcal C'\) is a \emph{triangulated sub-\(S\)-category} of \(\mathcal C\) if, for every \(X\in\Ob S\), \(\mathcal C'_X\) is a triangulated subcategory of \(\mathcal C_X\) \([V]\), Chap. I, §1, 2.3.

% SGA-SRC-001475 — reconstructed by mathematical inference; scan inconclusive
% The malformed documentary abbreviation footnote is not reproduced in this English text.
\sgaStableTarget{sga6.e01.para.1-1-5}\textbf{1.1.5.} Let \(\mathcal C\) be an additive \(S\)-category. One defines (SGA 4 XVII 2.2, 2.4) an additive (respectively triangulated) \(S\)-category
\[
  \Cpx(\mathcal C)\qquad\text{respectively}\qquad \Kpx(\mathcal C),
\]
whose fiber over \(X\in\Ob S\) is the category \(\Cpx(\mathcal C_X)\) of complexes with differential of degree \(+1\) (respectively the category \(\Kpx(\mathcal C_X)\) of complexes of \(\mathcal C_X\) ``up to homotopy''). If \(\mathcal C\) is a flat abelian \(S\)-category, one defines, loc. cit., a triangulated \(S\)-category \(\Dpx(\mathcal C)\), whose fiber over \(X\in\Ob S\) is the derived category \(\Dpx(\mathcal C_X)\); if \(f:X\to Y\) is an arrow of \(S\), we shall write
\[
  f^*:\Dpx(\mathcal C_Y)\longrightarrow \Dpx(\mathcal C_X)
\]
for the inverse-image functor, obtained from
\[
  f^*:\Kpx(\mathcal C_Y)\longrightarrow \Kpx(\mathcal C_X)
\]
by passage to the quotient. One defines full triangulated sub-\(S\)-categories
\[
  \Kpx^*(\mathcal C)\qquad\text{respectively}\qquad \Dpx^*(\mathcal C)
  \quad (*=+,-,b)
\]
by the usual degree conditions.

\sgaStableTarget{sga6.e01.sec.1.2}\subsection*{1.2. Terminology}

Let \(S\) be a site, let \(\mathcal C\) be a flat abelian \(S\)-category \hyperlink{sga6.e01.sec.1.1}{(1.1)}, and let \(\mathcal C_0\) be a strictly full fibered sub-\(S\)-category.

Let \(F\in\Ob \mathcal C_X\), for \(X\in\Ob S\). The set of arrows \(f:U\to X\) in \(S/X\) such that there exists an epimorphism
\[
  E\longrightarrow f^*F,\qquad E\in\Ob \mathcal C_{0,U},
\]
is a sieve on \(X\). If this sieve is covering, we shall say that \(F\) is of \(\mathcal C_0\)-\emph{finite type} (or simply of finite type when no confusion is to be feared).

We shall say that \(\mathcal C_0\) is \emph{locally stable under kernels of epimorphisms} if the following condition is satisfied: for every \(X\in\Ob S\) and every epimorphism
\[
  u:E\longrightarrow F
\]
of \(\mathcal C_X\), with \(E,F\in\Ob\mathcal C_{0,X}\), the kernel of \(u\) is locally in \(\mathcal C_0\); in other words the set of \(f:U\to X\) in \(S/X\) such that
\[
  f^*\Ker(u)\in\Ob\mathcal C_{0,U}
\]
is a refinement of \(X\).

By the standard process of decomposing a long exact sequence into short exact sequences, one sees that the preceding condition is equivalent to the following apparently stronger one: for every \(X\in\Ob S\) and every exact sequence in \(\mathcal C_X\)
\sgaStableTarget{sga6.e01.eq.unnamed}\begin{equation}\tag{*}
  0\longrightarrow E^0\longrightarrow E^1\longrightarrow\cdots\longrightarrow E^n\longrightarrow 0
\end{equation}
with \(E^i\in\Ob\mathcal C_{0,X}\) for \(1\leq i\leq n\), the object \(E^0\) is locally in \(\mathcal C_0\).

We shall say that \(\mathcal C_0\) is \emph{locally quasi-stable under kernels of epimorphisms} if, for every exact sequence such as \((*)\), \(E^0\) is of \(\mathcal C_0\)-finite type.

On the other hand, we shall say that \(\mathcal C_0\) is \emph{locally liftable in} \(\mathcal C\) if, for every \(X\in\Ob S\), every epimorphism
\[
  u:E\longrightarrow F
\]
of \(\mathcal C_X\) with \(F\in\Ob\mathcal C_{0,X}\) locally admits a section. Equivalently, for \(\mathcal C_0\) to be locally liftable in \(\mathcal C\), it is necessary and sufficient that, for every diagram of \(\mathcal C_X\), \(X\in\Ob S\),
\[
\begin{tikzcd}
& H \arrow[d]\\
E \arrow[r,"u"'] & F
\end{tikzcd}
\]
where \(u\) is an epimorphism and \(H\in\Ob\mathcal C_0\), the set of \(U\) over \(X\) for which there exists a commutative diagram
\[
\begin{tikzcd}
& H_U \arrow[d]\\
E_U \arrow[r,"u_U"'] \arrow[ur] & F_U
\end{tikzcd}
\]
is a refinement of \(X\).

We shall say that \(\mathcal C_0\) is \emph{locally quasi-liftable in} \(\mathcal C\) if the following weaker condition is satisfied: for every \(X\in\Ob S\), and every epimorphism
\[
  u:E\longrightarrow F
\]
of \(\mathcal C_X\), with \(F\in\Ob\mathcal C_{0,X}\), the set of arrows \(f:U\to X\) of \(S/X\) for which there exists
\[
  v:F'\longrightarrow f^*E,\qquad F'\in\Ob\mathcal C_{0,U},
\]
such that \(f^*(u)v\) is an epimorphism, is a refinement of \(X\). This condition is equivalent to the following: for every diagram in \(\mathcal C_X\)
\[
\begin{tikzcd}
& H \arrow[d]\\
E \arrow[r,"u"'] & F ,
\end{tikzcd}
\]
where \(u\) is an epimorphism and \(H\in\Ob\mathcal C_0\), the set of \(U\) over \(X\) such that there exists a commutative diagram
\[
\begin{tikzcd}
G \arrow[r,"v"] \arrow[d] & H_U \arrow[d]\\
E_U \arrow[r,"u_U"'] & F_U
\end{tikzcd}
\]
with \(G\in\Ob\mathcal C_0\) and \(v\) an epimorphism, is a refinement of \(X\).

If \(S\) is a chaotic site, for example the punctual site (one object, one morphism), we shall omit the adverb ``locally'' in the preceding definitions.

\sgaStableTarget{sga6.e01.sec.1.3}\subsection*{1.3. Examples}

\textbf{a)} Let \(S\) be the punctual site, let \(A\) be a ring, and let \(\mathcal C\) be the category of left \(A\)-modules. The full subcategory \(\mathcal C_0\) of \(\mathcal C\) formed by the projective \(A\)-modules of finite type is stable under kernels of epimorphisms and is liftable in \(\mathcal C\); an \(A\)-module is of \(\mathcal C_0\)-finite type if and only if it is of finite type in the ordinary sense. The full subcategory \(\mathcal C_1\) of \(\mathcal C\) formed by the \(A\)-modules of finite type is quasi-liftable in \(\mathcal C\), but is not liftable in general; if \(A\) is noetherian, \(\mathcal C_1\) is stable under kernels of epimorphisms. The full subcategory \(\mathcal C'_0\) of \(\mathcal C\) formed by the free \(A\)-modules of finite type is liftable in \(\mathcal C\) and quasi-stable, though not stable in general, under kernels of epimorphisms.

\textbf{b)} Let \((S,\mathcal O_S)\) be a space ringed by local rings, let \(S\) be the site of open subsets of \(S\), and let \(\mathcal C\) be the fibered \(S\)-category, even split (SGA 1 VI 9), defined by the functor
\[
  U\longmapsto \{\text{category of left }\mathcal O_U\text{-modules}\}.
\]
Then \(\mathcal C\) is a flat abelian \(S\)-category \hyperlink{sga6.e01.sec.1.1}{(1.1)}. Let \(\mathcal C_0\) be the strictly full additive sub-\(S\)-category of \(\mathcal C\) defined by
\[
  U\longmapsto \{\text{category of finite-type free }\mathcal O_U\text{-modules}\}.
\]
An \(\mathcal O_S\)-module \(M\) is locally of \(\mathcal C_0\)-finite type if and only if \(M\) is of finite type in the ordinary sense (EGA \(0_{\mathrm I}\), 5.2). From EGA \(0_{\mathrm I}\), 5.2.7 and 5.4.9, it follows that \(\mathcal C_0\) is locally stable under kernels of epimorphisms and locally liftable in \(\mathcal C\).

\textbf{c)} Let \(G\) be a group, let \(S\) be the topos of \(G\)-sets (SGA 4 V 3.11), and let \(k\) be a field, regarded as the constant ring on \(S\). Let \(\mathcal C\) be the split \(S\)-category defined by
\[
  U\longmapsto \{\text{category of left }k_U\text{-modules}\};
\]
then \(\mathcal C\) is a flat abelian \(S\)-category. Let \(\mathcal C_0\) be the strictly full additive sub-\(S\)-category of \(\mathcal C\) defined by
\[
  U\longmapsto \{\text{category of finite-type free }k_U\text{-modules}\}.
\]
Since localization in \(S\) is equivalent to forgetting the operations of \(G\), it is clear that a \(G\)-\(k\)-module \(M\) is locally of \(\mathcal C_0\)-finite type if and only if \(M\) is finite-dimensional as a vector space over \(k\), and that \(\mathcal C_0\) is locally stable under kernels of epimorphisms and locally liftable in \(\mathcal C\).

The reader is left to look at the ``Galois'' variant of c), where \(G\) is a profinite group and \(S\) is the topos of sets on which \(G\) operates continuously.

In fact, examples b) and c) are particular cases of the following example.

\textbf{d)} Let \((S,\mathcal O_S)\) be a ringed topos (cf. SGA 4 IV 2.10), and let \(\mathcal C\) be the split \(S\)-category defined by
\[
  U\longmapsto \{\text{category of left }\mathcal O_U\text{-modules}\};
\]
then \(\mathcal C\) is a flat abelian \(S\)-category. Let \(\mathcal C_0\) be the strictly full additive sub-\(S\)-category of \(\mathcal C\) defined by
\[
  U\longmapsto \{\text{category of finite-type free }\mathcal O_U\text{-modules}\}.
\]
An \(\mathcal O_S\)-module locally of \(\mathcal C_0\)-finite type will simply be called \emph{of finite type}. The category \(\mathcal C_0\) is locally quasi-stable under kernels of epimorphisms and locally liftable in \(\mathcal C\), as follows from:

\sgaStableTarget{sga6.e01.lem.1-3-1}\textbf{Lemma 1.3.1.} Let \(P\) be an \(\mathcal O_S\)-module. The following conditions are equivalent:
\begin{enumerate}[label=(\roman*)]
\item \(P\) is of finite type and the functor \(\Hom(P,\cdot)\) is exact;
\item \(P\) is of finite type and every epimorphism \(M\to P\) locally admits a section;
\item \(P\) is locally a direct factor of a free module of finite type.
\end{enumerate}
The proof is immediate and is left as an exercise to the reader.

\sgaStableTarget{sga6.e01.sec.1.4}\subsection*{1.4. An ``extension'' lemma (cf. EGA \(0_{\mathrm{III}}\), 11.9.1)}

\sgaStableTarget{sga6.e01.prop.1-4-1}\textbf{Proposition 1.4.1.} Let \(S\) be a site, let \(\mathcal C\) be a flat abelian \(S\)-category \hyperlink{sga6.e01.sec.1.1}{(1.1)}, and let \(\mathcal C_0\) be a strictly full additive sub-\(S\)-category of \(\mathcal C\). Assume that \(\mathcal C_0\) is locally quasi-liftable in \(\mathcal C\) \hyperlink{sga6.e01.sec.1.2}{(1.2)}. Let \(S\in\Ob S\), let
\[
  u:E\longrightarrow F
\]
be a morphism, of degree \(0\), of complexes of \(\mathcal C_S\) (differentials of degree \(+1\)), let \(G\) be the cone of \(u\) (cf. [V], Chap. I, §1, 2.1), let \(n\in\mathbb Z\), and assume that \(\Ho^n(G)\) is of \(\mathcal C_0\)-finite type \hyperlink{sga6.e01.sec.1.2}{(1.2)}. Then the set of \(X\in\Ob(S/S)\) satisfying property \((P)\) below is a refinement of \(S\).

\((P)\): there exist an object \(M\in\mathcal C_{0,X}\) and arrows of \(\mathcal C_X\)
\[
  r:M\longrightarrow Z^{n+1}(E_X)=
  \Ker\!\left(E_X^{n+1}\xrightarrow{d^{n+1}}E_X^{n+2}\right),
  \qquad
  s:M\longrightarrow F_X^n
\]
such that the following diagram of \(\mathcal C_X\)
\[
\begin{tikzcd}[column sep=tiny]
\cdots \arrow[r] & E_X^{n-1}
 \arrow[r,"{\binom{d_E^{n-1}}{0}}"]
 \arrow[d,"u_X^{n-1}"'] &
E_X^n\oplus M
 \arrow[r,"{(d_E^n,r)}"]
 \arrow[d,"{(u_X^n,s)}"'] &
E_X^{n+1} \arrow[r] \arrow[d,"u_X^{n+1}"] &
\cdots\\
\cdots \arrow[r] & F_X^{n-1} \arrow[r,"d_F^{n-1}"'] &
F_X^n \arrow[r,"d_F^n"'] &
F_X^{n+1} \arrow[r] &
\cdots
\end{tikzcd}
\]
is a morphism of complexes \(u':E'\to F_X\), that is,
\[
  d_F^n s=u_X^{n+1}r,
\]
and such that
\begin{enumerate}[label=(\roman*)]
\item \(\Ho^{n+1}(u')\) is a monomorphism, with
\[
  \Im\Ho^{n+1}(u')=\Im\Ho^{n+1}(u_X);
\]
\item \(\Ho^n(u')\) is an epimorphism.
\end{enumerate}

In pictorial terms, one may say that, after modifying \(E\) and \(u\) locally in degree \(n\) only by adding, in degree \(n\), an object of \(\mathcal C_0\), one can locally make \(\Ho^{n+1}(u)\) injective, without changing its image, and \(\Ho^n(u)\) surjective.

\emph{Proof.} The procedure is as follows: first one shows that \(r\) and \(s\) may be chosen locally so as to make \(\Ho^{n+1}(u)\) injective; one verifies that the finiteness hypothesis on \(\Ho^n(G)\) is preserved; a new local correction in degree \(n\) then makes it possible to obtain \(\Ho^n(u)\) surjective.

The exact cohomology sequence
\[
  \cdots\longrightarrow \Ho^n(G)\longrightarrow
  \Ho^{n+1}(E)\longrightarrow
  \Ho^{n+1}(F)\longrightarrow \cdots
\]
shows that the kernel of \(\Ho^{n+1}(u)\) is of \(\mathcal C_0\)-finite type. In other words, the set of \(X\) over \(S\) for which there is an exact sequence
\[
  M\longrightarrow \Ho^{n+1}(E_X)
  \xrightarrow{\Ho^{n+1}(u)}
  \Ho^{n+1}(F_X),
  \qquad M\in\Ob\mathcal C_{0,X},
\]
is a refinement of \(S\). Since \(\mathcal C_0\) is locally quasi-liftable in \(\mathcal C\), one may moreover impose that the arrow
\[
  M\longrightarrow \Ho^{n+1}(E_X)
\]
factor through \(Z^{n+1}(E_X)\):
\[
  M\xrightarrow{r} Z^{n+1}(E_X)\longrightarrow \Ho^{n+1}(E_X).
\]
The composite
\[
  M\xrightarrow{u_X^{n+1}r} Z^{n+1}(F_X)\longrightarrow \Ho^{n+1}(F_X)
\]
is zero by construction, so \(u_X^{n+1}r\) factors through
\[
  B^{n+1}(F_X)=\Im\!\left(F_X^n\xrightarrow{d_F^n}F_X^{n+1}\right).
\]
Applying quasi-liftability once more, one may impose that \(u_X^{n+1}r\) factor through \(F_X^n\):
\[
  M\xrightarrow{s}F_X^n\longrightarrow B^{n+1}(F_X)\longrightarrow Z^{n+1}(F_X).
\]
Then \(u'\), defined as in \hyperlink{sga6.e01.prop.1-4-1}{(1.4.1)}, is plainly a morphism of complexes satisfying (i).

Let \(G'\) be the cone of \(u':E'\to F_X\). We show that \(\Ho^n(G')\) is of \(\mathcal C_0\)-finite type. By construction there is an exact sequence
\[
  0\longrightarrow E_X\xrightarrow{v}E'\longrightarrow M[-n]\longrightarrow 0,
\]
with \(u'v=u_X\). To compare \(G_X\) and \(G'\), one may, for instance, form the octahedron in \(\Kpx(\mathcal C_X)\):
\[
\begin{tikzpicture}[>=stealth,baseline=(GX.base),node distance=1.0cm]
  \node (EX) at (0,0) {$E_X$};
  \node (Ep) at (2.6,0) {$E'$};
  \node (FX) at (5.2,0) {$F_X$};
  \node (M)  at (1.35,1.25) {$M[-n]$};
  \node (GX) at (2.65,2.55) {$G_X$};
  \node (Gp) at (3.95,1.25) {$G'$};
  \node at (3.22,1.72) {$(*)$};
  \draw[->] (EX) -- node[below] {$v$} (Ep);
  \draw[->] (Ep) -- node[below] {$u'$} (FX);
  \draw[->] (Ep) -- (M);
  \draw[->] (M) -- (GX);
  \draw[->] (GX) -- (Gp);
  \draw[->] (FX) -- (Gp);
  \draw[->] (M) to[bend left=38] node[left] {$+1$} (EX);
  \draw[->] (Gp) to[bend left=28] node[above] {$+1$} (M);
  \draw[->] (Gp) to[bend left=28] node[right] {$+1$} (Ep);
  \draw[->] (GX) to[bend left=55] node[left] {$+1$} (EX);
  \draw[->] (FX) to[bend right=42] (GX);
\end{tikzpicture}
\]
The exact cohomology sequence of the distinguished triangle marked by the octahedron shows that
\[
  \Ho^n(G_X)\longrightarrow \Ho^n(G')
\]
is an epimorphism; hence \(\Ho^n(G')\) is indeed of \(\mathcal C_0\)-finite type.

The preceding correction therefore reduces us to the case where \(\Ho^{n+1}(u)\) is a monomorphism. We suppose this from now on. The exact cohomology sequence of \(E\to F\to G\) then gives an exact sequence
\[
  \Ho^n(E)\xrightarrow{\Ho^n(u)}
  \Ho^n(F)\longrightarrow \Ho^n(G)\longrightarrow 0.
\]
Since \(\Ho^n(G)\) is of \(\mathcal C_0\)-finite type, the set of \(X\) over \(S\) for which there exists an epimorphism
\[
  M\longrightarrow \Ho^n(G_X),\qquad M\in\Ob\mathcal C_{0,X},
\]
is, by definition, a refinement of \(S\). Using the quasi-liftability of \(\mathcal C_0\) in \(\mathcal C\), one may impose that this epimorphism factor through \(Z^n(F_X)\):
\[
  M\xrightarrow{s}Z^n(F_X)\longrightarrow \Ho^n(F_X)\longrightarrow \Ho^n(G_X).
\]
It is clear that \(u'\), defined as in \hyperlink{sga6.e01.prop.1-4-1}{(1.4.1)} with \(r=0\), is a morphism of complexes and that
\[
  \Ho^n(u'):\Ho^n(E_X)\oplus M\longrightarrow \Ho^n(F_X)
\]
is an epimorphism. This completes the proof of \hyperlink{sga6.e01.prop.1-4-1}{(1.4.1)}.

\sgaStableTarget{sga6.e01.lem.1-4-2}\textbf{Lemma 1.4.2.} Let \(\mathcal C\) be an abelian category, let \(u:E\to F\) be a morphism of complexes of \(\mathcal C\) (respectively an arrow of \(\Dpx(\mathcal C)\)), and let \(G\) be the cone of \(u\). For a given \(n\in\mathbb Z\), the following conditions are equivalent:
\begin{enumerate}[label=(\roman*)]
\item \(\Ho^i(G)=0\) for \(i\geq n\);
\item \(\Ho^i(u)\) is an isomorphism for \(i>n\) and an epimorphism for \(i=n\).
\end{enumerate}
The proof is immediate from the exact cohomology sequence.

\sgaStableTarget{sga6.e01.def.1-4-3}\textbf{Definition 1.4.3.} We shall say that \(u\) is an \(n\)-\emph{quasi-isomorphism} (respectively an \(n\)-\emph{isomorphism}) if \(u\) satisfies the equivalent conditions of \hyperlink{sga6.e01.lem.1-4-2}{(1.4.2)}.

That being so, one deduces at once from \hyperlink{sga6.e01.prop.1-4-1}{(1.4.1)} the following:

\sgaStableTarget{sga6.e01.cor.1-4-4}\textbf{Corollary 1.4.4.} Let \(S,S,\mathcal C,\mathcal C_0\) be as in \hyperlink{sga6.e01.prop.1-4-1}{(1.4.1)}. Let \(u:E\to F\) be an \((n+1)\)-quasi-isomorphism of complexes of \(\mathcal C_S\), with \(n\in\mathbb Z\) fixed, such that \(\Ho^n(G)\) is of \(\mathcal C_0\)-finite type. Then the set of \(X\) over \(S\) for which there exist \((M,r,s)\) as in \hyperlink{sga6.e01.prop.1-4-1}{(1.4.1)} such that \(u'\) is an \(n\)-quasi-isomorphism is a refinement of \(S\).

In pictorial terms, one may, using the finiteness condition on \(\Ho^n(G)\), locally ``gain one notch,'' that is, locally modify \(u\) in degree \(n\) by adding to \(E^n\) an object of \(\mathcal C_0\) in such a way as to make \(u\) an \(n\)-quasi-isomorphism.

\sgaStableTarget{sga6.e01.cor.1-4-5}\textbf{Corollary 1.4.5.} Let \(\mathcal C\) be an abelian category, let \(n\in\mathbb Z\), and let \(u:E\to F\) be an \(n\)-quasi-isomorphism of complexes of \(\mathcal C\). There exists a commutative diagram in \(\Cpx(\mathcal C)\)
\[
\begin{tikzcd}
E \arrow[r,"v"] \arrow[d,"u"'] & E' \arrow[dl,"u'"]\\
F &
\end{tikzcd}
\]
where \(u'\) is a quasi-isomorphism and \(v\) is a monomorphism inducing the identity in degrees \(\geq n\), that is,
\[
  E'^i=E^i,\qquad v^i=\id(E^i)\qquad\text{for }i\geq n.
\]

\emph{Proof.} Applying \hyperlink{sga6.e01.cor.1-4-4}{(1.4.4)} in the case where \(S\) is the punctual site and \(\mathcal C_0=\mathcal C\), one obtains a commutative triangle in \(\Cpx(\mathcal C)\)
\[
\begin{tikzcd}
E \arrow[r,"v_{n-1}"] \arrow[d,"u"'] & E'_{\,n-1} \arrow[dl,"u'_{n-1}"]\\
F &
\end{tikzcd}
\]
where \(u'_{n-1}\) is an \((n-1)\)-quasi-isomorphism and \(v_{n-1}\) is a monomorphism inducing the identity in degrees \(\geq n\). Iterating the process, one finds, for \(i\leq n-1\), a sequence of commutative triangles in \(\Cpx(\mathcal C)\)
\[
\begin{tikzcd}
E'_i \arrow[r,"v_{i-1}"] \arrow[d,"u'_i"'] & E'_{\,i-1} \arrow[dl,"u'_{i-1}"]\\
F &
\end{tikzcd}
\]
where \(u'_i\) is an \(i\)-quasi-isomorphism and \(v_{i-1}\) is a monomorphism inducing the identity in degrees \(\geq i\). The corollary follows by passing to the limit. More precisely, writing \(d'_i\) for the differential of \(E'_i\), define \(E'\) by
\[
  E'^i=E^i,\quad d'^i=d^i\quad (i\geq n),
\]
and
\[
  E'^i=(E'_i)^i,\quad d'^i=(d'_i)^i\quad (i<n),
\]
and then define \(v\) and \(u'\) by
\[
  v^j=(v_i)^j,\qquad u'^j=(u'_i)^j\qquad (i\leq j).
\]
One immediately verifies that \(E'\), \(v\), and \(u'\), so defined, answer the question.

In abbreviated form: one can extend an \(n\)-quasi-isomorphism into a quasi-isomorphism ``without changing anything given in degree \(>n\).''

\sgaStableTarget{sga6.e01.sec.2}\section*{2. Pseudo-coherent complexes}

\sgaStableTarget{sga6.e01.para.2-0}\textbf{2.0.} In this section, \(S\) denotes a site, \(S\) an object of \(S\), \(\mathcal C\) a flat abelian \(S\)-category \hyperlink{sga6.e01.sec.1.1}{(1.1)}, and \(\mathcal C_0\) a strictly full additive sub-\(S\)-category of \(\mathcal C\). We assume that \(\mathcal C_0\) is locally quasi-stable under kernels of epimorphisms \hyperlink{sga6.e01.sec.1.2}{(1.2)} and locally quasi-liftable in \(\mathcal C\) \hyperlink{sga6.e01.sec.1.2}{(1.2)}.

\sgaStableTarget{sga6.e01.def.2-1}\textbf{Definition 2.1.} Let \(X\in\Ob S\), let \(E\) be a complex of \(\mathcal C_X\), and let \(n\in\mathbb Z\). We shall say that \(E\) is \emph{strictly \(\mathcal C_0\)-\(n\)-pseudo-coherent} (respectively \(\mathcal C_0\)-pseudo-coherent, respectively \(\mathcal C_0\)-perfect) if \(E^i=0\) for \(i\) sufficiently large and \(E^i\in\Ob\mathcal C_{0,S}\) for \(i\geq n\) (respectively for every \(i\), respectively for every \(i\) and \(E^i=0\) for almost all \(i\)).

When no confusion is to be feared, we shall sometimes say ``strictly \(n\)-pseudo-coherent'' (respectively, and so on) instead of ``strictly \(\mathcal C_0\)-\(n\)-pseudo-coherent''.

\sgaStableTarget{sga6.e01.lem.2-1-1}\textbf{Lemma 2.1.1.} Let \(E\in\Ob\Cpx(\mathcal C_S)\) and \(n\in\mathbb Z\). If \(E\) is strictly \(n\)-pseudo-coherent and if \(\Ho^i(E)=0\) for \(i>n\), then \(\Ho^n(E)\) is of finite type.

\emph{Proof.} The sequence
\[
  0\longrightarrow Z^n(E)\longrightarrow E^n\longrightarrow E^{n+1}\longrightarrow\cdots\longrightarrow 0
\]
is exact, and \(E^i\in\Ob\mathcal C_0\) for \(i\geq n\). Since \(\mathcal C_0\) is locally quasi-stable under kernels of epimorphisms, it follows that \(Z^n(E)\) is of finite type, hence so is \(\Ho^n(E)\), which is a quotient of it.

\sgaStableTarget{sga6.e01.prop.2-2}\textbf{Proposition 2.2.} Let \(F\in\Ob\Dpx(\mathcal C_S)\) and \(n\in\mathbb Z\). Then:

\textbf{a)} If \(F\) is strictly \(n\)-pseudo-coherent (relative to \(\mathcal C_0\)) and if
\[
  u:E\longrightarrow F
\]
is a quasi-isomorphism, then the set of \(X\) over \(S\) such that there exists a quasi-isomorphism
\[
  v:F'\longrightarrow E_X
\]
with \(F'\) strictly \(n\)-pseudo-coherent is a refinement of \(S\).

\textbf{b)} The following conditions are equivalent:
\begin{enumerate}[label=(\roman*)]
\item the set of \(X\) over \(S\) such that there exists a quasi-isomorphism
\[
  u:E\longrightarrow F_X
\]
of \(\Cpx(\mathcal C_X)\), with \(E\) strictly \(n\)-pseudo-coherent, is a refinement of \(S\);
\item[\((\mathrm{i}')\)] the same condition as (i), but with \(u\) an isomorphism of \(\Dpx(\mathcal C_X)\);
\item the set of \(X\) over \(S\) such that there exists an \(n\)-quasi-isomorphism
\[
  u:E\longrightarrow F_X
\]
of \(\Cpx(\mathcal C_X)\), with \(E\) strictly perfect, is a refinement of \(S\);
\item[\((\mathrm{ii}')\)] the same condition as (ii), but with \(u\) an \(n\)-isomorphism of \(\Dpx(\mathcal C_X)\).
\end{enumerate}

\emph{Proof.} For a), there exists, for \(i\) sufficiently large, an \(i\)-quasi-isomorphism
\[
  v_i:F'_i\longrightarrow E
\]
with \(F'_i\) strictly \(n\)-pseudo-coherent (take \(F'_i=0\)). We show that, if \(i>n\), there exists, after refining \(S\), an \((i-1)\)-quasi-isomorphism
\[
  v_{i-1}:F'_{i-1}\longrightarrow E
\]
with \(F'_{i-1}\) strictly \(n\)-pseudo-coherent. By \hyperlink{sga6.e01.cor.1-4-4}{(1.4.4)}, it is enough to show that \(\Ho^{i-1}(M')\), where \(M'\) is the cone of \(v_i\), is of \(\mathcal C_0\)-finite type. Now if \(M''\) denotes the cone of \(uv_i\), the hypothesis that \(u\) is a quasi-isomorphism and \hyperlink{sga6.e01.lem.1-4-2}{(1.4.2)} imply that
\[
  \Ho^j(M')\qis \Ho^j(M'')\qquad\text{for all }j,
\]
and
\[
  \Ho^j(M')=0\qquad\text{for }j\geq i.
\]
The cone formula
\[
  M''{}^j=F'_i{}^{\,j+1}\oplus F^j
\]
shows, on the other hand, that \(M''\) is strictly \((i-1)\)-pseudo-coherent. Hence, by \hyperlink{sga6.e01.lem.2-1-1}{(2.1.1)},
\[
  \Ho^{i-1}(M'')=\Ho^{i-1}(M')
\]
is of finite type, whence the assertion. It follows, by iteration and after refining \(S\), that there is an \(n\)-quasi-isomorphism
\[
  v_n:F'_n\longrightarrow E
\]
where \(F'_n\) is strictly \(n\)-pseudo-coherent. But, by \hyperlink{sga6.e01.cor.1-4-5}{(1.4.5)}, one can extend \(v_n\) into a quasi-isomorphism
\[
  v:F'\longrightarrow E
\]
without changing anything in degrees \(\geq n\); in particular \(F'\) is strictly \(n\)-pseudo-coherent. This proves a).

For b), the equivalence of (i) and \((\mathrm{i}')\) follows at once from a) and from the definition of isomorphisms of \(\Dpx^-(\mathcal C_X)\). If \(E\in\Ob\Dpx^-(\mathcal C_S)\) is strictly \(n\)-pseudo-coherent, there is an exact sequence
\[
  0\longrightarrow E'\longrightarrow E\longrightarrow E''\longrightarrow 0
\]
where \(E'\) is strictly perfect and \(E''{}^i=0\) for \(i\geq n\). Consequently \(E'\to E\) is an \(n\)-quasi-isomorphism, which gives the implication \((\mathrm{i})\Rightarrow(\mathrm{ii})\). The implication \((\mathrm{ii})\Rightarrow(\mathrm{ii}')\) is trivial, so it remains to prove \((\mathrm{ii}')\Rightarrow(\mathrm{i})\). Let
\[
  u:E\longrightarrow F
\]
be an \(n\)-isomorphism of \(\Dpx^-(\mathcal C_S)\), with \(E\) strictly perfect. Then \(u\) is represented by a diagram in \(\Cpx^-(\mathcal C_S)\)
\[
\begin{tikzcd}
& E' \arrow[dl,"s"'] \arrow[dr,"u'"]\\
E && F ,
\end{tikzcd}
\]
where \(s\) is a quasi-isomorphism and \(u'\) an \(n\)-quasi-isomorphism. By a), after refining \(S\), there exists a quasi-isomorphism
\[
  t:E''\longrightarrow E',
\]
with \(E''\) strictly \(n\)-pseudo-coherent. Applying \hyperlink{sga6.e01.cor.1-4-5}{(1.4.5)} to \(u't:E''\to F\), we extend it to a quasi-isomorphism
\[
  E'''\longrightarrow F
\]
with \(E'''\) strictly \(n\)-pseudo-coherent. This proves \((\mathrm{ii}')\Rightarrow(\mathrm{i})\) and completes the proof of \hyperlink{sga6.e01.prop.2-2}{(2.2)}.

\sgaStableTarget{sga6.e01.def.2-3}\textbf{Definition 2.3.} Under the hypotheses of \hyperlink{sga6.e01.prop.2-2}{(2.2)}, we shall say that \(F\in\Ob\Dpx^-(\mathcal C_S)\) is \(n\)-\(\mathcal C_0\)-pseudo-coherent if \(F\) satisfies the equivalent conditions b) of \hyperlink{sga6.e01.prop.2-2}{(2.2)}. We shall say that \(F\) is \(\mathcal C_0\)-pseudo-coherent if \(F\) is \(n\)-\(\mathcal C_0\)-pseudo-coherent for every \(n\).

When this cannot cause confusion, we shall sometimes allow ourselves to say \(n\)-pseudo-coherent, pseudo-coherent, instead of \(n\)-\(\mathcal C_0\)-pseudo-coherent, \(\mathcal C_0\)-pseudo-coherent.

\sgaStableTarget{sga6.e01.rem.2-3-1}\textbf{Remark 2.3.1.} Let \(F\in\Ob\Dpx(\mathcal C_S)\) be a pseudo-coherent complex. The set of \(X\) over \(S\) such that there exists a quasi-isomorphism
\[
  E\longrightarrow F_X
\]
with \(E\) strictly pseudo-coherent \hyperlink{sga6.e01.def.2-1}{(2.1)} is not in general a refinement of \(S\). However, we shall see later (\hyperlink{sga6.e01.prop.2-7}{2.7} a) and \hyperlink{sga6.e02.expose}{Exposé II}) important cases in which this is so.

The proof of (\hyperlink{sga6.e01.prop.2-2}{2.2} a) gives the following result.

\sgaStableTarget{sga6.e01.prop.2-4}\textbf{Proposition 2.4.} Let \(n,i\in\mathbb Z\), and let
\[
  u:E\longrightarrow F
\]
be an \(i\)-quasi-isomorphism of \(\Cpx(\mathcal C_S)\), with \(F\) \(n\)-pseudo-coherent and \(E\) strictly \(n\)-pseudo-coherent. Then, after refining \(S\), there exists a commutative triangle in \(\Cpx(\mathcal C_S)\)
\[
\begin{tikzcd}
E \arrow[r,"v"] \arrow[d,"u"'] & E' \arrow[dl,"u'"]\\
F &
\end{tikzcd}
\]
where \(E'\) is strictly \(n\)-pseudo-coherent, \(u'\) is a quasi-isomorphism, and \(v\) is a monomorphism inducing the identity in degree \(\geq i\).

\sgaStableTarget{sga6.e01.prop.2-5}\textbf{Proposition 2.5.} \textbf{a)} If \(F\in\Ob\Dpx(\mathcal C_S)\) is \(n\)-pseudo-coherent (respectively pseudo-coherent), then \(F[m]\) is \((n-m)\)-pseudo-coherent (respectively pseudo-coherent).

\textbf{b)} Let
\sgaStableTarget{sga6.e01.eq.unnamed.occ2}\[
\begin{tikzcd}
E \arrow[rr] && F \arrow[dl]\\
& G \arrow[ul,"+1"] &
\end{tikzcd}
\tag{*}
\]
be a distinguished triangle of \(\Dpx(\mathcal C_S)\). If \(E\) and \(G\) are \(n\)-pseudo-coherent (respectively pseudo-coherent), then \(F\) is \(n\)-pseudo-coherent (respectively pseudo-coherent).

\emph{Proof.} a) is trivial from the definitions.

b) It is enough, of course, to prove the non-parenthesized assertion. The triangle \((*)\) can also be written
\[
\begin{tikzcd}
G[-1] \arrow[rr] && E \arrow[dl]\\
& F \arrow[ul,"+1"] &
\end{tikzcd}
\]
By a), \(G[-1]\) is \((n+1)\)-pseudo-coherent. Thus we are reduced to proving that, if \(E\) is \((n+1)\)-pseudo-coherent and \(F\) is \(n\)-pseudo-coherent, then \(G\) is \(n\)-pseudo-coherent. After localizing, we may assume \(E\) (respectively \(F\)) strictly \((n+1)\)-pseudo-coherent (respectively \(n\)-pseudo-coherent). The arrow
\[
  E\xrightarrow{u}F
\]
of \(\Dpx^-(\mathcal C_S)\) is represented by a pair of arrows of \(\Cpx^-(\mathcal C_S)\)
\[
\begin{tikzcd}
& E' \arrow[dl,"s"'] \arrow[dr,"u'"]\\
E && F ,
\end{tikzcd}
\]
where \(s\) is a quasi-isomorphism. By (\hyperlink{sga6.e01.prop.2-2}{2.2} a)), after localizing, there exists a quasi-isomorphism
\[
  t:E''\longrightarrow E',
\]
where \(E''\) is strictly \((n+1)\)-pseudo-coherent. We are therefore reduced to the case where \(u\) is defined by a true arrow of \(\Cpx^-(\mathcal C_S)\). The assertion then follows from the definitions and from the formula
\[
  C(u)^i=E^{i+1}\oplus F^i,
\]
where \(C(u)\) denotes the cone of \(u\).

\sgaStableTarget{sga6.e01.schol.2-6}\textbf{Scholium 2.6.} By rotation of triangles, one sees that pseudo-coherence properties for two of the objects of the triangle \((*)\) imply another such property for the third. The various cases may be summarized in the table below; the lower-case letters denote degrees of pseudo-coherence, and in each row the conclusion is underlined:
\[
\begin{array}{ccc}
E&F&G\\[0.2em]
n+1&n&\underline n\\
n&\underline n&n\\
\underline n&n&n-1 .
\end{array}
\]

We shall write
\[
  \Dpx(\mathcal C)_{\mathcal C_0\text{-}n\mathrm{coh}}
  \quad\text{respectively}\quad
  \Dpx(\mathcal C)_{\mathcal C_0\text{-}\mathrm{coh}}
\]
for the strictly full fibered subcategory of \(\Dpx(\mathcal C)\) whose fiber over each \(X\in\Ob S\) is formed by the objects of \(\Dpx(\mathcal C_X)\) that are \(n\)-pseudo-coherent (respectively pseudo-coherent) relative to \(\mathcal C_0\). By \hyperlink{sga6.e01.prop.2-5}{(2.5)}, the category
\[
  \Dpx(\mathcal C)_{\mathcal C_0\text{-}\mathrm{coh}}
\]
is a triangulated sub-\(S\)-category of \(\Dpx(\mathcal C)\) \hyperlink{sga6.e01.sec.1.1}{(1.1)}.

\sgaStableTarget{sga6.e01.prop.2-7}\textbf{Proposition 2.7.} Suppose that \(S\) is the punctual site, so that
\[
  \mathcal C=\mathcal C_S,\qquad \mathcal C_0=\mathcal C_{0,S}.
\]
Then:
\begin{enumerate}[label=\alph*)]
\item If \(F\in\Ob\Dpx(\mathcal C)\) is pseudo-coherent, there exists a quasi-isomorphism
\[
  E\longrightarrow F
\]
with \(E\) strictly pseudo-coherent \hyperlink{sga6.e01.def.2-1}{(2.1)}.
\item Let \(\Kpx^{-,\mathrm{ac}}(\mathcal C_0)\) denote the full triangulated subcategory of \(\Kpx^-(\mathcal C_0)\) formed by acyclic complexes. The canonical functor
\[
  \Kpx^-(\mathcal C_0)/\Kpx^{-,\mathrm{ac}}(\mathcal C_0)
  \longrightarrow \Dpx(\mathcal C)
\]
is fully faithful and has essential image \(\Dpx(\mathcal C)_{\mathcal C_0\text{-}\mathrm{coh}}\).
\end{enumerate}

\emph{Proof.} a) Let \(F\in\Ob\Dpx(\mathcal C)_{\mathcal C_0\text{-}\mathrm{coh}}\), and suppose constructed an \(i\)-quasi-isomorphism
\[
  u_i:E_i\longrightarrow F
\]
of \(\Cpx(\mathcal C)\), with \(E_i\) strictly pseudo-coherent; this is possible for \(i\) sufficiently large, taking \(E_i=0\). Let \(G_i\) be the cone of \(u_i\). It follows from \hyperlink{sga6.e01.lem.1-4-2}{(1.4.2)} and \hyperlink{sga6.e01.lem.2-1-1}{(2.1.1)} that \(\Ho^{i-1}(G_i)\) is of \(\mathcal C_0\)-finite type. Hence, by \hyperlink{sga6.e01.cor.1-4-4}{(1.4.4)}, there exists a commutative triangle in \(\Cpx(\mathcal C)\)
\[
\begin{tikzcd}
E_i \arrow[r,"v_{i-1}"] \arrow[d,"u_i"'] & E_{i-1} \arrow[dl,"u_{i-1}"]\\
F &
\end{tikzcd}
\]
where \(E_{i-1}\) is strictly pseudo-coherent, \(u_{i-1}\) an \((i-1)\)-quasi-isomorphism, and \(v_{i-1}\) a monomorphism inducing the identity in degrees \(\geq i\). Passing to the limit (cf. the proof of \hyperlink{sga6.e01.cor.1-4-5}{(1.4.5)}) gives a quasi-isomorphism
\[
  u:E\longrightarrow F
\]
with \(E\) strictly pseudo-coherent.

b) Apply the criterion of \([V]\), Chap. I, §2, 4.2 b)(i), together with a).

We note that the proof of a) gives the following more precise result, which generalizes \hyperlink{sga6.e01.cor.1-4-5}{(1.4.5)}.

\sgaStableTarget{sga6.e01.prop.2-7-1}\textbf{Proposition 2.7.1.} Under the hypotheses of \hyperlink{sga6.e01.prop.2-7}{(2.7)}, let \(u:E\to F\) be an \(n\)-quasi-isomorphism of \(\Cpx(\mathcal C)\), with \(F\) pseudo-coherent and \(E\) strictly pseudo-coherent. Then there exists a commutative triangle in \(\Cpx(\mathcal C)\)
\[
\begin{tikzcd}
E \arrow[r,"v"] \arrow[d,"u"'] & E' \arrow[dl,"u'"]\\
F &
\end{tikzcd}
\]
where \(E'\) is strictly pseudo-coherent, \(u'\) is a quasi-isomorphism, and \(v\) is a monomorphism inducing the identity in degree \(\geq n\).

\sgaStableTarget{sga6.e01.def.2-8}\textbf{Definition 2.8.} Let \(n\in\mathbb N\) and \(F\in\Ob\mathcal C_S\). We shall say that \(F\) is of \(n\)-\(\mathcal C_0\)-finite presentation if the set of \(X\) over \(S\) such that there exists an exact sequence
\[
  E^{-n}\longrightarrow\cdots\longrightarrow E^0\longrightarrow F_X\longrightarrow 0,
\qquad E^i\in\Ob\mathcal C_{0,X}\quad(-n\leq i\leq 0),
\]
is a refinement of \(S\). We shall say that \(F\) is of \(\infty\)-\(\mathcal C_0\)-finite presentation if \(F\) is of \(n\)-\(\mathcal C_0\)-finite presentation for every \(n\in\mathbb N\).

Of course, we shall sometimes simply say \(n\)-finite presentation when there is no possible ambiguity about \(\mathcal C_0\). We shall say ``finite presentation'' for ``1-finite presentation,'' and ``finite type'' for ``0-finite presentation,'' the latter notion coinciding with the one already introduced in \hyperlink{sga6.e01.sec.1.2}{(1.2)}.

\sgaStableTarget{sga6.e01.prop.2-9}\textbf{Proposition 2.9.} Let \(n\in\mathbb N\) and \(F\in\Ob\mathcal C_S\). The following conditions are equivalent:
\begin{enumerate}[label=(\roman*)]
\item \(F\) is of \(n\)-finite presentation (respectively of \(\infty\)-finite presentation);
\item \(F\) is \((-n)\)-pseudo-coherent (respectively pseudo-coherent).
\end{enumerate}
In the case where \(S\) is the punctual site, \(F\) is pseudo-coherent if and only if there exists an exact sequence
\[
  \cdots\longrightarrow E^i\longrightarrow\cdots\longrightarrow E^0\longrightarrow F\longrightarrow 0,
\qquad E^i\in\Ob\mathcal C_0\quad\text{for all }i.
\]

\emph{Proof.} For the equivalence \((\mathrm{i})\Leftrightarrow(\mathrm{ii})\), it is plainly enough to treat the non-parenthesized case. The implication \((\mathrm{i})\Rightarrow(\mathrm{ii})\) is immediate from the definitions, and \((\mathrm{ii})\Rightarrow(\mathrm{i})\) follows at once from \hyperlink{sga6.e01.prop.2-4}{(2.4)}: the arrow \(0\to F\) is a \(1\)-quasi-isomorphism, and from this one deduces, after localizing, a quasi-isomorphism
\[
  E\longrightarrow F
\]
with \(E\) strictly \((-n)\)-pseudo-coherent and \(E^i=0\) for \(i\geq 1\), that is, an \(n\)-finite presentation of \(F\). The second assertion is proved in the same way, using \hyperlink{sga6.e01.prop.2-7-1}{(2.7.1)}.

\sgaStableTarget{sga6.e01.prop.2-10}\textbf{Proposition 2.10.} \textbf{a)} Let \(m,n\in\mathbb Z\), with \(n\leq m\). If \(F\in\Ob\Dpx^-(\mathcal C_S)\) is \(n\)-pseudo-coherent and acyclic in degrees \(\geq m\), then there exists locally a quasi-isomorphism
\[
  E\longrightarrow F
\]
with \(E\) strictly \(n\)-pseudo-coherent and \(E^i=0\) for \(i\geq m\). If \(S\) is punctual, and if \(F\) is pseudo-coherent and acyclic in degrees \(\geq m\), there exists a quasi-isomorphism
\[
  E\longrightarrow F
\]
with \(E\) strictly pseudo-coherent and \(E^i=0\) for \(i\geq m\).

\textbf{b)} Let \(F\in\Ob\Dpx^-(\mathcal C_S)\) be acyclic in degrees \(>n\). For \(F\) to be \(n\)-pseudo-coherent, it is necessary and sufficient that \(\Ho^n(F)\) be of finite type.

\emph{Proof.} a) Since \(F\) is acyclic in degrees \(\geq m\), the arrow \(0\to F\) is an \(m\)-quasi-isomorphism. The assertion follows from \hyperlink{sga6.e01.prop.2-4}{(2.4)} and \hyperlink{sga6.e01.prop.2-7-1}{(2.7.1)}.

b) There exists a distinguished triangle in \(\Dpx^-(\mathcal C_S)\):
\[
\begin{tikzcd}
F \arrow[rr] && F' \arrow[dl]\\
& \Ho^n(F)[-n] \arrow[ul,"+1"] &
\end{tikzcd}
\]
(to see this, begin by truncating \(F\), replacing \(F^n\) by \(Z^n(F)\) and \(F^i\) by \(0\) for \(i>n\)). The exact cohomology sequence implies
\[
  \Ho^i(F')=0\quad\text{for }i\geq n-1,
\]
so \(F'\) is \((n-1)\)-pseudo-coherent. We conclude by \hyperlink{sga6.e01.prop.2-5}{(2.5)} and \hyperlink{sga6.e01.schol.2-6}{(2.6)}.



% Additional notation for Expose I continuation
\providecommand{\RHom}{\mathrm R\!\Hom}
\providecommand{\Ext}{\operatorname{Ext}}
\providecommand{\coh}{\operatorname{coh}}
\providecommand{\PreCoh}{\operatorname{precoh}}
\providecommand{\Tr}{\operatorname{Tr}}
\providecommand{\Qis}{\operatorname{Qis}}
\providecommand{\Rad}{\operatorname{rad}}
\providecommand{\Spec}{\operatorname{Spec}}
\providecommand{\End}{\operatorname{End}}
\providecommand{\Id}{\operatorname{Id}}
\providecommand{\varinjlim}{\mathop{\mathrm{lim}}\limits_{\longrightarrow}}
\providecommand{\varprojlim}{\mathop{\mathrm{lim}}\limits_{\longleftarrow}}
\providecommand{\qcoh}{\operatorname{qcoh}}
\providecommand{\amp}{\operatorname{amp}}


\sgaStableTarget{sga6.e01.expose.occ2}\section*{Exposé I, continued: pseudo-coherent complexes and perfect complexes}

\sgaStableTarget{sga6.e01.prop.2-11}\textbf{Proposition 2.11.} Let \(n\in\mathbb Z\).

\textbf{a)} Let \(F\in\Ob\Cpx^-(\mathcal C_S)\). If, for every \(i\geq n\) (respectively for every \(i\)), \(F^i\) is of finite \((i-n)\)-presentation (respectively of finite \(\infty\)-presentation), then \(F\) is \(n\)-pseudo-coherent (respectively pseudo-coherent).

\textbf{b)} Let \(F\in\Ob\Dpx^-(\mathcal C_S)\). If, for every \(i\geq n\) (respectively for every \(i\)), \(\Ho^i(F)\) is of finite \((i-n)\)-presentation (respectively of finite \(\infty\)-presentation), then \(F\) is pseudo-coherent.

\emph{Proof.} It is enough to treat the non-parenthesized assertions.

\textbf{a)} We argue by induction on \(N\) such that \(F^i\neq0\) only for \(i\leq n+N\). If \(N=0\), the assertion follows at once from (\hyperlink{sga6.e01.prop.2-10}{2.10} b). Suppose \(N\geq1\). We have an exact sequence
\[
  0\longrightarrow F^{n+N}[-n-N]\longrightarrow F\longrightarrow F'\longrightarrow 0 .
\]
By the induction hypothesis, \(F'\) is \(n\)-pseudo-coherent; on the other hand \(F^{n+N}[-n-N]\) is \(n\)-pseudo-coherent by hypothesis, using (\hyperlink{sga6.e01.prop.2-9}{2.9} and \hyperlink{sga6.e01.prop.2-5}{2.5} a). Hence, by (\hyperlink{sga6.e01.prop.2-5}{2.5} b), \(F\) is \(n\)-pseudo-coherent.

\textbf{b)} We argue by induction on \(N\) such that \(\Ho^i(F)=0\) for \(i>n+N\). If \(N=0\), the assertion follows from (\hyperlink{sga6.e01.prop.2-10}{2.10} b). Suppose \(N\geq1\). As in the proof of (\hyperlink{sga6.e01.prop.2-10}{2.10} b), we have a distinguished triangle in \(\Dpx^-(\mathcal C_S)\)
\[
  F'\longrightarrow F\longrightarrow \Ho^{n+N}(F)[-n-N]\longrightarrow F'[1].
\]
more plainly, with cone \(\Ho^{n+N}(F)[-n-N]\). The exact cohomology sequence implies
\[
  \Ho^i(F')=0\quad\text{for } i>n+N-1,
\]
that is, for \(i>(n-1)+(N-1)\), and
\[
  \Ho^i(F')\qis \Ho^{i+1}(F)\quad\text{for }i<n+N-1 .
\]
By the induction hypothesis, applied with \(n-1\) in place of \(n\), \(F'\) is therefore \((n-1)\)-pseudo-coherent. On the other hand, \(\Ho^{n+N}(F)[-n-N]\) is \(n\)-pseudo-coherent by (\hyperlink{sga6.e01.prop.2-9}{2.9} and \hyperlink{sga6.e01.prop.2-5}{2.5} a). Hence, by \hyperlink{sga6.e01.schol.2-6}{(2.6)}, \(F\) is \(n\)-pseudo-coherent.

\sgaStableTarget{sga6.e01.rem.2-11-1}\textbf{Remark 2.11.1.} If \(F\) is pseudo-coherent, the \(F^i\) (respectively the \(\Ho^i(F)\)) are not in general pseudo-coherent, nor even of finite type. Nevertheless, in the following paragraph, in the excursions on the notions of coherence, we shall see important examples where pseudo-coherence of \(F\) implies that of the \(\Ho^i(F)\).

\sgaStableTarget{sga6.e01.prop.2-12}\textbf{Proposition 2.12.} Let \(F',F''\in\Ob\Dpx^-(\mathcal C_S)\), and let \(F=F'\oplus F''\). For \(F\) to be \(n\)-pseudo-coherent (respectively pseudo-coherent), it is necessary and sufficient that \(F'
\) and \(F''\) be \(n\)-pseudo-coherent (respectively pseudo-coherent).

\emph{Proof.} It is enough to prove the non-parenthesized assertion. Sufficiency is obvious; let us prove necessity. By induction we may suppose that it has already been proved that \(F'\) and \(F''\) are \((n+1)\)-pseudo-coherent. After localizing, we may even suppose that \(F'\) and \(F''\) are strictly \((n+1)\)-pseudo-coherent and that \(F=F'\oplus F''\) in \(\Cpx^-(\mathcal C_S)\). Let \(F'_{>n}\) (respectively \(F'_{\leq n}\)) denote the complex obtained from \(F'\) by replacing \(F'^i\) by \(0\) for \(i\leq n\) (respectively for \(i>n\)), and define \(F''_{>n}\) and \(F''_{\leq n}\) similarly. We have an exact sequence
\[
  0\longrightarrow F'_{>n}\oplus F''_{>n}\longrightarrow F
  \longrightarrow F'_{\leq n}\oplus F''_{\leq n}\longrightarrow 0 .
\]
The complex \(F'_{>n}\oplus F''_{>n}\) is strictly perfect; since \(F\) is \(n\)-pseudo-coherent, \hyperlink{sga6.e01.schol.2-6}{(2.6)} implies that \(F'_{\leq n}\oplus F''_{\leq n}\) is \(n\)-pseudo-coherent. This means, by (\hyperlink{sga6.e01.prop.2-10}{2.10} b), that
\[
  \Ho^n(F'_{\leq n}\oplus F''_{\leq n})
\]
is of finite type. But
\[
  \Ho^n(F'_{\leq n}\oplus F''_{\leq n})
  =\Ho^n(F'_{\leq n})\oplus \Ho^n(F''_{\leq n}),
\]
so \(\Ho^n(F'_{\leq n})\) (respectively \(\Ho^n(F''_{\leq n})\)) is of finite type, and again by (\hyperlink{sga6.e01.prop.2-10}{2.10} b) \(F'_{\leq n}\) (respectively \(F''_{\leq n}\)) is \(n\)-pseudo-coherent. Applying \hyperlink{sga6.e01.prop.2-5}{(2.5)} to the exact sequence
\[
  0\longrightarrow F'_{>n}\longrightarrow F'\longrightarrow F'_{\leq n}\longrightarrow 0
\]
(and similarly for \(F''\)), we deduce that \(F'\) and \(F''\) are \(n\)-pseudo-coherent.

\sgaStableTarget{sga6.e01.prop.2-13}\textbf{Proposition 2.13.} Let \((S',\mathcal C',\mathcal C'_0)\) be a triple satisfying the same hypotheses \hyperlink{sga6.e01.para.2-0}{(2.0)} as \((S,\mathcal C,\mathcal C_0)\). Let \(f:S'\to S\) be a morphism of sites, corresponding to an underlying functor \(f^{-1}:S\to S'\), and let
\[
  u:\Dpx^*(\mathcal C)\longrightarrow \Dpx^*(\mathcal C')\qquad (*=+,-,b)
\]
be a Cartesian functor above \(f\), exact on each fiber. Suppose that there exist \(a,b\in\mathbb Z\) such that, for every \(S\in\Ob S\), the following conditions hold:

\begin{enumerate}[label=(\roman*)]
\item if \(E\in\Ob\mathcal C_{0,S}\), then \(u(E)\) is \(a\)-\(\mathcal C'_0\)-pseudo-coherent (respectively \(\mathcal C'_0\)-pseudo-coherent);
\item if \(F\in\Ob\Dpx^*(\mathcal C_S)\) is acyclic in degrees \(>0\), then \(u(F)\) is acyclic in degrees \(>b\).
\end{enumerate}

Then, if \(F\in\Ob\Dpx^*(\mathcal C_S)\) is \(m\)-pseudo-coherent relative to \(\mathcal C_0\) and acyclic in degrees \(>n\), the object
\[
  u(F)\in\Ob\Dpx^*(\mathcal C'_{S'}),\qquad S'=f^{-1}(S),
\]
is \(\sup(m+b,n+a)\)-pseudo-coherent relative to \(\mathcal C'_0\) (respectively \((m+b)\)-pseudo-coherent), and is acyclic in degrees \(>n+b\). In particular, in the parenthesized case, if \(F\in\Ob\Dpx^*(\mathcal C_S)\) is pseudo-coherent relative to \(\mathcal C_0\), then \(u(F)\) is pseudo-coherent relative to \(\mathcal C'_0\).

\emph{Proof.} We may restrict to the non-parenthesized case. First suppose that \(F\in\Ob\Dpx^*(\mathcal C_S)\) is strictly perfect \hyperlink{sga6.e01.def.2-1}{(2.1)}, with \(F^i=0\) for \(i>n\). Then \(u(F)\) is \((n+a)\)-\(\mathcal C'_0\)-pseudo-coherent. We argue by induction on the number \(N\) of indices \(i\) such that \(F^i\neq0\). If \(N=1\), then \(F=F^p[-p]\) for some \(p\leq n\), so \(u(F)\) is \((p+a)\)-\(\mathcal C'_0\)-pseudo-coherent, hence a fortiori \((n+a)\)-pseudo-coherent. If \(N>1\), there is an exact sequence
\[
  0\longrightarrow F'\longrightarrow F\longrightarrow F''\longrightarrow 0,
\]
where \(F'\) and \(F''\) are strictly perfect of length \(<N-1\), and where \(F'^i=F''^i=0\) for \(i>n\). Applying (\hyperlink{sga6.e01.prop.2-5}{2.5} b) to the distinguished triangle
\[
  u(F')\longrightarrow u(F)\longrightarrow u(F'')\longrightarrow u(F')[1].
\]
and using the induction hypothesis, we get that \(u(F)\) is \((n+a)\)-pseudo-coherent.

Now pass to the general case: \(F\) is \(m\)-\(\mathcal C_0\)-pseudo-coherent and acyclic in degrees \(>n\). After localizing, we may, by (\hyperlink{sga6.e01.prop.2-10}{2.10} a), suppose \(F\) strictly \(m\)-pseudo-coherent and such that \(F^i=0\) for \(i>n\). There is then an exact sequence
\[
  0\longrightarrow F'\longrightarrow F\longrightarrow F''\longrightarrow 0,
\]
where \(F''\) is zero in degrees \(\geq m\), and \(F'\) is strictly perfect and zero in degrees \(>n\). In the distinguished triangle
\[
  u(F')\longrightarrow u(F)\longrightarrow u(F'')\longrightarrow u(F')[1].
\]
the object \(u(F'')\) is acyclic in degrees \(>m+b\), hence \((m+b)\)-pseudo-coherent, while \(u(F')\) is \((n+a)\)-pseudo-coherent by what was just proved. We conclude by (\hyperlink{sga6.e01.prop.2-5}{2.5} b).

\sgaStableTarget{sga6.e01.cor.2-14}\textbf{Corollary 2.14.} Let \(\mathcal C_1\) be a strictly full additive sub-\(S\)-category of \(\mathcal C\) such that
\[
  \mathcal C_0\subset \mathcal C_1\subset \mathcal C.
\]
Assume that every object of \(\mathcal C_1\) is \(\mathcal C_0\)-pseudo-coherent. Then:

\begin{enumerate}[label=(\roman*)]
\item \(\mathcal C_1\) is locally quasi-liftable in \(\mathcal C\) \hyperlink{sga6.e01.sec.1.2}{(1.2)}, and locally quasi-stable under kernels of epimorphisms \hyperlink{sga6.e01.sec.1.2}{(1.2)};
\item for \(S\in\Ob S\) and \(F\in\Ob\Dpx(\mathcal C_S)\), in order that \(F\) be \(n\)-pseudo-coherent (respectively pseudo-coherent) relative to \(\mathcal C_1\), it is necessary and sufficient that it be so relative to \(\mathcal C_0\).
\end{enumerate}

\emph{Proof.} For (i), let \(E\to F\) be an epimorphism in \(\mathcal C_S\), with \(F\in\Ob\mathcal C_{1,S}\). Since \(F\) is \(\mathcal C_0\)-pseudo-coherent, hence a fortiori of \(\mathcal C_0\)-finite type, there exists, after localizing, an epimorphism \(E'\to F\) with \(E'\in\Ob\mathcal C_{0,S}\). Using the quasi-liftability of \(\mathcal C_0\) in \(\mathcal C\), we may suppose, after localizing again, that this epimorphism factors through \(E\); this gives the quasi-liftability of \(\mathcal C_1\). On the other hand, if \(E\in\Ob\mathcal C_S\) admits a finite right resolution by objects of \(\mathcal C_1\), then \(E\) is \(\mathcal C_0\)-pseudo-coherent by (\hyperlink{sga6.e01.prop.2-11}{2.11} a), hence a fortiori of \(\mathcal C_1\)-finite type.

For (ii), it is enough to prove necessity. For this one may apply, at choice, (\hyperlink{sga6.e01.prop.2-11}{2.11} a) or \hyperlink{sga6.e01.prop.2-13}{(2.13)}, with \(S'=S\), \(\mathcal C'=\mathcal C\), \(f=\Id\), \(u=\Id\), and with \(\mathcal C'_0\) and \(\mathcal C_0\) replaced respectively by \(\mathcal C_0\) and \(\mathcal C_1\).

\sgaStableTarget{sga6.e01.schol.2-14-1}\textbf{Scholium 2.14.1.} There is therefore a largest strictly full additive sub-\(S\)-category of \(\mathcal C\) giving the same notion of \(n\)-pseudo-coherence (respectively pseudo-coherence) as \(\mathcal C_0\), namely the category formed by the \(\mathcal C_0\)-pseudo-coherent objects. This latter category is locally stable under kernels of epimorphisms.

\sgaStableTarget{sga6.e01.ex.2-15}\textbf{Example 2.15.} Let \((S,A)\) be a ringed topos, and let \(\mathcal C,\mathcal C_0\) be the \(S\)-categories considered in (\hyperlink{sga6.e01.sec.1.3}{1.3} d). The fibered category \(\Dpx(\mathcal C)\) will often be denoted by \(\Dpx(S,A)\), or \(\Dpx_A(S)\), when one wants to suggest that the modules are left \(A\)-modules, or even simply \(\Dpx(S)\) if there is no ambiguity about \(A\). We denote by \(S\) the final object of \(S\), and put \(\Dpx(\mathcal C_S)=\Dpx_A(S)\) (or \(\Dpx(S)\)). An object of \(\Dpx(S)\) will be said to be \(n\)-pseudo-coherent (respectively pseudo-coherent) if it is \(n\)-\(\mathcal C_0\)-pseudo-coherent (respectively \(\mathcal C_0\)-pseudo-coherent), and similarly for the strict notions. The category fibered in \(\Dpx(\mathcal C)\), denoted \(\Dpx(\mathcal C)_{n-\mathcal C_0-\coh}\) (respectively \(\Dpx(\mathcal C)_{\mathcal C_0-\coh}\)) in \hyperlink{sga6.e01.schol.2-6}{(2.6)}, will be denoted \(\Dpx_A(S)_{n-\coh}\) (respectively \(\Dpx_A(S)_{\coh}\)). One observes that, by \hyperlink{sga6.e01.cor.2-14}{(2.14)}, the same notion of pseudo-coherence is obtained if \(\mathcal C_0\) is replaced by the \(S\)-category \(\mathcal C_1\) defined by
\[
  U\longmapsto \text{the category of left }A_U\text{-modules locally direct summands of free finite-type modules}.
\]
If an \(A\)-module \(M\) is a direct summand of a free finite-type \(A\)-module, it does not follow in general that \(M\) is locally free of finite type. One does, however, have the following result.

\sgaStableTarget{sga6.e01.lem.2-15-1}\textbf{Lemma 2.15.1.} Suppose that \(S\) satisfies one of the following two hypotheses:
\begin{enumerate}[label=(\roman*)]
\item \(S\) has enough points (SGA 4 IV 6.4), and for every point \(s\) of \(S\), the left \(A_s\)-modules which are projective of finite type are free; this last condition is satisfied, for example, if \(A_s/\Rad(A_s)\) is a field or if \(\Rad(A_s)\) is nilpotent;
\item \((S,A)\) is a ``locally ringed'' topos in the sense of [1], which means that \(A\) is commutative and that, for every \(X\in\Ob S\) and every \(f\in\Gamma(X,A)\), one has
\[
  X=X_f\cup X_{1-f},
\]
where \(X_f\) (respectively \(X_{1-f}\)) denotes the largest subobject of \(X\) on which \(f\) (respectively \(1-f\)) is invertible.
\end{enumerate}
Then every direct summand of a free finite-type left \(A\)-module is locally free of finite type.

\emph{Proof.} Suppose first that (i) holds. We note that, if \(E\) and \(F\) are two \(A\)-modules and \(E\) is of finite presentation, the canonical arrow
\[
  \Hom(E,F)_s\longrightarrow \Hom(E_s,F_s)
\]
is an isomorphism for every point \(s\) of \(S\). It follows that, if \(E\) and \(F\) are both of finite presentation and if \(E_s\qis F_s\) at a point \(s\) of \(S\), then there is an \(s\)-pointed object \(U\) of \(S\) such that \(E|U\qis F|U\). The conclusion follows at once.

Suppose next that (ii) holds. Let \(n\in\mathbb N\), and let \(\mathbb D^n\) denote the functor on the category of locally ringed topoi belonging to a fixed universe, defined by
\[
  T\longmapsto \Gamma(T,\mathcal O_T)^n.
\]
It is known [1] that the functor \(\End_{\mathbb D}(\mathbb D^n)\) is represented by
\[
  M_n=\Spec\mathbb Z[T_{ij}]\quad(1\leq i\leq n,
  1\leq j\leq n),
\]
endowed with the endomorphism \(U_n\) of \((\mathcal O_{M_n})^n\) defined by the matrix \((T_{ij})\). Now let
\[
  p:A^n\longrightarrow A^n
\]
be a projector. We show that \(p(A^n)\) is locally free of finite type. From what has just been said, there exists a morphism
\[
  s:S\longrightarrow M_n
\]
of locally ringed topoi such that \(s^*(u_n)=p\). We may interpret \(s\) as a point of \(M_n\) with values in \(\Gamma(S,A)\), that is, as a matrix of order \(n\) with coefficients in \(\Gamma(S,A)\). Since this matrix is a projector, \(s\) factors through the closed subscheme
\[
  i:P_n\hookrightarrow M_n
\]
corresponding to projectors; say \(S\xrightarrow{s'}P_n\xrightarrow{i}M_n\). Hence
\[
  p=s'^* i^*(u_n),
\]
and we are done, because, by the assertion under hypothesis (i), the image of \(i^*(u_n)\) is a locally free finite-type sheaf.

\sgaStableTarget{sga6.e01.cor.2-16}\textbf{Corollary 2.16.} Let
\[
  f:(S',A')\longrightarrow (S,A)
\]
be a morphism of ringed topoi, defined by a functor \(f^{-1}:S\to S'\) and a morphism of sheaves of rings \(f^{-1}(A)\to A'\). Let \(\Dpx(_{A'}S'_A)\) denote the derived category of \(A'\)-\(A\)-bimodules, left over \(A'\) and right over \(A\), and let
\[
  \otimes_A^{\mathrm L}:\Dpx^-(_{A'}S'_A)\times_{S'}\Dpx^-(_A S)
  \longrightarrow \Dpx^-(_{A'}S')
\]
be the derived functor of the tensor product
\[
  (E',E)\longmapsto E'\otimes_A f^{-1}(E).
\]
Let \(X\in\Ob S\), put \(X'=u(X)\), and let \(E'\in\Ob\Dpx^-(_{A'}X'_A)\), \(E\in\Ob\Dpx^-(_A X)\). If \(E'\) is \(m'\)-pseudo-coherent (as an object of \(\Dpx^-(_A S')\)) and acyclic in degrees \(>n'\), and if \(E\) is \(m\)-pseudo-coherent and acyclic in degrees \(>n\), then
\[
  E'\otimes_A^{\mathrm L}E
\]
is \(\sup(m+n',m'+n)\)-pseudo-coherent and acyclic in degrees \(>n+n'\).

\emph{Proof.} After replacing \(S\) by \(S/X\), we may suppose \(X=S\), and apply \hyperlink{sga6.e01.prop.2-13}{(2.13)} with
\[
  u=E'\otimes_A^{\mathrm L}:\Dpx^-(_A S)\longrightarrow \Dpx^-(_{A'}S').
\]
Let \(U\in\Ob S\), let \(U'=f^{-1}(U)\), and let \(E\) be a complex of left \(A_U\)-modules. If \(E\) is acyclic in degrees \(>n\), then \(E'\otimes_A^{\mathrm L}E\) is acyclic in degrees \(>n+n'\), by the definition of \(\otimes_A^{\mathrm L}\). On the other hand, if \(E=(A_U)^p\), then
\[
  E'\otimes_A^{\mathrm L}E=(E'|U')^p,
\]
which is \(m'\)-pseudo-coherent by hypothesis and \hyperlink{sga6.e01.prop.2-12}{(2.12)}. The conditions of \hyperlink{sga6.e01.prop.2-13}{(2.13)} are therefore satisfied, with \(a=m'\) and \(b=n'\), whence the conclusion.

In particular:

\sgaStableTarget{sga6.e01.cor.2-16-1}\textbf{Corollary 2.16.1.} Under the hypotheses of \hyperlink{sga6.e01.cor.2-16}{(2.16)}, one has:

\textbf{a)} The functor
\[
  \mathrm L f^*: \Dpx^-(_A S)\longrightarrow \Dpx^-(_{A'}S')
\]
(which is none other than the functor \(A'\otimes_A^{\mathrm L}\)) induces a functor
\[
  \mathrm L f^*: \Dpx^-(_A S)_{\coh}\longrightarrow \Dpx^-(_{A'}S')_{\coh}
\]
(with the notation of \hyperlink{sga6.e01.schol.2-6}{(2.6)}).

\textbf{b)} If \(A\) is commutative, the functor
\[
  \otimes_A^{\mathrm L}:\Dpx^-(S)\times_S\Dpx^-(S)
  \longrightarrow \Dpx^-(S)
\]
induces a functor
\[
  \otimes_A^{\mathrm L}:\Dpx^-(S)_{\coh}\times_S\Dpx^-(S)_{\coh}
  \longrightarrow \Dpx^-(S)_{\coh} .
\]

\sgaStableTarget{sga6.e01.exercise.2-16-2}\textbf{Exercise 2.16.2.} Let \(S\) be a topos, and let \(A\to A'\) be a morphism of rings of \(S\). Assume one of the following hypotheses:
\begin{enumerate}[label=(\roman*)]
\item \(A\to A'\) is surjective, with nilpotent ideal as kernel;
\item the functor \(A'\otimes_A-\), on the category of left \(A\)-modules, is exact and faithful.
\end{enumerate}
Show that \(E\in\Ob\Dpx^-(_A S)\) is \(n\)-pseudo-coherent (respectively pseudo-coherent) if and only if
\[
  A'\otimes_A E\in\Ob\Dpx^-(_{A'}S)
\]
is \(n\)-pseudo-coherent (respectively pseudo-coherent).

Hints: first show that a left \(A\)-module \(L\) is of finite type if and only if \(A'\otimes_A L\) is of finite type (cf. Bourbaki, Commutative Algebra, Chap. I, §3, Prop. 11 and Chap. II, §3, Prop. 4), and then reduce to this case by truncation, using \hyperlink{sga6.e01.prop.2-5}{(2.5)} and (\hyperlink{sga6.e01.prop.2-10}{2.10} b).

\sgaStableTarget{sga6.e01.sec.3}\section*{3. Link with the classical notion of coherence}

\sgaStableTarget{sga6.e01.para.3-0}\textbf{3.0.} The hypotheses and notation are those of \hyperlink{sga6.e01.para.2-0}{(2.0)}, except that \(\mathcal C_0\) is no longer assumed locally quasi-stable under kernels of epimorphisms.

\sgaStableTarget{sga6.e01.def.3-1}\textbf{Definition 3.1.} Let \(F\in\Ob\mathcal C_S\). We shall say that \(F\) is \emph{precoherent} relative to \(\mathcal C_0\) if, for every \(U\) over \(S\) and every morphism
\[
  u:E\longrightarrow F|U,
\]
where \(E\in\Ob\mathcal C_{0,U}\), \(\Ker(u)\) is of \(\mathcal C_0\)-finite type \hyperlink{sga6.e01.sec.1.2}{(1.2)}. We shall say that \(F\) is \emph{coherent} relative to \(\mathcal C_0\) if \(F\) is precoherent and of finite type.

\emph{Technical note.} We introduce this notion here only as a technical intermediary; the only notion to retain is that of coherence.

\sgaStableTarget{sga6.e01.rem.3-2}\textbf{Remark 3.2.} Let \(F\in\Ob\mathcal C_S\). It is clear from the definitions that, if \(F\) is precoherent (respectively of finite type), then every subobject (respectively quotient) of \(F\) is precoherent (respectively of finite type).

\sgaStableTarget{sga6.e01.lem.3-3}\textbf{Lemma 3.3.}

\textbf{a)} Let
\[
  0\longrightarrow F\longrightarrow G\longrightarrow H
\]
be an exact sequence of \(\mathcal C_S\). If \(G\) is of finite type and \(H\) is precoherent, then \(F\) is of finite type.

\textbf{b)} Let
\[
  F\longrightarrow G\longrightarrow H\longrightarrow 0
\]
be an exact sequence of \(\mathcal C_S\). If \(F\) is of finite type and \(G\) is precoherent, then \(H\) is precoherent.

\textbf{c)} Let
\[
  0\longrightarrow F\longrightarrow G\longrightarrow H\longrightarrow 0
\]
be an exact sequence of \(\mathcal C_S\). If \(F\) and \(H\) are precoherent (respectively of finite type), then so is \(G\).

\emph{Proof.} a) After localizing, there is an epimorphism \(G'\to G\) with \(G'\in\Ob\mathcal C_0\). Hence there is a commutative diagram with exact rows
\[
\begin{tikzcd}
0\arrow[r]&F'\arrow[r]\arrow[d]&G'\arrow[r]\arrow[d]&H\arrow[d,"\Id"]\\
0\arrow[r]&F\arrow[r]&G\arrow[r]&H .
\end{tikzcd}
\]
Since \(H\) is precoherent, \(F'\) is of finite type; but \(F'\to F\) is an epimorphism, so \(F\) is of finite type.

b) Let \(U\) be over \(S\), and let \(H'\xrightarrow{f}H|U\) be a morphism of \(\mathcal C_U\), with \(H'\in\Ob\mathcal C_{0,U}\). After a change of notation, we may suppose \(U=S\). By the quasi-liftability of \(\mathcal C_0\) in \(\mathcal C\), there exists, after localizing, a commutative square
\[
\begin{tikzcd}
G' \arrow[r,"p"] \arrow[d] & H' \arrow[d,"f"]\\
G \arrow[r] & H,
\end{tikzcd}
\]
where the horizontal arrows are epimorphisms and \(G'\in\Ob\mathcal C_0\). Since \(\Ker(f)\) is a quotient of \(\Ker(fp)\), it is enough to show that \(\Ker(fp)\) is of finite type; hence we may suppose \(H'=G'\). The arrows \(F\to G\) and \(H'\to G\) define
\[
  g:F\oplus H'\longrightarrow G.
\]
Putting \(K=\Ker(G\to H)\), we have a commutative diagram, with exact rows and with \(q\) an epimorphism,
\[
\begin{tikzcd}
0\arrow[r]&F\arrow[r]\arrow[d,"q"]&F\oplus H'\arrow[r]\arrow[d,"g"]&H'\arrow[r]\arrow[d,"f"]&0\\
0\arrow[r]&K\arrow[r]&G\arrow[r]&H\arrow[r]&0 .
\end{tikzcd}
\]
The snake diagram shows that \(\Ker(g)\to\Ker(f)\) is an epimorphism. But, since \(F\oplus H'\) is of finite type and \(G\) is precoherent, \(\Ker(g)\) is of finite type by a); hence we are done.

c) First prove the non-parenthesized assertion. Up to a change of notation, it is a matter of showing that, if \(g:G'\to G\) is a morphism of \(\mathcal C_S\), with \(G'\in\Ob\mathcal C_0\), then \(\Ker(g)\) is of finite type. Let \(H'\) be the image of the composed arrow \(G'\to H\). We have a commutative diagram with exact rows
\[
\begin{tikzcd}
0\arrow[r]&F'\arrow[r]\arrow[d,"f"]&G'\arrow[r]\arrow[d,"g"]&H'\arrow[r]\arrow[d,"h"]&0\\
0\arrow[r]&F\arrow[r]&G\arrow[r]&H\arrow[r]&0 .
\end{tikzcd}
\]
Since \(h\) is a monomorphism, the snake diagram gives an isomorphism \(\Ker(f)\qis\Ker(g)\). But \(F'\) is of finite type by a), hence \(\Ker(f)\) is of finite type, again by a). This proves the non-parenthesized assertion. For the parenthesized assertion, strictly speaking one cannot invoke (\hyperlink{sga6.e01.prop.2-5}{2.5} b), since \(\mathcal C_0\) is not assumed locally quasi-stable under kernels of epimorphisms. After localizing, there are epimorphisms
\[
  F'\xrightarrow{u}F,
  \qquad H'\xrightarrow{v}H,
\]
with \(F',H'\in\Ob\mathcal C_0\). After localizing further, we may suppose that \(v\) factors through \(G\), whence a commutative diagram with exact rows
\[
\begin{tikzcd}
0\arrow[r]&F'\arrow[r]\arrow[d,"u"]&F'\oplus H'\arrow[r]\arrow[d,"w"]&H'\arrow[r]\arrow[d,"v"]&0\\
0\arrow[r]&F\arrow[r]&G\arrow[r]&H\arrow[r]&0 .
\end{tikzcd}
\]
Since \(u\) and \(v\) are epimorphisms, so is \(w\), and this proves the assertion. The proof of \hyperlink{sga6.e01.lem.3-3}{(3.3)} is complete.

It follows immediately that:

\sgaStableTarget{sga6.e01.cor.3-4}\textbf{Corollary 3.4.} The full (respectively strictly full) subcategory of \(\mathcal C_S\) formed by coherent objects is stable under finite projective (respectively inductive) limits. In particular, the kernel, cokernel, and image of an arrow between two coherent objects are coherent. Let
\[
  E^0\longrightarrow E^1\longrightarrow E^2\longrightarrow E^3\longrightarrow E^4
\]
be an exact sequence of \(\mathcal C_S\). If \(E^0,E^1,E^3,E^4\) are coherent, then \(E^2\) is coherent as well.

\sgaStableTarget{sga6.e01.cor.3-5}\textbf{Corollary 3.5.} Suppose that the objects of \(\mathcal C_0\) are \(\mathcal C_0\)-coherent (which in particular implies that \(\mathcal C_0\) is locally quasi-stable under kernels of epimorphisms, that is, that we are under the conditions of \hyperlink{sga6.e01.para.2-0}{(2.0)}).

\textbf{a)} For \(F\in\Ob\mathcal C_S\), the following conditions are equivalent:
\begin{enumerate}[label=(\roman*)]
\item \(F\) is of finite presentation;
\item \(F\) is coherent;
\item \(F\) is of finite \(\infty\)-presentation \hyperlink{sga6.e01.def.2-8}{(2.8)}.
\end{enumerate}

\textbf{b)} Let \(F\in\Ob\Dpx^-(\mathcal C_S)\), and let \(n\in\mathbb Z\). Then
\[
\begin{split}
F\text{ is }n\text{-pseudo-coherent}
\quad\Longleftrightarrow\quad&
\Ho^i(F)\text{ is coherent for }i\geq n+1,\\
&\text{and }\Ho^n(F)\text{ is of finite type.}
\end{split}
\]
In particular, for \(F\) to be pseudo-coherent it is necessary and sufficient that \(\Ho^n(F)\) be coherent for every \(n\).

\emph{Proof.} Assertion a) and the implication \(\Rightarrow\) in b) follow trivially from \hyperlink{sga6.e01.cor.3-4}{(3.4)}. The implication \(\Leftarrow\) in b) follows at once from a) and (\hyperlink{sga6.e01.prop.2-11}{2.11} b).

\sgaStableTarget{sga6.e01.ex.3-6}\textbf{Examples 3.6.} Let \((S,A)\) be a ringed topos, and let \(\mathcal C,\mathcal C_0\) be the \(S\)-categories of (\hyperlink{sga6.e01.sec.1.3}{1.3} d), also as in \hyperlink{sga6.e01.ex.2-15}{(2.15)}. In order that the objects of \(\mathcal C_0\) be coherent, it is necessary and sufficient that \(A\) be coherent, that is, precoherent. Here are examples where this is the case:
\begin{itemize}
\item \(A\) is constant with value a left noetherian ring (SGA 4 IX 2.1);
\item \(S\) is the topos of sheaves of sets, for the Zariski topology, on a locally noetherian scheme \(S\), and \(A=\mathcal O_S\);
\item \(S\) is the topos of sheaves of sets on a complex analytic space \(S\), and \(A=\mathcal O_S\).
\end{itemize}
On the other hand, the sheaf of continuous complex-valued functions on a topological space is not coherent in general.

Of course, if \((S,A)\) is defined by a ringed space, the notion of coherence considered here coincides with the usual notion of Serre (EGA \(0_I\), 5.3.1).

\sgaStableTarget{sga6.e01.prop.3-7}\textbf{Proposition 3.7.} \textbf{a)} Let
\[
  f:(S',A')\longrightarrow(S,A)
\]
be a morphism of ringed topoi. Assume that \(f\) is flat on the right, that is, that the functor \(f^*=A'\otimes_A-\) is exact. There is, for \(F\in\Ob\Dpx_A^-(S)\) and \(G\in\Ob\Dpx_A^+(S)\), a functorial canonical arrow
\[
  f^*\RHom(F,G)\longrightarrow \RHom(f^*(F),f^*(G)).
\]
This arrow is an isomorphism when \(F\) is pseudo-coherent.

\textbf{b)} Suppose \(A\) is commutative and coherent, and let \(F\in\Ob\Dpx^-(S)\), \(G\in\Ob\Dpx^+(S)\). If \(\Ho^i(F)\) and \(\Ho^i(G)\) are coherent for all \(i\), then the same is true of \(\Ext^i(F,G)\).

\emph{Proof.} a) To define the arrow (cf. SGA 4 XVII), one may suppose that \(G\) has injective components, and one takes the composite arrow
\[
  f^*\Hom(F,G)\longrightarrow \Hom(f^*F,f^*G)
  \longrightarrow \RHom(f^*F,f^*G).
\]
Let us show that it is an isomorphism for \(F\) pseudo-coherent. For brevity write \(T\) (respectively \(T'\)) for the functor \(f^*\RHom(\cdot,G)\) (respectively \(\RHom(f^*\cdot,f^*G)\)), and let \(a\in\mathbb Z\) be such that \(\Ho^i(G)=0\) for \(i<a\). It is clear that one has \(TA\qis T'A\), hence by dévissage \(TF\qis T'F\) for \(F\) strictly perfect \hyperlink{sga6.e01.def.2-1}{(2.1)}. On the other hand, if \(F\) is acyclic in degrees \(>n\), then \(TF\) and \(T'F\) are acyclic in degrees \(<a-n\). To verify that \(TF\qis T'F\) for \(F\) pseudo-coherent, the question being local, one may suppose \(F\) strictly \(n\)-pseudo-coherent; that is, there is a distinguished triangle
\[
  F'\longrightarrow F\longrightarrow F''\longrightarrow F'[1].
\]
with \(F'\) strictly perfect and \(F''\) acyclic in degrees \(>n-1\). From this one obtains a morphism of distinguished triangles
\[
\begin{tikzcd}[column sep=tiny]
TF'' \arrow[r] \arrow[d] & TF \arrow[r] \arrow[d] & TF' \arrow[d] \\
T'F'' \arrow[r] & T'F \arrow[r] & T'F' .
\end{tikzcd}
\]
Now \(TF'\qis T'F'\), since \(F'\) is strictly perfect, and \(TF''\) and \(T'F''\) are acyclic in degrees \(\leq a-n\). Writing the induced morphism on the exact cohomology sequences of the preceding triangles, one deduces that \(TF\to T'F\) induces an isomorphism
\[
  \Ho^i(TF)\qis \Ho^i(T'F)
\]
for \(i\leq a-n\). Since \(n\) may be chosen arbitrarily negative, this proves the assertion.

b) The proof is left as an exercise to the reader; see, for example, [H] II 3.3.

\sgaStableTarget{sga6.e01.rem.3-7-1}\textbf{Remark 3.7.1.} We shall see below \hyperlink{sga6.e01.prop.7-1-2}{(7.1.2)} another case of compatibility of \(\RHom\) with \(\mathrm L f^*\).

\sgaStableTarget{sga6.e01.cor.3-7-2}\textbf{Corollary 3.7.2.} (cf. EGA \(0_I\), 5.2.7.) Let \((S,A)\) be a ringed topos, let \(s\) be a point of \(S\) (SGA 4 VIII 7.8), and let
\[
  E,F\in\Ob\D^b_{A,\coh}(S).
\]
If \(E_s\qis F_s\), then there exists an \(s\)-pointed object \(U\) of \(S\) such that
\[
  E|U\qis F|U.
\]

\emph{Proof.} Let \(\alpha:E_s\to F_s\) and \(\beta:F_s\to E_s\) be reciprocal isomorphisms. By (\hyperlink{sga6.e01.prop.3-7}{2.22} a), one has
\[
  \Hom(E,F)_s\qis\Hom(E_s,F_s),
  \qquad
  \Hom(F,E)_s\qis\Hom(F_s,E_s).
\]
There are therefore a neighbourhood \(U\) of \(s\) and a section \(u\) (respectively \(v\)) of \(\Hom(E,F)\) (respectively of \(\Hom(F,E)\)) over \(U\) such that \(u_s=\alpha\) (respectively \(v_s=\beta\)). The morphisms \(vu:E|U\to E|U\) and \(uv:F|U\to F|U\) induce the identity at \(s\); hence on a suitable \(V\) over \(U\), as required.

\sgaStableTarget{sga6.e01.sec.4}\section*{4. Perfect complexes}

\sgaStableTarget{sga6.e01.para.4-0}\textbf{4.0.} In this section, \(S\) denotes a site, \(S\) an object of \(S\), \(\mathcal C\) a flat abelian \(S\)-category, and \(\mathcal C_0\) a strictly full additive sub-\(S\)-category of \(\mathcal C\) \hyperlink{sga6.e01.sec.1.1}{(1.1)}. Unless otherwise mentioned, we suppose that \(\mathcal C_0\) is locally liftable in \(\mathcal C\) and locally stable under kernels of epimorphisms \hyperlink{sga6.e01.sec.1.2}{(1.2)}.

\sgaStableTarget{sga6.e01.lem.4-1}\textbf{Lemma 4.1.} Let
\[
  u:E\longrightarrow F
\]
be an arrow of \(\Cpx(\mathcal C_S)\). If \(E\) is strictly \(\mathcal C_0\)-perfect \hyperlink{sga6.e01.def.2-1}{(2.1)} and \(F\) is acyclic, then \(u\) is locally homotopic to zero.

\emph{Proof.} One has to construct locally a homotopy \(k\) such that
\[
  u=dk+kd.
\]
Since \(E\) is strictly perfect, there are \(a,b\in\mathbb Z\), with \(a\leq b\), such that \(E^i=0\) for \(i\notin[a,b]\). Thus one may take \(k^i=0\) for \(i\notin[a,b]\). We construct the \(k^i\), for \(i>a\), by descending induction on \(i\). Suppose \(k^i:E^i\to F^{i-1}\) has been constructed so that
\[
  u^i=k^{i+1}d_E^i+d_F^{i-1}k^i
\]
for \(i>n>a\). We show that, after localizing, there exists
\[
  k^{n-1}:E^{n-1}\longrightarrow F^{n-2}
\]
such that
\[
  u^{n-1}=k^n d_E^{n-1}+d_F^{n-2}k^{n-1}.
\]
By the induction hypothesis, the arrow
\[
  v^{n-1}=u^{n-1}-k^n d_E^{n-1}:E^{n-1}\longrightarrow F^{n-1}
\]
factors through \(Z^{n-1}(F)\). But \(Z^{n-1}(F)=B^{n-1}(F)\), since \(F\) is acyclic. Since \(\mathcal C_0\) is locally liftable in \(\mathcal C\), it follows, after localizing, that there exists \(k^{n-1}:E^{n-1}\to F^{n-2}\) such that \(v^{n-1}=d_F^{n-2}k^{n-1}\), as required.

\sgaStableTarget{sga6.e01.cor.4-2}\textbf{Corollary 4.2.} Let
\[
\begin{tikzcd}
P \arrow[d,"v"'] \\
F & E \arrow[l,"f"]
\end{tikzcd}
\]
be a diagram of \(\Cpx(\mathcal C_S)\), where \(P\) is strictly \(\mathcal C_0\)-perfect and \(f\) is a quasi-isomorphism. The set of \(X\) over \(S\) for which there exists an arrow
\[
  u:P_X\longrightarrow E_X
\]
of \(\Cpx(\mathcal C_X)\) such that the diagram
\[
\begin{tikzcd}
& P_X \arrow[d,"v_X"]\\
E_X \arrow[r,"f_X"'] \arrow[ur,"u"] & F_X
\end{tikzcd}
\]
is commutative in \(\Kpx(\mathcal C_X)\), is a refinement of \(S\).

\emph{Proof.} Let \(G\) be the cone of \(f\). There is an exact sequence
\[
  \longrightarrow \Hom_{\Kpx(\mathcal C_S)}(P,E)
  \longrightarrow \Hom_{\Kpx(\mathcal C_S)}(P,F)
  \longrightarrow \Hom_{\Kpx(\mathcal C_S)}(P,G)
  \longrightarrow .
\]
Since \(G\) is acyclic, the image of \(v\) in \(\Hom_{\Kpx(\mathcal C_S)}(P,G)\) is locally zero by \hyperlink{sga6.e01.lem.4-1}{(4.1)}. Therefore, after localizing, there exists \(u:P\to E\) such that \(v=fu\) in \(\Kpx(\mathcal C_S)\), as required.

\sgaStableTarget{sga6.e01.cor.4-3}\textbf{Corollary 4.3.} Let \(f:E\to F\) be a quasi-isomorphism of \(\Cpx(\mathcal C_S)\), with \(F\) strictly perfect relative to \(\mathcal C_0\). The set of \(X\) over \(S\) for which there exists a quasi-isomorphism
\[
  g:F_X\longrightarrow E_X
\]
such that \(f_Xg\) is homotopic to \(\Id_{F_X}\), is a refinement of \(S\).

\emph{Proof.} This is a special case of \hyperlink{sga6.e01.cor.4-2}{(4.2)}.

\sgaStableTarget{sga6.e01.schol.4-4}\textbf{Scholium 4.4.} Let \(F\in\Ob\Kpx(\mathcal C_S)\). One may define, for example with the aid of a cleavage of \(\mathcal C\) (SGA 1 VI 7.1), a category \((\Qis/F)\) fibered over \(S/S\), whose fiber over \(X\) over \(S\) is the category \((\Qis/F_X)^\circ\), opposite to the category of quasi-isomorphisms with target \(F_X\). One may then express \hyperlink{sga6.e01.cor.4-3}{(4.3)} by saying that, if \(F\) is strictly perfect, the arrow \(\Id_F\) is locally cofinal in \((\Qis/F)\).

\sgaStableTarget{sga6.e01.notation.4-5}\textbf{Notation 4.5.} Let \(\mathcal M\) be a fibered \(S\)-category. Let \(X\in\Ob S\), and let \(E,F\in\Ob\mathcal M_X\). We shall denote by
\[
  \underline{\Hom}_{\mathcal M}(E,F)
\]
the sheaf on \(S/X\) associated to the presheaf of \(S\)-morphisms from \(E\) to \(F\) in \(\mathcal M\). For the definition of this presheaf, see [G] I 2.6. One must be careful that, even if \(\mathcal C\) is a stack ([G]), the corresponding presheaf for \(\mathcal M=\Kpx(\mathcal C)\) or \(\Dpx(\mathcal C)\) is not in general a sheaf.

\sgaStableTarget{sga6.e01.cor.4-6}\textbf{Corollary 4.6.} Let \(E,F\in\Ob\Cpx(\mathcal C_S)\), with \(E\) strictly perfect. The canonical arrow
\[
  \underline{\Hom}_{\Kpx(\mathcal C)}(E,F)
  \longrightarrow
  \underline{\Hom}_{\Dpx(\mathcal C)}(E,F)
\]
is an isomorphism.

\emph{Proof.} This is an immediate consequence of \hyperlink{sga6.e01.schol.4-4}{(4.4)} and of the formula
\[
  \underline{\Hom}_{\Dpx(\mathcal C_X)}(E_X,F_X)
  = \lim_{(\Qis/E_X)^\circ}\underline{\Hom}_{\Kpx(\mathcal C_X)}(\,\cdot\,,F_X)
\]
for \(X\) over \(S\).

\sgaStableTarget{sga6.e01.def.4-7}\textbf{Definition 4.7.} Let \(F\in\Ob\Dpx(\mathcal C_S)\), and let \(a,b\in\mathbb Z\), with \(a\leq b\). We shall say that \(F\) has \(\mathcal C_0\)-perfect amplitude contained in the interval \([a,b]\), and we shall write
\[
  \mathcal C_0\text{-}\Parf\text{-}\amp(F)\subset[a,b],
\]
if the set of \(X\) over \(S\) for which there exists a quasi-isomorphism
\[
  F'\longrightarrow F
\]
in \(\Cpx(\mathcal C_X)\), with \(F'\) strictly \(\mathcal C_0\)-perfect (2.1) and such that \(F'^i=0\) for \(i\notin[a,b]\), is a refinement of \(S\). By \hyperlink{sga6.e01.cor.4-3}{(4.3)}, this condition is equivalent to the one obtained by replacing ``quasi-isomorphism in \(\Cpx(\mathcal C_X)\)'' by ``isomorphism in \(\Dpx(\mathcal C_X)\).'' We shall say that \(F\) has finite \(\mathcal C_0\)-perfect amplitude if there are integers \(a\) and \(b\) such that \(\mathcal C_0\)-\(\Parf\)-\(\amp(F)\subset[a,b]\). We shall say that \(F\) is \(\mathcal C_0\)-perfect if \(F\) is locally of finite \(\mathcal C_0\)-perfect amplitude.

\sgaStableTarget{sga6.e01.ex.4-8}\textbf{Example 4.8.} Let \((S,A)\) be a ringed topos, let \(\mathcal C\) be the \(S\)-category defined by
\[
  U\longmapsto \text{the category of left }A_U\text{-modules},
\]
and let \(\mathcal C_0\) be the sub-\(S\)-category defined by
\[
  U\longmapsto \text{the category of left }A_U\text{-modules which are direct summands of free finite-type modules}.
\]
The pair \((\mathcal C,\mathcal C_0)\) satisfies the hypotheses \hyperlink{sga6.e01.para.4-0}{(4.0)} by \hyperlink{sga6.e01.lem.1-3-1}{(1.3.1)}. An object \(F\) of \(\Dpx_A(S)\) (cf. \hyperlink{sga6.e01.ex.2-15}{(2.15)}) will be said to have perfect amplitude contained in \([a,b]\) (respectively finite perfect amplitude, respectively to be perfect) if \(F\) has \(\mathcal C_0\)-perfect amplitude contained in \([a,b]\) (respectively the corresponding property). When the direct summands of free finite-type modules are locally free (cf. \hyperlink{sga6.e01.lem.2-15-1}{(2.15.1)}), to say that \(F\) is perfect means that \(F\) is locally isomorphic, in \(\Dpx(S)\), to a bounded complex whose components are free of finite type.

\sgaStableTarget{sga6.e01.notation.4-9}\textbf{Notation 4.9.} We shall denote by \(\Dpx(\mathcal C)_{\mathcal C_0\text{-}\Parf}\), or \(\Parf(\mathcal C,\mathcal C_0)\), the strictly full sub-\(S\)-category of \(\Dpx(\mathcal C)\) formed by the \(\mathcal C_0\)-perfect objects. In the case of \hyperlink{sga6.e01.ex.4-8}{Example (4.8)}, we shall write \(\Dpx_A(S)_{\Parf}\), or \(\Parf_A(S)\), or simply \(\Parf(S)\) if no ambiguity about \(A\) is possible.

\sgaStableTarget{sga6.e01.prop.4-10}\textbf{Proposition 4.10.} The category \(\Parf(\mathcal C,\mathcal C_0)\) is a triangulated sub-\(S\)-category of \(\Dpx(\mathcal C)\) \hyperlink{sga6.e01.para.1-1-4}{(1.1.4)}. The canonical functor
\[
  \Kpx^b(\mathcal C_0)\longrightarrow \Parf(\mathcal C,\mathcal C_0)
\]
(respectively
\[
  \Tr(\Kpx^b(\mathcal C_0))\longrightarrow \Tr(\Parf(\mathcal C,\mathcal C_0))
\]
\hyperlink{sga6.e01.para.1-1-3}{(1.1.3)}) induces an \(S\)-equivalence on the associated stacks ([G] II 2.2.4), which means:

\begin{enumerate}[label=(\roman*)]
\item for every \(X\in\Ob S\) and all \(E,F\in\Ob\Kpx^b(\mathcal C_{0,X})\) (respectively \(\Tr(\Kpx^b(\mathcal C_{0,X}))\)), the canonical arrow
\[
  \Hom_{\Kpx^b(\mathcal C_0)}(E,F)
  \longrightarrow
  \Hom_{\Dpx(\mathcal C)}(E,F)
\]
(respectively
\[
  \Hom_{\Tr(\Kpx^b(\mathcal C_0))}(E,F)
  \longrightarrow
  \Hom_{\Tr(\Dpx(\mathcal C))}(E,F)
\]) is an isomorphism; and
\item every object of \(\Parf(\mathcal C,\mathcal C_0)\) (respectively of \(\Tr(\Parf(\mathcal C,\mathcal C_0))\)) is locally isomorphic in \(\Dpx(\mathcal C)\) (respectively in \(\Tr(\Dpx(\mathcal C))\)) to an object of \(\Kpx^b(\mathcal C_0)\) (respectively of \(\Tr(\Kpx^b(\mathcal C_0))\)).
\end{enumerate}

\emph{Proof.} The non-parenthesized assertion (i) is a special case of \hyperlink{sga6.e01.cor.4-6}{(4.6)}, and the non-parenthesized assertion (ii) is true by the \hyperlink{sga6.e01.def.4-7}{definition (4.7)} and \hyperlink{sga6.e01.notation.4-9}{(4.9)}. It is clear, moreover, that \(\Parf(\mathcal C,\mathcal C_0)\) is stable under translation of degrees. To prove that \(\Parf(\mathcal C,\mathcal C_0)\) is a triangulated sub-\(S\)-category of \(\Dpx(\mathcal C)\), it remains to show that, if
\sgaStableTarget{sga6.e01.eq.unnamed.occ3}\[
  E\xrightarrow{u}F\longrightarrow G\longrightarrow E[1].
  \tag{*}
\]
is a distinguished triangle of \(\Dpx(\mathcal C)\) with \(E\) and \(F\) perfect, then \(G\) is perfect. After localizing, we may suppose that \(E\) and \(F\) are objects of \(\Kpx^b(\mathcal C_0)\). After localizing further, we may suppose, by the non-parenthesized assertion (i), that \(u\) is an arrow of \(\Kpx^b(\mathcal C_0)\); then \(G\) is perfect, since it is isomorphic to the cone of \(u\), which is strictly perfect. Hence \(\Parf(\mathcal C,\mathcal C_0)\) is indeed a triangulated sub-\(S\)-category of \(\Dpx(\mathcal C)\). The preceding argument also proves the parenthesized assertion (ii). It remains to prove the parenthesized assertion (i). Let \(X\in\Ob S\), and let
\[
E=(E_1\xrightarrow{a_1}E_2\xrightarrow{a_2}E_3\xrightarrow{a_3}E_1[+1]),
\quad
F=(F_1\xrightarrow{b_1}F_2\xrightarrow{b_2}F_3\xrightarrow{b_3}F_1[+1])
\]
be distinguished triangles of \(\Kpx^b(\mathcal C_{0,X})\). We must show that
\sgaStableTarget{sga6.e01.eq.unnamed.occ4}\[
  \Hom_{\Tr(\Kpx^b(\mathcal C_0))}(E,F)
  \longrightarrow
  \Hom_{\Tr(\Dpx(\mathcal C))}(E,F)
  \tag{**}
\]
is an isomorphism. By the non-parenthesized assertion (i), we already know that \((**)\) is injective. Let \((u_i):E\to F\) be a morphism of triangles in \(\Dpx(\mathcal C)\). After localizing, the non-parenthesized assertion (i) permits us to suppose that \(u_i\), for \(1\leq i\leq3\), is an arrow of \(\Kpx^b(\mathcal C_0)\). The arrow
\[
  u_{i+1}a_i-b_iu_i,
\]
for \(1\leq i\leq3\) (with the convention \(u_4=u_1\)), is zero in \(\Dpx(\mathcal C)\), hence locally zero in \(\Kpx^b(\mathcal C_0)\) by the non-parenthesized assertion (i). In other words, \((u_i)\) locally defines a morphism of triangles in \(\Kpx^b(\mathcal C_0)\). This proves that \((**)\) is surjective and completes the proof of \hyperlink{sga6.e01.prop.4-10}{(4.10)}.

\sgaStableTarget{sga6.e01.rem.4-10-1}\textbf{Remark 4.10.1.} In \hyperlink{sga6.e01.prop.4-10}{(4.10)}, one may replace \(\Parf(\mathcal C,\mathcal C_0)\) by
\[
  \Dpx^*(\mathcal C)_{\mathcal C_0\text{-}\Parf}
  =\Dpx^*(\mathcal C)\cap\Parf(\mathcal C,\mathcal C_0)
  \qquad(*=-\text{ or }b),
\]
and nothing changes in the proof.

\sgaStableTarget{sga6.e01.complement.4-11}\textbf{Complement 4.11.} The triangulated structure of \(\Parf(\mathcal C,\mathcal C_0)\) can be made more precise in terms of perfect amplitude \hyperlink{sga6.e01.def.4-7}{(4.7)}. If \(E\in\Ob\Dpx(\mathcal C)\) has perfect amplitude contained in \([a,b]\), then \(E[n]\) has perfect amplitude contained in \([a-n,b-n]\).

Let
\[
  E\longrightarrow F\longrightarrow G\longrightarrow E[1].
\]
be a distinguished triangle of \(\Dpx(\mathcal C)\). One may draw up a table analogous to \hyperlink{sga6.e01.schol.2-6}{(2.6)}, where the entries indicate perfect amplitudes and each row corresponds to an implication whose target is underlined:
\[
\begin{array}{ccc}
E&F&G\\[0.3em]
{}[a+1,b+1]&{}[a,b]&\underline{[a,b]}\\
{}[a,b]&\underline{[a,b]}&{}[a,b]\\
\underline{[a,b]}&{}[a,b]&{}[a-1,b-1].
\end{array}
\]
From \hyperlink{sga6.e01.prop.4-10}{(4.10)} one immediately deduces:

\sgaStableTarget{sga6.e01.cor.4-12}\textbf{Corollary 4.12.} Suppose that \(S\) is a punctual site. For \(F\in\Ob\Dpx(\mathcal C)\) to have perfect amplitude relative to \(\mathcal C_0\) contained in \([a,b]\), it is necessary and sufficient that there exist a quasi-isomorphism \(F'\to F\), with \(F'\) strictly perfect and such that \(F'^i=0\) for \(i\notin[a,b]\). The canonical functors
\[
  \Kpx^b(\mathcal C_0)\longrightarrow\Parf(\mathcal C,\mathcal C_0)
\]
and
\[
  \Tr(\Kpx^b(\mathcal C_0))\longrightarrow\Tr(\Parf(\mathcal C,\mathcal C_0))
\]
are equivalences of categories.

\sgaStableTarget{sga6.e01.rem.4-12-1}\textbf{Remark 4.12.1.} When \(S\) is a punctual site, the objects of \(\mathcal C_0\) are projective. When \(\mathcal C_0\) contains all the projective objects of \(\mathcal C\), one says ``projective amplitude'' instead of ``\(\mathcal C_0\)-perfect amplitude.''



% Additional notation for Expose I continuation
\providecommand{\Cpx}{\mathrm C}
\providecommand{\Kpx}{\mathrm K}
\providecommand{\Dpx}{\mathrm D}
\providecommand{\Ho}{\mathrm H}
\providecommand{\RHom}{\mathrm R\!\operatorname{Hom}}
\providecommand{\Ob}{\operatorname{ob}}
\providecommand{\qis}{\xrightarrow{\sim}}
\providecommand{\RHom}{\mathrm R\!\Hom}
\providecommand{\Ext}{\operatorname{Ext}}
\providecommand{\Hom}{\operatorname{Hom}}
\providecommand{\Tr}{\operatorname{Tr}}

\providecommand{\Add}{\operatorname{Add}}
\providecommand{\Cart}{\operatorname{Cart}}
\providecommand{\perfamp}{\operatorname{perf.amp}}
\providecommand{\perfdim}{\operatorname{perf.dim}}
\providecommand{\toramp}{\operatorname{tor.amp}}
\providecommand{\torf}{\operatorname{torf}}
\providecommand{\injf}{\operatorname{injf}}
\providecommand{\rg}{\operatorname{rg}}
\providecommand{\is}{\mathrm{is}}
\providecommand{\Mod}{\operatorname{Mod}}


\sgaStableTarget{sga6.e01.expose.occ3}\section*{Exposé I, continued: perfect amplitude, finite Tor-dimension, and rank}

\sgaStableTarget{sga6.e01.rem.4-12-2}\textbf{Remark 4.12.2.} Let \(\mathcal C\) be an abelian category (which may be regarded as fibered over a punctual site), and let \(\mathcal C_0\) be a strictly full additive subcategory. Suppose that \(\mathcal C_0\) is stable under kernels of epimorphisms and is quasi-liftable in \(\mathcal C\). Then an acyclic object of \(\Kpx^b(\mathcal C_0)\) is not in general zero; hence there is no hope that the functor
\[
  \Kpx^b(\mathcal C_0)\longrightarrow \Dpx(\mathcal C)
\]
should be faithful. However, if \(\Kpx^{b,\mathrm{ac}}(\mathcal C_0)\) denotes the thick subcategory of \(\Kpx^b(\mathcal C_0)\) formed by the objects which are acyclic in \(\mathcal C\), the canonical functor
\[
  \Kpx^b(\mathcal C_0)/\Kpx^{b,\mathrm{ac}}(\mathcal C_0)
  \longrightarrow \Dpx(\mathcal C)
\]
is fully faithful.

Indications: one already knows from \hyperlink{sga6.e01.prop.2-7}{(2.7)} that the functor
\[
  \Kpx^-(\mathcal C_0)/\Kpx^{-,\mathrm{ac}}(\mathcal C_0)
  \longrightarrow \Dpx^-(\mathcal C)
\]
is fully faithful. Show, for example with the aid of the criterion of [V], 4.2 b)(ii), that the functor
\[
  \Kpx^b(\mathcal C_0)/\Kpx^{b,\mathrm{ac}}(\mathcal C_0)
  \longrightarrow
  \Kpx^-(\mathcal C_0)/\Kpx^{-,\mathrm{ac}}(\mathcal C_0)
\]
is fully faithful.

\sgaStableTarget{sga6.e01.lem.4-13}\textbf{Lemma 4.13.} Let \(F\in\Ob \Dpx(\mathcal C_S)\), and let \(a,b,c\in\mathbb Z\) be such that \(a\leq b\) and \(a\leq c\). If \(F\) is of perfect amplitude contained in \([a,c]\) \hyperlink{sga6.e01.def.4-7}{(4.7)}, and if \(\Ho^i(F)=0\) for \(i>b\), then \(F\) is of perfect amplitude contained in \([a,d]\), where \(d=\inf(b,c)\).

\emph{Proof.} We may suppose \(b\leq c\), and \(F\) strictly perfect with \(F^i=0\) for \(i\notin [a,c]\). Then \(F\) is isomorphic to its truncation \(F'\) defined by
\[
  F'^i=F^i\quad (i<b),\qquad F'^b=Z^b(F),\qquad F'^i=0\quad (i>b).
\]
But \(Z^b(F)\) admits a finite right resolution by objects of \(\mathcal C_0\), hence is locally an object of \(\mathcal C_0\), whence the assertion.

\sgaStableTarget{sga6.e01.cor.4-14}\textbf{Corollary 4.14.} Let \(F\in\Ob\mathcal C_S\) and \(n\in\mathbb N\). The following conditions are equivalent:
\begin{enumerate}[label=(\roman*)]
\item \(\perfamp(F)\subset[-n,0]\);
\item \(F\) locally admits a left resolution of length \(n\) by objects of \(\mathcal C_0\), that is, the set of objects \(X\) over \(S\) such that there exists an exact sequence
\[
  0\longrightarrow L^{-n}\longrightarrow L^{-n+1}\longrightarrow \cdots
  \longrightarrow L^0\longrightarrow F_X\longrightarrow 0,
\]
with \(L^i\in\Ob\mathcal C_{0,X}\) for all \(i\), is a refinement of \(S\).
\end{enumerate}

\emph{Proof.} This is an immediate consequence of \hyperlink{sga6.e01.def.4-7}{(4.7)} and \hyperlink{sga6.e01.lem.4-13}{(4.13)}.

\sgaStableTarget{sga6.e01.def.4-14-1}\textbf{Definition 4.14.1.} We shall say that \(F\) is of perfect dimension \(\leq n\), and shall write \(\perfdim(F)\leq n\), if \(F\) satisfies the equivalent conditions of \hyperlink{sga6.e01.cor.4-14}{(4.14)}. We shall call the perfect dimension of \(F\) the smallest integer \(n\) such that \(\perfdim(F)\leq n\); if there is no such integer, we shall say that \(F\) is of infinite perfect dimension.

When \(S\) is the punctual site and \(\mathcal C_0\) is the full subcategory of \(\mathcal C\) formed by the projective objects, the perfect dimension is just the usual projective dimension.

\sgaStableTarget{sga6.e01.prop.4-15}\textbf{Proposition 4.15.} Let \(a,b,n\in\mathbb Z\), with \(a\leq b\) and \(n\geq0\), and let \(F\in\Ob\Dpx(\mathcal C)\) be such that \(F^i=0\) (respectively \(\Ho^i(F)=0\)) for \(i\notin[a,b]\). If, for every \(i\in[a,b]\), \(F^i\) (respectively \(\Ho^i(F)\)) is of perfect dimension \(\leq n+(i-a)\), then \(F\) is of perfect amplitude contained in \([a-n,b]\). Consequently, if for every \(i\in[a,b]\), \(F^i\) (respectively \(\Ho^i(F)\)) is perfect, then \(F\) is perfect.

\emph{Proof.} Proceed by induction on \(b\) (with \(a\) and \(n\) fixed), by truncating \(F\) and applying \hyperlink{sga6.e01.complement.4-11}{(4.11)}, cf. the proof of \hyperlink{sga6.e01.prop.2-11}{(2.11)}.

\sgaStableTarget{sga6.e01.lem.4-16}\textbf{Lemma 4.16.} Let \(a,b\in\mathbb Z\), \(a\leq b\), and let \(F\in\Ob\Dpx(\mathcal C)\) be such that \(F^i=0\) for \(i\notin[a,b]\). Suppose that \(\perfamp(F)\subset[a,b]\), and that \(F^i\in\Ob\mathcal C_0\) for \(i>a\). Then \(F^a\) is locally an object of \(\mathcal C_0\).

\emph{Proof.} There is an exact sequence, obtained by truncating \(F\),
\[
  0\longrightarrow F'\longrightarrow F\longrightarrow F^a[-a]\longrightarrow 0,
\]
with \(\perfamp(F')\subset[a+1,b+1]\). From \hyperlink{sga6.e01.complement.4-11}{(4.11)} and \hyperlink{sga6.e01.lem.4-13}{(4.13)} one deduces that \(\perfdim(F^a)=0\), i.e. that \(F^a\) is locally an object of \(\mathcal C_0\).

\sgaStableTarget{sga6.e01.prop.4-17}\textbf{Proposition 4.17.} Let \(F',F''\in\Ob\Dpx(\mathcal C_S)\), and let \(a,b\in\mathbb Z\), with \(a\leq b\). Consider the conditions:
\begin{enumerate}[label=(\roman*)]
\item \(\perfamp(F')\subset[a,b]\) and \(\perfamp(F'')\subset[a,b]\);
\item \(\perfamp(F'\oplus F'')\subset[a,b]\).
\end{enumerate}
Then (i) implies (ii). If \(\mathcal C_0\) is locally stable under direct factors, i.e. if the relation \(E'\oplus E''\in\Ob\mathcal C_0\), for \(E',E''\in\Ob\mathcal C_X\) and \(X\in\Ob S\), implies that locally \(E'\) and \(E''\) are objects of \(\mathcal C_0\), then (i) and (ii) are equivalent.

\emph{Proof.} The implication (i) \(\Rightarrow\) (ii) is a particular case of \hyperlink{sga6.e01.complement.4-11}{(4.11)}. Let us prove the converse. Put \(F=F'\oplus F''\). Since \(F\) is pseudo-coherent, \(F'\) and \(F''\) are also pseudo-coherent \hyperlink{sga6.e01.prop.2-12}{(2.12)}. On the other hand, the relation \(\Ho^i(F)=0\) for \(i\notin[a,b]\) implies \(\Ho^i(F')=\Ho^i(F'')=0\) for \(i\notin[a,b]\). By (\hyperlink{sga6.e01.prop.2-10}{2.10} a), we may therefore suppose, after localizing, that \(F'^i\) and \(F''^i\) are in \(\Ob\mathcal C_0\) for \(i\geq a+1\), and that \(F'^i=F''^i=0\) for \(i>b\). Then, by truncating, we get \(F'^i=F''^i=0\) for \(i<a\). Since \(\perfamp(F)\subset[a,b]\), it follows from \hyperlink{sga6.e01.lem.4-16}{(4.16)} that \(F^a\) is locally an object of \(\mathcal C_0\). But \(F^a=F'^a\oplus F''^a\), hence \(F'^a\) and \(F''^a\) are locally objects of \(\mathcal C_0\), since \(\mathcal C_0\) is locally stable under direct factors.

\sgaStableTarget{sga6.e01.prop.4-18}\textbf{Proposition 4.18.} Let \((S',\mathcal C',\mathcal C'_0)\) be a triple satisfying the same hypotheses \hyperlink{sga6.e01.para.4-0}{(4.0)} as \((S,\mathcal C,\mathcal C_0)\). Let \(f:S'\to S\) be a morphism of sites, defined by \(f^{-1}:S\to S'\), and let
\[
  u:\Dpx^*(\mathcal C)\longrightarrow \Dpx^*(\mathcal C')\qquad (*=-,+,b)
\]
be a functor over \(f\) (cf. \hyperlink{sga6.e01.prop.2-13}{(2.13)}). If \(u\) sends objects of \(\mathcal C_0\) to \(\mathcal C'_0\)-perfect complexes, then \(u\) defines a functor
\[
  \Dpx^*(\mathcal C)_{\mathcal C_0\text{-}\mathrm{parf}}
  \longrightarrow
  \Dpx^*(\mathcal C')_{\mathcal C'_0\text{-}\mathrm{parf}}.
\]

\emph{Proof.} Let \(F\in\Ob\Dpx^*(\mathcal C)_{\mathcal C_0\text{-}\mathrm{parf}}\). We show that \(u(F)\in\Ob\Dpx^*(\mathcal C')_{\mathcal C'_0\text{-}\mathrm{parf}}\). Since the question is local, one may suppose \(F\) strictly perfect. It remains only to use dévissage, taking into account the fact that \(\Dpx^*(\mathcal C')_{\mathcal C'_0\text{-}\mathrm{parf}}\) is a triangulated subcategory of \(\Dpx^*(\mathcal C')\) \hyperlink{sga6.e01.prop.4-10}{(4.10)}.

\sgaStableTarget{sga6.e01.rem.4-18-1}\textbf{Remark 4.18.1.} Suppose that there exist \(m,n\in\mathbb Z\), \(m\leq n\), such that \(F\in\Ob\mathcal C_0\) implies
\[
  \mathcal C'_0\text{-}\perfamp(u(F))\subset[m,n].
\]
Then, if \(F\in\Ob\Dpx^*(\mathcal C)\), one has
\[
  \mathcal C_0\text{-}\perfamp(F)
  \subset [a,b]
  \quad\Rightarrow\quad
  \mathcal C'_0\text{-}\perfamp(u(F))
  \subset [a+m,b+n].
\]
The proof is left as an exercise to the reader (use \hyperlink{sga6.e01.complement.4-11}{(4.11)}).

\sgaStableTarget{sga6.e01.cor.4-19}\textbf{Corollary 4.19.} Under the hypotheses of \hyperlink{sga6.e01.cor.2-16}{(2.16)}, if \(E'\in\Ob\Dpx^-(A,S')\) is perfect as an object of \(\Dpx^-(A,S')\), and \(E\in\Ob\Dpx^-(A,S)\) is perfect, then
\[
  E'\overset{\mathrm L}{\otimes}_A E
\]
is perfect.

\emph{Proof.} We may suppose that \(E'\in\Ob\Dpx^-(A,S'_A)\), where \(S'\) is the final object of \(S'\), and apply \hyperlink{sga6.e01.prop.4-18}{(4.18)} to
\[
  u=E'\overset{\mathrm L}{\otimes}_A:
  \Dpx^-(A,S)\longrightarrow \Dpx^-(A,S').
\]
In particular:

\sgaStableTarget{sga6.e01.cor.4-19-1}\textbf{Corollary 4.19.1.} \textbf{a)} The functor
\[
  \mathrm L f^*:\Dpx^-(A,S)\longrightarrow \Dpx^-(A,S')
\]
induces a functor
\[
  \mathrm L f^*:\Dpx^-(A,S)_{\mathrm{parf}}
  \longrightarrow
  \Dpx^-(A,S')_{\mathrm{parf}}.
\]

\textbf{b)} If \(A\) is commutative, the functor
\[
  \overset{\mathrm L}{\otimes}_A:
  \Dpx^-(S)\times_S\Dpx^-(S)
  \longrightarrow \Dpx^-(S)
\]
induces a functor
\[
  \overset{\mathrm L}{\otimes}_A:
  \Dpx^-(S)_{\mathrm{parf}}\times_S\Dpx^-(S)_{\mathrm{parf}}
  \longrightarrow \Dpx^-(S)_{\mathrm{parf}}.
\]

\sgaStableTarget{sga6.e01.rem.4-19-2}\textbf{Remark 4.19.2.} One can refine \hyperlink{sga6.e01.cor.4-19}{(4.19)} in terms of perfect amplitude:
\[
  \perfamp(E)\subset[a,b],\qquad
  \perfamp(E')\subset[a',b']
  \quad\Rightarrow\quad
  \perfamp(E'\overset{\mathrm L}{\otimes}_A E)\subset[a+a',b+b'].
\]
In particular,
\[
  \perfamp(E)\subset[a,b]
  \quad\Rightarrow\quad
  \perfamp(\mathrm L f^*(E))\subset[a,b].
\]

\sgaStableTarget{sga6.e01.sec.5}\section*{5. Finite Tor-dimension and perfection}

\sgaStableTarget{sga6.e01.para.5-0}\textbf{5.0.} The reader will surely have asked the question: what is missing from a pseudo-coherent complex \(E\) in order that it be perfect? Is it enough, for example, that \(E\) have bounded cohomology? The answer is no. Indeed, take for \(\mathcal C\) the category of modules over a commutative noetherian ring \(A\), and for \(\mathcal C_0\) the full subcategory of projective \(A\)-modules of finite type. Any \(A\)-module of finite type \(E\) is pseudo-coherent and has bounded cohomology, but \(E\) is perfect only if, in addition, it is of finite projective dimension (or, equivalently, of finite Tor-dimension). This example suggests that the difference between a pseudo-coherent complex and a perfect complex lies in a question of cohomological dimension. This is what we shall now examine in the case of the category of modules of a ringed topos (cf. \hyperlink{sga6.e01.lem.3-3}{(3.3)}). We shall need a few preliminary sorties on the notion of finite Tor-dimension. We shall be brief, since this question already appears in the literature: cf. [H] II 4 and SGA 4 XVII 4.1.9.

\sgaStableTarget{sga6.e01.prop.5-1}\textbf{Proposition 5.1.} Let \((S,A)\) be a ringed topos, \(X\in\Ob S\), \(F\in\Ob\Dpx^-(A,X)\), and let \(m,n\in\mathbb Z\), with \(m\leq n\). The following conditions are equivalent:
\begin{enumerate}[label=(\roman*)]
\item \(F\) is isomorphic, in \(\Dpx^-(A,X)\), to a complex \(F'\) such that \(F'^i\) is flat for every \(i\), and zero for \(i\notin[m,n]\);
\item for every right \(A_X\)-module \(E\), one has
\[
  \Ho^i(E\overset{\mathrm L}{\otimes}_A F)=0\quad\text{for }i\notin[m,n].
\]
\end{enumerate}

\emph{Proof.} The implication (i) \(\Rightarrow\) (ii) is trivial; the implication (ii) \(\Rightarrow\) (i) follows, by truncation, from the following lemma.

\sgaStableTarget{sga6.e01.lem.5-1-1}\textbf{Lemma 5.1.1.} (cf. \hyperlink{sga6.e01.lem.4-16}{(4.16)}). Let \(F\in\Ob\Dpx^-(A,X)\). Suppose that \(F\) satisfies hypothesis (ii) of (5.1), that \(F^i=0\) for \(i\notin[m,n]\), and that \(F^i\) is flat for \(i\ne m\). Then \(F^m\) is flat.

\emph{Proof.} There is a distinguished triangle
\[
\begin{tikzcd}
F' \arrow[rr] && F \arrow[dl]\\
& F^m[-m] \arrow[ul,"+1"] &
\end{tikzcd}
\]
where \(F'\) has flat components and is zero outside the interval \([m+1,n]\) (suppose \(m<n\)). Apply the functor \(E\overset{\mathrm L}{\otimes}_A\cdot\), for an arbitrary right \(A_X\)-module \(E\), and write the exact cohomology sequence.

\sgaStableTarget{sga6.e01.def.5-2}\textbf{Definition 5.2.} An object \(F\) of \(\Dpx^-(A,S)\) satisfying the equivalent conditions of \hyperlink{sga6.e01.prop.5-1}{(5.1)} will be said to be of flat amplitude contained in the interval \([m,n]\), and we shall write
\[
  \toramp(F)\subset[m,n].
\]
We shall say that \(F\) is of finite flat amplitude, or of finite Tor-dimension, if there exist integers \(m\) and \(n\) such that the preceding condition is satisfied.

\sgaStableTarget{sga6.e01.para.5-3}\textbf{5.3.} Let \(X\in\Ob S\). We shall write \(\Dpx^-(A,X)_{\torf}\) for the full subcategory of \(\Dpx^-(A,X)\) formed by the objects of finite Tor-dimension. By criterion (ii) of \hyperlink{sga6.e01.prop.5-1}{(5.1)}, this is a triangulated subcategory of \(\Dpx^-(A,X)\). As \(X\) varies, the categories \(\Dpx^-(A,X)_{\torf}\) form a triangulated sub-\(S\)-category \hyperlink{sga6.e01.para.1-1-4}{(1.1.4)} of \(\Dpx^-(A,S)\), which we shall denote by \(\Dpx^-(S)_{A\text{-}\torf}\). The functor \(\overset{\mathrm L}{\otimes}_A\) induces an exact \(S\)-functor
\[
  \overset{\mathrm L}{\otimes}_A:
  \Dpx^b(S_A)\times_S \Dpx(A/S)_{\torf}
  \longrightarrow \Dpx^b(S_{\mathbb Z}).
\]

\sgaStableTarget{sga6.e01.para.5-4}\textbf{5.4.} In order that a left \(A\)-module \(F\) be of flat amplitude contained in \([-n,0]\), \(n\in\mathbb N\), it is necessary and sufficient that \(F\) be of Tor-dimension \(\leq n\) in the ordinary sense, that is, that \(F\) admit a left resolution of length \(n\) by flat \(A\)-modules.

\sgaStableTarget{sga6.e01.para.5-5}\textbf{5.5.} Let \(X\in\Ob S\), and let \(m,n\in\mathbb Z\), with \(m\leq n\). In order that \(F\in\Ob\Dpx^-(A,X)\) be of flat amplitude contained in \([m,n]\), it is necessary and sufficient that there exist a covering family \((X_i\to X)\) such that \(F|X_i\) is of flat amplitude contained in \([m,n]\). In particular, if \(S/X\) is of finite type (SGA 4 VI 1.7), and if \(F\) is locally of finite Tor-dimension, then \(F\) is of finite Tor-dimension.

\sgaStableTarget{sga6.e01.prop.5-6}\textbf{Proposition 5.6.} Under the hypotheses of \hyperlink{sga6.e01.cor.2-16}{(2.16)}, if \(E'\) is of flat amplitude contained in \([m',n']\) as an object of \(\Dpx^-(A,S')\), and if \(E\) is of flat amplitude contained in \([m,n]\), then
\[
  E'\overset{\mathrm L}{\otimes}_A E
\]
is of flat amplitude contained in \([m+m',n+n']\).

\emph{Proof.} This is immediate: use form \hyperlink{sga6.e01.prop.5-1}{(5.1 (i))} of the notion of finite Tor-dimension and the fact that, if \(E\) is a flat left \(A\)-module, then \(f^{-1}(E)\) is a flat left \(f^{-1}(A)\)-module (SGA 4 IV, App.).

In particular:

\sgaStableTarget{sga6.e01.cor.5-6-1}\textbf{Corollary 5.6.1.} \textbf{a)} The functor
\[
  \mathrm L f^*:\Dpx^-(A,S)\longrightarrow \Dpx^-(A,S')
\]
induces a functor
\[
  \mathrm L f^*:\Dpx(A,S)_{\torf}\longrightarrow \Dpx(A,S')_{\torf}.
\]

\textbf{b)} If \(A\) is commutative, the functor
\[
  \overset{\mathrm L}{\otimes}_A:
  \Dpx^-(S)\times_S\Dpx^-(S)\longrightarrow \Dpx^-(S)
\]
induces a functor
\[
  \overset{\mathrm L}{\otimes}_A:
  \Dpx(S)_{\torf}\times_S\Dpx(S)_{\torf}
  \longrightarrow \Dpx(S)_{\torf}.
\]

\sgaStableTarget{sga6.e01.rem.5-6-2}\textbf{Remark 5.6.2.} Let \(E\in\Ob\Dpx^-(A,S)\), where \(S\) is the final object of \(S\), and let \(m,n\in\mathbb Z\), with \(m\leq n\). If \(E\) is of flat amplitude contained in \([m,n]\), then so is \(E_s\), for every point \(s\) of \(S\). The converse is true if \(S\) has enough points; use form \hyperlink{sga6.e01.prop.5-1}{(5.1(ii))} of the notion of flat amplitude and the fact that \(\overset{\mathrm L}{\otimes}_A\) commutes with passage to the fibers.

\sgaStableTarget{sga6.e01.exercise.5-7}\textbf{Exercises 5.7.} \textbf{a)} Suppose that \(A\) is commutative and that \((S',A')\) is faithfully flat over \((S,A)\) (cf. \hyperlink{sga6.e01.exercise.2-16-2}{(2.16.2)}). In order that \(E\in\Ob\Dpx^-(A,S)\)\footnote{Here \(S\) denotes the final object of the topos \(S\).} be of flat amplitude contained in \([m,n]\), it is necessary and sufficient that \(f^*(E)\) have that property.

\textbf{b)} Let \(F\in\Ob\Dpx^+(A,S)\)\textsuperscript{\thefootnote}. Show the equivalence of the following conditions:
\begin{enumerate}[label=(\roman*)]
\item \(F\) is isomorphic, in \(\Dpx^+(A,S)\), to a complex \(F'\) such that \(F'^i\) is injective for every \(i\) and zero for \(i\notin[m,n]\);
\item for every left \(A_S\)-module \(E\), one has \(\Ext^i(E,F)=0\) for \(i\notin[m,n]\).
\end{enumerate}
If \(F\) satisfies the preceding conditions, one says that \(F\) is of injective amplitude contained in \([m,n]\). The full subcategory of \(\Dpx^+(A,S)\) formed by the objects of finite injective amplitude is a triangulated subcategory, denoted \(\Dpx^+(A,S)_{\injf}\). Show that, if \(A\) is commutative, the functor
\[
  \mathrm R\Hom:\Dpx(S)^\circ\times\Dpx^+(S)\longrightarrow \Dpx(S)
\]
induces a functor
\[
  \mathrm R\Hom:\Dpx(S)^\circ_{\torf}\times\Dpx(S)_{\injf}
  \longrightarrow \Dpx(S)_{\injf}.
\]

We can now answer the question of \hyperlink{sga6.e01.para.5-0}{(5.0)} and state the main result of this section.

\sgaStableTarget{sga6.e01.thm.5-8}\textbf{Theorem 5.8.} Let \((S,A)\) be a ringed topos, \(E\in\Ob\Dpx(A,S)\), and let \(m,n\in\mathbb Z\), with \(m\leq n\). The following conditions are equivalent:
\begin{enumerate}[label=(\roman*)]
\item \(\perfamp(E)\subset[m,n]\) \hyperlink{sga6.e01.ex.4-8}{(4.8)};
\item \(E\) is \((m-1)\)-pseudo-coherent \hyperlink{sga6.e01.ex.2-15}{(2.15)} and \(\toramp(E)\subset[m,n]\) \hyperlink{sga6.e01.def.5-2}{(5.2)}.
\end{enumerate}
Since ``perfect'' plainly implies ``pseudo-coherent'', one deduces:

\sgaStableTarget{sga6.e01.cor.5-8-1}\textbf{Corollary 5.8.1.} In order that \(E\) be perfect \hyperlink{sga6.e01.ex.4-8}{(4.8)}, it is necessary and sufficient that \(E\) be pseudo-coherent \hyperlink{sga6.e01.ex.2-15}{(2.15)} and locally of finite Tor-dimension \hyperlink{sga6.e01.def.5-2}{(5.2)}.

For the proof, we shall need a few lemmas.

\sgaStableTarget{sga6.e01.lem.5-8-2}\textbf{Lemma 5.8.2.} Let \(B\) be a ring of \(S\). If \(E\) is a left \(A\)-module, \(G\) a right \(B\)-module, and \(F\) an \(A\)-\(B\)-bimodule (left over \(A\), right over \(B\)), there is a canonical functorial homomorphism
\sgaStableTarget{sga6.e01.eq.5-8-2-1}\begin{equation}\tag{5.8.2.1}
  \Hom_B(F,G)\otimes_A E
  \longrightarrow
  \Hom_B(\Hom_A(E,F),G).
\end{equation}
It is an isomorphism for \(G\) injective and \(E\) of finite presentation.

\emph{Proof.} First define the arrow. By the adjunction formula dear to Cartan, it amounts to defining a homomorphism of \(B\)-modules:
\stepcounter{footnote}
\sgaStableTarget{sga6.e01.eq.unnamed.occ5}\begin{equation}\tag{*}
  \Hom_B(F,G)\otimes_A E\otimes_{\mathbb Z_S}\Hom_A(E,F)
  \longrightarrow G.\text{\textsuperscript{\thefootnote}}
\end{equation}
\begin{NoHyper}
\footnotetext[\value{footnote}]{We write \(S\) for the final object of the topos \(S\), and \(\mathbb Z_S\) for the constant sheaf with value \(\mathbb Z\).}
\end{NoHyper}
There is a canonical homomorphism of \(A\)-\(B\)-bimodules
\sgaStableTarget{sga6.e01.eq.unnamed.occ6}\begin{equation}\tag{**}
  E\otimes_{\mathbb Z_S}\Hom_A(E,F)
  \longrightarrow F
\end{equation}
coming, by adjunction, from the identity arrow of \(\Hom_A(E,F)\). Similarly, there is a canonical homomorphism of \(B\)-modules
\sgaStableTarget{sga6.e01.eq.unnamed.occ7}\begin{equation}\tag{***}
  \Hom_B(F,G)\otimes_A F\longrightarrow G.
\end{equation}
We then define (*) as the composite of (**) and (***). To see that the arrow \hyperlink{sga6.e01.eq.5-8-2-1}{(5.8.2.1)} is an isomorphism when \(G\) is injective and \(E\) is of finite presentation, it is enough to note that, if \(G\) is injective, the source and the target of the arrow are right-exact functors in \(E\), and that the arrow induces the identity for \(E=A\).

Here is now a particular case of \hyperlink{sga6.e01.thm.5-8}{(5.8)}.

\sgaStableTarget{sga6.e01.lem.5-8-3}\textbf{Lemma 5.8.3.} Let \(E\) be a left \(A\)-module. The following conditions are equivalent:
\begin{enumerate}[label=(\roman*)]
\item \(E\) is locally a direct factor of a free module of finite type;
\item \(E\) is flat and of finite presentation.
\end{enumerate}

\emph{Proof} (after Bourbaki, Alg. comm., Chap. I, §2, Exer. 15). It is clear that (i) implies (ii). Let us prove (ii) \(\Rightarrow\) (i). By \hyperlink{sga6.e01.lem.1-3-1}{(1.3.1)}, it is enough to show that the functor \(\Hom_A(E,\cdot)\) is exact. Let
\[
  F\longrightarrow F''\longrightarrow 0
\]
be an exact sequence of left \(A\)-modules. We must show that the sequence
\[
  \Hom_A(E,F)\longrightarrow \Hom_A(E,F'')\longrightarrow 0
\]
is exact. For this it is enough to prove that, for every injective \(\mathbb Z_S\)-module \(I\), the sequence obtained by applying the functor \(\Hom_{\mathbb Z_S}(\cdot,I)\) is exact. But there is a commutative diagram, where the vertical arrows are the canonical arrows of \hyperlink{sga6.e01.eq.5-8-2-1}{(3.17.1)}:
\[
\begin{array}{ccccccccc}
0&\to&\Hom_{\mathbb Z_S}(F'',I)\otimes_A E&\to&\Hom_{\mathbb Z_S}(F,I)\otimes_A E\\
&&\downarrow&&\downarrow\\
0&\to&\Hom_{\mathbb Z_S}(\Hom_A(E,F''),I)&\to&\Hom_{\mathbb Z_S}(\Hom_A(E,F),I).
\end{array}
\]
The vertical arrows are isomorphisms by \hyperlink{sga6.e01.lem.5-8-2}{(5.8.2)}. The top row is exact because \(E\) is flat. Hence the bottom row is exact, q.e.d.

\emph{Proof of \hyperlink{sga6.e01.thm.5-8}{(5.8)}.} The implication (i) \(\Rightarrow\) (ii) follows trivially from the definitions. Let us prove (ii) \(\Rightarrow\) (i). Since the question is local, one may suppose \(E\) strictly \((m-1)\)-pseudo-coherent. Since, on the other hand, \(E\) is acyclic in degrees \(<m\), the truncation arrow
\[
\begin{array}{ccccccc}
\cdots&\to&E^{m-1}&\to&E^m&\to&E^{m+1}\to\cdots\\
&&\downarrow&&\downarrow&&\downarrow\\
\cdots&\to&0&\to&E^m/B^m&\to&E^{m+1}\to\cdots
\end{array}
\]
is a quasi-isomorphism. Let \(E'\) be the truncated complex. By construction, \(E'\) is of flat amplitude contained in \([m,n]\), and \(E'^i\) is free of finite type for \(i>m\). Therefore, by \hyperlink{sga6.e01.lem.5-1-1}{(5.1.1)}, \(E'^m\) is flat. On the other hand there is a distinguished triangle
\[
\begin{tikzcd}
& E'^m[-m] \arrow[dl,"+1"'] & \\
E'' \arrow[rr] && E' \arrow[ul]
\end{tikzcd}
\]
where \(E''\) is strictly perfect. By \hyperlink{sga6.e01.schol.2-6}{(2.6)}, \(E'^m[-m]\) is \((m-1)\)-pseudo-coherent; i.e. by \hyperlink{sga6.e01.prop.2-9}{(2.9)}, \(E'^m\) is of finite presentation. We conclude from \hyperlink{sga6.e01.lem.5-8-3}{(5.8.3)} that \(E'^m\) is locally a direct factor of a free module of finite type, hence that \(\perfamp(E)\subset[m,n]\).

\sgaStableTarget{sga6.e01.rem.5-8-4}\textbf{Remarks 5.8.4.} \textbf{a)} Using \hyperlink{sga6.e01.prop.2-5}{(2.5)}, one recovers the fact that perfect complexes form a triangulated subcategory of \(\Dpx(A,S)\).

\textbf{b)} A perfect complex, even with bounded cohomology, is not necessarily of finite Tor-dimension. It is so, however, if \(S\) is of finite type \hyperlink{sga6.e01.para.5-5}{(5.5)}. Write \(\Dpx(A,S)_{\mathrm{parf}}\) for the full sub-\(S\)-category of \(\Dpx^b(A,S)\) formed by perfect complexes of finite Tor-dimension. It follows at once from (\hyperlink{sga6.e01.cor.2-16-1}{2.16.1} b) and \hyperlink{sga6.e01.para.5-3}{(5.3)} that, if \(A\) is commutative, the functor \(\overset{\mathrm L}{\otimes}_A\) induces a functor
\[
  \overset{\mathrm L}{\otimes}_A:
  \Dpx(S)_{\mathrm{parf}}\times_S \Dpx^b(S)_{\mathrm{coh}}
  \longrightarrow \Dpx^b(S)_{\mathrm{coh}}.
\]
This formula is the starting point for the theory of intersections on schemes. Note, on the other hand, that the derived tensor product of two pseudo-coherent complexes with bounded cohomology is not, in general, of bounded cohomology.

\sgaStableTarget{sga6.e01.cor.5-9}\textbf{Corollary 5.9.} Suppose that \(S=(\mathrm{Ens})\) (so that \(A\) ``is'' an ordinary ring) and that every left \(A\)-module is of finite Tor-dimension; this is the case, for example, if \(A\) is of finite global cohomological dimension, cf. (MULT VI 2.6) and (EGA \(0_{\mathrm{IV}}\) 17.2.8). Then
\[
\text{(i)}\quad \Dpx^b(A,S)=\Dpx(A,S)_{\torf},
\qquad
\text{and}
\qquad
\text{(ii)}\quad \Dpx^b(A,S)_{\mathrm{coh}}=\Dpx(A,S)_{\mathrm{parf}}.
\]
In other words, in order that \(E\in\Ob\Dpx(A,S)\) be of finite Tor-dimension (respectively finite Tor-dimension and perfect), it is necessary and sufficient that \(E\) have bounded cohomology (respectively bounded and pseudo-coherent cohomology).

\emph{Proof.} By \hyperlink{sga6.e01.thm.5-8}{(5.8)}, it is enough to prove (i). Let \(E\in\Ob\Dpx^b(A,S)\). We may suppose, after extending \(E\) by \(0\), that \(E\in\Ob\Dpx^b(A,S)\), where \(S\) is the final object of \(S\). By induction on the number of non-zero cohomology objects of \(E\), one is reduced, by dévissage, to the case where \(E\) is an ordinary \(A\)-module; then one wins.

More generally, one may state:

\sgaStableTarget{sga6.e01.cor.5-10}\textbf{Corollary 5.10.} Let \(X\in\Ob S\). Suppose that \(S/X\) has enough points (SGA 4 IV 6.4.1) and that, for every point \(s\) of \(S/X\), every left \(A_s\)-module is of finite Tor-dimension. Then
\[
  \Dpx^b(A,X)_{\mathrm{coh}}=\Dpx^b(A,X)_{\mathrm{parf}}.
\]
Equivalently, in order that \(E\in\Ob\Dpx^b(A,X)\) be pseudo-coherent, it is necessary and sufficient that \(E\) be perfect.

\emph{Proof.} Let \(E\in\Ob\Dpx^b(A,X)_{\mathrm{coh}}\) and let \(s\) be a point of \(S/X\). One must show that \(E\) is perfect in a neighbourhood of \(s\), that is, that there exists an \(s\)-pointed object \(U\) of \(S\) over \(X\) such that \(E|U\) is perfect. Now, since the functor \(s^*\) is exact, (\hyperlink{sga6.e01.cor.2-16-1}{2.16.1} a) gives \(E_s\in\Ob\Dpx^b(A_s)_{\mathrm{coh}}\). Hence, by \hyperlink{sga6.e01.cor.5-9}{(5.9)}, \(E_s\) is perfect. The assertion is therefore a consequence of \hyperlink{sga6.e01.cor.3-7-2}{(3.7.2)} and of the following lemma.

\sgaStableTarget{sga6.e01.lem.5-10-1}\textbf{Lemma 5.10.1.} Let \(s\) be a point of \(S\), and let \(F\in\Ob\Dpx(A_s)_{\mathrm{parf}}\). There exist an \(s\)-pointed object \(U\) of \(S\) and \(F'\in\Ob\Dpx(A_U)_{\mathrm{parf}}\) such that \(F'_s=F\).

\emph{Proof.} This is evident if \(F\) is a free \(A_s\)-module of finite type. Suppose now that \(F\) is projective of finite type, hence the image of a projector
\[
  p:L\longrightarrow L,
\]
where \(L\) is free of finite type. There is then a neighbourhood \(U\) of \(s\), a free \(A_U\)-module \(L'\), and a homomorphism \(p':L'\to L'\) such that \(p'_s=p\). Since \(p_s^2=p_s\), we may suppose, after shrinking \(U\), that \(p'^2=p'\). Thus \(F'=\operatorname{Im}(p')\) is a direct factor of \(L'\) and satisfies \(F'_s=F\). The case where \(F\) is strictly perfect, that is, a bounded complex made up of projective modules of finite type, follows immediately.

\sgaStableTarget{sga6.e01.ex.5-11}\textbf{Examples 5.11.} \hyperlink{sga6.e01.cor.5-10}{Corollary (5.10)} applies in particular when \((S,A)\) is defined by a locally ringed space in regular local rings, for example a regular scheme or a smooth complex analytic space over \(\mathbb C\).

\sgaStableTarget{sga6.e01.sec.6}\section*{6. Rank of a perfect complex}

\sgaStableTarget{sga6.e01.para.6-0}\textbf{6.0.} Let \((X,\mathcal O_X)\) be a ringed space in local rings. It is well known that one can attach, to every locally free sheaf of finite type \(L\) on an open set \(U\) of \(X\), a function \(\rg_U(L)\), locally constant on \(U\) and with integral values, called the rank of \(L\). The family of functions \(\rg_U(L)\) is characterized by the following properties:
\begin{enumerate}[label=(\roman*)]
\item if \(V\subset U\) is an inclusion of open sets, then \(\rg_V(L|V)=\rg_U(L)|V\);
\item for every exact sequence of locally free \(\mathcal O_U\)-modules of finite type
\[
  0\longrightarrow L'\longrightarrow L\longrightarrow L''\longrightarrow 0,
\]
one has \(\rg_U(L)=\rg_U(L')+\rg_U(L'')\);
\item \(\rg_X(\mathcal O_X)=1\).
\end{enumerate}
We propose in this section to extend the notion of rank to perfect complexes.

\sgaStableTarget{sga6.e01.def.6-1}\textbf{Definition 6.1} (cf. SGA 5 VIII 2). Let \(\mathcal C\) be an additive (respectively triangulated) category, and let \(F\) be an abelian group. A function
\[
  f:\Ob(\mathcal C)\longrightarrow F
\]
is said to be an additive function from \(\mathcal C\) to \(F\) if, for all \(L,L',L''\in\Ob(\mathcal C)\) such that \(L\simeq L'\oplus L''\) (respectively such that there is a distinguished triangle
\[
\begin{tikzcd}
& L'' \arrow[dl,dashed] & \\
L' \arrow[rr] && L \arrow[ul]
\end{tikzcd}
\]
), one has
\[
  f(L)=f(L')+f(L'').
\]
The additive functions from \(\mathcal C\) to \(F\) form an abelian group, denoted \(\Add(\mathcal C,F)\).

\sgaStableTarget{sga6.e01.para.6-2}\textbf{6.2.} Let \(\mathcal C\) be a triangulated category, and let \(f\) be an additive function from \(\mathcal C\) to an abelian group \(F\). Then \(f\) induces an additive function on the underlying additive category \(\mathcal C^{\mathrm{ad}}\) of \(\mathcal C\); hence an inclusion
\[
  \Add(\mathcal C,F)\subset \Add(\mathcal C^{\mathrm{ad}},F).
\]
In particular, \(f(0)=0\) and \(f(L)=f(M)\) if \(L\simeq M\). On the other hand,
\[
  f(L[n])=(-1)^n f(L)
\]
for all \(L\in\Ob(\mathcal C)\) and \(n\in\mathbb Z\).

\sgaStableTarget{sga6.e01.para.6-3}\textbf{6.3.} Let \(\mathcal C\) be an additive (respectively triangulated) category. The group \(\Add(\mathcal C,F)\) depends functorially on \(F\). The functor \(\Add(\mathcal C,\cdot)\) is representable (cf. SGA 5 VIII 2) by an abelian group, denoted \(k(\mathcal C)\)\footnote{We shall not use here the usual notation \(K(\mathcal C)\), since \(K(\mathcal C)\) already denotes, when \(\mathcal C\) is additive, the category of complexes of \(\mathcal C\) ``up to homotopy.''}, endowed with an additive application
\[
  \cl_{\mathcal C}:\Ob(\mathcal C)\longrightarrow k(\mathcal C).
\]
If \(\mathcal C\) is additive, one takes for \(k(\mathcal C)\) the symmetrized monoid of isomorphism classes of objects of \(\mathcal C\). If \(\mathcal C\) is triangulated, one takes for \(k(\mathcal C)\) the quotient of the free abelian group generated by \(\Ob(\mathcal C)\) by the relations identifying \(L\) with \(L'+L''\) whenever there is a distinguished triangle
\[
  L'\longrightarrow L\longrightarrow L''\longrightarrow L'[1].
\]

The group \(k(\mathcal C)\) depends functorially on \(\mathcal C\), with respect to additive (respectively exact) functors. If \(u,u':\mathcal C\to\mathcal D\) are isomorphic additive (respectively exact) functors, then \(k(u)=k(u')\).

If \(\mathcal C\) is a triangulated category, one has a canonical epimorphism
\[
  k(\mathcal C^{\mathrm{ad}})\longrightarrow k(\mathcal C)
\]
(cf. \hyperlink{sga6.e01.para.6-2}{(6.2)}).

\sgaStableTarget{sga6.e01.lem.6-4}\textbf{Lemma 6.4.} Let \(\mathcal C\) be an additive category. The canonical functor
\[
  \mathcal C\longrightarrow \Kpx^b(\mathcal C)
\]
induces an isomorphism, functorial in \(\mathcal C\),
\[
  k(\mathcal C)\xrightarrow{\sim} k(\Kpx^b(\mathcal C)).
\]

\emph{Proof.} This is an immediate consequence of the following formula, which follows trivially from \hyperlink{sga6.e01.para.6-2}{(6.2)}:
\sgaStableTarget{sga6.e01.eq.6-4-1}\begin{equation}\tag{6.4.1}
  f(L)=\sum_i (-1)^i f(L^i)
\end{equation}
for \(L\in\Ob\Kpx^b(\mathcal C)\) and \(f\) an additive function on \(\Kpx^b(\mathcal C)\).

\sgaStableTarget{sga6.e01.def.6-5}\textbf{Definition 6.5.} Let \(S\) be a category, \(\mathcal C\) an additive (respectively triangulated) \(S\)-category \hyperlink{sga6.e01.para.1-1-1}{(1.1.1)}, and \(F\) an abelian presheaf on \(S\). A family of functions
\[
  f_X:\Ob(\mathcal C_X)\longrightarrow F(X),
  \qquad X\in\Ob S,
\]
is said to be an additive \(S\)-function from \(\mathcal C\) to \(F\) if the following conditions are satisfied:
\begin{enumerate}[label=(\roman*)]
\item for every \(X\in\Ob S\), \(f_X\) is an additive function \hyperlink{sga6.e01.def.6-1}{(6.1)};
\item for every arrow \(u:X\to Y\) of \(S\), the following diagram is commutative:
\[
\begin{tikzcd}
\Ob\mathcal C_Y \arrow[r,"f_Y"] \arrow[d,"u^*"'] & F(Y) \arrow[d,"F(u)"]\\
\Ob\mathcal C_X \arrow[r,"f_X"'] & F(X).
\end{tikzcd}
\]
\end{enumerate}
The additive \(S\)-functions from \(\mathcal C\) to \(F\) form an abelian group, denoted \(\Add_S(\mathcal C,F)\).

\sgaStableTarget{sga6.e01.rem.6-6}\textbf{Remarks 6.6.} \textbf{a)} Let \(\mathcal C^{\is}\) denote the fibered \(S\)-category whose fiber over \(X\in\Ob S\) is the category with the same objects as \(\mathcal C_X\) but whose arrows are the isomorphisms of \(\mathcal C_X\). Let \(\underline F\) denote the fibered \(S\)-category whose fiber over \(X\in\Ob S\) is the discrete category (i.e. the only arrows are identities) \(F(X)\), the inverse-image functor by \(u\in\mathrm{Fl}(S)\) being \(F(u)\). Then an additive \(S\)-function from \(\mathcal C\) to \(F\) is simply a Cartesian \(S\)-functor from \(\mathcal C^{\is}\) to \(\underline F\), inducing on each fiber an additive function in the sense of \hyperlink{sga6.e01.def.6-1}{(6.1)}. It will be convenient for us to interpret additivity as follows. Suppose, for definiteness, that \(\mathcal C\) is triangulated. The category
\[
  \Cart_S(\mathcal C^{\is},\underline F)
\]
(respectively \(\Cart_S(\Tr(\mathcal C)^{\is},\underline F)\) \hyperlink{sga6.e01.para.1-1-2}{(1.1.2)}) of Cartesian \(S\)-functors from \(\mathcal C^{\is}\) (respectively \(\Tr(\mathcal C)^{\is}\)) to \(\underline F\) is the discrete category associated to an abelian group, and there is a homomorphism
\[
  \rho:\Cart_S(\mathcal C^{\is},\underline F)
  \longrightarrow
  \Cart_S(\Tr(\mathcal C)^{\is},\underline F),
\]
defined by
\[
  \rho(f)(T)=f(L)-f(L')-f(L'')
\]
for \(f\in\Cart_S(\mathcal C^{\is},\underline F)\) and \(T=(L'\to L\to L''\to L'[1])\in\Ob\Tr(\mathcal C)\). To say that \(f\in\Cart_S(\mathcal C^{\is},\underline F)\) is additive means that \(\rho(f)=0\), or equivalently that one has a functorial canonical exact sequence
\sgaStableTarget{sga6.e01.eq.6-6-1}\begin{equation}\tag{6.6.1}
  0\longrightarrow \Add_S(\mathcal C,F)
  \longrightarrow \Cart_S(\mathcal C^{\is},\underline F)
  \xrightarrow{\rho}
  \Cart_S(\Tr(\mathcal C)^{\is},\underline F).
\end{equation}

\textbf{b)} Let \(u:X\to Y\) be an arrow of \(S\). The inverse-image functors by \(u\), being all isomorphic to one another, define by \hyperlink{sga6.e01.para.6-3}{(6.3)} a unique homomorphism
\[
  k(Y)\longrightarrow k(X).
\]
Consequently, \(X\mapsto k(X)\) defines an abelian presheaf on \(S\), denoted \(k(\mathcal C)\). An additive \(S\)-function from \(\mathcal C\) to \(F\) is then simply interpreted as a homomorphism of presheaves
\[
  k(\mathcal C)\longrightarrow F,
\]
or in other words
\[
  \Add_S(\mathcal C,F)=\Hom(k(\mathcal C),F).
\]
The presheaf \(k(\mathcal C)\) of course depends functorially on \(\mathcal C\) with respect to additive (respectively exact) \(S\)-functors.

\sgaStableTarget{sga6.e01.prop.6-7}\textbf{Proposition 6.7.} Let \(S\) be a site, \(\mathcal C\) a flat abelian \(S\)-category, and \(\mathcal C_0\) a strictly full additive sub-\(S\)-category of \(\mathcal C\) satisfying the hypotheses \hyperlink{sga6.e01.para.4-0}{(4.0)}. Then the morphism of abelian presheaves
\[
  k(\mathcal C_0)\longrightarrow k(\Parf(\mathcal C,\mathcal C_0))\qquad(\text{cf. }\hyperlink{sga6.e01.notation.4-9}{(4.9)})
\]
induced by the inclusion of \(\mathcal C_0\) in \(\Parf(\mathcal C,\mathcal C_0)\), defines an isomorphism on the associated sheaves. In other words, for every abelian sheaf \(F\) on \(S\), the restriction homomorphism
\[
  \Add_S(\Parf(\mathcal C,\mathcal C_0),F)
  \longrightarrow
  \Add_S(\mathcal C_0,F)
\]
is an isomorphism.

\emph{Proof.} Let \(F\) be an abelian sheaf on \(S\). We show that the arrow
\[
  \Add_S(\Parf(\mathcal C,\mathcal C_0),F)
  \longrightarrow
  \Add_S(\mathcal C_0,F)
\]
is an isomorphism. It factors as
\[
\Add_S(\Parf(\mathcal C,\mathcal C_0),F)
 \xrightarrow{a}
\Add_S(\Kpx^b(\mathcal C_0),F)
 \xrightarrow{b}
\Add_S(\mathcal C_0,F).
\]
By \hyperlink{sga6.e01.lem.6-4}{(6.4)}, the arrow \(b\) is an isomorphism. We are reduced to proving that \(a\) is an isomorphism.

Consider the following commutative diagram, whose rows are exact by \hyperlink{sga6.e01.eq.6-6-1}{(6.6.1)}, and whose vertical arrows are the restriction arrows:
\[
\begin{array}{ccccccccc}
0&\to&\Add_S(\Parf(\mathcal C),F)&\to&\Cart_S(\Parf(\mathcal C)^{\is},\underline F)&\to&\Cart_S(\Tr(\Parf(\mathcal C))^{\is},\underline F)\\
&&\downarrow a&&\downarrow&&\downarrow\\
0&\to&\Add_S(\Kpx^b(\mathcal C_0),F)&\to&\Cart_S(\Kpx^b(\mathcal C_0)^{\is},\underline F)&\to&\Cart_S(\Tr(\Kpx^b(\mathcal C_0))^{\is},\underline F).
\end{array}
\]
Since \(F\) is a sheaf, \(\underline F\) is a stack (G \hyperlink{sga6.e02.def.2-2-4}{II 2.2.4}); hence, by \hyperlink{sga6.e01.prop.4-10}{(4.10)}, the two right vertical arrows are isomorphisms. The same is therefore true of \(a\), q.e.d.

\sgaStableTarget{sga6.e01.rem.6-7-1}\textbf{Remark 6.7.1} (cf. \hyperlink{sga6.e01.rem.4-10-1}{4.10.1}). In \hyperlink{sga6.e01.prop.6-7}{(6.7)}, one may replace \(\Parf(\mathcal C,\mathcal C_0)\) by \(\Dpx^*(\mathcal C)_{\mathcal C_0\text{-}\mathrm{parf}}\) (\(*=-\) or \(b\)).

\sgaStableTarget{sga6.e01.para.6-8}\textbf{6.8.} Let \((S,A)\) be a ringed topos, and let \(\mathcal C\) and \(\mathcal C_0\) be the fibered \(S\)-categories considered in \hyperlink{sga6.e01.ex.4-8}{(4.8)}. There are homomorphisms of abelian presheaves
\[
  k(\mathcal C_0)
  \xrightarrow{\alpha}
  k(\Dpx^-(A,S)_{\mathrm{parf}})
  \xrightarrow{\beta}
  k(\Parf(A,S)),
\]
which, by \hyperlink{sga6.e01.prop.6-7}{(6.7)}, induce isomorphisms on the associated sheaves. The homomorphism \(\alpha\) is compatible with inverse images: that is, if
\[
  u:(S',A')\longrightarrow(S,A)
\]
is a morphism of ringed topoi, there is a commutative diagram, whose vertical arrows are induced by the functor \(\mathrm L u^*\) (cf. \hyperlink{sga6.e01.cor.4-19-1}{(4.19.1)}):
\[
\begin{tikzcd}
 k(\mathcal C_0) \arrow[r] \arrow[d] & k(\Dpx^-(A,S)_{\mathrm{parf}}) \arrow[d]\\
 k(\mathcal C'_0) \arrow[r] & k(\Dpx^-(A',S')_{\mathrm{parf}}),
\end{tikzcd}
\]
where \(\mathcal C'_0\) is defined analogously to \(\mathcal C_0\). If \(A\) is commutative, the derived tensor product endows \(k(\mathcal C_0)\) and \(k(\Dpx^-(A,S)_{\mathrm{parf}})\) with structures of presheaves of rings (cf. \hyperlink{sga6.e01.cor.4-19-1}{(4.19.1)}), and \(\alpha\) is a homomorphism of presheaves of rings.

\sgaStableTarget{sga6.e01.para.6-9}\textbf{6.9.} In the situation of \hyperlink{sga6.e01.para.6-8}{(6.8)}, suppose that \((S,A)\) is locally ringed (cf. \hyperlink{sga6.e01.lem.2-15-1}{(2.15.1)}). Then, for every \(X\in\Ob S\), \(\mathcal C_{0,X}\) is the category of locally free \(A_X\)-modules of finite type. Let \(\mathcal C_1\) be the fibered \(S\)-category defined by
\[
  X\longmapsto\text{the category of free }A_X\text{-modules of finite type}.
\]
It is clear that \(\mathcal C_1\) satisfies the hypotheses \hyperlink{sga6.e01.para.4-0}{(4.0)}. An object of \(\Dpx(S)\) is perfect if and only if it is perfect relative to \(\mathcal C_1\). There are homomorphisms of abelian presheaves
\[
  k(\mathcal C_1)\longrightarrow k(\mathcal C_0)\longrightarrow
  k(\Dpx^-(S)_{\mathrm{parf}})\longrightarrow k(\Parf(S))
\]
(the two on the left being homomorphisms of presheaves of rings), which, by \hyperlink{sga6.e01.prop.6-7}{(6.7)}, induce isomorphisms on the associated sheaves. We shall show that the associated sheaves are in fact canonically isomorphic to the constant sheaf \(\mathbb Z_S\).

\sgaStableTarget{sga6.e01.lem.6-10}\textbf{Lemma 6.10.} Let \((S,A)\) be a locally ringed topos, let \(U\) be an object of \(S\) distinct from the empty object, and let \(p,q\in\mathbb N\). If
\[
  A_U^p\simeq A_U^q,
\]
then \(p=q\).

\emph{Proof.} By an argument similar to that given in the proof of \hyperlink{sga6.e01.lem.2-15-1}{(2.15.1)}, one is reduced to the case where \((S,A)\) is the locally ringed topos defined by a prescheme; in this case the assertion is trivial: take a point in \(U\).

\sgaStableTarget{sga6.e01.def.6-11}\textbf{Definition 6.11.} The rank, denoted \(\rg\), is the additive \(S\)-function from \(\mathcal C_1\) to \(\mathbb Z_S\) defined by
\[
  \rg(L)=p\quad\text{if } L\simeq A_X^p,
  \qquad X\ne\varnothing,
\]
\[
  \rg(L)=0\quad\text{if } L\in\Ob\mathcal C_{1,X},
  \qquad X=\varnothing.
\]
By \hyperlink{sga6.e01.prop.6-7}{(6.7)}, the rank function extends uniquely to an additive \(S\)-function from \(\Parf(S)\) to \(\mathbb Z_S\), which we shall again call rank and denote by \(\rg\).

\sgaStableTarget{sga6.e01.prop.6-12}\textbf{Proposition 6.12.} \textbf{a)} For \(L,M\in\Ob\Dpx^-(X)_{\mathrm{parf}}\), with \(X\in\Ob S\), one has
\[
  \rg(L\overset{\mathrm L}{\otimes} M)=\rg(L)\,\rg(M).
\]

\textbf{b)} Let
\[
  u:(S',A')\longrightarrow(S,A)
\]
be a morphism of locally ringed topoi. For \(M\in\Ob\Dpx^-(S)_{\mathrm{parf}}\), one has
\[
  \rg(\mathrm L u^*(M))=u^*(\rg(M)).
\]
In other words, the rank homomorphism
\[
  \rg:k(\Dpx^-(S)_{\mathrm{parf}})\longrightarrow \mathbb Z_S
\]
is a homomorphism of presheaves of rings, compatible with inverse images.



% Additional notation for Exposé I
\providecommand{\Cpx}{\mathrm C}
\providecommand{\Kpx}{\mathrm K}
\providecommand{\Dpx}{\mathrm D}
\providecommand{\Ho}{\mathrm H}
\providecommand{\RHom}{\mathrm R\!\operatorname{Hom}}
\providecommand{\Ob}{\operatorname{ob}}
\providecommand{\Hom}{\operatorname{Hom}}
\providecommand{\RHom}{\mathrm R\!\operatorname{Hom}}
\providecommand{\Ext}{\operatorname{Ext}}
\providecommand{\Tor}{\operatorname{Tor}}
\providecommand{\Tr}{\operatorname{Tr}}
\providecommand{\Add}{\operatorname{Add}}
\providecommand{\Cart}{\operatorname{Cart}}
\providecommand{\Parf}{\operatorname{Parf}}
\providecommand{\coh}{\operatorname{coh}}
\providecommand{\amp}{\operatorname{amp}}
\providecommand{\cl}{\operatorname{cl}}
\providecommand{\rg}{\operatorname{rg}}
\providecommand{\Id}{\operatorname{Id}}
\providecommand{\Ker}{\operatorname{Ker}}
\providecommand{\Im}{\operatorname{Im}}
\providecommand{\cH}{\mathcal H}
\providecommand{\Zs}{\mathbb Z_S}
\providecommand{\A}{\mathcal A}
\providecommand{\cC}{\mathcal C}
\providecommand{\cCzero}{\mathcal C_0}
\providecommand{\otL}{\mathop{\otimes}^{\mathrm L}}
\providecommand{\tamp}{\operatorname{tor.amp}}
\providecommand{\pamp}{\operatorname{parf.amp}}
\providecommand{\pdim}{\operatorname{parf.dim}}
\providecommand{\bon}{\mathrm b}
\providecommand{\parf}{\operatorname{parf}}
\providecommand{\parff}{\operatorname{parf}}
\providecommand{\ac}{\mathrm{ac}}


% Additional notation for Exposé I, completion pages 154--166
\providecommand{\End}{\operatorname{End}}
\providecommand{\tr}{\operatorname{tr}}

\emph{Proof.} It is enough to show that the homomorphism induced by rank on the sheaf associated to \(k(\Dpx^-(S)_{\Parf})\) is a homomorphism of sheaves of rings compatible with inverse images. By \hyperlink{sga6.e01.para.6-9}{(6.9)}, one may replace \(k(\Dpx^-(S)_{\Parf})\) by \(k(\mathcal C_1)\). We are thus reduced to verifying a) and b) for free modules of finite type, which is trivial.

\sgaStableTarget{sga6.e01.para.6-13}\textbf{6.13.} With \((S,A)\) still denoting a locally ringed topos, let \(X\in\Ob S\) and \(n\in H^0(X,\mathbb Z_S)\). We define an object of \(\Parf(X)\), denoted \(A_X^n\), as follows. If \(n\) is constant of value an integer \(n\geq0\) (respectively \(n<0\)), put \(A_X^n=A_X^{\oplus n}\) (respectively \(A_X^n=A_X^{-n}[1]\)). In the general case, decompose \(X\) into disjoint pieces \(X_i\) on which \(n|X_i\) is constant, and define \(A_X^n\) as the unique object of \(\Parf(X)\) inducing \(A_{X_i}^{n|X_i}\) on each \(X_i\). More generally, for \(L\in\Ob\Parf(X)\), one may define \(L^n\) by putting \(L^n=A_X^n\otimes_A L\).

The map
\[
  i_X:H^0(X,\mathbb Z_S)\longrightarrow k(\Parf(X)),
  \qquad i_X(n)=\cl(A_X^n),
\]
is a homomorphism of rings, with \(i_X(1)=1\), and defines, as \(X\) varies, a homomorphism of presheaves of rings
\[
  i:\mathbb Z_S\longrightarrow k(\Parf(S)).
\]

\sgaStableTarget{sga6.e01.prop.6-14}\textbf{Proposition 6.14.} Under the hypotheses of \hyperlink{sga6.e01.para.6-13}{(6.13)}, one has
\[
  \rg\circ i=\Id(\mathbb Z_S).
\]
In other words, \(k(\Parf(S))\), considered as a presheaf of \(\mathbb Z_S\)-algebras by means of \(i\), is augmented toward \(\mathbb Z_S\) by the rank. Moreover, the rank
\[
  \rg:k(\Parf(S))\longrightarrow \mathbb Z_S
\]
defines an isomorphism on the associated sheaves.

\emph{Proof.} The first assertion follows trivially from the relations
\[
  \cl(L[1])=-\cl(L),\qquad \cl(L^n)=n\cl(L),
  \qquad L\in\Ob\Parf(S).
\]
Let us prove the second. First observe that \(i\) factors in the evident way through \(k(\mathcal C_1)\). Since \(\rg\circ i\) is the identity, we are reduced, by \hyperlink{sga6.e01.para.6-9}{(6.9)}, to proving that the homomorphism
\[
  i\circ\rg:k(\mathcal C_1)\longrightarrow k(\mathcal C_1)
\]
induces an isomorphism on the associated sheaves; equivalently, for \(L\in\Ob\mathcal C_1\), that \(i(\rg(L))\) is locally isomorphic to \(L\). This is immediate.

One deduces, by (\hyperlink{sga6.e01.rem.6-6}{6.6} b):

\sgaStableTarget{sga6.e01.cor.6-15}\textbf{Corollary 6.15.} Under the hypotheses of \hyperlink{sga6.e01.para.6-13}{(6.13)}, for every abelian sheaf \(F\) on \(S\), the group of \(S\)-additive functions from \(\Parf(S)\) to \(F\) is canonically identified with \(H^0(S,F)\), where \(S\) is the final object of \(S\). The rank function corresponds, for \(F=\mathbb Z_S\), to the unit section of \(\mathbb Z_S\).

\sgaStableTarget{sga6.e01.sec.7}\section*{7. Duality of perfect complexes}

\sgaStableTarget{sga6.e01.para.7-0}\textbf{7.0.} In this section, \((S,A)\) denotes a topos ringed in commutative rings, and \(S\) also denotes the final object of \(S\). We propose to generalize to perfect complexes the duality notions that are well known for locally free sheaves of finite type.

\sgaStableTarget{sga6.e01.prop.7-1}\textbf{Proposition 7.1.} Let \(E\in\Ob\Dpx^-(S)_{\Parf}\) and \(F\in\Ob\Dpx^+(S)\).

\textbf{a)} If \(F\) has locally bounded cohomology, then so does \(\RHom(E,F)\). If \(F\) is pseudo-coherent (respectively locally of finite Tor-dimension, respectively perfect), then the same is true of \(\RHom(E,F)\).

\textbf{b)} Suppose, moreover, that \(E\) is of finite Tor-dimension, a condition satisfied if \(S\) is of finite type. Then, if \(F\) has bounded cohomology (respectively finite Tor-dimension), the same is true of \(\RHom(E,F)\). More precisely, let \(a,b,m,n\) be integers such that
\[
  \tamp(E)\subset[a,b]
\]
(or, equivalently by \hyperlink{sga6.e01.thm.5-8}{(5.8)}, \(\pamp(E)\subset[a,b]\)), and such that \(\Ho^i(F)=0\) for \(i\notin[m,n]\) (respectively \(\tamp(F)\subset[m,n]\)). Then
\[
  \Ho^i(\RHom(E,F))=0\quad(i\notin[m-b,n-a]),
\]
respectively
\[
  \tamp(\RHom(E,F))\subset[m-b,n-a].
\]

\emph{Proof.} Assertion a) is local, so we may suppose \(E\) strictly perfect, and then conclude by devissage on \(E\), reducing to \(E=A\). For assertion b), we may suppose \(E\) strictly perfect with \(E^i=0\) for \(i\notin[a,b]\), and \(F^i=0\) for \(i\notin[m,n]\) (respectively \(F^i\) flat for every \(i\) and \(F^i=0\) for \(i\notin[m,n]\)). It is then enough to apply the following lemma.

\sgaStableTarget{sga6.e01.lem.7-1-1}\textbf{Lemma 7.1.1.} Let \(E\in\Ob\Dpx^-(S)\) and \(F\in\Ob\Dpx^+(S)\). If \(E\) is strictly pseudo-coherent, that is, in the sense of \hyperlink{sga6.e01.def.2-1}{(2.1)}, with components direct factors of free finite-type modules, then the canonical arrow of \(\Dpx^+(S)\)
\[
  \Hom^\bullet(E,F)\longrightarrow \RHom(E,F)
\]
is an isomorphism.

\emph{Proof.} The spectral sequence
\[
  E_2^{p0}=\Ho^p(\Hom^\bullet(E,F))
  \Longrightarrow \Ext^p(E,F)
\]
where \(\Ext^q(E^i,F^j)=0\) for every \(q>0\) by \hyperlink{sga6.e01.lem.1-3-1}{(1.3.1)}, degenerates. One obtains isomorphisms
\[
  E_2^{p0}=\Ho^p(\Hom^\bullet(E,F))
  \xrightarrow{\sim}\Ext^p(E,F)=\Ho^p(\RHom(E,F)),
\]
which prove the assertion.

\sgaStableTarget{sga6.e01.prop.7-1-2}\textbf{Proposition 7.1.2.} Let
\[
  f:(S',A')\longrightarrow(S,A)
\]
be a morphism of ringed topoi. Let \(E\in\Ob\Dpx^-(S)\) and \(F\in\Ob\Dpx^b(S)\) be such that \(\RHom(E,F)\in\Ob\Dpx^b(S)\) and \(\mathrm Lf^*(F)\in\Ob\Dpx^b(S')\). There is a functorial canonical arrow
\[
  \mathrm Lf^*\RHom(E,F)
  \longrightarrow
  \RHom(\mathrm Lf^*E,\mathrm Lf^*F).
\]
This is an isomorphism if \(E\) is perfect.

\emph{Proof.} By the adjunction of \(\mathrm Lf^*\) and \(\mathrm Rf_*\), it is enough to define an arrow
\[
  \RHom(E,F)
  \longrightarrow
  \mathrm Rf_*\RHom(\mathrm Lf^*E,\mathrm Lf^*F).
\]
For this we use the adjunction arrow \(F\to\mathrm Rf_*\mathrm Lf^*F\) and the duality isomorphism
\[
  \mathrm Rf_*\RHom(\mathrm Lf^*E,\mathrm Lf^*F)
  \xrightarrow{\sim}
  \RHom(E,\mathrm Rf_*\mathrm Lf^*F).
\]
To see that the arrow of \hyperlink{sga6.e01.prop.7-1-2}{(7.1.2)} is an isomorphism for \(E\) perfect, one reduces trivially, by localization and devissage, to the case \(E=A\), where it is a tautology.

\sgaStableTarget{sga6.e01.prop.7-2}\textbf{Proposition 7.2.} There exists (cf. SGA 5 I 1.6), for \(E\in\Ob\Dpx(S)\) and \(F\in\Ob\Dpx^+(S)\), a functorial canonical arrow
\[
  E\longrightarrow\RHom(\RHom(E,F),F).
\]
This is an isomorphism if \(E\in\Ob\Dpx(S)_{\Parf}\) and \(F=A\).

\emph{Proof.} Define the arrow by resolving \(F\) injectively and using the canonical arrow in \(\Kpx(S)\)
\[
  E\longrightarrow \Hom^\bullet(\Hom^\bullet(E,F),F).
\]
To see that the arrow in question is an isomorphism if \(E\) is perfect and \(F=A\), one reduces, by localization and devissage, to \(E=A\), where it is evident.

\sgaStableTarget{sga6.e01.def.7-3}\textbf{Definition 7.3.} Let \(E\in\Ob\Dpx(S)\). We shall put
\[
  \RHom(E,A)=E^\vee,
\]
and we shall say that \(E^\vee\) is the dual of \(E\).

By \hyperlink{sga6.e01.prop.7-1}{(7.1)} and \hyperlink{sga6.e01.prop.7-2}{(7.2)}, if \(E\in\Ob\Dpx(S)\) is perfect (respectively perfect and of finite Tor-dimension), then so is \(E^\vee\), and the canonical arrow \(E\to E^{\vee\vee}\) is an isomorphism. In other words, the functor \(\RHom(\,\cdot\,,A)\) defines an anti-auto-equivalence of the category \(\Dpx(S)_{\Parf}\) (respectively \(\Dpx(S)_{\parf}\), cf. (\hyperlink{sga6.e01.rem.5-8-4}{5.8.4} b)) with itself.

\sgaStableTarget{sga6.e01.lem.7-4}\textbf{Lemma 7.4} (Cartan isomorphism) (cf. ([H] II 5.15)). For \(E\in\Ob\Dpx(S)\), \(F\in\Ob\Dpx^-(S)\), and \(G\in\Ob\Dpx^+(S)\), there are functorial canonical isomorphisms
\[
\begin{aligned}
\RHom(E\otL F,G)&\xrightarrow{\sim}\RHom(E,\RHom(F,G)),\\
\RHom(E\otimes F,G)&\xrightarrow{\sim}\RHom(E,\RHom(F,G)),\\
\Hom(E\otimes F,G)&\xrightarrow{\sim}\Hom(E,\RHom(F,G)).
\end{aligned}
\]

\emph{Proof.} This is standard: resolve \(G\) injectively and \(F\) flatly.

\sgaStableTarget{sga6.e01.lem.7-5}\textbf{Lemma 7.5.} For \(E\in\Ob\Dpx^-(S)\) and \(F\in\Ob\Dpx^+(S)\), there is a functorial canonical homomorphism
\[
  E\otL \RHom(E,F)\longrightarrow F.
\]

\sgaStableTarget{sga6.e01.prop.7-6}\textbf{Proposition 7.6} (cf. ([H] II 5.13)). Let \(E\in\Ob\Dpx(S)\), \(F\in\Ob\Dpx^+(S)\), and \(G\in\Ob\Dpx^+(S)\). Under either of the following hypotheses:
\begin{enumerate}[label=(\roman*)]
\item \(G\) is of finite Tor-dimension;
\item \(E\in\Ob\Dpx^-(S)\), \(F\) is of finite Tor-dimension, and \(G\in\Ob\Dpx^b(S)\),
\end{enumerate}
there is a functorial canonical homomorphism
\[
  \RHom(E,F)\otL G
  \longrightarrow
  \RHom(E,F\otL G).
\]
This is an isomorphism when \(E\) is perfect.

\emph{Proof.} Let us define the arrow. Under hypothesis (i), it is enough to resolve \(F\) injectively and \(G\) flatly. Under hypothesis (ii), by \hyperlink{sga6.e01.lem.7-4}{(7.4)} it is equivalent to define an arrow
\[
  E\otL\RHom(E,F)\otL G\longrightarrow F\otL G.
\]
Apply the functor \(-\otL G\) to the arrow of \hyperlink{sga6.e01.lem.7-5}{(7.5)}.

To see that the arrow is an isomorphism when \(E\) is perfect, use the standard method (cf. the proof of \hyperlink{sga6.e01.prop.7-2}{(7.2)}): reduce, by localization and devissage, to the case \(E=A\).

\sgaStableTarget{sga6.e01.cor.7-7}\textbf{Corollary 7.7.} There exists, for \(E\in\Ob\Dpx(S)_{\Parf}\) and \(F\in\Ob\Dpx^+(S)\), a functorial isomorphism
\[
  E^\vee\otL F\xrightarrow{\sim}\RHom(E,F),
\]
with the notation of \hyperlink{sga6.e01.def.7-3}{(7.3)}.

\emph{Proof.} Apply \hyperlink{sga6.e01.prop.7-6}{(7.6)}, with \(F\) and \(G\) replaced respectively by \(A\) and \(F\), and use the fact \hyperlink{sga6.e01.prop.7-1}{(7.1)} that \(E^\vee\in\Ob\Dpx(S)_{\Parf}\).

\sgaStableTarget{sga6.e01.cor.7-8}\textbf{Corollary 7.8.} Let \(E,F\in\Ob\Dpx(S)_{\Parf}\). There is a natural arrow
\[
  \RHom(E,F)\otL \RHom(F,E)\longrightarrow A,
\]
which is a ``perfect duality,'' that is, modulo \hyperlink{sga6.e01.lem.7-4}{(7.4)}, it identifies each of the complexes \(\RHom(E,F)\) and \(\RHom(F,E)\) with the dual of the other, in the sense of \hyperlink{sga6.e01.def.7-3}{(7.3)}.

\emph{Proof.} By \hyperlink{sga6.e01.lem.7-5}{(7.5)}, one has natural pairings \(E^\vee\otL E\to A\) and \(F^\vee\otL F\to A\). Hence, by tensor product, one obtains a natural pairing
\[
  E^\vee\otL F\otL F^\vee\otL E\longrightarrow A,
\]
which, in view of \hyperlink{sga6.e01.cor.7-7}{(7.7)}, gives the announced pairing on the \(\RHom\)'s. Let us show that it is a perfect duality. There are canonical isomorphisms
\[
\begin{aligned}
\RHom(F^\vee\otL E,A)&\xrightarrow{\sim}\RHom(F^\vee,E^\vee) &&\hyperlink{sga6.e01.lem.7-4}{(7.4)}\\
&\xrightarrow{\sim} F^{\vee\vee}\otL E^\vee &&\hyperlink{sga6.e01.prop.7-6}{(7.6)}\\
&\xrightarrow{\sim} F\otL E^\vee &&\hyperlink{sga6.e01.prop.7-2}{(7.2)}.
\end{aligned}
\]
Thus, up to a compatibility verification, the result follows. One could also localize and devissage to reduce to \(E=F=A\).


\sgaStableTarget{sga6.e01.exercise.7-9}\textbf{Exercise 7.9.} Let \(E,F\in\Ob\Dpx(S)_{\Parf}\). Show that there is a canonical isomorphism
\[
  E^\vee\otL F^\vee \xrightarrow{\sim} (E\otL F)^\vee .
\]

\sgaStableTarget{sga6.e01.sec.8}\section*{8. Traces and cup-products}

\sgaStableTarget{sga6.e01.para.8-0}\textbf{8.0.} In this section we generalize to perfect complexes the notion of the trace of an endomorphism of a projective module of finite type (see SGA 5 III for applications). The hypotheses and notation are those of \hyperlink{sga6.e01.para.7-0}{(7.0)}.

\sgaStableTarget{sga6.e01.para.8-1}\textbf{8.1.} Let \(E\in\Ob\Dpx(S)_{\parf}\) (\hyperlink{sga6.e01.rem.5-8-4}{5.8.4} b). There is a natural arrow
\[
  \tau:\RHom(E,E)\longrightarrow A,
\]
obtained by composing the inverse of the isomorphism \hyperlink{sga6.e01.cor.7-7}{(7.7)}
\[
  \RHom(E,E)\xrightarrow{\sim}E^\vee\otL E
\]
with the canonical arrow \hyperlink{sga6.e01.lem.7-5}{(7.5)} \(E^\vee\otL E\to A\).

Let \(U\in\Ob S\). Applying to \(\tau\) the functor \(H^0(U,\cdot)\), and using the fact that \(H^0(U,\RHom(E,E))\) identifies with \(\Hom_{\Dpx(U)}(E|U,E|U)\), one obtains a homomorphism of \(H^0(U,A)\)-modules
\[
  \Tr_{E|U}:\Hom_{\Dpx(U)}(E|U,E|U)\longrightarrow H^0(U,A),
\]
which we shall call the trace homomorphism. The finite Tor-dimension hypothesis on \(E\) is in fact superfluous for the definition of \(\Tr_{E|U}\); see \hyperlink{sga6.e01.rem.8-5}{(8.5)}.

Before going further, we give some properties of the trace that follow immediately from the definition.

\sgaStableTarget{sga6.e01.para.8-1-1}\textbf{8.1.1.} If \(E=A\), then
\[
  \Tr_{A|U}:\Hom(A|U,A|U)\longrightarrow H^0(U,A)
\]
is just the well-known isomorphism. If \(E\) is a direct factor of a free module of finite type, an endomorphism of \(E\) over \(U\) is locally defined by
\[
  e\longmapsto \sum_i\langle e,x_i^\vee\rangle x_i,
\]
where \(x_i\) (respectively \(x_i^\vee\)) is a local section of \(E\) (respectively \(E^\vee\)); the corresponding trace is \(\sum_i\langle x_i,x_i^\vee\rangle\). It follows easily that, if \(E\) and \(F\) are direct factors of free modules of finite type, and if
\[
  u:(E+F)|U\longrightarrow(E+F)|U
\]
is defined by a matrix \(\begin{pmatrix}a&b\\ c&d\end{pmatrix}\), then
\[
  \Tr_{(E+F)|U}(u)=\Tr_{E|U}(a)+\Tr_{F|U}(d).
\]

\sgaStableTarget{sga6.e01.para.8-1-2}\textbf{8.1.2.} Let \(E\) be a strictly perfect complex, let \(U\in\Ob S\), and let
\[
  u:E|U\longrightarrow E|U
\]
be a homomorphism of complexes, with components \(u^i:E^i|U\to E^i|U\). Then
\[
  \Tr_{E|U}(u)=\sum_i(-1)^i\Tr_{E^i|U}(u^i).
\]
Indeed, the question being local, one may suppose that \(u\) is defined by a section of \(H^0(E^\vee\otimes E)\), and even that this section is of the form
\[
  \sum_{i,j}x_{ij}^\vee\otimes x_{ij},
\]
where \(x_{ij}\) (respectively \(x_{ij}^\vee\)) is a section of \(E^i\) (respectively \(E^{i\vee}\)). Making explicit the arrow \hyperlink{sga6.e01.lem.7-5}{(7.5)} \(E^\vee\otL E\to A\), one finds
\[
  \Tr_{E|U}(u)=\sum_{i,j}(-1)^i\langle x_{ij},x_{ij}^\vee\rangle
  =\sum_i(-1)^i\Tr_{E^i|U}(u^i).
\]

\sgaStableTarget{sga6.e01.para.8-2}\textbf{8.2.} Let \(E,F\in\Ob\Dpx(S)_{\Parf}\). We have seen in \hyperlink{sga6.e01.cor.7-8}{(7.8)} that there is a canonical homomorphism
\[
  \chi:\RHom(E,F)\otL\RHom(F,E)\longrightarrow A.
\]
Let \(U\in\Ob S\), and let
\[
  u\in H^0(U,\RHom(E,F)),\qquad v\in H^0(U,\RHom(F,E)).
\]
One may interpret \(u\) (respectively \(v\)) as a homomorphism, in \(\Dpx(U)\), from \(A|U\) to \(\RHom(E,F)|U\) (respectively \(\RHom(F,E)|U\)), and form
\[
  u\otL v:A|U\longrightarrow(\RHom(E,F)\otL\RHom(F,E))|U.
\]
The map \((u,v)\mapsto\langle u,v\rangle=\chi\circ(u\otL v)\) defines an \(H^0(U,A)\)-bilinear pairing
\[
  \Hom_{\Dpx(U)}(E|U,F|U)\otimes_{H^0(U,A)}\Hom_{\Dpx(U)}(F|U,E|U)
  \longrightarrow H^0(U,A),
\]
which we shall call the cup-product. The same comment as for the definition of \(\Tr_{E|U}\) applies here.

The relation between trace and cup-product is expressed in the following proposition.

\sgaStableTarget{sga6.e01.prop.8-3}\textbf{Proposition 8.3.} In the situation of \hyperlink{sga6.e01.para.8-2}{(8.2)}, for \(u\in\Hom(E|U,F|U)\) and \(v\in\Hom(F|U,E|U)\), one has
\[
  \Tr_{E|U}(vu)=\Tr_{F|U}(uv)=\langle u,v\rangle.
\]
In particular, if \(E=F\), then
\[
  \Tr_{E|U}(u)=\langle u,\Id\rangle=\langle\Id,u\rangle.
\]

\emph{Proof.} For simplicity we suppose \(U=S\). We therefore have two arrows
\[
  u:A\longrightarrow E^\vee\otL F,
  \qquad
  v:A\longrightarrow F^\vee\otL E,
\]
and hence
\[
  u\otL v:A\longrightarrow E^\vee\otL F\otL F^\vee\otL E.
\]
On the other hand, we have seen that \(vu:A\to E^\vee\otL E\) and \(uv:A\to F^\vee\otL F\). The proposition follows from the commutativity of the following diagrams, whose verification is purely mechanical:
\[
\begin{tikzcd}[column sep=large,row sep=large]
 A \arrow[r,"u\otL v"] \arrow[dr,"vu"'] & E^\vee\otL F\otL F^\vee\otL E \arrow[d,"\Id\otimes\tau_F"] \\
 & E^\vee\otL E
\end{tikzcd}
\qquad
\begin{tikzcd}[column sep=large,row sep=large]
 A \arrow[r,"u\otL v"] \arrow[dr,"uv"'] & E^\vee\otL F\otL F^\vee\otL E \arrow[d,"\tau_E\otimes\Id"] \\
 & F^\vee\otL F
\end{tikzcd}
\]
\[
\begin{tikzcd}[column sep=large,row sep=large]
E^\vee\otL F\otL F^\vee\otL E \arrow[r,"\Id\otimes\tau_E"] \arrow[d,"\Id\otimes\tau_F"'] & F^\vee\otL F \arrow[d,"\tau_F"]\\
E^\vee\otL E \arrow[r,"\tau_E"'] & A .
\end{tikzcd}
\]

\sgaStableTarget{sga6.e01.cor.8-4}\textbf{Corollary 8.4.} Let \(E,F,U\) be as in \hyperlink{sga6.e01.para.8-2}{(8.2)}, let \(u\in\Hom(E|U,E|U)\), and let \(s:E|U\to F|U\) be an isomorphism. Then \(sus^{-1}\in\Hom(F|U,F|U)\), and
\[
  \Tr_{F|U}(sus^{-1})=\Tr_{E|U}(u).
\]

\emph{Proof.} Immediate.

\sgaStableTarget{sga6.e01.rem.8-5}\textbf{Remark 8.5.} If \(\mathcal C\) is a fibered \(S\)-category, denote by \(\mathcal E(\mathcal C)\) the fibered \(S\)-category whose fiber over \(U\in\Ob S\) is the category whose objects are pairs \((L,u)\), with \(L\in\Ob\mathcal C_U\) and \(u\in\operatorname{End}_S(L)\), and whose arrows from \((L,u)\) to \((M,v)\) are the arrows \(f:L\to M\) such that \(fu=vf\). By \hyperlink{sga6.e01.cor.8-4}{(8.4)} one may regard the trace homomorphism as an \(S\)-functor from \(\mathcal E(\Dpx(S)_{\Parf})^{\mathrm{is}}\) to \(\mathcal A\), where the exponent \(\mathrm{is}\) means that only isomorphisms are considered, and where \(\mathcal A\) is the \(S\)-category with discrete fibers defined by \(A\) (cf. (\hyperlink{sga6.e01.rem.6-6}{6.6} a)). Since every perfect complex is locally of finite Tor-dimension, the forgetful \(S\)-functor
\[
  \mathcal E(\Dpx(S)_{\Parf})^{\mathrm{is}}
  \longrightarrow
  \mathcal E(\Dpx(S)_{\parf})^{\mathrm{is}}
\]
induces an isomorphism on the associated stacks. Consequently, the trace homomorphism extends uniquely to an \(S\)-functor from \(\mathcal E(\Dpx(S)_{\Parf})^{\mathrm{is}}\) to \(\mathcal A\).

For \(E\in\Ob\Parf(U)\), \(U\in\Ob S\), we again write
\[
  \Tr_E:\Hom(E,E)\longrightarrow H^0(U,A)
\]
for the corresponding homomorphism, which is an \(H^0(U,A)\)-module homomorphism. One sees in the same way that the definition of the cup-product \hyperlink{sga6.e01.para.8-2}{(8.2)}, as well as assertions \hyperlink{sga6.e01.prop.8-3}{(8.3)} and \hyperlink{sga6.e01.cor.8-4}{(8.4)}, generalize to perfect complexes, not necessarily of finite Tor-dimension.

\sgaStableTarget{sga6.e01.prop.8-6}\textbf{Proposition 8.6.} Let \(U\in\Ob S\), let \(E,F\in\Parf(U)\), and let \(u\in\Hom(E,E)\), \(v\in\Hom(F,F)\); hence \(u\otL v\in\Hom(E\otL F,E\otL F)\). Then
\[
  \Tr_{E\otL F}(u\otL v)=\Tr_E(u)\,\Tr_F(v).
\]

\emph{Proof.} The question being local, by \hyperlink{sga6.e01.cor.8-4}{(8.4)} and \hyperlink{sga6.e01.prop.4-10}{(4.10)} one may suppose that \(u\) and \(v\) are arrows of \(K^b(\mathcal C_0)\). Using \hyperlink{sga6.e01.para.8-1-2}{(8.1.2)}, one is reduced to the case where \(u\) and \(v\) are arrows of \(\mathcal C_0\). One finishes by \hyperlink{sga6.e01.para.8-1-1}{(8.1.1)}.

\sgaStableTarget{sga6.e01.bibliography}\section*{Bibliography}
\begin{description}[leftmargin=2.5cm,style=nextline]
\item[{[H]}] Hartshorne, R. \emph{Residues and Duality}, Lecture Notes no. 20, Springer, 1966.
\item[{[V]}] Verdier, J.-L. \emph{Catégories dérivées}, I.H.E.S., 1963.
\item[{[G]}] Giraud, J. \emph{Cohomologie non abélienne}, to appear in North-Holland Publishing Company.
\item[{SGA 4}] \emph{Cohomologie étale des schémas}, by M. Artin, A. Grothendieck, and J.-L. Verdier, to appear in North-Holland Publishing Company.
\item[{SGA 5}] \emph{Cohomologie \(\ell\)-adique et fonctions L}, by A. Grothendieck, to appear in North-Holland Publishing Company.
\item[{[1]}] Hakim, M. \emph{Topos annelés et schémas relatifs}, thesis (in preparation).
\item[{[2]}] Verdier, J.-L. \emph{Les catégories dérivées des catégories abéliennes}, to appear in North-Holland Publishing Company.
\item[{[MUL]}] Cartan, H. and Eilenberg, S. \emph{Homological Algebra}, Princeton, 1956.
\end{description}



\providecommand{\cS}{\mathcal S}
\providecommand{\cCz}{\mathcal C_0}
\providecommand{\cCone}{\mathcal C_1}
\providecommand{\cA}{\mathcal A}
\providecommand{\cB}{\mathcal B}
\providecommand{\cL}{\mathcal L}
\providecommand{\cI}{\mathcal I}
\providecommand{\cJ}{\mathcal J}
\providecommand{\Db}{\mathrm D}
\providecommand{\Kb}{\mathrm K}
\providecommand{\Dminus}{\mathrm D^{-}}
\providecommand{\Kminus}{\mathrm K^{-}}
\providecommand{\Kbd}{\mathrm K^{b}}
\providecommand{\Mod}{\operatorname{Mod}}
\providecommand{\Lib}{\operatorname{Lib}}
\providecommand{\Loclib}{\operatorname{Loclib}}
\providecommand{\Qcoh}{\operatorname{Qcoh}}
\providecommand{\Coh}{\operatorname{Coh}}
\providecommand{\Ext}{\operatorname{Ext}}
\providecommand{\Hom}{\operatorname{Hom}}
\providecommand{\Ob}{\operatorname{ob}}
\providecommand{\im}{\operatorname{Im}}
\providecommand{\coker}{\operatorname{Coker}}
\providecommand{\coh}{\operatorname{coh}}
\providecommand{\parf}{\operatorname{parf}}
\providecommand{\amp}{\operatorname{amp}}
\providecommand{\RHom}{\operatorname{RHom}}


\sgaStableTarget{sga6.e02.expose}\section*{Exposé II. Existence of global resolutions}
\begin{center}
by L. Illusie
\end{center}

\sgaStableTarget{sga6.e02.sec.1}\section*{1. General globalization criteria}

\sgaStableTarget{sga6.e02.sec.1.0}\subsection*{1.0}
In this section we are given a site $S$, a flat abelian $S$-category $\mathcal C$ (\hyperlink{sga6.e01.sec.1.1}{I 1.1}), and a strictly full additive sub-$S$-category $\mathcal C_0$ (\hyperlink{sga6.e01.sec.1.1}{I 1.1}). We assume that $\mathcal C_0$ is locally quasi-liftable in $\mathcal C$ (\hyperlink{sga6.e01.sec.1.2}{I 1.2}) and locally quasi-stable under kernels of epimorphisms (\hyperlink{sga6.e01.sec.1.2}{I 1.2}).

\sgaStableTarget{sga6.e02.prop.1-1}\textbf{Proposition 1.1.} Let $S\in \Ob(S)$, and suppose that:
\begin{enumerate}[label=(\roman*)]
\item $\mathcal C_{0S}$ is quasi-liftable in $\mathcal C_S$ (\hyperlink{sga6.e01.sec.1.2}{I 1.2}), where $\mathcal C_S$ is regarded as a category fibered over a punctual site;
\item every object of $\mathcal C_S$ which is of finite $\mathcal C_0$-type (\hyperlink{sga6.e01.sec.1.2}{I 1.2}), respectively of finite $\mathcal C_0$-presentation (\hyperlink{sga6.e01.def.2-8}{I 2.8}), is of finite $\mathcal C_{0S}$-type.
\end{enumerate}
Then:

\begin{enumerate}[label=\alph*)]
\item $\mathcal C_{0S}$ is quasi-stable under kernels of epimorphisms (\hyperlink{sga6.e01.sec.1.2}{I 1.2}), where $\mathcal C_{0S}$ is regarded as a category fibered over a punctual site.
\item Let $E\in\Ob \mathrm D^{-}(\mathcal C_S)$ and $n\in\mathbb Z$. If $E$ is $n$-$\mathcal C_0$-pseudo-coherent (\hyperlink{sga6.e01.prop.2-2}{I 2.2}), then $E$ is $n$-$\mathcal C_{0S}$-pseudo-coherent, respectively $(n+1)$-$\mathcal C_{0S}$-pseudo-coherent. In other words,
\[
\mathrm D^{-}(\mathcal C_S)_{n-\mathcal C_0\text{-coh}}
=
\mathrm D^{-}(\mathcal C_S)_{n-\mathcal C_{0S}\text{-coh}},
\]
and, respectively,
\[
\mathrm D^{-}(\mathcal C_S)_{n-\mathcal C_0\text{-coh}}
\subset
\mathrm D^{-}(\mathcal C_S)_{(n+1)-\mathcal C_{0S}\text{-coh}}
\subset
\mathrm D^{-}(\mathcal C_S)_{(n+1)-\mathcal C_0\text{-coh}},
\]
inclusions of strictly full subcategories of $\mathrm D(\mathcal C_S)$, the inclusion on the right being trivial.
\item For an object $E\in\Ob\mathrm D(\mathcal C_S)$ to be pseudo-coherent relative to $\mathcal C_0$ (\hyperlink{sga6.e01.def.2-3}{I 2.3}), it is necessary and sufficient that it be pseudo-coherent relative to $\mathcal C_{0S}$; equivalently,
\[
\mathrm D^{-}(\mathcal C_S)_{\mathcal C_0\text{-coh}}
=
\mathrm D^{-}(\mathcal C_S)_{\mathcal C_{0S}\text{-coh}}.
\]
Moreover, the canonical functor
\[
\mathrm K^{-}(\mathcal C_{0S})/\mathrm K^{-,\delta}(\mathcal C_{0S})
\longrightarrow \mathrm D(\mathcal C_S),
\]
where $\mathrm K^{-,\delta}(\mathcal C_{0S})$ denotes the thick subcategory of $\mathrm K^{-}(\mathcal C_{0S})$ formed by acyclic complexes, is fully faithful and has as essential image $\mathrm D^{-}(\mathcal C_S)_{\mathcal C_0\text{-coh}}$. If $\mathcal C_{0S}$ is liftable in $\mathcal C_S$ (\hyperlink{sga6.e01.sec.1.2}{I 1.2}), that is, if the objects of $\mathcal C_{0S}$ are projective, then $\mathrm K^{-,\delta}(\mathcal C_{0S})$ is reduced to zero objects, and consequently the functor
\[
\mathrm K^{-}(\mathcal C_{0S})\longrightarrow \mathrm D(\mathcal C_S)
\]
is fully faithful and has essential image $\mathrm D^{-}(\mathcal C_S)_{\mathcal C_0\text{-coh}}$.
\end{enumerate}

\emph{Proof.} a) Let
\[
0\longrightarrow L\longrightarrow L^0\longrightarrow\cdots\longrightarrow L^m\longrightarrow0
\]
be an exact sequence in $\mathcal C_S$, with $L^i\in\Ob\mathcal C_{0S}$ for all $i$. It is enough to prove that $L$ is of finite $\mathcal C_{0S}$-type. By (\hyperlink{sga6.e01.prop.2-2}{I 2.2}), $L$ is $\mathcal C_0$-pseudo-coherent, hence, by (\hyperlink{sga6.e01.prop.2-9}{I 2.9}), of finite $\mathcal C_0$-presentation; the assertion follows from (ii).

b) We prove by descending induction on $i$ that $E$ is $i$-$\mathcal C_{0S}$-pseudo-coherent for $i\ge n$, respectively for $i\ge n+1$. One may already suppose that $E$ is strictly $(n+1)$-$\mathcal C_{0S}$-pseudo-coherent, respectively strictly $(n+2)$-$\mathcal C_{0S}$-pseudo-coherent (\hyperlink{sga6.e01.def.2-1}{I 2.1}); the point is then to prove that $E$ is $n$-$\mathcal C_{0S}$-pseudo-coherent, respectively $(n+1)$-$\mathcal C_{0S}$-pseudo-coherent. Thus there is an exact sequence
\[
0\longrightarrow E'\longrightarrow E\longrightarrow E''\longrightarrow0,
\]
where $E'$ is strictly $\mathcal C_{0S}$-perfect (\hyperlink{sga6.e01.def.2-1}{I 2.1}) and $E''{}^i=0$ for $i>n$, respectively for $i>n+1$. Replacing $E$ by $E''$, one may therefore suppose (\hyperlink{sga6.e01.schol.2-6}{I 2.6}) that $E^i=0$ for $i>n$, respectively for $i>n+1$. Then, by (\hyperlink{sga6.e01.prop.2-10}{I 2.10} a)), $H^n(E)$, respectively $H^{n+1}(E)$, is of finite $\mathcal C_0$-type, respectively of finite $\mathcal C_0$-presentation. Hence, by (ii), $H^n(E)$, respectively $H^{n+1}(E)$, is of finite $\mathcal C_{0S}$-type. Applying (\hyperlink{sga6.e01.prop.2-10}{I 2.10} b)) again, one obtains that $E$ is $n$-$\mathcal C_{0S}$-pseudo-coherent, respectively $(n+1)$-$\mathcal C_{0S}$-pseudo-coherent. The hypothesis (i) and part a) have been used implicitly in the references to (\hyperlink{sga6.e01.schol.2-6}{I 2.6} and \hyperlink{sga6.e01.prop.2-10}{2.10}).

c) The first assertion follows from b), the second from (\hyperlink{sga6.e01.prop.2-7}{I 2.7}), and the last is evident.

\sgaStableTarget{sga6.e02.prop.1-2}\textbf{Proposition 1.2.} Suppose that $\mathcal C_0$ is locally liftable in $\mathcal C$ (\hyperlink{sga6.e01.sec.1.2}{I 1.2}) and locally stable under kernels of epimorphisms (\hyperlink{sga6.e01.sec.1.2}{I 1.2}), and make the hypotheses (\hyperlink{sga6.e02.prop.1-1}{1.1 (i)} and (ii), respectively). Let $\mathcal C_{1S}$ denote the strictly full subcategory of $\mathcal C_S$ formed by the objects which are locally in $\mathcal C_0$. Then:
\begin{enumerate}[label=\alph*)]
\item $\mathcal C_{1S}$ is stable under kernels of epimorphisms (\hyperlink{sga6.e01.sec.1.2}{I 1.2}) and is quasi-liftable in $\mathcal C_S$ (\hyperlink{sga6.e01.sec.1.2}{I 1.2}), as a category fibered over a punctual site.
\item Let $E\in\Ob\mathrm D(\mathcal C_S)$ and let $[m,n]$ be an interval of $\mathbb Z$. For $E$ to have $\mathcal C_0$-perfect amplitude contained in $[m,n]$, it is necessary and sufficient that $E$ be isomorphic, in $\mathrm D(\mathcal C_S)$, to a complex $F$ such that $F^i=0$ for $i\notin[m,n]$, $F^i\in\Ob\mathcal C_{0S}$ for $i>m$, and $F^m\in\Ob\mathcal C_{1S}$.
\item The canonical functor
\[
\mathrm K^b(\mathcal C_{1S})/\mathrm K^{b,\delta}(\mathcal C_{1S})
\longrightarrow \mathrm D(\mathcal C_S)
\]
is fully faithful and has as essential image the full subcategory of $\mathrm D(\mathcal C_S)$ formed by complexes of finite $\mathcal C_0$-perfect amplitude.
\item If $\mathcal C_{1S}$ is not merely quasi-liftable but liftable in $\mathcal C_S$, one may replace, in b), the expression ``$E$ is isomorphic to $F$ in $\mathrm D(\mathcal C_S)$'' by ``there exists a quasi-isomorphism $F\to E$''; and, in c), $\mathrm K^{b,\delta}(\mathcal C_{1S})$ is reduced to zero objects.
\end{enumerate}

\emph{Proof.} a) It is clear that $\mathcal C_{1S}$ is stable under kernels of epimorphisms. On the other hand, the objects of $\mathcal C_{1S}$ are $\mathcal C_0$-pseudo-coherent, hence $\mathcal C_{0S}$-pseudo-coherent by (\hyperlink{sga6.e02.prop.1-1}{1.1} c)); the second assertion follows from (\hyperlink{sga6.e01.cor.2-14}{I 2.14}), taking $S$ to be the punctual site.

b) It suffices to prove necessity. By (\hyperlink{sga6.e02.prop.1-1}{1.1} c)), $E$ is already $\mathcal C_{0S}$-pseudo-coherent. Since $E$ is acyclic in degree $>n$, one may, after replacing $E$ by an isomorphic complex in $\mathrm D(\mathcal C_S)$, suppose (\hyperlink{sga6.e01.prop.2-10}{I 2.10}) that $E$ is strictly $(m+1)$-$\mathcal C_{0S}$-pseudo-coherent and that $E^i=0$ for $i>n$. Since $E$ is acyclic in degree $<m$, one may also replace $E$ by its truncation
\[
0\longrightarrow E^m/B^m\longrightarrow E^{m+1}\longrightarrow\cdots\longrightarrow E^n\longrightarrow0,
\]
so that $E^i=0$ for $i<m$. Then, by (\hyperlink{sga6.e01.lem.4-16}{I 4.16}), $E^m$ is locally an object of $\mathcal C_0$, and therefore an object of $\mathcal C_{1S}$.

c) In view of (\hyperlink{sga6.e02.prop.1-2}{1.2} a)), the assertion follows from (\hyperlink{sga6.e02.prop.1-2}{1.2} b)) and (\hyperlink{sga6.e01.rem.4-12-2}{I 4.12.2}). d) The fact that $\mathrm K^{b,\delta}(\mathcal C_{1S})$ is reduced to zero objects is evident. The other assertion follows from (\hyperlink{sga6.e01.cor.4-6}{I 4.6}), applied over the punctual site.

\sgaStableTarget{sga6.e02.sec.2}\section*{2. Application to certain categories of modules}

\sgaStableTarget{sga6.e02.sec.2.0}\subsection*{2.0. Notation}
Let $(S,A)$ be a ringed topos. For $U\in\Ob S$, write $\Mod(U)$ for the category of left $A_U$-modules, and $\Lib(U)$, respectively $\Loclib(U)$, for the full subcategory of $\Mod(U)$ formed by free $A_U$-modules, respectively modules locally direct factors of free modules, of finite type.

\sgaStableTarget{sga6.e02.sec.2.1}\subsection*{2.1. Lifting lemmas}

\sgaStableTarget{sga6.e02.lem.2-1-1}\textbf{Lemma 2.1.1} (cf. EGA $0_{\mathrm I}$, 5.2.3). Let $(S,A)$ be a ringed topos, and let $S\in\Ob S$. Suppose that $S$ is quasi-compact, that is, every covering family $(S_i\to S)$ is dominated by a finite covering family.
\begin{enumerate}[label=\alph*)]
\item Let $F$ be a left $A_S$-module which is the filtered inductive limit, over a filtered category $\Lambda$, of $A_S$-modules $F_\lambda$, and let
\[
u=\varinjlim u_\lambda:F\longrightarrow G
\]
be an epimorphism, where $G$ is an $A_S$-module of finite type. Then there exists $\lambda\in\Lambda$ such that $u_\lambda$ is an epimorphism.
\item Let $F$ be a left $A_S$-module generated by its sections, that is, such that there exists an epimorphism $A_S^{(I)}\to F$ for a suitable set of indices $I$. If $F$ is of finite type, then there exists an epimorphism $A_S^p\to F$ for some $p\in\mathbb N$.
\end{enumerate}
The proof of a) is analogous to the proof of the cited result and is left to the reader; b) follows immediately from a).

\sgaStableTarget{sga6.e02.lem.2-1-2}\textbf{Lemma 2.1.2.} Let $(S,A)$ be a ringed topos, let $\mathcal C$ be the $S$-category fibered in left $A$-modules, let $\mathcal C_0$ be the fibered subcategory of free $A$-modules of finite type (\hyperlink{sga6.e01.sec.1.3}{I 1.3} d)), let $S\in\Ob S$, and let $\mathcal C'_S$ be a strictly full additive subcategory of $\mathcal C_S$ containing $\mathcal C_{0S}$ and stable under kernels, cokernels, and extensions. Suppose that
\[
H^1(S,F)=0\quad\text{for every }F\in\Ob\mathcal C'_S.
\]
Then:
\begin{enumerate}[label=\alph*)]
\item Every exact sequence in $\mathcal C_S$
\[
0\longrightarrow E\longrightarrow F\longrightarrow G\longrightarrow0,
\]
such that $\operatorname{Hom}(G,E)\in\Ob\mathcal C'_S$ and such that $G$ is locally a direct factor of a free module of finite type, is split. The category $\mathcal C_{0S}$ is liftable in $\mathcal C'_S$ (\hyperlink{sga6.e01.prop.2-2}{I 2.2}).
\item Let $\Mod(A(S))$ denote the category of $A(S)$-modules. The functor
\[
\Gamma_S:\mathcal C'_S\longrightarrow\Mod(A(S))
\]
is exact. Moreover:
\begin{enumerate}[label=(\roman*)]
\item if every object $F$ of $\mathcal C'_S$ admits a finite global presentation
\[
A_S^q\longrightarrow A_S^p\longrightarrow F\longrightarrow0,
\]
then $\Gamma_S$ is fully faithful and has as essential image the full subcategory of $\Mod(A(S))$ formed by $A(S)$-modules of finite presentation;
\item if $\mathcal C'_S$ contains the free $A_S$-modules, if every object $F$ of $\mathcal C'_S$ admits a global presentation
\[
A_S^{(J)}\longrightarrow A_S^{(I)}\longrightarrow F\longrightarrow0,
\]
and if $\Gamma_S$ commutes with filtered inductive limits, then $\Gamma_S$ is an equivalence of categories.
\end{enumerate}
\end{enumerate}

\emph{Proof.} a) Since $G$ is locally a direct factor of a free module of finite type, the spectral sequence
\[
E_2^{pq}=H^p(S,\mathcal{E}xt^q(G,E))\Longrightarrow \Ext^{p+q}(G,E)
\]
degenerates (\hyperlink{sga6.e01.lem.1-3-1}{I 1.3.1}) and gives an isomorphism
\[
H^1(S,\operatorname{Hom}(G,E))\simeq \Ext^1(G,E),
\]
whence the first assertion. The second is a particular case, since for $G\in\Ob\mathcal C_{0S}$ and $E\in\Ob\mathcal C'_S$ one has $\operatorname{Hom}(G,E)\in\Ob\mathcal C'_S$.

b) The first assertion is trivial. Let us prove, in case (i) or (ii), the full faithfulness of $\Gamma_S:\mathcal C'_S\to\Mod(A(S))$. Let $F,G\in\Ob\mathcal C'_S$ and let
\sgaStableTarget{sga6.e02.eq.unnamed}\[
A_S^{(J)}\longrightarrow A_S^{(I)}\longrightarrow F\longrightarrow0
\tag{$*$}
\]
be a global presentation of $F$, with $I$ and $J$ finite in case (i). Applying $\Gamma_S$, which is exact on $\mathcal C'_S$ and, in case (ii), commutes with filtered inductive limits, one obtains an exact sequence
\sgaStableTarget{sga6.e02.eq.unnamed.occ2}\[
A(S)^{(J)}\longrightarrow A(S)^{(I)}\longrightarrow\Gamma_S(F)\longrightarrow0.
\tag{$**$}
\]
Applying the left exact functors $\operatorname{Hom}_{\mathcal C_S}(-,G)$ and $\operatorname{Hom}_{A(S)}(-,\Gamma_S(G))$ to $(*)$ and $(**)$ gives a commutative diagram with exact rows:
\[
\begin{tikzcd}[column sep=small]
0 \arrow[r] & \operatorname{Hom}_{\mathcal C_S}(F,G) \arrow[r] \arrow[d] &
\operatorname{Hom}_{\mathcal C_S}(A_S^{(I)},G) \arrow[r] \arrow[d,"\sim"] &
\operatorname{Hom}_{\mathcal C_S}(A_S^{(J)},G) \arrow[d,"\sim"] \\
0 \arrow[r] & \operatorname{Hom}_{A(S)}(\Gamma_SF,\Gamma_SG) \arrow[r] &
\operatorname{Hom}_{A(S)}(A(S)^{(I)},\Gamma_SG) \arrow[r] &
\operatorname{Hom}_{A(S)}(A(S)^{(J)},\Gamma_SG).
\end{tikzcd}
\]
The two vertical arrows on the right are visibly isomorphisms, hence so is the arrow on the left. The description of the essential image in cases (i) and (ii) is immediate and is left to the reader.

\sgaStableTarget{sga6.e02.lem.2-1-3}\textbf{Lemma 2.1.3.} Let $\mathcal A$ and $\mathcal B$ be abelian categories, and let $\varphi:\mathcal A\to\mathcal B$ be an exact functor.
\begin{enumerate}[label=\alph*)]
\item For $M\in\Ob\mathcal A$, the following conditions are equivalent:
\begin{enumerate}[label=(\roman*)]
\item for every $L\in\Ob\mathcal A$, the canonical map induced by $\varphi$
\[
\Ext^n_{\mathcal A}(L,M)\longrightarrow\Ext^n_{\mathcal B}(\varphi L,\varphi M)
\]
is an isomorphism for all $n$;
\item for every $L\in\Ob\mathcal A$, the canonical map
\[
\operatorname{Hom}_{\mathcal A}(L,M)\longrightarrow\operatorname{Hom}_{\mathcal B}(\varphi L,\varphi M)
\]
is an isomorphism, and for every $n\ge1$ and every $x\in\Ext^n_{\mathcal B}(\varphi L,\varphi M)$ there exists an epimorphism $L'\to L$ in $\mathcal A$ which kills $x$.
\end{enumerate}
\item The following conditions are equivalent:
\begin{enumerate}[label=(\roman*)]
\item for all $L,M\in\Ob\mathcal A$ and all $n\in\mathbb Z$, the canonical map
\[
\Ext^n_{\mathcal A}(L,M)\longrightarrow\Ext^n_{\mathcal B}(\varphi L,\varphi M)
\]
is an isomorphism;
\item the canonical functor induced by $\varphi$
\[
\mathrm D^b(\mathcal A)\longrightarrow\mathrm D^b(\mathcal B)
\]
is fully faithful and has as essential image the full subcategory of $\mathrm D^b(\mathcal B)$ formed by complexes $F$ such that $H^i(F)$ is in the essential image of $\varphi$ for every $i$.
\end{enumerate}
\end{enumerate}

\emph{Proof.} For (i)$\Rightarrow$(ii) in a), it is enough to show that, if $n\ge1$, $L\in\Ob\mathcal A$, and $y\in\Ext^n_{\mathcal A}(L,M)$, then there exists an epimorphism $L'\to L$ killing $y$. Using Yoneda's interpretation ([18], III 3.2), $y$ is the class of an exact sequence
\[
0\longrightarrow M\longrightarrow E_{n-1}\longrightarrow\cdots\longrightarrow E_0\longrightarrow L\longrightarrow0,
\]
and the epimorphism $E_0\to L$ kills $y$ ([18], III 3.2.5 and 3.2.8). Conversely, put
\[
F^n(L)=\Ext^n_{\mathcal A}(L,M),\qquad
G^n(L)=\Ext^n_{\mathcal B}(\varphi L,\varphi M).
\]
The effaceability hypothesis says, for $n\ge1$ and $L\in\Ob\mathcal A$, that
\[
G^n(L)=\varinjlim\operatorname{Coker}igl(G^{n-1}(L_0)\to G^{n-1}(L_1)\bigr),
\]
the limit being taken over the exact sequences $L_1\to L_0\to L\to0$. As above, one has the same formula for $F^n(L)$. Since $F^0\to G^0$ is an isomorphism, induction on $n$ shows that the morphism $F^n\to G^n$ induced by $\varphi$ is an isomorphism for every $n$. The equivalence in b) follows by an immediate dévissage, left to the reader.

\sgaStableTarget{sga6.e02.rem.2-1-3-1}\textbf{Remarks 2.1.3.1.} There is of course a dual version of a), which we leave to the reader to state. On the other hand, if $\mathcal A$ and $\mathcal B$ have enough injectives, the conditions (i) and (ii) of b) are equivalent to:
\begin{enumerate}[label=(\roman*),start=3]
\item for $L\in\Ob\mathrm D^{-}(\mathcal A)$ and $M\in\Ob\mathrm D^{+}(\mathcal A)$, the canonical map
\[
\mathrm R\!\operatorname{Hom}(L,M)\longrightarrow
\mathrm R\!\operatorname{Hom}(\varphi L,\varphi M)
\]
is an isomorphism,
\end{enumerate}
as is seen at once by a spectral sequence argument.

\sgaStableTarget{sga6.e02.sec.2.2}\subsection*{2.2. Noetherian schemes and divisorial schemes}

\sgaStableTarget{sga6.e02.notation.2-2-1}\textbf{Notation 2.2.1.} Let $S$ be a scheme. We write $\Qcoh(S)$, respectively $\Coh(S)$, for the abelian category of quasi-coherent, respectively coherent, $\mathcal O_S$-modules. Let $[m,n]$ be an interval of $\mathbb Z$. We shall say that an object $E$ of $\mathrm D(\Qcoh(S))$ is $n$-pseudo-coherent, respectively pseudo-coherent, respectively of perfect amplitude contained in $[m,n]$, respectively perfect, if it is so relative to the category fibered over the Zariski site of $S$ defined by
\[
U\longmapsto \Loclib(U)
\]
((\hyperlink{sga6.e01.def.2-3}{I 2.3}), (\hyperlink{sga6.e01.def.4-7}{I 4.7}), and \hyperlink{sga6.e02.sec.2.0}{(2.0)}). We denote by $\mathrm D^{-}(\Qcoh(S))_{\coh}$, respectively $\mathrm D(\Qcoh(S))_{\parf}$, the full triangulated subcategory of $\mathrm D(\Qcoh(S))$ formed by pseudo-coherent, respectively perfect, complexes.

\sgaStableTarget{sga6.e02.prop.2-2-2}\textbf{Proposition 2.2.2.} Let $S$ be a noetherian scheme. The canonical functor
\[
\mathrm D^{-}(\Coh(S))\longrightarrow \mathrm D(\Qcoh(S))
\]
is fully faithful and has as essential image $\mathrm D^{-}(\Qcoh(S))_{\coh}$.

\emph{Proof.} Apply \hyperlink{sga6.e02.prop.1-1}{(1.1)}, taking for the site the Zariski site of $S$, and for $\mathcal C$, respectively $\mathcal C_0$, the fibered categories $U\mapsto\Qcoh(U)$, respectively $U\mapsto\Coh(U)$. Since $\mathcal O_S$ is coherent, the objects of $\mathcal C_0$ are pseudo-coherent (\hyperlink{sga6.e01.cor.3-5}{I 3.5}); hence an object of $\mathrm D(\Qcoh(U))$ is pseudo-coherent, in the sense of \hyperlink{sga6.e02.notation.2-2-1}{(2.2.1)}, if and only if it is pseudo-coherent relative to $\mathcal C_0$ (\hyperlink{sga6.e01.schol.2-14-1}{I 2.14.1}). By (\hyperlink{sga6.e02.prop.1-1}{1.1} c)), it remains to verify conditions (i) and (ii) of \hyperlink{sga6.e02.prop.1-1}{(1.1)}. For (ii) this is trivial, since a quasi-coherent module which is a quotient of a coherent module is coherent, and coherence is local. For (i), one uses (\hyperlink{sga6.e02.lem.2-1-1}{2.1.1} a)) and the fact that every quasi-coherent $\mathcal O_S$-module is a filtered inductive limit of coherent $\mathcal O_S$-modules (EGA I 9.4.9).

Granting \hyperlink{sga6.e02.prop.3-5}{(3.5)} and \hyperlink{sga6.e02.prop.3-7}{(3.7)}, which are independent of the present section, one easily deduces the following.

\sgaStableTarget{sga6.e02.cor.2-2-2-1}\textbf{Corollary 2.2.2.1.} Let $S$ be a noetherian scheme, respectively a noetherian scheme of finite dimension. The canonical functor
\[
\mathrm D^b(\Coh(S))\longrightarrow \mathrm D(S)
\]
respectively
\[
\mathrm D^{-}(\Coh(S))\longrightarrow \mathrm D(S)
\]
is fully faithful and has as essential image $\mathrm D^b(S)_{\coh}$, respectively $\mathrm D^{-}(S)_{\coh}$ (\hyperlink{sga6.e01.ex.2-15}{I 2.15}).

\sgaStableTarget{sga6.e02.prop.2-2-3}\textbf{Proposition 2.2.3} (cf. EGA II 4.5.2 and 4.5.5). Let $S$ be a quasi-compact and quasi-separated scheme, and let $(L_i)_{i\in I}$ be a family of invertible sheaves on $S$. For every $\mathcal O_S$-module $F$ and every pair $(i,n)\in I\times\mathbb Z$, let $F_{(i,n)}$ denote the sub-$\mathcal O_S$-module of $F\otimes L_i^{\otimes n}$ generated by the sections of $F\otimes L_i^{\otimes n}$ over $S$. The following conditions are equivalent:
\begin{enumerate}[label=(\roman*)]
\item the open sets $S_f$, for $f\in\Gamma(S,L_i^{\otimes n})$, $i\in I$ and $n\in\mathbb N$, form a base for the topology of $S$;
\item those of the $S_f$ which are affine cover $S$;
\item for every quasi-coherent $\mathcal O_S$-module $F$ of finite type, there exists a finite family $(i_\alpha,n_\alpha)$ with $n_\alpha\ge0$ and an epimorphism
\[
\bigoplus_\alpha L_{i_\alpha}^{\otimes(-n_\alpha)}\longrightarrow F;
\]
\item for every $F\in\Ob\Qcoh(S)$, one has
\[
F=\sum_{(i,n)\in I\times\mathbb N}F_{(i,n)}\otimes L_i^{\otimes(-n)},
\]
a sum, not necessarily direct, of sub-sheaves of $F$;
\item same condition as (iv), but with $F$ a quasi-coherent ideal of $\mathcal O_S$.
\end{enumerate}

\emph{Proof.} This is a simple paraphrase of the cited passages; we recall the argument briefly for the reader's convenience. For (i)$\Rightarrow$(ii), choose for $x\in S$ an affine open $U$ containing $x$. An open $S_f$ with $x\in S_f\subset U$, where $f\in\Gamma(S,L_i^{\otimes n})$ and $n\ge0$, is necessarily affine, by the following lemma.

\sgaStableTarget{sga6.e02.lem.2-2-3-1}\textbf{Lemma 2.2.3.1} (EGA II 5.5.8). Let $X$ be a scheme, $L$ an invertible $\mathcal O_X$-module, and $f\in\Gamma(X,L)$. The canonical immersion $X_f\to X$ is affine.

\emph{Proof.} This is trivial: the question is local on $X$, so one may suppose $X$ affine and $L\simeq\mathcal O_X$.

For (ii)$\Rightarrow$(iii), choose a finite covering of $S$ by affine opens $S_{f_{ij}}$, with $f_{ij}\in\Gamma(S,L_i^{\otimes m_{ij}})$. Replacing the $f_{ij}$ by suitable powers, one may suppose that the $m_{ij}$ with fixed $i$ are all equal to an integer $m_i$. If $F$ is quasi-coherent of finite type, then, for each $i,j$, there is an integer $q_i\ge0$ such that every section of $F$ over $S_{f_{ij}}$ extends to a section of $F\otimes L_i^{\otimes q_im_i}$ over $S$. Thus $F\otimes L_i^{\otimes q_im_i}$ is generated over $S_{f_{ij}}$ by finitely many global sections. Proceeding as in the proof of EGA II 4.5.5, one obtains, for each $i$, an integer $n_i\ge0$ such that $F\otimes L_i^{\otimes n_i}$ is generated over the union of the corresponding $S_{f_{ij}}$ by finitely many sections over $S$. Hence there are integers $r_i\ge0$ and a morphism
\[
L_i^{-n_i\, r_i}\longrightarrow F
\]
which induces an epimorphism over that union; the assertion follows.

For (iii)$\Rightarrow$(iv), since every quasi-coherent sheaf on $S$ is the inductive limit of its quasi-coherent subsheaves of finite type (EGA I 9.4.9 and IV 1.7.7), it is enough to treat the case where $F$ is of finite type. There exist morphisms
\[
u_i:L_i^{-n_i\, r_i}\longrightarrow F
\]
whose sum is an epimorphism. But
\[
\operatorname{Im}(u_i)=\operatorname{Im}\bigl(u_i\otimes\operatorname{Id}(L_i^{n_i})\bigr)\otimes L_i^{-n_i},
\]
and $\operatorname{Im}\bigl(u_i\otimes\operatorname{Id}(L_i^{n_i})\bigr)$ is a subsheaf of $F_{(i,n_i)}$, which gives the assertion. The implication (iv)$\Rightarrow$(v) is trivial. Finally, for (v)$\Rightarrow$(i), let $x\in S$, let $U$ be an open containing $x$, and let $F$ be a quasi-coherent ideal defining $S-U$. By (v), there are $(i,n)\in I\times\mathbb N$ and a section $f$ of $F\otimes L_i^{\otimes n}$ over $S$ with $f(x)\ne0$. Then $x\in S_f\subset U$, proving (i).

\sgaStableTarget{sga6.e02.def.2-2-4}\textbf{Definition 2.2.4.} Let $S$ be a quasi-compact and quasi-separated scheme, and let $(L_i)_{i\in I}$ be a family of invertible sheaves on $S$. We say that the $L_i$ form an \emph{ample family} if the equivalent conditions of \hyperlink{sga6.e02.prop.2-2-3}{(2.2.3)} hold. In particular, a family consisting of a single sheaf $L$ is ample if and only if $L$ is ample in the usual sense (EGA II 4.5.3).

\sgaStableTarget{sga6.e02.def.2-2-5}\textbf{Definition 2.2.5.} A scheme $S$ is called \emph{divisorial} if $S$ is quasi-compact and quasi-separated and if there exists an ample family \hyperlink{sga6.e02.def.2-2-4}{(2.2.4)} of invertible $\mathcal O_S$-modules.

\sgaStableTarget{sga6.e02.lem.2-2-6}\textbf{Lemma 2.2.6.} Let $X$ be a locally noetherian scheme. Let $U$ be an open of $X$ such that the inclusion $U\to X$ is affine, for example $U$ affine and $X$ separated. Then $X-U$ is purely of codimension $0$ or $1$; that is, at every maximal point $x$ of $X-U$ one has $\dim\mathcal O_{X,x}=0$ or $1$. It is purely of codimension $1$ if $U$ contains the maximal points of $X$.

\emph{Proof.} Put $Y=X-U$. Let $x$ be a maximal point of $Y$ and make the base change $X'=\operatorname{Spec}\mathcal O_{X,x}\to X$; one obtains a cartesian diagram
\[
\begin{tikzcd}
U' \arrow[r] \arrow[d] & U \arrow[d] \\
X' \arrow[r] & X .
\end{tikzcd}
\]
Let $Y'=X'-U'$. Since $x$ is a maximal point of $Y$, $Y'$ is just the closed point of $X'$. Since $X'$ is affine, $U'$ is an affine open of $X'$. We show that the dimension $d$ of $X'$ is $0$ or $1$. Indeed, one knows (SGA 2 V 5.7) that $H^d_{Y'}(X',\mathcal O_{X'})\ne0$. But if $d\ge2$, then
\[
H^d_{Y'}(X',\mathcal O_{X'})=H^{d-1}(U',\mathcal O_{X'}),
\]
and the latter group is zero because $U'$ is affine and $d-1\ge1$, a contradiction.

\sgaStableTarget{sga6.e02.rem.2-2-6-1}\textbf{Remark 2.2.6.1.} Without the separation hypothesis on $X$, \hyperlink{sga6.e02.lem.2-2-6}{(2.2.6)} would be false for an affine open $U$ of $X$, as shown by the example of the affine plane with the origin doubled (EGA I 5.5.11).

\sgaStableTarget{sga6.e02.prop.2-2-7}\textbf{Proposition 2.2.7.} Let $S$ be a noetherian separated locally factorial scheme, that is, a scheme whose local rings are factorial. Then $S$ is divisorial (2.2.5). In particular (EGA IV 21.11.1):

\sgaStableTarget{sga6.e02.cor.2-2-7-1}\textbf{Corollary 2.2.7.1.} Every noetherian regular separated scheme is divisorial.

\emph{Proof of 2.2.7.} One may suppose that $S$ is integral. Let $(S_i)_{i\in I}$ be a finite covering of $S$ by nonempty affine opens. By (2.2.6), $S-S_i$ is purely of codimension $1$; hence, by EGA IV 21.7.4, it is the support of a positive divisor. In other words, for each $i\in I$, there is an invertible sheaf $L_i$ and a section $f_i\in\Gamma(S,L_i)$ such that $S_i=S_{f_i}$. The family $(L_i)$ is then ample by \hyperlink{sga6.e02.def.2-2-4}{(2.2.4)}, using \hyperlink{sga6.e02.prop.2-2-3}{(2.2.3 (ii))}.

\sgaStableTarget{sga6.e02.rem.2-2-7-2}\textbf{Remark 2.2.7.2.} As in \hyperlink{sga6.e02.lem.2-2-6}{(2.2.6)}, the separation hypothesis on $S$ is essential. Let $S$ again be the affine plane with the origin doubled. Recall that $S$ is obtained by gluing two copies of the affine plane $\mathbb A^2_k$ along the identity on the complement $U$ of the origin $O$; hence there is a canonical projection
\[
p:S\longrightarrow\mathbb A^2_k
\]
which is an isomorphism over $U$, while $p^{-1}(O)$ consists of two points $O_1$ and $O_2$. An invertible sheaf on $S$ is obtained by gluing two invertible sheaves on $\mathbb A^2_k$ by an isomorphism over $U$. Since every invertible sheaf on $\mathbb A^2_k$ is trivial, an invertible sheaf on $S$ is defined by an element $g\in\Gamma(U,\mathcal O_U^*)$. But $g$ is necessarily constant, since the restriction $\Gamma(\mathbb A^2_k,\mathcal O_{\mathbb A^2})\to\Gamma(U,\mathcal O_U)$ is an isomorphism, the origin being of depth $2$ in the plane. Thus every invertible sheaf on $S$ is trivial. It is immediate, on the other hand, that $p$ induces an isomorphism
\[
\Gamma(\mathbb A^2_k,\mathcal O_{\mathbb A^2})\xrightarrow{\sim}\Gamma(S,\mathcal O_S).
\]
Consequently every global section of $\mathcal O_S$ which vanishes at $O_1$ also vanishes at $O_2$, and the opens $S_f$, for $f\in\Gamma(S,\mathcal O_S)$, do not form a base for the topology of $S$.

\sgaStableTarget{sga6.e02.thm.2-2-8}\textbf{Theorem 2.2.8.} Let $S$ be a divisorial scheme \hyperlink{sga6.e02.def.2-2-5}{(2.2.5)}, and let $(L_i)_{i\in I}$ be an ample family of invertible $\mathcal O_S$-modules. Let $L(S)$ be the strictly full additive subcategory of $\Qcoh(S)$ generated by the $L_i^{\otimes(-n)}$, for $(i,n)\in I\times\mathbb N$. Let $E\in\Ob\mathrm D(\Qcoh(S))$.
\begin{enumerate}[label=\alph*)]
\item The categories $L(S)$ and $\Loclib(S)$ are quasi-liftable in $\Qcoh(S)$ (\hyperlink{sga6.e01.sec.1.2}{I 1.2}) and quasi-stable under kernels of epimorphisms (\hyperlink{sga6.e01.sec.1.2}{I 1.2}).
\item Let $n\in\mathbb Z$. For $E$ to be $n$-pseudo-coherent, respectively pseudo-coherent, in the sense of \hyperlink{sga6.e02.notation.2-2-1}{(2.2.1)}, it is necessary and sufficient that $E$ be so relative to $L(S)$ or to $\Loclib(S)$ (\hyperlink{sga6.e01.def.2-3}{I 2.3}). The canonical functors
\[
\mathrm K^{-}(L(S))/\mathrm K^{-,\delta}(L(S))\longrightarrow \mathrm D(\Qcoh(S)),
\]
\[
\mathrm K^{-}(\Loclib(S))/\mathrm K^{-,\delta}(\Loclib(S))\longrightarrow \mathrm D(\Qcoh(S))
\]
are fully faithful and have as essential image $\mathrm D^{-}(\Qcoh(S))_{\coh}$.
\item Let $[m,n]$ be an interval of $\mathbb Z$. For $E$ to have perfect amplitude contained in $[m,n]$ \hyperlink{sga6.e02.notation.2-2-1}{(2.2.1)}, it is necessary and sufficient that $E$ be isomorphic, in $\mathrm D(\Qcoh(S))$, to a complex $F$ such that $F^i=0$ for $i\notin[m,n]$, $F^i\in\Ob L(S)$ for $i>m$, and $F^m$ is locally free of finite type. The canonical functor
\[
\mathrm K^b(\Loclib(S))/\mathrm K^{b,\delta}(\Loclib(S))\longrightarrow \mathrm D(\Qcoh(S))
\]
is fully faithful and has as essential image $\mathrm D(\Qcoh(S))_{\parf}$.
\end{enumerate}

\emph{Proof.} Let $F\to G$ be an epimorphism in $\Qcoh(S)$, with $G$ of finite type. Since $F$ is the inductive limit of its quasi-coherent subsheaves of finite type (EGA I 9.4.9), there exists, by (\hyperlink{sga6.e02.lem.2-1-1}{2.1.1} a)), a quasi-coherent subsheaf $F'$ of finite type in $F$ such that $F'\to F\to G$ is an epimorphism. By \hyperlink{sga6.e02.prop.2-2-3}{(2.2.3 (iii))}, there is an epimorphism $E\to F'$, with $E\in\Ob L(S)$; hence $E\to F'\to F\to G$ is an epimorphism. This proves that the category of quasi-coherent $\mathcal O_S$-modules of finite type is quasi-liftable in $\Qcoh(S)$, and a fortiori the same is true for $\Loclib(S)$ and for $L(S)$.

For the second assertion of a), apply \hyperlink{sga6.e02.prop.1-1}{(1.1)} on the Zariski site of $S$, with $\mathcal C$ the fibered category $U\mapsto\Qcoh(U)$ and $\mathcal C_0$ the fibered category $U\mapsto L(U)$. Since every object of $L$ is pseudo-coherent, the notions of $n$-pseudo-coherence and pseudo-coherence in the fibered category $\mathrm D(\Qcoh)$ coincide with the same notions relative to $L$ (\hyperlink{sga6.e01.cor.2-14}{I 2.14}). Condition (ii) of \hyperlink{sga6.e02.prop.1-1}{(1.1)} is satisfied by \hyperlink{sga6.e02.prop.2-2-3}{(2.2.3 (iii))}, and thus (\hyperlink{sga6.e02.prop.1-1}{1.1} a)) gives the part of a) concerning $L(S)$. The assertion for $\Loclib(S)$ is proved similarly. Part b) follows immediately from (\hyperlink{sga6.e02.prop.1-1}{1.1} b), c)), taking into account that, if $E\in\Ob\mathrm D(\Qcoh(S))$ is $n$-pseudo-coherent, then $E$ lies in $\mathrm D^{-}(\Qcoh(S))$, since $S$ is quasi-compact. An object of the fibered category $\Qcoh$ is locally in $L$ if and only if it is locally free of finite type; hence $L$ is locally liftable in $\Qcoh$ and locally stable under kernels of epimorphisms. Thus c) follows from \hyperlink{sga6.e02.prop.1-2}{(1.2)}, with $\mathcal C=\Qcoh$, $\mathcal C_0=L$, and $\mathcal C_1=\Loclib(S)$.

\sgaStableTarget{sga6.e02.prop.2-2-9}\textbf{Proposition 2.2.9.} The hypotheses and notation are those of \hyperlink{sga6.e02.thm.2-2-8}{(2.2.8)}, but assume moreover that $S$ is separated or noetherian.
\begin{enumerate}[label=\alph*)]
\item The canonical functors
\[
\mathrm K^{-}(L(S))/\mathrm K^{-,\delta}(L(S))\longrightarrow \mathrm D(S),
\qquad
\mathrm K^{-}(\Loclib(S))/\mathrm K^{-,\delta}(\Loclib(S))\longrightarrow \mathrm D(S)
\]
are fully faithful and have as essential image $\mathrm D^{-}(S)_{\coh}$. Moreover, if $S$ is of finite cohomological dimension in the sense of \hyperlink{sga6.e02.prop.3-7}{(3.7)}, for example if the underlying space of $S$ is noetherian of finite Krull dimension, then the functors
\[
\mathrm K(L(S))/\mathrm K^{\delta}(L(S))\longrightarrow
\mathrm K(\Loclib(S))/\mathrm K^{\delta}(\Loclib(S))\longrightarrow
\mathrm D^{-}(S)_{\coh}
\]
are equivalences of categories.
\item The canonical functor
\[
\mathrm K^b(\Loclib(S))/\mathrm K^{b,\delta}(\Loclib(S))\longrightarrow \mathrm D(S)
\]
is fully faithful and has as essential image $\mathrm D(S)_{\parf}$ (\hyperlink{sga6.e01.notation.4-9}{I 4.9}). For $E\in\Ob\mathrm D(S)$ to have perfect amplitude contained in $[m,n]$ (\hyperlink{sga6.e01.ex.4-8}{I 4.8}), it is necessary and sufficient that $E$ be isomorphic, in $\mathrm D(S)$, to a complex $F$ such that $F^i=0$ for $i\notin[m,n]$, $F^i\in\Ob L(S)$ for $i>m$, and $F^m$ is locally free of finite type.
\end{enumerate}

\emph{Proof.} In view of \hyperlink{sga6.e02.prop.3-5}{(3.5)} and \hyperlink{sga6.e02.prop.3-7}{(3.7)}, this is an immediate corollary of \hyperlink{sga6.e02.thm.2-2-8}{(2.2.8)}. The only point to check, and it follows trivially from \hyperlink{sga6.e02.prop.3-5}{(3.5)}, respectively \hyperlink{sga6.e02.prop.3-7}{(3.7)}, is that for an object $E$ of $\mathrm D^{-}(\Qcoh(S))$, respectively of $\mathrm D(\Qcoh(S))$ when $S$ has finite cohomological dimension in the sense of \hyperlink{sga6.e02.prop.3-7}{(3.7)}, the notions of pseudo-coherence and perfection of \hyperlink{sga6.e02.notation.2-2-1}{(2.2.1)} coincide with the usual notions (\hyperlink{sga6.e01.ex.2-15}{I 2.15} and \hyperlink{sga6.e01.notation.4-9}{4.9}) after regarding $E$ as an object of $\mathrm D(S)$.

\sgaStableTarget{sga6.e02.rem.2-2-10}\textbf{Remark 2.2.10.} a) If $S=\operatorname{Spec}A$ is affine, the family consisting of the single sheaf $\mathcal O_S$ is ample, and for it $L(S)=\Lib(S)$. It is also clear that the categories $\mathrm K^{-,\delta}(\Lib(S))$ and $\mathrm K^{-,\delta}(\Loclib(S))$ are reduced to zero objects. It follows from \hyperlink{sga6.e02.prop.2-2-9}{(2.2.9)} that, writing $\Mod(A)$ for the category of $A$-modules and $\mathrm D(A)=\mathrm D(\Mod(A))$, the functor $R\Gamma_S:\mathrm D(S)\to\mathrm D(A)$ induces fully faithful functors
\[
\mathrm D^{-}(S)_{\coh}\longrightarrow \mathrm D(A),
\qquad
\mathrm D(S)_{\parf}\longrightarrow \mathrm D(A),
\]
whose essential images are respectively the triangulated subcategories $\mathrm D^{-}(A)_{\coh}$ and $\mathrm D(A)_{\parf}$, formed by complexes with bounded-above cohomology and pseudo-coherent, respectively perfect, relative to the subcategory $\Lib(A)$, respectively $\Loclib(A)$, of free, respectively projective, modules of finite type. There are commutative diagrams of equivalences of categories
\[
\begin{tikzcd}
\mathrm K^{-,b}(\Lib(S)) \arrow[r] \arrow[d,"\Gamma_S"'] &
\mathrm K^{-,b}(\Loclib(S)) \arrow[r] \arrow[d,"\Gamma_S"'] &
\mathrm D^b(S)_{\coh} \arrow[d,"R\Gamma_S"] \\
\mathrm K^{-,b}(\Lib(A)) \arrow[r] &
\mathrm K^{-,b}(\Loclib(A)) \arrow[r] &
\mathrm D^b(A)_{\coh},
\end{tikzcd}
\]
\[
\begin{tikzcd}
\mathrm K^b(\Loclib(S)) \arrow[r] \arrow[d,"\Gamma_S"'] &
\mathrm D(S)_{\parf} \arrow[d,"R\Gamma_S"] \\
\mathrm K^b(\Loclib(A)) \arrow[r] &
\mathrm D(A)_{\parf}.
\end{tikzcd}
\]
In the upper diagram, one may replace $b$ by $-$ when $S$ has finite cohomological dimension in the sense of \hyperlink{sga6.e02.prop.3-7}{(3.7)}.

b) For an arbitrary scheme $S$, it follows from a) that a complex of $\mathcal O_S$-modules which is pseudo-coherent and has bounded cohomology is locally isomorphic, in $\mathrm D(S)$, to a strictly pseudo-coherent complex, which was by no means evident a priori (cf. \hyperlink{sga6.e01.rem.2-3-1}{I 2.3.1}).

\sgaStableTarget{sga6.e02.sec.2.3}\subsection*{2.3. Soft sheaves}

\sgaStableTarget{sga6.e02.para.2-3-1}\textbf{2.3.1.} Recall a few definitions and results from [T.F.]. Let $X$ be a paracompact space. A sheaf of sets $F$ on $X$ is called \emph{soft} if, for every closed subset $Y$ of $X$, every section of $F$ over $Y$ extends to a section of $F$ over $X$. Soft abelian sheaves are acyclic: if $F$ is soft, then $H^i(X,F)=0$ for $i>0$ (T.F. II 4.4.3).

Let $A$ be a sheaf of rings on $X$. If $F$ is a soft $A$-module, then $F$ is generated by its sections (cf. \hyperlink{sga6.e02.lem.2-1-1}{2.1.1}): indeed, the homomorphism $A_X^{(I)}\to F$ defined by $I=H^0(X,F)$ is an epimorphism. Moreover, if $A$ is soft, every $A$-module is soft (T.F. II 3.7.1). Examples of soft sheaves of rings include the sheaf of continuous real or complex-valued functions on a paracompact space and the sheaf of $C^\infty$ functions on a $C^\infty$ manifold.

\sgaStableTarget{sga6.e02.prop.2-3-2}\textbf{Proposition 2.3.2.} Let $S$ be a compact space equipped with a soft sheaf of rings $A$. Write $\Mod(S)$ for the category of left $A$-modules and $\mathrm D(S)=\mathrm D(\Mod(S))$.
\begin{enumerate}[label=\alph*)]
\item The category $\Lib(S)$, respectively $\Loclib(S)$ \hyperlink{sga6.e02.sec.2.0}{(2.0)}, is liftable in $\Mod(S)$ and quasi-stable, respectively stable, under kernels of epimorphisms (\hyperlink{sga6.e01.sec.1.2}{I 1.2}). This assertion holds more generally for $S$ paracompact.
\item Let $E\in\Ob\mathrm D(S)$ and $n\in\mathbb Z$. For $E$ to be $n$-pseudo-coherent, respectively pseudo-coherent, it is necessary and sufficient that it be so relative to $\Lib(S)$ or to $\Loclib(S)$. The canonical functors
\[
\mathrm K^{-}(\Lib(S))\longrightarrow \mathrm D(S),
\qquad
\mathrm K^{-}(\Loclib(S))\longrightarrow \mathrm D(S)
\]
are fully faithful and have as essential image $\mathrm D^{-}(S)_{\coh}$ (\hyperlink{sga6.e01.ex.2-15}{I 2.15}).
\item Let $E\in\Ob\mathrm D(S)$ and let $[m,n]$ be an interval of $\mathbb Z$. For $E$ to have perfect amplitude contained in $[m,n]$ (\hyperlink{sga6.e01.ex.4-8}{I 4.8}), it is necessary and sufficient that there exist a quasi-isomorphism $F\to E$, where $F$ is a complex such that $F^i=0$ for $i\notin[m,n]$, $F^i\in\Ob\Lib(S)$ for $i>m$, and $F^m\in\Ob\Loclib(S)$. The canonical functor
\[
\mathrm K^b(\Loclib(S))\longrightarrow \mathrm D(S)
\]
is fully faithful and has as essential image $\mathrm D(S)_{\parf}$ (\hyperlink{sga6.e01.notation.4-9}{I 4.9}).
\item Let $\Mod(A(S))$ be the category of left $A(S)$-modules and $\mathrm D(A(S))$ the corresponding derived category. The functor
\[
H^0(S,-):\Mod(S)\longrightarrow\Mod(A(S))
\]
is exact and is an equivalence of categories. Moreover it defines equivalences of categories
\[
\mathrm D(S)\xrightarrow{\sim}\mathrm D(A(S)),\qquad
\mathrm D^{-}(S)_{\coh}\xrightarrow{\sim}\mathrm D^{-}(A(S))_{\coh},\qquad
\mathrm D(S)_{\parf}\xrightarrow{\sim}\mathrm D(A(S))_{\parf},
\]
where pseudo-coherence, respectively perfection, in $\mathrm D(A(S))$ is understood relative to the category of free, respectively projective, modules of finite type. It also gives an equivalence from the full subcategory of $\mathrm D(S)$ formed by $n$-pseudo-coherent objects, respectively by objects of perfect amplitude contained in $[m,n]$, onto the analogous full subcategory of $\mathrm D(A(S))$.
\end{enumerate}

\emph{Proof.} The first assertion of a) follows from (\hyperlink{sga6.e02.lem.2-1-2}{2.1.2} a)), applied with $S$ the topos of sheaves of sets on $S$ and $\mathcal C'_S=\Mod(S)$, using \hyperlink{sga6.e02.para.2-3-1}{(2.3.1)}. The second follows from (\hyperlink{sga6.e02.prop.1-1}{1.1} a)), respectively (\hyperlink{sga6.e02.prop.1-2}{1.2} a)), with $\mathcal C=\Mod$ and $\mathcal C_0=\Lib$, respectively $\Loclib$; the condition (ii) of \hyperlink{sga6.e02.prop.1-1}{(1.1)} is verified by \hyperlink{sga6.e02.para.2-3-1}{(2.3.1)} and (\hyperlink{sga6.e02.lem.2-1-1}{2.1.1} b)). Part b) follows from (\hyperlink{sga6.e02.prop.1-1}{1.1} b), c)) with the same $\mathcal C$ and $\mathcal C_0$. Part c) follows from (\hyperlink{sga6.e02.prop.1-2}{1.2} b), c), d)), with $\mathcal C=\Mod$ and $\mathcal C_0=\Lib$. The first assertion of d) follows, using \hyperlink{sga6.e02.para.2-3-1}{(2.3.1)}, from (\hyperlink{sga6.e02.lem.2-1-2}{2.1.2} b) (ii)) with $\mathcal C'_S=\mathcal C_S=\Mod(S)$. The other assertions in d) follow from it, together with b) and c), and from the fact that $\Gamma_S$ carries free modules, respectively direct factors of free modules, of finite type to free, respectively projective, modules of finite type.

\sgaStableTarget{sga6.e02.sec.2.4}\subsection*{2.4. Stein spaces}

\sgaStableTarget{sga6.e02.def.2-4-1}\textbf{Definition 2.4.1.} Let $S$ be a locally ringed topos (\hyperlink{sga6.e01.lem.2-15-1}{I 2.15.1 (ii)}), with final object also denoted $S$. We say that $S$ is \emph{Stein} if the structural sheaf $\mathcal O_S$ is coherent (\hyperlink{sga6.e01.def.3-1}{I 3.1}) and if the following properties hold:

\emph{Theorem A:} every coherent $\mathcal O_S$-module (\hyperlink{sga6.e01.def.3-1}{I 3.1}) is generated by its sections;

\emph{Theorem B:} for every coherent $\mathcal O_S$-module $F$, one has $H^i(S,F)=0$ for $i>0$.

\sgaStableTarget{sga6.e02.rem.2-4-1-1}\textbf{Remark 2.4.1.1.} If $S$ has enough closed points $s$ such that the ideal $\mathfrak m_s$ defined by $s$ is coherent, then the relation $H^1(S,F)=0$ for every coherent sheaf $F$ implies Theorem A. Indeed, let $F$ be a coherent $\mathcal O_S$-module and let $s$ be a closed point of $S$. There is an exact sequence
\[
0\longrightarrow \mathfrak m_sF\longrightarrow F\longrightarrow F/\mathfrak m_sF\longrightarrow0,
\]
where $F/\mathfrak m_sF$ is a coherent sheaf concentrated at $s$. Applying $H^0(S,-)$ and writing the cohomology exact sequence, one finds that
\[
H^0(S,F)\longrightarrow H^0(S,F/\mathfrak m_sF)
\]
is an epimorphism. Since
\[
H^0(S,F/\mathfrak m_sF)=F_s/\mathfrak m_sF_s,
\]
where $\mathfrak m_s$ is the maximal ideal of $\mathcal O_s$, Nakayama's lemma gives the conclusion. The editor does not know whether the preceding implication is true in general when one assumes only that $\mathcal O_S$ is a coherent sheaf of local rings. The editor also does not know whether the relation $H^1(S,F)=0$ for every coherent sheaf $F$ implies Theorem B.

\sgaStableTarget{sga6.e02.ex.2-4-2}\textbf{Examples 2.4.2.} A complex analytic Stein space, a compact convex subset of $\mathbb C^n$ equipped with the sheaf induced by the sheaf of holomorphic functions, the spectrum of a noetherian ring, and an affine rigid analytic space define Stein topoi.

\sgaStableTarget{sga6.e02.lem.2-4-3}\textbf{Lemma 2.4.3.} Let $S$ be a Stein topos \hyperlink{sga6.e02.def.2-4-1}{(2.4.1)}. Suppose that the final object $S$ is quasi-compact \hyperlink{sga6.e02.lem.2-1-1}{(2.1.1)}, and let $\Coh(S)$ be the abelian category of coherent $\mathcal O_S$-modules, that is, modules of finite presentation. Then:
\begin{enumerate}[label=\alph*)]
\item $\Lib(S)$, respectively $\Loclib(S)$, is liftable in $\Coh(S)$ and quasi-stable, respectively stable, under kernels of epimorphisms (\hyperlink{sga6.e01.sec.1.2}{I 1.2}).
\item The canonical functors
\[
\mathrm K^{-}(\Lib(S))\longrightarrow \mathrm K^{-}(\Loclib(S))\longrightarrow \mathrm D^{-}(\Coh(S))
\]
are equivalences of categories.
\item Put $A=\Gamma(S,\mathcal O_S)$, and write $\Mod(A)$ for the category of $A$-modules and $\mathrm D(A)$ for the corresponding derived category. The functor
\[
\Gamma_S:\Coh(S)\longrightarrow\Mod(A)
\]
is exact, fully faithful, and has as essential image the strictly full additive subcategory of $\Mod(A)$ formed by $A$-modules of finite presentation. It induces a fully faithful functor
\[
\mathrm D^{-}(\Coh(S))\longrightarrow\mathrm D(A),
\]
whose essential image is the full subcategory $\mathrm D^{-}(A)_{\coh}$ formed by complexes pseudo-coherent relative to the subcategory of free modules of finite type.
\end{enumerate}
Here ``enough closed points'' means enough fiber functors (SGA 4 IV 6.4.1) of the form $i^*:S\to s$, with $i_*$ the inclusion of a closed subtopos of $S$; and $\mathfrak m_s$ denotes the sections of $\mathcal O_S$ whose fiber at $s$ belongs to the maximal ideal of $\mathcal O_s$.

\emph{Proof.} For the first assertion of a), apply (\hyperlink{sga6.e02.lem.2-1-2}{2.1.2} a)) with $\mathcal C'_S=\Coh(S)$; the hypothesis of \hyperlink{sga6.e02.lem.2-1-2}{(2.1.2)} is verified by Theorem B, since if $E$ is coherent and $G$ is locally free of finite type, then $\operatorname{Hom}(G,E)$ is coherent. Thus $\Loclib(S)$, and a fortiori $\Lib(S)$, is liftable in $\Coh(S)$. The second assertion of a) follows from (\hyperlink{sga6.e02.prop.1-1}{1.1} a)), respectively (\hyperlink{sga6.e02.prop.1-2}{1.2} a)), applied with $\mathcal C_0=\Lib$, respectively $\Loclib$; condition (ii) of \hyperlink{sga6.e02.prop.1-1}{(1.1)} is verified by Theorem A and (\hyperlink{sga6.e02.lem.2-1-1}{2.1.1} b)). Part b) follows from (\hyperlink{sga6.e02.prop.1-1}{1.1} c)), applied successively with $\mathcal C_0=\Lib$ and $\mathcal C_0=\Loclib$. The first assertion of c) follows from (\hyperlink{sga6.e02.lem.2-1-2}{2.1.2} b)(i)), applied with $\mathcal C'_S=\Coh(S)$. To prove that $\Gamma_S:\mathrm D^{-}(\Coh(S))\to\mathrm D(A)$ is fully faithful, it is enough, by b), to prove that the composite
\[
\mathrm K^{-}(\Lib(S))\longrightarrow\mathrm D^{-}(\Coh(S))\longrightarrow\mathrm D(A)
\]
is so. But this composite factors as
\[
\mathrm K^{-}(\Lib(S))\xrightarrow{\Gamma_S}\mathrm K^{-}(\Lib(A))\longrightarrow\mathrm D(A),
\]
where $\Lib(A)$ denotes the category of free $A$-modules of finite type; the first arrow is an equivalence of categories, and the second is fully faithful by (\hyperlink{sga6.e01.prop.2-7}{I 2.7} b)), applied with $\mathcal C=\Mod(A)$ and $\mathcal C_0=\Lib(A)$, since $\mathrm K^{-,\delta}(\Lib(A))=0$. This proves the last assertion of c).

\sgaStableTarget{sga6.e02.lem.2-4-4}\textbf{Lemma 2.4.4.} Under the hypotheses of \hyperlink{sga6.e02.lem.2-4-3}{(2.4.3)}, the canonical functor
\[
\mathrm D^b(\Coh(S))\longrightarrow\mathrm D(S)
\]
is fully faithful and has as essential image $\mathrm D^b(S)_{\coh}$ (\hyperlink{sga6.e01.ex.2-15}{I 2.15}).

\emph{Proof.} Apply \hyperlink{sga6.e02.lem.2-1-3}{(2.1.3)} with $\mathcal A=\Coh(S)$, $\mathcal B=\Mod(S)$, and $\varphi:\mathcal A\to\mathcal B$ the inclusion functor. Since this functor is fully faithful, it is enough, by (\hyperlink{sga6.e02.lem.2-1-3}{2.1.3} a)(ii)), to check that for fixed $M\in\Ob\Coh(S)$ and $n>0$ the functor
\[
L\longmapsto \Ext^n_{\Mod(S)}(L,M)
\]
is coeffaceable on $L\in\Ob\Coh(S)$. By Theorem A and (\hyperlink{sga6.e02.lem.2-1-1}{2.1.1} b)), every coherent $\mathcal O_S$-module is a quotient of a free module of finite type; and if $L$ is free of finite type, then $\Ext^n_{\Mod(S)}(L,M)=0$ by Theorem B. The assertion follows. Hence the functor $\mathrm D^b(\Coh(S))\to\mathrm D(S)$ is fully faithful and has as essential image the full subcategory of $\mathrm D(S)$ formed by complexes with bounded coherent cohomology, which is precisely $\mathrm D^b(S)_{\coh}$ since $\mathcal O_S$ is coherent (\hyperlink{sga6.e01.cor.3-5}{I 3.5}).

\sgaStableTarget{sga6.e02.prop.2-4-5}\textbf{Proposition 2.4.5.} With the hypotheses and notation of \hyperlink{sga6.e02.lem.2-4-3}{(2.4.3)}, one has:
\begin{enumerate}[label=\alph*)]
\item The canonical functors
\[
\mathrm K^{-,b}(\Lib(S))\longrightarrow\mathrm K^{-,b}(\Loclib(S))\longrightarrow \mathrm D^b(S)_{\coh}
\]
are equivalences of categories. Moreover, the functor $R\Gamma_S:\mathrm D^b(S)\to\mathrm D^b(A)$ induces a fully faithful functor
\[
R\Gamma_S:\mathrm D^b(S)_{\coh}\longrightarrow \mathrm D^b(A),
\]
whose essential image is $\mathrm D^b(A)_{\coh}$, the triangulated subcategory of $\mathrm D(A)$ formed by complexes with bounded cohomology and pseudo-coherent relative to the category of free modules of finite type.
\item Let $E\in\Ob\mathrm D(S)$, and let $[m,n]$ be an interval of $\mathbb Z$. If $E$ has perfect amplitude contained in $[m,n]$ (\hyperlink{sga6.e01.notation.4-5}{I 4.5}), then $E$ is isomorphic to a complex $F$ such that $F^i=0$ for $i\notin[m,n]$, $F^i\in\Ob\Lib(S)$ for $i>m$, and $F^m\in\Ob\Loclib(S)$.
\item The canonical functor
\[
\mathrm K^b(\Loclib(S))\longrightarrow \mathrm D(S)_{\parf}
\]
is an equivalence of categories (\hyperlink{sga6.e01.notation.4-9}{I 4.9}). Moreover, $R\Gamma_S:\mathrm D^b(S)\to\mathrm D^b(A)$ induces a fully faithful functor
\[
R\Gamma_S:\mathrm D(S)_{\parf}\longrightarrow\mathrm D^b(A),
\]
whose essential image is $\mathrm D^b(A)_{\parf}$, the triangulated subcategory of $\mathrm D(A)$ formed by complexes perfect relative to the category of projective modules of finite type.
\end{enumerate}

\emph{Proof.} a) The functors
\[
\mathrm K^{-,b}(\Lib(S))\longrightarrow\mathrm K^{-,b}(\Loclib(S))\longrightarrow
\mathrm D^b(\Coh(S))\longrightarrow\mathrm D^b(A)_{\coh}
\]
are equivalences of categories by (\hyperlink{sga6.e02.lem.2-4-3}{2.4.3} b), c). Since the fully faithful functor $\Gamma_S:\mathrm D^b(\Coh(S))\to\mathrm D(A)$ factors as
\[
\mathrm D^b(\Coh(S))\longrightarrow\mathrm D^b(S)_{\coh}\xrightarrow{R\Gamma_S}\mathrm D(A),
\]
the conclusion follows from \hyperlink{sga6.e02.lem.2-4-4}{(2.4.4)}.

b) Suppose $\operatorname{parf.amp}(E)\subset[m,n]$. In particular $E$ is pseudo-coherent; by a), after replacing $E$ by an isomorphic object, we may assume that $E^i\in\Ob\Lib(S)$ for $i>m$ and $E^i=0$ for $i>n$. Truncating, we may also suppose that $E^i=0$ for $i<m$. By (\hyperlink{sga6.e01.lem.4-16}{I 4.16}), applied with $\mathcal C=\Mod$ and $\mathcal C_0=\Lib$, one then has $E^m\in\Ob\Loclib(S)$.

c) The hypotheses \hyperlink{sga6.e02.prop.1-1}{(1.1 (i))} and, respectively, \hyperlink{sga6.e02.prop.1-1}{(1.1 (ii))} are verified for $\mathcal C=\Coh(S)$ and $\mathcal C_0=\Lib$, by (\hyperlink{sga6.e02.lem.2-4-3}{2.4.3} a)) and by (\hyperlink{sga6.e02.lem.2-1-1}{2.1.1} b)) together with Theorem A. Hence (\hyperlink{sga6.e02.prop.1-2}{1.2} c), d)) shows that the functor
\[
\mathrm K^b(\Loclib(S))\longrightarrow\mathrm D^b(\Coh(S))
\]
is fully faithful. Therefore, by \hyperlink{sga6.e02.lem.2-4-4}{(2.4.4)}, the same is true for
\[
\mathrm K^b(\Loclib(S))\longrightarrow\mathrm D(S),
\]
and its essential image is $\mathrm D(S)_{\parf}$ by b). By an argument analogous to that used in the proof of (\hyperlink{sga6.e02.lem.2-4-3}{2.4.3} c)), one sees that the functor $\mathrm K^b(\Loclib(S))\to\mathrm D(A)$ induced by $\Gamma_S$ is fully faithful and has essential image $\mathrm D^b(A)_{\parf}$; one finishes the proof as in a).

\sgaStableTarget{sga6.e02.rem.2-4-6}\textbf{Remark 2.4.6.} The editor does not know whether in \hyperlink{sga6.e02.lem.2-4-4}{(2.4.4)} one may replace $\mathrm D^b$ by $\mathrm D^{-}$, and consequently whether (\hyperlink{sga6.e02.prop.2-4-5}{2.4.5} a)) remains true when one replaces $\mathrm K^{-,b}$ by $\mathrm K^{-}$ and $\mathrm D^b$ by $\mathrm D^{-}$.

We finish with an interesting complement to \hyperlink{sga6.e02.prop.2-4-5}{(2.4.5)}. First we give a definition.

\sgaStableTarget{sga6.e02.def.2-4-7}\textbf{Definition 2.4.7.} A \emph{complex pseudo-analytic space} is a ringed space locally isomorphic to a ringed space of the form $(Y,\mathcal O_X|Y)$, where $X$ is a complex analytic space and $Y$ is a closed subset of $X$.

The structural sheaf of a complex pseudo-analytic space is therefore coherent. We say that a complex pseudo-analytic space is Stein if the corresponding topos is Stein \hyperlink{sga6.e02.def.2-4-1}{(2.4.1)}.



% Exposé II continuation macros
\providecommand{\Qcoh}{\operatorname{Qcoh}}
\providecommand{\Mod}{\operatorname{Mod}}
\providecommand{\Coh}{\operatorname{Coh}}
\providecommand{\Ob}{\operatorname{Ob}}
\providecommand{\Hom}{\operatorname{Hom}}
\providecommand{\Ext}{\operatorname{Ext}}
\providecommand{\Spec}{\operatorname{Spec}}
\providecommand{\qcoh}{\mathrm{qcoh}}
\providecommand{\parf}{\mathrm{parf}}
\providecommand{\coh}{\mathrm{coh}}
\providecommand{\Loclib}{\operatorname{Loclib}}
\providecommand{\Lib}{\operatorname{Lib}}
\providecommand{\im}{\operatorname{Im}}

\sgaStableTarget{sga6.e02.expose.occ2}\section*{Exposé II. Global resolutions (continued)}

\sgaStableTarget{sga6.e02.prop.2-4-8}\textbf{Proposition 2.4.8.} Let $S$ be a compact Stein complex pseudo-analytic space \hyperlink{sga6.e02.def.2-4-7}{(2.4.7)}. Suppose that, for every open $U$ of $S$ and every reduced closed subspace $X$ of $U$, the normalization of $X$ is locally connected. This condition is verified, for example, when $S$ is locally defined by real analytic equations or inequalities in a space $\mathbb R^n$. Then
\[
 A=\Gamma(S,\mathcal O_S)
\]
is a noetherian ring.

\emph{Proof.} It is enough to show that, if $M$ is an $A$-module of finite presentation, then the set of submodules of $M$, or equivalently the set of submodules of finite type of $M$, satisfies the maximal condition. Indeed, if $N$ is a submodule of finite type of $M$, then $M/N$ is of finite presentation by \hyperlink{sga6.e01.schol.2-6}{I 2.6} and \hyperlink{sga6.e01.prop.2-9}{2.9}, so that $N$ is itself of finite presentation by the equivalence (\hyperlink{sga6.e02.lem.2-4-3}{2.4.3} c). Using again the equivalence (\hyperlink{sga6.e02.lem.2-4-3}{2.4.3} c) between the category of coherent $\mathcal O_S$-modules and the category of $A$-modules of finite presentation, together with the compactness of $S$, one is reduced to proving the following lemma.

\sgaStableTarget{sga6.e02.lem.2-4-8-1}\textbf{Lemma 2.4.8.1.} Let $S$ be a complex pseudo-analytic space satisfying the local connectedness condition of \hyperlink{sga6.e02.prop.2-4-8}{(2.4.8)}. Let $F$ be a coherent $\mathcal O_S$-module, let $(E_n)$ be an increasing sequence of coherent sub-$\mathcal O_S$-modules of $F$, and let $s$ be a point of $S$. Then there exists an open neighbourhood $U$ of $s$ such that the sequence $(E_n|U)$ is stationary.

\emph{Proof.} Since the ring $\mathcal O_{S,s}$ is noetherian, there exists an integer $n_0$ such that $(E_n)_s=(E_{n_0})_s$ for $n\geq n_0$. Replacing the sequence $(E_n)$ by the sequence $(E_n/E_{n_0})$, one may therefore suppose that $(E_n)_s=0$ for every $n$. On the other hand $F_s=M$ is a finite $\mathcal O_{S,s}$-module; hence it has a finite filtration
\[
 M=M_0\supset M_1\supset\cdots\supset M_q=0
\]
such that
\[
 M_i/M_{i+1}\simeq \mathcal O_s/\mathfrak p_i,
\]
where $\mathfrak p_i$ is a prime ideal of $\mathcal O_s$. Since the category of finite-type $\mathcal O_s$-modules is equivalent to the category of germs at $s$ of coherent $\mathcal O_S$-modules, after shrinking around $s$ one may suppose that the $M_i$, respectively the $\mathfrak p_i$, are the fibres at $s$ of coherent sub-$\mathcal O_S$-modules $F_i$, respectively coherent ideals $\mathfrak p_i$, of $F$, respectively of $\mathcal O_S$, and that one has
\[
 F=F_0\supset F_1\supset\cdots\supset F_q=0,
 \qquad F_i/F_{i+1}\simeq\mathcal O_S/\mathfrak p_i.
\]
To prove that the sequence $(E_n)$ is stationary in a neighbourhood of $s$ is therefore equivalent to proving the same assertion for the induced sequences on the graded pieces of this filtration. Thus one may suppose that $F=\mathcal O_S$ and that $\mathcal O_{S,s}$ is an integral domain; the $E_n$ then form an increasing sequence of ideals of $\mathcal O_S$ with $(E_n)_s=0$.

Since $\mathcal O_{S,s}$ is an integral domain, hence in particular reduced, one may moreover suppose, by Oka's theorem, after localizing at $s$, that $S$ is reduced. We have to show that there is an open neighbourhood $U$ of $s$ such that $E_n|U=0$ for $n$ large enough. Consider, to do this, the normalization $\widetilde S$ of $S$. It is the analytic spectrum of a finite $\mathcal O_S$-algebra $\widetilde{\mathcal O}_S$ whose fibre at each $t\in S$ is the integral closure of $\mathcal O_{S,t}$. Since $E_n\subset E_n\widetilde{\mathcal O}_S$, it is enough to prove that $E_n\widetilde{\mathcal O}_S=0$ for $n$ large. This reduces us to the case where $S$ is normal.

In that case, since $(E_n)_s=0$, the closed immersion
\[
 i_n:S_n\longrightarrow S
\]
defined by $E_n$ is étale at $s$. But, $S$ being normal, the set of points at which $i_n$ is étale is both open and closed by the lemma below. Hence its image in $S$ is simultaneously open and closed. The restriction of $E_n$ to the connected component of $s$ in $S$ is therefore zero; this component is open because $S$ is locally connected. The lemma is proved.

\sgaStableTarget{sga6.e02.lem.2-4-8-2}\textbf{Lemma 2.4.8.2} (cf. EGA IV 18.10.3). Let $Y$ be a normal pseudo-analytic space, and let $f:X\to Y$ be a non-ramified morphism. The set of points of $X$ at which $f$ is étale is both open and closed.

\emph{Proof.} It is immediate that the set of points at which $f$ is étale is open. It remains to show that the set of points at which $f$ is not étale is open. Since a non-ramified morphism is locally an immersion, one may suppose that $f$ is a closed immersion defined by a coherent ideal $I$. Let $x$ be a point of $X$ such that $I_x\ne0$. Let $\mathfrak p$ be a prime ideal of $\mathcal O_{Y,x}$ containing $I_x$. One has
\[
 \dim(\mathcal O_{Y,x}/\mathfrak p)>0,
\]
for otherwise $\mathfrak p=0$, since $\mathcal O_{Y,x}$ is integral. Thus $Y$ has positive codimension at $x$, hence positive codimension in a neighbourhood of $x$; the assertion follows.

\sgaStableTarget{sga6.e02.sec.3}\section*{3. Complements on quasi-coherent sheaves on schemes}

\sgaStableTarget{sga6.e02.para.3-0}\textbf{3.0.} In this section we fix a quasi-compact and quasi-separated scheme $S$. The canonical, fully faithful inclusion functor
\[
 \iota:\Qcoh(S)\longrightarrow\Mod(S)
\]
extends to a functor, again denoted $\iota$, from $\mathrm D(\Qcoh(S))$ to $\mathrm D(S)=\mathrm D(\Mod(S))$, which we shall study.

\sgaStableTarget{sga6.e02.lem.3-1}\textbf{Lemma 3.1.} The category $\Qcoh(S)$ is an abelian category satisfying AB5, and the quasi-coherent sheaves of finite type form a small family of generators. In particular $\Qcoh(S)$ has enough injectives.

\emph{Proof.} The assertion about generators follows at once from the fact that every quasi-coherent sheaf on $S$ is the filtered inductive limit of its quasi-coherent subsheaves of finite type (EGA I 9.4.9 and IV 1.7.7). The other assertions are evident.

\sgaStableTarget{sga6.e02.rem.3-1-1}\textbf{Remark 3.1.1.} The editor does not know whether $\Qcoh(S)$ still has enough injectives when one assumes only that $S$ is quasi-separated.

\sgaStableTarget{sga6.e02.lem.3-2}\textbf{Lemma 3.2.} The functor
\[
 \iota:\Qcoh(S)\longrightarrow\Mod(S)
\]
has a right adjoint
\[
 Q:\Mod(S)\longrightarrow\Qcoh(S),
\]
called the coherer.

\emph{Proof.} If $S$ is affine, the existence of $Q$ is immediate: one takes $Q=\widetilde{\Gamma_S}$ (EGA I 1.3.7). Let now $f:U\to S$ be an affine open of $S$. It is clear that, for $E\in\Ob\Qcoh(S)$ and $F\in\Ob\Mod(U)$, one has a canonical isomorphism
\sgaStableTarget{sga6.e02.eq.1}\begin{equation}\tag{1}
 \Hom(E,f_*F)\simeq \Hom(E,f_*Q(F)).
\end{equation}
Let $(f_i:U_i\to S)$ be a finite covering of $S$ by affine opens, and, for each pair $(i,j)$, let $(U_{ijk}\to U_i\cap U_j)_k$ be a finite covering of $U_i\cap U_j$ by affine opens, which is possible since $S$ is quasi-separated. For $F\in\Ob\Mod(S)$ one has an exact sequence
\sgaStableTarget{sga6.e02.eq.2}\begin{equation}\tag{2}
 0\longrightarrow F\longrightarrow \prod_i f_{i*}(F|U_i)
 \longrightarrow \prod_{i,j,k}f_{ijk*}(F|U_{ijk}),
\end{equation}
where $f_{ijk}$ denotes the inclusion of $U_{ijk}$ into $S$. If $g:V\to U$ is an inclusion of affine opens of $S$, there is an evident natural arrow
\[
 Q(F|U)\longrightarrow g_*Q(F|V).
\]
Consequently (2) gives an arrow
\sgaStableTarget{sga6.e02.eq.3}\begin{equation}\tag{3}
 \prod_i f_{i*}Q(F|U_i)\longrightarrow
 \prod_{i,j,k}f_{ijk*}Q(F|U_{ijk}),
\end{equation}
whose kernel we denote by $QF$. If $E\in\Ob\Qcoh(S)$, applying the functor $\Hom(E,\cdot)$ to (2), and using (1), one obtains a functorial isomorphism in $E$,
\[
 \Hom(E,F)\simeq\Hom(E,QF).
\]
This proves the existence of $Q$; its uniqueness is evident, since it is an adjoint.

\sgaStableTarget{sga6.e02.lem.3-3}\textbf{Lemma 3.3.} \begin{enumerate}[label=\alph*)]
\item The coherer is an additive functor and transforms injectives into injectives.
\item Let
\[
 RQ:\mathrm D^+(S)\longrightarrow \mathrm D^+(\Qcoh(S))
\]
be the right derived functor of the coherer. For $E\in\Ob\mathrm D(\Qcoh(S))$ and $F\in\Ob\mathrm D^+(S)$ there are canonical functorial isomorphisms
\sgaStableTarget{sga6.e02.eq.i}\sgaStableTarget{sga6.e02.eq.ii}\begin{align}
 R\Hom_{\Mod}(E,F)&\simeq R\Hom_{\Qcoh}(E,RQ(F)), \tag{i}\\
 \Hom_{\Mod}(E,F)&\simeq \Hom_{\Qcoh}(E,RQ(F)). \tag{ii}
\end{align}
\end{enumerate}

\emph{Proof.} Assertion a) is an immediate consequence of the definition. For b), it is enough to establish the isomorphism (i), since (ii) follows by passing to cohomology. For this purpose one may suppose that $F$ has injective components. Then $QF$ has injective components by a), and the isomorphism (i) follows directly from the definition of $Q$.

\sgaStableTarget{sga6.e02.lem.3-4}\textbf{Lemma 3.4.} Let $f:X\to Y$ be a morphism of quasi-compact and quasi-separated schemes. For $E\in\Ob\Mod(X)$ there is a canonical functorial isomorphism
\[
 Q_Y f_*(E)\simeq f_*Q_X(E).
\]

\emph{Proof.} Since $f$ is quasi-compact and quasi-separated (EGA IV 1.2.3 and 1.2.4), $f_*$ transforms quasi-coherent sheaves into quasi-coherent sheaves (EGA III 1.4.10 and IV 1.7.21). The functors $Q_Yf_*$ and $f_*Q_X$ from $\Mod(X)$ to $\Qcoh(Y)$ are therefore both right adjoints to the functor
\[
 f^*: \Qcoh(Y)\longrightarrow \Mod(X),
\]
and hence are canonically isomorphic.

\sgaStableTarget{sga6.e02.prop.3-5}\textbf{Proposition 3.5.} If $S$ is noetherian, or quasi-compact and separated, the functor
\[
 \iota:\mathrm D^+(\Qcoh(S))\longrightarrow \mathrm D^+(S)
\]
is fully faithful and has as essential image the full subcategory $\mathrm D^+(S)_{\qcoh}$ of $\mathrm D^+(S)$ formed by complexes with quasi-coherent cohomology.

\emph{Proof.} One may restrict to proving the assertion in the case where $S$ is separated, since the noetherian case is already known. It is plainly equivalent to prove the following two assertions.

\begin{enumerate}[label=(3.5.\arabic*)]
\item For $E\in\Ob\mathrm D^+(\Qcoh(S))$, the adjunction arrow
\[
 E\longrightarrow RQ(E)
\]
is an isomorphism.
\item For $E\in\Ob\mathrm D^+(S)_{\qcoh}$, the adjunction arrow
\[
 RQ(E)\longrightarrow E
\]
is an isomorphism.
\end{enumerate}

\emph{Proof of \hyperlink{sga6.e02.prop.3-5}{(3.5.1)}.} By an immediate dévissage one is reduced to proving that quasi-coherent sheaves are acyclic for the coherer, that is, for $E\in\Ob\Qcoh(S)$, one has $R^iQ(E)=0$ for $i>0$, and that, moreover, the adjunction arrow $E\to QE$ is an isomorphism. If $S$ is affine, this follows from the formula $Q=\widetilde{\Gamma_S}$ used in the proof of \hyperlink{sga6.e02.lem.3-2}{(3.2)}. In the general case, cover $S$ by a finite number of affine opens $U_i$. Since $S$ is separated, the inclusions
\[
 f_{i_0\ldots i_p}:U_{i_0}\cap\cdots\cap U_{i_p}\longrightarrow S
\]
are affine morphisms. Let $E\in\Ob\Qcoh(S)$, and consider the resolution of $E$ defined by this covering,
\sgaStableTarget{sga6.e02.eq.unnamed.occ3}\begin{equation}\tag{*}
 0\longrightarrow E\longrightarrow \prod_i f_{i*}f_i^*E\longrightarrow\cdots\longrightarrow
 \prod_{i_0\ldots i_p}f_{i_0\ldots i_p*}f_{i_0\ldots i_p}^*E\longrightarrow\cdots .
\end{equation}
It gives a spectral sequence
\sgaStableTarget{sga6.e02.eq.unnamed.occ4}\begin{equation}\tag{**}
 E_1^{pq}=\prod_{i_0\ldots i_p}R^qQ\bigl(f_{i_0\ldots i_p*}f_{i_0\ldots i_p}^*E\bigr)
 \Longrightarrow R^{p+q}Q(E).
\end{equation}
Since $f_{i_0\ldots i_p}$ is affine, \hyperlink{sga6.e02.lem.3-4}{(3.4)} gives, for every $q$,
\sgaStableTarget{sga6.e02.eq.unnamed.occ5}\begin{equation}\tag{***}
 R^qQ\bigl(f_{i_0\ldots i_p*}f_{i_0\ldots i_p}^*E\bigr)
 \simeq f_{i_0\ldots i_p*}R^qQ\bigl(f_{i_0\ldots i_p}^*E\bigr).
\end{equation}
In particular $E_1^{pq}=0$ for $q>0$, by the affine case already treated. From (**) one deduces isomorphisms
\[
 R^nQ(E)\simeq H^n(QE')\quad(n\geq0),
\]
where $E'$ is the resolution of $E$ defined by the exact sequence (*). But by (***) one has $QE^n\simeq E^n$ for $n\geq0$; hence finally
\[
 R^nQ(E)=0\quad(n>0),
\]
and $Q(E)\simeq E$. This proves \hyperlink{sga6.e02.prop.3-5}{(3.5.1)}.

\emph{Proof of \hyperlink{sga6.e02.prop.3-5}{(3.5.2)}.} The spectral sequence
\[
 R^pQ(H^q(E))\Longrightarrow R^{p+q}Q(E)
\]
reduces us immediately to the case where $E\in\Ob\Qcoh(S)$. The composite
\[
 E\longrightarrow RQE\longrightarrow E
\]
of the adjunction arrows is the identity, and the assertion follows from \hyperlink{sga6.e02.prop.3-5}{(3.5.1)}. This completes the proof of \hyperlink{sga6.e02.prop.3-5}{(3.5)}.

\sgaStableTarget{sga6.e02.rem.3-6}\textbf{Remark 3.6.} The conclusion of \hyperlink{sga6.e02.prop.3-5}{(3.5)} would be false if one omitted the hypothesis ``separated or noetherian'', as is shown by a neat counterexample of J.-L. Verdier; see Appendix I to the present exposé.

\sgaStableTarget{sga6.e02.prop.3-7}\textbf{Proposition 3.7.} Under the hypothesis of \hyperlink{sga6.e02.prop.3-5}{(3.5)}, suppose moreover that there is an integer $N$ such that, for every $F\in\Ob\Mod(S)$ and every quasi-compact open $U$ of $S$, one has
\[
 H^i(U,F)=0\quad\text{for }i>N.
\]
Then:
\begin{enumerate}[label=\alph*)]
\item The coherer has finite cohomological dimension.
\item The functor
\[
 \iota:\mathrm D(\Qcoh(S))\longrightarrow \mathrm D(S)
\]
is fully faithful and has as essential image the full subcategory $\mathrm D(S)_{\qcoh}$ of $\mathrm D(S)$ formed by complexes with quasi-coherent cohomology.
\end{enumerate}

\emph{Proof of a).} If $S$ is affine, the formula $Q=\widetilde{\Gamma_S}$ implies that $Q$ is of cohomological dimension $\leq N$. We pass to the general case. Let us show that there is an integer $a$ such that $R^qQ(E)=0$ for every $E\in\Ob\Mod(S)$ and every $q>a$. The hypothesis implies that every $E\in\Ob\Mod(S)$ admits a finite right resolution $E'$ of length $N$ such that the $E^i$ are acyclic, that is, $H^q(U,E^i)=0$ for every open $U$ of $S$ and every $q>0$. Thus it is enough to prove that there is an integer $b$ such that
\[
 R^qQ(E)=0
\]
for every acyclic $\mathcal O_S$-module $E$ and every $q>b$; one may then take $a=b+N$.

Let $(U_i)$ be a finite covering of $S$ by $N'$ affine opens. Let $E$ be an acyclic $\mathcal O_S$-module and let $C'(E)$ be the resolution of $E$ defined by the covering $(U_i)$. There are two cases.

If $S$ is separated, the components of $C'(E)$ are sums of terms of the form $f_*f^*E$, where $f:U\to S$ is the inclusion of an affine open, hence an affine morphism. One has
\[
 RQ(Q_Xf_*)f^*E\simeq RQ\,Rf_*f^*E\simeq RQ\,f_*f^*E,
\]
since $f^*E$ is acyclic. On the other hand, by \hyperlink{sga6.e02.lem.3-4}{(3.4)},
\[
 R(Q_Xf_*)f^*E\simeq R(f_*Q_X)f^*E\simeq f_*RQf^*E,
\]
since $f$ is affine. Hence
\[
 R^nQ(f_*f^*E)\simeq f_*R^nQ(f^*E)\quad(n\geq0).
\]
By the affine case,
\[
 R^nQ(f_*f^*E)=0\quad(n>N).
\]
Since $C'(E)$ has length $N'$, it follows that
\[
 R^nQ(E)=0\quad(n>N+N').
\]

If $S$ is noetherian, the components of $C'(E)$ are sums of terms of the same form $f_*f^*E$, where now $f:U\to S$ is the inclusion of a quasi-compact and separated open. As in the first case, one has
\[
 R(Q_Xf_*)f^*E\simeq RQ\,f_*f^*E.
\]
Since $U$ is noetherian, $\Qcoh(U)$ has enough injective objects in $\Mod(U)$, and hence the right derived functor of $f_*:\Qcoh(U)\to\Qcoh(S)$ is induced by $Rf_*:\mathrm D^+(U)\to\mathrm D^+(S)$ when $\mathrm D^+(\Qcoh(U))$ and $\mathrm D^+(\Qcoh(S))$ are identified with full subcategories of $\mathrm D^+(U)$ and $\mathrm D^+(S)$ by \hyperlink{sga6.e02.prop.3-5}{(3.5)}. Thus
\[
 R(f_*Q_X)f^*E\simeq Rf_*RQf^*E.
\]
One deduces a spectral sequence
\[
 E_2^{pq}=R^pf_*R^qQ(f^*E)\Longrightarrow R^{p+q}(f_*Q)f^*E.
\]
By the first case there is an integer $n(U)$ depending only on $U$ such that $Q$ is of cohomological dimension $<n(U)$ on $\Mod(U)$; hence $E_2^{pq}=0$ for $q>n(U)$. Also $E_2^{pq}=0$ for $p>N$, by the cohomological-dimension hypothesis on $S$. Thus
\[
 R^n(f_*Q)f^*E=0\quad(n>N+N''),
\]
where $N''$ is an integer greater than all the $n(U)$, with $U$ running over the intersections $U_{i_0}\cap\cdots\cap U_{i_p}$. Taking account of the preceding displayed isomorphism, one obtains
\[
 R^nQ(f_*f^*E)=0\quad(n>N+N''),
\]
and hence finally
\[
 R^nQ(E)=0\quad(n>N+N'+N'').
\]
This proves a).

\emph{Proof of b).} The coherer, being of finite cohomological dimension by a), admits a right derived functor
\[
 RQ:\mathrm D(S)\longrightarrow\mathrm D(\Qcoh(S))
\]
whose restriction to $\mathrm D^+(S)$ coincides with the functor $RQ$ defined above. For $E\in\Ob\mathrm D(\Qcoh(S))$ there is a canonical functorial isomorphism
\sgaStableTarget{sga6.e02.eq.1.occ2}\begin{equation}\tag{1}
 E\xrightarrow{\sim}RQ(E),
\end{equation}
since, by \hyperlink{sga6.e02.prop.3-5}{(3.5.1)}, the $E^i$ are acyclic for $Q$ and satisfy $E^i\xrightarrow{\sim}Q(E^i)$. On the other hand, for $E\in\Ob\mathrm D(S)$ one has a canonical functorial homomorphism
\sgaStableTarget{sga6.e02.eq.2.occ2}\begin{equation}\tag{2}
 RQ(E)\longrightarrow E,
\end{equation}
which, in the case where the $E^i$ are $Q$-acyclic, comes from the ordinary adjunction homomorphism between $\iota$ and $Q$. It is checked directly that these homomorphisms make
\[
 RQ:\mathrm D(S)\longrightarrow\mathrm D(\Qcoh(S))
\]
a right adjoint of
\[
 \iota:\mathrm D(\Qcoh(S))\longrightarrow\mathrm D(S).
\]
The fact that (1) is an isomorphism implies that $\iota$ is fully faithful. Finally, using the spectral sequence
\[
 R^pQ(H^q(E))\Longrightarrow R^{p+q}Q(E),
\]
which converges since $Q$ has finite cohomological dimension, one sees that (2) is an isomorphism for $E\in\Ob\mathrm D(S)_{\qcoh}$; hence $\iota$ has image $\mathrm D(S)_{\qcoh}$. This proves b) and completes the proof of \hyperlink{sga6.e02.prop.3-7}{(3.7)}.

\sgaStableTarget{sga6.e02.appendix-i-a-counterexample-of-verdier}\section*{Appendix I. A counterexample of Verdier}

\sgaStableTarget{sga6.e02.sec.0}\subsection*{0. Introduction}
If $S$ is a scheme, write $\Mod(S)$, respectively $\Qcoh(S)$, for the category of $\mathcal O_S$-modules, respectively of quasi-coherent $\mathcal O_S$-modules. The purpose of this appendix is to draw attention to certain pathological properties of the category $\mathrm D(\Qcoh(S))$, and in particular to prove the following facts.

\sgaStableTarget{sga6.e02.para.0-1}\textbf{0.1.} Let $X$ be an affine scheme with ring $A$, and let $I$ be an injective $A$-module. We know that, if $A$ is noetherian, the sheaf $\widetilde I$ is an injective $\mathcal O_X$-module. If $A$ is not noetherian, $\widetilde I$ is not necessarily injective, nor even flasque, even if the underlying space of $X$ is noetherian of finite Krull dimension. Moreover, the restriction of $\widetilde I$ to an open $U$ of $X$ is not, in general, an injective object of $\Qcoh(U)$.

\sgaStableTarget{sga6.e02.para.0-2}\textbf{0.2.} Let $f:X\to Y$ be a morphism of quasi-compact and quasi-separated schemes. The functor
\[
 Rf_*:\mathrm D^+(X)\longrightarrow\mathrm D^+(Y)
\]
does not in general induce, on $\mathrm D^+(\Qcoh(X))$ and $\mathrm D^+(\Qcoh(Y))$, the right derived functor of
\[
 f_*:\Qcoh(X)\longrightarrow\Mod(Y).
\]
More precisely, an injective object of $\Qcoh(X)$ is not necessarily acyclic for $f_*:\Mod(X)\to\Mod(Y)$.

\sgaStableTarget{sga6.e02.para.0-3}\textbf{0.3.} Let $S$ be a quasi-compact and quasi-separated scheme. An object of $\Qcoh(S)$, even an injective object, is not necessarily acyclic for the coherer (\hyperlink{sga6.e02.lem.3-2}{II 3.2}). The functor
\[
 \iota:\mathrm D^b(\Qcoh(S))\longrightarrow\mathrm D^b(S)
\]
is not in general fully faithful.

\sgaStableTarget{sga6.e02.sec.1.occ2}\subsection*{1. The counterexamples}
We first recall the following lemma (SGA 2 II 9).

\sgaStableTarget{sga6.e02.lem.1-1}\textbf{Lemma 1.1.} Let $A$ be a ring, $X=\Spec(A)$, let $\mathbf f=(f_\alpha)$ be a finite system of elements of $A$, and let $Y$ be the closed subset defined by $\mathbf f(x)=0$. For $i>0$, the following conditions are equivalent:
\begin{enumerate}[label=(\roman*)]
\item for every injective $A$-module $I$, one has $H^i_Y(X,\widetilde I)=0$;
\item the pro-object $n\mapsto H_i(\mathbf f^n,A)$ is zero, that is, for every $n$ there exists $n'\geq n$ such that
\[
 H_i(\mathbf f^{n'},A)\longrightarrow H_i(\mathbf f^n,A)
\]
is zero.
\end{enumerate}

\sgaStableTarget{sga6.e02.para.1-2}\textbf{1.2.} We shall use \hyperlink{sga6.e02.lem.1-1}{(1.1)} to prove \hyperlink{sga6.e02.para.0-1}{(0.1)}. Let $k$ be a field,
\[
 B=k[x,y],\qquad \mathbf f=(x,y),\qquad
 M=\bigoplus_{n>0}B/(\mathbf f^n),
\]
and let $A$ be the $B$-algebra $D_B(M)$ (EGA $0_{\mathrm{IV}}$ 18.2.3). One has trivially
\[
 H_2^A(\mathbf f^n,A)\simeq H_2^B(\mathbf f^n,A).
\]
On the other hand,
\[
 H_2^B(\mathbf f^n,A)\simeq H_2^B(\mathbf f^n,B)\oplus H_2^B(\mathbf f^n,M),
\]
and
\[
 H_2^B(\mathbf f^n,B)=0,
\]
since $\mathbf f$ is a regular sequence. Thus
\[
 H_2^B(\mathbf f^n,A)\simeq H_2^B(\mathbf f^n,M).
\]
Finally,
\[
 H_2^B(\mathbf f^n,M)\simeq H_B^0(\mathbf f^n,M),
\]
where $H_B^0(\mathbf f^n,M)$ is canonically identified with the submodule $(\mathbf f^n,M)$ of $M$ annihilated by $\mathbf f^n$. Hence
\sgaStableTarget{sga6.e02.eq.unnamed.occ6}\begin{equation}\tag{*}
 H_2^A(\mathbf f^n,A)\simeq (\mathbf f^n,M).
\end{equation}
Since the transition morphism
\[
 H_2^A(\mathbf f^{n'},A)\longrightarrow H_2^A(\mathbf f^n,A)
\]
is induced by multiplication by $\mathbf f^{\,n'-n}$, one checks immediately from (*) that the pro-object $n\mapsto H_2^A(\mathbf f^n,A)$ is not zero. Hence, by \hyperlink{sga6.e02.lem.1-1}{(1.1)}, there is an injective $A$-module $I$ such that
\[
 H^2_Y(X,\widetilde I)\ne0,
\]
and therefore $\widetilde I$ is not flasque. Notice also that the projection $\Spec(A)\to\Spec(B)$ is a homeomorphism, so that the underlying space of $X$ is noetherian of dimension $2$. If $U$ denotes the complement of $Y$ in $X$, then by SGA 2 II 4,
\[
  H^2_Y(X,\widetilde I)\simeq H^1(U,\widetilde I)\ne0.
\]
But, since $U$ is separated, the functor $\mathrm D^+(\Qcoh(U))\to\mathrm D^+(U)$ is fully faithful by \hyperlink{sga6.e02.prop.3-5}{II 3.5}, so $\widetilde I|U$ is not injective in $\Qcoh(U)$. This proves assertion \hyperlink{sga6.e02.para.0-1}{(0.1)}.

Let $Z$ be the quasi-compact and quasi-separated scheme obtained by gluing two copies of $X$ along $U$; let $j_1,j_2$ be the two canonical open immersions of $X$ into $Z$. Thus one has a Cartesian square
\[
\begin{tikzcd}
 U \arrow[r,"i"] \arrow[d,"i"'] & X \arrow[d,"j_1"]\\
 X \arrow[r,"j_2"'] & Z .
\end{tikzcd}
\]
One has
\[
 j_2^*R^1j_{1*}(\widetilde I)\simeq R^1i_*(i^*\widetilde I),
\]
and hence
\[
 H^0(X,R^1j_{1*}(\widetilde I))\simeq H^1(U,\widetilde I)\ne0,
\]
so that $R^1j_{1*}(\widetilde I)\ne0$, which proves \hyperlink{sga6.e02.para.0-2}{(0.2)}.

On the other hand, by \hyperlink{sga6.e02.lem.3-4}{II 3.4} and \hyperlink{sga6.e02.prop.3-5}{3.5} one has
\[
 RQ\,Rj_{1*}(\widetilde I)\simeq j_{1*}(\widetilde I).
\]
It follows from this a spectral sequence whose exact sequence in low degrees gives the isomorphism
\[
 R^1j_{1*}(\widetilde I)\simeq R^2Q(j_{1*}\widetilde I).
\]
Since the sheaf $j_{1*}\widetilde I$ is injective in $\Qcoh(Z)$, the first part of \hyperlink{sga6.e02.para.0-3}{(0.3)} is proved.

Finally, let $E\in\Ob\Qcoh(Z)$. We have the duality isomorphism
\[
 R\Hom_{\Mod}(E,Rj_{1*}\widetilde I)
 \simeq R\Hom_{\Mod}(j_1^*E,\widetilde I),
\]
but, since $X$ is affine, \hyperlink{sga6.e02.prop.3-5}{II 3.5} gives
\[
 R\Hom_{\Mod}(j_1^*E,\widetilde I)
  \simeq R\Hom_{\Qcoh}(j_1^*E,\widetilde I)
  \simeq \Hom(j_1^*E,\widetilde I).
\]
Consequently
\[
 R\Hom_{\Mod}(E,Rj_{1*}\widetilde I)
 \simeq \Hom(j_1^*E,\widetilde I).
\]
A spectral sequence then gives, in low degree, the isomorphism
\[
 \Ext^2_{\Mod}(E,j_{1*}\widetilde I)
 \simeq \Hom(E,R^1j_{1*}\widetilde I).
\]
Thus the identity morphism of $E=R^1j_{1*}\widetilde I$ supplies a non-zero element of
\[
 \Ext^2_{\Mod}(E,j_{1*}\widetilde I).
\]
It follows that the functor
\[
 \mathrm D^b(\Qcoh(Z))\longrightarrow \mathrm D^b(Z)
\]
is not fully faithful. This completes the proof of \hyperlink{sga6.e02.para.0-3}{(0.3)}.




% Macros for Exposé II, Appendix II
\providecommand{\Diff}{\operatorname{Diff}}
\providecommand{\Hom}{\operatorname{Hom}}
\providecommand{\im}{\operatorname{Im}}
\providecommand{\ind}{\operatorname{ind}}
\providecommand{\Cone}{\operatorname{Cone}}
\providecommand{\Loclib}{\operatorname{Loclib}}
\providecommand{\ob}{\operatorname{ob}}
\providecommand{\Parf}{\operatorname{Parf}}
\providecommand{\Mod}{\operatorname{Mod}}
\providecommand{\Ob}{\operatorname{Ob}}
\providecommand{\Coker}{\operatorname{Coker}}
\providecommand{\Ker}{\operatorname{Ker}}
\providecommand{\cH}{\mathcal H}
\providecommand{\cO}{\mathcal O}
\providecommand{\cP}{\mathcal P}
\providecommand{\cE}{\mathcal E}
\providecommand{\cF}{\mathcal F}
\providecommand{\cG}{\mathcal G}
\providecommand{\cD}{\mathcal D}
\providecommand{\cT}{\mathcal T}
\providecommand{\cL}{\mathcal L}
\providecommand{\RHom}{\mathrm R\!\operatorname{Hom}}
\providecommand{\RGamma}{\mathrm R\Gamma}
\providecommand{\dlim}{\operatorname*{lim}}
\providecommand{\plim}{\operatorname*{lim}}


\newpage
\sgaStableTarget{sga6.e02.appendix-ii-definition-of-the-analytic-index-of}\section*{Appendix II. Definition of the analytic index of a relative elliptic complex}

In this appendix we propose to show that the notion of a perfect complex leads to a simple and natural definition of the ``analytic index'' of a continuous family of elliptic differential operators, or more generally of a relative elliptic complex. For the initiated reader, let us say that the definition given here has the advantage, compared with the classical definitions (see for example [2] and [17]), of not requiring any deformation of the elliptic complex under consideration.

Sections 1 and 2 consist of reminders on the classical notion of a ``mixed structure''.
\footnote{The informed reader is advised to skip them.}
We give, of course, only the minimum of definitions and results required in order to state and prove, in Section 3, the central theorem we have in view. It says, more precisely, that the direct image, by the projection (assumed proper) onto the base (assumed locally compact), of a relative elliptic complex is a perfect complex. When the base $S$ is compact, the definition of the analytic index as an element of $K(S)$, the Grothendieck group of complex vector bundles of finite rank on $S$, then follows from the globalization sorite of \hyperlink{sga6.e02.prop.2-3-2}{II 2.3.2}.

Throughout what follows, if $S$ is a topological space, $\mathcal O_S$ will denote the sheaf of continuous real-valued functions on $S$.

\sgaStableTarget{sga6.e02.sec.1.occ3}\subsection*{1. The sorite of mixed structures}

In this section a topological space $S$ is fixed.

\sgaStableTarget{sga6.e02.sec.1.1}\subsubsection*{1.1. Definition of mixed varieties}

Let $V$ be an open subset of $\mathbb R^n$. For $r\in\mathbb N$ one writes $C^r(V)$ for the Fréchet space of real-valued functions of class $C^r$ on $V$. Put
\[
 C^\infty(V)=\plim_r C^r(V).
\]
Let $U$ be an open subset of $S$. For $r\in\mathbb N\cup\{\infty\}$, denote by ${}^r\mathcal O_{U\times V}$ the sheaf of rings on $U\times V$ defined by
\[
 \Gamma(U'\times V',{}^r\mathcal O_{U\times V})
 =C(U',C^r(V')),
\]
where $U'$ (respectively $V'$) is an open subset of $U$ (respectively of $V$), and $C(U',C^r(V'))$ denotes the space of continuous maps from $U'$ into $C^r(V')$. There is a canonical injection
\[
 \operatorname{pr}_1^{-1}(\mathcal O_U)\longrightarrow {}^r\mathcal O_{U\times V}
\]
which makes $(U\times V,{}^r\mathcal O_{U\times V})$ into a ringed space over $U$.

A ringed space $X$ over $S$ is said to be a \emph{mixed variety of class $C^r$} ($r\in\mathbb N\cup\{\infty\}$) if $X$ is locally isomorphic, as a ringed space over $S$, to a space of the type just described. If $S$ is reduced to one point, a mixed variety of class $C^r$ over $S$ is just a variety of class $C^r$ in the usual sense.

Let $f:X\to S$ be a mixed variety of class $C^r$. For $s\in S$, the fiber $X_s=f^{-1}(s)$, endowed with the sheaf $i_s^{-1}(\mathcal O_X)/I_s$, where $i_s:X_s\to X$ is the injection and $I_s$ is the ideal of sections which vanish on $X_s$, is a variety of class $C^r$.

\sgaStableTarget{sga6.e02.sec.1.2}\subsubsection*{1.2. Relative differential operators}

It is easy, on the model of algebraic geometry (EGA IV 16, SGA 3 VII), to develop an infinitesimal study of mixed varieties. We shall only review rapidly the notions we shall need.

\sgaStableTarget{sga6.e02.para.1-2-1}\textbf{1.2.1.} Let $f:X\to S$ be a mixed variety of class $C^\infty$. For each $n\in\mathbb N$ one defines the sheaf $P^n_{X/S}$ of principal parts of order $n$ on $X$. In two different ways, corresponding to the two projections $\operatorname{pr}_1$ and $\operatorname{pr}_2$ of $X\times_SX$ onto $X$, it is an $\mathcal O_X$-algebra augmented toward $\mathcal O_X$, and is locally free of finite type as an $\mathcal O_X$-module. When the $\mathcal O_X$-module structure on $P^n_{X/S}$ is not specified, it will be understood to be the one coming from $\operatorname{pr}_1$. One has the canonical identifications
\[
 P^0_{X/S}=\mathcal O_X,
 \qquad
 P^1_{X/S}/\mathcal O_X=\Omega^1_{X/S}.
\]
The sheaf $\Omega^1_{X/S}$ is the sheaf of relative differentials of $X$ over $S$, or relative cotangent sheaf. Its dual is denoted $T_{X/S}$ and called the relative tangent sheaf. The sheaf $T_{X/S}$, respectively $\Omega^1_{X/S}$, is the sheaf of sections of a vector bundle of finite rank on $X$, called the relative tangent, respectively cotangent, bundle; these bundles are mixed varieties over $X$.

The $f^{-1}(\mathcal O_S)$-linear homomorphism, corresponding to the $\mathcal O_X$-algebra structure defined by $\operatorname{pr}_2$,
\[
 d^n_{X/S}:\mathcal O_X\longrightarrow P^n_{X/S},
\]
which associates to a section of $\mathcal O_X$ its principal part of order $n$, is called the universal differential operator of order $n$ on $X$ relative to $S$.

\sgaStableTarget{sga6.e02.para.1-2-2}\textbf{1.2.2.} More generally, if $E$ is a locally free $\mathcal O_X$-module of finite type, one may consider the sheaf of principal parts of order $n$ of $E$,
\[
 P^n_{X/S}(E)=P^n_{X/S}\otimes_{\mathcal O_X}E,
\]
where the tensor product is defined by means of the second $\mathcal O_X$-module structure on $P^n_{X/S}$. One has a universal differential operator
\[
 d^n_{X/S}(E):E\longrightarrow P^n_{X/S}(E).
\]
Let $E,F$ be locally free $\mathcal O_X$-modules of finite type. By definition, a differential operator from $E$ to $F$ of order $\le n$ is an $f^{-1}(\mathcal O_S)$-linear homomorphism from $E$ to $F$ which factors through $d^n_{X/S}(E)$ as an $\mathcal O_X$-linear homomorphism $P^n_{X/S}(E)\to F$; the factorization is then unique. In other words, $d^n_{X/S}(E)$ gives a bijection
\sgaStableTarget{sga6.e02.eq.1-2-2-1}\begin{equation}\tag{1.2.2.1}
 \Hom_{\mathcal O_X}(P^n_{X/S}(E),F)\xrightarrow{\sim}\Diff^n(E,F),
\end{equation}
where $\Diff^n(E,F)$ denotes the set of differential operators from $E$ to $F$ of order $\le n$.

\sgaStableTarget{sga6.e02.para.1-2-3}\textbf{1.2.3.} We shall also have to consider differential operators with complex coefficients. Let $\mathcal O_X\otimes\mathbb C$ (respectively $\mathcal O_S\otimes\mathbb C$) be the complexification of $\mathcal O_X$ (respectively $\mathcal O_S$). Let $E,F$ be locally free $(\mathcal O_X\otimes\mathbb C)$-modules of finite type. By a differential operator from $E$ to $F$ of order $\le n$ one means an $f^{-1}(\mathcal O_S\otimes\mathbb C)$-linear homomorphism from $E$ to $F$ which factors through $d^n_{X/S}(E)$ as an $(\mathcal O_X\otimes\mathbb C)$-linear homomorphism $P^n_{X/S}(E)\to F$. If $\Diff^n_{\mathbb C}(E,F)$ denotes the set of such operators, $d^n_{X/S}(E)$ therefore defines a bijection
\sgaStableTarget{sga6.e02.eq.1-2-3-1}\begin{equation}\tag{1.2.3.1}
 \Hom_{\mathcal O_X\otimes\mathbb C}(P^n_{X/S}(E),F)
 \xrightarrow{\sim}\Diff^n_{\mathbb C}(E,F).
\end{equation}

\sgaStableTarget{sga6.e02.para.1-2-4}\textbf{1.2.4.} For $n\ge1$ there is a canonical exact sequence
\[
 0\longrightarrow S^n(\Omega^1_{X/S})\longrightarrow P^n_{X/S}
 \longrightarrow P^{n-1}_{X/S}\longrightarrow 0.
\]
Hence, tensoring over $\mathcal O_X$ with a locally free $\mathcal O_X$-module, respectively $(\mathcal O_X\otimes\mathbb C)$-module, $E$ of finite type, one obtains an exact sequence
\sgaStableTarget{sga6.e02.eq.1-2-4-1}\begin{equation}\tag{1.2.4.1}
0\longrightarrow S^n(\Omega^1_{X/S})\otimes E\longrightarrow P^n_{X/S}(E)
\longrightarrow P^{n-1}_{X/S}(E)\longrightarrow0.
\end{equation}
Here $S^n(\Omega^1_{X/S})$ denotes the $n$-th symmetric power.

Let $F$ be a locally free $\mathcal O_X$-module, respectively $(\mathcal O_X\otimes\mathbb C)$-module, of finite type. Applying $\Hom(\ ,F)$, respectively $\Hom_{\mathcal O_X\otimes\mathbb C}(\ ,F)$, to \hyperlink{sga6.e02.eq.1-2-4-1}{(1.2.4.1)}, and using \hyperlink{sga6.e02.eq.1-2-2-1}{(1.2.2.1)} and \hyperlink{sga6.e02.eq.1-2-3-1}{(1.2.3.1)}, one obtains an exact sequence
\sgaStableTarget{sga6.e02.eq.1-2-4-2}\begin{equation}\tag{1.2.4.2}
0\longrightarrow \Diff^{n-1}(E,F)\longrightarrow \Diff^n(E,F)
\xrightarrow{\sigma_n}
\Hom_{\mathcal O_X}\bigl(S^n(\Omega^1_{X/S})\otimes E,F\bigr),
\end{equation}
and analogously in the complex case. If $X$ is paracompact and $\mathcal O_X$ is soft, the sequence may be completed on the right by adding a zero. The homomorphism $\sigma_n$ is called the \emph{symbol}. It associates to a differential operator from $E$ to $F$ of order $\le n$ a homogeneous polynomial map of degree $n$ from $T^*_{X/S}$ into $\Hom(E,F)$, or equivalently a homomorphism $p^*E\to p^*F$, where $p$ is the projection of the relative cotangent bundle onto $X$, satisfying the usual homogeneity condition.

Let $E,F,G$ be locally free $\mathcal O_X$-modules of finite type, $u\in\Diff^n(E,F)$ and $v\in\Diff^m(F,G)$. Then $vu\in\Diff^{m+n}(E,G)$ and one checks at once the relation
\sgaStableTarget{sga6.e02.eq.1-2-4-3}\begin{equation}\tag{1.2.4.3}
 \sigma_{m+n}(vu)=\sigma_m(v)\sigma_n(u).
\end{equation}
The same relation holds for differential operators with complex coefficients. In particular, if the composite of two differential operators is zero, then the composite of their symbols is zero.

\sgaStableTarget{sga6.e02.para.1-2-5}\textbf{1.2.5.} Let $E^\bullet$ be a bounded complex of differential operators with complex coefficients on $X$ relative to $S$, and let $d^i:E^i\to E^{i+1}$ be of order $n_i$. By \hyperlink{sga6.e02.eq.1-2-4-3}{(1.2.4.3)},
\[
 \sigma_{n_{i+1}}(d^{i+1})\sigma_{n_i}(d^i)=0.
\]
The complex $E^\bullet$ is called an \emph{elliptic complex of order $(n_i)$} if the symbol complex on the relative cotangent bundle, pulled back to that bundle, is acyclic outside the zero section.

When $S$ is reduced to one point, this notion of elliptic complex on $X$ coincides with the usual one, studied in detail in [4] and [16].

\sgaStableTarget{sga6.e02.para.1-2-6}\textbf{1.2.6.} Let $s\in S$, and let $i_s$ be the inclusion of the fiber $X_s$ in $X$, considered as a monomorphism of ringed spaces \hyperlink{sga6.e02.sec.1.1}{(1.1)}. For $n\in\mathbb N$ one has a canonical isomorphism
\[
 i_s^*P^n_{X/S}\simeq P^n_{X_s},
\]
and hence a restriction homomorphism
\[
 i_s^*:\Diff^n(E,F)\longrightarrow \Diff^n(i_s^*E,i_s^*F),
\]
for $E,F$ locally free $\mathcal O_X$-modules of finite type, and similarly for operators with complex coefficients. If $E^\bullet$ is a bounded complex of differential operators with complex coefficients on $X$, then $E^\bullet$ is elliptic if and only if, for every $s\in S$, $i_s^*(E^\bullet)$ is an elliptic complex. The verification is left to the reader.

\sgaStableTarget{sga6.e02.sec.2.occ2}\subsection*{2. Rigidity lemmas}

\sgaStableTarget{sga6.e02.para.2-0}\textbf{2.0.} It is ``well known'' that compact differentiable varieties are rigid, that is, they admit no continuous deformation other than the trivial one. Since this result will be used essentially in Section 3, and references seem rather scarce, it is perhaps useful to give precise statements and a few proofs.

A topological space $S$ is fixed. The mixed varieties under consideration are assumed to be of class $C^r$, with $r\gg1$; possibly $r=\infty$.

\sgaStableTarget{sga6.e02.para.2-1}\textbf{2.1. Notation.} Let $X,Y$ be mixed $S$-varieties. We denote by $\Hom(X,Y)$ the set of $S$-morphisms from $X$ to $Y$, and by $\mathcal Hom(X,Y)$ the sheaf on $X$ of germs of $S$-morphisms from $X$ to $Y$, defined by
\[
 \mathcal Hom(X,Y)(U)=\Hom(U,Y)
\]
for $U$ open in $X$.

Let $s\in S$, and write $i_s$ for the inclusion of $X_s$ in $X$. There is a canonical induction homomorphism
\sgaStableTarget{sga6.e02.eq.2-1-1}\begin{equation}\tag{2.1.1}
 i_s^{-1}\mathcal Hom(X,Y)\longrightarrow \mathcal Hom(X_s,Y_s)
\end{equation}
which is an epimorphism, as is immediate. In the case $Y=S\times\mathbb R$, the sheaf $\mathcal Hom(X,Y)$ is canonically identified with $\mathcal O_X$, and the epimorphism \hyperlink{sga6.e02.eq.2-1-1}{(2.1.1)} with the epimorphism
\[
 i_s^{-1}(\mathcal O_X)\longrightarrow i_s^*(\mathcal O_X)
\]
considered in \hyperlink{sga6.e02.sec.1.1}{(1.1)}.

\sgaStableTarget{sga6.e02.prop.2-2}\textbf{Proposition 2.2.} Let $X,Y$ be mixed $S$-varieties, and let $s\in S$. Suppose that the fiber $X_s$ is paracompact and that $s$ has a fundamental system of paracompact neighborhoods. Let $A$ be a closed subset of $X_s$, and let
\[
 f\in H^0(A,\mathcal Hom(X_s,Y_s)).
\]
The subsheaf $\mathcal F$ of $i_s^{-1}\mathcal Hom(X,Y)$ consisting of the germs which extend $f$, that is, whose image by \hyperlink{sga6.e02.eq.2-1-1}{(2.1.1)} coincides with $f$ over $A$, is soft.

\emph{Proof.} Let $X'$ be an open neighborhood in $X$ which is isomorphic, as a mixed variety, to $S'\times U$, where $S'$ is a paracompact neighborhood of $s$ and $U$ is a relatively compact open subset of $\mathbb R^n$, and such that the image by $f$ of $X'_s\cap A$ is contained in an open subset of $Y$ isomorphic to $S'\times V$, where $V$ is a convex open subset of $\mathbb R^q$ containing the origin. By T.F. II 3.4.1 it is enough to show that every section of $\mathcal F$ over a closed subset of $X_s$ contained in $X'_s$ extends to $X'_s$. Since $X'_s=\{s\}\times U$ is relatively compact, a closed subset of $X_s$ contained in $X'_s$ is compact. Identifying $X'$ with $S'\times U$ and $Y'$ with $S'\times V$, putting $A'=A\cap X'_s$, and using T.F. II 3.3.1, the assertion is reduced to the following elementary lemma.\footnote{Communicated to the writer by A. Douady. A closed subset of a paracompact space has a fundamental system of paracompact neighborhoods.}

\sgaStableTarget{sga6.e02.lem.2-2-1}\textbf{Lemma 2.2.1.} Let $U$ be an open subset of $\mathbb R^n$, and let $V$ be a convex open subset of $\mathbb R^q$ containing $0$. Let $A_1,A_2$ be closed subsets of $\{s\}\times U$. Let $U_1$, respectively $U_2$, be an open subset of $S\times U$, respectively of $\{s\}\times U$, containing $A_1$, respectively $A_2$. Let $f_1:U_1\to V$ be a morphism and let $f_2:U_2\to S\times V$ be an $S$-morphism which, on $U_1\cap(\{s\}\times U)$, defines the same germ as $f_1$ near $A_1$. Then there exists an open subset $U'$ of $S\times U$ containing $\{s\}\times U$ and an $S$-morphism
\[
 F:U'\longrightarrow S\times V
\]
such that the germ of $F$ along $\{s\}\times U$ extends the germ defined by $f_1$, and the morphism induced by $F$ on $\{s\}\times U$ defines the same germ as $f_2$ near $A_2$.

\emph{Proof.} Replacing $f_1$ by $af_1$, where $a$ is a $C^\infty$ function on $U$ equal to $1$ near $A_1$ and with support in $U_1$, one may suppose $U_1=U$. Let $b$ be a $C^\infty$ function on $U$, equal to $1$ on $A_2$ and with support in $U_2\cap(\{s\}\times U)$. Then
\[
 U' = U_1\cup\bigl(S\times (U-\operatorname{Supp}b)\bigr)
\]
is an open subset of $S\times U$ containing $\{s\}\times U$. The morphism $F:U'\to S\times V$ defined by
\[
 F(t,x)=\bigl(t,\,b(x)\operatorname{pr}_2 f_2(x)+(1-b(x))f_1(t,x)\bigr)
\]
visibly satisfies the required conditions.

\sgaStableTarget{sga6.e02.cor.2-3}\textbf{Corollary 2.3.} Under the hypotheses of \hyperlink{sga6.e02.prop.2-2}{(2.2)}, the sheaves $i_s^{-1}\mathcal Hom(X,Y)$ and $\mathcal Hom(X_s,Y_s)$ are soft. In particular, the sheaves $i_s^{-1}(\mathcal O_X)$ and $\mathcal O_{X_s}$ are soft. For every closed subset $A$ of $X_s$, the map
\[
 H^0(X_s,i_s^{-1}\mathcal Hom(X,Y))\longrightarrow H^0(A,\mathcal Hom(X_s,Y_s))
\]
deduced from \hyperlink{sga6.e02.eq.2-1-1}{(2.1.1)} is surjective.

\emph{Proof.} The softness of $i_s^{-1}\mathcal Hom(X,Y)$ follows from \hyperlink{sga6.e02.prop.2-2}{(2.2)} applied with $A=\varnothing$. Taking $S=\{s\}$, one obtains the softness of $\mathcal Hom(X_s,Y_s)$. The last assertion follows from \hyperlink{sga6.e02.prop.2-2}{(2.2)}.\footnote{For $\mathcal O_{X_s}$ this is very well known. The writer does not know whether $\mathcal O_X$ is soft when the hypotheses of \hyperlink{sga6.e02.prop.2-2}{(2.2)} are satisfied for every $s\in S$ and $X$ is paracompact. This would trivially imply, under those hypotheses, the softness of $i_s^{-1}(\mathcal O_X)$ as well as \hyperlink{sga6.e02.cor.2-5}{(2.5)} below.}

\sgaStableTarget{sga6.e02.cor.2-4}\textbf{Corollary 2.4.} Under the hypotheses of \hyperlink{sga6.e02.prop.2-2}{(2.2)}, let $E$ be an $\mathcal O_X$-module. For every closed subset $A$ of $X_s$, the canonical homomorphism
\[
 H^0(X_s,i_s^{-1}E)\longrightarrow H^0(A,i_s^*E)
\]
deduced from the epimorphism
\[
 i_s^{-1}E\longrightarrow i_s^*E=i_s^{-1}E\otimes_{i_s^{-1}\mathcal O_X}\mathcal O_{X_s}
\]
is an epimorphism.

\emph{Proof.} Since $i_s^{-1}\mathcal O_X$ is soft by \hyperlink{sga6.e02.cor.2-3}{(2.3)}, the functor $H^0(A,\ )$ is exact on the category of $i_s^{-1}\mathcal O_X$-modules (T.F. II 3.7.1, 3.4.2, 4.4.3). Hence
\[
 H^0(A,i_s^{-1}E)\longrightarrow H^0(A,i_s^*E)
\]
is an epimorphism. Since $i_s^{-1}E$ is soft, the homomorphism
\[
 H^0(X_s,i_s^{-1}E)\longrightarrow H^0(A,i_s^{-1}E)
\]
is an epimorphism, and this proves the assertion.

\sgaStableTarget{sga6.e02.cor.2-5}\textbf{Corollary 2.5.} Let $p:X\to S$ be a mixed variety and let $s\in S$. Suppose that $s$ has a fundamental system of paracompact neighborhoods $S'$ such that the $p^{-1}(S')$ form a fundamental system of paracompact neighborhoods of $X_s$. This condition is satisfied, for example, if $S$ is locally compact and $p$ is a proper map of topological spaces. Then, for every $\mathcal O_X$-module $E$,
\[
 (R^ip_*E)_s=0\qquad(i>0).
\]

\emph{Proof.} By T.F. II 4.17.1,
\[
 (R^ip_*E)_s=H^i(X_s,i_s^{-1}E).
\]
But $i_s^{-1}\mathcal O_X$ is soft by \hyperlink{sga6.e02.cor.2-3}{(2.3)}, hence $H^i(X_s,i_s^{-1}E)=0$ for $i>0$.

\sgaStableTarget{sga6.e02.cor.2-6}\textbf{Corollary 2.6.} Let $X,Y$ be mixed $S$-varieties, and let $s\in S$. Suppose that $X$ satisfies the hypotheses of \hyperlink{sga6.e02.cor.2-5}{(2.5)}. For every morphism $f:X_s\to Y_s$, there exists an open neighborhood $S'$ of $s$ and a morphism
\[
 F:X|S'\longrightarrow Y|S'
\]
such that $F_s=f$.

\emph{Proof.} This is an immediate consequence of \hyperlink{sga6.e02.cor.2-3}{(2.3)}, with $A=X_s$, and of the fact that the hypotheses on $X$ give
\[
 H^0(X_s,i_s^{-1}\mathcal Hom(X,Y))
 =\dlim_{S'\ni s}\Hom(X|S',Y|S')
\]
for open $S'$ containing $s$ (T.F. II 8.5.5).

\sgaStableTarget{sga6.e02.prop.2-7}\textbf{Proposition 2.7.} Let $X,Y$ be mixed $S$-varieties, let $f:X\to Y$ be a morphism, let $s\in S$, and let $x\in X_s$. If the tangent linear map of
\[
 f_s:X_s\longrightarrow Y_s
\]
is invertible at $x$, then $f$ is a local isomorphism at $x$; that is, there is an open neighborhood $U$ of $x$ in $X$ and an open subset $V$ of $Y$ such that $f$ induces an isomorphism $U\simeq V$.

\emph{Proof.} This is an easy consequence of the implicit function theorem. The details are left to the reader.

\sgaStableTarget{sga6.e02.cor.2-8}\textbf{Corollary 2.8.} Let $X,Y$ be mixed $S$-varieties and let $s\in S$. Suppose that $s$ has a fundamental system of neighborhoods $S'$ such that the $X|S'$, respectively $Y|S'$, form a fundamental system of neighborhoods of $X_s$, respectively $Y_s$; this is the case, for example, if $S$ is locally compact and $X$, respectively $Y$, is proper over $S$. Let $f:X\to Y$ be an $S$-morphism such that
\[
 f_s:X_s\longrightarrow Y_s
\]
is an isomorphism. Then there exists a neighborhood $S'$ of $s$ such that
\[
 f|X|S':X|S'\longrightarrow Y|S'
\]
is an isomorphism.

\emph{Proof.} By \hyperlink{sga6.e02.prop.2-7}{(2.7)}, there exists an open neighborhood $U$ of $X_s$ in $X$ such that $f|U$ is étale, that is, a local isomorphism. We may suppose that $U=X|S'$ for a neighborhood $S'$ of $s$; after shrinking $S$, we may therefore assume that $f$ is étale. Then the fiber product $Z=X\times_YX$ is naturally a mixed $S$-variety, étale over $X$ for each projection, and the diagonal morphism $\Delta:X\to Z$ is an open immersion. Moreover the sets $Z|S'$ form a fundamental system of neighborhoods of $Z_s$. Since
\[
 \Delta_s:X_s\longrightarrow Z_s
\]
is an isomorphism, it follows, after shrinking $S$, that $\Delta$ is an isomorphism. Thus $f$ is an open immersion, and after shrinking $S$ again it is an isomorphism.

\sgaStableTarget{sga6.e02.cor.2-9}\textbf{Corollary 2.9.} Let $X$ be a mixed $S$-variety and let $s\in S$. Suppose that the hypothesis of \hyperlink{sga6.e02.cor.2-5}{(2.5)} is satisfied. Then there exists a neighborhood $S'$ of $s$ in $S$ such that $X|S'$ is $S'$-isomorphic to $S'\times X_s$.

\emph{Proof.} Apply \hyperlink{sga6.e02.cor.2-6}{(2.6)} and \hyperlink{sga6.e02.cor.2-8}{(2.8)}.

\sgaStableTarget{sga6.e02.cor.2-10}\textbf{Corollary 2.10.} Let $s\in S$ and let $V$ be a paracompact variety. Put $X=S\times V$, and let $i_s$ be the injection of $V$ into $X$ defined by $s$. Suppose that the hypothesis of \hyperlink{sga6.e02.cor.2-5}{(2.5)} is satisfied. Let $E$ be a locally free $\mathcal O_X$-module, respectively $\mathcal O_X\otimes\mathbb C$-module, of finite type. Then there exists a neighborhood $S'$ of $s$ in $S$ such that $E|X|S'$ is isomorphic to $\operatorname{pr}_2^*(i_s^*E)$, where $\operatorname{pr}_2:S'\times V\to V$ is the second projection.

\emph{Proof.} Put $F=\operatorname{pr}_2^*(i_s^*E)$. By \hyperlink{sga6.e02.cor.2-4}{(2.4)}, applied to $\mathcal Hom(E,F)$, the induction homomorphism
\[
 \Hom(i_s^{-1}E,i_s^{-1}F)\longrightarrow \Hom(i_s^*E,i_s^*F)
\]
is an epimorphism. On the other hand, by T.F. II 3.3.1,
\[
 \Hom(i_s^{-1}E,i_s^{-1}F)
 =\dlim_{S'\ni s}\Hom(E|X|S',F|X|S') .
\]
Thus the canonical isomorphism $i_s^*E\simeq i_s^*F$ is induced, after shrinking $S$, by a homomorphism $f:E\to F$. Since $f_s$ is an isomorphism, $f$ is an isomorphism near $X_s$ by Nakayama, hence over $X|S'$ for a suitable neighborhood $S'$ of $s$.

\sgaStableTarget{sga6.e02.sec.3.occ2}\subsection*{3. The finiteness theorem}

\sgaStableTarget{sga6.e02.sec.3.0}\subsubsection*{3.0. Conventions and notation}

\sgaStableTarget{sga6.e02.para.3-0-1}\textbf{3.0.1.} Let $X$ and $Y$ be topological spaces. We write $C(X,Y)$ for the set of continuous maps from $X$ to $Y$, and $\mathcal C(X,Y)$ for the sheaf on $X$ defined by $U\mapsto C(U,Y)$.

\sgaStableTarget{sga6.e02.para.3-0-2}\textbf{3.0.2.} Let $E,F$ be Banach spaces, real or complex. We write $L(E,F)$ for the Banach space of continuous linear maps from $E$ to $F$.

\sgaStableTarget{sga6.e02.para.3-0-3}\textbf{3.0.3.} The mixed varieties occurring in this section are assumed to be of class $C^\infty$. If $X$ is a mixed $S$-variety, we shall often identify a locally free $\mathcal O_X$-module, respectively $\mathcal O_X\otimes\mathbb C$-module, of finite type with the corresponding real, respectively complex, vector bundle.

\sgaStableTarget{sga6.e02.para.3-0-4}\textbf{3.0.4.} Let $X$ be a compact ordinary variety, and let $E$ be a real or complex vector bundle of finite rank on $X$. Write $C^\infty(X,E)$ for the Fréchet space of $C^\infty$ sections of $E$ on $X$. For $r\in\mathbb Z$, write $C_r(X,E)$ for the Hilbert space of sections of $E$ which are $r$ times differentiable in $L^2$ ([16] IX). For $s\le r$, $C_r(X,E)$ maps to $C_s(X,E)$ by a continuous linear injective map with dense image, and
\[
 C^\infty(X,E)=\plim_r C_r(X,E).
\]

\sgaStableTarget{sga6.e02.para.3-0-5}\textbf{3.0.5.} Let $S$ be a topological space. We shall need the notion of a Hilbert bundle over $S$, for which we refer the reader to [12] or [14]. There is an additive functor from the category of real, respectively complex, Hilbert bundles over $S$ to the category of $\mathcal O_S$-, respectively $\mathcal O_S\otimes\mathbb C$-, modules: it sends a Hilbert bundle $E$ to the sheaf of continuous sections of $E$.

\sgaStableTarget{sga6.e02.sec.3.1}\subsubsection*{3.1. Auxiliary technical results}

\sgaStableTarget{sga6.e02.lem.3-1-1}\textbf{Lemma 3.1.1.} Let $X,Y$ be topological spaces, with $X$ locally compact. Let $\mathcal C'(X,Y)$ be the presheaf, on the basis of relatively compact open subsets of $X$, defined by
\[
 \mathcal C'(X,Y)(U)=C(\overline U,Y).
\]
The canonical restriction morphism
\[
 \mathcal C'(X,Y)\longrightarrow \mathcal C(X,Y)
\]
induces an isomorphism on the associated sheaves.

\emph{Proof.} Immediate.

\sgaStableTarget{sga6.e02.lem.3-1-2}\textbf{Lemma 3.1.2.} Let $E,F$ be locally convex topological vector spaces over $\mathbb C$, and let $u:E\to F$ be a continuous linear map, respectively a continuous linear map with dense image. Let $X$ be compact. Endow $C(X,E)$ and $C(X,F)$ with the topology of uniform convergence. Then the induced linear map
\[
 u_X:C(X,E)\longrightarrow C(X,F)
\]
is continuous, respectively continuous with dense image.

\emph{Proof.} It is clear that $u_X$ is continuous. On the other hand, by Bourbaki, \emph{Intégration}, chapter II, §2, lemmas 1 and 2, the subspace $C(X,\mathbb C)\otimes E$ of $C(X,E)$, formed by the continuous maps from $X$ to $E$ whose values lie in a finite-dimensional vector subspace, is dense; and similarly $C(X,\mathbb C)\otimes F$ is dense in $C(X,F)$. If $u$ has dense image, then $u_X(C(X,E))$ is dense in $C(X,\mathbb C)\otimes F$, and taking closures gives the assertion.

\sgaStableTarget{sga6.e02.para.3-1-3}\textbf{3.1.3.} Let $V$ be a compact variety, let $E$ and $F$ be complex vector bundles of finite rank on $V$, and let $n\in\mathbb N$. The space $\Diff^n_{\mathbb C}(E,F)$, identified by \hyperlink{sga6.e02.eq.1-2-3-1}{(1.2.3.1)} with
\[
 \Hom(P^n(E),F)\simeq C^\infty(V,\Hom(P^n(E),F)),
\]
is naturally a Fréchet space. Let $u\in\Diff^n_{\mathbb C}(E,F)$. One checks easily that the linear map $C^\infty(V,E)\to C^\infty(V,F)$ defined by $u$ is continuous. If $L(C^\infty(V,E),C^\infty(V,F))$ denotes the space of continuous linear maps, one therefore has a linear map
\[
 \delta:\Diff^n_{\mathbb C}(E,F)
 \longrightarrow L(C^\infty(V,E),C^\infty(V,F)).
\]
Moreover, for $r\in\mathbb Z$, $u$ acts by extension of scalars on the $C_r$ sections of $E$ and defines a linear map
\[
 \delta_r:\Diff^n_{\mathbb C}(E,F)
 \longrightarrow L(C_r(V,E),C_{r-n}(V,F)).
\]
The maps $\delta_r$ are continuous, and for $u\in\Diff^n_{\mathbb C}(E,F)$, $\delta(u)$ is the inverse-limit element of the projective system of the $\delta_r(u)$. These assertions follow easily from [16] IX: one reduces immediately to $n=0$; the statement on the $C_r$ follows directly from the definitions, while the statement for $C^\infty$ follows, after localizing by a partition of unity, from the formula for successive derivatives of a product. The analogous real statements are also valid.

\sgaStableTarget{sga6.e02.lem.3-1-4}\textbf{Lemma 3.1.4.} Let $S$ be a topological space, let $V$ be a compact variety, let $E,F$ be complex vector bundles of finite rank on $V$, and let $n\in\mathbb N$. Put $X=S\times V$. For
\[
 u\in\Diff^n_{\mathbb C}(\operatorname{pr}_2^*E,\operatorname{pr}_2^*F),
\]
the map $s\mapsto i_s^*u$ from $S$ into $\Diff^n_{\mathbb C}(E,F)$ is continuous. The linear map
\[
 \Diff^n_{\mathbb C}(\operatorname{pr}_2^*E,\operatorname{pr}_2^*F)
 \longrightarrow C(S,\Diff^n_{\mathbb C}(E,F))
\]
so defined is an isomorphism.

\emph{Proof.} By \hyperlink{sga6.e02.eq.1-2-3-1}{(1.2.3.1)} one may restrict to the case $n=0$. Putting $G=\Hom(E,F)$, one is reduced to showing that for $u\in\Gamma(X,\operatorname{pr}_2^*G)$, the map $s\mapsto u_s=i_s^*u$ from $S$ into $C^\infty(V,G)$ is continuous and that the resulting map
\[
 \Gamma(X,\operatorname{pr}_2^*G)\longrightarrow C(S,C^\infty(V,G))
\]
is an isomorphism. This follows immediately from the definition of mixed varieties.

\sgaStableTarget{sga6.e02.para.3-1-5}\textbf{3.1.5.} Under the hypotheses of \hyperlink{sga6.e02.lem.3-1-4}{(3.1.4)}, let
\[
 u\in\Diff^n_{\mathbb C}(\operatorname{pr}_2^*E,\operatorname{pr}_2^*F).
\]
Identifying $u$, by \hyperlink{sga6.e02.lem.3-1-4}{(3.1.4)}, with a continuous map from $S$ into $\Diff^n(E,F)$, one sees that, for $U$ open in $S$, the morphism
\[
 \operatorname{pr}_{1*}(u)(U):C^\infty(U,E)\longrightarrow C^\infty(U,F)
\]
is canonically identified with the linear map
\[
 C(U,C^\infty(V,E))\longrightarrow C(U,C^\infty(V,F))
\]
deduced from
\[
 \delta u:U\longrightarrow L(C^\infty(V,E),C^\infty(V,F)).
\]
Suppose $S$ locally compact. By \hyperlink{sga6.e02.lem.3-1-1}{(3.1.1)}, the morphism of sheaves
\[
 \operatorname{pr}_{1*}(u):
 \operatorname{pr}_{1*}\operatorname{pr}_2^*E
 \longrightarrow
 \operatorname{pr}_{1*}\operatorname{pr}_2^*F
\]
is the morphism associated to the morphism of presheaves
\[
 \mathcal C'(S,C^\infty(V,E))\longrightarrow \mathcal C'(S,C^\infty(V,F))
\]
defined by $\delta u$. On the other hand, for $r\in\mathbb Z$,
\[
 \delta_r u:S\longrightarrow L(C_r(V,E),C_{r-n}(V,F))
\]
defines a morphism of presheaves
\[
 u_r:\mathcal C'(S,C_r(V,E))\longrightarrow \mathcal C'(S,C_{r-n}(V,F)).
\]
By \hyperlink{sga6.e02.para.3-1-3}{(3.1.3)}, the $u_r$ form a projective system whose limit is the previous morphism. It will be convenient to interpret $\delta_r u$ as a morphism of trivial Hilbert bundles
\[
 S\times C_r(V,E)\longrightarrow S\times C_{r-n}(V,F),
\]
inducing on sheaves of sections the sheaf morphism associated to $u_r$.

\sgaStableTarget{sga6.e02.sec.3.2}\subsubsection*{3.2. Definition of the index}

\sgaStableTarget{sga6.e02.thm.3-2-1}\textbf{Theorem 3.2.1.} Let $S$ be a locally compact space, let $p:X\to S$ be a proper mixed variety, and let $E^\bullet$ be an elliptic complex of order $(n_i)$ of complex vector bundles on $X$ \hyperlink{sga6.e02.para.1-2-5}{(1.2.5)}. Then the complex $Rp_*(E^\bullet)$, which by \hyperlink{sga6.e02.cor.2-5}{(2.5)} is isomorphic to $p_*(E^\bullet)$ in the derived category of $\mathcal O_S\otimes\mathbb C$-modules, is perfect.

\emph{Proof.} The question being local on $S$, one may suppose by \hyperlink{sga6.e02.cor.2-9}{(2.9)} that $X=S\times V$, where $V$ is a compact variety, and that $p=\operatorname{pr}_1$. Moreover, by \hyperlink{sga6.e02.cor.2-10}{(2.10)}, one may suppose that for all $i$,
\[
 E^i=\operatorname{pr}_2^*(F^i),
\]
where $F^i$ is a complex vector bundle of finite rank on $V$. By \hyperlink{sga6.e02.para.3-1-5}{(3.1.5)}, $p_*(E^\bullet)$ is associated to the complex of presheaves
\[
 \mathcal C'(S,C^\infty(V,F^\bullet)),
\]
where the differentials are induced by the given differential operators. Put, for each $i\in\mathbb Z$,
\[
 m_i=\sum_{j<i}n_j,
\]
with the convention $n_i=0$ if $E^i=0$. For each $r\in\mathbb Z$ one has, again by \hyperlink{sga6.e02.para.3-1-5}{(3.1.5)}, a complex of presheaves
\[
 \mathcal C'(S,C_{r-m_i}(V,F^i)),
\]
which we denote by $\mathcal C'(S,C_r(V,F^\bullet))$. These complexes form a projective system, and
\sgaStableTarget{sga6.e02.eq.3-2-1-1}\begin{equation}\tag{3.2.1.1}
 \mathcal C'(S,C^\infty(V,F^\bullet))
 =\plim_r \mathcal C'(S,C_r(V,F^\bullet)).
\end{equation}
Finally, denote by $F_r^\bullet$ the complex of trivial Hilbert bundles over $S$ defined by the corresponding $\delta_r(d^i)$, and by $\mathcal F_r^\bullet$ the complex of sheaves of sections; by \hyperlink{sga6.e02.para.3-1-5}{(3.1.5)}, it is just the complex of sheaves associated to $\mathcal C'(S,C_r(V,F^\bullet))$. To prove the theorem it is enough to prove the following assertions:

(a) For $r\in\mathbb Z$, the canonical morphism
\[
 \mathcal C'(S,C^\infty(V,F^\bullet))\longrightarrow \mathcal C'(S,C_r(V,F^\bullet))
\]
is a quasi-isomorphism of complexes of presheaves, that is, it induces isomorphisms on the presheaves of cohomology.

(b) For $r\in\mathbb Z$, the complex $\mathcal F_r^\bullet$ is perfect.

Proof of (b). The ellipticity of $E^\bullet$ implies ([16] XI, theorem 6) that, for every $s\in S$, the differential of $(F_r^\bullet)_s$ has closed image and the spaces $H^i((F_r^\bullet)_s)$ are finite-dimensional vector spaces. In the terminology of [11], $F_r^\bullet$ is a direct quasi-acyclic complex of Hilbert bundles. Hence, by [12] 2.3, $F_r^\bullet$ is locally homotopically equivalent to a bounded complex of vector bundles of finite rank. Since the sheaf-of-sections functor is additive \hyperlink{sga6.e02.para.3-0-5}{(3.0.5)}, $\mathcal F_r^\bullet$ is locally homotopically equivalent to a bounded complex of locally free sheaves of finite type; hence it is perfect.

For the proof of (a) we shall need the following intermediate assertion:

(a') For $p\gg q$, the canonical morphism
\[
 F_p^\bullet\longrightarrow F_q^\bullet
\]
is a quasi-isomorphism.

Proof of (a'). Let $M_{pq}^\bullet$ be the complex of Hilbert bundles which is the cone of the morphism $F_p^\bullet\to F_q^\bullet$. The ellipticity of $E^\bullet$ implies that, for every $s\in S$, the morphism $(F_p^\bullet)_s\to(F_q^\bullet)_s$ is a quasi-isomorphism ([16] XI, corollary 2 and theorem 6). In other words, $(M_{pq}^\bullet)_s$ is a direct acyclic complex. By [11] 2.3, $M_{pq}^\bullet$ is locally homotopically trivial; hence its sheaf complex, the cone of $\mathcal F_p^\bullet\to\mathcal F_q^\bullet$, is acyclic.

Proof of (a). It is enough to show that, for every relatively compact open subset $U$ of $S$, the morphism
\[
 C^\infty(U,C^\infty(V,F^\bullet))\longrightarrow C^\infty(U,C_r(V,F^\bullet))
\]
is a quasi-isomorphism. Replacing $S$ by $\overline U$ and $E^\bullet$ by the induced complex over $\overline U$, one is reduced to proving that
\[
 C(S,C^\infty(V,F^\bullet))\longrightarrow C(S,C_r(V,F^\bullet))
\]
is a quasi-isomorphism, with $S$ compact. By \hyperlink{sga6.e02.eq.3-2-1-1}{(3.2.1.1)},
\[
 C(S,C^\infty(V,F^\bullet))=\plim_i C(S,C_i(V,F^\bullet)).
\]
For each $i$, $C(S,C_i(V,F^j))$ is a Banach space for the topology of uniform convergence, and the transition morphisms
\[
 C(S,C_i(V,F^j))\longrightarrow C(S,C_{i-1}(V,F^j))
\]
are continuous linear maps with dense image by \hyperlink{sga6.e02.para.3-0-4}{(3.0.4)} and \hyperlink{sga6.e02.lem.3-1-2}{(3.1.2)}. On the other hand, applying to the quasi-isomorphisms (a') the functor ``sections'', which is exact since $\mathcal O_S$ is soft, one sees that the transition morphisms are quasi-isomorphisms. Assertion (a) follows from EGA $0_{\mathrm I}$ 13.2.3, topological variant. This completes the proof of \hyperlink{sga6.e02.thm.3-2-1}{(3.2.1)}.

\sgaStableTarget{sga6.e02.para.3-2-2}\textbf{3.2.2.} If $S$ is a topological space, we denote by $\Parf(S)$ the triangulated category of perfect complexes of $(\mathcal O_S\otimes\mathbb C)$-modules, and by $K^0(S)$ the corresponding Grothendieck group.\footnote{See \hyperlink{sga6.e01.para.6-3}{I 6.3}. For a general review of $K$-functors, see \hyperlink{sga6.e04.expose}{Exposé IV}.}

Let $S,X,E^\bullet$ be as in \hyperlink{sga6.e02.thm.3-2-1}{(3.2.1)}. The \emph{analytic index} of $E^\bullet$ is the element
\[
 \ind_{\mathrm{an}}(E^\bullet)=[Rp_*(E^\bullet)]\in K^0(S).
\]

If $S$ is compact, $\ind_{\mathrm{an}}(E^\bullet)$ canonically defines an element in the usual group $K(S)$ of topologists. Indeed, in general, for every topological space $S$, let $\Loclib(S)$ be the additive category of locally free $(\mathcal O_S\otimes\mathbb C)$-modules of finite type, and let $K(S)$ be the Grothendieck group of the triangulated category
\[
 K^b(\Loclib(S))/K^b_{\mathrm{ac}}(\Loclib(S)).
\]
This category is in fact identified with $K^b(\Loclib(S))$ when $S$ is paracompact (\hyperlink{sga6.e02.prop.2-2-2}{II 2.2.2(a)}), so in this case, by \hyperlink{sga6.e01.lem.6-4}{I 6.4}, $K(S)$ is canonically isomorphic to the Grothendieck group of the additive category $\Loclib(S)$. There is a natural homomorphism
\[
 K(S)\longrightarrow K^0(S).
\]
By \hyperlink{sga6.e02.prop.2-3-2}{II 2.3.2} this is an isomorphism for compact $S$, proving the assertion. More generally, $\ind_{\mathrm{an}}(E^\bullet)$ defines an element in the completed group
\[
 \widehat K(S)=\plim_A K(A),
\]
the projective limit being taken over the compact subsets $A$ of $S$.

\sgaStableTarget{sga6.e02.para.3-2-3}\textbf{3.2.3.} The writer does not know whether the finiteness \hyperlink{sga6.e02.thm.3-2-1}{theorem (3.2.1)} remains valid if, instead of assuming $S$ locally compact, one merely assumes that each point of $S$ has a fundamental system of paracompact neighborhoods. Moreover, the local compactness hypothesis on $S$ does not seem very natural. The space of ``universal local parameters'' for elliptic complexes of fixed order is in fact a subspace of infinite dimension in a Fréchet space. More precisely, let $V$ be a compact variety, let $F=(F^i)$ be a sequence of complex vector bundles of finite rank on $V$, zero except for finitely many indices, and let $n=(n_i)$ be a sequence of integers $\ge0$. Let $\operatorname{Comp.diff}_n(F)$ be the space of complexes of differential operators of base $F$ and order $n$, that is, the closed subspace of
\[
 \prod_i \Diff^{n_i}(F^i,F^{i+1})
\]
defined by the equations $d^{i+1}d^i=0$. This product is a Fréchet space. The subset $\operatorname{Ell}_n(F)$ formed by the elliptic complexes of order $n$ is open. It is this space, in general not locally compact, which plays the role of a space of universal local parameters.

\sgaStableTarget{sga6.e02.para.3-2-4}\textbf{3.2.4.} Nevertheless, even if the generalization of \hyperlink{sga6.e02.thm.3-2-1}{(3.2.1)} just mentioned is impossible, one can always associate to $E^\bullet$, under the sole hypothesis that the fibration $p$ is proper and locally trivial, an element of $K^0(S)$ deserving the name analytic index of $E^\bullet$. Indeed, suppose, for simplicity, that we are in the local situation where $X=S\times V$ with $V$ compact, and $E^\bullet$ is given by a continuous map from $S$ into $\operatorname{Ell}_n(F)$ \hyperlink{sga6.e02.para.3-2-3}{(3.2.3)}. Assertions (a') and (b) in the proof of \hyperlink{sga6.e02.thm.3-2-1}{(3.2.1)} are in fact valid without any hypothesis on $S$; it is therefore reasonable to put, with the notation of that proof,
\[
 \ind_{\mathrm{an}}(E^\bullet)=[\mathcal F_r^\bullet]\in K^0(S),
\]
this element being independent of $r$.

\sgaStableTarget{sga6.e02.para.3-2-5}\textbf{3.2.5.} Suppose $S$ paracompact. Let $H$ be a complex Hilbert space of infinite Hilbert dimension, and let $I(H)$ be the subspace of $L(H,H)$ formed by the Fredholm operators. The complexes $F_r^\bullet$ define [12] a canonical element in the group $[S,I(H)]$ of homotopy classes of continuous maps from $S$ to $I(H)$. But the writer does not know, for non-compact $S$ (even when $S$ is locally compact), what the relations are between $[S,I(H)]$ and $K^0(S)$: is there a natural arrow from one to the other? are they isomorphic? A fortiori, the relation between the two indices thus defined remains unclear.

For a generalization, under the hypotheses of \hyperlink{sga6.e02.thm.3-2-1}{(3.2.1)}, of the Atiyah--Singer index theorem, we refer to the recent work of Shih [17], or to work in preparation of M. F. Atiyah and I. M. Singer.\footnote{For the general case, use the notion of ``weak Hilbert bundle'', or continuous field of Hilbert spaces, of [12].}

\sgaStableTarget{sga6.e02.bibliography}\subsection*{Bibliography}

\begin{enumerate}[label={[\,\arabic*\,]}]
\item Artin, M., Grothendieck, A., and Verdier, J.-L., \emph{Cohomologie étale des schémas}, to appear with North-Holland Publishing Company (cited as SGA 4).
\item Atiyah, M., and Segal, G., secret papers.
\item Cartan, H., \emph{Funzioni e varietà complesse. Faisceaux analytiques cohérents}, course at the Centro Internazionale Matematico Estivo, 1964, Rome.
\item Cartan, H., and Schwartz, L., \emph{Théorème d'Atiyah--Singer}, Seminar 1963--64, Institut Henri Poincaré, Paris.
\item Frisch, J., Fonctions analytiques sur un ensemble semi-analytique, note in C.R.A.S. Paris, t. 260, pp. 2974--2976 (1965).
\item Godement, R., \emph{Théorie des faisceaux}, Hermann (1958), cited as T.F.
\item Grothendieck, A., \emph{Cohomologie locale des faisceaux cohérents et théorèmes de Lefschetz locaux et globaux}, Séminaire de Géométrie Algébrique de l'I.H.E.S. 1962, to appear with North-Holland Publishing Company, cited as SGA 2.
\item Grothendieck, A., \emph{Fondements de la géométrie analytique}, Séminaire Cartan 1959--60, Institut Henri Poincaré, Paris.
\item Hartshorne, R., \emph{Residues and Duality}, Lecture Notes in Mathematics 20, Springer, cited as [H].
\item Houzel, C., \emph{Géométrie analytique locale}, Séminaire Cartan 1959--60, Institut Henri Poincaré, Paris.
\item Illusie, L., \emph{Complexes de fibrés et définitions du foncteur $K$}, Séminaire Shih Weishu, I.H.E.S. 1964--65.
\item Illusie, L., Complexes quasi-acycliques directs de fibrés banachiques, note in C.R.A.S. Paris, t. 260, pp. 6499--6502.
\item Kiehl, R., Theorem A und Theorem B in der nichtarchimedischen Funktionentheorie, \emph{Inventiones Mathematicae} 2, pp. 256--273 (1967).
\item Lang, S., \emph{Introduction to differentiable manifolds}.
\item Łojasiewicz, S., \emph{Annali della Scuola Normale Superiore di Pisa}, Serie III, 18, fasc. IV, 1964.
\item Palais, R., lectures at the Princeton Seminar 1963--64 on the index theorem.
\item Shih, W., \emph{Fiber cobordism and the index of a family of elliptic differential operators}.
\item Verdier, J.-L., \emph{Les catégories dérivées des catégories abéliennes}, to appear with North-Holland Publishing Company.
\end{enumerate}




\providecommand{\Ob}{\operatorname{Ob}}
\providecommand{\qcoh}{\mathrm{qcoh}}
\providecommand{\coh}{\mathrm{coh}}
\providecommand{\torf}{\mathrm{torf}}
\providecommand{\toramp}{\operatorname{tor.amp}}
\providecommand{\amp}{\operatorname{amp}}
\providecommand{\Id}{\operatorname{Id}}
\providecommand{\pr}{\operatorname{pr}}
\providecommand{\coker}{\operatorname{coker}}
\providecommand{\im}{\operatorname{Im}}
\providecommand{\Spec}{\operatorname{Spec}}


\sgaStableTarget{sga6.e03.expose}\section*{Exposé III. Relative finiteness conditions}
\begin{center}
by L. Illusie
\end{center}

\sgaStableTarget{sga6.e03.sec.1}\section*{1. Relative pseudo-coherence}

\sgaStableTarget{sga6.e03.thm.1-1}\textbf{Theorem 1.1.} Let \(n\in \mathbb Z\). For each locally finite type morphism of schemes
\[
  f:X\longrightarrow Y
\]
there exists one and only one full subcategory \(D(f)_{n\text{-}\coh}\) of \(D(X)\) such that the following conditions are satisfied:

\begin{enumerate}[label=(\roman*)]
\item If \(E\in \Ob D(f)_{n\text{-}\coh}\), then \(i^*(E)\in \Ob D(fi)_{n\text{-}\coh}\) for every open immersion \(i:U\to X\).

\item If \(E\in \Ob D(X)\), and if there exists an open covering \((i_\alpha:U_\alpha\to X)_{\alpha\in A}\) such that \(i_\alpha^*(E)\in \Ob D(fi_\alpha)_{n\text{-}\coh}\) for every \(\alpha\in A\), then \(E\in \Ob D(f)_{n\text{-}\coh}\).

\item If
\[
\begin{tikzcd}
X \arrow[r,"f"] \arrow[dr,"h"'] & Y \arrow[d,"g"]\\
& Z
\end{tikzcd}
\]
is a commutative triangle of locally finite type morphisms, where \(f\) is a closed immersion and \(g\) is smooth, then one has
\[
  E\in \Ob D(h)_{n\text{-}\coh}
  \quad\Longleftrightarrow\quad
  f_*(E)\in D(Y)_{n\text{-}\coh}.
\]
Here \(D(X)\) denotes the derived category of \(\mathcal O_X\)-modules, and \(D(Y)_{n\text{-}\coh}\) denotes the full subcategory of \(D(Y)\) formed by the \(n\)-pseudo-coherent objects.
\end{enumerate}

\emph{Proof.} Uniqueness is evident. Let indeed \(f:X\to Y\) be a locally finite type morphism. Choose an open covering \((i_\alpha:X_\alpha\to X)_{\alpha\in A}\) such that, for every \(\alpha\), \(fi_\alpha\) factors as
\[
  X_\alpha \xrightarrow{i'_\alpha} Y_\alpha \xrightarrow{f'_\alpha} Y,
\]
where \(i'_\alpha\) is a closed immersion and \(f'_\alpha\) is smooth. It follows at once from (i), (ii), (iii) that the objects of \(D(f)_{n\text{-}\coh}\) are precisely the objects \(E\in D(X)\) such that
\[
  i'_{\alpha *}i_\alpha^*(E)\in \Ob D(Y_\alpha)_{n\text{-}\coh};
\]
therefore \(D(f)_{n\text{-}\coh}\) is well determined, since it is a full subcategory of \(D(X)\). To prove existence we shall need several lemmas.

\sgaStableTarget{sga6.e03.lem.1-1-1}\textbf{Lemma 1.1.1.} Let \(i:X\to Y\) be a closed immersion of schemes such that \(i_*(\mathcal O_X)\) is pseudo-coherent, and let \(E\in \Ob D(X)\). In order that \(E\) be \(n\)-pseudo-coherent, it is necessary and sufficient that \(i_*(E)\) be so.

\emph{Proof.} The necessity follows from the general stability criterion for pseudo-coherence under application of a functor. Conversely, suppose that \(i_*(E)\) is \(n\)-pseudo-coherent, and prove that the same is true of \(E\). Since \(i_*(E)\) has locally bounded-above cohomology, the same is true of \(E\). The question being local, one may suppose that \(E\) is \(m\)-pseudo-coherent for \(m\) sufficiently large. We prove by descending induction on \(m\), with \(E\) assumed \(m\)-pseudo-coherent for some \(m\geq n\). Suppose already proved that \(E\) is \((n+1)\)-pseudo-coherent. After localization, there is a distinguished triangle
\[
E'\longrightarrow E\longrightarrow E''\longrightarrow E'[1]
\]
where \(E'\) is perfect and \(H^i(E'')=0\) for \(i>n\). By \hyperlink{sga6.e01.schol.2-6}{Exposé I, Proposition 2.6}, it remains to show that \(E''\) is \(n\)-pseudo-coherent. Applying \(i_*\) to the preceding triangle, we get a distinguished triangle
\[
i_*(E')\longrightarrow i_*(E)\longrightarrow i_*(E'')\longrightarrow i_*(E')[1]
\]
and, by the first part of the proof, \(i_*(E')\) is pseudo-coherent. Hence, again by \hyperlink{sga6.e01.schol.2-6}{Exposé I, Proposition 2.6}, \(i_*(E'')\) is \(n\)-pseudo-coherent. Since \(i_*(E'')\) is zero in degrees \(>n\), this means, by \hyperlink{sga6.e01.prop.2-10}{Exposé I, Proposition 2.10} b), that \(H^n(i_*(E''))\) is of finite type. But
\[
  H^n(i_*(E''))=i_*H^n(E''),
\]
so the assertion is reduced to showing, by \hyperlink{sga6.e01.prop.2-10}{Exposé I, Proposition 2.10} b), that if \(M\) is an \(\mathcal O_X\)-module such that \(i_*(M)\) is of finite type, then \(M\) is of finite type. This is immediate and independent of the hypothesis on \(i_*(\mathcal O_X)\).

\sgaStableTarget{sga6.e03.rem.1-1-2}\textbf{Remark 1.1.2.} The preceding lemma applies in particular to the case where \(i\) is a regular immersion. Indeed, by means of the Koszul complex, \(i_*(\mathcal O_X)\) has locally a finite left resolution by finite-type free \(\mathcal O_Y\)-modules; hence \(i_*(\mathcal O_X)\) is perfect, and a fortiori pseudo-coherent.

\sgaStableTarget{sga6.e03.lem.1-1-3}\textbf{Lemma 1.1.3.} Let
\[
\begin{tikzcd}
X' \arrow[r,"f"] \arrow[d,"i'"'] & X \arrow[d,"i"]\\
Y' \arrow[r,"g"'] & Y
\end{tikzcd}
\]
be a Cartesian square of schemes, where \(i\) is a closed immersion, and let \(E\) be a complex of \(\mathcal O_X\)-modules. Then the canonical base-change arrow
\[
  g^*i_*(E)\longrightarrow i'_*f^*(E)
\]
is an isomorphism.

If \(E'=f^*(E)\), if \(g\) is flat and \(i_*(E)\) is \(n\)-pseudo-coherent, then \(i'_*(E')\) is \(n\)-pseudo-coherent. Conversely, if \(g\) is faithfully flat and \(i'_*(E')\) is \(n\)-pseudo-coherent, then \(i_*(E)\) is \(n\)-pseudo-coherent.

\emph{Proof.} The assertion is verified trivially on the fibers.

\sgaStableTarget{sga6.e03.lem.1-1-4}\textbf{Lemma 1.1.4.} Let
\[
\begin{tikzcd}
X' & X \arrow[l,"i'"'] \arrow[r,"i''"] & X''\\
& Y \arrow[ul,"f'"] \arrow[ur,"f''"'] &
\end{tikzcd}
\]
be a commutative diagram of schemes, where \(i'\) (respectively \(i''\)) is a closed immersion and \(f'\) (respectively \(f''\)) is smooth, and let \(E\in \Ob D(X)\). Then
\[
  i'^*(E)\text{ is }n\text{-pseudo-coherent}
  \quad\Longleftrightarrow\quad
  i''^*(E)\text{ is }n\text{-pseudo-coherent}.
\]

\emph{Proof.} Following a device due to Lichtenbaum, send \(X\) diagonally into \(X'\times_YX''\). One is at once reduced to proving the assertion in the case where one of the morphisms \(f'\), \(f''\) is the identity. Changing notation, one has a commutative triangle
\[
\begin{tikzcd}
& Y \arrow[d,"f"]\\
X \arrow[r,"i"'] \arrow[ur,"j"] & Z
\end{tikzcd}
\]
where \(i\) and \(j\) are closed immersions and \(f\) is smooth, and one must show that \(i_*(E)\) is \(n\)-pseudo-coherent if and only if \(j_*(E)\) is \(n\)-pseudo-coherent. The question being local on \(Y\) and on \(Z\) in a neighborhood of a point of \(X\), one may suppose that \(X,Y,Z\) are affine and that there is a factorization
\[
\begin{tikzcd}
Y \arrow[r,"i'"] \arrow[d,"f"'] & Z[T_1,\ldots,T_r]\arrow[dl]\\
Z&
\end{tikzcd}
\]
where \(i'\) is a closed immersion and the oblique arrow is the canonical projection. Then, by \hyperlink{sga6.e07.prop.1-10}{Exposé VII, 1.10}, \(i'\) is a regular immersion; hence by \hyperlink{sga6.e03.lem.1-1-1}{Lemma 1.1.1} an object \(F\) of \(D(Y)\) is \(n\)-pseudo-coherent if and only if \(i'_*(F)\) is so.

Thus one may suppose that \(Y=Z[T_1,\ldots,T_r]\), \(f\) being the canonical projection. The preceding diagram then corresponds to a diagram of rings
\[
\begin{tikzcd}
A & C \arrow[l,"v"'] \arrow[r,"u"] & B=C[T_1,\ldots,T_r]
\end{tikzcd}
\]
where \(u\) and \(v\) are epimorphisms and \(c:C\to B\) is the canonical inclusion. For \(1\leq i\leq r\), put \(u(T_i)=x_i\), and let \(z_i\in C\) be such that \(v(z_i)=x_i\). The retraction
\[
  \sigma:B\longrightarrow C,\qquad \sigma(T_i)=z_i,
\]
is defined by \(c(T_i)=z_i\), is such that \(v\sigma=u\). Geometrically, it corresponds to a section \(s\) of \(f\) such that \(i=sj\). By \hyperlink{sga6.e07.prop.1-10}{Exposé VII, 1.10}, \(s\) is a regular immersion; the assertion therefore follows from \hyperlink{sga6.e03.lem.1-1-1}{Lemma 1.1.1}.

\emph{End of the proof of \hyperlink{sga6.e03.thm.1-1}{Theorem 1.1}.} Let \(f:X\to Y\) be a locally finite type morphism. There exists an open covering \((X_\alpha)\) of \(X\) such that each \(X_\alpha\) embeds by a closed immersion into a scheme \(Y_\alpha\) smooth over \(Y\). Put \(\coprod X_\alpha=X'\), \(\coprod Y_\alpha=Y'\). We have a commutative diagram
\[
\begin{tikzcd}
X' \arrow[r,"i'"] \arrow[d,"p'"'] & Y' \arrow[dl,"f'"]\\
X \arrow[d,"f"']\\
Y
\end{tikzcd}
\]
where \(p'\) is the morphism which is the sum of a surjective family of open immersions (we shall say \emph{special}, for short), \(i'\) is a closed immersion, and \(f'\) is smooth. We define \(D(f)_{n\text{-}\coh}\) as the full subcategory of \(D(X)\) formed by the objects \(E\) such that \(i'_*p'^*(E)\) is \(n\)-pseudo-coherent. First we show that this definition makes sense, that is, that \(D(f)_{n\text{-}\coh}\) is independent of the choice of \(i'\), \(p'\), \(f'\). Let
\[
\begin{tikzcd}
X'' \arrow[r,"i''"] \arrow[d,"p''"'] & Y'' \arrow[dl,"f''"]\\
X \arrow[d,"f"']\\
Y
\end{tikzcd}
\]
be an analogous commutative diagram, with \(p''\) special, \(i''\) a closed immersion, and \(f''\) smooth. From the two diagrams one deduces a commutative diagram
\[
\begin{tikzcd}
X'\times_X X'' \arrow[r] \arrow[d] & Y'\times_Y Y'' \arrow[d]\\
X' \arrow[r,"i'"] \arrow[d] & Y' \arrow[d,"f'"]\\
X'' \arrow[r,"i''"'] \arrow[d] & Y'' \arrow[d,"f''"']\\
X \arrow[r,"f"'] & Y
\end{tikzcd}
\]
where the horizontal arrows, except for \(f\), are closed immersions, and the arrows of the left (respectively right) Cartesian square are special morphisms (respectively smooth morphisms). Thus one is reduced to proving the following assertion.

\[
\begin{tikzcd}
X_1 \arrow[r,"i_1"] \arrow[d,"p"'] & Y_1 \arrow[d,"q"]\\
X_2 \arrow[r,"i_2"'] & Y_2
\end{tikzcd}
\]
Let this be a commutative square of schemes, where \(p\) is special, \(q\) is smooth, \(i_1\) and \(i_2\) are closed immersions, and let \(E\in \Ob D(X_2)\). Then \(i_{2*}(E)\) is \(n\)-pseudo-coherent if and only if \(i_{1*}p^*(E)\) is so.

\emph{Proof of the assertion.} By definition, \(p\) is the sum of a family of opens \(U_\alpha\) of \(X_2\) covering \(X_2\). Each \(U_\alpha\) is induced by an open \(V_\alpha\) of \(Y_2\), and the \(V_\alpha\) cover an open neighborhood \(X'_2\) of \(X_2\) in \(Y_2\). There is therefore a Cartesian commutative diagram
\[
\begin{tikzcd}
X_1 \arrow[r,"i'_1"] \arrow[d,"p"'] & X'_1 \arrow[d,"p'"]\\
X_2 \arrow[r,"i'_2"'] & X'_2 \arrow[r,"j"'] & Y_2
\end{tikzcd}
\]
where \(i'_1\) and \(i'_2\) are closed immersions, \(p\) and \(p'\) are special morphisms, and \(j\) is an open immersion. The condition that \(i_{2*}(E)\) be \(n\)-pseudo-coherent is plainly equivalent to the condition that \(i'_{2*}(E)\) be \(n\)-pseudo-coherent. On the other hand, by \hyperlink{sga6.e03.lem.1-1-3}{Lemma 1.1.3}, the condition that \(i'_{2*}(E)\) be \(n\)-pseudo-coherent is equivalent to the condition that \(i'_{1*}p^*(E)\) be \(n\)-pseudo-coherent. Applying \hyperlink{sga6.e03.lem.1-1-4}{Lemma 1.1.4} to the factorizations \((jp')i'_1\) and \(qi_1\) of \(i_2p\), and to \(p^*(E)\in \Ob D(X_1)\), one finally concludes that \(i'_{2*}(E)\) is \(n\)-pseudo-coherent if and only if \(i_{1*}p^*(E)\) is \(n\)-pseudo-coherent, as required.

We have thus shown that the category \(D(f)_{n\text{-}\coh}\) defined above by means of a diagram as above depends in fact only on \(f\). Moreover it is trivial from the definition that it satisfies (i) to (iii). This proves \hyperlink{sga6.e03.thm.1-1}{Theorem 1.1}.

\sgaStableTarget{sga6.e03.def.1-2}\textbf{Definition 1.2.} Let \(f:X\to Y\) be a locally finite type morphism, let \(E\in \Ob D(X)\), and let \(n\in\mathbb Z\). We say that \(E\) is \(n\)-pseudo-coherent relative to \(f\) if \(E\in \Ob D(f)_{n\text{-}\coh}\) (cf. \hyperlink{sga6.e03.thm.1-1}{1.1}). We say that \(E\) is pseudo-coherent, or \((-\infty)\)-pseudo-coherent, relative to \(f\) if \(E\) is \(n\)-pseudo-coherent relative to \(f\) for every \(n\in\mathbb Z\). We write \(D(f)_{\mathrm{coh}}\) for the full subcategory of \(D(X)\) formed by the objects pseudo-coherent relative to \(f\).

We shall say that \(f\) is \(n\)-pseudo-coherent (respectively pseudo-coherent) if \(\mathcal O_X\) is \(n\)-pseudo-coherent (respectively pseudo-coherent) relative to \(f\). When this cannot lead to confusion, we shall sometimes say ``\(n\)-pseudo-coherent (respectively pseudo-coherent) relative to \(Y\)'' instead of ``relative to \(f\),'' and we shall also write
\[
  D(X)_{Y,n\text{-}\coh}\quad(\text{respectively }D(X)_{Y,\coh})
\]
instead of
\[
  D(f)_{n\text{-}\coh}\quad(\text{respectively }D(f)_{\mathrm{coh}}).
\]
Finally, for \(*=+,-,b\), put
\[
  D^*(f)_{n\text{-}\coh}=D^*(X)\cap D(f)_{n\text{-}\coh},
  \qquad
  D^*(f)_{\mathrm{coh}}=D^*(X)\cap D(f)_{\mathrm{coh}}.
\]

\sgaStableTarget{sga6.e03.rem.1-2-1}\textbf{Remarks 1.2.1.}
\begin{enumerate}[label=\alph*)]
\item The categories \(D(f)_{n\text{-}\coh}\), \(D(f)_{\mathrm{coh}}\) are not in general contained in \(D^{-}(X)\). This is the case, however, if \(X\) is quasi-compact, since an object of \(D(X)\) with locally bounded-above cohomology then has bounded-above cohomology.

\item The objects of \(D(f)_{\mathrm{coh}}\) have quasi-coherent cohomology. Indeed, if \(E\in \Ob D(f)_{\mathrm{coh}}\), the question being local, one may suppose that \(f\) is a closed immersion; then \(f_*(E)\) is pseudo-coherent, hence has quasi-coherent cohomology, and one uses the fact that an \(\mathcal O_X\)-module \(M\) is quasi-coherent if and only if \(f_*(M)\) is.

\item When it is a matter of an \(\mathcal O_X\)-module, we shall use the terminology ``of finite presentation relative to \(f\)'' in preference to ``\((-1)\)-pseudo-coherent relative to \(f\).'' For an \(\mathcal O_X\)-module \(E\) to be of finite presentation relative to \(f\), it is necessary and sufficient, by \hyperlink{sga6.e01.prop.2-9}{Exposé I, Proposition 2.9}, that there exist locally a closed immersion \(i:X\to X'\), with \(X'\) smooth over \(Y\), such that \(i_*(E)\) is of finite presentation in the ordinary sense. In particular, to say that \(f\) is \((-1)\)-pseudo-coherent is equivalent to saying that \(f\) is locally of finite presentation in the sense of EGA IV 1.4.
\end{enumerate}

\sgaStableTarget{sga6.e03.ex.1-3}\textbf{Examples 1.3.} A smooth morphism and a regular immersion are pseudo-coherent morphisms. The same is true of an arbitrary locally finite type morphism \(f:X\to Y\) if \(Y\) is locally noetherian.

\sgaStableTarget{sga6.e03.prop.1-4}\textbf{Proposition 1.4.} Let \(f:X\to Y\) be a locally finite type morphism and let
\[
E\longrightarrow F\longrightarrow G\longrightarrow E[1]
\]
be a distinguished triangle of \(D(X)\). The following table, where \(n\in \mathbb Z\cup\{-\infty\}\) indicates a degree of pseudo-coherence relative to \(f\), has the property that each line represents an implication whose conclusion is underlined:
\[
\begin{array}{ccc}
E&F&G\\
n+1&n&\underline n\\
n&\underline n&n\\
\underline n&n&n-1.
\end{array}
\]
In particular, \(D(f)_{\mathrm{coh}}\) is a triangulated subcategory of \(D(X)\).

\emph{Proof.} By the properties (i) to (iii) of \hyperlink{sga6.e03.thm.1-1}{Theorem 1.1}, one may suppose that \(f\) is a closed immersion. We then have a distinguished triangle
\[
f_*(E)\longrightarrow f_*(F)\longrightarrow f_*(G)\longrightarrow f_*(E)[1]
\]
and the assertion follows from \hyperlink{sga6.e01.schol.2-6}{Exposé I, Proposition 2.6}.

\sgaStableTarget{sga6.e03.def.1-5}\textbf{Definition 1.5.} Let \(X\) and \(Y\) be two \(S\)-schemes. We shall say that \(X\) and \(Y\) are \emph{Tor-independent over \(S\)} if
\[
  \Tor_i^S(\mathcal O_X,\mathcal O_Y)=0\quad\text{for }i\ne 0,
\]
where the \(\Tor_i^S\) are the local hyper-Tor sheaves defined in EGA III 6.5. This is a condition of local nature. If \(X,Y,S\) are affine, with respective rings \(B,C,A\), it means that the canonical map
\[
  B\otimes_A^{\mathbf L}C\longrightarrow B\otimes_A C
\]
is an isomorphism. It is therefore satisfied in particular if \(X\) or \(Y\) is flat over \(S\).

\sgaStableTarget{sga6.e03.lem.1-5-1}\textbf{Lemma 1.5.1.} Let
\[
\begin{tikzcd}
X & X' \arrow[l] & X'' \arrow[l]\\
Y \arrow[u] & Y' \arrow[l] \arrow[u] & Y'' \arrow[l] \arrow[u]
\end{tikzcd}
\]
be Cartesian squares of schemes. If \(X\) and \(Y'\) are Tor-independent over \(Y\), and if \(X'\) and \(Y''\) are Tor-independent over \(Y'\), then \(X\) and \(Y''\) are Tor-independent over \(Y\).

\emph{Proof.} This is an easy exercise on associativity of the tensor product.

\sgaStableTarget{sga6.e03.notation.1-6}\textbf{Notation 1.6.} Let \(X\) and \(Y\) be two \(S\)-schemes, \(E\in \Ob D^{-}(X)\), \(F\in \Ob D^{-}(Y)\). We put
\[
  E\otimes_S^{\mathbf L}F
  =
  L\pr_1^*(E)\otimes^{\mathbf L}L\pr_2^*(F),
\]
where \(\pr_1,\pr_2\) are the two projections from \(X\times_SY\).

\sgaStableTarget{sga6.e03.lem.1-7}\textbf{Lemma 1.7.} Let \(S\) be a scheme, and for \(i=1,2\) let
\[
  f_i:X_i\longrightarrow Y_i
\]
be a closed \(S\)-immersion. Suppose that \(X_1\) and \(X_2\), and also \(Y_1\) and \(Y_2\), are Tor-independent over \(S\). Then, for \(E_i\in \Ob D^{-}(X_i)_{\mathrm{qcoh}}\), there is a canonical isomorphism
\[
  f_{1*}(E_1)\otimes_S^{\mathbf L} f_{2*}(E_2)
  \xrightarrow{\sim}
  (f_1\times_S f_2)_*(E_1\otimes_S^{\mathbf L}E_2).
\]

\emph{Proof.} Put \(f_1\times_Sf_2=f\). One easily defines, by adjunction, a canonical arrow
\[
  f_{1*}(E_1)\otimes_S^{\mathbf L} f_{2*}(E_2)
  \longrightarrow
  f_*(E_1\otimes_S^{\mathbf L}E_2).
\]
We show that it is an isomorphism. The question being local, one may suppose \(S,X_i,Y_i\) affine with respective rings \(C,A_i,B_i\). By dévissage, one is reduced to \(E_i=\mathcal O_{X_i}\). One must then verify that
\[
  A_1\otimes^{\mathbf L}_{B_1}B\otimes^{\mathbf L}_{B_2}A_2
  \longrightarrow A
\]
is an isomorphism, where \(A=A_1\otimes_C A_2\), \(B=B_1\otimes_C B_2\). In view of the Tor-independence hypotheses, this follows at once from associativity of the tensor product.

\sgaStableTarget{sga6.e03.prop.1-8}\textbf{Proposition 1.8.} Let \(S\) be a scheme, and for \(i=1,2\) let
\[
  f_i:X_i\longrightarrow Y_i
\]
be a locally finite type \(S\)-morphism. Suppose that \(X_1\) and \(X_2\), and also \(Y_1\) and \(Y_2\), are Tor-independent over \(S\). Let \([a_i,b_i]\) be intervals of \(\mathbb Z\), and let
\[
  E_i\in \Ob D^{-}(X_i)_{\mathrm{qcoh}}
\]
be a complex which is \(a_i\)-pseudo-coherent relative to \(f_i\) and acyclic in degree \(>b_i\). Then
\[
  E_1\otimes_S^{\mathbf L}E_2
\]
is \(a\)-pseudo-coherent relative to \(f=f_1\times_S f_2\) and acyclic in degree \(>b\), where
\[
  a=\sup(a_1+b_2,a_2+b_1),\qquad b=b_1+b_2.
\]

\emph{Proof.} The question being local on \(X_1\times_SX_2\), one may suppose that \(f_i\) factors as
\[
  X_i\xrightarrow{i_i}Y'_i\xrightarrow{p_i}Y_i,
\]
where \(i_i\) is a closed immersion and \(p_i\) a smooth morphism. Put \(f'=f'_1\times_Sf'_2\). We must show that \(f'_*(E_1\otimes_S^{\mathbf L}E_2)\) is \(a\)-pseudo-coherent and acyclic in degree \(>b\). By \hyperlink{sga6.e03.lem.1-5-1}{Lemma 1.5.1}, \(Y'_1\) and \(Y'_2\) are Tor-independent over \(S\). After changing notation, one may suppose that \(f_i=f'_i\), in other words that the \(f_i\) are closed immersions. By hypothesis, \(f_{i*}(E_i)\) is \(a_i\)-pseudo-coherent and acyclic in degree \(>b_i\); hence by \hyperlink{sga6.e01.cor.2-16}{Exposé I, Proposition 2.16},
\[
  f_{1*}(E_1)\otimes_S^{\mathbf L}f_{2*}(E_2)
\]
is \(a\)-pseudo-coherent and acyclic in degree \(>b\). One concludes by \hyperlink{sga6.e03.lem.1-7}{Lemma 1.7}.

\sgaStableTarget{sga6.e03.cor.1-9}\textbf{Corollary 1.9.} Under the hypotheses of \hyperlink{sga6.e03.prop.1-8}{Proposition 1.8}, the functor
\[
  \otimes_S^{\mathbf L}:D^{-}(X_1)\times D^{-}(X_2)\longrightarrow D^{-}(X)
\]
induces a functor
\[
  \otimes_S^{\mathbf L}:D^{-}(f_1)_{\mathrm{coh}}\times D^{-}(f_2)_{\mathrm{coh}}
  \longrightarrow D^{-}(f)_{\mathrm{coh}} .
\]
In particular, if \(f_1\) and \(f_2\) are pseudo-coherent, then so is \(f\).

\sgaStableTarget{sga6.e03.cor.1-10}\textbf{Corollary 1.10.} Let
\[
\begin{tikzcd}
X' \arrow[r,"g"] \arrow[d,"f'"'] & X \arrow[d,"f"]\\
Y' \arrow[r,"h"'] & Y
\end{tikzcd}
\]
be a Cartesian square, with \(f\) locally of finite type. Then:
\begin{enumerate}[label=\alph*)]
\item Suppose that \(X\) and \(Y'\) are Tor-independent over \(Y\). For \(n\in\mathbb Z\cup\{-\infty\}\), the functor \(Lg^*:D^{-}(X)\to D^{-}(X')\) induces a functor
\[
  Lg^*:D^{-}(f)_{n\text{-}\coh}\longrightarrow D^{-}(f')_{n\text{-}\coh}.
\]
In particular, if \(f\) is \(n\)-pseudo-coherent, then so is \(f'\).

\item Suppose that \(h\) is faithfully flat. Let \(n\in\mathbb Z\cup\{-\infty\}\), and let \(E\in\Ob D^{-}(X)\) be a complex with quasi-coherent cohomology. In order that \(E\) be \(n\)-pseudo-coherent relative to \(f\), it is necessary and sufficient that \(g^*(E)\) be \(n\)-pseudo-coherent relative to \(f'\). In particular, in order that \(f\) be \(n\)-pseudo-coherent, it is necessary and sufficient that \(f'\) be so.
\end{enumerate}

\emph{Proof.} Taking \hyperlink{sga6.e03.rem.1-2-1}{Remark 1.2.1} b) into account, all assertions of Corollaries \hyperlink{sga6.e03.cor.1-9}{1.9} and \hyperlink{sga6.e03.cor.1-10}{1.10} follow trivially from \hyperlink{sga6.e03.prop.1-8}{Proposition 1.8}, except the necessity in \hyperlink{sga6.e03.cor.1-10}{1.10} b). For this one may, the question being local on \(X\), suppose that \(f\) factors as a closed immersion followed by a smooth morphism. One is then reduced by \hyperlink{sga6.e03.lem.1-1-3}{Lemma 1.1.3} to the case where \(f\) (respectively \(f'\)) is the identity of \(X\) (respectively \(X'\)). The assertion then follows, by a standard dévissage, from the fact that a quasi-coherent sheaf \(L\) on \(X\) is of finite type if and only if \(g^*L\) is of finite type.

\sgaStableTarget{sga6.e03.prop.1-11}\textbf{Proposition 1.11.} Let
\[
\begin{tikzcd}
X \arrow[r,"f"] \arrow[dr,"h"'] & Y \arrow[d,"g"]\\
& Z
\end{tikzcd}
\]
be a commutative triangle of locally finite type morphisms. Let \(E\in\Ob D(X)\), \(F\in\Ob D(Y)\), and let \([a,b]\), \([c,d]\) be intervals of \(\mathbb Z\).

\begin{enumerate}[label=\alph*)]
\item If \(E\) is \(a\)-pseudo-coherent relative to \(f\) and acyclic in degree \(>b\), and if \(F\) is \(c\)-pseudo-coherent relative to \(g\) and acyclic in degree \(>d\), then
\[
  E\otimes_{\mathcal O_Y}^{\mathbf L}F
\]
is \(m\)-pseudo-coherent relative to \(h\) and acyclic in degree \(>n\), where
\[
  m=\sup(a+d,b+c),\qquad n=b+d.
\]

\item If \(E\) is \(a\)-pseudo-coherent relative to \(h\) and acyclic in degree \(>b\), and if \(g\) is \((a+1-b)\)-pseudo-coherent, then \(E\) is \(a\)-pseudo-coherent relative to \(f\).
\end{enumerate}

\emph{Proof.} The question being local on \(X\), one may suppose that there is a Cartesian commutative diagram
\[
\begin{tikzcd}[column sep=small]
X \arrow[r,"i"] \arrow[d,"f"'] & E^r_Y \arrow[r,"j"] \arrow[d,"p"] & E^{r+s}_Z \arrow[d,"q"]\\
Y \arrow[r,"k"] & E^s_Z \arrow[r] & Z
\end{tikzcd}
\]
where \(i,j,k\) are closed immersions and \(p,q\) are the canonical projections of affine spaces. If \(F\) is \(c\)-pseudo-coherent relative to \(Z\), then \(k_*(F)\) is \(c\)-pseudo-coherent by \hyperlink{sga6.e03.thm.1-1}{Theorem 1.1 (iii)}, and \(j_*p^*(F)\) is \(c\)-pseudo-coherent by \hyperlink{sga6.e03.lem.1-1-3}{Lemma 1.1.3}. Similarly, if \(g\) is \(x\)-pseudo-coherent, then \(j_*(\mathcal O_{E^r_Y})\) is \(x\)-pseudo-coherent. By the definition of relative pseudo-coherence, one is thus reduced to proving a) and b) when \(f\) and \(g\) are closed immersions.

\emph{Proof of a).} By \hyperlink{sga6.e03.thm.1-1}{Theorem 1.1 (iii)}, one has to show the following. Suppose \(f_*(E)\) is \(a\)-pseudo-coherent and acyclic in degree \(>b\), and \(g_*(F)\) is \(c\)-pseudo-coherent and acyclic in degree \(>d\). Then
\[
  h_*(E\otimes_Y^{\mathbf L}F)
\]
is \(m\)-pseudo-coherent and acyclic in degree \(>n\). There is a canonical isomorphism
\[
  f_*(E)\otimes_Y^{\mathbf L}F\xrightarrow{\sim} f_*(E\otimes_Y^{\mathbf L}F),
\]
a special case of the isomorphism used in the proof of \hyperlink{sga6.e03.prop.1-8}{Proposition 1.8}. Since the functor \(g_*(-\otimes^{\mathbf L}F)\) takes finite-type locally free sheaves to \(c\)-pseudo-coherent complexes, and complexes acyclic in degree \(>0\) to complexes acyclic in degree \(>d\), the assertion follows from \hyperlink{sga6.e01.prop.2-13}{Exposé I, Proposition 2.13}.

\emph{Proof of b).} One must prove the following assertion: suppose \(f_*(E)\) is \(a\)-pseudo-coherent and acyclic in degree \(>b\), and \(g_*(\mathcal O_Y)\) is \((a+1-b)\)-pseudo-coherent. Then \(f_*(E)\) is \(a\)-pseudo-coherent. Since the assertion concerns \(f_*(E)\), one may suppose \(X=Y\) and \(f=\Id_X\). One then proceeds as in the proof of \hyperlink{sga6.e03.lem.1-1-1}{Lemma 1.1.1}, proving by descending induction that \(E\) is pseudo-coherent in degree \(a\). Suppose already proved that \(E\) is \((a+1)\)-pseudo-coherent. By \hyperlink{sga6.e01.prop.2-10}{Exposé I, Proposition 2.10}, after localizing there is a distinguished triangle
\[
E'\longrightarrow E\longrightarrow E''\longrightarrow E'[1]
\]
where \(E'^i\) is locally free for \(i\in [a+1,b]\), \(E'^i=0\) for \(i\notin [a+1,b]\), and \(E''^i=0\) for \(i>a\). The hypothesis that \(g_*(\mathcal O_Y)\) is \((a+1-b)\)-pseudo-coherent implies, by \hyperlink{sga6.e01.prop.2-11}{Exposé I, Proposition 2.11} a), that \(g_*(E')\) is \((a+1)\)-pseudo-coherent. Since \(g_*(E)\) is \(a\)-pseudo-coherent, it follows from \hyperlink{sga6.e01.schol.2-6}{Exposé I, Proposition 2.6}, that \(g_*(E'')\) is \(a\)-pseudo-coherent. We conclude by \hyperlink{sga6.e01.prop.2-10}{Exposé I, Proposition 2.10} b), and by the fact that an \(\mathcal O_Y\)-module \(M\) is of finite type if and only if \(g_*(M)\) is. Corollaries \hyperlink{sga6.e03.cor.1-12}{1.12}--\hyperlink{sga6.e03.cor.1-14}{1.14} follow immediately.

\sgaStableTarget{sga6.e03.cor.1-12}\textbf{Corollary 1.12.} In the situation of \hyperlink{sga6.e03.prop.1-11}{Proposition 1.11}, if \(g\) is pseudo-coherent, then, for \(r\in\mathbb Z\cup\{-\infty\}\),
\[
  D(f)_{r\text{-}\coh}=D(h)_{r\text{-}\coh}.
\]
In particular, taking \(X=Y\), the notions of \(r\)-pseudo-coherence over \(Y\), in the absolute sense or relative to \(g\), coincide.

\sgaStableTarget{sga6.e03.cor.1-13}\textbf{Corollary 1.13.} In the situation of \hyperlink{sga6.e03.prop.1-11}{Proposition 1.11}, if \(f\) is pseudo-coherent, the functor
\[
  \mathcal O_X\otimes_Y^{\mathbf L}-=Lf^*:D^{-}(Y)\longrightarrow D^{-}(X)
\]
induces, for \(r\in\mathbb Z\cup\{-\infty\}\), a functor
\[
  Lf^*:D^{-}(g)_{r\text{-}\coh}\longrightarrow D^{-}(h)_{r\text{-}\coh}.
\]

\sgaStableTarget{sga6.e03.cor.1-14}\textbf{Corollary 1.14.} In the situation of \hyperlink{sga6.e03.prop.1-11}{Proposition 1.11}, if \(f\) is \(r\)-pseudo-coherent and \(g\) is \(s\)-pseudo-coherent, with \(r,s\in\mathbb Z\cup\{-\infty\}\), then \(h\) is \(\sup(r,s)\)-pseudo-coherent. If \(h\) is \(r\)-pseudo-coherent and \(g\) is \((r+1)\)-pseudo-coherent, then \(f\) is \(r\)-pseudo-coherent.

\sgaStableTarget{sga6.e03.rem.1-14-1}\textbf{Remark 1.14.1.} In particular, for \(r=s=-1\), one recovers the assertions (ii) and (v) of EGA IV 1.4.3.

\sgaStableTarget{sga6.e03.sec.2}\section*{2. The finiteness theorem}

\sgaStableTarget{sga6.e03.conj.2-1}\textbf{Conjecture 2.1.} Let \(S\) be a scheme, let \(f:X\to Y\) be a proper \(S\)-morphism of schemes locally of finite type over \(S\), and let \(E\in \Ob D^+(X)\). If \(E\) is pseudo-coherent relative to \(S\), then the same is true of \(\mathbf R f_*(E)\).

\sgaStableTarget{sga6.e03.thm.2-2}\textbf{Theorem 2.2.} The preceding conjecture is true in the following cases:
\[
  \text{(i) } S\text{ is locally noetherian};\qquad
  \text{(ii) } f\text{ is projective}.
\]

\emph{Proof in case (i).} Since \(S\) is locally noetherian, the projection of \(X\) (respectively \(Y\)) onto \(S\) is a pseudo-coherent morphism, and the notions of pseudo-coherence over \(X\) (respectively \(Y\)) and relative to \(S\) coincide (\hyperlink{sga6.e03.cor.1-12}{Corollary 1.12)}. On the other hand, since \(X\) (respectively \(Y\)) is locally noetherian, hence \(\mathcal O_X\) (respectively \(\mathcal O_Y\)) coherent, an object \(M\) of \(D(X)\) (respectively \(D(Y)\)) is pseudo-coherent if and only if its \(H^i(M)\) are coherent and \(M\) is locally in \(D^{-}\) (\hyperlink{sga6.e01.cor.3-5}{Exposé I, 3.5)}. The theorem therefore follows from:

\sgaStableTarget{sga6.e03.thm.2-2-1}\textbf{Theorem 2.2.1.} Let \(f:X\to Y\) be a proper morphism of locally noetherian schemes, and let \(E\in \Ob D^+(X)\). If the \(H^i(E)\) are coherent, then the same is true of \(R^if_*(E)\). If moreover \(E\) is locally in \(D^{-}\), then the same is true of \(\mathbf Rf_*(E)\).

\emph{Proof.} The first assertion follows from the spectral sequence
\[
  R^pf_*(H^q(E)) \Longrightarrow R^{p+q}f_*(E)
\]
and from the finiteness theorem EGA III 3.2.1. For the second, one may suppose \(Y\) affine. Then \(X\) is quasi-compact and separated; suppose it covered by \(r\) affine opens. One knows (EGA III 1.4.12) that under these conditions, for a quasi-coherent \(\mathcal O_X\)-module \(M\), one has \(R^if_*(M)=0\) for every \(i>r\). If \(E\) is locally in \(D^{-}\), then \(E\in D^{-}\), since \(X\) is quasi-compact; say \(H^i(E)=0\) for \(i>N\). The preceding spectral sequence then gives \(R^if_*(E)=0\) for \(i>N+r\), which proves the assertion.

\emph{Proof in case (ii).} The question being local on \(Y\), one may suppose \(S\) and \(Y\) affine. There is then a Cartesian commutative diagram
\[
\begin{tikzcd}
X \arrow[r,"i"] \arrow[d,"f"'] & \mathbb P^r_Y \arrow[d,"p"] \arrow[r,"j"] & \mathbb P^r_{Y'} \arrow[d,"p'"]\\
Y \arrow[r] & Y' \arrow[r,"g"'] & S
\end{tikzcd}
\]
where \(i,j,j'\) are closed immersions, \(g'\) is smooth and affine, and \(p,p'\) are projections of projective spaces. Suppose that \(E\) is pseudo-coherent relative to \(S\). By \hyperlink{sga6.e03.thm.1-1}{Theorem 1.1 (iii)}, this means that \((ji)_*(E)\) is pseudo-coherent. We must show that \(\mathbf Rf_*(E)\) is pseudo-coherent relative to \(S\), that is, by \hyperlink{sga6.e03.thm.1-1}{Theorem 1.1 (iii)}, that
\[
  j'_*\mathbf Rf_*(E)
\]
is pseudo-coherent. But
\[
  j'_*\mathbf Rf_*(E)\simeq \mathbf Rp'_*\bigl((ji)_*(E)\bigr),
\]
so the assertion is a consequence of the following more precise result.

\sgaStableTarget{sga6.e03.thm.2-2-2}\textbf{Theorem 2.2.2.} Let \(Y\) be an affine scheme, let
\[
  f:X=\mathbb P^r_Y\longrightarrow Y
\]
be the projection of a projective space of finite type, and let \(E\in \Ob D^b(X)_{\mathrm{coh}}\). Then \(\mathbf Rf_*(E)\) is pseudo-coherent. Moreover, if \(E\) is acyclic in degree \(>a\), then \(\mathbf Rf_*(E)\) is acyclic in degree \(>a+r\), and there exists an integer \(N\) such that, for all \(n\geq N\),
\[
  \mathbf Rf_*(E(n)),\qquad E(n)=E\otimes\mathcal O_X(n),
\]
is acyclic in degree \(>a\).

\emph{Proof.} Since \(Y\) is affine, \(X\) is quasi-compact and separated. By \hyperlink{sga6.e02.prop.2-2-9}{Exposé II, 2.2.9} a), \(E\) is isomorphic, in \(D(X)\), to a complex \(F\) bounded above such that, for every \(i\), \(F^i\in \Ob L(X)\), that is, \(F^i\) is a finite direct sum of sheaves of the type \(\mathcal O_X(-n)\), \(n\geq 0\). If \(E\) is acyclic in degree \(>a\), \(F\) may be chosen so that \(F^i=0\) for \(i>a\); this follows formally from the preceding assertion, taking account of \hyperlink{sga6.e01.prop.2-10}{Exposé I, 2.10} a), and \hyperlink{sga6.e02.thm.2-2-8}{Exposé II, 2.2.8} a). Replacing \(E\) by \(F\), one may suppose that \(E^i\in \Ob L(X)\) for every \(i\), and that \(E^i=0\) for \(i>a\). Since \(X\) is covered by \(r+1\) affine opens, EGA III 1.4.10 and the hypercohomology spectral sequence give
\[
  R^if_*(E)=0\qquad\text{for }i>a+r.
\]
Let \(t\leq a\) be an integer and put
\[
  E'=(0\to E^t\to\cdots\to E^a\to 0).
\]
There is an exact sequence
\[
  0\longrightarrow E'\longrightarrow E\longrightarrow E''\longrightarrow 0,
\]
where \(E''\) satisfies the same conditions as \(E\), but with \(a\) replaced by \(t-1\). Applying \(\mathbf Rf_*\), one obtains a distinguished triangle
\[
\mathbf Rf_*(E')\longrightarrow \mathbf Rf_*(E)\longrightarrow \mathbf Rf_*(E'')\longrightarrow \mathbf Rf_*(E')[1]
\]
Since \(t\) can be chosen arbitrarily negative, it is enough, for the first assertion, to prove that \(\mathbf Rf_*(E')\) is \((t+r)\)-pseudo-coherent. But \(\mathbf Rf_*(E'')\) is acyclic in degree \(\geq t+r\), hence \((t+r)\)-pseudo-coherent. Thus one wins by \hyperlink{sga6.e01.schol.2-6}{Exposé I, Proposition 2.6}, if one proves that \(\mathbf Rf_*(E')\) is perfect.

For this one is reduced by dévissage to the case where \(E'\) is concentrated in degree \(0\) and equal to \(\mathcal O_X(-n)\), \(n>0\). Then, by EGA III 2.1.16,
\[
  \mathbf Rf_*(\mathcal O_X(-n))\simeq R^rf_*(\mathcal O_X(-n))[-r].
\]
Therefore \(\mathbf Rf_*(\mathcal O_X(-n))\) is perfect, and this ends the proof of the first assertion of \hyperlink{sga6.e03.thm.2-2-2}{Theorem 2.2.2}. We have already seen that \(\mathbf Rf_*(E)\) is acyclic in degree \(>a+r\). It remains to prove the last assertion. For \(t<-r+a\), \(\mathbf Rf_*(E'')\) is acyclic in degree \(>a\). Taking account of the preceding triangle, one is reduced to doing the proof for \(E'\) instead of for \(E\). Since \(E'\) is bounded above and has components in \(L(X)\), there is an integer \(N\) such that, for all \(n\geq N\), the components of \(E'(n)\) are finite direct sums of sheaves of type \(\mathcal O_X(p)\) with \(p\geq 0\). Then, by EGA III 2.1.16 and the hypercohomology spectral sequence, \(\mathbf Rf_*(E'(n))\) is acyclic in degree \(>a\) for all \(n\geq N\), which completes the proof of \hyperlink{sga6.e03.thm.2-2-2}{Theorem 2.2.2} and hence of \hyperlink{sga6.e03.thm.2-2}{Theorem 2.2}.

We also have the following classical complement to \hyperlink{sga6.e03.thm.2-2}{2.2 (ii)}, which will be given for later reference.

\sgaStableTarget{sga6.e03.cor.2-3}\textbf{Corollary 2.3} (EGA III 2.2.1). Let \(S\) be a quasi-compact scheme, \(f:X\to Y\) a projective \(S\)-morphism of finite type schemes over \(S\), and let
\[
  L=i^*\mathcal O_P(1)
\]
be the invertible sheaf defined by an extension \(i:X\to P=P(F)\), where \(F\) is a quasi-coherent \(\mathcal O_Y\)-module of finite type. For \(M\in \Ob D(X)\) and \(n\in\mathbb Z\), put
\[
  M(n)=M\otimes L^{\otimes n}.
\]
If \(E\in \Ob D^+(X)\) is pseudo-coherent relative to \(S\) and acyclic in degree \(>a\), then there exists an integer \(N\) such that, for all \(n\geq N\), \(\mathbf Rf_*(E(n))\) has the same properties.

\emph{Proof.} One reduces trivially, as in the proof of \hyperlink{sga6.e03.thm.2-2}{Theorem 2.2 (ii)}, to the situation of \hyperlink{sga6.e03.thm.2-2-2}{Theorem 2.2.2}: \(Y=S\), \(f:X=\mathbb P^r_Y\to Y\) the canonical projection, \(L=\mathcal O_X(1)\) the canonical sheaf, \(E\in\Ob D^b(X)\) a pseudo-coherent complex acyclic in degree \(>a\). The conclusion then follows from \hyperlink{sga6.e03.thm.2-2-2}{Theorem 2.2.2}.

\sgaStableTarget{sga6.e03.cor.2-4}\textbf{Corollary 2.4.} Let \(S\) be a quasi-compact scheme and let \(f:X\to Y\) be a proper \(S\)-morphism of schemes of finite type over \(S\). If \(S\) is noetherian, or if \(f\) is projective, the functor
\[
  \mathbf Rf_*:D^+(X)\longrightarrow D^+(Y)
\]
induces a functor
\[
  \mathbf Rf_*:D^b(X)_{S\text{-}\coh}\longrightarrow D^b(Y)_{S\text{-}\coh}.
\]

\emph{Proof.} Consequence of \hyperlink{sga6.e03.thm.2-2}{Theorem 2.2} and of the fact that there is then an integer \(N\) such that, for every quasi-coherent \(\mathcal O_X\)-module \(E\), \(R^if_*(E)=0\) for all \(i>N\).

\sgaStableTarget{sga6.e03.cor.2-5}\textbf{Corollary 2.5.} Let \(f:X\to Y\) be a proper and pseudo-coherent morphism, and let \(E\in \Ob D^+(X)\). Suppose \(Y\) locally noetherian or \(f\) projective. If \(E\) is pseudo-coherent in the absolute sense, then the same is true of \(\mathbf Rf_*(E)\). If \(Y\) is quasi-compact, the functor \(\mathbf Rf_*\) induces a functor
\[
  \mathbf Rf_*:D^b(X)_{\mathrm{coh}}\longrightarrow D^b(Y)_{\mathrm{coh}}.
\]

\emph{Proof.} Consequence of \hyperlink{sga6.e03.thm.2-2}{Theorem 2.2}, \hyperlink{sga6.e03.cor.2-4}{Corollary 2.4}, and \hyperlink{sga6.e03.cor.1-12}{Corollary 1.12}.

\sgaStableTarget{sga6.e03.schol.2-6}\textbf{Scholie 2.6.} To prove \hyperlink{sga6.e03.conj.2-1}{Conjecture 2.1} in general, new techniques will surely have to be introduced. One may hope that they will one day make it possible to prove the analogous conjecture in analytic geometry, for a proper morphism of analytic spaces over a suitably topologized locally ringed topos in \(\mathbb C\)-algebras. For the present this conjecture has been proved only over a complex analytic space, by Grauert's theorem.\footnote{Added in April 1971: see a forthcoming article of R. Kiehl and J.-L. Verdier, where the finiteness theorem is proved by different, simpler techniques and under more general hypotheses.}

\sgaStableTarget{sga6.e03.sec.3}\section*{3. Relative finite Tor-dimension}

\sgaStableTarget{sga6.e03.def.3-1}\textbf{Definition 3.1.} Let
\[
  f:(X,\mathcal O_X)\longrightarrow (Y,\mathcal O_Y)
\]
be a morphism of ringed topoi, let \(E\in\Ob D^{-}(X)\), and let \([m,n]\) be an interval of \(\mathbb Z\). We say that \(E\) is of flat amplitude relative to \(f\) (or to \(Y\), if there is no possible confusion) contained in \([m,n]\), and write
\[
  \toramp_f(E)\subset [m,n]\qquad(\text{or }\toramp_Y(E)\subset [m,n]),
\]
if \(E\), considered as a complex of \(f^{-1}(\mathcal O_Y)\)-modules, is of flat amplitude contained in \([m,n]\) in the sense of \hyperlink{sga6.e01.def.5-2}{Exposé I, 5.2}. We say that \(f\) is of flat amplitude contained in \([-a,0]\), or of Tor-dimension \(\leq a\), if \(\mathcal O_X\) is of flat amplitude relative to \(f\) contained in \([-a,0]\). We say that \(E\) is of finite flat amplitude relative to \(f\) if there exists an interval \([m,n]\) such that \(\toramp_f(E)\subset[m,n]\).

\sgaStableTarget{sga6.e03.notation.3-2}\textbf{Notation 3.2.} We write
\[
  D(f)_{\mathrm{torf}},\qquad D(X)_{f\text{-}\torf},\qquad D(X)_{Y\text{-}\torf}
\]
if no confusion is possible, for the full subcategory of \(D^{-}(X)\) formed by complexes of finite flat amplitude relative to \(f\). By \hyperlink{sga6.e01.para.5-3}{Exposé I, 5.3}, this is a triangulated subcategory of \(D^b(X)\).

% SGA-SRC-001581
\sgaStableTarget{sga6.e03.prop.3-3}\textbf{Proposition 3.3.} With the notation of Definition 3.1, suppose that \(X\) has enough points (SGA 4 IV 6.4). The following conditions are then equivalent:
\begin{enumerate}[label=(\roman*)]
\item \(\toramp_f(E)\subset[m,n]\).
\item For every point \(x\) of \(X\), one has
\[
  \toramp_y(E_x)\subset[m,n],
\]
where \(y=f(x)\), ringed by \(\mathcal O_{Y,y}\).
\item For every object \(V\) of \(Y\) and every \(\mathcal O_V\)-module \(M\), one has
\[
  H^i\left(E\vert_{f^{-1}(V)}
  \otimes_{\mathcal O_V}^{\mathbf L}M\right)=0
  \quad\text{for }i\notin[m,n].
\]
Equivalently, for every \(\mathcal O_V\)-module \(M\),
\[
  E\vert_{f^{-1}(V)}\otimes_{\mathcal O_V}^{\mathbf L}M
  \simeq E\vert_{f^{-1}(V)}\otimes_{\mathcal O_V}^{\mathbf L}L_!f^*(M).
\]
\item The same condition as (iii), but with \(M\) the cokernel of a homomorphism
\[
  \mathcal O_V^q\longrightarrow \mathcal O_V^p,
  \qquad p,q\in\mathbb N.
\]
\end{enumerate}

\emph{Proof.} The equivalence (i)\(\Longleftrightarrow\)(ii) follows from \hyperlink{sga6.e01.rem.5-6-2}{Exposé I, 5.6.2}, and the implications (ii)\(\Rightarrow\)(iii)\(\Rightarrow\)(iv) are trivial. We prove (iv)\(\Rightarrow\)(ii). Let \(x\) be a point of \(X\), \(y=f(x)\). One has to show that
\[
  H^i(E_x\otimes_{\mathcal O_{Y,y}}^{\mathbf L}M)=0
\]
for every \(\mathcal O_{Y,y}\)-module \(M\) and every \(i\notin[m,n]\). Since every \(\mathcal O_{Y,y}\)-module is an inductive filtered limit of modules of finite presentation, one may restrict to \(M\) of finite presentation. If \(M\) is a finitely presented \(\mathcal O_{Y,y}\)-module, the same argument as for ordinary ringed spaces (EGA I, 5.3.9) shows that there exists a \(y\)-pointed object \(V\) of \(Y\) and an \(\mathcal O_V\)-module \(P\), cokernel of a homomorphism
\[
  \mathcal O_V^q\longrightarrow\mathcal O_V^p,
  \qquad p,q\in\mathbb N,
\]
such that \(P_y\simeq M\). One then has
\[
  \bigl(E\vert_{f^{-1}(V)}\otimes_{\mathcal O_V}^{\mathbf L}P\bigr)_x
  \simeq
  E_x\otimes_{\mathcal O_{Y,y}}^{\mathbf L}M,
\]
by the commutation of tensor product with passage to fibers, and hypothesis (iv) implies the desired vanishing.

\sgaStableTarget{sga6.e03.rem.3-3-1}\textbf{Remark 3.3.1.} The editor does not know whether the equivalences (i)\(\Longleftrightarrow\)(iii)\(\Longleftrightarrow\)(iv) are true without the hypothesis that \(X\) has enough points.

\sgaStableTarget{sga6.e03.cor.3-4}\textbf{Corollary 3.4.} Let
\[
\begin{tikzcd}
X \arrow[r,"f"] \arrow[dr,"h"'] & Y \arrow[d,"g"]\\
& Z
\end{tikzcd}
\]
be a commutative triangle of ringed topoi, where \(X\) is assumed to have enough points. Let \(E\in\Ob D(X)\), \(F\in\Ob D(Y)\), and let \([a,b]\), \([c,d]\) be intervals of \(\mathbb Z\). Then
\[
  \toramp_f(E)\subset[a,b]\quad\text{and}\quad \toramp_g(F)\subset[c,d]
\]
imply
\[
  \toramp_h(E\otimes_Y^{\mathbf L}F)\subset[a+c,b+d].
\]

\emph{Proof.} By \hyperlink{sga6.e03.prop.3-3}{Proposition 3.3}, one must show that the relation on the left implies
\[
  H^i\left(\bigl((E\otimes_Y^{\mathbf L}F)\vert_{h^{-1}(Z')}\bigr)
    \otimes_{\mathcal O_{Z'}}^{\mathbf L}M\right)=0
\]
for every \(i\notin[a+c,b+d]\) and every \(\mathcal O_{Z'}\)-module \(M\), where \(Z'\) is an object of \(Z\). Put \(Y'=g^{-1}(Z')\), \(X'=f^{-1}(Y')=h^{-1}(Z')\). Then
\[
  \bigl((E\otimes_Y^{\mathbf L}F)\vert_{X'}\bigr)\otimes_{Z'}^{\mathbf L}M
  \simeq
  (E\vert_{X'})\otimes_{Y'}^{\mathbf L}\bigl((F\vert_{Y'})\otimes_{Z'}^{\mathbf L}M\bigr),
\]
and the assertion follows at once from the hypotheses.

In particular:

\sgaStableTarget{sga6.e03.cor.3-4-1}\textbf{Corollary 3.4.1.} Under the hypotheses of \hyperlink{sga6.e03.cor.3-4}{Corollary 3.4}:
\begin{enumerate}[label=\alph*)]
\item If \(g\) is of Tor-dimension \(\leq n\), then
\[
  \toramp_f(E)\subset[a,b]\Longrightarrow
  \toramp_h(E)\subset[a-n,b].
\]
\item If \(f\) is of Tor-dimension \(\leq m\), then
\[
  \toramp_g(F)\subset[c,d]\Longrightarrow
  \toramp_h(Lf^*(F))\subset[c-m,d].
\]
\end{enumerate}

\sgaStableTarget{sga6.e03.rem.3-4-2}\textbf{Remark 3.4.2.} Let \(f:X\to Y\) be a morphism of finite Tor-dimension and \(E\in \Ob D(X)\). By \hyperlink{sga6.e03.cor.3-4-1}{Corollary 3.4.1} a), if \(E\) is of finite flat amplitude in the absolute sense, then \(E\) is of finite flat amplitude relative to \(f\). But the converse is false, even if \(f\) is flat, as is shown by the example of a singularity of a scheme over a field.

\sgaStableTarget{sga6.e03.prop.3-5}\textbf{Proposition 3.5.} Let \(S\) be a scheme, and for \(i=1,2\) let
\[
  f_i:X_i\longrightarrow Y_i
\]
be an \(S\)-morphism, with \(E_i\in \Ob D^{-}(X_i)\), and let \([a_i,b_i]\) be an interval of \(\mathbb Z\). Suppose that \(X_1\) and \(X_2\), and also \(Y_1\) and \(Y_2\), are Tor-independent over \(S\). Put
\[
  a=a_1+a_2,\qquad b=b_1+b_2,\qquad
  f=f_1\times_S f_2:X\to Y.
\]
Then:
\begin{enumerate}[label=\alph*)]
\item If \(\toramp_{f_1}(E_1)\subset[a_1,b_1]\), and if \(E_2\) is acyclic outside \([a_2,b_2]\), then
\[
  E_1\otimes_S^{\mathbf L}E_2
\]
is acyclic outside \([a,b]\).

\item If \(\toramp_{f_i}(E_i)\subset[a_i,b_i]\) for \(i=1,2\), then
\[
  \toramp_f(E_1\otimes_S^{\mathbf L}E_2)\subset[a,b].
\]
\end{enumerate}

\emph{Proof.} By \hyperlink{sga6.e03.prop.3-3}{Proposition 3.3}, it is enough to verify the assertion on the fibers. Let \(s\) be a point of \(S\), \(x\) a point of \(X\) above \(s\), \(y=f(x)\), and let \(x_i\) (respectively \(y_i\)) be the canonical projection of \(x\) (respectively \(y\)) onto \(X_i\) (respectively \(Y_i\)). Put
\[
  (E_i)_{x_i}=M_i,\quad
  \mathcal O_{X_i,x_i}=A_i,\quad
  \mathcal O_{Y_i,y_i}=B_i,\quad
  \mathcal O_{S,s}=C,
\]
and
\[
  \mathcal O_{X,x}=A=A_1\otimes_C A_2,\qquad
  \mathcal O_{Y,y}=B=B_1\otimes_C B_2.
\]
Then
\[
  (E_1\otimes_S^{\mathbf L}E_2)_x
  \simeq
  M_1\otimes_{A_1}^{\mathbf L} A
  \otimes_{A_2}^{\mathbf L} M_2.
\]
The Tor-independence hypotheses imply on the other hand that
\[
  A\simeq A_1\otimes_{B_1}^{\mathbf L}B
       \otimes_{B_2}^{\mathbf L}A_2.
\]
Consequently, the second member of the displayed isomorphism is canonically isomorphic, in the derived category \(D(B)\) of \(B\)-modules, to
\[
  M_1\otimes_{B_1}^{\mathbf L}B
       \otimes_{B_2}^{\mathbf L}M_2.
\]
The relations \(\toramp_{f_i}(E_i)\subset[a_i,b_i]\), \(i=1,2\), imply
\[
  \toramp_{B_i}(M_i)\subset[a_i,b_i],
\]
hence, by \hyperlink{sga6.e01.para.5-5}{Exposé I, 5.5},
\[
  \toramp_B\bigl(M_1\otimes_{B_1}^{\mathbf L}B
  \otimes_{B_2}^{\mathbf L}M_2\bigr)\subset[a,b],
\]
which proves b). If, on the other hand, \(\toramp_{f_1}(E_1)\subset[a_1,b_1]\), \(M_1\) is isomorphic to a complex of flat \(B_1\)-modules, zero outside \([a_1,b_1]\), and if \(E_2\), hence \(M_2\), is acyclic outside \([a_2,b_2]\), then
\[
  M_1\otimes_{B_1}^{\mathbf L}B\otimes_{B_2}^{\mathbf L}M_2
\]
is acyclic outside \([a,b]\). This proves a) and completes the proof.

\sgaStableTarget{sga6.e03.cor.3-5-1}\textbf{Corollary 3.5.1.} Under the hypotheses of \hyperlink{sga6.e03.prop.3-5}{Proposition 3.5}, if \(f_1\) and \(f_2\) are of finite Tor-dimension, then so is \(f\).

\sgaStableTarget{sga6.e03.cor.3-5-2}\textbf{Corollary 3.5.2.} Let
\[
\begin{tikzcd}
X' \arrow[r,"g"] \arrow[d,"f'"'] & X \arrow[d,"f"]\\
Y' \arrow[r,"h"'] & Y
\end{tikzcd}
\]
be a Cartesian square of schemes, where \(X\) and \(Y'\) are assumed Tor-independent over \(Y\), and let \(E\in \Ob D(X)\). Then one has
\[
  \toramp_f(E)\subset[a,b]\Longrightarrow
  \toramp_{f'}(Lg^*(E))\subset[a,b],
\]
and consequently the functor \(Lg^*:D^{-}(X)\to D^{-}(X')\) induces a functor
\[
  Lg^*:D(f)_{\mathrm{torf}}\longrightarrow D(f')_{\mathrm{torf}} .
\]



\providecommand{\Ob}{\operatorname{Ob}}
\providecommand{\D}{\mathrm D}
\providecommand{\Ho}{\mathrm H}
\providecommand{\RHom}{\mathrm R\!\operatorname{Hom}}
\providecommand{\parfamp}{\operatorname{parf.amp}}
\providecommand{\tamp}{\operatorname{tor.amp}}
\providecommand{\Qcoh}{\operatorname{Qcoh}}
\providecommand{\Mod}{\operatorname{Mod}}
\providecommand{\Coh}{\operatorname{Coh}}
\providecommand{\Lib}{\operatorname{Lib}}
\providecommand{\Loclib}{\operatorname{Loclib}}
\providecommand{\qcoh}{\mathrm{qcoh}}
\providecommand{\coh}{\mathrm{coh}}
\providecommand{\parf}{\mathrm{parf}}
\providecommand{\torf}{\mathrm{torf}}
\providecommand{\id}{\operatorname{id}}
\providecommand{\Coker}{\operatorname{Coker}}
\providecommand{\rank}{\operatorname{rank}}
\providecommand{\rg}{\operatorname{rg}}
\providecommand{\Hom}{\operatorname{Hom}}
\providecommand{\Ext}{\operatorname{Ext}}
\providecommand{\Spec}{\operatorname{Spec}}
\providecommand{\Cone}{\operatorname{Cone}}
\providecommand{\pt}{\operatorname{Pt}}
\providecommand{\ind}{\operatorname{ind}}
\hbadness=10000
\hfuzz=20pt
\emergencystretch=3em

\sgaStableTarget{sga6.e03.prop.3-6}\textbf{Proposition 3.6.} Let
\[
\begin{tikzcd}
X \arrow[r,"i"] \arrow[d,"f"'] & X' \arrow[dl,"f'"]\\
Y &
\end{tikzcd}
\]
be a commutative triangle of schemes, where \(i\) is a closed immersion and \(f'\) is smooth, and let \(E\in\Ob D(X)\). For \(E\) to be locally of finite flat amplitude relative to \(f\), it is necessary and sufficient that \(i_*(E)\) be locally of finite flat amplitude in the absolute sense. More precisely, if \(f'\) is of relative dimension \(\leq d\), then
\begin{align*}
\text{(i)}\quad &\tamp(i_*E)\subset[a,b] \Longrightarrow \tamp_f(E)\subset[a,b],\\
\text{(ii)}\quad &\tamp_f(E)\subset[a,b] \Longrightarrow \tamp(i_*E)\subset[a-d,b].
\end{align*}

\emph{Proof.} If \([m,n]\) is an interval in \(\mathbb Z\), one has evidently
\[
 \tamp(E)\subset[m,n] \Longleftrightarrow \tamp(i_*E)\subset[m,n]
\]
in the reduction in which one tests by modules after the closed immersion. Thus it remains to treat the case \(X=X'\), \(i=\id_X\), \(f=f'\). Assertion (i) is immediate from (\hyperlink{sga6.e03.cor.3-4-1}{3.4.1} a). For (ii), consider the diagonal section
\[
 s:X\longrightarrow X\times_YX
\]
of the two projections. If \(f\) is smooth of relative dimension \(d\), then, by \hyperlink{sga6.e07.prop.1-10}{VII 1.10}, \(s\) is a regular immersion of codimension \(\leq d\), hence of Tor-dimension \(\leq d\) (as is seen from the Koszul complex). By \hyperlink{sga6.e03.cor.3-5-2}{(3.5.2)},
\[
 \tamp_f(E)\subset[a,b]
 \Longleftrightarrow
 \tamp_{\operatorname{pr}_2}(\operatorname{pr}_1^*E)\subset[a,b].
\]
By (\hyperlink{sga6.e03.cor.3-4-1}{3.4.1} b) it follows that
\[
 \tamp(s^*\operatorname{pr}_1^*E)
 \subset[a-d,b],
\]
and the implication follows because \(s^*\operatorname{pr}_1^*\simeq\id\).

\sgaStableTarget{sga6.e03.rem.3-6-1}\textbf{Remark 3.6.1.} The statement \hyperlink{sga6.e03.prop.3-6}{(3.6)} would be false if one supposed only that \(f'\) were flat; see \hyperlink{sga6.e03.rem.3-4-2}{(3.4.2)}.

\sgaStableTarget{sga6.e03.prop.3-7}\textbf{Proposition 3.7 (projection formula).} Let \(S\) be a quasi-compact scheme, let \(f:X\to Y\) be a quasi-compact and quasi-separated \(S\)-morphism of quasi-compact schemes over \(S\), let \(E\in\Ob D^{-}(X)_{\qcoh}\), and let \(M\in\Ob D^b(S)_{\qcoh}\). If \(M\) has finite flat amplitude, or if \(E\) has finite flat amplitude relative to \(S\), there is a functorial canonical isomorphism in \(D(Y)_{\qcoh}\)
\[
 Rf_*(E)\otimes_S^{\mathbf L}M
 \xrightarrow{\sim}
 Rf_*(E\otimes_S^{\mathbf L}M).
\]

\emph{Proof.} First, by EGA III 1.4.12 and IV 1.7.21, there is an integer \(N\) such that \(R^if_*(F)=0\) for \(i>N\) and for every quasi-coherent \(\mathcal O_X\)-module \(F\). Hence \(Rf_*\) sends \(D^{-}(X)_{\qcoh}\) into \(D^{-}(Y)_{\qcoh}\). The adjunction arrow
\[
 Lf^*Rf_*(E)\longrightarrow E
\]
yields, after tensoring with \(M\), an arrow
\[
 Lf^*(Rf_*(E)\otimes_S^{\mathbf L}M)
 \longrightarrow E\otimes_S^{\mathbf L}M,
\]
and therefore, by adjunction, the canonical arrow
\sgaStableTarget{sga6.e03.eq.unnamed}\[
 Rf_*(E)\otimes_S^{\mathbf L}M
 \longrightarrow Rf_*(E\otimes_S^{\mathbf L}M).\tag{*}
\]
Both sides have quasi-coherent cohomology. The assertion is local on \(Y\), so assume \(S\) and \(Y\) affine. If \(M\) has finite flat amplitude, devissage and the way-out property of the two functors in \((*)\) reduce successively to \(M\) an ordinary quasi-coherent module, then to \(M\) free, then to \(M\) free of finite type, and finally to \(M=\mathcal O_S\), where \((*)\) is tautological. If instead \(E\) has finite flat amplitude relative to \(S\), devissage reduces to the case \(M=P\), with \(P\) a flat \(A\)-module and \(S=\Spec A\). Put \(A'=D_A(P)\) (EGA \(0_{\mathrm{IV}}\), 18.2.3) and \(S'=\Spec A'\). Then the assertion is the flat base-change assertion that
\[
 q^*Rf_*(E)\longrightarrow Rf'_*p^*(E)
\]
is an isomorphism for the base change to \(S'\), which is EGA III 1.4.15 together with IV 1.7.21.

\sgaStableTarget{sga6.e03.cor.3-7-1}\textbf{Corollary 3.7.1.} Let \(S\) be a scheme, let \(f:X\to Y\) be a quasi-compact and quasi-separated \(S\)-morphism of quasi-compact schemes over \(S\), and let \(E\in\Ob D^{-}(X)_{\qcoh}\). If \(E\) is locally of finite flat amplitude relative to \(S\), then the same is true of \(Rf_*(E)\). More precisely, if \(S\) is quasi-compact and if \(N\) is an integer such that \(R^if_*(F)=0\) for \(i>N\) and every quasi-coherent \(\mathcal O_X\)-module \(F\), then
\[
 \tamp_S(E)\subset[a,b]
 \Longrightarrow
 \tamp_S(Rf_*E)\subset[a,b+N].
\]

\emph{Proof.} By (\hyperlink{sga6.e03.prop.3-3}{3.3} (1) (iv)), and using that \(Rf_*\) commutes with restriction to an open of \(S\), it suffices to prove that the hypothesis implies
\[
 H^i(Rf_*E\otimes_S^{\mathbf L}M)=0
\]
for \(i\notin[a,b+N]\) and every quasi-coherent \(\mathcal O_S\)-module \(M\). But the hypothesis says that \(H^i(E\otimes_S^{\mathbf L}M)=0\) for \(i\notin[a,b]\), and \hyperlink{sga6.e03.prop.3-7}{(3.7)} gives the result.

\sgaStableTarget{sga6.e03.cor.3-7-2}\textbf{Corollary 3.7.2.} Let \(f:X\to Y\) be quasi-compact, quasi-separated, and locally of finite Tor-dimension, and let \(E\in\Ob D^{-}(X)_{\qcoh}\). If \(E\) is locally of finite flat amplitude in the absolute sense, then so is \(Rf_*(E)\).

\emph{Proof.} Combine (\hyperlink{sga6.e03.cor.3-4-1}{3.4.1} a) and \hyperlink{sga6.e03.cor.3-7-1}{(3.7.1)}.

\sgaStableTarget{sga6.e03.sec.4}\section*{4. Relative perfection}

\sgaStableTarget{sga6.e03.def.4-1}\textbf{Definition 4.1.} Let \(f:X\to Y\) be a morphism locally of finite type of schemes, let \(E\in\Ob D(X)\), and let \([a,b]\) be an interval of \(\mathbb Z\). We say that \(E\) is of perfect amplitude relative to \(f\) (or to \(Y\), if there is no ambiguity), contained in \([a,b]\), and write
\[
 \parfamp_f(E)\subset[a,b]
\]
if \(E\), relative to \(f\), is pseudo-coherent \hyperlink{sga6.e03.def.1-2}{(1.2)} and of flat amplitude contained in \([a,b]\) \hyperlink{sga6.e03.def.3-1}{(3.1)}. We say that \(f\) is of perfect amplitude contained in \([-n,0]\), or of perfect dimension \(\leq n\), if \(\mathcal O_X\) has perfect amplitude contained in \([-n,0]\). We say that \(E\) is of finite perfect amplitude relative to \(f\) if such an interval \([a,b]\) exists. We say that \(E\) is perfect relative to \(f\) if it is locally of finite perfect amplitude relative to \(f\). We say that \(f\) is perfect if \(\mathcal O_X\) is locally of finite perfect amplitude relative to \(f\).

\sgaStableTarget{sga6.e03.ex.4-1-1}\textbf{Example 4.1.1.} A regular immersion of codimension \(d\) (\hyperlink{sga6.e07.def.1-4}{VII 1.4}) is of perfect dimension \(\leq d\), as follows from the Koszul complex. A morphism locally a complete intersection, i.e. locally the composite of a regular immersion and a smooth morphism, is a perfect morphism.

\sgaStableTarget{sga6.e03.notation.4-2}\textbf{Notation 4.2.} Let \(f:X\to Y\) be locally of finite type. We write
\[
 D(f)_{\parf},\qquad D(f)_{\mathrm{parf}},\qquad D(f)_{\mathrm{parf}}^{-}
\]
when no confusion is possible, for the full subcategories of \(D(X)\) formed by complexes perfect, respectively of finite perfect amplitude, relative to \(f\). By (\hyperlink{sga6.e01.para.5-3}{I 5.3}) and \hyperlink{sga6.e03.prop.1-4}{(1.4)}, these are triangulated subcategories of \(D(X)\). If \(X\) is quasi-compact, the two finite-amplitude/perfect notations coincide in the evident bounded sense.

\sgaStableTarget{sga6.e03.prop.4-3}\textbf{Proposition 4.3} (cf. \hyperlink{sga6.e01.thm.5-8}{I 5.8}). Let \(f:X\to Y\) be locally of finite type, let \(E\in\Ob D(X)\), let \([a,b]\) be an interval of \(\mathbb Z\), and let \(x\in X\), \(y=f(x)\). The following conditions are equivalent:
\begin{enumerate}[label=\roman*)]
\item \(\parfamp_f(E)\subset[a,b]\) in a neighbourhood of \(x\);
\item in a neighbourhood of \(x\), \(E\) is \((a-1)\)-pseudo-coherent relative to \(f\), is acyclic outside \([a,b]\), and
\[
 \tamp_{\mathcal O_{Y,y}}(E_x)\subset[a,b].
\]
\end{enumerate}

\emph{Proof.} The implication i) \(\Rightarrow\) ii) is immediate. For the converse, the question is local around \(x\), and by the definitions one reduces to the case where \(f\) is smooth. Then \(E\) is \((a-1)\)-pseudo-coherent near \(x\); after localizing, one may assume that \(E\) is strictly \((a-1)\)-pseudo-coherent and zero in degrees \(>b\). After truncating, one may represent it by a complex with \(E^i\) free of finite type for \(i>a\), \(E^a\) of finite presentation, and zero outside \([a,b]\). By (\hyperlink{sga6.e01.lem.5-1-1}{I 5.1.1}), the term in degree \(a\) modulo boundaries is flat relative to \(f\) near \(x\). Since \(f\) is smooth, one embeds locally into an affine space over \(Y\) and descends finite presentation and flatness from a noetherian approximation; pseudo-coherence follows from (\hyperlink{sga6.e03.cor.1-10}{1.10} a), and the desired perfect-amplitude assertion follows.

\sgaStableTarget{sga6.e03.cor.4-3-1}\textbf{Corollary 4.3.1.} Let \(f:X\to Y\) be locally of finite type and let \(E\) be an \(\mathcal O_X\)-module. One has \(\parfamp_f(E)\subset[0,0]\) if and only if \(E\) is flat and of finite presentation relative to \(f\) (\hyperlink{sga6.e03.rem.1-2-1}{1.2.1} c). In particular, \(f\) is of perfect dimension zero if and only if \(f\) is flat and locally of finite presentation.

\sgaStableTarget{sga6.e03.cor.4-3-2}\textbf{Corollary 4.3.2.} Let \(f:X\to Y\) be locally of finite type, let \(E\) be a complex of \(\mathcal O_X\)-modules flat and of finite presentation relative to \(f\), and let \([a,b]\) be an interval. If \(E^i=0\) for \(i\notin[a,b]\), then \(\parfamp_f(E)\subset[a,b]\). If \(E\) is bounded above, then \(E\) is pseudo-coherent relative to \(f\).

\emph{Proof.} This follows from \hyperlink{sga6.e03.cor.4-3-1}{(4.3.1)} and from \hyperlink{sga6.e01.prop.4-15}{I 4.15} and \hyperlink{sga6.e01.ex.4-8}{I 4.8}.

\sgaStableTarget{sga6.e03.prop.4-4}\textbf{Proposition 4.4.} Let
\[
\begin{tikzcd}
X \arrow[r,"i"] \arrow[d,"f"'] & X' \arrow[dl,"g"]\\
Y &
\end{tikzcd}
\]
be a commutative triangle of schemes, where \(i\) is a closed immersion and \(g\) is smooth, and let \(E\in\Ob D(X)\). Then \(E\) is perfect relative to \(f\) if and only if \(i_*(E)\) is perfect in the absolute sense. More precisely, if \(g\) has relative dimension \(\leq d\), then
\begin{align*}
\parfamp(i_*E)
  \subset[a,b]&\Longrightarrow \parfamp_f(E)
  \subset[a,b],\\
\parfamp_f(E)
  \subset[a,b]&\Longrightarrow \parfamp(i_*E)
  \subset[a-d,b].
\end{align*}

\emph{Proof.} This is an immediate consequence of \hyperlink{sga6.e03.prop.3-6}{(3.6)}.

\sgaStableTarget{sga6.e03.prop.4-5}\textbf{Proposition 4.5.} Let
\[
\begin{tikzcd}
X \arrow[r,"f"] \arrow[dr,"h"'] & Y \arrow[d,"g"]\\
& S
\end{tikzcd}
\]
be a commutative triangle of morphisms locally of finite type, let \(E\in\Ob D(X)\), \(F\in\Ob D(Y)\), and let \([a,b]\), \([c,d]\) be intervals of \(\mathbb Z\). Then
\[
 \parfamp_f(E)\subset[a,b],\quad
 \parfamp_g(F)\subset[c,d]
 \Longrightarrow
 \parfamp_h(E\otimes_Y^{\mathbf L}F)
     \subset[a+c,b+d].
\]
In particular, if \(E\) is perfect relative to \(f\) and \(F\) is perfect relative to \(g\), then \(E\otimes_Y^{\mathbf L}F\) is perfect relative to \(h\).

\emph{Proof.} Combine \hyperlink{sga6.e03.thm.1-1}{(1.1)} and \hyperlink{sga6.e03.cor.3-4}{(3.4)}.

\sgaStableTarget{sga6.e03.cor.4-5-1}\textbf{Corollary 4.5.1.} Under the hypotheses of \hyperlink{sga6.e03.prop.4-5}{(4.5)}:
\begin{enumerate}[label=\alph*)]
\item if \(g\) is of perfect dimension \(\leq n\), then
\[
 \parfamp_f(E)
  \subset[a,b]
 \Longrightarrow
 \parfamp_h(E)
  \subset[a-n,b];
\]
therefore, if \(g\) is perfect, \(E\) perfect relative to \(f\) implies \(E\) perfect relative to \(h\);
\item if \(f\) is of perfect dimension \(\leq m\), then
\[
 \parfamp_g(F)
  \subset[c,d]
 \Longrightarrow
 \parfamp_h(Lf^*F)
  \subset[c-m,d];
\]
therefore, if \(f\) is perfect, \(F\) perfect relative to \(g\) implies \(Lf^*F\) perfect relative to \(h\).
\end{enumerate}

\sgaStableTarget{sga6.e03.rem.4-5-2}\textbf{Remark 4.5.2.} With the preceding notation, if \(g\) is perfect, \(E\) perfect relative to \(h\) does not necessarily imply that \(E\) is perfect relative to \(f\); the usual example of a singular point of a scheme over a field already shows this. The following criterion is available.

\sgaStableTarget{sga6.e03.prop.4-6}\textbf{Proposition 4.6} (criterion of perfection by fibers). Let
\[
\begin{tikzcd}
X \arrow[r,"f"] \arrow[dr,"h"'] & Y \arrow[d,"g"]\\
& S
\end{tikzcd}
\]
be a commutative triangle of morphisms of schemes. Suppose that \(f\) and \(h\) are locally of finite type, and that \(g\) is locally of finite presentation and flat. Let \(x\in X\), \(y=f(x)\), \(s=h(x)=g(y)\), and let \(E\in\Ob D(X)\). Let \([a,b]\) be an interval of \(\mathbb Z\), and let \(f_s:X_s\to Y_s\) be the fiber of \(f\) over \(s\), i.e. the morphism obtained from \(f\) by the base change \(\Spec k(s)\to S\). Then the following conditions are equivalent:
\begin{enumerate}[label=\roman*)]
\item \(\parfamp_h(E)\subset[a,b]\) near \(x\), and
\[
 \tamp_{\mathcal O_{Y_s,y}}(E_x\otimes_{\mathcal O_{Y,y}}^{\mathbf L}\mathcal O_{Y_s,y})
 \subset[a,b];
\]
\item \(\parfamp_f(E)\subset[a,b]\) near \(x\).
\end{enumerate}

\emph{Proof.} The implication ii) \(\Rightarrow\) i) follows because \(g\) has perfect dimension zero by \hyperlink{sga6.e03.cor.4-3-1}{(4.3.1)}, hence (\hyperlink{sga6.e03.cor.4-5-1}{4.5.1} a) gives \(\parfamp_h(E)\subset[a,b]\) near \(x\); the asserted Tor-amplitude condition follows from \hyperlink{sga6.e03.prop.4-3}{(4.3)}. Conversely, the question is local around \(x\). We may factor \(f\) as a closed immersion followed by a smooth morphism. Replacing \(E\) by its direct image under the immersion and replacing the smooth factor accordingly does not change the conditions; hence one may suppose \(f\) smooth. Since \(g\) is perfect of dimension zero and \(h\) is perfect by hypothesis, \(E\) is pseudo-coherent. After localizing, use (\hyperlink{sga6.e02.rem.2-2-10}{II 2.2.10} b) and (\hyperlink{sga6.e01.prop.2-10}{I 2.10} a) to suppose that \(E\) is strictly pseudo-coherent, zero in degrees \(>b\), and quasi-isomorphic to the truncation
\[
 t:
 E\longrightarrow
 (0\to E^a/B^a\to E^{a+1}\to\cdots\to E^b\to0).
\]
It remains to prove that \(E^a/B^a\) is \(f\)-flat near \(x\). Let \(i:X_s\to X\) be the canonical map. One has
\[
 E_x\otimes_{\\mathcal O_{Y,y}}^{\mathbf L}\mathcal O_{Y_s,y}
 \simeq i^*(E)_x.
\]
Using \hyperlink{sga6.e03.cor.3-5-2}{(3.5.2)}, the Tor-amplitude hypothesis on the fiber shows that the truncated complex \(i^*(t)\) is a quasi-isomorphism, so \(i^*(E^a/B^a)_x\) is flat over \(\mathcal O_{Y_s,y}\). On the other hand, the Tor-amplitude of \(E\) relative to \(h\) implies that \(E^a/B^a\) is \(h\)-flat. The usual criterion of flatness by fibers (EGA IV 11.3.10) then gives that \(E^a/B^a\) is \(f\)-flat near \(x\).

\sgaStableTarget{sga6.e03.cor.4-6-1}\textbf{Corollary 4.6.1.} Let \(f:X\to S\) be locally of finite presentation and flat, let \(x\in X\), \(s=f(x)\), let \(E\in\Ob D(X)\), and let \([a,b]\) be an interval. If \(i_s:X_s=X\times_S\Spec k(s)\to X\) is the canonical projection and \(E_s=Li_s^*E\), then
\[
 \parfamp_S(E)
  \subset[a,b]\text{ near }x,
 \quad
 \tamp(E_s)
  \subset[a,b]
\]
if and only if
\[
 \parfamp(E)
  \subset[a,b]\text{ near }x.
\]
Consequently \(E\) is perfect if and only if \(E\) is perfect relative to \(f\) and \(E_s\) is perfect for every \(s\in S\).

\emph{Proof.} Apply \hyperlink{sga6.e03.prop.4-6}{(4.6)} with \(X=Y\) and \(f=\id_X\).

\sgaStableTarget{sga6.e03.rem.4-6-2}\textbf{Remark 4.6.2.} From \hyperlink{sga6.e03.cor.4-6-1}{(4.6.1)} and \hyperlink{sga6.e01.cor.5-10}{I 5.10} it follows that, if \(f\) is locally of finite presentation, flat, and has regular schemes as fibers, then \(E\) perfect relative to \(f\) is equivalent to \(E\) perfect. This equivalence was already obtained in \hyperlink{sga6.e03.prop.4-4}{(4.4)} when \(f\) is smooth.

\sgaStableTarget{sga6.e03.prop.4-7}\textbf{Proposition 4.7.} Let \(S\) be a scheme, and for \(i=1,2\) let
\[
 f_i:X_i\longrightarrow Y_i
\]
be an \(S\)-morphism locally of finite type, with \(E_i\in\Ob D(X_i)\) and intervals \([a_i,b_i]\). Suppose that \(X_1\) and \(X_2\), and also \(Y_1\) and \(Y_2\), are Tor-independent over \(S\). Put \(a=a_1+a_2\), \(b=b_1+b_2\), and \(f=f_1\times_S f_2\). Then:
\begin{enumerate}[label=\alph*)]
\item if \(\parfamp_{f_1}(E_1)\subset[a_1,b_1]\), and if \(E_2\) is pseudo-coherent relative to \(f_2\) and acyclic outside \([a_2,b_2]\), then \(E_1\otimes_S^{\mathbf L}E_2\) is pseudo-coherent relative to \(f\) and acyclic outside \([a,b]\);
\item if \(\parfamp_{f_i}(E_i)\subset[a_i,b_i]\) for \(i=1,2\), then
\[
 \parfamp_f(E_1\otimes_S^{\mathbf L}E_2)
 \subset[a,b].
\]
\end{enumerate}

\emph{Proof.} Combine \hyperlink{sga6.e03.cor.1-9}{(1.9)} and \hyperlink{sga6.e03.prop.3-5}{(3.5)}.

\sgaStableTarget{sga6.e03.cor.4-7-1}\textbf{Corollary 4.7.1.} Under the hypotheses of \hyperlink{sga6.e03.prop.4-7}{(4.7)}, if \(f_i\) has perfect dimension \(\leq n_i\) for \(i=1,2\), then \(f\) has perfect dimension \(\leq n_1+n_2\). Consequently, if \(f_1\) and \(f_2\) are perfect, then so is \(f\).

\sgaStableTarget{sga6.e03.cor.4-7-2}\textbf{Corollary 4.7.2.} In a Cartesian square
\[
\begin{tikzcd}
X' \arrow[r,"g"] \arrow[d,"f'"'] & X \arrow[d,"f"]\\
Y' \arrow[r,"h"'] & Y,
\end{tikzcd}
\]
where \(f\) is locally of finite type and \(X\) and \(Y'\) are Tor-independent over \(Y\), one has
\[
 \parfamp_f(E)
 \subset[a,b]
 \Longrightarrow
 \parfamp_{f'}(Lg^*E)
 \subset[a,b].
\]
Consequently \(Lg^*:D^{-}(X)\to D^{-}(X')\) sends \(D(f)_{\parf}\), respectively \(D(f)_{\mathrm{parf}}\), into the corresponding categories for \(f'\).

\sgaStableTarget{sga6.e03.prop.4-8}\textbf{Proposition 4.8.} Let \(S\) be a scheme, let \(f:X\to Y\) be a proper \(S\)-morphism of \(S\)-schemes locally of finite type, and let \(E\in\Ob D^{-}(X)\). Under the proviso that \hyperlink{sga6.e03.conj.2-1}{(2.1)} is valid for \(f\) (hence at least if \(S\) is locally noetherian or if \(f\) is projective):
\begin{enumerate}[label=\alph*)]
\item if \(E\) is perfect relative to \(S\), then so is \(Rf_*(E)\);
\item if \(S\) is quasi-compact, \(X\) and \(Y\) are of finite type over \(S\), \([a,b]\) is an interval, and \(N\) is such that \(R^if_*(M)=0\) for \(i>N\) and every quasi-coherent \(\mathcal O_X\)-module \(M\), then
\[
 \parfamp_S(E)\subset[a,b]
 \Longrightarrow
 \parfamp_S(Rf_*E)
 \subset[a,b+N].
\]
\end{enumerate}
Consequently \(Rf_*\) sends perfect complexes relative to \(S\) to perfect complexes relative to \(S\).

\emph{Proof.} Assertion a) follows from b), and b) is obtained by combining \hyperlink{sga6.e03.conj.2-1}{(2.1)} and \hyperlink{sga6.e03.prop.3-7}{(3.7)}.

\sgaStableTarget{sga6.e03.cor.4-8-1}\textbf{Corollary 4.8.1.} Let \(f:X\to Y\) be proper and perfect, and let \(E\in\Ob D^{-}(X)\). Under the proviso that \hyperlink{sga6.e03.conj.2-1}{(2.1)} holds for \(f\), if \(E\) is perfect then \(Rf_*(E)\) is perfect. If moreover \(Y\) is quasi-compact, \(Rf_*\) sends \(D(X)_{\parf}\) to \(D(Y)_{\parf}\).

\emph{Proof.} If \(E\) is perfect, then \(E\) is perfect relative to \(f\) by (\hyperlink{sga6.e03.cor.4-5-1}{4.5.1} a).

\sgaStableTarget{sga6.e03.prop.4-9}\textbf{Proposition 4.9.} Let \(S\) be noetherian, let \(X\) and \(Y\) be \(S\)-schemes of finite type, let \(f:X\to Y\) be a compactifiable \(S\)-morphism, i.e. the composite of an open immersion and a proper morphism, and let \(E\in\Ob D(X)\), \(F\in\Ob D^{-}(Y)\). If \(E\) is perfect relative to \(f\) and \(F\) is perfect relative to \(S\), then
\[
 \RHom(E,f^!F)
\]
is perfect relative to \(S\).

\emph{Proof.} The question is local on \(X\). One may suppose that \(f\) factors as
\[
 X\xrightarrow{i}X'\xrightarrow{f'}Y,
\]
where \(i\) is a closed immersion and \(f'\) is smooth. If \(E\) is perfect relative to \(f\), then it is perfect relative to \(i\) by \hyperlink{sga6.e03.prop.4-4}{(4.4)}; if \(F\) is perfect relative to \(S\), then \(f'^!F\) is perfect relative to \(S\) by (\hyperlink{sga6.e03.cor.4-5-1}{4.5.1} b), since \(f'^!F\) is locally isomorphic to \(f'^*F\) up to translation. As \(f^!=i^!f'^!\), it remains to prove the assertion when \(f\) is a closed immersion. Then \(f_*E\) is perfect and, by global duality ([H] III 6.7),
\[
 f_*\RHom(E,f^!F)\simeq\RHom(f_*E,F).
\]
Since \(f_*E\) is perfect, (\hyperlink{sga6.e01.cor.7-7}{I 7.7}) gives
\[
 \RHom(f_*E,F)\simeq(f_*E)^{\vee}\otimes^{\mathbf L}F,
\]
and the result follows from \hyperlink{sga6.e03.prop.4-5}{(4.5)} with the identity morphism on \(Y\).

\sgaStableTarget{sga6.e03.cor.4-9-1}\textbf{Corollary 4.9.1.} Under the hypotheses of \hyperlink{sga6.e03.prop.4-8}{(4.8)}, if \(f\) is perfect, then \(f^!\) sends \(D(Y)_{\parf}\) into \(D(X)_{\parf}\).

\emph{Proof.} Take \(E=\mathcal O_X\) in \hyperlink{sga6.e03.prop.4-9}{(4.9)}.

\sgaStableTarget{sga6.e03.cor.4-9-2}\textbf{Corollary 4.9.2.} Let \(S\) be noetherian and let \(f:X\to S\) be compactifiable. Put \(K_X=f^!(\mathcal O_S)\), and \(D_X(\cdot)=\RHom(\cdot,K_X)\). For \(E\in\Ob D(f)_{\parf}\), one has \(D_X(E)\in\Ob D(f)_{\parf}\), and the canonical arrow
\[
 E\longrightarrow D_XD_X(E)
\]
is an isomorphism.

\emph{Proof.} For the first assertion, apply \hyperlink{sga6.e03.prop.4-9}{(4.9)} with \(Y=S\) and \(F=\mathcal O_S\). For the second, one reduces easily to the case where \(f\) is a closed immersion. It is enough then to apply \(f_*\) to the canonical arrow. By global duality one has \(f_*D_X=D_Sf_*\), and one is reduced to ordinary biduality for perfect complexes (\hyperlink{sga6.e01.prop.7-2}{I 7.2}).

\sgaStableTarget{sga6.e03.sec.5}\section*{5. Applications: exchange and semi-continuity theorems}

\sgaStableTarget{sga6.e03.lem.5-1}\textbf{Lemma 5.1.} Let \((X,\mathcal O_X)\) be a commutatively ringed topos, let \(\Mod(X)\) be the category of \(\mathcal O_X\)-modules, and write \(D(X)=D(\Mod(X))\). Let \(E\in\Ob D(X)\) and let \(p\in\mathbb Z\). Consider the following conditions:
\begin{enumerate}[label=\alph*)]
\item the functor \(M\mapsto H^p(E\otimes^{\mathbf L}M)\) from \(\Mod(X)\) to \(\Mod(X)\) is right exact;
\item[\(\mathrm a'\)] the canonical map
\[
 H^p(E)\otimes M\longrightarrow H^p(E\otimes^{\mathbf L}M)
\]
is an isomorphism for all \(M\in\Ob\Mod(X)\);
\item[\(\mathrm b\)] the functor \(M\mapsto H^{p+1}(E\otimes^{\mathbf L}M)\) is left exact;
\item[\(\mathrm b'\)] the canonical map defined below
\[
 H^{p+1}(E\otimes^{\mathbf L}M)
 \longrightarrow
 \Hom(H^{-p-1}(E^{\vee}),M)
\]
is an isomorphism for all \(M\in\Ob\Mod(X)\).
\end{enumerate}
Then
\[
 \mathrm a\Longleftrightarrow \mathrm a'
 \Longleftrightarrow \mathrm b
 \Longleftarrow \mathrm b'.
\]
If \(E\) is pseudo-coherent, then a), a'), b), and b') are equivalent.

The map in b') is obtained as follows. The canonical evaluation arrow \(E^{\vee}\otimes^{\mathbf L}E\to\mathcal O_X\), tensorized with \(M\), gives an arrow
\[
 E\otimes^{\mathbf L}M\longrightarrow\RHom(E^{\vee},M).
\]
On \(H^{p+1}\), and then by the canonical map
\[
 \Ext^{p+1}(E^{\vee},M)	o\Hom(H^{-p-1}(E^{\vee}),M),
\]
one obtains the arrow b').

\emph{Proof.} The implications \(\mathrm a'\Rightarrow\mathrm a\Leftrightarrow\mathrm b\Leftarrow\mathrm b'\) are formal. To prove \(\mathrm a\Rightarrow\mathrm a'\), present \(M\) by a sequence \(M^0\to M^1\to M\to0\) in which \(M^0\) and \(M^1\) are direct sums of sheaves of the form \(\mathcal O_{U,X}\), with \(U\) an object of \(X\). Right exactness and compatibility with direct sums reduce the assertion to \(M=\mathcal O_{U,X}\). Then both sides are identified with extension by zero from \(U\), and exactness of that extension gives the result.

Assume now that \(E\) is pseudo-coherent and prove \(\mathrm b\Rightarrow\mathrm b'\). It is enough to show that the preceding comparison arrow
\[
 E\otimes^{\mathbf L}M\longrightarrow\RHom(E^{\vee},M)
\]
is an isomorphism when \(M\) is injective. Let \(T(F)=F\otimes^{\mathbf L}M\) and \(T'(F)=\RHom(F^{\vee},M)\). Both functors have the same boundedness behaviour under truncation; the arrow is local and is compatible with restriction to opens, and an injective remains injective after such a restriction. After localizing and truncating, one may suppose \(E\) strictly perfect. A devissage then reduces to \(E=\mathcal O_X\), where the assertion is immediate.

\sgaStableTarget{sga6.e03.rem.5-2}\textbf{Remarks 5.2.}
\begin{enumerate}[label=\roman*)]
\item Conditions a), a'), b), and b') are local on the topos: if \((X_\alpha)\) is an open covering of the final object and the condition holds for each \(X_\alpha\) and \(E|X_\alpha\), then it holds on \(X\).
\item If \(x\) is a point of \(X\), write \(a_x\), \(a'_x\), \(b_x\), etc. for the corresponding conditions for \(E_x\) and \(\Mod(\mathcal O_{X,x})\). If \(X\) has enough points, condition a), respectively a'), b), is equivalent to the corresponding stalkwise condition for every \(x\); if \(E\) is pseudo-coherent, this is also equivalent to b') stalkwise. Moreover, a') may be tested on quotients of modules of the form \(\mathcal O_{U,x}^{(I)}\to\mathcal O_{U,x}\).
\item The criterion EGA III 7.4.2 extends word for word: if \(E\) is represented by flat terms, condition a) says that \(\Coker(E^p\to E^{p+1})\) is flat, hence locally a direct factor of a free module of finite type when \(E^p\) and \(E^{p+1}\) are of finite presentation. In particular, for \(E\) perfect, condition a) for degree \(p\) is equivalent to the corresponding condition for \(E^{\vee}\) in degree \(-p-1\).
\item Let \(X\) be a quasi-separated scheme and let \(E\in\Ob D(X)\) have locally bounded quasi-coherent cohomology. Let \(a_{\qcoh}\), \(a'_{\qcoh}\), etc. be the analogous conditions with \(\Mod(X)\) replaced by \(\Qcoh(X)\). Then
\[
 a\Rightarrow a_{\qcoh},\qquad
 a'\Rightarrow a'_{\qcoh},\qquad
 b\Rightarrow b_{\qcoh},
\]
and these conditions are all equivalent when \(E\) is pseudo-coherent. This follows by reducing to affine opens, using \hyperlink{sga6.e02.prop.3-5}{II 3.5} and \hyperlink{sga6.e02.rem.2-2-10}{II 2.2.10} a), and the fact that quasi-coherent sheaves extend from affine opens on a quasi-separated scheme.
\end{enumerate}

\sgaStableTarget{sga6.e03.prop.5-3}\textbf{Proposition 5.3} (exchange property; EGA III 7.7.5 II). Let \(f:X\to Y\) be a morphism of quasi-compact and quasi-separated schemes, let \(E\in\Ob D^{-}(X)\) have quasi-coherent cohomology and finite flat amplitude relative to \(f\), and let \(p\in\mathbb Z\). Consider:
\begin{enumerate}[label=\alph*)]
\item \(M\mapsto R^pf_*(E\otimes_Y^{\mathbf L}M)\), from \(\Qcoh(Y)\) to \(\Qcoh(Y)\), is right exact;
\item[\(\mathrm a'\)] the canonical arrow
\[
 R^pf_*(E)\otimes M
 \longrightarrow
 R^pf_*(E\otimes_Y^{\mathbf L}M)
\]
is an isomorphism for every \(M\in\Ob\Qcoh(Y)\);
\item[\(\mathrm b\)] \(M\mapsto R^{p+1}f_*(E\otimes_Y^{\mathbf L}M)\) is left exact;
\item[\(\mathrm b'\)] the canonical arrow
\[
 R^{p+1}f_*(E\otimes_Y^{\mathbf L}M)
 \longrightarrow
 \Hom(H^{-p-1}((Rf_*E)^{\vee}),M)
\]
is an isomorphism for every \(M\in\Ob\Qcoh(Y)\);
\item[\(\mathrm c\)] for every base change \(Y'\to Y\), the canonical arrow
\[
 R^pf_*(E)\otimes_Y\mathcal O_{Y'}
 \longrightarrow
 H^p(Rf_*E\otimes_Y^{\mathbf L}\mathcal O_{Y'})
\]
is an isomorphism.
\end{enumerate}
Then
\[
 a\Longleftrightarrow a'\Longleftrightarrow c\Longleftrightarrow b\Longleftarrow b'.
\]
If, moreover, \(f\) is proper, \(E\) is pseudo-coherent relative to \(f\) (hence perfect relative to \(f\)), and \hyperlink{sga6.e03.conj.2-1}{(2.1)} is valid for \(f\), for instance if \(f\) is projective or \(Y\) is noetherian, then the preceding conditions are all equivalent.

\emph{Proof.} The implications among a), a'), b), and b') follow at once from \hyperlink{sga6.e03.prop.3-7}{(3.7)}, \hyperlink{sga6.e03.lem.5-1}{(5.1)}, and (\hyperlink{sga6.e03.rem.5-2}{5.2} iv); the implication \(b\Rightarrow b'\) in the proper pseudo-coherent case follows from \hyperlink{sga6.e03.conj.2-1}{(2.1)}. It remains to compare a') and c). This is supplied by the following lemma.

\sgaStableTarget{sga6.e03.lem.5-3-1}\textbf{Lemma 5.3.1.} Let \(X\) be a quasi-separated scheme, and let \(E\in\Ob D^{-}(X)\) have locally bounded quasi-coherent cohomology. Then condition (\hyperlink{sga6.e03.lem.5-1}{5.1} a) is equivalent to:
\begin{enumerate}[label=\alph*)]
\setcounter{enumi}{2}
\item for every \(X'\to X\), the canonical map
\[
 H^p(E)\otimes_X\mathcal O_{X'}
 \longrightarrow
 H^p(E\otimes_X^{\mathbf L}\mathcal O_{X'})
\]
is an isomorphism.
\end{enumerate}
If \(E\) is pseudo-coherent, condition (\hyperlink{sga6.e03.lem.5-1}{5.1} a) also implies that, for every \(X'\to X\), the map
\sgaStableTarget{sga6.e03.eq.5-3-1-1}\[
 H^{-p-1}(E^{\vee})\otimes_X\mathcal O_{X'}
 \longrightarrow
 H^{-p-1}((E\otimes_X^{\mathbf L}\mathcal O_{X'})^{\vee})
\tag{5.3.1.1}
\]
is an isomorphism.

\emph{Proof of the lemma.} For \(a\Rightarrow c\), localize on \(X'\), and assume \(X=\Spec A\), \(X'=\Spec A'\), and \(E=\widetilde F\). By the affine case the hypothesis gives
\[
 H^p(F)\otimes_AA'
 \longrightarrow
 H^p(F\otimes_A^{\mathbf L}A')
\]
an isomorphism, which is exactly c). Conversely, by (\hyperlink{sga6.e03.rem.5-2}{5.2} iv) it suffices to prove \(c\Rightarrow a'_{\qcoh}\). Again localizing, write \(X=\Spec A\), \(E=\widetilde F\), and let \(M=\widetilde N\) be quasi-coherent. Taking \(A'=D_A(N)\) and applying c) to \(\Spec A'\to X\) gives the required map
\[
 H^p(F)\otimes_AN
 \longrightarrow
 H^p(F\otimes_A^{\mathbf L}N)
\]
as an isomorphism.

For the last assertion, the question is local on \(X'\). By (\hyperlink{sga6.e02.rem.2-2-10}{II 2.2.10} b) and truncation, one may suppose \(E\) strictly perfect. Then \((E\otimes_X^{\mathbf L}\mathcal O_{X'})^{\vee}\) identifies canonically with \(E^{\vee}\otimes_X^{\mathbf L}\mathcal O_{X'}\), and \hyperlink{sga6.e03.eq.5-3-1-1}{(5.3.1.1)} is the canonical map
\sgaStableTarget{sga6.e03.eq.unnamed.occ2}\[
 H^{-p-1}(E^{\vee})\otimes_X\mathcal O_{X'}
 \longrightarrow
 H^{-p-1}(E^{\vee}\otimes_X^{\mathbf L}\mathcal O_{X'}).
\tag{*}
\]
If \(E\) satisfies (\hyperlink{sga6.e03.lem.5-1}{5.1} a) for \(p\), then by (\hyperlink{sga6.e03.rem.5-2}{5.2} iii) \(E^{\vee}\) satisfies the same condition for \(-p-1\), and \((*)\) is an isomorphism by the first part of the lemma. The lemma, and therefore \hyperlink{sga6.e03.prop.5-3}{(5.3)}, follow.

\sgaStableTarget{sga6.e03.rem.5-3-2}\textbf{Remark 5.3.2.} \hyperlink{sga6.e03.prop.5-3}{Proposition (5.3)} applies in particular when \(E\) is a bounded complex of quasi-coherent flat \(\mathcal O_X\)-modules of finite presentation relative to \(f\): by \hyperlink{sga6.e03.cor.4-3-2}{(4.3.2)}, such a complex has finite perfect amplitude relative to \(f\).

\sgaStableTarget{sga6.e03.def.5-4}\textbf{Definition 5.4.} Let \(X\) be a topos, let \(\pt(X)\) be its set of points (SGA 4 IV 6.1), let \(M\) be an ordered set, and let \(\varphi:\pt(X)\to M\) be a function. We say that \(\varphi\) is upper semi-continuous at \(x\) (respectively lower semi-continuous, respectively locally constant) if there exists an \(x\)-pointed object \(U\) of \(X\) such that for every point \(y\) of \(U\), one has \(\varphi(y)\leq\varphi(x)\) (respectively \(\varphi(y)\geq\varphi(x)\), respectively \(\varphi(y)=\varphi(x)\)). We say that \(\varphi\) is upper semi-continuous (respectively lower semi-continuous, respectively locally constant) if it is so at every point.

\sgaStableTarget{sga6.e03.ex.5-4-1}\textbf{Examples 5.4.1.}
\begin{enumerate}[label=\alph*)]
\item If \(\underline M\) is the constant sheaf with value \(M\) and \(s\in H^0(X,\underline M)\), then \(x\mapsto x^*s\) is locally constant. In particular, if \(X\) is locally ringed and \(E\) is a perfect complex on \(X\), the rank of \(E\) (\hyperlink{sga6.e01.def.6-11}{I 6.11}) defines a locally constant function on \(\pt(X)\) with values in \(\mathbb Z\).
\item If \(X\) is locally ringed and \(P\) is an \(\mathcal O_X\)-module of finite type, then for \(x\in\pt(X)\), \(P_x\otimes_{\mathcal O_{X,x}}k(x)\) is a finite-dimensional vector space over \(k(x)\), and
\[
 x\longmapsto \dim_{k(x)}(P_x\otimes k(x))
\]
is upper semi-continuous. This is the usual Nakayama argument: an epimorphism \(\mathcal O_{X,x}^{n}\to P_x\) lifts over a pointed neighbourhood.
\end{enumerate}

\sgaStableTarget{sga6.e03.lem.5-5}\textbf{Lemma 5.5.} Let \((X,\mathcal O_X)\) be a locally ringed topos, let \(E\in\Ob D(X)\), and let \(p\in\mathbb Z\). Then:
\begin{enumerate}[label=\roman*)]
\item If \(E\) is perfect, the vector spaces
\[
 H^i(E_x\otimes_{\mathcal O_{X,x}}^{\mathbf L}k(x))
\]
are finite-dimensional over \(k(x)\), zero except for finitely many \(i\), and the function
\[
 x\longmapsto
 \sum_i(-1)^i\dim_{k(x)}H^i(E_x\otimes^{\mathbf L}k(x))
\]
is locally constant on \(\pt(X)\).
\item If \(E\) is \((p-1)\)-pseudo-coherent, then
\[
 H^p(E_x\otimes_{\mathcal O_{X,x}}^{\mathbf L}k(x))
\]
is finite-dimensional over \(k(x)\), and its dimension is upper semi-continuous on \(\pt(X)\). If moreover \(M\mapsto H^p(E\otimes^{\mathbf L}M)\) is exact on \(\Mod(X)\), then \(H^p(E)\) is locally free of finite type and the preceding function is locally constant.
\item If \(E\) is \(p\)-pseudo-coherent (respectively \((p-1)\)-pseudo-coherent), and if for a point \(x\) the functor
\[
 M\longmapsto H^p(E_x\otimes_{\mathcal O_{X,x}}^{\mathbf L}M)
\]
is right exact (respectively left exact) on \(\Mod(\mathcal O_{X,x})\), then there exists an \(x\)-pointed object \(U\) such that
\[
 M\longmapsto H^p(E|_U\otimes^{\mathbf L}M)
\]
is right exact (respectively left exact) on \(\Mod(U)\).
\end{enumerate}

\emph{Proof.} For i), by \hyperlink{sga6.e01.cor.4-19-1}{I 4.19.1} the complex \(E_x\otimes^{\mathbf L}k(x)\) is a perfect complex of \(k(x)\)-modules; hence its cohomology spaces have finite rank and vanish except for finitely many degrees. The identity
\[
 \rank(E_x\otimes^{\mathbf L}k(x))
 =\sum_i(-1)^i\dim_{k(x)}H^i(E_x\otimes^{\mathbf L}k(x))
\]
follows by devissage from the additivity of rank. The assertion follows from the compatibility of the rank with inverse images (\hyperlink{sga6.e01.prop.6-12}{I 6.12}).

For ii), after truncation and translation, one may assume for the first assertion that \(E\) is strictly perfect, \(p=0\), and \(E^i=0\) for \(i\notin[-1,1]\). Put \(E(x)=E_x\otimes^{\mathbf L}k(x)\). Then
\[
 0=\rank(E(x))
 =\dim H^{-1}(E(x))-\dim H^0(E(x))+\dim H^1(E(x))
\]
up to the corresponding alternating-rank convention; hence upper semi-continuity of \(\dim H^0(E(x))\) follows from that of the neighbouring cohomology dimensions and from local constancy of the rank of the terms. For \(H^1\) this is immediate from (\hyperlink{sga6.e03.ex.5-4-1}{5.4.1} b). For \(H^{-1}\), replacing \(E\) by the truncated complex \((0\to E^{-1}\to E^0\to0)\) reduces to the same argument.

For the second assertion in ii), after truncating one may suppose \(E\) strictly perfect and \(E^i=0\) for \(i<p-1\). Then \(Z'^{p+1}(E)=\Coker(E^p\to E^{p+1})\) is of finite presentation and one has an exact sequence
\[
0\to H^p(E)\to Z'^p(E)\to E^{p+1}\to Z'^{p+1}(E)\to0.
\]
Exactness of \(M\mapsto H^p(E\otimes^{\mathbf L}M)\) implies, by (\hyperlink{sga6.e03.rem.5-2}{5.2} iii), that \(Z'^{p}(E)\) and \(Z'^{p+1}(E)\) are flat, hence locally free of finite type by \hyperlink{sga6.e01.lem.5-8-3}{I 5.8.3}. The exact sequence then shows that \(H^p(E)\) is locally free of finite type, and the function \(x\mapsto\dim_{k(x)}H^p(E)_x\) is locally constant.

For iii), one may assume that \(E\) is strictly perfect and that \(E^i=0\) for \(i<p\) (respectively for \(i<p-1\)). Saying that the functor
\[
 M\longmapsto H^p\bigl(E_x\otimes_{\mathcal O_{X,x}}^{\mathbf L}M\bigr)
\]
is right exact (respectively left exact) on \(\Mod(\mathcal O_{X,x})\) is equivalent, by (\hyperlink{sga6.e03.rem.5-2}{5.2} iii), to saying that \(Z'^{p+1}(E_x)\) (respectively \(Z'^p(E_x)\)) is flat, hence free of finite type since it is of finite presentation. Since
\[
 Z'^{p+1}(E_x)\simeq Z'^{p+1}(E)_x
 \qquad\text{(respectively }Z'^p(E_x)\simeq Z'^p(E)_x\text{)},
\]
the usual argument (EGA \(0_I\), 5.2.7) gives an \(x\)-pointed object \(U\) of \(X\) on which \(Z'^{p+1}(E|_U)\) (respectively \(Z'^p(E|_U)\)) is locally free of finite type. This proves the exactness assertion.

\sgaStableTarget{sga6.e03.cor.5-6}\textbf{Corollary 5.6} (EGA III 7.6.9). Let \(X\) be a scheme, let \(p\in\mathbb Z\), and let \(E\in\Ob D^{-}(X)\) be \((p-1)\)-pseudo-coherent. Then:
\begin{enumerate}[label=\roman*)]
\item For \(x\in X\), \(H^p(E\otimes_X^{\mathbf L}k(x))\) has finite rank over \(k(x)\), and
\[
 x\longmapsto \dim_{k(x)}H^p(E\otimes_X^{\mathbf L}k(x))
\]
is constructible and upper semi-continuous on \(X\).
\item Suppose \(X\) is quasi-separated and \(E\) has locally bounded quasi-coherent cohomology. If \(M\mapsto H^p(E\otimes^{\mathbf L}M)\) is exact on the category of quasi-coherent \(\mathcal O_X\)-modules, then the preceding function is locally constant and \(H^p(E)\) is locally free of finite type. Conversely, if \(X\) is locally noetherian and reduced and that function is locally constant, then the same functor is exact on quasi-coherent modules.
\end{enumerate}

\emph{Proof.} If \(X\) is integral with generic point \(x\), the functor \(M\mapsto H^p(E_x\otimes M)\) is exact on \(\mathcal O_{X,x}\)-modules, since \(\mathcal O_{X,x}=k(x)\). By (\hyperlink{sga6.e03.lem.5-5}{5.5} iii) there is a non-empty open neighbourhood \(U\) of \(x\) on which the corresponding functor is exact, and by (\hyperlink{sga6.e03.lem.5-5}{5.5} ii) the rank function is locally constant on \(U\). This gives constructibility. The other assertions follow from (\hyperlink{sga6.e03.lem.5-5}{5.5} ii, iii) and (\hyperlink{sga6.e03.rem.5-2}{5.2} iv), except for the converse in ii), for which one reduces locally and by truncation to the strictly perfect case and applies EGA III 7.6.2 and 7.6.5 to the functor \(T_M=H^p(E\otimes^{\mathbf L}M)\).

\sgaStableTarget{sga6.e03.prop.5-7}\textbf{Proposition 5.7} (EGA III 7.7.5 and 7.9.4). Let \(f:X\to Y\) be a proper morphism of quasi-compact schemes, let \(E\in\Ob D(X)_{\parf}\), and let \(p\in\mathbb Z\). Assuming \hyperlink{sga6.e03.conj.2-1}{(2.1)} is valid for \(f\), for instance if \(f\) is projective or if \(Y\) is noetherian, the following hold:
\begin{enumerate}[label=\roman*)]
\item For \(y\in Y\), the vector spaces \(H^i(Rf_*E\otimes_Y^{\mathbf L}k(y))\) are finite-dimensional over \(k(y)\), zero except for finitely many \(i\), and their alternating rank is locally constant on \(Y\).
\item The function
\[
 y\longmapsto \dim_{k(y)}H^p(Rf_*E\otimes_Y^{\mathbf L}k(y))
\]
is constructible and upper semi-continuous on \(Y\).
\item Suppose \(Y\) is quasi-separated. If \(M\mapsto R^pf_*(E\otimes_Y^{\mathbf L}M)\) is exact on the category of quasi-coherent \(\mathcal O_Y\)-modules, then \(R^pf_*(E)\) is locally free of finite type and the function in ii) is locally constant. Conversely, if \(Y\) is noetherian and reduced and the function in ii) is locally constant, then the same functor is exact on quasi-coherent modules.
\end{enumerate}

\emph{Proof.} By \hyperlink{sga6.e03.prop.4-8}{(4.8)}, \(Rf_*(E)\) is a perfect complex. Assertion i), respectively ii), follows from (\hyperlink{sga6.e03.lem.5-5}{5.5} i), respectively (\hyperlink{sga6.e03.cor.5-6}{5.6} i). If \(M\) is quasi-coherent on \(Y\), the canonical arrow
\[
 Rf_*(E)\otimes_Y^{\mathbf L}M
 \xrightarrow{\sim}
 Rf_*(E\otimes_Y^{\mathbf L}M)
\]
is an isomorphism by \hyperlink{sga6.e03.prop.3-7}{(3.7)}, and iii) follows from (\hyperlink{sga6.e03.cor.5-6}{5.6} ii).

\sgaStableTarget{sga6.e03.rem.5-8}\textbf{Remarks 5.8.}
\begin{enumerate}[label=\alph*)]
\item In the situation of \hyperlink{sga6.e03.prop.5-7}{(5.7)}, let \(i_y:X_y\to X\) be the inclusion of the fiber over a point \(y\in Y\). One may show, under the usual flatness hypotheses, the Künneth formula
\[
 H^i(Rf_*E\otimes_Y^{\mathbf L}k(y))
 \simeq H^i(X_y,E_y),
\]
where \(E_y=Li_y^*E\), or \(E_y=i_y^*E\) in the flat case.
\item One can develop for complex analytic spaces a theory analogous to that of this section: exchange, semi-continuity, and continuity properties. Its essential inputs are Grauert's finiteness theorem for proper morphisms and the analytic analogue of the projection formula. More precisely, for a morphism \(f:X\to Y\) of analytic spaces, with \(X\) of bounded dimension, and for \(E\in\Ob D(X)\), \(M\in\Ob D(Y)_{\coh}\), the canonical arrow
\[
 Rf_*(E)\otimes_Y^{\mathbf L}M
 \longrightarrow
 Rf_*(E\otimes_Y^{\mathbf L}M)
\]
is an isomorphism. Unlike Grauert's theorem, this assertion is essentially formal; it follows by the same devissage as in the proof of \hyperlink{sga6.e03.prop.3-7}{(3.7)}, and requires no hypothesis on \(E\) beyond belonging to the appropriate derived category.
\end{enumerate}

\sgaStableTarget{sga6.e03.bibliography}\subsection*{Bibliography}
\begin{enumerate}[label={[\arabic*]}]
\item Grauert, H., \emph{Ein Theorem der Analytischen Garbentheorie}, Publ. IHES, no. 5.
\item Hartshorne, R., \emph{Residues and Duality}, Lecture Notes in Mathematics, no. 20, Springer, 1966.
\item Godement, R., \emph{Théorie des faisceaux}, Hermann, 1958.
\end{enumerate}



% Exposé IV macros
\providecommand{\Ob}{\operatorname{Ob}}
\providecommand{\Mod}{\operatorname{Mod}}
\providecommand{\Coh}{\operatorname{Coh}}
\providecommand{\Loclib}{\operatorname{Loclib}}
\providecommand{\coh}{\mathrm{coh}}
\providecommand{\parf}{\mathrm{parf}}
\providecommand{\naif}{\mathrm{naif}}
\providecommand{\loc}{\mathrm{loc}}
\providecommand{\fin}{\mathrm{fin}}
\providecommand{\clA}{\operatorname{cl}}
\providecommand{\Rg}{\operatorname{rg}}
\providecommand{\varinjlim}{\varinjlim}
\providecommand{\varprojlim}{\varprojlim}
\providecommand{\id}{\operatorname{id}}
\providecommand{\Th}{\operatorname{Th}}
\providecommand{\Hom}{\operatorname{Hom}}
\providecommand{\Dcat}{\mathrm D}
\providecommand{\Kcat}{\mathrm K}
\providecommand{\Kdot}{K_{\bullet}}
\providecommand{\Kprime}{K'}
\providecommand{\naif}{\mathrm{na\ddot{\i}f}}


\sgaStableTarget{sga6.e04.expose}\section*{Exposé IV. Grothendieck groups of ringed topoi}
\begin{center}
by L. Illusie
\end{center}

\sgaStableTarget{sga6.e04.sec.1}\section*{1. Reminders and generalities on Grothendieck groups}

\sgaStableTarget{sga6.e04.para.1-1}\textbf{1.1.} Let \(C\) be a triangulated category. Recall that a map \(f\) from \(\Ob C\) to an abelian group \(G\) is called \emph{additive} if
\[
 f(E)=f(E')+f(E'')
\]
for every distinguished triangle
\[
 E'\longrightarrow E\longrightarrow E''\longrightarrow E'[1].
\]
The additive maps from \(\Ob C\) to \(G\) form an abelian group, depending functorially on \(G\). The functor so obtained is represented by an abelian group \(k(C)\) and a universal additive map
\[
 \cl_C:\Ob C\longrightarrow k(C),
\]
which we shall write simply as \(\cl\) when no confusion is possible.

The group \(k(C)\) depends functorially on \(C\) with respect to exact functors. If \(C,C'\) are triangulated categories, two isomorphic exact functors from \(C\) to \(C'\) induce the same homomorphism \(k(C)\to k(C')\); in particular, if \(u:C\to C'\) is an exact equivalence of categories, then \(k(u):k(C)\to k(C')\) is an isomorphism.

\sgaStableTarget{sga6.e04.lem.1-2}\textbf{Lemma 1.2.} Let \(C\) be a triangulated category and let \(L\in\Ob C\).
\begin{enumerate}[label=\alph*)]
\item For every \(n\in\mathbb Z\),
\[
 \cl(L[n])=(-1)^n\cl(L).
\]
If \(L',L''\) are objects of \(C\) such that \(L\simeq L'\oplus L''\), then
\[
 \cl(L)=\cl(L')+\cl(L'').
\]
\item Suppose that there exists \(n\in\mathbb Z\) such that
\[
 \coprod_{i\geq0} L[2ni]
\]
is representable. Then \(\cl(L)=0\).
\end{enumerate}

\emph{Proof.} Assertion a) follows directly from the definitions. For b), it is enough to observe that
\[
 \coprod_{i\geq0} L[2ni]\simeq
 L\oplus\left(\coprod_{i\geq0}L[2ni]\right)[2n],
\]
and to apply a).

\sgaStableTarget{sga6.e04.lem.1-3}\textbf{Lemma 1.3.} Let \(C\) be a triangulated category and let \(A\) be a thick subcategory of \(C\) \(([\mathrm V]\ \mathrm{I}\ 2.1)\). The inclusion and quotient functors define an exact sequence
\[
 k(A)\longrightarrow k(C)\longrightarrow k(C/A)\longrightarrow 0.
\]

\emph{Proof.} See SGA 5 VIII 3.1.

\sgaStableTarget{sga6.e04.lem.1-4}\textbf{Lemma 1.4.} Let \(A\) be an additive category, respectively an abelian category, and let \(f\) be an additive function on \(\Ob K^b(A)\), respectively on \(\Ob D^b(A)\). For \(E\in\Ob K^b(A)\), respectively \(E\in\Ob D^b(A)\), one has
\[
 f(E)=\sum_i(-1)^i f(E^i),
\]
respectively
\[
 f(E)=\sum_i(-1)^i f(H^i(E)).
\]

\emph{Proof.} Left as an exercise to the reader.

\sgaStableTarget{sga6.e04.para.1-5}\textbf{1.5.} Let \(C\) be an abelian category and let \(A\) be a full subcategory of \(C\), stable under finite sums. A map \(f\) from \(\Ob A\) to an abelian group \(G\) is called \emph{additive} if
\[
 \sum_i(-1)^i f(E^i)=0
\]
for every acyclic bounded complex \(E\) of \(C\) whose components belong to \(A\). When \(A\) is stable under kernels of epimorphisms (\hyperlink{sga6.e01.sec.1.2}{I 1.2}), this amounts to saying that
\[
 f(L)=f(L')+f(L'')
\]
for every exact sequence
\[
 0\longrightarrow L'\longrightarrow L\longrightarrow L''\longrightarrow0
\]
in \(C\) whose terms belong to \(A\). The additive maps from \(\Ob A\) to \(G\) form an abelian group depending functorially on \(G\). Denote by \(K^{b,\emptyset}(A)\) the thick subcategory of \(K^b(A)\) consisting of acyclic complexes. It follows easily from \hyperlink{sga6.e04.lem.1-3}{(1.3)} and \hyperlink{sga6.e04.lem.1-4}{(1.4)} that the map
\sgaStableTarget{sga6.e04.eq.1-5-1}\begin{equation}\tag{1.5.1}
 \cl_A:\Ob A\longrightarrow k\bigl(K^b(A)/K^{b,\emptyset}(A)\bigr)
\end{equation}
induced by \(\cl_{K^b(A)/K^{b,\emptyset}(A)}\) is a universal additive map, i.e. that the pair
\[
 \left(k\bigl(K^b(A)/K^{b,\emptyset}(A)\bigr),\cl_A\right)
\]
represents the functor which associates to each abelian group \(G\) the abelian group of additive maps from \(\Ob A\) to \(G\). We put
\sgaStableTarget{sga6.e04.eq.1-5-2}\begin{equation}\tag{1.5.2}
 k\bigl(K^b(A)/K^{b,\emptyset}(A)\bigr)=k(A),
\end{equation}
and write \(k(A,C)\) instead of \(k(A)\) when it is necessary to avoid ambiguity about the ambient abelian category.\footnote{The notation \(k(A)\) had been introduced in \hyperlink{sga6.e01.para.6-3}{I 6.3} for the group \(k(K^b(A))\), which is not in general isomorphic to the group defined by \hyperlink{sga6.e04.eq.1-5-2}{(1.5.2)}. We therefore change convention here: in this exposé the notation \(k(A)\) will mean only the group defined by \hyperlink{sga6.e04.eq.1-5-2}{(1.5.2)}.}
In particular, for \(A=C\), one has
\sgaStableTarget{sga6.e04.eq.1-5-3}\begin{equation}\tag{1.5.3}
 k(D^b(C))=k(C).
\end{equation}

\sgaStableTarget{sga6.e04.lem.1-6}\textbf{Lemma 1.6.} Let \(C\) be an abelian category and let \(A\) be a thick abelian subcategory, i.e. a full subcategory stable under kernels, cokernels, and extensions. Let \(D_A(C)\) be the full subcategory of \(D(C)\) formed by the objects \(E\) such that \(H^i(E)\in\Ob A\) for every \(i\). Then \(D_A(C)\) is a triangulated subcategory \(([\mathrm V]\ \mathrm I\ 2.3)\) of \(D(C)\), and the canonical functor
\[
 D^b(A)\longrightarrow D_A^b(C)=D^b(C)\cap D_A(C)
\]
induces an isomorphism
\[
 k(A)\xrightarrow{\sim} k(D_A^b(C)),
\]
where \(k(A)\) is defined by \hyperlink{sga6.e04.eq.1-5-2}{(1.5.2)}.

\emph{Proof.} The first assertion is an immediate consequence of the long exact cohomology sequence. For the second, consider the map
\[
 f:\Ob D_A^b(C)\longrightarrow k(A)
\]
defined by
\[
 f(E)=\sum_i(-1)^i\cl_A(H^i(E)).
\]
By the long exact cohomology sequence, \(f\) is additive in the sense of \hyperlink{sga6.e04.para.1-1}{(1.1)}, hence defines a homomorphism
\[
 \varphi:k(D_A^b(C))\longrightarrow k(A).
\]
It follows at once from \hyperlink{sga6.e04.lem.1-4}{(1.4)} that \(\varphi\) is inverse to the canonical homomorphism \(k(A)\to k(D_A^b(C))\).

\sgaStableTarget{sga6.e04.rem.1-6-1}\textbf{Remark 1.6.1.} In the situation of \hyperlink{sga6.e04.lem.1-6}{(1.6)}, the canonical functor
\[
 D^b(A)\longrightarrow D_A^b(C)
\]
is not in general an equivalence of categories. This is shown by the example where \(S\) is a scheme, \(C\) is the category of \(\mathcal O_S\)-modules, and \(A\) is the category of quasi-coherent \(\mathcal O_S\)-modules (\hyperlink{sga6.e02.para.0-3}{II App. I 0.3}). See, however, the criterion (\hyperlink{sga6.e02.lem.2-1-3}{II 2.1.3}).

\sgaStableTarget{sga6.e04.def.1-7}\textbf{Definition 1.7.} Let \(C\) be a triangulated category, respectively an abelian category, and let \(C'\) be a subset of \(\Ob C\). We shall say that \(C'\) is an \emph{exact subset} if, for every distinguished triangle
\[
 E'\longrightarrow E\longrightarrow E''\longrightarrow E'[1],
\]
respectively for every exact sequence
\[
 0\longrightarrow E'\longrightarrow E\longrightarrow E''\longrightarrow0,
\]
such that two of the terms \(E',E,E''\) belong to \(C'\), the third also belongs to \(C'\). We shall say that \(C'\) is \emph{dense} in \(\Ob C\) if \(\Ob C\) is the smallest exact subset of \(\Ob C\) containing \(C'\).

\sgaStableTarget{sga6.e04.lem.1-8}\textbf{Lemma 1.8.} Let \(C\) be a triangulated category, respectively an abelian category, and let \(C'\) be a dense subset of \(\Ob C\). Then the elements \(\cl_C(E)\), for \(E\in C'\), generate \(k(C)\).

\emph{Proof.} Put \(C'=C_0\), and for \(i\geq1\) define \(C_i\) recursively as the set of objects of \(C\) occurring in a distinguished triangle, respectively a short exact sequence, two of whose terms belong to \(C_{i-1}\). The \(C_i\) form an increasing filtration of \(\Ob C\), which is exhaustive by hypothesis. For \(i\geq0\), let \(k(C)_i\) be the subgroup of \(k(C)\) generated by the elements \(\cl_C(E)\), \(E\in C_i\). It is immediate from the definitions that \(k(C)_i=k(C)_{i+1}\) for every \(i\geq0\), hence \(k(C)_i=k(C)_0\). If \(x\in k(C)\), one can write \(x=\cl(E')-\cl(E'')\) for \(E',E''\in\Ob C\) (in the non-respectable case one may even arrange \(E''=0\)). Since the filtration \((C_i)\) is exhaustive, there exists \(i\geq0\) such that \(E'\) and \(E''\) belong to \(C_i\); hence \(x\in k(C)_i=k(C)_0\), as required.

\sgaStableTarget{sga6.e04.exercise.1-9}\textbf{Exercise 1.9.} Let \(C\) be an abelian category and \(A\) a full subcategory stable under finite sums, subobjects, and quotients, such that \(\Ob A\) is dense in \(\Ob C\). Then the homomorphism
\[
 k(A,C)\longrightarrow k(C)
\]
defined by the additive function \(E\mapsto\cl_C(E)\) on \(\Ob A\) (cf. \hyperlink{sga6.e04.para.1-5}{(1.5)}) is an isomorphism.

\emph{Hint.} Put \(A=A_0\), and define \(A_i\), for \(i\geq1\), as the full subcategory of \(C\) consisting of objects of \(C\) which are extensions of two objects of \(A_{i-1}\). Prove that \(A_i\) satisfies the same hypotheses as \(A\), and then prove that the canonical homomorphism
\[
 \varinjlim_i k(A_i,C)\longrightarrow k(C)
\]
is an isomorphism. The desired conclusion follows at once.

\sgaStableTarget{sga6.e04.sec.2}\section*{2. The functors \(K_{\bullet}\) and \(K'\) of a ringed topos}

\sgaStableTarget{sga6.e04.para.2-1}\textbf{2.1.} Let \((X,A)\) be a ringed topos. Recall the following notation:
\[
 D^b(A_X)_{\coh}
\]
is the category of pseudo-coherent complexes with bounded cohomology of left \(A\)-modules (\hyperlink{sga6.e01.ex.2-15}{I 2.15});
\[
 D(A_X)_{\parf}
\]
is the category of complexes of left \(A\)-modules of finite perfect amplitude (\hyperlink{sga6.e01.rem.5-8-4}{I 5.8.4});
\[
 \Loclib(A_X)
\]
is the category of left \(A\)-modules locally direct factors of finite free modules (\hyperlink{sga6.e02.para.2-0}{II 2.0});
\[
 \Coh(A_X)
\]
is the category of coherent left \(A\)-modules, relative to the category of finite-type free \(A\)-modules (\hyperlink{sga6.e01.def.3-1}{I 3.1}).

The last two categories are full subcategories of the category \(\Mod(A_X)\) of left \(A\)-modules, and the first two are full triangulated subcategories of
\[
 D(A_X)=D(\Mod(A_X)).
\]

\sgaStableTarget{sga6.e04.def.2-2}\textbf{Definition 2.2.} Let \((X,A)\) be a ringed topos. Unless otherwise stated, we shall put
\[
 K_{\bullet}(A_X)=k(D^b(A_X)_{\coh}),\qquad
 K'(A_X)=k(D(A_X)_{\parf}),
\]
where \(k(\ )\) is defined in \hyperlink{sga6.e04.para.1-1}{(1.1)}, and
\[
 K_{\bullet}(A_X)_{\naif}=k(\Coh(A_X)),\qquad
 K'(A_X)_{\naif}=k(\Loclib(A_X)),
\]
where \(k(\ )\) is defined by \hyperlink{sga6.e04.eq.1-5-2}{(1.5.2)}, the ambient abelian category being \(\Mod(A_X)\).\footnote{We shall allow ourselves to omit the index \(A\) when no confusion can result. We shall sometimes also drop the index ``naive'', but will then say so explicitly.} We say that \(K_{\bullet}(A_X)\), respectively \(K'(A_X)\), is the Grothendieck group of pseudo-coherent, respectively perfect, \(A\)-module complexes on \(X\), and that \(K_{\bullet}(A_X)_{\naif}\), respectively \(K'(A_X)_{\naif}\), is the Grothendieck group of coherent, respectively locally free, left \(A\)-modules.

\sgaStableTarget{sga6.e04.rem.2-3}\textbf{Remarks 2.3.}
\begin{enumerate}[label=\alph*)]
\item One can define variants of \(K_{\bullet}(\ )\) and \(K'(\ )\) by putting
\[
 K_{\bullet}^{\loc}(A_X)=k(D^{\loc b}(A_X)_{\coh}),
\]
where ``loc \(b\)'' means ``with locally bounded cohomology'', and
\[
 K_{\loc}'(A_X)=k(D(A_X)_{\parf}).
\]
These groups coincide with those defined in \hyperlink{sga6.e04.def.2-2}{(2.2)} when the final object of \(X\) is quasi-compact, but in general they lend themselves less well to computations in intersection theory.
\item The group \(K_{\bullet}(A_X)_{\naif}\) has practically no interest unless \(A\) is coherent as a left \(A\)-module; otherwise there is a strong risk that very few coherent sheaves exist.
\item Considering the \(k\) of triangulated categories which are ``too large'' is of no interest. Thus, from (\hyperlink{sga6.e04.lem.1-2}{1.2} b), one immediately deduces the relations
\[
 k(D(A_X))=
 k(D^-(A_X))=
 k(D^+(A_X))=
 k(D^b(A_X))=
 k(D^-(A_X)_{\coh})=0,
\]
for the last equality observing that \(L\in\Ob D^-(A_X)_{\coh}\) implies
\[
 \coprod_{i\geq0}L[2i]\in\Ob D^-(A_X)_{\coh}.
\]
\end{enumerate}

\sgaStableTarget{sga6.e04.para.2-4}\textbf{2.4. Case of a coherent sheaf of rings.} Let \((X,A)\) be a ringed topos, with \(A\) assumed coherent as a left \(A\)-module (\hyperlink{sga6.e01.def.3-1}{I 3.1} and \hyperlink{sga6.e01.ex.3-6}{3.6}). Every coherent left \(A\)-module is a pseudo-coherent complex (\hyperlink{sga6.e01.cor.3-5}{I 3.5}), and one has a fully faithful functor
\sgaStableTarget{sga6.e04.eq.2-4-1}\begin{equation}\tag{2.4.1}
 \Coh(A_X)\longrightarrow D^b(A_X)_{\coh}.
\end{equation}
Moreover \(\Coh(A_X)\) is a thick subcategory of \(\Mod(A_X)\) (\hyperlink{sga6.e01.cor.3-4}{I 3.4}), and, by (\hyperlink{sga6.e01.cor.3-5}{I 3.5}), \(D^b(A_X)_{\coh}\) is the full subcategory of \(D^b(A_X)\) formed by the objects \(L\) such that \(H^i(L)\in\Ob\Coh(A_X)\) for all \(i\). Consequently, it follows from \hyperlink{sga6.e04.lem.1-6}{(1.6)} that the homomorphism
\sgaStableTarget{sga6.e04.eq.2-4-2}\begin{equation}\tag{2.4.2}
 K_{\bullet}(A_X)_{\naif}\longrightarrow K_{\bullet}(A_X)
\end{equation}
defined by \hyperlink{sga6.e04.eq.2-4-1}{(2.4.1)} is an isomorphism.

\sgaStableTarget{sga6.e04.para.2-5}\textbf{2.5. The homomorphism \(K'\to K_{\bullet}\).} Let \((X,A)\) be a ringed topos. The inclusion functor
\[
 D(A_X)_{\parf}\longrightarrow D^b(A_X)_{\coh}
\]
defines a homomorphism
\sgaStableTarget{sga6.e04.eq.2-5-1}\begin{equation}\tag{2.5.1}
 K'(A_X)\longrightarrow K_{\bullet}(A_X).
\end{equation}
It is not in general injective or surjective. However, assume the following hypotheses:
\begin{enumerate}[label=(\roman*)]
\item the final object of \(X\) is quasi-compact;
\item \(X\) has enough points, and for every point \(x\) of \(X\), every left \(A_x\)-module has finite Tor-dimension.
\end{enumerate}
Then, by (\hyperlink{sga6.e01.cor.5-10}{I 5.10}),
\[
 D(A_X)_{\parf}=D^b(A_X)_{\coh},
\]
and \hyperlink{sga6.e04.eq.2-5-1}{(2.5.1)} is an isomorphism. The hypotheses (i) and (ii) are satisfied, for example, by a compact space ringed by regular local rings.

\sgaStableTarget{sga6.e04.para.2-6}\textbf{2.6. Pairings between \(K'\) and \(K'\) in the non-commutative case.} Let
\[
 f:(X',A')\longrightarrow (X,A)
\]
be a morphism of ringed topoi, and denote by \(\Mod(A',X',A)\) the category of bimodules which are left \(A'\)-modules and right \(f^{-1}(A)\)-modules. Let
\[
 \Loclib_A^b(A',X',A)
\]
be the full subcategory formed by the objects which belong, as objects over \(A'\), to \(\Loclib(A',X')\). By analogy with \hyperlink{sga6.e04.def.2-2}{(2.2)}, put
\[
 K'_{\naif}(A',X',A)_{A'}=k(\Loclib_A^b(A',X',A)).
\]
Tensor product
\[
 \otimes_A:\Mod(A',X',A)\times\Mod(A_X)\longrightarrow \Mod(A',X')
\]
sends
\[
 (E',E)\longmapsto E'\otimes_{f^{-1}(A)}f^{-1}(E),
\]
and induces a functor
\sgaStableTarget{sga6.e04.eq.2-6-1}\begin{equation}\tag{2.6.1}
 \Loclib_A^b(A',X',A)\times\Loclib(A_X)\longrightarrow\Loclib(A',X').
\end{equation}
The map \((E',E)\mapsto\cl(E'\otimes_A E)\) deduced from \hyperlink{sga6.e04.eq.2-6-1}{(2.6.1)} is biadditive in the sense of \hyperlink{sga6.e04.para.1-5}{(1.5)}, and therefore defines a homomorphism
\sgaStableTarget{sga6.e04.eq.2-6-2}\begin{equation}\tag{2.6.2}
 K'_{\naif}(A',X',A)_{A'}\otimes K'_{\naif}(A_X)\longrightarrow K'_{\naif}(A',X').
\end{equation}

Let \(D^-(A',X',A)\) be the derived category of \(\Mod(A',X',A)\), and let
\[
 D^-(A',X',A)_{A'\text{-}\parf}
\]
be the full subcategory formed by the complexes of finite perfect amplitude relative to \(A'\), i.e. as objects of \(D^-(A',X')\). Put, as in \hyperlink{sga6.e04.def.2-2}{(2.2)},
\[
 K'(A',X',A)_{A'}=k(D^-(A',X',A)_{A'\text{-}\parf}).
\]
The inclusion functors define evident homomorphisms
\sgaStableTarget{sga6.e04.eq.2-6-3-1}\sgaStableTarget{sga6.e04.eq.2-6-3-2}\sgaStableTarget{sga6.e04.eq.2-6-3-3}\begin{align}
 K'_{\naif}(A',X',A)_{A'}&\longrightarrow K'(A',X',A)_{A'}, \tag{2.6.3.1}\\
 K'_{\naif}(A_X)&\longrightarrow K'(A_X), \tag{2.6.3.2}\\
 K'_{\naif}(A',X')&\longrightarrow K'(A',X'). \tag{2.6.3.3}
\end{align}
On the other hand, one knows (\hyperlink{sga6.e01.prop.4-18}{I 4.18}, \hyperlink{sga6.e01.cor.4-19}{4.19}) that the derived tensor functor
\[
 \otimes_A^{\mathbf L}:D^-(A',X',A)\times D^-(A_X)\longrightarrow D^-(A',X')
\]
induces an exact functor
\sgaStableTarget{sga6.e04.eq.2-6-4}\begin{equation}\tag{2.6.4}
 D^-(A',X',A)_{A'\text{-}\parf}\times D(A_X)_{\parf}
 \longrightarrow D(A',X')_{\parf}.
\end{equation}
The latter defines a homomorphism
\sgaStableTarget{sga6.e04.eq.2-6-5}\begin{equation}\tag{2.6.5}
 K'(A',X',A)_{A'}\otimes K'(A_X)\longrightarrow K'(A',X').
\end{equation}
Since \hyperlink{sga6.e04.eq.2-6-1}{(2.6.1)} is induced by \hyperlink{sga6.e04.eq.2-6-4}{(2.6.4)}, the homomorphisms \hyperlink{sga6.e04.eq.2-6-2}{(2.6.2)} and \hyperlink{sga6.e04.eq.2-6-5}{(2.6.5)} are compatible with the homomorphisms \hyperlink{sga6.e04.eq.2-6-3-1}{(2.6.3.1)}, \hyperlink{sga6.e04.eq.2-6-3-2}{(2.6.3.2)}, and \hyperlink{sga6.e04.eq.2-6-3-3}{(2.6.3.3)}, in other words one has a commutative diagram
\sgaStableTarget{sga6.e04.eq.2-6-6}\[
\begin{tikzcd}[column sep=tiny]
K'_{\naif}(A',X',A)_{A'}\otimes K'_{\naif}(A_X)\arrow[r]\arrow[d]
&K'_{\naif}(A',X')\arrow[d]\\
K'(A',X',A)_{A'}\otimes K'(A_X)\arrow[r]
&K'(A',X'),
\end{tikzcd}
\tag{2.6.6}
\]
where the upper and lower horizontal arrows are respectively \hyperlink{sga6.e04.eq.2-6-2}{(2.6.2)} and \hyperlink{sga6.e04.eq.2-6-5}{(2.6.5)}, and the vertical arrows are defined by \hyperlink{sga6.e04.eq.2-6-3-1}{(2.6.3.1)}, \hyperlink{sga6.e04.eq.2-6-3-2}{(2.6.3.2)}, and \hyperlink{sga6.e04.eq.2-6-3-3}{(2.6.3.3)}.

\sgaStableTarget{sga6.e04.para.2-7}\textbf{2.7. Contravariant law and ring structure on \(K'\).} We spell out some important particular cases of the pairings \hyperlink{sga6.e04.para.2-6}{(2.6)}, although their direct study is immediate.
\begin{enumerate}[label=\alph*)]
\item Let
\[
 f:(X',A')\longrightarrow (X,A)
\]
be a morphism of ringed topoi. The functor
\[
 Lf^*:D^-(A_X)\longrightarrow D^-(A',X')
\]
sends \(D(A_X)_{\parf}\) to \(D(A',X')_{\parf}\) (\hyperlink{sga6.e01.cor.4-19-1}{I 4.19.1}), and likewise sends \(\Loclib(A_X)\) to \(\Loclib(A',X')\). Thus it defines a homomorphism
\sgaStableTarget{sga6.e04.eq.2-7-1}\begin{equation}\tag{2.7.1}
 f^*:K'(A_X)\longrightarrow K'(A',X')
\end{equation}
and, respectively,
\sgaStableTarget{sga6.e04.eq.2-7-2}\begin{equation}\tag{2.7.2}
 f^*:K'(A_X)_{\naif}\longrightarrow K'(A',X')_{\naif}.
\end{equation}
This can also be interpreted as multiplication by \(\cl(A')\) in \hyperlink{sga6.e04.eq.2-6-2}{(2.6.2)}, respectively \hyperlink{sga6.e04.eq.2-6-5}{(2.6.5)}. One has a commutative diagram, a special case of \hyperlink{sga6.e04.eq.2-6-6}{(2.6.6)},
\[
\begin{tikzcd}
K'(A_X)_{\naif}\arrow[r,"f^*"]\arrow[d]
&K'(A',X')_{\naif}\arrow[d]\\
K'(A_X)\arrow[r,"f^*"]
&K'(A',X'),
\end{tikzcd}
\]
where the left and right vertical arrows are respectively \hyperlink{sga6.e04.eq.2-6-3-2}{(2.6.3.2)} and \hyperlink{sga6.e04.eq.2-6-3-3}{(2.6.3.3)}.

Let \(g:(X,A)\to(X'',A'')\) be a morphism of ringed topoi. The transitivity of inverse-image functors implies that the homomorphisms of type \hyperlink{sga6.e04.eq.2-7-1}{(2.7.1)}, respectively \hyperlink{sga6.e04.eq.2-7-2}{(2.7.2)}, satisfy
\sgaStableTarget{sga6.e04.eq.2-7-3}\begin{equation}\tag{2.7.3}
 (gf)^*=f^*g^*.
\end{equation}
Thus \(K'(A_X)\), respectively \(K'(A_X)_{\naif}\), is a contravariant functor of \((X,A)\).

\item Let \((X,\mathcal O_X)\) be a commutatively ringed topos. Thanks to the associativity and commutativity of tensor product, \hyperlink{sga6.e04.eq.2-6-5}{(2.6.5)}, respectively \hyperlink{sga6.e04.eq.2-6-2}{(2.6.2)}, with \(A=\mathcal O_X\), \((X',A')=(X,A)\), and \(f=\id\), endows \(K'(X)\), respectively \(K'(X)_{\naif}\), with the structure of a commutative ring with unit
\[
 1=\cl(\mathcal O_X).
\]
The canonical homomorphism
\[
 K'(X)_{\naif}\longrightarrow K'(X)
\]
is a homomorphism of unital rings.

If \(f:(X,\mathcal O_X)\to(Y,\mathcal O_Y)\) is a morphism of commutatively ringed topoi, compatibility of tensor product with inverse image shows that
\[
 f^*:K'(Y)\longrightarrow K'(X),
\]
respectively
\[
 f^*:K'(Y)_{\naif}\longrightarrow K'(X)_{\naif},
\]
is a homomorphism of unital rings.
\end{enumerate}

\sgaStableTarget{sga6.e04.para.2-8}\textbf{2.8. Augmentation of \(K'\) toward \(H^0(\ ,\mathbb Z)\).} Let \((X,\mathcal O_X)\) be a locally ringed topos (\hyperlink{sga6.e01.lem.2-15-1}{I 2.15.1}). Recall that one has defined (\hyperlink{sga6.e01.def.6-11}{I 6.11} to \hyperlink{sga6.e01.prop.6-14}{6.14}) a homomorphism of unital rings
\[
 i:H^0(X,\mathbb Z)\longrightarrow K'(X)_{\naif},
\]
thereby endowing \(K'(X)_{\naif}\), and then \(K'(X)\) through the canonical homomorphism \(K'(X)_{\naif}\to K'(X)\), with the structure of a unital algebra over \(H^0(X,\mathbb Z)\). Moreover, \(K'(X)_{\naif}\), respectively \(K'(X)\), is augmented toward \(H^0(X,\mathbb Z)\) by the rank, and the canonical homomorphism \(K'(X)_{\naif}\to K'(X)\) is a homomorphism of augmented algebras.

If \(f:(X,\mathcal O_X)\to(Y,\mathcal O_Y)\) is a morphism of locally ringed topoi, the homomorphisms
\[
 f^*:K'(Y)\longrightarrow K'(X),\qquad
 f^*:K'(Y)_{\naif}\longrightarrow K'(X)_{\naif}
\]
are compatible (\hyperlink{sga6.e01.prop.6-12}{I 6.12}, \hyperlink{sga6.e01.para.6-13}{6.13}) with the preceding augmented-algebra structures; this is expressed by the commutativity of the diagram, where \(\Rg\) denotes the rank,
\[
\begin{tikzcd}[column sep=tiny]
H^0(Y,\mathbb Z)\arrow[r,"i"]\arrow[d]
&K'(Y)_{\naif}\arrow[r]\arrow[d]
&K'(Y)\arrow[r,"\Rg"]\arrow[d]
&H^0(Y,\mathbb Z)\arrow[d]\\
H^0(X,\mathbb Z)\arrow[r,"i"]
&K'(X)_{\naif}\arrow[r]
&K'(X)\arrow[r,"\Rg"]
&H^0(X,\mathbb Z).
\end{tikzcd}
\]
Later in the Seminar (\hyperlink{sga6.e05.ex.3-9-1}{V 3.9.1} and VI) one will see that \(K'(X)_{\naif}\) can be endowed with a structure richer than that of an algebra augmented toward \(H^0(X,\mathbb Z)\), namely with a structure of a \(\lambda\)-algebra augmented toward \(H^0(X,\mathbb Z)\), viewed as a binomial ring (\hyperlink{sga6.e05.ex.2-8}{V 2.8}), this structure being compatible with inverse-image homomorphisms. Deligne has recently proved that \(K'(X)\) can be endowed with an analogous structure, at least when \(X\) is of characteristic \(p\geq0\), and that the canonical homomorphism
\[
 K'(X)_{\naif}\longrightarrow K'(X)
\]
is a homomorphism of augmented \(\lambda\)-algebras, thereby answering affirmatively a conjecture of Grothendieck (XIV 1).

\sgaStableTarget{sga6.e04.para.2-9}\textbf{2.9. Comparison of \(K'\) and \(K'_{\naif}\).} Let \((X,A)\) be a ringed topos. There is in general no reason for the canonical homomorphism
\[
 K'(X)_{\naif}\longrightarrow K'(X)
\]
to be injective or surjective. However, the ``global resolution existence criteria'' of \hyperlink{sga6.e02.expose}{Exposé II} show that it is an isomorphism in the following cases:
\begin{enumerate}[label=(\roman*)]
\item \(X\) is the Zariski topos of a separated or noetherian divisorial scheme, and \(A=\mathcal O_X\) (\hyperlink{sga6.e02.prop.2-2-9}{II 2.2.9});
\item \(X\) is the topos defined by a compact space and \(A\) is a soft sheaf (\hyperlink{sga6.e02.prop.2-3-2}{II 2.3.2});
\item \(X\) is a Stein topos whose final object is quasi-compact and \(A=\mathcal O_X\) (\hyperlink{sga6.e02.prop.2-4-5}{II 2.4.5}).
\end{enumerate}

\sgaStableTarget{sga6.e04.para.2-10}\textbf{2.10. Pairing between \(K'\) and \(K_{\bullet}\).} Let \((X,\mathcal O_X)\) be a commutatively ringed topos. One knows (\hyperlink{sga6.e01.rem.5-8-4}{I 5.8.4} b) that there is a pairing
\[
 \otimes_{\mathcal O_X}^{\mathbf L}:D(X)_{\parf}\times D^b(X)_{\coh}
 \longrightarrow D^b(X)_{\coh}.
\]
It defines a homomorphism
\sgaStableTarget{sga6.e04.eq.2-10-1}\begin{equation}\tag{2.10.1}
 K'(X)\otimes K_{\bullet}(X)\longrightarrow K_{\bullet}(X).
\end{equation}
Thanks to associativity of tensor product, \hyperlink{sga6.e04.eq.2-10-1}{(2.10.1)} endows \(K_{\bullet}(X)\) with the structure of a \(K'(X)\)-module. More generally, if
\[
 f:(X,\mathcal O_X)\longrightarrow (Y,\mathcal O_Y)
\]
is a morphism of commutatively ringed topoi, \(K_{\bullet}(X)\) may be regarded as a \(K'(Y)\)-module through the homomorphism
\[
 f^*:K'(Y)\longrightarrow K'(X).
\]

\sgaStableTarget{sga6.e04.para.2-11}\textbf{2.11. Covariant law on \(K_{\bullet}\) and projection formula.} Let \((X,A)\) be a ringed topos. In contrast with \(K'(A_X)\), the group \(K_{\bullet}(A_X)\) has no natural variance with respect to \((X,A)\). Nevertheless, in algebraic geometry and analytic geometry, \(K_{\bullet}(X)\) tends to be covariant in \(X\) for proper morphisms. We now make this precise.

\sgaStableTarget{sga6.e04.para.2-11-1}\textbf{2.11.1.} Let \(f:X\to Y\) be a proper and pseudo-coherent morphism (\hyperlink{sga6.e03.def.1-2}{III 1.2}) of quasi-compact schemes. Then, subject to the validity of the finiteness conjecture for \(f\) (\hyperlink{sga6.e03.conj.2-1}{III 2.1}), hence in particular if \(Y\) is noetherian (in which case the pseudo-coherence hypothesis on \(f\) is automatically verified (\hyperlink{sga6.e03.ex.1-3}{III 1.3})) or if \(f\) is projective, the functor \(Rf_*\) induces, by (\hyperlink{sga6.e03.schol.2-6}{III 2.6}), a functor
\[
 Rf_*:D^b(X)_{\coh}\longrightarrow D^b(Y)_{\coh}.
\]
It therefore defines a homomorphism
\sgaStableTarget{sga6.e04.eq.2-11-1-1}\begin{equation}\tag{2.11.1.1}
 f_*:K_{\bullet}(X)\longrightarrow K_{\bullet}(Y).
\end{equation}
This homomorphism is \(K'(Y)\)-linear (cf. \hyperlink{sga6.e04.para.2-10}{(2.10)}); in other words, for \(y\in K'(Y)\) and \(x\in K_{\bullet}(X)\), one has the \emph{projection formula}
\sgaStableTarget{sga6.e04.eq.2-11-1-2}\begin{equation}\tag{2.11.1.2}
 f_*(yx)=y f_*(x).
\end{equation}
Indeed, if \(x=\cl(E)\) for \(E\in\Ob D^b(X)_{\coh}\) and \(y=\cl(M)\) for \(M\in\Ob D(Y)_{\parf}\), then by (\hyperlink{sga6.e03.prop.3-7}{III 3.7}) one has
\[
 Rf_*(E)\otimes_Y^{\mathbf L}M\simeq
 Rf_*(E\otimes_Y^{\mathbf L}M).
\]
Finally, if \(g:Y\to Z\) is a proper morphism satisfying the same hypotheses as \(f\), and for which the finiteness conjecture (\hyperlink{sga6.e03.conj.2-1}{III 2.1}) is true, the transitivity of direct-image functors gives, for the homomorphisms \hyperlink{sga6.e04.eq.2-11-1-1}{(2.11.1.1)}, the formula
\sgaStableTarget{sga6.e04.eq.2-11-1-3}\begin{equation}\tag{2.11.1.3}
 g_*f_*=(gf)_*.
\end{equation}

\sgaStableTarget{sga6.e04.para.2-11-2}\textbf{2.11.2.} Let \(f:X\to Y\) be a proper morphism of finite-dimensional analytic spaces. Then, by Grauert's finiteness theorem [2], \(Rf_*\) sends \(D^b(X)_{\coh}\) into \(D^b(Y)_{\coh}\), and hence defines a homomorphism
\sgaStableTarget{sga6.e04.eq.2-11-2-1}\begin{equation}\tag{2.11.2.1}
 f_*:K_{\bullet}(X)\longrightarrow K_{\bullet}(Y).
\end{equation}
Like its algebraic analogue \hyperlink{sga6.e04.eq.2-11-1-1}{(2.11.1.1)}, this homomorphism is \(K'(Y)\)-linear; equivalently, formula \hyperlink{sga6.e04.eq.2-11-1-2}{(2.11.1.2)} is valid, as follows from the projection formula at the level of derived categories (\hyperlink{sga6.e03.rem.5-8}{III 5.9} b). One also has the transitivity formula \hyperlink{sga6.e04.eq.2-11-1-3}{(2.11.1.3)}.

\sgaStableTarget{sga6.e04.para.2-11-3}\textbf{2.11.3.} Outside the algebraic or analytic framework, this covariance of \(K_{\bullet}\) for proper morphisms disappears. For instance, let
\[
 f:X\to Y
\]
be a proper morphism of topological spaces ringed by the sheaves of complex continuous functions. Then, if \(E\in\Ob D^b(X)_{\coh}\), in general \(Rf_*(E)\) need not belong to \(D^b(Y)_{\coh}\), as is already evident for \(E=\mathcal O_X\) and \(Y\) reduced to a point. In fact, the \hyperlink{sga6.e04.def.2-2}{definition (2.2)} of \(K_{\bullet}\) is especially adapted to the needs of algebraic and analytic geometry, and a good definition of \(K_{\bullet}\) for a topological space remains to be found. It seems, moreover, that one will have to leave the framework of \(\mathcal O_X\)-linear complexes, and define, for a morphism \(f:X\to Y\) of topological spaces, or of mixed varieties (II App. II), a suitable class of complexes of topological vector-sheaves on \(X\), meant to contain, when \(f\) is the structural morphism of a mixed variety, the elliptic complexes on \(X\) relative to \(Y\). This is suggested by the finiteness theorem (II App. II), as well as by the analogous theorem, where \(Y\) is assumed reduced to a point, proved by Grothendieck in [3].

\sgaStableTarget{sga6.e04.para.2-12}\textbf{2.12. Unusual variances of \(K'\) and \(K_{\bullet}\).} Let
\[
 f:(X,A)\longrightarrow(Y,B)
\]
be a morphism of finite Tor-dimension of ringed topoi (\hyperlink{sga6.e03.def.3-1}{III 3.1}). The functor
\[
 Lf^*:D^-(B_Y)\longrightarrow D^-(A_X)
\]
sends \(D^b(B_Y)_{\coh}\) into \(D^b(A_X)_{\coh}\) (\hyperlink{sga6.e01.cor.2-16-1}{I 2.16.1}), and therefore defines a homomorphism
\sgaStableTarget{sga6.e04.eq.2-12-1}\begin{equation}\tag{2.12.1}
 f^*:K_{\bullet}(B_Y)\longrightarrow K_{\bullet}(A_X).
\end{equation}
If \(A\) and \(B\) are commutative, commutation of inverse image with tensor product shows that \hyperlink{sga6.e04.eq.2-12-1}{(2.12.1)} is compatible with the module structures of \(K_{\bullet}(Y)\) and \(K_{\bullet}(X)\) over \(K'(Y)\) and \(K'(X)\) respectively \hyperlink{sga6.e04.para.2-10}{(2.10)}. In other words, for \(a\in K'(Y)\) and \(y\in K_{\bullet}(Y)\), one has
\sgaStableTarget{sga6.e04.eq.2-12-2}\begin{equation}\tag{2.12.2}
 f^*(ay)=a f^*(y).
\end{equation}
If \(g:(Y,B)\to(Z,C)\) is a morphism of finite Tor-dimension of ringed topoi, transitivity of inverse-image functors implies that the homomorphisms \hyperlink{sga6.e04.eq.2-12-1}{(2.12.1)} satisfy
\sgaStableTarget{sga6.e04.eq.2-12-1-1}\begin{equation}\tag{2.12.1.1}
 f^*g^*=(gf)^*.
\end{equation}

Now suppose that we are in the situation of \hyperlink{sga6.e04.para.2-11-1}{(2.11.1)}, respectively \hyperlink{sga6.e04.para.2-11-2}{(2.11.2)}, with \(f\) of finite Tor-dimension. Then by (\hyperlink{sga6.e03.cor.4-8-1}{III 4.8.1} and \hyperlink{sga6.e03.rem.5-8}{5.9} b), the functor \(Rf_*\) sends \(D(X)_{\parf}\) into \(D(Y)_{\parf}\), and hence defines a homomorphism
\sgaStableTarget{sga6.e04.eq.2-12-3}\begin{equation}\tag{2.12.3}
 f_*:K'(X)\longrightarrow K'(Y).
\end{equation}
This is not generally a ring homomorphism. Nevertheless, for \(y\in K_{\bullet}(Y)\), respectively \(K'(Y)\), and \(x\in K'(X)\), one has
\sgaStableTarget{sga6.e04.eq.2-12-4}\begin{equation}\tag{2.12.4}
 f_*(f^*(y)x)=y f_*(x),
\end{equation}
where \(f^*\) is the homomorphism \hyperlink{sga6.e04.eq.2-12-1}{(2.12.1)}, respectively \hyperlink{sga6.e04.eq.2-7-1}{(2.7.1)}. The proof is analogous to that of \hyperlink{sga6.e04.eq.2-11-1-2}{(2.11.1.2)}. If \(g:Y\to Z\) is a morphism satisfying the same hypotheses as \(f\), the homomorphisms \hyperlink{sga6.e04.eq.2-12-3}{(2.12.3)} satisfy the transitivity formula
\sgaStableTarget{sga6.e04.eq.2-12-3-1}\begin{equation}\tag{2.12.3.1}
 g_*f_*=(gf)_*.
\end{equation}

The homomorphism \hyperlink{sga6.e04.eq.2-12-3}{(2.12.3)} is one of the essential ingredients of the Riemann--Roch theorem, both in algebraic geometry and in analytic geometry, for the indications in the analytic case see \hyperlink{sga6.e00.sec.5}{Exposé 0, 5}.

Suppose that \(X\) and \(Y\) are separated and divisorial schemes (\hyperlink{sga6.e02.def.2-2-5}{II 2.2.5}), for example schemes possessing ample invertible sheaves. Then \hyperlink{sga6.e04.para.2-9}{(2.9 (i))} identifies \(K'(X)\), respectively \(K'(Y)\), with \(K'(X)_{\naif}\), respectively \(K'(Y)_{\naif}\), and \hyperlink{sga6.e04.eq.2-12-3}{(2.12.3)} therefore defines a homomorphism
\sgaStableTarget{sga6.e04.eq.2-12-3.occ2}\begin{equation}\tag{2.12.3'}
 f_*:K'(X)_{\naif}\longrightarrow K'(Y)_{\naif}.
\end{equation}
As was already explained in \hyperlink{sga6.e00.sec.3.1}{Exposé 0, 3.1}, the definition of (\hyperlink{sga6.e04.eq.2-12-3}{2.12.3}') was far from evident a priori. In practice, it is precisely one of the main motivations of this Seminar, from \hyperlink{sga6.e01.expose}{Exposé I} through \hyperlink{sga6.e04.expose}{Exposé IV}.

\sgaStableTarget{sga6.e04.sec.3}\section*{3. Complements on Grothendieck groups of schemes}

The results of this section will not be used in \hyperlink{sga6.e05.expose}{Exposés V} to VIII.

\sgaStableTarget{sga6.e04.prop.3-1-0}\textbf{Proposition 3.1.0.} Let
\[
\begin{tikzcd}
X' \arrow[r,"g"]\arrow[d,"f'"'] & X\arrow[d,"f"]\\
Y' \arrow[r,"h"'] & Y
\end{tikzcd}
\]
be a cartesian square of schemes, where \(Y\) is quasi-compact and \(f\) is quasi-compact and quasi-separated. Let \(E\in\Ob D^b(X)\) be a complex with quasi-coherent cohomology. Suppose that \(X\) and \(Y'\) are Tor-independent over \(Y\) (\hyperlink{sga6.e03.def.1-5}{III 1.5}). Suppose also either that \(h\) has finite Tor-dimension, or that \(E\) has finite flat amplitude relative to \(f\). Under these conditions, the base-change arrow
\[
 Lh^*\,Rf_*E\longrightarrow Rf'_*\,Lg^*E
\]
is defined and is an isomorphism.

\emph{Proof.} To define the arrow, one uses the trivial duality formula (SGA 4 XVII 2.3) for \(h\), which says that there is a canonical functorial isomorphism
\[
 \Hom(Lh^*L,M)\xrightarrow{\sim}\Hom(L,Rh_*M)
\]
for \(L\in\Ob D^-(Y)\) and \(M\in\Ob D^+(Y')\). Using (\hyperlink{sga6.e03.cor.3-5-2}{III 3.5.2}) and (\hyperlink{sga6.e03.cor.3-7-1}{III 3.7.1}), one easily verifies that one may take \(L=Rf_*E\) and \(M=Rf'_*Lg^*E\), so that applying this formula reduces the problem to defining an arrow
\[
 Rf_*E\longrightarrow Rf_*Rg_*Lg^*E.
\]
For this arrow we take the image under \(Rf_*\) of the adjunction map
\[
 E\longrightarrow Rg_*Lg^*E,
\]
which comes from the trivial duality formula for \(g\). The base-change arrow is therefore defined.

It remains to show that it is an isomorphism. For this, one may first suppose that \(Y\) and \(Y'\) are affine. Then, by the Leray spectral sequence of an open covering of \(X\), one is reduced to the case where \(X\) and \(X'\) are affine. One thus has a cartesian square of rings
\[
\begin{tikzcd}
B\arrow[r] & B'\\
A\arrow[u]\arrow[r] & A'\arrow[u],
\end{tikzcd}
\]
and one may suppose that \(E\) is defined by a bounded-above complex of \(B\)-modules. The Tor-independence of \(A'\) and \(B\) over \(A\), written
\[
 B\otimes_A^{\mathbf L}A'\xrightarrow{\sim}B',
\]
then implies, in \(D(A')\),
\[
 E\otimes_A^{\mathbf L}A'\xrightarrow{\sim}E\otimes_B^{\mathbf L}B',
\]
which proves the proposition.

Translating \hyperlink{sga6.e04.prop.3-1-0}{(3.1.0)} into terms of Grothendieck groups gives:

\sgaStableTarget{sga6.e04.prop.3-1-1}\textbf{Proposition 3.1.1.} In the cartesian square of schemes
\[
\begin{tikzcd}
X' \arrow[r,"g"]\arrow[d,"f'"'] & X\arrow[d,"f"]\\
Y' \arrow[r,"h"'] & Y,
\end{tikzcd}
\]
suppose that \(Y\) and \(Y'\) are quasi-compact, that \(X\) and \(Y'\) are Tor-independent over \(Y\) (\hyperlink{sga6.e03.def.1-5}{III 1.5}), that \(f\) is proper and pseudo-coherent, and that the finiteness conjecture (\hyperlink{sga6.e03.conj.2-1}{III 2.1}) is valid for \(f\) and \(f'\), for example if \(Y\) and \(Y'\) are noetherian, or if \(f\) is projective. Then, if \(f\) is perfect, respectively if \(h\) has finite Tor-dimension, for \(x\in K'(X)\), respectively \(x\in K_{\bullet}(X)\), one has
\[
 h^* f_*(x)=f'_* g^*(x),
\]
where \(h^*,g^*\) are defined by \hyperlink{sga6.e04.eq.2-7-1}{(2.7.1)}, respectively \hyperlink{sga6.e04.eq.2-12-1}{(2.12.1)}, and \(f_*,f'_*\) by \hyperlink{sga6.e04.eq.2-12-3}{(2.12.3)}, respectively \hyperlink{sga6.e04.eq.2-11-1-1}{(2.11.1.1)}. Notice that by (\hyperlink{sga6.e03.cor.4-7-1}{III 4.7.1}) \(f'\) is perfect if \(f\) is, and by (\hyperlink{sga6.e03.cor.3-5-1}{III 3.5.1}) \(g\) has finite Tor-dimension if \(h\) does.

\sgaStableTarget{sga6.e04.para.3-2}\textbf{3.2. Passage to the limit on \(K'_{\naif}\) and \(K_{\bullet}\).} Let \((X_i,f_{ij})_{i\in I}\) be a filtered projective system of quasi-compact and quasi-separated schemes whose transition morphisms are affine (cf. EGA IV 8.2.2). Put
\[
 X=\varprojlim_i X_i,
\]
and let \(f_i:X\to X_i\) denote the canonical morphism.

\sgaStableTarget{sga6.e04.para.3-2-1}\textbf{3.2.1.} The homomorphisms \hyperlink{sga6.e04.eq.2-7-2}{(2.7.2)}
\[
 f_i^*:K'(X_i)_{\naif}\longrightarrow K'(X)_{\naif}
\]
define, by the transitivity formula \hyperlink{sga6.e04.eq.2-7-3}{(2.7.3)}, a homomorphism
\sgaStableTarget{sga6.e04.eq.3-2-1-1}\begin{equation}\tag{3.2.1.1}
 \varinjlim_i K'(X_i)_{\naif}\longrightarrow K'(X)_{\naif}.
\end{equation}
I say that this homomorphism is an isomorphism. Indeed, let \(E\) be a locally free \(\mathcal O_X\)-module of finite type. By EGA IV 8.5.2 and 8.5.5, there exists an index \(i\) and a locally free \(\mathcal O_{X_i}\)-module of finite type \(E_i\) such that
\[
 E\simeq f_i^*(E_i).
\]
By EGA IV 8.5.2.5, the image of \(\cl(E_i)\) in \(\varinjlim_i K'(X_i)_{\naif}\) does not depend on the choice of \(E_i\); denote it by \(u(E)\). By EGA IV 8.5.2 and 8.5.6, the map
\[
 u:\Ob\Loclib(X)\longrightarrow \varinjlim_i K'(X_i)_{\naif}
\]
so defined is additive in the sense of \hyperlink{sga6.e04.para.1-5}{(1.5)}. It therefore defines a homomorphism
\[
 K'(X)_{\naif}\longrightarrow \varinjlim_i K'(X_i)_{\naif},
\]
and one verifies immediately, using EGA IV 8.5.2, that it is inverse to (3.2.1.1). This proves the assertion.

\sgaStableTarget{sga6.e04.para.3-2-2}\textbf{3.2.2.} Likewise one defines a canonical homomorphism
\sgaStableTarget{sga6.e04.eq.3-2-2-1}\begin{equation}\tag{3.2.2.1}
 \varinjlim_i K'(X_i)\longrightarrow K'(X).
\end{equation}
Using Deligne's localization technique (SGA 4 XVII), which implies that the category of perfect complexes on \(X\) is equivalent to the category
\[
 \varinjlim_i \Parf(X_i)
\]
of the analogous categories of perfect complexes on the \(X_i\), one proves again that \hyperlink{sga6.e04.eq.3-2-2-1}{(3.2.2.1)} is an isomorphism.

\sgaStableTarget{sga6.e04.para.3-2-3}\textbf{3.2.3.} Suppose that the \(X_i\), and \(X\), are noetherian, and that the morphisms \(f_i\) are flat (then so are the \(f_{ij}\)). Then the homomorphisms \hyperlink{sga6.e04.eq.2-12-1}{(2.12.1)}
\[
 f_i^*:K_{\bullet}(X_i)\longrightarrow K_{\bullet}(X)
\]
define an isomorphism
\[
 \varinjlim_i K_{\bullet}(X_i)\xrightarrow{\sim}K_{\bullet}(X).
\]
The proof is analogous to that of \hyperlink{sga6.e04.para.3-2-1}{(3.2.1)}, using the fact \hyperlink{sga6.e04.para.2-4}{(2.4)} that \(K_{\bullet}(X)\simeq K_{\bullet}(X)_{\naif}\) is computed in terms of coherent sheaves on \(X\).

Moreover, the module structures of the \(K_{\bullet}(X_i)\) over the \(K'(X_i)_{\naif}\) are compatible with the inverse images \(f_{ij}^*\) by \hyperlink{sga6.e04.eq.2-7-1}{(2.7.1)} and \hyperlink{sga6.e04.eq.2-12-2}{(2.12.2)}. Passing to the limit, one finds that \(\varinjlim_i K_{\bullet}(X_i)\) is a module over \(\varinjlim_i K'(X_i)_{\naif}\); this is compatible, through the isomorphisms \hyperlink{sga6.e04.eq.3-2-1-1}{(3.2.1.1)} and \hyperlink{sga6.e04.eq.3-2-2-1}{(3.2.2.1)}, with the module structure of \(K_{\bullet}(X)\) over \(K'(X)_{\naif}\). In other words, one has a commutative diagram
\[
\begin{tikzcd}[column sep=tiny]
\varinjlim K'(X_i)_{\naif}\otimes \varinjlim K_{\bullet}(X_i)
\arrow[r]\arrow[d,"\sim"']
&\varinjlim K_{\bullet}(X_i)\arrow[d,"\sim"]\\
K'(X)_{\naif}\otimes K_{\bullet}(X)\arrow[r]
&K_{\bullet}(X).
\end{tikzcd}
\]

\sgaStableTarget{sga6.e04.para.3-3}\textbf{3.3. Relative Grothendieck groups.}

\sgaStableTarget{sga6.e04.para.3-3-1}\textbf{3.3.1.} Let
\[
 f:X\longrightarrow Y
\]
be a morphism of schemes locally of finite type. Put
\[
 K'(f)=k(D(f)_{\parf}),\qquad
 K_{\bullet}(f)=k(D^b(f)_{\coh}),
\]
where \(D^b(f)_{\coh}\) denotes the category of complexes which are pseudo-coherent relative to \(f\) and have bounded cohomology (\hyperlink{sga6.e03.def.1-2}{III 1.2}), and \(D(f)_{\parf}\) the category of complexes of finite perfect amplitude relative to \(f\) (\hyperlink{sga6.e03.notation.4-2}{III 4.2}). When there is no possible ambiguity, we shall sometimes write \(K'(X/Y)\), respectively \(K_{\bullet}(X/Y)\), instead of \(K'(f)\), respectively \(K_{\bullet}(f)\).

The groups \(K'(f)\) and \(K_{\bullet}(f)\) do not possess the rich structure of their absolute analogues \(K'(X)\) and \(K_{\bullet}(X)\), for tensor product does not in general define homomorphisms
\[
 K'(f)\otimes K'(f)\longrightarrow K'(f),
 \qquad
 K'(f)\otimes K_{\bullet}(f)\longrightarrow K_{\bullet}(f).
\]
On the other hand, they satisfy interesting variance properties, which we now review.

\sgaStableTarget{sga6.e04.para.3-3-2}\textbf{3.3.2.} Let
\[
\begin{tikzcd}
X\arrow[rr,"f"]\arrow[dr,"h"'] &&Y\arrow[dl,"g"]\\
&S&
\end{tikzcd}
\]
be a commutative triangle of morphisms locally of finite type of schemes.
\begin{enumerate}[label=\textbf{3.3.2.\arabic*.}, leftmargin=4em]
\item If \(g\) is pseudo-coherent, one has
\[
 K_{\bullet}(f)=K_{\bullet}(h)
\]
(\hyperlink{sga6.e03.cor.1-12}{III 1.12}). If \(g\) is smooth and \(X\) quasi-compact, one has
\[
 K'(f)=K'(h)
\]
(\hyperlink{sga6.e03.prop.4-4}{III 4.4.1}).

\item By (\hyperlink{sga6.e03.prop.1-11}{III 1.11} and \hyperlink{sga6.e03.prop.4-5}{4.5}), the functor \(\otimes_Y^{\mathbf L}\) defines pairings
\[
 D(f)_{\parf}\times D^b(g)_{\coh}\longrightarrow D^b(h)_{\coh},
\]
\[
 D(f)_{\parf}\times D(g)_{\parf}\longrightarrow D(h)_{\parf},
\]
which induce homomorphisms
\[
 K'(f)\otimes K_{\bullet}(g)\longrightarrow K_{\bullet}(h),
\]
\[
 K'(f)\otimes K'(g)\longrightarrow K'(h).
\]
In particular, taking \(X=Y\) and \(f=\id\), one deduces that \(K_{\bullet}(g)\) is a \(K'(Y)\)-module. On the other hand, if \(f\) is perfect, product by
\[
 \cl(\mathcal O_X)\in K'(f)
\]
defines \(K'(Y)\)-linear homomorphisms
\[
 f^*:K_{\bullet}(Y/S)\longrightarrow K_{\bullet}(X/S),
\]
and
\[
 f^*:K'(Y/S)\longrightarrow K'(X/S).
\]

\item Suppose \(S\) quasi-compact, \(X\) and \(Y\) of finite type over \(S\), and \(f\) proper. Then, subject to the validity of the finiteness conjecture (\hyperlink{sga6.e03.conj.2-1}{III 2.1}), in particular if \(Y\) is noetherian or if \(f\) is projective, the functor \(Rf_*\) sends \(D^b(h)_{\coh}\) into \(D^b(g)_{\coh}\), and therefore defines a homomorphism
\[
 f_*:K_{\bullet}(X/S)\longrightarrow K_{\bullet}(Y/S).
\]
If in addition \(f\) is perfect, \(Rf_*\) sends \(D(h)_{\parf}\) into \(D(g)_{\parf}\) (\hyperlink{sga6.e03.prop.4-8}{III 4.8}), and defines a homomorphism
\[
 f_*:K'(X/S)\longrightarrow K'(Y/S).
\]
From the projection formula (\hyperlink{sga6.e03.prop.3-7}{III 3.7}) one deduces
\[
 f_*(a x)=a f_*(x),
\]
for \(f\) proper, \(a\in K'(S)\), \(x\in K_{\bullet}(X)\), respectively for \(f\) proper and perfect, \(a\in K_{\bullet}(S)\), \(x\in K'(X/S)\).

\item Suppose \(S\) noetherian, \(X\) and \(Y\) of finite type over \(S\), and \(f\) compactifiable and perfect. Then by (\hyperlink{sga6.e03.cor.4-9-1}{III 4.9.1}), the functor \(f^!\) sends \(D(g)_{\parf}\) into \(D(h)_{\parf}\), and defines a homomorphism
\[
 f^!:K'(Y/S)\longrightarrow K'(X/S).
\]
\end{enumerate}

\sgaStableTarget{sga6.e04.para.3-3-3}\textbf{3.3.3.} Under the hypotheses of (\hyperlink{sga6.e03.prop.4-7}{III 4.7}), the functor \(\otimes_S^{\mathbf L}\) defines pairings
\[
 D(f_1)_{\parf}\times D^b(f_2)_{\coh}\longrightarrow D^b(f)_{\coh},
\]
\[
 D(f_1)_{\parf}\times D(f_2)_{\parf}\longrightarrow D(f)_{\parf},
\]
which induce homomorphisms
\[
 K'(f_1)\otimes K_{\bullet}(f_2)\longrightarrow K_{\bullet}(f),
\]
\[
 K'(f_1)\otimes K'(f_2)\longrightarrow K'(f).
\]
Of course, for \(X_i=Y_i=S\), one recovers the \(K'(S)\)-module structure on \(K_{\bullet}(S)\).

If \(X\) is flat over \(S\), one deduces, for \(n\geq1\), homomorphisms
\[
 K'(X/S)^{\otimes n}\longrightarrow K'\bigl((X/S)^n\bigr),
\]
which should, in a sense, replace the absent ring structure on \(K'(X/S)\).

\sgaStableTarget{sga6.e04.bibliography}\subsection*{Bibliography}
\begin{enumerate}[label={[\arabic*]}]
\item Deligne, P., papiers secrets.
\item Grauert, H., \emph{Ein Theorem der Analytischen Garbentheorie}, Publ. Math. IHES, no. 5 (1960).
\item Grothendieck, A., \emph{Théorèmes de finitude pour la cohomologie des faisceaux}, Bull. Soc. Math. France, 84, 1956, pp. 1--7.
\item[{[V]}] Verdier, J.-L., \emph{Catégories dérivées}, IHES (1964), miméographié.
\end{enumerate}



% Expose V macros
\providecommand{\Chat}{\widehat C}
\providecommand{\Ghat}{\widehat G}
\providecommand{\Gg}{\widehat G_g}
\providecommand{\Ga}{\widehat G_a}
\providecommand{\Grhat}{\widehat{\operatorname{Gr}}}
\providecommand{\Abfil}{\mathrm{Ab\,fil}}
\providecommand{\Kalgfil}{K\text{-}\mathrm{Alg.fil.}}
\providecommand{\Hom}{\operatorname{Hom}}
\providecommand{\Sym}{\mathfrak S}
\providecommand{\id}{\operatorname{Id}}
\providecommand{\Log}{\operatorname{Log}}
\providecommand{\Toddc}{\operatorname{Todd}}
\providecommand{\Ob}{\operatorname{Ob}}
\providecommand{\cG}{\mathcal G}
\providecommand{\cCzero}{\mathcal C_0}
\providecommand{\Fil}{\operatorname{Fil}}


\sgaStableTarget{sga6.e05.expose}\section*{Exposé V. Generalities on \(\lambda\)-rings}
\begin{center}
by P. Berthelot
\end{center}

The theory of \(\lambda\)-rings is useful whenever one uses rings such as the ring of classes of vector bundles on a ringed site, the representation ring of a group, and so on, for which the exterior-power operations define maps \(\lambda^i\). We shall use it abundantly later in the Seminar. It is therefore necessary first to develop the general formalism. This exposé is accordingly entirely independent of the rest of the Seminar, and may be useful to users of \(\lambda\)-rings outside algebraic geometry.

The first paragraph essentially gives the theory of Hirzebruch's ``multiplicative sequences of polynomials'' (cf. [2]). Besides its generality, the functorial viewpoint adopted here has the advantage of considerably reducing the calculations. The following paragraphs will show how, to every augmented \(\lambda\)-ring, one can associate a ``Chern ring'' which permits the development of a theory of Chern classes. This is the theory which will allow us to state the Riemann--Roch theorem in the general form of \hyperlink{sga6.e08.expose}{Exposé VIII}.

\sgaStableTarget{sga6.e05.sec.1}\section*{1. Universal polynomials}

\sgaStableTarget{sga6.e05.para.1-1}\textbf{1.1.} Let \(K\) be a commutative ring with unit, fixed throughout this paragraph. Let \(\mathcal C\) denote the category of filtered commutative \(K\)-algebras with unit, whose filtration is decreasing and such that the group of elements of filtration zero is the whole algebra. We define on \(\mathcal C\) a functor with values in the category of abelian groups as follows. If \(A\in\operatorname{Ob}(\mathcal C)\), associate to \(A\) the set of formal series \(\psi\in A[[t]]\) satisfying the following conditions: if
\[
  \psi=a_0+a_1t+\cdots+a_nt^n+\cdots ,
\]
then
\sgaStableTarget{sga6.e05.eq.1-1-1}\begin{equation}\tag{1.1.1}
  a_0=1,\qquad \forall i\geq 1,\quad a_i\in \Fil^i A .
\end{equation}
This set will be denoted \(\Ghat_K(A)\), or \(\Ghat(A)\) when no confusion is possible. It is clearly an abelian group for multiplication of formal series, and this defines a functor
\[
  \Ghat_K:K\text{-Alg.fil.}\longrightarrow \mathrm{Ab}.
\]
If \(A\) has the gross filtration, then \(\Ghat(A)\) is the group of formal series with augmentation \(1\), which will also be denoted
\[
  1+A[[t]]^+.
\]

\sgaStableTarget{sga6.e05.lem.1-2}\textbf{Lemma 1.2.} The functor \(\Ghat_K\) is representable.

Let \((X_i)_{i\in\mathbb N,\ i\geq 1}\) be a family of indeterminates. Assign to \(X_i\) the weight \(i\), and filter the ring \(K[X_i]_{i\geq 1}\) by weights with respect to the \(X_i\). One then obtains a bijection
\[
  \Ghat_K(A)\xrightarrow{\sim}
  \Hom_{K\text{-Alg.fil.}}\bigl(K[X_i]_{i\geq 1},A\bigr)
\]
by associating to every series
\[
  \psi=1+a_1t+\cdots+a_nt^n+\cdots\in\Ghat_K(A)
\]
the morphism \(f_\psi\) determined by \(f_\psi(X_i)=a_i\). This homomorphism is compatible with the filtrations, and one thus obtains an isomorphism of functors.

Put
\[
  \varphi=1+X_1t+\cdots+X_nt^n+\cdots .
\]
For every filtered \(K\)-algebra \(A\) and every series \(\psi\in\Ghat(A)\), there exists a unique homomorphism
\[
  f_\psi\in K[X_i]_{i\geq 1}\longrightarrow A
\]
such that
\[
  \Ghat(f_\psi)(\varphi)=\psi .
\]

\sgaStableTarget{sga6.e05.para.1-3}\textbf{1.3.} For every \(A\in\operatorname{Ob}(\mathcal C)\) we define a filtration on \(\Ghat(A)\) as follows: a series
\[
  \psi=1+a_1t+\cdots
\]
is said to be of filtration \(\geq k\) if \(a_i=0\) for \(1\leq i<k\). Let \(F^k\Ghat(A)\) be the subgroup so obtained. This is a functor on \(\mathcal C\), and from now on we shall regard \(\Ghat\) as a functor with values in the category \(\Abfil\) of filtered abelian groups by means of this filtration, with \(F^0\Ghat=\Ghat\).

Let \(H\) be a functor on \(\mathcal C\) with values in \(\Abfil\). We shall suppose that \(H\) satisfies the following two conditions:

\[
\begin{array}{ll}
\text{i)} & \forall A\in\operatorname{Ob}(\mathcal C),\ H(A)\text{ is separated and complete},\\[0.3em]
\text{ii)} & \text{if }A\in\operatorname{Ob}(\mathcal C)\text{ is a graded }K\text{-algebra, with the filtration associated}\\
&\text{to its grading, and if a finite group }G\text{ acts on }A\text{ by graded morphisms, then}\\
& H(A^G)\simeq H(A)^G .
\end{array}
\]
Here \(H(A)^G\) denotes the invariants under \(G\).

In fact we shall need these conditions only when \(A\) is a polynomial ring and \(G\) is a finite product of symmetric groups acting by permutation of the indeterminates.

\sgaStableTarget{sga6.e05.ex.unnamed}\textbf{Examples.}
\[
\begin{array}{ll}
1)& H(A)=\Ghat(A),\\
\sgaStableTarget{sga6.e05.ex.1-3-2}2)& H(A)=1+A[[t]]^+.
\end{array}
\]
The latter functor will be denoted \(\Gg\), with \(g\) standing for ``gross''. Finally,
\[
3)\quad H(A)=\Ga(A),
\]
where \(\Ga(A)\) denotes the group, for addition, of formal series
\[
  a=a_0+a_1t+\cdots
\]
such that \(a_0=0\) and \(a_i\in\Fil^i A\) for \(i\geq 1\), endowed with the filtration analogous to that on \(\Ghat\).

\sgaStableTarget{sga6.e05.para.1-4}\textbf{1.4.} For every \(n\), denote by \(f_n\) the endomorphism of \(K[X_1,\ldots,X_n]\) defined by
\[
  f_n(X_i)=X_i\quad (i<n),\qquad f_n(X_n)=0.
\]
It factors through the inclusion of \(K[X_1,\ldots,X_{n-1}]\) in \(K[X_1,\ldots,X_n]\). Hence the homomorphism
\[
  H(K[X_1,\ldots,X_{n-1}])\longrightarrow H(K[X_1,\ldots,X_n])
\]
is injective. This will allow us to identify the first group with a subgroup of the second; similarly, we shall regard \(H(K[X_1,\ldots,X_n])\) as contained in \(H(K[X_i]_{i\geq 1})\).

One should note that \(f_1\) is none other than the augmentation of \(K[X]\).

\sgaStableTarget{sga6.e05.para.1-5}\textbf{1.5.} Embed \(K[X_1,\ldots,X_n]\) in \(K[T_1,\ldots,T_n]\) by sending \(X_i\) to the \(i\)-th elementary symmetric function in the \(T_j\). Let the symmetric group \(\Sym_n\) act by permutation of the \(T_j\). Using the theorem on symmetric polynomials and (1.3.ii), one sees that the induced homomorphism
\[
  H(K[X_1,\ldots,X_n])\longrightarrow H(K[T_1,\ldots,T_n])
\]
is injective, which again permits us to identify the former with a subgroup of the latter:
\[
  H(K[X_1,\ldots,X_n])\simeq H(K[T_1,\ldots,T_n])^{\Sym_n}
  \subset H(K[T_1,\ldots,T_n]).
\]
Let
\[
  u_i^n:K[X]\longrightarrow K[T_1,\ldots,T_n]
\]
be the homomorphism defined by \(u_i^n(X)=T_i\), and put
\[
  u_H^n=\sum_{i=1}^n H(u_i^n):H(K[X])\longrightarrow H(K[T_1,\ldots,T_n]).
\]
The symmetric group permutes the morphisms \(H(u_i^n)\), and so leaves \(u_H^n\) invariant; thus \(u_H^n\) takes its values in
\[
  H(K[T_1,\ldots,T_n])^{\Sym_n}
   \simeq H(K[X_1,\ldots,X_n])
\]
by (1.3.ii). Notice, finally, that if \(H=\Ghat\), then
\[
  u_H^n(1+Xt)=\prod_{i=1}^n(1+T_it)=1+X_1t+\cdots+X_nt^n.
\]

\sgaStableTarget{sga6.e05.thm.1-6}\textbf{Theorem 1.6.} The map from \(\Hom_{\Abfil}(\Ghat,H)\) to \(H(K[X])\) which, to a morphism of functors, associates its value on the element \(1+Xt\) of \(\Ghat(K[X])\), is injective. Its image is the subgroup of \(H(K[X])\) formed by the elements \(\theta\) such that, with the notation of \hyperlink{sga6.e05.para.1-4}{1.4} for the \(f_n\),
\[
\begin{array}{ll}
\text{i)}& H(f_1)(\theta)=0,\\[0.3em]
\text{ii)}& \forall n>1,\quad
 H(f_n)\bigl(u_H^n(\theta)\bigr)\equiv u_H^n(\theta)
 \pmod{F^nH(K[X_1,\ldots,X_n])}.
\end{array}
\]

First prove injectivity. By \hyperlink{sga6.e05.lem.1-2}{Lemma 1.2} it is enough to see that a morphism of functors \(\xi\) which is zero on \(1+Xt\) is zero on
\[
  \varphi=1+X_1t+\cdots+X_nt^n+\cdots .
\]
The commutative diagram
\sgaStableTarget{sga6.e05.eq.1-6-1}\[
\begin{tikzcd}
\Ghat(K[X]) \arrow[r,"u^n_{\Ghat}"] \arrow[d,"\xi"'] &
\Ghat(K[X_1,\ldots,X_n]) \arrow[d,"\xi"]\\
H(K[X]) \arrow[r,"u_H^n"'] &
H(K[X_1,\ldots,X_n])
\end{tikzcd}
\tag{1.6.1}
\]
shows that, for every \(n\),
\[
  \xi(1+X_1t+\cdots+X_nt^n)=0.
\]
But
\[
  \varphi\equiv 1+X_1t+\cdots+X_nt^n
  \pmod{F^{n+1}\Ghat},
\]
hence
\[
  \xi(\varphi)\equiv \xi(1+X_1t+\cdots+X_nt^n)
  \pmod{F^{n+1}H};
\]
therefore \(\xi(\varphi)\equiv 0\pmod{F^{n+1}H}\) for every \(n\). Since \(H(K[X_i]_{i\geq 1})\) is separated and complete, \(\xi(\varphi)=0\).

Now let \(\xi\) be any element of \(\Hom(\Ghat,H)\), and set
\[
  \theta=\xi(1+Xt).
\]
We show that \(\theta\) satisfies conditions i) and ii). First,
\[
  H(f_1)(\theta)=H(f_1)\xi(1+Xt)
  =\xi\,\Ghat(f_1)(1+Xt)=\xi(1)=0.
\]
Second,
\[
  u_H^n(\theta)=\xi(1+\cdots+X_nt^n)
\]
by \hyperlink{sga6.e05.eq.1-6-1}{(1.6.1)}, while
\[
  1+\cdots+X_nt^n\equiv 1+\cdots+X_{n-1}t^{\,n-1}
  \pmod{F^n\Ghat},
\]
and
\[
  1+\cdots+X_{n-1}t^{\,n-1}=\Ghat(f_n)(1+\cdots+X_nt^n).
\]
Since \(H\circ \Ghat(f_n)=H(f_n)\circ \xi\), and since \(\xi\) is a morphism of filtered functors, condition ii) follows.

It remains to show that if \(\theta\in H(K[X])\) satisfies i) and ii), then there exists \(\xi\in\Hom(\Ghat,H)\) with
\[
  \theta=\xi(1+Xt).
\]
First determine the value of \(\xi\) on the element \(1+\cdots+X_nt^n\) of \(\Ghat(K[X_1,\ldots,X_n])\). In view of \hyperlink{sga6.e05.eq.1-6-1}{(1.6.1)}, it must be, by definition,
\[
  \xi(1+\cdots+X_nt^n)
  =
  \sum_{i=1}^n \xi(1+T_it),
\]
where the \(\xi(1+T_it)\) are known by functoriality from the morphisms \(u_i^n\).

One then glues the \(\xi(1+\cdots+X_nt^n)\), for variable \(n\), by means of the commutative diagram
\sgaStableTarget{sga6.e05.eq.1-6-2}\[
\begin{tikzcd}
K[X_1,\ldots,X_n] \arrow[r,"f_n"] \arrow[d,"\sigma_n"'] &
K[X_1,\ldots,X_n] \arrow[d,"\sigma'_n"]\\
K[T_1,\ldots,T_n] \arrow[r,"f'_n"'] &
K[T_1,\ldots,T_n],
\end{tikzcd}
\tag{1.6.2}
\]
where \(f'_n\) is determined by \(f'_n(T_i)=T_i\) if \(i<n\), and \(f'_n(T_n)=0\); where \(\sigma_n\) sends \(X_i\) to the \(i\)-th elementary symmetric function in \(T_1,\ldots,T_n\); and where \(\sigma'_n\) sends \(X_i\) to the \(i\)-th elementary symmetric function in \(T_1,\ldots,T_{n-1}\) for \(i<n\), and sends \(X_n\) to \(0\).

Indeed one obtains from this
\[
\begin{aligned}
H(f_n)\bigl(\xi(1+\cdots+X_nt^n)\bigr)
&=H(f'_n)\left(\sum_{i=1}^n \xi(1+T_it)\right)\\
&=\sum_{i=1}^{n-1}\xi(1+T_it)
   =\xi(1+\cdots+X_{n-1}t^{\,n-1}),
\end{aligned}
\]
inside \(H(K[T_1,\ldots,T_{n-1}])\) by the identifications of 1.4. By 1.6.ii the elements \(\xi(1+\cdots+X_nt^n)\) define an element \(\psi\) of \(H(K[X_i]_{i\geq 1})\), since this is a separated complete group. Put
\[
  \xi(\varphi)=\psi .
\]

Finally, if \(A\in\operatorname{Ob}(\mathcal C)\) and \(x\in\Ghat(A)\), define
\[
  \xi_A(x)=H(f_x)(\psi),
\]
where \(f_x:K[X_i]_{i\geq 1}\to A\) is the homomorphism defined in 1.2. Since \(\xi\) has already been defined on the elements \(1+\cdots+X_nt^n=\varphi_n\), one must verify that the two definitions give the same result, i.e.
\[
  H(f_{\varphi_n})(\psi)=\xi(1+\cdots+X_nt^n).
\]
But for every \(p\) one has
\[
  \psi\equiv \xi(1+\cdots+X_pt^p)
  \pmod{F^{p+1}H(K[X_i]_{i\geq 1})}.
\]
Thus, for \(p\geq n\),
\[
  H(f_{\varphi_n})(\psi)
  \equiv
  H(f_{\varphi_n})\bigl(\xi(1+\cdots+X_pt^p)\bigr)
  \pmod{F^{p+1}H(K[X_i]_{i\geq 1})}.
\]
Using the identifications of \hyperlink{sga6.e05.para.1-4}{1.4} and the gluing of the \(\xi(\varphi_m)\) for variable \(m\), one obtains
\[
  H(f_{\varphi_n})(\xi(\varphi_p))=\xi(\varphi_n).
\]
Hence finally
\[
  \forall p\geq n,\quad
  H(f_{\varphi_n})(\psi)\equiv \xi(1+\cdots+X_nt^n)
  \pmod{F^pH(K[X_1,\ldots,X_n])}.
\]
Since \(H\) is separated, the required equality follows.

We have thus defined a system of maps
\[
  \xi_A:\Ghat(A)\longrightarrow H(A).
\]
To verify that they constitute a morphism of functors with values in groups, one must prove first that, for every morphism \(h:A\to A'\), the square
\[
\begin{tikzcd}
\Ghat(A) \arrow[r,"\Ghat(h)"] \arrow[d,"\xi_A"'] &
\Ghat(A') \arrow[d,"\xi_{A'}"]\\
H(A) \arrow[r,"H(h)"'] & H(A')
\end{tikzcd}
\]
is commutative. If \(x\in\Ghat(A)\), then \(h\circ f_x=f_{\Ghat(h)(x)}\), whence
\[
  H(h)(\xi_A(x))=H(h)H(f_x)(\psi)
  =H(f_{\Ghat(h)(x)})(\psi)=\xi_{A'}(\Ghat(h)(x)).
\]
It remains only to show that the morphism of functors so constructed is compatible with the structures of filtered groups.

Suppose \(x\in F^n\Ghat(A)\). To say that \(x\) is of filtration \(>n\) is to say that
\(f_x:K[X_i]_{i\geq 1}\to A\) factors through the endomorphism \(g_n\) of \(K[X_i]_{i\geq 1}\) defined by
\[
  g_n(X_i)=0\quad (1\leq i\leq n-1),\qquad
  g_n(X_i)=X_i\quad (i\geq n).
\]
To show that \(\xi\) is compatible with filtrations, it is enough to show that
\[
  H(g_n)(\xi(\varphi))\in F^nH.
\]
But by construction
\[
  \xi(\varphi)\equiv \xi(1+\cdots+X_{n-1}t^{\,n-1})
  \pmod{F^nH}.
\]
Therefore
\[
  H(g_n)(\xi(\varphi))
  =H(g_n)\bigl(\xi(1+\cdots+X_{n-1}t^{\,n-1})\bigr)
  \equiv 0\pmod{F^nH}.
\]

To verify the additivity of \(\xi\), it suffices to do this for the elements of finite degree of \(\Ghat(A)\), and then pass to the limit. Let \(X_1,X_2\) be two elements of degree \(n\) of \(\Ghat(A)\), and let \(X\) be their product. Consider the commutative diagram
\[
\begin{tikzcd}[column sep=tiny]
K[X_1,\ldots,X_{2n}] \arrow[rr,"f_X"] \arrow[d] && A\\
\bigl(K[T_1,\ldots,T_{2n}]\bigr)^{\Sym_n\times\Sym_n}
\arrow[urr,"f'_X"'] &&
\end{tikzcd}
\]
where the first factor of \(\Sym_n\times\Sym_n\) permutes \(T_1,\ldots,T_n\) and the second permutes \(T_{n+1},\ldots,T_{2n}\). The ring \(K[X_1,\ldots,X_{2n}]\) is embedded into the ring of invariants by sending \(X_j\) to the \(j\)-th elementary symmetric function of the \(T_i\); \(f'\) sends the \(i\)-th elementary symmetric function of \(T_1,\ldots,T_n\), respectively of \(T_{n+1},\ldots,T_{2n}\), to the \(i\)-th coefficient of \(X_1\), respectively \(X_2\). Then
\[
\begin{aligned}
\xi(X)
&=H(f'_X)\left(\xi\left(\prod_{i=1}^{2n}(1+T_it)\right)\right),\\
\xi(X_1)
&=H(f'_X)\left(\xi\left(\prod_{i=1}^{n}(1+T_it)\right)\right),
\end{aligned}
\]
and similarly for \(X_2\), using \(i=n+1,\ldots,2n\). Since \(H(f'_X)\) is a homomorphism of groups, we are reduced to the case
\[
  X=\prod_{i=1}^{2n}(1+T_it),
\]
and in this case the additivity of \(\xi\) follows from the definition of the \(\xi(1+X_1t+\cdots+X_k t^k)\).

\sgaStableTarget{sga6.e05.para.1-7}\textbf{1.7.} One can generalize the preceding considerations to the case of a product of functors \(\Ghat\). The functor \((\Ghat)^k\) is representable by the ring
\[
  K[X_i^{(j)}]_{j=1,\ldots,k;\ i\geq 1},
\]
where \(X_i^{(j)}\) has weight \(i\). We denote by \(\varphi^{(j)}\) the series
\[
  1+X_1^{(j)}t+\cdots+X_n^{(j)}t^n+\cdots.
\]
The element \((\varphi^{(j)})_j\) of
\[
  \bigl(\Ghat(K[X_i^{(j)}])\bigr)^k
\]
has the universal property analogous to that of \(\varphi\) (cf. \hyperlink{sga6.e05.lem.1-2}{1.2}).

One can introduce the morphisms \(f_n^{(j)}\) and \(u_H^{n,(j)}\) as above, by repeating the construction for each family \((X_i^{(j)})_{i\geq 1}\) separately. One obtains a result analogous to \hyperlink{sga6.e05.thm.1-6}{Theorem 1.6} by considering morphisms of functors from \((\Ghat)^k\) to \(H\) which are \(\mathbb Z\)-multilinear and compatible with the filtrations in each of the \(k\) arguments separately. They are determined by their value on
\[
  (1+X^{(j)}t)_{j=1,\ldots,k}
  \in \bigl(\Ghat(K[X^{(j)}]_j)\bigr)^k.
\]
An element
\[
  \theta\in H(K[X^{(j)}]_j)
\]
determines such a morphism of functors if and only if it satisfies
\[
\begin{array}{ll}
\text{i)} & \forall j,\quad H(f_1^{(j)})(\theta)=0,\\[0.3em]
\text{ii)}& \forall n,j,\quad
 H(f_n^{(j)})\bigl(u_H^{n,(j)}(\theta)\bigr)
 \equiv u_H^{n,(j)}(\theta)\pmod{F^nH}.
\end{array}
\]
The reader will easily convince himself that, apart from notational complications, it is enough to copy the proof of 1.6.

\sgaStableTarget{sga6.e05.para.1-8}\textbf{1.8.} Let \(\mathcal C_0\) be the category of \(K\)-algebras, and let
\[
  i:\mathcal C_0\longrightarrow\mathcal C
\]
be the functor which assigns to a \(K\)-algebra \(A\) the \(K\)-algebra \(A\) endowed with the gross filtration. Let
\[
  g:\mathcal C\longrightarrow\mathcal C_0
\]
be the forgetful functor which forgets the filtration. These functors satisfy the following relations:
\[
  g\circ i=\id_{\mathcal C_0},\qquad
  \Ghat\circ i\circ g=\Gg,\qquad
  \Gg\circ i\circ g=\Gg
\]
(cf. \hyperlink{sga6.e05.para.1-3}{1.3}, \hyperlink{sga6.e05.ex.1-3-2}{example 2)}.

\sgaStableTarget{sga6.e05.lem.1-9}\textbf{Lemma 1.9.} Let \(H:\mathcal C_0\to\Abfil\) be a functor on \(\mathcal C_0\). Then the canonical homomorphism
\[
  \Hom_{\Abfil}(\Ghat,H\circ g)
  \longrightarrow
  \Hom_{\Abfil}(\Ghat\circ i,H)
\]
is an isomorphism.

To see that it is injective, one uses the sequence of canonical homomorphisms
\[
 \Hom(\Ghat,H\circ g)\to
 \Hom(\Ghat\circ i,H)\to
 \Hom(\Gg,H\circ g)\to
 \Hom(\Ghat,H\circ g),
\]
the last arrow coming from the fact that \(\Ghat\) is a subfunctor of \(\Gg\). To verify that the composite is the identity, observe that a morphism of functors
\[
  \xi:\Ghat\longrightarrow H\circ g
\]
is determined by its values on objects of \(\mathcal C\) whose filtration is gross. Indeed, if \(A\in\operatorname{Ob}(\mathcal C)\), and if \(A_g\) is \(A\) endowed with the gross filtration, the canonical morphism \(A\to A_g\) defines a commutative diagram
\[
\begin{tikzcd}
\Ghat(A) \arrow[r,"\xi_A"] \arrow[d] &
(H\circ g)(A) \arrow[d]\\
\Ghat(A_g) \arrow[r,"\xi_{A_g}"'] &
(H\circ g)(A_g),
\end{tikzcd}
\]
which shows that \(\xi_A\) is the restriction of \(\xi_{A_g}\) to \(\Ghat(A)\).

Next consider the composite homomorphism
\[
  \Hom(\Ghat\circ i,H)\longrightarrow
  \Hom(\Gg,H\circ g)\longrightarrow
  \Hom(\Ghat,H\circ g)\longrightarrow
  \Hom(\Ghat\circ i,H).
\]
It is clear that this is the identity, whence the asserted result.

\sgaStableTarget{sga6.e05.thm.1-10}\textbf{Theorem 1.10.} Let \(H\) be a functor on the category \(\mathcal C_0\) of \(K\)-algebras, with values in the category of separated complete filtered abelian groups, and commuting with invariants under a finite group. Then the map from
\[
  \Hom_{\Abfil}(\Ghat\circ i,H)
\]
to \(H(K[X])\) which, to a morphism of functors, associates its value on the element \(1+Xt\) of \(\Ghat\circ i(K[X])\), is injective. Its image is the subgroup of \(H(K[X])\) formed by the elements \(\theta\) such that, with notation analogous to that of \hyperlink{sga6.e05.para.1-4}{1.4} and \hyperlink{sga6.e05.para.1-5}{1.5},
\[
\begin{array}{ll}
\text{i)}& H(f_1)(\theta)=0,\\[0.3em]
\text{ii)}& \forall n>1,\quad
 H(f_n)\bigl(u_H^n(\theta)\bigr)\equiv u_H^n(\theta)
 \pmod{F^nH(K[X_1,\ldots,X_n])}.
\end{array}
\]
The theorem follows immediately from \hyperlink{sga6.e05.thm.1-6}{Theorem 1.6} and \hyperlink{sga6.e05.lem.1-9}{Lemma 1.9}, since \(H\circ g(K[X])\) does not depend on the chosen filtration on \(K[X]\).

The analogous theorem obtained by replacing \(\Ghat\) by \((\Ghat)^k\) is also true (cf. \hyperlink{sga6.e05.para.1-7}{1.7}).

When there is no risk of confusion, we write the functor \(\Ghat\circ i\) simply as \(\Ghat\). This is legitimate because \(i\) identifies \(\mathcal C_0\) with a subcategory of \(\mathcal C\).

\sgaStableTarget{sga6.e05.para.1-11}\textbf{1.11.} Let \(\mathcal C_g\) denote the category of graded \(K\)-algebras, and let
\[
  j:\mathcal C_g\longrightarrow\mathcal C
\]
be the functor which assigns to a graded \(K\)-algebra \(B\) the \(K\)-algebra \(B\) endowed with the filtration canonically defined by its grading. The associated graded functor
\[
  \operatorname{Gr}:\mathcal C\longrightarrow \mathcal C_g
\]
will be denoted \(\operatorname{Gr}\). There is an isomorphism of functors
\[
  \varepsilon:\id_{\mathcal C_g}\xrightarrow{\sim}\operatorname{Gr}\circ j.
\]

If \(B\in\operatorname{Ob}(\mathcal C_g)\), denote by
\[
  1+\widehat{B}^{+}
\]
the multiplicative group of elements of
\[
  \prod_{i=0}^{\infty}B^i
\]
having augmentation \(1\). With the filtration analogous to that on \(\Ghat\), one obtains a functor
\[
  \mathcal C_g\longrightarrow\Abfil.
\]
Furthermore, for every \(A\in\operatorname{Ob}(\mathcal C)\), one defines a homomorphism of filtered groups
\[
  \Ghat(A)\longrightarrow 1+\widehat{\operatorname{Gr}A}^{+}
\]
by sending the series \(1+\cdots+a_nt^n+\cdots\) to the series \(1+\cdots+\bar a_n+\cdots\), where \(\bar a_n\) is the class of \(a_n\) in \(\operatorname{Gr}^nA\). Its kernel is the group \(\Ghat'(A)\) of formal series such that
\[
  a_0=1,\qquad \forall i\geq 1,\quad a_i\in\Fil^{i+1}A.
\]
Thus one has a functorial isomorphism in \(A\):
\[
  \Ghat(A)/\Ghat'(A)\xrightarrow{\sim} 1+\widehat{\operatorname{Gr}A}^{+}.
\]

\sgaStableTarget{sga6.e05.lem.1-12}\textbf{Lemma 1.12.} Let \(H:\mathcal C_g\to\Abfil\) be a functor on \(\mathcal C_g\). Then the canonical homomorphism
\[
  \rho:\Hom_{\Abfil}(1+\widehat{\,\cdot\,}^{+},H)
  \longrightarrow
  \Hom_{\Abfil}(\Ghat,H\circ\operatorname{Gr})
\]
is an isomorphism.

The homomorphism \(\rho\) is defined as the composite
\[
 \Hom(1+\widehat{\,\cdot\,}^{+},H)
 \longrightarrow
 \Hom(1+\widehat{\operatorname{Gr}(\,\cdot\,)}^{+},H\circ\operatorname{Gr})
 \longrightarrow
 \Hom(\Ghat,H\circ\operatorname{Gr}),
\]
the last arrow being defined by the homomorphism
\[
  \Ghat\longrightarrow 1+\widehat{\operatorname{Gr}(\,\cdot\,)}^{+}.
\]
We construct an inverse morphism
\[
  \rho':\Hom(\Ghat,H\circ\operatorname{Gr})
  \longrightarrow
  \Hom(1+\widehat{\,\cdot\,}^{+},H).
\]
For this, we show that every morphism of functors from \(\Ghat\) to \(H\circ\operatorname{Gr}\) factors through the canonical morphism
\[
  \Ghat\longrightarrow 1+\widehat{\operatorname{Gr}(\,\cdot\,)}^{+}.
\]
Equivalently, it suffices to show that it vanishes on \(\Ghat'\), by 1.11. Let \(A\in\operatorname{Ob}(\mathcal C)\) and let \(x\in\Ghat'(A)\). Define a homomorphism of filtered \(K\)-algebras
\[
  K[X_i]_{i\geq 1}\longrightarrow A
\]
by sending \(X_1\) to \(0\) and \(X_i\) to the \((i-1)\)-st coefficient of \(x\) for \(i>1\). Using this homomorphism, it is enough to verify that the morphism of functors considered vanishes on the element
\[
  1+X_2t+\cdots+X_{n+1}t^n+\cdots .
\]
Use the endomorphism \(f\) of \(K[X_i]_{i\geq 1}\) defined by
\[
  f(X_i)=X_{i+1}\quad\text{for all }i,
\]
and the commutative diagram
\[
\begin{tikzcd}
\Ghat(K[X_i]) \arrow[r] \arrow[d,"\Ghat(f)"'] &
(H\circ\operatorname{Gr})(K[X_i]) \arrow[d,"H(\operatorname{Gr}(f))"]\\
\Ghat(K[X_i]) \arrow[r] &
(H\circ\operatorname{Gr})(K[X_i]).
\end{tikzcd}
\]
Now \(\operatorname{Gr}(f)=\operatorname{Gr}(f_1)\), where \(f_1\) is the augmentation homomorphism of \(K[X_i]_{i\geq 1}\). Since the latter annihilates \(\varphi\), the analogous diagram written for \(f_1\) shows that \(H(\operatorname{Gr}(f))\) annihilates the image of \(\varphi\) in \(H\circ\operatorname{Gr}(K[X_i]_{i\geq 1})\), and the result follows.

Passing to the quotient, one associates to every morphism of functors
\[
  \Ghat\longrightarrow H\circ\operatorname{Gr}
\]
a morphism of functors
\[
  1+\widehat{\operatorname{Gr}}^{+}\longrightarrow H\circ\operatorname{Gr}.
\]
By composing with \(j\) and using the isomorphism \(\varepsilon\), one obtains a morphism
\[
  1+\widehat{\,\cdot\,}^{+}\longrightarrow H,
\]
which defines the desired homomorphism \(\rho'\). It is immediate to verify that \(\rho\) and \(\rho'\) are inverses of one another.

\sgaStableTarget{sga6.e05.thm.1-13}\textbf{Theorem 1.13.} Let \(H\) be a functor on \(\mathcal C_g\), with values in the category of separated complete filtered abelian groups, and commuting with invariants under a finite group. Then the map from
\[
  \Hom_{\Abfil}(1+\widehat{\,\cdot\,}^{+},H)
\]
to \(H(K[X])\), which to a morphism of functors associates its value on the element \(1+X\) of \(1+\widehat{K[X]}^{+}\), is injective. Its image is the subgroup of \(H(K[X])\) formed by the elements \(\theta\) such that, with notation as in \hyperlink{sga6.e05.para.1-4}{1.4} and \hyperlink{sga6.e05.para.1-5}{1.5},
\[
\begin{array}{ll}
\text{i)}&H(f_1)(\theta)=0,\\[0.3em]
\text{ii)}&\forall n>1,\quad
H(f_n)\bigl(u_H^n(\theta)\bigr)\equiv u_H^n(\theta)
 \pmod{F^nH(K[X_1,\ldots,X_n])}.
\end{array}
\]
This is immediate from \hyperlink{sga6.e05.thm.1-6}{1.6} and 1.12. There is an analogous statement replacing \(\Ghat\) by \((\Ghat)^k\).

\sgaStableTarget{sga6.e05.rem.1-14}\textbf{Remark 1.14.} If \(H\) is a functor on \(\mathcal C_g\), with values in abelian groups, commuting with invariants under a finite group, and satisfying the condition
\[
  \forall B\in\operatorname{Ob}(\mathcal C_g),\qquad
  H(B)\simeq \varprojlim_n H(B/B^{\geq n}),
  \qquad B^{\geq n}=\prod_{i\geq n}B^i,
\]
then condition ii) of \hyperlink{sga6.e05.thm.1-13}{1.13} is always satisfied if one takes for the filtration on \(H\) the filtration by the kernels of
\[
  H(B)\longrightarrow H(B/B^{\geq n}),
\]
which is separated and complete. Indeed, if \(p_n\) is the canonical homomorphism
\[
  K[X_1,\ldots,X_n]\longrightarrow K[X_1,\ldots,X_n]/K[X_1,\ldots,X_n]^{\geq n},
\]
then
\[
  p_n\circ f_n=p_n,
\]
and hence
\[
  H(p_n)\bigl[H(f_n)\circ u_H^n(\theta)-u_H^n(\theta)\bigr]=0.
\]
This says precisely that
\[
  H(f_n)\circ u_H^n(\theta)-u_H^n(\theta)
  \in F^nH(K[X_1,\ldots,X_n]).
\]
Thus, in order to give a morphism of functors from \(1+\widehat{\,\cdot\,}^{+}\) to \(H\), it is enough to give an element of \(H(K[X])\) satisfying only condition i) of 1.13.

\sgaStableTarget{sga6.e05.para.1-15}\textbf{1.15.} Let \(H:\mathcal C\to\Abfil\) (respectively \(\mathcal C_0\to\Abfil\), \(\mathcal C_g\to\Abfil\)) be a functor satisfying the conditions of \hyperlink{sga6.e05.para.1-3}{1.3} (respectively \hyperlink{sga6.e05.thm.1-10}{1.10}, \hyperlink{sga6.e05.thm.1-13}{1.13}), and let
\[
  \xi:(\Ghat)^k\to H
\]
(respectively \((\Ghat\circ i)^k\to H\), \((1+\widehat{\,\cdot\,}^{+})^k\to H\)) be a morphism of functors. We have seen in \hyperlink{sga6.e05.para.1-7}{1.7} that \(\xi\) is determined by its value on the element
\[
  (\varphi^{(j)})_{j=1,\ldots,k}
\]
of
\[
  \bigl(\Ghat(K[X_i^{(j)}]_{j=1,\ldots,k;\ i\geq 1})\bigr)^k
\]
(respectively similarly). Suppose \(H\) is one of the functors
\[
  \Ghat,\quad \Gg,\quad \Ga,\quad 1+\widehat{\operatorname{Gr}}^{+},\quad\text{etc.}
\]
Then the element \(\xi((\varphi^{(j)}))\) is a formal series with coefficients in
\[
  K[X_i^{(j)}]_{i\geq 1,\ j=1,\ldots,k}.
\]
The \(n\)-th coefficient of this series is called the \(n\)-th universal polynomial of \(\xi\). The knowledge of these polynomials is equivalent to the knowledge of \(\xi\).

According to the proof of \hyperlink{sga6.e05.thm.1-6}{1.6}, these polynomials are themselves determined by the element
\[
  \xi(1+X^{(1)}t,\ldots,1+X^{(k)}t).
\]
The calculation is performed by the method of the proof: one computes the elements
\[
  \xi(1+X_1^{(1)}t+\cdots+X_n^{(1)}t^n,\ldots,
        1+X_1^{(k)}t+\cdots+X_n^{(k)}t^n),
\]
by decomposing
\[
  1+X_1^{(1)}t+\cdots+X_n^{(1)}t^n
  =\prod_{j=1}^n(1+T_j^{(1)}t)
\]
and by using the additivity of \(\xi\). The polynomials thus obtained are symmetric with respect to the \(T_j^{(i)}\), for fixed \(i\) and variable \(j\), and can be expressed as functions of the \(X_j^{(i)}\). Gluing by means of the filtration of \(H\), one thus determines the universal polynomials of \(\xi\). If
\[
  (\psi^{(1)},\ldots,\psi^{(k)})\in(\Ghat(A))^k,
\]
the calculation of \(\xi(\psi^{(1)},\ldots,\psi^{(k)})\) is made by substituting the coefficients of the \(\psi^{(i)}\) for the \(X_j^{(i)}\) in the universal polynomials.

Suppose that there exists an integer \(m\) such that the coefficient of \(t^n\) in
\[
  \theta=\xi(1+X^{(1)}t,\ldots,1+X^{(k)}t)
\]
is a polynomial whose isobaric components have weights between \(n\) and \(m\). Then the \(k\)-th universal polynomial of \(\xi\) has isobaric components whose weights range between \(k\) and \(km\). This follows immediately from the calculation just described and from the theorem on symmetric functions.

In particular, if \(H=\Ghat\), then the \(n\)-th universal polynomial of \(\xi\) has no component of weight \(<n\). Now consider the canonical morphism
\[
  H\longrightarrow 1+\widehat{\operatorname{Gr}}^{+}.
\]
Composing it with \(\xi\) gives a morphism
\[
  (\Ghat)^k\longrightarrow 1+\widehat{\operatorname{Gr}}^{+},
\]
which, by the proof of \hyperlink{sga6.e05.lem.1-12}{1.12}, factors through
\[
  (1+\widehat{\operatorname{Gr}}^{+})^k
\]
and defines a morphism
\[
  \bar\xi:(1+\widehat{\,\cdot\,}^{+})^k\longrightarrow 1+\widehat{\,\cdot\,}^{+}.
\]
One sees at once that the \(n\)-th universal polynomial of \(\bar\xi\) is the same as that of
\[
  (\Ghat)^k\longrightarrow 1+\widehat{\operatorname{Gr}}^{+},
\]
that is, the isobaric component of weight \(n\) of the \(n\)-th universal polynomial of \(\xi\).

The same remarks apply to the case \(H=\Ga\). Then one may define a morphism of functors
\[
  H\longrightarrow \widehat{\operatorname{Gr}}^{+}
\]
analogously to \hyperlink{sga6.e05.para.1-11}{1.11}, by putting
\[
  \widehat B^{+}=\prod_{i\geq 1}B^i
\]
for a graded algebra \(B\); one repeats the same reasoning and reaches the conclusions of the preceding paragraph.

\sgaStableTarget{sga6.e05.para.1-16}\textbf{1.16.} We finish with an example. Suppose \(K=\mathbb Q\). Define an endomorphism of the functor \(1+\widehat{\,\cdot\,}^{+}\), denoted \(\Toddc\) and called the Todd character, by putting
\[
  \Toddc(1+X)=\frac{X}{1-e^{-X}}.
\]
The conditions i) and ii) of \hyperlink{sga6.e05.thm.1-13}{Theorem 1.13} are immediately verified, which shows that an endomorphism of \(1+\widehat{\,\cdot\,}^{+}\) has indeed been defined.

\sgaStableTarget{sga6.e05.sec.2}\section*{2. Definition of \(\lambda\)-rings; examples}

In this paragraph the base ring of paragraph 1 will be \(\mathbb Z\), except in 2.9. We place ourselves in the category \(\mathcal C_0\) of \hyperlink{sga6.e05.para.1-8}{1.8}, and write \(\Ghat\) for the restriction of \(\Ghat\) to \(\mathcal C_0\), as in 1.10.

\sgaStableTarget{sga6.e05.def.2-1}\textbf{Definition 2.1.} A \emph{pre-\(\lambda\)-ring} is a commutative ring \(K\) with unit, endowed with a homomorphism of abelian groups
\[
  \lambda_t:K\longrightarrow 1+K[[t]]^+
\]
such that, for every \(x\in K\),
\[
  \lambda_t(x)=1+xt+\cdots.
\]
We shall write
\[
  \lambda_t(x)=\sum_{i\geq 0}\lambda^i(x)t^i.
\]
Thus giving a pre-\(\lambda\)-ring structure on a ring \(K\) is equivalent to giving a family of maps
\[
  \lambda^i:K\longrightarrow K\qquad (i\in\mathbb N)
\]
such that
\[
\begin{array}{ll}
\text{i)}&\forall x\in K,\quad \lambda^0(x)=1,\\[0.3em]
\text{ii)}&\forall x\in K,\quad \lambda^1(x)=x,\\[0.3em]
\text{iii)}&\forall x,y\in K,\ \forall n\in\mathbb N,\quad
\lambda^n(x+y)=\displaystyle\sum_{p+q=n}\lambda^p(x)\lambda^q(y).
\end{array}
\]
If \(K,K'\) are two pre-\(\lambda\)-rings and \(f:K\to K'\) is a ring homomorphism, we say that \(f\) is a homomorphism of pre-\(\lambda\)-rings, or a \(\lambda\)-homomorphism, if the diagram
\[
\begin{tikzcd}
K \arrow[r,"\lambda_t"] \arrow[d,"f"'] &
1+K[[t]]^+ \arrow[d,"\Ghat(f)"]\\
K' \arrow[r,"\lambda_t"'] &
1+K'[[t]]^+
\end{tikzcd}
\]
is commutative. This is equivalent to the relations
\[
  \forall x\in K,\ \forall i\in\mathbb N,\qquad
  \lambda^i(f(x))=f(\lambda^i(x)).
\]

\sgaStableTarget{sga6.e05.ex.2-2}\textbf{Examples 2.2.}
\begin{enumerate}[label=\alph*)]
\item One can give all pre-\(\lambda\)-ring structures on \(\mathbb Z\). It is enough to give \(\lambda_t(1)\), which is then a formal series with integral coefficients subject to the sole condition that it begin with \(1+t\).

\item Let \((T,A)\) be a commutatively ringed topos, and let \(K\) be the ``Grothendieck group'' of \(A\)-modules locally isomorphic to a direct factor of a free \(A\)-module of finite type. We know that \(K\) can be given a ring structure by the tensor product (\hyperlink{sga6.e04.expose}{Exposé IV}), and we now endow it with a pre-\(\lambda\)-ring structure as follows. For every \(A\)-module \(E\) locally a direct factor of a free module of finite type, let \(\lambda^i(E)\) denote the class in \(K\) of the \(i\)-th exterior power of \(E\), which is again locally a direct factor of a free module of finite type. This defines a function from the category of such \(A\)-modules to the abelian group \(1+K[[t]]^+\). It remains to show that this function is additive, which will imply that it factors through \(K\) and hence gives the desired structure.
\end{enumerate}

\sgaStableTarget{sga6.e05.lem.2-2-1}\textbf{Lemma 2.2.1.} Let
\[
  0\longrightarrow E' \xrightarrow{u} E \xrightarrow{v} E''\longrightarrow 0
\]
be an exact sequence of \(A\)-modules which is locally split, for example with \(E''\) locally a direct factor of a free module of finite type. There exists on \(\Lambda^n(E)\) a finite filtration and isomorphisms
\[
  \Lambda^{n-p}(E')\otimes_A\Lambda^p(E'')
  \xrightarrow{\sim}
  \operatorname{Gr}^p\Lambda^n(E).
\]
Take for \(F^p\Lambda^n(E)\) the image of the canonical morphism
\[
  \Lambda^{n-p}(E')\otimes_A\Lambda^p(E)\longrightarrow \Lambda^n(E)
  \quad (0\leq p\leq n).
\]
If \(p>n\), put \(F^p\Lambda^n(E)=\Lambda^n(E)\), and if \(p<0\), put \(F^p\Lambda^n(E)=0\). This does indeed define a finite filtration on \(\Lambda^n(E)\).

Define a homomorphism
\[
  f_p:\Lambda^{n-p}(E')\otimes_A\Lambda^p(E'')
  \longrightarrow
  F^p\Lambda^n(E)/F^{p-1}\Lambda^n(E)
\]
by
\[
  f_p(x'_1\wedge\cdots\wedge x'_{n-p}\otimes x''_1\wedge\cdots\wedge x''_p)
  =
  \operatorname{cl}\bigl(u(x'_1)\wedge\cdots\wedge u(x'_{n-p})\wedge
  x_1\wedge\cdots\wedge x_p\bigr),
\]
where the \(x_i\) are lifts in \(E\) of the \(x''_i\) in \(E''\). This visibly does not depend on the choice of the \(x_i\). If the exact sequence is locally split, then \(f_p\) is an isomorphism, since locally
\[
  \Lambda^n(E)\simeq\bigoplus_{p+q=n}
  \Lambda^p(E')\otimes_A\Lambda^q(E'').
\]

When \(E',E,E''\) are locally direct factors of free modules of finite type, the \(\operatorname{Gr}^p\Lambda^n(E)\) are locally direct factors of free modules of finite type, and hence so are the \(F^p\Lambda^n(E)\). Thus one may write in \(K\)
\[
  \lambda^n(E)=\sum_{p=0}^{n}\operatorname{cl}(\operatorname{Gr}^p\Lambda^n(E))
  =
  \sum_{p=0}^{n}\lambda^{n-p}(E')\lambda^p(E''),
\]
which gives the required additivity.

Let \((T,A)\) and \((T',A')\) be two commutatively ringed topoi, and let
\[
  f:(T,A)\longrightarrow (T',A')
\]
be a morphism of ringed topoi. We know (\hyperlink{sga6.e04.para.2-7}{Exposé IV, 2.7)} that \(f\) defines a morphism of rings on the corresponding \(K\)'s, namely
\[
  f^*:K^*_{\mathrm{naif}}(T',A')\longrightarrow K^*_{\mathrm{naif}}(T,A).
\]
Since inverse image commutes with exterior powers, the corresponding morphism \(f^*\) on \(K\) is a \(\lambda\)-homomorphism for the structure just defined.

We point out two special cases:
\begin{itemize}
\item If \(X\) is a scheme, then \(K^*_{\mathrm{naif}}(X)\) is thus endowed with a pre-\(\lambda\)-ring structure. This is the case which will be used later in the Seminar, and it may be regarded as the model for the general theory of \(\lambda\)-rings.
\item If \(A\) is a ring, \(X=\operatorname{Spec}A\), and \(G\) is a group acting trivially on \(A\), then, taking for \((T,A)\) the topos of sheaves of sets on \(X\) on which \(G\) acts, the ring \(K\) thus obtained is the representation ring of \(G\) in the projective \(A\)-modules of finite type.
\end{itemize}

\sgaStableTarget{sga6.e05.para.2-3}\textbf{2.3.} Let \(K\) be a commutative ring with unit. We propose to endow the abelian group
\[
  \Ghat(K)=1+K[[t]]^+
\]
with a pre-\(\lambda\)-ring structure, functorially in \(K\). To this end, we apply the results of \hyperlink{sga6.e05.thm.1-10}{1.10} to the functor \(\Ghat\).

First define multiplication by giving a \(\mathbb Z\)-bilinear morphism of functors
\[
  \Ghat\times\Ghat\longrightarrow\Ghat
\]
through the element \(1+XYt\) of \(\Ghat(\mathbb Z[X,Y])\). This element satisfies 1.7.i; to see that it satisfies 1.7.ii, one must show that in
\[
  \prod_{i=1}^{n}(1+T_iYt),
\]
when expressed by means of the elementary symmetric functions \(X_j\) in the \(T_i\), the variable \(X_n\) does not occur in the coefficients of \(t^k\) for \(k<n\). This is evident by degree considerations.

The law so defined will be denoted by the sign \(\circ\). Thus, by definition,
\sgaStableTarget{sga6.e05.eq.2-3-1}\begin{equation}\tag{2.3.1}
  (1+Xt)\circ(1+Yt)=1+XYt .
\end{equation}
It is distributive by construction, and it is associative and commutative by the diagrams
\[
\begin{tikzcd}[column sep=tiny]
\Ghat\times\Ghat\times\Ghat \arrow[r] \arrow[d] &
\Ghat\times\Ghat \arrow[d]\\
\Ghat\times\Ghat \arrow[r] & \Ghat
\end{tikzcd}
\qquad
\begin{tikzcd}[column sep=tiny]
\Ghat\times\Ghat \arrow[r] \arrow[d] &
\Ghat\times\Ghat \arrow[d]\\
\Ghat \arrow[r,"\id"'] & \Ghat,
\end{tikzcd}
\]
where the meaning of the arrows is evident, and where it is enough to verify commutativity on the elements \(1+Xt,1+Yt,1+Zt\), respectively \(1+Xt,1+Xt\), of \(\Ghat(\mathbb Z[X,Y,Z])\), respectively \(\Ghat(\mathbb Z[X,Y])\), by \hyperlink{sga6.e05.thm.1-10}{1.10}; this verification is immediate.

By \hyperlink{sga6.e05.para.1-8}{1.8}, the product under this law of two series is calculated by substituting their coefficients into the universal polynomials of the law. The \(n\)-th universal polynomial will be denoted \(P_n\). It is a polynomial in
\[
  \mathbb Z[X_i,Y_j]_{1\leq i,j\leq n}
\]
which is isobaric of weight \(n\) with respect to each of the two groups of variables (see \hyperlink{sga6.e05.para.1-15}{1.15}), and symmetric with respect to these two groups. For example,
\[
  P_2(X_1,X_2;Y_1,Y_2)=X_2Y_1^2+Y_2X_1^2-2X_2Y_2 .
\]
One also has the following relation: if \(K\in\operatorname{Ob}(\mathcal C_0)\), \(\psi\in\Ghat(K)\), and \(a\in K\), then
\sgaStableTarget{sga6.e05.eq.2-3-2}\begin{equation}\tag{2.3.2}
  \psi(t)\circ(1+at)=\psi(at).
\end{equation}
As usual, by the standard functorial argument this reduces to the case \(K=\mathbb Z[X,Y]\), \(\psi=1+Xt\), \(a=Y\), where it is nothing other than \hyperlink{sga6.e05.eq.2-3-1}{(2.3.1)}.

Applying \hyperlink{sga6.e05.eq.2-3-2}{(2.3.2)} to \(\Ghat(K)=1+K[[t]]^+\) with \(a=1\), one sees that \(1+t\) is the unit for this multiplication. To avoid confusion, we sometimes denote it by \(1\), or by \(1_{\circ}\).

The functor \(\Ghat\) can now be considered as a functor from \(\mathcal C_0\) to the category of commutative rings with unit, and hence can be composed with itself. To avoid confusion, choose different indeterminates for each of the two factors in the product of \(\Ghat\) with itself; we indicate the indeterminate by an exponent when necessary:
\[
  \Ghat^t,\quad \Ghat^u,\quad\text{etc.}
\]
For \(A\in\operatorname{Ob}(\mathcal C_0)\) one has
\[
  \Ghat^u\circ \Ghat^t(A)=1+\{1+A[[t]]^+\}[[u]]^+.
\]
To define the operations \(\lambda^i\) on \(1+K[[t]]^+\), we are therefore going to give a morphism of functors
\[
  \Ghat\longrightarrow \Ghat^u\circ\Ghat^t.
\]
For this, one must first define a filtration of the functor \(\Ghat^u\circ\Ghat^t\) which permits the gluing of 1.6.

First note that \(F^k\Ghat\) is an ideal for the law \(\circ\). For this follows, for the product
\[
  \left(1+\sum_{i\geq k}X_i t^i\right)\circ(1+Yt),
\]
by reducing by functoriality to the case \(A=\mathbb Z[X_i,Y]_{i\geq 1}\), and using \hyperlink{sga6.e05.eq.2-3-2}{(2.3.2)}. Therefore \(F^n(\Ghat^u\circ\Ghat^t)\) is defined to be the set of series of the form
\[
  1+\psi_1u+\cdots+\psi_iu^i+\cdots,\qquad
  \psi_i\in F^{[n/i]}\Ghat,
\]
where \([n/i]\) denotes the integral part of \(n/i\).

The element defining the morphism \(\lambda_u\) is
\[
  1+\{1+Xt\}u.
\]
The reader will verify as an exercise that the conditions of \hyperlink{sga6.e05.thm.1-6}{1.6} are indeed fulfilled; in fact the filtration was chosen ad hoc.

Again, the operations \(\lambda^i\) are calculated by substitution from
\[
  \lambda_u(1+\cdots+X_nt^n+\cdots),
\]
and so are determined by universal polynomials. Denote by
\[
  Q_{i,j}(X_1,\ldots,X_{ij})
\]
the coefficient of \(t^j\) in
\[
  \lambda^i(1+\cdots+X_nt^n+\cdots).
\]
It is an isobaric polynomial of weight \(ij\) with respect to the variables \(X_1,\ldots,X_{ij}\) (by a verification analogous to that of \hyperlink{sga6.e05.para.1-15}{1.15}). It is calculated by symmetric functions from the defining element
\sgaStableTarget{sga6.e05.eq.2-3-3}\begin{equation}\tag{2.3.3}
  \lambda_u(1+Xt)=1+\{1+Xt\}u .
\end{equation}
For example,
\[
  Q_{2,2}(X_1,X_2,X_3,X_4)=X_1X_3-X_4.
\]

Finally, we indicate the origin of formulas \hyperlink{sga6.e05.eq.2-3-1}{(2.3.1)} and \hyperlink{sga6.e05.eq.2-3-3}{(2.3.3)}. Consider the \hyperlink{sga6.e05.ex.2-2}{example 2.2}.b, where the operations \(\lambda^i\) are defined by exterior powers of \(A\)-modules. We want
\[
  \lambda_t:K\longrightarrow \Ghat(K)
\]
to be a \(\lambda\)-homomorphism. Thus on \(\Ghat\) one must define laws such that, for all \(A\)-modules \(E,F\),
\[
  \lambda_t(EF)=\lambda_t(E)\circ\lambda_t(F),\qquad
  \lambda_u(\lambda_t(E))=\Ghat^u_{\lambda_t}\bigl(\lambda_t(\lambda_u(E))\bigr).
\]
The relations \hyperlink{sga6.e05.eq.2-3-1}{(2.3.1)} and \hyperlink{sga6.e05.eq.2-3-3}{(2.3.3)} are precisely the preceding relations applied to the case where \(E\) and \(F\) are invertible modules.

We shall see in \hyperlink{sga6.e06.expose}{Exposé VI} to what extent the desired property is verified.

\sgaStableTarget{sga6.e05.def.2-4}\textbf{Definition 2.4.} A pre-\(\lambda\)-ring \(K\) will be called a \emph{\(\lambda\)-ring} if the homomorphism of abelian groups
\[
  \lambda_t:K\longrightarrow 1+K[[t]]^+
\]
is a \(\lambda\)-homomorphism for the pre-\(\lambda\)-ring structure on the second member defined in 2.3. Thus the following relations hold for all \(x,y\in K\):
\sgaStableTarget{sga6.e05.eq.2-4-1}\begin{align}
  \lambda_t(1)&=1+t,&
  \lambda_t(xy)&=\lambda_t(x)\circ\lambda_t(y),&
  \lambda_u(\lambda_t(x))&=\Ghat(\lambda_t)\bigl(\lambda_u(x)\bigr).
  \tag{2.4.1}
\end{align}
In terms of the maps \(\lambda^i\), the last two relations are equivalent, for \(x,y\in K\) and \(i,j\in\mathbb N\), to
\sgaStableTarget{sga6.e05.eq.2-4-2}\sgaStableTarget{sga6.e05.eq.2-4-3}\begin{align}
  \lambda^i(xy)
  &=P_i(\lambda^1(x),\ldots,\lambda^i(x);
       \lambda^1(y),\ldots,\lambda^i(y)), \tag{2.4.2}\\
  \lambda^j(\lambda^i(x))
  &=Q_{i,j}(\lambda^1(x),\ldots,\lambda^{ij}(x)). \tag{2.4.3}
\end{align}

We note that here we depart from the usual terminology (cf. [1]), where what we call a pre-\(\lambda\)-ring was called a \(\lambda\)-ring, and our \(\lambda\)-rings were called ``special'' \(\lambda\)-rings. The reason is that only ``special'' \(\lambda\)-rings are useful in practice.

\sgaStableTarget{sga6.e05.ex.2-5}\textbf{Examples 2.5.}
\begin{enumerate}[label=\alph*)]
\item We have determined all pre-\(\lambda\)-ring structures on \(\mathbb Z\). It follows from \hyperlink{sga6.e05.eq.2-4-1}{(2.4.1)} that only the structure defined by
\[
  \lambda_t(1)=1+t
\]
is capable of being a \(\lambda\)-ring structure. One then has
\[
  \lambda_t(n)=(1+t)^n,
\]
hence
\[
  \forall n,n'\in\mathbb Z,\quad
  \lambda_t(n)\circ\lambda_t(n')=(1+t)^n\circ(1+t)^{n'}
  =(1+t)^{nn'}=\lambda_t(nn'),
\]
and
\[
  \lambda_u((1+t)^n)
  =
  (\lambda_u(1+t))^n
  =
  (1+\{1+t\}u)^n
  =
  \Ghat(\lambda_t)((1+u)^n)
  =
  \Ghat(\lambda_t)(\lambda_u(n)).
\]
Thus this is indeed a \(\lambda\)-ring structure. If \(\binom{X}{k}\) denotes the polynomial with rational coefficients
\[
  \frac{1}{k!}X(X-1)\cdots(X-k+1),
\]
then
\sgaStableTarget{sga6.e05.eq.2-5-1}\begin{equation}\tag{2.5.1}
  \lambda^i(n)=\binom{n}{i}.
\end{equation}

\item For every ring \(K\), the pre-\(\lambda\)-ring \(1+K[[t]]^+\) is a \(\lambda\)-ring. First,
\[
  \lambda_u(1)=\lambda_u(1+t)=1+\{1+t\}u=1+1\cdot u.
\]
The other relations are verified by the commutativity of the diagrams
\[
\begin{tikzcd}[column sep=tiny]
\Ghat^t\times\Ghat^t \arrow[r,"\lambda_u\times\lambda_u"] \arrow[d,"\circ"'] &
(\Ghat^u\circ\Ghat^t)\times(\Ghat^u\circ\Ghat^t) \arrow[d,"\circ"]\\
\Ghat^t \arrow[r,"\lambda_u"'] &
\Ghat^u\circ\Ghat^t,
\end{tikzcd}
\qquad
\begin{tikzcd}[column sep=tiny]
\Ghat^t \arrow[r,"\lambda_v"] \arrow[d,"\lambda_u"'] &
\Ghat^v\circ\Ghat^t \arrow[d,"\Ghat^v(\lambda_u)"]\\
\Ghat^u\circ\Ghat^t \arrow[r,"\lambda_v"'] &
\Ghat^v\circ\Ghat^u\circ\Ghat^t,
\end{tikzcd}
\]
which are checked by \hyperlink{sga6.e05.thm.1-10}{1.10} on the elements \(1+Xt\) and \(1+Yt\). For the first,
\[
\begin{aligned}
\lambda_u(1+Xt)\circ\lambda_u(1+Yt)
&=(1+\{Xt\}u)\circ(1+\{1+Yt\}u)\\
&=1+\{(1+Xt)\circ(1+Yt)\}u\\
&=1+\{1+XYt\}u\\
&=\lambda_u((1+Xt)\circ(1+Yt)).
\end{aligned}
\]
For the second,
\[
  \lambda_v(\lambda_u(1+Xt))
  =
  1+\{1+\{1+Xt\}u\}v
  =
  \Ghat^v(\lambda_u)\lambda_v(1+Xt).
\]

\item The geometric examples of \hyperlink{sga6.e05.ex.2-2}{2.2} will be treated later.
\end{enumerate}

\sgaStableTarget{sga6.e05.para.2-6}\textbf{2.6.} One can generalize \hyperlink{sga6.e05.ex.2-5}{Example 2.5}.a as follows. Let \(K\) be a commutative ring with unit such that:
\[
\begin{array}{ll}
\text{i)}& K\text{ is torsion-free},\\
\text{ii)}& \forall x\in K,\ \forall n\in\mathbb N,\quad \binom{x}{n}\in K .
\end{array}
\]
Then, for all \(x\in K\) and all \(n\in\mathbb N\), put
\[
  \lambda^n(x)=\binom{x}{n},
\]
that is,
\sgaStableTarget{sga6.e05.eq.2-6-1}\begin{equation}\tag{2.6.1}
  \lambda_t(x)=\sum_{n=0}^{\infty}\binom{x}{n}t^n
  =(1+t)^x=\exp(x\Log(1+t)),
\end{equation}
where the last identity is considered in \(K\otimes_{\mathbb Z}\mathbb Q[[t]]\). By multiplicativity of the exponential, this defines a pre-\(\lambda\)-ring structure on \(K\).

To see that it is a \(\lambda\)-ring structure, one must prove the relations \hyperlink{sga6.e05.eq.2-4-2}{(2.4.2)}, and it is enough to prove them in \(K\otimes\mathbb Q\), by i). One is then reduced by functoriality to the case \(K=\mathbb Q[X,Y]\) and to the elements \(X,Y\). In view of the definition of the \(\lambda^i\), the two sides of the equalities \hyperlink{sga6.e05.eq.2-4-2}{(2.4.2)} are polynomials in \(X,Y\), with rational coefficients. It is therefore enough, in order to prove their equality, to show that they take the same value for all integral values of \(X,Y\). But the relations obtained by substituting integral values in \(X,Y\) are precisely those expressing that the pre-\(\lambda\)-ring \(\mathbb Z\) of 2.5.a is a \(\lambda\)-ring. They are therefore satisfied.

\sgaStableTarget{sga6.e05.def.2-7}\textbf{Definition 2.7.} A \emph{binomial ring} is a ring satisfying the conditions 2.6.i and ii), endowed with the \(\lambda\)-ring structure defined by \hyperlink{sga6.e05.eq.2-6-1}{(2.6.1)}.

\sgaStableTarget{sga6.e05.ex.2-8}\textbf{Example 2.8.} Let \(X\) be a site, \(K\) a binomial ring, and \(\mathcal K\) the sheaf associated to the constant presheaf with value \(K\) on \(X\). Then \(H^0(X,\mathcal K)\) is a binomial ring. Indeed, the conditions \hyperlink{sga6.e05.para.2-6}{2.6} are preserved under passage to an inductive limit. In particular, \(H^0(X,\mathbb Z)\) is a binomial ring.

\sgaStableTarget{sga6.e05.para.2-9}\textbf{2.9.} Let \(K\) be a binomial ring and \(A\) a \(K\)-algebra. For every series
\[
  \varphi\in\Ghat(A)
\]
and every \(a\in K\), one can define a series \(\varphi^a\) in the usual way. Writing \(\varphi=1+\psi\), put
\sgaStableTarget{sga6.e05.eq.2-9-1}\begin{equation}\tag{2.9.1}
  \varphi^a=(1+\psi)^a
  =
  \sum_{n=0}^{\infty}\binom{a}{n}\psi^n
  \left(=\exp(a\Log\varphi)\right),
\end{equation}
where the last equality can be written only modulo torsion. These formulas are obtained by substitution from the analogous formulas for \(\varphi=1+t\). Therefore the following relations, the classical properties of the exponential, remain valid for \(a,b\in K\) and \(\varphi,\psi\in\Ghat(A)\):
\sgaStableTarget{sga6.e05.eq.2-9-2}\begin{equation}\tag{2.9.2}
  \varphi^a\varphi^b=\varphi^{a+b},\qquad
  \varphi^{ab}=(\varphi^a)^b,\qquad
  \varphi^a\psi^a=(\varphi\psi)^a .
\end{equation}



% Exposé V continuation macros
\providecommand{\Ghat}{\widehat G}
\providecommand{\Gg}{\widehat G_g}
\providecommand{\Gcirc}{\widehat G_{\circ}}
\providecommand{\Gnabla}{\widehat G_{\nabla}}
\providecommand{\Grhat}{\widehat{\operatorname{Gr}}}
\providecommand{\Log}{\operatorname{Log}}
\providecommand{\Fil}{\operatorname{Fil}}
\providecommand{\eps}{\varepsilon}
\providecommand{\id}{\operatorname{id}}
\providecommand{\Sym}{\operatorname{Sym}}
\providecommand{\Mon}{\operatorname{Mon}}


\sgaStableTarget{sga6.e05.continuation-of-2-9}\subsection*{Continuation of 2.9}

We still suppose that \(K\) is binomial. Multiplication, for the law \(\circ\), by \(\lambda_t(a)\), where \(a\in K\), defines a morphism of functors
\[
  \Ghat_K\longrightarrow \Ghat_K
\]
which is characterized by its value on the element \(1+Xt\) of \(G(K[X])\). Thus
\[
  \lambda_t(a)\circ(1+Xt)=\lambda_{Xt}(a),
\]
and hence
\sgaStableTarget{sga6.e05.eq.2-9-3}\begin{equation}\tag{2.9.3}
  \lambda_t(a)\circ(1+Xt)
  =(1+Xt)^a
  =\exp(a\Log(1+Xt))
  =\sum_{n=0}^{\infty}\binom{a}{n}X^nt^n .
\end{equation}
It follows that
\sgaStableTarget{sga6.e05.eq.2-9-4}\begin{equation}\tag{2.9.4}
  \lambda_t(a)\circ\varphi=\varphi^a
  \quad(=\exp(a\Log\varphi)).
\end{equation}
Indeed, the two members are additive functors which coincide on \(1+Xt\). From \hyperlink{sga6.e05.eq.2-9-3}{(2.9.3)} and \hyperlink{sga6.e05.para.1-15}{1.15} one obtains that the universal polynomial of degree \(n\) defining multiplication by \(\lambda_t(a)\) is isobaric of weight \(n\). We denote it by \(R_{n,a}\).

Finally, if \(A\) is a \(K\)-\(\lambda\)-algebra and \(b\in A\), then
\sgaStableTarget{sga6.e05.eq.2-9-5}\begin{equation}\tag{2.9.5}
  \lambda_t(ab)=(\lambda_t(b))^a .
\end{equation}
This follows from \(\lambda_t(ab)=\lambda_t(a)\circ\lambda_t(b)\) and from \hyperlink{sga6.e05.eq.2-9-4}{(2.9.4)}.

\sgaStableTarget{sga6.e05.sec.3}\section*{3. The operations \(\gamma\)}

We still work over the category \(\mathscr C_0\).

\sgaStableTarget{sga6.e05.para.3-1}\textbf{3.1.} Let \(K\) be a \(\lambda\)-ring. We shall define from the operations \(\lambda^i\) new operations denoted \(\gamma^i\). They are the operations which are adapted to Chern classes. If \(E\) is a vector bundle of rank \(d\), the top Chern class of \(E\) is represented, in \(K\)-theoretic language, by the alternating sum of the exterior powers. To obtain the \(n\)-th class, one first replaces \(E\) by \(E\oplus\mathbf1^{\,n-d}\), so that the same top-class recipe applies. This motivates the formula below.

\sgaStableTarget{sga6.e05.prop.3-2}\textbf{Proposition 3.2.} Let \(K\) be a \(\lambda\)-ring.
\begin{enumerate}[label=\roman*)]
\item For every \(x\in K\) and every \(n\in\mathbb N\),
\[
  \lambda^n(x+n-1)=(-1)^n\sum_{i=0}^n(-1)^i\lambda^i(x+n).
\]
\item If
\[
  \gamma^n(x)=\lambda^n(x+n-1),\qquad
  \gamma_t(x)=\sum_{n\ge0}\gamma^n(x)t^n,
\]
then
\[
  \gamma_t(x)=\lambda_{t/(1-t)}(x),
  \qquad\text{hence}\qquad
  \lambda_t(x)=\gamma_{t/(1+t)}(x).
\]
\end{enumerate}
Thus knowledge of the operations \(\gamma\) is equivalent to knowledge of the operations \(\lambda\).

\emph{Proof.} For i),
\[
  \lambda^n(x+n-1)=\sum_{i=0}^{n}\lambda^i(x+n)\lambda^{n-i}(-1).
\]
Since \(\lambda_t(1)=1+t\), one has \(\lambda_t(-1)=1/(1+t)\), and hence \(\lambda^k(-1)=(-1)^k\), giving the desired relation. For ii),
\[
  \lambda_{t/(1-t)}(x)=\sum_{n\ge0}\lambda^n(x)t^n(1-t)^{-n},
\]
and
\[
  t^n(1-t)^{-n}=\sum_{p\ge n}\binom{p-1}{p-n}t^p.
\]
The coefficient of \(t^p\) is therefore
\[
  \sum_{n=0}^{p}\lambda^n(x)\binom{p-1}{p-n}
  =\sum_{n=0}^{p}\lambda^n(x)\lambda^{p-n}(p-1)
  =\lambda^p(x+p-1).
\]

\sgaStableTarget{sga6.e05.para.3-3}\textbf{3.3.} If \(f:K\to K'\) is a homomorphism of \(\lambda\)-rings, then \(f\) is a \(\lambda\)-homomorphism if and only if, for every \(i\in\mathbb N\),
\[
  f(\gamma^i(x))=\gamma^i(f(x)),
\]
or equivalently
\[
  \gamma_t(f(x))=\Ghat(f)(\gamma_t(x)).
\]

\sgaStableTarget{sga6.e05.para.3-4}\textbf{3.4.} From \hyperlink{sga6.e05.prop.3-2}{3.2} one obtains
\sgaStableTarget{sga6.e05.eq.3-4-1}\begin{equation}\tag{3.4.1}
  \gamma^0(x)=1,
  \qquad \gamma^1(x)=x,
  \qquad \gamma^n(x+y)=\sum_{p+q=n}\gamma^p(x)\gamma^q(y).
\end{equation}
Thus \(\gamma_t\) defines a new pre-\(\lambda\)-ring structure. It is not, in general, a \(\lambda\)-ring structure. Let \(\rho_t\) denote the automorphism of the functor \(\Ghat\) induced by the change of variable
\[
  t\longmapsto \frac{t}{1-t}.
\]
Then \(\gamma_t=\rho_t\lambda_t\). Transporting the \(\lambda\)-ring functor structure of \(\Ghat\) by \(\rho_t\) gives a second structure on \(\Ghat\), for which \(\gamma_t:K\to\Ghat(K)\) is a \(\lambda\)-homomorphism whenever \(K\) is a \(\lambda\)-ring.

\sgaStableTarget{sga6.e05.para.3-5}\textbf{3.5.} The multiplication in this transported structure is the composite
\[
  \Ghat\times\Ghat\xrightarrow{\rho_t^{-1}\times\rho_t^{-1}}
  \Ghat\times\Ghat\xrightarrow{\circ}\Ghat\xrightarrow{\rho_t}\Ghat .
\]
It will be denoted \(\nabla\). It is determined by
\sgaStableTarget{sga6.e05.eq.3-5-1}\begin{align}
 (1+Xt)\nabla(1+Yt)
 &=\rho_t\!\left(\rho_t^{-1}(1+Xt)\circ\rho_t^{-1}(1+Yt)\right) \notag\\
 &=\rho_t\!\left(\frac{1+(X+1)t}{1+t}\circ
                 \frac{1+(Y+1)t}{1+t}\right) \notag\\
 &=\rho_t\!\left(\frac{(1+((X+1)(Y+1))t)(1+t)}
                      {(1+(X+1)t)(1+(Y+1)t)}\right) \notag\\
 &=\frac{1+(1+X+Y+XY)t}{(1+Xt)(1+Yt)} . \tag{3.5.1}
\end{align}
Thus, for \(x,y\in K\),
\sgaStableTarget{sga6.e05.eq.3-5-2}\begin{equation}\tag{3.5.2}
  \gamma_t(xy)=\gamma_t(x)\nabla\gamma_t(y).
\end{equation}
The universal polynomial of degree \(n\) for this law will be denoted \(P'_n\). Its isobaric components have weights between \(n\) and \(2n\) in each of the two groups of variables, and
\sgaStableTarget{sga6.e05.eq.3-5-3}\begin{equation}\tag{3.5.3}
  \gamma^n(xy)
  =P'_n(\gamma^1(x),\ldots,\gamma^n(x),
       \gamma^1(y),\ldots,\gamma^n(y)).
\end{equation}
The unit for \(\nabla\) is \(\gamma_t(1)=1/(1-t)\).

\sgaStableTarget{sga6.e05.para.3-6}\textbf{3.6.} The operations \(\lambda'\) of the transported structure are defined by
\sgaStableTarget{sga6.e05.eq.3-6-1}\[
\begin{tikzcd}[column sep=tiny]
\Ghat \\[-0.6em]
1+K[[t]]^+ \arrow[r,"\lambda_u"] \arrow[d,"\rho_t"'] &
1+\{1+K[[t]]^+\}[[u]]^+ \arrow[d,"\Ghat(\rho_t)"]\\
1+K[[t]]^+ \arrow[r,"\lambda'_u"'] &
1+\{1+K[[t]]^+\}[[u]]^+ .
\end{tikzcd}
\tag{3.6.1}
\]
If \(\gamma'\) denotes the \(\gamma\)-operations attached to these operations \(\lambda'\), then
\sgaStableTarget{sga6.e05.eq.3-6-2}\begin{equation}\tag{3.6.2}
  \gamma'_u(1+Xt)=\frac1{1-u}+\{1+Xt\}u .
\end{equation}
Let \(Q'_{i,j}\) be the coefficient of \(t^j\) in \(\gamma'^i(1+X_1t+X_2t^2+\cdots)\). It is a polynomial in \(\mathbb Z[X_1,\ldots,X_{ij}]\), whose isobaric components have weights from \(j\) to \(ij\). For a \(\lambda\)-ring \(K\) and \(x\in K\),
\sgaStableTarget{sga6.e05.eq.3-6-3}\begin{equation}\tag{3.6.3}
  \gamma_t(\lambda^i(x))=\lambda'^i(\gamma_t(x)),
  \qquad
  \gamma_t(\gamma^i(x))=\gamma'^i(\gamma_t(x)),
\end{equation}
and therefore
\sgaStableTarget{sga6.e05.eq.3-6-4}\begin{equation}\tag{3.6.4}
  \gamma^j(\gamma^i(x))
  =Q'_{i,j}(\gamma^1(x),\ldots,\gamma^{ij}(x)).
\end{equation}
The transported pre-\(\lambda\)-ring structure on \(\Ghat\) is itself a \(\lambda\)-ring structure.

\sgaStableTarget{sga6.e05.para.3-7}\textbf{3.7.} We denote by \(\Gcirc\) the functor \(\Ghat\) with the structure of §2, and by \(\Gnabla\) the same functor with the structure of §3. Giving a \(\lambda\)-ring structure on a ring \(A\) is equivalent either to giving a ring homomorphism \(A\to\Gcirc(A)\) satisfying the usual diagram for \(\lambda_t\), or to giving a ring homomorphism \(A\to\Gnabla(A)\) satisfying the corresponding diagram for \(\gamma_t\). Explicitly, the two diagrams are
\sgaStableTarget{sga6.e05.eq.3-7-1}\[
\begin{tikzcd}[column sep=tiny]
A \arrow[r,"\lambda_u"] \arrow[d,"\lambda_t"'] & \Gcirc^u(A) \arrow[d,"\Gcirc^u(\lambda_t)"]\\
\Gcirc^t(A) \arrow[r,"\lambda_u"'] & \Gcirc^u(\Gcirc^t(A))
\end{tikzcd}
\tag{3.7.1}
\]
and
\sgaStableTarget{sga6.e05.eq.3-7-2}\[
\begin{tikzcd}[column sep=tiny]
A \arrow[r,"\gamma_u"] \arrow[d,"\gamma_t"'] & \Gnabla^u(A) \arrow[d,"\Gnabla^u(\gamma_t)"]\\
\Gnabla^t(A) \arrow[r,"\gamma'_u"'] & \Gnabla^u(\Gnabla^t(A)).
\end{tikzcd}
\tag{3.7.2}
\]

\sgaStableTarget{sga6.e05.para.3-8}\textbf{3.8.} Let \(K\) be a binomial ring. From \hyperlink{sga6.e05.eq.2-6-1}{(2.6.1)}, for \(x\in K\),
\sgaStableTarget{sga6.e05.eq.3-8-1}\begin{equation}\tag{3.8.1}
  \gamma_t(x)=\left(\frac1{1-t}\right)^x
  =\exp(-x\Log(1-t)).
\end{equation}
For \(a\in K\), multiplication by \(\gamma_t(a)\) for the law \(\nabla\) is characterized by its value on \(1+Xt\):
\sgaStableTarget{sga6.e05.eq.3-8-2}\begin{equation}\tag{3.8.2}
  \gamma_t(a)\nabla(1+Xt)
  =(1+Xt)^a
  =\exp(a\Log(1+Xt))
  =\sum_{n=0}^{\infty}\binom{a}{n}X^nt^n .
\end{equation}
Thus the universal polynomial of degree \(n\) defining this multiplication is again \(R_{n,a}\).

\sgaStableTarget{sga6.e05.def.3-9}\textbf{Definition 3.9.} Let \(K\) be a \(\lambda\)-ring. An \emph{augmented \(K\)-\(\lambda\)-algebra} is a \(K\)-\(\lambda\)-algebra \(A\), together with a \(\lambda\)-homomorphism of \(K\)-algebras
\[
  \epsilon:A\longrightarrow K.
\]
Morphisms are \(\lambda\)-homomorphisms compatible with the augmentations.

\sgaStableTarget{sga6.e05.ex.3-9-1}\textbf{Example 3.9.1.} Let \((X,A)\) be a locally ringed topos. The Grothendieck group \(K^\circ(X)\) of locally free finite-type \(A\)-modules has a natural pre-\(\lambda\)-ring structure by exterior powers. There is a natural \(\lambda\)-homomorphism from \(H^0(X,\mathbb Z)\) to \(K^\circ(X)\): a section \(a\) of the constant sheaf \(\mathbb Z\) partitions the final object into opens \(U_n\) on which \(a=n\), and is sent to the class of the locally free module which is \(A^n\) on \(U_n\) for \(n\ge0\), minus the analogous class on the negative pieces. The rank map is the augmentation. This is the augmented structure used below.

\sgaStableTarget{sga6.e05.para.3-10}\textbf{3.10.} Let \(A\) be an augmented \(K\)-\(\lambda\)-algebra. For \(n>0\), define \(\Fil^n A\) to be the ideal generated by products
\sgaStableTarget{sga6.e05.eq.3-10-1}\begin{equation}\tag{3.10.1}
  \gamma^{i_1}(x_1)\cdots\gamma^{i_k}(x_k),
  \qquad i_1+\cdots+i_k\ge n,
  \qquad \epsilon(x_i)=0.
\end{equation}
If \(A\) is generated as an abelian group by \((x_\alpha)\), then \(\Fil^nA\) is generated over \(K\) by
\sgaStableTarget{sga6.e05.eq.3-10-2}\begin{equation}\tag{3.10.2}
  \gamma^{i_1}(x_{\alpha_1}-\epsilon(x_{\alpha_1}))\cdots
  \gamma^{i_k}(x_{\alpha_k}-\epsilon(x_{\alpha_k})),
  \qquad i_1+\cdots+i_k\ge n.
\end{equation}
The filtration is decreasing and multiplicative, and
\[
  \Fil^0A=A,
  \qquad \Fil^1A=\Ker\epsilon.
\]
It is functorial for homomorphisms of augmented \(K\)-\(\lambda\)-algebras.

\sgaStableTarget{sga6.e05.prop.3-11}\textbf{Proposition 3.11.} Let \(A\) be a \(\lambda\)-ring and let \(y_1,\ldots,y_n\) be elements of \(A\) such that \(\gamma^j(y_i)=0\) for \(j\ge2\). Put \(x=\sum_i y_i\). Then \(\gamma^k(x-n)=0\) for \(k>n\), and, for \(1\le k\le n\), \(\gamma^k(x-n)\) is the \(k\)-th elementary symmetric function of the \(y_i-1\). This follows from
\[
  \gamma_t(x-n)=\prod_{i=1}^{n}\gamma_t(y_i-1).
\]

\sgaStableTarget{sga6.e05.sec.4}\section*{4. \(\lambda\)-rings generated by generators and relations}

\sgaStableTarget{sga6.e05.para.4-1}\textbf{4.1.} Let \(K\) be a \(\lambda\)-ring and let \((X_i)_{i\in I}\) be indeterminates. Put
\[
  A=K[X_i,{}^nX_i]_{i\in I,\ n\ge2}.
\]
Define
\sgaStableTarget{sga6.e05.eq.4-1-1}\begin{equation}\tag{4.1.1}
  \lambda_t(X_i)=1+X_it+\sum_{n\ge2}{}^nX_i t^n,
  \qquad
  \lambda_t({}^nX_i)=\lambda^n(\lambda_t(X_i)).
\end{equation}
This gives a \(\lambda\)-ring structure on \(A\) extending that of \(K\).

\sgaStableTarget{sga6.e05.prop.4-2}\textbf{Proposition 4.2.} For every \(K\)-\(\lambda\)-algebra \(B\) and every family \((b_i)_{i\in I}\) in \(B\), there exists a unique \(K\)-\(\lambda\)-homomorphism \(A\to B\) carrying \(X_i\) to \(b_i\). Necessarily it carries \({}^nX_i\) to \(\lambda^n(b_i)\), and this condition is sufficient.

\sgaStableTarget{sga6.e05.para.4-3}\textbf{4.3.} The same construction may be made using \(\gamma\)'s. Let
\[
  A'=K[X_i,{}^\gamma_nX_i]_{i\in I,\ n\ge2},
\]
and define
\sgaStableTarget{sga6.e05.eq.4-3-1}\begin{equation}\tag{4.3.1}
  \gamma_t(X_i)=1+X_it+\sum_{n\ge2}{}^\gamma_nX_i t^n,
  \qquad
  \gamma_t({}^\gamma_nX_i)=\gamma'^n(\gamma_t(X_i)).
\end{equation}
This gives a \(\lambda\)-ring satisfying the same universal property, hence canonically isomorphic to \(A\).

\sgaStableTarget{sga6.e05.def.4-4}\textbf{Definition 4.4.} The preceding \(\lambda\)-ring is the \emph{free \(\lambda\)-ring} over \(K\) generated by the family \((X_i)\).

\sgaStableTarget{sga6.e05.def.4-5}\textbf{Definition 4.5.} An ideal \(J\) of a \(\lambda\)-ring \(A\) is a \emph{\(\lambda\)-ideal} if \(\lambda^n(x)\in J\) for every \(x\in J\) and \(n>1\). Kernels of \(\lambda\)-homomorphisms are \(\lambda\)-ideals; conversely, if \(J\) is a \(\lambda\)-ideal, \(A/J\) receives a unique \(\lambda\)-ring structure. The intersection of \(\lambda\)-ideals is a \(\lambda\)-ideal. Thus every family of elements has a smallest containing \(\lambda\)-ideal, generated by the elements and all their iterated \(\lambda^n\)'s.

\sgaStableTarget{sga6.e05.def.4-6}\textbf{Definition 4.6.} The quotient of the free \(\lambda\)-ring by the \(\lambda\)-ideal generated by prescribed relations \((P_i)\) is called the \(\lambda\)-ring generated by the generators \((X_i)\) subject to the relations \((P_i)\). Equivalently, one may formulate the same construction in terms of \(\gamma\)-operations.

\sgaStableTarget{sga6.e05.prop.4-7}\textbf{Proposition 4.7.} This quotient has the expected universal property: every family \((b_i)\) in a \(K\)-\(\lambda\)-algebra satisfying the relations \((P_i)\) determines a unique \(K\)-\(\lambda\)-homomorphism from the quotient to that algebra.

\sgaStableTarget{sga6.e05.prop.4-8}\textbf{Proposition 4.8.} Let \(n\ge0\). In the free \(\lambda\)-ring over \(K\) on one generator \(X\), the ideal generated by the \(\lambda^kX\), \(k>n\), is a \(\lambda\)-ideal. Indeed, after quotienting, one has only to check that if \(\varphi\in\Ghat(C)\) has no terms of degree \(>n\), then \(\lambda^k(\varphi)=0\) for \(k>n\). By the splitting principle, it is enough to write
\[
  \varphi=\prod_{i=1}^{n}(1+T_it),
\]
and then
\[
  \lambda_u(\varphi)=\prod_{i=1}^{n}\bigl(1+\{1+T_it\}u\bigr),
\]
which has degree \(n\) in \(u\).

\sgaStableTarget{sga6.e05.cor.4-9}\textbf{Corollary 4.9.} The \(\lambda\)-ring generated over \(K\) by one element \(X\), subject to \(\lambda^k(X)=0\) for \(k>n\), is
\[
  K[X,\lambda^2X,\ldots,\lambda^nX].
\]
More generally, for a family \((X_i)\) with bounds \((n_i)\), the quotient subject to \(\lambda^k(X_i)=0\) for \(k>n_i\) is the polynomial ring generated by \(X_i,\lambda^2X_i,\\ldots,\lambda^{n_i}X_i\).

\sgaStableTarget{sga6.e05.para.4-10}\textbf{4.10.} Let \(K\) be binomial and let \(A\) be an augmented \(K\)-\(\lambda\)-algebra. Put
\[
  B=A[X,\gamma^2X,\gamma^3X,\ldots],
\]
the free \(A\)-\(\lambda\)-algebra on \(X\), written in \(\gamma\)-coordinates. Extend the augmentation by \(\epsilon(X)=0\), hence \(\epsilon(\gamma^iX)=0\) for \(i>0\). Let \(M_k\) be the set of monomials in the variables \(\gamma^jX\), counted with weight \(j\), of total weight \(k\), and set \(M_0=\{1\}\).

\sgaStableTarget{sga6.e05.lem.4-11}\textbf{Lemma 4.11.} If \(x\in\Fil^pA\), then every coefficient of \(\gamma^n(x+X)\) which is homogeneous of weight \(k\) in the \(\gamma^jX\) belongs to \(\Fil^{p+n-k}A\).

\sgaStableTarget{sga6.e05.prop.4-12}\textbf{Proposition 4.12.} For every \(k\ge0\),
\sgaStableTarget{sga6.e05.eq.4-12-1}\begin{equation}\tag{4.12.1}
  \Fil^kB=
  \Fil^kA+\Fil^{k-1}A\,M_1+\cdots+\Fil^1A\,M_{k-1}
  +\Fil^0A\,M_k+B\,M_{k+1}.
\end{equation}
The inclusion from right to left follows from \hyperlink{sga6.e05.lem.4-11}{Lemma 4.11}. The reverse inclusion is obtained by writing the generators of \(\Fil^kB\) as products of \(\gamma\)-operations and reducing, term by term, to the displayed weights.

\sgaStableTarget{sga6.e05.rem.4-13}\textbf{Remarks 4.13.} The proof gives the same assertion for a free \(\lambda\)-algebra on any finite or infinite family of generators, with the evident multi-weight. It is also unchanged when one uses the \(\lambda\)-coordinates, since \(\lambda\) and \(\gamma\) coordinates differ by triangular universal formulas.

\sgaStableTarget{sga6.e05.cor.4-14}\textbf{Corollary 4.14.} Passing to the associated graded ring, one has
\sgaStableTarget{sga6.e05.eq.4-14-1}\begin{equation}\tag{4.14.1}
  \Gr^kB\simeq
  \Gr^kA\oplus \Gr^{k-1}A\,M_1\oplus\cdots\oplus\Gr^0A\,M_k.
\end{equation}

\sgaStableTarget{sga6.e05.cor.4-15}\textbf{Corollary 4.15.} If \(B'=K[X,\gamma^2X,\ldots,\gamma^nX]\) is the \(K\)-\(\lambda\)-algebra generated by one element satisfying \(\gamma^k(X)=0\) for \(k>n\), then the \(\lambda\)-ring filtration of \(B'\) is exactly the weight filtration in the variables \(\gamma^jX\).

\sgaStableTarget{sga6.e05.sec.5}\section*{5. The elements \(\lambda^p(N,x)\)}

This section will be used only in \hyperlink{sga6.e07.expose}{Exposé VII}, §4, in the study of regular immersions; it may be omitted on first reading.

\sgaStableTarget{sga6.e05.para.5-1}\textbf{5.1.} Let \(d\ge1\). Let \(A\) be the \(\mathbb Z\)-\(\lambda\)-algebra generated by \(N\) and \(x\), subject to the relations \(\lambda^i(N)=0\) for \(i>d\). Thus
\[
  A=\mathbb Z[N,\ldots,\lambda^dN;
       x,\ldots,\lambda^jx,\ldots].
\]
Put
\sgaStableTarget{sga6.e05.eq.5-1-1}\begin{equation}\tag{5.1.1}
  \lambda_{-1}(N)=\sum_{i=0}^{d}(-1)^i\lambda^i(N).
\end{equation}

\sgaStableTarget{sga6.e05.lem.5-2}\textbf{Lemma 5.2.} For every \(p\ge1\), the element \(\lambda^p(x\lambda_{-1}(N))\) is divisible by \(\lambda_{-1}(N)\).

It is enough, by isobaricity, to prove this for \(\lambda^p(\lambda_{-1}(N))\). Let \(B\) be the \(\lambda\)-ring generated by \(N_1,\ldots,N_d\), with \(\lambda^i(N_j)=0\) for \(i\ge2\). The map \(N\mapsto\sum_iN_i\) identifies \(A\) with the symmetric polynomials in the \(N_i\), and
\[
  \lambda_{-1}(N)=\prod_{i=1}^{d}(1-N_i).
\]
If some \(N_i=1\), then \(\lambda_t(\lambda_{-1}(N))=\lambda_t(0)=1\), so every positive coefficient vanishes. Thus each \(\lambda^p(\lambda_{-1}(N))\) is divisible by every \(1-N_i\), hence by their product.

\sgaStableTarget{sga6.e05.def.5-3}\textbf{Definition 5.3.} Since \(A\) is integral, define \(\lambda^p(N,x)\) by
\sgaStableTarget{sga6.e05.eq.5-3-1}\begin{equation}\tag{5.3.1}
  \lambda^p(N,x)\lambda_{-1}(N)=\lambda^p(x\lambda_{-1}(N)),
  \qquad p\ge1.
\end{equation}
Thus \(\lambda^p(N,x)\) is a polynomial in the \(\lambda^i(x)\) and \(\lambda^j(N)\), isobaric of weight \(p\) in the \(\lambda^i(x)\). If a \(\lambda\)-ring \(B\) contains an element \(N'\) with \(\lambda^k(N')=0\) for \(k>d\), then \(\lambda^p(N',x')\) is defined by substitution; it is independent of the chosen bound \(d\).

\sgaStableTarget{sga6.e05.para.5-4}\textbf{5.4.} In the \(\lambda\)-ring generated over \(\mathbb Z\) by \(N,x,y\), subject to \(\lambda^k(N)=0\) for \(k>d\), formula \hyperlink{sga6.e05.eq.5-3-1}{(5.3.1)} gives
\begin{align*}
 \lambda^p(N,x+y)\lambda_{-1}(N)
 &=\lambda^p\bigl((x+y)\lambda_{-1}(N)\bigr)\\
 &=\sum_{i+j=p}
   \lambda^i\bigl(x\lambda_{-1}(N)\bigr)
   \lambda^j\bigl(y\lambda_{-1}(N)\bigr)\\
 &=\sum_{\substack{i+j=p\\ i,j\ne0}}
   \lambda^i(N,x)\lambda^j(N,y)\bigl(\lambda_{-1}(N)\bigr)^2
   +\bigl(\lambda^p(N,x)+\lambda^p(N,y)\bigr)\lambda_{-1}(N).
\end{align*}
Since the ring under consideration is an integral domain, this yields the relation, valid by substitution in every \(\lambda\)-ring,
\sgaStableTarget{sga6.e05.eq.5-4-1}\begin{equation}\tag{5.4.1}
 \lambda^p(N,x+y)
 =\sum_{\substack{i+j=p\\ i,j\ne0}}
   \lambda^i(N,x)\lambda^j(N,y)\lambda_{-1}(N)
   +\lambda^p(N,x)+\lambda^p(N,y).
\end{equation}
Now consider the \(\lambda\)-ring generated over \(\mathbb Z\) by \(N,N',x\), subject to \(\lambda^k(N)=0\) for \(k>d\) and \(\lambda^k(N')=0\) for \(k>d'\). Then \(\lambda^k(N+N')=0\) for \(k>d+d'\), and
\begin{align*}
 \lambda^p(N+N',x)\lambda_{-1}(N)\lambda_{-1}(N')
 &=\lambda^p\bigl(x\lambda_{-1}(N)\lambda_{-1}(N')\bigr)\\
 &=\lambda^p\bigl(N',x\lambda_{-1}(N)\bigr)\lambda_{-1}(N').
\end{align*}
Here we used \(\lambda_{-1}(N+N')=\lambda_{-1}(N)\lambda_{-1}(N')\). Cancelling in the integral domain gives the relation, again valid by substitution,
\sgaStableTarget{sga6.e05.eq.5-4-2}\begin{equation}\tag{5.4.2}
  \lambda^p(N+N',x)\lambda_{-1}(N)
  =\lambda^p\bigl(N',x\lambda_{-1}(N)\bigr).
\end{equation}

\sgaStableTarget{sga6.e05.para.5-5}\textbf{5.5.} Let \(A\) be a \(\lambda\)-ring and let \(N\in A\) satisfy \(\lambda^k(N)=0\) for \(k>d\). Define on the underlying abelian group of \(A\) a new multiplication by
\[
  x\cdot_N y=xy\lambda_{-1}(N).
\]
After adjoining a unit, one obtains a ring \(\mathbb Z\times A\) with
\[
  (n,x)(m,y)=(nm,ny+mx+x\cdot_N y).
\]
The formulas
\sgaStableTarget{sga6.e05.eq.5-5-1}\begin{equation}\tag{5.5.1}
  \lambda_t(n,x)=(1+t)^n\left(1+
  \sum_{p\ge1}\lambda^p(N,x)t^p\right)
\end{equation}
define a pre-\(\lambda\)-ring structure. In fact, when \(A\) is a \(\lambda\)-ring, this is a \(\lambda\)-ring structure; it is the one transported from the original structure by multiplication by \(\lambda_{-1}(N)\). We denote this augmented \(\lambda\)-ring by \(\widehat A_N\). Its \(\gamma\)-operations satisfy
\sgaStableTarget{sga6.e05.eq.5-5-2}\begin{equation}\tag{5.5.2}
  \gamma^p(N,x)\lambda_{-1}(N)=\gamma^p(x\lambda_{-1}(N)).
\end{equation}

\sgaStableTarget{sga6.e05.para.5-6}\textbf{5.6.} Suppose that, in addition, \(L\) is an element such that \(\lambda^k(L)=0\) for \(k>1\). Then \(\lambda^p(N+L,x)\) is the coefficient of \(t^p\) in the series
\sgaStableTarget{sga6.e05.eq.5-6-1}\begin{equation}\tag{5.6.1}
  \frac{\displaystyle\sum_{k\ge1}\lambda^k(N,x)
        (1+L+\cdots+L^{k-1})t^k}
       {\displaystyle 1+\sum_{k\ge1}\lambda^k(N,x)
        \lambda_{-1}(N)L^k t^k}.
\end{equation}
This is obtained by applying \hyperlink{sga6.e05.eq.5-4-2}{(5.4.2)} to \(N+L\) and using \(\lambda_{-1}(L)=1-L\).

\sgaStableTarget{sga6.e05.prop.5-7}\textbf{Proposition 5.7.} Let \(K\) be a binomial ring, \(A\) an augmented \(K\)-\(\lambda\)-algebra, and \(N\in A\) such that \(\lambda^k(N)=0\) for \(k>d\). Then, for every \(x\in A\) and every \(p\ge1\),
\sgaStableTarget{sga6.e05.eq.5-7-1}\begin{equation}\tag{5.7.1}
  \gamma^p(N,x)\in\Fil^{p-d}A.
\end{equation}
Indeed, let \(A_0\) be the \(K\)-\(\lambda\)-algebra generated by \(N'_0,X'_0\), subject to \(\gamma^k(N'_0)=0\) for \(k>d\), and augmented by \(\epsilon(N'_0)=\epsilon(X'_0)=0\). This is a polynomial algebra over \(K\) in the variables \(\gamma^i(N'_0)\), \(1\le i\le d\), and \(\gamma^j(X'_0)\), \(j\ge1\), and its \(\lambda\)-filtration is the weight filtration, with weight \(i\) assigned to \(\gamma^i(N'_0)\) and \(\gamma^i(X'_0)\) (\hyperlink{sga6.e05.prop.4-12}{4.12} and \hyperlink{sga6.e05.cor.4-15}{4.15}). Put
\[
  N_0=N'_0+d,
  \qquad X_0=X'_0+\epsilon(x).
\]
For \(N_0,X_0\), the assertion is equivalent to
\[
  \gamma^p(N_0,X_0)\gamma^d(N'_0)\in\Fil^pA_0.
\]
Since \(\gamma^d(N'_0)=(-1)^d\lambda_{-1}(N_0)\), this becomes
\[
  (-1)^d\gamma^p\bigl(X_0\gamma^d(N'_0)\bigr)\in\Fil^pA_0,
\]
by \hyperlink{sga6.e05.eq.5-5-2}{(5.5.4)}. This holds because \(\gamma^d(N'_0)\) has zero augmentation. By \hyperlink{sga6.e05.prop.4-7}{4.7} there is a \(K\)-\(\lambda\)-homomorphism \(A_0\to A\), compatible with augmentations and sending \(N'_0\) to \(N-d\) and \(X'_0\) to \(x-\epsilon(x)\), hence sending \(N_0\) to \(N\) and \(X_0\) to \(x\). Compatibility with the filtrations gives the result.

\sgaStableTarget{sga6.e05.sec.6}\section*{6. Chern ring: opening}

Throughout this section the base ring is a fixed binomial ring \(K\). We use the functor \(G\) of \hyperlink{sga6.e05.para.1-1}{1.1}, and return to the notation of §1: \(G\) denotes the functor there defined and \(G_i\) its restriction to \(\mathscr C_0\).

\sgaStableTarget{sga6.e05.para.6-1}\textbf{6.1.} In \hyperlink{sga6.e05.para.3-5}{3.5} and \hyperlink{sga6.e05.para.3-6}{3.6} we defined morphisms of functors
\[
  G_i\times G_i\xrightarrow{\star}G_i,
  \qquad
  G_i^\circ\xrightarrow{\gamma'}G_i^\circ\circ G_S,
  \qquad
  G_i\xrightarrow{a}G_i,
\]
where the last map is multiplication, for the law \(\nabla\), by \(\gamma_t(a)\) with \(a\in K\). By \hyperlink{sga6.e05.lem.1-9}{1.9} these give morphisms
\[
  G\times G\xrightarrow{\star}G,
  \qquad
  G\xrightarrow{\gamma'}G\circ G_S,
  \qquad
  G\xrightarrow{a}G.
\]
The isobaricity of the universal polynomials involved implies that the first and third images lie in \(G\), while the second lies in the subgroup of formal series in \(u\) whose non-zero terms have the prescribed degrees.



% Local macros for Exposé V, §§6--7
\providecommand{\Ghat}{\widehat G}
\providecommand{\Aqhat}{\widehat A_{\mathbb Q}}
\providecommand{\Ch}{\operatorname{Ch}}
\providecommand{\Log}{\operatorname{Log}}
\providecommand{\id}{\operatorname{id}}
\providecommand{\gr}{\operatorname{gr}}
\providecommand{\Hom}{\operatorname{Hom}}
\providecommand{\epsi}{\varepsilon}
\providecommand{\clgr}{\operatorname{cl}}
\providecommand{\Fil}{\operatorname{Fil}}
\providecommand{\Gr}{\operatorname{Gr}}
\providecommand{\Sym}{\operatorname{Sym}}


\sgaStableTarget{sga6.e05.expose.occ2}\section*{Exposé V, 6. Chern ring: completion}

\sgaStableTarget{sga6.e05.def.6-2}\textbf{Definition 6.2.} The functor defined in \hyperlink{sga6.e05.para.6-1}{6.1} on the category of graded \(K\)-algebras, with values in the category of \(K\)-\(\lambda\)-algebras, is called the \emph{Chern ring}; its value on \(A\in\operatorname{Ob}(\mathscr C)\) is denoted \(\Ch(A)\), or, if there is a risk of confusion, \(\Ch_K(A)\).

From formulas \hyperlink{sga6.e05.para.6-1}{(6.1.1)}, \hyperlink{sga6.e05.para.6-1}{(6.1.2)}, \hyperlink{sga6.e05.para.6-1}{(6.1.4)}, and \hyperlink{sga6.e05.eq.3-6-2}{(3.6.2)}, one deduces the relations
\sgaStableTarget{sga6.e05.eq.6-2-1}\begin{equation}\tag{6.2.1}
 (1,1+X)(1,1+Y)=(1,1+X+Y),
 \qquad
 \lambda^i(1,1+X)=0\quad\text{for }i\ge2,
\end{equation}
the relations which served in [1] to define the Chern ring. Conversely, from these formulas one can recover the relations \hyperlink{sga6.e05.para.6-1}{(6.1.1)} and \hyperlink{sga6.e05.para.6-1}{(6.1.4)}.

Let us indicate why the Chern ring is introduced. We place ourselves in the situation of \hyperlink{sga6.e05.ex.3-9-1}{Example 3.9}, and write \(A\) for the graded ring associated, by the filtration defined in \hyperlink{sga6.e05.para.3-10}{3.10}, with the \(K^\bullet\) of a commutatively locally ringed topos; for simplicity suppose that \(K^\bullet\) is a \(\lambda\)-ring. The augmentation ring is \(K\). We shall see below that for every element of an augmented \(\lambda\)-ring one can define Chern classes with values in the associated graded ring \hyperlink{sga6.e05.para.6-7}{(6.7)}. To every locally free sheaf of finite type one can therefore associate an element of \(\Ch(A)\), whose component on the factor \(K\) is the rank of the sheaf and whose component on the factor \(1+\widehat A^+\) is the sum of the Chern classes. Formulas \hyperlink{sga6.e05.eq.6-2-1}{(6.2.1)} are then the usual formulas giving the Chern classes and the rank of the tensor product of two invertible sheaves and of the exterior powers of an invertible sheaf; in other words, the map just defined from \(K^\bullet\) to \(\Ch(A)\) is a \(\lambda\)-homomorphism.

\sgaStableTarget{sga6.e05.para.6-3}\textbf{6.3.} For every graded \(K\)-algebra \(A\), denote by \(A_{\mathbb Q}\) the \(K\)-algebra \(A\otimes\mathbb Q\), graded in the evident way. By \hyperlink{sga6.e05.thm.1-13}{1.13} one defines a functorial homomorphism
\[
  1+\widehat A^+\xrightarrow{\ \ch\ }\widehat A_{\mathbb Q}^{+}
\]
by giving its value on \(1+X\):
\sgaStableTarget{sga6.e05.eq.6-3-1}\begin{equation}\tag{6.3.1}
  \ch(1+X)=\exp(X)-1.
\end{equation}
This homomorphism is compatible with the multiplication defined in \hyperlink{sga6.e05.para.6-1}{6.1} on \(1+\widehat A^+\). By \hyperlink{sga6.e05.thm.1-13}{1.13} it is enough to show this for the elements \(1+X\), \(1+Y\) of the universal ring \(1+K[X,Y]^+\). One has
\[
\begin{aligned}
 \ch\bigl((1+X)*(1+Y)\bigr)
 &= \ch\!\left(\frac{1+X+Y}{(1+X)(1+Y)}\right)\\
 &= \ch(1+X+Y)-\ch(1+X)-\ch(1+Y)\\
 &= \exp(X+Y)-\exp(X)-\exp(Y)+1\\
 &= \ch(1+X)\,\ch(1+Y).
\end{aligned}
\]

Let \(\eta\) be the endomorphism of \(\widehat A_{\mathbb Q}^{+}\), functorial in \(A\), defined by
\[
 \eta\!\left(\sum_{k\ge1} a_k t^k\right)
 =
 \sum_{k\ge1}\frac{(-1)^{k-1}}{(k-1)!}\,a_k t^k.
\]
Then, for every \(\varphi\in1+\widehat A^+\),
\sgaStableTarget{sga6.e05.eq.6-3-2}\begin{equation}\tag{6.3.2}
  \ch(\varphi)=\eta(\Log\varphi).
\end{equation}
Indeed, \(\ch\) and \(\eta\circ\Log\) are two functorial homomorphisms of abelian groups, and it is enough for them to agree on \(1+X\), which is evident.

It follows that \(\ch\) commutes with the operations of \(K\): again it is enough to show this for \(1+X\), and for \(a\in K\) one has
\[
\begin{aligned}
 \ch((1+X)^a)
 &=\ch(\exp(a\Log(1+X)))\\
 &=\eta(a\Log(1+X))\\
 &=a\,\eta(\Log(1+X))\\
 &=a\,\ch(1+X).
\end{aligned}
\]
Thus one can extend \(\ch\) in the evident way to a homomorphism of \(K\)-algebras
\[
  K\times(1+\widehat A^+)\longrightarrow K\oplus\widehat A_{\mathbb Q}^{+},
\]
namely
\sgaStableTarget{sga6.e05.eq.6-3-3}\begin{equation}\tag{6.3.3}
  \ch(a,\varphi)=a+\eta(\Log\varphi).
\end{equation}

\sgaStableTarget{sga6.e05.def.6-4}\textbf{Definition 6.4.} The homomorphism
\[
  \ch:\Ch A\longrightarrow K\oplus\widehat A_{\mathbb Q}^{+}
\]
just defined is called the \emph{Chern character}.

If \(A\) is of characteristic zero, that is, is a \(\mathbb Q\)-algebra, the Chern character is an isomorphism. Indeed, if \(\psi\in K\oplus\widehat A^+\), one may write \(\psi=a+\psi_1\), where \(\psi_1\) is of positive degree, and then an inverse morphism is defined by
\[
  \psi\longmapsto (a,\exp(\eta^{-1}(\psi_1))).
\]

\sgaStableTarget{sga6.e05.prop.6-5}\textbf{Proposition 6.5.} Let \(A\) be a graded \(K\)-algebra. On \(\Ch(A)\) consider the filtration defined as follows, the degree filtration:
\sgaStableTarget{sga6.e05.eq.6-5-1}\begin{equation}\tag{6.5.1}
  (\Ch(A))_0=\Ch(A),\qquad
  (\Ch(A))_n=0\times(1+\widehat A^{\ge n})\quad(n\ge1),
\end{equation}
where \(\widehat A^{\ge n}=\prod_{i\ge n}A^i\). The following congruences hold.

\begin{enumerate}[label=\roman*)]
\item If \(n\ge1\) and \(\varphi\in(\Ch A)_n\), writing \(\varphi=1+a_n+\cdots\), then
\sgaStableTarget{sga6.e05.eq.6-5-2}\begin{equation}\tag{6.5.2}
  \ch(\varphi)=\frac{(-1)^{n-1}}{(n-1)!}\,a_n+\cdots .
\end{equation}

\item If \(\varphi\in(\Ch A)_m\), \(\psi\in(\Ch A)_n\), and if
\(\varphi=1+a_m+\cdots\), \(\psi=1+b_n+\cdots\), then
\sgaStableTarget{sga6.e05.eq.6-5-3}\begin{equation}\tag{6.5.3}
  \varphi*\psi
  =
  1-\frac{(m+n-1)!}{(m-1)!(n-1)!}\,a_m b_n+\cdots .
\end{equation}
Consequently the degree filtration is a ring filtration.
\end{enumerate}

For (i), the Chern character is by construction a morphism of filtered functors; from the congruence \(\varphi\equiv1+a_n\pmod{(\Ch A)_{n+1}}\) one obtains
\[
  \ch(\varphi)=\ch(1+a_n)\pmod{\widehat A^{\ge n+1}},
\]
and, since \(\ch(1+a_n)=\eta(\Log(1+a_n))\), this is
\[
  \frac{(-1)^{n-1}}{(n-1)!}a_n+\cdots .
\]
For (ii), by substitution it is enough to prove \hyperlink{sga6.e05.eq.6-5-3}{(6.5.3)} for the torsion-free universal polynomial algebra
\[
 A=K[X_i,Y_j]_{i,j\in\mathbb N}
\]
and for \(\varphi=1+\sum_{k\ge m}X_k\), \(\psi=1+\sum_{k\ge n}Y_k\). Since \(K\) is torsion-free in the universal case, one can embed this algebra in its tensor product with \(\mathbb Q\), so that \(\ch\) is injective. Put \(\varphi*\psi=1+Z_p+\cdots\), where \(Z_p\) is the first non-zero term different from \(1\). The multiplicativity of \(\ch\), together with \hyperlink{sga6.e05.eq.6-5-2}{(6.5.2)}, gives
\[
  1+\frac{(-1)^{p-1}}{(p-1)!}Z_p+\cdots
  =
  \left(1+\frac{(-1)^{m-1}}{(m-1)!}X_m+\cdots\right)
  \left(1+\frac{(-1)^{n-1}}{(n-1)!}Y_n+\cdots\right),
\]
which gives the stated value of \(Z_p\) by comparing the terms of lowest degree.

\sgaStableTarget{sga6.e05.para.6-6}\textbf{6.6.} The \(\lambda\)-ring \(\Ch A\) is augmented over \(K\): one defines a \(\lambda\)-homomorphism \(\Ch A\to K\) by \((a,\varphi)\mapsto a\). It is therefore equipped with the canonical filtration of an augmented \(\lambda\)-ring. We compare it with the filtration just defined.

\sgaStableTarget{sga6.e05.prop.6-6-1}\textbf{Proposition 6.6.1.} Let \(x\in\Ch A\) be an element of the form \(x=(n,1+\cdots+a_n)\). Then:
\[
\begin{array}{ll}
\text{i)}& \lambda^i(x)=0\quad\text{for }i>n,\\[0.3em]
\text{ii)}& \lambda^n(x)=(1,1+a_1),\\[0.3em]
\text{iii)}& \displaystyle\sum_{i=0}^n(-1)^i\lambda^i(x)
       =(0,1-(n-1)!a_n+\cdots),
\end{array}
\]
where the unwritten terms have degree \(>n\).

By functoriality, it is enough to prove this in the case where
\(A=K[X_1,\ldots,X_n]\), filtered by weight with respect to the \(X_i\), \(X_i\) having weight \(i\), and for the element \((n,1+X_1+\cdots+X_n)\). Embed \(A\) in \(A'=K[T_1,\ldots,T_n]\) by the symmetric functions. Then
\[
 (n,1+X_1+\cdots+X_n)
 =
 \left(n,\prod_{i=1}^n(1+T_i)\right)
 =
 \sum_{i=1}^n(1,1+T_i).
\]
Hence, by \hyperlink{sga6.e05.eq.6-2-1}{(6.2.1)},
\[
  \lambda_t(x)=\prod_{i=1}^n\bigl(1+(1,1+T_i)t\bigr),
\]
so that
\[
  \lambda^k(x)=0\quad(k>n),\qquad
  \lambda^n(x)=\prod_{i=1}^n(1,1+T_i)
   =\left(1,1+\sum_{i=1}^nT_i\right)=(1,1+X_1).
\]
Finally,
\[
  \sum_{i=1}^n(-1)^i\lambda^i(x)=\lambda_{-1}(x)
  =\prod_{i=1}^n\bigl(1-(1,1+T_i)\bigr)
  =(-1)^n\prod_{i=1}^n(0,1+T_i).
\]
By \hyperlink{sga6.e05.eq.6-5-3}{(6.5.3)},
\[
  \prod_{i=1}^n(0,1+T_i)
  =
  \left(0,1+(-1)^{n-1}(n-1)!\prod_{i=1}^nT_i+\cdots\right),
\]
hence
\[
  \lambda_{-1}(x)=(0,1-(n-1)!X_n+\cdots).
\]

\sgaStableTarget{sga6.e05.cor.6-6-2}\textbf{Corollary 6.6.2.} Relation \hyperlink{sga6.e05.prop.6-6-1}{6.6.1(iii)} is true for every element \(X\) of augmentation \(n\).

For this we shall use the following lemma.

\sgaStableTarget{sga6.e05.lem.6-6-3}\textbf{Lemma 6.6.3.} For every \(n\), \((\Ch A)_n\) is a \(\lambda\)-ideal of \(\Ch A\).

We may suppose \(n>1\). Let \(x=(0,\varphi)\in(\Ch A)_n\), and write
\[
  \lambda^i(0,\varphi)=(0,1+Q^i_{i,1}+\cdots+Q^i_{i,j}+\cdots),
\]
where the \(Q^i_{i,j}\) are the isobaric components of weight \(j\) of the \(Q^i_i\) of \hyperlink{sga6.e05.para.3-6}{(3.6)}. For every \(j\), the polynomial \(Q^i_{i,j}\) is a polynomial in the first \(j\) coefficients of \(\varphi\). If \(j<n\), these coefficients are zero; hence the \(Q^i_{i,j}\) are zero, and so the \(\lambda^i(x)\) belong to \((\Ch A)_n\). Thus \((\Ch A)_n\) is a \(\lambda\)-ideal.

Let \(x=(n,1+a_1+\cdots+a_n+\cdots)\) be an element of augmentation \(n\) in \(\Ch A\). Then
\[
  x\equiv (n,1+\cdots+a_n)\pmod{(\Ch A)_{n+1}},
\]
and therefore
\[
  \lambda_{-1}(x)\equiv \lambda_{-1}(n,1+\cdots+a_n)
  =(0,1-(n-1)!a_n+\cdots),
\]
which proves the corollary.

\sgaStableTarget{sga6.e05.cor.6-6-4}\textbf{Corollary 6.6.4.} Let \(\epsilon\) be the augmentation homomorphism of \(\Ch A\). Then, for every \(x\in\Ch A\), one has
\[
  \gamma^n(x-\epsilon(x))
  =
  (0,1+(-1)^{n-1}(n-1)!a_n+\cdots).
\]
By definition,
\[
  \gamma^n(x-\epsilon(x))
  =
  (-1)^n\sum_{i=0}^n(-1)^i\lambda^i(x+n-\epsilon(x)).
\]
Since \(x+n-\epsilon(x)\) has augmentation \(n\), \hyperlink{sga6.e05.cor.6-6-2}{Corollary 6.6.2} applies to it, giving the result.

\sgaStableTarget{sga6.e05.cor.6-6-5}\textbf{Corollary 6.6.5.} With the notation of \hyperlink{sga6.e05.para.3-10}{3.10},
\[
  \Fil^n\Ch A\subset(\Ch A)_n.
\]
This follows from the preceding corollary, which gives \(\gamma^n(x-\epsilon(x))\in(\Ch A)_n\), and from the fact that the two filtrations are ring filtrations and that the elements \(\gamma^n(x-\epsilon(x))\), for \(n\) and \(x\) variable, generate the \(\lambda\)-ring filtration.

\sgaStableTarget{sga6.e05.cor.6-6-7}\textbf{Corollary 6.6.7.} Let \(A\) be a graded \(K\)-algebra of characteristic zero, bounded by degree. Then the two filtrations defined on \(\Ch A\) are identical.

Suppose that \(\Fil^p A=0\), and let \(x=(0,1+a_n+\cdots+a_p)\) be an element of \((\Ch A)_n\). Put
\[
  x'=\left(0,1+\frac{(-1)^{n-1}}{(n-1)!}a_n+a_{n+1}+\cdots+a_p\right).
\]
Then \(\gamma^n(x')=(0,1+a_n+\cdots+a'_p)\), since \(x'\) has zero augmentation. Hence \(x-\gamma^n(x')\in(\Ch A)_{n+1}\). Iterating the operation \(p-n\) times, one reaches \((\Ch A)_{p+1}=0\), and therefore \(x\) is written as an element of \(\Fil^n\Ch A\). In view of \hyperlink{sga6.e05.cor.6-6-5}{(6.6.5)}, the two filtrations are identical.

\sgaStableTarget{sga6.e05.para.6-7}\textbf{6.7.} Now suppose that \(A\) is a \(\lambda\)-ring augmented over \(K\), and let \(\epsilon\) be the augmentation homomorphism. Then \(A\) is equipped with its canonical filtration of augmented \(\lambda\)-ring, and one can associate to \(\Gr A\) its Chern ring.

If \(x\in A\), then \(\gamma_t(x-\epsilon(x))\in\Ghat(A)\) by definition of the filtration of an augmented \(\lambda\)-ring. We write \(c(x)\) for the image of \(\gamma_t(x-\epsilon(x))\) in
\[
  \Ghat(A)/\Ghat'(A)=1+\widehat{\Gr A}^{+}.
\]
It is called the \emph{total Chern class} of \(x\). We write \(c^i(x)\), and call it the \(i\)-th Chern class of \(x\), for the \(i\)-th coefficient of \(c(x)\). Thus
\sgaStableTarget{sga6.e05.eq.6-7-1}\begin{equation}\tag{6.7.1}
  c^i(x)=\cl_{\Gr^iA}\!\left(\gamma^i(x-\epsilon(x))\right),
  \qquad
  c(x)=1+c^1(x)+\cdots .
\end{equation}
Put
\sgaStableTarget{sga6.e05.eq.6-7-2}\begin{equation}\tag{6.7.2}
  \widetilde c(x)=(\epsilon(x),c(x)).
\end{equation}
This is called the \emph{completed Chern class} of \(x\).

\sgaStableTarget{sga6.e05.prop.6-8}\textbf{Proposition 6.8.} The completed Chern class
\[
  \widetilde c:A\longrightarrow\Ch(\Gr A)
\]
is a \(K\)-\(\lambda\)-homomorphism of augmented \(K\)-\(\lambda\)-algebras.

It is clear that the homomorphism so defined is additive. To verify that it is multiplicative, write
\[
\begin{aligned}
\gamma_t(xy-\epsilon(xy))
&=\gamma_t\bigl(\epsilon(y)(x-\epsilon(x))\bigr)
  \gamma_t\bigl(\epsilon(x)(y-\epsilon(y))\bigr)
  \gamma_t\bigl((x-\epsilon(x))(y-\epsilon(y))\bigr)\\
&=\bigl(\gamma_t(x-\epsilon(x))\bigr)^{\epsilon(y)}
  \bigl(\gamma_t(y-\epsilon(y))\bigr)^{\epsilon(x)}
  \bigl(\gamma_t(x-\epsilon(x))\circ\gamma_t(y-\epsilon(y))\bigr),
\end{aligned}
\]
which gives the result by \hyperlink{sga6.e05.para.6-1}{(6.1.2)} after passage to the quotient. If \(a\in K\), then \(\widetilde c(a)=(a,1)\), so \(\widetilde c\) is a homomorphism of \(K\)-algebras; moreover it is compatible with augmentations.

It remains to see that it is a \(\lambda\)-homomorphism. One must prove the commutativity of the diagram
\[
\begin{tikzcd}[column sep=large,row sep=large]
A \arrow[r,"\gamma_t"] \arrow[d,"\widetilde c"'] &
\Ghat(A) \arrow[d,"\Ghat(\widetilde c)"]\\
\Ch A \arrow[r,"\gamma_t"'] &
\Ghat(\Ch A).
\end{tikzcd}
\]
By additivity, it is enough to prove it for the elements of \(K\), where it is evident, and for the elements of zero augmentation in \(A\), where it follows, by the definition of \(c\) and \hyperlink{sga6.e05.para.6-1}{(6.1.4)}, from diagram \hyperlink{sga6.e05.eq.3-7-2}{(3.7.2)}.

\sgaStableTarget{sga6.e05.prop.6-9}\textbf{Proposition 6.9.} Let \(A\) be an augmented \(K\)-\(\lambda\)-algebra and let \(z\in\Fil^mA\). Then
\[
  \gamma^m(z)=(-1)^{m-1}(m-1)!z\pmod{\Fil^{m+1}A},
\]
that is,
\[
  c^m(z)=\cl_{\Gr^mA}\!\left((-1)^{m-1}(m-1)!z\right).
\]

One must see that these two elements give the same image in \(\Gr^mA\), that is, that the \(m\)-th Chern class of \(z\) is the image of \((-1)^{m-1}(m-1)!z\).

One may write \(z\) as a sum
\[
  \sum a\,\gamma^{i_1}(x_1)\cdots\gamma^{i_k}(x_k),
\]
with \(a\in A\), \(\epsilon(x_i)=0\) for all \(i\), and \(i_1+\cdots+i_k\ge m\). We then have
\[
  \widetilde c(z)=
  \sum \widetilde c(a)\,
        \widetilde c(\gamma^{i_1}(x_1))\cdots
        \widetilde c(\gamma^{i_k}(x_k)).
\]
Since \(\widetilde c\) is a \(\lambda\)-homomorphism, one may write
\[
  \widetilde c(\gamma^i(x_j))=\gamma^i(\widetilde c(x_j)).
\]
Taking account of the definition of \(c\) and of \hyperlink{sga6.e05.cor.6-6-4}{(6.6.4)}, one obtains
\[
  \widetilde c(\gamma^{i_j}(x_j))
  =
  \left(0,1+(-1)^{i_j-1}(i_j-1)!\,\gamma^{i_j}(x_j)+\cdots\right).
\]
By \hyperlink{sga6.e05.eq.6-5-3}{(6.5.3)}, and since \(\widetilde c(a)=(a,1)\), one finally gets
\[
\begin{aligned}
\widetilde c(z)
&=\sum
\left(0,\left(1+
(-1)^{i_1+\cdots+i_k-1}
(i_1+\cdots+i_k-1)!
\,\gamma^{i_1}(x_1)\cdots\gamma^{i_k}(x_k)
+\cdots\right)^a\right)\\
&=
\left(0,1+
\sum_{\substack{(x_1,\ldots,x_k)\\ i_1+\cdots+i_k=m}}
(-1)^{m-1}(m-1)!a\,\gamma^{i_1}(x_1)\cdots\gamma^{i_k}(x_k)
+\cdots\right)\\
&=
(0,1+(-1)^{m-1}(m-1)!z+\cdots),
\end{aligned}
\]
which proves the assertion.

\sgaStableTarget{sga6.e05.prop.6-10}\textbf{Proposition 6.10.} Let \(K\) be a binomial ring, \(R\) an augmented \(K\)-\(\lambda\)-algebra, \(N\in R\) an element such that \(\lambda^k(N)=0\) for \(k>d\), and \(x\in\Fil^iR\). Then
\[
  \gamma^{d+i}(N,x)-(-1)^{d+i-1}(d+i-1)!x\in\Fil^{i+1}R.
\]

One can find elements \(z_{\alpha,j}\in R\) and \(a_\alpha\in K\) such that
\[
  x=\sum_\alpha a_\alpha\,
       \gamma^{i_1}(x_{\alpha,1})\cdots
       \gamma^{i_k}(x_{\alpha,k}),
\]
with \(i_1+\cdots+i_k\ge i\) and \(\epsilon(x_{\alpha,j})=0\) for every \(\alpha\) and every \(j\). Let \(R_0\) be the \(K\)-\(\lambda\)-algebra generated by \(N'_0\) and \(X_{\alpha,j}\), subject to the relations
\[
  \gamma^k(N'_0)=0\quad\text{for }k>d.
\]
Define an augmentation homomorphism \(R_0\to K\) by \(\epsilon(N'_0)=\epsilon(X_{\alpha,j})=0\). By \hyperlink{sga6.e05.cor.4-9}{(4.9.1)}, \(R_0\) is the ring of polynomials with coefficients in \(K\) in the indeterminates \(\gamma^k(N'_0)\), \(1\le k\le d\), and \(\gamma^p(X_{\alpha,j})\), \(p\ge1\); the \(\lambda\)-filtration of \(R_0\) is its weight filtration, where \(\gamma^k(N'_0)\) and \(\gamma^k(X_{\alpha,j})\) have weight \(k\) (\hyperlink{sga6.e05.prop.4-12}{4.12} and \hyperlink{sga6.e05.cor.4-15}{4.15}).

Put
\[
  X_0=\sum_\alpha a_\alpha\,
       \gamma^{i_1}(X_{\alpha,1})\cdots
       \gamma^{i_k}(X_{\alpha,k}).
\]
We shall prove \hyperlink{sga6.e05.prop.6-10}{6.10} first for \(X_0\) and \(N'_0=N_0+d\), which do satisfy the hypotheses of 6.10. In view of the preceding remark about the filtration of \(R_0\), it is enough to show that
\[
  \gamma^{d+i}(N_0,X_0)\gamma^d(N'_0)
  -
  (-1)^{d+i-1}(d+i-1)!X_0\gamma^d(N'_0)
  \in\Fil^{d+i+1}R_0 .
\]
But \(\lambda_{-1}(N_0)=(-1)^d\gamma^d(N'_0)\). Thus this relation becomes
\[
  (-1)^d\gamma^{d+i}(X_0,(-1)^d\gamma^d(N'_0))
  -
  (-1)^{d+i-1}(d+i-1)!X_0\gamma^d(N'_0)
  \in\Fil^{d+i+1}R_0 .
\]
Putting \(z=(-1)^dX_0\gamma^d(N'_0)\), we have \(z\in\Fil^{d+i}R_0\), and the statement is reduced to
\[
  \gamma^{d+i}(z)-(-1)^{d+i-1}(d+i-1)!z\in\Fil^{d+i+1}R_0,
\]
which is just 6.9.

To deduce \hyperlink{sga6.e05.prop.6-10}{6.10} from this particular case, consider the \(K\)-\(\lambda\)-homomorphism \(R_0\to R\) which sends \(N'_0\) to \(N-d\) and \(X_{\alpha,j}\) to \(x_{\alpha,j}\), for every \(\alpha\) and \(j\). This homomorphism is compatible with the \(\lambda\)-filtrations and sends \(N_0\) to \(N\) and \(X_0\) to \(x\). Hence the result follows.

\sgaStableTarget{sga6.e05.prop.6-11}\textbf{Proposition 6.11.} Suppose that \(K\) is of characteristic zero, and let \(n\ge0\) be an integer. Then the functor \(\Ch\) is an equivalence from the category of graded \(K\)-algebras \(A\) such that \(A_0=K\) and \(A_k=0\) for \(k>n\), to the category of augmented \(K\)-\(\lambda\)-algebras \(\Lambda\) such that \(\Fil^{n+1}\Lambda=0\). A quasi-inverse functor is given by \(\Lambda\mapsto\Gr\Lambda\).

Let \(A\) be a graded \(K\)-algebra such that \(A_k=0\) for \(k>n\), equipped with the filtration associated to its grading. Then \((\Ch A)_{n+1}=0\), and hence \(\Fil^{n+1}\Ch A=0\) by \hyperlink{sga6.e05.cor.6-6-5}{(6.6.5)}. Conversely, if \(\Lambda\) is an augmented \(K\)-\(\lambda\)-algebra such that \(\Fil^{n+1}\Lambda=0\), then \(\Gr^0\Lambda=K\) and \(\Gr^k\Lambda=0\) for \(k>n\).

We shall define isomorphisms of functors
\[
  A\longrightarrow \Gr(\Ch A),
  \qquad
  \Lambda\longrightarrow\Ch(\Gr\Lambda).
\]
Since \(A\) is of finite grading, the associated gradeds of the \(\lambda\)-ring filtration on \(\Ch A\) and of the degree filtration are the same by \hyperlink{sga6.e05.cor.6-6-7}{(6.6.7)}. Define
\[
  u:A\longrightarrow \Gr_{\deg}(\Ch A)
\]
to be the identity on \(A_0\), and, for \(i\ge1\) and \(a_i\in A_i\),
\[
  u(a_i)=
  \cl_{\Gr^i}\left(0,1+(-1)^{i-1}(i-1)!a_i\right),
\]
then extend by additivity. By relation \hyperlink{sga6.e05.eq.6-5-3}{(6.5.3)} this is a homomorphism of graded rings, and it is easy to see that it is compatible with the operations of \(K\). The inverse morphism \(u^{-1}\) may be defined as follows: for \(i\ge1\) and \(\bar\varphi\in\Gr^i(\Ch A)\), the \(i\)-th coefficient \(a_i\) of \(\varphi\) is independent of the lift of \(\bar\varphi\) in \((\Ch A)_i\), and one puts
\[
  u^{-1}(\bar\varphi)=\frac{(-1)^{i-1}}{(i-1)!}a_i.
\]
Thus \(u\) is an isomorphism.

The isomorphism \(\Lambda\to\Ch(\Gr\Lambda)\) is defined by the completed Chern class \(\widetilde c\). To verify that it really defines an isomorphism, it is enough to verify that the homomorphism obtained by passing to associated graded rings is one, the filtrations being discrete. One gets
\[
  \Gr\Lambda
   \xrightarrow{\Gr(\widetilde c)}
  \Gr_{\lambda}(\Ch(\Gr\Lambda))
   \simeq
  \Gr_{\deg}(\Ch(\Gr\Lambda))
   \xrightarrow{u^{-1}}
  \Gr\Lambda .
\]
It remains to show that the composite is the identity. Since the homomorphisms are homomorphisms of graded \(K\)-algebras, it is enough, by the definition of the \(\lambda\)-ring filtration of \(\Lambda\), to show that it is the identity on the \(\gamma^i(x)\), where \(\epsilon(x)=0\). One then obtains
\[
  \Gr(\widetilde c)\bigl(\gamma^i(x)\bigr)
  =
  \cl_{\Gr^i}\bigl(\gamma^i(c(x))\bigr).
\]
But \(\gamma^i(\widetilde c(x))=(0,1+(-1)^{i-1}(i-1)!c^i(x)+\cdots)\), by \hyperlink{sga6.e05.cor.6-6-4}{(6.6.4)}. Therefore
\[
  \cl_{\Gr^i}\bigl(\gamma^i(\widetilde c(x))\bigr)=u(c^i(x)),
\]
which shows that \(u^{-1}\circ\Gr(\widetilde c)=\id\).

\sgaStableTarget{sga6.e05.cor.6-12}\textbf{Corollary 6.12.} If \(K\) is of characteristic zero, and if \(\Lambda\) is a \(K\)-\(\lambda\)-algebra with discrete filtration, then the completed Chern class
\[
  \widetilde c:\Lambda\longrightarrow \Ch(\Gr\Lambda)
\]
is an isomorphism.

The proof is contained in that of 6.11.

\sgaStableTarget{sga6.e05.def.6-13}\textbf{Definition 6.13.} Let \(A\) be a graded \(K\)-algebra and \(\Lambda\) an augmented \(K\)-\(\lambda\)-algebra. A \emph{Chern theory for \(\Lambda\) with values in \(A\)} is a \(\lambda\)-homomorphism
\[
  \Lambda\longrightarrow \Ch(A)
\]
compatible with augmentations.

Thus one defines a functor on the category of graded \(K\)-algebras, with values in the category of sets, by associating with a graded \(K\)-algebra \(A\) the set of Chern theories of \(\Lambda\) with values in \(A\), namely
\[
  \Hom_\lambda(\Lambda,\Ch(A)).
\]

\sgaStableTarget{sga6.e05.prop.6-14}\textbf{Proposition 6.14.} The functor defined in \hyperlink{sga6.e05.def.6-13}{6.13} is representable. If \(K\) is of characteristic zero and \(\Lambda\) has discrete filtration, it is represented by the graded ring associated with \(\Lambda\).

First prove the first assertion. Consider the graded \(K\)-algebra generated by generators \(c^i(x)\), with \(i\ge1\), \(x\in\Fil^1\Lambda\), \(c^i(x)\) of degree \(i\), and take its quotient by the following relations:
\[
\begin{array}{ll}
\text{i)}&
 c^k(x+y)=\displaystyle\sum_{i+j=k}c^i(x)c^j(y),
 \quad k\ge1,\ x,y\in\Fil^1\Lambda,\\[0.8em]
\text{ii)}&
 c^k(xy)=P''_k(c^i(x),c^j(y))
 \quad(\text{cf. }\hyperlink{sga6.e05.para.6-1}{6.1}),
 \quad k\ge1,\ x,y\in\Fil^1\Lambda,\\[0.4em]
\text{iii)}&
 c^k(ax)=R_{k,a}(c^i(x))
 \quad(\text{cf. }\hyperlink{sga6.e05.para.6-1}{6.1}),
 \quad k\ge1,\ a\in K,\ x\in\Fil^1\Lambda,\\[0.4em]
\text{iv)}&
 c^k(\gamma^j(x))=Q''_{j,k}(c^i(x))
 \quad(\text{cf. }\hyperlink{sga6.e05.para.6-1}{6.1}),
 \quad j,k\ge1,\ x\in\Fil^1\Lambda.
\end{array}
\]
These relations are isobaric, so the quotient is a graded \(K\)-algebra \(\Omega\). Define a morphism
\[
  \widetilde c_0:\Lambda\longrightarrow \Ch(\Omega)
\]
by
\[
  x\longmapsto
  \bigl(\epsilon(x),1+c^1(x-\epsilon(x))+\cdots\bigr).
\]
The imposed relations show that this is a Chern theory for \(\Lambda\) with values in \(\Omega\). Conversely, if \(f\) is a Chern theory with values in a graded \(K\)-algebra \(A\), then there is a unique homomorphism \(\Omega\to A\) through which \(f\) factors through \(\Ch(\Omega)\): it sends, for every \(i\ge1\) and every \(x\in\Fil^1\Lambda\), the element \(c^i(x)\) to the \(i\)-th coefficient of \(f(x)\). This is well defined because \(f\) is a \(\lambda\)-homomorphism.

Now suppose that \(K\) is of characteristic zero and \(\Lambda\) has discrete filtration. Let \(n\) be such that \(\Fil^n\Lambda=0\). Then the graded algebra \(\Omega\) is zero in degrees \(>n\). To show this, it is enough to prove that, for
\[
  x_1,\ldots,x_p\in\Fil^1\Lambda
  \quad\text{and}\quad
  k_1+\cdots+k_p\ge n,
\]
one has
\[
  c^{k_1}(x_1)\cdots c^{k_p}(x_p)=0.
\]
By hypothesis,
\[
  \gamma^{k_1}(x_1)\cdots\gamma^{k_p}(x_p)=0.
\]
Applying the \(\lambda\)-homomorphism \(\widetilde c_0\), one gets
\[
  \gamma^{k_1}(\widetilde c_0(x_1))\cdots
  \gamma^{k_p}(\widetilde c_0(x_p))=(0,1).
\]
By \hyperlink{sga6.e05.cor.6-6-4}{(6.6.4)} and \hyperlink{sga6.e05.eq.6-5-3}{(6.5.3)}, this implies
\[
  (k_1+\cdots+k_p)!\,
  c^{k_1}(x_1)\cdots c^{k_p}(x_p)=0.
\]
Since \(K\) is of characteristic zero, the desired result follows.

The universal property of \(\Omega\), applied to the canonical homomorphism \(\widetilde c\), gives a commutative diagram
\[
\begin{tikzcd}[column sep=large,row sep=large]
\Lambda \arrow[r,"\widetilde c"] \arrow[d,"\widetilde c_0"'] &
\Ch(\Gr\Lambda)\\
\Ch(\Omega) \arrow[ur] &
\end{tikzcd}
\]
By \hyperlink{sga6.e05.prop.6-11}{Proposition 6.11}, the homomorphism \(\Omega\to\Gr\Lambda\) is an isomorphism if and only if the homomorphism
\[
  \Ch(\Omega)\longrightarrow \Ch(\Gr\Lambda)
\]
is one. Since it is visibly surjective, the latter is also surjective; it remains to show that
\(\Ch(\Omega)\to\Ch(\Gr\Lambda)\) is injective. By functoriality of \(\Ch\), one obtains the diagram
\[
\begin{tikzcd}[column sep=large,row sep=large]
\Lambda \arrow[r,"\widetilde c_0"] \arrow[d,"\widetilde c"'] &
\Ch(\Omega) \arrow[d,"\widetilde c"]\\
\Ch(\Gr\Lambda) \arrow[r,"\Ch(\widetilde c_0)"'] &
\Ch\bigl(\Gr_\Lambda(\Ch(\Omega))\bigr).
\end{tikzcd}
\]
Since \(\Omega\) has bounded degrees, the right vertical arrow is an isomorphism by \hyperlink{sga6.e05.cor.6-12}{(6.12)}. Using \hyperlink{sga6.e05.prop.6-11}{6.11}, one therefore obtains a morphism
\[
  \Gr\Lambda\longrightarrow\Omega
\]
whose two composites with \(\Omega\to\Gr\Lambda\) are the identity, by the universal property of \(\Omega\).

\sgaStableTarget{sga6.e05.sec.7}\section*{7. Appendix: the Adams operations \(\psi^k\)}

We follow here a presentation due to J.-P. Serre.

\sgaStableTarget{sga6.e05.para.7-1}\textbf{7.1.} Let \(K\) be a pre-\(\lambda\)-ring. For \(x\in K\), write
\[
  \frac{d}{dt}(\lambda_t(x))
\]
for the derivative with respect to \(t\) of \(\lambda_t(x)\). Define elements of \(K\), denoted \(\psi^k(x)\), by putting
\sgaStableTarget{sga6.e05.eq.7-1-1}\begin{equation}\tag{7.1.1}
  \frac{d}{dt}(\lambda_t(x))/\lambda_t(x)
  =
  \sum_{k=1}^{\infty}(-1)^{k-1}\psi^k(x)t^{k-1}.
\end{equation}
The maps \(\psi^k:K\to K\) are additive. Indeed,
\[
\begin{aligned}
 \frac{d}{dt}(\lambda_t(x+y))/\lambda_t(x+y)
 &=
 \frac{d}{dt}(\lambda_t(x)\lambda_t(y))/\lambda_t(x)\lambda_t(y)\\
 &=
 \frac{d}{dt}(\lambda_t(x))/\lambda_t(x)
 +\frac{d}{dt}(\lambda_t(y))/\lambda_t(y),
\end{aligned}
\]
and therefore, for every \(k\),
\[
  \psi^k(x+y)=\psi^k(x)+\psi^k(y).
\]
By \hyperlink{sga6.e05.eq.7-1-1}{(7.1.1)}, the \(\psi^k\) are isobaric polynomials of weight \(k\) in the \(\lambda^i\), \(i=1,\ldots,k\). For example,
\[
  \psi^1(x)=\lambda^1(x)=x,
  \qquad
  \psi^2(x)=x^2-2\lambda^2(x).
\]

\sgaStableTarget{sga6.e05.para.7-2}\textbf{7.2.} Let \(K\) be a ring of characteristic zero. Suppose given additive maps \(\psi^k:K\to K\) such that \(\psi^1=\id_K\). Define a pre-\(\lambda\)-structure on \(K\) by
\sgaStableTarget{sga6.e05.eq.7-2-1}\begin{equation}\tag{7.2.1}
  \lambda_t(x)=
  \exp\left(\sum_{k=1}^{\infty}(-1)^{k-1}\psi^k(x)t^k/k\right).
\end{equation}
Multiplicativity follows from the additivity of the \(\psi^k\). Thus there is an equivalence between giving a pre-\(\lambda\)-structure on \(K\) and giving the corresponding operations \(\psi^k\).

\sgaStableTarget{sga6.e05.para.7-3}\textbf{7.3.} It is clear from formula \hyperlink{sga6.e05.eq.7-1-1}{(7.1.1)} that the \(\psi^k\) commute with \(\lambda\)-homomorphisms. If \(K,K'\) are pre-\(\lambda\)-rings of characteristic zero, it is necessary and sufficient for a homomorphism \(f:K\to K'\) to be a \(\lambda\)-homomorphism that it commute with the \(\psi^k\). This follows from \hyperlink{sga6.e05.para.7-2}{7.2}, since the \(\lambda^i\) are expressed as polynomials in the \(\psi^k\).

\sgaStableTarget{sga6.e05.ex.7-4}\textbf{Examples 7.4.} i) For \(K=\mathbb Z\), with its \(\lambda\)-ring structure, one has
\[
  \frac{d}{dt}(\lambda_t(n))/\lambda_t(n)
  =
  n(1+t)^{n-1}/(1+t)^n
  =
  n\sum_{k=1}^{\infty}(-1)^{k-1}t^{k-1}.
\]
Consequently the \(\psi^k\) are all equal to the identity map of \(\mathbb Z\).

ii) Work in the category \(\mathscr C_0\) of \(\mathbb Z\)-algebras \hyperlink{sga6.e05.para.1-8}{(1.8)}. By \hyperlink{sga6.e05.para.7-3}{7.3}, the operations \(\psi^k\) define additive endomorphisms of the functor \(G\) on \(\lambda\)-rings. By \hyperlink{sga6.e05.thm.1-10}{1.10} they are characterized by their value on \(1+Xt\). One has
\[
  \frac{d}{dt}(\lambda_u(1+Xt))/\lambda_u(1+Xt)
  =
  \frac{1+Xt}{1+\{1+Xt\}u}
  =
  \sum_{k=1}^{\infty}(-1)^{k-1}\{1+Xt\}^{\circ k}u^{k-1}.
\]
Thus
\sgaStableTarget{sga6.e05.eq.7-4-1}\begin{equation}\tag{7.4.1}
  \psi^k(1+Xt)=1+X^k t.
\end{equation}

\sgaStableTarget{sga6.e05.prop.7-5}\textbf{Proposition 7.5.} i) If \(K\) is a \(\lambda\)-ring, the operations \(\psi^k\) satisfy the identities
\sgaStableTarget{sga6.e05.eq.7-5-1}\begin{equation}\tag{7.5.1}
  \psi^k(1)=1,\qquad
  \psi^k(xy)=\psi^k(x)\psi^k(y),\qquad
  \psi^k\circ\psi^{k'}=\psi^{kk'}.
\end{equation}
ii) Conversely, if the \(\psi^k\) satisfy these identities and if \(K\) is of characteristic zero, then \(K\) is a \(\lambda\)-ring.

For (i), if \(K\) is a \(\lambda\)-ring, \(\lambda_t(1)=1+t\), hence
\[
  \frac{d}{dt}(\lambda_t(1))/\lambda_t(1)
  =
  \frac1{1+t}
  =
  \sum_{k=1}^{\infty}(-1)^{k-1}t^{k-1},
\]
and therefore \(\psi^k(1)=1\).

To prove the second relation of \hyperlink{sga6.e05.eq.7-5-1}{(7.5.1)}, consider the diagram
\sgaStableTarget{sga6.e05.eq.7-5-2}\begin{equation}\tag{7.5.2}
\begin{tikzcd}[column sep=large,row sep=large]
K\times K \arrow[r,"\lambda_t\times\lambda_t"] \arrow[d] &
\widehat G(K)\times\widehat G(K) \arrow[r,"g^k\times g^k"] \arrow[d,"\circ"] &
K\times K \arrow[d]\\
K \arrow[r,"\lambda_t"'] &
\widehat G(K) \arrow[r,"g^k"'] &
K ,
\end{tikzcd}
\end{equation}
where the vertical arrows are defined by multiplication, and the maps \(g^k\) as follows: the logarithmic derivative is a group homomorphism from \(G(K)\) to \(K[[t]]\), and the map which to a series assigns its \((k-1)\)-st coefficient, multiplied by \((-1)^{k-1}\), is an additive homomorphism from \(K[[t]]\) to \(K\); \(g^k\) is their composite. Clearly \(g^k\circ\lambda_t=\psi^k\) by \hyperlink{sga6.e05.eq.7-1-1}{definition (7.1.1)}.

The second of formulas \hyperlink{sga6.e05.eq.7-5-1}{(7.5.1)} therefore expresses the commutativity of the diagram just written. If \(K\) is a \(\lambda\)-ring, the left square is commutative, and it is enough to verify the commutativity of the right square. This is a diagram of functor homomorphisms on \(\mathscr C_0\), with source \(\widehat G\times\widehat G\). By \hyperlink{sga6.e05.thm.1-10}{1.10} it is enough to verify commutativity on \((1+Xt,1+Yt)\). Then
\[
  g^k(1+Xt)=X^k,\qquad
  g^k(1+Yt)=Y^k,
\]
and
\[
  g^k\bigl((1+Xt)\circ(1+Yt)\bigr)
  =
  g^k(1+XYt)
  =
  (XY)^k
  =
  g^k(1+Xt)g^k(1+Yt),
\]
which gives the result.

To prove the last relation, first observe that the diagram
\sgaStableTarget{sga6.e05.eq.7-5-3}\begin{equation}\tag{7.5.3}
\begin{tikzcd}[column sep=large,row sep=large]
\widehat G_t(K) \arrow[r,"\lambda_u"] \arrow[d,"g^k"'] &
\widehat G_u\widehat G_t(K) \arrow[d,"g^k_u"]\\
K \arrow[r,"\lambda_u"'] &
\widehat G^u(K)
\end{tikzcd}
\end{equation}
is commutative. As above, it is enough to see this for \(1+Xt\), and one has
\[
  g^{k'}_u\lambda_u(1+Xt)
  =
  g^{k'}_u(1+\{1+Xt\}u)
  =
  (1+Xt)^{\circ k'}
  =
  1+X^{k'}t,
\]
and then
\[
  g^k(1+X^{k'}t)=X^{kk'}=g^k(1+Xt).
\]
Consequently
\[
  \psi^k\circ\psi^{k'}
  =
  g^k\circ g^{k'}_u\circ\lambda_u\circ\lambda_t
  =
  g^k\circ\lambda_t\circ\psi^{k'}.
\]
Moreover
\[
  \psi^k\circ\psi^{k'}
  =
  g^k\circ\lambda_t\circ\psi^{k'}
  =
  \psi^k\circ\lambda_t\circ\psi^{k'}.
\]
If \(K\) is a \(\lambda\)-ring, the diagram
\sgaStableTarget{sga6.e05.eq.7-5-4}\begin{equation}\tag{7.5.4}
\begin{tikzcd}[column sep=large,row sep=large]
K \arrow[r,"\lambda_t"] \arrow[d,"\psi^{k'}"'] &
\widehat G(K) \arrow[d,"\psi^{k'}"]\\
K \arrow[r,"\lambda_t"'] &
\widehat G(K)
\end{tikzcd}
\end{equation}
is commutative by \hyperlink{sga6.e05.para.7-3}{7.3}, since \(\lambda_t\) is a \(\lambda\)-homomorphism. The desired relations follow at once.

For (ii), if \(K\) is of characteristic zero, the logarithmic derivative is an isomorphism from \(\widehat G(K)\) onto \(K[[t]]\), since one can integrate. Therefore two elements of \(\widehat G(K)\) are equal if and only if, for every \(k\), they have the same image by \(g^k\).

To show that \(\lambda_t\) is a ring homomorphism, it is enough to show that the left square in \hyperlink{sga6.e05.eq.7-5-2}{(7.5.2)} is commutative. By the preceding remark, it is enough to show that, for every \(k\), diagram \hyperlink{sga6.e05.eq.7-5-2}{(7.5.2)} is commutative; the right square always is. But this is exactly the relation
\[
  \psi^k(xy)=\psi^k(x)\psi^k(y).
\]
To show that \(\lambda_t\) is a \(\lambda\)-homomorphism, it is enough, by \hyperlink{sga6.e05.para.7-3}{7.3}, to show that diagram \hyperlink{sga6.e05.eq.7-5-4}{(7.5.4)} is commutative for every \(k'\). By the preceding remark, this is equivalent to showing, for all \(k,k'\), that
\[
  g^k\circ\psi^{k'}\circ\lambda_t
  =
  g^k\circ\lambda_t\circ\psi^{k'},
\]
that is, the relations \(\psi^{kk'}=\psi^k\circ\psi^{k'}\). Finally, the relation \(\lambda_t(1)=1+t\) is obtained by integration from the relations \(\psi^k(1)=1\).

\sgaStableTarget{sga6.e05.bibliography}\subsection*{Bibliography}
\begin{enumerate}[label={[\arabic*]}]
\item A. Grothendieck, \emph{La théorie des classes de Chern}, Bull. Soc. math. France, 86, 1958, p. 137 à 154.
\item F. Hirzebruch, \emph{Neue topologische Methoden in der algebraischen Geometrie}.
\end{enumerate}



\providecommand{\Sym}{\operatorname{Sym}}
\providecommand{\Proj}{\operatorname{Proj}}
\providecommand{\Projz}{\operatorname{Proj}_0}
\providecommand{\Spec}{\operatorname{Spec}}
\providecommand{\Comp}{\operatorname{Comp}}
\providecommand{\Dplus}{D_{+}}
\providecommand{\id}{\operatorname{id}}
\providecommand{\cG}{\mathcal G}
\providecommand{\cA}{\mathcal A}
\providecommand{\cB}{\mathcal B}
\providecommand{\cL}{\mathscr L}
\providecommand{\cM}{\mathscr M}
\providecommand{\cN}{\mathscr N}
\providecommand{\cE}{\mathscr E}
\providecommand{\cK}{\mathscr K}
\providecommand{\cP}{\mathscr P}
\providecommand{\cQ}{\mathscr Q}
\providecommand{\OO}{\mathcal O}
\providecommand{\rk}{\operatorname{rk}}
\providecommand{\colim}{\operatorname*{colim}}


\sgaStableTarget{sga6.e06.expose}\section*{Exposé VI. The \(K^\bullet\) of a projective bundle: computations and consequences}
\begin{center}
Pierre Berthelot
\end{center}

\sgaStableTarget{sga6.e06.sec.1}\section*{1. Computation of \(K^\bullet\) of a projective bundle: the case of locally free sheaves of finite type}

In this section \(K^\bullet(X)\) denotes the Grothendieck group of locally free sheaves of finite type on a scheme \(X\). The aim is to prove the following structure theorem.

\sgaStableTarget{sga6.e06.thm.1-1}\textbf{Theorem 1.1.} Let \(S\) be a quasi-compact scheme, let \(E\) be a locally free \(\mathcal O_S\)-module of rank \(r+1\), and let
\[
  X=\mathbf P(E)
\]
be the associated projective bundle. The structural morphism \(f:X\to S\) gives \(K^\bullet(X)\) the structure of a \(K^\bullet(S)\)-algebra. If \(\xi\in K^\bullet(X)\) denotes the class of \(\mathcal O_X(1)\), then \(K^\bullet(X)\) is a free \(K^\bullet(S)\)-module with basis
\[
  1,\xi,\ldots,\xi^r .
\]

\sgaStableTarget{sga6.e06.sec.1.2}\subsection*{1.2. Plan of the proof}
The proof is divided into three parts.

First, one proves that the powers of \(\xi\) generate \(K^\bullet(X)\) as a \(K^\bullet(S)\)-module. For this it is enough to establish the following facts. If \(F\) is locally free of finite type on \(X\), one can find a graded \(\Sym(E)\)-module \(M\), of finite presentation and flat over \(\mathcal O_S\), such that \(F=\Proj_0(M)\). Such an \(M\) admits a finite graded resolution
\[
  0\longrightarrow L_{r+1}\longrightarrow \cdots \longrightarrow L_0\longrightarrow M\longrightarrow 0
\]
by graded \(\Sym(E)\)-modules locally free of finite type. Finally, each \(L_i\) has a finite filtration whose quotients are of the form
\[
  N\otimes_{\mathcal O_S}\Sym(E)(n),
\]
with \(N\) locally free of finite type on \(S\). Applying \(\Proj_0\) then gives the required generation by the powers of \(\xi\).

Second, one proves that \(1,\xi,\ldots,\xi^{r+1}\) are linearly dependent over \(K^\bullet(S)\). This is obtained from the canonical epimorphism \(f^*E\to\mathcal O_X(1)\).

Third, one proves that \(1,\xi,\ldots,\xi^r\) are linearly independent over \(K^\bullet(S)\).

Let \(S\) be a scheme, let \(B=\bigoplus_{n\geq 0}B_n\) be a graded \(\mathcal O_S\)-algebra, and let \(M\) be a graded \(B\)-module. If \(N\) is a graded \(B\)-module, \(N(n)\) denotes the graded module defined by
\[
  N(n)_k=N_{n+k}.
\]
A graded \(B\)-module is said to be free, respectively locally free, respectively locally free of finite type, if locally on \(S\) it is a direct sum, respectively locally a direct sum, respectively locally a finite direct sum, of modules \(B(n)\). It is said to be of finite presentation if, locally on \(S\), it admits a presentation by free graded modules of finite type.

For every integer \(n\) there is a canonical homomorphism of graded \(B\)-modules
\sgaStableTarget{sga6.e06.eq.1-2-1}\begin{equation}\tag{1.2.1}
  M_n\otimes_{\mathcal O_S}B(-n)\longrightarrow M,
\end{equation}
defined by \(x\otimes a\mapsto ax\). A free graded module of finite type admits a decomposition
\sgaStableTarget{sga6.e06.eq.1-2-2}\begin{equation}\tag{1.2.2}
  M\simeq \bigoplus_{n=n_0}^{n_1} B(-n)^{k_n}.
\end{equation}

\sgaStableTarget{sga6.e06.prop.1-2}\textbf{Proposition 1.2.} The integers \(k_n\) occurring in a decomposition \hyperlink{sga6.e06.eq.1-2-2}{(1.2.2)} are independent of the chosen decomposition.

Indeed, after restricting to an open set where \(B_0=\mathcal O_S\), one compares the first non-zero homogeneous component. This determines \(k_{n_0}\); passing to the quotient and arguing by induction gives all the \(k_n\).

\sgaStableTarget{sga6.e06.prop.1-3}\textbf{Proposition 1.3.} Let \(S\) be quasi-compact, let \(E\) be a locally free \(\mathcal O_S\)-module of finite type, let \(X=\mathbf P(E)\), and let \(F\) be an \(\mathcal O_X\)-module of finite presentation which is flat over \(S\). Then there exists an integer \(n_0\) such that, for all \(n\geq n_0\),
\[
\begin{array}{ll}
\text{\rm (i)}& R^i f_*(F(n))=0\quad (i\geq 1),\\[2mm]
\text{\rm (ii)}& f_*(F(n))\text{ is locally free of finite type on }S,\\[2mm]
\text{\rm (iii)}& \displaystyle M=\bigoplus_{n\geq n_0} f_*(F(n))
\text{ is a graded }\Sym(E)\text{-module of finite presentation.}
\end{array}
\]

The statement is local on \(S\). Since \(S\) is quasi-compact, one is reduced to the affine case. If \(S\) is noetherian, (i) is Serre's theorem, (ii) follows from EGA III, 7.9.10, and (iii) from the usual comparison between a coherent sheaf on a projective bundle and its associated graded module.

In the general affine case one uses passage to the limit. Write \(S=\Spec(A)\) as a filtered projective limit of noetherian affine schemes \(S_\lambda=\Spec(A_\lambda)\), with affine transition morphisms. The locally free module \(E\), the scheme \(X\), and the sheaf \(F\) descend, after increasing \(\lambda\), to corresponding data \(E_\lambda,X_\lambda,F_\lambda\), with \(F_\lambda\) flat over \(S_\lambda\). The noetherian case applies to \(F_\lambda\), and the theorem of change of base gives the assertions after pullback to \(S\).

In this reduction one uses the Cartesian square
\[
\begin{tikzcd}[column sep=large,row sep=large]
X \arrow[r,"g"] \arrow[d,"f"'] & X_0 \arrow[d,"f_0"]\\
S \arrow[r,"g_0"'] & S_0 .
\end{tikzcd}
\]
Here \(S_0\) is noetherian, \(X_0=\mathbf P(E_0)\), and \(F=g^*(F_0)\). The previous paragraph proves the assertion for \(F_0\), and EGA III, 7.7.5 and 7.8.5 give the required base-change identifications for \(F\).

\sgaStableTarget{sga6.e06.cor.1-4}\textbf{Corollary 1.4.} Under the hypotheses of \hyperlink{sga6.e06.prop.1-3}{Proposition 1.3} there exists a graded \(\Sym(E)\)-module \(M\), of finite presentation and flat over \(\mathcal O_S\), such that
\[
  F=\Proj_0(M).
\]

Choose an integer \(n_0\) for which the properties of \hyperlink{sga6.e06.prop.1-3}{Proposition 1.3} hold. For
\[
 M=\bigoplus_{n\ge n_0}f_*(F(n)),
\]
we have seen that \(M\) is of finite presentation; it is flat over \(S\), since the \(f_*(F(n))\), for \(n\ge n_0\), are locally free on \(S\). Moreover, \(F=\Proj_0(M)\) by EGA II 3.4. We now turn to part b) of the proof.

\sgaStableTarget{sga6.e06.lem.1-5-1}\textbf{Lemma 1.5.1.} Let \(A\) be a ring, let \(B\) be a positively graded \(A\)-algebra with \(B_0=A\), and let \(M\) be a graded \(B\)-module. Suppose that \(M\), regarded as an ungraded \(B\)-module, is projective of finite type. Let \(\widetilde B\) be the sheaf of algebras on \(\Spec(A)\) defined by \(B\), and let \(\widetilde M\) be the associated graded \(\widetilde B\)-module. Then, locally on \(\Spec(A)\), \(\widetilde M\) is free of finite type over \(\widetilde B\) as a graded module \hyperlink{sga6.e06.prop.1-2}{(1.2)}.

Indeed, since \(M\) is of finite type, its degrees are bounded below. Further, \(M\otimes_B A\) is projective of finite type over \(A\), and therefore has only finitely many non-zero homogeneous components, each projective of finite type over \(A\). After localizing on \(\Spec(A)\), one may thus suppose that \(M\otimes_B A\) is free of finite type over \(A\). Finally, projectivity of \(M\) gives
\[
 \Tor_1^B(M,A)=0.
\]
These properties imply that \(M\) is free of finite type over \(B\) as a graded module (Bourbaki, \emph{Algèbre}, ch.~II, \S~11, prop.~7), and hence \(\widetilde M\) is free of finite type over \(\widetilde B\), as required.

\sgaStableTarget{sga6.e06.lem.1-5-2}\textbf{Lemma 1.5.2.} Let \(A\) be a ring, let \(B\) be an \(A\)-algebra of finite presentation, and let \(M\) be a \(B\)-module of finite presentation. Suppose that \(B\) and \(M\) are flat over \(A\), and that, for every residue field \(k\) of \(A\), the \(B\otimes_A k\)-module \(M\otimes_A k\) has projective dimension at most \(n\) (as is the case if \(B=A[t_0,\ldots,t_r]\)). Then \(M\) has a resolution of length \(n\) by projective \(B\)-modules of finite type:
\[
  0\longrightarrow L_n\longrightarrow\cdots\longrightarrow L_0\longrightarrow M\longrightarrow 0.
\]
This is the local criterion of EGA IV, 12.3.2.

\sgaStableTarget{sga6.e06.prop.1-6}\textbf{Proposition 1.6.} Let \(S\) be quasi-compact, let \(E\) be locally free of rank \(r+1\), let \(B=\Sym(E)\) with its usual grading, and let \(M\) be a graded \(B\)-module of finite presentation, flat over \(S\). Then \(M\) admits a graded resolution
\[
  0\longrightarrow L_{r+1}\longrightarrow\cdots\longrightarrow L_0\longrightarrow M\longrightarrow 0,
\]
where the \(L_i\) are graded \(B\)-modules locally free of finite type.

First, \(M\) admits a graded epimorphism \(L\to M\) with \(L\) locally free of finite type: choose an integer \(N\) such that \(M\) is generated in degrees \(\leq N\), and take
\[
  L=\bigoplus_{i\leq N} M_i\otimes_{\mathcal O_S} B(-i),
\]
with the morphism \hyperlink{sga6.e06.eq.1-2-1}{(1.2.1)}. Repeating the construction for successive kernels gives a global graded resolution. After \(r+1\) steps, \hyperlink{sga6.e06.lem.1-5-2}{Lemma 1.5.2} shows locally on affine opens that the final kernel is projective over \(B\), and \hyperlink{sga6.e06.lem.1-5-1}{Lemma 1.5.1} then shows that it is locally free as a graded \(B\)-module; truncation gives the asserted resolution.

\sgaStableTarget{sga6.e06.prop.1-7}\textbf{Proposition 1.7.} Let \(S\) be quasi-compact, let \(E\) be locally free of finite type, let \(B=\Sym(E)\), and let \(L\) be a graded \(B\)-module locally free of finite type. Then \(L\) admits a finite decreasing filtration by graded submodules \(L^{(m)}\) such that
\[
\begin{array}{ll}
\text{\rm (i)}& L^{(m)}=L\quad\text{for }m\leq m_0,\\
\text{\rm (ii)}& L^{(m)}=0\quad\text{for }m\geq m_1,\\
\text{\rm (iii)}& \operatorname{gr}^m L
  \simeq N_m\otimes_{\mathcal O_S}B(m),
\end{array}
\]
where \(N_m\) is locally free of finite type on \(S\).

The question is local on \(S\). Since \(S\) is quasi-compact, the degrees of homogeneous bases that occur locally lie between two fixed integers. One argues by induction on the length of this interval. The first step is the canonical map from the lowest non-zero homogeneous component,
\[
  L_{n_0}\otimes_{\mathcal O_S}B(-n_0)\longrightarrow L,
\]
which locally identifies the corresponding direct summand. The quotient has local homogeneous bases in higher degrees, and induction applies.

\sgaStableTarget{sga6.e06.sec.1.8}\subsection*{1.8. Generation by the powers of \(\xi\)}
Let \(F\) be locally free of finite type on \(X\). By \hyperlink{sga6.e06.cor.1-4}{Corollary 1.4} and \hyperlink{sga6.e06.prop.1-6}{Proposition 1.6}, \(F\) comes from a graded module \(M\) admitting a finite graded resolution by \(L_i\)'s. Applying \(\Proj_0\) gives an exact sequence of locally free sheaves on \(X\), and therefore in \(K^\bullet(X)\)
\[
  \cl(F)=\sum_i(-1)^i\cl(\Proj_0(L_i)).
\]
Using the filtration of \hyperlink{sga6.e06.prop.1-7}{Proposition 1.7}, each term is a sum of classes
\[
  \cl\bigl(f^*N_{ij}(j)\bigr)=f^*\cl(N_{ij})\,\xi^j .
\]
Thus \(K^\bullet(X)\) is generated as a \(K^\bullet(S)\)-module by the powers of \(\xi\).

\sgaStableTarget{sga6.e06.sec.1.9}\subsection*{1.9. The relation}
For linear dependence, use the canonical epimorphism \(f^*E\to\mathcal O_X(1)\). Its kernel is locally free of rank \(r\). Hence
\[
  \lambda^{r+1}(f^*E-\xi)=0.
\]
Since \(\lambda_t(-\xi)=1/(1+\xi t)\), one obtains the relation
\sgaStableTarget{sga6.e06.eq.1-9-1}\begin{equation}\tag{1.9.1}
  \sum_{i=0}^{r+1}(-1)^i\,\lambda^i(E)\,\xi^{r+1-i}=0,
\end{equation}
equivalently the characteristic equation of \(\xi\) associated with the locally free module \(E\).

\sgaStableTarget{sga6.e06.sec.1.10}\subsection*{1.10--1.11. Linear independence}
The proof of the linear-independence assertion is due to J.-P. Jouanolou.

\sgaStableTarget{sga6.e06.lem.1-10}\textbf{Lemma 1.10.} With the hypotheses of \hyperlink{sga6.e06.thm.1-1}{Theorem 1.1}, let \(F\) and \(G\) be locally free sheaves of finite type on \(X\), and suppose that they have the same class in \(K^\bullet(X)\). Then there is an integer \(n_0\) such that, for all \(n\geq n_0\),
\[
\begin{array}{ll}
\text{\rm (i)}& f_*(F(n))\text{ and }f_*(G(n))\text{ are locally free of finite type on }S,\\
\text{\rm (ii)}& \cl f_*(F(n))=\cl f_*(G(n))\text{ in }K^\bullet(S).
\end{array}
\]

An equality of classes in \(K^\bullet(X)\) is a finite integral combination of relations coming from short exact sequences of locally free sheaves. Twisting by \(\mathcal O_X(n)\) and applying \(f_*\), for \(n\) large enough all higher direct images vanish and all direct images are locally free of finite type; the same relations then hold in \(K^\bullet(S)\).

Suppose that
\sgaStableTarget{sga6.e06.eq.1-11-1}\begin{equation}\tag{1.11.1}
  f^*(x_0)+f^*(x_1)\xi+\cdots+f^*(x_r)\xi^r=0 .
\end{equation}
Writing each \(x_i\) as a difference of classes of locally free sheaves and adding suitable summands, \hyperlink{sga6.e06.eq.1-11-1}{(1.11.1)} is converted into an equality \(\cl(F)=\cl(G)\) for locally free sheaves of finite type on \(X\). By \hyperlink{sga6.e06.lem.1-10}{Lemma 1.10}, for all sufficiently large \(n\) one has
\sgaStableTarget{sga6.e06.eq.1-11-2}\begin{equation}\tag{1.11.2}
  \sum_{i=0}^{r} x_i\,\sigma_{n+i}(E)=0
  \quad\text{in }K^\bullet(S),
\end{equation}
where \(\sigma_p(E)\) denotes the class of \(\Sym^p(E)\), and \(\sigma_p(E)=0\) for \(p<0\).

The Koszul complex of the symmetric algebra gives, for \(p\geq 1\),
\sgaStableTarget{sga6.e06.eq.1-11-3}\begin{equation}\tag{1.11.3}
  \sum_{i=0}^{r+1}(-1)^i\lambda^i(E)\,\sigma_{p-i}(E)=0 .
\end{equation}
Using this relation, one descends by induction from large \(n\) in \hyperlink{sga6.e06.eq.1-11-2}{(1.11.2)}. It follows that \hyperlink{sga6.e06.eq.1-11-2}{(1.11.2)} holds for all \(n\geq -r\). Taking \(n=-r\) gives \(x_r=0\); repeating the argument gives successively all \(x_i=0\). This proves the linear independence and hence \hyperlink{sga6.e06.thm.1-1}{Theorem 1.1}.

\sgaStableTarget{sga6.e06.rem.1-12}\textbf{Remark 1.12.} Let
\[
  \sigma_t(E)=\sum_{i\geq 0}\sigma_i(E)t^i .
\]
The relation \hyperlink{sga6.e06.eq.1-11-3}{(1.11.3)} is equivalent to
\[
  \lambda_t(E)\,\sigma_{-t}(E)=1.
\]
It follows in particular that \(\sigma_t\) is additive and extends to a homomorphism
\[
  K^\bullet(S)\longrightarrow 1+K^\bullet(S)[[t]]_+ .
\]

\sgaStableTarget{sga6.e06.rem.1-13}\textbf{Remark 1.13.} Any sequence of \(r+1\) consecutive powers of \(\xi\) is a basis of \(K^\bullet(X)\) over \(K^\bullet(S)\). If
\[
  x=1-\xi^{-1}
\]
is the class of a hyperplane section, then \(1,x,\ldots,x^r\) is also a basis. The dependence relation becomes
\sgaStableTarget{sga6.e06.eq.1-13-1}\begin{equation}\tag{1.13.1}
  \sum_{k=0}^{r+1}\lambda^{r+1-k}(\check E-r-1)\,x^k=0.
\end{equation}
For \(E=\mathcal O_S^{r+1}\), so that \(X=\mathbf P^r_S\), this reduces to
\[
  x^{r+1}=0.
\]
Using relative schemes in Hakim's sense, \hyperlink{sga6.e06.thm.1-1}{Theorem 1.1} also makes sense when \(S\) is a locally ringed topos; if the final object of the topos is quasi-compact, the proof transports without essential change.

\sgaStableTarget{sga6.e06.sec.2}\section*{2. Computation of \(K^\bullet\) of a projective bundle: the case of perfect complexes}

The present section will not be used in the rest of the Seminar. In this section the notation \(K^\bullet(X)\) denotes the Grothendieck group of perfect complexes on \(X\) (\hyperlink{sga6.e04.def.2-2}{IV 2.2}). We propose to prove the following theorem, analogous to \hyperlink{sga6.e06.thm.1-1}{1.1}:

\sgaStableTarget{sga6.e06.thm.2-1}\textbf{Theorem 2.1.} Let \(S\) be quasi-compact, let \(E\) be a locally free \(\mathcal O_S\)-module of rank \(r+1\), let \(X=\mathbf P(E)\) be the associated projective bundle; let \(\xi\) be the class in \(K^\bullet(X)\) of the complex equal to \(\mathcal O_X(1)\) in degree zero, and zero in the other degrees. Then \(K^\bullet(X)\), endowed with its natural \(K^\bullet(S)\)-algebra structure, is a free \(K^\bullet(S)\)-module with basis
\[
  1,\xi,\ldots,\xi^r .
\]

\sgaStableTarget{sga6.e06.sec.2.2}\subsection*{2.2. Notation}
Fix the notation: \(B\) will denote the symmetric algebra of \(E\); write \(\mathcal A\) for the category of \(\mathcal O_X\)-modules, \(\mathcal G\) for the category of graded \(B\)-modules with graded morphisms as morphisms, \(\mathcal G_0\) for the full subcategory of \(\mathcal G\) whose objects are the \(B\)-modules locally free of finite type as graded \(B\)-modules, and \(\mathcal G'\) for the category of \(\mathcal O_S\)-modules. For every \(n\), and every complex \(L^\bullet\) of \(D^+(\mathcal A)\), we shall write \(L^\bullet(n)\) for the complex
\[
  L^\bullet(n)=L^\bullet\otimes_{\mathcal O_X}\mathcal O_X(n).
\]

\sgaStableTarget{sga6.e06.sec.2.3}\subsection*{2.3. The functor \(\Proj_0\)}
We shall need to apply the functor \(\Proj_0\) to arbitrary graded \(B\)-modules, and not only to quasi-coherent graded \(B\)-modules as in the construction of EGA II 3.2. We shall therefore indicate briefly how this construction can be made in the general case.

\sgaStableTarget{sga6.e06.para.2-3-1}\textbf{2.3.1.} First let \(S\) be a scheme, let \(A\) be a quasi-coherent \(\mathcal O_S\)-algebra, let \(Y=\Spec(A)\), and let \(f:Y\to S\) be the structural morphism. From the isomorphism
\[
  A\simeq f_*\mathcal O_Y
\]
one obtains a canonical homomorphism
\sgaStableTarget{sga6.e06.eq.2-3-1-1}\begin{equation}\tag{2.3.1.1}
  f^{-1}(A)\longrightarrow \mathcal O_Y,
\end{equation}
where \(f^{-1}\) denotes inverse image in the sense of sheaves of sets; in other words, one has a morphism of ringed spaces
\sgaStableTarget{sga6.e06.eq.2-3-1-2}\begin{equation}\tag{2.3.1.2}
  \bar f:(Y,\mathcal O_Y)\longrightarrow (S,A).
\end{equation}
It gives an inverse-image functor \(\bar f^{\,*}\) from the category of \(A\)-modules to the category of \(\mathcal O_Y\)-modules, left adjoint to the direct-image functor \(\bar f_*\) from the category of \(\mathcal O_Y\)-modules to the category of \(f_*(\mathcal O_Y)\)-modules. Thus
\sgaStableTarget{sga6.e06.eq.2-3-1-3}\begin{equation}\tag{2.3.1.3}
  \bar f^{\,*}(M)=f^{-1}(M)\otimes_{f^{-1}A}\mathcal O_Y
\end{equation}
for every \(A\)-module \(M\).

Since the homomorphism \hyperlink{sga6.e06.eq.2-3-1-1}{(2.3.1.1)} is flat (in each fiber it corresponds to a localization), the functor \(\bar f^{\,*}\) just defined is exact (cf. [4], p. 164). If \(M\) is a quasi-coherent \(A\)-module, this construction gives the module \(\widetilde M\) defined in EGA II 1.4.3. Indeed, the isomorphism
\[
  M\simeq f_*(\widetilde M)=\bar f_*(\widetilde M)
\]
gives, by adjunction, a homomorphism
\[
  \bar f^{\,*}(M)\longrightarrow \widetilde M
\]
and one checks immediately on each fiber that it is an isomorphism.

\sgaStableTarget{sga6.e06.para.2-3-2}\textbf{2.3.2.} Now let \(B\) be a quasi-coherent graded \(\mathcal O_S\)-algebra, let \(M\) be a graded \(B\)-module, and let \(X=\Proj(B)\). For every affine open \(U\) of \(S\), and every section \(s\) of \(B\) over \(U\), homogeneous of degree \(d>0\), write \(M_{(s),U}\) for the homogeneous component of degree \(0\) of the module of fractions \(M_s\); thus
\[
  M_{(s),U}=
  \bigl((M|_U)\otimes_{B|_U}(B|_U)_s\bigr)_0 .
\]
If \(V\) is an affine open of \(U\), then \(M_{(s),V}\) is the restriction to \(V\) of \(M_{(s),U}\).

Consider the base of opens of \(X\) formed by the affine opens
\[
  D_+(U,s)\simeq \Spec\bigl((B|_U)_{(s)}\bigr)
\]
Write \(f\) for the structural morphism \(X\to S\), and \(f_{U,s}\) for the affine morphism
\[
  f_{U,s}:D_+(U,s)\longrightarrow U .
\]
Put
\[
  \Proj_0(M)_{U,s}
  =
  \bar f_{U,s}^{\,*}\bigl(M_{(s),U}\bigr).
\]

With the notation of \hyperlink{sga6.e06.para.2-3-1}{2.3.1}, we shall show that the \(\Proj_0(M)_{U,s}\) glue as \(U\) and \(s\) vary, thereby defining an \(\mathcal O_X\)-module denoted \(\Proj_0(M)\). The gluing for variable \(U\) and fixed \(s\) being evident, one is reduced (cf. EGA II 2.5.2) to showing that, if \(t\) is a section of \(B\) on \(U\), homogeneous of degree \(e>0\), then the restriction of \(\Proj_0(M)_{U,s}\) to \(D_+(U,st)\) is canonically isomorphic to \(\Proj_0(M)_{U,st}\). Consider
\[
\begin{tikzcd}[column sep=large,row sep=large]
D_+(st) \arrow[r,"i"] \arrow[dr,"f_{U,st}"'] &
D_+(s) \arrow[d,"f_{U,s}"]\\
& U
\end{tikzcd}
\]
It is a matter of showing that the modules
\[
  i^*\bar f_{U,s}^{\,*}\bigl(M_{(s),U}\bigr)
  \quad\text{and}\quad
  \bar f_{U,st}^{\,*}\bigl(M_{(st),U}\bigr)
\]
are canonically isomorphic. On the one hand, one has the canonical isomorphisms
\[
\begin{aligned}
i^*\bar f_{U,s}^{\,*}\bigl(M_{(s),U}\bigr)
&\simeq
i^*\Bigl(f_{U,s}^{-1}\bigl(M_{(s),U}\bigr)
 \otimes_{f_{U,s}^{-1}((B|_U)_{(s)})}\mathcal O_{D_+(s)}\Bigr)\\
&\simeq
i^{-1}\Bigl(f_{U,s}^{-1}\bigl(M_{(s),U}\bigr)
 \otimes_{f_{U,s}^{-1}((B|_U)_{(s)})}\mathcal O_{D_+(s)}\Bigr)
 \otimes_{i^{-1}(\mathcal O_{D_+(s)})}\mathcal O_{D_+(st)}\\
&\simeq
f_{U,st}^{-1}\bigl(M_{(s),U}\bigr)
 \otimes_{f_{U,st}^{-1}((B|_U)_{(s)})}\mathcal O_{D_+(st)} .
\end{aligned}
\]
On the other hand,
\[
  \bar f_{U,st}^{\,*}\bigl(M_{(st),U}\bigr)
  \simeq
  f_{U,st}^{-1}\bigl(M_{(st),U}\bigr)
  \otimes_{f_{U,st}^{-1}((B|_U)_{(st)})}\mathcal O_{D_+(st)} .
\]
Finally, it follows from EGA II 2.2.2, translated into sheaf language, that one has the canonical isomorphisms
\[
  M_{(st),U}\simeq
  \bigl(M_{(s),U}\bigr)_{t^d/s^e}
  \simeq
  M_{(s),U}\otimes_{(B|_U)_{(s)}}(B|_U)_{(st)},
\]
and hence the desired isomorphism.

It is clear that, if \(M\) is quasi-coherent, the module \(\Proj_0(M)\) thus constructed is indeed the one defined in EGA II.

\sgaStableTarget{sga6.e06.para.2-3-3}\textbf{2.3.3.} It follows from \hyperlink{sga6.e06.para.2-3-1}{2.3.1} and \hyperlink{sga6.e06.para.2-3-2}{2.3.2} that the functor \(\Proj_0\) is exact. Returning to the notation of \hyperlink{sga6.e06.sec.2.2}{2.2}, by passage to derived categories one therefore obtains a functor, still denoted \(\Proj_0\):
\[
  \Proj_0:D^+(\mathcal G)\longrightarrow D^+(\mathcal A).
\]
It is also clear from the preceding constructions that the functor \(\Proj_0\), on the category of graded \(B\)-modules, is compatible with tensor product; the same remains true of the functor obtained on the derived categories.

\sgaStableTarget{sga6.e06.sec.2.4}\subsection*{2.4. The functor \(F_k\)}
For an \(\mathcal O_X\)-module \(L\) and an integer \(k\), set
\sgaStableTarget{sga6.e06.eq.2-4-1}\begin{equation}\tag{2.4.1}
  F_k(L)=\bigoplus_{n\geq k} f_*(L(n)).
\end{equation}
The functor \(F_k\) is left exact and therefore has a right derived functor
\sgaStableTarget{sga6.e06.eq.2-4-2}\begin{equation}\tag{2.4.2}
  RF_k:D^+(\mathcal A)\longrightarrow D^+(\mathcal G).
\end{equation}
If \(L\) is locally free of finite type, then for \(k\) sufficiently large it is \(F_k\)-acyclic, by the vanishing theorem used above in 1.3.

For every \(k\) there is a functorial morphism
\sgaStableTarget{sga6.e06.eq.2-5-1}\begin{equation}\tag{2.5.1}
  \beta:\Proj_0(F_k(F))\longrightarrow F,
\end{equation}
constructed on each \(D_+(U,s)\) by the usual localization formula: a fraction represented by a homogeneous section of sufficiently high degree is sent to the corresponding local section of \(F\). In the quasi-coherent case this is the canonical morphism of EGA. Passing to derived categories gives
\sgaStableTarget{sga6.e06.eq.2-5-2}\begin{equation}\tag{2.5.2}
  \Proj_0\circ RF_k \longrightarrow \operatorname{id}_{D^+(\mathcal A)}.
\end{equation}

\sgaStableTarget{sga6.e06.prop.2-6}\textbf{Proposition 2.6.} Let \(S\) be quasi-compact, let \(E\) be locally free of finite type, let \(X=\mathbf P(E)\), and let \(L^\bullet\) be a perfect complex on \(X\). Then, for \(k\) sufficiently large,
\[
\begin{array}{ll}
\text{\rm (i)}& RF_k(L^\bullet)\text{ is }\mathcal G_0\text{-perfect},\\[1mm]
\text{\rm (ii)}& \Proj_0 RF_k(L^\bullet)\longrightarrow L^\bullet
\text{ is an isomorphism in }D^+(\mathcal A).
\end{array}
\]

The assertion is local on \(S\). One may therefore assume \(S\) affine. Since \(\mathcal O_X(1)\) is ample, the perfect complex \(L^\bullet\) is locally, and then globally after refining, represented by a bounded complex whose terms are locally free of finite type. For \(k\) sufficiently large all those finitely many terms are \(F_k\)-acyclic, and the graded modules \(F_k(L^i)\) are \(\mathcal G_0\)-perfect. This gives (i), and the morphism \hyperlink{sga6.e06.eq.2-5-2}{(2.5.2)} is then an isomorphism term by term, giving (ii).

\sgaStableTarget{sga6.e06.sec.2.7}\subsection*{2.7. Decomposition of a \(\mathcal G_0\)-perfect complex}
On \(\mathcal G\), the functor which associates to a graded \(B\)-module its \(n\)-th homogeneous component is exact with values in \(\mathcal G'\); it defines a functor
\sgaStableTarget{sga6.e06.eq.2-7-1}\begin{equation}\tag{2.7.1}
  \Comp_n:D^+(\mathcal G)\longrightarrow D^+(\mathcal G').
\end{equation}
The homomorphism \hyperlink{sga6.e06.eq.1-2-1}{(1.2.1)} gives a morphism of functors
\sgaStableTarget{sga6.e06.eq.2-7-2}\begin{equation}\tag{2.7.2}
  \Comp_n\otimes_{\mathcal O_S}B(-n)\longrightarrow
  \operatorname{id}_{D^+(\mathcal G)}.
\end{equation}

Let \(M^\bullet\) be a \(\mathcal G_0\)-perfect complex. We shall show that there exists a finite sequence of distinguished triangles
\[
\begin{array}{ccccccccccc}
& P_1^\bullet &&&& P_2^\bullet &&&& P_k^\bullet &\\
\swarrow && \nwarrow && \swarrow && \nwarrow && \swarrow && \nwarrow\\
N_1^\bullet & \longrightarrow & M^\bullet, & \qquad &
N_2^\bullet & \longrightarrow & P_1^\bullet, & \cdots, &
N_k^\bullet & \longrightarrow & P_{k-1}^\bullet,
\end{array}
\]
where each \(P_k^\bullet\) at the last stage is acyclic, the terms \(N_i^\bullet\) are of the form
\[
  N_i^\bullet=K_i^\bullet\otimes_{\mathcal O_S}B(d_i),
\]
and the \(K_i^\bullet\) are perfect complexes on \(S\) in the ordinary sense.

Indeed, because \(M^\bullet\) is \(\mathcal G_0\)-perfect and \(S\) is quasi-compact, there exists a finite covering of \(S\) and an integer \(d_0\) such that, on every open of the covering, \(M^\bullet\) is represented by a bounded complex of locally free \(B\)-modules of finite type, with no homogeneous generators in degrees \(<d_0\). Taking \(d_0\) minimal, the morphism \hyperlink{sga6.e06.eq.2-7-2}{(2.7.2)} extracts the part of the lowest degree. Its cone has all homogeneous generators in strictly higher degrees. Repeating the process, and using the finiteness of the degree interval from \hyperlink{sga6.e06.prop.1-7}{Proposition 1.7}, one obtains the desired finite sequence of triangles.

\sgaStableTarget{sga6.e06.sec.2.8}\subsection*{2.8. Proof of Theorem 2.1}
Let \(L^\bullet\) be a perfect complex on \(X\). For \(k\) sufficiently large put
\[
  M^\bullet=RF_k(L^\bullet).
\]
By \hyperlink{sga6.e06.prop.2-6}{Proposition 2.6}, \(M^\bullet\) is \(\mathcal G_0\)-perfect and \(\Proj_0(M^\bullet)\simeq L^\bullet\). Applying the decomposition of \hyperlink{sga6.e06.sec.2.7}{2.7} and then \(\Proj_0\), one obtains in \(K^\bullet(X)\)
\[
  \cl(L^\bullet)
  =
  \sum_i \cl\bigl(f^*K_i^\bullet(d_i)\bigr)
  =
  \sum_i f^*\cl(K_i^\bullet)\,\xi^{d_i}.
\]
Thus the powers of \(\xi\) generate \(K^\bullet(X)\) as a \(K^\bullet(S)\)-module.

The dependence relation is the same relation \hyperlink{sga6.e06.eq.1-9-1}{(1.9.1)}, now interpreted in the Grothendieck group of perfect complexes. For linear independence one uses the \(K^\bullet(S)\)-linear homomorphism
\[
  f_*:K^\bullet(X)\longrightarrow K^\bullet(S).
\]
For the structure sheaf and negative twists of the projective bundle one has
\[
  f_*(1)=1,\qquad f_*(\xi^n)=0\quad(-r\leq n<0).
\]
Multiplying a relation among \(1,\xi,\ldots,\xi^r\) by suitable negative powers of \(\xi\) and applying \(f_*\) successively gives all its coefficients equal to zero. This completes the proof of \hyperlink{sga6.e06.thm.2-1}{Theorem 2.1}.



% Continuing macros for Expose VI
\providecommand{\Sym}{\operatorname{Sym}}
\providecommand{\GL}{\operatorname{GL}}
\providecommand{\Grass}{\operatorname{Grass}}
\providecommand{\gr}{\operatorname{gr}}
\providecommand{\ctop}{c_{\mathrm{top}}}
\providecommand{\Rep}{\operatorname{Rep}}
\providecommand{\RZ}{R_{\mathbb Z}}
\providecommand{\Dflag}{D}
\providecommand{\Gmn}{G}


This completes the induction, for if \(0\le s\le r\) and if \(x_j=0\) for \(s<j\le r\), then from a relation
\[
  \sum_{i=0}^{s} f^*(x_i)\xi^i=0
\]
after multiplication by \(\xi^{r-s}\) one gets
\[
  \sum_{i=0}^{s} f^*(x_i)\xi^{i+r-s}=0,
\]
and hence \(x_s=0\) by what precedes.

\sgaStableTarget{sga6.e06.sec.3}\section*{3. Consequences of the structure theorem for the \(K^\bullet\) of a projective bundle}

\sgaStableTarget{sga6.e06.prop.3-1}\textbf{Proposition 3.1.} Let \(S\) be a quasi-compact scheme, let \(E_1,\ldots,E_n\) be \(\mathcal O_S\)-modules locally free of finite type, and let \(D(E_i)\) be the flag scheme of \(E_i\). Put
\[
  X=D(E_1)\times_S\cdots\times_S D(E_n).
\]
Then the canonical homomorphism
\[
  K^\bullet(S)\longrightarrow K^\bullet(X)
\]
is injective, where \(K^\bullet\) is defined either from \(\mathcal O_S\)-modules locally free of finite type or from perfect complexes.

One verifies this immediately by using the universal property of flag schemes. The scheme \(X\) can be constructed as follows: one considers a sequence of \(S\)-schemes \(X_1,\ldots,X_n\), defined by induction by taking \(X_k\) to be the flag scheme of the inverse image of \(E_k\) on \(X_{k-1}\); then \(X=X_n\). Thus, to prove the proposition, it is enough to prove it when \(n=1\), the general case following by induction. We may therefore suppose that \(X\) is the flag scheme of an \(\mathcal O_S\)-module \(E\), locally free of finite type.

% SGA-SRC-001528 — reconstructed by mathematical inference; scan inconclusive
After decomposing \(S\) into finitely many open-and-closed subschemes, we may suppose that \(E\) has constant rank \(r\). The scheme \(X\) may then be constructed as follows. Put \(E=F_0\), \(S=X_0\), and define successively schemes \(X_1,\ldots,X_r\) and, for each \(i\), an \(\mathcal O_{X_i}\)-module locally free of finite type \(F_i\), by
\[
  X_i=\mathbb P(F_{i-1}),\qquad
  F_i=\Ker\bigl(f_i^*(F_{i-1})\longrightarrow \mathcal O_{X_i}(1)\bigr),
\]
where \(f_i:X_i\to X_{i-1}\) is the structural morphism. It follows that it is enough to prove the proposition for a projective bundle. This is precisely the content of \hyperlink{sga6.e06.thm.1-1}{1.1}, respectively of \hyperlink{sga6.e06.thm.2-1}{2.1}, according to which definition of \(K^\bullet\) is being used.

\sgaStableTarget{sga6.e06.thm.3-2}\textbf{Theorem 3.2.} Let \(S\) be a quasi-compact scheme. The pre-\(\lambda\)-ring \(K^\bullet(S)\), defined from \(\mathcal O_S\)-modules locally free of finite type, is a \(\lambda\)-ring.

The homomorphism
\[
  \lambda_t:K^\bullet(S)\longrightarrow 1+K^\bullet(S)[[t]]^+
\]
being additive, it is enough, in order to verify that it is a \(\lambda\)-homomorphism, to check the universal relations on elements of the form \(\cl(E)\), where \(E\) is an \(\mathcal O_S\)-module locally free of finite type. Thus one has to show that for every pair \(E,F\) of such modules, with classes again denoted \(E,F\), the relations \hyperlink{sga6.e05.eq.2-4-2}{V 2.4.2} are satisfied:
\[
  \lambda^i(EF)=P_i\bigl(\lambda^1(E),\ldots,\lambda^i(E);\lambda^1(F),\ldots,\lambda^i(F)\bigr),
\]
\[
  \lambda^j(\lambda^i(E))=Q_{i,j}\bigl(\lambda^1(E),\ldots,\lambda^{ij}(E)\bigr).
\]
Since \(S\) is quasi-compact, it admits a finite covering by disjoint opens such that on each of them \(E\) and \(F\) have constant rank. Since the \(K^\bullet\) of a finite disjoint union is the product of the corresponding \(K^\bullet\)'s, it is enough to check the relations on the members of this covering; hence we may suppose that \(E\) and \(F\) have constant ranks on \(S\).

Put then
\[
  X=D(E)\times_S D(F).
\]
By \hyperlink{sga6.e06.prop.3-1}{Proposition 3.1} the \(\lambda\)-homomorphism \(K^\bullet(S)\to K^\bullet(X)\) is injective. Consequently, in order to prove the relations \hyperlink{sga6.e05.eq.2-4-2}{V 2.4.2}, we may work in \(K^\bullet(X)\). But in \(K^\bullet(X)\) the elements \(E\) and \(F\) split as sums of classes of invertible sheaves. By additivity it is therefore enough to prove the same relations for elements \(\xi,\eta\) which are classes of invertible sheaves. Written in the form \hyperlink{sga6.e05.eq.2-4-1}{V 2.4.1}, the relations are
\[
  1+\xi\eta t=(1+\xi t)\circ(1+\eta t),
  \qquad
  \lambda_u(1+\xi t)=1+\{1+\xi t\}u.
\]
These are just the relations \hyperlink{sga6.e05.eq.2-3-1}{V 2.3.1} and \hyperlink{sga6.e05.eq.2-3-3}{V 2.3.3} defining the \(\lambda\)-ring structure on \(1+K^\bullet(S)[[t]]^+\).

Let us also record the following more general result.

\sgaStableTarget{sga6.e06.thm.3-3}\textbf{Theorem 3.3.} Let \((T,A)\) be a commutatively ringed topos, and let \(K\) be the Grothendieck group of \(A\)-modules locally free of bounded rank on \(T\). Then \(K\), endowed with its canonical pre-\(\lambda\)-structure, is a \(\lambda\)-ring.

To prove that \(K\) is a \(\lambda\)-ring, it is necessary to prove the relations \hyperlink{sga6.e05.eq.2-4-1}{V 2.4.1}; by additivity, it is enough to prove them for classes of \(A\)-modules \(E,F\) locally free of bounded rank. Since the ranks of \(E\) and \(F\) are bounded, one may decompose \(T\) as a finite disjoint sum of opens over which \(E\) and \(F\) have constant rank. Using the compatibility of \(K\) with finite disjoint sums, and writing the relations \hyperlink{sga6.e05.eq.2-4-1}{V 2.4.1} in the form \hyperlink{sga6.e05.eq.2-4-2}{V 2.4.2}, we are reduced to the case of modules of constant rank.

Let \(E\) be an \(A\)-module locally free of finite rank \(n\). There is a torsor \(P\) under \(\GL(n,A)\) such that
\[
  E=P\times_{\GL(n,A)} A^n,
\]
where \(A^n\) is considered as a \(\GL(n,A)\)-module by the standard representation. One can then define a \(\lambda\)-homomorphism from the representation ring \(R_{\mathbb Z}(\GL(n))\) to \(K\) by associating to a representation \(\rho\) of \(\GL(n)\) in a free \(\mathbb Z\)-module of rank \(p\) the class of
\[
  P\times_{\GL(n,A)} A^p.
\]
Under this homomorphism, the class of \(E\) is the image of the standard representation. Thus any polynomial relation among the \(\lambda^i(E)\)'s follows from the corresponding relation for the standard representation in \(R_{\mathbb Z}(\GL(n))\).

Similarly, for two locally free modules \(E,F\) of ranks \(n,p\), one uses a torsor under \(\GL(n,A)\times\GL(p,A)\), with \(\GL(n,A)\) acting through the standard representation and \(\GL(p,A)\) through the trivial one on the first factor, and conversely for \(F\). We are reduced to proving that, for every group scheme \(G\) over \(\Spec\mathbb Z\) which is a finite product of linear group schemes, the pre-\(\lambda\)-ring \(R_{\mathbb Z}(G)\) is a \(\lambda\)-ring. This may be obtained by the explicit computation of \(R_{\mathbb Z}(G)\) in [5]; alternatively, one can reason as in \hyperlink{sga6.e06.thm.3-2}{3.2}, using 1.14.

\sgaStableTarget{sga6.e06.sec.4}\section*{4. Computation of the \(K^\bullet\) of a flag bundle}

\sgaStableTarget{sga6.e06.sec.4.1}\subsection*{4.1}
Let \(S\) be a scheme and let \(E\) be an \(\mathcal O_S\)-module locally free of rank \(n\). Let
\[
  \pi=(p_1,\ldots,p_k),\qquad p_i\geq 1,
  \qquad \sum_{i=1}^k p_i=n.
\]
The flag functor of type \(\pi\) of \(E\) is the functor on the category of \(S\)-schemes, with values in sets, which to an \(S\)-scheme \(f:X\to S\) associates the set of chains of quotients of \(f^*E\)
\[
  (L_{k-1},\ldots,L_1),
\]
where, for each \(i\), \(L_i\) is a quotient of \(L_{i+1}\), locally free of rank \(p_1+\cdots+p_i\). This functor is representable by an \(S\)-scheme \(D_\pi(E)\), endowed with a canonical sequence of quotients of the inverse image of \(E\), locally free of ranks \(p_1,p_1+p_2,\ldots\). If \(\pi=(1,n-1)\), then \(D_\pi(E)=\mathbb P(E)\); if \(\pi=(r,n-r)\), then \(D_\pi(E)=G_{r,n}(E)\), the Grassmannian of type \((r,n-r)\) of \(E\). If all \(p_i\)'s are equal to \(1\), one obtains the usual flag bundle of \(E\), denoted \(D(E)\).

\sgaStableTarget{sga6.e06.lem.4-2}\textbf{Lemma 4.2.} Let \(S\) be a scheme, let \(E\) be an \(\mathcal O_S\)-module locally free of rank \(n\), and let
\[
  \pi=(p_1,\ldots,p_i,\ldots,p_k),\qquad p_1+\cdots+p_k=n.
\]
Let \(p',p''\geq 1\) be such that \(p'+p''=p_i\), and put
\[
  \pi'=(p_1,\ldots,p_{i-1},p',p'',p_{i+1},\ldots,p_k),
  \qquad \pi''=(p',p'').
\]
Let \(L_1,\ldots,L_{k-1}\) be the canonical quotients of the inverse image of \(E\) on \(D_\pi(E)\), and let \(F\) be the kernel of \(L_i\to L_{i-1}\). Then there is a canonical \(S\)-isomorphism
\[
  D_{\pi'}(E)\simeq D_{\pi''}(F).
\]
The proof, immediate from the functorial definition of flag schemes, is left to the reader.

\sgaStableTarget{sga6.e06.lem.4-3}\textbf{Lemma 4.3.} Let \(A\) be a commutative ring with unit, let \(n>0\), and let \((c_i)_{1\leq i\leq n}\) be elements of \(A\). Let \(C\) be the \(A\)-algebra generated by \(\xi_1,\ldots,\xi_n\), subject to the relations
\[
  \sigma_i(\xi_1,\ldots,\xi_n)=c_i,
  \qquad i=1,\ldots,n,
\]
where \(\sigma_i\) denotes the elementary symmetric functions. Let \(p_1,\ldots,p_k\) be integers \(\geq 1\) whose sum is \(n\). For \(1\leq j\leq k\) put
\[
  \xi^{(j)}=(\xi^{(j)}_1,\ldots,\xi^{(j)}_{p_j}),
  \qquad
  \xi_i^{(j)}=\xi_{\sum_{\alpha<j}p_\alpha+i}.
\]
Put
\[
  c_i^{(j)}=\sigma_i(\xi^{(j)}_1,\ldots,\xi^{(j)}_{p_j}),
  \qquad 1\leq j\leq k,
  \quad 1\leq i\leq p_j,
\]
and let \(B\) be the subalgebra of \(C\) generated by the \(c_i^{(j)}\)'s. Then:

\begin{enumerate}[label=\textup{\roman*)}]
\item The elements
\[
  (\xi^{(1)})^{\alpha(1)}\cdots(\xi^{(k)})^{\alpha(k)},
  \qquad 0\leq \alpha_i^{(j)}\leq p_j-i,
\]
form a basis of \(C\), considered as a \(B\)-module.
\item The generators \(c_i^{(j)}\) of the \(A\)-algebra \(B\) satisfy the \(n\) isobaric relations deduced from
\[
% SGA-SRC-001520 — reconstructed by mathematical inference; scan inconclusive
  1+\sum_{i=1}^n c_i
  =
  \prod_{j=1}^{k}\left(1+\sum_{i=1}^{p_j} c_i^{(j)}\right),
\]
where \(c_i\) and \(c_i^{(j)}\) are assigned weight \(i\); these \(n\) relations generate the ideal of relations among the \(c_i^{(j)}\)'s.
\end{enumerate}

\sgaStableTarget{sga6.e06.cor.4-4}\textbf{Corollary 4.4.} The elements of \(C\) of the form
\[
  \xi_1^{\alpha_1}\cdots \xi_{n-1}^{\alpha_{n-1}},
  \qquad 0\leq \alpha_i\leq n-i,
\]
form a basis of \(C\) as an \(A\)-module.

For the proof of the lemma and corollary we refer to [2], Exposé 4, pages 4.19, from which they are borrowed.

\sgaStableTarget{sga6.e06.sec.4.5}\subsection*{4.5}
Let \(S\) be quasi-compact, let \(E\) be an \(\mathcal O_S\)-module locally free of rank \(n\), let \(\pi=(p_1,\ldots,p_k)\) be a sequence of integers \(\geq 1\) with sum \(n\), and put \(X=D_\pi(E)\). The scheme \(X\) is equipped with canonical quotients \(L_i\). Put \(L_0=0\), \(L_k=f^*E\), where \(f:X\to S\) is the structural morphism. For \(1\leq i\leq k\), let \(F_i\) be the kernel of \(L_i\to L_{i-1}\). Let \(K^\bullet(X)\) be the Grothendieck group of locally free finite-type sheaves on \(X\), respectively of perfect complexes. Denote by \(F_j\) the class of \(F_j\), and by \(\lambda_i^{(j)}\) the class of \(\lambda^i(F_j)\). In the case where \(K^\bullet(X)\) is defined by perfect complexes, one gives the same evident meaning to the series \(\lambda_t(E)\) and \(\lambda_t(F_j)\), where \(E\) is the class of \(E\). Finally,
\[
  \sum_{j=1}^{k}F_j=E.
\]

\sgaStableTarget{sga6.e06.thm.4-6}\textbf{Theorem 4.6.} With the hypotheses and notations of \hyperlink{sga6.e06.sec.4.5}{4.5}, the \(K^\bullet(S)\)-algebra \(K^\bullet(X)\) is generated by the \(\lambda_i^{(j)}\), and the ideal of relations among the \(\lambda_i^{(j)}\)'s is generated by the relations deduced from
\[
  \lambda_t(E)=\prod_{j=1}^{k}\lambda_t(F_j).
\]
Equivalently, \(K^\bullet(X)\) is the \(K^\bullet(S)\)-\(\lambda\)-algebra generated by \(F_1,\ldots,F_k\), subject to the relations
\[
  \text{(i)}\quad \lambda^p(F_i)=0\quad(p>p_i),
  \qquad
  \text{(ii)}\quad \sum_i F_i=E.
\]

We first treat the case where all \(p_i\)'s are equal to \(1\). Then the theorem reads as follows.

\sgaStableTarget{sga6.e06.cor.4-7}\textbf{Corollary 4.7.} Let \(S\) be quasi-compact, let \(E\) be an \(\mathcal O_S\)-module locally free of rank \(n\), let \(X=D(E)\) be the flag bundle of \(E\), and let \(\xi_1,\ldots,\xi_n\) be the classes of the invertible sheaves \(F_1,\ldots,F_n\). Then \(K^\bullet(X)\) is the \(K^\bullet(S)\)-algebra generated by the \(\xi_i\)'s, subject to the relations
\[
  \sigma_i(\xi_1,\ldots,\xi_n)=\lambda^i(E).
\]

We argue by induction on \(n\). For \(n=1\) the assertion is trivial. Suppose it known for \(n-1\). Consider \(\mathbb P(E)\), and on \(\mathbb P(E)\) the exact sequence
\[
  0\longrightarrow F\longrightarrow g^*E\longrightarrow \mathcal O_{\mathbb P(E)}(1)\longrightarrow 0,
\]
where \(g\) is the structural morphism of \(\mathbb P(E)\). By \hyperlink{sga6.e06.lem.4-2}{Lemma 4.2}, \(X\) identifies with the usual flag bundle \(D(F)\) over \(\mathbb P(E)\), and the canonical sheaves correspond under this identification. The diagram
\[
\begin{tikzcd}[column sep=large,row sep=large]
D(E)=D(F) \arrow[r,"h"] \arrow[dr,"f"'] & \mathbb P(E) \arrow[d,"g"]\\
& S
\end{tikzcd}
\]
shows that \(K^\bullet(S)\to K^\bullet(X)\) factors through \(K^\bullet(\mathbb P(E))\). By the induction hypothesis applied to \(D(F)\), and by \hyperlink{sga6.e06.cor.4-4}{Corollary 4.4}, \(K^\bullet(X)\) has over \(K^\bullet(\mathbb P(E))\) a basis formed by
\[
  \xi_2^{\alpha_2}\cdots \xi_n^{\alpha_{n-1}},
  \qquad 0\leq \alpha_i\leq n-i.
\]
Moreover \(K^\bullet(\mathbb P(E))\) has over \(K^\bullet(S)\) the basis \(1,\xi_1,\ldots,\xi_1^{n-1}\). Hence \(K^\bullet(X)\) has over \(K^\bullet(S)\) the basis
\[
  \xi_1^{\alpha_1}\cdots \xi_{n-1}^{\alpha_{n-1}},
  \qquad 0\leq \alpha_i\leq n-i.
\]
In addition, the additivity of \(\lambda_t\) gives \(\prod_i(1+\xi_i t)=\lambda_t(E)\), so the displayed relations are satisfied. If \(C\) is the \(K^\bullet(S)\)-algebra generated by \(\xi_i\)'s and subjected to these relations, then there is a surjective homomorphism \(C\to K^\bullet(X)\). By \hyperlink{sga6.e06.cor.4-4}{Corollary 4.4} this homomorphism sends a basis of \(C\) to a basis of \(K^\bullet(X)\), hence it is an isomorphism.

For a general flag bundle of type \(\pi\), the usual flag bundle \(D(E)\) is obtained from \(D_\pi(E)\) by constructing successively the flag bundles of the \(F_j\)'s. It follows that \(K^\bullet(D(E))\), considered as a module over \(K^\bullet(X)\), has a basis formed by the elements appearing in 4.3.i), the \(\xi_i^{(j)}\)'s being the classes of the sheaves analogous to the \(F_j\)'s in the splitting of the inverse images of the \(F_j\)'s on \(D(E)\). Thus \(K^\bullet(X)\) is a subring of \(K^\bullet(D(E))\). Moreover, additivity gives
\[
  \lambda_i^{(j)}=\sigma_i(\xi_1^{(j)},\ldots,\xi_{p_j}^{(j)}),
\]
and therefore \(K^\bullet(X)\) contains the sub-\(K^\bullet(S)\)-algebra of \(K^\bullet(D(E))\) generated by the \(\lambda_i^{(j)}\)'s. If \(B\) is this subalgebra, then by 4.3.i) the same family of elements is a basis both for the \(B\)-module structure and for the \(K^\bullet(X)\)-module structure. Hence \(K^\bullet(X)=B\), proving generation. The ideal of relations among the \(\lambda_i^{(j)}\)'s is then given by 4.3.ii), taking into account the structure of \(K^\bullet(D(E))\) over \(K^\bullet(S)\) determined by \hyperlink{sga6.e06.cor.4-7}{4.7}, and the fact that the \(\xi_i^{(j)}\)'s are precisely the classes involved in the splitting of the inverse image of \(E\) on \(D(E)\). This proves the first form of 4.6.

For the second form, one observes that the elements \(F_i\) of the \(K^\bullet(S)\)-\(\lambda\)-algebra \(K^\bullet(X)\) satisfy (i) and (ii). If \(A\) is the \(K^\bullet(S)\)-\(\lambda\)-algebra generated by elements \(F_i'\) satisfying (i) and (ii), there is therefore a \(K^\bullet(S)\)-\(\lambda\)-homomorphism \(A\to K^\bullet(X)\) sending \(F_i'\) to \(F_i\). By \hyperlink{sga6.e06.thm.4-6}{4.6} this homomorphism is surjective. Since the \(\lambda\)-ideal generated by (i) and (ii) in \(K^\bullet(S)[F_1,\ldots,F_k]\) contains the usual ideal generated by the first form of the relations, the homomorphism is also injective. Thus it is an isomorphism.

\sgaStableTarget{sga6.e06.cor.4-8}\textbf{Corollary 4.8.} With the hypotheses and notations of \hyperlink{sga6.e06.sec.4.5}{4.5}, put
\[
  c_i^{(j)}=\gamma^i(F_j-p_j),
  \qquad 1\leq i\leq p_j.
\]
Then the \(K^\bullet(S)\)-algebra \(K^\bullet(X)\) is generated by the \(c_i^{(j)}\)'s, and the ideal of relations among them is generated by the relations deduced from
\[
  \gamma_t(E-n)=\prod_{j=1}^{k}\gamma_t(F_j-p_j).
\]
This is merely another formulation of 4.6. Indeed, the \(c_i^{(j)}\)'s may be expressed as functions of the \(\lambda_i^{(j)}\)'s and conversely. Moreover the relation \(\lambda_t(E)=\prod\lambda_t(F_j)\) is equivalent to \(\gamma_t(E-n)=\prod\gamma_t(F_j-p_j)\).

\sgaStableTarget{sga6.e06.sec.4.9}\subsection*{4.9}
Let \(S\) be a scheme. Denote by \(G_S^{n,r}\) the Grassmannian of type \((n,r-n)\) of \(\mathcal O_S^r\), with \(r\ge n\). For variable \(r\), the \(G_S^{n,r}\) form an inductive system, the morphism
\[
  G_S^{n,r}\longrightarrow G_S^{n,r+1}
\]
being defined by considering \(\mathcal O_{G_S^{n,r}}^r\) as the quotient of \(\mathcal O_{G_S^{n,r}}^{r+1}\) by the \((r+1)\)-st factor. In this system of morphisms the canonical quotients therefore correspond by inverse image. The inductive system so defined will be denoted \(G_S^{n,\infty}\).

We define its \(K^\bullet\) by the formula
\sgaStableTarget{sga6.e06.eq.4-9-1}\begin{equation}\tag{4.9.1}
  K^\bullet(G_S^{n,\infty})
  = \varprojlim_r K^\bullet(G_S^{n,r}) .
\end{equation}
Let \(F_1,F_2\) be the canonical sheaves on \(G_S^{n,r}\) \hyperlink{sga6.e06.sec.4.5}{(4.5)}, and \(F_1,F_2\) their classes in \(K^\bullet(G_S^{n,r})\). Put
\[
  c_i=\gamma^i(F_1-n).
\]
When \(r\) varies, the \(c_i\) correspond under the transition morphisms and consequently define elements, still denoted \(c_i\), of \(K^\bullet(G_S^{n,\infty})\).

\sgaStableTarget{sga6.e06.prop.4-10}\textbf{Proposition 4.10.} Let \(S\) be a quasi-compact scheme. Then, with the notation of \hyperlink{sga6.e06.sec.4.9}{4.9}, one has an isomorphism
\[
  K^\bullet(G_S^{n,\infty})\simeq K^\bullet(S)[[c_1,\ldots,c_n]] .
\]

By \hyperlink{sga6.e06.cor.4-8}{4.8}, \(K^\bullet(G_S^{n,r})\) is the \(K^\bullet(S)\)-algebra generated by
\[
  c_i=\gamma^i(F_1-n),\quad 1\le i\le n,
\]
and by
\[
  \gamma^j(F_2-r+n),\quad 1\le j\le r-n,
\]
subject to the relations deduced from
\[
  \gamma_t(F_1-n)\cdot\gamma_t(F_2-r+n)
  = \gamma_t(r-r)=1,
\]
or equivalently
\[
  \gamma_t(F_2-r+n)=1/\gamma_t(F_1-n).
\]
This relation makes it possible to express the \(\gamma^j(F_2-r+n)\) in terms of the \(c_i\), and gives a series of isobaric relations among the \(c_i\) (where \(c_i\) is assumed to have weight \(i\)) by writing that the coefficients of \(1/\gamma_t(F_1-n)\) vanish in degree \(>r-n\); moreover, these relations generate the ideal of relations among the \(c_i\). We may therefore write
\[
  K^\bullet(G_S^{n,r})=K^\bullet(S)[c_1,\ldots,c_n]/I_r,
\]
where \(I_r\) is a homogeneous ideal of \(K^\bullet(S)[c_1,\ldots,c_n]\), of degree \(>r-n\). It then follows from the \hyperlink{sga6.e06.eq.4-9-1}{definition (4.9.1)} that \(K^\bullet(G_S^{n,\infty})\) is the separated completion of \(K^\bullet(S)[c_1,\ldots,c_n]\) for the topology defined by the \(I_r\). This topology is finer than the topology defined by weight; we shall show that it is equivalent to it, which will imply the proposition.

Put \(A=K^\bullet(S)[c_1,\ldots,c_n]/I_r\), and let \(C\) be the \(A\)-algebra generated by generators \(x_1,\ldots,x_n\) subject to the relations
\[
  \sigma_j(x_1,\ldots,x_n)=c_j,
\]
where the \(\sigma_j\) denote the elementary symmetric functions. By \hyperlink{sga6.e06.cor.4-4}{4.4}, the canonical homomorphism \(A\to C\) is injective. We obtain
\[
  \gamma_t(F_1-n)
  =1+c_1t+\cdots+c_nt^n
  =\prod_{i=1}^n(1+x_it).
\]
Since \(1/\gamma_t(F_1-n)\) is a polynomial, so is \(1/(1+x_it)\), for every \(i\). Hence the \(x_i\) are nilpotent, and there exists an integer \(N\) such that every polynomial of degree \(>N\) in the \(x_i\), with coefficients in \(K^\bullet(S)\), is zero. Consequently every polynomial of \(A\) of weight \(>N\) is zero, i.e. every polynomial of \(K^\bullet(S)[c_1,\ldots,c_n]\) of weight \(>N\) belongs to \(I_r\). It follows that the two topologies under consideration are equivalent, completing the proof.

\sgaStableTarget{sga6.e06.sec.5}\section*{5. Applications to projective bundles; study of \(f_*\)}

\sgaStableTarget{sga6.e06.sec.5.1}\subsection*{5.1}
Let \(S\) be quasi-compact, let \(E\) be locally free of rank \(r+1\), and put \(X=\mathbb P(E)\), with structural morphism \(f:X\to S\). On \(X\) one has the canonical exact sequence
\[
  0\longrightarrow F\longrightarrow f^*E\longrightarrow \mathcal O_X(1)\longrightarrow 0.
\]
Let \(E,F,\xi\) denote the corresponding classes in \(K^\bullet\), and put \(x=\xi-1\).

\sgaStableTarget{sga6.e06.sec.5.2}\subsection*{5.2}
For a locally free sheaf \(G\) on \(X\), the higher direct images \(R^if_*G\) need not be locally free of finite type on \(S\). Thus the direct image \(f_*:K^\bullet(X)\to K^\bullet(S)\) is not defined here by alternating sums of higher direct images. \hyperlink{sga6.e06.thm.1-1}{Theorem 1.1} gives a substitute: every \(y\in K^\bullet(X)\) is uniquely expressible in the form
\[
  y=a_0+a_1\xi+\cdots+a_r\xi^r,
  \qquad a_i\in K^\bullet(S),
\]
and we put
\[
  f_*(y)=a_0.
\]
This additive homomorphism satisfies the projection formula. Moreover,
\[
  f_*(1)=1,
  \qquad
  f_*(\xi^n)=0\quad(-r\leq n<0),
  \qquad
  f_*(\xi^n)=s_n\quad(n\geq 0),
\]
where \(s_n\) is the \(n\)-th complete symmetric function in the roots of \(E\). The last relations follow by induction from the dependence relation and the Koszul relation.

\sgaStableTarget{sga6.e06.prop.5-3}\textbf{Proposition 5.3.} For every \(k\geq 0\),
\[
  \Fil^k K^\bullet(X)
  =
  \Fil^k K^\bullet(S)
  +\Fil^{k-1}K^\bullet(S)x+\cdots+\Fil^{k-r}K^\bullet(S)x^r,
\]
as a direct sum of \(K^\bullet(S)\)-modules, where the filtration is the canonical filtration of augmented \(\lambda\)-rings.

Here \(x\) has augmentation zero and
\[
  \gamma_t(x)=\lambda_{t/(1-t)}(1-\xi^{-1})=\frac{1}{1-xt},
\]
so that \(\gamma^i(x)=x^i\). The result follows from the computation of the filtration for free \(\lambda\)-rings in \hyperlink{sga6.e05.prop.4-12}{V, 4.12} and \hyperlink{sga6.e05.rem.4-13}{V, 4.13}, together with the relation \hyperlink{sga6.e06.eq.1-9-1}{(1.9.1)}. Directness follows from the independence of \(1,x,\ldots,x^r\).

\sgaStableTarget{sga6.e06.cor.5-4}\textbf{Corollary 5.4.} The filtration induced on \(K^\bullet(S)\) by \(K^\bullet(X)\) is the original filtration of \(K^\bullet(S)\).

\sgaStableTarget{sga6.e06.cor.5-5}\textbf{Corollary 5.5.} If \(X\) is a product over \(S\) of flag bundles of locally free sheaves, then
\[
  \Gr K^\bullet(S)\longrightarrow \Gr K^\bullet(X)
\]
is injective.

\sgaStableTarget{sga6.e06.cor.5-6}\textbf{Corollary 5.6.} With the notation of the projective bundle above, \(\Gr K^\bullet(X)\) is a free \(\Gr K^\bullet(S)\)-module with basis \(1,x,\ldots,x^r\), and
\[
  \Gr^k K^\bullet(X)
  =\Gr^kK^\bullet(S)\oplus\Gr^{k-1}K^\bullet(S)x\oplus\cdots\oplus\Gr^{k-r}K^\bullet(S)x^r.
\]
The elements \(1,x,\ldots,x^{r+1}\) satisfy
\[
  \sum_{i=0}^{r+1}(-1)^i c_i(E)x^{r+1-i}=0,
\]
where \(c_i(E)\) is the \(i\)-th Chern class of \(E\).

\sgaStableTarget{sga6.e06.cor.5-7}\textbf{Corollary 5.7.} For a flag bundle as in \hyperlink{sga6.e06.sec.4.5}{4.5}, the \(\Gr K^\bullet(S)\)-algebra \(\Gr K^\bullet(X)\) is generated by the classes \(\xi_i^{(\alpha)}\) of \(\xi_i^{(\alpha)}-1\), with relations generated by
\[
  c_t(E)=\prod_{i,\alpha}c_t(\xi_i^{(\alpha)}),
\]
where \(c_t\) denotes the total Chern class.

\sgaStableTarget{sga6.e06.cor.5-8}\textbf{Corollary 5.8.} For every \(k\),
\[
  f_*(\Fil^k K^\bullet(X))\subset \Fil^k K^\bullet(S).
\]
This follows immediately from \hyperlink{sga6.e06.prop.5-3}{Proposition 5.3} and the projection formula.

\sgaStableTarget{sga6.e06.prop.5-9}\textbf{Proposition 5.9.} With the notation of \hyperlink{sga6.e06.sec.5.1}{5.1}, for every \(y\in K^\bullet(X)\), one has
\[
  f_*(\lambda_{-1}(F)y)=0,
  \qquad
  f_*(\gamma^*(F)y)=0.
\]
In particular,
\[
  f_*(\lambda_{-1}(F))=1.
\]
If
\[
  \xi_t:1+K^\bullet(X)[[t]]^+\longrightarrow 1+K^\bullet(S)[[t]]^+
\]
is defined coefficientwise by
\[
  \xi_t\left(\sum a_i t^i\right)=\sum f_*(a_i)t^i,
\]
then, although it is not in general a homomorphism of groups, it satisfies
\[
  \xi_t(mf^*y)=\xi_t(m)y.
\]
Using \(E=F+\xi\) and \(\lambda_t(E)=\lambda_t(F)(1+\xi t)\), one obtains the asserted identities.

\sgaStableTarget{sga6.e06.prop.5-10}\textbf{Proposition 5.10.} Let \(z\in K^\bullet(X)\). The condition
\[
  z(1-\xi)=0
\]
is equivalent to the existence of an element \(y\in K^\bullet(S)\) such that
\[
  \text{(i)}\quad z=y\lambda_{-1}(F),
  \qquad
  \text{(ii)}\quad y\lambda_{-1}(E)=0.
\]
Sufficiency follows from
\[
  \lambda_{-1}(E)=\lambda_{-1}(F)\lambda_{-1}(\xi)=\lambda_{-1}(F)(1-\xi).
\]
Conversely, write \(z=\sum_{i=0}^r a_i\xi^i\). Combining the dependence relation among the powers of \(\xi\) with \(z(1-\xi)=0\), and using the independence of \(1,\ldots,\xi^r\), expresses \(z\) as a multiple of \(\lambda_{-1}(F)\). The top relation gives \(y\lambda_{-1}(E)=0\).

\sgaStableTarget{sga6.e06.prop.5-11}\textbf{Proposition 5.11.} Let \(z\in \Gr K^\bullet(X)\). The condition \(zx=0\) is equivalent to the existence of \(y\in \Gr K^\bullet(S)\) such that
\[
  \text{(i)}\quad z=y\,c_{\mathrm{top}}(F),
  \qquad
  \text{(ii)}\quad y\,c_{\mathrm{top}}(E)=0.
\]
The proof is the same as that of \hyperlink{sga6.e06.prop.5-10}{5.10}, after passing to associated graded rings and using 5.6.

\sgaStableTarget{sga6.e06.sec.6}\section*{6. Study of the filtration of \(K^\bullet(X)\) when \(X\) has an ample sheaf}

Let \(X\) be a scheme possessing an ample sheaf. Then the two definitions of \(K^\bullet(X)\), by locally free finite-type sheaves and by perfect complexes, coincide. Let
\[
  \varepsilon:K^\bullet(X)\longrightarrow H^0(X,\mathbb Z)
\]
be the augmentation.

\sgaStableTarget{sga6.e06.prop.6-1}\textbf{Proposition 6.1.} If \(X\) is quasi-compact and has an ample sheaf, then \(\Fil^1 K^\bullet(X)\) is a nilideal.

It is enough to show that \(E-\varepsilon(E)\) is nilpotent, where \(E\) is the class of a locally free sheaf. By decomposing \(X\) into finitely many disjoint opens and then passing to the flag bundle of \(E\), one reduces to the case where \(E\) is the class of an invertible sheaf. Let \(L\) be the class of an ample invertible sheaf. For \(n\) sufficiently large, \(E^\vee(n)\) is generated by global sections, giving a surjection \((E(-n))^m\to \mathcal O_X\) with locally free kernel. Thus in \(K^\bullet(X)\)
\[
  mEL^{-n}=1+F.
\]
% SGA-SRC-001543
The corresponding \(\gamma\)-relation shows that \(1-EL^{-n}\), and therefore \(L^n-E\), is nilpotent. Since \(E-1=(L^n-1)-(L^n-E)\), and since \(L^n-1\) is also nilpotent for \(n\) large, the assertion follows.

\sgaStableTarget{sga6.e06.prop.6-2}\textbf{Proposition 6.2.} Let \(R\) be an augmented sub-\(\lambda\)-algebra of \(K^\bullet(X)\), generated by the classes of finitely many sheaves. Then the \(\lambda\)-ring filtration of \(R\) is discrete and \(\Fil^1R\) is nilpotent. Moreover, \(R\) is a finite module over its image in \(H^0(X,\mathbb Z)\).

After reducing to constant ranks, \(R\) is generated by finitely many elements \(x_i=E_i-\varepsilon(E_i)\) and their \(\gamma\)-operations. \hyperlink{sga6.e06.prop.6-1}{Proposition 6.1} shows that all these elements are nilpotent. Hence the augmentation ideal of \(R\) is nilpotent; only finitely many monomials in the \(x_i\)'s and the \(\gamma^j(x_i)\)'s remain. This gives both discreteness and finiteness.

\sgaStableTarget{sga6.e06.cor.6-3}\textbf{Corollary 6.3.} Let \(A\) be a commutative ring. An element \(y\) of \(K^\bullet(X)\otimes_{\mathbb Z}A\) is invertible if and only if its augmentation is invertible. In particular, the class of every nowhere-zero locally free sheaf of finite type becomes invertible in \(K^\bullet(X)\otimes_{\mathbb Z}A\).

Indeed, the augmentation ideal is nil.

\sgaStableTarget{sga6.e06.cor.6-4}\textbf{Corollary 6.4.} Let \(X\) be quasi-compact with an ample sheaf, and let \(f:X'\to X\) be proper and perfect. Assume either that \(f\) is projective or that \(X\) is noetherian, so that the finiteness theorem applies. If there is an element \(x'\in K^\bullet(X')\otimes A\) such that the augmentation of \(f_*(x')\) is invertible, then
\[
  f^*:K^\bullet(X)\otimes A\longrightarrow K^\bullet(X')\otimes A
\]
is injective.
Indeed, by the projection formula,
\[
  f_*(f^*y\cdot x')=y\cdot f_*(x'),
\]
and \(f_*(x')\) is invertible by 6.3.

\sgaStableTarget{sga6.e06.sec.6.5}\subsection*{6.5. The codimension filtration}
Assume now that \(X\) is noetherian of dimension \(d\), and has an ample sheaf. Define a filtration \(K^\bullet(X)^i\) as follows: an element \(x\in K^\bullet(X)\) lies in \(K^\bullet(X)^i\) if, for every closed subscheme \(Y\subset X\), it can be represented by a perfect complex \(M^\bullet(Y)\) such that
\begin{enumerate}[label=\textup{\alph*)}]
\item \(\cl(M^\bullet)=x\);
\item \(\codim(\Supp H^j(M^\bullet)\cap Y,Y)\geq i+j\), and \(H^j(M^\bullet)=0\) for \(j\gg 0\).
\end{enumerate}
These \(K^\bullet(X)^i\)'s are additive subgroups.

\sgaStableTarget{sga6.e06.prop.6-6}\textbf{Proposition 6.6.} The filtration of \hyperlink{sga6.e06.sec.6.5}{6.5} is a ring filtration, and \(K^\bullet(X)^{d+1}=0\).

If \(x\) and \(x'\) are represented by perfect complexes \(M^\bullet\) and \(M'{}^\bullet\) satisfying the above codimension estimates, then \(xx'\) is represented by \(M^\bullet\otimes M'{}^\bullet\). The support of \(H^n(M^\bullet\otimes M'{}^\bullet)\) is contained in the union of the intersections of supports with total degree \(n\), and the codimension estimate follows from EGA \(0_{\mathrm{IV}}\), 14.2.2.3. Finally \(K^\bullet(X)^{d+1}=0\), because a support in codimension greater than \(d\) is empty.

\sgaStableTarget{sga6.e06.lem.6-7}\textbf{Lemma 6.7.} Let \(X=\Spec A\) be affine noetherian, let \(x_1,\ldots,x_n\) be finitely many points of \(X\), and let \(E\) be invertible. There exists a section \(s\in \Gamma(X,E)\) which is a basis of \(E_{x_i}\) for every \(i\).

The proof uses closed specializations \(x_i'\) with distinct maximal ideals \(\mathfrak m_i\), the isomorphisms \(M/\mathfrak m_iM\simeq A/\mathfrak m_i\), for \(M=\Gamma(X,E)\), and Nakayama's lemma.

% SGA-SRC-001548 — reconstructed by mathematical inference; scan inconclusive
\sgaStableTarget{sga6.e06.lem.6-8}\textbf{Lemma 6.8.} Let \(X\) be noetherian with an ample sheaf \(L\), let \(E\) be invertible, and let \(x_1,\ldots,x_n\) be points of \(X\). There exists an integer \(N\) such that, for every \(k\geq1\), \(E(kN)\) has a global section which is a basis at every \(x_i\).

Choose a power of the ample sheaf and a section whose non-vanishing locus is affine and contains all the \(x_i\)'s; apply \hyperlink{sga6.e06.lem.6-7}{Lemma 6.7} there, and extend after multiplying by a further power.

\sgaStableTarget{sga6.e06.thm.6-9}\textbf{Theorem 6.9.} Let \(X\) be a noetherian scheme of dimension \(d\) with an ample sheaf. Then
\[
  \Fil^{d+1}K^\bullet(X)=0,
\]
and therefore \((\Fil^1K^\bullet(X))^{d+1}=0\).

First, if \(E\) and \(E'\) are invertible, with classes still denoted \(E,E'\), then \(E-E'\in K^\bullet(X)^1\). For a closed \(Y\subset X\), choose one point in each irreducible component. \hyperlink{sga6.e06.lem.6-8}{Lemma 6.8} gives morphisms from suitable negative powers of the ample sheaf to \(E\) and \(E'\) which are isomorphisms at those points. The two-term complex built from these morphisms has class \(E-E'\), and is acyclic at the chosen generic points of \(Y\); hence its cohomology supports have codimension at least \(1\) in \(Y\).

Now take \(x\in \Fil^{d+1}K^\bullet(X)\). It is enough to treat products
\[
  \gamma^{i_1}(E_1-F_1)\cdots\gamma^{i_k}(E_k-F_k),
  \qquad \sum i_j\geq d+1,
\]
where the \(E_i,F_i\) are represented by locally free sheaves of the same constant rank. Pull back to the product of the flag bundles of these sheaves. The pullback is injective on \(K^\bullet\), and each class splits into sums of invertible classes. For invertible classes \(\xi,\eta\),
\[
  \gamma_t(\xi-\eta)
  =\frac{\gamma_t(\xi-1)}{\gamma_t(\eta-1)}
  =\frac{1+(\xi-1)t}{1+(\eta-1)t},
\]
so
\[
  \gamma^i(\xi-\eta)=(-1)^{i-1}(\xi-\eta)(\eta-1)^{i-1}\in K^\bullet(X')^i.
\]
Thus the pullback of \(x\) belongs to \(K^\bullet(X')^{d+1}\). Multiplication by a basis element of \(K^\bullet(X')\) over \(K^\bullet(X)\) of relative codimension \(r\) gives an element of \(K^\bullet(X')^{d+1+r}=0\), since \(\dim X'=d+r\). The basis property and injectivity then give \(x=0\).

\sgaStableTarget{sga6.e06.problem.6-10}\textbf{Problem 6.10.} Under the hypotheses of \hyperlink{sga6.e06.thm.6-9}{Theorem 6.9}, is it true that for every \(k\) one has
\[
  K^\bullet(X)^k=\Fil^kK^\bullet(X)?
\]

\sgaStableTarget{sga6.e06.bibliography-for-expose-vi}\subsection*{Bibliography for Exposé VI}
\begin{enumerate}[label={\textup{[\arabic*]}}]
\item H. Bass, \emph{K-Theory and stable algebra}, Publications Mathématiques de l'IHÉS, no. 22.
\item C. Chevalley, Séminaire 1958: \emph{Anneaux de Chow et applications}.
\item M. Hakim, \emph{Topos localement annelés et schémas relatifs}, thesis in preparation.
\item R. Hartshorne, \emph{Residues and duality}, Lecture Notes in Mathematics 20, Springer-Verlag.
\item J.-P. Serre, \emph{Le groupe de Grothendieck des schémas en groupes réductifs}, Publications Mathématiques de l'IHÉS, no. 34.
\end{enumerate}



\providecommand{\Kosz}{\mathrm K}
\providecommand{\Proj}{\operatorname{Proj}}
\providecommand{\Sym}{\operatorname{Sym}}
\providecommand{\gr}{\operatorname{gr}}
\providecommand{\Tor}{\operatorname{Tor}}
\providecommand{\Spec}{\operatorname{Spec}}
\providecommand{\Imf}{\operatorname{Im}}
\providecommand{\Ker}{\operatorname{Ker}}
\providecommand{\Id}{\operatorname{Id}}
\providecommand{\rad}{\operatorname{rad}}
\providecommand{\cO}{\mathcal O}
\providecommand{\cS}{\mathcal S}
\providecommand{\cJ}{\mathcal J}
\providecommand{\cN}{\mathcal N}
\providecommand{\bfP}{\mathbf P}
\providecommand{\bfA}{\mathbf A}


\sgaStableTarget{sga6.e07.expose}\section*{Exposé VII. Regular immersions and the computation of \texorpdfstring{$K^\bullet$}{K} of a blow-up scheme}
\begin{center}
by Pierre Berthelot
\end{center}

\sgaStableTarget{sga6.e07.sec.1}\section*{1. Generalities on regular immersions}

The classical notion of regular immersion (EGA IV 16.9.2) proves inadequate in the non-noetherian case. We shall define a weaker notion, one which gives the desired properties, and which is equivalent to the usual definition in the noetherian case.

\sgaStableTarget{sga6.e07.def.1-1}\textbf{Definition 1.1.} Let \(A\) be a commutative ring, let \(E\) be a projective \(A\)-module of finite type, and let \(u:E\to A\) be an \(A\)-linear homomorphism. We say that \(u\) is \emph{regular} if the Koszul complex \(\Kosz(u)\) defined by \(u\) is acyclic in dimensions \(>0\), and hence is a resolution of \(B=A/u(E)\).

Recall that the Koszul complex associated with \(u\) is the graded \(A\)-module \(\bigwedge(E)\), with boundary the interior multiplication by \(u\), regarded as an element of the dual \(\check E\). This boundary is an antiderivation of degree \(-1\), equal to \(u\) on \(E\).

\sgaStableTarget{sga6.e07.lem.1-2}\textbf{Lemma 1.2.} Let \(u:E\to A\) be regular. Let \(A'\) be an \(A\)-algebra and let \(u':E'\to A'\) be obtained from \(u\) by base change; suppose that \(u'\) is regular. Put \(J=u(E)\), \(J'=JA'=u(E')\). For every \(n\ge0\) one has
\sgaStableTarget{sga6.e07.eq.1-2-1}\sgaStableTarget{sga6.e07.eq.1-2-2}\sgaStableTarget{sga6.e07.eq.1-2-3}\sgaStableTarget{sga6.e07.eq.1-2-4}\begin{align}
 J^n/J^{n+1}\otimes_AA'&\xrightarrow{\sim}J'^n/J'^{n+1}, \tag{1.2.1}\\
 \Tor_i^A(J^n/J^{n+1},A')&=0\quad(i\ge1), \tag{1.2.2}\\
 J^n\otimes_AA'&\xrightarrow{\sim}J'^n, \tag{1.2.3}\\
 \Tor_i^A(J^n,A')&=0\quad(i\ge1). \tag{1.2.4}
\end{align}

First assume that the modules \(J^n/J^{n+1}\) are projective over \(B=A/J\). Since \(J\) is of finite type they are locally free of finite type over \(B\). The assertions are checked pointwise on \(\Spec(A)\), so that one may suppose \(A\) local and \(u(E)\subset\rad(A)\). Then \(J^n/J^{n+1}\) is free over \(B\), and \hyperlink{sga6.e07.eq.1-2-2}{(1.2.2)} is equivalent to \(\Tor_i^A(B,A')=0\) for \(i\ge1\). This follows by computing with the Koszul resolution \(\Kosz(u)\), since after tensoring with \(A'\) it becomes \(\Kosz(u')\), acyclic in positive degrees.

The exact sequences
\[
0\to J^n/J^{n+1}\to A/J^{n+1}\to A/J^n\to0,
\qquad
0\to J^n\to A\to A/J^n\to0
\]
then give \hyperlink{sga6.e07.eq.1-2-1}{(1.2.1)}, \hyperlink{sga6.e07.eq.1-2-3}{(1.2.3)}, and \hyperlink{sga6.e07.eq.1-2-4}{(1.2.4)} by induction on \(n\). It remains to know that \(J^n/J^{n+1}\) is projective over \(B\). Locally one reduces to \(E=A^d\), with \(u\) given by \(f_1,\ldots,f_d\), and to the universal homomorphism \(\mathbb Z[T_1,\ldots,T_d]^d\to \mathbb Z[T_1,\ldots,T_d]\). Since the sequence \((T_1,\ldots,T_d)\) is regular, the conclusion follows by base change.

\sgaStableTarget{sga6.e07.prop.1-3}\textbf{Proposition 1.3.} Let \(A\) be a ring, \(E\) a projective \(A\)-module of finite type, \(u:E\to A\) a regular homomorphism, and \(J=u(E)\). Then \(J\) is quasi-regular: \(J\) is of finite type, \(J/J^2\) is a projective \(A/J\)-module of finite type, and the canonical homomorphism
\sgaStableTarget{sga6.e07.eq.1-3-1}\begin{equation}\tag{1.3.1}
 \Sym_{A/J}(J/J^2)\longrightarrow \gr_J(A)=\bigoplus_{n\ge0}J^n/J^{n+1}
\end{equation}
is an isomorphism. This is obtained from the preceding proof, again by reduction to the universal polynomial case.

\sgaStableTarget{sga6.e07.def.1-4}\textbf{Definition 1.4.} Let \((X,A)\) be a commutatively locally ringed site, let \(E\) be an \(A\)-module locally free of finite type, and let \(u:E\to A\) be \(A\)-linear. We say that \(u\) is regular if \(\Kosz(u)\) is acyclic in degrees \(>0\). If \(J\) is an ideal of \(A\), we say that \(J\) is regular if, locally on \(X\), it is the image of a surjective regular homomorphism from a free \(A\)-module of finite type.

Let \(i:Y\to X\) be an immersion of preschemes. If \(U\) is an open subset of \(X\) containing \(i(Y)\), and if \(i\) is a closed immersion into \(U\), then \(i\) is said to be regular if the ideal defining \(i(Y)\) in \(\cO_U\) is regular. This condition is independent of \(U\), and regularity may be tested in a neighbourhood of every point.

The definition replaces here that of EGA IV 16.9.2, and in the noetherian case the two definitions are equivalent. A closed immersion locally defined by a regular sequence is regular, and a regular immersion is quasi-regular by \hyperlink{sga6.e07.prop.1-3}{Proposition 1.3}. A minimal system of generators of an ideal \(J\) at a point is a family whose images form a basis of \(J_x\otimes k(x)\); regularity is equivalently the local acyclicity of the corresponding Koszul complex. If \(J\) is of finite type and \(J/J^2\) is locally free, these systems correspond to local bases of \(J/J^2\).

\sgaStableTarget{sga6.e07.prop.1-5}\textbf{Proposition 1.5.} If \(i:Y\to X\) is regular, it remains regular after every flat base change. Conversely, \(i\) is regular if and only if it is so locally for the faithfully flat quasi-compact topology.

\sgaStableTarget{sga6.e07.def.1-6}\textbf{Definition 1.6.} If \(i:Y\to X\) is defined by an ideal \(J\), the \(\cO_Y\)-module \(J/J^2\) is the conormal sheaf of \(Y\) in \(X\), denoted \(N_{X/Y}\). If \(i\) is regular it is locally free of finite type, and its rank is the codimension. For closed immersions \(Z\xrightarrow{j}Y\xrightarrow{i}X\) one has the exact conormal sequence
\sgaStableTarget{sga6.e07.eq.1-6-1}\begin{equation}\tag{1.6.1}
 j^*N_{X/Y}\longrightarrow N_{X/Z}\longrightarrow N_{Y/Z}\longrightarrow0.
\end{equation}

\sgaStableTarget{sga6.e07.prop.1-7}\textbf{Proposition 1.7.} Let \(i:Y\to X\) and \(j:Z\to Y\) be immersions. If \(i\) and \(j\) are regular, then \(i\circ j\) is regular and the conormal sequence, completed by \(0\) on the left, is exact. If \(i\circ j\) and \(i\) are regular and if this sequence is exact and locally split, then \(j\) is regular. If \(X\) is noetherian and if \(i\circ j\) and \(j\) are regular, then \(i\) is regular at the points of \(Z\).

Locally, write \(X=\Spec(A)\), \(Y=\Spec(B)\), \(Z=\Spec(C)\), with \(B=A/I\) and \(C=B/K\). Choosing liftings \(f_1,\ldots,f_n\) of a basis of \(I/I^2\) and \(g_1,\ldots,g_m\) of a basis of \(K/K^2\), one obtains the diagram
\sgaStableTarget{sga6.e07.eq.1-7-1}\begin{equation}\tag{1.7.1}
\begin{tikzcd}[column sep=large,row sep=small]
A_0=\mathbb Z[T_1,\ldots,T_{n+m}] \arrow[r] \arrow[d] & A \arrow[d]\\
B_0=\mathbb Z[T_1,\ldots,T_m] \arrow[r] \arrow[d] & B \arrow[d]\\
C_0=\mathbb Z \arrow[r] & C,
\end{tikzcd}
\end{equation}
Therefore:
\sgaStableTarget{sga6.e07.eq.1-7-2}\begin{equation}\tag{1.7.2}
\begin{aligned}
 i\ \text{regular} &\Longleftrightarrow \forall i\ge1,\quad
 \Tor_i^{A_0}(B_0,A)=0,\\
 i\circ j\ \text{regular} &\Longleftrightarrow \forall i\ge1,\quad
 \Tor_i^{A_0}(C_0,A)=0,\\
 j\ \text{regular} &\Longleftrightarrow \forall i\ge1,\quad
 \Tor_i^{B_0}(C_0,B)=0.
\end{aligned}
\end{equation}
One then uses the transitivity spectral sequence
\sgaStableTarget{sga6.e07.eq.1-7-3}\begin{equation}\tag{1.7.3}
E^2_{pq}=\Tor_p^{B_0}\bigl(C_0,\Tor_q^{A_0}(B_0,A)\bigr)\Rightarrow\Tor_n^{A_0}(C_0,A).
\end{equation}
This gives the first two assertions; the noetherian assertion is that of EGA IV 19.1.5, with the definitions identified in the noetherian case.

\sgaStableTarget{sga6.e07.prop.1-8}\textbf{Proposition 1.8.} Let \(X\) be a scheme, \(E\) a locally free \(\cO_X\)-module of finite type, \(u:E\to\cO_X\) a regular homomorphism, and \(J=u(E)\). Let \(Y\) denote the closed subscheme of \(X\) defined by \(J\), let \(X'\) be the scheme obtained by blowing up \(Y\), let \(P=\bfP(E)\) be the projective bundle associated with \(E\), and let \(p:P\to X\) be the structural morphism. On \(P\) consider the canonical exact sequence
\sgaStableTarget{sga6.e07.eq.1-8-1}\begin{equation}\tag{1.8.1}
0\to F\to p^*E\to\cO_P(1)\to0,
\end{equation}
and the composite morphism
\[
  v:F\longrightarrow p^*E\xrightarrow{p^*(u)}\cO_P .
\]
Then:
\begin{enumerate}
\item the formation of \(X'\) commutes with every base change \(f:T\to X\) such that \(f^*(u)\) is regular;
\item the closed immersion \(j:X'\to P\) defined by the surjective homomorphism
\[
  \cS(E)\longrightarrow\bigoplus_{n\ge0}J^n
\]
defined by \(u\) in degree \(1\) is a regular immersion, its ideal being the image of the regular morphism \(v\). In particular, if \(\operatorname{codim}(i)\) is constant, then \(\operatorname{codim}(j)=\operatorname{codim}(i)-1\);
\item the \(\cS(E)\)-module \(\bigoplus_{n\ge0}J^n\) is perfect as a graded \(\cS(E)\)-module, i.e. locally there is a finite graded resolution of \(\bigoplus_{n\ge0}J^n\) by graded \(\cS(E)\)-modules locally free of finite type.
\end{enumerate}

For (i), since \(X'=\Proj(\bigoplus_{n\ge0}J^n)\), the result is an immediate consequence of the relations (1.2.3) and of EGA II.

For (ii), the asserted properties being local on \(X\), one may suppose that \(X=\Spec(A)\), and moreover that \(E\) is free on \(X\). We first prove the assertion in the standard case
\[
  A=\mathbb Z[T_1,\ldots,T_d],
\]
where \(E\) is given the basis \(e_1,\ldots,e_d\), with \(u(e_i)=T_i\). Then
\[
  P=\Proj\bigl(\mathbb Z[T_1,\ldots,T_d,X_1,\ldots,X_d]\bigr),
\]
the grading being taken with respect to the \(X_i\). To show that \(v\) is a regular morphism whose image is the ideal of \(j\), it suffices to do so locally on \(P\), hence on the open sets \(D_+(X_i)\); let us place ourselves, for instance, on \(D_+(X_1)\). Then \(p^*(E)\) is defined by the module
\[
  \mathbb Z[T_1,\ldots,T_d,X_2/X_1,\ldots,X_d/X_1]\otimes E,
\]
and \(\cO_P(1)\) by
\[
  \mathbb Z[T_1,\ldots,T_d,X_2/X_1,\ldots,X_d/X_1],
\]
the canonical morphism sending \(e_i\) to \(X_i/X_1\) for every \(i\). Thus \(F\) has as a basis the elements
\[
  e_i-e_1X_i/X_1\qquad (i>1).
\]
One has
\[
  v(e_i-e_1X_i/X_1)=T_i-T_1X_i/X_1;
  \qquad T'_i=T_i-T_1X_i/X_1.
\]
Since the sequence of the \(T_i\) is regular in
\[
\mathbb Z[T_1,\ldots,T_d,X_2/X_1,\ldots,X_d/X_1],
\]
the same is true of the sequence \(T_1,T'_2,\ldots,T'_d\), which is obtained from it by an invertible matrix, and therefore also of the extracted sequence \(T'_2,\ldots,T'_d\).

It remains to show that the \(T'_i\) generate the ideal of \(j\). In \(\bigoplus_{n\ge0}J^n\), for \(m\le p\), we shall write \(x^{(m)}\) for the element \(x\) of \(J^p\), considered as being of degree \(m\). Over \(D_+(X_1)\),
\[
  X'=\Spec\left(\left(\bigoplus_{n\ge0}J^n\right)_{(T_1^{(1)})}\right),
\]
and \(j\) is defined by the morphism sending \(X_i\) to \(T_i^{(1)}\). In this morphism the image of \(T'_i\) is
\[
T_i^{(0)}-T_1^{(0)}T_i^{(1)}/T_1^{(1)}
=\bigl((T_iT_1)^{(1)}-(T_1T_i)^{(1)}\bigr)/T_1^{(1)}=0.
\]
Thus the \(T'_i\) belong to the ideal of \(j\). Since every element of
\[
\mathbb Z[T_1,\ldots,T_d,X_2/X_1,\ldots,X_d/X_1]
\]
can be expressed as a polynomial in \(T_1,T'_2,\ldots,T'_d,X_2/X_1,\ldots,X_d/X_1\), it suffices, in order to show that the \(T'_i\) generate the ideal of \(j\), to show that no non-zero element of
\[
\mathbb Z[T_1,X_2/X_1,\ldots,X_d/X_1]
\]
belongs to the ideal of \(j\). This amounts to showing that the homomorphism
\[
\mathbb Z[T_1,X_1,\ldots,X_d]\longrightarrow\bigoplus_{n\ge0}J^n
\]
is injective, since \(T_1^{(1)}\) is not a zero divisor in \(\bigoplus_{n\ge0}J^n\). If one grades \(\mathbb Z[T_1,X_1,\ldots,X_d]\) by assigning degree \(0\) to \(T_1\) and degree \(1\) to the \(X_i\), the kernel of the homomorphism in question is a graded ideal. But it is contained in the kernel of the composite homomorphism
\[
\mathbb Z[T_1,X_1,\ldots,X_d]\longrightarrow
\bigoplus_{n\ge0}J^n\longrightarrow A,
\]
the last homomorphism being defined in degree \(n\) by the injection of \(J^n\) into \(A\). The kernel of the composite homomorphism is the principal ideal generated by \(X_1-T_1\), which, by integrality, contains no non-zero homogeneous element. Therefore the first homomorphism is injective, which completes the proof in this case.

Return to the general case where \(A\) is arbitrary, with \(E\) still assumed free on \(X\). Put index \(0\) on the special case just treated. By part (i), the situation under consideration is obtained from the preceding one by the base change \(A_0\to A\), where the homomorphism \(A_0\to A\) is defined by choosing a basis of \(E\), its image \(f_1,\ldots,f_d\) in \(J\), and sending \(T_i\) to \(f_i\). Put
\[
  B_0=\left(\bigoplus_{n\ge0}J_0^n\right)_{(T_1^{(1)})},\qquad
  B=\left(\bigoplus_{n\ge0}J^n\right)_{(f_1^{(1)})}.
\]
Then \(B=B_0\otimes_{A_0}A\), and the Koszul complex defined by \(v:F\to\cO_P\) is obtained by base change from the Koszul complex defined by \(v_0\). Since the latter is a projective resolution of \(B_0\) on a suitable open set, the homology groups of \(K_\bullet(v)\) are, on this open set, the groups
\[
\Tor_i^{A_0}(A,B_0).
\]
These are the components of degree \(0\) of
\[
\Tor_i^{A_0}\left(A,\left(\bigoplus_{n\ge0}J_0^n\right)_{T_1^{(1)}}\right)
=\left(\bigoplus_{n\ge0}\Tor_i^{A_0}(A,J_0^n)\right)_{T_1^{(1)}}.
\]
But the \(\Tor_i^{A_0}(A,J_0^n)\) vanish for \(i\ge1\) by the relations \hyperlink{sga6.e07.lem.1-2}{1.2}; hence \(K_\bullet(v)\) is a resolution of \(B\), which shows that \(j\) is regular.

For (iii), the property being local on \(X\), one may again suppose \(X\) affine and \(E\) free of finite type. We then resume the notation of (ii). First consider the standard case of (ii). The ring \(A_0\) is noetherian and of finite cohomological dimension, and so is the algebra \(\cS(E_0)\). One may then construct, step by step, a partial graded resolution of \(\bigoplus_{n\ge0}J_0^n\) by graded free \(\cS(E_0)\)-modules of finite type, of length greater than the cohomological dimension of \(A_0\). The last kernel is then a graded \(\cS(E_0)\)-module, projective of finite type as an ungraded \(\cS(E_0)\)-module, hence locally free of finite type on \(\Spec(A_0)\) as a graded \(\cS(E_0)\)-module by \hyperlink{sga6.e06.lem.1-5-1}{VI 1.5.1}. In this way one obtains the announced finite resolution, replacing the last module of the partial resolution by the corresponding kernel.

Tensoring the resolution thus constructed by \(A\), one obtains a graded complex of graded \(\cS(E)\)-modules locally free of finite type, and it remains to see that this is a resolution of \(\bigoplus_{n\ge0}J^n\). But the \(i\)-th homology module of the resulting complex is a graded \(\cS(E)\)-module whose \(n\)-th component is \(\Tor_i^{A_0}(A,J_0^n)\), which is zero for \(i\ge1\) by \hyperlink{sga6.e07.eq.1-2-4}{(1.2.4)}. Consequently, by \hyperlink{sga6.e07.eq.1-2-3}{(1.2.3)}, one has thus constructed a resolution of the required type.

\sgaStableTarget{sga6.e07.prop.1-9}\textbf{Proposition 1.9.} If \(i:Y\to X\) is a regular closed immersion, if \(X'\) is the blow-up of \(Y\), and if \(f:X'\to X\) is structural, then \(i\) and \(f\) are perfect morphisms: \(\Kosz\) resolves \(\cO_Y\) by finite locally free \(\cO_X\)-modules, and \(f\) factors through a regular immersion into a projective bundle followed by a smooth morphism.

\sgaStableTarget{sga6.e07.prop.1-10}\textbf{Proposition 1.10.} If \(X\) and \(Y\) are smooth over \(S\), every \(S\)-immersion from \(X\) into \(Y\) is regular. In particular, every section of a smooth morphism is a regular immersion.

\sgaStableTarget{sga6.e07.sec.2}\section*{2. Computations on regular immersions}

\sgaStableTarget{sga6.e07.para.2-1}\textbf{2.1--2.3.} For two \(A\)-algebras \(B,C\), the graded module \(\bigoplus_{i\ge0}\Tor_i^A(B,C)\) is given a natural structure of associative graded \(B\otimes_A C\)-algebra by choosing a projective resolution \(P_\bullet\to B\), using the Kunneth morphisms, and extending the augmentation \(B\otimes_A B\to B\). If the resolution is an alternating differential graded algebra, the Tor algebra is alternating. Globally this gives, for two \(X\)-schemes \(X',X''\), a graded algebra structure on
\[
\bigoplus_i\Tor_i^{\cO_X}(\cO_{X'},\cO_{X''}).
\]
The multiplication is represented by the commutative square of complexes
\sgaStableTarget{sga6.e07.eq.2-2-1}\begin{equation}\tag{2.2.1}
\begin{tikzcd}[column sep=large,row sep=large]
P_i\otimes P_j \arrow[r] \arrow[d] & (P_{i-1}\otimes P_j)\oplus(P_i\otimes P_{j-1}) \arrow[d]\\
P_{i+j} \arrow[r] & P_{i+j-1}.
\end{tikzcd}
\end{equation}

\sgaStableTarget{sga6.e07.lem.2-4}\textbf{Lemma 2.4.} For two \(X\)-schemes \(X',X''\) and a flat \(\cO_{X'}\)-module \(M\), one has canonically
\sgaStableTarget{sga6.e07.eq.2-4-1}\begin{equation}\tag{2.4.1}
\Tor_i^{\cO_X}(M,\cO_{X''})\simeq
M\otimes_{\cO_{X'}}\Tor_i^{\cO_X}(\cO_{X'},\cO_{X''}).
\end{equation}
This is immediate locally from a projective resolution, and then glues.

% SGA-SRC-001489 — reconstructed by mathematical inference; scan inconclusive
\sgaStableTarget{sga6.e07.prop.2-5}\textbf{Proposition 2.5.} Let \(Y\) and \(Z\) be closed subschemes of \(X\), with \(T=Y\times_XZ\). Under the local Koszul hypotheses of this proposition, one has
\sgaStableTarget{sga6.e07.eq.2-5-1}\begin{equation}\tag{2.5.1}
\Tor_i^{\cO_X}(\cO_Y,\cO_Z)\simeq \bigwedge^i\Tor_1^{\cO_X}(\cO_Y,\cO_Z),
\end{equation}
and
\sgaStableTarget{sga6.e07.eq.2-5-2}\begin{equation}\tag{2.5.2}
\Tor_1^{\cO_X}(\cO_Y,\cO_Z)
\simeq I\cap J/IJ
\simeq \Ker\bigl(i^*N_{X/Y}\times j^*N_{X/Z}\to N_{X/T}\bigr).
\end{equation}
In particular, for a regular closed immersion \(Y\to X\),
\sgaStableTarget{sga6.e07.eq.2-5-3}\begin{equation}\tag{2.5.3}
\Tor_i^{\cO_X}(\cO_Y,\cO_Y)\simeq
\bigwedge^i\Tor_1^{\cO_X}(\cO_Y,\cO_Y),
\qquad
\Tor_1^{\cO_X}(\cO_Y,\cO_Y)\simeq N_{X/Y}.
\end{equation}
The proof is the standard comparison of Koszul complexes, together with the bicomplex spectral sequence and the preceding multiplicative construction.

\sgaStableTarget{sga6.e07.para.2-6}\textbf{2.6--2.7.} Assume that \(X\) has an ample sheaf, so that the two definitions of \(K^\bullet\) coincide. If \(i:Y\to X\) is a regular closed immersion, one defines
\[
i_*:K^\bullet(Y)\to K^\bullet(X)
\]
using the class of the direct image, and one has, for every \(x\in K^\bullet(Y)\),
\sgaStableTarget{sga6.e07.eq.2-7-1}\begin{equation}\tag{2.7.1}
i^*i_*(x)=x\lambda_{-1}(N),
\qquad
\lambda_{-1}(N)=\sum_{r\ge0}(-1)^r\lambda^r(N),
\end{equation}
where \(N\) is the conormal sheaf. Indeed, if \(x\) is represented by a locally free finite type \(\cO_Y\)-module \(M\), the homology sheaves of the derived inverse image of \(i_*M\) are \(M\otimes\Tor_i^{\cO_X}(\cO_Y,\cO_Y)\), and \hyperlink{sga6.e07.prop.2-5}{Proposition 2.5} gives the displayed formula.

\sgaStableTarget{sga6.e07.cor.2-8}\textbf{Corollary 2.8.} Under the hypotheses of \hyperlink{sga6.e07.para.2-6}{2.7}, for \(x,y\in K^\bullet(Y)\),
\sgaStableTarget{sga6.e07.eq.2-8-1}\begin{equation}\tag{2.8.1}
i_*(xy\lambda_{-1}(N))=i_*(x)i_*(y).
\end{equation}
This follows from \hyperlink{sga6.e07.eq.2-7-1}{(2.7.1)} and the projection formula.

\sgaStableTarget{sga6.e07.cor.2-9}\textbf{Corollary 2.9.} Under the same hypotheses, if \(i_*(1)\) is invertible in \(K^\bullet(X)\), then \(i_*\) is an isomorphism of \(K^\bullet(Y)\) onto a direct summand.

\sgaStableTarget{sga6.e07.sec.3}\section*{3. Computation of \texorpdfstring{$K^\bullet$}{K} of a blow-up scheme}

\sgaStableTarget{sga6.e07.para.3-1}\textbf{3.1.} Let \(i:Y\to X\) be a regular closed immersion. Let \(X'\) be the blow-up of \(Y\), and let \(Y'=Y\times_X X'\). By EGA II 8.1.8, \(Y'\) is a closed subscheme of \(X'\) defined by an ideal canonically isomorphic to \(\cO_{X'}(1)\). Thus \(j:Y'\to X'\) is a regular closed immersion of codimension \(1\). Let \(J\) be the ideal of \(Y\) in \(X\), and \(J'\) the ideal of \(Y'\) in \(X'\); then
\[
J'\simeq J\cO_{X'}\simeq\cO_{X'}(1).
\]
Writing \(f:X'\to X\) and \(g:Y'\to Y\) for the structural morphisms, one has the commutative diagram
\sgaStableTarget{sga6.e07.eq.3-1-1}\begin{equation}\tag{3.1.1}
\begin{tikzcd}[column sep=large,row sep=large]
Y' \arrow[r,"j"] \arrow[d,"g"'] & X' \arrow[d,"f"]\\
Y \arrow[r,"i"'] & X.
\end{tikzcd}
\end{equation}
Moreover,
\[
X'=\Proj(S),\qquad S=\bigoplus_{n\ge0}J^n,
\]
and
\[
Y'=\Proj\bigl(i^*S\bigr)
=\Proj\bigl(\bigoplus_{n\ge0}J^n/J^{n+1}\bigr).
\]
Since \(i\) is regular, \(\bigoplus J^n/J^{n+1}\simeq\Sym_{\cO_Y}(J/J^2)\). If \(N=J/J^2\), then \(Y'=\bfP(N)\). On \(Y'\) there is the exact sequence
\sgaStableTarget{sga6.e07.eq.3-1-2}\begin{equation}\tag{3.1.2}
0\to F\to g^*N\to\cO_{Y'}(1)\to0.
\end{equation}
The sheaf \(\cO_{Y'}(1)=j^*\cO_{X'}(1)=J'/J'^2\) is the conormal sheaf of \(j\); denote it by \(L\).

\sgaStableTarget{sga6.e07.lem.3-2}\textbf{Lemma 3.2.} With the hypotheses and notation of \hyperlink{sga6.e07.para.3-1}{3.1}, one has
\sgaStableTarget{sga6.e07.eq.3-2-i}\sgaStableTarget{sga6.e07.eq.3-2-ii}\begin{align}
\Tor_i^{\cO_X}(\cO_{X'},\cO_Y)&\simeq
\bigwedge^i\Tor_1^{\cO_X}(\cO_{X'},\cO_Y),\tag{3.2.i}\\
\Tor_1^{\cO_X}(\cO_{X'},\cO_Y)&\simeq F.\tag{3.2.ii}
\end{align}

For (i), since \(\cO_Y\) can locally be resolved by a Koszul complex, the algebra \(\bigoplus_i\Tor_i^{\cO_X}(\cO_{X'},\cO_Y)\) is alternating. Hence there is a unique algebra homomorphism
\sgaStableTarget{sga6.e07.eq.3-2-1}\begin{equation}\tag{3.2.1}
\bigwedge \Tor_1^{\cO_X}(\cO_{X'},\cO_Y)
\longrightarrow
\bigoplus_{i\ge0}\Tor_i^{\cO_X}(\cO_{X'},\cO_Y)
\end{equation}
extending the injection of degree one; it is enough to verify locally that it is an isomorphism. On an affine open, choose a regular homomorphism \(E\to A\) defining \(J\), with basis images \(f_1,\ldots,f_d\). On the open of \(X'\) where \(f_1^{(1)}\) is invertible, change the basis by
\[
e'_1=e_1,
\qquad
e'_i=e_i-(f_i^{(1)}/f_1^{(1)})e_1\quad(i\ne1).
\]
The Koszul differential then gives local isomorphisms
\sgaStableTarget{sga6.e07.eq.3-2-2}\begin{equation}\tag{3.2.2}
\Tor_i^{\cO_X}(\cO_{X'},\cO_Y)
\simeq
(\cO_{X'}/f_1^{(0)}\cO_{X'})\otimes_{\cO_{X'}}
\bigwedge^i(\cO_{X'}^{d-1})
\simeq
\bigwedge^i(\cO_{Y'}^{d-1}).
\end{equation}
In particular, for \(i=1\) one obtains \(\Tor_1^{\cO_X}(\cO_{X'},\cO_Y)\simeq \cO_{Y'}^{d-1}\), and hence
\sgaStableTarget{sga6.e07.eq.3-2-3}\begin{equation}\tag{3.2.3}
\Tor_i^{\cO_X}(\cO_{X'},\cO_Y)
\simeq
\bigwedge^i\Tor_1^{\cO_X}(\cO_{X'},\cO_Y).
\end{equation}
For (ii), tensor the exact sequence \(0\to J\to\cO_X\to\cO_Y\to0\) with \(\cO_{X'}\). Since the image of \(J\otimes\cO_{X'}\to\cO_{X'}\) is \(J\cO_{X'}=\cO_{X'}(1)\), and then tensoring with \(\cO_{Y'}\), one obtains
\sgaStableTarget{sga6.e07.eq.3-2-4}\begin{equation}\tag{3.2.4}
0\to \Tor_1^{\cO_X}(\cO_{X'},\cO_Y)\to g^*N\to\cO_{Y'}(1)\to0,
\end{equation}
which identifies this Tor sheaf with \(F\) in \hyperlink{sga6.e07.eq.3-1-2}{(3.1.2)}.

\sgaStableTarget{sga6.e07.para.3-3}\textbf{3.3.} Suppose now that \(X\) has an ample sheaf. Then the same is true for \(X'\), which is projective over \(X\), and for \(Y\) and \(Y'\). Moreover \(i\) and \(j\) are perfect morphisms, \(g\) is smooth hence perfect because \(Y'=\bfP(N)\), and \(f\) is perfect by 1.9. Since these morphisms are projective, one has additive homomorphisms
\[
i_*:K^\bullet(Y)\to K^\bullet(X),\quad
j_*:K^\bullet(Y')\to K^\bullet(X'),\quad
f_*:K^\bullet(X')\to K^\bullet(X),\quad
g_*:K^\bullet(Y')\to K^\bullet(Y).
\]

\sgaStableTarget{sga6.e07.prop.3-4}\textbf{Proposition 3.4.} With the hypotheses and notation of \hyperlink{sga6.e07.para.3-1}{3.1} and \hyperlink{sga6.e07.para.3-3}{3.3}, for every \(x\in K^\bullet(Y)\) one has
\sgaStableTarget{sga6.e07.eq.3-4-1}\begin{equation}\tag{3.4.1}
f^*(i_*(x))=j_*\bigl(g^*(x)\lambda_{-1}(F)\bigr).
\end{equation}
By additivity, one may assume \(x\) represented by a locally free finite type \(\cO_Y\)-module \(M\). The class of \(f^*i_*(x)\) is represented by the complex obtained from a projective resolution of \(i_*M\) and application of \(f^*\). Its homology sheaves are \(\Tor_i^{\cO_X}(M,\cO_{X'})\). By \hyperlink{sga6.e07.lem.2-4}{2.4} and \hyperlink{sga6.e07.lem.3-2}{3.2},
\[
\Tor_i^{\cO_X}(M,\cO_{X'})
\simeq M\otimes_{\cO_Y}\Tor_i^{\cO_X}(\cO_{X'},\cO_Y)
\simeq M\otimes_{\cO_Y}\bigwedge^iF,
\]
and the alternating sum of their classes is exactly the right hand side of \hyperlink{sga6.e07.eq.3-4-1}{(3.4.1)}.

\sgaStableTarget{sga6.e07.lem.3-5}\textbf{Lemma 3.5.} With the hypotheses and notation of \hyperlink{sga6.e07.para.3-1}{3.1}, one has the following isomorphisms:
\begin{enumerate}[label=\roman*)]
\item for all \(n\ge0\) and \(i\ge1\),
\[
R^if_*\bigl(\cO_{X'}(n)\bigr)=0;
\]
\item for all \(n\ge0\),
\[
f_*\bigl(\cO_{X'}(n)\bigr)=J^n.
\]
\end{enumerate}
One may suppose \(X\) affine, hence quasi-compact. For (i) one uses the following lemma: if \(f:X'\to X\) is projective and perfect, with \(X\) quasi-compact, and if \(L'\) is a perfect complex on \(X'\) acyclic in degrees \(>0\), then there exists \(n_0\) such that, for \(n\ge n_0\), \(Rf_*(L'\otimes\cO_{X'}(n))\) is acyclic in degrees \(>0\). Applied to \(\cO_{X'}\), this gives the vanishing for large \(n\). The exact sequence
\[
0\to \cO_{X'}(n+1)\to\cO_{X'}(n)\to\cO_{Y'}(n)\to0
\]
and the fact that \(Y'=\bfP(N)\) then allow descending induction to \(n=0\).

For (ii), the canonical graded homomorphism
\[
\bigoplus_{n\ge0}J^n\longrightarrow\bigoplus_{n\ge0}f_*\bigl(\cO_{X'}(n)\bigr)
\]
is shown to be an isomorphism locally. Write \(S=\bigoplus J^n\) and use a surjective graded homomorphism \(\cO_X[T_1,\ldots,T_d]\to S\) given by a set of generators of \(J\). This realizes \(X'\) as a regular closed subscheme of \(\bfP_X^{d-1}\). With a finite graded presentation \(L''\to L'\to S\to0\), one obtains on \(\bfP_X^{d-1}\) a presentation \(M''\to M'\to\cO_{X'}\to0\). For each \(n\) there is a commutative diagram
\[
\begin{tikzcd}[column sep=large,row sep=large]
L''_n \arrow[r] \arrow[d,"\sim"'] & L'_n \arrow[r] \arrow[d,"\sim"'] & J^n \arrow[r] \arrow[d] & 0 \arrow[d]\\
p_*M''(n) \arrow[r] & p_*M'(n) \arrow[r] & p_*\cO_{X'}(n) \arrow[r] & 0,
\end{tikzcd}
\]
where \(p:\bfP_X^{d-1}\to X\). For \(n\) sufficiently large the bottom row is exact; hence the middle vertical arrow is an isomorphism for large \(n\). Descending induction, using
\[
0\to J^{n+1}\to J^n\to J^n/J^{n+1}\to0
\]
and the analogous direct-image exact sequence, completes the proof.

\sgaStableTarget{sga6.e07.prop.3-6}\textbf{Proposition 3.6.} With the hypotheses and notation of \hyperlink{sga6.e07.para.3-1}{3.1} and \hyperlink{sga6.e07.para.3-3}{3.3},
\sgaStableTarget{sga6.e07.eq.3-6-1}\begin{equation}\tag{3.6.1}
f_*f^*=\Id_{K^\bullet(X)}.
\end{equation}
Indeed, by the projection formula, \(f_*f^*(x)=x f_*(1)\), and \hyperlink{sga6.e07.lem.3-5}{Lemma 3.5} says that \(Rf_*\cO_{X'}\) has \(\cO_X\) in degree \(0\) and no higher cohomology.

\sgaStableTarget{sga6.e07.thm.3-7}\textbf{Theorem 3.7.} Let \(i:Y\to X\) be a regular closed immersion, let \(X'\) be the blow-up of \(Y\), and let \(Y'=X'\times_XY\), so that one has
\[
\begin{tikzcd}[column sep=large,row sep=large]
Y' \arrow[r,"j"] \arrow[d,"g"'] & X' \arrow[d,"f"]\\
Y \arrow[r,"i"'] & X.
\end{tikzcd}
\]
Assume that \(X\) is quasi-compact and has an ample sheaf. Then the sequence
\sgaStableTarget{sga6.e07.eq.3-7-0}\begin{equation}\tag{3.7.0}
0\to K^\bullet(Y)\xrightarrow{u}K^\bullet(Y')\times K^\bullet(X)\xrightarrow{v}K^\bullet(X')
\end{equation}
is exact, where
\sgaStableTarget{sga6.e07.eq.3-7-u}\begin{equation}\tag{3.7.u}
u(y)=\bigl(g^*(y)\lambda_{-1}(F),-i_*(y)\bigr),
\end{equation}
and, for \(y'\in K^\bullet(Y')\), \(x\in K^\bullet(X)\),
\sgaStableTarget{sga6.e07.eq.3-7-v}\begin{equation}\tag{3.7.v}
v(y',x)=j_*(y')+f^*(x).
\end{equation}
Moreover \(u\) has a left inverse \(u'\), defined by \(u'(y',x)=g_*(y')\). If, in addition, \(X\) is noetherian and \(Y\) regular, then \(v\) is surjective and
\sgaStableTarget{sga6.e07.eq.3-7-2}\begin{equation}\tag{3.7.2}
0\to K^\bullet(Y)\xrightarrow{u}K^\bullet(Y')\times K^\bullet(X)\xrightarrow{v}K^\bullet(X')\to0
\end{equation}
is exact.

First, \(u'\circ u=1\), because \(g_*(g^*y\lambda_{-1}(F))=y\,g_*(\lambda_{-1}(F))=y\). Also \(v\circ u=0\), by \hyperlink{sga6.e07.prop.3-4}{Proposition 3.4}. If \((y',x)\) belongs to the kernel of \(v\), then \(j_*(y')+f^*(x)=0\). One must find \(y\in K^\bullet(Y)\) such that
\begin{equation}\tag{\sgaStableTargetInline{sga6.e07.eq.3-7-1}3.7.1}
x=-i_*(y),\qquad y'=g^*(y)\lambda_{-1}(F).
\end{equation}
Necessarily \(y=u'(y',x)=g_*(y')\). Applying \(f_*\) to the relation \(f^*(x)=-j_*(y')\), and using \hyperlink{sga6.e07.prop.3-6}{Proposition 3.6}, gives \(x=-i_*(y)\). To prove the second equality in \hyperlink{sga6.e07.eq.3-7-1}{(3.7.1)}, put \(z=y'-g^*(g_*y')\lambda_{-1}(F)\). Then \(g_*z=0\), and the preceding formulas give \(j_*z=0\). Since \(\Ker j_*\cap\Ker g_*=0\), it follows that \(z=0\). This proves exactness at the middle term.

For the final surjectivity assertion, assume \(X\) noetherian and \(Y\) regular. Then \(Y'\) is regular, being a projective bundle over \(Y\), and local rings of \(X'\) along \(Y'\) are regular because \(Y'\) is cut out by one regular element. Therefore coherent modules on \(X'\) with support in \(Y'\) are perfect. If such a module \(E\) has class in \(K^\bullet(X')\), filtration by the powers \(J'^nE\) shows that its class is in the image of \(j_*\). For arbitrary \(x\in K^\bullet(X')\), choose a perfect complex \(M'\) representing \(x\), and form the distinguished triangle
\[
\begin{tikzcd}[column sep=large,row sep=large]
& N' \arrow[dl] \arrow[dr] &\\
M' \arrow[rr] && Lf^*Rf_*(M').
\end{tikzcd}
\]
The horizontal arrow is the canonical morphism. It is an isomorphism off \(Y'\), so the cohomology sheaves of \(N'\) are coherent sheaves supported on \(Y'\). Their class lies in \(\Imf(j_*)\). Hence
\[
x-f^*f_*(x)\in\Imf(j_*),
\]
and \(v\) is surjective.



% Macros for Exposé VII continuation
\providecommand{\Q}{\mathbb Q}
\providecommand{\bfP}{\mathbf P}
\providecommand{\gr}{\operatorname{gr}}
\providecommand{\Id}{\operatorname{Id}}
\providecommand{\Fil}{\operatorname{Fil}}
\providecommand{\Imf}{\operatorname{Im}}
\providecommand{\rank}{\operatorname{rank}}
\providecommand{\codim}{\operatorname{codim}}
\providecommand{\cJ}{\mathcal J}
\providecommand{\cI}{\mathcal I}
\providecommand{\cN}{\mathcal N}
\providecommand{\cE}{\mathcal E}
\providecommand{\cO}{\mathcal O}


\sgaStableTarget{sga6.e07.remark-3-8}\subsection*{Remark 3.8}

In the case where \(X\) is noetherian and \(Y\) regular, \hyperlink{sga6.e07.thm.3-7}{3.7} determines the abelian group \(K^\bullet(X')\) entirely from the \(\lambda\)-rings \(K^\bullet(X)\), \(K^\bullet(Y)\), the maps \(i_*,i^*\) between these rings, and the element \(N\in K^\bullet(Y)\). Indeed, \(K^\bullet(Y')\) can be computed from \(K^\bullet(Y)\) and \(N\) by \hyperlink{sga6.e06.thm.1-1}{VI 1.1}. We shall see in the following paragraph that the multiplicative structure of \(K^\bullet(X')\) is determined by the same data, and so is the \(\lambda\)-structure after tensoring with \(\mathbb Q\) (this restriction is probably unnecessary; see \hyperlink{sga6.e07.rem.4-5}{4.5}).

\sgaStableTarget{sga6.e07.sec.4}\section*{4. Regular immersions and the filtration of \(K^\bullet\)}

Throughout this paragraph we keep the notation and hypotheses of 3.1.

\sgaStableTarget{sga6.e07.lem.4-1}\textbf{Lemma 4.1.} Suppose that there exists on \(Y\) a locally free sheaf \(N'\) of finite type such that, if \(N,N'\) denote the classes in \(K^\bullet(Y)\) of \(N,N'\), then \(N=N'+2\). Then the following congruence is verified:
\[
  \lambda_{-1}(F)\equiv 0 \pmod{1-L}.
\]

Let \(d\) be the rank of \(N\), which we may suppose constant; the rank of \(F\) is then \(d-1\). By definition,
\[
\begin{aligned}
\lambda_{-1}(F)
 &= \sum_{i\ge0}(-1)^i\lambda^i(F)\\
 &= (-1)^{d-1}\lambda^{d-1}(F-1)\\
 &= (-1)^{d-1}\lambda^{d-1}(F+L-1-L)\\
 &= (-1)^{d-1}\lambda^{d-1}(N-2+1-L).
\end{aligned}
\]
Now
\[
  \lambda_t(1-L)=\frac{\lambda_t(1)}{\lambda_t(L)}
  =\frac{1+t}{1+Lt}.
\]
Thus, for every \(i\ge1\), \(\lambda^i(1-L)\) is divisible by \(1-L\). Consequently
\[
\begin{aligned}
\lambda_{-1}(F)
  &\equiv (-1)^{d-1}\lambda^{d-1}(N-2) \pmod{1-L}\\
  &=(-1)^{d-1}\lambda^{d-1}(N') \pmod{1-L}.
\end{aligned}
\]
Since \(N'\) has rank \(d-2\), \(\lambda^{d-1}(N')=0\), proving the result.

\sgaStableTarget{sga6.e07.sec.4.2}\subsection*{4.2}

Let \(i:Y\to X\) be a regular immersion. Let \(r\ge1\) be an integer and let
\[
  s:X\longrightarrow \mathbf P_X^r
\]
be the section defined by a homomorphism \((\mathcal O_X)^{r+1}\to\mathcal O_X\), after passing to the quotient by the sub-bundle generated by the corresponding sections. If \(\mathbf P_X^r\) is defined by the graded \(\mathcal O_X\)-algebra \(\mathcal O_X[T_0,\ldots,T_r]\), the ideal \(I\) defining the section \(s\) is generated by the sections \(T_1,\ldots,T_r\). Therefore \(I/I^2\) is isomorphic to \((\mathcal O_X[T_0])^r\), so the ideal corresponding to \(\mathcal O_{\mathbf P^r}\) is \((\mathcal O_X)^r\).

Since the section \(s\) is evidently a regular immersion, the same is true for \(s\circ i\), and one may write the corresponding exact sequence of conormal sheaves \hyperlink{sga6.e07.prop.1-7}{(1.7)}. Thus, if \(N\) denotes the conormal sheaf of \(i\), and \(N_1\) the conormal sheaf of \(s\circ i\), one obtains the equality
\[
  N_1=N+r.
\]
It follows that every regular immersion \(i\) can be made to satisfy the property of \hyperlink{sga6.e07.lem.4-1}{Lemma 4.1} by composing it with a section \(X\to\mathbf P_X^2\), or with a section \(X\to\mathbf P_X^1\).

\sgaStableTarget{sga6.e07.thm.4-3}\textbf{Theorem 4.3.} Let \(i:Y\to X\) be a regular closed immersion, and let \(N\) be the class in \(K^\bullet(Y)\) of the conormal sheaf of \(Y\) in \(X\). Suppose that \(X\) is quasi-compact and has an ample sheaf. Then, with the notation of \hyperlink{sga6.e07.para.3-3}{3.3}, for every \(y\in K^\bullet(Y)\) and every \(p\ge1\),
\sgaStableTarget{sga6.e07.eq.4-3-1}\begin{equation}\tag{4.3.1}
 i_*\bigl(\lambda^p(N,y)\bigr)\equiv \lambda^p(i_*(y))
 \pmod{\text{torsion}},
\end{equation}
where \(\lambda^p(N,y)\) is the element defined in \hyperlink{sga6.e05.def.5-3}{V 5.3}.

First suppose that the hypothesis of \hyperlink{sga6.e07.lem.4-1}{4.1} is satisfied; in that case we shall show
\[
  i_*\bigl(\lambda^p(N,y)\bigr)=\lambda^p(i_*(y)).
\]
Since \(f_*f^*=\Id\), this relation is equivalent to
\[
  f^*\bigl(i_*(\lambda^p(N,y))\bigr)=
  f^*\bigl(\lambda^p(i_*(y))\bigr).
\]
Now one has
\sgaStableTarget{sga6.e07.eq.3-4}\[
\begin{aligned}
 f^*\bigl(\lambda^p(i_*(y))\bigr)
 &=\lambda^p\bigl(f^*(i_*(y))\bigr)\\
 &=\lambda^p\bigl(j_*g^*(y)\lambda_{-1}(F)\bigr),
\end{aligned}
\tag{3.4}
\]
whereas
\sgaStableTarget{sga6.e07.eq.3-4.occ2}\[
  f^*\bigl(i_*(\lambda^p(N,y))\bigr)
  =
  j_*\bigl(g^*(\lambda^p(N,y))\lambda_{-1}(F)\bigr)
  \tag{3.4}
\]
by \hyperlink{sga6.e07.prop.3-4}{Proposition 3.4}. Since \(g^*(N)=F+L\), formula \hyperlink{sga6.e05.eq.5-4-2}{V 5.4.2} gives
\[
  g^*(\lambda^p(N,y))\lambda_{-1}(F)
  =
  \lambda^p(g^*(N),g^*(y))\lambda_{-1}(F)
  =
  \lambda^p(L,g^*(y)\lambda_{-1}(F)).
\]
Putting \(z=g^*(y)\lambda_{-1}(F)\), we are reduced to showing
\[
  j_*\bigl(\lambda^p(L,z)\bigr)=\lambda^p(j_*(z)).
\]
This is the analogous relation for the immersion \(j\) and the element \(z\). By hypothesis, \(\lambda_{-1}(F)\) is a multiple of \(1-L\), so one can find \(z'\) such that \(z=z'(1-L)\). Then
\[
  z=j_*(u')\quad\text{with }u'=g_*(z')
\]
by \hyperlink{sga6.e07.para.2-6}{(2.7)}, and since \(L=j^*(J')\), it remains only to verify
\[
  j_*\bigl(j^*(\lambda^p(J',u'))\bigr)=
  \lambda^p\bigl(j_*j^*(u')\bigr).
\]
The exact sequence
\[
  0\longrightarrow J'\longrightarrow \mathcal O_{X'}
  \longrightarrow \mathcal O_{Y'}\longrightarrow 0
\]
gives
\sgaStableTarget{sga6.e07.eq.4-3-2}\begin{equation}\tag{4.3.2}
  j_*j^*(1)=1-J'.
\end{equation}
The desired relation is therefore, after applying the projection formula,
\[
  \lambda^p(J',u')(1-J')=\lambda^p(u'(1-J')).
\]
Since \(1-J'=\lambda_{-1}(J')\), this is precisely the formula which serves to define \(\lambda^p(J',u')\).

We now pass to the general case. Compose \(i\) first with a section \(X\to\mathbf P_X^1\), and then with a section \(\mathbf P_X^1\to\mathbf P_{\mathbf P_X^1}^1\), in such a way that the composed immersion satisfies the condition of \hyperlink{sga6.e07.lem.4-1}{4.1} (see \hyperlink{sga6.e07.sec.4.2}{4.2}). \hyperlink{sga6.e07.thm.4-3}{Proposition 4.3} is then true for the immersion \(s\circ i\) which precedes it; it remains to descend to \(X\). We descend first to \(\mathbf P_X^1\), and then to \(X\), so that the proof is made by applying twice the following lemma, due to J.-P. Jouanolou.

\sgaStableTarget{sga6.e07.lem.unnamed}\textbf{Lemma.} Under the conditions of \hyperlink{sga6.e07.thm.4-3}{4.3}, let \(s:X\to\mathbf P_X^1\) be the section defined in 4.2. If \hyperlink{sga6.e07.thm.4-3}{Proposition 4.3} is true for \(s\circ i\) and for \(s\), then it is true for \(i\).

Suppose \hyperlink{sga6.e07.thm.4-3}{4.3} true for \(s\circ i\). Then, for every \(p\ge1\) and every \(y\in K^\bullet(Y)\),
\sgaStableTarget{sga6.e07.eq.4-3-3}\begin{equation}\tag{4.3.3}
 (s\circ i)_*\bigl(\lambda^p(N+1,y)\bigr)
 \equiv
 \lambda^p\bigl((s\circ i)_*(y)\bigr)
 \pmod{\text{torsion}}.
\end{equation}
The section \(s\) is defined by an ideal isomorphic to \(\mathcal O_{\mathbf P}(-1)\) \hyperlink{sga6.e07.sec.4.2}{(4.2)}. Let \(\xi\) be its class in \(K^\bullet(\mathbf P_X^1)\). Then, for every \(x\in K^\bullet(X)\),
\[
  s_*(x)=p^*(x)(1-\xi),
\]
where \(p\) is the structural morphism \(\mathbf P_X^1\to X\). Indeed, since \(p\circ s=\Id_X\), one has \(x=s^*(p^*(x))\), whence \(s_*(x)=s_*(s^*p^*(x))=p^*(x)s_*(1)\); and \(s_*(1)=1-\xi\). Thus
\[
\begin{aligned}
(s\circ i)_*\bigl(\lambda^p(N+1,y)\bigr)
&\equiv \lambda^p\bigl(p^*(i_*(y))(1-\xi)\bigr)\pmod{\text{torsion}}\\
&= \lambda^p(\xi,p^*(i_*(y)))(1-\xi)\pmod{\text{torsion}}.
\end{aligned}
\]
Now \(\lambda^p(\xi,p^*(i_*(y)))-\lambda^p(1,p^*(i_*(y)))\) is divisible by \(1-\xi\). To see this, one may work in the universal \(\lambda\)-ring used to define \(\lambda^p(\xi,x)\) (V 5). Then \(\lambda^p(\xi,x)-\lambda^p(1,x)\) is a polynomial in \(\xi\) which vanishes for \(\xi=1\), hence is divisible by \(\xi-1\). By the independence equation for \(\xi\) on \(K^\bullet(X)\), one has \((1-\xi)^2=0\) (\hyperlink{sga6.e06.eq.1-13-1}{VI 1.13.1}). Hence
\[
\begin{aligned}
(s\circ i)_*\bigl(\lambda^p(N+1,y)\bigr)
&\equiv \lambda^p(1,p^*(i_*(y)))(1-\xi) \pmod{\text{torsion}}\\
&=p^*(\lambda^p(1,i_*(y)))(1-\xi) \pmod{\text{torsion}}.
\end{aligned}
\]
Applying \(p_*\), we obtain
\sgaStableTarget{sga6.e07.eq.4-3-4}\begin{equation}\tag{4.3.4}
  i_*\bigl(\lambda^p(N+1,y)\bigr)
  \equiv \lambda^p(1,i_*(y)) \pmod{\text{torsion}},
\end{equation}
since \(p_*(\xi)=0\).

By formula \hyperlink{sga6.e05.para.5-6}{V 5.6}, applied with \(L=1\), \(\lambda^p(N+1,y)\) is the coefficient of \(t^p\) in the series
\[
  \frac{\sum_{k\ge1}\lambda^k(N,y)t^k}
       {1+\sum_{k\ge1}\lambda^k(N,y)\lambda_{-1}(N)t^k}.
\]
Similarly, the same formula applied to \(N=0\) and \(L=1\) shows that \(\lambda^p(1,y)\) is the coefficient of \(t^p\) in
\[
  \frac{\sum_{k\ge1}\lambda^k(y)t^k}
       {\sum_{k\ge0}\lambda^k(y)t^k}.
\]
The congruences \hyperlink{sga6.e07.eq.4-3-4}{(4.3.4)} can therefore be written
\[
i_*\left(
  \frac{\sum_{k\ge1}k\lambda^k(N,y)t^k}
       {1+\sum_{k\ge1}\lambda^k(N,y)\lambda_{-1}(N)t^k}
\right)
=
  \frac{\sum_{k\ge1}k\lambda^k(i_*(y))t^k}
       {\sum_{k\ge0}\lambda^k(i_*(y))t^k}
       \pmod{\text{torsion}}.
\]
Applying formula \hyperlink{sga6.e07.cor.2-8}{2.8}, one obtains
\[
  \frac{\sum_{k\ge0}k\,i_*(\lambda^k(N,y))t^k}
       {1+\sum_{k\ge1} i_*(\lambda^k(N,y))t^k}
  =
  \frac{\sum_{k\ge0}k\lambda^k(i_*(y))t^k}
       {\sum_{k\ge0}\lambda^k(i_*(y))t^k}
       \pmod{\text{torsion}}.
\]
If
\[
  u(t)=1+\sum_{k\ge1}i_*(\lambda^k(N,y))t^k,\qquad
  v(t)=\sum_{k\ge0}\lambda^k(i_*(y))t^k,
\]
this last relation can also be written
\[
  \frac{d}{dt}\!\left(\frac{u(t)}{v(t)}\right)\equiv0
  \pmod{\text{torsion}}.
\]
Since the constant term of \(u(t)/v(t)\) is \(1\), and since the congruences are modulo torsion, one may integrate and write the congruence \(u(t)=v(t)\). This gives the announced congruences.

\sgaStableTarget{sga6.e07.cor.4-4}\textbf{Corollary 4.4.} With the notation and hypotheses of \hyperlink{sga6.e07.thm.4-3}{4.3}, for every \(p\ge1\) and every \(y\in K^\bullet(Y)\), one has
\[
  i_*\bigl(\psi^p(N,y)\bigr)\equiv \psi^p(i_*(y))
  \pmod{\text{torsion}}.
\]
This follows immediately from \hyperlink{sga6.e07.thm.4-3}{4.3}, using the expression of the \(\psi^p\) as integral linear combinations of the \(\lambda^p\).

\sgaStableTarget{sga6.e07.rem.4-5}\textbf{Remark 4.5.} The proof of \hyperlink{sga6.e07.thm.4-3}{4.3} itself suggests that \hyperlink{sga6.e07.eq.4-3-1}{(4.3.1)} should be true not merely modulo torsion, but also as an equality. In \(\RRR\), the equality was proved in characteristic \(0\) by means of a geometric characterization of \(\lambda^p(N,x)\). An extension of this method to all characteristics would probably be possible if one knew how to define exterior-power operations on perfect complexes (cf. XIV).

\sgaStableTarget{sga6.e07.thm.4-6}\textbf{Theorem 4.6.} Let \(X\) be a quasi-compact scheme having an ample sheaf, let \(d>0\) be an integer, and let \(i:Y\to X\) be a regular immersion of codimension \(d\). Then, for every \(k\),
\[
 i_*\bigl((\Fil^kK^\bullet(Y))\otimes\mathbb Q\bigr)
 \subset
 (\Fil^{k+d}K^\bullet(X))\otimes\mathbb Q.
\]

First note that \(i_*(K^\bullet(Y))\subset \Fil^1K^\bullet(X)\). Indeed, \(X-Y\) is dense in \(X\); otherwise \(Y\) would contain an affine open, and on this open the ideal \(J\) defining \(Y\) would be a nilideal, hence a locally nilpotent ideal, since it is of finite type. This is impossible, since the \(J^n/J^{n+1}\) are non-zero locally free modules. The augmentation in \(K^\bullet(X)\) of an element of \(i_*(K^\bullet(Y))\) is a locally constant function on \(X\); since it vanishes on \(X-Y\), it vanishes on all of \(X\).

For every abelian group \(A\), write \(A_{\mathbb Q}=A\otimes\mathbb Q\). Let \(x\in\Fil^kK^\bullet(Y)_{\mathbb Q}\). Then \(x\) can be written in the form
\[
  x=\sum a\,\gamma^{i_1}(x_1)\cdots\gamma^{i_n}(x_n),
\]
with \(\epsilon(x_i)=0\) for all \(i\), \(a\in\mathbb Q\), and \(i_1+\cdots+i_n\ge k\). The \(x_i\) can be expressed as integral linear combinations of the classes \(E_\alpha\) of finitely many locally free finite-type \(\mathcal O_Y\)-modules \(E_\alpha\). Let \(R\) be the sub-\(\mathbb Q\)-\(\lambda\)-algebra of \(K^\bullet(Y)_{\mathbb Q}\) generated by the \(E_\alpha\) and by \(N\). Then \(x\in\Fil^kR\).

The filtration of the \(\lambda\)-ring \(R\) is discrete (\hyperlink{sga6.e06.prop.6-2}{VI 6.2}); choose \(m\) such that \(\Fil^mR=0\). Finally, after decomposing \(X\) as a finite sum of disjoint open subsets and proving the assertion on each such open, one may suppose that the \(E_\alpha\) have constant rank \(r_\alpha\). The augmentation ring of \(R\) is then \(\mathbb Q\).

We shall show that, for every \(i\),
\[
  i_*(\Fil^iR)\subset \Fil^{i+d}K^\bullet(X)_{\mathbb Q}.
\]
This will imply in particular
\[
  i_*(x)\in\Fil^{k+d}K^\bullet(X),
\]
and so prove the theorem.

The proof is by descending induction on \(i\), the case \(i=m\) being trivial. Let \(y\in\Fil^iR\), and set, for \(j\),
\[
  a_j=(-1)^{j-1}(j-1)!.
\]
Then, by \hyperlink{sga6.e05.prop.6-10}{V 6.10},
\[
  \gamma^{d+i}(N,y)-a_{d+i}y\in\Fil^{i+1}R.
\]
The induction hypothesis can therefore be applied, giving
\[
  i_*(\Fil^{i+1}R)\subset\Fil^{i+d+1}K^\bullet(X)_{\mathbb Q}.
\]
Thus
\[
  i_*\bigl(\gamma^{d+i}(N,y)-a_{d+i}y\bigr)
  \in\Fil^{i+d+1}K^\bullet(X)_{\mathbb Q}.
\]
By \hyperlink{sga6.e07.cor.4-4}{4.4} this implies
\[
  \gamma^{d+i}\bigl(i_*(y)\bigr)-a_{d+i}i_*(y)
  \in\Fil^{i+d+1}K^\bullet(X)_{\mathbb Q}.
\]
Since \(i_*(y)\) has zero augmentation,
\[
  \gamma^{d+i}\bigl(i_*(y)\bigr)\in\Fil^{i+d+1}K^\bullet(X).
\]
As we are working modulo torsion, this gives
\[
  i_*(y)\in\Fil^{i+d}K^\bullet(X)_{\mathbb Q}.
\]

\sgaStableTarget{sga6.e07.rem.4-6-1}\textbf{Remark 4.6.1.} One sees that, even if \hyperlink{sga6.e07.thm.4-3}{4.3} were true without tensoring by \(\mathbb Q\), the preceding proof works only modulo torsion. For a counterexample, see \hyperlink{sga6.e14.para.4-6}{XIV 4.6}.

\sgaStableTarget{sga6.e07.sec.4.7}\subsection*{4.7}

We return to the situation of \hyperlink{sga6.e07.thm.3-7}{3.7}, where one had the commutative diagram
\[
\begin{tikzcd}[column sep=large,row sep=large]
Y' \arrow[r,"j"] \arrow[d,"g"'] & X' \arrow[d,"f"]\\
Y \arrow[r,"i"'] & X,
\end{tikzcd}
\]
where \(i\) is a regular immersion, \(X'\) is the \(X\)-scheme obtained by blowing up \(Y\), and \(Y'=X'\times_XY\). Suppose \(X\) noetherian with an ample sheaf, and \(Y\) regular. We saw in \hyperlink{sga6.e07.thm.3-7}{3.7} that there is then the exact sequence
\[
0\longrightarrow K^\bullet(Y)\xrightarrow{u}
K^\bullet(Y')\times K^\bullet(X)\xrightarrow{v}
K^\bullet(X')\longrightarrow0.
\]
We shall show that this exact sequence permits one to describe the whole situation formed by the \(K^\bullet\) of \(X,Y,X',Y'\), from the data of the \(\lambda\)-rings \(K^\bullet(X)\), \(K^\bullet(Y)\), the maps \(i^*\) and \(i_*\), and the element \(N\in K^\bullet(Y)\), at least up to torsion, up to one further point.

Let \(K,K'\) be two \(\lambda\)-rings, let \(i_*:K\to K'\) be an additive homomorphism, let \(i^*:K'\to K\) be a \(\lambda\)-homomorphism, let \(d\) be an integer, and let \(N\in K\) be such that \(\lambda^d(N)\ne0\) for \(i>d\) and \(\lambda^d(N)\) is invertible. Suppose, moreover, that the following properties hold:
\begin{enumerate}[label=\roman*)]
\item \(i_*\) is \(K'\)-linear (the ``projection formula'');
\item \(i^*i_*\) is multiplication by \(\lambda_{-1}(N)\);
\item for all \(n\ge1\) and all \(y\in K\),
\[
 i_*\bigl(\lambda^n(N,y)\bigr)=\lambda^n(i_*(y))
\]
(relation \hyperlink{sga6.e07.eq.4-3-1}{(4.3.1)}).
\end{enumerate}

We first use \hyperlink{sga6.e06.thm.1-1}{VI 1.1} to construct a \(K\)-\(\lambda\)-algebra \(K_1\) which plays the role of \(K^\bullet(Y')\). Let \(L\) be an indeterminate, and let \(K[L]\) be the \(K\)-\(\lambda\)-algebra generated by \(L\), subject to the relations \(\lambda^i(L)=0\) for \(i\ge2\) (\hyperlink{sga6.e05.cor.4-9}{V 4.9}). Put
\[
  P(L)=\lambda^d(N-L)=\sum_{i=0}^{d}(-1)^i\lambda^{d-i}(N)L^i.
\]
The ideal generated by \(P\) is a \(\lambda\)-ideal. To see this, first observe that the relations \(\lambda^i(N)=0\) for \(i>d\) imply that the \(\lambda^j(N-L)\), for \(j\ge d\), are multiples of \(\lambda^d(N-L)\). If one works in \(1+K[(P)][[t]]\), the image \(\lambda_t(N-L)\) is a polynomial of degree \(<d\), and a universal argument shows that its \(\lambda^i\) for \(i\ge d\) are equal to \(1\) (cf. \hyperlink{sga6.e05.prop.4-8}{V 4.8}). Since in
\[
1+K[[t]]^+[L]/(L^N)
\]
one has \(\lambda_t(\lambda^d(N-L))=\lambda(\lambda_t(N-L))\), it follows that the \(\lambda^i(\lambda^d(N-L))\), \(i\ge1\), are multiples of \(P\). By the relations \hyperlink{sga6.e05.eq.2-4-2}{V 2.4.2}, the ideal generated by \(P\) is therefore a \(\lambda\)-ideal.

Put \(K_1=K/(P)\). The natural homomorphism \(K\to K_1\) will be denoted \(g^*\); it is a \(\lambda\)-homomorphism. Moreover \(K_1\) is a free \(K\)-module with basis
\[
  1,L,\ldots,L^{d-1}.
\]
Indeed, from
\[
  \lambda^d(N)=\lambda^{d-1}(N-L)L
\]
and the invertibility hypothesis on \(\lambda^d(N)\), it follows that \(L\) is invertible. Thus \(K_1\) admits as a basis over \(K\) the elements \(1,L^{-1},\ldots,L^{-d+1}\). Define an application \(g_*:K_1\to K\) by assigning to
\[
y'=y_0+y_1L+\cdots+y_{d-1}L^{d-1}
\]
the element \(y_{d-1}\) of \(K\). It is clear that the application so defined is an additive, \(K\)-linear homomorphism, and that \(g_*g^*=\Id\).

We next equip \(K_1\times K'\) with a \(\lambda\)-ring structure by putting
\sgaStableTarget{sga6.e07.eq.4-7-1}\sgaStableTarget{sga6.e07.eq.4-7-2}\begin{align}
 (y'_1,x_1)*(y'_2,x_2)
   &=(y'_1y'_2(1-L)+x_1y'_2+x_2y'_1,\; x_1x_2), \tag{4.7.1}\\
 \lambda^n_*(y',x)
   &=\sum_{p=1}^{n}\lambda^{\,n-p}(x)\lambda^p(L,y')
     +\lambda^n(x). \tag{4.7.2}
\end{align}
In these formulas \(K_1\) is regarded as a \(K'\)-algebra by the homomorphism
\[
  K'\xrightarrow{i^*}K\xrightarrow{g^*}K_1.
\]
Next, for \(y_1,y_2\in K\) and \(n\ge1\), put
\sgaStableTarget{sga6.e07.eq.4-7-3}\sgaStableTarget{sga6.e07.eq.4-7-4}\begin{align}
 y_1*_y y_2&=-\,y_1y_2\lambda_{-1}(N), \tag{4.7.3}\\
 \lambda^n_*(y_1)&=-\,\lambda^n(N,-y_1), \tag{4.7.4}
\end{align}
and define a homomorphism
\[
  u:K\longrightarrow K_1\times K'
\]
by
\[
  u(y)=\bigl(g^*(y)\lambda_{-1}(F),-i_*(y)\bigr),
  \qquad F=N-L.
\]

\sgaStableTarget{sga6.e07.prop.4-7-5}\textbf{Proposition 4.7.5.} The operations defined by \hyperlink{sga6.e07.eq.4-7-1}{(4.7.1)} and \hyperlink{sga6.e07.eq.4-7-2}{(4.7.2)} furnish \(K_1\times K'\) with a \(\lambda\)-ring structure, and \(u\) is an injective homomorphism which identifies \(K\), equipped with the operations defined by \hyperlink{sga6.e07.eq.4-7-3}{(4.7.3)} and \hyperlink{sga6.e07.eq.4-7-4}{(4.7.4)}, with a \(\lambda\)-ideal of \(K_1\times K'\).

It is immediate to verify that \(K_1\times K'\) is a \(\lambda\)-ring; for example, one may refer to \hyperlink{sga6.e05.para.5-5}{V 5.5}, where an entirely analogous construction is made. The verification of the second part of the proposition is also straightforward from the preceding relations, the relation
\[
  i_*(y_1y_2\lambda_{-1}(N))=i_*(y_1)\,i_*(y_2),
\]
which follows immediately from the hypotheses, the relation \(g_*(\lambda_{-1}(F))=1\) (\hyperlink{sga6.e06.prop.5-9}{VI 5.9}), and the relations of V 5.

Put
\[
  K'_1=(K_1\times K')/K,
\]
and let \(v:K_1\times K'\to K'_1\) be the canonical homomorphism. Regard \(K'_1\) as a \(\lambda\)-ring via the quotient structure. Define additive homomorphisms
\[
\begin{array}{llll}
j_*:K_1\to K'_1, & j^*:K'_1\to K_1, &
f_*:K'_1\to K', & f^*:K'\to K'_1 .
\end{array}
\]
The homomorphism \(j_*\) is the composite \(K_1\to K_1\times K'\to K'_1\). Similarly, \(f^*\) is the composite \(K'\to K_1\times K'\to K'_1\). Thus
\[
  v=j_*+f^*.
\]
Define next a homomorphism \(K_1\times K'\to K_1\) by assigning to \((y',x)\) the element
\[
  y'(1-L)+g^*i^*(x).
\]
This homomorphism vanishes on \(K\), and hence defines a homomorphism
\[
  j^*:K'_1\longrightarrow K_1.
\]
Finally, define a homomorphism \(K_1\times K'\to K'\) by assigning to \((y',x)\) the element
\[
  i_*g_*(y')+x.
\]
This homomorphism vanishes on \(K\), and hence defines a homomorphism
\[
  f_*:K'_1\longrightarrow K'.
\]

\sgaStableTarget{sga6.e07.prop.4-7-6}\textbf{Proposition 4.7.6.}
\begin{enumerate}[label=\roman*)]
\item \(j^*\) and \(f^*\) are \(\lambda\)-homomorphisms;
\item \(j_*\) is \(K'_1\)-linear and \(f_*\) is \(K'\)-linear;
\item \(j^*j_*\) is the homomorphism ``multiplication by \(1-L\)'';
\item for every \(y'\in K_1\) and every \(p\ge1\),
\[
  j_*\bigl(\lambda^p(L,y')\bigr)=\lambda^p(j_*(y'));
\]
\item \(f_*f^*=\Id\).
\end{enumerate}
It suffices to let the machine run.

If one returns to the geometric situation recalled above, one sees that \(K^\bullet(X')\) is constructed from \(K^\bullet(Y)\), \(K^\bullet(X)\), the maps \(i^*\) and \(i_*\), and the element \(N\in K^\bullet(Y)\), by the preceding method. However, the \(\lambda\)-operations on \(K^\bullet(X')\) agree with the \(\lambda\)-operations defined above on \(K'_1\) only up to torsion, so long as relation \hyperlink{sga6.e07.eq.4-3-1}{(4.3.1)} has been proved only modulo torsion.

\sgaStableTarget{sga6.e07.thm.4-8}\textbf{Theorem 4.8.} Let \(i:Y\to X\) be a regular closed immersion of constant codimension \(d\), let \(X'\) be the scheme obtained by blowing up \(Y\), and let \(Y'=X'\times_XY\), so that one has the commutative diagram
\[
\begin{tikzcd}[column sep=large,row sep=large]
Y' \arrow[r,"j"] \arrow[d,"g"'] & X' \arrow[d,"f"]\\
Y \arrow[r,"i"'] & X.
\end{tikzcd}
\]
Assume that \(X\) is quasi-compact and has an ample sheaf. Then, for every \(k\), there is an exact sequence
\[
0\longrightarrow
\operatorname{Gr}^{k-d}K^\bullet(Y)_{\mathbb Q}
\xrightarrow{u^{\gr}}
\operatorname{Gr}^{k-1}K^\bullet(Y')_{\mathbb Q}
\times
\operatorname{Gr}^{k}K^\bullet(X)_{\mathbb Q}
\xrightarrow{v^{\gr}}
\operatorname{Gr}^{k}K^\bullet(X')_{\mathbb Q},
\]
where, for \(y\in \operatorname{Gr}^{k-d}K^\bullet(Y)_{\mathbb Q}\),
\[
  u^{\gr}(y)=
  \bigl(g^{\gr}(y)(-1)^{d-1}c^{d-1}(F),-\,i_{\gr}(y)\bigr),
\]
and, for \(y'\in \operatorname{Gr}^{k-1}K^\bullet(Y')_{\mathbb Q}\) and \(x\in \operatorname{Gr}^{k}K^\bullet(X)_{\mathbb Q}\),
\[
  v^{\gr}(y',x)=j_{\gr}(y')+f^{\gr}(x).
\]
The notation is that of \hyperlink{sga6.e07.para.3-1}{3.1} and \hyperlink{sga6.e07.para.3-3}{3.3}, and \(c^{d-1}(F)\) denotes the \((d-1)\)-st Chern class of \(F\) (\hyperlink{sga6.e05.para.6-7}{V 6.7}). Moreover, \(u^{\gr}\) has a left inverse \(u'^{\gr}\), defined by
\[
  u'^{\gr}(y',x)=g_{\gr}(y').
\]
If, in addition, \(X\) is noetherian and \(Y\) regular, one has the exact sequence
\[
0\longrightarrow
\operatorname{Gr}^{k-d}K^\bullet(Y)_{\mathbb Q}
\xrightarrow{u^{\gr}}
\operatorname{Gr}^{k-1}K^\bullet(Y')_{\mathbb Q}
\times
\operatorname{Gr}^{k}K^\bullet(X)_{\mathbb Q}
\xrightarrow{v^{\gr}}
\operatorname{Gr}^{k}K^\bullet(X')_{\mathbb Q}
\longrightarrow0.
\]

\emph{Proof.} Notice that one could give, in the case where \(Y\) is regular, a purely algebraic proof in terms of the theory of \(\lambda\)-rings, using the description \hyperlink{sga6.e07.sec.4.7}{4.7} of the situation. We shall instead compute it from the proof of 3.7.

We know that the inverse-image homomorphisms \(f^*,g^*\) are \(\lambda\)-homomorphisms, are compatible with the filtrations of the \(\lambda\)-rings, and therefore define homomorphisms on the associated gradeds:
\[
  f^{\gr}:\operatorname{Gr}^kK^\bullet(X)\to\operatorname{Gr}^kK^\bullet(X'),
  \qquad
  g^{\gr}:\operatorname{Gr}^{k-d}K^\bullet(Y)\to\operatorname{Gr}^{k-d}K^\bullet(Y'),
\]
and similarly after tensoring with \(\mathbb Q\). Moreover, \hyperlink{sga6.e07.thm.4-6}{4.6} shows that, after tensoring with \(\mathbb Q\), the homomorphisms \(i_*,j_*\) pass to the associated gradeds, giving
\[
  i_{\gr}:\operatorname{Gr}^{k-d}K^\bullet(Y)_{\mathbb Q}
  \to \operatorname{Gr}^{k}K^\bullet(X)_{\mathbb Q},
  \qquad
  j_{\gr}:\operatorname{Gr}^{k-1}K^\bullet(Y')_{\mathbb Q}
  \to \operatorname{Gr}^{k}K^\bullet(X')_{\mathbb Q}.
\]
Finally, since
\[
  \lambda_{-1}(F)=(-1)^{d-1}\gamma^{d-1}(F),
\]
the morphisms \(u^{\gr}\) and \(v^{\gr}\) are obtained by passage to the quotient from the morphisms \(u\) and \(v\) of 3.7. Thus \(v^{\gr}u^{\gr}=0\).

By \hyperlink{sga6.e06.cor.5-8}{VI 5.8},
\[
  g_*(\Fil^{k-1}K^\bullet(Y'))\subset \Fil^{k-d}K^\bullet(Y)
\]
for every \(k\), since \(Y'=\mathbf P(N)\) and \(N\) has rank \(d\). We therefore obtain a homomorphism
\[
  g_{\gr}:\operatorname{Gr}^{k-1}K^\bullet(Y')
  \to\operatorname{Gr}^{k-d}K^\bullet(Y),
\]
and an analogous homomorphism after tensoring with \(\mathbb Q\). The homomorphism \(u'^{\gr}\) is obtained from \(u'\) \hyperlink{sga6.e07.thm.3-7}{(3.7)} by passage to the associated graded. Since \(u'\) is a left inverse of \(u\), \(u'^{\gr}\) is a left inverse of \(u^{\gr}\), proving at the same time the injectivity of \(u^{\gr}\).

To show that \(\Ker(v^{\gr})=\Im(u^{\gr})\), suppose that
\[
y'\in\Fil^{k-1}K^\bullet(Y')_{\mathbb Q},\qquad
x\in\Fil^kK^\bullet(X)_{\mathbb Q}
\]
are such that
\sgaStableTarget{sga6.e07.eq.4-6-1}\begin{equation}\tag{4.6.1}
  j_*(y')+f^*(x)\in\Fil^{k+1}K^\bullet(X')_{\mathbb Q}.
\end{equation}
We must show that there exists \(y\in\Fil^{k-d}K^\bullet(Y)_{\mathbb Q}\) such that
\[
\begin{aligned}
\text{i)}\quad &x+i_*(y)\in\Fil^{k+1}K^\bullet(X)_{\mathbb Q},\\
\text{ii)}\quad &y'-g^*(y)\lambda_{-1}(F)\in\Fil^{k}K^\bullet(Y')_{\mathbb Q}.
\end{aligned}
\]
Put
\[
  y=g_*(y').
\]
By \hyperlink{sga6.e06.cor.5-8}{VI 5.8}, \(y\in\Fil^{k-d}K^\bullet(Y)\). Since \(f_*\) preserves the filtration (using results of \hyperlink{sga6.e08.expose}{Exposé VIII}, which is inconvenient here, for \hyperlink{sga6.e07.thm.4-8}{4.8} is not used to prove this fact), it follows from \hyperlink{sga6.e07.prop.1-8}{1.8} ii) that \(f\) is a complete-intersection morphism of virtual relative dimension \(0\) (\hyperlink{sga6.e08.prop.4-2}{VIII 4.2}). Since \(f\) is projective, the result follows from \hyperlink{sga6.e08.prop.3-2}{VIII 3.2}. Applying \(f_*\) to \hyperlink{sga6.e07.eq.4-6-1}{(4.6.1)}, we get
\[
  f_*(j_*(y'))+f_*f^*(x)\in\Fil^{k+1}K^\bullet(X)_{\mathbb Q}.
\]
Thus, by \hyperlink{sga6.e07.prop.3-6}{3.6},
\[
  i_*(y)+x\in\Fil^{k+1}K^\bullet(X)_{\mathbb Q}.
\]

To prove (ii), one must show that
\[
  y'-g^*(g_*y')\lambda_{-1}(F)\in\Fil^kK^\bullet(Y')_{\mathbb Q}.
\]
Put
\[
  z=y'-g^*(g_*y')\lambda_{-1}(F).
\]
Then
\[
\begin{aligned}
j_*(z)
  &=j_*(y')-j_*\bigl(g^*(g_*y')\lambda_{-1}(F)\bigr)\\
  &=j_*(y')-f^*(i_*g_*y')\qquad\text{by 3.4.}
\end{aligned}
\]
But relation (i) gives
\[
  f^*(i_*g_*y')+f^*(x)\in\Fil^{k+1}K^\bullet(X')_{\mathbb Q}.
\]
Together with \hyperlink{sga6.e07.eq.4-6-1}{(4.6.1)}, this gives
\[
  j_*(z)\in\Fil^{k+1}K^\bullet(X')_{\mathbb Q}.
\]
Consequently
\[
  j^*(j_*(z))=z(1-L)\in\Fil^{k+1}K^\bullet(Y')_{\mathbb Q}.
\]
Let \(z'\) be the image of \(z\) in
\(\operatorname{Gr}^{k-1}K^\bullet(Y')_{\mathbb Q}\), and let \((1-L)\) denote the class of \(1-L\) in \(\operatorname{Gr}^{1}K^\bullet(Y')_{\mathbb Q}\). We know that \(z'(1-L)=0\). We may then apply \hyperlink{sga6.e06.prop.5-11}{VI 5.11} (which remains valid after tensoring with \(\mathbb Q\), since its proof uses only the structure of \(K^\bullet(Y')\) as a free \(K^\bullet(Y)\)-algebra) and find \(z''\in K^\bullet(Y)_{\mathbb Q}\) such that
\[
  z=g^*(z'')\lambda_{-1}(F)\in\Fil^kK^\bullet(Y')_{\mathbb Q}.
\]
Applying \(g_*\), one obtains, since \(g_*(z)=0\),
\[
  z''\in\Fil^{k+1-d}K^\bullet(Y)_{\mathbb Q},
\]
and hence, since \(\lambda_{-1}(F)\in\Fil^{d-1}K^\bullet(Y')_{\mathbb Q}\),
\[
  z\in\Fil^kK^\bullet(Y')_{\mathbb Q},
\]
which finishes the proof.

Assume now that \(X\) is noetherian and \(Y\) regular. Then \(v\) is surjective \hyperlink{sga6.e07.thm.3-7}{(3.7)}. We saw in the proof of \hyperlink{sga6.e07.thm.3-7}{3.7} that for every \(x'\in K^\bullet(X')\) there exists \(y'\in K^\bullet(Y')\) such that
\[
  x'=f^*(f_*x')+j_*(y').
\]
The same property remains true after tensoring with \(\mathbb Q\). Hence, if \(x'\in\Fil^kK^\bullet(X')_{\mathbb Q}\), then there exist \((y',x)\in K^\bullet(Y')_{\mathbb Q}\times K^\bullet(X)_{\mathbb Q}\) such that
\[
  x'=v(y',x),\qquad x\in\Fil^kK^\bullet(X)_{\mathbb Q}.
\]
Indeed it suffices to take \(x=f_*(x')\).

Filter \(K^\bullet(Y')_{\mathbb Q}\times K^\bullet(X)_{\mathbb Q}\), as in \hyperlink{sga6.e07.sec.4.7}{4.7}, by putting
\[
  \Fil^k\bigl(K^\bullet(Y')_{\mathbb Q}\times K^\bullet(X)_{\mathbb Q}\bigr)
  =
  \Fil^{k-1}K^\bullet(Y')_{\mathbb Q}\times \Fil^kK^\bullet(X)_{\mathbb Q}.
\]
It follows immediately from \hyperlink{sga6.e07.eq.4-7-1}{4.7.1} that this is a ring filtration. To show that for every \(k\) the homomorphism
\[
  \Fil^1\bigl(K^\bullet(Y')_{\mathbb Q}\times K^\bullet(X)_{\mathbb Q}\bigr)
  \longrightarrow
  \Fil^1K^\bullet(X')_{\mathbb Q}
\]
is surjective, it is enough to show that, if \(x'=y(x'_1)\) with \(\epsilon(x'_1)=0\), there is \((y',x)\in\Fil^1(K^\bullet(Y')_{\mathbb Q}\times K^\bullet(X)_{\mathbb Q})\) such that \(x'=v(y',x)\).

By what precedes, there exists \((y'_1,x_1)\in K^\bullet(Y')_{\mathbb Q}\times K^\bullet(X)_{\mathbb Q}\) such that \(x'_1=v(y'_1,x_1)\) and \(x_1\in\Fil^1K^\bullet(X)_{\mathbb Q}\). By \hyperlink{sga6.e07.sec.4.7}{4.7},
\[
  x'=v\bigl(\gamma^k_*(y'_1,x_1)\bigr).
\]
But
\[
  \gamma^k_*(y'_1,x_1)
  =
  \sum_{p=1}^{k}\gamma^{k-p}(x_1)\gamma^p(L,y'_1),
\]
which is an immediate consequence of \hyperlink{sga6.e07.eq.4-7-2}{(4.7.2)}. Since \(x_1\) has zero augmentation, \(\gamma^k(x_1)\in\Fil^1K^\bullet(X)_{\mathbb Q}\), and by \hyperlink{sga6.e05.prop.5-7}{V 5.7}
\[
  \sum_{p=1}^{k}\gamma^{k-p}(x_1)\gamma^p(L,y'_1)
  \in\Fil^{k-1}K^\bullet(Y')_{\mathbb Q}.
\]
This proves the assertion.

\sgaStableTarget{sga6.e07.rem.4-9}\textbf{Remark 4.9.} In the case where \(X\) is noetherian and \(Y\) regular, one can, proceeding as in \hyperlink{sga6.e07.sec.4.7}{4.7}, describe the structure of graded ring of \(\operatorname{Gr}^*K^\bullet(X')_{\mathbb Q}\) in terms of the graded rings \(\operatorname{Gr}^*K^\bullet(X)_{\mathbb Q}\) and \(\operatorname{Gr}^*K^\bullet(Y)_{\mathbb Q}\), of the additive homomorphisms \(i_{\gr}\) and \(i^{\gr}\) induced by \(i^*\) and \(i_*\), and of the element \(c(N)\in\operatorname{Gr}^*K^\bullet(Y)\). Similarly, one can describe, from the same data, the homomorphisms \(f^{\gr}, f_{\gr}, j^{\gr}, j_{\gr}\) defined from \(f^*, f_*, j^*, j_*\) by passage to the associated gradeds.

\sgaStableTarget{sga6.e07.sec.4.10}\subsection*{4.10}

% SGA-SRC-001500
Let \(X\) be a regular scheme having an ample sheaf. We know (\hyperlink{sga6.e04.para.2-12}{IV 2.12}) that the Grothendieck groups formed with locally free finite-type \(\mathcal O_X\)-modules, and with coherent \(\mathcal O_X\)-modules, are isomorphic. We shall use this fact to define on \(K^\bullet(X)\) a ``topological'' filtration: take \(\Fil_{\mathrm{top}}^iK^\bullet(X)\) to be the subgroup generated by the classes of coherent sheaves whose support has codimension at least \(i\). One shows easily that \(\Fil_{\mathrm{top}}^iK^\bullet(X)\) is generated by the classes of the sheaves \(\mathcal O_Z\), where \(Z\) ranges over the closed integral subschemes of \(X\) of codimension \(\ge i\).

\sgaStableTarget{sga6.e07.prop.4-11}\textbf{Proposition 4.11.} Let \(X\) be a regular scheme having an ample sheaf. Then, for every \(k\ge0\),
\[
  \Fil_{\mathrm{top}}^kK^\bullet(X)_{\mathbb Q}
  \subset
  \Fil^kK^\bullet(X)_{\mathbb Q}.
\]
The reverse inclusion, valid without tensoring by \(\mathbb Q\), will be proved in X.

By \hyperlink{sga6.e07.sec.4.10}{4.10}, it is enough to show that, if \(Z\) is an integral closed subscheme of \(X\), of codimension \(k\), then
\[
  \operatorname{cl}(\mathcal O_Z)\in\Fil^kK^\bullet(X)_{\mathbb Q}.
\]
To this end we proceed by noetherian induction on the integral closed subschemes of \(X\). Suppose the property already proved for the proper integral closed subschemes of \(Z\), and let us prove it for \(Z\). Let \(J\) be the ideal of \(\mathcal O_X\) defining \(Z\), and let \(z\) be the generic point of \(Z\). Since \(Z\) is integral, \(J_z\) is the maximal ideal of \(\mathcal O_{X,z}\); and since \(X\) is regular, \(J_z\) can be generated by an \(\mathcal O_{X,z}\)-regular sequence. One can therefore find a sequence of sections of \(J\), in a neighbourhood of \(z\), forming an \(\mathcal O_X\)-regular sequence. Consequently there is an open neighbourhood \(U\) of \(z\), such that \(Z'=X-U\) is contained in \(Z\), and such that \(J\mathcal O_U\) is a regular ideal of \(\mathcal O_U\). The immersion \(Z\cap U\to U\) is then a regular immersion of codimension
\[
  \dim(\mathcal O_{X,z})=\codim(Z,X)
\]
(EGA IV 5.1.2). Moreover, one has the exact sequence (IV)
\[
  K^\bullet(Z')\xrightarrow{j_*}K^\bullet(X)
  \xrightarrow{u^*}K^\bullet(U)\to0,
\]
where \(j\) is the immersion \(Z'\to X\) and \(u\) is the immersion \(U\to X\). After tensoring with \(\mathbb Q\), one obtains
\[
  \operatorname{cl}(\mathcal O_Z)
  \equiv i_*(1_{\mathcal O_{Z\cap U}}),
\]
and by \hyperlink{sga6.e07.thm.4-6}{4.6} the image under \(u^*\) of the class of \(\mathcal O_Z\) belongs to \(\Fil^kK^\bullet(U)_{\mathbb Q}\). Since \(u^*\) is a surjective \(\lambda\)-homomorphism, it induces a surjection
\[
  \Fil^kK^\bullet(X)\longrightarrow \Fil^kK^\bullet(U).
\]
Thus \(u^*(1_Z)=u^*(x)\), with \(x\in\Fil^kK^\bullet(X)_{\mathbb Q}\). Therefore there exists
\[
  y\in j_*(K^\bullet(Z'))
\]
such that
\[
  \operatorname{cl}(\mathcal O_Z)-y\in\Fil^kK^\bullet(X)_{\mathbb Q}.
\]
But \(Z'\subset Z\), and since \(z\notin Z'\), \(\codim(Z',X)>k\). Thus \(y\) can be expressed as a combination of classes of sheaves \(\mathcal O_{Z''_\alpha}\), where the \(Z''_\alpha\) are integral closed subschemes of \(Z\), different from \(Z\). By the induction hypothesis, \(y\in\Fil^kK^\bullet(X)_{\mathbb Q}\), and this completes the proof.

\sgaStableTarget{sga6.e07.bibliography}\subsection*{Bibliography}
\begin{enumerate}[label={[\arabic*]}]
\item Cartan--Eilenberg, \emph{Homological Algebra}.
\end{enumerate}



\providecommand{\etale}{\'etale}
\providecommand{\Ch}{\operatorname{Ch}}


\sgaStableTarget{sga6.e08.expose}\section*{Expos\'e VIII. The Riemann--Roch theorem}
\begin{center}
by Pierre Berthelot
\end{center}

\sgaStableTarget{sga6.e08.sec.1}\subsection*{1. Complete-intersection morphisms}

\sgaStableTarget{sga6.e08.def.1-1}\textbf{Definition 1.1.} Let $f:X\to Y$ be a morphism of schemes. One says that $f$ is a \emph{complete-intersection morphism} if, for every point $x$ of $X$, there exist an open neighbourhood $U$ of $x$ and a smooth $Y$-scheme $V$ such that the restriction of $f$ to $U$ factors as
\[
\begin{tikzcd}[column sep=tiny]
U \arrow[r,"i"] \arrow[dr,"f"'] & V \arrow[d,"g"]\\
& Y,
\end{tikzcd}
\]
where $i$ is a regular immersion (\hyperlink{sga6.e07.def.1-4}{VII 1.4}). After replacing $V$ by an open subset, one may suppose that $i$ is a regular closed immersion.

In this definition one may replace $V$ by a vector bundle $Y[T_1,\ldots,T_n]$. Indeed, since the notion is local, one may suppose that $V$ is affine and that its image in $Y$ is contained in an affine open $W$ of $Y$. There then exists an integer $n$ and a closed $W$-immersion
\[
V\longrightarrow W[T_1,\ldots,T_n],
\]
hence a $Y$-immersion $V\to Y[T_1,\ldots,T_n]$, which is regular by \hyperlink{sga6.e07.prop.1-10}{VII 1.10}. Since the composite of two regular immersions is regular, one obtains a regular $Y$-immersion
\[
U\longrightarrow Y[T_1,\ldots,T_n]
\]
factorizing $f$.

\sgaStableTarget{sga6.e08.prop.1-2}\textbf{Proposition 1.2.} Let $f:X\to Y$ be a complete-intersection morphism. Then every $Y$-immersion $i:X\to Z$ of $X$ in a smooth $Y$-scheme is regular.

We first prove the following corollary.

\sgaStableTarget{sga6.e08.cor.1-3}\textbf{Corollary 1.3.} Let
\[
\begin{tikzcd}[column sep=large,row sep=large]
& V' \arrow[d,"p"]\\
X \arrow[ur,"j"] \arrow[r,"i"'] & V
\end{tikzcd}
\]
be a commutative diagram in which $i$ and $j$ are immersions, and $p$ is a smooth morphism. Then $i$ is a regular immersion if and only if $j$ is a regular immersion.

(i) Suppose $i$ regular. Put $X'=V'\times_V X$. The datum of $j$ defines a section $j'$ of $X'$ over $X$, such that $j=i'j'$:
\[
\begin{tikzcd}[column sep=large,row sep=large]
X' \arrow[r,"i'"] \arrow[d,"p'"'] & V' \arrow[d,"p"]\\
X \arrow[r,"i"'] \arrow[ur,"j"] & V .
\end{tikzcd}
\]
Now $j'$ is a regular immersion, being a section of a smooth morphism (\hyperlink{sga6.e07.prop.1-10}{VII 1.10}); $i'$ is a regular immersion because it is obtained from $i$ by the smooth, hence flat, base change (\hyperlink{sga6.e07.prop.1-5}{VII 1.5}). Therefore $j$ is a regular immersion (\hyperlink{sga6.e07.prop.1-7}{VII 1.7}).

(ii) Suppose $j$ regular and $p$ \etale. Then $j'$ is an open immersion (EGA IV 17.9.3), hence an isomorphism of $X$ onto $j'(X)$; put $U=j'(X)$. Let $i''$ be the restriction of $i'$ to $U$, and let $U'$ be an open subset of $V'$ such that $U=U'\cap X'$. Then
\[
i''=j\circ(j')^{-1},
\]
so that $i''$ is a regular immersion of $U$ into $U'$. But the morphism $U'\to V$, being smooth, is flat and locally of finite presentation, hence open. Consequently it is faithfully flat and quasi-compact from $U'$ onto an open neighbourhood of $X$ in $V$. It follows from \hyperlink{sga6.e07.prop.1-5}{VII 1.5} that $i$ is a regular immersion.

(iii) Suppose $j$ regular, and suppose that $V'$ is a vector bundle over $V$. Since the property to be proved is local on $V$, one may suppose $V$ affine, $i$ closed, and $V'=V[T_1,\ldots,T_n]$. Put $V=\Spec(A)$, $X=\Spec(B)$. One has the commutative diagram
\[
\begin{tikzcd}[column sep=large,row sep=large]
& A[T_1,\ldots,T_n] \arrow[d,"v"]\\
A \arrow[r] \arrow[ur,"u"] & B,
\end{tikzcd}
\]
where the morphisms $u$ and $v$ are surjective. Choose elements $t_i\in A$, $1\leq i\leq n$, such that $u(t_i)=v(T_i)$ for every $i$. Define a homomorphism
\[
w:A[T_1,\ldots,T_n]\longrightarrow A,
\qquad w(T_i)=t_i.
\]
This gives a closed immersion $s:V\to V'$ which is a section of $p$, and such that $j=s\circ i$. Moreover, being a section of a smooth morphism, $s$ is a regular immersion (\hyperlink{sga6.e07.prop.1-10}{VII 1.10}). Since $s$ and $s\circ i$ are regular, it will follow that $i$ is regular if one proves that the sequence
\[
0\longrightarrow i^*(K/K^2)\longrightarrow J/J^2\longrightarrow I/I^2\longrightarrow 0
\]
is exact and locally split; here $I,J,K$ denote respectively the ideals defining $i,j,s$ (\hyperlink{sga6.e07.prop.1-7}{VII 1.7}). It is enough to show that $i^*(K/K^2)$ is a direct factor of $J/J^2$. But $j'$ is obtained from $s$ by the base change $i':X'\to V'$. Moreover $j'$ and $s$ are regular immersions (\hyperlink{sga6.e07.prop.1-10}{VII 1.10}). Thus, by \hyperlink{sga6.e07.lem.1-2}{VII 1.2}, the ideal defining $j'$ is $i'^*(K)$, and the conormal sheaf of $j'$ is $i'^*(K/K^2)$. The exact sequence of conormal sheaves attached to the decomposition $j=i'j'$ therefore gives in particular a homomorphism
\[
J/J^2\longrightarrow i^*(K/K^2).
\]
It is immediate to verify that the composite
\[
i^*(K/K^2)\longrightarrow J/J^2\longrightarrow i^*(K/K^2)
\]
is the identity. In affine terms, this corresponds to the decomposition of $J$ as the direct sum of an ideal of $A$ and the ideal generated by the $T_i-t_i$. This gives the required splitting.

(iv) Finally suppose $j$ regular and $p$ smooth. For every $x\in X$, one can find an open neighbourhood $W$ of $i(x)$, an open neighbourhood $W'$ of $j(x)$ lying over $W$, and a commutative diagram
\[
\begin{tikzcd}[column sep=large,row sep=large]
W' \arrow[r,"q"] \arrow[d,"p"'] & W[T_1,\ldots,T_n] \arrow[d,"p_0"]\\
W \arrow[r,equal] & W,
\end{tikzcd}
\]
where $p_0$ is the structural morphism, and $q$ is an \etale morphism (EGA IV 17.11.4). One may then suppose $V=W$ and $V'=W'$. It is easy to see that $qj$ is an immersion, by introducing the fiber product $X[T_1,\ldots,T_n]$. Applying (ii) and (iii), one deduces that $qj$ is regular, and then that $i$ is regular.

We now prove \hyperlink{sga6.e08.prop.1-2}{Proposition 1.2}. Since regularity is local on $Z$, we may work in a neighbourhood of a point $x\in X$. There is then an open neighbourhood $U$ of $x$ in $X$, a smooth $Y$-scheme $V$, and a regular $Y$-immersion $i':U\to V$. We have a commutative diagram
\[
\begin{tikzcd}[column sep=large,row sep=large]
U \arrow[r,"i'"] \arrow[d,"i"'] & V \arrow[d]\\
Z \arrow[r] & Y,
\end{tikzcd}
\]
where the unlabeled vertical arrows are the structural morphisms. From it one deduces a closed immersion
\[
i'':U\longrightarrow Z\times_Y V,
\]
and a commutative diagram
\[
\begin{tikzcd}[column sep=large,row sep=large]
& U \arrow[dl,"i"'] \arrow[dr,"i'"] &\\
Z & Z\times_Y V \arrow[l,"p"'] \arrow[r,"p'"] & V.
\end{tikzcd}
\]
Since $i'$ is regular and $p,p'$ are smooth, the corollary implies that $i''$ and $i$ are regular.

The following proposition will not be used in the sequel of the Seminar.

\sgaStableTarget{sga6.e08.prop.1-4}\textbf{Proposition 1.4.} Let $f:X\to Y$ be a flat morphism locally of finite presentation. For $f$ to be a complete-intersection morphism in the sense of \hyperlink{sga6.e08.def.1-1}{1.1}, it is necessary and sufficient that it be a complete-intersection morphism in the sense of EGA IV 19.3.6.

Suppose that $f$ is a complete-intersection morphism in the sense of \hyperlink{sga6.e08.def.1-1}{1.1}, and let $x\in X$. Then there are an open neighbourhood $U$ of $x$ and a factorization
\[
\begin{tikzcd}[column sep=large,row sep=large]
U \arrow[r,"i"] \arrow[dr,"f"'] & V \arrow[d,"g"]\\
& Y,
\end{tikzcd}
\]
where $i$ is a regular immersion and $g$ is smooth. Let $y=f(x)$. Then the immersion $i_y:U_y\to V_y$ is regular at $x$. Indeed, after possibly restricting $V$, one may suppose that there is a finite-type free $\mathcal O_V$-module $E$ and a regular homomorphism $u:E\to\mathcal O_V$ such that the Koszul complex associated to $u$ is a resolution of $\mathcal O_U$. Let $u_y$ be the homomorphism induced on the fiber over $y$, and let $K_\bullet(u_y)=K_\bullet(u)_y$ be the associated Koszul complex. Since $g$ is smooth, hence flat, $K_\bullet(u)$ is a $Y$-flat resolution of $\mathcal O_U$. Consequently the homology groups of $K_\bullet(u)_y$ are the groups
\[
\Tor_i^{\mathcal O_Y}(\mathcal O_U,k(y)),
\]
where $k(y)$ is the residue field of $\mathcal O_{Y,y}$. Since $f$ is flat, these groups are zero; hence $i_y$ is regular at $x$. Since $g$ is smooth, the local ring of $V_y$ at $x$ is regular. The local ring of $U_y$ at $y$ is therefore the quotient of a regular ring by a regular ideal, hence a complete-intersection ring (EGA IV 19.3.2). It follows that $f$ is a complete intersection at $x$ in the sense of EGA IV 19.3.6.

Conversely, suppose that $f$ is a complete-intersection morphism in the sense of EGA IV 19.3.6, and let $x\in X$. Let $U$ be an affine neighbourhood of $x$ lying over an affine neighbourhood $U'$ of $y=f(x)$. One can then find a smooth $U'$-scheme $V$, and a closed immersion $U\to V$ (for instance by taking for $V$ a vector bundle over $U'$). We shall show that $i$ is regular. By EGA IV 19.3.7, $i$ is transversally regular at $x$, since $V$ is smooth over $Y$, hence $V_y$ is regular at $x$, and $U$ is flat over $Y$. It follows from EGA IV 19.2.4, a)3c), that $i$ is regular at $x$. Therefore $f$ is a complete-intersection morphism in the sense of 1.1.

\sgaStableTarget{sga6.e08.prop.1-5}\textbf{Proposition 1.5.} Let $f:X\to Y$ and $g:Y\to Z$ be two locally complete-intersection morphisms. Then $gf$ is locally complete-intersection.

Let $x\in X$. One can find open neighbourhoods $U,V$ of $x,f(x)$ and regular immersions $i,j$ such that one has factorizations
\[
\begin{tikzcd}[column sep=large,row sep=large]
U \arrow[r,"i"] \arrow[dr,"f"'] & V[T_1,\ldots,T_n] \arrow[d]\\
& V
\end{tikzcd}
\qquad
\begin{tikzcd}[column sep=large,row sep=large]
V \arrow[r,"j"] \arrow[dr,"g"'] & Z[T'_1,\ldots,T'_m] \arrow[d]\\
& Z.
\end{tikzcd}
\]
This diagram can be completed by the commutative square
\[
\begin{tikzcd}[column sep=large,row sep=large]
V[T_1,\ldots,T_n] \arrow[r,"j'"] \arrow[d] & Z[T'_1,\ldots,T'_m,T_1,\ldots,T_n] \arrow[d]\\
V \arrow[r,"j"'] & Z[T'_1,\ldots,T'_m],
\end{tikzcd}
\]
where $j'$ is obtained from $j$ by base change. This base change is flat; hence $j'$ is a regular immersion, and so $j'i$ is a regular immersion of $U$ in a scheme smooth over $W$, where $W$ is the relevant open of $Z$. This proves the assertion.

% SGA-SRC-001509
\sgaStableTarget{sga6.e08.prop.1-6}\textbf{Proposition 1.6.} Let $f:X\to Y$ be a morphism and let $g:Y'\to Y$ be a flat base change. If $f$ is complete-intersection, then the morphism
\[
f':X\times_Y Y'\longrightarrow Y'
\]
deduced from $f$ is complete-intersection. The converse is true if, moreover, $g$ is quasi-compact and surjective.

One can find an open covering of $X$ by open subsets $U_i$ such that, for every $i$, the restriction of $f$ to $U_i$ factors as
\[
U_i\longrightarrow V_i\longrightarrow Y,
\]
where $U_i\to V_i$ is a regular immersion and $V_i\to Y$ is smooth. The $U_i\times_Y Y'$ form an open covering of $X\times_Y Y'$, the immersions
\[
U_i\times_Y Y'\longrightarrow V_i\times_Y Y'
\]
are regular by \hyperlink{sga6.e07.prop.1-5}{VII 1.5}, and the morphisms $V_i\times_Y Y'\to Y'$ are smooth, giving the first assertion. For the converse, EGA IV 2.7.1 (iv) first implies that $f$ is locally of finite presentation; this reduces us easily, using \hyperlink{sga6.e08.prop.1-2}{1.2}, to the case where $f$ is a closed immersion. In that case the assertion is just \hyperlink{sga6.e07.prop.1-5}{VII 1.5}.

\sgaStableTarget{sga6.e08.prop.1-7}\textbf{Proposition 1.7.} Let $f:X\to Y$ be a complete-intersection morphism. Then $f$ is a perfect morphism.

Since the property is local on $X$, one may work in a neighbourhood of a point $x\in X$ and suppose that there is an open subset $U$ containing $x$ and a factorization
\[
\begin{tikzcd}[column sep=large,row sep=large]
U \arrow[r,"i"] \arrow[dr,"f"'] & V \arrow[d,"g"]\\
& Y,
\end{tikzcd}
\]
where $i$ is a regular immersion and $g$ is smooth. Since $g$ is smooth, it is a perfect morphism (\hyperlink{sga6.e03.def.4-1}{III 4.1}). By \hyperlink{sga6.e07.prop.1-9}{VII 1.9}, $i$ is also a perfect morphism. Since the composite of two perfect morphisms is perfect (\hyperlink{sga6.e03.prop.4-5}{III 4.5}), the proposition follows.

\sgaStableTarget{sga6.e08.prop.1-8}\textbf{Proposition 1.8.} Let $f:X\to Y$ be a complete-intersection morphism, let $x$ be a point of $X$, and let
\[
\begin{tikzcd}[column sep=large,row sep=large]
U \arrow[r,"i"] \arrow[dr,"f"'] & V \arrow[d,"g"]\\
& Y
\end{tikzcd}
\]
be a factorization in which $U$ is an open neighbourhood of $x$, $i$ is a regular immersion, and $g$ is smooth. Let $J$ be the ideal of $\mathcal O_V$ defining $U$. Then the integer
\[
d_x=\operatorname{rank}\bigl((\Omega^1_{V/Y})_{i(x)}\bigr)-\operatorname{rank}\bigl((J/J^2)_x\bigr)
\]
is independent of the chosen factorization of $f$.

Suppose that the restriction of $f$ to $U$ factors through two smooth $Y$-schemes $V$ and $V'$, and put $V''=V\times_Y V'$. Then $V''$ is smooth over $Y$, and also over $V$ and $V'$. Moreover, the immersions $X\to V$ and $X\to V'$ define an immersion $X\to V''$, which is regular by 1.2. It is therefore enough to compare the integers $d_x$ attached to two factorizations through smooth $Y$-schemes $V$ and $V'$ such that there exists a commutative diagram over $Y$
\[
\begin{tikzcd}[column sep=large,row sep=large]
& V' \arrow[d,"p"]\\
X \arrow[ur,"j"] \arrow[r,"i"'] & V,
\end{tikzcd}
\]
where $p$ is smooth. We may also suppose that $\Omega^1_{V/Y}$, $\Omega^1_{V'/Y}$, and $\Omega^1_{V'/V}$ have constant rank, as do the conormal sheaves of $i$ and $j$.

From the commutative diagram, all morphisms being smooth,
\[
\begin{tikzcd}[column sep=large,row sep=large]
V' \arrow[r,"p"] \arrow[dr] & V \arrow[d]\\
& Y,
\end{tikzcd}
\]
one deduces
\sgaStableTarget{sga6.e08.eq.1-8-1}\begin{equation}\tag{1.8.1}
\operatorname{rank}(\Omega^1_{V'/Y})=
\operatorname{rank}(\Omega^1_{V/Y})+
\operatorname{rank}(\Omega^1_{V'/V}).
\end{equation}
Let $X'=V'\times_V X$; then the rank of $\Omega^1_{V'/V}$ is equal to that of $\Omega^1_{X'/X}$. Further, $j$ defines a section $j':X\to X'$, which is a regular immersion of codimension equal to the rank of $\Omega^1_{X'/X}$ (EGA IV 17.10.4). If $i'$ is the immersion obtained from $i$ by the base change $V'\to V$, then $j=i'j'$, and consequently (\hyperlink{sga6.e07.prop.1-7}{VII 1.7})
\sgaStableTarget{sga6.e08.eq.1-8-2}\begin{equation}\tag{1.8.2}
\operatorname{rank}(N_{V'/X})=
\operatorname{rank}(N_{V/X})+
\operatorname{rank}(\Omega^1_{V'/V}).
\end{equation}
Comparing \hyperlink{sga6.e08.eq.1-8-1}{(1.8.1)} and \hyperlink{sga6.e08.eq.1-8-2}{(1.8.2)} gives the proposition.

\sgaStableTarget{sga6.e08.def.1-9}\textbf{Definition 1.9.} The integer $d_x$ of \hyperlink{sga6.e08.prop.1-8}{Proposition 1.8} is called the \emph{relative virtual dimension} of $X$ over $Y$ (or of $f$) at the point $x$.

The relative virtual dimension of $X$ over $Y$ is therefore a locally constant function on $X$, with integral values, positive or negative. If $f$ is a smooth morphism of relative dimension $n$, its relative virtual dimension is $n$. If $f$ is a regular immersion of codimension $n$, its relative virtual dimension is $-n$.

\sgaStableTarget{sga6.e08.prop.1-10}\textbf{Proposition 1.10.} Let $f:X\to Y$ and $g:Y\to Z$ be two complete-intersection morphisms, and let $x\in X$. Then
\[
\operatorname{dim.rel.virt.}_x(gf)=
\operatorname{dim.rel.virt.}_x(f)+
\operatorname{dim.rel.virt.}_{f(x)}(g).
\]
This follows immediately from \hyperlink{sga6.e08.def.1-9}{1.9} and from the proof of 1.5.

\sgaStableTarget{sga6.e08.sec.2}\subsection*{2. The relative cotangent complex}

\sgaStableTarget{sga6.e08.para.2-1}\textbf{2.1.} We give here a definition of the cotangent complex which, without having the greatest possible generality (see \hyperlink{sga6.e00.sec.4.4}{Expos\'e 0, 4.4)}, will be sufficient for the needs of the Riemann--Roch theorem. Of course the definition given in \hyperlink{sga6.e00.expose}{Expos\'e 0} coincides, in the case considered there, with the definition given here.

Let $f:X\to Y$ be a smoothable morphism, that is, a morphism admitting a factorization
\[
\begin{tikzcd}[column sep=large,row sep=large]
X \arrow[r,"i"] \arrow[dr,"f"'] & V \arrow[d,"g"]\\
& Y,
\end{tikzcd}
\]
where $g$ is smooth and $i$ is an immersion, which may be supposed closed. Let $J$ be the ideal of $X$ in $V$. The canonical derivation
\[
d:\mathcal O_V\longrightarrow\Omega^1_{V/Y}
\]
defines a homomorphism
\[
J\longrightarrow\mathcal O_V\xrightarrow{d}\Omega^1_{V/Y}\longrightarrow
\Omega^1_{V/Y}\otimes_{\mathcal O_V}\mathcal O_X
\]
which vanishes on $J^2$, and hence defines a homomorphism, still denoted $d$,
\sgaStableTarget{sga6.e08.eq.2-1-1}\begin{equation}\tag{2.1.1}
J/J^2\xrightarrow{d} i^*(\Omega^1_{V/Y}).
\end{equation}
One checks easily that this homomorphism is $\mathcal O_X$-linear. One then defines a chain complex $L^{X/Y}_\bullet$ by putting
\[
L^{X/Y}_0=i^*(\Omega^1_{V/Y}),\qquad
L^{X/Y}_1=J/J^2,
\qquad
L^{X/Y}_i=0\quad(i\neq 0,1),
\]
and taking as homomorphism $L^{X/Y}_1\to L^{X/Y}_0$ the homomorphism \hyperlink{sga6.e08.eq.2-1-1}{(2.1.1)}.

\sgaStableTarget{sga6.e08.prop.2-2}\textbf{Proposition 2.2.} The complex $L^{X/Y}_\bullet$ is independent, up to canonical isomorphism in the derived category $D(X)$, of the factorization chosen through a smooth $Y$-scheme.

Let $V,V'$ be two smooth $Y$-schemes through which $f$ factors. Put $V''=V\times_Y V'$. Then $V''$ is smooth over $Y$, and also over $V$ and $V'$. Moreover the immersions $X\to V$ and $X\to V'$ define an immersion $X\to V''$ which factors $f$. It is therefore enough to compare the complexes $L^{X/Y}_\bullet$ attached to two factorizations through smooth $Y$-schemes $V$ and $V'$ for which there is a commutative diagram over $Y$
\[
\begin{tikzcd}[column sep=large,row sep=large]
& V' \arrow[d,"p"]\\
X \arrow[ur,"j"] \arrow[r,"i"'] & V,
\end{tikzcd}
\]
where $i$ and $j$ are closed immersions, and $p$ is smooth. Functoriality then gives homomorphisms
\[
I/I^2\longrightarrow J/J^2,
\qquad
p^*(\Omega^1_{V/Y})\longrightarrow\Omega^1_{V'/Y},
\]
hence a diagram
\sgaStableTarget{sga6.e08.eq.2-2-1}\begin{equation}\tag{2.2.1}
\begin{tikzcd}[column sep=large,row sep=large]
I/I^2 \arrow[r] \arrow[d,"d"'] & J/J^2 \arrow[d,"d"]\\
i^*(\Omega^1_{V/Y}) \arrow[r] & j^*(\Omega^1_{V'/Y}).
\end{tikzcd}
\end{equation}
Since $V$ and $V'$ are smooth over $Y$, one obtains the exact sequence
\[
0\longrightarrow p^*(\Omega^1_{V/Y})\longrightarrow\Omega^1_{V'/Y}
\longrightarrow\Omega^1_{V'/V}\longrightarrow0.
\]
Since the modules of differentials considered are locally free on $\mathcal O_V$, applying the functor $j^*$ gives the exact sequence
\sgaStableTarget{sga6.e08.eq.2-2-2}\begin{equation}\tag{2.2.2}
0\longrightarrow i^*(\Omega^1_{V/Y})\longrightarrow j^*(\Omega^1_{V'/Y})
\longrightarrow j^*(\Omega^1_{V'/V})\longrightarrow0.
\end{equation}
Furthermore, $j$ defines a section $j'$ of $X'=V'\times_V X$:
\[
\begin{tikzcd}[column sep=large,row sep=large]
X' \arrow[r,"i'"] \arrow[d,"p'"'] & V' \arrow[d,"p"]\\
X \arrow[r,"i"'] \arrow[ur,"j"] & V .
\end{tikzcd}
\]
% SGA-SRC-001513
Let $I,J,K$ be the ideals defining $i,j,j'$. Then $j=i'j'$, and the immersion $j'$ is regular; moreover $i'$ is defined by the ideal $p^*(I)$, because $p$ is flat, and its conormal sheaf is therefore
\[
i'^*(p^*(I))=p'^*(I/I^2).
\]
One obtains the exact sequence
\sgaStableTarget{sga6.e08.eq.2-2-3}\begin{equation}\tag{2.2.3}
0\longrightarrow I/I^2\longrightarrow J/J^2\longrightarrow K/K^2\longrightarrow0.
\end{equation}
We shall show: (a) the diagram \hyperlink{sga6.e08.eq.2-2-1}{(2.2.1)} is commutative; this will give a homomorphism of complexes $\theta$ from the complex $L^{X/Y}_\bullet$ defined by $V$ to the complex $L^{X/Y}_\bullet$ defined by $V'$; (b) the homomorphism
\[
K/K^2\longrightarrow j^*(\Omega^1_{V'/V})
\]
defined by \hyperlink{sga6.e08.eq.2-2-2}{(2.2.2)}, \hyperlink{sga6.e08.eq.2-2-3}{(2.2.3)}, and \hyperlink{sga6.e08.eq.2-2-1}{(2.2.1)} is an isomorphism. This will imply that the mapping cylinder of $\theta$ is acyclic, hence that $\theta$ is a quasi-isomorphism, and therefore defines an isomorphism in the derived category.

The properties (a) and (b) are local on $X$. One may therefore suppose $X,Y,V,V'$ affine. Put $Y=\Spec(A)$, $V=\Spec(B)$, $V'=\Spec(B')$, $X=\Spec(C)$, and keep the same letters for the corresponding homomorphisms and ideals. One has the diagram
\[
\begin{tikzcd}[column sep=large,row sep=large]
B\otimes_A B \arrow[r,"p\otimes p"] \arrow[d,"\Delta"'] & B'\otimes_A B' \arrow[d,"\Delta"] & B'\otimes_B B' \arrow[l] \arrow[d]\\
B \arrow[r,"p"] \arrow[dr,"i"'] & B' \arrow[d,"j"] & B'\otimes_B C = B'/IB' \arrow[dl,"p'" ]\\
{} & C & {}
\end{tikzcd}
\]
where the left and right displayed squares are denoted (1) and (2), respectively.
The homomorphism $I/I^2\to J/J^2$ is deduced from the homomorphism $IB'\to J$ and coincides with that deduced from $p:I\to J$; the homomorphism
\[
i^*(\Omega^1_{V/Y})\longrightarrow j^*(\Omega^1_{V'/Y})
\]
is deduced from the commutative square (1). Hence property (a) is evident.

The homomorphism $K/K^2\to j^*(\Omega^1_{V'/V})$ considered in (b) may be defined as follows. Let $\overline x\in K/K^2$, and let $x$ be an element of $J$ lifting $\overline x$ (such an element exists because $K=J/IB'$). To $\overline x$ one associates the image of $x$ by the composite
\[
J\xrightarrow{d}\Omega^1_{B'/A}\otimes_{B'}C
\longrightarrow\Omega^1_{B'/B}\otimes_{B'}C.
\]
On the other hand, let $H$ be the augmentation ideal of
\[
B'\otimes_B B'\longrightarrow B'.
\]
The commutative square (2) gives a homomorphism $H\to K$, hence a homomorphism
\[
H/H^2\otimes_{B'}C\longrightarrow K/K^2.
\]
But $H/H^2=\Omega^1_{B'/B}$. It remains only to check that the homomorphisms thus constructed are inverse to each other, which is straightforward.

\sgaStableTarget{sga6.e08.def.2-3}\textbf{Definition 2.3.} Let $f:X\to Y$ be a smoothable morphism. The complex $L^{X/Y}_\bullet$, defined up to canonical isomorphism in the derived category $D(X)$, is called the \emph{relative cotangent complex} of $X$ over $Y$ (or of $f$).

\sgaStableTarget{sga6.e08.prop.2-4}\textbf{Proposition 2.4.} Let $f:X\to Y$ be a smoothable complete-intersection morphism. Then $L^{X/Y}_\bullet$ is perfect. More precisely, for every factorization of $f$ as the composite of an immersion and a smooth morphism, the complex $L^{X/Y}_\bullet$ defined in \hyperlink{sga6.e08.para.2-1}{2.1} is strictly perfect.

Let $f=g i$, where $g$ is smooth and $i$ is an immersion. By \hyperlink{sga6.e08.prop.1-2}{1.2}, $i$ is a regular immersion. If $J$ is the ideal defining $i$, then $J/J^2$ is a locally free finite-type $\mathcal O_X$-module. Likewise, since $g:V\to Y$ is smooth, $i^*(\Omega^1_{V/Y})$ is locally free of finite type. This proves 2.4.

\sgaStableTarget{sga6.e08.cor.2-5}\textbf{Corollary 2.5.} Let $f:X\to Y$ be a smoothable complete-intersection morphism, and let $f=gi$ be a factorization of $f$ through a smooth $Y$-scheme $V$. Let $J$ be the ideal of $\mathcal O_V$ defining $i$. If $K^{\bullet}(X)$ denotes the Grothendieck group of locally free finite-type $\mathcal O_X$-modules (respectively of perfect complexes on $X$), the element $T_f\in K^{\bullet}(X)$ defined by
\[
T_f=\operatorname{cl}(i^*\Omega^1_{V/Y})-\operatorname{cl}(J/J^2)
\]
respectively
\[
T_f=\operatorname{cl}(L^{X/Y}_\bullet)
\]
does not depend on the chosen factorization of $f$. The value at a point $x$ of $X$ of the augmentation of $T_f$ is precisely the relative virtual dimension of $X$ over $Y$ at $x$ \hyperlink{sga6.e08.def.1-9}{(1.9)}.

The element $T_f$ is well defined in both cases by \hyperlink{sga6.e08.prop.2-4}{2.4}; its class in $K^{\bullet}(X)$ is independent of the factorization, by \hyperlink{sga6.e08.prop.2-2}{2.2} in the second case and by the proof of \hyperlink{sga6.e08.prop.2-2}{2.2} in the first. The last assertion is trivial.

\sgaStableTarget{sga6.e08.prop.2-6}\textbf{Proposition 2.6.} Let $f:X\to Y$ and $g:Y\to Z$ be two smoothable complete-intersection morphisms, and suppose in addition that the following condition is satisfied:

(G) There exists a smooth $Y$-scheme $V$ factorizing $f$ such that, for every closed immersion $Y\to V'$ into a smooth $Z$-scheme, one can find a smooth $V'$-scheme $V''$ with
\[
V \simeq V''\times_{V'}Y.
\]
This condition is in particular satisfied if $f$ is quasi-projective and if $Y$ has an ample sheaf, by taking for $V$ some $\mathbb P^r_Y$.

Then $gf$ is complete-intersection and smoothable, and there is a distinguished triangle
\[
\begin{tikzcd}[column sep=large,row sep=large]
& L^{X/Y}_\bullet \arrow[dl] \\
f^*L^{Y/Z}_\bullet \arrow[rr] && L^{X/Z}_\bullet \arrow[ul]
\end{tikzcd}
\]
in the derived category.

Let $V'$ be a smooth $Z$-scheme factorizing $g$. Consider the commutative diagram
\[
\begin{tikzcd}[column sep=large,row sep=large]
X \arrow[r,"i"] \arrow[dr,"f"'] & V \arrow[r,"j'"] \arrow[d,"p"] & V'' \arrow[d,"p''"]\\
& Y \arrow[r,"j"] \arrow[dr,"g"'] & V' \arrow[d,"p'"]\\
& & Z .
\end{tikzcd}
\]
Then the immersion $j'$ is regular, since it is obtained from $j$ by a flat base change. Hence $j'i$ is a regular immersion factorizing $gf$, which is therefore smoothable and complete-intersection. Moreover, one has the exact sequence of conormal sheaves
\[
0\longrightarrow i^*(J'/J'^2)\longrightarrow K/K^2\longrightarrow I/I^2\longrightarrow0,
\]
where $I,J,J',K$ denote the sheaves defining respectively the immersions $i,j,j',j'i$. Since $p''$ is flat, $J'=p''^*(J)$, and the sequence can be written
\sgaStableTarget{sga6.e08.eq.2-6-1}\begin{equation}\tag{2.6.1}
0\longrightarrow f^*(J/J^2)\longrightarrow K/K^2\longrightarrow I/I^2\longrightarrow0.
\end{equation}
Furthermore, since $V'$ and $V''$ are smooth over $Z$, there is an exact sequence
\[
0\longrightarrow p''^*(\Omega^1_{V'/Z})\longrightarrow\Omega^1_{V''/Z}\longrightarrow\Omega^1_{V''/V'}\longrightarrow0.
\]
Applying $(j'i)^*$, one obtains an exact sequence
\[
0\longrightarrow (p''ji)^*(\Omega^1_{V'/Z})\longrightarrow (j'i)^*(\Omega^1_{V''/Z})\longrightarrow (j'i)^*(\Omega^1_{V''/V'})\longrightarrow0.
\]
Since
\[
j'^*(\Omega^1_{V''/V'})=\Omega^1_{V/Y}
\]
by hypothesis (G), one obtains
\sgaStableTarget{sga6.e08.eq.2-6-2}\begin{equation}\tag{2.6.2}
0\longrightarrow f^*(j^*\Omega^1_{V'/Z})\longrightarrow (j'i)^*(\Omega^1_{V''/Z})\longrightarrow i^*(\Omega^1_{V/Y})\longrightarrow0.
\end{equation}
Combining \hyperlink{sga6.e08.eq.2-6-1}{(2.6.1)} and \hyperlink{sga6.e08.eq.2-6-2}{(2.6.2)}, one is reduced to checking that the diagram
\[
\begin{tikzcd}[column sep=large,row sep=large]
0 \arrow[r] & f^*(J/J^2) \arrow[r] \arrow[d,"f^*(d)"'] & K/K^2 \arrow[r] \arrow[d,"d"] & I/I^2 \arrow[r] \arrow[d,"d"] & 0\\
0 \arrow[r] & f^*(j^*\Omega^1_{V'/Z}) \arrow[r] & (j'i)^*\Omega^1_{V''/Z} \arrow[r] & i^*\Omega^1_{V/Y} \arrow[r] & 0
\end{tikzcd}
\]
is commutative. This is verified locally without difficulty by returning to the definitions of the homomorphisms involved. The announced distinguished triangle follows from the resulting exact sequence of complexes.

\sgaStableTarget{sga6.e08.cor.2-7}\textbf{Corollary 2.7.} With the notation and hypotheses of \hyperlink{sga6.e08.prop.2-6}{2.6}, one has in $K^{\bullet}(X)$ the equality
\[
T_{gf}=T_f+f^*(T_g).
\]

\sgaStableTarget{sga6.e08.sec.3}\subsection*{3. Riemann--Roch theorem: statement}

\sgaStableTarget{sga6.e08.para.3-1}\textbf{3.1.} In no. 3, $Y$ is a quasi-compact scheme having an ample sheaf.

Let $f:X\to Y$ be a projective complete-intersection morphism \hyperlink{sga6.e08.def.1-1}{(1.1)}; suppose moreover that the relative virtual dimension of $f$ \hyperlink{sga6.e08.def.1-9}{(1.9)} is constant on $X$. Since $f$ is projective, there is a closed immersion $X\to P$, where $P$ is a projective bundle over $Y$; since $Y$ has an ample sheaf, one may suppose that $P=\mathbb P^r_Y$. Thus there is a factorization
\[
\begin{tikzcd}[column sep=large,row sep=large]
X \arrow[r,"i"] \arrow[dr,"f"'] & \mathbb P^r_Y \arrow[d,"g"]\\
& Y.
\end{tikzcd}
\]
Since $f$ is complete-intersection, the immersion $i$ is regular \hyperlink{sga6.e08.prop.1-2}{(1.2)}.

By \hyperlink{sga6.e08.prop.1-7}{1.7}, $f$ is a perfect morphism. Moreover $Y$ has an ample sheaf, and therefore so do $X$ and $P$. One can therefore define an additive homomorphism
\[
f_*:K^{\bullet}(X)\longrightarrow K^{\bullet}(Y)
\]
(\hyperlink{sga6.e04.para.2-12}{IV 2.12}), the groups $K^{\bullet}$ defined in terms of locally free sheaves of finite type, or of perfect complexes, being identical. One can also define additive homomorphisms
\[
i_*:K^{\bullet}(X)\longrightarrow K^{\bullet}(P),
\qquad
 g_*:K^{\bullet}(P)\longrightarrow K^{\bullet}(Y),
\]
and one has $f_*=g_*i_*$.

Let $d$ be the relative virtual dimension of $X$ over $Y$. The codimension of $i$ is then $r-d$. By \hyperlink{sga6.e07.thm.4-6}{VII 4.6}, the homomorphism $i_*$ gives, after tensoring with $\mathbb Q$, inclusions
\[
i_*\bigl(\Fil^kK^{\bullet}(X)_{\mathbb Q}\bigr)
\subset
\Fil^{k+r-d}K^{\bullet}(P)_{\mathbb Q}.
\]
By \hyperlink{sga6.e06.cor.5-8}{VI 5.8}, the homomorphism $g_*$ gives inclusions
\[
g_*\bigl(\Fil^{k'}K^{\bullet}(P)\bigr)
\subset
\Fil^{k'-r}K^{\bullet}(Y).
\]
Composing, one obtains the following.

\sgaStableTarget{sga6.e08.prop.3-2}\textbf{Proposition 3.2.} Let $f:X\to Y$ be a projective complete-intersection morphism whose relative virtual dimension $d$ is constant. Suppose that $Y$ is quasi-compact and has an ample sheaf. Then for every $k\in\mathbb Z$,
\[
f_*\bigl(\Fil^kK^{\bullet}(X)_{\mathbb Q}\bigr)
\subset
\Fil^{k-d}K^{\bullet}(Y)_{\mathbb Q}.
\]
Consequently one obtains an additive homomorphism
\[
f^{\mathrm{gr}}_*:\Gr K^{\bullet}(X)_{\mathbb Q}\longrightarrow
\Gr K^{\bullet}(Y)_{\mathbb Q}
\]
sending $\Gr^kK^{\bullet}(X)_{\mathbb Q}$ into $\Gr^{k-d}K^{\bullet}(Y)_{\mathbb Q}$.

\sgaStableTarget{sga6.e08.para.3-3}\textbf{3.3.} Recall that in \hyperlink{sga6.e05.para.6-1}{V 6.1} we defined a functor
\[
A\longmapsto \operatorname{Ch}(A)
\]
from the category of graded $K$-algebras to the category of augmented $K$-$\lambda$-algebras, where $K$ is a binomial ring (\hyperlink{sga6.e05.def.2-7}{V 2.7}). If $A$ is a graded $K$-algebra, then
\[
\operatorname{Ch}(A)=K\times(1+\widehat A^+).
\]
For every augmented $K$-$\lambda$-algebra $\Lambda$ we defined a $K$-$\lambda$-homomorphism, called the completed Chern class (\hyperlink{sga6.e05.para.6-7}{V 6.7}),
\[
\widetilde c=(\varepsilon,c):\Lambda\longrightarrow\operatorname{Ch}(\Gr_{\Lambda}).
\]
Finally, we defined a ring homomorphism, functorial in $A$, called the Chern character (\hyperlink{sga6.e05.def.6-4}{V 6.4}),
\[
\ch:\operatorname{Ch}(A)\longrightarrow K\oplus\widehat A^+_{\mathbb Q},
\]
and a group endomorphism, called the Todd character (\hyperlink{sga6.e05.para.1-16}{V 1.16}),
\[
\Todd:1+\widehat A^+\longrightarrow 1+\widehat A^+.
\]
We shall apply these definitions to the augmented $\lambda$-rings $K^{\bullet}(X)$ and $K^{\bullet}(Y)$; the augmentation rings $H^0(X,\mathbb Z)$ and $H^0(Y,\mathbb Z)$ are binomial (\hyperlink{sga6.e05.ex.2-8}{V 2.8}). To simplify notation, for every scheme $Z$ and every $z\in K^{\bullet}(Z)$ we write
\[
\ch(z)=\ch(\widetilde c(z)),
\qquad
\Todd(z)=\Todd(c(z)).
\]
For every graded algebra $A$ we regard the group $1+\widehat A^+$ as contained in $\widehat A$. Finally, if $f:A\to B$ is additive and compatible with the gradings up to a shift, we shall still write $f$ for the homomorphism $\widehat A\to\widehat B$ deduced from it.

\sgaStableTarget{sga6.e08.para.3-4}\textbf{3.4.} Let $X$ be a scheme. For every locally free finite-type $\mathcal O_X$-module $E$, denote by $\check E$ the class in $K^{\bullet}(X)$ of the dual $\check E$ of $E$. It is clear that the function $E\mapsto\check E$ is additive on the category of locally free finite-type $\mathcal O_X$-modules, and therefore defines an endomorphism of $K^{\bullet}(X)$, denoted $x\mapsto\check x$.

\sgaStableTarget{sga6.e08.prop.3-5}\textbf{Proposition 3.5.} Let $X$ be a quasi-compact scheme. Then, for all $x\in K^{\bullet}(X)$ and all $i\geq0$, one has
\[
c^i(\check x)=(-1)^i c^i(x).
\]
Equivalently,
\[
c_t(\check x)=c_{-t}(x),
\qquad
c_t(x)=\sum_{i\geq0}c^i(x)t^i.
\]
By additivity one may suppose that $x$ is the class of a locally free finite-type $\mathcal O_X$-module $E$. Moreover, after replacing $X$ by the flag scheme of $E$, which gives an injection on the groups $\Gr K^{\bullet}$ (\hyperlink{sga6.e06.cor.5-5}{VI 5.5}), one may suppose that $E$ is an invertible sheaf. Then $\check E=E^{-1}$, so $\check x=x^{-1}$.

The $c^i(\check x)$ and $c^j(x)$ are zero for $i,j\geq2$. It is enough to prove
\[
c^1(x^{-1})=-c^1(x).
\]
Now $c^1(x)$ is the class in $\Gr^1K^{\bullet}(X)$ of $x-1$, while $c^1(x^{-1})$ is that of $x^{-1}-1$. The equality
\[
x^{-1}-1=1-x+(x-1)(1-x^{-1})
\]
gives
\[
x^{-1}-1=1-x\pmod{\Fil^2K^{\bullet}(X)},
\]
and the result follows.

\sgaStableTarget{sga6.e08.thm.3-6}\textbf{Theorem 3.6 (Riemann--Roch).} Let $f:X\to Y$ be a projective complete-intersection morphism \hyperlink{sga6.e08.def.1-1}{(1.1)}, whose relative virtual dimension $d$ is constant \hyperlink{sga6.e08.def.1-9}{(1.9)}. Suppose that $Y$ is quasi-compact and admits an ample sheaf. With the notation of \hyperlink{sga6.e08.para.3-3}{3.3}, for every $x\in K^{\bullet}(X)$ one has
\[
\ch(f_*(x))=f^{\mathrm{gr}}_*\bigl(\ch(x)\Todd(\check T_f)\bigr),
\]
where $\check T_f$ denotes the dual element of the class of the relative cotangent complex $L^{X/Y}_\bullet$ (\hyperlink{sga6.e08.cor.2-5}{2.5} and \hyperlink{sga6.e08.def.2-3}{2.3}).

In other words, the diagram
\[
\begin{tikzcd}[column sep=huge,row sep=large]
K^{\bullet}(X) \arrow[r,"\ch\circ\widetilde c"] \arrow[d,"f_*"'] & \Gr K^{\bullet}(X)_{\mathbb Q} \arrow[d,"x\mapsto f^{\mathrm{gr}}_*(x\Todd(\check T_f))"]\\
K^{\bullet}(Y) \arrow[r,"\ch\circ\widetilde c"'] & \Gr K^{\bullet}(Y)_{\mathbb Q}
\end{tikzcd}
\]
is commutative.



\sgaStableTarget{sga6.e08.continuation-of-the-proof-of-theorem-3-6}\subsection*{Continuation of the proof of Theorem 3.6}

The proof of \hyperlink{sga6.e08.thm.3-6}{Theorem 3.6} will occupy the end of this exposé. Let us give some preliminaries.

\sgaStableTarget{sga6.e08.prop.3-7}\textbf{Proposition 3.7.} For arbitrary
\[
x\in \Gr K^{\bullet}(X)_{\mathbb Q},\qquad
y\in \Gr K^{\bullet}(Y)_{\mathbb Q},
\]
one has
\[
 f^{\mathrm{gr}}_*\bigl(x\cdot f^{\mathrm{gr}}(y)\bigr)
 =
 f^{\mathrm{gr}}_*(x)\,y .
\]

By additivity, one may suppose \(x\) and \(y\) homogeneous; say
\[
x\in \Gr^iK^{\bullet}(X)_{\mathbb Q},\qquad
y\in \Gr^jK^{\bullet}(Y)_{\mathbb Q}.
\]
Let
\[
x'\in \Fil^iK^{\bullet}(X)_{\mathbb Q},\qquad
y'\in \Fil^jK^{\bullet}(Y)_{\mathbb Q}
\]
be elements lifting \(x\) and \(y\). Then the projection formula gives
\[
 f_*\bigl(x'\cdot f^*(y')\bigr)=f_*(x')\,y'.
\]
By \hyperlink{sga6.e08.prop.3-2}{3.2}, this equality is written in
\(\Fil^{i+j-d}K^{\bullet}(Y)_{\mathbb Q}\). One deduces from it an equality in
\(\Gr^{i+j-d}K^{\bullet}(Y)_{\mathbb Q}\), and one verifies immediately that this is the required equality.

\sgaStableTarget{sga6.e08.lem.3-8}\textbf{Lemma 3.8.} Let \(f:X\to Y\) and \(g:Y\to Z\) be two morphisms satisfying the hypotheses of 3.6. Then \(g\circ f\) also satisfies them. If \(x\) is an element of \(K^{\bullet}(X)\), and if \hyperlink{sga6.e08.thm.3-6}{3.6} is verified by \(f\) for the element \(x\), and by \(g\) for the element \(f_*(x)\), then it is verified by \(g\circ f\) for the element \(x\).

By hypothesis, one has
\sgaStableTarget{sga6.e08.eq.3-8-1}\begin{equation}\tag{3.8.1}
 \ch(f_*(x))=
 f^{\mathrm{gr}}_*\bigl(\ch(x)\Todd(\check T_f)\bigr),
\end{equation}
and
\sgaStableTarget{sga6.e08.eq.3-8-2}\begin{equation}\tag{3.8.2}
 \ch(g_*(f_*(x)))=
 g^{\mathrm{gr}}_*\bigl(\ch(f_*(x))\Todd(\check T_g)\bigr).
\end{equation}
Combining \hyperlink{sga6.e08.eq.3-8-1}{(3.8.1)} and \hyperlink{sga6.e08.eq.3-8-2}{(3.8.2)}, one obtains
\[
\ch(g_*(f_*(x)))=
g^{\mathrm{gr}}_*\!\left(
 f^{\mathrm{gr}}_*\bigl(\ch(x)\Todd(\check T_f)\bigr)\cdot \Todd(\check T_g)
\right).
\]
Applying \hyperlink{sga6.e08.prop.3-7}{3.7}, one finds
\[
\ch(g_*(f_*(x)))=
g^{\mathrm{gr}}_*\!\left(
 f^{\mathrm{gr}}_*\bigl(\ch(x)\Todd(\check T_f)\cdot
 f^{\mathrm{gr}}(\Todd(\check T_g))\bigr)
\right).
\]
The functoriality of \(c\) and of \(\Todd\) implies
\[
 f^{\mathrm{gr}}(\Todd(\check T_g))=\Todd(f^*(\check T_g)).
\]
By \hyperlink{sga6.e08.cor.2-7}{2.7}, one has
\[
 \check T_{g\circ f}=\check T_f+f^*(\check T_g).
\]
The multiplicativity of the Todd character therefore gives
\[
 \Todd(\check T_f)\,f^{\mathrm{gr}}(\Todd(\check T_g))
 =
 \Todd(\check T_{g\circ f}).
\]
Thus finally
\[
 \ch(g_*(f_*(x)))=
 g^{\mathrm{gr}}_* f^{\mathrm{gr}}_*\bigl(\ch(x)\Todd(\check T_{g\circ f})\bigr),
\]
which is the announced result.

\sgaStableTarget{sga6.e08.para.3-9}\textbf{3.9.} If one applies \hyperlink{sga6.e08.lem.3-8}{3.8} to the factorization of \(f\) given in \hyperlink{sga6.e08.para.3-1}{3.1}, one sees that, to prove the Riemann--Roch theorem, it is enough to prove it in the case of a regular immersion of constant codimension, and in the case of the structural morphism
\[
 \mathbb P^r_Y\longrightarrow Y .
\]
Moreover, one may suppose that the immersion \(i\) satisfies the condition of \hyperlink{sga6.e07.lem.4-1}{VII 4.1}. Indeed, let \(s\) be the section of
\[
 \mathbb P^2_{\mathbb P^r_Y}
\]
introduced in \hyperlink{sga6.e07.sec.4.2}{VII 4.2}. Then \(s\circ i\) is a regular immersion satisfying the condition of \hyperlink{sga6.e07.lem.4-1}{VII 4.1} (\hyperlink{sga6.e07.sec.4.2}{VII 4.2}), and by \hyperlink{sga6.e08.lem.3-8}{3.8} it is enough to prove \hyperlink{sga6.e08.thm.3-6}{3.6} for
\[
s\circ i,\qquad
\mathbb P^2_{\mathbb P^r_Y}\longrightarrow\mathbb P^r_Y,\qquad
\mathbb P^r_Y\longrightarrow Y .
\]

We shall therefore prove \hyperlink{sga6.e08.thm.3-6}{3.6} by proving the following two particular cases:
\begin{enumerate}[label=\roman*)]
\item \(f\) is a regular immersion of constant codimension, satisfying the condition of \hyperlink{sga6.e07.lem.4-1}{VII 4.1};
\item \(f\) is the structural morphism \(\mathbb P^r_Y\to Y\).
\end{enumerate}

\sgaStableTarget{sga6.e08.sec.4}\section*{4. Riemann--Roch theorem: case of a regular closed immersion}

\sgaStableTarget{sga6.e08.para.4-1}\textbf{4.1.} Let \(i:Y\to X\) be a regular closed immersion. The Riemann--Roch theorem for \(i\) is essentially relation \hyperlink{sga6.e07.cor.4-4}{VII 4.4}, which makes it possible to compute the Chern classes of \(i_*(x)\) from the \(\gamma^p(N,x)\). More precisely, if one introduces the \(\lambda\)-ring \(K^{\bullet}(Y)_{\mathbb Q N}\) defined in \hyperlink{sga6.e05.para.5-5}{V 5.5}, then the homomorphism
\[
K^{\bullet}(Y)_{\mathbb Q N}\longrightarrow K^{\bullet}(X)_{\mathbb Q}
\]
deduced from \(i_*\) is a \(\lambda\)-homomorphism by \hyperlink{sga6.e07.cor.2-8}{VII 2.8} and \hyperlink{sga6.e07.thm.4-3}{VII 4.3}. The Riemann--Roch theorem is then reduced to an analogous formula concerning the additive homomorphism
\[
K^{\bullet}(Y)\longrightarrow K^{\bullet}(Y)_{\mathbb Q N}
\]
defined by \(y\mapsto(0,y)\), that is, to a formal relation valid for every \(\lambda\)-ring. We shall not develop this proof here, in order to avoid introducing still more formalism on \(\lambda\)-rings (cf. \(\RRR\), I 1.5), and we shall resume the method of [1], which used blow-up schemes. This is the method already employed for the proof of \hyperlink{sga6.e07.thm.4-3}{VII 4.3}, and also the method at the end of \(\RRR\).

As indicated in \hyperlink{sga6.e08.para.3-9}{3.9}, one may restrict to the case of a regular immersion of constant codimension \(d\) satisfying the condition of \hyperlink{sga6.e07.lem.4-1}{VII 4.1}; the case of a general regular immersion is then deduced from 3.8. In this paragraph we resume the notation of \hyperlink{sga6.e07.para.3-1}{VII 3.1}. Thus one has the Cartesian square
\[
\begin{tikzcd}[column sep=large,row sep=large]
Y' \arrow[r,"j"] \arrow[d,"g"'] & X' \arrow[d,"f"]\\
Y \arrow[r,"i"'] & X ,
\end{tikzcd}
\]
where \(X'\) is the \(X\)-scheme obtained by blowing up \(Y\). We suppose moreover that \(X\) is quasi-compact and has an ample sheaf.

\sgaStableTarget{sga6.e08.prop.4-2}\textbf{Proposition 4.2.} The morphism \(f\) is projective, a complete-intersection morphism, and of virtual relative dimension zero. In addition one has the relation
\[
 f^{\mathrm{gr}}_* f^{\mathrm{gr}}=\operatorname{Id}_{\Gr K^{\bullet}(X)_{\mathbb Q}} .
\]
The first assertion follows immediately from \hyperlink{sga6.e07.prop.1-8}{VII 1.8}. The second relation is a consequence of the relation
\[
 f_*f^*=\operatorname{Id}_{K^{\bullet}(X)}
\]
proved in \hyperlink{sga6.e07.prop.3-6}{VII 3.6}.

\sgaStableTarget{sga6.e08.prop.4-3}\textbf{Proposition 4.3.} For every \(y\in K^{\bullet}(Y)_{\mathbb Q}\), one has the relation
\[
 f^{\mathrm{gr}}_* i^{\mathrm{gr}}_*(y)
 =
 j^{\mathrm{gr}}_*\bigl(g^{\mathrm{gr}}(y)\,c^{d-1}(\check F)\bigr),
\]
where \(c^{d-1}\) denotes the \((d-1)\)-st Chern class.

This relation is contained in \hyperlink{sga6.e07.thm.4-8}{VII 4.8}, for it is nothing other than
\[
 v^{\mathrm{gr}}\circ u^{\mathrm{gr}}=0 .
\]

\sgaStableTarget{sga6.e08.prop.4-4}\textbf{Proposition 4.4.} Let \(X\) be a quasi-compact scheme. If \(E\) is a locally free sheaf of finite type on \(X\), of class \(E\) in \(K^{\bullet}(X)\), and of rank \(n\), then
\[
 \ch(\lambda_{-1}(E))=c^n(\check E)\Todd(\check E)^{-1}.
\]

The two sides of the relation to be proved are multiplicative. This is clear for \(\ch\circ\lambda_{-1}\) and for \(\Todd\); it is also true for \(c^n\), which is the class of \((-1)^nE\lambda_{-1}(E)\). Consequently, by passing to the flag bundle of \(E\), one may suppose that \(E\) is a sum of classes of invertible sheaves, and reduce to proving the formula for an invertible sheaf. Let \(\xi\) be the class of such a sheaf. Then
\[
 \ch(\lambda_{-1}(\xi))=\ch(1-\xi)=1-\ch(\xi).
\]
Since \(\widetilde c(\xi)=(1,1+c^1(\xi))\), the definition of \(\ch\) gives
\[
 \ch(\lambda_{-1}(\xi))=1-\exp(c^1(\xi)).
\]
Moreover, \(\xi\) has augmentation \(1\); the Chern class which intervenes in formula \hyperlink{sga6.e08.prop.4-4}{4.4} relative to \(\xi\) is therefore \(c^1\). Taking into account the definition of \(\Todd\), and \hyperlink{sga6.e08.prop.3-5}{3.5}, the relation to be proved is written
\[
 1-\exp(c^1(\xi))
 =
 -c^1(\xi)\left(\frac{-c^1(\xi)}{1-\exp(c^1(\xi))}\right)^{-1},
\]
and is therefore verified.

\sgaStableTarget{sga6.e08.rem.4-4-1}\textbf{Remark 4.4.1.} It is also possible to deduce \hyperlink{sga6.e08.prop.4-4}{4.4} from a corresponding purely algebraic statement; cf. \(\RRR\), \hyperlink{sga6.e01.sec.1.3}{I 1.3}.

\sgaStableTarget{sga6.e08.para.4-5}\textbf{4.5.} We shall now show that, in order to prove the Riemann--Roch theorem in the case considered in this paragraph, one may reduce to the case where \(Y\) is a divisor of \(X\), and where the element to which the formula is applied is the inverse image of an element of \(K^{\bullet}(X)\).

By \hyperlink{sga6.e08.prop.3-2}{3.2}, the formula to be proved,
\[
 \ch(i_*(x))=
 i^{\mathrm{gr}}_*\bigl(\ch(x)\Todd(\check T_i)\bigr),
\]
is equivalent to the formula
\[
 f^{\mathrm{gr}}\bigl(\ch(i_*(x))\bigr)
 =
 f^{\mathrm{gr}}i^{\mathrm{gr}}_*\bigl(\ch(x)\Todd(\check T_i)\bigr).
\]
Taking account of the functoriality of \(\ch\), and of \hyperlink{sga6.e08.prop.4-3}{4.3}, it may also be written
\[
 \ch\bigl(f^*(i_*(x))\bigr)
 =
 j^{\mathrm{gr}}_*\bigl(\ch(g^*(x))\Todd(g^*(\check T_i))\,c^{d-1}(\check F)\bigr),
\]
or, by \hyperlink{sga6.e07.prop.3-4}{VII 3.4},
\[
 \ch\bigl(j_*(g^*(x)\lambda_{-1}(F))\bigr)
 =
 j^{\mathrm{gr}}_*\bigl(\ch(g^*(y))\Todd(g^*(\check T_i))\,c^{d-1}(\check F)\bigr).
\]

By definition,
\[
 T_i=-c^1(J/J^2)=-N,\qquad
 T_j=-c^1(J'/J'^2)=-L.
\]
Moreover, the exact sequence
\[
0\longrightarrow F\longrightarrow g^*(N)\longrightarrow \mathcal O_{Y'}(1)\longrightarrow0
\]
gives the relation
\[
 g^*(T_i)=-(F+L),
\]
hence
\[
 \Todd(g^*(\check T_i))=\Todd(\check T_j)\Todd(\check F)^{-1}.
\]
Taking account of \hyperlink{sga6.e08.prop.4-4}{4.4}, the relation to be proved is therefore written
\[
 \ch\bigl(j_*(g^*(x)\lambda_{-1}(F))\bigr)
 =
 j^{\mathrm{gr}}_*\bigl(\ch(g^*(x)\lambda_{-1}(F))\Todd(\check T_j)\bigr).
\]
Putting
\[
z=g^*(x)\lambda_{-1}(F),
\]
this relation is nothing but the Riemann--Roch theorem applied to the regular immersion of codimension \(1\), namely \(j\), and to the element \(z\in K^{\bullet}(Y')\).

Furthermore, since \(i\) was supposed to satisfy the hypothesis of \hyperlink{sga6.e07.lem.4-1}{VII 4.1}, \(\lambda_{-1}(F)\) is divisible by \(1-L\). Thus one can find \(z'\) such that \(z=z'(1-L)\), or still
\[
 z=j^*(j_*(z'))
\]
(\hyperlink{sga6.e07.para.2-6}{Exp. VII 2.7)}. We are therefore reduced to the case where \(Y\) is a divisor, and to an element \(x\) for which there exists \(y\), with \(x=i^*(y)\).

\sgaStableTarget{sga6.e08.para.4-6}\textbf{4.6.} It remains to show the relation
\[
 \ch\bigl(i_*(i^*(y))\bigr)
 =
 i^{\mathrm{gr}}_*\bigl(\ch(i^*(y))\Todd(\check T_i)\bigr).
\]
The projection formula gives
\[
 i_*(i^*(y))=y\,i_*(1),
\]
and from the exact sequence
\[
0\longrightarrow J\longrightarrow\mathcal O_X\longrightarrow\mathcal O_Y\longrightarrow0,
\]
where \(J\) is here invertible, one obtains
\[
 i_*(1)=1-J.
\]
Thus
\[
 \ch(i_*(i^*(y)))=\ch(y)\,\ch(1-J).
\]
Since \(J\) is invertible, one has
\[
 \widetilde c(J)=(1,1+c^1(J)),\qquad \ch(J)=\exp(c^1(J)).
\]
Consequently
\[
 \ch(i_*(i^*(y)))=\ch(y)(1-\exp(c^1(J))).
\]

On the other hand, the functoriality of \(\ch\) and the projection formula \hyperlink{sga6.e08.prop.3-7}{3.7} imply
\[
 i^{\mathrm{gr}}_*\bigl(\ch(i^*(y))\Todd(\check T_i)\bigr)
 =
 \ch(y)\,i^{\mathrm{gr}}_*(\Todd(\check T_i)).
\]
Now
\[
T_i=-c^1(J/J^2)=-i^*(J).
\]
Hence
\[
 i^{\mathrm{gr}}_*(\Todd(\check T_i))
 =
 i^{\mathrm{gr}}_*\,i^{\mathrm{gr}}(\Todd(-\check J))
 =
 i^{\mathrm{gr}}_*(1)\Todd(-\check J).
\]
Since \(i_*(1)=1-J\), \(i^{\mathrm{gr}}_*(1)\) is the image in \(\Gr^1K^{\bullet}(X)\) of \(1-J\), hence it is \(-c^1(J)\). Moreover
\[
 \Todd(-\check J)=\Todd(J)^{-1}
 =
 \frac{1-\exp(c^1(J))}{-c^1(J)}.
\]
Thus finally
\[
 i^{\mathrm{gr}}_*\bigl(\ch(i^*(y))\Todd(\check T_i)\bigr)
 =
 \ch(y)(1-\exp(c^1(J))),
\]
which proves the relation to be shown.

\sgaStableTarget{sga6.e08.sec.5}\section*{5. Riemann--Roch theorem: case of the structural morphism of a projective bundle}

\sgaStableTarget{sga6.e08.para.5-1}\textbf{5.1.} By \hyperlink{sga6.e08.para.3-9}{3.9}, in order to prove the Riemann--Roch theorem, it remains to prove it for the structural morphism
\[
f:\mathbb P^r_Y\longrightarrow Y,
\]
where \(Y\) is a quasi-compact scheme admitting an ample sheaf. Put
\[
P=\mathbb P^r_Y,\qquad \xi=c^1(\mathcal O_P(1)),\qquad E=(\mathcal O_Y)^{r+1}.
\]

\sgaStableTarget{sga6.e08.prop.5-2}\textbf{Proposition 5.2.} On \(P\) one has the exact sequence
\[
0\longrightarrow\Omega^1_{P/Y}\longrightarrow f^*(E)(-1)\longrightarrow\mathcal O_P\longrightarrow0,
\]
i.e. one has an isomorphism
\[
\Omega^1_{P/Y}\simeq F(-1),
\qquad
F=\Ker\bigl(f^*(E)\longrightarrow\mathcal O_P(1)\bigr).
\]

Consider the vector bundle over \(Y\) defined by \(E\), and denote by \(Z\) the complement of the zero section of this bundle. One has the commutative diagram
\[
\begin{tikzcd}[column sep=large,row sep=large]
P \arrow[d,"f"'] & Z \arrow[l,"h"'] \arrow[dl,"g"]\\
Y &
\end{tikzcd}
\]
The morphism \(h\) identifies \(Z\) with the affine \(P\)-scheme defined by the sheaf of algebras
\[
B=\bigoplus_{n\in\mathbb Z}\mathcal O_P(n);
\]
\(B\) is a graded \(\mathcal O_P\)-algebra.

Since \(f\), \(g\), and \(h\) are smooth morphisms, one may write the exact sequence of sheaves of relative differentials. Since \(Z\) is affine over \(P\), one may write this sequence in terms of \(B\)-modules; thus one obtains
\sgaStableTarget{sga6.e08.eq.5-2-1}\begin{equation}\tag{5.2.1}
0\longrightarrow
 \Omega^1_{P/Y}\otimes_{\mathcal O_P}B
\longrightarrow
 \Omega^1_{B/Y}
\longrightarrow
 \Omega^1_{B/\mathcal O_P}
\longrightarrow0 .
\end{equation}
This exact sequence is an exact sequence of graded \(B\)-modules. Indeed, it follows from the commutative diagram
\[
\begin{tikzcd}[column sep=large,row sep=large]
S(E) \arrow[rr] && \displaystyle\bigoplus_{n\in\mathbb Z}S(E)(n)\simeq S(E)[T,T^{-1}]\\
& \mathcal O_X \arrow[ul] \arrow[ur] &
\end{tikzcd}
\]
where the arrows are algebra homomorphisms, compatible with the grading of \(S(E)[T,T^{-1}]\) (grading defined by the degree in the indeterminate \(T\)), the two other algebras being supposed to be reduced to degree \(0\).

The module of differentials \(\Omega^1_{B/\mathcal O_P}\) is isomorphic to
\[
 B(-1)\otimes_{\mathcal O_P}\mathcal O_P(1),
\]
for one may define an \(S(E)\)-derivation
\[
S(E)[T,T^{-1}]
\longrightarrow
S(E)[T,T^{-1}]\otimes_{S(E)}S(E)\,dT
\]
satisfying the universal property which defines \(\Omega^1_{B/\mathcal O_P}\), by setting
\[
 d(aT^n)=naT^{n-1}\otimes dT;
\]
moreover, \(d\) is compatible with the gradings if the degree \(1\) is attributed to the factor \(S(E)\,dT\).

Let \(V\) be the vector bundle defined by \(E\) over \(Y\). Then \(\Omega^1_{Z/Y}\) is the restriction to \(Z\) of \(\Omega^1_{V/Y}\). But \(V\) is the affine \(Y\)-scheme defined by \(S(E)\), so \(\Omega^1_{V/Y}\) is defined by
\[
S(E)\otimes_{\mathcal O_Y}E
\]
(since \(E\) is locally free of finite type), and therefore \(\Omega^1_{V/Y}\) is equal to \(g^*(E)\). Consequently
\[
\Omega^1_{B/Y}=f^*(E)\otimes_{\mathcal O_P}B.
\]
To see its grading, one looks locally, and one sees at once that \(f^*(E)\) must be considered as being of degree \(1\). Taking the terms of degree \(0\) in the exact sequence \hyperlink{sga6.e08.eq.5-2-1}{(5.2.1)}, one obtains the announced exact sequence.

\sgaStableTarget{sga6.e08.rem.5-2-1}\textbf{Remark 5.2.1.} Using the general results on the differential study of Hilbert functors ([2], 5), one directly obtains an isomorphism
\[
\Omega^1_{P/Y}\xrightarrow{\sim}F(-1)
\]
from the functorial characterization given of \(g^*(\Omega^1_{P/Y})\) for every \(Y\)-morphism \(g:Y'\to P\).

\sgaStableTarget{sga6.e08.para.5-3}\textbf{5.3.} By \hyperlink{sga6.e08.prop.5-2}{5.2}, one has
\[
 T_f=c_1(\Omega^1_{P/Y})=f^*(E)\xi^{-1}-1.
\]
Consequently
\[
 \widetilde c(T_f)=\widetilde c(E)\,\widetilde c(\xi^{-1})-\widetilde c(1).
\]
Since \(E=(\mathcal O_Y)^{r+1}\), one has \(E=r+1\), and hence
\[
 \widetilde c(E)=(r+1,1).
\]
Put
\[
 x=c^1_{\Gr^1K^{\bullet}(P)}(\xi-1)
 =
 c^1_{\Gr^1K^{\bullet}(P)}(1-\xi^{-1})
 =
 c^1(\xi).
\]
Since \(\xi^{-1}=\check \xi\), one has
\[
 \widetilde c(\xi^{-1})=(1,1-x).
\]
By \hyperlink{sga6.e05.eq.6-2-1}{V (6.2.1)}, one deduces
\[
 \widetilde c(T_f)=(r,(1-x)^{r+1}).
\]
Taking \hyperlink{sga6.e08.prop.3-5}{3.5} into account, one finally obtains
\[
 \widetilde c(\check T_f)=(r,(1+x)^{r+1}).
\]
By the definition of \(\Todd\), one therefore has
\sgaStableTarget{sga6.e08.eq.5-3-1}\begin{equation}\tag{5.3.1}
 \Todd(\check T_f)=
 \left(\frac{x}{1-\exp(-x)}\right)^{r+1}.
\end{equation}

\sgaStableTarget{sga6.e08.para.5-4}\textbf{5.4.} Using the basis of \(K^{\bullet}(P)\) over \(K^{\bullet}(Y)\) formed by the elements \(\xi^k\), with \(-r\leq k\leq0\), one sees by additivity that it is enough to prove the theorem for elements of the form
\[
 f^*(a)\xi^k .
\]
Then
\[
 \ch\bigl(f_*(f^*(a)\xi^k)\bigr)
 =
 \ch(a f_*(\xi^k))
 =
 \ch(a)\,\ch(f_*(\xi^k)).
\]
On the other hand,
\[
 f^{\mathrm{gr}}_*\bigl(\ch(f^*(a)\xi^k)\Todd(\check T_f)\bigr)
 =
 f^{\mathrm{gr}}_*\bigl(f^{\mathrm{gr}}(\ch(a))\ch(\xi^k)\Todd(\check T_f)\bigr)
 =
 \ch(a)\,f^{\mathrm{gr}}_*\bigl(\ch(\xi^k)\Todd(\check T_f)\bigr).
\]
Thus it is enough to prove the theorem for the \(\xi^k\), with \(-r\leq k\leq0\). By \hyperlink{sga6.e06.sec.5.2}{VI (5.2.3)}, one has
\[
 f_*(\xi^k)=0\quad\text{for }-r\leq k<0,\qquad f_*(1)=1.
\]
Consequently
\sgaStableTarget{sga6.e08.eq.5-4-1}\begin{equation}\tag{5.4.1}
 \ch(f_*(\xi^k))=0\quad\text{if }-r\leq k<0,\qquad
 \ch(f_*(1))=1.
\end{equation}

We now calculate the second member of the Riemann--Roch formula. For every \(k\), one has
\[
 \widetilde c(\xi^k)=(\widetilde c(\xi))^k=(1,1+x)^k=(1,1+kx)
\]
by \hyperlink{sga6.e05.eq.6-2-1}{V (6.2.1)}. Hence
\sgaStableTarget{sga6.e08.eq.5-4-2}\begin{equation}\tag{5.4.2}
 \ch(\xi^k)=\exp(kx).
\end{equation}
The second member of the Riemann--Roch formula is therefore
\sgaStableTarget{sga6.e08.eq.5-4-3}\begin{equation}\tag{5.4.3}
 f^{\mathrm{gr}}_*\left(\exp(kx)\,
 \frac{x^{r+1}}{(1-\exp(-x))^{r+1}}\right).
\end{equation}
Since \(x\) is the class in \(\Gr^1K^{\bullet}(P)\) of \(1-\xi^{-1}\), \(f^{\mathrm{gr}}_*(x^n)\) is, for \(0\leq n\leq r\), the class in
\(\Gr^{n-r}K^{\bullet}(Y)\) of \(f_*((1-\xi^{-1})^n)=1\). Thus
\sgaStableTarget{sga6.e08.eq.5-4-4}\begin{equation}\tag{5.4.4}
 f^{\mathrm{gr}}_*(x^n)=0\quad\text{if }0\leq n<r,\qquad
 f^{\mathrm{gr}}_*(x^r)=1.
\end{equation}
Finally recall that the elements
\[
1,x,\ldots,x^{r+1}
\]
are linearly dependent over \(\Gr K^{\bullet}(Y)\) (\hyperlink{sga6.e06.cor.5-6}{VI 5.6}), and satisfy the equation
\sgaStableTarget{sga6.e08.eq.5-4-5}\begin{equation}\tag{5.4.5}
 x^{r+1}=0,
\end{equation}
and that \(1,\ldots,x^r\) are linearly independent over \(\Gr K^{\bullet}(Y)\).

One can therefore write the series
\[
\varphi_k=\exp(kx)\,
\frac{x^{r+1}}{(1-\exp(-x))^{r+1}}
\]
in a unique way as a linear combination, with coefficients in
\(\Gr K^{\bullet}(Y)_{\mathbb Q}\), of the elements \(1,\ldots,x^r\). The image under \(f^{\mathrm{gr}}_*\) of \(\varphi_k\) is then, by \hyperlink{sga6.e08.eq.5-4-4}{(5.4.4)} and \hyperlink{sga6.e08.prop.3-7}{3.7}, its coefficient of \(x^r\). Taking account of \hyperlink{sga6.e08.eq.5-4-1}{(5.4.1)}, the Riemann--Roch theorem for \(f\) is therefore equivalent to
\sgaStableTarget{sga6.e08.eq.rr}\[
\left\{
\begin{array}{ll}
\text{for }-r\leq k<0,\text{ the coefficient of }x^r\text{ in }\varphi_k\text{ is }0,\\
\text{for }k=0,\text{ the coefficient of }x^r\text{ in }\varphi_k\text{ is }1.
\end{array}
\right.
\tag{RR}
\]
This coefficient is the residue of the differential form
\[
 \frac{\exp(kx)\,dx}{(1-\exp(-x))^{r+1}}.
\]
Putting \(y=1-\exp(-x)\), one is reduced to calculating the residue of
\[
 \frac{(1-y)^{-k-1}\,dy}{y^{r+1}},
\]
that is, the coefficient of \(y^r\) in \((1-y)^{-k-1}\). For \(-r\leq k<0\), it is zero, and for \(k=0\), it is \(1\).

This completes the proof of the Riemann--Roch theorem.

\sgaStableTarget{sga6.e08.bibliography}\subsection*{Bibliography}
\begin{enumerate}[label={[\arabic*]}]
\item A. Borel and J.-P. Serre, \emph{Le théorème de Riemann--Roch}, Bull. Soc. Math. de France, t. 86, 1958.
\item A. Grothendieck, \emph{Techniques de construction et théorèmes d'existence en Géométrie Algébrique}, Sém. Bourbaki, 1960--61.
\end{enumerate}



\sgaStableTarget{sga6.e09.expose}\section*{Exposé IX. Some computations of $K_\bullet$-groups}
\begin{center}
by P. Berthelot
\end{center}

This exposé corresponds to no oral exposé and will not be used in the sequel of the seminar. It contains a certain number of computations of the groups $K_\bullet$ for various fibrations: vector bundles, projective bundles, torsors under a linear group, and so on. Throughout the exposé, the base scheme will be assumed noetherian, and the notation $K_\bullet(X)$ will denote the Grothendieck group of coherent sheaves on a scheme $X$. Recall that, if $X$ is noetherian, the homomorphism from the $K$ of the category of coherent $\mathcal O_X$-modules to the $K$ of the category of pseudo-coherent complexes on $X$ is an isomorphism (\hyperlink{sga6.e04.para.2-4}{IV 2.4}).

The results set out here are the analogues of those one obtains for the Chow ring, and which are to be found in Grothendieck's exposés in the Chevalley seminar of 1958. The question also arises of determining to what extent statements of the same type can be obtained for the graded group associated with $K_\bullet$ for the filtration defined by the dimension of the support (cf. X). It does not seem that the answer is known, even for \hyperlink{sga6.e09.prop.1-1}{Proposition 1.1}, the homotopy exact sequence, which is one of the properties most used in the computations which follow.

\sgaStableTarget{sga6.e09.sec.1}\subsection*{1. Vector bundles}

\sgaStableTarget{sga6.e09.prop.1-1}\textbf{Proposition 1.1 (homotopy exact sequence).} Let $X$ be a noetherian scheme, let $X'$ be a closed subscheme of $X$, let $U$ be the open complement of $X'$ in $X$, let $i$ be the immersion of $X'$ in $X$, and let $j$ be the immersion of $U$ in $X$. Then the sequence
\[
 K_\bullet(X')\xrightarrow{i_*}K_\bullet(X)\xrightarrow{j^*}K_\bullet(U)\longrightarrow 0
\]
is exact.

If $F'$ is a coherent $\mathcal O_{X'}$-module, $i_*(F')$ is an $\mathcal O_X$-module whose support is contained in $X'$; consequently its restriction to $U$ is zero, and $j^*i_*(\cl(F'))=0$. Since the elements of the form $\cl(F')$ generate $K_\bullet(X')$, $j^*i_*=0$. We therefore obtain a homomorphism
\[
 \rho:K_\bullet(X)/\operatorname{Im}K_\bullet(X')\longrightarrow K_\bullet(U).
\]
We shall define an inverse homomorphism to $\rho$. Let $G$ be a coherent $\mathcal O_U$-module. Since $X$ is noetherian, there exists a coherent $\mathcal O_X$-module $F$ whose restriction to $U$ is isomorphic to $G$. If $F'$ is a second coherent $\mathcal O_X$-module whose restriction to $U$ is isomorphic to $G$, then in $K_\bullet(X)$ one has
\[
 \cl(F)-\cl(F')\in \operatorname{Im}K_\bullet(X').
\]
To see this, consider $G$ as a subsheaf of $j^*(F\oplus F')$ by means of the diagonal map. Then $G$ extends to a subsheaf $F''$ of $F\oplus F'$, and one is reduced to proving that
\[
 \cl(F)-\cl(F'')\in \operatorname{Im}K_\bullet(X')
\]
and likewise for $F'$. Let $H$ and $K$ be the kernel and cokernel of the composite homomorphism
\[
 F''\longrightarrow F\oplus F'\longrightarrow F.
\]
This homomorphism restricts on $U$ to the identity of $G$, and hence $H$ and $K$ have support contained in $X'$. If $I$ is the ideal of $X'$ in $X$, it follows that there exist integers $m$ and $n$ such that $I^mH=0$ and $I^nK=0$ (EGA I 5.3.4). Thus one may dévissser $H$ and $K$ into extensions of modules annihilated by $I$, hence of $\mathcal O_{X'}$-modules, and
\[
 \cl(H)\in\operatorname{Im}K_\bullet(X'),\qquad \cl(K)\in\operatorname{Im}K_\bullet(X'),
\]
which proves the assertion. Thus, to every coherent $\mathcal O_U$-module, assigning the class of one of its extensions to $X$ defines a function on the category of coherent $\mathcal O_U$-modules with values in $K_\bullet(X)/\operatorname{Im}K_\bullet(X')$.

This function is additive. Indeed, let
\[
0\longrightarrow G'\longrightarrow G\longrightarrow G''\longrightarrow 0
\]
be an exact sequence of coherent $\mathcal O_U$-modules. Let $F$ be an extension of $G$ to $X$. There exists a submodule $F'$ of $F$ extending $G'$, and $F/F'$ extends $G''$. Hence the function factors through $K_\bullet(U)$, defining the desired homomorphism $K_\bullet(U)\to K_\bullet(X)/\operatorname{Im}K_\bullet(X')$. One verifies at once that it is inverse to $\rho$.

\sgaStableTarget{sga6.e09.prop.1-2}\textbf{Proposition 1.2.} Let $X$ and $S$ be noetherian schemes, and let $f:X\to S$ be a flat morphism. Let $(\xi_i)_{i\in I}$ be a family of elements of $K^\circ(X)$, and consider the homomorphism
\[
 \varphi:K_\bullet(S)^{(I)}\longrightarrow K_\bullet(X)
\]
defined by
\[
 \varphi((x_i)_{i\in I})=\sum_{i\in I}\xi_i\,f^*(x_i),
\]
where $f^*$ is the inverse-image homomorphism from $K_\bullet(S)$ to $K_\bullet(X)$ defined in \hyperlink{sga6.e04.para.2-12}{IV 2.12}, and where $K_\bullet(X)$ is considered as a $K^\circ(X)$-module by \hyperlink{sga6.e04.para.2-10}{IV 2.10}. Suppose that, for every point $s\in S$, the images of the $\xi_i$ in $K_\bullet(X_s)$ by the homomorphism
\[
 K^\circ(X)\longrightarrow K^\circ(X_s)=K_\bullet(X_s)
\]
generate $K_\bullet(X_s)$ as an abelian group. Then $\varphi$ is surjective.

Let $S'$ be a closed subscheme of $S$, and let $U$ be the open complement of $S'$ in $S$. Write $X_U$ for the restriction of $X$ over $U$, and $X_{S'}$ for the restriction of $X$ over $S'$. We again denote by $\xi_i$ the images of the $\xi_i$ in $K^\circ(X_{S'})$ and $K^\circ(X_U)$, and write
\[
 \varphi_{S'}:K_\bullet(S')^{(I)}\to K_\bullet(X_{S'}),\qquad
 \varphi_U:K_\bullet(U)^{(I)}\to K_\bullet(X_U)
\]
for the homomorphisms defined in the same way as $\varphi$.

Consider the diagram
\sgaStableTarget{sga6.e09.eq.1-2-1}\[
\begin{tikzcd}[column sep=tiny]
K_\bullet(S')^{(I)} \arrow[r,"i_*"] \arrow[d,"\varphi_{S'}"'] &
K_\bullet(S)^{(I)} \arrow[r,"j^*"] \arrow[d,"\varphi"] &
K_\bullet(U)^{(I)} \arrow[r] \arrow[d,"\varphi_U"] & 0\\
K_\bullet(X_{S'}) \arrow[r,"u_*"] & K_\bullet(X) \arrow[r,"v^*"] & K_\bullet(X_U) \arrow[r] & 0,
\end{tikzcd}\tag{1.2.1}
\]
where $i$ and $j$ denote the immersions of $S'$ and $U$ into $S$, and $u$ and $v$ the immersions of $X_{S'}$ and $X_U$ into $X$. We show that, if $\varphi_{S'}$ and $\varphi_U$ are surjective, then so is $\varphi$.

The two horizontal rows are exact by 1.1. To prove that $\varphi$ is surjective, it is therefore enough to prove that the two squares commute. For $x_i\in K_\bullet(S')$, one has
\[
 \varphi(i_*(x_i))=\sum_i \xi_i f^*i_*(x_i),
\]
whereas
\[
 u_*\varphi_{S'}((x_i))=u_*\left(\sum_i\xi_i u^*f^*(x_i)\right)=\sum_i \xi_i u_*(u^*f^*(x_i)),
\]
the last equality following from the projection formula \hyperlink{sga6.e04.eq.2-11-1-2}{IV 2.11.1.2}. The equality $\varphi i_*=u_*\varphi_{S'}$ therefore follows from $f^*i_*=u_*u^*f^*$, which is true by \hyperlink{sga6.e04.prop.3-1-1}{IV 3.1.1} since $i$ is proper and $f$ is flat. Similarly, for $(x_i)\in K_\bullet(S)^{(I)}$,
\[
 \varphi_U(j^*(x_i))=\sum_i\xi_i f_U^*j^*(x_i),
\]
and
\[
 v^*\varphi((x_i))=v^*\left(\sum_i\xi_i f^*(x_i)\right)=\sum_i\xi_i v^*f^*(x_i)=\sum_i\xi_i f_U^*j^*(x_i),
\]
the middle equality being \hyperlink{sga6.e04.eq.2-12-2}{IV 2.12.2}.

This result permits one to prove \hyperlink{sga6.e09.prop.1-2}{1.2} by noetherian induction on the set of closed subschemes of $S$. If one assumes that, for every closed subscheme $S'$ of $S$ distinct from $S$, $\varphi_{S'}$ is surjective, the preceding argument shows that one may reduce to proving \hyperlink{sga6.e09.prop.1-2}{1.2} for an irreducible affine base scheme, after possibly replacing $S$ by an irreducible affine open. Moreover, the diagram
\[
\begin{tikzcd}
K_\bullet(S_{\mathrm{red}})^{(I)} \arrow[r] \arrow[d,"\varphi_{S_{\mathrm{red}}}"'] & K_\bullet(S)^{(I)} \arrow[d,"\varphi"]\\
K_\bullet(X_{S_{\mathrm{red}}}) \arrow[r] & K_\bullet(X)
\end{tikzcd}
\]
where $K_\bullet(X_{S_{\mathrm{red}}})\to K_\bullet(X)$ is surjective by \hyperlink{sga6.e09.prop.1-1}{1.1}, shows that one may also assume $S$ reduced.

We are thus on an integral affine base scheme. Put $S=\Spec(A)$, and let $K$ denote the fraction field of $A$. If $s$ is the generic point of $S$, then $\mathcal O_{S,s}=K$. Let $y\in K_\bullet(X)$. By hypothesis, $K_\bullet(X_s)$ is generated by the images of the $\xi_i$; if $y_s$ is the image of $y$ in $K_\bullet(X_s)$, then one can find integers $n_i$ such that
\[
 y_s=\sum_i \xi_i n_i.
\]
On the other hand $K=\varinjlim_f A_f$, where $f$ runs through the elements of $A$. Hence $\Spec(K)=\varprojlim_f \Spec(A_f)$ and
\[
 X_s=X\times_S\Spec(K)=\varprojlim_f X\times_S\Spec(A_f)
\]
by EGA IV 8.2.3 and 8.2.5. Since the schemes considered are noetherian and the transition morphisms affine and flat, one has, by \hyperlink{sga6.e04.para.3-2-3}{IV 3.2.3},
\[
 K_\bullet(X_s)=\varinjlim_f K_\bullet(X\times_S\Spec(A_f)).
\]
As inverse-image homomorphisms on the $K_\bullet$ are linear with respect to $K^\circ$ (\hyperlink{sga6.e04.eq.2-12-2}{IV 2.12.2}), there exists a non-empty open $U=\Spec(A_f)$ of $S$ such that in $K_\bullet(X_U)$ one has
\[
 y_U=\sum_i \xi_i n_i.
\]
Let $S'$ be a closed subscheme of $S$ with support $S-U$. Then $\varphi_{S'}$ is surjective by the induction hypothesis, and consequently diagram \hyperlink{sga6.e09.eq.1-2-1}{(1.2.1)} shows that
\[
 y-
\sum_i\xi_i n_i\in \operatorname{Im}(\varphi).
\]
Since $n_i=f^*(n_i)$, $\varphi$ is surjective.

\sgaStableTarget{sga6.e09.cor.1-3}\textbf{Corollary 1.3.} Let $f:X\to S$ be a flat morphism of noetherian schemes. Suppose that, for every point $s$ of $S$, the homomorphism
\[
 K^\circ(X)\longrightarrow K_\bullet(X_s)
\]
is surjective. Then the homomorphism deduced from $f^*$,
\[
  K_\bullet(S)\otimes_{K^\circ(S)}K^\circ(X)\longrightarrow K_\bullet(X),
\]
is surjective.

Take for the family $(\xi_i)$ the family of all elements of $K^\circ(X)$. The hypotheses of \hyperlink{sga6.e09.prop.1-2}{1.2} are then satisfied, and hence
\[
 \varphi:K_\bullet(S)^{(K^\circ(X))}\longrightarrow K_\bullet(X)
\]
is surjective. By definition of $\varphi$, this is equivalent to the surjectivity of the displayed tensor-product homomorphism.

\sgaStableTarget{sga6.e09.cor.1-4}\textbf{Corollary 1.4.} Let $f:X\to S$ be a flat morphism of noetherian schemes. Suppose that, for every point $s$ of $S$, the canonical homomorphism $\mathbb Z\to K_\bullet(X_s)$ is an isomorphism. Then
\[
 f^*:K_\bullet(S)\longrightarrow K_\bullet(X)
\]
is surjective. If, moreover, $f$ has a section $g$ of finite Tor-dimension, then $f^*$ is an isomorphism.

Take for the family $(\xi_i)$ the family reduced to the single element $1\in K^\circ(X)$. Then $\varphi=f^*$, and the conditions of \hyperlink{sga6.e09.prop.1-2}{1.2} being fulfilled, $f^*$ is surjective. If $f$ admits a section $g$ of finite Tor-dimension, one may define a homomorphism
\[
 g^*:K_\bullet(X)\longrightarrow K_\bullet(S)
\]
(\hyperlink{sga6.e04.para.2-12}{IV 2.12}). Since $g^*f^*=(fg)^*=\operatorname{Id}$ by \hyperlink{sga6.e04.eq.2-12-1-1}{IV 2.12.1.1}, $f^*$ is injective, hence an isomorphism.

\sgaStableTarget{sga6.e09.rem.1-5}\textbf{Remark 1.5.} For $g$ to be of finite Tor-dimension, it is necessary and sufficient that the fibres of $X$ be regular at the points of $g(S)$, i.e. that $f$ be regular at the points of $g(S)$, or smooth at these points if one assumes $f$ of finite presentation.

These properties being local on $X$ and $S$, one may suppose them affine: $S=\Spec(A)$, $X=\Spec(B)$. Then
\[
 \operatorname{tor.dim}_B(A)=\operatorname{dim.proj}_B(A)=\sup_{s\in S}\operatorname{dim.proj}_{B\otimes_A k(s)}(k(s)),
\]
by EGA IV 12.3.2, where the finite-presentation hypothesis on $B$ may be replaced by noetherian hypotheses. But $k(s)$ is isomorphic to the residue field of $\mathcal O_{X_s,g(s)}$, and by EGA $0_{\mathrm{IV}}$ 17.2.5 one has
\[
 \operatorname{dim.proj}_{B\otimes_A k(s)}(k(s))=
 \operatorname{dim.proj}_{\mathcal O_{X_s,g(s)}}(k(s)).
\]
It follows from EGA $0_{\mathrm{IV}}$ 17.2.11 that the latter is finite if and only if $\mathcal O_{X_s,g(s)}$ has finite cohomological dimension, i.e. is regular.

\sgaStableTarget{sga6.e09.prop.1-6}\textbf{Proposition 1.6.} Let $S$ be a noetherian scheme, let $E$ be a locally free $\mathcal O_S$-module of finite type, let $X$ be the vector bundle associated with $E$, and let $f$ be the structural morphism of $X$. Then
\[
 f^*:K_\bullet(S)\longrightarrow K_\bullet(X)
\]
is an isomorphism.

Since $f$ is smooth and admits a section, namely the zero section, it is enough, to prove \hyperlink{sga6.e09.prop.1-6}{1.6}, to show that for every $s\in S$ one has $K_\bullet(X_s)\simeq \mathbb Z$ \hyperlink{sga6.e09.cor.1-4}{(1.4)}. Thus one is reduced to the case where the base scheme is a field. In this case $E$ is a direct sum of rank-one submodules, which permits one to decompose $f$ into a sequence of rank-one fibrations, and hence to reduce to the case where $X$ has rank one, namely $X=\Spec(k[T])$. The proposition then follows from the next lemma.

\sgaStableTarget{sga6.e09.lem.1-7}\textbf{Lemma 1.7.} Let $A$ be a principal ideal domain and let $X=\Spec(A)$. Then the canonical homomorphism
\[
 \mathbb Z\longrightarrow K_\bullet(X)
\]
is an isomorphism.

Since $X$ is regular, the homomorphism $K^\circ(X)\to K_\bullet(X)$ is an isomorphism. Hence $K_\bullet(X)$ is generated by the classes of locally free $\mathcal O_X$-modules of finite type, i.e. by the classes of projective $A$-modules of finite type. Since $A$ is principal, every projective module of finite type is free, which gives surjectivity. Injectivity follows from the fact that $K^\circ(X)$ is $\mathbb Z$-augmented.

\sgaStableTarget{sga6.e09.prop.1-8}\textbf{Proposition 1.8.} Let $S$ be a noetherian scheme, let $E$ be a locally free $\mathcal O_S$-module of rank $r$, let $X$ be the vector bundle associated with $E$, let $f$ be the structural morphism of $X$, and let $g$ be the zero section of $X$. Then the composite homomorphism
\[
 K_\bullet(S)\xrightarrow{g_*}K_\bullet(X)\xrightarrow{g^*}K_\bullet(S)
\]
where $g^*$ is defined by \hyperlink{sga6.e04.eq.2-12-1}{IV 2.12.1} and \hyperlink{sga6.e04.para.1-5}{1.5}, is multiplication by the element $\lambda_{-1}(E)$ of $K^\circ(S)$, writing $E$ for the class of $E$ in $K^\circ(S)$ and putting
\[
 \lambda_{-1}(E)=\sum_{i=0}^r(-1)^i\lambda^i(E).
\]

This statement is the analogue, for $K_\bullet$, of \hyperlink{sga6.e07.para.2-6}{VII 2.7}. Let $F$ be a coherent $\mathcal O_S$-module, of class $F$ in $K_\bullet(S)$. Then $g_*(F)$ is the class in $K_\bullet(X)$ of the complex $Rg_*(F)=g_*(F)$, hence of $F$ considered as an $\mathcal O_X$-module through $g$. To compute $g^*(g_*(F))$, one must take the class of
\[
 Lg^*(F)=F\otimes^L_{\mathcal O_X}\mathcal O_S.
\]
For this, consider the Koszul complex
\[
0\longrightarrow \Lambda^r(E)\otimes_{\mathcal O_S}S(E)\longrightarrow
\Lambda^{r-1}(E)\otimes_{\mathcal O_S}S(E)\longrightarrow \cdots\longrightarrow
E\otimes_{\mathcal O_S}S(E)\longrightarrow S(E)\longrightarrow 0,
\]
where $S(E)$ is the symmetric algebra of $E$. This is a resolution of $\mathcal O_S$ over $S(E)$ (\hyperlink{sga6.e06.sec.1.10}{VI 1.11}). The associated complex on $X$, denoted $L_\bullet$, has $k$-th term $\Lambda^k(f^*E)$ and is therefore a resolution of $\mathcal O_S$ by locally free $\mathcal O_X$-modules of finite type. Thus
\[
 L_\bullet\otimes_{\mathcal O_X}F=\mathcal O_S\otimes_{\mathcal O_X}F,
\]
and hence
\[
\begin{aligned}
 g^*(g_*(F))
 &=\cl(L_\bullet\otimes_{\mathcal O_X}F)
   =\sum_{i=0}^r(-1)^i\cl(L_i\otimes_{\mathcal O_X}F)\\
 &=\sum_{i=0}^r(-1)^i\cl(\Lambda^i(E)\otimes_{\mathcal O_S}F)
   =F\lambda_{-1}(E).
\end{aligned}
\]
The proposition follows by additivity.

\sgaStableTarget{sga6.e09.sec.2}\subsection*{2. Principal bundles under split tori}

\sgaStableTarget{sga6.e09.para.2-1}\textbf{2.1.} Let $S$ be a noetherian scheme, and let $T$ be a torus over $S$. Suppose $T$ is split; thus one has an isomorphism $T\simeq D_S(M)$, where $M$ is a free abelian group of rank $r$, or equivalently $T\simeq \mathbb G_{m,S}^r$. Let $P$ be a torsor under $T$. Then $P$ is affine over $S$, defined by a graded $\mathcal O_S$-algebra $A$ of type $M$. If
\[
 A=\bigoplus_{m\in M}A_m,
\]
then $A_m$ is an invertible $\mathcal O_S$-module for every $m$, and, for $m,m'\in M$, the homomorphism
\[
 A_m\otimes_{\mathcal O_S}A_{m'}\longrightarrow A_{m+m'}
\]
defined by the multiplication of $A$ is an isomorphism (SGA 3 VIII 4.1). One therefore defines a multiplicative homomorphism
\[
 \eta:M\longrightarrow K^\circ(S)
\]
by assigning to $m\in M$ the class of $A_m$ in $K^\circ(S)$.

\sgaStableTarget{sga6.e09.prop.2-2}\textbf{Proposition 2.2.} With the notation and hypotheses of \hyperlink{sga6.e09.para.2-1}{2.1}, let $f:P\to S$ be the structural morphism of $P$. Then the homomorphism
\[
 f^*:K_\bullet(S)\longrightarrow K_\bullet(P)
\]
is surjective, and its kernel is the subgroup of $K_\bullet(S)$ generated by the subgroups $(\eta(m)-1)K_\bullet(S)$ for $m\in M$, or, equivalently, for $m$ running through a system of generators of $M$.

First suppose $r=1$, i.e. $T=\mathbb G_{m,S}$. Then $M=\mathbb Z$, and for every $m\in M$ one has
\[
 A_m\simeq A_1^{\otimes m}.
\]
If $A_1=L$, then $A=\bigoplus_{n\in\mathbb Z}L^{\otimes n}$, which shows that $P$ identifies with the complement of the zero section in the bundle $X=\mathbb V(L)$. Considering $S$ as a closed subscheme of $X$ by the zero section, one has, by \hyperlink{sga6.e09.prop.1-1}{1.1}, the exact sequence
\[
 K_\bullet(S)\xrightarrow{g_*}K_\bullet(X)\longrightarrow K_\bullet(P)\longrightarrow 0,
\]
where $g:S\to X$ is the zero section of $X$. Moreover $f^*$ factors as
\[
 K_\bullet(S)\xrightarrow{f'^*}K_\bullet(X)\longrightarrow K_\bullet(P),
\]
where $f'$ is the structural morphism of $X$. By \hyperlink{sga6.e09.prop.1-6}{1.6}, $f'^*$ is an isomorphism, inverse to $g^*$. Thus we have the exact sequence
\[
 K_\bullet(S)\xrightarrow{g^*g_*}K_\bullet(S)\xrightarrow{f^*}K_\bullet(P)\longrightarrow 0,
\]
and by \hyperlink{sga6.e09.prop.1-8}{1.8} one obtains
\[
 K_\bullet(P)\simeq K_\bullet(S)/\lambda_{-1}(L)K_\bullet(S)=K_\bullet(S)/(1-L)K_\bullet(S),
\]
where $L$ is the class of $L$ in $K^\circ(S)$. Since, for every $n\in\mathbb Z$, $L^n-1$ is a multiple of $1-L$, this proves \hyperlink{sga6.e09.prop.2-2}{2.2} when $P$ has rank one.

Continue the proof by induction on the rank $r$ of $P$. Suppose the result proved for torsors under split tori of rank $<r$. Write $M=M_1\oplus M_2$, with $M_1\simeq\mathbb Z^{r-1}$ and $M_2\simeq\mathbb Z$. Put $A'=\bigoplus_{m\in M_1}A_m$ and $A''=\bigoplus_{m\in M_2}A_m$, so that $A=A'\otimes_{\mathcal O_S}A''$.
Let $P_1=\Spec(A')$. Then $P_1$ is a torsor under $D_S(M_1)$, and $P$ is, over $P_1$, a torsor under the group $\mathbb G_{m,P_1}$ with algebra the inverse image of $A''$ on $P_1$. By the preceding rank-one case, $K_\bullet(P)$ is the quotient of $K_\bullet(P_1)$ by the subgroup generated by the $(\eta(m_2)-1)K_\bullet(P_1)$ for $m_2\in M_2$. By the induction hypothesis, $K_\bullet(P_1)$ is the quotient of $K_\bullet(S)$ by the subgroup generated by the $(\eta(m_1)-1)K_\bullet(S)$ for $m_1\in M_1$. Taking account of the relation
\[
 \eta(m_1+m_2)-1=\eta(m_1)\eta(m_2)-1
 =\eta(m_1)(\eta(m_2)-1)+\eta(m_1)-1,
\]
we conclude that $K_\bullet(P)$ is the quotient of $K_\bullet(S)$ by the subgroup generated by the $(\eta(m)-1)K_\bullet(S)$ for $m\in M$, or equivalently for $m$ running through a system of generators of $M$.

\sgaStableTarget{sga6.e09.cor.2-3}\textbf{Corollary 2.3.} With the notation and hypotheses of \hyperlink{sga6.e09.para.2-1}{2.1}, suppose $S$ regular. Then one has a canonical isomorphism
\[
 K^\circ(P)\simeq K^\circ(S)/\bigl((\eta(m)-1)_{m\in M}\bigr),
\]
where $((\eta(m)-1)_{m\in M})$ is the ideal of $K^\circ(S)$ generated by the $\eta(m)-1$, for $m$ running through a system of generators of $M$.

If $S$ is regular, the canonical homomorphism $K^\circ(S)\to K_\bullet(S)$ is an isomorphism (\hyperlink{sga6.e04.para.2-5}{IV 2.5}). Likewise, $P$ is then regular, and $K^\circ(P)\to K_\bullet(P)$ is an isomorphism. The corollary therefore follows immediately from 2.2.

\sgaStableTarget{sga6.e09.sec.3}\subsection*{3. Projective bundles and flag bundles}

\sgaStableTarget{sga6.e09.prop.3-1}\textbf{Proposition 3.1.} Let $S$ be a noetherian scheme, let $E$ be a locally free $\mathcal O_S$-module of rank $n$, let $\pi=(p_1,\ldots,p_k)$ be a sequence of integers such that $\sum p_i=n$, and let $X=D_\pi(E)$ be the scheme of flags of type $\pi$ in $E$. Then the canonical homomorphism
\[
 K_\bullet(S)\otimes_{K^\circ(S)}K^\circ(X)\longrightarrow K_\bullet(X)
\]
is an isomorphism.

We resume the notation of \hyperlink{sga6.e06.sec.4.5}{VI 4.5}. Thus $L_i$ denotes the canonical quotients of the inverse image of $E$ on $X$, $F_j$ denotes the kernel of $L_j\to L_{j-1}$, and $\lambda_i^{(j)}$ denotes the class in $K^\circ(X)$ of $\Lambda^i(F_j)$. Then the $\lambda_i^{(j)}$ generate $K^\circ(X)$ as a $K^\circ(S)$-algebra (\hyperlink{sga6.e06.thm.4-6}{VI 4.6}). Likewise, the images of the $\lambda_i^{(j)}$ in $K^\circ(X_s)$ generate $K^\circ(X_s)$ as an algebra over $K^\circ(k(s))=\mathbb Z$ for every point $s$ of $S$. Since $X_s$ is a flag scheme over a field, it is regular, and hence $K^\circ(X_s)\to K_\bullet(X_s)$ is an isomorphism. It follows that $K^\circ(X)\to K_\bullet(X_s)$ is surjective. By \hyperlink{sga6.e09.cor.1-3}{1.3}, this implies that
\[
 K_\bullet(S)\otimes_{K^\circ(S)}K^\circ(X)\longrightarrow K_\bullet(X)
\]
with $K^\circ(X)$ in the second factor is surjective.

To prove injectivity, one may first reduce to the case of the complete flag bundle $D(E)$. Indeed, suppose that
\[
 K_\bullet(S)\otimes_{K^\circ(S)}K^\circ(D(E))\longrightarrow K_\bullet(D(E))
\]
is injective. Since $K_\bullet(S)\to K_\bullet(D(E))$ factors through $K_\bullet(X)$, the corresponding homomorphism factors as
\[
 K_\bullet(S)\otimes_{K^\circ(S)}K^\circ(X)\otimes_{K^\circ(X)}K^\circ(D(E))
 \longrightarrow
 K_\bullet(X)\otimes_{K^\circ(X)}K^\circ(D(E))
 \longrightarrow K_\bullet(D(E)).
\]
The first homomorphism is injective. Since, by \hyperlink{sga6.e06.lem.4-3}{VI 4.3}, \hyperlink{sga6.e06.thm.4-6}{4.6} and \hyperlink{sga6.e06.cor.4-7}{4.7}, $K^\circ(D(E))$ is free over $K^\circ(X)$, this implies injectivity for
\[
 K_\bullet(S)\otimes_{K^\circ(S)}K^\circ(X)\to K_\bullet(X).
\]
Thus assume $X=D(E)$. One proceeds by induction on the rank of $E$. The case $n=1$ is trivial. Suppose the property proved for $n-1$. Introduce $P=\mathbb P(E)$, and let $F$ be the kernel of the canonical homomorphism $p^*(E)\to\mathcal O_P(1)$, where $p$ is the structural morphism of $P$. Then $X$ identifies with the scheme of flags of $F$ over $P$, and by the induction hypothesis it is enough to prove \hyperlink{sga6.e09.prop.3-1}{3.1} for $P$. Taking \hyperlink{sga6.e06.thm.1-1}{VI 1.1} into account, the result to prove may therefore be stated as the following corollary.

\sgaStableTarget{sga6.e09.cor.3-2}\textbf{Corollary 3.2.} Let $S$ be a noetherian scheme, let $E$ be a locally free $\mathcal O_S$-module of rank $n$, let $X=\mathbb P(E)$ be the associated projective bundle, let $f$ be the structural morphism of $X$, and let $\xi$ be the class in $K^\circ(X)$ of the sheaf $\mathcal O_X(1)$. Then the homomorphism
\[
 \varphi:(K_\bullet(S))^n\longrightarrow K_\bullet(X)
\]
defined by
\[
 \varphi((x_i))=\sum_{i=0}^{n-1}\xi^{-i}f^*(x_i)
\]
is an isomorphism.

We have already seen that this homomorphism is surjective, since the elements $\xi^i$, $i=0,\ldots,-n+1$, form a basis of $K^\circ(X)$ over $K^\circ(S)$. To see that it is injective, use the homomorphism $f_*$, defined because $f$ is projective. By \hyperlink{sga6.e04.eq.2-12-4}{IV 2.12.4},
\[
 f_*\left(\sum_{i=0}^{n-1}\xi^{-i}f^*(x_i)\right)
 =\sum_{i=0}^{n-1}x_i f_*(\xi^{-i})=x_0,
\]
the last equality following from the fact that $f_*(\xi^i)=0$ for $-n+1\leq i<0$ (\hyperlink{sga6.e06.sec.2.8}{VI 2.8}). If $\varphi((x_i))=0$, what precedes shows that $x_0=0$. Applying the same reasoning to the elements $\xi^i\varphi((x_i))$, for $i=1,\ldots,n-1$, one obtains by induction that all the $x_i$ are zero. Thus $\varphi$ is injective.

\sgaStableTarget{sga6.e09.para.3-3}\textbf{3.3.} Let $S$ still be a noetherian scheme, and let $E$ be a locally free $\mathcal O_S$-module of rank $n$. Write $V=\mathbb V(E)$ for the vector bundle associated with $E$, contravariantly in $E$ (SGA 3 I 4.6), considered as a group scheme over $S$. Let $P$ be a torsor under $V$; we wish to compute $K_\bullet(P)$.

Recall how one embeds $P$ in a projective bundle over $S$. Torsors under $V$ correspond to extensions $F$ of $E$ by $\mathcal O_S$. Indeed, if one has an extension
\sgaStableTarget{sga6.e09.eq.3-3-1}\[
0\longrightarrow\mathcal O_S\xrightarrow{s}F\xrightarrow{u}E\longrightarrow 0,
\tag{3.3.1}
\]
one obtains a torsor under $V$ by considering the homomorphism $\mathbb V(F)\to\mathbb V(\mathcal O_S)$ defined by $s$, the section of $\mathbb V(\mathcal O_S)=\Spec(\mathcal O_S[T])$ defined by $T=1$, and taking
\[
 P=\mathbb V(F)\times_{\mathbb V(\mathcal O_S)}S,
\]
the operations of $V$ on $P$ being induced by those of $V$ on $\mathbb V(F)$.

Conversely, let $P$ be a torsor under $V$. Denote by $\underline P$ the sheaf of germs of homomorphisms from $S$ to $P$. It is a principal homogeneous sheaf under the sheaf of germs of homomorphisms from $S$ to $V$, which is identified with the dual $\check E$ of $E$. In the classical way this yields an extension of $\mathcal O_S$ by $\check E$, and by passing to duals one obtains \hyperlink{sga6.e09.eq.3-3-1}{(3.3.1)}.

Thus let $P$ be a torsor under $V$, corresponding to an extension $F$ of $E$ by $\mathcal O_S$. The algebra $A(P)$ of $P$ over $\mathcal O_S$ is
\[
 S(F)\otimes_{S(\mathcal O_S)}\mathcal O_S,
\]
that is,
\[
 S(F)/(s-1)S(F),
\]
which identifies with the ring of homogeneous fractions of degree $0$, $S(F)_{(s)}$. Consequently $P$ identifies with the open $D_+(s)$ of $\mathbb P(F)$. The complement of $P$ in $\mathbb P(F)$ is then $\operatorname{Proj}(S(F)/sS(F))$, which is identified, by \hyperlink{sga6.e09.eq.3-3-1}{(3.3.1)}, with $\mathbb P(E)$. Thus one has immersions
\[
\begin{tikzcd}[column sep=large,row sep=small]
\mathbb P(E) \arrow[r,hook] \arrow[dr] & \mathbb P(F) & P \arrow[l,hook'] \arrow[dl]\\
& S &
\end{tikzcd}
\]
where $\mathbb P(E)$ is closed in $\mathbb P(F)$, and $P$ is open and complementary to $\mathbb P(E)$.

\sgaStableTarget{sga6.e09.prop.3-4}\textbf{Proposition 3.4.} Let $S$ be a noetherian scheme, let $E$ be a locally free $\mathcal O_S$-module of rank $n$, let $V=\mathbb V(E)$ be the associated vector bundle, let $P$ be a torsor under $V$, and let $f$ be the structural morphism of $P$. Then the canonical homomorphism
\[
 f^*:K_\bullet(S)\longrightarrow K_\bullet(P)
\]
is an isomorphism.

With the identifications of \hyperlink{sga6.e09.para.3-3}{3.3}, one has the exact sequence
\sgaStableTarget{sga6.e09.eq.3-4-1}\[
 K_\bullet(\mathbb P(E))\xrightarrow{i_*}K_\bullet(\mathbb P(F))\longrightarrow K_\bullet(P)\longrightarrow 0,
\tag{3.4.1}
\]
where $i$ denotes the closed immersion $\mathbb P(E)\to\mathbb P(F)$ defined by $u$. By \hyperlink{sga6.e09.cor.3-2}{3.2},
\[
 K_\bullet(\mathbb P(E))=K_\bullet(S)\otimes_{K^\circ(S)}K^\circ(\mathbb P(E)),
\qquad
 K_\bullet(\mathbb P(F))=K_\bullet(S)\otimes_{K^\circ(S)}K^\circ(\mathbb P(F)).
\]
Since $\mathbb P(E)$ and $\mathbb P(F)$ are smooth, $i$ is a regular immersion, and hence we have a homomorphism
\[
 i_*:K^\circ(\mathbb P(E))\longrightarrow K^\circ(\mathbb P(F)).
\]
The homomorphism between the $K_\bullet$ groups is then obtained by tensoring the homomorphism between the $K^\circ$ groups by $K_\bullet(S)$. It is enough, to see this, to show that for $a\in K_\bullet(S)$ and $b\in K^\circ(\mathbb P(E))$ one has
\[
 i_*(b\,h^*(a))=i_*(b)g^*(a),
\]
where $g$ and $h$ are the structural morphisms of $\mathbb P(F)$ and $\mathbb P(E)$. Since $h=gi$,
\[
 i_*(b\,h^*(a))=i_*(b\,i^*g^*(a))=g^*(a)i_*(b)
\]
by \hyperlink{sga6.e04.eq.2-12-4}{IV 2.12.4}.

From \hyperlink{sga6.e09.eq.3-4-1}{(3.4.1)} one therefore gets
\[
 K_\bullet(P)\simeq K_\bullet(S)\otimes_{K^\circ(S)}
 \bigl(K^\circ(\mathbb P(F))/\operatorname{Im}K^\circ(\mathbb P(E))\bigr).
\]
Let $\xi_F$ and $\xi_E$ be the classes of the canonical invertible sheaves on $\mathbb P(F)$ and $\mathbb P(E)$. By \hyperlink{sga6.e06.thm.1-1}{VI 1.1}, $K^\circ(\mathbb P(F))$ (respectively $K^\circ(\mathbb P(E))$) has as basis over $K^\circ(S)$ the elements $\xi_F^k$ for $k=0,\ldots,-n$ (respectively $\xi_E^k$ for $k=0,\ldots,-n+1$). The projection formula gives
\[
 i_*(\xi_E^k)=i_*(i^*(\xi_F^k))=\xi_F^k i_*(1).
\]
If $I$ is the ideal of $\mathbb P(E)$ in $\mathbb P(F)$, then $I$ is invertible and $i_*(1)=1-I$. Moreover $I$ is defined by the ideal $J$ of $S(F)$, kernel of $S(F)\to S(E)$, which is generated by a submodule of $F$ isomorphic to $\mathcal O_S$. Hence
\[
 I\simeq \mathcal O_{\mathbb P(F)}(-1).
\]
The image of $K^\circ(\mathbb P(E))$ is thus the $K^\circ(S)$-submodule generated by the $\xi_F^k-\xi_F^{k-1}$, for $k=0,\ldots,-n+1$. Therefore
\[
 K^\circ(S)\longrightarrow K^\circ(\mathbb P(F))/\operatorname{Im}K^\circ(\mathbb P(E))
\]
is an isomorphism, and consequently so is $K_\bullet(S)\to K_\bullet(P)$.

\sgaStableTarget{sga6.e09.sec.4}\subsection*{4. Principal bundles under the groups $\mathrm{GL}(n)_S$}

\sgaStableTarget{sga6.e09.prop.4-1}\textbf{Proposition 4.1.} Let $S$ be a noetherian scheme, let $G=\mathrm{GL}(n)_S$, let $P$ be a torsor under $G$, and let $E$ be the locally free $\mathcal O_S$-module of rank $n$ corresponding to $P$. Denote by $E$ the class of $E$ in $K^\circ(S)$, and let $J$ be the $\lambda$-ideal generated by $E-n$. If $f$ is the structural morphism of $P$, then $f^*$, defined because $f$ is flat, is surjective and defines an isomorphism
\[
 K_\bullet(S)/J K_\bullet(S)\xrightarrow{\sim}K_\bullet(P).
\]

Let $E_0=(\mathcal O_S)^n$, with evident basis $e_1,\ldots,e_n$, and let $E_0^{(i)}$ be the submodule of $E_0$ generated by $e_1,\ldots,e_i$. The subfunctor of $G$ which to every $S$-scheme $S'$ assigns the subgroup of automorphisms of $E_{0,S'}$ leaving invariant the submodules $E_{0,S'}^{(i)}$ for all $i$ is represented by a group subscheme $B$ of $G$, a Borel subgroup of $G$. The quotient sheaf $G/B$ for the fpqc topology is represented by the flag scheme of $E_0$, denoted $X_0$. Let $L_0^{(i)}$ be the canonical quotients of the inverse image of $E_0$ on $X_0$, let $G_0^{(i)}$ be the kernel of $E_{0,X_0}\to L_0^{(i)}$, and put
\[
 F_0^{(i)}=G_0^{(i-1)}/G_0^{(i)}\simeq \ker(L_0^{(i)}\to L_0^{(i-1)}).
\]

% SGA-SRC-001466
Moreover, $B$ is the semi-direct product of its unipotent radical $U$ and of its maximal subtorus $T$ leaving the $e_i$ invariant. Let $Y_0=G/U$. Then $Y_0$ is a torsor over $X_0$ under the group $T_{X_0}$. If $S'$ is an $S$-scheme, the datum of an $S'$-point of $Y_0$ is equivalent to the datum, on some faithfully flat quasi-compact $S''$ over $S'$, of an automorphism of $E_{0,S''}$ modulo the unipotents of $U_{S''}$; that is, of a flag of $E_{0,S''}$ with free submodules, and, for each $i$, of a class of free elements $f_i^\alpha$ on $\mathcal O_{S''}$ such that $f_i^\alpha-f_i^\beta\in G_{0,S''}^{(i)}$. Equivalently, it is the datum of a flag of $E_{0,S''}$ and, for every $i$, of a basis of $F_{0,S''}^{(i)}$, this datum being subject to the gluing condition on $S''\times_{S'}S''$. By fpqc descent, this amounts to the same datum on $S'$. Thus $Y_0$ is the affine $X_0$-scheme defined by the algebra
\[
 A_0=\bigoplus_{k_i\in\mathbb Z}(F_0^{(1)})^{\otimes k_1}\otimes\cdots\otimes(F_0^{(n)})^{\otimes k_n}.
\]

The group $G$ operates on $E_0$, on $X_0$ provided with the sheaves $L_0^{(i)}$, and on $Y_0$. Taking the product of these objects with $P$ modulo $G$, one obtains the sheaf $E$ associated with $P$, the scheme $X=D(E)=P/B$ of flags of $E$, and the scheme $Y=P/U$. We denote the canonical sheaves on $X$ by the same letters as their analogues on $X_0$. Thus one has the diagram
\[
\begin{array}{ccccc}
&& Y=P/U &&\\[-0.2em]
& \nearrow && \searrow &\\[-0.2em]
P &&&& X=P/B=D(E)\\[-0.2em]
& \searrow && \swarrow &\\[-0.2em]
&& S &&
\end{array}
\]
We note that $P$ is a torsor over $Y$ with structural group $U_Y$, and that $Y$ is a torsor over $X$ with structural group $T_X$, defined by the $\mathcal O_X$-algebra
\[
 A=\bigoplus_{k_i\in\mathbb Z}(F^{(1)})^{\otimes k_1}\otimes\cdots\otimes(F^{(n)})^{\otimes k_n}.
\]

The homomorphism $K_\bullet(S)\to K_\bullet(P)$ therefore factors as
\[
 K_\bullet(S)\longrightarrow K_\bullet(X)\longrightarrow K_\bullet(Y)\longrightarrow K_\bullet(P).
\]
By \hyperlink{sga6.e09.prop.3-1}{3.1}, $K_\bullet(X)=K_\bullet(S)\otimes_{K^\circ(S)}K^\circ(X)$. Let $\xi_i$ be the class of $F^{(i)}$ in $K^\circ(X)$. By \hyperlink{sga6.e09.prop.2-2}{2.2}, $K_\bullet(Y)$ is the quotient of $K_\bullet(X)$ by the submodule generated by the $(\xi_i-1)K_\bullet(X)$. If $I$ is the ideal of $K^\circ(X)$ generated by the $\xi_i-1$, then
\[
 K_\bullet(Y)=K_\bullet(S)\otimes_{K^\circ(S)}(K^\circ(X)/I).
\]
By \hyperlink{sga6.e06.cor.4-8}{VI 4.8}, $K^\circ(X)$ is the $K^\circ(S)$-algebra generated by the $\xi_i-1$, subject to the relations deduced from
\[
 \gamma_t(E-n)=\prod_{i=1}^n\gamma_t(\xi_i-1).
\]
Therefore $K^\circ(X)/I=K^\circ(S)/J$, where $J$ is the ideal generated by the $\gamma^j(E-n)$ for $j=1,\ldots,n$, hence the $\lambda$-ideal generated by $E-n$. Thus
\[
 K_\bullet(Y)\simeq K_\bullet(S)/J K_\bullet(S).
\]
The demonstration is completed by the following lemma.

\sgaStableTarget{sga6.e09.lem.4-2}\textbf{Lemma 4.2.} Let $S$ be a noetherian scheme, let $B$ be a parabolic subgroup of a reductive $S$-group, and let $P$ be a torsor under the unipotent radical $U$ of $B$. Then one has an isomorphism
\[
 K_\bullet(S)\xrightarrow{\sim}K_\bullet(P).
\]

Indeed, there exists a finite sequence of subgroup schemes of $U$,
\[
 U_0=U\supset U_1\supset U_2\supset\cdots\supset U_p=e,
\]
closed in $B$, finite because $B$ is noetherian, and such that, for every $i$, $U_i/U_{i+1}$ is isomorphic to an $S$-vector group (SGA 3 XXVI 2.1). Consequently $P\to S$ factors as a sequence of torsors under vector groups, and the result follows from 3.4.

\sgaStableTarget{sga6.e09.cor.4-3}\textbf{Corollary 4.3.} Let $S$ be a noetherian scheme and let $G=\mathrm{GL}(n)_S$. Then the canonical homomorphism
\[
 K_\bullet(S)\longrightarrow K_\bullet(G)
\]
is an isomorphism.

This follows immediately from \hyperlink{sga6.e09.prop.4-1}{4.1}, applied to $P=G$, since then $E=\mathcal O_S^n$.

\sgaStableTarget{sga6.e09.cor.4-4}\textbf{Corollary 4.4.} Under the hypotheses of \hyperlink{sga6.e09.prop.4-1}{4.1}, suppose $S$ regular. Then one has an isomorphism
\[
 K^\circ(S)/J\xrightarrow{\sim}K^\circ(P).
\]

If $S$ is regular, $G$ is also regular, and $P$ as well. Thus $K^\circ(S)\simeq K_\bullet(S)$ and $K^\circ(P)\simeq K_\bullet(P)$, whence the result follows from 4.1.

% END INCLUDED SGA6_English_current_buildable_source_20260823/SGA6_sourcePDF001_525_English_Inherited_PartiallySourceSynchronized_fragment.tex


% Exact structural seam: Exposé IX Corollary 4.4 -> Exposé X / idx532 / sourcePDF526.
\clearpage
% BEGIN INCLUDED SGA6_English_current_buildable_source_20260823/SGA6_tail_macro_state.tex
% Exact macro state used by the standalone idx532--702 tail master.
% This file is loaded only at the sourcePDF525/526 seam.  Defining these
% commands in the global preamble would pre-empt Exposés I--IX local
% \providecommand choices and silently change the inherited prefix.

\providecommand{\Filt}{}\renewcommand{\Filt}{\operatorname{Filt}}
\providecommand{\Coh}{}\renewcommand{\Coh}{\operatorname{Coh}}
\providecommand{\Supp}{}\renewcommand{\Supp}{\operatorname{Supp}}
\providecommand{\longg}{}\renewcommand{\longg}{\operatorname{length}}
\providecommand{\cl}{}\renewcommand{\cl}{\operatorname{cl}}
\providecommand{\Ker}{}\renewcommand{\Ker}{\operatorname{Ker}}
\providecommand{\codim}{}\renewcommand{\codim}{\operatorname{codim}}
\providecommand{\degtr}{}\renewcommand{\degtr}{\operatorname{deg.tr}}
\providecommand{\Pic}{}\renewcommand{\Pic}{\operatorname{Pic}}
\providecommand{\Spec}{}\renewcommand{\Spec}{\operatorname{Spec}}
\providecommand{\Gr}{}\renewcommand{\Gr}{\operatorname{Gr}}
\providecommand{\Gro}{}\renewcommand{\Gro}{\operatorname{Gr}^{\circ}}
\providecommand{\Grb}{}\renewcommand{\Grb}{\operatorname{Gr}_{\bullet}}
\providecommand{\fpqc}{}\renewcommand{\fpqc}{\mathrm{fpqc}}
\providecommand{\Sch}{}\renewcommand{\Sch}{\operatorname{Sch}}
\providecommand{\Ob}{}\renewcommand{\Ob}{\operatorname{Ob}}
\providecommand{\Tor}{}\renewcommand{\Tor}{\operatorname{Tor}}
\providecommand{\red}{}\renewcommand{\red}{\mathrm{red}}
\providecommand{\Mod}{}\renewcommand{\Mod}{\operatorname{Mod}}
\providecommand{\Aut}{}\renewcommand{\Aut}{\operatorname{Aut}}
\providecommand{\GL}{}\renewcommand{\GL}{\operatorname{GL}}
\providecommand{\SL}{}\renewcommand{\SL}{\operatorname{SL}}
\providecommand{\Gm}{}\renewcommand{\Gm}{\mathbb G_{\mathrm m}}
\providecommand{\Der}{}\renewcommand{\Der}{\operatorname{Der}}
\providecommand{\id}{}\renewcommand{\id}{\operatorname{id}}
\providecommand{\Ch}{}\renewcommand{\Ch}{\operatorname{Ch}}
\providecommand{\cO}{}\renewcommand{\cO}{\mathcal O}
\providecommand{\C}{}\renewcommand{\C}{\mathbb C}
\providecommand{\Q}{}\renewcommand{\Q}{\mathbb Q}
\providecommand{\Z}{}\renewcommand{\Z}{\mathbb Z}
\providecommand{\Zell}{}\renewcommand{\Zell}{\mathbf Z_\ell}
\providecommand{\Qell}{}\renewcommand{\Qell}{\mathbf Q_\ell}
\providecommand{\ch}{}\renewcommand{\ch}{\operatorname{ch}}
\providecommand{\Todd}{}\renewcommand{\Todd}{\operatorname{Todd}}
\providecommand{\Gtop}{}\renewcommand{\Gtop}{\operatorname{Gr}_{\mathrm{top}}}
\providecommand{\Fil}{}\renewcommand{\Fil}{\operatorname{Fil}}
\providecommand{\Qcoh}{}\renewcommand{\Qcoh}{\operatorname{Qcoh}}
\providecommand{\Parf}{}\renewcommand{\Parf}{\operatorname{Parf}}
\providecommand{\Lib}{}\renewcommand{\Lib}{\operatorname{Lib}}
\providecommand{\Loclib}{}\renewcommand{\Loclib}{\operatorname{Loclib}}
\providecommand{\Proj}{}\renewcommand{\Proj}{\operatorname{Proj}}
\providecommand{\Hom}{}\renewcommand{\Hom}{\operatorname{Hom}}
\providecommand{\Diff}{}\renewcommand{\Diff}{\operatorname{Diff}}
\providecommand{\Dcat}{}\renewcommand{\Dcat}{\mathrm D}
\providecommand{\Tr}{}\renewcommand{\Tr}{\operatorname{Tr}}
\providecommand{\toramp}{}\renewcommand{\toramp}{\operatorname{tor.amp}}
\providecommand{\parfamp}{}\renewcommand{\parfamp}{\operatorname{par.amp}}
\providecommand{\coh}{}\renewcommand{\coh}{\operatorname{coh}}
\providecommand{\qcoh}{}\renewcommand{\qcoh}{\operatorname{qcoh}}
\providecommand{\parf}{}\renewcommand{\parf}{\operatorname{parf}}
\providecommand{\torf}{}\renewcommand{\torf}{\operatorname{torf}}
\providecommand{\amp}{}\renewcommand{\amp}{\operatorname{amp}}
\providecommand{\num}{}\renewcommand{\num}{\operatorname{num}}
\providecommand{\naif}{}\renewcommand{\naif}{\operatorname{naif}}
\providecommand{\alg}{}\renewcommand{\alg}{\operatorname{alg}}

% repair108 suppresses diagnostics globally.  Restore the ordinary TeX
% thresholds used by the clean standalone tail build.
\hbadness=1000
\hfuzz=0.1pt
\allowdisplaybreaks

% END INCLUDED SGA6_English_current_buildable_source_20260823/SGA6_tail_macro_state.tex

% BEGIN INCLUDED SGA6_English_current_buildable_source_20260823/SGA6_idx532_702_and_backmatter_English_Synchronized_body.tex

% Current rescribe index 532; printed volume page 519; source-PDF page 526.
\sgaStableTarget{sga6.e10.expose}\section*{Expos\'e X. Formalism of Intersections on Proper Algebraic Schemes}
\begin{center}
by O. Jussila\\
(with an appendix by A. Grothendieck)
\end{center}

\sgaStableTarget{sga6.e10.sec.1}\section*{1. Compatibility of the filtrations with the composition law}
\[
K^\circ(X)\times K_\bullet(X)\longrightarrow K_\bullet(X).
\]

\sgaStableTarget{sga6.e10.sec.1.1}\subsection*{1.1}
% SGA-SRC-001616
First let $X$ be a noetherian scheme, and let $K_\bullet(X)$ be the group of
classes of coherent $\mathcal O_X$-modules. For every coherent
$\mathcal O_X$-module $F$, we denote its class in $K_\bullet(X)$ by
$\cl_\bullet(F)$.

\sgaStableTarget{sga6.e10.def.1-1-1}\textbf{Definition 1.1.1.}
For every integer $n\in\mathbb Z$, let $\Filt_n(X)$ denote the subgroup of
$K_\bullet(X)$ generated by the classes $\cl_\bullet(F)$ of coherent
$\mathcal O_X$-modules $F$ such that
\[
\dim(\Supp(F))\leq n.
\]
This defines an increasing filtration $(\Filt_n(X))_{n\in\mathbb Z}$ of
$K_\bullet(X)$ such that $\Filt_n(X)=0$ for $n<0$.

\sgaStableTarget{sga6.e10.prop.1-1-2}\textbf{Proposition 1.1.2.}
Let $X$ be a noetherian scheme, let $F$ be a coherent $\mathcal O_X$-module
such that $\dim(\Supp(F))\leq n$, and, for $n\in\mathbb Z$, let $X(n)$ be the
subset
\[
\{x\in X\,;\,\dim(\overline{\{x\}})=n\}
\]
of $X$. For every $x\in X$, let $Z(x)$ denote the integral closed subscheme
associated with $\overline{\{x\}}$. Then the length of the
$\mathcal O_{X,x}$-module $F_x$ is finite for every $x\in X(n)$ and is
nonzero for only finitely many $x\in X(n)$, and one has
% Current rescribe index 533; printed volume page 520; source-PDF page 527.
\sgaStableTarget{sga6.e10.eq.1-1-3}\begin{equation}\tag{1.1.3}
\cl_\bullet(F)=
\sum_{x\in X(n)}\longg(F_x)\cl_\bullet(\mathcal O_{Z(x)})+w,
\end{equation}
where $w\in\Filt_{n-1}(X)$.

By immediate induction on $n$, one obtains:

\sgaStableTarget{sga6.e10.cor.1-1-4}\textbf{Corollary 1.1.4.}
For every noetherian scheme $X$ and every integer $n$, $\Filt_n(X)$ is the
subgroup of $K_\bullet(X)$ generated by the classes
$\cl_\bullet(\mathcal O_Z)$, where $Z$ ranges over the integral closed
subschemes of $X$ of dimension at most $n$.

\emph{Proof of 1.1.2.}
Let $X'=\Supp(F)$. Since $X'$ is a noetherian space of dimension at most $n$,
the elements of $X'\cap X(n)$ are maximal in $X'$ and hence finite in number;
moreover, $\longg(F_x)$ is finite for every $x\in X'\cap X(n)$. This proves
the first assertion.

Let $\mathcal K=\Coh(X)$ be the category of coherent $\mathcal O_X$-modules,
and let $\mathcal K'$ be the subset of $\operatorname{Ob}(\mathcal K)$ formed
by the coherent $\mathcal O_X$-modules concentrated on $X'$ that satisfy
condition \hyperlink{sga6.e10.eq.1-1-3}{(1.1.3)}. Plainly $\mathcal K'$ contains the sheaves
$\mathcal O_{Z(x)}$ for $x\in X'$. It is therefore enough to show that
$\mathcal K'$ is an ``exact'' subset of $\operatorname{Ob}(\mathcal K)$,
for then the devissage lemma (EGA III 3.1.2) gives $F\in\mathcal K'$.

Thus let
\sgaStableTarget{sga6.e10.eq.1-1-5}\begin{equation}\tag{1.1.5}
0\longrightarrow G'\longrightarrow G\longrightarrow G''\longrightarrow0
\end{equation}
be an exact sequence of $\mathcal O_X$-modules having two of its terms in
$\mathcal K'$. The sheaves $G'$, $G$, and $G''$ are then all concentrated on
$X'$. Hence the lengths $\ell'_x$, $\ell_x$, and $\ell''_x$ of the
$\mathcal O_{X,x}$-modules $G'_x$, $G_x$, and $G''_x$ are finite for every
% Current rescribe index 534; printed volume page 521; source-PDF page 528.
$x\in X(n)$, and
\sgaStableTarget{sga6.e10.eq.1-1-6}\begin{equation}\tag{1.1.6}
\ell_x=\ell'_x+\ell''_x\qquad(x\in X(n)).
\end{equation}
Assertion \hyperlink{sga6.e10.eq.1-1-3}{(1.1.3)} for all the terms of \hyperlink{sga6.e10.eq.1-1-5}{(1.1.5)} now follows immediately from
the hypothesis and the equation
\sgaStableTarget{sga6.e10.eq.1-1-7}\begin{equation}\tag{1.1.7}
\cl_\bullet(G)=\cl_\bullet(G')+\cl_\bullet(G'').
\end{equation}

It is known that $K_\bullet$ is a covariant functor on the category of locally
noetherian schemes with respect to proper morphisms. The following proposition
asserts that our filtration is functorial.

\sgaStableTarget{sga6.e10.prop.1-1-8}\textbf{Proposition 1.1.8.}
Let $f:X'\to X$ be a proper morphism of locally noetherian schemes, and let
\[
f_*:K_\bullet(X')\longrightarrow K_\bullet(X)
\]
be the morphism induced by $f$. Then
\sgaStableTarget{sga6.e10.eq.1-1-9}\begin{equation}\tag{1.1.9}
f_*(\Filt_n(X'))\subset\Filt_n(X)
\end{equation}
for every integer $n$.

\emph{Proof.}
If $F$ is a coherent $\mathcal O_{X'}$-module, then by definition (\hyperlink{sga6.e04.para.2-11}{IV 2.11}),
\[
f_*(\cl_\bullet(F))=
\sum_{i\in\mathbb N}(-1)^i\cl_\bullet(R^if_*(F)).
\]
The $\mathcal O_X$-modules $R^if_*(F)$ are concentrated on
$f(\Supp(F))$, and by EGA IV 5.4.1(ii),
\[
\dim(\Supp(F))\geq\dim(f(\Supp(F))),
\]
which proves the assertion.

% Current rescribe index 535; printed volume page 522; source-PDF page 529.
\sgaStableTarget{sga6.e10.sec.1.2}\subsection*{1.2}
Let now $X$ be an arbitrary scheme, and consider the group $K^\circ(X)$ of
classes of locally free $\mathcal O_X$-modules (henceforth understood to be of
finite type). It is an augmented pre-$\lambda$-algebra over the ring
$H^0(X,\mathbb Z)$ (\hyperlink{sga6.e05.ex.3-9-1}{V 3.9.1}). The map
$f\rightsquigarrow f\cdot1\in K^\circ(X)$ allows us henceforth to identify
$H^0(X,\mathbb Z)$ with a subring of $K^\circ(X)$. If $E$ is a locally free
$\mathcal O_X$-module, we denote its class in $K^\circ(X)$ by
$\cl^\circ(E)$. The operators $\lambda^k$ ($k\in\mathbb N$) are defined by
\sgaStableTarget{sga6.e10.eq.1-2-1}\begin{equation}\tag{1.2.1}
\lambda^k(\cl^\circ(E))=\cl^\circ(\wedge_{\mathcal O}^kE)
\end{equation}
and by
\sgaStableTarget{sga6.e10.eq.1-2-2}\begin{equation}\tag{1.2.2}
\lambda^k(x+y)=
\sum_{\substack{i+j=k\\ i,j\in\mathbb N}}\lambda^i(x)\lambda^j(y).
\end{equation}

(\hyperlink{sga6.e05.ex.2-2}{V 2.2}). Recall that in the $\lambda$-filtration
$(\Filt^i(X))_{i\in\mathbb Z}$ of $K^\circ(X)$, one has
$\Filt^i(X)=K^\circ(X)$ for $i\leq0$, while for $i\geq1$, $\Filt^i(X)$ is
the ideal generated by the products
\sgaStableTarget{sga6.e10.eq.1-2-3}\begin{equation}\tag{1.2.3}
P=\gamma^{i_1}(x_1)\cdots\gamma^{i_k}(x_k)
\end{equation}
with $1\leq i_1,\ldots,i_k$; $i_1+\cdots+i_k\geq i$; and
$x_1,\ldots,x_k\in I=\Ker(\varepsilon)$, where $\varepsilon$ is the
augmentation homomorphism (cf. \hyperlink{sga6.e05.para.3-10}{Expos\'e V 3.10)}. It was proved in \hyperlink{sga6.e05.para.3-10}{V 3.10} that
$\Filt^i(X)$ is even the $H^0(X,\mathbb Z)$-submodule, hence, when $X$ is
connected, the abelian subgroup generated by the products \hyperlink{sga6.e10.eq.1-2-3}{(1.2.3)}, where one
may moreover suppose that the elements $x_p$ ($p=1,\ldots,k$) have the form
\sgaStableTarget{sga6.e10.eq.1-2-4}\begin{equation}\tag{1.2.4}
x_p=\cl^\circ(E_p)-n_p,
\end{equation}
with $E_p$ a locally free $\mathcal O_X$-module of rank
$n_p\in H^0(X,\mathbb Z)$.

% Current rescribe index 536; printed volume page 523; source-PDF page 530.
\sgaStableTarget{sga6.e10.sec.1.3}\subsection*{1.3}
Recall (\hyperlink{sga6.e04.para.2-10}{IV 2.10}) that the bifunctor
\[
(E,F)\rightsquigarrow E\otimes_{\mathcal O_X}F,
\]
where $E$ is a locally free $\mathcal O_X$-module and $F$ is a coherent
$\mathcal O_X$-module, induces a composition law
\begin{equation}\tag{\sgaStableTargetInline{sga6.e10.eq.1-3-1}1.3.1}
K^\circ(X)\times K_\bullet(X)\longrightarrow K_\bullet(X)
\end{equation}
that makes $K_\bullet(X)$ a unital $K^\circ(X)$-module. The key result of the
present expos\'e is the following.

\sgaStableTarget{sga6.e10.thm.1-3-2}\textbf{Theorem 1.3.2.}
Let $X$ be a noetherian scheme. For all $i,j\in\mathbb Z$, the composition law
\hyperlink{sga6.e10.eq.1-3-1}{(1.3.1)} induces by restriction a composition law
\sgaStableTarget{sga6.e10.eq.1-3-3}\begin{equation}\tag{1.3.3}
\Filt^i(X)\times\Filt_j(X)\longrightarrow\Filt_{j-i}(X).
\end{equation}

\sgaStableTarget{sga6.e10.cor.1-3-3}\textbf{Corollary 1.3.3.}
If $X$ is a noetherian scheme of dimension at most $j$, then the canonical
homomorphism
\[
\Theta_X:K^\circ(X)\longrightarrow K_\bullet(X),
\]
defined by $\Theta_X(x)=x\cdot\cl_\bullet(\mathcal O_X)$
($x\in K^\circ(X)$), induces for every $i\in\mathbb Z$ a homomorphism
\[
\Filt^i(X)\longrightarrow\Filt_{j-i}(X).
\]

\emph{Proof.}
First let us show that \hyperlink{sga6.e10.cor.1-3-3}{(1.3.3)}, for $i+j\leq k\in\mathbb Z$ and every integral
closed subscheme of dimension at most $j$, implies \hyperlink{sga6.e10.thm.1-3-2}{Theorem 1.3.2} for
$i+j\leq k$. By \hyperlink{sga6.e10.cor.1-1-4}{(1.1.4)}, it is enough to show that, for every
$x\in\Filt^i(X)$ and every integral closed subscheme $Z$ of dimension at most
$j$, one has
% Current rescribe index 537; printed volume page 524; source-PDF page 531.
\[
x\cl_\bullet(\mathcal O_Z)\in\Filt_{j-i}(X).
\]
To avoid confusion, let $1_Z$ denote the class of $\mathcal O_Z$ in
$K_\bullet(Z)$. If $\alpha:Z\to X$ is the canonical immersion, the projection
formula (IV \ldots) gives
\[
x\cl_\bullet(\mathcal O_Z)
=x\alpha_*(1_Z)
=\alpha_*(\alpha^*(x)1_Z).
\]
Since $\alpha^*(x)\in\Filt^i(Z)$, \hyperlink{sga6.e10.prop.1-1-8}{(1.1.8)} and \hyperlink{sga6.e10.cor.1-3-3}{(1.3.3)}, applied with $X=Z$
and $i+j\leq k$, give
\[
\alpha_*(\alpha^*(x)1_Z)\in\Filt_{j-i}(X),
\]
which proves the assertion.

\hyperlink{sga6.e10.thm.1-3-2}{Theorem 1.3.2} is trivial for $i+j\leq0$. Suppose by induction that it has been
proved for $i+j<k$ and for every noetherian scheme $X$. By the preceding
remark, it remains to prove that, for every integral noetherian scheme $X$ of
dimension $j$ and every product \hyperlink{sga6.e10.eq.1-2-3}{(1.2.3)} satisfying \hyperlink{sga6.e10.eq.1-2-4}{(1.2.4)}, one has
\sgaStableTarget{sga6.e10.eq.1-3-4}\begin{equation}\tag{1.3.4}
P\cl_\bullet(\mathcal O_X)\in\Filt_{j-i}(X).
\end{equation}
If $k>1$ in \hyperlink{sga6.e10.eq.1-2-3}{(1.2.3)}, the assertion follows from the induction hypothesis and
the associativity of the law \hyperlink{sga6.e10.eq.1-3-1}{(1.3.1)}. We may therefore suppose, since $X$ is
connected, that
\sgaStableTarget{sga6.e10.eq.1-3-5}\begin{equation}\tag{1.3.5}
x=P=\gamma^m(\cl^\circ(E)-n),
\end{equation}
where $m\geq i$ and $E$ is a locally free $\mathcal O_X$-module of rank
$n>0$. We distinguish two cases.

I. Suppose that there exist invertible $\mathcal O_X$-modules
$L_1,\ldots,L_n$ such that
\[
\cl^\circ(E)=\cl^\circ(L_1)+\cdots+\cl^\circ(L_n)\in K^\circ(X).
\]

% Current rescribe index 538; printed volume page 525; source-PDF page 532.
It follows from (V) that $x$ is the $m$-th elementary symmetric function in
the $\cl^\circ(L_r)-1$ ($r=1,\ldots,n$), and hence is a homogeneous
polynomial of degree $m$ with integral coefficients in the variables
$(\cl^\circ(L_r)-1)$ ($r=1,\ldots,n$). Applying again the induction
hypothesis and the associativity of the law \hyperlink{sga6.e10.eq.1-3-1}{(1.3.1)}, we are reduced to proving
that, for every invertible $\mathcal O_X$-module $L$, one has
\sgaStableTarget{sga6.e10.eq.1-3-6}\begin{equation}\tag{1.3.6}
(\cl^\circ(L)-1)\cl_\bullet(\mathcal O_X)\in\Filt_{j-1}(X).
\end{equation}
Since $X$ is integral, $L$ and $\mathcal O_X$ are isomorphic to submodules of
the constant $\mathcal O_X$-module $R_X$ of rational functions on $X$. The
$\mathcal O_X$-module sum $F=\mathcal O_X+L$ is coherent, and one obtains
\begin{align*}
(\cl^\circ(L)-1)\cl_\bullet(\mathcal O_X)
&=\cl_\bullet(L)-\cl_\bullet(\mathcal O_X)\\
&=(\cl_\bullet(F)-\cl_\bullet(\mathcal O_X))
  -(\cl_\bullet(F)-\cl_\bullet(L))\\
&=\cl_\bullet(F/\mathcal O_X)-\cl_\bullet(F/L).
\end{align*}
But the fibers at the generic point of the sheaves $F/\mathcal O_X$ and
$F/L$ vanish; their supports in $X$ are therefore nowhere dense (since $X$
is irreducible), and hence have dimension at most $j-1$. Assertion \hyperlink{sga6.e10.eq.1-3-6}{(1.3.6)}
follows immediately.

II. It remains to prove the general case. Introduce the flag bundle
\[
g:D\longrightarrow X
\]
of $E$ over $X$, on which there exist invertible $\mathcal O_D$-modules
$L_1,\ldots,L_n$ such that
\sgaStableTarget{sga6.e10.eq.1-3-7}\begin{equation}\tag{1.3.7}
g^*(\cl^\circ(E))=\cl^\circ(L_1)+\cdots+\cl^\circ(L_n)\in K^\circ(D).
\end{equation}
Let $e$ be the generic point of $X$, and let $k(e)$ be its residue field.
There exists a rational point $u$ of the generic fiber $D_e$ over $k(e)$,

% Current rescribe index 539; printed volume page 526; source-PDF page 533.
because these points are in one-to-one correspondence with flags of type
$(1,\ldots,1)$ ($n$ terms) in the vector space $E_e$ of dimension $n>0$ over
$k(e)$. These rational points also correspond to germs of sections of $g$ in
a neighborhood of $e$, so there exists such a section $s:U\to D$ over an
open subset $U$ of $X$. Let $X'=\overline{s(U)}$, endowed with its reduced
induced structure, and let $f:X'\to X$ be the restriction of $g$ to $X'$.
The morphism $f$ is projective and, since it is an isomorphism over $s(U)$,
it is birational.

Thus $\dim(X')\leq\dim(X)$ (EGA IV 5.6.6.1). On the other hand,
$f(X')=g(X')$ is closed and contains $U$, so $f$ is surjective and
$\dim(X')\geq\dim(X)$ (EGA IV 5.4.1(ii)). The hypotheses of case I are
therefore satisfied for $X'$ and $x'=f^*(x)\in\Filt^i(X')$, and one obtains
\sgaStableTarget{sga6.e10.eq.1-3-7.occ2}\begin{equation}\tag{1.3.7'}
f^*(x)\cl_\bullet(\mathcal O_{X'})\in\Filt_{j-i}(X').
\end{equation}
Consequently, by the projection formula and \hyperlink{sga6.e10.prop.1-1-8}{(1.1.8)},
\sgaStableTarget{sga6.e10.eq.1-3-8}\begin{equation}\tag{1.3.8}
x f_*(\cl_\bullet(\mathcal O_{X'}))\in\Filt_{j-i}(X).
\end{equation}
On the other hand, the canonical homomorphism
\[
\mathcal O_X\longrightarrow f_*(\mathcal O_{X'})
\]
is an isomorphism over the dense open subset $U$, and the
$\mathcal O_X$-modules $R^if_*(\mathcal O_{X'})$ ($i\geq1$) are concentrated
on the nowhere-dense closed subset $X-U$. It follows that
\sgaStableTarget{sga6.e10.eq.1-3-9}\begin{equation}\tag{1.3.9}
f_*(\cl_\bullet(\mathcal O_{X'}))=\cl_\bullet(\mathcal O_X)+w,
\end{equation}

% Current rescribe index 540; printed volume page 527; source-PDF page 534.
where $w\in\Filt_{j-1}(X)$. By the induction hypothesis,
$x\cdot w\in\Filt_{j-i-1}(X)$; hence, by \hyperlink{sga6.e10.eq.1-3-8}{(1.3.8)} and \hyperlink{sga6.e10.eq.1-3-9}{(1.3.9)},
\[
x\cl_\bullet(\mathcal O_X)\in\Filt_{j-i}(X).
\]
This proves the assertion.

\sgaStableTarget{sga6.e10.rem.1-4}\textbf{Remark 1.4.}
Let $X$ be a locally noetherian scheme. Observe that, in this section, we
have used only the following properties of the ``dimension function''
\[
D(x)=D_Y(x)=\dim(\overline{\{x\}})
\]
on proper $X$-schemes $Y$:
\begin{enumerate}[label=(\roman*)]
\item For every proper $X$-scheme $Y$ and every $x\in Y$, one has
  $D_Y(x)\in\mathbb Z$.
\item If $y\ne x$ is a specialization of $x$, then $D(y)<D(x)$.
\item For every $X$-morphism $f:Y_1\to Y_2$ of proper $X$-schemes and every
  $y\in Y_1$, one has
  \[
  D_{Y_2}(f(y))\leq D_{Y_1}(y).
  \]
\item Equality holds in (iii) if $f$ is an immersion, or if $Y_1$ is
  irreducible, $f$ is birational and surjective, and $y$ is the generic point
  of $Y_1$.
\end{enumerate}
Thus all the results of this section, in particular \hyperlink{sga6.e10.thm.1-3-2}{(1.3.2)} and \hyperlink{sga6.e10.cor.1-3-3}{(1.3.3)},
remain valid unchanged if ``dimension'' is replaced by a function $D$
satisfying conditions (i)--(iv) above and the filtration of $K_\bullet(X)$ is
defined using such a function $D$.

% Current rescribe index 541; printed volume page 528; source-PDF page 535.
\sgaStableTarget{sga6.e10.ex.1-5}\textbf{Example 1.5.}
Let $S$ be a locally noetherian universally catenary scheme (for example, a
subscheme of a Cohen--Macaulay scheme). Put
\sgaStableTarget{sga6.e10.eq.1-5-1}\begin{equation}\tag{1.5.1}
D_S(x)=n-\codim(\overline{\{x\}},S)
      =n-\dim\mathcal O_{S,x}\qquad(x\in S),
\end{equation}
and, for every proper $S$-scheme $s_Y:Y\to S$ and every $y\in Y$, also put
\sgaStableTarget{sga6.e10.eq.1-5-2}\begin{equation}\tag{1.5.2}
D_Y(y)=D_S(s_Y(y))+\degtr_{k(s_Y(y))}k(y).
\end{equation}
If $f:Y_1\to Y_2$ is an $S$-morphism of proper $S$-schemes, the additivity of
transcendence degrees and \hyperlink{sga6.e10.eq.1-5-2}{(1.5.2)} immediately give
\sgaStableTarget{sga6.e10.eq.1-5-3}\begin{equation}\tag{1.5.3}
D_{Y_1}(y)=D_{Y_2}(f(y))+\degtr_{k(f(y))}k(y)
\qquad(y\in Y_1).
\end{equation}
We show that the functions $D_Y$ so defined satisfy conditions (i)--(iv) of
1.4. By \hyperlink{sga6.e10.eq.1-5-3}{(1.5.3)}, condition (ii) is the only nontrivial one. In view of (iv),
we may plainly suppose that $Y=\overline{\{x\}}$. Factor the structural
morphism $s_Y$ through its image $s_Y(Y)=Z\hookrightarrow S$, say
$s_Y=g\circ f$:
\[
\begin{tikzcd}
Y \arrow[r,"f"] & Z \arrow[r,hook,"g"] & S,
\end{tikzcd}
\]
where $Y$ and $Z$ are irreducible, $f$ is proper and surjective, and $g$ is a
closed immersion. Put $x'=s_Y(x)$ and $y'=s_Y(y)$. From \hyperlink{sga6.e10.eq.1-5-2}{(1.5.2)} and \hyperlink{sga6.e10.eq.1-5-1}{(1.5.1)}
one obtains
\[
\begin{split}
D_Y(x)-D_Y(y)={}&\dim\mathcal O_{S,y'}-\dim\mathcal O_{S,x'}\\
&+\degtr_{k(x')}k(x)-\degtr_{k(y')}k(y).
\end{split}
\]
Since $S$ is catenary, one has

% Current rescribe index 542; printed volume page 529; source-PDF page 536.
\[
\dim\mathcal O_{S,y'}-\dim\mathcal O_{S,x'}
=\dim\mathcal O_{Z,y'}-\dim\mathcal O_{Z,x'}
\]
(EGA $0_{\mathrm{IV}}$ 14.3.2.1), whence
\sgaStableTarget{sga6.e10.eq.1-5-4}\begin{equation}\tag{1.5.4}
\begin{split}
D_Y(x)-D_Y(y)={}&
[\dim\mathcal O_{Z,y'}-\degtr_{k(y')}k(y)]\\
&-[\dim\mathcal O_{Z,x'}-\degtr_{k(x')}k(x)].
\end{split}
\end{equation}
Since $Z$ is universally catenary, EGA IV 5.6.5.1 applies and gives
\[
C_1=\dim\mathcal O_{Y,y}-e,
\qquad
C_2=\dim\mathcal O_{Y,x}-e,
\]
where $C_1$ and $C_2$ are the two terms on the right of \hyperlink{sga6.e10.eq.1-5-4}{(1.5.4)}. It follows
immediately from \hyperlink{sga6.e10.eq.1-5-4}{(1.5.4)} that
\[
D_Y(x)-D_Y(y)
=\dim\mathcal O_{Y,y}-\dim\mathcal O_{Y,x}>0,
\]
as required.

Thus \hyperlink{sga6.e10.thm.1-3-2}{(1.3.2)} and \hyperlink{sga6.e10.cor.1-3-3}{(1.3.3)} remain valid for a noetherian universally catenary
scheme $X$ of dimension $n$ if the filtration of $K_\bullet(X)$ is defined by
means of \hyperlink{sga6.e10.eq.1-5-1}{(1.5.1)}, according to the codimension of supports. This filtration
coincides with that of \hyperlink{sga6.e10.def.1-1-1}{(1.1.1)} when, in addition, $X$ is biequidimensional.

\sgaStableTarget{sga6.e10.sec.1.6}\subsection*{1.6}
One may ask whether it is always possible, on a noetherian scheme $X$ of
finite dimension, to define a function $D$ possessing property (i) of \hyperlink{sga6.e10.rem.1-4}{1.4}
and also the following property, stronger than (ii):
\begin{quote}
(ii bis) If $x\in X$ and $y\ne x$ is an immediate specialization of $x$,
then $D(x)=D(y)+1$.
\end{quote}

% Current rescribe index 543; printed volume page 530; source-PDF page 537.
For this it is plainly necessary that $X$ be catenary, and sufficient that
$X$ be biequidimensional (EGA $0_{\mathrm{IV}}$ 14.3.3). There nevertheless
exist noetherian schemes of finite dimension that are catenary and
equidimensional but admit no such function $D$. To see this, modify the
example of EGA IV 5.6.11 as follows. Take two copies $Y_1$ and $Y_2$ of the
spectrum of the ring $E$ of loc. cit., and let $y_i$ and $z_i$ ($i=1,2$) be
the closed points of $Y_i$ corresponding respectively to the maximal ideals
$\mathfrak m$ and $\mathfrak m'$ of $E$, of heights $2$ and $1$. Recall that
the residue fields of $\mathfrak m$ and $\mathfrak m'$ are isomorphic. One
can therefore construct a scheme $X$ by gluing $Y_1$ and $Y_2$, identifying
$y_1$ with $z_2$ and $y_2$ with $z_1$. The schemes $Y_1$ and $Y_2$ may then
be identified with the two irreducible components of $X$; inside $Z$,
$Y_1\cap Y_2=\{y_1,y_2\}$ has dimension zero. Thus $X$ is a noetherian,
catenary, equidimensional scheme of dimension $2$. If (ii bis) held on $X$,
one would have
\[
D(y_1)=D(z_1)-1=D(y_2)-1,
\]
and, on the other hand,
\[
D(y_1)=D(z_2)=D(y_2)+1,
\]
which is a contradiction.

% Current rescribe index 544; printed volume page 531; source-PDF page 538.
\sgaStableTarget{sga6.e10.sec.2}\section*{2. Snapper polynomials}

\sgaStableTarget{sga6.e10.sec.2.1}\subsection*{2.1}
We first recall the definition and a few elementary properties of polynomial
maps.

\sgaStableTarget{sga6.e10.def.2-1-1}\textbf{Definition 2.1.1.}
Let $f$ be a map from an abelian group $A$ to an abelian group $B$. We say
that $f$ is a polynomial map if, for every sequence $x_1,\ldots,x_r$
($r\geq1$) of elements of $A$, there is a family
\[
\bigl(y(x_1,\ldots,x_r;k_1,\ldots,k_r)\bigr)_{k_1,\ldots,k_r\geq0}
\]
of elements of $B$, almost all zero, satisfying the following condition: for
every sequence $n_1,\ldots,n_r$ of integers, one has
\[
f(n_1x_1+\cdots+n_rx_r)=
\sum_{k_1,\ldots,k_r\geq0}
n_1^{k_1}\cdots n_r^{k_r}
y(x_1,\ldots,x_r;k_1,\ldots,k_r).
\]
We say that $f$ has degree $\leq m$ if, for every integer $r\geq1$ and all
$x_1,\ldots,x_r\in A$, one has
\[
y(x_1,\ldots,x_r;k_1,\ldots,k_r)=0
\]
whenever $k_1+\cdots+k_r>m$.

\sgaStableTarget{sga6.e10.lem.2-1-2}\textbf{Lemma 2.1.2.}
Every homomorphism of abelian groups is a polynomial map of degree $\leq1$.
If $h=f\circ g$, where $f$ is a polynomial map of degree $p$ and $g$ is a
polynomial map of degree $q$, then $h$ is a polynomial map of degree
$\leq pq$.

The proof is trivial.

\sgaStableTarget{sga6.e10.lem.2-1-3}\textbf{Lemma 2.1.3.}
Let
\[
0\longrightarrow G'\longrightarrow G\xrightarrow{j}G''\longrightarrow0
\]
be an exact sequence of abelian groups such that $G'$ is $p$-divisible for
some prime $p$ and $G''$ is a torsion group. Then every polynomial map
$f:G\to\mathbb Q$ that factors through $\mathbb Z$ is constant.

% Current rescribe index 545; printed volume page 532; source-PDF page 539.
\emph{Proof.}
It is enough to show that $f(x)=f(0)$ for every $x\in G$.

(a) Suppose first that $x$ is a torsion element, say $mx=0$
($m\in\mathbb Z$, $m\ne0$). Then the polynomial
$P_x(n)=f(nx)-f(0)$ has infinitely many roots, namely all integers divisible
by $m$. Thus $P_x$ is identically zero, and $f(x)=f(0)$.

(b) Suppose now that $x$ is an element of infinite order in $G$, divisible by every power
$p^k$ ($k\in\mathbb Z$). Consider the polynomials $P_k$ defined by
\[
P_k(n)=f(nx/p^k)=\sum_{i\in\mathbb N}n^iy(x/p^k;i)\in\mathbb Z
\qquad(n\in\mathbb Z;\ y(x/p^k;i)\in\mathbb Q).
\]
For every $n\in\mathbb Z$, one has the identity
\[
P_k(n)=P_{k+1}(pn),
\]
that is,
\[
\sum_{i\in\mathbb N}n^iy(x/p^k;i)=
\sum_{i\in\mathbb N}p^in^iy(x/p^{k+1};i).
\]
Comparing coefficients gives
\[
y(x/p^k;i)=p^iy(x/p^{k+1};i)\qquad(i\in\mathbb N).
\]

% Current rescribe index 546; printed volume page 533; source-PDF page 540.
Every polynomial $P:\mathbb Q\to\mathbb Q$ is continuous for the topology
induced by $\mathbb R$. Thus, for $k$ sufficiently large,
\[
\lvert P_x(n/p^k)-P_x(0)\rvert<1.
\]
But, by definition, $P_x(n/p^k)$ is an integer for all
$n,k\in\mathbb Z$. It follows that $P_x(n/p^k)-P_x(0)=0$ for sufficiently
large $k$. Hence $P_x$ is constant and, in particular,
$P_x(1)=P_x(0)$, i.e. $f(x)=f(0)$.

The two cases above show that the restriction of $f$ to $G'$ is constant.
We now show that $f$ is constant on each coset of $G'$ in $G$. For every
$b\in G$, let $g_b:G'\to\mathbb Q$ be defined by
\[
g_b(y)=f(b+y)\qquad(y\in G').
\]
The map $g_b$ is a polynomial map from $G'$ to $\mathbb Q$ that factors
through $\mathbb Z$, and is therefore constant by the preceding argument.
Thus $f$ is constant on the cosets of $G'$ in $G$. Consequently one may
write $f=h\circ j$, where $h:G''\to\mathbb Q$ is a polynomial map that
factors through $\mathbb Z$. Since $G''$ is a torsion group, $h$ is constant.
Therefore $f$ is constant.

\sgaStableTarget{sga6.e10.lem.2-1-4}\textbf{Lemma 2.1.4.}
Let $A$ be a unital commutative ring and $I$ a nilpotent ideal of $A$. Let
$P$ be the multiplicative subgroup of $A^*$ consisting of the invertible elements
of $A$ of the form $x=1+y$ with $y\in I$, and let $i:P\to A$ be the natural
injection.

% Current rescribe index 547; printed volume page 534; source-PDF page 541.
Then the composite map
\[
\alpha:P\xrightarrow{i}A\xrightarrow{j}A\otimes_{\mathbb Z}\mathbb Q,
\]
where $j(x)=x\otimes1$, is a polynomial map. If $I^{m+1}=0$, then $\alpha$
has degree $\leq m$.

\emph{Proof.}
Let $x_1,\ldots,x_r$ be a sequence of elements of $P$, with
$x_i=1+y_i$ ($i=1,\ldots,r$). The binomial formula gives, for every
$n_1,\ldots,n_r\in\mathbb Z$,
\sgaStableTarget{sga6.e10.eq.2-1-5}\begin{equation}\tag{2.1.5}
x_1^{n_1}\cdots x_r^{n_r}=
\sum_{k_1,\ldots,k_r\geq0}
\binom{n_1}{k_1}\cdots\binom{n_r}{k_r}
y_1^{k_1}\cdots y_r^{k_r},
\end{equation}
where the sum on the right is finite because the elements
$y_1,\ldots,y_r$ are nilpotent. The coefficients
$\binom{n_1}{k_1}\cdots\binom{n_r}{k_r}$ are polynomials with rational
coefficients in $n_1,\ldots,n_r$, of degree $k_1+\cdots+k_r$; hence
$\alpha$ is indeed a polynomial map. If $I^{m+1}=0$, then
$y_1^{k_1}\cdots y_r^{k_r}=0$ whenever $k_1+\cdots+k_r>m$. Thus $\alpha$
has degree $\leq m$.

\sgaStableTarget{sga6.e10.sec.2.2}\subsection*{2.2}
Now let $X$ be a noetherian scheme and let $y\in\Filt_j(X)$ \hyperlink{sga6.e10.def.1-1-1}{(1.1.1)}. The
assignment $x\mapsto xy$ defines a homomorphism of $K^\circ(X)$-modules
\[
\Theta_y:K^\circ(X)\longrightarrow K_\bullet(X),
\]
which, by \hyperlink{sga6.e10.thm.1-3-2}{(1.3.2)}, factors through the canonical epimorphism
\[
j:K^\circ(X)\longrightarrow
K^\circ(X)/\Filt^{j+1}(X)=A.
\]

% Current rescribe index 548; printed volume page 535; source-PDF page 542.
Let $\beta_y:A\to K_\bullet(X)$ be the homomorphism so defined. If $I$
denotes the ideal
\[
I=\Filt^1(X)/\Filt^{j+1}(X)
\]
of $A$, then $I^{j+1}=0$. Keeping the notation of \hyperlink{sga6.e10.lem.2-1-4}{Lemma 2.1.4}, one obtains a
homomorphism of abelian groups
\[
\gamma:\Pic(X)\longrightarrow P
\]
that sends the class $\cl_P(L)$ of an invertible $\mathcal O_X$-module $L$
in $\Pic(X)$ to the element $j(\cl^\circ(L))$ of $P$.

Lemmas \hyperlink{sga6.e10.lem.2-1-2}{2.1.2} and \hyperlink{sga6.e10.lem.2-1-4}{2.1.4} immediately give:

\sgaStableTarget{sga6.e10.prop.2-2-1}\textbf{Proposition 2.2.1.}
For every noetherian scheme $X$ and every $y\in\Filt_j(X)$
($j\in\mathbb Z$), the composite map
\[
\Pic(X)\longrightarrow K_\bullet(X)\otimes_{\mathbb Z}\mathbb Q,
\]
defined by
\[
\cl_P(L)\longmapsto(\cl^\circ(L)y)\otimes1,
\]
is a polynomial map of degree $\leq j$.

\sgaStableTarget{sga6.e10.sec.2.3}\subsection*{2.3}
If $X$ is a proper scheme over a field $k$ and
$g:X\to\Spec(k)$ is the structural morphism, then the homomorphism
\[
g_*:K_\bullet(X)\longrightarrow K_\bullet(\Spec(k))\simeq\mathbb Z
\]
sends the class $\cl_\bullet(F)$ of every coherent $\mathcal O_X$-module
$F$ to the Euler--Poincar\'e characteristic $\chi(F)\in\mathbb Z$ of $F$
on $X$. \hyperlink{sga6.e10.prop.2-2-1}{Proposition 2.2.1} immediately gives:

% Current rescribe index 549; printed volume page 536; source-PDF page 543.
\sgaStableTarget{sga6.e10.prop.2-3-1}\textbf{Proposition 2.3.1.}
For every scheme $X$ proper over $k$ and every $y\in\Filt_j(X)$, the map
\[
S_y:\Pic(X)\longrightarrow\mathbb Z\hookrightarrow\mathbb Q,
\]
defined by
\[
\cl_P(L)\longmapsto
g_*(\cl^\circ(L)y)=\chi(\cl^\circ(L)y),
\]
is a polynomial map of degree $\leq j$.

When $L_1,\ldots,L_n$ are invertible $\mathcal O_X$-modules and $F$ is a
coherent $\mathcal O_X$-module such that $\dim(\Supp(F))=j$, \hyperlink{sga6.e10.prop.2-3-1}{Proposition
2.3.1} implies in particular that the Euler--Poincar\'e characteristic
\sgaStableTarget{sga6.e10.eq.2-3-2}\begin{equation}\tag{2.3.2}
\chi(L_1^{\otimes m_1}\otimes\cdots\otimes
L_n^{\otimes m_n}\otimes F)
\end{equation}
is a numerical polynomial in $m_1,\ldots,m_n\in\mathbb Z$, of degree
$\leq\dim(\Supp(F))$, sometimes called the Snapper polynomial of $F$ and
the sequence $L_1,\ldots,L_n$.

\sgaStableTarget{sga6.e10.sec.2.4}\subsection*{2.4}
Recall [2] that, for the $\fpqc$ topology on $\Sch/k$, the Picard functor
$\Pic_{X/k}$ is the sheaf of groups associated with the presheaf
\[
T\longmapsto\Pic(X\times_kT)
\qquad(T\in\Ob(\Sch/k)).
\]
When $k$ is algebraically closed, the canonical homomorphism
\[
\Pic(X)\longrightarrow\Pic_{X/k}(\Spec(k))
\]
is an isomorphism.

% Current rescribe index 550; printed volume page 537; source-PDF page 544.
When $X$ is proper over $k$, it is known [4] that the Picard functor is
representable by a group scheme, locally of finite type over $k$, also denoted
$\Pic_{X/k}$. The neutral component $\Pic^0_{X/k}$ of $\Pic_{X/k}$ is a
connected commutative algebraic group over $k$.

\sgaStableTarget{sga6.e10.def.2-4-1}\textbf{Definition 2.4.1.}
Let $X$ be a scheme proper over a field $k$, and let $L$ be an invertible
$\mathcal O_X$-module (respectively, let $x\in\Pic(X)$). We say that $L$
(respectively, $x$) is $\tau$-equivalent to zero if there is an integer
$n\ne0$ such that $L^{\otimes n}$ (respectively, $nx$) is algebraically
equivalent to zero. The elements of $\Pic(X)$ algebraically equivalent
(respectively, $\tau$-equivalent) to zero form a subgroup $\Pic^0(X)$
(respectively, $\Pic^\tau(X)$) of $\Pic(X)$.

\sgaStableTarget{sga6.e10.lem.2-4-3}\textbf{Lemma 2.4.3.}
For every proper scheme $X$ over an algebraically closed field $k$, the group
$\Pic^0(X)$ is divisible by every integer $n$ prime to the characteristic of
$k$.

\emph{Proof.}
Since $k$ is algebraically closed, one has
\[
\Pic^0(X)\simeq\Pic^0_{X/k}(k).
\]
Since $\Pic^0_{X/k}$ is connected and of finite type, the $n$-th-power
endomorphism is surjective when $n$ is prime to the characteristic of $k$
(SGA 3 XV 1.3 and VI 1.3.2). Since $k$ is algebraically closed, the
corresponding endomorphism of $\Pic^0_{X/k}(k)$ is surjective.

It follows from \hyperlink{sga6.e10.lem.2-4-3}{2.4.3} that, in the exact sequence
\[
0\longrightarrow\Pic^0(X)\longrightarrow\Pic^\tau(X)
\longrightarrow T\longrightarrow0,
\]
$\Pic^0(X)$ is $p$-divisible for some prime $p$ and $T$ is a torsion group.
Thus \hyperlink{sga6.e10.lem.2-1-3}{Lemma 2.1.3} applies to the restriction of the polynomial map \hyperlink{sga6.e10.prop.2-3-1}{(2.3.1)} to
$\Pic^\tau(X)$. One obtains:

% Current rescribe index 551; printed volume page 538; source-PDF page 545.
\sgaStableTarget{sga6.e10.prop.2-4-4}\textbf{Proposition 2.4.4.}
The polynomial map $S_y$ of \hyperlink{sga6.e10.prop.2-3-1}{(2.3.1)} is constant on the cosets of
$\Pic^\tau(X)$ in $\Pic(X)$.

\emph{Proof.}
If $K$ is an algebraically closed extension of $k$, one has a commutative
diagram
\[
\begin{tikzcd}[column sep=large,row sep=large]
\Pic(X) \arrow[r,"S_y"] \arrow[d] &
K_\bullet(\Spec(k)) \arrow[r] &
\mathbb Z \arrow[r,hook] & \mathbb Q\\
\Pic(X_K) \arrow[r,"\overline S_y"'] &
K_\bullet(\Spec(K)) \arrow[ur] &&
\end{tikzcd}
\]
where $\overline S_y$ is the polynomial map on $\Pic(X_K)$ corresponding to
$S_y$. Hence one may suppose that $k$ is algebraically closed. The
assertion then follows immediately from Lemmas \hyperlink{sga6.e10.lem.2-1-3}{2.1.3} and 2.4.3.

\sgaStableTarget{sga6.e10.rem.2-4-5}\textbf{Remark 2.4.5.}
In \hyperlink{sga6.e13.expose}{Expos\'e XIII} this result will be proved without using representability of
the functor $\Pic_{X/k}$.

\sgaStableTarget{sga6.e10.sec.3}\section*{3. Projection formulas for the associated graded groups}

\sgaStableTarget{sga6.e10.sec.3.1}\subsection*{3.1}
Let $X$ be a noetherian scheme. Denote by $\Gro(X)$ the graded group associated
with the \mbox{$\lambda$-filtration} of $K^\circ(X)$, and by $\Grb(X)$ the graded
group associated with the filtration \hyperlink{sga6.e10.def.1-1-1}{1.1.1} of $K_\bullet(X)$. By \hyperlink{sga6.e10.thm.1-3-2}{Theorem
1.3.2}, the composition law \hyperlink{sga6.e10.eq.1-3-1}{(1.3.1)}
\[
K^\circ(X)\times K_\bullet(X)\longrightarrow K_\bullet(X)
\]
induces, on passing to the associated graded groups, composition laws
\sgaStableTarget{sga6.e10.eq.3-1-1}\begin{equation}\tag{3.1.1}
\Gr^i(X)\times\Gr_j(X)\longrightarrow\Gr_{j-i}(X)
\qquad(i,j\in\mathbb Z).
\end{equation}

% Current rescribe index 552; printed volume page 539; source-PDF page 546.
Thus $\Grb(X)$ becomes a module over the ring $\Gro(X)$, and even a graded
module over the graded ring $\Gro(X)$, provided that the sign of the grading
of $\Grb(X)$ is changed by putting
\[
\Gr^i_\bullet(X)=\Gr_{-i}(X)\qquad(i\in\mathbb Z).
\]

Now let
\[
f:X'\longrightarrow X
\]
be a proper morphism of noetherian schemes. Then $f$ induces a homomorphism
\[
f^*:\Gro(X)\longrightarrow\Gro(X')
\]
of graded rings (\hyperlink{sga6.e05.expose}{Expos\'e V}), and, by \hyperlink{sga6.e10.prop.1-1-8}{Proposition 1.1.8}, a homomorphism
\[
f_*:\Grb(X')\longrightarrow\Grb(X)
\]
of graded abelian groups. Passing the projection formula (\hyperlink{sga6.e04.para.2-11}{IV 2.11}) and
\hyperlink{sga6.e10.eq.3-1-1}{(3.1.1)} to the associated graded groups gives the graded projection formula
\sgaStableTarget{sga6.e10.eq.3-1-2}\begin{equation}\tag{3.1.2}
x\cdot f_*(y)=f_*\bigl(f^*(x)\cdot y\bigr)
\qquad(x\in\Gro(X),\ y\in\Grb(X')).
\end{equation}

\sgaStableTarget{sga6.e10.sec.3.2}\subsection*{3.2}
Let $f:X'\to X$ be a projective complete-intersection morphism (\hyperlink{sga6.e08.expose}{Expos\'e
VIII}), of virtual dimension $r$, between noetherian schemes, with $X$ carrying
an ample invertible $\mathcal O_X$-module. In \hyperlink{sga6.e08.expose}{Expos\'e VIII} it was proved that
$f$ induces a homomorphism
\sgaStableTarget{sga6.e10.eq.3-2-1}\begin{equation}\tag{3.2.1}
f_*:\Filt^k(X')_{\mathbb Q}\longrightarrow
\Filt^{k-r}(X)_{\mathbb Q}\qquad(k\in\mathbb Z),
\end{equation}
which moreover satisfies the projection formula (\hyperlink{sga6.e04.expose}{Expos\'e IV}):
\sgaStableTarget{sga6.e10.eq.3-2-2}\begin{equation}\tag{3.2.2}
f_*\bigl(f^*(x)\cdot y\bigr)=x\cdot f_*(y)
\qquad(x\in K^\circ(X)_{\mathbb Q},\ y\in K^\circ(X')_{\mathbb Q}).
\end{equation}

% Current rescribe index 553; printed volume page 540; source-PDF page 547.
Passing to the associated graded groups in \hyperlink{sga6.e10.eq.3-2-1}{(3.2.1)} and \hyperlink{sga6.e10.eq.3-2-2}{(3.2.2)}, one obtains a
homomorphism of degree $-r$:
\sgaStableTarget{sga6.e10.eq.3-2-3}\begin{equation}\tag{3.2.3}
f_*:\Gro(X')_{\mathbb Q}\longrightarrow\Gro(X)_{\mathbb Q},
\end{equation}
and, by \hyperlink{sga6.e10.thm.1-3-2}{Theorem 1.3.2}, a graded projection formula
\[
f_*\bigl(f^*(x)\cdot y\bigr)=x\cdot f_*(y)
\qquad(x\in\Gro(X)_{\mathbb Q},\ y\in\Gro(X')_{\mathbb Q}).
\]

\sgaStableTarget{sga6.e10.sec.3.3}\subsection*{3.3}
Finally, let $f:X'\to X$ be a morphism of finite Tor-dimension between
noetherian schemes. It is known (\hyperlink{sga6.e04.expose}{Expos\'e IV}) that $f$ induces a homomorphism
\[
f^*:K_\bullet(X)\longrightarrow K_\bullet(X')
\]
such that, for every coherent $\mathcal O_X$-module $F$,
\[
f^*(\cl_\bullet(F))=
\sum_i(-1)^i\cl_\bullet\bigl(
\Tor_i^{\mathcal O_X}(F,\mathcal O_{X'})\bigr).
\]
The $\mathcal O_{X'}$-modules
$\Tor_i^{\mathcal O_X}(F,\mathcal O_{X'})$ are supported on
$f^{-1}(\Supp(F))$. If the fibers of $f$ are everywhere of dimension
$\leq d$, then (EGA IV 5.6.7)
\[
\dim\bigl(f^{-1}(\Supp(F))\bigr)\leq\dim(\Supp(F))+d.
\]
It follows that $f$ induces, for every $k\in\mathbb Z$, a homomorphism
\[
f^*:\Filt_k(X)\longrightarrow\Filt_{k+d}(X'),
\]
hence a homomorphism of degree $d$ of graded abelian groups.

% Current rescribe index 554; printed volume page 541; source-PDF page 548.
Thus one has
\sgaStableTarget{sga6.e10.eq.3-3-1}\begin{equation}\tag{3.3.1}
f^*:\Grb(X)\longrightarrow\Grb(X').
\end{equation}

\sgaStableTarget{sga6.e10.sec.3.4}\subsection*{3.4}
Suppose finally that the morphisms \hyperlink{sga6.e10.eq.3-2-3}{(3.2.3)} and \hyperlink{sga6.e10.eq.3-3-1}{(3.3.1)} are both defined for a
morphism $f:X'\to X$ and that $r=d$. This is the case, for example, if $X$
is quasi-compact with an ample invertible $\mathcal O_X$-module and if
$X'=\mathbf P_X(E)$, where $E$ is a locally free $\mathcal O_X$-module of
rank $r+1$. The projection formula in the derived category then gives, for
every $x\in\Filt^i(X')_{\mathbb Q}$ and $y\in\Filt_j(X)_{\mathbb Q}$,
\[
f_*\bigl(x\cdot f^*(y)\bigr)=f_*(x)\cdot y
\in\Filt_{j+r-i}(X)_{\mathbb Q}.
\]
When $X$ is noetherian, \hyperlink{sga6.e10.thm.1-3-2}{Theorem 1.3.2} yields the corresponding graded formula.

\sgaStableTarget{sga6.e10.sec.4}\section*{4. Intersection numbers}

\sgaStableTarget{sga6.e10.sec.4.1}\subsection*{4.1}
Let $X$ be a proper scheme over a field $k$. Composing the product \hyperlink{sga6.e10.eq.3-1-1}{(3.1.1)}
\[
\Gr^i(X)\times\Gr_i(X)\longrightarrow\Gr_0(X)=\Filt_0(X)
\qquad(i\in\mathbb Z)
\]
with the Euler--Poincar\'e characteristic
\[
\chi:\Filt_0(X)\longrightarrow\mathbb Z
\]
gives a pairing called the \emph{intersection number}:
\sgaStableTarget{sga6.e10.eq.4-1-1}\begin{equation}\tag{4.1.1}
\Gr^i(X)\times\Gr_i(X)\longrightarrow\mathbb Z.
\end{equation}

% Current rescribe index 555; printed volume page 542; source-PDF page 549.
We denote this pairing by
\[
(x,y)\longmapsto\langle x,y\rangle\in\mathbb Z
\qquad(i\in\mathbb Z,\ x\in\Gr^i(X),\ y\in\Gr_i(X)).
\]

\sgaStableTarget{sga6.e10.ex.4-2}\textbf{Example 4.2.}
Suppose that $X$ has dimension $n$. Let $y\in\Gr_n(X)$, for example
$y=\cl_\bullet(\mathcal O_Z)$, where $Z$ is an integral closed subscheme of
$X$ of dimension $n$. Let $P$ also be a homogeneous polynomial of weight $n$,
\[
P\in\mathbb Z[T_1,\ldots,T_n],
\]
where $T_i$ is assigned weight $i$ ($i=1,\ldots,n$); for example, $P$ may
contain a monomial $T_{i_1}\cdots T_{i_r}$ with
$i_1+\cdots+i_r=n$. Substituting for each variable $T_i$ the Chern class
\[
c_i(x)=\cl^i\bigl(\gamma^i(x-\varepsilon(x))\bigr)\in\Gr^i(X)
\qquad(x\in K^\circ(X))
\]
gives an element
\[
P(c_1(x),\ldots,c_n(x))\in\Gr^n(X).
\]
The integer
\[
P_y(x)=\bigl\langle P(c_1(x),\ldots,c_n(x)),y\bigr\rangle
\]
is called the \emph{Chern number} of $x$ relative to $y$. When
$y=\cl_\bullet(\mathcal O_X)$, the reference to $y$ is omitted from this
terminology.

\sgaStableTarget{sga6.e10.sec.4.3}\subsection*{4.3}
Without any hypothesis on the dimension of $X$, consider invertible
$\mathcal O_X$-modules $\mathcal L_1,\ldots,\mathcal L_d$ and a coherent
$\mathcal O_X$-module $\mathcal F$ such that
$\dim(\Supp(\mathcal F))\leq d$. Put
\[
L_i=\cl^\circ(\mathcal L_i)\quad(i=1,\ldots,d),
\qquad F=\cl_\bullet(\mathcal F).
\]

% Current rescribe index 556; printed volume page 543; source-PDF page 550.
\sgaStableTarget{sga6.e10.prop.4-3-1}\textbf{Proposition 4.3.1.}
% SGA-SRC-001331
With the preceding notation, for every sequence of nonnegative integers
$\alpha_1,\ldots,\alpha_d$ such that
$\sum_{i=1}^d\alpha_i=d$, the coefficient of
\[
\binom{n_1}{\alpha_1}\cdots\binom{n_d}{\alpha_d}
\]
in the Snapper polynomial
\[
\chi(\mathcal L_1^{\otimes n_1}\otimes\cdots\otimes
\mathcal L_d^{\otimes n_d}\otimes\mathcal F)
\]
is equal to
\sgaStableTarget{sga6.e10.eq.4-3-2}\begin{equation}\tag{4.3.2}
\bigl\langle c_1(L_1)^{\alpha_1}\cdots
c_1(L_d)^{\alpha_d},\lvert F\rvert\bigr\rangle,
\end{equation}
where $\lvert F\rvert$ is the class of $F$ in $\Gr_d(X)$. Thus, in the ordinary
expansion of this polynomial in monomials in the $n_i$ ($i=1,\ldots,d$),
the coefficient of $n_1^{\alpha_1}\cdots n_d^{\alpha_d}$ is
\[
% SGA-SRC-001330
  \frac{1}{\alpha_1!\cdots\alpha_d!}
\bigl\langle c_1(L_1)^{\alpha_1}\cdots
c_1(L_d)^{\alpha_d},\lvert F\rvert\bigr\rangle.
\]

In particular, one obtains immediately:

\sgaStableTarget{sga6.e10.cor.4-3-3}\textbf{Corollary 4.3.3.}
The coefficient of the monomial $n_1n_2\cdots n_d$ in the Snapper polynomial
\[
\chi(\mathcal L_1^{\otimes n_1}\otimes\cdots\otimes
\mathcal L_d^{\otimes n_d}\otimes\mathcal F)
\]
is equal to
\sgaStableTarget{sga6.e10.eq.4-3-4}\begin{equation}\tag{4.3.4}
\bigl\langle c_1(L_1)\cdots c_1(L_d),\lvert F\rvert\bigr\rangle
=\bigl(c_1(L_1)\cdots c_1(L_d),\lvert F\rvert\bigr).
\end{equation}

\emph{Proof.}
Put $L_i=1+M_i$ for $i=1,\ldots,d$. By the definitions, the integer \hyperlink{sga6.e10.eq.4-3-2}{(4.3.2)}
is equal to
\[
\chi(M_1^{\alpha_1}\cdots M_d^{\alpha_d}F).
\]
Substituting
\[
A=K^\circ(X)/\Filt^{d+1}(X),
\qquad I=\Filt^1(X)/\Filt^{d+1}(X)
\]
in \hyperlink{sga6.e10.lem.2-1-4}{Lemma 2.1.4}, formula \hyperlink{sga6.e10.eq.2-1-5}{(2.1.5)} gives, since
$M_i\in\Filt^1(X)$ for every $i=1,\ldots,d$, the congruence
% SGA-SRC-001332
\sgaStableTarget{sga6.e10.eq.4-3-5}\begin{equation}\tag{4.3.5}
L_1^{n_1}\cdots L_d^{n_d}\equiv
\sum_{\substack{\alpha_1,\ldots,\alpha_d\geq0\\
\alpha_1+\cdots+\alpha_d\leq d}}
\binom{n_1}{\alpha_1}\cdots\binom{n_d}{\alpha_d}
M_1^{\alpha_1}\cdots M_d^{\alpha_d}
\pmod{\Filt^{d+1}(X)}.
\end{equation}
Hence the coefficient of
$\binom{n_1}{\alpha_1}\cdots\binom{n_d}{\alpha_d}$ in
$\chi(\mathcal L_1^{\otimes n_1}\otimes\cdots\otimes
\mathcal L_d^{\otimes n_d}\otimes\mathcal F)$ is
$\chi(M_1^{\alpha_1}\cdots M_d^{\alpha_d}F)$, as required.

% Current rescribe index 557; printed volume page 544; source-PDF page 551.
\sgaStableTarget{sga6.e10.sec.4.4}\subsection*{4.4}
The integer \hyperlink{sga6.e10.eq.4-3-4}{(4.3.4)} is sometimes called the \emph{intersection number} of the
invertible $\mathcal O_X$-modules $\mathcal L_1,\ldots,\mathcal L_d$ with
$\mathcal F$, and is denoted by
\[
(\mathcal L_1,\ldots,\mathcal L_d;\mathcal F)
\qquad(\emph{cf.}~[3]).
\]
Expression \hyperlink{sga6.e10.eq.4-3-4}{(4.3.4)} immediately gives the following properties.

\begin{enumerate}[label=\arabic*)]
\item The integer \hyperlink{sga6.e10.eq.4-3-4}{(4.3.4)} is a symmetric $d$-linear function of
  the classes $L_1,\ldots,L_d\in\Pic(X)$, is zero if
  $\dim(\Supp(\mathcal F))<d$, and is additive in $\mathcal F$, that is, in
  its class $\cl_\bullet(\mathcal F)=F$.

\item By \hyperlink{sga6.e10.prop.1-1-2}{Proposition 1.1.2}, one obtains
\sgaStableTarget{sga6.e10.eq.4-4-1}\begin{equation}\tag{4.4.1}
(\mathcal L_1,\ldots,\mathcal L_d;\mathcal F)=
\sum_{x\in X(d)}
\operatorname{length}_{\mathcal O_{X,x}}(\mathcal F_x)
(\mathcal L_1,\ldots,\mathcal L_d;\mathcal O_{Z(x)}).
\end{equation}

\item Let $Y$ be a closed subscheme of $X$, and let $\mathcal G$ be a coherent
  $\mathcal O_Y$-module such that $\mathcal F=f_*(\mathcal G)$, where
  $f:Y\to X$ is the
  canonical immersion. Then
\[
(f^*\mathcal L_1,\ldots,f^*\mathcal L_d;\mathcal G)
=(\mathcal L_1,\ldots,\mathcal L_d;\mathcal F).
\]
More generally, the projection formula gives
\sgaStableTarget{sga6.e10.eq.4-4-2}\begin{equation}\tag{4.4.2}
\chi(f^*(x)\cdot y)
=\chi\bigl(f_*(f^*(x)\cdot y)\bigr)
=\chi(x\cdot f_*(y))
\qquad(x\in K^\circ(X),\ y\in K_\bullet(Y))
\end{equation}
for every $k$-morphism $f:Y\to X$ of proper $k$-schemes.

\item Suppose that, in the morphism $f:Y\to X$ above, $X$ and $Y$ are
  irreducible of dimension $\leq d$. If $f$ has degree $n$, then

% Current rescribe index 558; printed volume page 545; source-PDF page 552.
\sgaStableTarget{sga6.e10.eq.4-4-3}\begin{equation}\tag{4.4.3}
(f^*\mathcal L_1,\ldots,f^*\mathcal L_d;\mathcal O_Y)
=n(\mathcal L_1,\ldots,\mathcal L_d;\mathcal O_X).
\end{equation}
Recall the definition of the degree $n$ of $f$ [3]: it is zero unless
\[
\dim(X)=\dim(Y)=\dim(f(Y)).
\]
In the latter case, one puts
\sgaStableTarget{sga6.e10.eq.4-4-4}\begin{equation}\tag{4.4.4}
n=
\frac{\operatorname{length}_{\mathcal O_{X,x}}(\mathcal O_{Y,y})}
{\operatorname{length}_{\mathcal O_{X,x}}(\mathcal O_{X,x})},
\end{equation}
where $x$ is the generic point of $X$ and $y$ is the generic point of $Y$.
\end{enumerate}

In the proof of \hyperlink{sga6.e10.eq.4-4-3}{(4.4.3)}, the case $\dim(X)=\dim(Y)$ is the only nontrivial
one. If $\dim(f(Y))<d$, then
\sgaStableTarget{sga6.e10.eq.4-4-5}\begin{equation}\tag{4.4.5}
R^if_*(\mathcal O_Y)=0
\end{equation}
for $i>0$ in a neighborhood of $x$. Condition \hyperlink{sga6.e10.eq.4-4-5}{(4.4.5)} also holds when
$\dim(f(Y))=d$, since the generic fiber then has dimension zero (EGA III
4.2.2 and EGA IV 5.5.6). Proposition 1.1.2 therefore gives the congruences
\[
f_*(\cl_\bullet(\mathcal O_Y))
\equiv\cl_\bullet(f_*\mathcal O_Y)
\equiv\operatorname{length}_{\mathcal O_{X,x}}(\mathcal O_{Y,y})
\cl_\bullet(\mathcal O_{X_{\red}})
\pmod{\Filt_{d-1}(X)}.
\]
By \hyperlink{sga6.e10.eq.4-4-2}{(4.4.2)}, it follows that
\[
(f^*\mathcal L_1,\ldots,f^*\mathcal L_d;\mathcal O_Y)
=\operatorname{length}_{\mathcal O_{X,x}}(\mathcal O_{Y,y})
(\mathcal L_1,\ldots,\mathcal L_d;\mathcal O_{X_{\red}}).
\]
On the other hand, item 2 gives
\[
(\mathcal L_1,\ldots,\mathcal L_d;\mathcal O_X)
=\operatorname{length}_{\mathcal O_{X,x}}(\mathcal O_{X,x})
(\mathcal L_1,\ldots,\mathcal L_d;\mathcal O_{X_{\red}}),
\]
which proves \hyperlink{sga6.e10.eq.4-4-3}{(4.4.3)}.

% Current rescribe index 559; printed volume page 546; source-PDF page 553.
\sgaStableTarget{sga6.e10.sec.4.5}\subsection*{4.5}
The intersection product $\langle\, ,\,\rangle$ defines an equivalence
relation called \emph{numerical equivalence}.

\sgaStableTarget{sga6.e10.def.4-5-1}\textbf{Definition 4.5.1.}
Two elements $x$ and $y$ of $\Gr^i(X)$ ($i\in\mathbb Z$), respectively of
$\Gr_i(X)$, are said to be \emph{numerically equivalent} if they define the
same homomorphism
\[
\Gr_i(X)\longrightarrow\mathbb Z,
\qquad\text{respectively}\qquad
\Gr^i(X)\longrightarrow\mathbb Z.
\]
Two invertible $\mathcal O_X$-modules are said to be numerically equivalent
if their first Chern classes in $\Gr^1(X)$ are numerically equivalent.

\sgaStableTarget{sga6.e10.prop.4-5-2}\textbf{Proposition 4.5.2.}
An invertible $\mathcal O_X$-module $\mathcal L$ is numerically equivalent to
zero if and only if
\sgaStableTarget{sga6.e10.eq.4-5-3}\begin{equation}\tag{4.5.3}
\chi(\mathcal L^{\otimes n}\otimes_{\mathcal O_X}\mathcal O_Z)
=\chi(\mathcal O_Z)
\end{equation}
for every $n\in\mathbb Z$ and every integral closed curve $Z\subset X$.

This follows from Propositions \hyperlink{sga6.e10.prop.1-1-2}{1.1.2} and 4.3.1.

Let $\Pic^N(X)$ denote the subgroup of $\Pic(X)$ consisting of the classes of
invertible $\mathcal O_X$-modules numerically equivalent to zero.
% SGA-SRC-001335
\hyperlink{sga6.e10.prop.2-4-4}{Proposition 2.4.4} gives:

% SGA-SRC-001336
\sgaStableTarget{sga6.e10.cor.4-5-4}\textbf{Corollary 4.5.4.}
\[
\Pic^\tau(X)\subset\Pic^N(X).
\]
It will be proved in \hyperlink{sga6.e13.expose}{Expos\'e XIII} that in fact
\[
\Pic^N(X)=\Pic^\tau(X).
\]

\sgaStableTarget{sga6.e10.sec.5}\section*{5. The isomorphism $\Pic(X)\simeq\Gr^1(X)$}

\sgaStableTarget{sga6.e10.sec.5.1}\subsection*{5.1}
Let $X$ be an arbitrary ringed topos. The ring $K^\circ(X)$ and the
filtration $(\Filt^i(X))_{i\in\mathbb Z}$ can be defined exactly as for
schemes. Likewise, define $\Pic(X)$ to be the group of isomorphism classes of
invertible $\mathcal O_X$-modules. Let

% Current rescribe index 560; printed volume page 547; source-PDF page 554.
\sgaStableTarget{sga6.e10.eq.5-1-1}\begin{equation}\tag{5.1.1}
i:\Pic(X)\longrightarrow K^\circ(X)
\end{equation}
be the map sending the class $\cl_P(\mathcal L)\in\Pic(X)$ of an invertible
$\mathcal O_X$-module $\mathcal L$ to
$\cl^\circ(\mathcal L)\in K^\circ(X)$. This map is plainly a homomorphism
from $\Pic(X)$ to the multiplicative group $K^\circ(X)^*$ of units of
$K^\circ(X)$. Henceforth the group law of $\Pic(X)$ is written
multiplicatively.

The map $i$ has a left inverse defined by the determinant functor. If
$\mathcal E$ is a locally free $\mathcal O_X$-module of rank $n$, put
\[
\det(\mathcal E)=\bigwedge\nolimits^n\mathcal E.
\]
If
\sgaStableTarget{sga6.e10.eq.5-1-2}\begin{equation}\tag{5.1.2}
0\longrightarrow\mathcal E'\longrightarrow\mathcal E
\longrightarrow\mathcal E''\longrightarrow0
\end{equation}
is an exact sequence of locally free $\mathcal O_X$-modules, there is a
canonical isomorphism
\[
\det(\mathcal E)\xrightarrow{\sim}
\det(\mathcal E')\otimes_{\mathcal O_X}\det(\mathcal E''),
\]
arising from \hyperlink{sga6.e10.eq.6-2-9}{(6.2.9)} and the fact that \hyperlink{sga6.e10.eq.5-1-2}{(5.1.2)} is locally split. One therefore
obtains a homomorphism of abelian groups
\sgaStableTarget{sga6.e10.eq.5-1-3}\begin{equation}\tag{5.1.3}
\det:K^\circ(X)\longrightarrow\Pic(X)
\end{equation}
which is the required left inverse of $i$. In particular, $i$ is injective,
and henceforth $\Pic(X)$ is identified with its image in $K^\circ(X)^*$.

\sgaStableTarget{sga6.e10.sec.5.2}\subsection*{5.2}
Define another map
\sgaStableTarget{sga6.e10.eq.5-2-1}\begin{equation}\tag{5.2.1}
\delta:\Pic(X)\longrightarrow\Filt^1(X)
\end{equation}

% Current rescribe index 561; printed volume page 548; source-PDF page 555.
by putting, for every $L\in\Pic(X)$,
\[
\delta(L)=1-L^{-1}.
\]

\sgaStableTarget{sga6.e10.rem.5-2-2}\textbf{Remarks 5.2.2.}
\begin{enumerate}[label=\alph*)]
% SGA-SRC-001337
\item Suppose that $X$ is the locally ringed topos of Zariski sheaves on a
scheme $S$. If $D$ is an effective Cartier divisor on $S$ and
\[
L=\cl^\circ(\mathcal O_S(D)),
\]
then
\sgaStableTarget{sga6.e10.eq.5-2-3}\begin{equation}\tag{5.2.3}
\Theta_S(\delta(L))=\cl_\bullet(\mathcal O_D),
\end{equation}
where $\Theta_S:K^\circ(X)\to K_\bullet(S)$ is the canonical homomorphism of
1.2.3. This follows immediately from the exact sequence
\[
0\longrightarrow\mathcal O_S(-D)\longrightarrow\mathcal O_S
\longrightarrow\mathcal O_D\longrightarrow0.
\]

\item Under the hypotheses of \hyperlink{sga6.e10.sec.5.1}{5.1} and \hyperlink{sga6.e10.sec.5.2}{5.2}, one has, for every
$L\in\Pic(X)$,
\sgaStableTarget{sga6.e10.eq.5-2-4}\begin{equation}\tag{5.2.4}
\delta(L)\equiv L-1\pmod{\Filt^2(X)},
\end{equation}
which justifies the notation $c_1(L)$ introduced below. Indeed,
\[
(L-1)-(1-L^{-1})=(L-1)(1-L^{-1})\in\Filt^2(X).
\]
\end{enumerate}

\sgaStableTarget{sga6.e10.sec.5.3}\subsection*{5.3}
Composing $\delta$ with the canonical epimorphism
\[
\Filt^1(X)\longrightarrow\Gr^1(X)
\]
gives the \emph{first Chern class} on $\Pic(X)$:
\sgaStableTarget{sga6.e10.eq.5-3-1}\begin{equation}\tag{5.3.1}
c_1:\Pic(X)\longrightarrow\Gr^1(X).
\end{equation}

% Current rescribe index 562; printed volume page 549; source-PDF page 556.
\sgaStableTarget{sga6.e10.thm.5-3-2}\textbf{Theorem 5.3.2.}
The map $c_1$ in \hyperlink{sga6.e10.eq.5-3-1}{(5.3.1)} is an injective homomorphism, and the homomorphism
\[
\det:\Filt^1(X)\longrightarrow\Pic(X)
\]
induces on passing to the quotient a left inverse
\[
\alpha:\Gr^1(X)\longrightarrow\Pic(X)
\]
of $c_1$. If every section $n$ of the constant sheaf $\mathbb Z_X$ is
bounded---for example, if $X$ is quasi-compact, or if the final object $e$ of
$X$ is a finite disjoint sum of connected objects---then $c_1$ is an
isomorphism and $\alpha$ is its inverse.

\emph{Proof.}
The map $c_1$ is a homomorphism, since for all $L,L'\in\Pic(X)$ one has
\[
(1-(LL')^{-1})-(1-L^{-1})-(1-{L'}^{-1})
=-(1-L^{-1})(1-{L'}^{-1})\in\Filt^2(X).
\]
For the moment, admit that
\[
\det(\Filt^2(X))=\{1\}\subset\Pic(X),
\]
so that $\alpha$ is well defined. Then
\[
\alpha\circ c_1=\id_{\Pic(X)},
\]
because, by \hyperlink{sga6.e10.rem.5-2-2}{5.2.2} b),
\[
\det(1-L^{-1})=\det(L^{-1})^{-1}=L
\]
for every $L\in\Pic(X)$. Next we show that
\[
c_1\circ\alpha=\id_{\Gr^1(X)}
\]
when every section $n$ of $\mathbb Z_X$ is bounded. Since $\Filt^1(X)$ is
the abelian group generated by the elements $\cl^\circ(\mathcal E)-n$, where
$\mathcal E$ is a locally free $\mathcal O_X$-module of rank $n$, it is
enough, in view of \hyperlink{sga6.e10.rem.5-2-2}{5.2.2} b), to show that
% Current rescribe index 563; printed volume page 550; source-PDF page 557.
\sgaStableTarget{sga6.e10.eq.5-3-3}\begin{equation}\tag{5.3.3}
\det(\cl^\circ(\mathcal E))-\cl^\circ(\mathcal E)+n-1\in\Filt^2(X).
\end{equation}
Since $n$ is bounded, the final object $e$ of $X$ decomposes as a finite
disjoint sum $\coprod_{i=1}^p e_i$ of subobjects such that $n$ is constant on
each subtopos $X/e_i$. Since the functor
\[
Y\longmapsto K^\circ(X/Y)\qquad(Y\in\Ob(X))
\]
takes finite disjoint sums to finite direct products, one may suppose that
$n$ is constant on $X$, say $n\in\mathbb Z$. Put
$E=\cl^\circ(\mathcal E)$. Then (cf. \hyperlink{sga6.e05.expose}{Expos\'e V}):
\[
\det(E)=\lambda^n(E)=\gamma^n((E-n)+1)
=\sum_{0\leq i\leq n}\gamma^i(E-n).
\]
This proves \hyperlink{sga6.e10.eq.5-3-3}{(5.3.3)}, since its left-hand side is
\[
\sum_{2\leq i\leq n}\gamma^i(E-n).
\]

It remains to prove the following lemma.

\sgaStableTarget{sga6.e10.lem.5-3-4}\textbf{Lemma 5.3.4.}
For every ringed topos $X$ and every $x\in\Filt^2(X)$,
\[
\det(x)=1.
\]

\emph{Proof.}
Let $\varepsilon:K^\circ(X)\to\Gamma(\mathbb Z_X)$ be the augmentation
homomorphism. Formula \hyperlink{sga6.e10.eq.6-2-8}{(6.2.8)}, together with \hyperlink{sga6.e10.para.6-5-3}{6.5.3} below, gives by linearity
\sgaStableTarget{sga6.e10.eq.5-3-5}\begin{equation}\tag{5.3.5}
\det(xy)=\det(x)^{\varepsilon(y)}\det(y)^{\varepsilon(x)}
\qquad(x,y\in K^\circ(X)).
\end{equation}

% Current rescribe index 564; printed volume page 551; source-PDF page 558.
It is therefore enough to show that, for every $k>1$,
\[
\det\bigl(\gamma^k(E-n)\bigr)=1,
\]
where, as above, $E=\cl^\circ(\mathcal E)$ and
$n=\varepsilon(E)$. For every $x\in K^\circ(X)$,
\[
\gamma^k(x)=\lambda^k(x+k-1),
\]
and hence
\[
\begin{aligned}
% SGA-SRC-001340
\det\bigl(\gamma^k(E-n)\bigr)
&=\det\bigl(\lambda^k(E-n+k-1)\bigr)\\
&=\det\left(
\sum_{p+q=k}\binom{-n+k-1}{p}\lambda^q(E)
\right).
\end{aligned}
\]
But \hyperlink{sga6.e10.eq.6-2-10}{(6.2.10)} and \hyperlink{sga6.e10.para.6-5-3}{6.5.3} give
\[
\det\bigl(\lambda^q(E)\bigr)
=\det(E)^{\binom{n-1}{q-1}}
\qquad(q>0).
\]
It follows that
\[
\det\bigl(\gamma^k(E-n)\bigr)=\det(E)^r,
\]
where
\[
r=\sum_{\substack{p+q=k\\q\geq1}}
\binom{-n+k-1}{p}\binom{n-1}{q-1}.
\]
Now $r$ is a section of $\mathbb Z_X$, everywhere equal to the coefficient of
$t^{k-1}$ in $(1+t)^{k-2}$, and is therefore zero. This proves the lemma.

\sgaStableTarget{sga6.e10.sec.6}\section*{6. Appendix: Calculation of Determinants of Locally Free Sheaves}

\sgaStableTarget{sga6.e10.sec.6.1}\subsection*{6.1}
First let $X=(\mathcal C,\mathcal O)$ be an arbitrary ringed topos. For every
positive section $n\in\Gamma(X,\mathbb Z_X)$ of the constant sheaf and every
$\mathcal O$-module $M\in\Mod_{\mathcal O}(X)$, one can define the
$\mathcal O$-modules

% Current rescribe index 565; printed volume page 552; source-PDF page 559.
\sgaStableTarget{sga6.e10.eq.6-1-1}\begin{equation}\tag{6.1.1}
M^n,\qquad \bigotimes\nolimits^n M,\qquad
\bigwedge\nolimits^n M,\quad\text{etc.}
\end{equation}
The functors $(\ )^n$, $\bigotimes^n$, and $\bigwedge^n$ carry free modules,
respectively locally free modules, to free modules, respectively locally free
modules.

More generally, let $n_1,\ldots,n_r$ be sections of $\mathbb Z_X$ and let
$T$ be an $r$-variable multifunctor, covariant for definiteness,
\[
T:\Mod_{\mathcal O}^{n_1}(X)\times\cdots\times
\Mod_{\mathcal O}^{n_r}(X)\longrightarrow\Mod_{\mathcal O}(X),
\]
where $\Mod_{\mathcal O}^n(X)$ denotes the category fibered over $X$ of
locally free $\mathcal O$-modules of rank $n$. Put
\[
M=T(\mathcal O^{n_1},\ldots,\mathcal O^{n_r})
\]
and
\[
G=\GL(n_1)\times\cdots\times\GL(n_r),
\qquad
\GL(n_i)=\Aut_{\mathcal O}(\mathcal O^{n_i}).
\]
It follows immediately that $T$ defines a representation
\sgaStableTarget{sga6.e10.eq.6-1-2}\begin{equation}\tag{6.1.2}
\alpha_T:G\longrightarrow\Aut_{\mathcal O}(M).
\end{equation}
If $M$ is itself locally free of rank $m$, then for every composite functor
$T'\circ T$, where
\[
T':\Mod_{\mathcal O}^m(X)\longrightarrow\Mod_{\mathcal O}(X),
\]
one has
\sgaStableTarget{sga6.e10.eq.6-1-3}\begin{equation}\tag{6.1.3}
\alpha_{T'\circ T}=\alpha_{T'}\circ\alpha_T.
\end{equation}

Conversely, from a representation
$\alpha=\alpha_T:G\to\Aut_{\mathcal O}(M)$ one recovers
$T(E_1,\ldots,E_r)$

% Current rescribe index 566; printed volume page 553; source-PDF page 560.
for every sequence $E_1,\ldots,E_r$ of locally free $\mathcal O$-modules of
ranks $n_1,\ldots,n_r$, by gluing the free module
$T(\mathcal O^{n_1},\ldots,\mathcal O^{n_r})$ on local trivializing charts
through $\alpha$. More precisely, the sequence $E_1,\ldots,E_r$ determines a
right $G$-torsor $P=P_{E_1,\ldots,E_r}$ and canonical isomorphisms
\sgaStableTarget{sga6.e10.eq.6-1-3.occ2}\begin{equation}\tag{6.1.3}
E_i\xrightarrow{\sim}P\mathbin{\wedge^G}\mathcal O^{n_i}
\qquad(i=1,\ldots,r),
\end{equation}
where $\wedge^G$ denotes the contracted product and $G$ acts on
$\mathcal O^{n_i}$ through the $i$-th projection $G\to\GL(n_i)$. Every
representation
\[
\alpha:G\longrightarrow\Aut_{\mathcal O}(M)
\]
therefore defines a multifunctor $T=T_\alpha$ such that
\sgaStableTarget{sga6.e10.eq.6-1-4}\begin{equation}\tag{6.1.4}
T_\alpha(E_1,\ldots,E_r)
=P_{E_1,\ldots,E_r}\mathbin{\wedge^G}M.
\end{equation}
Moreover, for every composite representation
$\alpha=\alpha_1\circ\alpha_2$, one has
\sgaStableTarget{sga6.e10.eq.6-1-5}\begin{equation}\tag{6.1.5}
T_\alpha=T_{\alpha_1}\circ T_{\alpha_2}.
\end{equation}

\sgaStableTarget{sga6.e10.sec.6.2}\subsection*{6.2}
For every locally free $\mathcal O$-module $E$ of rank $n$, write
\[
\det(E)=\bigwedge\nolimits^n E.
\]
With the notation of \hyperlink{sga6.e10.sec.6.1}{6.1}, the functor $\det(T(E_1,\ldots,E_r))$ is
determined, up to canonical isomorphism, by the corresponding representation
\sgaStableTarget{sga6.e10.eq.6-2-1}\begin{equation}\tag{6.2.1}
\det\circ\alpha_T:G\longrightarrow\Gm=\GL(1).
\end{equation}

Suppose now that $X$ is the ringed topos of sheaves on the ringed site
$(\Sch/S)$ endowed with the \'etale topology. Then
$\det(T(E_1,\ldots,E_r))$, or equivalently the representation
$\det\circ\alpha_T$, can be calculated by means of the following proposition.

% Current rescribe index 567; printed volume page 554; source-PDF page 561.
\sgaStableTarget{sga6.e10.prop.6-2-2}\textbf{Proposition 6.2.2.}
For every positive section $n$ of $\mathbb Z$, the \'etale sheaf of groups
$\Der(\GL(n))$ of commutators of $\GL(n)$ is represented by
\[
\SL(n)=\Ker\bigl(\det:\GL(n)\longrightarrow\Gm\bigr).
\]

\sgaStableTarget{sga6.e10.cor.6-2-3}\textbf{Corollary 6.2.3.}
Every representation
\sgaStableTarget{sga6.e10.eq.6-2-4}\begin{equation}\tag{6.2.4}
\alpha:G=\GL(n_1)\times\cdots\times\GL(n_r)\longrightarrow\Gm
\end{equation}
is of the form
\sgaStableTarget{sga6.e10.eq.6-2-5}\begin{equation}\tag{6.2.5}
\alpha(Y)(s)=\det(s_1)^{m_1}\cdots\det(s_r)^{m_r},
\end{equation}
where $Y\in\Ob(\Sch/S)$, $s=(s_1,\ldots,s_r)\in G(Y)$, and
$m_1,\ldots,m_r$ are sections of $\mathbb Z$ uniquely determined by
$\alpha$.

\sgaStableTarget{sga6.e10.cor.6-2-6}\textbf{Corollary 6.2.6.}
With the notation of \hyperlink{sga6.e10.sec.6.1}{6.1},
\sgaStableTarget{sga6.e10.eq.6-2-6}\begin{equation}\tag{6.2.6}
\det\bigl(T_\alpha(E_1,\ldots,E_r)\bigr)
=\det(E_1)^{\otimes m_1}\otimes\cdots\otimes
\det(E_r)^{\otimes m_r},
\end{equation}
where $m_1,\ldots,m_r$ are the sections of $\mathbb Z$ uniquely determined
by $\alpha$.

\sgaStableTarget{sga6.e10.rem.6-2-7}\textbf{Remark 6.2.7.}
The exponents $m_1,\ldots,m_r$ in \hyperlink{sga6.e10.eq.6-2-5}{(6.2.5)} and \hyperlink{sga6.e10.eq.6-2-6}{(6.2.6)} can be determined by
successively substituting in \hyperlink{sga6.e10.eq.6-2-5}{(6.2.5)}
\[
s_i=\lambda\,\id_{\mathcal O_Y^{n_i}}\quad(\lambda\in\Gm(Y)),
\qquad
s_j=\id_{\mathcal O_Y^{n_j}}\quad(j\ne i),
\]
for $i=1,\ldots,r$. Then
\[
\det(s_i)=\lambda^{n_i},
\qquad
\alpha(Y)(s)=\lambda^{n_i m_i}.
\]
For example, \hyperlink{sga6.e10.eq.6-2-6}{(6.2.6)} gives
\sgaStableTarget{sga6.e10.eq.6-2-8}\sgaStableTarget{sga6.e10.eq.6-2-9}\sgaStableTarget{sga6.e10.eq.6-2-10}\begin{align}
% SGA-SRC-001344
\det(E_1\otimes E_2)
&=\det(E_1)^{\otimes n_2}\otimes\det(E_2)^{\otimes n_1},
\tag{6.2.8}\\
\det(E_1\oplus E_2)
&=\det(E_1)\otimes\det(E_2),
\tag{6.2.9}\\
\det\bigl(\bigwedge\nolimits^m E_1\bigr)
&=\det(E_1)^{\otimes\binom{n_1-1}{m-1}}.
\tag{6.2.10}
\end{align}

% Current rescribe index 568; printed volume page 555; source-PDF page 562.
\sgaStableTarget{sga6.e10.sec.6.3}\subsection*{6.3}
Let us first show how \hyperlink{sga6.e10.cor.6-2-3}{6.2.3} follows from 6.2.2. It is plainly enough to prove
\hyperlink{sga6.e10.eq.6-2-5}{(6.2.5)} for
\[
s_{(i)}=(1,\ldots,s_i,\ldots,1)\qquad(i=1,\ldots,r).
\]
Thus one may suppose that $r=1$ and $G=\GL(n_1)$. By \hyperlink{sga6.e10.prop.6-2-2}{6.2.2}, the
representation $\alpha$ factors through $\GL(n_1)/\SL(n_1)$; equivalently,
there is a commutative diagram
\[
\begin{tikzcd}[column sep=large,row sep=large]
\GL(n_1) \arrow[r,"\alpha"] \arrow[d,"j"'] & \Gm\\
\GL(n_1)/\SL(n_1) \arrow[ur,"\beta"'] &
\end{tikzcd}
\]
in which $j$ is the canonical epimorphism and $\beta$ is a representation
of $\GL(n_1)/\SL(n_1)$ in $\Gm$. The homomorphism
\[
\det:\GL(n_1)\longrightarrow\Gm
\]
has a section. Consequently, the monomorphism
\[
\GL(n_1)/\SL(n_1)\longrightarrow\Gm
\]
induced by $\det$ is an isomorphism. On the other hand, it is known
(SGA 3 VIII 1.3) that every representation of $\Gm$ in itself is of the
form
\[
s\longmapsto s^m,
\]
where $m$ is a section of $\mathbb Z$. The assertion follows.

\hyperlink{sga6.e10.cor.6-2-6}{Corollary 6.2.6} is simply a reformulation of 6.2.3. Notice that \hyperlink{sga6.e10.cor.6-2-3}{6.2.3} also
follows immediately from SGA 3 XXII 6.2.1, but we have wished to give an
elementary proof here by means of 6.2.2.

% Current rescribe index 569; printed volume page 556; source-PDF page 563.
\sgaStableTarget{sga6.e10.sec.6.4}\subsection*{6.4. Proof of 6.2.2}
It is enough to prove that, for every $Y\in\Ob(\Sch/S)$, every $y\in Y$,
and every $s\in\SL(n)(Y)$, there is an \'etale neighborhood $Y'\to Y$ of
$y$ such that the restriction of $s$ to $Y'$ is a product of commutators
in $\GL(n)(Y')$. Since the question is local, one may suppose that $Y$ is
affine with ring $A$ and that the section $n$ is constant on $Y$. The case
$n=1$ is trivial; hence assume $n\geq2$. Let
$\mathfrak p$ be the prime ideal corresponding to $y$, and write
$s=(s_{ij})_{i,j=1,\ldots,n}$. For every $\lambda\in A$ and every
$h,k\in\{1,\ldots,n\}$ with $h\ne k$, let $c_{h,k,\lambda}$ be the matrix
$C=(c_{ij})$ defined by
\[
c_{hk}=\lambda,\qquad
c_{ij}=\delta_{ij}\quad\text{for }(i,j)\ne(h,k),
\]
where $\delta_{ij}$ is the Kronecker symbol.

% SGA-SRC-001349 — reconstructed by mathematical inference; scan inconclusive
Thus $c_{h,k,\lambda}=I+\lambda E_{hk}$. For every matrix $a$, left
multiplication by $c_{h,k,\lambda}$ adds $\lambda$ times row $k$ of $a$ to
row $h$, whereas right multiplication adds $\lambda$ times column $h$ of
$a$ to column $k$. In particular,
\[
\det(c_{h,k,\lambda})=1,
\]
and
% SGA-SRC-001348
\sgaStableTarget{sga6.e10.eq.6-4-1}\begin{equation}\tag{6.4.1}
c_{h,k,\lambda}c_{h,k,\mu}=c_{h,k,\lambda+\mu}
\end{equation}
for all $\lambda,\mu\in A$. Hence one obtains a representation
\[
c_{h,k}:\mathbb G_{\mathrm a}\longrightarrow\SL(n)
\qquad(h,k=1,\ldots,n,\ h\ne k).
\]
No invertibility assumption is needed in the row-reduction step. We reduce
to invertible parameters before the commutator step.

The proof has two steps. First, we show that $s$ is a product of matrices
$c_{h,k,\lambda}$ in a sufficiently small Zariski neighborhood of $y$.
Then we show that these matrices are products of commutators in $\GL(n)$ in
a sufficiently small \'etale neighborhood of $y$.

% Current rescribe index 570; printed volume page 557; source-PDF page 564.
For the first step, one must transform $s$ into the identity matrix
$I=(\delta_{ij})$ in a Zariski neighborhood of $y$ by multiplying on the
left and right by suitable matrices $c_{h,k,\lambda}$. We first make
$s_{11}=1$. If $s_{11}\ne1$ and $s_{21}\in\mathfrak p$, then, since
$\det(s)=1\notin\mathfrak p$, there is an index $i$ for which
$s_{i1}\notin\mathfrak p$. Replacing $s$ by $c_{2,i,1}s$ replaces
$s_{21}$ by $s_{21}+s_{i1}$, which is not in $\mathfrak p$. Thus, after
localizing, one may suppose that $s_{21}$ is invertible.

If $s_{11}-1\in\mathfrak p$, replacing $s$ by $c_{1,2,1}s$ replaces
$s_{11}-1$ by
\[
(s_{11}+s_{21})-1\notin\mathfrak p.
\]
After a further localization, one may therefore suppose that $s_{11}-1$ is
invertible. Then
\[
\lambda=\frac{1-s_{11}}{s_{21}}
\]
is invertible, and left multiplication by $c_{1,2,\lambda}$ makes the new
entry $s_{11}$ equal to $1$.

Assume now that $s_{11}=1$. Left multiplication by
$c_{i,1,-s_{i1}}$ for $i>1$ clears the remaining entries in the first
column, and right multiplication by $c_{1,i,-s_{1i}}$ clears the remaining
entries in the first row. Repeating this procedure on the lower-right block
transforms $s$, in a Zariski neighborhood $U$ of $y$, into a matrix
$s'=(s'_{ij})$ satisfying
\[
s'_{ij}=\delta_{ij}\qquad\text{whenever }i<n\text{ or }j<n.
\]
By construction, $s'=c_1sc_2$ on $U$, where $c_1$ and $c_2$ are products
of matrices $c_{h,k,\lambda}$. Hence
\[
\det(s')=\det(s)=1,
\]
so $s'_{nn}=1$ and therefore $s'=I$. Since
$c_{h,k,\lambda}^{-1}=c_{h,k,-\lambda}$, this proves that $s$ is itself a
product of elementary matrices of this form.

After shrinking about $y$, a parameter $a\notin\mathfrak p$ is a unit. If
$a\in\mathfrak p$, then $a=1+(a-1)$, and both $1$ and $a-1$ are units near
$y$. By \hyperlink{sga6.e10.eq.6-4-1}{(6.4.1)},
\[
c_{h,k,a}=c_{h,k,1}c_{h,k,a-1}
\]
Thus it remains only to treat invertible parameters.

It remains to prove that these elementary matrices are products of
commutators in an \'etale neighborhood of $y$. For invertible $\lambda$ and
$\mu$, the two matrices

% Current rescribe index 571; printed volume page 558; source-PDF page 565.
$c_{h,k,\lambda}$ and $c_{h,k,\mu}$ are conjugate. More precisely, there is
an invertible matrix $a=(a_{ij})$ such that
\[
a\,c_{h,k,\lambda}=c_{h,k,\mu}\,a.
\]
% SGA-SRC-001350
Indeed, the diagonal matrix defined by
\[
a_{hh}=\lambda^{-1},\qquad a_{kk}=\mu^{-1},\qquad
a_{ii}=1\quad(i\ne h,k)
\]
has this property; recall that $\lambda$ and $\mu$ are invertible and that
$h\ne k$. Thus, for every homomorphism
\[
\beta:\GL(n)(U)\longrightarrow B
\]
to an abelian group $B$, written additively, one has
% SGA-SRC-001351
\[
\beta(c_{h,k,\lambda})=\beta(c_{h,k,\mu}).
\]
If $\lambda+\mu$ is invertible, this equality and \hyperlink{sga6.e10.eq.6-4-1}{(6.4.1)} imply
\[
\beta(c_{h,k,\lambda})=0.
\]
It follows that $c_{h,k,\lambda}$ belongs to the commutator subgroup. If the
residue field $k(y)$ is not the prime field $\mathbb F_2=\{0,1\}$, then,
for every invertible $\lambda\in A$, there are a Zariski neighborhood
$U=\Spec(B)$ of $y$ and an invertible $\mu\in B$ such that
$\lambda+\mu$ is invertible in $B$. Hence $c_{h,k,\lambda}$ is a product
of commutators on $U$.

If $k(y)=\mathbb F_2$, then
\[
V=\Spec\bigl(A[T]/(T^2+T+1)\bigr)
\]
is an \'etale neighborhood of $y$ whose residue field is different from
$\mathbb F_2$. This completes the proof.

\sgaStableTarget{sga6.e10.sec.6.5}\subsection*{6.5. Remarks}
\sgaStableTarget{sga6.e10.para.6-5-1}\textbf{6.5.1.}
If no residue field of $S$ is isomorphic to $\mathbb F_2$, the preceding
proof shows that \hyperlink{sga6.e10.prop.6-2-2}{6.2.2} remains valid for the ringed topos of Zariski sheaves
on $S$. By contrast, already for $S=\Spec(\mathbb F_2)$ one has
\[
\SL(2,\mathbb F_2)=\GL(2,\mathbb F_2),
\]
which is different from the commutator subgroup of
$\GL(2,\mathbb F_2)$. The latter has order $3$, whereas
$\SL(2,\mathbb F_2)$ has order $6$ (cf. [1]).

\sgaStableTarget{sga6.e10.para.6-5-2}\textbf{6.5.2.}
On an arbitrary locally ringed topos $X$, it is not true in general that
every representation $\Gm\to\Gm$ is given by a

% Current rescribe index 572; printed volume page 559; source-PDF page 566.
power. For example, on the topos $(\mathrm{Ens})$ ringed by the field
$\mathbb Q$, the representation $\Gm\to\Gm$ defined by
\[
r\longmapsto |r|\qquad(r\in\mathbb Q^*)
\]
is not a power. In particular, \hyperlink{sga6.e10.cor.6-2-3}{6.2.3} does not remain valid for an arbitrary
locally ringed topos $X$ and an arbitrary representation
$\alpha:G\to\Gm$.

\sgaStableTarget{sga6.e10.para.6-5-3}\textbf{6.5.3.}
On the other hand, \hyperlink{sga6.e10.cor.6-2-3}{6.2.3} remains valid, on an arbitrary ringed topos
$X=(\mathcal C,\mathcal O)$, for the representations associated with
formulas \hyperlink{sga6.e10.eq.6-2-8}{(6.2.8)}, \hyperlink{sga6.e10.eq.6-2-9}{(6.2.9)}, and \hyperlink{sga6.e10.eq.6-2-10}{(6.2.10)}, because these representations are
\emph{algebraic} in the following sense.

Let $e$ be the final object of $X$, put
$S=\Spec(\mathcal O(e))$, and consider the functor
\[
u:\mathcal C\longrightarrow(\Sch/S)
\]
which sends $Y\in\Ob(\mathcal C)$ to the affine $S$-scheme
$\Spec(\mathcal O(Y))$. For every positive integer $n$, the presheaf
$\GL(n)_X$ on $\mathcal C$ is obtained, by composition with $u$, from the
corresponding presheaf $\GL(n)_S$ on $(\Sch/S)$. A representation on $X$
\sgaStableTarget{sga6.e10.eq.6-5-4}\begin{equation}\tag{6.5.4}
\alpha:\GL(n)\longrightarrow\Gm
\end{equation}
is said to be \emph{algebraic} if it is induced via $u$ by a representation
\sgaStableTarget{sga6.e10.eq.6-5-5}\begin{equation}\tag{6.5.5}
\alpha_S:\GL(n)_S\longrightarrow\mathbb G_{\mathrm m,S}
\end{equation}
of group presheaves on $(\Sch/S)$; that is, if
\[
\alpha(s)=\alpha_S(s)
\]
for every $Y\in\Ob(\mathcal C)$ and every
$s\in\GL(n)(Y)=\GL(n)_S(u(Y))$. The representation $\alpha$ in \hyperlink{sga6.e10.eq.6-5-4}{(6.5.4)} is
of course completely determined by $\alpha_S$. It follows immediately that
\hyperlink{sga6.e10.cor.6-2-3}{6.2.3} remains valid in an arbitrary ringed topos whenever $\alpha$ is
algebraic. In particular, formulas \hyperlink{sga6.e10.eq.6-2-8}{(6.2.8)}, \hyperlink{sga6.e10.eq.6-2-9}{(6.2.9)}, and \hyperlink{sga6.e10.eq.6-2-10}{(6.2.10)} remain
valid on every ringed topos.

% Current rescribe index 573; printed volume page 560; source-PDF page 567.
\sgaStableTarget{sga6.e10.sec.7}\section*{7. Appendix: Specialization in Intersection Theory}
\begin{center}
by A. Grothendieck
\end{center}

\sgaStableTarget{sga6.e10.lem.7-1}\textbf{Lemma 7.1.}
Let $i:Y\to X$ be a regular closed immersion (\hyperlink{sga6.e07.def.1-4}{VII 1.4}) of bounded
codimension. Then $i$ is a perfect morphism (\hyperlink{sga6.e07.prop.1-9}{VII 1.9}), so the inverse-image
homomorphism (\hyperlink{sga6.e04.para.2-12}{IV 2.12})
\[
i^*:K_\bullet(X)\longrightarrow K_\bullet(Y)
\]
is defined. Let $N$ be the conormal sheaf of $i$. It is a locally free
sheaf of finite type on $Y$ (\hyperlink{sga6.e07.def.1-6}{VII 1.6}), and hence one can form
\[
\lambda_{-1}(N)\in K^\circ(Y).
\]
Consider also the direct-image homomorphism (\hyperlink{sga6.e04.para.2-11}{IV 2.11})
\[
i_*:K_\bullet(Y)\longrightarrow K_\bullet(X).
\]
Then
\[
i^*i_*(x)=x\,\lambda_{-1}(N)\qquad(x\in K_\bullet(Y)),
\]
where the product on the right is taken with respect to the
$K^\circ(Y)$-module structure on $K_\bullet(Y)$ defined in \hyperlink{sga6.e04.para.2-10}{IV 2.10}.

The proof is, essentially without change, the same as that of the
corresponding formula for $K^\circ(X)$ and $K^\circ(Y)$ established in
\hyperlink{sga6.e07.para.2-6}{VII 2.7}.

\sgaStableTarget{sga6.e10.cor.7-2}\textbf{Corollary 7.2.}
Let $f:X\to S$ be flat, where $S$ is a trait, that is, the spectrum of a
discrete valuation ring $A$. Let $s$ be the closed point of $S$, let $t$ be
its generic point, and let $i:X_s\to X$ be the canonical inclusion. Then
$i$ has finite Tor-dimension, so (\hyperlink{sga6.e04.para.2-12}{IV 2.12}) the homomorphism
\[
i^*:K_\bullet(X)\longrightarrow K_\bullet(X_s)
\]
is defined.

% Current rescribe index 574; printed volume page 561; source-PDF page 568.
On the other hand, (\hyperlink{sga6.e04.para.2-11}{IV 2.11}) gives a homomorphism
\[
i_*:K_\bullet(X_s)\longrightarrow K_\bullet(X).
\]
The composite
\[
i^*i_*:K_\bullet(X_s)\longrightarrow K_\bullet(X)
\longrightarrow K_\bullet(X_s)
\]
is zero.

Indeed, $i:X_s\to X$ is a regular immersion of codimension $1$. Its
conormal sheaf $N$ is free of rank $1$, because $X_s$ is defined by the
global equation $f^*(g)=0$, where $g$ is a uniformizer of $A$. Thus
$\lambda_{-1}(N)=0$, and the assertion follows from 7.1.

For the rest of this appendix, suppose that $X$ is noetherian and flat over
$S$. Apply the exact homotopy sequence (\hyperlink{sga6.e09.prop.1-1}{IX 1.1}) to $X$, with the open
subscheme $X_t$ and complementary closed subscheme $X_s$:
\[
K_\bullet(X_s)\xrightarrow{i_*}K_\bullet(X)
\xrightarrow{j^*}K_\bullet(X_t)\longrightarrow0,
\]
where $j:X_t\to X$ is the open immersion. \hyperlink{sga6.e10.cor.7-2}{Corollary 7.2} gives:

\sgaStableTarget{sga6.e10.prop.7-3}\textbf{Proposition 7.3.}
With the notation of \hyperlink{sga6.e10.cor.7-2}{7.2}, and assuming $X$ noetherian, there is a unique
group homomorphism
\sgaStableTarget{sga6.e10.eq.7-3-1}\begin{equation}\tag{7.3.1}
\sigma_\bullet:K_\bullet(X_t)\longrightarrow K_\bullet(X_s),
\end{equation}
called the \emph{specialization homomorphism}, for which the diagram
\sgaStableTarget{sga6.e10.eq.7-3-2}\begin{equation}\tag{7.3.2}
\begin{tikzcd}[column sep=large,row sep=large]
K_\bullet(X) \arrow[r,"j^*"] \arrow[d,"i^*"'] &
K_\bullet(X_t) \arrow[dl,"\sigma_\bullet"]\\
K_\bullet(X_s) &
\end{tikzcd}
\end{equation}
commutes.

% Current rescribe index 575; printed volume page 562; source-PDF page 569.
Regard the two sides of \hyperlink{sga6.e10.eq.7-3-1}{(7.3.1)} as $K^\circ(X)$-modules via the ring
homomorphisms
\[
K^\circ(X)\xrightarrow{j^*}K^\circ(X_t),
\qquad
K^\circ(X)\xrightarrow{i^*}K^\circ(X_s)
\qquad\text{(\hyperlink{sga6.e04.para.2-7}{IV 2.7} b)}
\]
and via the $K^\circ(X_t)$- and $K^\circ(X_s)$-module structures on
$K_\bullet(X_t)$ and $K_\bullet(X_s)$, respectively (\hyperlink{sga6.e04.para.2-10}{IV 2.10}). Then
$\sigma_\bullet$ is a homomorphism of $K^\circ(X)$-modules.

This last assertion follows because
$i^*:K_\bullet(X)\to K_\bullet(X_s)$ is $K^\circ(X)$-linear, as is the
epimorphism $j^*:K_\bullet(X)\to K_\bullet(X_t)$, by \hyperlink{sga6.e04.eq.2-12-2}{IV 2.12.2}.

We now review other properties of the specialization homomorphism.

\sgaStableTarget{sga6.e10.sec.7.4}\subsection*{7.4. Compatibility with filtrations}
Endow the groups $K_\bullet$ with the filtration $\Filt_\bullet K_\bullet$
by the dimension of supports \hyperlink{sga6.e10.def.1-1-1}{(1.1.1)}. If $f$ is of finite type, and more
generally if $X$ is catenary, then the specialization homomorphism is
compatible with these filtrations and consequently induces a homomorphism,
again called specialization,
\sgaStableTarget{sga6.e10.eq.7-4-1}\begin{equation}\tag{7.4.1}
\sigma_\bullet:\Gr_\bullet(X_t)\longrightarrow\Gr_\bullet(X_s).
\end{equation}

To prove this, by \hyperlink{sga6.e10.cor.1-1-4}{1.1.4} one is reduced to showing that, if $Y_t$ is an
integral closed subscheme of dimension $d$ of $X_t$, then
$\sigma_\bullet(\cl_\bullet(\mathcal O_{Y_t}))$ has filtration at most
$d$. Let $Y$ be the schematic closure of $Y_t$ in $X$. Then $Y$ is flat
over $S$ (EGA IV 2.8.1), and one readily obtains
\[
i^*(\cl_\bullet(\mathcal O_Y))
=\cl_\bullet(\mathcal O_{Y_s}).
\]
It follows that
\[
\sigma_\bullet(\cl_\bullet(\mathcal O_{Y_t}))
=\cl_\bullet(\mathcal O_{Y_s}).
\]
Since $X$ is catenary, and hence $Y$ is catenary, one checks that
\[
\dim Y_s\leq\dim Y_t.
\]
When $f$ is of finite type, this follows immediately, after reducing to the
case in which $X$ is integral, from EGA IV 7.1.13. This proves the assertion.

% Current rescribe index 576; printed volume page 563; source-PDF page 570.
\sgaStableTarget{sga6.e10.sec.7.5}\subsection*{7.5. ``Conservation of number''}
Suppose now that $f$ is proper. Then $X_t$ and $X_s$ are proper over
$k(t)$ and $k(s)$, respectively, and one has the homomorphisms defined in
\hyperlink{sga6.e10.sec.2.3}{2.3},
\[
\chi_t:K_\bullet(X_t)\longrightarrow\mathbb Z,
\qquad
\chi_s:K_\bullet(X_s)\longrightarrow\mathbb Z.
\]
For every $x\in K_\bullet(X_t)$ one has
\sgaStableTarget{sga6.e10.eq.7-5-1}\begin{equation}\tag{7.5.1}
\chi_s\bigl(\sigma_\bullet(x)\bigr)=\chi_t(x).
\end{equation}
This is the \emph{conservation-of-number theorem}; equivalently, the diagram
\sgaStableTarget{sga6.e10.eq.7-5-2}\begin{equation}\tag{7.5.2}
\begin{tikzcd}[column sep=large,row sep=large]
K_\bullet(X_t) \arrow[r,"\sigma_\bullet"] \arrow[dr,"\chi_t"'] &
K_\bullet(X_s) \arrow[d,"\chi_s"]\\
& \mathbb Z
\end{tikzcd}
\end{equation}
commutes.

To prove \hyperlink{sga6.e10.eq.7-5-1}{(7.5.1)}, one is again reduced to the case
$x=\cl_\bullet(\mathcal O_{Y_t})$, where $Y_t$ is an integral closed
subscheme of $X_t$. With the notation of the proof of \hyperlink{sga6.e10.sec.7.4}{7.4}, it remains to
prove
\[
\chi(Y_s)=\chi(Y_t),
\]
which is well known (EGA III 7.9.4, or Expos\'e III 5.7).

\sgaStableTarget{sga6.e10.sec.7.6}\subsection*{7.6}
It is immediate that
\sgaStableTarget{sga6.e10.eq.7-6-1}\begin{equation}\tag{7.6.1}
\sigma_\bullet\bigl(\cl_\bullet(\mathcal O_{X_t})\bigr)
=\cl_\bullet(\mathcal O_{X_s}).
\end{equation}
Since $\sigma_\bullet$ is $K^\circ(X)$-linear, the following diagram
\sgaStableTarget{sga6.e10.eq.7-6-2}\begin{equation}\tag{7.6.2}
\begin{tikzcd}[column sep=large,row sep=large]
& K^\circ(X) \arrow[dl,"j^*"'] \arrow[dr,"i^*"] &\\
K^\circ(X_t) \arrow[d] && K^\circ(X_s) \arrow[d]\\
K_\bullet(X_t) \arrow[rr,"\sigma_\bullet"'] && K_\bullet(X_s)
\end{tikzcd}
\end{equation}
commutes.

% Current rescribe index 577; printed volume page 564; source-PDF page 571.
The vertical arrows are the canonical homomorphisms
$K^\circ\to K_\bullet$ (\hyperlink{sga6.e04.eq.2-5-1}{IV 2.5.1}).

\sgaStableTarget{sga6.e10.sec.7.7}\subsection*{7.7. Specialization for $K^\bullet$ and $\Gr^\bullet$}
When $X_t$ and $X_s$ are regular (for example, when $f$ is smooth), the two
preceding vertical arrows are isomorphisms (\hyperlink{sga6.e04.para.2-5}{IV 2.5}). The specialization
homomorphism can then be interpreted as a homomorphism
\sgaStableTarget{sga6.e10.eq.7-7-1}\begin{equation}\tag{7.7.1}
\sigma^\bullet:K^\bullet(X_t)\longrightarrow K^\bullet(X_s),
\end{equation}
which is equivalently the unique homomorphism making the diagram
\sgaStableTarget{sga6.e10.eq.7-7-2}\begin{equation}\tag{7.7.2}
\begin{tikzcd}[column sep=large,row sep=large]
K^\bullet(X) \arrow[r,"j^*"] \arrow[d,"i^*"'] &
K^\bullet(X_t) \arrow[dl,"\sigma^\bullet"]\\
K^\bullet(X_s) &
\end{tikzcd}
\end{equation}
commute.

In this diagram $i^*$ and $j^*$ are ring homomorphisms (\hyperlink{sga6.e04.para.2-7}{IV 2.7} b), and
$j^*$ is surjective. Hence \hyperlink{sga6.e10.eq.7-7-1}{(7.7.1)} is a ring homomorphism. For the same
reason, if $X$ is separated, it is even a homomorphism of $\lambda$-rings.
Indeed, a regular separated noetherian scheme is divisorial (\hyperlink{sga6.e02.cor.2-2-7-1}{II 2.2.7.1}),
so its $K^\bullet$ is canonically identified with its ``naive''
$K^\bullet$ (\hyperlink{sga6.e04.para.2-9}{IV 2.9(ii)}), which carries a natural $\lambda$-ring structure
(\hyperlink{sga6.e06.thm.3-2}{VI 3.2}). Moreover, the regularity of $X_t$ and $X_s$ implies that $X$ is
regular (EGA $0_{\mathrm{IV}}$ 17.1.8), and
\[
i^*:K^\bullet(X)\longrightarrow K^\bullet(X_s),
\qquad
j^*:K^\bullet(X)\longrightarrow K^\bullet(X_t)
\]
are $\lambda$-homomorphisms. Therefore the same is true of \hyperlink{sga6.e10.eq.7-7-1}{(7.7.1)}, which
is obtained from $i^*$ by passage to the quotient.

% Current rescribe index 578; printed volume page 565; source-PDF page 572.
Since $X$ is regular, hence normal, its irreducible components are its
connected components, and they dominate $S$. Consequently the homomorphism
\[
H^0(X,\mathbb Z)\longrightarrow H^0(X_t,\mathbb Z)
\]
induced by $X_t\to X$ is bijective. The homomorphism induced by
$X_s\to X$ can therefore also be regarded as a specialization homomorphism
\sgaStableTarget{sga6.e10.eq.7-7-3}\begin{equation}\tag{7.7.3}
H^0(X_t,\mathbb Z)\longrightarrow H^0(X_s,\mathbb Z).
\end{equation}
It follows immediately from the definitions that \hyperlink{sga6.e10.eq.7-7-1}{(7.7.1)} is compatible,
via \hyperlink{sga6.e10.eq.7-7-3}{(7.7.3)}, with the augmented-algebra structures of its source and target
over $H^0(X_t,\mathbb Z)$ and $H^0(X_s,\mathbb Z)$, respectively, defined
in \hyperlink{sga6.e04.para.2-8}{IV 2.8}.

Still assuming that $X$ is separated, \hyperlink{sga6.e10.eq.7-7-1}{(7.7.1)} induces a homomorphism of
graded rings
\sgaStableTarget{sga6.e10.eq.7-7-4}\begin{equation}\tag{7.7.4}
\sigma^\bullet:\Gr^\bullet(X_t)\longrightarrow\Gr^\bullet(X_s).
\end{equation}
Using the Chern-class construction of V 6, one obtains the commutative
diagram
\[
\begin{tikzcd}[column sep=large,row sep=large]
K^\bullet(X_t) \arrow[r,"c_{X_t}"] \arrow[d] &
\Ch\,\Gr^\bullet(X_t) \arrow[d]\\
K^\bullet(X_s) \arrow[r,"c_{X_s}"'] &
\Ch\,\Gr^\bullet(X_s).
\end{tikzcd}
\]

\sgaStableTarget{sga6.e10.sec.7.8}\subsection*{7.8. Specialization for $\Pic$}
Retain the hypotheses of \hyperlink{sga6.e10.sec.7.7}{7.7}: $X$ is separated, and $X_t$ and $X_s$ are
regular. Using the canonical isomorphism \hyperlink{sga6.e10.thm.5-3-2}{(5.3.2)}
\[
\Gr^1\simeq\Pic,
\]
the homomorphism \hyperlink{sga6.e10.eq.7-7-4}{(7.7.4)} gives a canonical homomorphism
\sgaStableTarget{sga6.e10.eq.7-8-1}\begin{equation}\tag{7.8.1}
\Pic(X_t)\longrightarrow\Pic(X_s).
\end{equation}

% Current rescribe index 579; printed volume page 566; source-PDF page 573.
An equivalent description is as follows. Since $X$ is regular, its local
rings are factorial by the Auslander--Buchsbaum theorem (EGA IV 21.11.1).
Thus every invertible sheaf $L_t$ on $X_t$ extends to an invertible sheaf
$L$ on $X$ (EGA IV 21.6.11). The image of
$\cl(L_t)\in\Pic(X_t)$ under (7.8.1) is
$\cl(L_s)\in\Pic(X_s)$, where $L_s=L|_{X_s}$.

In fact, assume only that the local rings of $X$ are factorial and that
$X_s$ is locally integral, without assuming regularity or separatedness.
Then $\cl(L_s)$ is independent of the chosen extension $L$ of $L_t$ to an
invertible sheaf on $X$, and this gives \hyperlink{sga6.e10.eq.7-8-1}{(7.8.1)} directly. To prove the
independence, it is enough to show that
\[
L_t\simeq\mathcal O_{X_t}\quad\Longrightarrow\quad
L_s\simeq\mathcal O_{X_s}.
\]
The hypothesis means that $L\simeq\mathcal O_X(D)$ for a divisor $D$ on
$X$ whose support is contained in $X_s$ (EGA IV 21.2.11). One reduces at
once to the case in which $X_s$ is integral. Then $D=nX_s$, and
$D=\operatorname{div}(g^n)$, where $g$ is a uniformizer of $A$. Hence
$\mathcal O_X(D)\simeq\mathcal O_X$, and a fortiori
$L_s\simeq\mathcal O_{X_s}$.

\sgaStableTarget{sga6.e10.sec.7.9}\subsection*{7.9. Specialization and numerical equivalence}
Suppose that $f$ is proper and that $X_s$ and $X_t$ are regular. Consider
the two specialization homomorphisms \hyperlink{sga6.e10.eq.7-4-1}{(7.4.1)} and \hyperlink{sga6.e10.eq.7-7-4}{(7.7.4)}:
\sgaStableTarget{sga6.e10.eq.7-9-1}\begin{equation}\tag{7.9.1}
\begin{gathered}
\sigma_\bullet:\Gr_\bullet(X_t)\longrightarrow\Gr_\bullet(X_s),\\
\sigma^\bullet:\Gr^\bullet(X_t)\longrightarrow\Gr^\bullet(X_s).
\end{gathered}
\end{equation}
They are induced by the homomorphism
\sgaStableTarget{sga6.e10.eq.7-9-2}\begin{equation}\tag{7.9.2}
\sigma^\bullet:K^\bullet(X_t)\longrightarrow K^\bullet(X_s).
\end{equation}

% Current rescribe index 580; printed volume page 567; source-PDF page 574.
Since \hyperlink{sga6.e10.eq.7-9-2}{(7.9.2)} is compatible with the ring structures, formula \hyperlink{sga6.e10.eq.7-5-1}{(7.5.1)}
gives
\sgaStableTarget{sga6.e10.eq.7-9-3}\begin{equation}\tag{7.9.3}
\left\langle\sigma^\bullet(y_t),\sigma_\bullet(x_t)\right\rangle
=\left\langle y_t,x_t\right\rangle,
\qquad
x_t\in\Gr_\bullet(X_t),\quad y_t\in\Gr^\bullet(X_t),
\end{equation}
where the pairings between $\Gr^\bullet$ and $\Gr_\bullet$ are those of
4.1. It follows formally that, if $x\in\Gr_\bullet(X_t)$ (respectively,
$x\in\Gr^\bullet(X_t)$) and $\sigma_\bullet(x)$ (respectively,
$\sigma^\bullet(x)$) is numerically equivalent to zero \hyperlink{sga6.e10.def.4-5-1}{(4.5.1)}, then $x$
is numerically equivalent to zero. The converse holds if \hyperlink{sga6.e10.eq.7-9-2}{(7.9.2)}, or,
equivalently, either of the homomorphisms in \hyperlink{sga6.e10.eq.7-9-1}{(7.9.1)} obtained from it by passage to
the associated graded groups, is surjective modulo torsion. This condition
occurs only in exceptional cases.

\sgaStableTarget{sga6.e10.para.7-9-4}\textbf{7.9.4.}
It is plausible, at least when $f$ is proper and smooth, that the converse
just considered always holds: if $x$ is numerically equivalent to zero,
then so is $\sigma_\bullet(x)$ (respectively, $\sigma^\bullet(x)$). This
would imply that the homomorphisms \hyperlink{sga6.e10.eq.7-9-1}{(7.9.1)} induce monomorphisms on the
corresponding groups modulo numerical equivalence,
\[
\Gr_\bullet^{\mathrm{num}}
\qquad\text{and}\qquad
\Gr_{\mathrm{num}}^\bullet.
\]
At least the preceding result shows that
$\Gr_\bullet^{\mathrm{num}}(X_t)$ is isomorphic to a quotient of a subgroup
of $\Gr_\bullet^{\mathrm{num}}(X_s)$; in particular, its rank is at most
that of $\Gr_\bullet^{\mathrm{num}}(X_s)$. The same holds for
$\Gr_{\mathrm{num}}^\bullet$, and also for $K_{\mathrm{num}}^\bullet$.

\sgaStableTarget{sga6.e10.sec.7.10}\subsection*{7.10. Compatibility with direct and inverse images}
Let
\[
g:X\longrightarrow Y
\]
be a morphism of noetherian schemes flat over $S$. Denote by
$\sigma_{X,\bullet}$ and $\sigma_{Y,\bullet}$ the specialization
homomorphisms for $X$ and $Y$, and by
$i_X,j_X$ (respectively, $i_Y,j_Y$)

% Current rescribe index 581; printed volume page 568; source-PDF page 575.
the inclusions of the special and generic fibers into the ambient schemes.
Write
\[
g_t:X_t\longrightarrow Y_t,
\qquad
g_s:X_s\longrightarrow Y_s
\]
for the induced morphisms. They give the Cartesian squares
\sgaStableTarget{sga6.e10.eq.7-10-0}\begin{equation}\tag{7.10.0}
\begin{tikzcd}[column sep=large,row sep=large]
X_t \arrow[r,"g_t"] \arrow[d,"j_X"'] & Y_t \arrow[d,"j_Y"'] &
X_s \arrow[r,"g_s"] \arrow[d,"i_X"'] & Y_s \arrow[d,"i_Y"']\\
X \arrow[r,"g"] & Y & X \arrow[r,"g"] & Y.
\end{tikzcd}
\end{equation}

First suppose that $g$ is proper. Then the diagram
\sgaStableTarget{sga6.e10.eq.7-10-1}\begin{equation}\tag{7.10.1}
\begin{tikzcd}[column sep=large,row sep=large]
K_\bullet(X_t) \arrow[r,"g_{t*}"] \arrow[d,"\sigma_{X,\bullet}"'] &
K_\bullet(Y_t) \arrow[d,"\sigma_{Y,\bullet}"]\\
K_\bullet(X_s) \arrow[r,"g_{s*}"'] & K_\bullet(Y_s)
\end{tikzcd}
\end{equation}
commutes. This generalizes \hyperlink{sga6.e10.sec.7.5}{7.5}: in the present notation take $Y=S$
(equivalently, take the scheme denoted by $X$ in \hyperlink{sga6.e10.eq.7-3-2}{(7.3.2)} to be $S$); the
specialization homomorphism for $S$ is the identity on $\mathbb Z$. If $X$
and $Y$ satisfy the conditions of \hyperlink{sga6.e10.sec.7.7}{7.7}, that is, if their fibers are
regular, then $g_t$ and $g_s$ automatically have finite Tor-dimension, and
\hyperlink{sga6.e10.eq.7-10-1}{(7.10.1)} may be read as a commutative square whose four terms are
$K^\bullet$-groups. We leave this reformulation to the reader.

Using the definitions of $\sigma_{X,\bullet}$ and $\sigma_{Y,\bullet}$,
the commutativity of \hyperlink{sga6.e10.eq.7-10-1}{(7.10.1)} reduces to that of the two front faces of the
prism
\sgaStableTarget{sga6.e10.eq.unnamed}\[
\begin{tikzcd}[column sep=large,row sep=large]
K_\bullet(X_t) \arrow[rr,"g_{t*}"] && K_\bullet(Y_t)\\
& K_\bullet(X) \arrow[ul,"j_X^*" description]
  \arrow[dl,"i_X^*" description] \arrow[rr,"g_*"] &&
K_\bullet(Y) \arrow[ul,"j_Y^*" description]
  \arrow[dl,"i_Y^*" description]\\
K_\bullet(X_s) \arrow[rr,"g_{s*}"'] && K_\bullet(Y_s).
\end{tikzcd}
\tag{*}
\]

% Current rescribe index 582; printed volume page 569; source-PDF page 576.
These two faces are the base-change squares obtained from \hyperlink{sga6.e10.eq.7-10-0}{(7.10.0)}, so IV 3
applies. It remains only to verify that they correspond to the pairs
$(X,Y_t)$ and $(X,Y_s)$ of schemes over $Y$, and that these pairs are
Tor-independent over $Y$. This is clear for the first pair because
$Y_t\to Y$ is flat. It is also immediate for the second, because
$Y_s\to Y$ is a regular immersion of codimension $1$ and remains so after
base change by $X\to Y$.

Now drop the properness assumption and instead suppose that $g$ has finite
Tor-dimension. Then $g_t$ and $g_s$ also have finite Tor-dimension. This is
immediate for $g_t$; for $g_s$ it follows from the Tor-independence of
$X$ and $Y_s$ over $Y$ just noted. Hence one can form the square
\sgaStableTarget{sga6.e10.eq.7-10-2}\begin{equation}\tag{7.10.2}
\begin{tikzcd}[column sep=large,row sep=large]
K_\bullet(Y_t) \arrow[r,"g_t^*"] \arrow[d,"\sigma_{Y,\bullet}"'] &
K_\bullet(X_t) \arrow[d,"\sigma_{X,\bullet}"]\\
K_\bullet(Y_s) \arrow[r,"g_s^*"'] & K_\bullet(X_s),
\end{tikzcd}
\end{equation}
whose horizontal arrows are those defined in \hyperlink{sga6.e04.para.2-12}{IV 2.12}. If $X$ and $Y$
satisfy the conditions of \hyperlink{sga6.e10.sec.7.7}{7.7}, that is, if their fibers are regular, then
$Y$ is regular and $g$ automatically has finite Tor-dimension; the square
may then be interpreted as a commutative square of $K^\bullet$-rings.

To prove the commutativity of \hyperlink{sga6.e10.eq.7-10-2}{(7.10.2)}, it is again enough to prove that of
the two front faces of the prism analogous to $(*)$, with the direct-image
homomorphisms $g_{t*},g_*,g_{s*}$ replaced by the inverse-image
homomorphisms $g_t^*,g^*,g_s^*$. Their commutativity follows

% Current rescribe index 583; printed volume page 570; source-PDF page 577.
immediately from that of the corresponding squares in \hyperlink{sga6.e10.eq.7-10-0}{(7.10.0)}, using the
transitivity of inverse images (\hyperlink{sga6.e04.para.2-12}{IV 2.12}).

\sgaStableTarget{sga6.e10.rem.7-10-3}\textbf{Remark 7.10.3.}
When $g$ is proper, the commutativity of \hyperlink{sga6.e10.eq.7-10-1}{(7.10.1)} immediately implies that
of the following diagram, obtained by passage to the associated graded
groups:
\sgaStableTarget{sga6.e10.eq.7-10-1bis}\begin{equation}\tag{7.10.1 bis}
\begin{tikzcd}[column sep=large,row sep=large]
\Gr_\bullet(X_t) \arrow[r,"g_{t*}"]
  \arrow[d,"\sigma_{X,\bullet}"'] &
\Gr_\bullet(Y_t) \arrow[d,"\sigma_{Y,\bullet}"]\\
\Gr_\bullet(X_s) \arrow[r,"g_{s*}"'] & \Gr_\bullet(Y_s).
\end{tikzcd}
\end{equation}
Likewise, if $X$ and $Y$ are separated and have regular fibers, the
commutativity of \hyperlink{sga6.e10.eq.7-10-2}{(7.10.2)} implies that of
\sgaStableTarget{sga6.e10.eq.7-10-2bis}\begin{equation}\tag{7.10.2 bis}
\begin{tikzcd}[column sep=large,row sep=large]
\Gr^\bullet(Y_t) \arrow[r,"g_t^*"] \arrow[d,"\sigma_Y^\bullet"'] &
\Gr^\bullet(X_t) \arrow[d,"\sigma_X^\bullet"]\\
\Gr^\bullet(Y_s) \arrow[r,"g_s^*"'] & \Gr^\bullet(X_s).
\end{tikzcd}
\end{equation}
If $X$ and $Y$ are moreover quasi-projective and smooth over $S$, then
$g_t^*$ and $g_s^*$ are compatible with the filtrations by codimension of
supports. Thus one obtains a commutative diagram of the same form as
(\hyperlink{sga6.e10.eq.7-10-2}{7.10.2} bis), with $\Gr^\bullet$ replaced by
$\Gr_{\mathrm{top}}^\bullet$ (\hyperlink{sga6.e00.sec.2.1}{Expos\'e 0, 2.1)}.

Finally, if $X$ and $Y$ are quasi-projective and separated over $S$ and have
regular fibers, then $g_t$ and $g_s$ are automatically complete-intersection
morphisms (\hyperlink{sga6.e08.def.1-1}{VIII 1.1}). Using \hyperlink{sga6.e08.prop.3-2}{VIII 3.2}, one obtains a commutative diagram
analogous to (\hyperlink{sga6.e10.eq.7-10-1}{7.10.1} bis), with $\Gr_\bullet$ replaced by
$\Gr^\bullet(-)_{\mathbb Q}$.

% Current rescribe index 584; printed volume page 571; source-PDF page 578.
\sgaStableTarget{sga6.e10.para.7-11}\textbf{7.11.}
For a scheme $T$ over an algebraically closed field $k$, one can define the
subgroup of $K^\bullet(T)$ consisting of elements algebraically equivalent to
zero as the subgroup of elements of the form
\sgaStableTarget{sga6.e10.eq.7-11-1}\begin{equation}\tag{7.11.1}
E(x_1)-E(x_0),
\end{equation}
where $x_0$ and $x_1$ are two $k$-rational points of a proper, regular,
connected algebraic curve $C$ over $k$, where
$E\in K^\bullet(T\times_k C)$, and where, for every $k$-rational point $x$ of
$C$, one sets
\[
E(x)=i_x^*(E)\in K^\bullet(T),
\qquad\text{where}\qquad
i_x:T\longrightarrow T\times_k C
\]
is the morphism $t\mapsto(t,x)$ defined by $x$. That the elements of the form
\hyperlink{sga6.e10.eq.7-11-1}{(7.11.1)} constitute a subgroup of $K^\bullet(T)$ follows by a standard
argument, using a product $C\times_k C'$, the two points
$(x_1,x'_1)$ and $(x_0,x'_0)$ of that product, and the fact that one can
always find an irreducible algebraic curve through any two points of an
irreducible quasi-projective algebraic scheme over $k$ (here
$C\times_k C'$). (A reader
unwilling to admit this fact may define algebraic equivalence by the group
generated by the elements of the form \hyperlink{sga6.e10.eq.7-11-1}{(7.11.1)}.)

The elements algebraically equivalent to zero form an ideal of
$K^\bullet(T)$, as one sees at once. This ideal is moreover stable under the
operations $\lambda^i$ whenever they are defined, that is, when one works
with the ``naive'' $K^\bullet$; it is also contained in the augmentation
ideal of $K^\bullet(T)$, as is immediate. One can therefore define a quotient
ring
\sgaStableTarget{sga6.e10.eq.7-11-2}\begin{equation}\tag{7.11.2}
K^\bullet_{\mathrm{alg}}(T),
\end{equation}
which is even a $\lambda$-ring when one works with naive $K^\bullet$-groups.

In the same way one can define a subgroup of elements algebraically
equivalent to zero in $K_\bullet(T)$, taking the element $E$ above in
$K_\bullet(T\times_k C)$ and observing that

% Current rescribe index 585; printed volume page 572; source-PDF page 579.
$i_x:T\to T\times_k C$ is a perfect morphism: it is obtained, by flat base
change along $T\times_k C\to C$, from the morphism $j_x:x\to C$, which is
perfect because $C$ is regular. One thereby obtains a quotient group
\sgaStableTarget{sga6.e10.eq.7-11-3}\begin{equation}\tag{7.11.3}
K_\bullet^{\mathrm{alg}}(T)
\end{equation}
of $K_\bullet(T)$. This group is naturally a module over
$K^\bullet_{\mathrm{alg}}(T)$, since the definitions show at once that the
product of an element of $K^\bullet(T)$ with an element of $K_\bullet(T)$ is
algebraically equivalent to zero if either factor is, using (\hyperlink{sga6.e04.eq.2-12-2}{IV 2.12.2}) for
$i_x$.

When one works with naive $K^\bullet$, define the ideal in $\Gr^\bullet(T)$
of elements algebraically equivalent to zero to be the ideal generated by the
$c^i(x)$, where $i\geq1$ and $x\in K^\bullet(T)$ is algebraically equivalent
to zero. The quotient graded ring
\sgaStableTarget{sga6.e10.eq.7-11-4}\begin{equation}\tag{7.11.4}
\Gr^\bullet_{\mathrm{alg}}(T)
\end{equation}
is then precisely the Chern ring associated with the $\lambda$-ring
$K^\bullet_{\mathrm{alg}}(T)$.

Finally, when $T$ is locally noetherian, define the graded subgroup of
$\Gr_\bullet(T)$ consisting of elements algebraically equivalent to zero by
requiring its component of dimension $i$ to be the image in $\Gr_i(T)$ of
the subgroup of $\operatorname{Fil}_iK_\bullet(T)$ formed by the elements algebraically
equivalent to zero in $K_\bullet(T)$. It follows at once that the quotient
graded group
\sgaStableTarget{sga6.e10.eq.7-11-5}\begin{equation}\tag{7.11.5}
\Gr_\bullet^{\mathrm{alg}}(T)
\end{equation}
is a module over the ring $\Gr^\bullet_{\mathrm{alg}}(T)$, by passage to the
quotient from the module structure of $\Gr_\bullet(T)$ over
$\Gr^\bullet(T)$ \hyperlink{sga6.e10.sec.3.1}{(3.1)}.

\sgaStableTarget{sga6.e10.para.7-11-6}\textbf{7.11.6.}
The preceding definitions would still make sense without assuming $k$
algebraically closed, provided that $C$ above is taken geometrically
connected over $k$. It is not clear, however, that these definitions would
be very useful. In particular, it does not appear that, for the group
$\Pic(T)$, one would

% Current rescribe index 586; printed volume page 573; source-PDF page 580.
recover the classical ``geometric'' notion of algebraic equivalence
considered in 2.4.1. We therefore prefer to say that an element of
$K^\bullet(T)$ is algebraically equivalent to zero if its image in
$K^\bullet(T\otimes_k\bar k)$ is so, where $\bar k$ is an algebraic closure
of $k$, and to proceed in the same way when defining algebraic equivalence
in $K_\bullet(T)$, $\Gr^\bullet(T)$, and $\Gr_\bullet(T)$.

\sgaStableTarget{sga6.e10.para.7-12}\textbf{7.12.}
With these definitions in place, I claim that the \emph{specialization homomorphism}
\hyperlink{sga6.e10.eq.7-3-1}{(7.3.1)}
\emph{carries elements algebraically equivalent to zero to elements
algebraically equivalent to zero}, and consequently induces, by passage to
the quotient, a homomorphism
\sgaStableTarget{sga6.e10.eq.7-12-1}\begin{equation}\tag{7.12.1}
K_\bullet^{\mathrm{alg}}(X_t)
\longrightarrow K_\bullet^{\mathrm{alg}}(X_s).
\end{equation}

To see this, one reduces immediately to the case in which $A$ is complete
and $k(s)$ is algebraically closed, using the straightforward compatibility
made explicit in 7.15. It is enough to show that $\sigma_\bullet$ carries an
element of the form \hyperlink{sga6.e10.eq.7-11-1}{(7.11.1)} in $K_\bullet(X_t)$ to an element algebraically
equivalent to zero in $K_\bullet(X_s)$. Here $T=X_t$, and accordingly we
write $C_t,E_t$ in place of $C,E$ above. Since $C_t$ is projective over
$k(t)$, EGA IV 2.8.5 gives a projective flat scheme $C$ over $S$ whose
generic fiber is $C_t$. Using Abhyankar's resolution of singularities for
arithmetic surfaces [6], one sees that $C$ can even be chosen regular.
(N.B. Flatness is preserved under resolution, by EGA IV 14.3.8.) Since $C$ is
proper over $S$, the $k(t)$-rational points $x_{0t}$ and $x_{1t}$ of $C_t$
extend to sections $x_0,x_1$ of $C$ over $S$. Since $C$ is regular, it is
well known that $C$ is smooth over $S$ at the points of $x_0(S)$ and
$x_1(S)$ (EGA IV 17.12.1(c)$\Rightarrow$(a), and 19.1.1). Consequently,

% Current rescribe index 587; printed volume page 574; source-PDF page 581.
the immersions $x_i:S\to C$ $(i=0,1)$ are regular and remain regular
immersions $x_i:s\to C_s$ after base change by $s\to S$. Note that, by the
exact homotopy sequence (\hyperlink{sga6.e09.prop.1-1}{IX 1.1}), the element
\[
E_t\in K_\bullet(X_t\times_{k(t)}C_t)
\]
comes from an element $E\in K_\bullet(X\times_SC)$. Thus the element
$E_t(x_{1t})-E_t(x_{0t})$ of $K_\bullet(X_t)$ comes from the element
$E(x_1)-E(x_0)$ of $K_\bullet(X)$, and its image in $K_\bullet(X_s)$ under
\hyperlink{sga6.e10.eq.7-3-1}{(7.3.1)} is therefore
\sgaStableTarget{sga6.e10.eq.unnamed.occ2}\[
i^*\bigl(E(x_1)-E(x_0)\bigr).
\tag{*}
\]

For a section $x$ of $C$ over $S$, of course, we set
\[
E(x)=i_x^*(E),
\qquad\text{where}\qquad
i_x:X\longrightarrow X\times_SC
\]
is defined as above. It may be viewed as obtained from $x$ by the flat base
change $X\times_SC\to C$; it is therefore a regular immersion and, a
fortiori, a perfect morphism, so that the expression $E(x)$ is well defined.
For the same reason $E_s(x_s)$ is defined. By transitivity of inverse images
(\hyperlink{sga6.e04.para.2-12}{IV 2.12}) in the following commutative diagram of perfect morphisms,
\[
\begin{tikzcd}[column sep=large,row sep=large]
X \arrow[r,"i_x"] & X\times_SC\\
X_s \arrow[u,"i"] \arrow[r,"i_{x_s}"] &
X_s\times_{k(s)}C_s \arrow[u]
\end{tikzcd}
\;,
\]
the expression $(*)$ is equal to
\sgaStableTarget{sga6.e10.eq.unnamed.occ3}\[
E_s(x_{1s})-E_s(x_{0s}).
\tag{**}
\]
Because $k(s)$ is algebraically closed, the reduced irreducible components
of $C_s$ are geometrically integral over $k(s)$, and their intersection
points are rational over $k(s)$. It follows readily that $(**)$ is
algebraically equivalent to zero in $K_\bullet(X_s)$, using the connectedness

% Current rescribe index 588; printed volume page 575; source-PDF page 582.
of $C_s$, which follows from the geometric connectedness of $C_t$.

From \hyperlink{sga6.e10.eq.7-12-1}{(7.12.1)}, using \hyperlink{sga6.e10.sec.7.4}{7.4}, one obtains a homomorphism
\sgaStableTarget{sga6.e10.eq.7-12-2}\begin{equation}\tag{7.12.2}
\Gr_\bullet^{\mathrm{alg}}(X_t)
\longrightarrow \Gr_\bullet^{\mathrm{alg}}(X_s).
\end{equation}
If $X_s$ and $X_t$ are moreover regular, then \hyperlink{sga6.e10.eq.7-12-1}{(7.12.1)} can also be
interpreted as a ring homomorphism
\sgaStableTarget{sga6.e10.eq.7-12-3}\begin{equation}\tag{7.12.3}
K^\bullet_{\mathrm{alg}}(X_t)
\longrightarrow K^\bullet_{\mathrm{alg}}(X_s),
\end{equation}
which, if $X$ is also separated, is a homomorphism of $\lambda$-rings and
therefore induces a homomorphism
\sgaStableTarget{sga6.e10.eq.7-12-4}\begin{equation}\tag{7.12.4}
\Gr^\bullet_{\mathrm{alg}}(X_t)
\longrightarrow \Gr^\bullet_{\mathrm{alg}}(X_s).
\end{equation}

\sgaStableTarget{sga6.e10.para.7-12-5}\textbf{7.12.5.}
Let $T$ be a smooth proper algebraic scheme over a field $k$. It is easy to
verify that an element $x\in K^\bullet(T)$ for which there is an integer
$n\geq1$ such that $nx$ is algebraically equivalent to zero is numerically
equivalent to zero. It is plausible that the converse is always true, at
least when $T$ is projective. In considering this question, one may plainly
assume that $k$ is algebraically closed and replace $K^\bullet(T)$ by
$\Gr^\bullet(T)$ or $\Gr_\bullet(T)$ without changing the conjecture. One
obtains an equivalent conjecture by working with the Chow ring $A(T)$
(cf. \hyperlink{sga6.e14.para.4-1}{XIV 4.1}), namely the more or less classical conjecture that, if $Z$ is an
algebraic cycle on $T$ numerically equivalent to zero, then there is an
integer $n\geq1$ such that $nZ$ is algebraically equivalent to zero. If this
conjecture holds for $X_t$, the compatibility of \hyperlink{sga6.e10.eq.7-3-1}{(7.3.1)} with algebraic
equivalence would imply that this specialization homomorphism, when
$f:X\to S$ is projective and smooth, is also compatible with numerical
equivalence, as conjectured in 7.9.

% Current rescribe index 589; printed volume page 576; source-PDF page 583.
\sgaStableTarget{sga6.e10.sec.7.13}\subsection*{7.13. Relations with $\ell$-adic cohomology}
Recall (SGA 5 VII) that, for every locally free module of finite type $E$ on
a scheme $T$, and every prime number $\ell$ distinct from the residue
characteristics of $T$, one defines $\ell$-adic Chern classes (for $i\geq0$)
\[
c^i(E)\in H^{2i}(T,\mathbb Z_\ell(i)),
\]
satisfying the additivity formula
\[
c(E)=c(E')c(E'')
\]
when $E$ is an extension of a locally free module $E'$ by a locally free
module $E''$. It follows from this additivity that the total Chern class
$c(E)=\sum_i c^i(E)$ arises from a homomorphism
\sgaStableTarget{sga6.e10.eq.7-13-0}\begin{equation}\tag{7.13.0}
c:K^\bullet(T)\longrightarrow
\widehat G\bigl(H^{2*}(T,\mathbb Z_\ell(*))\bigr)
=1+\prod_{i\geq1}H^{2i}(T,\mathbb Z_\ell(i)),
\end{equation}
where the target is a group under the product explained in V 1, and where
$K^\bullet(T)$ is the ``naive'' $K^\bullet$. One obtains canonical maps
\sgaStableTarget{sga6.e10.eq.7-13-1}\begin{equation}\tag{7.13.1}
c^i:K^\bullet(T)\longrightarrow H^{2i}(T,\mathbb Z_\ell(i)).
\end{equation}

Now let $X$ be a proper flat scheme over $S$, with regular fibers, and suppose
that $S$ is henselian. Assume that $\ell$ is different from the residue
characteristic of $S$. One can then form the diagram
\sgaStableTarget{sga6.e10.eq.7-13-2}\begin{equation}\tag{7.13.2}
\begin{tikzcd}[column sep=large,row sep=large]
K^\bullet(X_t) \arrow[r,"c^i"] \arrow[d,"\sigma"'] &
H^{2i}(X_t,\mathbb Z_\ell(i))\\
K^\bullet(X_s) \arrow[r,"c^i"'] &
H^{2i}(X_s,\mathbb Z_\ell(i)) \arrow[u,"\sigma^{-1}"']
\end{tikzcd}
\;,
\end{equation}
where the horizontal arrows are the Chern-class maps for $X_t$ and $X_s$,
respectively, and where $\sigma^{-1}$ is the composite homomorphism
\[
H(X_s)\xleftarrow[\sim]{i^*}H(X)\xrightarrow{j^*}H(X_t)\,;
\]
% Current rescribe index 590; printed volume page 577; source-PDF page 584.
the arrows displayed here are restriction homomorphisms, and one uses the
fact that $i^*$ is an isomorphism by the proper base change theorem
(SGA 4 XII 5.5). I claim that the square (7.13.2) is commutative. The proof is
essentially the same as the proof of the commutativity of \hyperlink{sga6.e10.eq.7-10-2}{(7.10.2)}: it is
enough to use the compatibility of the formation of Chern classes with
inverse images for the morphisms $j:X_t\to X$ and $i:X_s\to X$.

\sgaStableTarget{sga6.e10.para.7-13-3}\textbf{7.13.3.}
The square \hyperlink{sga6.e10.eq.7-13-2}{(7.13.2)} is obtained by passage to the projective limit from the
analogous squares in which cohomology with coefficients in the $\ell$-adic
sheaf $\mathbb Z_\ell(i)$ is replaced by cohomology with coefficients in the
torsion sheaves $\mu_n^{\otimes i}$, where $n$ runs through the powers of
$\ell$. Thus the commutativity of \hyperlink{sga6.e10.eq.7-13-2}{(7.13.2)} merely expresses the
commutativity of the squares involving Chern classes modulo $n$, for every
power $n$ of $\ell$.

Choose an algebraic closure $\bar K$ of $K=k(t)$, and hence an algebraic
closure $\bar k$ of $k=k(s)$, namely the residue field of the integral closure
of $A$ in $\bar K$. For every finite subextension $K_\alpha$ of $\bar K$, one
can consider the analogous commutative diagram, for fixed $n$, in the
situation obtained from $f:X\to S$ by the base change $S_\alpha\to S$, where
$S_\alpha$ is the normalization of $S$ in $K_\alpha$. By Krull--Akizuki and
the henselian hypothesis on $S$, the scheme $S_\alpha$ is indeed a trait.
Here one must moreover assume that the fibers of $X$ over $S$ are
geometrically regular, that is, that $X$ is smooth over $S$, so that this
hypothesis is preserved by base change. As $\alpha$ varies, one obtains an
inductive system of commutative squares. Passing to the inductive limit and
using SGA 4 VII 5.8 and IV 3.2, one obtains a commutative square (cf. \hyperlink{sga6.e10.sec.7.16}{7.16}):
% Current rescribe index 591; printed volume page 578; source-PDF page 585.
\[
\begin{tikzcd}[column sep=large,row sep=large]
K^\bullet(X_{\bar t}) \arrow[r,"c^i"] \arrow[d,"\sigma"'] &
H^{2i}(X_{\bar t},\mu_n^{\otimes i}) \arrow[d,"\sigma'"]\\
K^\bullet(X_{\bar s}) \arrow[r,"c^i"'] &
H^{2i}(X_{\bar s},\mu_n^{\otimes i})
\end{tikzcd}
\;,
\]
where $X_{\bar t}$ and $X_{\bar s}$ are the geometric fibers of $X$.
Moreover, the specialization homomorphism $\sigma^{-1}$ for cohomology is
now known to be an isomorphism (SGA 4 XVI 2.2). Hence the commutativity of the
preceding square can also be expressed as the commutativity of the square
obtained by reversing that arrow. Passing to the projective limit over $n$ in
the resulting commutative squares gives a commutative square
\sgaStableTarget{sga6.e10.eq.7-13-4}\begin{equation}\tag{7.13.4}
\begin{tikzcd}[column sep=large,row sep=large]
K^\bullet(X_{\bar t}) \arrow[r] \arrow[d,"\sigma"'] &
H^{2i}(X_{\bar t},\mathbb Z_\ell(i)) \arrow[d,"\sim"']\\
K^\bullet(X_{\bar s}) \arrow[r] &
H^{2i}(X_{\bar s},\mathbb Z_\ell(i))
\end{tikzcd}
\;.
\end{equation}

\sgaStableTarget{sga6.e10.para.7-13-5}\textbf{7.13.5.}
Another way to express this commutativity is to say that, if
$x_t\in K^\bullet(X_t)$ and $x_s$ is its image in $K^\bullet(X_s)$, then
$c^i(x_{\bar t})$ and $c^i(x_{\bar s})$ are the values at the geometric
points $\bar t$ and $\bar s$, respectively, of a section of
$R^{2i}f_*(\mathbb Z_\ell(i))$. This section is easily interpreted: take an
element $x\in K^\bullet(X)$ inducing $x_t$. Then
$c^i(x)\in H^{2i}(X,\mathbb Z_\ell(i))$ defines the required section of
$R^{2i}f_*(\mathbb Z_\ell(i))$.

\sgaStableTarget{sga6.e10.para.7-13-6}\textbf{7.13.6.}
Let $T$ be a scheme. By passage to the associated graded groups, the
homomorphism \hyperlink{sga6.e10.eq.7-13-0}{(7.13.0)} induces homomorphisms
\[
\psi_i:\Gr^i(T)\longrightarrow H^{2i}(T,\mathbb Z_\ell(i)),
\]
the image under $\psi_i$ of an element $x\in\Filt^iK^\bullet(T)$ being, by
definition, $c^i(x)$. (The compatibility of \hyperlink{sga6.e10.eq.7-13-0}{(7.13.0)} with the
$\lambda$-filtration on
% Current rescribe index 592; printed volume page 579; source-PDF page 586.
$K^\bullet(T)$ and the degree filtration on the target follows readily from
\hyperlink{sga6.e05.cor.6-6-4}{V 6.6.4} and from the fact that the $c^i$ define a
Chern theory on $K^\bullet(T)$ in the sense of \hyperlink{sga6.e05.def.6-13}{V 6.13}, when $K^\bullet(T)$ is
regarded as a $\lambda$-algebra augmented over the binomial ring
$H^0(T,\mathbb Z)$.) One thereby obtains maps
\sgaStableTarget{sga6.e10.eq.7-13-7}\begin{equation}\tag{7.13.7}
\operatorname{cl}:\Gr^i(T)\longrightarrow H^{2i}(T,\mathbb Q_\ell(i))
\stackrel{\mathrm{dfn}}{=}
H^{2i}(T,\mathbb Z_\ell(i))\otimes_{\mathbb Z_\ell}\mathbb Q_\ell,
\end{equation}
by the formula
\sgaStableTarget{sga6.e10.eq.7-13-8}\begin{equation}\tag{7.13.8}
\operatorname{cl}(x^{(i)})=
\frac{1}{(-1)^{i-1}(i-1)!}\,\psi_i(x^{(i)}),
\end{equation}
compare \hyperlink{sga6.e14.para.5-1}{XIV 5.1} and 5.2. The commutativity of \hyperlink{sga6.e10.eq.7-13-2}{(7.13.2)} and \hyperlink{sga6.e10.eq.7-13-4}{(7.13.4)}
immediately implies that of the analogous diagrams
\sgaStableTarget{sga6.e10.eq.7-13-9}\begin{equation}\tag{7.13.9}
\begin{tikzcd}[column sep=large,row sep=large]
\Gr^i(X_t)_{\mathbb Q} \arrow[r,"\operatorname{cl}"]
  \arrow[d,"\sigma_{\mathbb Q}"'] &
H^{2i}(X_t,\mathbb Q_\ell(i))\\
\Gr^i(X_s)_{\mathbb Q} \arrow[r,"\operatorname{cl}"'] &
H^{2i}(X_s,\mathbb Q_\ell(i)) \arrow[u,"\sigma'"]
\end{tikzcd}
\;,
\end{equation}
and
\sgaStableTarget{sga6.e10.eq.7-13-10}\begin{equation}\tag{7.13.10}
\begin{tikzcd}[column sep=large,row sep=large]
\Gr^i(X_{\bar t})_{\mathbb Q} \arrow[r,"\operatorname{cl}"] \arrow[d] &
H^{2i}(X_{\bar t},\mathbb Q_\ell(i)) \arrow[d,"\sim"]\\
\Gr^i(X_{\bar s})_{\mathbb Q} \arrow[r,"\operatorname{cl}"'] &
H^{2i}(X_{\bar s},\mathbb Q_\ell(i))
\end{tikzcd}
\;.
\end{equation}

\sgaStableTarget{sga6.e10.rem.7-13-11}\textbf{Remark 7.13.11.}
The commutativity of the last two diagrams can be expressed by saying that
specialization of algebraic cycles is compatible with formation of the
associated $\ell$-adic cohomology classes. This statement is not very
satisfactory in the form just given, because we had to replace the
coefficients $\mathbb Z_\ell(i)$ by $\mathbb Q_\ell(i)$, that is, to neglect
torsion. To obtain a statement that does not neglect torsion, it seems that
% Current rescribe index 593; printed volume page 580; declared 702-page scan: ABSENT.
% Supplemental 360-dpi witness: zero-based page index 593 (one-based PDF page 594).
one will be forced to work with the Chow ring, since it is on that ring that
the homomorphism $\operatorname{cl}$ is naturally defined with values in
$H^{2*}(-,\mathbb Z_\ell(*))$; compare \hyperlink{sga6.e14.para.4-1}{XIV 4.1}. Unfortunately, no theory of
specialization for the Chow ring has yet been developed (cf. \hyperlink{sga6.e10.sec.7.14}{7.14}) that would
allow one to write commutative squares analogous to \hyperlink{sga6.e10.eq.7-13-9}{(7.13.9)} and \hyperlink{sga6.e10.eq.7-13-10}{(7.13.10)} in
terms of Chow rings and cohomology with coefficients in
$\mathbb Z_\ell(*)$. All one can say when $f:X\to S$ is smooth and proper is
that there is a commutative square
\sgaStableTarget{sga6.e10.eq.7-13-12}\begin{equation}\tag{7.13.12}
\begin{tikzcd}[column sep=large,row sep=large]
Z^i(X_t) \arrow[r,"\operatorname{cl}"] \arrow[d,"\sigma"'] &
H^{2i}(X_t,\mathbb Z_\ell(i))\\
Z^i(X_s) \arrow[r,"\operatorname{cl}"'] &
H^{2i}(X_s,\mathbb Z_\ell(i)) \arrow[u]
\end{tikzcd}
\;,
\end{equation}
analogous to \hyperlink{sga6.e10.eq.7-13-9}{(7.13.9)}, together with a corresponding square analogous to
\hyperlink{sga6.e10.eq.7-13-10}{(7.13.10)}, which we leave to the reader to write down. Here $Z^i$ denotes the
group of algebraic cycles of codimension $i$, modulo no relation at all; the
horizontal arrows are the associated-cohomology-class homomorphisms defined
in SGA 5 IV; and $\sigma$ is the specialization homomorphism for algebraic
cycles, associating with a cycle $Z_t$ the cycle $Z_s$, where $Z$ is the
schematic closure of $Z_t$ in $X$ (cf. \hyperlink{sga6.e10.sec.7.14}{7.14}). This commutativity follows from
the following fact, proved in SGA 5 IV. If $Z$ is a closed subscheme of $X$
of codimension at least $i$ on every fiber, one can associate with it a
cohomology class
\[
\operatorname{cl}(Z)\in H^{2i}(X,\mathbb Z_\ell(i))
\]
whose formation commutes with every base change and whose restriction to
$X_t$ (respectively, $X_s$) is therefore
$\operatorname{cl}(Z_t)$ (respectively, $\operatorname{cl}(Z_s)$). Of course,
once it is known that the homomorphism $\sigma$ in \hyperlink{sga6.e10.eq.7-13-12}{(7.13.12)} carries cycles
rationally equivalent to zero to cycles rationally equivalent to zero, with
$X$ projective over $S$,
% Current rescribe index 594; printed volume page 581; source-PDF page 587.
one can replace the $Z^i$ by the $A^i$ in this diagram and obtain a diagram
of the desired kind in terms of the Chow rings of $X_t$ and $X_s$.

\sgaStableTarget{sga6.e10.rem.7-13-13}\textbf{Remark 7.13.13.}
Let $T$ be a smooth projective scheme over an algebraically closed field. It
follows immediately from the formalism developed in SGA 5 IV that, if $Z$ is
a cycle of codimension $i$ on $T$ whose image in
$H^{2i}(T,\mathbb Q_\ell(i))$ is zero, then $Z$ is numerically equivalent to
zero. One of the standard conjectures on algebraic cycles asserts that the
converse is also true. By \hyperlink{sga6.e07.prop.4-11}{VII 4.11}, one obtains an equivalent conjecture by
replacing the group of algebraic cycles with $\Gr^\bullet(T)_{\mathbb Q}$,
using the \hyperlink{sga6.e10.eq.7-13-8}{definition (7.13.8)}, or with $K^\bullet(T)$, by stating that an
element of $K^\bullet(T)$ is numerically equivalent to zero if and only if
its augmentation and its Chern classes in the groups
$H^{2i}(T,\mathbb Q_\ell(i))$ vanish.

It follows immediately from the commutativity of \hyperlink{sga6.e10.eq.7-13-10}{(7.13.10)} that, if
$X_{\bar t}$ satisfies the preceding conjecture, then the specialization
homomorphism
\[
\Gr^\bullet(X_{\bar t})\longrightarrow\Gr^\bullet(X_{\bar s})
\]
(and, a fortiori, $\Gr^\bullet(X_t)\to\Gr^\bullet(X_s)$) is compatible with
numerical equivalence, as conjectured in 7.9. The conjecture used here
(numerical equivalence implies $\mathbb Q_\ell$-cohomological equivalence) is
weaker than the conjecture considered in 7.12.5.

\sgaStableTarget{sga6.e10.sec.7.14}\subsection*{7.14. Specialization homomorphism for groups of cycles and cycle classes}
If $T$ is a noetherian scheme, denote by $Z_i(T)$ (respectively, $Z^i(T)$)
the group of algebraic cycles on $T$ that are purely of dimension $i$
(respectively, purely of codimension $i$). When $f:X\to S$ is of finite type,
one defines
% Current rescribe index 595; printed volume page 582; declared 702-page scan: ABSENT.
% Supplemental 360-dpi witness: zero-based page index 595 (one-based PDF page 596).
a specialization homomorphism
\sgaStableTarget{sga6.e10.eq.7-14-1}\begin{equation}\tag{7.14.1}
\sigma_i:Z_i(X_t)\longrightarrow Z_i(X_s)
\end{equation}
as follows. For every prime cycle $Z_t$ of $X_t$, purely of dimension $i$ and
identified with an integral subscheme of $X_t$, define $\sigma_i(Z_t)$ to be
the cycle associated with the closed subscheme $Z_s$ of $X_s$, where $Z$ is
the schematic closure of $Z_t$ in $X$. This $Z$ is a closed subscheme of $X$
flat over $S$ (EGA IV 2.8), with generic fiber equal to the original $Z_t$,
and it satisfies $\dim Z_s=\dim Z_t$ (EGA IV 7.1.13). When $X_s$ is locally
equidimensional, that is, when all irreducible components passing through a
given point have the same dimension, the specialization homomorphism just
defined immediately induces homomorphisms
\sgaStableTarget{sga6.e10.eq.7-14-1bis}\begin{equation}\tag{7.14.1 bis}
\sigma^i:Z^i(X_t)\longrightarrow Z^i(X_s).
\end{equation}

Now suppose that $f$ is quasi-projective, so that rational equivalence of
cycles is available on $X_t$ and $X_s$ (cf. [5], where the theory is
developed over an algebraically closed field). It seems very plausible that
\hyperlink{sga6.e10.eq.7-14-1}{(7.14.1)} carries cycles rationally equivalent to zero to cycles rationally
equivalent to zero and therefore induces a corresponding homomorphism on the
``Chow'' groups of cycle classes,
\sgaStableTarget{sga6.e10.eq.7-14-2}\begin{equation}\tag{7.14.2}
\sigma_i:A_i(X_t)\longrightarrow A_i(X_s),
\end{equation}
or, when $X_s$ is locally equidimensional,
\sgaStableTarget{sga6.e10.eq.7-14-2bis}\begin{equation}\tag{7.14.2 bis}
\sigma^i:A^i(X_t)\longrightarrow A^i(X_s).
\end{equation}

One might try to prove this by a method paralleling that of \hyperlink{sga6.e10.para.7-12}{7.12}, but here
resolution of singularities and reduction
% Current rescribe index 596; printed volume page 583; declared 702-page scan: ABSENT.
% Supplemental 360-dpi witness: zero-based page index 596 (one-based PDF page 597).
to an algebraically closed residue field would be unnecessary, since
$C=\mathbb P^1_S$ is available as a
``parameter curve'' along which to deform cycles. A difficulty nevertheless
arises. With the notation of \hyperlink{sga6.e10.para.7-12}{7.12}, $E$ must now denote a cycle on
$\mathbb P^1_S\times_SX$ for which
$E_t(x_{1t})-E_t(x_{0t})$ is ``defined,'' but it does not follow that the
cycle $E_s(x_{1s})-E_s(x_{0s})$ is also defined. As was to be expected, one
is thus reduced to formulating a statement of the kind of Chow's moving
lemma in order to define \hyperlink{sga6.e10.eq.7-14-2}{(7.14.2)}.
Observe, however, that the existence of \hyperlink{sga6.e10.eq.7-4-1}{(7.4.1)} shows that \hyperlink{sga6.e10.eq.7-14-1}{(7.14.1)} is
compatible with the equivalence relation on cycles defined by the surjective
homomorphism
\[
Z_\bullet(T)\longrightarrow\Gr_\bullet(T),
\]
% SGA-SRC-001374
which is only slightly coarser than linear equivalence. Moreover, if $T$ is
quasi-projective and smooth over a field, this relation differs from linear
equivalence only by torsion (\hyperlink{sga6.e14.para.4-1}{XIV 4.1}, \hyperlink{sga6.e14.para.4-2}{4.2}). In that case one therefore obtains
a homomorphism analogous to \hyperlink{sga6.e10.eq.7-14-2}{(7.14.2)}, or better to (\hyperlink{sga6.e10.eq.7-14-2}{7.14.2} bis), with both
sides tensored by $\mathbb Q$.

The most natural way to define \hyperlink{sga6.e10.eq.7-14-2}{(7.14.2)} directly would seem to be to proceed
as in 7.3. This would require a theory of Chow groups $A_i$ (or of the Chow
ring in the regular case) for schemes more general than schemes
quasi-projective over a field, broad enough to include $X$, together with a
contravariant law for these groups making it possible to define the
restriction homomorphism to $X_s$,
\[
A_i(X)\longrightarrow A_i(X_s).
\]
Compare XIV 8. Once such a sufficiently general theory of Chow groups and
rings is available, it should be possible to transcribe into Chow theory all
the results stated above for the specialization homomorphism in $K$-theory.

% Current rescribe index 597; printed volume page 584; declared 702-page scan: ABSENT.
% Supplemental 360-dpi witness: zero-based page index 597 (one-based PDF page 598).
\sgaStableTarget{sga6.e10.sec.7.15}\subsection*{7.15. Varying $S$}
Let $S'$ be a second trait, and let
\[
S'\longrightarrow S
\]
be a morphism of traits, that is, a dominant local morphism, carrying the
closed point to the closed point and the generic point to the generic point.
Denote the closed point (respectively, the generic point) of $S'$ by $s'$
(respectively, $t'$), and set $X'=X\times_SS'$. Suppose that $X'$ is
noetherian, as is the case, for example, when $X$ is of finite type over $S$.
There is a commutative diagram
\[
% SGA-SRC-001375
\begin{tikzcd}[column sep=large,row sep=large]
X & X_t \arrow[l,hook,"j_X"'] & & \\
& & X' \arrow[ull] & X'_{t'} \arrow[l,hook,"j_{X'}"'] \arrow[ull] \\
X_s \arrow[uu,hook,"i_X"] & & & \\
& & X'_{s'} \arrow[uu,hook,"i_{X'}"'] \arrow[ull] &
\end{tikzcd}
\;,
\]
in which the oblique arrows are flat morphisms. One can therefore take the
inverse-image homomorphisms on $K_\bullet$ for every arrow of the diagram
(\hyperlink{sga6.e04.para.2-12}{IV 2.12}). Applying \hyperlink{sga6.e10.prop.7-3}{Definition 7.3}, one finds that the diagram
% SGA-SRC-001376
\sgaStableTarget{sga6.e10.eq.7-15-1}\begin{equation}\tag{7.15.1}
\begin{tikzcd}[column sep=large,row sep=large]
K_\bullet(X_t) \arrow[r,"\sigma_X"] \arrow[d] &
K_\bullet(X_s) \arrow[d]\\
K_\bullet(X'_{t'}) \arrow[r,"\sigma_{X'}"'] &
K_\bullet(X'_{s'})
\end{tikzcd}
\end{equation}
is commutative, where the vertical arrows are the inverse-image
homomorphisms. One immediately obtains an analogous commutative diagram with
$\Gr_\bullet$ in place of $K_\bullet$. When $X_s$ and $X_t$ are regular, one
may likewise write $K^\bullet$ and $\Gr^\bullet$ in place of $K_\bullet$ and
$\Gr_\bullet$.

% BEGIN INCLUDED SGA6_English_current_buildable_source_20260823/SGA6_Expose_X_idx598_607_English_SourceChecked_fragment.tex
% Promoted cumulative-tranche extension: current-rescribe indices 598--607.
% Source checked against the declared scan and independently against the
% validated high-resolution supplemental witness on 2026-07-17/18.

% Current rescribe index 598; printed volume page 585; source-PDF page 588.
% Supplemental witness: zero-based page index 598 (one-based PDF page 599).
\sgaStableTarget{sga6.e10.sec.7.16}\subsection*{7.16. Passage to geometric fibers}
Let $\bar K$ be an algebraic closure of $K=k(t)$, expressed as the inductive
limit of its finite subextensions $K_i$, and let
% SGA-SRC-001377
$\bar A=\varinjlim_i\bar A_i$ be the normalization of $A$ in $\bar K$, the
inductive limit of its normalizations $\bar A_i$ in the $K_i$. Choose a
maximal ideal $\bar{\mathfrak m}$ of $\bar A$; it induces a maximal ideal
$\mathfrak m_i$ on each $\bar A_i$. Set
$B_i=(\bar A_i)_{\mathfrak m_i}$. By Krull--Akizuki, $B_i$ is again a
discrete valuation ring. Denote its spectrum by $S_i$ and its closed and
generic points by $s_i$ and $t_i$, respectively. The residue field
$\bar k=\bar A/\bar{\mathfrak m}$ is an algebraic closure of $k=k(s)$ and is
the inductive limit of the residue fields
$k_i=\bar A_i/\mathfrak m_i=k(s_i)$ of the $B_i$. Set
$\bar t=\Spec(\bar K)$ and $\bar s=\Spec(\bar k)$.

Assume now that $X$ is of finite type over $S$. Then $X_{\bar t}$ and
$X_{\bar s}$ are of finite type over $\bar t$ and $\bar s$, respectively,
and hence are noetherian. By \hyperlink{sga6.e04.para.3-2-3}{IV 3.2.3},
\[
K_\bullet(X_{\bar t})\xleftarrow{\sim}
  \varinjlim_i K_\bullet(X_{t_i}),
\qquad
K_\bullet(X_{\bar s})\xleftarrow{\sim}
  \varinjlim_i K_\bullet(X_{s_i}).
\]
On the other hand, \hyperlink{sga6.e10.sec.7.15}{7.15} gives a homomorphism from the inductive system of the
$K_\bullet(X_{t_i})$ to that of the $K_\bullet(X_{s_i})$. Passing to the
limit gives a canonical homomorphism
\sgaStableTarget{sga6.e10.eq.7-16-1}\begin{equation}\tag{7.16.1}
K_\bullet(X_{\bar t})\longrightarrow K_\bullet(X_{\bar s}),
\end{equation}
again, as expected, called the \emph{specialization homomorphism}. By
\hyperlink{sga6.e10.sec.7.4}{7.4}, it induces a homomorphism
$\Gr_\bullet(X_{\bar t})\to\Gr_\bullet(X_{\bar s})$, using the fact (whose
verification is left to the reader) that the functor $\Gr_\bullet$, like
$K_\bullet$, commutes with the limits occurring here. Similarly, when $f$
is smooth, one may interpret \hyperlink{sga6.e10.eq.7-16-1}{(7.16.1)} as a homomorphism on $K^\bullet$; it is
then a ring homomorphism, and even a $\lambda$-ring homomorphism if $X$ is
separated. In that case, it therefore induces a homomorphism on
$\Gr^\bullet$.

% Current rescribe index 599; printed volume page 586; source-PDF page 589.
% Supplemental witness: zero-based page index 599 (one-based PDF page 600).
\sgaStableTarget{sga6.e10.para.7-16-2}\textbf{7.16.2.}
We shall not review here, for geometric fibers, the extension of all the
properties of specialization homomorphisms proved in the preceding sections
for the ``arithmetic'' fibers; they all follow immediately from the latter
by passage to the limit. Let us mention only the following interesting
consequence of 7.9.4. When $f$ is smooth and proper,
$\Gr_\bullet^{\mathrm{num}}(X_{\bar t})$ is isomorphic to a quotient of a
subgroup of $\Gr_\bullet^{\mathrm{num}}(X_{\bar s})$; in particular, its
rank is bounded above by that of the latter group. This is therefore a
semi-continuity result for the rank of the group of cycles modulo numerical
equivalence.

\sgaStableTarget{sga6.e10.sec.7.17}\subsection*{7.17. Case of an arbitrary base}
Let
\[
f:X\longrightarrow S
\]
be a flat and quasi-compact morphism of noetherian preschemes, and let $s,t$
be two points of $S$ such that $s$ is a specialization of $t$. We wish to
relate the groups $K_\bullet$, etc., for the two fibers $X_t$ and $X_s$, as
well as for the corresponding geometric fibers. For this, note (EGA II
7.1.9) that there exists a morphism
\[
g:S'\longrightarrow S,
\]
where $S'$ is a trait and $g$ carries the closed point $s'$ to $s$ and the
generic point $t'$ to $t$; one may even assume that
$k(t)\xrightarrow{\sim}k(t')$. Suppose that $X'=X\times_SS'$ is
noetherian. This is the case if $f$ is of finite type or if $S'$ is
essentially of finite type over $S$ (EGA IV 1.3.8), as one may assume when
$S$ itself is universally Japanese (EGA $0_{\mathrm{IV}}$ 23.1.1). The
specialization homomorphism for $X'$ then defines a homomorphism
\sgaStableTarget{sga6.e10.eq.7-17-1}\begin{equation}\tag{7.17.1}
\sigma_\bullet:K_\bullet(X_t)\xrightarrow{\sim}K_\bullet(X_{t'})
  \longrightarrow K_\bullet(X_{s'}).
\end{equation}
% Current rescribe index 600; printed volume page 587; source-PDF page 590.
% Supplemental witness: zero-based page index 600 (one-based PDF page 601).
Thus one obtains a specialization homomorphism on $K_\bullet(X_t)$, but with
values not in $K_\bullet(X_s)$, but in $K_\bullet(X_{s'})$, where $X_{s'}$
is obtained from $X_s$ by an extension $k(s)\to k(s')$ of the base fields
(which may be assumed to be of finite type if $S$ is universally Japanese).
Passing to geometric fibers as explained in \hyperlink{sga6.e10.sec.7.16}{7.16}, one likewise obtains a
homomorphism
\sgaStableTarget{sga6.e10.eq.7-17-2}\begin{equation}\tag{7.17.2}
\sigma_\bullet:K_\bullet(X_{\bar t})
  \longrightarrow K_\bullet(X_{\overline{s'}}).
\end{equation}
Here again, the ``geometric'' specialization homomorphism takes its values
not in $K_\bullet(X_{\bar s})$, but in
$K_\bullet(X_{\overline{s'}})$, where $X_{\overline{s'}}$ is obtained from
$X_{\bar s}$ by the extension of algebraically closed fields
$k(\bar s)\to k(\overline{s'})$.

\sgaStableTarget{sga6.e10.para.7-17-3}\textbf{7.17.3.}
In general, the homomorphism
$K_\bullet(X_{\bar s})\to K_\bullet(X_{\overline{s'}})$ is not an
isomorphism, because of the possible presence of ``continuous'' components
in $K_\bullet(X_{\bar s})$. It is, however, easy to verify that it is an
isomorphism modulo algebraic equivalence \hyperlink{sga6.e10.para.7-11}{(7.11)}, and also modulo numerical
equivalence \hyperlink{sga6.e10.def.4-5-1}{(4.5.1)} if $f$ is smooth and proper. Consequently, using \hyperlink{sga6.e10.para.7-12}{7.12},
one finds that \hyperlink{sga6.e10.eq.7-17-2}{(7.17.2)} induces a specialization homomorphism
\sgaStableTarget{sga6.e10.eq.7-17-3-1}\begin{equation}\tag{7.17.3.1}
K_\bullet^{\mathrm{alg}}(X_{\bar t})
  \longrightarrow K_\bullet^{\mathrm{alg}}(X_{\bar s});
\end{equation}
and, if $f$ is proper and smooth, provided that the specialization
homomorphism \hyperlink{sga6.e10.eq.7-17-2}{(7.17.2)} is compatible with numerical equivalence (cf. \hyperlink{sga6.e10.para.7-9-4}{7.9.4}),
one likewise obtains a homomorphism
\sgaStableTarget{sga6.e10.eq.7-17-3-2}\begin{equation}\tag{7.17.3.2}
K_\bullet^{\mathrm{num}}(X_{\bar t})
  \longrightarrow K_\bullet^{\mathrm{num}}(X_{\bar s}).
\end{equation}
We leave to the reader the task of developing, if needed, the corresponding
constructions for $\Gr_\bullet$, $K^\bullet$, and $\Gr^\bullet$, and the
properties of these specialization homomorphisms that follow immediately
from those developed in the preceding sections. Let us mention only
% Current rescribe index 601; printed volume page 588; source-PDF page 591.
% Supplemental witness: zero-based page index 601 (one-based PDF page 602).
that, by \hyperlink{sga6.e10.para.7-16-2}{7.16.2}, $\Gr_\bullet^{\mathrm{num}}(X_{\bar t})$ is isomorphic to a
quotient of a subgroup of $\Gr_\bullet^{\mathrm{num}}(X_{\bar s})$ when $f$
is proper and smooth; in particular, its rank is again bounded above by the
rank of the latter group.

\sgaStableTarget{sga6.e10.para.7-17-4}\textbf{7.17.4.}
When $S$ is regular at $s$ and $t$ is a maximal point, one can construct
$S'$ by
blowing up $s$ in the localization of $S$ at $s$, and then localizing at the
generic point of the fiber of $s$ in the blow-up. One then obtains a
\emph{pure extension} $k(s')$ of $k(s)$. Now it is not difficult to show
that, if $k'$ is a pure extension of a field $k$, then for every noetherian
scheme $T$ over $k$ the canonical homomorphism
\[
K_\bullet(T)\longrightarrow K_\bullet(T\otimes_k k')
\]
is an isomorphism: represent $k'$ as the inductive limit of the coordinate
rings of nonempty affine open subsets of schemes of the form
$\Spec k[T_1,\ldots,T_n]$, and use \hyperlink{sga6.e09.prop.1-6}{IX 1.6}. It follows that \hyperlink{sga6.e10.eq.7-17-1}{(7.17.1)} may be
viewed as a canonical homomorphism
\sgaStableTarget{sga6.e10.eq.7-17-4-1}\begin{equation}\tag{7.17.4.1}
K_\bullet(X_t)\longrightarrow K_\bullet(X_s).
\end{equation}
One should note, however, that this does not yield a homomorphism
$K_\bullet(X_{\bar t})\to K_\bullet(X_{\bar s})$: after passing to a finite
extension $K_i$ of $K=k(t)$ and normalizing $S$ in $K_i$, the regularity
hypothesis on $S$ is lost.

Another way to define \hyperlink{sga6.e10.eq.7-17-4-1}{(7.17.4.1)} is to reduce to the case in which $S$ is
local with closed point $s$ (which is immediate), and to consider an
increasing sequence of regular closed subschemes $S_i$ of $S$, with generic
points $t_i$, such that
\[
\dim S_i=i,\qquad 0\leq i\leq n=\dim S.
\]
For every $i$, $0\leq i\leq n-1$, the local ring
$\mathcal O_{S_{i+1},t_i}$ is a discrete valuation ring with residue field
$k(t_i)$ and fraction field $k(t_{i+1})$. There is therefore a
specialization homomorphism
\[
K_\bullet(X_{t_{i+1}})\longrightarrow K_\bullet(X_{t_i}),
\]
% Current rescribe index 602; printed volume page 589; source-PDF page 592.
% Supplemental witness: zero-based page index 602 (one-based PDF page 603).
and the composite of these homomorphisms for $0\leq i\leq n-1$ defines
\hyperlink{sga6.e10.eq.7-17-4-1}{(7.17.4.1)}. By an argument analogous to that of \hyperlink{sga6.e10.prop.7-3}{7.3}, one can show that this
is the same homomorphism as the one defined above by means of a blow-up. The
advantage of the second method is that one even has homomorphisms of the
form \hyperlink{sga6.e10.eq.7-16-1}{(7.16.1)}:
\[
K_\bullet(X_{\bar t_{i+1}})\longrightarrow
K_\bullet(X_{\bar t_i}),
\]
with suitable coherent choices of algebraic closures $k(\bar t_i)$ of the
$k(t_i)$. Their composite defines a homomorphism
\sgaStableTarget{sga6.e10.eq.7-17-4-2}\begin{equation}\tag{7.17.4.2}
K_\bullet(X_{\bar t})\longrightarrow K_\bullet(X_{\bar s}),
\end{equation}
which probably depends essentially on the choice of the $S_i$, that is, of
the $t_i$.

\sgaStableTarget{sga6.e10.para.7-17-5}\textbf{7.17.5.}
The construction of a specialization homomorphism \hyperlink{sga6.e10.eq.7-17-4-2}{(7.17.4.2)} can be
generalized to the case in which $S$ is no longer assumed regular at $s$
and $t$ is no longer assumed to be a maximal point, provided that resolution
of singularities is available for the integral closed subscheme $T$ of $S$
with generic point $t$, or merely for its localization at $s$. This
hypothesis holds, in particular, if $\mathcal O_{S,s}$ is excellent of
characteristic zero, by Hironaka. Indeed, one can then find a morphism
\[
S'\longrightarrow\Spec(\mathcal O_{T,s})
\]
with $S'$ regular. Taking a closed point of the fiber $S'_s$, one can apply
the preceding construction to $X'=X\times_SS'$ over $S'$, its generic point
$t'$, and its closed point $s'$. This yields a homomorphism of the form
\hyperlink{sga6.e10.eq.7-17-4-2}{(7.17.4.2)}, since $k(s')$ is a finite extension of $k(s)$ and therefore has
the same algebraic closure.

\sgaStableTarget{sga6.e10.para.7-17-6}\textbf{7.17.6.}
In all the constructions considered in \hyperlink{sga6.e10.sec.7.17}{7.17}, one has reduced the
construction of a specialization homomorphism to the
% Current rescribe index 603; printed volume page 590; source-PDF page 593.
% Supplemental witness: zero-based page index 603 (one-based PDF page 604).
case in which the base is a trait. It follows that the properties of the
specialization homomorphism reviewed in the preceding sections for a trait
as base extend automatically to the case of a more general base. Of course,
the methods of definition considered apply equally well to $\Gr_\bullet$,
to $K^\bullet$ when $f$ is smooth, and to $\Gr^\bullet$ when $f$ is also
separated.

\sgaStableTarget{sga6.e10.para.7-17-7}\textbf{7.17.7.}
Finally, note that the definition of a homomorphism \hyperlink{sga6.e10.eq.7-17-3-1}{(7.17.3.1)} or
\hyperlink{sga6.e10.eq.7-17-3-2}{(7.17.3.2)}, or of a homomorphism \hyperlink{sga6.e10.eq.7-17-4-2}{(7.17.4.2)}, depends essentially on certain
arbitrary choices: a suitable trait over $S$ in \hyperlink{sga6.e10.para.7-17-3}{7.17.3}, a sequence $(S_i)$
in \hyperlink{sga6.e10.para.7-17-4}{7.17.4}, and, in addition, the choice of a resolution of singularities in
7.17.5. The same is true of the variants for $\Gr_\bullet$, etc. Obvious
examples can be constructed with $s$, say, an ordinary double point of an
algebraic curve and $X$ projective and smooth over $S$ (for example, an
abelian scheme of relative dimension $2$), in which the specialization
homomorphism on $\Gr^1=\Pic$, and even on
$\Gr^1_{\mathrm{alg}}=\mathrm{NS}$ (the Néron--Severi group), actually
depends on the trait $S'$ chosen over $S$, or, what amounts to the same thing
here, on the choice of a point above $s$ in the normalization of $S$.

\sgaStableTarget{sga6.e10.sec.7.18}\subsection*{7.18. Appendix to the Appendix: base-field extension for certain
functors with values in commutative groups}
We spell out here a technical point raised in 7.17.4.

\sgaStableTarget{sga6.e10.para.7-18-1}\textbf{7.18.1.}
Let $k$ be a field, let $C$ be the category of integral $k$-algebras and
flat homomorphisms of such algebras, and let $F$ be a functor on $C$ with
values in the category $(\mathrm{Ens})$ of sets. Suppose that $F$ commutes
with filtered inductive limits. Given an extension $K$ of $k$, we seek
criteria
% Current rescribe index 604; printed volume page 591; source-PDF page 594.
% Supplemental witness: zero-based page index 604 (one-based PDF page 605).
for
\[
F(k)\longrightarrow F(K)
\]
to be injective (respectively, surjective).

\sgaStableTarget{sga6.e10.para.7-18-2}\textbf{7.18.2.}
Suppose that one of the following two conditions holds:
\begin{enumerate}[label=\alph*)]
\item The field $k$ is separably closed, and $K/k$ is a separable extension.
\item The field $k$ is infinite, and $K$ is a pure extension of $k$.
\end{enumerate}
Then the map $F(k)\to F(K)$ is injective.

Indeed, $K$ is the filtered inductive limit of its finite-type
$k$-subalgebras, with flat transition morphisms. Using this fact and the
commutation of $F$ with filtered inductive limits, one is reduced in case
a) to proving the following: if $A$ is an integral $k$-algebra of finite
type whose fraction field $K$ is separable over $k$, then
$F(k)\to F(A)$ is injective. The hypotheses on $k$ and $K$ imply that there
is a $k$-homomorphism $A\to k$; the composite
\[
F(k)\longrightarrow F(A)\longrightarrow F(k)
\]
is therefore the identity, and the assertion follows. In case b), passage
to the limit similarly reduces the question to an integral finite-type
$k$-algebra $A$ whose fraction field is a pure extension of $k$. Since $k$
is infinite, the hypothesis again implies the existence of a
$k$-homomorphism $A\to k$, and one concludes as above.

\sgaStableTarget{sga6.e10.para.7-18-3}\textbf{7.18.3.}
Suppose now that $F$ takes values in the category $(\mathrm{Ab})$ of
abelian groups, and suppose that for every finite extension $k'$ of $k$ a
norm homomorphism
\[
N_{k'/k}:F(k')\longrightarrow F(k)
\]
has been defined, such that the composite
\[
F(k)\longrightarrow F(k')\xrightarrow{N_{k'/k}}F(k)
\]
is multiplication by $[k':k]$. Under these conditions:
% Current rescribe index 605; printed volume page 592; source-PDF page 595.
% Supplemental witness: zero-based page index 605 (one-based PDF page 606).
\begin{enumerate}[label=(\roman*)]
\item for every extension $K$ of $k$, the kernel of
  $F(k)\to F(K)$ is a torsion group;
\item for every pure extension $K$ of $k$, the homomorphism
  $F(k)\to F(K)$ is injective. (Note that $k$ is not assumed infinite.)
\end{enumerate}

For (i), it is enough to prove that every element $x$ in the kernel of
$F(k)\to F(A)$, where $A$ is an integral $k$-algebra of finite type, is a
torsion element. Using a $k$-homomorphism $A\to k'$, where $k'$ is a finite
$k$-algebra, one reduces to the case in which $A=k'$ is a finite extension
of $k$. The norm property above then gives $[k':k]x=0$, proving that $x$ is
torsion. One even has $x=0$ when the greatest common divisor of the degrees
of the finite extensions $k'$ of $k$ that are quotients of $A$ is $1$. This
is the case when the fraction field $K$ of $A$ is a pure extension of $k$,
as is easily checked (only the case in which $k$ is finite requires a new
verification). This proves (ii).

\sgaStableTarget{sga6.e10.para.7-18-4}\textbf{7.18.4.}
Suppose that, for every nonempty affine open subset $U$ of the affine space
$\Spec k[T_1,\ldots,T_n]$, if $A$ is the ring of $U$, the homomorphism
$F(k)\to F(A)$ is surjective (compare \hyperlink{sga6.e09.prop.1-6}{IX 1.6}). Using the commutation of $F$
with filtered inductive limits, one then finds that
$F(k)\to F(K)$ is surjective for every pure extension $K$ of $k$. If, in
addition, $k$ is infinite, or if norm homomorphisms can be defined as in
\hyperlink{sga6.e10.para.7-18-3}{7.18.3}, this homomorphism is therefore bijective.

\sgaStableTarget{sga6.e10.para.7-18-5}\textbf{7.18.5.}
As an example, let $X$ be a scheme of finite type over $k$, and set
\[
F(A)=K_\bullet(X\otimes_k A),
\]
% Current rescribe index 606; printed volume page 593; source-PDF page 596.
% Supplemental witness: zero-based page index 606 (one-based PDF page 607).
which is a functor on the category $C$ by \hyperlink{sga6.e04.para.2-12}{IV 2.12}. By \hyperlink{sga6.e04.para.3-2-3}{IV 3.2.3}, it commutes
with filtered inductive limits. Moreover, if $k'$ is a finite extension of
$k$, the projection
\[
p:X\otimes_k k'\longrightarrow X
\]
is finite, flat, and of finite presentation, and therefore defines a direct
image homomorphism
\[
p_*=N_{k'/k}:K_\bullet(X\otimes_k k')\longrightarrow K_\bullet(X).
\]
The condition on the norm homomorphisms in \hyperlink{sga6.e10.para.7-18-3}{7.18.3} is immediately verified.
Finally, if $A$ is the affine algebra of an open subset $U$ of the affine
space $\Spec k[T_1,\ldots,T_n]$, then $F(k)\to F(A)$ identifies with the map
\[
K_\bullet(X)\longrightarrow K_\bullet(X\times_k U),
\]
which is surjective for $X$ noetherian, by \hyperlink{sga6.e09.prop.1-6}{IX 1.6}. The results proved above
therefore give:
\begin{enumerate}[label=(\roman*)]
\item If $K/k$ is a pure extension of $k$, then
  $K_\bullet(X)\to K_\bullet(X_K)$ is injective, and is even bijective if
  $X$ is noetherian.
\item If $k$ is separably closed and $K/k$ is a separable extension, then
  $K_\bullet(X)\to K_\bullet(X_K)$ is injective.
\item For every extension $K$ of $k$, the kernel of
  $K_\bullet(X)\to K_\bullet(X_K)$ is a torsion group.
\end{enumerate}
% SGA-SRC-001379
% SGA-SRC-001390
These results remain valid if one assumes only that $X$ is noetherian,
provided one restricts to extensions $K$ of $k$ of finite type (so that
$X_K$ is still noetherian). One obtains essentially the same results for
$K^\bullet$ in place of the functor $K_\bullet$, with the sole difference
that in (i) bijectivity cannot be asserted in general, because there is no
homotopy theorem for $K^\bullet$ in place of $K_\bullet$. To conclude
bijectivity, one would have to assume that $X$ is noetherian and regular;
then the groups $K^\bullet$ and $K_\bullet$ agree for $X$ and, respectively,
for $X_K$.

% Current rescribe index 607; printed volume page 594; source-PDF page 597.
% Supplemental witness: zero-based page index 607 (one-based PDF page 608).
One obtains essentially the same results for $\Gr^\bullet$ in place of
$K_\bullet$ or $K^\bullet$, with the same reservation concerning the
bijectivity assertion in (i). If one restricts to $\Gr^1=\Pic$, that
assertion is valid when $X$ is noetherian and normal: in that case, for every
open subset $U$ of an affine space over $k$, one sees at once that
\[
\Pic(X)\longrightarrow\Pic(X\times_k U)
\]
is surjective (EGA $\mathrm{Err}_{\mathrm{IV}}$ 53, 21.4.13).

\sgaStableTarget{sga6.e10.bibliography}\subsection*{Bibliography}
\begin{enumerate}[label={[\arabic*]}]
\item Dickson, E. \emph{Linear groups}, Dover, New York, 1958.
\item Grothendieck, A. TDTE, V, VI; S\'eminaire Bourbaki 232, 236.
\item Kleiman, S. \emph{Towards a numerical theory of ampleness}, Ann. of
  Math. 84 (1966), pp. 293--344.
\item Murre, J. \emph{On contravariant functors from the category of
  preschemes over a field into the category of abelian groups}, Publications
  Math\'ematiques de l'IHES 23 (1964).
\item Chevalley, C. \emph{Les classes d'\'equivalence rationnelles, I et II},
  S\'eminaire \`a l'ENS 1958, expos\'es 2 et 3.
\item S. Abhyankar, \emph{Resolution of singularities of arithmetical
  surfaces}, in: \emph{Arithmetical Algebraic Geometry}, Purdue University,
  Harper--Row, New York, 1965.
\end{enumerate}

% END INCLUDED SGA6_English_current_buildable_source_20260823/SGA6_Expose_X_idx598_607_English_SourceChecked_fragment.tex


% BEGIN INCLUDED SGA6_English_current_buildable_source_20260823/SGA6_idx608_646_English_SourceChecked_fragment.tex
% Cumulative source-checked continuation after idx607.
% Exact page ranges remain in separately reviewable fragments.
% BEGIN INCLUDED SGA6_English_current_buildable_source_20260823/SGA6_idx608_615_English_SourceChecked_fragment.tex
% Source-checked English fragment: current-rescribe idx608--615.
% Printed pages 595--602; declared source-PDF pages 598--605.
% Continues Expose X bibliography at a new page and ends mid-item in Lemma 2.4(ii).
\providecommand{\Hom}{\operatorname{Hom}}
\providecommand{\Isom}{\operatorname{Isom}}
\providecommand{\Coker}{\operatorname{Coker}}
\providecommand{\Ass}{\operatorname{Ass}}
\providecommand{\sgacircled}[1]{\tikz[baseline=(char.base)]\node[shape=circle,draw,inner sep=.35pt](char){\mbox{\scriptsize #1}};}
\newpage
% Current rescribe index 608; printed volume page 595; source-PDF page 598.
\sgaStableTarget{sga6.e12.expose}\section*{Expos\'e XII. A relative representability theorem for the Picard functor}
\begin{center}
by M. Raynaud (written up by S. Kleiman)
\end{center}

\sgaStableTarget{sga6.e12.warning}\subsection*{Warning}

The purpose of \hyperlink{sga6.e12.expose}{Expos\'es XII} and XIII is to prove the nine finiteness theorems
for the Picard functor stated by A. Grothendieck in the comments
(pp. C 07--C 11) to his \emph{Fondements de la G\'eom\'etrie Alg\'ebrique}
(cited as FGA). In the first expos\'e, the first finiteness theorem is proved
in the following form, which is more precise than the statement in FGA. Let
$X$ and $Y$ be proper schemes over an integral noetherian base $S$, and let
$f:X\to Y$ be a surjective $S$-morphism. Then there is a nonempty open subset
$V$ of $S$ over which the induced morphism
$f^*:\Pic_{Y/S}\to\Pic_{X/S}$ is representable by a quasi-affine morphism of
finite presentation. In the second expos\'e, the remaining assertions are
proved.

The two expos\'es differ notably in their proof techniques. In the first,
Oort's d\'evissage, non-flat descent, and the representability of unramified
functors are used in an essential way. In the second, one first reduces to
the case in which the schemes $X/S$ under consideration are projective and
flat with geometrically integral fibers, or, when necessary, normal or
smooth, always by applying the first finiteness theorem to a suitable
morphism $X'\to X$. One then uses the tool of $(b)$-sheaves, forged by
I. Castelnuovo and sharpened and used by T. Matsusaka, D. Mumford, Y. Nakai,
and S. Kleiman. In particular, this tool makes it possible to dispense with
Chow coordinates, the Frobenius morphism, and the Lefschetz theorems used in
Grothendieck's original proofs.

The present \hyperlink{sga6.e12.expose}{Expos\'e XII} was presented orally by M. Raynaud following
Grothendieck; the editor wishes to thank him for communicating his personal
notes.

% Current rescribe index 609; printed volume page 596; source-PDF page 599.
\sgaStableTarget{sga6.e12.sec.1}\section*{1. Statement of the main theorem and applications}

\sgaStableTarget{sga6.e12.thm.1-1}\textbf{Theorem 1.1.} Let $S$ be a noetherian scheme, let $X$ and $Y$ be two
proper $S$-schemes, and let $f:X\to Y$ be a morphism. Suppose that $S$ is
integral and that $f$ is surjective. Then there is a nonempty open subset $V$
of $S$ such that the morphism
\[
  f^*:\Pic_{Y/S}|_V\longrightarrow\Pic_{X/S}|_V
\]
between the Picard functors (localized for the fppf topology) is representable
by a quasi-affine morphism of finite presentation.

\sgaStableTarget{sga6.e12.cor.1-2}\textbf{Corollary 1.2.} Let $X$ be a proper scheme over a noetherian scheme
$S$, and suppose that $S$ is integral. Then there is a nonempty open subset
$V$ of $S$ such that the Picard functor $\Pic_{X/S}|_V$ is representable by a
$V$-group scheme that is a sum of quasi-projective open subschemes of finite
type, and hence, \emph{a fortiori}, is separated over $V$.

Indeed, recall the following fact (FGA no.~232). Let $S$ be a scheme and let
$X$ be a projective $S$-scheme, of finite presentation and integral over $S$
(that is, flat with geometrically integral fibers). Then the Picard functor is
represented by an $S$-group scheme $\Pic_{X/S}$, an increasing union of
quasi-projective open subschemes of finite presentation. Consideration of the
``Hilbert polynomials'' decomposes $\Pic_{X/S}$ here as a sum of open
subschemes $\Pic^q_{X/S}$, $q\in\mathbb Q[t]$. In the following expos\'e it is
proved, under the additional hypothesis that $S$ is noetherian, that the
$\Pic^q_{X/S}$ are of finite type over $S$. Thus \hyperlink{sga6.e12.cor.1-2}{1.2} is known when $X$ is in
addition projective and integral over $S$; the general case will be deduced
using the following lemma.

\sgaStableTarget{sga6.e12.lem.1-3}\textbf{Lemma 1.3.} Let $X$ be of finite type over an integral noetherian
scheme $S$. Then there are a nonempty open subset $V$ of $S$, an integral
scheme $S'$, and a finite
% Current rescribe index 610; printed volume page 597; source-PDF page 600.
flat morphism $S'\to V$ such that every irreducible component $X'_i$ of
$X_{S'}$, endowed with the induced reduced structure, is integral over $S'$.

Indeed, there is a finite extension $L$ of the function field $k(S)$ such
that every irreducible component of $(X_L)_{\red}$ is geometrically integral
(EGA IV 4.6.8). Let $S_1$ be an integral $S$-scheme such that $k(S_1)=L$ and
such that there are a nonempty open subset $V_1$ of $S$ and a finite flat
morphism $h:S_1\to V_1$. There is a nonempty open subset $S_2$ of $S_1$ such
% SGA-SRC-001353
that every irreducible component of $X_{S_2}$, endowed with the induced
reduced structure, is integral over $S_2$ (EGA IV 9.7.7). Set
\[
  V=V_1\setminus h(S_1\setminus S_2),\qquad S'=h^{-1}(V).
\]

We now prove 1.2. Let $h:S'\to V_1$ and the $X'_i$ be as in 1.3. For each
$i$, let $f_i:Z_i\to X'_i$ be a Chow presentation, with $Z_i$ integral and
$f_i$ birational. After shrinking $V$ and $S'$, one may suppose that $Z_i$
% SGA-SRC-001354
is integral over $S'$ (EGA IV 9.7.7). Let
\[
  f:Z=\coprod_i Z_i\longrightarrow X'=X_{S'}
\]
be the canonical morphism. Then
$\Pic_{Z/S'}=\prod_i\Pic_{Z_i/S'}$ is representable by a group scheme that is
a sum of quasi-projective open subschemes of finite type. Finally, \hyperlink{sga6.e12.thm.1-1}{1.1} says
that, after shrinking $V$ and $S'$, the morphism
\[
  f^*:\Pic_{X'/S'}\longrightarrow\Pic_{Z/S'}
\]
is representable by a quasi-affine morphism. Consequently,
$\Pic_{X'/S'}$ is representable by a group scheme that is a sum of
quasi-projective open subschemes of finite type; by finite flat descent
(SGA 1 VIII 7.6 and EGA II 6.6.1), the same is true of
$\Pic_{X/S}|_V$.

\sgaStableTarget{sga6.e12.lem.1-4}\textbf{Lemma 1.4.} Let $k$ be a field, let $P$ and $Q$ be $k$-group schemes
locally of finite type, and let $f:P\to Q$ be a quasi-affine morphism. Then
$f$ is affine.

Indeed, let $R=\Ker(f)$. Then $R$ is a quasi-affine $k$-group scheme of
finite type; by SGA 3 $\mathrm{VI}_{\mathrm B}$ 11.11, $R$ is affine.
% Current rescribe index 611; printed volume page 598; source-PDF page 601.
Thus $P\to P/R$, which makes $P$ an $R$-torsor over $P/R$, is affine.
Moreover, by SGA 3 I 4.2, $P/R\to Q$ is a closed immersion, hence affine;
therefore the composite $P\to P/R\to Q$ is affine.

\sgaStableTarget{sga6.e12.cor.1-5}\textbf{Corollary 1.5.} Let $k$ be a field.
\begin{enumerate}[label=\alph*)]
\item If $X$ is a proper $k$-scheme, then $\Pic_{X/k}$ is representable
  (J. P. Murre).
\item If $f:X\to Y$ is a surjective morphism between two proper $k$-schemes,
  then the induced morphism
  \[
    f^*:\Pic_{Y/k}\longrightarrow\Pic_{X/k}
  \]
  is affine.
\end{enumerate}

\sgaStableTarget{sga6.e12.cor.1-6}\textbf{Corollary 1.6.} Let $S$ be an integral noetherian scheme, and let
$f:X\to Y$ be a surjective $S$-morphism between two proper $S$-schemes.
\begin{enumerate}[label=\alph*)]
\item If $Y$ is normal over $S$ (that is, flat with geometrically normal
  fibers), then, after shrinking $S$,
  $f^*:\Pic_{Y/S}\to\Pic_{X/S}$ is quasi-finite.
\item If $\Pic_{Y/S}$ is essentially proper over $S$ (that is, if it
  satisfies the valuative criterion of properness; for example, this holds
  if $Y$ is smooth over $S$), then, after shrinking $S$,
  $f^*:\Pic_{Y/S}\to\Pic_{X/S}$ is finite (and, in particular, affine).
\end{enumerate}
Indeed, in (a), for every $s\in S$, $R_s=\Ker(f_s^*)$ is an affine algebraic
group; its identity component is a proper algebraic group over $k(s)$
(FGA 236, Theorem 2.1 (iii)), so $R_s$ is finite. In (b), $f^*$ is both
proper and quasi-affine, hence finite.

\sgaStableTarget{sga6.e12.rem.1-7}\textbf{Remark 1.7.} It is not known whether \hyperlink{sga6.e12.thm.1-1}{1.1} remains valid when
``quasi-affine'' is replaced by ``affine,'' as \hyperlink{sga6.e12.cor.1-5}{1.5} and \hyperlink{sga6.e12.cor.1-6}{1.6} might suggest.

% Current rescribe index 612; printed volume page 599; source-PDF page 602.
\sgaStableTarget{sga6.e12.sec.2}\section*{2. First reductions}

\sgaStableTarget{sga6.e12.para.2-1}\textbf{2.1.} Apply \hyperlink{sga6.e12.lem.1-3}{1.3} to $X/S$ and $Y/S$. One obtains a nonempty open
subset $V$ of $S$, an integral scheme $S'$, and a finite flat morphism
$S'\to V$ such that every irreducible component $X'_i$ of $X_{S'}$ and
$Y'_j$ of $Y_{S'}$, endowed with the induced reduced structure, is integral
over $S'$. By fpqc descent (EGA IV 2.7.1 (xiv) and SGA 1 VIII 7.9), one may
replace $S$ by $S'$. It then follows from the commutative diagram
\[
\begin{tikzcd}
X_{\red} \arrow[r,"f_{\red}"] \arrow[d,"i"'] & Y_{\red} \arrow[d,"j"]\\
X \arrow[r,"f"'] & Y
\end{tikzcd}
\]
which implies the commutative diagram
\[
\begin{tikzcd}
\Pic_{X_{\red}/S} & \Pic_{Y_{\red}/S} \arrow[l]\\
\Pic_{X/S} \arrow[u] & \Pic_{Y/S} \arrow[l] \arrow[u]
\end{tikzcd}
\]
and from elementary sorites concerning morphisms of functors representable
by quasi-affine morphisms (EGA II 5.1.10 (ii), (v)) that it is enough to prove
the following two assertions:
\begin{enumerate}[label=\Roman*.]
\item Theorem 1.1 when $f$ is an immersion (necessarily nilpotent).
\item Theorem 1.1 when $X$ and $Y$ are reduced and every irreducible component
  $X_i$ of $X$ and $Y_i$ of $Y$, endowed with the induced reduced structure,
  is integral over $S$.
\end{enumerate}

\sgaStableTarget{sga6.e12.para.2-2}\textbf{2.2.} Consider assertion II. For every $i$, let $j=\varphi(i)$ be an
index such that $f(X_i)\subset Y_j$, and put
$\Sigma_j=\coprod_{\varphi(i)=j}X_i$. Then one has the two diagrams
% Current rescribe index 613; printed volume page 600; source-PDF page 603.
\[
\begin{tikzcd}
Y & \coprod_jY_j \arrow[l]\\
X \arrow[u] & \coprod_j\Sigma_j \arrow[l] \arrow[u]
\end{tikzcd}
\]
and
\[
\begin{tikzcd}[column sep=large,row sep=large]
\Pic_{Y/S} \arrow[r,"\sgacircled{1}"] \arrow[d,"\sgacircled{2}"'] &
\Pic_{\coprod_jY_j/S}=\prod_j\Pic_{Y_j/S}
  \arrow[d,"\sgacircled{4}"]\\
\Pic_{X/S} \arrow[r,"\sgacircled{3}"'] &
\Pic_{\coprod_j\Sigma_j/S}=\prod_j\Pic_{\Sigma_j/S}.
\end{tikzcd}
\]
To establish II, it is enough to see that, after shrinking $S$,
$\sgacircled{1}$, $\sgacircled{3}$, and $\sgacircled{4}$ are quasi-affine.

Consider $\sgacircled{4}$ first. Replacing $Y$ by $Y_j$ and $X$ by
$\Sigma_j$, one is reduced to the case in which $Y$ is integral over $S$.
Since $f$ is surjective, one of the $X_i$, say $X_0$, dominates $Y$, and one
has the two sequences
\[
% SGA-SRC-001357
  X_0\hookrightarrow\coprod_iX_i\longrightarrow Y
\]
and
\[
  \Pic_{X_0/S}\longleftarrow\Pic_{\coprod_iX_i/S}
  \longleftarrow\Pic_{Y/S}.
\]
Thus, if the $\Pic_{X_i/S}$, and hence also
$\Pic_{\coprod_iX_i/S}$, are known to be representable and the composite
morphism is quasi-affine, then $\sgacircled{4}$ is quasi-affine. In view of
the proof of \hyperlink{sga6.e12.cor.1-2}{1.2}, proving $\sgacircled{4}$ is therefore reduced to proving
the following assertion:

\medskip
\noindent
$\mathrm{II}'$. \hyperlink{sga6.e12.thm.1-1}{Theorem 1.1} when $X$ and $Y$ are integral over $S$.
\medskip

For $\sgacircled{1}$ and $\sgacircled{3}$, one must treat the case of a
morphism $\coprod_iX_i\to X$, where $X$ is reduced and the $X_i$ are the
  irreducible components of $X$ and are integral over $S$. One may suppose that
  $X$ has connected fibers. Indeed, let $\eta$ be a geometric generic point of
  $S$.
Then the irreducible components of $X_\eta$ are the $(X_i)_\eta$, and, after
shrinking $S$, one may suppose that $X_i\cap X_j=\varnothing$ if and only if
% Current rescribe index 614; printed volume page 601; source-PDF page 604.
$(X_i)_\eta\cap(X_j)_\eta=\varnothing$. Let
$Y_k=\bigcup_{i\in I(k)}X_i$ be the connected components of $X$. Then $f$
decomposes as
\[
  f=\coprod_k f_k:
  \coprod_k\left(\coprod_{i\in I(k)}X_i\right)
% SGA-SRC-001358
  \longrightarrow\coprod_kY_k,
\]
and the assertion follows. Finally, an immediate induction reduces us to the
case of a morphism $X_1\amalg X_2\to X$. After shrinking $S$,
$\sgacircled{1}$ and $\sgacircled{3}$ then follow from EGA III 7.8.6 and the
following proposition.

\sgaStableTarget{sga6.e12.prop.2-3}\textbf{Proposition 2.3.} Let $S$ be a scheme, let $g:Y\to S$ be a proper
morphism, and let $Y_1,Y_2$ be two closed subschemes defined by the ideals
$I_1,I_2$. Suppose that:
\begin{enumerate}[label=\alph*)]
\item $Y,Y_1,Y_2,Y_1\cap Y_2$ are faithfully flat and of finite presentation
  over $S$;
\item $g_*\mathcal O_Y=\mathcal O_S$ and
  $(g_i)_*\mathcal O_{Y_i}=\mathcal O_S$ universally, where
  $g_i=g|_{Y_i}$, $i=1,2$;
\item $I_1\cap I_2=0$.
\end{enumerate}
Then
\[
  \Pic_{Y/S}\longrightarrow\Pic_{(Y_1\amalg Y_2)/S}
\]
is representable by an affine morphism of finite presentation.

First note that the hypotheses are stable under every base change. For (c),
this follows from the two exact sequences
\[
  0\longrightarrow I_1+I_2\longrightarrow\mathcal O_Y
  \longrightarrow\mathcal O_{Y_1\cap Y_2}\longrightarrow0
\]
and
\[
  0\longrightarrow I_1\cap I_2\longrightarrow I_1\oplus I_2
  \longrightarrow I_1+I_2\longrightarrow0.
\]
The first shows that $I_1+I_2$ is flat over $S$, and the second then shows
that $I_1\cap I_2$ is stable under base change.

The assertion is local for the fppf topology. After the base change
$Y_1\cap Y_2\to S$, one may therefore suppose that $Y_1\cap Y_2$ has a
section $\alpha$, which defines sections $\alpha_i$ of $Y_i$ and $\beta$ of
$Y$. After a suitable base change, one is reduced to showing that every point
$\mu:S\to\Pic_{(Y_1\amalg Y_2)/S}$ has as inverse image a subfunctor of
% Current rescribe index 615; printed volume page 602; source-PDF page 605.
$\Pic_{Y/S}$ that is representable by an affine scheme over $S$. Now $\mu$
is defined by an invertible sheaf $L_i$ on $Y_i$, furnished with a
rigidification
\[
  u_i:\alpha_i^*L_i\xrightarrow{\sim}\mathcal O_S,
  \qquad i=1,2.
\]
Likewise, a point $\lambda:S\to\Pic_{Y/S}$ corresponds to an invertible
sheaf $L$ on $Y$ furnished with a rigidification
% SGA-SRC-001360
$u:\beta^*L\xrightarrow{\sim}\mathcal O_S$. The image of $\lambda$ in
$\Pic_{(Y_1\amalg Y_2)/S}$ is $\mu$ if and only if there are isomorphisms of
rigidified sheaves $f_i^*L\simeq L_i$, where $f_i:Y_i\hookrightarrow Y$ is
the inclusion, for $i=1,2$. These isomorphisms define an isomorphism
\[
  \varphi:L_1|_{Y_1\cap Y_2}\xrightarrow{\sim}
  L_2|_{Y_1\cap Y_2}
\]
such that $u_1=u_2\circ\varphi$.

Conversely, if $\varphi$ is such an isomorphism, there is a rigidified
invertible sheaf $L$ on $Y$ obtained by gluing along $\varphi$. Explicitly,
if $f_{12}:Y_1\cap Y_2\hookrightarrow Y$ denotes the common restriction of
$f_1$ and $f_2$, then $L$ is the kernel of the canonical homomorphism
\[
  (f_1)_*L_1\oplus(f_2)_*L_2
  \longrightarrow
  (f_{12})_*\bigl(L_2|_{Y_1\cap Y_2}\bigr)
\]
defined by $\varphi$. This $L$ then corresponds to a point $\lambda$ of
$\Pic_{Y/S}$ mapping to $\mu$.

Finally, let $I$ be the functor that assigns, to an $S$-scheme $T$, the set
of isomorphisms
\[
  \varphi:(L_1)_T|_{(Y_1\cap Y_2)_T}\xrightarrow{\sim}
  (L_2)_T|_{(Y_1\cap Y_2)_T}
\]
such that $(u_1)_T=(u_2)_T\circ\varphi$. We have just established an
isomorphism between the inverse image of $\mu$ in $\Pic_{Y/S}$ and the
functor $I$. It is therefore enough to apply the following lemma (with
$X=Y_1\cap Y_2$).

\sgaStableTarget{sga6.e12.lem.2-4}\textbf{Lemma 2.4.} Let $X$ be a flat, proper scheme of finite presentation
over $S$. Let $L,M$ be two invertible sheaves on $X$, and let $\Isom(L,M)$ be
the functor of isomorphisms $L_T\to M_T$, for $T/S$. Then:
\begin{enumerate}[label=\roman*)]
\item $\Isom(L,M)$ is representable by an affine $S$-scheme of finite
  presentation.
\item If, moreover, $L$ and $M$ are furnished with rigidifications (relative
  to a

% END INCLUDED SGA6_English_current_buildable_source_20260823/SGA6_idx608_615_English_SourceChecked_fragment.tex

% BEGIN INCLUDED SGA6_English_current_buildable_source_20260823/SGA6_idx616_624_English_SourceChecked_fragment.tex
% Current rescribe index 616; printed volume page 603; source-PDF page 606.
given section of $X$), the subfunctor of $\Isom(L,M)$ formed by the
isomorphisms compatible with these rigidifications is representable by a
closed subscheme of $\Isom(L,M)$.
\end{enumerate}

Indeed, (ii) is clear. For (i), recall that the functors
$T\mapsto\operatorname{Hom}(L_T,M_T)$ and
$T\mapsto\operatorname{Hom}(M_T,L_T)$ are representable by
vector bundles $\mathbb V(E)$ and $\mathbb V(F)$ (EGA III 7.7.8). Thus, if
$u\in\operatorname{Hom}(L_T,M_T)$ and
$v\in\operatorname{Hom}(M_T,L_T)$, the conditions
$u\circ v=1$ and $v\circ u=1$ define a closed subscheme $Z$ of
$\mathbb V(E)\times_S\mathbb V(F)$, and it is clear that
% SGA-SRC-001363
$Z\simeq\Isom(L,M)$.

\sgaStableTarget{sga6.e12.para.2-5}\textbf{2.5.} Finally, let us reduce assertion $\mathrm{II}'$ to the following
assertion:

\medskip
\noindent
$\mathrm{II}^*$. \hyperlink{sga6.e12.thm.1-1}{Theorem 1.1} when $X$ and $Y$ are integral over $S$ and the
sequence
\sgaStableTarget{sga6.e12.eq.unnamed}\begin{equation}\tag{*}
  \mathcal O_Y\longrightarrow f_*\mathcal O_X
  \rightrightarrows g_*\mathcal O_{X\times_YX}
\end{equation}
(where $g:X\times_YX\to Y$ is the structural morphism) is exact.
\medskip

For this, it is enough to prove that there is a factorization of $f$
\[
  X\xrightarrow{f_0}Y_0\xrightarrow{f_1}Y_1
  \longrightarrow\cdots\xrightarrow{f_n}Y_n=Y
\]
such that every $f_i$ satisfies (*) and, after possibly restricting $S$,
every $Y_i$ is integral over $S$.

For this purpose, put $Y_0=\operatorname{Spec}(f_*\mathcal O_X)$, with $f_0$ arising from
the Stein factorization of $f$ (EGA III 4.3.3), so that
$f_{0*}\mathcal O_X\xleftarrow{\sim}\mathcal O_{Y_0}$; hence $f_0$ does
satisfy (*). To factor $Y_0\to Y$ in the same way, one uses the following
lemma.

\sgaStableTarget{sga6.e12.lem.2-6}\textbf{Lemma 2.6.} Let $Y$ be a noetherian scheme, and let $f:X\to Y$ be a
finite epimorphism (that is, $\mathcal O_Y\to f_*\mathcal O_X$ is
injective). Then $f$ factors into a finite sequence of finite morphisms
\[
  X\xrightarrow{\sim}Y_0\xrightarrow{f_1}Y_1
  \longrightarrow\cdots\longrightarrow Y_n=Y
\]
that are effective epimorphisms (that is, for each $i$ the sequence

% Current rescribe index 617; printed volume page 604; source-PDF page 607.
\[
  \mathcal O_{Y_i}\longrightarrow(f_i)_*\mathcal O_{Y_{i-1}}
  \rightrightarrows
  (g_i)_*\bigl(
    \mathcal O_{Y_{i-1}}\otimes_{\mathcal O_{Y_i}}
    \mathcal O_{Y_{i-1}}
  \bigr)
\]
is exact).

Put $A=\mathcal O_Y$, $B_0=f_*\mathcal O_X$, and define recursively
\[
  B_i=\Ker\bigl(
    B_{i-1}\rightrightarrows B_{i-1}\otimes_A B_{i-1}
  \bigr).
\]
One checks at once that
\[
  B_{i-1}\otimes_A B_{i-1}
  \xrightarrow{\sim}
  B_{i-1}\otimes_{B_i}B_{i-1}.
\]
It remains to prove that $B_i=A$ for all sufficiently large $i$. Indeed,
% SGA-SRC-001365 — reconstructed by mathematical inference; scan inconclusive
  consider the decreasing sequence of $A$-modules $M_i=B_i/A$, and put
$S_i=\operatorname{Supp}(M_i)$. The decreasing sequence of closed subsets $S_i$ of $Y$
is stationary; let $S_\infty$ be its stationary value. We show that
$S_\infty=\varnothing$. Otherwise, let $s$ be a maximal point of
$S_\infty$. After replacing $Y$ by $\operatorname{Spec}\mathcal O_{Y,s}$, we may suppose
that $Y$ is local with closed point $s$. By the choice of $s$, $M_i$ is
artinian for all sufficiently large $i$, and therefore the sequence $(M_i)$
is stationary. Thus $B_i=B_{i+1}$ for some $i$. It follows that
$\operatorname{Spec}(B_i)\to Y$ is a monomorphism, hence a closed immersion (Nakayama).
Since $A\subset B_i$, one has $A=B_i$, contrary to $s\in S_\infty$.

\sgaStableTarget{sga6.e12.sec.3}\section*{3. Proof of I: Oort's d\'evissage}

\sgaStableTarget{sga6.e12.prop.3-1}\textbf{Proposition 3.1.} (Oort, \emph{Bull. soc. math. France}, vol.~90
(1962)). Let $S$ be a locally noetherian scheme, let $f:X\to S$ be a proper
morphism, and let $X'\subset X$ be a closed subscheme. Suppose that:
\begin{enumerate}[label=\roman*)]
\item $X'$ is defined by an ideal $I$ contained in the nilradical $N$ of
  $X$, and $NI=0$.
\item $X'$ is flat over $S$ and cohomologically flat over $S$ in dimension
  $0$ (EGA III 7.8.1).
\item $I$ is flat over $S$ and cohomologically flat over $S$ in dimensions
  $0$ and $1$, so that (by EGA III 7.4.6) there are a coherent
  $\mathcal O_S$-module $Q$ and, functorially in the coherent
  $\mathcal O_S$-module $M$, an isomorphism

% Current rescribe index 618; printed volume page 605; source-PDF page 608.
\[
  R^2f_*\bigl(I\otimes f^*M\bigr)
  \simeq\operatorname{Hom}_{\mathcal O_S}(Q,M).
\]
\item The cokernel $E$ of
  $f_*\mathcal O_{X'}\to R^1f_*(I)$ is a locally free
  $\mathcal O_S$-module.
\item The homomorphism
  \[
    (f_*\mathcal O_X)_{\red}
    \longrightarrow(f_*\mathcal O_{X'})_{\red}
  \]
  is surjective.
\end{enumerate}
Then the canonical morphism
\[
  \Pic_{X/S}\longrightarrow\Pic_{X'/S}
\]
is representable by an affine morphism of finite presentation. More
precisely, there is a canonical exact sequence of fppf sheaves
\[
  0\longrightarrow\mathbb V(E^\vee)\longrightarrow\Pic_{X/S}
  \longrightarrow\Pic_{X'/S}\xrightarrow{u}\mathbb V(Q),
\]
and, if $P=\Ker(u)$, the monomorphism
$P\hookrightarrow\Pic_{X'/S}$ is representable by a closed immersion and
$\Pic_{X/S}$ is a torsor over $P$ under the group $\mathbb V(E^\vee)$
(that is, the sequence of fppf sheaves
\[
  0\longrightarrow\mathbb V(E^\vee)\longrightarrow\Pic_{X/S}
  \longrightarrow P\longrightarrow0
\]
is exact).

For an $S$-scheme $T$, consider the following two diagrams with exact rows:
\[
\begin{tikzcd}
0 \arrow[r] & I_T \arrow[r] & \mathcal O_{X_T} \arrow[r] &
  \mathcal O_{X'_T}^{*} \arrow[r] & 0\\
0 \arrow[r] & I_T \arrow[r,"\alpha"] \arrow[u,equal] &
  \mathcal O_{X_T}^{*} \arrow[r] \arrow[u] &
  \mathcal O_{X'_T}^{*} \arrow[r] \arrow[u] & 0,
\end{tikzcd}
\qquad \alpha(i)=1+i,
\]
and
\[
\begin{tikzcd}[column sep=tiny]
(f_T)_*\mathcal O_{X'_T} \arrow[r,"\delta_T"] &
  R^1(f_T)_*I_T \arrow[d,equal]\\
(f_T)_*\mathcal O_{X'_T}^{*} \arrow[r,"\delta_T^*"'] \arrow[u] &
  R^1(f_T)_*I_T \arrow[r] &
  R^1(f_T)_*\mathcal O_{X_T}^{*} \arrow[r] &
  R^1(f_T)_*\mathcal O_{X'_T}^{*} \arrow[r] &
  R^2(f_T)_*I_T.
\end{tikzcd}
\]
Note carefully that the left-hand square in this diagram is \emph{not
necessarily commutative}. Nevertheless,
\[
  \sgaStableTarget{sga6.e12.para.3-2}\textbf{3.2.}\qquad \operatorname{im}\delta_T^*
  =\operatorname{im}\delta_T.
\]

% Current rescribe index 619; printed volume page 606; source-PDF page 609.
Indeed, we may plainly suppose that $S$ and $T$ are affine. Choose an affine
covering $\{V_i\}$ of $X_T$. View $\mathcal O_{X'_T}$ as its direct image
under the closed immersion $X'_T\hookrightarrow X_T$. Let
$g\in\Gamma(X_T,\mathcal O_{X'_T}^{*})$, and let $g_i$ be a lift of
$g|_{V_i}$ in $\Gamma(V_i,\mathcal O_{X_T}^{*})$. Then $\delta_T^*(g)$ is
the class of the $1$-cocycle
\[
  \delta_{ij}^*=g_i g_j^{-1}-1,
\]
whereas $\delta_T(g)$ is the class of
\[
  \delta_{ij}=g_i-g_j=g_j\delta_{ij}^*=g\delta_{ij}^*.
\]
Consequently,
\[
  \delta_T(g)=g\,\delta_T^*(g),
\]
taking into account the action of
$\Gamma(X_T,\mathcal O_{X'_T})$ on $I_T$, and hence on $H^1(X_T,I_T)$.

By (ii),
\[
  \Gamma(X_T,\mathcal O_{X'_T})
  =\Gamma(X,\mathcal O_{X'})\otimes_S\Gamma(T,\mathcal O_T),
\]
and hence $g=\sum g_\beta\otimes t_\beta$, with
$g_\beta\in\Gamma(X,\mathcal O_{X'})$ and
$t_\beta\in\Gamma(T,\mathcal O_T)$. By (v), the map
\[
  (f_*\mathcal O_X)_{\red}\longrightarrow
  (f_*\mathcal O_{X'})_{\red}
\]
is surjective. Thus
$g_\beta=v(h_\beta)+n_\beta$, where
$h_\beta\in\Gamma(X,\mathcal O_X)$,
$v:\Gamma(X,\mathcal O_X)\to\Gamma(X,\mathcal O_{X'})$, and
$n_\beta\in\operatorname{nilrad}\bigl(\Gamma(X,\mathcal O_{X'})\bigr)$.

By (i), $(n_\beta\otimes t_\beta)I_T=0$. Therefore
$\delta_T(g)=h\,\delta_T^*(g)$, where
\[
  h=\sum h_\beta\otimes t_\beta
  \in\Gamma(X_T,\mathcal O_{X_T}).
\]
It follows that
\[
  \delta_T^*(g)=h^{-1}\delta_T(g)=\delta_T(h^{-1}g),
\]
and hence $\operatorname{im}\delta_T^*\subset
\operatorname{im}\delta_T$.

Finally, let
$g=\sum g_\beta\otimes t_\beta\in\Gamma(X_T,\mathcal O_{X'_T})$.
Write $g_\beta=v(h_\beta)+n_\beta$, with
$n_\beta\in\operatorname{nilrad}\bigl(\Gamma(X,\mathcal O_{X'})\bigr)$. Then
\[
\begin{split}
  \delta_T(g)
  &=\delta_T\Bigl(\sum n_\beta\otimes t_\beta\Bigr)\\
  &=\delta_T\Bigl(1+\sum n_\beta\otimes t_\beta\Bigr)\\
  &=\delta_T^*\Bigl(1+\sum n_\beta\otimes t_\beta\Bigr),
\end{split}
\]
whence $\operatorname{im}\delta_T\subset
\operatorname{im}\delta_T^*$.

By \hyperlink{sga6.e12.para.3-2}{3.2}, (iii), and (iv), there is an exact sequence, functorial in $T$,
\[
  0\longrightarrow\mathbb V(E^\vee)(T)
  \longrightarrow R^1(f_T)_*\mathcal O_{X_T}^{*}
  \longrightarrow R^1(f_T)_*\mathcal O_{X'_T}^{*}
  \longrightarrow\mathbb V(Q)(T),
\]
which, on passing to the associated fppf sheaves, gives an exact sequence
\[
  0\longrightarrow\mathbb V(E^\vee)\longrightarrow\Pic_{X/S}
  \longrightarrow\Pic_{X'/S}\xrightarrow{u}\mathbb V(Q).
\]
If $P=\Ker(u)$, the monomorphism $P\hookrightarrow\Pic_{X'/S}$ is obtained
by base change from the unit section $S\hookrightarrow\mathbb V(Q)$.
Since $\mathbb V(Q)$ is representable and separated, it follows that
$P\hookrightarrow\Pic_{X'/S}$ is representable by a

% Current rescribe index 620; printed volume page 607; source-PDF page 610.
closed immersion.

Finally,
\[
  0\longrightarrow\mathbb V(E^\vee)\longrightarrow\Pic_{X/S}
  \xrightarrow{w}P\longrightarrow0
\]
is a torsor for the fppf topology. A fortiori, by fpqc descent from a
``splitting'' of the torsor, $w$ is representable by a faithfully flat
affine morphism of finite presentation (see Oort, \emph{Commutative group
schemes}, Proposition 17.4).

\sgaStableTarget{sga6.e12.cor.3-3}\textbf{Corollary 3.3.} Under the hypotheses of \hyperlink{sga6.e12.prop.3-1}{3.1}:
\begin{enumerate}[label=\roman*)]
\item If
  \[
    \Pic_{(\mathrm{zar})(X'/S)}\longrightarrow
    \Pic_{(\mathrm{fppf})(X'/S)}
  \]
  is an isomorphism, then
  \[
    \Pic_{(\mathrm{zar})(X/S)}\longrightarrow
    \Pic_{(\mathrm{fppf})(X/S)}
  \]
  is an isomorphism.
\item If the $S$-scheme $T$ is affine and
  \[
    \Pic(X'_T)\longrightarrow\Pic_{(\mathrm{zar})(X'/S)}^{(T)}
  \]
  is surjective, then
  \[
    \Pic(X_T)\longrightarrow\Pic_{(\mathrm{zar})(X/S)}^{(T)}
  \]
  is surjective.
\end{enumerate}

Indeed, (i) follows by applying the five lemma to the following diagram of
sheaves for the Zariski topology:
\[
\begin{tikzcd}[column sep=small]
0 \arrow[r] \arrow[d,"S"] &
  \mathbb V(E^\vee) \arrow[r] \arrow[d,"S"] &
  \Pic_{(\mathrm{zar})(X/S)} \arrow[r] \arrow[d] &
  \Pic_{(\mathrm{zar})(X'/S)} \arrow[r] \arrow[d,"S"] &
  \mathbb V(Q) \arrow[d,"S"]\\
0 \arrow[r] & \mathbb V(E^\vee) \arrow[r] &
  \Pic_{(\mathrm{fppf})(X/S)} \arrow[r] &
  \Pic_{(\mathrm{fppf})(X'/S)} \arrow[r] & \mathbb V(Q).
\end{tikzcd}
\]

Assertion (ii) follows from the following diagram with exact rows:
\[
\begin{tikzcd}[column sep=small]
H^1(X_T,I_T) \arrow[r] \arrow[d] &
  H^1(X_T,\mathcal O_{X_T}^{*}) \arrow[r] \arrow[d] &
  H^1(X_T,\mathcal O_{X'_T}^{*}) \arrow[r] \arrow[d] &
  H^2(X_T,I_T) \arrow[d]\\
H^0(T,R^1(f_T)_*I_T) \arrow[r] &
  H^0(T,R^1(f_T)_*\mathcal O_{X_T}^{*}) \arrow[r] &
  H^0(T,R^1(f_T)_*\mathcal O_{X'_T}^{*}) \arrow[r] &
  H^0(T,R^2(f_T)_*I_T),
\end{tikzcd}
\]
because the outside vertical arrows are isomorphisms when $T$ is affine.

% Current rescribe index 621; printed volume page 608; source-PDF page 611.
\sgaStableTarget{sga6.e12.cor.3-4}\textbf{Corollary 3.4.} Let $k$ be a field, and let $X$ be a proper
$k$-scheme. Suppose that $\Pic_{X_{\red}/k}$ possesses a universal
invertible sheaf. (For example, descent theory shows that this is the case
if every connected component $X'_i$ of $X_{\red}$ has a rational point; in
particular, if $k$ is algebraically closed, or more generally if the greatest
common divisor of the degrees of the $0$-cycles on every $X'_i$ is $1$.)
Then, over every affine open subset $V$ of $\Pic_{X/k}$, there is a universal
invertible sheaf on $X\times_kV$.

Finally, we prove assertion I in the following more precise form.

\sgaStableTarget{sga6.e12.prop.3-5}\textbf{Proposition 3.5.} Let $S$ be an integral noetherian scheme, and let
$f:X\hookrightarrow Y$ be a nilpotent immersion of proper $S$-schemes.
Then there is a nonempty open subset $V$ of $S$ such that the morphism
\[
  f^*:\Pic_{Y/S}|_V\longrightarrow\Pic_{X/S}|_V
\]
is representable by an affine morphism of finite presentation. In fact,
$N=\Ker(f^*)$ is an affine group scheme with unipotent fibers (in the sense
of SGA 3 XVII).

Indeed, $X$ is defined by a nilpotent ideal $I$. The ideal $I$ has a
filtration by ideals
\[
  I_0=I\supset I_1\supset\cdots\supset I_n=0
\]
such that $NI_{i-1}\subset I_i$ for $1\leq i\leq n$, where $N$ is the
nilradical of $\mathcal O_Y$ (for example, $I_i=N^iI$). Let $X_i$ be the
closed subscheme of $Y$ defined by $I_i$. By induction it is enough to study
the immersions $X_{i-1}\hookrightarrow X_i$. We may therefore suppose that
$NI=0$; this is condition (i) of \hyperlink{sga6.e12.prop.3-1}{3.1}, and we are reduced finally to the case
in which the other conditions (ii)--(v) are satisfied as well. By the generic
% SGA-SRC-001315 — reconstructed by mathematical inference; scan inconclusive
flatness theorem (EGA IV 6.9.1), applied in the notation of \hyperlink{sga6.e12.prop.3-1}{3.1} to
$\mathcal O_{X'}$, the $R^if_*(I)$, and $E$, and taking account of the
K\"unneth relations of EGA III 6.9.9, after restricting $S$ we may suppose
that conditions (ii), (iii), and (iv) hold.

% Current rescribe index 622; printed volume page 609; source-PDF page 612.
There is a finite extension $K$ of $k(S)$ such that
$(f_K)_*\mathcal O_{X'_K}$ has all its residue field extensions trivial.
If $S'$ is an integral scheme, finite and flat over an open subset $V$ of
$S$, with function field $k(S')=K$, there is an open subset $V'$ of $S'$
such that
\[
  \bigl[(f_{S'})_*\mathcal O_{X'_{S'}}\bigr]_{\red}
  \simeq
  \bigl[(f_{S'})_*\mathcal O_{X_{S'}}\bigr]_{\red}
\]
over $V'$. After restricting $S$ and $S'$ further, we may suppose that
$S=V$ and $S'=V'$. Since the assertion to be proved is local for the fpqc
topology, we may replace $S$ by $S'$. Finally, apply 3.1.

\sgaStableTarget{sga6.e12.rem.3-6}\textbf{Remark 3.6.} It follows in the same way, taking account of the exact
sequences
\[
  \Pic_{X_i/S}\longrightarrow\Pic_{X_{i-1}/S}
  \longrightarrow\mathbb V(Q_i),
\]
that, for every $s\in S$,
\[
  \operatorname{coker}\bigl(
    \Pic^0_{Y_s/s}\longrightarrow\Pic^0_{X_s/s}
  \bigr)
\]
is a unipotent group.

\sgaStableTarget{sga6.e12.sec.4}\section*{4. Proof of $\mathrm{II}^*$: the most delicate part of the proof}

Recall assertion $\mathrm{II}^*$: \hyperlink{sga6.e12.thm.1-1}{Theorem 1.1} when $X$ and $Y$ are integral
over $S$ and the sequence
\[
  \mathcal O_Y\longrightarrow f_*\mathcal O_X
  \rightrightarrows g_*\mathcal O_{X\times_YX}
\]
is exact. Applying the generic flatness theorem, EGA IV 6.9.1, and EGA III
7.8.8, we are reduced to proving the following proposition.

\sgaStableTarget{sga6.e12.prop.4-1}\textbf{Proposition 4.1.} Let $S$ be an integral noetherian scheme, let
$\lambda:X\to S$ and $\mu:Y\to S$ be proper morphisms of finite
presentation, let $f:X\to Y$ be a surjective $S$-morphism, and let
$g:X\times_YX\to Y$ be the canonical morphism. Suppose that:
\begin{enumerate}[label=\arabic*)]
\item $X$, $Y$, and $X\times_YX$ are flat over $S$.
\item $\lambda_*\mathcal O_X=\mathcal O_S$ and
  $\mu_*\mathcal O_Y=\mathcal O_S$ universally.

% Current rescribe index 623; printed volume page 610; source-PDF page 613.
\item For all $i,j\geq0$, the modules $R^if_*\mathcal O_X$ and
  $R^jg_*\mathcal O_{X\times_YX}$ are flat over $S$.
\item The sequence
  \sgaStableTarget{sga6.e12.eq.unnamed.occ2}\begin{equation}\tag{*}
    \mathcal O_Y\longrightarrow f_*\mathcal O_X
    \mathrel{\substack{\xrightarrow{\pi_1}\\[-0.7ex]
                       \xrightarrow[\pi_2]{}}}
    g_*\mathcal O_{X\times_YX}
  \end{equation}
  is exact, and $\operatorname{coker}(\pi_1-\pi_2)$ is an
  $\mathcal O_Y$-module flat over $S$.
\end{enumerate}
Then
\[
  f^*:\Pic_{Y/S}\longrightarrow\Pic_{X/S}
\]
is representable by a quasi-affine morphism of finite presentation.

First note that (1)--(4) are stable under every base change $T\to S$. Only
the stability of (3) and (4) is not evident. By the K\"unneth formula (EGA III
6.9.9), (1) and (3) imply that the formation of $R^if_*\mathcal O_X$ and
$R^ig_*\mathcal O_{X\times_YX}$ commutes with base change $T\to S$; this
immediately implies the stability of condition (3). Moreover, the formation
of the cokernel of
\[
  \pi_1-\pi_2:f_*\mathcal O_X
  \longrightarrow g_*\mathcal O_{X\times_YX}
\]
also commutes with base change. Since it is flat over $S$ by (4),
$\operatorname{im}(\pi_1-\pi_2)$ is flat over $S$ and commutes with base
change. Hence $\Ker(\pi_1-\pi_2)$ commutes with base change $T\to S$ and is
therefore universally equal to $\mathcal O_Y$.

The stability of the hypotheses implies two things:
\begin{enumerate}[label=\alph*.]
\item Since the assertion to be proved is local for the fppf topology, by
  descent we may suppose that $X$ has a section $\alpha$; then $Y$ has the
  section $\beta=f\circ\alpha$.
\item To prove the assertion, it is enough to consider an $S$-point of
  $\Pic_{X/S}$, that is, by (a) and (2), an invertible sheaf $L$ on $X$
  endowed with a rigidification $u$ relative to $\alpha$, and to show that
  the fppf sheaf $Z$, the inverse image of this point in $\Pic_{Y/S}$, is
  representable by a quasi-affine $S$-scheme of finite presentation.
\end{enumerate}

% Current rescribe index 624; printed volume page 611; source-PDF page 614.
\sgaStableTarget{sga6.e12.lem.4-2}\textbf{Lemma 4.2.} Under the hypotheses of \hyperlink{sga6.e12.prop.4-1}{4.1}, $f$ is a descent morphism
for the category of locally free sheaves; that is, for every pair $L,M$ of
locally free sheaves on $Y$, the diagram
\[
\begin{tikzcd}[column sep=tiny]
\operatorname{Hom}_Y(L,M) \arrow[r] &
  \operatorname{Hom}_X(f^*L,f^*M)
    \arrow[r,shift left=.7ex] \arrow[r,shift right=.7ex] &
  \operatorname{Hom}_{X\times_YX}(g^*L,g^*M)
\end{tikzcd}
\]
is exact.

Indeed, if $N=\operatorname{Hom}_{\mathcal O_Y}(L,M)$, one must show that the sequence
\[
  \Gamma(Y,N)\longrightarrow\Gamma(X,f^*N)
  \rightrightarrows\Gamma(X\times_YX,g^*N)
\]
is exact, or equivalently (by EGA $0_{\mathrm I}$ 5.4.10.1) that the sequence
\[
  \Gamma(Y,N)
  \longrightarrow
  \Gamma\bigl(Y,f_*\mathcal O_X\otimes_{\mathcal O_Y}N\bigr)
  \rightrightarrows
  \Gamma\bigl(Y,g_*\mathcal O_{X\times_YX}
    \otimes_{\mathcal O_Y}N\bigr)
\]
is exact. But the sequence (*)
\[
  \mathcal O_Y\longrightarrow f_*\mathcal O_X
  \rightrightarrows g_*\mathcal O_{X\times_YX}
\]
is exact by hypothesis. Applying the left exact functor
$\Gamma(N\otimes -)$ gives the lemma, since $N$ is flat over $Y$.

Let $L$ be an invertible sheaf on $X$, endowed with a rigidification $u$
relative to $\alpha$. By \hyperlink{sga6.e12.lem.4-2}{4.2} there is an equivalence between the category
$C$ of pairs $(M,\varphi)$ consisting of an invertible sheaf $M$ on $Y$ and
an isomorphism $\varphi:f^*M\xrightarrow{\sim}L$, and the category
$\mathrm{Deff}(S)$ of effective descent data on $L$ relative to $f$.
If $f^*M\simeq L$, then, since $L$ is rigidified relative to $\alpha$, $M$
is canonically endowed with a rigidification relative to $\beta$.
Consequently, $C$ is equivalent to the discrete category defined by the set
of classes of $\beta$-rigidified invertible sheaves $(M,v)$ on $Y$ such that
\[
  f^*(M,v)\simeq(L,u).
\]
But this is precisely the set $Z(S)$ of $S$-points of the inverse image in
$\Pic_{Y/S}$ of the $S$-point of $\Pic_{X/S}$ defined by $(L,u)$. Moreover,
this description of $Z(S)$ is functorial in $S$. Hence $Z$ identifies with
$\mathrm{Deff}$, and we may forget the rigidifications from now on.

% Exact continuation boundary: idx624 ends after the preceding complete
% paragraph; idx625 begins with the canonical monomorphism
% \mathrm{Deff}\hookrightarrow\mathrm{Rec}.

% END INCLUDED SGA6_English_current_buildable_source_20260823/SGA6_idx616_624_English_SourceChecked_fragment.tex

% BEGIN INCLUDED SGA6_English_current_buildable_source_20260823/SGA6_idx625_635_English_SourceChecked_fragment.tex
% SGA 6 English synchronization fragment.
% Deliverable range only: current-rescribe idx625--635.
% Printed volume pages 612--622; declared source-PDF pages 615--625.
% French control: sga6_fr_workpass.tex (Claude workpass, source-checked through idx646).
% Scan authority: local sga6.pdf; high-resolution supplement used only to
% disambiguate symbols and source-level typographical defects.

\providecommand{\Isom}{\operatorname{Isom}}
\providecommand{\Coker}{\operatorname{Coker}}
\providecommand{\Ker}{\operatorname{Ker}}
\providecommand{\Supp}{\operatorname{Supp}}
\providecommand{\Ass}{\operatorname{Ass}}

% Current rescribe index 625; printed volume page 612; source-PDF page 615.
There is a canonical monomorphism from \(\mathrm{Deff}\) into the functor
\(\mathrm{Rec}\) of gluing data on \(L\). By definition,
\(\mathrm{Rec}\) is the functor
\(\Isom(p_1^*L,p_2^*L)\), where the
\(p_i:X\times_YX\to X\) are the projections; by \hyperlink{sga6.e12.lem.2-4}{2.4} this functor is
represented by an affine scheme \(I\) of finite presentation over \(S\). It
therefore remains to represent the monomorphism
\[
  i:\mathrm{Deff}\hookrightarrow \Isom(p_1^*L,p_2^*L)
\]
by a morphism of finite presentation, which will necessarily be quasi-affine
(by EGA IV 8.11.2).

Let us note in passing that, if desired, one may restrict
\(\Isom(p_1^*L,p_2^*L)\) by imposing the cocycle condition, but this is not
needed below.

After the base change \(I\to S\), we are thus reduced to the following
problem: if \(L\) is an invertible sheaf on \(X\), endowed with a gluing datum
\[
  \varphi:p_1^*L\xrightarrow{\sim}p_2^*L,
\]
represent the subfunctor of \(S\) on which this datum becomes effective. We
have the following lemma.

\sgaStableTarget{sga6.e12.lem.4-3}\textbf{Lemma 4.3.} Under the hypotheses of \hyperlink{sga6.e12.prop.4-1}{4.1}, let \(L\) be an invertible
sheaf on \(X\), endowed with a gluing datum
\[
  \varphi:p_1^*L\xrightarrow{\sim}p_2^*L.
\]
This datum is effective if and only if the following five conditions hold.
\begin{enumerate}[label=\arabic*)]
\item For every \(i\geq0\), \(R^if_*(L)\) is flat over \(S\).
\item For every \(j\geq0\), \(R^jg_*(p_1^*L)\) is flat over \(S\).
\item The cokernel of
\[
  \pi_1-\pi_2:f_*(L)\longrightarrow
  g_*(p_1^*L)\simeq g_*(p_2^*L)
\]
is flat over \(S\), where the latter isomorphism is \(g_*(\varphi)\).
\item The kernel \(N\) of
\(\pi_1-\pi_2:f_*(L)\to g_*(p_1^*L)\) is invertible on \(Y\).
\item The composite morphism
\[
  f^*N\longrightarrow f^*f_*L\longrightarrow L
\]
is an isomorphism.
\end{enumerate}

\newpage
% Current rescribe index 626; printed volume page 613; source-PDF page 616.
Indeed, these conditions are sufficient, since (4) and (5) alone imply
effectivity. Conversely, let \(M\) be invertible on \(Y\), and suppose that
there is an isomorphism \(f^*M\simeq L\) compatible with the gluing datum. We
may cover \(Y\) by open subsets \(V\) over which \(M\) is isomorphic to
\(\mathcal O_Y\). The descent datum on \(L\), restricted over
\(f^{-1}(V)\), is therefore isomorphic to the trivial descent datum on
\(\mathcal O_X|_{f^{-1}(V)}\). Finally, conditions (1)--(5) are local on
\(Y\), and the hypotheses of \hyperlink{sga6.e12.prop.4-1}{4.1} show that they hold for
\(L=\mathcal O_X\); hence they hold for \(L\).

It remains to represent conditions (1)--(5), one after another, by
monomorphisms of finite presentation. For (1) and (2), we apply the following
lemma.

\sgaStableTarget{sga6.e12.lem.4-4}\textbf{Lemma 4.4.} Let \(S\) be a scheme, let
\(\lambda:X\to S\) and \(\mu:Y\to S\) be proper morphisms of finite
presentation, and let \(f:X\to Y\) be an \(S\)-morphism. Let \(F\) be a
sheaf on \(X\), of finite presentation and flat over \(S\). For an integer
\(n\), let \(S_n\) be the subfunctor of \(S\) on which \(R^if_*(F)\) is
flat over \(S\) for every \(i\geq n\); more precisely, for a morphism
\(\phi:T\to S\),
\[
  S_n(T)=
  \begin{cases}
    \{\phi\},&\text{if \(R^i(f_T)_*(F_T)\) is flat for every \(i\geq n\),}\\
    \varnothing,&\text{otherwise.}
  \end{cases}
\]
Then \(S_n\hookrightarrow S\) is representable by a surjective monomorphism
of finite presentation; moreover, for \(i\geq n\), the formation of
\(R^i(f_{S_n})_*(F_{S_n})\) commutes with base changes \(T\to S_n\). In
addition, \(S_n\) coincides with the functor \(S'_n\) on which the functors
\[
  M\longmapsto
  R^i(f_V)_*\bigl(F_V\otimes\lambda_V^*M\bigr),
  \qquad i\geq n,
\]
are exact, where \(M\) ranges over the quasi-coherent
\(\mathcal O_V\)-modules, for every affine open subset \(V\) of \(S\).

Indeed, the assertion is local on \(S\), so we may suppose that \(S\) is
quasi-compact. The relative dimension of \(X/Y\) is then bounded, say by
\(r\), and, for \(n>r\),
\[
  R^nf_*(F_V\otimes\lambda_V^*M)=0.
\]
The assertion therefore holds for \(n>r\).

\newpage
% Current rescribe index 627; printed volume page 614; source-PDF page 617.
Proceeding by descending induction, suppose it is known for \(n+1\). After
replacing \(S\) by \(S_{n+1}=S'_{n+1}\), we may suppose that
\(R^if_*(F)\) is flat over \(S\) for \(i\geq n+1\). Since \(F\) is flat
over \(S\), EGA III 6.9.8 then shows that the formation of \(R^nf_*(F)\)
commutes with every base change. Consequently, the representability of
\(S_n\) follows from the next lemma (see J. P. Murre,
\emph{Representation of unramified functors}, S\'eminaire Bourbaki, May
1965, p.~294-11, Theorem~2).

\sgaStableTarget{sga6.e12.lem.4-5}\textbf{Lemma 4.5.} Let \(X\to S\) be a proper morphism of finite
presentation, and let \(F\) be a sheaf of finite presentation on \(X\). Then
the subfunctor \(S_F\) of \(S\) on which \(F\) is flat over \(S\) is
representable by a surjective monomorphism of finite presentation (which, of
course, is not an immersion in general).

To represent condition (3) of \hyperlink{sga6.e12.lem.4-3}{4.3}, note that, by (1) and (2), the formation
of \(f_*(L)\) and \(g_*(p_1^*L)\) commutes with the base changes
\(I\to S\); hence so does the formation of
\(\Coker(\pi_1-\pi_2)\). We then apply 4.5.

To represent (4), note that, by (3), and because \(f_*(L)\) and
\(g_*(p_1^*L)\) are flat and commute with base change,
\[
  N=\Ker(\pi_1-\pi_2)
\]
also commutes with base change. We then apply the following lemma.

\sgaStableTarget{sga6.e12.lem.4-6}\textbf{Lemma 4.6.} Let \(\lambda:X\to S\) be a proper morphism of finite
presentation, and let \(u:F\to G\) be a homomorphism of
\(\mathcal O_X\)-modules of finite presentation, with \(G\) flat over
\(S\). Then the subfunctor \(S_1\) of \(S\) on which \(u\) is an
isomorphism is representable by an open subscheme of \(S\), of finite
presentation.

As usual, we reduce to the case where \(S\) is noetherian, using EGA IV 8
and 11.2.6; this relieves us of the need to consider finite presentation.
First let us represent the functor on which \(u\) is an epimorphism,

\newpage
% Current rescribe index 628; printed volume page 615; source-PDF page 618.
that is, on which \(\Coker u\) is zero. Since the formation of
\(\Coker u\) commutes with base change, this functor is represented by the
open subset
\[
  S-\lambda\bigl(\Supp(\Coker u)\bigr)
\]
of \(S\). We may therefore suppose that \(u\) is already an epimorphism.
Using now the hypothesis that \(G\) is flat over \(S\), we find that the
formation of \(\Ker u\) commutes with base change; consequently, the functor
under consideration is represented by the open subset
\[
  S-\lambda\bigl(\Supp(\Ker u)\bigr)
\]
of \(S\).

\newpage
% Current rescribe index 629; printed volume page 616; source-PDF page 619.
\sgaStableTarget{sga6.e13.expose}\section*{Expos\'e XIII. Finiteness Theorems for the Picard Functor}
\begin{center}
by S. Kleiman
\end{center}

\sgaStableTarget{sga6.e13.sec.1}\subsection*{1. The \texorpdfstring{\((b)\)}{(b)}-sheaves}

Let \(k\) be an algebraically closed field and \(X\) a projective
\(k\)-scheme, endowed with an ample invertible sheaf \(\mathcal O_X(1)\).

\sgaStableTarget{sga6.e13.def.1-1}\textbf{Definition 1.1 (Mumford).} Let \(m\in\mathbb Z\). A coherent sheaf
\(F\) on \(X\) is said to be \emph{\(m\)-regular} if, at every point of the
support of \(F\), \(\mathcal O_X(1)\) is generated by its global sections
and if, for every \(q\geq1\),
\[
  H^q(F(m-q))=0.
\]

\sgaStableTarget{sga6.e13.lem.1-2}\textbf{Lemma 1.2.} Let
\[
  0\longrightarrow F(-1)\longrightarrow F\longrightarrow G\longrightarrow0
\]
be an exact sequence of coherent sheaves on \(X\).\footnote{Here, as below
in analogous situations, it is understood that \(F(-1)\to F\) is defined
by tensoring with a given section of \(\mathcal O_X(1)\).} If \(F\) is
\(m\)-regular, then so is \(G\).

Indeed, we have the corresponding exact sequence
\[
  H^q(F(m-q))\longrightarrow H^q(G(m-q))
  \longrightarrow H^{q+1}(F(m-q-1)).
\]

\sgaStableTarget{sga6.e13.prop.1-3}\textbf{Proposition 1.3.} If a sheaf \(F\) on \(X\) is \(m\)-regular,
then, for every \(n\geq m\),
\begin{enumerate}[label=(\roman*)]
\item \(F\) is \(n\)-regular.
\item The map
\[
  H^0(F(n))\otimes H^0(\mathcal O_X(1))
  \longrightarrow H^0(F(n+1))
\]
is surjective.
\item \(F(n)\) is generated by \(H^0(F(n))\).
\end{enumerate}

\newpage
% Current rescribe index 630; printed volume page 617; source-PDF page 620.
The proof is by induction on the dimension \(s\) of the support of \(F\); if
\(s=-1\), the assertion is trivial. Choose a section
\(\sigma\in H^0(\mathcal O_X(1))\) which vanishes at no point of the finite
set \(\Ass(F)\). Then \(\sigma\) defines an exact sequence
\[
  0\longrightarrow F(-1)\longrightarrow F\longrightarrow G\longrightarrow0,
\]
and \(G\) is \(m\)-regular by \hyperlink{sga6.e13.lem.1-2}{1.2}, with support of dimension \(s-1\).

Consider the corresponding exact sequence
\[
  H^q(F(n-q-1))\longrightarrow H^q(F(n-q))
  \longrightarrow H^q(G(n-q)).
\]
By the induction hypothesis, \(H^q(G(n-q))=0\) for \(n\geq m\), so the
vanishing of \(H^q(F(n-q-1))\) implies that of \(H^q(F(n-q))\). But
\(H^q(F(m-q))=0\), which proves (i).

To prove (ii), we chase through the following commutative diagram:
\[
\begin{tikzcd}[column sep=1.7em,row sep=2.4em]
  & & H^0(F(n))\otimes H^0(\mathcal O_X(1))
      \arrow[rr] \arrow[d]
  & & H^0(G(n))\otimes H^0(\mathcal O_X(1))
      \arrow[r] \arrow[d]
  & 0 \\
  0 \arrow[r]
  & H^0(F(n)) \arrow[ur] \arrow[r]
  & H^0(F(n+1)) \arrow[rr]
  & & H^0(G(n+1)) \arrow[d] \\
  & & & & 0
\end{tikzcd}
\]
Here the oblique arrow is defined by
\(\varphi\mapsto\varphi\otimes\sigma\); the horizontal rows are exact, the
first by the \(n\)-regularity of \(F\) and by (i), already proved, and the
second column is exact by the induction hypothesis.

Finally, (iii) follows from the commutative diagram
\[
\begin{tikzcd}[column sep=large,row sep=large]
H^0(F(n))_X\otimes H^0(\mathcal O_X(1))_X
  \arrow[r,"\alpha_n"]
  \arrow[d]
& H^0(F(n)(1))_X
  \arrow[d,"\beta_{n+1}"] \\
{[H^0(F(n))_X](1)}
  \arrow[r,"\beta_n\otimes\mathrm{id}"]
& F(n)(1).
\end{tikzcd}
\]

\newpage
% Current rescribe index 631; printed volume page 618; source-PDF page 621.
Here, if \(M\) is a vector space over \(k\), we write \(M_X\) for
\(M\otimes_k\mathcal O_X\). For \(n\geq m\), \(\alpha_n\) is surjective
by (ii); therefore, if \(\beta_{n+1}\) is surjective, so is \(\beta_n\). By
Serre's theorem, \(\beta_{n+1}\) is surjective for \(n\gg0\); this proves
(iii).

\sgaStableTarget{sga6.e13.prop.1-4}\textbf{Proposition 1.4.} Let
\[
  0\longrightarrow F(-1)\longrightarrow F\longrightarrow G\longrightarrow0
\]
be an exact sequence on \(X\). If \(G\) is \(m\)-regular, then:
\begin{enumerate}[label=(\roman*)]
% SGA-SRC-001406
\item \(H^q(F(n))=0\) for \(q\geq2\) and \(n\geq m-q\).
\item \(h^1(F(n-1))\geq h^1(F(n))\) for \(n\geq m-1\).
\item \(H^1(F(n))=0\) for
\(n\geq(m-1)+h^1(F(m-1))\).
\end{enumerate}
In particular, \(F\) is \([m+h^1(F(m-1))]\)-regular.

Indeed, in the exact sequence
\[
  H^{q-1}(G(n))\longrightarrow H^q(F(n-1))
  \longrightarrow H^q(F(n))\longrightarrow H^q(G(n)),
\]
the two outer groups are zero for \(q\geq2\) and
\(n\geq m-(q-1)\), by \hyperlink{sga6.e13.prop.1-3}{1.3(i)}. Thus
\[
  H^q(F(m-q))\xrightarrow{\sim}H^q(F(m-q+1))
  \xrightarrow{\sim}H^q(F(m-q+2))\xrightarrow{\sim}\cdots;
\]
but, by Serre's theorem, \(H^q(F(n))=0\) for \(n\gg0\), proving (i).

For \(n\geq m-1\), consider the exact sequence
\[
  0\to H^0(F(n-1))\to H^0(F(n))
  \xrightarrow{\alpha_n}H^0(G(n))\to H^1(F(n-1))
  \to H^1(F(n))\to0.
\]
It implies (ii). Now suppose that \(\alpha_n\) is surjective, and consider
the commutative diagram
\[
\begin{tikzcd}[column sep=large,row sep=large]
H^0(F(n))\otimes H^0(\mathcal O_X(1))
  \arrow[r,"\alpha_n\otimes\mathrm{id}"]
  \arrow[d]
& H^0(G(n))\otimes H^0(\mathcal O_X(1))
  \arrow[r]
  \arrow[d,"\beta_n"]
& 0 \\
H^0(F(n+1))
  \arrow[r,"\alpha_{n+1}"]
& H^0(G(n+1)).
\end{tikzcd}
\]

\newpage
% Current rescribe index 632; printed volume page 619; source-PDF page 622.
% SGA-SRC-001407
If \(n\geq m\), then \(\beta_n\) is surjective by \hyperlink{sga6.e13.prop.1-3}{1.3(ii)}, and hence
\(\alpha_{n+1}\) is surjective. Consequently,
\[
  H^1(F(n-1))\xrightarrow{\sim}H^1(F(n))
  \xrightarrow{\sim}\cdots\xrightarrow{\sim}0.
\]
On the other hand, if \(\alpha_n\) is not surjective, then
\(h^1(F(n-1))>h^1(F(n))\). Thus the function
% SGA-SRC-001408
\(n\longmapsto h^1(F(n))\) decreases \emph{strictly}, starting from
\(h^1(F(m-1))\), until it becomes zero; this proves (iii).

\sgaStableTarget{sga6.e13.def.1-5}\textbf{Definition 1.5.} Let \(F\) be a coherent sheaf on \(X\), let
\(r\geq\dim\Supp(F)\), and let \((b)=(b_0,\ldots,b_r)\) be a sequence of
\(r+1\) integers. Then \(F\) is said to be a \emph{\((b)\)-sheaf} if it
satisfies either of the following two equivalent sets of conditions.

\noindent\sgaStableTarget{sga6.e13.para.1-5-1}\textbf{(1.5.1)}
\begin{enumerate}[label=(\roman*)]
\item At every point of the support of \(F\), \(\mathcal O_X(1)\) is
generated by \(H^0(\mathcal O_X(1))\).
\item \(h^0(F(-1))\leq b_0\).
\item If \(r\geq1\), there is a section
\(\sigma\in H^0(\mathcal O_X(1))\) which defines an exact sequence
\[
  0\longrightarrow F(-1)\longrightarrow F\longrightarrow G\longrightarrow0
\]
such that \(G\) is a \((b_1,\ldots,b_r)\)-sheaf.
\end{enumerate}

\noindent\sgaStableTarget{sga6.e13.para.1-5-2}\textbf{(1.5.2)} There is an \(F\)-regular sequence of \(r\)
sections \(\sigma_1,\ldots,\sigma_r\in H^0(\mathcal O_X(1))\) such that,
for \(i=0,\ldots,r\), if \(F_i\) denotes the restriction of \(F\) to the
zero scheme of \(\sigma_1,\ldots,\sigma_i\), then
\[
  h^0(F_i(-1))\leq b_i.
\]

\sgaStableTarget{sga6.e13.prop.1-6}\textbf{Proposition 1.6.} Let \(F\) be a \((b)\)-sheaf on \(X\). Then:
\begin{enumerate}[label=(\roman*)]
\item For every sufficiently general sequence \(\sigma\) of \(r\) sections
\(\sigma_1,\ldots,\sigma_r\in H^0(\mathcal O_X(1))\), if
\(F_{\sigma,i}\) denotes the restriction of \(F\) to the zero scheme of
\(\sigma_1,\ldots,\sigma_i\), then
\(h^0(F_{\sigma,i}(-1))\leq b_i\), for \(0\leq i\leq r\).
\item Every subsheaf \(G\) of \(F\) is a \((b)\)-sheaf.
\end{enumerate}

\newpage
% Current rescribe index 633; printed volume page 620; source-PDF page 623.
% SGA-SRC-001410
Indeed, let \(S'=\mathbb A_k^N\) be the affine space whose \(k\)-points
correspond to the sequences \(\sigma\), and let \(T\) be the open subset of
\(S'\) corresponding to the \(F\)-regular sequences. For fixed \(i\), the
sheaves \(F_{\sigma,i}(-1)\) corresponding to the \(k\)-points of \(T\) are
contained in a flat family over \(T\), and \(T\) is nonempty by hypothesis.
Thus (i) follows from upper semicontinuity of the function
\[
  \sigma\longmapsto h^0(F_{\sigma,i}(-1))
\]
(EGA III 7.7.5). Finally, since every sufficiently general sequence is
\(G\)-regular and is such that \(G_{\sigma,i}\) is a subsheaf of
\(F_{\sigma,i}\), for every \(i\), (ii) follows from (i).

\sgaStableTarget{sga6.e13.lem.1-7}\textbf{Lemma 1.7.} Let
\[
% SGA-SRC-001411
  0\longrightarrow F(-1)\longrightarrow F\longrightarrow G\longrightarrow0
\]
be an exact sequence on \(X\). If the Hilbert polynomial of \(F\) is written
in the form
\[
  \chi(F(n))=\sum_{i=0}^{r}a_i\binom{n+i}{i},
\]
then
\[
  \chi(G(n))=\sum_{i=0}^{r-1}a_{i+1}\binom{n+i}{i}.
\]
Indeed,
\[
\begin{aligned}
  \chi(G(n))
  &=\chi(F(n))-\chi(F(n-1))\\
  &=\sum_{i=0}^{r}a_i
    \left[\binom{n+i}{i}-\binom{n-1+i}{i}\right]\\
  &=\sum_{i=1}^{r}a_i\binom{n+i-1}{i-1},
\end{aligned}
\]
which proves the assertion.

\sgaStableTarget{sga6.e13.prop.1-8}\textbf{Proposition 1.8.} Let \(F\) be a \((b)\)-sheaf on \(X\), let
\(s=\dim\Supp(F)\), and let
\[
  \chi(F(n))=\sum_{i=0}^{s}a_i\binom{n+i}{i}
\]
be the Hilbert polynomial of \(F\). Then:
\begin{enumerate}[label=(\roman*)]
\item For \(n\geq-1\),
\[
  h^0(F(n))\leq\sum_{i=0}^{s}b_i\binom{n+i}{i}.
\]
\item \(a_s\leq b_s\), and \(F\) is also a
\((b_0,\ldots,b_{s-1},a_s)\)-sheaf.
\end{enumerate}
Indeed, we argue by induction on \(s\). If \(s=0\), the assertion is clear,
since \(a_0=h^0(F)=h^0(F(-1))\leq b_0\). If \(s\geq1\), we have an exact
sequence
\[
  0\longrightarrow F(-1)\longrightarrow F\longrightarrow G\longrightarrow0
\]

\newpage
% Current rescribe index 634; printed volume page 621; source-PDF page 624.
with \(G\) a \((b_1,\ldots,b_r)\)-sheaf and
\(\dim\Supp(G)=s-1\). Then
\[
  h^0(F(n))-h^0(F(n-1))\leq h^0(G(n)),
\]
whereas
\[
  h^0(G(n))\leq
  \sum_{i=0}^{s-1}b_{i+1}\binom{n+i}{i}
\]
by induction on \(s\), and \(h^0(F(-1))\leq b_0\). Assertion (i) follows
by induction on \(n\geq-1\). Finally, by \hyperlink{sga6.e13.lem.1-7}{1.7} and induction on \(s\),
\(a_s\leq b_s\), and \(G\) is a
\((b_1,\ldots,b_{s-1},a_s)\)-sheaf; this proves (ii).

\sgaStableTarget{sga6.e13.def.1-9}\textbf{Definition 1.9.} The following polynomials with nonnegative rational
coefficients, defined recursively, are called the \emph{\((b)\)-polynomials}:
\[
\begin{cases}
P_{-1}=0,\\
P_r(X_0,\ldots,X_r)=P_{r-1}(X_1,\ldots,X_r)
+\displaystyle\sum_{i=0}^{r}X_i
\binom{P_{r-1}(X_1,\ldots,X_r)-1+i}{i}.
\end{cases}
\]

\sgaStableTarget{sga6.e13.rem.1-10}\textbf{Remark 1.10.} For the \((b)\)-polynomials, induction on \(r\geq t\)
immediately gives
\[
  P_r(X_0,\ldots,X_t,0,\ldots,0)=P_t(X_0,\ldots,X_t).
\]

\sgaStableTarget{sga6.e13.thm.1-11}\textbf{Theorem 1.11.} \emph{(The main theorem on \((b)\)-sheaves.)} Let
\(F\) be a \((b)\)-sheaf on \(X\), and let
\[
  \chi(F(n))=\sum_{i=0}^{r}a_i\binom{n+i}{i}
\]
be the Hilbert polynomial of \(F\). Let
\((c)=(c_0,\ldots,c_r)\) be a sequence of integers satisfying
\((c)\geq(b)-(a)\), and let \(m=P_r(c)\) be the integer given by the
\(r\)-th \((b)\)-polynomial \hyperlink{sga6.e13.def.1-9}{(1.9)}. Then \(m\geq0\), and \(F\) is
\(m\)-regular. In particular, if \(s=\dim\Supp(F)\), then \(F\) is
\[
  P_{s-1}(c_0,\ldots,c_{s-1})\text{-regular}.
\]

\newpage
% Current rescribe index 635; printed volume page 622; source-PDF page 625.
Indeed, we argue by induction on \(r\). If \(r=0\), the assertion is clear,
% SGA-SRC-001412 — reconstructed by mathematical inference; scan inconclusive
since \(m=c_0\geq0\) and \(F\) is \(n\)-regular for every \(n\). If \(r\geq1\),
there is an exact sequence
\[
  0\longrightarrow F(-1)\longrightarrow F\longrightarrow G\longrightarrow0
\]
in which \(G\) is a \((b_1,\ldots,b_r)\)-sheaf. By \hyperlink{sga6.e13.lem.1-7}{1.7} and the induction
hypothesis, if \(n=P_{r-1}(c_1,\ldots,c_r)\), then \(n\geq0\) and \(G\) is
\(n\)-regular. Hence, by \hyperlink{sga6.e13.prop.1-4}{1.4}, \(F\) is
\([n+h^1(F(n-1))]\)-regular, and \(h^q(F(n-1))=0\) for \(q\geq2\).
Therefore
\[
  h^1(F(n-1))=h^0(F(n-1))-\chi(F(n-1)),
\]
which, by \hyperlink{sga6.e13.prop.1-8}{1.8(i)}, is at most
\[
  \sum_{i=0}^{r}(b_i-a_i)\binom{n-1+i}{i},
\]
since \(b_i\geq0\).
% SGA-SRC-001413 — reconstructed by mathematical inference; scan inconclusive
Since \(c_i\geq b_i-a_i\), this quantity is at most
\[
  \sum_{i=0}^{r}c_i\binom{n-1+i}{i}.
\]
Finally, by \hyperlink{sga6.e13.prop.1-3}{1.3(i)}, \(F\) is
\[
  \left[n+\sum_{i=0}^{r}c_i\binom{n-1+i}{i}\right]\text{-regular}.
\]
This proves the first assertion; the second follows from \hyperlink{sga6.e13.prop.1-8}{1.8(ii)} and 1.10.

\sgaStableTarget{sga6.e13.para.1-12}\textbf{1.12.} Let \(S\) be a noetherian scheme, and let \(X\) be an
\(S\)-scheme of finite type. Let \(\mathcal F\) be a family of classes of
coherent sheaves on the fibers of \(X/S\). In other words, for every point
\(s\in S\) and every extension \(K\) of \(k(s)\), coherent sheaves \(F_K\)
are specified on the algebraic scheme \(X_K\); and \(F_K\) and
\(F'_{K'}\) are declared to lie in the same class if there are
\(k(s)\)-homomorphisms from \(K\) and \(K'\) into a common extension
\(K''\) of \(k(s)\) such that
\[
  F_{K''}=F_K\otimes_KK''
  \quad\text{and}\quad
  F'_{K''}=F'_{K'}\otimes_{K'}K''
\]
are isomorphic on \(X_{K''}\). The family \(\mathcal F\) is said to be
\emph{bounded by the coherent sheaf} \(F\) \emph{on}
\(X_T=X\times_ST\), where \(T\) is of finite type over \(S\), if
\(\mathcal F\) is contained in the family of classes of the sheaves
\(F_{k(t)}\), with \(t\in T\). The family \(\mathcal F\) is said to be
\emph{bounded} if such \(T\) and \(F\) exist.

Suppose in addition that \(X/S\) is projective and endowed with a
(relatively) ample invertible sheaf \(\mathcal O_X(1)\). Then
\(\mathcal F\) is called a \((b)\)-\emph{family}, for a sequence of integers
\((b)=(b_0,\ldots,b_r)\), if every class in \(\mathcal F\) is represented
by an \(F_K\), with \(K\) algebraically closed, which is a
\((b)\)-sheaf.

% END INCLUDED SGA6_English_current_buildable_source_20260823/SGA6_idx625_635_English_SourceChecked_fragment.tex

% BEGIN INCLUDED SGA6_English_current_buildable_source_20260823/SGA6_idx636_646_English_SourceChecked_fragment.tex
% SGA 6 English source-checked fragment: current-rescribe idx636--646.
% French control: sga6_fr_workpass.tex. Ultimate check: sga6.pdf.
% Printed pp. 623--633; source-PDF pp. 626--636.
% idx646 closes TeX syntax but not the sentence in Definition XIII.3.3.

% --- current-rescribe idx636 | printed p. 623 | source-PDF p. 626 ---
\sgaStableTarget{sga6.e13.thm.1-13}\textbf{Theorem 1.13.} \emph{Let \(S\) be a noetherian scheme, and let
\(X\) be a projective \(S\)-scheme endowed with an ample sheaf
\(\mathcal O_X(1)\) such that the induced sheaves \(\mathcal O_{X_s}(1)\),
for \(s\in S\), are generated by their \(H^0(\mathcal O_{X_s}(1))\). Let
\(\mathcal F\) be a family of classes of coherent sheaves on the fibers of
\(X/S\). Then the following conditions are equivalent:}
\begin{enumerate}[label=(\roman*)]
\item \emph{\(\mathcal F\) is bounded. Moreover, if all the
\(F_K\in\mathcal F\) are locally free of rank \(p\), then \(\mathcal F\)
may be assumed to be bounded by a sheaf \(F\) on some \(X_T\), with \(F\)
locally free of rank \(p\).}
\item \emph{The set of Hilbert polynomials \(\chi(F_K(n))\), for
\(F_K\in\mathcal F\), is finite, and there exists a sequence of integers
\((b)\) such that \(\mathcal F\) is a \((b)\)-family.}
\item \emph{The set of Hilbert polynomials \(\chi(F_K(n))\), for
\(F_K\in\mathcal F\), is finite, and there exists an integer \(m\) such
that all the \(F_K\in\mathcal F\) are \(m\)-regular.}
\item \emph{The set of Hilbert polynomials \(\chi(F_K(n))\), for
\(F_K\in\mathcal F\), is finite, and \(\mathcal F\) is contained in the
family of classes of quotients of sheaves of the form \(E_K\), where \(E\)
is a coherent sheaf on some \(X_T\). Moreover, one may assume that \(T=S\)
and that \(E\) has the form \(\mathcal O_X(-m)^{\oplus M}\).}
\item \emph{\(\mathcal F\) is contained in the family of classes of
cokernels of homomorphisms of the form \(E'_K\to E_K\), where \(E'\) and
\(E\) are coherent sheaves on some \(X_T\). Moreover, one may assume that
\(T=S\) and that \(E'\) and \(E\) have the forms
\(\mathcal O_X(-m')^{\oplus M'}\) and \(\mathcal O_X(-m)^{\oplus M}\),
respectively.}
\end{enumerate}

Indeed, suppose that \(\mathcal F\) is bounded by a sheaf \(F\) on some
\(X_T\). After partitioning \(T\) into pieces and applying the generic
flatness theorem (EGA IV 6.9.1), one may assume that \(F\) is flat over
\(T\). Then,
% --- current-rescribe idx637 | printed p. 624 | source-PDF p. 627 ---
% SGA-SRC-001425
the number of polynomials \(\chi(F_t(n))\) is at most the number of connected
components of \(T\) (EGA III 7.9.5). It is also an elementary lemma that,
if \(t\in T\) and \(F_t\) is locally free of rank \(p\), then \(F\) is
locally free of rank \(p\) over an open neighborhood of \(t\) in \(T\)
(see EGA IV 11.3.10, with \(X=Y\) and \(f=\mathrm{id}\)). The hypotheses
and EGA III 7.7.6 show, moreover, that after partitioning \(T\) still more
finely one can find a sequence \(\sigma\) of sections
\[
\sigma_1,\ldots,\sigma_r\in H^0(\mathcal O_X(1))
\]
that is \(F\)-regular. Thus (ii) follows from the upper semicontinuity of
the function
\[
t\longmapsto h^0(F_{t,i}(-1)),
\]
where, for \(t\in T\) and \(0\le i\le r\), \(F_{t,i}\) denotes the
restriction of \(F_t\) to the zero scheme of
\(\sigma_1,\ldots,\sigma_i\) (EGA III 7.7.5).

The implication (ii) \(\Rightarrow\) (iii) follows immediately from 1.11.
The implication (iii) \(\Rightarrow\) (iv) follows from \hyperlink{sga6.e13.prop.1-3}{1.3 (iii)}, on
taking
\[
M=\max\bigl(\chi(F_K(m))\bigr),
\qquad E=\mathcal O_X(-m)^{\oplus M}.
\]

Suppose that \(\mathcal F\) satisfies (iv). There are exact sequences
\[
0\longrightarrow F'_K\longrightarrow E_K\longrightarrow F_K
\longrightarrow0,
\]
and the set of polynomials \(\chi(F_K(n))\) is finite. By hypothesis, the
family of classes of the \(E_K\) is bounded. Hence, by the already proved
implication (i) \(\Rightarrow\) (ii), the set of polynomials
\(\chi(E_K(n))\) is finite, and there exists a sequence of integers \((b)\)
such that all the \(E_K\), with \(K\) algebraically closed, are
\((b)\)-sheaves. It follows that the set of polynomials
\(\chi(F'_K(n))\) is finite and, by \hyperlink{sga6.e13.prop.1-6}{1.6}, all the \(F'_K\) are
\((b)\)-sheaves. Applying the already proved implication (ii)
\(\Rightarrow\) (iv) to the family of classes of the \(F'_K\), one obtains
(v).

Finally, suppose that \(\mathcal F\) satisfies (v), and let us prove that
\(\mathcal F\) is bounded. One may assume that \(E\) (respectively \(E'\))
has the form \(\mathcal O_X(-m)^{\oplus M}\). Indeed,
% --- current-rescribe idx638 | printed p. 625 | source-PDF p. 628 ---
by (i) \(\Rightarrow\) (iv), take a surjection
\[
L=\mathcal O_{X_T}(-m)^{\oplus M}\longrightarrow E.
\]
After partitioning \(T\) into pieces so that \(E\) becomes flat over \(T\),
formation of the exact sequence
\[
0\longrightarrow I\xrightarrow{u}L\longrightarrow E\longrightarrow0
\]
commutes with passage to the fibers. By (i) \(\Rightarrow\) (iii) and
(i) \(\Rightarrow\) (iv), choose \(m_1\) sufficiently large that all the
groups
\[
\operatorname{Ext}^1\bigl(\mathcal O_{X_K}(-m_1),I_K\bigr)
=H^1(I_K(m_1))
\]
vanish and that there exist surjections
\[
L_1=\mathcal O_{X_T}(-m_1)^{\oplus M_1}\longrightarrow E',
\qquad
\alpha:L_2=\mathcal O_{X_T}(-m_1)^{\oplus M_2}\longrightarrow I.
\]
The maps
\[
\operatorname{Hom}(L_{1,K},L_K)\longrightarrow
\operatorname{Hom}(L_{1,K},E_K)
\]
are then surjective. Let \(\beta:E'_K\to E_K\) be a homomorphism, and let
\(\gamma\) be the composite
% SGA-SRC-001426
\[
L_{1,K}\longrightarrow E'_K\xrightarrow{\beta}E_K.
\]
Then \(\gamma\) comes from a homomorphism
\(\delta:L_{1,K}\to L_K\), and
\[
(\delta,u_K\alpha):L_{1,K}\oplus L_{2,K}\longrightarrow L_K
\]
and \(\beta:E'_K\to E_K\) have the same cokernel. Consequently,
\(\mathcal F\) is contained in the family of classes of cokernels of
homomorphisms of the form
\[
\mathcal O_{X_K}(-m_1)^{\oplus(M_1+M_2)}
\longrightarrow
\mathcal O_{X_K}(-m)^{\oplus M}.
\]

Still after partitioning \(T\) into pieces, one may assume in addition
that \(X_T\) (and hence
\(\mathcal H\!om_{\mathcal O_{X_T}}(E',E)\)) and all the
\[
R^q(f_T)_*\mathcal H\!om_{\mathcal O_{X_T}}(E',E),
\]
where \(f_T:X_T\to T\) is the structural morphism, are flat over \(T\),
so that formation of
\[
G=(f_T)_*\mathcal H\!om_{\mathcal O_{X_T}}(E',E)
\]
commutes with every base change \(T'\to T\) (EGA III 6.9.9.2). Let
\[
R=\mathbb V(G^*).
\]
Then \(X_R\) canonically carries a \emph{universal} exact sequence
\[
E'_R\longrightarrow E_R\longrightarrow F\longrightarrow0,
\]
and \(F\) bounds \(\mathcal F\).

\sgaStableTarget{sga6.e13.sec.2}\subsection*{2. Several technical lemmas}

\sgaStableTarget{sga6.e13.para.2-1}\textbf{2.1.} Let \(k\) be a field and \(X\) a proper \(k\)-scheme. Let
\(L_1,\ldots,L_s\) be invertible sheaves on \(X\), let \(c_1(L_i)\) be the
Chern class of \(L_i\), and let \(Y\) be a closed subscheme of \(X\) of
dimension \(\le s\). Recall that the notation
% --- current-rescribe idx639 | printed p. 626 | source-PDF p. 629 ---
\[
\bigl\langle c_1(L_1)\cdots c_1(L_s)\cdot Y\bigr\rangle
\]
denotes the intersection number, which may be calculated as the
coefficient of the monomial \(n_1\cdots n_s\) in the Snapper polynomial
\[
\chi\!\left(
L_1^{\otimes n_1}\otimes\cdots\otimes L_s^{\otimes n_s}
\otimes\mathcal O_Y
\right)
\]
(\hyperlink{sga6.e10.prop.4-3-1}{X 4.3.1}).

\sgaStableTarget{sga6.e13.para.2-2}\textbf{2.2.} From now on, suppose that \(k\) is algebraically closed,
that \(X\) is projective and endowed with a very ample invertible sheaf
\(H=\mathcal O_X(1)\), and that \(X\) is integral of dimension \(r\).

Let \(L\) be an invertible sheaf on \(X\). The integer
\[
% SGA-SRC-001427
\bigl\langle c_1(L)\cdot c_1(H)^{r-1}\bigr\rangle
\]
is called the \emph{degree} of \(L\) (relative to \(H\)) and is denoted
\(d(L)\).

In what follows we consider an invertible sheaf \(L\) and the Hilbert
polynomials
\[
\chi(L(n))=\sum_{i=0}^{r}a_i\binom{n+i}{i},
\qquad
\chi(\mathcal O_X(n))=\sum_{i=0}^{r}e_i\binom{n+i}{i}.
\]

\sgaStableTarget{sga6.e13.lem.2-3}\textbf{Lemma 2.3.} \emph{One has:}
\begin{enumerate}[label=(\roman*)]
\item \emph{\(a_r=e_r=d(H)\).}
\item \emph{If \(r\ge1\), then \(a_{r-1}-e_{r-1}=d(L)\).}
\end{enumerate}

Indeed, (i) and (ii) follow immediately by setting \(m=0,1\) in the
following polynomial:
\[
\chi\bigl(L^{\otimes m}\otimes H^{\otimes n}\bigr)
=\sum_{i+j\le r}f_{ij}
\binom{m+i}{i}\binom{n+j}{j}.
\]

\sgaStableTarget{sga6.e13.lem.2-4}\textbf{Lemma 2.4.} \emph{If \(d(L)\le0\) and \(h^0(L)\ne0\), then
\(L\simeq\mathcal O_X\); hence \(d(L)=0\) and \(h^0(L)=1\).}

Indeed, since \(X\) is integral, there is a divisor \(Y\ge0\) such that
\(L\simeq\mathcal O_X(Y)\). If \(Y>0\), one would have
\[
0\ge
\bigl\langle c_1(L)\cdot c_1(H)^{r-1}\bigr\rangle
=\bigl\langle c_1(H)^{r-1}\cdot Y\bigr\rangle
\ge1.
\]

\sgaStableTarget{sga6.e13.lem.2-5}\textbf{Lemma 2.5.} \emph{If \(d(L)\le0\), then \(L\) is a
\((0,\ldots,0,a_r)\)-sheaf (with \(r-1\) zeros).}
% --- current-rescribe idx640 | printed p. 627 | source-PDF p. 630 ---

Indeed, if \(r=0\), then \(L\) is plainly an \(a_0\)-sheaf. If
\(r\ge1\), then by \hyperlink{sga6.e13.lem.2-4}{2.4} one has \(h^0(L(-1))=0\), since
\(d(L(-1))<0\). Let \(Y\) be a hyperplane section of \(X\), and let
\[
0\longrightarrow L(-1)\longrightarrow L\longrightarrow L_Y
\longrightarrow0
\]
be an exact sequence defined by \(Y\). If \(r=1\), then by \hyperlink{sga6.e13.lem.1-7}{1.7} one has
\(h^0(L_Y(-1))=a_1\); hence \(L\) is a \((0,a_1)\)-sheaf. If \(r\ge2\),
one may take \(Y\) integral by Bertini's theorem and may thus assume
inductively that \(L_Y\) is a \((0,\ldots,0,a_r)\)-sheaf; the assertion
follows.

% SGA-SRC-001428
\sgaStableTarget{sga6.e13.lem.2-6}\textbf{Lemma 2.6.} \emph{If \(r=2\) and \(d(L)\le0\), then:}
\begin{enumerate}[label=(\roman*)]
\item \emph{For every \(n\),
\[
h^1(L(n))\le a_1(a_1+1)a_2-a_0.
\]}
\item \emph{\(a_0\le a_1(a_1+1)\).}
\end{enumerate}

Let \(Y\) be an integral hyperplane section of \(X\). By \hyperlink{sga6.e13.lem.1-7}{1.7} and \hyperlink{sga6.e13.lem.2-6}{2.6},
\(L_Y\) is a \((0,a_2)\)-sheaf with Hilbert polynomial
\[
\chi(L_Y(n))=a_1+a_2\binom{n+1}{1}.
\]
Thus, by \hyperlink{sga6.e13.thm.1-11}{1.11}, \(-a_1\ge0\), and \(L_Y\) is
\([-a_1]\)-regular.

Consider the following exact sequence:
\[
H^0(L_Y(n))\longrightarrow H^1(L(n-1))
\longrightarrow H^1(L(n))\longrightarrow H^1(L_Y(n)).
\]
For \(n\ge-a_1-1\), one has
\(h^1(L(-a_1-2))\ge h^1(L(n))\), since \(h^1(L_Y(n))=0\) by \hyperlink{sga6.e13.prop.1-3}{1.3 (i)}.
For \(n\le-1\), one has
\(h^1(L(n-1))\le h^1(L(-1))\), since \(h^0(L_Y(n))=0\) by 2.4.
Finally, for every \(n\),
\(h^0(L_Y(n-1))\le h^0(L_Y(n))\), since there is an injection
\(L_Y(n-1)\to L_Y(n)\), and
\[
h^0(L_Y(-a_1-1))=a_2\binom{-a_1}{1}+a_1,
\]
since \(h^1(L_Y(-a_1-1))=0\). Hence, for \(n\le-a_1-2\),
\[
h^1(L(n))-h^1(L(n+1))\le -a_2a_1+a_1.
\]
It follows that, for every \(n\),
\sgaStableTarget{sga6.e13.eq.2-6-1}\begin{equation}
h^1(L(n))\le
h^1(L(-a_1-2))+(-a_1-1)(-a_2a_1+a_1).
\tag{2.6.1}
\end{equation}
% --- current-rescribe idx641 | printed p. 628 | source-PDF p. 631 ---

Since \(L_Y\) is \([-a_1]\)-regular, one has
\(h^2(L(-a_1-2))=0\) by \hyperlink{sga6.e13.prop.1-4}{1.4}; hence
\[
h^1(L(-a_1-2))
=h^0(L(-a_1-2))-\chi(L(-a_1-2)).
\]
But
\[
h^0(L(-a_1-2))\le a_2\binom{-a_1}{2}
\]
by \hyperlink{sga6.e13.prop.1-8}{1.8 (i)}, since \(L\) is a \((0,0,a_2)\)-sheaf by 2.5. Therefore
\sgaStableTarget{sga6.e13.eq.2-6-2}\begin{equation}
0\le h^1(L(-a_1-2))\le -a_1(-a_1-1)-a_0.
\tag{2.6.2}
\end{equation}
Finally, \hyperlink{sga6.e13.eq.2-6-2}{2.6.2} gives (ii), and \hyperlink{sga6.e13.eq.2-6-1}{2.6.1} together with \hyperlink{sga6.e13.eq.2-6-2}{2.6.2} gives (i).

\sgaStableTarget{sga6.e13.lem.2-8}\textbf{Lemma 2.8.} \emph{Suppose that \(r=2\) and \(d(L)\le0\). Put
\(\delta=d(L)\) and \(\beta=\delta(e_2-1)+e_1\). Then:}

% SOURCE-SIC / PENDING CLAUDE: the printed source has a plus sign before
% beta(beta+1), but the displayed calculation below appears to yield a minus.
% The translation preserves the printed plus sign; no silent emendation.
\[
a_0\ge
\bigl\langle c_1(L)^2\bigr\rangle+2e_0
+\sum_{i=1}^{2}(2e_i-a_i)\binom{\delta-1+i}{i}
+\beta(\beta+1).
\]

Indeed, write
\[
\chi\bigl(L^{\otimes m}(\delta-1)\bigr)
=\alpha_2m^2+\alpha_1m+\alpha_0.
\]
Comparison of polynomials gives
\[
2\alpha_2=\bigl\langle c_1(L)^2\bigr\rangle,
\qquad
\alpha_0=\sum_{i=0}^{2}e_i\binom{\delta-1+i}{i},
\]
and
\[
\alpha_2+\alpha_1+\alpha_0
=\sum_{i=0}^{2}a_i\binom{\delta-1+i}{i}.
\]
Write
\[
\chi\bigl(L^{-1}(\delta)(n)\bigr)
=\sum_{i=0}^{2}\beta_i\binom{n+i}{i}.
\]
Then \(\beta_0=\alpha_2-\alpha_1+\alpha_0\), and by \hyperlink{sga6.e13.lem.2-3}{2.3},
\[
\beta_1=d\bigl(L^{-1}(\delta)\bigr)+e_1=\beta.
\]
Thus
\[
a_0=2\alpha_2+2\alpha_0
-\sum_{i=1}^{2}a_i\binom{\delta-1+i}{i}-\beta_0.
\]
Since \(d(L^{-1}(\delta))\le0\), one has
\(\beta_0\le\beta(\beta+1)\) by \hyperlink{sga6.e13.lem.2-6}{2.6 (ii)}; whence the assertion.

\sgaStableTarget{sga6.e13.lem.2-9}\textbf{Lemma 2.9.} \emph{Suppose that \(r\ge2\), that \(d(L)\le0\),
and that \(X\) is normal. Then:}
\begin{enumerate}[label=(\roman*)]
\item \emph{\(h^1(L(-n))=0\) for
\(n\ge2+h_2(a_{r-2},a_{r-1},a_r)\);}
\item \emph{\(h^1(L(n))\le h_r(a_{r-2},a_{r-1},a_r)\) for every
\(n\);}
\item \emph{\(|a_i|\le g_{r-i}(a_{r-2},a_{r-1},a_r)\) for
\(i=0,\ldots,r\), where
\(h_2(X_0,X_1,X_2),h_3(X_0,X_1,X_2),\ldots\) and
\(g_0(X_0,X_1,X_2),g_1(X_0,X_1,X_2),\ldots\) are universal
polynomials with coefficients in \(\mathbb Q\), independent of both
\(k\) and \(X\).}
\end{enumerate}
% --- current-rescribe idx642 | printed p. 629 | source-PDF p. 632 ---

Indeed, first consider the case \(r=2\). Put
\[
g_i(X_0,X_1,X_2)=X_{2-i}^2\quad(i=0,1,2),
\qquad
h_2(X_0,X_1,X_2)=X_1(X_1+1)X_2-X_0.
\]
Then (iii) is clear, and (ii) is \hyperlink{sga6.e13.lem.2-6}{2.6 (i)}.

Let \(X\subset P=\mathbb P_k^N\), and put
\[
\omega=
\mathcal E\!xt_{\mathcal O_P}^{N-2}
\bigl(\mathcal O_X,\Omega^N_{P/k}\bigr).
\]
Then
\[
h^1(L(-n))=h^1(\omega\otimes L^{-1}(n)),
\]
and \(\omega\) is invertible outside the finite set \(Z\) of singular
points of \(X\) (see FGA 149 \S~6). Let \(Y\) be a normal hyperplane
section of \(X\) not containing \(Z\). Then
\(\omega_Y(1)\simeq\Omega^1_{Y/k}\), and there is an exact sequence
\[
0\longrightarrow\omega\otimes L^{-1}(-1)
\longrightarrow\omega\otimes L^{-1}
\longrightarrow\Omega^1_{Y/k}\otimes L_Y^{-1}(-1)
\longrightarrow0.
\]
Now
\[
h^1\bigl(\Omega^1_{Y/k}\otimes L_Y^{-1}(n)\bigr)
=h^0(L_Y(-n))
\]
for every \(n\), and \(h^0(L_Y(-n))=0\) for \(n\ge1\) by 2.5. Thus
\(\Omega^1_{Y/k}\otimes L_Y^{-1}(-1)\) is 3-regular and, by \hyperlink{sga6.e13.prop.1-4}{1.4},
\(\omega\otimes L^{-1}\) is
\[
\bigl[3+h^1(\omega\otimes L^{-1}(2))\bigr]\text{-regular}.
\]
Consequently, \(h^1(L(-n))=0\) for
\(n\ge2+h^1(L(-2))\); this proves (i).

Now suppose that \(r\ge3\), and argue by induction on \(r\). Let \(Y\)
be a normal hyperplane section of \(X\), and consider the exact sequence
\sgaStableTarget{sga6.e13.eq.2-9-1}\begin{equation}
H^0(L_Y(-n))\longrightarrow H^1(L(-n-1))
\longrightarrow H^1(L(-n))\longrightarrow H^1(L_Y(-n)).
\tag{2.9.1}
\end{equation}
By \hyperlink{sga6.e13.lem.1-7}{1.7},
\[
\chi(L_Y(n))=
\sum_{i=0}^{r-1}a_{i+1}\binom{n+i}{i}.
\]
Thus, for \(n\ge2+h_2(a_{r-2},a_{r-1},a_r)\), the induction
hypothesis gives \(h^1(L_Y(-n))=0\); and since \(n\ge2\) by (ii) for
\(r=2\), \hyperlink{sga6.e13.lem.2-4}{2.4} gives \(h^0(L_Y(-n))=0\). Hence
\[
h^1(L(-n-1))=h^1(L(-n)).
\]
% SGA-SRC-001430 — reconstructed by mathematical inference; scan inconclusive
But for \(n\gg0\), one has \(h^1(L(-n))=0\) (see FGA 149 \S~6); hence
(i).

As induction hypothesis, suppose that the polynomials \(g_i\) and \(h_i\)
have been defined for \(i\le r-1\). For every \(n\), \hyperlink{sga6.e13.eq.2-9-1}{2.9.1} gives
\[
h^1(L(n))\le h^1(L(n-1))
+h_{r-1}(a_{r-2},a_{r-1},a_r).
\]
Therefore, by (i)
% --- current-rescribe idx643 | printed p. 630 | source-PDF p. 633 ---
already proved, it follows that for every \(m\ge-2\),
\sgaStableTarget{sga6.e13.eq.2-9-2}\begin{equation}
h^1(L(n))\le
\bigl[2+h_2(a_{r-2},a_{r-1},a_r)+m\bigr]
h_{r-1}(a_{r-2},a_{r-1},a_r)
\tag{2.9.2}
\end{equation}
for every \(n\le m\).

By \hyperlink{sga6.e13.lem.2-5}{2.5}, \(L_Y\) is a \((0,\ldots,0,a_r)\)-sheaf. Put
\[
c_i=g_{r-i}(a_{r-2},a_{r-1},a_r)
\quad(i=1,\ldots,r),
\qquad
m=P_{r-2}(c_1,\ldots,c_{r-1}).
\]
By \hyperlink{sga6.e13.thm.1-11}{1.11}, \(m\ge0\), and \(L_Y\) is \(m\)-regular. Hence, by \hyperlink{sga6.e13.prop.1-4}{1.4 (ii)},
\[
h^1(L(m-2))\ge h^1(L(n))
\qquad\text{for every }n\ge m-1.
\]
Thus, if polynomials are defined by
\[
M=P_{r-2}(g_{r-1},\ldots,g_1),
\qquad
h_r=(h_2+M)h_{r-1},
\]
then (ii) follows from 2.9.2.

Finally, since \(L_Y\) is \(m\)-regular, one has \(h^q(L(m))=0\) for
\(q\ge2\) by \hyperlink{sga6.e13.prop.1-4}{1.4 (i)}. Hence
\[
\sum_{i=0}^{r}a_i\binom{m+i}{i}
=h^0(L(m))-h^1(L(m)).
\]
But, by \hyperlink{sga6.e13.prop.1-8}{1.8 (i)},
\[
h^0(L(m))\le a_r\binom{m+r}{r}.
\]
Therefore, if one defines
\[
g_r=
\sum_{i=1}^{r}g_{r-i}\binom{M+i}{i}
+a_r\binom{M+r}{r}+h_r,
\]
then (iii) follows.

\sgaStableTarget{sga6.e13.rem.2-10}\textbf{Remark 2.10.} For fixed \(m\), consider an equality of
polynomials in \(n\):
\[
\sum_{i=0}^{r}a_{m,i}\binom{n+i}{i}
=\sum_{j=0}^{r}a_j\binom{m+n+j}{j}.
\]
One readily finds the formula
\[
a_{m,i}=
% SGA-SRC-001431 — reconstructed by mathematical inference; scan inconclusive
\sum_{j=0}^{r-i}a_{j+i}\binom{m-1+j}{j}
\]
by induction on \(r\), reducing \(a_{m,i}\) to \(a_{m,0}\), and then
setting \(n=-1\).

Similarly, one finds a polynomial with rational coefficients
\[
P_i^{(m,r)}(X_r,\ldots,X_i),
\]
increasing in \(X_i\), such that in the equality
\[
\sum_{i=0}^{r}a_i^{(m)}\binom{n+i}{i}
=\sum_{j=0}^{r}a_j\binom{mn+j}{j},
\]
% --- current-rescribe idx644 | printed p. 631 | source-PDF p. 634 ---
one has
\[
a_i^{(m)}=P_i^{(m,r)}(a_r,\ldots,a_i).
\]

\sgaStableTarget{sga6.e13.lem.2-11}\textbf{Lemma 2.11.} \emph{Let \(S\) be a noetherian scheme and \(X\)
a projective \(S\)-scheme endowed with a very ample sheaf
\(H=\mathcal O_X(1)\). Suppose that the fibers of \(X/S\) are
geometrically integral and all have the same dimension \(r\). Let
\(\mathcal L\) be a family of classes of invertible sheaves \(L_K\) on the
fibers of \(X/S\). For \(\mathcal L\) to be bounded, it is necessary and
sufficient that all the coefficients \(a_i\) in the Hilbert polynomials}
\[
\chi(L_K(n))=\sum_{i=0}^{r}a_i\binom{n+i}{i}
\]
\emph{of the \(L_K\in\mathcal L\) remain bounded. Moreover, if the
fibers of \(X/S\) are also geometrically normal, then it is enough that
the three coefficients \(a_r,a_{r-1},a_{r-2}\) remain bounded.}

Indeed, let
% SGA-SRC-001432
\[
\delta=\sup\{d(L_K)\}.
\]
If \(\delta>0\), one may replace \(\mathcal L\) by the family of classes
of the \(L_K(-\delta)\), by \hyperlink{sga6.e13.lem.2-3}{2.3} and \hyperlink{sga6.e13.rem.2-10}{2.10}; thus one may assume
\(\delta\le0\). Then \(\mathcal L\) is a \((b)\)-family by \hyperlink{sga6.e13.lem.2-5}{2.5}, and is
therefore bounded if all the coefficients \(a_i\) remain bounded, by
1.13. If the fibers of \(X/S\) are geometrically normal, it is enough
for this purpose that the three coefficients \(a_r,a_{r-1},a_{r-2}\)
remain bounded, by 2.9.

\sgaStableTarget{sga6.e13.sec.3}\subsection*{3. General finiteness theorems}

\sgaStableTarget{sga6.e13.lem.3-1}\textbf{Lemma 3.1.} \emph{Let \(S\) be a scheme and \(X\) a
quasi-compact and quasi-separated \(S\)-scheme. Let
\(\underline{\Pic}_{X/S}\) be the Picard functor, localized for the fppf
topology.}\footnote{Cf. SGA 3 IV 6.3 and FGA no.~232 for the definition of
\(\underline{\Pic}_{X/S}\); here, as indicated, the fpqc topology should be
replaced by the fppf topology.}
% --- current-rescribe idx645 | printed p. 632 | source-PDF p. 635 ---
\emph{Then, for every filtered inverse system \((Z_\alpha)\) of
quasi-compact and quasi-separated \(S\)-schemes with affine transition
morphisms, the canonical map}
\sgaStableTarget{sga6.e13.eq.3-1-1}\begin{equation}
\varinjlim_\alpha\underline{\Pic}_{X/S}(Z_\alpha)
\longrightarrow
\underline{\Pic}_{X/S}\!\left(\varprojlim_\alpha Z_\alpha\right)
\tag{3.1.1}
\end{equation}
\emph{is bijective.}

\emph{Consequently} (EGA IV 8.14.2), \emph{if the Picard functor is
represented by an \(S\)-scheme \(\Pic_{X/S}\), then \(\Pic_{X/S}\) is
locally of finite presentation.}

Indeed, an element \(\lambda_\alpha\) of
\(\underline{\Pic}_{X/S}(Z_\alpha)\) is represented by an invertible sheaf
% SGA-SRC-001433
\(L'_\alpha\) on \(X_{Z'_\alpha}\), where
\(Z'_\alpha\to Z_\alpha\) is an fppf morphism. Put
\[
Z=\varprojlim_\alpha Z_\alpha.
\]
The image \(\lambda\) of \(\lambda_\alpha\) in
\(\underline{\Pic}_{X/S}(Z)\) is represented by the invertible sheaf
\(L'\) on \(X_{Z'}\) obtained by inverse image from \(L'_\alpha\), where
\[
Z'=Z'_\alpha\times_{Z_\alpha}Z.
\]

If \(\lambda\) is the neutral element of
\(\underline{\Pic}_{X/S}(Z)\), there exists an fppf morphism
\(Z''\to Z'\) such that the inverse image of \(L'\) on \(X_{Z''}\) is
isomorphic to \(\mathcal O_{X_{Z''}}\). By EGA IV 8.8.2, 8.10.5 (vi),
11.2.6 (ii), 8.5.2 (i), and 8.5.2.4, there exist
\(\beta\ge\alpha\) and an fppf morphism
\[
Z''_\beta\longrightarrow Z'_\beta,
\qquad
Z'_\beta=Z'_\alpha\times_{Z_\alpha}Z_\beta,
\]
such that the inverse image of \(L'_\alpha\) on \(X_{Z''_\beta}\) is
isomorphic to \(\mathcal O_{X_{Z''_\beta}}\). The image of
\(\lambda_\alpha\) is therefore already the neutral element in
\(\underline{\Pic}_{X/S}(Z_\beta)\). Thus \hyperlink{sga6.e13.eq.3-1-1}{3.1.1} is injective.

The surjectivity of \hyperlink{sga6.e13.eq.3-1-1}{3.1.1} is proved similarly.

\sgaStableTarget{sga6.e13.prop.3-2}\textbf{Proposition 3.2.} \emph{Let \(S\) be a noetherian scheme and
\(X\) a proper \(S\)-scheme. Suppose that the Picard functor is
represented by an \(S\)-scheme \(\Pic_{X/S}\). Then:}
\begin{enumerate}[label=(\roman*)]
\item \emph{\(\Pic_{X/S}\) is locally of finite type over \(S\).}
% --- current-rescribe idx646 | printed p. 633 | source-PDF p. 636 ---
\item \emph{Subsets \(\Lambda\) of \(\Pic_{X/S}\) correspond bijectively
to families \(\mathscr L\) of classes of invertible sheaves on the fibers
of \(X/S\).}
\item \emph{A subset \(\Lambda\) of \(\Pic_{X/S}\) is quasi-compact if
% SOURCE-NORMALIZATION / PENDING CLAUDE: print and workpass say "a family
% Lambda of Pic" here, but item (ii) and the proof both treat Lambda as a subset.
and only if the corresponding family \(\mathscr L\) is bounded.}
\end{enumerate}

Indeed, (i) follows from 3.1. Assertion (ii) follows from the following
fact: if \(s\) is a point of \(S\) and \(K\) is an algebraically closed
extension of \(k(s)\), then the \(K\)-points of \(\Pic_{X/S}\) correspond
bijectively to the isomorphism classes of invertible sheaves on \(X_K\).

For (iii), suppose that the subset \(\Lambda\) of \(\Pic_{X/S}\) is
quasi-compact. Then \(\Lambda\) is contained in an open subscheme \(U\)
of finite type, and the inclusion \(U\to\Pic_{X/S}\) corresponds to an
invertible sheaf \(L\) on some \(X_T\), where \(T\to U\) is a surjective
(flat) morphism of finite type. Finally \(T\to S\) is of finite type, so
\(L\) bounds \(\mathscr L\).

Conversely, suppose that \(\mathscr L\) is bounded by a coherent sheaf
\(L\) on some \(X_T\), with \(T\) of finite type over \(S\). One may
assume that \(L\) is invertible \hyperlink{sga6.e13.thm.1-13}{(1.13 (i))}, so that \(L\) defines a
morphism
\[
f:T\longrightarrow\Pic_{X/S}.
\]
The image of \(f\) contains \(\Lambda\) and is quasi-compact, hence
noetherian by (i). It follows that \(\Lambda\) is quasi-compact, proving
(iii).

By \hyperlink{sga6.e13.prop.3-2}{3.2}, the usual definition of morphisms of finite type between two
Picard schemes generalizes as follows.

\sgaStableTarget{sga6.e13.def.3-3}\textbf{Definition 3.3.} \emph{Let \(X\) and \(Y\) be proper over a
noetherian scheme \(S\). A morphism between the Picard functors}
\[
\begin{tikzcd}
\underline{\Pic}_{Y/S}\arrow[r]&\underline{\Pic}_{X/S}
\end{tikzcd}
\]

% End idx646 / printed p. 633 / source-PDF p. 636.
% Definition 3.3 continues at idx647: the morphism is said to be of finite
% type if the inverse image of every bounded family on X/S is a bounded
% family on Y/S. No idx647 source text is promoted in this fragment.

% END INCLUDED SGA6_English_current_buildable_source_20260823/SGA6_idx636_646_English_SourceChecked_fragment.tex


% END INCLUDED SGA6_English_current_buildable_source_20260823/SGA6_idx608_646_English_SourceChecked_fragment.tex


% BEGIN INCLUDED SGA6_English_current_buildable_source_20260823/SGA6_idx647_665_English_SourceChecked_fragment.tex
% SGA 6 English source-checked fragment: current-rescribe idx647--665.
% French control for idx647--662: predecessor workpass commit 8ccdcf8ee.
% French control for idx663--665: source-rescribe workpass commit
% 859bd5abc3a813ee51e7cbcc6d623e08c1336235, release-TeX SHA-256
% 7F32C2080A78A2746CBE52DCC1EC43A8505269F25518FA7B9A86E4E89AF858AC.
% English idx663--665 is reconciled to that working control and the scan;
% this is not an unqualified whole-volume certification.
% Ultimate check: sga6.pdf. Printed pp. 634--652; source-PDF pp. 637--655.
% Begins in Definition XIII.3.3; ends cleanly after XIII.5.4.

% --- current-rescribe idx647 | printed p. 634 | source-PDF p. 637 ---
\emph{is said to be of finite type if every bounded family of classes of
invertible sheaves on the fibers of \(X/S\) has as inverse image a bounded
family of classes of invertible sheaves on the fibers of \(Y/S\).}

\sgaStableTarget{sga6.e13.lem.3-4}\textbf{Lemma 3.4.} \emph{Let \(X,Y,X',Y'\) be proper over a noetherian
scheme \(S\), and let \(T\to S\) be surjective and of finite type. Then:}
\begin{enumerate}[label=(\roman*)]
\item \emph{A morphism
\(\Phi:\underline{\Pic}_{Y/S}\to\underline{\Pic}_{X/S}\)
is of finite type if and only if
\(\Phi_T:\underline{\Pic}_{Y_T/T}\to\underline{\Pic}_{X_T/T}\)
is of finite type.}
\item \emph{In a commutative diagram}
\[
\begin{tikzcd}
\underline{\Pic}_{Y'/S}\arrow[r,"\Phi'"]&
\underline{\Pic}_{X'/S}\\
\underline{\Pic}_{Y/S}\arrow[r,"\Phi"']\arrow[u,"\Phi_Y"]&
\underline{\Pic}_{X/S}\arrow[u,"\Phi_X"']
\end{tikzcd}
\]
\emph{in which \(\Phi'\) and \(\Phi_Y\) are of finite type, \(\Phi\) is
of finite type.}
\end{enumerate}

Suppose that \(\Phi_T\) is of finite type. Let \(\mathscr L\) be a family
on the fibers of \(X/S\), bounded by an invertible sheaf \(L\) on some
\(X_R\). The family \(\mathscr L_T=\mathscr L\times_S T\) on the fibers of
\(X_T/T\) is bounded by the inverse image of \(L\) on
\(X_{R\times_S T}\). If \(\Phi_T^{-1}(\mathscr L_T)\) is bounded by an
invertible sheaf \(M\) on some \(Y_Q\), then \(M\) also bounds
\(\Phi^{-1}(\mathscr L)\). The converse in (i), and assertion (ii), are
proved similarly.

\sgaStableTarget{sga6.e13.thm.3-5}\textbf{Theorem 3.5.} \emph{Let \(S\) be a noetherian scheme, let \(X\)
and \(Y\) be proper \(S\)-schemes, and let \(f:X\to Y\) be a morphism. If
\(f\) is surjective, then the induced morphism between the Picard
functors}
\[
f^*:\underline{\Pic}_{Y/S}\longrightarrow\underline{\Pic}_{X/S}
\]

% --- current-rescribe idx648 | printed p. 635 | source-PDF p. 638 ---
\emph{is of finite type.}

By \hyperlink{sga6.e13.lem.3-4}{3.4 (i)}, one may assume that \(S\) is reduced and irreducible. There
is a nonempty open subset \(U\subset S\) such that
\(\underline{\Pic}_{Y/S}\) and \(\underline{\Pic}_{X/S}\) are
representable and the restriction of \(f^*\) is quasi-affine (\hyperlink{sga6.e12.thm.1-1}{XII 1.1}),
hence of finite type by \hyperlink{sga6.e13.prop.3-2}{3.2 (i)}. Give \(T=S-U\) its induced reduced
structure. By noetherian induction, one may suppose that
\[
\underline{\Pic}_{Y_T/T}\longrightarrow\underline{\Pic}_{X_T/T}
\]
induced by \(f^*\) is of finite type. The assertion follows from \hyperlink{sga6.e13.lem.3-4}{3.4 (i)},
applied to \(U\amalg T\to S\).

\sgaStableTarget{sga6.e13.thm.3-6}\textbf{Theorem 3.6.} \emph{Let \(S\) be a noetherian scheme, let \(X\)
be a proper \(S\)-scheme, and let \(\underline{\Pic}_{X/S}\) be its
Picard functor. For every nonzero integer \(n\), the \(n\)-th power
homomorphism}
\[
\Phi_n:\underline{\Pic}_{X/S}\longrightarrow\underline{\Pic}_{X/S}
\]
\emph{is a morphism of finite type.}

Let \(f:X'\to X\) be a Chow presentation: \(X'\) is projective over \(S\),
and \(f\) is surjective. If the theorem is known for \(X'\), then \hyperlink{sga6.e13.thm.3-5}{3.5} and
\hyperlink{sga6.e13.lem.3-4}{3.4 (ii)} imply it for \(X\). We may suppose \(X/S\) projective.

There are a nonempty open \(U\subset S\), an integral \(S'\), and a
finite surjective morphism \(S'\to U\), such that every irreducible
component \(X'_i\) of \(X_{S'}\), with its induced reduced structure, is
integral over \(S'\) (\hyperlink{sga6.e12.lem.1-3}{XII 1.3}). Give \(T=S-U\) its induced reduced
structure. By noetherian induction, the theorem may be assumed for
\(X_T/T\). If it is assumed for all \(X_i/S'\), it follows for
\(\coprod_iX_i/S'\) by the

% --- current-rescribe idx649 | printed p. 636 | source-PDF p. 639 ---
following easy lemma, then for \(X_{S'}/S'\) by \hyperlink{sga6.e13.thm.3-5}{3.5} and \hyperlink{sga6.e13.lem.3-4}{3.4 (ii)}, and
finally for \(X/S\) by \hyperlink{sga6.e13.lem.3-4}{3.4 (i)}.

\sgaStableTarget{sga6.e13.lem.3-7}\textbf{Lemma 3.7.} \emph{Let \(X=\coprod_iX_i\) be of finite type over a
noetherian scheme \(S\), and let \(\mathcal F\) be a family of classes of
coherent sheaves on the fibers of \(X/S\). Then \(\mathcal F\) is bounded
if and only if, for every \(i\), the induced family \(\mathcal F_i\) on
the fibers of \(X_i/S\) is bounded.}

In the proof of \hyperlink{sga6.e13.thm.3-6}{3.6} we reduced to \(X\) projective and integral over \(S\);
endow \(X/S\) with a very ample sheaf \(\mathcal O_X(1)\). Let
\(\mathscr L\) be bounded on the fibers of \(X/S\), and let \(L\) on some
\(X_T\), with \(T\) of finite type over \(S\), bound it. After
partitioning \(T\), suppose \(L\) flat over \(T\); after taking connected
components, suppose \(T\) connected. Then
\[
P(q,m)=\chi\bigl(L_t^{\otimes q}(m)\bigr)
\]
is independent of \(t\in T\) (EGA III 7.9.4).

Put \(\mathscr L_n=\Phi_n^{-1}(\mathscr L)\). It consists of classes of
invertible sheaves \(M_K\) whose \(n\)-th power class belongs to
\(\mathscr L\). Write
\[
\chi\bigl(M_K^{\otimes p}(m)\bigr)
=\sum_{i=0}^{r}a_i(p)\binom{m+i}{i},
\]
where \(a_i(p)\) has degree at most \(r-i\) in \(p\). Since
\(\chi(M_K^{\otimes qn}(m))=P(q,m)\), every \(a_i(qn)\), and hence every
polynomial \(a_i(p)\), is independent of \(M_K\). Thus \(\mathscr L_n\)
has one Hilbert polynomial \(\chi(M_K(m))\), and is bounded by 2.11.

\sgaStableTarget{sga6.e13.thm.3-8}\textbf{Theorem 3.8.} \emph{Let \(S\) be noetherian, let \(Y\) be a
projective \(S\)-scheme with a relatively ample invertible sheaf
\(\mathcal O_Y(1)\), let \(X\) be the zero scheme of a section \(\varphi\)
of \(\mathcal O_Y(1)\), and let \(f:X\to Y\) be the inclusion. Assume}

% --- current-rescribe idx650 | printed p. 637 | source-PDF p. 640 ---
\emph{that for every \(s\in S\) and every irreducible component \(Y_s^i\)
of \(Y_s\), the irreducible components of
\(X_s\cap Y_s^i=V(\varphi|_{Y_s^i})\) have dimension at least \(2\)
(as happens when the irreducible components of the fibers of \(Y\) have
dimension at least \(3\)). Then}
\[
f^*:\underline{\Pic}_{Y/S}\longrightarrow\underline{\Pic}_{X/S}
\]
\emph{is of finite type. Moreover, if \(S\) is integral, there is a
nonempty open \(U\subset S\) such that the restricted Picard functors
are representable and}
\[
\underline{\Pic}_{Y/S}|_U\longrightarrow\underline{\Pic}_{X/S}|_U
\]
\emph{is quasi-affine (and of finite type).}

For the first assertion, noetherian induction permits replacing \(S\) by
any nonempty open. One reduces to \(S\) integral and \(Y\) integral over
\(S\), replacing \(X\) by \(X\times_Y Y'\) when \(Y\) is replaced by
\(Y'\). One then reduces to \(Y\) normal over \(S\), using the following
lemma in place of \hyperlink{sga6.e12.lem.1-3}{XII 1.3}.

\sgaStableTarget{sga6.e13.lem.3-9}\textbf{Lemma 3.9.} \emph{Let \(S\) be integral noetherian and \(X\) an
\(S\)-scheme integral over \(S\). There are a nonempty open \(U\subset S\),
an integral scheme \(S'\) finite faithfully flat over \(U\), an
\(S'\)-scheme \(X'\) normal and integral over \(S'\), and a finite
surjective \(S'\)-morphism \(X'\to X_{S'}\).}

% --- current-rescribe idx651 | printed p. 638 | source-PDF p. 641 ---
Let \(\eta\) be the generic point of \(S\), and let \(K\) be an algebraic
closure of \(k(\eta)\). The scheme \(\operatorname{Spec}(K)\) is the
limit of the filtered inverse system \((S_\alpha)\) of integral affine
schemes with finite faithfully flat morphisms
\(S_\alpha\to U_\alpha\), where \(U_\alpha\) is a nonempty open of \(S\).
The normalization of \(X_K\) descends to an \(S_\alpha\)-scheme \(X'\) for
some \(\alpha\), and the finite surjective map \(X'_K\to X_K\) descends
% SOURCE-CORRECTION / PENDING CLAUDE: print and workpass have the undefined
% base S'_\alpha at this point; context requires the already defined S_\alpha.
to a finite surjective \(S_\alpha\)-morphism \(X'\to X_{S_\alpha}\)
(EGA IV 8.8.2, 8.10.5). Since \(X'_K\) is normal and integral, there is
an open \(S'\subset S_\alpha\) such that \(X'_{S'}\) is normal and
integral over \(S'\) (EGA IV 9.9.5).

\sgaStableTarget{sga6.e13.rem.3-10}\textbf{Remark 3.10.} Retain \hyperlink{sga6.e13.lem.3-9}{3.9} and suppose there is a resolving
morphism \(X'_K\to X_K\) (EGA IV 7.9.1), as happens when
% WORKPASS-CORRECTION / PENDING CLAUDE: commit 8ccdcf8ee changes the scan's
% spelling "Abhyankar" to the erroneous "Abhyankhar"; English follows print.
\(\dim(X_K)\le2\), by Abhyankar. A similar argument gives a finite
faithfully flat \(S'\to U\), a scheme \(X'\) smooth over \(S'\), and a
proper surjective birational morphism \(X'\to X_{S'}\).

In the proof of \hyperlink{sga6.e13.thm.3-8}{3.8} we have reduced to \(S\) integral and \(Y\) normal and
integral over \(S\). If \(\varphi=0\), then \(X=Y\). Otherwise, after
restricting \(S\), suppose that \(\varphi_s\ne0\) for every \(s\in S\).
Thus \(\varphi_s\) is \(\mathcal O_{Y_s}\)-regular, and the dimension
\(r\) of \(Y_s\) is at least \(3\).

Let \(\mathscr L\) be bounded on the fibers of \(X/S\), and let \(M_K\)
be invertible on a geometric fiber \(Y_K\), with
\[
\chi(M_K(n))=\sum_{i=0}^{r}a_i\binom{n+i}{i}.
\]
The Hilbert polynomial of \(f^*M_K\) is
\[
\sum_{i=0}^{r-1}a_{i+1}\binom{n+i}{i}.
\]
As the class of \(M_K\) varies through \((f^*)^{-1}(\mathscr L)\), that
is, as

% --- current-rescribe idx652 | printed p. 639 | source-PDF p. 642 ---
the class of \(f^*M_K\) varies through \(\mathscr L\), the coefficients
\(a_r,a_{r-1},a_{r-2}\) remain bounded by 1.13.

Let \(H=\mathcal O_Y(m)\) be a very ample multiple of
\(\mathcal O_Y(1)\). In the Hilbert polynomials relative to \(H\),
\[
\chi(M_K(mn))=\sum_{i=0}^{r}a_i^{(m)}\binom{n+i}{i},
\]
the coefficients \(a_r^{(m)},a_{r-1}^{(m)},a_{r-2}^{(m)}\) remain bounded
by 2.10. Consequently \((f^*)^{-1}(\mathscr L)\) is bounded by \hyperlink{sga6.e13.lem.2-11}{2.11},
% SOURCE-CORRECTION / PENDING CLAUDE: print and workpass say merely that
% \mathscr L is bounded; the finite-type claim requires its inverse image.
proving the first part of 3.8.

For the second part, generic representability and separatedness of the
Picard functors were established in \hyperlink{sga6.e12.cor.1-2}{XII 1.2}. An argument like XII 2,
using \hyperlink{sga6.e13.lem.3-9}{3.9} and \hyperlink{sga6.e12.thm.1-1}{XII 1.1}, reduces to \(Y\) normal over \(S\). It is enough
for \(\underline{\Pic}_{Y/S}\to\underline{\Pic}_{X/S}\) to be separated
and quasi-finite (EGA IV 8.12.8), which can be checked on geometric
fibers and follows from:

\sgaStableTarget{sga6.e13.lem.3-11}\textbf{Lemma 3.11.} \emph{Let \(k\) be algebraically closed, let \(Y\)
be normal and projective with an ample sheaf \(\mathcal O_Y(1)\), and let
\(X\) be the zero scheme of a section \(\varphi\). Suppose
\(\dim(Y)\ge3\) (respectively, \(\dim(Y)\ge2\)). Then}
\[
\Pic_{Y/k}\longrightarrow\Pic_{X/k}
\quad
\text{(respectively, }\Pic^\tau_{Y/k}\longrightarrow\Pic^\tau_{X/k}
\text{; see \hyperlink{sga6.e13.para.4-1}{4.1})}
\]
\emph{is of finite type and has a finite kernel \(N\), in general
nonreduced.}

Assume \(\varphi\ne0\); it is regular because \(Y\) is integral. The map
is of finite type by \hyperlink{sga6.e13.thm.3-8}{3.8} (respectively,\hyperlink{sga6.e13.thm.4-7}{ 4.7 (iii)}). Its kernel \(N\) is
an algebraic group contained in \(\Pic^\tau_{Y/k}\), hence is the same

% --- current-rescribe idx653 | printed p. 640 | source-PDF p. 643 ---
in both cases. Since \(\Pic^\tau_{Y/k}\) is proper because \(Y\) is
normal (FGA no.~236, 2.13), and \(N\) is closed (SGA 3 VI\(_B\) 1.9.2),
it is enough to show that \(N\) contains no subgroup of type
\(\mu_\ell\;(=\mathbb Z/\ell\mathbb Z)\), with
\((\ell,\operatorname{char}(k))=1\).

Consider Kummer theory, arising from
\[
0\longrightarrow\mu_\ell\longrightarrow\mathbb G_m
\xrightarrow{x\mapsto x^\ell}\mathbb G_m\longrightarrow0
\]
on \(Y\) (respectively, on \(X\)):
\[
\begin{tikzcd}[column sep=1.6em,row sep=1.2em]
&&H^1(Y,\mu_\ell)\arrow[r]\arrow[d]&
H^1(Y,\mathbb G_m)\arrow[r,"\ell"]\arrow[d]&
H^1(Y,\mathbb G_m)\arrow[d]\\
H^0(X,\mathbb G_m)\arrow[r,"\ell"]&
H^0(X,\mathbb G_m)\arrow[r]&
H^1(X,\mu_\ell)\arrow[r]&
H^1(X,\mathbb G_m)\arrow[r,"\ell"]&
H^1(X,\mathbb G_m).
\end{tikzcd}
\]
By Hilbert's Theorem 90, \(H^1(Y,\mathbb G_m)=\Pic_{Y/k}(k)\)
(respectively, \(H^1(X,\mathbb G_m)=\Pic_{X/k}(k)\)). Since \(X\) is
proper and integral, \(H^0(X,\mathbb G_m)=k^*\), so
\(H^1(X,\mu_\ell)\to H^1(X,\mathbb G_m)\) is injective.

It remains to show that
\[
H^1(Y,\mathbb Z/\ell\mathbb Z)\longrightarrow
H^1(X,\mathbb Z/\ell\mathbb Z)
\]
is injective. This follows from \(\pi_1(X,\xi)\to\pi_1(Y,\xi)\)
surjective; equivalently (SGA 1 V 6.9), every connected \'etale cover
\(\widetilde Y\) of \(Y\) with group \(\mathbb Z/\ell\mathbb Z\) has
connected restriction \(\widetilde X=X\times_Y\widetilde Y\). This is a
consequence of Bertini and Zariski connectedness (SGA 1 X 2.10).

\sgaStableTarget{sga6.e13.rem.3-12}\textbf{Remark 3.12.} Retain the hypotheses of 3.11.
\begin{enumerate}[label=(\roman*)]
\item The kernel \(N\) is finite \emph{unipotent} (SGA 3 XVII 1.1),
since it contains no subgroup of multiplicative type \(\mu_n\), for any
\(n\) (SGA 3 XVII 4.6.1 (vi)). This is proved similarly, using instead

% --- current-rescribe idx654 | printed p. 641 | source-PDF p. 644 ---
of Kummer theory the following fact of Grothendieck (SGA 1 XI, last
page): if \(X\) is proper over \(k\), with \(H^0(X,\mathcal O_X)=k\),
then there is an isomorphism of \(k\)-group schemes
\[
\underline{\operatorname{Hom}}_{\mathrm{grp}}
(\mu_n,\underline{\Pic}_{X/k})
\simeq \underline{H^1(X,\mathbb Z/n\mathbb Z)}.
\]
Consequently, when \(k\) is algebraically closed, its \(k\)-points give
\[
\operatorname{Hom}_{\mathrm{grp}}(\mu_n,\underline{\Pic}_{X/k})
\simeq H^1(X,\mathbb Z/n\mathbb Z).
\]
\item If \(k\) has characteristic zero, then \(N=0\), since every finite
unipotent group is zero (SGA 3 XVII 1.7). This approximately recovers
Kodaira's vanishing theorem in Mumford's form [2].
\end{enumerate}

\sgaStableTarget{sga6.e13.thm.3-13}\textbf{Theorem 3.13.} \emph{Let \(S\) be noetherian and \(X\) a
projective \(S\)-scheme with a relatively ample invertible sheaf
\(H=\mathcal O_X(1)\). Suppose the geometric fibers of \(X/S\) are
integral of dimension \(r\). For a family \(\mathscr L\) of classes of
invertible sheaves on the fibers, the following are equivalent:}
\begin{enumerate}[label=(\roman*)]
\item \emph{\(\mathscr L\) is bounded.}
\item \emph{In
\(\chi(L_K(n))=\sum_{i=0}^{r}a_i\binom{n+i}{i}\),
the coefficient \(a_{r-1}\) remains bounded and \(a_{r-2}\) remains
bounded below.}
\item \emph{The degree
\[
d(L_K)=\langle c_1(L_K)c_1(H_K)^{r-1}\rangle
\]
remains bounded and
\(\langle c_1(L_K)^2c_1(H_K)^{r-2}\rangle\) remains bounded below.}
\item \emph{There is an integer \(N\) such that
\[
H^{-N}<\mathscr L<H^N;
\]
that is, \(H_K^{\otimes N}\otimes L_K\) and
\(H_K^{\otimes N}\otimes L_K^{-1}\) are ample for every
\(L_K\in\mathscr L\).}
\end{enumerate}

% --- current-rescribe idx655 | printed p. 642 | source-PDF p. 645 ---
Suppose \(\mathscr L\) is bounded by \(L\) on some \(X_T\). Then (iv)
holds by EGA II 4.6.12. Taking \(L\) flat over \(T\), the polynomial
\(\chi(L_t^{\otimes m}(n))\) depends only on the connected component of
\(T\) (EGA III 7.9.4). Hence (ii) holds, and \hyperlink{sga6.e13.para.2-1}{2.1} gives (iii).

The families of classes of \(H_K\) and \(\mathcal O_{X_K}\) are bounded,
so they satisfy (iii) and (ii), respectively. If \(\mathscr L\) satisfies
(iv), then
\[
\langle[Nc_1(H_K)+c_1(L_K)]c_1(H_K)^{r-1}\rangle>0,
\]
so \(-Nd(H_K)<d(L_K)<Nd(H_K)\). Moreover,
\[
\langle[Nc_1(H_K)+c_1(L_K)]^2c_1(H_K)^{r-2}\rangle>0,
\]
whence
\[
\langle c_1(L_K)^2c_1(H_K)^{r-2}\rangle
>-2Nd(L_K)-N^2d(H_K).
\]
Thus (iii) holds.

Finally, suppose (ii), respectively (iii), holds. By \hyperlink{sga6.e13.rem.2-10}{2.10} and \hyperlink{sga6.e13.lem.2-3}{2.3},
replace \(H\) by a multiple and assume it very ample. Partition \(S\) so
there is a linear section \(Y\) of \(X\), integral over \(S\), with
two-dimensional fibers. By \hyperlink{sga6.e13.thm.3-8}{3.8}, replace \(X,\mathscr L\) by
\(Y,\mathscr L|_Y\); \hyperlink{sga6.e13.lem.1-7}{1.7}, respectively the projection formula, preserves
the relevant condition. Thus \(r=2\). Put
\[
\delta=\sup\{d(L_K)\}.
\]
If \(\delta>0\), replace \(\mathscr L\) by the classes of
\(L_K(-\delta)\), using \hyperlink{sga6.e13.lem.2-3}{2.3} and \hyperlink{sga6.e13.rem.2-10}{2.10}, so that \(\delta\le0\).

If \(a_1\) remains bounded and \(a_0\) bounded below in
\(\sum_{i=0}^{2}a_i\binom{n+i}{i}\), all three coefficients remain
bounded by \hyperlink{sga6.e13.lem.2-3}{2.3} and 2.7. Thus \(\mathscr L\) is bounded by \hyperlink{sga6.e13.lem.2-11}{2.11}, so (ii)
implies

% --- current-rescribe idx656 | printed p. 643 | source-PDF p. 646 ---
(i). On the other hand, (iii) implies (ii) by \hyperlink{sga6.e13.lem.2-3}{2.3} and 2.8.

\sgaStableTarget{sga6.e13.sec.4}\subsection*{4. Finiteness theorems for \(\Pic^\tau_{X/S}\)}

\sgaStableTarget{sga6.e13.para.4-1}\textbf{4.1.} Let \(S\) be a scheme and \(G\) a commutative group functor
over \(S\). Denote by \(G^0\), respectively \(G^\tau\), the subgroup
functor whose \(T\)-points \(f:T\to G\) satisfy: for every algebraically
closed field \(K\) and \(K\)-point \(t\) of \(T\), there are a connected
% SGA-SRC-001414
\(K\)-scheme \(Z\), two \(K\)-points \(z,z_0\), and \(g:Z\to G\), with
\(g(z)=f(t)\) and \(g(z_0)=0\) in \(G(K)\); respectively, there is a
nonzero \(n\) such that \(nf(t)\in G^0(K)\).

Formation of \(G^0\), respectively \(G^\tau\), commutes with base change,
and \(G\to G'\) induces \(G^0\to G'^0\), respectively
\(G^\tau\to G'^\tau\). If \(S\) is a field and \(G\) is represented by a
group scheme locally of finite type, \(G^0\) is represented by the
connected component of the identity, an open and closed finite-type
geometrically irreducible subscheme (SGA 3 VI\(_A\) 2.4).

\sgaStableTarget{sga6.e13.lem.4-2}\textbf{Lemma 4.2.} \emph{Let \(\Phi:G\to G'\) be a morphism of
commutative group functors over \(S\). Suppose that for every
algebraically closed \(K\) over \(S\), \(G_K,G'_K\) are represented by
group schemes locally of finite type and \(\Phi_K\) by a finite-type
(equivalently quasi-compact) morphism. Then}
\[
\Phi^{-1}(G'^\tau)=G^\tau.
\]

% --- current-rescribe idx657 | printed p. 644 | source-PDF p. 647 ---
% WORKPASS-CORRECTION / PENDING CLAUDE: print has Phi^{-1}(G'^T) superset
% G^T; the certified workpass instead has the ill-typed Phi(G'^tau).
The inclusion \(\Phi^{-1}(G'^\tau)\supset G^\tau\) has already been
noted. Conversely, \(g\in G(K)\) belongs to \(G^\tau(K)\) if and only if
the set of multiples \(ng\), \(n\in\mathbb N\), is quasi-compact. Since
\(\Phi_K\) is quasi-compact, it is enough that the set
\(\Phi(ng)=n\Phi(g)\) have this property. This proves the lemma.

\sgaStableTarget{sga6.e13.lem.4-3}\textbf{Lemma 4.3.} \emph{Let \(X,Y\) be proper over \(S\), and let
\(f:X\to Y\) be surjective (respectively, the inclusion of the zero
scheme of a section \(\varphi\) of an ample \(\mathcal O_Y(1)\)). Then}
\[
(f^*)^{-1}\underline{\Pic}^\tau_{X/S}
=\underline{\Pic}^\tau_{Y/S}.
\]
This follows formally from \hyperlink{sga6.e13.lem.4-2}{4.2}, \hyperlink{sga6.e12.cor.1-2}{XII 1.2}, \hyperlink{sga6.e13.lem.3-1}{3.1}, and \hyperlink{sga6.e13.thm.3-5}{3.5} (respectively,
\hyperlink{sga6.e13.thm.3-8}{3.8}).

\sgaStableTarget{sga6.e13.para.4-4}\textbf{4.4.} Let \(k\) be algebraically closed, \(X\) proper over \(k\),
and \(L\) invertible on \(X\). The sheaf \(L\) is \emph{algebraically
equivalent} to zero if there are a connected finite-type \(k\)-scheme
\(T\), an invertible sheaf \(M\) on \(X_T=X\times_kT\), and two
\(k\)-points \(t,t_0\), such that \(L\simeq M_t\) and
\(\mathcal O_X\simeq M_{t_0}\). It is \emph{\(\tau\)-equivalent} to zero
if some nonzero tensor power is algebraically equivalent to zero.
Equivalently, the point \(\lambda\) corresponding to \(L\) belongs to
\(\underline{\Pic}^0_{X/k}(k)\), respectively
\(\underline{\Pic}^\tau_{X/k}(k)\).

Recall that \(L\) is \emph{numerically equivalent} to zero if
\[
\langle c_1(L)\cdot Y\rangle=0
\]
for every closed integral curve \(Y\subset X\) (see \hyperlink{sga6.e13.para.2-1}{2.1} or \hyperlink{sga6.e10.sec.4.5}{X 4.5}).

% --- current-rescribe idx658 | printed p. 645 | source-PDF p. 648 ---
\sgaStableTarget{sga6.e13.prop.4-5}\textbf{Proposition 4.5 (Functorial behavior).} \emph{Let
\(f:X'\to X\) be a morphism between proper \(k\)-schemes, and let \(L\)
be invertible on \(X\). Then:}
\begin{enumerate}[label=(\roman*)]
\item \emph{If \(L\) is \(\tau\)-equivalent to zero (respectively,
numerically equivalent to zero), then so is \(f^*L\).}
\item \emph{Conversely, if \(f\) is surjective and \(f^*L\) is
\(\tau\)-equivalent to zero (respectively, numerically equivalent to
zero), then so is \(L\).}
\end{enumerate}

The assertions for \(\tau\)-equivalence follow from \hyperlink{sga6.e13.para.4-1}{4.1} and 4.3. Let
\(Y'\) be a closed integral curve in \(X'\), and let \(Y=f(Y')\), with
its induced reduced structure. If \(\dim(Y)=1\), the projection formula
(\hyperlink{sga6.e10.eq.4-4-4}{X 4.4.4}) gives
\sgaStableTarget{sga6.e13.eq.4-5-1}\begin{equation}
\langle c_1(f^*L)\cdot Y'\rangle
=[k(Y'):k(Y)]\,\langle c_1(L)\cdot Y\rangle.
\tag{4.5.1}
\end{equation}
If \(\dim(Y)=0\), then \((f^*L)_{Y'}\simeq\mathcal O_{Y'}\). This proves
(i) for numerical equivalence.

Conversely, for a closed integral curve \(Y\subset X\), choose a point
\(y'\in f^{-1}(y)\) algebraic over the function field \(k(y)\) of the
generic point \(y\), and let \(Y'\) be its closure. Then \(f(Y')=Y\), and
\hyperlink{sga6.e13.eq.4-5-1}{4.5.1} proves (ii).

\sgaStableTarget{sga6.e13.thm.4-6}\textbf{Theorem 4.6 (The numerical characterization of
\(\Pic^\tau_{X/k}\)).} \emph{Let \(k\) be algebraically closed, \(X\)
proper over \(k\), and \(L\) invertible on \(X\). The following are
equivalent:}
\begin{enumerate}[label=(\roman*)]
\item \emph{\(L\) is \(\tau\)-equivalent to zero.}
\item \emph{For every coherent \(F\) on \(X\),}

% --- current-rescribe idx659 | printed p. 646 | source-PDF p. 649 ---
\[
\chi(F\otimes L)=\chi(F).
\]
\item \emph{\(L\) is numerically equivalent to zero.}
\end{enumerate}
\emph{If \(X/k\) is projective with an ample
\(H=\mathcal O_X(1)\), these are also equivalent to:}
\begin{enumerate}[label=(\roman*),start=4]
\item \emph{\(L^{\otimes m}(1)\) is ample for every integer \(m\).}
\item \emph{If \(X\) is irreducible, then
\[
\chi(L^{\otimes m}(n))=\chi(\mathcal O_X(n))
\]
for every \(m,n\).}
\item \emph{If \(X\) is irreducible of dimension \(r\ge2\), then
\[
\langle c_1(L)c_1(H)^{r-1}\rangle=0,\qquad
\langle c_1(L)^2c_1(H)^{r-2}\rangle=0.
\]}
\end{enumerate}

For (i) \(\Rightarrow\) (ii), devissage (EGA III 3.1) and \hyperlink{sga6.e13.prop.4-5}{4.5} reduce to
\(X\) integral and \(F=\mathcal O_X\). It remains to show that
\(\chi(L^{\otimes m})\) is constant. Replace \(L\) by a nonzero tensor
power and suppose it algebraically equivalent to zero. If \(M\) on
\(X\times_kT\), with \(T\) connected of finite type, deforms \(L\) to
\(\mathcal O_X\), then \(\chi(M_t^{\otimes m})\) is independent of
\(t\in T\) (EGA III 7.9.5).

For a closed integral curve \(Y\), take \(F=\mathcal O_Y\) in (ii).
The Riemann formula
\[
\chi(L_Y)=\langle c_1(L)\cdot Y\rangle+\chi(\mathcal O_Y)
\]
shows that (ii) implies (iii).

Let \(f:X'\to X\) be a Chow presentation. If (iii) implies (i) for
\(X'\), then \hyperlink{sga6.e13.prop.4-5}{4.5} gives it for \(X\). Hence

% --- current-rescribe idx660 | printed p. 647 | source-PDF p. 650 ---
assume \(X\) projective with an ample \(H\). Passing through the finite
surjection \(\coprod_iX_i\to X\), where the \(X_i\) are the reduced
irreducible components, reduces (iii) \(\Rightarrow\) (i),(iv) to \(X\)
integral (\hyperlink{sga6.e13.prop.4-5}{4.5}, EGA III 2.6.2).

For irreducible \(X\), (ii) implies (v), (iii) implies (vi), and \hyperlink{sga6.e13.para.2-1}{2.1}
gives (v) implies (vi); condition (vi) is unchanged on
\(X_{\mathrm{red}}\). It remains to prove (vi) implies (i),(iv), and
(iv) implies (i), for \(X\) integral. Under (vi), respectively (iv), the
family of classes \(L^{\otimes m}\), \(m\in\mathbb N\), satisfies \hyperlink{sga6.e13.thm.3-13}{3.13
(iii)}, respectively (iv), and is bounded. Hence some invertible
\(\mathcal M\) on \(X_T\), with \(T\) of finite type over \(k\), bounds
it. Choose distinct \(p,q\) such that \(L^{\otimes p}\) and
% SOURCE-CORRECTION / PENDING CLAUDE: print/workpass use exponent r in the
% first power despite having just introduced p,q; the argument requires p.
\(L^{\otimes q}\) correspond to the same connected component of \(T\).
Then \(L^{\otimes(p-q)}\) is algebraically equivalent to zero, proving
(i).

Condition (vi) also gives an \(N\) such that
\(L^{\otimes m}\otimes H^{\otimes N}\) is ample for every \(m\).
Therefore \((L^{\otimes m}\otimes H)^{\otimes N}\) is ample, so
\(L^{\otimes m}\otimes H\) is ample. Thus (vi) implies (iv).

\sgaStableTarget{sga6.e13.thm.4-7}\textbf{Theorem 4.7.} \emph{Let \(S\) be noetherian and \(X\) proper over
\(S\). Then:}
\begin{enumerate}[label=(\roman*)]
\item \emph{\(\underline{\Pic}^{\tau}_{X/S}\to
\underline{\Pic}_{X/S}\) is representable by an open immersion.}
\item \emph{If \(X\) is projective and integral over \(S\), this map}

% --- current-rescribe idx661 | printed p. 648 | source-PDF p. 651 ---
\emph{is also representable by a closed immersion.}
\item \emph{\(\underline{\Pic}^{\tau}_{X/S}\) is of finite type over
\(S\); that is, the family of classes \(\tau\)-equivalent to zero on the
fibers is bounded.}
\end{enumerate}

For (i), respectively (ii), take an arbitrary
\(T\to\underline{\Pic}_{X/S}\) and show that the inverse image of
\(\underline{\Pic}^{\tau}_{X/S}\) is an open, respectively open and
closed, subscheme of \(T\). Locally, \(T=\operatorname{Spec}(B)\) over
\(\operatorname{Spec}(A)\subset S\), and \(T\) is the limit of
\(T_\alpha=\operatorname{Spec}(B_\alpha)\), for the finite-type
\(A\)-subalgebras \(B_\alpha\subset B\). By \hyperlink{sga6.e13.lem.3-1}{3.1} the morphism descends to
some \(T_\alpha\), so assume \(T\) of finite type over \(S\).

The morphism corresponds to an fppf \(R\to T\) and an invertible sheaf
\(L\) on \(X_R\). By fppf descent replace \(T\) by \(R\), then base-change
along \(R\to S\). Compatibility with base change \hyperlink{sga6.e13.para.4-1}{(4.1)} reduces to a
section \(S\to\underline{\Pic}_{X/S}\) defined by \(L\) on \(X\).

Let \(S_0\) be the set of \(s\in S\) such that \(L_s\) is
\(\tau\)-equivalent to zero. It must be open, and under (ii) also closed.
Since \(S\) is noetherian, for \(s\in S_0\) it suffices that
\begin{enumerate}[label=(\alph*)]
\item every generalization of \(s\) lies in \(S_0\);
\item every specialization sufficiently near \(s\) lies in \(S_0\);
\item under (ii), every specialization of \(s\) lies in \(S_0\).
\end{enumerate}

% --- current-rescribe idx662 | printed p. 649 | source-PDF p. 652 ---
Let \(f:X'\to X\) be a Chow presentation. If (a),(b),(c), and (iii) hold
for \(X'\), then \hyperlink{sga6.e13.lem.4-3}{4.3} and \hyperlink{sga6.e13.thm.3-5}{3.5} give them for \(X\); assume \(X\) projective
over \(S\).

Let \(H\) be ample for \(X/S\). If
\((L^{\otimes m}\otimes H)_s\) is ample, then it remains ample at every
generalization of \(s\) (EGA III 4.7.1). The equivalence \hyperlink{sga6.e13.thm.4-6}{4.6
(i)} \(\Longleftrightarrow\) (iv) gives (a).

For (b),(c), and (iii), reduce to \(S\) integral, with \(s\) generic in
(b),(c); for (iii), apply \hyperlink{sga6.e13.lem.3-4}{3.4 (i)} to
\(\coprod_iS_i\to S\), where \(S_i\) are the reduced irreducible
components.

There are an integral \(S'\) and an fppf \(S'\to S\) such that every
reduced irreducible component \(X'_i\) of \(X_{S'}\) is integral over
\(S'\) (\hyperlink{sga6.e12.lem.1-3}{XII 1.3}). It is enough to prove (b),(iii) for \(X_{S'}/S'\), by
descent and by noetherian induction with \hyperlink{sga6.e13.lem.3-4}{3.4 (i)}. The surjection
\(\coprod_iX'_i\to X_{S'}\), together with \hyperlink{sga6.e13.lem.4-3}{4.3}, \hyperlink{sga6.e13.thm.3-5}{3.5}, and \hyperlink{sga6.e13.lem.3-7}{3.7}, reduces
(b),(iii) to (c),(iii) for each \(X'_i/S'\).

Thus we are reduced to proving (c) and (iii), assuming that \(S\) is
integral with generic point \(s\), and that \(X\) is projective and
integral over \(S\). For a very ample \(\mathcal O_X(1)\) for \(X/S\),
the values \(\chi(L_t^{\otimes m}(n))\) and
\(\chi(\mathcal O_{X_t}(n))\), for every \(m,n\), are independent of
\(t\in S\) (EGA III 7.9.4). Hence (c) follows from \hyperlink{sga6.e13.thm.4-6}{4.6
(i)} \(\Longleftrightarrow\) (v), which also says that the family of
classes of invertible sheaves \(\tau\)-equivalent to zero on the fibers
of \(X/S\) has only one

% --- current-rescribe idx663 | printed p. 650 | source-PDF p. 653 ---
% Reconciled to French source-rescribe commit 859bd5abc3a813ee51e7cbcc6d623e08c1336235
% and both source scans; current French-control cursor after this tranche is idx685.
Hilbert polynomial. Thus (iii) follows from 2.11.

\sgaStableTarget{sga6.e13.rem.4-8}\textbf{Remark 4.8.} Let \(X\) be proper and integral over a noetherian
scheme \(S\). Without the projectivity hypothesis on \(X/S\), it is
unknown whether
\(\underline{\Pic}^{\tau}_{X/S}\to\underline{\Pic}_{X/S}\) remains a
\emph{closed} immersion, except when \(X\) is normal over \(S\) and
\(\underline{\Pic}_{X/S}\) is \emph{representable} (FGA 236, 2.3).

\sgaStableTarget{sga6.e13.sec.5}\subsection*{5. Finiteness theorems of ``N\'eron--S\'everi'' type}

\sgaStableTarget{sga6.e13.thm.5-1}\textbf{Theorem 5.1.} \emph{Let \(X\) be proper over noetherian \(S\).
Then the N\'eron--S\'everi groups
\(\underline{\Pic}_{X_K/K}(K)/\underline{\Pic}^0_{X_K/K}(K)\) of the
geometric fibers \(X_K/K\) of \(X/S\) are finitely generated, and their
ranks and the orders of their torsion subgroups
\(\underline{\Pic}^\tau_{X_K/K}(K)/\underline{\Pic}^0_{X_K/K}(K)\)
remain bounded.}

By noetherian induction, replace \(S\) by an arbitrary nonempty open of
\(S_{\mathrm{red}}\), and suppose
\(\underline{\Pic}^{\tau}_{X/S}\) represented by a finite-type
\(S\)-scheme (\hyperlink{sga6.e12.cor.1-2}{XII 1.2},\hyperlink{sga6.e13.thm.4-7}{ 4.7 (iii)}). The order of
\(\underline{\Pic}^\tau_{X_K/K}(K)/\underline{\Pic}^0_{X_K/K}(K)\)
is the number of connected components of
\((\underline{\Pic}^{\tau}_{X/S})_K\), constant over a nonempty open of
\(S\) (EGA IV 9.7.9).

It remains to bound the rank of
\(\underline{\Pic}_{X_K/K}(K)/\underline{\Pic}^\tau_{X_K/K}(K)\).
Reduce as in \hyperlink{sga6.e13.thm.4-7}{4.7 (iii)} to
\(S\) integral and \(X\) projective and integral. Let \(r\) be the common
fiber dimension of \(X/S\). For \(r=0\), respectively \(1\), the rank is \(0\),
respectively \(1\). For \(r\ge2\), reduce to \(r=2\) and \(X\) smooth.
After restricting \(S\), take a linear section \(Y\) of \(X\), still

% --- current-rescribe idx664 | printed p. 651 | source-PDF p. 654 ---
integral over \(S\), with two-dimensional fibers, and replace \(X\) by
\(Y\) using 4.3. Then use \hyperlink{sga6.e13.lem.4-3}{4.3} and \hyperlink{sga6.e13.rem.3-10}{3.10}, together with resolution of
surface singularities, and replace \(S\) by a flat cover of a nonempty
open, to make \(X\) smooth over \(S\).

By \hyperlink{sga6.e13.thm.4-6}{4.6 (i)} \(\Longleftrightarrow\) (iii) and finiteness and specialization
for \'etale cohomology (SGA 4 XVI 2.2), Theorem 5.1 follows from:

\sgaStableTarget{sga6.e13.prop.5-2}\textbf{Proposition 5.2.} \emph{Let \(k\) be algebraically closed and
\(X\) proper smooth over \(k\). The group \(C^p_{\mathrm{num}}(X)\) of
codimension-\(p\) algebraic cycles modulo numerical equivalence is free
of rank at most \(b^{2p}\), the \(2p\)-th Betti number.}

Disregarding the action of roots of unity, \'etale cohomology
\(H^*(X,\mathbb Q_\ell)=_{\mathrm{dfn}}\allowbreak
\left[\varprojlim_\nu H^*(X,\mathbb Z/\ell^\nu\mathbb Z)\right]
\allowbreak\otimes_{\mathbb Z_\ell}\mathbb Q_\ell\), where
\((\ell,\operatorname{char}(k))=1\), is a Weil cohomology: it satisfies
Poincar\'e duality and the K\"unneth formula, and receives a functorial,
nonzero homomorphism from the ring of algebraic cycles modulo rational
equivalence \(C^*_{\mathrm{rat}}(X)\) (see \hyperlink{sga6.e10.sec.7.13}{X 7.13}).

Choose \(a_1,\ldots,a_m\in C^{r-p}_{\mathrm{rat}}(X)\), \(r=\dim X\),
whose images form a basis for the \(\mathbb Q_\ell\)-subspace of
\(H^{2(r-p)}(X,\mathbb Q_\ell)\) spanned by algebraic cycles. Define
\[
\alpha:C^p_{\mathrm{rat}}(X)\longrightarrow\mathbb Z^m.
\]
It is defined by
\(\alpha(b)=(\langle b\cdot a_1\rangle,\ldots,\langle b,a_m\rangle)\),
where \(\langle\ \rangle\) denotes the intersection number. Its image
identifies with \(C^p_{\mathrm{num}}(X)\).

% --- current-rescribe idx665 | printed p. 652 | source-PDF p. 655 ---
\sgaStableTarget{sga6.e13.rem.5-3}\textbf{Remark 5.3.} If \(X\) is a proper smooth surface over
algebraically closed \(k\), then its second Betti number is
\(b^2(X)=c_2(X)-2+2b^1(X)\), where
\(b^1(X)=2\dim\Pic^0_{X/k}\), and \(c_2(X)\), the second Chern
class, is the Euler-Poincar\'e characteristic, equivalently the
self-intersection of the diagonal in \(X\times X\). Thus \(b^2(X)\) has
a cohomology-independent interpretation. In characteristic \(>0\), the
inequality \(\operatorname{rank}C^1_{\mathrm{num}}(X)\le b^2(X)\) in
this form is due, before \'etale cohomology, to Igusa, who used deep
results on the tame fundamental group of an algebraic curve minus some
points.

\sgaStableTarget{sga6.e13.thm.5-4}\textbf{Theorem 5.4.} \emph{Let \(k\) be algebraically closed and \(X\)
proper over \(k\). There are finitely many closed integral curves
\(Y_1,\ldots,Y_p\) in \(X\) such that a subset \(\Lambda\) of
\(\Pic_{X/k}\) is quasi-compact if and only if the intersection numbers
\(\langle\lambda\cdot Y_i\rangle\), for \(\lambda\in\Lambda\), remain
bounded. Moreover, \(p\) may be taken to be the rank of the
N\'eron--S\'everi group \(\Pic_{X/k}(k)/\Pic^0_{X/k}(k)\).}

Indeed, \(\Pic^\tau_{X/k}\) is an open quasi-compact subgroup of
\(\Pic_{X/k}\) \hyperlink{sga6.e13.thm.4-7}{(4.7)}. Thus \(\Lambda\) is quasi-compact if and only if
the corresponding set of points in
\(\bigl[\Pic_{X/k}(k)/\Pic^\tau_{X/k}(k)\bigr]
\otimes_{\mathbb Z}\mathbb Q\) is finite. By \hyperlink{sga6.e13.thm.4-6}{4.6 (i)}
\(\Longleftrightarrow\) (iii),
\(\tau\)-equivalence equals numerical equivalence. Choose curves whose
images form a basis of the dual vector space; the assertion follows.

% End idx665 / printed p. 652 / source-PDF p. 655. Theorem 5.4 is complete.

% END INCLUDED SGA6_English_current_buildable_source_20260823/SGA6_idx647_665_English_SourceChecked_fragment.tex


% Reconciled French-backed prefix through idx684, followed by the two sealed
% July 22 scan-backed English terminal successors and the pre-existing index
% back matter.  The terminal source-note blocks are moved after the complete
% Expose XIV body so that the idx692/idx693 sentence remains continuous.
% BEGIN INCLUDED SGA6_English_current_buildable_source_20260823/fragments/SGA6_idx666_684_English_Reconciled_Prefix.tex
% SGA 6 English integration fragment: current-rescribe idx666--702 plus
% the unindexed volume back matter on source-PDF pages 693--702.
% English idx666--684 is reconciled to French source-rescribe commit
% 859bd5abc3a813ee51e7cbcc6d623e08c1336235 (release-TeX SHA-256
% 7F32C2080A78A2746CBE52DCC1EC43A8505269F25518FA7B9A86E4E89AF858AC)
% and the source scans. Idx685--702 and terminal matter remain scan-checked
% working text pending later French-control reconciliation.
% This file contains no document wrapper and is suitable for \input into the
% cumulative SGA 6 English document when its required packages/macros are loaded.

% [SGA6 current-rescribe idx666 | printed 653 | source-PDF 656 | high-res page 667]
\sgaStableTarget{sga6.e13.sec.6}\subsection*{6. Appendix: Study of $(b)$-sheaves on $P=\mathbf P^N_k$.}

\sgaStableTarget{sga6.e13.para.6-1}\textbf{6.1.} Let $k$ be an algebraically closed field and let $P=\mathbf P^N_k$. Let $F$ be a coherent sheaf on $P$, with Hilbert polynomial
\[
  \chi(F(n))=\sum_{i=0}^{r}a_i\binom{n+i}{i},
\]
and let
\[
  (c)=(c_0,\ldots,c_r)
\]
be a sequence of integers with $(c)\geq(a)$. Put
\[
  a_{m,i}=\sum_{j=0}^{r-i}a_{j+i}\binom{m-1+j}{j}
  \quad\text{and}\quad
  c_{m,i}=\sum_{j=0}^{r-i}c_{j+i}\binom{m-1+j}{j}.
\]
Finally, for every integer $q\geq1$, denote by $N_q$ the largest coherent subsheaf of $F$ whose support has dimension at most $q-1$, and put $F_q=F/N_q$.

\sgaStableTarget{sga6.e13.lem.6-2}\textbf{Lemma 6.2.} \emph{Under the conditions of \hyperlink{sga6.e13.para.6-1}{6.1}, consider an exact sequence
\[
  0\longrightarrow I\longrightarrow\cO_P^{\oplus M}
  \longrightarrow F(m)\longrightarrow0
\]
with $m\geq0$. Then $I$ is $p$-regular, where
\(p=P_r(c_{m,0},\ldots,c_{m,r})\), and $P_r$ is the $r$-th
$(b)$-polynomial of 1.9.}

Indeed, $\cO_P$ is plainly a $(0,\ldots,0,1)$-sheaf ($N+1$ terms), so
$I$ is a $(0,\ldots,0,M)$-sheaf by 1.6. By \hyperlink{sga6.e13.rem.2-10}{2.10},
\(\chi(I(n))=M\binom{n+N}{N}
-\sum_{i=0}^{r}a_{m,i}\binom{n+i}{i}\). The assertion follows from
\hyperlink{sga6.e13.thm.1-11}{1.11} and 1.10.

\sgaStableTarget{sga6.e13.prop.6-3}\textbf{Proposition 6.3.} \emph{Under the conditions of \hyperlink{sga6.e13.para.6-1}{6.1}, suppose that $F$ is a $(b)$-sheaf. Then, by means of universal polynomials in the $a_i$, the $b_i$, and $N$, one can calculate integers $m,M,m_1,M_1$ for which there is an exact sequence
\[
  \cO_P(-m_1)^{\oplus M_1}\longrightarrow
  \cO_P(-m)^{\oplus M}\longrightarrow F\longrightarrow0.
\]
Moreover, $N$ enters only into the calculation of $M_1$.}

% [SGA6 current-rescribe idx667 | printed 654 | source-PDF 657 | high-res page 668]
Indeed, take $m=P_{r-1}(c_0,\ldots,c_{r-1})$, with $c_i=b_i-a_i$,
and $M=\sum_i a_i\binom{m+i}{i}$, so that, by \hyperlink{sga6.e13.thm.1-11}{1.11} and \hyperlink{sga6.e13.prop.1-3}{1.3}, there is
an exact sequence
\[
  0\longrightarrow I\longrightarrow\cO_P^{\oplus M}
  \longrightarrow F(m)\longrightarrow0.
\]
Likewise, taking account of \hyperlink{sga6.e13.lem.6-2}{6.2}, take $m_1=m+p$, where
$p=P_r(a_{m,0},\ldots,a_{m,r})$, and
$M_1=M\binom{p+N}{N}-\sum_i a_i\binom{m_1+i}{i}
\quad(=\chi(I(p)))$.

\sgaStableTarget{sga6.e13.thm.6-4}\textbf{Theorem 6.4.} \emph{Under the conditions of \hyperlink{sga6.e13.para.6-1}{6.1}, suppose that
$F$ is a quotient of $\cO_P(-m)^{\oplus M}$, with $m\geq0$. Put
$b_i=P_{r-i}(c_{m,i},\ldots,c_{m,r})$ for $i=0,\ldots,r$, where $P_s$
is the $s$-th $(b)$-polynomial of \hyperlink{sga6.e13.def.1-9}{1.9}; put $b=b_1+m-1$, and put
$B=\sum_{i=1}^{r}a_i\binom{b+i}{i}$. Then:}
\begin{enumerate}[label=(\roman*)]
\item \emph{$F$ is $b$-regular.}
\item \emph{$-B\leq a_0=h^0(F(b))-B$}
\item \emph{$F$ is a $(b_0,\ldots,b_r)$-sheaf.}
\end{enumerate}

Indeed, choose a section $\sigma$ of $\cO_P(1)$ for which one obtains the following commutative diagram with exact rows and columns:
\[
\begin{tikzcd}[column sep=large,row sep=large]
&0 \arrow[d]&0 \arrow[d]&0 \arrow[d]&\\
0 \arrow[r]&I(-1) \arrow[r] \arrow[d]&I \arrow[r] \arrow[d]&J \arrow[r] \arrow[d]&0\\
0 \arrow[r]&\cO_P(-1)^{\oplus M} \arrow[r] \arrow[d]&\cO_P^{\oplus M} \arrow[r] \arrow[d]&\cO_Y^{\oplus M} \arrow[r] \arrow[d]&0\\
0 \arrow[r]&F(m-1) \arrow[r] \arrow[d]&F(m) \arrow[r] \arrow[d]&G(m) \arrow[r] \arrow[d]&0\\
&0&0&0&
\end{tikzcd}
\]

% [SGA6 current-rescribe idx668 | printed 655 | source-PDF 658 | high-res page 669]
Since $\chi(G(n))=\sum_{i=0}^{r-1}a_{i+1}\binom{n+i}{i}$ by \hyperlink{sga6.e13.lem.1-7}{1.7},
$J$ is $b_1$-regular by \hyperlink{sga6.e13.lem.6-2}{6.2}; hence, by \hyperlink{sga6.e13.prop.1-4}{1.4(i)}, $H^q(I(p))=0$ for
$q\geq2$ and $p\geq b_1-q$. The exact sequence
$H^q\bigl(\cO_P^{\oplus M}(p)\bigr)\longrightarrow
H^q(F(m+p))\longrightarrow H^{q+1}(I(p))$ then shows that $F(m)$ is
$(b_1-1)$-regular; this proves (i).

Assertion (i) implies
$0\leq h^0(F(b))=\chi(F(b))
=\sum_{i=0}^{r}a_i\binom{b+i}{i}$, which proves (ii). By \hyperlink{sga6.e13.rem.2-10}{2.10} and \hyperlink{sga6.e13.def.1-9}{1.9},
\[
  \chi(F(b))
  =\sum_{i=0}^{r}a_{m,i}\binom{b_1-1+i}{i}
  \leq b_1+\sum_{i=0}^{r}c_{m,i}\binom{b_1-1+i}{i}
  =b_0.
\]
Thus $h^0(F(-1))\leq h^0(F(b))\leq b_0$. By induction on $r$, one may suppose that $G$ is a $(b_1,\ldots,b_r)$-sheaf, so (iii) follows from 1.6.2.

\sgaStableTarget{sga6.e13.lem.6-5}\textbf{Lemma 6.5.} \emph{Under the conditions of \hyperlink{sga6.e13.para.6-1}{6.1}, suppose that $F$
has no associated prime cycle of dimension $0$.}
\begin{enumerate}[label=(\roman*)]
\item \emph{If $h^0(F)\geq1$, then $h^0(F(-1))\leq h^0(F)-1$.}
\item \emph{$H^0(F(-n))=0$ for $n\geq h^0(F)$.}
\item \emph{Let $0\to F(-1)\to F\to G\to0$ be exact, and let
$n_0\geq0$ satisfy $H^0(G(-n_0))=0$. Then $H^0(F(-n_0))=0$ and
$h^0(F)\leq n_0\cdot h^0(G)$.}
\end{enumerate}

For (i), let $0\ne\sigma\in H^0(F)$. The section $\sigma$ defines a nonzero subsheaf $F'=\cO_X\cdot\sigma$ of $F$, so $s=\dim\Supp(F')\geq1$. Consider the commutative diagram with exact rows and columns
\[
\begin{tikzcd}[column sep=large,row sep=large]
&0 \arrow[d]&0 \arrow[d]&0 \arrow[d]&\\
0 \arrow[r]&F'(-1) \arrow[r] \arrow[d]&F' \arrow[r] \arrow[d]&G' \arrow[r] \arrow[d]&0\\
0 \arrow[r]&F(-1) \arrow[r]&F \arrow[r]&G \arrow[r]&0.
\end{tikzcd}
\]

% [SGA6 current-rescribe idx669 | printed 656 | source-PDF 659 | high-res page 670]
Since $s\geq1$, one has $G'\ne0$; hence $F'$ is not contained in $F(-1)$. In other words, the section $\sigma$ is not in the image of the injection $H^0(F(-1))\to H^0(F)$, which proves (i).

Assertion (ii) follows immediately from (i), bearing in mind that
$h^0(F)=0$ plainly implies $h^0(F(-1))=0$, since
$h^0(F(-1))\leq h^0(F)$. In (iii) one obtains
$0\leq h^0(F(n))-h^0(F(n-1))\leq h^0(G(n))$.
Thus, for $n\geq n_0$, $h^0(F(-n_0))=h^0(F(-n))$, since $h^0(G(-n_0))=0$; hence $h^0(F(-n_0))=0$ by (ii). Consequently, for $p\geq1$,
$h^0(F(p-n_0))
\leq h^0(G(1-n_0))+\cdots+h^0(G(p-n_0))
\leq p\cdot h^0(G(p-n_0))$,
which gives (iii).

\sgaStableTarget{sga6.e13.prop.6-6}\textbf{Proposition 6.6.} \emph{Under the conditions of \hyperlink{sga6.e13.para.6-1}{6.1}, suppose
that $F$ is a $(b_0,\ldots,b_r)$-sheaf. Then $F_q$ is a
$\bigl((b_{q-1})^q,\ldots,(b_{q-1})^2,b_{q-1},b_q,\ldots,b_r\bigr)$-sheaf.}

Indeed, let $\sigma$ be a ``general'' section of $\cO_X(1)$. It defines a commutative diagram with exact rows and columns:
\[
\begin{tikzcd}[column sep=large,row sep=large]
&0 \arrow[d]&0 \arrow[d]&0 \arrow[d]&\\
0 \arrow[r]&N_q(-1) \arrow[r] \arrow[d]&N_q \arrow[r] \arrow[d]&G' \arrow[r] \arrow[d]&0\\
0 \arrow[r]&F(-1) \arrow[r] \arrow[d]&F \arrow[r] \arrow[d]&G \arrow[r] \arrow[d]&0\\
0 \arrow[r]&F_q(-1) \arrow[r] \arrow[d]&F_q \arrow[r] \arrow[d]&G'' \arrow[r] \arrow[d]&0\\
&0&0&0&
\end{tikzcd}
\]
where $G$ is a $(b_1,\ldots,b_r)$-sheaf by \hyperlink{sga6.e13.prop.1-6}{1.6(i)}, and where
$G''=G_{q-1}$ (a depth question).

% [SGA6 current-rescribe idx670 | printed 657 | source-PDF 660 | high-res page 671]
Suppose first that $q=1$. Then $\dim\Supp(N_q)=0$, so
$H^1(N_q(-1))=0$ and $h^0(F_q(-1))\leq h^0(F(-1))\leq b_0$.
Moreover, $G'=0$, and hence $G\simeq G''$. Thus $F_1$ is indeed a $(b_0,\ldots,b_r)$-sheaf. In addition,\hyperlink{sga6.e13.lem.6-5}{ 6.5(ii)} shows that
$H^0(F_1(-b_0))=0$. If $q\geq2$, argue by induction on $q$. One may
suppose that $G_{q-1}$ is a
$\bigl((b_{q-1})^{q-1},\ldots,b_{q-1},\ldots,b_r\bigr)$-sheaf and that
$H^0(G_{q-1}(-b_{q-1}))=0$. \hyperlink{sga6.e13.lem.6-5}{Lemma 6.5} then gives
$H^0(F_q(-b_{q-1}))=0$ and $h^0(F_q(-1))\leq(b_{q-1})^q$,
which proves the assertion.

\sgaStableTarget{sga6.e13.thm.6-7}\textbf{Theorem 6.7.} \emph{There are two sequences
$\{A_i(X_0,\ldots,X_i;Y)\}$ and
$\{A_i^{(q)}(X_0,\ldots,X_q;Y)\}$ of polynomials having the following
properties. Assume the conditions of \hyperlink{sga6.e13.para.6-1}{6.1} and suppose that $F$ is a
quotient of $\cO_P(-m)^{\oplus M}$, with $m\geq0$.}
\begin{enumerate}[label=(\roman*)]
\item \emph{If $F$ is a $(b_0,\ldots,b_r)$-sheaf, then, for
$i=0,\ldots,r$,}
\[
  |a_i|\leq A_{r-i}(b_i,\ldots,b_r;m).
\]
\item \emph{If
$\chi(F_q(n))=\sum_{i=0}^{r}a_i^{(q)}\binom{n+i}{i}$ is the Hilbert
polynomial of $F_q$, then, for $i=0,\ldots,q-1$,}
\[
  |a_i^{(q)}|\leq A_{r-i}^{(r-q)}(c_{q-1},\ldots,c_r;m).
\]
\end{enumerate}
\emph{(Notice that $a_{q-1}^{(q)}\leq a_{q-1}$ and
$a_q^{(q)}=a_q,\ldots,a_r^{(q)}=a_r$, since
$\dim\Supp(N_q)\leq q-1$.)}

To prove (i), use induction on $r$ and apply \hyperlink{sga6.e13.thm.6-4}{6.4} to an exact sequence
$0\longrightarrow F(-1)\longrightarrow F\longrightarrow G\longrightarrow0$
in which $G$ is a $(b_1,\ldots,b_r)$-sheaf. By induction on $r$, one may
suppose that $A_0,\ldots,A_{r-1}$ have been defined, taking \hyperlink{sga6.e13.lem.1-7}{1.7} into
account. Then \hyperlink{sga6.e13.thm.6-4}{6.4(ii)}, together with \hyperlink{sga6.e13.prop.1-8}{1.8}, displays the definition of
$A_r$. Finally, (ii) follows formally from \hyperlink{sga6.e13.thm.6-4}{6.4(iii)}, \hyperlink{sga6.e13.prop.6-6}{6.6}, and assertion
(i), applied to $F_q$ in place of $F$.

% [SGA6 current-rescribe idx671 | printed 658 | source-PDF 661 | high-res page 672]
\sgaStableTarget{sga6.e13.cor.6-8}\textbf{Corollary 6.8. (Grothendieck).} \emph{Let $X$ be projective over
a noetherian scheme $S$, let $\cO_X(1)$ be very ample for $X/S$, and let
$\mathcal J$ be a family of classes of coherent sheaves on the fibers of
$X/S$. Suppose that:}
\begin{enumerate}[label=(\alph*)]
\item \emph{There is a coherent sheaf $E$ on $X$ such that $\mathcal J$
is contained in the family of classes of quotients of sheaves of the
form $E_K$.}
\item[$(b_q)$] \emph{In the Hilbert polynomials $\chi(F_K(n))$, for
$F_K\in\mathcal J$, the coefficients in degrees at least $q-1$ remain
bounded above.}
\end{enumerate}
\emph{Then the $(F_K)_q$ (as $F_K$ ranges over $\mathcal J$) form a
bounded family; moreover, the coefficients of $\chi(F_K(n))$ in degrees
at least $q-2$ remain bounded below.}

Indeed, one may plainly suppose that $X=\mathbf P^N_S$ and
$E=\cO_X(-m)^{\oplus M}$. The first assertion follows formally from\hyperlink{sga6.e13.thm.6-7}{
6.7(ii)} and 1.13. The second follows by induction, using an exact sequence
$0\longrightarrow F(-1)\longrightarrow F\longrightarrow G\longrightarrow0$
and \hyperlink{sga6.e13.thm.6-4}{6.4(ii)}.

\sgaStableTarget{sga6.e13.def.6-9}\textbf{Definition 6.9.} \emph{Let $k$ be a field. We call a
\emph{special positive $k$-cycle of dimension $r$} any projective
$k$-scheme $X$, equipped with a very ample sheaf $\cO_X(1)$, which is a
union of closed subschemes $X_j$ of dimension $r$, each $X_j$ being
obtained by base change $\Spec(k)\to\Spec(k_j)$ as the inverse image of
an integral $k_j$-scheme $X'_j$ (equipped with $\cO_{X'_j}(1)$). We call
the \emph{degree} of $X$ the coefficient $a_r$ in the Hilbert polynomial
$\chi(\cO_X(n))=\sum_i a_i\binom{n+i}{i}$.}

\sgaStableTarget{sga6.e13.lem.6-10}\textbf{Lemma 6.10.} \emph{Let $k$ be an algebraically closed field, and let $X$ be a special positive $k$-cycle of dimension $r$ and degree $d$. Then $\cO_X$ is a $(0,\ldots,0,d)$-sheaf.}

% [SGA6 current-rescribe idx672 | printed 659 | source-PDF 662 | high-res page 673]
Indeed, view $X$ as a subscheme of $\mathbf P^N_k$. After replacing $k$ by the algebraic closure of a purely transcendental extension, which is permissible by \hyperlink{sga6.e13.prop.1-6}{1.6(i)}, one may suppose that $k$ contains $r(N+1)$ elements algebraically independent over all the $k_j$. The $r$ successive ``general'' hyperplane sections of $X$ defined by these elements are then always special positive $k$-cycles of degree $d$.

Arguing by induction on $r$, it remains only to observe that, when $r=0$,
one has $h^0(\cO_X(-1))=d$, while, when $r\geq1$, one has
$h^0(\cO_X(-1))=0$. For the latter point, the injection
$\cO_X\longrightarrow\coprod_j\cO_{X_j}$ allows one to suppose that $X$
itself comes from an integral scheme $X'$. But then
$h^0(\cO_X(-1))=h^0(\cO_{X'}(-1))=0$, because $H^0(\cO_{X'})$ is a
field.

\sgaStableTarget{sga6.e13.cor.6-11}\textbf{Corollary 6.11.} (i) \emph{Let $r,d$ be two integers, put
$c_i=A_{r-i}(0,\ldots,0,d;0)$ for $i=0,\ldots,r$ \hyperlink{sga6.e13.thm.6-7}{(6.7(i))}, and put
$p=P_r(c_0,\ldots,c_r)$ \hyperlink{sga6.e13.def.1-9}{(1.9)}. Let $k$ be an algebraically closed field,
and let $X$ be a special positive $k$-cycle of dimension at most $r$ and
degree at most $d$, with Hilbert polynomial
$\chi(\cO_X(n))=\sum_{i=0}^{r}a_i\binom{n+i}{i}$. Then, for
$i=0,\ldots,r$,}
\[
  |a_i|\leq c_i.
\]
\emph{Moreover, $X$ embeds as a subscheme of $\mathbf P^N_k$, with
$N=d(r+1)-1$, defined by (at most) $\binom{N+p}{N}$ equations of degree
$p$.}

(ii) (Chow) \emph{Let $S$ be a noetherian scheme. As $K$ ranges over
the algebraically closed extensions of the $k(s)$, for $s\in S$, the
special positive $K$-cycles of bounded dimension and degree form a
``bounded family'' (either with the $X$ regarded as subschemes of a fixed
$\mathbf P^N_S$, or in an evident abstract sense).}

% [SGA6 current-rescribe idx673 | printed 660 | source-PDF 663 | high-res page 674]
Indeed, in (i), the first assertion follows from \hyperlink{sga6.e13.thm.6-7}{6.7(i)}; the second
follows from \hyperlink{sga6.e13.lem.6-10}{6.10} and \hyperlink{sga6.e13.prop.1-8}{1.8(i)}, and from \hyperlink{sga6.e13.lem.6-2}{6.2}, applied to
$0\longrightarrow I\longrightarrow\cO_P\longrightarrow\cO_X\longrightarrow0$,
$P=\mathbf P^N_k$, together with 1.3. Finally, (ii) follows from (i) and
1.13.

\sgaStableTarget{sga6.e13.rem.6-12}\textbf{Remark 6.12.} Let $k$ be an algebraically closed field and let
$X$ be a projective $k$-scheme, purely of dimension $r$, with no embedded
associated prime cycle. If arbitrary nilpotent elements are allowed in
$\cO_X$, the coefficients in the Hilbert polynomial
$\chi(\cO_X(n))=\sum_{i=0}^{r}a_i\binom{n+i}{i}$ can no longer be bounded
in terms of $a_r$ alone.

For example, let $Z$ be a smooth curve of degree $d$ in $\mathbf P^3_k$.
For every $n\gg0$, there is a smooth surface $Y$ of degree $n$ containing
$Z$. Then $Z$ is a Cartier divisor on $Y$; put
$\cO_{X_n}=\cO_Y/\cO_Y(-2Z)$. There is an exact sequence
% SGA-SRC-001437 — reconstructed by mathematical inference; scan inconclusive
$0\longrightarrow\cO_Y(-Z)\longrightarrow
\cO_Y\longrightarrow\cO_Z\longrightarrow0$ and the formula
$\langle Z^2\rangle_Y=2p_a(Z)-2-(n-4)d$. It follows that
$a_0=\chi(\cO_{X_n}(-1))\longrightarrow\infty$ as $n\longrightarrow
\infty$, although $a_1=2d$.

\sgaStableTarget{sga6.e13.cor.6-13}\textbf{Corollary 6.13.} \emph{Under the conditions of \hyperlink{sga6.e13.para.6-1}{6.1}, consider an exact sequence
\[
  \cO_P(-m_1)^{\oplus M_1}\longrightarrow
  \cO_P(-m)^{\oplus M}\xrightarrow{\alpha}F\longrightarrow0
\]
with $m_1\geq m\geq0$. Then, for $i=0,\ldots,N$,
\[
% SGA-SRC-001438
  |a_i|\leq A_{N-i}(0,\ldots,0,M;m_1)
  +M\binom{m}{N-i}.
\]}

Indeed, let $I=\Ker(\alpha)$. Plainly, $I$ is a
$(0,\ldots,0,M)$-sheaf, and
$\chi(I(n))=\allowbreak\sum_{i=0}^{N}
\bigl[M\tbinom{-m-1+N-i}{N-i}-a_i\bigr]\allowbreak\tbinom{n+i}{i}$.
The assertion follows from \hyperlink{sga6.e13.thm.6-7}{6.7(i)} and the formula
$\tbinom{-m-1+Q}{Q}=(-1)^Q\tbinom{m}{Q}$.

% [SGA6 current-rescribe idx674 | printed 661 | source-PDF 664 | high-res page 675]
\sgaStableTarget{sga6.e13.cor.6-14}\textbf{Corollary 6.14. (I. Hermann).} \emph{For every $N\geq0$, there
is a polynomial $R_N(m)$ such that, for every field $k_0$ and every ideal
$I$ of the polynomial ring $k_0[T_1,\ldots,T_N]$ generated by elements
of degree at most $m$, the radical $\sqrt I$ is generated by elements of
degree at most $R_N(m)$.}

Indeed, let $k$ be the algebraic closure of $k_0$. After introducing an
auxiliary indeterminate, let $Y$ (respectively $X$) be the closed
subscheme of $P=\mathbf P^N_k$ defined by the homogenized ideal of $I$
(respectively of $\sqrt I$). In the Hilbert polynomial
$\chi(\cO_Y(n))=\sum_{i=0}^{N}b_i\binom{n+i}{i}$, one has
$|b_i|\leq c_i$, where $c_i=A_{N-i}(0,\ldots,0,1;m)$, by 6.13.

Let $X=\bigcup X^q$ be the decomposition by dimension. Then
$\deg(X^q)\leq e_q$, where $e_q=P_{N-q}(c_q,\ldots,c_N)$.
Indeed, after cutting $X$ (respectively $Y$) by a ``general'' linear subspace of codimension $q$, one may suppose that $q=0$. Then
$\deg(X^q)=h^0(\cO_{X^q})$ and
$h^0(\cO_{X^q})\leq h^0(\cO_Y)$, and \hyperlink{sga6.e13.thm.6-4}{6.4(iii)} gives
$h^0(\cO_Y)\leq e_q$.

Consequently, $\cO_{X^q}$ is a $(0,\ldots,0,e_q)$-sheaf by 6.10. It
follows from the injection $\cO_X\longrightarrow\prod_q\cO_{X^q}$ that
$\cO_X$ is an $(e_0,\ldots,e_N)$-sheaf. Hence, in the Hilbert polynomial
% SOURCE-EMENDATION: the print has upper bound n; N is required by the
% ambient projective dimension and is retained deliberately.
$\chi(\cO_X(n))=\sum_{i=0}^{N}a_i\binom{n+i}{i}$, one has
$|a_i|\leq f_i$, where $f_i=A_{N-i}(e_i,\ldots,e_N;0)$, by 6.7.
Finally, the ideal $\sqrt I$ of $X$ is $R_N(m)$-regular, where
$R_N(m)=P_N(f_0,\ldots,f_N)$, and the assertion follows from 1.3.

% [SGA6 current-rescribe idx675 | printed 662 | source-PDF 665 | high-res page 676]
\sgaStableTarget{sga6.e13.sec.7}\subsection*{7. Appendix: The Hodge index theorem}

Here we recover and complete some earlier results by a direct argument.

\sgaStableTarget{sga6.e13.thm.7-1}\textbf{Theorem 7.1.} \emph{Let $X$ be an irreducible surface, proper
over an algebraically closed field. Suppose that there is an invertible
sheaf $H$ on $X$ such that $\langle c_1(H)^2\rangle>0$ (for example,
$H$ ample). Then an invertible sheaf $L$ on $X$ satisfying
\[
  \langle c_1(L)\cdot c_1(H)\rangle=0,
  \qquad
  \langle c_1(L)^2\rangle\geq0,
\]
is numerically equivalent to $0$. Consequently, by \hyperlink{sga6.e13.thm.4-6}{4.6}, the intersection
form on $A^1_{\mathrm{num}}(X)_{\Q}
=\bigl[\Pic(X)/\Pic^\tau(X)\bigr]\otimes\Q$ is nondegenerate of type
$(1,\rho-1)$.}

\sgaStableTarget{sga6.e13.para.7-1-1}\textbf{7.1.1.} First suppose that $X$ is reduced (hence integral), that
$H$ is ample, that $\langle c_1(L)\cdot c_1(H)\rangle=0$, and that
$\langle c_1(L)^2\rangle>0$. Arguing by contradiction, choose $m>0$
such that $H_1=L\otimes H^{\otimes m}$ is ample, and choose $n$ such
that, on putting $N=L^{\otimes n}\otimes H^{-1}$, one has
$\langle c_1(N)\cdot c_1(H_1)\rangle
=n\langle c_1(L)^2\rangle-m\langle c_1(H)^2\rangle>0$. One also has
$\langle c_1(N)\cdot c_1(H)\rangle=-\langle c_1(H)^2\rangle<0$ and
$\langle c_1(N)^2\rangle
=n^2\langle c_1(L)^2\rangle+\langle c_1(H)^2\rangle>0$, contradicting
the following lemma.

\sgaStableTarget{sga6.e13.lem.7-1-2}\textbf{Lemma 7.1.2.} \emph{Let $X$ be an integral surface, projective over an algebraically closed field. For an invertible sheaf $N$ on $X$ such that $\langle c_1(N)^2\rangle>0$, the following conditions are equivalent:
\begin{enumerate}[label=(\roman*)]
\item for every ample sheaf $H$, one has $\langle c_1(N)\cdot c_1(H)\rangle>0$;
\item[(i$'$)] there is an ample sheaf $H$ such that $\langle c_1(N)\cdot c_1(H)\rangle>0$;
\item there is an integer $n>0$ such that $H^0(X,N^{\otimes n})\ne0$.
\end{enumerate}}

% [SGA6 current-rescribe idx676 | printed 663 | source-PDF 666 | high-res page 677]
Indeed, (ii) implies that there is an effective divisor $D$ such that
$N^{\otimes n}\simeq\cO_X(D)$, and $\langle c_1(N)^2\rangle>0$ implies
$D>0$; hence (i).

To see that (i$'$) implies (ii), one may plainly suppose that
$H=\cO_X(1)$ is very ample. There is then a closed integral locally
principal curve $Y$ such that $H=\cO_X(Y)$. By the Riemann theorem,
(i$'$) implies $H^1\bigl(N_Y^{\otimes n}(p)\bigr)=0$ for $p\geq0$ and
$n\geq n_0$, where $n_0=2h^1(\cO_Y)-1$.
Since $H^2(N^{\otimes n}(p))=0$ for $p\gg0$, the exact sequence
\[
  H^1\bigl(N_Y^{\otimes n}(p)\bigr)\longrightarrow
  H^2\bigl(N^{\otimes n}(p-1)\bigr)\longrightarrow
  H^2\bigl(N^{\otimes n}(p)\bigr)\longrightarrow0
\]
gives $H^2(N^{\otimes n})=0$ for $n\geq n_0$. Thus, for $n\gg0$,
\[
  h^0(N^{\otimes n})
  =\chi(N^{\otimes n})+h^1(N^{\otimes n})
  \geq\chi(N^{\otimes n})
  =\frac12\langle c_1(N)^2\rangle n^2+\cdots>0,
\]
which completes the proof of the lemma.

\sgaStableTarget{sga6.e13.para.7-1-3}\textbf{7.1.3.} Now suppose in \hyperlink{sga6.e13.thm.7-1}{7.1} that there is an invertible sheaf $M$
on $X$ such that $\langle c_1(L)\cdot c_1(M)\rangle\ne0$, and argue by
contradiction. Replace $X$ by the underlying reduced scheme of a Chow
presentation of $X$, which is permissible by the projection formula \hyperlink{sga6.e10.eq.4-4-4}{X
4.4.4}, and let $G$ be ample on $X$. Choose $m>0$ such that
$H_1=G^{\otimes m}\otimes M$ is ample and
$\langle c_1(L)\cdot c_1(H_1)\rangle\ne0$. Choose $p,q$, with $q\ne0$,
such that, on putting $L_1=L^{\otimes p}\otimes H^{\otimes q}$, one has
$\langle c_1(L_1)\cdot c_1(H_1)\rangle
=p\langle c_1(L)\cdot c_1(H_1)\rangle
+q\langle c_1(H)\cdot c_1(H_1)\rangle=0$. Then one also has
$\langle c_1(L_1)^2\rangle=p^2\langle c_1(L)^2\rangle
+q^2\langle c_1(H)^2\rangle>0$. This is the case \hyperlink{sga6.e13.para.7-1-1}{7.1.1}, with $L=L_1$
and $H=H_1$.

\sgaStableTarget{sga6.e13.para.7-1-4}\textbf{7.1.4.} Finally, consider the general case of 7.1. Still arguing by contradiction, suppose that $L$ is not numerically equivalent to $0$. Replace $X$ by the normalization of $X_{\red}$, so that $X$ has

% [SGA6 current-rescribe idx677 | printed 664 | source-PDF 667 | high-res page 678]
only finitely many singular points; this reduction is justified by the projection formula \hyperlink{sga6.e10.eq.4-4-4}{X 4.4.4} and by \hyperlink{sga6.e13.prop.4-5}{4.5(ii)}. There is therefore a closed integral curve $Y$ on $X$ such that
$\langle c_1(L)\cdot Y\rangle\ne0$, and $Y$ is locally principal
outside a finite set $Z$. Let $f:X'\to X$ be the blow-up of the ideal
% SOURCE-EMENDATION: the print says I is an ideal of O_Y; O_X is required
% for the blow-up of X and is retained deliberately.
$I\subset\cO_X$ defining $Y$. The ideal $\mathcal M=I\cO_{X'}$ is
invertible and defines a subscheme of $X'$ which is the union of two
subschemes $Y'$ and $Y''$ such that $f(Y')=Y$ and $f(Y'')\subset Z$.
By the projection formula \hyperlink{sga6.e10.eq.4-4-4}{X 4.4.4},
$\langle c_1(f^*L)\cdot c_1(\mathcal M)\rangle
=-\langle c_1(f^*L)\cdot Y'\rangle
-\langle c_1(f^*L)\cdot Y''\rangle
=-\langle c_1(L)\cdot Y\rangle-0\ne0$. Thus one may replace $X,L,H$
by $X',f^*L,f^*H$. This reduces the problem to the preceding case and
completes the proof of 7.1.

\sgaStableTarget{sga6.e13.rem.7-2}\textbf{Remark 7.2.} It is not known whether an irreducible proper surface over an algebraically closed field which possesses an invertible sheaf $H$ satisfying $\langle c_1(H)^2\rangle>0$ is necessarily projective. Nor is it known whether, for every irreducible proper surface, the intersection form on $A^1_{\mathrm{num}}(X)_{\Q}$ is nondegenerate. It follows from \hyperlink{sga6.e13.thm.7-1}{7.1} that, when it is degenerate, it is necessarily negative (non-definite).

% SGA-SRC-001442 — reconstructed by mathematical inference; scan inconclusive
\sgaStableTarget{sga6.e13.cor.7-3}\textbf{Corollary 7.3.} \emph{Under the conditions of \hyperlink{sga6.e13.thm.7-1}{7.1}, let $M$ be
an invertible sheaf on $X$ which is not numerically proportional to $H$,
and put $d=\langle c_1(M)\cdot c_1(H)\rangle$ and
$h=\langle c_1(H)^2\rangle$. Then
$\langle c_1(M)^2\rangle<d^2/h$.}

% SGA-SRC-001441
Indeed, let $L=M^{\otimes h}\otimes H^{\otimes(-d)}$. Since
$\langle c_1(L)\cdot c_1(H)\rangle=0$, \hyperlink{sga6.e13.thm.7-1}{Theorem 7.1} gives
$\langle c_1(L)^2\rangle<0$, which proves the assertion.

% [SGA6 current-rescribe idx678 | printed 665 | source-PDF 668 | high-res page 679]
\sgaStableTarget{sga6.e13.cor.7-4}\textbf{Corollary 7.4.} \emph{Let $X$ be an irreducible scheme of
dimension $r\geq2$, projective over an algebraically closed field, and
let $H$ be an ample sheaf on $X$. Then:}
\begin{enumerate}[label=(\roman*)]
\item \emph{An invertible sheaf $L$ on $X$ is numerically equivalent to
$0$ if and only if}
\[
  \langle c_1(L)\cdot c_1(H)^{r-1}\rangle=0,
  \qquad
  \langle c_1(L)^2\cdot c_1(H)^{r-2}\rangle\geq0;
\]
\item \emph{For every invertible sheaf $M$ on $X$ which is not
numerically equivalent to $H$,}
\[
  \langle c_1(M)^2\cdot c_1(H)^{r-2}\rangle
  <\frac{\langle c_1(M)\cdot c_1(H)^{r-1}\rangle^2}
  {\langle c_1(H)^r\rangle}.
\]
\end{enumerate}

Assertion (ii) follows immediately from (i), by the argument of 7.3. To
prove (i), let $Y$ be a closed integral curve in $X$; we wish to prove
that $\langle c_1(L)\cdot Y\rangle=0$. Embed $X$ in $\mathbf P^N_k$ by
means of $H^{\otimes n}$, with $n>0$. Let $I=\bigoplus I_\nu$ be the
homogeneous ideal of $Y$, and let $\nu_0$ be an integer such that
$I_{\nu_0}$ generates $I_\nu$ for $\nu\geq\nu_0$. Then the
hypersurfaces $Z$ of degree $\nu_0$ containing $Y$ have no common points
outside $Y$. By induction on $r=\dim X$ and Bertini's theorem, one finds
$r-2$ hypersurfaces $Z_1,\ldots,Z_{r-2}$ of degree $\nu_0$ containing
$Y$ such that $X'=Z_1\cap\cdots\cap Z_{r-2}\cap X$ is irreducible.
Since $\langle c_1(L)\cdot Y\rangle
=\langle c_1(L|_{X'})\cdot Y\rangle$, it is enough to prove that
$L|_{X'}$ is numerically equivalent to $0$ on $X'$. By the projection
formula, $\langle c_1(L|_{X'})\cdot c_1(H|_{X'})\rangle=0$ and
$\langle c_1(L|_{X'})^2\rangle\geq0$, so the assertion follows from 7.1.

\sgaStableTarget{sga6.e13.cor.7-5}\textbf{Corollary 7.5.} \emph{Let $k$ be an algebraically closed field,
let $Y$ be a projective $k$-scheme equipped with an ample sheaf $H$, and
let $X$ be the zero scheme of a section $\varphi$ of $H$. Suppose that,
for every irreducible component $Y^i$ of $Y$, the irreducible components
of $X\cap Y^i=V(\varphi|_{Y^i})$
% [SGA6 current-rescribe idx679 | printed 666 | source-PDF 669 | high-res page 680]
have dimension at least $2$. Then the morphism
$A^1_{\mathrm{num}}(Y)\longrightarrow A^1_{\mathrm{num}}(X)$ is
injective.}

Indeed, one may plainly suppose that $Y$ is integral, and the assertion follows from \hyperlink{sga6.e13.cor.7-4}{7.4(i)} by the projection formula.

\sgaStableTarget{sga6.e13.bibliography}\subsection*{Bibliography}
\begin{enumerate}[label={[\arabic*]}]
\item S. Abhyankar, ``Resolution of singularities of arithmetical surfaces,'' in \emph{Arithmetical Algebraic Geometry}, Purdue University, Harper \& Row, New York, 1965.
\item D. Mumford, ``Pathologies III,'' \emph{Amer. J. Math.}, 1967, pp. 94--103. \emph{Amer. J. Math.}, 1967, pp. 94--103.
\item A. Grothendieck and J. Dieudonn\'e, \emph{EGA, \'El\'ements de G\'eom\'etrie Alg\'ebrique}, Publ. Math. Inst. des Hautes \'Etudes Sci., 1960.
\item A. Grothendieck, \emph{FGA, Fondements de la g\'eom\'etrie alg\'ebrique}, extracts from the S\'eminaire Bourbaki 1957--1962.
% SGA-SRC-001447
\item A. Grothendieck et al., \emph{SGA, S\'eminaire de g\'eom\'etrie alg\'ebrique}, 1960--66, I.H.E.S.
\item O. Jussila, ``Formalisme des intersections sur les sch\'emas alg\'ebriques propres,'' Expos\'e X in this seminar.
\item M. Raynaud, ``Un th\'eor\`eme de repr\'esentabilit\'e relative sur le foncteur de Picard,'' Expos\'e XII in this seminar.
\end{enumerate}

% [SGA6 current-rescribe idx680 | printed 667 | source-PDF 670 | high-res page 681]
\sgaStableTarget{sga6.e14.expose}\section*{Expos\'e XIV. Open problems in intersection theory}
\begin{center}
by A. Grothendieck
\end{center}

In the present expos\'e we review some of the problems connected with questions encountered in this seminar which deserve to be studied. We leave aside the problems connected with Serre's \emph{local intersection theory}~[14], which we have not had to use in the seminar, and refer in particular to \emph{loc. cit.} for the two fundamental conjectures which remain open in that direction (\emph{loc. cit.}, Chapter V, \S4). We also pass over the problems, very important for arithmetic applications, connected with the \emph{structure of groups of cycle classes}, especially on a smooth projective algebraic scheme. For discussion of these problems, we refer to S. Kleiman, ``Algebraic cycles and the Weil conjectures,'' in \emph{Dix expos\'es sur la cohomologie des sch\'emas} (Masson--North Holland Publishing Co.). In no. 4 of the present expos\'e we also give some supplements to the seminar, concerning in particular the \emph{Chow ring}, which had not found a place in the preceding expos\'es.

\sgaStableTarget{sga6.e14.sec.1}\section*{1. The operations $\Lambda^i$ in the derived category $D(X)$\footnote{Cf. the note on p. 23.}}

Let $X$ be a locally ringed topos, and let $K^\bullet(X)$ be the ring of classes of perfect complexes on $X$ (\hyperlink{sga6.e04.expose}{Expos\'e IV}). One would like to define on $K^\bullet(X)$ a $\lambda$-ring structure (\hyperlink{sga6.e05.expose}{Expos\'e V}), functorial in $X$, such that the natural homomorphism
\[
  K^\circ_{\mathrm{naive}}(X)\longrightarrow K^\bullet(X)
\]
(where the first member is the ring of classes of locally free modules on $X$, equipped with the $\lambda$-structure defined in \hyperlink{sga6.e05.expose}{Expos\'e V}) is a $\lambda$-homomorphism. For this purpose, one would have to define functors, in general nonadditive, of ``exterior power''
\[
  \Lambda^i:D^-(X)\longrightarrow D^-(X),
\]
transforming pseudo-coherent (respectively perfect) complexes into pseudo-coherent (respectively perfect) complexes and having the formal properties with respect to exact triangles needed to define a $\lambda$-structure on $K^\bullet(X)$ by the formula
\sgaStableTarget{sga6.e14.eq.1-1}\begin{equation}\tag{1.1}
  \lambda^i\bigl(\cl(L^\bullet)\bigr)
  =\cl\bigl(\Lambda^iL^\bullet\bigr).
\end{equation}

% [SGA6 current-rescribe idx681 | printed 668 | source-PDF 671 | high-res page 682]
When $X$ is of characteristic zero, that is, when its structure sheaf is a sheaf of $\Q$-algebras, the problem is solved simply by putting
\sgaStableTarget{sga6.e14.eq.1-2}\begin{equation}\tag{1.2}
  \Lambda^iL^\bullet
  =\left({}^{\otimes i}L^\bullet\right)^{\mathrm{alt}},
\end{equation}
where the $i$-th tensor power, in the sense of derived categories, of
$L^\bullet$ is regarded as an object of $D^-(\mathcal O_i)$, with
$\mathcal O_i=\mathcal O_X[\mathfrak S_i]$ the algebra of the symmetric
group $\mathfrak S_i$ over $\mathcal O_X$, and where the superscript
\(\mathrm{alt}\) denotes passage to the ``alternating part under the
action of $\mathfrak S_i$,'' an operation which makes sense because of
the characteristic-zero hypothesis. A formula of this type was already
used in [8] (cited below as [RRR]) in the first proof of the
Riemann--Roch theorem.

Without a characteristic hypothesis, when $L^\bullet$ is zero in degrees $>0$, Dold--Puppe~[6] gives a canonical homotopy from $L^\bullet$ to the complex associated with a semisimplicial module $L^{\prime\bullet}$ canonically attached to $L^\bullet$. For $L^\bullet$ flat, say, a fairly natural candidate for $\Lambda^iL^\bullet$ would be the complex associated with the semisimplicial module $\Lambda^iL^{\prime\bullet}$, whose degree-$n$ component is $\Lambda^iL'_n$. But to extend the definition to complexes of $D^-(X)$ which are no longer assumed to be zero in degrees $>0$, one would have to study the behavior of the preceding operation with respect to translation $L^\bullet\mapsto L^\bullet[1]$, or \emph{suspension} in Dold--Puppe terminology~[6]. Unfortunately, there does not seem to be an isomorphism $\Lambda^i\bigl(L^\bullet[1]\bigr)\simeq\Lambda^i(L^\bullet)[i]$ which would make the desired extension trivial.

\sgaStableTarget{sga6.e14.sec.2}\section*{2. The Riemann--Roch formula without projectivity hypotheses}

\sgaStableTarget{sga6.e14.para.2-1}\textbf{2.1.} Once one has a natural $\lambda$-structure on the $K^\bullet(X)$, it is possible to formulate the Riemann--Roch theorem, in the form contemplated in \hyperlink{sga6.e08.expose}{Expos\'e VIII}, for any proper complete-intersection morphism $f:X\to Y$, and even, more generally, for a relative scheme $X$ over a locally ringed topos $Y$, with $X$ still subject to the corresponding relative conditions. The Todd class of $f$ is defined by means of the class in $K^\bullet(X)$ of the relative cotangent complex $T_f$ constructed in \hyperlink{sga6.e00.sec.4.4}{Expos\'e 0, 4.4}; that complex is perfect here by the local hypothesis on $f$. It seems clear that a proof of the Riemann--Roch formula in this general case will require essentially new ideas, and it is possible that these same ideas could provide the key

% [SGA6 current-rescribe idx682 | printed 669 | source-PDF 672 | high-res page 683]
to a proof of the analogous formula in the complex analytic (or rigid-analytic!) case. Recall that, in the complex analytic case, only the case of a locally free sheaf $F$ on a compact nonsingular variety $X$, over a one-point $Y$, is currently known, by the Atiyah--Singer index formula; the case in which $F$ is only assumed coherent still seems to escape presently known techniques.

\sgaStableTarget{sga6.e14.para.2-2}\textbf{2.2.} As an essential ingredient of the Riemann--Roch formula,
recall also the \emph{finiteness problem} raised in \hyperlink{sga6.e03.expose}{Expos\'e III}: if
$f:X\to Y$ is a proper morphism and $L^\bullet$ is a complex on $X$
which is pseudo-coherent relative to $Y$, is $Rf_*(L^\bullet)$
pseudo-coherent? A general proof of this conjecture will probably use
algebraic ideas which will yield a proof of Grauert's ``finiteness
theorem''~[7] under appreciably more general conditions than those of
\emph{loc. cit.}, working from the outset over topologically ringed
spaces (or topoi) in $\mathbb C$-algebras $Y$ more general than analytic
spaces. This would include the cases in which $Y$ is a locally compact
space with the sheaf of continuous complex-valued functions, a
$C^m$-manifold with the sheaf of complex-valued $C^m$-functions, a real
analytic space, and so forth.

\sgaStableTarget{sga6.e14.para.2-3}\textbf{2.3.} Finally, in these questions the properness hypothesis on $f$ should be replaced by the hypothesis that the supports of the cohomology sheaves $H^i(L^\bullet)$ of the complexes under consideration on $X$ are proper over $Y$.
% SOURCE-EMENDATION: the print says that their ``groups'' are proper over Y;
% ``supports'' is the mathematically meaningful reading retained here.

\sgaStableTarget{sga6.e14.sec.3}\section*{3. A Riemann--Roch formula ``without denominators'' for an immersion}

\sgaStableTarget{sga6.e14.para.3-1}\textbf{3.1.} Let $f:Y\to X$ be a regular closed immersion (\hyperlink{sga6.e07.expose}{Expos\'e VII}), with conormal sheaf $\underline N$. For simplicity, suppose that $X$ admits an ample invertible sheaf. The Riemann--Roch formula for $y\in K^\bullet(Y)$,
\[
  \ch_X(f_*(y))
  =f_*\bigl(\ch_Y(y)\Todd(-\check{\underline N})\bigr),
\]
together with the general formula (\hyperlink{sga6.e07.expose}{Expos\'e VII})
\[
  f_*(y')f_*(y'')
  =f_*\bigl(y'y''c_d(\underline N)\bigr),
  \qquad y',y''\in\Gr^\bullet(Y)_{\Q},
\]
allows one, assuming $f$ has codimension $d$, to express the Chern classes of $f_*(y)$ in $\Gr^\bullet(X)_{\Q}$ in the form

% [SGA6 current-rescribe idx683 | printed 670 | source-PDF 673 | high-res page 684]
\sgaStableTarget{sga6.e14.eq.3-1}\begin{equation}\tag{3.1}
  c_i(f_*(y))
  =f_*\bigl(P_i((c_\alpha(y))_\alpha,(c_\beta(\underline N))_\beta)\bigr),
\end{equation}
where $P_i$ is a universal polynomial in the $c_\alpha(y)$ and the $c_\beta(\underline N)$, a polynomial which a priori has rational coefficients. In fact, one finds in [RRR] that these polynomials have \emph{integer coefficients}. This raises the question of the validity of a Riemann--Roch formula ``without denominators.''

In fact, such a formula for $c_i(f_*(y))$ cannot in general be sought in $\Gr^\bullet(X)$ itself, since the homomorphism $f_*:K^\bullet(Y)\longrightarrow K^\bullet(X)$ is not itself compatible with the $\lambda$-filtrations, modulo a shift by $d$ \hyperlink{sga6.e14.para.4-6}{(4.6)}, but is compatible only up to torsion. Thus it is not possible to define a homomorphism $f_*:\Gr^\bullet(Y)\longrightarrow\Gr^\bullet(X)$ inducing the natural homomorphism $f_*:\Gr^\bullet(Y)_{\Q}\longrightarrow\Gr^\bullet(X)_{\Q}$ and making it possible to interpret \hyperlink{sga6.e14.eq.3-1}{(3.1)} as a formula in $\Gr^\bullet(X)$ rather than only in $\Gr^\bullet(X)_{\Q}$.

When $X$ is regular, however, so that $K^\bullet(X)\simeq K_\bullet(X)$ (\hyperlink{sga6.e04.para.2-5}{IV 2.5}), one can equip $K^\bullet(X)$ with the \emph{topological filtration} by codimension of supports, inherited from the corresponding filtration of $K_\bullet(X)$. It is plausible that this filtration is compatible with the ring structure. This is proved at least in [RRR] when $X$ is smooth over a field $k$, using the theory of the Chow ring [4]. When this is so, the associated graded
\[
  \Gtop^\bullet(X)
\]
is a \emph{graded ring}; and, if $Y$ is also regular, one has precisely what is needed to obtain a natural homomorphism of degree $d$ of graded groups
\[
  f_*=f_*^{\mathrm{top}}:\Gtop^\bullet(Y)
  \longrightarrow\Gtop^\bullet(X).
\]
Since the $\lambda$-filtration of $K^\bullet(X)$ is finer than the topological filtration by codimension just considered (\hyperlink{sga6.e10.cor.1-3-3}{X 1.3.3}), one obtains a canonical homomorphism of graded groups (and even of graded rings),
\[
  \Gr^\bullet(X)\longrightarrow\Gtop^\bullet(X),
\]
which is, moreover, an isomorphism modulo torsion (Expos\'e VIII). This makes it possible, as in [RRR], to consider a theory of Chern classes on $X$ with values in $\Gtop^\bullet(X)$, and likewise on $Y$, and thus to interpret \hyperlink{sga6.e14.eq.3-1}{(3.1)} as a formula in $\Gtop^\bullet(X)$. In this form, the formula is proved in [RRR] when $X$ is of \emph{characteristic zero}, at least if $X$ is smooth over a field (an assumption used only to ensure that the topological filtration of $K^\bullet(X)$ is compatible with the ring structure). The characteristic-zero hypothesis is used essential%
% [SGA6 current-rescribe idx684 | printed 671 | source-PDF 674 | high-res page 685]
ly to have formulas of the type \hyperlink{sga6.e14.eq.1-1}{(1.1)}, \hyperlink{sga6.e14.eq.1-2}{(1.2)}, and essentially the same proof should apply without a characteristic hypothesis once the question raised in no. 1 has been answered affirmatively.

\sgaStableTarget{sga6.e14.para.3-2}\textbf{3.2.} When $X$ is smooth over a field, one also has the Chow ring
\[
  A(X),
\]
and a theory of Chern classes on $X$ with values in $A(X)$~[9], finer than the theory with values in $\Gtop^\bullet(X)$ just introduced. The latter Chern classes are the images of the Chow-style Chern classes under the canonical surjective homomorphism (cf. \hyperlink{sga6.e14.para.4-1}{4.1})
\[
  \varphi:A(X)\longrightarrow\Gtop^\bullet(X).
\]
This homomorphism is an isomorphism modulo torsion, but not an isomorphism in general \hyperlink{sga6.e14.para.4-7}{(4.7)}. That said, if $X$ and $Y$ are smooth, the inclusion $f:Y\to X$ still defines a homomorphism of degree $d$ of graded groups
\[
  f_*:A(Y)\longrightarrow A(X),
\]
which makes it possible to interpret \hyperlink{sga6.e14.eq.3-1}{(3.1)} as a formula in $A^i(X)$.
This formula reduces to $0=0$ for $i<d$ and is therefore of interest
% SGA-SRC-001404
only for $i\geq d$. Even for $i=d=2$, this formula has not been proved at
present (cf. \hyperlink{sga6.e14.para.4-3}{4.3}), nor has the less precise formula deduced from it by
replacing rational equivalence of algebraic cycles by algebraic
equivalence. Unlike the analogous formula in $\Gtop^\bullet(X)$, there
does not seem to be any convincing reason to suppose that the formula
without denominators should be true in $A(X)$. A good way to test it,
given our almost total ignorance of the nature of the groups
$A^i(X)$\footnote{The formula, as well as the variant mentioned in \hyperlink{sga6.e14.para.3-3}{3.3},
has since been proved by P. Jouanolou (Nov. 1969). Cf.} for $i>2$, seems
to be to test the formulas that can be deduced from it in groups
$A^1(Z)=\Pic(Z)$, for example by applying to both sides of \hyperlink{sga6.e14.eq.3-1}{(3.1)}
$g_*:A(X)\longrightarrow A(Z)$,
where $g:X\to Z$ is a proper morphism to a smooth quasi-projective $Z$ of dimension $\dim X-i+1$.

\sgaStableTarget{sga6.e14.para.3-3}\textbf{3.3.} Finally, to give a meaning to \hyperlink{sga6.e14.eq.3-1}{(3.1)}, one may also use the cohomological theory of Chern classes (SGA 5 VII), whenever one has \emph{Gysin homomorphisms}
\[
  f_*:H^p(Y,\mu_n^{\otimes i})
  \longrightarrow H^{p+2d}(X,\mu_n^{\otimes(i+d)}).
\]

% END INCLUDED SGA6_English_current_buildable_source_20260823/fragments/SGA6_idx666_684_English_Reconciled_Prefix.tex

% BEGIN INCLUDED SGA6_English_current_buildable_source_20260823/fragments/SGA6_XIV_idx685_692_English_SourceChecked_body.tex
This holds at least when $Y$ and $X$ are both smooth over a scheme $S$, by the
theorem of relative cohomological purity (SGA 4 XVI 3.7); conjecturally, it
also holds when $Y$ and $X$ are both regular, by the theorem of absolute
cohomological purity, proved at least for excellent $X$ of characteristic
zero (SGA 4 XIX 3.2).

When $X$ and $Y$ are smooth over the field of complex numbers, one may also
use the integral Chern class on the locally compact space $X(\C)$ of complex
points of $X$. Formula \hyperlink{sga6.e14.eq.3-1}{(3.1)} was proved in this case by Atiyah and
Hirzebruch~[1], who more generally treat embeddings of complex analytic
varieties (working with naive $K'$). It follows, by the Lefschetz principle
and the general results of SGA 4, that the $\ell$-adic formula (3.1) holds
whenever $Y$ and $X$ are smooth over an algebraically closed field of
characteristic zero. This makes the cohomological formula \hyperlink{sga6.e14.eq.3-1}{(3.1)} extremely
plausible. Nevertheless, it should be noted that, at present, it has not been
proved even when $X$ and $Y$ are smooth over a field $k$ of characteristic
zero, if $k$ is not algebraically closed.

\noindent\sgaStableTarget{sga6.e14.para.3-4}\textbf{3.4.}
One may also work with the cohomological theory of Chern classes, in the de
Rham or Hodge style~[13]. The question of the validity of \hyperlink{sga6.e14.eq.3-1}{(3.1)} in this
setting arises only when $X$ is not of characteristic zero, since one knows
in any case (no.~6) that the formula in question is true modulo torsion. For
smooth quasi-projective schemes over a field $k$, the de Rham or Hodge
version of \hyperlink{sga6.e14.eq.3-1}{(3.1)} would moreover follow from the version using cycle classes
modulo algebraic equivalence. Thus the de Rham or Hodge cohomological version
provides another accessible way to test the validity of the cycle-class
version of \hyperlink{sga6.e14.eq.3-1}{(3.1)}.

\noindent\sgaStableTarget{sga6.e14.para.3-5}\textbf{3.5.}
The proof~[1] of the transcendental cohomological version of \hyperlink{sga6.e14.eq.3-1}{(3.1)} is
essentially transcendental in nature and uses tubular neighborhoods of $Y$
in $X$. This suggests that a proof of a formula of type \hyperlink{sga6.e14.eq.3-1}{(3.1)} may be tied to
the systematic consideration of $K^\bullet$-groups and cohomology groups with
support in $Y$, which might provide an adequate technical means of
\emph{localization around $Y$}.

\sgaStableTarget{sga6.e14.sec.4}\section*{4. Relations between $K^\bullet(X)$ and the Chow ring $A(X)$}

\noindent\sgaStableTarget{sga6.e14.para.4-1}\textbf{4.1.}
Suppose for simplicity that $X$ is a smooth quasi-projective scheme over a
field $k$ (see no.~8 for a more general case), so that the Chow ring $A(X)$
of cycle classes modulo rational equivalence is available~[4]. Two variants
of this ring, defined under appreciably more general conditions, are
$\operatorname{Gr}^\bullet(X)$, the graded ring associated with the
$\lambda$-filtration of $K^\bullet(X)$, and $\Gtop^\bullet(X)$, the graded
ring associated with the filtration of
$K^\bullet(X)\simeq K_\bullet(X)$ by codimension of supports. There are
theories of Chern classes with values in each of these three rings (see [9]
for the case of $A(X)$, and this seminar for
$\operatorname{Gr}^\bullet(X)$, which maps into $\Gtop^\bullet(X)$). There is
also an $\ell$-adic cohomological theory (SGA 5 VII). Thus, when $\ell$ is
prime to the characteristic of $k$,
\[
 c_i(E)\in H^{2i}(X,\Zell(i)),
\]
where $E$ is a vector bundle on $X$, or more generally any element of
$K^\bullet(X)$. These theories are related by the following canonical
homomorphisms, which are compatible with the multiplicative structures:
\sgaStableTarget{sga6.e14.eq.4-1}\begin{equation*}\tag{4.1}
\begin{tikzcd}[column sep=large,row sep=large]
 H^{2*}(X,\Zell(*)) & {} \\
 A(X) \arrow[u,"\gamma"] \arrow[r,"\varphi"] & \Gtop^\bullet(X) \\
 {} & \operatorname{Gr}^\bullet(X) \arrow[u,"j"']
\end{tikzcd}
\end{equation*}
where
\[
 H^{2*}(X,\Zell(*))=\bigoplus_i H^{2i}(X,\Zell(i)).
\]

The homomorphism $j$ comes from the fact that the $\lambda$-filtration of
$K^\bullet(X)$ is finer than the topological filtration by codimension of
support (\hyperlink{sga6.e10.expose}{Expose X}), and it is an isomorphism modulo torsion (Expose VII
4.11). The homomorphism $\varphi$, defined for example in [RRR], associates
to the class of the codimension-$i$ cycle carried by a closed subscheme
$Y\subset X$ the element $\cl^i(\mathcal O_Y)$ of
\[
 \Filt_{\mathrm{top}}^i K^\bullet(X)/
 \Filt_{\mathrm{top}}^{i+1}K^\bullet(X).
\]
It is surjective, and we shall see in \hyperlink{sga6.e14.para.4-2}{4.2} that it is an isomorphism modulo
torsion. Hence, modulo torsion, the three rings $A(X)$,
$\Gtop^\bullet(X)$, and $\operatorname{Gr}^\bullet(X)$ are canonically
isomorphic; consequently, the Riemann--Roch formulas expressed using any
one of these rings after tensoring with $\Q$ are equivalent.

The graded ring $\operatorname{Gr}^\bullet(X)$, used in this seminar chiefly
as a substitute for the Chow ring, has the considerable advantage of being
defined for every locally ringed topos, and in particular for every scheme,
whether regular or not. Its disadvantage \hyperlink{sga6.e14.para.4-6}{(4.6)} with respect to $A(X)$ and
$\Gtop^\bullet(X)$ is that it is not covariant as it stands for proper
morphisms, even between smooth quasi-projective schemes over a field; it
acquires that covariance only after tensoring with $\Q$.

Another advantage of $A(X)$ over its competitors $\Gtop^\bullet(X)$ and
$\operatorname{Gr}^\bullet(X)$ is the homomorphism
\[
 A(X)\longrightarrow H^{2*}(X,\Zell(*)),
\]
which permits cohomological Chern classes to be interpreted as images of
Chern classes with values in $A(X)$. In general this homomorphism does not
factor through $\Gtop^\bullet(X)$ \hyperlink{sga6.e14.para.4-7}{(4.7)}, and it is at least doubtful that
there is in general a homomorphism
\[
 \operatorname{Gr}^\bullet(X)\longrightarrow H^{2*}(X,\Zell(*))
\]
carrying Chern classes to Chern classes (see no.~5). These difficulties
disappear after tensoring with $\Q$, and one obtains a homomorphism
\sgaStableTarget{sga6.e14.eq.4-2}\begin{equation*}\tag{4.2}
 \Gtop^\bullet(X)\longrightarrow H^{2*}(X,\Qell(*))
 =H^{2*}(X,\Zell(*))\otimes_{\Zell}\Qell
\end{equation*}
that carries Chern classes to Chern classes.

\noindent\sgaStableTarget{sga6.e14.para.4-2}\textbf{4.2.}
To show that the surjective homomorphism $\varphi$ in \hyperlink{sga6.e14.eq.4-1}{(4.1)} is an
isomorphism modulo torsion, one uses the Chern homomorphism obtained from the
theory of Chern classes~[9],
\[
 K^\bullet(X)\longrightarrow \Ch(A(X)),
\]
where the target is the Chern ring of $A(X)$ as defined in \hyperlink{sga6.e05.expose}{Expose V}. This
homomorphism is compatible with the topological filtration of $K^\bullet(X)$
and the natural filtration of $\Ch(A(X))$, and therefore induces a
homomorphism of graded groups
\sgaStableTarget{sga6.e14.eq.4-3}\begin{equation*}\tag{4.3}
 \psi:\Gtop^\bullet(X)\longrightarrow A(X).
\end{equation*}
One proves the formulas
\sgaStableTarget{sga6.e14.eq.4-4}\begin{equation*}\tag{4.4}
\left\{
\begin{aligned}
 \psi\varphi(x)&=(-1)^{i-1}(i-1)!\,x
   &&\pmod{\text{torsion}},\quad x\in A^i(X),\\
 \varphi\psi(x)&=(-1)^{i-1}(i-1)!\,x
   &&\pmod{\text{torsion}},\quad x\in \Gtop^i(X).
\end{aligned}
\right.
\end{equation*}
The first formula implies the second because $\varphi$ is surjective. These
formulas show that $\varphi$ and $\psi$ are isomorphisms modulo torsion and,
up to the factor $(-1)^{i-1}(i-1)!$ in degree $i$, are inverse to one
another.

To prove the first formula in \hyperlink{sga6.e14.eq.4-4}{(4.4)}, one may suppose that $x$ is represented
by an irreducible closed subscheme $Y$ of $X$ of codimension $i$. The formula
is then equivalent to
\sgaStableTarget{sga6.e14.eq.4-5}\begin{equation*}\tag{4.5}
 c_i(\mathcal O_Y)=(-1)^{i-1}(i-1)!\,\cl(Y)
   \quad\pmod{\text{torsion}}.
\end{equation*}
On the left, $\mathcal O_Y$ is regarded as a coherent sheaf on $X$ and hence
defines an element of $K^\bullet(X)$; on the right, $\cl(Y)$ denotes the
class of $Y$ in $A^i(X)$. Since $A^i(X)$ is unchanged upon deleting from
$X$ a closed subset of codimension $>i$, one reduces to the case where $Y$
is regular, and therefore smooth over the base field $k$ if $k$ is perfect.
In that case, \hyperlink{sga6.e14.eq.4-5}{(4.5)} is precisely the Riemann--Roch formula for the inclusion
$Y\to X$ and the element $1\in K^\bullet(Y)$, written using Chern classes
with values in the Chow ring $A(X)$. Under the present hypotheses, this
formula is proved in [2] by a method very close to the one used in this
seminar for the analogous formula with values in $\Gtop^\bullet(X)$.

Thus, at least when $k$ is perfect, \hyperlink{sga6.e14.eq.4-4}{(4.4)} follows, and so does the fact that
$\varphi$ is an isomorphism modulo torsion. This in turn implies, a
posteriori, that Riemann--Roch with denominators for a proper morphism of
smooth algebraic schemes over $k$ is essentially the same whether it is
written using $A(X)$, as in [2], or using $\Gtop^\bullet(X)$, as in this
seminar. The case of an arbitrary $k$ reduces easily to the algebraically
closed case: for a finite extension $k'/k$ of degree $n$, the kernel of the
natural map
\[
 A(X)\xrightarrow{f_X^*}A(X\otimes_k k')
\]
is annihilated by $n$. Indeed, the trace homomorphism
\[
 f_{X*}:A(X\otimes_k k')\longrightarrow A(X)
\]
satisfies $f_{X*}f_X^*=n\,\id$.\footnote{Thus the formula is established,
thanks to P.~Jouanolou, for $X$ smooth and quasi-projective over a field (see
the note on p.~5).}

% Current-rescribe idx688 | printed volume page 675 | source-PDF page 678
\noindent\sgaStableTarget{sga6.e14.para.4-3}\textbf{4.3.}
It seems quite plausible that formulas \hyperlink{sga6.e14.eq.4-4}{(4.4)} hold exactly as stated, and not
merely modulo torsion; equivalently, that formula \hyperlink{sga6.e14.eq.4-5}{(4.5)} holds not merely
modulo torsion. This question is a very special case of the question whether
the denominator-free Riemann--Roch formula \hyperlink{sga6.e14.eq.3-1}{(3.1)} is valid in $A(X)$ for an
immersion $Y\to X$, namely the case $i=d$. This special case may reasonably
be expected to hold. Indeed, we already noted in no.~3 that, when $k$ has
characteristic zero, the denominator-free Riemann--Roch formula is valid at
least in $\Gtop^\bullet(X)$. Consequently the two sides of \hyperlink{sga6.e14.eq.4-5}{(4.5)} have the
same image in $\Gtop^\bullet(X)$; equivalently, the second formula in \hyperlink{sga6.e14.eq.4-4}{(4.4)}
holds. This result will extend automatically to arbitrary characteristic
once

% Current-rescribe idx689 | printed volume page 676 | source-PDF page 679
the question raised in no.~1 has been satisfactorily resolved. As evidence
in support of formula \hyperlink{sga6.e14.eq.4-5}{(4.5)} in $A(X)$ itself, let us note that it is clear a
priori that the restriction of the left-hand side to $X-Y$ is zero. Thus,
since $Y$ is assumed irreducible, the left-hand side must a priori be an
integral multiple of $\cl(Y)$:
\[
 c_i(Y)=n(Y)\cl(Y).
\]
This already proves the exact formula \hyperlink{sga6.e14.eq.4-5}{(4.5)} whenever $\cl(Y)$ is not a
torsion element of $A^i(X)$. It is immediate that this is the case if $X$ is
projective (and not merely quasi-projective) over the field $k$: indeed, if
$\xi\in A^1(X)$ is a polarization of $X$, then
% SGA-SRC-001450 — reconstructed by mathematical inference; scan inconclusive
$\deg(Y\mathbin{\cdot}\xi^{\dim Y})>0$. This proves the desired formula in that
case. It follows that the formula is still valid whenever $X$ is isomorphic
to an open subscheme of a smooth projective scheme over $k$, as will be the
case when an adequate solution of the problem of resolution of
singularities is available. In particular, this holds when $k$ has
characteristic zero, by Hironaka~[11].

The validity of the exact formula \hyperlink{sga6.e14.eq.4-5}{(4.5)}, and hence of \hyperlink{sga6.e14.eq.4-4}{(4.4)}, would imply that
the kernel of
\[
 \varphi^i:A^i(X)\longrightarrow\Gtop^i(X)
\]
is annihilated by $(i-1)!$, and in particular \emph{that
$\varphi^2:A^2(X)\to\Gtop^2(X)$ is an isomorphism}. (That $\varphi^1$ is an
isomorphism already follows immediately from \hyperlink{sga6.e10.thm.5-3-2}{X~5.3.2}.) Unless I am
mistaken, the assertion that $\varphi^2$ is an isomorphism has not yet been
proved, even in the special case $\dim X=3$. It was therefore in any event
premature for the editor to assert the validity of (4.4) and (4.5) in
[9, no.~4, $3^\circ$].

\noindent\sgaStableTarget{sga6.e14.para.4-4}\textbf{4.4.}
At present, the principal advantage of $\Gtop^\bullet(X)$ (or of
$\operatorname{Gr}^\bullet(X)$) over the Chow ring $A(X)$ is that it is
better suited to explicit calculations of certain intersection formulas
between cycles in cases of excess intersection, that is, intersections of
excessively large dimension. Indeed, the definition of the intersection
class by means of Chow's lemma, which allows a cycle to be moved sufficiently
within its class, is poorly suited to calculation. A striking example in
this regard is the following fundamental formula [\hyperlink{sga6.e07.expose}{Exposé VII}~\hyperlink{sga6.e07.para.2-6}{2.7}], for a
closed immersion $f:Y\to X$ of smooth quasi-projective schemes over $k$:
\sgaStableTarget{sga6.e14.eq.4-6}\begin{equation*}\tag{4.6}
 f^*f_*(y)=y\,c_d(\underline{N}),
\end{equation*}
a formula valid in $\Gtop^\bullet(Y)$ for $y\in\Gtop^\bullet(Y)$, where
$\underline{N}$ is the conormal sheaf

% Current-rescribe idx690 | printed volume page 677 | source-PDF page 680
of $Y$ in $X$, assumed everywhere to have rank $d$. In the cited passage,
this formula follows from a very simple computation of the sheaves
$\underline{\operatorname{Tor}}$. By \hyperlink{sga6.e14.para.4-2}{4.2}, the preceding formula implies the
same formula in $A(Y)$, but only modulo torsion. It is quite remarkable that
formula \hyperlink{sga6.e14.eq.4-6}{(4.6)} in $A(Y)$ itself has not yet been proved,\footnote{Cf. the note
on p.~23.} nor even the special case obtained by taking $y=1$, which gives
the self-intersection of $Y$:
\sgaStableTarget{sga6.e14.eq.4-7}\begin{equation*}\tag{4.7}
 f^*(\cl(Y))=c_d(\underline{N}),
\end{equation*}
even when $k$ is the field of complex numbers.

An analogous problem arises when one attempts to compute the Chow ring of
the variety $X'$ obtained by blowing up $X$ along $Y$, by imitating the
method used to compute $\operatorname{Gr}^\bullet(X')_{\Q}$
[\hyperlink{sga6.e07.expose}{Exposé VII}~\hyperlink{sga6.e07.thm.4-8}{4.8}], which also works (without tensoring with $\Q$) for the
computation of $\Gtop^\bullet(X')$. With the notation of the cited passage,
the key formula for this computation is
\sgaStableTarget{sga6.e14.eq.4-8}\begin{equation*}\tag{4.8}
 f^*(i_*(y))
   =j_*\bigl(g^*(y)c_{d-1}(\check{\underline{F}})\bigr).
\end{equation*}
At present, this formula has been established in $\Gtop^\bullet(X')$ (for
$y\in\Gtop^\bullet(Y)$) by a computation of
$\underline{\operatorname{Tor}}$ (\hyperlink{sga6.e07.expose}{Exposé VII}~\hyperlink{sga6.e07.prop.3-4}{3.4}), while its validity in
$A(X')$ (for $y\in A(Y)$) remains to be proved.

Let us note that even the $\ell$-adic analogues of formulas \hyperlink{sga6.e14.eq.4-6}{(4.6)}, \hyperlink{sga6.e14.eq.4-7}{(4.7)},
and \hyperlink{sga6.e14.eq.4-8}{(4.8)} have not yet been proved; when $y$ is the cohomology class of an
algebraic cycle, formulas \hyperlink{sga6.e14.eq.4-6}{(4.6)} and \hyperlink{sga6.e14.eq.4-8}{(4.8)} are in any event true modulo
torsion.

\noindent\sgaStableTarget{sga6.e14.para.4-5}\textbf{4.5.}
In this section and the following ones, we shall give examples for which,
in diagram \hyperlink{sga6.e14.eq.4-1}{(4.1)}, the homomorphism $j$, respectively the homomorphism
$\varphi$, is not an isomorphism. For a given $X$, it is plainly equivalent
to say that
\[
 j:\operatorname{Gr}^\bullet(X)\longrightarrow\Gtop^\bullet(X)
\]
is an isomorphism, a monomorphism, or an epimorphism. Moreover, the image of
$j$ is the subring of $\Gtop^\bullet(X)$ generated by the Chern classes of
locally free sheaves on $X$. If $K^\bullet(X)$ is generated as a
$\Z$-algebra by the classes of invertible sheaves on $X$, one sees at once
that the subring in question is also the subring generated by
$\Gtop^1(X)$, the degree-one component of $\Gtop^\bullet(X)$.

% Current-rescribe idx691 | printed volume page 678 | source-PDF page 681
Now take for $X$ a flag variety
\[
 X=G/B,
\]
where $G$ is a simply connected semisimple algebraic group over the
algebraically closed field $k$, and $B$ is a Borel subgroup~[5]. When
$k=\C$, the Bruhat cell decomposition of $X$ [5, Exposé 13] implies that
the canonical homomorphism
\sgaStableTarget{sga6.e14.eq.unnamed}\begin{equation*}\tag{*}
 K^\bullet(X)\longrightarrow K(X(\C))
\end{equation*}
is an isomorphism, as is the canonical homomorphism
\sgaStableTarget{sga6.e14.eq.unnamed.occ2}\begin{equation*}\tag{**}
 A(X)\longrightarrow H^*(X(\C),\Z),
\end{equation*}
where $X(\C)$ is the compact space of complex points of $X$, endowed with
its usual topology induced by that of $\C$. On the other hand, L.~Hodgkin
[12] proved that $K(X(\C))$ is generated by the classes of invertible
sheaves on $X$. The same is therefore true of $K^\bullet(X)$ itself;
moreover, it is not difficult to deduce, by a specialization argument using
semisimple groups split over $\Z$ (SGA 3), that the result remains valid
over every algebraically closed field $k$. The cell decomposition of $X$
(or, when $k=\C$, the bijectivity of $(**)$) also readily implies that
$A(X)$ is torsion-free. Hence, by \hyperlink{sga6.e14.para.4-2}{4.2}, the homomorphism
\[
 \varphi:A(X)\longrightarrow\Gtop^\bullet(X)
\]
is an isomorphism. It follows that $j$ is an isomorphism if and only if
$A(X)$ is generated by its elements of degree $1$, or, when $k=\C$, if and
only if the integral cohomology of $X(\C)$ is generated by its elements of
degree $2$, i.e. if the classifying space $B_{G(\C)}$ is without torsion. It
is known [15, p.~217, Theorem B, $(1)\Leftrightarrow(2)$, and 2.5] that this
holds if and only if $G$ is without torsion, equivalently if $G$ is
isomorphic to a product of groups $\operatorname{SL}(n)$ and
$\operatorname{Sp}(m)$. Thus it suffices to take for $G$ a group
$\operatorname{SO}(n)$ with $n\geq7$, or a simple exceptional group, to
obtain an example in which $j$ is not an isomorphism.

\noindent\sgaStableTarget{sga6.e14.para.4-6}\textbf{4.6.}
In this example there can be no doubt that the good ring is the torsion-free
ring
\[
 A(X)\simeq\Gtop^\bullet(X),
\]
whose basis consists of the cells in the Bruhat decomposition and which,
when $k=\C$, is also isomorphic to the integral cohomology of $X(\C)$. By
contrast, $\operatorname{Gr}^\bullet(X)$ appears as a ring having torsion
of a rather pathological nature. This example therefore confirms the view
that the invariant $\operatorname{Gr}^\bullet(X)$ is hardly interesting

% Current-rescribe idx692 | printed volume page 679 | source-PDF page 682
except modulo torsion, i.e. after tensoring with $\Q$. At the same time,
this example shows that $\operatorname{Gr}^\bullet(X)$ is covariant with
respect to proper morphisms only after tensoring with $\Q$, even if one
restricts to closed immersions of smooth quasi-projective schemes. Indeed,
if $X$ is such that $j$ is not an isomorphism, there is an integral closed
subscheme $Y$ of $X$ such that
% SGA-SRC-001452
\[
 \cl^\bullet(Y)\in K^\bullet(X)
\]
does not lie in $\Filt^i(X)$, where $i=\operatorname{codim}(Y,X)$. We may
choose $Y$ of minimal dimension, so that
\[
 \Filt^{i+1}(X)=\Filt_{\mathrm{top}}^{i+1}(X).
\]
If $Z$ is the singular locus of $Y$, and $X'=X\setminus Z$,
$Y'=Y\setminus Z$, it follows that
\[
 \cl^\bullet(Y')\in K^\bullet(X')
\]
does not lie in $\Filt^i(X')$. Here one uses the fact that
$K^\bullet(X)\to K^\bullet(X')$ is surjective and compatible with the
augmentations, and that its kernel is
\[
% SGA-SRC-001453
 \operatorname{Im}\bigl(K^\bullet(Z)\longrightarrow K^\bullet(X)\bigr)
 \subset
 \Filt^{i+1}(X)=\Filt_{\mathrm{top}}^{i+1}(X).
\]
Now consider the inclusion
\[
 f:Y'\longrightarrow X',
\]
which is a closed immersion of codimension $i$ between smooth
quasi-projective schemes over $k$, \emph{such that}
\[
 f_*:K^\bullet(Y')\longrightarrow K^\bullet(X')
\]
\emph{is not compatible with the $\lambda$-filtrations, even up to a shift
by $i$}, since $f_*(1)\notin\Filt^i(X')$.

\noindent\sgaStableTarget{sga6.e14.para.4-7}\textbf{4.7.}
Retaining the notation of \hyperlink{sga6.e14.para.4-5}{4.5}, note that \hyperlink{sga6.e09.expose}{Exposé IX} shows that Hodgkin's
theorem quoted above is equivalent to the isomorphism
\[
 K^\bullet(G)\simeq\Z,
 \qquad\text{hence}\qquad
 \Gtop^\bullet(G)\simeq\Z.
\]
The homomorphism $K^\bullet(G/B)\to K^\bullet(G)$ is a priori surjective,
and its kernel is the ideal generated by the elements
$\cl^\bullet(L)-1$, where $L$ ranges over the invertible sheaves on $X=G/B$
associated with representations $B\to\Gm$ (which in fact yield all
invertible sheaves on $X$ [5, Exposé 15]). An analogous argument shows that
\[
 A(G/B)\longrightarrow A(G)
\]
is also surjective, and that its kernel is the ideal generated by the
elements of $A^1(G/B)$ defined by representations $B\to\Gm$, hence by all
of $A^1(G/B)$.\footnote{Cf. \hyperlink{sga6.e09.prop.3-4}{IX 3.4}, 2.2.} Thus $A(G)$ is reduced to
% SGA-SRC-001454
$A^0(G)=\Z$ if and

% END INCLUDED SGA6_English_current_buildable_source_20260823/fragments/SGA6_XIV_idx685_692_English_SourceChecked_body.tex

% BEGIN INCLUDED SGA6_English_current_buildable_source_20260823/fragments/SGA6_XIV_idx693_702_English_SourceChecked_body.tex
only if $G$ is ``without torsion'', i.e. if and only if $G$ is isomorphic to
a product of groups $\operatorname{Sl}(n)$ and $\operatorname{Sp}(m)$. Thus,
in this case and only in this case, the homomorphism
\[
 \varphi:A(G)\longrightarrow\Gtop^\bullet(G)
\]
is an isomorphism. Consequently, if $G$ has torsion, the underlying scheme
provides an example of a non-bijective $\varphi$.

\noindent\sgaStableTarget{sga6.e14.para.4-8}\textbf{4.8.}
One likewise finds an example in which the homomorphism $\gamma$ of \hyperlink{sga6.e14.eq.4-1}{(4.1)}
\sgaStableTarget{sga6.e14.eq.4-9}\begin{equation}\tag{4.9}
 \gamma:A(G)\longrightarrow H^{2*}(G,\Zell(*))
\end{equation}
does not factor through $\Gtop^\bullet(G)$, or, what amounts to the same
thing here, an example in which this homomorphism is nonzero in degrees
$>0$. Restricting (as one may, thanks to suitable specialization theorems)
to the case $k=\C$, one is reduced to studying the homomorphism
\sgaStableTarget{sga6.e14.eq.4-10}\begin{equation}\tag{4.10}
 A(G)\longrightarrow H^*(G(\C),\Z),
\end{equation}
which is moreover obtained, by passage to the quotient, from the canonical
homomorphism
\[
 A(G/B)\simeq H^*(G/B(\C),\Z)
 \longrightarrow H^*(G(\C),\Z).
\]
It is tempting to conjecture that \hyperlink{sga6.e14.eq.4-10}{(4.10)} is in fact injective, since it is
also known that the torsion primes for $G(\C)$ are exactly those dividing
the order of the finite group
\[
 A^+(G)=\coprod_{i>0}A^i(G).
\]
If this were so, whenever $G$ had torsion, \hyperlink{sga6.e14.eq.4-10}{(4.10)} would not factor through
$\varphi$; hence, for a suitable $\ell$ (more precisely, for $\ell$ dividing
the order of $A^+(G)$, i.e. for $\ell$ a torsion prime of $G$), the same
would hold for \hyperlink{sga6.e14.eq.4-9}{(4.9)}. We can at least assert that, for the group $G=G_2$
over $\C$, \hyperlink{sga6.e14.eq.4-10}{(4.10)} is nonzero in degree $>0$: more precisely, there exists an
element of $A^3(G)$ whose image in $H^6(G(\C),\Z)$ has order $2$. It follows
that \hyperlink{sga6.e14.eq.4-9}{(4.9)} does not factor through $\varphi$ when $\ell=2$. The fact used
about $G_2$, which amounts to saying that there exists in
$H^6(G/B(\C),\Z)$ an element whose image in $H^6(G(\C),\Z)$ is nonzero, is
found in [3], via the familiar dictionary between complex semisimple
algebraic groups and compact Lie groups; it was kindly communicated to me
by J.-P. Serre.

% Current-rescribe idx694 | printed volume page 681 | source-PDF page 684
\sgaStableTarget{sga6.e14.sec.5}\section*{5. Relations between $\operatorname{Gr}^\bullet(X)$ and
$H^{2*}(X,\Zell(*))$}

\noindent\sgaStableTarget{sga6.e14.para.5-1}\textbf{5.1.}
The example given in \hyperlink{sga6.e14.para.4-8}{4.8} does not exclude the possibility that, for an
arbitrary scheme $X$, one might find a ring homomorphism
\[
 g':\operatorname{Gr}^\bullet(X)
 \longrightarrow H^{2*}(X,\Zell(*)),
\]
such that, for every locally free module $E$ on $X$, $g^i$ sends the Chern
class $c_i(E)$ of Exp.~V to the cohomological Chern class of SGA 5 VII. Of
course, $\ell$ here denotes a prime distinct from the residual
characteristics of $X$. Since the ring $\operatorname{Gr}^\bullet(X)$ is
% SGA-SRC-001421
generated by the $c_i(E)$, one sees in any case that $g'$ is uniquely
determined by the requirement that it be compatible with Chern classes.

\noindent\sgaStableTarget{sga6.e14.para.5-2}\textbf{5.2.}
Using the Chern homomorphism
\[
 K^\bullet(X)\longrightarrow
 \Ch\bigl(H^{2*}(X,\Zell(*))\bigr),
\]
one obtains, by passing to the associated gradeds respectively for the
filtration on the first member and the filtration by degree on the second,
a homomorphism of graded groups (but not of rings!)
\[
 \psi:\operatorname{Gr}^\bullet(X)
 \longrightarrow H^{2*}(X,\Zell(*)),
\]
which satisfies
\[
 \psi(c_i(E))=(-1)^{i-1}(i-1)!\,c_i^\ell(E),
\]
where $c_i^\ell(E)$ denotes the $\ell$-adic Chern class. Dividing the
components $\psi^i$ by the indicated factors therefore gives a ring
homomorphism
\[
 g_\Q:\operatorname{Gr}^\bullet(X)
 \longrightarrow H^{2*}(X,\Qell(*))
 =H^{2*}(X,\Zell(*))\otimes_{\Zell}\Qell,
\]
sending Chern classes to Chern classes, and such that the image of the first
member is contained in that of $H^2(X,\Zell(*))$. Consequently, the question
of the existence of $g'$ has an affirmative answer provided
$H^{2*}(X,\Zell(*))$ is replaced by its quotient by its torsion subgroup.
The first doubtful case for defining the components $g^i$ of the
hypothetical homomorphism $g$ is that of $g^3$; a priori one succeeds only
in defining $2g^3=\psi^3$. The answer to the question of the existence of
$g^3$ would be negative if, for example, one could find three invertible
sheaves $L_1,L_2,L_3$ on $X$ satisfying
\[
 (L_1-1)(L_2-1)(L_3-1)=0
 \quad\text{in }K^\bullet(X),
\]

% Current-rescribe idx695 | printed volume page 682 | source-PDF page 685
but such that
\[
 c_1(L_1)c_1(L_2)c_1(L_3)\ne0
 \quad\text{in }H^6(X,\mu_2^{\otimes3}).
\]

\noindent\sgaStableTarget{sga6.e14.para.5-3}\textbf{5.3.}
When $X$ is smooth and quasi-projective over a field $k$, one may similarly
ask whether there exists a ring homomorphism
\[
 g:\operatorname{Gr}^\bullet(X)\longrightarrow A(X)
\]
taking values in the Chow ring $A(X)$ and sending Chern classes to Chern
classes. The preceding observations may be repeated in this connection;
they show in particular how to construct $(-1)^{i-1}(i-1)!g^i$ directly
when $g^i$ itself is unavailable.

\noindent\sgaStableTarget{sga6.e14.para.5-4}\textbf{5.4.}
Let us note, to conclude this section, that essentially the same question
arises if, instead of a scheme $X$, one considers a compact topological
space (say) $X$. One can still consider the associated graded
$\operatorname{Gr}^\bullet(X)$ of $K^\bullet(X)$ endowed with its
$\lambda$-filtration, and seek a ring homomorphism
\[
 \operatorname{Gr}^\bullet(X)\longrightarrow H^{2*}(X,\Z)
\]
sending Chern classes to Chern classes.

\sgaStableTarget{sga6.e14.sec.6}\section*{6. Cohomological Riemann--Roch theorem and Gysin homomorphism}

\noindent\sgaStableTarget{sga6.e14.para.6-1}\textbf{6.1.}
For simplicity, we shall confine ourselves here to schemes $X$ admitting an
ample invertible module, so that there is no need to distinguish between the
naive $K^\bullet(X)$ and the other one. Let us begin with the Riemann--Roch
formula ``with denominators'', for a proper complete-intersection morphism
$f:X\longrightarrow Y$. For a prime $\ell$ prime to the residual
characteristics of $Y$, one can define the $\ell$-adic Chern classes of a
locally free module $E$,
\[
 c_i(E)\in H^{2i}(X,\Zell(i))
\]
(SGA 5 VII), corresponding to a homomorphism
\sgaStableTarget{sga6.e14.eq.6-1}\begin{equation}\tag{6.1}
 K^\bullet(X)\longrightarrow \Ch\bigl(H^{2*}(X,\Zell(*))\bigr).
\end{equation}
The question therefore arises of giving a cohomological version of the
Riemann--Roch formula. The only ingredient missing in order to give meaning
to the formula is a Gysin homomorphism

% Current-rescribe idx696 | printed volume page 683 | source-PDF page 686
\sgaStableTarget{sga6.e14.eq.6-2}\begin{equation}\tag{6.2}
 H^n(X,\Qell(i))\longrightarrow H^{n-2d}(Y,\Qell(i-d)),
\end{equation}
where $d$ is the virtual relative dimension of $f$. Assuming this
homomorphism to be defined, and taking into account the Riemann--Roch formula
in the form of \hyperlink{sga6.e08.expose}{Exp.~VIII}, with values in
$\operatorname{Gr}^\bullet(Y)_{\Q}$, one sees that the validity of the
$\ell$-adic Riemann--Roch formula for the morphism $f$ is exactly equivalent
to the commutativity of the square
\sgaStableTarget{sga6.e14.eq.6-3}\begin{equation}\tag{6.3}
\begin{tikzcd}[column sep=large,row sep=large]
 \operatorname{Gr}^\bullet(X)_{\Q} \arrow[r] \arrow[d]
   & H^{2*}(X,\Qell(*)) \arrow[d] \\
 \operatorname{Gr}^\bullet(Y)_{\Q} \arrow[r]
   & H^{2*}(Y,\Qell(*)),
\end{tikzcd}
\end{equation}
where the horizontal arrows are obtained from \hyperlink{sga6.e14.eq.6-1}{(6.1)} by passage to the
associated gradeds and division, so as to send Chern classes to Chern
classes, while the vertical arrows are respectively the arrow of \hyperlink{sga6.e08.expose}{Exp.~VIII}
and the Gysin homomorphism in cohomology (assumed given).

\noindent\sgaStableTarget{sga6.e14.para.6-2}\textbf{6.2.}
Note that the Gysin homomorphism is defined in all cases if $Y$ is smooth
over a field $k$ (SGA 5 IV), using the theorem of relative cohomological
purity over $k$. The same construction applies when $Y$ is excellent,
regular, and of characteristic zero, thanks to the theorem of absolute
cohomological purity of SGA 4 XIX; and it will automatically apply to every
excellent regular $Y$, without restriction on the characteristic, once the
theorem of absolute cohomological purity has been proved in full generality.
Finally, the same construction applies, thanks to the theorem of relative
cohomological purity (SGA 4 XVI 3.7), when $X$ and $Y$ are both smooth over
the same scheme $S$ and $f$ is an $S$-morphism. It is natural to conjecture
that, in each of these cases in which the Gysin homomorphism is defined,
diagram \hyperlink{sga6.e14.eq.6-3}{(6.3)} is indeed commutative. The only case in which this assertion
is currently verified is that in which $X$ and $Y$ are both smooth over a
field $k$; it then follows from the commutativity of the analogous diagram
\sgaStableTarget{sga6.e14.eq.6-4}\begin{equation}\tag{6.4}
\begin{tikzcd}[column sep=large,row sep=large]
 A(X) \arrow[r] \arrow[d]
   & H^{2*}(X,\Zell(*)) \arrow[d] \\
 A(Y) \arrow[r]
   & H^{2*}(Y,\Zell(*)),
\end{tikzcd}
\end{equation}

% Current-rescribe idx697 | printed volume page 684 | source-PDF page 687
involving the rings $A(X)$ and $A(Y)$ of cycle classes modulo rational
equivalence, whose commutativity was established in SGA 5 IV.

\noindent\sgaStableTarget{sga6.e14.para.6-3}\textbf{6.3.}
Note that in each of the cases considered above in which one succeeds in
defining a Gysin homomorphism, one in fact obtains a construction for
cohomology with coefficients in $\Zell(*)$, and not only in $\Qell(*)$.
This makes it possible, when $f$ is an \emph{immersion}, to pose the question
of the validity of the Riemann--Roch formula ``without denominators'', as
explained in no.~3. When $X$ and $Y$ are both smooth over a field $k$, the
commutativity of \hyperlink{sga6.e14.eq.6-4}{(6.4)} shows that it would suffice to prove the formula
without denominators in the Chow ring $A(Y)$; cf.~\hyperlink{sga6.e14.para.3-3}{(3.3)}.

\noindent\sgaStableTarget{sga6.e14.para.6-4}\textbf{6.4.}
When no regularity hypothesis is made on $Y$, and even if $X$ and $Y$ are of
finite type over $\C$, one no longer sees a plausible definition of a
cohomological Gysin homomorphism \hyperlink{sga6.e14.para.6-2}{(6.2)}, even for a regular immersion of
codimension $1$. Can one nevertheless find one with reasonable properties,
in particular making square \hyperlink{sga6.e14.para.6-3}{(6.3)} commutative? The question is evidently
connected with that of defining, for $X$ a complete intersection of
codimension $i$ in $Y$, ``its cohomology class''
\[
 \gamma(X)\in H^{2i}(Y,\Zell(i)),
\]
in such a way as to turn intersections (without excess components) into cup
products. This latter question has an affirmative solution at least modulo
torsion, i.e. in $H^{2*}(Y,\Qell(*))$, by defining
\[
 \gamma(X)=\frac{1}{(-1)^{i-1}(i-1)!}\,c_i(\mathcal O_X).
\]
But I do not know whether $c_i(\mathcal O_X)$ is divisible by $(i-1)!$ in
$H^{2i}(Y,\Zell(i))$. In particular, if $Y$ is, say, a projective scheme of
dimension $3$ over the field of complex numbers and $X$ is reduced to a
(singular) point of $Y$, is $c_3(\mathcal O_X)$ divisible by $2$ in
$H^6(Y,\mathbb Z_2(3))$, and in particular is the Chern class zero in
$H^6(Y,\Z/2\Z)$? Let us also mention, along the same lines, the following
question: let $Y$ be a projective (singular) scheme of dimension $4$, and
$X$ a closed subscheme of dimension $1$ which is globally an intersection
of three positive divisors $D_1,D_2,D_3$. Consider the product of their
Chern classes modulo $2$,
\[
 c_1(D_1)c_1(D_2)c_1(D_3)\in H^6(Y,\Z/2\Z);
\]

% Current-rescribe idx698 | printed volume page 685 | source-PDF page 688
is this element of $H^6$ independent of the particular choice of the $D_i$
whose intersection is $X$?

\sgaStableTarget{sga6.e14.sec.7}\section*{7. Chern classes of perfect complexes}

\noindent\sgaStableTarget{sga6.e14.para.7-1}\textbf{7.1.}
When one poses the question of developing a cohomological Riemann--Roch
formula, with or without denominators, outside the rather artificial
hypothesis that an ample invertible module exists on the schemes under
consideration, one is necessarily faced with the question of defining, for
a perfect complex of modules $L^\bullet$ on a scheme $X$, the Chern classes
\[
 c_i(L^\bullet)\in H^{2i}(X,\Zell(i)),
\]
generalizing the Chern classes of locally free modules, SGA 5 VII. The same
question arises, moreover, for general locally ringed topoi (compare [10]),
as well as for perfect complexes on paracompact spaces ringed by continuous
complex-valued functions, working then with ordinary integral cohomology.
For a compact space, it is known that the $K^\bullet(X)$ constructed with
perfect complexes is the same as the ``naive'' one constructed with locally
free modules (\hyperlink{sga6.e04.expose}{Exp.~IV}), and in that case one is trivially reduced to the
classical definition. But already for a locally compact space the problem
posed seems nontrivial, since the integral cohomology of such a space is not
always the inverse limit of the cohomologies of its compact subsets.

\noindent\sgaStableTarget{sga6.e14.para.7-2}\textbf{7.2.}
More precisely, when $Y$ is a closed subset of $X$ such that
$L^\bullet|_{X-Y}$ is zero (in the sense of derived categories), i.e. such
that the cohomology sheaves of $L^\bullet$ are supported in $Y$, one would
like to define Chern classes \emph{with support in $Y$}
\[
 c_i^Y(L^\bullet)\in H_Y^{2i}(X,\Zell(i))
 \quad\text{respectively}\quad H_Y^{2i}(X,\Z),
\]
which would recover the ``ordinary'' (sic) Chern classes of \hyperlink{sga6.e14.para.7-1}{7.1} by means of
the canonical homomorphism
\[
 H_Y^*(X,-)\longrightarrow H^*(X,-).
\]

% Current-rescribe idx699 | printed volume page 686 | source-PDF page 689
Even in the case where the latter reduce to the Chern classes of locally
free modules ($X$ a scheme with an ample invertible module, or a compact
space), the construction of these Chern classes ``with supports'' has not
been carried out and will probably only be possible together with that of
the general Chern classes postulated in 7.1.\footnote{Added in October 1968.
The problems raised here were solved in the framework of Hodge cohomology by
L. Illusie; cf.~[13].}

It is probable that Chern classes with supports will have a role to play in
the proof of the cohomological Riemann--Roch formula without denominators
(cf.~\hyperlink{sga6.e14.para.3-5}{3.5}). On the other hand, they should make it possible to give an
improved version of the Riemann--Roch formula (with or without denominators)
by introducing families of supports into that formula.

\sgaStableTarget{sga6.e14.sec.8}\section*{8. Chow ring of regular schemes}

The problem is to define rational equivalence of cycles and to study the
resulting ring of classes (in particular its functorial properties) for a
regular noetherian scheme $X$ admitting an ample invertible module, thereby
generalizing the known theory when $X$ is smooth over a field [4]. The key
technical result will again be a ``moving lemma'', ensuring that a cycle can
be varied within its class so as to meet a given cycle in the proper
dimension.

The problem under consideration is quite natural a priori, especially since
we have seen (\hyperlink{sga6.e14.para.4-5}{4.5} to \hyperlink{sga6.e14.para.4-8}{4.8}) that the substitutes
$\operatorname{Gr}^\bullet(X)$ and $\Gtop^\bullet(X)$ for the Chow ring
$A(X)$ are hardly satisfactory when one does not confine oneself to a theory
``modulo torsion''. When $X$ is obtained from a smooth projective algebraic
scheme $Y$ over a field $k$ by taking the local scheme at the vertex of a
% SGA-SRC-001456
projective cone $C$ and removing the origin, one finds (assuming that $A(X)$
is indeed the inductive limit of the Chow rings of $U-\{O\}$, where $U$
ranges over the open neighborhoods in $C$ of the vertex $O$) that $A(X)$ is
canonically isomorphic to the quotient ring
\[
 A(Y)/\xi A(Y),
\]
where $\xi\in A^1(Y)$ is induced by the chosen projective embedding of $Y$.
This shows that knowing the Chow ring $A(X)$ (which is of a nonclassical
type) is equivalent to knowing $A(Y)/\xi A(Y)$, which might be called the
``primitive part'' of the Chow ring $A(Y)$ (in a sense suggested by the
Lefschetz--Hodge cohomological theory of nonsingular projective algebraic
varieties). This example shows that the ``local Chow rings'' of local rings
with isolated

% Current-rescribe idx700 | printed volume page 687 | source-PDF page 690
singularities are geometric invariants that are by no means trivial and will
undoubtedly have a role to play in the ``arithmetic'' study of these local
rings.

Of course, once one has a reasonable notion of a Chow ring for a regular
noetherian scheme admitting an ample invertible module, one should
generalize to this case the relations established or conjectured in no.~4
between this ring and $\Gtop^\bullet(X)$, and formulate in terms of this ring
the question of the ``Riemann--Roch formula without denominators'' stated in
no.~3.

\sgaStableTarget{sga6.e14.bibliography-of-expose-xiv}\subsection*{Bibliography of Expos\'e XIV}

\begin{enumerate}[label={[\arabic*]},leftmargin=3.2em]
\item M. F. Atiyah and F. Hirzebruch,
  \emph{The Riemann--Roch theorem for analytic embeddings},
  Topology, vol.~1, pp.~151--166 (1962).

\item A. Borel and J.-P. Serre,
  \emph{Le th\'eor\`eme de Riemann--Roch},
  Bull. Soc. Math. France, vol.~86, pp.~97--136 (1958).

\item R. Bott and H. Samelson,
  Amer. Journ. of Math., vol.~80, p.~1004 (1958).

\item C. Chevalley,
  \emph{Les classes d'\'equivalence rationelle}, I and II,
  S\'eminaire \`a l'ENS 1958, Expos\'es 2 and 3.

\item C. Chevalley,
  \emph{Classification des groupes de Lie alg\'ebriques},
  S\'eminaire \`a l'ENS 1956/58.

\item A. Dold and D. Puppe,
  \emph{Homologie nicht-additiver Funktoren. Anwendungen},
  Annales de l'Inst. Fourier, t.~XI, pp.~201--312 (1961).

\item H. Grauert,
  \emph{Ein Theorem der analytischen Garbentheorie},
  Pub. Math. no.~5, pp.~5--64 (1960).

\item A. Grothendieck,
  \emph{Classes de faisceaux et th\'eor\`eme de Riemann--Roch},
  mimeographed notes prepared by the Institute for Advanced Study, 1957
  (the work cited as [RRR] in the present Expos\'e and reproduced in our
  seminar as an appendix to Chapter~0).

\item A. Grothendieck,
  \emph{La th\'eorie des classes de Chern},
  Bull. Soc. Math. France, vol.~86, pp.~137--154 (1958).

% Current-rescribe idx701 | printed volume page 688 | source-PDF page 691
\item A. Grothendieck,
  \emph{Classes de Chern et repr\'esentations lin\'eaires des groupes
  discrets}, in: \emph{Dix expos\'es sur la cohomologie des sch\'emas},
  Masson--North Holland Pub. Cie.

\item H. Hironaka,
  \emph{Resolution of singularities of an algebraic variety over a field of
  characteristic zero}, Ann. of Math. 79, pp.~109--326 (1964).

\item L. Hodgkin, thesis (forthcoming?).

\item L. Illusie, work in preparation.

\item J.-P. Serre, \emph{Alg\`ebre Locale, Multiplicit\'es},
  Lecture Notes in Math. no.~11 (1965), Springer.

\item A. Borel,
  \emph{Sous-groupes commutatifs et torsion des groupes de Lie compacts
  connexes}, Tohoku Math. Journal, vol.~13, pp.~216--240 (1961).

\item[{[16]}(*)] P. Deligne, Expos\'e XVIII in
  \emph{Th\'eorie des topos et cohomologie \'etale des sch\'emas (SGA 6)},
  by M. Artin, A. Grothendieck, and J. L. Verdier.

% SGA-SRC-001458
\item[{[17]}(*)] Yu. I. Manin,
  \emph{Correspondances, motifs et transformations mono\"idales}
  (in Russian). Mat. Sborik, t.~77 (119), no.~4, 1968.
\end{enumerate}

\noindent{\footnotesize (*) Added in February 1969.}

% Current-rescribe idx702 | printed volume page 689 | source-PDF page 692
\sgaStableTarget{sga6.e14.note-for-page-1}\subsection*{Note for page 1}

The theory of the operations $\Lambda^i$ in $D^{-}(X)$, and of the
operations $\lambda^i$ in $K^\bullet(X)$, was written down by P. Deligne in
1968. He proves that $K^\bullet(X)$ thereby becomes a pre-$\lambda$-ring and,
when $\mathcal O_X$ is a sheaf of algebras over a field, that
$K^\bullet(X)$ is even a $\lambda$-ring. It is plausible that a refinement
of Deligne's method would make it possible to dispense with the restrictive
hypothesis on $\mathcal O_X$.

\sgaStableTarget{sga6.e14.notes-for-page-11}\subsection*{Notes for page 11}

\noindent(*) (April 1969) D. Mumford has pointed out to us a proof of
\hyperlink{sga6.e14.para.4-6}{(4.6)}\SGAArchiveOnly{\footnote{\textbf{Source note
[SGA6-XIV-IDX702-XREF96-VS-XREF46-SRCDEF-001].} The printed French note
reads ``(9.6)''; the contextual reference is formula \hyperlink{sga6.e14.eq.4-6}{(4.6)}.}}
in the \emph{Chow ring} (without disregarding torsion), which he obtained in
a seminar at Harvard in 1959! The reader will find the proof of this
beautiful result in SGA 5 VII 9. This result obviously contains the
corresponding cohomological result mentioned in note (**). The cohomological
result, however, is proved by J.-P. Jouanolou under more general conditions
% SGA-SRC-001459
(without the quasi-projectivity hypothesis, and in the relative smooth case over
an arbitrary base).

\medskip

\noindent(**) (January 1969) J.-P. Jouanolou has just proved these formulas
for every $y$, without disregarding torsion; cf. SGA 5 VII \S~4.9. Deligne
% SGA-SRC-001460
and Manin [17] had previously observed independently, and by (easy)
arguments of a rather different nature, that it was enough to establish the
formulas for ``algebraic'' $y$, and even for $y=1$ in [16], in order to
conclude the same relation in the general case.

% END INCLUDED SGA6_English_current_buildable_source_20260823/fragments/SGA6_XIV_idx693_702_English_SourceChecked_body.tex

\SGAArchiveOnly{\input{fragments/SGA6_XIV_idx685_702_Editorial_Source_Notes.tex}}
% BEGIN INCLUDED SGA6_English_current_buildable_source_20260823/SGA6_BackMatter_sourcePDF693_702_English_SourceChecked_fragment.tex
% Unindexed SGA 6 back matter after current-rescribe idx702.
% The current-rescribe has no idx numbers for these ten source pages.
% Printed page 690 is absent from the 702-page scan; the indexes begin at 691.

\newpage
\sgaStableTarget{sga6.back.terminological-index}\section*{Terminological index}
\begingroup\small
\begin{longtable}{@{}>{\raggedright\arraybackslash}p{0.64\linewidth}>{\raggedright\arraybackslash}p{0.14\linewidth}>{\raggedright\arraybackslash}p{0.12\linewidth}@{}}
\toprule
 & Exposé & §\\
\midrule
\endhead
% Unindexed back matter | printed page 691 | source-PDF page 693
additive (map) & I & \hyperlink{sga6.e01.def.6-1}{6.1}\\
additive (map) & IV & \hyperlink{sga6.e04.para.1-5}{1.5}\\
augmented \(K\)-\(\lambda\)-algebra & V & \hyperlink{sga6.e05.def.3-9}{3.9}\\
algebraically equivalent to zero (in \(K^\bullet\)) & X App. & \hyperlink{sga6.e10.para.7-11}{7.11}\\
algebraically equivalent to zero (invertible sheaf) & XIII & \hyperlink{sga6.e13.para.4-4}{4.4}\\
of \(\mathcal C_0\)-perfect amplitude contained in \([a,b]\) (complex) & I & \hyperlink{sga6.e01.def.4-7}{4.7}\\
of finite \(\mathcal C_0\)-perfect amplitude (complex) & I & \hyperlink{sga6.e01.def.4-7}{4.7}\\
of perfect amplitude contained in \([a,b]\) (complex) & I & \hyperlink{sga6.e01.ex.4-8}{4.8}\\
of finite perfect amplitude (complex) & I & \hyperlink{sga6.e01.ex.4-8}{4.8}\\
of flat amplitude contained in \([a,b]\) (complex) & I & \hyperlink{sga6.e01.def.5-2}{5.2}\\
of finite flat amplitude (complex) & I & \hyperlink{sga6.e01.def.5-2}{5.2}\\
of perfect amplitude relative to \(f\) (or \(Y\)) contained in \([a,b]\) (complex) & III & \hyperlink{sga6.e03.def.4-1}{4.1}\\
of perfect amplitude contained in \([-n,0]\) (morphism) & III & \hyperlink{sga6.e03.def.4-1}{4.1}\\
of finite perfect amplitude relative to \(f\) (complex) & III & \hyperlink{sga6.e03.def.4-1}{4.1}\\
of flat amplitude relative to \(f\) (or \(Y\)) contained in \([a,b]\) (complex) & III & \hyperlink{sga6.e03.def.3-1}{3.1}\\
of flat amplitude contained in \([-n,0]\) (morphism) & III & \hyperlink{sga6.e03.def.3-1}{3.1}\\
of finite flat amplitude relative to \(f\) (complex) & III & \hyperlink{sga6.e03.def.3-1}{3.1}\\
projective amplitude & I & \hyperlink{sga6.e01.rem.4-12-1}{4.12.1}\\
\(\lambda\)-ring & V & \hyperlink{sga6.e05.def.2-4}{2.4}\\
\(\lambda\)-ring generated by generators and relations & V & \hyperlink{sga6.e05.def.4-6}{4.6}\\
free \(\lambda\)-ring generated by a family of generators & V & \hyperlink{sga6.e05.def.4-4}{4.4}\\
Chern ring & V & \hyperlink{sga6.e05.def.6-2}{6.2}\\
polynomial map & X & \hyperlink{sga6.e10.def.2-1-1}{2.1.1}\\
binomial (ring) & V & \hyperlink{sga6.e05.def.2-7}{2.7}\\
% Unindexed back matter | printed page 692 | source-PDF page 694
Chern character & V & \hyperlink{sga6.e05.def.6-4}{6.4}\\
Todd character & V & \hyperlink{sga6.e05.para.1-16}{1.16}\\
abelian \(S\)-category & I & \hyperlink{sga6.e01.para.1-1-1}{1.1.1}\\
flat abelian \(S\)-category & I & \hyperlink{sga6.e01.para.1-1-1}{1.1.1}\\
additive \(S\)-category & I & \hyperlink{sga6.e01.para.1-1-1}{1.1.1}\\
triangulated \(S\)-category & I & \hyperlink{sga6.e01.para.1-1-1}{1.1.1}\\
Chern class & V & \hyperlink{sga6.e05.para.6-7}{6.7}\\
completed Chern class & V & \hyperlink{sga6.e05.para.6-7}{6.7}\\
total Chern class & V & \hyperlink{sga6.e05.para.6-7}{6.7}\\
codimension (of a regular immersion) & VII & \hyperlink{sga6.e07.def.1-6}{1.6}\\
coherent (complex) & I & \hyperlink{sga6.e01.def.3-1}{3.1}\\
coherer & II & \hyperlink{sga6.e02.sec.3.2}{3.2}\\
relative cotangent complex & VIII & \hyperlink{sga6.e08.def.2-3}{2.3}\\
elliptic complex & II App. II & \hyperlink{sga6.e02.para.1-2-5}{1.2.5}\\
cup-product & I & \hyperlink{sga6.e01.para.8-2}{8.2}\\
special positive \(k\)-cycle & XIII & \hyperlink{sga6.e13.def.6-9}{6.9}\\
degree (of a polynomial map) & X & \hyperlink{sga6.e10.def.2-1-1}{2.1.1}\\
degree (of an invertible sheaf) & XIII & \hyperlink{sga6.e13.para.2-2}{2.2}\\
degree (of a morphism of schemes) & X & \hyperlink{sga6.e10.sec.4.4}{4.4}\\
dense (subset of an abelian or triangulated category) & IV & \hyperlink{sga6.e04.def.1-7}{1.7}\\
perfect dimension & I & \hyperlink{sga6.e01.def.4-14-1}{4.14.1}\\
virtual relative dimension & VIII & \hyperlink{sga6.e08.def.1-9}{1.9}\\
divisorial & II & \hyperlink{sga6.e02.def.2-2-5}{2.2.5}\\
dual & I & \hyperlink{sga6.e01.def.7-3}{7.3}\\
perfect duality & I & \hyperlink{sga6.e01.cor.7-8}{7.8}\\
numerical equivalence (in \(\operatorname{Gr}^\bullet(X)\) and \(\operatorname{Gr}_\bullet(X)\)) & X & \hyperlink{sga6.e10.sec.4.5}{4.5}\\
\(\tau\)-equivalence & X & \hyperlink{sga6.e10.def.2-4-1}{2.4.1}\\
\(\tau\)-equivalence & XIII & \hyperlink{sga6.e13.para.4-4}{4.4}\\
complex pseudo-analytic space & II & \hyperlink{sga6.e02.def.2-4-7}{2.4.7}\\
complex Stein pseudo-analytic space & II & \hyperlink{sga6.e02.def.2-4-7}{2.4.7}\\
exact (subset of an abelian or triangulated category) & IV & \hyperlink{sga6.e04.def.1-7}{1.7}\\
\((b)\)-sheaf & XIII & \hyperlink{sga6.e13.def.1-5}{1.5}\\
conormal sheaf & VII & \hyperlink{sga6.e07.def.1-6}{1.6}\\
relative cotangent sheaf & II App. II & \hyperlink{sga6.e02.para.1-2-1}{1.2.1}\\
% Unindexed back matter | printed page 693 | source-PDF page 695
relative tangent sheaf & II App. II & \hyperlink{sga6.e02.para.1-2-1}{1.2.1}\\
\((b)\)-family & XIII & \hyperlink{sga6.e13.para.1-12}{1.12}\\
ample family & II & \hyperlink{sga6.e02.def.2-2-4}{2.2.4}\\
family of classes of coherent sheaves & XIII & \hyperlink{sga6.e13.para.1-12}{1.12}\\
limited family of classes of coherent sheaves & XIII & \hyperlink{sga6.e13.para.1-12}{1.12}\\
family of classes of coherent sheaves limited by a coherent sheaf & XIII & \hyperlink{sga6.e13.para.1-12}{1.12}\\
relative cotangent bundle & II App. II & \hyperlink{sga6.e02.para.1-2-1}{1.2.1}\\
relative tangent bundle & II App. II & \hyperlink{sga6.e02.para.1-2-1}{1.2.1}\\
faithfully flat (morphism of ringed topoi) & I & \hyperlink{sga6.e01.exercise.2-16-2}{2.16.2}\\
filtration of an augmented \(\lambda\)-algebra & V & \hyperlink{sga6.e05.para.3-10}{3.10}\\
additive \(S\)-functor & I & \hyperlink{sga6.e01.para.1-1-3}{1.1.3}\\
exact \(S\)-functor & I & \hyperlink{sga6.e01.para.1-1-3}{1.1.3}\\
projection formula & IV & \hyperlink{sga6.e04.para.2-11}{2.11}\\
Grothendieck group & IV & \hyperlink{sga6.e04.def.2-2}{2.2}\\
\(\lambda\)-homomorphism & V & \hyperlink{sga6.e05.def.2-1}{2.1}\\
specialization homomorphism & X App. & \hyperlink{sga6.e10.prop.7-3}{7.3}\\
specialization homomorphism & X App. & \hyperlink{sga6.e10.sec.7.16}{7.16}\\
\(\lambda\)-ideal & V & \hyperlink{sga6.e05.def.4-5}{4.5}\\
analytic index & II App. II & \hyperlink{sga6.e02.para.3-2-2}{3.2.2}\\
complete-intersection (morphism) & VIII & \hyperlink{sga6.e08.def.1-1}{1.1}\\
\(n\)-isomorphism & I & \hyperlink{sga6.e01.def.1-4-3}{1.4.3}\\
homogeneous free (graded module) & VI & \hyperlink{sga6.e06.eq.1-2-2}{1.2.2}\\
locally stable under kernels of epimorphisms & I & \hyperlink{sga6.e01.sec.1.2}{1.2}\\
locally quasi-liftable & I & \hyperlink{sga6.e01.sec.1.2}{1.2}\\
locally quasi-stable under kernels of epimorphisms & I & \hyperlink{sga6.e01.sec.1.2}{1.2}\\
locally liftable & I & \hyperlink{sga6.e01.sec.1.2}{1.2}\\
Chern number & X & \hyperlink{sga6.e10.ex.4-2}{4.2}\\
intersection number & X & \hyperlink{sga6.e10.sec.4.1}{4.1}\\
intersection number (for a family of invertible sheaves) & X & \hyperlink{sga6.e10.sec.4.4}{4.4}\\
numerically equivalent to zero (invertible sheaf) & XIII & \hyperlink{sga6.e13.para.4-4}{4.4}\\
universal differential operator of order \(n\) & II App. II & \hyperlink{sga6.e02.para.1-2-1}{1.2.1}\\
% Unindexed back matter | printed page 694 | source-PDF page 696
Adams operations & V & \hyperlink{sga6.e05.para.7-1}{7.1}\\
\(\mathcal C_0\)-perfect (complex) & I & \hyperlink{sga6.e01.def.4-7}{4.7}\\
perfect (complex) & I & \hyperlink{sga6.e01.ex.4-8}{4.8}\\
perfect (morphism) & III & \hyperlink{sga6.e03.def.4-1}{4.1}\\
perfect relative to \(f\) (complex) & III & \hyperlink{sga6.e03.def.4-1}{4.1}\\
precoherent & I & \hyperlink{sga6.e01.def.3-1}{3.1}\\
of finite \(n\)-\(\mathcal C_0\)-presentation & I & \hyperlink{sga6.e01.def.2-8}{2.8}\\
of finite presentation & I & \hyperlink{sga6.e01.def.2-8}{2.8}\\
of finite \(\infty\)-\(\mathcal C_0\)-presentation & I & \hyperlink{sga6.e01.def.2-8}{2.8}\\
\(n\)-\(\mathcal C_0\)-pseudo-coherent (complex) & I & \hyperlink{sga6.e01.prop.2-2}{2.2}\\
\(\mathcal C_0\)-pseudo-coherent (complex) & I & \hyperlink{sga6.e01.prop.2-2}{2.2}\\
\(n\)-pseudo-coherent (complex) & I & \hyperlink{sga6.e01.ex.2-15}{2.15}\\
\(n\)-pseudo-coherent (morphism) & III & \hyperlink{sga6.e03.def.1-2}{1.2}\\
pseudo-coherent (complex) & I & \hyperlink{sga6.e01.ex.2-15}{2.15}\\
pseudo-coherent (morphism) & III & \hyperlink{sga6.e03.def.1-2}{1.2}\\
\(n\)-pseudo-coherent relative to \(f\) (complex) & III & \hyperlink{sga6.e03.def.1-2}{1.2}\\
pseudo-coherent relative to \(f\) (complex) & III & \hyperlink{sga6.e03.def.1-2}{1.2}\\
\((b)\)-polynomial & XIII & \hyperlink{sga6.e13.def.1-9}{1.9}\\
Snapper polynomial & X & \hyperlink{sga6.e10.prop.2-3-1}{2.3.1}\\
universal polynomial & V & \hyperlink{sga6.e05.para.1-15}{1.15}\\
pre-\(\lambda\)-ring & V & \hyperlink{sga6.e05.def.2-1}{2.1}\\
\(n\)-quasi-isomorphism & I & \hyperlink{sga6.e01.def.1-4-3}{1.4.3}\\
rank (of a perfect complex) & I & \hyperlink{sga6.e01.def.6-11}{6.11}\\
\(m\)-regular (sheaf) & XIII & \hyperlink{sga6.e13.def.1-1}{1.1}\\
regular (homomorphism) & VII & \hyperlink{sga6.e07.def.1-1}{1.1}\\
regular (ideal) & VII & \hyperlink{sga6.e07.def.1-4}{1.4}\\
regular (immersion) & VII & \hyperlink{sga6.e07.def.1-4}{1.4}\\
triangulated sub-\(S\)-category & I & \hyperlink{sga6.e01.para.1-1-4}{1.1.4}\\
special (\(\lambda\)-ring) & V & \hyperlink{sga6.e05.def.2-4}{2.4}\\
special (morphism) & III & \hyperlink{sga6.e03.lem.1-1-4}{1.1.4}\\
Stein (locally ringed topos) & II & \hyperlink{sga6.e02.def.2-4-1}{2.4.1}\\
strictly \(\mathcal C_0\)-perfect & I & \hyperlink{sga6.e01.def.2-1}{2.1}\\
strictly \(\mathcal C_0\)-\(n\)-pseudo-coherent & I & \hyperlink{sga6.e01.def.2-1}{2.1}\\
strictly \(\mathcal C_0\)-pseudo-coherent & I & \hyperlink{sga6.e01.def.2-1}{2.1}\\
minimum system of generators & VII & \hyperlink{sga6.e07.def.1-4}{1.4.1}\\
% Unindexed back matter | printed page 695 | source-PDF page 697
of finite Tor-dimension & I & \hyperlink{sga6.e01.def.5-2}{5.2}\\
Tor-independence & III & \hyperlink{sga6.e03.def.1-5}{1.5}\\
trace & I & \hyperlink{sga6.e01.para.8-1}{8.1}\\
of finite \(\mathcal C_0\)-type & I & \hyperlink{sga6.e01.sec.1.2}{1.2}\\
mixed variety of class \(C^r\) & II App. II & \hyperlink{sga6.e02.sec.1.1}{1.1}\\
\bottomrule
\end{longtable}
\endgroup

\newpage
\sgaStableTarget{sga6.back.index-of-notations}\section*{Index of notations}
\begingroup\small
\begin{longtable}{@{}>{\raggedright\arraybackslash}p{0.50\linewidth}>{\raggedright\arraybackslash}p{0.17\linewidth}>{\raggedright\arraybackslash}p{0.15\linewidth}@{}}
\toprule
 & Exposé & §\\
\midrule
\endhead
% Unindexed back matter | printed page 696 | source-PDF page 698
\(A_N\) & V & \hyperlink{sga6.e05.para.5-5}{5.5}\\
\(\ch\) & V & \hyperlink{sga6.e05.para.6-3}{6.3}\\
\(\Ch(A),\ \Ch_K(A)\) & V & \hyperlink{sga6.e05.def.6-2}{6.2}\\
\((\Ch(A))_n\) & V & \hyperlink{sga6.e05.prop.6-5}{6.5}\\
\(\cl_{\mathcal C}\) & IV & \hyperlink{sga6.e04.para.1-1}{1.1}\\
\(\Coh(S)\) & II & \hyperlink{sga6.e02.notation.2-2-1}{2.2.1}\\
\(\mathcal C_0\text{-}\parf\text{-}\amp(F)\) & I & \hyperlink{sga6.e01.def.4-7}{4.7}\\
\(C^p_{\num}(X)\) & XIII & \hyperlink{sga6.e13.prop.5-2}{5.2}\\
\(C^r(V),\ C^\infty(V)\) & II App. II & \hyperlink{sga6.e02.sec.1.1}{1.1}\\
\(c(x),\ c^i(x),\ \widetilde c(x)\) & V & \hyperlink{sga6.e05.para.6-7}{6.7}\\
\(\Dcat(A,S),\ \Dcat(A_{\underline S})\) & I & \hyperlink{sga6.e01.ex.2-15}{2.15}\\
\(\Dcat(A_S)_{\coh},\ \Dcat(A_S)_{n\text{-}\coh}\) & I & \hyperlink{sga6.e01.ex.2-15}{2.15}\\
\(\Dcat(A_S)_{\parf}\) & I & \hyperlink{sga6.e01.notation.4-9}{4.9}\\
\(\Dcat^{-}(A_S)_{\torf},\ \Dcat^{-}(A_X)_{\torf}\) & I & \hyperlink{sga6.e01.para.5-3}{5.3}\\
\(\Dcat(\mathcal C)_{\mathcal C_0\text{-}\coh},\ \Dcat(\mathcal C)_{\mathcal C_0\text{-}n\text{-}\coh}\) & I & \hyperlink{sga6.e01.schol.2-6}{2.6}\\
\(\Dcat(\mathcal C)_{\mathcal C_0\text{-}\parf}\) & I & \hyperlink{sga6.e01.notation.4-9}{4.9}\\
\(\Dcat(f)_{\coh}\) & III & \hyperlink{sga6.e03.def.1-2}{1.2}\\
\(\Dcat(f)_{n\text{-}\coh}\) & III & \hyperlink{sga6.e03.def.1-2}{1.2}\\
\(\Dcat(f)_{\parf},\ \Dcat(f)_{\underline{\mathrm{parf}}}\) & III & \hyperlink{sga6.e03.notation.4-2}{4.2}\\
\(\Dcat(f)_{\torf}\) & III & \hyperlink{sga6.e03.notation.3-2}{3.2}\\
\(\Diff^n(E,F)\) & II App. II & \hyperlink{sga6.e02.para.1-2-2}{1.2.2}\\
\(d(L)\) & XIII & \hyperlink{sga6.e13.para.2-2}{2.2}\\
\(d^n_{X/S}\) & II App. II & \hyperlink{sga6.e02.para.1-2-1}{1.2.1}\\
\(d^n_{X/S}(E)\) & II App. II & \hyperlink{sga6.e02.para.1-2-2}{1.2.2}\\
% Unindexed back matter | printed page 697 | source-PDF page 699
\(\Dcat^{-}(\Qcoh(S))_{\coh},\ \Dcat(\Qcoh(S))_{\parf}\) & II & \hyperlink{sga6.e02.notation.2-2-1}{2.2.1}\\
\(\Dcat(S),\ \Dcat(\underline S),\ \Dcat(\underline S,A)\) & I & \hyperlink{sga6.e01.ex.2-15}{2.15}\\
\(\Dcat(X)_{f\text{-}\parf},\ \Dcat(X)_{\underline{f\text{-}\mathrm{parf}}}\) & III & \hyperlink{sga6.e03.notation.4-2}{4.2}\\
\(\Dcat(X)_{f\text{-}\torf}\) & III & \hyperlink{sga6.e03.notation.3-2}{3.2}\\
\(\Dcat(X)_{\qcoh}\) & III & \hyperlink{sga6.e03.lem.1-7}{1.7}\\
\(\Dcat(X)_{Y\text{-}\coh},\ \Dcat(X)_{Y\text{-}n\text{-}\coh}\) & III & \hyperlink{sga6.e03.def.1-2}{1.2}\\
\(\Dcat(X)_{Y\text{-}\parf},\ \Dcat(X)_{\underline{Y\text{-}\mathrm{parf}}}\) & III & \hyperlink{sga6.e03.notation.4-2}{4.2}\\
\(\Dcat(X)_{Y\text{-}\torf}\) & III & \hyperlink{sga6.e03.notation.3-2}{3.2}\\
\(\Dcat(Y)_{n\text{-}\coh}\) & III & \hyperlink{sga6.e03.thm.1-1}{1.1}\\
\(E^\vee\) & I & \hyperlink{sga6.e01.def.7-3}{7.3}\\
\(*\) & V & \hyperlink{sga6.e05.para.6-1}{6.1}\\
\(E|\) & X & \hyperlink{sga6.e10.prop.4-3-1}{4.3.1}\\
\(f^*\) (between the \(K^\bullet\)) & IV & \hyperlink{sga6.e04.para.2-7}{2.7}\\
\(f^*\) (between the \(K_\bullet\)) & IV & \hyperlink{sga6.e04.para.2-12}{2.12}\\
\(f_*\) (between the \(K_\bullet\)) & IV & \hyperlink{sga6.e04.para.2-11}{2.11}\\
\(f_*\) (between the \(K^\bullet\)) & IV & \hyperlink{sga6.e04.para.2-12}{2.12}\\
\(f_*\) (projective-bundle case) & VI & \hyperlink{sga6.e06.sec.5.2}{5.2}\\
\(f_{\operatorname{gr}}\) & VIII & \hyperlink{sga6.e08.prop.3-2}{3.2}\\
\(\Fil_n(X)\) & X & \hyperlink{sga6.e10.def.1-1-1}{1.1.1}\\
\(F^{\prime >n},\ F'_{\leq n}\) & I & \hyperlink{sga6.e01.prop.2-12}{2.12}\\
\(\varphi^a\) & V & \hyperlink{sga6.e05.para.2-9}{2.9}\\
\(\circ,\ \widehat G_K\) & V & \hyperlink{sga6.e05.para.1-1}{1.1}\\
\(\circ,\ \widehat G_{\nabla}\) & V & \hyperlink{sga6.e05.para.3-7}{3.7}\\
\(\widehat G_a\) & V & \hyperlink{sga6.e05.para.1-3}{1.3.3}\\
\(\widehat G_g\) & V & \hyperlink{sga6.e05.para.1-3}{1.3.2}\\
\(\operatorname{Gr}^\bullet(X),\ \operatorname{Gr}_\bullet(X)\) & X & \hyperlink{sga6.e10.sec.3.1}{3.1}\\
\(\operatorname{Gr}^\bullet_{\alg},\ \operatorname{Gr}^{\alg}_\bullet\) & X App. & \hyperlink{sga6.e10.para.7-11}{7.11}\\
% Unindexed back matter | printed page 698 | source-PDF page 700
\(\operatorname{Gr}^\bullet_{\num},\ \operatorname{Gr}^{\num}_\bullet\) & X App. & \hyperlink{sga6.e10.para.7-9-4}{7.9.4}\\
\(\gamma^n,\ \gamma_t\) & V & \hyperlink{sga6.e05.prop.3-2}{3.2}\\
\(\gamma^p(N,x)\) & V & \hyperlink{sga6.e05.para.5-5}{5.5}\\
\(\widehat G^+,\ \widehat G^u\) & V & \hyperlink{sga6.e05.para.2-3}{2.3}\\
\(\underline{\Hom}_{\Lambda}(E,F)\) & I & \hyperlink{sga6.e01.notation.4-5}{4.5}\\
\(k(A)\) & IV & \hyperlink{sga6.e04.para.1-5}{1.5}\\
\(K^\bullet_{\alg},\ K^{\alg}_\bullet\) & X App. & \hyperlink{sga6.e10.para.7-11}{7.11}\\
\(K^\bullet(A_X),\ K_\bullet(A_X)\) & IV & \hyperlink{sga6.e04.def.2-2}{2.2}\\
\(K^\bullet(A_X)_{\naif},\ K_\bullet(A_X)_{\naif}\) & IV & \hyperlink{sga6.e04.def.2-2}{2.2}\\
\(k(C)\) & I & \hyperlink{sga6.e01.para.6-3}{6.3}\\
\(K^\bullet(f),\ K_\bullet(f)\) & IV & \hyperlink{sga6.e04.para.3-3-1}{3.3.1}\\
\(K^\bullet_{\num},\ K^{\num}_\bullet\) & X App. & \hyperlink{sga6.e10.para.7-9-4}{7.9.4}\\
\(K^\bullet(X),\ K_\bullet(X)\) & IV & \hyperlink{sga6.e04.def.2-2}{2.2}\\
\(K^\bullet(X)_{\naif},\ K_\bullet(X)_{\naif}\) & IV & \hyperlink{sga6.e04.def.2-2}{2.2}\\
\(K^\bullet(X/Y),\ K_\bullet(X/Y)\) & IV & \hyperlink{sga6.e04.para.3-3-1}{3.3.1}\\
\(\Lib(U)\) & II & \hyperlink{sga6.e02.sec.2.0}{2.0}\\
\(\Loclib(U)\) & II & \hyperlink{sga6.e02.sec.2.0}{2.0}\\
\(L^{X/Y}\) & VIII & \hyperlink{sga6.e08.para.2-1}{2.1}\\
\(\lambda^i\) & V & \hyperlink{sga6.e05.def.2-1}{2.1}\\
\(\lambda_{-1}(N)\) & V & \hyperlink{sga6.e05.para.5-1}{5.1}\\
\(\lambda_{-1}(N)\) & VII & \hyperlink{sga6.e07.para.2-6}{2.7}\\
\(\lambda^p(N,x)\) & V & \hyperlink{sga6.e05.def.5-3}{5.3}\\
\(\lambda_t\) & V & \hyperlink{sga6.e05.def.2-1}{2.1}\\
\(\Mod(U)\) & II & \hyperlink{sga6.e02.sec.2.0}{2.0}\\
% Unindexed back matter | printed page 699 | source-PDF page 701
\(\parfamp_f(E),\ \parfamp_Y(E)\) & III & \hyperlink{sga6.e03.def.4-1}{4.1}\\
\(\Parf(A_S),\ \Parf(\mathcal C,\mathcal C_0),\ \Parf(S)\) & I & \hyperlink{sga6.e01.notation.4-9}{4.9}\\
\(\Pic^N(X)\) & X & \hyperlink{sga6.e10.prop.4-5-2}{4.5.2}\\
\(\Pic^\circ(X)\) & X & \hyperlink{sga6.e10.def.2-4-1}{2.4.1}\\
\(\underline{\Pic}^{\circ}_{X/k}\) & X & \hyperlink{sga6.e10.sec.2.4}{2.4}\\
\(\Pic^\tau(X)\) & X & \hyperlink{sga6.e10.def.2-4-1}{2.4.1}\\
\(\Pic(X)\) & X & \hyperlink{sga6.e10.sec.2.4}{2.4}\\
\(\underline{\Pic}_{X/k}\) & X & \hyperlink{sga6.e10.sec.2.4}{2.4}\\
\(P^n_{X/S}\) & II App. II & \hyperlink{sga6.e02.para.1-2-1}{1.2.1}\\
\(P^n_{X/S}(E)\) & II App. II & \hyperlink{sga6.e02.para.1-2-2}{1.2.2}\\
\(\Proj\) & VI & \hyperlink{sga6.e06.para.2-3-2}{2.3.2}\\
\(\psi_k(x)\) & V & \hyperlink{sga6.e05.para.7-1}{7.1}\\
\(\Qcoh(S)\) & II & \hyperlink{sga6.e02.notation.2-2-1}{2.2.1}\\
\(r^0_{\mathcal O_{U\times V}}\) & II App. II & \hyperlink{sga6.e02.sec.1.1}{1.1}\\
\(\rho_t\) & V & \hyperlink{sga6.e05.para.3-4}{3.4}\\
\(\circ\) & V & \hyperlink{sga6.e05.para.2-3}{2.3}\\
\(T_f\) & VIII & \hyperlink{sga6.e08.cor.2-5}{2.5}\\
\(\Todd\) & V & \hyperlink{sga6.e05.para.1-16}{1.16}\\
\(\toramp(F)\) & I & \hyperlink{sga6.e01.def.5-2}{5.2}\\
\(\toramp_f(E),\ \toramp_Y(E)\) & III & \hyperlink{sga6.e03.def.3-1}{3.1}\\
\(\Tr_{E|U}\) & I & \hyperlink{sga6.e01.para.8-1}{8.1}\\
\(T_{X/S}\) & II App. II & \hyperlink{sga6.e02.para.1-2-1}{1.2.1}\\
\(\tau\) & I & \hyperlink{sga6.e01.para.8-1}{8.1}\\
\(\nabla\) & V & \hyperlink{sga6.e05.para.3-5}{3.5}\\
% Unindexed back matter | printed page 700 | source-PDF page 702
\(1+A[[t]]^+\) & V & \hyperlink{sga6.e05.para.1-1}{1.1}\\
\(1+\widehat B^+\) & V & \hyperlink{sga6.e05.para.1-11}{1.11}\\
\(X(n)\) & X & \hyperlink{sga6.e10.prop.1-1-2}{1.1.2}\\
\(x_N\) & V & \hyperlink{sga6.e05.para.5-5}{5.5}\\
\(\langle x,y\rangle\) & X & \hyperlink{sga6.e10.sec.4.1}{4.1}\\
\(Z(x)\) & X & \hyperlink{sga6.e10.prop.1-1-2}{1.1.2}\\
\bottomrule
\end{longtable}
\endgroup

% END INCLUDED SGA6_English_current_buildable_source_20260823/SGA6_BackMatter_sourcePDF693_702_English_SourceChecked_fragment.tex


% END INCLUDED SGA6_English_current_buildable_source_20260823/SGA6_idx532_702_and_backmatter_English_Synchronized_body.tex


\end{document}

% SGA-CUMULATIVE-END 07_SGA6_EN_CURRENT_LAYERED_READER.tex

% SGA-CUMULATIVE-BEGIN 08_SGA7I_EN_EXPOSES_I_II_VI_IX_READER.tex
\documentclass[11pt,a4paper,leqno]{article}

\usepackage[T1]{fontenc}
\usepackage[utf8]{inputenc}
\usepackage[english,french]{babel}
\usepackage{amsmath,amssymb,amsfonts,mathrsfs}
\usepackage{mathtools}
\usepackage{tikz}
\usetikzlibrary{cd,arrows.meta}
\usepackage{geometry}
\geometry{margin=2.4cm}
\usepackage{enumitem}
\usepackage[normalem]{ulem}
\usepackage{hyperref}
\hypersetup{
  colorlinks=true,
  linkcolor=blue,
  citecolor=blue,
  urlcolor=blue,
  hypertexnames=false,
  pdftitle={SGA 7-I: Monodromy Groups in Algebraic Geometry - English Reader},
  pdfauthor={Interlanguage Mathematical Translation Project},
  pdfsubject={Source-first English translation from the complete French working transcription of LNM 288},
  pdfkeywords={SGA 7; monodromy; algebraic geometry; English translation}
}
\let\mathunderline\underline
\renewcommand{\underline}[1]{\ifmmode\mathunderline{#1}\else\uline{#1}\fi}
\emergencystretch=3em

\DeclareUnicodeCharacter{2192}{\ensuremath{\rightarrow}}
\DeclareUnicodeCharacter{2297}{\ensuremath{\otimes}}
\DeclareUnicodeCharacter{2295}{\ensuremath{\oplus}}
\DeclareUnicodeCharacter{2228}{\ensuremath{\vee}}
\DeclareUnicodeCharacter{223C}{\ensuremath{\sim}}
\DeclareUnicodeCharacter{2248}{\ensuremath{\approx}}
\DeclareUnicodeCharacter{00D7}{\ensuremath{\times}}
\DeclareUnicodeCharacter{03B1}{\ensuremath{\alpha}}
\DeclareUnicodeCharacter{03B2}{\ensuremath{\beta}}
\DeclareUnicodeCharacter{00A7}{\S}
\DeclareUnicodeCharacter{010C}{\v{C}}
\DeclareUnicodeCharacter{00C3}{\~{A}}
\DeclareUnicodeCharacter{02D9}{\.{}}

\providecommand{\Zl}{\mathbb{Z}_{\ell}}
\providecommand{\Ql}{\mathbb{Q}_{\ell}}
\providecommand{\Gal}{\operatorname{Gal}}
\providecommand{\Spec}{\operatorname{Spec}}
\providecommand{\Isom}{\operatorname{Isom}}
\providecommand{\Aut}{\operatorname{Aut}}
\providecommand{\Hom}{\operatorname{Hom}}
\providecommand{\Ext}{\operatorname{Ext}}
\providecommand{\Biext}{\operatorname{Biext}}
\providecommand{\Pic}{\operatorname{Pic}}
\providecommand{\Coker}{\operatorname{Coker}}
\providecommand{\Ker}{\operatorname{Ker}}
\providecommand{\Im}{\operatorname{Im}}
\providecommand{\Tor}{\operatorname{Tor}}
\providecommand{\prof}{\operatorname{prof}}
\providecommand{\codim}{\operatorname{codim}}
\providecommand{\pr}{\operatorname{pr}}
\providecommand{\id}{\mathrm{id}}
\providecommand{\sym}{\mathrm{sym}}
\providecommand{\Corr}{\operatorname{Corr}}
\providecommand{\Proj}{\operatorname{Proj}}
\providecommand{\Comp}{\operatorname{Comp}}
\providecommand{\cC}{\mathcal{C}}
\providecommand{\cY}{\mathcal{Y}}
\providecommand{\cT}{\mathcal{T}}
\providecommand{\cPsi}{\Psi}
\providecommand{\Dcat}{D}
\providecommand{\La}{\Lambda}
\providecommand{\Ghat}{\widehat{G}}
\providecommand{\uHom}{\underline{\operatorname{Hom}}}
\providecommand{\uExt}{\underline{\operatorname{Ext}}}
\providecommand{\uTor}{\underline{\operatorname{Tor}}}
\providecommand{\cl}{\operatorname{c\ell}}
\providecommand{\e}{\mathrm{e}}

\begin{document}
\selectlanguage{english}
\begin{titlepage}
\centering
\vspace*{2.2in}
{\LARGE\textbf{SGA 7-I}}\\[0.6em]
{\Large Monodromy Groups in Algebraic Geometry}\\[1.5em]
{\large Complete English reader}\\[1.2em]
{\normalsize Exposés I, II, VI, VII, VIII, and IX}
\vfill
\end{titlepage}
\tableofcontents
\clearpage
\phantomsection\addcontentsline{toc}{section}{Exposé I}
% BEGIN INCLUDED SGA7I_English_current_buildable_source_20260823/components/01_expose_I_opening_through_0_0.tex
% SOURCE AUTHORITY:
% expose_I_body.tex lines 1--54; scan indices 12--13; source folios 1--2.
% Hard stop before source line 57, paragraph 0.1.

%%%%%%%%%%  scan idx 12  |  folio 1  %%%%%%%%%%

\SGARefAliasTarget{sga7.sga7i.ei.nroot}SGA 7

\underline{Expos\'e I}

\underline{Summary of the First Lectures by A. Grothendieck}

\underline{written up by P. Deligne}

\underline{Contents}

0. Preliminaries

1. Arithmetic proof of the monodromy theorem

2. Vanishing cycles

3. Geometric proof of the monodromy theorem

4. Vanishing criteria for the sheaves of vanishing cycles

5. Action of monodromy on the $\pi_1$

6. Appendix by P. Deligne: arithmetic proof of the stable reduction theorem

0. \underline{Preliminaries}

\SGARefTarget{sga7.sga7i.ei.n0-0}{0.0}. \underline{Terminology}. Recall the following definitions:

(\SGARefTarget{sga7.sga7i.ei.n0-0-1}{0.0.1}) \underline{trait}: the spectrum of a discrete valuation ring. The generic and closed points of a trait are often denoted by $\eta$ and $s$.

(\SGARefTarget{sga7.sga7i.ei.n0-0-2}{0.0.2}) \underline{strictly local}: for a ring, local henselian with separably closed residue field; for a scheme, the spectrum of such a ring.

(\SGARefTarget{sga7.sga7i.ei.n0-0-3}{0.0.3}) \underline{geometric point of $S$} (in this expos\'e): a morphism $\bar{x}:\operatorname{Spec}(\bar{k})\longrightarrow S$ with image $x\in S$, where $k(\bar{x})\underset{\mathrm{dfn}}{=}\bar{k}$ is a separable closure of $k(x)$. For a henselian trait $S=\operatorname{Spec}(V)$, a generic geometric point of $S$ (that is, one with image $\eta$) is often denoted by $\bar{\eta}$. The residue field of the valuation ring $\bar{V}$, the integral closure of $V$ in $k(\bar{\eta})$, is a separably closed algebraic extension of $k(s)$ and defines a geometric point $\bar{s}$ of $S$ centered at $s$.

%%%%%%%%%%  scan idx 13  |  folio 2  %%%%%%%%%%

(\SGARefTarget{sga7.sga7i.ei.n0-0-4}{0.0.4}) \underline{strict localization of $X$ at the geometric point $\bar{x}$}: SGA 4 VIII 4.4.

(\SGARefTarget{sga7.sga7i.ei.n0-0-5}{0.0.5}) \underline{essentially of finite type over $S$}: for a local (respectively, henselian local; respectively, strictly local) $S$-scheme, the spectrum of a local ring (respectively, the henselization of a local ring; respectively, a strict localization) of a scheme of finite type over $S$.

(\SGARefTarget{sga7.sga7i.ei.n0-0-6}{0.0.6}) $F\longmapsto F(n)$: Tate twist.

(\SGARefTarget{sga7.sga7i.ei.n0-0-7}{0.0.7}) An endomorphism $T$ of a finite-dimensional vector space is \underline{quasi-unipotent} if its eigenvalues are roots of unity, that is, if there exist $N\geqslant 1$ and $M\geqslant 1$ such that $(T^N-1)^M=0$.

% END INCLUDED SGA7I_English_current_buildable_source_20260823/components/01_expose_I_opening_through_0_0.tex

% BEGIN INCLUDED SGA7I_English_current_buildable_source_20260823/components/02_expose_I_0_1_through_1_4.tex
% SOURCE AUTHORITY:
% expose_I_body.tex lines 57--246; scan indices 13--18; source folios 2--7.
% Hard stop before source line 249, section 2.

\SGARefTarget{sga7.sga7i.ei.n0-1}{0.1}. \underline{Review of $\ell$-adic cohomology}

Let $X$ be a scheme of finite type over an algebraically closed field $k$ of characteristic exponent $p$, and let $\ell$ be a prime different from $p$.

The $\ell$-adic cohomology groups of $X$ were defined in SGA 4 and SGA 5.

a) $H^i(X,\mathbb{Z}/\ell^n)$ is the $i$-th cohomology group of the \`etale site $X_{\mathrm{et}}$ of $X$ with coefficients in $\mathbb{Z}/\ell^n$;

b) by definition,
\[
H^i(X,\mathbb{Z}_\ell)=\varprojlim_n H^i(X,\mathbb{Z}/\ell^n);
\]

c) by definition,
\[
H^i(X,\mathbb{Q}_\ell)=H^i(X,\mathbb{Z}_\ell)\otimes_{\mathbb{Z}_\ell}\mathbb{Q}_\ell.
\]

The $\ell$-adic cohomology groups of $X$ with proper supports are defined in the same way from the groups $H^i_c(X,\mathbb{Z}/\ell^n)$ of SGA 4 XVII.

In each of the following cases, one knows how to prove that these groups $H^i(X)$ have reasonable finiteness properties:

(\SGARefTarget{sga7.sga7i.ei.n0-1-1}{0.1.1}) for cohomology with proper supports;

(\SGARefTarget{sga7.sga7i.ei.n0-1-2}{0.1.2}) when $X$ is proper over $k$;

(\SGARefTarget{sga7.sga7i.ei.n0-1-3}{0.1.3}) when $X$ is smooth over $k$ (by (\SGARefLink{sga7.sga7i.ei.n0-1-1}{0.1.1}) and Poincar\'e duality, SGA 4 XVIII);

(\SGARefTarget{sga7.sga7i.ei.n0-1-4}{0.1.4}) when Hironaka's resolution of singularities is available for schemes of finite type over $k$ of dimension at most $\dim(X)$;

(\SGARefTarget{sga7.sga7i.ei.n0-1-5}{0.1.5}) when $i=0,1$.

%%%%%%%%%%  scan idx 14  |  folio 3  %%%%%%%%%%

In all these cases, the groups $H^i(X,\mathbb{Z}/\ell^n)$ are finite, the groups $H^i(X,\mathbb{Z}_\ell)$ are finitely generated, the groups $H^i(X,\mathbb{Q}_\ell)$ are finite-dimensional, and (except perhaps in case (\SGARefLink{sga7.sga7i.ei.n0-1-5}{0.1.5})) there are split short exact sequences

\begin{equation*}
0 \longrightarrow H^i(X,\mathbb{Z}_\ell)\otimes \mathbb{Z}/\ell^n
\longrightarrow H^i(X,\mathbb{Z}/\ell^n)
\longrightarrow \Tor_1(H^{i+1}(X,\mathbb{Z}_\ell),\mathbb{Z}/\ell^n)
\longrightarrow 0.
\end{equation*}
Moreover, in each case the proof of the finiteness theorem also proves a generic constructibility theorem: if $k$ is the field of fractions of an integral noetherian scheme $S$, and $X$ is the generic fibre of a morphism $f_1:X_1\longrightarrow S$ of finite type, then there exists a nonempty open subset $U\subset S$ such that over $U$ the sheaves $R^i f_*\mathbb{Z}/\ell^n$ are locally constant and their formation commutes with every base change.

\SGARefTarget{sga7.sga7i.ei.n0-2}{0.2}. We no longer suppose that $k$ is algebraically closed. Let $\bar{k}$ be an algebraic closure of $k$ (or a separable closure, which amounts to the same thing here). We consider the ``geometric'' cohomology groups $H^i(X_{\bar{k}})$, where $X_{\bar{k}}$ is obtained from $X$ by extension of scalars from $k$ to $\bar{k}$. By transport of structure, the Galois group $\Gal(\bar{k}/k)$ acts on these groups. In $\mathbb{Q}_\ell$-cohomology, in any of the cases (\SGARefLink{sga7.sga7i.ei.n0-1-1}{0.1.1})--(\SGARefLink{sga7.sga7i.ei.n0-1-5}{0.1.5}), this gives continuous representations of $\Gal(\bar{k}/k)$ in a linear group $\mathrm{GL}(n,\mathbb{Q}_\ell)$.

\SGARefTarget{sga7.sga7i.ei.n0-3}{0.3}. \underline{Review of traits} (\SGARefLink{sga7.sga7i.ei.n0-0-1}{0.0.1})

For the proofs, we refer to [4]. Let $S$ be a henselian trait, and let $\eta$, $s$, $V$, $\bar V$, $\bar\eta$, and $\bar s$ be as in (\SGARefLink{sga7.sga7i.ei.n0-0-1}{0.0.1}) and (\SGARefLink{sga7.sga7i.ei.n0-0-3}{0.0.3}). Let $p$ denote the characteristic exponent of $k(s)$, and let $\Gal(\bar\eta/\eta)$ (respectively, $\Gal(\bar s/s)$) denote the Galois group $\Gal(k(\bar\eta)/k(\eta))$ (respectively, the corresponding group for $\bar s/s$). By transport of structure, $\Gal(\bar\eta/\eta)$ acts on $k(\bar s)$; hence there is a homomorphism from $\Gal(\bar\eta/\eta)$ onto $\Gal(\bar s/s)$ whose kernel is the \underline{inertia group} $I$. The subfield of $k(\bar\eta)$ fixed by $I$ is a maximal unramified extension of $k(\eta)$.

The profinite group $I$ has a largest pro-$p$ subgroup $P$. The quotient $I/P$ is the \underline{tame inertia group}. Canonically,

\begin{equation*}
I/P\simeq \varprojlim_n\mu_n(k(\bar s))
\underset{\mathrm{dfn}}{=}\widehat{\mathbb{Z}}(1)(k(\bar s)).
\end{equation*}

%%%%%%%%%%  scan idx 15  |  folio 4  %%%%%%%%%%

If $\phi_n$ is the canonical map from $I/P$ to $\mu_n(\bar k)$ and $u$ is an element of $\bar V$ of valuation $1/n$, then, for $\sigma\in I$,

\begin{equation*}
\sigma(u)\equiv\phi_n(\sigma)\,u
\qquad (\text{modulo elements of valuation greater than }1/n).
\end{equation*}
% Printed source: \phi_n(u)u. The argument must be \sigma because \phi_n is
% a character of I/P; the English text transparently corrects this source typo.
In summary,

\begin{equation}\tag{\SGARefTarget{sga7.sga7i.ei.n0-3-1}{0.3.1}}
\begin{aligned}
\Gal(\bar\eta/\eta)&\supset I\supset P,\\
\Gal(\bar\eta/\eta)/I&\simeq\Gal(\bar s/s),\\
I/P&\simeq\widehat{\mathbb{Z}}(1)(k(\bar s)),\\
P&:\ \text{a pro-$p$ group}\quad(\{e\}\text{ if }p=1),
\end{aligned}
\end{equation}
and the action (by inner automorphisms in $\Gal(\bar\eta/\eta)$) of $\Gal(\bar\eta/\eta)/I$ on $I/P$ is identified with the natural action of $\Gal(\bar s/s)$ on $\widehat{\mathbb{Z}}(1)(k(\bar s))$.

\SGARefTarget{sga7.sga7i.ei.n0-4}{0.4}. More generally, let $S$ be strictly local and regular, let $D$ be a regular divisor in $S$, and let $p'$ be the characteristic exponent of the generic point of $D$. Let $\bar\eta$ be a geometric generic point of $S$, and let $\bar s$ be the resulting geometric point centered at the closed point $s$. By Abhyankar's lemma, the fundamental group $\pi_1(S-D,\bar\eta)$ has a filtration analogous to (\SGARefLink{sga7.sga7i.ei.n0-3-1}{0.3.1}):

\begin{equation}\tag{\SGARefTarget{sga7.sga7i.ei.n0-4-1}{0.4.1}}
\begin{aligned}
\pi_1(S-D,\bar\eta)&\supset I\supset P,\\
\pi_1(S-D,\bar\eta)/I&\simeq\Gal(\bar s/s),\\
I/P&\simeq\widehat{\mathbb{Z}}(1)(k(\bar s)),\\
P&:\ \text{with no quotient of order prime to }p'.
\end{aligned}
\end{equation}
The action by inner automorphisms of $\pi_1(S-D,\bar\eta)/I$ on $I/P$ is identified with the natural action of $\Gal(\bar s/s)$ on $\widehat{\mathbb{Z}}(1)(k(\bar s))$.

\SGARefTarget{sga7.sga7i.ei.n0-5}{0.5}. \underline{N\'eron desingularization}

The method of N\'eron desingularization gives the following result.

\underline{Lemma} \SGARefTarget{sga7.sga7i.ei.n0-5-1}{0.5.1}. \underline{Let $S\longrightarrow S_o$ be a morphism of henselian traits, with $S=\Spec(V)$ and $S_o=\Spec(V_o)$, and let $\pi$ be a uniformizer of $V_o$. Suppose that $\pi$ is a uniformizer of $V$, and that the extensions $k(s)/k(s_o)$ and $k(\eta)/k(\eta_o)$ are separable. Then $V$ is a filtered direct limit of smooth henselian $V_o$-algebras $B_i$ essentially of finite type over $V_o$.}

%%%%%%%%%%  scan idx 16  |  folio 5  %%%%%%%%%%

To construct the desired $V_o$-algebras, one takes $V_o\subset B\subset V$, with $B$ of finite type, and desingularizes $\Spec(B)$ by N\'eron's method ([1], \S~4).

Let $S=\Spec(V)$ be a henselian trait.

(\SGARefTarget{sga7.sga7i.ei.n0-5-2}{0.5.2}) If $S$ is of equal characteristic, with uniformizer $\pi$, one may apply (\SGARefLink{sga7.sga7i.ei.n0-5-1}{0.5.1}) with $S_o=\mathbf{F}_p[\pi]_{(\pi)}$. The extension $s/s_o$ is separable because $\mathbf{F}_p$ is perfect; the extension $\eta/\eta_o$ is separable because $\pi$, being a uniformizer, is not a $p$-th power (if $p\neq1$), and hence belongs to a $p$-basis. It follows that $S$ is a limit of spectra of smooth henselian $\mathbf{F}_p[\pi]_{(\pi)}$-algebras essentially of finite type.

(\SGARefTarget{sga7.sga7i.ei.n0-5-3}{0.5.3}) If $S$ is complete and of mixed characteristic, with perfect residue field $k$, then $V$ is an Eisenstein extension of the ring of Witt vectors (that is, the Cohen ring) $W(k)$:

\begin{equation*}
V\simeq W(k)[\pi]/\left(\pi^n+\sum_{i=0}^{n-1}a_i\pi^i\right).
\end{equation*}
If the $a_i$ are perturbed slightly, one obtains an isomorphic extension. We may therefore suppose that the $a_i$ belong to $W(k')$, where $k'$ is of finite type over $\mathbf{F}_p$. Let $k_o$ be the perfect closure of $k'$ in $k$, and set

\begin{equation*}
V_o=W(k_o)[\pi]/\left(\pi^n+\sum_{i=0}^{n-1}a_i\pi^i\right).
\end{equation*}
Lemma \SGARefLink{sga7.sga7i.ei.n0-5-1}{0.5.1} applies to $V/V_o$, and the residue field of $V_o$ is a radicial extension of a field of finite type over $\mathbf{F}_p$.

1. \underline{Arithmetic proof of the monodromy theorem}

The reader will find in [5], Appendix, a proof of the following result of A. Grothendieck.

\underline{Proposition} \SGARefTarget{sga7.sga7i.ei.n1-1}{1.1}. \underline{Let $S$ be a henselian trait, let $\ell$ be a prime different from the characteristic exponent $p$ of $k(s)$, and let $\rho$ be a continuous representation of $\Gal(\bar\eta/\eta)$ in $\mathrm{GL}(n,\mathbb{Q}_\ell)$. Suppose that the following condition holds:}

$(*_\ell)$ \quad \underline{No finite extension of $k(s)$ contains all roots of unity of $\ell$-power order.}

%%%%%%%%%%  scan idx 17  |  folio 6  %%%%%%%%%%

\underline{Then there exists an open subgroup $I_1$ of the inertia group $I$ such that $\rho(\sigma)$ is unipotent for every $\sigma\in I_1$.}

Condition $(*_\ell)$ is automatically satisfied when $k(s)$ is of finite type over the prime field (or is radicial over such a field).

The following theorem follows immediately from \SGARefLink{sga7.sga7i.ei.n1-1}{1.1}.

\underline{Monodromy theorem} \SGARefTarget{sga7.sga7i.ei.n1-2}{1.2}. \underline{With the preceding notation, let $X$ be a scheme of finite type over $\eta$, and let $H$ be a cohomology space of $X_{\bar\eta}$ with coefficients in $\mathbb{Q}_\ell$, of one of the types (\SGARefLink{sga7.sga7i.ei.n0-1-1}{0.1.1})--(\SGARefLink{sga7.sga7i.ei.n0-1-5}{0.1.5}). If condition $(*_\ell)$ is satisfied, the action on $H$ of every element of $I$ is quasi-unipotent.}
% Printed source refers to (1.1.1)--(1.1.5); these do not exist. The intended
% reference is the list (0.1.1)--(0.1.5) immediately above.

The limit arguments of \SGARefLink{sga7.sga7i.ei.n0-5}{0.5} allow us to dispense with hypothesis $(*_\ell)$.

\underline{Variant} \SGARefTarget{sga7.sga7i.ei.n1-3}{1.3}. \underline{Condition $(*_\ell)$ in \SGARefLink{sga7.sga7i.ei.n1-2}{1.2} is superfluous.}

Let $H'_o$ be a $\mathbb{Z}_\ell$-cohomology space of $X_{\bar\eta}$ giving rise to $H$, and put $H_o=H'_o/\text{torsion}$. By hypothesis,
$H=H_o\otimes_{\mathbb{Z}_\ell}\mathbb{Q}_\ell$, and the action of $\Gal(\bar\eta/\eta)$ is induced by

\begin{equation*}
\rho:\Gal(\bar\eta/\eta)\longrightarrow\mathrm{GL}(H_o).
\end{equation*}
Let $I_{\ell'}$ be the largest prime-to-$\ell$ subgroup of $I$. Its image in $\mathrm{GL}(H_o)$ is an $\ell$-adic Lie group of order prime to $\ell$, and is therefore finite.

Make an extension of traits $S'\longrightarrow S$. This does not change $H_o$ (SGA 4 XVI 6.6), and $\rho(I')$ has finite index in $\rho(I)$, because $I'/P'$ has finite index in $I/P$, by what precedes. It therefore suffices to prove the theorem over $S'$; we may suppose that

a) $\Gal(\bar\eta/\eta)$ acts trivially on $H_o/\ell H_o$;

b) if $S$ is of mixed characteristic, then $S$ is complete with perfect residue field (EGA $\mathrm{O}_{III}$ 10.3.1).

Put $S=\Spec(V)$, define $V_o$ as in (\SGARefLink{sga7.sga7i.ei.n0-5-2}{0.5.2}) and (\SGARefLink{sga7.sga7i.ei.n0-5-3}{0.5.3}), and apply (\SGARefLink{sga7.sga7i.ei.n0-5-1}{0.5.1}). We have

\begin{equation*}
V=\varinjlim_i B_i.
\end{equation*}

%%%%%%%%%%  scan idx 18  |  folio 7  %%%%%%%%%%

Here $S_i=\Spec(B_i)$ is essentially of finite type and smooth over $V_o$. Let $D_i\subset S_i$ be defined by the equation $\pi=0$. For sufficiently large $i$, $H_o$ comes from a twisted constant $\mathbb{Z}_\ell$-sheaf $\mathcal H_i$ on $S_i-D_i$ (by ``generic constructibility,'' (\SGARefLink{sga7.sga7i.ei.n0-1}{0.1})), and $\mathcal H_i/\ell\mathcal H_i$ is constant.

The filtrations (\SGARefLink{sga7.sga7i.ei.n0-3-1}{0.3.1}) and (\SGARefLink{sga7.sga7i.ei.n0-4-1}{0.4.1}) give rise to commutative diagrams:

\begin{equation}\tag{\SGARefTarget{sga7.sga7i.ei.n1-3-1}{1.3.1}}
\begin{tikzcd}[column sep=small,row sep=large]
\Gal(\bar\eta/\eta) \arrow[d] & \supset & I \arrow[d] & \supset & P \arrow[d] \\
\pi_1(S_i-D_i,\bar\eta) & \supset & I_i & \supset & P_i
\end{tikzcd}
\qquad\text{and}\qquad
\begin{tikzcd}[row sep=large]
I/P \arrow[r,dash,"\sim"] \arrow[d] & \widehat{\mathbb Z}(1) \arrow[d,equal] \\
I_i/P_i \arrow[r,dash,"\sim"] & \widehat{\mathbb Z}(1)
\end{tikzcd}
\end{equation}

% The authority has two distinct commutative diagrams separated by the word
% ``and.'' The horizontal isomorphisms in the right-hand diagram are rules
% labelled \sim, not arrows; its right vertical connector is equality.
The representation on $H_o$ of $\Gal(\bar\eta/\eta)$ is induced by an action of $\pi_1(S_i-D_i,\bar\eta)$ on $H_o$ that is trivial on $H_o/\ell H_o$. The group
\[
\Ker\bigl(\mathrm{GL}(H_o)\longrightarrow\mathrm{GL}(H_o/\ell H_o)\bigr)
\]
is a pro-$\ell$ group. Thus $P_i$ acts trivially, and we apply the following variant of \SGARefLink{sga7.sga7i.ei.n1-1}{1.1}.
% Printed source writes GL(H) -> GL(H/ell H), but the lattice being acted on
% and reduced here is H_o. The English text makes this evident subscript repair.

\underline{Variant} \SGARefTarget{sga7.sga7i.ei.n1-4}{1.4}. \underline{Retain the notation of \SGARefLink{sga7.sga7i.ei.n0-4}{0.4}, and let $\rho$ be a continuous representation of $\Gal(\bar\eta/\eta)$ in $\mathrm{GL}(n,\mathbb{Q}_\ell)$. Suppose that $k(s)$ satisfies $(*_\ell)$ and that $\rho(P)$ is finite. Then the action of $I$ is quasi-unipotent.}
% Printed source has Gal(n,Q_l), an unambiguous typographical error for GL.

% END INCLUDED SGA7I_English_current_buildable_source_20260823/components/02_expose_I_0_1_through_1_4.tex

% BEGIN INCLUDED SGA7I_English_current_buildable_source_20260823/components/03_expose_I_section_2.tex
% SOURCE AUTHORITY:
% expose_I_body.tex lines 249--366; scan indices 18--21; source folios 7--10.
% Hard stop before source line 368, section 3.

2. \underline{Vanishing cycles}

\SGARefTarget{sga7.sga7i.ei.n2-1}{2.1}. Let $S$ be a strictly local trait, and retain the notation of \SGARefLink{sga7.sga7i.ei.n0-3}{0.3}. Thus, by hypothesis, $s=\bar s$. Let $f:X\longrightarrow S$ be a scheme of finite type over $S$. The formalism of vanishing cycles is a tool for studying the relation between the cohomology of $X_s$ and that of $X_{\bar\eta}$, as well as the action of the inertia group on the latter, using information about $f$ that is \underline{local} on $X$.

A key result of this kind is the specialization theorem (SGA 4 XVI 2.2), according to which $X_s$ and $X_{\bar\eta}$ have the same cohomology when $f$ is proper and smooth.

In the transcendental theory, if $f:X\longrightarrow D$ is a proper morphism from an analytic space $X$ to the disk $D$, vanishing cycles arise as follows:

a) After shrinking $D$ if necessary, one may suppose that $X$ is a topological fibre bundle over the punctured disk $D^*$ (Thom's isotopy theorem). In particular, the fibres

%%%%%%%%%%  scan idx 19  |  folio 8  %%%%%%%%%%

$X_t$ $(t\in D^*)$ are all isomorphic.

b) It is not unreasonable to picture the special fibre $X_s$ as obtained from \lq\lq the\rq\rq\ general fibre $X_t$ by contracting polyhedra in $X_t$; one would then have a map $\pi:X_t\longrightarrow X_s$. The method of vanishing cycles consists in studying the difference between the cohomology of $X_s$ and that of $X_t$ by means of the Leray spectral sequence for $\pi$.

The geometric theory explained below is very close to this transcendental theory: the general point $t\in D$ must be replaced by $\bar\eta$, the geometric generic point of $S$. Conversely, the transcendental theory can be modeled on the geometric theory by replacing $t$ or $\bar\eta$ with the universal cover $\widetilde D^*$ of $D^*$, and $X_t$ or $X_{\bar\eta}$ with $X\times_D\widetilde D^*$. This eliminates many $\varepsilon$'s and $\eta$'s, but sometimes loses information.

\SGARefTarget{sga7.sga7i.ei.n2-2}{2.2}. Fix a torsion ring $\Lambda$ in which $p$ is invertible (for example, $\Lambda=\mathbb Z/\ell^n\mathbb Z$, with $\ell$ prime to $p$). Consider the morphisms

\begin{center}
\begin{tikzcd}[column sep=large,row sep=huge]
X_s \arrow[r,"i"] \arrow[d,"f_s"] & X \arrow[d,"f"] & & X_{\bar\eta} \arrow[ll,"\bar j"] \arrow[d,"f_{\bar\eta}"] \\
s \arrow[r] & S & \eta \arrow[l] & \bar\eta \arrow[l]
\end{tikzcd}
\end{center}
% Authority geometry: seven nodes and eight plain arrows. The top row is
% X_s -> X <- X_{bar eta}; the bottom row is s -> S <- eta <- bar eta;
% the three vertical arrows are f_s, f, and f_{bar eta}.
The sheaves of vanishing cycles $\psi^i$ on $X_s$ are

\begin{equation*}\tag{\SGARefTarget{sga7.sga7i.ei.n2-2-1}{2.2.1}}
\psi^i=i^*R^i\bar j_*\Lambda.
\end{equation*}
If $f$ is \underline{proper}, then (SGA 4 XII 5.5)

\begin{equation*}\tag{\SGARefTarget{sga7.sga7i.ei.n2-2-2}{2.2.2}}
i^*:H^p(X,R^q\bar j_*\Lambda)\xrightarrow{\ \sim\ }
H^p(X_s,i^*R^q\bar j_*\Lambda),
\end{equation*}
% The fresh French TeX drops the bar from j on the right of (2.2.2); the
% printed authority has bar-j on both sides.
and the Leray spectral sequence for $\bar j$ becomes

\begin{equation*}\tag{\SGARefTarget{sga7.sga7i.ei.n2-2-3}{2.2.3}}
H^p(X_s,\psi^q)\Longrightarrow H^{p+q}(X_{\bar\eta},\Lambda).
\end{equation*}
When nothing happens at infinity, this spectral sequence can sometimes be obtained without a properness hypothesis. That some hypothesis at infinity is necessary is already clear from the trivial case $X_s=\varnothing$.

%%%%%%%%%%  scan idx 20  |  folio 9  %%%%%%%%%%

The definitions immediately give:

\underline{Proposition} \SGARefTarget{sga7.sga7i.ei.n2-3}{2.3}. \underline{Let $\bar x$ be a geometric point of $X_s$ (for example, a closed point), let $X_{(\bar x)}$ be the strict henselization of $X$ at $\bar x$, and let $X_{(\bar x)\bar\eta}$ be the geometric generic fibre of the projection $X_{(\bar x)}\longrightarrow S$. Then}

\begin{equation*}
\psi^i_{\bar x}=H^i(X_{(\bar x)\bar\eta},\Lambda).
\end{equation*}
% The subscript bar-x is visible in the printed formula but absent from the
% fresh French TeX; it is restored here.

\underline{Corollary} \SGARefTarget{sga7.sga7i.ei.n2-4}{2.4}. \underline{Wherever $f$ is smooth, $\psi^i=0$ for $i>0$ and $\psi^0=\Lambda$.}

This is SGA 4 XV 2.1.

The group $\Gal(\bar\eta/\eta)$ (which, by hypothesis, equals the inertia group) acts on the $\psi^i$ by transport of structure, and

\underline{Scholium} \SGARefTarget{sga7.sga7i.ei.n2-5}{2.5}. \underline{The spectral sequence (\SGARefLink{sga7.sga7i.ei.n2-2-3}{2.2.3}) is $\Gal(\bar\eta/\eta)$-equivariant.}

\underline{Variant} \SGARefTarget{sga7.sga7i.ei.n2-6}{2.6}. Put $\bar S=\Spec(\bar V)$ and $\bar X=X\times_S\bar S$. The scheme $X_{\bar\eta}$ is the complement of $X_s=X_{\bar s}$ in $\bar X$. Put

\begin{equation*}
\phi^i=\underline{H}^{\,i+1}_{X_s}(\bar X,\Lambda)
\end{equation*}
(the local-cohomology sheaves). There is an exact sequence of sheaves

\begin{equation*}
0\longrightarrow\phi^{-1}\longrightarrow\Lambda\longrightarrow\psi^0
\longrightarrow\phi^0\longrightarrow0,
\end{equation*}
and, for $i>0$, isomorphisms $\psi^i\xrightarrow{\ \sim\ }\phi^i$. If $f$ is smooth, the $\phi^i$ vanish.

Put $\phi^i_{\mathrm{gl}}=H^i_{X_s}(\bar X,\Lambda)$, the global analogue of $\phi^i$. There is a spectral sequence

\begin{equation*}\tag{\SGARefTarget{sga7.sga7i.ei.n2-6-1}{2.6.1}}
H^p(X_s,\phi^q)\Longrightarrow\phi^{p+q}_{\mathrm{gl}},
\end{equation*}
and a long exact cohomology sequence which, when $f$ is proper, becomes

\begin{equation*}\tag{\SGARefTarget{sga7.sga7i.ei.n2-6-2}{2.6.2}}
\cdots\longrightarrow H^i(X_s,\Lambda)\longrightarrow H^i(X_{\bar\eta},\Lambda)
\longrightarrow\phi^i_{\mathrm{gl}}\longrightarrow\cdots.
\end{equation*}

\underline{Variant} \SGARefTarget{sga7.sga7i.ei.n2-7}{2.7}. Let $K_t$ be the maximal tamely ramified extension of $K$ in $\bar K$. Put $X_t=X\times_S\Spec(K_t)$, and let $\bar j_t$ be the projection from $X_t$ to $X$. Define tame variants of the $\psi^i$ by

%%%%%%%%%%  scan idx 21  |  folio 10  %%%%%%%%%%

\begin{equation*}
\psi^i_t=R^i\bar j_{t*}\Lambda.
\end{equation*}
Since $P$ is a pro-$p$ group and $p$ is invertible in $\Lambda$, one obtains

\begin{equation}\tag{\SGARefTarget{sga7.sga7i.ei.n2-7-1}{2.7.1}}
H^*(X_t,\Lambda)=H^*(X_{\bar\eta},\Lambda)^P,
\end{equation}
\begin{equation}\tag{\SGARefTarget{sga7.sga7i.ei.n2-7-2}{2.7.2}}
\psi^i_t=(\psi^i)^P.
\end{equation}
When $f$ is proper, the Leray spectral sequence for $\bar j_t$ is obtained from (\SGARefLink{sga7.sga7i.ei.n2-2-3}{2.2.3}) by passing to $P$-invariants:

\begin{equation}\tag{\SGARefTarget{sga7.sga7i.ei.n2-7-3}{2.7.3}}
H^p(X_s,\psi^q_t)\Longrightarrow H^{p+q}(X_t,\Lambda)
=H^{p+q}(X_{\bar\eta},\Lambda)^P.
\end{equation}

% END INCLUDED SGA7I_English_current_buildable_source_20260823/components/03_expose_I_section_2.tex

% BEGIN INCLUDED SGA7I_English_current_buildable_source_20260823/components/04_expose_I_section_3.tex
% SOURCE AUTHORITY:
% expose_I_body.tex lines 368--457; scan indices 21--24; source folios 10--13.
% Hard stop before source line 459, section 4.

3. \underline{Geometric proof of the monodromy theorem}

Let $\Lambda$ be as in \SGARefLink{sga7.sga7i.ei.n2-2}{2.2}.

\underline{Conjecture} \SGARefTarget{sga7.sga7i.ei.n3-1}{3.1} (\underline{purity}). \underline{Let $A$ be an
excellent regular strictly local ring} (\SGARefLink{sga7.sga7i.ei.n0-0-2}{0.0.2}) \underline{and let $D$ be a
regular divisor in $\Spec(A)$ defined by a parameter $x$. Put
$U=\Spec(A)-D=\Spec(A[1/x])$. Then}
% The French TeX/print has the unrelated cross-reference (0.2) and A[1/X].
% The definition of "strictly local" is (0.0.2), and the named parameter is x.

\begin{equation}\tag{\SGARefTarget{sga7.sga7i.ei.n3-1-1}{3.1.1}}
H^{q}(U,\Lambda)=
\begin{cases}
\Lambda & \text{if }q=0,\\[4pt]
\Lambda(-1) & \text{if }q=1\quad\text{(notation (\SGARefLink{sga7.sga7i.ei.n0-0-6}{0.0.6})),}\\[4pt]
0 & \text{if }q>1.
\end{cases}
\end{equation}

Here are some cases in which this conjecture is known (SGA 4 XIX): for
$q=0$ or $1$; when $A$ has characteristic $0$; when $A$ has equal
characteristic $p$ and resolution of singularities is available in
characteristic $p$ in dimensions $\leqslant\dim(A)$; or when
$(\Spec(A),D)$ is a smooth pair (the relative purity theorem, SGA 4 XVI
3.7).

When purity is available and
$D=\sum_{i\in B}D_i$ is a normal-crossings divisor in $\Spec(A)$ defined
by part of a regular system of parameters, the cohomology of
$U=\Spec(A)-D$ is given by

%%%%%%%%%%  scan idx 22  |  folio 11  %%%%%%%%%%

\begin{equation}\tag{\SGARefTarget{sga7.sga7i.ei.n3-1-2}{3.1.2}}
\begin{aligned}
H^0(U,\Lambda)&=\Lambda,\\[4pt]
H^1(U,\Lambda)&=\Lambda(-1)^B,\\[4pt]
H^q(U,\Lambda)&=\bigwedge\nolimits^q H^1(U,\Lambda)
  \qquad\text{(via the cup product).}
\end{aligned}
\end{equation}
These formulas are identical to those which, over $\mathbf C$, give the
cohomology of $D^{*B}\times D^k$.

\SGARefTarget{sga7.sga7i.ei.n3-2}{3.2}. Let $S$ be a strictly local trait, and retain the notation of
(\SGARefLink{sga7.sga7i.ei.n0-0-1}{0.0.1}) and (\SGARefLink{sga7.sga7i.ei.n0-0-3}{0.0.3}). Let $f:X\longrightarrow S$ be a morphism of finite
type. Suppose that $X$ is regular, that $X_\eta$ is smooth, and that
$X_s$ is a normal-crossings divisor in $X$. We shall suppose that this
divisor is, globally and not merely locally for the \'{e}tale topology, a
sum of regular divisors:
\[
(X_s)_{\mathrm{red}}=\sum_{i\in B}D_i
\]
(the general case can be recovered by localization). Denote their
multiplicities by $m_i$, so that
$X_s=\sum_{i\in B}m_iD_i$.

Let $x$ be a point of $X_s$, let $\bar x$ be a geometric point of $X$
lying over $x$ (for example, $x$ may be closed and equal to $\bar x$),
put $C=\{i\mid x\in D_i\}$, and let $K$ be the complex concentrated in
degrees $0$ and $-1$
\begin{equation*}
K:\mathbf Z^C\xrightarrow{\ d\ }\mathbf Z,
\qquad d((n_j))=\sum_jm_jn_j.
\end{equation*}

\underline{Theorem} \SGARefTarget{sga7.sga7i.ei.n3-3}{3.3} (\underline{uses purity}). \underline{The groups
of tame vanishing cycles $\psi^i_{t\bar x}$ are as follows:}

(i) \underline{There is a canonical isomorphism of graded
$\Lambda$-algebras}
\begin{equation*}
\psi^*_{t\bar x}\simeq
\bigwedge\nolimits^*\bigl(H_1(K)\otimes\Lambda(-1)\bigr)
\otimes_\Lambda\psi^0_{t\bar x}.
\end{equation*}

(ii) $\psi^0_{t\bar x}$ \underline{is the} $\Lambda$-\underline{algebra}
$\Lambda^E$, \underline{where $E$ is a set on which the tame inertia
group} $I_t=I/P=\widehat{\mathbf Z}(1)(k(s))$
\underline{acts transitively, and whose cardinality is the largest
prime-to-$p$ divisor of $\#H_0(K)$.}

Here is the method of proof.

a) Let $X_{(\bar x)}$ be the strict localization of $X$ at $\bar x$,
let $U=X_{(\bar x)}-X_{(\bar x)s}$ (the complement of a
normal-crossings divisor), let $t_j$ be a local equation for $D_j$, put
\[
X_n=X_{(\bar x)}\bigl[(t_j^{1/n})\bigr]
\qquad ((n,p)=1),
\]
which is a regular scheme, and let $U_n$ be the inverse image of $U$ in
$X_n$. By (\SGARefLink{sga7.sga7i.ei.n3-1-2}{3.1.2}),
\[
H^q(U_n,\Lambda)=\bigwedge\nolimits^q\bigl(\Lambda(-1)^C\bigr).
\]
For $m=nd$, the pullback map from $H^q(U_n,\Lambda)$ to
$H^q(U_m,\Lambda)$ is multiplication
% The French TeX/print has m=xd. In this inverse system x is a point while
% m,n,d are the indices; the pullback relation is necessarily m=nd.

%%%%%%%%%%  scan idx 23  |  folio 12  %%%%%%%%%%

by $d^q$. Putting
$\widetilde U=\varprojlim_{(n,p)=1}U_n$, one therefore has
\begin{equation}\tag{\SGARefTarget{sga7.sga7i.ei.n3-3-1}{3.3.1}}
\begin{cases}
H^0(\widetilde U,\Lambda)=\Lambda,\\[6pt]
H^q(\widetilde U,\Lambda)=\displaystyle\varinjlim_n
  H^q(U_n,\Lambda)=0 & \text{if }q\ne0.
\end{cases}
\end{equation}

b) $\widetilde U$ is a pro-covering of $U$ with group
$\widehat{\mathbf Z}(1)(k(s))^C$. By definition,
\[
\psi^*_{t\bar x}=H^*(U_{\eta_t},\Lambda).
\]
% The French working TeX misreads the printed Greek eta in U_{eta_t} as n.
$\widetilde U$ is a pro-covering, with group
$H_1(K)\otimes\widehat{\mathbf Z}(1)(k(s))$, of each connected component
of $U_{\eta_t}$, and these components are permuted transitively by the
tame inertia group. Formula (\SGARefLink{sga7.sga7i.ei.n3-3-1}{3.3.1}) and the Hochschild--Serre spectral
sequence now give the formulas in \SGARefLink{sga7.sga7i.ei.n3-3}{3.3}.

Pictorially, one may say that $U_{\bar\eta}$ is the fibre of a morphism
$\delta$ from the cohomological
$K(\widehat{\mathbf Z}(1)(k(s))^C,1)$ represented by $U$ to the
cohomological $K(\widehat{\mathbf Z}(1)(k(s)),1)$ represented by $\eta$;
the morphism $\delta$ induces
\[
d:(n_j)\longmapsto\sum_jm_jn_j:
\widehat{\mathbf Z}(1)(k(s))^C
\longrightarrow\widehat{\mathbf Z}(1)(k(s)).
\]

\underline{Corollary} \SGARefTarget{sga7.sga7i.ei.n3-4}{3.4} (\underline{uses purity}). \underline{Suppose
that $f$ is proper. Let $N$ be the least common multiple of the
multiplicities $m_j$ $(j\in B)$. Let $q\geqslant0$, and let $q'$ be the
minimum of $q$ and the greatest number of divisors $D_j$ passing through
one point. For $T$ in the tame inertia group $I_t=I/P$, one has}
\begin{equation*}
(T^N-1)^{q'+1}=0
\quad\underline{\text{on}}\quad
H^q(X_{\bar\eta},\Lambda)^P.
\end{equation*}
% The French print has exponent q'. For q=0 that asserts 1=0; the
% filtration argument and the bound restated in 3.8 require q'+1.
This follows at once from \SGARefLink{sga7.sga7i.ei.n3-3}{3.3}, Scholium \SGARefLink{sga7.sga7i.ei.n2-5}{2.5}, and Variant \SGARefLink{sga7.sga7i.ei.n2-7}{2.7}.

The required purity results are available for $H^1$, in characteristic
$0$, or when $S$ is essentially of finite type over a field.

\underline{Theorem} \SGARefTarget{sga7.sga7i.ei.n3-5}{3.5}. \underline{Let $C$ be a smooth projective curve
over the fraction field $k(\eta)$ of a discrete valuation ring $V$. The
inertia group $I$ acts on $H^1(C_{\bar\eta},\mathbf Q_\ell)$ by
quasi-unipotent automorphisms of level $2$: there exists $N$ such that,
for $T\in I$,}
\begin{equation*}
(T^N-1)^2=0
\quad\text{on}\quad H^1(C_{\bar\eta},\mathbf Q_\ell).
\end{equation*}
We reduce to the case in which $V$ is strictly local. The assertion is
that there exists an open subgroup $I_1$ of $I$ such that
$(T-1)^2=0$ on $H^1$ for $T\in I_1$.

%%%%%%%%%%  scan idx 24  |  folio 13  %%%%%%%%%%

We may extend the scalars of $k(\eta)$ to a finite extension, replacing
$V$ by its normalization in that extension. Since the image of the
pro-$p$ group $P$ in an $\ell$-adic Lie group
$\mathrm{GL}(n,\mathbf Q_\ell)$ is finite, we may in particular suppose
that $P$ acts trivially.

Let $C_o$ be a minimal model of $C$ over $V$ (for the existence of
$C_o$, which is based on Abhyankar's theorem, see the discussion in [3],
\S~2, p.~87; one could instead reduce to the case where $V$ is excellent
by passing to its completion and invoking SGA 4 XVI 1.6).

Let $C_1$ be a regular model of $C$, obtained from $C_o$ by blowups, whose
special fibre is a normal-crossings divisor. The arguments of \SGARefLink{sga7.sga7i.ei.n3-3}{3.3} and
\SGARefLink{sga7.sga7i.ei.n3-4}{3.4} apply to $C_1$, $\psi_t^0$, $\psi_t^1$, and
\[
H^1(C_{\bar\eta},\mathbb Z/\ell^n)^P
=H^1(C_{\bar\eta},\mathbb Z/\ell^n),
\]
with $N$ independent of $n$; Theorem \SGARefLink{sga7.sga7i.ei.n3-5}{3.5} follows.
% The print uses a letter o in psi_t^o, but this is cohomological degree zero.

\underline{Remark} \SGARefTarget{sga7.sga7i.ei.n3-6}{3.6}. Let $A$ be an abelian variety over $k(\eta)$.
Since $A$ is a quotient of a Jacobian, it follows again that the action
of $I$ on $H^1(A_{\bar\eta},\mathbb Q_\ell)$, or on $T_\ell(A)$, is
quasi-unipotent of level $2$.
% The print omits the geometric bar from the fibre in this sentence.

\underline{Remark} \SGARefTarget{sga7.sga7i.ei.n3-7}{3.7}. The action of $I$ on
$H^1(X_{\bar\eta},\mathbb Q_\ell)$ is in fact quasi-unipotent of level
$2$ for every scheme $X$ of finite type over $\eta$.
% The print has Q here; the surrounding l-adic statement requires Q_ell.

Theorem \SGARefLink{sga7.sga7i.ei.n3-5}{3.5} (or Remark \SGARefLink{sga7.sga7i.ei.n3-6}{3.6}) is the basis for the proof of the stable
reduction theorem for abelian varieties that will be given in Expos\'{e}
IX, \S~3.

\SGARefTarget{sga7.sga7i.ei.n3-8}{3.8}. It is clear that whenever purity and resolution of singularities
are available (for example, in characteristic $0$), the preceding method
proves the monodromy theorem and gives bounds for the nilpotence
exponent. Thus, in equal characteristic $0$, if $X$ is proper and smooth
over $\eta$ and $T$ belongs to an open subgroup $I_1$ of $I$, one has
\begin{equation*}
(T-1)^{q+1}=0
\quad\text{on}\quad H^q(X_{\bar\eta},\mathbb Z_\ell)
\qquad(T\in I_1\subset I).
\end{equation*}

% END INCLUDED SGA7I_English_current_buildable_source_20260823/components/04_expose_I_section_3.tex

% BEGIN INCLUDED SGA7I_English_current_buildable_source_20260823/components/05_expose_I_section_4.tex
% SOURCE AUTHORITY:
% expose_I_body.tex lines 459--589; scan indices 24--28; source folios 13--17.
% Hard stop before source line 591, section 5.

4. \underline{Vanishing criteria for the sheaves of vanishing cycles}

\SGARefTarget{sga7.sga7i.ei.n4-1}{4.1}. Let $S$ be a strictly local trait and let
$f:X\longrightarrow S$ be a scheme of finite type over $S$. Put
$n=\dim(X_\eta)$.

If $X'$ is the scheme-theoretic closure of $X_\eta$ in $X$, then:

a) On $X_s-X'_s$, one has $\phi^{-1}=\Lambda$, and the other
$\phi^i$ vanish;

%%%%%%%%%%  scan idx 25  |  folio 14  %%%%%%%%%%

b) on $X'_s$, the $\phi^i$ and $\psi^i$ relative to $X/S$ or to
$X'/S$ coincide, and $\phi^{-1}=0$.

Since $X'$ is flat over $S$, this often permits a reduction to the case
in which $X$ is flat over $S$.

\underline{Theorem} \SGARefTarget{sga7.sga7i.ei.n4-2}{4.2}.

(i) \underline{One has $\psi^i=0$ for $i>n$.}

(ii) \underline{More precisely, let $x$ be a point of $X_s$ whose
closure has dimension $k$}
($k=\operatorname{trdeg}_{k(s)}k(x)$),
\underline{and let $\bar x$ be a geometric point lying over $x$. Then}
\begin{equation*}
\psi^i_{\bar x}=0
\quad\underline{\text{for}}\quad i>n-k,
\qquad\text{i.e.}\qquad
d(\psi^i)\leqslant n-i
\quad(\underline{\text{notation of SGA 4 XIV 2.1}}).
\end{equation*}
Indeed,
\[
\psi^i_{\bar x}=H^i(X_{(\bar x)\bar\eta},\Lambda),
\]
and $X_{(\bar x)\bar\eta}$ is an inverse limit of affine schemes of
dimension at most $n$ over $k(\bar\eta)$. Assertion (i) therefore follows
from the affine Lefschetz theorem (SGA 4 XIV 3.2) by passage to the
limit.

To prove (ii), we may suppose that $X$ is flat over $S$ (\SGARefLink{sga7.sga7i.ei.n4-1}{4.1}) and that
there is a quasi-finite flat $S$-morphism
\[
g:X\longrightarrow\mathbb A^k_S
\]
such that $g(x)$ is the generic point of $\mathbb A^k_s$. Let $S'$ be
the strict localization of $\mathbb A^k_S$ at $\bar x$. Then:

a) $S'_{\bar\eta}$ is the spectrum of a field, and
$\Gal(\bar\eta'/s'_{\bar\eta})$ is a $p$-group $Q$;

b) by (i), applied to $g$,
\[
H^i(X_{(\bar x)\bar\eta'},\Lambda)=0
\qquad\text{for }i>n-k;
\]

c) one has
\[
H^i(X_{(\bar x)\bar\eta},\Lambda)
=H^i(X_{(\bar x)\bar\eta'},\Lambda)^Q,
\]
which proves the theorem.

\underline{Corollary} \SGARefTarget{sga7.sga7i.ei.n4-3}{4.3}. \underline{If $f$ is proper and flat and the
non-smooth locus of $f$ in $X_s$ has dimension at most $d$, the
specialization morphism}
\[
H^i(X_s)\longrightarrow H^i(X_{\bar\eta})
\]
\underline{is an isomorphism for $i>n+d+1$ and an epimorphism for
$i=n+d+1$.}

\SGARefTarget{sga7.sga7i.ei.n4-4}{4.4}. To prove, in certain cases, that the $\phi^i$ vanish for small
$i$, we shall use local duality, either directly or through the local
Lefschetz theorems of SGA 2 XIV. This tool is available only in
characteristic $0$, or in equal characteristic $p$ when resolution of
singularities \`a la Hironaka is available in the dimensions under
consideration.

%%%%%%%%%%  scan idx 26  |  folio 15  %%%%%%%%%%

\underline{Theorem} \SGARefTarget{sga7.sga7i.ei.n4-5}{4.5}. \underline{Suppose that $S$ has characteristic
$0$. Let $f:X\longrightarrow S$ be a morphism of finite type and let
$x$ be a closed point of $X_s$. Suppose that:}

(a) \underline{$x$ is an isolated non-smooth point in its fibre;}

(b) \underline{$X$ has relative geometric depth at least $n$ at $x$:
in a neighbourhood of $x$, the scheme $X$ is identified with a
subscheme defined by $k$ equations in a smooth scheme of relative
dimension $N$ at $x$, with $n\leqslant N-k$.}

\underline{Then $\phi^i_x=0$ for $i<n$.}

Let $X_x$ be the strict localization of $X$ at $x$, and put
$V=X_x-\{x\}$. For every trait $S'$ finite over $S$, put similarly
\[
X'_x=X_x\times_S S',\qquad V'=V\times_S S'.
\]
Let $\bar S$ be the normalization of $S$ in $\bar\eta$, and put
\[
\bar X_x=X_x\times_S\bar S,
\qquad \bar V=V\times_S\bar S.
\]
We have:

(1) If $\phi^i=0$ on $(X_x)_s$ (i.e. at every generization of $x$ in
$X_s$) for $i\leqslant m$, then
\begin{equation*}
H^{i-1}_{\{x\}}(\bar X_x)
=\widetilde H^{,i}(\bar V)
\xrightarrow{\ \sim\ }
\widetilde H^{,i}(\bar V_{\bar\eta})
=\phi^i
\qquad\text{for }i\leqslant m.
\end{equation*}
The $\phi^i$ are indeed the cohomology sheaves of $\bar V$ with support
in $V_s$; under the hypotheses of \SGARefLink{sga7.sga7i.ei.n4-5}{4.5}, condition (a) makes (1)
applicable with $m=\infty$ (SGA 4 XV 2.1).

(2) Hypothesis (b) of \SGARefLink{sga7.sga7i.ei.n4-5}{4.5} is stable under base change. By SGA 2 XIV
5.6, it implies
\begin{equation*}
\widetilde H^{,i}(V')
=H^{i-1}_{\{x\}}(X'_x)=0
\qquad\text{for }i<n.
\end{equation*}
We conclude by observing that $\widetilde H^{,i}(\bar V)$ is the
filtered colimit of the $\widetilde H^{,i}(V')$.

Combining \SGARefLink{sga7.sga7i.ei.n4-2}{4.2} and \SGARefLink{sga7.sga7i.ei.n4-5}{4.5} gives:

\underline{Corollary} \SGARefTarget{sga7.sga7i.ei.n4-6}{4.6}. \underline{Let $S$ be a henselian trait of
characteristic $0$, let $X/S$ be a relative complete intersection of
relative dimension $n$, and let $x\in X_s$ be an isolated non-smooth
point in its fibre. Then $\phi^i_x=0$ for $i\ne n$.}

\underline{Corollary} \SGARefTarget{sga7.sga7i.ei.n4-7}{4.7}. \underline{Under the hypotheses of \SGARefLink{sga7.sga7i.ei.n4-6}{4.6},
$\phi^n_x$ is a flat $\Lambda$-module.}

Indeed, if $\Lambda'$ is a finite $\Lambda$-algebra, a general nonsense
(SGA 4 XVII 5.2.11) gives, for the $\phi^i_{\Lambda'}$,

%%%%%%%%%%  scan idx 27  |  folio 16  %%%%%%%%%%

\begin{equation*}
\phi^{n-i}_{\Lambda'}
=\Tor_i^\Lambda(\Lambda',\phi^n_\Lambda),
\end{equation*}
and we apply \SGARefLink{sga7.sga7i.ei.n4-6}{4.6} to $\Lambda'$.

\underline{Variant} \SGARefTarget{sga7.sga7i.ei.n4-8}{4.8}. In \SGARefLink{sga7.sga7i.ei.n4-5}{4.5}, omit condition (a), and let $Z$ be the
non-smooth locus of $f$. One may still conclude that
\begin{equation*}
\phi^i_x=0\qquad\text{for }i<n-\dim Z_s.
\end{equation*}
We proceed by descending induction on $\dim\overline{\{x\}}$. The
induction hypothesis permits us to apply \SGARefLink{sga7.sga7i.ei.n4-5}{4.5} (1) with
$m=n-\dim Z_s$, and the proof is completed as before.

\SGARefTarget{sga7.sga7i.ei.n4-9}{4.9}. It is possible to obtain a vanishing criterion using only
information about the geometric fibres of $f$. For complete
intersections, however, the result is one degree weaker than \SGARefLink{sga7.sga7i.ei.n4-5}{4.5}.

\underline{Proposition} \SGARefTarget{sga7.sga7i.ei.n4-10}{4.10}. \underline{Under the general hypotheses
of \SGARefLink{sga7.sga7i.ei.n4-5}{4.5}, suppose that:}

(a) \underline{$x$ is an isolated non-smooth point in its fibre;}

(b) \underline{for every point $y$ of $X_{\bar\eta}$ whose closure in
$X_{\bar\eta}$ has dimension $d(y)$ and which is a generization of $x$,
the \'{e}tale depth satisfies}
\begin{equation*}
\operatorname{depth}_{\mathrm{\acute{e}t},y}(X_{\bar\eta})
\geqslant n-d(y),
\quad\text{i.e.}\quad
H^i_{\{y\}}(X_{\bar\eta})=0
\quad\underline{\text{for }}i\leqslant n-d(y);
\end{equation*}

(c) \underline{$\operatorname{depth}_{\mathrm{\acute{e}t},x}(X_s)
\geqslant n$.}

\underline{Then $\phi^i_x=0$ for $i<n-1$.}

Retain the notation of \SGARefLink{sga7.sga7i.ei.n4-5}{4.5}. We may apply \SGARefLink{sga7.sga7i.ei.n4-5}{4.5} (1) with $m=\infty$.
By the local Lefschetz theorem (SGA 2 XIV), hypothesis (b) gives
\begin{center}
\begin{tikzcd}[row sep=large, column sep=huge]
H^i(V') \arrow[r, "\sim"] & H^i(V_s) & {(\ i<n\text{-}1\ )} \\
H^i(V') \arrow[r, hook] & H^i(V_s) & {(\ i=n\text{-}1\ )}
\end{tikzcd}
\end{center}
% Direct authority comparison: the first arrow is an isomorphism arrow
% with the tilde above; the second is a rightward hook arrow, not a
% leftward arrowhead. There are no vertical or diagonal arrows.
Condition (c) then gives
\[
H^i(\bar V)=\varinjlim H^i(V')=0
\qquad\text{for }i\leqslant n-1,
\]
which proves \SGARefLink{sga7.sga7i.ei.n4-10}{4.10}.

\SGARefTarget{sga7.sga7i.ei.n4-11}{4.11}. As in \SGARefLink{sga7.sga7i.ei.n4-8}{4.8}, one can give a variant of \SGARefLink{sga7.sga7i.ei.n4-10}{4.10} in which (a) is not
assumed.

%%%%%%%%%%  scan idx 28  |  folio 17  %%%%%%%%%%

% END INCLUDED SGA7I_English_current_buildable_source_20260823/components/05_expose_I_section_4.tex

% BEGIN INCLUDED SGA7I_English_current_buildable_source_20260823/components/06_expose_I_section_5.tex
% SOURCE AUTHORITY:
% expose_I_body.tex lines 591--688; scan indices 28--31; source folios 17--20.
% Hard stop before source line 689, section 6.

5. \underline{Action of monodromy on $\pi_1$}

In this section (Grothendieck's expos\'{e} of 19 November 1968), we
shall use the semistable reduction theorem for curves. Artin and Winters
[2] gave a direct proof of this theorem; it can also be deduced [3] from
the analogous theorem for abelian varieties (IX \S~3, using \SGARefLink{sga7.sga7i.ei.n3-6}{3.6}). We
shall also use Schlessinger's theory (Rim's Expos\'{e} VI).

\SGARefTarget{sga7.sga7i.ei.n5-1}{5.1}. Let $k$ be an algebraically closed field, let $C$ be a projective
curve over $k$, and let $x_1,\ldots,x_n$ be a finite set of closed points
of $C$. The pointed curve
\[
X=(C;x_1,\ldots,x_n)
\]
is called \underline{stable} if:

(a) $C$ is reduced, and its only singularities are double points with
distinct tangents;

(b) the $x_i$ are distinct smooth points of $C$;

(c) if $C_1$ is an irreducible component of $C$, and $E$ is the set of
points of the normalization $C'_1$ of $C_1$ lying above a singular point
of $C$ or above one of the $x_i$, then, if $C'_1$ has genus $0$
(respectively $1$), the set $E$ has at least $3$ (respectively $1$)
elements.
% The print says "a singular point of C_1". Stability requires all points
% above the singular locus of the whole curve C, including intersections
% with other irreducible components.

Condition (c) is equivalent to each of the following:

(c') $X$ has no nontrivial infinitesimal automorphisms.

(c'') If $\omega$ is the invertible dualizing sheaf, the invertible
sheaf
\[
\omega'=\omega\Bigl(\sum_i x_i\Bigr)
\]
is ample.

An $S$-scheme $C$ equipped with sections $x_1,\ldots,x_n$ is a
\underline{stable curve} over $S$ if its geometric fibres are stable
curves. In [3], \S\S~1--2, only the case $n=0$ is considered. The
extension of the results of loc. cit. to the general case is immediate:

(1) If $X/k$ is stable, then
$H^1(C,\omega'^{\otimes n})=0$ for $n\geqslant2$, and
$\omega'^{\otimes n}$ is very ample for $n\geqslant3$ (cf. [3] 1.2).

(2) The formal moduli scheme $M$ of $X$ is smooth, and $C$ remains
singular over a normal-crossings divisor of $M$.

%%%%%%%%%%  scan idx 29  |  folio 18  %%%%%%%%%%

(3) If $X$ and $Y$ are stable over $S$, then $\Isom_S(X,Y)$ is finite
and unramified over $S$ (isomorphisms are required to respect the marked
points).

(4) If $X/k$ is stable, the map
\[
\Aut_k(X)\longrightarrow\Aut_k\Pic^\circ(C)
\]
is injective.

Finally, we have:

\underline{Proposition} \SGARefTarget{sga7.sga7i.ei.n5-1-1}{5.1.1}. \underline{Let
$X_\eta=(C_\eta;x_1,\ldots,x_n)$ be stable over the generic point
$\eta$ of a trait $S$, with $C_\eta$ smooth. There exist a finite
separable extension $\eta'$ of $\eta$ and a stable curve $X'$ over the
normalization $S'$ of $S$ in $\eta'$ such that}
\begin{equation*}
X'\otimes_{S'}\eta'=X_\eta\otimes_\eta\eta'.
\end{equation*}
% The print omits the prime from the second occurrence of eta in the
% phrase defining S'. The extension and the displayed base change make
% eta' mathematically forced.

This follows from the usual semistable reduction theorem, which permits
one to choose $\eta'$ so that the geometric special fibre of the minimal
model $C''$ of $C_{\eta'}$ over $S'$ satisfies (a). One then blows up,
iteratively if necessary, \underline{smooth} points of $C''$ so that
(b) is satisfied. Next one contracts the chains of smooth rational
components of self-intersection $-2$ which contain none of the
$\bar x_i$ in the special fibre; the result is $C'$ (cf. [3], proof of
2.3, p.~88).
% The print says "self intersection two"; the omitted minus sign is the
% standard and mathematically necessary value for the contracted chains.

As in loc. cit., one can check that $C'$ exists as soon as the Jacobian
of $C_\eta$ has semistable reduction. The smoothness hypothesis on
$C_\eta$ is moreover unnecessary.

\SGARefTarget{sga7.sga7i.ei.n5-2}{5.2}. Let $X_\eta$ be a geometrically connected scheme of finite type
over the fraction field of a strictly local trait $S$. For a geometric
point $\xi$ of $X_{\bar\eta}$, there is an exact sequence
\begin{equation}\tag{\SGARefTarget{sga7.sga7i.ei.n5-2-1}{5.2.1}}
\begin{tikzcd}[column sep=large, row sep=tiny]
0 \arrow[r] & \pi_1(X_{\bar\eta},\xi) \arrow[r]
  & \pi_1(X,\xi) \arrow[r]
  & \Gal(\bar\eta/\eta) \arrow[r] & 0 .
\end{tikzcd}
\end{equation}
% The source typescript omits the arrowhead on the first exact-sequence
% map. The English reader supplies the mathematically required arrow.
It induces a homomorphism from $\Gal(\bar\eta/\eta)$ to the group of
outer automorphisms of $\pi_1(X_{\bar\eta})$. Passing to the largest
prime-to-$p$ quotient $\pi_1^{(p)}(X_{\bar\eta})$ of
$\pi_1(X_{\bar\eta})$, one obtains
\begin{equation}\tag{\SGARefTarget{sga7.sga7i.ei.n5-2-2}{5.2.2}}
\Gal(\bar\eta/\eta)
\longrightarrow
\Aut_{\mathrm{ext}}\bigl(\pi_1^{(p)}(X_{\bar\eta})\bigr).
\end{equation}

Suppose that $X_\eta$ has a rational point $x$, and let $\bar x$ be the
geometric point
\[
\bar\eta\longrightarrow\eta\xrightarrow{\ x\ }X_\eta.
\]
For $\xi=\bar x$, the exact sequence (\SGARefLink{sga7.sga7i.ei.n5-2-1}{5.2.1}) splits canonically, giving

%%%%%%%%%%  scan idx 30  |  folio 19  %%%%%%%%%%

\begin{equation}\tag{\SGARefTarget{sga7.sga7i.ei.n5-2-3}{5.2.3}}
[x]:\Gal(\bar\eta/\eta)
\longrightarrow\Aut\bigl(\pi_1(X_{\bar\eta},\bar x)\bigr)
\longrightarrow\Aut\bigl(\pi_1^{(p)}(X_{\bar\eta},\bar x)\bigr).
\end{equation}
% The print refers here to the "exact sequence (5.2.2)". The only exact
% sequence is (5.2.1); (5.2.2) is the induced outer action.

\underline{Theorem} \SGARefTarget{sga7.sga7i.ei.n5-3}{5.3}. \underline{The image of $P$ under (\SGARefLink{sga7.sga7i.ei.n5-2-2}{5.2.2})
(or under (\SGARefLink{sga7.sga7i.ei.n5-2-3}{5.2.3}), which amounts to the same thing) is finite.}

\SGARefTarget{sga7.sga7i.ei.n5-3-1}{5.3.1}. \underline{Reduction to the case in which $X_\eta$ is
geometrically normal and integral.}

Let $X'_\eta$ be the normalization of $X_\eta$, let
$(X'_{i\eta})_i$ be the family of connected components of $X'_\eta$,
put
\[
X''_\eta=X'_\eta\times_{X_\eta}X'_\eta,
\]
and let $(X''_{j\eta})_j$ be the family of connected components of
$X''_\eta$. After a finite extension of $k(\eta)$, we may suppose that
the $X'_{i\eta}$ are geometrically normal and integral and have rational
points $x_i$, and that the $X''_{j\eta}$ are geometrically connected and
have rational points $y_j$.

For $k=1$ or $2$, and whenever $p_k(y_j)\in X'_{i\eta}$, choose a class
of paths $\ell_{kij}$ from $\operatorname{pr}_k(\bar y_j)$ to $\bar x_i$.
It is an element of a torsor under
$\pi_1(X'_{i\bar\eta},\bar x_i)$. Let $\bar\ell_{kij}$ be the
corresponding prime-to-$p$ path (the same path modulo
$\Ker(\pi\longrightarrow\pi^{(p)})$). We may choose
$\bar\ell_{kij}$ to be invariant under $P$: the set of prime-to-$p$
paths is an inverse limit of finite sets of cardinality prime to $p$ on
which the pro-$p$ group $P$ acts.

Consequently, if a finite-index subgroup $P'$ of $P$ acts trivially on
the groups $\pi_1^{(p)}(X'_{i\bar\eta},\bar x_i)$, the Van Kampen
formulas---equivalently, the fact that the normalization map
$X'_{\bar\eta}\longrightarrow X_{\bar\eta}$ is effective descent for
\'{e}tale coverings---show that $P'$ acts trivially on
$\pi_1^{(p)}(X_{\bar\eta},\bar x)$.
% The print carries a spurious component index on the target of the
% normalization map; effective descent concerns X'_{bar eta} -> X_{bar eta}.

\SGARefTarget{sga7.sga7i.ei.n5-3-2}{5.3.2}. If $X_{\bar\eta}$ is normal and $U$ is an open subset, then
$\pi_1(U)$ maps onto $\pi_1(X_{\bar\eta})$; we may therefore suppose
that $X_{\bar\eta}$ is smooth and affine. By cutting with
\underline{generic} hyperplane sections (which replaces $S$ by its
strict localization at the generic point of the special fibre of an
$\mathbb A^n_S$) and applying Bertini, we reduce to the case in which
$X_\eta$ is \underline{smooth, geometrically connected, and of
dimension one}.

After a finite extension of $k(\eta)$, we may suppose that $X_\eta$ is
the complement in a smooth projective curve $\bar X_\eta$ of a finite
set of rational points $x_i$, whose number we are free to increase. By
the semistable reduction theorem (\SGARefLink{sga7.sga7i.ei.n5-1-1}{5.1.1}), we reduce to the case in which
there is a stable curve $X=(C;x_1,\ldots,x_n)$ over $S$ such that
\[
X_\eta=C_\eta-\{x_1,\ldots,x_n\}.
\]
We then apply the following result.

%%%%%%%%%%  scan idx 31  |  folio 20  %%%%%%%%%%

\underline{Theorem} \SGARefTarget{sga7.sga7i.ei.n5-4}{5.4}. \underline{Let $S$ be a strictly local scheme,
let $f:\bar X\longrightarrow S$ be a proper flat morphism whose special
fibre is a curve with ordinary double points as its only non-smooth
points, and let $Y$ be a subscheme of $\bar X$ which is the union of a
finite number of disjoint sections contained in the smooth locus. Let
$x$ be a rational point of $X=\bar X-Y$. Let $U$ be the open subset of
$S$ over which $X$ is smooth, and let $\xi$ be a geometric point of $U$.
Then the image of $\pi_1(U,\xi)$ in
$\Aut(\pi_1^{(p)}(X_\xi,x\xi))$ is abelian and prime to $p$.}

Observe that $\pi_1(U,\xi)$ can act on
$\pi_1^{(p)}(X_\xi,x\xi)$ only because the groups $\pi_1^{(p)}$
satisfy the specialization theorem.

We reduce first to the case in which $S$ is complete, and then to the
universal case in which $S$ is a formal moduli scheme for the special
fibre equipped with its sections. Then $S$ is the spectrum of a formal
power-series ring over a Cohen ring:
\[
S=\Spec\bigl(W[[t_1,\ldots,t_n]]\bigr),
\]
and $U$ is the complement of a normal-crossings divisor with equation
$t_1\cdots t_m=0$. In this universal case, Abhyankar's lemma shows that
$\pi_1(U,\xi)$ is abelian and prime to $p$, which proves the theorem.

\underline{Remark} \SGARefTarget{sga7.sga7i.ei.n5-5}{5.5}. One can prove a result analogous to \SGARefLink{sga7.sga7i.ei.n5-4}{5.4} in
arbitrary relative dimension by taking a generic plane section of
dimension one and applying Bertini.


% END INCLUDED SGA7I_English_current_buildable_source_20260823/components/06_expose_I_section_5.tex

% BEGIN INCLUDED SGA7I_English_current_buildable_source_20260823/components/07_expose_I_section_6.tex
6. \underline{Appendix}: \underline{arithmetic proof of the stable reduction theorem}

(by P. Deligne)

Let $S$ be a trait and let $A_\eta$ be an abelian variety over $k(\eta)$,
of dimension $d\neq 0$. For a local morphism of traits
$S'\longrightarrow S$, denote by $A_{\eta'}$ the abelian variety over
$k(\eta')$ obtained from $A_\eta$ by extension of scalars. Let
$\mathcal A_{S'}$ be the N\'eron model of $A_{\eta'}$ over $S'$, let
$\mathcal A_{S',s'}$ be its special fibre, and let
$\mathcal A^\circ_{S',s'}$ be the identity component of
$\mathcal A_{S',s'}$. The stable reduction theorem is as follows.

%%%%%%%%%%  scan idx 32  |  folio 21  %%%%%%%%%%

\underline{Theorem \SGARefTarget{sga7.sga7i.ei.n6-1}{6.1}. For a suitable $S'$, $A_{\eta'}$ has stable
reduction, i.e. $\mathcal A^\circ_{S',s'}$ is an extension of an abelian
variety by a torus.}

It follows fairly readily from \SGARefLink{sga7.sga7i.ei.n6-1}{6.1} that one may choose $S'$ such that
$k(\eta')$ is a finite separable extension of $k(\eta)$. Suppose first
that the residue field of $S$ is of finite type over the prime field (the
same argument would work if $k(s)$ were radicial over such a field).

Let $\ell$ be a prime different from the residual characteristic exponent
$p$, let $T_\ell(A)$ be the Tate module, let
$V_\ell(A)=T_\ell(A)\otimes\mathbf Q_\ell$, and let
\[
 \rho:\Gal(\bar\eta/\eta)\longrightarrow
 \mathrm{GL}(T_\ell(A))
\]
be the natural representation. After a preliminary extension of scalars,
we are reduced to the case in which the inertia group $I$ acts
unipotently (\SGARefLink{sga7.sga7i.ei.n1-2}{1.2}), and hence through its largest pro-$\ell$ quotient
$\mathbf Z_\ell(1)$ (\SGARefLink{sga7.sga7i.ei.n0-3}{0.3}). Let $T$ be the image in
$\mathrm{GL}(T_\ell(A))$ of a generator of $\mathbf Z_\ell(1)$, and let
$r$ be the largest integer $\geq 0$ such that $(T-1)^r\neq 0$. We have:

\underline{Lemma \SGARefTarget{sga7.sga7i.ei.n6-2}{6.2}. The map}
\[
 (\sigma,v)\in I\times V_\ell(A)
 \longmapsto (\rho(\sigma)-1)^r(v)
\]
\underline{induces the morphism and the isomorphism in the commutative
diagram}
\begin{equation}\tag{\SGARefTarget{sga7.sga7i.ei.n6-2-1}{6.2.1}}
\begin{tikzcd}[column sep=huge, row sep=huge]
V_\ell(A)_I\otimes\mathbf Q_\ell(r)
  \arrow[r,"M"] \arrow[d]
& V_\ell(A)^I \\
\bigl(V_\ell(A)/\Ker(T-1)^r\bigr)\otimes\mathbf Q_\ell(r)
  \arrow[r,"\sim"]
& \operatorname{Im}(T-1)^r \arrow[u]
\end{tikzcd}
\end{equation}
$V_\ell(A)_I$ \underline{and} $V_\ell(A)^I$ \underline{are
$\Gal(\bar s/s)$-modules, and $M$ is Galois-equivariant.}

\SGARefTarget{sga7.sga7i.ei.n6-3}{6.3}. An $\ell$-adic representation
\[
 \tau:\Gal(\bar s/s)\longrightarrow\mathrm{GL}(V)
\]
is said to be of weight $n$ ($n\in\mathbf Z$) if the following condition
is satisfied:

(*) For a suitable model $X$ of $k(s)$ (that is, an irreducible variety
of finite type over the prime field such that $k(s)=k(X)$), $\tau$
factors through $\pi_1(X,\bar s)$ and, for every closed point $x$ of
$X$, the eigenvalues of $\tau(F_x)$ are algebraic numbers all of whose
complex conjugates have absolute value $|k(x)|^{n/2}$ (here $F_x$
denotes the geometric Frobenius $\phi_x^{-1}$).

%%%%%%%%%%  scan idx 33  |  folio 22  %%%%%%%%%%

If $V$ (respectively $W$) has weight $n$ (respectively $m$) and
$n\neq m$, then $\Hom_{\Gal}(V,W)=0$.

\SGARefTarget{sga7.sga7i.ei.n6-4}{6.4}. It follows from the universal property of the N\'eron model that
$V_\ell(A)^I=V_\ell(\mathcal A^\circ_{S,s})$. If $a$ is the dimension
of the largest abelian quotient of $\mathcal A^\circ_{S,s}$ and $m$ is
the dimension of its multiplicative part, then, by Weil,
$V_\ell(A)^I$ is an extension of a representation of dimension $2a$ and
weight $-1$ by a representation of dimension $m$ and weight $-2$.

\SGARefTarget{sga7.sga7i.ei.n6-5}{6.5}. Let $A^*$ be the abelian variety dual to $A$. Recall that
$V_\ell(A^*)$ and $V_\ell(A)$ (and hence $V_\ell(A^*)^I$ and
$V_\ell(A)_I$) are in duality with values in $\mathbf Q_\ell(1)$. If
$a^*$ and $m^*$ are defined for $A^*$ as $a$ and $m$ are for $A$, it
follows from \SGARefLink{sga7.sga7i.ei.n6-4}{6.4} that $V_\ell(A)_I$ is an extension of a representation
of dimension $m^*$ and weight $0$ by a representation of dimension
$2a^*$ and weight $-1$.

\SGARefTarget{sga7.sga7i.ei.n6-6}{6.6}. The Galois module $\mathbf Q_\ell(r)$ has weight $-2r$. Since
$\operatorname{Im}(T-1)^r\neq 0$, it follows from \SGARefLink{sga7.sga7i.ei.n6-4}{6.4}, \SGARefLink{sga7.sga7i.ei.n6-5}{6.5}, and the isomorphism in \SGARefLink{sga7.sga7i.ei.n6-2}{6.2}
that $r\leq 1$. If $r=0$, then
\[
\left\{
\begin{array}{l}
a+m\leq d,\\[1ex]
2a+m=2d,
\end{array}
\right.
\]
so $m=0$, $a=d$, and $A_\eta$ has good reduction. Suppose that $r=1$.
It then follows from \SGARefLink{sga7.sga7i.ei.n6-2}{6.2} that $\operatorname{Im}(T-1)$ has weight $-2$ and that
$V_\ell(A)/V_\ell(A)^I$ has weight $0$; these spaces have the same
dimension $m_1\leq m$.

We then have
\[
 \dim V_\ell(A)^I=2d-m_1=2a+m,
\]
and
\[
 2\dim\mathcal A_{S,s}\geq 2(a+m)\geq
 2a+m+m_1=2d=2\dim\mathcal A_{S,s},
\]
so $A_\eta$ has stable reduction
($\dim\mathcal A_{S,s}=a+m$). At the same time, we have obtained
$m_1=m$; that is,
\[
 \operatorname{Im}(T-1)\subset V_\ell(A)^I\simeq
 V_\ell(\mathcal A^\circ_{S,s})
\]
is the $V_\ell$ of the multiplicative part of
$\mathcal A^\circ_{S,s}$.

\SGARefTarget{sga7.sga7i.ei.n6-7}{6.7}. To treat the general case, we shall use (\SGARefLink{sga7.sga7i.ei.n0-5}{0.5}) (cf. \SGARefLink{sga7.sga7i.ei.n1-3}{1.3}). We reduce
to the case in which:

%%%%%%%%%%  scan idx 34  |  folio 23  %%%%%%%%%%

(a) the conditions of (\SGARefLink{sga7.sga7i.ei.n0-5-2}{0.5.2}) or (\SGARefLink{sga7.sga7i.ei.n0-5-3}{0.5.3}) are satisfied;

(b) $\Gal(\bar\eta/\eta)$ acts trivially on $A_\ell$;

(c) the action of $I$ on $T_\ell(A)$ is unipotent.

Under these hypotheses, we shall prove that $A_\eta$ has stable
reduction. Hypothesis (a) allows us to apply (\SGARefLink{sga7.sga7i.ei.n0-5-1}{0.5.1}) and write
$S=\varprojlim S_i$ as in \SGARefLink{sga7.sga7i.ei.n1-3}{1.3}, whose notation we retain. Let $s_i$ be
the closed point of $S_i$.

For sufficiently large $i$, the identity component $\mathcal A^\circ$
of the N\'eron model $\mathcal A_S$ descends to a smooth group scheme
$\mathcal A_i$ with connected fibres over $S_i$, and
$(\mathcal A_i)_\ell$ is constant over $S_i-D_i$. The group
$\pi_1(S_i-D_i,\bar\eta)$ acts on $T_\ell(A)$ and acts trivially on
$T_\ell(A)/\ell T_\ell(A)\simeq A_\ell$. As in \SGARefLink{sga7.sga7i.ei.n1-3}{1.3}, one sees that
$P_i$ acts trivially. We then have, by (\SGARefLink{sga7.sga7i.ei.n1-3-1}{1.3.1}),
\[
 T_\ell(A)_I=T_\ell(A)_{I_i},\qquad
 T_\ell(\mathcal A_{i,s_i})=T_\ell(\mathcal A^\circ_s)
 =T_\ell(A)^I=T_\ell(A)^{I_i},
\]
and (\SGARefLink{sga7.sga7i.ei.n6-2-1}{6.2.1}) is a diagram of $\Gal(\bar s_i/s_i)$-modules to which the
arguments of \SGARefLink{sga7.sga7i.ei.n6-4}{6.4}--\SGARefLink{sga7.sga7i.ei.n6-6}{6.6} may be applied.

%%%%%%%%%%  scan idx 35  |  folio 24  %%%%%%%%%%

\begin{center}
\underline{BIBLIOGRAPHY}
\end{center}

\begin{enumerate}[label={[\arabic*]},leftmargin=3em]
\item M. Artin, ``Algebraic approximation of structures over complete
local rings,'' \emph{Publ. Math. I.H.E.S.} \underline{36} (1969),
pp. 23--58.

\item M. Artin and G. Winters, ``Degenerate fibres and reduction of
curves.''

\item P. Deligne and D. Mumford, ``The irreducibility of the space of
curves of a given genus,'' \emph{Publ. Math. I.H.E.S.} \underline{36}
(1969), pp. 75--110.

\item J.-P. Serre, \emph{Corps locaux}, Publ. Math. Univ.
Nancago--Hermann, 1962.

\item J.-P. Serre and J. Tate, ``Good reduction of abelian varieties,''
\emph{Ann. of Math.} \underline{88}, no. 3 (1968), pp. 492--517.
\end{enumerate}

% END INCLUDED SGA7I_English_current_buildable_source_20260823/components/07_expose_I_section_6.tex

\clearpage
\phantomsection\addcontentsline{toc}{section}{Exposé II}
% BEGIN INCLUDED SGA7I_English_current_buildable_source_20260823/components/08_expose_II_complete.tex
%%%%%%%%%%  scan idx 36  |  folio 25  %%%%%%%%%%

\begin{center}
\underline{\SGARefAliasTarget{sga7.sga7i.eii.nroot}Expos\'e II}

\vspace{1em}
\underline{\underline{Finiteness properties of the fundamental group}}

\vspace{1em}
\underline{by Mich\`ele Raynaud}
\end{center}

Let $k$ be a separably closed field of characteristic $p\geq 0$, let
$X$ be a connected scheme of finite type over $k$, and let
$\Pi_1^{(p)}(X)$ be the largest ``prime-to-$p$'' quotient of $\Pi_1(X)$.
In this expos\'e we prove that, when $X$ is a surface,
$\Pi_1^{(p)}(X)$ is topologically finitely presented, and that the same
holds for arbitrary $X$ if resolution of singularities is assumed.

For a surface, the proof method, due to J.-P. Murre, consists in reducing
to a fibration over $P^1$ having properties (i), (ii), and (iii) of the
following proposition.

\underline{Proposition \SGARefTarget{sga7.sga7i.eii.n2-1}{2.1}. Let $k$ be an algebraically closed field and
let $S$ be a projective, irreducible, smooth surface over $k$. Then there
exists a blow-up}
\[
 g:S'\longrightarrow S,
\]
\underline{where $S'$ is a regular surface equipped with a fibration}
\[
 f:S'\longrightarrow P^1,
\]
\underline{such that $f$ has a section $\sigma$ and the following
conditions hold:}

(i) \underline{$f$ is smooth at the points of $\sigma(P^1)$.}

(ii) \underline{The fibres of $f$ are geometrically integral curves
having only ordinary singularities.}

(iii) \underline{There exists a nonempty open subset $U$ of $P^1$ such
that the fibres of $f$ over the points of $U$ are smooth.}

The proposition follows from the existence of a projective embedding
$S\longrightarrow P^r$ for which a Lefschetz pencil exists (XVII 2.5).
Indeed, let $D=\{H_t\}_{t\in P^1}$ be such a pencil of hyperplanes. Recall
that it has the following properties (loc. cit. \SGARefLink{sga7.sga7i.eii.n2-2}{2.2}):

%%%%%%%%%%  scan idx 37  |  folio 26  %%%%%%%%%%

(a) The axis $\Delta$ of $D$ meets $S$ transversely; the closed
subscheme $\Delta\cap S$ is therefore reduced and consists of finitely
many closed points $x_1,\ldots,x_n$ of $S$.

(b) There exists a nonempty open subset $U$ of $P^1$ such that, for
every $t\in U$, $H_t$ meets $S$ transversely.

(c) For every $t\in P^1-U$, $H_t$ meets $S$ transversely except at one
point, which is an ordinary quadratic singularity.

Through every point $x$ of $S$ not belonging to $\Delta$ there passes a
unique hyperplane $H_t$. The map assigning $t\in P^1$ to $x$ is a
rational map $a:S\dashrightarrow P^1$ with domain of definition
$S-(\Delta\cap S)$.

Consider the blow-up $g:S'\longrightarrow S$ of the closed subscheme
$S\cap\Delta$. The scheme $S'$ is regular, and the fibre $S'_{x_i}$ over
$x_i$ is the projective bundle defined by the tangent space of $S$ at
$x_i$, hence is isomorphic to $P^1$ (EGA IV 19.4.2, 19.4.4). Moreover,
the rational map $f=ag$ is everywhere defined (XVIII 3.1). Finally, to
each $x_i$ one can associate the section of $f$ that sends
$t\in P^1$ to the tangent vector at $x_i$ of $H_t\cap S$.

Let us verify conditions (i), (ii), and (iii). For each $t\in P^1$, the
curves $S'_t$ and $S_t=H_t\cap S$ are isomorphic. Indeed, the morphism
$S'_t\longrightarrow S_t$ is an isomorphism except possibly at the
points $x_i$, where $S_t$ is regular; moreover, since $S'_t$ is a
complete intersection in $S'$, it is integral, and is therefore
isomorphic to $S_t$.

The asserted properties of the fibres follow from (a)--(c). Finally,
$f$ is smooth at the points of $\sigma$, since $f$ is flat and the
fibres $S'_t$ are geometrically regular at $\sigma(t)$.

%%%%%%%%%%  scan idx 38  |  folio 27  %%%%%%%%%%

\underline{Corollary \SGARefTarget{sga7.sga7i.eii.n2-2}{2.2}. With the notation of \SGARefLink{sga7.sga7i.eii.n2-1}{2.1}, let $Y$ be a
normal-crossings divisor in $S$, and let $Y'=Y\times_S S'$. Then $f$,
$g$, and $U$ may be chosen so that $Y'\times_{P^1}U$ is a divisor with
normal crossings relative to $U$.}

It is known that, if a Lefschetz pencil exists for an embedding
$S\longrightarrow P^r$, the Lefschetz pencils form an open subset of the
variety of lines in $P^r$ (XVII 3.2.1). The same is true of the set of
Lefschetz pencils whose axis does not meet $Y$ and for which there is a
nonempty open subset $U_1$ of $P^1$ such that $Y\cap H_t$ is smooth for
$t\in U_1$. It is then enough to construct $f$ and $g$ from a pencil
satisfying these two conditions and replace $U$ by $U\cap U_1$.

\SGARefTarget{sga7.sga7i.eii.n2-3-0}{2.3.0}. Let $G$ be a profinite group. Recall that $G$ is said to be
topologically finitely generated if there is an integer $n$ such that
$G$ is isomorphic to a quotient of the free pro-group $F(n)$ on $n$
generators. Put $K=\Ker(F(n)\longrightarrow G)$. The group $G$ is said
to be topologically finitely presented if, in addition, there are
finitely many elements $k_1,\ldots,k_r$ of $K$ such that the smallest
closed normal subgroup of $K$ containing $k_1,\ldots,k_r$ is $K$.

\underline{Theorem \SGARefTarget{sga7.sga7i.eii.n2-3-1}{2.3.1}. Let $k$ be a separably closed field of
characteristic $p\geq 0$, and let $X$ be a connected scheme of finite
type over $k$. Suppose that schemes of finite type and dimension at most
$\dim X$ over an algebraic closure of $k$ admit strong resolution of
singularities (SGA 5 I 3.1.5). This condition is satisfied when
$\dim X=2$, by [1], or when $\operatorname{char}k=0$, by [2]. Then the
fundamental group $\Pi_1^{(p)}(X)$ is topologically finitely presented.}

By SGA 1 IX 4.10, we may suppose that $k$ is algebraically closed.

1) \underline{Reduction to the case in which $X$ is a regular surface.}

If $f:X'\longrightarrow X$ is a morphism of finite type, put
$X''=X'\times_X X'$, and let $(X'_a)_{a\in A}$ and
$(X''_b)_{b\in B}$ be the families of connected components of $X'$ and
$X''$, respectively. When $f$ is an effective descent morphism for the
category of finite \'etale coverings, the fundamental group $\Pi_1(X)$
can be computed explicitly

%%%%%%%%%%  scan idx 39  |  folio 28  %%%%%%%%%%

in terms of the fundamental groups $\Pi_1(X'_a)$ and $\Pi_1(X''_b)$.
It follows (SGA 1 IX 5.2, \SGARefLink{sga7.sga7i.eix.n5-3}{5.3}) that if the $\Pi_1(X'_a)$ are
topologically finitely generated, then so is $\Pi_1(X)$; and if the
$\Pi_1(X'_a)$ are topologically finitely presented while the
$\Pi_1(X''_b)$ are topologically finitely generated, then $\Pi_1(X)$ is
topologically finitely presented.

Since $X$ admits resolution of singularities, there is a regular scheme
$X'$ and a proper birational morphism $f:X'\longrightarrow X$. Since
$f$ is an effective descent morphism for finite \'etale coverings, it
suffices to prove the theorem under the additional assumption that $X$
is regular. One first deduces that $\Pi_1(X)$ is topologically finitely
generated; knowing that $\Pi_1(X')$ is topologically finitely presented
and that the $\Pi_1(X''_b)$ are topologically finitely generated then
shows that $\Pi_1(X)$ is topologically finitely presented.

Now let $(U_i)_{i\in I}$ be a finite affine open covering of $X$, and
let $f:\bigsqcup U_i\longrightarrow X$ be the canonical morphism. Since
$f$ is an effective descent morphism for finite \'etale coverings, it is
enough to prove the theorem for the $U_i$. We are thus reduced to the
case in which $X$ is affine and smooth.

Assume henceforth that $X$ is affine and smooth over $k$, with
$\dim X>2$. The Lefschetz theorem for quasi-projective schemes
([3] V 7.4) applies: if $Y$ is a sufficiently general hyperplane section
of $X$, then
\[
 \Pi_1(Y)\simeq\Pi_1(X).
\]
It therefore suffices to prove the theorem for $Y$, reducing us to the
case $\dim X=2$.

2) \underline{The case in which $X$ is a smooth surface over $k$.}

By the strong resolution theorem for surfaces, there exist a smooth,
integral, projective surface $S$ and an open immersion
$i:X\longrightarrow S$ such that $Y=S-X$ is a normal-crossings divisor.
We apply Corollary \SGARefLink{sga7.sga7i.eii.n2-2}{2.2} and retain its notation.

%%%%%%%%%%  scan idx 40  |  folio 29  %%%%%%%%%%

We show that it is enough to prove that the fundamental group of
$X'=X\times_S S'$ is topologically finitely presented. Since $X$ and
$X'$ are regular and
\[
 \operatorname{codim}\!\left(\bigcup_{1\leq i\leq n}\{x_i\},X\right)
 \geq 2,
\]
every \'etale covering of $X'$ induces on
$X-\bigcup_{1\leq i\leq n}\{x_i\}$ an \'etale covering which extends
uniquely to an \'etale covering of $X$ (SGA 2 1.11). This gives an
isomorphism
\[
 \Pi_1(X')\simeq\Pi_1(X).
\]

Consider the cartesian diagram
\[
\begin{tikzcd}[column sep=large,row sep=large]
X' \arrow[d,"h"'] & X'_U \arrow[l] \arrow[d,"h_U"] \\
P^1 & U \arrow[l]
\end{tikzcd}
\]
where $h=f|_{X'}$. Let $L$ be the set of primes different from $p$, and
let $\bar\eta$ be a geometric point over the generic point $\eta$ of
$P^1$. We verify the hypotheses of SGA 1 XIII 4.8. The morphism $h$ has
a section $\sigma$, because $f$ does. To prove condition (a) of
SGA 1 XIII 4.7, note first that, since the fibres of $h$ are
geometrically integral, $h$ is $0$-acyclic and locally $0$-acyclic
(SGA 4 XV 4.1, 1.16). It remains to show that, for every locally constant
sheaf of sets $F$ on $X'$, $(F,h)$ is cohomologically proper in
dimensions $\leq 0$. If $\xi$ is a closed point of $P^1$ and $Z$ denotes
the integral closure of $\Spec\mathcal O_{P^1,\xi}$ in $\bar\eta$, then
$X'_Z=X'\times_{P^1}Z$ satisfies condition $(S_2)$ at its closed points
and is regular at every other point, hence is normal. Consequently, an
\'etale covering of $X'_Z$ whose restriction to $X'_{\bar\eta}$ is
connected is necessarily connected. This proves the assertion, and
hence condition (a).

Since $X'_U$ is the complement in $S'_U$ of a divisor with normal
crossings relative to $U$, $h_U$ is locally $1$-aspherical for $L$
(SGA 4 XV 2.1); and, for every finite constant sheaf of $L$-groups $F$
on $X'_U$, $(F,h_U)$ is cohomologically proper in dimensions $\leq 1$
(SGA 1 XIII 2.4), which gives condition (b). Since $P^1$ is normal and
$T=P^1-U$ has codimension $1$, we have

%%%%%%%%%%  scan idx 41  |  folio 30  %%%%%%%%%%

\[
 \operatorname{prof\,\acute et}_T(P^1)\geq 2,
\]
which gives condition (c). The condition
\[
 \operatorname{prof\,hop}_x(X')\geq 3
\]
holds at every closed point of $X'$ by the purity theorem, which gives
condition (e). Finally, for every $t\in P^1$, $h$ is smooth at
$\sigma(t)$, and hence is locally $1$-aspherical for $L$ at $\sigma(t)$,
which gives condition (d).

By SGA 1 XIII 4.7, if $a$ is a geometric point of
$X'_{\bar\eta}$, there is an exact sequence
\[
 1\longrightarrow\Pi_1^{(p)}(X'_{\bar\eta},a)
 \longrightarrow\Pi'_1(X'_U,a)
 \longrightarrow\Pi_1(U,a)\longrightarrow 1,
\]
and the choice of the section $\sigma$ defines a splitting of this exact
sequence. Let $K$ be the image of $\Pi_1(U,a)$ in
$\Aut(\Pi_1^{(p)}(X'_{\bar\eta},a))$. The group of coinvariants under
$K$,
\[
 \Pi_1^{(p)}(X'_{\bar\eta},a)_K,
\]
is the quotient of $\Pi_1^{(p)}(X'_{\bar\eta},a)$ by the closed normal
subgroup generated by the elements $x^{-1}q(x)$, where
$x\in\Pi_1^{(p)}(X'_{\bar\eta},a)$ and $q\in K$. By loc. cit. (4.7.4),
there is an isomorphism
\[
 \Pi_1^{(p)}(X',a)\simeq
 \Pi_1^{(p)}(X'_{\bar\eta},a)_K.
\]

Write $P^1-U=\coprod_{1\leq i\leq n}\{t_i\}$. For each $i$, let
$\bar T_i$ be the strict localization of $P^1$ at a geometric point
$\bar t_i$ over $t_i$, and let $s_i$ be the generic point of $\bar T_i$.
An \'etale covering of $U$ whose restriction to every $s_i$ is trivial
extends to an \'etale covering of $P^1$, and is therefore trivial. For
each $i$, let $\overline{k(s_i)}$ be an algebraic closure of $k(s_i)$.
The preceding assertion says that $\Pi_1(U,a)$ is the smallest closed
normal subgroup containing the images of the Galois groups
$\Gal(\overline{k(s_i)}/k(s_i))$ in $\Pi_1(U,a)$. By I \SGARefLink{sga7.sga7i.ei.n5-3}{5.3}, the image
$K_i$ of $\Gal(\overline{k(s_i)}/k(s_i))$ in
$\Aut(\Pi_1^{(p)}(X'_{\bar\eta},a))$ is finite. Since $K$ is the
smallest closed normal subgroup containing all the $K_i$, and since
$\Pi_1^{(p)}(X'_{\bar\eta},a)$ is topologically finitely presented
(SGA 1 XIII 3.2), this proves that
$\Pi_1^{(p)}(X'_{\bar\eta},a)_K$, and hence also the isomorphic group
$\Pi_1^{(p)}(X)$, is topologically finitely presented.

%%%%%%%%%%  scan idx 42  |  folio 31  %%%%%%%%%%

\underline{Remark \SGARefTarget{sga7.sga7i.eii.n2-3-2}{2.3.2}.} Resolution of singularities is not needed to
prove Theorem \SGARefLink{sga7.sga7i.eii.n2-3-1}{2.3.1} when $X$ is the open complement of a
normal-crossings divisor in a projective scheme smooth over $k$.

\begin{center}\rule{1.5cm}{0.4pt}\end{center}

\begin{center}\underline{BIBLIOGRAPHY}\end{center}

\begin{enumerate}[label={[\arabic*]},leftmargin=3em]
\item S. Abhyankar, ``Local uniformization on algebraic surfaces over
ground fields of characteristic $p\neq0$,'' \emph{Ann. of Math.},
vol. 63, no. 3, 1956, pp. 491--526.

\item H. Hironaka, ``Resolution of singularities of an algebraic variety
over a field of characteristic zero,'' \emph{Ann. of Math.}, vol. 79,
1964.

\item M. Raynaud, \emph{Th\'eor\`eme de Lefschetz en cohomologie des
faisceaux coh\'erents et en cohomologie \'etale}, thesis, forthcoming.
\end{enumerate}

% END INCLUDED SGA7I_English_current_buildable_source_20260823/components/08_expose_II_complete.tex

\clearpage
\phantomsection\addcontentsline{toc}{section}{Exposé VI}
% BEGIN INCLUDED SGA7I_English_current_buildable_source_20260823/components/09_expose_VI_opening_through_1_1.tex
%%%%%%%%%%  scan idx 43  |  folio 32  %%%%%%%%%%

\begin{center}
\underline{\SGARefAliasTarget{sga7.sga7i.evi.nroot}Expos\'e VI}

\vspace{1em}
\underline{\underline{FORMAL DEFORMATION THEORY}}

\vspace{1em}
\underline{by D. S. Rim}
\end{center}

1. \underline{Formal existence theorem}

We recall the notion of a cofibered category (S.G.A. 1 IV). Let
$\underline C$ be a fixed category. A cofibered category $\underline A$
over $\underline C$ is, by definition, a category $\underline A$
together with a functor $P:\underline A\longrightarrow\underline C$
such that:

\underline{Ax. 1.} For every map $f$ in $\underline C$ and every object
$a$ in $\underline A$ for which $P(a)$ is the source of $f$, there is a
cocartesian $f$-morphism $\alpha:a\longrightarrow b$. Thus, for every
$f$-morphism $\alpha':a\longrightarrow b'$, there is a unique
$T$-morphism $\tau:b\longrightarrow b'$ in $\underline A$ such that
$\alpha'=\tau\alpha$, where $T$ is the target of $f$.

\underline{Ax. 2.} A composite of cocartesian morphisms in
$\underline A$ is again cocartesian.

The following properties of a cofibered category
$P:\underline A\longrightarrow\underline C$ follow immediately from
the definitions.

(Cancellation) Let
$a\overset{\alpha}{\underset{\beta}{\rightrightarrows}}a'$ be two
cocartesian maps such that $P(\alpha)=P(\beta)$. If there is a
cocartesian map $b\xrightarrow{\gamma}a$ such that
$\alpha\gamma=\beta\gamma$, then $\alpha=\beta$.

(Factorization) Given cocartesian maps
$a\xrightarrow{\alpha}a'$ and $a\xrightarrow{\beta}a''$, there is a
cocartesian map $a'\xrightarrow{\gamma}a''$ such that
$\gamma\alpha=\beta$ if and only if there is a map
$f:P(a')\longrightarrow P(a'')$ in $\underline C$ such that
$fP(\alpha)=P(\beta)$. Moreover, $\gamma$ is unique.

Let $\underline A$ be a cofibered category over $\underline C$. For each
object $S$ of $\underline C$, denote by $\underline A(S)$ the
subcategory of $\underline A$ whose objects are the objects $a$ for
which $P(a)=S$, and in which
\[
 \Hom_{\underline A(S)}(a,a')=
 \{\alpha\in\Hom_{\underline A}(a,a')\mid P(\alpha)=\id_S\}.
\]
A cofibered groupoid over $\underline C$ is a cofibered category
$\underline A$ over $\underline C$ such that $\underline A(S)$

%%%%%%%%%%  scan idx 44  |  folio 33  %%%%%%%%%%

is a groupoid for every object $S$ of $\underline C$. Recall that a
groupoid is a category in which every morphism is an isomorphism. Thus a
cofibered groupoid over $\underline C$ is a cofibered category
$P:\underline A\longrightarrow\underline C$ in which every map of
$\underline A$ is cocartesian. Furthermore, for every
$\alpha\in\operatorname{Fl}(\underline A)$, $\alpha$ is an isomorphism
in $\underline A$ if and only if $P(\alpha)$ is an isomorphism.

Throughout the rest of this section, fix a complete noetherian local ring
$\Lambda$ with residue field $k$, and let $\underline C_\Lambda$ be the
category of artinian local $\Lambda$-algebras with residue field $k$,
whose maps are local homomorphisms over $\Lambda$. From \SGARefLink{sga7.sga7i.evi.n1-2}{1.2} onward, a
cofibered groupoid $\underline A$ over $\underline C_\Lambda$ is always
assumed to satisfy $\underline A(k)=\{e\}$, where $\{e\}$ denotes the
category in which $\Hom_{\{e\}}(a,b)$ consists of exactly one element
for every two objects $a,b$ of $\{e\}$. A map
$a'\xrightarrow{\alpha}a$ in $\underline A$ will simply be called a
$\underline C_\Lambda$-map if $P(a')=P(a)$ and
$P(\alpha)=1_{P(a)}$. By abuse of language, a map
$a'\xrightarrow{\alpha}a$ in $\underline A$ will be called surjective
if $P(\alpha):P(a')\longrightarrow P(a)$ is surjective in
$\underline C_\Lambda$.

\underline{Ex. \SGARefTarget{sga7.sga7i.evi.n1-1.item-a}{\SGARefAliasTarget{sga7.sga7i.evi.n1-1}1.1}(a).} Let
$F:\underline C\longrightarrow(\text{categories})$ be a functor.
Define the category $\underline F$ as follows. An object of
$\underline F$ is a pair $(S,a)$ with
$S\in\mathrm{ob}(\underline C)$ and $a\in\mathrm{ob}F(S)$; a map
$(S,a)\longrightarrow(S',a')$ is a pair $(f,u')$, where
$f:S\longrightarrow S'$ is a map in $\underline C$ and
$u':F(f)a\longrightarrow a'$ is a map in $F(S')$. The composite of
\[
 (S,a)\xrightarrow{(f,u')}(S',a')
 \xrightarrow{(g,u'')}(S'',a'')
\]
is
\[
 (gf,\,u''\mathbin{\cdot}F(g)u'):(S,a)\longrightarrow(S'',a'').
\]
With respect to the functor
$P:\underline F\longrightarrow\underline C$ given by
$(S,a)\longmapsto S$, $\underline F$ becomes a cofibered category over
$\underline C$. If $F(S)$ is a groupoid for every
$S\in\mathrm{ob}(\underline C)$, then $\underline F$ is a cofibered
groupoid over $\underline C$. In particular, a functor
$F:\underline C\longrightarrow(\text{groups})$ yields a cofibered
group, and a functor $F:\underline C\longrightarrow(\mathrm{Ens})$
yields a cofibered discrete groupoid. A natural transformation
$\varphi:F\longrightarrow G$, where $F,G:\underline C\longrightarrow
(\text{categories})$, induces a cocartesian functor
$\varphi:\underline F\longrightarrow\underline G$ over $\underline C$.

\underline{Ex. \SGARefTarget{sga7.sga7i.evi.n1-1.item-b}{1.1}(b).} Let $X$ and $Y$ be schemes over $\Lambda$
together with a fixed

%%%%%%%%%%  scan idx 45  |  folio 34  %%%%%%%%%%

isomorphism
\[
 X\otimes_\Lambda k\xleftrightarrow{\ \tau\ }Y\otimes_\Lambda k.
\]
Then the functor
$\Isom_{X,Y}:\underline C_\Lambda\longrightarrow(\text{groupoids})$
given by
\[
 \Isom_{X,Y}(S)=
 \left\{X\otimes_\Lambda S\xleftrightarrow{\ \sigma\ }
 Y\otimes_\Lambda S\;\middle|\;
 \sigma\otimes_\Lambda k=\tau\right\}
\]
yields a cofibered groupoid
$P:\underline{\Isom}_{X,Y}\longrightarrow\underline C_\Lambda$.

\underline{Ex. \SGARefTarget{sga7.sga7i.evi.n1-1.item-c}{1.1}(c).} Let $X_0$ be a fixed scheme over $k$, and let
$\underline M_{X_0}$ be the category of deformations of $X_0$ over
$\underline C_\Lambda$ (see \S4). Thus an object of
$\underline M_{X_0}$ is a pair $(R,X)$, where
$R\in\underline C_\Lambda$ and $X$ is a deformation of $X_0$ to $R$.
Then the functor
$P:\underline M_{X_0}\longrightarrow\underline C_\Lambda$ given by
$(R,X)\longmapsto R$ defines a cofibered groupoid over
$\underline C_\Lambda$.

\underline{Ex. \SGARefTarget{sga7.sga7i.evi.n1-1.item-d}{1.1}(d).} For any groupoid $X$, denote by $[X]$ the
discrete groupoid consisting of the isomorphism classes of objects of
$X$. If $P:\underline F\longrightarrow\underline C_\Lambda$ is a
cofibered groupoid, we obtain a functor
$[\underline F]:\underline C_\Lambda\longrightarrow(\mathrm{Ens})$
given by $[\underline F](S)=[\underline F(S)]$, and hence a cofibered
discrete groupoid which, by abuse of notation, is again denoted by
$[\underline F]$.

% END INCLUDED SGA7I_English_current_buildable_source_20260823/components/09_expose_VI_opening_through_1_1.tex

% BEGIN INCLUDED SGA7I_English_current_buildable_source_20260823/components/10_expose_VI_1_2_through_1_3.tex
\underline{Definition \SGARefTarget{sga7.sga7i.evi.n1-2}{1.2}. A cofibered groupoid $\underline A$ over
$\underline C_\Lambda$ is called semi-homogeneous if the following
properties hold:}

(1) \underline{Every diagram}
\[
\begin{tikzcd}[row sep=large,column sep=large]
a'' \arrow[r,dashed] \arrow[d,dotted] & a' \arrow[d] \\
a_1 \arrow[r] & a
\end{tikzcd}
\]
\underline{of solid arrows in $\underline A$, where
$a'\longrightarrow a$ is surjective, can be completed to a commutative
diagram by including the dotted arrows.}

(2) \underline{Let}
\[
\begin{tikzcd}[column sep=large]
& R \\
R\underset{k}{\times}k[\varepsilon]
 \arrow[ur,"\pi_1"] \arrow[dr,"\pi_2"'] & \\
& k[\varepsilon]
\end{tikzcd}
\]
\underline{be the projection maps, where $k[\varepsilon]$ is the ring
of dual numbers over $k$, and let}
\[
\begin{tikzcd}[column sep=huge]
a' \arrow[dr,"\alpha"] & \\
& a \\
b' \arrow[ur,"\beta"'] &
\end{tikzcd}
\]
\underline{be $\pi_1$-maps in $\underline A$. Then there is a
$\underline C_\Lambda$-map $\gamma:a'\longrightarrow b'$ such that
$\alpha=\beta\gamma$ if and only if
$(\pi_2)_*(a')\simeq(\pi_2)_*(b')$.}

\underline{Remark \SGARefTarget{sga7.sga7i.evi.n1-3.item-a}{\SGARefAliasTarget{sga7.sga7i.evi.n1-3}1.3}(a).} A reformulation of the notion in terms of
$2$-categories will be given later in \SGARefLink{sga7.sga7i.evi.n2-6.item-b}{2.6}(b). If $\underline A$ is
semi-homogeneous, then so is $[\underline A]$. In particular, \SGARefLink{sga7.sga7i.evi.n1-2}{1.2}(2)
implies that the canonical map

%%%%%%%%%%  scan idx 46  |  folio 35  %%%%%%%%%%

\[
 [\underline A(R\times_k k[\varepsilon])]
 \longrightarrow
 [\underline A(R)]\times[\underline A(k[\varepsilon])]
\]
is bijective.

\underline{Remark \SGARefTarget{sga7.sga7i.evi.n1-3.item-b}{1.3}(b).} Consider a commutative diagram
\[
\begin{tikzcd}[row sep=large,column sep=large]
& a' \arrow[dl,"v_1"'] \arrow[dr,"v_2"] & \\
a_1 \arrow[dr,"u_1"'] & & a_2 \arrow[dl,"u_2"] \\
& a &
\end{tikzcd}
\]
in a cofibered category $\underline A$ over $\underline C_\Lambda$.
Put $R=P(a)$, $R'=P(a')$, and $R_i=P(a_i)$ for $i=1,2$. If $a''$ is
the direct image under the map
\[
 (P(v_1),P(v_2)):R'\longrightarrow R_1\times_RR_2,
\]
we obtain a commutative diagram
\[
\begin{tikzcd}[row sep=large,column sep=large]
& a' \arrow[d] \arrow[ddl,"v_1"'] \arrow[ddr,"v_2"] & \\
& a'' \arrow[dl] \arrow[dr] & \\
a_1 \arrow[dr,"u_1"'] & & a_2 \arrow[dl,"u_2"] \\
& a &
\end{tikzcd}
\]
Thus \SGARefLink{sga7.sga7i.evi.n1-2}{1.2}(1) may be replaced by:

(1)$'$ Every diagram
\[
\begin{tikzcd}[cramped,column sep=small,row sep=small]
a_1 \arrow[dr] & & a_2 \arrow[dl] \\
& a &
\end{tikzcd}
\qquad\text{in }\underline A\text{ above}\qquad
\begin{tikzcd}[cramped,column sep=small,row sep=small]
R_1 \arrow[dr] & & R_2 \arrow[dl] \\
& R &
\end{tikzcd}
\quad\text{in }\underline C_\Lambda,
\]
where $R_2\longrightarrow R$ is surjective, can be completed to a
commutative diagram
\[
\begin{tikzcd}[row sep=large,column sep=large]
& a' \arrow[dl] \arrow[dr] & \\
a_1 \arrow[dr] & & a_2 \arrow[dl] \\
& a &
\end{tikzcd}
\qquad\text{above}\qquad
\begin{tikzcd}[row sep=large,column sep=large]
& R_1\times_RR_2 \arrow[dl,"\pi_1"'] \arrow[dr,"\pi_2"] & \\
R_1 \arrow[dr] & & R_2 \arrow[dl] \\
& R &
\end{tikzcd}
\]
where $\pi_i$ ($i=1,2$) are the projections.

\underline{Remark \SGARefTarget{sga7.sga7i.evi.n1-3.item-c}{1.3}(c).} If $\underline A$ is a semi-homogeneous
cofibered groupoid over $\underline C_\Lambda$, then
$[\underline A(k[\varepsilon])]$ is canonically endowed with a vector
space structure over $k$, where $k[\varepsilon]$ is the ring of dual
numbers over $k$. Indeed, let $\underline{\underline V}$ be the category
of

%%%%%%%%%%  scan idx 47  |  folio 36  %%%%%%%%%%

finite-dimensional vector spaces over $k$. Every product-preserving
functor $H:\underline{\underline V}\longrightarrow(\mathrm{Ens})$
factors canonically as
\[
\begin{tikzcd}[row sep=large,column sep=large]
\underline{\underline V} \arrow[rr,"H"]
 \arrow[dr,"\overline H"'] & & (\mathrm{Ens}) \\
& \underline{\underline V} \arrow[ur] &
\end{tikzcd}
\]
where $\underline{\underline V}\longrightarrow(\mathrm{Ens})$ is the
forgetful functor. On the other hand, the functor
\[
 G:\underline{\underline V}\longrightarrow\underline C_\Lambda,
 \qquad W\longmapsto k[W]=k\oplus W,\qquad W^2=0,
\]
takes products in $\underline{\underline V}$ to fibre products over $k$
in $\underline C_\Lambda$. If $\underline A$ is a semi-homogeneous
cofibered groupoid over $\underline C_\Lambda$, then
\[
 [\underline A(k[V]\times_k k[W])]
 \longrightarrow
 [\underline A(k[V])]\times[\underline A(k[W])]
\]
is bijective by \SGARefLink{sga7.sga7i.evi.n1-3.item-a}{1.3}(a). Hence the functor
$V\longmapsto[\underline A(k[V])]$ commutes with products, proving the
assertion.

% END INCLUDED SGA7I_English_current_buildable_source_20260823/components/10_expose_VI_1_2_through_1_3.tex

% BEGIN INCLUDED SGA7I_English_current_buildable_source_20260823/components/11_expose_VI_1_4_through_1_6.tex
%%%%%%%%%%  scan idx 47  |  folio 36 (continued)  %%%%%%%%%%

\underline{Definition \SGARefTarget{sga7.sga7i.evi.n1-4}{1.4}.} (1) \underline{Let $R$ be an object of
$\underline C_\Lambda$. A small extension of $R$ in
$\underline C_\Lambda$ is a surjective map
$R'\xrightarrow{f}R$ in $\underline C_\Lambda$ such that
$(\ker f)^2=0$ and $\ker f\simeq k$. In other words,
$R'\xrightarrow{f}R$ is a small extension of $R$ if}
\[
0\longrightarrow k\xrightarrow{t}R'\xrightarrow{f}R
 \longrightarrow0
\]
\underline{is exact for some $t\in R'$ with $t^2=0$.}

(2) \underline{Let $\underline A$ be a cofibered groupoid over
$\underline C_\Lambda$, and let $a$ be an object of $\underline A$.
A small prolongation of $a$ is an arrow
$a'\xrightarrow{\alpha}a$ in $\underline A$ such that $P(\alpha)$ is
a small extension in $\underline C_\Lambda$. An arrow
$a'\xrightarrow{\alpha}a$ will be called versal for small
prolongations of $a$ if, for every small prolongation
$a'_1\longrightarrow a$, there exists a commutative diagram}
\[
\begin{tikzcd}[row sep=large,column sep=large]
a' \arrow[rr] \arrow[dr,"\alpha"'] & & a'_1 \arrow[dl] \\
& a &
\end{tikzcd}
\]

\underline{Proposition \SGARefTarget{sga7.sga7i.evi.n1-5}{1.5}. Let $\underline A$ be a semi-homogeneous
cofibered groupoid over $\underline C_\Lambda$.}

(1) \underline{Let}
\[
\begin{tikzcd}[column sep=large]
& R \\
R\times_k k[\varepsilon]
 \arrow[ur,"\pi_1"] \arrow[dr,"\pi_2"'] & \\
& k[\varepsilon]
\end{tikzcd}
\]
\underline{be the projection maps in $\underline C_\Lambda$, and let
$a'\xrightarrow{\alpha}a$ be a $\pi_1$-map. Then there exists a map
$a\xrightarrow{\beta}a'$ such that $\alpha\beta=1_a$ if and only if
there exists a map $a\longrightarrow(\pi_2)_*a'$.}

%%%%%%%%%%  scan idx 48  |  folio 37  %%%%%%%%%%

(2) \underline{Let $R'\xrightarrow{f}R$ be a small extension in
$\underline C_\Lambda$, and consider a commutative diagram}
\[
\begin{tikzcd}[row sep=large,column sep=small]
& a'' \arrow[dl,"\alpha'"'] \arrow[dr,"\alpha'_1"]
  & & & & & R'\times_RR' \arrow[dl] \arrow[dr] & \\
a' \arrow[dr,"\alpha"'] & & a' \arrow[dl,"\alpha_1"]
  & & \underline{\text{above}} & R' \arrow[dr] & & R' \arrow[dl] \\
& a & & & & & R &
\end{tikzcd}
\]
\underline{Define the canonical map
$\widetilde f:R'\times_RR'\longrightarrow k[\ker f]$ by}
\[
 \widetilde f(r'_1,r'_2)=r'_1(0)+(r'_1-r'_2).
\]
\underline{If there exists a map
$a'\longrightarrow(\widetilde f)_*(a'')$, then there exists a map
$\gamma:a'\longrightarrow a'$ such that
$\alpha=\alpha_1\gamma$.}

\underline{Proof.} (1) The ``only if'' part is clear. Assume conversely
that there exists a map
$\varphi:a\longrightarrow(\pi_2)_*a'$. Set
\[
 h:R\longrightarrow R\times_k k[\varepsilon],
 \qquad h=1\times P(\varphi).
\]
The composite $\pi_1h$ is the identity. We therefore obtain maps
\[
 a\mathrel{\mathop{\rightleftarrows}^{\gamma}_{\alpha_1}}h_*a,
 \qquad \alpha_1\gamma=1_a,
 \qquad P(\gamma)=h,
 \qquad P(\alpha_1)=\pi_1.
\]
On the other hand, $\pi_2h=P(\varphi)$, so there exists a
$\pi_2$-map
\[
 h_*a\longrightarrow(\pi_2)_*a'.
\]
Thus $(\pi_2)_*(h_*a)=(\pi_2)_*a'$, and \SGARefLink{sga7.sga7i.evi.n1-2}{1.2}(2) gives a
$\underline C_\Lambda$-map $u:h_*a\longrightarrow a'$ such that
$\alpha u=\alpha_1$. Define $\beta=u\gamma$. Then
\[
 \alpha\beta=\alpha u\gamma=\alpha_1\gamma=1_a.
\]

(2) The map
\[
 1\times\widetilde f:R'\times_RR'
 \longrightarrow R'\times_k k[\ker f]
\]
is an isomorphism compatible with the projections onto the first
factor, and $\pi_2(1\times\widetilde f)=\widetilde f$, where
$\pi_2:R'\times_k k[\ker f]\longrightarrow k[\ker f]$ is the second
projection. Since $R'\xrightarrow fR$ is a small extension,
$k[\ker f]\simeq k[\varepsilon]$. Hence, if there exists a map
\[
 a'\longrightarrow(\widetilde f)_*a''
   =(\pi_2)_*\bigl((1\times\widetilde f)_*a''\bigr),
\]
part (1) yields a map $\gamma':a'\longrightarrow a''$ such that
$\alpha'\gamma'=1_{a'}$. Define
$\gamma=\alpha'_1\gamma':a'\longrightarrow a'$. Then
$\alpha_1\gamma=\alpha$.

%%%%%%%%%%  scan idx 49  |  folio 38  %%%%%%%%%%

\underline{Corollary \SGARefTarget{sga7.sga7i.evi.n1-6}{1.6}. Let $\underline A$ be a semi-homogeneous
cofibered groupoid over $\underline C_\Lambda$, and consider a diagram}
\[
\begin{tikzcd}[column sep=large]
a' \arrow[dr,"\alpha"'] & & a_1 \arrow[dl,"\alpha_1"] \\
& a &
\end{tikzcd}
\]
\underline{in $\underline A$, where $\alpha_1$ is a small
prolongation. Assume that $a$ is versal for small prolongations of
$e$. Then there exists a map $\rho:a'\longrightarrow a_1$ such that
$\alpha=\alpha_1\rho$ if and only if there exists a map
$f:P(a')\longrightarrow P(a_1)$ such that
$P(\alpha)=P(\alpha_1)f$ in $\underline C_\Lambda$.}

\underline{Proof.} Assume that $P(\alpha)=P(\alpha_1)f$, where
$f:P(a')\longrightarrow P(a_1)$. Replacing $a'$ by its direct image
under $f$, we may and shall assume that
$P(\alpha)=P(\alpha_1)$. Put $R'=P(a_1)$ and $R=P(a)$. Since
$\underline A$ is semi-homogeneous, we obtain a commutative diagram
\[
\begin{tikzcd}[row sep=large,column sep=small]
& a'' \arrow[dl,"\beta"'] \arrow[dr,"\beta_1"]
  & & & & & R'\times_RR' \arrow[dl] \arrow[dr] & \\
a' \arrow[dr,"\alpha"'] & & a_1 \arrow[dl,"\alpha_1"]
  & & \text{above} & R' \arrow[dr] & & R' \arrow[dl] \\
& a & & & & & R &
\end{tikzcd}
\]
Since $a$ is versal for small prolongations of $e$, so is $a'$;
the assertion therefore follows from \SGARefLink{sga7.sga7i.evi.n1-5}{1.5}(2).


% END INCLUDED SGA7I_English_current_buildable_source_20260823/components/11_expose_VI_1_4_through_1_6.tex

% BEGIN INCLUDED SGA7I_English_current_buildable_source_20260823/components/12_expose_VI_1_7_through_1_8.tex
%%%%%%%%%%  scan idx 49  |  folio 38 (continued)  %%%%%%%%%%

\underline{Definition \SGARefTarget{sga7.sga7i.evi.n1-7}{1.7}. For each object $R$ of
$\underline C_\Lambda$, put}
\[
 \overline\Omega_R=M/(m_\Lambda R+M^2),
\]
\underline{where $m_\Lambda$ and $M$ are the maximal ideals of
$\Lambda$ and $R$, respectively. This is a finite-dimensional vector
space over $k$, and $R\longmapsto\overline\Omega_R$ defines a functor}
\[
 \overline\Omega:\underline C_\Lambda
   \longrightarrow\underline{\underline V},
\]
\underline{where $\underline{\underline V}$ is the category of
finite-dimensional vector spaces over $k$. If $\underline A$ is a
cofibered category over $\underline C_\Lambda$, the composite functor
$\underline A\xrightarrow P\underline C_\Lambda
\longrightarrow\underline{\underline V}$ is again denoted
$\overline\Omega$, by abuse of notation.}

\underline{Proposition \SGARefTarget{sga7.sga7i.evi.n1-8}{1.8}. Let $\underline A$ be a semi-homogeneous
cofibered groupoid over $\underline C_\Lambda$. Given an object $a$ of
$\underline A$ that is versal for small prolongations of $e$, there
exists an arrow $a'\longrightarrow a$ such that:}

%%%%%%%%%%  scan idx 50  |  folio 39  %%%%%%%%%%

(i) \underline{$a'$ is versal for small prolongations of $a$;}

(ii) \underline{$P(a')\longrightarrow P(a)$ is surjective, and its
kernel is annihilated by the maximal ideal of $P(a)$;}

(iii) \underline{$\overline\Omega_{a'}\longrightarrow
\overline\Omega_a$ is an isomorphism.}

\underline{Proof.} Put $R=P(a)$, and consider an exact sequence
\[
 0\longrightarrow J\longrightarrow S\longrightarrow R
 \longrightarrow0,
\]
where $S$ is a formal power-series ring over $\Lambda$ in finitely
many variables and $J\subset M^2$, with $M$ the maximal ideal of $S$.
Let $\Sigma$ be the set of ideals $J_1$ of $S$ such that

(i) $J\supset J_1\supset MJ$, and

(ii) there exists a $\pi_1$-morphism $a_1\longrightarrow a$ in
$\underline A$, where $\pi_1:S/J_1\longrightarrow R$ is the canonical
map.

Property \SGARefLink{sga7.sga7i.evi.n1-2}{1.2}(1) of $\underline A$, together with the fact that $J/MJ$
is a finite-dimensional vector space over $k$, implies that $\Sigma$
has a smallest ideal. Denote it by $J'$, and choose a
$\pi'$-morphism $a'\xrightarrow g a$, where
$\pi':S/J'\longrightarrow R$ is the canonical map. Since
$J'\subset J\subset M^2$, the map
$\overline\Omega_{a'}\longrightarrow\overline\Omega_a$ is an
isomorphism.

We claim that $a'\xrightarrow g a$ is versal for small prolongations
of $a$. Indeed, let $a_1\longrightarrow a$ be a small prolongation,
and put $R_1=P(a_1)$. We obtain an exact commutative diagram
\[
\begin{tikzcd}[column sep=small]
0 \arrow[r] & J/MJ \arrow[r] \arrow[d,"h"]
 & S/MJ \arrow[r] \arrow[d,"h_1"]
 & R \arrow[r] \arrow[d,equal] & 0 \\
0 \arrow[r] & k \arrow[r] & R_1 \arrow[r] & R \arrow[r] & 0.
\end{tikzcd}
\]
If $h=0$, then $\pi_1:R_1\longrightarrow R$ admits a cross-section,
so $h_1$ induces a map $f:S/J'\longrightarrow R_1$. If $h\ne0$,
then $h_1:S/MJ\longrightarrow R_1$ is surjective and hence, by the
definition of $J'$, again induces $f:S/J'\longrightarrow R_1$.
In either case we obtain a commutative diagram
\[
\begin{tikzcd}[column sep=large]
S/J' \arrow[rr,"f"] \arrow[dr,"\pi'"']
 & & R_1 \arrow[dl,"\pi_1"] \\
& R &
\end{tikzcd}
\]
Since $a$ is versal for small prolongations of $e$, Corollary \SGARefLink{sga7.sga7i.evi.n1-6}{1.6}
gives a commutative diagram

%%%%%%%%%%  scan idx 51  |  folio 40  %%%%%%%%%%

\[
\begin{tikzcd}[column sep=large]
a' \arrow[rr,"\rho"] \arrow[dr] & & a_1 \arrow[dl] \\
& a &
\end{tikzcd}
\]
as required.


% END INCLUDED SGA7I_English_current_buildable_source_20260823/components/12_expose_VI_1_7_through_1_8.tex

% BEGIN INCLUDED SGA7I_English_current_buildable_source_20260823/components/13_expose_VI_pro_categories_through_1_10.tex
%%%%%%%%%%  scan idx 51  |  folio 40 (continued)  %%%%%%%%%%

Let $P:\underline A\longrightarrow\underline C_\Lambda$ be a
cofibered category in groupoids. It then gives a cofibered category in
groupoids
\[
 \operatorname{pro}(P):\operatorname{pro}(\underline A)
 \longrightarrow\operatorname{pro}(\underline C_\Lambda)
\]
(for the definition of the category of pro-objects
$\operatorname{pro}(\underline A)$, see A. Grothendieck,
S\'eminaire Bourbaki 195), fitting into the commutative diagram
\[
\begin{tikzcd}[column sep=large,row sep=large]
\underline A \arrow[r] \arrow[d,"P"']
 & \operatorname{pro}(\underline A)
   \arrow[d,"\operatorname{pro}(P)"] \\
\underline C_\Lambda \arrow[r]
 & \operatorname{pro}(\underline C_\Lambda).
\end{tikzcd}
\]

Let $\widehat{\underline C}_\Lambda$ be the full subcategory of
$\operatorname{pro}(\underline C_\Lambda)$ whose objects are systems
$(R_i)_{i\in\mathbb N}$ with $R_i\in\underline C_\Lambda$ such that
each map $\varphi_i:R_i\longrightarrow R_{i-1}$ is surjective and,
for all sufficiently large $i$, induces an isomorphism
\[
 \frac{m_i}{m_\Lambda R_i+m_i^2}
 \xrightarrow{\ \sim\ }
 \frac{m_{i-1}}{m_\Lambda R_{i-1}+m_{i-1}^2},
\]
where $\mathbb N$ denotes the ordered category of natural numbers and
$m_i$ is the maximal ideal of $R_i$.

Let $\widehat{\underline A}$ be the full subcategory of
$\operatorname{pro}(\underline A)$ defined by
\[
 \operatorname{ob}(\widehat{\underline A})=
 \left\{a\in\operatorname{ob}(\operatorname{pro}(\underline A))
 \mathrel{\Big|}
 \operatorname{pro}(P)(a)\in
 \operatorname{ob}(\widehat{\underline C}_\Lambda)\right\}.
\]
We obtain a cofibered category in groupoids
$\widehat P:\widehat{\underline A}\longrightarrow
\widehat{\underline C}_\Lambda$ and a commutative diagram
\[
\begin{tikzcd}[column sep=large,row sep=large]
\underline A \arrow[r,hook] \arrow[d,"P"']
 & \widehat{\underline A} \arrow[d,"\widehat P"] \\
\underline C_\Lambda \arrow[r,hook]
 & \widehat{\underline C}_\Lambda.
\end{tikzcd}
\]
The category $\widehat{\underline C}_\Lambda$ is canonically
equivalent to the category of complete local $\Lambda$-algebras of
formally finite type over $\Lambda$ with fixed residue field $k$; we
henceforth make this identification.

\underline{Definition \SGARefTarget{sga7.sga7i.evi.n1-9}{1.9}. Let $\underline A$ be a cofibered groupoid
over $\underline C_\Lambda$. An object $\xi$ of
$\widehat{\underline A}$ is called quasi-initial if, for every object
$a$ of $\widehat{\underline A}$, there exists a map $\xi\to a$ in
$\widehat{\underline A}$. A quasi-initial object $\xi$ is called
minimal if the morphism
$\xi:h_R\longrightarrow[\widehat{\underline A}]$ induces a bijection}

%%%%%%%%%%  scan idx 52  |  folio 41  %%%%%%%%%%

\[
 \xi(k[\varepsilon]):
 \operatorname{Hom}_{\Lambda\text{-alg}}(R,k[\varepsilon])
 \xrightarrow{\ \sim\ }
 [\underline A(k[\varepsilon])],
 \qquad R=P(\xi).
\]
\underline{An object $\xi$ of $\widehat{\underline A}$ is called a
versal formal object of $\underline A$ if every surjective map
$\alpha:a'\to a$ in $\underline A$ induces a surjective map}
\[
 \operatorname{Hom}_{\widehat{\underline A}}(\xi,a')
 \longrightarrow
 \operatorname{Hom}_{\widehat{\underline A}}(\xi,a).
\]
\underline{A versal formal object $\xi$ is called minimal if
$\xi:h_R\longrightarrow[\widehat{\underline A}]$ induces the above
bijection on $k[\varepsilon]$, where $R=P(\xi)$.}

\underline{Proposition \SGARefTarget{sga7.sga7i.evi.n1-10}{1.10}. Let $\underline A$ be a cofibered
groupoid over $\underline C_\Lambda$.}

(i) \underline{Every (minimally) versal formal object of
$\underline A$ is a quasi-initial (minimal) object of
$\widehat{\underline A}$.}

(ii) \underline{Suppose that $\xi$ is an object of
$\widehat{\underline A}$ for which}
\[
 \widetilde\xi(k[\varepsilon]):
 \operatorname{Hom}_{\Lambda\text{-alg}}(R,k[\varepsilon])
 \longrightarrow[\underline A(k[\varepsilon])]
\]
\underline{is injective, where $P(\xi)=R$. Then every endomorphism of
$\xi$ in $\widehat{\underline A}$ is an automorphism. In particular,
a quasi-initial minimal object of $\widehat{\underline A}$, if it
exists, is unique up to noncanonical isomorphism.}

\underline{Proof.} (i) Let $\xi$ be a versal formal object of
$\underline A$. For an object $a$ of $\widehat{\underline A}$, choose
a presentation by surjective maps in $\underline A$,
\[
 \cdots\longrightarrow a_n\xrightarrow{\alpha_n}a_{n-1}
 \longrightarrow\cdots\longrightarrow
 a_1\xrightarrow{\alpha_1}a_0,
 \qquad a=\varprojlim_n a_n.
\]
Suppose maps $\varphi_i:\xi\longrightarrow a_i$ have been chosen with
$\alpha_i\varphi_i=\varphi_{i-1}$ for all $i<n$. Since $\xi$ is
versal and $\alpha_n$ is surjective, there is a map
$\varphi_n:\xi\longrightarrow a_n$ with
$\alpha_n\varphi_n=\varphi_{n-1}$. Thus
$\varphi=\varprojlim_n\varphi_n$ defines a map $\xi\longrightarrow a$.

(ii) Let $\alpha:\xi\longrightarrow\xi$ be an endomorphism in
$\widehat{\underline A}$. For every
$\lambda\in\operatorname{Hom}_{\Lambda\text{-alg}}
(R,k[\varepsilon])$, we have
\[
 \lambda_*\xi\simeq
 (\lambda\circ P(\alpha))_*\xi.
\]
Injectivity therefore gives
$\lambda=\lambda\circ P(\alpha)$ for every $\lambda$. It follows that
$P(\alpha):R\longrightarrow R$ is surjective. Every surjective
endomorphism of a noetherian ring is an automorphism, so $P(\alpha)$,
and hence $\alpha$, is an automorphism.


% END INCLUDED SGA7I_English_current_buildable_source_20260823/components/13_expose_VI_pro_categories_through_1_10.tex

% BEGIN INCLUDED SGA7I_English_current_buildable_source_20260823/components/14_expose_VI_1_11_through_1_14.tex
%%%%%%%%%%  scan idx 52  |  folio 41 (continued)  %%%%%%%%%%

\underline{Theorem \SGARefTarget{sga7.sga7i.evi.n1-11}{1.11} (Schlessinger). Let $\underline A$ be a
cofibered groupoid over $\underline C_\Lambda$. Then $\underline A$
admits a minimally versal formal object if and only if:}

%%%%%%%%%%  scan idx 53  |  folio 42  %%%%%%%%%%

(i) \underline{$\underline A$ is semi-homogeneous;}

(ii) \underline{$[\underline A(k[\varepsilon])]$ is a
finite-dimensional vector space over $k$ (cf. \SGARefLink{sga7.sga7i.evi.n1-3.item-c}{1.3}(c)).}

\underline{Proof. Necessity.} Let $\xi\in\widehat{\underline A}$ be a
minimally versal formal object of $\underline A$, and consider a
diagram
\[
\begin{tikzcd}[column sep=large,row sep=large]
a_1 \arrow[dr] & & a_2 \arrow[dl,"\alpha"] \\
& a &
\end{tikzcd}
\]
in $\underline A$, where $\alpha$ is surjective. Since $\xi$ is
quasi-initial, there is a map $\xi\longrightarrow a_1$. Since
$\alpha$ is surjective, the composite
$\xi\longrightarrow a_1\longrightarrow a$ lifts to $a_2$; hence we
obtain a commutative diagram
\[
\begin{tikzcd}[column sep=large,row sep=large]
& \xi \arrow[dl] \arrow[dr] & \\
a_1 \arrow[dr] & & a_2 \arrow[dl,"\alpha"] \\
& a &
\end{tikzcd}
\]
which proves condition \SGARefLink{sga7.sga7i.evi.n1-2}{1.2}(1).

We now prove condition \SGARefLink{sga7.sga7i.evi.n1-2}{1.2}(2). Let
\[
\begin{tikzcd}[column sep=large]
& S \\
S\times_k k[\varepsilon]
 \arrow[ur,"\pi_1"] \arrow[dr,"\pi_2"'] & \\
& k[\varepsilon]
\end{tikzcd}
\]
be the projections, and let
\[
\begin{tikzcd}[column sep=large]
a' \arrow[dr,"\alpha"'] & & b' \arrow[dl,"\beta"] \\
& a &
\end{tikzcd}
\]
be $\pi_1$-maps in $\underline A$. By the definition of $\xi$, there
is a commutative diagram
\[
\begin{tikzcd}[column sep=large,row sep=large]
& \xi \arrow[dl,"u"'] \arrow[dr,"v"] & \\
a' \arrow[dr,"\alpha"'] & & b' \arrow[dl,"\beta"] \\
& a &
\end{tikzcd}
\]
and therefore $\pi_1P(u)=\pi_1P(v)$. Suppose now that
$(\pi_2)_*a'\simeq(\pi_2)_*b'$. Minimality of $\xi$ gives
$\pi_2P(u)=\pi_2P(v)$, whence

%%%%%%%%%%  scan idx 54  |  folio 43  %%%%%%%%%%

\[
 P(u)=\pi_1P(u)\times\pi_2P(u)
     =\pi_1P(v)\times\pi_2P(v)=P(v).
\]
Since $\underline A$ is a cofibered groupoid, there is a unique
isomorphism $\sigma:a'\xrightarrow{\sim}b'$ in
$\underline A(S\times_k k[\varepsilon])$ such that $\sigma u=v$.
The equality $\beta\sigma u=\alpha u$ then implies
$\beta\sigma=\alpha$.

\underline{Sufficiency.} Let $c_1,\ldots,c_n$ be objects of
$\underline A(k[\varepsilon])$ whose classes form a $k$-basis of
$[\underline A(k[\varepsilon])]$. Choose
\[
 a_1\in\underline A(
 k[\varepsilon]\times_k\cdots\times_k k[\varepsilon])
\]
such that $(\pi_j)_*a_1\simeq c_j$. Then
\[
 \widetilde a_1(k[\varepsilon]):
 \operatorname{Hom}_{\Lambda\text{-alg}}(R_1,k[\varepsilon])
 \longrightarrow[\underline A(k[\varepsilon])]
\]
is bijective, where $R_1=P(a_1)$. In particular,
$a_1\longrightarrow e$ is versal for small prolongations of $e$.
Choose $a_2\longrightarrow a_1$ versal for small prolongations of
$a_1$ such that
$\overline\Omega_{a_2}\longrightarrow\overline\Omega_{a_1}$ is
bijective. Continue inductively, choosing
$a_{n+1}\longrightarrow a_n$ versal for small prolongations and
inducing a bijection on $\overline\Omega$, and set
\[
 \xi=\varprojlim_n a_n.
\]
Then $\overline\Omega_\xi\longrightarrow\overline\Omega_{a_1}$ is
bijective, so
\[
 \xi(k[\varepsilon]):
 \operatorname{Hom}_{\Lambda\text{-alg}}(R,k[\varepsilon])
 \xrightarrow{\ \sim\ }
 [\underline A(k[\varepsilon])]
\]
is bijective.

It remains to prove versality. Consider a diagram
\[
\begin{tikzcd}
& a' \arrow[d,"\alpha"] \\
\xi \arrow[r,"h"] & a
\end{tikzcd}
\]
where $\alpha$ is surjective. We must lift $h$ to $a'$. By induction
we may assume that $a'$ is a small prolongation of $a$. Since $P(a)$
is artinian, $h$ factors through some $a_q$,
\[
 \xi\longrightarrow a_q\longrightarrow a.
\]
Semi-homogeneity gives a commutative diagram
\[
\begin{tikzcd}
& a'' \arrow[r] \arrow[d,"\alpha'"] & a' \arrow[d,"\alpha"] \\
\xi \arrow[r] & a_q \arrow[r] & a,
\end{tikzcd}
\]
where $\alpha'$ is also a small prolongation. Since
$a_{q+1}\longrightarrow a_q$ is versal for small prolongations of
$a_q$, there is a commutative diagram
\[
\begin{tikzcd}
& & a'' \arrow[d,"\alpha'"] \\
\xi \arrow[r] & a_{q+1} \arrow[r] \arrow[ur] & a_q,
\end{tikzcd}
\]
which supplies the required lift.

\underline{Remark \SGARefTarget{sga7.sga7i.evi.n1-12}{1.12}.} Let

%%%%%%%%%%  scan idx 55  |  folio 44  %%%%%%%%%%

\[
 \cdots\longrightarrow a_n\longrightarrow a_{n-1}
 \longrightarrow\cdots\longrightarrow a_2\longrightarrow a_1
 \longrightarrow a_0=e
\]
be an infinite sequence of arrows in $\underline A$ such that
\[
\left\{
\begin{aligned}
&a_i\text{ is versal for small prolongations of }a_{i-1},\\
&P(a_i)\longrightarrow P(a_{i-1})\text{ is surjective},\\
&\overline\Omega_{a_i}\longrightarrow
 \overline\Omega_{a_{i-1}}\text{ is bijective for }i\gg0.
\end{aligned}
\right.
\]
The preceding proof shows that $\varprojlim_i a_i$ is a versal formal
object of $\underline A$.

\underline{Remark \SGARefTarget{sga7.sga7i.evi.n1-13}{1.13}.} Let $\underline A$ be a cofibered groupoid
over $\underline C_\Lambda$. If $\underline A$ satisfies \SGARefLink{sga7.sga7i.evi.n1-11}{1.11}(i) and
(ii), then so does $[\underline A]$; moreover,
$\xi\in\operatorname{ob}(\widehat{\underline A})$ is a minimally
versal formal object of $\underline A$ if and only if it is so for
$[\underline A]$, by 1.10(ii).

\underline{Remark \SGARefTarget{sga7.sga7i.evi.n1-14}{1.14}.} Let $\underline A$ be a cofibered groupoid
over $\underline C_\Lambda$ admitting a minimally versal formal
object $\xi$. Any invariant of $R=P(\xi)$, a complete local
$\Lambda$-algebra of formally finite type over $\Lambda$, is
automatically an invariant of $\underline A$ over
$\underline C_\Lambda$. Thus, for example, we say that $\underline A$
is (formally) smooth over $\underline C_\Lambda$ if $R$ is
(formally) smooth over $\Lambda$, that is, a formal power-series
algebra over $\Lambda$; and we say that $\underline A$ is irreducible over
$\underline C_\Lambda$ if $\operatorname{Spec}(R)$ is irreducible.

In the next section we introduce smoothness for a functor
$\underline A\longrightarrow\underline B$ over
$\underline C_\Lambda$, where $\underline A$ and $\underline B$ are
cofibered categories in groupoids over $\underline C_\Lambda$.
It follows from \SGARefLink{sga7.sga7i.evi.n2-4}{2.4} that $\underline A$ being (formally) smooth over
$\underline C_\Lambda$ is equivalent to the functor
$P:\underline A\longrightarrow\underline C_\Lambda$ being smooth,
so no ambiguity arises. We define
\[
 \dim_{\underline C_\Lambda}\underline A
   =\dim(k\otimes_\Lambda R).
\]

For this, let $\underline C_k$ be the category of artinian local
$k$-algebras with residue field $k$. The functor
$\underline C_k\longrightarrow\underline C_\Lambda$ is fully
faithful, and its left adjoint is
\[
 \underline C_\Lambda\longrightarrow\underline C_k,
 \qquad R\longmapsto k\otimes_\Lambda R.
\]
Let $\underline A|_{\underline C_k}$ be the pullback of
$P:\underline A\longrightarrow\underline C_\Lambda$ along
$\underline C_k\longrightarrow\underline C_\Lambda$. If $\xi$ is a
minimally versal formal object for $\underline A$ over
$\underline C_\Lambda$, then $k\otimes_\Lambda\xi$ is minimally
versal for $\underline A|_{\underline C_k}$ over $\underline C_k$;
the arrow $\xi\longrightarrow k\otimes_\Lambda\xi$ is cocartesian
over $R\longrightarrow k\otimes_\Lambda R$. Consequently,
\[
 \dim_{\underline C_\Lambda}\underline A
 =\dim_{\underline C_k}(\underline A|_{\underline C_k}).
\]

% END INCLUDED SGA7I_English_current_buildable_source_20260823/components/14_expose_VI_1_11_through_1_14.tex

% BEGIN INCLUDED SGA7I_English_current_buildable_source_20260823/components/15_expose_VI_1_15_through_1_17.tex
%%%%%%%%%%  scan idx 56  |  folio 45  %%%%%%%%%%

\textbf{\SGARefTarget{sga7.sga7i.evi.n1-15}{1.15}.} In the remainder of this paragraph, we briefly explain
how to modify the preceding theory to allow residue-field extensions.

Let $\Lambda$ and $K$ be complete noetherian local rings, with $K$ a
$\Lambda$-algebra. Assume that:

(a) the natural map $\Lambda\longrightarrow K$ is a local
homomorphism;

(b) the residue field $K'$ of $K$ is an extension of finite type of
the residue field $k$ of $\Lambda$.

The most interesting case is $K=K'$ with $K$ a finite inseparable
extension of $k$.

Denote by $\widehat C_{\Lambda,K}$ the category whose objects are
complete noetherian local rings $R$ equipped with a
$\Lambda$-algebra structure and a $\Lambda$-epimorphism
$s:R\longrightarrow K$. Let $C_{\Lambda,K}$ be its full subcategory
consisting of pairs $(R,s)$ for which $\ker(s)$ has finite length as
an $R$-module. The category $\widehat C_{\Lambda,K}$ identifies with
a full subcategory of $\operatorname{Pro}(C_{\Lambda,K})$. If $K$ is
artinian, then every object of $C_{\Lambda,K}$ is artinian.

For a finite-dimensional vector space $V$ over $K'$, write $D_V(K)$
for the algebra of dual numbers over $K$ defined by $V$. Its elements
have the form $x+v\varepsilon$, where $x\in K$, $v\in V$, and
$\varepsilon^2=0$, with multiplication
\[
 (x+v\varepsilon)(x'+v'\varepsilon)
 =xx'+(xv'+x'v)\varepsilon.
\]
The map
\[
 s:D_V(K)\longrightarrow K,
 \qquad s(x+v\varepsilon)=x,
\]
makes $D_V(K)$ an object of $C_{\Lambda,K}$. The role previously
played by $k[\varepsilon]$ is now played by $D_{K'}(K)$.

As in \SGARefLink{sga7.sga7i.evi.n1-8}{1.8}, every cofibered groupoid $\underline A$ over
$C_{\Lambda,K}$ extends naturally to a cofibered category
$\widehat{\underline A}$ over $\widehat C_{\Lambda,K}$. A map in
$\underline A$ or $\widehat{\underline A}$ is called
\underline{surjective} when the corresponding algebra homomorphism
in $C_{\Lambda,K}$ or $\widehat C_{\Lambda,K}$ is surjective.

%%%%%%%%%%  scan idx 57  |  folio 46  %%%%%%%%%%

\underline{We assume that $\underline A(K)=\{e\}$.}

\underline{Definition \SGARefTarget{sga7.sga7i.evi.n1-16}{1.16}. A cofibered groupoid $\underline A$ over
$C_{\Lambda,K}$ is called semi-homogeneous if the following
properties hold:}

(I) \underline{Every diagram of solid arrows}
\[
\begin{tikzcd}[row sep=large,column sep=large]
a'' \arrow[r,dashed] \arrow[d,dashed]
 & a' \arrow[d] \\
a_1 \arrow[r] & a
\end{tikzcd}
\]
\underline{in $\underline A$, where $a'\longrightarrow a$ is
surjective, can be completed to a commutative diagram by including the
dotted arrows.}

(II) \underline{For every $R$, the natural map}
\[
 [\underline A(R\times_KD_{K'}(K))]
 \longrightarrow
 [\underline A(R)]\times[\underline A(D_{K'}(K))]
\]
\underline{is bijective.}

If $\underline A$ is semi-homogeneous, then so is $[\underline A]$.

\textbf{\SGARefTarget{sga7.sga7i.evi.n1-17}{1.17}.} First restrict $\underline A$ to the full subcategory
of $C_{\Lambda,K}$ whose objects are the $D_V(K)$. Put
\[
 \Omega=\Omega^1_{K/\Lambda}\otimes_KK',
 \qquad d:K\longrightarrow\Omega.
\]
A map $f:D_V(K)\longrightarrow D_W(K)$ can be written
\[
 f(x+v\varepsilon)
 =x+\bigl(u(v)+D(dx)\bigr)\varepsilon,
\]
with $u\in\operatorname{Hom}_{K'}(V,W)$ and
$D\in\operatorname{Hom}_{K'}(\Omega,W)$. Identifying $f$ with
$(u,D)$, composition is given by
\[
 (u,D)\circ(u',D')=(uu',D+uD').
\]
Condition (II) implies:

(II$'$) the functor from $K'$-vector spaces to sets
\[
 V\longmapsto[\underline A(D_V(K))]
\]
commutes with products.


% END INCLUDED SGA7I_English_current_buildable_source_20260823/components/15_expose_VI_1_15_through_1_17.tex

% BEGIN INCLUDED SGA7I_English_current_buildable_source_20260823/components/16_expose_VI_lemma_1_18.tex
%%%%%%%%%%  scan idx 58  |  folio 47  %%%%%%%%%%

\underline{Lemma \SGARefTarget{sga7.sga7i.evi.n1-18}{1.18}. Assume that condition (II$'$) is satisfied by
$\underline A$. Then:}

(1) $[\underline A(D_{K'}(K))]$ carries a unique vector-space
structure such that, functorially in $V$,
\[
 [\underline A(D_V(K))]
 \simeq
 V\otimes_{K'}[\underline A(D_{K'}(K))].
\]

(2) There is a linear map
\[
 \gamma:\operatorname{Hom}(\Omega,K')
 \longrightarrow [\underline A(D_{K'}(K))]
\]
such that, for every $D\in\operatorname{Hom}(\Omega,K')$, the map
\[
 [\underline A((1,D))]:
 [\underline A(D_{K'}(K))]
 \longrightarrow [\underline A(D_{K'}(K))]
\]
is $x\longmapsto x+\gamma(D)$.

(3) For every
$(u,D):D_V(K)\longrightarrow D_W(K)$, the induced map
\[
 (u,D):[\underline A(D_V(K))]
 \longrightarrow[\underline A(D_W(K))],
\]
that is,
\[
 (u,D):
 V\otimes[\underline A(D_{K'}(K))]
 \longrightarrow
 W\otimes[\underline A(D_{K'}(K))],
\]
is given by
\[
 x\longmapsto u(x)+\gamma(D).
\]

Introduce the notation
\[
\begin{aligned}
 F(V)&=[\underline A(D_V(K))],\\
 A&=F(K'),\\
 F(u,D):F(V)&\longrightarrow F(W)
   &&\text{for the map induced by }
     (u,D):D_V(K)\longrightarrow D_W(K),\\
 F(D):F(V)&\longrightarrow F(V),
   &F(D)&=F(1,D).
\end{aligned}
\]
By (II$'$), and because $K'$ is a $K'$-vector-space object in the
category of $K'$-vector spaces, $F(K')$ is a $K'$-vector-space object
in the category of sets, hence a $K'$-vector space. This proves (1).
We henceforth identify $F(V)$ with $V\otimes A$.

The identity
\[
 (u,D)=(1,D)\circ(u,0)
\]
gives
\[
 F(u,D)=F(D)\circ(u\otimes1_A).
\]
The functoriality of $F$ amounts to the identities
\[
\begin{aligned}
 F(0)&=\operatorname{Id},\\
 F(D)\circ F(D')&=F(D+D'),\\
 (u\otimes1)\circ F(D)
   &=F(uD)\circ(u\otimes1_A).
\end{aligned}
\]

For any $D:\Omega\longrightarrow V$ and any $n\in\mathbb N$, let
\[
 i:V\lhook\joinrel\longrightarrow V\oplus K'^{\,n},
 \qquad
 p:V\oplus K'^{\,n}\longrightarrow K'^{\,n}
\]
be the canonical inclusion and projection. Then $piD=0$, and the
following diagram is commutative:

%%%%%%%%%%  scan idx 59  |  folio 48  %%%%%%%%%%

\[
\begin{tikzcd}[column sep=huge,row sep=large]
 V\otimes A
   \arrow[r]
   \arrow[d,"{F(D)}"]
 & (V\oplus K'^{\,n})\otimes A
   \arrow[r]
   \arrow[d,"{F(i,D)}"]
 & K'^{\,n}\otimes A
   \arrow[d,equal]\\
 V\otimes A
   \arrow[r]
 & (V\oplus K'^{\,n})\otimes A
   \arrow[r]
 & K'^{\,n}\otimes A.
\end{tikzcd}
\]

Let $x:K'^{\,n}\longrightarrow V$ be a map, put
$x_i=x(e_i)$, and let $(a_i)$ be a family of elements of $A$. We
have
\[
 (1_V+x)\circ i\circ D=D.
\]
Consequently, if
\[
 F(i,D)\Bigl(\sum_i e_i a_i\Bigr)
 =\sum_i e_i a_i+a,
 \qquad a\in V\otimes A,
\]
then
\[
 F(D)\Bigl(\sum_i x_i a_i\Bigr)
 =\sum_i x_i a_i+a,
\]
and $a$ is independent of the $x_i$. Hence
\[
\begin{aligned}
 F(D)(y)&=y+G(D),\\
 G(D+D')&=G(D)+G(D'),\\
 uG(D)&=G(uD)
   \qquad\text{for every }u:V\longrightarrow W.
\end{aligned}
\]
Every such $G$ is given, for some $\gamma\in\Omega\otimes A$, by
\[
 G(D)=(D\otimes1_A)(\gamma).
\]
This proves (2) and (3).

% END INCLUDED SGA7I_English_current_buildable_source_20260823/components/16_expose_VI_lemma_1_18.tex

% BEGIN INCLUDED SGA7I_English_current_buildable_source_20260823/components/17_expose_VI_1_19_through_1_20.tex
%%%%%%%%%%  scan idx 59  |  folio 48 (continued)  %%%%%%%%%%

\underline{Definition \SGARefTarget{sga7.sga7i.evi.n1-19}{1.19}.} (1) An object
$\xi\in\widehat{\underline A}$ is called a \underline{versal formal
object} if every surjective map $\alpha:a'\longrightarrow a$ in
$\underline A$ induces a surjective map
\[
 \operatorname{Hom}_{\widehat{\underline A}}(\xi,a')
 \longrightarrow
 \operatorname{Hom}_{\widehat{\underline A}}(\xi,a).
\]

(2) A versal formal object $\xi\in\widehat{\underline A}$, with
image $(R,s)$ in $\widehat C_{\Lambda,K}$, is called
\underline{minimal} if the following condition holds. For any two
maps
\[
 \varphi_1,\varphi_2:R\rightrightarrows D_{K'}(K),
\]
if $\varphi_1(\xi)$ is isomorphic to $\varphi_2(\xi)$, then there is
a derivation $D:K\longrightarrow K'$ such that
\[
 (1,D)(\varphi_1)=\varphi_2,
 \qquad\text{that is,}\qquad
 \varphi_2-\varphi_1=Ds.
\]

\underline{Theorem \SGARefTarget{sga7.sga7i.evi.n1-20}{1.20}. If $\underline A$ is semi-homogeneous and}
\[
 \dim_{K'}[\underline A(D_{K'}(K))]<\infty,
\]
\underline{then $\underline A$ admits a minimally versal formal
object, and any two such objects are isomorphic.}

%%%%%%%%%%  scan idx 60  |  folio 49  %%%%%%%%%%

With the notation of \SGARefLink{sga7.sga7i.evi.n1-18}{1.18}, let $H^*$ be a subspace of
$[\underline A(D_{K'}(K))]$ complementary to the image of $\gamma$,
and let $H$ be the dual of $H^*$. Set
\[
 R_0=D_H(K).
\]
Choose $\xi_0\in\underline A(R_0)$ whose class, under the
identifications
\[
 [\underline A(R_0)]
 \simeq H\otimes[\underline A(D_{K'}(K))]
 \simeq
 \operatorname{Hom}\bigl(H^*,[\underline A(D_{K'}(K))]\bigr),
\]
is the inclusion
$H^*\hookrightarrow[\underline A(D_{K'}(K))]$. Then:

(1) for every $V$, the maps
\[
 \operatorname{Hom}_{C_{\Lambda,K}}
   \bigl(D_H(K),D_V(K)\bigr)
 \longrightarrow[\underline A(D_V(K))]
\]
are surjective;

(2) the minimality condition of \SGARefLink{sga7.sga7i.evi.n1-19}{1.19}(2) holds for
$(R_0,\xi_0)$.

Choose $S\in\widehat C_{\Lambda,K}$ formally smooth over $\Lambda$,
together with an epimorphism $t:S\longrightarrow R_0$, such that
\[
 \Omega_{R_0/\Lambda}\otimes K'
 \xrightarrow{\ \sim\ }
 \Omega_{S/\Lambda}\otimes K'.
\]
If $M$ is the maximal ideal of $S$, define inductively, for
$n\geq1$,
\[
 R_n=S/I_n,
 \qquad
 \xi_n\in\operatorname{Ob}(\underline A(R_n)),
\]
by the following conditions:

(a) $I_{n+1}$ is the intersection of the ideals $I$ such that
\[
 I_nM\subset I\subset I_n
\]
and $\xi_n$ extends to an object $\xi\in\underline A(S/I)$;

(b) $\xi_{n+1}$ is an extension of $\xi_n$.

If
\[
 R=\varprojlim_nR_n,
 \qquad
 \xi=\varprojlim_n\xi_n,
\]
then $(R,\xi)$ is a minimally versal formal object, and it is unique
up to isomorphism.


% END INCLUDED SGA7I_English_current_buildable_source_20260823/components/17_expose_VI_1_19_through_1_20.tex

% BEGIN INCLUDED SGA7I_English_current_buildable_source_20260823/components/18_expose_VI_section_2_opening_through_2_2.tex
\section*{\SGARefTarget{sga7.sga7i.evi.n2}{2}. Prorepresentable cofibered groupoids}

Let $\underline A$ be a cofibered groupoid over
$\underline C_\Lambda$ admitting a minimally versal formal object
$\xi$. In general, $\xi$ alone does not determine $\underline A$ up
to $\underline C_\Lambda$-equivalence. One must also know when two
maps
\[
 R\mathrel{\substack{\xrightarrow{\ f\ }\\[-.6ex]
                      \xrightarrow[\ g\ ]{}}}S
\]
give rise to a $\underline C_\Lambda$-isomorphism
$\xi(f)\longrightarrow\xi(g)$, where $R=P(\xi)$. As an application
of \SGARefLink{sga7.sga7i.evi.n1-11}{1.11} and \SGARefLink{sga7.sga7i.evi.n1-12}{1.12}, we now consider the problem of recovering
$\underline A$, up to $\underline C_\Lambda$-equivalence. We obtain
a satisfactory answer for a broad class of cofibered groupoids.

%%%%%%%%%%  scan idx 61  |  folio 50  %%%%%%%%%%

Many cofibered groupoids that arise in practice satisfy a property
slightly stronger than semi-homogeneity; it will be defined
explicitly in \SGARefLink{sga7.sga7i.evi.n2-5}{2.5}. To investigate them, it is convenient to recast
the preceding discussion in terms of ``smooth functors.'' All the
notation and terminology of the preceding section remain in force
throughout this section.

\underline{Definition \SGARefTarget{sga7.sga7i.evi.n2-1}{2.1}.} Let
$\varphi:\underline A\longrightarrow\underline B$ be a functor of
cofibered groupoids over $\underline C_\Lambda$. We say that
$\varphi$ is \underline{smooth} if every surjective map
$\beta:b'\longrightarrow\varphi(a)$ in $\underline B$ can be
completed to a commutative diagram
\begin{equation}\tag{*}
\begin{tikzcd}[row sep=large,column sep=huge]
 \varphi(a') \arrow[r] \arrow[dr,"{\varphi(\alpha)}"']
 & b' \arrow[d,"\beta"]\\
 & \varphi(a)
\end{tikzcd}
\end{equation}
for some $\alpha:a'\longrightarrow a$ in $\underline A$. We say
that $\varphi$ is \underline{minimally smooth} if it is smooth and
\[
 \varphi(k[\varepsilon]):
 [\underline A(k[\varepsilon])]
 \longrightarrow
 [\underline B(k[\varepsilon])]
\]
is bijective.

\underline{Remark \SGARefTarget{sga7.sga7i.evi.n2-2.item-a}{\SGARefAliasTarget{sga7.sga7i.evi.n2-2}2.2}(a).} Let
$\varphi:\underline A\longrightarrow\underline B$ be a functor of
cofibered groupoids over $\underline C_\Lambda$. The following
conditions are equivalent:

(i) $\varphi$ is smooth;

(ii) every surjective map
$\beta:b'\longrightarrow\varphi(a)$ in $\underline B$ can be
completed to a diagram $(*)$ in which
$\varphi(a')\longrightarrow b'$ is a
$\underline C_\Lambda$-isomorphism;

(iii) every small prolongation
$\beta:b'\longrightarrow\varphi(a)$ in $\underline B$ can be
completed to a commutative diagram $(*)$.

\underline{Remark \SGARefTarget{sga7.sga7i.evi.n2-2.item-b}{2.2}(b).} For every cofibered groupoid
$\underline A$ over $\underline C_\Lambda$, the natural functor
\[
 \underline A\longrightarrow[\underline A]
\]
is minimally smooth.

\underline{Remark \SGARefTarget{sga7.sga7i.evi.n2-2.item-c}{2.2}(c).} For each
$R\in\widehat{\underline C}_\Lambda$, let
\[
 h_R:\underline C_\Lambda\longrightarrow(\mathrm{Ens})
\]
be the functor defined by
\[
 h_R(S)=\operatorname{Hom}_{\Lambda\text{-alg}}(R,S)
 \qquad(S\in\underline C_\Lambda).
\]
It determines a cofibered discrete groupoid $\underline h_R$ over
$\underline C_\Lambda$. A map $R\longrightarrow S$ in
$\widehat{\underline C}_\Lambda$ induces a functor
\[
 \underline h_S\longrightarrow\underline h_R.
\]
This functor is smooth if and only if $S$ is a formal power-series
ring over $R$.

%%%%%%%%%%  scan idx 62  |  folio 51  %%%%%%%%%%

\underline{Remark \SGARefTarget{sga7.sga7i.evi.n2-2.item-d}{2.2}(d).} Given
$R\in\widehat{\underline C}_\Lambda$, denote by
$(R,\underline C_\Lambda)$ the category of arrows in
$\widehat{\underline C}_\Lambda$ with source $R$ and target in
$\underline C_\Lambda$. Let $\underline A$ be a cofibered groupoid
over $\underline C_\Lambda$. Given
$\xi\in\operatorname{Ob}\widehat{\underline A}$ with $P(\xi)=R$,
let
\[
 \underline A'\longrightarrow(R,\underline C_\Lambda)
\]
be the pullback of
$\underline A\longrightarrow\underline C_\Lambda$ along the
canonical functor
$(R,\underline C_\Lambda)\longrightarrow\underline C_\Lambda$.
Choose a section $\xi^\#$ such that
$\xi^\#(1_R)=\xi$. Thus, for every map
$f:R\longrightarrow S$, we choose an $f$-map
$\xi\longrightarrow\xi(f)$.

We obtain a functor
\[
 \xi^\#:\underline h_R\longrightarrow\underline A
\]
over $\underline C_\Lambda$ by sending
$(R\xrightarrow{f}S)$ to $\xi(f)$. Up to natural equivalence, this
functor does not depend on the choice of $\xi^\#$.

The assertion that
$\xi^\#:\underline h_R\longrightarrow\underline A$ is
(minimally) smooth means that every surjective map
$a'\xrightarrow{\alpha}\xi(f)$ in $\underline A$ can be completed
to a commutative diagram
\[
\begin{tikzcd}[row sep=large,column sep=huge]
 \xi(f') \arrow[r] \arrow[dr,"{\xi(\,)}"']
 & a' \arrow[d,"\alpha"]\\
 & \xi(f)
\end{tikzcd}
\]
(and, in the minimal case, that
$\underline h_R(k[\varepsilon])\longrightarrow
\underline A(k[\varepsilon])$ is bijective). This means precisely
that $\xi$ is a (minimally) versal formal object of $\underline A$
over $\underline C_\Lambda$.

Thus the existence of a (minimally) versal formal object of
$\underline A$ is equivalent to the existence of a (minimally)
smooth functor
\[
 \underline h_R\longrightarrow\underline A
\]
over $\underline C_\Lambda$ for some
$R\in\operatorname{Ob}(\widehat{\underline C}_\Lambda)$.


% END INCLUDED SGA7I_English_current_buildable_source_20260823/components/18_expose_VI_section_2_opening_through_2_2.tex

% BEGIN INCLUDED SGA7I_English_current_buildable_source_20260823/components/19_expose_VI_2_3_through_2_7.tex
The following general properties of smoothness are immediate.

\underline{Proposition \SGARefTarget{sga7.sga7i.evi.n2-3}{2.3}.} Let
\[
 \underline A\xrightarrow{\ \varphi\ }\underline B
 \xrightarrow{\ \psi\ }\underline C
\]
be functors of cofibered groupoids over $\underline C_\Lambda$.

(a) If $\varphi$ and $\psi$ are both smooth, then so is
$\psi\circ\varphi$.

(b) If $\psi\circ\varphi$ is smooth and $\varphi$ is surjective on
isomorphism classes of objects, then $\psi$ is smooth.

%%%%%%%%%%  scan idx 63  |  folio 52  %%%%%%%%%%

\underline{Remark \SGARefTarget{sga7.sga7i.evi.n2-4}{2.4}.} Let $\underline A$ be a cofibered groupoid
over $\underline C_\Lambda$ admitting a versal formal object $\xi$.
By \SGARefLink{sga7.sga7i.evi.n2-2.item-d}{2.2}(d), this means that
\[
 \xi:\underline h_R\longrightarrow\underline A
\]
is smooth, where $R=P(\xi)$. Therefore the following assertions are
equivalent:

\begin{itemize}
\item $\underline A$ is formally smooth over $\underline C_\Lambda$
      (cf. \SGARefLink{sga7.sga7i.evi.n1-13}{1.13}), that is, $R$ is a formal power-series algebra over
      $\Lambda$;
\item the composite functor
      \[
       \underline h_R\xrightarrow{\ \xi\ }\underline A
       \xrightarrow{\ P\ }\underline C_\Lambda
       \;(=\underline h_\Lambda)
      \]
      is smooth;
\item $P:\underline A\longrightarrow\underline C_\Lambda$ is
      smooth.
\end{itemize}
The last equivalence follows from Proposition \SGARefLink{sga7.sga7i.evi.n2-3}{2.3}(b).

We now consider a stronger property than semi-homogeneity.

\underline{Definition \SGARefTarget{sga7.sga7i.evi.n2-5}{2.5}.} A cofibered groupoid $\underline A$ over
$\underline C_\Lambda$ is called \underline{homogeneous} if every
diagram
\[
\begin{tikzcd}[row sep=large,column sep=huge]
 & a' \arrow[d]\\
 b' \arrow[r] & a
\end{tikzcd}
\]
in which $a'\longrightarrow a$ is surjective admits a fibre product
$a'\times_a b'$ in $\underline A$ such that
\[
 P(a'\times_a b')=
 P(a')\times_{P(a)}P(b').
\]

\underline{Remark \SGARefTarget{sga7.sga7i.evi.n2-6.item-a}{\SGARefAliasTarget{sga7.sga7i.evi.n2-6}2.6}(a).} Let $\underline A$ be a cofibered
groupoid over $\underline C_\Lambda$, and consider a diagram
\[
(*)\qquad
\begin{tikzcd}[row sep=small,column sep=small]
 & a' \arrow[dl] \arrow[dr] & \\
 a_1 \arrow[dr] & & a_2 \arrow[dl]\\
 & a &
\end{tikzcd}
\qquad\text{in $\underline A$ above the diagram}\qquad
\begin{tikzcd}[row sep=small,column sep=small]
 & R_1\times_R R_2 \arrow[dl,"{\pi_1}"'] \arrow[dr,"{\pi_2}"] & \\
 R_1 \arrow[dr] & & R_2 \arrow[dl]\\
 & R &
\end{tikzcd}
\]
in $\underline C_\Lambda$, where the $\pi_i$ are the projections.
If a fibre product $a_1\times_a a_2$ exists in $\underline A$,
there is a unique map
\[
 \varphi:a'\longrightarrow a_1\times_a a_2
\]
compatible with the projections. If
\[
 P(a_1\times_a a_2)=R_1\times_R R_2,
\]
then $P(\varphi)$ is the identity map of $R_1\times_R R_2$;
hence $\varphi$ is a $\underline C_\Lambda$-isomorphism, that is,
$a'=a_1\times_a a_2$. Thus, if $\underline A$ is homogeneous,
every diagram $(*)$ defines a fibre product $a_1\times_a a_2$
whenever $R_1\longrightarrow R$ is surjective. It follows at once

%%%%%%%%%%  scan idx 64  |  folio 53  %%%%%%%%%%

that every homogeneous cofibered groupoid is, a fortiori,
semi-homogeneous.

\underline{Remark \SGARefTarget{sga7.sga7i.evi.n2-6.item-b}{2.6}(b).} Consider a diagram of categories and
functors
\[
\begin{tikzcd}[row sep=large,column sep=large]
 X \arrow[dr,"{\varphi}"'] & & Y \arrow[dl,"{\psi}"]\\
 & Z &
\end{tikzcd}
\]
The ordinary fibre product $X\times_ZY$ is the category whose
objects are pairs $(x,y)$ satisfying $\varphi(x)=\psi(y)$ and whose
maps $(x,y)\longrightarrow(x',y')$ are pairs
\[
 \alpha:x\longrightarrow x',\qquad
 \beta:y\longrightarrow y'
\]
such that $\varphi(\alpha)=\psi(\beta)$. The fibre product
$X\times_Z^2Y$ in the sense of $2$-categories is the category whose
objects are triples $(x,y;\gamma)$, where
\[
 \varphi(x)\xrightarrow{\ \gamma\ }\psi(y)
\]
is an isomorphism. A map
$(x,y;\gamma)\longrightarrow(x',y';\gamma')$ is a pair
$\alpha:x\longrightarrow x'$ and $\beta:y\longrightarrow y'$ for
which the diagram
\[
\begin{tikzcd}[row sep=huge,column sep=huge]
 \varphi(x) \arrow[r,"{\varphi(\alpha)}"] \arrow[d,"{\gamma}"']
 & \varphi(x') \arrow[d,"{\gamma'}"]\\
 \psi(y) \arrow[r,"{\psi(\beta)}"'] & \psi(y')
\end{tikzcd}
\]
commutes.

The notions of a (semi-)homogeneous cofibered groupoid and of a
smooth functor can be rephrased in terms of fibre products of
categories in this $2$-categorical sense. Let
$P:\underline A\longrightarrow\underline C_\Lambda$ be a
cofibered category, and consider $R_1\times_RR_2$ in
$\underline C_\Lambda$. Choose direct-image functors
\[
 (\pi_i)_*:\underline A(R_1\times_RR_2)
 \longrightarrow\underline A(R_i),\qquad i=1,2.
\]
They induce a functor
\[
 \underline A(R_1\times_RR_2)\longrightarrow
 \underline A(R_1)\times^2_{\underline A(R)}\underline A(R_2).
\]
Up to canonical natural equivalence, it is independent of the
chosen direct-image functors. The groupoid $\underline A$ is
homogeneous if and only if this functor is an equivalence whenever
$R_1\longrightarrow R$ is surjective; it is semi-homogeneous if and
only if
\[
 [\underline A(R_1\times_RR_2)]\longrightarrow
 [\underline A(R_1)]
 \times_{[\underline A(R)]}[\underline A(R_2)]
\]

%%%%%%%%%%  scan idx 65  |  folio 54  %%%%%%%%%%

is surjective whenever $R_1\longrightarrow R$ is surjective, and is
bijective whenever $R=k$ and $R_1=k[\varepsilon]$.

Now let $\varphi:\underline A\longrightarrow\underline B$ be a
functor of cofibered groupoids over $\underline C_\Lambda$. For
every map $R'\longrightarrow R$ in $\underline C_\Lambda$,
$\varphi$ induces, up to canonical natural equivalence, a functor
\[
 \underline A(R')\longrightarrow
 \underline A(R)\times^2_{\underline B(R)}\underline B(R').
\]
The functor $\varphi$ is smooth if and only if this functor is
surjective on isomorphism classes of objects whenever
$R'\longrightarrow R$ is surjective.

\underline{Remark \SGARefTarget{sga7.sga7i.evi.n2-6.item-c}{2.6}(c).} If $\underline A$ is homogeneous, then
for any finite-dimensional $k$-vector spaces $V$ and $W$ the functor
\[
 \underline A(k[V\times W])\longrightarrow
 \underline A(k[V])\times\underline A(k[W])
\]
is an equivalence of categories. Recall our convention
$\underline A(k)=\{e\}$. In particular,
\[
 \operatorname{Aut}_{\underline A(k[V\times W])}(1_{V\times W})
 \xrightarrow{\ \sim\ }
 \operatorname{Aut}_{\underline A(k[V])}(1_V)
 \times
 \operatorname{Aut}_{\underline A(k[W])}(1_W),
\]
where $1_V$ is the direct image of $e$ under the unique map
$k\longrightarrow k[V]$ in $\underline C_\Lambda$. Consequently,
if $\underline A$ is homogeneous, both
$\operatorname{Aut}(1_\varepsilon)$ and
$[\underline A(k[\varepsilon])]$ carry canonical $k$-vector-space
structures.

\underline{Remark \SGARefTarget{sga7.sga7i.evi.n2-6.item-d}{2.6}(d).} For every
$R\in\widehat{\underline C}_\Lambda$, the cofibered discrete
groupoid $\underline h_R$ over $\underline C_\Lambda$ is
homogeneous.

If $\underline A$ is semi-homogeneous, then so is $[\underline A]$.
Homogeneity of $\underline A$, however, does not in general imply
homogeneity of $[\underline A]$.

\underline{Proposition \SGARefTarget{sga7.sga7i.evi.n2-7}{2.7}.} Let $\underline A$ be a homogeneous
cofibered groupoid over $\underline C_\Lambda$. The following
conditions are equivalent.

(i) $[\underline A]$ is homogeneous; that is,
\[
 [\underline A(R'\times_RS)]\longrightarrow
 [\underline A(R')]
 \times_{[\underline A(R)]}[\underline A(S)]
\]
is bijective for every small extension $R'\longrightarrow R$.

(ii) For every small prolongation $a'\longrightarrow a$ in
$\underline A$, the map
\[
 \operatorname{Aut}_{\underline A(R')}(a')
 \longrightarrow
 \operatorname{Aut}_{\underline A(R)}(a)
\]

%%%%%%%%%%  scan idx 66  |  folio 55  %%%%%%%%%%

is surjective, where $R'=P(a')$ and $R=P(a)$.

\underline{Proof.} (i)$\Rightarrow$(ii). Let
$f:R'\longrightarrow R$ be a small extension in
$\underline C_\Lambda$, and let $a'\longrightarrow a$ be an
$f$-map in $\underline A$. Given
$\tau\in\operatorname{Aut}_{\underline A(R)}(a)$, consider
$a'\times_a a'$ and $a'\times_a a'_\tau$, where
$a'_\tau\longrightarrow a$ is the composite
\[
 a'\longrightarrow a\xrightarrow{\ \tau\ }a.
\]
Since
\[
 [\underline A(R'\times_RR')]
 \longrightarrow
 [\underline A(R')]
 \times_{[\underline A(R)]}[\underline A(R')]
\]
is bijective, there is an isomorphism
\[
 a'\times_a a'\xrightarrow{\ u\ }a'\times_a a'_\tau.
\]
Projection onto the second factor then gives an isomorphism
$\tau':a'\longrightarrow a'$ lifting $\tau$.

(ii)$\Rightarrow$(i). Let $a''_i$ $(i=1,2)$ be objects of
$\underline A(R'\times_RS)$ with the same image under
\[
 [\underline A(R'\times_RS)]\longrightarrow
 [\underline A(R')]
 \times_{[\underline A(R)]}[\underline A(S)].
\]
We obtain a diagram of solid arrows
\[
\begin{tikzcd}[column sep=small,row sep=small]
 && a''_2 \arrow[ddll] \arrow[ddrr] &&\\
 && a''_1 \arrow[u,dashed] \arrow[dl] \arrow[dr] &&\\
 a'_2 \arrow[ddrr] & a'_1 \arrow[l,"{\sim}"'] \arrow[dr]
 && b_1 \arrow[r,"{\tau}"] \arrow[dl] & b_2 \arrow[ddll]\\
 && a_1 \arrow[d,"{\bar\tau}"] &&\\
 && a_2 &&
\end{tikzcd}
\qquad\text{above}\qquad
\begin{tikzcd}[column sep=small,row sep=small]
 & R'\times_RS \arrow[dl] \arrow[dr] &\\
 R' \arrow[dr] && S \arrow[dl]\\
 & R &
\end{tikzcd}
\]
where $\bar\tau:a_1\longrightarrow a_2$ is the isomorphism induced
by $\tau$, so the lower-right square commutes. Under condition (ii),
the given isomorphism $a'_1\xrightarrow{\sim}a'_2$ can be replaced
so that the left square also commutes. Since
\[
 a''_i=a'_i\times_{a_i}b_i\qquad(i=1,2)
\]
by Remark \SGARefLink{sga7.sga7i.evi.n2-6.item-a}{2.6}(a), we obtain an isomorphism
$a''_1\xrightarrow{\sim}a''_2$. Hence
\[
 [\underline A(R'\times_RS)]\longrightarrow
 [\underline A(R')]
 \times_{[\underline A(R)]}[\underline A(S)]
\]
is injective, and therefore bijective because $\underline A$ is
homogeneous by hypothesis. \hfill$\square$


% END INCLUDED SGA7I_English_current_buildable_source_20260823/components/19_expose_VI_2_3_through_2_7.tex

% BEGIN INCLUDED SGA7I_English_current_buildable_source_20260823/components/20_expose_VI_2_8_through_2_9.tex
Let $\underline A$ be a homogeneous cofibered groupoid over
$\underline C_\Lambda$, and let $f:R'\longrightarrow R$ be a small
extension in $\underline C_\Lambda$. There is a canonical map

%%%%%%%%%%  scan idx 67  |  folio 56  %%%%%%%%%%

\[
 \widetilde f:R'\times_RR'\longrightarrow k[\operatorname{Ker}f],
 \qquad
 (r'_1,r'_2)\longmapsto r'_1(0)+(r'_1-r'_2),
\]
and
\[
 1\times\widetilde f:R'\times_RR'
 \longrightarrow R'\times_k k[\operatorname{Ker}f]
\]
is an isomorphism. Given two $f$-maps
$a_i\longrightarrow a$ $(i=1,2)$ in $\underline A$, the fibre
product $a_1\times_a a_2$ is an object over $R'\times_RR'$. We may
therefore form
\[
 \tau(a_1,a_2)=
 (\widetilde f)_*(a_1\times_a a_2)
 \in\underline A(k[\operatorname{Ker}f]).
\]
Since $1\times\widetilde f$ is an isomorphism, there is an
isomorphism
\[
 a_1\times_a a_2\simeq a_1\times\tau(a_1,a_2)
\]
compatible with projection onto the first factor. Statement \SGARefLink{sga7.sga7i.evi.n1-5}{1.5}(2)
now takes the following more definite form.

\underline{Lemma \SGARefTarget{sga7.sga7i.evi.n2-8}{2.8}.} Let $\underline A$ be a homogeneous
cofibered groupoid over $\underline C_\Lambda$, and let
\[
 a_i\xrightarrow{\ \alpha_i\ }a\qquad(i=1,2)
\]
be $f$-maps in $\underline A$.

(1) There exists $\rho:a_1\longrightarrow a_2$ in $\underline A$
with $\alpha_2\circ\rho=\alpha_1$ if and only if there is an arrow
$a_1\longrightarrow\tau(a_1,a_2)$ in $\underline A$.

(2) There exists such a $\rho$ with $P(\rho)=\operatorname{id}_{R'}$
if and only if
\[
 \tau(a_1,a_2)=0
 \quad\text{in}\quad
 [\underline A(k[\operatorname{Ker}f])].
\]

\underline{Proof.} We prove both assertions simultaneously. The
existence of $\rho:a_1\longrightarrow a_2$ satisfying
$\alpha_2\circ\rho=\alpha_1$ (and $P(\rho)=\operatorname{id}_{R'}$)
is equivalent successively to the existence of
\[
 h:a_1\longrightarrow a_1\times_a a_2,
 \qquad \operatorname{pr}_1\circ h=\operatorname{id}_{a_1},
\]
with $P(h)$ the diagonal map; to the existence of
\[
 g:a_1\longrightarrow a_1\times\tau(a_1,a_2),
 \qquad \operatorname{pr}_1\circ g=\operatorname{id}_{a_1},
\]
with
\[
 P(g):R'\longrightarrow R'\times_k k[\operatorname{Ker}f],
 \qquad r'\longmapsto(r',r'(0));
\]
and to the existence of an arrow
$\gamma:a_1\longrightarrow\tau(a_1,a_2)$ in $\underline A$, with
\[
 P(\gamma):R'\longrightarrow k[\operatorname{Ker}f],
 \qquad r'\longmapsto r'(0).
\]
Such an arrow factors through
$R'\longrightarrow k\longrightarrow k[\operatorname{Ker}f]$ if and
only if $\tau(a_1,a_2)=0$. \hfill$\square$

\underline{Proposition \SGARefTarget{sga7.sga7i.evi.n2-9}{2.9}.} (a) Let
\[
 \underline A\xrightarrow{\ \varphi\ }\underline B
 \xrightarrow{\ \psi\ }\underline C
\]
be functors of cofibered groupoids over $\underline C_\Lambda$.
Assume that

(i) $\underline B$ is semi-homogeneous and $\underline C$ is
homogeneous;

(ii) the map
\[
 \psi(k[\varepsilon]):[\underline B(k[\varepsilon])]
 \longrightarrow[\underline C(k[\varepsilon])]
\]
is injective.

Then $\psi\circ\varphi$ is smooth if and only if $\varphi$ and
$\psi$ are both

%%%%%%%%%%  scan idx 68  |  folio 57  %%%%%%%%%%

smooth.

(b) Let $\varphi:\underline A\longrightarrow\underline B$ be a
smooth functor of cofibered groupoids over $\underline C_\Lambda$.
If $\underline A$ is semi-homogeneous and $\underline B$ is
homogeneous, then
\[
 \widehat\varphi:\widehat{\underline A}
 \longrightarrow\widehat{\underline B}
\]
is smooth. Explicitly, for every surjective map
$b'\xrightarrow{\beta}\widehat\varphi(a)$ in
$\widehat{\underline B}$, there are a map
$\alpha:a'\longrightarrow a$ in $\widehat{\underline A}$ and a
$\widehat{\underline C}_\Lambda$-isomorphism
\[
 \rho:\widehat\varphi(a')\xrightarrow{\ \sim\ }b'
\]
such that $\beta\circ\rho=\varphi(\alpha)$.

\underline{Proof.} (a) Suppose $\psi\circ\varphi$ is smooth. If
$\varphi$ is smooth, then $\psi$ is smooth by Proposition \SGARefLink{sga7.sga7i.evi.n2-3}{2.3}(b),
so it remains to prove that $\varphi$ is smooth. Let
$b'\xrightarrow{\beta}\varphi(a)$ be a small prolongation in
$\underline B$. Since $\psi\circ\varphi$ is smooth, there are an
arrow $a'\xrightarrow{\alpha}a$ in $\underline A$ and a
$\underline C_\Lambda$-isomorphism
\[
 \rho:\psi(b')\longrightarrow\psi\varphi(a')
\]
for which
\[
\begin{tikzcd}[column sep=huge,row sep=huge]
 \psi(b') \arrow[rr,"{\rho}"] \arrow[dr,"{\psi(\beta)}"']
 && \psi\varphi(a') \arrow[dl,"{\psi\varphi(\alpha)}"]\\
 & \psi\varphi(a) &
\end{tikzcd}
\]
commutes. Since $\underline B$ is semi-homogeneous, we can choose a
commutative diagram
\[
\begin{tikzcd}[column sep=small,row sep=small]
 & b'' \arrow[dl] \arrow[dr] &
 &&& R'\times_RR' \arrow[dl] \arrow[dr] &\\
 b' \arrow[dr,"{\beta}"'] &&
 \varphi(a') \arrow[dl,"{\varphi(\alpha)}"]
 & \text{above} &
 R' \arrow[dr] && R' \arrow[dl]\\
 & \varphi(a) &&&& R &
\end{tikzcd}
\]
where $R=P(a)$ and $R'=P(b')=P(a')$. Homogeneity of
$\underline C$ gives
\[
 \psi(b'')=
 \psi(b')\times_{\psi\varphi(a)}\psi\varphi(a')
\]
by Remark \SGARefLink{sga7.sga7i.evi.n2-6.item-a}{2.6}(a). Lemma \SGARefLink{sga7.sga7i.evi.n2-8}{2.8}(2) then gives
\[
 0=\tau(\psi\varphi(a'),\psi(b'))
   =(\widetilde f)_*\psi(b'').
\]
By hypothesis (ii), it follows that $(\widetilde f)_*(b'')=0$.
Statement \SGARefLink{sga7.sga7i.evi.n1-5}{1.5}(2) now provides a map
$\gamma:b'\longrightarrow\varphi(a')$ satisfying
$\beta=\varphi(\alpha)\circ\gamma$. Thus $\varphi$ is smooth, and
then so is $\psi$. The converse follows from Proposition \SGARefLink{sga7.sga7i.evi.n2-3}{2.3}(a).

(b) Let a surjective map
$\beta:b'\longrightarrow\widehat\varphi(a)$ in
$\widehat{\underline B}$ be given. It is represented by a sequence
of commutative diagrams

%%%%%%%%%%  scan idx 69  |  folio 58  %%%%%%%%%%

\[
\begin{tikzcd}[column sep=large]
 {} \arrow[r,dashed] & b'_3 \arrow[r,"{\beta'_2}"] \arrow[d,"{\beta_3}"]
 & b'_2 \arrow[r,"{\beta'_1}"] \arrow[d,"{\beta_2}"]
 & b'_1 \arrow[d,"{\beta_1}"]\\
 {} \arrow[r,dashed] & \varphi(a_3) \arrow[r,"{\varphi(\alpha_2)}"']
 & \varphi(a_2) \arrow[r,"{\varphi(\alpha_1)}"']
 & \varphi(a_1)
\end{tikzcd}
\]
in $\underline B$, where every $\alpha_v$, $\beta_v$, and
$\beta'_v$ is surjective and
$\beta=\varprojlim_v\beta_v$. It is enough to show inductively that
every diagram of solid arrows below, with
$\alpha_n$, $\beta_n$, $\beta_{n+1}$, and $\beta'_n$ surjective,
can be completed by the dotted arrows:
\[
\begin{tikzcd}[column sep=large,row sep=large]
 \varphi(a'_{n+1})
   \arrow[rrrr,dashed,"{\varphi(\,)}"]
   \arrow[dr,dashed,leftrightarrow]
   \arrow[ddr,dashed,"{\varphi(\,)}"']
 &&&& \varphi(a'_n)
   \arrow[dl,leftrightarrow]
   \arrow[ddl,"{\varphi(\alpha'_n)}"]\\
 & b'_{n+1} \arrow[rr,"{\beta'_n}"] \arrow[d,"{\beta_{n+1}}"']
 && b'_n \arrow[d,"{\beta_n}"'] &\\
 & \varphi(a_{n+1}) \arrow[rr,"{\varphi(\alpha_n)}"']
 && \varphi(a_n) &
\end{tikzcd}
\]
By hypothesis $\underline A$ is semi-homogeneous and
$\alpha_n:a'_{n}\longrightarrow a_n$ is surjective. Hence there is
a commutative diagram
\[
\begin{tikzcd}[column sep=small,row sep=small]
 & a' \arrow[dl,"{\gamma}"'] \arrow[dr,"{\gamma'}"] &
 &&& R_{n+1}\times_{R_n}R'_n \arrow[dl] \arrow[dr] &\\
 a_{n+1} \arrow[dr,"{\alpha_n}"'] &&
 a'_n \arrow[dl,"{\alpha'_n}"]
 & \text{in $\underline A$ over} &
 R_{n+1} \arrow[dr] && R'_n \arrow[dl]
 & \text{in $\underline C_\Lambda$}\\
 & a_n &&&& R_n &&
\end{tikzcd}
\]
Homogeneity of $\underline B$ implies that
\[
\begin{tikzcd}[column sep=huge,row sep=huge]
 & \varphi(a') \arrow[dl,"{\varphi(\gamma)}"']
                  \arrow[dr,"{\varphi(\gamma')}"] &\\
 \varphi(a_{n+1}) \arrow[dr,"{\varphi(\alpha_n)}"']
 && \varphi(a'_n) \arrow[dl,"{\varphi(\alpha'_n)}"]\\
 & \varphi(a_n) &
\end{tikzcd}
\]

%%%%%%%%%%  scan idx 70  |  folio 59  %%%%%%%%%%

is a fibre-product diagram in $\underline B$. Thus there is a unique
map $\delta:b'_{n+1}\longrightarrow\varphi(a')$ with the required
commutativity. Since $\beta_{n+1}$ and $\beta'_n$ are surjective, so
is $\delta$. Smoothness of $\varphi$ therefore gives a map
$a'_{n+1}\xrightarrow{\alpha'}a'$ in $\underline A$ and a
$\underline C_\Lambda$-isomorphism
\[
 \rho:\varphi(a'_{n+1})\xrightarrow{\ \sim\ }b'_{n+1}
\]
such that $\delta\circ\rho=\varphi(\alpha')$. We obtain the
commutative diagram
\[
\begin{tikzcd}[row sep=huge,column sep=huge]
 \varphi(a'_{n+1})
   \arrow[rrrr,"{\varphi(\gamma'\circ\alpha')}"]
   \arrow[dr,"{\rho}"]
   \arrow[ddr,"{\varphi(\gamma\circ\alpha')}"']
 &&&& \varphi(a'_n)
   \arrow[dl,"{\sim}"']
   \arrow[ddl,"{\varphi(\alpha'_n)}"]\\
 & b'_{n+1} \arrow[rr] \arrow[d]
 && b'_n \arrow[d] &\\
 & \varphi(a_{n+1}) \arrow[rr]
 && \varphi(a_n) &
\end{tikzcd}
\]
which completes the induction and proves (b). \hfill$\square$


% END INCLUDED SGA7I_English_current_buildable_source_20260823/components/20_expose_VI_2_8_through_2_9.tex

% BEGIN INCLUDED SGA7I_English_current_buildable_source_20260823/components/21_expose_VI_corollary_2_10.tex
\underline{Corollary \SGARefTarget{sga7.sga7i.evi.n2-10}{2.10}.} (a) Let
$\varphi:\underline A\longrightarrow\underline B$ be a smooth functor
of cofibered groupoids over $\underline C_\Lambda$. Let $\xi$ be a
versal formal object of $\underline A$, and let $\eta$ be a minimally
versal formal object of $\underline B$. If $\underline B$ is
homogeneous, then for every map
$\beta:\eta\longrightarrow\varphi(\xi)$ in
$\widehat{\underline B}$, the ring $P(\xi)$ is a formal power-series
ring over $P(\eta)$ via
\[
 P(\beta):P(\eta)\longrightarrow P(\varphi(\xi))=P(\xi).
\]

(b) Let $\underline A$ be a homogeneous cofibered groupoid over
$\underline C_\Lambda$, and let $\xi$ and $\eta$ be versal formal
objects of $\underline A$. If $\eta$ is minimal, then for every map
$\alpha:\eta\longrightarrow\xi$ in $\widehat{\underline A}$, the
ring $P(\xi)$ is a formal power-series ring over $P(\eta)$ via
\[
 P(\alpha):P(\eta)\longrightarrow P(\xi).
\]

\underline{Proof.} (a) A map
$\beta:\eta\longrightarrow\varphi(\xi)$ induces a commutative
diagram
\[
\begin{tikzcd}[row sep=huge,column sep=huge]
 \underline h_{P(\xi)} \arrow[r,"{\widetilde\beta}"]
                           \arrow[d,"{\widetilde\xi}"']
 & \underline h_{P(\eta)} \arrow[d,"{\widetilde\eta}"]\\
 \underline A \arrow[r,"{\varphi}"'] & \underline B.
\end{tikzcd}
\]
Since $\widetilde\xi$ and $\varphi$ are smooth, the composite
$\widetilde\eta\circ\widetilde\beta
=\varphi\circ\widetilde\xi$ is smooth. On the other hand,

%%%%%%%%%%  scan idx 71  |  folio 60  %%%%%%%%%%

\[
 \widetilde\eta(k[\varepsilon]):
 \underline h_{P(\eta)}(k[\varepsilon])
 \longrightarrow[\underline B(k[\varepsilon])]
\]
is an isomorphism. Proposition \SGARefLink{sga7.sga7i.evi.n2-9}{2.9}(a) therefore shows that
$\widetilde\beta$ is smooth, which means precisely that $P(\xi)$ is
a formal power-series ring over $P(\eta)$.

(b) Apply (a) with $\underline A=\underline B$. \hfill$\square$

% END INCLUDED SGA7I_English_current_buildable_source_20260823/components/21_expose_VI_corollary_2_10.tex

% BEGIN INCLUDED SGA7I_English_current_buildable_source_20260823/components/22_expose_VI_representable_cogroupoids_intro.tex
We conclude by considering the ``representability'' of a cofibered
groupoid. Let $\underline C$ be a fixed category. Every object $R$ of
$\underline C$ defines a functor
\[
 h_R:\underline C\longrightarrow(\mathrm{Sets}),
 \qquad h_R(S)=\operatorname{Hom}_{\underline C}(R,S),
\]
and hence a fully faithful functor
\[
 h:\underline C^{\mathrm{op}}
 \longrightarrow
 \underline{\operatorname{Hom}}(\underline C,(\mathrm{Sets})).
\]
A \emph{cocategory in $\underline C$} is a diagram
\[
\begin{tikzcd}[column sep=huge]
 h_{R_1}\times_{h_{R_0}}h_{R_1} \arrow[r]
 & h_{R_1}
   \arrow[r,shift left=3,"{s_1}"]
   \arrow[r,shift right=3,"{s_2}"']
 & h_{R_0} \arrow[l,"{\epsilon}" description]
\end{tikzcd}
\qquad
\text{with }s_i\circ\epsilon=\operatorname{id}\quad(i=1,2),
\]
where the first arrow is denoted by $c$ and
\[
 h_{R_1}\times_{h_{R_0}}h_{R_1}
 = (h_{R_1})_{s_1}\times_{h_{R_0}}{}_{s_2}(h_{R_1}),
\]
such that

(i) $c$ is associative and compatible with $s_i$ for $i=1,2$;

(ii) both composites
\begin{align*}
 h_{R_1}&\xrightarrow{\ \sim\ }
 h_{R_0}\times_{h_{R_0}}h_{R_1}
 \xrightarrow{\ \epsilon\times1\ }
 h_{R_1}\times_{h_{R_0}}h_{R_1}
 \xrightarrow{\ c\ }h_{R_1},\\[1ex]
 h_{R_1}&\xrightarrow{\ \sim\ }
 h_{R_1}\times_{h_{R_0}}h_{R_0}
 \xrightarrow{\ 1\times\epsilon\ }
 h_{R_1}\times_{h_{R_0}}h_{R_1}
 \xrightarrow{\ c\ }h_{R_1}
\end{align*}
are identities.

If $(R_0,R_1,s_1,s_2,\epsilon,c)$ is a cocategory in
$\underline C$, then for every $S\in\underline C$ we define a
category $h_{(R_0,R_1,\ldots)}(S)$ by
\[
 \operatorname{ob}h_{(R_0,R_1,\ldots)}(S)
   =\operatorname{Hom}_{\underline C}(R_0,S),
 \qquad
 \operatorname{Fl}h_{(R_0,R_1,\ldots)}(S)
   =\operatorname{Hom}_{\underline C}(R_1,S),
\]
with composition supplied by
\[
 c:h_{R_1}\times_{h_{R_0}}h_{R_1}\longrightarrow h_{R_1}.
\]
The categories so obtained depend functorially on $S$ and define a
cofibered category
$\underline h_{(R_0,R_1,\ldots)}$ over $\underline C$. For example,
every object $R$ of $\underline C$ defines a discrete cocategory
\[
 \underline h_R
 =\underline h_{(R,R,s_1,s_2,\epsilon,c)}
\]
by taking $s_1,s_2,\epsilon$, and $c$ all to be identity maps. A
homomorphism
\[
 (R_0,R_1,\ldots)\longrightarrow(R'_0,R'_1,\ldots)
\]

%%%%%%%%%%  scan idx 72  |  folio 61  %%%%%%%%%%

of cocategories is a pair of maps of represented functors
$h_{R_0}\longrightarrow h_{R'_0}$ and
$h_{R_1}\longrightarrow h_{R'_1}$ whose induced map
\[
 h_{(R_0,R_1,\ldots)}(S)
 \longrightarrow h_{(R'_0,R'_1,\ldots)}(S)
\]
is a functor for every $S\in\underline C$. Every cocategory
$(R_0,R_1,\ldots)$ has, in particular, a homomorphism
\[
 \epsilon:\underline h_{R_0}
 \longrightarrow\underline h_{(R_0,R_1,\ldots)}.
\]
A \emph{cogroupoid in $\underline C$} is a cocategory
$(R_0,R_1,\ldots)$ for which $h_{(R_0,R_1,\ldots)}(S)$ is a groupoid
for every $S\in\underline C$; equivalently,
$\underline h_{(R_0,R_1,\ldots)}$ is a category cofibered in
groupoids over $\underline C$.

In practice, cocategories and cogroupoids arise as follows. Let
$P:\underline F\longrightarrow\underline C$ be a cofibered category,
and fix an object $\xi$ of $\underline F$. Put $R_0=P(\xi)$, and let
$(R_0,\underline C)$ denote the category of arrows in $\underline C$
with source $R_0$. Let
$P:\underline F'\longrightarrow(R_0,\underline C)$ be the pullback of
$\underline F$ along the canonical functor
$(R_0,\underline C)\longrightarrow\underline C$. Choose a
cocartesian section $\xi^{\#}$ of $\underline F'$ over
$(R_0,\underline C)$ such that
$\xi^{\#}(\operatorname{id}_{R_0})=\xi$. Thus, for every arrow
$R_0\xrightarrow fS$ in $\underline C$, we choose a cocartesian
$f$-map $\xi\longrightarrow\xi^{\#}(f)$.

Define
\[
 \operatorname{Hom}_{\xi^{\#}}:
 (R_0\amalg R_0,\underline C)\longrightarrow(\mathrm{Sets})
\]
by
\[
 X=(f_1,f_2)\longmapsto
 \operatorname{Hom}_{\underline F(S)}
   (\xi^{\#}(f_1),\xi^{\#}(f_2)),
\]
where $S$ is the common target of $f_1$ and $f_2$. Up to canonical
natural equivalence, this functor is independent of the choice of
$\xi^{\#}$; henceforth we write simply $\xi$ for $\xi^{\#}$.

Suppose the functor is represented by an object
\[
 X_0=
 \left(R_0\mathrel{\substack{\xrightarrow{\ s_1\ }\\[-.6ex]
                              \xrightarrow[\ s_2\ ]{}}}R_1\right)
\]
of $(R_0\amalg R_0,\underline C)$. Thus there is an $R_0$-map
$\varphi:\xi(s_1)\longrightarrow\xi(s_2)$ in $\underline F$ such
that
\[
 \widetilde\varphi:
 \operatorname{Hom}_{(R_0\amalg R_0,\underline C)}(X_0,X)
 \longrightarrow\operatorname{Hom}_\xi(X)
\]
is bijective for every $X$ in $(R_0\amalg R_0,\underline C)$. Then
$X_0$ naturally defines a cocategory $(R_0,R_1,\ldots)$ in
$\underline C$. Indeed, a triple of arrows

%%%%%%%%%%  scan idx 73  |  folio 62  %%%%%%%%%%

\[
\begin{tikzcd}[column sep=huge]
 R_0 \arrow[r,shift left=3,"u"]
     \arrow[r,"v"]
     \arrow[r,shift right=3,"w"']
 & S
\end{tikzcd}
\]
in $\underline C$ defines a composition
\[
 \operatorname{Hom}_{\underline F(S)}(\xi(v),\xi(w))
 \times
 \operatorname{Hom}_{\underline F(S)}(\xi(u),\xi(v))
 \longrightarrow
 \operatorname{Hom}_{\underline F(S)}(\xi(u),\xi(w)),
\]
and hence a composition rule
\[
 h_{R_1}(S)\times_{h_{R_0}(S)}h_{R_1}(S)
 \longrightarrow h_{R_1}(S),
\]
functorial in $S$. Consequently, $\operatorname{Hom}_\xi$ defines the
cocategory $(R_0,R_1,\ldots)$ in $\underline C$ and the cofibered
category $\underline h_{(R_0,R_1,\ldots)}$ over $\underline C$.
Explicitly,
\[
 \operatorname{ob}\underline h_{(R_0,R_1,\ldots)}
 =\coprod_{S\in\operatorname{ob}(\underline C)}
   \operatorname{Hom}_{\underline C}(R_0,S).
\]
An arrow
\[
 (R_0\xrightarrow fS)\longrightarrow(R_0\xrightarrow{f'}S')
\]
is a pair of arrows $S\xrightarrow hS'$ and $R_1\xrightarrow uS'$ in
$\underline C$ such that
\[
 u\circ s_1=h\circ f,
 \qquad u\circ s_2=f'.
\]
The structural functor
$P:\underline h_{(R_0,R_1,\ldots)}\longrightarrow\underline C$ sends
$(R_0\xrightarrow fS)$ to $S$.

The pair $(\xi,\varphi)$ now defines a functor
\[
 (\xi,\varphi):\underline h_{(R_0,R_1,\ldots)}
 \longrightarrow\underline F
 \quad\text{over }\underline C
\]
by
\[
 (R_0\xrightarrow fS)\longmapsto\xi(f)
\]
on objects and by
\[
 \left[(R_0\xrightarrow fS)
 \xrightarrow{(h,u)}(R_0\xrightarrow{f'}S')\right]
 \longmapsto
 \left[\xi(f)\xrightarrow{h_*}\xi(hf)
 \xrightarrow{u(\varphi)}\xi(f')\right]
\]
on arrows. For every $S\in\underline C$, the functor
\[
 (\xi,\varphi)(S):
 \underline h_{(R_0,R_1,\ldots)}(S)\longrightarrow\underline F(S)
\]
is fully faithful by the representability of
$\operatorname{Hom}_\xi$. It follows that
\[
 (\xi,\varphi):\underline h_{(R_0,R_1,\ldots)}
 \longrightarrow\underline F
\]
is a $\underline C$-equivalence if and only if $\xi$ is quasi-initial
in $\underline F$ over $\underline C$, that is, if every object $a$
of $\underline F$ admits a cocartesian arrow $\xi\longrightarrow a$.

One may instead consider the functor
\[
 \operatorname{Isom}_\xi:
 (R_0\amalg R_0,\underline C)\longrightarrow(\mathrm{Sets}),
\]

%%%%%%%%%%  scan idx 74  |  folio 63  %%%%%%%%%%

given by
\[
 X=(f_1,f_2)\longmapsto
 \operatorname{Isom}_{\underline F(S)}(\xi(f_1),\xi(f_2)),
\]
where $S$ is the common target of $f_1$ and $f_2$. If this functor is
represented by an object $(R_0\rightrightarrows R_1)$ of
$(R_0\amalg R_0,\underline C)$, then it defines a cogroupoid
$(R_0,R_1,\ldots)$ in $\underline C$, together with a functor
\[
 (\xi,\varphi):\underline h_{(R_0,R_1,\ldots)}
 \longrightarrow\underline F.
\]
When $\underline F$ is a cofibered groupoid, this functor is a
$\underline C$-equivalence if and only if $\xi$ is quasi-initial in
$\underline F$ over $\underline C$.

% END INCLUDED SGA7I_English_current_buildable_source_20260823/components/22_expose_VI_representable_cogroupoids_intro.tex

% BEGIN INCLUDED SGA7I_English_current_buildable_source_20260823/components/23_expose_VI_2_11_and_2_12_ab.tex
\underline{Definition and Proposition \SGARefTarget{sga7.sga7i.evi.n2-11}{2.11}.} A cogroupoid
$(R_0,R_1,\ldots)$ in $\widehat{\underline C}_\Lambda$ is called
\underline{smooth} if either, and hence both, of the following
equivalent conditions holds:

(i) $R_1$ is a formal power-series ring over $R_0$ via $s_1$ (and
via $s_2$);

(ii) the functor
\[
 \epsilon:\underline h_{R_0}
 \longrightarrow\underline h_{(R_0,R_1,\ldots)}
\]
is smooth over $\underline C_\Lambda$.

A \underline{minimally smooth cogroupoid}, also called a
\underline{normalized smooth cogroupoid}, in
$\widehat{\underline C}_\Lambda$ is a smooth cogroupoid
$(R_0,R_1,\ldots)$ for which
$\underline h_{(R_0,R_1,\ldots)}(k[\varepsilon])$ is totally
disconnected; equivalently,
\[
 s_1(k[\varepsilon])=s_2(k[\varepsilon]).
\]

\underline{Proof.} Smoothness of
$\epsilon:\underline h_{R_0}\to
\underline h_{(R_0,R_1,\ldots)}$ means that, for every surjective
map
\[
 (S',f')\xrightarrow{(h,u)}(S,f)
\]
in $\underline h_{(R_0,R_1,\ldots)}$, there are a lift $f''$ of $f$
to $S'$ and an isomorphism
\[
 (S',f'')\xrightarrow{(1,u')}(S',f').
\]
Equivalently, given $u\in h_{R_1}(S)$ and a lift $f'$ of $s_1(u)$
to $S'$, there is a lift $u'$ of $u$ to $S'$ such that
$s_1(u')=f'$, for every surjection $S'\to S$ in
$\underline C_\Lambda$. This is precisely the smoothness of
\[
 s_1:\underline h_{R_1}\longrightarrow\underline h_{R_0}
\]
over $\underline C_\Lambda$, and that in turn says that $R_1$ is a
formal power-series ring over $R_0$ via $s_1$. \hfill$\square$

\underline{Proposition \SGARefTarget{sga7.sga7i.evi.n2-12}{2.12}.} Let $(R_0,R_1,\ldots)$ be a cogroupoid
in $\widehat{\underline C}_\Lambda$.

(a) If $(R_0,R_1,\ldots)$ is smooth, then
$\underline h_{(R_0,R_1,\ldots)}$ is homogeneous.

(b) If $s_1(k[\varepsilon])=s_2(k[\varepsilon])$, then the following
conditions are equivalent:

(i) $(R_0,R_1,\ldots)$ is smooth;

(ii) $\underline h_{(R_0,R_1,\ldots)}$ is homogeneous;

(iii) $\underline h_{(R_0,R_1,\ldots)}$ is semi-homogeneous.

%%%%%%%%%%  scan idx 75  |  folio 64  %%%%%%%%%%

(c) If $\underline h_{(R_0,R_1,\ldots)}$ is semi-homogeneous, then
it is equivalent to $\underline h_{(S_0,S_1,\ldots)}$ for some
minimally smooth cogroupoid $S_*$. 

\underline{Proof.} (a) Consider a diagram
\[
\begin{tikzcd}[column sep=huge,row sep=large]
 (S_1,f_1) \arrow[dr,"{(h_1,u_1)}"']
 && (S_2,f_2) \arrow[dl,"{(h_2,u_2)}"]\\
 & (S,f) &
\end{tikzcd}
\]
in $\underline h_{(R_0,R_1,\ldots)}$, where
$h_2:S_2\longrightarrow S$ is surjective in
$\underline C_\Lambda$. Since
$s_1:\underline h_{R_1}\longrightarrow\underline h_{R_0}$ is
smooth, $u_1^{-1}u_2$ lifts to an element $v\in h_{R_1}(S_2)$ with
$s_1(v)=f_2$. We have
\[
 h_2\circ s_2(v)=h_1\circ f_1,
\]
so $f_1\times s_2(v)$ is an element of
$h_{R_0}(S_1\times_SS_2)$ and gives a commutative diagram
\[
\begin{tikzcd}[column sep=huge,row sep=large]
 & (S_1\times_SS_2,,f_1\times s_2(v))
   \arrow[dl,"{\pi_1}"']
   \arrow[dr,"{(\pi_2,v^{-1})}"] &\\
 (S_1,f_1) \arrow[dr,"{(h_1,u_1)}"']
 && (S_2,f_2) \arrow[dl,"{(h_2,u_2)}"]\\
 & (S,f) &
\end{tikzcd}
\]
We claim that its top object is the fibre product
\[
 (S_1,f_1)\times_{(S,f)}(S_2,f_2)
\]
in $\underline h_{(R_0,R_1,\ldots)}$. Indeed, for every object
$(T,g)$, a map to the top object is equivalently a tuple
$(h'_1\times h'_2,w_1,w_2)$ satisfying
\begin{align*}
 h_1w_1&=h_2w_2,&
 s_1(w_1)&=h'_1g,&s_2(w_1)&=f_1,\\
 s_1(w_2)&=h'_2g,&
 s_2(w_2)&=s_2(v).
\end{align*}
Putting $w'_2=v^{-1}w_2$ gives a bijection with the tuples
$(h'_1\times h'_2,w_1,w'_2)$ satisfying
\begin{align*}
 u_2h_2w'_2&=u_1h_1w_1,&
 s_1(w_1)&=h'_1g,&s_2(w_1)&=f_1,\\
 s_1(w'_2)&=h'_2g,&
 s_2(w'_2)&=f_2.
\end{align*}
These are exactly the pairs of arrows from $(T,g)$ to $(S_1,f_1)$
and $(S_2,f_2)$ whose composites to $(S,f)$ coincide. Thus
\begin{multline*}
 \operatorname{Hom}((T,g),(S_1\times_SS_2,
                 f_1\times s_2(v)))\\
 \simeq
 \operatorname{Hom}((T,g),(S_1,f_1))
 \times_{\operatorname{Hom}((T,g),(S,f))}
 \operatorname{Hom}((T,g),(S_2,f_2)),
\end{multline*}
which proves (a).

%%%%%%%%%%  scan idx 76  |  folio 65  %%%%%%%%%%

(b) It remains only to prove (iii)$\Rightarrow$(i). If
$\underline h_{(R_0,R_1,\ldots)}$ is semi-homogeneous, then
\[
 \dim[\underline h_{(R_0,R_1,\ldots)}(k[\varepsilon])]
 =\dim h_{R_0}(k[\varepsilon])<\infty,
\]
so it admits a minimally versal formal object. On the other hand,
$(R_0,1_{R_0})$ is a quasi-initial minimal object of
$\widehat{\underline h}_{(R_0,R_1,\ldots)}$, because
\[
 h_{R_0}(k[\varepsilon])
 \longrightarrow
 [\underline h_{(R_0,R_1,\ldots)}(k[\varepsilon])]
\]
is bijective. Proposition \SGARefLink{sga7.sga7i.evi.n1-10}{1.10}(ii) therefore makes
$(R_0,1_{R_0})$ a minimally versal formal object of
$\underline h_{(R_0,R_1,\ldots)}$. Equivalently,
\[
 \epsilon:\underline h_{R_0}
 \longrightarrow\underline h_{(R_0,R_1,\ldots)}
\]
is minimally smooth. This proves (i).

% END INCLUDED SGA7I_English_current_buildable_source_20260823/components/23_expose_VI_2_11_and_2_12_ab.tex

% BEGIN INCLUDED SGA7I_English_current_buildable_source_20260823/components/24_expose_VI_2_12_c_through_2_14.tex
(c) Let $(S_0,\xi_0)$ be a minimally versal formal object of
$\underline h_{(R_0,R_1,\ldots)}$ (Theorem \SGARefLink{sga7.sga7i.evi.n1-11}{1.11}). For any object
\[
 \xi:R_0\longrightarrow T
 \qquad\text{of }\underline h_{(R_0,R_1,\ldots)}(T),
\]
consider the functor
\[
 \operatorname{Isom}(\xi,\xi):
 \underline C_\Lambda/(T\widehat\otimes_\Lambda T)
 \longrightarrow(\mathrm{Sets}).
\]
It assigns to $R$ the set of isomorphisms between the two images of
$\xi$ on $R$ obtained from
\[
 R_0\xrightarrow{\ \xi\ }T
 \mathrel{\substack{\xrightarrow{}\\[-.6ex]\xrightarrow[]{}}}
 T\widehat\otimes_\Lambda T\longrightarrow R.
\]
This functor is represented by
\[
 R_1\otimes_{R_0\widehat\otimes_\Lambda R_0}
       (T\widehat\otimes_\Lambda T).
\]
Passing to the inverse limit shows that
$\operatorname{Isom}(\xi_0,\xi_0)$ is prorepresented by an object
$S_1$, and that
\[
 \underline h_{(R_0,R_1,\ldots)}
 \simeq\underline h_{(S_0,S_1,\ldots)}.
\]
By (b), the cogroupoid $S_*$ is minimally smooth. \hfill$\square$

\underline{Definition \SGARefTarget{sga7.sga7i.evi.n2-13}{2.13}.} A cofibered groupoid $\underline A$
over $\underline C_\Lambda$ is called \underline{(minimally)
prorepresentable} if it is $\underline C_\Lambda$-equivalent to
$\underline h_{(R_0,R_1,\ldots)}$ for some (normalized) cogroupoid
$(R_0,R_1,\ldots)$ in $\widehat{\underline C}_\Lambda$.

\underline{Proposition \SGARefTarget{sga7.sga7i.evi.n2-14}{2.14}.} If $(R_0,R_1,\ldots)$ is a normalized
cogroupoid in $\widehat{\underline C}_\Lambda$, then every
homomorphism
\[
 (\alpha_0,\alpha_1):(R_0,R_1,\ldots)
 \longrightarrow(R_0,R_1,\ldots)
\]
which induces a $\underline C_\Lambda$-equivalence
\[
 \underline h_{(\alpha_0,\alpha_1)}:
 \underline h_{(R_0,R_1,\ldots)}
 \longrightarrow\underline h_{(R_0,R_1,\ldots)}
\]
is an isomorphism. In

%%%%%%%%%%  scan idx 77  |  folio 66  %%%%%%%%%%

particular, a normalized prorepresentation of a cofibered groupoid
over $\underline C_\Lambda$, if it exists, is unique up to
isomorphism.

\underline{Proof.} Consider the commutative diagram
\[
\begin{tikzcd}[column sep=huge,row sep=huge]
 \underline h_{(R_0,R_1,\ldots)}
   \arrow[r,"{\underline h_{(\alpha_0,\alpha_1)}}"]
 & \underline h_{(R_0,R_1,\ldots)}\\
 \underline h_{R_0} \arrow[u,"\epsilon"']
   \arrow[r,"{\underline h_{\alpha_0}}"']
 & \underline h_{R_0} \arrow[u,"\epsilon"]
\end{tikzcd}
\]
in which $\underline h_{(\alpha_0,\alpha_1)}$ is a
$\underline C_\Lambda$-equivalence. Since
$\underline h_{(R_0,R_1,\ldots)}(k[\varepsilon])$ is totally
disconnected, both
\[
 [\underline h_{(\alpha_0,\alpha_1)}(k[\varepsilon])]
 \quad\text{and}\quad
 [\epsilon(k[\varepsilon])]
\]
are bijections, and hence so is
$\underline h_{\alpha_0}(k[\varepsilon])$. Thus
$\alpha_0:R_0\longrightarrow R_0$ is a surjective endomorphism and
therefore an automorphism. It follows that $\alpha_1$ is also an
automorphism of $R_1$. \hfill$\square$

% END INCLUDED SGA7I_English_current_buildable_source_20260823/components/24_expose_VI_2_12_c_through_2_14.tex

% BEGIN INCLUDED SGA7I_English_current_buildable_source_20260823/components/25_expose_VI_2_15_through_2_17.tex
\underline{Lemma \SGARefTarget{sga7.sga7i.evi.n2-15}{2.15}.} Let $A$ and $B$ be discrete groupoids over
$\underline C_\Lambda$, and let $\varphi:A\longrightarrow B$ be a
minimally smooth functor. If $A$ and $B$ are homogeneous, then
$\varphi$ is a $\underline C_\Lambda$-equivalence.

\underline{Proof.} We must show that
$\varphi(S):A(S)\longrightarrow B(S)$ is bijective for every
$S\in\operatorname{ob}(\underline C_\Lambda)$. It is surjective for
every $S$ and bijective for $S=k[\varepsilon]$. We argue by induction
on the length of $S$. Let $S'\xrightarrow fS$ be a small extension
in $\underline C_\Lambda$. Since $A$ and $B$ are homogeneous
discrete groupoids, the canonical isomorphism
\[
 1\times\widetilde f:S'\times_SS'
 \xrightarrow{\ \sim\ }S'\times_k k[\operatorname{Ker}f]
\]
induces a commutative diagram

%%%%%%%%%%  scan idx 78  |  folio 67  %%%%%%%%%%

\[
\begin{tikzcd}[column sep=huge,row sep=huge]
 A(S')\times_{A(S)}A(S') \arrow[r,"\sim"] \arrow[d]
 & A(S')\times A(k[\operatorname{Ker}f]) \arrow[d]\\
 B(S')\times_{B(S)}B(S') \arrow[r,"\sim"']
 & B(S')\times B(k[\operatorname{Ker}f]).
\end{tikzcd}
\]
Thus $A(S')$ and $B(S')$ are principal homogeneous sets over $A(S)$
and $B(S)$, with structural groups identified by
\[
 A(k[\operatorname{Ker}f])\xrightarrow{\ \sim\ }
 B(k[\operatorname{Ker}f]).
\]
If $\varphi(S)$ is bijective, it follows that $\varphi(S')$ is
injective, and hence bijective. This completes the induction.
\hfill$\square$

\underline{Corollary \SGARefTarget{sga7.sga7i.evi.n2-16}{2.16}.} Let $\underline A$ be a homogeneous
groupoid over $\underline C_\Lambda$ such that
\[
 \operatorname{Aut}_{\underline A(k[\varepsilon])}(1_\varepsilon)
\]
is finite-dimensional. For every
$\xi\in\operatorname{ob}(\widehat{\underline A})$, the functor
\[
 \operatorname{Isom}_\xi:
 \widehat{\underline C}_{R\widehat\otimes_\Lambda R}
 \longrightarrow(\mathrm{Sets})
\]
is prorepresentable, where $R=P(\xi)$.

\underline{Proof.} If $\operatorname{Isom}_\xi$ is homogeneous,
Theorem \SGARefLink{sga7.sga7i.evi.n1-11}{1.11} supplies a minimally smooth functor
\[
 \Phi:\underline h_{R_1}\longrightarrow
 \underline{\operatorname{Isom}}_\xi.
\]
Lemma \SGARefLink{sga7.sga7i.evi.n2-15}{2.15} makes $\Phi$ a
$\widehat{\underline C}_{R\widehat\otimes_\Lambda R}$-equivalence,
and hence $\operatorname{Isom}_\xi$ is prorepresentable. It remains
to prove homogeneity. Consider a diagram
\[
\begin{tikzcd}[column sep=large,row sep=large]
 T_1 \arrow[dr] && T_2 \arrow[dl]\\
 & T &
\end{tikzcd}
\]
in $\widehat{\underline C}_{R\widehat\otimes_\Lambda R}$, with
$T_2\longrightarrow T$ surjective. Since $\underline A$ is
homogeneous, the fibre products
\[
 \xi_i(T_1)\times_{\xi_i(T)}\xi_i(T_2)
\]
exist in $\underline A$, and there are canonical isomorphisms
\[
 \xi_i(T_1\times_TT_2)\xrightarrow{\ \sim\ }
 \xi_i(T_1)\times_{\xi_i(T)}\xi_i(T_2)
 \qquad(i=1,2)
\]
(cf. \SGARefLink{sga7.sga7i.evi.n2-6.item-a}{2.6}(a)). Therefore
\[
 \operatorname{Isom}_\xi(T_1\times_TT_2)
 \longrightarrow
 \operatorname{Isom}_\xi(T_1)
 \times_{\operatorname{Isom}_\xi(T)}
 \operatorname{Isom}_\xi(T_2)
\]
is an isomorphism. \hfill$\square$

\underline{Theorem \SGARefTarget{sga7.sga7i.evi.n2-17}{2.17}.} Let $\underline A$ be a cofibered groupoid
over $\underline C_\Lambda$. Then it is prorepresentable by a
(normalized) smooth cogroupoid if and only if

%%%%%%%%%%  scan idx 79  |  folio 68  %%%%%%%%%%

(1) $\underline A$ is homogeneous;

(2) $[\underline A(k[\varepsilon])]$ and
$\operatorname{Aut}_{\underline A(k[\varepsilon])}(1_\varepsilon)$
are finite-dimensional vector spaces over $k$.

\underline{Proof.} Suppose first that $\underline A$ is
prorepresented by a smooth cogroupoid. Thus $\underline A$ is
$\underline C_\Lambda$-equivalent to
\[
 \underline h=\underline h_{(R_0,R_1,\ldots)}
\]
for a smooth cogroupoid $(R_0,R_1,\ldots)$ in
$\widehat{\underline C}_\Lambda$. Proposition \SGARefLink{sga7.sga7i.evi.n2-12}{2.12}(a) shows that
$\underline h$ is homogeneous. Moreover,
\begin{align*}
 [\underline h(k[\varepsilon])]
 &=\operatorname{Coker}\!\left(
   \operatorname{Hom}_{\Lambda\text{-alg}}(R_1,k[\varepsilon])
   \mathrel{\substack{\xrightarrow{\ s_1\ }\\[-.6ex]
                       \xrightarrow[\ s_2\ ]{}}}
   \operatorname{Hom}_{\Lambda\text{-alg}}(R_0,k[\varepsilon])
   \right),\\[1ex]
 \operatorname{Aut}_{\underline h}(1_\varepsilon)
 &=\operatorname{Ker}\!\left(
   \operatorname{Hom}_{\Lambda\text{-alg}}(R_1,k[\varepsilon])
   \xrightarrow{\ s_1\ }
   \operatorname{Hom}_{\Lambda\text{-alg}}(R_0,k[\varepsilon])
   \right)\\
 &=\operatorname{Ker}\!\left(
   \operatorname{Hom}_{\Lambda\text{-alg}}(R_1,k[\varepsilon])
   \xrightarrow{\ s_2\ }
   \operatorname{Hom}_{\Lambda\text{-alg}}(R_0,k[\varepsilon])
   \right).
\end{align*}
Both are finite-dimensional vector spaces over $k$, so the same is
true for $\underline A$.

Conversely, assume (1) and (2). Since $\underline A$ is homogeneous
and $[\underline A(k[\varepsilon])]$ is finite-dimensional, it has a
minimally versal formal object $\xi$. Put $P(\xi)=R_0$ and consider
\[
 \operatorname{Isom}_\xi:
 \widehat{\underline C}_{R_0\widehat\otimes_\Lambda R_0}
 \longrightarrow(\mathrm{Sets}),
\]
given by
\[
 T\longmapsto
 \operatorname{Isom}_{\underline A(T)}(\xi_1(T),\xi_2(T)),
\]
where $\xi_i(T)=\xi(f_i)$ and $f_i:R_0\longrightarrow T$ are the
composite maps in
\[
\begin{tikzcd}[column sep=large]
 R_0 \arrow[r,shift left=2,"s_1"]
     \arrow[r,shift right=2,"s_2"']
 & R_0\widehat\otimes_\Lambda R_0 \arrow[r] & T.
\end{tikzcd}
\]
Corollary \SGARefLink{sga7.sga7i.evi.n2-16}{2.16} prorepresents this functor by an object
\[
 R_1\in
 \operatorname{ob}(\widehat{\underline C}_{R_0\widehat\otimes_\Lambda R_0})
\]
and an isomorphism
$\xi_1(R_1)\xrightarrow{\varphi}\xi_2(R_1)$ for which
\[
 \Phi:\underline h_{R_1}\longrightarrow
 \underline{\operatorname{Isom}}_\xi
\]
is an isomorphism over
$\widehat{\underline C}_{R_0\widehat\otimes_\Lambda R_0}$. The pair
$(R_0,R_1)$ therefore defines a cogroupoid, and
\[
 (\xi,\varphi):\underline h_{(R_0,R_1,\ldots)}
 \longrightarrow\underline A
\]
is a $\underline C_\Lambda$-equivalence. Because $\xi$ is minimally
versal, the composite
\[
 \underline h_{(R_0,R_0,\ldots)}
 \xrightarrow{\ \epsilon\ }
 \underline h_{(R_0,R_1,\ldots)}
 \xrightarrow{\ (\xi,\varphi)\ }
 \underline A
\]
is minimally smooth. Since $(\xi,\varphi)$ is an equivalence,
$\epsilon$ is minimally smooth as well. Thus
$(R_0,R_1,\ldots)$ is a normalized smooth cogroupoid in
$\widehat{\underline C}_\Lambda$. \hfill$\square$

% END INCLUDED SGA7I_English_current_buildable_source_20260823/components/25_expose_VI_2_15_through_2_17.tex

% BEGIN INCLUDED SGA7I_English_current_buildable_source_20260823/components/26_expose_VI_2_18_through_2_22.tex
%%%%%%%%%%  scan idx 80  |  folio 69  %%%%%%%%%%

\underline{Corollary \SGARefTarget{sga7.sga7i.evi.n2-18}{2.18}.} A functor
$F:\underline C_\Lambda\longrightarrow(\mathrm{Sets})$ is
prorepresentable if and only if it is homogeneous and
$\dim F(k[\varepsilon])<\infty$.

\underline{Corollary \SGARefTarget{sga7.sga7i.evi.n2-19}{2.19}.} Every smooth cogroupoid in
$\widehat{\underline C}_\Lambda$ can be normalized up to isomorphism.

Given a cogroupoid $(R_0,R_1,\ldots)$ in
$\widehat{\underline C}_\Lambda$, set $\bar R_1=R_1/J$, where $J$ is
the ideal of $R_1$ generated by
\[
 \{s_1(r)-s_2(r)\mid r\in R_0\}.
\]
Then $(R_0,\bar R_1,\ldots)$ defines a formal group scheme over $R_0$
with the induced composition law.

Now let $\underline A$ be a cofibered groupoid over
$\underline C_\Lambda$, prorepresented by a normalized smooth
cogroupoid $(R_0,R_1,\ldots)$ via $(\xi,\varphi)$. Let
\[
 \underline C_k\lhook\joinrel\longrightarrow
 \underline C_{R_0}\lhook\joinrel\longrightarrow
 \underline C_{R_0\widehat\otimes_\Lambda R_0}
\]
be the embeddings corresponding to the surjections
\[
 R_0\widehat\otimes_\Lambda R_0\longrightarrow R_0\longrightarrow k,
\]
and let $\operatorname{Aut}_\xi$ and $\operatorname{Aut}^{o}$ be the
restrictions of $\operatorname{Isom}_\xi$ to
$\underline C_{R_0}$ and $\underline C_k$, respectively. Thus
$\operatorname{Aut}_\xi(S)$ is the automorphism group of $\xi(S)$ for
$S\in\operatorname{ob}(\underline C_{R_0})$, while
$\operatorname{Aut}^{o}(S)$ is the automorphism group of $1_S$ for
$S\in\operatorname{ob}(\underline C_k)$, where $1_S$ is the direct
image of $1\in\underline A(k)$ under $k\longrightarrow S$.

Since $R_1$ prorepresents $\operatorname{Isom}_\xi$ on
$\underline C_{R_0\widehat\otimes_\Lambda R_0}$, it follows that
$\bar R_1$ prorepresents $\operatorname{Aut}_\xi$ on
$\underline C_{R_0}$, and that
$k\otimes_{R_0}\bar R_1$ prorepresents $\operatorname{Aut}^{o}$ on
$\underline C_k$. Although $\underline A$ is prorepresented by the
smooth cogroupoid $(R_0,R_1,\ldots)$, neither
$(R_0,\bar R_1,\ldots)$ nor
$(k,k\otimes_{R_0}\bar R_1,\ldots)$ is smooth in general. More
precisely:

\underline{Corollary \SGARefTarget{sga7.sga7i.evi.n2-20}{2.20}.} Let $\underline A$ be a homogeneous
cofibered groupoid over $\underline C_\Lambda$, prorepresented by a
normalized smooth cogroupoid $(R_0,R_1,\ldots)$ via
$(\xi,\varphi)$.

\begin{enumerate}[label=(\alph*)]
\item $(R_0,\bar R_1,\ldots)$ and
$(k,k\otimes_{R_0}\bar R_1,\ldots)$ prorepresent
$\operatorname{Aut}_\xi$ and $\operatorname{Aut}^{o}$, respectively.

\item The following statements are equivalent:
\begin{enumerate}[label=(\roman*)]
\item the formal group scheme
$\underline h_{(R_0,\bar R_1,\ldots)}$ is formally smooth over $R_0$;
\item for every surjective arrow $a'\longrightarrow a$ in
$\underline A$, the map
\[
 \operatorname{Aut}_{P(a')}(a')\longrightarrow
 \operatorname{Aut}_{P(a)}(a)
\]
is surjective;
\item $R_0$ prorepresents the functor $[\underline A]$ over
$\underline C_\Lambda$.
\end{enumerate}
\end{enumerate}

%%%%%%%%%%  scan idx 81  |  folio 70  %%%%%%%%%%

\underline{Proof.} We have already noted (a). For (b), the
equivalence (i)$\Leftrightarrow$(ii) is clear, whereas
(ii)$\Leftrightarrow$(iii) follows from \SGARefLink{sga7.sga7i.evi.n2-7}{2.7} and \SGARefLink{sga7.sga7i.evi.n2-15}{2.15}, since
\[
 \widetilde\xi:\underline h_{R_0}\longrightarrow[\underline A]
\]
is minimally smooth. \hfill$\square$

\underline{Remark \SGARefTarget{sga7.sga7i.evi.n2-21}{2.21}.} Let $\underline A$ be a cofibered groupoid
over $\underline C_\Lambda$, and let
$\underline A|\underline C_k$ be the pullback of
$p:\underline A\longrightarrow\underline C_\Lambda$ along
$\underline C_k\lhook\joinrel\longrightarrow\underline C_\Lambda$.
If $\underline A$ is prorepresented by $(R_0,R_1,\ldots)$ via
$(\xi,\varphi)$, then $\underline A|\underline C_k$ is
prorepresented by
\[
 (k\otimes_\Lambda R_0,k\otimes_\Lambda R_1,\ldots)
\]
via $(k\otimes_\Lambda\xi,k\otimes_\Lambda\varphi)$. Here
$k\otimes_\Lambda\xi$ is the direct image of $\xi$ under
$R_0\longrightarrow k\otimes_\Lambda R_0$, and
$k\otimes_\Lambda\varphi$ is the image of $\varphi$ under
\[
 \operatorname{Isom}_{\underline A(R_0)}
   \bigl(\xi(s_1),\xi(s_2)\bigr)
 \longrightarrow
 \operatorname{Isom}_{\underline A(k\otimes_\Lambda R_1)}
   \bigl(k\otimes_\Lambda\xi(s_1),
         k\otimes_\Lambda\xi(s_2)\bigr).
\]
If $(R_0,R_1,\ldots)$ is normalized, then so is
$(k\otimes_\Lambda R_0,k\otimes_\Lambda R_1,\ldots)$; moreover, the
relative dimension of $R_1$ over $R_0$ equals that of
$k\otimes_\Lambda R_1$ over $k\otimes_\Lambda R_0$, since $R_1$ is
smooth over $R_0$.

\underline{Remark \SGARefTarget{sga7.sga7i.evi.n2-22}{2.22}.} Suppose that $\underline A$ is
prorepresented by smooth cogroupoids $(R_0,R_1,\ldots)$ and
$(R'_0,R'_1,\ldots)$, with $(R_0,R_1,\ldots)$ normalized. We obtain a
homomorphism
\[
 (i_0,i_1):(R_0,R_1,\ldots)\longrightarrow(R'_0,R'_1,\ldots).
\]
By \SGARefLink{sga7.sga7i.evi.n2-10}{2.10}(b), $R'_0$ is a formal power-series ring over $R_0$, and
$(i_0,i_1)$ induces an isomorphism
\[
 R'_0\otimes_{R_0}(R_0,R_1,\ldots)
 \xrightarrow{\ \sim\ }(R'_0,R'_1,\ldots).
\]

% END INCLUDED SGA7I_English_current_buildable_source_20260823/components/26_expose_VI_2_18_through_2_22.tex

% BEGIN INCLUDED SGA7I_English_current_buildable_source_20260823/components/27_expose_VI_section_3_opening_through_3_1_b.tex
%%%%%%%%%%  scan idx 81  |  folio 70  %%%%%%%%%%

\section*{\SGARefTarget{sga7.sga7i.evi.n3}{3}. The coherent sheaves $\underline D_X^i$}

One important functor which is usually not prorepresentable, but
nevertheless admits a formal versal object, is the deformation
functor. With this in view, we recall some basic facts about
extensions of commutative rings and their obstruction theory. This
section contains nothing new, except perhaps for its systematic use
of Grothendieck cohomology on a site. Its purpose is to make this
expos\'e as self-contained as possible; in fact, we prove somewhat
more than will be needed later. The basic results are stated in
Theorem \SGARefLink{sga7.sga7i.evi.n3-12}{3.12}. Throughout, all rings are commutative and have a unit.

%%%%%%%%%%  scan idx 82  |  folio 71  %%%%%%%%%%

Given a ring $R$, let $\underline C_R$ denote the category of
$R$-algebras. For every $A\in\operatorname{Ob}(\underline C_R)$, put
\[
 \operatorname{Cov}(A)=
 \{\text{surjective arrows }\varphi:A'\longrightarrow A
   \text{ in }\underline C_R\}.
\]
This defines a pretopology, and hence a site, on $\underline C_R$.
For every $A\in\operatorname{Ob}(\underline C_R)$, let
$\underline C_{A|R}$ be the category of arrows in $\underline C_R$
with target $A$. We endow $\underline C_{A|R}$ with the topology
induced by the forgetful functor
\[
 j_A:\underline C_{A|R}\longrightarrow\underline C_R.
\]
The category $\underline C_{A|R}$ has a final object
$(1_A:A\longrightarrow A)$ and an initial object
$(i_A:R\longrightarrow A)$, and it admits both fibre products and
coproducts. When there is no risk of confusion, we write simply $B$
for an object $(B\longrightarrow A)$ of $\underline C_{A|R}$.

Let $\mathscr C$ be an abelian category. A $\mathscr C$-valued sheaf
on $\underline C_{A|R}$ is, by definition, a functor
\[
 F:\underline C_{A|R}^{\mathrm{op}}\longrightarrow\mathscr C
\]
such that, for every surjective arrow $B'\longrightarrow B$ in
$\underline C_{A|R}$, the sequence
\[
 0\longrightarrow F(B)\longrightarrow F(B')
 \mathrel{\substack{\longrightarrow\\[-.6ex]\longrightarrow}}
 F(B'\mathbin{\times_B}B')
\]
is exact. The \underline{additive} presheaves are those which
preserve coproducts; thus
\[
 F(B_1\otimes_R B_2)=F(B_1)\oplus F(B_2).
\]
For example, an $A$-module $M$ defines the functor
\[
 \operatorname{Der}_R(-,M):
 \underline C_{A|R}^{\mathrm{op}}\longrightarrow(A\text{-mod}),
\]
which sends $B$ to the $R$-linear derivations
$\operatorname{Der}_R(B,M)$, with $M$ regarded as a $B$-module via
the structural map $B\longrightarrow A$. This is plainly an
additive sheaf on $\underline C_{A|R}$. Here $(A\text{-mod})$ denotes
the category of $A$-modules.

A presheaf $F$ on $\underline C_{A|R}$ with values in
$(A\text{-mod})$ is said to be of \underline{finite type} if $F(B)$
is an $A$-module of finite type whenever $B$ is an $R$-algebra of
finite type. For example, if $A$ is noetherian and $M$ is an
$A$-module of finite type, then $\operatorname{Der}_R(-,M)$ is an
additive sheaf of finite type on $\underline C_{A|R}$.

We record a little more notation and a few basic facts from
homological algebra.

\underline{\SGARefTarget{sga7.sga7i.evi.n3-1.item-a}{\SGARefAliasTarget{sga7.sga7i.evi.n3-1}3.1}(a).} Throughout the remainder of this section,
$\widetilde{\underline C}_{A|R}$ and
$\overset{\vee}{\underline C}_{A|R}$ denote, respectively, the
categories of sheaves and presheaves with values in
$(A\text{-mod})$. There is a left-exact inclusion functor
\[
 i_{A|R}:\widetilde{\underline C}_{A|R}
 \longrightarrow\overset{\vee}{\underline C}_{A|R}
\]
and an exact associated-sheaf functor
\[
 G_{A|R}:\overset{\vee}{\underline C}_{A|R}
 \longrightarrow\widetilde{\underline C}_{A|R},
\]
left adjoint to $i_{A|R}$.

%%%%%%%%%%  scan idx 83  |  folio 72  %%%%%%%%%%

An object $B\in\operatorname{Ob}(\underline C_{A|R})$ induces a
continuous functor
$j:\underline C_{B|R}\longrightarrow\underline C_{A|R}$ and hence
a commutative diagram
\begin{equation}\tag{*}
\begin{tikzcd}[column sep=huge,row sep=huge]
 \widetilde{\underline C}_{A|R}
   \arrow[r,"\widetilde j"] \arrow[d,"i_{A|R}"']
 & \widetilde{\underline C}_{B|R} \arrow[d] \\
 \overset{\vee}{\underline C}_{A|R}
   \arrow[r,"j^{\vee}"']
 & \overset{\vee}{\underline C}_{B|R}.
\end{tikzcd}
\end{equation}
The restriction functor $\widetilde j$, as well as $j^{\vee}$, is
exact, because the site $\underline C_{B|R}$ is simply the category
of arrows in $\underline C_{A|R}$ with target $B$, endowed with the
induced topology.

A ring homomorphism $S\longrightarrow R$ induces a continuous
functor
$\beta:\underline C_{A|R}\longrightarrow\underline C_{A|S}$ and
hence a commutative diagram
\begin{equation}\tag{**}
\begin{tikzcd}[column sep=huge,row sep=huge]
 \widetilde{\underline C}_{A|S}
   \arrow[r,"\widetilde\beta"] \arrow[d,"i_{A|S}"']
 & \widetilde{\underline C}_{A|R} \arrow[d] \\
 \overset{\vee}{\underline C}_{A|S}
   \arrow[r,"\beta^{\vee}"']
 & \overset{\vee}{\underline C}_{A|R}.
\end{tikzcd}
\end{equation}
Notice that $\widetilde\beta$ need not be exact.

\underline{\SGARefTarget{sga7.sga7i.evi.n3-1.item-b}{3.1}(b).} The right-derived functor of the inclusion
\[
 i_{A|R}:\widetilde{\underline C}_{A|R}
 \longrightarrow\overset{\vee}{\underline C}_{A|R}
\]
will be denoted by $\underline H^{\bullet}_{A|R}$. For every sheaf
$F$ on $\underline C_{A|R}$, put
\[
 H^{\bullet}(A|R,F)=
 [\underline H^{\bullet}_{A|R}(F)](A).
\]

For $B'\longrightarrow B$ in $\operatorname{Cov}(B)$, define
\[
 H^0(\{B'\longrightarrow B\},F)=
 \ker\!\left(
 F(B')\mathrel{\substack{\longrightarrow\\[-.6ex]\longrightarrow}}
 F(B'\mathbin{\times_B}B')
 \right).
\]
We denote the right-derived functors of
$H^0(\{B'\longrightarrow B\},-)$ by
$H^{\bullet}(\{B'\longrightarrow B\},-)$. These are the cohomology
groups of the \v{C}ech semisimplicial module

%%%%%%%%%%  scan idx 84  |  folio 73  %%%%%%%%%%

\[
\begin{tikzcd}[column sep=large]
 F(B')
   \arrow[r,shift left=2] \arrow[r,shift right=2]
 & F(B'\times_BB')
   \arrow[r,shift left=3] \arrow[r] \arrow[r,shift right=3]
 & F(B'\times_BB'\times_BB')
   \arrow[r,shift left=4.5] \arrow[r,shift left=1.5]
   \arrow[r,shift right=1.5] \arrow[r,shift right=4.5]
 & \cdots .
\end{tikzcd}
\]

The right-derived functor of
\[
 \check{\underline H}^{,0}_{A|R}:
 \overset{\vee}{\underline C}_{A|R}
 \longrightarrow\overset{\vee}{\underline C}_{A|R}
\]
will be denoted by $\check{\underline H}^{\,\bullet}_{A|R}$, where
\[
 [\check{\underline H}^{,0}_{A|R}(F)](B)=
 \varinjlim_{\{B'\to B\}\in\operatorname{Cov}(B)}
 H^0(\{B'\longrightarrow B\},F).
\]
For every presheaf $F$ on $\underline C_{A|R}$, put
\[
 \check H^{\bullet}(A|R,F)=
 [\check{\underline H}^{\,\bullet}_{A|R}(F)](A)=
 \varinjlim_{\{A'\to A\}\in\operatorname{Cov}(A)}
 H^{\bullet}(\{A'\longrightarrow A\},F).
\]
For every sheaf $F$ on $\underline C_{A|R}$ there is the usual
spectral sequence
\[
 \check H^p\bigl(A|R,\underline H^q_{A|R}(F)\bigr)
 \Longrightarrow H^{p+q}(A|R,F).
\]
Here
$\check H^0(A|R,\underline H^q_{A|R}(F))=0$ for every $q>0$.
Consequently,
\[
 \check H^1(A|R,F)=H^1(A|R,F),
\]
and the sequence
\[
 0\longrightarrow\check H^2(A|R,F)
 \longrightarrow H^2(A|R,F)
 \longrightarrow
 \check H^1\bigl(A|R,\underline H^1_{A|R}(F)\bigr)
 \longrightarrow\check H^3(A|R,F)
\]
is exact.

% END INCLUDED SGA7I_English_current_buildable_source_20260823/components/27_expose_VI_section_3_opening_through_3_1_b.tex

% BEGIN INCLUDED SGA7I_English_current_buildable_source_20260823/components/28_expose_VI_3_1_c_through_3_4.tex
\underline{\SGARefTarget{sga7.sga7i.evi.n3-1.item-c}{3.1}(c).} If $B$ is an object of
$\underline C_{A|R}$, diagram $(*)$ in \SGARefLink{sga7.sga7i.evi.n3-1.item-a}{3.1}(a), together with the
exactness of $\widetilde j$ and $j^{\vee}$, gives a canonical
isomorphism
\[
 H^{\bullet}(B|R,\widetilde jF)=
 [\underline H^{\bullet}_{A|R}(F)](B).
\]
On the other hand, a ring homomorphism $S\longrightarrow R$ induces
diagram $(**)$ in \SGARefLink{sga7.sga7i.evi.n3-1.item-a}{3.1}(a). The functor $\widetilde\beta$ need not be
exact; however, $\beta$ preserves fibre products, and we obtain a
spectral sequence
\[
 \underline H^p_{A|R}(R^q\widetilde\beta F)
 \Longrightarrow
 \beta^{\vee}\underline H^{p+q}_{A|S}(F),
\]
or, on evaluating at $A$,
\[
 H^p(A|R,R^q\widetilde\beta F)
 \Longrightarrow H^{p+q}(A|S,F).
\]
Since there is a canonical isomorphism
\[
 \widetilde\beta=
 G_{A|R}\circ\beta^{\vee}\circ i_{A|S}
\]
and the functor
$\beta^{\#}=G_{A|R}\circ\beta^{\vee}$ is exact, we have

%%%%%%%%%%  scan idx 85  |  folio 74  %%%%%%%%%%

\[
 R^{\bullet}\widetilde\beta F=
 \beta^{\#}\underline H^{\bullet}_{A|S}(F).
\]
Thus the preceding spectral sequence may be written
\[
 H^p\bigl(A|R,\beta^{\#}\underline H^q_{A|S}(F)\bigr)
 \Longrightarrow H^{p+q}(A|S,F).
\]

\underline{Lemma \SGARefTarget{sga7.sga7i.evi.n3-2}{3.2}.} Let
$B\in\operatorname{Ob}(\underline C_{A|R})$, and choose a
surjective arrow $P\longrightarrow B$ in $\underline C_{B|R}$,
where $P$ is a polynomial algebra over $R$. Then, for every
$F\in\operatorname{Ob}(\overset{\vee}{\underline C}_{A|R})$, the
natural map
\[
 H^{\bullet}(\{P\longrightarrow B\},F)
 \longrightarrow[\check{\underline H}^{\,\bullet}(F)](B)
\]
is an isomorphism.

\underline{Proof.} It suffices to prove that the natural map
\[
 H^0(\{P\longrightarrow B\},F)
 \longrightarrow[\check{\underline H}^{,0}(F)](B)
 =\varinjlim_{\{B'\to B\}\in\operatorname{Cov}(B)}
   H^0(\{B'\longrightarrow B\},F)
\]
is an isomorphism. This is clear because, for every
$\{B'\longrightarrow B\}\in\operatorname{Cov}(B)$, there is a map
of covers
\[
 \{P\longrightarrow B\}\longrightarrow
 \{B'\longrightarrow B\}.
\]
\hfill$\square$

\underline{Corollary \SGARefTarget{sga7.sga7i.evi.n3-3}{3.3}.}

\begin{enumerate}[label=(\alph*)]
\item Suppose that $A$ is noetherian. If
$F\in\operatorname{Ob}(\widetilde{\underline C}_{A|R})$ is of
finite type, then $\underline H^n(F)$ is of finite type for every
$n$.

\item If $P$ is a polynomial algebra over $R$, then
\[
 H^n(P|R,F)=0
\]
for every $n>0$ and every
$F\in\operatorname{Ob}(\widetilde{\underline C}_{P|R})$.

\item A sequence
\[
 0\longrightarrow F'\longrightarrow F\longrightarrow F''
 \longrightarrow0
\]
in $\widetilde{\underline C}_{A|R}$ is exact if and only if
\[
 0\longrightarrow F'(P)\longrightarrow F(P)
 \longrightarrow F''(P)\longrightarrow0
\]
is exact for every polynomial $R$-algebra
$P\in\operatorname{Ob}(\underline C_{A|R})$.
\end{enumerate}

\underline{Proof.}

(a) In view of the spectral sequence
\[
 \check H^p\bigl(B|R,\underline H^q(F)\bigr)
 \Longrightarrow H^{p+q}(B|R,F)
\]
and the equality
$\check H^0(\underline H^q(F))=0$ for $q>0$, it is enough to show
that $\check{\underline H}^{,n}(G)$ is of finite type for every
$n$ and every presheaf $G$ of finite type. By Lemma \SGARefLink{sga7.sga7i.evi.n3-2}{3.2},
\[
 \check H^{\bullet}(B|R,G)=
 H^{\bullet}(\{P\longrightarrow B\},G),
\]
where $P$ is a polynomial algebra over $R$ and
$\{P\longrightarrow B\}\in\operatorname{Cov}(B)$. The latter is
the cohomology of the complex

%%%%%%%%%%  scan idx 86  |  folio 75  %%%%%%%%%%

\[
 \cdots
 G(P\times_BP\times_BP)
 \mathrel{\substack{\longleftarrow\\[-.6ex]\longleftarrow\\[-.6ex]\longleftarrow}}
 G(P\times_BP)
 \mathrel{\substack{\longleftarrow\\[-.6ex]\longleftarrow}}
 G(P).
\]
If $B$ is an $R$-algebra of finite type, $P$ may be chosen of
finite type over $R$. All the iterated fibre products
$P\times_B\cdots\times_BP$ are then $R$-algebras of finite type.
If $G$ is of finite type, their images under $G$ are $A$-modules of
finite type. Hence
$H^{\bullet}(\{P\longrightarrow B\},G)$ is an $A$-module of finite
type because $A$ is noetherian.

(b) Let $P$ be a polynomial algebra over $R$. Every surjective
$R$-algebra map $B\longrightarrow P$ admits a section, so
\[
 H^i(\{B\longrightarrow P\},G)=0\qquad(i>0).
\]
It follows that $\check H^i(P|R,G)=0$ for every $i>0$ and every
presheaf $G$ on $\underline C_{P|R}$. Consequently the spectral
sequence
\[
 \check H^p\bigl(P|R,\underline H^q_{P|R}(F)\bigr)
 \Longrightarrow H^{p+q}(P|R,F)
\]
degenerates, and
\[
 \check H^0\bigl(P|R,\underline H^n(F)\bigr)
 \simeq H^n(P|R,F)
\]
for every $n$. For a sheaf $F$, however,
$\check H^0(P|R,\underline H^n(F))=0$ whenever $n>0$. Thus
$H^n(P|R,F)=0$ for every $n>0$.

(c) If
$0\longrightarrow F'\longrightarrow F\longrightarrow F''
\longrightarrow0$ is exact in
$\widetilde{\underline C}_{A|R}$, then the sequence
\[
 0\longrightarrow F'(B)\longrightarrow F(B)
 \longrightarrow F''(B)
 \longrightarrow H^1(B|R,F')\longrightarrow H^1(B|R,F)
 \longrightarrow\cdots
\]
is exact. If $P\in\operatorname{Ob}(\underline C_{A|R})$ is a
polynomial algebra over $R$, then $H^1(P|R,F')=0$ by (b), and hence
\[
 0\longrightarrow F'(P)\longrightarrow F(P)
 \longrightarrow F''(P)\longrightarrow0
\]
is exact.

Conversely, let $\underline P_{A|R}$ be the full subcategory of
$\underline C_{A|R}$ with
\[
 \operatorname{Ob}(\underline P_{A|R})=
 \{P\in\operatorname{Ob}(\underline C_{A|R})\mid
   P\text{ is a polynomial algebra over }R\},
\]
endowed with the induced topology, and let
$p:\underline P_{A|R}\longrightarrow\underline C_{A|R}$ be the
inclusion. Every object of $\underline C_{A|R}$ is covered by an
object of $\underline P_{A|R}$; the comparison lemma therefore
shows that
\[
 \widetilde p:\widetilde{\underline C}_{A|R}
 \longrightarrow\widetilde{\underline P}_{A|R}
\]
is an equivalence of categories. If the sequence obtained by
evaluating $F'\longrightarrow F\longrightarrow F''$ at every
$P\in\operatorname{Ob}(\underline P_{A|R})$ is exact, then
\[
 0\longrightarrow\widetilde pF'\longrightarrow\widetilde pF
 \longrightarrow\widetilde pF''\longrightarrow0
\]
is exact in $\widetilde{\underline P}_{A|R}$, and therefore the
original sequence is exact in
$\widetilde{\underline C}_{A|R}$. \hfill$\square$

%%%%%%%%%%  scan idx 87  |  folio 76  %%%%%%%%%%

\underline{Remark \SGARefTarget{sga7.sga7i.evi.n3-4}{3.4}.} Every surjective arrow in
$\underline P_{A|R}$ admits a section. It follows at once that every
presheaf on $\underline P_{A|R}$ is a sheaf; that is,
\[
 \widetilde{\underline P}_{A|R}=
 \overset{\vee}{\underline P}_{A|R}.
\]
Consequently, the restriction functor
\[
 \widetilde{\underline C}_{A|R}
 \longrightarrow\overset{\vee}{\underline P}_{A|R}
\]
is an equivalence of categories.

% END INCLUDED SGA7I_English_current_buildable_source_20260823/components/28_expose_VI_3_1_c_through_3_4.tex

% BEGIN INCLUDED SGA7I_English_current_buildable_source_20260823/components/29_expose_VI_3_5_through_3_6.tex
\underline{Lemma \SGARefTarget{sga7.sga7i.evi.n3-5}{3.5}.} Consider a diagram of $R$-algebras
\[
\begin{tikzcd}[column sep=large,row sep=large]
 & C \arrow[d]\\
 B' \arrow[r,"\text{surjective}"'] & B
\end{tikzcd}
\]
and let $S$ be an $R$-algebra.

\begin{enumerate}[label=(\roman*)]
\item If $\operatorname{Tor}_1^R(B,S)=0$, then the canonical map
\[
 (B'\times_BC)\otimes_RS
 \longrightarrow
 (B'\otimes_RS)\times_{B\otimes_RS}(C\otimes_RS)
\]
is an isomorphism.

\item Suppose that $B'$ is flat over $R$. If
\[
 \operatorname{Tor}_i^R(C,S)=0\quad(0<i<n)
 \qquad\text{and}\qquad
 \operatorname{Tor}_i^R(B,S)=0\quad(0<i\leq n),
\]
then
\[
 \operatorname{Tor}_i^R(B'\times_BC,S)=0\qquad(0<i<n).
\]
\end{enumerate}

\underline{Proof.} Let
$0\longrightarrow J\longrightarrow B'\longrightarrow B
\longrightarrow0$ be exact.

(i) If $\operatorname{Tor}_1^R(B,S)=0$, then
\[
 0\longrightarrow J\otimes_RS\longrightarrow B'\otimes_RS
 \longrightarrow B\otimes_RS\longrightarrow0
\]
is exact, and hence so is
\[
 0\longrightarrow J\otimes_RS
 \longrightarrow
 (B'\otimes_RS)\times_{B\otimes_RS}(C\otimes_RS)
 \longrightarrow C\otimes_RS\longrightarrow0.
\]
We therefore obtain a commutative diagram with exact rows:
\begin{center}
\scriptsize
\begin{tikzcd}[column sep=small,row sep=large]
 & J\otimes_RS \arrow[r] \arrow[d,equal]
 & (B'\times_BC)\otimes_RS \arrow[r] \arrow[d]
 & C\otimes_RS \arrow[r] \arrow[d,equal]
 & 0\\
 0 \arrow[r] & J\otimes_RS \arrow[r]
 & (B'\otimes_RS)\times_{B\otimes_RS}(C\otimes_RS) \arrow[r]
 & C\otimes_RS \arrow[r] & 0.
\end{tikzcd}
\end{center}

%%%%%%%%%%  scan idx 88  |  folio 77  %%%%%%%%%%

It follows that
\[
 (B'\times_BC)\otimes_RS
 \longrightarrow
 (B'\otimes_RS)\times_{B\otimes_RS}(C\otimes_RS)
\]
is an isomorphism.

(ii) Since $B'$ is flat over $R$ and
$\operatorname{Tor}_i^R(B,S)=0$ for $0<i\leq n$, the long exact
Tor sequence gives
\[
 \operatorname{Tor}_i^R(J,S)=0\qquad(0<i<n).
\]
Applying the same sequence to
\[
 0\longrightarrow J\longrightarrow B'\times_BC
 \longrightarrow C\longrightarrow0
\]
and using the hypothesis on $C$ gives
\[
 \operatorname{Tor}_i^R(B'\times_BC,S)=0\qquad(0<i<n).
\]
\hfill$\square$

\underline{Corollary \SGARefTarget{sga7.sga7i.evi.n3-6}{3.6}.}

\begin{enumerate}[label=(\alph*)]
\item Let $A$ and $S$ be $R$-algebras. If
$\operatorname{Tor}_1^R(A,S)=0$, then for every
$G\in\operatorname{Ob}
(\overset{\vee}{\underline C}_{A\otimes_RS|S})$ the canonical map
\[
 \check H^{\bullet}(A\otimes_RS|S,G)
 \longrightarrow\check H^{\bullet}(A|R,\check SG)
\]
is an isomorphism. Here
\[
 S:\underline C_{A|R}\longrightarrow
 \underline C_{A\otimes_RS|S},
 \qquad X\longmapsto X\otimes_RS.
\]
In particular, if $S$ is flat over $R$, then
\[
 \check S\,\check{\underline H}^{\,\bullet}(G)
 \longrightarrow
 \check{\underline H}^{\,\bullet}(\check SG)
\]
is an isomorphism for every such $G$.

\item Let $P\longrightarrow B$ be a surjective map of $R$-algebras,
where $P$ is a polynomial algebra over $R$. If
\[
 \operatorname{Tor}_i^R(B,S)=0\qquad(0<i\leq n),
\]
then, for every iterated fibre product of copies of $P$ over $B$,
\[
 \operatorname{Tor}_i^R
 (P\times_B P\times_B\cdots\times_B P,S)=0
 \qquad(0<i<n).
\]
\end{enumerate}

\underline{Proof.}

(a) Given $B\in\operatorname{Ob}(\underline C_{A|R})$, choose a
cover $\{P\longrightarrow B\}\in\operatorname{Cov}(B)$ with $P$ a
polynomial algebra over $R$. If
$\operatorname{Tor}_1^R(B,S)=0$, Lemma \SGARefLink{sga7.sga7i.evi.n3-5}{3.5}(i) shows that the natural
map
\[
 \left(\overbrace{P\times_B\cdots\times_BP}^{n}\right)\otimes_RS
 \longrightarrow
 \overbrace{(P\otimes_RS)\times_{B\otimes_RS}\cdots
 \times_{B\otimes_RS}(P\otimes_RS)}^{n}
\]
is an isomorphism for every $n$. Consequently,
\begin{align*}
 \check H^{\bullet}(B|R,\check SG)
 &=H^{\bullet}(\{P\longrightarrow B\},\check SG)\\
 &=H^{\bullet}(\{P\otimes_RS\longrightarrow B\otimes_RS\},G)\\
 &=\check H^{\bullet}(B\otimes_RS|S,G).
\end{align*}
If $S$ is flat over $R$, then
$\operatorname{Tor}_1^R(X,S)=0$ for every
$X\in\operatorname{Ob}(\underline C_{A|R})$. The preceding
isomorphism is therefore natural in $X$, and gives
\[
 \check S\,\check{\underline H}^{\,\bullet}(G)
 \xrightarrow{\ \sim\ }
 \check{\underline H}^{\,\bullet}(\check SG).
\]

(b) Put
\[
 P^{(m)}=\overbrace{P\times_B\cdots\times_BP}^{m+1}.
\]
For $m=0$ there is nothing to prove, since $P$ is flat over $R$.
Suppose inductively that
$\operatorname{Tor}_i^R(P^{(m-1)},S)=0$ for $0<i<n$. Since
$P\longrightarrow B$ is surjective and
$P^{(m)}=P\times_BP^{(m-1)}$, Lemma \SGARefLink{sga7.sga7i.evi.n3-5}{3.5}(ii) gives
\[
 \operatorname{Tor}_i^R(P^{(m)},S)=0\qquad(0<i<n).
\]
\hfill$\square$

% END INCLUDED SGA7I_English_current_buildable_source_20260823/components/29_expose_VI_3_5_through_3_6.tex

% BEGIN INCLUDED SGA7I_English_current_buildable_source_20260823/components/30_expose_VI_proposition_3_7.tex
%%%%%%%%%%  scan idx 89  |  folio 78  %%%%%%%%%%

\underline{Proposition \SGARefTarget{sga7.sga7i.evi.n3-7}{3.7}.}

\begin{enumerate}[label=(\alph*)]
\item Let $A$ and $S$ be $R$-algebras. If
\[
 \operatorname{Tor}_i^R(A,S)=0\qquad(0<i\leq n),
\]
then, for every sheaf $F$ on
$\underline C_{A\otimes_RS|S}$, the canonical map
\[
 H^i(A\otimes_RS|S,F)
 \longrightarrow H^i(A|R,\widetilde SF)
\]
is an isomorphism for every $i\leq n$. Here
\[
 S:\underline C_{A|R}\longrightarrow
 \underline C_{A\otimes_RS|S},
 \qquad X\longmapsto X\otimes_RS.
\]

\item Let $P$ be a polynomial algebra over $R$. Then, for every
additive sheaf $F$ on $\underline C_{A\otimes_RP|R}$, the
canonical map
\[
 H^n(A\otimes_RP|R,F)\longrightarrow H^n(A|R,F)
\]
is an isomorphism for every $n>0$.

\item Let $S\longrightarrow R\longrightarrow A$ be ring
homomorphisms. Then, for every additive sheaf $F$ on
$\underline C_{A|S}$, there is an exact sequence
\begin{multline*}
0\longrightarrow H^0(A|R,F_R)\longrightarrow H^0(A|S,F)
\longrightarrow H^0(R|S,F)\longrightarrow H^1(A|R,F_R)\\
\longrightarrow H^1(A|S,F)\longrightarrow H^1(R|S,F)
\longrightarrow\cdots\longrightarrow H^n(A|R,F_R)\\
\longrightarrow H^n(A|S,F)\longrightarrow H^n(R|S,F)
\longrightarrow H^{n+1}(A|R,F_R)\longrightarrow\cdots,
\end{multline*}
where $F_R$ is the sheaf on $\underline C_{A|R}$ defined by
\[
 F_R(X)=\ker\bigl(F(X)\longrightarrow F(R)\bigr)
\]
for every $X\in\operatorname{Ob}(\underline C_{A|R})$, via the
forgetful functor
$\underline C_{A|R}\longrightarrow\underline C_{A|S}$.
\end{enumerate}

\underline{Proof.}

(a) Consider the continuous functor
\[
 S:\underline C_{A|R}\longrightarrow
 \underline C_{A\otimes_RS|S},
 \qquad S(X)=X\otimes_RS.
\]
Since $S$ sends polynomial algebras over $R$ to polynomial
algebras over $S$, the functor
\[
 \widetilde S:\widetilde{\underline C}_{A\otimes_RS|S}
 \longrightarrow\widetilde{\underline C}_{A|R}
\]
is exact by \SGARefLink{sga7.sga7i.evi.n3-3}{3.3}(c). We therefore obtain the canonical map
\[
 \check S\,\underline H^{\bullet}_{A\otimes_RS|S}(F)
 \longrightarrow
 \underline H^{\bullet}_{A|R}(\widetilde SF)
\]
and the following commutative diagram of spectral sequences:
\begin{center}
\small
\begin{tikzcd}[column sep=huge,row sep=large]
E_S^{p,q}(S(X))=
 \check H^p\bigl(X\otimes_RS|S,
 \underline H^q_{A\otimes_RS|S}(F)\bigr)
 \arrow[r,Rightarrow] \arrow[d]
& H^{p+q}(X\otimes_RS|S,F) \arrow[dd]\\
\check H^p\bigl(X|R,
 \check S\,\underline H^q_{A\otimes_RS|S}(F)\bigr)
 \arrow[d] & \\
E_R^{p,q}(X)=
 \check H^p\bigl(X|R,
 \underline H^q_{A|R}(\widetilde SF)\bigr)
 \arrow[r,Rightarrow]
& H^{p+q}(X|R,\widetilde SF).
\end{tikzcd}
\end{center}

%%%%%%%%%%  scan idx 90  |  folio 79  %%%%%%%%%%

By \SGARefLink{sga7.sga7i.evi.n3-6}{3.6}(a), the map
\[
 E_S^{p,q}(S(X))\longrightarrow
 \check H^p\bigl(X|R,
 \check S\,\underline H^q_{A\otimes_RS|S}(F)\bigr)
\]
is an isomorphism whenever
$\operatorname{Tor}_1^R(X,S)=0$. Suppose inductively that
\[
 [\check S\,\underline H^q_{A\otimes_RS|S}(F)](X)
 \longrightarrow
 [\underline H^q_{A|R}(\widetilde SF)](X)
\]
is an isomorphism for every $q<n$ and every
$X\in\operatorname{Ob}(\underline C_{A|R})$ satisfying
\[
 \operatorname{Tor}_i^R(X,S)=0\qquad(0<i\leq n).
\]
It follows from \SGARefLink{sga7.sga7i.evi.n3-6}{3.6}(b) that
$E_S^{p,q}(S(X))\longrightarrow E_R^{p,q}(X)$ is an isomorphism
for every $q<n$ and every $p$, under the same Tor hypothesis.
Moreover,
$E_S^{0,n}(S(X))\longrightarrow E_R^{0,n}(X)$ is an isomorphism,
since both sides are zero. Hence
\[
 H^n(X\otimes_RS|S,F)
 \longrightarrow H^n(X|R,\widetilde SF)
\]
is an isomorphism whenever
$\operatorname{Tor}_i^R(X,S)=0$ for $0<i\leq n$. Equivalently,
\[
 [\check S\,\underline H^i_{A\otimes_RS|S}(F)](X)
 \longrightarrow
 [\underline H^i_{A|R}(\widetilde SF)](X)
\]
is an isomorphism for every $i\leq n$ under that hypothesis.

(b) Let
\[
 f:\underline C_{A|R}\longrightarrow
 \underline C_{A\otimes_RP|R},
 \qquad f(X)=X\otimes_RP.
\]
Choose a surjective arrow $P'\longrightarrow X$ in
$\underline C_{A|R}$, with $P'$ polynomial over $R$. Then
\[
 P'\otimes_RP\longrightarrow X\otimes_RP
\]
is a surjective arrow in $\underline C_{A\otimes_RP|R}$, and
$P'\otimes_RP$ is again polynomial over $R$. Since $P$ is flat
over $R$, \SGARefLink{sga7.sga7i.evi.n3-5}{3.5}(i) gives
\begin{align*}
 \check H^{\bullet}(X|R,\check fG)
 &=H^{\bullet}(\{P'\longrightarrow X\},\check fG)\\
 &=H^{\bullet}(\{P'\otimes_RP\longrightarrow X\otimes_RP\},G)\\
 &=\check H^{\bullet}(X\otimes_RP|R,G)
\end{align*}
for every
$G\in\operatorname{Ob}(\overset{\vee}{\underline C}_{A\otimes_RP|R})$.
Thus $\widetilde f$ sends acyclic sheaves to acyclic sheaves. On
the other hand, $f$ sends polynomial algebras over $R$ to
polynomial algebras over $R$, so
\[
 \widetilde f:\widetilde{\underline C}_{A\otimes_RP|R}
 \longrightarrow\widetilde{\underline C}_{A|R}
\]
is exact by \SGARefLink{sga7.sga7i.evi.n3-3}{3.3}(c). It follows that
\[
 H^{\bullet}(A\otimes_RP|R,F)
 \longrightarrow H^{\bullet}(A|R,\widetilde fF)
\]
is an isomorphism for every
$F\in\operatorname{Ob}(\widetilde{\underline C}_{A\otimes_RP|R})$.

For each $X\in\operatorname{Ob}(\underline C_{A|R})$ there is a
canonical map
\[
 (\widetilde fF)(X)=F(X\otimes_RP)
 \longrightarrow F(X)\oplus F(P),
\]
and hence a canonical map
$\widetilde fF\longrightarrow F\oplus F(P)$, where $F(P)$ denotes
the constant sheaf and $F$ on the right denotes the restriction of
$F$ to $\underline C_{A|R}$. If $F$ is additive, this map is an
isomorphism. Therefore, for every $n>0$,
\[
 H^n(A|R,\widetilde fF)
 \xrightarrow{\ \sim\ }
 H^n(A|R,F\oplus F(P))=H^n(A|R,F),
\]
because the constant sheaf $F(P)$ is flasque. This proves (b).

%%%%%%%%%%  scan idx 91  |  folio 80  %%%%%%%%%%

(c) This argument was communicated by Rinehart. Consider the
forgetful functor
$\beta:\underline C_{A|R}\longrightarrow\underline C_{A|S}$.
There is a spectral sequence
\[
 E^{p,q}=H^p\bigl(A|R,
 \beta^{\#}\underline H^q_{A|S}(F)\bigr)
 \Longrightarrow H^{p+q}(A|S,F),
\]
where
\[
 \beta^{\#}=G_{A|R}\circ\beta^{\vee},
 \qquad
 G_{A|R}:\overset{\vee}{\underline C}_{A|R}
 \longrightarrow\widetilde{\underline C}_{A|R}
\]
is the associated-sheaf functor.

Suppose now that $F$ is additive. We claim that
$\beta^{\#}\underline H^q_{A|S}(F)$ is a constant sheaf for every
$q>0$. Indeed, if $P$ is a polynomial algebra over $R$, then
$P=R\otimes_SP'$ for a polynomial algebra $P'$ over $S$, and
\[
 [\beta^{\#}\underline H^q_{A|S}(F)](P)
 =H^q(P|S,F)
 =H^q(R\otimes_SP'|S,F)
 \longrightarrow H^q(R|S,F)
\]
is an isomorphism for $q>0$ by (b), and is surjective for $q=0$.
By \SGARefLink{sga7.sga7i.evi.n3-3}{3.3}(c), the canonical map
\[
 \beta^{\#}\underline H^q_{A|S}(F)
 \longrightarrow H^q(R|S,F)
\]
to the corresponding constant sheaf is therefore an isomorphism
for $q>0$ and a surjection for $q=0$. Hence $E^{p,q}=0$ whenever
$p,q>0$, and the spectral sequence yields the exact sequence
\[
 (*)\quad
 0\longrightarrow E^{1,0}\longrightarrow H^1
 \longrightarrow E^{0,1}\longrightarrow E^{2,0}
 \longrightarrow H^2\longrightarrow E^{0,2}
 \longrightarrow E^{3,0}\longrightarrow H^3\longrightarrow\cdots.
\]
Here
\[
 E^{0,q}=[\beta^{\#}\underline H^q_{A|S}(F)](A)
 \xrightarrow{\ \sim\ }H^q(R|S,F)\quad(q>0),
 \qquad
 E^{p,0}=H^p(A|R,\widetilde\beta F).
\]
On the other hand, the exact sequence
\[
 0\longrightarrow F_R\longrightarrow\widetilde\beta F
 \longrightarrow F(R)\longrightarrow0
\]
gives the exact diagram
\begin{center}
\scriptsize
\begin{tikzcd}[column sep=small,row sep=large]
0 \arrow[r]
& H^0(A|R,F_R) \arrow[r]
& H^0(A|R,\widetilde\beta F) \arrow[r] \arrow[d,equal]
& H^0(A|R,F(R)) \arrow[r] \arrow[d,equal]
& H^1(A|R,F_R) \arrow[r]
& H^1(A|R,\widetilde\beta F) \arrow[r]
& 0\\
&& H^0(A|S,F) & H^0(R|S,F) &&&
\end{tikzcd}
\end{center}
labelled $(**)$ in the source. Combining $(*)$ and $(**)$ gives
the desired exact sequence. \hfill$\square$

% END INCLUDED SGA7I_English_current_buildable_source_20260823/components/30_expose_VI_proposition_3_7.tex

% BEGIN INCLUDED SGA7I_English_current_buildable_source_20260823/components/31_expose_VI_corollaries_3_8_and_3_9.tex
%%%%%%%%%%  scan idx 92  |  folio 81  %%%%%%%%%%

\underline{Corollary \SGARefTarget{sga7.sga7i.evi.n3-8}{3.8}.}

\begin{enumerate}[label=(\alph*)]
\item Let $A$ and $B$ be $R$-algebras such that
\[
 \operatorname{Tor}_i^R(A,B)=0\qquad(0<i\leq n).
\]
Then, for every additive sheaf
$F\in\operatorname{Ob}(\widetilde{\underline C}_{A\otimes_RB|R})$,
the canonical map
\[
 H^i(A\otimes_RB|R,F)
 \longrightarrow H^i(A|R,F)\oplus H^i(B|R,F)
\]
is an isomorphism for every $i\leq n$.

\item Let $A$ be an $R$-algebra such that
$\operatorname{Tor}_i^R(A,A)=0$ for every $i>0$. If the map
$A\otimes_RA\longrightarrow A$ is an isomorphism, then, for every
additive sheaf
$F\in\operatorname{Ob}(\widetilde{\underline C}_{A|R})$, one has
\[
 H^i(A|R,F)=0
\]
for every $i$.
\end{enumerate}

\underline{Proof.}

(a) Let
$F\in\operatorname{Ob}(\widetilde{\underline C}_{A\otimes_RB|R})$
be additive. Applied to
$R\longrightarrow A\longrightarrow A\otimes_RB$, Proposition
\SGARefLink{sga7.sga7i.evi.n3-7}{3.7}(c) gives the exact sequence
\[
 (*)\quad\cdots\longrightarrow
 H^i(A\otimes_RB|A,F_A)
 \longrightarrow H^i(A\otimes_RB|R,F)
 \longrightarrow H^i(A|R,F)\longrightarrow\cdots.
\]
Consider the commutative diagram
\[
\begin{tikzcd}[column sep=huge,row sep=huge]
H^i(A\otimes_RB|A,F_A) \arrow[r] \arrow[d]
& H^i(A\otimes_RB|R,F) \arrow[d]\\
H^i(B|R,\widetilde A F_A) \arrow[r,"\approx"]
& H^i(B|R,F).
\end{tikzcd}
\]
Here
\[
 A:\underline C_{B|R}\longrightarrow
 \underline C_{A\otimes_RB|A},
 \qquad X\longmapsto A\otimes_RX,
\]
and the bottom map is induced by
$\widetilde A F_A\longrightarrow F$, which is an isomorphism
because $F$ is additive. By functoriality, the commutativity of the
diagram reduces to the case $i=0$.

If $\operatorname{Tor}_i^R(A,B)=0$ for $0<i\leq n$, then
Proposition \SGARefLink{sga7.sga7i.evi.n3-7}{3.7}(a) shows that the left vertical map is an
isomorphism for $i\leq n$. Consequently, both composites
\[
 H^i(A\otimes_RB|A,F_A)
 \longrightarrow H^i(A\otimes_RB|R,F)
 \longrightarrow H^i(B|R,F)
\]
and, by symmetry,

%%%%%%%%%%  scan idx 93  |  folio 82  %%%%%%%%%%

\[
 H^i(A\otimes_RB|B,F_B)
 \longrightarrow H^i(A\otimes_RB|R,F)
 \longrightarrow H^i(A|R,F)
\]
are isomorphisms for $i\leq n$. It follows that $(*)$ breaks into
split short exact sequences
\[
 0\longrightarrow H^i(A\otimes_RB|A,F_A)
 \longrightarrow H^i(A\otimes_RB|R,F)
 \longrightarrow H^i(A|R,F)\longrightarrow0
\]
for $i\leq n$, and hence that
\[
 H^i(A\otimes_RB|R,F)
 \longrightarrow H^i(A|R,F)\oplus H^i(B|R,F)
\]
is an isomorphism in that range.

(b) By (a), the map
\[
 H^i(A\otimes_RA|R,F)
 \longrightarrow H^i(A|R,F)\oplus H^i(A|R,F)
\]
is an isomorphism for every $i$. On the other hand, the assumed
isomorphism $A\otimes_RA\longrightarrow A$ induces an isomorphism
\[
 H^i(A|R,F)\longrightarrow H^i(A\otimes_RA|R,F).
\]
Their composite is the diagonal map
\[
 H^i(A|R,F)\longrightarrow
 H^i(A|R,F)\oplus H^i(A|R,F).
\]
Since it is an isomorphism, $H^i(A|R,F)=0$ for every $i$.
\hfill$\square$

\underline{Corollary \SGARefTarget{sga7.sga7i.evi.n3-9}{3.9}.} Let $S$ be a multiplicative subset of
an $R$-algebra $A$. Suppose that
\[
 F\in\operatorname{Ob}
 (\widetilde{\underline C}_{S^{-1}A|R})
\]
is additive and that $F_A$ is also additive. Then the canonical map
\[
 H^i(S^{-1}A|R,F)\longrightarrow H^i(A|R,F)
\]
is an isomorphism for every $i$. Here
$F_A\in\operatorname{Ob}
(\widetilde{\underline C}_{S^{-1}A|A})$ is defined by
\[
 F_A(X)=\ker\bigl(F(X)\longrightarrow F(A)\bigr)
\]
for every
$X\in\operatorname{Ob}(\underline C_{S^{-1}A|A})$.

\underline{Proof.} Applied to
$R\longrightarrow A\longrightarrow S^{-1}A$, Proposition \SGARefLink{sga7.sga7i.evi.n3-7}{3.7}(c)
gives the exact sequence
\[
 \cdots\longrightarrow H^n(S^{-1}A|A,F_A)
 \longrightarrow H^n(S^{-1}A|R,F)
 \longrightarrow H^n(A|R,F)\longrightarrow\cdots.
\]
The map
\[
 S^{-1}A\otimes_AS^{-1}A\longrightarrow S^{-1}A
\]
is an isomorphism, and $F_A$ is additive by hypothesis. Corollary
\SGARefLink{sga7.sga7i.evi.n3-8}{3.8}(b) therefore gives
$H^n(S^{-1}A|A,F_A)=0$ for every $n$. The middle map in the exact
sequence is consequently an isomorphism for every $n$.
\hfill$\square$

% END INCLUDED SGA7I_English_current_buildable_source_20260823/components/31_expose_VI_corollaries_3_8_and_3_9.tex

% BEGIN INCLUDED SGA7I_English_current_buildable_source_20260823/components/32_expose_VI_definition_and_lemma_3_10.tex
%%%%%%%%%%  scan idx 94  |  folio 83  %%%%%%%%%%

Let $A$ be an $R$-algebra. An $A$-module $M$ defines an additive
sheaf
\[
 \operatorname{Der}_R(-,M)
 \in\operatorname{Ob}(\widetilde{\underline C}_{A|R}),
 \qquad X\longmapsto\operatorname{Der}_R(X,M).
\]
We make the following definition.

\underline{Definition \SGARefTarget{sga7.sga7i.evi.n3-10}{3.10}.} Let $A$ be an $R$-algebra. For every
$A$-module $M$, set
\[
 D^i(A|R,M)=H^i(A|R,\operatorname{Der}_R(-,M)).
\]

We now compute the groups of low degree, namely $D^i(A|R,M)$ for
$i=1,2$. Choose a surjective arrow
$P\longrightarrow A$ in $\underline C_{A|R}$ with $P$ a
polynomial algebra over $R$, and write
\[
 0\longrightarrow I\longrightarrow P\longrightarrow A
 \longrightarrow0.
\]

\underline{Lemma \SGARefTarget{sga7.sga7i.evi.n3-10.occ-2}{3.10}.}
\begin{align*}
 D^1(A|R,M)
 &=\operatorname{Coker}\bigl(
   \operatorname{Der}_R(P,M)\longrightarrow
   \operatorname{Hom}_P(I,M)\bigr),\\
 D^2(A|R,M)
 &=H^1\bigl(A|R,
   \underline H^1_{A|R}(\operatorname{Der}_R(-,M))\bigr).
\end{align*}

\underline{Proof.} Consider the simplicial algebra
\[
 (*)\qquad
 \cdots\ P_A^{(2)}
 \mathrel{\substack{\longrightarrow\\[-.6ex]\longrightarrow\\[-.6ex]\longrightarrow}}
 P_A^{(1)}
 \mathrel{\substack{\longrightarrow\\[-.6ex]\longrightarrow}}
 P_A^{(0)},
\]
where
\[
 P_A^{(n)}=
 \overbrace{P\times_A\cdots\times_AP}^{n+1}.
\]
For every $n$ there is an exact sequence
\[
 0\longrightarrow I^{(n)}\longrightarrow P_A^{(n)}
 \longrightarrow A\longrightarrow0,
 \qquad
 I^{(n)}=\overbrace{I\times\cdots\times I}^{n+1}.
\]
Put $\overline I=I/I^2$, and let $\Delta$ denote the diagonal
map. It is straightforward to verify that, for every $n$, the
sequence
\[
 0\longrightarrow
 \operatorname{Hom}_A
 \bigl(\overline I^{(n)}/\Delta(\overline I),M\bigr)
 \longrightarrow\operatorname{Der}_R(P_A^{(n)},M)
 \xrightarrow{\operatorname{Der}(\Delta,M)}
 \operatorname{Der}_R(P,M)\longrightarrow0
 \tag{a}
\]
is exact. We also have the exact sequence
\[
 0\longrightarrow
 \operatorname{Hom}_A
 \bigl(\overline I^{(n)}/\Delta(\overline I),M\bigr)
 \longrightarrow\operatorname{Hom}_A(\overline I^{(n)},M)
 \xrightarrow{\operatorname{Hom}(\Delta,M)}
 \operatorname{Hom}_A(\overline I,M)\longrightarrow0.
 \tag{b}
\]
The sequences (a) and (b) yield the corresponding exact sequences
of simplicial $A$-modules. Since $\{I\longrightarrow0\}$ and
$\{I\longrightarrow I\}$ admit sections,
\[
 H^i(\{I\longrightarrow0\},\operatorname{Hom}_A(-,M))
 =H^i(\{I\longrightarrow I\},\operatorname{Hom}_A(-,M))=0
 \qquad(i>0).
\]

%%%%%%%%%%  scan idx 95  |  folio 84  %%%%%%%%%%

It follows that
\[
\left\{
\begin{aligned}
0\longrightarrow\operatorname{Der}_R(A,M)
&\longrightarrow\operatorname{Der}_R(P,M)
\longrightarrow\operatorname{Hom}_A(\overline I,M)\\
&\longrightarrow
H^1(\{P\longrightarrow A\},\operatorname{Der}_R(-,M))
\longrightarrow0
&&\text{is exact},\\
H^i(\{P\longrightarrow A\},\operatorname{Der}_R(-,M))
&=0 &&(i>1).
\end{aligned}
\right.
\]
By \SGARefLink{sga7.sga7i.evi.n3-1.item-b}{3.1}(b), therefore,
\begin{align*}
D^1(A|R,M)
 &=\check H^1(A|R,\operatorname{Der}_R(-,M))\\
 &=H^1(\{P\longrightarrow A\},\operatorname{Der}_R(-,M))\\
 &=\operatorname{Coker}\bigl(
   \operatorname{Der}_R(P,M)\longrightarrow
   \operatorname{Hom}_A(\overline I,M)\bigr)\\
 &=\operatorname{Coker}\bigl(
   \operatorname{Der}_R(P,M)\longrightarrow
   \operatorname{Hom}_P(I,M)\bigr),\\
D^2(A|R,M)
 &=H^1\bigl(A|R,
   \underline H^1_{A|R}(\operatorname{Der}_R(-,M))\bigr).
\end{align*}
\hfill$\square$

Choose a free $P$-module $F$ and a $P$-linear map
$\gamma:F\longrightarrow P$ such that $\gamma(F)=I$. Associated
with $\gamma$ is a Koszul complex, denoted by $K_I$ by abuse of
notation.

% END INCLUDED SGA7I_English_current_buildable_source_20260823/components/32_expose_VI_definition_and_lemma_3_10.tex

% BEGIN INCLUDED SGA7I_English_current_buildable_source_20260823/components/33_expose_VI_lemma_3_11.tex
\underline{Lemma \SGARefTarget{sga7.sga7i.evi.n3-11}{3.11}.} There is an isomorphism
\[
 D^2(A|R,M)\simeq
 \operatorname{Coker}\bigl(
 H^1(K_I,M)\longrightarrow
 \operatorname{Hom}_A(H_1(K_I),M)\bigr).
\]

\underline{Proof.} We have
\begin{align*}
D^2(A|R,M)
 &=H^1\bigl(A|R,
   \underline H^1_{A|R}(\operatorname{Der}_R(-,M))\bigr)\\
 &=\ker\left(
   H^1(P_A^{(1)}|R,M)
   \mathrel{\substack{\longrightarrow\\[-.6ex]\longrightarrow\\[-.6ex]\longrightarrow}}
   H^1(P_A^{(2)}|R,M)\right),
\end{align*}
because
\[
 [\underline H^1_{A|R}(\operatorname{Der}_R(-,M))](P_A^{(0)})
 =[\underline H^1_{A|R}(\operatorname{Der}_R(-,M))](P)=0.
\]

Let
\[
 0\longrightarrow K\longrightarrow F\xrightarrow{\gamma}I
 \longrightarrow0
\]
be exact. Write $P[F]$ for the symmetric algebra over $P$
generated by $F$. Then $P[F]$ and
$P[F]\otimes_PP[F]$ are polynomial algebras over $R$, and there
are surjective maps
\begin{align*}
P[F]&\longrightarrow P_A^{(1)}=P\times_AP,
&X&\longmapsto(0,\gamma(X)),\\
P[F]\otimes_PP[F]&\longrightarrow
P_A^{(2)}=P\times_AP\times_AP,
&X\otimes1&\longmapsto(0,\gamma(X),0),\\[-.3ex]
&&1\otimes X&\longmapsto(0,0,\gamma(X)),
\end{align*}
for $X\in F$. Put
\[
 J=\ker(P[F]\longrightarrow P_A^{(1)}),
 \qquad
 J'=\ker(P[F]\otimes_PP[F]\longrightarrow P_A^{(2)}).
\]
If $F_\gamma$ denotes the ideal generated by
$\{X-\gamma(X)\mid X\in F\}$, then
\begin{align*}
J&=(F)\cap(F_\gamma)=(K)+(F\cdot F_\gamma),\\
J'&=(F\otimes1+1\otimes F)
 \cap(F_\gamma\otimes1+1\otimes F)
 \cap(F\otimes1+1\otimes F_\gamma)\\
&=(K\otimes1+1\otimes K)
 +(F\cdot F_\gamma\otimes1+1\otimes F\cdot F_\gamma)
 +(F\otimes F).
\end{align*}

%%%%%%%%%%  scan idx 96  |  folio 85  %%%%%%%%%%

The three maps
$\pi_i:P_A^{(2)}\longrightarrow P_A^{(1)}$ admit lifts
$\widehat\pi_i:P[F]\otimes_PP[F]\longrightarrow P[F]$ given, for
$X\in F$, by
\[
\begin{array}{c@{\qquad}c@{\qquad}c}
\widehat\pi_0:
\begin{cases}
 X\otimes1\longmapsto\gamma(X)-X,\\
 1\otimes X\longmapsto X,
\end{cases}
&
\widehat\pi_1:
\begin{cases}
 X\otimes1\longmapsto0,\\
 1\otimes X\longmapsto X,
\end{cases}
&
\widehat\pi_2:
\begin{cases}
 X\otimes1\longmapsto X,\\
 1\otimes X\longmapsto0.
\end{cases}
\end{array}
\]
A direct computation shows that
$\widehat\pi_0-\widehat\pi_1+\widehat\pi_2$ annihilates
\[
 K\otimes1+1\otimes K+1\otimes F\cdot F_\gamma
\]
and maps
$F\cdot F_\gamma\otimes1+F\otimes F$ onto
$F\cdot F_\gamma$. For every
\[
 D\in\operatorname{Der}_R(P[F]\otimes_PP[F],M)
\]
one has
$D(F\cdot F_\gamma\otimes1+F\otimes F)=0$, because $I\cdot M=0$.
Consequently,
\begin{align*}
D^2(A|R,M)
 &=\ker\left(
   D^1(P_A^{(1)}|R,M)
   \mathrel{\substack{\longrightarrow\\[-.6ex]\longrightarrow\\[-.6ex]\longrightarrow}}
   D^1(P_A^{(2)}|R,M)\right)\\
 &=\{[f]\in D^1(P_A^{(1)}|R,M)
       \mid f(F\cdot F_\gamma)=0\}\\
 &=\operatorname{Coker}\left(
   \operatorname{Der}_R(P[F],M)\longrightarrow
   \operatorname{Hom}_{P[F]}
   \left(\frac{(F)\cap(F_\gamma)}{F\cdot F_\gamma},M\right)
   \right)\\
 &=\operatorname{Coker}\left(
   \operatorname{Der}_P(P[F],M)\longrightarrow
   \operatorname{Hom}_{P[F]}
   \left(\frac{(F)\cap(F_\gamma)}{(F)(F_\gamma)},M\right)
   \right).
\end{align*}

The sequences
\[
 0\longrightarrow(F)\longrightarrow P[F]
 \xrightarrow{P[0]}P\longrightarrow0,
 \qquad
 0\longrightarrow(F_\gamma)\longrightarrow P[F]
 \xrightarrow{P[\gamma]}P_\gamma\longrightarrow0
\]
are exact. Hence
\[
 \frac{(F)\cap(F_\gamma)}{(F)(F_\gamma)}
 =\operatorname{Tor}_1^{P[F]}(P,P_\gamma)
 =H_1(K_I),
\]
where $K_I$ is the Koszul complex associated with
$F\xrightarrow{\gamma}I\subset P$. We therefore obtain
\begin{align*}
D^2(A|R,M)
 &=\operatorname{Coker}\bigl(
   \operatorname{Hom}_P(F,M)\longrightarrow
   \operatorname{Hom}_P(H_1(K_I),M)\bigr)\\
 &=\operatorname{Coker}\bigl(
   H^1(K_I,M)\longrightarrow
   \operatorname{Hom}_A(H_1(K_I),M)\bigr).
\end{align*}
\hfill$\square$

We can now state the basic facts needed below, summarizing the
results already established.

% END INCLUDED SGA7I_English_current_buildable_source_20260823/components/33_expose_VI_lemma_3_11.tex

% BEGIN INCLUDED SGA7I_English_current_buildable_source_20260823/components/34_expose_VI_theorem_3_12_statement.tex
%%%%%%%%%%  scan idx 97  |  folio 86  %%%%%%%%%%

\underline{Theorem \SGARefTarget{sga7.sga7i.evi.n3-12}{3.12}.}

\begin{enumerate}[label=(\arabic*)]
\item If
\[
 0\longrightarrow M'\longrightarrow M\longrightarrow M''
 \longrightarrow0
\]
is an exact sequence of $A$-modules, then the sequence
\begin{multline*}
0\longrightarrow D^0(A|R,M')\longrightarrow D^0(A|R,M)
\longrightarrow D^0(A|R,M'')\longrightarrow D^1(A|R,M')\\
\longrightarrow D^1(A|R,M)\longrightarrow D^1(A|R,M'')
\longrightarrow\cdots\longrightarrow D^n(A|R,M')\\
\longrightarrow D^n(A|R,M)\longrightarrow D^n(A|R,M'')
\longrightarrow\cdots
\end{multline*}
is exact.

\item Let $S\longrightarrow R\longrightarrow A$ be ring
homomorphisms. Then, for every $A$-module $M$, the sequence
\begin{multline*}
0\longrightarrow D^0(A|R,M)\longrightarrow D^0(A|S,M)
\longrightarrow D^0(R|S,M)\longrightarrow D^1(A|R,M)\\
\longrightarrow D^1(A|S,M)\longrightarrow D^1(R|S,M)
\longrightarrow\cdots\longrightarrow D^n(A|R,M)\\
\longrightarrow D^n(A|S,M)\longrightarrow D^n(R|S,M)
\longrightarrow\cdots
\end{multline*}
is exact.

\item Let $A$ and $S$ be $R$-algebras such that
$\operatorname{Tor}_i^R(A,S)=0$ for every $i>0$. Then, for every
$A\otimes_RS$-module $M$, the canonical map
\[
 D^i(A\otimes_RS|S,M)\longrightarrow D^i(A|R,M)
\]
is an isomorphism for every $i$. In particular, if $A$ is flat
over $R$, then
\[
 D^i(A\otimes_RS|S,M)\xrightarrow{\ \sim\ }D^i(A|R,M)
\]
for every $i$, every $R$-algebra $S$, and every
$A\otimes_RS$-module $M$.

\item Let $S$ be a multiplicative subset of $A$. For every
$S^{-1}A$-module $N$, the canonical map
\[
 D^i(S^{-1}A|R,N)\longrightarrow D^i(A|R,N)
\]
is an isomorphism for every $i$. Consequently, for every
$A$-module $M$, there is a canonical isomorphism
\[
 D^i(S^{-1}A|R,S^{-1}M)
 \xrightarrow{\ \sim\ }S^{-1}D^i(A|R,M)
\]
for every $i$.

\item Let $A$ and $B$ be $R$-algebras such that
$\operatorname{Tor}_i^R(A,B)=0$ for $0<i\leq n$. Then, for every
$A\otimes_RB$-module $M$, the canonical map
\[
 D^i(A\otimes_RB|R,M)
 \longrightarrow D^i(A|R,M)\oplus D^i(B|R,M)
\]
is an isomorphism for $i\leq n$. In particular, suppose that
$\operatorname{Tor}_i^R(A,A)=0$ for every $i>0$ and that the
diagonal
\[
 \Delta:\operatorname{Spec}(A)\longrightarrow
 \operatorname{Spec}(A)\times_{\operatorname{Spec}(R)}
 \operatorname{Spec}(A)
\]
is an open immersion---for example, suppose that $A$ is \emph{étale}
over $R$. Then
\[
 D^i(A|R,M)=0
\]
for every $i>0$ and every $A$-module $M$.

%%%%%%%%%%  scan idx 98  |  folio 87  %%%%%%%%%%

\item Let $R$ be noetherian, let $A$ be an $R$-algebra essentially
of finite type, and let $E$ be a flat $A$-module. Then the canonical
map
\[
 E\otimes_AD^i(A|R,M)
 \xrightarrow{\ \sim\ }
 D^i(A|R,E\otimes_AM)
\]
is an isomorphism for every $i\geq0$ and every $A$-module $M$.

\item If $R$ is noetherian, $A$ is an $R$-algebra essentially of
finite type, and $M$ is an $A$-module of finite type, then
$D^i(A|R,M)$ is an $A$-module of finite type for every $i$.

\item Let $P$ be a polynomial algebra over $R$, and let
\[
 0\longrightarrow I\longrightarrow P\longrightarrow A
 \longrightarrow0
\]
be exact. Then
\[
 D^0(A|R,M)=\operatorname{Der}_R(A,M),
\]
\begin{align*}
D^1(A|R,M)
 &=\operatorname{Coker}\bigl(
   \operatorname{Der}_R(P,M)\longrightarrow
   \operatorname{Hom}_A(I/I^2,M)\bigr)\\
 &=\left\{
   \begin{array}{c}
   \text{isomorphism classes of commutative}\
   R\text{-algebra extensions of }A\text{ by }M
   \end{array}\right\},\\
D^2(A|R,M)
 &=\operatorname{Coker}\bigl(
   H^1(K_I,M)\longrightarrow
   \operatorname{Hom}_A(H_1(K_I),M)\bigr).
\end{align*}
\end{enumerate}

% END INCLUDED SGA7I_English_current_buildable_source_20260823/components/34_expose_VI_theorem_3_12_statement.tex

% BEGIN INCLUDED SGA7I_English_current_buildable_source_20260823/components/35_expose_VI_theorem_3_12_proof.tex
\underline{Proof.}

\begin{enumerate}[label=(\arabic*)]
\item If
\[
 0\longrightarrow M'\longrightarrow M\longrightarrow M''
 \longrightarrow0
\]
is an exact sequence of $A$-modules, then, for every polynomial
algebra $P$ over $R$, the sequence
\[
 0\longrightarrow\operatorname{Der}_R(P,M')
 \longrightarrow\operatorname{Der}_R(P,M)
 \longrightarrow\operatorname{Der}_R(P,M'')
 \longrightarrow0
\]
is exact. Hence, by \SGARefLink{sga7.sga7i.evi.n3-3}{3.3}(c),
\[
 0\longrightarrow\operatorname{Der}_R(-,M')
 \longrightarrow\operatorname{Der}_R(-,M)
 \longrightarrow\operatorname{Der}_R(-,M'')
 \longrightarrow0
\]
is an exact sequence of sheaves.

\item The functor $\operatorname{Der}_S(-,M)$ is an additive sheaf
on $\underline C_{A|S}$, and
\[
 \operatorname{Der}_R(B,M)
 =\ker\bigl(\operatorname{Der}_S(B,M)
 \longrightarrow\operatorname{Der}_S(R,M)\bigr)
\]
for every $B\in\operatorname{Ob}(\underline C_{A|R})$.
The assertion therefore follows from \SGARefLink{sga7.sga7i.evi.n3-7}{3.7}(c).

\item Let
\[
 S:\underline C_{A|R}\longrightarrow
 \underline C_{A\otimes_RS|S},
 \qquad X\longmapsto X\otimes_RS.
\]
Then
$\widetilde S\operatorname{Der}_S(-,M)
=\operatorname{Der}_R(-,M)$, so the assertion follows from \SGARefLink{sga7.sga7i.evi.n3-7}{3.7}(a).

%%%%%%%%%%  scan idx 99  |  folio 88  %%%%%%%%%%

\item This follows from \SGARefLink{sga7.sga7i.evi.n3-9}{3.9}. Indeed,
$F=\operatorname{Der}_R(-,N)$ is an additive sheaf on
$\underline C_{S^{-1}A|R}$, and
$F_A\in\operatorname{Ob}
(\widetilde{\underline C}_{S^{-1}A|A})$ is also additive, since
\begin{align*}
 F_A(X)
 &=\ker\bigl(F(X)\longrightarrow F(A)\bigr)\\
 &=\ker\bigl(\operatorname{Der}_R(X,N)
   \longrightarrow\operatorname{Der}_R(A,N)\bigr)\\
 &=\operatorname{Der}_A(X,N).
\end{align*}
Thus $F_A=\operatorname{Der}_A(-,N)$.

\item The first assertion follows from \SGARefLink{sga7.sga7i.evi.n3-8}{3.8}(a). Now suppose that
$\operatorname{Tor}_i^R(A,A)=0$ for every $i>0$ and that
\[
 \Delta:\operatorname{Spec}(A)\longrightarrow
 \operatorname{Spec}(A)\times_{\operatorname{Spec}(R)}
 \operatorname{Spec}(A)
\]
is an open immersion. Let
$\varphi:A\otimes_RA\longrightarrow A$ be the multiplication map.
For every maximal ideal $\mathfrak m$ of $A$, the induced map
\[
 (A\otimes_RA)_{\varphi^{-1}(\mathfrak m)}
 \longrightarrow A_{\mathfrak m}
\]
is an isomorphism. The first assertion, together with \SGARefLink{sga7.sga7i.evi.n3-12}{3.12}(4),
therefore implies that the diagonal map
\[
 \Delta_{\mathfrak m}:D^i(A|R,M)_{\mathfrak m}
 \longrightarrow D^i(A|R,M)_{\mathfrak m}
 \oplus D^i(A|R,M)_{\mathfrak m}
\]
is an isomorphism for every maximal ideal $\mathfrak m$ of $A$.
Consequently,
\[
 \Delta:D^i(A|R,M)\longrightarrow
 D^i(A|R,M)\oplus D^i(A|R,M)
\]
is an isomorphism, and hence $D^i(A|R,M)=0$ for every $i>0$.

\item In view of \SGARefLink{sga7.sga7i.evi.n3-12}{3.12}(4), we may assume that $A$ is an
$R$-algebra of finite type. Consider the morphism of sheaves
\[
 E\otimes_A\operatorname{Der}_R(-,M)
 \longrightarrow
 \operatorname{Der}_R(-,E\otimes_AM)
\]
on $\widetilde{\underline C}_{A|R}$. Since $E$ is flat over $A$,
the map
\[
 E\otimes_A\operatorname{Der}_R(B,M)
 \xrightarrow{\ \sim\ }
 \operatorname{Der}_R(B,E\otimes_AM)
\]
is an isomorphism whenever
$B\in\operatorname{Ob}(\underline C_{A|R})$ is of finite type
over $R$. By \SGARefLink{sga7.sga7i.evi.n3-3}{3.3}(b), it follows that
\[
 D^i\bigl(A|R,E\otimes_A\operatorname{Der}_R(-,M)\bigr)
 \longrightarrow
 D^i\bigl(A|R,\operatorname{Der}_R(-,E\otimes_AM)\bigr)
\]
is an isomorphism. On the other hand,
$E\otimes_A-$ is exact, and therefore
\[
 E\otimes_AD^i(A|R,M)
 =D^i\bigl(A|R,E\otimes_A\operatorname{Der}_R(-,M)\bigr).
\]
This proves the assertion.

\item This follows from \SGARefLink{sga7.sga7i.evi.n3-3}{3.3}(a).

%%%%%%%%%%  scan idx 100  |  folio 89  %%%%%%%%%%

\item In view of \SGARefLink{sga7.sga7i.evi.n3-10.occ-2}{3.10} and \SGARefLink{sga7.sga7i.evi.n3-11}{3.11}, it suffices to show that
\[
 \operatorname{Coker}\bigl(
 \operatorname{Der}_R(P,M)\longrightarrow
 \operatorname{Hom}_A(I/I^2,M)\bigr)
 =\operatorname{Exalcomm}(A|R,M),
\]
where $\operatorname{Exalcomm}(A|R,M)$ denotes the set of
isomorphism classes of $R$-algebra extensions of $A$ by $M$.
Let
\[
 0\longrightarrow M\longrightarrow A'\longrightarrow A
 \longrightarrow0
\]
be such an extension. Since $P$ is a polynomial algebra over $R$,
we obtain a commutative diagram
\[
\begin{tikzcd}[column sep=large]
0 \arrow[r] & I \arrow[r] \arrow[d]
& P \arrow[r] \arrow[d]
& A \arrow[r] \arrow[d,equal] & 0 \\
0 \arrow[r] & M \arrow[r]
& A' \arrow[r]
& A \arrow[r] & 0.
\end{tikzcd}
\]
The map $I\longrightarrow M$ is determined by the extension $A'$,
up to the action of $\operatorname{Der}_R(P,M)$. We thus obtain a
map
\[
 \Phi:\operatorname{Exalcomm}(A|R,M)
 \longrightarrow
 \operatorname{Coker}\bigl(
 \operatorname{Der}_R(P,M)\longrightarrow
 \operatorname{Hom}_A(I/I^2,M)\bigr).
\]
Conversely, a $P$-linear map $I\longrightarrow M$ yields, by
pushout, an extension fitting into a diagram
\[
\begin{tikzcd}[column sep=large]
0 \arrow[r] & I \arrow[r] \arrow[d]
& P \arrow[r] \arrow[d]
& A \arrow[r] \arrow[d,equal] & 0 \\
0 \arrow[r] & M \arrow[r]
& A' \arrow[r]
& A \arrow[r] & 0.
\end{tikzcd}
\]
Two $P$-linear maps $I\rightrightarrows M$ that differ by an
element of $\operatorname{Der}_R(P,M)$ yield isomorphic
extensions. Hence there is a map
\[
 \Psi:\operatorname{Coker}\bigl(
 \operatorname{Der}_R(P,M)\longrightarrow
 \operatorname{Hom}_A(I/I^2,M)\bigr)
 \longrightarrow\operatorname{Exalcomm}(A|R,M).
\]
It is straightforward to verify that $\Phi$ and $\Psi$ are inverse
to each other.
\end{enumerate}
\hfill$\square$

% END INCLUDED SGA7I_English_current_buildable_source_20260823/components/35_expose_VI_theorem_3_12_proof.tex

% BEGIN INCLUDED SGA7I_English_current_buildable_source_20260823/components/36_expose_VI_corollary_3_13.tex
%%%%%%%%%%  scan idx 100  |  folio 89  %%%%%%%%%%

\underline{Corollary \SGARefTarget{sga7.sga7i.evi.n3-13}{3.13}.}

\begin{enumerate}[label=(\alph*)]
\item The $R$-algebra $A$ is formally smooth over $R$ if and only if
$D^1(A|R,M)=0$ for every $A$-module $M$.

\item Let $R$ be noetherian and let $A$ be an $R$-algebra of finite
type. Then $A$ is a relative complete intersection over $R$ if and
only if $D^2(A|R,M)=0$ for every $A$-module $M$.

\item Let $R$ be noetherian and let $A$ be an $R$-algebra of finite
type. For every multiplicative subset $S$ of $A$, there is a canonical
isomorphism
\[
 D^i(S^{-1}A|R,S^{-1}M)
 \xrightarrow{\ \sim\ }S^{-1}D^i(A|R,M)
 \qquad(i\geq0).
\]

\item Let $R$ be noetherian and let $A$ be a local $R$-algebra
essentially of finite type. Write $\widehat A$ for the
$\mathfrak m$-adic completion of $A$, where $\mathfrak m$ is the
maximal ideal of $A$. Then, for every $A$-module $E$ of finite type,
there is a canonical isomorphism
\[
 D^1(\widehat A|R,\widehat A\otimes_AE)
 \xrightarrow{\ \sim\ }
 \widehat A\otimes_AD^1(A|R,E).
\]
\end{enumerate}

\underline{Proof.}

Part (a) follows immediately from
$D^1(A|R,M)=\operatorname{Exalcomm}(A|R,M)$.

For (b), let $P$ be a polynomial algebra over $R$ in finitely many
variables, and let
\[
 0\longrightarrow I\longrightarrow P\longrightarrow A
 \longrightarrow0
\]
be exact. After localization, we may suppose that $I$ is generated by
a $P$-regular sequence. Then $H_1(K_I)=0$, and hence
\[
 D^2(A|R,M)
 =\operatorname{Coker}\bigl(
 H^1(K_I,M)\longrightarrow
 \operatorname{Hom}_A(H_1(K_I),M)\bigr)=0.
\]
Conversely, suppose that $D^2(A|R,M)=0$ for every $A$-module $M$.
Choose an exact sequence
\[
 0\longrightarrow K\longrightarrow F\longrightarrow I
 \longrightarrow0,
\]
where $F$ is a free $P$-module of finite type. If
\[
 K_0=\operatorname{Im}\bigl(\textstyle\bigwedge^2F
 \longrightarrow\bigwedge^1F\bigr),
\]
then $H_1(K_I)=K/K_0$. Since
$D^2(A|R,H_1(K_I))=0$, the map
\[
 H^1(K_I,H_1(K_I))
 \longrightarrow
 \operatorname{Hom}_A(H_1(K_I),H_1(K_I))
\]
is surjective. Consequently, the exact sequence
\[
 0\longrightarrow K/K_0\longrightarrow F/IF
 \longrightarrow I/I^2\longrightarrow0
\]
splits as a sequence of $A$-modules. In particular, $I/I^2$ is
projective. After localizing $A$ once more, we may suppose that
$I/I^2$ is free and that $F/IF\longrightarrow I/I^2$ is an
isomorphism. The split exact sequence then gives
$H_1(K_I)=K/K_0=0$. Thus $I$ is generated by a $P$-regular sequence,
so $A$ is a relative complete intersection over $R$.

Part (c) follows immediately from \SGARefLink{sga7.sga7i.evi.n3-12}{3.12}(4) and (6).

%%%%%%%%%%  scan idx 101  |  folio 90  %%%%%%%%%%

For (d), \SGARefLink{sga7.sga7i.evi.n3-12}{3.12}(6) first gives canonical isomorphisms
\[
 D^i(A|R,\widehat A\otimes_AE)
 \xrightarrow{\ \sim\ }
 \widehat A\otimes_AD^i(A|R,E)
\]
for every $i$. It therefore remains to prove that
\[
 D^1(\widehat A|R,\widehat A\otimes_AE)
 \longrightarrow D^1(A|R,\widehat A\otimes_AE)
\]
is an isomorphism. We use the description of $D^1$ in terms of
algebra extensions.

Let
\[
 0\longrightarrow\widehat A\otimes_AE\longrightarrow A'
 \xrightarrow{\ \varphi\ }\widehat A\longrightarrow0
\]
be an $R$-algebra extension whose restriction to $A$ splits. Since
$\widehat A\otimes_AE$ is an $\widehat A$-module of finite type,
$A'$ is a local $R$-algebra complete for its $\mathfrak m'$-adic
topology, where $\mathfrak m'$ is the maximal ideal of $A'$. If
$\psi:A\longrightarrow A'$ is an $R$-algebra map lifting the
canonical map $A\longrightarrow\widehat A$, then $\psi$ extends
uniquely to a map
$\widehat\psi:\widehat A\longrightarrow A'$ satisfying
$\varphi\circ\widehat\psi=\operatorname{id}_{\widehat A}$.
This proves that
\[
 D^1(\widehat A|R,\widehat A\otimes_AE)
 \longrightarrow D^1(A|R,\widehat A\otimes_AE)
\]
is injective.

%%%%%%%%%%  scan idx 102  |  folio 91  %%%%%%%%%%

Now let
\[
 0\longrightarrow\widehat A\otimes_AE\longrightarrow B
 \xrightarrow{\ \varphi\ }A\longrightarrow0
\]
be an $R$-algebra extension, and put
$\widehat E=\widehat A\otimes_AE$. We first claim that the topology
on $\widehat E$ induced by the $\mathfrak m_B$-adic topology on $B$
coincides with its topology as an $\widehat A$-module. Indeed, since
$A$ is essentially of finite type over $R$, there is a subalgebra
$B_0\subset B$ such that $B_0$ is a local $R$-algebra essentially
of finite type and $\varphi(B_0)=A$. Put $E_0=\widehat E\cap B_0$,
and denote the maximal ideals of $B$, $B_0$, and $A$ by
$\mathfrak m_B$, $\mathfrak m_0$, and $\mathfrak m$, respectively.
Then
\[
 \mathfrak m_B=\mathfrak m_0+\widehat E,
 \qquad
 \mathfrak m_B^{n+1}=\mathfrak m_0^{n+1}
   +\mathfrak m^n\widehat E
\]
for every $n$. Since $B_0$ is noetherian, the Artin--Rees theorem
gives an integer $r>0$ such that, for every $n\geq r$,
\begin{align*}
 \widehat E\cap\mathfrak m_B^{n+1}
 &=\mathfrak m^n\widehat E
   +\widehat E\cap\mathfrak m_0^{n+1}\\
 &=\mathfrak m^n\widehat E
   +E_0\cap\mathfrak m_0^{n+1}\\
 &=\mathfrak m^n\widehat E
   +\mathfrak m^{n+1-r}(E_0\cap\mathfrak m_0^r)
 \subset\mathfrak m^{n+1-r}\widehat E.
\end{align*}
This proves the assertion about the two topologies. In particular,
$\widehat E$ is complete for the topology induced by the
$\mathfrak m_B$-adic topology on $B$, and hence completion gives an
exact commutative diagram
\[
\begin{tikzcd}[column sep=large]
0 \arrow[r] & \widehat E \arrow[r] \arrow[d,equal]
& B \arrow[r] \arrow[d]
& A \arrow[r] \arrow[d] & 0 \\
0 \arrow[r] & \widehat E \arrow[r]
& \widehat B \arrow[r]
& \widehat A \arrow[r] & 0.
\end{tikzcd}
\]
Here $\widehat B$ denotes the $\mathfrak m_B$-adic completion of
$B$. Thus
\[
 D^1(\widehat A|R,\widehat E)
 \longrightarrow D^1(A|R,\widehat E)
\]
is surjective, and the proof is complete.
\hfill$\square$

% END INCLUDED SGA7I_English_current_buildable_source_20260823/components/36_expose_VI_corollary_3_13.tex

% BEGIN INCLUDED SGA7I_English_current_buildable_source_20260823/components/37_expose_VI_3_14_through_3_16.tex
%%%%%%%%%%  scan idx 103  |  folio 92  %%%%%%%%%%

\underline{Corollary \SGARefTarget{sga7.sga7i.evi.n3-14}{3.14}.}
Let $R$ be noetherian, and let $A$ be an $R$-algebra of finite type
that is smooth everywhere except at one closed point
$x\in\operatorname{Spec}(A)$. Then there are isomorphisms
\[
 D^1(A|R,A)\simeq D^1(A_x|R,A_x)
 \simeq D^1(\widehat A_x|R,\widehat A_x).
\]

\underline{Proof.}
Since $A$ is smooth over $R$ away from $x$,
$D^1(A|R,A)$ is an $A$-module of finite type supported on
$\{x\}$. Hence
\[
 D^1(A|R,A)=D^1(A|R,A)_x
 =\widehat{D^1(A|R,A)_x}.
\]
The assertion now follows from \SGARefLink{sga7.sga7i.evi.n3-13}{3.13}.
\hfill$\square$

\underline{Remark \SGARefTarget{sga7.sga7i.evi.n3-15}{3.15}.}
Let $R$ be a topological ring and $A$ a topological $R$-algebra.
For every $A$-module $E$, define
\[
 D^i_{\mathrm{top}}(A|R,E)
 =\varinjlim_\alpha
 D^i(A_\alpha|R_\alpha,A_\alpha\otimes_AE),
\]
where $A_\alpha=A/J_\alpha$, $R_\alpha=R/I_\alpha$, and
$J_\alpha$ and $I_\alpha$ are open in $A$ and $R$, respectively.
Thus $A$ is formally smooth as a topological $R$-algebra precisely
when $D^1_{\mathrm{top}}(A|R,E)=0$ for every $A$-module $E$.

Now let $X$ be a scheme over a ring $R$, and let $\underline E$ be a
quasi-coherent $\mathcal O_X$-module. We define sheaves
\[
 \underline D^i_{X|R}(\underline E)
 =\underline D^i(X|R,\underline E)
\]
on $X$ from the presheaves on affine open subsets $U\subset X$
given by
\[
 \bigl[\underline D^i_{X|R}(\underline E)\bigr](U)
 =D^i\bigl(\Gamma(U,\mathcal O_X)|R,
            \Gamma(U,\underline E)\bigr).
\]
If $R$ is noetherian and $X$ is of finite type over $R$, the
localization property \SGARefLink{sga7.sga7i.evi.n3-13}{3.13}(c) shows that these presheaves canonically
form quasi-coherent $\mathcal O_X$-modules.

When no confusion is possible, we abbreviate
\[
 \underline D^i_{X|R}
 =\underline D^i_{X|R}(\mathcal O_X).
\]
All the assertions of \SGARefLink{sga7.sga7i.evi.n3-13}{3.13} have corresponding formulations for the
quasi-coherent sheaves $\underline D^i_{X|R}(\underline E)$. We leave
those formulations to the reader and record below the properties
most relevant to our purpose.

\underline{Corollary \SGARefTarget{sga7.sga7i.evi.n3-15.occ-2}{3.15}.}
Let $R$ be noetherian and let $X$ be an $R$-scheme of finite type.

%%%%%%%%%%  scan idx 104  |  folio 93  %%%%%%%%%%

\begin{enumerate}[label=(\alph*)]
\item For every coherent $\mathcal O_X$-module $\underline E$,
$\underline D^i_{X|R}(\underline E)$ is again coherent for every
$i$. Moreover,
\[
 \operatorname{Supp}\underline D^i_{X|R}(\underline E)
 \subset X^{\mathrm{Sing}}
 \qquad(i>0),
\]
where
\[
 X^{\mathrm{Sing}}
 =\{x\in X\mid X\text{ is not smooth over }R\text{ at }x\}.
\]

\item If $X$ is flat over $R$, then for every $R$-algebra $S$ there
is a canonical isomorphism
\[
 \underline D^i_{X\otimes_RS\,|\,S}(\underline E)
 \xrightarrow{\ \sim\ }
 \underline D^i_{X|R}(\underline E)
\]
for every $i$.

\item Suppose that $X$ is a relative complete intersection over
$R$. For every surjective ring homomorphism $S\longrightarrow R$,
there is an exact sequence
\[
0\longrightarrow\underline D^1_{X|R}(\underline E)
\longrightarrow\underline D^1_{X|S}(\underline E)
\longrightarrow
\underline{\operatorname{Hom}}_{\mathcal O_X}
 \bigl((I/I^2)\otimes_R\mathcal O_X,\underline E\bigr)
\longrightarrow0,
\]
where $I=\ker(S\longrightarrow R)$. If, moreover, $R$ is local
with residue field $k$ and $X$ is flat over $R$, then
\[
0\longrightarrow
 \underline D^1_{X_0|k}(\mathcal O_{X_0})
\longrightarrow
 \underline D^1_{X|S}(\mathcal O_{X_0})
\longrightarrow
 D^1(R|S,k)\otimes_k\mathcal O_{X_0}
\longrightarrow0,
\]
where $X_0=X\otimes_Rk$.
\end{enumerate}

\underline{Proof.}
Part (a) follows from the finiteness and low-degree descriptions in
\SGARefLink{sga7.sga7i.evi.n3-12}{3.12}, while (b) follows from \SGARefLink{sga7.sga7i.evi.n3-12}{3.12}(3).

For (c), because $X$ is a relative complete intersection over $R$,
we have $\underline D^i_{X|R}(\underline E)=0$ for every $i\geq2$.
The transitivity sequence \SGARefLink{sga7.sga7i.evi.n3-12}{3.12}(2) therefore gives an exact sequence
\begin{multline*}
0\longrightarrow\underline D^0_{X|R}(\underline E)
\longrightarrow\underline D^0_{X|S}(\underline E)
\longrightarrow\underline D^0_{R|S}(\underline E)
\longrightarrow\underline D^1_{X|R}(\underline E)\\
\longrightarrow\underline D^1_{X|S}(\underline E)
\longrightarrow\underline D^1_{R|S}(\underline E)
\longrightarrow0,
\end{multline*}
where the quasi-coherent sheaves
$\underline D^i_{R|S}(\underline E)$ on $X$ are defined on every
affine open subset $U\subset X$ by
\[
 \bigl[\underline D^i_{R|S}(\underline E)\bigr](U)
 =D^i(R|S,\underline E(U)).
\]
Since $S\longrightarrow R$ is surjective,
$\underline D^0_{R|S}(\underline E)=0$; and the low-degree
description in \SGARefLink{sga7.sga7i.evi.n3-12}{3.12}(8) gives
\[
 \underline D^1_{R|S}(\underline E)
 =\underline{\operatorname{Hom}}_{\mathcal O_X}
  \bigl((I/I^2)\otimes_R\mathcal O_X,\underline E\bigr).
\]
This proves the first exact sequence.

%%%%%%%%%%  scan idx 105  |  folio 94  %%%%%%%%%%

If $X$ is flat over $R$, then, for every quasi-coherent
$\mathcal O_{X_0}$-module $\underline E$, there are canonical
isomorphisms
\[
 \underline D^i_{X_0|k}(\underline E)
 \xrightarrow{\ \sim\ }
 \underline D^i_{X|R}(\underline E)
\]
for every $i$. In particular,
\[
 \underline D^1_{X_0|k}(\mathcal O_{X_0})
 \xrightarrow{\ \sim\ }
 \underline D^1_{X|R}(\mathcal O_{X_0}),
\]
from which the second exact sequence follows.
\hfill$\square$

\underline{Remark \SGARefTarget{sga7.sga7i.evi.n3-16}{3.16}.}
The quasi-coherent sheaves $\underline D^i_{X|R}(\underline E)$
defined above should occur as the terms $E_2^{0,i}$ of a spectral
sequence relating the local theory to the global theory. L.~Illusie
has obtained such a theory in \emph{Lecture Notes in Mathematics}
239. His method uses ``homotopical algebra,'' generalizing the methods
of M.~André and D.~Quillen, rather than ``cohomological algebra.'' It
seems that the approach taken here can be globalized by means of
Grothendieck cohomology on the category of schemes, equipped with an
appropriate topology generalizing the one on the affine category.

% END INCLUDED SGA7I_English_current_buildable_source_20260823/components/37_expose_VI_3_14_through_3_16.tex

% BEGIN INCLUDED SGA7I_English_current_buildable_source_20260823/components/38_expose_VI_section_4_opening_through_4_1.tex
%%%%%%%%%%  scan idx 105  |  folio 94  %%%%%%%%%%

\section*{\SGARefTarget{sga7.sga7i.evi.n4}{4}. Formal moduli of deformations}

Among the natural cofibered categories whose isomorphism classes of
objects are generally not prorepresentable but nevertheless admit a
formal versal object is the cofibered category of deformations. We
return to the situation and notation of section 1. Thus, throughout
this section, $\Lambda$ is a fixed complete noetherian local ring with
residue field $k$, and $\underline C_\Lambda$ is the category of
artinian local $\Lambda$-algebras with residue field $k$.

Let $X_0$ be a fixed scheme over $k$. By an (infinitesimal)
deformation of $X_0$ we mean a pair $(R,X)$, where
$R\in\operatorname{Ob}(\underline C_\Lambda)$ and $X$ is a flat
$R$-scheme, together with a commutative diagram
\[
\begin{tikzcd}[column sep=huge,row sep=huge]
X_0 \arrow[r,"i_X"] \arrow[d] & X \arrow[d,"\mathrm{flat}"] \\
\operatorname{Spec}(k) \arrow[r] & \operatorname{Spec}(R)
\end{tikzcd}
\]
such that the induced morphism
\[
 X_0\longrightarrow
 X\times_{\operatorname{Spec}(R)}\operatorname{Spec}(k)
\]
is an isomorphism. Since $R$ is artinian local,
$i_X:X_0\longrightarrow X$ is necessarily a closed immersion and a
homeomorphism on the underlying topological spaces.

%%%%%%%%%%  scan idx 106  |  folio 95  %%%%%%%%%%

If $(R,X)$ and $(R',X')$ are deformations of $X_0$ over $R$ and
$R'$, respectively, an arrow
\[
 (R',X')\longrightarrow(R,X)
\]
consists of a map $\varphi:R'\longrightarrow R$ in
$\underline C_\Lambda$ and a morphism of schemes
$\Phi:X\longrightarrow X'$ fitting into a commutative diagram
\[
\begin{tikzcd}[column sep=huge,row sep=large]
X \arrow[rr,"\Phi"] \arrow[ddd] && X' \arrow[ddd] \\
& X_0 \arrow[ul,"i_X"'] \arrow[ur,"i_{X'}"] \arrow[d] & \\
& \operatorname{Spec}(k) \arrow[dl] \arrow[dr] & \\
\operatorname{Spec}(R)
 \arrow[rr,"\operatorname{Spec}(\varphi)"']
&& \operatorname{Spec}(R').
\end{tikzcd}
\]
This defines the category $\underline M_{X_0}$ of infinitesimal
deformations of $X_0$. It is equipped with the functor
\[
 p:\underline M_{X_0}\longrightarrow\underline C_\Lambda,
 \qquad (R,X)\longmapsto R.
\]

Let $(R,X)\in\operatorname{Ob}(\underline M_{X_0})$ and let
$\varphi:R\longrightarrow S$ be a map in $\underline C_\Lambda$.
There is a commutative diagram
\[
\begin{tikzcd}[column sep=huge,row sep=huge]
X\times_{\operatorname{Spec}(R)}\operatorname{Spec}(S)
 \arrow[r,"p_1"] \arrow[d,"\mathrm{flat}"']
& X \arrow[d,"\mathrm{flat}"] \\
\operatorname{Spec}(S)
 \arrow[r,"\operatorname{Spec}(\varphi)"']
& \operatorname{Spec}(R),
\end{tikzcd}
\]
and the $\varphi$-arrow
\[
 (\varphi,p_1):(R,X)\longrightarrow
 \bigl(S,X\times_{\operatorname{Spec}(R)}\operatorname{Spec}(S)\bigr)
\]
is cocartesian in $\underline M_{X_0}$. Hence
$\underline M_{X_0}$ is a cofibered category over
$\underline C_\Lambda$. We record the following facts.

%%%%%%%%%%  scan idx 107  |  folio 96  %%%%%%%%%%

\medskip
\noindent\textbf{\SGARefTarget{sga7.sga7i.evi.n4-1.item-a}{\SGARefAliasTarget{sga7.sga7i.evi.n4-1}4.1}(a).}
Let
\[
 (1_R,\Phi):(R,X_1)\longrightarrow(R,X_2)
\]
be an arrow in $\underline M_{X_0}$. Thus
$\Phi:X_2\longrightarrow X_1$ is an $R$-morphism whose reduction
\[
 X_2\otimes_Rk\longrightarrow X_1\otimes_Rk
\]
is an isomorphism. Since $R$ is artinian and the $X_i$ are flat over
$R$, $\Phi$ is an isomorphism. Consequently,
$\underline M_{X_0}$ is a cofibered category in groupoids over
$\underline C_\Lambda$.

\medskip
\noindent\textbf{\SGARefTarget{sga7.sga7i.evi.n4-1.item-b}{4.1}(b).}
Suppose that $X_0$ is affine, say
$X_0=\operatorname{Spec}(A_0)$. Then every
$(R,X)\in\operatorname{Ob}(\underline M_{X_0})$ has $X$ affine as
well. Consequently, for every
$R\in\operatorname{Ob}(\underline C_\Lambda)$, taking global
sections identifies $\underline M_{X_0}(R)$ with the following
category. An object is a flat $R$-algebra $A$ together with an
$R$-algebra map $j_A:A\longrightarrow A_0$ inducing an isomorphism
\[
 A\otimes_Rk\xrightarrow{\ \sim\ }A_0.
\]
A morphism $\Phi:A_1\longrightarrow A_2$ is an $R$-algebra map for
which the diagram
\[
\begin{tikzcd}[column sep=huge,row sep=huge]
A_1 \arrow[r,"\Phi"] \arrow[d,"j_{A_1}"']
& A_2 \arrow[d,"j_{A_2}"] \\
A_0 \arrow[r,equal] & A_0
\end{tikzcd}
\]
commutes.

\medskip
\noindent\textbf{\SGARefTarget{sga7.sga7i.evi.n4-1.item-c}{4.1}(c).}
Consider the cofibered category $\widehat{\underline M}_{X_0}$ in
groupoids over $\widehat{\underline C}_\Lambda$. By definition, for
every $R\in\operatorname{Ob}(\widehat{\underline C}_\Lambda)$, the
category $\widehat{\underline M}_{X_0}(R)$ is canonically equivalent
to the category of formal schemes $\mathfrak X$, adic over
$\operatorname{Spf}(R)$, such that
\[
 \mathfrak X_n
 =\mathfrak X\otimes_R R/\mathfrak m^{n+1}
\]
is flat over $R/\mathfrak m^{n+1}$ for every $n\geq0$. If $X_0$ is
noetherian, this is equivalent to requiring $\mathfrak X$ to be flat
over $\operatorname{Spf}(R)$. Thus, in that case,
$\widehat{\underline M}_{X_0}$ is canonically equivalent over
$\widehat{\underline C}_\Lambda$ to the category of deformations of
$X_0$ over the objects of $\widehat{\underline C}_\Lambda$, in the
sense of formal schemes.

\medskip
\noindent\textbf{\SGARefTarget{sga7.sga7i.evi.n4-1.item-d}{4.1}(d).}
Let $(R,X)\in\operatorname{Ob}(\underline M_{X_0})$, so that there is
a commutative diagram
\[
\begin{tikzcd}[column sep=huge,row sep=huge]
X_0 \arrow[r,"i_X"] \arrow[d] & X \arrow[d,"\mathrm{flat}"] \\
\operatorname{Spec}(k) \arrow[r] & \operatorname{Spec}(R)
\end{tikzcd}
\]
whose induced map
$X_0\longrightarrow X\times_{\operatorname{Spec}(R)}
\operatorname{Spec}(k)$ is an isomorphism.

%%%%%%%%%%  scan idx 108  |  folio 97  %%%%%%%%%%

Let $U_0\subset X_0$ be an open subscheme. Since
$i_X:X_0\longrightarrow X$ is a closed immersion and a homeomorphism
on underlying spaces, restriction gives a commutative diagram
\[
\begin{tikzcd}[column sep=huge,row sep=huge]
U_0 \arrow[r,"i_X|_{U_0}"] \arrow[d]
& X|_{U_0} \arrow[d,"\mathrm{flat}"] \\
\operatorname{Spec}(k) \arrow[r] & \operatorname{Spec}(R).
\end{tikzcd}
\]
The induced map
\[
 U_0\longrightarrow
 (X|_{U_0})\times_{\operatorname{Spec}(R)}\operatorname{Spec}(k)
 =(X\otimes_Rk)|_{U_0}
\]
is an isomorphism, so $(R,X|_{U_0})$ is a deformation of $U_0$ over
$R$. Moreover, every arrow
$(\varphi,\Phi):(R,X)\longrightarrow(R',X')$ in
$\underline M_{X_0}$ restricts to an arrow
\[
 (\varphi,\Phi|_{U_0}):(R,X|_{U_0})
 \longrightarrow(R',X'|_{U_0})
\]
in $\underline M_{U_0}$. We therefore have a restriction functor
\[
 \underline M_{X_0}\longrightarrow\underline M_{U_0}
\]
over $\underline C_\Lambda$.

One may likewise localize and complete at a point. Let
$x_0\in X_0$. If $(R,X)\in\operatorname{Ob}(\underline M_{X_0})$,
then $\operatorname{Spec}(\mathcal O_{X,x_0})$ is a deformation of
$\operatorname{Spec}(\mathcal O_{X_0,x_0})$ over $R$, while
$\operatorname{Spec}(\widehat{\mathcal O}_{X,x_0})$ is a deformation
of $\operatorname{Spec}(\widehat{\mathcal O}_{X_0,x_0})$ over $R$;
the hats denote completion for the topology defined by powers of the
maximal ideal. Thus there are functors
\[
 \underline M_{X_0}
 \longrightarrow
 \underline M_{\operatorname{Spec}(\mathcal O_{X_0,x_0})}
 \longrightarrow
 \underline M_{\operatorname{Spec}(\widehat{\mathcal O}_{X_0,x_0})}.
\]

We first prove that $\underline M_{X_0}$ is a homogeneous cofibered
category over $\underline C_\Lambda$.

% END INCLUDED SGA7I_English_current_buildable_source_20260823/components/38_expose_VI_section_4_opening_through_4_1.tex

% BEGIN INCLUDED SGA7I_English_current_buildable_source_20260823/components/39_expose_VI_lemma_4_2_and_corollary_4_3.tex
%%%%%%%%%%  scan idx 108  |  folio 97  %%%%%%%%%%

\medskip
\noindent\textbf{Lemma \SGARefTarget{sga7.sga7i.evi.n4-2}{4.2}.}
Consider a diagram
\[
\begin{tikzcd}[column sep=huge,row sep=large]
& R' \arrow[d] \\
R \arrow[r,"\mathrm{surjective}"'] & \bar R
\end{tikzcd}
\]
in $\underline C_\Lambda$, and let $E$ and $E'$ be flat modules over
$R$ and $R'$, respectively. Then, for every $R'$-homomorphism
\[
 E'\longrightarrow \bar E=E\otimes_R\bar R,
\]
the fiber product $E\times_{\bar E}E'$ is flat as a module over
$R\times_{\bar R}R'$.

%%%%%%%%%%  scan idx 109  |  folio 98  %%%%%%%%%%

\medskip
\noindent\textit{Proof.}
Let
\[
 0\longrightarrow I\longrightarrow R\longrightarrow\bar R
 \longrightarrow0
\]
be exact, so that
\[
 0\longrightarrow I\longrightarrow R\times_{\bar R}R'
 \longrightarrow R'\longrightarrow0
\]
is exact. Since $E$ is flat over $R$,
\[
 0\longrightarrow I\otimes_RE\longrightarrow E
 \longrightarrow\bar E\longrightarrow0
\]
is exact, and hence we obtain the exact sequence
\begin{equation}\tag{*}
0\longrightarrow I\otimes_RE\longrightarrow
E\times_{\bar E}E'\longrightarrow E'\longrightarrow0.
\end{equation}
Set $S=R\times_{\bar R}R'$. Since $I$ is nilpotent,
$E\times_{\bar E}E'$ is flat as an $S$-module if and only if
$(E\times_{\bar E}E')\otimes_SR'$ is flat as an $R'$-module and
\[
 \operatorname{Tor}_1^S(E\times_{\bar E}E',R')=0.
\]
But
\[
 I\otimes_RE=I\otimes_S(E\times_{\bar E}E'),
 \qquad
 E'=R'\otimes_S(E\times_{\bar E}E'),
\]
and therefore the exact sequence $(*)$ may be rewritten as
\[
0\longrightarrow I\otimes_S(E\times_{\bar E}E')
\longrightarrow E\times_{\bar E}E'
\longrightarrow R'\otimes_S(E\times_{\bar E}E')
\longrightarrow0.
\]
Thus
\[
 \operatorname{Tor}_1^S(R',E\times_{\bar E}E')=0,
 \qquad
 R'\otimes_S(E\times_{\bar E}E')=E',
\]
and $E'$ is flat over $R'$. It follows that
$E\times_{\bar E}E'$ is flat over $S$. \hfill$\square$

\medskip
\noindent\textbf{Corollary \SGARefTarget{sga7.sga7i.evi.n4-3}{4.3}.}
Given a diagram
\[
\begin{tikzcd}[column sep=huge,row sep=large]
(R,X) \arrow[dr,"{(\varphi,\Phi)}"']
&& (R',X') \arrow[dl,"{(\varphi',\Phi')}"] \\
& (\bar R,\bar X) &
\end{tikzcd}
\]
in $\underline M_{X_0}$, the amalgamated sum
$X\amalg_{\bar X}X'$ (in the category of schemes) is a flat scheme
over $R\times_{\bar R}R'$, provided that
$\varphi:R\longrightarrow\bar R$ is surjective. Consequently,
$\underline M_{X_0}$ is a homogeneous cofibered category in groupoids
over $\underline C_\Lambda$.

\medskip
\noindent\textit{Proof.}
For every affine open $U_0\subset X_0$, by definition we have
\[
 (X\amalg_{\bar X}X')|_{U_0}
 \simeq
 \operatorname{Spec}\!\left(
   \Gamma(U_0,X)
   \times_{\Gamma(U_0,\bar X)}
   \Gamma(U_0,X')
 \right).
\]
Consequently, if $\varphi:R\longrightarrow\bar R$ is surjective,
$X\amalg_{\bar X}X'$ is flat over $R\times_{\bar R}R'$ by Lemma
\SGARefLink{sga7.sga7i.evi.n4-2}{4.2}. Moreover, the morphism
\[
 i_X\amalg_{i_{\bar X}}i_{X'}:
 X_0\longrightarrow X\amalg_{\bar X}X'
\]
induces an isomorphism
\[
 X_0\longrightarrow (X\amalg_{\bar X}X')\otimes_Sk,
 \qquad S=R\times_{\bar R}R'.
\]
Therefore
\[
 \bigl(R\times_{\bar R}R',X\amalg_{\bar X}X'\bigr)
 \in\operatorname{Ob}(\underline M_{X_0}),
\]
and, in $\underline M_{X_0}$, we have
\[
 \bigl(R\times_{\bar R}R',X\amalg_{\bar X}X'\bigr)
 =
 (R,X)\times_{(\bar R,\bar X)}(R',X').
\]
Thus $\underline M_{X_0}$ is homogeneous over
$\underline C_\Lambda$. \hfill$\square$

%%%%%%%%%%  scan idx 110  |  folio 99  %%%%%%%%%%

% END INCLUDED SGA7I_English_current_buildable_source_20260823/components/39_expose_VI_lemma_4_2_and_corollary_4_3.tex

% BEGIN INCLUDED SGA7I_English_current_buildable_source_20260823/components/40_expose_VI_lemma_4_4.tex
%%%%%%%%%%  scan idx 110  |  folio 99  %%%%%%%%%%

\medskip
\noindent\textbf{Lemma \SGARefTarget{sga7.sga7i.evi.n4-4}{4.4}.}
For every scheme $X_0$ over $k$, there is a canonical exact sequence
\[
0\longrightarrow H^1(X_0,\underline D^0_{X_0})
\longrightarrow[\underline M_{X_0}(k[\varepsilon])]
\longrightarrow H^0(X_0,\underline D^1_{X_0}),
\]
where the $\underline D^i_{X_0}$ are as defined in section 3.

\medskip
\noindent\textit{Proof.}
First suppose that $X_0$ is affine, and put
$\Gamma(X_0,\mathcal O_{X_0})=A_0$. As noted in \SGARefLink{sga7.sga7i.evi.n4-1.item-b}{4.1}(b), the category
$\underline M_{X_0}(k[\varepsilon])$ is canonically equivalent to a
category $\underline E$ whose objects are flat $k[\varepsilon]$-algebras
$A$ equipped with a $k[\varepsilon]$-algebra map
$j_A:A\longrightarrow A_0$ inducing an isomorphism
\[
k\otimes_{k[\varepsilon]}A\xrightarrow{\ \sim\ }A_0.
\]
Equivalently, an object of $\underline E$ is a $k$-algebra $A$
equipped with an exact sequence
\[
0\longrightarrow A_0\xrightarrow{\ \varepsilon\ }A
\xrightarrow{\ j_A\ }A_0\longrightarrow0
\]
in which $\varepsilon(A_0)^2=0$ as an ideal of $A$.

An arrow $A_1\longrightarrow A_2$ in $\underline E$, in the first
description, is a $k[\varepsilon]$-algebra map
$\Phi:A_1\longrightarrow A_2$ such that
$j_{A_1}=j_{A_2}\circ\Phi$. In the second description it is a
$k$-algebra map $\Phi:A_1\longrightarrow A_2$ yielding a commutative
diagram
\[
\begin{tikzcd}[column sep=large,row sep=large]
0 \arrow[r]
& A_0 \arrow[r,"\varepsilon"] \arrow[d,equal]
& A_1 \arrow[r,"j_{A_1}"] \arrow[d,"\Phi"]
& A_0 \arrow[r] \arrow[d,equal]
& 0 \\
0 \arrow[r]
& A_0 \arrow[r]
& A_2 \arrow[r,"j_{A_2}"]
& A_0 \arrow[r]
& 0.
\end{tikzcd}
\]
It follows that, via $\Gamma(X_0,-)$,
$[\underline M_{X_0}(k[\varepsilon])]$ is canonically isomorphic to
the set of isomorphism classes of $k$-algebra extensions of $A_0$ by
$A_0$. In other words, there is a canonical bijection
\[
[\underline M_{X_0}(k[\varepsilon])]
\xrightarrow{\ \sim\ }H^0(X_0,\underline D^1_{X_0})
\]
which plainly respects the vector-space structures on both sides.

%%%%%%%%%%  scan idx 111  |  folio 100  %%%%%%%%%%

In general, for every affine open $U_0\subset X_0$, restriction and
the affine case give the composite map
\[
[\underline M_{X_0}(k[\varepsilon])]
\longrightarrow[\underline M_{U_0}(k[\varepsilon])]
\xrightarrow{\ \sim\ }H^0(U_0,\underline D^1_{X_0}).
\]
These maps give a canonical map
\[
[\underline M_{X_0}(k[\varepsilon])]
\longrightarrow H^0(X_0,\underline D^1_{X_0}).
\]
Its kernel consists of the isomorphism classes of locally trivial
deformations of $X_0$ over $k[\varepsilon]$, that is, of deformations
locally isomorphic to $X_0\otimes_k k[\varepsilon]$. The sheaf of
automorphisms of $X_0\otimes_k k[\varepsilon]$ inducing the identity
on $X_0$ is
\[
\underline D^0_{X_0}
=\underline{\operatorname{Der}}_k
  (\mathcal O_{X_0},\mathcal O_{X_0}).
\]
The isomorphism classes of such locally trivial deformations are
therefore canonically identified with
$H^1(X_0,\underline D^0_{X_0})$, which gives the asserted exact
sequence. \hfill$\square$

% END INCLUDED SGA7I_English_current_buildable_source_20260823/components/40_expose_VI_lemma_4_4.tex

% BEGIN INCLUDED SGA7I_English_current_buildable_source_20260823/components/41_expose_VI_theorem_4_5.tex
%%%%%%%%%%  scan idx 111  |  folio 100  %%%%%%%%%%

\medskip
\noindent\textbf{Theorem \SGARefTarget{sga7.sga7i.evi.n4-5}{4.5}.}
Let $X_0$ be a scheme of finite type over $k$ for which either

\smallskip
\noindent(i) $X_0$ is proper over $k$, or

\noindent(ii) $X_0$ is affine and
\[
X_0^{\mathrm{sing}}
=\{x\in X_0\mid X_0\text{ is not smooth over }k\text{ at }x\}
\]
is finite.

\smallskip
\noindent\textbf{(a)}
The category $\underline M_{X_0}$ admits a formal versal object over
$\underline C_\Lambda$, called a formal versal deformation of $X_0$.
If $(R,\mathcal X)$ is a formal versal deformation of $X_0$, then for
every $(S,\mathcal Y)\in\operatorname{Ob}(\widehat{\underline M}_{X_0})$
there is a map
\[
(\varphi,\Phi):(R,\mathcal X)\longrightarrow(S,\mathcal Y)
\]
in $\widehat{\underline M}_{X_0}$ such that the following diagram is
cartesian:
\[
\begin{tikzcd}[column sep=huge,row sep=large]
\mathcal Y \arrow[r,"\Phi"] \arrow[d]
& \mathcal X \arrow[d] \\
\operatorname{Spf}(S)
 \arrow[r,"\operatorname{Spf}(\varphi)"']
& \operatorname{Spf}(R).
\end{tikzcd}
\]
Furthermore, a formal versal deformation $(R,\mathcal X)$ of $X_0$
prorepresents the functor $[\underline M_{X_0}]$ if and only if, for
every surjective map $R'\longrightarrow R$ in $\underline C_\Lambda$
and every deformation $(R',X')$ of $X_0$ over $R'$, the homomorphism
\[
\operatorname{Aut}_{R'}(X'\mid X_0)
\longrightarrow
\operatorname{Aut}_R(X'\otimes_{R'}R\mid X_0)
\]
is surjective.

\smallskip
\noindent\textbf{(b)}
If $X_0$ is proper over $k$, then $\underline M_{X_0}$ is smoothly
pseudo-representable over $\underline C_\Lambda$.

%%%%%%%%%%  scan idx 112  |  folio 101  %%%%%%%%%%

\medskip
\noindent\textit{Proof.}
\textbf{(a)} By Corollary 4.3, $\underline M_{X_0}$ is a homogeneous
cofibered category in groupoids over $\underline C_\Lambda$. By \SGARefLink{sga7.sga7i.evi.n1-11}{1.11}
and \SGARefLink{sga7.sga7i.evi.n2-9}{2.9}(b), it therefore suffices to prove that
$[\underline M_{X_0}(k[\varepsilon])]$ is finite-dimensional over
$k$.

Consider the exact sequence
\begin{equation}\tag{*}
0\longrightarrow H^1(X_0,\underline D^0_{X_0})
\longrightarrow[\underline M_{X_0}(k[\varepsilon])]
\longrightarrow H^0(X_0,\underline D^1_{X_0}).
\end{equation}
If $X_0$ is proper over $k$, both
$H^1(X_0,\underline D^0_{X_0})$ and
$H^0(X_0,\underline D^1_{X_0})$ are finite-dimensional, since the
$\underline D^i_{X_0}$ are coherent for all $i\geq0$ by \SGARefLink{sga7.sga7i.evi.n2-13}{2.13}(a).

Now suppose that $X_0$ is affine and $X_0^{\mathrm{sing}}$ is finite.
The sheaf $\underline D^1_{X_0}$ is coherent and
\[
\operatorname{Supp}\underline D^1_{X_0}
\subset X_0^{\mathrm{sing}}
\]
by \SGARefLink{sga7.sga7i.evi.n2-13}{2.13}(a). Hence $H^0(X_0,\underline D^1_{X_0})$ is
finite-dimensional over $k$, whereas
$H^1(X_0,\underline D^0_{X_0})=0$ because $X_0$ is affine. In either
case, the exact sequence $(*)$ shows that
$[\underline M_{X_0}(k[\varepsilon])]$ is finite-dimensional over
$k$. The second assertion of part (a) follows from \SGARefLink{sga7.sga7i.evi.n2-7}{2.7}.

\smallskip
\noindent\textbf{(b)} If $X_0$ is proper, then
\[
\operatorname{Aut}(X_0\otimes_k k[\varepsilon]\mid X_0)
=H^0(X_0,\underline D^0_{X_0})
\]
is finite-dimensional, and the assertion follows from \SGARefLink{sga7.sga7i.evi.n2-13}{2.13}.
\hfill$\square$

% END INCLUDED SGA7I_English_current_buildable_source_20260823/components/41_expose_VI_theorem_4_5.tex

% BEGIN INCLUDED SGA7I_English_current_buildable_source_20260823/components/42_expose_VI_remarks_4_6_and_pair_deformations.tex
%%%%%%%%%%  scan idx 112  |  folio 101  %%%%%%%%%%

\medskip
\noindent\textbf{Remark \SGARefTarget{sga7.sga7i.evi.n4-6.item-a}{\SGARefAliasTarget{sga7.sga7i.evi.n4-6}4.6}(a).}
Let $X_0$ be affine. A formal versal deformation of $X_0$ exists if
and only if the support of $\underline D^1_{X_0}$ is a finite set.
This does not imply, however, that all nonsmooth points of $X_0$ are
isolated: there are affine schemes $X_0$ over $k$ having a
nonisolated nonsmooth point for which
$\operatorname{Supp}\underline D^1_{X_0}$ is finite.

\medskip
\noindent\textbf{Remark \SGARefTarget{sga7.sga7i.evi.n4-6.item-b}{4.6}(b).}
Henceforth, a deformation of $X_0$ means an object of
$\underline M_{X_0}$. Where confusion is possible, such an object
will be called an infinitesimal deformation of $X_0$.

\medskip
\noindent\textbf{Remark \SGARefTarget{sga7.sga7i.evi.n4-6.item-c}{4.6}(c).}
Let $X_0$ be a scheme over $k$. We have fixed a complete noetherian
local ring $\Lambda$ with residue field $k$ and considered
deformations of $X_0$ over $\underline C_\Lambda$.

%%%%%%%%%%  scan idx 113  |  folio 102  %%%%%%%%%%

In practice one usually takes $\Lambda=k$, or $\Lambda=W(k)$ when
$k$ is a perfect field of positive characteristic. The inclusion
$\underline C_k\subset\underline C_\Lambda$ is fully faithful, and
$\underline M_{X_0}|_{\underline C_k}$ is simply the category of
deformations of $X_0$ over the objects of $\underline C_k$.
Consequently, if $(R,\mathfrak X)$ is a formal versal deformation of
$X_0$ over $\underline C_\Lambda$, then
\[
 (k\otimes_\Lambda R,\;k\otimes_\Lambda\mathfrak X)
\]
is a formal versal deformation of $X_0$ over $\underline C_k$. It
follows, for example, that
\[
 \dim_{\underline C_\Lambda}\underline M_{X_0}
 =\dim(k\otimes_\Lambda R)
\]
does not depend on the choice of $\Lambda$, but only on $X_0$.

The preceding theorem admits many generalizations. For example,
consider a pair $Y_0\longrightarrow X_0$, where $X_0$ is a scheme of
finite type over $k$ and $Y_0$ is a fixed closed subscheme of $X_0$.
For $R\in\operatorname{Ob}(\underline C_\Lambda)$, a deformation of
$(X_0,Y_0)$ over $R$ is a commutative diagram
\[
\begin{tikzcd}[row sep=large,column sep=huge]
Y_0 \arrow[rrr] \arrow[dr,hook] \arrow[ddr]
&&& Y \arrow[dl,hook] \arrow[ddl,"\mathrm{flat}"] \\
& X_0 \arrow[r] \arrow[d]
& X \arrow[d,"\mathrm{flat}"'] & \\
& \operatorname{Spec}(k) \arrow[r]
& \operatorname{Spec}(R) &
\end{tikzcd}
\]
such that the induced map
\[
(X_0,Y_0)\longrightarrow
\left(
 X\times_{\operatorname{Spec}(R)}\operatorname{Spec}(k),
 Y\times_{\operatorname{Spec}(R)}\operatorname{Spec}(k)
\right)
\]
is an isomorphism. Defining morphisms
\[
(R';X',Y')\longrightarrow(R;X,Y)
\]
of deformations of $(X_0,Y_0)$ in the evident way gives a cofibered
category $\underline M_{X_0,Y_0}$ in groupoids over
$\underline C_\Lambda$. This construction is frequently used to
rigidify automorphisms in deformation-moduli problems. The argument
of Corollary \SGARefLink{sga7.sga7i.evi.n4-3}{4.3}, based on Lemma \SGARefLink{sga7.sga7i.evi.n4-2}{4.2}, immediately shows that
$\underline M_{X_0,Y_0}$ is a homogeneous cofibered category in
groupoids over $\underline C_\Lambda$.

% END INCLUDED SGA7I_English_current_buildable_source_20260823/components/42_expose_VI_remarks_4_6_and_pair_deformations.tex

% BEGIN INCLUDED SGA7I_English_current_buildable_source_20260823/components/43_expose_VI_lemma_4_7_through_theorem_4_9.tex
%%%%%%%%%%  scan idx 113  |  folio 102  %%%%%%%%%%

\medskip
\noindent\textbf{Lemma \SGARefTarget{sga7.sga7i.evi.n4-7}{4.7}.}
The sequence
\[
0\longrightarrow
H^0\!\left(
Y_0,
\underline{\operatorname{Hom}}_{\mathcal O_{Y_0}}
(\underline I/\underline I^2,\mathcal O_{Y_0})
\right)
\longrightarrow[\underline M_{X_0,Y_0}(k[\varepsilon])]
\longrightarrow[\underline M_{X_0}(k[\varepsilon])]
\]
is exact, where $\underline I$ is the ideal sheaf defining $Y_0$ in
$X_0$; thus
\[
0\longrightarrow\underline I\longrightarrow\mathcal O_{X_0}
\longrightarrow\mathcal O_{Y_0}\longrightarrow0
\]
is exact.

%%%%%%%%%%  scan idx 114  |  folio 103  %%%%%%%%%%

\medskip
\noindent\textit{Proof.}
Let $\underline B$ be the full subcategory of
$\underline M_{X_0,Y_0}(k[\varepsilon])$ whose objects have the form
\[
\bigl(k[\varepsilon];
X_0\times_{\operatorname{Spec}(k)}\operatorname{Spec}(k[\varepsilon]),
Y\bigr).
\]
Then
\[
[\underline B]
=\ker\!\left(
[\underline M_{X_0,Y_0}(k[\varepsilon])]
\longrightarrow[\underline M_{X_0}(k[\varepsilon])]
\right).
\]
The category $\underline B$ is the groupoid of deformations of $Y_0$
over $k[\varepsilon]$ inside
$X_0\times_{\operatorname{Spec}(k)}\operatorname{Spec}(k[\varepsilon])$.
Thus an object is a commutative diagram
\[
\begin{tikzcd}[column sep=huge,row sep=large]
Y \arrow[rr,"i_Y"] \arrow[dr,"\mathrm{flat}"']
&& X_0\times_{\operatorname{Spec}(k)}
   \operatorname{Spec}(k[\varepsilon]) \arrow[dl] \\
& \operatorname{Spec}(k[\varepsilon]) &
\end{tikzcd}
\]
whose special fiber is the fixed closed immersion $i_{Y_0}$. An
arrow $(Y_1,i_{Y_1})\longrightarrow(Y_2,i_{Y_2})$ is a
$k[\varepsilon]$-morphism $\Phi:Y_1\longrightarrow Y_2$ (necessarily
an isomorphism) whose special fiber is $\operatorname{id}_{Y_0}$ and
for which $i_{Y_2}\circ\Phi=i_{Y_1}$.

Equivalently, an object of $\underline B$ is a pair $(Y,j_Y)$, where
$Y$ is a deformation of $Y_0$ over $k[\varepsilon]$ and
$j_Y:Y\longrightarrow X_0$ is a $k$-morphism whose composite
$Y_0\longrightarrow Y\xrightarrow{j_Y}X_0$ is $i_{Y_0}$. An arrow
$(Y_1,j_{Y_1})\longrightarrow(Y_2,j_{Y_2})$ is a morphism
$\Phi:Y_1\longrightarrow Y_2$ whose special fiber is
$\operatorname{id}_{Y_0}$ and for which
$j_{Y_2}\circ\Phi=j_{Y_1}$.

The assignment sending $(Y,j_Y)$ to the diagram
\[
\begin{tikzcd}[column sep=large,row sep=large]
0 \arrow[r]
& \underline I \arrow[r] \arrow[d]
& \mathcal O_{X_0} \arrow[r] \arrow[d,"j_Y"]
& \mathcal O_{Y_0} \arrow[r] \arrow[d,equal]
& 0 \\
0 \arrow[r]
& \mathcal O_{Y_0} \arrow[r,"\varepsilon"]
& \mathcal O_Y \arrow[r]
& \mathcal O_{Y_0} \arrow[r]
& 0
\end{tikzcd}
\]
gives an isomorphism
\[
\begin{aligned}
[\underline B]
&\xrightarrow{\ j_Y|_{\underline I}\ }
\operatorname{Hom}_{\mathcal O_{X_0}}
 (\underline I,\mathcal O_{Y_0}) \\
&=
\operatorname{Hom}_{\mathcal O_{Y_0}}
 (\underline I/\underline I^2,\mathcal O_{Y_0}) \\
&=
H^0\!\left(
Y_0,
\underline{\operatorname{Hom}}_{\mathcal O_{Y_0}}
(\underline I/\underline I^2,\mathcal O_{Y_0})
\right).
\end{aligned}
\]
This proves exactness. \hfill$\square$

%%%%%%%%%%  scan idx 115  |  folio 104  %%%%%%%%%%

\medskip
\noindent\textbf{Remark \SGARefTarget{sga7.sga7i.evi.n4-8}{4.8}.}
For a more general formula implying Lemma \SGARefLink{sga7.sga7i.evi.n4-7}{4.7}, see SGA 3, III, 4.4.
One may also note the following. For every $S$-scheme $X$ and every
quasi-coherent $\mathcal O_X$-module $E$, one can define
quasi-coherent $\mathcal O_X$-modules $\underline D^i_{X|S}(E)$ for
all $i\geq0$, generalizing the case in which $S$ is affine. As usual,
we write
\[
\underline D^i_{X|S}(\mathcal O_X)=\underline D^i_{X|S}.
\]
For the closed immersion $Y_0\hookrightarrow X_0$, we then have
\[
H^0\!\left(
Y_0,
\underline{\operatorname{Hom}}_{\mathcal O_{Y_0}}
(\underline I/\underline I^2,\mathcal O_{Y_0})
\right)
=H^0(Y_0,\underline D^1_{Y_0|X_0}),
\]
where
$0\longrightarrow\underline I\longrightarrow\mathcal O_{X_0}
\longrightarrow\mathcal O_{Y_0}\longrightarrow0$ is exact.

\medskip
\noindent\textbf{Theorem \SGARefTarget{sga7.sga7i.evi.n4-9}{4.9}.}
Let $X_0$ be a $k$-scheme of finite type and $Y_0$ a closed
subscheme of $X_0$. Suppose either

\smallskip
\noindent(i) $X_0$ is proper over $k$, or

\noindent(ii) $X_0$ is affine and smooth over $k$ outside finitely
many points, and $Y_0$ is proper over $k$.

\smallskip
\noindent\textbf{(a)}
The category $\underline M_{X_0,Y_0}$ admits a formal versal
deformation over $\underline C_\Lambda$. It prorepresents the functor
$[\underline M_{X_0,Y_0}]$ if and only if, for every surjective map
$R'\longrightarrow R$ in $\underline C_\Lambda$ and every deformation
$(R';X',Y')$, the map
\[
\operatorname{Aut}_{R'}((X',Y')\mid X_0)
\longrightarrow
\operatorname{Aut}_{R}\!\left(
(X'\otimes_{R'}R,Y'\otimes_{R'}R)\mid X_0
\right)
\]
is surjective.

\smallskip
\noindent\textbf{(b)}
If, in addition,
\[
H^0\!\left(
X_0,
\underline D^0_{Y_0|k}
\times_{\underline D^0_{X_0|k}(\mathcal O_{Y_0})}
\underline D^0_{X_0|k}
\right)
\]
is finite-dimensional, then $\underline M_{X_0,Y_0}$ is smoothly
pseudo-representable.

\medskip
\noindent\textit{Proof.}
This follows immediately from Lemmas \SGARefLink{sga7.sga7i.evi.n4-4}{4.4} and \SGARefLink{sga7.sga7i.evi.n4-7}{4.7} together with \SGARefLink{sga7.sga7i.evi.n1-11}{1.11},
\SGARefLink{sga7.sga7i.evi.n2-7}{2.7}, and \SGARefLink{sga7.sga7i.evi.n2-13}{2.13}. \hfill$\square$

% END INCLUDED SGA7I_English_current_buildable_source_20260823/components/43_expose_VI_lemma_4_7_through_theorem_4_9.tex

% BEGIN INCLUDED SGA7I_English_current_buildable_source_20260823/components/44_expose_VI_localization_and_completion_intro.tex
%%%%%%%%%%  scan idx 115  |  folio 104  %%%%%%%%%%

We now consider localization and completion in the deformation theory
of schemes. We first fix notation. Let $X_0$ be a $k$-scheme of finite
type and let $x_0\in X_0$ be a fixed point. There are natural functors
\[
\underline M_{X_0}
\longrightarrow
\underline M_{\operatorname{Spec}(\mathcal O_{X_0,x_0})}
\longrightarrow
\underline M_{\operatorname{Spec}(\widehat{\mathcal O}_{X_0,x_0})},
\]

%%%%%%%%%%  scan idx 116  |  folio 105  %%%%%%%%%%

and hence functors
\[
\widehat{\underline M}_{X_0}
\longrightarrow
\widehat{\underline M}_{\operatorname{Spec}(\mathcal O_{X_0,x_0})}
\longrightarrow
\widehat{\underline M}_{\operatorname{Spec}
 (\widehat{\mathcal O}_{X_0,x_0})}.
\]
As observed in \SGARefLink{sga7.sga7i.evi.n4-1.item-c}{4.1}(c), $\widehat{\underline M}_{X_0}$ is canonically
equivalent to the category whose objects are pairs $(R,\mathfrak X)$,
where $R\in\operatorname{Ob}(\widehat{\underline C}_\Lambda)$ and
$\mathfrak X$ is a flat $R$-adic formal scheme equipped with a fixed
isomorphism $\mathfrak X_0\xrightarrow{\sim}X_0$.

When there is no danger of confusion, we denote the images of
$(R,\mathfrak X)$ under localization and completion by
$(R,\mathfrak X_{(x_0)})$ and
$(R,\mathfrak X_{(\widehat{x}_0)})$, respectively. Thus, if
\[
\mathfrak X=\varinjlim_\nu\mathfrak X_\nu,
\qquad
\mathfrak X_\nu
=\mathfrak X\times_{\operatorname{Spf}(R)}
 \operatorname{Spec}(R/\mathfrak m^{\nu+1}),
\]
then
\[
\mathfrak X_{(x_0)}
=\varinjlim_\nu
 \operatorname{Spec}(\mathcal O_{\mathfrak X_\nu,x_0})
\]
and
\[
\mathfrak X_{(\widehat{x}_0)}
=\varinjlim_\nu
 \operatorname{Spec}(\widehat{\mathcal O}_{\mathfrak X_\nu,x_0}).
\]

% END INCLUDED SGA7I_English_current_buildable_source_20260823/components/44_expose_VI_localization_and_completion_intro.tex

% BEGIN INCLUDED SGA7I_English_current_buildable_source_20260823/components/45_expose_VI_theorem_4_10.tex
%%%%%%%%%%  scan idx 116  |  folio 105  %%%%%%%%%%

\medskip
\noindent\textbf{Theorem \SGARefTarget{sga7.sga7i.evi.n4-10}{4.10}.}
Let $X_0$ be an affine scheme of finite type over $k$.

\smallskip
\noindent\textbf{(a)}
If $X_0$ is locally a complete intersection, then
$\underline M_{X_0}$ is smooth over $\underline C_\Lambda$.

\smallskip
\noindent\textbf{(b)}
Suppose that, for a fixed closed point $x\in X_0$, the open
$X_0-\{x\}$ is smooth over $k$. Then both functors
\[
\underline M_{X_0}
\longrightarrow
\underline M_{\operatorname{Spec}(\mathcal O_{X_0,x})}
\longrightarrow
\underline M_{\operatorname{Spec}(\widehat{\mathcal O}_{X_0,x})}
\]
are minimally smooth. In particular, if $(R,\mathfrak X)$ is a formal
versal deformation of $X_0$, then $(R,\mathfrak X_{(x)})$ and
$(R,\mathfrak X_{(\widehat x)})$ are formal versal deformations of
$\operatorname{Spec}(\mathcal O_{X_0,x})$ and
$\operatorname{Spec}(\widehat{\mathcal O}_{X_0,x})$, respectively.

\medskip
\noindent\textit{Proof.}
We use the notation for the cohomology of commutative algebras from
section 3, and write $X_0=\operatorname{Spec}(A_0)$.

%%%%%%%%%%  scan idx 117  |  folio 106  %%%%%%%%%%

\smallskip
\noindent\textbf{(a)}
It suffices to show that
$P:[\underline M_{X_0}]\longrightarrow\underline C_\Lambda$ is
smooth. Let $(R,\operatorname{Spec}(A))$ be a deformation of
$\operatorname{Spec}(A_0)$ over
$R\in\operatorname{Ob}(\underline C_\Lambda)$. Given a small
extension $R'\longrightarrow R$ in $\underline C_\Lambda$, so that
\[
0\longrightarrow k\longrightarrow R'\longrightarrow R
\longrightarrow0
\]
is exact, we must lift $(R,\operatorname{Spec}(A))$ to $R'$. The
transitivity sequence for $D^i$ with coefficients in $A_0$, applied
to $R'\longrightarrow R\longrightarrow A$, gives
\[
\cdots\longrightarrow D^1(A|R,A_0)
\longrightarrow D^1(A|R',A_0)
\longrightarrow D^1(R|R',k)\otimes_kA_0
\longrightarrow D^2(A|R,A_0)\longrightarrow\cdots .
\]
Because $R'\longrightarrow R$ is a small extension,
$D^1(R|R',k)$ is one-dimensional over $k$, generated by the class
$[R']$ represented by
$0\longrightarrow k\longrightarrow R'\longrightarrow R
\longrightarrow0$. Since $A$ is flat over $R$,
\[
D^i(A|R,A_0)=D^i(A_0|k,A_0)
\qquad(i>0).
\]
In particular, $D^2(A|R,A_0)=0$ because $A_0$ is locally a complete
intersection, by the low-degree description in \SGARefLink{sga7.sga7i.evi.n3-12}{3.12}(8). We therefore
obtain an exact sequence
\[
\longrightarrow D^1(A_0|k,A_0)
\longrightarrow D^1(A|R',A_0)
\longrightarrow D^1(R|R',k)\otimes_kA_0
\longrightarrow0.
\]
Hence there is an $R'$-algebra extension $A'$ of $A$ by $A_0$ whose
class maps to $[R']\otimes1$. Equivalently,
$(R,\operatorname{Spec}(A))$ lifts to the deformation
$(R',\operatorname{Spec}(A'))$.

\smallskip
\noindent\textbf{(b)}
By \SGARefLink{sga7.sga7i.evi.n3-14}{3.14} there are isomorphisms
\[
D^1(A_0|k,A_0)
\xrightarrow{\ \sim\ }
D^1((A_0)_x|k,(A_0)_x)
\xrightarrow{\ \sim\ }
D^1(\widehat{(A_0)_x}|k,\widehat{(A_0)_x}),
\]
and hence
\[
[\underline M_{X_0}(k[\varepsilon])]
\xrightarrow{\ \sim\ }
[\underline M_{\operatorname{Spec}(\mathcal O_{X_0,x})}
 (k[\varepsilon])]
\xrightarrow{\ \sim\ }
[\underline M_{\operatorname{Spec}(\widehat{\mathcal O}_{X_0,x})}
 (k[\varepsilon])].
\]
By \SGARefLink{sga7.sga7i.evi.n2-7}{2.7}(a), it remains to show that
\[
[\underline M_{X_0}]
\longrightarrow
[\underline M_{\operatorname{Spec}(\mathcal O_{X_0,x})}]
\longrightarrow
[\underline M_{\operatorname{Spec}(\widehat{\mathcal O}_{X_0,x})}]
\]
are smooth functors.

Let $(R,\operatorname{Spec}(A))$ be a deformation of
$X_0=\operatorname{Spec}(A_0)$, and let
$0\longrightarrow k\longrightarrow R'\longrightarrow R
\longrightarrow0$ be a small extension. As above, the homomorphisms
$R'\longrightarrow R\longrightarrow A$ give an exact commutative
diagram

%%%%%%%%%%  scan idx 118  |  folio 107  %%%%%%%%%%

\begin{center}
\scriptsize
\begin{tikzcd}[column sep=tiny,row sep=large]
{} \arrow[r]
& D^1(A_0|k,A_0) \arrow[r] \arrow[d,"\sim"']
& D^1(A|R',A_0) \arrow[r,"p"] \arrow[d]
& D^1(R|R',k)\otimes_kA_0 \arrow[r] \arrow[d]
& D^2(A_0|k,A_0) \arrow[r] \arrow[d,"\sim"]
& \cdots \\
{} \arrow[r]
& D^1(A_0|k,(A_0)_x) \arrow[r]
& D^1(A|R',(A_0)_x) \arrow[r]
& D^1(R|R',k)\otimes_k(A_0)_x \arrow[r]
& D^2(A_0|k,(A_0)_x) \arrow[r]
& \cdots \\
{} \arrow[r]
& D^1((A_0)_x|k,(A_0)_x) \arrow[r] \arrow[u]
& D^1(A_x|R',(A_0)_x) \arrow[r,"p_x"] \arrow[u]
& D^1(R|R',k)\otimes_k(A_0)_x \arrow[r] \arrow[u]
& D^2((A_0)_x|k,(A_0)_x) \arrow[r] \arrow[u]
& \cdots
\end{tikzcd}
\end{center}
The indicated isomorphisms follow from the fact that $A_0$ is smooth
over $k$ away from the single closed point $x$, so that
\[
D^i(A_0|k,(A_0)_x)
=D^i(A_0|k,A_0)_x
=D^i(A_0|k,A_0)
\qquad(i>0).
\]

Suppose we are given an arrow
\[
(R',\operatorname{Spec}(B))
\longrightarrow(R,\operatorname{Spec}(A_x))
\]
in $[\underline M_{\operatorname{Spec}((A_0)_x)}]$. This amounts to
a class $[B]\in D^1(A_x|R',(A_0)_x)$ satisfying
$p_x([B])=[R']$. A diagram chase gives a class
$[A']\in D^1(A|R',A_0)$ such that $p([A'])=[R']$ and $A'_x=B$.
Thus $(R',\operatorname{Spec}(A'))$ is a deformation of
$\operatorname{Spec}(A_0)$ over $R'$ whose localization at $x$ is
$\operatorname{Spec}(B)$. This proves that
\[
\underline M_{X_0}\longrightarrow
\underline M_{\operatorname{Spec}(\mathcal O_{X_0,x})}
\]
is smooth.

The smoothness of
\[
\underline M_{\operatorname{Spec}((A_0)_x)}
\longrightarrow
\underline M_{\operatorname{Spec}(\widehat{(A_0)_x})}
\]
follows from the analogous diagram and the identities
\[
D^i((A_0)_x|k,\widehat{(A_0)_x})
=\widehat{D^i((A_0)_x|k,(A_0)_x)}
=D^i((A_0)_x|k,(A_0)_x)
\qquad(i>0),
\]
since $x$ is the only possible nonsmooth point. \hfill$\square$

% END INCLUDED SGA7I_English_current_buildable_source_20260823/components/45_expose_VI_theorem_4_10.tex

% BEGIN INCLUDED SGA7I_English_current_buildable_source_20260823/components/46_expose_VI_theorem_4_11_and_remark_4_12.tex
%%%%%%%%%%  scan idx 118  |  folio 107  %%%%%%%%%%

\medskip
\noindent\textbf{Theorem \SGARefTarget{sga7.sga7i.evi.n4-11}{4.11}.}
Let $X_0$ be a separated scheme of finite type over $k$, smooth away
from a finite closed subset
\[
X_0^{\mathrm{Sing}}=\{x_1,\ldots,x_r\}.
\]
Let $U_0^1,\ldots,U_0^r$ be affine open neighbourhoods of
$x_1,\ldots,x_r$ such that
$U_0^i\cap X_0^{\mathrm{Sing}}=\{x_i\}$. If
$H^2(X_0,\underline D^0_{X_0})=0$, then the natural functor
\[
\underline M_{X_0}\longrightarrow
\underline M_{U_0^1}\times\cdots\times\underline M_{U_0^r}
\]
is smooth.

\smallskip
\noindent\textit{Proof.}
Complete $U_0^1,\ldots,U_0^r$ to an affine open covering of $X_0$
by choosing $U_0^{r+1},\ldots,U_0^N$ such that
$U_0^j\cap X_0^{\mathrm{Sing}}=\varnothing$ for every $j>r$.
Let $R'\longrightarrow R$ be a small extension in
$\underline C_\Lambda$, and let
$(R',U'^1),\ldots,(R',U'^r)$ be deformations of
$U_0^1,\ldots,U_0^r$ over $R'$ such that
\[
U'^i\otimes_{R'}R=X|_{U_0^i},\qquad 1\leq i\leq r,
\]
for some deformation $(R,X)$ of $X_0$ over $R$. We must find a
deformation $(R',X')$ of $X_0$ over $R'$ such that
\[
X'\otimes_{R'}R=X,
\qquad
X'|_{U_0^i}=U'^i\quad(1\leq i\leq r).
\]

%%%%%%%%%%  scan idx 119  |  folio 108  %%%%%%%%%%

For $j>r$, the scheme $U_0^j$ is smooth over $k$; hence, by \SGARefLink{sga7.sga7i.evi.n4-9}{4.9}(a),
$X|_{U_0^j}$ lifts to a deformation $(R',U'^j)$ of $U_0^j$ over
$R'$. Thus
\[
U'^j\otimes_{R'}R=X|_{U_0^j}\qquad(j>r).
\]
For every pair $i\neq j$, the intersection $U_0^i\cap U_0^j$ is
affine and smooth over $k$. We therefore obtain an isomorphism
\[
\tau_{ij}:U'^i|_{U_0^i\cap U_0^j}
\xrightarrow{\ \sim\ }
U'^j|_{U_0^i\cap U_0^j}
\]
such that
$\tau_{ij}\otimes_{R'}R=\operatorname{id}_{X|_{U_0^i\cap U_0^j}}$.
On triple intersections, the automorphisms
\[
\tau_{ij}\tau_{jk}\tau_{ik}^{-1}
\]
define a cohomology class in
$H^2(X_0,\underline D^0_{X_0})$. This class is the obstruction to
gluing the $U'^i$ to a global deformation. It vanishes by hypothesis,
so the $U'^i$ glue to a deformation $(R',X')$ with
$X'\otimes_{R'}R=X$ and $X'|_{U_0^i}\simeq U'^i$. \hfill$\square$

\medskip
\noindent\textbf{Remark \SGARefTarget{sga7.sga7i.evi.n4-12}{4.12}.}
The preceding proof assumes that $X_0$ is separated, but this
assumption is unnecessary. Indeed, fix a deformation $(R,X)$ of
$X_0$ and deformations $(R',U'^i)$ of the members $U_0^i$ of an open
covering. The deformations $(R',V')$ of open subsets $V$ of $X$ over
$R'$ that are compatible with the already-given deformations
$(R',U'^i)$ form a ``gerbe'' (in Giraud's terminology) whose
structure sheaf is $\underline D^0_{X_0}$. The obstruction to the
existence of a section of this gerbe---that is, to a global
deformation $(R',X')$ of $X$ compatible with the given data
$(R',U'^i)$---lies in $H^2(X_0,\underline D^0_{X_0})$, and hence
vanishes under the hypothesis of the theorem.


% END INCLUDED SGA7I_English_current_buildable_source_20260823/components/46_expose_VI_theorem_4_11_and_remark_4_12.tex

% BEGIN INCLUDED SGA7I_English_current_buildable_source_20260823/components/47_expose_VI_corollary_4_13.tex
\medskip
\noindent\textbf{Corollary \SGARefTarget{sga7.sga7i.evi.n4-13}{4.13}.}
Let $X_0$ be a scheme of finite type over $k$, smooth away from a
finite subset
\[
X_0^{\mathrm{Sing}}=\{x_1,\ldots,x_r\},
\]
and suppose that $H^2(X_0,\underline D^0_{X_0})=0$. Then:

\smallskip
\noindent\textbf{(a)} The functor
\[
\underline M_{X_0}\longrightarrow
\prod_{i=1}^r
\underline M_{\operatorname{Spec}(\widehat{\mathcal O}_{X_0,x_i})}
\]
is smooth. In particular,
\[
\dim\underline M_{X_0}
=\dim_k H^1(X_0,\underline D^0_{X_0})
+\sum_{i=1}^r
\dim\underline M_{\operatorname{Spec}(\widehat{\mathcal O}_{X_0,x_i})}.
\]

\smallskip
\noindent\textbf{(b)} If, in addition, $X_0$ is locally a complete
intersection, then $\underline M_{X_0}$ is smooth.

%%%%%%%%%%  scan idx 120  |  folio 109  %%%%%%%%%%

\smallskip
\noindent\textit{Proof.}
\textbf{(a)} The first assertion follows from \SGARefLink{sga7.sga7i.evi.n4-10}{4.10} and \SGARefLink{sga7.sga7i.evi.n4-9}{4.9}(b).
Let $(R,\mathfrak X)$ and
$(R_1,\mathfrak X_1),\ldots,(R_r,\mathfrak X_r)$ be formal versal
deformations of $X_0$ and of
\[
\operatorname{Spec}(\mathcal O_{X_0,x_1}),\ldots,
\operatorname{Spec}(\mathcal O_{X_0,x_r}),
\]
respectively. We obtain a commutative diagram
\[
\begin{tikzcd}[column sep=large,row sep=large]
\underline h_R
  \arrow[r]
  \arrow[d]
& \underline h_{R_1\widehat\otimes_\Lambda\cdots
  \widehat\otimes_\Lambda R_r}
  \arrow[d] \\
\underline M_{X_0}
  \arrow[r]
& \displaystyle\prod_{i=1}^r
  \underline M_{\operatorname{Spec}(\widehat{\mathcal O}_{X_0,x_i})}.
\end{tikzcd}
\]
The composite functor
\[
\underline h_R\longrightarrow\underline M_{X_0}
\longrightarrow
\prod_{i=1}^r
\underline M_{\operatorname{Spec}(\widehat{\mathcal O}_{X_0,x_i})}
\]
is smooth, whereas
\[
\underline h_{R_1\widehat\otimes_\Lambda\cdots
\widehat\otimes_\Lambda R_r}
\longrightarrow
\prod_{i=1}^r
\underline M_{\operatorname{Spec}(\widehat{\mathcal O}_{X_0,x_i})}
\]
is minimally smooth. It follows from \SGARefLink{sga7.sga7i.evi.n2-7}{2.7}(a) that
\[
\underline h_R\longrightarrow
\underline h_{R_1\widehat\otimes_\Lambda\cdots
\widehat\otimes_\Lambda R_r}
\]
is smooth; equivalently, $R$ is a formal power-series ring over
$R_1\widehat\otimes_\Lambda\cdots\widehat\otimes_\Lambda R_r$.

By \SGARefLink{sga7.sga7i.evi.n4-4}{4.4}, the sequence
\[
0\longrightarrow H^1(X_0,\underline D^0_{X_0})
\longrightarrow[\underline M_{X_0}(k[\varepsilon])]
\longrightarrow
\prod_{i=1}^r
[\underline M_{\operatorname{Spec}(\widehat{\mathcal O}_{X_0,x_i})}
  (k[\varepsilon])]
\]
is exact. Equivalently,
\[
0\longrightarrow H^1(X_0,\underline D^0_{X_0})
\longrightarrow\operatorname{Hom}_{\Lambda\text{-alg}}(R,k[\varepsilon])
\longrightarrow
\operatorname{Hom}_{\Lambda\text{-alg}}
 (R_1\widehat\otimes_\Lambda\cdots\widehat\otimes_\Lambda R_r,
  k[\varepsilon])
\]
is exact. Consequently, $R$ is a formal power-series ring over
$R_1\widehat\otimes_\Lambda\cdots\widehat\otimes_\Lambda R_r$ in
$d$ variables, where
\[
d=\dim_k H^1(X_0,\underline D^0_{X_0}).
\]
Therefore
\begin{align*}
\dim\underline M_{X_0}
 &=\dim(k\otimes_\Lambda R) \\
 &=d+\dim\bigl(k\otimes_\Lambda
   (R_1\widehat\otimes_\Lambda\cdots\widehat\otimes_\Lambda R_r)
   \bigr) \\
 &=d+\dim\bigl((k\otimes_\Lambda R_1)
   \widehat\otimes_k\cdots\widehat\otimes_k
   (k\otimes_\Lambda R_r)\bigr) \\
 &=d+\sum_{i=1}^r\dim(k\otimes_\Lambda R_i) \\
 &=d+\sum_{i=1}^r
   \dim\underline M_{\operatorname{Spec}
   (\widehat{\mathcal O}_{X_0,x_i})}.
\end{align*}

%%%%%%%%%%  scan idx 121  |  folio 110  %%%%%%%%%%

\noindent\textbf{(b)} This follows from \SGARefLink{sga7.sga7i.evi.n4-10}{4.10} and \SGARefLink{sga7.sga7i.evi.n4-9}{4.9}(a).
\hfill$\square$


% END INCLUDED SGA7I_English_current_buildable_source_20260823/components/47_expose_VI_corollary_4_13.tex

% BEGIN INCLUDED SGA7I_English_current_buildable_source_20260823/components/48_expose_VI_4_14_and_4_15.tex
\medskip
\noindent\textbf{\SGARefTarget{sga7.sga7i.evi.n4-14}{4.14}.}
We close this section by constructing explicitly a formal versal
deformation in a simple case. When $X_0$ is affine and a complete
intersection, such a deformation is immediate to write down. Let
\[
A_0=k[x_1,\ldots,x_n]/(f_1,\ldots,f_r)=S/J,
\]
where $S=k[x_1,\ldots,x_n]$, $J=(f_1,\ldots,f_r)$, and
$f_1,\ldots,f_r$ is an $S$-regular sequence. We have an exact
sequence
\[
\cdots\longrightarrow\operatorname{Der}_k(S,A_0)
\longrightarrow\operatorname{Hom}_{A_0}(J/J^2,A_0)
\longrightarrow D^1(A_0|k,A_0)\longrightarrow0.
\]
The $A_0$-module $J/J^2$ is free of rank $r$, generated by the
classes of $f_1,\ldots,f_r$, because these elements form an
$S$-regular sequence.

Assume now that $D^1(A_0|k,A_0)$ is finite-dimensional over $k$;
for example, assume that $A_0$ has only isolated nonsmooth points.
Put
\[
d=\dim_kD^1(A_0|k,A_0),
\]
and choose maps
\[
p_i:J/J^2\longrightarrow A_0\qquad(1\leq i\leq d)
\]
whose classes form a $k$-basis of $D^1(A_0|k,A_0)$. For
$1\leq i\leq d$ and $1\leq j\leq r$, choose
$P_{ij}\in\Lambda[x_1,\ldots,x_n]$ such that
\[
P_{ij}\equiv p_i(f_j)\pmod{\mathfrak m_\Lambda+J},
\]
where $\mathfrak m_\Lambda$ is the maximal ideal of $\Lambda$. Set
\[
R=\Lambda[[t_1,\ldots,t_d]],
\qquad
B=R[x_1,\ldots,x_n]/(G_1,\ldots,G_r),
\]
where
\[
G_j=F_j-\sum_{i=1}^d t_iP_{ij},
\]
and $F_j\in\Lambda[x_1,\ldots,x_n]$ is obtained from $f_j$ by
lifting its coefficients from $k$ to $\Lambda$.

The algebra $R[x_1,\ldots,x_n]$ is flat over $R$, and the images of
$G_1,\ldots,G_r$ modulo $\mathfrak m_R$ form the regular sequence
$f_1,\ldots,f_r$ in
$k\otimes_RR[x_1,\ldots,x_n]$. Hence $B$ is $R$-flat and
$k\otimes_RB=A_0$. By construction, the map
\[
\operatorname{Hom}_{\Lambda\text{-alg}}(R,k[\varepsilon])
\longrightarrow D^1(A_0|k,A_0)
=[\underline M_{\operatorname{Spec}(A_0)}(k[\varepsilon])]
\]
is bijective. Thus, if

%%%%%%%%%%  scan idx 122  |  folio 111  %%%%%%%%%%

\[
\operatorname{Spf}(B)
=\varinjlim_\nu
\operatorname{Spec}\bigl((R/\mathfrak m_R^{\nu+1})\otimes_RB\bigr),
\]
then $(R,\operatorname{Spf}(B))$ is an object of
$\widehat{\underline M}_{\operatorname{Spec}(A_0)}$. The functor
\[
\underline h_R\longrightarrow\underline M_{\operatorname{Spec}(A_0)}
\]
that it induces is smooth by \SGARefLink{sga7.sga7i.evi.n1-6}{1.6}, since $R$ is a formal power-series
ring over $\Lambda$. In other words,
$(R,\operatorname{Spf}(B))$ is a formal versal object of
$\underline M_{\operatorname{Spec}(A_0)}$.

By \SGARefLink{sga7.sga7i.evi.n4-10}{4.10}(b), if instead we put
\[
B=R[[x_1,\ldots,x_n]]/(G_1,\ldots,G_r),
\]
then $(R,\operatorname{Spf}(B))$ is a formal versal deformation of
$\operatorname{Spec}(A_0)$ for
\[
A_0=k[[x_1,\ldots,x_n]]/(f_1,\ldots,f_r).
\]

\medskip
\noindent\textbf{Remark \SGARefTarget{sga7.sga7i.evi.n4-15}{4.15}.}
Let $X_0$ be a scheme of finite type over $k$ with only isolated
nonsmooth points, and let $(R,\mathfrak X)$ be a formal versal
deformation of $X_0$. The preceding construction shows that, when
$X_0$ is affine and locally a complete intersection, $\mathfrak X$
is the formalization of an $R$-scheme of finite type. The same fact
extends readily to the nonaffine case provided that
$H^2(X_0,\underline D^0_{X_0})=0$.


% END INCLUDED SGA7I_English_current_buildable_source_20260823/components/48_expose_VI_4_14_and_4_15.tex

% BEGIN INCLUDED SGA7I_English_current_buildable_source_20260823/components/49_expose_VI_remark_4_16.tex
\medskip
\noindent\textbf{Remark \SGARefTarget{sga7.sga7i.evi.n4-16}{4.16}.}
The global theory of the relative cotangent complex---here reduced to
the truncated complex of length one, which nevertheless suffices for
many purposes in deformation theory---generalizes and simplifies
results such as \SGARefLink{sga7.sga7i.evi.n4-10}{4.10}. Let $S$ be a scheme, and let
$S_0\subset S'\subset S$ be subschemes defined by ideals
$J\supset I$ with $JI=0$. Let $X'$ be flat over $S'$, and write
\[
L_\bullet=L^{X_0/S_0}
\]
for the relative cotangent complex of $X_0$ over $S_0$, regarded as
an object of $D(X_0,\mathcal O_{X_0})$. Let $\underline F$ be the
category of flat $S$-schemes $X$ extending $X'/S'$; thus an object is
equipped with an $S'$-isomorphism
\[
X\times_SS'\xrightarrow{\ \sim\ }X'.
\]
Then:

\smallskip
\noindent\textbf{(a)} For an object $X$ of $\underline F$, its group
of automorphisms inducing the identity on $X'$ is canonically
isomorphic to
\[
\operatorname{Ext}^0(X_0;L_\bullet,I_{X_0}),
\]
where $I_{X_0}$ is the inverse image on $X_0$ of the
$\mathcal O_{S_0}$-module $I$.

%%%%%%%%%%  scan idx 123  |  folio 112  %%%%%%%%%%

\smallskip
\noindent\textbf{(b)} The set of isomorphism classes of objects of
$\underline F$ is either empty or a torsor (a principal homogeneous
space) under
$\operatorname{Ext}^1(X_0;L_\bullet,I_{X_0})$.

\smallskip
\noindent\textbf{(c)} There is a canonical element
\[
c\in\operatorname{Ext}^2(X_0;L_\bullet,I_{X_0})
\]
whose vanishing is necessary and sufficient for $\underline F$ to be
nonempty.

To study these groups, it is often convenient to write
\[
\operatorname{Ext}^i(X_0;L_\bullet,I_{X_0})
=H^i\bigl(X_0,R\underline{\operatorname{Hom}}
  (L_\bullet,I_{X_0})\bigr).
\]
The cohomology sheaves of
$R\underline{\operatorname{Hom}}(L_\bullet,I_{X_0})$ are precisely
the sheaves $\underline D^i_{X_0}(I_{X_0})$ introduced in section 3.
For $i=0,1,2$, they may be viewed respectively as the sheaf of
automorphisms of an object $X$ of $\underline F$ (automorphisms that
induce the identity on $X'$), the sheaf measuring the indeterminacy
of isomorphism classes of local solutions to the flat-deformation
problem, and the sheaf of obstructions to the existence of local
solutions.

The standard spectral sequence is
\[
E_2^{p,q}
=H^p\bigl(X_0,\underline D^q_{X_0}(I_{X_0})\bigr)
\Longrightarrow
\operatorname{Ext}^{p+q}(X_0;L_\bullet,I_{X_0}).
\]
It gives the familiar isomorphisms
\begin{align*}
\operatorname{Aut}(X/X')
&\simeq H^0\bigl(X_0,\underline D^0_{X_0}(I_{X_0})\bigr)\\
&\simeq H^0\bigl(X_0,
  \underline{\operatorname{Hom}}
  (\Omega^1_{X_0/S_0},I_{X_0})\bigr)\\
&=\operatorname{Hom}(\Omega^1_{X_0/S_0},I_{X_0}),
\end{align*}
and the low-degree exact sequence
\[
0\longrightarrow H^1(X_0,\underline D^0_{X_0}(I_{X_0}))
\longrightarrow\operatorname{Ext}^1
\longrightarrow H^0(X_0,\underline D^1_{X_0}(I_{X_0}))
\longrightarrow H^2(X_0,\underline D^0_{X_0}(I_{X_0}))
\longrightarrow\operatorname{Ext}^2.
\]
In particular, if
$H^2(X_0,\underline D^0_{X_0}(I_{X_0}))=0$ and local flat
deformations of $X'/S'$ exist, then compatible local deformations
glue to a global flat deformation. More generally, for a prescribed
section of the sheaf of isomorphism classes of local flat
deformations, the argument of \SGARefLink{sga7.sga7i.evi.n4-10}{4.10} places the obstruction to a
compatible global deformation in
$H^2(X_0,\underline D^0_{X_0}(I_{X_0}))$.

%%%%%%%%%%  scan idx 124  |  folio 113  %%%%%%%%%%

The same spectral sequence shows that $\operatorname{Ext}^2$ vanishes
whenever
\begin{equation}\tag{*}
H^0(X_0,\underline D^2_{X_0})=0,
\qquad
H^1(X_0,\underline D^1_{X_0})=0,
\qquad
H^2(X_0,\underline D^0_{X_0})=0.
\end{equation}
The first equality always holds when $X_0$ is a relative complete
intersection over $S_0$, since then
$\underline D^2_{X_0}(M)=0$ for every $\mathcal O_{X_0}$-module $M$.
The second holds when the support of $\underline D^1_{X_0}$ is
discrete, as happens when $X_0$ is smooth over $S_0$ away from a
discrete subset. If all three equalities in $(*)$ hold, then
$\underline F$ is nonempty by (c), which generalizes \SGARefLink{sga7.sga7i.evi.n4-13}{4.13}(b).

There is also a geometric interpretation, requiring only the
globalized functors $\underline D^i_{X_0}$ rather than the global
theory of the relative cotangent complex. The three successive
obstructions to extending $X'/S'$ flatly over $S$ lie in the three
groups displayed in $(*)$. The local theory defines the first
obstruction
\[
c_1\in H^0(X_0,\underline D^2_{X_0})
\]
as the obstruction to the existence of local solutions. If $c_1=0$,
the sheaf $\underline F$ of isomorphism classes of local solutions is
canonically a torsor under $\underline D^1_{X_0}$. Its obstruction to
having a section is
\[
c_2\in H^1(X_0,\underline D^1_{X_0});
\]
such a section is a coherent system of isomorphism classes of local
solutions. If $c_2=0$, choose a section of $\underline F$. The
obstruction to a global solution realizing the chosen local classes
is then
\[
c_3\in H^2(X_0,\underline D^0_{X_0}).
\]
Changing the section of $\underline F$ by a section of
$\underline D^1_{X_0}$ gives a map
\[
d:H^0(X_0,\underline D^1_{X_0})
\longrightarrow H^2(X_0,\underline D^0_{X_0}).
\]
A solution exists if and only if $c_1=0$, $c_2=0$, and the class of
$c_3$ in $\operatorname{Coker}d$ is zero (each $c_i$ being defined
only after the preceding obstructions vanish).

%%%%%%%%%%  scan idx 125  |  folio 114  %%%%%%%%%%

The three obstruction groups
\[
H^0(X_0,\underline D^2_{X_0}),
\qquad
H^1(X_0,\underline D^1_{X_0}),
\qquad
H^2(X_0,\underline D^0_{X_0})/\operatorname{Im}d
\]
are the initial terms $E_2^{0,2}$, $E_2^{1,1}$, and $E_3^{2,0}$ of
the preceding spectral sequence; up to sign, $d$ is its differential
\[
d_2:E_2^{0,1}\longrightarrow E_2^{2,0}.
\]


% END INCLUDED SGA7I_English_current_buildable_source_20260823/components/49_expose_VI_remark_4_16.tex

% BEGIN INCLUDED SGA7I_English_current_buildable_source_20260823/components/50_expose_VI_section_5_opening_through_5_2.tex
\bigskip
\noindent\textbf{5. Formal Jacobian subschemes}

\medskip
Let $X_0$ be a scheme of finite type over a field $k$, and let
$(R,\mathfrak X)$ be a formal versal deformation of $X_0$. The space
$\mathfrak X$ is not an ordinary $R$-scheme, but an $R$-adic formal
scheme. The purpose of this section is to clarify the meaning of the
``nonsmooth locus of $\mathfrak X$ over $R$.''

We first recall the Fitting ideals of a finite-type module. Let $A$ be
a commutative ring and let $E$ be a finite-type $A$-module. Choose a
presentation
\[
K\longrightarrow F\longrightarrow E\longrightarrow0,
\]
where $F$ is locally free of constant rank $n$. For
$p=0,1,2,\ldots$, define
\[
I_A^p(E)=\operatorname{Im}\left(
 \bigwedge^{n-p}K\otimes_A\bigwedge^{n-p}F^*
 \longrightarrow A\right).
\]
These ideals are independent of the chosen presentation. They are
invariants of the $A$-module $E$, called its \emph{Fitting ideals}.
The following properties follow directly from the definition.

\medskip
\noindent\textbf{\SGARefTarget{sga7.sga7i.evi.n5-1}{5.1}.}
\textbf{(a)} For every ring homomorphism $A\longrightarrow B$, the
canonical map
\[
B\otimes_A A/I_A^p(E)
\xrightarrow{\ \sim\ }
B/I_B^p(B\otimes_AE)
\]
is an isomorphism for every $p$. In particular,
\[
I_{S^{-1}A}^p(S^{-1}E)=S^{-1}I_A^p(E)
\]
for every multiplicative subset $S\subset A$.

%%%%%%%%%%  scan idx 126  |  folio 115  %%%%%%%%%%

\smallskip
\noindent\textbf{(b)} There are inclusions
\[
I_A^0(E)\subset I_A^1(E)\subset I_A^2(E)\subset\cdots.
\]
For every $x\in\operatorname{Spec}(A)$,
\[
I_A^p(E)_x=A_x
\qquad\text{for all }p\geq\dim_{\kappa(x)}E(x),
\]
where $E(x)=\kappa(x)\otimes_AE$.

\smallskip
\noindent\textbf{(c)} A point
$x\in\operatorname{Spec}(A)-\operatorname{Supp}(A/I_A^p(E))$ if and
only if, for any (equivalently, for some) presentation
\[
K\longrightarrow F\longrightarrow E\longrightarrow0,
\qquad F\simeq A^n,
\]
the module $K_x$ contains a copy of $A_x^{n-p}$ that is a direct
summand of $F_x$. In particular,
\[
\operatorname{Supp}(A/I_A^0(E))=\operatorname{Supp}(E).
\]

\smallskip
\noindent\textbf{(d)} For any two finite-type $A$-modules,
\[
I_A^r(E_1\oplus E_2)
=\sum_{p+q=r}I_A^p(E_1)I_A^q(E_2).
\]

\smallskip
\noindent\textbf{(e)} If $A$ is a noetherian adic ring, then
\[
A/I_A^p(E)=\varprojlim_\nu A_\nu/I_{A_\nu}^p(E_\nu),
\]
where $A=\varprojlim_\nu A_\nu$ and $E_\nu=E\otimes_AA_\nu$.

Property \SGARefLink{sga7.sga7i.evi.n5-1}{5.1}(a) globalizes the definition. If $X$ is a scheme and
$\underline E$ is a quasi-coherent $\mathcal O_X$-module of finite
presentation, there are quasi-coherent ideal sheaves
$\underline I_X^p(\underline E)$, $p\geq0$, characterized on every
affine open $U\subset X$ by
\[
\Gamma(U,\underline I_X^p(\underline E))
\xrightarrow{\ \sim\ }
I_{\Gamma(U,\mathcal O_X)}^p(\Gamma(U,\underline E)).
\]

For a morphism of schemes $X\longrightarrow Y$, the relative module
of K\"ahler differentials
$\Omega_{X|Y}=\Omega^1_{X|Y}$ is quasi-coherent. If the morphism is
essentially of finite presentation, then $\Omega_{X|Y}$ is of finite
presentation, so the ideal sheaves
$\underline I_X^p(\Omega_{X|Y})$ are defined.

\medskip
\noindent\textbf{Definition \SGARefTarget{sga7.sga7i.evi.n5-2}{5.2}.}
Let $X\longrightarrow Y$ be a morphism essentially of finite
presentation. For every integer $p\geq0$, define $J_X^p(X|Y)$ to be
the closed subscheme of $X$

%%%%%%%%%%  scan idx 127  |  folio 116  %%%%%%%%%%

\noindent defined by the ideal sheaf
\[
\underline I_X^p(X|Y)=\underline I_X^p(\Omega_{X|Y}).
\]
It is called the $p$th \emph{Jacobian subscheme} of $X$ over $Y$.
For every $x\in X$, put
\[
X_x=\operatorname{Spec}(\mathcal O_{X,x}).
\]


% END INCLUDED SGA7I_English_current_buildable_source_20260823/components/50_expose_VI_section_5_opening_through_5_2.tex

% BEGIN INCLUDED SGA7I_English_current_buildable_source_20260823/components/51_expose_VI_theorem_5_3.tex
\medskip
\noindent\textbf{Theorem \SGARefTarget{sga7.sga7i.evi.n5-3}{5.3}.}
Let $X\longrightarrow Y$ be a morphism essentially of finite
presentation. Then:

\smallskip
\noindent\textbf{(1)} For every base change $Y'\longrightarrow Y$,
the canonical map
\[
J_{X'}^p(X'|Y')\xrightarrow{\ \sim\ }
J_X^p(X|Y)\times_YY'
\]
is an isomorphism for every $p$, where $X'=X\times_YY'$.

\smallskip
\noindent\textbf{(2)} For every $x\in X$, there is a canonical
isomorphism
\[
J_{X_x}^p(X_x|Y)\xrightarrow{\ \sim\ }J_X^p(X|Y)_x.
\]

\smallskip
\noindent\textbf{(3)} There are inclusions
\[
J_X^0(X|Y)\supset J_X^1(X|Y)\supset J_X^2(X|Y)\supset\cdots,
\]
and
\[
x\notin J_X^p(X|Y)
\qquad\text{for every }p\geq
\dim_{\kappa(x)}\Omega_{X|Y}(x).
\]

\smallskip
\noindent\textbf{(4)} A point $x$ does not belong to
$J_X^p(X|Y)$ if and only if it has an open neighbourhood $U$ that
admits a closed immersion $i:U\hookrightarrow U'$ into a $Y$-scheme
$U'$ such that $U'$ is smooth over $Y$ at $x'=i(x)$ and
\[
\dim_{x'}U'_{f(x)}=p.
\]

\smallskip
\noindent\textbf{(5)} For the projections
$\pi_i:X_1\times_YX_2\longrightarrow X_i$,
\[
\underline I_{X_1\times_YX_2}^r(X_1\times_YX_2|Y)
=\sum_{p+q=r}
 \pi_1^*\underline I_{X_1}^p(X_1|Y)\,
 \pi_2^*\underline I_{X_2}^q(X_2|Y).
\]

\smallskip
\noindent\textit{Proof.}
Assertion (3) follows from \SGARefLink{sga7.sga7i.evi.n5-1}{5.1}(b). Assertion (1) follows from \SGARefLink{sga7.sga7i.evi.n5-1}{5.1}(a)
and the canonical isomorphism
\[
P^*\Omega_{X|Y}\xrightarrow{\ \sim\ }\Omega_{X'|Y'},
\qquad P:X\times_YY'\longrightarrow X.
\]
Likewise, (2) follows from \SGARefLink{sga7.sga7i.evi.n5-1}{5.1}(a) and
\[
\mathcal O_{X,x}\otimes_{\mathcal O_X}\Omega_{X|Y}
\xrightarrow{\ \sim\ }\Omega_{X_x|Y}.
\]
Finally, (5) follows from \SGARefLink{sga7.sga7i.evi.n5-1}{5.1}(d) and
\[
\pi_1^*\Omega_{X_1|Y}\oplus\pi_2^*\Omega_{X_2|Y}
\xrightarrow{\ \sim\ }
\Omega_{X_1\times_YX_2|Y}.
\]

For (4), first suppose that there is a diagram
\[
\begin{tikzcd}[column sep=huge,row sep=large]
U \arrow[r,hook,"i"] \arrow[d,"f"']
  & U' \arrow[dl] \\
Y & {}
\end{tikzcd}
\]
where $U'$ is smooth over $Y$ at $x'=i(x)$ and
$\dim_{x'}U'_{f(x)}=p$. Then
\[
\dim_{\kappa(x')}\Omega_{U'|Y}(x')=p,
\]
and the sequence
\[
i^*\Omega_{U'|Y}\longrightarrow\Omega_{U|Y}\longrightarrow0
\]
is exact. Hence
$\dim_{\kappa(x)}\Omega_{U|Y}(x)\leq p$, so
$x\notin J_U^p(U|Y)$ by (3). Since
\[
J_U^p(U|Y)=J_X^p(X|Y)|_U,
\]
we have $x\notin J_X^p(X|Y)$.

%%%%%%%%%%  scan idx 128  |  folio 117  %%%%%%%%%%

Conversely, suppose that $x\notin J_X^p(X|Y)$. The question is
local, so write $X=\operatorname{Spec}(B)$ and
$Y=\operatorname{Spec}(A)$. Choose an $A$-presentation
\[
0\longrightarrow J\longrightarrow P\longrightarrow B
\longrightarrow0,
\qquad P=A[x_1,\ldots,x_n].
\]
There is an exact sequence
\[
J/J^2\xrightarrow{d}B\otimes_P\Omega_{P|A}
\longrightarrow\Omega_{B|A}\longrightarrow0,
\]
and $B\otimes_P\Omega_{P|A}$ is free of rank $n$ over $B$. Thus
\[
I_B^p(\Omega_{B|A})
=\operatorname{Im}\left(
 \bigwedge^{n-p}(J/J^2)\otimes_B
 (B\otimes_P\bigwedge^{n-p}\Omega_{P|A})^*
 \longrightarrow B\right).
\]
Consequently, $x\notin J_X^p(X|Y)$ precisely when there are
$f_1,\ldots,f_{n-p}\in J$ and a choice of $n-p$ among the variables
$x_1,\ldots,x_n$ for which the corresponding Jacobian minor
\[
\det\left(\frac{\partial f_i}{\partial x_j}\right)(x)
\]
is nonzero. Equivalently,
\[
\operatorname{Spec}\bigl(P/(f_1,\ldots,f_{n-p})\bigr)
\longrightarrow\operatorname{Spec}(A)
\]
is smooth at the image $x'$ of $x$. Taking
\[
U'=\operatorname{Spec}\bigl(P/(f_1,\ldots,f_{n-p})\bigr)
\]
gives the required closed immersion locally. \hfill$\square$


% END INCLUDED SGA7I_English_current_buildable_source_20260823/components/51_expose_VI_theorem_5_3.tex

% BEGIN INCLUDED SGA7I_English_current_buildable_source_20260823/components/52_expose_VI_theorem_5_4.tex
\medskip
\noindent\textbf{Theorem \SGARefTarget{sga7.sga7i.evi.n5-4}{5.4}.}
Let $X\xrightarrow{f}Y$ be a morphism essentially of finite
presentation. Suppose that there is a fixed integer $n\geq0$ for
which one of the following conditions holds:

\smallskip
\noindent\textbf{(1)} $f$ is flat and
$\dim_xX_{f(x)}=n$ for every $x\in X$;

\smallskip
\noindent\textbf{(2)} $f$ is a relative complete intersection of
virtual dimension $n$ (SGA 6, VIII, 1.9).

\smallskip
Then $x\notin J_X^n(X|Y)$ if and only if $f$ is smooth at $x$.

\smallskip
\noindent\textit{Proof.}
By \SGARefLink{sga7.sga7i.evi.n5-3}{5.3}(4), $x\notin J_X^n(X|Y)$ if and only if there is a diagram
\[
\begin{tikzcd}[column sep=huge,row sep=large]
U \arrow[r,hook,"i"] \arrow[d,"f"']
  & U' \arrow[dl,"f'"] \\
Y & {}
\end{tikzcd}
\]
such that $f'$ is smooth at $x'$, with
$\dim_{x'}U'_{f(x)}=n$, while $i$ is a closed immersion and
$x'=i(x)$. It therefore suffices to show that under either (1) or
(2), $i$ is an isomorphism at $x$.

Set
\[
A=\mathcal O_{Y,f(x)},\qquad
B'=\mathcal O_{U',x'},\qquad B=\mathcal O_{U,x},
\]
and let
\[
0\longrightarrow I\longrightarrow B'\longrightarrow B
\longrightarrow0
\]
be exact. We must show that $I=0$.

Assume (1). Since $B$ is flat over $A$, the sequence
\[
0\longrightarrow I/\mathfrak mI
\longrightarrow B'/\mathfrak mB'
\longrightarrow B/\mathfrak mB\longrightarrow0
\]
is exact, where $\mathfrak m$ is the maximal ideal of $A$. By
hypothesis,
\[
\dim_xU_{f(x)}=\dim_{x'}U'_{f(x)},
\]
and hence
$\dim(B'/\mathfrak mB')=\dim(B/\mathfrak mB)$. Since
$B'/\mathfrak mB'$ is a regular local ring, it follows that
$I/\mathfrak mI=0$, and therefore $I=0$.

%%%%%%%%%%  scan idx 129  |  folio 118  %%%%%%%%%%

Now assume (2), and consider the exact sequence
\[
I/I^2\xrightarrow{d}B\otimes_{B'}\Omega_{B'|A}
\longrightarrow\Omega_{B|A}\longrightarrow0.
\]
Since $X$ is a relative complete intersection over $Y$ and
$U'\longrightarrow Y$ is smooth at $x'$, the map
\[
d:I/I^2\longrightarrow B\otimes_{B'}\Omega_{B'|A}
\]
is injective at $x$, and both $I/I^2$ and
$B\otimes_{B'}\Omega_{B'|A}$ are free at $x$. By the definition of
the virtual dimension of a relative complete intersection, we must
therefore have
\[
n=\dim_{\kappa(x)}\Omega_{B'|A}(x')
  -\dim_{\kappa(x)}(I/I^2)(x)
 =n-\dim_{\kappa(x)}(I/I^2)(x).
\]
Thus $I/I^2=0$, and hence $I=0$. \hfill$\square$

% END INCLUDED SGA7I_English_current_buildable_source_20260823/components/52_expose_VI_theorem_5_4.tex

% BEGIN INCLUDED SGA7I_English_current_buildable_source_20260823/components/53_expose_VI_5_5_through_5_6.tex
\medskip
\noindent\textbf{Definition \SGARefTarget{sga7.sga7i.evi.n5-5}{5.5}.}
Let $X\xrightarrow{f}Y$ be a morphism essentially of finite
presentation. If there is a fixed integer $n$ satisfying (1) or (2)
of \SGARefLink{sga7.sga7i.evi.n5-4}{5.4}, set
\[
\underline I_X(X|Y)=\underline I_X^n(X|Y),
\qquad
J_X(X|Y)=J_X^n(X|Y).
\]

%%%%%%%%%%  scan idx 130  |  folio 119  %%%%%%%%%%

This notion extends naturally to formal schemes. Let $\mathfrak X$
be a noetherian formal scheme and $\mathcal E$ a coherent
$\mathfrak X$-module. For every affine open $U\subset\mathfrak X$,
define
\[
\underline I_{\mathfrak X}^p(\mathcal E)(U)
=I_{\mathfrak X(U)}^p\bigl(\mathcal E(U)\bigr).
\]
If $\mathfrak X=\varinjlim_\nu X_\nu$, then
$\mathcal E=\varprojlim_\nu\mathcal E_\nu$, where
\[
\mathcal E_\nu
=\mathcal E\otimes_{\mathcal O_{\mathfrak X}}\mathcal O_{X_\nu},
\]
and noetherianness gives
\[
\underline I_{\mathfrak X}^p(\mathcal E)
=\varprojlim_\nu\underline I_{X_\nu}^p(\mathcal E_\nu).
\]

Now let $S$ be a noetherian formal scheme and $\mathfrak X$ an
$S$-adic formal scheme essentially of finite type. Let
$\underline I_{\mathfrak X}$ be the ideal sheaf defining the closed
diagonal immersion
\[
\Delta:\mathfrak X\longrightarrow
\mathfrak X\times_S\mathfrak X
\]
(where the fibre product is taken in the category of formal schemes),
so that
\[
0\longrightarrow\underline I_{\mathfrak X}
\longrightarrow\mathcal O_{\mathfrak X\times_S\mathfrak X}
\longrightarrow\mathcal O_{\mathfrak X}\longrightarrow0
\]
is exact. Define
\[
\Omega_{\mathfrak X|S}
=\underline I_{\mathfrak X}/\underline I_{\mathfrak X}^{,2}.
\]
Writing
\[
\mathfrak X=\varinjlim_\nu X_\nu,
\qquad X_\nu=\mathfrak X\times_SS_\nu,
\]
we have
\[
\Omega_{\mathfrak X|S}
\simeq\varprojlim_\nu\Omega_{X_\nu|S_\nu},
\]
so $\Omega_{\mathfrak X|S}$ is a coherent $\mathfrak X$-module.
For every formal $S$-scheme $S'$ of finite type there is a canonical
isomorphism
\[
\Omega_{\mathfrak X\times_SS'|S'}
\simeq
\mathcal O_{S'}\otimes_{\mathcal O_S}\Omega_{\mathfrak X|S}.
\]

\medskip
\noindent\textbf{Definition \SGARefTarget{sga7.sga7i.evi.n5-6}{5.6}.}
Let $S$ be a noetherian formal scheme and $\mathfrak X$ an $S$-adic
formal scheme essentially of finite type. Define
\[
\underline I_{\mathfrak X}^p(\mathfrak X|S)
=\underline I_{\mathfrak X}^p(\Omega_{\mathfrak X|S}),
\]
and let $J_{\mathfrak X}^p(\mathfrak X|S)$ be the closed formal
subscheme of $\mathfrak X$ defined by this ideal sheaf.

Since
$\Omega_{\mathfrak X|S}=\varprojlim_\nu\Omega_{X_\nu|S_\nu}$,
we have
\[
\underline I_{\mathfrak X}^p(\mathfrak X|S)
=\varprojlim_\nu\underline I_{X_\nu}^p(X_\nu|S_\nu)
\]
and
\[
J_{\mathfrak X}^p(\mathfrak X|S)
=\varinjlim_\nu J_{X_\nu}^p(X_\nu|S_\nu).
\]

%%%%%%%%%%  scan idx 131  |  folio 120  %%%%%%%%%%

The formal scheme $\mathfrak X_{(x)}$ is again $S$-adic and
essentially of finite type, and there is a commutative diagram
\[
\begin{tikzcd}[column sep=huge,row sep=large]
\mathfrak X_{(\widehat{x})}
  \arrow[rr,"i_{(\widehat{x})}"] \arrow[dr]
  && \mathfrak X \\
& \mathfrak X_{(x)} \arrow[ur,"i_{(x)}"'] & {}
\end{tikzcd}
\]
% Direct 1800-dpi source check: the upper-left subscript and its arrow
% label are both \widehat{x}; the lower subscript and label are plain x.
If $x$ is an isolated point of $\mathfrak X$, then
$i_{(x)}$ and $i_{(\widehat{x})}$ are both isomorphisms onto the
connected component of $\mathfrak X$ containing $x$.

In particular, if
\[
|\mathfrak X|=\{x_1,\ldots,x_m\}
\]
is a finite discrete space, then there is a canonical identification
\[
\mathfrak X=\coprod_{1\leq i\leq m}\mathfrak X_{(x_i)}
\qquad\text{(disjoint union)}.
\]

% END INCLUDED SGA7I_English_current_buildable_source_20260823/components/53_expose_VI_5_5_through_5_6.tex

% BEGIN INCLUDED SGA7I_English_current_buildable_source_20260823/components/54_expose_VI_corollary_5_7.tex
\medskip
\noindent\textbf{Corollary \SGARefTarget{sga7.sga7i.evi.n5-7}{5.7}.}
Let $S$ be a noetherian formal scheme and $\mathfrak X$ an $S$-adic
formal scheme essentially of finite type. Then:

\smallskip
\noindent\textbf{(1)} Formation of
$J_{\mathfrak X}^p(\mathfrak X|S)$ commutes with base change. More
precisely, if
\[
\mathfrak X'=\mathfrak X\times_SS',
\]
then
\[
J_{\mathfrak X'}^p(\mathfrak X'|S')
=J_{\mathfrak X}^p(\mathfrak X|S)\times_SS'.
\]

\smallskip
\noindent\textbf{(2)} The connected components of
$J_{\mathfrak X}^p(\mathfrak X|S)$ are in bijective correspondence
with those of $J_{X_0}^p(X_0|S_0)$. For every
\[
x\in\bigl|J_{X_0}^p(X_0|S_0)\bigr|
\]
there is a canonical isomorphism
\[
J_{\mathfrak X}^p(\mathfrak X|S)_{(x)}
\simeq
J_{\mathfrak X_{(x)}}^p(\mathfrak X_{(x)}|S).
\]
In particular, if
\[
\bigl|J_{X_0}^p(X_0|S_0)\bigr|=\{x_1,\ldots,x_m\}
\]
is discrete, then
\[
J_{\mathfrak X}^p(\mathfrak X|S)
=\coprod_{1\leq i\leq m}
 J_{\mathfrak X_{(x_i)}}^p(\mathfrak X_{(x_i)}|S)
\qquad\text{(disjoint union)}.
\]

%%%%%%%%%%  scan idx 132  |  folio 121  %%%%%%%%%%

\smallskip
\noindent\textbf{(3)} Let $S$ be an ordinary noetherian scheme and
$X$ an $S$-scheme essentially of finite type. Given a closed
subscheme $S_0\subset S$, let $\widehat S$ be the formal completion
of $S$ along $S_0$. If
\[
\widehat X=X\times_S\widehat S
\]
is the formal completion of $X$ along
$X_0=X\times_SS_0$, then
\[
J_{\widehat X}^p(\widehat X|\widehat S)
=J_X^p(X|S)\times_S\widehat S.
\]
Equivalently, it is the formal completion of $J_X^p(X|S)$ along
$J_{X_0}^p(X_0|S_0)$.

\smallskip
\noindent\textbf{(4)} Suppose that a fixed integer $n\geq0$
satisfies condition \SGARefLink{sga7.sga7i.evi.n5-4}{5.4}(1) or \SGARefLink{sga7.sga7i.evi.n5-4}{5.4}(2) for every
$X_\nu\longrightarrow S_\nu$. Then
\[
x\notin J_{\mathfrak X}^n(\mathfrak X|S)
\]
if and only if $\mathfrak X$ is formally smooth over $S$ at $x$.

\smallskip
\noindent\textit{Proof.}
Assertion (1) follows from
\[
\Omega_{\mathfrak X'|S'}
\simeq
\mathcal O_{S'}\widehat\otimes_{\mathcal O_S}
\Omega_{\mathfrak X|S},
\]
whereas (2) follows from
\[
\mathcal O_{\mathfrak X_{(x)}}
\widehat\otimes_{\mathcal O_{\mathfrak X}}
\Omega_{\mathfrak X|S}
\simeq\Omega_{\mathfrak X_{(x)}|S}.
\]
For (3),
\[
\begin{aligned}
J_{\widehat X}^p(\widehat X|\widehat S)
&=\varinjlim_\nu
 J_{X\times_SS_\nu}^p(X\times_SS_\nu|S_\nu) \\
&=\varinjlim_\nu\bigl(J_X^p(X|S)\times_SS_\nu\bigr)
 =J_X^p(X|S)\times_S\widehat S.
\end{aligned}
\]
Finally, since $\mathfrak X$ is $S$-adic,
\[
\begin{aligned}
\mathfrak X\to S\text{ is formally smooth at }x
&\Longleftrightarrow
 X_\nu\to S_\nu\text{ is formally smooth at }x
 \text{ for every }\nu \\
&\Longleftrightarrow
 x\notin J_{X_\nu}^n(X_\nu|S_\nu)
 \text{ for every }\nu \\
&\Longleftrightarrow
 x\notin\varinjlim_\nu J_{X_\nu}^n(X_\nu|S_\nu)
 =J_{\mathfrak X}^n(\mathfrak X|S).
\end{aligned}
\]
\hfill$\square$

% END INCLUDED SGA7I_English_current_buildable_source_20260823/components/54_expose_VI_corollary_5_7.tex

% BEGIN INCLUDED SGA7I_English_current_buildable_source_20260823/components/55_expose_VI_remark_5_8.tex
\medskip
\noindent\textbf{Remark \SGARefTarget{sga7.sga7i.evi.n5-8}{5.8}.}
Let $\mathfrak X$ be an $S$-adic formal scheme essentially of finite
type. For a point $x\in\mathfrak X$, consider also the $S$-adic
formal scheme
\[
\mathfrak X_{(\widehat{x})}
=\varinjlim_\nu
 \operatorname{Spec}\bigl(\widehat{\mathcal O_{X_\nu,x}}\bigr).
\]
There are natural morphisms
\[
\mathfrak X_{(\widehat{x})}
\longrightarrow\mathfrak X_{(x)}
\longrightarrow\mathfrak X.
\]
In general, $\widehat{\mathcal O_{X_0,x}}$ is not essentially of
finite type over $S_{(y)}$, where $y$ is the image of $x$ under
$\mathfrak X\longrightarrow S$.

For simplicity, put
\[
A_\nu=\mathcal O_{X_\nu,x},
\qquad R_\nu=\mathcal O_{S_\nu,y}.
\]
The natural map
\[
\Omega_{A_\nu|R_\nu}
\longrightarrow\widehat\Omega_{A_\nu|R_\nu}
\]
induces an isomorphism
\[
\widehat{\Omega_{A_\nu|R_\nu}}
\xrightarrow{\ \sim\ }
\widehat{\widehat\Omega_{A_\nu|R_\nu}},
\]
and therefore
\[
\varprojlim_\nu\widehat{\Omega_{A_\nu|R_\nu}}
\xrightarrow{\ \sim\ }
\varprojlim_\nu
 \widehat{\widehat\Omega_{A_\nu|R_\nu}}.
\]
Here the outer hat denotes the separated completion of a module for
the topology induced by the ring. Consequently, if we define
\[
\widehat\Omega_{\mathfrak X_{(\widehat{x})}|S}
=\varprojlim_\nu
 \widehat{\widehat\Omega_{
   \mathcal O_{X_\nu,x}|\mathcal O_{S_\nu,y}}},
\]
then $\widehat\Omega_{\mathfrak X_{(\widehat{x})}|S}$ is a coherent
$\mathfrak X_{(\widehat{x})}$-module.

%%%%%%%%%%  scan idx 133  |  folio 122  %%%%%%%%%%

Define
$J_{\mathfrak X_{(\widehat{x})}}^p
 (\mathfrak X_{(\widehat{x})}|S)$
to be the closed formal subscheme of
$\mathfrak X_{(\widehat{x})}$ defined by
\[
\underline I_{\mathfrak X_{(\widehat{x})}}^p
 (\mathfrak X_{(\widehat{x})}|S)
=\underline I_{\mathfrak X_{(\widehat{x})}}^p
 (\widehat\Omega_{\mathfrak X_{(\widehat{x})}|S}).
\]
Then
\[
J_{\mathfrak X_{(x)}}^p(\mathfrak X_{(x)}|S)_{(\widehat{x})}
\simeq
J_{\mathfrak X_{(\widehat{x})}}^p
 (\mathfrak X_{(\widehat{x})}|S).
\]
If $x$ is an isolated point of
$J_{\mathfrak X}^p(\mathfrak X|S)$, the natural morphism
\[
J_{\mathfrak X_{(\widehat{x})}}^p
 (\mathfrak X_{(\widehat{x})}|S)
\longrightarrow
J_{\mathfrak X_{(x)}}^p(\mathfrak X_{(x)}|S)
\]
is an isomorphism.

To form the ``image'' of $J_{\mathfrak X}^p(\mathfrak X|S)$ under
$f:\mathfrak X\longrightarrow S$, one could simply take the formal
subscheme of $S$ defined by the ideal sheaf
\[
\ker\left(
 \mathcal O_S\longrightarrow
 \mathcal O_{\mathfrak X}/
 \underline I_{\mathfrak X}^p(\mathfrak X|S)
\right).
\]
For our purposes, however, it is more convenient to use the Fitting
ideal once again.

% END INCLUDED SGA7I_English_current_buildable_source_20260823/components/55_expose_VI_remark_5_8.tex

% BEGIN INCLUDED SGA7I_English_current_buildable_source_20260823/components/56_expose_VI_5_9_through_5_11.tex
\medskip
\noindent\textbf{Definition \SGARefTarget{sga7.sga7i.evi.n5-9}{5.9}.}
Let $S$ be a noetherian formal scheme and
$f:\mathfrak X\longrightarrow S$ an $S$-adic formal scheme
essentially of finite type. Suppose that
$J_{\mathfrak X}^p(\mathfrak X|S)$ is proper over $S$, so that
\[
f_*\mathcal O_{J_{\mathfrak X}^p(\mathfrak X|S)}
\]
is a coherent $\mathcal O_S$-module. Define
\[
\underline I_S^p(\mathfrak X|S)
=\underline I_S^0\!\left(
 f_*\mathcal O_{J_{\mathfrak X}^p(\mathfrak X|S)}
\right),
\]
and let $J_S^p(\mathfrak X|S)$ be the closed formal subscheme of $S$
defined by $\underline I_S^p(\mathfrak X|S)$.

If there is a fixed integer $n$ satisfying condition \SGARefLink{sga7.sga7i.evi.n5-4}{5.4}(1) or
\SGARefLink{sga7.sga7i.evi.n5-4}{5.4}(2) for every $X_\nu\longrightarrow S_\nu$, set
\[
\begin{aligned}
\underline I_{\mathfrak X}(\mathfrak X|S)
 &=\underline I_{\mathfrak X}^n(\mathfrak X|S),
&\qquad
J_{\mathfrak X}(\mathfrak X|S)
 &=J_{\mathfrak X}^n(\mathfrak X|S),\\
\underline I_S(\mathfrak X|S)
 &=\underline I_S^n(\mathfrak X|S),
&
J_S(\mathfrak X|S)
 &=J_S^n(\mathfrak X|S).
\end{aligned}
\]

\medskip
\noindent\textbf{Corollary \SGARefTarget{sga7.sga7i.evi.n5-10}{5.10}.}
Let $S$ be a noetherian formal scheme and
$f:\mathfrak X\longrightarrow S$ an $S$-adic formal scheme
essentially of finite type such that
$J_{\mathfrak X}^p(\mathfrak X|S)$ is proper over $S$. Then:

\smallskip
\noindent\textbf{(1)}
\[
\bigl|J_S^p(\mathfrak X|S)\bigr|
=f\bigl(\,|J_{\mathfrak X}^p(\mathfrak X|S)|\,\bigr).
\]

%%%%%%%%%%  scan idx 134  |  folio 123  %%%%%%%%%%

\smallskip
\noindent\textbf{(2)} Formation of $J_S^p(\mathfrak X|S)$ commutes
with base change: if $\mathfrak X'=\mathfrak X\times_SS'$, then
\[
J_{S'}^p(\mathfrak X'|S')
=J_S^p(\mathfrak X|S)\times_SS'.
\]

\smallskip
\noindent\textbf{(3)} If
\[
J_{\mathfrak X}^p(\mathfrak X|S)
=\coprod_{1\leq i\leq m}\mathfrak Y^{(i)}
\]
is a disjoint union, then
\[
\underline I_S^p(\mathfrak X|S)
=\prod_{1\leq i\leq m}
 \underline I_S^0\bigl(f_*\mathcal O_{\mathfrak Y^{(i)}}\bigr).
\]
In particular, if
\[
\bigl|J_{X_0}^p(X_0|S_0)\bigr|=\{x_1,\ldots,x_m\}
\]
is discrete, then
\[
\underline I_S^p(\mathfrak X|S)
=\prod_{1\leq i\leq m}
 \underline I_S^p(\mathfrak X_{(x_i)}|S)
=\prod_{1\leq i\leq m}
 \underline I_S^p(\mathfrak X_{(\widehat{x_i})}|S).
\]

\smallskip
\noindent\textbf{(4)} Let $S$ be an ordinary noetherian scheme,
$X$ an $S$-scheme essentially of finite type, and
$S_0\hookrightarrow S$ a closed subscheme. Let $\widehat S$ be the
formal completion of $S$ along $S_0$. If
\[
\widehat X=X\times_S\widehat S
\]
is the formal completion of $X$ along
$X_0=X\times_SS_0$, then
\[
J_{\widehat S}^p(\widehat X|\widehat S)
=J_S^p(X|S)\times_S\widehat S.
\]
Thus it is the formal completion of $J_S^p(X|S)$ along
$J_{S_0}^p(X_0|S_0)$.

\smallskip
\noindent\textit{Proof.}
By \SGARefLink{sga7.sga7i.evi.n5-1}{5.1}(c),
\[
\operatorname{Supp}
 f_*\mathcal O_{J_{\mathfrak X}^p(\mathfrak X|S)}
=\bigl|J_S^p(\mathfrak X|S)\bigr|,
\]
which proves (1). Assertions (2) and (4) follow from \SGARefLink{sga7.sga7i.evi.n5-1}{5.1}(a), while
(3) follows from \SGARefLink{sga7.sga7i.evi.n5-1}{5.1}(d) and, for an isolated point
$x\in J_{\mathfrak X}^p(\mathfrak X|S)$, the isomorphism
\[
J_{\mathfrak X_{(\widehat{x})}}^p
 (\mathfrak X_{(\widehat{x})}|S)
\xrightarrow{\ \sim\ }
J_{\mathfrak X_{(x)}}^p(\mathfrak X_{(x)}|S).
\]
\hfill$\square$

\medskip
\noindent\textbf{Remark \SGARefTarget{sga7.sga7i.evi.n5-11}{5.11}.}
Our definition of the Jacobian ideal, adopted from the classical
definition over a base field, commutes with base change and, under the
mild hypothesis of \SGARefLink{sga7.sga7i.evi.n5-4}{5.4}, gives a smoothness criterion. In a general
setting, however, one can probably define in many ways a canonical
subscheme whose support is the nonsmooth locus. We outline one such
definition.

For an $A$-algebra $B$, put
\[
I_{B|A}=\bigcap_E\operatorname{Ann}D^1(B|A,E),
\]
where $E$ ranges over all finite-type $B$-modules. By \SGARefLink{sga7.sga7i.evi.n3-13}{3.13}(c), the
ideal $I_{B|A}$ commutes with localization. Thus a morphism
$f:X\longrightarrow Y$ determines an ideal sheaf
$\underline I_{X|Y}$. If $f$ is essentially of finite type, then by
\SGARefLink{sga7.sga7i.evi.n3-13}{3.13}(a), $f$ is smooth at $x$ if and only if
\[
x\notin\operatorname{Supp}
 \bigl(\mathcal O_X/\underline I_{X|Y}\bigr).
\]

%%%%%%%%%%  scan idx 135  |  folio 124  %%%%%%%%%%

This construction carries over naturally to adic formal schemes
essentially of finite type, and is therefore equally suitable for
deformation theory.

% END INCLUDED SGA7I_English_current_buildable_source_20260823/components/56_expose_VI_5_9_through_5_11.tex

% BEGIN INCLUDED SGA7I_English_current_buildable_source_20260823/components/57_expose_VI_section_6_opening_and_lemma_6_1.tex
\par\bigskip
\noindent\textbf{\SGARefAliasTarget{sga7.sga7i.evi.n6-0}6. Non-degenerate quadratic singularities}

\medskip
The purpose of this section is to study deformations of a $k$-scheme
whose only singularities are non-degenerate quadratic singularities.
In view of \SGARefLink{sga7.sga7i.evi.n4-13}{4.13}, this becomes, in a suitable situation, a purely local
problem.

We return to the notation of \S4. Thus $\Lambda$ is a fixed complete
noetherian local ring with residue field $k$, and
$\underline C_\Lambda$ denotes the category of local artinian
$\Lambda$-algebras having residue field $k$.

Let $X$ be a $k$-scheme of finite type, and let $x$ be a closed point
of $X$. Recall that $x$ is called a \emph{non-degenerate quadratic
singularity rational over $k$} if
\[
\widehat{\mathcal O}_{X,x}
 \simeq k[[x_1,\ldots,x_n]]/(f)
\]
and
\[
\left(\frac{\partial f}{\partial x_1},\ldots,
      \frac{\partial f}{\partial x_n}\right)
=(x_1,\ldots,x_n).
\]
One may always take $f$ to be a quadratic form. Indeed:

\medskip
\noindent\textbf{Lemma \SGARefTarget{sga7.sga7i.evi.n6-1}{6.1}.}
Let $f\in k[[x_1,\ldots,x_n]]$, and put
$M=(x_1,\ldots,x_n)$. If
\[
\left(\frac{\partial f}{\partial x_1},\ldots,
      \frac{\partial f}{\partial x_n}\right)=M,
\]
then there are elements $x_1',\ldots,x_n'$ satisfying
$x_i'\equiv x_i\pmod{M^2}$ such that
\[
f=q(x_1',\ldots,x_n')
\]
for a quadratic form $q$.

\smallskip
\noindent\textit{Proof.}
We use successive approximation. Write
\[
f=f_2+h_3,
\]
where $f_2$ is quadratic and $h_3\in M^3$. Since the partial
derivatives of $f$ generate $M$, there are
$\varepsilon_1^{(1)},\ldots,\varepsilon_n^{(1)}\in M^2$ such that
\[
\sum_i\varepsilon_i^{(1)}\frac{\partial f}{\partial x_i}
 \equiv h_3\pmod{M^4}.
\]
Put $x_i^{(1)}=x_i-\varepsilon_i^{(1)}$. Then
\[
\begin{aligned}
f(x_1^{(1)},\ldots,x_n^{(1)})
&=f(x_1,\ldots,x_n)
 -\sum_i\varepsilon_i^{(1)}\frac{\partial f}{\partial x_i}
 +\cdots\\
&\equiv f_2(x_1,\ldots,x_n)\pmod{M^4}.
\end{aligned}
\]
Thus
\[
f(x_1^{(1)},\ldots,x_n^{(1)})
=f_2(x_1,\ldots,x_n)+h_4,
\qquad h_4\in M^4.
\]
Moreover, the derivatives with respect to the new variables still
generate $M$. Hence one can choose
$\varepsilon_1^{(2)},\ldots,\varepsilon_n^{(2)}\in M^3$ such that
\[
\sum_i\varepsilon_i^{(2)}
 \frac{\partial f}{\partial x_i^{(1)}}
 \equiv h_4\pmod{M^5}.
\]

%%%%%%%%%%  scan idx 136  |  folio 125  %%%%%%%%%%

It follows that
\[
\begin{aligned}
&f(x_1^{(1)}-\varepsilon_1^{(2)},\ldots,
  x_n^{(1)}-\varepsilon_n^{(2)})\\
&\qquad =f(x_1^{(1)},\ldots,x_n^{(1)})
 -\sum_i\varepsilon_i^{(2)}
       \frac{\partial f}{\partial x_i^{(1)}}+\cdots\\
&\qquad \equiv f_2(x_1,\ldots,x_n)\pmod{M^5}.
\end{aligned}
\]
Continuing in this way, one finds elements
$\varepsilon_i^{(\nu)}\in M^{\nu+1}$ such that, on setting
\[
x_i^{(\nu)}
=x_i-\bigl(\varepsilon_i^{(1)}+\cdots+
            \varepsilon_i^{(\nu)}\bigr),
\]
one has
\[
f(x_1^{(\nu)},\ldots,x_n^{(\nu)})
\equiv f_2(x_1,\ldots,x_n)\pmod{M^{\nu+3}}.
\]
Let $\varepsilon_i=\sum_\nu\varepsilon_i^{(\nu)}$. Completeness then
gives
\[
f(x_1-\varepsilon_1,\ldots,x_n-\varepsilon_n)
=f_2(x_1,\ldots,x_n),
\]
or, equivalently,
\[
f(x_1,\ldots,x_n)
=f_2(x_1+\varepsilon_1,\ldots,x_n+\varepsilon_n).
\]
Taking $x_i'=x_i+\varepsilon_i$ proves the assertion.
\hfill$\square$

% END INCLUDED SGA7I_English_current_buildable_source_20260823/components/57_expose_VI_section_6_opening_and_lemma_6_1.tex

% BEGIN INCLUDED SGA7I_English_current_buildable_source_20260823/components/58_expose_VI_lemma_6_2.tex
\par\medskip
\noindent\textbf{Lemma \SGARefTarget{sga7.sga7i.evi.n6-2}{6.2}.}
Let $A_0$ be a local $k$-algebra essentially of finite type such that
\[
\widehat A_0\simeq k[[x_1,\ldots,x_n]]/(f),
\]
where $f$ is a non-degenerate quadratic form, and let
$(R,\mathfrak X)$ be a formal versal deformation of
$X_0=\operatorname{Spec}(A_0)$. Then:

\smallskip
\noindent\textbf{(1)} The canonical maps
\[
J_{\mathfrak X}(\mathfrak X|R)
 \longrightarrow J_R(\mathfrak X|R)
 \longrightarrow \operatorname{Spec}(\Lambda)
\]
are both isomorphisms.

\smallskip
\noindent\textbf{(2)} The ideal
$I_R(\mathfrak X|R)$ is principal and generated by a free formal
generator of $R$ over $\Lambda$; in other words,
\[
I_R(\mathfrak X|R)=(t),
\qquad R=\Lambda[[t]].
\]

\smallskip
\noindent\textbf{(3)} The morphism
$\mathfrak X\longrightarrow\operatorname{Spec}(\Lambda)$ is formally
smooth.

\smallskip
\noindent\textit{Proof.}
Write $\mathfrak X=\operatorname{Spf}(A)$. By \SGARefLink{sga7.sga7i.evi.n4-10}{4.10}(b),
$(R,\operatorname{Spf}(\widehat A))$ is a formal versal deformation
of $\widehat X_0=\operatorname{Spec}(\widehat A_0)$, where
\[
\widehat A=\varprojlim_\nu A\otimes_RR_\nu.
\]
Now
\[
\widehat A_0\simeq k[[x_1,\ldots,x_n]]/(f),
\qquad
\left(\frac{\partial f}{\partial x_1},\ldots,
      \frac{\partial f}{\partial x_n}\right)
=(x_1,\ldots,x_n).
\]
It follows immediately that
\[
\bigl[\underline M_{\widehat X_0}(k[\varepsilon])\bigr]
\simeq k.
\]
Consequently, by \SGARefLink{sga7.sga7i.evi.n4-14}{4.14},
\[
R\simeq\Lambda[[t]],
\qquad
\widehat A\simeq R[[x_1,\ldots,x_n]]/(F-t),
\]
where $F$ is the quadratic polynomial in
$\Lambda[x_1,\ldots,x_n]\subset R[[x_1,\ldots,x_n]]$ obtained from
$f$ by lifting its coefficients to $\Lambda$.

The sequence
\[
0\longrightarrow\widehat A
 \xrightarrow{\ dF\ }
 \widehat A\,dx_1\oplus\cdots\oplus\widehat A\,dx_n
 \longrightarrow\Omega_{\widehat A|R}\longrightarrow0
\]
is exact, where
\[
1\longmapsto
 \frac{\partial F}{\partial x_1}dx_1+\cdots+
 \frac{\partial F}{\partial x_n}dx_n.
\]

%%%%%%%%%%  scan idx 137  |  folio 126  %%%%%%%%%%

Therefore
\[
I_{\widehat A}(\widehat A|R)
=\left(\frac{\partial F}{\partial x_1},\ldots,
       \frac{\partial F}{\partial x_n}\right)
=(x_1,\ldots,x_n),
\]
and hence
\[
\begin{aligned}
\widehat A/I_{\widehat A}(\widehat A|R)
&=\widehat A/(x_1,\ldots,x_n)\\
&=R[[x_1,\ldots,x_n]]/(x_1,\ldots,x_n,F-t)
 \simeq\Lambda.
\end{aligned}
\]
Thus the canonical map
$\Lambda\longrightarrow
\widehat A/I_{\widehat A}(\widehat A|R)$ is an isomorphism. The
completion comparison identifies this quotient with
$A/I_A(A|R)$, and consequently
\[
\Lambda\xrightarrow{\ \sim\ }R/I_R(A|R).
\]
This proves (1) and (2). Finally, (3) follows from the isomorphisms
\[
\begin{aligned}
\widehat A
&\simeq R[[x_1,\ldots,x_n]]/(F-t)\\
&\simeq\Lambda[[t,x_1,\ldots,x_n]]/(F-t)
 \simeq\Lambda[[x_1,\ldots,x_n]]
\end{aligned}
\]
of $\Lambda$-algebras.
\hfill$\square$

% END INCLUDED SGA7I_English_current_buildable_source_20260823/components/58_expose_VI_lemma_6_2.tex

% BEGIN INCLUDED SGA7I_English_current_buildable_source_20260823/components/59_expose_VI_corollary_6_3.tex
\par\medskip
\noindent\textbf{Corollary \SGARefTarget{sga7.sga7i.evi.n6-3}{6.3}.}
Let $X_0$ be a $k$-scheme of finite type, and let $x\in X_0$ be a
non-degenerate quadratic singularity rational over $k$.

\smallskip
\noindent\textbf{(1)} Let $(R,\mathfrak X)$ and
$(R,\mathfrak Y)$ be two deformations of $X_0$, with
$R\in\operatorname{ob}(\widehat{\underline C}_\Lambda)$. Assume that
$k=k^2$. Then
\[
\mathfrak X_{(\widehat x)}\simeq
\mathfrak Y_{(\widehat x)}
\]
as $R$-adic formal schemes if and only if
\[
I_R(\mathfrak X_{(x)}|R)=I_R(\mathfrak Y_{(x)}|R).
\]

\smallskip
\noindent\textbf{(2)} Let
$R\in\operatorname{ob}(\widehat{\underline C}_\Lambda)$ be a regular
local ring, and let $(R,\mathfrak X)$ be a deformation of $X_0$.
Then $\mathcal O_{\mathfrak X,x}$ is regular if and only if
\[
I_R(\mathfrak X_{(x)}|R)\not\subset\mathfrak m_R^2;
\]
equivalently, $I_R(\mathfrak X_{(x)}|R)$ is a regular divisor in
$\operatorname{Spf}(R)$.

\smallskip
\noindent\textit{Proof.}
Let $(R_0,Z)$ be a formal versal deformation of
$\operatorname{Spec}(\widehat{\mathcal O}_{X_0,x})$. Since $x$ is a
non-degenerate quadratic singularity rational over $k$, Lemma \SGARefLink{sga7.sga7i.evi.n6-2}{6.2}
allows us to take
\[
R_0\simeq\Lambda[[t]],
\qquad
Z\simeq\operatorname{Spf}
 \bigl(R_0[[X_1,\ldots,X_n]]/(F-t)\bigr).
\]

For (1), the deformations
$(R,\mathfrak X_{(\widehat x)})$ and
$(R,\mathfrak Y_{(\widehat x)})$ are induced from $(R_0,Z)$ by
continuous local $\Lambda$-algebra homomorphisms
\[
\varphi,\psi:R_0\longrightarrow R.
\]
Put $a=\varphi(t)$ and $b=\psi(t)$. Then
\[
\begin{aligned}
\mathfrak X_{(\widehat x)}
&\simeq\operatorname{Spf}
 \bigl(R[[X_1,\ldots,X_n]]/(F-a)\bigr),\\
\mathfrak Y_{(\widehat x)}
&\simeq\operatorname{Spf}
 \bigl(R[[X_1,\ldots,X_n]]/(F-b)\bigr).
\end{aligned}
\]

%%%%%%%%%%  scan idx 138  |  folio 127  %%%%%%%%%%

Suppose that
$I_R(\mathfrak X_{(x)}|R)=I_R(\mathfrak Y_{(x)}|R)$; equivalently,
$(a)=(b)$ in $R$. Then $a=ub$ for some unit $u\in R$. Since
$k=k^2$ and $R$ is complete, $u=v^2$ for some unit $v\in R$, and
hence $a=v^2b$. Thus
\[
(F-a)\longmapsto v^2(F-b)
\]
under $X_i\mapsto vX_i$, and this substitution induces an
$R$-algebra isomorphism
\[
R[[X_1,\ldots,X_n]]/(F-a)
 \xrightarrow{\ \sim\ }
R[[X_1,\ldots,X_n]]/(F-b),
\]
because $F$ is quadratic. The converse is immediate.

For (2), there is a continuous local $\Lambda$-algebra homomorphism
\[
\varphi:R_0=\Lambda[[t]]\longrightarrow R
\]
such that
\[
\mathfrak X_{(\widehat x)}
\simeq\operatorname{Spf}
 \bigl(R[[X_1,\ldots,X_n]]/(F-\varphi(t))\bigr).
\]
Moreover,
\[
I_R(\mathfrak X_{(x)}|R)
=I_R(\mathfrak X_{(\widehat x)}|R)
=(\varphi(t)),
\]
and
\[
R[[X_1,\ldots,X_n]]/(F-\varphi(t))
=\widehat{\mathcal O}_{\mathfrak X,x}
\]
is regular if and only if $\mathcal O_{\mathfrak X,x}$ is regular.
Since $R$ is regular,
$R[[X_1,\ldots,X_n]]$ is a regular local ring. If $M$ denotes its
maximal ideal, then its quotient by $F-\varphi(t)$ is regular if and
only if
\[
F-\varphi(t)\notin M^2,
\]
which is equivalent to
$\varphi(t)\notin\mathfrak m_R^2$, and hence to
\[
I_R(\mathfrak X_{(x)}|R)\not\subset\mathfrak m_R^2.
\]
\hfill$\square$

% END INCLUDED SGA7I_English_current_buildable_source_20260823/components/59_expose_VI_corollary_6_3.tex

% BEGIN INCLUDED SGA7I_English_current_buildable_source_20260823/components/60_expose_VI_theorem_6_4_through_part_4.tex
\par\medskip
\noindent\textbf{Theorem \SGARefTarget{sga7.sga7i.evi.n6-4}{6.4}.}
Let $X_0$ be a $k$-scheme of finite type whose nonsmooth points are
isolated, and let $(R,\mathfrak X)$ be a formal versal deformation of
$X_0$. Denote the set of nonsmooth points of $X_0$ by
$\{z_1,\ldots,z_m\}$. Assume that
\[
H^2(X_0,\underline D^0_{X_0})=0
\]
and that every $z_i$ is a non-degenerate quadratic singularity
rational over $k$. Then:

\smallskip
\noindent\textbf{(1)} $R$ is formally smooth over $\Lambda$; that is,
$R$ is a formal power-series ring over $\Lambda$.

\smallskip
\noindent\textbf{(2)} There is a disjoint-union decomposition
\[
J_{\mathfrak X}(\mathfrak X|R)
=\coprod_{1\leq i\leq m}
 J_{\mathfrak X_{(z_i)}}(\mathfrak X_{(z_i)}|R),
\]
and the canonical morphism
\[
J_{\mathfrak X_{(z_i)}}(\mathfrak X_{(z_i)}|R)
 \xrightarrow{\ \sim\ }
J_R(\mathfrak X_{(z_i)}|R)
\]
is an isomorphism.

\smallskip
\noindent\textbf{(3)} Each ideal
$I_R(\mathfrak X_{(z_i)}|R)$ is principal, as is
\[
I_R(\mathfrak X|R)
=\prod_{1\leq i\leq m}I_R(\mathfrak X_{(z_i)}|R).
\]
If $I_R(\mathfrak X_{(z_i)}|R)=(t_i)$, then
$t_1,\ldots,t_m$ can be extended to a system of free formal
generators of $R$ over $\Lambda$.

\smallskip
\noindent\textbf{(4)} The morphism
$\mathfrak X\longrightarrow\operatorname{Spec}(\Lambda)$ is formally
smooth; equivalently,
\[
J_{\mathfrak X}(\mathfrak X|\Lambda)=\varnothing.
\]

%%%%%%%%%%  scan idx 139  |  folio 128  %%%%%%%%%%

\smallskip
\noindent\textbf{(5)} If $X_0$ is a complete irreducible curve, then
\[
\dim\underline M_{X_0}=3g-3+a,
\]
where
\[
g=\dim_kH^1(X_0,\mathcal O_{X_0}),
\qquad
a=\dim_kH^0(X_0,\underline D^0_{X_0}).
\]

\smallskip
\noindent\textit{Proof.}
Assertion (1) follows from (2) and (3). The first assertion in (2)
follows from \SGARefLink{sga7.sga7i.evi.n5-7}{5.7}(2). For each $i$, let $(R_i,\mathfrak Y_i)$ be a
formal versal deformation of
$\operatorname{Spec}(\widehat{\mathcal O}_{X_0,z_i})$. Since
$(R,\mathfrak X_{(\widehat z_i)})$ is such a deformation, there is a
continuous local $\Lambda$-algebra homomorphism
\[
\varphi_i:R_i\longrightarrow R
\]
and a cartesian diagram
\[
\begin{tikzcd}[column sep=huge,row sep=large]
\mathfrak X_{(\widehat z_i)}
  \arrow[r] \arrow[d]
& \mathfrak Y_i \arrow[d]\\
\operatorname{Spf}(R)
  \arrow[r,"{\operatorname{Spf}(\varphi_i)}"']
& \operatorname{Spf}(R_i).
\end{tikzcd}
\]
Consequently,
\[
\begin{aligned}
J_{\mathfrak X_{(z_i)}}(\mathfrak X_{(z_i)}|R)
&=J_{\mathfrak X_{(\widehat z_i)}}
   (\mathfrak X_{(\widehat z_i)}|R)\\
&\simeq
J_{\mathfrak Y_i}(\mathfrak Y_i|R_i)
 \times_{\operatorname{Spf}(R_i)}\operatorname{Spf}(R),
\end{aligned}
\]
and
\[
\begin{aligned}
J_R(\mathfrak X_{(z_i)}|R)
&=J_R(\mathfrak X_{(\widehat z_i)}|R)\\
&\simeq
J_{R_i}(\mathfrak Y_i|R_i)
 \times_{\operatorname{Spf}(R_i)}\operatorname{Spf}(R).
\end{aligned}
\]
By Lemma \SGARefLink{sga7.sga7i.evi.n6-2}{6.2}(1),
\[
J_{\mathfrak Y_i}(\mathfrak Y_i|R_i)
 \longrightarrow J_{R_i}(\mathfrak Y_i|R_i)
\]
is an isomorphism, proving the remaining assertion in (2).

By Lemma \SGARefLink{sga7.sga7i.evi.n6-2}{6.2}(2),
\[
I_{R_i}(\mathfrak Y_i|R_i)=(u_i),
\qquad R_i=\Lambda[[u_i]].
\]
It follows that
\[
I_R(\mathfrak X_{(z_i)}|R)
=I_R(\mathfrak X_{(\widehat z_i)}|R)
=(t_i),
\qquad t_i=\varphi_i(u_i),
\]
and therefore
\[
I_R(\mathfrak X|R)
=\prod_{1\leq i\leq m}I_R(\mathfrak X_{(z_i)}|R)
=\left(\prod_{i=1}^m t_i\right).
\]
By \SGARefLink{sga7.sga7i.evi.n4-13}{4.13}, the natural morphism
\[
\underline M_{X_0}\longrightarrow
\underline M_{X_0,(\widehat z_1)}\times\cdots\times
\underline M_{X_0,(\widehat z_m)}
\]
is smooth. Hence
\[
R_1\widehat\otimes_\Lambda\cdots
 \widehat\otimes_\Lambda R_m\longrightarrow R
\]
is formally smooth. Thus $t_1,\ldots,t_m$ extend to a system of free
formal generators of $R$ over $\Lambda$, proving (1) and (3).

%%%%%%%%%%  scan idx 140  |  folio 129  %%%%%%%%%%

For (4), the projection
\[
\mathfrak X_{(\widehat z_i)}
 \xrightarrow{\ \sim\ }
\mathfrak Y_i\times_{\operatorname{Spf}(R_i)}
 \operatorname{Spf}(R)
 \longrightarrow\mathfrak Y_i
\]
is formally smooth because
$\operatorname{Spf}(R)\longrightarrow\operatorname{Spf}(R_i)$ is
formally smooth. Lemma \SGARefLink{sga7.sga7i.evi.n6-2}{6.2}(3) shows that
$\mathfrak Y_i\longrightarrow\operatorname{Spec}(\Lambda)$ is
formally smooth. Hence every
$\mathfrak X_{(\widehat z_i)}\longrightarrow
\operatorname{Spec}(\Lambda)$ is formally smooth. Since the $z_i$
are precisely the nonsmooth points of $X_0$, it follows that
$\mathfrak X\longrightarrow\operatorname{Spec}(\Lambda)$ is formally
smooth.

% END INCLUDED SGA7I_English_current_buildable_source_20260823/components/60_expose_VI_theorem_6_4_through_part_4.tex

% BEGIN INCLUDED SGA7I_English_current_buildable_source_20260823/components/61_expose_VI_theorem_6_4_part_5_and_remark_6_5.tex
For (5), $R$ is formally smooth over $\Lambda$, and by \SGARefLink{sga7.sga7i.evi.n4-4}{4.4} and \SGARefLink{sga7.sga7i.evi.n4-11}{4.11}
the sequence
\[
0\longrightarrow H^1(X_0,\underline D^0_{X_0})
 \longrightarrow[\underline M_{X_0}(k[\varepsilon])]
 \longrightarrow H^0(X_0,\underline D^1_{X_0})
 \longrightarrow0
\]
is exact. Thus
\[
\begin{aligned}
\dim\underline M_{X_0}
&=\dim_kH^1(X_0,\underline D^0_{X_0})
 +\dim_kH^0(X_0,\underline D^1_{X_0})\\
&=\dim_kH^1(X_0,\underline D^0_{X_0})+m.
\end{aligned}
\]
It is therefore enough to show that
\[
\dim_kH^1(X_0,\underline D^0_{X_0})+m
=3g-3+\dim_kH^0(X_0,\underline D^0_{X_0}),
\]
or equivalently
\[
\chi(\underline D^0_{X_0})=3-3g+m,
\]
where, for a coherent $\mathcal O_{X_0}$-module $\mathcal E$, we put
\[
\chi(\mathcal E)
=\dim_kH^0(X_0,\mathcal E)-\dim_kH^1(X_0,\mathcal E).
\]

The Euler characteristic of $\underline D^0_{X_0}$ is unchanged by
extension of the base field to its algebraic closure, so we may assume
that $k$ is algebraically closed. Let $z$ be a nonsmooth point of
$X_0$, put
\[
A=\widehat{\mathcal O}_{X_0,z},
\]
and let $A'$ be the normalization of $A$. By hypothesis,
\[
A\simeq k[[x,y]]/(f),
\]
where $f$ is a non-degenerate quadratic form. Over the algebraically
closed field $k$ we may take $f=xy$, so that
\[
A\simeq k[[x,y]]/(xy).
\]
It follows at once that
\[
\ell(A'/A)=1,
\qquad
\operatorname{Der}_k(A,\mathfrak m_A)
=\operatorname{Der}_k(A,A).
\]
Let $C$ be the conductor of $A$ in $A'$. Since $A$ is a complete
intersection,
\[
\ell(A/C)=\ell(A'/A)=1,
\]
and hence $C=\mathfrak m_A$. Therefore
\[
A'=\{a'\in K\mid a'\mathfrak m_A\subset\mathfrak m_A\},
\]
where $K$ is the total fraction ring of $A$.

We claim that
\[
\operatorname{Der}_k(A,A)
 \subset\operatorname{Der}_k(A',A')
\]
inside $\operatorname{Der}_k(K,K)$. Indeed, let
$D\in\operatorname{Der}_k(A,A)$. For $a'\in A'$ and
$x\in\mathfrak m_A$, one has $xa'\in\mathfrak m_A$, so
$D(xa')\in\mathfrak m_A$ and
$a'D(x)\in\mathfrak m_A$. The identity
\[
D(xa')=a'D(x)+xD(a')
\]
therefore gives $xD(a')\in\mathfrak m_A$ for every
$x\in\mathfrak m_A$, whence $D(a')\in A'$.

%%%%%%%%%%  scan idx 141  |  folio 130  %%%%%%%%%%

Moreover, the exact sequence
\[
0\longrightarrow\operatorname{Der}_k(A,A)
 \longrightarrow\operatorname{Der}_k(A',A')
 \longrightarrow A'/\mathfrak m_AA'\longrightarrow0
\]
may be written
\[
0\longrightarrow\operatorname{Der}_k(A,A)
 \longrightarrow\operatorname{Der}_k(A',A')
 \longrightarrow A'/C\longrightarrow0.
\]
If $X_0'$ is the normalization of $X_0$, these local sequences give
an exact sequence
\[
0\longrightarrow\underline D^0_{X_0}
 \longrightarrow\underline D^0_{X_0'}
 \longrightarrow\mathcal O_{X_0'}/C\longrightarrow0.
\]
Since $X_0$ is locally a complete intersection,
\[
\ell(\mathcal O_{X_0'}/C)
=2\ell(\mathcal O_{X_0'}/\mathcal O_{X_0})=2m,
\]
whereas the Riemann--Roch theorem gives
\[
\chi(\underline D^0_{X_0'})=3-3(g-m).
\]
Consequently,
\[
\begin{aligned}
\chi(\underline D^0_{X_0})
&=\chi(\underline D^0_{X_0'})
 -\ell(\mathcal O_{X_0'}/C)\\
&=3-3(g-m)-2m=3-3g+m.
\end{aligned}
\]
This completes the proof.
\hfill$\square$

\par\medskip
\noindent\textbf{Remark \SGARefTarget{sga7.sga7i.evi.n6-5}{6.5}.}
In the formula of Theorem \SGARefLink{sga7.sga7i.evi.n6-4}{6.4}(5), the number
$a=\dim_kH^0(X_0,\underline D^0_{X_0})$ is easily computed by the
Riemann--Roch theorem:
\[
a=
\begin{cases}
0,&g\geq2,\\
1,&g=1,\\
3,&g=0.
\end{cases}
\]

We finish with a remark on characteristic $2$. If
$\operatorname{char}k=2$, there is no non-degenerate quadratic form
in an odd number of variables. Equivalently, there is no
non-degenerate quadratic singularity of even dimension in
characteristic $2$. In this case one may instead consider a quadratic
form of maximal non-degeneracy.

% END INCLUDED SGA7I_English_current_buildable_source_20260823/components/61_expose_VI_theorem_6_4_part_5_and_remark_6_5.tex

% BEGIN INCLUDED SGA7I_English_current_buildable_source_20260823/components/62_expose_VI_6_6_through_6_8.tex
%%%%%%%%%%  scan idx 142  |  folio 131  %%%%%%%%%%

\par\medskip
\noindent\textbf{Definition \SGARefTarget{sga7.sga7i.evi.n6-6}{6.6}.}
Let $X$ be a $k$-scheme of finite type. A closed point $z\in X$ is
called an \emph{ordinary quadratic singularity rational over $k$} if
the following conditions hold.

\smallskip
\noindent\textbf{(i)} If either $\operatorname{char}k\neq2$, or
$\operatorname{char}k=2$ and $\dim\mathcal O_{X,z}$ is odd, then $z$
is a non-degenerate quadratic singularity rational over $k$.

\smallskip
\noindent\textbf{(ii)} If $\operatorname{char}k=2$ and
$\dim\mathcal O_{X,z}$ is even, then
\[
\overline k\otimes_k\widehat{\mathcal O}_{X,z}
 \simeq
\overline k[[X_0,X_1,\ldots,X_{2n}]]/(f),
\]
where $\overline k$ is an algebraic closure of $k$ and
\[
f=X_0^2+X_1X_2+X_3X_4+\cdots+X_{2n-1}X_{2n}.
\]

\par\medskip
\noindent\textbf{Lemma \SGARefTarget{sga7.sga7i.evi.n6-7}{6.7}.}
Let $k$ be a field of characteristic $2$, and put
\[
A_0=k[[X_0,X_1,\ldots,X_{2n}]]/(f),
\qquad
f=X_0^2+X_1X_2+X_3X_4+\cdots+X_{2n-1}X_{2n}.
\]
Let $(R,\mathfrak X)$ be a formal versal deformation of
$X_0=\operatorname{Spec}(A_0)$. Then
\[
R=\Lambda[[t,t_1]],
\]
where $t,t_1$ are free formal generators of $R$ over $\Lambda$, and
\[
I_R(\mathfrak X|R)=(t^2).
\]
Moreover,
$\mathfrak X\longrightarrow\operatorname{Spec}(\Lambda)$ is formally
smooth.

\smallskip
\noindent\textit{Proof.}
Write $\mathfrak X=\operatorname{Spf}(A)$. Since
\[
\left(f,\frac{\partial f}{\partial X_0},\ldots,
         \frac{\partial f}{\partial X_{2n}}\right)
=(X_0^2,X_1,\ldots,X_{2n}),
\]
we have
\[
D^1(A_0|k,A_0)\simeq k[X_0]/(X_0^2).
\]
Thus $[\underline M_{X_0}(k[\varepsilon])]$ is a two-dimensional
$k$-vector space. By \SGARefLink{sga7.sga7i.evi.n4-14}{4.14},
\[
R\simeq\Lambda[[t,t_1]],
\qquad
A\simeq R[[X_0,X_1,\ldots,X_{2n}]]/(F),
\]
where
\[
F=f+t_1+tX_0.
\]
The exact sequence
\[
0\longrightarrow A\xrightarrow{\ dF\ }
A\,dX_0\oplus\cdots\oplus A\,dX_{2n}
\longrightarrow\widehat\Omega_{A|R}\longrightarrow0,
\]
in which
\[
dF:\ 1\longmapsto
\sum_{0\leq i\leq2n}\frac{\partial f}{\partial X_i}\,dX_i
+t\,dX_0,
\]
gives
\[
I_A(A|R)
=(t,X_1,\ldots,X_{2n})
=(X_0^2+t_1,t,X_1,\ldots,X_{2n}).
\]
% The printed formula omits the coefficient t before dX_0, although
% F=f+t_1+tX_0 and the ensuing ideal both force this derivative term.

%%%%%%%%%%  scan idx 143  |  folio 132  %%%%%%%%%%

Hence
\[
\begin{aligned}
A/I_A(A|R)
&\simeq A/(X_0^2+t_1,t,X_1,\ldots,X_{2n})\\
&\simeq\Lambda[[t_1,X_0]]/(X_0^2+t_1).
\end{aligned}
\]
As an $R$-module this last module is isomorphic to
$R/(t)\oplus R/(t)$. It follows that
\[
I_R(A|R)=(t^2),
\]
which proves the first assertion. On the other hand,
\[
\begin{aligned}
A
&\simeq
\Lambda[[t,t_1,X_0,X_1,\ldots,X_{2n}]]/
 (X_0^2+X_1X_2+\cdots+X_{2n-1}X_{2n}+tX_0+t_1)\\
&\simeq\Lambda[[t,X_0,X_1,\ldots,X_{2n}]],
\end{aligned}
\]
so $A$ is formally smooth over $\Lambda$.
\hfill$\square$

Lemma \SGARefLink{sga7.sga7i.evi.n6-7}{6.7} gives the following analogue of Theorem \SGARefLink{sga7.sga7i.evi.n6-4}{6.4}; we leave its
proof to the reader.

\par\medskip
\noindent\textbf{Corollary \SGARefTarget{sga7.sga7i.evi.n6-8}{6.8}.}
Let $X_0$ be a $k$-scheme of finite type whose nonsmooth points are
isolated, assume that
\[
H^2(X_0,\underline D^0_{X_0})=0,
\]
and let $(R,\mathfrak X)$ be a formal versal deformation of $X_0$.
Denote the set of nonsmooth points by $\{z_1,\ldots,z_m\}$. Suppose
that $\operatorname{char}k=2$, that $\dim X_0$ is even, and that all
the $z_i$ are ordinary quadratic singularities. Then:

\smallskip
\noindent\textbf{(1)} $R$ is formally smooth over $\Lambda$.

\smallskip
\noindent\textbf{(2)} Each ideal
$I_R(\mathfrak X_{(z_i)}|R)$ is principal, as is
\[
I_R(\mathfrak X|R)
=\prod_{1\leq i\leq m}I_R(\mathfrak X_{(z_i)}|R).
\]
More precisely,
\[
I_R(\mathfrak X_{(z_i)}|R)=(t_i^2),
\]
and $t_1,\ldots,t_m$ can be extended to a system of free formal
generators of $R$ over $\Lambda$.

\smallskip
\noindent\textbf{(3)} The morphism
$\mathfrak X\longrightarrow\operatorname{Spec}(\Lambda)$ is formally
smooth.

\smallskip
\noindent\textbf{(4)} If $k$ is algebraically closed and $X_0$ is a
complete irreducible curve, then
\[
\dim\underline M_{X_0}=3g-3+a,
\]
where
\[
g=\dim_kH^1(X_0,\mathcal O_{X_0}),
\qquad
a=\dim_kH^0(X_0,\underline D^0_{X_0}).
\]

% END INCLUDED SGA7I_English_current_buildable_source_20260823/components/62_expose_VI_6_6_through_6_8.tex

\clearpage
\phantomsection\addcontentsline{toc}{section}{Exposé VII}
% BEGIN INCLUDED SGA7I_English_current_buildable_source_20260823/components/63_expose_VII_opening_through_1_1_2_1.tex
% SOURCE AUTHORITY:
% expose_VII_body.tex lines 5--58; scan indices 144--147; source folios 133--136.
% Hard stop before source line 60, paragraph 1.1.3.

%%%%%%%%%%  scan idx 144  |  folio 133  %%%%%%%%%%

\begin{center}
\SGARefAliasTarget{sga7.sga7i.evii.nroot}SGA 7

\vspace{1em}
\underline{Expos\'e VII}

\vspace{1em}
\underline{BIEXTENSIONS OF SHEAVES OF GROUPS}

\vspace{1em}
\underline{by A. Grothendieck}
\end{center}

0. \underline{Introduction}

The notion of a \underline{biextension} was introduced by D.~Mumford
[6], in the context of formal groups, as a convenient way of expressing
the relations between the formal groups associated with an abelian scheme
and the dual abelian scheme. In the present expos\'e we undertake a
systematic study of this notion in the general context of sheaves of
groups on a given site. In particular, we show that it fits remarkably
simply into the classical cohomological formalism of
$\overset{\mathrm L}{\otimes}$ and $\underline{\mathrm{RHom}}$
(\SGARefLink{sga7.sga7i.evii.n3-6-4}{3.6.4} and \SGARefLink{sga7.sga7i.evii.n3-6-5}{3.6.5}). In the following expos\'e we shall make explicit the
most interesting special cases, notably that of biextensions of group
schemes by the multiplicative group $\underline G_m$ (working with such
sites as the Zariski site, the fpqc site, or an intermediate site). We
shall see, in particular, that if $A$ and $B$ are two abelian schemes over
a scheme $S$, biextensions of $(A,B)$ by $\underline G_m$ admit a
remarkable interpretation as, essentially, the classical divisorial
correspondences on $A,B$. This gives a particularly convenient explicit
cohomological meaning to the notion of divisorial correspondence (already
implicit in Tate duality for abelian varieties over $p$-adic fields).
This interpretation of divisorial correspondences is also convenient for
following the fate of such a correspondence when $A$ and $B$ (assumed to
be defined over the fraction field of a discrete valuation ring) are made
to degenerate by means of N\'eron theory.

%%%%%%%%%%  scan idx 145  |  folio 134  %%%%%%%%%%

The results of Expos\'e VIII will show, for example, that one obtains a
biextension of the pair of identity components of the N\'eron models,
hence a biextension of their connected special fibres, which plays the
role of a specialization of the original divisorial correspondence.
These results will be used in Expos\'e IX, in conjunction with the
monodromy theorem (I~3), to prove certain results about the N\'eron models
of abelian varieties, the most important being the ``semistable reduction
theorem'' for the N\'eron model.

Throughout what follows, unless expressly stated otherwise, we work in a
fixed $\underline U$-topos denoted by $\underline T$; that is (SGA~4,
Expos\'e~IV), $\underline T$ is a category equivalent to the category of
sheaves on a suitable $\underline U$-site (which may, if desired, be taken
to be $\underline T$ itself, endowed with its canonical topology). In our
terminology we shall often identify the objects of $\underline T$ with
the sheaves of sets on $\underline T$. All groups, extensions of groups,
and so forth considered below are relative to this topos $\underline T$.

1. \underline{Complements on extensions of groups}

\SGARefTarget{sga7.sga7i.evii.n1-1}{1.1}. Let $P$ and $G$ be two groups of $\underline T$, and let

\begin{equation}\tag{\SGARefTarget{sga7.sga7i.evii.n1-1-1}{1.1.1}}
\begin{tikzcd}[column sep=large]
1 \arrow[r] & G \arrow[r, "i"] & E \arrow[r, "j"] & P \arrow[r] & 1
\end{tikzcd}
\end{equation}

be an extension of $P$ by $G$. Thus $E$ is a group,
$j:E\longrightarrow P$ is an epimorphism of groups, and $i$ induces an
isomorphism of $G$ with $\ker j$. We shall reinterpret the datum of such
an extension in different terms, which will be convenient for the
definition and study of biextensions.

%%%%%%%%%%  scan idx 146  |  folio 135  %%%%%%%%%%

\SGARefTarget{sga7.sga7i.evii.n1-1-2}{1.1.2}. First, the homomorphism $i$ makes $G$ act on $E$ in two ways, on
the left and on the right, by setting
\[
 gxg'=i(g)xi(g')
\]
for $g,g'\in G(S)$ and $x\in E(S)$, where $S$ is any object of
$\underline T$. If, by means of $j$, we regard $E$ as an object ``over
$P$'', that is, as an object of the induced topos $\underline T_{/P}$,
this amounts to saying that the group $G_P=G\times P$ induced by $G$ in
$\underline T_{/P}$ acts on $E$ on the left and on the right. Moreover,
it is well known that each of these actions makes $E$ a
\underline{torsor} (a left torsor and a right torsor, respectively) over
$P$, with group $G_P$; more precisely still, it makes $E$ a
\underline{bitorsor} under $(G_P,G_P)$. Recall [2] that this means that
$E$ is a right torsor $E_d$ under the right action of $G_P$, and that the
left action induces an isomorphism of $G_P$ with the group of
automorphisms of $E_d$ (the sheaf of groups of automorphisms of $E$ that
commute with the right action of $G_P$ on $E$). This condition is
equivalent to the symmetric condition: the left action of $G_P$ on $E$
makes $E$ a left torsor $E_s$, and the right action induces an
isomorphism of the opposite group $G_P^{\mathrm o}$ with the group of
automorphisms of $E_s$.

\SGARefTarget{sga7.sga7i.evii.n1-1-2-1}{1.1.2.1}. When $\underline T$ is the ``point topos'', that is, the
category of sets (equivalent to the category of sheaves on the one-point
topological space), so that $P$ and $G$ are ordinary groups, the datum of
a bitorsor $E$ under $(G_P,G_P)$ is plainly equivalent to the datum, for
each $p\in P$, of a bitorsor $E_p$ under $(G,G)$, namely the fibre of $E$
at $p$. An analogous interpretation is available in the general case.
Indeed, giving a bitorsor $E$ under $(G_P,G_P)$ is equivalent to giving,
for every ``point'' $p\in P(S)$, where $S$ is any object of
$\underline T$, a bitorsor $E_p$ under $(G_S,G_S)$ over $S$, in such a
way that the formation of $E_p$ is ``compatible with base change along
every morphism $S'\longrightarrow S$''. Thus, if
$p':S'\longrightarrow P$ is the composite morphism, one is also given an
isomorphism of bitorsors from $E_{p'}$ to the inverse image of $E_p$ under
$S'\longrightarrow S$, and these isomorphisms satisfy the usual
transitivity condition for a composite
$S''\longrightarrow S'\longrightarrow S$.

%%%%%%%%%%  scan idx 147  |  folio 136  %%%%%%%%%%

When stated in canonical language (SGA~1, Expos\'e~VI), this reduces to
the following tautology: giving a bitorsor under $(G_P,G_P)$ over $P$ is
equivalent to giving a cartesian section, over the category
$\underline T_{/P}$, of the fibred category of bitorsors under
$(G_S,G_S)$ over a variable base $S\in\operatorname{ob}\underline
T_{/P}$. More precisely, the ``value at $P$'' functor induces an
equivalence between the category of the cartesian sections in question
and the category of bitorsors under $(G_P,G_P)$. This follows trivially
from the fact that $P$ is a final object of the base category
$\underline T_{/P}$ of our fibred category: for \underline{every}
category fibred over $\underline T_{/P}$, the ``value at $P$'' functor is
an equivalence between the category of cartesian sections of that fibred
category and its fibre category over $P$. Despite its tautological
nature, however, this interpretation will be useful for expressing
certain compatibilities in language close to the set-theoretic case,
using points with values in an object $S$ in place of ordinary points.


% END INCLUDED SGA7I_English_current_buildable_source_20260823/components/63_expose_VII_opening_through_1_1_2_1.tex

% BEGIN INCLUDED SGA7I_English_current_buildable_source_20260823/components/64_expose_VII_1_1_3_through_1_1_4_2.tex
% SOURCE AUTHORITY:
% expose_VII_body.tex lines 60--137; scan indices 147--151; source folios 136--140.
% Hard stop before source line 139, paragraph 1.1.5.
% The print intermittently writes P_p where its own definitions and the
% associativity diagram require E_p; the coherent E_p notation is retained.

\SGARefTarget{sga7.sga7i.evii.n1-1-3}{1.1.3}. In terms of the extension structure (\SGARefLink{sga7.sga7i.evii.n1-1-1}{1.1.1}), the bitorsor
structure of \SGARefLink{sga7.sga7i.evii.n1-1-2}{1.1.2} on $E$ under $(G_P,G_P)$ used only part of the
multiplicative structure of $E$: namely, multiplication of pairs of
``points'' of which at least one comes from a point of $G$ (via $i$). We
shall now interpret the multiplicative structure of $E$ as an additional
structure on the bitorsor under consideration.

%%%%%%%%%%  scan idx 148  |  folio 137  %%%%%%%%%%

For this purpose, first consider two sections $p,p'$ of $P$. With the
notation of \SGARefLink{sga7.sga7i.evii.n1-1-2-1}{1.1.2.1}, these give two bitorsors $E_p,E_{p'}$ under $(G,G)$,
namely the inverse images in $E$ of the sections
$p,p':e\longrightarrow P$ (where $e$ denotes the final object of
$\underline T$). It is then clear that multiplication in $E$ induces a
morphism of objects of $\underline T$
\[
 \varphi:E_p\times E_{p'}\longrightarrow E_{pp'},
\]
which, in terms of the left and right structures of the three bitorsors
involved (and, as usual, using points with values in an arbitrary object
$S\in\operatorname{ob}\underline T$), satisfies
\begin{equation}\tag{\SGARefTarget{sga7.sga7i.evii.n1-1-3-0}{1.1.3.0}}
\left\{
\begin{aligned}
 \varphi(gx,x')&=g\varphi(x,x'),
 &\qquad \varphi(x,x'g')&=\varphi(x,x')g',\\
 \varphi(xg,x')&=\varphi(x,gx').
\end{aligned}
\right.
\end{equation}
Here $g,g'\in G(S)$, $x\in E_p(S)$, and $x'\in E_{p'}(S)$. In terms of
the contracted product [2]
\[
 E_p\mathbin{\overset{G}{\times}}E_{p'},
\]
the third relation says that $\varphi$ factors as
\[
 E_p\times E_{p'}\longrightarrow
 E_p\mathbin{\overset{G}{\times}}E_{p'}
 \xrightarrow{\ \varphi_{p,p'}\ }E_{pp'},
\]
where the first arrow is the canonical quotient morphism. The first two
relations, on the other hand, say that the morphism thus obtained,
\begin{equation}\tag{\SGARefTarget{sga7.sga7i.evii.n1-1-3-1}{1.1.3.1}}
 \varphi_{p,p'}:
 E_p\mathbin{\overset{G}{\times}}E_{p'}\longrightarrow E_{pp'},
\end{equation}
is a homomorphism of \underline{bitorsors}, where the bitorsor structure
on the source is the one for which the left (respectively, right)
operations of $G$ come from the left (respectively, right) operations of
$G$ on the first factor $E_p$ (respectively, on the second factor
$E_{p'}$).

%%%%%%%%%%  scan idx 149  |  folio 138  %%%%%%%%%%

Notice that, like every homomorphism of bitorsors under a fixed pair of
groups, $\varphi_{p,p'}$ is necessarily an \underline{isomorphism}.

The definition of the preceding isomorphism (\SGARefLink{sga7.sga7i.evii.n1-1-3-1}{1.1.3.1}) immediately extends
to the case in which $p,p'$ are points of $P$ with values in an arbitrary
object $S$ of $\underline T$, rather than sections of $P$: one simply
applies the preceding discussion to the induced topos
$\underline T_{/S}$ and to the extension of groups over $S$ obtained from
(\SGARefLink{sga7.sga7i.evii.n1-1-1}{1.1.1}) by base change. The formation of the isomorphisms
$\varphi_{p,p'}$ is then compatible with every base change
$S'\longrightarrow S$. In other words, it corresponds to a homomorphism
of cartesian sections of the fibred category over the category
$\underline T_{/P\times P}$ of pairs of arrows
$p,p':S\rightrightarrows P$ with target $P$.

Giving a system of $\varphi_{p,p'}$ satisfying this condition is, of
course, equivalent to giving $\varphi_{p,p'}$ in the ``universal'' case
$S=P\times P$, where $p,p'$ are the two projections
$\operatorname{pr}_i$ ($i=1,2$) from $P\times P$ to $P$. It is then an
isomorphism of bitorsors under
$(G_{P\times P},G_{P\times P})$:
\begin{equation}\tag{\SGARefTarget{sga7.sga7i.evii.n1-1-3-2}{1.1.3.2}}
 \varphi:\operatorname{pr}_1^*(E)\,\operatorname{pr}_2^*(E)
 \overset{\sim}{\longrightarrow}\pi^*(E),
\end{equation}
where
\[
 \pi:P\times P\longrightarrow P
\]
is the morphism of objects of $\underline T$ defining the multiplicative
structure of the group $P$. The compatibility properties below are most
conveniently stated, however, in terms of the data (\SGARefLink{sga7.sga7i.evii.n1-1-3-1}{1.1.3.1}), avoiding
diagrams that are less close to intuition.

%%%%%%%%%%  scan idx 150  |  folio 139  %%%%%%%%%%

\SGARefTarget{sga7.sga7i.evii.n1-1-4}{1.1.4}. In writing the relations (\SGARefLink{sga7.sga7i.evii.n1-1-3-0}{1.1.3.0}), which enabled us to express
the given composition law on $E$ in terms of isomorphisms of bitorsors
(\SGARefLink{sga7.sga7i.evii.n1-1-3-1}{1.1.3.1}), we used associativity of this composition law only for a
product of three ``points'' of $E$ of which at least one comes from a
point of $G$. To express full associativity in terms of the
$\varphi_{p,p'}$, one must require, for three points $p,p',p''$ of $P$
with values in an object $S$ of $\underline T$, that the following
diagram be commutative:
\begin{equation}\tag{\SGARefTarget{sga7.sga7i.evii.n1-1-4-1}{1.1.4.1}}
\begin{tikzcd}[column sep=huge, row sep=large]
 & E_{pp'}E_{p''} \arrow[dr, "{\varphi_{pp',p''}}"] & \\
 E_pE_{p'}E_{p''}
   \arrow[ur, "{\varphi_{p,p'}\wedge\id}"]
   \arrow[dr, "{\id\wedge\varphi_{p',p''}}"']
 && E_{pp'p''} \\
 & E_pE_{p'p''} \arrow[ur, "{\varphi_{p,p'p''}}"'] &
\end{tikzcd}
\end{equation}
In this diagram we have written the contracted product of bitorsors
without the signs $\mathbin{\overset{G_S}{\times}}$. It is understood
that the contracted product $Q.R$ of two bitorsors is formed using, in
passing to the quotient in $Q\times R$, the action of the structural
group $G$ \underline{on the right of} $Q$ \underline{and on the left of}
$R$. This contracted-product operation is associative up to a canonical
isomorphism [2], and we make the identification
\[
 (E_pE_{p'})E_{p''}\simeq E_p(E_{p'}E_{p''}),
\]
writing both sides simply as $E_pE_{p'}E_{p''}$. Obvious conventions of
this kind will be used throughout the rest of this expos\'e without being
expressly noted each time.

\SGARefTarget{sga7.sga7i.evii.n1-1-4-2}{1.1.4.2}. Just as the isomorphisms $\varphi_{p,p'}$ of (\SGARefLink{sga7.sga7i.evii.n1-1-3-1}{1.1.3.1}) are
determined by the ``universal'' case $S=P\times P$,
$p=\operatorname{pr}_1$, $p'=\operatorname{pr}_2$ considered in
(\SGARefLink{sga7.sga7i.evii.n1-1-3-2}{1.1.3.2}), so the associativity condition (\SGARefLink{sga7.sga7i.evii.n1-1-4-1}{1.1.4.1}) may likewise be
checked only in the universal case
$S=P\times P\times P$, $p=\operatorname{pr}_1$,
$p'=\operatorname{pr}_2$, $p''=\operatorname{pr}_3$. We leave it to the
reader to write the diagram in whatever notation is preferred.


% END INCLUDED SGA7I_English_current_buildable_source_20260823/components/64_expose_VII_1_1_3_through_1_1_4_2.tex

% BEGIN INCLUDED SGA7I_English_current_buildable_source_20260823/components/65_expose_VII_1_1_5_through_1_1_6.tex
% SOURCE AUTHORITY:
% expose_VII_body.tex lines 139--235; scan indices 151--154; source folios 140--143.
% Hard stop before source line 237, paragraph 1.2.

\SGARefTarget{sga7.sga7i.evii.n1-1-5}{1.1.5}. Conversely, suppose that the groups $P$ and $G$ of
$\underline T$ are given, together with a bitorsor
\begin{equation}\tag{\SGARefTarget{sga7.sga7i.evii.n1-1-5-1}{1.1.5.1}}
 j:E\longrightarrow P
\end{equation}
under $(G_P,G_P)$ (the two actions of $G_P$ on $E$, or equivalently of
$G$ on $E$, being understood in the notation), and a system of
isomorphisms of bitorsors over a variable base $S$,
\begin{equation}\tag{\SGARefTarget{sga7.sga7i.evii.n1-1-5-2}{1.1.5.2}}
 \varphi_{p,p'}:E_pE_{p'}\overset{\sim}{\longrightarrow}E_{pp'}
 \qquad (p,p'\in P(S)),
\end{equation}
compatible with base change $S'\longrightarrow S$; in other words, an
isomorphism of bitorsors (\SGARefLink{sga7.sga7i.evii.n1-1-3-2}{1.1.3.2}). Suppose that the corresponding
diagrams (\SGARefLink{sga7.sga7i.evii.n1-1-4-1}{1.1.4.1}) commute. I claim that \underline{there is a unique
extension structure on $E$ of $P$ by $G$ such that the preceding data
are obtained from it by the constructions of \SGARefLink{sga7.sga7i.evii.n1-1-2}{1.1.2} and \SGARefLink{sga7.sga7i.evii.n1-1-3}{1.1.3}}.

First observe that, on taking $p,p'$ in (\SGARefLink{sga7.sga7i.evii.n1-1-5-2}{1.1.5.2}) to be the identity
section $1$ of $P$, one obtains an isomorphism
\begin{equation}\tag{\SGARefTarget{sga7.sga7i.evii.n1-1-5-3}{1.1.5.3}}
 \varphi_{1,1}:E_1\overset{\sim}{\longleftarrow}
 E_1\mathbin{\overset{G}{\times}}E_1.
\end{equation}
Taking the contracted product, on the left say, of both sides with the
inverse bitorsor $E_1^{-1}$ of $E_1$ [2] therefore gives a canonical
isomorphism
\begin{equation}\tag{\SGARefTarget{sga7.sga7i.evii.n1-1-5-4}{1.1.5.4}}
 E_1\simeq\text{the unit bitorsor }{}_sG_d
\end{equation}
(that is, $G$ acting on itself by left and right translations), in view
of the canonical isomorphism

%%%%%%%%%%  scan idx 152  |  folio 141  %%%%%%%%%%

\[
 E_1^{-1}.E_1\simeq\text{the unit bitorsor }{}_sG_d
\]
of loc.~cit., valid for every bitorsor. Equivalently, one has a section
$1_E$ of $E_1$, hence of $E$,
\begin{equation}\tag{\SGARefTarget{sga7.sga7i.evii.n1-1-5-5}{1.1.5.5}}
 1_E\in\Gamma(E_1)=\operatorname{Hom}(e,E_1)\subset\Gamma(E),
\end{equation}
namely the image under (\SGARefLink{sga7.sga7i.evii.n1-1-5-4}{1.1.5.4}) of the distinguished section of
${}_sG_d$, satisfying
\[
 g.1_E=1_E.g\qquad(g\in G(S)).
\]
By construction this section is characterized by
\begin{equation}\tag{\SGARefTarget{sga7.sga7i.evii.n1-1-5-6}{1.1.5.6}}
 \varphi_{1,1}(1_E,1_E)=1_E.
\end{equation}

On the other hand, giving the morphism of torsors $\varphi$ of
(\SGARefLink{sga7.sga7i.evii.n1-1-3-2}{1.1.3.2}) is plainly equivalent to giving a morphism
\begin{equation}\tag{\SGARefTarget{sga7.sga7i.evii.n1-1-5-7}{1.1.5.7}}
 \pi_E:E\times E\longrightarrow E,
 \qquad\text{also denoted }(x,y)\longmapsto xy,
\end{equation}
satisfying
\begin{equation}\tag{\SGARefTarget{sga7.sga7i.evii.n1-1-5-8}{1.1.5.8}}
 j(xy)=j(x)j(y)\qquad(x,y\in E(S)),
\end{equation}
so that $j$ is compatible with the multiplicative structures $\pi_E$
on $E$ and $\pi$ on $P$, together with the three relations (\SGARefLink{sga7.sga7i.evii.n1-1-3-0}{1.1.3.0})
(preferably beginning with the last of these). Finally define a morphism
of objects of $\underline T$,
\[
 i:G\longrightarrow E,
\]
by the formula
\begin{equation}\tag{\SGARefTarget{sga7.sga7i.evii.n1-1-5-9}{1.1.5.9}}
 i(g)=g.1_E=1_E.g.
\end{equation}

%%%%%%%%%%  scan idx 153  |  folio 142  %%%%%%%%%%

The equality of the last two terms was already observed in constructing
$1_E$. I now claim that \underline{the multiplication law $\pi_E$ on
$E$ makes $E$ a group with identity section $1_E$, and that
$j:E\longrightarrow P$ and $i:G\longrightarrow E$ are homomorphisms of
groups making $E$ an extension of $P$ by $G$}.

Indeed, associativity of the composition law on $E$ is precisely the
commutativity of the diagrams (\SGARefLink{sga7.sga7i.evii.n1-1-4-1}{1.1.4.1}). The construction of $\pi_E$ by
means of the isomorphisms $\varphi_{p,p'}$ of (\SGARefLink{sga7.sga7i.evii.n1-1-3-1}{1.1.3.1}) also shows that,
for every $x\in E(S)$, left and right translation by $x$ in $E(S)$ are
injective. Since $1_E1_E=1_E$ by (\SGARefLink{sga7.sga7i.evii.n1-1-5-6}{1.1.5.6}), it follows at once that
$1_E$ is a two-sided identity of $E$. We leave to the reader the
verification that an $x\in E(S)$ lying over $p\in P(S)$ has an inverse;
one uses the canonical isomorphism deduced from $\varphi_{p,p^{-1}}$,
\begin{equation}\tag{\SGARefTarget{sga7.sga7i.evii.n1-1-5-10}{1.1.5.10}}
 E_{p^{-1}}\simeq(E_p)^{\mathrm o}.
\end{equation}
Thus $E$ is indeed a \underline{group}. We have already observed that
$j$ is compatible with the composition laws (\SGARefLink{sga7.sga7i.evii.n1-1-5-8}{1.1.5.8}), and the same is
true of $i$, since, for $g,g'\in G(S)$,
\begin{align*}
 i(g)i(g')
 &= (g.1_E)(g'.1_E)\\
 &= (g.1_E.g')1_E &&\text{(the last relation of (\SGARefLink{sga7.sga7i.evii.n1-1-3-0}{1.1.3.0}))}\\
 &= g.1_E.g' &&\text{(since $1_E$ is an identity)}\\
 &= g.(1_E.g')=g.(g'.1_E)\\
 &= (gg').1_E &&\text{by (\SGARefLink{sga7.sga7i.evii.n1-1-5-9}{1.1.5.9})}\\
 &= i(gg').
\end{align*}
Since $j$ is the structural projection of a torsor, it is an
epimorphism of objects of $\underline T$. It remains only to see that
$i$ induces an isomorphism of $G$ with $\ker j=E_1$, which is clear from
the definition (\SGARefLink{sga7.sga7i.evii.n1-1-5-9}{1.1.5.9}) and the construction of $1_E$.

%%%%%%%%%%  scan idx 154  |  folio 143  %%%%%%%%%%

\SGARefTarget{sga7.sga7i.evii.n1-1-6}{1.1.6}. We shall accept that the preceding arguments prove the following.
For a fixed topos $\underline T$, the preceding constructions give an
\underline{equivalence between the category of extensions of groups of
$\underline T$} as in (\SGARefLink{sga7.sga7i.evii.n1-1-1}{1.1.1}), where $P,G,E$ are all variable,
\underline{and the category of systems $(P,G,E,\varphi)$ in which $P$
and $G$ are groups of $\underline T$, $E$ is a bitorsor under
$(G_P,G_P)$ over $P$, and $\varphi$ is an isomorphism of bitorsors}
(\SGARefLink{sga7.sga7i.evii.n1-1-3-2}{1.1.3.2}) (or, equivalently, a system of base-change-compatible
isomorphisms of bitorsors (\SGARefLink{sga7.sga7i.evii.n1-1-3-1}{1.1.3.1})), \underline{satisfying the
associativity condition expressed by commutativity of diagram (\SGARefLink{sga7.sga7i.evii.n1-1-4-2}{1.1.4.2})}
(or, equivalently, of the diagrams (\SGARefLink{sga7.sga7i.evii.n1-1-4-1}{1.1.4.1})).

We leave it to the reader to spell out, using the preceding dictionary,
the evident notion of inverse image (respectively, direct image) of an
extension (\SGARefLink{sga7.sga7i.evii.n1-1-1}{1.1.1}) under a homomorphism of groups
\[
 P'\longrightarrow P
 \qquad\bigl(\text{respectively, }G\longrightarrow G'\bigr),
\]
or under both operations simultaneously. We likewise leave to the reader
the compatibility of the preceding equivalence of categories with the
inverse-image functor $f^*$ associated with a morphism of topoi
\[
 f:\underline T'\longrightarrow\underline T.
\]
For this one should use the datum $\varphi$ in the form (\SGARefLink{sga7.sga7i.evii.n1-1-3-2}{1.1.3.2}), in
preference to the $\varphi_{p,p'}$ of (\SGARefLink{sga7.sga7i.evii.n1-1-3-1}{1.1.3.1}); it then defines the
corresponding datum $f^*(\varphi)$ on the topos $\underline T'$.


% END INCLUDED SGA7I_English_current_buildable_source_20260823/components/65_expose_VII_1_1_5_through_1_1_6.tex

% BEGIN INCLUDED SGA7I_English_current_buildable_source_20260823/components/66_expose_VII_section_1_2.tex
% SOURCE AUTHORITY:
% expose_VII_body.tex lines 237--261; scan indices 154--155; source folios 143--144.
% Hard stop before source line 264, section 1.3.

\SGARefTarget{sga7.sga7i.evii.n1-2}{1.2}. \underline{The commutative case}. From \SGARefLink{sga7.sga7i.evii.n1-4}{1.4} onward we shall be
concerned only with \underline{commutative} extensions of groups. We
therefore assume that $P$ and $G$ are commutative, and spell out here,
in terms of the bitorsor description $(P,G,E,\varphi)$ of an extension,
the condition that $E$ be a commutative group.

%%%%%%%%%%  scan idx 155  |  folio 144  %%%%%%%%%%

An obviously necessary first condition---which in fact expresses that
$i(G)$ be a \underline{central} subgroup of $E$---is that the bitorsors
$E_p$, for $p\in P(S)$, be ordinary torsors; that is, the left and right
actions of $G$ on $E_p$ must coincide. Equivalently, $E$ over $P$ is an
ordinary torsor under $G_P$. (Thus there will henceforth be no need to
speak of bitorsors, and we shall consider only ordinary torsors.) Assuming
this preliminary condition, commutativity of $E$ is expressed by
commutativity of the following diagrams, for $p,p'\in P(S)$:
\begin{equation}\tag{\SGARefTarget{sga7.sga7i.evii.n1-2-1}{1.2.1}}
\begin{tikzcd}[column sep=large, row sep=large]
 E_pE_{p'} \arrow[r, "{\varphi_{p,p'}}"] \arrow[d, "\mathrm{sym}"']
   & E_{pp'} \arrow[d, "\id"] \\
 E_{p'}E_p \arrow[r, "{\varphi_{p',p}}"]
   & E_{p'p}.
\end{tikzcd}
\end{equation}
Here the first vertical arrow is the canonical commutativity isomorphism
(defined because $G$ is now commutative); moreover $pp'=p'p$ because
$P$ is commutative. This condition is plainly equivalent to the
corresponding condition in the universal situation
$S=P\times P$, $p=\operatorname{pr}_1$,
$p'=\operatorname{pr}_2$, which we need not write out again.


% END INCLUDED SGA7I_English_current_buildable_source_20260823/components/66_expose_VII_section_1_2.tex

% BEGIN INCLUDED SGA7I_English_current_buildable_source_20260823/components/67_expose_VII_1_3_1_through_1_3_3.tex
% SOURCE AUTHORITY:
% expose_VII_body.tex lines 264--286; scan indices 155--158; source folios 144--147.
% Hard stop before source line 288, paragraph 1.3.4.

\SGARefTarget{sga7.sga7i.evii.n1-3}{1.3}. \underline{Relations with rigidified or birigidified torsors}

\SGARefTarget{sga7.sga7i.evii.n1-3-1-0}{\SGARefAliasTarget{sga7.sga7i.evii.n1-3-1}1.3.1.0}. Let $P$ be an object of $\underline T$ equipped with a section
$e_P$, let $G$ be a group of $\underline T$, and let $E$ be a torsor
under $G_P$. A \underline{rigidification of $E$}, relative to the given
section $e_P$ of $P$, is a trivialization of the torsor
$E_1=e_P^*(E)$ under $G$;

%%%%%%%%%%  scan idx 156  |  folio 145  %%%%%%%%%%

that is, a section of this torsor. This terminology is inspired by the
case where $G$ is commutative and
\begin{equation}\tag{\SGARefTarget{sga7.sga7i.evii.n1-3-1-1}{1.3.1.1}}
 \operatorname{Hom}_{\mathrm{pt}}(P,G)\simeq(e),
\end{equation}
where the left-hand side denotes the set of morphisms carrying $e_P$ to
the identity section $e_G$ of $G$. This relation is equivalent to saying
that the natural homomorphism
\begin{equation}\tag{\SGARefTarget{sga7.sga7i.evii.n1-3-1-2}{1.3.1.2}}
 \Gamma(G)\xrightarrow{\ \sim\ }\operatorname{Hom}(P,G),
\end{equation}
which associates with a section $g$ of $G$ the constant morphism
$P\longrightarrow G$ with value $g$, is an isomorphism. It is immediate
that these equivalent conditions are also equivalent to saying that
every automorphism of the torsor $E$ preserving the given
rigidification is the identity automorphism (such automorphisms
correspond precisely to the elements of the left-hand side of
(\SGARefLink{sga7.sga7i.evii.n1-3-1-1}{1.3.1.1})). In particular, under condition (\SGARefLink{sga7.sga7i.evii.n1-3-1-1}{1.3.1.1}), the category of
torsors under $G_P$ rigidified relative to $e_P$ is ``rigid''.

\SGARefTarget{sga7.sga7i.evii.n1-3-2}{1.3.2}. Let, in addition, $Q$ be an object of $\underline T$ equipped
with a section $e_Q$, and now let $E$ denote a torsor under
$G_{P\times Q}$. Since $e_Q$ defines a section
$e_{P\times Q/P}$ of $P\times Q$ over $P$, one may consider
rigidifications $\alpha$ of $E$ relative to this section, working in the
induced topos $\underline T_{/P}$; explicitly, these are sections of the
torsor induced by $E$ on $P\times e_Q$. Interchanging the roles of $P$
and $Q$, one may likewise consider rigidifications $\beta$ of $E$
relative to the section $e_{P\times Q/Q}$ of $P\times Q$ over $Q$
defined by $e_P$. Two such partial rigidifications are said to be
\underline{compatible} if the sections $\alpha$ and $\beta$ agree on
$e_P\times e_Q$, that is, if they define the same section of the torsor
$E_{1,1}$ under $G$, the inverse image

%%%%%%%%%%  scan idx 157  |  folio 146  %%%%%%%%%%

of $E$ under the section $(e_P,e_Q)$ of $P\times Q$. A
\underline{birigidification} of the torsor $E$ relative to the pair of
sections $(e_P,e_Q)$ is a pair of compatible partial rigidifications.

If $G$ is commutative, then $E$ admits a birigidification
$(\alpha,\beta)$ if and only if it admits partial rigidifications
$\alpha$ and $\beta$; moreover one may find a birigidification in which
$\alpha$ (or $\beta$) is prescribed in advance. Indeed, it suffices to
correct $\beta$ by the morphism $Q\longrightarrow G$ obtained as the
composite
\[
 Q\longrightarrow e\xrightarrow{\lambda}G,
\]
where $\lambda$ is defined by
$\beta e_Q=(\alpha e_P)\lambda$.

Still in the commutative case, if $E,E'$ are two birigidified torsors
under $G_{P\times Q}$, saying that they are isomorphic merely as
$G_{P\times Q}$-torsors is equivalent to saying that they are
isomorphic as birigidified $G_{P\times Q}$-torsors. In other words, if
there is an isomorphism $u:E\longrightarrow E'$ of
$G_{P\times Q}$-torsors, then there is one
$v:E\longrightarrow E'$ that preserves the birigidifications. Indeed,
measuring the discrepancy of $u$ from the given birigidifications gives
morphisms $f:P\longrightarrow G$ and $g:Q\longrightarrow G$ such that
$f(e_P)=g(e_Q)=a\in G(e)$. There is then a morphism
$P\times Q\longrightarrow G$ whose restrictions along
$e_{P\times Q/P}$ and $e_{P\times Q/Q}$ are respectively $f$ and $g$,
namely
\[
 h(p,q)=f(p)g(q)a^{-1},
\]
and one may take $v=uh$.

\SGARefTarget{sga7.sga7i.evii.n1-3-3}{1.3.3}. The terminology of \SGARefLink{sga7.sga7i.evii.n1-3-1}{1.3.1} extends immediately to the case of a
\underline{bitorsor} $E$ under $(G_P,G_P)$. A rigidification of the
bitorsor $E$ is then required to be an isomorphism from $e_P^*(E)$ to
the trivial bitorsor ${}_sG_d$; equivalently, it is defined by a
\underline{central} section $e_E$ of $e_P^*(E)$, one satisfying
\[
 g.e_E=e_E.g\qquad(g\in G(S)).
\]
The terminology of birigidifications extends in the same way to
bitorsors.


% END INCLUDED SGA7I_English_current_buildable_source_20260823/components/67_expose_VII_1_3_1_through_1_3_3.tex

% BEGIN INCLUDED SGA7I_English_current_buildable_source_20260823/components/68_expose_VII_1_3_4.tex
% SOURCE AUTHORITY:
% expose_VII_body.tex lines 288--303; scan indices 158--159; source folios 147--148.
% Hard stop before source line 305, Proposition 1.3.5.

\SGARefTarget{sga7.sga7i.evii.n1-3-4}{1.3.4}. Return to the case in which two groups $P$ and $G$ are given,
and choose the identity section $e_P$ in order to define
rigidifications of bitorsors over $P$ and birigidifications of bitorsors
over $P\times P$. There are natural functors
\begin{equation}\tag{\SGARefTarget{sga7.sga7i.evii.n1-3-4-1}{1.3.4.1}}
 \mathbf{EXT}(P;G)\xrightarrow{\ i\ }
 \mathrm{BITORSRIG}(P,G)\xrightarrow{\ \delta\ }
 \mathbf{BITORSBIRIG}(P,P;G),
\end{equation}
where the three expressions denote, respectively, the category of
extensions of the group $P$ by $G$, the category of bitorsors under
$(G_P,G_P)$ rigidified at $e_P$, and the category of bitorsors under
$(G_{P\times P},G_{P\times P})$ birigidified at $(e_P,e_P)$. The functor
$i$ associates with an extension $E$ of $P$ by $G$ the corresponding
bitorsor over $P$ (\SGARefLink{sga7.sga7i.evii.n1-1-6}{1.1.6}), rigidified by the identity section of $E$.
The functor $\delta$ is defined by
\begin{equation}\tag{\SGARefTarget{sga7.sga7i.evii.n1-3-4-2}{1.3.4.2}}
 \delta(E)=\pi^*(E)\,\pr_1^*(E)^{-1}\,\pr_2^*(E)^{-1},
\end{equation}
where $\pi:P\times P\longrightarrow P$ is the composition law of $P$,
and $\pr_1,\pr_2$ are the two projections. There is a
\underline{canonical birigidification of $\delta(E)$}. Indeed, the
restriction of $\delta(E)$ to $P\times e_P$ is canonically isomorphic
to
\[
 \id_P^*(E)\,\id_P^*(E)^{-1}\,e^*(E)^{-1},
\]
where $e=\pr_2\operatorname{inj}_1$ is the constant morphism
$P\longrightarrow P$ with value $e_P$. The given rigidification of $E$
provides an isomorphism from $e^*(E)$ to the trivial bitorsor, and hence
an isomorphism of the restriction of $\delta(E)$ to $P\times e_P$ with
the trivial bitorsor. Symmetrically one defines a partial
rigidification along $e_P\times P$, and the two partial
rigidifications are immediately seen to be compatible.

%%%%%%%%%%  scan idx 159  |  folio 148  %%%%%%%%%%

\SGARefTarget{sga7.sga7i.evii.n1-3-4-3}{1.3.4.3}. Observe that the composite $\delta i$ of the functors in
(\SGARefLink{sga7.sga7i.evii.n1-3-4-1}{1.3.4.1}) is canonically isomorphic to the trivial functor, that is, to
the constant functor whose value is the trivial birigidified bitorsor
over $P\times P$. Indeed, for every extension $E$ of $P$ by $G$, the
isomorphism $\varphi$ of (\SGARefLink{sga7.sga7i.evii.n1-1-3-2}{1.1.3.2}) plainly defines a trivialization of
$\delta i(E)$ compatible with its canonical birigidification.


% END INCLUDED SGA7I_English_current_buildable_source_20260823/components/68_expose_VII_1_3_4.tex

% BEGIN INCLUDED SGA7I_English_current_buildable_source_20260823/components/69_expose_VII_proposition_1_3_5.tex
% SOURCE AUTHORITY:
% expose_VII_body.tex lines 305--358; scan indices 159--162; source folios 148--151.
% Source-backed corrections: the proposition's opening Q is G throughout its statement;
% E in part (b) is a bitorsor, as required by delta(E) and the ensuing sentence;
% the final associativity specialization follows the three partial products stated in (b).
% Hard stop before source line 360, Corollary 1.3.6.

\underline{Proposition \SGARefTarget{sga7.sga7i.evii.n1-3-5}{1.3.5}. Let $P$ and $G$ be two groups of
$\underline T$.}

a) \underline{Suppose that every morphism of sheaves of sets
$P\times P\longrightarrow G$ which is trivial on $P\times e_P$ and on
$e_P\times P$ is trivial. Then the functor $i$ of (\SGARefLink{sga7.sga7i.evii.n1-3-4-1}{1.3.4.1}) is fully
faithful. In particular, on a bitorsor $E$ under $(G_P,G_P)$ rigidified
relative to $e_P$, that is, equipped with a central section $e_E$ of
$E_1=e_P^*(E)$, there is at most one extension structure of $P$ by $G$
compatible with its bitorsor structure and having $e_E$ as identity
section.}

b) \underline{Suppose in addition that every morphism of sheaves of
sets $P\times P\times P\longrightarrow G$ which is trivial on the
partial products $P\times P\times e_P$, $P\times e_P\times P$, and
$e_P\times P\times P$ is trivial. Let $E$ be a bitorsor under
$(G_P,G_P)$. In order that $E$ admit an extension structure of $P$ by
$G$ compatible with its bitorsor structure, it is necessary and
sufficient that the bitorsor under
$(G_{P\times P},G_{P\times P})$ underlying}
\[
 \delta(E)=\pi^*(E)\,\pr_1^*(E)^{-1}\,\pr_2^*(E)^{-1}
\]
\underline{be trivial. When this is so, for every central section
$e_E$ of $E_1=e_P^*(E)$, that is, every rigidification of $E$ relative
to $e_P$, there is a unique extension structure on $E$ compatible with
its bitorsor structure and having $e_E$ as identity section. In other
words, $E$, regarded as a bitorsor rigidified relative to $e_P$, lies
in the essential image of $i$. Moreover, the birigidified bitorsor
$\delta(E)$ has one and only one trivialization.}

%%%%%%%%%%  scan idx 160  |  folio 149  %%%%%%%%%%

Thus, in terms of diagram (\SGARefLink{sga7.sga7i.evii.n1-3-4-1}{1.3.4.1}), the essential image of the fully
faithful functor $i$ is ``the kernel'' of $\delta$. More picturesquely,
diagram (\SGARefLink{sga7.sga7i.evii.n1-3-4-1}{1.3.4.1}) is an ``exact sequence of pointed categories,'' and
even of categories endowed with pseudo-group structures (in the
terminology introduced in \SGARefLink{sga7.sga7i.evii.n2-5-1}{2.5.1} below).

Let us prove \SGARefLink{sga7.sga7i.evii.n1-3-5}{1.3.5}. For a), let $E,E'$ be two extensions of $P$ by $G$,
and let $f:E\longrightarrow E'$ be a morphism of the underlying
rigidified bitorsors. We show that $f$ is a morphism of extensions, or,
what plainly amounts to the same thing, a homomorphism of groups. We
must prove $f(xy)=f(x)f(y)$. If $j':E'\longrightarrow P$ denotes the
structural morphism of the torsor $E'$ over $P$, then
$j'f(xy)=j'(f(x)f(y))$: the two sides are respectively $j(xy)$ and
$j(x)j(y)$, which are equal because $j$ is multiplicative. Hence
\begin{equation*}
 f(xy)=u(j(x),j(y))f(x)f(y),
\end{equation*}
where
\begin{equation*}
 u:P\times P\longrightarrow G
\end{equation*}
is a uniquely determined morphism. Setting first $x=1$ and then $y=1$
in the preceding relation, and using $f(1)=1$, gives
\begin{equation*}
 u(p,1)=u(1,p)=1.
\end{equation*}
Thus $u$ is trivial on $P\times e_P$ and on $e_P\times P$; by
hypothesis it is trivial, and therefore $f(xy)=f(x)f(y)$, as required.

For b), necessity of the stated condition for $E$ to admit an
extension structure compatible with its bitorsor structure is immediate
from \SGARefLink{sga7.sga7i.evii.n1-3-4-3}{1.3.4.3}. It remains to prove sufficiency for the first assertion.
Pulling a trivialization $\varphi$ of $\delta(E)$ back along the
identity section of $P\times P$ gives a trivialization of
$E_1=e_P^*(E)$ (compare 1.5.1), that is, a rigidification of $E$
relative to $e_P$. Fix this rigidification. A given trivialization
$\varphi$ of $\delta(E)$ need not be compatible with the corresponding
birigidification. Relative to the natural partial rigidifications of
$\delta(E)$, its restrictions to $P\times e_P$ and $e_P\times P$ are
described by central morphisms
\begin{equation*}
 f:P\longrightarrow G,qquad g:P\longrightarrow G,
\end{equation*}
whose values at the identity section of $P$ coincide:
$f(1)=g(1)=a\in\Gamma G$. The central morphism
\begin{equation*}
 h:P\times P\longrightarrow G,qquad h(p,q)=f(p)g(q)a^{-1},
\end{equation*}
has restrictions $f$ and $g$ to $P\times e_P$ and $e_P\times P$.
Correcting $\varphi$ by $h$, that is, replacing it by
$\varphi(p,q)h(p,q)^{-1}$, therefore gives a trivialization of
$\delta(E)$ compatible with its birigidification. Such a trivialization
is unique: two of them differ by a central morphism
$P\times P\longrightarrow G$ trivial on $P\times e_P$ and
$e_P\times P$, hence trivial by hypothesis. This proves the last
assertion of \SGARefLink{sga7.sga7i.evii.n1-3-5}{1.3.5} b).

It remains to prove that $\varphi$, viewed as an isomorphism of the form
(\SGARefLink{sga7.sga7i.evii.n1-1-3-2}{1.1.3.2}), satisfies the associativity condition (\SGARefLink{sga7.sga7i.evii.n1-1-4-1}{1.1.4.1}); by our
choice of $\varphi$, $e_E$ is already a two-sided identity for the
composition law defined by $\varphi$. The two composites in (\SGARefLink{sga7.sga7i.evii.n1-1-4-1}{1.1.4.1})
are isomorphisms of bitorsors under
$(G_{P\times P\times P},G_{P\times P\times P})$, so they differ by a
central morphism
\begin{equation*}
 u:P\times P\times P\longrightarrow G
\end{equation*}
through the formula
\begin{equation*}
 (xy)z=u(jx,jy,jz)\,x(yz).
\end{equation*}

%%%%%%%%%%  scan idx 162  |  folio 151  %%%%%%%%%%

Setting successively $z=1$, $x=1$, and $y=1$, and using the fact that
$1$ is a two-sided identity for the composition law, shows that $u$ is
trivial on the three partial products $P\times P\times e_P$,
$e_P\times P\times P$, and $P\times e_P\times P$. The hypothesis then
implies that $u$ is trivial, so $(xy)z=x(yz)$. This completes the proof.


% END INCLUDED SGA7I_English_current_buildable_source_20260823/components/69_expose_VII_proposition_1_3_5.tex

% BEGIN INCLUDED SGA7I_English_current_buildable_source_20260823/components/70_expose_VII_1_3_6_through_1_3_7.tex
% SOURCE AUTHORITY:
% expose_VII_body.tex lines 360--398; scan indices 162--164; source folios 151--153.
% Source-backed correction: the action in 1.3.6 is defined by the extension E, not G.
% Hard stop before source line 400, paragraph 1.3.8.

\underline{Corollary \SGARefTarget{sga7.sga7i.evii.n1-3-6}{1.3.6}. Suppose that every morphism of sheaves of
sets $P\times G\longrightarrow G$ which is trivial on $P\times e_G$
and on $e_P\times G$ is trivial. Then, if $G$ is commutative, $G$ is
central in every extension of $P$ by $G$. If $P$ is also commutative
and the hypothesis of \SGARefLink{sga7.sga7i.evii.n1-3-5}{1.3.5} a) holds, then every extension $E$ of $P$
by $G$ is commutative.}

For the first assertion, we must prove that the action of $P$ on $G$
defined by the extension $E$ is trivial. This action is expressed by a
morphism $u:P\times G\longrightarrow G$. The restrictions to
$P\times e_G$ and $e_P\times G$ of the corresponding morphism
\[
 (p,g)\longmapsto u(p,g)g^{-1}
\]
are plainly trivial. Hence this last morphism is trivial, so
$u(p,g)=g$, as required. Now, if $E$ is an extension of $P$ by $G$, its
opposite group $E^{\mathrm o}$ is an extension of $P^{\mathrm o}$ by
$G^{\mathrm o}$, and the bitorsor structure underlying
$i(E^{\mathrm o})$ is obtained from $i(E)$ by interchanging the left
and right actions. When $E$ is a central extension, one therefore has
$i(E)=i(E^{\mathrm o})$; if, in addition, $P$ is commutative, so that
$P=P^{\mathrm o}$, the uniqueness assertion \SGARefLink{sga7.sga7i.evii.n1-3-5}{1.3.5} a) gives
$E=E^{\mathrm o}$. Thus $E$ is commutative, proving the corollary.

%%%%%%%%%%  scan idx 163  |  folio 152  %%%%%%%%%%

\SGARefTarget{sga7.sga7i.evii.n1-3-7}{1.3.7}. One interesting special case in which all the hypotheses of
\SGARefLink{sga7.sga7i.evii.n1-3-5}{1.3.5} and \SGARefLink{sga7.sga7i.evii.n1-3-6}{1.3.6} hold is
\begin{equation}\tag{\SGARefTarget{sga7.sga7i.evii.n1-3-7-1}{1.3.7.1}}
 \underline{\operatorname{Hom}}_{\mathrm{pt}}(P,G)=e
 \quad\text{(the final sheaf)},
\end{equation}
where the left-hand side denotes the sheaf of morphisms of pointed
sheaves of sets from $P$ to $G$. Independently of the marked sections,
this condition may also be written in the form, analogous to (\SGARefLink{sga7.sga7i.evii.n1-3-1-2}{1.3.1.2}),
\begin{equation}\tag{\SGARefTarget{sga7.sga7i.evii.n1-3-7-2}{1.3.7.2}}
 G\simeq\underline{\operatorname{Hom}}_{\mathrm{sets}}(P,G),
\end{equation}
where the second member denotes the sheaf of morphisms between the
underlying sheaves of sets of $P$ and $G$. These conditions also mean
that, for every object $Q$ of $\underline T$, the map
$u\longmapsto u\circ\pr_2$
\begin{equation}\tag{\SGARefTarget{sga7.sga7i.evii.n1-3-7-3}{1.3.7.3}}
 \operatorname{Hom}_{\mathrm{sets}}(Q,G)
 \xrightarrow{\ \sim\ }
 \operatorname{Hom}_{\mathrm{sets}}(P\times Q,G)
\end{equation}
is bijective. This property of $P$ (with $G$ fixed) is stable under
Cartesian products; in particular it implies the same property for
$P\times P$, $P\times P\times P$, and so forth. It follows immediately
that it implies the hypotheses made in \SGARefLink{sga7.sga7i.evii.n1-3-5}{1.3.5} and \SGARefLink{sga7.sga7i.evii.n1-3-6}{1.3.6}. It also implies
that all the categories occurring in (\SGARefLink{sga7.sga7i.evii.n1-3-4-1}{1.3.4.1}) are \underline{rigid},
that is, every automorphism of an object of any one of them is the
identity.

Moreover, since condition (\SGARefLink{sga7.sga7i.evii.n1-3-7-1}{1.3.7.1}) on $P,G$ is stable under every base
change $S\longrightarrow e$, descent together with the uniqueness
result \SGARefLink{sga7.sga7i.evii.n1-3-5}{1.3.5} a) shows that a rigidified bitorsor $E$ over $P$ comes from
an extension of $P$ by $G$ as soon as it does so \underline{locally over
$e$}. By \SGARefLink{sga7.sga7i.evii.n1-3-5}{1.3.5} b), this may also be expressed by the vanishing of the
section
\begin{equation*}
 \underline{\delta}(E)\in\Gamma R^1h_*(G_{P\times P})
\end{equation*}
defined by $\delta(E)$, where
\begin{equation*}
 h:P\times P\longrightarrow e
\end{equation*}
is the structural morphism.


% END INCLUDED SGA7I_English_current_buildable_source_20260823/components/70_expose_VII_1_3_6_through_1_3_7.tex

% BEGIN INCLUDED SGA7I_English_current_buildable_source_20260823/components/71_expose_VII_1_3_8.tex
% SOURCE AUTHORITY:
% expose_VII_body.tex lines 400--418; scan indices 164--165; source folios 153--154.
% Source-backed correction: final references 3.5 and 3.6 are 1.3.5 and 1.3.6.
% Hard stop before source line 421, section 1.4.

\SGARefTarget{sga7.sga7i.evii.n1-3-8}{1.3.8}. Let $p:P\longrightarrow S$ be a morphism of schemes which is
\underline{coherent} (that is, quasi-compact and quasi-separated) and
\underline{flat}, and suppose that
\begin{equation}\tag{\SGARefTarget{sga7.sga7i.evii.n1-3-8-1}{1.3.8.1}}
 p_*(\mathcal O_P)\xleftarrow{\ \sim\ }\mathcal O_S.
\end{equation}
Since the formation of $p_*(\mathcal O_P)$ commutes with every flat base
change (by coherence of $p$), the stated conditions on the relative
scheme $P/S$ are stable under Cartesian products and hence hold, in
particular, for the Cartesian powers $P\times P$,
$P\times P\times P$, and so forth. They are also stable under every flat
base change.

Now let $G$ be an \underline{affine} scheme over $S$. Relation
(\SGARefLink{sga7.sga7i.evii.n1-3-7-2}{1.3.7.2}) follows at once from (\SGARefLink{sga7.sga7i.evii.n1-3-8-1}{1.3.8.1}), since
\[
 \operatorname{Hom}_S(P,G)
 =\operatorname{Hom}_{\mathrm{Alg}}(\underline A,p_*(\mathcal O_P))
 \simeq\operatorname{Hom}_{\mathrm{Alg}}(\underline A,\mathcal O_S),
\]
where $\underline A$ is the quasi-coherent algebra over $S$ such that
$G=\operatorname{Spec}(\underline A)$. By the preceding observation,
the same conclusion applies to $P\times P$, $P\times P\times P$, and
so on. Likewise, if $G$ is flat over $S$, the same conclusion applies
to the relative scheme $P_G$ over $G$ with respect to the relative
scheme $G_G$ over $G$.

When $P$ and $G$ are group schemes over $S$, and one works with a
topology on $(\mathrm{Sch})_{/S}$ coarser than the canonical topology,
so that $P$ and $G$ may be identified with sheaves on this site, every
morphism from $P$, $P\times P$, $P\times P\times P$, and so on to $G$
which is compatible with the marked sections is trivial. Similarly, if
$G$ is flat, every $S$-morphism $P\times G\longrightarrow G$ compatible
with the marked sections is trivial, as one sees by interpreting it as
a $G$-morphism from $P_G$ to $G_G$. Thus the hypotheses of \SGARefLink{sga7.sga7i.evii.n1-3-5}{1.3.5} hold,
as do those of \SGARefLink{sga7.sga7i.evii.n1-3-6}{1.3.6} when $G$ is flat. Moreover, the three categories
occurring in (\SGARefLink{sga7.sga7i.evii.n1-3-4-1}{1.3.4.1}) are rigid.

%%%%%%%%%%  scan idx 165  |  folio 154  %%%%%%%%%%

\SGARefTarget{sga7.sga7i.evii.n1-3-8-2}{1.3.8.2}. The most interesting case in which (\SGARefLink{sga7.sga7i.evii.n1-3-8-1}{1.3.8.1}) holds, and
continues to hold after every extension of the base, is that in which
$P$ is an \underline{abelian scheme} over $S$; then the conditions of
\SGARefLink{sga7.sga7i.evii.n1-3-7}{1.3.7} themselves hold. In this case the conclusions of \SGARefLink{sga7.sga7i.evii.n1-3-5}{1.3.5} and \SGARefLink{sga7.sga7i.evii.n1-3-6}{1.3.6}
are therefore valid for every affine group $G$ over $S$. This is the
case originally considered by J.~P. Serre, and the preceding
developments were plainly inspired by it. An interesting case in which
(\SGARefLink{sga7.sga7i.evii.n1-3-8-1}{1.3.8.1}) holds without holding universally is that in which $P$ is the
\underline{N\'eron model} of an abelian scheme over the fraction field
of a discrete valuation ring with spectrum $S$.


% END INCLUDED SGA7I_English_current_buildable_source_20260823/components/71_expose_VII_1_3_8.tex

% BEGIN INCLUDED SGA7I_English_current_buildable_source_20260823/components/72_expose_VII_section_1_4.tex
% SOURCE AUTHORITY:
% expose_VII_body.tex lines 421--460; scan indices 165--167; source folios 154--156.
% Hard stop before source line 465, section 2.

\SGARefTarget{sga7.sga7i.evii.n1-4}{1.4}. \underline{Extensions of a free commutative group}

Let $I$ be an object of $\underline T$, and put
\begin{equation}\tag{\SGARefTarget{sga7.sga7i.evii.n1-4-1}{1.4.1}}
 P=\mathbb Z[I],
\end{equation}
the free $\mathbb Z$-module generated by $I$ (SGA 4, IV, 11). With
every extension $E$ of $P$ by a given $\mathbb Z$-module $G$ (from now
on only commutative group extensions will occur), associate the
corresponding torsor under $G_P$, and then restrict it to $I$ along the
canonical morphism
\[
 I\longrightarrow P=\mathbb Z[I].
\]

%%%%%%%%%%  scan idx 166  |  folio 155  %%%%%%%%%%

This gives a canonical functor
\begin{equation}\tag{\SGARefTarget{sga7.sga7i.evii.n1-4-2}{1.4.2}}
 \operatorname{EXT}(P,G)\longrightarrow
 \operatorname{TORS}(I,G_I).
\end{equation}
I claim that this functor is an \underline{equivalence of categories}.
Equivalently, for every $G$, the maps
\begin{equation}\tag{\SGARefTarget{sga7.sga7i.evii.n1-4-3}{1.4.3}}
 \operatorname{Hom}_{\mathbb Z}(P,G)\longrightarrow
 H^0(I,G_I)=G(I)
\end{equation}
and
\begin{equation}\tag{\SGARefTarget{sga7.sga7i.evii.n1-4-4}{1.4.4}}
 \operatorname{Ext}^1_{\mathbb Z}(P,G)\longrightarrow H^1(I,G_I)
\end{equation}
are isomorphisms. More generally, if $A$ is a ring (not necessarily
commutative) of $\underline T$ and $P=A[I]$ is the free $A$-module
generated by $I$, there are canonical isomorphisms, functorial in the
variable $A$-module $G$,
\begin{equation}\tag{\SGARefTarget{sga7.sga7i.evii.n1-4-5}{1.4.5}}
 \operatorname{Ext}^i_A(A[I],G)\xrightarrow{\ \sim\ }H^i(I,G_I).
\end{equation}
For $A=\mathbb Z$ and $i=0,1$, these coincide with (\SGARefLink{sga7.sga7i.evii.n1-4-3}{1.4.3}) and
(\SGARefLink{sga7.sga7i.evii.n1-4-4}{1.4.4}). To define (\SGARefLink{sga7.sga7i.evii.n1-4-5}{1.4.5}), recall that its left-hand side is the
derived functor of
\[
 G\longmapsto\operatorname{Hom}_A(A[I],G),
\]
which, by the definition of $A[I]$, is isomorphic to
$G\longmapsto G(I)=H^0(I,G_I)$. Since $G\longmapsto G_I$ carries
injective modules to injective modules of $\underline T_{/I}$ (SGA 4,
V), the assertion follows. We leave to the reader the verification of
the stated compatibility with (\SGARefLink{sga7.sga7i.evii.n1-4-3}{1.4.3}) and (\SGARefLink{sga7.sga7i.evii.n1-4-4}{1.4.4}).

Finally, the functor (\SGARefLink{sga7.sga7i.evii.n1-4-2}{1.4.2}) is compatible with inverse images under a
morphism of topoi
\[
 f:\underline T'\longrightarrow\underline T.
\]


% END INCLUDED SGA7I_English_current_buildable_source_20260823/components/72_expose_VII_section_1_4.tex

% BEGIN INCLUDED SGA7I_English_current_buildable_source_20260823/components/73_expose_VII_section_2_0_phi.tex
% SOURCE AUTHORITY:
% expose_VII_body.tex lines 465--540; scan indices 167--169; source folios 156--158.
% Source-backed correction: the printed prose interchanges the P- and Q-direction extension names;
% the displayed phi laws unambiguously multiply the P-coordinate, hence define P_Q by G_Q.
% The missing prime in the upper-left label of (2.0.5) is restored from its nodes and associativity.
% Hard stop before the psi-law continuation on source line 540.

2. \underline{The notion of a biextension of abelian sheaves}

\SGARefTarget{sga7.sga7i.evii.n2-0}{2.0}. Throughout this section, $P,Q,G$ denote commutative groups of the
topos $\underline T$. We shall frequently use the Cartesian diagram
\begin{equation}\tag{\SGARefTarget{sga7.sga7i.evii.n2-0-1}{2.0.1}}
\begin{tikzcd}[column sep=4em,row sep=2.6em]
 P \arrow[d] & P\times Q \arrow[l] \arrow[d] \\
 e & Q \arrow[l]
\end{tikzcd}
\end{equation}
and the groups $G_P$ (respectively $G_Q$, $G_{P\times Q}$) over $P$
(respectively $Q$, $P\times Q$), obtained by inverse image from the
group $G$ over $e$. We identify $G_{P\times Q}$ either with the inverse
image $(G_P)_{P\times Q}$ under $\pr_1$, or with the inverse image
$(G_Q)_{P\times Q}$ under $\pr_2$. In this notation the group
structures of $P,Q,P\times Q$ play no role; indeed, we shall almost
never use the product-group structure on $P\times Q$. We shall instead
use on $P\times Q$ the structure of a group over $P$ obtained by inverse
image from $Q$ over $e$,
\begin{equation}\tag{\SGARefTarget{sga7.sga7i.evii.n2-0-2}{2.0.2}}
 P\times Q=Q_P,
\end{equation}
and, similarly, the structure of a group over $Q$ obtained by inverse
image from $P$ over $e$,
\begin{equation}\tag{\SGARefTarget{sga7.sga7i.evii.n2-0-3}{2.0.3}}
 P\times Q=P_Q.
\end{equation}
Thus over $P$ we have the two groups $G_P,Q_P$, and over $Q$ the two
groups $G_Q,P_Q$.

If $E$ is a torsor over $P\times Q$ under $G_{P\times Q}$, it may be
viewed either as a torsor over $Q_P$ under
$G_{P\times Q}=(G_P)_{Q_P}$, or as a torsor over $P_Q$ under
$G_{P\times Q}=(G_Q)_{P_Q}$. It therefore makes sense to equip $E$ with
a commutative extension structure in either variable. Applied in the
induced topoi $\underline T_{/P}$ and $\underline T_{/Q}$, \SGARefLink{sga7.sga7i.evii.n1-1-6}{1.1.6} gives
a convenient explicit description, using the canonical isomorphisms
\[
 \underline T_{/P\times Q}\simeq
 (\underline T_{/P})_{/P\times Q},\qquad
 \underline T_{/P\times Q}\simeq
 (\underline T_{/Q})_{/P\times Q}.
\]

An extension structure in the $P$-direction---that is, of $P_Q$ by
$G_Q$---is equivalent to giving, for every object $S$ of $\underline T$
and elements $p,p'\in P(S)$, $q\in Q(S)$, an isomorphism of torsors
under $G_S$
\begin{equation}\tag{\SGARefTarget{sga7.sga7i.evii.n2-0-4}{2.0.4}}
 \varphi_{p,p';q}:E_{p,q}E_{p',q}
 \xrightarrow{\ \sim\ }E_{pp',q}.
\end{equation}
As usual (cf. \SGARefLink{sga7.sga7i.evii.n1-1-2-1}{1.1.2.1}), for $(p,q)\in(P\times Q)(S)$, $E_{p,q}$ denotes
the torsor under $G_S$ obtained from $E$ by inverse image along
$(p,q):S\longrightarrow P\times Q$. The data (\SGARefLink{sga7.sga7i.evii.n2-0-4}{2.0.4}) must be
\underline{functorial in $S$}, associative, and commutative. The
associativity condition is the commutativity of
\begin{equation}\tag{\SGARefTarget{sga7.sga7i.evii.n2-0-5}{2.0.5}}
\begin{tikzcd}[column sep=5.1em,row sep=3.4em]
 & E_{pp',q}E_{p'',q}
   \arrow[rd,"{\varphi_{pp',p'';q}}"] & \\
 E_{p,q}E_{p',q}E_{p'',q}
   \arrow[ru,"{\varphi_{p,p';q}\wedge\mathrm{id}}"]
   \arrow[rd,"{\mathrm{id}\wedge\varphi_{p',p'';q}}"']
 && E_{pp'p'',q} \\
 & E_{p,q}E_{p'p'',q}
   \arrow[ru,"{\varphi_{p,p'p'';q}}"'] &
\end{tikzcd}
\end{equation}
for $p,p',p''\in P(S)$ and $q\in Q(S)$. Commutativity is expressed by
\begin{equation}\tag{\SGARefTarget{sga7.sga7i.evii.n2-0-6}{2.0.6}}
\begin{tikzcd}[column sep=5em,row sep=3.1em]
 E_{p,q}E_{p',q} \arrow[r,"{\varphi_{p,p';q}}"]
   \arrow[d,"\mathrm{sym}"']
 & E_{pp',q} \arrow[d,"\mathrm{id}"] \\
 E_{p',q}E_{p,q} \arrow[r,"{\varphi_{p',p;q}}"']
 & E_{p'p,q}.
\end{tikzcd}
\end{equation}
Here $p,p'\in P(S)$ and $q\in Q(S)$.


% END INCLUDED SGA7I_English_current_buildable_source_20260823/components/73_expose_VII_section_2_0_phi.tex

% BEGIN INCLUDED SGA7I_English_current_buildable_source_20260823/components/74_expose_VII_psi_and_definition_2_1.tex
% SOURCE AUTHORITY:
% expose_VII_body.tex lines 540--658; scan indices 169--171; source folios 158--160.
% Source-backed corrections: (2.0.7) targets E_{p,qq'}, not the impossible E_{pp',q};
% the second structure is Q_P by G_P, not the printed reversal G_Q by P_Q.
% Hard stop before source line 660, section 2.2.

Symmetrically, an extension structure in the $Q$-direction---that is,
of $Q_P$ by $G_P$---is described by isomorphisms of torsors under $G_S$
\begin{equation}\tag{\SGARefTarget{sga7.sga7i.evii.n2-0-7}{2.0.7}}
 \psi_{p;q,q'}:E_{p,q}E_{p,q'}
 \xrightarrow{\ \sim\ }E_{p,qq'},
\end{equation}
for $p\in P(S)$ and $q,q'\in Q(S)$. These are functorial in $S$ and
satisfy the symmetric associativity and commutativity conditions
\begin{equation}\tag{\SGARefTarget{sga7.sga7i.evii.n2-0-8}{2.0.8}}
\begin{tikzcd}[column sep=5.1em,row sep=3.4em]
 & E_{p,qq'}E_{p,q''}
   \arrow[rd,"{\psi_{p;qq',q''}}"] & \\
 E_{p,q}E_{p,q'}E_{p,q''}
   \arrow[ru,"{\psi_{p;q,q'}\wedge\mathrm{id}}"]
   \arrow[rd,"{\mathrm{id}\wedge\psi_{p;q',q''}}"']
 && E_{p,qq'q''} \\
 & E_{p,q}E_{p,q'q''}
   \arrow[ru,"{\psi_{p;q,q'q''}}"'] &
\end{tikzcd}
\end{equation}
for $p\in P(S)$ and $q,q',q''\in Q(S)$, and
\begin{equation}\tag{\SGARefTarget{sga7.sga7i.evii.n2-0-9}{2.0.9}}
\begin{tikzcd}[column sep=5em,row sep=3.1em]
 E_{p,q}E_{p,q'} \arrow[r,"{\psi_{p;q,q'}}"]
   \arrow[d,"\mathrm{sym}"']
 & E_{p,qq'} \arrow[d,"\mathrm{id}"] \\
 E_{p,q'}E_{p,q} \arrow[r,"{\psi_{p;q',q}}"']
 & E_{p,q'q}.
\end{tikzcd}
\end{equation}
Here $p\in P(S)$ and $q,q'\in Q(S)$.

%%%%%%%%%%  scan idx 170  |  folio 159  %%%%%%%%%%

\underline{Definition \SGARefTarget{sga7.sga7i.evii.n2-1}{2.1}. Let $P,Q,G$ be abelian groups of
$\underline T$, and let $E$ be a torsor under $G_{P\times Q}$ equipped
with an extension structure in the $P$-direction, defined by the
isomorphisms (\SGARefLink{sga7.sga7i.evii.n2-0-4}{2.0.4}), and an extension structure in the $Q$-direction,
defined by the isomorphisms (\SGARefLink{sga7.sga7i.evii.n2-0-7}{2.0.7}), each functorial in $S$ and subject
to the associativity and commutativity conditions above. The two
structures are called compatible if, for every object $S$ of
$\underline T$ and $p,p'\in P(S)$, $q,q'\in Q(S)$, the following
diagram commutes:}
\begin{equation}\tag{\SGARefTarget{sga7.sga7i.evii.n2-1-1}{2.1.1}}
\begin{tikzcd}[column sep=2.7em,row sep=3.2em,font=\scriptsize]
 && E_{p,qq'}E_{p',qq'}
    \arrow[rrdd,"{\varphi_{p,p';qq'}}"] && \\
 E_{p,q}E_{p,q'}E_{p',q}E_{p',q'}
    \arrow[rru,"{\psi_{p;q,q'}\wedge\psi_{p';q,q'}}"]
    \arrow[d,"{}"'] &&&& \\
 E_{p,q}E_{p',q}E_{p,q'}E_{p',q'}
    \arrow[rrd,"{\varphi_{p,p';q}\wedge\varphi_{p,p';q'}}"']
 &&&& E_{pp',qq'} \\
 && E_{pp',q}E_{pp',q'}
    \arrow[rru,"{\psi_{pp';q,q'}}"'] &&
\end{tikzcd}
\end{equation}
\underline{where the vertical arrow is induced by the symmetry
isomorphism}
\[
 E_{p,q'}E_{p',q}\xrightarrow{\ \sim\ }E_{p',q}E_{p,q'}.
\]
\underline{A \emph{biextension} of $(P,Q)$ by $G$ is a torsor $E$
under $G_{P\times Q}$ equipped with compatible extension structures in
the $P$- and $Q$-directions in the preceding sense.}

%%%%%%%%%%  scan idx 171  |  folio 160  %%%%%%%%%%

\SGARefLink{sga7.sga7i.evii.n2-1-1}{2.1.1}. Equivalently, a biextension structure on the
$G_{P\times Q}$-torsor $E$ is given by two partially defined composition
laws $\varphi(x,y)$ and $\psi(x,y)$ on $E$. The first is defined when
$\pr_2j(x)=\pr_2j(y)$, and the second when
$\pr_1j(x)=\pr_1j(y)$, where $j:E\longrightarrow P\times Q$ is the
structural morphism of the torsor. Besides associativity relations of
the type
\begin{equation}\tag{\SGARefTarget{sga7.sga7i.evii.n2-1-1-1}{2.1.1.1}}
 \varphi(gx,y)=\varphi(x,gy)=g\varphi(x,y)
 \qquad(x,y\in E(S),\ g\in G(S))
\end{equation}
and the analogous relations for $\psi$, the axioms say that these
partial composition laws are associative and commutative and satisfy
\begin{equation}\tag{\SGARefTarget{sga7.sga7i.evii.n2-1-1-2}{2.1.1.2}}
 \varphi(\psi(x,y),\psi(z,t))
 =\psi(\varphi(x,z),\varphi(y,t)),
\end{equation}
where $x,y,z,t$ lie respectively in
$E_{p,q}(S),E_{p,q'}(S),E_{p',q}(S),E_{p',q'}(S)$.


% END INCLUDED SGA7I_English_current_buildable_source_20260823/components/74_expose_VII_psi_and_definition_2_1.tex

% BEGIN INCLUDED SGA7I_English_current_buildable_source_20260823/components/75_expose_VII_section_2_2.tex
% SOURCE AUTHORITY:
% expose_VII_body.tex lines 660--683; scan indices 171--173; source folios 160--162.
% Source-backed corrections: e_{E/P} trivializes over P x 1_Q and e_{E/Q} over 1_P x Q;
% the final restricted torsor is an inverse image in E, not in G.
% Hard stop before source line 688, section 2.3.

\SGARefTarget{sga7.sga7i.evii.n2-2}{2.2}. \underline{Identity sections}. Consider first the identity section
of $E$ regarded as an extension of $Q_P$ by $G_P$,
\begin{equation}\tag{\SGARefTarget{sga7.sga7i.evii.n2-2-1}{2.2.1}}
 e_{E/P}\in\Gamma(E/P)=\operatorname{Hom}_P(P,E),
\end{equation}
and, symmetrically, the identity section of $E$ regarded as an extension
of $P_Q$ by $G_Q$,
\begin{equation}\tag{\SGARefTarget{sga7.sga7i.evii.n2-2-2}{2.2.2}}
 e_{E/Q}\in\Gamma(E/Q)=\operatorname{Hom}_Q(Q,E).
\end{equation}
In particular, \underline{the torsor $E$ over $P\times Q$ under
$G_{P\times Q}$ is canonically trivialized over $P\times1_Q$ by
$e_{E/P}$, and over $1_P\times Q$ by $e_{E/Q}$. These trivializations
coincide over the identity section $1_P\times1_Q$ of $P\times Q$}, and
hence define the same section
\begin{equation}\tag{\SGARefTarget{sga7.sga7i.evii.n2-2-3}{2.2.3}}
 e_E\in\Gamma(E),
\end{equation}
which, by abuse of language, will be called the \underline{identity
section} (or \underline{neutral section}) of $E$.

%%%%%%%%%%  scan idx 172  |  folio 161  %%%%%%%%%%

To prove the asserted equality, let $e,e'$ be the two sections of $E$
deduced respectively from $e_{E/P}$ and $e_{E/Q}$. They are
characterized as the identity sections in the inverse image
$E_{1_P,1_Q}$ of $E$ along the identity section of $P\times Q$, for the
two composition laws $\varphi$ and $\psi$, respectively. Applying
(\SGARefLink{sga7.sga7i.evii.n2-1-1-2}{2.1.1.2}) to the four sections $e,e',e',e$ of $E_{1,1}$ gives $e$ on
the left and $e'$ on the right; hence $e=e'$. It follows moreover that
\underline{the two composition laws $\varphi$ and $\psi$ on $E_{1,1}$
are identical}. Indeed, under transport of structure by the morphisms
$g\mapsto g.e$ and $g\mapsto g.e'$, respectively, both correspond to
the group law of $G$, and $e=e'$.

The extension of $Q_P$ by $G_P$ over the base $P$, restricted over the
identity section $1_P$ of $P$, gives an extension of $Q$ by $G$ on the
inverse image of $1_P\times Q$ in $E$. The homomorphism $e_{E/Q}$ is
immediately seen to be multiplicative for this extension structure, so
the extension is canonically split. Symmetrically, the extension of $P$
by $G$ obtained by restricting the extension of $P_Q$ by $G_Q$ over
$Q$ to the identity section $1_Q$ is canonically split by $e_{E/P}$.


% END INCLUDED SGA7I_English_current_buildable_source_20260823/components/75_expose_VII_section_2_2.tex

% BEGIN INCLUDED SGA7I_English_current_buildable_source_20260823/components/76_expose_VII_section_2_3.tex
% SOURCE AUTHORITY:
% expose_VII_body.tex lines 688--725; scan indices 173--174; source folios 162--163.
% The relative homomorphisms are retained exactly with the two bases made explicit.
% Hard stop before source line 728, section 2.4.

\SGARefTarget{sga7.sga7i.evii.n2-3}{2.3}. \underline{Morphisms of biextensions}. Let $E$ be a biextension
of $(P,Q)$ by $G$, and let $E'$ likewise be a biextension of
$(P',Q')$ by $G'$, where $P',Q',G'$ are again commutative groups of
$\underline T$. A \underline{morphism of biextensions} from $E$ to
$E'$ is a system $(u,v,w,f)$, where
\begin{equation}\tag{\SGARefTarget{sga7.sga7i.evii.n2-3-1}{2.3.1}}
 u:P\longrightarrow P',\qquad
 v:Q\longrightarrow Q',\qquad
 w:G\longrightarrow G'
\end{equation}
are morphisms of groups, and
\begin{equation}\tag{\SGARefTarget{sga7.sga7i.evii.n2-3-2}{2.3.2}}
 f:E\longrightarrow E'
\end{equation}
is a morphism of the sheaves of sets underlying $E$ and $E'$, such
that $f$ is simultaneously a homomorphism of group extensions over
$P$ and $P'$, respectively, associated with the relative group
homomorphisms
\[
 u\times v:Q_P\longrightarrow Q'_{P'},
 \qquad
 u\times w:G_P\longrightarrow G'_{P'},
\]
and a homomorphism of group extensions over $Q$ and $Q'$, respectively,
associated with the relative group homomorphisms
\[
 u\times v:P_Q\longrightarrow P'_{Q'},
 \qquad
 v\times w:G_Q\longrightarrow G'_{Q'}.
\]
Notice moreover that $u,v,w$ are determined by $f$ alone. Indeed,
passing to the quotient recovers $u\times v$ from $f$, and then
recovers $u$ and $v$ by restriction to $P\times1_Q$ and
$1_P\times Q$, respectively; restricting $f$ to $G.e_E$ recovers
$w$. Thus no ambiguity arises from identifying the morphism
$(u,v,w,f)$ with $f$.

After taking the usual precaution of adjoining their sources and
targets to morphisms (the canonical genuflection), one sees that
\underline{the biextensions of the topos $\underline T$ are the
objects of a category}, whose objects and morphisms are those defined
in \SGARefLink{sga7.sga7i.evii.n2-1}{2.1} and \SGARefLink{sga7.sga7i.evii.n2-3}{2.3} and whose composition of morphisms is defined in the
evident way. Moreover, $(P,Q,G,E)\mapsto(P,Q,G)$ gives a canonical
functor
\begin{equation}\tag{\SGARefTarget{sga7.sga7i.evii.n2-3-3}{2.3.3}}
 \operatorname{BIEXT}(\underline T)
 \longrightarrow
 \operatorname{GRAB}(\underline T)\times
 \operatorname{GRAB}(\underline T)\times
 \operatorname{GRAB}(\underline T),
\end{equation}
where $\operatorname{BIEXT}(\underline T)$ denotes the category of
biextensions over $\underline T$, and
$\operatorname{GRAB}(\underline T)$ the category of abelian groups
over $\underline T$. It is useful to decompose this canonical functor
into the two functors $E\mapsto(P,Q)$,
\begin{equation}\tag{\SGARefTarget{sga7.sga7i.evii.n2-3-4}{2.3.4}}
 \operatorname{BIEXT}(\underline T)
 \longrightarrow
 \operatorname{GRAB}(\underline T)\times
 \operatorname{GRAB}(\underline T),
\end{equation}
and $E\mapsto G$,
\begin{equation}\tag{\SGARefTarget{sga7.sga7i.evii.n2-3-5}{2.3.5}}
 \operatorname{BIEXT}(\underline T)
 \longrightarrow
 \operatorname{GRAB}(\underline T).
\end{equation}


% END INCLUDED SGA7I_English_current_buildable_source_20260823/components/76_expose_VII_section_2_3.tex

% BEGIN INCLUDED SGA7I_English_current_buildable_source_20260823/components/77_expose_VII_section_2_4.tex
% SOURCE AUTHORITY:
% expose_VII_body.tex lines 728--773; scan indices 174--176; source folios 163--165.
% Source-backed corrections: the covariant construction begins with a biextension of
% (P,Q), not the printed (P,G); (2.4.2) has BIEXT(P,Q;G x G'), not BIEXT(P,G;...).
% The scanned source confirms the stray <<CONF marker is not mathematical content.
% Hard stop before source line 776, section 2.5.

\SGARefTarget{sga7.sga7i.evii.n2-4}{2.4}. \underline{The category $\operatorname{BIEXT}(P,Q;G)$ for
variable $P,Q,G$}.

I claim that \underline{the canonical functor (\SGARefLink{sga7.sga7i.evii.n2-3-4}{2.3.4}) is fibred}, and
that \underline{the canonical functor (\SGARefLink{sga7.sga7i.evii.n2-3-5}{2.3.5}) is cofibred}. The first
assertion means, in particular, that if $E$ is a biextension of
$(P,Q)$ by $G$, and if
\[
 u:P'\longrightarrow P,\qquad v:Q'\longrightarrow Q
\]
are morphisms of abelian groups, then there is a biextension $E'$ of
$(P',Q')$ by $G$ and a morphism of biextensions
\[
 f:E'\longrightarrow E
\]
compatible with $u,v,\mathrm{id}_G$, having the usual universal
property for morphisms of biextensions into $E$ compatible with
$u,v,\mathrm{id}_G$.

The second assertion is expressed similarly. In particular, if $E$ is
a biextension of $(P,Q)$ by $G$, and if
\[
 w:G\longrightarrow G'
\]
is a morphism of abelian groups, then there is a biextension $E'$ of
$(P,Q)$ by $G'$ and a morphism of biextensions
\[
 g:E\longrightarrow E'
\]
compatible with $\mathrm{id}_P,\mathrm{id}_Q,w$, having the usual
universal property for morphisms of biextensions from $E$ compatible
with $\mathrm{id}_P,\mathrm{id}_Q,w$. We leave it to the reader to
spell out the two constructions whose existence has just been
asserted; they are obtained easily by paraphrasing the familiar
analogous constructions for extensions of ordinary groups.

Thus, for three abelian groups $P,Q,G$ of $\underline T$, write
\begin{equation}\tag{\SGARefTarget{sga7.sga7i.evii.n2-4-1}{2.4.1}}
 \operatorname{BIEXT}(P,Q;G)
\end{equation}
for the category of biextensions of $(P,Q)$ by $G$, defined as the
fibre category of (\SGARefLink{sga7.sga7i.evii.n2-3-3}{2.3.3}) over $(P,Q,G)$ (so its morphisms are the
morphisms of biextensions compatible with the identity morphisms of
$P,Q,G$). Then \underline{this category depends covariantly on $G$ and
contravariantly on $P,Q$}. Moreover, with these variances,
\underline{the category $\operatorname{BIEXT}(P,Q;G)$ depends
additively on each of the three arguments $P,Q,G$}. For example, for
the argument $G$, this means precisely:

\begin{enumerate}[label=\alph*)]
\item $\operatorname{BIEXT}(P,Q;0)$ is equivalent to the one-point
category;
\item for two abelian groups $G,G'$ of $\underline T$, the functor
\begin{equation}\tag{\SGARefTarget{sga7.sga7i.evii.n2-4-2}{2.4.2}}
 \operatorname{BIEXT}(P,Q;G\times G')
 \longrightarrow
 \operatorname{BIEXT}(P,Q;G)\times
 \operatorname{BIEXT}(P,Q;G')
\end{equation}
deduced from the two projections $G\times G'\longrightarrow G$ and
$G\times G'\longrightarrow G'$ is an
\underline{equivalence of categories}.
\end{enumerate}
In the language of [4], one also says that, for fixed $P,Q$, the
category of biextensions of $(P,Q)$ by variable groups $G$ is, via the
functor $E\mapsto G$ induced by (\SGARefLink{sga7.sga7i.evii.n2-3-5}{2.3.5}), an
\underline{additive cofibred category} over
$\operatorname{GRAB}(\underline T)$.

\underline{Remark \SGARefTarget{sga7.sga7i.evii.n2-4-3}{2.4.3}.} In the terminology of \emph{loc.\ cit.}, the
preceding cofibred category is not merely additive but even
\underline{left exact} [4, \SGARefLink{sga7.sga7i.evii.n2-2}{2.2}]. We shall not verify this property here,
since it will reappear below in \SGARefLink{sga7.sga7i.evii.n3-6}{3.6} as a consequence of more precise
facts. We only point out that the general theory of \emph{loc.\ cit.}
would then describe the structure of this cofibred category, up to
equivalence of fibred categories, in terms of a chain complex
$L_\bullet$ in $\operatorname{GRAB}(\underline T)$ depending only on
$P$ and $Q$. We shall also give below a direct construction of this
complex as
\[
 P\overset{\mathbb L}{\otimes}Q,
\]
the tensor product taken in a derived category.


% END INCLUDED SGA7I_English_current_buildable_source_20260823/components/77_expose_VII_section_2_4.tex

% BEGIN INCLUDED SGA7I_English_current_buildable_source_20260823/components/78_expose_VII_section_2_5.tex
% SOURCE AUTHORITY:
% expose_VII_body.tex lines 776--844; scan indices 176--179; source folios 165--168.
% The scan confirms all displayed Biext groups and the derived tensor notation.
% Hard stop before source line 846, section 2.6.

\SGARefTarget{sga7.sga7i.evii.n2-5}{2.5}. \underline{Structure of the category
$\operatorname{BIEXT}(P,Q;G)$}.

The general formal considerations of [4, no.~1] give several useful
consequences of the additive dependence of the categories
$\operatorname{BIEXT}(P,Q;G)$ on $G$.

\SGARefTarget{sga7.sga7i.evii.n2-5-1}{2.5.1}. First, for fixed $G$, the category
$\operatorname{BIEXT}(P,Q;G)$ has a composition law which assigns
functorially to two biextensions $E,E'$ of $(P,Q)$ by $G$ a
\underline{product biextension}, denoted $E.E'$. This construction is
simply the analogue of the familiar product of two extensions of one
commutative group by another. It is compatible with that construction
whether $E,E'$ are viewed as group extensions of $Q_P$ by $G_P$, or
as group extensions of $P_Q$ by $G_Q$. In particular, as a torsor
under $G_{P\times Q}$, $EE'$ is the product of the torsors under
$G_{P\times Q}$ defined by $E$ and $E'$.

The formation of $E.E'$ is moreover \underline{associative and
commutative} up to canonical isomorphisms, and this composition law on
$\operatorname{BIEXT}(P,Q;G)$ has a two-sided neutral object, called
the \underline{trivial biextension of $(P,Q)$ by $G$}. It may be
constructed, as in \emph{loc.\ cit.}, from the unique biextension
$E_0$ of $(P,Q)$ by the zero group $0$, up to isomorphism, by means of
the homomorphism $0\longrightarrow G$. From this description one
immediately obtains an alternative description of the trivial
biextension as the trivial torsor
\[
 E=P\times Q\times G.
\]
Its two group structures over $P$ and over $Q$ are obtained,
respectively, from the product group structures on $Q\times G$ and
$P\times G$ by the base changes $P\longrightarrow e$ and
$Q\longrightarrow e$.

Finally, for every biextension $E$ of $(P,Q)$ by $G$, one can define
the \underline{inverse biextension}, denoted $E^{-1}$, together with
an isomorphism from $EE^{-1}$ to the trivial biextension. This
determines $E^{-1}$ up to unique isomorphism and makes its dependence
on $E$ functorial.

In short, the category $\operatorname{BIEXT}(P,Q;G)$ behaves like a
group object in the category $\mathbf{Cat}$ of categories, except that
associativity, commutativity, unit, and inverse relations hold only up
to specified isomorphisms. One may call it a
\underline{pseudo-group} in the 2-category $\mathbf{Cat}$.

\SGARefTarget{sga7.sga7i.evii.n2-5-2}{2.5.2}. For fixed $P,Q,G$, the preceding properties imply the
following facts. For every biextension $E$ of $(P,Q)$ by $G$, the
monoid of endomorphisms of $E$ is canonically isomorphic to the monoid
of endomorphisms of the trivial biextension $\underline 1$: an
endomorphism $u$ of $\underline 1$ is sent to
$u.\mathrm{id}_E$ on $\underline 1.E\simeq E$. The endomorphism monoid
of $\underline 1$ is in fact a \underline{commutative group}, denoted
\begin{equation}\tag{\SGARefTarget{sga7.sga7i.evii.n2-5-2-1}{2.5.2.1}}
 \operatorname{Biext}^{0}(P,Q;G).
\end{equation}
Its composition law is also given by
$(u,v)\mapsto u.v\in
\operatorname{End}(\underline 1.\underline 1)
\simeq\operatorname{End}(\underline 1)$.

Likewise, the composition law $(E,E')\mapsto EE'$ equips the set
\begin{equation}\tag{\SGARefTarget{sga7.sga7i.evii.n2-5-2-2}{2.5.2.2}}
 \operatorname{Biext}^{1}(P,Q;G)
\end{equation}
of isomorphism classes of biextensions of $(P,Q)$ by $G$ with the
structure of a \underline{commutative group}; the inverse of the class
of $E$ is the class of $E^{-1}$.

\SGARefTarget{sga7.sga7i.evii.n2-5-3}{2.5.3}. Given a group homomorphism $G\longrightarrow G'$, the
corresponding functor
\begin{equation}\tag{\SGARefTarget{sga7.sga7i.evii.n2-5-3-1}{2.5.3.1}}
 \operatorname{BIEXT}(P,Q;G)
 \longrightarrow
 \operatorname{BIEXT}(P,Q;G')
\end{equation}
is compatible, up to canonical isomorphisms, with the operations
$E.E'$ and $E^{-1}$ of \SGARefLink{sga7.sga7i.evii.n2-5-1}{2.5.1}. Consequently, the two evident maps
\begin{equation}\tag{\SGARefTarget{sga7.sga7i.evii.n2-5-3-2}{2.5.3.2}}
 \operatorname{Biext}^{i}(P,Q;G)
 \longrightarrow
 \operatorname{Biext}^{i}(P,Q;G')
 \qquad(i=0,1)
\end{equation}
are \underline{group homomorphisms}. Moreover, if $E$ is a
biextension of $(P,Q)$ by $G$ and $E'$ the corresponding
biextension by $G'$, the natural homomorphism
\begin{equation}\tag{\SGARefTarget{sga7.sga7i.evii.n2-5-3-3}{2.5.3.3}}
 \operatorname{Hom}(E,E)=\operatorname{Aut}(E)
 \longrightarrow
 \operatorname{Hom}(E',E')=\operatorname{Aut}(E')
\end{equation}
is precisely (\SGARefLink{sga7.sga7i.evii.n2-5-3-2}{2.5.3.2}) for $i=0$.

\SGARefTarget{sga7.sga7i.evii.n2-5-4}{2.5.4}. The group $\operatorname{Biext}^{0}(P,Q;G)$ admits a simple
explicit description. If $E$ is a biextension of $(P,Q)$ by $G$, an
automorphism $u$ of $E$ is first of all an automorphism of the
underlying torsor under $G_{P\times Q}$, and hence is given by a
morphism
\[
 u_0:P\times Q\longrightarrow G
\]
through the formula
\[
 u(x)=u_0(j(x))x.
\]
Compatibility of $u$ with the first, respectively the second,
extension structure on $E$ is equivalent to additivity of $u_0$ in
the second, respectively the first, argument. Thus
\underline{$u$ is an automorphism of the biextension structure of $E$
if and only if $u_0$ is bilinear}. This gives a
\underline{canonical bijection}
\begin{equation}\tag{\SGARefTarget{sga7.sga7i.evii.n2-5-4-1}{2.5.4.1}}
 \operatorname{Biext}^{0}(P,Q;G)
 \simeq
 \operatorname{Hom}(P\otimes Q,G).
\end{equation}
It is readily checked that this bijection is independent of $E$,
respects the group structures, and is functorial in $G$ for the
functoriality on the left described in (\SGARefLink{sga7.sga7i.evii.n2-5-3-2}{2.5.3.2}).

Below we shall determine $\operatorname{Biext}^{1}(P,Q;G)$ and find
that, canonically and functorially in $G$, it is isomorphic to
\[
 \operatorname{Ext}^{1}
 \bigl(P\overset{\mathbb L}{\otimes}Q,G\bigr).
\]


% END INCLUDED SGA7I_English_current_buildable_source_20260823/components/78_expose_VII_section_2_5.tex

% BEGIN INCLUDED SGA7I_English_current_buildable_source_20260823/components/79_expose_VII_sections_2_6_and_2_7.tex
% SOURCE AUTHORITY:
% expose_VII_body.tex lines 846--870; scan indices 179--181; source folios 168--170.
% Source-backed corrections: the fibre in 2.6 is BIEXT(P,Q;G), not the printed
% BIEXT(P,P';G); the symmetric biextension has structural target Q x P.
% Hard stop before source line 873, section 2.8.

\SGARefTarget{sga7.sga7i.evii.n2-6}{2.6}. In \SGARefLink{sga7.sga7i.evii.n2-5}{2.5} only $G$ was varied, while $P$ and $Q$ were held fixed.
We now say a few words about variation in $P,Q$. The essential fact,
from which the others follow formally, is that the contravariant
dependence of $\operatorname{BIEXT}(P,Q;G)$ on $P,Q$
\underline{commutes} with its covariant dependence on $G$. Thus the
categories in question may be viewed as the fibre categories of a
\underline{cofibred category}
\begin{equation}\tag{\SGARefTarget{sga7.sga7i.evii.n2-6-1}{2.6.1}}
 \operatorname{BIEXT}'(\underline T)
 \longrightarrow
 \operatorname{GRAB}(\underline T)^{\mathrm{op}}\times
 \operatorname{GRAB}(\underline T)^{\mathrm{op}}\times
 \operatorname{GRAB}(\underline T),
\end{equation}
where the superscript $\mathrm{op}$ denotes the opposite category and
the structural functor is $E\mapsto(P,Q,G)$. This cofibred category is
additive in each of its three arguments, or, as we shall say,
\underline{tri-additive}.

By paraphrasing standard arguments, which we shall not reproduce, it
follows that the composition $E.E'$ in the fibre category
$\operatorname{BIEXT}(P,Q;G)$ may be described using the additive
structure of any one of the three arguments: the three resulting
pseudo-group structures are canonically isomorphic. Consequently, the
abelian groups
\begin{equation}\tag{\SGARefTarget{sga7.sga7i.evii.n2-6-2}{2.6.2}}
 \operatorname{Biext}^{i}(P,Q;G),\qquad i=0,1,
\end{equation}
may be regarded as two trilinear functors in $P,Q,G$, respectively in
$\operatorname{GRAB}(\underline T)^{\mathrm{op}}$,
$\operatorname{GRAB}(\underline T)^{\mathrm{op}}$, and
$\operatorname{GRAB}(\underline T)$. Equivalently, they are
\underline{trilinear functors} of variable $P,Q,G$ in
$\operatorname{GRAB}(\underline T)$, contravariant in $P,Q$ and
covariant in $G$, with values in the category of abelian groups.
For this variance, the isomorphism (\SGARefLink{sga7.sga7i.evii.n2-5-4-1}{2.5.4.1}) is functorial in all of
$P,Q,G$, and not merely in $G$. The analogous assertion will hold
below for the description of $\operatorname{Biext}^{1}$ by an
$\operatorname{Ext}^{1}$.

\SGARefTarget{sga7.sga7i.evii.n2-7}{2.7}. \underline{The symmetric biextension of a given biextension}.

If $E$ is a biextension of $(P,Q)$ by $G$, define its
\underline{symmetric biextension} $\,{}^{s}E$, which is a biextension
of $(Q,P)$ by $G$. The underlying sheaf of sets of $\,{}^{s}E$ is that
of $E$, and the actions of $G$ on $E$ and $\,{}^{s}E$ are the same.
The structural morphism
\[
 {}^{s}E\longrightarrow Q\times P
\]
is the composite of $E\longrightarrow P\times Q$ with the symmetry
$P\times Q\longrightarrow Q\times P$; finally, the left and right
extension structures of $\,{}^{s}E$ are, respectively, the right and
left extension structures of $E$. Evidently $E\mapsto{}^{s}E$ is an
automorphism of the category $\operatorname{BIEXT}(\underline T)$ of
biextensions, and indeed an \underline{involutive automorphism}:
\[
 {}^{s}({}^{s}E)=E.
\]
Every statement about bitorsors $E,\ldots$ therefore gives a symmetric
statement upon being applied to $\,{}^{s}E,\ldots$. We may call it the
statement obtained by \emph{exchanging left and right}. In general we
shall spell out only one of the two symmetric statements, choosing
between left and right, for example, according to our political
preferences.


% END INCLUDED SGA7I_English_current_buildable_source_20260823/components/79_expose_VII_sections_2_6_and_2_7.tex

% BEGIN INCLUDED SGA7I_English_current_buildable_source_20260823/components/80_expose_VII_section_2_8.tex
% SOURCE AUTHORITY:
% expose_VII_body.tex lines 873--899; scan indices 181--182; source folios 170--171.
% The conventional inverse-image notation f^* is used for the scan's typewritten f*.
% Page-header tokens at source lines 889 and 904 are excluded as running matter.
% Hard stop before source line 907, section 2.9.

\SGARefTarget{sga7.sga7i.evii.n2-8}{2.8}. \underline{Change of topos. Localization.}

Let
\[
 f:\underline T'\longrightarrow\underline T
\]
be a morphism of topoi. The inverse-image functor $f^*$ transforms a
biextension of $(P,Q)$ by $G$ into a biextension $f^*(E)$ of
$(f^*P,f^*Q)$ by $f^*G$, as always for a structure expressible in
terms of finite projective limits and arbitrary inductive limits. To
prove compatibility of the two extension structures on $f^*(E)$
induced by those on $E$, one expresses compatibility by commutativity
of (\SGARefLink{sga7.sga7i.evii.n2-1-1}{2.1.1}) in the universal case
\[
 S=P\times P\times Q\times Q,
\]
where $p,p',q,q'$ are the four projections. As expected,
$E\mapsto f^*(E)$ gives a canonical functor
\[
 f^*:\operatorname{BIEXT}(\underline T)
 \longrightarrow
 \operatorname{BIEXT}(\underline T').
\]

In what follows we shall use without further mention the fact that
every construction made with biextensions is compatible with inverse
images under morphisms of topoi. This applies in particular to the
product $EE'$ and inverse $E^{-1}$ considered in \SGARefLink{sga7.sga7i.evii.n2-5}{2.5}.

One may take for $\underline T'$ an \underline{induced topos}
$\underline T_{/S'}$, and for $f$ the canonical morphism, called the
localization morphism (SGA 4 IV 5),
\[
 f=i_{S'}:\underline T_{/S'}\longrightarrow\underline T.
\]
Then $f^*(E)$ is called the \underline{biextension induced} by $E$ on
$S'$, and is also denoted, as usual, by $E|_{S'}$ or $E_{S'}$.

Since restriction has the usual transitivity properties, the
categories $\operatorname{BIEXT}(\underline T_{/S})$, for variable
$S\in\underline T$, form the fibres of a \underline{fibred category}
over $\underline T$ (SGA 1 VI 8). In fact, this category is even a
\underline{stack} over $\underline T$ in the sense of Giraud [2]:
morphisms and objects in its fibres can be glued. For example, to give
a biextension over $S$ up to canonical isomorphism, it suffices to give
a covering family (that is, an epimorphic family) $S_i\longrightarrow
S$, a biextension $E_i$ over each $S_i$ (an object of
$\underline T_{/S_i}$), and, for every pair $(i,j)$, an isomorphism
between $E_i|_{S_j}$ and $E_j|_{S_i}$, subject to the usual cocycle
condition.

We call this the \underline{stack of biextensions over $\underline T$}.
Replacing $\underline T$ in (\SGARefLink{sga7.sga7i.evii.n2-3-3}{2.3.3}), and hence also in (\SGARefLink{sga7.sga7i.evii.n2-3-4}{2.3.4}) and
(\SGARefLink{sga7.sga7i.evii.n2-3-5}{2.3.5}), by $\underline T_{/S}$ for variable $S$ gives morphisms of
stacks over $\underline T$. The same applies to the operations $E.E'$
and $E^{-1}$ of (\SGARefLink{sga7.sga7i.evii.n2-5-1}{2.5.1}) in the categories
$\operatorname{BIEXT}(P_S,Q_S;G_S)$ for variable $S$, where $P,Q,G$
are fixed commutative groups over $\underline T$.


% END INCLUDED SGA7I_English_current_buildable_source_20260823/components/80_expose_VII_section_2_8.tex

% BEGIN INCLUDED SGA7I_English_current_buildable_source_20260823/components/81_expose_VII_section_2_9_opening_and_proposition.tex
% SOURCE AUTHORITY:
% expose_VII_body.tex lines 907--939; scan indices 183--185; source folios 172--174.
% Source-backed corrections: e_p is normalized to e_P; the second extension in
% the proof is of P_Q by G_Q, not the printed G_P.
% Hard stop before source line 941, Corollary 2.9.4.

\SGARefTarget{sga7.sga7i.evii.n2-9}{2.9}. \underline{Relations with birigidified torsors}.

In the terminology of \SGARefLink{sga7.sga7i.evii.n1-3-2}{1.3.2}, we saw in \SGARefLink{sga7.sga7i.evii.n2-2}{2.2} that, for every
biextension $E$ of $(P,Q)$ by $G$, the underlying torsor structure
under $G_{P\times Q}$ is canonically birigidified relative to the pair
of identity sections of $P$ and $Q$. This gives a canonical functor
\begin{equation}\tag{\SGARefTarget{sga7.sga7i.evii.n2-9-1}{2.9.1}}
 \operatorname{BIEXT}(P,Q;G)
 \xrightarrow{\ i\ }
 \operatorname{TORSBIRIG}(P,Q;G),
\end{equation}
where the category on the right consists of birigidified torsors over
$P\times Q$ under $G_{P\times Q}$.

Using the identity sections of $P$ and $Q$, for every product of $n$
objects each isomorphic to $P$ or $Q$ one obtains a canonical
morphism $\operatorname{inj}_i$ from the $i$-th factor into that
product. Thus, for a morphism from such a product to $G$, one may speak
of its restriction to the $i$-th factor. Consider the products
\begin{equation}\tag{\SGARefTarget{sga7.sga7i.evii.n2-9-2}{2.9.2}}
 \left\{
 \begin{array}{l}
 P\times Q,\\[1mm]
 P\times Q\times Q,\quad P\times Q\times Q\times Q,\\[1mm]
 P\times P\times Q,\quad P\times P\times P\times Q,\\[1mm]
 P\times P\times Q\times Q.
 \end{array}
 \right.
\end{equation}

\underline{Proposition \SGARefTarget{sga7.sga7i.evii.n2-9-3}{2.9.3}. Suppose that every morphism from one of
the products (\SGARefLink{sga7.sga7i.evii.n2-9-2}{2.9.2}) to $G$ which is trivial on each factor of that
product (embedded by the identity sections as just explained) is
trivial. Then:}

\medskip
\noindent
a) \underline{The canonical functor (\SGARefLink{sga7.sga7i.evii.n2-9-1}{2.9.1}) is fully faithful, and
both sides of (\SGARefLink{sga7.sga7i.evii.n2-9-1}{2.9.1}) are rigid categories. In particular, if $E$ is a
torsor under $G_{P\times Q}$ birigidified relative to $(e_P,e_Q)$,
there is at most one biextension structure on $E$ compatible with its
birigidified torsor structure.}

\medskip
\noindent
b) \underline{Such a biextension structure exists on $E$ if and only
if $E$ admits an extension structure of $Q_P$ by $G_P$ compatible with
its torsor structure under $G_P$, and an extension structure of $P_Q$
by $G_Q$ compatible with its torsor structure under $G_Q$
(cf.\ \SGARefLink{sga7.sga7i.evii.n1-3-5}{1.3.5}(b)). In that case there is a unique extension structure of
$Q_P$ by $G_P$ compatible with the torsor structure under $G_P$ and
the given rigidification over $P\times e_Q$, and a unique extension
structure of $P_Q$ by $G_Q$ compatible with the torsor structure under
$G_Q$ and the given rigidification over $e_P\times Q$; these two
extension structures are compatible in the sense of \SGARefLink{sga7.sga7i.evii.n2-1}{2.1}.}

\underline{Proof.} a) If $E$ is a birigidified torsor over
$P\times Q$, its automorphism group identifies with the group of
morphisms $P\times Q\longrightarrow G$ trivial on the two partial
factors. By the hypothesis applied to $P\times Q$, this group contains
only the identity automorphism, proving rigidity. Now let $E,E'$ be
biextensions of $(P,Q)$ by $G$. To see that every morphism
$f:E\longrightarrow E'$ of birigidified torsors is a morphism of
biextensions, apply \SGARefLink{sga7.sga7i.evii.n1-3-5}{1.3.5}(a) to the two extension structures defined
by $E,E'$ over $P$ and over $Q$, respectively. This uses the
hypothesis for $P\times Q\times Q$ and $P\times P\times Q$ in
(\SGARefLink{sga7.sga7i.evii.n2-9-2}{2.9.2}).

\noindent
b) Necessity is clear. For sufficiency, \SGARefLink{sga7.sga7i.evii.n1-3-5}{1.3.5}(b) first shows that the
stated condition gives unique extension structures of $Q_P$ by $G_P$
and of $P_Q$ by $G_Q$, compatible with the torsor structure and the
rigidifications. Here the hypothesis is used for
$P\times Q\times Q\times Q$ and $P\times P\times P\times Q$.
It remains to prove that the two extension structures are compatible,
that is, that diagram (\SGARefLink{sga7.sga7i.evii.n2-1-1}{2.1.1}) commutes. Its two composites are
isomorphisms of torsors over $P\times P\times Q\times Q$, and hence
differ by a morphism
\[
 u:P\times P\times Q\times Q\longrightarrow G.
\]
Because the partial trivializations of $E$ are identity sections for
the laws $\varphi$ and $\psi$, respectively, $u$ is trivial on each of
the four factors. The hypothesis for
$P\times P\times Q\times Q$ therefore makes $u$ trivial. This proves
the proposition.


% END INCLUDED SGA7I_English_current_buildable_source_20260823/components/81_expose_VII_section_2_9_opening_and_proposition.tex

% BEGIN INCLUDED SGA7I_English_current_buildable_source_20260823/components/82_expose_VII_corollary_2_9_4_through_remark_2_9_6.tex
% SOURCE AUTHORITY:
% expose_VII_body.tex lines 941--984; scan indices 185--187; source folios 174--176.
% Source-backed corrections: references to "2.9.2 b)" are to Proposition 2.9.3(b);
% Biext in (2.9.6.1) is written Biext^1 to identify the group of isomorphism classes.
% Hard stop before source line 986, section 2.10.

\underline{Corollary \SGARefTarget{sga7.sga7i.evii.n2-9-4}{2.9.4}. Let $P,Q,G$ be commutative groups
satisfying}
\begin{equation}\tag{\SGARefTarget{sga7.sga7i.evii.n2-9-4-1}{2.9.4.1}}
 \underline{\operatorname{Hom}}_{\mathrm{pt}}(P,G)=e,
 \qquad
 \underline{\operatorname{Hom}}_{\mathrm{pt}}(Q,G)=e
\end{equation}
\underline{(cf.\ \SGARefLink{sga7.sga7i.evii.n1-3-7}{1.3.7}). Then the hypothesis of Proposition 2.9.3 is
satisfied, and conclusions (a) and (b) of that proposition hold.
Moreover, if $E$ is a torsor under $G_{P\times Q}$ birigidified
relative to $(e_P,e_Q)$, it admits a biextension structure of $(P,Q)$
by $G$ compatible with its given structures if and only if the two
morphisms of pointed sheaves which it defines,}
\begin{equation}\tag{\SGARefTarget{sga7.sga7i.evii.n2-9-4-2}{2.9.4.2}}
 Q\longrightarrow R^1p_*(G_P),
 \qquad
 P\longrightarrow R^1q_*(G_Q),
\end{equation}
\underline{where $p:P\longrightarrow e$ and $q:Q\longrightarrow e$
are the structural morphisms, are morphisms of groups.}

The first assertion was proved in \SGARefLink{sga7.sga7i.evii.n1-3-7}{1.3.7}. For the second, necessity is
clear. To prove sufficiency, apply Proposition \SGARefLink{sga7.sga7i.evii.n2-9-3}{2.9.3}(b) and \SGARefLink{sga7.sga7i.evii.n1-3-5}{1.3.5}(b).
Multiplicativity of the first morphism in (\SGARefLink{sga7.sga7i.evii.n2-9-4-2}{2.9.4.2}) says that the
rigidified torsor $E$ under $G_P$ satisfies the condition of \SGARefLink{sga7.sga7i.evii.n1-3-5}{1.3.5}(b)
\underline{locally over $P$}. As observed in \SGARefLink{sga7.sga7i.evii.n1-3-7}{1.3.7}, that condition is
then true globally over $P$, so $E$ admits an extension structure of
$Q_P$ by $G_P$ compatible with its rigidified torsor structure.
Symmetrically, $E$ admits an extension structure of $P_Q$ by $G_Q$
compatible with its rigidified torsor structure under $G_Q$. The
conclusion follows from Proposition \SGARefLink{sga7.sga7i.evii.n2-9-3}{2.9.3}(b).

\underline{Example \SGARefTarget{sga7.sga7i.evii.n2-9-5}{2.9.5}.} Let $S$ be a scheme, let $P$ and $Q$ be
abelian schemes over $S$, and let $G=\underline G_{m,S}$ be the
multiplicative group scheme over $S$. We noted in \SGARefLink{sga7.sga7i.evii.n1-3-8}{1.3.8} that
(\SGARefLink{sga7.sga7i.evii.n2-9-4-1}{2.9.4.1}) then holds. A birigidified torsor $E$ over $P\times Q$ is
equivalently a birigidified invertible sheaf over $P\times Q$. Recall
that a \underline{divisorial correspondence on $P,Q$} is the
isomorphism class of such a birigidified invertible sheaf (see, for
example, [1]). In the fppf topology, the morphisms (\SGARefLink{sga7.sga7i.evii.n2-9-4-2}{2.9.4.2}) are
\begin{equation}\tag{\SGARefTarget{sga7.sga7i.evii.n2-9-5-1}{2.9.5.1}}
 Q\longrightarrow\underline{\operatorname{Pic}}_{P/S},
 \qquad
 P\longrightarrow\underline{\operatorname{Pic}}_{Q/S},
\end{equation}
defined by the divisorial correspondence. The theorem of the square
(\emph{loc.\ cit.}) says that these morphisms necessarily respect the
group structures. Thus \SGARefLink{sga7.sga7i.evii.n2-9-4}{2.9.4} shows that
\underline{the category of biextensions of $(P,Q)$ by
$G=\underline G_{m,S}$ is equivalent to the category of birigidified
invertible sheaves on $P\times Q$}; it is therefore a rigid category,
and its set of isomorphism classes identifies with the set of
divisorial correspondences on $(P,Q)$. This identification plainly
respects the natural group laws on the two sets, the first being that
of (\SGARefLink{sga7.sga7i.evii.n2-5-2-2}{2.5.2.2}).

\underline{Remark \SGARefTarget{sga7.sga7i.evii.n2-9-6}{2.9.6}.} Additivity of the first morphism in
(\SGARefLink{sga7.sga7i.evii.n2-9-4-2}{2.9.4.2}) is equivalent to factorization of the second through the
subsheaf $\underline{\operatorname{Ext}}^1(Q,G)$; note that, by \SGARefLink{sga7.sga7i.evii.n1-3-5}{1.3.5},
the canonical morphism
\[
 \underline{\operatorname{Ext}}^1(Q,G)
 \longrightarrow R^1q_*(G_Q)
\]
is injective. Likewise, additivity of the second morphism in
(\SGARefLink{sga7.sga7i.evii.n2-9-4-2}{2.9.4.2}) means that the first factors through
$\underline{\operatorname{Ext}}^1(P,G)$. This gives group
isomorphisms
\begin{equation}\tag{\SGARefTarget{sga7.sga7i.evii.n2-9-6-1}{2.9.6.1}}
 \operatorname{Biext}^{1}(P,Q;G)
 \simeq
 \operatorname{Hom}\bigl(Q,\underline{\operatorname{Ext}}^1(P,G)\bigr)
 \simeq
 \operatorname{Hom}\bigl(P,\underline{\operatorname{Ext}}^1(Q,G)\bigr).
\end{equation}
In the situation of \SGARefLink{sga7.sga7i.evii.n2-9-5}{2.9.5} these reduce to the familiar isomorphisms
\begin{equation}\tag{\SGARefTarget{sga7.sga7i.evii.n2-9-6-2}{2.9.6.2}}
 \operatorname{Corr}(P,Q)
 \simeq \operatorname{Hom}(P,Q')
 \simeq \operatorname{Hom}(Q,P'),
\end{equation}
where
\[
 P'=\underline{\operatorname{Ext}}^1(P,\underline G_m),
 \qquad
 Q'=\underline{\operatorname{Ext}}^1(Q,\underline G_m)
\]
are the dual abelian schemes of $P$ and $Q$, identified respectively
with $\underline{\operatorname{Pic}}^0_{P/S}$ and
$\underline{\operatorname{Pic}}^0_{Q/S}$. In both formulas,
$\operatorname{Hom}$ denotes the group of \underline{group
homomorphisms}.


% END INCLUDED SGA7I_English_current_buildable_source_20260823/components/82_expose_VII_corollary_2_9_4_through_remark_2_9_6.tex

% BEGIN INCLUDED SGA7I_English_current_buildable_source_20260823/components/83_expose_VII_section_2_10.tex
% SOURCE AUTHORITY:
% expose_VII_body.tex lines 986--1007; scan indices 187--190; source folios 176--179.
% Source-backed corrections: in 2.10.2 the multiplicative group is G_{m,S};
% in the A-linear dictionary the torsor is E_p, not the printed P_p.
% Hard stop before source line 1010, section 3.

\SGARefTarget{sga7.sga7i.evii.n2-10}{2.10}. \underline{Various generalizations of the notion of
biextension}. We mention them for completeness, although we shall not
use them.

\SGARefTarget{sga7.sga7i.evii.n2-10-1}{2.10.1}. One can define a biextension of $(P,Q)$ by $G$ without
assuming $P$ or $Q$ commutative, only $G$ being commutative. It is a
torsor under $G_{P\times Q}$ equipped with central extension
structures of $Q_P$ by $G_P$ and of $P_Q$ by $G_Q$, each compatible
with the torsor structure and compatible with one another, meaning
that they make (\SGARefLink{sga7.sga7i.evii.n2-1-1}{2.1.1}) commute. Most of the preceding developments
remain valid, \emph{mutatis mutandis}, for this more general notion.
Even when $P,Q$ are commutative, it is more general than the notion of
\SGARefLink{sga7.sga7i.evii.n2-1}{2.1}, which is the case in which the two underlying relative group
structures on $E$ are \underline{commutative}. In the noncommutative
case, the category of biextensions of $(P,Q)$ by $G$ still depends
additively on $G$, but not on $P$ and $Q$.

\SGARefTarget{sga7.sga7i.evii.n2-10-2}{2.10.2}. Paraphrasing \SGARefLink{sga7.sga7i.evii.n2-1}{2.1}, for a fixed integer $n\geq1$ one defines an
\underline{$n$-extension of a sequence of $n$ commutative groups}
$(P_1,P_2,\ldots,P_n)$ by a commutative group $G$. It consists of a
torsor $E$ under $G_{P_1\times\cdots\times P_n}$, equipped with $n$
partial extension structures over the partial products of $n-1$
factors, with a compatibility relation of type (\SGARefLink{sga7.sga7i.evii.n2-1-1}{2.1.1}) for every pair
of distinct indices between $1$ and $n$.

If $n\geq3$, the $P_i$ are abelian schemes over a scheme $S$, and
$G=\underline G_{m,S}$, every $n$-extension of
$(P_1,\ldots,P_n)$ by $G$ is \underline{trivial}.

\SGARefTarget{sga7.sga7i.evii.n2-10-3}{2.10.3}. The notion of biextension in \SGARefLink{sga7.sga7i.evii.n2-1}{2.1} may be viewed as a variant of
the notion of a commutative extension of commutative groups, that is,
of an extension of $\mathbb Z$-modules. Let $A$ be a ring of
$\underline T$, assumed commutative for simplicity, and let $P,Q,G$
be three $A$-modules. More generally one can define an
\underline{$A$-module biextension} of $(P,Q)$ by $G$, a variant of an
extension of $A$-modules.

Such a biextension is a biextension in the sense of \SGARefLink{sga7.sga7i.evii.n2-1}{2.1} whose
underlying partial extension structures are enriched to extension
structures of $A_P$-modules and $A_Q$-modules, respectively,
compatible with the relative module structures already present on
$Q_P,G_P$ and $P_Q,G_Q$. The two relative module structures on $E$
may be interpreted as two actions of the underlying multiplicative
monoid of $A$ on the sheaf of sets underlying $E$, and these two
actions are required to commute. This condition is added to the
compatibility expressed by (\SGARefLink{sga7.sga7i.evii.n2-1-1}{2.1.1}), and is itself expressed by an
analogous diagram, which we leave to the reader. Finally, the first
action of $A$ on $E$ is required to be distributive with respect not
only to the first composition law on $E$ but also to the second; the
same is required of the second action.

The dictionary of \SGARefLink{sga7.sga7i.evii.n1-1-6}{1.1.6} also has an $A$-linear variant. An extension
of the $A$-module $P$ by the $A$-module $G$ amounts to a family of
torsors $E_p$ under $G_p$, parametrized by points $p\in P(S)$ for
variable $S\in\operatorname{ob}\underline T$, with $E_p$ varying
\underline{$A$-linearly}, rather than merely additively, in $p$.
Besides the isomorphisms $\varphi_{p,p'}$ of (\SGARefLink{sga7.sga7i.evii.n1-1-3-1}{1.1.3.1}), this is
expressed by isomorphisms
\[
 \mu_{a,p}:\,{}^{a}E_p\xrightarrow{\ \sim\ }E_{ap}
 \qquad(p\in P(S),\ a\in A(S)),
\]
where $\,{}^{a}E_p$ denotes the image of the $G_S$-torsor $E_p$ under
extension of the structure group along
\[
 g\longmapsto ag:G_S\longrightarrow G_S.
\]
We leave to the reader the conditions on $\varphi$ and $\mu$ which
make these data correspond to an extension of $A$-modules $E$ of $P$
by $G$ (cf.\ \SGARefLink{sga7.sga7i.evii.n3-5-4}{3.5.4}(b)).


% END INCLUDED SGA7I_English_current_buildable_source_20260823/components/83_expose_VII_section_2_10.tex

% BEGIN INCLUDED SGA7I_English_current_buildable_source_20260823/components/84_expose_VII_section_3_1.tex
% SOURCE AUTHORITY:
% expose_VII_body.tex lines 1010--1110; scan indices 190--194; source folios 179--183.
% Direct scan comparison: physical PDF pages 191--195 at 600 dpi; diagram 3.1.3 checked
% node-by-node and arrow-by-arrow against physical page 192.
% Transparent notation corrections: printed/typewritten subscript letter o is rendered as zero
% throughout L_0, d_0, H_0 and degree 0; the erroneous prose reference 3.1.7.1 is 3.1.8.
% Hard stop before source line 1116, section 3.2.

\section*{\SGARefTarget{sga7.sga7i.evii.n3}{3}. Supplements on homological algebra}

The principal purpose of this section is to define the isomorphism
\[
 \Biext^1(P,Q;G)\simeq
 \Ext^1(P\overset{\mathbb L}{\otimes}Q,G)
\]
announced above. For this we shall need certain developments in
homological algebra which are of independent interest and will occupy
sections \SGARefLink{sga7.sga7i.evii.n3-1}{3.1} through \SGARefLink{sga7.sga7i.evii.n3-4}{3.4}. The developments in \SGARefLink{sga7.sga7i.evii.n3-1}{3.1} and \SGARefLink{sga7.sga7i.evii.n3-3}{3.3} may be
regarded as supplements to [4], especially to section 6.9 therein.

\subsection*{\SGARefTarget{sga7.sga7i.evii.n3-1}{3.1}. Construction of an additive cofibred category
$\underline\Psi_{L_\bullet}$ in terms of a chain complex $L_\bullet$
in an abelian category $\underline C$}

Let $\underline C$ be an abelian category and let $L_\bullet$ be a
chain complex in $\underline C$,
\begin{equation}\tag{\SGARefTarget{sga7.sga7i.evii.n3-1-1}{3.1.1}}
 L_2\longrightarrow L_1\longrightarrow L_0\longrightarrow0,
\end{equation}
written with lower-degree conventions, so that the differential has
degree $-1$. Thus $L_i=0$ for $i<0$. We shall define a cofibred
category
\[
 \underline\Psi=\underline\Psi_{L_\bullet}
\]
over $\underline C$; up to canonical equivalence of cofibred
categories it is the category of the same name in [4, 6.9]. The
objects of $\underline\Psi_{L_\bullet}$ are pairs
\[
 (X,\alpha),
\]
where $X$ is an extension of $L_0$ by an object $A$ of
$\underline C$, depending on $X$,
\begin{equation}\tag{\SGARefTarget{sga7.sga7i.evii.n3-1-2}{3.1.2}}
 0\longrightarrow A\xrightarrow{i}X\xrightarrow{j}L_0
 \longrightarrow0,
\end{equation}
and $\alpha$ is a trivialization of the inverse image $d_0^*(X)$ of
this extension by $d_0:L_1\longrightarrow L_0$. It is required that
the corresponding trivialization of the inverse image
$d_1^*d_0^*(X)$ by $d_1:L_2\longrightarrow L_1$ be the evident one,
obtained from the transitivity isomorphism
\[
 d_1^*d_0^*(X)\simeq(d_0d_1)^*(X)
\]
and the relation $d_0d_1=0$.

Equivalently, $\alpha$ is a lifting of $d_0:L_1\longrightarrow L_0$
to a morphism $\alpha:L_1\longrightarrow X$, subject to
$\alpha d_1=0$:
\begin{equation}\tag{\SGARefTarget{sga7.sga7i.evii.n3-1-3}{3.1.3}}
\begin{tikzcd}[row sep=large,column sep=large]
 &&0\arrow[d]\\
 &&A\arrow[d,"i"]\\
 &&X\arrow[d,"j"]\\
 L_2\arrow[r,"d_1"']&L_1\arrow[r,"d_0"']\arrow[ru,"\alpha"]
   &L_0\arrow[d]\\
 &&0
\end{tikzcd}
\qquad,\qquad \alpha d_1=0.\quad(*)
\end{equation}
Here is a simpler description, pointed out by Deligne:
$\underline\Psi_{L_\bullet}(A)$ is the category of extensions of the
complex $L_\bullet$ by $A$, regarded as a complex concentrated in
degree zero.

In this form one sees that the description of the objects
$(X,\alpha)$ depends only on the chain complex $[L_\bullet]$
\underline{truncated in degree 1}, obtained from $L_\bullet$ by
``killing $H_i(L_\bullet)$ for $i\geq2$'': its terms are zero in
degrees at least $2$, agree with those of $L_\bullet$ in degrees at
most $0$, and its degree-$1$ term is the cokernel of
$d_1:L_2\longrightarrow L_1$.

A homomorphism from $(X,\alpha)$ to $(X',\alpha')$ is a homomorphism
of extensions $u:X\longrightarrow X'$, compatible with the identity
on $L_0$, such that $u\alpha=\alpha'$. Homomorphisms are composed by
composing the corresponding homomorphisms of extensions. This gives
the announced category $\underline\Psi_{L_\bullet}$. Up to canonical
isomorphism, it depends only on the degree-$1$ truncation
$[L_\bullet]$ of $L_\bullet$.

There is an evident functor $(X,\alpha)\longmapsto A$,
\begin{equation}\tag{\SGARefTarget{sga7.sga7i.evii.n3-1-4}{3.1.4}}
 \pi:\underline\Psi_{L_\bullet}\longrightarrow\underline C.
\end{equation}
\underline{This functor is cofibred}. Indeed, let
$u:A\longrightarrow B$ be a morphism of $\underline C$, and let
$(X,\alpha)$ be an object of the fibre
$\underline\Psi_{L_\bullet}(A)$. Thus $X$ is an extension of $L_0$ by
$A$. Form the extension $u_*(X)$ of $L_0$ by $B$, the direct image of
$X$ under $u$, and equip it with the partial trivialization
$\beta=f\alpha$, where $f:X\longrightarrow u_*(X)$ is the canonical
morphism. It follows immediately that
\[
 f:(X,\alpha)\longrightarrow(u_*(X),\beta)
\]
is cocartesian for (\SGARefLink{sga7.sga7i.evii.n3-1-4}{3.1.4}): it has the usual universal property for
$u$-morphisms in $\underline\Psi_{L_\bullet}$ with source
$(X,\alpha)$. This is just the familiar and elementary corresponding
universal property of $f:X\longrightarrow u_*(X)$.

Finally, the cofibred category $\underline\Psi_{L_\bullet}$ is
\underline{additive}; that is, it satisfies
\begin{equation}\tag{\SGARefTarget{sga7.sga7i.evii.n3-1-5}{3.1.5}}
 \underline\Psi_{L_\bullet}(0)
 \quad\text{is equivalent to the one-point category},
\end{equation}
and, for arbitrary $A,B$ in $\underline C$,
\begin{equation}\tag{\SGARefTarget{sga7.sga7i.evii.n3-1-6}{3.1.6}}
 \underline\Psi_{L_\bullet}(A\times B)
 \longrightarrow
 \underline\Psi_{L_\bullet}(A)\times
 \underline\Psi_{L_\bullet}(B)
 \quad\text{is an equivalence}.
\end{equation}
This follows immediately from the corresponding properties of the
cofibred category of extensions of $L_0$ by variable objects of
$\underline C$. We shall freely use the consequences of the
additivity of $\underline\Psi_{L_\bullet}$ established in [4, 1] and
recalled above in \SGARefLink{sga7.sga7i.evii.n2-5}{2.5}.

\SGARefTarget{sga7.sga7i.evii.n3-1-7}{3.1.7}. Among these consequences, for every object $A$ of
$\underline C$ the category $\underline\Psi_{L_\bullet}(A)$ has a
\underline{zero}, or \underline{trivial}, object (\SGARefLink{sga7.sga7i.evii.n2-5-1}{2.5.1}). It is the
trivial extension $X=L_0+A$ of $L_0$ by $A$, equipped with the
partial trivialization $\alpha:L_1\longrightarrow X$ whose components
are $d_0$ and $0$.

Moreover, by \SGARefLink{sga7.sga7i.evii.n2-5-2}{2.5.2}, for every $(X,\alpha)$ in
$\underline\Psi_{L_\bullet}(A)$ the automorphism group of this object
is canonically isomorphic to the automorphism group of the trivial
object. Denote the latter commutative group by
$\Psi^0_{L_\bullet}(A)$. The explicit description above gives the
canonical isomorphisms
\begin{multline}\tag{\SGARefTarget{sga7.sga7i.evii.n3-1-8}{3.1.8}}
 \Psi^0_{L_\bullet}(A)
 =\Ker\bigl(\Hom(L_0,A)\longrightarrow\Hom(L_1,A)\bigr)
 \simeq\Hom(H_0(L_\bullet),A)\\
 \simeq\Ext^0(L_\bullet,A),
\end{multline}
induced by the usual identification of the automorphism group of the
extension $X$ with $\Hom(L_0,A)$. The isomorphism (\SGARefLink{sga7.sga7i.evii.n3-1-8}{3.1.8}) is plainly
functorial in $A$.

\SGARefTarget{sga7.sga7i.evii.n3-1-9}{3.1.9}. Finally, recall from \SGARefLink{sga7.sga7i.evii.n2-5-2}{2.5.2} that the natural multiplicative
structure on the fibre category $\underline\Psi_{L_\bullet}(A)$
endows its set of isomorphism classes with a natural group structure.
We denote this group by $\Psi^1_{L_\bullet}(A)$; it is an additive
functor of $A$.


% END INCLUDED SGA7I_English_current_buildable_source_20260823/components/84_expose_VII_section_3_1.tex

% BEGIN INCLUDED SGA7I_English_current_buildable_source_20260823/components/85_expose_VII_section_3_2.tex
% SOURCE AUTHORITY:
% expose_VII_body.tex lines 1116--1243; scan indices 194--197; source folios 183--186.
% Direct scan comparison: physical PDF pages 195--198 at 600 dpi.  Diagrams (3.2.3),
% (*), and (3.2.5.2) were checked node-by-node, arrow-by-arrow, and label-by-label.
% Transparent notation corrections: printed/typewritten subscript letter o is rendered as zero;
% the source's erroneous ``j : L_0 -> X'' in the injectivity conclusion is rendered as the
% mathematically required ``s : L_0 -> X'', consistently with (3.2.5.1).
% The source's two distinct displays both numbered (3.2.3) are preserved as printed.
% Hard stop before source line 1246, section 3.3.

\subsection*{\SGARefTarget{sga7.sga7i.evii.n3-2}{3.2}. Determination of $\Psi^1_{L_\bullet}(A)$}

Consider an object $(X,\alpha)$ (\SGARefLink{sga7.sga7i.evii.n3-1-3}{3.1.3}) of
$\underline\Psi_{L_\bullet}(A)$. Consider the cochain complex of
length one
\begin{equation}\tag{\SGARefTarget{sga7.sga7i.evii.n3-2-1}{3.2.1}}
 [X\longrightarrow L_0]:\quad
 0\longrightarrow X\longrightarrow L_0\longrightarrow0\longrightarrow
 \cdots
\end{equation}
defined by the morphism $j:X\longrightarrow L_0$, with $X$ as its
degree-zero component. It is a resolution of $A$, and hence may be
identified with $A$ as an object of the derived category
$D(\underline C)$ [5]. The diagram (\SGARefLink{sga7.sga7i.evii.n3-1-3}{3.1.3}) then gives a degree-one
morphism of complexes
\begin{equation}\tag{\SGARefTarget{sga7.sga7i.evii.n3-2-2}{3.2.2}}
 L_\bullet\longrightarrow[X\longrightarrow L_0],
\end{equation}
made explicit by the following diagram with anticommutative squares:
\begin{equation}\tag{\SGARefTarget{sga7.sga7i.evii.n3-2-3}{3.2.3}}
\begin{tikzcd}[column sep=large,row sep=huge]
 \cdots & L_2\arrow[r,"d_1"]
 & L_1\arrow[r,"d_0"]\arrow[d,"u^{-1}=-\alpha"']
 & L_0\arrow[r]\arrow[d,"u^0=\mathrm{id}"'] & 0 & \cdots\\
 \cdots & 0\arrow[r] & X\arrow[r,"j"] & L_0\arrow[r] & 0 & \cdots
\end{tikzcd}
\end{equation}
(We use the usual convention which associates with the chain complex
$L_\bullet$ a complex whose differential has degree $1$ and whose
degree-$i$ component is $L_{-i}$.) We thus obtain a degree-one
morphism in the derived category
\begin{equation}\tag{\SGARefTarget{sga7.sga7i.evii.n3-2-3.occ-2}{3.2.3}}
 c(X,\alpha):L_\bullet\longrightarrow A
 \qquad
 \bigl(\text{a degree-one morphism in }D(\underline C)\bigr),
\end{equation}
which plainly depends only on the isomorphism class of $(X,\alpha)$.
We have thereby defined a canonical map
\begin{equation}\tag{\SGARefTarget{sga7.sga7i.evii.n3-2-4}{3.2.4}}
 \Psi^1_{L_\bullet}(A)\longrightarrow\Ext^1(L_\bullet,A)
 \overset{\mathrm{def}}{=}
 \Hom_{D(\underline C)}(L_\bullet,A[1]).
\end{equation}

\underline{Theorem} \SGARefTarget{sga7.sga7i.evii.n3-2-5}{3.2.5}. \underline{The map} (\SGARefLink{sga7.sga7i.evii.n3-2-4}{3.2.4})
\underline{is an isomorphism of groups, functorial in the object}
$A$ \underline{of} $\underline C$.

We leave to the reader the verification that (\SGARefLink{sga7.sga7i.evii.n3-2-4}{3.2.4}) is additive and
functorial in $A$, and shall confine ourselves to proving that it is
bijective. For injectivity, let us first show that if the morphism of
complexes $u$ in (\SGARefLink{sga7.sga7i.evii.n3-2-3}{3.2.3}) is zero in $D(\underline C)$, then it is
already zero in the homotopy category $K(\underline C)$; that is,
there exists a morphism
\[
 s:L_0\longrightarrow X
\]
such that
\begin{equation}\tag{\SGARefTarget{sga7.sga7i.evii.n3-2-5-1}{3.2.5.1}}
 js=u^0=\mathrm{id}_{L_0},
 \qquad sd_0=-u^{-1}=\alpha.
\end{equation}
Indeed, the hypothesis says that the conclusion becomes true after
the morphism of complexes under consideration is composed with a
suitable quasi-isomorphism $v$
\begin{equation}\tag{*}
\begin{tikzcd}[column sep=large,row sep=huge]
 0\arrow[r] & X\arrow[r,"j"]\arrow[d]
 & L_0\arrow[r]\arrow[d] & 0\\
 0\arrow[r] & C^0\arrow[r,"j'"] & C^1\arrow[r] & C^2\arrow[r] & \cdots
\end{tikzcd}
\end{equation}
(and, after replacing $C^i$ by $0$ for $i\geq2$ and $C^1$ by
$\operatorname{Im}(C^0\longrightarrow C^1)
=\operatorname{Ker}(C^1\longrightarrow C^2)$ if necessary, we may
suppose that $C^i=0$ for $i\geq2$). But the morphism (*) induces the
identity on $A$, and therefore identifies $X$ with the fibre product
of $L_0$ and $C^0$ over $C^1$. A morphism
$s':L_0\longrightarrow C^0$ defining a homotopy $vu\sim0$, and hence
satisfying $j's'=v^1u^0$, consequently factors through a morphism
$s:L_0\longrightarrow X$ satisfying
$js=u^0=\mathrm{id}_{L_0}$. Thus $s$ splits the extension $X$ of
$L_0$ by $A$, so the complex $X\longrightarrow L_0$ is isomorphic
\underline{in} $K(\underline C)$ to the complex $K^\bullet$
concentrated in $A$. For the latter complex it is clear that
\[
 \Hom_{K(\underline C)}(L_\bullet,K^\bullet)
 \longrightarrow
 \Hom_{D(\underline C)}(L_\bullet,K^\bullet)
\]
is an isomorphism, proving the assertion.

Now the choice of a morphism $s:L_0\longrightarrow X$ satisfying
(\SGARefLink{sga7.sga7i.evii.n3-2-5-1}{3.2.5.1}) is precisely equivalent to a trivialization of the
extension $X$ of $L_0$ by $A$ compatible with the partial
trivialization $\alpha$. This proves the injectivity of (\SGARefLink{sga7.sga7i.evii.n3-2-4}{3.2.4}).

For surjectivity, start with an element of the right-hand side. By
definition it may be described by a degree-one morphism of complexes
\begin{equation}\tag{\SGARefTarget{sga7.sga7i.evii.n3-2-5-2}{3.2.5.2}}
\begin{tikzcd}[column sep=large,row sep=huge]
 \ldots & L_2\arrow[r,"d_1"]
 & L_1\arrow[r,"d_0"]\arrow[d,"u^{-1}"']
 & L_0\arrow[r]\arrow[d,"u^0"'] & 0\\
 & 0\arrow[r] & C^0\arrow[r] & C^1\arrow[r] & \ldots
\end{tikzcd}
\end{equation}
with anticommutative squares, where $C^\bullet$ is a resolution of
$A$ and, as above, we may suppose that $C^i=0$ for $i\geq2$. Thus
$C^0$ is an extension of $C^1$ by $A$, and we may form the extension
$X$ obtained by pulling it back along $L_0\longrightarrow C^1$.
Anticommutativity in (\SGARefLink{sga7.sga7i.evii.n3-2-5-2}{3.2.5.2}) shows that $-u^{-1}$ factors through a
morphism $\alpha:L_1\longrightarrow X$ which is a partial splitting,
that is, $j\alpha=d_0$ (with $j$ as in (\SGARefLink{sga7.sga7i.evii.n3-1-2}{3.1.2})). Moreover, the
relation $u^{-1}d_1=0$ in (\SGARefLink{sga7.sga7i.evii.n3-2-5-2}{3.2.5.2}) implies
$\alpha d_1=0$. Hence $(X,\alpha)$ is indeed an object of
$\underline\Psi_{L_\bullet}(A)$. By construction, the morphism
(\SGARefLink{sga7.sga7i.evii.n3-2-5-2}{3.2.5.2}) is the composite of the morphism (\SGARefLink{sga7.sga7i.evii.n3-2-3.occ-2}{3.2.3}) associated with
$(X,\alpha)$ and the morphism
\[
 [X\longrightarrow L_0]\longrightarrow C^\bullet
 =[C^0\longrightarrow C^1]
\]
whose components are $(f,u^0)$, where $f:X\longrightarrow C^0$ is
the canonical morphism obtained from the definition of $X$ as the
pullback of the extension $C^0$.

The preceding morphism is a morphism of resolutions of $A$ and
induces the identity on $A$. It follows, by the definition
$\Ext^1=\Hom^1_{D(\underline C)}$, that the given element of
$\Ext^1(L_\bullet,A)$ described by (\SGARefLink{sga7.sga7i.evii.n3-2-5-2}{3.2.5.2}) is precisely
$c(X,\alpha)$. This proves the surjectivity of (\SGARefLink{sga7.sga7i.evii.n3-2-4}{3.2.4}) and completes
the proof of \SGARefLink{sga7.sga7i.evii.n3-2-5}{3.2.5}.


% END INCLUDED SGA7I_English_current_buildable_source_20260823/components/85_expose_VII_section_3_2.tex

% BEGIN INCLUDED SGA7I_English_current_buildable_source_20260823/components/86_expose_VII_section_3_3.tex
% SOURCE AUTHORITY:
% expose_VII_body.tex lines 1246--1352; scan indices 197--203; source folios 186--192.
% Direct scan comparison: physical PDF pages 198--204 at 600 dpi.
% Transparent source corrections: the two references printed as 3.2.3 in 3.3.2 are 3.3.2.1;
% the quasi-maximal construction is written with the induced map \bar d_0:L'_1->L_0 rather
% than the internally inconsistent printed d'_1/d_1 notation; the canonical map A->X is i,
% not the printed j; and the boundary arrow in (3.3.6.1) is \partial, as the following prose
% and (3.3.6.2) require.  The typo ``auticommutative'' is translated as anticommutative.
% Hard stop before source line 1355, section 3.4.

\subsection*{\SGARefTarget{sga7.sga7i.evii.n3-3}{3.3}. Properties of the cofibred category
$\underline\Psi_{L_\bullet}$}

In this section, which is not essential for understanding what
follows, we freely use the notions of [4], our purpose being to
compare the preceding results with those of that work.

\underline{Theorem} \SGARefTarget{sga7.sga7i.evii.n3-3-1}{3.3.1}. \underline{The additive cofibred category}
$\underline\Psi_{L_\bullet}$ \underline{over} $\underline C$
\underline{is left exact} ([4, \SGARefLink{sga7.sga7i.evii.n2-2}{2.2}]), \underline{admits a
quasi-maximal object} ([4, 4.4]), \underline{and, for every object}
$E=(X,\alpha)$ \underline{of this category, there exists a typical
complex} $L^E_\bullet$ \underline{for} $E$ ([4, \SGARefLink{sga7.sga7i.evii.n3-1}{3.1}]).
\underline{Moreover, the typical pro-complex of}
$\underline\Psi_{L_\bullet}$ ([4, \SGARefLink{sga7.sga7i.evii.n3-7}{3.7} and 4.4])
\underline{is canonically isomorphic in} $D(\underline C)$
\underline{to the truncated complex} $[L_\bullet]$
\underline{considered in} \SGARefLink{sga7.sga7i.evii.n3-1}{3.1}.

We now spell out the meaning of the several assertions in \SGARefLink{sga7.sga7i.evii.n3-3-1}{3.3.1}.

\SGARefTarget{sga7.sga7i.evii.n3-3-2}{3.3.2}. To say that the cofibred category
$\underline\Psi_{L_\bullet}$ is \underline{left exact} means that it
is additive and also satisfies the following condition:

\SGARefTarget{sga7.sga7i.evii.n3-3-2-1}{3.3.2.1}. For every exact sequence
\[
 0\longrightarrow A'\xrightarrow{u}A\xrightarrow{v}A''
 \longrightarrow0
\]
in $\underline C$, the natural functor (in the notation of [4, 1.8])
\[
 \underline\Psi_{L_\bullet}(A)
 \longrightarrow
 \underline\Psi_{L_\bullet}([A\longrightarrow A''])
\]
is an equivalence of categories.

Recall that if $\underline\Psi$ is a cofibred category over
$\underline C$ and $v:A\longrightarrow A''$ is a morphism in
$\underline C$, then $\underline\Psi([A\longrightarrow A''])$ is the
category of pairs $(E,\rho)$ in which $E$ is an object of
$\underline\Psi(A)$ and $\rho$ is a ``trivialization'' of the
corresponding object $v_*(E)$ of $\underline\Psi(A'')$; that is, an
isomorphism from it to the zero object (\SGARefLink{sga7.sga7i.evii.n2-5-1}{2.5.1}) of
$\underline\Psi(A'')$. The definition of the functor
\[
 \underline\Psi(A')\longrightarrow
 \underline\Psi([A\longrightarrow A''])
\]
in \SGARefLink{sga7.sga7i.evii.n3-3-2-1}{3.3.2.1} is evident. When
$\underline\Psi=\underline\Psi_{L_\bullet}$, the objects of
$\underline\Psi_{L_\bullet}([A\longrightarrow A''])$ are explicitly
the triples $(X,\alpha,\rho)$, where $(X,\alpha)$ is as in (\SGARefLink{sga7.sga7i.evii.n3-1-3}{3.1.3})
and $\rho:X\longrightarrow A''$ is a morphism in $\underline C$
satisfying
\begin{equation}\tag{*}
 \rho i=v,\qquad \rho\alpha=0.
\end{equation}
If $v:A\longrightarrow A''$ belongs to a short exact sequence as in
\SGARefLink{sga7.sga7i.evii.n3-3-2-1}{3.3.2.1}, then, as is well known, giving a morphism satisfying the
first relation (*) is equivalent, by the left exactness of the
cofibred category of extensions of $L_0$ by variable objects of
$\underline C$, to giving an extension $X'$ of $L_0$ by
$A'\simeq\Ker v$ which is a subextension of $X$: one takes
$X'=\Ker\rho$. The second relation (*) says that $\alpha$ factors
through a morphism $\alpha':L_1\longrightarrow X'$. Up to the
compatibilities whose verification is immediate, this proves the
left exactness stated in \SGARefLink{sga7.sga7i.evii.n3-3-2-1}{3.3.2.1}.

\SGARefTarget{sga7.sga7i.evii.n3-3-3}{3.3.3}. Recall that an object $E_0$ of a cofibred category
$\underline\Psi$ over $\underline C$ is called \underline{quasi-maximal}
if, for every object $E$ of $\underline\Psi$ lying over an object
$A$ of $\underline C$, there is a morphism $E\longrightarrow E'$ in
$\underline\Psi$ inducing a \underline{monomorphism}
$A\longrightarrow A'$ in $\underline C$, such that $E$ is
\underline{dominated by} $E_0$, meaning that there is a morphism
$E_0\longrightarrow E'$. If $\underline C$ has enough injective
objects, this is equivalent to saying that every object of
$\underline\Psi$ over an injective object of $\underline C$ is
dominated by $E_0$.

For $\underline\Psi=\underline\Psi_{L_\bullet}$ there is the
following canonical quasi-maximal object. Put
\[
 L'_1=\Coker(d_1:L_2\longrightarrow L_1)
\]
and let $\bar d_0:L'_1\longrightarrow L_0$ be the morphism induced
by $d_0:L_1\longrightarrow L_0$. Let $E_0=(X_0,\alpha_0)$ be the
object of $\underline\Psi_{L_\bullet}(L'_1)$ for which
$X_0=L_0\oplus L'_1$ is the trivial extension of $L_0$ by $L'_1$
and $\alpha_0$ is the composite
\[
 L_1\twoheadrightarrow L'_1
 \xrightarrow{(\bar d_0,\,\mathrm{id}_{L'_1})}
 L_0\oplus L'_1.
\]

To see that $E_0$ is quasi-maximal, observe that for every object
$(X,\alpha)$ of $\underline\Psi_{L_\bullet}$ one can find a morphism
from it to an object $(X',\alpha')$, inducing a
\underline{monomorphism} $A\longrightarrow A'$ in $\underline C$,
such that the extension $X'$ of $L_0$ by $A'$ is \underline{split}.
Indeed, it suffices to take $A'=X$ and
$(X',\alpha')=i_*(X,\alpha)$, with the notation of (\SGARefLink{sga7.sga7i.evii.n3-1-2}{3.1.2}). On the
other hand, it is immediate that any object $(X',\alpha')$ of
$\underline\Psi$ for which the extension $X'$ is split is dominated
by $E_0$.

\SGARefTarget{sga7.sga7i.evii.n3-3-4}{3.3.4}. Let $E$ be an object of an additive cofibred category
$\underline\Psi$ over $\underline C$, lying above an object $A$ of
$\underline C$, and consider the following functor of the variable
object $B$ of $\underline C$:
\begin{equation}\tag{\SGARefTarget{sga7.sga7i.evii.n3-3-4-1}{3.3.4.1}}
 B\longmapsto\Hom(E,\theta_B),
\end{equation}
where $\theta_B$ denotes the trivial object of $\underline\Psi(B)$.
We say that $E$ \underline{admits a typical complex} if this functor
is representable by an object $\Omega_E$ of $\underline C$. Since
the functor in question maps canonically to the functor
$B\longmapsto\Hom(A,B)$ represented by $A$, one obtains a canonical
morphism
\begin{equation}\tag{\SGARefTarget{sga7.sga7i.evii.n3-3-4-2}{3.3.4.2}}
 A\longrightarrow\Omega_E,
\end{equation}
and the corresponding length-one \underline{chain complex}
$A\longrightarrow\Omega_E$ in $\underline C$ is called the
\underline{typical complex of} $E$.

When $\underline\Psi=\underline\Psi_{L_\bullet}$ and
$E=(X,\alpha)$, one checks immediately that the functor (\SGARefLink{sga7.sga7i.evii.n3-3-4-1}{3.3.4.1}) is
represented by $\Coker\alpha$, and that the canonical morphism
(\SGARefLink{sga7.sga7i.evii.n3-3-4-2}{3.3.4.2}) is simply the composite
\[
 A\xrightarrow{i}X\xrightarrow{\mathrm{can}}\Coker\alpha.
\]

\SGARefTarget{sga7.sga7i.evii.n3-3-5}{3.3.5}. If an additive cofibred category $\underline\Psi$ over
$\underline C$ admits a quasi-maximal object (\SGARefLink{sga7.sga7i.evii.n3-3-3}{3.3.3}), and the typical
complexes of its objects exist, then the typical complex $T_\bullet$
of a quasi-maximal object $E_0$, regarded as an object of the
\underline{derived category} $D(\underline C)$, is independent, up
to canonical isomorphism, of the choice of $E_0$ ([4, 4.4]). It is
called the \underline{typical complex of the additive cofibred
category} $\underline\Psi$. If, in addition, $\underline\Psi$ is
left exact (\SGARefLink{sga7.sga7i.evii.n3-3-2}{3.3.2}), then $\underline\Psi$ can be reconstructed from
the typical complex $T_\bullet$, up to an equivalence of fibred
categories over $\underline C$ unique up to unique isomorphism
([4, 6.10]). There are also canonical isomorphisms, functorial in
$A\in\underline C$ ([4, 5.4]),
\begin{equation}\tag{\SGARefTarget{sga7.sga7i.evii.n3-3-5-1}{3.3.5.1}}
 \Psi^0(A)\xrightarrow{\sim}\Ext^0(T_\bullet,A),
 \qquad
 \Psi^1(A)\xrightarrow{\sim}\Ext^1(T_\bullet,A).
\end{equation}

To complete the comparison with [4], we shall prove that when
$\underline\Psi=\underline\Psi_{L_\bullet}$ its typical complex is
the truncation $[L_\bullet]$. This gives another definition of the
isomorphisms (\SGARefLink{sga7.sga7i.evii.n3-1-8}{3.1.8}) and (\SGARefLink{sga7.sga7i.evii.n3-2-4}{3.2.4}) in terms of (\SGARefLink{sga7.sga7i.evii.n3-3-5-1}{3.3.5.1}), and one
checks that it is compatible with the definitions in \SGARefLink{sga7.sga7i.evii.n3-1}{3.1} and \SGARefLink{sga7.sga7i.evii.n3-2}{3.2}.

For this assertion we may plainly restrict to the case where
$L_\bullet$ is already truncated in degree $1$, since \SGARefLink{sga7.sga7i.evii.n3-1}{3.1} noted that
truncation does not change $\underline\Psi_{L_\bullet}$ up to an
isomorphism of cofibred categories. The explicit construction in
\SGARefLink{sga7.sga7i.evii.n3-3-4}{3.3.4} then shows immediately that the typical complex of the
quasi-maximal object described in \SGARefLink{sga7.sga7i.evii.n3-3-3}{3.3.3} is canonically isomorphic to
$L_\bullet$.

\underline{Remark} \SGARefTarget{sga7.sga7i.evii.n3-3-6}{3.3.6}. Let
\[
 0\longrightarrow A'\longrightarrow A\longrightarrow A''
 \longrightarrow0
\]
be an exact sequence in $\underline C$. It gives an exact sequence
of $\Ext^i$ groups
\[
 0\longrightarrow\Ext^0(L_\bullet,A')
 \longrightarrow\Ext^0(L_\bullet,A)
 \longrightarrow\cdots\longrightarrow\Ext^1(L_\bullet,A)
 \longrightarrow\Ext^1(L_\bullet,A''),
\]
which could of course be continued to the right by introducing
$\Ext^i$ for $i\geq2$. Using the isomorphisms (\SGARefLink{sga7.sga7i.evii.n3-1-8}{3.1.8}) and (\SGARefLink{sga7.sga7i.evii.n3-2-4}{3.2.4}),
we obtain the exact sequence
\begin{multline}\tag{\SGARefTarget{sga7.sga7i.evii.n3-3-6-1}{3.3.6.1}}
 0\longrightarrow\Psi^0_{L_\bullet}(A')
 \longrightarrow\Psi^0_{L_\bullet}(A)
 \longrightarrow\Psi^0_{L_\bullet}(A'')
 \xrightarrow{\partial}\Psi^1_{L_\bullet}(A')\\
 \longrightarrow\Psi^1_{L_\bullet}(A)
 \longrightarrow\Psi^1_{L_\bullet}(A'').
\end{multline}
All arrows other than $\partial$ are defined simply by functoriality
of $\Psi^0_{L_\bullet}$ and $\Psi^1_{L_\bullet}$. On the other hand,
[4, \SGARefLink{sga7.sga7i.evii.n2-5-2}{2.5.2}] defines an exact sequence (3.3.6.1), for every left-exact
cofibred category $\underline\Psi$ over $\underline C$, through a
morphism
\begin{equation}\tag{\SGARefTarget{sga7.sga7i.evii.n3-3-6-2}{3.3.6.2}}
 \partial:\Psi^0(A'')\longrightarrow\Psi^1(A').
\end{equation}
It associates with an element $u$ of the left-hand side the class,
in the right-hand side, of the object of $\underline\Psi(A')$,
determined up to canonical isomorphism, whose image in
$\underline\Psi([A\longrightarrow A''])$ (notation of 3.3.2) is the
pair $(X_0,u)$, where $X_0$ is the trivial object of
$\underline\Psi(A)$.

For $\underline\Psi=\underline\Psi_{L_\bullet}$, the morphism
(\SGARefLink{sga7.sga7i.evii.n3-3-6-2}{3.3.6.2}) defined by this \underline{``geometric'' construction is
the negative of the boundary morphism} $\partial$ \underline{in}
(\SGARefLink{sga7.sga7i.evii.n3-3-6-1}{3.3.6.1}). Thus, after the identifications above and a change of
sign in $\partial$, the latter exact sequence is precisely that of
[4, \SGARefLink{sga7.sga7i.evii.n2-5-2}{2.5.2}]. This follows from the final assertion of [4, 5.4], in
view of the compatibilities noted in \SGARefLink{sga7.sga7i.evii.n3-3-5}{3.3.5}.

\underline{Remark} \SGARefTarget{sga7.sga7i.evii.n3-3-6-3}{3.3.6.3}. We take this opportunity to clarify a
point in [4, no.~3.11] concerning the fundamental isomorphisms
(3.11.3) and (3.11.4), which play an essential role in the proof of
[4, 5.4]. The diagram on page 35 of [4] is commutative, not
anticommutative, and therefore does not as it stands define a
morphism
\[
 u:[L^X_1\longrightarrow L^X_0]\longrightarrow K^\bullet[1].
\]
It must be made anticommutative by replacing $u^1$ by $-u^1$. One
could instead adopt the opposite convention and replace $u^0$ by
$-u^0$; however, the stated convention, in which $u^1$ becomes
$-u^1$, is the one for which (3.11.3) generalizes the isomorphism
(\SGARefLink{sga7.sga7i.evii.n3-5-3}{3.5.3}) of [4], as claimed there. With this clarification, the claim
made in [4, 5.4] must also be corrected: the two boundary morphisms
considered there \underline{differ by a sign}.

\SGARefTarget{sga7.sga7i.evii.n3-3-7}{3.3.7}. In [4, 6.13] the variance of the cofibred category
$\underline\Psi_{L_\bullet}$ over $\underline C$ as $L_\bullet$
varies was also made explicit, in particular by reconstructing the
category of cocartesian functors from $\underline\Psi_{L_\bullet}$
to $\underline\Psi_{L'_\bullet}$ in terms of the chain complexes
$L_\bullet$ and $L'_\bullet$. The set of isomorphism classes of such
cocartesian functors is canonically in bijection with
\[
 \Hom_{D(\underline C)}(L'_\bullet,L_\bullet).
\]
We shall use only a small part of these results, which we develop in
the next section using the simple description of
$\underline\Psi_{L_\bullet}$ given in \SGARefLink{sga7.sga7i.evii.n3-1}{3.1}.

% END INCLUDED SGA7I_English_current_buildable_source_20260823/components/86_expose_VII_section_3_3.tex

% BEGIN INCLUDED SGA7I_English_current_buildable_source_20260823/components/87_expose_VII_section_3_4_opening_through_3_4_7.tex
% SOURCE AUTHORITY:
% expose_VII_body.tex lines 1355--1434; scan indices 203--205; source folios 192--194.
% Direct scan comparison: physical PDF pages 204--206 at 600 dpi.  Diagrams (3.4.2)
% and (3.4.6) were checked node-by-node, arrow-by-arrow, and label-by-label.
% Transparent corrections: typewritten subscript letter o is rendered as zero; the arrow
% omitted by the working transcription in (3.4.5) is restored from the scan; and the proof's
% unindexed f is rendered as the required f^1.
% Hard stop before source line 1436, 3.4.8.

\subsection*{\SGARefTarget{sga7.sga7i.evii.n3-4}{3.4}. Variance of $\underline\Psi_{L_\bullet}$ with
$L_\bullet$}

Consider a morphism of complexes
\begin{equation}\tag{\SGARefTarget{sga7.sga7i.evii.n3-4-1}{3.4.1}}
 f:L'_\bullet\longrightarrow L_\bullet
\end{equation}
given by a commutative diagram
\begin{equation}\tag{\SGARefTarget{sga7.sga7i.evii.n3-4-2}{3.4.2}}
\begin{tikzcd}[column sep=large,row sep=large]
 {}\arrow[r]&L'_2\arrow[r]\arrow[d,"f_2"']
 &L'_1\arrow[r]\arrow[d,"f_1"']
 &L'_0\arrow[r]\arrow[d,"f_0"']&0\\
 {}\arrow[r]&L_2\arrow[r]&L_1\arrow[r]&L_0\arrow[r]&0.
\end{tikzcd}
\end{equation}
We associate with it a canonical functor
\begin{equation}\tag{\SGARefTarget{sga7.sga7i.evii.n3-4-3}{3.4.3}}
 f^*:\underline\Psi_{L_\bullet}\longrightarrow
 \underline\Psi_{L'_\bullet}.
\end{equation}
It is in fact a cocartesian functor of cofibred categories over
$\underline C$: it is compatible with the structural functors to
$\underline C$ and carries cocartesian morphisms to cocartesian
morphisms; see SGA 1, Exposé VI.

If $(X,\alpha)$ is an object of $\underline\Psi_{L_\bullet}$, define
$f^*(X,\alpha)=(X',\alpha')$ as follows. The extension $X'$ of
$L'_0$ by $A$ is the pullback, along $f_0$, of the given extension
$X$ of $L_0$ by $A$, and $\alpha'$ is the partial trivialization
deduced from $\alpha$ by transitivity in the second square of
(\SGARefLink{sga7.sga7i.evii.n3-4-2}{3.4.2}). The compatibility of $\alpha'$ with the relation
$d'_0d'_1=0$ follows immediately from the analogous compatibility
for $\alpha$. By construction, (\SGARefLink{sga7.sga7i.evii.n3-4-3}{3.4.3}) commutes with the structural
functors to $\underline C$. Its cocartesian character follows at
once from the fact that, in the category of extensions of objects of
$\underline C$, direct images of extensions commute with inverse
images of extensions.

The functor (\SGARefLink{sga7.sga7i.evii.n3-4-3}{3.4.3}) induces morphisms of functors
\begin{equation}\tag{\SGARefTarget{sga7.sga7i.evii.n3-4-4}{3.4.4}}
 f^i:\Psi^i_{L_\bullet}\longrightarrow\Psi^i_{L'_\bullet},
 \qquad i=0,1
\end{equation}
(compare \SGARefLink{sga7.sga7i.evii.n2-6}{2.6}). The morphism (\SGARefLink{sga7.sga7i.evii.n3-4-1}{3.4.1}), on the other hand, induces
morphisms
\begin{equation}\tag{\SGARefTarget{sga7.sga7i.evii.n3-4-5}{3.4.5}}
 f^i:\Ext^i(L_\bullet,-)\longrightarrow\Ext^i(L'_\bullet,-),
 \qquad i\in\mathbb Z.
\end{equation}
For $i=0,1$, the morphisms (\SGARefLink{sga7.sga7i.evii.n3-4-4}{3.4.4}) and (\SGARefLink{sga7.sga7i.evii.n3-4-5}{3.4.5}) are compatible with
the isomorphisms (\SGARefLink{sga7.sga7i.evii.n3-1-8}{3.1.8}) and (\SGARefLink{sga7.sga7i.evii.n3-2-4}{3.2.4}). In other words, for
$A\in\operatorname{Ob}\underline C$ and $i=0,1$, the following
diagram is commutative:
\begin{equation}\tag{\SGARefTarget{sga7.sga7i.evii.n3-4-6}{3.4.6}}
\begin{tikzcd}[column sep=large,row sep=large]
 \Psi^i_{L_\bullet}(A)\arrow[r,"\sim"]\arrow[d]
 &\Ext^i(L_\bullet,A)\arrow[d]\\
 \Psi^i_{L'_\bullet}(A)\arrow[r,"\sim"]
 &\Ext^i(L'_\bullet,A).
\end{tikzcd}
\end{equation}
The verification of this compatibility is essentially immediate and
is left to the reader.

\underline{Proposition} \SGARefTarget{sga7.sga7i.evii.n3-4-7}{3.4.7}. \underline{The functor} (\SGARefLink{sga7.sga7i.evii.n3-4-3}{3.4.3})
\underline{is faithful} (\underline{respectively fully faithful,
respectively an equivalence of categories}) \underline{if and only
if}
\[
 H_0(f):H_0(L'_\bullet)\longrightarrow H_0(L_\bullet)
\]
\underline{is an epimorphism} (\underline{respectively, if}
$H_0(f)$ \underline{is an isomorphism and} $H_1(f)$
\underline{is an epimorphism; respectively, if} $H_0(f)$
\underline{and} $H_1(f)$ \underline{are isomorphisms}).

Indeed, with $f^i$ denoting the functorial morphism (\SGARefLink{sga7.sga7i.evii.n3-4-4}{3.4.4}), the
required condition is that $f^0$ be a monomorphism (respectively,
that $f^0$ be an isomorphism and $f^1$ a monomorphism; respectively,
that $f^0$ and $f^1$ be isomorphisms). By the commutativity of
(\SGARefLink{sga7.sga7i.evii.n3-4-6}{3.4.6}), these conditions may be expressed by the analogous
conditions on the functors $\Ext^i$, $i=0,1$. The proposition then
reduces to a straightforward statement of homological algebra,
which we leave to the reader as an exercise.


% END INCLUDED SGA7I_English_current_buildable_source_20260823/components/87_expose_VII_section_3_4_opening_through_3_4_7.tex

% BEGIN INCLUDED SGA7I_English_current_buildable_source_20260823/components/88_expose_VII_sections_3_4_8_and_3_4_9.tex
% SOURCE AUTHORITY:
% expose_VII_body.tex lines 1436--1542; scan indices 205--209; source folios 194--198.
% Direct scan comparison: physical PDF pages 206--210 at 600 dpi, with the dense
% diagram inspected separately at original detail.  The apparent second arrowhead
% on the X'--L'_0 diagonal belongs to the crossing alpha' arrow: the source has the
% single morphism j':X'->L'_0 required by the extension, not a bidirectional arrow.
% Transparent corrections: 0_{L,A'} is rendered 0_{L,A}; unprimed j in the four
% defining relations is rendered j'; the concluding Psi_L is rendered Psi_{L''};
% and the missing comma in Ext^0(L',A) is restored.
% Hard stop before source line 1545, section 3.5.

\SGARefTarget{sga7.sga7i.evii.n3-4-8}{3.4.8}. Consider now an exact sequence of chain complexes
\begin{equation}\tag{\SGARefTarget{sga7.sga7i.evii.n3-4-8-1}{3.4.8.1}}
 0\longrightarrow L_\bullet\xrightarrow{f}L'_\bullet
 \xrightarrow{g}L''_\bullet\longrightarrow0.
\end{equation}
It gives functors of the form (\SGARefLink{sga7.sga7i.evii.n3-4-3}{3.4.3}),
\[
 \underline\Psi_{L''_\bullet}\xrightarrow{g^*}
 \underline\Psi_{L'_\bullet}\xrightarrow{f^*}
 \underline\Psi_{L_\bullet}.
\]
The transitivity property implicit in functors of the form (\SGARefLink{sga7.sga7i.evii.n3-4-3}{3.4.3})
gives an isomorphism of functors
\[
 f^*g^*\simeq(gf)^*=(0_{L_\bullet\to L''_\bullet})^*.
\]
The last member is evidently canonically isomorphic to the
cocartesian zero functor from $\underline\Psi_{L''_\bullet}$ to
$\underline\Psi_{L_\bullet}$.  The preceding isomorphism therefore
defines a canonical functor
\begin{equation}\tag{\SGARefTarget{sga7.sga7i.evii.n3-4-8-2}{3.4.8.2}}
 \underline\Psi_{L''_\bullet}\longrightarrow
 \underline\Psi_{[L_\bullet\longrightarrow L'_\bullet]}.
\end{equation}
Here $\underline\Psi_{[L_\bullet\longrightarrow L'_\bullet]}$
denotes the category of pairs $(X',t)$ in which $X'$ is an object of
$\underline\Psi_{L'_\bullet}$ and $t$ is a trivialization of its
inverse image $f^*(X')$, that is, an isomorphism
\[
 f^*(X')\simeq0_{L_\bullet,A},
\]
where $0_{L_\bullet,A}$ denotes the zero object of
$\underline\Psi_{L_\bullet}(A)$ and $A$ is the object of
$\underline C$ underlying $X'$.  We claim that \underline{the
functor} (\SGARefLink{sga7.sga7i.evii.n3-4-8-2}{3.4.8.2}) \underline{is an equivalence of categories}.
Together with the left exactness in $A$ established in \SGARefLink{sga7.sga7i.evii.n3-3-1}{3.3.1}, this
completes the exactness property in $L_\bullet$ of the expression
$\underline\Psi_{L_\bullet}(A)$.

Let us sketch the proof.  To give an object of the target of
(\SGARefLink{sga7.sga7i.evii.n3-4-8-2}{3.4.8.2}) is to give an extension
\[
 0\longrightarrow A\xrightarrow{i'}X'\xrightarrow{j'}L'_0
 \longrightarrow0
\]
of $L'_0$ by an object $A$ of $\underline C$, together with a
morphism $\alpha':L'_1\longrightarrow X'$ satisfying
$j'\alpha'=d'_0$ and $\alpha'd'_1=0$, and a morphism
$t:L_0\longrightarrow X'$ satisfying $j't=f_0$ and
$td_0=\alpha'f_1$.  These data are displayed in the following
diagram:
\begin{center}
\begin{tikzpicture}[
  x=1cm,y=1cm,>=Latex,
  every node/.style={font=\small,inner sep=1.2pt}
]
  \node (z00) at (0,0) {$0$};
  \node (l0)  at (1.3,0) {$L_0$};
  \node (lp0) at (2.8,0) {$L'_0$};
  \node (lpp0) at (4.3,0) {$L''_0$};
  \node (z05) at (5.5,0) {$0$};

  \node (z10) at (0,-1.35) {$0$};
  \node (l1)  at (1.3,-1.35) {$L_1$};
  \node (lp1) at (2.8,-1.35) {$L'_1$};
  \node (lpp1) at (4.3,-1.35) {$L''_1$};
  \node (z15) at (5.5,-1.35) {$0$};

  \node (z20) at (0,-2.70) {$0$};
  \node (l2)  at (1.3,-2.70) {$L_2$};
  \node (lp2) at (2.8,-2.70) {$L'_2$};
  \node (lpp2) at (4.3,-2.70) {$L''_2$};
  \node (z25) at (5.5,-2.70) {$0$};

  \node (xp) at (3.25,1.25) {$X'$};
  \node (ap) at (4.0,2.25) {$A$};
  \node (zap) at (4.75,3.25) {$0$};
  \node (xpp) at (5.85,1.25) {$X''$};
  \node (app) at (6.60,2.25) {$A$};
  \node (zapp) at (7.35,3.25) {$0$};

  \draw[->] (z00)--(l0);
  \draw[->] (l0)--node[above]{$f_0$}(lp0);
  \draw[->] (lp0)--node[above]{$g_0$}(lpp0);
  \draw[->] (lpp0)--(z05);
  \draw[->] (z10)--(l1);
  \draw[->] (l1)--node[above]{$f_1$}(lp1);
  \draw[->] (lp1)--node[above]{$g_1$}(lpp1);
  \draw[->] (lpp1)--(z15);
  \draw[->] (z20)--(l2);
  \draw[->] (l2)--(lp2);
  \draw[->] (lp2)--(lpp2);
  \draw[->] (lpp2)--(z25);

  \draw[->] (l1)--node[right]{$d_0$}(l0);
  \draw[->] (lp1)--node[right]{$d'_0$}(lp0);
  \draw[->] (lpp1)--node[right]{$d''_0$}(lpp0);
  \draw[->] (l2)--(l1);
  \draw[->] (lp2)--node[right]{$d'_1$}(lp1);
  \draw[->] (lpp2)--node[right]{$d''_1$}(lpp1);

  \draw[->] (zap)--(ap);
  \draw[->] (ap)--node[left]{$i'$}(xp);
  \draw[->,dashed] (zapp)--(app);
  \draw[->,dashed] (app)--node[left]{$i''$}(xpp);
  \draw[->] (l0) to[bend left=22] node[above]{$t$} (xp);
  \draw[->] (xp)--node[left]{$j'$}(lp0);
  \draw[->] (lp1)--node[pos=.82,right]{$\alpha'$}(xp);
  \draw[->,dashed] (xpp)--node[above]{$j''$}(lpp0);
  \draw[->,dashed] (lpp1) to[bend right=25]
    node[pos=.82,right]{$\alpha''$} (xpp);

  \node[anchor=west] at (7.55,-0.70) {$
    \left\{
    \begin{aligned}
      j'\alpha'&=d'_0,\\
      \alpha'd'_1&=0,\\
      j't&=f_0,\\
      td_0&=\alpha'f_1.
    \end{aligned}
    \right.$};
\end{tikzpicture}
\end{center}

The familiar left exactness, in the first argument, of extensions
of an object $L$ by an object $A$, applied to the first row of the
diagram, shows that a morphism $t$ satisfying $j't=f_0$ allows one
to interpret $X'$ as the inverse image of an extension $X''$ of
$L''_0$ by $A$.  Under this interpretation, $t$ corresponds to the
natural trivialization of the inverse image of that extension along
$g_0f_0=0$.

The morphism $\alpha'$ satisfying $j'\alpha'=d'_0$ and
$\alpha'f_1=td_0$ is interpreted by essentially the same dictionary,
now applied to the second row of the diagram, as arising from a
morphism $\alpha'':L''_1\longrightarrow X''$ satisfying
$j''\alpha''=d''_0$.  Here $\alpha'$ and $\alpha''$ are regarded as
trivializations of the inverse images of $X'$ and $X''$ along
$d'_0$ and $d''_0$, respectively.  Finally, the relation
$\alpha'd'_1=0$ is then equivalent to $\alpha''d''_1=0$.  Thus the
data $(X',\alpha',t)$, an object of
$\underline\Psi_{[L_\bullet\longrightarrow L'_\bullet]}$, are
equivalent to the data $(X'',\alpha'')$, an object of
$\underline\Psi_{L''_\bullet}$.  This proves, in substance, that
(\SGARefLink{sga7.sga7i.evii.n3-4-8-2}{3.4.8.2}) is indeed an equivalence of categories.

\SGARefTarget{sga7.sga7i.evii.n3-4-9}{3.4.9}. Continue to consider the exact sequence (\SGARefLink{sga7.sga7i.evii.n3-4-8-1}{3.4.8.1}), and the
exact sequence of $\Ext^i$ groups that it defines:
\begin{multline}\tag{\SGARefTarget{sga7.sga7i.evii.n3-4-9-1}{3.4.9.1}}
 0\longrightarrow\Ext^0(L''_\bullet,A)
 \xrightarrow{g^0}\Ext^0(L'_\bullet,A)
 \xrightarrow{f^0}\Ext^0(L_\bullet,A)
 \xrightarrow{\partial^0}\Ext^1(L''_\bullet,A)\\
 \xrightarrow{g^1}\Ext^1(L'_\bullet,A)
 \xrightarrow{f^1}\Ext^1(L_\bullet,A)\longrightarrow\cdots.
\end{multline}
The canonical isomorphisms of the forms (\SGARefLink{sga7.sga7i.evii.n3-1-8}{3.1.8}) and (\SGARefLink{sga7.sga7i.evii.n3-2-4}{3.2.4}) give a
geometric interpretation of the first six terms of this exact
sequence in terms of the categories $\underline\Psi_{L_\bullet}$,
$\underline\Psi_{L'_\bullet}$, and
$\underline\Psi_{L''_\bullet}$.  Compatibility (\SGARefLink{sga7.sga7i.evii.n3-4-6}{3.4.6}) likewise
gives a geometric interpretation of the arrows
$f^0,g^0,f^1,g^1$.  We shall now give such an interpretation of
$\partial^0$, using the result just proved that (\SGARefLink{sga7.sga7i.evii.n3-4-8-2}{3.4.8.2}) is an
equivalence of categories.

Let
\[
 u\in\Ext^0(L_\bullet,A),
 \qquad\text{i.e.}\qquad
 u:L_0\longrightarrow A,\quad ud_0=0.
\]
Consider the object $(0,\overline u)$ of
$\underline\Psi_{[L_\bullet\to L'_\bullet]}$.  Here $0$ is the
zero object of $\underline\Psi_{L'_\bullet}$, whose inverse image
along $f_0$ is therefore identified with the zero object of
$\underline\Psi_{L_\bullet}$, and $\overline u$ is the automorphism
of the latter defined by $u$ through the canonical isomorphism
(\SGARefLink{sga7.sga7i.evii.n3-1-8}{3.1.8}).  Since (\SGARefLink{sga7.sga7i.evii.n3-4-8-2}{3.4.8.2}) is an equivalence, $(0,\overline u)$ comes
from an object of $\underline\Psi_{L''_\bullet}(A)$, determined up
to unique isomorphism.  Its isomorphism class is an element of
\[
 \Psi^1_{L''_\bullet}(A)\simeq\Ext^1(L''_\bullet,A).
\]
\underline{We claim that this element is precisely}
$\partial^0(u)$.

We prove this compatibility by reducing it to the analogous one in
[4, Cor.~5.5], where the sign must be corrected as indicated in
3.3.6.3.  The argument also gives another proof of the equivalence
(\SGARefLink{sga7.sga7i.evii.n3-4-8-2}{3.4.8.2}), by reducing it to [4, Prop.~\SGARefLink{sga7.sga7i.evii.n2-7}{2.7}].  Consider the opposite
category $\underline C^{\mathrm o}$ and interpret the construction
of $\underline\Psi_{L_\bullet}(A)$, for fixed $A$, together with
the constructions of \SGARefLink{sga7.sga7i.evii.n3-4-8}{3.4.8} and the present section, in terms of
$\underline C^{\mathrm o}$ and the category $\underline E$ cofibred
over $\underline C^{\mathrm o}$ whose objects are extensions of the
fixed object $A$ of $\underline C^{\mathrm o}$ by a variable object
$L$ of $\underline C^{\mathrm o}$.  This is an additive cofibred
category over $\underline C^{\mathrm o}$, and is even left exact
there (in the sense of [4, \SGARefLink{sga7.sga7i.evii.n2-2}{2.2}]).

A chain complex $L_\bullet$ of $\underline C$ may be interpreted as
a cochain complex of $\underline C^{\mathrm o}$.  Under this
interpretation, the construction of $\underline\Psi_{L_\bullet}(A)$
coincides with the opposite of the category
$\underline{\widehat E}(L_\bullet)$ in the sense of [4, 1.8].  The
functor (\SGARefLink{sga7.sga7i.evii.n3-4-8-2}{3.4.8.2}) is then identified with the opposite of the functor
of the form [4, (2.1.2)], with $E$ replaced by its extension
$\underline{\widehat E}$ to the category of cochain complexes.
Consequently, the assertion that this functor is an equivalence is
the special case [4, \SGARefLink{sga7.sga7i.evii.n2-7}{2.7}, a)$\Rightarrow$d)].

On the other hand, the exact sequence (\SGARefLink{sga7.sga7i.evii.n3-4-9-1}{3.4.9.1}) is identified with
the exact sequence of the groups $\Ext^i(A,L''_\bullet)$,
$\Ext^i(A,L'_\bullet)$, and $\Ext^i(A,L_\bullet)$, apart from a
factor $(-1)^i$ on $\partial^i$ (Verdier dixit).  Similarly, the
isomorphisms (\SGARefLink{sga7.sga7i.evii.n3-1-8}{3.1.8}) and (\SGARefLink{sga7.sga7i.evii.n3-2-4}{3.2.4}) are identified with the analogous
isomorphisms (3.11.5) and (3.11.4).  The asserted compatibility is
therefore the special case of [4, Cor.~5.5] corrected as above.

\underline{Remark} \SGARefTarget{sga7.sga7i.evii.n3-4-9-2}{3.4.9.2}. The editor, being very afraid of signs,
confesses that he does not have complete confidence in the preceding
proof: he fears that signs may have been omitted in passing from
$C$ to $C^{\mathrm o}$.  He recommends that a courageous reader
verify the asserted compatibility directly (or, respectively, its
opposite).  Fortunately, these sign questions will play no role in
the following expos\'es.

% END INCLUDED SGA7I_English_current_buildable_source_20260823/components/88_expose_VII_sections_3_4_8_and_3_4_9.tex

% BEGIN INCLUDED SGA7I_English_current_buildable_source_20260823/components/89_expose_VII_section_3_5.tex
% SOURCE AUTHORITY:
% expose_VII_body.tex lines 1545--1662; scan indices 209--215; source folios 198--204.
% Direct scan comparison: physical PDF pages 210--216 at 600 dpi.
% Transparent corrections: the first overbrace in (3.5.1) is L_2 as printed, not
% the working transcription's L_o; chain subscripts o are rendered as zero; the
% chain-complex relation is written in compositional order d_0d_1=0 throughout;
% and the display used later as (*) receives the tag visible in the scan.
% Hard stop before source line 1668, section 3.6.

\subsection*{\SGARefTarget{sga7.sga7i.evii.n3-5}{3.5}. A canonical partial flat resolution of an
arbitrary $\mathbb Z$-module}

We apply the preceding results when $\underline C$ is the category
of commutative groups of the topos $\underline T$, that is, the
category of $\mathbb Z$-modules of $\underline T$.  Let $P$ be an
object of $\underline C$.  We shall define, functorially in $P$, a
chain complex $L_\bullet(P)$ augmented to $P$:
\begin{equation}\tag{\SGARefTarget{sga7.sga7i.evii.n3-5-1}{3.5.1}}
\begin{tikzcd}[column sep=small,row sep=large]
 0\arrow[r]&\overbrace{\mathbb Z[P\times P]\times
   \mathbb Z[P\times P\times P]}^{L_2}
   \arrow[r,"d_1"]
 &\overbrace{\mathbb Z[P\times P]}^{L_1}
   \arrow[r,"d_0"]
 &\overbrace{\mathbb Z[P]}^{L_0}
   \arrow[r]\arrow[d,"\varepsilon"']&0\\
 &&&P&{}
\end{tikzcd}
\end{equation}
with $L_i(P)=0$ for $i\geq3$.  The components are displayed in
(\SGARefLink{sga7.sga7i.evii.n3-5-1}{3.5.1}), using the notation of \SGARefLink{sga7.sga7i.evii.n1-4}{1.4}.  To describe the differentials
$d_0,d_1$ and the augmentation $\varepsilon$, write $[p]$ for the
section of $\mathbb Z[P](S)$ defined by $p\in P(S)$, and similarly
write $[p,q]$ and $[p,q,r]$ for the corresponding sections of
$\mathbb Z[P\times P](S)$ and
$\mathbb Z[P\times P\times P](S)$.  Then set
\begin{equation}\tag{\SGARefTarget{sga7.sga7i.evii.n3-5-2}{3.5.2}}
\left\{
\begin{aligned}
 \varepsilon[p]&=p,\\
 d_0[p,q]&=[p+q]-[p]-[q],\\
 d_1[p,q]&=[p,q]-[q,p],\\
 d_1[p,q,r]&=[p+q,r]-[p,q+r]+[p,q]-[q,r].
\end{aligned}
\right.
\end{equation}
These homomorphisms of commutative groups do define a complex
$L_\bullet(P)$ augmented to $P$:
\[
 \varepsilon d_0=0,\qquad d_0d_1=0.
\]
The first relation is immediate.  The second splits into two
relations according to the two factors of $L_2(P)$; they express,
respectively, commutativity and associativity of the composition law
on $P$.

The augmented complex depends functorially on $P$.  A morphism
$u:P\longrightarrow P'$ gives a commutative diagram
\[
\begin{tikzcd}[column sep=huge,row sep=large]
 L_\bullet(P)\arrow[r,"L_\bullet(u)"]\arrow[d]
 &L_\bullet(P')\arrow[d]\\
 P\arrow[r,"u"]&P'.
\end{tikzcd}
\]
Its components are flat, as is every free $\mathbb Z$-module.

\underline{Proposition} \SGARefTarget{sga7.sga7i.evii.n3-5-3}{3.5.3}. \underline{One has}
$H_1(L_\bullet(P))=0$, \underline{and the augmentation}
$\varepsilon$ \underline{induces an isomorphism}
\[
 H_0(L_\bullet(P))\xrightarrow{\sim}P.
\]

For the proof one can make the standard reduction to the case in
which $\underline T$ is the category of sets and $P$ is an ordinary
$\mathbb Z$-module, where (\SGARefLink{sga7.sga7i.evii.n3-5-1}{3.5.1}) is likely familiar.  We instead
give a direct proof that does not use this reduction, applying \SGARefLink{sga7.sga7i.evii.n3-4-7}{3.4.7}
to the augmentation
\[
 \varepsilon:L_\bullet(P)\longrightarrow P,
\]
where $P$ is regarded, as usual, as a chain complex concentrated in
degree zero.  It remains to prove that, for every object $G$ of
$\underline C$, the functor
\begin{equation}\tag{*}
 \varepsilon^*:\underline\Psi_P(G)\longrightarrow
 \underline\Psi_{L_\bullet(P)}(G)
\end{equation}
is an equivalence of categories.

The source of (*) is plainly isomorphic to the category of
extensions of the $\mathbb Z$-module $P$ by $G$.  To describe its
target, use \SGARefLink{sga7.sga7i.evii.n1-4}{1.4} to identify the categories of extensions of
$L_0,L_1,L_2$ by $G$, respectively, with the category of torsors
under $G_P$, the category of torsors under $G_{P\times P}$, and the
category of pairs consisting of a torsor under $G_{P\times P}$ and
a torsor under $G_{P\times P\times P}$.

Thus an object of the target of (*) is a pair $(E,\varphi)$, where
$E$ is a torsor under $G_P$ and $\varphi$ is a trivialization of the
torsor $E'$ on $P\times P$ corresponding to the inverse image along
$d_0$ of the extension of $\mathbb Z[P]$ by $G$ defined by $E$.
The trivialization is subject to the compatibility with
$d_0d_1=0$ described below.  Formula (\SGARefLink{sga7.sga7i.evii.n3-5-2}{3.5.2}) immediately gives
\[
 E'=\pi^*E\,\pr_1^*(E)^{-1}\,\pr_2^*(E)^{-1},
\]
and a trivialization of this torsor is the same as an isomorphism,
also denoted by $\varphi$,
\[
 \varphi:\pr_1^*E\,\pr_2^*E\xrightarrow{\sim}\pi^*E.
\]
Expanding the compatibility of $\varphi$ with $d_0d_1=0$ gives the
commutativity condition (\SGARefLink{sga7.sga7i.evii.n1-2-1}{1.2.1}) and the associativity condition
(\SGARefLink{sga7.sga7i.evii.n1-1-4-1}{1.1.4.1}), corresponding to the two factors of $L_2$.  The assertion
that (*) is an equivalence is therefore precisely the dictionary of
\SGARefLink{sga7.sga7i.evii.n1-1-6}{1.1.6}.  This proves \SGARefLink{sga7.sga7i.evii.n3-5-3}{3.5.3}.

\underline{Remark} \SGARefTarget{sga7.sga7i.evii.n3-5-4}{3.5.4}. a) It would be interesting to strengthen
\SGARefLink{sga7.sga7i.evii.n3-5-3}{3.5.3} by constructing a resolution $L_\bullet(P)$ of $P$, functorial
in $P$, whose components are functorially isomorphic to sums of free
$\mathbb Z$-modules generated by cartesian powers of $P$.  In view
of isomorphisms of the form (\SGARefLink{sga7.sga7i.evii.n1-4-5}{1.4.5}), this would give a method for
partially computing $\Ext^i(P,A)$ as the abutment of a spectral
sequence whose initial term is expressed through groups
\[
 H^j(P\times P\times\cdots\times P,A).
\]
Here we obtain only very partial information about $\Ext^1$, and no
information about the higher $\Ext^i$.  There is, in any case, a
natural way to continue the partial resolution $L_\bullet(P)$ by one
degree.\footnote{L. Breen constructed such a resolution (editor's
note).}

b) There is a variant of (\SGARefLink{sga7.sga7i.evii.n3-5-1}{3.5.1}) and \SGARefLink{sga7.sga7i.evii.n3-5-3}{3.5.3} in the category of
$A$-modules, where $A$ is a fixed ring of $\underline T$.  The form
of $L_\bullet(P)$ is more complicated because the description \SGARefLink{sga7.sga7i.evii.n1-1-6}{1.1.6}
must be enlarged (compare \SGARefLink{sga7.sga7i.evii.n2-10-3}{2.10.3}).  One obtains
\begin{equation}\tag{\SGARefTarget{sga7.sga7i.evii.n3-5-5}{3.5.5}}
\left\{
\begin{aligned}
 L_0={}&A[P],\\
 L_1={}&A[P\times P]+A[A\times P],\\
 L_2={}&A[P\times P\times P]+A[P\times P]
       +A[A\times A\times P]+A[A\times P\times P].
\end{aligned}
\right.
\end{equation}
The augmented-complex structure is given by (\SGARefLink{sga7.sga7i.evii.n3-5-2}{3.5.2}), supplemented by
\begin{equation}\tag{\SGARefTarget{sga7.sga7i.evii.n3-5-6}{3.5.6}}
\left\{
\begin{aligned}
 d_0[\lambda,p]={}&[\lambda p]-\lambda[p],\\
 d_1[\lambda,\mu,p]={}&[\lambda,\mu p]-[\lambda\mu,p]
                         +\lambda[\mu,p],\\
 d_1[\lambda,p,q]={}&[\lambda,p+q]-[\lambda,p]-[\lambda,q]
                       +\lambda[p,q]-[\lambda p,\lambda q].
\end{aligned}
\right.
\end{equation}
The relation $d_0d_1=0$ follows from the analogous relation on the
components appearing in (\SGARefLink{sga7.sga7i.evii.n3-5-2}{3.5.2}), which, as we saw, is equivalent to
associativity and commutativity of addition on $P$, together with
the relations
\[
 d_0d_1[\lambda,\mu,p]=0,
 \qquad d_0d_1[\lambda,p,q]=0.
\]
The latter express, respectively, associativity of the given action
of $A$ on $P$,
\[
 \lambda(\mu p)=(\lambda\mu)p,
\]
and its distributivity over addition,
\[
 \lambda(p+q)=\lambda p+\lambda q.
\]

For a torsor $E$ under $G_P$, a trivialization of $d_0^*E$ amounts
to an internal composition law on $E$ (satisfying the compatibility
conditions already described for the internal law on $E$ and its
torsor structure), together with an external composition law of $A$
on $E$, compatible with the action of $A$ on $P$ and on $G$.  The
compatibility of this trivialization with the evident trivialization
of $d_1^*d_0^*E$ is equivalent to the conditions already described
above, which say that the internal law on $E$ is a commutative-group
law extending $P$ by $G$, together with associativity of the external
law and distributivity of the latter with respect to the internal
law.  These conditions say exactly that $E$, with its two
composition laws and its torsor structure, is an $A$-module extension
of the $A$-modules $P$ and $G$.

It follows, by essentially the same method as in \SGARefLink{sga7.sga7i.evii.n3-5-3}{3.5.3}, that
$L_\bullet$ is a partial resolution of the module $P$, that is,
$H_1(L_\bullet)=0$.  It would again be interesting to extend this
partial resolution to an infinite resolution of the same kind, with
components of the form
\[
 A[A\times A\times\cdots\times P\times P\times\cdots].
\]
Compare the analogous question in a) for commutative groups.

% END INCLUDED SGA7I_English_current_buildable_source_20260823/components/89_expose_VII_section_3_5.tex

% BEGIN INCLUDED SGA7I_English_current_buildable_source_20260823/components/90_expose_VII_section_3_6.tex
% SOURCE AUTHORITY:
% expose_VII_body.tex lines 1668--1726; scan indices 215--219; source folios 204--208.
% Direct scan comparison: physical PDF pages 216--220 at 600 dpi.
% Transparent corrections: (3.6.3) and Theorem 3.6.4 use L'_\bullet(P,Q), not
% the impossible L_\bullet(P',Q'); extension notation uses L_0(P,Q); the induced
% maps from P x Q x Q have target P x Q; and the extension of P_Q is by G_Q.
% Typewritten letter-o subscripts are rendered as zero.
% Hard stop before source line 1729, section 3.7.

\subsection*{\SGARefTarget{sga7.sga7i.evii.n3-6}{3.6}. Application to the description of the cofibred
category of biextensions of a given pair $(P,Q)$ of commutative
groups}

Let $P$ and $Q$ be two $\mathbb Z$-modules of the topos
$\underline T$.  Choose a flat resolution $L'_\bullet(P)$ of $P$
that agrees in degrees at most $2$ with the augmented complex
$L_\bullet(P)$ of \SGARefLink{sga7.sga7i.evii.n3-5}{3.5}.  This is possible by \SGARefLink{sga7.sga7i.evii.n3-5-3}{3.5.3} and because every
$\mathbb Z$-module on $\underline T$ is a quotient of a free, hence
flat, $\mathbb Z$-module.  Similarly choose a flat resolution
$L'_\bullet(Q)$ of $Q$ agreeing with $L_\bullet(Q)$ in degrees at
most $2$.  We have canonical homomorphisms
\begin{equation}\tag{\SGARefTarget{sga7.sga7i.evii.n3-6-1}{3.6.1}}
 L_\bullet(P)\longrightarrow L'_\bullet(P),
 \qquad
 L_\bullet(Q)\longrightarrow L'_\bullet(Q),
\end{equation}
and hence a canonical homomorphism
\begin{equation}\tag{\SGARefTarget{sga7.sga7i.evii.n3-6-2}{3.6.2}}
\begin{split}
 L_\bullet(P,Q)
 &\overset{\mathrm{def}}{=}L_\bullet(P)\otimes L_\bullet(Q)
 \longrightarrow L'_\bullet(P,Q)\\
 &\overset{\mathrm{def}}{=}L'_\bullet(P)\otimes L'_\bullet(Q)
 \simeq P\mathbin{\overset{\mathbb L}{\otimes}}Q.
\end{split}
\end{equation}
It is an isomorphism on components in degrees at most $2$.
Consequently, by definition,
\begin{equation}\tag{\SGARefTarget{sga7.sga7i.evii.n3-6-3}{3.6.3}}
 \underline\Psi_{L_\bullet(P,Q)}
 =
 \underline\Psi_{L'_\bullet(P,Q)}.
\end{equation}

\underline{Theorem} \SGARefTarget{sga7.sga7i.evii.n3-6-4}{3.6.4}. \underline{The additive cofibred category
over the category}
\[
 \underline C=\operatorname{GRAB}(\underline T)
\]
\underline{of commutative groups of $\underline T$, consisting of
biextensions of $(P,Q)$ by a variable abelian group $G$} (\SGARefLink{sga7.sga7i.evii.n2-4}{2.4}),
\underline{is canonically equivalent to the additive cofibred
category}
\[
 \underline\Psi_{L_\bullet(P,Q)}
 =
 \underline\Psi_{L'_\bullet(P,Q)}
\]
\underline{defined from}
$L_\bullet(P,Q)=L_\bullet(P)\otimes L_\bullet(Q)$
\underline{as in \SGARefLink{sga7.sga7i.evii.n3-1}{3.1}}.

This and \SGARefLink{sga7.sga7i.evii.n3-2-5}{3.2.5} give the principal result of the present section,
announced above:

\underline{Corollary} \SGARefTarget{sga7.sga7i.evii.n3-6-5}{3.6.5}. \underline{For fixed $P$ and $Q$ there
is a canonical isomorphism, functorial in $G$,}
\begin{equation}\tag{\SGARefTarget{sga7.sga7i.evii.n3-6-5-1}{3.6.5.1}}
 \Biext^1(P,Q;G)
 \xrightarrow{\sim}
 \Ext^1(P\mathbin{\overset{\mathbb L}{\otimes}}Q,G).
\end{equation}

To prove \SGARefLink{sga7.sga7i.evii.n3-6-5}{3.6.5}, first make explicit the components of degrees
$0,1,2$ of $L_\bullet(P,Q)$:
\begin{equation}\tag{\SGARefTarget{sga7.sga7i.evii.n3-6-6}{3.6.6}}
\left\{
\begin{aligned}
 L_0(P,Q)
  ={}&L_0(P)\otimes L_0(Q)
     =\mathbb Z[P\times Q],\\
 L_1(P,Q)
  ={}&L_0(P)\otimes L_1(Q)+L_1(P)\otimes L_0(Q)\\
  ={}&\mathbb Z[P\times Q\times Q]
     +\mathbb Z[P\times P\times Q],\\
 L_2(P,Q)
  ={}&L_0(P)\otimes L_2(Q)+L_1(P)\otimes L_1(Q)
     +L_2(P)\otimes L_0(Q)\\
  ={}&\mathbb Z[P\times Q\times Q]
     +\mathbb Z[P\times Q\times Q\times Q]\\
   &+\mathbb Z[P\times P\times Q\times Q]\\
   &+\mathbb Z[P\times P\times Q]
     +\mathbb Z[P\times P\times P\times Q].
\end{aligned}
\right.
\end{equation}

To describe the cofibred category
$\underline\Psi_{L_\bullet(P,Q)}$ through its fibres
$\underline\Psi_{L_\bullet(P,Q)}(G)$, use \SGARefLink{sga7.sga7i.evii.n1-4}{1.4} to express extensions
of $L_i(P,Q)$, for $i=0,1,2$, by $G$ in terms of torsors.  An
extension of $L_0(P,Q)$ by $G$ is a torsor under $G_{P\times Q}$.
An extension of $L_1(P,Q)$ by $G$ is a pair consisting of a torsor
under $G_{P\times Q\times Q}$ and one under
$G_{P\times P\times Q}$.  An extension of $L_2(P,Q)$ by $G$ is
similarly a system of five torsors, under the groups obtained from
$G$ by base change to
\[
\begin{gathered}
 P\times Q\times Q,\quad
 P\times Q\times Q\times Q,\quad
 P\times P\times Q\times Q,\\
 P\times P\times Q,\quad
 P\times P\times P\times Q.
\end{gathered}
\]

Under this dictionary, a pair $(X,\alpha)$ consisting of an
extension of $L_0(P,Q)$ by $G$ and a trivialization of
$d_0^*X$ on $L_1(P,Q)$ becomes a torsor $E$ on $P\times Q$ equipped
with two trivializations $(\varphi,\psi)$ of the torsors
$(E_1,E_2)$ on $P\times Q\times Q$ and
$P\times P\times Q$, respectively, obtained from $E$ by inverse
image along
\[
 d_0:L_1(P,Q)\longrightarrow L_0(P,Q).
\]
Expanding $d_0$, as in the proof of \SGARefLink{sga7.sga7i.evii.n3-5-3}{3.5.3}, identifies
$\varphi$ and $\psi$ with isomorphisms
\[
 \varphi:q_1^*E\,q_2^*E\xrightarrow{\sim}\pi_Q^*E,
 \qquad
 \psi:p_1^*E\,p_2^*E\xrightarrow{\sim}\pi_P^*E.
\]
Here $\pi_Q,q_1,q_2:P\times Q\times Q\longrightarrow P\times Q$
are induced by the composition law and the first and second
projections $Q\times Q\longrightarrow Q$; the maps
\[
 \pi_P,p_1,p_2:P\times P\times Q\longrightarrow P\times Q
\]
are defined symmetrically.  Thus $\varphi$ and $\psi$ are data of
the forms (\SGARefLink{sga7.sga7i.evii.n2-0-4}{2.0.4}) and (\SGARefLink{sga7.sga7i.evii.n2-0-7}{2.0.7}) appearing in \SGARefLink{sga7.sga7i.evii.n2-1}{2.1}.

Finally, the condition that $(X,\alpha)$ be an object of
$\underline\Psi_{L_\bullet(P,Q)}(G)$, namely compatibility of the
partial trivialization $\alpha$ with $d_0d_1=0$, becomes five
compatibility relations for the structure $(E,\varphi,\psi)$ on
torsors over, respectively,
\[
\begin{gathered}
 P\times Q\times Q,\quad
 P\times Q\times Q\times Q,\quad
 P\times P\times Q\times Q,\\
 P\times P\times Q,\quad
 P\times P\times P\times Q.
\end{gathered}
\]
They correspond to the five terms in the expression (\SGARefLink{sga7.sga7i.evii.n3-6-6}{3.6.6}) for
$L_2(P,Q)$.  As in the proof of \SGARefLink{sga7.sga7i.evii.n3-5-3}{3.5.3}, the first two are precisely
the commutativity relation (\SGARefLink{sga7.sga7i.evii.n2-0-6}{2.0.6}) and the associativity relation
(\SGARefLink{sga7.sga7i.evii.n2-0-5}{2.0.5}); together they say that $\varphi$ gives $E$ the structure
of a commutative extension of $Q_P$ by $G_P$.  Symmetrically, the
last two say that $\psi$ gives $E$ the structure of a commutative
extension of $P_Q$ by $G_Q$.  The middle relation, on torsors over
$P\times P\times Q\times Q$, is exactly the commutativity of
diagram (\SGARefLink{sga7.sga7i.evii.n2-1-1}{2.1.1}).  It says that $\varphi$ and $\psi$ are compatible,
hence, in view of the other conditions, define on $E$ a
biextension structure of $(P,Q)$ by $G$.

The verification of these interpretations is immediate and is left
to the reader.  The preceding discussion indeed defines an
equivalence of categories cofibred over
$\underline C=\operatorname{GRAB}(\underline T)$, as asserted in
\SGARefLink{sga7.sga7i.evii.n3-6-4}{3.6.4}.

\underline{Remark} \SGARefTarget{sga7.sga7i.evii.n3-6-7}{3.6.7}. Each of the three generalizations of
biextensions mentioned in \SGARefLink{sga7.sga7i.evii.n2-10}{2.10} has a corresponding variant of \SGARefLink{sga7.sga7i.evii.n3-6-4}{3.6.4}
and hence of \SGARefLink{sga7.sga7i.evii.n3-6-5}{3.6.5}.  The generalization to any number of factors
(\SGARefLink{sga7.sga7i.evii.n2-10-2}{2.10.2}) is immediate, as is the generalization to biextensions of
$A$-modules (\SGARefLink{sga7.sga7i.evii.n2-10-3}{2.10.3}), once the variant of $L_\bullet(P)$ described
in \SGARefLink{sga7.sga7i.evii.n3-5-4}{3.5.4} b) has been constructed.  From the viewpoint of
homological algebra, the case in which $P$ and $Q$ are not assumed
commutative (\SGARefLink{sga7.sga7i.evii.n2-10-1}{2.10.1}) is more interesting.

Associate with a not necessarily commutative group $P$ of
$\underline T$ a complex $K_\bullet(P)$ truncated in degree $2$.
It is described as $L_\bullet(P)$ was in \SGARefLink{sga7.sga7i.evii.n3-5}{3.5}, except that
\[
 K_2(P)=\mathbb Z[P\times P\times P],
\]
the first factor occurring in $L_2(P)$ being omitted; $d_0$ and
$d_1$ are defined as in (\SGARefLink{sga7.sga7i.evii.n3-5-2}{3.5.2}).  Then
\[
 H_0(K_\bullet(P))=P/[P,P]
 \simeq\underline H_1(P,\mathbb Z),
\]
and similarly
\[
 H_1(K_\bullet(P))\simeq\underline H_2(P,\mathbb Z).
\]
Indeed, $K_\bullet(P)$ is precisely the truncation in degree $2$ of
the chain complex $C_\bullet(P,\mathbb Z)$ of $P$ with coefficients
in $\mathbb Z$.

The argument proving \SGARefLink{sga7.sga7i.evii.n3-5-3}{3.5.3} also shows that the cofibred category
$\underline\Psi_{K_\bullet(P)}$ over
$\underline C=\operatorname{GRAB}(\underline T)$ is canonically
equivalent to the category of central extensions of $P$ by variable
commutative groups $G$.  Likewise, for two groups $P,Q$, the
cofibred category of biextensions of $(P,Q)$ in the sense of \SGARefLink{sga7.sga7i.evii.n2-10-1}{2.10.1}
is equivalent to $\underline\Psi_{K_\bullet(P,Q)}$, where
\[
 K_\bullet(P,Q)
 =
 K_\bullet(P)\mathbin{\overset{\mathbb L}{\otimes}}K_\bullet(Q).
\]

These considerations suggest that, when $P$ is commutative, the
canonical resolution sought in \SGARefLink{sga7.sga7i.evii.n3-5-4}{3.5.4} a) might be obtained by a
suitable modification of $C_\bullet(P,\mathbb Z)$.

% END INCLUDED SGA7I_English_current_buildable_source_20260823/components/90_expose_VII_section_3_6.tex

% BEGIN INCLUDED SGA7I_English_current_buildable_source_20260823/components/91_expose_VII_section_3_7.tex
% SOURCE AUTHORITY:
% expose_VII_body.tex lines 1729--1804; scan indices 219--223; source folios 208--212.
% Direct scan comparison: physical PDF pages 220--224 at 600 dpi.
% Transparent normalization: typewritten letter-o superscripts are rendered as zero;
% the third factor of the base category in Proposition 3.7.6 is covariant, exactly as
% printed; and the chain-complex subscripts use zero.  Hard stop before source line
% 1810, section 3.8.

\subsection*{\SGARefTarget{sga7.sga7i.evii.n3-7}{3.7}. Various compatibilities}

In this section we examine the behaviour, for variable $P$ and $Q$,
of the equivalence of cofibred categories in \SGARefLink{sga7.sga7i.evii.n3-6-4}{3.6.4}.  The variance of
the cofibred category of biextensions of $(P,Q)$ by a variable $G$
with respect to variable $P,Q$ is defined in principle in \SGARefLink{sga7.sga7i.evii.n2-4}{2.4}; it may
also be expressed by saying that one has a cofibred functor (\SGARefLink{sga7.sga7i.evii.n2-6-1}{2.6.1}).
Notice that the variance in $P,Q$ is \underline{contravariant}.  The
variance of the cofibred category
$\underline\Psi_{L_\bullet(P,Q)}$ with respect to variable $P,Q$ is
also clear, since
$L_\bullet(P,Q)=L_\bullet(P)\otimes L_\bullet(Q)$ depends
functorially on the pair $(P,Q)$, while $\underline\Psi_{L_\bullet}$
depends contravariantly on $L_\bullet$ (\SGARefLink{sga7.sga7i.evii.n3-4}{3.4}).  With this understood,
a first compatibility is expressed, for every triple of morphisms
\begin{equation}\tag{\SGARefTarget{sga7.sga7i.evii.n3-7-1}{3.7.1}}
 P'\longrightarrow P,\qquad
 Q'\longrightarrow Q,\qquad
 G\longrightarrow G',
\end{equation}
by the fact that the diagram of functors
\begin{equation}\tag{\SGARefTarget{sga7.sga7i.evii.n3-7-2}{3.7.2}}
\begin{tikzcd}[column sep=huge,row sep=large]
 \underline\Psi_{L_\bullet(P,Q)}(G)
   \arrow[r]\arrow[d]
 &\underline\Psi_{L_\bullet(P',Q')}(G')
   \arrow[d]\\
 \operatorname{BIEXT}(P,Q;G)
   \arrow[r]
 &\operatorname{BIEXT}(P',Q';G')
\end{tikzcd}
\end{equation}
is commutative up to a canonical isomorphism.  We shall admit this
compatibility: making the commutativity isomorphism explicit could
only be as obvious as it is tedious (as, indeed, could the complete
explication it presupposes of the equivalence of cofibred categories
in \SGARefLink{sga7.sga7i.evii.n3-6-4}{3.6.4}).

\SGARefTarget{sga7.sga7i.evii.n3-7-3}{3.7.3}. It follows in particular from the preceding compatibility
that the isomorphism (\SGARefLink{sga7.sga7i.evii.n3-6-5-1}{3.6.5.1}) is functorial not only in $G$, but also
in $P$ and $Q$.  One may state, similarly and more trivially, the
analogous functoriality of the isomorphism defined through (\SGARefLink{sga7.sga7i.evii.n3-1-8}{3.1.8}),
just as (\SGARefLink{sga7.sga7i.evii.n3-6-5-1}{3.6.5.1}) was defined through (\SGARefLink{sga7.sga7i.evii.n3-2-4}{3.2.4}):
\begin{equation}\tag{\SGARefTarget{sga7.sga7i.evii.n3-7-4}{3.7.4}}
 \Biext^0(P,Q;G)
 \xrightarrow{\sim}
 \Ext^0(P\mathbin{\overset{\mathbb L}{\otimes}}Q,G)
 \simeq\Hom(P\otimes Q,G).
\end{equation}
But such an isomorphism was already obtained in (\SGARefLink{sga7.sga7i.evii.n2-5-4-1}{2.5.4.1}), and was
observed in \SGARefLink{sga7.sga7i.evii.n2-6}{2.6} to be functorial in its three arguments.  It is
immediate to verify that the two isomorphisms (\SGARefLink{sga7.sga7i.evii.n3-7-4}{3.7.4}) and (\SGARefLink{sga7.sga7i.evii.n2-5-4-1}{2.5.4.1})
are equal.

\SGARefTarget{sga7.sga7i.evii.n3-7-5}{3.7.5}. The value of the isomorphism in \SGARefLink{sga7.sga7i.evii.n3-6-5}{3.6.5} is that it expresses the
remarkable group of geometric origin $\Biext^1(P,Q;G)$ in terms of
the very simple homological-algebra invariant
$\Ext^1(P\mathbin{\overset{\mathbb L}{\otimes}}Q,G)$.  The latter
belongs to the sequence of invariants
$\Ext^i(P\mathbin{\overset{\mathbb L}{\otimes}}Q,G)$, to which the
usual formalism of homological algebra applies.  In particular, the
sequence of these $\Ext^i$ is a \underline{cohomological functor} in
the sense of [3] (or a \underline{connected sequence of functors} in
the terminology of [5]) in each of the three arguments $P,Q,G$.
Thus an exact sequence
\begin{equation}\tag{\SGARefTarget{sga7.sga7i.evii.n3-7-5-1}{3.7.5.1}}
 0\longrightarrow G'\longrightarrow G\longrightarrow G''
 \longrightarrow0
\end{equation}
(respectively,
\begin{equation}\tag{\SGARefTarget{sga7.sga7i.evii.n3-7-5-2}{3.7.5.2}}
 0\longrightarrow P'\longrightarrow P\longrightarrow P''
 \longrightarrow0,\qquad\text{or}\qquad
 0\longrightarrow Q'\longrightarrow Q\longrightarrow Q''
 \longrightarrow0)
\end{equation}
gives rise to an infinite exact sequence of the $\Ext^i$.  Its first
six terms, which involve only $\Ext^0$ and $\Ext^1$, may be
interpreted in terms of biextensions by \SGARefLink{sga7.sga7i.evii.n3-6-5}{3.6.5} and \SGARefLink{sga7.sga7i.evii.n3-7-4}{3.7.4}.  We have
already indicated in \SGARefLink{sga7.sga7i.evii.n3-3-6}{3.3.6} a geometric interpretation of the exact
sequence deduced from (\SGARefLink{sga7.sga7i.evii.n3-7-5-1}{3.7.5.1}), and especially of the boundary map
\begin{equation}\tag{\SGARefTarget{sga7.sga7i.evii.n3-7-5-3}{3.7.5.3}}
 \partial:\Biext^0(P,Q;G'')\longrightarrow
 \Biext^1(P,Q;G'),
\end{equation}
within the framework of a \underline{left exact cofibred category}.
One may seek an analogous interpretation of the six-term exact
sequence associated with either short exact sequence in (\SGARefLink{sga7.sga7i.evii.n3-7-5-2}{3.7.5.2}),
say the first; equivalently, in view of \SGARefLink{sga7.sga7i.evii.n3-7-3}{3.7.3}, one may seek a
geometric interpretation of the boundary homomorphism
\begin{equation}\tag{\SGARefTarget{sga7.sga7i.evii.n3-7-5-4}{3.7.5.4}}
 \partial:\Biext^0(P',Q;G)\longrightarrow
 \Biext^1(P'',Q;G).
\end{equation}
Indeed, we have the following result.

\underline{Proposition} \SGARefTarget{sga7.sga7i.evii.n3-7-6}{3.7.6}. \underline{The cofibred category over}
\[
 \operatorname{GRAB}(\underline T)^{\mathrm{op}}
 \times\operatorname{GRAB}(\underline T)^{\mathrm{op}}
 \times\operatorname{GRAB}(\underline T)
\]
\underline{of biextensions of $(P,Q)$ by $G$, with $P,Q,G$
variable} (cf. (\SGARefLink{sga7.sga7i.evii.n2-6-1}{2.6.1})), \underline{is left exact not only in the
argument $G$, but also in each of the arguments $P$ and $Q$.
Moreover, the boundary homomorphism (\SGARefLink{sga7.sga7i.evii.n3-7-5-4}{3.7.5.4}) deduced from the exact
sequence of the $\Ext^i$ associated with (\SGARefLink{sga7.sga7i.evii.n3-7-5-2}{3.7.5.2}), by means of the
isomorphisms \SGARefLink{sga7.sga7i.evii.n3-6-5}{3.6.5} and \SGARefLink{sga7.sga7i.evii.n3-7-4}{3.7.4}, is the boundary homomorphism in the
theory of left exact cofibred categories over}
$\underline C=\operatorname{GRAB}(\underline T)^{\mathrm{op}}$
\underline{recalled in \SGARefLink{sga7.sga7i.evii.n3-3-6}{3.3.6}, applied to the cofibred category of
biextensions of $(P,Q)$ by $G$, where $Q$ and $G$ are fixed and $P$
varies} (subject to a possible sign error in \SGARefLink{sga7.sga7i.evii.n3-4-9}{3.4.9}; cf. \SGARefLink{sga7.sga7i.evii.n3-4-9-2}{3.4.9.2}).

We sketch the proof of \SGARefLink{sga7.sga7i.evii.n3-7-6}{3.7.6}.  By \SGARefLink{sga7.sga7i.evii.n3-6-4}{3.6.4} and compatibility (\SGARefLink{sga7.sga7i.evii.n3-7-2}{3.7.2}),
we interpret biextensions of $(P,Q)$ by $G$ as objects of
$\underline\Psi_{L_\bullet(P,Q)}(G)$.  Consider the homomorphism
\[
 L_\bullet(P,Q)
 =L_\bullet(P)\otimes L_\bullet(Q)
 \longrightarrow P\otimes L_\bullet(Q)
\]
defined by the augmentation $L_\bullet(P)\longrightarrow P$.
It follows from \SGARefLink{sga7.sga7i.evii.n3-5-3}{3.5.3} and the flatness of $L_\bullet(Q)$ that this
homomorphism induces isomorphisms on $H_0$ and $H_1$, and hence by
\SGARefLink{sga7.sga7i.evii.n3-4-7}{3.4.7} an equivalence
\[
 \underline\Psi_{P\otimes L_\bullet(Q)}
 \longrightarrow
 \underline\Psi_{L_\bullet(P,Q)}.
\]
We may therefore interpret biextensions of $(P,Q)$ by $G$ as objects
of $\underline\Psi_{P\otimes L_\bullet(Q)}(G)$.  The exact sequence
(\SGARefLink{sga7.sga7i.evii.n3-7-5-2}{3.7.5.2}) now gives an exact sequence of chain complexes
\begin{equation}\tag{\SGARefTarget{sga7.sga7i.evii.n3-7-6-1}{3.7.6.1}}
 0\longrightarrow P'\otimes L_\bullet(Q)
 \longrightarrow P\otimes L_\bullet(Q)
 \longrightarrow P''\otimes L_\bullet(Q)
 \longrightarrow0.
\end{equation}
The asserted left exactness follows from \SGARefLink{sga7.sga7i.evii.n3-4-8}{3.4.8} applied to this exact
sequence of complexes.  (A direct verification would in any event
be immediate.)  Moreover, the homomorphism (\SGARefLink{sga7.sga7i.evii.n3-7-5-4}{3.7.5.4}) is precisely the
boundary homomorphism $\Ext^0\longrightarrow\Ext^1$ deduced from this
exact sequence by taking the $\Ext^i$ with values in $G$.  The last
assertion of \SGARefLink{sga7.sga7i.evii.n3-7-6}{3.7.6} is consequently a special case of \SGARefLink{sga7.sga7i.evii.n3-4-9}{3.4.9}.

% END INCLUDED SGA7I_English_current_buildable_source_20260823/components/91_expose_VII_section_3_7.tex

% BEGIN INCLUDED SGA7I_English_current_buildable_source_20260823/components/92_expose_VII_section_3_8_and_bibliography.tex
% SOURCE AUTHORITY:
% expose_VII_body.tex lines 1810--1977; scan indices 223--228; source folios 212--217.
% Direct scan comparison: physical PDF pages 224--229 at 600 dpi.
% Transparent normalization: typewritten letter-o subscripts are rendered as zero.
% The headless diagonal in (3.8.2) is reproduced as printed, while the prose fixes
% its mathematical role as the trivializing homomorphism.  The derived-category
% arrow in (3.8.10) follows the French authority, from P derived-tensor Q to
% L_\bullet tensor M_\bullet; the inherited abridgment reversed it.  Hard stop at
% the end of the Expose VII bibliography.

\subsection*{\SGARefTarget{sga7.sga7i.evii.n3-8}{3.8}. Biextensions of $(P,Q)$ trivialized over extensions of $P$ and $Q$}

First consider two exact sequences of abelian groups
\begin{equation}\tag{\SGARefTarget{sga7.sga7i.evii.n3-8-1}{3.8.1}}
 0\longrightarrow L_1\xrightarrow{\ i\ }L_0\xrightarrow{\ u\ }P
 \longrightarrow0,\qquad
 0\longrightarrow M_1\xrightarrow{\ j\ }M_0\xrightarrow{\ v\ }Q
 \longrightarrow0.
\end{equation}
We shall describe the category of biextensions $E$ of $(P,Q)$ by an
abelian group $G$ for which the inverse image $(u,v)^*(E)$ is trivial;
that is, those for which there exists a homomorphism $s$ making the
following triangle commutative:
% SGA-SRC-001262 — reconstructed by mathematical inference; scan inconclusive
\begin{equation}\tag{\SGARefTarget{sga7.sga7i.evii.n3-8-2}{3.8.2}}
\begin{tikzcd}[column sep=huge,row sep=huge]
 E\arrow[d]&\\
 P\times Q&L_0\times M_0\arrow[l,"u\times v"']\arrow[ul,"s"]
\end{tikzcd}
\end{equation}
and satisfying the additivity condition required in each of its two
arguments.  The choice of such an $s$ therefore identifies
$(u,v)^*(E)$ with the trivial biextension
$E_0=L_0\times M_0\times G$ of $(L_0,M_0)$ by $G$.  The left
exactness in $P,Q$ of the cofibred category of biextensions of
$(P,Q)$ by fixed $G$ (\SGARefLink{sga7.sga7i.evii.n3-7-6}{3.7.6}) then shows that to give the biextension
$E$ of $(P,Q)$, together with the trivialization $s$ of its inverse
image, is essentially equivalent to giving two partial
trivializations of the trivial biextension $E_0$ of $(L_0,M_0)$ by
$G$, over $(L_1,M_0)$ and $(L_0,M_1)$ respectively, which agree over
$(L_1,M_1)$.  Since the biextensions induced by $E_0$ over
$(L_1,M_0)$, $(L_0,M_1)$, and $(L_1,M_1)$ identify with the trivial
biextensions, a partial trivialization over $(L_1,M_0)$ (respectively
over $(L_0,M_1)$) is equivalent to giving a homomorphism (or pairing)
\begin{equation}\tag{\SGARefTarget{sga7.sga7i.evii.n3-8-3}{3.8.3}}
 f:L_1\otimes M_0\longrightarrow G
 \qquad
 \bigl(\text{respectively }g:L_0\otimes M_1\longrightarrow G\bigr).
\end{equation}
The compatibility in question simply means that the two
homomorphisms agree over $L_1\otimes M_1$, namely
\begin{equation}\tag{\SGARefTarget{sga7.sga7i.evii.n3-8-4}{3.8.4}}
 f(x_1\otimes y_1)=g(x_1\otimes y_1),
\end{equation}
for sections $x_1,y_1$ of $L_1,M_1$ over an object of
$\underline T$.  In terms of the homomorphism $s$ of (\SGARefLink{sga7.sga7i.evii.n3-8-2}{3.8.2}), the
forms $f$ and $g$ are characterized by
\begin{equation}\tag{\SGARefTarget{sga7.sga7i.evii.n3-8-5}{3.8.5}}
 s(x_1,y_0)=f(x_1\otimes y_0)e_{E/Q}(y_0),\qquad
 s(x_0,y_1)=g(x_0\otimes y_1)e_{E/P}(x_0),
\end{equation}
which again makes clear the necessity of condition (\SGARefLink{sga7.sga7i.evii.n3-8-4}{3.8.4}).

Once the trivialization $s$ has been chosen, the other
trivializations are the morphisms
$s':L_0\times M_0\longrightarrow E$ of the form
\begin{equation}\tag{\SGARefTarget{sga7.sga7i.evii.n3-8-6}{3.8.6}}
 s'(x_0,y_0)=s(x_0,y_0)+h(x_0\otimes y_0),
\end{equation}
where
\[
 h:L_0\otimes M_0\longrightarrow G
\]
is an arbitrary form.  It is then clear from (\SGARefLink{sga7.sga7i.evii.n3-8-5}{3.8.5}) that the forms
$f,g$ defined by $s$ are changed into
\begin{equation}\tag{\SGARefTarget{sga7.sga7i.evii.n3-8-7}{3.8.7}}
\begin{aligned}
 f'(x_1\otimes y_0)&=f(x_1\otimes y_0)+h(x_1\otimes y_0),\\
 g'(x_0\otimes y_1)&=g(x_0\otimes y_1)+h(x_0\otimes y_1).
\end{aligned}
\end{equation}
Thus we obtain:

\underline{Proposition} \SGARefTarget{sga7.sga7i.evii.n3-8-8}{3.8.8}.
\underline{The subgroup of $\Biext^1(P,Q;G)$ formed by the classes
of biextensions whose inverse image under $(u,v)$ is a trivial
biextension of $(L_0,M_0)$} (that is, the kernel of
$\Biext^1(P,Q;G)\longrightarrow\Biext^1(L_0,M_0;G)$)
\underline{is canonically isomorphic, by the preceding description,
to the quotient of the group of pairs of forms $(f,g)$ in (\SGARefLink{sga7.sga7i.evii.n3-8-3}{3.8.3})
satisfying compatibility condition (\SGARefLink{sga7.sga7i.evii.n3-8-4}{3.8.4}), by the pairs of forms
which arise by restriction from a form}
$h:L_0\otimes M_0\longrightarrow G$.

We now interpret this isomorphism in terms of homological algebra,
taking into account the canonical isomorphism
\[
 \Biext^1(P,Q;G)\simeq
 \Ext^1(P\mathbin{\overset{\mathbb L}{\otimes}}Q,G)
\]
of \SGARefLink{sga7.sga7i.evii.n3-6-5}{3.6.5}.  To this end, interpret the data (\SGARefLink{sga7.sga7i.evii.n3-8-1}{3.8.1}) as two complexes
resolving $P$ and $Q$, respectively:
\begin{equation}\tag{\SGARefTarget{sga7.sga7i.evii.n3-8-9}{3.8.9}}
 L_\bullet:\ 0\longrightarrow L_1\longrightarrow L_0
 \longrightarrow0,\qquad
 M_\bullet:\ 0\longrightarrow M_1\longrightarrow M_0
 \longrightarrow0,
\end{equation}
and consider the tensor-product complex $L_\bullet\otimes M_\bullet$
\[
 0\longrightarrow L_1\otimes M_1
 \xrightarrow{d_2}
 L_1\otimes M_0+L_0\otimes M_1
 \xrightarrow{d_1}
 L_0\otimes M_0\longrightarrow0.
\]
There is, moreover, a canonical homomorphism in the derived category
\[
 P\mathbin{\overset{\mathbb L}{\otimes}}Q
 \longrightarrow L_\bullet\otimes M_\bullet,
\]
and hence an equally canonical homomorphism
\begin{equation}\tag{\SGARefTarget{sga7.sga7i.evii.n3-8-10}{3.8.10}}
 H^1\!\left(\Hom^{\cdot}(L_\bullet\otimes M_\bullet,G)\right)
 \longrightarrow
 \Ext^1(P\mathbin{\overset{\mathbb L}{\otimes}}Q,G).
\end{equation}
The group of $1$-cocycles of
$\Hom^{\cdot}(L_\bullet\otimes M_\bullet,G)$ identifies with the
group of pairs $(f,g)$ as in (\SGARefLink{sga7.sga7i.evii.n3-8-3}{3.8.3}) which satisfy compatibility
condition (\SGARefLink{sga7.sga7i.evii.n3-8-4}{3.8.4}): to such a pair one associates the homomorphism
$F:L_1\otimes M_0+L_0\otimes M_1\longrightarrow G$ with components
$(f,g)$,
\begin{equation}\tag{\SGARefTarget{sga7.sga7i.evii.n3-8-11}{3.8.11}}
 F(x_1\otimes y_0+x_0\otimes y_1)
 =f(x_1\otimes y_0)+g(x_0\otimes y_1).
\end{equation}
Moreover, the coboundary of the $0$-cochain
$h:L_0\otimes M_0\longrightarrow G$ is precisely the pair of its
restrictions to $L_1\otimes M_0$ and $L_0\otimes M_1$.  Consequently,
the group of pairs $(f,g)$ modulo the pairs $(h,h)$ described in
\SGARefLink{sga7.sga7i.evii.n3-8-8}{3.8.8} is canonically isomorphic to the first term of (\SGARefLink{sga7.sga7i.evii.n3-8-10}{3.8.10}).
With this understood, one has the following compatibility.

\underline{Proposition} \SGARefTarget{sga7.sga7i.evii.n3-8-12}{3.8.12}.
\underline{If}
\[
 \xi\in\Ker\!\left(
 \Biext^1(P,Q;G)\longrightarrow\Biext^1(L_0,M_0;G)\right)
\]
\underline{is described by the pair $(f,g)$ as in \SGARefLink{sga7.sga7i.evii.n3-8-8}{3.8.8}, then the
image of $\xi$ in}
$\Ext^1(P\mathbin{\overset{\mathbb L}{\otimes}}Q,G)$
\underline{under the isomorphism (\SGARefLink{sga7.sga7i.evii.n3-6-5-1}{3.6.5.1}) is the image under
(\SGARefLink{sga7.sga7i.evii.n3-8-10}{3.8.10}) of the element of its first term defined by the cocycle
(\SGARefLink{sga7.sga7i.evii.n3-8-11}{3.8.11})}.

The verification of this compatibility is left to the reader.

\SGARefTarget{sga7.sga7i.evii.n3-8-13}{3.8.13}. We shall apply this especially in the case
\[
 L_0=\mathbb Z[P],\qquad M_0=\mathbb Z[Q],
\]
where $L_1$ and $M_1$ are respectively the kernels of the canonical
homomorphisms $\mathbb Z[P]\longrightarrow P$ and
$\mathbb Z[Q]\longrightarrow Q$.  Notice the canonical isomorphisms
\[
 L_0\mathbin{\overset{\mathbb L}{\otimes}}M_0
 \simeq L_0\otimes M_0
 \simeq\mathbb Z[P\times Q].
\]
Since $L_0$ and $M_0$ are flat, (\SGARefLink{sga7.sga7i.evii.n1-4-5}{1.4.5}) gives canonical isomorphisms
\[
 \Ext^i(L_0\mathbin{\overset{\mathbb L}{\otimes}}M_0,G)
 \simeq H^i(P\times Q,G_{P\times Q}).
\]
It follows from this and \SGARefLink{sga7.sga7i.evii.n3-6-5}{3.6.5}, or directly from the definitions and
\SGARefLink{sga7.sga7i.evii.n1-4}{1.4}, that the functor which assigns to each biextension of
$(L_0,M_0)$ by $G$ the $G_{P\times Q}$-torsor over $P\times Q$ that
it induces is an equivalence of categories.  (This remains valid
without using an additive structure on $P,Q$.)  Consequently, if
$E$ is a biextension of $(P,Q)$ by $G$, giving a trivialization of
its inverse image over $(L_0,M_0)$ is equivalent to giving a section
of $E$ considered only as an object over $P\times Q$, that is, a
trivialization of $E$ considered as a torsor under $G_{P\times Q}$.

Biextensions equipped with this additional structure are therefore
described by the pairs $(f,g)$ of (\SGARefLink{sga7.sga7i.evii.n3-8-3}{3.8.3}) satisfying compatibility
(\SGARefLink{sga7.sga7i.evii.n3-8-4}{3.8.4}), and their classification up to isomorphism is the one given
in \SGARefLink{sga7.sga7i.evii.n3-8-8}{3.8.8}.

With the notation of \SGARefLink{sga7.sga7i.evii.n3-5}{3.5}, one may also regard $L_1$ as the cokernel
of $L_2(P)\longrightarrow L_1(P)$, and similarly for $M_1$.  This
makes it possible to express the data $(f,g)$ in terms of a system
of homomorphisms
\[
\left\{
\begin{aligned}
 f'&:L_1(P)\otimes L_0(Q)
     =\mathbb Z(P\times P\times Q)\longrightarrow G,\\
 g'&:L_0(P)\otimes L_1(Q)
     =\mathbb Z(P\times Q\times Q)\longrightarrow G,
\end{aligned}
\right.
\]
that is, homomorphisms of sheaves of \underline{sets}
\[
 P\times P\times Q\longrightarrow G,\qquad
 P\times Q\times Q\longrightarrow G,
\]
satisfying five compatibility conditions.  These express that $f'$
and $g'$ factor through $f$ and $g$---which gives $2\times2$
equations, in view of the expression of $L_2(P)$ and $L_2(Q)$ as
sums of two terms---together with condition (\SGARefLink{sga7.sga7i.evii.n3-8-4}{3.8.4}).  One then sees
that $f'$ and $g'$ are precisely the partial multiplication data
defining a biextension, and that the five conditions in question are
exactly the five axioms for biextensions considered in \SGARefLink{sga7.sga7i.evii.n2-1}{2.1} (cf. the
proof of \SGARefLink{sga7.sga7i.evii.n3-6-4}{3.6.4}).

\subsection*{Bibliography to Expos\'e VII}

\begin{enumerate}[label={[\arabic*]}]
\item M. Demazure, J. Giraud, and M. Raynaud,
\emph{Schémas abéliens}, Séminaire de la Faculté des Sciences
d'Orsay 1967/68 (to appear).
\item J. Giraud, \emph{Cohomologie non abélienne}, thesis, Paris,
1967 (to appear).
\item A. Grothendieck, \emph{Sur quelques points d'Algèbre
homologique}, Tohoku Mathematical Journal 9 (1956), 119--221.
\item A. Grothendieck, \emph{Catégories cofibrées additives et
Complexe cotangent relatif}, Lecture Notes in Mathematics 79,
Springer, 1968.
\item J.-L. Verdier, \emph{Catégories dérivées des catégories
abéliennes} (to appear).
\item D. Mumford, \emph{Biextensions of formal groups}, Tata
Institute of Fundamental Research, 1968 (to appear).
\end{enumerate}

% END INCLUDED SGA7I_English_current_buildable_source_20260823/components/92_expose_VII_section_3_8_and_bibliography.tex

\clearpage
\phantomsection\addcontentsline{toc}{section}{Exposé VIII}
% BEGIN INCLUDED SGA7I_English_current_buildable_source_20260823/components/93_expose_VIII_opening_and_introduction.tex
% SOURCE AUTHORITY:
% expose_VIII_body.tex lines 5--52; scan indices 229--231; source folios 218--220.
% Direct scan comparison: physical PDF pages 230--232 at 1100 dpi.
% Hard stop before section 1, source line 55.

\begin{center}
\SGARefAliasTarget{sga7.sga7i.eviii.nroot}SGA 7

\vspace{1em}
\underline{Expos\'e VIII}

\vspace{1em}
\underline{COMPLEMENTS ON BIEXTENSIONS. GENERAL PROPERTIES}

\underline{OF BIEXTENSIONS OF GROUP SCHEMES}

\vspace{1em}
\underline{by A. Grothendieck}
\end{center}

\subsection*{Contents}

\begin{enumerate}[label=\arabic*.]
\setcounter{enumi}{-1}
\item Introduction
\item Various special cases over an arbitrary topos
\item Pairings defined by a biextension
\item Biextensions of group schemes $(P,Q)$ by $\underline G_m$: generalities
\item Extensions and biextensions of smooth connected group schemes over a field
\item Extensions and biextensions by constant truncated torsion-free group schemes
\item Extensions and biextensions by the N\'eron model of $\underline G_m$
\item Canonical prolongations of extensions and biextensions by $\underline G_m$
\end{enumerate}

0. \underline{Introduction}

\SGARefTarget{sga7.sga7i.eviii.n0-1}{0.1}. We continue here the general study of biextensions begun in the
preceding expos\'e.  First we give some complements in the case of an
arbitrary base topos, developing in particular, with the detail it
deserves, the formalism of the pairings associated with biextensions
(Section 2), whose importance in the theory of abelian schemes is well
known.  We then begin the study of biextensions of general group
schemes, the case of greatest interest for us being that of
biextensions by the multiplicative group $\underline G_m$.  We shall
nevertheless be obliged, as a technical intermediate step, also to
study biextensions by other types of group schemes (Sections 5 and 6).

The most interesting result for what follows is given in \SGARefLink{sga7.sga7i.eviii.n7-1}{7.1}, which
gives conditions under which, for $P$ and $Q$ smooth over a trait with
generic point $\eta$, a biextension of $(P_\eta,Q_\eta)$ by
$\underline G_{m\eta}$ can be prolonged to a biextension of $(P,Q)$ by
$\underline G_m$, determined up to a unique isomorphism.  This result
will be very convenient in the following expos\'e, devoted to the
theory of the N\'eron model of abelian varieties.  It is likely that
one could manage there at less expense, without the imposing
paperweight of Expos\'es VII and VIII (which will no longer be used in
the remainder of the Seminar); it nevertheless seems likely to us that
the developments given in these expos\'es are destined to be useful
elsewhere as well.

\SGARefTarget{sga7.sga7i.eviii.n0-2}{0.2}. When we work with group schemes over a scheme $S$, it will be
understood that we work with the fppf topos of $S$.  This is the topos
associated with the $\underline U$-site of schemes locally of finite
presentation over $S$ (and belonging to $\underline U$), endowed with
the fppf topology (SGA 3 IV 6.3).  The reader will readily verify that
most of the statements we shall give would also be valid (and
essentially equivalent to the fppf statements) for other topologies on
$S$, whose underlying site contains at least the schemes locally of
finite presentation over $S$, such as the Zariski, \'etale, and fpqc
topologies.  One reason is that, by descent theory, the category of
extensions of a group scheme $P$ locally of finite presentation over
$S$ by the group scheme $\underline G_m$ does not depend, up to
equivalence, on the chosen topology; the same statement holds for the
notion of biextension.

The technical reason for preferring the fppf topology (or a finer
topology, such as fpqc) to coarser topologies such as the Zariski and
\'etale topologies is that, for every integer $n>0$, the homomorphism
$n\cdot\id$ from $\underline G_m$ to itself is an \underline{fppf
epimorphism} (because it is faithfully flat and of finite
presentation), but in general is not a Zariski or an \'etale
epimorphism.  Notice that it is an \'etale epimorphism if and only if
$n$ is prime to the residue characteristics of $S$, and a Zariski
epimorphism if and only if $S$ is empty.

\SGARefTarget{sga7.sga7i.eviii.n0-3}{0.3}. Finally, most of the statements concerning extensions or
biextensions by $\underline G_m$ would remain valid if that group were
replaced by a \underline{torus} $G$.  Since such a $G$ is locally
isomorphic (for the fppf, and even for the \'etale topology) to a group
$\underline G_m^r$, this is an essentially trivial generalization.

% END INCLUDED SGA7I_English_current_buildable_source_20260823/components/93_expose_VIII_opening_and_introduction.tex

% BEGIN INCLUDED SGA7I_English_current_buildable_source_20260823/components/94_expose_VIII_section_1.tex
% SOURCE AUTHORITY:
% expose_VIII_body.tex lines 55--276; scan indices 231--237; source folios 220--226.
% Direct scan comparison: physical PDF pages 232--238 at 1100 dpi.
% Transparent source dispositions:
% - the sentence before (1.1.3) is stated in standard truncation language rather
%   than repeating the printed impossible claim that all Tor_i vanish for i >= 1;
% - the fourth term of (1.1.4) uses Hom(Q,G), as required by the displayed
%   spectral sequence, rather than the printed Hom(Q,P);
% - broken parentheses in (1.3.5), the transition exponent in part (c), and the
%   references to the final member of (1.3.5) are restored from the formulas;
% - in 1.6 the parenthetical Biext argument and the terminal Ext coefficient are
%   corrected to Q and G respectively.
% Hard stop before section 2, source line 279.

\subsection*{\SGARefTarget{sga7.sga7i.eviii.n1}{1}. Various special cases over an arbitrary topos}

\SGARefTarget{sga7.sga7i.eviii.n1-1}{1.1}. In this section and the next, as in Expos\'e VII, $T$ denotes a
topos, and $P,Q,G$ denote three abelian groups of $T$.  Thus there is a
canonical isomorphism (VII \SGARefLink{sga7.sga7i.evii.n3-6-5}{3.6.5})
\begin{equation}\tag{\SGARefTarget{sga7.sga7i.eviii.n1-1-1}{1.1.1}}
 \Biext^1(P,Q;G)
 \simeq
 \Ext^1\!\left(P\mathbin{\overset{\mathbb L}{\otimes}}Q,G\right).
\end{equation}
Recall the canonical isomorphisms (called the Cartan isomorphisms)
\begin{equation}\tag{\SGARefTarget{sga7.sga7i.eviii.n1-1-2}{1.1.2}}
\begin{split}
 \Ext^1\!\left(P\mathbin{\overset{\mathbb L}{\otimes}}Q,G\right)
 &\simeq \Ext^1\!\left(P,R\uHom(Q,G)\right)\\
 &\simeq \Ext^1\!\left(Q,R\uHom(P,G)\right).
\end{split}
\end{equation}
All these isomorphisms are functorial in the three arguments $P,Q,G$.

Taking the standard truncation of the derived tensor product, whose
homology objects are the $\Tor_i(P,Q)$, one obtains from (\SGARefLink{sga7.sga7i.eviii.n1-1-1}{1.1.1}) the
following infinite exact sequence:
\begin{equation}\tag{\SGARefTarget{sga7.sga7i.eviii.n1-1-3}{1.1.3}}
\begin{aligned}
0\longrightarrow{}&
 \Ext^1(P\otimes Q,G)
 \longrightarrow \Biext^1(P,Q;G)
 \longrightarrow \Hom(\Tor_1(P,Q),G)\\
 \longrightarrow{}&
 \Ext^2(P\otimes Q,G)
 \longrightarrow
 \Ext^2\!\left(P\mathbin{\overset{\mathbb L}{\otimes}}Q,G\right)
 \longrightarrow\cdots .
\end{aligned}
\end{equation}
On the other hand, since the cohomology objects of $R\uHom(Q,G)$ are
the $\uExt^i(Q,G)$, the first equality in (\SGARefLink{sga7.sga7i.eviii.n1-1-2}{1.1.2}) gives rise to the
five-term exact sequence
\begin{equation}\tag{\SGARefTarget{sga7.sga7i.eviii.n1-1-4}{1.1.4}}
\begin{aligned}
0\longrightarrow{}&
 \Ext^1(P,\uHom(Q,G))
 \longrightarrow \Biext^1(P,Q;G)
 \longrightarrow \Hom(P,\uExt^1(Q,G))\\
 \longrightarrow{}&
 \Ext^2(P,\uHom(Q,G))
 \longrightarrow \Ext^2(P,R\uHom(Q,G)),
\end{aligned}
\end{equation}
associated with the usual spectral sequence
\[
 \Ext^*(P,R\uHom(Q,G))
 \Longleftarrow
 \Ext^p(P,\uExt^q(Q,G)).
\]
We leave it to the reader, should the need arise, to make the first
arrow of (\SGARefLink{sga7.sga7i.eviii.n1-1-3}{1.1.3}), respectively of (\SGARefLink{sga7.sga7i.eviii.n1-1-4}{1.1.4}), explicit by means of the
\underline{natural fully faithful functor}
\begin{equation}\tag{\SGARefTarget{sga7.sga7i.eviii.n1-1-5}{1.1.5}}
 \operatorname{EXT}(P\otimes Q,G)
 \longrightarrow
 \operatorname{BIEXT}(P,Q;G),
\end{equation}
respectively
\begin{equation}\tag{\SGARefTarget{sga7.sga7i.eviii.n1-1-6}{1.1.6}}
 \operatorname{EXT}(P,\uHom(Q,G))
 \longrightarrow
 \operatorname{BIEXT}(P,Q;G),
\end{equation}
where $\operatorname{EXT}$ and $\operatorname{BIEXT}$ denote the
evident categories of extensions and biextensions.

Of course, the isomorphism
\[
 \Biext^1(P,Q;G)\simeq\Ext^1(Q,R\uHom(P,G))
\]
gives rise to an exact sequence analogous to (\SGARefLink{sga7.sga7i.eviii.n1-1-3}{1.1.3}), with the roles
of $P$ and $Q$ interchanged.  Its first arrow can be made explicit by
means of the fully faithful functor
\begin{equation}\tag{\SGARefTarget{sga7.sga7i.eviii.n1-1-7}{1.1.7}}
 \operatorname{EXT}(Q,\uHom(P,G))
 \longrightarrow
 \operatorname{BIEXT}(P,Q;G)
\end{equation}
analogous to (\SGARefLink{sga7.sga7i.eviii.n1-1-6}{1.1.6}).

\SGARefTarget{sga7.sga7i.eviii.n1-2}{1.2}. Suppose that
\begin{equation}\tag{\SGARefTarget{sga7.sga7i.eviii.n1-2-1}{1.2.1}}
 \Tor_1(P,Q)=0.
\end{equation}
Then (\SGARefLink{sga7.sga7i.eviii.n1-1-3}{1.1.3}) supplies a canonical isomorphism
\begin{equation}\tag{\SGARefTarget{sga7.sga7i.eviii.n1-2-2}{1.2.2}}
 \Biext^1(P,Q;G)\simeq\Ext^1(P\otimes Q,G).
\end{equation}
This holds, for example, if $P$ or $Q$ is a \underline{flat}
$\mathbb Z$-module.  More precisely, under condition (\SGARefLink{sga7.sga7i.eviii.n1-2-1}{1.2.1}), the
functor (\SGARefLink{sga7.sga7i.eviii.n1-1-5}{1.1.5}) is an equivalence of categories.

\SGARefTarget{sga7.sga7i.eviii.n1-3}{1.3}. Suppose that
\begin{equation}\tag{\SGARefTarget{sga7.sga7i.eviii.n1-3-1}{1.3.1}}
 P\otimes Q=0.
\end{equation}
Then (\SGARefLink{sga7.sga7i.eviii.n1-1-3}{1.1.3}) supplies a canonical isomorphism
\begin{equation}\tag{\SGARefTarget{sga7.sga7i.eviii.n1-3-2}{1.3.2}}
 \Biext^1(P,Q;G)\simeq\Hom(\Tor_1(P,Q),G).
\end{equation}
The case of greatest interest to us in which (\SGARefLink{sga7.sga7i.eviii.n1-3-1}{1.3.1}) holds is that in
which there is an integer $n>0$ such that
\begin{equation}\tag{\SGARefTarget{sga7.sga7i.eviii.n1-3-3}{1.3.3}}
 nP=P,\qquad
 Q=\varinjlim_i {}_{n^i}Q;
\end{equation}
that is, $P$ is divisible by $n$, and $Q$ is the limit of subgroups
annihilated by powers of $n$.  Put, as usual,
\[
 T_n(P)=({}_{n^i}P)_{i\geq0}\quad\text{(a projective system)},\qquad
 {}_{n^\infty}P=\varinjlim_i{}_{n^i}P.
\]
From (\SGARefLink{sga7.sga7i.eviii.n1-3-3}{1.3.3}) one obtains the following formulas, more precise than
(\SGARefLink{sga7.sga7i.eviii.n1-3-1}{1.3.1}):
\begin{equation}\tag{\SGARefTarget{sga7.sga7i.eviii.n1-3-4}{1.3.4}}
\begin{split}
 P\otimes Q&=0,\\
 \Tor_1(P,Q)
 &\simeq\varinjlim_i
   \bigl({}_{n^i}P\otimes{}_{n^i}Q\bigr)
 \simeq T_n(P)\otimes Q\\
 &\simeq{}_{n^\infty}P\otimes T_n(Q).
\end{split}
\end{equation}
Thus the isomorphism (\SGARefLink{sga7.sga7i.eviii.n1-3-2}{1.3.2}) here assumes the more concrete form
\begin{equation}\tag{\SGARefTarget{sga7.sga7i.eviii.n1-3-5}{1.3.5}}
\begin{split}
 \Biext^1(P,Q;G)
 &\simeq\Hom(T_n(P),D(Q))\\
 &\simeq\Hom(Q,D(T_n(P)))\\
 &\simeq\Hom(T_n(P)\otimes T_n(Q),T_n(G)).
\end{split}
\end{equation}
Let us specify the meaning of these formulas.

\begin{enumerate}[label=\alph*)]
\item In the inductive system occurring in (\SGARefLink{sga7.sga7i.eviii.n1-3-4}{1.3.4}), the transition
homomorphism
\[
 {}_{n^i}P\otimes{}_{n^i}Q
 \longrightarrow
 {}_{n^j}P\otimes{}_{n^j}Q
 \qquad(i\leq j)
\]
is described by identifying the first member with
$\,{}_{n^j}P\otimes{}_{n^i}Q$ by means of
$n^{j-i}\otimes\id$, and then using the homomorphism
\[
 \id\otimes\operatorname{inj}:
 {}_{n^j}P\otimes{}_{n^i}Q
 \longrightarrow
 {}_{n^j}P\otimes{}_{n^j}Q,
\]
where
$\operatorname{inj}:{}_{n^i}Q\longrightarrow{}_{n^j}Q$
denotes the inclusion.

\item By definition, put
\[
 T_n(P)\otimes Q
 \overset{\mathrm{def}}{=}
 \varinjlim_i\bigl(T_n(P)\otimes{}_{n^i}Q\bigr),
\]
where
\[
 T_n(P)\otimes M
 \overset{\mathrm{def}}{=}
 {}_{n^i}P\otimes M
 \qquad\text{if }n^iM=0.
\]
The meaning of the third term written in (\SGARefLink{sga7.sga7i.eviii.n1-3-4}{1.3.4}) is specified in the
same way when $nQ=Q$ is also assumed, so that $T_n(Q)$, like
$T_n(P)$, is a strict projective system.  The last term of (\SGARefLink{sga7.sga7i.eviii.n1-3-5}{1.3.5}) is
a Hom between projective systems of groups.

\item The letter $D$ denotes $\uHom(-,G)$.  The object $D(Q)$ is
defined as the projective system of groups
\[
 D(Q)=\bigl(D({}_{n^i}Q)\bigr)_{i\geq0}
 =\bigl(\uHom({}_{n^i}Q,G)\bigr)_{i\geq0},
\]
the second member of (\SGARefLink{sga7.sga7i.eviii.n1-3-5}{1.3.5}) being the group of homomorphisms of
projective systems of groups.  The following definition gives a
simpler interpretation of the next member:
\[
 D(T_n(P))=\varinjlim_iD({}_{n^i}P),
\]
where the transition morphisms of this inductive system are obtained
by transposition from the transition morphisms
\[
 n^{j-i}:{}_{n^j}P\longrightarrow{}_{n^i}P
\]
of $T_n(P)$.
\end{enumerate}

We leave it to the reader to make explicit the calculation of
$\Tor_1(P,Q)$ indicated in (\SGARefLink{sga7.sga7i.eviii.n1-3-4}{1.3.4}), and to derive (\SGARefLink{sga7.sga7i.eviii.n1-3-5}{1.3.5}) from it.
We note only that one uses the relations, deduced from the second
hypothesis in (\SGARefLink{sga7.sga7i.eviii.n1-3-3}{1.3.3}),
\[
 \Tor_p(P,Q)\simeq
 \varinjlim_i\Tor_p(P,{}_{n^i}Q).
\]
For $p=0$ this immediately gives $P\otimes Q=0$, and it reduces the
calculation of $\Tor_1(P,Q)$ to that of the
$\Tor_1(P,{}_{n^i}Q)$ and their transition morphisms.  Thus, when
$Q$ is annihilated by an integer $m=n^i$ such that $mP=P$, it remains
to make explicit an isomorphism
\begin{equation}\tag{\SGARefTarget{sga7.sga7i.eviii.n1-3-6}{1.3.6}}
 \Tor_1(P,Q)
 \xleftarrow{\ \sim\ }
 {}_mP\otimes{}_mQ
 \qquad(mP=P,\ mQ=0).
\end{equation}
There are several ways to do this.  One consists in using the exact
Tor sequence associated with
\[
 0\longrightarrow{}_mP\longrightarrow P
 \xrightarrow{\ m\ }P\longrightarrow0.
\]
One may also use the homomorphism (\SGARefLink{sga7.sga7i.eviii.n2-1-1}{2.1.1}) below; the fact that the
transition morphisms are those specified in part (a) is then precisely
\SGARefLink{sga7.sga7i.eviii.n2-1-11}{2.1.11}.  In any case, the choice of (\SGARefLink{sga7.sga7i.eviii.n1-3-6}{1.3.6}) amounts to the choice of a
sign; see \SGARefLink{sga7.sga7i.eviii.n2-1-10}{2.1.10}.

\SGARefTarget{sga7.sga7i.eviii.n1-4}{1.4}. Suppose that
\begin{equation}\tag{\SGARefTarget{sga7.sga7i.eviii.n1-4-1}{1.4.1}}
 \uHom(Q,G)=0.
\end{equation}
Then (\SGARefLink{sga7.sga7i.eviii.n1-1-4}{1.1.4}) supplies a canonical isomorphism
\begin{equation}\tag{\SGARefTarget{sga7.sga7i.eviii.n1-4-2}{1.4.2}}
 \Biext^1(P,Q;G)
 \xrightarrow{\ \sim\ }
 \Hom(P,\uExt^1(Q,G)).
\end{equation}
This case occurs in particular when $Q$ is an abelian scheme and
$G=\underline G_m$; see VII \SGARefLink{sga7.sga7i.evii.n1-3-8}{1.3.8}.  Hence the functor in $P$,
\[
 P\longmapsto\Biext^1(P,Q;G),
\]
is \underline{representable} by $Q^*=\uExt^1(Q,G)$.  In particular,
there is a canonical biextension $E^Q$ of $(Q^*,Q)$ by $G$, determined
up to a unique isomorphism, since its automorphism group is
\[
 \Hom(Q^*,\uHom(Q,G))=0.
\]
This biextension, or its symmetric biextension $\,{}^sE^Q$ of
$(Q,Q^*)$ by $G$, may be called the \underline{Weil biextension} of
$(Q^*,Q)$, respectively $(Q,Q^*)$, by $G$.

\SGARefTarget{sga7.sga7i.eviii.n1-5}{1.5}. Suppose that
\begin{equation}\tag{\SGARefTarget{sga7.sga7i.eviii.n1-5-1}{1.5.1}}
 \Hom(P,\uExt^1(Q,G))=0.
\end{equation}
This holds in particular if $\uExt^1(Q,G)=0$; for example, if $Q$ is
a finite locally free group scheme or a group of multiplicative type
(SGA 3 IX), and $G=\underline G_m$; see \SGARefLink{sga7.sga7i.eviii.n3-4}{3.4} below.  Then (\SGARefLink{sga7.sga7i.eviii.n1-1-4}{1.1.4})
supplies a canonical isomorphism
\begin{equation}\tag{\SGARefTarget{sga7.sga7i.eviii.n1-5-2}{1.5.2}}
 \Biext^1(P,Q;G)
 \xleftarrow{\ \sim\ }
 \Ext^1(P,\uHom(Q,G)),
\end{equation}
where $\uHom(Q,G)$ plays the role of a "dual group" of $Q$ (relative
to $G$).  More precisely, it follows from (\SGARefLink{sga7.sga7i.eviii.n1-5-2}{1.5.2}) that the fully
faithful functor (\SGARefLink{sga7.sga7i.eviii.n1-1-6}{1.1.6}) is an \underline{equivalence of categories}.

\SGARefTarget{sga7.sga7i.eviii.n1-6}{1.6}. Suppose that the two relations
\begin{equation}\tag{\SGARefTarget{sga7.sga7i.eviii.n1-6-1}{1.6.1}}
 \uHom(P,G)=0,\qquad
 \Hom(P,\uExt^1(Q,G))=0
\end{equation}
hold simultaneously.  By \SGARefLink{sga7.sga7i.eviii.n1-4}{1.4} and \SGARefLink{sga7.sga7i.eviii.n1-5}{1.5}, respectively, these two
relations give isomorphisms
\[
 \Biext^1(P,Q;G)
 \xrightarrow{\ \sim\ }
 \Hom(Q,\uExt^1(P,G)),
\]
\[
 \Biext^1(P,Q;G)
 \xleftarrow{\ \sim\ }
 \Ext^1(P,\uHom(Q,G)).
\]
Their composite is the canonical isomorphism
\begin{equation}\tag{\SGARefTarget{sga7.sga7i.eviii.n1-6-2}{1.6.2}}
 \Ext^1(P,\uHom(Q,G))
 \xrightarrow{\ \sim\ }
 \Hom(Q,\uExt^1(P,G))
 \quad\bigl(\simeq\Biext^1(P,Q;G)\bigr),
\end{equation}
which plays the role of a duality formula.

Notice also that the first relation in (\SGARefLink{sga7.sga7i.eviii.n1-6-1}{1.6.1}) plainly implies
$\Hom(P\otimes Q,G)=0$.  Hence the category
$\operatorname{EXT}(P,\uHom(Q,G))$ of commutative extensions of $P$
by $\uHom(Q,G)$ is rigid (every automorphism of an object is the
identity), as is $\operatorname{BIEXT}(P,Q;G)$.  The structure of
these categories is therefore known once the group of isomorphism
classes of objects is known; by (\SGARefLink{sga7.sga7i.eviii.n1-6-2}{1.6.2}), this group is canonically
isomorphic to
\[
 \Hom(Q,\uExt^1(P,G)).
\]

The map (\SGARefLink{sga7.sga7i.eviii.n1-6-2}{1.6.2}) is in fact defined in all cases, without assuming
(\SGARefLink{sga7.sga7i.eviii.n1-6-1}{1.6.1}), and can also be described directly, without passing through
$\Biext^1(P,Q;G)$, in terms of the natural pairing
\[
 Q\times Q'\longrightarrow G,
\qquad Q'=\uHom(Q,G).
\]
Interpreting it as a homomorphism
$Q\longrightarrow\uHom(Q',G)$ gives a composite
\begin{equation}\tag{\SGARefTarget{sga7.sga7i.eviii.n1-6-3}{1.6.3}}
\begin{split}
 \Ext^1(P,Q')\times Q
 &\longrightarrow
 \Ext^1(P,Q')\times\Hom(Q',G)\\
 &\longrightarrow
 \Ext^1(P,G),
\end{split}
\end{equation}
which recovers (\SGARefLink{sga7.sga7i.eviii.n1-6-2}{1.6.2}).

For a useful case in which the considerations of \SGARefLink{sga7.sga7i.eviii.n1-6}{1.6} apply, see \SGARefLink{sga7.sga7i.eviii.n3-7}{3.7}
below.

% END INCLUDED SGA7I_English_current_buildable_source_20260823/components/94_expose_VIII_section_1.tex

% BEGIN INCLUDED SGA7I_English_current_buildable_source_20260823/components/95_expose_VIII_section_2_1_opening_through_2_1_8.tex
% SOURCE AUTHORITY:
% expose_VIII_body.tex lines 279--382; scan indices 237--241; source folios 226--230.
% Direct scan comparison: physical PDF pages 238--242 at 1100 dpi.
% Transparent source dispositions:
% - typewritten letter-o homological indices are normalized to zero;
% - after (2.1.5), the second map is induced by the inclusion {}_nP -> P,
%   not the authority's impossible {}_nP -> Q;
% - the malformed two-argument Tor display in 2.1.8 is set with both arguments.
% Hard stop before 2.1.9, source line 384.

\subsection*{\SGARefTarget{sga7.sga7i.eviii.n2}{2}. Pairings defined by a biextension}

\SGARefTarget{sga7.sga7i.eviii.n2-1}{2.1}. Let $P$ and $Q$ be two abelian groups of $T$, and let $n>0$ be
an integer.  Recall that $\,{}_nP$ denotes the kernel of multiplication
by $n\cdot\id_P$.  We shall define a canonical homomorphism
\begin{equation}\tag{\SGARefTarget{sga7.sga7i.eviii.n2-1-1}{2.1.1}}
 \alpha=\alpha^n_{P,Q}:
 {}_nP\otimes{}_nQ\longrightarrow\uTor_1(P,Q),
\end{equation}
which will be the composite
\begin{equation}\tag{\SGARefTarget{sga7.sga7i.eviii.n2-1-2}{2.1.2}}
 {}_nP\otimes{}_nQ
 \longrightarrow \uTor_1({}_nP,{}_nQ)
 \longrightarrow \uTor_1(P,Q),
\end{equation}
where the first arrow is $\alpha^n_{{}_nP,{}_nQ}$ and the second is
induced by the inclusions $\,{}_nP\longrightarrow P$ and
$\,{}_nQ\longrightarrow Q$.  We are thus reduced to defining (\SGARefLink{sga7.sga7i.eviii.n2-1-1}{2.1.1})
when $nP=0$ and $nQ=0$, that is, when $P$ and $Q$ are
$\mathbb Z/n\mathbb Z$-modules.

Suppose only that $nQ=0$.  Taking a flat resolution $L_\bullet$ of
$P$ (as a $\mathbb Z$-module), one obtains canonical isomorphisms
\[
 P\mathbin{\overset{\mathbb L}{\otimes}_{\mathbb Z}}Q
 \simeq L_\bullet\otimes_{\mathbb Z}Q
 \simeq
 (L_\bullet\otimes_{\mathbb Z}\mathbb Z/n\mathbb Z)
 \otimes_{\mathbb Z/n\mathbb Z}Q,
\]
which may be written as an isomorphism in the derived category
\begin{equation}\tag{\SGARefTarget{sga7.sga7i.eviii.n2-1-3}{2.1.3}}
 P\mathbin{\overset{\mathbb L}{\otimes}_{\mathbb Z}}Q
 \simeq
 \left(P\mathbin{\overset{\mathbb L}{\otimes}_{\mathbb Z}}
       \mathbb Z/n\mathbb Z\right)
 \mathbin{\overset{\mathbb L}{\otimes}_{\mathbb Z/n\mathbb Z}}Q.
\end{equation}
This gives rise to the homological K\"unneth-type spectral sequence
(cf.\ SGA 1 IV 5.2)
\begin{equation}\tag{\SGARefTarget{sga7.sga7i.eviii.n2-1-4}{2.1.4}}
 \uTor^{\mathbb Z}_*(P,Q)
 \Longleftarrow
 E^2_{p,q}
 =
 \uTor^{\mathbb Z/n\mathbb Z}_p
 \bigl(\uTor^{\mathbb Z}_q(P,\mathbb Z/n\mathbb Z),Q\bigr).
\end{equation}
Recall also the canonical isomorphisms
\[
 \uTor^{\mathbb Z}_0(P,\mathbb Z/n\mathbb Z)
 =P_n\mathrel{\mathop:}=P/nP,
 \qquad
 \uTor^{\mathbb Z}_1(P,\mathbb Z/n\mathbb Z)\simeq{}_nP,
\]
and
\[
 \uTor^{\mathbb Z}_i(P,\mathbb Z/n\mathbb Z)=0
 \qquad\text{for }i\geq2.
\]
They follow by computing these objects from the canonical flat
resolution of $(\mathbb Z/n\mathbb Z)_T$:
\[
\begin{tikzcd}[column sep=large,row sep=large]
 0\arrow[r]&
 \mathbb Z_T\arrow[r,"n"]\arrow[d,equal]&
 \mathbb Z_T\arrow[r]\arrow[d,equal]&
 0\\
 &L_1&L_0&
\end{tikzcd}
\]
whose augmentation
$L_0\longrightarrow(\mathbb Z/n\mathbb Z)_T$ sends $1$ to $1$.
Consequently, the spectral sequence (\SGARefLink{sga7.sga7i.eviii.n2-1-4}{2.1.4}) reduces essentially to
the exact sequence
\begin{equation}\tag{\SGARefTarget{sga7.sga7i.eviii.n2-1-5}{2.1.5}}
\begin{split}
0\longrightarrow{}&
 \uTor^{\mathbb Z/n\mathbb Z}_2(P_n,Q)
 \longrightarrow{}_nP\otimes Q
 \xrightarrow{\alpha}\uTor^{\mathbb Z}_1(P,Q)\\
\longrightarrow{}&
 \uTor^{\mathbb Z/n\mathbb Z}_1(P_n,Q)
 \longrightarrow0,
\end{split}
\end{equation}
where $P_n=P\otimes_{\mathbb Z}\mathbb Z/n\mathbb Z=P/nP$.
The desired homomorphism $\alpha$ is the edge homomorphism
$E^2_{0,1}\longrightarrow\uTor^{\mathbb Z}_1$ in this exact sequence.
Since the sequence is plainly functorial in $P$ and $Q$, functoriality
in $P$ shows that the homomorphism in question is indeed the composite
\[
 {}_nP\otimes Q
 \longrightarrow\uTor_1({}_nP,Q)
 \longrightarrow\uTor_1(P,Q),
\]
where the first homomorphism is the $\alpha$ just defined, with
$(P,Q)$ replaced by $({}_nP,Q)$, and the second is induced by the
inclusion $\,{}_nP\longrightarrow P$.

\SGARefTarget{sga7.sga7i.eviii.n2-1-6}{2.1.6}. This therefore defines, for arbitrary $P$ and $Q$, a
homomorphism (\SGARefLink{sga7.sga7i.eviii.n2-1-1}{2.1.1}) satisfying the factorization condition (\SGARefLink{sga7.sga7i.eviii.n2-1-2}{2.1.2}).
Moreover, when $nQ=0$, the exact sequence (\SGARefLink{sga7.sga7i.eviii.n2-1-5}{2.1.5}) shows that $\alpha$
is an isomorphism whenever $P_n$ is flat over
$\mathbb Z/n\mathbb Z$.  This holds in particular if $P_n=0$, that
is, if $nP=P$ and $P$ is $n$-divisible, or if $P$ itself is a flat
$\mathbb Z/n\mathbb Z$-module.  It also holds if $Q$ is a flat
$\mathbb Z/n\mathbb Z$-module.

\SGARefTarget{sga7.sga7i.eviii.n2-1-7}{2.1.7}. Let us make the calculation of the homomorphism $\alpha$ in
(\SGARefLink{sga7.sga7i.eviii.n2-1-5}{2.1.5}) explicit.  To determine it, it suffices, for every section
$q$ of $Q$ (over an object $S$ of $T$; take the final object for
simplicity), to determine the corresponding homomorphism
\[
 {}_nP\longrightarrow\uTor^{\mathbb Z}_1(P,Q).
\]
The section $q$ may be regarded as the image of the section $1$ of
$(\mathbb Z/n\mathbb Z)_T$ under the uniquely determined
homomorphism
\[
 u_q:(\mathbb Z/n\mathbb Z)_T\longrightarrow Q.
\]
Functoriality of $\alpha$ in $Q$ therefore reduces us to the case
$Q=(\mathbb Z/n\mathbb Z)_T$, with $q$ its unit section.  In this
case, identify
\[
 \uTor^{\mathbb Z}_1(P,Q)
 =
 \uTor^{\mathbb Z}_1(P,\mathbb Z/n\mathbb Z)
\]
with $\,{}_nP$ as explained after (\SGARefLink{sga7.sga7i.eviii.n2-1-4}{2.1.4}).  The homomorphism in question
is then the identity.  Equivalently, the edge homomorphism
\[
 E^2_{0,q}
 =
 \uTor^{\mathbb Z}_q(P,\mathbb Z/n\mathbb Z)
 \otimes_{\mathbb Z/n\mathbb Z}\mathbb Z/n\mathbb Z
 \longrightarrow
 \uTor^{\mathbb Z}_q(P,\mathbb Z/n\mathbb Z)
\]
is the identity.  This follows immediately from the definitions.
The analogous assertion remains valid when $\mathbb Z$ and its
quotient $\mathbb Z/n\mathbb Z$ are replaced by any ring $A$ on $T$
and any quotient ring $A/I$.

\SGARefTarget{sga7.sga7i.eviii.n2-1-8}{2.1.8}. The homomorphism $\alpha$ of (\SGARefLink{sga7.sga7i.eviii.n2-1-1}{2.1.1}) may also be determined by
describing, for every pair of sections $p$ of $\,{}_nP$ and $q$ of
$\,{}_nQ$ (over an object $S$ of $T$, again taken to be final to
simplify the notation), the section of
$\uTor^{\mathbb Z}_1(P,Q)$ determined by $(p,q)$.  By functoriality in
$P$ and $Q$, as above, we are reduced to the case
\[
 P=Q=(\mathbb Z/n\mathbb Z)_T,
\]
with $p$ and $q$ the unit section.  Thus we need only determine
\[
 \alpha:
 (\mathbb Z/n\mathbb Z)_T\otimes(\mathbb Z/n\mathbb Z)_T
 \longrightarrow
 \uTor^{\mathbb Z}_1
 \bigl((\mathbb Z/n\mathbb Z)_T,(\mathbb Z/n\mathbb Z)_T\bigr).
\]
The description in \SGARefLink{sga7.sga7i.eviii.n2-1-7}{2.1.7} shows that this is the isomorphism
\begin{equation}\tag{\SGARefTarget{sga7.sga7i.eviii.n2-1-8-1}{2.1.8.1}}
 (\mathbb Z/n\mathbb Z)_T
 \xrightarrow{\sim}
 \uTor^{\mathbb Z}_1
 \bigl((\mathbb Z/n\mathbb Z)_T,(\mathbb Z/n\mathbb Z)_T\bigr)
\end{equation}
deduced from the canonical flat resolution
\begin{equation}\tag{\SGARefTarget{sga7.sga7i.eviii.n2-1-8-2}{2.1.8.2}}
 0\longrightarrow\mathbb Z_T
 \xrightarrow{n}\mathbb Z_T
 \longrightarrow0
\end{equation}
of the \underline{first} factor
$P=(\mathbb Z/n\mathbb Z)_T$ occurring in $\uTor_1$.

% END INCLUDED SGA7I_English_current_buildable_source_20260823/components/95_expose_VIII_section_2_1_opening_through_2_1_8.tex

% BEGIN INCLUDED SGA7I_English_current_buildable_source_20260823/components/96_expose_VIII_section_2_1_9_through_2_1_12.tex
% SOURCE AUTHORITY:
% expose_VIII_body.tex lines 384--547; scan indices 241--244; source folios 230--233.
% Direct scan comparison: physical PDF pages 242--245 at 1100 dpi.
% Transparent source dispositions:
% - typewritten letter-o complex indices are normalized to zero;
% - broken internal references 2.11.1/2.11.2 are restored as 2.1.11.1/2.1.11.2;
% - the upper vertex in the proof of 2.1.11 is set as P/nP, as required by
%   the identification in the prose and the four maps.
% Hard stop before 2.1.13, source line 552.

\SGARefTarget{sga7.sga7i.eviii.n2-1-9}{2.1.9}. The definition given in (\SGARefLink{sga7.sga7i.eviii.n2-1-1}{2.1.1}) is not symmetric.  By
interchanging the roles of $P$ and $Q$ in the preceding construction,
one may define an analogous homomorphism
\begin{equation}\tag{\SGARefTarget{sga7.sga7i.eviii.n2-1-9-1}{2.1.9.1}}
 \beta^n_{P,Q}:{}_nP\otimes{}_nQ\longrightarrow\uTor_1(P,Q),
\end{equation}
first reducing to the case $nP=0$ and then using a flat resolution
\underline{of} $Q$.  Equivalently, $\beta$ may be defined in terms of
$\alpha$ by requiring the following diagram to commute:
\begin{equation}\tag{\SGARefTarget{sga7.sga7i.eviii.n2-1-9-2}{2.1.9.2}}
\begin{tikzcd}[column sep=huge,row sep=huge]
 {}_nP\otimes{}_nQ
   \arrow[r,"\sim"]
   \arrow[d,"{\beta^n_{P,Q}}"']&
 {}_nQ\otimes{}_nP
   \arrow[d,"{\alpha^n_{Q,P}}"]\\
 \uTor_1(P,Q)\arrow[r,"\sim"]&
 \uTor_1(Q,P).
\end{tikzcd}
\end{equation}
Here the horizontal arrows are the symmetry isomorphisms.  With this
convention, the relation between $\alpha$ and $\beta$ is
\begin{equation}\tag{\SGARefTarget{sga7.sga7i.eviii.n2-1-9-3}{2.1.9.3}}
 \beta^n_{P,Q}=-\alpha^n_{P,Q}.
\end{equation}
In other words:

\medskip
\noindent\underline{Lemma \SGARefTarget{sga7.sga7i.eviii.n2-1-10}{2.1.10}.}
\underline{Let $P$ and $Q$ be two commutative groups and let $n>0$ be
an integer.  Then the following diagram is anticommutative, where
$\alpha$ is defined in (\SGARefLink{sga7.sga7i.eviii.n2-1-1}{2.1.1}):}
\begin{equation}\tag{\SGARefTarget{sga7.sga7i.eviii.n2-1-10-1}{2.1.10.1}}
\begin{tikzcd}[column sep=huge,row sep=huge]
 {}_nP\otimes{}_nQ
   \arrow[r,"\sim"]
   \arrow[d,"{\alpha^n_{P,Q}}"']&
 {}_nQ\otimes{}_nP
   \arrow[d,"{\alpha^n_{Q,P}}"]\\
 \uTor_1(P,Q)
   \arrow[r,"\sim"]
   \arrow[ur,phantom,"-" description]&
 \uTor_1(Q,P).
\end{tikzcd}
\end{equation}
\underline{The horizontal arrows are the symmetry isomorphisms.}

\medskip
\noindent\underline{Proof.}
The procedure of \SGARefLink{sga7.sga7i.eviii.n2-1-8}{2.1.8} immediately reduces us to the case
$P=Q=(\mathbb Z/n\mathbb Z)_T$.  We must then prove that the two
isomorphisms (\SGARefLink{sga7.sga7i.eviii.n2-1-8-1}{2.1.8.1}), obtained by resolving respectively the first
and second arguments of $\uTor_1$ by (\SGARefLink{sga7.sga7i.eviii.n2-1-8-2}{2.1.8.2}), differ by a sign
(equivalently, that the symmetry isomorphism on
$\uTor^{\mathbb Z}_1(\mathbb Z/n\mathbb Z,\mathbb Z/n\mathbb Z)$ is
$-\id$).  We may plainly suppose that $T$ is the punctual topos, so
that we work with ordinary commutative groups.  Let $L_\bullet$ be the
complex
\[
 L_\bullet=\bigl(0\longrightarrow\mathbb Z
 \xrightarrow{n}\mathbb Z\longrightarrow0\bigr),
\]
regarded as a flat resolution of
$P=Q=\mathbb Z/n\mathbb Z$.  Use
$L_\bullet\otimes L_\bullet$ to compute
\[
 \Tor_1(P,Q)=\Tor_1(\mathbb Z/n\mathbb Z,\mathbb Z/n\mathbb Z).
\]
The isomorphisms with
$H_1(L_\bullet\otimes\mathbb Z/n\mathbb Z)$, which defines $\alpha$,
and with
$H_1(\mathbb Z/n\mathbb Z\otimes L_\bullet)$, which defines $\beta$,
are induced respectively by the augmentations
$L_\bullet\longrightarrow\mathbb Z/n\mathbb Z$ of the second and
first factors of $L_\bullet\otimes L_\bullet$.  It follows at once
that $\alpha(1,1)$ is represented by the cycle
\[
 -1_0\otimes1_1+1_1\otimes1_0
 \in(L_\bullet\otimes L_\bullet)_1
 =L_0\otimes L_1+L_1\otimes L_0,
\]
where $1_0$ and $1_1$ denote the element $1$ of $L_0$ and $L_1$
respectively, whereas $\beta(1,1)$ is represented by
$1_0\otimes1_1-1_1\otimes1_0$.  This proves \SGARefLink{sga7.sga7i.eviii.n2-1-10}{2.1.10}.

To finish these generalities on the pairings $\alpha$, we must still
specify the relations among the homomorphisms $\alpha^n_{P,Q}$ as
$n$ varies.  The main result is the following.

\medskip
\noindent\underline{Lemma \SGARefTarget{sga7.sga7i.eviii.n2-1-11}{2.1.11}.}
\underline{Let $P,Q$ be two commutative groups of $T$, and let $n,n'$
be positive integers such that $n'$ is a multiple of $n$, say
$n'=mn$.  Then the following diagram is commutative:}
\begin{equation}\tag{\SGARefTarget{sga7.sga7i.eviii.n2-1-11-1}{2.1.11.1}}
\begin{tikzcd}[column sep=large,row sep=large]
 &{}_{n'}P\otimes{}_nQ
   \arrow[ld,"P_m\otimes\id_{{}_nQ}"']
   \arrow[rd,"\id_{{}_{n'}P}\otimes i_m"]&\\
 {}_nP\otimes{}_nQ
   \arrow[rd,"{\alpha^n_{P,Q}}"']&&
 {}_{n'}P\otimes{}_{n'}Q
   \arrow[ld,"{\alpha^{n'}_{P,Q}}"]\\
 &\uTor_1(P,Q)&
\end{tikzcd}
\end{equation}
\underline{Here $P_m:{}_{n'}P\longrightarrow{}_nP$ is the homomorphism
induced by $m\cdot\id_P$, and
$i_m:{}_nQ\longrightarrow{}_{n'}Q$ is the inclusion.  Likewise, the
diagram (\SGARefTarget{sga7.sga7i.eviii.n2-1-11-2}{2.1.11.2}) symmetric to (\SGARefLink{sga7.sga7i.eviii.n2-1-11-1}{2.1.11.1}) is commutative; it has the
same vertices except that the upper vertex
$\,{}_{n'}P\otimes{}_nQ$ is replaced by
$\,{}_nP\otimes{}_{n'}Q$.}

\medskip
It suffices to prove the first assertion; the second follows by
symmetry in view of \SGARefLink{sga7.sga7i.eviii.n2-1-10}{2.1.10}.  By the procedure of \SGARefLink{sga7.sga7i.eviii.n2-1-8}{2.1.8}, we are reduced
to the punctual topos, where $P$ and $Q$ are ordinary commutative
groups and
\[
 P=\mathbb Z/n'\mathbb Z,\qquad Q=\mathbb Z/n'\mathbb Z.
\]
In fact, for $P$ we use only the hypothesis $n'P=0$.  Using the
resolution
\[
 0\longrightarrow\mathbb Z\xrightarrow{n'}\mathbb Z
 \longrightarrow0
\]
of $Q$, identify $\Tor_1(P,Q)$ with $\,{}_{n'}P=P$.  Likewise, using
\[
 0\longrightarrow\mathbb Z\xrightarrow n\mathbb Z
 \longrightarrow0
\]
to resolve
$\,{}_nQ\simeq\mathbb Z/n\mathbb Z$, identify
$\Tor_1({}_nP,{}_nQ)$ with $\,{}_nP$.  The augmentation in the latter
resolution sends $1_0$ to $m\cdot1_Q$, where $1_Q$ denotes the element
$1$ of $Q=\mathbb Z/n'\mathbb Z$.  Under these identifications the
homomorphism
\[
\begin{tikzcd}[column sep=large,row sep=large]
 \Tor_1({}_nP,{}_nQ)\arrow[r]\arrow[d,no head,"\wr"']&
 \Tor_1(P,Q)\arrow[d,no head,"\wr"']\\
 {}_nP&{}_{n'}P=P
\end{tikzcd}
\]
induced by the inclusions
$\,{}_nP\hookrightarrow P$ and $\,{}_nQ\hookrightarrow Q$ is simply
the natural inclusion $\,{}_nP\hookrightarrow{}_{n'}P$.

Consequently, identify the upper vertex of (\SGARefLink{sga7.sga7i.eviii.n2-1-11-1}{2.1.11.1}) with $P/nP$
using the generator $m\cdot1_Q$ of $\,{}_nQ$, which identifies
$\,{}_nQ$ with $\mathbb Z/n\mathbb Z$.  The diagram becomes
\[
\begin{tikzcd}[column sep=large,row sep=large]
 &P/nP
   \arrow[ld,"x\longmapsto mx"']
   \arrow[rd,"x\longmapsto mx"]&\\
 {}_nP\arrow[rd,"x\longmapsto x"']&&
 P\arrow[ld,"x\longmapsto x"]\\
 &{}_{n'}P&
\end{tikzcd}
\]
which is plainly commutative.

\medskip
\noindent\underline{Corollary \SGARefTarget{sga7.sga7i.eviii.n2-1-12}{2.1.12}.}
\underline{With the notation of \SGARefLink{sga7.sga7i.eviii.n2-1-11}{2.1.11}, the following diagram is
commutative:}
\begin{equation}\tag{\SGARefTarget{sga7.sga7i.eviii.n2-1-12-1}{2.1.12.1}}
\begin{tikzcd}[column sep=huge,row sep=huge]
 {}_{n'}P\otimes{}_{n'}Q
   \arrow[r,"m\otimes m"]
   \arrow[d,"{\alpha^{n'}_{P,Q}}"']&
 {}_nP\otimes{}_nQ
   \arrow[d,"{\alpha^n_{P,Q}}"]\\
 {}_{n'}\uTor_1(P,Q)
   \arrow[r,"m"]&
 {}_n\uTor_1(P,Q).
\end{tikzcd}
\end{equation}
\underline{In the description of the horizontal arrows, $m$ denotes
the homomorphisms $\,{}_{n'}E\longrightarrow{}_nE$ induced by
$m\cdot\id_E$, for $E=P,Q$, and $\uTor_1(P,Q)$ respectively.}

\medskip
Indeed, by \SGARefLink{sga7.sga7i.eviii.n2-1-11}{2.1.11}, for
$x\in{}_{n'}P(S)$ and $y\in{}_{n'}Q(S)$ one has
\[
 \alpha^n(mx\otimes my)=\alpha^{n'}(x\otimes my).
\]
The right-hand side equals
$m\alpha^{n'}(x\otimes y)$, since
$x\otimes my=m(x\otimes y)$.  This proves the commutativity of
(\SGARefLink{sga7.sga7i.eviii.n2-1-12-1}{2.1.12.1}).

% END INCLUDED SGA7I_English_current_buildable_source_20260823/components/96_expose_VIII_section_2_1_9_through_2_1_12.tex

% BEGIN INCLUDED SGA7I_English_current_buildable_source_20260823/components/97_expose_VIII_section_2_1_13_and_remark_2_1_14.tex
% SOURCE AUTHORITY:
% expose_VIII_body.tex lines 552--675; scan indices 245--248; source folios 234--237.
% Direct scan comparison: physical PDF pages 246--249 at 1100 dpi.
% Diagram 2.1.14.6: six targeted vector crops at 5000 dpi, covering every
% node, arrow, label, attachment, and terminal punctuation.
% Transparent source dispositions:
% - the authority's self-contradictory phrase "injective isomorphism" in
%   2.1.14 is rendered as "injective homomorphism";
% - in the multiplication-by-pq map following (2.1.14.6), the authority
%   omits pq before the second ideal in the target numerator; it is restored
%   as pq a \cap pq b, as required by multiplication by pq and by the diagram.
% Hard stop before 2.2, source line 676.

\SGARefTarget{sga7.sga7i.eviii.n2-1-13}{2.1.13}. Let $\ell$ be a prime number.  For every commutative group
$E$, denote by $T_\ell(E)$ the inverse system
\begin{equation}\tag{\SGARefTarget{sga7.sga7i.eviii.n2-1-13-1}{2.1.13.1}}
 T_\ell(E)=({}_{\ell^\nu}E)_{\nu\geq0},
\end{equation}
which plainly depends functorially on $E$.  The compatibilities of
\SGARefLink{sga7.sga7i.eviii.n2-1-12}{2.1.12}, for $n,n'$ of the form $\ell^\nu,\ell^{\nu'}$ with
$\nu'\geq\nu$, show that the homomorphisms
$\alpha^{\ell^\nu}_{P,Q}$, as $\nu$ varies, define a homomorphism of
inverse systems, functorial in $P,Q$:
\begin{equation}\tag{\SGARefTarget{sga7.sga7i.eviii.n2-1-13-2}{2.1.13.2}}
 \alpha^{(\ell)}_{P,Q}:
 T_\ell(P)\otimes T_\ell(Q)
 \longrightarrow T_\ell\bigl(\uTor_1(P,Q)\bigr).
\end{equation}
Knowledge of these homomorphisms for every prime $\ell$ is plainly
equivalent to knowledge of the $\alpha^n_{P,Q}$ for every integer
$n$.  The antisymmetry relations of \SGARefLink{sga7.sga7i.eviii.n2-1-10}{2.1.10} then become the analogous
$\ell$-adic antisymmetry relations, for every $\ell$, expressed by
the anticommutativity of the diagrams
\begin{equation}\tag{\SGARefTarget{sga7.sga7i.eviii.n2-1-13-3}{2.1.13.3}}
\begin{tikzcd}[column sep=huge,row sep=huge]
 T_\ell(P)\otimes T_\ell(Q)
   \arrow[r]
   \arrow[d,"{\alpha^{(\ell)}_{P,Q}}"']
   \arrow[rd,phantom,"-" description]&
 T_\ell(Q)\otimes T_\ell(P)
   \arrow[d,"{\alpha^{(\ell)}_{Q,P}}"]\\
 T_\ell\bigl(\uTor_1(P,Q)\bigr)
   \arrow[r]&
 T_\ell\bigl(\uTor_1(Q,P)\bigr).
\end{tikzcd}
\end{equation}
Here the horizontal arrows are the symmetry arrows.

\medskip
\noindent\underline{Remark \SGARefTarget{sga7.sga7i.eviii.n2-1-14}{2.1.14}.}
The construction of the homomorphism (\SGARefLink{sga7.sga7i.eviii.n2-1-1}{2.1.1}), and the results \SGARefLink{sga7.sga7i.eviii.n2-1-10}{2.1.10}
and \SGARefLink{sga7.sga7i.eviii.n2-1-11}{2.1.11} concerning it, may be generalized to the situation in
which one is given a commutative ring $\underline A$ on $T$ and two
ideals $\underline a$ and $\underline b$ of $\underline A$.  In the
case treated above, both ideals are the ideal
$n\mathbb Z_T$ of the constant ring $\mathbb Z_T$ on $T$.  One then
defines a homomorphism, functorial in the $\underline A$-modules
$P,Q$,
\begin{equation}\tag{\SGARefTarget{sga7.sga7i.eviii.n2-1-14-1}{2.1.14.1}}
 {}_{\underline a}P\otimes{}_{\underline b}Q
 \xrightarrow{\ \alpha\ }
 \uHom_{\underline A}
 \left(
   \frac{\underline a\cap\underline b}
        {\underline a\mathbin{\cdot}\underline b},
   \uTor^{\underline A}_1(P,Q)
 \right),
\end{equation}
as the composite of the homomorphisms
\[
\begin{split}
 {}_{\underline a}P\otimes{}_{\underline b}Q
 &\longrightarrow
 \uHom(\underline A/\underline a,P)
 \otimes
 \uHom(\underline A/\underline b,Q)\\
 &\longrightarrow
 \uHom\!\left(
   \uTor_1(\underline A/\underline a,\underline A/\underline b),
   \uTor_1(P,Q)
 \right),
\end{split}
\]
using the well-known canonical isomorphism
\begin{equation}\tag{\SGARefTarget{sga7.sga7i.eviii.n2-1-14-2}{2.1.14.2}}
 \uTor_1(\underline A/\underline a,\underline A/\underline b)
 \simeq
 \frac{\underline a\cap\underline b}
      {\underline a\mathbin{\cdot}\underline b}.
\end{equation}
In fact, the sign in \SGARefLink{sga7.sga7i.eviii.n2-1-10}{2.1.10} is ``hidden'' in the choice of this
isomorphism, according as one works with the resolution
$0\longrightarrow\underline a\longrightarrow\underline A
\longrightarrow0$ of $\underline A/\underline a$, or with the
resolution $0\longrightarrow\underline b\longrightarrow\underline A
\longrightarrow0$ of $\underline A/\underline b$.  Here it becomes
the anticommutativity of the diagram
\begin{equation}\tag{\SGARefTarget{sga7.sga7i.eviii.n2-1-14-3}{2.1.14.3}}
\begin{tikzcd}[column sep=huge,row sep=huge]
 {}_{\underline a}P\otimes{}_{\underline b}Q
   \arrow[r]
   \arrow[d]
   \arrow[rd,phantom,"-1" description]&
 {}_{\underline b}Q\otimes{}_{\underline a}P
   \arrow[d]\\
 \uHom\!\left(
   \dfrac{\underline a\cap\underline b}
         {\underline a\mathbin{\cdot}\underline b},
   \uTor_1(P,Q)
 \right)
   \arrow[r]&
 \uHom\!\left(
   \dfrac{\underline a\cap\underline b}
         {\underline a\mathbin{\cdot}\underline b},
   \uTor_1(Q,P)
 \right).
\end{tikzcd}
\end{equation}
The horizontal arrows are the symmetry arrows.  When
$\underline a=\underline b$, this anticommutativity may also be
expressed by saying that the symmetry automorphism on
$\uTor_1(\underline A/\underline a,\underline A/\underline a)$ is
$-\id$.

In the special case where
$\underline a=\underline b=n\underline A$, with $n$ a section of
$\underline A$, multiplication by $n$ induces an epimorphism
\[
 \underline A/\underline a
 \longrightarrow
 \frac{\underline a\cap\underline a}
      {\underline a\mathbin{\cdot}\underline a}
 =\underline a/\underline a^2,
\]
and hence a functorial injective homomorphism, in the
$\underline A$-module $T$,
\[
 \uHom(\underline a/\underline a^2,T)
 \longrightarrow{}_{\underline a}T.
\]
Thus in this case the homomorphism (\SGARefLink{sga7.sga7i.eviii.n2-1-14-1}{2.1.14.1}) may be interpreted as a
homomorphism
\begin{equation}\tag{\SGARefTarget{sga7.sga7i.eviii.n2-1-14-4}{2.1.14.4}}
 {}_nP\otimes{}_nQ
 \longrightarrow{}_n\uTor_1(P,Q),
\end{equation}
which is closer in form to (\SGARefLink{sga7.sga7i.eviii.n2-1-1}{2.1.1}).

Continue to suppose that
$\underline a=\underline b=n\underline A$, and that $n$ is a regular
section of $\underline A$.  The spectral sequence (\SGARefLink{sga7.sga7i.eviii.n2-1-4}{2.1.4}) still makes
sense in the present setting and yields the following exact sequence
of low-degree terms, generalizing (\SGARefLink{sga7.sga7i.eviii.n2-1-5}{2.1.5}), where $Q$ is now an
$\underline A$-module annihilated by $\underline a$:
\begin{equation}\tag{\SGARefTarget{sga7.sga7i.eviii.n2-1-14-5}{2.1.14.5}}
\begin{split}
 \uTor^{\underline A}_2(P,Q)
 &\longrightarrow
 \uTor^{\underline A/n\underline A}_2(P_n,Q)
 \longrightarrow{}_nP\otimes Q
 \xrightarrow{\ \alpha\ }
 \uTor^{\underline A}_1(P,Q)\\
 &\longrightarrow
 \uTor^{\underline A/n\underline A}_1(P_n,Q)
 \longrightarrow0.
\end{split}
\end{equation}

Return now to arbitrary $\underline a$ and $\underline b$, and let
$p$ and $q$ be two sections of $\underline A$.  We then obtain the
commutative diagram below; (\SGARefLink{sga7.sga7i.eviii.n2-1-11-1}{2.1.11.1}) is essentially the special case
obtained by taking $q=1$:
\begin{equation}\tag{\SGARefTarget{sga7.sga7i.eviii.n2-1-14-6}{2.1.14.6}}
{\footnotesize
\begin{tikzcd}[row sep=large,column sep=tiny]
 &{}_{p\underline a}P\otimes{}_{q\underline b}Q
   \arrow[dl,"p\otimes q"']
   \arrow[dr,"i\otimes j"]&\\
 {}_{\underline a}P\otimes{}_{\underline b}Q
   \arrow[d,"\alpha"']&&
 {}_{pq\underline a}P\otimes{}_{pq\underline b}Q
   \arrow[d,"\alpha"]\\
 \uHom\!\left(
   \dfrac{\underline a\cap\underline b}
         {\underline a\mathbin{\cdot}\underline b},
   \uTor_1(P,Q)
 \right)
   \arrow[dr,"\mathrm{incl.}"']&&
 \uHom\!\left(
   \dfrac{pq\underline a\cap pq\underline b}
         {pq\underline a\mathbin{\cdot}pq\underline b},
   \uTor_1(P,Q)
 \right)
   \arrow[dl,"(pq)^*"]\\
 &\uHom\!\left(
   \dfrac{\underline a\cap\underline b}
         {pq\underline a\mathbin{\cdot}\underline b},
   \uTor_1(P,Q)
 \right)&
\end{tikzcd}}
\end{equation}
Here
\[
 i:{}_{p\underline a}P\longrightarrow{}_{pq\underline a}P,
 \qquad
 j:{}_{q\underline b}Q\longrightarrow{}_{pq\underline b}Q
\]
are the canonical inclusions, and $(pq)^*$ denotes the injective
homomorphism obtained, by passing to $\uHom$, from the surjective
homomorphism
\[
 \frac{\underline a\cap\underline b}
      {pq\underline a\mathbin{\cdot}\underline b}
 \longrightarrow
 \frac{pq\underline a\cap pq\underline b}
      {pq\underline a\mathbin{\cdot}pq\underline b}
\]
induced by multiplication by $pq$.  When
$\underline a=\underline b=n\underline A$, this diagram may be
rewritten in a simpler form, given in (\SGARefLink{sga7.sga7i.eviii.n2-1-11-1}{2.1.11.1}) when $q=1$.

% END INCLUDED SGA7I_English_current_buildable_source_20260823/components/97_expose_VIII_section_2_1_13_and_remark_2_1_14.tex

% BEGIN INCLUDED SGA7I_English_current_buildable_source_20260823/components/98_expose_VIII_section_2_2.tex
% SOURCE AUTHORITY:
% expose_VIII_body.tex lines 676--834; scan indices 248--252; source folios 237--241.
% Direct scan comparison: physical PDF pages 249--253 at 1100 dpi.
% Transparent source dispositions:
% - in 2.2.7, the restriction is to the image of alpha^n_{P,Q} in
%   Tor_1(P,Q), not the authority's impossible phi-map to Tor_1(P,Q;G);
% - the cited pairings are (2.2.5) and (2.2.6), not (2.2.4) and (2.2.6);
% - typewritten letter-o complex subscripts are normalized to zero;
% - the missing closing parenthesis in the lower-left node of (2.2.11.2)
%   is restored.
% Hard stop before 2.3, source line 836.

\SGARefTarget{sga7.sga7i.eviii.n2-2}{2.2}. Consider now, in addition to $P$ and $Q$, a commutative group
$G$.  In (\SGARefLink{sga7.sga7i.eviii.n1-1-3}{1.1.3}) we defined a canonical homomorphism, functorial in
$P,Q$, and $G$,
\begin{equation}\tag{\SGARefTarget{sga7.sga7i.eviii.n2-2-1}{2.2.1}}
 \mathrm{Biext}^1(P,Q;G)
 \longrightarrow
 \operatorname{Hom}\bigl(\operatorname{Tor}_1(P,Q),G\bigr).
\end{equation}
On the other hand, if $n$ is a positive integer (respectively, if
$\ell$ is a prime number), the homomorphism $\alpha^n_{P,Q}$ of
(\SGARefLink{sga7.sga7i.eviii.n2-1-1}{2.1.1}) (respectively, the homomorphism $\alpha^{(\ell)}_{P,Q}$ of
(\SGARefLink{sga7.sga7i.eviii.n2-1-13-2}{2.1.13.2})) gives a homomorphism
\[
 \operatorname{Hom}\bigl(\operatorname{Tor}_1(P,Q),G\bigr)
 \longrightarrow
 \operatorname{Hom}({}_nP\otimes{}_nQ,G)
\]
(respectively, a homomorphism
\[
 \operatorname{Hom}\bigl(\operatorname{Tor}_1(P,Q),G\bigr)
 \longrightarrow
 \operatorname{Hom}\bigl(
   T_\ell(P)\otimes T_\ell(Q),T_\ell(G)
 \bigr)).
\]
These last homomorphisms are also functorial in $P,Q,G$.
Composing them with (\SGARefLink{sga7.sga7i.eviii.n2-2-1}{2.2.1}), we obtain canonical homomorphisms,
functorial in $P,Q,G$,
\begin{equation}\tag{\SGARefTarget{sga7.sga7i.eviii.n2-2-2}{2.2.2}}
 \varphi^n=\varphi^n_{P,Q}:
 \mathrm{Biext}^1(P,Q;G)
 \longrightarrow
 \operatorname{Hom}({}_nP\otimes{}_nQ,G),
\end{equation}
\begin{equation}\tag{\SGARefTarget{sga7.sga7i.eviii.n2-2-3}{2.2.3}}
 \varphi^{(\ell)}=\varphi^{(\ell)}_{P,Q}:
 \mathrm{Biext}^1(P,Q;G)
 \longrightarrow
 \operatorname{Hom}\bigl(
   T_\ell(P)\otimes T_\ell(Q),T_\ell(G)
 \bigr).
\end{equation}
For
\begin{equation}\tag{\SGARefTarget{sga7.sga7i.eviii.n2-2-4}{2.2.4}}
 \xi\in\mathrm{Biext}^1(P,Q;G),
\end{equation}
denote its image under (\SGARefLink{sga7.sga7i.eviii.n2-2-2}{2.2.2}) (respectively, under (\SGARefLink{sga7.sga7i.eviii.n2-2-3}{2.2.3})) by
$\varphi^n_\xi$ (respectively, by $\varphi^{(\ell)}_\xi$):
\begin{equation}\tag{\SGARefTarget{sga7.sga7i.eviii.n2-2-5}{2.2.5}}
 \varphi^n_\xi:{}_nP\otimes{}_nQ
 \longrightarrow G
 \qquad\text{(or, equivalently, }{}_nG\text{)},
\end{equation}
\begin{equation}\tag{\SGARefTarget{sga7.sga7i.eviii.n2-2-6}{2.2.6}}
 \varphi^{(\ell)}_\xi:
 T_\ell(P)\otimes T_\ell(Q)
 \longrightarrow T_\ell(G).
\end{equation}

\SGARefTarget{sga7.sga7i.eviii.n2-2-7}{2.2.7}. Plainly, knowledge of $\varphi^{(\ell)}_\xi$ is equivalent
to knowledge of the $\varphi^n_\xi$ for all $n=\ell^\nu$, with
$\nu\geq0$.  For fixed $n$, knowledge of $\varphi^n_\xi$ is
equivalent to knowledge, on the image of
\[
 \alpha^n_{P,Q}:{}_nP\otimes{}_nQ\longrightarrow\uTor_1(P,Q),
\]
of the homomorphism
\begin{equation}\tag{\SGARefTarget{sga7.sga7i.eviii.n2-2-8}{2.2.8}}
 \uTor_1(P,Q)\longrightarrow G
\end{equation}
which is the image of $\xi$ under (\SGARefLink{sga7.sga7i.eviii.n2-2-1}{2.2.1}).

Observe that if $P=nP$ and $Q=nQ$, then the preceding image is
$\,{}_n\uTor_1(P,Q)$.  This follows readily by considering the exact
sequence (\SGARefLink{sga7.sga7i.eviii.n1-3-5}{1.3.5}) and the analogous exact sequence for $Q$.
Furthermore, if $P$ and $Q$ are divisible, that is, if $nP=P$ and
$nQ=Q$ for every positive integer $n$, one readily checks that
$\uTor_1(P,Q)$ is torsion, that is,
\[
 \uTor_1(P,Q)=\varinjlim_n{}_{n}\uTor_1(P,Q).
\]
Thus, in this case, knowledge of the $\varphi^n_\xi$ for all $n$ is
equivalent to knowledge of (\SGARefLink{sga7.sga7i.eviii.n2-2-8}{2.2.8}), whereas knowledge of
$\varphi^{(\ell)}_\xi$ for a given prime $\ell$ is equivalent to
knowledge of the restriction of (\SGARefLink{sga7.sga7i.eviii.n2-2-8}{2.2.8}) to the $\ell$-primary
component of $\uTor_1(P,Q)$.

We now make explicit the symmetry properties of the pairings (\SGARefLink{sga7.sga7i.eviii.n2-2-5}{2.2.5})
and (\SGARefLink{sga7.sga7i.eviii.n2-2-6}{2.2.6}) associated with biextensions.  First observe that there is
a commutative diagram
\begin{equation}\tag{\SGARefTarget{sga7.sga7i.eviii.n2-2-9}{2.2.9}}
\begin{tikzcd}[column sep=large,row sep=huge]
 \mathrm{Biext}^1(P,Q;G)
   \arrow[r,"\sim"]
   \arrow[d]&
 \mathrm{Biext}^1(Q,P;G)
   \arrow[d]\\
 \operatorname{Hom}\bigl(\uTor_1(P,Q),G\bigr)
   \arrow[r,"\sim"]&
 \operatorname{Hom}\bigl(\uTor_1(Q,P),G\bigr),
\end{tikzcd}
\end{equation}
where the horizontal arrows are the symmetry arrows (the first of
which was defined in VII \SGARefLink{sga7.sga7i.evii.n2-7}{2.7}), and the vertical arrows are (\SGARefLink{sga7.sga7i.eviii.n2-2-1}{2.2.1}).
The commutativity follows from that of the two constituent squares in
the fuller diagram
\begin{equation}\tag{\SGARefTarget{sga7.sga7i.eviii.n2-2-10}{2.2.10}}
\begin{tikzcd}[column sep=large,row sep=large]
 \mathrm{Biext}^1(P,Q;G)
   \arrow[r,"\sim"]
   \arrow[d]&
 \mathrm{Biext}^1(Q,P;G)
   \arrow[d]\\
 \operatorname{Ext}^1\!\left(
   P\mathbin{\overset{\mathbb L}{\otimes}}Q,G
 \right)
   \arrow[r,"\sim"]
   \arrow[d]&
 \operatorname{Ext}^1\!\left(
   Q\mathbin{\overset{\mathbb L}{\otimes}}P,G
 \right)
   \arrow[d]\\
 \operatorname{Hom}\bigl(\uTor_1(P,Q),G\bigr)
   \arrow[r,"\sim"]&
 \operatorname{Hom}\bigl(\uTor_1(Q,P),G\bigr).
\end{tikzcd}
\end{equation}
The first vertical arrow in each column is the canonical isomorphism
of VII \SGARefLink{sga7.sga7i.evii.n3-6-5}{3.6.5}, and the second is the evident homomorphism.  The second
square in (\SGARefLink{sga7.sga7i.eviii.n2-2-10}{2.2.10}) is trivially commutative: by definition, the
symmetry isomorphism
\[
 P\mathbin{\overset{\mathbb L}{\otimes}}Q
 \xrightarrow{\ \sim\ }
 Q\mathbin{\overset{\mathbb L}{\otimes}}P
\]
induces the symmetry isomorphism
$\uTor_1(P,Q)\simeq\uTor_1(Q,P)$, and, for a variable complex
$L_\bullet$, the canonical isomorphism
\[
 \operatorname{Ext}^1(L_\bullet,G)
 \simeq
 \operatorname{Hom}\bigl(H_1(L_\bullet),G\bigr)
\]
is functorial in $L_\bullet$.

As for the commutativity of the first square of (\SGARefLink{sga7.sga7i.eviii.n2-2-10}{2.2.10}), it follows
by making explicit simultaneously the symmetry isomorphisms and the
vertical isomorphisms defined in VII \SGARefLink{sga7.sga7i.evii.n3-6-5}{3.6.5}.  We leave the details of
the verification to the reader.  One uses the fact that the symmetry
isomorphism
\[
 L_\bullet(P)\otimes L_\bullet(Q)
 \simeq
 L_\bullet(Q)\otimes L_\bullet(P)
\]
induces in degree $1$ the isomorphism
\[
 L_0(P)\otimes L_1(Q)+L_1(P)\otimes L_0(Q)
 \simeq
 L_0(Q)\otimes L_1(P)+L_1(Q)\otimes L_0(P),
\]
which, on corresponding summands, is given by the ordinary symmetry
isomorphism of the underlying ungraded modules (and hence introduces
no sign).

In view of the commutativity of (\SGARefLink{sga7.sga7i.eviii.n2-2-9}{2.2.9}), the anticommutativity of
(\SGARefLink{sga7.sga7i.eviii.n2-1-10-1}{2.1.10.1}), in terms of the $\varphi^n$ or
$\varphi^{(\ell)}$, becomes the following relation.

\medskip
\noindent\underline{Proposition \SGARefTarget{sga7.sga7i.eviii.n2-2-11}{2.2.11}.}
\underline{The following diagrams are anticommutative:}
\begin{equation}\tag{\SGARefTarget{sga7.sga7i.eviii.n2-2-11-1}{2.2.11.1}}
\begin{tikzcd}[row sep=large,column sep=huge]
 \mathrm{Biext}^1(P,Q;G)
   \arrow[rr,"\sim"]
   \arrow[dd,"{\varphi^n_{P,Q}}"']&&
 \mathrm{Biext}^1(Q,P;G)
   \arrow[dd,"{\varphi^n_{Q,P}}"]\\
 &- &\\
 \operatorname{Hom}({}_nP\otimes{}_nQ,G)
   \arrow[rr,"\sim"]&&
 \operatorname{Hom}({}_nQ\otimes{}_nP,G),
\end{tikzcd}
\end{equation}
\begin{equation}\tag{\SGARefTarget{sga7.sga7i.eviii.n2-2-11-2}{2.2.11.2}}
\begin{tikzcd}[row sep=large,column sep=huge]
 \mathrm{Biext}^1(P,Q;G)
   \arrow[rr,"\sim"]
   \arrow[dd,"{\varphi^{(\ell)}_{P,Q}}"']&&
 \mathrm{Biext}^1(Q,P;G)
   \arrow[dd,"{\varphi^{(\ell)}_{Q,P}}"]\\
 &- &\\
 \operatorname{Hom}\bigl(
   T_\ell(P)\otimes T_\ell(Q),T_\ell(G)
 \bigr)
   \arrow[rr,"\sim"]&&
 \operatorname{Hom}\bigl(
   T_\ell(Q)\otimes T_\ell(P),T_\ell(G)
 \bigr),
\end{tikzcd}
\end{equation}
\underline{where the horizontal arrows are the symmetry arrows.}

We record the important corollary.

\medskip
\noindent\underline{Corollary \SGARefTarget{sga7.sga7i.eviii.n2-2-12}{2.2.12}.}
\underline{Suppose that $P=Q$, and suppose that $\xi$ is the class of
a symmetric (respectively, antisymmetric) biextension of $(P,P)$ by
$G$; that is, suppose ${}^s\xi=\xi$ (respectively,
${}^s\xi=-\xi$), where ${}^s\xi$ denotes the image of $\xi$ under
symmetry.  Then the bilinear maps $\varphi^n_\xi$, on
${}_nP\times{}_nP$, and $\varphi^{(\ell)}_\xi$, on
$T_\ell(P)\times T_\ell(P)$, are antisymmetric (respectively,
symmetric).}

\medskip
\noindent\underline{Remark \SGARefTarget{sga7.sga7i.eviii.n2-2-13}{2.2.13}.}
Suppose that $\xi$ is symmetric.  If $n$ (respectively, $\ell$) is
odd, then \SGARefLink{sga7.sga7i.eviii.n2-2-12}{2.2.12} shows that $\varphi^n_\xi$ (respectively,
$\varphi^{(\ell)}_\xi$) is in fact an alternating form.  This
conclusion is not valid in general without the restriction that $n$
(respectively, $\ell$) be odd, as is already seen when $P$ is the
ordinary group $\mathbb Z/2\mathbb Z$ and
$G=\mathbb Z/2\mathbb Z$.

Suppose, however, that $2P=P$.  This is the case in particular when
$P$ is an abelian scheme, working as usual in the fppf or fpqc
topology.  The preceding forms are then still alternating.  It is
enough to prove that $\varphi^{2^\nu}_\xi(x,x)=0$.  Since $2P=P$, a
given element $x\in{}_{2^\nu}P(S)$ may be written as $x=2y$, with
$y\in{}_{2^{\nu+1}}P(S)$.  Hence
\[
 \varphi^{2^\nu}_\xi(x,x)
 =2\varphi^{2^{\nu+1}}_\xi(y,y)=0,
\]
because $\varphi^{2^{\nu+1}}_\xi$ is antisymmetric.

% END INCLUDED SGA7I_English_current_buildable_source_20260823/components/98_expose_VIII_section_2_2.tex

% BEGIN INCLUDED SGA7I_English_current_buildable_source_20260823/components/99_expose_VIII_section_2_3_opening_through_proposition_2_3_11.tex
% SOURCE AUTHORITY:
% expose_VIII_body.tex lines 836--923; scan indices 252--254; source folios 241--243.
% Direct scan comparison: physical PDF pages 253--255 at 1100 dpi.
% Transparent source disposition:
% - the right-hand vertex of (2.3.10) is restored as
%   Hom({}_nP,{}_n Ext^1(nQ,G)); the printed/transcribed "nQnG" is a
%   comma-key slip and is incompatible with (2.3.5), (2.3.9), and the prose.
% The authority's numbering jump from (2.3.2) to (2.3.4) is preserved.
% Hard stop before the proof of 2.3.11, source line 924.

\SGARefTarget{sga7.sga7i.eviii.n2-3}{2.3}. \underline{Relation with ``Cartier duality''.}
The description given in \SGARefLink{sga7.sga7i.eviii.n2-2}{2.2} of the pairings defined by a
biextension, through the isomorphism theorem VII \SGARefLink{sga7.sga7i.evii.n3-6-5}{3.6.5} and the
homomorphism (\SGARefLink{sga7.sga7i.eviii.n2-2-1}{2.2.1}), is plainly not the definition formerly used for
abelian schemes, at a time when the very notion of biextension did not
yet exist.  That definition proceeds through Cartier duality, and we
now give its relation to the present one.

Consider again the homomorphism
\begin{equation}\tag{\SGARefTarget{sga7.sga7i.eviii.n2-3-1}{2.3.1}}
 \mathrm{Biext}^1(P,Q;G)
 \longrightarrow
 \operatorname{Hom}\bigl(P,\uExt^1(Q,G)\bigr)
\end{equation}
arising from the exact sequence (\SGARefLink{sga7.sga7i.eviii.n1-4}{1.4}).  It is in fact an isomorphism
when
\begin{equation}\tag{\SGARefTarget{sga7.sga7i.eviii.n2-3-2}{2.3.2}}
 \uHom(Q,G)=0.
\end{equation}
This condition holds in particular when $Q$ is an abelian scheme over
a scheme $S$ and $G$ is the multiplicative group $\underline G_m$
(VII \SGARefLink{sga7.sga7i.evii.n1-3-8}{1.3.8}).

There is also an evident ``restriction to $nQ$'' homomorphism
\[
 \uExt^1(Q,G)\longrightarrow\uExt^1(nQ,G),
\]
and hence a composite homomorphism
\begin{equation}\tag{\SGARefTarget{sga7.sga7i.eviii.n2-3-4}{2.3.4}}
\begin{split}
 \operatorname{Hom}\bigl(P,\uExt^1(Q,G)\bigr)
 &\longrightarrow
 \operatorname{Hom}\bigl(P,\uExt^1(nQ,G)\bigr)\\
 &\longrightarrow
 \operatorname{Hom}\bigl({}_nP,{}_n\uExt^1(nQ,G)\bigr),
\end{split}
\end{equation}
where the second arrow is restriction to $\,{}_nP$.  Composing with
(\SGARefLink{sga7.sga7i.eviii.n2-3-1}{2.3.1}), we obtain a homomorphism
\begin{equation}\tag{\SGARefTarget{sga7.sga7i.eviii.n2-3-5}{2.3.5}}
 \mathrm{Biext}^1(P,Q;G)
 \longrightarrow
 \operatorname{Hom}\bigl({}_nP,{}_n\uExt^1(nQ,G)\bigr).
\end{equation}

On the other hand, the exact sequence
\begin{equation}\tag{\SGARefTarget{sga7.sga7i.eviii.n2-3-6}{2.3.6}}
 0\longrightarrow{}_nQ\longrightarrow Q
 \xrightarrow{\ n\ }nQ\longrightarrow0
\end{equation}
gives rise to a long exact sequence of the $\uExt^i(-,G)$, and hence
to the short exact sequence
\begin{equation}\tag{\SGARefTarget{sga7.sga7i.eviii.n2-3-7}{2.3.7}}
 0\longrightarrow\uHom(Q,G)_n
 \longrightarrow\uHom({}_nQ,G)
 \xrightarrow{\ \partial\ }
 {}_n\uExt^1(nQ,G)
 \longrightarrow0,
\end{equation}
where, as usual, $E_n$ denotes $E/nE$.  We retain from this sequence
the canonical surjective homomorphism
\begin{equation}\tag{\SGARefTarget{sga7.sga7i.eviii.n2-3-8}{2.3.8}}
 \uHom({}_nQ,G)
 \xrightarrow{\hspace{3em}\partial\hspace{3em}}
 {}_n\uExt^1(nQ,G),
\end{equation}
which is an isomorphism if $\uHom(Q,G)_n=0$, and a fortiori if (\SGARefLink{sga7.sga7i.eviii.n2-3-2}{2.3.2})
holds.  From (\SGARefLink{sga7.sga7i.eviii.n2-3-8}{2.3.8}) we deduce a homomorphism
\begin{equation}\tag{\SGARefTarget{sga7.sga7i.eviii.n2-3-9}{2.3.9}}
\begin{split}
 \operatorname{Hom}\bigl({}_nP,\uHom({}_nQ,G)\bigr)
 &\simeq
 \operatorname{Hom}({}_nP\otimes{}_nQ,G)\\
 &\longrightarrow
 \operatorname{Hom}\bigl({}_nP,{}_n\uExt^1(nQ,G)\bigr),
\end{split}
\end{equation}
which is an isomorphism when (\SGARefLink{sga7.sga7i.eviii.n2-3-2}{2.3.2}) holds.

All the homomorphisms under consideration are functorial in $P,Q,G$.
The homomorphisms (\SGARefLink{sga7.sga7i.eviii.n2-3-5}{2.3.5}), (\SGARefLink{sga7.sga7i.eviii.n2-3-9}{2.3.9}), and (\SGARefLink{sga7.sga7i.eviii.n2-2-2}{2.2.2}) form the diagram
\begin{equation}\tag{\SGARefTarget{sga7.sga7i.eviii.n2-3-10}{2.3.10}}
\begin{tikzcd}[column sep=large,row sep=large]
 &\mathrm{Biext}^1(P,Q;G)
   \arrow[ddl,"\varphi^n"']
   \arrow[ddr,"(2.3.5)"]&\\
 &-&\\
 \operatorname{Hom}({}_nP\otimes{}_nQ,G)
   \arrow[rr,"(2.3.9)"]&&
 \operatorname{Hom}\bigl({}_nP,{}_n\uExt^1(nQ,G)\bigr),
\end{tikzcd}
\end{equation}
where the horizontal arrow (\SGARefLink{sga7.sga7i.eviii.n2-3-9}{2.3.9}) is an isomorphism if (\SGARefLink{sga7.sga7i.eviii.n2-3-2}{2.3.2})
holds.

\medskip
\noindent\underline{Proposition \SGARefTarget{sga7.sga7i.eviii.n2-3-11}{2.3.11}.}
\underline{The diagram (\SGARefLink{sga7.sga7i.eviii.n2-3-10}{2.3.10}) is anticommutative.}

When $\uHom(Q,G)=0$, this statement therefore recovers a
characterization of the homomorphism $\varphi^n$ of (\SGARefLink{sga7.sga7i.eviii.n2-2-2}{2.2.2}) in terms
of the homomorphisms (\SGARefLink{sga7.sga7i.eviii.n2-3-5}{2.3.5}) and (\SGARefLink{sga7.sga7i.eviii.n2-3-9}{2.3.9}).  This is the
characterization formerly used in the theory of abelian schemes,
where moreover $nQ=Q$.

% END INCLUDED SGA7I_English_current_buildable_source_20260823/components/99_expose_VIII_section_2_3_opening_through_proposition_2_3_11.tex

% BEGIN INCLUDED SGA7I_English_current_buildable_source_20260823/components/100_expose_VIII_proof_2_3_11_parts_a_and_b.tex
% SOURCE AUTHORITY:
% expose_VIII_body.tex lines 924--1021; scan indices 254--257; source folios 243--246.
% Direct scan comparison: physical PDF pages 255--258 at 1100 dpi.
% Transparent source dispositions:
% - the proof heading's printed 2.5.11 is restored as 2.3.11;
% - typewritten letter-o indices and cochain degrees are normalized to zero;
% - the arrow-key slip inside the defining formula (2.3.14) is restored as
%   the minus sign already printed in (2.3.13) and required by the cocycle.
% Hard stop before part c), source line 1022.

\medskip
\noindent\underline{Proof of \SGARefLink{sga7.sga7i.eviii.n2-3-11}{2.3.11}.}\footnote{For a more elegant
proof, see \SGARefLink{sga7.sga7i.eviii.n2-3-19}{2.3.19} below.}
Let $\xi$ be an element of $\mathrm{Biext}^1(P,Q;G)$, let
\[
 \varphi_\xi:{}_nP\otimes{}_nQ\longrightarrow G
\]
be the associated pairing, and let
\[
 u_\xi:{}_nP\longrightarrow{}_n\uExt^1(nQ,G)
\]
be the image of $\xi$ under (\SGARefLink{sga7.sga7i.eviii.n2-3-5}{2.3.5}).  We must prove that $-u_\xi$ is
obtained from $\varphi_\xi$ by the arrow (\SGARefLink{sga7.sga7i.eviii.n2-3-9}{2.3.9}), that is, that it is
the composite
\begin{equation}\tag{*}
 {}_nP
 \xrightarrow{\ \varphi_\xi\ }
 \uHom({}_nQ,G)
 \xrightarrow{\ (\SGARefLink{sga7.sga7i.eviii.n2-3-8}{2.3.8})\ }
 {}_n\uExt^1(nQ,G).
\end{equation}
To prove this, we may take a section $x$ of $\,{}_nP$ and prove that
its image under $-u_\xi$ equals its image under the composite (*).
If desired, this would allow us to reduce to the case
$P=(\mathbb Z/n\mathbb Z)_T$, but we shall not need that reduction.

\medskip
\noindent a) Choose flat resolutions $L_\bullet$ and $M_\bullet$,
\[
 0\longrightarrow L_1\longrightarrow L_0\longrightarrow0,
 \qquad
 0\longrightarrow M_1\longrightarrow M_0\longrightarrow0,
\]
of $P$ and $Q$, respectively.  A flat resolution of $P$, for example,
is obtained by writing $P$ as a quotient of a flat
$\mathbb Z_T$-module $L_0$ and taking $L_1$ to be the kernel of
$L_0\longrightarrow P$; this kernel is flat because it is
torsion-free.  We obtain the complex $L_\bullet\otimes M_\bullet$
\[
 0\longrightarrow L_1\otimes M_1
 \xrightarrow{d_2}
 L_1\otimes M_0+L_0\otimes M_1
 \xrightarrow{d_1}
 L_0\otimes M_0
 \longrightarrow0,
\]
whose sheaf $\underline H_1$ will serve as $\uTor_1(P,Q)$.

Let $L'_0$ (respectively, $M'_0$) be the subsheaf of $L_0$
(respectively, $M_0$) which is the inverse image of $\,{}_nP$
(respectively, $\,{}_nQ$).  Thus
\[
 0\longrightarrow L_1\longrightarrow L'_0\longrightarrow0
 \qquad
 \left(\text{respectively, }
 0\longrightarrow M_1\longrightarrow M'_0\longrightarrow0\right)
\]
is a flat resolution of $\,{}_nP$ (respectively, of $\,{}_nQ$).  The
homomorphism
\[
 {}_nP\otimes{}_nQ
 \longrightarrow
 \uTor_1(P,Q)=\underline H_1(L_\bullet\otimes M_\bullet)
\]
may then be described, by passage to the quotient, from the
homomorphism
\[
 L'_0\otimes M'_0
 \longrightarrow
 \underline H_1(L_\bullet\otimes M_\bullet)
\]
given by
\[
 x_0\otimes y_0
 \longmapsto
 \text{the class of the $1$-cycle }
 [nx_0]\otimes y_0-x_0\otimes[ny_0].
\]
Here the brackets indicate that the element in question is regarded
as a section of the subsheaf $L_1$ of $L_0$ (respectively, $M_1$ of
$M_0$).  This description follows from the calculation made in the
proof of \SGARefLink{sga7.sga7i.eviii.n2-1-10}{2.1.10}.

\medskip
\noindent b) Take an injective resolution
\[
 C^\bullet=\bigl(0\longrightarrow C^0
 \xrightarrow{d^0}C^1\longrightarrow\cdots\bigr)
\]
of $G$, and identify
\[
 \mathrm{Biext}^1(P,Q;G)
 \simeq
 \operatorname{Ext}^1\!\left(
   P\mathbin{\overset{\mathbb L}{\otimes}}Q,G
 \right)
 =H^1\!\left(
   \operatorname{Hom}^{\boldsymbol\cdot}
   (L_\bullet\otimes M_\bullet,C^{\boldsymbol\cdot})
 \right).
\]
The element $\xi$ may therefore be represented by a pair of
homomorphisms $f_1,f_0$, giving the anticommutative diagram
\begin{equation}\tag{\SGARefTarget{sga7.sga7i.eviii.n2-3-12}{2.3.12}}
\begin{tikzcd}[column sep=huge,row sep=huge]
 L_1\otimes M_0+L_0\otimes M_1
   \arrow[r,"d_1"]
   \arrow[d,"f_1"']
   \arrow[rd,phantom,"-" description]&
 L_0\otimes M_0
   \arrow[d,"f_0"']\\
 C^0\arrow[r,"d^0"]&C^1.
\end{tikzcd}
\end{equation}
By definition, the pairing
\[
 \varphi_\xi:{}_nP\otimes{}_nQ\longrightarrow G
\]
is obtained, by passage to the quotient, from the homomorphism
\[
 \varphi:L'_0\otimes M'_0\longrightarrow G
\]
given by
\begin{equation}\tag{\SGARefTarget{sga7.sga7i.eviii.n2-3-13}{2.3.13}}
 \varphi(x_0\otimes y_0)
 =f_1\bigl([nx_0]\otimes y_0-x_0\otimes[ny_0]\bigr)
 \in G=\ker(C^0\longrightarrow C^1).
\end{equation}
Since assertion \SGARefLink{sga7.sga7i.eviii.n2-3-11}{2.3.11} is local on $T$, we may suppose that the given
section $x$ of $\,{}_nP$ lifts to a section $x_0$ of $L'_0$.
Regarding $y_0$ as variable, the preceding formula then defines a
homomorphism
\begin{equation}\tag{\SGARefTarget{sga7.sga7i.eviii.n2-3-14}{2.3.14}}
 \varphi':M'_0\longrightarrow G,
 \qquad
 y_0\longmapsto
 f_1\bigl([nx_0]\otimes y_0-x_0\otimes[ny_0]\bigr),
\end{equation}
which passes to the quotient to give
\[
 {}_nQ\longrightarrow G,
 \qquad
 y\longmapsto\varphi_\xi(x,y),
\]
the homomorphism associated with $\varphi_\xi$ and $x$.

% END INCLUDED SGA7I_English_current_buildable_source_20260823/components/100_expose_VIII_proof_2_3_11_parts_a_and_b.tex

% BEGIN INCLUDED SGA7I_English_current_buildable_source_20260823/components/101_expose_VIII_proof_2_3_11_parts_c_through_e_and_remark_2_3_16.tex
% SOURCE AUTHORITY:
% expose_VIII_body.tex lines 1022--1221; scan indices 257--260;
% source folios 246--249.
% Direct scan comparison: physical PDF pages 258--261 in overlapping
% 1100-dpi bands.  The ambiguous right-hand component in the first diagram
% of part d was additionally examined at 5000-dpi equivalent magnification.
% Transparent source dispositions:
% - typewritten letter-o indices and cochain degrees are normalized to zero;
% - the source prints f_1 on the right vertical arrow of the degree-one
%   diagram in part d; its codomain is C^1, and the immediately following
%   formula reads g_0=n f_0, so the English diagram restores the type-correct
%   f_0;
% - the uppercase X in Remark 2.3.16 is restored to the section x used in
%   the preceding and following clauses.
% Hard stop before 2.3.17, source line 1222.

\medskip
\noindent c) We must take the image of this homomorphism under the
connecting operator $\partial$ of (\SGARefLink{sga7.sga7i.eviii.n2-3-7}{2.3.7}), associated with the exact
sequence (\SGARefLink{sga7.sga7i.eviii.n2-3-6}{2.3.6}).  Now $nQ\simeq Q/{}_nQ$ is plainly also isomorphic
to $M_0/M'_0$, so that we have a homomorphism of exact sequences
\[
\begin{tikzcd}[column sep=large,row sep=large]
 0\arrow[r]&M'_0\arrow[r]\arrow[d]&M_0\arrow[r]\arrow[d]
   &M_0/M'_0\arrow[r]\arrow[d,"\wr"']&0\\
 0\arrow[r]&{}_nQ\arrow[r]&Q\arrow[r]&nQ\arrow[r]&0&{,}
\end{tikzcd}
\]
where the last vertical arrow is an isomorphism.  This shows that the
desired section of $\uExt^1(nQ,G)$ may also be computed as the image
of $\varphi'$ in (\SGARefLink{sga7.sga7i.eviii.n2-3-14}{2.3.14}) under the connecting operator
\[
 \uHom(M'_0,G)\xrightarrow{\ \partial\ }
 \uExt^1(M_0/M'_0,G),
\]
or, equivalently, as the section of the second term defined by the
element of global $\mathrm{Ext}^1$ which is the image of $\varphi'$
under the connecting operator
\[
 \mathrm{Hom}(M'_0,G)\longrightarrow
 \mathrm{Ext}^1(M_0/M'_0,G).
\]
The usual construction of the connecting operator gives the following
recipe.  Extend the composite (\SGARefLink{sga7.sga7i.eviii.n2-3-13}{2.3.13})
\[
 M'_0\xrightarrow{\ \varphi'\ }G\longrightarrow C^0
\]
to a homomorphism
\[
 \psi:M_0\longrightarrow C^0,
\]
which is possible since $C^0$ is injective, and consider the unique
homomorphism
\begin{equation}\tag{\SGARefTarget{sga7.sga7i.eviii.n2-3-15}{2.3.15}}
 \lambda:M_0/M'_0\longrightarrow C^1
\end{equation}
which makes the diagram
\[
\begin{tikzcd}[column sep=large,row sep=large]
 0\arrow[r]&M'_0\arrow[r]\arrow[d,"\varphi'"']
   &M_0\arrow[r]\arrow[d,"\psi"']
   &M_0/M'_0\arrow[r]\arrow[d,"\lambda"']&0\\
 0\arrow[r]&G\arrow[r]&C^0\arrow[r]&C^1&&{.}
\end{tikzcd}
\]
commutative.  It is clear that (\SGARefLink{sga7.sga7i.eviii.n2-3-15}{2.3.15}) takes values in the
$1$-cocycles---that is, it is a $1$-cocycle---and that this cocycle
represents the desired element of
\[
 \mathrm{Ext}^1(nQ,G)=\mathrm{Ext}^1(M_0/M'_0,G).
\]

\medskip
\noindent d) It remains to describe explicitly the section of
$\uExt^1(nQ,G)$ obtained from $\xi$ and $x$ by means of arrow (\SGARefLink{sga7.sga7i.eviii.n2-3-5}{2.3.5})
in diagram (\SGARefLink{sga7.sga7i.eviii.n2-3-10}{2.3.10}).  For this purpose we may disregard the fact that
the section $x$ of $P$ is annihilated by $n$.  In terms of the chosen
lift $x_0$ of $x$ to a section of $L_0$, the description is as follows.
This lift defines an augmentation homomorphism
\[
 u:M_\bullet\longrightarrow L_\bullet\otimes M_\bullet,
 \qquad
 y_0\longmapsto x_0\otimes y_0,
 \quad
 y_1\longmapsto x_0\otimes y_1.
\]
The section of $\uExt^1(Q,G)$ obtained from (\SGARefLink{sga7.sga7i.eviii.n2-3-1}{2.3.1}) is therefore the
one obtained from the element $\xi$ of
$\mathrm{Ext}^1(L_\bullet\otimes M_\bullet,G)$ by composition with
the preceding homomorphism.  Thus it is represented by the
degree-one homomorphism
\[
\begin{tikzcd}[column sep=large,row sep=huge]
 0\arrow[r]&M_1\arrow[r,"d_1"]
   \arrow[d,"{y_1\mapsto f_1(x_0\otimes y_1)}"']
   \arrow[rd,phantom,"-" description]
   &M_0\arrow[r]\arrow[d,"{y_0\mapsto f_0(x_0\otimes y_0)}"]&0\\
 0\arrow[r]&C^0\arrow[r,"d^0"]&C^1&\cdots&{.}
\end{tikzcd}
\]
To restrict this class to $nQ\hookrightarrow Q$, use the isomorphism
$nQ\simeq M_0/M'_0$, which gives a flat resolution
\[
 0\longrightarrow M'_0\longrightarrow M_0\longrightarrow0
\]
of $nQ$.  The inclusion $nQ\hookrightarrow Q$ lifts to the following
homomorphism of resolutions:
\[
\begin{tikzcd}[column sep=large,row sep=huge]
 0\arrow[r]&M'_0\arrow[r]
   \arrow[d,"{y'_0\mapsto[ny'_0]}"']
   \arrow[rd,phantom,"+" description]
   &M_0\arrow[r]\arrow[d,"{y_0\mapsto ny_0}"]&0\\
 0\arrow[r]&M_1\arrow[r]&M_0\arrow[r]&0&{.}
\end{tikzcd}
\]
By composition, the desired section of $\uExt^1(nQ,G)$ is therefore
associated with the element of $\mathrm{Ext}^1(nQ,G)$ represented by
the degree-one homomorphism of complexes
\[
\begin{tikzcd}[column sep=large,row sep=large]
 0\arrow[r,dash]&M'_0\arrow[r,dash,"d_1"]\arrow[d,"g_1"']
   &M_0\arrow[r,dash]\arrow[d,"g_0"]&0&{}\\
 0\arrow[r,dash]&C^0\arrow[r,dash]&C^1\arrow[r,dash]&\dots&{,}
\end{tikzcd}
\]
where
\[
 g_1(y'_0)=f_1\bigl(x_0\otimes[ny'_0]\bigr),
 \qquad
 g_0(y_0)=n f_0(x_0\otimes y_0).
\]

\medskip
\noindent e) Using now the fact that $x$ is a section of $\,{}_nP$,
we show that the preceding element of $\mathrm{Ext}^1(nQ,G)$ is the
negative of the element defined by (\SGARefLink{sga7.sga7i.eviii.n2-3-15}{2.3.15}).  For any homomorphism
\[
 h:M_0\longrightarrow C^0,
\]
the element of $\mathrm{Ext}^1$ considered in d) is unchanged if
$g_1,g_0$ are replaced, respectively, by
\[
 g'_1=g_1+h d_1,
 \qquad
 g'_0=g_0-d^0h.
\]
We have at our disposal the homomorphism
$\psi:M_0\longrightarrow C^0$ constructed in c), which by definition
satisfies
\[
 \psi(y'_0)=
 f_1\bigl([nx_0]\otimes y'_0-x_0\otimes[ny'_0]\bigr)
\]
for every section $y'_0$ of $M'_0$.  Equivalently,
\[
 f_1\bigl(x_0\otimes[ny'_0]\bigr)
 \bigl(=g_1(y'_0)\bigr)
 =f_1\bigl([nx_0]\otimes y'_0\bigr)-\psi(y'_0).
\]
Define $h:M_0\longrightarrow C^0$ by
\[
 h(y_0)=-f_1\bigl([nx_0]\otimes y_0\bigr)+\psi(y_0).
\]
By construction, then,
\[
 g'_1=0,
 \qquad
 g'_0(y_0)=n f_0(x_0\otimes y_0)
 +d^0f_1\bigl([nx_0]\otimes y_0\bigr)-d^0\psi(y_0).
\]
Since $d^0f_1=-f_0d_1$, we have
\[
 d^0f_1\bigl([nx_0]\otimes y_0\bigr)
 =-f_0(nx_0\otimes y_0),
\]
and hence
\[
 g'_1=0,
 \qquad
 g'_0=-d^0\psi.
\]
Thus the element in question is the one defined by the homomorphism
\[
 M_0/M'_0\simeq nQ\longrightarrow C^1
\]
induced by $-d^0\psi$ upon passage to the quotient.  We therefore
obtain the negative of the element defined by (\SGARefLink{sga7.sga7i.eviii.n2-3-15}{2.3.15}).

\medskip
\noindent\underline{Remarks \SGARefTarget{sga7.sga7i.eviii.n2-3-16}{2.3.16}.}  The preceding proof actually
gives a slightly more precise compatibility than the one stated in
\SGARefLink{sga7.sga7i.eviii.n2-3-11}{2.3.11}, obtained by choosing a section $x$ of $\,{}_nP$.  This section
then defines a homomorphism
\[
 {}_nQ\longrightarrow G,
 \qquad
 y\longmapsto\varphi_\xi(x,y),
\]
and hence, by the connecting homomorphism associated with (\SGARefLink{sga7.sga7i.eviii.n2-3-6}{2.3.6}), an
element of $\mathrm{Ext}^1(nQ,G)$.  On the other hand, if we forget
that $x$ is annihilated by $n$, we consider the homomorphism
$\mathbb Z_T\longrightarrow P$ defined by $x$.  Its inverse image is a
biextension of $(\mathbb Z_T,Q)$ by $G$, whose class may be interpreted
as an element of $\mathrm{Ext}^1(Q,G)$ and therefore induces an element
of $\mathrm{Ext}^1(nQ,G)$.  With these conventions, the two elements
obtained in $\mathrm{Ext}^1(nQ,G)$ are \underline{opposite}.

% END INCLUDED SGA7I_English_current_buildable_source_20260823/components/101_expose_VIII_proof_2_3_11_parts_c_through_e_and_remark_2_3_16.tex

% BEGIN INCLUDED SGA7I_English_current_buildable_source_20260823/components/102_expose_VIII_section_2_3_17.tex
% SOURCE AUTHORITY:
% expose_VIII_body.tex lines 1222--1241; scan indices 260--262;
% source folios 249--251.
% Direct scan comparison: physical PDF pages 261--263 in overlapping
% 1100-dpi bands.
% Transparent source dispositions:
% - typewritten letter-o indices are normalized to zero;
% - the final reference to phi^n is rendered as phi_xi^n, the pairing that
%   is explicitly being identified throughout 2.3.17.
% Hard stop before 2.3.19, source line 1242.

\medskip
\noindent \SGARefTarget{sga7.sga7i.eviii.n2-3-17}{2.3.17}.  Let us use the explicit cocycle description of the
pairing
\[
 \varphi_\xi^n:{}_nP\otimes{}_nQ\longrightarrow G
\]
given in (\SGARefLink{sga7.sga7i.eviii.n2-3-12}{2.3.12}) to obtain another expression for it.  First suppose
that the biextension $E$ of class $\xi$ admits a section $s$ over
$P\times Q$.  As in VII \SGARefLink{sga7.sga7i.evii.n3-8-13}{3.8.13}, take for the flat resolution of $P$
the one with $L_0=\mathbb Z[P]$ and $L_1$ the kernel of the canonical
homomorphism $\mathbb Z[P]\longrightarrow P$, and define
$M_\bullet$ similarly.  There is then a unique biextension
trivialization of the inverse image of $E$ on $(L_0,M_0)$ whose
restriction to $P\times Q$ is given by $s$; denote this trivialization
again by $s$ (cf. VII \SGARefLink{sga7.sga7i.evii.n3-8-2}{3.8.2}).  By VII \SGARefLink{sga7.sga7i.evii.n3-8-12}{3.8.12}, the class of $\xi$ in
\[
 \mathrm{Ext}^1\!\left(
 P\mathbin{\overset{\mathbb L}{\otimes}}Q,G\right)
\]
is represented by the cocycle with components $(f,g)$ characterized by
\[
 s(x_1,y_0)=f(x_1\otimes y_0)e_{E/Q}(y_0),
 \qquad
 s(x_0,y_1)=g(x_0\otimes y_1)e_{E/P}(x_0).
\]
If $x$ and $y$ are sections of $\,{}_nP$ and $\,{}_nQ$, lift them to
the corresponding sections $x_0=[x]$ and $y_0=[y]$ of $L_0$ and
$M_0$.  Applying formula (\SGARefLink{sga7.sga7i.eviii.n2-3-13}{2.3.13}), and putting $z=s(x,y)$, gives
\begin{equation}\tag{\SGARefTarget{sga7.sga7i.eviii.n2-3-18}{2.3.18}}
 \varphi_\xi^n(x,y)
 ={}^n z/e_{E/Q}(y)-z^n/e_{E/P}(x).
\end{equation}
Here the exponent $n$ written on the left, respectively on the right,
of $z$ denotes the $n$th power for the left, respectively right,
multiplicative structure on $E$.  It therefore gives a section of $E$
over $(0,y)$, respectively $(x,0)$; the meaning of the ``quotient'' of
such an element by $e_{E/Q}(y)$, respectively $e_{E/P}(x)$, is clear.

The preceding formula (\SGARefLink{sga7.sga7i.eviii.n2-3-18}{2.3.18}) still makes sense whenever $z$ is a
global section of $E$ lifting the sections $x,y$ of
$\,{}_nP,\,{}_nQ$.  I claim that (\SGARefLink{sga7.sga7i.eviii.n2-3-18}{2.3.18}) is valid for every such
section, without assuming that a section of $E$ over $P\times Q$
exists.  The section $(x,y)$ of $P\times Q$ always lifts locally to a
section of $E$, since $E\longrightarrow P\times Q$ is an epimorphism;
thus the formula computes $\varphi_\xi^n(x,y)$ in every case.
Moreover, if $z$ is replaced by another section $z'$ of $E$ lifting
$(x,y)$, then $z'=z+g$ for a section $g$ of $G$.  The two terms on the
right-hand side of (\SGARefLink{sga7.sga7i.eviii.n2-3-18}{2.3.18}) corresponding to $z'$ are the terms
corresponding to $z$, each increased by $ng$.  Their difference is
therefore independent of $z$, showing that (\SGARefLink{sga7.sga7i.eviii.n2-3-18}{2.3.18}) defines a
well-determined homomorphism
\[
 {}_nP\otimes{}_nQ\longrightarrow G.
\]

It remains to show that this homomorphism is indeed $\varphi_\xi^n$.
For this we reduce at once to the case
\[
 P=Q=(\mathbb Z/n\mathbb Z)_T.
\]
Then $P\times Q$ is of the form $I_T$, where
$I=\mathbb Z/n\mathbb Z\times\mathbb Z/n\mathbb Z$ is a finite set.
In this case the torsor $E$ under $G_{P\times Q}$---which identifies
with a family of torsors $E_i$ under $G$, for $i\in I$---plainly
becomes trivial locally on $T$.  We are thus reduced to the already
treated case in which $E$ has a global section over $P\times Q$.

% END INCLUDED SGA7I_English_current_buildable_source_20260823/components/102_expose_VIII_section_2_3_17.tex

% BEGIN INCLUDED SGA7I_English_current_buildable_source_20260823/components/103_expose_VIII_alternative_proof_2_3_19.tex
% SOURCE AUTHORITY:
% expose_VIII_body.tex lines 1242--1342; scan indices 262--263;
% source folios 251--252.
% Direct scan comparison: physical PDF pages 263--264 in overlapping
% 1100-dpi bands.  The three-square homotopy display and the full 2.3.10
% derived-Hom diagram were additionally inspected in five targeted crops at
% 5000-dpi equivalent magnification.
% Transparent source dispositions:
% - the source's typewritten [o] is normalized to the cohomological [0];
% - the circled source mark "- ?" in the large diagram is retained exactly
%   as an unresolved sign annotation, rather than silently adjudicated.
% Hard stop before section 2.4, source line 1346 / scan index 264.

\medskip
\noindent \SGARefTarget{sga7.sga7i.eviii.n2-3-19}{2.3.19}.  \underline{Another proof of \SGARefLink{sga7.sga7i.eviii.n2-3-11}{2.3.11}} (by
Deligne).

\medskip
\noindent a) Reduction to $P=\mathbb Z/n\mathbb Z$ and to
$T=(\mathrm{Sets})$.

\medskip
\noindent b) Let $\partial$ be the following arrow from $nQ$ to
$\,{}_nQ[1]$:
\[
 nQ\xleftarrow[\text{``$n$''}]{\ \sim\ }
 [\,{}_nQ\longrightarrow Q\,]
 \longrightarrow{}_nQ[1].
\]
The following diagram is anticommutative:
\[
\begin{tikzcd}[row sep=huge,column sep=huge]
 nQ[0]\arrow[r,"1"]\arrow[d,"\partial"']
   &{[Q\xrightarrow{n}Q]}\\
 {}_nQ[1]\arrow[ru,"1"']&{.}
\end{tikzcd}
\]
Indeed, the arrows
\[
\begin{tikzcd}[column sep=small,row sep=large]
 {[{}_nQ}\arrow[r]\arrow[dd,"0"']&{Q]}\arrow[dd,"n"]
   &&{[{}_nQ}\arrow[r]\arrow[dd,"1"']&{Q]}\arrow[dd,"0"']
   &\text{are opposite}
   &{[{}_nQ}\arrow[r]&{Q]}\arrow[ddl,"1"]\\
 &&\text{and}&&&\text{by means of}\quad H\,:&&\\
 {[Q}\arrow[r,"n"]&{Q]}&&{[Q}\arrow[r,"n"]&{Q]}
   &&{[Q}\arrow[r,"n"]&{Q]}
\end{tikzcd}
\]

\medskip
\noindent c) Diagram \SGARefLink{sga7.sga7i.eviii.n2-3-10}{2.3.10} is
\[
\begin{tikzcd}[column sep=small,row sep=small]
 &&&H^1\mathrm{RHom}
   (\mathbb Z\mathbin{\overset{\mathbb L}{\otimes}}Q,G)
   \arrow[ddd]\\
 &&&+\\
 H^1\mathrm{RHom}
   (\mathbb Z/n\mathbb Z\mathbin{\overset{\mathbb L}{\otimes}}Q,G)
   \arrow[rr]\arrow[dd]\arrow[uurrr]
 &&H^1\mathrm{RHom}
   (\mathbb Z/n\mathbb Z,\mathrm{RHom}(Q,G))
   \arrow[dd]&\\
 &&&\mathrm{Ext}^1(Q,G)\arrow[ddl]\\
 \mathrm{Hom}(\mathrm{Tor}_1(\mathbb Z/n\mathbb Z,Q),G)\arrow[d]
 &|[draw,circle,inner sep=1pt]|{-\,?}
 &\mathrm{Hom}(\mathbb Z/n\mathbb Z,\mathrm{Ext}^1(Q,G))
   \arrow[d]&\\
 \mathrm{Hom}({}_nQ,G)\arrow[rr]
 &&\mathrm{Ext}^1(nQ,G)&
\end{tikzcd}
\]
and:

\smallskip
\noindent c1) We have an isomorphism
\[
 \mathbb Z/n\mathbb Z
 \mathbin{\overset{\mathbb L}{\otimes}}Q
 \simeq[Q\xrightarrow{n}Q],
\]
and, under this isomorphism, the first composite vertical arrow is
induced by
\[
 {}_nQ[1]\longrightarrow[Q\xrightarrow{n}Q].
\]

\smallskip
\noindent c2) The second horizontal arrow is the transpose of the
arrow $\partial$ displayed above in b).

\smallskip
\noindent c3) If the diagram is traversed clockwise, it is clear that
one obtains the arrow induced by
\[
 nQ[0]\longrightarrow[Q\xrightarrow{n}Q].
\]
This proves \SGARefLink{sga7.sga7i.eviii.n2-3-11}{2.3.11}.

% END INCLUDED SGA7I_English_current_buildable_source_20260823/components/103_expose_VIII_alternative_proof_2_3_19.tex

% BEGIN INCLUDED SGA7I_English_current_buildable_source_20260823/components/104_expose_VIII_section_2_4.tex
% SOURCE AUTHORITY:
% expose_VIII_body.tex lines 1346--1380; scan index 264; source folio 253.
% Direct scan comparison: physical PDF page 265 in four overlapping
% 1100-dpi bands.
% Hard stop before section 3, source line 1386 / scan index 265.

\medskip
\noindent \SGARefTarget{sga7.sga7i.eviii.n2-4}{2.4}.  Let
\[
 f:\underline T'\longrightarrow\underline T
\]
be a morphism of topoi.  There is then a canonical homomorphism
(VII \SGARefLink{sga7.sga7i.evii.n2-8}{2.8})
\[
 \Biext^1(P,Q;G)\longrightarrow\Biext^1(P',Q';G'),
\]
where primes denote inverse-image groups.  Since, for every
commutative group $E$ of $\underline T$, there is plainly a canonical
isomorphism
\[
 ({}_nE)'\simeq{}_n(E'),
\]
the homomorphisms (\SGARefLink{sga7.sga7i.eviii.n2-2-2}{2.2.2}) over $\underline T$ and $\underline T'$ give
a canonical diagram
\begin{equation}\tag{\SGARefTarget{sga7.sga7i.eviii.n2-4-1}{2.4.1}}
\begin{tikzcd}[column sep=huge,row sep=huge]
 \Biext^1(P,Q;G)\arrow[r]\arrow[d,"{\varphi^n_{P,Q}}"']
   &\Biext^1(P',Q';G')\arrow[d,"{\varphi^n_{P',Q'}}"]\\
 \Hom({}_nP\otimes{}_nQ,G)\arrow[r]
   &\Hom({}_nP'\otimes{}_nQ',G').
\end{tikzcd}
\end{equation}
Similarly, the homomorphisms (\SGARefLink{sga7.sga7i.eviii.n2-2-3}{2.2.3}) give an analogous diagram in
terms of the $T_\ell$.  I claim that these diagrams are
\underline{commutative}.  The verification is essentially immediate
from the definitions and is left to the reader.  More generally,
notice that \underline{all} the canonical homomorphisms considered in
this section and the preceding one are \underline{compatible with
change of topos}, in the evident sense.  We shall henceforth use
compatibilities of this kind without further mention.


% END INCLUDED SGA7I_English_current_buildable_source_20260823/components/104_expose_VIII_section_2_4.tex

% BEGIN INCLUDED SGA7I_English_current_buildable_source_20260823/components/105_expose_VIII_sections_3_1_and_3_2.tex
% SOURCE AUTHORITY:
% expose_VIII_body.tex lines 1386--1427; scan indices 265--266;
% source folios 254--255.
% Direct scan comparison: physical PDF pages 266--267 in four overlapping
% 1100-dpi bands per page.
% Hard stop before section 3.3, source line 1432 / scan index 267.

\subsection*{\SGARefTarget{sga7.sga7i.eviii.n3}{3}. Biextensions of group schemes $(P,Q)$ by
$\underline G_m$: generalities}

Throughout this section we fix a base scheme $S$ and, as announced in
Section 0, work in the fppf topos of $S$.

\medskip
\noindent \SGARefTarget{sga7.sga7i.eviii.n3-1}{3.1}.  Let us first recall the most important example of a
group $\Biext^1(P,Q;\underline G_m)$: the case where $P$ and $Q$ are
two abelian schemes over $S$, and where $\Biext^1$ is interpreted as
the group of divisorial correspondences on $(P,Q)$ (VII \SGARefLink{sga7.sga7i.evii.n2-9-4}{2.9.4}).  The
$\ell$-adic pairings of the preceding section are then the familiar
pairings [1] in the theory of abelian schemes.

\medskip
\noindent \SGARefTarget{sga7.sga7i.eviii.n3-2}{3.2}.  More generally, suppose only that $Q$ is an abelian
scheme.  It is well known (VII \SGARefLink{sga7.sga7i.evii.n1-3-8}{1.3.8}) that
\begin{equation}\tag{\SGARefTarget{sga7.sga7i.eviii.n3-2-1}{3.2.1}}
 \uHom(Q,\underline G_m)=0,
\end{equation}
and hence (\SGARefLink{sga7.sga7i.eviii.n1-4}{1.4}) there is a canonical isomorphism
\begin{equation}\tag{\SGARefTarget{sga7.sga7i.eviii.n3-2-2}{3.2.2}}
 \Biext^1(P,Q;\underline G_m)\simeq\Hom(P,Q^*),
\end{equation}
where
\begin{equation}\tag{\SGARefTarget{sga7.sga7i.eviii.n3-2-3}{3.2.3}}
 Q^*=\uExt^1(Q,\underline G_m)
\end{equation}
is the dual abelian scheme of $Q$.  (For the representability of
$Q^*$, due to M. Raynaud in the general case, see [1].)

When $P$ is also an abelian scheme, symmetrically to (\SGARefLink{sga7.sga7i.eviii.n3-2-2}{3.2.2}), there is
a canonical isomorphism
\begin{equation}\tag{\SGARefTarget{sga7.sga7i.eviii.n3-2-4}{3.2.4}}
 \Biext^1(P,Q;\underline G_m)\simeq\Hom(Q,P^*).
\end{equation}
Thus the datum of a biextension $\xi$ of $(P,Q)$ by $\underline G_m$
may be interpreted equivalently as the datum of a homomorphism
\[
 u:P\longrightarrow Q^*
\]
or of a homomorphism
\[
 v:Q\longrightarrow P^*.
\]

Recall that a homomorphism of abelian schemes $w:A\longrightarrow B$
is an isomorphism if and only if, for every integer $n>0$, the induced
homomorphism ${}_nw:{}_nA\longrightarrow{}_nB$ is an isomorphism.  In
the case of $u$ and $v$ above, ${}_nu:{}_nP\longrightarrow{}_nQ^*$ is
identified with the homomorphism
${}_nP\longrightarrow D({}_nQ)$ defined by the pairing
\[
 \varphi_\xi^n:{}_nP\times{}_nQ\longrightarrow\underline G_m
\]
(see \SGARefLink{sga7.sga7i.eviii.n2-3-11}{2.3.11}), whereas ${}_nv:{}_nQ\longrightarrow{}_nP^*$ is
identified with the homomorphism
${}_nQ\longrightarrow D({}_nP)$ defined by the pairing
\[
 \psi_\xi^n:{}_nP\times{}_nQ\longrightarrow\underline G_m.
\]
Now $\varphi_\xi^n=-\psi_\xi^n$ (\SGARefLink{sga7.sga7i.eviii.n2-1-10}{2.1.10}), so, up to sign,
${}_nu$ and ${}_nv$ are transposes of one another.  Since ${}_nP$ and
${}_nQ$ are finite and locally free over $S$, and therefore satisfy
Cartier duality, it follows that ${}_nu$ is an isomorphism if and only
if ${}_nv$ is an isomorphism.  Consequently $u$ is an isomorphism if
and only if $v$ is an isomorphism.

An equivalent, though slightly less symmetric, formulation is the
following.  Consider the canonical biextension (the Weil
biextension) of $(P,P^*)$ by $\underline G_m$, corresponding to the
identity homomorphism $v:P^*\longrightarrow P^*$.  Then the
corresponding canonical homomorphism
\[
 u:P\longrightarrow P^{**}
\]
is an \underline{isomorphism}.  This is the \underline{biduality
theorem} of Barsotti and Cartier.


% END INCLUDED SGA7I_English_current_buildable_source_20260823/components/105_expose_VIII_sections_3_1_and_3_2.tex

% BEGIN INCLUDED SGA7I_English_current_buildable_source_20260823/components/106_expose_VIII_section_3_3.tex
% SOURCE AUTHORITY:
% expose_VIII_body.tex lines 1432--1481; scan indices 267--270;
% source folios 256--259.
% Direct scan comparison: physical PDF pages 268--271 in four overlapping
% 1100-dpi bands per page.
% Transparent source dispositions:
% - dualizing 0 -> T -> M -> L -> 0 gives the type-correct terminal D(T),
%   not the repeated D(L) in the working transcription;
% - the visibly blank citation after "SGA 3 X 4.9 et" is not guessed; the
%   dangling conjunction is omitted in English.
% Hard stop before Proposition 3.4, source line 1483 / scan index 270.

\medskip
\noindent \SGARefTarget{sga7.sga7i.eviii.n3-3}{3.3}.  We now turn to the case where $Q$ is finite locally
free over $S$, or is a group scheme of multiplicative type over $S$.
The key lemma in this case is the following well-known result.

\medskip
\noindent\underline{Proposition \SGARefTarget{sga7.sga7i.eviii.n3-3-1}{3.3.1}.}  \underline{Let $Q$ be a
commutative group scheme over $S$, and suppose that $Q$ is finite
locally free over $S$, or that $Q$ is of finite type and of
multiplicative type over $S$ (SGA 3 IX).  Then}
\[
 \uExt^1(Q,\underline G_m)=0.
\]

\noindent\underline{Proof.}  We distinguish the two cases.

\medskip
\noindent a) Suppose that $Q$ is finite locally free over $S$.  We
must prove that every extension $E$ of $Q$ by $\underline G_m$ splits
locally for the fppf topology.  We may plainly suppose $S$ affine;
there is then an integer $n>0$ such that ${}_nQ=0$ (SGA 3 IX 6.1).
Moreover, $n\cdot\id$ is an epimorphism of $\underline G_m$ for the
fppf topology, and hence the Kummer exact sequence
\[
 0\longrightarrow\mu_n\longrightarrow\underline G_m
 \xrightarrow{\ n\ }\underline G_m\longrightarrow0
\]
gives an exact sequence of the $\Ext^i(Q,-)$ and, in particular, an
\underline{epimorphism}
\[
 \Ext^1(Q,\mu_n)\longrightarrow
 {}_n\Ext^1(Q,\underline G_m)=\Ext^1(Q,\underline G_m).
\]
The last equality follows because multiplication by $n$, being zero
on $Q$, is also zero on $\Ext^1(Q,\underline G_m)$.  Thus the
extension $E$ comes from an extension $F$ of $Q$ by $\mu_n$.

The group $F$ is represented by a finite locally free scheme over $S$
(fpqc descent theory, SGA 1 VII).  By the known properties of Cartier
duality for finite locally free groups (SGA 3 VII 3),
\[
 D(F)=\uHom(F,\underline G_m)
\]
is an extension of $D(\mu_n)$ by $D(Q)$.  Here one uses the following
fact: if $A\longrightarrow B$ is a monomorphism of commutative finite
locally free groups over $S$, then
$D(B)\longrightarrow D(A)$ is faithfully flat.  Since $D(B)$ and
$D(A)$ are themselves finite locally free over $S$, the fibrewise
criterion for faithful flatness reduces us to the case where $S$ is
the spectrum of a field $k$.  In that case the assertion says that
the homomorphism of affine coordinate rings
\[
 k(D(A))\longrightarrow k(D(B))
\]
is injective (SGA 3 $\mathrm{VI}_B$ 11.14).  This homomorphism is the
transpose of $k(B)\longrightarrow k(A)$, which is surjective because
$A\longrightarrow B$ is a closed immersion (SGA 3 $\mathrm{VI}_B$
1.4.2).  Hence $k(D(A))\longrightarrow k(D(B))$ is indeed injective.

The inclusion $\mu_n\longrightarrow\underline G_m$ identifies with a
section of
\[
 D(\mu_n)=\uHom(\mu_n,\underline G_m),
\]
and it remains only to prove that, locally for the fppf topology on
$S$, this inclusion extends to a homomorphism
$F\longrightarrow\underline G_m$; equivalently, that the section of
$D(\mu_n)$ just considered lifts fppf-locally to a section of $D(F)$.
This is clear, since $D(F)\longrightarrow D(\mu_n)$ is faithfully flat
and of finite presentation.

\medskip
\noindent b) By SGA 3 X 4.5, after replacing $S$ by an \`etale cover,
we may suppose that $Q$ is diagonalizable, say $Q=D(M)$, where $M$ is
a finitely generated $\mathbb Z$-module.  Let $T$ be the torsion
submodule of $M$, and put $L=M/T$, so that $L\simeq\mathbb Z^r$.  The
exact sequence
\[
 0\longrightarrow T\longrightarrow M\longrightarrow L
 \longrightarrow0
\]
then gives an exact sequence (SGA 3 IX 3.1)
\[
 0\longrightarrow D(L)\longrightarrow D(M)\longrightarrow D(T)
 \longrightarrow0.
\]
It is therefore enough to prove separately that the $\uExt^1$ of
$D(L)$ and of $D(T)$ by $\underline G_m$ vanish.  The case of $D(T)$
follows from a), and the case of $D(L)$ reduces to $L=\mathbb Z$.
Thus it remains to check that
\[
 \uExt^1(\underline G_m,\underline G_m)=0,
\]
that is, that every extension $E$ of $\underline G_m$ by
$\underline G_m$ is fppf-locally trivial.

Such an extension is represented by an affine flat group scheme over
$S$ (fpqc descent theory, SGA 1 VII) whose fibres are tori; it is
therefore a torus (SGA 3 X 4.9).  By SGA 3 X 4.5 we may suppose that
$E$ is a split torus
\[
 E=D(R)\simeq\underline G_m^{,2}.
\]
After a Zariski localization, we may also suppose that the
monomorphism $\underline G_m\hookrightarrow E=D(R)$ comes from a
constant morphism $R\longrightarrow\mathbb Z_S$ (SGA 3 VIII 1.4),
which is necessarily an epimorphism (SGA 3 VIII 3.2 b)).  Choosing
$r\in R$ above $1\in\mathbb Z$ gives a homomorphism
$E\longrightarrow\underline G_m$ splitting the extension, and proves
the assertion.

\medskip
\noindent\underline{Corollary \SGARefTarget{sga7.sga7i.eviii.n3-3-2}{3.3.2}.}  \underline{Let $Q$ be as in
\SGARefLink{sga7.sga7i.eviii.n3-3-1}{3.3.1}.  Then, for every $P$, there is a canonical isomorphism}
\begin{equation}\tag{\SGARefTarget{sga7.sga7i.eviii.n3-3-2-1}{3.3.2.1}}
 \Biext^1(P,Q;\underline G_m)
 \xleftarrow{\ \sim\ }\Ext^1(P,D(Q)),
\end{equation}
\underline{where $D(Q)$ denotes the Cartier dual
$\uHom(Q,\underline G_m)$ of $Q$.  More precisely, the canonical
functor (\SGARefLink{sga7.sga7i.eviii.n1-1-6}{1.1.6}) is an equivalence from the category of extensions of
$P$ by $D(Q)$ to the category of biextensions of $(P,Q)$ by
$\underline G_m$.}

This is the special case of \SGARefLink{sga7.sga7i.eviii.n1-5}{1.5} obtained by applying \SGARefLink{sga7.sga7i.eviii.n3-3-1}{3.3.1}.

\medskip
\noindent\underline{Remarks \SGARefTarget{sga7.sga7i.eviii.n3-3-3}{3.3.3}.}  Under the hypotheses of \SGARefLink{sga7.sga7i.eviii.n3-3-1}{3.3.1},
$D(Q)$ is represented by a separated group scheme, locally of finite
presentation and locally quasi-finite over $S$.  More precisely, if
$Q$ is finite locally free over $S$, then so is $D(Q)$, as is well
known (SGA 3 VII 3).  If $Q$ is of multiplicative type, then $D(Q)$
is a \underline{finitely generated twisted constant group} (SGA 3 X
5.1), as was seen in SGA 3 IX 4.5 and X 5.7.  If $P$ is a group
scheme, every extension of $P$ by $D(Q)$ is therefore representable,
by a familiar lemma from fpqc descent theory (SGA 3 X 5.4).

\medskip
\noindent \SGARefTarget{sga7.sga7i.eviii.n3-3-4}{3.3.4}.  It is false that
$\uExt^1(Q,\underline G_m)=0$ for every affine flat group scheme $Q$
of finite presentation over $S$.  This already fails for
\[
 Q=\underline G_{a_S},\qquad
 S=\Spec k[t]/(t^2),
\]
where $k$ is a field of characteristic $p>0$.


% END INCLUDED SGA7I_English_current_buildable_source_20260823/components/106_expose_VIII_section_3_3.tex

% BEGIN INCLUDED SGA7I_English_current_buildable_source_20260823/components/107_expose_VIII_section_3_4.tex
% SOURCE AUTHORITY:
% expose_VIII_body.tex lines 1483--1508; scan indices 270--273;
% source folios 259--262.
% Direct scan comparison: physical PDF pages 271--274 in four overlapping
% 1100-dpi bands per page.  The pages are clear at this resolution and no
% higher-resolution ambiguity pass was required.
% Transparent source disposition:
% - in the proof of Proposition 3.4, the printed phrase "une section f de P
%   sur P" is rendered as the mathematically required section of E over P;
%   the following displayed cocycle condition and Lemma 3.4.2 make the
%   intended object unambiguous.
% Hard stop before Corollary 3.5, source line 1510 / scan index 273.

\medskip
\noindent\underline{Proposition \SGARefTarget{sga7.sga7i.eviii.n3-4}{3.4}.}  \underline{Let $P$ be a smooth
group scheme of finite presentation over $S$, with connected fibres,
and let $Q$ be a torus over $S$.  Then the category of biextensions
of $(P,Q)$ by $\underline G_m$ is equivalent to the one-point
category; that is,}
\begin{equation}\tag{\SGARefTarget{sga7.sga7i.eviii.n3-4-1}{3.4.1}}
 \Biext^0(P,Q;\underline G_m)=
 \Biext^1(P,Q;\underline G_m)=0.
\end{equation}

\noindent\underline{Proof.}  Put $M=D(Q)$.  This is a twisted
constant group over $S$ whose fibres are isomorphic to groups
$\mathbb Z^r$ (\SGARefLink{sga7.sga7i.eviii.n3-3-3}{3.3.3}).  By \SGARefLink{sga7.sga7i.eviii.n3-3-2}{3.3.2}, the category in question is
equivalent to the category of extensions of $P$ by $M$.  This
category is \underline{rigid}: the automorphism group of any object
is the trivial group.  This is the first equality in (\SGARefLink{sga7.sga7i.eviii.n3-4-1}{3.4.1}), which
may also be written
\[
 \Hom(P,M)=0.
\]
It follows at once from the fact that $M$ is etale over $S$, whereas
$P$ has connected fibres.  It remains to prove
\[
 \Ext^1(P,M)=0,
\]
that is, that every extension $E$ of $P$ by $M$ is trivial.

By the uniqueness of a splitting when one exists, it is enough to
prove that $E$ is locally trivial.  We may therefore reduce to the
case $M=\mathbb Z_S^r$, and then to $M=\mathbb Z_S$.  When $S$ is the
spectrum of a field, the conclusion follows from \SGARefLink{sga7.sga7i.eviii.n5-5}{5.5} (i) below.
(Notice that we shall not use \SGARefLink{sga7.sga7i.eviii.n3-4}{3.4} before IX 6, so in particular there
is no vicious circle.)  We now reduce the general case to this
special case.

It is enough to find a section $f$ of $E$ over $P$ such that
$f(0)=0$, for such a section is necessarily multiplicative.  Indeed,
\[
 f(x+y)-f(x)-f(y)=0:
\]
the left-hand side defines a morphism
$P\times_S P\longrightarrow\mathbb Z_S$ which vanishes on the zero
section of $P\times_S P$.  It therefore vanishes identically, since
$P\times_S P$ has connected fibres (the geometric fibres of $P$
being connected).  The required conclusion now follows from the
following lemma.

\medskip
\noindent\underline{Lemma \SGARefTarget{sga7.sga7i.eviii.n3-4-2}{3.4.2}.}  \underline{Let $S$ be a scheme,
$P$ a smooth $S$-scheme of finite presentation with geometrically
connected fibres, $Z$ a group, and $E$ a $Z_P$-torsor (for the fppf,
etale, or fpqc topology; these notions are equivalent here, cf. SGA
3 X 5.4).  Let $e$ be a section of $P$ over $S$, and let $e'$ be a
section of $E$ over $S$ above $e$ (that is, a section of $e^{-1}(E)$
over $S$).}

\underline{The torsor $E$ is trivial if and only if, for every
$s\in S$, the induced torsor $E_s$ over $P_s$ is trivial.  In that
case there is a unique section $f$ of $E$ over $P$ such that
$f\circ e=e'$.}

\noindent\underline{Proof.}  Uniqueness is clear.  Two such sections
$f$ and $g$ ``differ'' by a morphism
\[
 h:P\longrightarrow Z_S,
 \qquad h\circ e=1,
\]
and such a morphism is constantly equal to $1$, since $P$ has
connected fibres.  Moreover, if a section $f$ of $E$ over $P$
exists, it can always be corrected by a morphism
$P\longrightarrow Z_S$ coming from a morphism $z:S\longrightarrow
Z_S$, so as to obtain a section $f'$ satisfying $f'\circ e=e'$.
One need only choose $z$ so that
\[
 f(e(t))=e'(t)z(t).
\]

Applying this observation fibrewise, for every $s\in S$ we obtain a
well-determined section $f_s$ of $E_s$ over $P_s$ taking $e_s$ to
$e'_s$.  It remains to prove that the $f_s$ arise from a section $f$
of $E$ over $P$.  Equivalently (EGA IV 17.9.3), we must prove that
\[
 U=\bigcup_{s\in S}f_s(P_s)\subset E
\]
is open in $E$.

Suppose first that $S$ is noetherian.  Then $U$ is constructible.  To
see this, for each $s\in S$ let $T$ be the closure of $s$, endowed
with its reduced induced structure.  There is an open neighbourhood
$T'$ of $s$ in $T$ and a section of $E_{T'}$ over $P_{T'}$ extending
$f_s$ and compatible with the marked sections $e$ and $e'$ (as
follows directly from EGA IV 8.8.2).  By the criterion of EGA IV
1.10.1, it remains to show that $U$ is stable under generization.
The usual argument using EGA II 7.1.9 reduces this to the case in
which $S$ is the spectrum of a discrete valuation ring.

Since $P$ is smooth over $S$, it is regular and hence normal; and
since it has connected fibres, it is connected.  The given section
of $E$ over the generic fibre of $P$ over $S$ therefore extends
uniquely to a section $f$ of $E$ over $P$ (EGA IV 18.10.19 and
18.10.20), necessarily compatible with $e$ and $e'$.  Thus in this
case $U=f(P)$ is open.  This proves \SGARefLink{sga7.sga7i.eviii.n3-4}{3.4} when $S$ is noetherian.  The
general case reduces readily to this one by the usual limit
arguments of EGA IV, 8 and 9; we leave the details to the reader.

\medskip
\noindent\underline{Remark \SGARefTarget{sga7.sga7i.eviii.n3-4-3}{3.4.3}.}  Lemma 3.4.2 can be ``crushed''
under learned references, namely SGA 4 XV 1.15 (the case $n=1$) and
\SGARefLink{sga7.sga7i.eviii.n4-1}{4.1}.


% END INCLUDED SGA7I_English_current_buildable_source_20260823/components/107_expose_VIII_section_3_4.tex

% BEGIN INCLUDED SGA7I_English_current_buildable_source_20260823/components/108_expose_VIII_sections_3_5_through_3_7.tex
% SOURCE AUTHORITY:
% expose_VIII_body.tex lines 1510--1553; scan indices 273--275;
% source folios 262--264.
% Direct scan comparison: physical PDF pages 274--276, with pages 275--276
% inspected in four overlapping 1100-dpi bands.  The typography is clear at
% that resolution; no 5000--9000-dpi ambiguity escalation was needed.
% Transparent source dispositions:
% - formula (3.6.2) has Biext^i on both sides.  The printed right-hand
%   exponent 1 is incompatible with the following qualifier i=0,1 and with
%   the asserted equivalence of categories;
% - in the proof of 3.7 the second cyclic case is M=Z/nZ, not the printed
%   phrase "Z=Z/nZ".
% Hard stop before section 4, source line 1559 / scan index 276.

\medskip
\noindent\underline{Corollary \SGARefTarget{sga7.sga7i.eviii.n3-5}{3.5}.}  \underline{Let $P$ be a smooth
group scheme of finite presentation over $S$, with connected fibres,
and let}
\[
 0\longrightarrow Q'\longrightarrow Q\longrightarrow Q''
 \longrightarrow0
\]
\underline{be an exact sequence of commutative groups over $S$, with
$Q'$ a torus.  Then the inverse-image functor on biextensions}
\begin{equation}\tag{\SGARefTarget{sga7.sga7i.eviii.n3-5-1}{3.5.1}}
 \operatorname{BIEXT}(P,Q'';\underline G_m)
 \xrightarrow{\ \sim\ }
 \operatorname{BIEXT}(P,Q;\underline G_m)
\end{equation}
\underline{is an equivalence of categories.  In particular, the
following natural homomorphisms are isomorphisms:}
\begin{equation}\tag{\SGARefTarget{sga7.sga7i.eviii.n3-5-2}{3.5.2}}
 \Biext^i(P,Q'';\underline G_m)
 \xrightarrow{\ \sim\ }
 \Biext^i(P,Q;\underline G_m),
 \qquad i=0,1.
\end{equation}

Indeed, up to equivalence, $\operatorname{BIEXT}(P,Q'';\underline
G_m)$ identifies with the category of biextensions $E$ of $(P,Q)$ by
$\underline G_m$ \underline{equipped} with a trivialization of the
induced biextension $E'$ of $(P,Q')$ by $\underline G_m$ (VII
\SGARefLink{sga7.sga7i.evii.n3-7-6}{3.7.6}).  By \SGARefLink{sga7.sga7i.eviii.n3-4}{3.4} there is exactly one such trivialization of $E'$ for
each given biextension $E$.  One can of course also prove directly
that (\SGARefLink{sga7.sga7i.eviii.n3-5-2}{3.5.2}) is an isomorphism for $i=0,1$ - which is equivalent to
the conclusion of \SGARefLink{sga7.sga7i.eviii.n3-5}{3.5} - by using the exact sequence of the
$\Biext^i$ in VII \SGARefLink{sga7.sga7i.evii.n3-7-5}{3.7.5} associated with
\[
 0\longrightarrow Q'\longrightarrow Q\longrightarrow Q''
 \longrightarrow0.
\]

\medskip
\noindent\underline{Remark \SGARefTarget{sga7.sga7i.eviii.n3-6}{3.6}.}  The symmetric statement obtained
from \SGARefLink{sga7.sga7i.eviii.n3-5}{3.5} by interchanging left and right is of course also valid; it
also follows formally from \SGARefLink{sga7.sga7i.eviii.n3-5}{3.5} by passing to symmetric biextensions
(VII \SGARefLink{sga7.sga7i.evii.n2-7}{2.7}).  Suppose, in addition to the data and hypotheses of \SGARefLink{sga7.sga7i.eviii.n3-5}{3.5},
that the symmetric data and hypotheses are given: namely, that $Q$
has connected fibres (equivalently, that $Q''$ has connected fibres,
since $Q'$ has connected fibres), and that there is an exact sequence
\[
 0\longrightarrow P'\longrightarrow P\longrightarrow P''
 \longrightarrow0
\]
with $P'$ a torus.  Successive application of \SGARefLink{sga7.sga7i.eviii.n3-5}{3.5} and of its
symmetric form, applied to this last sequence and to $Q''$, shows
that the functor
\begin{equation}\tag{\SGARefTarget{sga7.sga7i.eviii.n3-6-1}{3.6.1}}
 \operatorname{BIEXT}(P'',Q'';\underline G_m)
 \xrightarrow{\ \sim\ }
 \operatorname{BIEXT}(P,Q;\underline G_m)
\end{equation}
is an equivalence of categories.  In particular, the following
homomorphism is an isomorphism:
\begin{equation}\tag{\SGARefTarget{sga7.sga7i.eviii.n3-6-2}{3.6.2}}
 \Biext^i(P'',Q'';\underline G_m)
 \xrightarrow{\ \sim\ }
 \Biext^i(P,Q;\underline G_m),
 \qquad i=0,1.
\end{equation}

\medskip
\noindent\underline{Proposition \SGARefTarget{sga7.sga7i.eviii.n3-7}{3.7}.}  \underline{Let $A$ be an
abelian scheme over $S$, let $T$ be a group scheme of multiplicative
type (SGA 3 IX) and of finite type over $S$, and let}
\[
 \underline M=\uHom(T,\underline G_m)
\]
\underline{be the Cartier dual of $T$ (a finitely generated twisted
constant group over $S$, by \SGARefLink{sga7.sga7i.eviii.n3-3-3}{3.3.3}).  Then the categories
$\operatorname{EXT}(A,T)$ and
$\operatorname{BIEXT}(A,\underline M;\underline G_m)$, formed
respectively by commutative extensions of $A$ by $T$ and by
biextensions of $(A,\underline M)$ by $\underline G_m$, are rigid and
canonically equivalent.  Their groups of isomorphism classes are
canonically isomorphic to $\Hom(\underline M,A^*)$, where}
\[
 A^*=\uExt^1(A,\underline G_m)
\]
\underline{is the dual abelian scheme of $A$ (\SGARefLink{sga7.sga7i.eviii.n3-2}{3.2}).  In particular,
there is a canonical isomorphism}
\begin{equation}\tag{\SGARefTarget{sga7.sga7i.eviii.n3-7-1}{3.7.1}}
 \Ext^1(A,T)\xrightarrow{\ \sim\ }
 \Hom(\underline M,A^*).
\end{equation}

This is the special case of \SGARefLink{sga7.sga7i.eviii.n1-6}{1.6} obtained by taking $P=A$ and
$Q=\underline M$.  By the biduality theorem (SGA 3 X 5.7), we do
indeed have
\[
 T\simeq\uHom(\underline M,\underline G_m),
\]
and it remains only to verify the conditions (\SGARefLink{sga7.sga7i.eviii.n1-6-1}{1.6.1}), which here read
\[
 \uHom(A,\underline G_m)=0,
 \qquad
 \Hom\bigl(A,\uExt^1(\underline M,\underline G_m)\bigr)=0.
\]
The first equality was recalled in (\SGARefLink{sga7.sga7i.eviii.n3-2-1}{3.2.1}); the second may be
strengthened to
\[
 \uExt^1(\underline M,\underline G_m)=0.
\]
This assertion is local on $S$, so we may suppose that $\underline
M$ is constant, $\underline M=M_S$.  Decomposing $M$ as a direct sum
of cyclic $\mathbb Z$-modules reduces us to $M=\mathbb Z$, which is
trivial, and to $M=\mathbb Z/n\mathbb Z$, which follows from the
$n$-divisibility of $\underline G_m$ (or, more learnedly, from
\SGARefLink{sga7.sga7i.eviii.n3-3-1}{3.3.1}).

When $T=\underline G_m$, and hence $\underline M=\mathbb Z_S$,
(\SGARefLink{sga7.sga7i.eviii.n3-7-1}{3.7.1}) reduces to the identity isomorphism (\SGARefLink{sga7.sga7i.eviii.n3-2-3}{3.2.3}).  Of course, a
direct proof of (\SGARefLink{sga7.sga7i.eviii.n3-7-1}{3.7.1}) would be essentially trivial, without
introducing the group
\[
 \Ext^1\bigl(A\mathbin{\mathop{\otimes}^{\mathbb L}}\underline M,
              \underline G_m\bigr)
 =\Biext^1(A,\underline M;\underline G_m).
\]


% END INCLUDED SGA7I_English_current_buildable_source_20260823/components/108_expose_VIII_sections_3_5_through_3_7.tex

% BEGIN INCLUDED SGA7I_English_current_buildable_source_20260823/components/109_expose_VIII_sections_4_1_through_4_4.tex
% SOURCE AUTHORITY:
% expose_VIII_body.tex lines 1559--1575; scan indices 276--277;
% source folios 265--266.
% Direct scan comparison: physical PDF pages 277--278 in four overlapping
% 1100-dpi bands.  A targeted crop of the formula in Proposition 4.2 confirms
% that the printed argument is P in all three pullbacks, although the sentence
% has just introduced the line bundle L.  The mathematically type-correct L is
% therefore used below and the product is written explicitly as a tensor
% product.  No other high-resolution ambiguity remained.
% Hard stop before Lemma 4.5, source line 1581 / scan index 278.

\subsection*{\SGARefTarget{sga7.sga7i.eviii.n4}{4}. Extensions and biextensions of smooth connected group
schemes over a field}

We recall the following lemma due to Rosenlicht [3 bis, Proposition
3].

\medskip
\noindent\underline{Lemma \SGARefTarget{sga7.sga7i.eviii.n4-1}{4.1}.}  \underline{Let $k$ be a field, and
let $X$ and $Y$ be two separable, geometrically connected $k$-schemes
of finite type, equipped respectively with $k$-rational points $e_X$
and $e_Y$.  Then every morphism}
\[
 X\times Y\longrightarrow\underline G_m
\]
\underline{which is trivial on the two subschemes $X\times e_Y$ and
$e_X\times Y$ is trivial.}

We shall use this result only when $X$ and $Y$ are group schemes over
$k$; in that case the hypotheses imply that $X$ and $Y$ are in fact
smooth over $k$.  By induction, \SGARefLink{sga7.sga7i.eviii.n4-1}{4.1} also implies the analogous result
for any number of factors.  We may therefore apply the results of VII
\SGARefLink{sga7.sga7i.evii.n1-3}{1.3} and VII \SGARefLink{sga7.sga7i.evii.n2-9}{2.9}, and obtain the following two propositions.

\medskip
\noindent\underline{Proposition \SGARefTarget{sga7.sga7i.eviii.n4-2}{4.2}.}  \underline{Let $P$ be a
smooth connected group over the field $k$.  The natural functor from
the category of group extensions of $P$ by $\underline G_m$ to the
category of line bundles on $P$ rigidified at the identity element
$e_P$ is fully faithful.  Its essential image consists of the
rigidified line bundles $\mathbb L$ on $P$ for which}
\[
 \pi^*\mathbb L\simeq
 \pr_1^*\mathbb L\otimes\pr_2^*\mathbb L,
\]
\underline{where $\pi:P\times P\longrightarrow P$ is the group law
of $P$, and $\pr_1,\pr_2$ are the two projections.}

\medskip
\noindent\underline{Proposition \SGARefTarget{sga7.sga7i.eviii.n4-3}{4.3}.}  \underline{Let $P$ and $Q$ be
two smooth connected algebraic groups over the field $k$.  The
natural functor from the category of biextensions of $(P,Q)$ by
$\underline G_m$ to the category of line bundles on $P\times Q$
birigidified with respect to the identity sections $e_P$ and $e_Q$ of
$P$ and $Q$ is fully faithful; moreover, both categories are rigid.
A birigidified line bundle $\mathbb L$ on $P\times Q$ comes from a
biextension if and only if its pullbacks to $P\times P\times Q$ and
$P\times Q\times Q$ by $\pi_P\times\id_Q$ and
$\id_P\times\pi_Q$ satisfy the ``theorem of the cube'' ([1] or
[2]).}

This last condition is simply the translation of the conditions in
VII \SGARefLink{sga7.sga7i.evii.n2-9-3}{2.9.3} b), in the explicit form given in VII \SGARefLink{sga7.sga7i.evii.n3-5}{3.5} b).

\medskip
\noindent \SGARefTarget{sga7.sga7i.eviii.n4-4}{4.4}.  I do not know whether the cube condition in \SGARefLink{sga7.sga7i.eviii.n4-3}{4.3} is
automatically satisfied.  At least it is known [2, IV 2.6] that, for
a given $\mathbb L$, there is always an integer $n\geq1$, a power of
the characteristic exponent of $k$, such that
$\mathbb L^{\otimes n}$ satisfies that condition and hence comes from
a biextension of $(P,Q)$ by $\underline G_m$.  It is also always
possible to find a purely inseparable extension $k'$ of $k$ such that
the required condition becomes true after the base extension
$k\longrightarrow k'$; after this extension, therefore, $\mathbb L$
is a birigidified line bundle arising from a biextension of $(P,Q)$
by $\underline G_m$.

In particular, when $k$ is perfect, the condition in \SGARefLink{sga7.sga7i.eviii.n4-3}{4.3} is always
satisfied.  This last fact, which formally implies the preceding
statements by a trace argument, can also be seen by applying \SGARefLink{sga7.sga7i.eviii.n4-7}{4.7}
below; cf. \SGARefLink{sga7.sga7i.eviii.n4-9}{4.9}.


% END INCLUDED SGA7I_English_current_buildable_source_20260823/components/109_expose_VIII_sections_4_1_through_4_4.tex

% BEGIN INCLUDED SGA7I_English_current_buildable_source_20260823/components/110_expose_VIII_sections_4_5_and_4_6.tex
% SOURCE AUTHORITY:
% expose_VIII_body.tex lines 1581--1605; scan indices 278--280;
% source folios 267--269.
% Direct scan comparison: physical PDF pages 279--281 in four overlapping
% 1100-dpi bands per page.  Prose and formulas are fully legible at that
% resolution.
% Transparent source dispositions in the second case of Proposition 4.6:
% - the descent base is Y=P x B, not the printed Y=P x A;
% - after reducing to the case in which Q is abelian, the rigidifying
%   restriction is P x e_Q, not the printed P x e_A.
% Both readings are forced by the displayed morphism to Pic_{Q/k} and by the
% immediately preceding extension 1 -> M -> Q -> B -> 1.
% Hard stop before Proposition 4.7, source line 1607 / scan index 280.

We recall, in the form of a lemma, the following result from EGA
$\mathrm{Err}_{\mathrm{IV}}$ 53, 21.4.13.

\medskip
\noindent\underline{Lemma \SGARefTarget{sga7.sga7i.eviii.n4-5}{4.5}.}  \underline{Let
$f:X\longrightarrow Y$ be a smooth surjective morphism with integral
fibres, where $Y$ is normal and integral with generic point $y$.  Let
$\mathbb L$ be a line bundle on $X$ whose restriction to $X_y$ is
trivial.  Then there is a line bundle $\mathbb M$ on $Y$ such that}
\[
 \mathbb L\simeq f^*\mathbb M.
\]

\medskip
\noindent\underline{Proposition \SGARefTarget{sga7.sga7i.eviii.n4-6}{4.6}.}  \underline{Let $P,Q$ be two
smooth connected algebraic groups over the field $k$, with $P$
affine.  Suppose in addition either that $P$ admits a composition
series each of whose factors is a torus or is isomorphic to the
additive group $\underline G_a$ (a condition satisfied in particular
when $k$ is perfect [SGA 3, XVII 7.2.1, 4.1.1 (i)$\Longleftrightarrow$
(ii), and 4.1.5]), or that $Q$ is an extension of an abelian variety
$B$ by a smooth connected affine algebraic group satisfying the
preceding condition.  Under either hypothesis, every birigidified
line bundle $\mathbb L$ on $P\times Q$ is trivial; a fortiori, by
\SGARefLink{sga7.sga7i.eviii.n4-3}{4.3}, every biextension of $(P,Q)$ by $\underline G_m$ is trivial.}

\noindent\underline{Proof.}  First observe that every automorphism of
a birigidified line bundle is the identity (\SGARefLink{sga7.sga7i.eviii.n4-3}{4.3}).  Descent theory
therefore shows that it is enough to prove that $\mathbb L$ becomes
trivial after a finite separable extension of $k$.  Passing to the
limit immediately reduces us to the case where $k$ is separably
closed.

Every torus over such a field is split.  Hence, if $P$ admits a
composition series as in \SGARefLink{sga7.sga7i.eviii.n4-6}{4.6}, it also admits one whose quotients are
isomorphic to $\underline G_a$ or $\underline G_m$.  Repeated
application of \SGARefLink{sga7.sga7i.eviii.n4-5}{4.5}, together with
\[
 \Pic(\underline G_{a_K})=\Pic(\underline G_{m_K})=0
\]
for every field $K$, gives
\[
 \Pic(P)=0.
\]
This equality remains true after every extension of the base field
$k$.  Applying \SGARefLink{sga7.sga7i.eviii.n4-5}{4.5} to
\[
 \pr_2:P\times Q\longrightarrow Q
\]
and to $\mathbb L$, we find a line bundle $\mathbb M$ on $Q$ such
that $\mathbb L\simeq\pr_2^*\mathbb M$.  Since the restriction of
$\mathbb L$ to $e_P\times Q$ is trivial, $\mathbb M$ is trivial, and
so is $\mathbb L$.  This proves the first case.

In the second case, let $Q$ be an extension of an abelian scheme $B$
by a group $M$ admitting a composition series with quotients
isomorphic to $\underline G_m$ or $\underline G_a$.  Apply \SGARefLink{sga7.sga7i.eviii.n4-5}{4.5} to
\[
 f:X=P\times Q\longrightarrow Y=P\times B.
\]
The fibre over the generic point $y$ of $P\times B$ is an
$M_y$-torsor.  It is necessarily trivial by the particular structure
of $M$ and the classical equalities
\[
 H^1(K,\underline G_a)=0,
 \qquad
 H^1(K,\underline G_m)=0,
\]
valid for every field $K$, here taken to be $k(y)$.  Since
$\Pic(M)=0$, we have $\Pic(X_y)=0$; thus \SGARefLink{sga7.sga7i.eviii.n4-5}{4.5} applies and shows that
$\mathbb L$ is the pullback of a line bundle $\mathbb N$ on
$Y=P\times B$.

Let $\mathbb N_1$ and $\mathbb N_2$ be the restrictions of $\mathbb
N$ to $P\times e_B$ and $e_P\times B$, respectively.  Replacing
$\mathbb N$ by
\[
 \mathbb N\otimes\pr_1^*(\mathbb N_1)^{-1}
 \otimes\pr_2^*(\mathbb N_2)^{-1},
\]
we may suppose that $\mathbb N$ is birigidified and that $\mathbb L$
is its birigidified pullback.  We are thereby reduced to the case in
which $Q$ itself is an abelian scheme.

Then $\mathbb L$ defines a canonical morphism
\begin{equation}\tag{*}
 P\longrightarrow\underline{\Pic}_{Q/k}
\end{equation}
taking the identity to the identity.  Since $P$ is connected, this
morphism factors through the dual abelian scheme
\[
 Q'=\underline{\Pic}^{,0}_{Q/k}.
\]
It is well known that every morphism of schemes from a smooth
connected affine algebraic group to an abelian scheme is constant.
This is a geometric assertion, so we may suppose $k$ algebraically
closed; a devissage then reduces it immediately to $P=\underline
G_a$ or $P=\underline G_m$, and every rational map from
$\mathbb P^1_k$ to an abelian scheme over $k$ is constant.

The constancy, hence vanishing, of (*) means that $\mathbb L$ is the
pullback of a line bundle $\mathbb M$ on $P$.  Since $\mathbb L$ is
trivial on $P\times e_Q$, $\mathbb M$ is trivial, and therefore so is
$\mathbb L$.  This completes the proof.


% END INCLUDED SGA7I_English_current_buildable_source_20260823/components/110_expose_VIII_sections_4_5_and_4_6.tex

% BEGIN INCLUDED SGA7I_English_current_buildable_source_20260823/components/111_expose_VIII_sections_4_7_through_4_10.tex
% SOURCE AUTHORITY:
% expose_VIII_body.tex lines 1607--1696; scan indices 280--284;
% source folios 269--273.
% Direct scan comparison: physical PDF pages 281--284 in overlapping
% 1100-dpi bands.  The prose and displayed formulas are fully legible at
% that resolution; no 5000--9000-dpi ambiguity escalation was needed.
% Source disposition in the proof of Corollary 4.8: the print displays
% Biext^i(P,M;G_m).  This is retained.  Together with (4.7.1), applied with
% Q=M, those vanishings give the Biext^i(A,M;G_m) vanishings required by
% the exact sequence for 0 -> M -> Q -> B -> 0.
% Hard stop before 4.11, source line 1701 / scan index 284.

\medskip
\noindent\underline{Proposition \SGARefTarget{sga7.sga7i.eviii.n4-7}{4.7}.}  \underline{Let $P$ and $Q$
be two smooth connected algebraic groups over the field $k$.
Suppose that $P$ is an extension of an abelian scheme $A$ by a smooth
connected affine algebraic group $L$ admitting a composition series
each of whose factors is isomorphic to a torus or to $\underline
G_a$ (a condition always satisfied when $k$ is perfect).  Then the
functor considered in \SGARefLink{sga7.sga7i.eviii.n1-3}{1.3} is an equivalence of categories; that is,
every birigidified line bundle $\mathbb L$ on $P\times Q$ comes from
a biextension of $(P,Q)$.  Moreover, the functor between the rigid
categories}
\[
 \operatorname{BIEXT}(A,Q;\underline G_m)
 \longrightarrow
 \operatorname{BIEXT}(P,Q;\underline G_m)
\]
\underline{is an equivalence of categories; in other words, it
induces an isomorphism of groups}
\begin{equation}\tag{\SGARefTarget{sga7.sga7i.eviii.n4-7-1}{4.7.1}}
 \Biext^1(A,Q;\underline G_m)
 \xrightarrow{\ \sim\ }
 \Biext^1(P,Q;\underline G_m).
\end{equation}

For the first assertion, descent and \SGARefLink{sga7.sga7i.eviii.n4-3}{4.3} again reduce us to the case
where $k$ is separably closed.  Proceeding as in the proof of \SGARefLink{sga7.sga7i.eviii.n4-6}{4.6},
we deduce from \SGARefLink{sga7.sga7i.eviii.n4-5}{4.5} and the hypothesis on $L$ that $\mathbb L$ is
isomorphic to the pullback of a birigidified line bundle $\mathbb M$
on $A\times Q$.  For the first assertion of \SGARefLink{sga7.sga7i.eviii.n4-7}{4.7} we are thus reduced
to the case $A=P$.  The line bundle $\mathbb L$ then defines a
morphism of pointed schemes
\[
 f:Q\longrightarrow\underline{\Pic}_{P/k},
\]
which necessarily factors through the dual abelian scheme $P'$ of
$P$.  Let $\mathbb L'$ be the pullback of the Weil sheaf on
$P\times P'$ along
\[
 \id_P\times f:P\times Q\longrightarrow P\times P'.
\]
This is a birigidified line bundle which, by the definition of $f$,
differs from $\mathbb L$ only by the pullback of a line bundle
$\mathbb M$ on $P$.  Since both $\mathbb L'$ and $\mathbb L$ have
trivial restriction to $P\times e_Q$, it follows that $\mathbb M$ is
trivial and hence that $\mathbb L\simeq\mathbb L'$.  We noted in VII
\SGARefLink{sga7.sga7i.evii.n2-9-4}{2.9.4} that divisorial correspondences on a product of two abelian
schemes arise from biextensions.  Consequently $\mathbb L'$, and
therefore $\mathbb L$, comes from a biextension.

Returning to the general case, where $P$ is an extension of $A$ by
$L$, it remains to prove that (\SGARefLink{sga7.sga7i.eviii.n4-7-1}{4.7.1}) is an isomorphism.  This follows
from the exact sequence of the $\Biext^i$, for $i=0,1$, associated
with
\[
 0\longrightarrow L\longrightarrow P\longrightarrow A
 \longrightarrow0,
\]
and from the relations
\[
 \Biext^0(L,Q;\underline G_m)=0,
 \qquad
 \Biext^1(L,Q;\underline G_m)=0.
\]
The first is a special case of the rigidity assertion in \SGARefLink{sga7.sga7i.eviii.n4-3}{4.3}, and
the second is precisely the first case of \SGARefLink{sga7.sga7i.eviii.n4-6}{4.6}.

\medskip
\noindent\underline{Corollary \SGARefTarget{sga7.sga7i.eviii.n4-8}{4.8}.}  \underline{Under the hypotheses
of \SGARefLink{sga7.sga7i.eviii.n4-7}{4.7}, suppose that $Q$ is an extension of an algebraic group $B$
(necessarily smooth and connected) by a smooth connected affine
algebraic group $M$.  Then the functor between the rigid categories}
\[
 \operatorname{BIEXT}(A,B;\underline G_m)
 \longrightarrow
 \operatorname{BIEXT}(P,Q;\underline G_m)
\]
\underline{is an equivalence of categories; in other words, it
induces an isomorphism of groups}
\begin{equation}\tag{\SGARefTarget{sga7.sga7i.eviii.n4-8-1}{4.8.1}}
 \Biext^1(A,B;\underline G_m)
 \xrightarrow{\ \sim\ }
 \Biext^1(P,Q;\underline G_m).
\end{equation}

In view of (\SGARefLink{sga7.sga7i.eviii.n4-7-1}{4.7.1}), it remains to prove that
\[
 \Biext^1(A,B;\underline G_m)
 \xrightarrow{\ \sim\ }
 \Biext^1(A,Q;\underline G_m)
\]
is an isomorphism.  As in the proof of (\SGARefLink{sga7.sga7i.eviii.n4-7-1}{4.7.1}), this follows from the
exact sequence of the $\Biext^i$ associated with
\[
 0\longrightarrow M\longrightarrow Q\longrightarrow B
 \longrightarrow0,
\]
together with the relations
\[
 \Biext^0(P,M;\underline G_m)=0,
 \qquad
 \Biext^1(P,M;\underline G_m)=0.
\]
Indeed, (\SGARefLink{sga7.sga7i.eviii.n4-7-1}{4.7.1}), applied with $M$ in place of $Q$, identifies these
groups with the corresponding groups having $A$ in the first
variable, as required by the preceding exact sequence.  The second
vanishing displayed above is a special case of the second case of
\SGARefLink{sga7.sga7i.eviii.n4-6}{4.6}, after interchanging the two factors.

\medskip
\noindent\underline{Remark \SGARefTarget{sga7.sga7i.eviii.n4-9}{4.9}.}  Let $P$ and $Q$ be two smooth
connected algebraic groups over a \underline{perfect} field $k$.
It is then known [3] that each is an extension of an abelian scheme
by a smooth connected affine group, the latter being the product of
a torus with an iterated extension of copies of $\underline G_a$
(cf. the reference given in \SGARefLink{sga7.sga7i.eviii.n4-6}{4.6}):
\[
 0\longrightarrow L\longrightarrow P\longrightarrow A
 \longrightarrow0,
 \qquad
 0\longrightarrow M\longrightarrow Q\longrightarrow B
 \longrightarrow0.
\]
Thus \SGARefLink{sga7.sga7i.eviii.n4-8}{4.8} applies and shows that the theory of biextensions of $(P,Q)$
by $\underline G_m$ reduces to the analogous (and much more
classical!) theory for the pair $(A,B)$ of abelian varieties.

\medskip
\noindent\underline{Corollary \SGARefTarget{sga7.sga7i.eviii.n4-10}{4.10}.}  \underline{Let $P$ and $Q$ be
two smooth connected algebraic groups over $k$, let $L$ be a smooth
connected affine algebraic subgroup of $P$, and let $E$ be a
biextension of $(P,Q)$ by $\underline G_m$.  By \SGARefLink{sga7.sga7i.eviii.n2-2}{2.2}, $E$ induces a
pairing}
\[
 T_\ell(P)\times T_\ell(Q)\longrightarrow T_\ell(\underline G_m).
\]
\underline{Then the induced pairing}
\[
 T_\ell(L)\times T_\ell(Q)\longrightarrow T_\ell(\underline G_m)
\]
\underline{is zero.}

Indeed, we may plainly suppose that $k$ is algebraically closed.  We
are then under the hypotheses of \SGARefLink{sga7.sga7i.eviii.n4-6}{4.6}, which imply that the
restriction of $E$ to $(L,Q)$ is the trivial biextension.  The
asserted conclusion follows, since the induced pairing in question
is the pairing defined by this restriction.

% END INCLUDED SGA7I_English_current_buildable_source_20260823/components/111_expose_VIII_sections_4_7_through_4_10.tex

% BEGIN INCLUDED SGA7I_English_current_buildable_source_20260823/components/112_expose_VIII_section_4_11.tex
% SOURCE AUTHORITY:
% expose_VIII_body.tex lines 1701--1743; scan indices 284--286;
% source folios 273--275.
% Direct scan comparison: physical PDF pages 285--287 in overlapping
% 1100-dpi bands.  All formulas, primes, kernels, and arrow directions are
% legible at that resolution; no higher-resolution escalation was needed.
% Transparent source disposition: the reference printed as (1.11.1) in the
% middle of the discussion is the pairing (4.11.1), as the local argument
% and every surrounding reference require.
% Hard stop before 4.12, source line 1745 / scan index 286.

\medskip
\noindent \SGARefTarget{sga7.sga7i.eviii.n4-11}{4.11}.  Let $P$ and $Q$ be algebraic groups smooth over
$k$, let $E$ be a biextension of $(P,Q)$ by $\underline G_m$, and
let $T$ be the largest subtorus of $P$.  We shall determine the left
annihilator of the pairing
\begin{equation}\tag{\SGARefTarget{sga7.sga7i.eviii.n4-11-1}{4.11.1}}
 T_\ell(P)\times T_\ell(Q)
 \longrightarrow T_\ell(\underline G_{m,k}),
\end{equation}
where $\ell$ is again a fixed prime number.  By \SGARefLink{sga7.sga7i.eviii.n4-10}{4.10} this annihilator
always contains the subgroup $T_\ell(T)$ of $T_\ell(P)$.

Choose an extension $k'$ of $k$ over which $P$ and $Q$ can each be
written as an extension of an abelian scheme by an affine group
scheme that is an iterated extension of tori and copies of
$\underline G_a$.  If $P$ and $Q$ have unipotent rank zero, we may
take $k'=k$; in general, we may take for $k'$ the perfect closure of
$k$ (or a finite radicial extension).  Let $A'$ and $B'$ be the
abelian parts of $P_{k'}$ and $Q_{k'}$, respectively.  By \SGARefLink{sga7.sga7i.eviii.n4-7}{4.7}, the
biextension $E_{k'}$ of $(P_{k'},Q_{k'})$ comes from a biextension
$F'$ of $(A',B')$, uniquely determined up to unique isomorphism.
Consequently the pairing
\begin{equation}\tag{\SGARefTarget{sga7.sga7i.eviii.n4-11-2}{4.11.2}}
 T_\ell(P_{k'})\times T_\ell(Q_{k'})
 \longrightarrow T_\ell(\underline G_{m,k'}),
\end{equation}
obtained from (\SGARefLink{sga7.sga7i.eviii.n4-11-1}{4.11.1}) by the base-field extension $k\to k'$, and
also defined by the biextension $E_{k'}$, is obtained from the
pairing
\begin{equation}\tag{\SGARefTarget{sga7.sga7i.eviii.n4-11-3}{4.11.3}}
 T_\ell(A')\times T_\ell(B')
 \longrightarrow T_\ell(\underline G_{m,k'})
\end{equation}
defined by $F'$, by composing with the canonical homomorphisms
(which are in fact epimorphisms if $P$ and $Q$ are
$\ell$-divisible)
\begin{equation}\tag{\SGARefTarget{sga7.sga7i.eviii.n4-11-4}{4.11.4}}
 T_\ell(P_{k'})\longrightarrow T_\ell(A'),
 \qquad
 T_\ell(Q_{k'})\longrightarrow T_\ell(B').
\end{equation}

It follows that the left annihilator of (\SGARefLink{sga7.sga7i.eviii.n4-11-2}{4.11.2}), which is obtained
from the left annihilator of (\SGARefLink{sga7.sga7i.eviii.n4-11-1}{4.11.1}) by base change, is the inverse
image under the first homomorphism in (\SGARefLink{sga7.sga7i.eviii.n4-11-4}{4.11.4}) of the left
annihilator of (\SGARefLink{sga7.sga7i.eviii.n4-11-3}{4.11.3}).  Let $C'$ be the largest abelian subscheme
of $A'$ such that the restriction of $F'$ to $(C',B')$ is the
trivial biextension.  Equivalently, $C'$ is the reduced identity
component of the kernel of the homomorphism
\begin{equation}\tag{\SGARefTarget{sga7.sga7i.eviii.n4-11-5}{4.11.5}}
 A'\longrightarrow (B')^*
\end{equation}
defined by $F'$, where $(B')^*$ is the dual abelian scheme of $B'$.
The left annihilator of (\SGARefLink{sga7.sga7i.eviii.n4-11-3}{4.11.3}) identifies with $T_\ell(C')$;
therefore the left annihilator of (\SGARefLink{sga7.sga7i.eviii.n4-11-1}{4.11.1}) is precisely the inverse
image of
\[
 T_\ell(C')\lhook\joinrel\longrightarrow T_\ell(A')
\]
under the first morphism in (\SGARefLink{sga7.sga7i.eviii.n4-11-4}{4.11.4}).

In particular, to say that the left annihilator of (\SGARefLink{sga7.sga7i.eviii.n4-11-1}{4.11.1}) is
exactly $T_\ell(T)$ -- equivalently, that the left annihilator of
(\SGARefLink{sga7.sga7i.eviii.n4-11-2}{4.11.2}) is the kernel $T_\ell(T')$ of the first morphism in
(\SGARefLink{sga7.sga7i.eviii.n4-11-4}{4.11.4}) -- is to say that (\SGARefLink{sga7.sga7i.eviii.n4-11-3}{4.11.3}) is nondegenerate on the left,
that is, that $C'=0$, or equivalently that (\SGARefLink{sga7.sga7i.eviii.n4-11-5}{4.11.5}) has finite
kernel.  In this last form the condition is visibly independent of
$\ell$, and in any of its forms it is plainly independent of the
choice of $k'$.  When it holds, we shall simply say that the given
biextension $E$ of $(P,Q)$ by $\underline G_m$ is
\underline{nondegenerate on the left}.  By duality, this also means
that the homomorphism
\begin{equation}\tag{\SGARefTarget{sga7.sga7i.eviii.n4-11-6}{4.11.6}}
 B'\longrightarrow (A')^*
\end{equation}
defined by $F'$ is an epimorphism, that is, is surjective.

The notion of a biextension \underline{nondegenerate on the right}
is defined symmetrically: (\SGARefLink{sga7.sga7i.eviii.n4-11-6}{4.11.6}) must have finite kernel, or,
equivalently, (\SGARefLink{sga7.sga7i.eviii.n4-11-5}{4.11.5}) must be an epimorphism.  The biextension $E$
of $(P,Q)$ by $\underline G_{m,k}$ is called
\underline{nondegenerate} if it is nondegenerate on both the left
and the right.  This is equivalent to saying that (\SGARefLink{sga7.sga7i.eviii.n4-11-5}{4.11.5}), or
equivalently (\SGARefLink{sga7.sga7i.eviii.n4-11-6}{4.11.6}), is an isogeny.

It is immediate from the definition that these notions are
invariant under every extension of the base field $k$.

% END INCLUDED SGA7I_English_current_buildable_source_20260823/components/112_expose_VIII_section_4_11.tex

% BEGIN INCLUDED SGA7I_English_current_buildable_source_20260823/components/113_expose_VIII_sections_4_12_through_4_14.tex
% SOURCE AUTHORITY:
% expose_VIII_body.tex lines 1745--1799; scan indices 286--289;
% source folios 275--278.
% Direct scan comparison: physical PDF pages 287--290 in overlapping
% 1100-dpi bands.  The formulas and prose are fully legible at that
% resolution; no 5000--9000-dpi ambiguity escalation was needed.
% Transparent source dispositions in 4.12:
% - the canonical quotient is f:P -> A, not the printed f:A -> P; this is
%   forced by the displayed pullback Pic(A) -> Pic(P);
% - the descended relation is L \simeq f^*M, not the printed L^* relation;
% - M is a line bundle on A, not the printed "on M".
% In the last sentence of 4.14, the delta map on Pic(P) is (4.14.2), not
% the printed reference (4.14.1).
% Hard stop before section 5, source line 1800 / scan index 289.

\medskip
\noindent \SGARefTarget{sga7.sga7i.eviii.n4-12}{4.12}.  Let $P$ be a smooth connected group scheme, let
$\mathbb L$ be a line bundle on $P$, and put
\begin{equation}\tag{\SGARefTarget{sga7.sga7i.eviii.n4-12-1}{4.12.1}}
 \delta(\mathbb L)=
 \pi^*\mathbb L\otimes
 \pr_1^*(\mathbb L)^{-1}\otimes
 \pr_2^*(\mathbb L)^{-1},
\end{equation}
with the notation of \SGARefLink{sga7.sga7i.eviii.n4-2}{4.2}.  As already observed in VII \SGARefLink{sga7.sga7i.evii.n1-3-4}{1.3.4}, this is
a birigidified line bundle on $P\times P$.  By \SGARefLink{sga7.sga7i.eviii.n4-3}{4.3}, if
$\delta(\mathbb L)$ comes from a biextension, that biextension is
determined up to unique isomorphism.  In particular, by \SGARefLink{sga7.sga7i.eviii.n4-7}{4.7}, it is
defined when $P$ is an extension of an abelian scheme $A$ by an
iterated extension of tori and copies of $\underline G_a$.  Whenever
it is defined, this biextension is plainly symmetric.

In the preceding favourable case, and assuming that the maximal
torus of $P$ is split, the construction of $\delta(\mathbb L)$ can
be reduced to the case where $P$ is an abelian scheme.  Indeed, the
construction is functorial in $P$, while the homomorphism
\begin{equation}\tag{\SGARefTarget{sga7.sga7i.eviii.n4-12-2}{4.12.2}}
 \Pic(A)\longrightarrow\Pic(P)
\end{equation}
induced by the canonical quotient homomorphism
$f:P\longrightarrow A$ is surjective by \SGARefLink{sga7.sga7i.eviii.n4-5}{4.5} (cf. the beginning of
the proof of \SGARefLink{sga7.sga7i.eviii.n4-6}{4.6}).  We can therefore find an isomorphism
\[
 \mathbb L\simeq f^*\mathbb M,
\]
where $\mathbb M$ is a line bundle on $A$.  Then
$\delta(\mathbb L)$ is the pullback of the biextension
$\delta(\mathbb M)$ of $(A,A)$ by $\underline G_{m,k}$.

\medskip
\noindent\underline{Proposition \SGARefTarget{sga7.sga7i.eviii.n4-13}{4.13}.}  \underline{Under the
hypotheses of \SGARefLink{sga7.sga7i.eviii.n4-12}{4.12}, suppose that $\mathbb L$ is an ample line bundle
on $P$.  If the biextension $\delta(\mathbb L)$ of $(P,P)$ by
$\underline G_{m,k}$ is defined, then it is nondegenerate in the
sense of \SGARefLink{sga7.sga7i.eviii.n4-11}{4.11}.  If $P$ is an extension of an abelian scheme $A$ by a
smooth connected affine group $L$ satisfying the condition of \SGARefLink{sga7.sga7i.eviii.n4-7}{4.7},
then the biextension of $(A,A)$ from which $\delta(\mathbb L)$ comes
by \SGARefLink{sga7.sga7i.eviii.n4-8}{4.8} is of the form $\delta(\mathbb M)$, where $\mathbb M$ is an
ample line bundle on $A$.}

We reduce to the case where $k$ is algebraically closed.  By the
definition, it is then enough, with the notation of \SGARefLink{sga7.sga7i.eviii.n4-12}{4.12}, to prove
that the biextension $\delta(\mathbb M)$ of $(A,A)$ is
nondegenerate.  M. Raynaud proves in his thesis [2, XI 1.11] that if
$\mathbb L=f^*\mathbb M$ is ample, then $\mathbb M$ is ample as
well.  The theory of A. Weil then gives the familiar conclusion that
$\delta(\mathbb M)$ is nondegenerate [1, 2], that is, that the
homomorphism from $A$ to its dual abelian scheme is an isogeny.

\medskip
\noindent\underline{Remark \SGARefTarget{sga7.sga7i.eviii.n4-14}{4.14}.}  Suppose we are in the favourable
case where $P$ is an extension of an abelian scheme $A$ by a group
$L$ which is an iterated extension of tori and copies of
$\underline G_a$, and consider the canonical composite
\begin{equation}\tag{\SGARefTarget{sga7.sga7i.eviii.n4-14-1}{4.14.1}}
 \operatorname{NS}(A)
 \xrightarrow{\ \delta\ }
 \Biext^1(A,A;\underline G_m)
 \xrightarrow{\ \sim\ }
 \Biext^1(P,P;\underline G_m).
\end{equation}
By \SGARefLink{sga7.sga7i.eviii.n4-8}{4.8} the second arrow is bijective.  On the other hand, the theory
of abelian varieties shows that the first arrow identifies the
Neron--Severi group $\operatorname{NS}(A)$ of $A$ with the subgroup
of
\[
 \Biext^1(A,A;\underline G_m)=\operatorname{Corr}(A,A)
\]
consisting of the classes of symmetric biextensions, or equivalently
symmetric divisorial correspondences.  Thus (\SGARefLink{sga7.sga7i.eviii.n4-14-1}{4.14.1}) identifies
$\operatorname{NS}(A)$ with the group of symmetric biextensions of
$(P,P)$ by $\underline G_m$.  By \SGARefLink{sga7.sga7i.eviii.n4-7}{4.7} this group can equally be
described as the group of classes of symmetric divisorial
correspondences on $(P,P)$.

It is natural to call a divisorial correspondence on $(P,P)$
\underline{ample}, and likewise to call the corresponding
biextension of $(P,P)$ by $\underline G_m$ \underline{ample}, when
it comes from an ample element of $\operatorname{NS}(A)$.  This
condition is plainly invariant under extension of the base field;
the notion therefore still makes sense without the preceding
restrictive hypothesis on $P$, assuming only that $P$ is smooth and
connected.

Suppose further that $L$ is ``split solvable over $k$'', in the
sense that it admits a devissage by copies of $\underline G_a$ and
$\underline G_m$.  The preceding observation and the surjectivity of
(\SGARefLink{sga7.sga7i.eviii.n4-12-2}{4.12.2}) show that the image of
\begin{equation}\tag{\SGARefTarget{sga7.sga7i.eviii.n4-14-2}{4.14.2}}
 \Pic(P)\xrightarrow{\ \delta\ }
 \Biext^1(P,P;\underline G_m)
\end{equation}
consists of the symmetric biextensions, and that its kernel is the
subgroup $\Pic^0(P)$ of $\Pic(P)$ which is the image of $\Pic^0(A)$.
Consequently, if we put
\[
 \operatorname{NS}(P)=\Pic(P)/\Pic^0(P)
 \simeq\operatorname{NS}(A),
\]
we may also interpret (\SGARefLink{sga7.sga7i.eviii.n4-14-1}{4.14.1}) as a canonical isomorphism
\begin{equation}\tag{\SGARefTarget{sga7.sga7i.eviii.n4-14-3}{4.14.3}}
 \operatorname{NS}(P)
 \xrightarrow{\ \sim\ }
 \Biext^1(P,P;\underline G_m)^{\mathrm{sym}}.
\end{equation}

Finally, the theorem of Raynaud cited in \SGARefLink{sga7.sga7i.eviii.n4-13}{4.13} says that an element
$\xi$ of $\Pic(P)$ is ample if and only if its image
$\delta(\xi)$ under (\SGARefLink{sga7.sga7i.eviii.n4-14-2}{4.14.2}) is the class of an ample biextension.

% END INCLUDED SGA7I_English_current_buildable_source_20260823/components/113_expose_VIII_sections_4_12_through_4_14.tex

% BEGIN INCLUDED SGA7I_English_current_buildable_source_20260823/components/114_expose_VIII_section_5_opening_through_5_6.tex
% SOURCE AUTHORITY:
% expose_VIII_body.tex lines 1800--1916; scan indices 289--295;
% source folios 278--284.
% Direct scan comparison: physical PDF pages 290--296 in overlapping
% 1100-dpi bands.  The mathematical text and displayed formulas are fully
% legible at that resolution; no higher-resolution escalation was needed.
% Transparent source dispositions:
% - in the proof of 5.1, G_S and the final E are corrected to Z_S and Z;
% - the source repeats the target category in Corollary 5.2; the first
%   category is restored to torsors on P under Z_P;
% - the repeated P_i at the end of that proof is read as P_{i\eta};
% - in 5.5(ii), the second finite-type example is P, and every stray G_k or
%   G in the statement and proof is corrected to the stated coefficient Z;
% - in 5.6, Phi and Psi are groups of connected components (not identity
%   components), and the stray G_k in (5.6.1) is corrected to Z_k.
% Hard stop before 5.7, source line 1917 / scan index 295.

\medskip
\noindent\underline{5.  Extensions and biextensions by torsion-free
truncated constant group schemes.}

In this section and the next we give some auxiliary results that will
be used in section 7 in the case $Z=\mathbb Z$.

\medskip
\noindent\underline{Proposition \SGARefTarget{sga7.sga7i.eviii.n5-1}{5.1}.}  \underline{Let $S$ be an
irreducible geometrically unibranch scheme with generic point $\eta$,
and let $Z$ be a torsion-free group, hence giving a constant group
scheme $Z_S$ over $S$.  Then every torsor $E$ under $Z_S$ is trivial,
and the functor}
\[
 E\longmapsto E(\eta)
\]
\underline{from the category of torsors under $Z_S$ to the category of
torsors under $Z=Z_S(\eta)$ is an equivalence of categories.}

A familiar lemma of descent theory (SGA 3 X 5.4) shows that every
torsor $E$ under $Z_S$ is representable.  In fact, that result applies
to every group scheme over $S$ which is separated, locally of finite
presentation, and locally quasi-finite over $S$.  Since $Z_S$ is
etale over $S$, fpqc descent shows that the same is true of every
torsor $E$ under $Z_S$; such a torsor is therefore locally trivial for
the \underline{etale topology} (and not merely for the fppf topology).
To prove that $E$ is trivial, it is consequently enough to prove
\[
 H^1(S_{\mathrm{\acute{e}t}},Z_S)=e.
\]
Likewise, the full-faithfulness assertion of \SGARefLink{sga7.sga7i.eviii.n5-1}{5.1} reduces to the
relation
\[
 Z\xrightarrow{\ \sim\ }H^0(S_{\mathrm{\acute{e}t}},Z_S).
\]
The latter is immediate.  To prove the former, consider the inclusion
\[
 i:\eta\longrightarrow S.
\]
Because $S$ is geometrically unibranch, we have
\[
 Z_S\xrightarrow{\ \sim\ }i_*(Z_\eta)
\]
(SGA 4 IX 2.14.1).  The noncommutative low-degree exact sequence that
replaces the Leray spectral sequence for $i:\eta\to S$ (SGA 4 XII
3.2) therefore begins
\[
 e\longrightarrow H^1(S,i_*(Z_\eta))
   \longrightarrow H^1(\eta,Z_\eta)\longrightarrow\cdots,
\]
and shows that
\[
 H^1(S,Z_S)\longrightarrow H^1(\eta,Z_\eta)
\]
is injective.  We are thus reduced to the case in which $S$ is the
spectrum of a field $k$.  If $\pi$ is the Galois group of a separable
closure of $k$, then
\[
 H^1(k,Z)=\Hom_{\mathrm{cont}}(\pi,Z)/Z=e,
\]
because $Z$ is torsion-free.  This completes the proof.

\medskip
\noindent\underline{Corollary \SGARefTarget{sga7.sga7i.eviii.n5-2}{5.2}.}  \underline{With the notation of
\SGARefLink{sga7.sga7i.eviii.n5-1}{5.1}, let $P$ be a smooth scheme over $S$.  Then every torsor $E$ under
$Z_P$ is trivial, and the restriction functor}
\[
 E\longmapsto E_\eta
\]
\underline{from the category of torsors on $P$ under $Z_P$ to the
category of torsors on $P_\eta$ under $Z_{P_\eta}$ is an equivalence of
categories.}

Indeed, $P$ is geometrically unibranch (EGA IV 17.5.7).  Moreover,
every irreducible component of $P$ dominates $S$, because $P$ is flat
over $S$, and the set of irreducible components of $P$ is locally
finite.  Thus $P$ is the disjoint sum of irreducible geometrically
unibranch schemes $P_i$.  The assertion now follows immediately from
\SGARefLink{sga7.sga7i.eviii.n5-1}{5.1} applied to the various $P_i$ and their generic fibres $P_{i\eta}$,
noting that the generic point $p_i$ of $P_i$ is also the generic point
of $P_{i\eta}$.

\medskip
\noindent \SGARefTarget{sga7.sga7i.eviii.n5-2-1}{5.2.1}.  The preceding corollary immediately gives the
analogous statement for bitorsors on $P$ under $(Z_P,Z_P)$: restriction
to the generic fibre induces an equivalence from their category to the
category of bitorsors on $P_\eta$ under
$(Z_{P_\eta},Z_{P_\eta})$.  Indeed, it is enough to express the bitorsor
structure on $E$ in terms of its structure as a right torsor under
$Z_P$, together with an isomorphism from $Z_P$ to the commutant of
$Z_P$ for its action on $E$.  Since $E$ is necessarily trivial, that
commutant is itself isomorphic to $Z_P$.  We are thus reduced, on each
irreducible component $X$ of $P$, to the following immediate
observation: if $X$ is a nonempty connected scheme, the functor
$Z\mapsto Z_X$ from sets to $X$-schemes is fully faithful.

\medskip
\noindent\underline{Proposition \SGARefTarget{sga7.sga7i.eviii.n5-3}{5.3}.}  \underline{Let $S$ and $Z$ be
as in \SGARefLink{sga7.sga7i.eviii.n5-1}{5.1}, and let $P$ be a smooth group scheme over $S$.  Consider
extensions $E$ of group schemes of $P$ by $Z_S$.  Then restriction to
the generic fibre}
\[
 E\longmapsto E_\eta
\]
\underline{is an equivalence from the category of these extensions to
the category of extensions of $P_\eta$ by $Z_\eta$.  If $P$ and $Z$
are commutative, it also induces an equivalence between the
corresponding categories of commutative extensions.}

First observe that every extension of sheaves of groups of $P$ by
$Z_S$ (and similarly of $P_\eta$ by $Z_\eta$) gives, by \SGARefLink{sga7.sga7i.eviii.n5-2}{5.2}, a trivial
torsor, hence in particular a representable one.  The categories in
\SGARefLink{sga7.sga7i.eviii.n5-3}{5.3} may therefore equally well be interpreted as categories of
extensions of sheaves of groups.  The description in VII \SGARefLink{sga7.sga7i.evii.n1-1-6}{1.1.6} and
\SGARefLink{sga7.sga7i.evii.n1-2}{1.2} expresses these in terms of bitorsors on
$P$, $P\times P$, and $P\times P\times P$ (respectively on
$P_\eta$, $P_\eta\times P_\eta$, and
$P_\eta\times P_\eta\times P_\eta$), under the pairs of groups
obtained from $(Z,Z)$ by base change.  Applying \SGARefLink{sga7.sga7i.eviii.n5-2}{5.2} in the commutative
case and \SGARefLink{sga7.sga7i.eviii.n5-2-1}{5.2.1} in the general case to the smooth $S$-schemes
$P$, $P\times P$, and $P\times P\times P$ gives the asserted
conclusion.  The same argument gives the following result.

\medskip
\noindent\underline{Corollary \SGARefTarget{sga7.sga7i.eviii.n5-4}{5.4}.}  \underline{Let $S,Z$ be as in
\SGARefLink{sga7.sga7i.eviii.n5-1}{5.1}, and let $P,Q$ be two smooth group schemes over $S$.  Suppose that
$P,Q,$ and $Z$ are commutative, and consider biextensions $E$ of
$(P,Q)$ by $Z_S$.  Then restriction to the generic fibre}
\[
 E\longmapsto E_\eta
\]
\underline{is an equivalence from the category of these biextensions
to the category of biextensions of $(P_\eta,Q_\eta)$ by $Z_\eta$.}

The preceding results reduce the determination of categories of
extensions or biextensions by a torsion-free constant group to the
case in which the base scheme is the spectrum of a field.  In that
case we obtain the following.

\medskip
\noindent\underline{Proposition \SGARefTarget{sga7.sga7i.eviii.n5-5}{5.5}.}  \underline{Let $k$ be a field,
$Z$ a torsion-free group, $P$ a $k$-group scheme locally of finite
type, $P^\circ$ its identity component, and
$\Phi=P/P^\circ$ its group of connected components (SGA 3
$\mathrm{VI}_A$ 4).  Then:}

\smallskip
\noindent (i) \underline{The natural functor from the category of group
schemes which are extensions (respectively commutative extensions) of
$\Phi$ by $Z_k$ to the category of group schemes which are extensions
(respectively commutative extensions) of $P$ by $Z_k$ is an
equivalence.  In particular, if $P$ is connected, every extension of
$P$ by $Z_k$ is trivial, and the category of such extensions is
rigid.}

\smallskip
\noindent (ii) \underline{Suppose that $\Phi$ is a torsion group and
that $P$ and $Z$ are commutative.  Then the category of extensions of
$P$ by $Z_k$ is rigid: every automorphism of an object of this
category is trivial.  If, in addition, there is an integer $n>0$ such
that $n\id_\Phi=0$---for example, if $\Phi$ is finite, as happens when
$P$ is of finite type---then the group of classes of commutative
extensions of $P$ by $Z_k$ is given by the canonical isomorphism}
\begin{equation}\tag{\SGARefTarget{sga7.sga7i.eviii.n5-5-1}{5.5.1}}
 \Ext^1(P,Z_k)
 \xleftarrow{\ \sim\ }
 \Hom\bigl(\Phi,(Z\otimes_{\mathbb Z}\mathbb Q)/Z\bigr),
\end{equation}
\underline{defined by the boundary operator in the exact sequence of
$\Ext$ groups associated with the exact sequence}
\begin{equation}\tag{\SGARefTarget{sga7.sga7i.eviii.n5-5-2}{5.5.2}}
 0\longrightarrow Z\longrightarrow Z\otimes_{\mathbb Z}\mathbb Q
 \longrightarrow (Z\otimes_{\mathbb Z}\mathbb Q)/Z
 \longrightarrow 0.
\end{equation}

The canonical morphism
\[
 P\longrightarrow\Phi
\]
plainly induces an isomorphism on the sets of irreducible components,
since it is flat with irreducible fibres.  By \SGARefLink{sga7.sga7i.eviii.n5-1}{5.1}, the induced functor
from torsors under $Z_\Phi$ to torsors under $Z_P$ is therefore an
equivalence.  (Recall that a group scheme locally of finite type over
$k$ is geometrically unibranch.)  The same conclusion applies to
bitorsors and to the morphisms
\[
 P\times P\longrightarrow\Phi\times\Phi,
 \qquad
 P\times P\times P\longrightarrow\Phi\times\Phi\times\Phi.
\]
The description in VII \SGARefLink{sga7.sga7i.evii.n1-1-6}{1.1.6} of extensions of $P$ by $Z_k$ and of
$\Phi$ by $Z_k$ (and VII \SGARefLink{sga7.sga7i.evii.n1-3}{1.3} in the commutative case) now gives the
assertions in (i).

We restrict attention to commutative extensions.  The automorphism
group of an extension of $P$ by $Z_k$ is
\[
 \Hom(P,Z_k)\simeq\Hom(\Phi,Z_k),
\]
which is plainly zero when $\Phi$ is a torsion group.  Likewise, if
$V=Z\otimes_{\mathbb Z}\mathbb Q$ and $\Phi$ is annihilated by the
integer $n>0$, then
\[
 \Ext^i(\Phi,V_k)=0\qquad\text{for every }i\in\mathbb Z,
\]
because the groups in question are $\mathbb Q$-vector spaces
annihilated by $n$.  The isomorphism (\SGARefLink{sga7.sga7i.eviii.n5-5-1}{5.5.1}) follows at once.

The same method proves, essentially, the following.

\medskip
\noindent\underline{Corollary \SGARefTarget{sga7.sga7i.eviii.n5-6}{5.6}.}  \underline{Let $k$ be a field,
$Z$ a torsion-free commutative group, $P$ and $Q$ two commutative
group schemes locally of finite type over $k$, and let
$\Phi=P/P^\circ$ and $\Psi=Q/Q^\circ$ be their groups of connected
components.  Then:}

\smallskip
\noindent (i) \underline{The natural functor from the category of
biextensions of $(\Phi,\Psi)$ by $Z_k$ to the category of
biextensions of $(P,Q)$ by $Z_k$ is an equivalence.  In particular,
if $P$ or $Q$ is connected, the category of biextensions of $(P,Q)$
by $Z_k$ is equivalent to the terminal category.}

\smallskip
\noindent (ii) \underline{Suppose that $\Phi$ or $\Psi$ is a torsion
group.  Then the category of biextensions of $(P,Q)$ by $Z_k$ is
rigid.  If, in addition, there is an integer $n>0$ such that
$n\id_\Phi=0$ or $n\id_\Psi=0$, then the group of classes of
biextensions of $(P,Q)$ by $Z_k$ is given by the canonical
isomorphism}
\begin{equation}\tag{\SGARefTarget{sga7.sga7i.eviii.n5-6-1}{5.6.1}}
 \Biext^1(P,Q;Z_k)
 \xleftarrow{\ \sim\ }
 \Hom\bigl(\Phi\otimes\Psi,(Z\otimes_{\mathbb Z}\mathbb Q)/Z\bigr),
\end{equation}
\underline{defined by the boundary operator in the exact sequence of
$\Biext^i$ groups $(i=0,1)$ associated with (\SGARefLink{sga7.sga7i.eviii.n5-5-2}{5.5.2}).}


% END INCLUDED SGA7I_English_current_buildable_source_20260823/components/114_expose_VIII_section_5_opening_through_5_6.tex

% BEGIN INCLUDED SGA7I_English_current_buildable_source_20260823/components/115_expose_VIII_sections_5_7_through_5_10.tex
% SOURCE AUTHORITY:
% expose_VIII_body.tex lines 1917--1974; scan indices 295--298;
% source folios 284--287.
% Direct scan comparison: physical PDF pages 296--299 in overlapping
% 1100-dpi bands.  The formulas, subscripts, and arrow directions are fully
% legible at that resolution; no higher-resolution escalation was needed.
% Transparent source dispositions:
% - the reference printed as (2.7.2) in 5.7 is the local formula (5.7.2);
% - stray G_{S,T} occurrences in 5.8 and 5.9 are corrected to the group
%   Z_{S,T} fixed by the surrounding definitions and formulas.
% Hard stop before section 6, source line 1976 / scan index 298.

\medskip
\noindent \SGARefTarget{sga7.sga7i.eviii.n5-7}{5.7}.  Let $S$ be a scheme, $T$ a closed subscheme, and $Z$
a set.  The \underline{constant scheme over $S$, truncated along $T$,
with value $Z$}, denoted $Z_{S,T}$, is the scheme obtained by taking
the sum of copies $S_z$ of $S$, one for each $z\in Z$, and gluing all
these copies along the open subscheme $U=S\setminus T$.  Thus
\begin{equation}\tag{\SGARefTarget{sga7.sga7i.eviii.n5-7-1}{5.7.1}}
 Z_{S,T}\times_S U\simeq U,
 \qquad
 Z_{S,T}\times_S T\simeq Z_T,
\end{equation}
and $Z_{S,T}$ is etale over $S$.  For a variable $S$-scheme $S'$
there is an isomorphism, functorial in $S'$,
\begin{equation}\tag{\SGARefTarget{sga7.sga7i.eviii.n5-7-2}{5.7.2}}
 \Hom_S(S',Z_{S,T})\simeq\Hom_T(S'_T,Z_T).
\end{equation}
The map is restriction to $S'_T$, using the second isomorphism in
(\SGARefLink{sga7.sga7i.eviii.n5-7-1}{5.7.1}); we leave it to the reader to check that the map is bijective.
If
\[
 i:T\longrightarrow S
\]
denotes the inclusion, the functorial isomorphism (\SGARefLink{sga7.sga7i.eviii.n5-7-2}{5.7.2}) may also be
written as an isomorphism of sheaves
\begin{equation}\tag{\SGARefTarget{sga7.sga7i.eviii.n5-7-3}{5.7.3}}
 Z_{S,T}\simeq i_*(Z_T).
\end{equation}
This formula, or the definition directly, shows that as $Z$ varies,
$Z_{S,T}$ depends functorially on $Z$, and that this functor commutes
with finite inverse limits.  In particular, it takes groups to group
schemes and commutative groups to commutative group schemes, as is also
clear from the form (\SGARefLink{sga7.sga7i.eviii.n5-7-2}{5.7.2}) of the functor represented by $Z_{S,T}$.

It is also immediate from the definition, or from (\SGARefLink{sga7.sga7i.eviii.n5-7-3}{5.7.3}), that the
formation of $Z_{S,T}$ commutes with every base change
$f:S'\to S$.  This gives an isomorphism, functorial in $Z$,
\begin{equation}\tag{\SGARefTarget{sga7.sga7i.eviii.n5-7-4}{5.7.4}}
 Z_{S,T}\times_S S'\simeq Z_{S',T'},
 \qquad T'=T\times_S S'=T_{S'}.
\end{equation}

\medskip
\noindent \SGARefTarget{sga7.sga7i.eviii.n5-8}{5.8}.  We shall study torsors $E$ under a group scheme of the
form $Z_{S,T}$, where $Z$ is a group, assumed commutative here for
simplicity.  There is a natural functor
\begin{equation}\tag{\SGARefTarget{sga7.sga7i.eviii.n5-8-0}{5.8.0}}
 E\longmapsto E_T
\end{equation}
from the category of torsors under $Z_{S,T}$ to the category of
torsors under $Z_T$.  We claim that this functor is an
\underline{equivalence of categories}.  Its full faithfulness means
that for two torsors $E,F$ under $Z_{S,T}$, the map
\[
 \Hom(E,F)\longrightarrow\Hom(E_T,F_T)
\]
is bijective.  By fpqc descent this reduces immediately to the case in
which $E$ and $F$ are trivial, hence to $E=F=Z_{S,T}$, where the
assertion is precisely the bijectivity of (\SGARefLink{sga7.sga7i.eviii.n5-7-2}{5.7.2}) for $S'=S$.

It remains to prove that (\SGARefLink{sga7.sga7i.eviii.n5-8-0}{5.8.0}) is essentially surjective.  Let $E_T$
be a torsor under $Z_T$.  We already observed in the proof of \SGARefLink{sga7.sga7i.eviii.n5-1}{5.1} that
$E_T$ is representable and locally trivial for the etale topology.  It
therefore suffices to prove that
\[
 H^1(S_{\mathrm{\acute{e}t}},Z_{S,T})
 \longrightarrow H^1(T_{\mathrm{\acute{e}t}},Z_T)
\]
is bijective.  This follows from (\SGARefLink{sga7.sga7i.eviii.n5-7-3}{5.7.3}) and the relation
\[
 R^1i_{\mathrm{\acute{e}t},*}(Z_T)=0,
\]
which is the special case SGA 4 VIII 5.6.

\medskip
\noindent \SGARefTarget{sga7.sga7i.eviii.n5-8-1}{5.8.1}.  Let now $P$ be a commutative group scheme over $S$.
We wish to determine the category of extensions, as sheaves of
commutative groups, of $P$ by $Z_{S,T}$.  For this we use the
description in VII \SGARefLink{sga7.sga7i.evii.n1-1-6}{1.1.6} and \SGARefLink{sga7.sga7i.evii.n1-3}{1.3} in terms of torsors on
$P$, $P\times P$, and $P\times P\times P$, under the inverse images
of $Z_{S,T}$.  By (\SGARefLink{sga7.sga7i.eviii.n5-7-4}{5.7.4}), these inverse images are again of the same
form as $Z_{S,T}$.  Applying the preceding result to those torsors
shows that they correspond to torsors on
$P_T$, $P_T\times_TP_T$, and
$P_T\times_TP_T\times_TP_T$ under the inverse images of $Z_T$.
We therefore obtain the following.

\medskip
\noindent\underline{Proposition \SGARefTarget{sga7.sga7i.eviii.n5-9}{5.9}.}  \underline{Let $S$ be a scheme,
$T$ a closed subscheme, $Z$ a commutative group, $Z_{S,T}$ the
corresponding truncated constant group scheme over $S$ defined in
\SGARefLink{sga7.sga7i.eviii.n5-7}{5.7}, and $P$ a commutative group scheme over $S$.  Then restriction to
$T$ induces an equivalence between the category of sheaves of
commutative groups which are extensions of $P$ by $Z_{S,T}$ and the
category of sheaves of commutative groups which are extensions of
$P_T$ by $Z_T$.}

An essentially identical argument gives the following.

\medskip
\noindent\underline{Corollary \SGARefTarget{sga7.sga7i.eviii.n5-10}{5.10}.}  \underline{With the notation of
\SGARefLink{sga7.sga7i.eviii.n5-9}{5.9}, let $Q$ be another commutative group scheme over $S$.  Then
restriction to $T$ induces an equivalence between the category of
biextensions of $(P,Q)$ by $Z_{S,T}$ and the category of biextensions
of $(P_T,Q_T)$ by $Z_T$.}


% END INCLUDED SGA7I_English_current_buildable_source_20260823/components/115_expose_VIII_sections_5_7_through_5_10.tex

% BEGIN INCLUDED SGA7I_English_current_buildable_source_20260823/components/116_expose_VIII_section_6_1_and_lemma_6_2.tex
% SOURCE AUTHORITY:
% expose_VIII_body.tex lines 1976--2003; scan indices 298--300;
% source folios 287--289.
% Direct scan comparison: physical PDF pages 299--301 in overlapping
% 1100-dpi bands.  The prose, formulas, and subscripts are fully legible at
% that resolution; no higher-resolution escalation was needed in this unit.
% Transparent source dispositions:
% - the stray \mathbb Z_Z in the paragraph preceding Lemma 6.2 is corrected
%   to \mathbb Z_S, the free group scheme over the stated base S;
% - the corrupt insertion between "over S" and \overline G in Lemma 6.2 is
%   omitted; the authority clearly defines alpha:S->\mathbb Z_{S,T} and then
%   pulls back the extension \overline G along alpha.
% Hard stop before 6.2.1, source line 2004 / scan index 300.

\medskip
\noindent\underline{6.  Extensions and biextensions by the Neron model of
$G_m$.}

\medskip
\noindent \SGARefTarget{sga7.sga7i.eviii.n6-1}{6.1}.  Let $S$ be a scheme, $T$ a closed subscheme, and $Z$ a set,
giving a truncated constant scheme $Z_{S,T}$ over $S$ (\SGARefLink{sga7.sga7i.eviii.n5-7}{5.7}).  There is a
canonical homomorphism, functorial in $Z$, of etale schemes over $S$,
\begin{equation}\tag{\SGARefTarget{sga7.sga7i.eviii.n6-1-1}{6.1.1}}
 Z_S\longrightarrow Z_{S,T},
\end{equation}
which is a \underline{surjective} morphism, as follows at once from (\SGARefLink{sga7.sga7i.eviii.n5-7-1}{5.7.1}).
Its functoriality in $Z$ shows that when $Z$ is a group, (\SGARefLink{sga7.sga7i.eviii.n6-1-1}{6.1.1}) is a
homomorphism of etale group schemes over $S$.  Its kernel is therefore an
etale group scheme over $S$, denoted $Z_{S,U}$, and there is a canonical
exact sequence, functorial in $Z$,
\begin{equation}\tag{\SGARefTarget{sga7.sga7i.eviii.n6-1-2}{6.1.2}}
 0\longrightarrow Z_{S,U}\longrightarrow Z_S\longrightarrow Z_{S,T}
 \longrightarrow 0.
\end{equation}
It follows immediately from (\SGARefLink{sga7.sga7i.eviii.n5-7-1}{5.7.1}) that $Z_{S,U}$ is the sum of a family of
$S$-schemes $S_n$, $n\in Z$, where $S_n=S$ for $n=0$ and $S_n=U$ for
$n\ne0$.  This description of $Z_{S,U}$ readily shows that for every group
$F$ over $S$, restriction to $U$ induces a bijection
\begin{equation}\tag{\SGARefTarget{sga7.sga7i.eviii.n6-1-3}{6.1.3}}
 \mathrm{Hom}_{\mathrm{gr}}(Z_{S,U},F)
 \xrightarrow{\ \sim\ }
 \mathrm{Hom}_{\mathrm{gr}}(Z_U,F|U).
\end{equation}

\medskip
\noindent \SGARefTarget{sga7.sga7i.eviii.n6-1-4}{6.1.4}.  Henceforth we suppose that $Z$ is commutative, and we wish
to examine extensions of commutative groups of $Z_{S,T}$ by a given
commutative group $G$ over $S$.  By the exact sequence (\SGARefLink{sga7.sga7i.eviii.n6-1-2}{6.1.2}), the category
of these extensions is equivalent to the category of extensions $F$ of
$Z_S$ by $G$, furnished with a trivialization of the inverse image of the
extension along $Z_{S,U}\to Z_S$.  By (\SGARefLink{sga7.sga7i.eviii.n6-1-3}{6.1.3}), such a trivialization amounts
to a trivialization of the extension $F|U$ of $Z_U$ by $G|U$.

For what follows we take $Z$ to be the group $\mathbb Z$ of integers.  Thus
$\mathbb Z_S$ is the free group on one generator, and VII \SGARefLink{sga7.sga7i.evii.n1-4}{1.4} gives an
equivalence between the category of extensions of $\mathbb Z_S$ by $G$ and
the category of torsors under $G$.  Applying the same result to
$\mathbb Z_U$ and $G|U$ gives the following.

\medskip
\noindent\underline{Lemma \SGARefTarget{sga7.sga7i.eviii.n6-2}{6.2}.}  \underline{Let $S$ be a scheme, $T$ a
closed subscheme, $U=S\setminus T$, and $G$ a commutative group over $S$.
The category of commutative groups over $S$ which are extensions of
$\mathbb Z_{S,T}$ (defined in (\SGARefLink{sga7.sga7i.eviii.n5-7}{5.7})) by $G$ is equivalent to the category of
pairs $(F,s)$ consisting of a torsor $F$ under $G$ and a section $s$ of
$F|U$.  If $\alpha:S\to\mathbb Z_{S,T}$ is the section of
$\mathbb Z_{S,T}$ which is the image of the section $1$ of $\mathbb Z_S$
over $S$, then the torsor corresponding to an extension $\overline G$ of
$\mathbb Z_{S,T}$ by $G$ is $F=\alpha^*(\overline G)$ under $G$, furnished
with the section of $F|U$ given by restricting the unit section of
$\overline G$ to $U$.}


% END INCLUDED SGA7I_English_current_buildable_source_20260823/components/116_expose_VIII_section_6_1_and_lemma_6_2.tex

% BEGIN INCLUDED SGA7I_English_current_buildable_source_20260823/components/117_expose_VIII_sections_6_2_1_through_6_2_4.tex
% SOURCE AUTHORITY:
% expose_VIII_body.tex lines 2004--2029; scan indices 300--302;
% source folios 289--291.
% Direct scan comparison: physical PDF pages 301--303 in overlapping
% 1100-dpi bands.  The mathematical text is fully legible at that resolution.
% Transparent source disposition: the printed "description 3.2" at the end
% of 6.2.1 is corrected to the immediately preceding description in 6.2.
% Hard stop before 6.3, source line 2031 / scan index 302.

\medskip
\noindent \SGARefTarget{sga7.sga7i.eviii.n6-2-1}{6.2.1}.  If $\overline G$ corresponds to the pair $(F,s)$, then the
inverse image $f_n^*(\overline G)$ in $\overline G$ of the section $f_n$ of
$\mathbb Z_{S,T}$ defined by the constant section of $\mathbb Z_S$ with
value $n\in\mathbb Z$ is canonically identified, as a torsor under $G$, with
the $n$-th power $F^{(n)}$ of the $G$-torsor $F$.  The canonical isomorphism
of $G_U$-torsors
\[
 f_n^*(\overline G)|U\simeq G_U,
\]
defined by the first relation in (\SGARefLink{sga7.sga7i.eviii.n5-7-1}{5.7.1}), is determined by the section
$s^{(n)}$ of $F_U^{(n)}$, the $n$-th power of the given section $s$ of
$F_U$.  This makes the following description of the sections of
$\overline G$ over $S$ explicit.  A section $\overline g$ of $\overline G$
over $S$ is determined by:

\smallskip
\noindent a) a section $\underline n$ of $\mathbb Z_{S,T}$ over $S$, which
may be interpreted as a section of $\mathbb Z_T$, that is, as a locally
constant function on $T$ with values in $\mathbb Z$, or equivalently as a
partition
\[
 T=\coprod_{i\in\mathbb Z}T_i
\]
of $T$ into a disjoint sum of open subsets indexed by $i\in\mathbb Z$;

\smallskip
\noindent b) a family of sections
\[
 \begin{gathered}
  \overline g_i\in\Gamma(F^{(i)}/S_i'),\qquad i\in\mathbb Z,\\
  S_i'=S\setminus T_i',\qquad
  T_i'=\displaystyle\bigcup_{\substack{j\in\mathbb Z\\j\ne i}}T_j,
 \end{gathered}
\]
subject to the following compatibility condition.  For $i,j\in\mathbb Z$,
$\overline g_i$ and $\overline g_j$ agree on $S_i'\cap S_j'=U$ after
$F^{(i)}|U$ and $F^{(j)}|U$ have been identified with $G_U$ by means of the
sections $s^{(i)}$ and $s^{(j)}$, respectively.  In other words, there must
exist a necessarily unique section $h_{ij}\in\Gamma(G/U)$ such that
\[
 \overline g_i=h_{ij}s^{(i)},\qquad
 \overline g_j=h_{ij}s^{(j)}.
\]
Since the description in \SGARefLink{sga7.sga7i.eviii.n6-2}{6.2} of an extension $\overline G$ of
$\mathbb Z_{S,T}$ by $G$ is plainly compatible with every base change
$S'\to S$, the preceding gives a complete description of the sheaf
$\overline G$ as a functor of the argument
$S'\in\operatorname{Ob}(\mathrm{Sch}_{/S})$.

\medskip
\noindent \SGARefTarget{sga7.sga7i.eviii.n6-2-2}{6.2.2}.  If $T$ is a scheme which is a sum of nonempty connected
schemes $T_i$, $i\in I$, the description in \SGARefLink{sga7.sga7i.eviii.n6-2-1}{6.2.1} of the sections of
$\overline G$ over $S$ may be simplified.  Such a section is determined by
a family of integers $(n_i)_{i\in I}$ indexed by $I$, and a family of
sections
\[
 \overline g_i\in\Gamma(F^{(n_i)}/S_i'),\qquad i\in I,
\]
where
\[
 S_i'=S\setminus T_i',
 \qquad
 T_i'=\coprod_{\substack{j\in I\\j\ne i}}T_j.
\]
The $\overline g_i$ satisfy the following compatibility condition: for every
pair $i,j\in I$, the restrictions of $\overline g_i$ and $\overline g_j$ to
$S_i'\cap S_j'=U$ agree under the identifications of $F^{(n_i)}|U$ and
$F^{(n_j)}|U$ deduced from the respective sections $s^{(n_i)}$ and
$s^{(n_j)}$ of these torsors.

\medskip
\noindent \SGARefTarget{sga7.sga7i.eviii.n6-2-3}{6.2.3}.  Finally, if $T$ is assumed nonempty and connected, giving a
section $\overline g$ of $\overline G$ over $S$ amounts to giving an integer
$n$ and a section $\overline g$ of $F^{(n)}$ over $S$.  There is no longer
any compatibility condition to impose, so the section $s$ of $F|U$ does not
occur in the description of the group of sections of $\overline G$ over
$S$.  (Of course, $s$ occurs again as soon as one seeks to describe the
points of $\overline G$ with values in a general $S$-scheme $S'$.)

\medskip
\noindent \SGARefTarget{sga7.sga7i.eviii.n6-2-4}{6.2.4}.  Since $\overline G$ is the union of the ``open
subschemes'' $f_n^*(\overline G)$, $n\in\mathbb Z$, described in \SGARefLink{sga7.sga7i.eviii.n6-2-1}{6.2.1}, it
follows that $\overline G$ is representable if and only if the torsors
$F^{(n)}$ are representable.  For example, it suffices that $G$ be
representable and affine over $S$, by the descent theory of SGA 1 VIII 2.1.

% END INCLUDED SGA7I_English_current_buildable_source_20260823/components/117_expose_VIII_sections_6_2_1_through_6_2_4.tex

% BEGIN INCLUDED SGA7I_English_current_buildable_source_20260823/components/118_expose_VIII_sections_6_3_and_6_4.tex
% SOURCE AUTHORITY:
% expose_VIII_body.tex lines 2031--2063; scan indices 302--304;
% source folios 291--293.
% Direct scan comparison: physical PDF pages 303--305 in overlapping
% 1100-dpi bands.  The mathematical text and displayed formulas are fully
% legible at that resolution.
% Transparent source disposition: formula (6.3.2) prints s_D next to
% O_S(-D), while the immediately following defining sentence correctly
% identifies the canonical rational section as s_{-D}; the English formula
% uses s_{-D}.
% Hard stop before Theorem 6.5, source line 2064 / scan index 304.

\medskip
\noindent \SGARefTarget{sga7.sga7i.eviii.n6-3}{6.3}.  In what follows we suppose that
$G=\underline G_{mS}$.  Every extension $\overline G$ of
$\mathbb Z_{S,T}$ by $G=\underline G_{mS}$ is then representable (\SGARefLink{sga7.sga7i.eviii.n6-2-4}{6.2.4}),
and is of course smooth over $S$, being an extension of the etale group
scheme $\mathbb Z_{S,T}$ by the smooth group scheme $\underline G_{mS}$.
Moreover, giving a torsor $F$ under $\underline G_m$ is equivalent to giving
an invertible Module $\mathbb L$ over $S$ (by associating to the latter the
torsor of invertible sections of $\mathbb L$).  Thus $\overline G$ is
described by a pair $(\mathbb L,s)$, where $\mathbb L$ is an invertible
Module over $S$ and $s$ is a trivialization of $\mathbb L$ over
$U=S\setminus T$, that is, an invertible section of $\mathbb L|U$.

When $U$ is schematically dense in $S$, and $S$ is locally noetherian, or is
reduced with a locally finite set of irreducible components, giving such a
pair is equivalent to giving a \underline{divisor $D$ on $S$ such that}
\begin{equation}\tag{\SGARefTarget{sga7.sga7i.eviii.n6-3-1}{6.3.1}}
 \operatorname{supp}D\subset T,
\end{equation}
by associating to such a divisor the pair
\begin{equation}\tag{\SGARefTarget{sga7.sga7i.eviii.n6-3-2}{6.3.2}}
 \bigl(\underline O_S(-D),s_{-D}|U\bigr)
\end{equation}
formed by the invertible Module $\underline O_S(-D)$ associated with $-D$
(EGA IV 21.2.8.1) and the restriction to $U$ of its canonical rational
section $s_{-D}$ (EGA IV 21.2.11 and 20.2.11(ii)).  In any event, without
any hypothesis on $S$, a divisor $D$ on $S$ satisfying (\SGARefLink{sga7.sga7i.eviii.n6-3-1}{6.3.1}) defines by
(\SGARefLink{sga7.sga7i.eviii.n6-3-2}{6.3.2}) a pair $(\mathbb L,s)$ and hence an extension $\overline G$ of
$\mathbb Z_{S,T}$ by $\underline G_m$.  We call it the
\underline{Neron--Raynaud model of $\underline G_m$} associated with the
divisor $D$ and the closed subscheme $T$ of $S$ satisfying (\SGARefLink{sga7.sga7i.eviii.n6-3-1}{6.3.1}).  If
$T=\operatorname{supp}D$, we omit $T$ from the terminology.

It is immediate that the formation of this extension $\overline G$ commutes
with every base change $f:S'\to S$ for which the inverse-image divisor
$f^*(D)$ is defined (EGA IV 21.4.2), and in particular with every
\underline{flat} base change.

\medskip
\noindent \SGARefTarget{sga7.sga7i.eviii.n6-3-3}{6.3.3}.  The general description in \SGARefLink{sga7.sga7i.eviii.n6-2-1}{6.2.1} of the sections of
$\overline G$ immediately shows that they may be identified with the
sections $f$ of the sheaf $\mathcal M_S^*$ of regular meromorphic sections
on $S$ (EGA IV 20.1.8) whose divisor $\operatorname{div}(f)$ is locally, for
the Zariski topology, an integral multiple of $D$.  If $V$ is an open subset
of $S$ on which the restriction of $\operatorname{div}(f)$ equals that of
$nD$, where $n\in\mathbb Z$, then the section of $\mathbb Z_{S,T}$ defined
by $f$ is simply the image in (\SGARefLink{sga7.sga7i.eviii.n6-1-2}{6.1.2}) of the constant section of
$\mathbb Z_S$ with value $n$.  The preceding description of the group of
sections of $\overline G$ over $S$ remains valid over every open subset of
$S$, and more generally for the group of points of $\overline G$ with
values in any flat $S$-scheme $S'$.  Since $\overline G$ itself is flat
(indeed smooth) over $S$, this description again gives a complete
characterization of the Neron--Raynaud model of $\underline G_m$ defined by
$D$ and $T\supset\operatorname{supp}D$.

\medskip
\noindent \SGARefTarget{sga7.sga7i.eviii.n6-4}{6.4}.  Suppose now that $S$ is \underline{normal}, that its set of
irreducible components is locally finite (so that $S$ is a sum of normal
integral schemes), and that $D$ is a positive divisor with all
multiplicities equal to $1$ (EGA IV 21.6), that $T=\operatorname{supp}D$,
and that $T$ is locally irreducible, that is, its irreducible components are
disjoint.  Then for every rational function $f$ on $S$ which is regular on
$U=S\setminus T$, the divisor $\operatorname{div}(f)$ is locally a multiple
of $D$.  Thus in this case there is a canonical isomorphism of groups
\begin{equation}\tag{\SGARefTarget{sga7.sga7i.eviii.n6-4-1}{6.4.1}}
 \Gamma(S,\overline G)\xrightarrow{\ \sim\ }
 \Gamma(U,\underline G_{mU}).
\end{equation}
Suppose further that $T$ is geometrically unibranch and locally integral, a
hypothesis stable under smooth base change (EGA IV 17.5.7).  Since all the
hypotheses are then stable under base change, (\SGARefLink{sga7.sga7i.eviii.n6-4-1}{6.4.1}) implies that
\underline{for every $S'$ smooth over $S$ there is a canonical group
isomorphism, functorial in $S'$,}
\begin{equation}\tag{\SGARefTarget{sga7.sga7i.eviii.n6-4-2}{6.4.2}}
 \mathrm{Hom}_S(S',\overline G)\xrightarrow{\ \sim\ }
 \mathrm{Hom}_U(S'_U,\underline G_{mU})
 \simeq\Gamma(S'_U,\underline O_{S'_U}^*).
\end{equation}
Since $\overline G$ is itself smooth over $S$, this isomorphism characterizes
the group scheme $\overline G$ over $S$ up to unique isomorphism.  It also
shows that $\overline G$ plays the role of a ``Neron model'' for the group
scheme $\underline G_{mU}$ over the open subset $U$ of $S$.  Together with
the fact that it was first introduced and used by M.~Raynaud, this justifies
the terminology adopted in \SGARefLink{sga7.sga7i.eviii.n6-3}{6.3} for $\overline G$.


% END INCLUDED SGA7I_English_current_buildable_source_20260823/components/118_expose_VIII_sections_6_3_and_6_4.tex

% BEGIN INCLUDED SGA7I_English_current_buildable_source_20260823/components/119_expose_VIII_theorem_6_5_through_corollary_6_7.tex
% SOURCE AUTHORITY:
% expose_VIII_body.tex lines 2064--2102; scan indices 304--307;
% source folios 293--296.
% Direct scan comparison: physical PDF pages 305--308 in overlapping
% 1100-dpi bands.  The theorem statements, formulas, and proof are fully
% legible at that resolution.
% Transparent source disposition: the proof of essential surjectivity in
% Theorem 6.5 prints "local rings of X" once; the stated base is S, and the
% English corrects X to S.
% Hard stop before Remarks 6.8, source line 2103 / scan index 307.

\medskip
\noindent\underline{Theorem \SGARefTarget{sga7.sga7i.eviii.n6-5}{6.5}.}  \underline{Let $S$ be a normal locally
integral scheme and let $D$ be a positive divisor on $S$, without multiple
components, whose support is geometrically unibranch and locally integral.
Let $\overline G$ be the Neron--Raynaud model of $\underline G_m$ over $S$
relative to $D$ (\SGARefLink{sga7.sga7i.eviii.n6-3}{6.3}), and consider the category of torsors $E$ over $S$
under the group scheme $\overline G$.  The restriction functor}
\[
 E\longmapsto E_U=E|U
\]
\underline{from this category to the category of torsors over $U$ under the
group scheme $\underline G_{mU}$ is fully faithful.  If $S$ is regular (more
generally, if its local rings are factorial), this functor is an equivalence
of categories.}

Let us prove that the functor in question is fully faithful, that is, that
for two torsors $E,F$ under $\overline G$, the restriction map
\begin{equation}\tag{*}
 \mathrm{Hom}(E,F)\longrightarrow\mathrm{Hom}(E_U,F_U)
\end{equation}
is bijective, where $\mathrm{Hom}$ denotes the set of homomorphisms of
\underline{torsors}.  First observe that this question is plainly
\underline{local on $S$ for the Zariski topology}.  Indeed, suppose that
$S$ is the union of open subsets $S_i$ such that for every open subset $S'$
of $S$ contained in some $S_i$ (and hence in particular for $S'=S_i$ or
$S_i\cap S_j$), the map
\[
 \mathrm{Hom}(E_{S'},F_{S'})
 \longrightarrow
 \mathrm{Hom}(E_{S'\cap U},F_{S'\cap U})
\]
is bijective.  Then (*) is bijective.  On the other hand, it follows at once
from \SGARefLink{sga7.sga7i.eviii.n5-8}{5.8} applied to $S,T,\mathbb Z$, and from \SGARefLink{sga7.sga7i.eviii.n5-1}{5.1} applied to $T,\mathbb Z$,
that every torsor under $\overline G$ is locally trivial for the Zariski
topology.  These two observations reduce us to proving the bijectivity of
(*) when $E=F=\overline G$.  In that case the map is precisely (\SGARefLink{sga7.sga7i.eviii.n6-4-1}{6.4.1}),
whose bijectivity has already been established.

We now prove that the functor $E\mapsto E_U$ is essentially surjective when
the local rings of $S$ are factorial.  Since $\underline G_{mS}$ is
identified with an open subgroup scheme of $\overline G$, it is enough to
show that every torsor under $\underline G_{mU}$ is isomorphic to the
restriction to $U$ of a torsor under $\underline G_{mS}$.  Equivalently,
every invertible Module over $U$ must be isomorphic to the restriction of an
invertible Module over $S$.  This follows from the hypothesis on $S$
(EGA IV 21.6.11, where only the locally noetherian case is treated; that is
the only case needed below).  This proves \SGARefLink{sga7.sga7i.eviii.n6-5}{6.5}.

\medskip
\noindent \SGARefTarget{sga7.sga7i.eviii.n6-5-1}{6.5.1}.  Suppose now, in addition, that $S$ is locally noetherian
and \underline{regular}.  Since $S$ is then locally factorial, giving $D$
is equivalent to giving a reduced, purely $1$-codimensional closed
subscheme, again denoted $D$, which is geometrically unibranch.  Notice also
that the hypotheses on $S$ and $D$ are stable under every smooth base change
$S'\to S$, so that the conclusion of \SGARefLink{sga7.sga7i.eviii.n6-5}{6.5} applies to $S'$ with the subscheme
$D'=D\times_S S'$.  The description in VII \SGARefLink{sga7.sga7i.evii.n1-1-6}{1.1.6} of extension groups now
immediately gives the following consequence.

\medskip
\noindent\underline{Corollary \SGARefTarget{sga7.sga7i.eviii.n6-6}{6.6}.}  \underline{Let $S$ be a locally
noetherian regular scheme, let $D$ be a reduced, geometrically unibranch,
purely $1$-codimensional closed subscheme (thus defining a divisor on $S$),
let $\overline G$ be the Neron--Raynaud model of $\underline G_m$ relative
to $D$, and let $P$ be a smooth commutative group scheme over $S$.  Consider
the category of commutative extensions, as sheaves of groups, of $P$ by
$\overline G$.  Restriction to $U$, $E\mapsto E|U$, induces an equivalence
between this category and the category of extensions of $P_U$ by
$\underline G_{mU}$.}

\medskip
\noindent \SGARefTarget{sga7.sga7i.eviii.n6-6-1}{6.6.1}.  The same method gives the analogous statement for
extensions which are not necessarily commutative (with $P$ no longer
assumed commutative).  We leave the details of the proof to the reader; they
are slightly more involved because bitorsors must be considered in place of
torsors.

By essentially the same proof as for \SGARefLink{sga7.sga7i.eviii.n6-6}{6.6}, one likewise obtains:

\medskip
\noindent\underline{Corollary \SGARefTarget{sga7.sga7i.eviii.n6-7}{6.7}.}  \underline{Let $S,D,\overline G$ be as
in \SGARefLink{sga7.sga7i.eviii.n6-6}{6.6}, and let $P,Q$ be two smooth commutative group schemes over $S$.
Consider the category of biextensions of $(P,Q)$ by $\overline G$.  Then the
restriction functor to $U$, $E\mapsto E|U$, from this category to the
category of biextensions of $(P_U,Q_U)$ by $\underline G_{mU}$ is an
equivalence of categories.}


% END INCLUDED SGA7I_English_current_buildable_source_20260823/components/119_expose_VIII_theorem_6_5_through_corollary_6_7.tex

% BEGIN INCLUDED SGA7I_English_current_buildable_source_20260823/components/120_expose_VIII_remarks_6_8_and_autocritique_6_9.tex
% SOURCE AUTHORITY:
% expose_VIII_body.tex lines 2103--2170; scan indices 307--310;
% source folios 296--299.
% Direct scan comparison: physical PDF pages 308--311 in overlapping
% 1100-dpi bands.  The prose and formulas are fully legible at that
% resolution.  The hand-corrected diagram on physical page 310 was also
% inspected directly in a 5000-dpi targeted crop.
% Transparent source dispositions:
% - 6.9(a) begins with j:V->S but consistently defines U=S-Y; V is corrected
%   to U, and j_*G_{mS} is corrected to j_*G_{mU};
% - the source diagram retains a faint obsolete right arrowhead on the Y--S
%   stroke beneath a heavy hand-added left arrowhead.  The correction gives
%   Y->S.  The mark at U is an open-immersion hook, and the source's terminal
%   period is retained.
% Hard stop before section 7, source line 2171 / scan index 310.

\medskip
\noindent\underline{Remarks \SGARefTarget{sga7.sga7i.eviii.n6-8}{6.8}.}\quad a) We shall apply \SGARefLink{sga7.sga7i.eviii.n6-6}{6.6} and \SGARefLink{sga7.sga7i.eviii.n6-7}{6.7} when
$S$ is a trait and $D$ is the divisor defined by the inclusion of the closed
point $s$ of $S$ into $S$.  In this case $U$ is the spectrum of the fraction
field $K$ of $S$, and $\overline G$ is the usual Neron model of the
multiplicative group over $K$, introduced by Raynaud.  Notice that even in
this special case, the proof of \SGARefLink{sga7.sga7i.eviii.n6-6}{6.6} and \SGARefLink{sga7.sga7i.eviii.n6-7}{6.7} uses \SGARefLink{sga7.sga7i.eviii.n6-5}{6.5} over bases more general
than traits (such as $P$ and its Cartesian powers).  Nevertheless, in this
case it is enough to know the \underline{existence} of the Neron model
$\overline G$ in the special situation under consideration: \SGARefLink{sga7.sga7i.eviii.n6-5}{6.5} for base
schemes smooth over $S$, and hence \SGARefLink{sga7.sga7i.eviii.n6-6}{6.6} and \SGARefLink{sga7.sga7i.eviii.n6-7}{6.7}, then follow formally.  For
us, the introduction of $\overline G$ will chiefly serve as a convenient
intermediate step in studying, in the next section, the problem of extending
extensions and biextensions by the multiplicative group.

\smallskip
\noindent b) A Neron--Raynaud model $\overline G$ of
$\underline G_{mU}$, satisfying the fundamental condition (\SGARefLink{sga7.sga7i.eviii.n6-4-2}{6.4.2}) for every
$S'$ smooth over $S$, can be defined under much more general conditions than
those considered above, with $U$ merely a schematically dense open subset
of a normal locally noetherian scheme $S$.  In this case $\overline G$ is an
fppf \underline{sheaf}, not necessarily representable.  Let $\underline D$
denote the sheaf, on the small etale site of $S$, of divisors on $S$ whose
support is contained in $T=S\setminus U$, and, by abuse of language, also
denote by the same letter its canonical extension to a sheaf on
$(\mathrm{Sch})_{/S}$.  Its sections over a relative scheme
$f:S'\to S$ are the sections of $f_{\mathrm{\acute{e}t}}^*(\underline D)$
over $S'$.  The desired model is an extension of groups
\begin{equation}\tag{\SGARefTarget{sga7.sga7i.eviii.n6-8-1}{6.8.1}}
 0\longrightarrow\underline G_{mS}\longrightarrow\overline G
 \longrightarrow\underline D\longrightarrow0.
\end{equation}
The construction of a canonical extension (\SGARefLink{sga7.sga7i.eviii.n6-8-1}{6.8.1}), without any normality
hypothesis on $S$, presents no difficulty when one uses the general
description VII \SGARefLink{sga7.sga7i.evii.n1-1-6}{1.1.6}; at the same time one obtains (\SGARefLink{sga7.sga7i.eviii.n6-4-2}{6.4.2}) for $S'$ etale
over $S$.  For the same relation to hold for every $S'$ smooth over $S$, it
is equivalent to verify that the formation of the sheaf $\underline D$
commutes with every smooth base change $f:S'\to S$.  For normal $S$, this
follows from EGA $\mathrm{Err}_{\mathrm{IV}}$ 53, 21.4.9.

\medskip
\noindent \SGARefTarget{sga7.sga7i.eviii.n6-9}{6.9}.\quad ``\underline{Self-criticism}'' (observations of
P.~Deligne).\footnote{The author of the present Expose having admitted
defeat, the reader is invited, if necessary, to rewrite sections 5 and 6
according to the accompanying directions.}

Sections 5 and 6 are really rather uninspiring.  One periodically has the
impression that the formalism of $i_*,i^*$ is being redone ``by hand'',
without seeing what is geometric and what is ``general nonsense''.

I propose extracting the ``geometric'' statements a), b), c):

\smallskip
\noindent a) (cf. \SGARefLink{sga7.sga7i.eviii.n6-8}{6.8} b)) Let $j:U\hookrightarrow S$ be a schematically
dense open subset of a normal noetherian scheme $S$.  The etale sheaf
$\underline G_{mS}$ is then identified with a subsheaf of the etale sheaf
$j_*\underline G_{mU}$.  Let $\underline D$ be the quotient etale sheaf in
the exact sequence
\[
 0\longrightarrow\underline G_{mS}\longrightarrow
 j_*\underline G_{mU}\longrightarrow\underline D\longrightarrow0.
\]
\underline{Then} the formation of $\underline D$ commutes with every smooth
base change (EGA $\mathrm{Err}_{\mathrm{IV}}$ 53, 21.4.9); \underline{moreover},
$R^1j_*\underline G_m=0$ (Theorem 90) if the strictly local rings of $S$ are
factorial.

\smallskip
\noindent b) If $U=S\setminus Y$, where $Y$ is a geometrically unibranch
divisor, and if $i$ is the inclusion of $Y$ into $S$, then
\[
 \underline D=i_*\mathbb Z_Y.
\]

\smallskip
% SGA-SRC-001273
\noindent c) If $Y$ is a closed subscheme of a scheme $X$, with closed immersion
$i:Y\hookrightarrow X$, and if $E$ is a set,
then the etale sheaf $i_*\underline E_Y$ is represented by an etale scheme
over $X$, whose formation commutes with every base change.

\medskip
\noindent\underline{Corollary.}  \underline{Let $S$ be a normal noetherian
scheme, let $Y$ be a geometrically unibranch divisor, and let
$U=S\setminus Y$:}
\begin{center}
\begin{tikzcd}[column sep=huge]
 U \arrow[r, hook, "j"'] & S & Y \arrow[l, "i"']
\end{tikzcd}
\qquad .
\end{center}
\underline{There exists a unique smooth group scheme $\mathcal G_m$ over
$S$ such that after every smooth base change one has
$\mathcal G_m|S_{\mathrm{\acute{e}t}}=j_*(\underline G_m)_U$.}

Indeed, the preceding requirements determine $\mathcal G_m$ as a sheaf, for
the etale topology, on the category of schemes smooth over $S$.  By a) and
b), the sequence
\[
 0\longrightarrow\underline G_m\longrightarrow\mathcal G_m
 \xrightarrow{\ q\ }i_*\mathbb Z_Y\longrightarrow0
\]
is exact.  By c), $i_*\mathbb Z_Y$ is representable.  Finally, the morphism
$q$ is representable.

Use of the preceding ``small smooth site'' should make it possible to deduce
the statements of section 6 formally from a).

% END INCLUDED SGA7I_English_current_buildable_source_20260823/components/120_expose_VIII_remarks_6_8_and_autocritique_6_9.tex

% BEGIN INCLUDED SGA7I_English_current_buildable_source_20260823/components/121_expose_VIII_section_7_opening_and_theorem_7_1_statement.tex
% SOURCE AUTHORITY:
% expose_VIII_body.tex lines 2171--2203; scan indices 310--311;
% source folios 299--300.
% Direct scan comparison: physical PDF pages 311--312 in overlapping
% 1100-dpi bands.  The prose, component-group notation, and obstruction
% formulas are fully legible; no higher-resolution escalation was needed.
% Transparent source dispositions:
% - typewritten letter-o special-fibre subscripts are normalized to zero;
% - in 7.1(b), the source twice drops the eta from E_eta in the extension
%   criterion; the English retains the stated generic-fibre object E_eta;
% - the coefficient is typeset coherently as \underline G_{mS}.
% Hard stop before the proof, source line 2204 / scan index 311.

\medskip
\noindent\underline{7.  Canonical extensions of extensions and
biextensions by $\underline G_m$.}

\medskip
\noindent \SGARefTarget{sga7.sga7i.eviii.n7-0}{7.0}.  We now reap the fruit of our preparations in sections 5 and
6.  First let the base scheme $S$ be a \underline{trait}, that is, the
spectrum of a discrete valuation ring; let $\eta$ be its generic point,
$s$ its closed point, and $k=k(s)$.  Let $P$ be a smooth commutative group
scheme over $S$, and write $P_0=P_s$ for its special fibre.  In the present
section we consider only commutative groups; in particular, all group
extensions under consideration are commutative.

\medskip
\noindent\underline{Theorem \SGARefTarget{sga7.sga7i.eviii.n7-1}{7.1}.}  \underline{The notation is that recalled
in \SGARefLink{sga7.sga7i.eviii.n7-0}{7.0}.}

\smallskip
\noindent a) \underline{Suppose that the group of connected components of
the special fibre of $P$,}
\[
 \Phi=P_0/P_0^{\circ},
\]
\underline{is a torsion group---as is the case, in particular, if $P$ is of
finite type over $S$.  Then the functor $E\mapsto E_\eta$ from the category
of extensions of $P$ by $\underline G_{mS}$ to the category of extensions
of $P_\eta$ by $\underline G_{m\eta}$ is fully faithful.  Suppose that there
is an integer $n>0$ such that $n\Phi=0$.  Then every extension $E_\eta$ of
$P_\eta$ by $\underline G_{m\eta}$ has an associated canonical
homomorphism}
\begin{equation}\tag{\SGARefTarget{sga7.sga7i.eviii.n7-1-1}{7.1.1}}
 d(E_\eta):\Phi\longrightarrow(\mathbb Q/\mathbb Z)_k,
\end{equation}
\underline{whose vanishing is necessary and sufficient for $E_\eta$ to be
isomorphic to the generic fibre of an extension of $P$ by
$\underline G_{mS}$.  In particular, if $P_0$ is connected, the preceding
functor $E\mapsto E_\eta$ is an equivalence of categories.}

\smallskip
\noindent b) \underline{Let $P$ and $Q$ be two smooth group schemes over
$S$, put}
\[
 \Phi=P_0/P_0^{\circ},\qquad \Psi=Q_0/Q_0^{\circ},
\]
\underline{and suppose that either $\Phi$ or $\Psi$ is a torsion group (for
example, that either $P$ or $Q$ is of finite type over $S$).  Then the
functor $E\mapsto E_\eta$ from the category of biextensions of $(P,Q)$ by
$\underline G_{mS}$ to the category of biextensions of
$(P_\eta,Q_\eta)$ by $\underline G_{m\eta}$ is fully faithful.  Suppose that
there is an integer $n>0$ such that $n\Phi=0$ or $n\Psi=0$.  Then every
biextension $E_\eta$ of $(P_\eta,Q_\eta)$ by $\underline G_{m\eta}$ has an
associated canonical pairing}
\begin{equation}\tag{\SGARefTarget{sga7.sga7i.eviii.n7-1-2}{7.1.2}}
 d(E_\eta):\Phi\otimes_{\mathbb Z}\Psi
 \longrightarrow(\mathbb Q/\mathbb Z)_k,
\end{equation}
\underline{whose vanishing is necessary and sufficient for $E_\eta$ to be
isomorphic to the generic fibre of a biextension $E$ of $(P,Q)$ by
$\underline G_{mS}$.  In particular, if either $P_0$ or $Q_0$ is connected,
the preceding functor $E\mapsto E_\eta$ is an equivalence of categories.}


% END INCLUDED SGA7I_English_current_buildable_source_20260823/components/121_expose_VIII_section_7_opening_and_theorem_7_1_statement.tex

% BEGIN INCLUDED SGA7I_English_current_buildable_source_20260823/components/122_expose_VIII_theorem_7_1_proof.tex
% SOURCE AUTHORITY:
% expose_VIII_body.tex lines 2204--2242; scan indices 311--313;
% source folios 300--302.
% Direct scan comparison: physical PDF pages 312--314 in overlapping
% 1100-dpi bands.  Every displayed Ext/Biext formula and tensor product is
% fully legible at that resolution.
% Transparent source dispositions:
% - part (a) prints "biextensions by \overline G" where Theorem 6.6 and the
%   entire argument concern extensions; the English uses extensions;
% - the final target printed Q/Z is the same (Q/Z)_k coefficient used in
%   (7.1.2), 5.6, and the preceding Ext calculation, so the English restores
%   (\mathbb Q/\mathbb Z)_k.
% Hard stop before Remark 7.2, source line 2243 / scan index 313.

\medskip
\noindent\underline{Proof.}\quad a) Let $\overline G$ be the
Neron--Raynaud model of $\underline G_m$.  The functor under consideration
is the composite of the functor ``extension of the kernel
$\underline G_{mS}\to\overline G$'' with the functor ``restriction to
$\eta$'' defined on the category of extensions by $\overline G$.  The latter
functor is an equivalence of categories by \SGARefLink{sga7.sga7i.eviii.n6-6}{6.6}, so it remains to study the
first.  For this purpose we use the exact sequence of group schemes
\begin{equation}\tag{\SGARefTarget{sga7.sga7i.eviii.n7-1-3}{7.1.3}}
 0\longrightarrow\underline G_{mS}\longrightarrow\overline G
 \longrightarrow\mathbb Z_{S,s}\longrightarrow0,
\end{equation}
where the notation $\mathbb Z_{S,s}$ is that of \SGARefLink{sga7.sga7i.eviii.n5-7}{5.7}.  The exact sequence of
$\mathrm{Ext}^i$, $i=0,1$, associated with (\SGARefLink{sga7.sga7i.eviii.n7-1-3}{7.1.3}) shows that the functor in
question is fully faithful if and only if
\begin{equation}\tag{\SGARefTarget{sga7.sga7i.eviii.n7-1-4}{7.1.4}}
 \mathrm{Hom}(P,\mathbb Z_{S,s})=0,
\end{equation}
and that the obstruction for an extension $\overline E$ of $P$ by
$\overline G$ to come from an extension $E$ of $P$ by
$\underline G_{mS}$ lies in
\begin{equation}\tag{\SGARefTarget{sga7.sga7i.eviii.n7-1-5}{7.1.5}}
 \mathrm{Ext}^1(P,\mathbb Z_{S,s}).
\end{equation}
By (\SGARefLink{sga7.sga7i.eviii.n5-7-2}{5.7.2}), the first group in (\SGARefLink{sga7.sga7i.eviii.n7-1-4}{7.1.4}) is isomorphic to
$\mathrm{Hom}(P_0,\mathbb Z_k)$, which in turn is isomorphic to
$\mathrm{Hom}(\Phi,\mathbb Z_k)$.  Thus (\SGARefLink{sga7.sga7i.eviii.n7-1-4}{7.1.4}) is equivalent to
\begin{equation}\tag{7.1.4 bis}
 \mathrm{Hom}(\Phi,\mathbb Z_k)=0.
\end{equation}
This relation certainly holds if $\Phi$ is a torsion group.  Moreover, if
$n\Phi=0$, then \SGARefLink{sga7.sga7i.eviii.n5-9}{5.9} and \SGARefLink{sga7.sga7i.eviii.n5-5}{5.5} give canonical isomorphisms
\begin{equation}\tag{\SGARefTarget{sga7.sga7i.eviii.n7-1-6}{7.1.6}}
 \mathrm{Ext}^1(P,\mathbb Z_{S,s})
 \simeq\mathrm{Ext}^1(P_0,\mathbb Z_k)
 \simeq\mathrm{Hom}\bigl(\Phi,(\mathbb Q/\mathbb Z)_k\bigr),
\end{equation}
which completes the proof of \SGARefLink{sga7.sga7i.eviii.n7-1}{7.1}(a).

\smallskip
\noindent b) The proof is essentially the same.  By \SGARefLink{sga7.sga7i.eviii.n6-7}{6.7}, we reduce to
studying the functor ``extension of the kernel
$\underline G_{mS}\to\overline G$'' from the category of biextensions of
$(P,Q)$ by $\underline G_{mS}$ to the category of biextensions of $(P,Q)$ by
$\overline G$.  We use the exact sequence (\SGARefLink{sga7.sga7i.eviii.n7-1-3}{7.1.3}) and the associated exact
sequence of $\mathrm{Biext}^i$, $i=0,1$ (VII \SGARefLink{sga7.sga7i.evii.n3-7-5}{3.7.5}).  Full faithfulness of
the functor is equivalent to
\begin{equation}\tag{\SGARefTarget{sga7.sga7i.eviii.n7-1-7}{7.1.7}}
 \mathrm{Hom}(P\otimes Q,\mathbb Z_{S,s})=0.
\end{equation}
By (\SGARefLink{sga7.sga7i.eviii.n5-7-2}{5.7.2}), the left side is isomorphic to
$\mathrm{Hom}(P_0\otimes Q_0,\mathbb Z_k)$ and hence to
$\mathrm{Hom}(\Phi\otimes\Psi,\mathbb Z_k)$; this group is indeed zero if
either $\Phi$ or $\Psi$ is a torsion group.  On the other hand, for a given
biextension $\overline E$ of $(P,Q)$ by $\overline G$, the obstruction to
its coming from a biextension $E$ of $(P,Q)$ by $\underline G_{mS}$ lies in
\[
 \mathrm{Biext}^1(P,Q;\mathbb Z_{S,s}).
\]
By \SGARefLink{sga7.sga7i.eviii.n5-10}{5.10} this is isomorphic to
$\mathrm{Biext}^1(P_0,Q_0;\mathbb Z_k)$, which, when $n\Phi=0$ or
$n\Psi=0$, is itself isomorphic by \SGARefLink{sga7.sga7i.eviii.n5-6}{5.6} to
\[
 \mathrm{Hom}\bigl(\Phi\otimes\Psi,(\mathbb Q/\mathbb Z)_k\bigr).
\]
This completes the proof of \SGARefLink{sga7.sga7i.eviii.n7-1}{7.1}.


% END INCLUDED SGA7I_English_current_buildable_source_20260823/components/122_expose_VIII_theorem_7_1_proof.tex

% BEGIN INCLUDED SGA7I_English_current_buildable_source_20260823/components/123_expose_VIII_remark_7_2.tex
% SOURCE AUTHORITY:
% expose_VIII_body.tex lines 2243--2261; scan indices 313--314;
% source folios 302--303.
% Direct scan comparison: physical PDF pages 314--315 in overlapping
% 1100-dpi bands.  The obstruction map, codimension bound, and footnote are
% fully legible; no higher-resolution escalation was needed.
% Transparent source disposition: spacing/typographical slips in "for every"
% and the products P x P x P are normalized without changing content.
% Hard stop before 7.3, source line 2262 / scan index 314.

\medskip
\noindent\underline{Remark \SGARefTarget{sga7.sga7i.eviii.n7-2}{7.2}.}  Given a smooth group scheme $P$ over a
connected regular locally noetherian scheme $S$ with generic point $\eta$,
one may ask to study the functor $E\mapsto E_\eta$ from the category of
extensions of $P$ by $\underline G_{mS}$ to the category of extensions of
$P_\eta$ by $\underline G_{m\eta}$.  This problem is readily reduced to the
one treated in \SGARefLink{sga7.sga7i.eviii.n7-1}{7.1}(a).  For simplicity, suppose that $P$ is of finite type
over $S$.  It follows easily from \SGARefLink{sga7.sga7i.eviii.n7-1}{7.1} that this functor is fully
faithful.\footnote{For this it is enough that $S$ be normal rather than
regular.}  One readily deduces that if $E_\eta$ is an extension of $P_\eta$
by $\underline G_{m\eta}$, then it is isomorphic to the generic fibre of an
extension $E$ of $P$ by $\underline G_{mS}$ if and only if, for every
$s\in S$ such that
\[
 \dim\operatorname{Spec}(\underline O_{S,s})=1,
\]
the analogous extension problem obtained by base change
$\operatorname{Spec}(\underline O_{S,s})\to S$ has a solution.  By \SGARefLink{sga7.sga7i.eviii.n7-1}{7.1}(a),
this is expressed by the vanishing of a certain homomorphism
\begin{equation}\tag{\SGARefTarget{sga7.sga7i.eviii.n7-2-1}{7.2.1}}
 d(E_\eta,s):P_s/P_s^{\circ}
 \longrightarrow(\mathbb Q/\mathbb Z)_{k(s)}.
\end{equation}

The preceding assertion plainly reduces to the following one.  Let $U$ be
an open subset of $S$ whose complement $T$ satisfies
\[
 \operatorname{codim}(T,S)\geq2.
\]
Then the functor $E\mapsto E|U$ from the category of extensions of $P$ by
$\underline G_{mS}$ to the category of extensions of $P_U$ by
$\underline G_{mU}$ is an equivalence of categories.  In terms of the
description VII \SGARefLink{sga7.sga7i.evii.n1-1-6}{1.1.6} this assertion is essentially trivial: for every
smooth $S$-scheme $X$ (such as $P$, $P\times P$,
$P\times P\times P$, and so on), $X$ is locally factorial and
$X_T=X\setminus X_U$ has codimension at least $2$ in $X$.  Consequently,
the functor $F\mapsto F|X_U$ from the category of torsors under
$\underline G_{mX}$ to the category of torsors under
$\underline G_{mX_U}$ is an equivalence of categories.

There are entirely analogous results for the problem of extending
biextensions; we leave their formulation to the reader.


% END INCLUDED SGA7I_English_current_buildable_source_20260823/components/123_expose_VIII_remark_7_2.tex

% BEGIN INCLUDED SGA7I_English_current_buildable_source_20260823/components/124_expose_VIII_section_7_3_1.tex
% SOURCE AUTHORITY:
% expose_VIII_body.tex lines 2262--2314; scan indices 314--315;
% source folios 303--304.
% Direct scan comparison: physical PDF pages 315--316 in overlapping
% 1100-dpi bands. Both triangular diagrams were checked node by node and
% arrow by arrow; their directions, label sides, and terminal punctuation
% are clear at this resolution, so no higher-resolution escalation was needed.
% Transparent source dispositions:
% - typewritten letter-o special-fibre subscripts are normalized to zero;
% - the coefficient of the pulled-back extension in 7.3.1 is restored from
%   the printed bare G_m to the contextually required G_{m\eta};
% - the malformed phrase "si on des morphismes" is rendered "if we have".
% Hard stop before 7.3.2, source line 2316 / scan index 315.

\medskip
\noindent \SGARefTarget{sga7.sga7i.eviii.n7-3}{7.3}.  We shall make explicit some functorial properties of the
obstructions defined in (\SGARefLink{sga7.sga7i.eviii.n7-1-1}{7.1.1}) and (\SGARefLink{sga7.sga7i.eviii.n7-1-2}{7.1.2}), with $P,Q,S$ variable, and
some relations among them.  In what follows, we shall suppose that the
conditions allowing these obstructions to be defined are satisfied.  Since
the verifications present no difficulty, we leave them to the reader.

\medskip
\noindent \SGARefTarget{sga7.sga7i.eviii.n7-3-1}{7.3.1}.  Let
\[
 u:P'\longrightarrow P
\]
be a morphism of smooth group schemes over $S$, and let $E_\eta$ be an
extension of $P_\eta$ by $\underline G_{m\eta}$.  Pullback by
$u_\eta:P'_\eta\to P_\eta$ gives an extension $E'_\eta$ of $P'_\eta$ by
$\underline G_{m\eta}$.  With this notation, the following diagram is
commutative:
\begin{equation}\tag{\SGARefTarget{sga7.sga7i.eviii.n7-3-1-1}{7.3.1.1}}
\begin{tikzcd}[row sep=large,column sep=huge]
 P'_0/{P'_0}^{\circ}
   \arrow[dd,"u_0"']
   \arrow[rd,"d(E'_\eta)"] & \\
 & (\mathbb Q/\mathbb Z)_k \\
 P_0/P_0^{\circ}
   \arrow[ru,"d(E_\eta)"'] &
\end{tikzcd}
\qquad .
\end{equation}

Similarly, if we have morphisms of smooth group schemes over $S$
\[
 u:P'\longrightarrow P,\qquad v:Q'\longrightarrow Q,
\]
and a biextension $E_\eta$ of $(P_\eta,Q_\eta)$ by
$\underline G_{m\eta}$, whose pullback by $(u_\eta,v_\eta)$ is a
biextension $E'_\eta$ of $(P'_\eta,Q'_\eta)$ by
$\underline G_{m\eta}$, then the following diagram is commutative:
\begin{equation}\tag{\SGARefTarget{sga7.sga7i.eviii.n7-3-1-2}{7.3.1.2}}
\begin{tikzcd}[row sep=large,column sep=huge]
 P'_0/{P'_0}^{\circ}\otimes Q'_0/{Q'_0}^{\circ}
   \arrow[dd,"u_0\otimes v_0"']
   \arrow[rd,"d(E'_\eta)"] & \\
 & (\mathbb Q/\mathbb Z)_k \\
 P_0/P_0^{\circ}\otimes Q_0/Q_0^{\circ}
   \arrow[ru,"d(E_\eta)"'] &
\end{tikzcd}
\end{equation}
The essentially mechanical verification of these compatibilities is left
to the reader.

% END INCLUDED SGA7I_English_current_buildable_source_20260823/components/124_expose_VIII_section_7_3_1.tex

% BEGIN INCLUDED SGA7I_English_current_buildable_source_20260823/components/125_expose_VIII_sections_7_3_2_and_7_3_3.tex
% SOURCE AUTHORITY:
% expose_VIII_body.tex lines 2316--2361; scan indices 315--317;
% source folios 304--306.
% Direct scan comparison: physical PDF pages 316--318 in overlapping
% 1100-dpi bands. The kernel formulas and the symmetry diagram were checked
% directly; all nodes, arrows, label sides, and the closing comma are clear.
% Transparent source dispositions:
% - typewritten letter-o special-fibre subscripts are normalized to zero;
% - the biextension denoted once by E rather than E_eta in 7.3.2 is rendered
%   consistently as the stated generic-fibre biextension E_eta;
% - "biextension symetrique" in the construction of {}^sE is rendered
%   "transposed biextension" to distinguish that operation from the later
%   property E_eta \simeq {}^sE_eta of being symmetric.
% Hard stop before 7.3.4, source line 2367 / scan index 317.

\medskip
\noindent \SGARefTarget{sga7.sga7i.eviii.n7-3-2}{7.3.2}.  It follows from the compatibility (\SGARefLink{sga7.sga7i.eviii.n7-3-1-1}{7.3.1.1}) that, if
$E_\eta$ is an extension of $P_\eta$ by $\underline G_{m\eta}$, there is a
largest group subscheme $V$ of $P$ such that the restriction of $E_\eta$ to
$V_\eta$ extends to an extension of $V$ by $\underline G_{mS}$.  It is the
open subgroup $V$ of $P$ containing $P_\eta$ and defined by the condition
\begin{equation}\tag{\SGARefTarget{sga7.sga7i.eviii.n7-3-2-1}{7.3.2.1}}
 V_0/V_0^{\circ}=\ker d(E_\eta).
\end{equation}
If $u:P'\to P$ is then a morphism of smooth group schemes over $S$, the
extension $E'_\eta$, obtained by pulling $E_\eta$ back along $u_\eta$,
extends to an extension of $P'$ by $\underline G_{mS}$ if and only if
\begin{equation}\tag{\SGARefTarget{sga7.sga7i.eviii.n7-3-2-2}{7.3.2.2}}
 u(P')\subset V.
\end{equation}

In the case of a biextension $E_\eta$ of $(P_\eta,Q_\eta)$ by
$\underline G_{m\eta}$, there is likewise a largest group subscheme $V$ of
$P$ such that the restriction of $E_\eta$ to $(V_\eta,Q_\eta)$ extends to
a biextension of $(V,Q)$ by $\underline G_{mS}$.  It is the open subgroup
of $P$ containing the generic fibre $P_\eta$ and characterized by the fact
that $V_0/V_0^{\circ}$ is the left annihilator in $\Phi$ of the pairing
(\SGARefLink{sga7.sga7i.eviii.n7-1-2}{7.1.2}), that is, the kernel of the corresponding homomorphism from $\Phi$
to $\underline{\mathrm{Hom}}(\Psi,(\mathbb Q/\mathbb Z)_k)$:
\begin{equation}\tag{\SGARefTarget{sga7.sga7i.eviii.n7-3-2-3}{7.3.2.3}}
 V_0/V_0^{\circ}
 =\ker\!\left(
   \Phi\xrightarrow{\ d(E_\eta)\ }
   \underline{\mathrm{Hom}}(\Psi,(\mathbb Q/\mathbb Z)_k)
 \right).
\end{equation}
If $u:P'\to P$ is a morphism of smooth group schemes over $S$, the
biextension of $(P'_\eta,Q_\eta)$ by $\underline G_{m\eta}$ obtained by
pulling $E_\eta$ back along $u_\eta$ extends to a biextension of $(P',Q)$
by $\underline G_{mS}$ if and only if (\SGARefLink{sga7.sga7i.eviii.n7-3-2-2}{7.3.2.2}) holds.  Symmetric results
hold upon interchanging the roles of $P$ and $Q$.

Of course, in general there is no largest \underline{pair} of subgroups
$P'\to P$, $Q'\to Q$ for which the restriction $E'_\eta$ of $E_\eta$ to a
biextension of $(P'_\eta,Q'_\eta)$ by $\underline G_{m\eta}$ extends to a
biextension of $(P',Q')$ by $\underline G_{mS}$.  One can only say that,
for a given pair $(P',Q')$, or more generally for a given pair of
morphisms $P'\to P$, $Q'\to Q$ of smooth group schemes, the desired
extension exists if and only if the restriction of the form $d(E_\eta)$
to
\[
 P'_0/{P'_0}^{\circ}\otimes Q'_0/{Q'_0}^{\circ}
\]
is zero.

\medskip
\noindent \SGARefTarget{sga7.sga7i.eviii.n7-3-3}{7.3.3}.  \underline{Symmetry}.  Let $E_\eta$ be a biextension of
$(P_\eta,Q_\eta)$ by $\underline G_{m\eta}$, and consider the transposed
biextension ${}^{s}E_\eta$ (VII \SGARefLink{sga7.sga7i.evii.n2-7}{2.7}), which is a biextension of
$(Q_\eta,P_\eta)$ by $\underline G_{m\eta}$.  The following diagram is
commutative:
\begin{equation}\tag{\SGARefTarget{sga7.sga7i.eviii.n7-3-3-1}{7.3.3.1}}
\begin{tikzcd}[row sep=large,column sep=huge]
 \Phi\otimes\Psi
   \arrow[dd,"\mathrm{sym}"']
   \arrow[rd,"d(E_\eta)"] & \\
 & (\mathbb Q/\mathbb Z)_k \\
 \Psi\otimes\Phi
   \arrow[ru,"d({}^{s}E_\eta)"'] &
\end{tikzcd}
\qquad ,
\end{equation}
where the vertical arrow is the canonical symmetry.  (Compare this
statement with the \underline{antisymmetry} statement of \SGARefLink{sga7.sga7i.eviii.n2-2-10}{2.2.10}.)  It
follows in particular that, if $P=Q$ and $E_\eta$ is a
\underline{symmetric} biextension of $(P_\eta,Q_\eta)$ by
$\underline G_{m\eta}$, that is, if $E_\eta$ is isomorphic to
${}^{s}E_\eta$, then the bilinear form $d(E_\eta)$ on $\Phi\times\Phi$
with values in $(\mathbb Q/\mathbb Z)_k$ is \underline{symmetric} (and not
alternating, as in \SGARefLink{sga7.sga7i.eviii.n2-2-12}{2.2.12}!).  If it is zero, that is, if $E_\eta$ extends
to a biextension $E$ of $(P,P)$ by $\underline G_{mS}$, the latter is
necessarily symmetric as well: by \SGARefLink{sga7.sga7i.eviii.n7-1}{7.1}(a), every isomorphism
${}^{s}E_\eta\to E_\eta$ extends uniquely to an isomorphism
${}^{s}E\to E$.

% END INCLUDED SGA7I_English_current_buildable_source_20260823/components/125_expose_VIII_sections_7_3_2_and_7_3_3.tex

% BEGIN INCLUDED SGA7I_English_current_buildable_source_20260823/components/126_expose_VIII_sections_7_3_4_and_7_3_5.tex
% SOURCE AUTHORITY:
% expose_VIII_body.tex lines 2367--2446; scan indices 317--321;
% source folios 306--310.
% Direct scan comparison: physical PDF pages 318--322 in overlapping
% 1100-dpi bands. Formulas, primes, residue-field subscripts, and the
% ramification index are clear at this resolution; no higher escalation was
% needed.
% Transparent source dispositions:
% - formula (7.3.4.1) is restored from the transcription's impossible E_*
%   to the source scan's q_eta^*(E_eta);
% - the family P_{eta Q_eta} is written explicitly as
%   P_eta x_eta Q_eta;
% - typewritten letter-o special-fibre subscripts are normalized to zero;
% - the malformed phrase "a l'aise du changement de base" is rendered
%   "by base change".
% Hard stop before Proposition 7.4, source line 2449 / scan index 321.

\medskip
\noindent \SGARefTarget{sga7.sga7i.eviii.n7-3-4}{7.3.4}.  \underline{Reduction of obstructions of type (\SGARefLink{sga7.sga7i.eviii.n7-1-2}{7.1.2}) to
those of type (\SGARefLink{sga7.sga7i.eviii.n7-1-1}{7.1.1}).}

Suppose, for definiteness, that there is an integer $n>0$ such that
$n\Phi=0$, and let
\[
 q\in Q(S)
\]
be a section of $Q$.  By an abuse of notation, write
\[
 q_0\in\Psi(k)
\]
for the image in $\Psi(k)$ of the value of $q$ at $s$.  Regarding $E_\eta$
as an extension of $P_\eta\times_\eta Q_\eta$ by
$\underline G_{mQ_\eta}$ over $Q_\eta$, pullback gives an extension
\[
 q_\eta^*(E_\eta),\qquad
 \text{an extension of $P_\eta$ by $\underline G_{m\eta}$},
\]
and hence an obstruction
\begin{equation}\tag{\SGARefTarget{sga7.sga7i.eviii.n7-3-4-1}{7.3.4.1}}
 d_q\overset{\mathrm{dfn}}{=}
 d\bigl(q_\eta^*(E_\eta)\bigr):
 \Phi\longrightarrow(\mathbb Q/\mathbb Z)_k.
\end{equation}
With $d=d(E_\eta)$, one has the relation
\begin{equation}\tag{\SGARefTarget{sga7.sga7i.eviii.n7-3-4-2}{7.3.4.2}}
 d_q(x)=d(x,q_0),
\end{equation}
valid for every point $x$ of $\Phi$ with values in a $k$-scheme $T$.

\smallskip
\noindent \SGARefTarget{sga7.sga7i.eviii.n7-3-4-3}{7.3.4.3}.  For every $q\in Q(S)$, this formula determines $d_q$
in terms of the pairing $d$; conversely, it plainly recovers $d$ from the
invariants $d_q$, at least when every point of $\Psi$ comes from a section
of $Q$.  By Hensel's lemma, this condition holds automatically when $S$ is
strictly local.  The pairing $d$ is known once it is known over the
separable closure of $k$, and its formation commutes with etale localization
on $S$ by \SGARefLink{sga7.sga7i.eviii.n7-3-5-3}{7.3.5.3} below.  Thus, after replacing $S$ by its strict
localization if necessary, the determination of $d$ reduces to that of the
invariants $d_q$.

In particular, when every point of $\Psi$ comes from a section $q$ of $Q$,
the pairing $d$ is zero---that is, $E_\eta$ extends to a biextension $E$ of
$(P,Q)$ by $\underline G_m$---if and only if $d_q$ is zero for every
section $q$ of $Q$, that is, if and only if $q_\eta^*(E_\eta)$ extends to
an extension of $P$ by $\underline G_m$.  With no hypothesis on $\Psi$,
the analogous criterion remains valid provided that the $q$ are allowed to
be points with values in local etale $S$-schemes $S'$.

\smallskip
\noindent \SGARefTarget{sga7.sga7i.eviii.n7-3-4-4}{7.3.4.4}.  Conversely, starting from the situation of \SGARefLink{sga7.sga7i.eviii.n7-1}{7.1}(a), the
obstruction (\SGARefLink{sga7.sga7i.eviii.n7-1-1}{7.1.1}) may be interpreted as an obstruction of type (\SGARefLink{sga7.sga7i.eviii.n7-1-2}{7.1.2}):
take $Q=\mathbb Z_S$ and the biextension of $(P_\eta,\mathbb Z_\eta)$
deduced from $E_\eta$.  Indeed, it suffices to apply the preceding
construction with $q$ equal to the section $1$ of $\mathbb Z_S$.

\medskip
\noindent \SGARefTarget{sga7.sga7i.eviii.n7-3-5}{7.3.5}.  \underline{Effect of a base change $S'\to S$.}  Suppose
given a dominant local morphism of traits
\[
 f:S'\longrightarrow S,
\]
and denote by $\eta'$ the generic point of $S'$, by $s'$ its closed point,
and put $k'=k(s')$.  If $S=\operatorname{Spec}(V)$ and
$S'=\operatorname{Spec}(V')$, then $V'$ is a discrete valuation ring
dominating the discrete valuation ring $V$.  Put
\begin{equation}\tag{\SGARefTarget{sga7.sga7i.eviii.n7-3-5-1}{7.3.5.1}}
 e(S'/S)=e(V'/V)=\operatorname{length}_{V'}(V'/\mathfrak m V'),
\end{equation}
where $\mathfrak m$ is the maximal ideal of $V$.

Let now $P$ be as in \SGARefLink{sga7.sga7i.eviii.n7-1}{7.1}(a), let $E_\eta$ be an extension of $P_\eta$ by
$\underline G_{m\eta}$, put $P'=P\times_S S'$, and let $E'_{\eta'}$ be
the extension of $P'_{\eta'}$ by $\underline G_{m\eta'}$ obtained by
pulling back $E_\eta$.  With the notation of \SGARefLink{sga7.sga7i.eviii.n4-1}{4.1}(a), there is a canonical
isomorphism
\[
 \Phi'\overset{\mathrm{dfn}}{=}P'_0/{P'_0}^{\circ}
 \simeq\Phi\otimes_k k'\simeq\Phi_{k'},
\]
and hence one may consider
\[
 d(E_\eta)_{k'}:\Phi'\longrightarrow(\mathbb Q/\mathbb Z)_{k'}.
\]
Then
\begin{equation}\tag{\SGARefTarget{sga7.sga7i.eviii.n7-3-5-2}{7.3.5.2}}
 d(E'_{\eta'})=e(S'/S)\,d(E_\eta)_{k'}.
\end{equation}

Likewise, under the hypotheses of \SGARefLink{sga7.sga7i.eviii.n7-1}{7.1}(b) and with the evident notation,
(\SGARefLink{sga7.sga7i.eviii.n7-3-5-2}{7.3.5.2}) holds with $d(E_\eta)$ now denoting the pairing (\SGARefLink{sga7.sga7i.eviii.n7-1-2}{7.1.2}),
$d(E_\eta)_{k'}$ the pairing
\[
 \Phi'\otimes\Psi'\longrightarrow(\mathbb Q/\mathbb Z)_{k'}
\]
obtained from it by base change $k\to k'$, and $E'_{\eta'}$ the
biextension of $(P'_{\eta'},Q'_{\eta'})$ by
$\underline G_{m\eta'}$ obtained from the biextension $E_\eta$ of
$(P_\eta,Q_\eta)$ by $\underline G_{m\eta}$ through the base change
$\eta'\to\eta$.

\smallskip
\noindent \SGARefTarget{sga7.sga7i.eviii.n7-3-5-3}{7.3.5.3}.  An interesting case is $e(S'/S)=1$, that is,
\[
 \mathfrak m'=\mathfrak m V',
\]
where $\mathfrak m'$ denotes the maximal ideal of $V'$.  This holds in
particular if $V'$ is etale over $V$, or if $V'$ is the henselization,
strict henselization, or completion of $V$.  In this case (\SGARefLink{sga7.sga7i.eviii.n7-3-5-2}{7.3.5.2}) simply
says that \underline{the formation of obstruction (\SGARefLink{sga7.sga7i.eviii.n7-1-1}{7.1.1}), respectively
(\SGARefLink{sga7.sga7i.eviii.n7-1-2}{7.1.2}), commutes with the base change under consideration $S'\to S$}.

\smallskip
\noindent \SGARefTarget{sga7.sga7i.eviii.n7-3-5-4}{7.3.5.4}.  On the other hand, (\SGARefLink{sga7.sga7i.eviii.n7-3-5-2}{7.3.5.2}) in the general case shows
that, for every extension $E_\eta$ as in \SGARefLink{sga7.sga7i.eviii.n7-1}{7.1}(a) (respectively every
biextension $E_\eta$ as in \SGARefLink{sga7.sga7i.eviii.n7-1}{7.1}(b)), there is a \underline{finite} morphism
of traits $S'\to S$ such that $k(\eta')$ is a separable extension of
$k(\eta)$ and such that the extension (respectively biextension)
$E'_{\eta'}$ obtained from $E_\eta$ by base change $\eta'\to\eta$ extends
to an extension (respectively biextension) $E'$ over all of $S'$.  Indeed,
it is enough to choose $S'/S$ so that $e(S'/S)$ is a multiple of the
integer $n>0$ occurring in \SGARefLink{sga7.sga7i.eviii.n7-1}{7.1}, and such an extension can always be found
with $k(\eta')$ separable over $k(\eta)$.

% END INCLUDED SGA7I_English_current_buildable_source_20260823/components/126_expose_VIII_sections_7_3_4_and_7_3_5.tex

% BEGIN INCLUDED SGA7I_English_current_buildable_source_20260823/components/127_expose_VIII_proposition_7_4_corollary_7_5_and_bibliography.tex
% SOURCE AUTHORITY:
% expose_VIII_body.tex lines 2449--2495; scan indices 321--323;
% source folios 310--312.
% Direct scan comparison: physical PDF pages 322--324 in overlapping
% 1100-dpi bands. The category names, product bases, superscripts, theorem
% underlining, and bibliography were all clear; no higher-resolution
% escalation was needed.
% Transparent source dispositions:
% - source spacing around Z-algebra is normalized without changing meaning;
% - bibliographic accents and punctuation follow the printed authority where
%   the transcription contained obvious keyboard errors.
% This component closes Expose VIII at the end of its bibliography.

\medskip
\noindent\underline{Proposition \SGARefTarget{sga7.sga7i.eviii.n7-4}{7.4}.  Let $S$ be a reduced scheme.}

\smallskip
\noindent a) \underline{Let $P$ be a smooth group scheme over $S$ whose
fibres at the maximal points of $S$ are connected.  Then the natural
functor}
\begin{equation}\tag{\SGARefTarget{sga7.sga7i.eviii.n7-4-1}{7.4.1}}
 \mathrm{EXT}(P,\underline G_{mS})
 \longrightarrow
 \mathrm{TORSRIG}(P,e_P;\underline G_{mS})
\end{equation}
\underline{from the category of extensions of $P$ by $\underline G_{mS}$
to the category of torsors under $\underline G_{mP}$ rigidified along the
identity section of $P$ is fully faithful.  If $S$ is normal and $P$ has
connected fibres, an object $E$ of the target of (\SGARefLink{sga7.sga7i.eviii.n7-4-1}{7.4.1}) belongs to the
essential image of this functor if and only if, for every maximal point
$\eta$ of $S$, the rigidified torsor $E_\eta$ on $P_\eta$ comes from an
extension of $P_\eta$ by $\underline G_{m\eta}$.}

\smallskip
\noindent b) \underline{Let $P,Q$ be two smooth group schemes over $S$
whose fibres at the maximal points of $S$ are connected.  Then the natural
functor}
\begin{equation}\tag{\SGARefTarget{sga7.sga7i.eviii.n7-4-2}{7.4.2}}
 \mathrm{BIEXT}(P,Q;\underline G_{mS})
 \longrightarrow
 \mathrm{TORSBIRIG}(P,Q;\underline G_{mS})
\end{equation}
\underline{from the category of biextensions of $(P,Q)$ by
$\underline G_{mS}$ to the category of torsors under
$\underline G_{m,P\times_S Q}$ birigidified along the pair of identity
sections of $P$ and $Q$ is fully faithful.  If $S$ is normal and either $P$
or $Q$ has connected fibres, an object $E$ of the target of (\SGARefLink{sga7.sga7i.eviii.n7-4-2}{7.4.2}) belongs
to the essential image of the functor if and only if, for every maximal
point $\eta$ of $S$, $E_\eta$ comes from a biextension of
$(P_\eta,Q_\eta)$ by $\underline G_{m\eta}$.}

We prove (a); the essentially identical proof of (b) is left to the reader.
For the full-faithfulness assertion, note that \SGARefLink{sga7.sga7i.eviii.n4-1}{4.1} remains valid if $X$ and
$Y$ are flat schemes over a reduced scheme $S$ whose fibres at the maximal
points satisfy the conditions stated in \SGARefLink{sga7.sga7i.eviii.n4-1}{4.1}.  Indeed, the family of maximal
fibres of $X\times_S Y$ (which is flat over the reduced scheme $S$) is
schematically dense in $X\times_S Y$.  Hence a morphism from
$X\times_S Y$ to the separated $S$-scheme $\underline G_{mS}$ is trivial
as soon as it is trivial on the maximal fibres.  Thus the full faithfulness
of (\SGARefLink{sga7.sga7i.eviii.n7-4-1}{7.4.1}) is proved as in \SGARefLink{sga7.sga7i.eviii.n4-2}{4.2}, and the characterization of its essential
image given there remains valid.

It remains to prove that, if the rigidified torsor $E$ on $P$ is such that
the $E_\eta$ come from extensions, then so does $E$, assuming that $S$ is
normal and $P$ has connected fibres.  The latter condition implies that
$P$ is of finite presentation over $S$ (SGA 3 $\mathrm{VI}_B$ 5.5).
Full faithfulness also shows that the question is local on $S$, so that we
may suppose $S$ affine and local.  The usual limit argument, using
Nagata's theorem on the finiteness of the integral closure of a reduced
finitely generated $\mathbb Z$-algebra (EGA IV 7.8), reduces us further to
the case in which $S$ is noetherian.

Then \SGARefLink{sga7.sga7i.eviii.n7-1}{7.1}(a) shows that there is an open subset $U\subset S$ whose complement
has codimension at least $2$ and such that $E_U$ comes from an extension of
$P_U$ by $\underline G_{mU}$.  In other words, the desired isomorphism
\[
 \pi^*(E)\simeq
 \operatorname{pr}_1^*(E)\operatorname{pr}_2^*(E)
\]
can be constructed over $(P\times_S P)_U$.  This is an open subset of the
normal scheme $P\times_S P$ whose complement has codimension at least $2$.
The isomorphism on $(P\times_S P)_U$ therefore extends over all of
$P\times_S P$, completing the proof.

\medskip
\noindent\underline{Corollary \SGARefTarget{sga7.sga7i.eviii.n7-5}{7.5}.  Let $S$ be a normal scheme and let $P$
and $Q$ be two smooth group schemes over $S$, with connected fibres at the
maximal points, either $P$ or $Q$ having connected fibres.  Suppose in
addition that, for every maximal point $\eta$ of $S$, either $P_\eta$ or
$Q_\eta$ is an extension of an abelian scheme by a torus (or, more
generally, admits a composition series whose factors are abelian schemes,
tori, or groups $\underline G_a$).  Then the functor (\SGARefLink{sga7.sga7i.eviii.n7-4-2}{7.4.2}) is an
equivalence of categories.}

This is an immediate consequence of \SGARefLink{sga7.sga7i.eviii.n7-4}{7.4}(b) and \SGARefLink{sga7.sga7i.eviii.n4-7}{4.7}.

\bigskip
\noindent\underline{BIBLIOGRAPHY}

\begin{enumerate}
\item[{[1]}] M. Demazure, J. Giraud, M. Raynaud,
  \emph{Sch\'emas ab\'eliens}, S\'eminaire Orsay 1967/68, forthcoming.
\item[{[2]}] M. Raynaud, \emph{Faisceaux amples sur les sch\'emas en groupes
  et les espaces homog\`enes}, Paris thesis, 1968 (forthcoming in Lecture
  Notes, Springer).
\item[{[3]}] M. Rosenlicht, \emph{Some basic theorems on algebraic groups},
  Amer. Jour. Math. \textbf{78}, no.~2 (1956), pp.~401--443.
\item[{[3bis]}] M. Rosenlicht, \emph{Some rationality questions on algebraic
  groups}, Annali di Matematica, Series 4, no.~43 (1957), pp.~25--50.
\item[{[4]}] J.-P. Serre, \emph{Corps locaux}, Act. Sci. Ind. 1296 (1962),
  Hermann, Paris.
\end{enumerate}

% END INCLUDED SGA7I_English_current_buildable_source_20260823/components/127_expose_VIII_proposition_7_4_corollary_7_5_and_bibliography.tex

\clearpage
\phantomsection\addcontentsline{toc}{section}{Exposé IX}
% BEGIN INCLUDED SGA7I_English_current_buildable_source_20260823/components/128_expose_IX_title_and_contents.tex
% SOURCE AUTHORITY:
% expose_IX_body.tex lines 5--64; scan indices 324--325; source folios 313--314.
% Direct scan comparison: physical PDF pages 325--326 at 1100 dpi.
% Hard stop before Introduction 0, source line 65.

\begin{center}
\SGARefAliasTarget{sga7.sga7i.eix.nroot}SGA 7

\vspace{1em}
\underline{Expos\'e IX}

\vspace{1em}
\underline{\underline{N\'ERON MODELS AND MONODROMY}}

\vspace{1em}
\underline{by A. Grothendieck}

(with an appendix by M. Raynaud)
\end{center}

\subsection*{Contents}

\begin{enumerate}[label=\arabic*.]
\setcounter{enumi}{-1}
\item Introduction
\item N\'eron model of abelian schemes: notation; the canonical pairing
  $\Phi_o\times\Phi'_o\longrightarrow(\mathbb Q/\mathbb Z)_k$, and the
  canonical biextension $W^o$ of $(A^o,{A'}^o)$
\item Fixed part and toric part of $T_\ell(A_K)$. Criteria for good
  reduction. Orthogonality theorem for $\ell\ne p$
\item The semistable-reduction case. The semistable reduction theorem
\item Application to a conjecture of Serre--Tate and to the conductor
\item The orthogonality theorem for $\ell=p$ (semistable case). Duality
  of the abelian schemes $B_o$ and $B'_o$. Characterization of the fixed
  part. Criteria for good reduction, essentially good reduction, and
  semistable reduction
\item Remarks on the construction of the toric part and the fixed part
  of $T_p(A)$ in the non-semistable-reduction case
\item The Raynaud extension $G^\natural$ over $S$ attached to the
  abelian scheme $A_K$ with semistable reduction. Duality of the abelian
  schemes $B,B'$ over $S$
\item The Raynaud extension $A_K^{\natural\,o}$ in the semistable case,
  and the ind-group $A_K^\natural$
\item Definition of the monodromy pairing
  $\underline M_\ell\otimes\underline M'_\ell\longrightarrow
  \mathbb Z_{\ell S}$
\item Properties of the monodromy pairing. Integrality and positivity
  theorem
\item Connected components of the N\'eron model: duality relation,
  asymptotic behavior
\item N\'eron model of a Jacobian and the Picard--Lefschetz formula
\item Relations with transcendental theory: the complex-analytic case
  \begin{enumerate}[label=\arabic{enumi}.\arabic*.]
  \item Review of analytic groups over $S$
  \item Formal completion along a fiber
  \item The analytic Raynaud extension
  \item Interpretation in terms of Deligne's ``mixed Hodge structures''
  \item The relatively algebraic case
  \end{enumerate}
\item Relations with transcendental theory: the rigid-analytic case.
  The canonical homomorphism
  $\underline M_K\longrightarrow A_K^{\natural\,o}$
\end{enumerate}


% END INCLUDED SGA7I_English_current_buildable_source_20260823/components/128_expose_IX_title_and_contents.tex

% BEGIN INCLUDED SGA7I_English_current_buildable_source_20260823/components/129_expose_IX_introduction_and_leitfaden.tex
% SOURCE AUTHORITY:
% expose_IX_body.tex lines 65--140; scan indices 325--329; source folios 314--318.
% Direct scan comparison: physical PDF pages 326--330 at 1100 dpi.
% Leitfaden checked from four direct-source crops at 5000 dpi; all 18 strokes
% are plain, unlabelled, undirected segments. Hard stop before source line 144.

0. \underline{Introduction}

\SGARefTarget{sga7.sga7i.eix.n0-1}{0.1}. In the present expos\'e we apply the results of the preceding
expos\'es to the study of the \underline{N\'eron model} (\SGARefLink{sga7.sga7i.eix.n1-1}{1.1}) $A$ of an
abelian scheme $A_K$ defined over the function field $K$ of the
Henselian trait $S$. We shall see that the most important properties of
the N\'eron model can be interpreted in terms of the action of the local
monodromy group $I$ on the Tate module
$T_\ell(A_K(\overline K))$ (where $\ell$ is a prime number distinct
from $p$). Since the dual of the latter can be interpreted as the
$\ell$-adic cohomology group $H^1(A_{\overline K},\mathbb Z_\ell)$,
and since, for every smooth projective scheme $X$ over $K$,
$H^1(X_{\overline K},\mathbb Z_\ell)$ is identified with
$H^1(A_{\overline K},\mathbb Z_\ell)$, where $A_{\overline K}$ is the
Albanese variety of $X_{\overline K}$, the present expos\'e may also be
viewed as a detailed study of \underline{local monodromy phenomena for
$\ell$-adic $H^1$} (or, better still, for ``motivic'' $H^1$)
\underline{of smooth projective varieties over} $K$. From this point of
view, it seems clear that the principal results of the present expos\'e
are destined to be subsumed within a ``N\'eron theory'' for motives of
arbitrary weight, that is, for $H^i$ ($\ell$-adic, de Rham, Hodge, and
so forth) with arbitrary $i$, of which at present we are only beginning
to catch a glimpse. (See [9] in this connection, and more particularly
Deligne's Conjectures 9.8 through 9.13 in the cited report.)

\SGARefTarget{sga7.sga7i.eix.n0-2}{0.2}. The most important result of the present expos\'e is the
``\underline{semistable reduction theorem}'' \SGARefLink{sga7.sga7i.eix.n3-6}{3.6}, which asserts that,
after a finite separable extension $K'$ of $K$, the N\'eron model of
$A_{K'}$ has a special fiber whose identity component is an extension
of an abelian scheme by a torus. This result was conjectured by
J.-P. Serre in 1964, and the principle of the proof given here
\footnote{It is developed in a letter from the lecturer to J.-P. Serre
dated 30 October 1964.}
goes back to the same year. It appears in \SGARefLink{sga7.sga7i.eix.n3-5}{3.5} that this property of
``semistable reduction'' is equivalent to the action of the monodromy
group $I$ on $T_\ell(A_K(\overline K))$ being ``essentially unipotent
of step 2'' (where $\ell$ is a fixed prime distinct from the residual
characteristic $p$ of $S$). In this form it was proved, in terms of
$\ell$-adic $H^1$, in Expos\'e III as a special case of the
``\underline{monodromy theorem}'', using resolution of singularities
for excellent schemes of dimension $2$ (due to S. Abhyankar). To prove
the equivalence of the two properties, we shall also need an
``orthogonality theorem'' \SGARefLink{sga7.sga7i.eix.n2-4}{2.4}, due to Igusa [11] in the special case in
which $A_K$ is the Jacobian of a smooth curve $X_K$ over $K$, the
generic fiber of a regular connected projective scheme $X$ over $S$
whose geometric special fiber has no singular points other than
ordinary quadratic singularities. This is indeed a semistable-reduction
case, but it would not suffice for our purposes; it is chiefly in order
to establish the orthogonality theorem in its general form that we use
the convenient tool supplied by the theory of biextensions developed
in the two preceding expos\'es. Let us also point out that D. Mumford
has found a very different proof of the semistable reduction theorem,
one which does not use resolution of singularities and instead uses
his theory of theta functions (unpublished); unfortunately, he must
restrict himself to residual characteristic different from $2$.

\SGARefTarget{sga7.sga7i.eix.n0-3}{0.3}. The semistable reduction theorem is the principal tool used in the
subsequent sections to study N\'eron models, which we undertake here
from the \underline{``algebraic'' point of view}. It also seems
indispensable for studying abelian varieties over $K$ from the
\underline{``rigid-analytic'' point of view}, or from the point of view
of ``theta functions'' (these terms apparently being synonymous when
abelian varieties over complete valued fields are in question), as
developed recently by D. Mumford and M. Raynaud following work of Tate
and Morikawa. (In Sections 13 and 14 we shall say a few words about the
relations between the algebraic and transcendental points of view.)
Finally, the same theorem is a key ingredient in the proof of the
``\underline{semistable reduction theorem for algebraic curves}'' of
Artin--Deligne--Mumford [3],\footnote{P. Deligne gave the proof of this
theorem in the oral Seminar, but it was not considered useful to
reproduce it in the Seminar texts.} which we have already had occasion
to use in the study of the action of monodromy on the fundamental group
(the ``essentially tame action theorem'' of Expos\'e V), and which is
used again in the proof of the integrality and positivity theorem \SGARefLink{sga7.sga7i.eix.n10-4}{10.4}
in the present expos\'e.\footnote{There is another proof of the latter
theorem, by way of rigid-analytic theory, which does not use the
semistable reduction theorem.}

\SGARefTarget{sga7.sga7i.eix.n0-4}{0.4}. The present expos\'e differs from the other expos\'es of the
Seminar in that, in studying the Tate module $T_\ell(A_K)$, we have
systematically taken into account also the case in which $\ell$ equals
the residual characteristic $p$. It turns out that the language of
\underline{Barsotti--Tate groups} [33]\footnote{Or ``$p$-divisible
groups'' in Tate's terminology.} makes it possible to extend to this
case all the most important phenomena established for $\ell\ne p$,
subject on occasion to having available Tate's fundamental theorem
[33, Theorem 4], which at present has been established only when $K$
has characteristic zero. Of course, in the latter case the pro-$p$
Barsotti--Tate group $T_p(A_K)$ may still be interpreted in terms of
the action of the monodromy group $I$ on the $p$-adic module
$T_p(A_K(\overline K))$, but one should bear in mind that most of the
statements suggested by this superficial analogy with the case
$\ell\ne p$ are grossly false. (Thus it is no longer true that the
monodromy action is essentially unipotent when $\ell=p$.) By using more
judicious translations, however (such as \SGARefLink{sga7.sga7i.eix.n5-13}{5.13} below), it is already
clear that the study of pro-$p$ Barsotti--Tate groups defined
geometrically over the function field of a trait of residual
characteristic $p>0$, whether or not the trait has mixed
characteristic, must be regarded as part of the theory of ``local
monodromy'', in which these groups play the role of \underline{$p$-adic
local systems} on $\operatorname{Spec}(K)$. In fact, these are
essentially the ``local systems'' that enter into the \underline{study,
from the ``$p$-adic'' point of view, of $H^1$ of smooth projective
varieties over} $K$. For the still-to-be-developed study of higher
$H^i$ from the $p$-adic point of view, it appears that the category of
pro-$p$ Barsotti--Tate groups over $K$, respectively over $S$, will
have to be suitably enlarged, taking inspiration from the viewpoints
provided by ``crystalline cohomology''.

\SGARefTarget{sga7.sga7i.eix.n0-5}{0.5}. The expos\'e given here follows the oral lectures very closely up
to Section 4, but is markedly more detailed and complete from the
following sections onward, which had been only partially sketched. It
will not be used in the subsequent expos\'es. Some results of the
present expos\'e (\SGARefLink{sga7.sga7i.eix.n10-4}{10.4} and \SGARefLink{sga7.sga7i.eix.n12-5}{12.5}), on the other hand, use the
\underline{Picard--Lefschetz formula}, which will be established only
later.

\begin{center}
\underline{LEITFADEN}

\vspace{0.5em}
\begin{tikzpicture}[line width=0.35pt,every node/.style={fill=white,inner sep=1.2pt}]
  \coordinate (p1)  at ( 0.00, 0.00);
  \coordinate (p2)  at ( 0.00,-0.76);
  \coordinate (p3)  at ( 0.00,-1.53);
  \coordinate (p4)  at (-1.60,-1.92);
  \coordinate (p5)  at ( 0.00,-2.30);
  \coordinate (p6)  at (-1.62,-2.69);
  \coordinate (p7)  at ( 1.32,-2.87);
  \coordinate (p8)  at ( 0.04,-3.35);
  \coordinate (p9)  at ( 3.02,-2.69);
  \coordinate (p10) at ( 3.10,-3.35);
  \coordinate (p11) at ( 2.66,-4.21);
  \coordinate (p12) at ( 3.76,-4.21);
  \coordinate (p13) at ( 0.12,-4.37);
  \coordinate (p14) at ( 1.04,-4.37);
  \draw (p1)--(p2) (p2)--(p3) (p3)--(p4) (p3)--(p5)
        (p5)--(p6) (p5)--(p7) (p5)--(p8) (p5)--(p9)
        (p7)--(p8) (p7)--(p13) (p7)--(p14)
        (p8)--(p13) (p8)--(p14) (p9)--(p10)
        (p10)--(p11) (p10)--(p12) (p10)--(p13) (p10)--(p14);
  \node at (p1)  {1};  \node at (p2)  {2};
  \node at (p3)  {3};  \node at (p4)  {4};
  \node at (p5)  {5};  \node at (p6)  {6};
  \node at (p7)  {7};  \node at (p8)  {8};
  \node at (p9)  {9};  \node at (p10) {10};
  \node at (p11) {11}; \node at (p12) {12};
  \node at (p13) {13}; \node at (p14) {14};
\end{tikzpicture}
\end{center}


% END INCLUDED SGA7I_English_current_buildable_source_20260823/components/129_expose_IX_introduction_and_leitfaden.tex

% BEGIN INCLUDED SGA7I_English_current_buildable_source_20260823/components/130_expose_IX_section_1_0_and_1_1.tex
% Expose IX, 1.0--1.1.
% French transcription lines 144--260; authority scan physical pp. 331--334
% (source folios 319--322).  Ordinary pages were checked in four overlapping
% 1100-dpi bands.  The missing prime on A'_o in the final sentence of 1.1 is
% an ordinary printed/source slip and is restored below.

\medskip
\noindent
\underline{1. Neron models of abelian schemes: notation; the canonical pairing}
\[
 \Phi_o\times\Phi'_o\longrightarrow(\mathbb Q/\mathbb Z)_k,
\]
\underline{and the canonical biextension $W^o$ of $(A^o,{A'}^o)$.}

\medskip
\noindent
\SGARefTarget{sga7.sga7i.eix.n1-0}{1.0}. Throughout this expos\'e we fix a trait $S$, for which we use the
usual notation $s$, $\eta$, $k$, $K$, and so forth, of I~1.  Let $A_K$
be an abelian scheme over $K$, and denote its dual abelian scheme [5] by
$A'_K$.  The duality between $A_K$ and $A'_K$ is given by a divisorial
correspondence $W_K$ on $(A_K,A'_K)$, that is, an invertible sheaf on
$A_K\times A'_K$, birigidified with respect to the identity sections of
$A_K$ and $A'_K$.  It is determined up to unique isomorphism by its
class (cf. VIII~\SGARefLink{sga7.sga7i.eviii.n3-1}{3.1})
\begin{equation}\tag{\SGARefTarget{sga7.sga7i.eix.n1-0-1}{1.0.1}}
 w_K\in\mathrm{Biext}^1(A_K,A'_K;\mathbb G_{mK}).
\end{equation}
An important role will be played by the corresponding pairing
(VIII~\SGARefLink{sga7.sga7i.eviii.n2-2}{2.2})
\begin{equation}\tag{\SGARefTarget{sga7.sga7i.eix.n1-0-2}{1.0.2}}
 \varphi:T_\ell(A_K)\times T_\ell(A'_K)
 \longrightarrow\mathbb Z_\ell(1)_K
 \overset{\mathrm{def}}{=}T_\ell(\mathbb G_{mK}),
\end{equation}
where $\ell$ is an arbitrary prime number.  When
$\ell\ne p=\operatorname{char}k$, giving $\varphi$ is equivalent to
giving a pairing of finite free $\mathbb Z_\ell$-modules
\begin{equation}\tag{\SGARefTarget{sga7.sga7i.eix.n1-0-3}{1.0.3}}
 \varphi:T_\ell(A_K(\overline K))\times
 T_\ell(A'_K(\overline K))\longrightarrow T_\ell(\overline K^*),
\end{equation}
compatible with the action of
\[
 \pi=\operatorname{Gal}(\overline K/K)
\]
on the three modules in question, where $\overline K$ is a separable
closure of $K$.  It is known (VIII~\SGARefLink{sga7.sga7i.eviii.n3-2}{3.2}) that this pairing is a perfect
duality: it identifies each of the modules
$T_\ell(A_K(\overline K))$ and $T_\ell(A'_K(\overline K))$ with the
dual of the other, with values in the invertible module
$\mathbb Z_\ell(1)=T_\ell(\overline K^*)$.

At times we shall suppose that a polarization of $A_K$ has been chosen;
thus we are given an element
\begin{equation}\tag{\SGARefTarget{sga7.sga7i.eix.n1-0-4}{1.0.4}}
 \xi\in\operatorname{NS}(A_K/K)
 =\underline{\operatorname{NS}}_{A_K/K}(K),
\end{equation}
where
\[
 \underline{\operatorname{NS}}_{A_K/K}
 \overset{\mathrm{def}}{=}
 \underline{\operatorname{Pic}}_{A_K/K}/
 \bigl(\underline{\operatorname{Pic}}{}^o_{A_K/K}=A'_K\bigr).
\]
In the usual way it defines a homomorphism
\begin{equation}\tag{\SGARefTarget{sga7.sga7i.eix.n1-0-5}{1.0.5}}
 \psi_\xi:A_K\longrightarrow A'_K,
\end{equation}
hence a homomorphism
\begin{equation}\tag{\SGARefTarget{sga7.sga7i.eix.n1-0-6}{1.0.6}}
 T_\ell(\psi_\xi),\quad\text{or simply}\quad
 \psi_\xi:T_\ell(A_K)\longrightarrow T_\ell(A'_K),
\end{equation}
and a bilinear form $\varphi_\xi$
\begin{equation}\tag{\SGARefTarget{sga7.sga7i.eix.n1-0-7}{1.0.7}}
 \begin{split}
 \varphi_\xi(x,y)&=\varphi(x,\psi_\xi(y)),\\
 \varphi_\xi:T_\ell(A_K)\times T_\ell(A_K)
 &\longrightarrow\mathbb Z_\ell(1)_K,
 \end{split}
\end{equation}
which, for $\ell\ne p$, may also be interpreted as a pairing of
finite free $\mathbb Z_\ell$-modules with operators:
\begin{equation}\tag{\SGARefTarget{sga7.sga7i.eix.n1-0-8}{1.0.8}}
 \varphi_\xi:T_\ell(A_K(\overline K))\times
 T_\ell(A_K(\overline K))\longrightarrow T_\ell(\overline K^*).
\end{equation}
These constructions make sense for every element of the
Neron--Severi group of $A_K$, and yield an \underline{alternating} form
$\varphi_\xi$, in either interpretation (\SGARefLink{sga7.sga7i.eix.n1-0-7}{1.0.7}) or (\SGARefLink{sga7.sga7i.eix.n1-0-8}{1.0.8})
(VIII~\SGARefLink{sga7.sga7i.eviii.n2-2-11}{2.2.11}).  The fact that $\xi$ is a \underline{polarization}
implies that $\psi_\xi$ in (\SGARefLink{sga7.sga7i.eix.n1-0-5}{1.0.5}) is an isogeny; equivalently, since
$\dim A_K=\dim A'_K$, it is an epimorphism of groups.  This is also
equivalent to the injectivity of $T_\ell(\psi_\xi)$ in (\SGARefLink{sga7.sga7i.eix.n1-0-6}{1.0.6}) (or to
its having finite cokernel), and finally to the form $\varphi_\xi$ in
(\SGARefLink{sga7.sga7i.eix.n1-0-7}{1.0.7}) being \underline{separating}.

\medskip
\noindent
\SGARefTarget{sga7.sga7i.eix.n1-1}{1.1}. Recall that there exists a \underline{smooth scheme} $A$ over
$S$, called the \underline{Neron model of $A_K$}, which represents the
functor
\begin{equation}\tag{\SGARefTarget{sga7.sga7i.eix.n1-1-1}{1.1.1}}
 S'\longmapsto\operatorname{Hom}(S'_\eta,A_K)
\end{equation}
on the category of smooth $S$-schemes.  Thus, by definition, there is
an isomorphism, functorial in the smooth relative $S$-scheme $S'$,
\begin{equation}\tag{\SGARefTarget{sga7.sga7i.eix.n1-1-2}{1.1.2}}
 \operatorname{Hom}_S(S',A)\simeq
 \operatorname{Hom}_K(S'_\eta,A_K).
\end{equation}
Neron [19] constructed $A$ when the residue field $k$ of $S$ is
perfect.  Raynaud extended the construction to the general case in an
unwritten seminar held at IHES in 1966--67 (see [21] for the statement
of Raynaud's principal results, which complete those of Neron in
loc. cit.).  We shall admit this existence result here, together with
the fact that $A$ is a quasi-projective \underline{scheme} over $S$.
Applying (\SGARefLink{sga7.sga7i.eix.n1-1-2}{1.1.2}) when $S'$ is smooth \underline{over} $\eta$, one also
finds that the following canonical morphism is an isomorphism:
\begin{equation}\tag{\SGARefTarget{sga7.sga7i.eix.n1-1-3}{1.1.3}}
 A\times_S\operatorname{Spec}(K)\xrightarrow{\ \sim\ }A_K.
\end{equation}
This justifies the paired notation $A,A_K$; when convenient we shall
also write $A_\eta$ in place of $A_K$.  Using (\SGARefLink{sga7.sga7i.eix.n1-1-3}{1.1.3}) as an identifying
isomorphism, the canonical bijection (\SGARefLink{sga7.sga7i.eix.n1-1-2}{1.1.2}) is simply the map given by
\underline{restriction to the generic fibers}.

Since the functor (\SGARefLink{sga7.sga7i.eix.n1-1-1}{1.1.1}) plainly has the structure of a commutative
group, the Neron model $A$ representing it itself has the structure of
a \underline{smooth commutative group scheme over $S$}.  The remainder
of this expos\'e is devoted to the study of this group scheme.

An important invariant arising from the Neron model is, of course, the
smooth commutative group scheme of finite type over $k$
\begin{equation}\tag{\SGARefTarget{sga7.sga7i.eix.n1-1-4}{1.1.4}}
 A_o=A\times_Ss=A_k.
\end{equation}
Following the general notation, we denote its identity component by
$A_o^o$ and also consider the group scheme
\begin{equation}\tag{\SGARefTarget{sga7.sga7i.eix.n1-1-5}{1.1.5}}
 \Phi_o=A_o/A_o^o
\end{equation}
of connected components of $A_o$, which is finite and etale over $k$.
In terms of a separable closure $\overline k$ of $k$, it is therefore
determined by the ordinary finite group
\begin{equation}\tag{\SGARefTarget{sga7.sga7i.eix.n1-1-6}{1.1.6}}
 \Phi_o(\overline k),
\end{equation}
together with the natural action of
\[
 \pi_o=\operatorname{Gal}(\overline k/k)
\]
on this group.

Along with the Neron model $A$ of $A_K$, we shall study
\begin{equation}\tag{\SGARefTarget{sga7.sga7i.eix.n1-1-7}{1.1.7}}
 A'=\text{the Neron model of }A'_K,
\end{equation}
which is likewise a smooth commutative group scheme of finite type over
$S$.  We shall denote by
\begin{equation}\tag{\SGARefTarget{sga7.sga7i.eix.n1-1-8}{1.1.8}}
 \Phi'_o=A'_o/{A'}_o^o
\end{equation}
the group of connected components of its special fiber $A'_o$.

% END INCLUDED SGA7I_English_current_buildable_source_20260823/components/130_expose_IX_section_1_0_and_1_1.tex

% BEGIN INCLUDED SGA7I_English_current_buildable_source_20260823/components/131_expose_IX_sections_1_2_through_1_5.tex
% Expose IX, 1.2--1.5.
% French transcription lines 265--341; authority scan physical pp. 335--338
% (source folios 323--326).  Ordinary pages were checked in four overlapping
% 1100-dpi bands.  Formula (1.4.2) was checked in targeted 5000-dpi crops.

\medskip
\noindent
\SGARefTarget{sga7.sga7i.eix.n1-2}{1.2}. Apply VIII~\SGARefLink{sga7.sga7i.eviii.n7-1}{7.1}(b) to the pair $(A,A')$ and to the biextension
$W_K$ of $(A_K,A'_K)$ by $\mathbb G_{mK}$ in (\SGARefLink{sga7.sga7i.eix.n1-0-1}{1.0.1}).  We thereby
obtain a \underline{canonical pairing}
\begin{equation}\tag{\SGARefTarget{sga7.sga7i.eix.n1-2-1}{1.2.1}}
 \Phi_o\times\Phi'_o\longrightarrow(\mathbb Q/\mathbb Z)_k,
\end{equation}
that is, a pairing of $\pi_o$-modules
\begin{equation}\tag{\SGARefTarget{sga7.sga7i.eix.n1-2-2}{1.2.2}}
 \Phi_o(\overline k)\times\Phi'_o(\overline k)
 \longrightarrow\mathbb Q/\mathbb Z.
\end{equation}
Here the pairing appears as the obstruction to extending $W_K$ to a
biextension $W$ of $(A,A')$ by $\mathbb G_{mS}$; the meaning of this
point of view is explained in VIII~\SGARefLink{sga7.sga7i.eviii.n7-3-2}{7.3.2}.  We record at once the
following conjecture concerning this pairing.

\medskip
\noindent
\underline{Conjecture \SGARefTarget{sga7.sga7i.eix.n1-3}{1.3}.} \underline{The pairing} (\SGARefLink{sga7.sga7i.eix.n1-2-1}{1.2.1})
\underline{is a perfect duality}: \underline{it identifies each of the
component groups} $\Phi_o$ \underline{and} $\Phi'_o$
\underline{with the dual of the other, with values in}
$(\mathbb Q/\mathbb Z)_k$.

\medskip
\noindent
\SGARefTarget{sga7.sga7i.eix.n1-3-1}{1.3.1}. The formation of (\SGARefLink{sga7.sga7i.eix.n1-2-1}{1.2.1}) commutes with formally etale base
change (VIII~\SGARefLink{sga7.sga7i.eviii.n7-3-5-3}{7.3.5.3}), and the formation of Neron models has the
analogous property.  Consequently, to prove Conjecture~\SGARefLink{sga7.sga7i.eix.n1-3}{1.3} it would
suffice to treat the case in which $S$ is complete and has separably
closed residue field.  Artin and Mazur proved the conjecture when
$A_K$ is the Jacobian of a proper smooth curve over $K$, using their
general, unpublished theory of the self-duality of the relative
Jacobian (in the sense of derived categories) of a proper curve scheme
over $S$.  They also proved Conjecture~\SGARefLink{sga7.sga7i.eix.n1-3}{1.3} when the residue field is
\underline{finite}, using Tate duality [32].  We shall likewise see in
\SGARefLink{sga7.sga7i.eix.n11-4}{11.4} and \SGARefLink{sga7.sga7i.eix.n11-5}{11.5} (subject to compatibility checks) that \SGARefLink{sga7.sga7i.eix.n1-3}{1.3} holds on the
$\ell$-primary components for $\ell\ne p$, and that it holds in the
case of ``semistable reduction''.  These results make \SGARefLink{sga7.sga7i.eix.n1-3}{1.3} extremely
plausible.  Moreover, \SGARefLink{sga7.sga7i.eix.n1-3}{1.3} may also be regarded as a by-product of a
general duality theory for group schemes over fraction fields of
complete discrete valuation rings, formally analogous to SGA~5~I.
Such a theory remains heuristic at present; it would also contain
local class field theory in the form given to it by Serre [27].

\medskip
\noindent
\SGARefTarget{sga7.sga7i.eix.n1-3-2}{1.3.2}. For a given abelian scheme $A_K$, the validity of \SGARefLink{sga7.sga7i.eix.n1-3}{1.3} means that
the pairing (\SGARefLink{sga7.sga7i.eix.n1-2-1}{1.2.1}) is \underline{separating} on both the left and the
right.  By VIII~\SGARefLink{sga7.sga7i.eviii.n7-3-2}{7.3.2}, separation on the left says precisely that the
only open subgroup scheme $V$ of $A$ for which $W_K$ can be extended to
a biextension of $(V,A')$ by $\mathbb G_{mS}$ is the identity component
$A^o$ of $A$.  There is a symmetric interpretation of separation on
the right.

\medskip
\noindent
\SGARefTarget{sga7.sga7i.eix.n1-4}{1.4}. More generally, suppose that
\begin{equation}\tag{\SGARefTarget{sga7.sga7i.eix.n1-4-1}{1.4.1}}
 U\subset A\qquad\text{and}\qquad U'\subset A'
\end{equation}
are open subgroup schemes such that $W_K$ extends to a biextension of
$(U,U')$ by $\mathbb G_{mS}$.  By VIII~\SGARefLink{sga7.sga7i.eviii.n7-3-2}{7.3.2}, this says explicitly that
the subgroups
% SGA-SRC-001266
\begin{equation}\tag{\SGARefTarget{sga7.sga7i.eix.n1-4-2}{1.4.2}}
 F_o=U_o/U_o^o\subset\Phi_o,
 \qquad
 F'_o=U'_o/{U'}_o^o\subset\Phi'_o
\end{equation}
defined by $U$ and $U'$ are orthogonal under (\SGARefLink{sga7.sga7i.eix.n1-2-1}{1.2.1}).\footnote{The
printed formula, confirmed in a direct high-resolution check, reuses
$\Phi'_o$ for the first subgroup and writes $\Psi'_o\subset\Psi_o$ for
the second, colliding with the established notation for the component
group of $A'$.  The neutral local names $F_o,F'_o$ above make the
intended two subgroups unambiguous without changing the assertion.}
We denote this extension (unique up to unique isomorphism) of $W_K$ by
$W$.  It is determined up to unique isomorphism by its class
\begin{equation}\tag{\SGARefTarget{sga7.sga7i.eix.n1-4-3}{1.4.3}}
 w\in\mathrm{Biext}^1(U,U';\mathbb G_{mS}),
\end{equation}
which is itself characterized by the condition that its image in
$\mathrm{Biext}^1(A_K,A'_K;\mathbb G_{mK})$ be the class $w_K$ of
(\SGARefLink{sga7.sga7i.eix.n1-0-1}{1.0.1}).

For example, one may take for $U$ and $U'$ the identity components
$A^o$ and $A'{}^o$ of $A$ and $A'$.  One thereby obtains a
\underline{canonical biextension $W^o$, or $W$, of $(A^o,A'{}^o)$ by
$\mathbb G_{mS}$}, giving rise to a class
\begin{equation}\tag{\SGARefTarget{sga7.sga7i.eix.n1-4-4}{1.4.4}}
 w=w^o\in\mathrm{Biext}^1(A^o,A'{}^o;\mathbb G_{mS}),
\end{equation}
with which we shall work chiefly.  Evidently, the restrictions to
$(A^o,A'{}^o)$ of the classes $w$ in (\SGARefLink{sga7.sga7i.eix.n1-4-3}{1.4.3}), for the various pairs
$(U,U')$, are all equal to $w^o$ in (\SGARefLink{sga7.sga7i.eix.n1-4-4}{1.4.4}).

\medskip
\noindent
\underline{Proposition \SGARefTarget{sga7.sga7i.eix.n1-5}{1.5}.} \underline{Consider the special fiber
$W_s$ of the canonical biextension} $W$ \underline{of}
$(A^o,A'{}^o)$ \underline{by} $\mathbb G_{mS}$ \underline{in (\SGARefLink{sga7.sga7i.eix.n1-4-4}{1.4.4}),
so that $W_s$ is a biextension of $(A_o^o,{A'}_o^o)$ by}
$\mathbb G_{mk}$.  \underline{This biextension is separating}
(VIII~\SGARefLink{sga7.sga7i.eviii.n4-11}{4.11}).

Indeed, choose an ample invertible sheaf $\mathbb L$ on $A^o$.  Then
$\mathbb L_\eta$ defines a homomorphism (cf. \SGARefLink{sga7.sga7i.eix.n1-0-5}{1.0.5})
\[
 A_K\longrightarrow A'_K,
\]
which extends to a homomorphism $A\longrightarrow A'$, and hence gives
\[
 \psi:A^o\longrightarrow A'{}^o.
\]
We may consider the biextension $W_{\mathbb L}$ of $(A^o,A^o)$ by
$\mathbb G_{mS}$ obtained by pulling back $W$ along
$(\operatorname{id}_{A^o},\psi)$.  Its restriction to the generic
fibers is precisely the biextension $\delta(\mathbb L_\eta)$ of
VIII~\SGARefLink{sga7.sga7i.eviii.n4-12}{4.12}.  Hence, by VIII~\SGARefLink{sga7.sga7i.eviii.n7-4}{7.4}, $W_{\mathbb L}$ is associated with the
birigidified invertible sheaf $\delta(\mathbb L)$, and consequently
$(W_{\mathbb L})_s$ is defined by $\delta(\mathbb L_s)$.  On the other
hand, it may also be described as the pullback of $W_s$ by the
morphisms
\[
 \operatorname{id}_{A_o^o}:A_o^o\longrightarrow A_o^o,
 \qquad
 \psi_o:A_o^o\longrightarrow {A'}_o^o.
\]
Since $\mathbb L$, and hence $\mathbb L_s$, is ample, VIII~\SGARefLink{sga7.sga7i.eviii.n4-13}{4.13} shows
that $W_{\mathbb L_s}$ is separating, in particular separating on the
left.  A fortiori $W_s$ is therefore separating on the left.  By
symmetry it is also separating on the right, which proves \SGARefLink{sga7.sga7i.eix.n1-5}{1.5}.

We shall substantially sharpen \SGARefLink{sga7.sga7i.eix.n1-5}{1.5} when $A_o^o$ has unipotent rank
zero (\SGARefLink{sga7.sga7i.eix.n5-4}{5.4}).

% END INCLUDED SGA7I_English_current_buildable_source_20260823/components/131_expose_IX_sections_1_2_through_1_5.tex

% BEGIN INCLUDED SGA7I_English_current_buildable_source_20260823/components/132_expose_IX_section_2_1.tex
% Expose IX, section 2 title and 2.1--2.1.14.
% French transcription lines 344--441; authority physical pp. 338--341
% (source folios 326--329).  All pages were checked in four overlapping
% 1100-dpi bands.  Diagram (2.1.6) was checked node-by-node and arrow-by-arrow
% in four overlapping direct-source crops rendered at 5000 dpi.

\medskip
\noindent
\underline{2. Fixed and toric parts of $T_\ell(A_K)$.  Criteria for
good reduction.  Orthogonality theorem for $\ell\ne p$.}

\medskip
\noindent
\SGARefTarget{sga7.sga7i.eix.n2-1}{2.1}. Like every commutative algebraic group over $k$, $A_o$ has a
largest subtorus, whose formation commutes with every extension of the
base field [SGA~3 XVII~7.2.1, XII~1.12]:
\begin{equation}\tag{\SGARefTarget{sga7.sga7i.eix.n2-1-1}{2.1.1}}
 T_o\subset A_o.
\end{equation}
When $k$ is perfect, $T_o$ is characterized by the condition that
$A_o^o/T_o$ be an extension of an abelian scheme $B_o$ over $k$ by a
smooth connected unipotent group [25].  The latter necessarily has a
composition series all of whose quotients are isomorphic to
$\mathbb G_{ak}$ (SGA~3 XVII~4.1.3), and one has an exact sequence
\begin{equation}\tag{\SGARefTarget{sga7.sga7i.eix.n2-1-2}{2.1.2}}
 0\longrightarrow U_o\longrightarrow A_o^o/T_o
 \longrightarrow B_o\longrightarrow0.
\end{equation}
In any case, whether or not $k$ is perfect, \underline{if} $A_o^o$ is
an extension of an abelian scheme $B_o$ by a smooth connected affine
scheme $L_o$,
\begin{equation}\tag{\SGARefTarget{sga7.sga7i.eix.n2-1-3}{2.1.3}}
 0\longrightarrow L_o\longrightarrow A_o^o
 \longrightarrow B_o\longrightarrow0,
\end{equation}
this extension structure is known to be unique [25].  Moreover,
$T_o\subset L_o$, and $T_o$ is the largest subtorus of $L_o$.  Thus
one has the exact sequence (\SGARefLink{sga7.sga7i.eix.n2-1-2}{2.1.2}), with
\begin{equation}\tag{\SGARefTarget{sga7.sga7i.eix.n2-1-4}{2.1.4}}
 U_o=L_o/T_o.
\end{equation}

\medskip
\noindent
\SGARefTarget{sga7.sga7i.eix.n2-1-5}{2.1.5}. Let us note that $A_o^o$ has \underline{unipotent rank zero}
(that is, after passage to an algebraic closure of $k$, it has no
subgroup isomorphic to $\mathbb G_a$) if and only if $A_o^o/T_o$ is an
abelian scheme $B_o$.  In that case the preceding hypotheses apply
with $U_o=0$.

Let $\ell$ be a prime different from the characteristic of $k$, so that
$A_o^o$ is $\ell$-divisible.  The inclusion (\SGARefLink{sga7.sga7i.eix.n2-1-1}{2.1.1}) then gives an exact
sequence of pro-group schemes:
\begin{equation}\tag{\SGARefTarget{sga7.sga7i.eix.n2-1-6}{2.1.6}}
\begin{tikzcd}[row sep=large,column sep=large]
0 \arrow[r] & T_\ell(T_o) \arrow[r]
 & T_\ell(A_o^o) \arrow[r] \arrow[d,"\wr"]
 & T_\ell(A_o^o/T_o) \arrow[r] \arrow[d,"\wr"] & 0\\
 & & T_\ell(A_o) & T_\ell(A_o/T_o) & {,}
\end{tikzcd}
\end{equation}
whose terms may also be interpreted as twisted constant free
$\ell$-adic sheaves.  When an exact sequence (\SGARefLink{sga7.sga7i.eix.n2-1-2}{2.1.2}) is available---for
example, when $k$ is perfect or when $A_o^o$ has unipotent rank
zero---one has
\begin{equation}\tag{\SGARefTarget{sga7.sga7i.eix.n2-1-7}{2.1.7}}
 T_\ell(A_o^o/T_o)\xrightarrow{\ \sim\ }T_\ell(B_o),
\end{equation}
so that the exact sequence (\SGARefLink{sga7.sga7i.eix.n2-1-6}{2.1.6}) becomes
\begin{equation}\tag{\SGARefTarget{sga7.sga7i.eix.n2-1-8}{2.1.8}}
 0\longrightarrow T_\ell(T_o)\longrightarrow T_\ell(A_o^o)
 \longrightarrow T_\ell(B_o)\longrightarrow0.
\end{equation}
Put
\begin{equation}\tag{\SGARefTarget{sga7.sga7i.eix.n2-1-9}{2.1.9}}
 n=\dim A_K=\dim A_o,\qquad
 \mu=\dim T_o,\qquad
 \lambda=\dim U_o,\qquad
 \alpha=\dim B_o.
\end{equation}
Then
\begin{equation}\tag{\SGARefTarget{sga7.sga7i.eix.n2-1-10}{2.1.10}}
 n=\mu+\lambda+\alpha,
\end{equation}
and
\begin{equation}\tag{\SGARefTarget{sga7.sga7i.eix.n2-1-11}{2.1.11}}
 \left\{
 \begin{aligned}
 \operatorname{rank}T_\ell(T_o)&=\mu,&
 \operatorname{rank}T_\ell(A_o^o/T_o)
 &=\operatorname{rank}T_\ell(B_o)=2\alpha,\\
 \operatorname{rank}T_\ell(A_o^o)&=\mu+2\alpha.
 \end{aligned}
 \right.
\end{equation}
The second relation is a well-known theorem of Weil [5, 15].

\medskip
\noindent
\SGARefTarget{sga7.sga7i.eix.n2-1-12}{2.1.12}. The integers $n,\mu,\lambda,\alpha$ are plainly defined, after
passing to an algebraic closure of $k$, even without assuming the
existence of an exact sequence (\SGARefLink{sga7.sga7i.eix.n2-1-2}{2.1.2}).  They are respectively called
the dimension of $A_o$, the \underline{reductive rank of $A_o$}
(equal to the dimension of $T_o$), the \underline{unipotent rank of
$A_o$}, and the \underline{abelian rank of $A_o$}.  They are still
related by (\SGARefLink{sga7.sga7i.eix.n2-1-10}{2.1.10}) and (\SGARefLink{sga7.sga7i.eix.n2-1-11}{2.1.11}), with the term
$\operatorname{rank}T_\ell(B_o)$ omitted from the latter.

\medskip
\noindent
\SGARefTarget{sga7.sga7i.eix.n2-1-13}{2.1.13}. Suppose that $k$ has characteristic $p>0$ and that $A_o^o$ is
$p$-divisible, or equivalently that it has unipotent rank zero
(cf. \SGARefLink{sga7.sga7i.eix.n2-1-5}{2.1.5}).  One should then also consider the pro-group $T_p(A_o^o)$,
which is a \underline{Barsotti--Tate group} (or a
\underline{$p$-divisible group}, in the terminology of [33]).  The
exact sequence (\SGARefLink{sga7.sga7i.eix.n2-1-8}{2.1.8}) remains valid for $\ell=p$ and is then an exact
sequence of Barsotti--Tate groups.  Here $\lambda=0$, and hence
\begin{equation}\tag{\SGARefTarget{sga7.sga7i.eix.n2-1-13-1}{2.1.13.1}}
 n=\mu+\alpha.
\end{equation}
The formulas (\SGARefLink{sga7.sga7i.eix.n2-1-11}{2.1.11}) also remain valid, provided that rank is
interpreted as the rank (or ``height'') of the Barsotti--Tate groups in
question.

\medskip
\noindent
\SGARefTarget{sga7.sga7i.eix.n2-1-14}{2.1.14}. For ${A'}_o^o$ we likewise introduce the invariants
\[
 T'_o\subset {A'}_o^o,\qquad U'_o,\qquad B'_o,
\]
where $T'_o$ is the toric part of ${A'}_o^o$, $B'_o$ its abelian part
(defined in particular when $k$ is perfect or when ${A'}_o^o$ has
unipotent rank zero), and $U'_o$ its unipotent part (defined, for
example, when $k$ is perfect, or trivially when ${A'}_o^o$ has
unipotent rank zero).  There is a corresponding exact sequence of the
form (\SGARefLink{sga7.sga7i.eix.n2-1-6}{2.1.6}), or (\SGARefLink{sga7.sga7i.eix.n2-1-8}{2.1.8}), which we shall not write again.  One could
introduce integers $\mu',\lambda',\alpha'$ for ${A'}_o^o$ as in
(\SGARefLink{sga7.sga7i.eix.n2-1-9}{2.1.9}), but we shall see below (\SGARefLink{sga7.sga7i.eix.n2-2-7}{2.2.7}) that they are the same as those
for $A_o^o$; this spares us from introducing new notation for them.

% END INCLUDED SGA7I_English_current_buildable_source_20260823/components/132_expose_IX_section_2_1.tex

% BEGIN INCLUDED SGA7I_English_current_buildable_source_20260823/components/133_expose_IX_section_2_2.tex
% Expose IX, section 2.2--2.2.9.1.
% French transcription lines 443--599; authority physical pp. 341--348
% (source folios 329--336).  Physical pp. 342--348 were checked in four
% overlapping 1100-dpi bands; the opening on p. 341 was checked with the
% preceding component.  The printed cross-references 1.5.5 and 6.2.5.1 are
% transparently corrected below to 2.1.5 and 2.2.5.1 respectively.

\medskip
\noindent
\SGARefTarget{sga7.sga7i.eix.n2-2}{2.2}. We propose to make more precise the relations between
$T_\ell(A_K)$ and $T_\ell(A_o)=T_\ell(A_o^o)$ by bringing
$T_\ell(A^o)$ into play.  First recall the following.

\medskip
\noindent
\underline{Lemma} \SGARefTarget{sga7.sga7i.eix.n2-2-1}{2.2.1}. \underline{Let $m>0$ be an integer.  Then the
following conditions are equivalent:}

\smallskip
\noindent
(i) $m\,\mathrm{id}_{A^o}$ \underline{is flat}.

\noindent
(ii) $m\,\mathrm{id}_{A^o}$ \underline{is surjective}.

\noindent
(iii) $m\,\mathrm{id}_{A^o}$ \underline{is quasi-finite}.

\noindent
(iv) $m$ \underline{is prime to the characteristic of $k$}, or $A_o^o$
\underline{has unipotent rank zero} (cf. \SGARefLink{sga7.sga7i.eix.n2-1-5}{2.1.5}).

\smallskip
\noindent
\underline{Moreover, if these conditions hold for $m$, they hold for
$m^\nu$ for every $\nu>0$; the groups}
\[
 {}_{m^\nu}A^o\underset{\mathrm{dfn}}{=}
 \ker(m^\nu\,\mathrm{id}_{A^o})
\]
\underline{are quasi-finite, separated, and flat over $S$, and, as
$\nu$ varies, they form an inverse system of group schemes whose
transition morphisms}
\[
 x\longmapsto m^{\nu'-\nu}x:
 {}_{m^{\nu'}}A^o\longrightarrow{}_{m^\nu}A^o
 \qquad(\nu'\geq\nu)
\]
\underline{are faithfully flat}.

Indeed, the fiberwise criterion for flatness (EGA~IV~11.3.5) shows that
(i), like (ii) and (iii), is a condition on the fibers; these three
conditions are equivalent by SGA~3 $\mathrm{VI}_A$~5.6, since the
fibers in question are smooth and connected.  In the form (ii), a
classical theorem of Weil [5, 15] shows that the conditions hold on the
generic fiber, which is an abelian scheme; it is therefore enough to
state them on the special fiber.  Their equivalence with (iv) is then
also well known from the structure theory.  In any of these forms,
stability of the condition under $m\mapsto m^\nu$ is immediate.  It
follows that the ${}_{m^\nu}A^o$ are quasi-finite flat group schemes
over $S$, and they are plainly separated over $S$ because $A^o$ is.
Finally, faithful flatness of the transition morphisms---equivalently,
the fact that they are epimorphisms for the fppf topology---follows
formally from the fact that $m\,\mathrm{id}_{A^o}$ is faithfully flat,
hence an epimorphism for that topology.

\medskip
\noindent
\SGARefTarget{sga7.sga7i.eix.n2-2-2}{2.2.2}. Consequently, we shall take $m$ to be a prime $\ell$ satisfying
the conditions of \SGARefLink{sga7.sga7i.eix.n2-2-1}{2.2.1}.  We then denote by
\begin{equation}\tag{\SGARefTarget{sga7.sga7i.eix.n2-2-2-1}{2.2.2.1}}
 T_\ell(A^o)=({}_{\ell^\nu}A^o)_{\nu>0}
\end{equation}
the inverse system described in \SGARefLink{sga7.sga7i.eix.n2-2-1}{2.2.1}.  Under the two base changes
$\eta\longrightarrow S$ and $s\longrightarrow S$, it recovers the
inverse systems $T_\ell(A_K)$ and
$T_\ell(A_o^o)=T_\ell(A_o)$ considered in \SGARefLink{sga7.sga7i.eix.n1-0}{1.0} and \SGARefLink{sga7.sga7i.eix.n2-1}{2.1} respectively.

\medskip
\noindent
\SGARefTarget{sga7.sga7i.eix.n2-2-3}{2.2.3}. Suppose now that $S$ is henselian.  Every quasi-finite separated
scheme $X$ over $S$ then has a canonical decomposition as a disjoint
union
\begin{equation}\tag{\SGARefTarget{sga7.sga7i.eix.n2-2-3-1}{2.2.3.1}}
 X=X^f\coprod X',
\end{equation}
where $X^f$ is \underline{finite over $S$} and $X'_o=\varnothing$, that
is, $X'$ is confined to its generic fiber $X'_\eta$; thus $X^f$ has the
same special fiber as $X$.  This decomposition is plainly functorial in
$X$.  In particular, if $X$ is a group scheme over $S$, then $X^f$ is
a subgroup scheme.  If $X$ is flat, respectively \'{e}tale, over $S$,
then so is $X^f$.

We apply this to the components ${}_{\ell^\nu}A^o$ of $T_\ell(A^o)$.
This gives a projective subsystem of groups
\begin{equation}\tag{\SGARefTarget{sga7.sga7i.eix.n2-2-3-2}{2.2.3.2}}
 \bigl(({}_{\ell^\nu}A^o)^f\bigr)_{\nu>0}
 \underset{\mathrm{dfn}}{=}T_\ell(A^o)^f
 \subset T_\ell(A^o),
\end{equation}
called the \underline{fixed part} of the Tate module $T_\ell(A^o)$.
By construction, there are canonical isomorphisms, functorial in $A$,
\begin{equation}\tag{\SGARefTarget{sga7.sga7i.eix.n2-2-3-3}{2.2.3.3}}
 T_\ell(A^o)^f\times_S s\xrightarrow{\ \sim\ }
 T_\ell(A^o)\times_Ss\xrightarrow{\ \sim\ }T_\ell(A_o^o).
\end{equation}
On the other hand, the generic fiber of $T_\ell(A^o)^f$ defines a
canonical sub-pro-group of $T_\ell(A_K)$, again called the
``\underline{fixed part}'' of this Tate module:
\begin{equation}\tag{\SGARefTarget{sga7.sga7i.eix.n2-2-3-4}{2.2.3.4}}
 T_\ell(A_K)^f\subset T_\ell(A_K).
\end{equation}

\medskip
\noindent
\SGARefTarget{sga7.sga7i.eix.n2-2-4}{2.2.4}. When $\ell\ne p=\operatorname{char}k$, the fixed part
$T_\ell(A^o)^f$ may be interpreted as a twisted constant $\ell$-adic
sheaf on $S$ (the $({}_{\ell^\nu}A^o)^f$ being finite \'{e}tale over
$S$).  Since the functor
\[
 X\longmapsto X_o=X\times_Ss
\]
from finite \'{e}tale schemes over $S$ to finite \'{e}tale schemes
over $k$ is an equivalence of categories when $S$ is henselian, there
is an analogous equivalence for the categories of twisted constant
$\ell$-adic sheaves over $S$ and over $k$.  Hence (\SGARefLink{sga7.sga7i.eix.n2-2-3-3}{2.2.3.3}), which
gives the special fiber of $T_\ell(A^o)^f$, characterizes this fixed
part up to unique isomorphism.  On the other hand, as an $\ell$-adic
subsheaf of $T_\ell(A^o)$ it is plainly characterized by its generic
fiber (\SGARefLink{sga7.sga7i.eix.n2-2-3-4}{2.2.3.4}), or equivalently by the corresponding
$\mathbb Z_\ell$-submodule, stable under
$\pi=\operatorname{Gal}(\overline K/K)$:
\begin{equation}\tag{\SGARefTarget{sga7.sga7i.eix.n2-2-4-1}{2.2.4.1}}
 T_\ell(A_K)^f(\overline K)
 \underset{\mathrm{dfn}}{=}T_\ell(A(\overline K))^f
 \subset T_\ell(A(\overline K)).
\end{equation}
In fact, one has the following result.

\medskip
\noindent
\underline{Proposition} \SGARefTarget{sga7.sga7i.eix.n2-2-5}{2.2.5}. \underline{Suppose, as above, that
$\ell\ne p$ and that $S$ is henselian.  Then the submodule (\SGARefLink{sga7.sga7i.eix.n2-2-4-1}{2.2.4.1})
of $T_\ell(A(\overline K))=T_\ell(A_K)(\overline K)$ corresponding to
the fixed part $T_\ell(A_K)^f$ in (\SGARefLink{sga7.sga7i.eix.n2-2-3-4}{2.2.3.4}) is the submodule
$T_\ell(A(\overline K))^I$ of invariants under the inertia group
$I\subset\pi$.}

This follows, by passage to the inverse limit, from the analogous
relation valid for every integer $m>0$ prime to
$p=\operatorname{char}k$:
\begin{equation}\tag{\SGARefTarget{sga7.sga7i.eix.n2-2-5-1}{2.2.5.1}}
 ({}_mA)^f(\overline K)=({}_mA(\overline K))^I.
\end{equation}
To prove this relation, introduce the strict localization $\widetilde S$
of $S$, whose fraction field $\widetilde K$ is the maximal unramified
subextension of $\overline K/K$.  Since $({}_mA)^f$ is finite
\'{e}tale over $S$, its points with values in $\overline K$ already
have values in $\widetilde K$.  Thus the left-hand side of (\SGARefLink{sga7.sga7i.eix.n2-2-5-1}{2.2.5.1})
is
$\operatorname{Hom}_S(\widetilde S,({}_mA)^f)$, whereas the right-hand
side is, by definition of $\widetilde K$, equal to ${}_mA(\widetilde K)$.
By the definition of $({}_mA)^f$, and since $\widetilde S$ is an
inverse limit of finite \'{e}tale $S$-schemes $S_i$, the former is also
equal to
\[
 \operatorname{Hom}_S(\widetilde S,{}_mA)
 ={}_m\operatorname{Hom}_S(\widetilde S,A).
\]
By the defining property of the N\'{e}ron model,
$\operatorname{Hom}_S(\widetilde S,A)=A(\widetilde K)$, which proves
(\SGARefLink{sga7.sga7i.eix.n2-2-5-1}{2.2.5.1}).

To deduce \SGARefLink{sga7.sga7i.eix.n2-2-5}{2.2.5}, observe that (\SGARefLink{sga7.sga7i.eix.n2-2-5-1}{2.2.5.1}) gives
\begin{equation}\tag{*}
 T_\ell(A(\overline K))^I
 =\varprojlim_\nu({}_{\ell^\nu}A)^f(\overline K)
 \supset\varprojlim_\nu({}_{\ell^\nu}A^o)^f(\overline K).
\end{equation}
If $(x_\nu)$ belongs to the rightmost term, then for every
$\nu\geq0$, $x_\nu$ is a section of $A$ that is infinitely divisible
by $\ell$.  Its value at $s$ is therefore an element of $A_o(k)$
infinitely divisible by $\ell$, and consequently lies in $A_o^o(k)$.
This proves equality of the last two terms of (*) and completes the
proof of \SGARefLink{sga7.sga7i.eix.n2-2-5}{2.2.5}.

\medskip
\noindent
\underline{Proposition} \SGARefTarget{sga7.sga7i.eix.n2-2-6}{2.2.6}. \underline{Let $C_K$ be a second
abelian scheme over $K$, with N\'{e}ron model $C$, and consider a
homomorphism $u_K:C_K\longrightarrow A_K$, which therefore extends to
a homomorphism $u:C\longrightarrow A$.  Suppose that $u_K$ is an
isogeny.  Then, for every prime $\ell$,
$T_\ell(u_K):T_\ell(C_K)\longrightarrow T_\ell(A_K)$ induces an
isogeny of fixed parts}
\begin{equation}\tag{\SGARefTarget{sga7.sga7i.eix.n2-2-6-1}{2.2.6.1}}
 T_\ell(C_K)^f\longrightarrow T_\ell(A_K)^f,
\end{equation}
\underline{and likewise the map induced by
$u_o:C_o\longrightarrow A_o$}
\begin{equation}\tag{\SGARefTarget{sga7.sga7i.eix.n2-2-6-2}{2.2.6.2}}
 T_\ell(u_o^o):T_\ell(C_o^o)\longrightarrow T_\ell(A_o^o)
\end{equation}
\underline{is an isogeny.  Finally, $u_o:C_o\longrightarrow A_o$
induces an isogeny from the largest subtorus of $C_o$ to the largest
subtorus $T_o$ of $A_o$; and over an extension $k'/k$ on which the
``abelian parts'' of $C_o$ and $A_o$ are defined (for example, over
the perfect closure of $k$), $u_o$ induces an isogeny between the
abelian parts of $C_{ok'}$ and $A_{ok'}$.}

This statement follows formally from the fact that
$T_\ell(A_K)^f$, $T_\ell(A_o^o)$, the largest subtorus of $A_o$, and
the ``abelian part'' of $A_{ok'}$ are additive functors of the varying
$A_K$, and from the fact that $u_K$ is an isogeny precisely in the
sense that there is a morphism $v_K:A_K\longrightarrow C_K$ such that
\[
 u_Kv_K=n\,\mathrm{id},\qquad v_Ku_K=n\,\mathrm{id}
\]
for some integer $n>0$.

\medskip
\noindent
\underline{Corollary} \SGARefTarget{sga7.sga7i.eix.n2-2-7}{2.2.7}. \underline{If $A_K$ and $C_K$ are
isogenous abelian schemes over $K$, then the corresponding invariants
$\mu,\lambda,\alpha$ in (\SGARefLink{sga7.sga7i.eix.n2-1-11}{2.1.11}) are the same for $A_K$ and $C_K$.
In particular, they are the same for $A_K$ and its dual $A'_K$.}

\medskip
\noindent
\underline{Remarks} \SGARefTarget{sga7.sga7i.eix.n2-2-8}{2.2.8}. In the two preceding statements, \SGARefLink{sga7.sga7i.eix.n2-2-6}{2.2.6} and
\SGARefLink{sga7.sga7i.eix.n2-2-7}{2.2.7}, we assumed that $S$ is henselian.  It is immediate, however,
that \SGARefLink{sga7.sga7i.eix.n2-2-7}{2.2.7}, as well as the assertions of \SGARefLink{sga7.sga7i.eix.n2-2-6}{2.2.6} concerning the
properties of the induced homomorphism
$u_o:C_o\longrightarrow A_o$, holds without this hypothesis.  Indeed,
one reduces to the case already treated by the base change
$S^h\longrightarrow S$, where $S^h$ is the henselization of $S$, using
the fact---immediate from the definitions---that formation of N\'{e}ron
models commutes with this base change.

\medskip
\noindent
\underline{Corollary} \SGARefTarget{sga7.sga7i.eix.n2-2-9}{2.2.9}. (\underline{Good reduction.})
\underline{Let $S$ be a trait, let $A_K$ be an abelian scheme over its
fraction field $K$, let $A$ be the N\'{e}ron model of $A_K$, and let
$\ell$ be a prime different from $\operatorname{char}k$.  The following
conditions are equivalent.}

\smallskip
\noindent
(i) \underline{There exists an abelian scheme over $S$ whose generic
fiber is isomorphic to $A_K$.}

\noindent
(ii) $A$ \underline{is an abelian scheme}.

\noindent
(ii$^o$) $A^o$ \underline{is an abelian scheme}.

\noindent
(iii) $A_o$ \underline{is an abelian scheme}.

\noindent
(iii$^o$) $A_o^o$ \underline{is an abelian scheme}, \underline{that is,
with the notation of (\SGARefLink{sga7.sga7i.eix.n2-1-11}{2.1.11}), $\lambda=\mu=0$}.

\noindent
(iv) $A$ \underline{is proper over $S$}.

\noindent
(iv$^o$) $A^o$ \underline{is proper over $S$}.

\noindent
(v) $T_\ell(A_K)$ \underline{is ``unramified over $S$''}, that is,
\underline{it extends to a twisted constant $\ell$-adic sheaf over
$S$}, or equivalently \underline{the inertia group $I$ acts trivially
on $T_\ell(A_K(\overline K))$}.

\noindent
(v bis) \underline{For every integer $\nu\geq0$,
${}_{\ell^\nu}A_K$ is ``unramified over $S$''}, that is,
\underline{it extends to an \'{e}tale covering of $S$}, or equivalently
\underline{the inertia group $I$ acts trivially on
${}_{\ell^\nu}A_K(\overline K)$}.

\smallskip
\noindent
\underline{Finally, an abelian scheme over $S$ as in (i) is a N\'{e}ron
model of $A_K$, and in particular is determined up to unique
isomorphism.}

\medskip
\noindent
\SGARefTarget{sga7.sga7i.eix.n2-2-9-1}{2.2.9.1}. We say that $A_K$ has \underline{good reduction} over $S$ if
it satisfies condition (i) above.  Criterion (v), or (v bis), is known
as the \underline{N\'{e}ron--Ogg--Shafarevich criterion} for good
reduction.

Let us prove \SGARefLink{sga7.sga7i.eix.n2-2-9}{2.2.9}.  It is well known (cf., for example, [5]) that an
abelian scheme over $S$ satisfies the universal property of a
N\'{e}ron model; this reduces to the fact that a rational section of an
abelian scheme over a regular base is defined everywhere.  This proves
the equivalence
$(i)\Longleftrightarrow(ii)\Longleftrightarrow(ii^o)$, as well as the
last assertion of \SGARefLink{sga7.sga7i.eix.n2-2-9}{2.2.9}.  These conditions plainly imply (iii), (iv),
and (iv$^o$), each of which implies (iii$^o$).  Let us prove that
(iii$^o$)$\Longrightarrow$(ii$^o$), which will establish the
equivalence of (i) through (iv$^o$).  This follows from the fact that a
smooth group scheme over $S$ whose fibers are abelian schemes is proper
over $S$ (and hence is an abelian scheme), by EGA~IV~15.7.10.

On the other hand, in view of \SGARefLink{sga7.sga7i.eix.n2-2-5}{2.2.5} and the fact that
\[
 T_\ell(A_K(\overline K))/T_\ell(A_K(\overline K))^I
\]
is plainly a torsion-free $\mathbb Z_\ell$-module, condition (v) says
that $T_\ell(A_K(\overline K))$ and its fixed part have the same rank.
By (\SGARefLink{sga7.sga7i.eix.n2-1-11}{2.1.11}) and (\SGARefLink{sga7.sga7i.eix.n2-2-3-3}{2.2.3.3}), this is the condition
\[
 \mu+2\alpha=2(\lambda+\mu+\alpha),
\]
that is, $2\lambda+\mu=0$, or $\lambda=\mu=0$, which is precisely
(iii$^o$).  Finally, the equivalence of (v) and (v bis) is clear from
the definitions.  This completes the proof of \SGARefLink{sga7.sga7i.eix.n2-2-9}{2.2.9}.

% END INCLUDED SGA7I_English_current_buildable_source_20260823/components/133_expose_IX_section_2_2.tex

% BEGIN INCLUDED SGA7I_English_current_buildable_source_20260823/components/134_expose_IX_sections_2_3_and_2_4.tex
% Expose IX, sections 2.3 and 2.4.
% French transcription lines 601--685; authority physical pp. 348--352
% (source folios 336--340).  Physical pp. 349--352 were checked in four
% overlapping 1100-dpi bands; the opening on p. 348 was checked with the
% preceding component.  The printed reference (1.6.5) is transparently
% corrected to (2.2.5), the missing prime in the second factor of (2.4.5)
% is restored, and the impossible printed reference (1.10) is replaced by
% an unnumbered verbal reference to the displayed inclusion (2.4.7).

\medskip
\noindent
\SGARefTarget{sga7.sga7i.eix.n2-3}{2.3}. Suppose once again that $S$ is henselian and that
$\ell\ne p=\operatorname{char}k$.  We now use the inclusion
\begin{equation}\tag{\SGARefTarget{sga7.sga7i.eix.n2-3-1}{2.3.1}}
 T_\ell(T_o)\hookrightarrow T_\ell(A_o^o)
\end{equation}
of (\SGARefLink{sga7.sga7i.eix.n2-1-6}{2.1.6}), together with the equivalence of categories recalled in
\SGARefLink{sga7.sga7i.eix.n2-2-4}{2.2.4}.  We obtain a canonical twisted constant $\ell$-adic subsheaf
of $T_\ell(A^o)^f$, called the \underline{toric part of}
$T_\ell(A^o)$,
\begin{equation}\tag{\SGARefTarget{sga7.sga7i.eix.n2-3-2}{2.3.2}}
 T_\ell(A^o)^t\hookrightarrow T_\ell(A^o)^f
 \hookrightarrow T_\ell(A^o),
\end{equation}
characterized by the condition that its special fiber is the inclusion
(\SGARefLink{sga7.sga7i.eix.n2-3-1}{2.3.1}):
\begin{equation}\tag{\SGARefTarget{sga7.sga7i.eix.n2-3-3}{2.3.3}}
 T_\ell(A^o)^t\times_Ss=T_\ell(T_o)
 \hookrightarrow T_\ell(A_o^o).
\end{equation}
Later we shall give an interpretation of the toric part of
$T_\ell(A^o)$ as the Tate module $T_\ell(T)$ of a torus $T$ over $S$
lifting the torus $T_o$ over $k$.  That interpretation will allow the
definition of the toric part to be extended to the case $\ell=p$.
For the moment, and to fix ideas, we confine ourselves to the case
$\ell\ne p$.  In this case the toric part of $T_\ell(A^o)$ may be
described by the corresponding $\mathbb Z_\ell$-submodule of
$T_\ell(A(\overline K))$:
\begin{equation}\tag{\SGARefTarget{sga7.sga7i.eix.n2-3-4}{2.3.4}}
 T_\ell(A(\overline K))^t\hookrightarrow
 T_\ell(A(\overline K))^f\hookrightarrow
 T_\ell(A(\overline K)).
\end{equation}
It is stable under the action of $\pi$; its first term is defined as
$T_\ell(A^o)^t(\overline K)$ and is, as one would expect, again called
the \underline{toric submodule} of the Tate module
$T_\ell(A(\overline K))$.  Thus this Tate module carries the canonical
three-step filtration (\SGARefLink{sga7.sga7i.eix.n2-3-4}{2.3.4}).  Applying the same considerations to
the dual abelian scheme $A'_K$, we obtain the analogous three-step
filtration
\begin{equation}\tag{\SGARefTarget{sga7.sga7i.eix.n2-3-5}{2.3.5}}
 T_\ell(A'(\overline K))^t\hookrightarrow
 T_\ell(A'(\overline K))^f\hookrightarrow
 T_\ell(A'(\overline K)).
\end{equation}
The middle terms in (\SGARefLink{sga7.sga7i.eix.n2-3-4}{2.3.4}) and (\SGARefLink{sga7.sga7i.eix.n2-3-5}{2.3.5}), the ``fixed parts,'' are
described in terms of the $\pi$-module structures of the Tate modules
of $A(\overline K)$ and $A'(\overline K)$ as their submodules of
invariants under the inertia group $I$ (\SGARefLink{sga7.sga7i.eix.n2-2-5}{2.2.5}).  The toric parts are
then completely described in terms of the fixed parts and the pairing
$\varphi$ of (\SGARefLink{sga7.sga7i.eix.n1-0-2}{1.0.2}), by the following theorem (a special case of which
already appears in [11]).

\medskip
\noindent
\underline{Theorem} \SGARefTarget{sga7.sga7i.eix.n2-4}{2.4}. (\underline{Orthogonality theorem}.)
\underline{Let $S$ be a henselian trait, let $A_K$ be an abelian scheme
over its fraction field $K$, let $A'_K$ be the dual abelian scheme, and
let $\ell$ be a prime different from the residual characteristic of
$S$.  The Tate modules $T_\ell(A(\overline K))$ and
$T_\ell(A'(\overline K))$ then carry the canonical filtrations (\SGARefLink{sga7.sga7i.eix.n2-3-4}{2.3.4})
and (\SGARefLink{sga7.sga7i.eix.n2-3-5}{2.3.5}), respectively, by the ``fixed part'' and the ``toric
part.''  The toric part $T_\ell(A(\overline K))^t$ is the intersection
of the fixed part $T_\ell(A(\overline K))^f$ with the orthogonal of the
fixed part
$T_\ell(A'(\overline K))^f=T_\ell(A'(\overline K))^I$ with respect to
the canonical pairing $\varphi$ of (\SGARefLink{sga7.sga7i.eix.n1-0-3}{1.0.3}).  Symmetrically, the toric
part $T_\ell(A'(\overline K))^t$ is the intersection of
$T_\ell(A'(\overline K))^f$ with the orthogonal of the fixed part
$T_\ell(A(\overline K))^f=T_\ell(A(\overline K))^I$.  (Here the
superscript $I$ denotes the submodule of invariants under the action of
the inertia group $I$.)}

\smallskip
\noindent
\underline{Proof.}\quad a) Let us prove that
$T_\ell(A(\overline K))^t$ is orthogonal to
$T_\ell(A'(\overline K))^f$.  To do so, we interpret the restriction of
$\varphi$ to the fixed parts of the Tate modules using the canonical
biextension $W^o=W$ of (\SGARefLink{sga7.sga7i.eix.n1-4-4}{1.4.4}).  This is the canonical extension of the
biextension $W_K$ of $(A_K,A'_K)$ by $\mathbb G_{mK}$ to a biextension
of $(A^o,A'{}^o)$ by $\mathbb G_{mS}$.  By VIII~\SGARefLink{sga7.sga7i.eviii.n2-2}{2.2}, $W$ gives rise to
a pairing
\begin{equation}\tag{\SGARefTarget{sga7.sga7i.eix.n2-4-1}{2.4.1}}
 \varphi_W:T_\ell(A^o)\times T_\ell(A'{}^o)
 \longrightarrow T_\ell(\mathbb G_{mS}),
\end{equation}
from which the pairing $\varphi$ of (\SGARefLink{sga7.sga7i.eix.n1-0-3}{1.0.3}) is obtained by passing to
$\overline K$-valued points.  It follows that the restriction of
$\varphi$ to the fixed parts considered in \SGARefLink{sga7.sga7i.eix.n2-4}{2.4} is obtained, by passing
to $\overline K$-valued points, from the pairing
\begin{equation}\tag{\SGARefTarget{sga7.sga7i.eix.n2-4-2}{2.4.2}}
 T_\ell(A^o)^f\times T_\ell(A'{}^o)^f
 \longrightarrow T_\ell(\mathbb G_{mS}),
\end{equation}
the restriction of (\SGARefLink{sga7.sga7i.eix.n2-4-1}{2.4.1}).  Since the two factors on the left of
(\SGARefLink{sga7.sga7i.eix.n2-4-2}{2.4.2}) are now twisted constant $\ell$-adic sheaves, this pairing may
be studied on the special fibers.  The fibers of the two factors are
precisely the Tate modules $T_\ell(A_o^o)$ and
$T_\ell(A'{}_o^o)$ (\SGARefLink{sga7.sga7i.eix.n2-2-3-3}{2.2.3.3}), and the pairing induced on special fibers
by (\SGARefLink{sga7.sga7i.eix.n2-4-2}{2.4.2}) is
\begin{equation}\tag{\SGARefTarget{sga7.sga7i.eix.n2-4-3}{2.4.3}}
 T_\ell(A_o^o)\times T_\ell(A'{}_o^o)
 \longrightarrow T_\ell(\mathbb G_{mk}),
\end{equation}
obtained from the biextension $W_o$ of $(A_o^o,A'{}_o^o)$ by
$\mathbb G_{mk}$, the restriction of $W$ to the special fibers.

Recall from \SGARefLink{sga7.sga7i.eix.n2-3}{2.3} that the toric part of
$T_\ell(A(\overline K))$ is obtained by passing to
$\overline K$-valued points in the twisted constant $\ell$-adic
subsheaf $T_\ell(A^o)^t$ of $T_\ell(A^o)^f$.  Thus the orthogonality
relation to be proved says that the pairing induced by (\SGARefLink{sga7.sga7i.eix.n2-4-2}{2.4.2}),
\begin{equation}\tag{\SGARefTarget{sga7.sga7i.eix.n2-4-4}{2.4.4}}
 T_\ell(A^o)^t\times T_\ell(A'{}^o)^f
 \longrightarrow T_\ell(\mathbb G_{mS}),
\end{equation}
is identically zero.  It is enough to verify this for the corresponding
pairing on the special fibers, which by the definition of the toric
part is
\begin{equation}\tag{\SGARefTarget{sga7.sga7i.eix.n2-4-5}{2.4.5}}
 T_\ell(T_o)\times T_\ell(A'{}_o^o)
 \longrightarrow T_\ell(\mathbb G_{mk}),
\end{equation}
induced by (\SGARefLink{sga7.sga7i.eix.n2-4-3}{2.4.3}).  By VIII~\SGARefLink{sga7.sga7i.eviii.n4-10}{4.10} this pairing is necessarily zero,
which proves the asserted orthogonality.

\smallskip
\noindent
b) Let us prove that the inclusion
\begin{equation}\tag{\SGARefTarget{sga7.sga7i.eix.n2-4-6}{2.4.6}}
 T_\ell(A(\overline K))^t\subset
 T_\ell(A(\overline K))^f\cap
 \bigl(T_\ell(A'(\overline K))^f\bigr)^\perp
\end{equation}
is an equality, where $\perp$ denotes the orthogonal.  Equivalently, in
the pairing (\SGARefLink{sga7.sga7i.eix.n2-4-2}{2.4.2}), the annihilator of the second factor---which, as
we have just seen, contains the toric part $T_\ell(A^o)^t$---is equal
to that toric part.  Since these are again relations between twisted
constant $\ell$-adic sheaves, it is enough to check the assertion on
the special fibers: namely, that the inclusion
\begin{equation}\tag{\SGARefTarget{sga7.sga7i.eix.n2-4-7}{2.4.7}}
 T_\ell(T_o)\subset T_\ell(A'{}_o^o)^\perp
\end{equation}
is an equality.  In other words, we must verify that the biextension
$W_o$ of $(A_o^o,A'{}_o^o)$ is \underline{nondegenerate on the left}
(VIII~\SGARefLink{sga7.sga7i.eviii.n4-11}{4.11}).  But this is precisely what we already observed in \SGARefLink{sga7.sga7i.eix.n1-5}{1.5},
as a consequence of a result from Raynaud's thesis.

% END INCLUDED SGA7I_English_current_buildable_source_20260823/components/134_expose_IX_sections_2_3_and_2_4.tex

% BEGIN INCLUDED SGA7I_English_current_buildable_source_20260823/components/135_expose_IX_section_2_5_and_remarks_2_6.tex
% Expose IX, section 2.5 and Remarks 2.6.
% French transcription lines 687--757; authority physical pp. 352--354
% (source folios 340--342), with the section-3 boundary checked on p. 355.
% Physical pp. 353--354 were checked in four overlapping 1100-dpi bands;
% the opening on p. 352 was checked with the preceding component.  The
% six-arrow filtration diagram (2.5.4) was additionally checked directly
% in five overlapping 5000-dpi crops, including hook sides, label sides,
% arrow directions, and the terminal comma.

\medskip
\noindent
\SGARefTarget{sga7.sga7i.eix.n2-5}{2.5}. The orthogonality theorem has an equivalent formulation using a
polarization $\xi$ of $A_K$ and the alternating bilinear form (\SGARefLink{sga7.sga7i.eix.n1-0-8}{1.0.8}),
compatible with the operations of $\pi$,
\begin{equation}\tag{\SGARefTarget{sga7.sga7i.eix.n2-5-1}{2.5.1}}
 \varphi_\xi:T_\ell(A(\overline K))\times T_\ell(A(\overline K))
 \longrightarrow T_\ell(\overline K^*),
\end{equation}
obtained from the form $\varphi$ of (\SGARefLink{sga7.sga7i.eix.n1-0-3}{1.0.3}) by transporting structure in
the second argument along the isogeny of $\mathbb Z_\ell$-modules
\[
 T_\ell(A(\overline K))\longrightarrow T_\ell(A'(\overline K))
\]
induced by the isogeny $\psi_\xi:A_K\longrightarrow A'_K$ defined by
$\xi$.  The equality in (\SGARefLink{sga7.sga7i.eix.n2-4-6}{2.4.6}) then takes the simple form
\begin{equation}\tag{\SGARefTarget{sga7.sga7i.eix.n2-5-2}{2.5.2}}
 W=V\cap V^\perp,
\end{equation}
where we have set
\begin{equation}\tag{\SGARefTarget{sga7.sga7i.eix.n2-5-3}{2.5.3}}
 \left\{
 \begin{array}{l}
 U=T_\ell(A(\overline K)),\qquad
 V=U^I=\text{the ``fixed part'' of }T_\ell(A(\overline K)),\\[2mm]
 \hspace{4em}W=\text{the toric part of }T_\ell(A(\overline K)),
 \end{array}
 \right.
\end{equation}
and $\perp$ denotes the orthogonal with respect to the polarization form
(\SGARefLink{sga7.sga7i.eix.n2-5-1}{2.5.1}).  These modules give the following diagram of inclusions:
\begin{equation}\tag{\SGARefTarget{sga7.sga7i.eix.n2-5-4}{2.5.4}}
\begin{tikzcd}[column sep=large,row sep=large]
 && V\arrow[dl,hook',"2\lambda"'] && \\
 U & W^\perp=V+V^\perp\arrow[l,hook',"\mu"']
 && W=V\cap V^\perp
    \arrow[ul,hook',"2\alpha"']
    \arrow[dl,hook,"2\lambda"]
 & 0\arrow[l,hook',"\mu"'] \\
 && V^\perp\arrow[ul,hook,"2\alpha"] &&
\end{tikzcd}
\qquad ,
\end{equation}
where the integers $2\alpha$, $2\lambda$, and $\mu$ are the ranks of
the quotients of consecutive modules in the diagram, with the notation
of (\SGARefLink{sga7.sga7i.eix.n2-1-9}{2.1.9}).  (That $W$ has rank $\mu$ is the definition of $\mu$; that
$V/W$ has rank $2\alpha$ follows from the second relation in (\SGARefLink{sga7.sga7i.eix.n2-1-11}{2.1.11})
together with the exact sequence (\SGARefLink{sga7.sga7i.eix.n2-1-6}{2.1.6}).  Hence, by orthogonality,
$\operatorname{rank}(U/W^\perp)=\mu$ and
$\operatorname{rank}(W^\perp/V^\perp)=2\alpha$.  Finally, the formulas
% The print has 2\lambda here.  Diagram (2.5.4), the preceding rank
% computation, and Proposition 3.5 show that the quotient V/W has rank
% 2\alpha; the complementary quotient V^\perp/W has rank 2\lambda.
\begin{equation}\tag{\SGARefTarget{sga7.sga7i.eix.n2-5-5}{2.5.5}}
 \operatorname{rank}\bigl(V/(V\cap V^\perp)\bigr)
 =\operatorname{rank}(V/W)=2\alpha
\end{equation}
and the dual formula follow from the preceding ones, in view of
(\SGARefLink{sga7.sga7i.eix.n2-1-10}{2.1.10}) and $\operatorname{rank}U=n$.)

\medskip
\noindent
\underline{Remarks} \SGARefTarget{sga7.sga7i.eix.n2-6}{2.6}.\quad a) Here we obtained the orthogonality
theorem \SGARefLink{sga7.sga7i.eix.n2-4}{2.4} as an easy consequence of the theory of biextensions
developed in the preceding sections.  The \underline{inclusion}
(\SGARefLink{sga7.sga7i.eix.n2-4-6}{2.4.6}) can be proved without using the theory of biextensions, by an
arithmetic argument in the style of the proof of the ``monodromy
theorem'' given in Expos\'e I.  One reduces to a situation in which $S$
has finite residue field, where the orthogonality result follows from a
simple consideration of \underline{weights} for the Frobenius operation
(using A.~Weil's result on the eigenvalues of that Frobenius).  There is
no need to invoke resolution of singularities here: one may state the
desired result for group schemes $A$ over schemes $S$ under sufficiently
general conditions.  Unfortunately, it does not appear possible to
obtain by this method equality in the preceding inclusion.  That
equality will apparently be used in an essential way in the proof of
the ``semistable reduction theorem'' in the next section.

\smallskip
\noindent
b) The orthogonality theorem \SGARefLink{sga7.sga7i.eix.n2-4}{2.4} allows the most important properties
of the N\'{e}ron model of an abelian variety $A_K$ to be expressed in
terms of the representation of the inertia group $I$ on the Tate module
$T_\ell(A(\overline K))$, where $\ell$ is a fixed prime different from
the residual characteristic.  (The case $\ell=p$ will be considered
separately later, in section 5.)  As a first statement of this kind,
though one that does not use the orthogonality theorem, recall the
good-reduction criterion \SGARefLink{sga7.sga7i.eix.n2-2-9}{2.2.9}.  We shall give many further examples
below, beginning in the next section (\SGARefLink{sga7.sga7i.eix.n3-5}{3.5}).

% END INCLUDED SGA7I_English_current_buildable_source_20260823/components/135_expose_IX_section_2_5_and_remarks_2_6.tex

% BEGIN INCLUDED SGA7I_English_current_buildable_source_20260823/components/136_expose_IX_section_3_1.tex
% Expose IX, section 3.1.
% French transcription lines 763--839; authority physical pp. 355--359
% (source folios 343--347), stopping before Proposition 3.2.
% Physical pp. 355--358 and the opening/boundary on p. 359 were checked
% directly in four overlapping 1100-dpi bands per page.  The text and the
% single hooked-arrow display were unambiguous at that resolution; no
% higher-detail escalation was needed.  The inherited French transcription's
% impossible ``t\in T'' in Lemma 3.1.3 is corrected to the scan-confirmed
% ``t\in S'' in the English text.

\medskip
\noindent
3. \underline{The case of semistable reduction.  The semistable reduction
theorem.}

We first give an auxiliary result in the style of SGA 3.

\medskip
\noindent
\underline{Proposition} \SGARefTarget{sga7.sga7i.eix.n3-1}{3.1}. \underline{Let $S$ be a scheme, let $G$ be a
commutative $S$-group scheme that is flat, separated, and of finite
presentation over $S$, let $H$ be a quasi-separated group scheme locally of
finite type over $S$, let $u:G\longrightarrow H$ be a homomorphism of
$S$-group schemes, and put $K=\ker u$.  For every $s\in S$, let $\bar s$
denote the spectrum of an algebraic closure of $k(s)$, and denote by
$\mu(s)$ (respectively, $\alpha(s)$) the reductive rank (respectively, the
abelian rank) of $K_s$, that is, the dimension of the largest subtorus
(respectively, of the largest abelian-variety quotient) of
$(K_{\bar s})^o_{\mathrm{red}}$.  Let $s\in S$, let $t\in S$ be a
generalization of $s$, and let $p$ be the characteristic exponent of
$k(s)$.}

\smallskip
\noindent
a) \underline{For every integer $n>0$ prime to $p$, ${}_nK$ is etale, of
finite type, and separated over $S$ in a neighborhood of $s$, and one has
the divisibility relation}
\begin{equation}\tag{\SGARefTarget{sga7.sga7i.eix.n3-1-1}{3.1.1}}
 \operatorname{rank}({}_nK_s)\mid \operatorname{rank}({}_nK_t).
\end{equation}

\noindent
b) \underline{One has}
\begin{equation}\tag{\SGARefTarget{sga7.sga7i.eix.n3-1-2}{3.1.2}}
 \mu(s)+2\alpha(s)\leq \mu(t)+2\alpha(t).
\end{equation}

\noindent
c) \underline{If $(K_{\bar t})^o_{\mathrm{red}}$ is unipotent, then so is
$(K_{\bar s})^o_{\mathrm{red}}$.}

\smallskip
\noindent
d) \underline{Suppose that $G_s$ has unipotent rank zero, that $H$ is smooth
over $S$, and that $u_t$ is an isogeny.  Then $u$ is flat and quasi-finite
over an open neighborhood $U$ of $s$; a fortiori, $K|_U$ is flat and
quasi-finite.}

\smallskip
\noindent
e) \underline{Under the hypotheses of d), suppose in addition that $u_t$ is
a monomorphism (and hence an isomorphism from $G_t$ onto an open subgroup of
$H_t$).  Then $u_s$ is an isomorphism from $G_s$ onto an open subgroup of
$H_s$.  If, moreover, $S$ is irreducible with generic point $t$, then there
is an open neighborhood $U$ of $s$ such that $u$ induces an isomorphism from
$G|_U$ onto an open subgroup of $H|_U$.}

\medskip
\noindent
\underline{Proof.}

\smallskip
\noindent
a) We may suppose that $n$ is invertible on $S$.  Then $n\,\mathrm{id}_G$
is etale on every fibre (SGA 3, XV, 1.3); since $G$ is flat and of finite
presentation over $S$, the fibrewise criterion for etaleness (EGA IV,
17.8.2) shows that $n\,\mathrm{id}_G$ is etale.  Its kernel ${}_nG$ is
therefore etale; it is also separated and of finite presentation over $S$,
since $G$ is.  Similarly, ${}_nH$ is locally of finite type,
quasi-separated, and formally unramified over $S$,
\footnote{At least when $H$ is commutative.  In the general case, ${}_nH$
is still unramified over $S$ along its unit section, which is all that is
needed here.}
so its unit section is an open immersion (EGA IV, 17.4.2 b')).  Consequently,
the kernel of
\[
 {}_nu:{}_nG\longrightarrow {}_nH
\]
is an open subscheme of ${}_nG$, hence is etale over $S$; it is
quasi-compact in ${}_nG$ and therefore of finite presentation over $S$, and
it is separated over $S$ because $G$ is.  This kernel is plainly ${}_nK$.

Formula (\SGARefLink{sga7.sga7i.eix.n3-1-1}{3.1.1}) follows from these properties of ${}_nK$.  Indeed, after a
base change we may reduce to the case where $S$ is henselian local with
closed point $s$, and then use the decomposition of ${}_nK$ into
${}_nK^f$ and a scheme with empty special fibre (cf. \SGARefLink{sga7.sga7i.eix.n2-2-3}{2.2.3}).  More precisely,
the same argument gives a monomorphism on geometric fibres
\[
 {}_nK(\bar s)\lhook\joinrel\longrightarrow {}_nK(\bar t).
\]

\smallskip
\noindent
b) Choose an integer $n>0$ prime to the characteristic exponent of $k(s)$
and to the ranks of the quotients of the geometric fibres of $K$ at $s$ and
$t$ by their reduced identity components.  The two sides of (\SGARefLink{sga7.sga7i.eix.n3-1-1}{3.1.1}) are
then obtained by raising $n$, respectively, to the powers whose exponents
are the two sides of (\SGARefLink{sga7.sga7i.eix.n3-1-2}{3.1.2}), as follows from the known structure of
commutative algebraic groups.  Thus b) follows from a).

\smallskip
\noindent
c) The hypothesis that $(K_{\bar t})^o_{\mathrm{red}}$ is unipotent is
equivalent to the equalities $\mu(t)=\alpha(t)=0$, and likewise for
$K_{\bar s}$.  Assertion c) therefore follows formally from b).

\smallskip
\noindent
d) Since $u_t$ is an isogeny, $K_t$ is finite.  We are therefore under the
hypotheses of c), and $(K_{\bar s})^o_{\mathrm{red}}$ is unipotent.  As
$G_s$ has unipotent rank zero, it follows that $K_s$ is finite.  Moreover,
the fact that $u_t$ is an isogeny gives
$\dim G_t=\dim H_t$.  Since $G$ and $H$ are flat and locally of finite
presentation, one has
\[
 \dim G_s=\dim G_t,\qquad \dim H_s=\dim H_t
 \quad\text{(SGA 3, VI$_B$, 4.3)},
\]
and hence $\dim G_s=\dim H_s$.  Since $\dim\ker u_s=0$, it follows that
$\dim H_s=\dim u_s(G_s)$ (SGA 3, VI$_B$, 1.2), which, because $H_s$ is
smooth, is equivalent to $u_s$ being flat (SGA 3, VI$_A$, 5.6).  There is
therefore an open neighborhood $U$ of $s$ such that $u$ is flat and
quasi-finite at the points of the unit section over $U$ (SGA 3, VI$_B$,
2.2(i) and \SGARefLink{sga7.sga7i.eix.n2-5}{2.5}(ii)).  As for $u_s$, it follows that $u_{s'}$ is flat for
every $s'\in U$; hence $u|_{G_U}$ is flat and quasi-finite, since it has
these properties on every fibre.

\smallskip
\noindent
e) For the first assertion, a base change immediately reduces us to the
case where $S$ is integral with generic point $t$, and hence to proving the
second assertion of e).  That assertion follows at once from d) and the
following lemma.

\medskip
\noindent
\underline{Lemma} \SGARefTarget{sga7.sga7i.eix.n3-1-3}{3.1.3}. \underline{Let $S$ be a scheme, let $K$ be a
quasi-finite $S$-scheme of finite presentation that is separated and flat
over $S$, let $s\in S$, and let $t\in S$ be a generalization of $s$.  Then}
\[
 \operatorname{rank}K_s\leq \operatorname{rank}K_t.
\]
\underline{If $\operatorname{rank}K_\eta\leq1$ for every maximal point
$\eta$ of $S$, then $K\longrightarrow S$ is an open immersion (and hence
an isomorphism if it admits a section).}

For the first assertion, strict localization reduces us to the case where
$S$ is strictly local.  Decomposing $K$ as the disjoint union of a finite
scheme $K^f$ and a scheme with empty special fibre, we obtain
\[
 \operatorname{rank}K_s
 =\operatorname{rank}K^f_s
 =\operatorname{rank}K^f_t
 \leq\operatorname{rank}K_t.
\]
The second equality follows because $K^f$ is flat, finite, and of finite
presentation over $S$, and therefore locally free over $S$.

For the second assertion, the preceding argument already shows that
$\operatorname{rank}K_s\leq1$ for every $s\in S$.  Thus $K_s$ is empty or
isomorphic to $\operatorname{Spec}k(s)$.  Consequently,
$K\longrightarrow S$ is a monomorphism (SGA 3, VI$_B$, 2.11), and, being
flat and of finite presentation, it is an open immersion (EGA IV, 17.9.1).
This proves the lemma.

\medskip
\noindent
\underline{Remark} \SGARefTarget{sga7.sga7i.eix.n3-1-4}{3.1.4}.  Part \SGARefLink{sga7.sga7i.eix.n3-1}{3.1} a), and the remainder of Proposition
\SGARefLink{sga7.sga7i.eix.n3-1}{3.1}, give no information about ${}_nK_s$ when $n$ is a power of $p$ (for
$p>1$).  At least when $G$ is smooth (and $S$ is locally noetherian or $H$
is of finite presentation over $S$), essentially the same method as in the
proof of a) gives such information.  Indeed, reducing to the case where
$S$ is complete local noetherian and using SGA 3, XV, 1.6 c), X, 3.2, IX,
\SGARefLink{sga7.sga7i.eix.n3-6}{3.6}, and IX, \SGARefLink{sga7.sga7i.eix.n5-2}{5.2}, one obtains
\begin{equation}\tag{\SGARefTarget{sga7.sga7i.eix.n3-1-4-1}{3.1.4.1}}
 \operatorname{rank}_{\mathrm{mult}}({}_nK_s)
 \mid
 \operatorname{rank}_{\mathrm{mult}}({}_nK_t),
\end{equation}
where $\operatorname{rank}_{\mathrm{mult}}$ denotes the rank of the largest
subgroup of multiplicative type of the commutative algebraic group under
consideration (SGA 3, XVII, 7.2.1).  Notice that a torsion algebraic group
of multiplicative type is finite.  Formula (\SGARefLink{sga7.sga7i.eix.n3-1-4-1}{3.1.4.1}) is valid for every
integer $n\geq1$.  It follows from SGA 3, XVII, 4.6.1(vi) that, in c), if
$K_t$ is unipotent (SGA 3, VII, 1.1), then $K_s$ is unipotent as well.

% END INCLUDED SGA7I_English_current_buildable_source_20260823/components/136_expose_IX_section_3_1.tex

% BEGIN INCLUDED SGA7I_English_current_buildable_source_20260823/components/137_expose_IX_sections_3_2_through_3_4.tex
% Expose IX, sections 3.2--3.4.
% French transcription lines 841--875; authority physical pp. 359--361
% (source folios 347--349), stopping before Proposition 3.5.
% The opening on physical p. 359 was checked with the preceding component;
% physical pp. 360--361 were checked directly in four overlapping 1100-dpi
% bands per page.  The prose and formulas were unambiguous at that resolution,
% so no higher-detail escalation was needed.

\medskip
\noindent
\underline{Proposition} \SGARefTarget{sga7.sga7i.eix.n3-2}{3.2}. \underline{Let $S$ be a regular integral
locally noetherian scheme of dimension $1$, let $K$ be its field of rational
functions, let $A_K$ be an abelian scheme over $K$, and let $A$ be its Neron
model over $S$, which is a smooth separated group scheme over $S$ (and is
quasi-projective over $S$ if $S$ is noetherian [23, VIII, 2, 3\textdegree]).
The following conditions are equivalent:}

\smallskip
\noindent
(i) \underline{For every $s\in S$, $A_s^o$ has unipotent rank zero; that is,
it is an extension of an abelian scheme by a torus (\SGARefLink{sga7.sga7i.eix.n2-1-5}{2.1.5}).}

\smallskip
\noindent
(ii) \underline{There exists a smooth separated group scheme $G$ of finite
type over $S$, extending $A_K$, whose fibres have unipotent rank zero.}

\smallskip
\noindent
\underline{Moreover, if $G$ is as in (ii), then the unique morphism
$G\longrightarrow A$ inducing the given isomorphism on the generic fibres
is an isomorphism from $G$ onto an open subgroup of $A$; it therefore induces
an isomorphism of identity components}
\[
 G^o\xrightarrow{\ \sim\ }A^o.
\]

\medskip
\noindent
\SGARefTarget{sga7.sga7i.eix.n3-2-1}{3.2.1}.  Let us note at once that the definition of the Neron model $A$ of
$A_K$ is the same as in the local case recalled in \SGARefLink{sga7.sga7i.eix.n1-1}{1.1}.  The existence of
this Neron model, and the fact that it is separated and of finite type over
$S$, are well known and in fact reduce immediately to the local case.  Of
course, for every closed point $s\in S$,
\[
 S'=\operatorname{Spec}(\mathcal O_{S,s})
\]
is a trait, and $A\times_S S'$ is precisely the Neron model of $A_K$ over
$S'$ in the sense considered in \SGARefLink{sga7.sga7i.eix.n1-1}{1.1}.

Let us prove \SGARefLink{sga7.sga7i.eix.n3-2}{3.2}.  It is trivial that (i) implies (ii), since it suffices to
take $G=A$.  It remains to prove that if $G$ is as in (ii), then
$G\longrightarrow A$ is an open immersion.  This is indeed a special case
of \SGARefLink{sga7.sga7i.eix.n3-1}{3.1} e).

\medskip
\noindent
\underline{Corollary} \SGARefTarget{sga7.sga7i.eix.n3-3}{3.3}. \underline{With the notation of \SGARefLink{sga7.sga7i.eix.n3-2}{3.2}, suppose
that the identity components of the fibres of the Neron model have
unipotent rank zero.  Let $S'$ be a scheme satisfying the same hypotheses
as $S$, let $K'$ be its field of rational functions, let
$f:S'\longrightarrow S$ be a dominant morphism, let}
\[
 A_{K'}\simeq A_K\otimes_KK'
 \simeq \text{the generic fibre of }A\times_SS',
\]
\underline{let $A'$ be the Neron model of $A_{K'}$, and consider the unique
morphism}
\[
 u:A\times_SS'\longrightarrow A'
\]
\underline{inducing the identity on the generic fibres.  Then $u$ is an
open immersion and therefore induces an isomorphism on identity
components}
\begin{equation}\tag{\SGARefTarget{sga7.sga7i.eix.n3-3-1}{3.3.1}}
 A^o\times_SS'\xrightarrow{\ \sim\ }{A'}^o.
\end{equation}
Indeed, it suffices to apply \SGARefLink{sga7.sga7i.eix.n3-2}{3.2} over $S'$, taking
$G\simeq A\times_SS'$.

\medskip
\noindent
\SGARefTarget{sga7.sga7i.eix.n3-3-2}{3.3.2}.  Thus the equivalent properties of \SGARefLink{sga7.sga7i.eix.n3-2}{3.2} imply a partial
``stability'' property for Neron models.  One should bear in mind that when
$A_K$ does not have good reduction, that is, when $A$ is not an abelian
scheme, it is not true that, for every $S'$ as in \SGARefLink{sga7.sga7i.eix.n3-3}{3.3} (even for $S'$ finite
over $S$ and $K'/K$ separable), the morphism
\[
 A\times_SS'\longrightarrow A'
\]
itself is an isomorphism, or that it induces isomorphisms on the groups of
connected components of the fibres (see \SGARefLink{sga7.sga7i.eix.n11-10}{11.10} below for the behavior of
these groups of connected components under extension of the base).  These
observations, reinforced by \SGARefLink{sga7.sga7i.eix.n3-9}{3.9} below, justify the following terminology,
introduced by Mumford (and which also appears to agree with the notions of
the same name studied in his book [17]).

\medskip
\noindent
\underline{Definition} \SGARefTarget{sga7.sga7i.eix.n3-4}{3.4}. \underline{With the notation of \SGARefLink{sga7.sga7i.eix.n3-2}{3.2}, the
abelian scheme $A_K$ over $K$ is said to have semistable reduction over $S$
if the equivalent conditions (i) and (ii) of \SGARefLink{sga7.sga7i.eix.n3-2}{3.2} are satisfied.}

\medskip
\noindent
\SGARefTarget{sga7.sga7i.eix.n3-4-0}{3.4.0}.  Notice that in the form (ii), this condition does not invoke the
notion of a Neron model.  Moreover, by localization it reduces immediately
to the case where $S$ is a trait, as follows for example from \SGARefLink{sga7.sga7i.eix.n3-2-1}{3.2.1}.
Furthermore, since formation of the Neron model commutes with base change
from $S$ to its henselization or its strict henselization, the notion even
reduces to the case where $S$ is a strictly local trait.

% END INCLUDED SGA7I_English_current_buildable_source_20260823/components/137_expose_IX_sections_3_2_through_3_4.tex

% BEGIN INCLUDED SGA7I_English_current_buildable_source_20260823/components/138_expose_IX_sections_3_5_through_3_9.tex
% Expose IX, Proposition 3.5 through Corollary 3.9.
% French transcription lines 880--960; authority physical pp. 362--366
% (source folios 350--354); hard stop before section 4 on p. 366.
% All five pages were checked directly in four overlapping 1100-dpi bands.
% The empty parenthesis printed after Z_l(1) in 3.5.2 contains no recoverable
% text and is omitted rather than conjecturally supplied.
% Transparent source corrections:
% - Proposition 3.5(ii) prints V \subset V^\perp.  The filtration diagram
%   (2.5.4), condition (iii), and the proof require V^\perp \subset V.
% - In the proof, the print says tautologically that I acts trivially on
%   U^I=V; condition (iv) says that it acts trivially on U/U^I=U/V.
% - The proof of Corollary 3.8 prints (v) => (vi), whereas the argument proves
%   the required converse (vi) => (iv).

\medskip
\noindent
\underline{Proposition} \SGARefTarget{sga7.sga7i.eix.n3-5}{3.5}. \underline{(Galois criterion for semistable
reduction).  Suppose that $S$ is a trait.  The following conditions on the
abelian scheme $A_K$ are equivalent} (the notation being that of (\SGARefLink{sga7.sga7i.eix.n2-5-3}{2.5.3}),
applied to the situation deduced from $(A_K,S)$ by passing to the
henselization of $S$):

\smallskip
\noindent
(i) \underline{$A_K$ has semistable reduction over $S$} (\SGARefLink{sga7.sga7i.eix.n3-4}{3.4}).

\smallskip
\noindent
(ii) \underline{With the notation of (\SGARefLink{sga7.sga7i.eix.n2-5-3}{2.5.3}) and (\SGARefLink{sga7.sga7i.eix.n2-5-4}{2.5.4}), one has}
\[
 V^\perp\subset V.
\]

\noindent
(iii) \underline{For the toric part $W$ of the Tate module one has}
\[
 W=V^\perp.
\]

\noindent
(iv) \underline{The action of the inertia group $I$ on
$U=T_\ell(A(\overline K))$ is unipotent of level $2$} (that is, there is an
$I$-stable submodule $U'$ of $U$ such that $I$ acts trivially on $U'$ and
on $U/U'$).

We may suppose that $S$ is henselian (\SGARefLink{sga7.sga7i.eix.n3-4-0}{3.4.0}).  With the notation of \SGARefLink{sga7.sga7i.eix.n2-1-9}{2.1.9},
(i) means that
\[
 \lambda=0.
\]
By the rank formulas following (\SGARefLink{sga7.sga7i.eix.n2-5-4}{2.5.4}), this is equivalent to
$V^\perp=V\cap V^\perp$, that is, to (ii).  Since
$W=V\cap V^\perp$ by the orthogonality theorem \SGARefLink{sga7.sga7i.eix.n2-4}{2.4}, conditions (ii) and
(iii) are equivalent.  Moreover, (iv) also says that $I$ acts trivially
on $U/U^I=U/V$.  By duality this means that it acts trivially on
$V^\perp$, that is, that $V^\perp\subset U^I=V$, which is (ii).  This
completes the proof.

\medskip
\noindent
\underline{Remarks} \SGARefTarget{sga7.sga7i.eix.n3-5-1}{3.5.1}.  Conditions (ii) and (iv) plainly depend only
on the action of $I$ on the $\mathbb Q_\ell$-vector space
$U\otimes_{\mathbb Z_\ell}\mathbb Q_\ell$, and could have been formulated
in terms of that action.  Let us also note that these conditions are
manifestly stable under every base change by a morphism of traits
$S'\longrightarrow S$.  Finally, if (iv) holds for $A_K$, it plainly holds
for every abelian subscheme and every quotient abelian scheme of $A_K$,
and in particular for every abelian scheme isogenous to $A_K$ (for
example, for $A'_K$).

\medskip
\noindent
\underline{Corollary} \SGARefTarget{sga7.sga7i.eix.n3-5-2}{3.5.2}. \underline{Let $g$ be a topological generator
of $\mathbb Z_\ell(1)=T_\ell(\overline k^{\,*})$.  The equivalent conditions
of \SGARefLink{sga7.sga7i.eix.n3-5}{3.5} are also equivalent to the following condition:}

\smallskip
\noindent
% SGA-SRC-001269 — internal reference omitted in source; target not recoverable
(v) \underline{The action of the inertia group $I$ on
$U=T_\ell(A(\overline K))$ factors through the maximal pro-$\ell$ quotient
$I(\ell)$ of $I$, which is canonically isomorphic to $\mathbb Z_\ell(1)$,
and the action $g_U$ of the generator $g$ on $U$ satisfies}
\begin{equation}\tag{\SGARefTarget{sga7.sga7i.eix.n3-5-3}{3.5.3}}
 (1-g_U)^2=0.
\end{equation}
Indeed, for every continuous unipotent action of a profinite group $I$ on
an $\ell$-adic vector space, one sees at once that the image of $I$ is a
pro-$\ell$ group.  It follows that (v) is equivalent to (iv).

We can now prove the central result of the present expose:

\medskip
\noindent
\underline{Theorem} \SGARefTarget{sga7.sga7i.eix.n3-6}{3.6}. \underline{(Semistable reduction theorem).  Let
$S$ be a connected regular noetherian scheme of dimension $1$, let $K$ be
the field of rational functions of $S$, and let $A_K$ be an abelian scheme
over $K$.  Then there is a finite Galois extension $K'$ of $K$ such that
the abelian scheme $A_{K'}=A_K\otimes_KK'$ has semistable reduction} (\SGARefLink{sga7.sga7i.eix.n3-4}{3.4})
\underline{over the normalization $S'$ of $S$ in $K'$.}

Notice that $S'$ is connected, regular, noetherian, and of dimension $1$,
so that the stated conclusion makes sense.  To prove \SGARefLink{sga7.sga7i.eix.n3-6}{3.6}, choose a
separable closure $\overline K$ of $K$ and write $\overline K$ as the
filtered union of its finite Galois subextensions $K_i$.  Let $S_i$ be the
normalization of $S$ in $K_i$, and let $U$ be the largest open subset of
$S$ over which $A_K$ has good reduction.  Since $S$ is connected,
noetherian, and one-dimensional, and $U\ne\varnothing$, the complement
$T=S-U$ is a finite set of closed points of $S$, and
$\mathcal O_{S,t}$ is a discrete valuation ring for every $t\in T$.  It is
enough to find, for every $t\in T$, an index $i=i(t)$ such that $A_{K_i}$
has semistable reduction at the points of $S_i$ above $t$: one then takes
for $K'$ the compositum of the extensions $K_{i(t)}$, for $t\in T$.
Thus we are reduced to the case where $S$ is a trait.

Moreover, since the Galois group permutes the points of the normalization
lying above the closed point $s$ of $S$, it is enough to find $K'$ such
that $A_{K'}$ has semistable reduction at \emph{one} point of the
normalization $S'$.  Here it is unnecessary at first to require that
$K'/K$ be Galois: it is enough to find a finite etale extension $K'/K$
satisfying the preceding condition, since its Galois closure satisfies the
same condition.  Using criterion \SGARefLink{sga7.sga7i.eix.n3-5}{3.5}(iv), we see that Theorem \SGARefLink{sga7.sga7i.eix.n3-6}{3.6} reduces
to the following assertion.

\medskip
\noindent
\underline{Corollary} \SGARefTarget{sga7.sga7i.eix.n3-7}{3.7}. \underline{Suppose again that $S$ is a trait.
Then there is an open subgroup $\pi'$ of $\pi$ such that the action of
$I'=I\cap\pi'$ on $U=T_\ell(A(\overline K))$ is unipotent of level $2$;
equivalently, there is an open subgroup $I'$ of $I$ whose action on $U$ is
unipotent of level $2$.}

There is a canonical isomorphism, compatible with the actions of $\pi$,
\begin{equation}\tag{*}
 H^1(A_{\overline K},\mathbb Z_\ell)
 \simeq
 \operatorname{Hom}\bigl(T_\ell(A(\overline K)),\mathbb Z_\ell\bigr),
\end{equation}
which shows that assertion \SGARefLink{sga7.sga7i.eix.n3-7}{3.7} is equivalent to the same assertion for
the actions of $I$ on the left-hand side of (*).  But that assertion was
proved in Expose III (using resolution of singularities for excellent
schemes of dimension $2$).

\medskip
\noindent
\underline{Corollary} \SGARefTarget{sga7.sga7i.eix.n3-8}{3.8}. \underline{The equivalent conditions (i)--(iv)
of \SGARefLink{sga7.sga7i.eix.n3-5}{3.5} are also equivalent to the following condition:}

\smallskip
\noindent
(vi) \quad \underline{The action of $I$ on
$U=T_\ell(A(\overline K))$ is unipotent.}

Indeed, to prove that (vi) implies (iv), observe that if a group $I$
acting unipotently on a finite-dimensional vector space $E$ over a field
of characteristic zero (here $\mathbb Q_\ell$) has a subgroup $I'$ of
finite index acting with level $m$ (here $m=2$), then $I$ itself acts with
level $m$.  In fact, since $I'$ has finite index in $I$, the algebraic
group $G'\subset\operatorname{GL}(E)$ that it defines has finite index in
the group $G$ defined by $I$; but $G$ is unipotent and hence connected (in
characteristic zero), so $G=G'$.  Thus $I$ and $I'$ have the same
``level of unipotence.''

We also obtain a converse to \SGARefLink{sga7.sga7i.eix.n3-3}{3.3}:

\medskip
\noindent
\underline{Corollary} \SGARefTarget{sga7.sga7i.eix.n3-9}{3.9}. \underline{The notation is that of \SGARefLink{sga7.sga7i.eix.n3-2}{3.2}.  The
conditions (i) and (ii) of \SGARefLink{sga7.sga7i.eix.n3-2}{3.2} are also equivalent to the following
condition:}

\smallskip
\noindent
(iii) \quad \underline{For every finite etale extension $K'$ of $K$, let
$S'$ denote the normalization of $S$ in $K'$, and let $A'$ denote the
Neron model of
$A_{K'}=A\otimes_SK'=A_K\otimes_KK'$.  Then the canonical morphism}
\[
 A\times_SS'\longrightarrow A'
\]
\underline{induces an isomorphism on identity components.}

The necessity of the condition was proved, in a stronger form, in \SGARefLink{sga7.sga7i.eix.n3-3}{3.3}.
For sufficiency, choose $K'$ as in \SGARefLink{sga7.sga7i.eix.n3-6}{3.6}.  Since the fibres of ${A'}^o$ have
unipotent rank zero, the same is true of the fibres of
$(A\times_SS')^o=A^o\times_SS'$, and hence also of the fibres of $A^o$,
from which the former are obtained by extension of the base field.

% END INCLUDED SGA7I_English_current_buildable_source_20260823/components/138_expose_IX_sections_3_5_through_3_9.tex

% BEGIN INCLUDED SGA7I_English_current_buildable_source_20260823/components/139_expose_IX_sections_4_1_and_4_2.tex
% Expose IX, sections 4.1 and 4.2.
% French transcription lines 963--1093; authority physical pp. 366--370
% (source folios 354--358); hard stop before Theorem 4.3 on p. 370.
% All five pages were checked directly in four overlapping 1100-dpi bands.
% Diagram (4.2.4) was additionally checked in four overlapping 5000-dpi
% crops: all six arrowheads, all three unheaded extension-field segments,
% the two curved segments, node placement, and label sides were compared.
% Transparent source corrections:
% - The footnote to 4.1 prints ell != p, although the main text already
%   assumes ell != p and sections 5--6 treat precisely the case ell = p.
% - The printed reference to the semistable-reduction criterion as (4.5)
%   is corrected to (3.5), where that criterion is stated and proved.
% - After (4.2.11), the source prints the undefined Gamma^o, whereas the
%   exact sequence (4.2.3) makes Gamma the kernel of pi' -> pi_o; hence the
%   displayed character is trivial on Gamma.
% - The parenthetical reference following (4.2.6) is blank after "I" in
%   the source.  It is identified as omitted rather than conjecturally filled.

\medskip
\noindent
\underline{4. Application to a conjecture of Serre--Tate and to the
conductor}

\medskip
\noindent
\SGARefTarget{sga7.sga7i.eix.n4-1}{4.1}. We suppose that $S$ is henselian, and retain the notation of (\SGARefLink{sga7.sga7i.eix.n2-5-3}{2.5.3}),
with $\ell\ne p$.
\footnote{The case $\ell=p$ will be studied in \S\S5 and 6 below.}
To study the action of $\pi=\operatorname{Gal}(\overline K/K)$ and of its
inertia subgroup $I$ on $U=T_\ell(A(\overline K))$, we shall use the fact
that there exists an open subgroup $I'$ of $I$, which we may moreover take
to be normal in $\pi$, such that the action of $I'$ on $U$ is unipotent of
level two; that is,
\[
 U^{I'}\supset (U^{I'})^\perp.
\]
It is then clear that $U^{I'}$ does not depend on the particular choice of
$I'$.  We might call it the \underline{essentially fixed submodule} of
$U$, and shall denote it by $U^{\mathrm{ef}}$.  Its orthogonal will be
called the \underline{essentially toric submodule} of $U$, and will be
denoted by $U^{\mathrm{et}}$.  We thus have a three-step filtration
\begin{equation}\tag{\SGARefTarget{sga7.sga7i.eix.n4-1-1}{4.1.1}}
 U\supset U^{\mathrm{ef}}\supset U^{\mathrm{et}}\supset0,
\end{equation}
which is plainly invariant under the action of $\pi$.

Let us also note that there is a largest subgroup of $I$ acting
unipotently (or, equivalently here, unipotently of level two) on $U$: it
is the subgroup $I'$ consisting of the elements acting unipotently.  It is
also the kernel of the representation of $I$ on the graded module
associated with the filtration (\SGARefLink{sga7.sga7i.eix.n4-1-1}{4.1.1}), or, equivalently in the present
case, the kernel of the induced representation of $I$ on
$U^{\mathrm{ef}}$.  By the Galois criterion for semistable reduction
(\SGARefLink{sga7.sga7i.eix.n3-5}{3.5}), this subgroup $I'$ corresponds to \underline{the smallest Galois
extension} $\widetilde K'$ of $\widetilde K$ (the maximal unramified
extension of $K$) \underline{such that $A_{\widetilde K'}$ has
semistable reduction}.  This description also shows that the subgroup
$I'$ of $I$ is independent of the choice of the prime $\ell\ne p$.

Notice that, once the representation of $\pi$ on $U^{\mathrm{ef}}$, and
hence on $U^{\mathrm{et}}$, is known, the representation on
$U/U^{\mathrm{ef}}$ is known up to isogeny, by means of the isogeny
induced by the polarization pairing (\SGARefLink{sga7.sga7i.eix.n2-5-1}{2.5.1}):
\begin{equation}\tag{\SGARefTarget{sga7.sga7i.eix.n4-1-2}{4.1.2}}
 U/U^{\mathrm{ef}}\longrightarrow (U^{\mathrm{et}})^\vee(1)
 \qquad\text{an isogeny},
\end{equation}
where $(1)$ denotes the Tate twist by
$\mathbb Z_\ell(1)=T_\ell(\overline K^*)$; this homomorphism is plainly
compatible with the operations of $\pi$.  More precisely, there is a
canonical \underline{isomorphism}
\[
 U/U^{\mathrm{ef}}\simeq ((U')^{\mathrm{et}})^\vee(1),
 \qquad U'=T_\ell(A'(\overline K)).
\]
It therefore remains to study more closely the action of $\pi$ on
$U^{\mathrm{ef}}$.

By construction this action is trivial on the normal subgroup $I'$ of
$I$, and hence factors through a representation of the quotient group
\begin{equation}\tag{\SGARefTarget{sga7.sga7i.eix.n4-1-3}{4.1.3}}
 \pi'=\pi/I',
\end{equation}
an action of which we shall now give a description ``independent of
$\ell$.''

\medskip
\noindent
\SGARefTarget{sga7.sga7i.eix.n4-2}{4.2}. Let $K'$ be a Galois extension of $K$ corresponding to an open
normal subgroup $V$ of $\pi$ such that $V\cap I=I'$.  The Galois group
$G$ of $K'/K$ then occurs in an exact sequence
\begin{equation}\tag{\SGARefTarget{sga7.sga7i.eix.n4-2-1}{4.2.1}}
 1\longrightarrow\Gamma\longrightarrow G\longrightarrow G_o
 \longrightarrow1,
\end{equation}
where
\begin{equation}\tag{\SGARefTarget{sga7.sga7i.eix.n4-2-2}{4.2.2}}
 \Gamma=I/I'
\end{equation}
is the inertia subgroup, and $G_o$ is the Galois group of the largest
unramified subextension $K''$ of $K'/K$.  Thus (\SGARefLink{sga7.sga7i.eix.n4-2-1}{4.2.1}) is the image of the
analogous exact sequence
\begin{equation}\tag{\SGARefTarget{sga7.sga7i.eix.n4-2-3}{4.2.3}}
 1\longrightarrow\Gamma\longrightarrow\pi'\longrightarrow\pi_o
 \longrightarrow1,
\end{equation}
where $\pi'$ is the group (\SGARefLink{sga7.sga7i.eix.n4-1-3}{4.1.3}) acting on $U^{\mathrm{ef}}$.  We
therefore have the following diagram of field inclusions:
\begin{equation}\tag{\SGARefTarget{sga7.sga7i.eix.n4-2-4}{4.2.4}}
\begin{tikzpicture}[baseline=(current bounding box.center),
  every node/.style={inner sep=2pt}]
  \node (K) at (0,0) {$K$};
  \node (Ktwo) at (4.25,0) {$K''=\widetilde K\cap K'$};
  \node (Kt) at (8.5,0) {$\widetilde K$};
  \node (Kp) at (4.25,-1.8) {$K'$};
  \node (Ktp) at (8.5,-1.8) {$\widetilde K'$};
  \node (Kb) at (8.5,-3.6) {$\overline K$};
  \draw[->] (K) -- node[above] {$G_o$} (Ktwo);
  \draw[->] (Ktwo) -- (Kt);
  \draw[->] (Ktwo) -- node[right] {$\Gamma$} (Kp);
  \draw[->] (Kp) -- (Ktp);
  \draw[->] (Kt) -- node[right] {$\Gamma$} (Ktp);
  \draw[->] (Ktp) -- node[right] {$I'$} (Kb);
  \draw[-] (K) -- node[below left,pos=.58] {$G$} (Kp);
  \draw[-] (K) to[out=65,in=115,looseness=.85]
    node[above,pos=.78] {$\pi_o$} (Kt);
  \draw[-] (Kt) to[out=12,in=12,looseness=1.25]
    node[right,pos=.52] {$I$} (Kb);
\end{tikzpicture}
\end{equation}

Let $S'$ be the normalization of $S$ in $K'$, and let $A^*$ be the
Neron model of $A_{K'}$.  Then, by the Galois criterion for semistable
reduction (\SGARefLink{sga7.sga7i.eix.n3-5}{3.5}), $A_{K'}$ \underline{has semistable reduction over
$S'$}; that is, $A^{*o}_o$ \underline{is an extension of an abelian
scheme by a torus}.  On the other hand, by transport of structure, $G$
\underline{acts on $A^*$, and hence on $A^{*o}$, by automorphisms in a
manner compatible with the given action of $G$ on the base scheme $S'$
of $A^*$}.  It is then clear that \underline{the action of $\pi'$ on}
\[
 U^{\mathrm{ef}}=T_\ell(A^{*o})^f(\overline K)
\]
\underline{is obtained by transport of structure from the action of
$\pi'$ on $(A^{*o}/K',\widetilde K')$}, where $\pi'$ acts on $A^*$
through $G$.  Taking into account the canonical isomorphism
\begin{equation}\tag{\SGARefTarget{sga7.sga7i.eix.n4-2-6}{4.2.6}}
 U^{\mathrm{ef}}\simeq T_\ell(A^{*o})^f(\overline K)
 \simeq T_\ell(A^{*o}_o)(\overline k)
\end{equation}
% SGA-SRC-001270 — internal reference omitted in source; target not recoverable
(valid for every twisted constant $\ell$-adic sheaf over $S'$; cf.
Expos\'e I [reference omitted in the source]), we obtain the following
description of the action of $\pi'$ on $U^{\mathrm{ef}}$.

Consider the smooth connected group scheme $A^{*o}_o$ over the residue
field $k'\hookrightarrow\overline k$ of $S'$, which is a quasi-Galois
extension of $k$ with Galois group $G_o$.  The quotient group $G$ of
$\pi'$ acts on $A^{*o}_o$ compatibly with its action on $k'$ through
$G_o$.  Moreover, $\pi'$ acts on $\overline k$ through its quotient
$\pi_o$, and hence acts on the situation
$(A^{*o}_o/k',\overline k)$.  Thus, by transport of structure, $\pi'$
acts on
\[
 U^{\mathrm{ef}}=T_\ell(A^{*o}_o)(\overline k),
\]
and this is the desired action of $\pi'$.  In particular, the subgroup
$\Gamma$ of $G$, which acts trivially on $k'$, acts on $A^{*o}_o$ by
$k'$-automorphisms, and the corresponding representation of
$\Gamma=I/I'$ on $T_\ell(A^{*o}_o)(\overline k)$ is the representation
deduced from the natural representation of the inertia group $I$ on
$U^{\mathrm{ef}}$.

On the other hand, the essentially toric part $U^{\mathrm{et}}$ of $U$
corresponds simply to the submodule
\begin{equation}\tag{\SGARefTarget{sga7.sga7i.eix.n4-2-7}{4.2.7}}
 U^{\mathrm{et}}\simeq T_\ell(T^*_o)(\overline k)
 \subset U^{\mathrm{ef}}\simeq T_\ell(A^{*o}_o)(\overline k),
\end{equation}
where $T^*_o$ is the largest subtorus of $A^{*o}_o$.  Likewise, put
\begin{equation}\tag{\SGARefTarget{sga7.sga7i.eix.n4-2-8}{4.2.8}}
 B^*_o=A^{*o}_o/T^*_o
\end{equation}
for the abelian part of $A^{*o}_o$.  The action of $G$ on $A^{*o}_o$
induces separately actions of $G$ on $T^*_o$ and on $B^*_o$, and the
action of $\pi'$ on $U^{\mathrm{et}}$ and on
$U^{\mathrm{ef}}/U^{\mathrm{et}}$ is obtained from these actions by
means of the isomorphism (\SGARefLink{sga7.sga7i.eix.n4-2-7}{4.2.7}) and the corresponding isomorphism
\begin{equation}\tag{\SGARefTarget{sga7.sga7i.eix.n4-2-9}{4.2.9}}
 U^{\mathrm{ef}}/U^{\mathrm{et}}
 \simeq T_\ell(B^*_o)(\overline k).
\end{equation}
It follows from this and the isogeny (\SGARefLink{sga7.sga7i.eix.n4-1-2}{4.1.2}) that the character
$g\longmapsto\operatorname{Tr}g_U$ of the action of $\pi'$ (or of
$\pi$) on $U$ is given by
\begin{equation}\tag{\SGARefTarget{sga7.sga7i.eix.n4-2-10}{4.2.10}}
 \operatorname{Tr}g_{U_\ell}
 =\operatorname{Tr}g_{B^*_o,\ell}
 +\operatorname{Tr}g_{T^*_o,\ell}
 +\chi_\ell(g)\operatorname{Tr}g^{-1}_{T^*_o,\ell}.
\end{equation}
Here $U_\ell$ has been written for $U$ on the left to recall the chosen
prime $\ell\ne p$; $g_{B^*_o,\ell}$ and $g_{T^*_o,\ell}$ denote,
respectively, the actions of $g$ on the Tate modules
$T_\ell(B^*_o)(\overline k)$ and $T_\ell(T^*_o)(\overline k)$ induced by
its actions on $B^*_o$ and $T^*_o$ just described; and
\begin{equation}\tag{\SGARefTarget{sga7.sga7i.eix.n4-2-11}{4.2.11}}
 \chi_\ell:\pi'\longrightarrow\pi_o\longrightarrow\mathbb Z_\ell^*
\end{equation}
is the character of $\pi'$ (or of $\pi_o$) induced by the action on
$\mathbb Z_\ell(1)=T_\ell(\overline k^*)$---a character that is trivial
on $\Gamma$.

We are now in a position to prove a special case of a conjecture of
Serre--Tate [29, Appendix].
\footnote{This is essentially the case $i=1$ of those conjectures for
the $H^i$.}

% END INCLUDED SGA7I_English_current_buildable_source_20260823/components/139_expose_IX_sections_4_1_and_4_2.tex

% BEGIN INCLUDED SGA7I_English_current_buildable_source_20260823/components/140_expose_IX_theorem_4_3_and_corollary_4_4.tex
% Expose IX, Theorem 4.3 and Corollary 4.4.
% French transcription lines 1098--1171; authority physical pp. 371--375
% (source folios 359--363); hard stop before section 4.5 on p. 375.
% All five pages were checked directly in four overlapping 1100-dpi bands.
% The formulas, Frobenius powers, inverses, Tate twists, primes, and bars
% were unambiguous at that resolution; no 5000--9000-dpi escalation was
% warranted for this diagram-free block.
% In applying Lemma 4.3.2, the printed shorthand replacing k by a barred
% k-prime is rendered invariantly as an algebraic closure of k'.

\medskip
\noindent
\underline{Theorem} \SGARefTarget{sga7.sga7i.eix.n4-3}{4.3}. \underline{Let $S$ be a henselian trait with
residue field $k$ and function field $K$, and let $A_K$ be an abelian
scheme over $K$.  Consider the representation of}
$\pi=\operatorname{Gal}(\overline K/K)$ \underline{on the Tate module}
$U_\ell=T_\ell(A(\overline K))$, \underline{where $\ell$ is a prime
distinct from the residual characteristic.}

\smallskip
\noindent
a) \underline{There exists an open subgroup $I'$ of the inertia group
that acts unipotently on $U_\ell$, so that the character of the
representation of $\pi$ on $U_\ell$ is in fact a function on $\pi/I'$.}
\underline{Moreover, on $\Gamma=I/I'$ this function takes integral
values and is independent of $\ell$.  More generally, for every $g\in I$,
the coefficients of the characteristic polynomial of $g_{U_\ell}$ are
rational integers independent of $\ell$.}

\smallskip
\noindent
b) \underline{Suppose that $k$ is a finite field with $q$ elements, so
that}
\[
 \pi_o=\pi/I\simeq\operatorname{Gal}(\overline k/k)
\]
\underline{is identified with the profinite completion $\widehat{\mathbb
Z}$ of $\mathbb Z$, the generator $1$ of $\mathbb Z$ corresponding to
the Frobenius automorphism}
\[
 f_q:x\longmapsto x^q
\]
\underline{of $\overline k/k$.  Then, for every $g\in\pi$ whose image in
$\pi_o$ lies in $\mathbb Z$, the trace (and, more generally, every
coefficient of the characteristic polynomial) of $g_{U_\ell}$ is a
rational number independent of $\ell$.  It is even an integer if the
image of $g$ in $\pi_o$ is of the form $f_q^{,n}$, with $n\geq0$.}

Let us prove, for example, the assertions concerning traces.  The
assertions concerning the coefficients of the characteristic polynomial
could be proved in the same way.  One may also note that they follow
formally from the assertions about traces, by expressing these
coefficients in terms of the traces of the iterates (although those
formulas introduce denominators and therefore do not give the asserted
integrality), together with a lemma of Dwork--Fatou, for which a pleasant
proof is given by Kleiman in [14, 2.8].

\smallskip
\noindent
a) We use formula (\SGARefLink{sga7.sga7i.eix.n4-2-10}{4.2.10}), which here becomes
\begin{equation}\tag{\SGARefTarget{sga7.sga7i.eix.n4-3-0}{4.3.0}}
 \operatorname{Tr}g_{U_\ell}
 =\operatorname{Tr}g_{B^*_o,\ell}
 +\operatorname{Tr}g_{T^*_o,\ell}
 +\operatorname{Tr}g^{-1}_{T^*_o,\ell}.
\end{equation}
Here $g_{B^*_o}$ and $g_{T^*_o}$ are now $k'$-automorphisms of $B^*_o$
and $T^*_o$, respectively.  It therefore suffices to show that the two
corresponding traces are rational integers independent of $\ell$, which
is a general ``geometric'' assertion about algebraic groups.  For
$g_{B^*_o}$ this is a well-known result of Weil in the theory of abelian
varieties [5] [15].  For $g_{T^*_o}$ it is practically immediate: one
describes the representation of $\Gamma$ on $T=T^*_o$ by means of the
corresponding representation on the dual group $M'\simeq\mathbb Z^\mu$.
The required trace is that of the contragredient of the latter
representation, by the canonical isomorphism
\[
 T_\ell(T)\simeq\check M_\ell(1)
 \mathrel{\mathop{=}^{\mathrm{def}}}
 \check M'\otimes_{\mathbb Z}\mathbb Z_\ell(1).
\]

\smallskip
\noindent
b) If the image of $g$ in $\pi_o$ is $f_q^{,n}$, formula (\SGARefLink{sga7.sga7i.eix.n4-2-10}{4.2.10}) now
becomes
\begin{equation}\tag{\SGARefTarget{sga7.sga7i.eix.n4-3-1}{4.3.1}}
 \operatorname{Tr}g_{U,\ell}
 =\operatorname{Tr}g_{B^*_o,\ell}
 +\operatorname{Tr}g_{T^*_o,\ell}
 +q^n\operatorname{Tr}g^{-1}_{T^*_o,\ell}.
\end{equation}
We are thus reduced to proving the following lemma, in which $k$ is to be
replaced by an algebraic closure of $k'$, and $A$ by
$A^{*o}_o\otimes_{k'}\overline k$.

\medskip
\noindent
\underline{Lemma} \SGARefTarget{sga7.sga7i.eix.n4-3-2}{4.3.2}. \underline{Let $k$ be the algebraic closure of
a finite field, let $A$ be a commutative group scheme of finite type over
$k$, and let $g$ be an automorphism of this group scheme compatible with
the automorphism $f$ of $\operatorname{Spec}(k)$ obtained by transport of
structure from an iterated Frobenius automorphism}
\[
 x\longmapsto x^q,
 \qquad q=p^r,\quad r\geq0,
\]
\underline{where $q$ is a power of the characteristic.  Thus the action
of $g$ on functions $\varphi$ satisfies}
\[
 g^*(\lambda\varphi)=f^*(\lambda)g^*(\varphi)
 =\lambda^{q^{-1}}g^*(\varphi)
 =\lambda^{p^{-r}}g^*(\varphi)
 \qquad(\lambda\in k).
\]
\underline{Let $g_\ell$ be the automorphism induced by $g$ on
$T_\ell(A(k))$ by ``transport of structure,'' where $\ell$ is a prime
different from $p$.  Then the trace of $g_\ell$ (and, more generally,
every coefficient of its characteristic polynomial) is a rational
integer independent of $\ell$, and its eigenvalues are algebraic integers
of absolute value $q^{1/2}$ or $q$.  If $A$ is an abelian scheme, their
absolute values are $q^{1/2}$; if $A$ is affine, the eigenvalues are of
the form $q\varepsilon_i$, where the $\varepsilon_i$ are roots of unity.}

Indeed, consider the ``$q$-th power'' endomorphism $f_q^A$ of the
underlying absolute scheme of $A$, defined as the identity on points and
as the $q$-th power map on the structure sheaf.  It is an endomorphism of
the group scheme $A$ compatible with the endomorphism
$f_q^{\operatorname{Spec}(k)}=f^{-1}$ of the base
$\operatorname{Spec}(k)$.  Consequently the composite
\begin{equation}\tag{\SGARefTarget{sga7.sga7i.eix.n4-3-3}{4.3.3}}
 F=f_q^A g
\end{equation}
is an endomorphism of the group scheme $A$ compatible with the identity
on $\operatorname{Spec}(k)$; in other words, it is a $k$-endomorphism of
the group scheme $A$.  I claim that the automorphism $\Gamma_g$ of $A(k)$
defined by $g$ through transport of structure, namely
\begin{equation}\tag{\SGARefTarget{sga7.sga7i.eix.n4-3-4}{4.3.4}}
 \Gamma_g(u)=g\circ u\circ f^{-1}
 \qquad\text{for }u\in A(k),
\end{equation}
may also be described by
\begin{equation}\tag{\SGARefTarget{sga7.sga7i.eix.n4-3-5}{4.3.5}}
 \Gamma_g(u)=F(u)
 \qquad\text{for }u\in A(k).
\end{equation}
To verify this, it is enough to observe that the underlying set maps of
the morphisms $\operatorname{Spec}(k)\longrightarrow A$ represented by
the two sides are the same.  This reduces us to the equality obtained
from (\SGARefLink{sga7.sga7i.eix.n4-3-5}{4.3.5}) by composing its left-hand member on the left with $f_q^A$
and its right-hand member on the right with $f^{-1}$, since $f_q^A$ and
$f^{-1}$ induce the identity on the underlying sets of the source
schemes.  The resulting formula follows immediately from (\SGARefLink{sga7.sga7i.eix.n4-3-3}{4.3.3}) and
(\SGARefLink{sga7.sga7i.eix.n4-3-4}{4.3.4}).

Applying (\SGARefLink{sga7.sga7i.eix.n4-3-5}{4.3.5}) to $u\in{}_{\ell^\nu}A(k)$ and passing to the inverse
limit over $\nu$, we find that $g_\ell$ may also be interpreted as the
endomorphism $T_\ell(F)$ of $T_\ell(A(k))$ associated with the
endomorphism $F$ of $A$.  It is known, on the other hand, that
$T_\ell(F)$ has a characteristic polynomial with integral coefficients
independent of $\ell$.  By devissage, this reduces to the case in which
$A$ is either an abelian scheme or a torus, where the required result was
already used in part a) of the proof.  The same argument also proves the
assertions about the eigenvalues, in view, in the abelian-scheme case, of
the classical results of Weil [15].

This proves \SGARefLink{sga7.sga7i.eix.n4-3-2}{4.3.2} and completes the proof of \SGARefLink{sga7.sga7i.eix.n4-3}{4.3}.  At the same time it
proves the following result.

\medskip
\noindent
\underline{Corollary} \SGARefTarget{sga7.sga7i.eix.n4-4}{4.4}. \underline{Suppose that the hypotheses of
\SGARefLink{sga7.sga7i.eix.n4-3}{4.3} b) hold.  Then, for every $g\in\pi'$ whose image in $\pi_o$ is of the
form $f_q^n$ with $n\geq0$, the eigenvalues of the induced automorphisms
of}
\[
 U^{\mathrm{et}},\qquad
 U^{\mathrm{ef}}/U^{\mathrm{et}},\qquad
 U/U^{\mathrm{ef}}
\]
\underline{are algebraic integers whose absolute values are,
respectively, $q^n$, $q^{n/2}$, and $1$.  For the two extreme pieces one
may say more precisely that the corresponding eigenvalues are,
respectively,}
\[
 (q^n\varepsilon_i)_{1\leq i\leq\mu}
 \qquad\text{and}\qquad
 (\varepsilon_i)_{1\leq i\leq\mu},
\]
\underline{where the $\varepsilon_i$ $(1\leq i\leq\mu)$ are roots of
unity.}

Indeed, formula (\SGARefLink{sga7.sga7i.eix.n4-1-2}{4.1.2}) gives for the last piece the eigenvalues
\[
 \frac{q^n}{q^n\varepsilon_i}=\varepsilon_i^{-1}.
\]
But since $\varepsilon_i$ is a root of unity, one has
$\varepsilon_i^{-1}=\overline{\varepsilon_i}$ (complex conjugate), and
since the family of the $\varepsilon_i$ is stable under complex
conjugation (being the family of roots of an equation with integral
coefficients), the family of the $\varepsilon_i^{-1}$ is also the family
of the $\varepsilon_i$.  This proves the last assertion of \SGARefLink{sga7.sga7i.eix.n4-4}{4.4}.

% END INCLUDED SGA7I_English_current_buildable_source_20260823/components/140_expose_IX_theorem_4_3_and_corollary_4_4.tex

% BEGIN INCLUDED SGA7I_English_current_buildable_source_20260823/components/141_expose_IX_sections_4_5_through_4_7.tex
% Expose IX, sections 4.5--4.7.5.
% French transcription lines 1173--1270; authority physical pp. 375--380
% (source folios 363--368); hard stop before section 5 on physical p. 381.
% Physical pp. 375--380 were checked directly in four overlapping 1100-dpi
% bands. Two one-character/source-logic questions in 4.7.3 were additionally
% inspected in direct 5000-dpi crops.
%
% Transparent source corrections:
% - The source transcription reads I where the printed page has Gamma for
%   the group acting on the dual of L in the proof of Corollary 4.6.
% - In 4.7.3 the print says that E is generated by v but then gives it the
%   finite-order polynomial T^(p^N)-1 belonging to u. Since u=1-nv, the
%   corresponding polynomial in v is (1-nT)^(p^N)-1; that relation is used
%   below. The 5000-dpi crop confirms that the printed letter is v.
% - The print uses n both for the torsion integer and in the displayed root
%   order p^n, but immediately continues with "u != 1, i.e. m >= 1".
%   The root order is therefore rendered p^m, keeping n for the torsion
%   integer throughout.

\medskip
\noindent
\SGARefTarget{sga7.sga7i.eix.n4-5}{4.5}. We shall give an application of \SGARefLink{sga7.sga7i.eix.n4-3}{4.3} a) to the theory of the local
conductor of abelian schemes.  Suppose that $k$ is algebraically
closed.\footnote{It would in fact suffice for $k$ to be perfect, but,
after passing to the strict henselization of $S$, one is immediately
reduced to the case considered here.}  Recall (SGA 5 X 6.2, or [20])
that, for every commutative etale $\ell$-group scheme $F_K$ over $K$
($\ell\neq p$), one defines an integer
\begin{equation}\tag{\SGARefTarget{sga7.sga7i.eix.n4-5-1}{4.5.1}}
 \alpha_s(F_K)=\alpha_s(i_*F_K),
\end{equation}
where $i:\eta\longrightarrow S$ denotes the inclusion, as follows.  One
chooses a finite Galois extension $K'$ of $K$, with group $\Gamma$,
that splits $F_K$, so that $F_K$ is represented by a finite
$\ell$-group $M$ on which $\Gamma$ acts.  On the other hand, one
associates with the extension $K'/K$ a certain projective
$\mathbb Z_\ell[\Gamma]$-module, defined up to isomorphism,
\begin{equation}\tag{\SGARefTarget{sga7.sga7i.eix.n4-5-2}{4.5.2}}
 \operatorname{Sw}=\operatorname{Sw}(K'/K,\mathbb Z_\ell),
\end{equation}
and sets
\begin{equation}\tag{\SGARefTarget{sga7.sga7i.eix.n4-5-3}{4.5.3}}
 \alpha_s(F_K)
 =\operatorname{length}
   \operatorname{Hom}_{\mathbb Z_\ell(\Gamma)}(\operatorname{Sw},M).
\end{equation}
One checks (loc.\ cit.) that this expression is independent of the
choice of the extension $K'/K$.  Moreover, the projectivity of
$\operatorname{Sw}$ as a $\mathbb Z_\ell[\Gamma]$-module implies that
$\alpha_s(F_K)$ is additive in exact sequences.

We apply this definition to the sheaf ${}_\ell A_K$, where $A_K$ is our
abelian scheme over $K$.  In this case one therefore obtains an integer
\begin{equation}\tag{\SGARefTarget{sga7.sga7i.eix.n4-5-4}{4.5.4}}
 \alpha_s(A_K;\ell)
 =\operatorname{length}
   \operatorname{Hom}_{\mathbb Z_\ell(\Gamma)}
   \bigl(\operatorname{Sw},
         U_\ell\otimes_{\mathbb Z_\ell}\mathbb F_\ell\bigr),
 \qquad
 U_\ell=T_\ell(A(\overline K)),
\end{equation}
where $\Gamma$ is chosen as above (a choice that a priori depends on
$\ell$).  The integer (\SGARefLink{sga7.sga7i.eix.n4-5-4}{4.5.4}) is called the \emph{wild conductor
exponent} of the abelian scheme $A_K$ relative to the discrete valuation
ring $V\subset K$.\footnote{One should also define the \emph{tame
conductor exponent} of $A_K$ relative to $V$ as the difference between
the lengths over $\mathbb Z/\ell\mathbb Z$ of the geometric generic and
special fibres of ${}_\ell A^o$.  It is independent of $\ell$ and equals
$2\lambda+\mu$ in the notation of (1.1.11).  The (total) conductor
exponent of $A_K$ is the sum of the tame and wild exponents; it is the
local term occurring in the Euler--Poincare formula for ${}_\ell A^o$ in
SGA 5 X 6.2 and [20].}  After passage to the henselization it is of
course defined without assuming that $V$ is henselian, and one is
immediately reduced to the case in which $S$ is strictly local.  We
intend to prove the following result.

\medskip
\noindent
\underline{Corollary} \SGARefTarget{sga7.sga7i.eix.n4-6}{4.6}. \underline{The wild conductor exponent
(\SGARefLink{sga7.sga7i.eix.n4-5-4}{4.5.4}) is independent of the choice of the prime $\ell\neq p$.}

We may suppose that $S$ is strictly local.  We compute (\SGARefLink{sga7.sga7i.eix.n4-5-4}{4.5.4}), replacing
$\operatorname{Sw}$ by any other projective
$\mathbb Z_\ell(\Gamma)$-module $L$.  Observe that the filtration
(\SGARefLink{sga7.sga7i.eix.n4-1-1}{4.1.1}) of $U_\ell$ consists of $\mathbb Z_\ell$-submodules that are
manifestly direct summands.  It therefore induces a corresponding
filtration of
\[
 U(\ell)=U_\ell\otimes_{\mathbb Z_\ell}\mathbb F_\ell,
\]
whose successive factors are obtained, by tensoring with
$\mathbb F_\ell$, from the factors $V_i$ $(i=1,2,3)$
\[
 U^{\mathrm{et}},\qquad
 U^{\mathrm{ef}}/U^{\mathrm{et}},\qquad
 U/U^{\mathrm{ef}}
\]
appearing in (\SGARefLink{sga7.sga7i.eix.n4-1-1}{4.1.1}).  This filtration is, of course, stable under
$\Gamma$.  Additivity thus decomposes
\[
 \operatorname{length}
 \operatorname{Hom}_{\mathbb Z_\ell(\Gamma)}(L,U(\ell))
\]
as a sum of three terms, the lengths of the groups
\begin{equation}\tag{*}
 \operatorname{Hom}_{\mathbb Z_\ell(\Gamma)}
   (L,V_i\otimes_{\mathbb Z_\ell}\mathbb F_\ell)
 \simeq
 \operatorname{Hom}_{\mathbb Z_\ell(\Gamma)}(L,V_i)
   \otimes_{\mathbb Z_\ell}\mathbb F_\ell.
\end{equation}
The isomorphism in (*) again follows from the assumed projectivity of
$L$; it is only necessary to choose $K'$ large enough for $\Gamma$ to
act on the $V_i$.  Now
$\operatorname{Hom}_{\mathbb Z_\ell(\Gamma)}(L,V_i)$ is plainly a
torsion-free $\mathbb Z_\ell$-module, so its rank is the length of the
right-hand side of (*).  If $\varphi_i$ is the character of the
representation of $\Gamma$ on $V_i$, and $\psi$ the character of the
representation of $\Gamma$ on $\check L$, we obtain
\begin{equation}\tag{\SGARefTarget{sga7.sga7i.eix.n4-6-1}{4.6.1}}
 \operatorname{length}
 \operatorname{Hom}_{\mathbb Z_\ell(\Gamma)}(L,U(\ell))
 =
 \sum_{1\leq i\leq3}N^{-1}
 \sum_{g\in\Gamma}\psi(g)\varphi_i(g)
\end{equation}
\[
 =N^{-1}\sum_{g\in\Gamma}\psi(g)\varphi(g),
 \qquad N=\operatorname{card}\Gamma,
\]
where $\varphi(g)=\sum_i\varphi_i(g)$ is also the function on $\Gamma$
obtained, by passage to the quotient, from the character of the
monodromy representation of the inertia group $I$ on $U$.  By \SGARefLink{sga7.sga7i.eix.n4-3}{4.3} a),
this character is an integral-valued rational function independent of
$\ell$.  Taking $L=\operatorname{Sw}$, the same is true of the character
$\psi(g)$ of $L$, which is
\[
 a(g)-N\delta(g)+1,
\]
where $\delta$ is the Dirac function on $\Gamma$ and $a$ is the Artin
conductor of $K'/K$ [26].  This proves \SGARefLink{sga7.sga7i.eix.n4-6}{4.6} and at the same time gives the
following expression, independent of $\ell$, for the wild conductor
exponent:
\begin{equation}\tag{\SGARefTarget{sga7.sga7i.eix.n4-6-2}{4.6.2}}
 \alpha_s(A_K)
 =N^{-1}\sum_{g\in\Gamma}
   \bigl(a(g)-N\delta(g)+1\bigr)\varphi(g).
\end{equation}
Here $\Gamma$ is the Galois group of a finite Galois extension $K'$ of
$K$ that splits $U^{\mathrm{ef}}$ (for example the smallest such
extension; cf.\ \SGARefLink{sga7.sga7i.eix.n4-1}{4.1}), $a(g)$ is the Artin conductor on $\Gamma$, and
$\varphi(g)$ is the trace function on $\Gamma$ described in (\SGARefLink{sga7.sga7i.eix.n4-3-0}{4.3.0}).
In the semistable-reduction case one may take $K'=K$, hence
$\Gamma=\{e\}$, and obtains
\begin{equation}\tag{\SGARefTarget{sga7.sga7i.eix.n4-6-3}{4.6.3}}
 \alpha_s(A_K)=0
 \qquad\text{(semistable-reduction case).}
\end{equation}

To conclude, let us state explicitly a criterion for semistable
reduction that follows from the general considerations of \SGARefLink{sga7.sga7i.eix.n4-2}{4.2}.

\medskip
\noindent
\underline{Proposition} \SGARefTarget{sga7.sga7i.eix.n4-7}{4.7}. \underline{(M.\ Raynaud's criterion for
semistable reduction.)  Let $S$ be a trait, let $A_K$ be an abelian
scheme over its function field $K$, and let $n>1$ be an integer prime
to the residual characteristic $p$.  Suppose that ${}_nA_K$ is
unramified relative to $S$, that is, that the inertia group acts
trivially on ${}_nA(\overline K)$.  Then either $A_K$ has semistable
reduction over $S$, or $n=2$; in the latter case, if $S$ is strictly
local, $A_K$ acquires semistable reduction over a finite Galois extension
$K'$ of $K$ whose Galois group is isomorphic to
$(\mathbb Z/2\mathbb Z)^r$.}\footnote{As follows from an example
communicated to me by J.\ Tate, \SGARefLink{sga7.sga7i.eix.n4-7}{4.7} is no longer valid when $n$ equals
the residual characteristic $p$, even if $K$ has characteristic zero
and $A_K$ is an elliptic curve with complex multiplication.}

We may plainly suppose that $S$ is strictly local.  With the notation
of \SGARefLink{sga7.sga7i.eix.n4-1}{4.1} and \SGARefLink{sga7.sga7i.eix.n4-2}{4.2}, take $K'$ as small as possible; in other words, suppose
that $\Gamma$ acts faithfully on $A^{*o}_o$.  By hypothesis $\Gamma$
acts trivially on ${}_nA(\overline K)$, and hence a fortiori on
\[
 {}_nA(\overline K)^{\mathrm{ef}}
 \simeq{}_nA^*_o(k)\supset{}_nA^{*o}_o(k).
\]
The conclusion follows from the next two lemmas.

\medskip
\noindent
\underline{Lemma} \SGARefTarget{sga7.sga7i.eix.n4-7-1}{4.7.1}. \underline{(J.-P.\ Serre.)  Let $G$ be a
commutative group scheme over a field $k$, which is an extension of an
abelian scheme by a torus, and let $n\geq2$ be an integer.  Every
finite-order automorphism of $G$ that induces the identity on ${}_nG$ is
the identity if $n\neq2$, and is either the identity or has order two if
$n=2$.}

\medskip
\noindent
\underline{Lemma} \SGARefTarget{sga7.sga7i.eix.n4-7-2}{4.7.2}. \underline{A finite group in which every
element distinct from $1$ has order $2$ is isomorphic to
$(\mathbb Z/2\mathbb Z)^r$ for a suitable $r\geq0$.}

\smallskip
\noindent
\SGARefTarget{sga7.sga7i.eix.n4-7-3}{4.7.3}. Let us recall the well-known proof of \SGARefLink{sga7.sga7i.eix.n4-7-1}{4.7.1}.  Since
$n\,\operatorname{id}_G$ is an epimorphism, the hypothesis that $u$
induces the identity on ${}_nG$, or equivalently that $1-u$ induces zero
on ${}_nG$, is expressed by the condition
\[
 1-u=nv
\]
in the endomorphism ring of $G$.  To prove the asserted conclusion, we
reduce, by decomposing the finite cyclic group generated by $u$ into its
$p$-primary components for the various primes $p$ dividing the order of
$u$, to the case in which $u$ has order $p^N$.

Let $E$ be the subring of $\operatorname{End}(G)$ generated by $v$.
Since $u=1-nv$ and $u^{p^N}=1$, it is a quotient of
\[
 \mathbb Z[T]\big/\bigl((1-nT)^{p^N}-1\bigr);
\]
it is therefore of finite type over $\mathbb Z$.  It is moreover
torsion-free (a formal consequence of the divisibility of $G$), and thus
embeds in
\[
 E_{\mathbb Q}=E\otimes_{\mathbb Z}\mathbb Q.
\]
The latter is a quotient, after the change of variable $u=1-nT$, of
$\mathbb Q[u]/(u^{p^N}-1)$, which is a direct product of finite field
extensions of $\mathbb Q$; the same is therefore true of $E_{\mathbb Q}$.
Considering the images of $E$, $u$, and $v$ in the component fields
$K_i$ of $E_{\mathbb Q}$, it remains to prove the following assertion.
Let $K=\mathbb Q(u)$ be generated by a primitive $p^m$-th root of unity
$u$, and let $n\geq2$ be an integer such that
\[
 u-1\equiv0\pmod n
\]
in the ring of integers of $K$.  Then $u=1$, or $u=-1$ and $n=2$.

This follows at once from the familiar fact that, if $u\neq1$, that is,
$m\geq1$, the ideal $(u-1)$ is prime and $(p)$ is a power of it, with
exponent greater than $1$ except when $p=2$ and $m=1$, that is, when
$u=-1$.  Since $(u-1)$ is prime and divides $p$, an integer $n\geq2$ can
divide $u-1$ only if $n=p$; and $p$ can divide $u-1$ only if
$(u-1)=(p)$, that is, if $u=-1$.  Hence $n=2$.

\smallskip
\noindent
\SGARefTarget{sga7.sga7i.eix.n4-7-4}{4.7.4}. Let us prove \SGARefLink{sga7.sga7i.eix.n4-7-2}{4.7.2}.  It is enough to prove that $\Gamma$ is
commutative.  If $x,y\in\Gamma$, then $(xy)(xy)=1$; since $x^2=1$ and
$y^2=1$, this may also be written
\[
 xy=(xy)^{-1}=y^{-1}x^{-1}=yx,
\]
as required.

\medskip
\noindent
\underline{Remark} \SGARefTarget{sga7.sga7i.eix.n4-7-5}{4.7.5}. The proof given in \SGARefLink{sga7.sga7i.eix.n4-7-3}{4.7.3}, together with
\SGARefLink{sga7.sga7i.eix.n4-7-2}{4.7.2}, shows that if $G$ is a divisible commutative group in a topos,
$n\geq2$ is an integer, and $\Gamma$ is a finite group acting faithfully
on $G$ while inducing the identity on ${}_nG$, then $\Gamma$ is the
trivial group if $n\geq3$, and is isomorphic to a group
$(\mathbb Z/2\mathbb Z)^r$ if $n=2$.

% END INCLUDED SGA7I_English_current_buildable_source_20260823/components/141_expose_IX_sections_4_5_through_4_7.tex

% BEGIN INCLUDED SGA7I_English_current_buildable_source_20260823/components/142_expose_IX_section_5_1.tex
% Expose IX, section 5.1.
% French transcription lines 1276--1348; authority physical pp. 381--383
% (source folios 369--371); hard stop before Theorem 5.2.
% Physical pp. 381--383 were checked directly in four overlapping
% 1100-dpi bands per page.  The print is clear throughout this unit, so
% no higher-resolution ambiguity crop was needed.
%
% Transparent transcription correction:
% - Formula (5.1.7) in the French transcription has Lambda^o, whereas
%   the printed authority and the surrounding definition both have A^o.

\medskip
\noindent
\underline{5. The orthogonality theorem in the case $\ell=p$
(semistable case).  Duality of the abelian schemes $B_o$ and $B'_o$.
Characterization of the fixed part.  Criteria for good reduction,
essentially good reduction, and semistable reduction.}

\medskip
\noindent
\SGARefTarget{sga7.sga7i.eix.n5-1}{5.1}. Let us place ourselves in the situation of \SGARefLink{sga7.sga7i.eix.n2-2-3}{2.2.3}, where $S$ is
henselian and $A^o$ is $\ell$-divisible, so that the \underline{fixed
part}
\[
 T_\ell(A^o)^f\lhook\joinrel\longrightarrow T_\ell(A)
\]
is defined as the pro-$\ell$ Barsotti--Tate group formed by the
$({}_{\ell^\nu}A^o)^f$.  We shall now give a definition of the toric
part (\SGARefLink{sga7.sga7i.eix.n2-3-2}{2.3.2}) of $T_\ell(A^o)$ which is independent of the hypothesis
$\ell\neq p$ made in \SGARefLink{sga7.sga7i.eix.n2-3}{2.3}, but which instead assumes that $S$ is
\underline{complete}.\footnote{This hypothesis is in fact not
essential, as is shown in \SGARefLink{sga7.sga7i.eix.n6-1}{6.1} below.}  For this purpose, for every
integer $n\geq0$ introduce the local Artinian scheme
\begin{equation}\tag{\SGARefTarget{sga7.sga7i.eix.n5-1-1}{5.1.1}}
 S_n=\operatorname{Spec}(V/\mathfrak m^{n+1}),
\end{equation}
the $n$-th infinitesimal neighbourhood of $s$ in $S$, and, for every
$S$-scheme $X$, put $X_n=X\times_S S_n$.  Since
$(A^o)_n=(A_n)^o$ is smooth over $S_n$, it is known (SGA 3 IX 3.6 bis)
that $T_o$ is the special fibre of a unique subtorus
\begin{equation}\tag{\SGARefTarget{sga7.sga7i.eix.n5-1-2}{5.1.2}}
 T_n\subset A_n^o.
\end{equation}
By the uniqueness of the $T_n$, they are obtained from one another by
the base changes $S_{n'}\longrightarrow S_n$.  In other words, they
define a \underline{formal subtorus} $\widehat T$ of the group
$\widehat A^o$ (in the category of formal schemes over $S$), the formal
completion of $A^o$ along the special fibre:
\begin{equation}\tag{\SGARefTarget{sga7.sga7i.eix.n5-1-3}{5.1.3}}
 \widehat T\subset\widehat A^o.
\end{equation}

Since the functor $X\longmapsto\widehat X$ (formal completion along the
special fibre) induces an equivalence between the category of finite
schemes over $S$ and that of finite formal schemes over $S$ (EGA III
4.8), we shall simply identify a finite scheme over $S$ with the
corresponding formal scheme.  With this identification, the definition
of the \emph{fixed parts} immediately gives, for every $\nu\geq0$,
\begin{equation}\tag{\SGARefTarget{sga7.sga7i.eix.n5-1-4}{5.1.4}}
 ({}_{\ell^\nu}A^o)^f={}_{\ell^\nu}(\widehat A^o),
\end{equation}
which may also be written as an identity of projective systems:
\begin{equation}\tag{\SGARefTarget{sga7.sga7i.eix.n5-1-5}{5.1.5}}
 T_\ell(A^o)^f=T_\ell(\widehat A^o).
\end{equation}
We therefore obtain a Barsotti--Tate pro-subgroup of the fixed part of
$T_\ell(A^o)$ by taking
$T_\ell(\widehat T)\subset T_\ell(\widehat A^o)$.  We accordingly
define the \underline{toric part} of $T_\ell(A^o)$ by
\begin{equation}\tag{\SGARefTarget{sga7.sga7i.eix.n5-1-6}{5.1.6}}
 T_\ell(A^o)^t
 \stackrel{\mathrm{def}}{=}
 T_\ell(\widehat T)
 \subset T_\ell(\widehat A^o)
 =T_\ell(A^o)^f.
\end{equation}
Thus, by definition,
\begin{equation}\tag{\SGARefTarget{sga7.sga7i.eix.n5-1-7}{5.1.7}}
 T_\ell(A^o)^t=({}_{\ell^\nu}\widehat T)_{\nu\geq0},
\end{equation}
the finite flat $S$-group ${}_{\ell^\nu}\widehat T$ itself being
described as the inductive limit of the finite flat $S_n$-groups
\begin{equation}\tag{\SGARefTarget{sga7.sga7i.eix.n5-1-8}{5.1.8}}
 {}_{\ell^\nu}(T_n),\qquad n\ \text{variable};
\end{equation}
that is, its affine algebra is described as the projective limit of the
affine algebras of the groups (\SGARefLink{sga7.sga7i.eix.n5-1-8}{5.1.8}).  Of course,
\begin{equation}\tag{\SGARefTarget{sga7.sga7i.eix.n5-1-9}{5.1.9}}
 \bigl(T_\ell(A^o)^t\bigr)_n=T_\ell(T_n).
\end{equation}
For $n=0$ this is precisely formula (\SGARefLink{sga7.sga7i.eix.n2-3-3}{2.3.3}), by which we characterized
the toric part (as a profinite etale group) in the case $\ell\neq p$.
This shows that definition (\SGARefLink{sga7.sga7i.eix.n5-1-6}{5.1.6}) is compatible with that given in
\SGARefLink{sga7.sga7i.eix.n2-3}{2.3} when $\ell\neq p$.

Passing to generic fibres, the fixed and toric parts of $T_\ell(A^o)$
define fixed and toric parts in $T_\ell(A_K)$:
\begin{equation}\tag{\SGARefTarget{sga7.sga7i.eix.n5-1-10}{5.1.10}}
 T_\ell(A_K)^t\subset T_\ell(A_K)^f\subset T_\ell(A_K).
\end{equation}
These are inclusions of pro-$\ell$ Barsotti--Tate groups.  Consider the
analogous filtration corresponding to the dual abelian scheme $A'_K$:
\begin{equation}\tag{\SGARefTarget{sga7.sga7i.eix.n5-1-11}{5.1.11}}
 T_\ell(A'_K)^t\subset T_\ell(A'_K)^f\subset T_\ell(A'_K).
\end{equation}
We then have the following result, which generalizes \SGARefLink{sga7.sga7i.eix.n2-4}{2.4} (at least in
the case where $S$ is complete, to which one could easily have reduced
in \SGARefLink{sga7.sga7i.eix.n2-4}{2.4}).


% END INCLUDED SGA7I_English_current_buildable_source_20260823/components/142_expose_IX_section_5_1.tex

% BEGIN INCLUDED SGA7I_English_current_buildable_source_20260823/components/143_expose_IX_theorem_5_2.tex
% Expose IX, Theorem 5.2 and proof through 5.2.6.
% French transcription lines 1350--1448; authority physical pp. 383--387
% (source folios 371--375); hard stop before section 5.3.
% Physical p. 383 was checked with the preceding unit, and pp. 384--387
% were checked directly in four overlapping 1100-dpi bands per page.
% The formulas and typography are clear throughout; no higher-resolution
% ambiguity crop was needed.
%
% Transparent transcription corrections:
% - The theorem statement twice corrupts the dual abelian scheme A'_K,
%   once as A_{K'} and once as Lambda'_K.  The printed authority has A'_K.

\medskip
\noindent
\underline{Theorem} \SGARefTarget{sga7.sga7i.eix.n5-2}{5.2}. \underline{Suppose that $S$ is a complete
trait, and let $A_K$ be an abelian scheme over $K$, $A'_K$ the dual
abelian scheme, and $\ell$ a prime such that the connected Neron model
$A^o$ is $\ell$-divisible (cf.\ \SGARefLink{sga7.sga7i.eix.n2-2-1}{2.2.1}), so that one then has the
filtrations (\SGARefLink{sga7.sga7i.eix.n5-1-10}{5.1.10}) and (\SGARefLink{sga7.sga7i.eix.n5-1-11}{5.1.11}) of the pro-groups
$T_\ell(A_K)$ and $T_\ell(A'_K)$.  Then the toric parts are
recovered from the fixed parts by the formula}
\begin{equation}\tag{\SGARefTarget{sga7.sga7i.eix.n5-2-1}{5.2.1}}
 T_\ell(A_K)^t
 =
 T_\ell(A_K)^f\cap\bigl(T_\ell(A'_K)^f\bigr)^\perp,
\end{equation}
\underline{and by the analogous formula obtained by interchanging the
roles of $A_K$ and $A'_K$, where $\perp$ denotes the orthogonal
complement with respect to the canonical pairing (\SGARefLink{sga7.sga7i.eix.n1-0-2}{1.0.2}).}

\underline{Finally, when $A_K$ has semistable reduction over $S$
(which follows from the other hypotheses in the case $\ell=p$), one
even has the formula}
\begin{equation}\tag{\SGARefTarget{sga7.sga7i.eix.n5-2-2}{5.2.2}}
 T_\ell(A_K)^t=\bigl(T_\ell(A'_K)^f\bigr)^\perp.
\end{equation}

\smallskip
\noindent
\SGARefTarget{sga7.sga7i.eix.n5-2-3}{5.2.3}. The meaning of formula (\SGARefLink{sga7.sga7i.eix.n5-2-1}{5.2.1}) is quite clear when $\ell$ is
prime to the characteristic of $K$: as in \SGARefLink{sga7.sga7i.eix.n2-4}{2.4}, it can then be
interpreted by passing to the ordinary $\mathbb Z_\ell$-modules, with
an action of
$\pi=\operatorname{Gal}(\overline K/K)$, formed by the
$\overline K$-points of the pro-groups involved.  Let us spell it out
in the general case, using the notion of the dual
\begin{equation}\tag{\SGARefTarget{sga7.sga7i.eix.n5-2-3-1}{5.2.3.1}}
 D(G)=\operatorname{Hom}\bigl(G,T_\ell(\mathbb G_m)\bigr)
\end{equation}
of a pro-$\ell$ Barsotti--Tate group $G=(G(\nu))$.  By definition,
this is the projective system of Cartier duals $D(G(\nu))$
(SGA 3 VII 3.3.1 and VIII 1), whose transition morphisms are defined in
the evident way (for example, by transposition from the inclusion
morphisms among the $G(\nu)$).  The duality theory of abelian
varieties tells us that the pairing mentioned in \SGARefLink{sga7.sga7i.eix.n5-2}{5.2} actually defines
an \underline{isomorphism}
\begin{equation}\tag{\SGARefTarget{sga7.sga7i.eix.n5-2-3-2}{5.2.3.2}}
 T_\ell(A'_K)\simeq D\bigl(T_\ell(A_K)\bigr),
\end{equation}
and hence, by composition, a morphism
\begin{equation}\tag{\SGARefTarget{sga7.sga7i.eix.n5-2-3-3}{5.2.3.3}}
 T_\ell(A_K)^f\longrightarrow D\bigl(T_\ell(A'_K)^f\bigr).
\end{equation}
Formula (\SGARefLink{sga7.sga7i.eix.n5-2-1}{5.2.1}) says that the kernel of the preceding homomorphism
(which a priori is a pro-subgroup of its first term) is precisely the
toric part $T_\ell(A_K)^t$.  Notice that this is an
\underline{isomorphism of pro-objects, and not of projective systems}:
we do not assert (and it would in general be false) that, for every
integer $\nu$, $T_\ell(A_K)^t(\nu)$ equals the kernel of
\[
 T_\ell(A_K)^f(\nu)
 \longrightarrow
 D\bigl(T_\ell(A'_K)^f\bigr)(\nu);
\]
the latter can a priori be larger.  We shall see below, however, that
it is not larger if $A_K$ has semistable reduction; by the hypotheses
made in \SGARefLink{sga7.sga7i.eix.n5-2}{5.2}, this is in particular the case when $\ell=p$.

\smallskip
\noindent
\SGARefTarget{sga7.sga7i.eix.n5-2-4}{5.2.4}. The proof of (\SGARefLink{sga7.sga7i.eix.n5-2-1}{5.2.1}) merely repeats, with practically no change,
the proof of \SGARefLink{sga7.sga7i.eix.n2-4}{2.4}.  Observe that the induced pairing
\begin{equation}\tag{\SGARefTarget{sga7.sga7i.eix.n5-2-4-1}{5.2.4.1}}
 T_\ell(A_K)^f\times T_\ell(A'_K)^f
 \longrightarrow T_\ell(\mathbb G_{mK})
\end{equation}
which gives rise to (\SGARefLink{sga7.sga7i.eix.n5-2-3-3}{5.2.3.3}) is induced by a pairing
\begin{equation}\tag{\SGARefTarget{sga7.sga7i.eix.n5-2-4-2}{5.2.4.2}}
 T_\ell(A^o)^f\times T_\ell({A'}^o)^f
 \longrightarrow T_\ell(\mathbb G_{mS})
\end{equation}
of pro-Barsotti--Tate groups over $S$, itself induced by the pairing
of pro-groups
\begin{equation}\tag{\SGARefTarget{sga7.sga7i.eix.n5-2-4-3}{5.2.4.3}}
 T_\ell(A^o)\times T_\ell({A'}^o)
 \longrightarrow T_\ell(\mathbb G_{mS})
\end{equation}
deduced from the biextension $W$ of $(A^o,{A'}^o)$ by
$\mathbb G_{mS}$.  Consequently, the pairing
\begin{equation}\tag{\SGARefTarget{sga7.sga7i.eix.n5-2-4-4}{5.2.4.4}}
 T_\ell(A_o^o)\times T_\ell({A'}_o^o)
 \longrightarrow T_\ell(\mathbb G_{mk}),
\end{equation}
obtained from (\SGARefLink{sga7.sga7i.eix.n5-2-4-2}{5.2.4.2}) by passing to the special fibres, is precisely
the pairing deduced from the biextension $W_o$ of
$(A_o^o,{A'}_o^o)$ by $\mathbb G_{mk}$.  By VIII \SGARefLink{sga7.sga7i.eviii.n4-11}{4.11} and \SGARefLink{sga7.sga7i.eviii.n1-5}{1.5},
the left annihilator of the pairing (\SGARefLink{sga7.sga7i.eix.n5-2-4-4}{5.2.4.4}) is therefore
$T_\ell(T_o)$, which is precisely the special fibre of the toric part
$T_\ell(A^o)^t$.  Returning to the definition of left annihilators
given in \SGARefLink{sga7.sga7i.eix.n5-2-3}{5.2.3}, we conclude by means of the following lemma.

\medskip
\noindent
\underline{Lemma} \SGARefTarget{sga7.sga7i.eix.n5-2-5}{5.2.5}. \underline{Let $u:G\longrightarrow H$ be a
homomorphism of pro-$\ell$ Barsotti--Tate groups over $S$.  If
$u_o:G_o\longrightarrow H_o$ is zero, then $u$ is zero.  If
$u_o:G_o\longrightarrow H_o$ is an isogeny (respectively an
isomorphism), then $u$ is an isogeny (respectively an isomorphism).}

Indeed, apply the first assertion of \SGARefLink{sga7.sga7i.eix.n5-2-5}{5.2.5} to the composite morphism
\[
 T_\ell(A^o)^t
 \longrightarrow T_\ell(A^o)^f
 \longrightarrow D\bigl(T_\ell({A'}^o)^f\bigr),
\]
where the second morphism is deduced from the pairing (\SGARefLink{sga7.sga7i.eix.n5-2-4-2}{5.2.4.2}).  The
composite is zero, so the pairing under consideration vanishes on the
toric part on the left.  By symmetry it also vanishes on the toric part
on the right, and hence descends to a pairing on the quotients, or
\emph{abelian pieces},
\[
 T_\ell(A^o)^{ab}\times T_\ell({A'}^o)^{ab}
 \longrightarrow T_\ell(\mathbb G_m).
\]
These quotients are again pro-Barsotti--Tate groups, as one checks
readily.  It remains to prove that the morphism corresponding to this
pairing,
\[
 T_\ell(A^o)^{ab}
 \longrightarrow D\bigl(T_\ell({A'}^o)^{ab}\bigr),
\]
is an isogeny, which again follows from \SGARefLink{sga7.sga7i.eix.n5-2-5}{5.2.5}.

We leave to the reader the easy proof of \SGARefLink{sga7.sga7i.eix.n5-2-5}{5.2.5}, a result valid over any
connected Noetherian base.  Moreover, below in \SGARefLink{sga7.sga7i.eix.n7-4-7}{7.4.7}, by a somewhat
more precise analysis of the biextensions $W_n$ of
$(A_n^o,{A'}_n^o)$, we shall find a quite evident direct proof of
the two special cases in which we had to use \SGARefLink{sga7.sga7i.eix.n5-2-5}{5.2.5}.

\smallskip
\noindent
\SGARefTarget{sga7.sga7i.eix.n5-2-6}{5.2.6}. Suppose now that $A_K$ has semistable reduction over $S$,
so that $\ell$ may be any prime whatsoever (equal or unequal to
$p$).  Moreover, by the Galois criterion for semistable reduction
\SGARefLink{sga7.sga7i.eix.n3-5}{3.5}, for $\ell\neq p$ one has
\begin{equation}\tag{\SGARefTarget{sga7.sga7i.eix.n5-2-7}{5.2.7}}
 \bigl(T_\ell(A'_K)^f\bigr)^\perp
 \ \bigl(=T_\ell(A_K)^t\bigr)
 \subset T_\ell(A_K)^f.
\end{equation}
To finish the proof of \SGARefLink{sga7.sga7i.eix.n5-2}{5.2}, it remains to show that this inclusion is
valid without restriction on $\ell$.

Indeed, the left-hand term is canonically isomorphic to
\[
 D\bigl(T_\ell(A'_K)/T_\ell(A'_K)^f\bigr)
\]
and therefore has \emph{rank} (or \emph{height})
$2n-2\alpha-\mu=\mu$, in the notation of (\SGARefLink{sga7.sga7i.eix.n2-1-9}{2.1.9}), since by
hypothesis $\lambda=0$.  This is also the rank of
$T_\ell(A_K)^t$.  Since the inclusion
\begin{equation}\tag{*}
 T_\ell(A_K)^t\subset T_\ell(A_K)
\end{equation}
is \underline{strict}, in the sense that it already induces a
monomorphism on every term of the projective systems under
consideration, it follows readily that the left-hand term of (*) is
maximal among the pro-Barsotti--Tate subgroups of $T_\ell(A_K)$
having the same rank.  Thus equality holds in (\SGARefLink{sga7.sga7i.eix.n5-2-7}{5.2.7}); by the
orthogonality theorem this is equivalent to the inclusion between the
two extreme terms of (\SGARefLink{sga7.sga7i.eix.n5-2-7}{5.2.7}).


% END INCLUDED SGA7I_English_current_buildable_source_20260823/components/143_expose_IX_theorem_5_2.tex

% BEGIN INCLUDED SGA7I_English_current_buildable_source_20260823/components/144_expose_IX_sections_5_3_and_5_4.tex
% Expose IX, sections 5.3--5.4.
% French transcription lines 1450--1499; authority physical pp. 388--390
% (source folios 376--378); hard stop before section 5.5.
% All three pages were checked directly in four overlapping 1100-dpi
% bands.  The print is clear throughout; no higher-resolution ambiguity
% crop was needed.
%
% Transparent source correction:
% - The printed statement of Theorem 5.4 calls A_K an abelian scheme
%   over S.  Its subscript, the surrounding setup, and the reduction
%   language require an abelian scheme over K with semistable reduction
%   over S; the English states this type-correct formulation.

\medskip
\noindent
\SGARefTarget{sga7.sga7i.eix.n5-3}{5.3}. Using the orthogonality correspondence between strict
pro-Barsotti--Tate subgroups (in the sense specified above) of the dual
pro-groups $T_\ell(A_K)$ and $T_\ell(A'_K)$, one formally deduces
from (\SGARefLink{sga7.sga7i.eix.n5-2-2}{5.2.2}) that the pairing (\SGARefLink{sga7.sga7i.eix.n1-0-2}{1.0.2}) between these pro-groups induces a
pairing
\begin{equation}\tag{\SGARefTarget{sga7.sga7i.eix.n5-3-2}{5.3.2}}
 T_\ell(A_K)^{ab}\times T_\ell(A'_K)^{ab}
 \longrightarrow T_\ell(\mathbb G_{mK})
\end{equation}
which is a \underline{perfect duality}, that is, it identifies each of
the two factors on the left with the dual of the other (in the sense of
\SGARefLink{sga7.sga7i.eix.n5-2-3}{5.2.3}).\footnote{The superscripts $ab$ denote the \emph{abelian
parts}, the quotient of the fixed part by the toric part.}  It is also
clear that (\SGARefLink{sga7.sga7i.eix.n5-3-2}{5.3.2}) is induced by the pairing
\begin{equation}\tag{\SGARefTarget{sga7.sga7i.eix.n5-3-3}{5.3.3}}
 T_\ell(A^o)^{ab}\times T_\ell(A^{\prime o})^{ab}
 \longrightarrow T_\ell(\mathbb G_{mS})
\end{equation}
obtained from the pairing (\SGARefLink{sga7.sga7i.eix.n5-2-4-2}{5.2.4.2}) by passing to the quotients.  Since
this is a perfect duality on the generic fibres, that is, since the
corresponding homomorphism
\begin{equation}\tag{\SGARefTarget{sga7.sga7i.eix.n5-3-4}{5.3.4}}
 T_\ell(A^o)^{ab}
 \longrightarrow D\bigl(T_\ell(A^{\prime o})^{ab}\bigr)
\end{equation}
is an isomorphism on the generic fibres, one naturally wishes to deduce
that it is so over the whole of $S$, or equivalently on the special
fibres.  This deduction is trivially valid if $\ell\neq p$.  It is
also valid, at least in unequal characteristics when $\ell=p$, by a
deep theorem of Tate [33, Th.\ 4].  We shall nevertheless establish it
directly below in \SGARefLink{sga7.sga7i.eix.n7-4-6}{7.4.6}, without using Tate's theorem.

To make precise the geometric meaning of the result we wish to prove,
observe that the pairing obtained from (\SGARefLink{sga7.sga7i.eix.n5-3-3}{5.3.3}) by passage to the
special fibres is precisely the pairing
\begin{equation}\tag{\SGARefTarget{sga7.sga7i.eix.n5-3-5}{5.3.5}}
 T_\ell(B_o)\times T_\ell(B'_o)
 \longrightarrow T_\ell(\mathbb G_{mk})
\end{equation}
deduced from the biextension $W_o$ of
$(A_o^o,A_o^{\prime o})$ by passage to the quotient, where
$B_o$ and $B'_o$ denote respectively the abelian parts of
$A_o^o$ and $A_o^{\prime o}$.  By VIII \SGARefLink{sga7.sga7i.eviii.n4-8}{4.8}, the biextension
$W_o$ comes from a biextension $W_o^{ab}$ of
$(B_o,B'_o)$ by $\mathbb G_{mk}$, unique up to a unique
isomorphism.  Thus the pairing (\SGARefLink{sga7.sga7i.eix.n5-3-5}{5.3.5}) may also be interpreted as the
one deduced from the divisorial correspondence $W_o^{ab}$, and the
corresponding homomorphism
\begin{equation}\tag{\SGARefTarget{sga7.sga7i.eix.n5-3-6}{5.3.6}}
 T_\ell(B_o)
 \longrightarrow D\bigl(T_\ell(B'_o)\bigr)
 \simeq T_\ell\bigl(D(B'_o)\bigr)
\end{equation}
(where, in the last term, $D$ denotes the dual abelian scheme) is
precisely the morphism obtained by applying $T_\ell$ functorially to
the homomorphism
\begin{equation}\tag{\SGARefTarget{sga7.sga7i.eix.n5-3-7}{5.3.7}}
 B_o\longrightarrow D(B'_o)
\end{equation}
defined by $W_o^{ab}$.  Consequently, by the standard dictionary
(cf.\ [5]), saying that (5.3.5) is a perfect duality is also saying that
the kernel of (\SGARefLink{sga7.sga7i.eix.n5-3-7}{5.3.7}), which is a finite group over $k$ by \SGARefLink{sga7.sga7i.eix.n1-5}{1.5}, has
zero $\ell$-primary component.  Saying that this holds for every
$\ell$ therefore means that the isogeny (\SGARefLink{sga7.sga7i.eix.n5-3-7}{5.3.7}) is an isomorphism,
that is, that the divisorial correspondence $W_o^{ab}$ is a perfect
duality between $B_o$ and $B'_o$.  The preceding considerations
show only that the kernel of (\SGARefLink{sga7.sga7i.eix.n5-3-7}{5.3.7}) is a $p$-group.  We therefore
state the more complete result.

\medskip
\noindent
\underline{Theorem} \SGARefTarget{sga7.sga7i.eix.n5-4}{5.4}. \underline{Let $S$ be a trait, and let
$A_K$ be an abelian scheme over $K$ having semistable reduction
over $S$.  Let $B_o$ and $B'_o$ denote respectively the
abelian parts of the Neron models of $A_K$ and of the dual abelian
scheme $A'_K$.  Then the canonical divisorial correspondence
$W_o^{ab}$ on $(B_o,B'_o)$, deduced by VIII \SGARefLink{sga7.sga7i.eviii.n4-8}{4.8} from the
canonical biextension $W_o$ of $(A_o^o,A_o^{\prime o})$ (the
special fibre of the biextension $W$ of \SGARefLink{sga7.sga7i.eix.n1-4}{1.4}), is a perfect duality
between $B_o$ and $B'_o$; that is, it defines an isomorphism
of each of these abelian schemes with the dual of the other.}

The proof of \SGARefLink{sga7.sga7i.eix.n5-4}{5.4} in the general case will be given in section 7, as a
consequence of a very natural and important construction due to
Raynaud.  Let us merely observe here that, in proving \SGARefLink{sga7.sga7i.eix.n5-4}{5.4}, one may
assume that $S$ is complete, by \SGARefLink{sga7.sga7i.eix.n3-3}{3.3}.

% END INCLUDED SGA7I_English_current_buildable_source_20260823/components/144_expose_IX_sections_5_3_and_5_4.tex

% BEGIN INCLUDED SGA7I_English_current_buildable_source_20260823/components/145_expose_IX_sections_5_5_and_5_6.tex
\par\medskip
\noindent
\SGARefTarget{sga7.sga7i.eix.n5-5}{5.5}. To express conveniently the extreme factors in the filtrations
(\SGARefLink{sga7.sga7i.eix.n5-1-10}{5.1.10}) and (\SGARefLink{sga7.sga7i.eix.n5-1-11}{5.1.11}), in the case of semistable reduction, let us
introduce the twisted constant group schemes over $S$
\begin{equation}\tag{\SGARefTarget{sga7.sga7i.eix.n5-5-1}{5.5.1}}
 \underline M'_o=D(T_o),\qquad \underline M_o=D(T'_o),
\end{equation}
dual to the maximal tori $T_o,T'_o$ of the special fibres of $A,A'$
(SGA 3, X 5.7). Since $S$ is complete local, these group schemes may
be regarded as the special fibres of twisted constant schemes
\begin{equation}\tag{\SGARefTarget{sga7.sga7i.eix.n5-5-2}{5.5.2}}
 \underline M',\qquad \underline M\quad\hbox{over }S.
\end{equation}
Thus, if the torus $T_o$ is split (for example, if $k$ is separably
closed), that is, if $\underline M'_o$ is the constant group defined
by the ordinary group
\begin{equation}\tag{\SGARefTarget{sga7.sga7i.eix.n5-5-3}{5.5.3}}
 M'=\operatorname{Hom}_{k\text{-}\mathrm{gr}}(T_o,\mathbb G_{mk}),
\end{equation}
then $\underline M'$ is the constant group defined by $M'$:
\begin{equation*}
 \underline M'=M'_S.
\end{equation*}
One may moreover observe that, since $T_o$ and $T'_o$ are isogenous
(\SGARefLink{sga7.sga7i.eix.n2-2-6}{2.2.6}), $T'_o$ is split whenever $T_o$ is. In this case one should
therefore also introduce
\begin{equation}\tag{\SGARefTarget{sga7.sga7i.eix.n5-5-4}{5.5.4}}
 M=\operatorname{Hom}_{k\text{-}\mathrm{gr}}(T'_o,\mathbb G_{mk}),
\end{equation}
and then
\begin{equation*}
 \underline M=M_S.
\end{equation*}
With this notation, for every integer $n$ there are evidently
isomorphisms
\begin{equation}\tag{\SGARefTarget{sga7.sga7i.eix.n5-5-5}{5.5.5}}
 \underline M'_n\simeq D(T_n),\qquad
 \underline M_n\simeq D(T'_n),
\end{equation}
where $T'_n$ is defined in terms of $A'$ as $T_n$ is in terms of $A$
(\SGARefLink{sga7.sga7i.eix.n5-1-2}{5.1.2}). Hence there follow canonical isomorphisms (where
$\mathbb Z_\ell(1)=T_\ell(\mathbb G_{mS})$)
\begin{equation}\tag{\SGARefTarget{sga7.sga7i.eix.n5-5-6}{5.5.6}}
 T_\ell(T_n)\simeq
 \check{\underline M}'_n\otimes_{\mathbb Z}\mathbb Z_\ell(1),
 \qquad
 T_\ell(T'_n)\simeq
 \check{\underline M}_n\otimes_{\mathbb Z}\mathbb Z_\ell(1),
\end{equation}
and, on passing to the limit over $n$, canonical isomorphisms
\begin{equation}\tag{\SGARefTarget{sga7.sga7i.eix.n5-5-7}{5.5.7}}
 \left\{
 \begin{aligned}
  T_\ell(A^o)^t=T_\ell(\widehat T)
   &\simeq \check{\underline M}'\otimes\mathbb Z_\ell(1),\\[1ex]
  T_\ell(A^{\prime o})^t=T_\ell(\widehat T')
   &\simeq \check{\underline M}\otimes\mathbb Z_\ell(1).
 \end{aligned}
 \right.
\end{equation}
In the case of semistable reduction, duality, together with the
orthogonality theorem \SGARefLink{sga7.sga7i.eix.n5-2}{5.2} and relation (\SGARefLink{sga7.sga7i.eix.n5-2-2}{5.2.2}), then gives canonical
isomorphisms
\begin{equation}\tag{\SGARefTarget{sga7.sga7i.eix.n5-5-8}{5.5.8}}
 \left\{
 \begin{aligned}
  T_\ell(A_K)/T_\ell(A_K)^f
   &\simeq \underline M_K\otimes_{\mathbb Z}\mathbb Z_\ell,\\[1ex]
  T_\ell(A'_K)/T_\ell(A'_K)^f
   &\simeq \underline M'_K\otimes_{\mathbb Z}\mathbb Z_\ell.
 \end{aligned}
 \right.
\end{equation}
Likewise, the same formula (\SGARefLink{sga7.sga7i.eix.n5-2-2}{5.2.2}) gives canonical isomorphisms
\begin{equation}\tag{\SGARefTarget{sga7.sga7i.eix.n5-5-9}{5.5.9}}
 \left\{
 \begin{aligned}
  T_\ell(A_K)/T_\ell(A_K)^t
   &\simeq D\bigl(T_\ell(A'_K)^f\bigr),\\[1ex]
  T_\ell(A'_K)/T_\ell(A'_K)^t
   &\simeq D\bigl(T_\ell(A_K)^f\bigr),
 \end{aligned}
 \right.
\end{equation}
where one should recall that the groups $T_\ell(A_K)^f$ and
$T_\ell(A'_K)^f$ occurring on the right are the restrictions to
$\eta$ of the pro-$\ell$ Barsotti--Tate groups $T_\ell(A^o)^f$ and
$T_\ell(A^{\prime o})^f$, respectively, defined over $S$. Here $D$
denotes duality in the sense of pro-Barsotti--Tate groups, namely
$\underline{\operatorname{Hom}}(-,\mathbb Z_\ell(1))$. Thus we have:

\medskip
\noindent
\underline{Proposition} \SGARefTarget{sga7.sga7i.eix.n5-6}{5.6}. \underline{Suppose that $S$ is complete
and that $A_K$ has semistable reduction over $S$. Then, for every
prime number $\ell$ (which may equal the residue characteristic of
$S$), the pro-$\ell$ Barsotti--Tate group $T_\ell(A_K)^t$ is the
generic fibre of a pro-Barsotti--Tate group over $S$ ``of toric
type'', given by (\SGARefLink{sga7.sga7i.eix.n5-5-7}{5.5.7}), whereas
$T_\ell(A_K)/T_\ell(A_K)^f$ is the generic fibre of a
pro-Barsotti--Tate group over $S$ of pro-\'{e}tale type, given by
(\SGARefLink{sga7.sga7i.eix.n5-5-8}{5.5.8}). Finally, the groups $T_\ell(A_K)^f$ and
$T_\ell(A_K)/T_\ell(A_K)^t$ are the generic fibres of
pro-Barsotti--Tate groups defined over $S$, equal respectively to
$T_\ell(A^o)^f$ and to
$D(T_\ell(A^{\prime o})^f)=
\underline{\operatorname{Hom}}(T_\ell(A^{\prime o})^f,
\mathbb Z_\ell(1))$.}

\medskip
\noindent
\SGARefTarget{sga7.sga7i.eix.n5-6-1}{5.6.1}. Observe that, for an \'{e}tale covering $S'$ of $S$, splitting
$T_\ell(A_K)^t$ is equivalent to splitting
$T_\ell(A_K)/T_\ell(A_K)^f$: these conditions say respectively that
$S'$ splits the twisted constant groups $\underline M$ and
$\underline M'$, and these are isogenous because $T_o$ and $T'_o$
are. At the same time one sees that there is a smallest \'{e}tale
extension $S'$ of $S$ which does the job.

We also record the following easy consequence of the definitions in
\SGARefLink{sga7.sga7i.eix.n2-2-3}{2.2.3}.

% END INCLUDED SGA7I_English_current_buildable_source_20260823/components/145_expose_IX_sections_5_5_and_5_6.tex

% BEGIN INCLUDED SGA7I_English_current_buildable_source_20260823/components/146_expose_IX_proposition_5_7_through_corollary_5_10.tex
\par\medskip
\noindent
\underline{Proposition} \SGARefTarget{sga7.sga7i.eix.n5-7}{5.7}. \underline{Suppose that $A^o$ is
$\ell$-divisible and that $S$ is henselian. Let $F_K$ be a
constructible $\ell$-adic sheaf over $K$ which is ``unramified'', that
is, which extends to a twisted constant $\ell$-adic sheaf $F$ over
$S$, and let $u:F_K\longrightarrow T_\ell(A_K)$ be a homomorphism of
pro-groups. Then $u$ factors through a homomorphism from $F$ to
$T_\ell(A_K)^f$.}

Indeed, write
\begin{equation}\tag{\SGARefTarget{sga7.sga7i.eix.n5-7-1}{5.7.1}}
 u=(u(\nu))_{\nu\geqslant0},\qquad
 u(\nu):F_K(\nu)\longrightarrow T_\ell(A_K)(\nu)
 ={}_{\ell^\nu}A_K.
\end{equation}
By the universal property of $A$, since $F(\nu)$ is \'{e}tale and
hence smooth over $S$, the morphism $u(\nu)$ comes from a
homomorphism
\[
 v(\nu):F(\nu)\longrightarrow{}_{\ell^\nu}A.
\]
By definition of the fixed part, this homomorphism factors as
\[
 v(\nu):F(\nu)\longrightarrow({}_{\ell^\nu}A)^f.
\]
Since the image of the induced morphism on special fibres lies in the
infinitely $\ell$-divisible part of $A_o$, it follows that $v(\nu)$
in fact factors through $({}_{\ell^\nu}A^o)^f$. Passing to the limit
over $\nu$ gives a homomorphism
\[
 v:F\longrightarrow T_\ell(A^o)^f,
\]
and the induced morphism on generic fibres plainly gives the required
factorization.

\medskip
\noindent
\underline{Remarks} \SGARefTarget{sga7.sga7i.eix.n5-7-2}{5.7.2}. When $\ell\neq p$, property \SGARefLink{sga7.sga7i.eix.n5-7}{5.7} uniquely
characterizes the sub-pro-group $T_\ell(A_K)^f$ of $T_\ell(A_K)$ as
its largest unramified $\ell$-adic subsheaf. When $\ell=p$, the
statement of \SGARefLink{sga7.sga7i.eix.n5-7}{5.7} ceases to be very satisfactory. We shall give a
considerably stronger property, which can serve as an intrinsic
characterization of the fixed part of $T_p(A_K)$ (in terms of the
pro-Barsotti--Tate group structure of $T_p(A_K)$ alone), at least in
unequal characteristics.

\medskip
\noindent
\underline{Theorem} \SGARefTarget{sga7.sga7i.eix.n5-8}{5.8} (M. Raynaud). \underline{Let $A_K$ be an
abelian scheme with semistable reduction over the function field $K$
of a henselian trait $S$. Let $\ell$ be a prime number (the
interesting case being $\ell=p$), and consider the pro-$\ell$
Barsotti--Tate group $X_\eta=T_\ell(A_K)$ and its ``fixed part''
$X_\eta^f$, which is the generic fibre of the pro-$\ell$
Barsotti--Tate group $X^f=T_\ell(A^o)^f$. If $\ell=p$, suppose that
$\operatorname{char}K=0$ (so that Tate's theorem [33, Th.\ 4] is
available over $S$). Let $F$ be a pro-$\ell$ Barsotti--Tate group over
$S$. Then every homomorphism $F_\eta\longrightarrow X_\eta$ factors
through $X_\eta^f$.}

Passing to the ind-$\ell$ groups associated with the pro-$\ell$
Barsotti--Tate groups under consideration, the preceding statement
becomes equivalent, by the cited theorem of Tate, to the following.

\medskip
\noindent
\underline{Corollary} \SGARefTarget{sga7.sga7i.eix.n5-9}{5.9}. \underline{If $F$ is an ind-$\ell$
Barsotti--Tate group, then every homomorphism
$u_\eta:F_\eta\longrightarrow A_\eta$ extends uniquely to a
homomorphism $u:F\longrightarrow A$.}

In other words, the characteristic Neron property (\SGARefLink{sga7.sga7i.eix.n1-1-2}{1.1.2}) of the
Neron model remains valid for homomorphisms from ind-group schemes of
Barsotti--Tate type, with no smoothness hypothesis on them. (Only the
case of ind-\'{e}tale groups would follow directly from the
definition.)

\medskip
\noindent
\underline{Remark} \SGARefTarget{sga7.sga7i.eix.n5-9-1}{5.9.1}. In fact, M. Raynaud proves \SGARefLink{sga7.sga7i.eix.n5-9}{5.9} without
assuming that $A_K$ has semistable reduction, and even for every Neron
group scheme $A$ over $S$ (whose generic fibre need not be an abelian
scheme). Of course, the assumption that $S$ is henselian is also
unnecessary in the formulation of \SGARefLink{sga7.sga7i.eix.n5-9}{5.9}. Raynaud's proof of this more
general case reduces it to the special case considered in \SGARefLink{sga7.sga7i.eix.n5-8}{5.8}, using
an unpublished property in the construction of Neron models. Here we
shall give only Raynaud's proof of \SGARefLink{sga7.sga7i.eix.n5-9}{5.9}.

\medskip
\noindent
\underline{Proof of \SGARefLink{sga7.sga7i.eix.n5-9}{5.9}.} It is known [33] that $F$ is an extension
of an ind-\'{e}tale Barsotti--Tate group $F'$ by a Barsotti--Tate group
$F''$ with connected special fibre, so that for every integer $n$
there is an exact sequence
\[
 0\longrightarrow F''(n)\longrightarrow F(n)
 \longrightarrow F'(n)\longrightarrow0,
\]
with $F'(n)$ \'{e}tale and $F''(n)_o$ purely infinitesimal.
The given homomorphism $u_\eta$ is defined by a sequence of
homomorphisms
\[
 u_\eta(n):F(n)_\eta\longrightarrow A_\eta,
\]
and it is enough to extend these separately to homomorphisms
$u(n):F(n)\longrightarrow A$. Lemma \SGARefLink{sga7.sga7i.eix.n5-9-2}{5.9.2} below then reduces us to
proving that the homomorphism
$u''_\eta:F''_\eta\longrightarrow A_\eta$ induced by $u_\eta$
extends to $u'':F''\longrightarrow A$. Consequently, \underline{we
may suppose that $F$ has connected special fibre}.

We may regard $u_\eta$ as a homomorphism from $F_\eta$ to the
ind-$\ell$ Barsotti--Tate group $Y_\eta$ associated with
$X_\eta=T_\ell(A_K)$, and then consider the composite
\begin{equation}\tag{*}
 v_\eta:F_\eta\longrightarrow Y_\eta
 \longrightarrow Y_\eta/Y_\eta^f.
\end{equation}
By \SGARefLink{sga7.sga7i.eix.n5-6}{5.6}, the latter group extends to an ind-\'{e}tale
Barsotti--Tate group $Z$ over $S$, namely
\[
 \underline M\otimes_{\mathbb Z}(\mathbb Q_\ell/\mathbb Z_\ell).
\]
By Tate's theorem cited in \SGARefLink{sga7.sga7i.eix.n5-8}{5.8}, the homomorphism (*) then extends to a
homomorphism
\[
 v:F\longrightarrow Z.
\]
Since $F$ has connected special fibre and $Z$ is ind-\'{e}tale, this
morphism is zero. Hence $v_\eta$ is zero as well, and $u_\eta$
therefore factors through $Y_\eta^f$, as required.

It remains to prove the following lemma.

\medskip
\noindent
\underline{Lemma} \SGARefTarget{sga7.sga7i.eix.n5-9-2}{5.9.2}. Let $S$ be a trait, let $A$ be
a Neron group scheme over $S$ (\SGARefLink{sga7.sga7i.eix.n1-1-2}{1.1.2}), let
\[
 1\longrightarrow G'\longrightarrow G\longrightarrow G''
 \longrightarrow1
\]
be an extension of group schemes over $S$, with $G''$ smooth, let
$u:G'\longrightarrow A$ be a group homomorphism, and let
$v_\eta:G_\eta\longrightarrow A_\eta$ be a group homomorphism
extending $u_\eta$. Suppose that the extension of $G''$ by $A$
obtained from the extension $G$ by pushout along $u$ is representable
\footnote{Unpublished results of M. Raynaud imply, moreover, that this
condition is always satisfied.}---as is the case, for example, if
$G\longrightarrow G''$ admits a section locally for the ``finite
locally free'' topology (cf.\ SGA 3, IV 6.3) and if $A$ is
quasi-projective over $S$. Then there exists a unique group
homomorphism $v:G\longrightarrow A$ extending both $v_\eta$ and $u$.

Indeed, the last hypothesis reduces us to the case in which $G'=A$,
$u$ is the identity, and $v_\eta:G_\eta\longrightarrow A_\eta$ is a
projection which we seek to extend to a projection
$v:G\longrightarrow A$. But $G$ is smooth, being an extension of two
smooth group schemes, and the conclusion follows from the definition
of the Neron property of $A$.

\medskip
\noindent
\underline{Corollary} \SGARefTarget{sga7.sga7i.eix.n5-10}{5.10}. \underline{Let $A_K$ be an abelian
scheme over the function field $K$ of a trait $S$, and let $\ell$ be
a prime number. When $\ell=p$, suppose that
$\operatorname{char}K=0$, so that Tate's theorem already invoked in
\SGARefLink{sga7.sga7i.eix.n5-8}{5.8} is available. Then $A_K$ has good reduction over $S$ if and only
if the pro-$\ell$ Barsotti--Tate group $T_\ell(A_K)$ ``has good
reduction over $S$'', that is, extends to a Barsotti--Tate group over
$S$.}

(This is the result conjectured in [8, 4.9].)

Necessity being clear, it remains to prove sufficiency; observe that
the case $\ell\neq p$ has already been treated in \SGARefLink{sga7.sga7i.eix.n2-2-9}{2.2.9}. We may
suppose that $S$ is henselian (\SGARefLink{sga7.sga7i.eix.n2-2-9}{2.2.9}). First suppose that $A_K$ has
semistable reduction over $S$. Applying \SGARefLink{sga7.sga7i.eix.n5-8}{5.8} with $F=X$ gives
\[
 T_\ell(A_K)=T_\ell(A_K)^f.
\]
This also says that the maximal torus of $A_o$ is zero, that is, that
$A_o^o$ is an abelian scheme; the conclusion then follows from
\SGARefLink{sga7.sga7i.eix.n2-2-9}{2.2.9} (iii${}^o$).

We now treat the general case. By \SGARefLink{sga7.sga7i.eix.n2-2-9}{2.2.9} we may plainly assume that
$S$ is strictly local. Let $K'/K$ be a finite Galois extension such
that $A_{K'}$ has semistable reduction over the normalization $S'$ of
$S$ in $K'$ (\SGARefLink{sga7.sga7i.eix.n3-6}{3.6}). If $A^*$ denotes the corresponding Neron model,
$G=\operatorname{Gal}(K'/K)$ acts on the situation $(A^*,S')$, and
hence on $A_o^{*o}$. By \SGARefLink{sga7.sga7i.eix.n4-1}{4.1} it is enough to prove that $G$ acts
trivially on the latter; this indeed implies that $A_K$ has
semistable reduction, reducing us to the case already treated.

Let $X$ be a Barsotti--Tate group over $S$ extending
$X_\eta=T_\ell(A_K)$. By Tate's uniqueness theorem already invoked,
there is a unique isomorphism
\[
 X\times_S S'\simeq T_\ell(A^{*o})
\]
which induces the identity on generic fibres. By transport of
structure, this isomorphism is compatible with the actions of $G$.
Now $G$ plainly acts trivially on the special fibre of
$X\times_S S'$, and hence on $T_\ell(A_o^{*o})$. It follows at once
that $G$ acts trivially on $A_o^{*o}$, completing the proof.

% END INCLUDED SGA7I_English_current_buildable_source_20260823/components/146_expose_IX_proposition_5_7_through_corollary_5_10.tex

% BEGIN INCLUDED SGA7I_English_current_buildable_source_20260823/components/147_expose_IX_remark_5_11_through_corollary_5_15.tex
\medskip
\noindent
\underline{Remark} \SGARefTarget{sga7.sga7i.eix.n5-11}{5.11}. Once Tate's theorem has been proved without any
restrictive hypothesis on $\operatorname{char}K$, the preceding results
\SGARefLink{sga7.sga7i.eix.n5-8}{5.8}, \SGARefLink{sga7.sga7i.eix.n5-9}{5.9}, and \SGARefLink{sga7.sga7i.eix.n5-10}{5.10}, as well as \SGARefLink{sga7.sga7i.eix.n5-13}{5.13} below, will plainly hold without
such a restriction.

We record a definition which will also be useful elsewhere.

\medskip
\noindent
\underline{Definition} \SGARefTarget{sga7.sga7i.eix.n5-12}{5.12}. Let $S$ be a henselian trait, let $\ell$ be
a prime number, and let $X_\eta$ be a pro-$\ell$ Barsotti--Tate group
over the generic point $\eta$ of $S$. We say that $X_\eta$ has
\emph{good reduction over $S$} if it extends to a Barsotti--Tate group
over $S$ (cf. \SGARefLink{sga7.sga7i.eix.n5-10}{5.10}). We say that it has \emph{virtual good reduction
over $S$} if there exists a filtration
\[
 \bigl(\operatorname{Fil}^i(X_\eta)\bigr)_{i\geq0}
\]
by strict Barsotti--Tate subgroups\footnote{That is, subobjects as
projective systems, or equivalently subobjects term by term.} such that
the Barsotti--Tate groups
\[
 \operatorname{Gr}^i(X_\eta)=
 \operatorname{Fil}^i(X_\eta)/\operatorname{Fil}^{i+1}(X_\eta)
\]
have good reduction over $S$. If the filtration can be chosen to have
length at most $n$ (that is, if at most $n$ of the
$\operatorname{Gr}^i$ are nonzero), we say that $X_\eta$ has
\emph{virtual good reduction of level $n$ over $S$}.

Finally, we say that $X_\eta$ has \emph{essential good reduction over
$S$} (respectively, \emph{essential virtual good reduction},
respectively, \emph{essential virtual good reduction of level $n$}) if
there exists a finite extension $K'/K$ such that $X_{K'}$ has good
reduction (respectively, virtual good reduction, respectively, virtual
good reduction of level $n$) over the normalization $S'$ of $S$ in
$K'$.

\medskip
\noindent
\SGARefTarget{sga7.sga7i.eix.n5-12-1}{5.12.1}. When $\ell\neq p$, the categories of pro-$\ell$
Barsotti--Tate groups over $\eta$ and over $S$ may be interpreted in
terms of $\ell$-adic representations of
\[
 \pi=\pi_1(\eta), \qquad \pi_o=\pi_1(S),
\]
respectively. The six notions just introduced can then be expressed in
terms of the natural representation of the inertia group $I$ on
$X_\eta(\overline K)$: the representation is, respectively, trivial,
unipotent, unipotent of level $n$, essentially trivial, essentially
unipotent, or essentially unipotent of level $n$.

\medskip
\noindent
\SGARefTarget{sga7.sga7i.eix.n5-12-2}{5.12.2}. Take $X_\eta=T_\ell(A_K)$, where $A_K$ is an abelian scheme
over $K$. We saw in \SGARefLink{sga7.sga7i.eix.n5-6}{5.6} that, if $A_K$ has semistable reduction over
$S$, then $X_\eta$ has virtual good reduction of level 2 (and we shall
see the converse in a moment). Consequently, by the semistable
reduction theorem \SGARefLink{sga7.sga7i.eix.n3-6}{3.6}, $X_\eta$ has essential virtual good reduction of
level 2. This statement should be regarded as the ``right'' formulation
of the geometric monodromy theorem in Expos\'e III\footnote{The source
leaves the precise internal reference blank.} for $H^1$ of a proper
smooth scheme over $K$, without excluding the case $\ell=p$ (which was
not considered there). The geometric monodromy theorem---or, for the
present, conjecture---for the higher $H^i$ given in Expos\'e III should
likewise be reformulated so as to include ``$p$-adic'' coefficients, as
was observed in \SGARefLink{sga7.sga7i.eix.n0-4}{0.4}.

We now summarize the relations between the reduction properties of
abelian schemes over $K$ and those of the associated Barsotti--Tate
groups.

\medskip
\noindent
\underline{Proposition} \SGARefTarget{sga7.sga7i.eix.n5-13}{5.13}. Let $S$ be a henselian trait, let $A_K$
be an abelian scheme over its fraction field, and let $\ell$ be a prime
number. When $\ell=p$, assume that $K$ has characteristic zero
(cf. \SGARefLink{sga7.sga7i.eix.n5-11}{5.11}).

\smallskip
\noindent
a) The abelian scheme $A_K$ has good reduction over $S$ (\SGARefLink{sga7.sga7i.eix.n2-2-9}{2.2.9}) if and
only if $T_\ell(A_K)$ has good reduction over $S$ (\SGARefLink{sga7.sga7i.eix.n5-12}{5.12}).

\smallskip
\noindent
b) The abelian scheme $A_K$ has essential good reduction over $S$---that
is, there exists a finite extension $K'/K$ such that $A_{K'}$ has good
reduction over the normalization $S'$ of $S$ in $K'$---if and only if
$T_\ell(A_K)$ has essential good reduction over $S$ (\SGARefLink{sga7.sga7i.eix.n5-12}{5.12}). In that
case $K'$ may be chosen to be a finite separable extension of $K$.

\smallskip
\noindent
c) The following conditions are equivalent:
\begin{enumerate}
\item[1)] $A_K$ has semistable reduction over $S$ (3.4);
\item[2)] $T_\ell(A_K)$ has virtual good reduction of level 2 over $S$
  (\SGARefLink{sga7.sga7i.eix.n5-12}{5.12});
\item[3)] $T_\ell(A_K)$ has virtual good reduction over $S$.
\end{enumerate}

\noindent
\underline{Proof}. Part a) is just \SGARefLink{sga7.sga7i.eix.n5-10}{5.10} and has been included for
reference.

For b), the first assertion follows immediately from a) and the
definitions. To prove the final assertion, we may suppose that
$\ell\neq p$. In view of a), it then suffices to take for $K'$ the
finite separable extension of $K$ corresponding to an open subgroup
$\pi'$ of
\[
 \pi=\operatorname{Gal}(\overline K/K)
\]
whose intersection with the inertia group $I$ is contained in the
kernel of the representation of $I$ on
$T_\ell(A_K(\overline K))$; this kernel is open by hypothesis
(cf. \SGARefLink{sga7.sga7i.eix.n5-12-1}{5.12.1}).

For c), we have already noted that 1) implies 2) (\SGARefLink{sga7.sga7i.eix.n5-12-2}{5.12.2}), and it is
immediate that 2) implies 3). It remains to prove that 3) implies 1).
For $\ell\neq p$, this is precisely the Galois criterion for
semistable reduction \SGARefLink{sga7.sga7i.eix.n3-5}{3.5}. For $\ell=p$, argue as follows. Choose a
prime $\ell'\neq p$ and apply the considerations of \SGARefLink{sga7.sga7i.eix.n4-1}{4.1} to it. We may
suppose that $S$ is strictly local (\SGARefLink{sga7.sga7i.eix.n3-4-0}{3.4.0}), and choose a finite Galois
extension $K'/K$, with group $G$, such that $A_{K'}$ has semistable
reduction. We must prove that the representation of $G$ on each of the
three quotients occurring in (\SGARefLink{sga7.sga7i.eix.n4-1-1}{4.1.1}) is trivial. Equivalently, their
direct-sum representation is trivial, or again its trace is the
``trivial'' trace, that is, is constant. We saw in \SGARefLink{sga7.sga7i.eix.n4-3}{4.3}(a) that this
trace function is independent of the choice of $\ell'\neq p$, and in
\SGARefLink{sga7.sga7i.eix.n6-4}{6.4} below we indicate a method of calculating it in terms of
$T_p(A_K)$ as well. That calculation plainly gives a constant function
when $T_p(A_K)$ has virtual good reduction, and this completes the
proof.

\medskip
\noindent
\underline{Remark} \SGARefTarget{sga7.sga7i.eix.n5-13-1}{5.13.1}. It seems plausible that, if $X_\eta$ is a
pro-$\ell$ Barsotti--Tate group over $\eta$ having virtual good
reduction over $S$ and essential virtual good reduction of level $n$,
then it has virtual good reduction of level $n$. (This is true when
$\ell\neq p$.)

The following result, an easy consequence of \SGARefLink{sga7.sga7i.eix.n5-7}{5.7} and of the
semistable reduction theorem \SGARefLink{sga7.sga7i.eix.n3-6}{3.6}, was pointed out to me by
J.-P. Serre.

\medskip
\noindent
\underline{Proposition} \SGARefTarget{sga7.sga7i.eix.n5-14}{5.14} (Criterion for essential good reduction).
Let $S$ be a trait with function field $K$, let $A_K$ be an abelian
scheme over $K$, and let $\ell$ be a prime number distinct from
$\operatorname{char}K$.

\smallskip
\noindent
a) Suppose that $\ell\neq p=\operatorname{char}k$. Then the following
conditions are equivalent:
\begin{enumerate}
\item[(i)] $A_K$ has essential good reduction over $S$; that is, there
  exists a finite separable extension $K'/K$ such that $A_{K'}$ has
  good reduction over the normalization $S'$ of $S$ in $K'$;
\item[(ii)] the action of the inertia group $I$ on
  $U_\ell=T_\ell(A(\overline K))$ is essentially trivial; that is, it
  is trivial on an open subgroup of $I$, or equivalently factors
  through a finite quotient $\Gamma$ of $I$;
\item[(iii)] the action of $I$ on
  $U_\ell\otimes_{\mathbb Z_\ell}\mathbb Q_\ell$ is semisimple.
\end{enumerate}

\noindent
b) Suppose that $\ell=p$. If $n=\dim A_K\neq0$, then the image of $I$
in the automorphism group of $U_\ell$ is not finite; more precisely,
its image in the group $\mathbb Z_\ell^\times$ of automorphisms of
$\bigwedge^{2n}U_\ell$ has finite index. A sufficient (but by no means
necessary) condition for $A_K$ to have essential good reduction is
that the action of $I$ on
$U_\ell\otimes_{\mathbb Z_\ell}\mathbb Q_\ell$ be semisimple.

\noindent
\underline{Proof}. When $\ell\neq p$, the equivalence of (i) and (ii)
follows immediately from the good-reduction criterion \SGARefLink{sga7.sga7i.eix.n2-2-9}{2.2.9}, and (ii)
implies (iii) because $\mathbb Q_\ell$ has characteristic zero. The
homomorphism
\[
 \pi=\operatorname{Gal}(\overline K/K)
 \longrightarrow \mathbb Z_\ell^\times
 \simeq\operatorname{Aut}\!\left(\bigwedge^{2n}U_\ell\right)
\]
is $\chi^n$, where $\chi$ is induced by the character
\[
 \mathbb Z_\ell^\times=\operatorname{Aut}(\mathbb Z_\ell(1))
\]
arising from the action on the $\ell^\nu$-th roots of unity. It is
known that, when $\ell=p$, the image of $I$ under this character has
finite index [33]. The fact that the action of $I$ on
$U_p\otimes\mathbb Q_p$ need not be semisimple even when $A_K$ is an
elliptic curve with good reduction is exhibited in [28, IV A 2.4].

It remains to prove that (iii) implies (i), for every $\ell$. We use
the semistable reduction theorem \SGARefLink{sga7.sga7i.eix.n3-6}{3.6}, which immediately reduces us to
the case where $A_K$ has semistable reduction over $S$: hypothesis
(iii) is preserved upon replacing $K$ by a finite separable extension
$K'$, or equivalently $I$ by an open subgroup $I'$. When $\ell\neq p$,
the action of $I$ on $U_\ell\otimes\mathbb Q_\ell$ is unipotent
(\SGARefLink{sga7.sga7i.eix.n3-5}{3.5}(v)) and semisimple, hence trivial, and the good-reduction
criterion \SGARefLink{sga7.sga7i.eix.n2-2-9}{2.2.9} gives the result. We now give an argument valid with no
restriction on $\ell$, which yields a slightly sharper statement.

\medskip
\noindent
\underline{Corollary} \SGARefTarget{sga7.sga7i.eix.n5-15}{5.15}. With the notation of \SGARefLink{sga7.sga7i.eix.n5-14}{5.14}, suppose that
$A_K$ has semistable reduction over $S$, and that the fixed part
\[
 (U_{\ell\mathbb Q_\ell})^f
 \stackrel{\mathrm{def}}{=}(U_\ell^f)_{\mathbb Q_\ell}
\]
admits an $I$-stable complement in $U_{\ell\mathbb Q_\ell}$. Then
$A_K$ has good reduction.

It is enough to prove under these hypotheses that $U_\ell=U_\ell^f$.
By \SGARefLink{sga7.sga7i.eix.n5-6}{5.6} this says precisely that the toric part $T_o$ of $A_o^o$ is
zero, so that $A_o^o$ is an abelian scheme; hence $A$ is an abelian
scheme by \SGARefLink{sga7.sga7i.eix.n2-2-9}{2.2.9}. We may plainly suppose that $k$ is separably closed.
The intersection of the asserted complement in \SGARefLink{sga7.sga7i.eix.n5-15}{5.15} with $U_\ell$
then defines a pro-$\ell$ subgroup $F_K$ of $T_\ell(A_K)$ which is
isogenous to the cotoric part
\[
 T_\ell(A_K)/T_\ell(A_K)^f.
\]
The latter is pro-\'{e}tale and unramified by \SGARefLink{sga7.sga7i.eix.n5-6}{5.6}, and the same is
therefore true of $F_K$. It follows from \SGARefLink{sga7.sga7i.eix.n5-7}{5.7} that $F_K$ is contained
in $T_\ell(A_K)^f$; thus $F_K=0$, proving that $U_\ell$ is equal to
its fixed part.

% END INCLUDED SGA7I_English_current_buildable_source_20260823/components/147_expose_IX_remark_5_11_through_corollary_5_15.tex

% BEGIN INCLUDED SGA7I_English_current_buildable_source_20260823/components/148_expose_IX_section_6.tex
% Expose IX, section 6.
% French transcription lines 1783--1842; authority physical pp. 404--411
% (scan indices 403--410; source folios 392--399); hard stop before section 7.
% All eight physical pages were checked directly in four overlapping
% 1100-dpi bands per page.  The print, including its small subscripts,
% primes, infinity signs, and trace notation, is clear at that resolution;
% no unresolved feature warranted a higher-resolution ambiguity crop.
%
% Transparent transcription/source corrections:
% - Near the end of 6.1, the working transcription corrupts p^infinity as
%   p^(bar infinity); the print has an ordinary infinity sign.
% - The largest p-divisible subgroup is printed p^N A_o^o (the image of
%   multiplication by p^N), not the left-subscript p^N torsion group in the
%   working transcription.
% - Raynaud's example twice has type-level slips in print: it first calls
%   A_K simple over K' rather than K, and then calls K' a ramified quadratic
%   extension of S rather than of K.  The surrounding argument forces K.
% - Formula (6.3.2) clearly prints A'_{K'} in both orthogonal complements;
%   the working transcription has A'_k.  The following normalization is
%   likewise in K', not K.
% - The forward reference ``6.5'' in 6.4 cannot exist (section 7 follows
%   immediately); the notion of essential good reduction was defined in
%   5.12, which is the type-correct reference used below.
% - The final trace function is clearly on pi-prime; the working
%   transcription's pi-exclamation is a glyph error.

\medskip
\noindent
\underline{6. Remarks on the construction of the toric part and the fixed
part of} $T_p(A)$ \underline{in the case of non-semistable reduction}.

\medskip
\noindent
\SGARefTarget{sga7.sga7i.eix.n6-1}{6.1}. The construction of the toric part of $T_\ell(A_K)$ given in \SGARefLink{sga7.sga7i.eix.n5-1}{5.1}
still makes sense as soon as $S$ is complete, without assuming that
$A^o$ is $\ell$-divisible; that is (when $p=\ell$), without assuming
semistable reduction.  That hypothesis had been used only in order to
formulate the orthogonality theorem \SGARefLink{sga7.sga7i.eix.n5-2}{5.2}.  Moreover, one can construct
$T_p(A_K)^t$ while assuming only that $S$ is henselian: one shows that
the toric part $T_p(A_{\widehat K})^t$ over the completion, constructed
by the method of \SGARefLink{sga7.sga7i.eix.n5-1}{5.1}, comes from a Barsotti--Tate pro-subgroup of
$T_p(A_K)$, necessarily determined uniquely.

More precisely, one can construct a pro-$p$ Barsotti--Tate group
$T_p(A)^t$ \underline{over} $S$ and a homomorphism
\[
 T_p(A)^t\longrightarrow T_p(A)
\]
of projective systems which is a monomorphism term by term and whose
pullback to $\widehat S$ gives the analogous objects constructed by the
method of \SGARefLink{sga7.sga7i.eix.n5-1}{5.1}.  This follows by paraphrasing the construction of \SGARefLink{sga7.sga7i.eix.n5-1}{5.1},
except that, instead of lifting the whole maximal torus $T_o$ of $A_o$,
for each integer $\nu\geq0$ one merely lifts the subgroup
${}_{p^\nu}T_o$ to a multiplicative-type subgroup
$({}_{p^\nu}A^o)^t$ of ${}_{p^\nu}A^o$.

For the existence and uniqueness of such a lifting, note that the
functor
\[
 (\mathrm{Sch})^o_{/S}\longrightarrow\mathrm{Sets},
 \qquad
 S'\longmapsto
 \{\text{finite multiplicative-type subgroups of }A^o_{S'}\}
\]
is represented by an $S$-scheme $X$ of finite type (cf. the proof of
SGA 3 XI 3.12(a)), which is \'etale over $S$ by SGA 3 IX 3.6 bis.
Consequently, there is exactly one section of $X$ over $S$ extending a
given section of $X_s$ over $s$, which proves the assertion.

The same argument shows more generally that, for every
\underline{finite multiplicative-type subgroup} $H_o$ of $A_o$, there
is a unique finite multiplicative-type subgroup $H$ of $A$ extending
it.  Consequently, for every \underline{toric} Barsotti--Tate
$p$-subgroup $H_o$ of $A_o$ (that is, one whose components are of
multiplicative type), there is a unique Barsotti--Tate $p$-subgroup $H$
of $A$ extending it.  This in turn gives a Barsotti--Tate $p$-subgroup
of ${}_{p^\infty}A_K$, or, equivalently, a pro-$p$ Barsotti--Tate
subgroup of $T_p(A_K)$.

In particular, this procedure applies when $H_o$ is the largest toric
Barsotti--Tate $p$-subgroup of $A_o$.  It is obtained by taking the
largest $p$-divisible subgroup $\overline A_o$ of $A_o^o$
($\overline A_o=p^N A_o^o$ for $N$ large) and the ``toric part'' of
${}_{p^\infty}\overline A_o$ (defined, for example, by SGA 3 XVII
7.2.1).  We thus obtain a Barsotti--Tate subgroup of
${}_{p^\infty}A_K$ which might be called the \underline{toric part} of
the latter.

\medskip
\noindent
\SGARefTarget{sga7.sga7i.eix.n6-2}{6.2}. Again let $\overline A_o$ be the largest $p$-divisible subgroup
scheme of $A_o^o$, and consider the Barsotti--Tate $p$-group
\[
 F_o={}_{p^\infty}(\overline A_o).
\]
It is immediate that $F_o$ may be characterized as the largest
Barsotti--Tate $p$-subgroup of $A_o$ (or of $A_o^o$).  It is natural to
ask whether this Barsotti--Tate group can be extended to a
Barsotti--Tate subgroup $F$ of $A$, which would then play the role of
the \underline{fixed part} $T_\ell(A)^f$ of (\SGARefLink{sga7.sga7i.eix.n2-2-3-2}{2.2.3.2}) in the case in
which $A^o$ is $\ell$-divisible.

In the semistable case, where $\overline A_o=A_o^o$, the method of \SGARefLink{sga7.sga7i.eix.n2-2}{2.2}
answers the question affirmatively.  In the general case, whenever the
problem has a solution, it is readily seen to be unique.  The
corresponding pro-subgroups of $T_p(A_K)$ and $T_p(A'_K)$ (which again
deserve to be called the \underline{fixed parts} of the pro-groups in
question) still satisfy the orthogonality theorem \SGARefLink{sga7.sga7i.eix.n5-2}{5.2}, and the fixed
part $T_p(A_K)^f$ may still be characterized as in \SGARefLink{sga7.sga7i.eix.n5-8}{5.8}.

\medskip
\noindent
\SGARefTarget{sga7.sga7i.eix.n6-2-1}{6.2.1}. We shall nevertheless show that in general the question has no
affirmative answer; that is, \underline{the fixed part $T_p(A)^f$
cannot in general be defined} in the case of non-semistable reduction.
Indeed, by a lifting theorem of Serre--Tate [30], when $S$ is complete,
the existence of the fixed part is equivalent to the existence of
infinitesimal liftings $\overline A_n$ of $\overline A_o$ in the $A_n$,
or again to the existence of a smooth group scheme $\overline A$ over
$S$, an extension of an abelian scheme by a torus $T$, having special
fibre $\overline A_o$, together with a formal homomorphism
\[
 \widehat{\overline A}\longrightarrow\widehat A
\]
extending the inclusion on the special fibres.

Thus, when $\overline A_o$ is an abelian scheme, that is, when $T_o=0$,
this is equivalent to the existence of an abelian scheme $\overline A$
extending $\overline A_o$ and a homomorphism
$\overline A\longrightarrow A$ extending the inclusion of the special
fibres.  It will therefore suffice to give an example in which $T_o=0$,
$A_o^o$ is neither unipotent nor proper, and $A_K$ is simple over $K$.

\medskip
\noindent
\SGARefTarget{sga7.sga7i.eix.n6-2-2}{6.2.2}. Here is an example, due to M. Raynaud.  It is based on the
observation that, if $S$ is a trait, $S'$ is the normalization of $S$
in a finite separable extension $K'$ of $K$, and $A'$ is an abelian
scheme over $S'$, then the group scheme
\[
 A=\prod_{S'/S} A'/S'
\]
over $S$ is a N\'eron model: it identifies with the N\'eron model of its
generic fibre, which is indeed an abelian scheme.  This follows
immediately from the definitions and from the fact that $A'$ itself is
a N\'eron model over $S'$.  Then
\[
 A_o=\prod_{S'_s/s}(A'\times_{S'}S'_s)/S'_s
\]
is an extension of an abelian scheme, nonzero if $A'\neq0$, by a
connected unipotent scheme, nonzero if $A'\neq0$ and $S'$ is ramified
over $S$.

It therefore suffices to give an example in which $A'\neq0$, $S'$ is
ramified over $S$ (a complete local trait), and
\[
 A_K=\prod_{K'/K}A'_{K'}/K'
\]
is simple over $K$.  Take $K'/K$ to be a ramified quadratic extension
(such an extension always exists), and take $A'$ to be an elliptic
curve over $S'$.  Then $A$ has relative dimension two, and saying that
$A_K$ is simple over $K$ means that there is no nonzero homomorphism
\[
 B_K\longrightarrow A_K,
\]
where $B_K$ is an abelian scheme of dimension one over $K$.  Giving
such a homomorphism is equivalent to giving a nonzero homomorphism
\[
 B_{K'}\longrightarrow A'_{K'},
\]
which would necessarily be an isogeny.  Thus $A_K$ is simple over $K$
precisely when $A'_{K'}$ is not isogenous to an elliptic curve coming
from $K$.

If the contrary held, the $j$-invariant of $A'$ would be algebraically
dependent (over the prime field) on its conjugate $\overline j$ under
the nontrivial $K$-automorphism of $K'$.  If $j$ is chosen so that
\[
 a=j+\overline j=\operatorname{Tr}_{K'/K}j
\]
is transcendental over the prime field, this is equivalent to saying
that $j$ is algebraically dependent on $a$.  By cardinality
considerations, however, there exists an element $a\in V$
transcendental over the prime subfield $P$ of $K$, and then an element
$b\in V'$ (the normalization of $V$ in $K'$) which is transcendental
over $P(a)$ and satisfies
\[
 \operatorname{Tr}_{V'/V}b=0.
\]
It is enough to take $j=a+b\in V'$.  We may moreover suppose that
$b\in\mathfrak r(V')$, and that the image of $a$ in
$k=V/\mathfrak m$ is a prescribed element of $k$, chosen distinct from
the exceptional values $j_o=0$ and $j_o=12^3$.  This is possible even
when $k$ has only two elements, because then these two contemplated
values of $j_o$ coincide.  Thus the image $j_o$ of $j$ in $k(s')$ is
different from $0$ and $12^3$, and it is known that there then exists
an elliptic curve over $S'$ with $j$-invariant $j$.

\medskip
\noindent
\SGARefTarget{sga7.sga7i.eix.n6-3}{6.3}. Return to the case in which $S$ is merely assumed henselian.  In
analogy with the definitions of \SGARefLink{sga7.sga7i.eix.n4-1}{4.1} for $\ell\neq p$, one can then
define a filtration
\begin{equation}\tag{\SGARefTarget{sga7.sga7i.eix.n6-3-1}{6.3.1}}
 T_p(A_K)\supset T_p(A_K)^{\mathrm{ef}}
 \supset T_p(A_K)^{\mathrm{et}}\supset0
\end{equation}
by strict Barsotti--Tate pro-subgroups, called respectively the
\underline{essentially fixed part} and the
\underline{essentially toric part} of $T_p(A_K)$.

To do so, choose a finite Galois extension $K'$ of $K$ such that
$A_{K'}$ has semistable reduction over the normalization $S'$ of $S$
in $K'$.  The fixed part $T_p(A_{K'})^f$ of
$T_p(A_{K'})=(T_p(A_K))_{K'}$ is then defined (\SGARefLink{sga7.sga7i.eix.n2-2-3}{2.2.3}), as is the toric
part (\SGARefLink{sga7.sga7i.eix.n6-1}{6.1}).  The latter is moreover the orthogonal complement of the
fixed part $T_p(A'_{K'})^f$ in $T_p(A'_{K'})$ (\SGARefLink{sga7.sga7i.eix.n5-2}{5.2}).  Plainly
\[
 T_p(A_{K'})^t\subset T_p(A_{K'})^f
\]
(this follows from \SGARefLink{sga7.sga7i.eix.n5-3}{5.3}, or directly from the construction of the toric
part in \SGARefLink{sga7.sga7i.eix.n6-1}{6.1}), and the Galois group of $K'/K$ preserves these two strict
pro-subgroups of $(T_p(A_K))_{K'}$.  Descent therefore gives the desired
filtration (\SGARefLink{sga7.sga7i.eix.n6-3-1}{6.3.1}).  It is independent of the choice of $K'$, by the
semistable-reduction property \SGARefLink{sga7.sga7i.eix.n3-3}{3.3}.  Descent also plainly gives the
orthogonality property
\begin{equation}\tag{\SGARefTarget{sga7.sga7i.eix.n6-3-2}{6.3.2}}
 \begin{split}
 T_p(A_K)^{\mathrm{et}}
   &=\bigl(T_p(A'_{K'})^{\mathrm{ef}}\bigr)^\perp,\\
 T_p(A_K)^{\mathrm{ef}}
   &=\bigl(T_p(A'_{K'})^{\mathrm{et}}\bigr)^\perp.
 \end{split}
\end{equation}

One could give an intrinsic characterization of the filtration (\SGARefLink{sga7.sga7i.eix.n6-3-1}{6.3.1})
in terms of the Barsotti--Tate pro-group $X=T_p(A_K)$ over the fraction
field $K$ of $S$, in the spirit of \SGARefLink{sga7.sga7i.eix.n5-8}{5.8}---at least, for the moment, in
unequal characteristics.  The essentially fixed part is the largest
Barsotti--Tate pro-subgroup having ``essentially good reduction over
$S$'': after a finite separable extension $K'$ of $K$, it gives a
Barsotti--Tate pro-group $X_{K'}$ having good reduction over the
normalization $S'$ of $S$ in $K'$, that is, extending to a
Barsotti--Tate pro-group over $S'$.  Such an extension is necessarily
unique up to a unique isomorphism, by the already-cited theorem of Tate
[33, Th. 4].  The essentially toric part is then defined in terms of
the essentially fixed part by orthogonality (\SGARefLink{sga7.sga7i.eix.n6-3-2}{6.3.2}).

Let us also record the following immediate consequence of \SGARefLink{sga7.sga7i.eix.n5-8}{5.8} and the
definitions; for $\ell=p$ it reduces to \SGARefLink{sga7.sga7i.eix.n5-12}{5.12}(a).

\medskip
\noindent
\underline{Proposition} \SGARefTarget{sga7.sga7i.eix.n6-3-3}{6.3.3}. \underline{Let $\ell$ be a prime
number.  If $\ell=p$, assume that $\operatorname{char}K=0$.}
\footnote{This hypothesis is probably unnecessary; cf. \SGARefLink{sga7.sga7i.eix.n5-11}{5.11}.}
\underline{Then $A_K$ has essential good reduction over $S$ (\SGARefLink{sga7.sga7i.eix.n5-12}{5.12}(a))
if and only if its essentially fixed part}
$T_\ell(A_K)^{\mathrm{ef}}$
\underline{(\SGARefLink{sga7.sga7i.eix.n4-1-1}{4.1.1} and \SGARefLink{sga7.sga7i.eix.n6-3-1}{6.3.1}) is all of $T_\ell(A_K)$.}

\medskip
\noindent
\SGARefTarget{sga7.sga7i.eix.n6-4}{6.4}. Recall that, when $\ell\neq p$, in \SGARefLink{sga7.sga7i.eix.n4-3}{4.3}(a) we exhibited a
remarkable integer-valued function on the inertia group $I$ of $S$,
independent of the choice of $\ell$.  When the residue field is perfect,
knowledge of this function suffices to reconstruct the conductor
exponent of $A$ (\SGARefLink{sga7.sga7i.eix.n4-5}{4.5}): one takes, on a sufficiently large quotient
$\Gamma=I/I'$ of $I$, the inner product of this function with the Artin
conductor on $\Gamma$ (cf. (\SGARefLink{sga7.sga7i.eix.n4-6-2}{4.6.2})).  It is useful to observe that this
integer-valued function on $I$ can also be recovered from the knowledge
of $T_\ell(A_K)$ for $\ell=p$---at least, for the time being, in unequal
characteristics.

When the residue field $k$ is finite, an analogous observation applies
to the $\mathbb Q$-valued function considered in \SGARefLink{sga7.sga7i.eix.n4-3}{4.3}(b), defined on the
subgroup $\pi'$ of $\pi$ consisting of the elements whose image in
$\pi_o$ has the form $f_q^{\,n}$, $n\in\mathbb Z$.  We indicate the
principle of the verification.  Using the filtration (\SGARefLink{sga7.sga7i.eix.n6-3-1}{6.3.1}), one is
reduced to attaching a suitable function (on $I$, respectively on
$\pi'$) to a pro-$p$ Barsotti--Tate group $Y$ over $K$ having
``essentially good reduction over $S$'' in the sense of \SGARefLink{sga7.sga7i.eix.n5-12}{5.12}.  Here $Y$
will be one of the three quotients arising from the filtration (\SGARefLink{sga7.sga7i.eix.n6-3-1}{6.3.1}).

In the case of \SGARefLink{sga7.sga7i.eix.n4-3}{4.3}(a), a suitable residue-field extension within $S$
(EGA $0_{\mathrm{III}}$ 10.3.1) reduces us to the case of a perfect
residue field, or even an algebraically closed one if desired.  Choose
a finite Galois extension $K'$ of $K$ such that $Y_{K'}$ has good
reduction over $S'$.  We then have a Barsotti--Tate group $Y_{S'}$ over
$S'$ and an action of $G$ on $Y_{S'}$ compatible with its action on
$S'$.  If $k'$ is the residue field of $S'$, then $G$ acts, not
necessarily faithfully, on $k'$ (whose invariant field is $k$), and on
$Y_{k'}$ compatibly with its action on $k'$.

By Dieudonn\'e theory [16 bis], the Barsotti--Tate pro-$p$ group $Y_{k'}$
has a canonically associated finite free module $M$ (of rank equal to
the rank of $Y_{k'}$) over the ring $W(k')$ of Witt vectors of $k'$.
The group $G$ plainly acts on $M$ compatibly with its actions on
$W(k')$, through its actions on $k'$.  Hence the inertia subgroup
$\Gamma$ of $G$ acts $W(k')$-linearly on $M$; since $\Gamma$ is a
quotient of the inertia group $I$, this gives a trace function on
$\Gamma$, and hence on $I$, which is the desired function.

A priori this function takes values in $W(k')$,
\footnote{In fact it is immediate that this function always takes values
in $W(k)$, since it is invariant under $\operatorname{Gal}(k'/k)$.}
but it agrees with the analogous function defined in terms of the
$T_\ell$ for $\ell\neq p$ whenever the action of $G$ on $Y_{S'}$ is
obtained from an action of $G$ on a smooth $p$-divisible group scheme
$C'$ over $S'$ by taking $Y=T_p(C')$ (cf. [25 bis, Th. 7, p. 737]).
It is therefore integer-valued in this case.  The case of \SGARefLink{sga7.sga7i.eix.n4-3}{4.3}(b) is
handled similarly: to define a trace function on $\pi'$, one uses the
method of the proof of \SGARefLink{sga7.sga7i.eix.n4-3-2}{4.3.2} to ``linearize'' on $M$ the action defined
by an element $g\in\pi$ lying above an $f_q^n$, with $n\geq0$.

% END INCLUDED SGA7I_English_current_buildable_source_20260823/components/148_expose_IX_section_6.tex

% BEGIN INCLUDED SGA7I_English_current_buildable_source_20260823/components/149_expose_IX_section_7_1.tex
% Expose IX, section 7 title and 7.1--7.1.6.
% French transcription lines 1845--1887; authority scan physical pp. 411--414
% (scan indices 410--413; source folios 399--402).  All four pages were
% checked directly in four overlapping 1100-dpi bands; no higher-resolution
% escalation was needed.  In 7.1.3 the print itself changes incoherently from
% M on G to L on G_n.  The type-correct reading L on A, L_n on A_n, M_n on
% B_n is fixed by the surrounding construction and by Proposition 7.1.5.

\medskip
\noindent
7. \underline{The Raynaud extension $G^\natural$ over $S$ associated
with the abelian scheme $A_K$ having semistable reduction.  Duality of
the abelian schemes $B,B'$ over $S$.}

\medskip
\noindent
\SGARefTarget{sga7.sga7i.eix.n7-1}{7.1}. Let $S$ be a Noetherian scheme which is the spectrum of a ring
$V$, separated and complete for the topology defined by an ideal
$\underline m$ of $V$.  Let $A$ be a smooth group scheme over $S$, let
$A^o$ be its identity component, and suppose that $A_o/A_o^o$ is
finite (hence finite etale) over $S_o$, and that $A_o^o$ is an
extension of an abelian scheme by a torus.  As usual, we have put
$A_n=A\times_S S_n$, where
$S_n=\operatorname{Spec}(V/\underline m^{n+1})$.  We shall be
interested chiefly in the case where $S$ is a trait, the generic fiber
$A_K$ is an abelian scheme having semistable reduction over $S$, and
$A$ is the Neron model of $A_K$.

Let $T_o$ be the maximal torus of $A_o$, so that
$B_o=A_o^o/T_o$ is an abelian scheme.  The construction of \SGARefLink{sga7.sga7i.eix.n5-1}{5.1} then
applies and gives a formal subtorus $\widehat T$ of the formal
completion $\widehat A$.  For every $n$, the quotient
\begin{equation}\tag{\SGARefTarget{sga7.sga7i.eix.n7-1-1}{7.1.1}}
 B_n=A_n^o/T_n
\end{equation}
is representable (SGA~3 $\mathrm{VI}_A$~3.2); more precisely, it is a
flat scheme of finite type over $S_n$ whose pullback to $S_o$ is the
abelian scheme $B_o$.  Consequently $B_n$ is an \underline{abelian
scheme} over $S_n$.  As $n$ varies, the $B_n$ are obtained from one
another by base change and therefore define a \underline{formal
abelian scheme}
\begin{equation}\tag{\SGARefTarget{sga7.sga7i.eix.n7-1-2}{7.1.2}}
 \widehat B=(B_n)_{n\geqslant0}\,.
\end{equation}
We shall give conditions implying that this formal abelian scheme is
\underline{algebraizable} (EGA~III~5.4.1), more precisely, that it is
the formal completion of a projective abelian scheme over $S$, which
we shall then denote by $B$.  By loc. cit., $B$ is determined up to a
unique isomorphism (as an $S$-scheme; the group structure of
$\widehat B$ then comes uniquely from a group structure on $B$, again
by loc. cit.).  This abelian scheme $B$ will then be called the
\underline{abelian part} of the scheme $A$, and will sometimes be
denoted by $A^{\mathrm{ab}}$.

\medskip
\noindent
\SGARefTarget{sga7.sga7i.eix.n7-1-3}{7.1.3}. By EGA~III~5.4.5, projectivity of $\widehat B$ means that one
can find an invertible sheaf $\widehat{\mathbb M}$ on $\widehat B$,
that is, a sequence of invertible sheaves $\mathbb M_n$ on the $B_n$
which glue as $n$ varies, and which is relatively ample over $S$, that
is, such that $\mathbb M_o$ is ample.  The natural idea would be to
start with an ample invertible sheaf $\mathbb L$ on $A$ and show that,
for every $n$, its restriction $\mathbb L_n$ to $A_n$ is the pullback
of an invertible sheaf $\mathbb M_n$ on $B_n$.\footnote{The printed
text changes incoherently here from an invertible sheaf $\mathbb M$ on
$G$ to the restriction $\mathbb L_n$ of $\mathbb L$ to $G_n$.  The
notation above is the type-correct reading fixed by the surrounding
construction and by Proposition~\SGARefLink{sga7.sga7i.eix.n7-1-5}{7.1.5}.}  I do not know whether such
$\mathbb M_n$ necessarily exist (and strongly doubt it).  We shall
follow a slightly different method, constructing a
\underline{polarization} of $\widehat B$, that is, a projective system
of polarizations $\rho_n$ of the $B_n$ which glue.  It is known [5]
that the doubles of these polarizations canonically define
(symmetric) invertible sheaves $\mathbb M_n$ on the $B_n$.  By their
canonical nature these sheaves necessarily glue and furnish the
desired $\widehat{\mathbb M}$.

We shall interpret an element of $\operatorname{NS}(B_n/S_n)$ as a
symmetric divisorial correspondence on $(B_n,B_n)$ [5], that is, as a
biextension of $(B_n,B_n)$ by $\mathbb G_{mS_n}$.  By VIII~\SGARefLink{sga7.sga7i.eviii.n3-5}{3.5}, these
biextensions correspond bijectively to biextensions of $(A_n,A_n)$ by
$\mathbb G_{mS_n}$.  To obtain a coherent system of such biextensions
$W_n$, it suffices to start with
\begin{equation}\tag{\SGARefTarget{sga7.sga7i.eix.n7-1-4}{7.1.4}}
 W,\quad\text{a symmetric biextension of $(A,A)$ by $\mathbb G_{mS}$},
\end{equation}
and take the $W_n$ obtained from $W$ by the base changes
$S_n\longrightarrow S$.  We must in addition suppose that the element
of $\operatorname{NS}(B_o)$ defined by $W_o$ is a polarization of
$B_o$; we thereby obtain a polarization of $\widehat B$.

\medskip
\noindent
\underline{Proposition \SGARefTarget{sga7.sga7i.eix.n7-1-5}{7.1.5}.} \underline{With the preceding notation
for $S,A,\widehat B$, suppose that $S$ is normal.  Then $\widehat B$ is
algebraizable, and the abelian scheme $B$ over $S$ from which it comes
is projective over $S$.}

We may plainly suppose that $A=A^o$ and that $S$ is integral.  By
Raynaud [23, XI~1.13], the hypothesis that $S$ is normal implies that
$A$ is quasi-projective over $S$.  Choose a relatively ample
invertible sheaf $\mathbb L$ on $A$, and form the birigidified
invertible sheaf
\[
 \delta(\mathbb L)=
 \pi^*(\mathbb L)\operatorname{pr}_1^*(\mathbb L)^{-1}
 \operatorname{pr}_2^*(\mathbb L)^{-1}
\]
on $A\times_S A$, as in VIII~\SGARefLink{sga7.sga7i.eviii.n4-12}{4.12}.  Since the fibers of $A_o$ have
unipotent rank zero, the same is true of the fibers of $A$ (this
follows easily from \SGARefLink{sga7.sga7i.eix.n3-1}{3.1}).  Therefore, by VIII~\SGARefLink{sga7.sga7i.eviii.n7-5}{7.5},
$\delta(\mathbb L)$ defines a biextension $W$ of $(A,A)$, plainly
symmetric, which by VIII~\SGARefLink{sga7.sga7i.eviii.n4-13}{4.13} (that is, essentially by Raynaud
[23, XI~1.11]) satisfies the required ampleness condition at the
points of $S_o$.

\medskip
\noindent
\underline{Remark \SGARefTarget{sga7.sga7i.eix.n7-1-6}{7.1.6}.} Without assuming that $S$ is normal, one
can prove that $\widehat B$ is always algebraizable by reducing to the
normal case as in \SGARefLink{sga7.sga7i.eix.n13-5-3}{13.5.3}(a).

% END INCLUDED SGA7I_English_current_buildable_source_20260823/components/149_expose_IX_section_7_1.tex

% BEGIN INCLUDED SGA7I_English_current_buildable_source_20260823/components/150_expose_IX_lemma_7_2_1.tex
% Expose IX, 7.2 and Lemma 7.2.1 with proof 7.2.1.1--7.2.1.5.
% French transcription lines 1889--1938; authority scan physical pp. 414--418
% (scan indices 413--417; source folios 402--406).  All pages were checked
% directly in four overlapping 1100-dpi bands.  No unresolved feature needed
% higher-resolution escalation.
% Source notation corrected transparently: the adic base is Spf(V), not the
% printed Spf(A), since A is the varying group scheme; likewise Phi_o and
% Phi_n live over S_o and S_n, although 7.2.1.3 prints S.  The parenthetical
% source citation for representability in 7.2.1.3 is blank and is not guessed.

\medskip
\noindent
\SGARefTarget{sga7.sga7i.eix.n7-2}{7.2}. We now record the following lemma.

\medskip
\noindent
\underline{Lemma \SGARefTarget{sga7.sga7i.eix.n7-2-1}{7.2.1}.} Let $S$ be the spectrum of a Noetherian ring
$V$, separated and complete for the topology defined by an ideal
$\underline m$ of $V$.  Let $C$ be the category of group schemes
$\Lambda$ over $S$, smooth and of finite type over $S$, such that
$\Lambda/\Lambda^o$ is a finite etale group scheme over $S$ and
$\Lambda^o$ is an extension of an abelian scheme $B$ by an isotrivial
torus $T$ (SGA~3 IX~1.1).  Let $C'$ be the category of group objects
in the category of formal schemes of finite type (EGA~I~10.3.3) over
$\widehat S=\operatorname{Spf}(V)$, and consider the functor
\[
 A\longmapsto\widehat A
\]
from $C$ to $C'$ which takes the formal completion along
$A_o=A\times_S S_o$, where
$S_o=\operatorname{Spec}(V/\underline m)$.\footnote{At both occurrences
in this statement the print writes $\operatorname{Specf}(A)$, and it
also writes $\operatorname{Spec}(A/\underline m^{n+1})$ below.  Since
$V$ is the adic base ring and $A$ denotes the varying group scheme,
the type-correct base is $\operatorname{Spf}(V)$ and
$S_n=\operatorname{Spec}(V/\underline m^{n+1})$.} Then:

\smallskip
\noindent
\underline{(a) The preceding functor is fully faithful.}

\smallskip
\noindent
\underline{(b) Let $\mathcal A$ be an object of $C'$, smooth over
$\widehat S$}; that is, suppose that the ordinary schemes
\[
 A_n=\mathcal A\times_{\widehat S}S_n,
 \qquad
 S_n=\operatorname{Spec}(V/\underline m^{n+1}),
\]
are smooth over $S_n$.  Suppose that $A_o/A_o^o$ is a finite etale
group scheme over $S_o$, and that $A_o^o$ is an extension of an
abelian scheme $B_o$ by an isotrivial torus $T_o$.  Put
\[
 \mathcal A^o=(A_n^o)_{n\geq0},
 \qquad
 \mathcal A/\mathcal A^o=(A_n/A_n^o)_{n\geq0}.
\]
\underline{Then $\mathcal A/\mathcal A^o$ is finite etale over
$\widehat S$, and $\mathcal A^o$ is, in an essentially unique manner
and functorially in $\mathcal A$, an extension of a formal abelian
scheme $\mathcal B$ over $\widehat S$ by a formal torus $\mathcal T$
over $\widehat S$, reducing respectively to $B_o$ and $T_o$.  The
formal group scheme $\mathcal A$ belongs to the essential image of
$C\longrightarrow C'$---or, as we shall also say, is
``algebraizable''---if and only if $\mathcal B$ is algebraizable.}

\medskip
\noindent
\SGARefTarget{sga7.sga7i.eix.n7-2-1-1}{7.2.1.1}. First observe, over an arbitrary scheme $S$, that if $A$ is
a smooth group scheme such that $A^o$ is an extension of an abelian
scheme by a torus, this extension structure is essentially unique and
is functorial in $A$ in the evident sense.  Indeed, the assertion
reduces at once to the fact that every morphism from a torus $T$ to an
abelian scheme $B$ is zero.  This is well known on geometric fibers
[15], and SGA~3 IX~5.3 then gives the conclusion.  Applying the result
to the compatible systems of smooth group schemes $A_n$ over the
$S_n$ occurring in \SGARefLink{sga7.sga7i.eix.n7-2-1}{7.2.1} gives the analogous result for formal group
schemes over $S$.  Moreover, we already observed in \SGARefLink{sga7.sga7i.eix.n7-1}{7.1} that, if
$A_o^o$ is an extension of an abelian scheme by a torus, the same is
true of $A_n^o$.  This proves the first assertion of (b).

\medskip
\noindent
\SGARefTarget{sga7.sga7i.eix.n7-2-1-2}{7.2.1.2}. Suppose next that the formal group scheme $\mathcal A$ over
$S$ is of the form $\widehat A$, with $A\in\operatorname{ob}C$.  If
$A^o$ is an extension of an abelian scheme $B$ by a torus $T$, then
plainly $\mathcal B\simeq\widehat B$.  Thus algebraizability of
$\mathcal A$ implies algebraizability of $\mathcal B$.

Conversely, suppose that $\mathcal B$ is algebraizable, that is, that
it comes from an abelian scheme $B$, and let us prove that $\mathcal A$
is algebraizable.  We first prove that
$\mathcal A^o=(A_n^o)$ is algebraizable.  By VIII~\SGARefLink{sga7.sga7i.eviii.n3-7}{3.7}, the extensions
$A_n^o$ of $B_n$ by $T_n$ may be interpreted as homomorphisms
\begin{equation}\tag{*}
 \underline M_n\longrightarrow B'_n,
\end{equation}
where $\underline M_n$ is the twisted constant group over $S_n$ dual
to $T_n$ (SGA~3 X~5.7), and $B'_n$ is the abelian scheme dual to
$B_n$.  Restriction to $S_o$, or to any $S_n$, induces an equivalence
from the category of isotrivial tori over $S$ to the corresponding
category over $S_o$ (SGA~3 X~3.2).  Hence the $T_n$ come from a torus
$T$ over $S$, determined up to a unique isomorphism.  Denote its dual
group by $\underline M$; its restriction to $S_n$ is canonically
isomorphic to $\underline M_n$.

The isotriviality of $T$ means that $\underline M$ is a disjoint union
of $S$-schemes $\underline M^i$ which are finite over $S$
(SGA~3 X~5.11).  It follows at once that the homomorphisms $(*)$ come
from a unique homomorphism
\begin{equation}\tag{**}
 \underline M\longrightarrow B',
\end{equation}
where $B'$ is the abelian scheme dual to $B$.  In turn, this
homomorphism is interpreted as an extension of $B$ by $T$, say $A^o$,
which plainly supplies an algebraization of $\mathcal A^o$.

\medskip
\noindent
\SGARefTarget{sga7.sga7i.eix.n7-2-1-3}{7.2.1.3}. On the other hand, by hypothesis $A_o/A_o^o$ is a finite
group scheme $\Phi_o$ over $S_o$, and the $A_n/A_n^o$ are likewise
finite group schemes $\Phi_n$ over $S_n$; they are necessarily etale
over $S_n$.\footnote{The print writes $S$ rather than the fiber bases
$S_o$ and $S_n$ in these sentences, and leaves blank the
parenthetical citation from which representability of $\Phi_n$ is
said to follow.  The bases are restored from the statement of the
lemma; the missing citation is not guessed.}  There is therefore an
etale group scheme $\Phi$ over $S$, determined up to a unique
isomorphism, whose formal completion is the formal group scheme
$\mathcal A/\mathcal A^o$ defined by the $\Phi_n$.

Since $G=A^o$ is divisible, the category of extensions, as fppf
sheaves, of a finite locally constant group over $S$ such as $\Phi$
by the group scheme $A^o$ is canonically equivalent to the category
of torsors under the sheaf
\[
 H=\underline{\operatorname{Hom}}_{\mathrm{gr}}(\Phi,G),
\]
and this equivalence is compatible with every base change.  Applying
the same description to extensions of $\Phi_n$ by $G_n$ in terms of
torsors under
\[
 H_n\simeq
 \underline{\operatorname{Hom}}_{\mathrm{gr}}(\Phi_n,G_n),
\]
the $A_n$ define a coherent system of torsors $P_n$ under the $H_n$.
Because $G$ is an extension of an abelian scheme by a torus, it follows
readily that $H$ is a group scheme \underline{finite and locally free
over $S$}.  Consequently, the functor
\[
 P\longmapsto\widehat P=(P_n)
\]
is an equivalence between the category of torsors over $S$ under $H$
and the category of ``formal'' torsors $\widehat P$ under
$\widehat H$, that is, coherent systems of torsors $P_n$ under the
$H_n$.

Thus the $P_n$ defined by the extensions $A_n$ come from a single
torsor $P$ over $S$ under $H$.  We obtain an extension $A$ of $\Phi$
by $G=A^o$ which algebraizes $\mathcal A$, except that its
representability remains to be checked.  Regard $A$ as a torsor over
$\Phi$ under $A^o_\Phi$.  This torsor becomes trivial after base
change by the finite, surjective, locally free morphism
\[
 P\times_S\Phi\longrightarrow\Phi,
\]
since the base change $P\longrightarrow S$ splits the extension $A$
of $\Phi$ by $A^o$.  One may therefore conclude by one of Raynaud's
many criteria for effectiveness of descent [23, XI~3.1(6)].

\medskip
\noindent
\SGARefTarget{sga7.sga7i.eix.n7-2-1-4}{7.2.1.4}. This completes the proof of (b), and only the full
faithfulness assertion (a) remains.  We leave it to the reader to
verify that this assertion is essentially contained in the method of
proof of (b), or, if preferred, to find an independent direct proof.

\medskip
\noindent
\SGARefTarget{sga7.sga7i.eix.n7-2-1-5}{7.2.1.5}. Notice that the appeal to Raynaud's delicate descent result
becomes unnecessary if $A^o$ is assumed quasi-projective over $S$, or,
equivalently, if $B$ is projective over $S$.  This condition holds in
particular when one starts from the formal group scheme $\widehat A$
over $S$, with $A$ as in \SGARefLink{sga7.sga7i.eix.n7-1-5}{7.1.5}.  This is the only case in which we
shall use \SGARefLink{sga7.sga7i.eix.n7-2-1}{7.2.1}.

% END INCLUDED SGA7I_English_current_buildable_source_20260823/components/150_expose_IX_lemma_7_2_1.tex

% BEGIN INCLUDED SGA7I_English_current_buildable_source_20260823/components/151_expose_IX_sections_7_2_2_through_7_2_6.tex
% Expose IX, 7.2.2--7.2.6.
% French transcription lines 1940--1967; authority scan physical pp. 418--419
% (scan indices 417--418; source folios 406--407).  Both pages were checked
% directly in four overlapping 1100-dpi bands; no escalation was needed.

\medskip
\noindent
\SGARefTarget{sga7.sga7i.eix.n7-2-2}{7.2.2}. We apply \SGARefLink{sga7.sga7i.eix.n7-2-1}{7.2.1} to the formal group scheme $\widehat A$
considered in \SGARefLink{sga7.sga7i.eix.n7-1}{7.1}, in the favorable case treated there in which the
formal abelian scheme $\widehat B$ is algebraizable (cf. \SGARefLink{sga7.sga7i.eix.n7-1-5}{7.1.5}), and
we suppose in addition that the torus $T_o$ over $S_o$ is isotrivial.
We thereby obtain a group scheme $A^\natural$ over $S$, called the
\underline{Raynaud group} of $A$ over $S$, characterized up to
isomorphism by a canonical isomorphism of formal group schemes over
$\operatorname{Spf}(V)=\widehat S$:
\begin{equation}\tag{\SGARefTarget{sga7.sga7i.eix.n7-2-3}{7.2.3}}
 (A^\natural)^\wedge\simeq\widehat A.
\end{equation}
It follows from \SGARefLink{sga7.sga7i.eix.n7-2-1}{7.2.1} that
\begin{equation}\tag{\SGARefTarget{sga7.sga7i.eix.n7-2-4}{7.2.4}}
 \Phi=A^\natural/A^{\natural o}
\end{equation}
is a finite etale group over $S$, and that the identity component
\[
 A^{\natural o}=A^{o\natural}
\]
is endowed with the structure of an extension
\begin{equation}\tag{\SGARefTarget{sga7.sga7i.eix.n7-2-5}{7.2.5}}
 0\longrightarrow T\longrightarrow A^{\natural o}
 \longrightarrow B\longrightarrow0,
\end{equation}
where $B$ is the abelian scheme over $S$ whose formal completion is
the formal abelian scheme $\widehat B$ constructed in \SGARefLink{sga7.sga7i.eix.n7-1}{7.1}, and $T$ is
characterized, up to a unique isomorphism, as the isotrivial torus
over $S$ lifting $T_o$.  On restriction to the $S_n$, the extension
structure (\SGARefLink{sga7.sga7i.eix.n7-2-5}{7.2.5}) gives the extension structures on
$A_n^{\natural o}=A_n^o$ defined by (\SGARefLink{sga7.sga7i.eix.n7-1-1}{7.1.1}).

\medskip
\noindent
\SGARefTarget{sga7.sga7i.eix.n7-2-6}{7.2.6}. It also follows immediately from \SGARefLink{sga7.sga7i.eix.n7-2-1}{7.2.1} that the Raynaud group
$A^\natural$, as well as the exact sequence (\SGARefLink{sga7.sga7i.eix.n7-2-5}{7.2.5}), depends
functorially on $A$ (satisfying the conditions under consideration),
and that its formation commutes with every base change
$S'\longrightarrow S$ associated with a homomorphism of adic
topological rings $V\longrightarrow V'$.

% END INCLUDED SGA7I_English_current_buildable_source_20260823/components/151_expose_IX_sections_7_2_2_through_7_2_6.tex

% BEGIN INCLUDED SGA7I_English_current_buildable_source_20260823/components/152_expose_IX_section_7_3.tex
% Expose IX, section 7.3.
% French transcription lines 1969--2015; authority scan physical pp. 419--421
% (scan indices 418--420; source folios 407--409). All three pages were checked
% directly in four overlapping 1100-dpi bands; no escalation was needed.

\medskip
\noindent
\SGARefTarget{sga7.sga7i.eix.n7-3}{7.3}. Let us now consider, for a fixed prime number $\ell$, the
pro-$\ell$ group $T_{\ell}(A^{o})$. Identifying a finite scheme over
$S$ with the corresponding formal scheme, and defining the
``\underline{fixed part}'' $T_{\ell}(A^{o})^{f}$ as usual, we obtain
an isomorphism as in (\SGARefLink{sga7.sga7i.eix.n5-1-5}{5.1.5}):

\begin{equation}\tag{\SGARefTarget{sga7.sga7i.eix.n7-3-1}{7.3.1}}
T_{\ell}(A^{o})^{f} \simeq T_{\ell}(\widehat{A}^{o}) \quad ,
\end{equation}
and, defining the ``\underline{toric part}'' as in (\SGARefLink{sga7.sga7i.eix.n5-1-6}{5.1.6}),

\begin{equation}\tag{\SGARefTarget{sga7.sga7i.eix.n7-3-2}{7.3.2}}
T_{\ell}(A^{o})^{t} = T_{\ell}(\widehat{T}) \simeq T_{\ell}(T)
 \subset T_{\ell}(A^{o})^{f} \simeq T_{\ell}(\widehat{A}^{o}) \quad ,
\end{equation}
we immediately obtain, for the ``\underline{abelian part},'' an
isomorphism

\begin{equation}\tag{\SGARefTarget{sga7.sga7i.eix.n7-3-3}{7.3.3}}
T_{\ell}(A^{o})^{ab} \stackrel{\mathrm{def}}{=}
T_{\ell}(A^{o})^{f}/T_{\ell}(A^{o})^{t}
 \simeq T_{\ell}(\widehat{B}) \quad .
\end{equation}
We have not yet used here the hypothesis that $\widehat B$ is
algebraizable, the hypothesis that allowed us to construct the
abelian scheme $B$ and the Raynaud extension of $B$ by $T$. Since
there is, of course, a canonical isomorphism (the special case of
(\SGARefLink{sga7.sga7i.eix.n7-3-1}{7.3.1}) applied to the abelian scheme $B$!)
$T_{\ell}(\widehat B)\simeq T_{\ell}(B)$, the isomorphism (\SGARefLink{sga7.sga7i.eix.n7-3-3}{7.3.3})
also gives us

\begin{equation}\tag{\SGARefTarget{sga7.sga7i.eix.n7-3-4}{7.3.4}}
T_{\ell}(A^{o})^{ab} \simeq T_{\ell}(B) \quad .
\end{equation}
Let us also note in passing that (\SGARefLink{sga7.sga7i.eix.n7-2-3}{7.2.3}) and (\SGARefLink{sga7.sga7i.eix.n7-3-1}{7.3.1}) imply a canonical
isomorphism

\begin{equation}\tag{\SGARefTarget{sga7.sga7i.eix.n7-3-5}{7.3.5}}
T_{\ell}(A^{o})^{f} \simeq T_{\ell}(A^{\natural o}) \quad ,
\end{equation}
since the right-hand side is already a profinite group, that is,
identical to its fixed part. The canonical exact sequence

\begin{equation}\tag{\SGARefTarget{sga7.sga7i.eix.n7-3-6}{7.3.6}}
0 \longrightarrow T_{\ell}(A^{o})^{t}
 \longrightarrow T_{\ell}(A^{o})^{f}
 \longrightarrow T_{\ell}(B) \longrightarrow 0 \quad ,
\end{equation}
which can in any case (without the algebraizability hypothesis on
$\widehat B$) be interpreted as the exact sequence obtained by
applying the functor $T_{\ell}$ to the exact sequence

\begin{equation*}
0 \longrightarrow \widehat{T} \longrightarrow \widehat{A}^{o}
 \longrightarrow \widehat{B} \longrightarrow 0 \quad ,
\end{equation*}
can also be regarded as obtained by applying the functor $T_{\ell}$
to the canonical exact sequence of ordinary group schemes

\begin{equation*}
0 \longrightarrow T \longrightarrow A^{\natural o}
 \longrightarrow B \longrightarrow 0 \qquad .
\end{equation*}

% END INCLUDED SGA7I_English_current_buildable_source_20260823/components/152_expose_IX_section_7_3.tex

% BEGIN INCLUDED SGA7I_English_current_buildable_source_20260823/components/153_expose_IX_section_7_4.tex
% Expose IX, section 7.4.
% French transcription lines 2017--2080; authority scan physical pp. 421--424
% (scan indices 420--423; source folios 409--412). All four pages were
% checked directly in four overlapping 1100-dpi bands; no escalation was needed.

\medskip
\noindent
\SGARefTarget{sga7.sga7i.eix.n7-4}{7.4}. Suppose now that a second group scheme $A'$ over $S$ is given,
satisfying the same hypotheses as $A$, and suppose that $A$ and $A'$
have connected fibres. Let $W$ be a biextension of $(A,A')$ by
$\mathbb G_{mS}$. We thereby obtain a pairing (VIII \SGARefLink{sga7.sga7i.eviii.n2-2}{2.2})

\begin{equation}\tag{\SGARefTarget{sga7.sga7i.eix.n7-4-1}{7.4.1}}
T_{\ell}(A) \times T_{\ell}(A') \longrightarrow
T_{\ell}(\mathbb G_{mS}) = \mathbb Z_{\ell}(1)
\end{equation}
and, by restriction to the fixed parts, a pairing

\begin{equation}\tag{\SGARefTarget{sga7.sga7i.eix.n7-4-2}{7.4.2}}
T_{\ell}(A)^{f} \times T_{\ell}(A')^{f}
 = T_{\ell}(\widehat A) \times T_{\ell}(\widehat {A'})
 \longrightarrow \mathbb Z_{\ell}(1) \qquad .
\end{equation}
This can also be interpreted as the pairing defined by the
biextension $\widehat W$ of $(\widehat A,\widehat {A'})$ by
$\widehat{\mathbb G}_{mS}$ obtained from $W$ by completion (here, for
simplicity, defining a biextension of formal group schemes as a
coherent family of biextensions $W_n$ of $(A_n,A'_n)$ by
$\mathbb G_{mS_n}$). Now we know (VIII \SGARefLink{sga7.sga7i.eviii.n3-5}{3.5}) that $W_n$ comes from a
biextension $W_n^{ab}$ of $(B_n,B'_n)$, where $B'_n$ is the abelian
part of $A'_n$, so that (\SGARefLink{sga7.sga7i.eix.n7-4-2}{7.4.2}) factors through

\begin{equation}\tag{\SGARefTarget{sga7.sga7i.eix.n7-4-3}{7.4.3}}
T_{\ell}(A)^{ab} \times T_{\ell}(A')^{ab}
 = T_{\ell}(\widehat B) \times T_{\ell}(\widehat {B'})
 \longrightarrow \mathbb Z_{\ell}(1) \qquad .
\end{equation}
This is also the pairing defined by the biextension

\begin{equation*}
\widehat W^{ab}=(W_n^{ab})
\end{equation*}
of $(\widehat B,\widehat {B'})$ by $\widehat{\mathbb G}_{mS}$. When
$\widehat B$ and $\widehat {B'}$ are algebraizable, so that abelian
schemes $B,B'$ over $S$ are available, EGA IV 5.1.4 shows that
$\widehat W^{ab}$ comes from a biextension, defined up to a unique
isomorphism,

\begin{equation*}
W^{ab}\quad\text{of}\quad(B,B')\quad\text{by}\quad\mathbb G_{mS},
\end{equation*}
and the corresponding pairing

\begin{equation}\tag{\SGARefTarget{sga7.sga7i.eix.n7-4-4}{7.4.4}}
T_{\ell}(B) \times T_{\ell}(B')
 \longrightarrow T_{\ell}(\mathbb G_{mS})
\end{equation}
is canonically isomorphic to (\SGARefLink{sga7.sga7i.eix.n7-4-3}{7.4.3}).

Writing $D(B')$ for the abelian scheme dual to $B'$, $W^{ab}$ defines
a homomorphism

\begin{equation}\tag{\SGARefTarget{sga7.sga7i.eix.n7-4-5}{7.4.5}}
u:B\longrightarrow D(B') \qquad ,
\end{equation}
whose properties reflect those of the pairings (\SGARefLink{sga7.sga7i.eix.n7-4-4}{7.4.4}), that is,
(\SGARefLink{sga7.sga7i.eix.n7-4-3}{7.4.3}). For example, the pairing (\SGARefLink{sga7.sga7i.eix.n7-4-3}{7.4.3}) is a perfect duality at a
point $t\in S$ if and only if
$u_t:B_t\longrightarrow D(B')_t$ is an isogeny whose kernel has rank
prime to $\ell$. This implies that $u$ itself is an isogeny (since
$S$ is local, hence connected), and $\ker u$ is flat and finite over
$S$, so that $\ker u$ has the same rank at every point. Thus the
condition just considered is independent of the chosen point $t$.

\medskip
\noindent
\SGARefTarget{sga7.sga7i.eix.n7-4-6}{7.4.6}. This completes in particular the proof of \SGARefLink{sga7.sga7i.eix.n5-4}{5.4}. Indeed, under
the conditions of that theorem (where $S$ may be assumed complete),
applying the preceding discussion to $A^o$ and ${A'}^o$ shows that,
for every $\ell$, the pairing (\SGARefLink{sga7.sga7i.eix.n7-4-3}{7.4.3}) is a perfect duality at the
generic point $\eta$, and hence also at the point $s$.

\medskip
\noindent
\SGARefTarget{sga7.sga7i.eix.n7-4-7}{7.4.7}. The preceding considerations likewise make the orthogonality
theorem \SGARefLink{sga7.sga7i.eix.n5-2}{5.2} clear, without recourse to Lemma \SGARefLink{sga7.sga7i.eix.n5-2-4}{5.2.4}. Indeed, along the
way in \SGARefLink{sga7.sga7i.eix.n7-4}{7.4} we established that the pairing
$(\SGARefLink{sga7.sga7i.eix.n5-2-3-1}{5.2.3.1})=(\SGARefLink{sga7.sga7i.eix.n7-4-2}{7.4.2})$ vanishes on the toric parts on both sides, and
therefore factors through the pairing (\SGARefLink{sga7.sga7i.eix.n7-4-3}{7.4.3}). It remains only to
establish that this pairing is nondegenerate on the generic fibres of
the pro-groups under consideration, knowing that it is so on the
special fibres. This follows from the fact that if (\SGARefLink{sga7.sga7i.eix.n7-4-5}{7.4.5}) induces an
isogeny on the special fibres, then it is an isogeny.

\medskip
\noindent
\SGARefTarget{sga7.sga7i.eix.n7-4-8}{7.4.8}. Return to the general conditions of \SGARefLink{sga7.sga7i.eix.n7-4}{7.4}, and suppose that an
open subset $U\subset S-S_o$ is given such that $A_U$ and $A'_U$ are
\underline{abelian schemes}, dual to one another through the
biextension $W_U$. Restricting (\SGARefLink{sga7.sga7i.eix.n7-3-6}{7.3.6}), and the analogous exact
sequence for $A'$, to $U$, we obtain canonical filtrations on the
Barsotti--Tate pro-$\ell$ groups $T_\ell(A_U)$ and $T_\ell(A'_U)$:

\begin{equation}\tag{\SGARefTarget{sga7.sga7i.eix.n7-4-9}{7.4.9}}
\left\{
\begin{aligned}
T_\ell(A_U)&\supset T_\ell(A_U)^f\supset T_\ell(A_U)^t\supset0,\\[4ex]
T_\ell(A'_U)&\supset T_\ell(A'_U)^f\supset T_\ell(A'_U)^t\supset0
\qquad .
\end{aligned}
\right.
\end{equation}
I claim that \underline{these two filtrations are dual to one another}
with respect to the pairing (\SGARefLink{sga7.sga7i.eix.n7-4-1}{7.4.1}), which here is a perfect duality
(VIII \SGARefLink{sga7.sga7i.eviii.n3-2}{3.2}); that is, relation (\SGARefLink{sga7.sga7i.eix.n5-2-2}{5.2.2}) and its symmetric counterpart
still hold, with $A_K,A'_K$ replaced by $A_U,A'_U$. We saw above that
this orthogonality assertion is also equivalent to saying that the
biextension $W^{ab}$ of $(B,B')$ by $\mathbb G_{mS}$ is a perfect
duality over $U$ (or, equivalently when $U$ is dense in $S$, over
$S$). It is enough to verify this property at every point $\eta$ of
$U$; the usual argument then reduces us to the already known case in
which $S$ is a trait with generic point $\eta$.

The orthogonality relation also allows us to extend formulas (\SGARefLink{sga7.sga7i.eix.n5-5-7}{5.5.7}),
(\SGARefLink{sga7.sga7i.eix.n5-5-8}{5.5.8}), and (\SGARefLink{sga7.sga7i.eix.n5-5-9}{5.5.9}) to the present case, replacing restriction to the
generic point $\eta$ by restriction to the open subset $U$.
Consequently, \SGARefLink{sga7.sga7i.eix.n5-6}{5.6} likewise extends to the present case. We refrain
from rewriting the preceding formulas and statement here in this
context.

% END INCLUDED SGA7I_English_current_buildable_source_20260823/components/153_expose_IX_section_7_4.tex

% BEGIN INCLUDED SGA7I_English_current_buildable_source_20260823/components/154_expose_IX_proposition_7_5_and_section_7_6.tex
% Expose IX, Proposition 7.5 and section 7.6.
% French transcription lines 2082--2103; authority scan physical pp. 424--425
% (scan indices 423--424; source folios 412--413). Both pages were checked
% directly in four overlapping 1100-dpi bands; no escalation was needed.

Let us summarize, from the point of view of N\'{e}ron models, the
results obtained concerning the abelian parts and their pairings.

\medskip
\noindent
\underline{Proposition} \SGARefTarget{sga7.sga7i.eix.n7-5}{7.5}. \underline{Let $S$ be a complete trait
with function field $K$. To every abelian scheme $A_K$ over $K$
having semistable reduction over $S$, the procedure of \SGARefLink{sga7.sga7i.eix.n7-1}{7.1}, applied
to the identity component $A^o$ of the N\'{e}ron model $A$ of $A_K$,
associates its abelian part $B=A^{ab}$, which is an abelian scheme
over $S$ depending functorially on $A_K$. For every other abelian
scheme $A'_K$ with semistable reduction, one can associate to every
divisorial correspondence $W_K$ on $(A_K,A'_K)$ a unique divisorial
correspondence $W^{ab}$ on $(A^{ab},{A'}^{ab})$ such that, for every
integer $n\geq0$, the pullback of $W_n^{ab}$ along
$A_n^o\longrightarrow B_n=A_n^{oab}$ and
${A'}_n^o\longrightarrow B'_n={A'}_n^{oab}$ is the biextension $W_n$
of $(A_n^o,{A'}_n^o)$ obtained from $W$ by reduction, where $W$ is
the canonical extension (\SGARefLink{sga7.sga7i.eix.n1-4}{1.4}) of $W_K$ to $(A^o,{A'}^o)$.

The formation of $W^{ab}$ is compatible with the formation of
pullbacks of correspondences by morphisms of abelian schemes
$\overline A_K\longrightarrow A_K$ and
$\overline {A'}_K\longrightarrow A'_K$. If $W_K$ is a perfect
duality, then so is $W^{ab}$. If $W_K$ is the divisorial
correspondence on $(A_K,A_K)$ associated with a polarization of
$A_K$, then $W^{ab}$ is associated with a polarization of $A^{ab}$.}

The compatibility of the formation of the $W^{ab}$ with pullbacks is
immediate from the indicated characterization of the $W^{ab}$. The
additional assertion concerning polarizations was established in
\SGARefLink{sga7.sga7i.eix.n7-1-5}{7.1.5} at the same time as the algebraizability of $\widehat B$.

\medskip
\noindent
\SGARefTarget{sga7.sga7i.eix.n7-6}{7.6}. Continue to suppose that $S$ is a complete trait, with $A$ the
N\'{e}ron model of an abelian scheme $A_K$ over $K$. Set

\begin{equation}\tag{\SGARefTarget{sga7.sga7i.eix.n7-6-1}{7.6.1}}
A_K^{\natural}=(A^{\natural})_K \quad ,
\end{equation}
where $A^{\natural}$ is the Raynaud group scheme (\SGARefLink{sga7.sga7i.eix.n7-2-2}{7.2.2}): it is an
extension of the finite etale group $\Phi$ by
$A^{\natural o}$, while $A^{\natural o}$ is an extension of the
abelian scheme $B$ by the torus $T$. Thus $A_K^{\natural}$ is itself
an extension of $\Phi_K$ by
$(A_K^{\natural})^o=(A^{\natural o})_K=(A^{o\natural})_K$, which is an
extension of the abelian scheme $B_K$ by the torus $T_K$. For every
prime number $\ell$, (\SGARefLink{sga7.sga7i.eix.n7-3-5}{7.3.5}) and (\SGARefLink{sga7.sga7i.eix.n7-3-2}{7.3.2}) give

\begin{equation}\tag{\SGARefTarget{sga7.sga7i.eix.n7-6-2}{7.6.2}}
T_{\ell}(A_K^{\natural}) \simeq T_{\ell}(A_K^{\natural o})
 \simeq T_{\ell}(A_K)^f \; , \qquad
T_{\ell}(T_K) \simeq T_{\ell}(A_K)^t \quad .
\end{equation}

\medskip
\noindent
\SGARefTarget{sga7.sga7i.eix.n7-6-3}{7.6.3}. By \SGARefLink{sga7.sga7i.eix.n3-3}{3.3}, the formation of $A_K^{\natural o}$ commutes with
every base change associated with a morphism of complete traits
$S'\longrightarrow S$.

% END INCLUDED SGA7I_English_current_buildable_source_20260823/components/154_expose_IX_proposition_7_5_and_section_7_6.tex

% BEGIN INCLUDED SGA7I_English_current_buildable_source_20260823/components/155_expose_IX_section_8_1.tex
% Expose IX, section 8.1.
% French transcription lines 2109--2154; authority scan physical pp. 426--428
% (scan indices 425--427; source folios 414--416). All three pages were
% checked directly in four overlapping 1100-dpi bands; no escalation was
% needed.

\medskip
\noindent
8. \underline{The Raynaud extension $A_K^{\natural o}$ in the
non-semistable case, and the ind-group $A_K^{\flat}$}

\medskip
\noindent
\SGARefTarget{sga7.sga7i.eix.n8-1}{8.1}. Continue to suppose that the trait $S$ is complete, but no
longer suppose necessarily that $A_K$ has semistable reduction. As
usual, let $K'$ be a finite Galois extension of $K$ such that
$A_{K'}$ has semistable reduction relative to the normalization $S'$
of $S$ in $K'$. Consider the Raynaud extension
$A_{K'}^{\natural}$ associated with $A_{K'}$ (\SGARefLink{sga7.sga7i.eix.n7-6-1}{7.6.1}). It is clear
that the Galois group of $K'/K$ acts by transport of structure on
this group scheme; hence, by Galois descent, it gives a group scheme
over $K$, denoted $A_{K,K'}^{\natural}$, which depends on the choice
of the extension $K'$. Its identity component, however, is
independent of this choice up to canonical isomorphism, by \SGARefLink{sga7.sga7i.eix.n7-6-3}{7.6.3},
and will simply be denoted $A_K^{\natural o}$. Descent shows that
$A_K^{\natural o}$ is an extension

\begin{equation}\tag{\SGARefTarget{sga7.sga7i.eix.n8-1-1}{8.1.1}}
0 \longrightarrow T_K \longrightarrow A_K^{\natural o}
 \longrightarrow B_K \longrightarrow 0 \quad ,
\end{equation}
where $T_K$ is a torus and $B_K$ an abelian scheme over $K$.
Likewise, by the definition of the essentially fixed and essentially
toric parts of $T_{\ell}(A_K)$ (\SGARefLink{sga7.sga7i.eix.n4-1}{4.1} and 6.5), the isomorphisms
(\SGARefLink{sga7.sga7i.eix.n7-6-2}{7.6.2}) give canonical isomorphisms

\begin{equation}\tag{\SGARefTarget{sga7.sga7i.eix.n8-1-2}{8.1.2}}
T_{\ell}(A_K^{\natural}) \simeq T_{\ell}(A_K)^{ef} \quad ,
\qquad T_{\ell}(T_K) \simeq T_{\ell}(A_K)^{et} \quad ,
\end{equation}
and hence, on passing to the quotient,

\begin{equation}\tag{\SGARefTarget{sga7.sga7i.eix.n8-1-3}{8.1.3}}
T_{\ell}(B_K) \simeq
T_{\ell}(A_K)^{ef}/T_{\ell}(A_K)^{et} \quad .
\end{equation}

This therefore determines two of the quotients occurring in the
filtration

\begin{equation}\tag{\SGARefTarget{sga7.sga7i.eix.n8-1-4}{8.1.4}}
T_\ell(A_K) \supset T_\ell(A_K)^{ef}
 \supset T_\ell(A_K)^{et} \supset 0
\end{equation}
Thus the extension of group schemes (\SGARefLink{sga7.sga7i.eix.n8-1-1}{8.1.1}) describes two quotients
in this filtration. For the third quotient
$T_\ell(A_K)/T_\ell(A_K)^{ef}$, (6.5.2) gives a canonical
isomorphism

\begin{equation}\tag{\SGARefTarget{sga7.sga7i.eix.n8-1-5}{8.1.5}}
T_\ell(A_K)/T_\ell(A_K)^{ef}
 \simeq D(T_\ell(A'_K)^{et}) \quad ,
\end{equation}
where $D=\underline{\operatorname{Hom}}(-,T_\ell(\mathbb{G}_m))$
denotes Cartier duality (in the sense of Barsotti--Tate
pro-groups), and $A'_K$ denotes the abelian scheme dual to $A_K$.
Consider then the exact sequence of type (\SGARefLink{sga7.sga7i.eix.n8-1-1}{8.1.1})

\begin{equation}\tag{\SGARefTarget{sga7.sga7i.eix.n8-1-6}{8.1.6}}
0 \longrightarrow T'_K \longrightarrow {A'_K}^{\natural o}
 \longrightarrow B'_K \longrightarrow 0
\end{equation}
associated with the abelian scheme $A'_K$ over $K$, together with
the corresponding isomorphisms of type (\SGARefLink{sga7.sga7i.eix.n8-1-2}{8.1.2}). Formula (\SGARefLink{sga7.sga7i.eix.n8-1-5}{8.1.5})
then gives

\begin{equation}\tag{\SGARefTarget{sga7.sga7i.eix.n8-1-7}{8.1.7}}
T_\ell(A_K)/T_\ell(A_K)^{ef}
 \simeq D(T_\ell(T'_K))
 \simeq \underline{M}_K\otimes_{\mathbb{Z}}\mathbb{Z}_\ell \quad ,
\end{equation}
where, generalizing the notation introduced in (\SGARefLink{sga7.sga7i.eix.n5-5-2}{5.5.2}), we have set

\begin{equation}\tag{\SGARefTarget{sga7.sga7i.eix.n8-1-8}{8.1.8}}
\underline{M}_K=D(T'_K) \quad , \qquad
\underline{M}'_K=D(T_K) \quad ,
\end{equation}
and where $D$ now denotes Cartier duality
$\underline{\operatorname{Hom}}(-,\mathbb{G}_{mK})$ in the sense of
ordinary group schemes. Thus $\underline{M}_K$ and
$\underline{M}'_K$ are twisted constant groups over $K$ which become
constant over $K'$ and isomorphic to groups of the form
$\mathbb{Z}^{r}$, where $r=\dim T=\dim T'$.

% END INCLUDED SGA7I_English_current_buildable_source_20260823/components/155_expose_IX_section_8_1.tex

% BEGIN INCLUDED SGA7I_English_current_buildable_source_20260823/components/156_expose_IX_sections_8_2_and_8_3.tex
% Expose IX, sections 8.2--8.3.
% French transcription lines 2159--2188; authority scan physical pp. 428--429
% (scan indices 427--428; source folios 416--417). Both pages were checked
% directly in four overlapping 1100-dpi bands; physical p. 430 was checked
% for the exact section-9 boundary. The authority's evident A_K -> A'_K
% polarization is restored where the working transcription misplaces the
% prime as A_{K'}; its evident "extensions ... de K'" slip is rendered
% according to the mathematically required extensions of K.

\medskip
\noindent
\SGARefTarget{sga7.sga7i.eix.n8-2}{8.2}. The hypotheses are those of \SGARefLink{sga7.sga7i.eix.n8-1}{8.1}. It is clear that the extension
(\SGARefLink{sga7.sga7i.eix.n8-1-1}{8.1.1}) depends functorially on the abelian scheme $A_K$, and that
its formation commutes with every morphism of traits
$\widetilde S\longrightarrow S$. Choose a polarization of $A_K$,
hence an isogeny $A_K\longrightarrow A'_K$. It induces a
homomorphism from (\SGARefLink{sga7.sga7i.eix.n8-1-1}{8.1.1}) to (\SGARefLink{sga7.sga7i.eix.n8-1-6}{8.1.6}), which is immediately seen, by
looking at the $T_\ell$, to be an isogeny term by term. Likewise,
using \SGARefLink{sga7.sga7i.eix.n7-5}{7.5} and descent, one sees that every biextension by
$\mathbb{G}_{mK}$ of abelian schemes $A_K$ and $\overline A_K$ over
$K$ defines a biextension of $(B_K,\overline B_K)$ by
$\mathbb{G}_{mK}$, where $B_K$ and $\overline B_K$ denote the
abelian quotients of the Raynaud extensions corresponding to $A_K$
and $\overline A_K$, respectively. When $W$ is a perfect duality, so
is the corresponding biextension of $(B_K,\overline B_K)$. In
particular, with the notation of \SGARefLink{sga7.sga7i.eix.n7-6}{7.6}, the abelian schemes $B_K$ and
$B'_K$ are canonically dual to one another.

The correspondence between biextensions on $(A_K,\overline A_K)$
and biextensions on $(B_K,\overline B_K)$ has evident functoriality
properties as $A_K$ and $\overline A_K$ vary. These follow by descent
from the analogous properties stated in \SGARefLink{sga7.sga7i.eix.n7-5}{7.5}, and we leave their
formulation to the reader. We also leave to the reader the
compatibility of the preceding construction with every change of
traits $\widetilde S\longrightarrow S$.

\medskip
\noindent
\SGARefTarget{sga7.sga7i.eix.n8-3}{8.3}. Return to the group $A_{K,K'}^{\natural}$ and study its
variation with $K'$. It is clear that two isomorphic extensions
$K'$ and $K'_1$ of $K$ give canonically isomorphic groups, so it is
enough to consider the subextensions $K'$ of the separable closure
$\overline K$ of $K$. If $K'\hookrightarrow K''$ are two such finite
Galois subextensions of $\overline K$, it follows immediately from
\SGARefLink{sga7.sga7i.eix.n3-3}{3.3} that there is a canonical \underline{open immersion}

\begin{equation}\tag{\SGARefTarget{sga7.sga7i.eix.n8-3-1}{8.3.1}}
A^{\natural}_{K,K'} \hookrightarrow A^{\natural}_{K,K''}
 \qquad ,
\end{equation}
which identifies the first group with an open subgroup of the
second. As $K'$ varies, the $A^{\natural}_{K,K'}$ thus form an
inductive system of smooth groups of finite type over $K$, whose
transition morphisms are open immersions. If we set

\begin{equation}\tag{\SGARefTarget{sga7.sga7i.eix.n8-3-2}{8.3.2}}
\Phi_{K,K'}=A^{\natural}_{K,K'}
 /{A^{\natural}_{K,K'}}^{o} \qquad ,
\end{equation}
the $\Phi_{K,K'}$ themselves form an inductive system of finite
etale groups over $K$, which will be studied in greater detail
below.

We shall set

\begin{equation}\tag{\SGARefTarget{sga7.sga7i.eix.n8-3-3}{8.3.3}}
A^{\flat}_{K}=\varinjlim_{K'} A^{\natural}_{K,K'} \qquad ,
\end{equation}
where the limit is taken in the category of fppf sheaves over $K$.
Since the transition morphisms in this inductive system are open
immersions, it is clear moreover that this limit is represented by
a group scheme over $K$ (in general non-quasi-compact), which we
shall naturally denote by the same symbol $A^{\flat}_{K}$. For this
group scheme one has

\begin{equation}\tag{\SGARefTarget{sga7.sga7i.eix.n8-3-4}{8.3.4}}
{A_{K}^{\flat}}^{o}=A_{K}^{\natural o} \quad , \qquad
A^{\flat}_{K}/{A_{K}^{\flat}}^{o}
 \simeq \varinjlim_{K'}\Phi_{K,K'} \qquad .
\end{equation}
The latter group will be determined in \SGARefLink{sga7.sga7i.eix.n11-10}{11.10}: it is canonically
isomorphic to $\underline{M}_{K}\otimes\mathbb{Q}/\mathbb{Z}$, with
the notation of (\SGARefLink{sga7.sga7i.eix.n8-1-8}{8.1.8}).

% END INCLUDED SGA7I_English_current_buildable_source_20260823/components/156_expose_IX_sections_8_2_and_8_3.tex

% BEGIN INCLUDED SGA7I_English_current_buildable_source_20260823/components/157_expose_IX_section_9_1.tex
% Expose IX, section 9 heading and 9.1.
% French transcription lines 2194--2214; authority scan physical pp. 430--431
% (scan indices 429--430; source folios 418--419). Both pages were checked
% directly in four overlapping 1100-dpi bands; no escalation was needed.

\medskip
\noindent
9. \underline{Definition of the monodromy pairing}
$\underline{M}_{\ell}\otimes\underline{M}'_{\ell}
 \longrightarrow \mathbb{Z}_{\ell S}$.

\medskip
\noindent
\SGARefTarget{sga7.sga7i.eix.n9-1}{9.1}. In this section, suppose that $S$ is henselian and that $A_K$
has semistable reduction over $S$, or, equivalently, that $A'_K$
has semistable reduction over $S$ (\SGARefLink{sga7.sga7i.eix.n3-5-1}{3.5.1}). We then have over $S$ two
twisted constant $\mathbb{Z}$-modules $\underline{M}$ and
$\underline{M}'$ (\SGARefLink{sga7.sga7i.eix.n5-5-2}{5.5.2}). Over a suitable finite etale covering of
$S$ (and hence over $S$ itself if $S$ is strictly local, that is, if
$k$ is separably closed), they split into constant groups with
values $M$ and $M'$, respectively, noncanonically isomorphic to
$\mathbb{Z}^{r}$, where $r$ is the common dimension of the maximal
tori of $A_o$ and $A'_o$. For every prime number $\ell$, we shall
consider the corresponding $\ell$-adic sheaves

\begin{equation}\tag{\SGARefTarget{sga7.sga7i.eix.n9-1-1}{9.1.1}}
\underline{M}_{\ell}=\underline{M}\otimes_{\mathbb{Z}}\mathbb{Z}_{\ell}
 \quad , \qquad
\underline{M}'_{\ell}=\underline{M}'\otimes_{\mathbb{Z}}\mathbb{Z}_{\ell}
 \qquad ,
\end{equation}
whose geometric meaning is made explicit by formulas (\SGARefLink{sga7.sga7i.eix.n5-5-7}{5.5.7}) and
(\SGARefLink{sga7.sga7i.eix.n5-5-8}{5.5.8}). In this section we shall define a canonical pairing, called
the \underline{monodromy pairing},

\begin{equation}\tag{\SGARefTarget{sga7.sga7i.eix.n9-1-2}{9.1.2}}
u_{\ell}:\underline{M}_{\ell}\otimes\underline{M}'_{\ell}
 \longrightarrow\mathbb{Z}_{\ell} \qquad ,
\end{equation}
which will be defined in terms of the Barsotti--Tate pro-$\ell$-groups
$T_{\ell}(A^{o})$ and $T_{\ell}({A'}^{o})$, and even in terms of the
generic fiber $T_{\ell}(A_K)$ of the first of these groups (subject,
when $\ell=p$, to the availability of Tate's theorem [33, Th.~4],
that is, for the time being, provided we are then in unequal
characteristic). When $\ell\neq p$, this is equivalent to saying
that the pairing (\SGARefLink{sga7.sga7i.eix.n9-1-2}{9.1.2}) is described in terms of the Galois
$\pi$-module

\begin{equation}\tag{\SGARefTarget{sga7.sga7i.eix.n9-1-3}{9.1.3}}
U=T_{\ell}(A_K)(\overline K)=T_{\ell}(A_{\overline K})
\end{equation}
considered in (\SGARefLink{sga7.sga7i.eix.n2-5-3}{2.5.3}). More precisely, we shall see that knowing the
pairing is essentially equivalent to knowing the action of the
monodromy group $I$ on $U$.

% END INCLUDED SGA7I_English_current_buildable_source_20260823/components/157_expose_IX_section_9_1.tex

% BEGIN INCLUDED SGA7I_English_current_buildable_source_20260823/components/158_expose_IX_section_9_2.tex
% Expose IX, section 9.2.
% French transcription lines 2216--2276; authority scan physical pp. 431--433
% (scan indices 430--432; source folios 419--421). All three pages were
% checked directly in four overlapping 1100-dpi bands; physical p. 434 was
% checked for the exact 9.3 boundary. The authority controls the evident
% transcription repairs x (not u) in the generator phrase, V^\perp (not
% V^1) in (9.2.4), T'_o in the second line of (9.2.7), and balanced
% parentheses in the displayed formulas.

\medskip
\noindent
\SGARefTarget{sga7.sga7i.eix.n9-2}{9.2}. \underline{The case $\ell\neq p$}. Begin with the simpler case
$\ell\neq p$, using the notation introduced in (\SGARefLink{sga7.sga7i.eix.n2-5-3}{2.5.3}). Note that,
up to torsion, the relation $V=U^I$ also means that $V^{\perp}$ is
the module generated by the elements $gx-x$, where $x\in U$ and
$g\in I$, and in any case one has

\begin{equation}\tag{\SGARefTarget{sga7.sga7i.eix.n9-2-1}{9.2.1}}
gx-x\in V^{\perp}\quad\text{for }x\in U,\ g\in I \qquad .
\end{equation}
Thus, for $g\in I$,

\begin{equation}\tag{\SGARefTarget{sga7.sga7i.eix.n9-2-2}{9.2.2}}
gx=x+u(g)\mathbin{.}x \qquad ,
\end{equation}
where $g\longmapsto u(g)$ is a homomorphism from $U$ to
$V^{\perp}$ which vanishes on $V=U^I$, or equivalently

\begin{equation}\tag{\SGARefTarget{sga7.sga7i.eix.n9-2-3}{9.2.3}}
u\in\operatorname{Hom}(U/V,V^{\perp})
 \simeq (U/V)^{\vee}\otimes V^{\perp} \qquad .
\end{equation}
In the case of semistable reduction under consideration, that is,
when $V^{\perp}\subset V$, formula (\SGARefLink{sga7.sga7i.eix.n9-2-2}{9.2.2}) says that the action of
$g$ on $U$ is a \underline{transvection} parallel to
$V^{\perp}=W$, and composition of these transvections is simply
addition of the corresponding homomorphisms $u$ in (\SGARefLink{sga7.sga7i.eix.n9-2-3}{9.2.3}).
Consequently, the action of the inertia group $I$ on $U$ is
completely described by a group homomorphism

\begin{equation}\tag{\SGARefTarget{sga7.sga7i.eix.n9-2-4}{9.2.4}}
I\longrightarrow (U/V)^{\vee}\otimes V^{\perp} \qquad .
\end{equation}

The orthogonality theorem, applied both to $A_K$ and to $A'_K$,
gives

\begin{equation*}
V^{\perp}=W \quad , \qquad
(U/V)^{\vee}\simeq W'(-1) \qquad ,
\end{equation*}
where $W'$ is the toric part of
$U'=T_{\ell}(A'(\overline K))$, and $(-1)$ denotes the Tate twist.
Moreover, the homomorphism (\SGARefLink{sga7.sga7i.eix.n9-2-4}{9.2.4}) factors through
$I(\ell)\simeq\mathbb{Z}_{\ell}(1)$, so that (\SGARefLink{sga7.sga7i.eix.n9-2-4}{9.2.4}) may be
interpreted as a canonical homomorphism

\begin{equation}\tag{\SGARefTarget{sga7.sga7i.eix.n9-2-5}{9.2.5}}
u:\mathbb{Z}_{\ell}(1)\longrightarrow W'\otimes W(-1)
 \qquad ,
\end{equation}
or, equivalently, as an element

\begin{equation}\tag{\SGARefTarget{sga7.sga7i.eix.n9-2-6}{9.2.6}}
u\in W'\otimes W(-2) \qquad .
\end{equation}
Recall that $W$ and $W'$ are the groups of $\overline K$-valued
points of the two twisted constant $\ell$-adic sheaves
$T_{\ell}(A^o)^t$ and $T_{\ell}({A'}^o)^t$, whose special fibers are
$T_{\ell}(T_o)$ and $T_{\ell}(T'_o)$, respectively. We obtain the
canonical isomorphisms

\begin{equation}\tag{\SGARefTarget{sga7.sga7i.eix.n9-2-7}{9.2.7}}
\left\{
\begin{array}{l}
W\simeq T_{\ell}(T_o)(\overline k)
 \simeq\check M\otimes\mathbb{Z}_{\ell}(1),\\[3ex]
W'\simeq T_{\ell}(T'_o)(\overline k)
 \simeq\check M'\otimes\mathbb{Z}_{\ell}(1),
\end{array}
\right.
\end{equation}
where $M$ and $M'$ denote the character groups over $\overline k$
of $T_o$ and $T'_o$, respectively:

\begin{equation}\tag{\SGARefTarget{sga7.sga7i.eix.n9-2-8}{9.2.8}}
T_o=D(M) \quad , \qquad T'_o=D(M') \qquad .
\end{equation}
It follows that the element $u$ of (\SGARefLink{sga7.sga7i.eix.n9-2-6}{9.2.6}), which describes the
action of the inertia group $I$ on the Tate module $U$ of $A_K$, is
interpreted as a canonical element of
$(\check M'\otimes_{\mathbb{Z}}\check M)
 \otimes_{\mathbb{Z}}\mathbb{Z}_{\ell}$. It is moreover invariant
under $\pi_o=\pi/I=\operatorname{Gal}(\overline k/k)$, since by
definition it is invariant under transport of structure:

\begin{equation}\tag{\SGARefTarget{sga7.sga7i.eix.n9-2-9}{9.2.9}}
u\in
\bigl((\check M'\otimes_{\mathbb{Z}}\check M)
 \otimes_{\mathbb{Z}}\mathbb{Z}_{\ell}\bigr)^I
\simeq
\operatorname{Hom}_{\mathbb{Z}}
 (\underline M\otimes\underline M',\mathbb{Z}_{\ell})
\simeq
\operatorname{Hom}_{\mathbb{Z}_{\ell}}
 (\underline M_{\ell}\otimes\underline M'_{\ell},
  \mathbb{Z}_{\ell}) \quad .
\end{equation}
Up to writing $u$ in place of $u_{\ell}$, this is the element
(\SGARefLink{sga7.sga7i.eix.n9-1-2}{9.1.2}) that we wished to define.

To define (\SGARefLink{sga7.sga7i.eix.n9-1-2}{9.1.2}) in the general case, we shall need a few
preliminaries, which will be developed in the next two sections.

% END INCLUDED SGA7I_English_current_buildable_source_20260823/components/158_expose_IX_section_9_2.tex

% BEGIN INCLUDED SGA7I_English_current_buildable_source_20260823/components/159_expose_IX_section_9_3.tex
% Expose IX, section 9.3.
% French transcription lines 2279--2430; authority scan physical pp. 433--439
% (scan indices 432--438; source folios 421--427). The authority pages were
% checked directly in overlapping 1100-dpi bands. Higher-detail inspection was
% reserved for the functor diagram. Evident print defects are corrected and
% disclosed in reader notes: the P/Q order in the graded pieces; F for E twice
% in the Brauer-pushout paragraph; part (d) for part (c); Ext^2(R,Q) for
% Ext^2(Q,P); omega for omega-prime on the lower functor arrow; and Y for Y'
% in the change-of-base formula.

\medskip
\noindent
\SGARefTarget{sga7.sga7i.eix.n9-3}{9.3}. \underline{Mixed extensions.}

Let $\underline C$ be an abelian category, and let

\begin{equation}\tag{\SGARefTarget{sga7.sga7i.eix.n9-3-1}{9.3.1}}
P,\ R,\ Q
\end{equation}
be three objects of $\underline C$. We are interested in classifying the
objects $X$ of $\underline C$ endowed with a three-step filtration

\begin{equation}\tag{*}
X^0=X\supset X^1\supset X^2\supset X^3=0
\end{equation}
and with given isomorphisms\footnote{The printed source assigns $P,R,Q$,
in that order, to these three graded pieces.  The immediately following exact
sequences and every subsequent construction require $Q,R,P$; the evident
interchange of $P$ and $Q$ has therefore been corrected here.}

\begin{equation*}
X^0/X^1\xrightarrow{\ \sim\ }Q,
\qquad X^1/X^2\xrightarrow{\ \sim\ }R,
\qquad X^2/X^3=X^2\xrightarrow{\ \sim\ }P.
\end{equation*}
Putting

\begin{equation*}
E=X^0/X^2,
\qquad F=X^1/X^3=X^1,
\end{equation*}
we see that $X$ gives rise to two ordinary extensions

\begin{equation}\tag{\SGARefTarget{sga7.sga7i.eix.n9-3-2}{9.3.2}}
0\longrightarrow P\longrightarrow F\longrightarrow R\longrightarrow0,
\qquad
0\longrightarrow R\longrightarrow E\longrightarrow Q\longrightarrow0,
\end{equation}
of $R$ by $P$ and of $Q$ by $R$, respectively. Our point of view here will
be to start, in addition to the objects (\SGARefLink{sga7.sga7i.eix.n9-3-1}{9.3.1}), with the extension
structures (\SGARefLink{sga7.sga7i.eix.n9-3-2}{9.3.2}), which are likewise assumed to be given, and to consider
the corresponding three-step filtered structures $X$.

More precisely, a \underline{mixed extension} of the extension $E$ by the
extension $F$ will mean an object $X$ of $\underline C$, endowed with a
three-step filtration (*) and with given isomorphisms

\begin{equation}\tag{\SGARefTarget{sga7.sga7i.eix.n9-3-3}{9.3.3}}
E\simeq X^0/X^2,
\qquad F\simeq X^1/X^3=X^1,
\end{equation}
compatible with the filtrations: they carry $X^1/X^2$ to
$\operatorname{Im}(R\longrightarrow E)$ and $X^2$ to
$\operatorname{Im}(P\longrightarrow F)$, respectively, and the
isomorphism $R\simeq F/P$ (and hence $R\simeq X^1/X^2$) deduced from
(\SGARefLink{sga7.sga7i.eix.n9-3-3}{9.3.3}) is required to equal the one deduced from the first exact sequence
in (\SGARefLink{sga7.sga7i.eix.n9-3-2}{9.3.2}).

For fixed $P,Q,R$ and fixed extensions $E$ and $F$, the mixed extensions of
$E$ by $F$ form a category in the evident way; it is in fact a groupoid
(every arrow is invertible). The automorphism group of a mixed extension
$X$ of $E$ by $F$ is canonically isomorphic to $\operatorname{Hom}(Q,P)$
through

\begin{equation}\tag{\SGARefTarget{sga7.sga7i.eix.n9-3-4}{9.3.4}}
\operatorname{Hom}(Q,P)
 \xrightarrow{\ \sim\ }
\operatorname{Aut}_{\mathrm{extpan}}(X),
\qquad
u\longmapsto\operatorname{id}_X+\bar u,
\end{equation}
where $\bar u$ is the endomorphism of the object of $\underline C$
underlying $X$ obtained by composing
$X\longrightarrow Q\xrightarrow{u}P\longrightarrow X$; the outer arrows
are respectively the canonical epimorphism and the canonical monomorphism
coming from the given mixed-extension structure on $X$.

One could make the structure of the category
$\operatorname{EXTPAN}(E,F)$ of mixed extensions of $E$ by $F$ precise by
regarding it as a ``torsor'' under the ``group-category'' (VII \SGARefLink{sga7.sga7i.evii.n2-5}{2.5})
$\operatorname{EXT}(Q,P)$ of extensions of $Q$ by $P$, using a natural
functor analogous to the Brauer composition of ordinary extensions:

\begin{equation}\tag{\SGARefTarget{sga7.sga7i.eix.n9-3-5}{9.3.5}}
\begin{aligned}
\omega:
\operatorname{EXT}(Q,P)\times\operatorname{EXTPAN}(E,F)
 &\longrightarrow \operatorname{EXTPAN}(E,F),\\
\omega(G,X)&=G\wedge X.
\end{aligned}
\end{equation}
For every object $X$ of $\operatorname{EXTPAN}(E,F)$, the partial functor
$G\longmapsto G\wedge X$ is an equivalence from
$\operatorname{EXT}(Q,P)$ to $\operatorname{EXTPAN}(E,F)$.

We shall give only a brief description of $\omega$. Regard a mixed
extension $X$ of $E$ by $F$ as defining an ordinary extension, also denoted
$X$, of $Q$ by $F$ (viewed as an object of $\underline C$). For every
extension $G$ of $Q$ by $P$, let $\bar G$ denote the extension of $Q$ by
$F$ induced by the inclusion $P\longrightarrow F$. Set

\begin{equation}\tag{\SGARefTarget{sga7.sga7i.eix.n9-3-6}{9.3.6}}
G\wedge X=\omega(G,X)
 \overset{\mathrm{dfn}}{=}\bar G\wedge X,
\end{equation}
where $\wedge$ denotes Brauer composition in the category of extensions of
$Q$ by $F$. The three-step filtration of the right-hand side $X'$ of
(\SGARefLink{sga7.sga7i.eix.n9-3-6}{9.3.6}) is clear, as is the second isomorphism of type (\SGARefLink{sga7.sga7i.eix.n9-3-3}{9.3.3}) attached to
$X'$.

To define the first isomorphism of type (\SGARefLink{sga7.sga7i.eix.n9-3-3}{9.3.3}) for $X'$, observe that the
functor

\begin{equation*}
Y\longmapsto\widetilde Y,
\qquad
\operatorname{EXT}(Q,F)\longrightarrow\operatorname{EXT}(Q,R),
\end{equation*}
induced by the given morphism $F\longrightarrow R$ in (\SGARefLink{sga7.sga7i.eix.n9-3-2}{9.3.2}) is
compatible with Brauer composition. It therefore sends $X'$ to an
extension $\widetilde X'$ canonically isomorphic to
$\widetilde{\bar G}\wedge\widetilde X$. The first isomorphism (\SGARefLink{sga7.sga7i.eix.n9-3-3}{9.3.3})
gives an isomorphism $\widetilde X\simeq E$, while
$\widetilde{\bar G}$---the extension induced from $G$ by the composite
$P\longrightarrow F\longrightarrow R$---is canonically isomorphic to the
trivial extension, since that composite $P\longrightarrow R$ is zero.
Consequently there is a canonical isomorphism
$\widetilde X'\simeq E$, which specifies on the right-hand side of (\SGARefLink{sga7.sga7i.eix.n9-3-6}{9.3.6})
its mixed-extension structure of $E$ by $F$.
\footnote{At both occurrences in this paragraph the printed source has
$F$ where the objects involved are extensions of $Q$ by $R$ and hence must
be identified with $E$.}

We leave to the reader the definition of $\omega$ on the arrows of
$\operatorname{EXT}(Q,P)\times\operatorname{EXTPAN}(E,F)$, the
verification that (\SGARefLink{sga7.sga7i.eix.n9-3-5}{9.3.5}) has the asserted property---namely, for every
\underline{fixed} mixed extension $X$, the functor
$G\longmapsto G\wedge X$ is an equivalence from
$\operatorname{EXT}(Q,P)$ to $\operatorname{EXTPAN}(E,F)$---and the proof
of the associativity formula

\begin{equation}\tag{\SGARefTarget{sga7.sga7i.eix.n9-3-7}{9.3.7}}
G_1\wedge(G_2\wedge X)
 \simeq (G_1\wedge G_2)\wedge X,
\end{equation}
where $G_1\wedge G_2$ denotes Brauer composition in
$\operatorname{EXT}(Q,P)$.

The preceding discussion immediately gives parts (a) and (b) of the
following result.

\medskip
\noindent
\underline{Proposition \SGARefTarget{sga7.sga7i.eix.n9-3-8}{9.3.8}. Given two extensions $E$ and $F$ as in
(\SGARefLink{sga7.sga7i.eix.n9-3-2}{9.3.2}) in the abelian category $\underline C$, consider the category
$\operatorname{EXTPAN}(E,F)$ of mixed extensions of $E$ by $F$. Then:}

\smallskip
\noindent
a) \underline{$\operatorname{EXTPAN}(E,F)$ is a groupoid. For every one
of its objects the automorphism group is canonically isomorphic to}
$\operatorname{Ext}^0(Q,P)=\operatorname{Hom}(Q,P)$.

\smallskip
\noindent
b) \underline{The set $\operatorname{Extpan}(E,F)$ of isomorphism classes
of objects of $\operatorname{EXTPAN}(E,F)$ is empty, or is naturally a
torsor under the group} $\operatorname{Ext}^1(Q,P)$.

\smallskip
\noindent
c) \underline{One can canonically define a class}

\begin{equation}\tag{\SGARefTarget{sga7.sga7i.eix.n9-3-9}{9.3.9}}
c(E,F)\in\operatorname{Ext}^2(Q,P),
\end{equation}
\underline{whose vanishing is necessary and sufficient for}
$\operatorname{Extpan}(E,F)\neq\varnothing$, \underline{that is, for the
existence of a mixed extension of $E$ by $F$.}

It remains to prove (c).\footnote{The printed source says ``prove (d)'',
although Proposition \SGARefLink{sga7.sga7i.eix.n9-3-8}{9.3.8} has only parts (a), (b), and (c).} A mixed
extension of $E$ by $F$ may also be described as an extension $X$ of $E$
by $P$ together with an isomorphism of extensions of $R$ by $P$

\begin{equation*}
X^1\simeq F,
\end{equation*}
where $X^1$ denotes the image of the extension $X$ under the morphism
$R\longrightarrow E$ in (\SGARefLink{sga7.sga7i.eix.n9-3-2}{9.3.2}). Such a mixed extension therefore exists
if and only if there is an element of $\operatorname{Ext}^1(E,P)$ whose
image in $\operatorname{Ext}^1(R,P)$ is the class $\xi$ of $F$. Consider
the exact sequence of the $\operatorname{Ext}^i(-,P)$ associated with the
second exact sequence in (\SGARefLink{sga7.sga7i.eix.n9-3-2}{9.3.2})\footnote{The printed final term is
$\operatorname{Ext}^2(R,Q)$; applying $\operatorname{Ext}(-,P)$ to
$0\to R\to E\to Q\to0$ gives the displayed $\operatorname{Ext}^2(Q,P)$,
which is also the group specified in (\SGARefLink{sga7.sga7i.eix.n9-3-9}{9.3.9}).}:

\begin{equation*}
\cdots\longrightarrow
\operatorname{Ext}^1(E,P)\longrightarrow
\operatorname{Ext}^1(R,P)
\xrightarrow{\partial}
\operatorname{Ext}^2(Q,P)\longrightarrow\cdots .
\end{equation*}
It is enough to define $c(E,F)$ by

\begin{equation}\tag{\SGARefTarget{sga7.sga7i.eix.n9-3-10}{9.3.10}}
c(E,F)=\partial\xi,
\qquad
\xi=\operatorname{cl}(F)\in\operatorname{Ext}^1(R,P).
\end{equation}

Now suppose that we are given an exact functor

\begin{equation}\tag{\SGARefTarget{sga7.sga7i.eix.n9-3-11}{9.3.11}}
r:\underline C\longrightarrow\underline C',
\end{equation}
where $\underline C'$ is another abelian category. Write
$P',Q',R',E',F'$ for the objects of $\underline C'$ obtained from
$P,Q,R,E,F$ by applying $r$, so that $E'$ is an extension of $Q'$ by
$R'$ and $F'$ is an extension of $R'$ by $P'$. Up to canonical
isomorphism, the following diagram of functors is commutative:

\begin{equation}\tag{\SGARefTarget{sga7.sga7i.eix.n9-3-12}{9.3.12}}
\begin{tikzcd}[column sep=large,row sep=large]
\operatorname{EXT}(Q,P)\times\operatorname{EXTPAN}(E,F)
  \arrow[r,"\omega"]\arrow[d]
& \operatorname{EXTPAN}(E,F)\arrow[d]\\
\operatorname{EXT}(Q',P')\times\operatorname{EXTPAN}(E',F')
  \arrow[r,"\omega'"']
& \operatorname{EXTPAN}(E',F').
\end{tikzcd}
\end{equation}
The vertical arrows are the evident functors induced by the exact functor
$r$, and $\omega'$ is the arrow analogous to $\omega$ in (\SGARefLink{sga7.sga7i.eix.n9-3-5}{9.3.5}), with
$\underline C$ replaced by $\underline C'$. The lower arrow in the
printed diagram is labelled $\omega$ without a prime; the explanatory
sentence identifies it as $\omega'$, which is the notation used here.
We obtain analogous compatibilities, with the functor $r$, of the
structures considered in \SGARefLink{sga7.sga7i.eix.n9-3-8}{9.3.8}; their explicit formulation is left to the
reader.

We shall use the following consequence of this commutativity. Suppose that
$\operatorname{EXTPAN}(E,F)$ is nonempty and choose an object $X$ in it.
Then, by means of $G'\longmapsto G'\wedge X'$, where $X'=r(X)$, the
category $\operatorname{EXTPAN}(E',F')$ is equivalent to
$\operatorname{EXT}(Q',P')$. Thus an object $Y'$ of
$\operatorname{EXTPAN}(E',F')$ is known up to isomorphism when the
isomorphism class of the corresponding object of
$\operatorname{EXT}(Q',P')$ is known, that is, when the corresponding
element

\begin{equation*}
c_X(Y')\in\operatorname{Ext}^1(Q',P')
\end{equation*}
is known.

If $X$ is replaced by another object $X_1=H\wedge X$, where $H$ is an
extension of $Q$ by $P$, then one immediately obtains
\footnote{The printed right-hand side of the following formula has
$c_X(Y)$; the fixed object in this construction is $Y'$.}

\begin{equation*}
c_{X_1}(Y')=c_X(Y')-\operatorname{cl}(H')
\end{equation*}
from the associativity formula (\SGARefLink{sga7.sga7i.eix.n9-3-7}{9.3.7}). Hence, for fixed $Y'$ and varying
$X$, the elements $c_X(Y')$ run through exactly one coset

\begin{equation}\tag{\SGARefTarget{sga7.sga7i.eix.n9-3-13}{9.3.13}}
c(Y')\in
\operatorname{Ext}^1(Q',P')/
\operatorname{Im}\operatorname{Ext}^1(Q,P),
\end{equation}
which is therefore canonically associated with the mixed extension $Y'$ of
$E'$ by $F'$. By construction, its vanishing is necessary and sufficient
for $Y'$ to be isomorphic to the transform under $r$ of a mixed extension
$X$ of $E$ by $F$.

% END INCLUDED SGA7I_English_current_buildable_source_20260823/components/159_expose_IX_section_9_3.tex

% BEGIN INCLUDED SGA7I_English_current_buildable_source_20260823/components/160_expose_IX_section_9_4.tex
% Expose IX, section 9.4.
% French transcription lines 2435--2576; authority scan physical pp. 439--444
% (scan indices 438--443; source folios 427--432). The authority pages were
% checked directly in overlapping 1100-dpi bands. The scan is clear throughout;
% no higher-resolution escalation was required. Evident print/transcription
% defects are corrected and disclosed at the point of use: omitted checks and
% delimiters, Lambda_s misread as a wedge, Q/P reversals in Hom and Ext, and the
% nonexistent reference 9.4.6.2 for 9.4.7.2.

\medskip
\noindent
\SGARefTarget{sga7.sga7i.eix.n9-4}{9.4}. Let $S$ be a henselian trait with generic point $\eta$, let $n$ be an
integer $>0$, put $\Lambda=\mathbb Z/n\mathbb Z$, and let $\underline C$ be
the category of $\Lambda$-modules on $S$ (for the fppf site on $S$; cf.
VIII \SGARefLink{sga7.sga7i.eviii.n0-2}{0.2}) and $\underline C'$ the category of $\Lambda$-modules on $\eta$.
Let $P,R,Q$ be objects of $\underline C$, and let $E,F$ be extensions of
$Q$ by $R$ and of $R$ by $P$, respectively (\SGARefLink{sga7.sga7i.eix.n9-3-2}{9.3.2}). Suppose that $Q$ is a
locally free $\Lambda_S$-module of finite type and that $P$ is represented
by a finite group of multiplicative type over $S$ (SGA 3 IX), so that it is
of the form

\begin{equation*}
P=D(Q^*),
\end{equation*}
where $D=\underline{\operatorname{Hom}}(-,\mathbb G_{mS})$ denotes Cartier
duality. Let $P_\eta,R_\eta,Q_\eta,E_\eta,F_\eta$ denote the restrictions of
$P,R,Q,E,F$ to $\eta$, and suppose that a mixed extension $X_\eta$ of
$E_\eta$ by $F_\eta$ is given. Under these conditions we shall canonically
define a pairing

\begin{equation}\tag{\SGARefTarget{sga7.sga7i.eix.n9-4-1}{9.4.1}}
c(X_\eta):Q\otimes_\Lambda Q^*\longrightarrow\Lambda_S
\qquad(\Lambda=\mathbb Z/n\mathbb Z).
\end{equation}

For this purpose, first \underline{suppose} that $S$ is \underline{strictly local}. Then $Q$
and $Q^*$ are constant, say $Q=Q_{oS}$ and $Q^*=Q^*_{oS}$, and giving
(\SGARefLink{sga7.sga7i.eix.n9-4-1}{9.4.1}) amounts to giving an element

\begin{equation}\tag{\SGARefTarget{sga7.sga7i.eix.n9-4-2}{9.4.2}}
c(X_\eta)\in\check Q_o\otimes\check Q_o^*,
\end{equation}
where the check denotes the ordinary dual $\operatorname{Hom}(-,\Lambda)$.
We apply the considerations of \SGARefLink{sga7.sga7i.eix.n9-3}{9.3} with the restriction functor
$r:\underline C\longrightarrow\underline C'$. First observe that a mixed
extension of $E$ by $F$ exists. Indeed, the obstruction to its existence
lies in $\operatorname{Ext}^2(Q,P)$ (\SGARefLink{sga7.sga7i.eix.n9-3-8}{9.3.8}(c)), and it is enough to prove

\begin{equation*}
\operatorname{Ext}^2(Q,P)=0.
\end{equation*}

This is a special case of the following lemma.

\medskip
\noindent
\underline{Lemma \SGARefTarget{sga7.sga7i.eix.n9-4-3}{9.4.3}. Under the preceding hypotheses, with $S$ strictly
local, one has:}

\smallskip
\noindent
a) \underline{canonical isomorphisms}

\begin{equation*}
\operatorname{Ext}^i_\Lambda(Q,P)\simeq
H^i(S,\check Q\otimes P),
\qquad
\operatorname{Ext}^i_\Lambda(Q_\eta,P_\eta)\simeq
H^i(\eta,\check Q\otimes P);
\end{equation*}

\noindent
b) \underline{a canonical isomorphism}

\begin{equation}\tag{b)}
\check Q\otimes P\simeq D(Q\otimes Q^*);
\end{equation}

\noindent
c) \underline{for every finite constant group $H=H_{oS}$ over $S$,}

\begin{equation}\tag{c)}
H^i(S,D(H))\simeq H^i(\eta,D(H)_\eta)=0
\qquad\underline{\text{for }}i\geq2;
\end{equation}

\noindent
d) \underline{canonical isomorphisms}

\begin{equation}\tag{d)}
\begin{aligned}
H^1(S,D(H))&\simeq\mathbb G_m(S)\otimes\check H_o
              =V^*\otimes\check H_o,\\
H^1(\eta,D(H)_\eta)&\simeq\mathbb G_m(\eta)\otimes\check H_o
              =K^*\otimes\check H_o,
\end{aligned}
\end{equation}
\underline{where the check in (d) denotes $\operatorname{Hom}(-,\mathbb
Q/\mathbb Z)$.}

The formulas in (a) and (b) follow immediately from the fact that $Q$ is a
locally free $\Lambda_S$-module of finite type. To prove (c), we plainly
reduce to the case $H_o=\mathbb Z/m\mathbb Z$, that is,
$D(H)=\mu_m$. The cohomology sequence associated with the Kummer exact
sequence

\begin{equation*}
0\longrightarrow\mu_m\longrightarrow\mathbb G_m
 \xrightarrow{\ m\cdot\operatorname{id}\ }\mathbb G_m
 \longrightarrow0
\end{equation*}
then reduces us to proving

\begin{equation*}
H^i(S,\mathbb G_m)=H^i(\eta,\mathbb G_m)=0
\qquad\text{for }i\geq1.
\end{equation*}
These groups are the groups $H^i_{\mathrm{\acute et}}$ computed on the
étale site [6, Theorem \SGARefLink{sga7.sga7i.eix.n11-7}{11.7}], and therefore vanish for $S$, since $S$ is
strictly local. Their vanishing for $\eta$ is likewise well known
[6, Theorem \SGARefLink{sga7.sga7i.eix.n1-1}{1.1}].

It remains to describe the isomorphisms in (d). Over an arbitrary base $U$
there is a canonical homomorphism

\begin{equation}\tag{\SGARefTarget{sga7.sga7i.eix.n9-4-3-1}{9.4.3.1}}
\mathbb G_m(U)\otimes\check H_o
\longrightarrow H^1(U,D(H_{oU})).
\end{equation}
If $H_o$ is annihilated by an integer $m>0$, this homomorphism is described
using the boundary map
$\mathbb G_m(U)\xrightarrow{\partial}H^1(U,\mu_m)$ associated with the
Kummer sequence: it is the composite

\begin{equation*}
\begin{aligned}
\mathbb G_m(U)\otimes\check H_o
&\longrightarrow H^1(U,\mu_m)\otimes\check H_o\\
&\longrightarrow H^1(U,\mu_m\otimes\check H_o)
 \longrightarrow H^1(U,D(H_{oU})).
\end{aligned}
\end{equation*}
Here the second arrow is deduced from the evident homomorphism
$\check H_o\longrightarrow
\operatorname{Hom}(\mu_m,\mu_m\otimes\check H_o)$, and the third is a
special case of (b), with $m$ in place of $n$. We leave it to the reader to
check that (\SGARefLink{sga7.sga7i.eix.n9-4-3-1}{9.4.3.1}) is independent of the choice of the integer $m$ such
that $mH_o=0$. We also claim that (\SGARefLink{sga7.sga7i.eix.n9-4-3-1}{9.4.3.1}) is always injective and that its
cokernel vanishes if $\operatorname{Pic}(U)=0$. (In fact, the cokernel is
always isomorphic to $\operatorname{Hom}(H_o,\operatorname{Pic}(U))$, but
this is immaterial here.) Indeed, we reduce to the case
$H_o=\mathbb Z/n\mathbb Z$, and the assertion then follows immediately from
the cohomology sequence of the Kummer sequence above. This proves (d), and
immediately yields:

\medskip
\noindent
\underline{Corollary \SGARefTarget{sga7.sga7i.eix.n9-4-4}{9.4.4}. With the notation of \SGARefLink{sga7.sga7i.eix.n9-4-3}{9.4.3}, there is a canonical
isomorphism}

\noindent
e)
\begin{equation*}
H^1(\eta,D(H)_\eta)/\operatorname{Im}H^1(S,D(H))
\simeq\check H_o,
\end{equation*}
\underline{and hence, on combining (a) and (b), a canonical isomorphism}

\begin{equation}\tag{\SGARefTarget{sga7.sga7i.eix.n9-4-5}{9.4.5}}
\operatorname{Ext}^1(Q_\eta,P_\eta)/
\operatorname{Im}\operatorname{Ext}^1(Q,P)
\simeq\check Q_o\otimes\check Q_o^*.
\end{equation}

Returning to the definition of (\SGARefLink{sga7.sga7i.eix.n9-4-2}{9.4.2}) when $S$ is strictly local, the
relation $\operatorname{Ext}^2(Q,P)=0$ makes it possible to define the
canonical element (\SGARefLink{sga7.sga7i.eix.n9-3-13}{9.3.13}), taking for $Y'$ the given mixed extension
$X_\eta$ of $E_\eta$ by $F_\eta$. Formula (\SGARefLink{sga7.sga7i.eix.n9-4-5}{9.4.5}) identifies this element
$c(X_\eta)$ with an element of $\check Q_o\otimes\check Q_o^*$, thereby
completing the definition of (\SGARefLink{sga7.sga7i.eix.n9-4-2}{9.4.2}). The general considerations of \SGARefLink{sga7.sga7i.eix.n9-3}{9.3}
also show that this element is the obstruction to extending $X_\eta$ to a
mixed extension $X$ of $E$ by $F$.

\medskip
\noindent
\SGARefTarget{sga7.sga7i.eix.n9-4-6}{9.4.6}. The definition of (\SGARefLink{sga7.sga7i.eix.n9-4-1}{9.4.1}) when $S$ is henselian now reduces
immediately to the strictly local case by Galois descent from the strict
henselization of $S$.

\medskip
\noindent
\underline{Proposition \SGARefTarget{sga7.sga7i.eix.n9-4-7}{9.4.7}. Let $P,Q,R,E,F$ be as above, except that the
trait $S$ need not be henselian. Consider the functor}

\begin{equation}\tag{\SGARefTarget{sga7.sga7i.eix.n9-4-7-1}{9.4.7.1}}
X\longmapsto X_\eta
\end{equation}
\underline{from the category of mixed extensions of $E$ by $F$ to the
category of mixed extensions of $E_\eta$ by $F_\eta$. For every mixed
extension $X_\eta$ of $E_\eta$ by $F_\eta$, consider the pairing}

\begin{equation}\tag{\SGARefTarget{sga7.sga7i.eix.n9-4-7-2}{9.4.7.2}}
c_o(X_\eta):Q_o\otimes Q_o^*\longrightarrow\Lambda_s
\end{equation}
\underline{induced by the pairing of type (\SGARefLink{sga7.sga7i.eix.n9-4-1}{9.4.1}) defined over the
henselization $S^h$ of $S$. Then the functor (\SGARefLink{sga7.sga7i.eix.n9-4-7-1}{9.4.7.1}) is fully faithful,
and a mixed extension $X_\eta$ of $E_\eta$ by $F_\eta$ belongs to its
essential image if and only if the pairing (\SGARefLink{sga7.sga7i.eix.n9-4-7-2}{9.4.7.2}) is zero.}

The printed target of (\SGARefLink{sga7.sga7i.eix.n9-4-7-2}{9.4.7.2}) is $\Lambda_s$; the French transcription's
$\wedge_s$ is a character-recognition error.

\medskip
\noindent
\SGARefTarget{sga7.sga7i.eix.n9-4-7-3}{9.4.7.3}. This statement still makes sense and remains valid if one assumes
only that $S$ is a connected regular noetherian scheme of dimension $1$.
Instead of the single pairing (\SGARefLink{sga7.sga7i.eix.n9-4-7-2}{9.4.7.2}), one must then consider a pairing
$Q_s\otimes Q_s^*\longrightarrow\Lambda_s$ for each closed point $s$ of
$S$. It will be most convenient to prove \SGARefLink{sga7.sga7i.eix.n9-4-7}{9.4.7} in this more general form.

\smallskip
\noindent
a) For every mixed extension $X$, the map

\begin{equation*}
\operatorname{End}(X)\longrightarrow\operatorname{End}(X_\eta)
\end{equation*}
is bijective. The two sides are respectively the sections over $S$ and
$\eta$ of
$\underline{\operatorname{Hom}}(Q,P)=D(Q\otimes Q^*)$; hence the assertion
follows from the normality of $S$ and the finiteness of
$D(Q\otimes Q^*)$ over $S$.\footnote{The printed proof has
$\underline{\operatorname{Hom}}(P,Q)$ here. Proposition 9.3.8(a), formula
(\SGARefLink{sga7.sga7i.eix.n9-3-4}{9.3.4}), and the displayed identification with $D(Q\otimes Q^*)$ all
require $\underline{\operatorname{Hom}}(Q,P)$.} It follows at once that if $X$
and $Y$ are two mixed extensions already known to be isomorphic, then

\begin{equation*}
\operatorname{Hom}(X,Y)\longrightarrow
\operatorname{Hom}(X_\eta,Y_\eta)
\end{equation*}
is bijective.

\smallskip
\noindent
b) If $X$ and $Y$ are two mixed extensions, every isomorphism
$X_\eta\xrightarrow{\sim}Y_\eta$ comes from an isomorphism
$X\longrightarrow Y$. By (a), the question is local for the étale
topology, so we may reduce to the case in which $S$ is strictly local. In
that case \SGARefLink{sga7.sga7i.eix.n9-4-3}{9.4.3} shows that the map

\begin{equation*}
\operatorname{Ext}^1_\Lambda(Q,P)
\longrightarrow\operatorname{Ext}^1_\Lambda(Q_\eta,P_\eta)
\end{equation*}
identifies with

\begin{equation*}
V^*\otimes\check H_o\longrightarrow K^*\otimes\check H_o,
\qquad H_o=Q_o\otimes Q_o^*.
\end{equation*}
The latter is injective, since the exact sequence
$0\longrightarrow V^*\longrightarrow K^*\longrightarrow\mathbb Z
\longrightarrow0$ is split. By the discussion of \SGARefLink{sga7.sga7i.eix.n9-3}{9.3} this implies that if
$X_\eta$ and $Y_\eta$ are isomorphic, then so are $X$ and $Y$; the
conclusion now follows from (a).\footnote{The printed Ext groups in this
paragraph have the arguments $(P,Q)$. The action of
$\operatorname{Ext}^1(Q,P)$ on mixed extensions in Proposition \SGARefLink{sga7.sga7i.eix.n9-3-8}{9.3.8}(b),
and Lemma \SGARefLink{sga7.sga7i.eix.n9-4-3}{9.4.3} itself, require $(Q,P)$.}

Thus (a) and (b) show that $X\longmapsto X_\eta$ is \underline{fully} faithful.

\smallskip
\noindent
c) Let $X_\eta$ be a mixed extension over $\eta$. We prove that it extends
to $S$ if and only if the pairings (\SGARefLink{sga7.sga7i.eix.n9-4-7-2}{9.4.7.2}) vanish. Here again, the
uniqueness already proved shows that the question is local for the étale
topology, reducing us once more to the strictly local case. The conclusion
then holds by construction, as observed immediately before \SGARefLink{sga7.sga7i.eix.n9-4-6}{9.4.6}.
\footnote{The printed proof cites the nonexistent pairings
``(9.4.6.2)''; these are the pairings (\SGARefLink{sga7.sga7i.eix.n9-4-7-2}{9.4.7.2}).}

This completes the proof of \SGARefLink{sga7.sga7i.eix.n9-4-7}{9.4.7}.

% END INCLUDED SGA7I_English_current_buildable_source_20260823/components/160_expose_IX_section_9_4.tex

% BEGIN INCLUDED SGA7I_English_current_buildable_source_20260823/components/161_expose_IX_section_9_5.tex
% Expose IX, section 9.5.
% French transcription lines 2578--2653; authority scan physical pp. 444--448
% (scan indices 443--447; source folios 432--436). The authority pages were
% checked directly in four overlapping 1100-dpi bands. The scan is clear
% throughout. Four internally forced print slips are corrected and disclosed:
% 9.4.6 for 9.4.7, X_eta for its extension X over S, eta for S as the base of
% the etale group Q(n), and 9.5.6 for the homomorphism (9.5.5).

\medskip
\noindent
\SGARefTarget{sga7.sga7i.eix.n9-5}{9.5}. Fix a prime number $\ell$, pro-$\ell$ Barsotti--Tate groups over $S$

\begin{equation}\tag{\SGARefTarget{sga7.sga7i.eix.n9-5-1}{9.5.1}}
P=(P(n))_{n\geq0},
\qquad R=(R(n))_{n\geq0},
\qquad Q=(Q(n))_{n\geq0},
\end{equation}
and extensions of pro-$\ell$ Barsotti--Tate groups

\begin{equation}\tag{\SGARefTarget{sga7.sga7i.eix.n9-5-2}{9.5.2}}
0\longrightarrow P\longrightarrow F\longrightarrow R\longrightarrow0,
\qquad
0\longrightarrow R\longrightarrow E\longrightarrow Q\longrightarrow0.
\end{equation}
Thus, for every integer $n\geq0$, these define extensions of
$\Lambda_n$-modules, where
$\Lambda_n=\mathbb Z/\ell^{\,n+1}\mathbb Z$:

\begin{equation}\tag{\SGARefTarget{sga7.sga7i.eix.n9-5-3}{9.5.3}}
\begin{gathered}
0\longrightarrow P(n)\longrightarrow F(n)\longrightarrow R(n)
 \longrightarrow0,\\
0\longrightarrow R(n)\longrightarrow E(n)\longrightarrow Q(n)
 \longrightarrow0.
\end{gathered}
\end{equation}
We assume that $Q$ is of étale type and $P$ of multiplicative type.
Finally, suppose that a pro-$\ell$ Barsotti--Tate group $X_\eta$ is given
over $\eta$ which is a \underline{mixed extension} of $E_\eta$ by $F_\eta$. By this
we mean that it is endowed, as a pro-$\ell$ Barsotti--Tate group, with a
three-step filtration whose partial extensions are supplied with
isomorphisms to $E_\eta$ and $F_\eta$, respectively. It follows that, for
every $n$, $X_\eta(n)$ is endowed with the structure of a mixed extension
of $E(n)_\eta$ by $F(n)_\eta$. We may therefore associate with it a pairing
of type (\SGARefLink{sga7.sga7i.eix.n9-4-1}{9.4.1}):

\begin{equation}\tag{\SGARefTarget{sga7.sga7i.eix.n9-5-4}{9.5.4}}
c(X(n)_\eta):Q(n)\otimes Q^*(n)
\longrightarrow\Lambda_n=\mathbb Z/\ell^{\,n+1}\mathbb Z,
\end{equation}
where the $Q^*(n)$ are the finite étale groups over $S$ defined by

\begin{equation*}
P(n)\simeq D(Q^*(n)),
\qquad\text{i.e.}\qquad Q^*(n)=D(P(n)),
\end{equation*}
and $D$ again denotes Cartier duality. The $Q^*(n)$ form a pro-$\ell$ group
in the evident way, denoted by $Q^*$. We leave it to the reader to verify
that, as $n$ varies, the homomorphisms (\SGARefLink{sga7.sga7i.eix.n9-5-4}{9.5.4}) define a homomorphism of
projective systems, that is, a homomorphism of twisted constant
$\ell$-adic sheaves over $S$:

\begin{equation}\tag{\SGARefTarget{sga7.sga7i.eix.n9-5-5}{9.5.5}}
c(X_\eta):Q\otimes Q^*\longrightarrow\mathbb Z_{\ell S}.
\end{equation}

We mention in passing the following result, which will not be used below.

\medskip
\noindent
\underline{Proposition \SGARefTarget{sga7.sga7i.eix.n9-5-6}{9.5.6}. Let $P,Q,R,E,F$ be as above, except that the
trait $S$ need not be henselian. Consider the functor}

\begin{equation}\tag{\SGARefTarget{sga7.sga7i.eix.n9-5-6-1}{9.5.6.1}}
X\longmapsto X_\eta
\end{equation}
\underline{from the category of Barsotti--Tate mixed extensions of $E$ by
$F$ to the category of mixed extensions of $E_\eta$ by $F_\eta$. For every
Barsotti--Tate group $X_\eta$ which is a mixed extension of $E_\eta$ by
$F_\eta$, consider the corresponding pairing}

\begin{equation}\tag{\SGARefTarget{sga7.sga7i.eix.n9-5-6-2}{9.5.6.2}}
Q_o\otimes Q_o^*\longrightarrow\mathbb Z_{\ell s}
\end{equation}
\underline{induced by the pairing of type (\SGARefLink{sga7.sga7i.eix.n9-5-5}{9.5.5}) defined over the
henselization of $S$. Then the functor (\SGARefLink{sga7.sga7i.eix.n9-5-6-1}{9.5.6.1}) is fully faithful, and its
essential image consists of those Barsotti--Tate groups $X_\eta$ which are
mixed extensions of $E_\eta$ by $F_\eta$ and for which the corresponding
pairing (\SGARefLink{sga7.sga7i.eix.n9-5-6-2}{9.5.6.2}) is zero. Moreover, when $\operatorname{char}K=0$ (so that
Tate's theorem [33, Theorem 4] applies), $X_\eta$ belongs to the essential
image if and only if it extends to a Barsotti--Tate group over $S$.}

\medskip
\noindent
\SGARefTarget{sga7.sga7i.eix.n9-5-6-3}{9.5.6.3}. We leave it to the reader to extend this statement to the case in
which $S$ is assumed only to be a connected regular noetherian scheme of
dimension $1$.

The first assertion of \SGARefLink{sga7.sga7i.eix.n9-5-6}{9.5.6} follows immediately from \SGARefLink{sga7.sga7i.eix.n9-4-7}{9.4.7}, applied to the
components of the Barsotti--Tate groups under consideration.
\footnote{The printed proof cites \SGARefLink{sga7.sga7i.eix.n9-4-6}{9.4.6}, which only performs the henselian
descent construction; the full-faithfulness and essential-image statement
being applied componentwise is Proposition \SGARefLink{sga7.sga7i.eix.n9-4-7}{9.4.7}.}
It remains to prove that if a Barsotti--Tate group $X_\eta$, which is a
mixed extension of $E_\eta$ by $F_\eta$, extends to a Barsotti--Tate group
$X$ over $S$, then it also extends as a mixed extension.
\footnote{At this point the printed source gives the extension over $S$ the
name $X_\eta$ as well. We reserve $X_\eta$ for the generic fibre and write
$X$ for its extension over $S$.}
By Tate's theorem [33], the homomorphism
$X_\eta\longrightarrow Q_\eta$ extends uniquely to a homomorphism

\begin{equation*}
X\longrightarrow Q.
\end{equation*}
For every integer $n$, the corresponding homomorphism
$X(n)\longrightarrow Q(n)$ is surjective on the generic fibres and hence is
surjective. Since $Q(n)$ is étale over $S$, the map
$X(n)\longrightarrow Q(n)$ is flat, and consequently it is an epimorphism.
Thus $X\longrightarrow Q$ is an \underline{epimorphism component by component}.
\footnote{The printed sentence says that $Q(n)$ is étale over $\eta$.
The group $Q(n)$ was defined over $S$ and is assumed of étale type there;
the base required for this flatness argument is $S$.}
Let $\overline F$ be the kernel of $X\longrightarrow Q$. Then
$\overline F$ is a Barsotti--Tate group and $X$ is an extension of $Q$ by
$\overline F$. The isomorphism
$F_\eta\xrightarrow{\sim}\overline F_\eta$ extends, by Tate's theorem, to
an isomorphism $F\simeq\overline F$. Hence $X$ is an extension of $Q$ by
$F$. Consider now the quotient $\overline E=X/P$. A further application of
Tate's theorem gives an isomorphism $E\simeq\overline E$. It follows that
$X$ is a mixed extension of $E$ by $F$ extending the mixed extension
$X_\eta$, as required.

\medskip
\noindent
\underline{Remarks \SGARefTarget{sga7.sga7i.eix.n9-5-6-4}{9.5.6.4}.}

\smallskip
\noindent
a) Of course, the last assertion of \SGARefLink{sga7.sga7i.eix.n9-5-6}{9.5.6} should hold without restriction
on $K$, since the same should be true of the invoked theorem of Tate.

\smallskip
\noindent
b) One can give another proof of the good-reduction criterion \SGARefLink{sga7.sga7i.eix.n5-10}{5.10}. As in
\SGARefLink{sga7.sga7i.eix.n5-10}{5.10}, reduce to the case where $A_K$ has semistable reduction, and then
apply \SGARefLink{sga7.sga7i.eix.n9-5-6}{9.5.6} to $T_\ell(A_K)$ regarded as a mixed extension (cf. \SGARefLink{sga7.sga7i.eix.n9-6}{9.6}). This
implies $u_\ell=0$; using \SGARefLink{sga7.sga7i.eix.n11-6-2}{11.6.2}(a) below then gives
$T_\ell(A_K)^t=0$.

\medskip
\noindent
\SGARefTarget{sga7.sga7i.eix.n9-5-8}{9.5.8}. When $\ell\neq p$, the Barsotti--Tate groups under consideration
over $S$ and over $K$, respectively, may simply be interpreted in terms of
twisted constant $\ell$-adic sheaves over $S$, or equivalently in terms of
representations of
$\pi_o=\operatorname{Gal}(\overline k/k)$ and
$\pi=\operatorname{Gal}(\overline K/K)$, respectively, on free
$\mathbb Z_\ell$-modules of finite type. One then sees that the homomorphism
(\SGARefLink{sga7.sga7i.eix.n9-5-5}{9.5.5}) is the canonical homomorphism expressing the action of the inertia
subgroup $I\subset\pi$ on $X_\eta(\overline K)$, as defined by the method
of \SGARefLink{sga7.sga7i.eix.n9-2}{9.2}.
\footnote{The printed source refers here to ``the homomorphism (\SGARefLink{sga7.sga7i.eix.n9-5-6}{9.5.6})'';
the homomorphism constructed in this section is (\SGARefLink{sga7.sga7i.eix.n9-5-5}{9.5.5}), whereas \SGARefLink{sga7.sga7i.eix.n9-5-6}{9.5.6} is
the proposition.}
To verify this, one reduces immediately to the situation of \SGARefLink{sga7.sga7i.eix.n9-4}{9.4}, taking
$n$ there to be a power of $\ell$ and $R$ to be a locally constant
$\Lambda$-module (as are $Q$ and $P$); the same is then true of $E,F$, and
$X_\eta$. Verification that the two definitions agree is left to the
reader.

% END INCLUDED SGA7I_English_current_buildable_source_20260823/components/161_expose_IX_section_9_5.tex

% BEGIN INCLUDED SGA7I_English_current_buildable_source_20260823/components/162_expose_IX_sections_9_6_and_9_7.tex
% Expose IX, sections 9.6--9.7.
% French transcription lines 2655--2711; authority scan physical pp. 448--450
% (scan indices 447--449; source folios 436--438). The authority pages were
% checked directly in four overlapping 1100-dpi bands. The scan is clear.
% Four evident print slips affecting the English are corrected and disclosed:
% F_eta for E_eta in the quotient extension, (9.5.6) for (9.5.5), S for eta
% as the base of X_eta, and a missing closing parenthesis. A missing sentence
% stop in 9.7 is restored.

\medskip
\noindent
\SGARefTarget{sga7.sga7i.eix.n9-6}{9.6}. We are now in a position to define the announced canonical pairing
(\SGARefLink{sga7.sga7i.eix.n9-1-2}{9.1.2}), without any restriction on $\ell$, in such a way that, for
$\ell\neq p$, this definition agrees with the one given in \SGARefLink{sga7.sga7i.eix.n9-2}{9.2}. For this
purpose, we apply the considerations of \SGARefLink{sga7.sga7i.eix.n9-5}{9.5} to the following situation:

\begin{equation}\tag{\SGARefTarget{sga7.sga7i.eix.n9-6-1}{9.6.1}}
\left\{
\begin{aligned}
P&=T_\ell(A^o)^t=\check{\underline M}{}'_\ell(1),
& Q&=\underline M_\ell,
& R&=T_\ell(A^o)^f/T_\ell(A^o)^t,\\[1.25ex]
F&=T_\ell(A^o)^f,
& E&=D(F')=D\bigl(T_\ell({A'}^o)^f\bigr).
\end{aligned}
\right.
\end{equation}
In the last formula, $D$ denotes Cartier duality in the sense of
pro-$\ell$ Barsotti--Tate groups; that is, $D(Y)$ is the
\underline{projective} system of the $D(Y(n))$, with the usual transition
morphisms. The structure of $F$ as an extension of $R$ by $P$ is clear.
Likewise,

\begin{equation*}
F'=T_\ell({A'}^o)^f
\end{equation*}
has the structure of an extension of $R'$ by
$P'=T_\ell({A'}^o)^t=\check{\underline M}_\ell(1)$. Applying Cartier
duality, we conclude that $E=D(F')$ is endowed with the structure of an
extension of $D(P')=\underline M_\ell$ by $D(R')$. Now $D(R')$ is
canonically isomorphic to $R$ by \SGARefLink{sga7.sga7i.eix.n5-5}{5.5} (proved in \SGARefLink{sga7.sga7i.eix.n7-4-6}{7.4.6}); this specifies the
structure of $E$ as an extension of $Q$ by $R$.

As the mixed extension $X_\eta$ of $E_\eta$ by $F_\eta$, we take

\begin{equation}\tag{\SGARefTarget{sga7.sga7i.eix.n9-6-2}{9.6.2}}
X_\eta=T_\ell(A_K)=X^0_\eta
 \supset X^1_\eta=T_\ell(A_K)^f
 \supset X^2_\eta=T_\ell(A_K)^t
 \supset0.
\end{equation}
The isomorphism of extensions of $R_\eta$ by $P_\eta$

\begin{equation*}
X^1_\eta\simeq F_\eta
\end{equation*}
is clear. The isomorphism of extensions of $Q_\eta$ by $R_\eta$

\begin{equation*}
X_\eta/X^2_\eta\simeq E_\eta=D(F'_\eta)
\end{equation*}
is clear from the orthogonality theorem \SGARefLink{sga7.sga7i.eix.n5-2}{5.2}.
\footnote{The printed formula has $F_\eta=D(F'_\eta)$. Since
$E=D(F')$ and the displayed quotient is the extension of $Q_\eta$ by
$R_\eta$, the required term is $E_\eta$.}

We therefore have the data required to define a pairing (\SGARefLink{sga7.sga7i.eix.n9-5-5}{9.5.5}), which
here, since $Q^*=\underline M'_\ell$, takes the form (\SGARefLink{sga7.sga7i.eix.n9-1-2}{9.1.2}).
\footnote{The printed source calls this a pairing ``(\SGARefLink{sga7.sga7i.eix.n9-5-6}{9.5.6})''; the pairing
is the homomorphism (\SGARefLink{sga7.sga7i.eix.n9-5-5}{9.5.5}), whereas \SGARefLink{sga7.sga7i.eix.n9-5-6}{9.5.6} is the subsequent proposition.}
Its compatibility with the definition of \SGARefLink{sga7.sga7i.eix.n9-2}{9.2} follows from \SGARefLink{sga7.sga7i.eix.n9-5-8}{9.5.8}. The
considerations of \SGARefLink{sga7.sga7i.eix.n5-8}{5.8} also show, subject to the validity of Tate's theorem
[33] (hence at least when $\operatorname{char}K=0$), that the pairing in
question can be described in terms of the single pro-$\ell$
Barsotti--Tate group $X_\eta=T_\ell(A_K)$ over $\eta$: its three-step
filtration, and the extensions $F$ and $E$ extending $X^1_\eta$ and
$X_\eta/X^2_\eta$, can be defined canonically in terms of $X_\eta$ alone.
\footnote{The printed sentence says that $X_\eta$ is ``over $S$''; its
notation and definition make it the generic-fibre group over $\eta$. It
also lacks its closing parenthesis.}

\medskip
\noindent
\underline{Remark \SGARefTarget{sga7.sga7i.eix.n9-7}{9.7}.} We no longer necessarily assume that $A_K$ has
semistable reduction over the henselian trait $S$. Let $K'$ be a finite
Galois extension of $K$ such that $A_{K'}$ has semistable reduction over
the normalization $S'$ of $S$ in $K'$ (\SGARefLink{sga7.sga7i.eix.n3-6}{3.6}). With the notation of (\SGARefLink{sga7.sga7i.eix.n8-1-8}{8.1.8}),
Galois descent from the monodromy pairing for $A_{K'}$ then defines a
pairing

\begin{equation}\tag{*}
v_\ell^{K'}:\underline M_K\otimes\underline M'_K
 \longrightarrow\mathbb Z_{\ell K}.
\end{equation}
A priori this pairing depends on the choice of $K'$, but it follows from
(10.4.3) below that the pairing

\begin{equation}\tag{\SGARefTarget{sga7.sga7i.eix.n9-7-1}{9.7.1}}
u_\ell:\underline M_K\otimes\underline M'_K
 \longrightarrow\mathbb Q_{\ell K}
\end{equation}
given by

\begin{equation*}
u_\ell=\frac{1}{e(S'/S)}\,v_\ell^{K'}
\end{equation*}
is independent of it, where

\begin{equation*}
e(S'/S)=\operatorname{length}(V'/\mathfrak mV').
\end{equation*}
We may again call (\SGARefLink{sga7.sga7i.eix.n9-7-1}{9.7.1}) the \underline{monodromy pairing} for the
abelian scheme $A_K$ over $K$ and the prime number $\ell$ under
consideration: when $A_K$ itself has semistable reduction, it plainly
coincides with the restriction to $\eta$ of the monodromy pairing (\SGARefLink{sga7.sga7i.eix.n9-1-2}{9.1.2}).
It follows from \SGARefLink{sga7.sga7i.eix.n10-4}{10.4} below that this pairing in fact always has values in
$\mathbb Q_K$ and is independent of $\ell$.

One may ask whether it has values in $\mathbb Z_K$, which is equivalent to
saying that (\SGARefLink{sga7.sga7i.eix.n9-7-1}{9.7.1}) has values in $\mathbb Z_{\ell K}$ for every prime
number $\ell$. This is so at least when $A_K$ is an elliptic curve. The
only case requiring verification is then the one in which
$\underline M_K\simeq\underline M'_K$ has rank $1$, so that
$\underline M_K\otimes\underline M'_K$ is canonically isomorphic to
$\mathbb Z_K$ and the monodromy pairing (\SGARefLink{sga7.sga7i.eix.n9-7-1}{9.7.1}) can simply be interpreted
as an ordinary rational number. It follows from the results of Néron [19]
that this number is $v(j)$, where $j$ is the absolute invariant of $A_K$
and $v$ is the normalized valuation of $K$.


% END INCLUDED SGA7I_English_current_buildable_source_20260823/components/162_expose_IX_sections_9_6_and_9_7.tex

% BEGIN INCLUDED SGA7I_English_current_buildable_source_20260823/components/163_expose_IX_sections_10_0_through_10_2.tex
% Expose IX, section 10 opening and sections 10.1--10.2.
% French transcription lines 2717--2786; authority scan physical pp. 451--453
% (scan indices 450--452; source folios 439--441). Physical p. 454 was
% checked for the exact 10.3 boundary. All pages were compared in four
% overlapping 1100-dpi bands; the print is clear at that resolution.

\medskip
\noindent
10. \underline{Properties of the monodromy pairing}. \underline{Integrality
and positivity theorem}

In this section we shall first give certain essentially trivial formal
properties of the monodromy pairings, leaving the (somewhat tedious)
details of the verifications to the reader. We shall also state a deeper
theorem (\SGARefLink{sga7.sga7i.eix.n10-4}{10.4}) concerning these pairings and indicate how its proof may be
reduced to a special case of Jacobians, to be treated (by transcendental
methods) in section 12, which is in fact logically independent of section
11.

\medskip
\noindent
\SGARefTarget{sga7.sga7i.eix.n10-1}{10.1}. \underline{Transport of structure}. It is clear from the definition
of the monodromy pairing (\SGARefLink{sga7.sga7i.eix.n9-1-2}{9.1.2}) that this pairing is compatible with
transport of structure in the situation $(S,A_K,A'_K)$.

\medskip
\noindent
\SGARefTarget{sga7.sga7i.eix.n10-2}{10.2}. \underline{Symmetry}. Starting with the dual scheme $A'_K$ of $A_K$
and applying the construction of section 9 to it, we must introduce the
dual scheme $(A'_K)'$ of $A'_K$, which is canonically isomorphic to $A_K$
by the biduality theorem ([5] or VIII 3.2). Denote by
$u_\ell^{A_K}$ the monodromy pairing (\SGARefLink{sga7.sga7i.eix.n9-1-2}{9.1.2}) relative to $A_K$. In the same
way we obtain a pairing

\begin{equation}\tag{\SGARefTarget{sga7.sga7i.eix.n10-2-1}{10.2.1}}
u_\ell^{A'_K}:\underline M'_\ell\otimes\underline M_\ell
 \longrightarrow\mathbb Z_\ell.
\end{equation}
I claim that the latter is obtained from $u_\ell^{A_K}$ by the symmetry

\begin{equation*}
s:\underline M'_\ell\otimes\underline M_\ell
 \xrightarrow{\ \sim\ }
 \underline M_\ell\otimes\underline M'_\ell;
\end{equation*}
that is,

\begin{equation}\tag{\SGARefTarget{sga7.sga7i.eix.n10-2-2}{10.2.2}}
u_\ell^{A'_K}=u_\ell^{A_K}\circ s.
\end{equation}

Taking account of the known anticommutativity property of the pairings
$\varphi_\xi$ ([5] or VIII 2.2.11), the symmetry formula (\SGARefLink{sga7.sga7i.eix.n10-2-2}{10.2.2}) may be
reformulated as follows in terms of a fixed polarization $\xi$ of $A_K$.
It defines a homomorphism

\begin{equation*}
\widetilde\xi:A_K\longrightarrow A'_K
\end{equation*}
and hence a corresponding homomorphism ($t$ standing for ``toric'')

\begin{equation}\tag{\SGARefTarget{sga7.sga7i.eix.n10-2-3}{10.2.3}}
\xi_t:\underline M\longrightarrow\underline M'.
\end{equation}
We may therefore form the pairing

\begin{equation}\tag{\SGARefTarget{sga7.sga7i.eix.n10-2-4}{10.2.4}}
\begin{aligned}
\varphi^t_{\xi,\ell}:\underline M_\ell\otimes\underline M_\ell
 &\longrightarrow\mathbb Z_\ell,\\
\varphi^t_{\xi,\ell}(x,y)&=u_\ell^{A_K}\bigl(x,\xi_t(y)\bigr).
\end{aligned}
\end{equation}
The preceding form $\varphi^t_{\xi,\ell}$ is \underline{symmetric}:

\begin{equation}\tag{\SGARefTarget{sga7.sga7i.eix.n10-2-5}{10.2.5}}
\varphi^t_{\xi,\ell}(x,y)=\varphi^t_{\xi,\ell}(y,x).
\end{equation}
In fact, this assertion is valid for every element $\xi$ of the
Néron--Severi group of $A_K$, not necessarily a polarization. But only
when $\xi$ is nondegenerate---that is, defines an \underline{isogeny}
$A_K\longrightarrow A'_K$, and hence also an isogeny (\SGARefLink{sga7.sga7i.eix.n10-2-3}{10.2.3})---is the
symmetry relation (\SGARefLink{sga7.sga7i.eix.n10-2-5}{10.2.5}) equivalent to (\SGARefLink{sga7.sga7i.eix.n10-2-2}{10.2.2}). When $\xi$ is a
polarization, we shall moreover sharpen (\SGARefLink{sga7.sga7i.eix.n10-2-5}{10.2.5}) considerably in \SGARefLink{sga7.sga7i.eix.n10-4}{10.4}(b)
below.

Let us finally mention a third interpretation of the symmetry relation.
Consider a biextension (that is, a divisorial correspondence) $\xi$ on a
pair $(A_K,\overline A_K)$ of abelian schemes with semistable reduction
over $K$. It defines a homomorphism

\begin{equation*}
A_K\longrightarrow\overline A'_K.
\end{equation*}
Passing to the Néron models, this defines a homomorphism on the maximal
tori of the special fibres, and therefore, in the reverse direction, a
homomorphism on their character groups:

\begin{equation}\tag{\SGARefTarget{sga7.sga7i.eix.n10-2-6}{10.2.6}}
\xi_t:\overline{\underline M}\longrightarrow\underline M',
\end{equation}
where $\overline{\underline M}$ is defined from $\overline A_K$ as
$\underline M$ is defined from $A_K$. We can therefore form the pairing

\begin{equation}\tag{\SGARefTarget{sga7.sga7i.eix.n10-2-7}{10.2.7}}
\begin{aligned}
\varphi^t_\xi:\underline M_\ell\otimes
 \overline{\underline M}_\ell&\longrightarrow\mathbb Z_\ell,\\
\varphi^t_\xi(x,y)&=u_\ell^{A_K}\bigl(x,\xi_t(y)\bigr).
\end{aligned}
\end{equation}
\footnote{In the printed right-hand side the underlying pairing is denoted
only by $\varphi^t$, without a definition. It is the monodromy pairing
$u_\ell^{A_K}$ used in the identical construction (\SGARefLink{sga7.sga7i.eix.n10-2-4}{10.2.4}).}
Notice in passing that, for a variable abelian scheme $A_K$ with semistable
reduction, the twisted constant group $\underline M$ is a covariant
functor of $A_K$. Formation of the pairings (\SGARefLink{sga7.sga7i.eix.n10-2-7}{10.2.7}) associated with
biextensions of abelian schemes with semistable reduction over $(K,S)$ is
then compatible, in the evident sense, with taking inverse images of
biextensions.

When $\xi$ is the canonical (``Weil'') biextension of $(A_K,A'_K)$, the
pairing (\SGARefLink{sga7.sga7i.eix.n10-2-7}{10.2.7}) is of course precisely the monodromy pairing (\SGARefLink{sga7.sga7i.eix.n9-1-2}{9.1.2}).
\footnote{The printed source refers here to (\SGARefLink{sga7.sga7i.eix.n9-2-1}{9.2.1}), which is the inertia
orthogonality formula $gx-x\in V^\perp$, not the monodromy pairing. The
pairing is numbered (\SGARefLink{sga7.sga7i.eix.n9-1-2}{9.1.2}).}
Now consider the transpose ${}^s\xi$ of $\xi$, a biextension of
$(\overline A_K,A_K)$, and hence the pairing

\begin{equation*}
\varphi^t_{{}^s\xi}:\overline{\underline M}_\ell\otimes
 \underline M_\ell\longrightarrow\mathbb Z_\ell.
\end{equation*}
The symmetry relation is

\begin{equation}\tag{\SGARefTarget{sga7.sga7i.eix.n10-2-8}{10.2.8}}
\varphi^t_{{}^s\xi}=\varphi^t_\xi\circ s,
\qquad
s:\overline{\underline M}_\ell\otimes\underline M_\ell
 \xrightarrow{\ \sim\ }
 \underline M_\ell\otimes\overline{\underline M}_\ell,
\end{equation}
where $s$ is again the symmetry isomorphism. Compare this relation with
the antisymmetry relation of VIII \SGARefLink{sga7.sga7i.eviii.n2-2-11}{2.2.11}.


% END INCLUDED SGA7I_English_current_buildable_source_20260823/components/163_expose_IX_sections_10_0_through_10_2.tex

% BEGIN INCLUDED SGA7I_English_current_buildable_source_20260823/components/164_expose_IX_sections_10_3_and_10_4.tex
% Expose IX, sections 10.3--10.4.
% French transcription lines 2792--2883; authority scan physical pp. 454--456
% (scan indices 453--455; source folios 442--444). Physical p. 457 was
% checked for the exact 10.5 boundary. All pages were compared in four
% overlapping 1100-dpi bands; the text, formulas, and simple square are
% clear at that resolution.

\medskip
\noindent
\SGARefTarget{sga7.sga7i.eix.n10-3}{10.3}. \underline{Change of trait}. Let

\begin{equation}\tag{\SGARefTarget{sga7.sga7i.eix.n10-3-1}{10.3.1}}
g:S'\longrightarrow S
\end{equation}
be a dominant morphism of henselian traits, associated with an injective
local homomorphism of discrete valuation rings

\begin{equation*}
V\longrightarrow V'.
\end{equation*}
As usual, put

\begin{equation}\tag{\SGARefTarget{sga7.sga7i.eix.n10-3-2}{10.3.2}}
e(S'/S)=\operatorname{length}_V\bigl(V'/\underline mV'\bigr),
\end{equation}
where $\underline m$ is the maximal ideal of $V$. It is then known that
the homomorphism induced by (\SGARefLink{sga7.sga7i.eix.n10-3-1}{10.3.1}) on the ``tame'' parts of the inertia
groups,

\begin{equation}\tag{\SGARefTarget{sga7.sga7i.eix.n10-3-3}{10.3.3}}
g_*:I'_t\simeq\prod_{\ell\ne p}\mathbb Z_\ell(1)(k')
 \longrightarrow
 I_t=\prod_{\ell\ne p}\mathbb Z_\ell(1)(k),
\end{equation}
is given by

\begin{equation}\tag{\SGARefTarget{sga7.sga7i.eix.n10-3-4}{10.3.4}}
g_*(x)=e(S'/S)x,
\end{equation}
after the two sides of (\SGARefLink{sga7.sga7i.eix.n10-3-3}{10.3.3}) have been identified by means of the
evident isomorphisms
$\mathbb Z_\ell(1)(k)\xrightarrow{\ \sim\ }\mathbb Z_\ell(1)(k')$.
\footnote{The printed source leaves the internal reference following
``Expos\'e I'' blank.}

On the other hand, if $A_K$ is an abelian scheme over $K$ with
semistable reduction over $S$, formation of the twisted lattices
$\underline M$ and $\underline M'$ commutes with every base change
$S'\longrightarrow S$ as above, as follows immediately from the
definitions (especially \SGARefLink{sga7.sga7i.eix.n3-4}{3.4}). We may therefore regard the pairing
$u_\ell^{A_{K'}}$ defined by $A_{K'}$, where $K'$ is the fraction field
of $V'$, as a pairing

\begin{equation*}
u_\ell^{A_{K'}}:
 \underline M_{\ell S'}\otimes\underline M'_{\ell S'}
 \longrightarrow\mathbb Z_{\ell S'},
\end{equation*}
and compare it with the pairing $u_{\ell S'}^{A_K}$ obtained from
$u_\ell^{A_K}$ by base change. When $\ell\ne p$, the construction of the
pairings $u_\ell$ given in \SGARefLink{sga7.sga7i.eix.n9-2}{9.2}, together with formula (\SGARefLink{sga7.sga7i.eix.n10-3-4}{10.3.4}), shows at
once that\footnote{The printed source cites (\SGARefLink{sga7.sga7i.eix.n10-3-3}{10.3.3}), the displayed map
$g_*$; the multiplication formula it uses is (\SGARefLink{sga7.sga7i.eix.n10-3-4}{10.3.4}).}

\begin{equation}\tag{\SGARefTarget{sga7.sga7i.eix.n10-3-5}{10.3.5}}
u_\ell^{A_{K'}}=e(S'/S)u_\ell^{A_K}.
\end{equation}
In fact this formula remains valid for $\ell=p$. To see this, return to
the general construction of \SGARefLink{sga7.sga7i.eix.n9-3}{9.3}--\SGARefLink{sga7.sga7i.eix.n9-6}{9.6} and reduce to the analogous
compatibility assertion for the isomorphism \SGARefLink{sga7.sga7i.eix.n9-4-4}{9.4.4}(e), under a change of
strictly local traits $S'\longrightarrow S$. Notice that the isomorphism
in question is induced by the normalized valuation
$K^*/V^*\simeq\mathbb Z$, and therefore gives the commutative diagram
\footnote{The printed source says that $K^*/V^*\simeq\mathbb Z$ is
deduced from the ``norm.''  The displayed map is the normalized discrete
valuation.}

\begin{center}
\begin{tikzcd}
K^*/V^* \arrow[r, "\sim"] \arrow[d]
  & \mathbb Z \arrow[d, "e(S'/S)"] \\
K^{\prime *}/V^{\prime *} \arrow[r, "\sim"]
  & \mathbb Z .
\end{tikzcd}
\end{center}

Here is a deeper result concerning the monodromy pairings. We shall prove
it by transcendental methods in section 12, relying on results that will
be proved later in the Seminar. (We shall not, in fact, use \SGARefLink{sga7.sga7i.eix.n10-4}{10.4} in the
remainder of the Seminar, any more than we shall use the other results of
the present expos\'e.)

\medskip
\noindent
\underline{Theorem \SGARefTarget{sga7.sga7i.eix.n10-4}{10.4}.} \underline{Let $S$ be a henselian trait, and
let $A_K$ be an abelian scheme over $K$ with semistable reduction over
$S$. Consider the lattices $\underline M$ and $\underline M'$ associated
as in \SGARefLink{sga7.sga7i.eix.n9-1}{9.1} (cf. \SGARefLink{sga7.sga7i.eix.n5-5-2}{5.5.2}). Then the following assertions hold:}

\smallskip
\noindent
a) \underline{There exists a unique pairing}

\begin{equation}\tag{\SGARefTarget{sga7.sga7i.eix.n10-4-1}{10.4.1}}
u:\underline M\otimes\underline M'\longrightarrow\mathbb Z_S
\end{equation}
\underline{such that, for every prime number $\ell$,
$u\otimes_{\mathbb Z}\mathbb Z_\ell$ is the monodromy pairing $u_\ell$
of (\SGARefLink{sga7.sga7i.eix.n9-1-2}{9.1.2}). This pairing is separating.}

\smallskip
\noindent
b) \underline{Let $\xi$ be a polarization of $A_K$, thus defining a
homomorphism (\SGARefLink{sga7.sga7i.eix.n10-2-3}{10.2.3})}

\begin{equation*}
\xi_t:\underline M\longrightarrow\underline M',
\end{equation*}
\underline{and consider the form}

\begin{equation}\tag{\SGARefTarget{sga7.sga7i.eix.n10-4-2}{10.4.2}}
\begin{aligned}
\varphi_\xi^t:\underline M\otimes\underline M
  &\longrightarrow\mathbb Z_S,\\
\varphi_\xi^t(x,y)&=u\bigl(x,\xi_t(y)\bigr).
\end{aligned}
\end{equation}
\underline{This form is symmetric, positive, and nondegenerate.}

Of course, in terms of the geometric fibre $M_o$ of $\underline M$ at a
given geometric point of $S$, the preceding assertion says that the
corresponding form

\begin{equation*}
M_o\otimes M_o\longrightarrow\mathbb Z
\end{equation*}
is symmetric, positive, and nondegenerate. The pairing (\SGARefLink{sga7.sga7i.eix.n10-4-1}{10.4.1}) will
again be called the \underline{monodromy pairing}.

% END INCLUDED SGA7I_English_current_buildable_source_20260823/components/164_expose_IX_sections_10_3_and_10_4.tex

% BEGIN INCLUDED SGA7I_English_current_buildable_source_20260823/components/165_expose_IX_section_10_5.tex
% Expose IX, section 10.5.
% French transcription lines 2889--2955; authority scan physical pp. 457--460
% (scan indices 456--459; source folios 445--448). Physical p. 460 was
% checked through the exact section 11 boundary. All pages were compared in
% four overlapping 1100-dpi bands; the prose and formulas are clear at that
% resolution, so no higher-resolution escalation was necessary.

\medskip
\noindent
\SGARefTarget{sga7.sga7i.eix.n10-5}{10.5}. \underline{Reduction of \SGARefLink{sga7.sga7i.eix.n10-4}{10.4} to a particular case}

\smallskip
\noindent
\SGARefTarget{sga7.sga7i.eix.n10-5-1}{10.5.1}. The uniqueness of a $u$ satisfying the stated condition (even
for a single $\ell$) is clear. On the other hand, existence simply means
that every $u_\ell$ in fact maps $\underline M\otimes\underline M'$ into
$\mathbb Z_S$, and that the homomorphism thereby obtained

\begin{equation*}
\underline M\otimes\underline M'\longrightarrow\mathbb Z_S
\end{equation*}
is independent of $\ell$. In fact it even suffices to prove the analogous
assertion for the $u_\ell\otimes_{\mathbb Z_\ell}\mathbb Q_\ell$ and for
a homomorphism

\begin{equation*}
u_{\mathbb Q}:\underline M_{\mathbb Q}\otimes
 \underline M'_{\mathbb Q}\longrightarrow\mathbb Q_S.
\end{equation*}

Suppose, then, that existence has been established after a base change
$S'\longrightarrow S$ as in \SGARefLink{sga7.sga7i.eix.n9-3}{9.3}. It follows immediately from (\SGARefLink{sga7.sga7i.eix.n10-3-5}{10.3.5})
that there exists a homomorphism

\begin{equation*}
u:\underline M\otimes\underline M'
 \longrightarrow\frac{1}{e(S'/S)}\mathbb Z_S\subset\mathbb Q_S
\end{equation*}
compatible with the $u_\ell\otimes_{\mathbb Z_\ell}\mathbb Q_\ell$;
this proves a). As for b), it is then plain that, if it holds for
$A_{K'}$, it also holds for $A_K$.

Consequently, in proving \SGARefLink{sga7.sga7i.eix.n10-4}{10.4} we may make any dominant change of traits
$S'\longrightarrow S$. In particular, this allows us to suppose that
$S$ is complete with algebraically closed residue field. In that case
$\underline M$ and $\underline M'$ are constant group schemes, with
values $M$ and $M'$, finitely generated free $\mathbb Z$-modules, and a)
and b) may be reformulated more conveniently in terms of the latter.

\smallskip
\noindent
\SGARefTarget{sga7.sga7i.eix.n10-5-2}{10.5.2}. Suppose now that $A_K$ is isomorphic to a direct factor of an
abelian scheme $\overline A_K$ over $K$ with semistable reduction, and
that \SGARefLink{sga7.sga7i.eix.n10-4}{10.4} has been proved for the latter. It follows at once that \SGARefLink{sga7.sga7i.eix.n10-4}{10.4}
also holds for $A_K$, in view of the evident functorial behaviour of the
lattices $\underline M,\underline M'$ and of the monodromy pairings with
respect to direct sums of abelian schemes with semistable reduction.
\footnote{The printed source says that $\overline A_K$ is ``over $k$'';
the notation and the direct-factor argument require ``over $K$.''}

It is equally immediate from \SGARefLink{sga7.sga7i.eix.n10-5-1}{10.5.1} that the validity of \SGARefLink{sga7.sga7i.eix.n10-4}{10.4} depends
only on the isogeny class of $A_K$. Together with the preceding remark
and the ``Poincar\'e reducibility theorem'' [15], this shows that, if
$A_K$ is a quotient of an abelian scheme $\overline A_K$ for which \SGARefLink{sga7.sga7i.eix.n10-4}{10.4}
holds, then \SGARefLink{sga7.sga7i.eix.n10-4}{10.4} also holds for $A_K$.
\footnote{The printed source leaves the pinpoint following citation [15]
blank.}

\smallskip
\noindent
\SGARefTarget{sga7.sga7i.eix.n10-5-3}{10.5.3}. We now use the well-known fact [31, Th.~11 (p.~20)] that, after
making a finite extension of $K$ if necessary (which is permissible by
\SGARefLink{sga7.sga7i.eix.n10-5-1}{10.5.1}), every abelian scheme $A_K$ is a quotient of the Jacobian

\begin{equation*}
\underline{\operatorname{Pic}}^0_{X_K/K}
\end{equation*}
of a proper, smooth, geometrically connected curve $X_K$ over $K$.
Thus, by \SGARefLink{sga7.sga7i.eix.n10-5-2}{10.5.2}, we are reduced to proving \SGARefLink{sga7.sga7i.eix.n10-4}{10.4} for such a Jacobian.
\footnote{The printed source says that the curve $X_K$ is ``over $X$'';
the displayed relative Picard scheme and the argument require ``over
$K$.''}

\smallskip
\noindent
\SGARefTarget{sga7.sga7i.eix.n10-5-4}{10.5.4}. By a theorem of Artin--Mumford [3] (already cited and used in
Expos\'e V), after making a further finite extension $K'$ of $K$ if
necessary, we may suppose that $X_K$ extends to a projective and flat
scheme over $S$ which is regular and whose geometric special fibre
$X_{\bar s}$ has no singularities other than ordinary quadratic
singularities (that is, here, normal crossings of two branches).
\footnote{The printed source leaves the internal reference following
``Expos\'e V'' blank.}
In this case we shall prove \SGARefLink{sga7.sga7i.eix.n10-4}{10.4} a) in \SGARefLink{sga7.sga7i.eix.n12-5}{12.5} below, by describing the
monodromy pairing (\SGARefLink{sga7.sga7i.eix.n10-4-1}{10.4.1}) completely in terms of the geometric data of
the situation; at the same time we shall verify \SGARefLink{sga7.sga7i.eix.n10-4}{10.4} b) for the canonical
polarization of the Jacobian.

\smallskip
\noindent
\SGARefTarget{sga7.sga7i.eix.n10-5-5}{10.5.5}. It remains to see that, if for a given $A_K$ with semistable
reduction assertion \SGARefLink{sga7.sga7i.eix.n10-4}{10.4} a) holds and assertion \SGARefLink{sga7.sga7i.eix.n10-4}{10.4} b) holds for one
given polarization $\xi_o$ of $A_K$, then the same conclusion holds for
every polarization $\xi$ of $A_K$. Indeed, let

\begin{equation*}
\widetilde\xi_o,\widetilde\xi:A_K\longrightarrow A'_K
\end{equation*}
be the homomorphisms defined by $\xi_o$ and $\xi$; these are isogenies.
Then

\begin{equation*}
\widetilde\xi=\widetilde\xi_o\alpha
 ={}^t(\alpha^*)\widetilde\xi_o,
 \qquad
\alpha\in\operatorname{End}(A_K)\otimes_{\mathbb Z}\mathbb Q,
\end{equation*}
where $\alpha\longmapsto\alpha^*$ denotes the canonical involution of
the algebra $\operatorname{End}(A_K)\otimes_{\mathbb Z}\mathbb Q$
defined by the polarization $\xi_o$, and the superscript $t$ in
${}^t(\alpha^*)$ denotes transpose (associating with an endomorphism of
$A_K$ an endomorphism of $A'_K$). On the other hand, it is known
[15, Chap.~V, Th.~3] that

\begin{equation*}
\alpha=\sum_i\beta_i^*\beta_i,
 \qquad
\beta_i\in\operatorname{End}(A_K)\otimes_{\mathbb Z}\mathbb Q,
\end{equation*}
and hence

\begin{equation*}
\widetilde\xi=\sum_i{}^t\beta_i\,\widetilde\xi_o\,\beta_i.
\end{equation*}
Equivalently, in terms of the biextensions (that is, divisorial
correspondences) associated with the polarizations under consideration
(VIII \SGARefLink{sga7.sga7i.eviii.n4-14}{4.14}),

\begin{equation}\tag{\SGARefTarget{sga7.sga7i.eix.n10-5-5-1}{10.5.5.1}}
\delta(\xi)=\sum_i\delta(\xi_o)\circ(\beta_i,\beta_i).
\end{equation}
By what was said after (\SGARefLink{sga7.sga7i.eix.n10-2-7}{10.2.7}), this formula implies (with the notation
of (\SGARefLink{sga7.sga7i.eix.n10-4-2}{10.4.2}))

\begin{equation}\tag{\SGARefTarget{sga7.sga7i.eix.n10-5-5-2}{10.5.5.2}}
\varphi^t_\xi
 =\sum_i\varphi^t_{\xi_o}(\beta_i,\beta_i)
 \in\operatorname{Hom}
 \bigl(M_{\mathbb Q}\otimes_{\mathbb Q}M_{\mathbb Q},\mathbb Q\bigr).
\end{equation}
It indeed follows that, if $\varphi^t_{\xi_o}$ is symmetric and positive,
then so is $\varphi^t_\xi$.

Finally, the nondegeneracy of $\varphi^t_{\xi_o}$ is plainly equivalent
to the assertion that the pairing (\SGARefLink{sga7.sga7i.eix.n10-4-1}{10.4.1}) is separating (since
$\xi_{ot}:M\otimes_{\mathbb Z}\mathbb Q
\longrightarrow M'\otimes_{\mathbb Z}\mathbb Q$ is an isomorphism,
$\widetilde\xi_o$ being an isogeny), and is therefore likewise
independent of the chosen polarization. One may also observe that, in
order to verify that (\SGARefLink{sga7.sga7i.eix.n10-4-1}{10.4.1}) is separating, it suffices to verify the
same assertion for the monodromy pairing $u_\ell$ of (\SGARefLink{sga7.sga7i.eix.n9-1-2}{9.1.2}) for one
prime number $\ell$. Taking $\ell\ne p$, one may use the description of
$u_\ell$ in \SGARefLink{sga7.sga7i.eix.n9-2}{9.2}; the fact that $u_\ell$ is separating on the left (and
hence on both sides, for reasons of rank) is then immediate.

% END INCLUDED SGA7I_English_current_buildable_source_20260823/components/165_expose_IX_section_10_5.tex

% BEGIN INCLUDED SGA7I_English_current_buildable_source_20260823/components/166_expose_IX_sections_11_0_through_11_2.tex
% Expose IX, sections 11.0--11.2.
% French transcription lines 2956--3027; authority scan physical pp. 460--463
% (scan indices 459--462; source folios 448--451). Physical p. 463 was
% checked through the exact 11.3 boundary. All pages were compared in four
% overlapping 1100-dpi bands; the text and formulas are clear at that
% resolution, so no higher-resolution escalation was necessary.

\medskip
\noindent
11. \underline{Connected components of the N\'eron model: duality
relation and asymptotic behaviour}

\smallskip
\noindent
\SGARefTarget{sga7.sga7i.eix.n11-0}{11.0}. The purpose of the present section is to study the finite \'etale
$k$-group

\begin{equation}\tag{\SGARefTarget{sga7.sga7i.eix.n11-0-1}{11.0.1}}
\Phi_o=A_o/A_o^{,o}
\end{equation}
already considered in section 1. For every prime number $\ell$, denote
by

\begin{equation*}
\Phi_o(\ell)
\end{equation*}
its $\ell$-primary component. Thus knowledge of $\Phi_o$ is equivalent
to knowledge of the $\Phi_o(\ell)$, by the formula

\begin{equation}\tag{\SGARefTarget{sga7.sga7i.eix.n11-0-2}{11.0.2}}
\Phi_o=\coprod_\ell\Phi_o(\ell).
\end{equation}

Notice that the groups $\Phi_o$ and $\Phi_o(\ell)$ are determined,
respectively, by the ordinary finite groups $\Phi_o(\bar k)$ and
$\Phi_o(\ell)(\bar k)$ together with the natural action of

\begin{equation*}
\pi_o=\operatorname{Gal}(\bar k/k)
\end{equation*}
on them. On the other hand, it is clear that $\Phi_o(\bar k)$ is also
the group of connected components of the N\'eron model of
$A_{\widetilde K}$, where $\widetilde K$ is the fraction field of the
strict henselization of $S$ relative to the chosen separable closure
$\bar k$ of $k$. Thus, in studying $\Phi_o$, we may reduce to the case
where $S$ is strictly local; the actions of

\begin{equation*}
\pi_o\simeq\operatorname{Gal}(\widetilde K/K)
\end{equation*}
on $\Phi_o(\bar k)$ are then recovered automatically by transport of
structure. We could likewise reduce further to the case where $S$ is
complete, using the base change $\widehat S\longrightarrow S$ and the
fact that formation of the N\'eron model commutes with it
[21, Th.~\SGARefLink{sga7.sga7i.eix.n3-4}{3.4} b)].

\smallskip
\noindent
\SGARefTarget{sga7.sga7i.eix.n11-1}{11.1}. Suppose, then, that $S$ is strictly local, so that $\Phi_o$ may be
identified with the ordinary finite group $\Phi_o(k)$. We obtain

\begin{equation}\tag{\SGARefTarget{sga7.sga7i.eix.n11-1-1}{11.1.1}}
\Phi_o(k)\simeq A_o(k)/A_o^o(k)
 \simeq A(S)/A^o(S)
 \simeq A(K)/A(K)^o,
\end{equation}
where the first equality follows from the fact that $A_o$ is smooth over
$k$, the second from Hensel's lemma since $A$ is smooth over $S$ and
$S$ is henselian, and the third from the relation
$A(S)\xrightarrow{\sim}A(K)$ and from the convention

\begin{equation}\tag{\SGARefTarget{sga7.sga7i.eix.n11-1-2}{11.1.2}}
A(K)^o\mathrel{\mathop:}=A^o(S).
\end{equation}
Thus $A(K)^o$ appears as a natural subgroup of finite index in $A(K)$
(and is not to be confused with $A^o(K)=A(K)$). The isomorphisms
(\SGARefLink{sga7.sga7i.eix.n11-1-1}{11.1.1}) are plainly functorial in $A$ and compatible with transport of
structure in $(S,A_K)$, and in particular with the action of any Galois
groups.

Now let $\ell$ be a prime number different from $p$. Then $A^o$ is
$\ell$-divisible; more precisely, $\ell\operatorname{id}_{A^o}$ is an
\'etale surjective morphism. It consequently induces a surjective
endomorphism of $A^o(S)=A(K)^o$, which is therefore an
\underline{$\ell$-divisible group}. It follows at once that, for every
integer $n\geq 0$ prime to $p=\operatorname{char}(k)$,

\begin{equation}\tag{\SGARefTarget{sga7.sga7i.eix.n11-1-3}{11.1.3}}
{}_n\Phi_o(k)\simeq{}_nA(K)/{}_nA(K)^o,
 \qquad (n,p)=1.
\end{equation}
By the definition (\SGARefLink{sga7.sga7i.eix.n2-2-3}{2.2.3}) of the fixed part of ${}_nA_K$, we have

\begin{equation}\tag{\SGARefTarget{sga7.sga7i.eix.n11-1-4}{11.1.4}}
{}_nA(K)^o=({}_nA_K)^f(\bar K)
 \simeq T_n(A_K)^f(\bar K)
 \otimes_{\mathbb Z_n}\mathbb Z/n\mathbb Z.
\end{equation}
Taking $n$ to be a power of $\ell$, and putting as in (\SGARefLink{sga7.sga7i.eix.n2-5-3}{2.5.3})

\begin{equation}\tag{\SGARefTarget{sga7.sga7i.eix.n11-1-5}{11.1.5}}
U=T_\ell\bigl(A(\bar K)\bigr),
 \qquad
T_\ell\bigl(A(\bar K)\bigr)^f=U^I,
\end{equation}
where $U^I$ denotes the submodule of $U$ consisting of invariants under
the inertia group $I$, we may write (\SGARefLink{sga7.sga7i.eix.n11-1-3}{11.1.3}) in the form

\begin{equation}\tag{\SGARefTarget{sga7.sga7i.eix.n11-1-6}{11.1.6}}
{}_n\Phi_o(k)\simeq (U_n)^I/(U^I)_n,
 \qquad n=\ell^{r+1},
\end{equation}
where the subscript $n$ denotes reduction modulo $n$, that is, tensor
product over $\mathbb Z_\ell$ with
$\mathbb Z_\ell/n\mathbb Z_\ell$. The transition morphisms between the
groups occurring on the two sides of (\SGARefLink{sga7.sga7i.eix.n11-1-6}{11.1.6}) being evident, passage to
the inductive limit over the powers of $\ell$ gives the following
result.

\medskip
\noindent
\underline{Proposition \SGARefTarget{sga7.sga7i.eix.n11-2}{11.2}.} \underline{Suppose that $S$ is strictly
local, and let $\ell$ be a prime number different from $p$. Then there is
a canonical isomorphism}

\begin{equation}\tag{\SGARefTarget{sga7.sga7i.eix.n11-2-1}{11.2.1}}
\Phi_o(\ell)\simeq
 \frac{(U\otimes D_\ell)^I}{U^I\otimes D_\ell},
\end{equation}
\underline{where $U$ is defined in (\SGARefLink{sga7.sga7i.eix.n11-1-5}{11.1.5}) and
$D_\ell=\mathbb Q_\ell/\mathbb Z_\ell$.}

Thus the group $\Phi_o(\ell)$ can be described solely in terms of the
Tate module $U=T_\ell(A(\bar K))$, regarded as a Galois module under the
action of the inertia group $I$.

% END INCLUDED SGA7I_English_current_buildable_source_20260823/components/166_expose_IX_sections_11_0_through_11_2.tex

% BEGIN INCLUDED SGA7I_English_current_buildable_source_20260823/components/167_expose_IX_section_11_3.tex
% Expose IX, section 11.3.
% French transcription lines 3029--3147; authority scan physical pp. 463--466
% (scan indices 462--465; source folios 451--454). Physical p. 466 was
% checked through the exact 11.4 boundary. All pages were compared in four
% overlapping 1100-dpi bands. The text, formulas, arrow directions, labels,
% and two native diagrams are clear at that resolution.

\medskip
\noindent
\SGARefTarget{sga7.sga7i.eix.n11-3}{11.3}. We shall use the explicit expression (\SGARefLink{sga7.sga7i.eix.n11-2-1}{11.2.1}) to sketch the
definition of a canonical pairing

\begin{equation}\tag{\SGARefTarget{sga7.sga7i.eix.n11-3-1}{11.3.1}}
\Phi_o(\ell)\otimes\Phi'_o(\ell)\longrightarrow\mathbb Q/\mathbb Z
\end{equation}
which is a perfect duality; here $\Phi'_o$ denotes, as in section 1, the
group of connected components associated with the N\'eron model of the
\underline{dual} $A'_K$ of $A_K$. There is little doubt (although the
author has dispensed with the verification) that this pairing is none
other than, possibly up to sign, the one defined in (\SGARefLink{sga7.sga7i.eix.n1-2-1}{1.2.1}), and the
reader is invited, as an exercise, to establish the required
compatibility (with the correct sign!).

Put

\begin{equation*}
U'=T_\ell\bigl(A'(\bar K)\bigr),
 \qquad
T=T_\ell(1)(\bar K)=T_\ell(\bar K^*).
\end{equation*}
There is a \underline{perfect pairing} of finitely generated free
$\mathbb Z_\ell$-modules, compatible with the actions of $I$,

\begin{equation}\tag{\SGARefTarget{sga7.sga7i.eix.n11-3-2}{11.3.2}}
\varphi:U\otimes U'\longrightarrow T,
 \qquad (T\simeq\mathbb Z_\ell\text{ noncanonically}),
\end{equation}
and, by the known structure of $I$, an exact sequence
\footnote{The printed source leaves the internal reference following
the first occurrence of $I$ blank.}

\begin{equation}\tag{\SGARefTarget{sga7.sga7i.eix.n11-3-3}{11.3.3}}
1\longrightarrow R\longrightarrow I\longrightarrow T\longrightarrow1,
\end{equation}
where $R$ is a profinite group \underline{prime to $\ell$}. Starting now
from the data (\SGARefLink{sga7.sga7i.eix.n11-3-2}{11.3.2}) and (\SGARefLink{sga7.sga7i.eix.n11-3-3}{11.3.3}), introduce the groups

\begin{equation}\tag{\SGARefTarget{sga7.sga7i.eix.n11-3-4}{11.3.4}}
\begin{aligned}
\Phi&=\frac{(U\otimes D_\ell)^I}{U^I\otimes D_\ell},\qquad
\Phi'=\frac{(U'\otimes D_\ell)^I}{{U'}^I\otimes D_\ell},\qquad
D_\ell=\mathbb Q_\ell/\mathbb Z_\ell.
\end{aligned}
\end{equation}
We shall define a perfect duality

\begin{equation}\tag{\SGARefTarget{sga7.sga7i.eix.n11-3-5}{11.3.5}}
\Phi\otimes\Phi'\longrightarrow\mathbb Q/\mathbb Z,
\end{equation}
which will in particular give the announced pairing (\SGARefLink{sga7.sga7i.eix.n11-3-1}{11.3.1}). Consider
the exact sequence of $I$-modules

\begin{equation}\tag{\SGARefTarget{sga7.sga7i.eix.n11-3-6}{11.3.6}}
0\longrightarrow U\xrightarrow{\ i\ }
 U\otimes_{\mathbb Z_\ell}\mathbb Q_\ell
 \xrightarrow{\ j\ }U\otimes D_\ell\longrightarrow0
\end{equation}
and the associated exact cohomology sequence, where the cohomology of
$I$ is defined by \underline{continuous cochains} with coefficients in
the topological groups under consideration:

\begin{equation}\tag{\SGARefTarget{sga7.sga7i.eix.n11-3-7}{11.3.7}}
\resizebox{\textwidth}{!}{%
\begin{tikzcd}[ampersand replacement=\&, column sep=small, row sep=large]
0 \arrow[r]
 \& H^o(I,U) \arrow[r, "i^o"] \arrow[d, no head, "\wr"]
 \& H^o(I,U\otimes\mathbb Q_\ell) \arrow[r, "j^o"]
      \arrow[d, no head, "\wr"]
 \& H^o(I,U\otimes D_\ell) \arrow[r, "\partial"]
      \arrow[d, no head, "\wr"]
 \& H^1(I,U) \arrow[r, "i^1"]
 \& H^1(I,U\otimes\mathbb Q_\ell) \arrow[d, no head, "\wr"] \\
{}
 \& U^I
 \& (U\otimes\mathbb Q_\ell)^I\simeq U^I\otimes_{\mathbb Z}\mathbb Q
 \& (U\otimes D_\ell)^I
 \& {}
 \& H^1(I,U)\otimes_{\mathbb Z}\mathbb Q .
\end{tikzcd}}
\end{equation}
It is immediate that the group $\Phi$ defined by (\SGARefLink{sga7.sga7i.eix.n11-3-4}{11.3.4}) is precisely
$\operatorname{coker}j^o$, and is therefore canonically isomorphic to
$\ker i^1$. Thus we obtain a canonical exact sequence

\begin{equation}\tag{\SGARefTarget{sga7.sga7i.eix.n11-3-8}{11.3.8}}
0\longrightarrow\Phi\longrightarrow H^1(I,U)
 \longrightarrow H^1(I,U)\otimes_{\mathbb Z}\mathbb Q,
\end{equation}
which shows that $\Phi$ \underline{is canonically isomorphic to the
torsion subgroup of $H^1(I,U)$}.

On the other hand, from the exact sequence (\SGARefLink{sga7.sga7i.eix.n11-3-3}{11.3.3}) we deduce, for every
$n=\ell^r$, an isomorphism

\begin{equation}\tag{\SGARefTarget{sga7.sga7i.eix.n11-3-9}{11.3.9}}
H^1(I,T_n)\simeq\operatorname{Hom}(I,T_n)
 \xleftarrow{\sim}\operatorname{Hom}(T,T_n)
 \xleftarrow{\sim}\mathbb Z/n\mathbb Z.
\end{equation}
(Notice that $I$ acts trivially on $T$, and hence on $T_n=T/nT$.) It
follows that, for every finite $I$-$\mathbb Z_\ell$-module $M$, if

\begin{equation}\tag{\SGARefTarget{sga7.sga7i.eix.n11-3-10}{11.3.10}}
M'=\operatorname{Hom}(M,T\otimes D_\ell)
\end{equation}
denotes its ``Cartier dual,'' there is a canonical pairing

\begin{equation}\tag{\SGARefTarget{sga7.sga7i.eix.n11-3-11}{11.3.11}}
H^i(I,M)\otimes H^{1-i}(I,M')
 \longrightarrow H^1(I,T\otimes D_\ell)\simeq D_\ell,
\end{equation}
where the last isomorphism is deduced from the isomorphisms (\SGARefLink{sga7.sga7i.eix.n11-3-9}{11.3.9}) by
passage to the inductive limit over the powers $n$ of $\ell$. It is well
known that these pairings are \underline{perfect dualities} (see, for
example, SGA 5 I 5). Taking $M$ in these pairings to be of the form
$U_n$, so that $M'=U'_n$, and passing to the limit, we obtain a perfect
pairing

\begin{equation}\tag{\SGARefTarget{sga7.sga7i.eix.n11-3-12}{11.3.12}}
H^o(I,U\otimes D_\ell)\times H^1(I,U')\longrightarrow D_\ell,
\end{equation}
where the second factor on the left is a finitely generated
$\mathbb Z_\ell$-module. A further passage to the limit in the pairings
(\SGARefLink{sga7.sga7i.eix.n11-3-12}{11.3.12}), for transition morphisms of the form
$\ell^r\operatorname{id}$, gives a second perfect pairing

\begin{equation}\tag{\SGARefTarget{sga7.sga7i.eix.n11-3-13}{11.3.13}}
H^o(I,U\otimes\mathbb Q_\ell)
 \times H^1(I,U'\otimes\mathbb Q_\ell)
 \longrightarrow D_\ell
\end{equation}
of finite-dimensional vector spaces over $\mathbb Q_\ell$. Let $i'$ and
$j'$ denote the morphisms defined as in (\SGARefLink{sga7.sga7i.eix.n11-3-6}{11.3.6}), with $U$ replaced by
$U'$. It is then easy to see that the following homomorphisms

\begin{center}
\begin{tikzcd}[column sep=6em, row sep=5em]
H^o(I,U\otimes\mathbb Q_\ell) \arrow[r, "j^o"]
  & H^o(I,U\otimes D_\ell) \\
H^1(I,U'\otimes\mathbb Q_\ell)
  & H^1(I,U') \arrow[l, "{i'}^1"']
\end{tikzcd}
\end{center}
are \underline{transposes} of one another under the pairings (\SGARefLink{sga7.sga7i.eix.n11-3-12}{11.3.12})
and (\SGARefLink{sga7.sga7i.eix.n11-3-13}{11.3.13}). It follows immediately that
$\operatorname{coker}j^o$ and $\ker {i'}^1$ are dual to one another by a
canonical pairing with values in $D_\ell$; this is the announced pairing
(\SGARefLink{sga7.sga7i.eix.n11-3-5}{11.3.5}).

% END INCLUDED SGA7I_English_current_buildable_source_20260823/components/167_expose_IX_section_11_3.tex

% BEGIN INCLUDED SGA7I_English_current_buildable_source_20260823/components/168_expose_IX_sections_11_4_and_11_5.tex
% Expose IX, sections 11.4 and 11.5.
% French transcription lines 3148--3191; authority scan physical pp. 466--468
% (scan indices 465--467; source folios 454--456). Physical p. 468 was
% checked through the exact 11.6 boundary. All pages were compared in four
% overlapping 1100-dpi bands. The text, formulas, primes, checks, underlines,
% and displayed arrows are clear at that resolution; no escalation was needed.

\medskip
\noindent
\SGARefTarget{sga7.sga7i.eix.n11-4}{11.4}. The preceding description of a pairing (\SGARefLink{sga7.sga7i.eix.n11-3-1}{11.3.1}), by means of the
expression (\SGARefLink{sga7.sga7i.eix.n11-2-1}{11.2.1}) for the $\ell$-primary component of the group of
connected components, is valid only when $\ell\ne p$. Suppose now that
$A_K$ has \underline{semistable reduction} over $S$. We shall again
define pairings that are perfect dualities

\begin{equation}\tag{\SGARefTarget{sga7.sga7i.eix.n11-4-1}{11.4.1}}
\Phi_o(\ell)\otimes\Phi'_o(\ell)\longrightarrow\mathbb Q/\mathbb Z
\end{equation}
without any restriction on $\ell$. As in \SGARefLink{sga7.sga7i.eix.n11-3}{11.3},\footnote{The printed
source reads ``as in \SGARefLink{sga7.sga7i.eix.n11-4}{11.4}'' here, within section \SGARefLink{sga7.sga7i.eix.n11-4}{11.4} itself. The
reference is corrected to the immediately preceding construction in
\SGARefLink{sga7.sga7i.eix.n11-3}{11.3}.} one should establish that this recovers (up to a sign, perhaps,
which ought to be specified) the pairing defined in (\SGARefLink{sga7.sga7i.eix.n1-2-1}{1.2.1}). This would
therefore establish Conjecture \SGARefLink{sga7.sga7i.eix.n1-3}{1.3} in the case of semistable reduction,
just as it does when one restricts to the $\ell$-primary components for
$\ell\ne p$.

To define (\SGARefLink{sga7.sga7i.eix.n11-4-1}{11.4.1}), we shall give, essentially, an expression for
$\Phi_o(\ell)$ in terms of the Barsotti--Tate group $T_\ell(A_K)$ over
$K$; in the present case this will play the same role as \SGARefLink{sga7.sga7i.eix.n11-2}{11.2} did in
the case treated above.

\medskip
\noindent
\underline{Theorem} \SGARefTarget{sga7.sga7i.eix.n11-5}{11.5}. \underline{Suppose that $S$ is henselian and
that $A_K$ has semistable reduction over $S$. Let $\ell$ be a prime and
consider the monodromy pairing} (\SGARefLink{sga7.sga7i.eix.n9-1-2}{9.1.2}), \underline{and hence the
canonical homomorphism}

\begin{equation*}
\widetilde u_\ell:\underline M_\ell
 \longrightarrow\underline{\check M}'_\ell.
\end{equation*}
\underline{Put $\Phi(\ell)=\operatorname{coker}\widetilde u_\ell$.
Then there is a canonical isomorphism}

\begin{equation}\tag{\SGARefTarget{sga7.sga7i.eix.n11-5-1}{11.5.1}}
\Phi_o(\ell)\xrightarrow{\ \sim\ }\Phi(\ell)\times_S s.
\end{equation}

This statement immediately gives the definition of a perfect pairing
(\SGARefLink{sga7.sga7i.eix.n11-4-1}{11.4.1}). Indeed, applying \SGARefLink{sga7.sga7i.eix.n11-5}{11.5} to $A'_K$ in place of $A_K$ and using
(\SGARefLink{sga7.sga7i.eix.n10-2-2}{10.2.2}), we also obtain a canonical isomorphism

\begin{equation*}
\Phi'_o(\ell)\xrightarrow{\ \sim\ }\Phi'(\ell)\times_S s,
\qquad
\Phi'(\ell)=\operatorname{coker}\bigl(
 {}^t\!\widetilde u_\ell:\underline M'_\ell
 \longrightarrow\underline{\check M}_\ell\bigr).
\end{equation*}

\medskip
\noindent
\underline{Remarks} \SGARefTarget{sga7.sga7i.eix.n11-5-2}{11.5.2}. a) Theorem 11.5 implies in particular that
$\widetilde u_\ell$ is an isogeny, that is, that $u_\ell$ is
separating, including in the case $\ell=p$, where this assertion is not
immediate from the definition in \SGARefLink{sga7.sga7i.eix.n9-6}{9.6}. Notice that we do not use here
the (considerably deeper) Theorem \SGARefLink{sga7.sga7i.eix.n10-4}{10.4}, which will be proved in the next
section.

\smallskip
\noindent
b) Using \SGARefLink{sga7.sga7i.eix.n10-4}{10.4}, which supplies us with a monodromy pairing

\begin{equation*}
u:\underline M\otimes\underline M'\longrightarrow\mathbb Z_S,
\end{equation*}
and denoting by $\Phi$ the cokernel of the corresponding homomorphism

\begin{equation*}
\widetilde u:\underline M\longrightarrow\underline{\check M}',
\end{equation*}
we obtain an equivalent formulation of \SGARefLink{sga7.sga7i.eix.n11-5}{11.5}, independent of the choice
of $\ell$, in the canonical isomorphism

\begin{equation}\tag{\SGARefTarget{sga7.sga7i.eix.n11-5-3}{11.5.3}}
\Phi_o\xrightarrow{\ \sim\ }\Phi\times_S s.
\end{equation}

% END INCLUDED SGA7I_English_current_buildable_source_20260823/components/168_expose_IX_sections_11_4_and_11_5.tex

% BEGIN INCLUDED SGA7I_English_current_buildable_source_20260823/components/169_expose_IX_section_11_6_opening_through_lemma_11_6_13.tex
% Expose IX, section 11.6 opening through Lemma 11.6.13 (statement).
% French transcription lines 3192--3299; authority physical pp. 468--471
% (scan indices 467--470; source folios 456--459). Physical p. 471 was
% checked through the exact proof boundary. All pages were compared in four
% overlapping 1100-dpi bands. Checks, primes, subscripts, hooks, underlines,
% quotient bars, and formula punctuation are clear at that resolution.

\medskip
\noindent
\SGARefTarget{sga7.sga7i.eix.n11-6}{11.6}. To define the isomorphism (\SGARefLink{sga7.sga7i.eix.n11-5-1}{11.5.1}), we may restrict ourselves to
the case in which $S$ is strictly local (\SGARefLink{sga7.sga7i.eix.n11-0}{11.0}), so that $\underline M$
and $\underline M'$ are constant, with values

\begin{equation}\tag{\SGARefTarget{sga7.sga7i.eix.n11-6-1}{11.6.1}}
M=D(T'_o),\qquad M'=D(T_o),
\end{equation}
where $T_o$ and $T'_o$ are respectively the maximal tori of $A_o$ and
$A'_o$. Likewise identifying $\Phi_o$ with an ordinary finite group,
we must define a canonical isomorphism between its component
$\Phi_o(\ell)$ and

\begin{equation*}
\Phi(\ell)=\operatorname{coker}\bigl(
 \widetilde u_\ell:M_\ell\longrightarrow\check M'_\ell\bigr).
\end{equation*}
In fact, we shall directly define a canonical isomorphism

\begin{equation}\tag{\SGARefTarget{sga7.sga7i.eix.n11-6-2}{11.6.2}}
\Phi_o(\ell)\xrightarrow{\ \sim\ }
 \ker\bigl(\widetilde u_\ell\otimes D_\ell:
 M_\ell\otimes D_\ell\longrightarrow
 \check M'_\ell\otimes D_\ell\bigr).
\end{equation}
Since $\Phi_o(\ell)$ is finite and $M_\ell$ and $\check M'_\ell$ have
the same rank, it will first follow that
$\widetilde u_\ell:M_\ell\to\check M'_\ell$ is an isogeny. It then
follows that the right-hand side of (\SGARefLink{sga7.sga7i.eix.n11-6-2}{11.6.2}) is canonically isomorphic
to $\Phi(\ell)$:

\begin{equation}\tag{\SGARefTarget{sga7.sga7i.eix.n11-6-3}{11.6.3}}
\Phi(\ell)=\operatorname{coker}\widetilde u_\ell
 \xrightarrow{\ \sim\ }
 \ker(\widetilde u_\ell\otimes D_\ell).
\end{equation}
Indeed, identifying $M_\ell$ by means of $\widetilde u_\ell$ with a
finite-index submodule of $\check M'_\ell$, we may write the source and
target of $\widetilde u_\ell\otimes D_\ell$ as\footnote{At this point
the printed source labels $V/M_\ell$ as $\Phi(\ell)$ and
$V/\check M'_\ell$ as $\Phi'(\ell)$. Those quotients are instead the
source and target of the map whose kernel is the finite cokernel
$\check M'_\ell/M_\ell$; the contradictory labels are not reproduced.}

\begin{equation*}
M_\ell\otimes D_\ell=V/M_\ell,\qquad
\check M'_\ell\otimes D_\ell=V/\check M'_\ell,
\end{equation*}
where

\begin{equation*}
V=M_\ell\otimes\mathbb Q_\ell
 \xrightarrow{\ \sim\ }\check M'_\ell\otimes\mathbb Q_\ell.
\end{equation*}
Composing (\SGARefLink{sga7.sga7i.eix.n11-6-2}{11.6.2}) with (\SGARefLink{sga7.sga7i.eix.n11-6-3}{11.6.3}) now immediately gives the isomorphism
(\SGARefLink{sga7.sga7i.eix.n11-5-1}{11.5.1}) to be defined, so everything reduces to defining the canonical
isomorphism (\SGARefLink{sga7.sga7i.eix.n11-6-2}{11.6.2}).

Writing $M_\ell\otimes D_\ell$ (respectively,
$\check M'_\ell\otimes D_\ell$) as the inductive limit of the $M_n$
(respectively, the $\check M'_n$), where $n$ ranges over the powers of
$\ell$, it is equivalent to define canonical isomorphisms

\begin{equation}\tag{\SGARefTarget{sga7.sga7i.eix.n11-6-4}{11.6.4}}
{}_n\Phi_o\xrightarrow{\ \sim\ }
 \ker\bigl(\widetilde u_n:M_n\longrightarrow\check M'_n\bigr),
 \qquad n=\ell^{r+1},
\end{equation}
compatible with the evident transition morphisms as $n$ varies. Here,
by an abuse of notation, $\widetilde u_n$ now denotes the homomorphism
obtained by reducing
$\widetilde u_\ell:M_\ell\to\check M'_\ell$ modulo $n$.

For this purpose, consider the group scheme ${}_nA$ over $S$. By the
semistable-reduction hypothesis it is flat and quasi-finite over $S$
(\SGARefLink{sga7.sga7i.eix.n2-2-1}{2.2.1}). There is a canonical filtration

\begin{equation}\tag{\SGARefTarget{sga7.sga7i.eix.n11-6-5}{11.6.5}}
{}_nA\supset({}_nA^o)^f\supset({}_nA^o)^t\supset0,
\end{equation}
where the superscript $f$ (respectively, $t$) denotes the ``fixed part''
(\SGARefLink{sga7.sga7i.eix.n2-2-3}{2.2.3}) (respectively, the ``toric part'' (\SGARefLink{sga7.sga7i.eix.n5-1}{5.1})) of the group scheme to
which it is applied. In the definition of the toric part in \SGARefLink{sga7.sga7i.eix.n5-1}{5.1}, we had
assumed $S$ complete; this is permissible here (\SGARefLink{sga7.sga7i.eix.n11-0}{11.0}). (We saw moreover
in \SGARefLink{sga7.sga7i.eix.n6-1}{6.1} that the construction generalizes to the henselian case, if one
insists.) Put

\begin{equation}\tag{\SGARefTarget{sga7.sga7i.eix.n11-6-6}{11.6.6}}
{}_n\Psi={}_nA/({}_nA^o)^f.
\end{equation}
This is a flat, separated, quasi-finite group scheme over $S$
(SGA 3 $\mathrm{VI}_A$ 3.2 and $\mathrm{VI}_B$ \SGARefLink{sga7.sga7i.eix.n9-2}{9.2}), with generic fibre

\begin{equation*}
{}_n\Psi_K=({}_nA_K)/({}_nA_K)^f,
\end{equation*}
where, by definition,
$({}_nA_K)^f=({}_nA^o)^f_K$ denotes the $r$-th constituent
$(n=\ell^{r+1})$ of the Barsotti--Tate pro-group
$T_\ell(A_K)^f$. The first of relations (\SGARefLink{sga7.sga7i.eix.n5-5-8}{5.5.8}) therefore gives a
canonical isomorphism

\begin{equation}\tag{\SGARefTarget{sga7.sga7i.eix.n11-6-7}{11.6.7}}
{}_n\Psi_K\xrightarrow{\ \sim\ }\underline M_{nK},
\qquad
\underline M_n\mathrel{\mathop:}=\underline M
 \otimes_{\mathbb Z}(\mathbb Z/n\mathbb Z).
\end{equation}
In fact, we have:

\medskip
\noindent
\underline{Lemma} \SGARefTarget{sga7.sga7i.eix.n11-6-8}{11.6.8}. \underline{The group scheme ${}_n\Psi$ over
$S$ defined in (\SGARefLink{sga7.sga7i.eix.n11-6-6}{11.6.6}) is \emph{\'etale} over $S$, and the isomorphism
(\SGARefLink{sga7.sga7i.eix.n11-6-7}{11.6.7}) extends uniquely to a morphism}

\begin{equation}\tag{\SGARefTarget{sga7.sga7i.eix.n11-6-8-1}{11.6.8.1}}
{}_n\Psi\lhook\joinrel\longrightarrow\underline M_n
\end{equation}
\underline{which is an open immersion.}

To see that ${}_n\Psi$ is \'etale, in view of (\SGARefLink{sga7.sga7i.eix.n11-6-7}{11.6.7}) it remains only
to verify this at the points of its special fibre, or equivalently to
verify that this fibre is \'etale. This is immediate, since the fibre in
question is

\begin{equation*}
{}_nA_o/{}_nA_o^o\subset A_o/A_o^o=\Phi_o.
\end{equation*}
It follows that ${}_n\Psi$ is normal; and since $\underline M_n$ is
finite over $S$, the existence of the extension (\SGARefLink{sga7.sga7i.eix.n11-6-8-1}{11.6.8.1}) follows at
once. This is a morphism of \'etale group schemes over $S$, so its
kernel is \'etale over $S$. Since ${}_n\Psi$, and hence this kernel, is
separated, while its generic fibre is the unit group by (\SGARefLink{sga7.sga7i.eix.n11-6-7}{11.6.7}), the
kernel is the unit group (\SGARefLink{sga7.sga7i.eix.n3-1-3}{3.1.3}). Thus (\SGARefLink{sga7.sga7i.eix.n11-6-8-1}{11.6.8.1}) is an \'etale
monomorphism, hence an open immersion (EGA IV 17.9.1), proving the lemma.

Now put

\begin{equation}\tag{\SGARefTarget{sga7.sga7i.eix.n11-6-9}{11.6.9}}
{}_n\Phi=({}_nA)^f/({}_nA^o)^f\subset{}_n\Psi.
\end{equation}
Thus ${}_n\Phi$ is a finite open subscheme of ${}_n\Psi$, since
$({}_nA)^f$ is a finite open subgroup scheme of ${}_nA$. More
precisely, it is immediate that

\begin{equation}\tag{\SGARefTarget{sga7.sga7i.eix.n11-6-10}{11.6.10}}
{}_n\Phi\xrightarrow{\ \sim\ }({}_n\Psi)^f.
\end{equation}
Composing (\SGARefLink{sga7.sga7i.eix.n11-6-9}{11.6.9}) with (\SGARefLink{sga7.sga7i.eix.n11-6-8-1}{11.6.8.1}), we obtain an open immersion

\begin{equation}\tag{\SGARefTarget{sga7.sga7i.eix.n11-6-11}{11.6.11}}
{}_n\Phi\lhook\joinrel\longrightarrow\underline M_n.
\end{equation}
On the other hand,

\begin{equation*}
{}_n\Phi\times_Ss={}_n\Psi\times_Ss={}_nA_o/{}_nA_o^o,
\end{equation*}
and, since $A_o^o$ is $n$-divisible (\SGARefLink{sga7.sga7i.eix.n2-2-1}{2.2.1}), the last member is
immediately identified with ${}_n\Phi_o$. We therefore obtain an
isomorphism

\begin{equation}\tag{\SGARefTarget{sga7.sga7i.eix.n11-6-12}{11.6.12}}
{}_n\Phi\times_Ss\xrightarrow{\ \sim\ }{}_n\Phi_o,
\end{equation}
which shows that ${}_n\Phi$ is characterized, up to unique isomorphism,
as the finite \'etale scheme over $S$ whose special fibre is
${}_n\Phi_o$. Thus, to define the desired isomorphism (\SGARefLink{sga7.sga7i.eix.n11-6-4}{11.6.4}) by means
of the open immersion (\SGARefLink{sga7.sga7i.eix.n11-6-11}{11.6.11}), it is enough to prove the following.

\medskip
\noindent
\underline{Lemma} \SGARefTarget{sga7.sga7i.eix.n11-6-13}{11.6.13}. \underline{The image of the canonical
homomorphism (\SGARefLink{sga7.sga7i.eix.n11-6-11}{11.6.11}) is equal to the kernel of}

\begin{equation*}
\widetilde u_n:\underline M_n\longrightarrow\underline{\check M}'_n.
\end{equation*}

The required compatibility of the homomorphisms (\SGARefLink{sga7.sga7i.eix.n11-6-11}{11.6.11}) with the
transition homomorphisms as $n$ varies is immediate from the
definitions. Consequently Lemma \SGARefLink{sga7.sga7i.eix.n11-6-13}{11.6.13} will complete the proof of
Theorem \SGARefLink{sga7.sga7i.eix.n11-5}{11.5}.

% END INCLUDED SGA7I_English_current_buildable_source_20260823/components/169_expose_IX_section_11_6_opening_through_lemma_11_6_13.tex

% BEGIN INCLUDED SGA7I_English_current_buildable_source_20260823/components/170_expose_IX_lemma_11_6_13_proof_and_remarks_11_7.tex
% Expose IX, proof of Lemma 11.6.13 and Remarks 11.7.
% French transcription lines 3304--3391; authority physical pp. 472--474
% (scan indices 471--473; source folios 460--462). Physical p. 474 was
% checked through the exact 11.8 boundary. All pages were compared in four
% overlapping 1100-dpi bands. The two compact diagrams, their nodes, ten plus
% five arrows, equality/isomorphism marks, labels, bars, and punctuation are
% all decisive at that resolution; no higher-resolution escalation was needed.

To prove the lemma, note that

\begin{equation*}
K_n=\ker\widetilde u_n
\end{equation*}
is also the left annihilator of the canonical pairing

\begin{equation*}
u_n:\underline M_n\times\underline M'_n
 \longrightarrow(\mathbb Z/n\mathbb Z)_S,
\end{equation*}
whose explicit definition is given in \SGARefLink{sga7.sga7i.eix.n9-4}{9.4} and \SGARefLink{sga7.sga7i.eix.n9-6}{9.6} in terms of the
extensions

\begin{equation}\tag{\SGARefTarget{sga7.sga7i.eix.n11-6-14}{11.6.14}}
\resizebox{\textwidth}{!}{%
\begin{tikzcd}[ampersand replacement=\&, column sep=small, row sep=large]
0 \arrow[r]
 \& {\ensuremath{\check{\underline M}'_n(1)}} \arrow[r]
 \& F_n \arrow[r] \arrow[d, equal]
 \& R_n \arrow[r]
 \& 0
 \& {,}
 \& 0 \arrow[r]
 \& {\ensuremath{\check{\underline M}_n(1)}} \arrow[r]
 \& F'_n \arrow[r] \arrow[d, equal]
 \& R'_n \arrow[r]
 \& 0 \\
{}
 \& {}
 \& ({}_nA^o)^f
 \& {}
 \& {}
 \& {}
 \& {}
 \& {}
 \& ({}_nA'^o)^f
 \& {}
 \& {}
\end{tikzcd}}
\end{equation}
(where $R_n$ and $R'_n$ are Cartier dual to one another), and in terms
of the mixed extension $X_{n\eta}$ of $E_{n\eta}$ by $F_{n\eta}$,
which is the generic fibre of (\SGARefLink{sga7.sga7i.eix.n11-6-5}{11.6.5}),\footnote{The printed source
refers here to nonexistent formula (11.7.5); the filtration whose
generic fibre is being used is (\SGARefLink{sga7.sga7i.eix.n11-6-5}{11.6.5}).} where we put

\begin{equation}\tag{\SGARefTarget{sga7.sga7i.eix.n11-6-15}{11.6.15}}
\begin{tikzcd}[column sep=large, row sep=large]
E_n=D(F_n),
 & 0 \arrow[r]
 & R_n \arrow[r] \arrow[d, no head, "\wr"']
 & E_n \arrow[r]
 & \underline M_n \arrow[r]
 & 0
 & {.} \\
{}
 & {}
 & D(R'_n)
 & {}
 & {}
 & {}
 & {}
\end{tikzcd}
\end{equation}
It follows immediately from Proposition \SGARefLink{sga7.sga7i.eix.n9-4-7}{9.4.7}\footnote{The printed
source cites \SGARefLink{sga7.sga7i.eix.n9-4-6}{9.4.6}, but the required extension criterion is Proposition
\SGARefLink{sga7.sga7i.eix.n9-4-7}{9.4.7}.} that $K_n$ is the largest among the finite \'etale subgroups
$K'_n$ of $\underline M_n$ for which the inverse image
$\overline X_{n\eta}$ of $K'_n$ in the mixed extension $X_{n\eta}$,
regarded as a mixed extension of $\overline E_{n\eta}$ by $F_{n\eta}$,
extends to a mixed extension of $\overline E_n$ by $F_n$. Here
$\overline E_n$ is the evident subgroup extension of $K'_n$ by
$R_n$.\footnote{The printed source says ``of $K_n$'' in this general
criterion, although the subgroup being varied at this point is $K'_n$;
the specialization $K'_n=K_n$ is made only below.}

This condition is plainly satisfied for the subgroup ${}_n\Phi$ of
$\underline M_n$ in (\SGARefLink{sga7.sga7i.eix.n11-6-11}{11.6.11}), since the corresponding mixed extension
$({}_nA)^f$ over $S$ is available. Hence ${}_n\Phi\subset K_n$. To
prove the opposite inclusion, consider the mixed extension
$\overline X_n$ of $\overline E_n$ by $F_n$ extending
$\overline X_{n\eta}$, defined as above by taking $K'_n=K_n$. Thus it
is an extension

\begin{equation*}
0\longrightarrow F_n\longrightarrow\overline X_n
 \longrightarrow K_n\longrightarrow0,
\end{equation*}
and there is a canonical homomorphism

\begin{equation}\tag{*}
F_n=({}_nA^o)^f\lhook\joinrel\longrightarrow A,
\end{equation}
which on the generic fibre is induced by the canonical homomorphism

\begin{equation}\tag{**}
\overline X_{n\eta}\lhook\joinrel\longrightarrow X_{n\eta}
 ={}_nA_K\lhook\joinrel\longrightarrow A_K.
\end{equation}
It follows immediately from (\SGARefLink{sga7.sga7i.eix.n5-9-2}{5.9.2}) that the homomorphism $(**)$ extends
to a unique homomorphism

\begin{equation}\tag{***}
\overline X_n\longrightarrow A
\end{equation}
inducing $(*)$. Since this last homomorphism maps the generic fibre into
${}_nA$, which is a closed subscheme of $A$, it factors through
$\overline X_n\to{}_nA$; and since $\overline X_n$ is finite, it factors
through $({}_nA)^f$. Looking at the induced homomorphism on generic
fibres, we conclude that

\begin{equation*}
\overline X_{n\eta}\subset({}_nA)^f_\eta,
\quad\text{i.e.}\quad
K_{n\eta}\subset{}_n\Phi_\eta,
\quad\text{i.e.}\quad
K_n\subset{}_n\Phi.
\end{equation*}
This proves the lemma.

\medskip
\noindent
\underline{Remarks} \SGARefTarget{sga7.sga7i.eix.n11-7}{11.7}. The construction made in \SGARefLink{sga7.sga7i.eix.n11-6}{11.6}\footnote{The
printed source says ``the construction made in \SGARefLink{sga7.sga7i.eix.n11-7}{11.7}'' at the opening of
\SGARefLink{sga7.sga7i.eix.n11-7}{11.7}; the reference is to the construction just completed in \SGARefLink{sga7.sga7i.eix.n11-6}{11.6}.}
shows that, from the extensions (\SGARefLink{sga7.sga7i.eix.n11-6-14}{11.6.14}) over $S$ and the mixed
extension $X_{n\eta}={}_nA_K$ of $E_n$ by $F_n$, one can reconstruct,
up to unique isomorphism, the mixed extension (\SGARefLink{sga7.sga7i.eix.n11-6-5}{11.6.5}), and hence also
the subgroup

\begin{equation*}
({}_nA)^f\supset({}_nA^o)^f
\end{equation*}
of ${}_nA$. Indeed, by \SGARefLink{sga7.sga7i.eix.n9-4}{9.4} these data determine the pairing

\begin{equation*}
u_n:\underline M_n\times\underline M'_n
 \longrightarrow(\mathbb Z/n\mathbb Z)_S,
\end{equation*}
and the preceding construction gives $({}_nA)^f$ as the mixed extension
$\overline X_n$ of $\overline E_n$ by $F_n$, extending the mixed
extension $\overline X_{n\eta}$ induced by $X_{n\eta}$. Here
$\overline E_n$ is the inverse image of $\ker\widetilde u_n$ under
$E_n\to\underline M_n$. By (9.4.8), this mixed extension
$\overline X_n$ is thus determined \underline{up to unique
isomorphism}. We then recover ${}_nA$ by gluing $\overline X_n$ and
$X_{n\eta}$ along their common open subscheme $\overline X_{n\eta}$.

Thus, if the structural data (\SGARefLink{sga7.sga7i.eix.n9-6-1}{9.6.1}) and the mixed extension

\begin{equation*}
X_\eta=T_\ell(A_K)
\end{equation*}
of $E_\eta$ by $F_\eta$ are known---equivalently, if the extensions
(\SGARefLink{sga7.sga7i.eix.n11-6-14}{11.6.14}) and $X_{n\eta}$ are known for every $n=\ell^{r+1}$,
$r\geq0$---then one can reconstruct the inductive (or, if preferred,
projective) system of mixed extensions (\SGARefLink{sga7.sga7i.eix.n11-6-5}{11.6.5}).

% END INCLUDED SGA7I_English_current_buildable_source_20260823/components/170_expose_IX_lemma_11_6_13_proof_and_remarks_11_7.tex

% BEGIN INCLUDED SGA7I_English_current_buildable_source_20260823/components/171_expose_IX_sections_11_8_through_11_11.tex
% Expose IX, sections 11.8--11.11.
% French transcription lines 3393--3446; authority physical pp. 474--476
% (scan indices 473--475; source folios 462--464). Physical p. 477 was
% checked through the exact section-12 boundary. All pages were compared in
% four overlapping 1100-dpi bands; every formula, accent, subscript, arrow,
% underline, and punctuation mark was decisive without higher escalation.
% The page images restore two transcription slips: the quotient sign in
% (11.8.1), and Phi rather than Psi on the left of (11.8.3).

\SGARefTarget{sga7.sga7i.eix.n11-8}{11.8}. As in \SGARefLink{sga7.sga7i.eix.n8-3}{8.3}, regard the separable closure $\overline K$ of $K$ as the
inductive limit of its finite subextensions $K'$. For every such $K'$, let
$A[K']$ denote the N\'{e}ron model of

\begin{equation*}
A_{K'}=A_K\otimes_K K'.
\end{equation*}
It gives a finite \'{e}tale group scheme

\begin{equation}\tag{\SGARefTarget{sga7.sga7i.eix.n11-8-1}{11.8.1}}
\Phi_o[K']=A[K']_o/A[K']_o^o
\end{equation}
of connected components, defined over the residue field of the normalization
$S[K']$ of $S$ in $K'$. For simplicity, suppose that the residue field of
$S$ is separably closed, so that the same is true of the residue fields of
the $S[K']$. The groups (\SGARefLink{sga7.sga7i.eix.n11-8-1}{11.8.1}) may then be identified with ordinary finite
groups $\Phi[K']$, and it is immediate that, as $K'$ varies, these form an
inductive system, which we now propose to study. We shall be interested in
particular in the group

\begin{equation}\tag{\SGARefTarget{sga7.sga7i.eix.n11-8-2}{11.8.2}}
\Psi=\varinjlim_{K'}\Phi[K'].
\end{equation}
We know from \SGARefLink{sga7.sga7i.eix.n3-6}{3.6} that, for $K'$ sufficiently large, say
$K'\supset K'_o$, $A[K']$ has semistable reduction. Hence for
$K'\supset K'_o$ the transition morphisms are injective (\SGARefLink{sga7.sga7i.eix.n3-3}{3.3}). Moreover,
with the notation of \SGARefLink{sga7.sga7i.eix.n5-5}{5.5}, retained in the present section, there is for
$K'\supset K'_o$ a canonical isomorphism

\begin{equation}\tag{\SGARefTarget{sga7.sga7i.eix.n11-8-3}{11.8.3}}
\Phi[K'](\ell)\simeq
\ker\bigl(\widetilde u[K']_\ell\otimes D_\ell:
 M_\ell\otimes D_\ell\longrightarrow
 \check M'_\ell\otimes D_\ell\bigr),
\end{equation}
coming from (\SGARefLink{sga7.sga7i.eix.n11-6-2}{11.6.2}) applied to $A[K']$. One also verifies immediately that
the corresponding immersions

\begin{equation*}
\Phi[K'](\ell)\lhook\joinrel\longrightarrow M_\ell\otimes D_\ell
\end{equation*}
are compatible with the transition morphisms as $K'$ varies. Passing to the
limit therefore gives a canonical monomorphism

\begin{equation}\tag{\SGARefTarget{sga7.sga7i.eix.n11-8-4}{11.8.4}}
\Psi(\ell)\mathrel{\mathop{=}^{\mathrm{def}}}
 \varinjlim_{K'}\Phi[K'](\ell)
 \lhook\joinrel\longrightarrow M_\ell\otimes D_\ell
\end{equation}
from the $\ell$-primary component of $\Psi$; summing over all $\ell$ gives

\begin{equation}\tag{\SGARefTarget{sga7.sga7i.eix.n11-8-5}{11.8.5}}
\Psi\mathrel{\mathop{=}^{\mathrm{def}}}
 \varinjlim_{K'}\Phi[K']
 \longrightarrow M\otimes\mathbb Q/\mathbb Z.
\end{equation}

\medskip
\noindent
\underline{Theorem} \SGARefTarget{sga7.sga7i.eix.n11-9}{11.9}. \underline{The canonical homomorphism (\SGARefLink{sga7.sga7i.eix.n11-8-5}{11.8.5})
is an isomorphism}.

We already know that it is a monomorphism. To prove that it is an
isomorphism, we must use the characterization (\SGARefLink{sga7.sga7i.eix.n11-8-3}{11.8.3}) of the image of
$\Phi[K'](\ell)$ and the behaviour of $\widetilde u[K']$ as $K'$ varies.
That behaviour is given by (\SGARefLink{sga7.sga7i.eix.n10-3-4}{10.3.4}).

% SGA-SRC-001281 — internal reference omitted in source; target not recoverable
We may reduce to the case in which $S$ is strictly local. To prove \SGARefLink{sga7.sga7i.eix.n11-9}{11.9} we
are then reduced to the following familiar assertion, applied to $S[K'_o]$:
if $S$ is a strictly local trait with fraction field $K$, then for every
integer $n\geq1$ there is a separable extension $K'$ of $K$ whose degree is
divisible by $n$.
\footnote{The printed source appends ``cf. I'' followed by a blank internal
reference. No target is supplied, so none is invented here.}

\medskip
\noindent
\underline{Corollary} \SGARefTarget{sga7.sga7i.eix.n11-10}{11.10}. \underline{Under the hypotheses ($S$
complete) and with the notation of \SGARefLink{sga7.sga7i.eix.n8-3}{8.3}, the ind-finite group scheme over
$K$ of connected components of the group scheme $A_K^{\flat}$,
namely $A_K^{\flat}/{A_K^{\flat}}^o$, is canonically isomorphic to}
$\underline M_K\otimes\mathbb Q/\mathbb Z$, \underline{where
$\underline M_K$ is the twisted constant group over $K$ defined in
(\SGARefLink{sga7.sga7i.eix.n8-1-8}{8.1.8})}.

This follows immediately from (8.2.4) and \SGARefLink{sga7.sga7i.eix.n11-9}{11.9}, in view of the definition
of $\underline M_K$.

\medskip
\noindent
\underline{Remarks} \SGARefTarget{sga7.sga7i.eix.n11-11}{11.11}. a) If one restricts to the $\ell$-primary
components for $\ell\neq p$, Theorem \SGARefLink{sga7.sga7i.eix.n11-9}{11.9} is immediate in terms of \SGARefLink{sga7.sga7i.eix.n11-2}{11.2} and
does not require the more delicate considerations of \SGARefLink{sga7.sga7i.eix.n11-6}{11.6} (via \SGARefLink{sga7.sga7i.eix.n9-4}{9.4}).
Likewise, if $K$ is assumed to be \underline{of characteristic zero}, \SGARefLink{sga7.sga7i.eix.n11-9}{11.9}
can be proved without using (\SGARefLink{sga7.sga7i.eix.n11-8-3}{11.8.3}): it is immediate that the image of
${}_n\Phi[K'](\ell)$ in

\begin{equation*}
M_n={}_n(M_\ell\otimes D_\ell)
\end{equation*}
contains the image of the group ${}_nA(K')$ of $K'$-valued points of
${}_nA$. For $K'$ sufficiently large this group is equal to
${}_nA(\overline K)$, and therefore maps onto $M_n$. If
$\operatorname{char}K=p$, on the other hand, it appears that there is no
``trivial'' proof---that is, no proof avoiding \SGARefLink{sga7.sga7i.eix.n9-4}{9.4}---of \SGARefLink{sga7.sga7i.eix.n11-10}{11.10} for the
$p$-primary component.

b) The existence of an isomorphism

\begin{equation*}
\Psi\simeq(\mathbb Q/\mathbb Z)^r,
\end{equation*}
together with the correct value of $r$, was conjectured in 1964 by
J.-P. Serre.

% END INCLUDED SGA7I_English_current_buildable_source_20260823/components/171_expose_IX_sections_11_8_through_11_11.tex

% BEGIN INCLUDED SGA7I_English_current_buildable_source_20260823/components/172_expose_IX_section_12_1_through_12_2.tex
\medskip
\noindent
12. \underline{N\'eron model of a Jacobian and the Picard--Lefschetz
formula}

Let us first recall the following general results of M. Raynaud
[22, Ths.~5 and 7; 22 bis].

\medskip
\noindent
\underline{Theorem} \SGARefTarget{sga7.sga7i.eix.n12-1}{12.1} (M. Raynaud). \underline{Let $S$ be a strictly
local trait, and let $f:X\to S$ be a scheme proper and flat over $S$, with
one-dimensional fibres, such that}

\begin{equation}\tag{\SGARefTarget{sga7.sga7i.eix.n12-1-1}{12.1.1}}
 f_*\underline{\mathcal O}_X\simeq\underline{\mathcal O}_S,
\end{equation}

\noindent
\underline{and such that the local rings of $X$ are factorial. For every
maximal point $x_i$ of $X_o$, put}

\begin{equation}\tag{\SGARefTarget{sga7.sga7i.eix.n12-1-2}{12.1.2}}
 d_i=\operatorname{length}\underline{\mathcal O}_{X_o,x_i},
 \qquad d_i^t=\mu_i d_i,
\end{equation}

\noindent
\underline{where $\mu_i$ is the radicial multiplicity of $k(x_i)$ over
$k$} (EGA IV, 4.7.4) \underline{(thus a power of the characteristic
exponent of $k$). Let $d$ (respectively $d^t$) be the gcd of the integers
$d_i$ (respectively $d_i^t$); hence $d=d^t$ if $k$ is perfect. Put}

\begin{equation}\tag{\SGARefTarget{sga7.sga7i.eix.n12-1-3}{12.1.3}}
 P=\underline{\operatorname{Pic}}_{X/S}:
   (\operatorname{Sch})_{/S}^{,o}\longrightarrow(\operatorname{Ab}).
\end{equation}

\noindent
a) \underline{Consider the following conditions}:

\smallskip
\noindent
1) $d^t=1$.

\smallskip
\noindent
2) $X$ \underline{has a zero-cycle of degree $1$}.

\smallskip
\noindent
3) \underline{There exists an invertible sheaf $L$ on $X\times P^o$ such
that the corresponding homomorphism}
$P^o_\eta\to\underline{\operatorname{Pic}}_{X_\eta/\eta}=P_\eta$
\underline{is the canonical inclusion}. (It is known, by J.-P. Murre,
that $P_\eta$ is representable by a group scheme locally of finite type
over $\eta$; see, for example, SGA 6, XII, 1.5.)

\smallskip
\noindent
4) $X$ \underline{is cohomologically flat over $S$} (EGA III, 7.8.1)
(equivalently, cohomologically flat in dimension zero, or equivalently
$k\xrightarrow{\sim}H^o(X_o,\underline{\mathcal O}_{X_o})$),
\underline{and $d=1$}.

\smallskip
\noindent
5) $P$ \underline{is representable} (equivalently, representable by a
group scheme smooth over $S$).

\smallskip
\noindent
6) $P^o$ \underline{is representable} (equivalently, representable by a
group scheme smooth over $S$). (For the definition of $P^o$, see SGA 3,
$\mathrm{VI}_B$, \SGARefLink{sga7.sga7i.eix.n3-1}{3.1}.)

\smallskip
\noindent
7) $P^o$ \underline{is separated over $S$}, that is, it satisfies the
valuative criterion of separatedness of EGA II, 7.2.3 (equivalently,
every section $s$ of $P^o$ over $S$ for which $s(\eta)=0$ is zero).

\smallskip
\noindent
8) $d=1$.

\smallskip
\noindent
\underline{Then one has the implications}

\begin{equation}\tag{\SGARefTarget{sga7.sga7i.eix.n12-1-4}{12.1.4}}
 1)\Longleftrightarrow2)\Longrightarrow3)\Longrightarrow4)
 \Longleftrightarrow5)\Longleftrightarrow6)\Longleftrightarrow7)
 \Longrightarrow8).
\end{equation}

\noindent
\underline{In particular, if $d=d^t$, for example if $k$ is perfect, the
conditions 1) through 8) are equivalent}.

\medskip
\noindent
b) \underline{Suppose that $d^t=1$, so that $P$ is representable by a
group scheme smooth over $S$. Then the schematic closure of the identity
section of $P$ is a group subscheme $E$ which is \`etale over $S$. If $I$
is the set of reduced irreducible components $X_o^i$ of $X_o$, one obtains
an isomorphism}

\begin{equation}\tag{\SGARefTarget{sga7.sga7i.eix.n12-1-5}{12.1.5}}
 \mathbb Z^I\simeq\Gamma(S,E)
\end{equation}

\noindent
\underline{by associating to every $i\in I$ the section $S_i$ of $P$ over
$S$ (zero over $\eta$) defined by the invertible sheaf
$\underline{\mathcal O}_X(X_o^i)$ on $X$}\footnote{The printed source says
``on $S$'' here and later writes $X_i^o$ for the same divisor. The sheaf is
the invertible sheaf on $X$ associated with the component $X_o^i$; the
English text keeps that notation consistently.}. (Because the local rings of $X$
are factorial by hypothesis, the cycles $X_o^i$ are indeed divisors; thus
$\underline{\mathcal O}_X(X_o^i)$ has its stated meaning as the invertible
sheaf associated with that divisor; see EGA IV, 21.2.8.1.) \underline{One
has}

\begin{equation}\tag{\SGARefTarget{sga7.sga7i.eix.n12-1-6}{12.1.6}}
 P^o\cap E=\underline{\text{the zero section of }P}.
\end{equation}

\noindent
c) \underline{Under the conditions of b), the quotient}

\begin{equation}\tag{\SGARefTarget{sga7.sga7i.eix.n12-1-7}{12.1.7}}
 Q=P/E
\end{equation}

\noindent
\underline{is representable by a smooth separated group scheme over $S$,
and this group scheme is N\'eron, that is, it satisfies the universal
property (\SGARefLink{sga7.sga7i.eix.n1-1-2}{1.1.2}) in terms of its special fibre}.

\medskip
\noindent
d) \underline{Under the conditions of c), suppose in addition that the
geometric generic fibre of $X$ is integral. Then there is a canonical
exact sequence}

\begin{equation}\tag{\SGARefTarget{sga7.sga7i.eix.n12-1-8}{12.1.8}}
 0\longrightarrow A\longrightarrow Q
   \xrightarrow{\ \deg\ }\mathbb Z_S\longrightarrow0,
\end{equation}

\noindent
\underline{where $\deg$ is induced by the ``degree'' homomorphism
$P\to\mathbb Z_S$, and where $A$ is a smooth separated group scheme of
finite type over $S$, which is likewise N\'eron, that is, a N\'eron model
of}

\begin{equation}\tag{\SGARefTarget{sga7.sga7i.eix.n12-1-9}{12.1.9}}
 A_\eta=P_\eta^o=
 \underline{\operatorname{Pic}}^o_{X_\eta/\eta}.
\end{equation}

\noindent
\underline{Moreover, the composite homomorphism $P^o\to P\to Q$ induces
an isomorphism}

\begin{equation}\tag{\SGARefTarget{sga7.sga7i.eix.n12-1-10}{12.1.10}}
 P^o\xrightarrow{\ \sim\ }A^o.
\end{equation}

\medskip
\noindent
\SGARefTarget{sga7.sga7i.eix.n12-1-11}{12.1.11}. The case of interest to us is that in which $X_\eta$ is smooth
over $\eta$, so that, in view of (\SGARefLink{sga7.sga7i.eix.n12-1-1}{12.1.1}), the geometric generic fibre of
$X$ is smooth and connected (hence integral), and in which $X_\eta$ has a
zero-cycle of degree one. All the conditions considered in \SGARefLink{sga7.sga7i.eix.n12-1}{12.1} (and in
particular conditions 1) through 8) of a)) then hold. Moreover,
$A_\eta=P_\eta^o$ is then an abelian scheme, and d) gives us a very
explicit construction of its N\'eron model $A$. In particular, (\SGARefLink{sga7.sga7i.eix.n12-1-10}{12.1.10})
gives an isomorphism

\begin{equation}\tag{\SGARefTarget{sga7.sga7i.eix.n12-1-12}{12.1.12}}
 A_o^{,o}\simeq\underline{\operatorname{Pic}}^o_{X_o/k},
\end{equation}

\noindent
which shows in particular that $A$ \underline{has semistable reduction
over $S$ if and only if}
$\underline{\operatorname{Pic}}^o_{X_o/k}$ \underline{has unipotent rank
zero}. When $X$ is regular and $X_o$ is separable over $k$, this is known
to hold if and only if the geometric special fibre $\overline X_o$ has no
singular points other than ordinary double points (that is, points whose
completed local ring is the same as that of the union of the two
coordinate axes of affine $2$-space at the origin). Verification of this
fact, together with the ``well-known'' description of the Picard scheme
of a proper separable curve over a field $k$, using EGA IV, 21.8.5, is left
to the reader as an exercise. Notice also that the implication
``$\overline X_o$ has ordinary double points'' $\Longrightarrow$
``$A_\eta=\underline{\operatorname{Pic}}^o_{X_\eta/\eta}$ has semistable
reduction'' follows immediately from the ``geometric monodromy theorem''
III\footnote{The printed source leaves the internal locator following
``III'' blank; no target is invented here.}, in view of the Galois
criterion for semistable reduction in \SGARefLink{sga7.sga7i.eix.n3-5}{3.5}. This argument, however, does
not give the important isomorphism (\SGARefLink{sga7.sga7i.eix.n12-1-10}{12.1.10}) and its consequence
(\SGARefLink{sga7.sga7i.eix.n12-1-12}{12.1.12}).

\medskip
\noindent
\SGARefTarget{sga7.sga7i.eix.n12-2}{12.2}. For the proof of \SGARefLink{sga7.sga7i.eix.n12-1}{12.1}, we refer to [22 bis]; in fact we shall need
only the case where $X_\eta$ is smooth over $\eta$. Let us also note that
[22 bis] proves relations between Picard schemes and N\'eron models under
conditions substantially more general than those of \SGARefLink{sga7.sga7i.eix.n12-1}{12.1} c) and d). As
their statements are more technical, we have chosen not to include them
in this review. We mention only that b), c), and d) are easy consequences
of a), and that in a) the only delicate points are the implications
$3)\Longrightarrow4)\Longrightarrow5)$ and $7)\Longrightarrow4)$, the
first being the most difficult. This point is trivial, however, if one
assumes that $X_o$ is separable over $k$ (which implies that $X$ has a
section over $S$, since $S$ is strictly local, and hence that condition
1), $d^t=1$, holds). We advise the reader to give the proof in this
particularly simple case, assuming in addition, if necessary, that
$X_\eta$ is smooth, so that $P_\eta^o$ is an abelian scheme. One may then
use its N\'eron model $A$, and $A^o$ represents
$\underline{\operatorname{Pic}}^o_{X/S}$. We shall use \SGARefLink{sga7.sga7i.eix.n12-1}{12.1} only in this
particular case, and under the further assumption that $\overline X_o$
has no singularities other than ordinary double points, that is, by
(\SGARefLink{sga7.sga7i.eix.n12-1-11}{12.1.11}), that $A_\eta$ has semistable reduction. This is the case
considered in 12.5.4, to which we restrict ourselves from now on.

% END INCLUDED SGA7I_English_current_buildable_source_20260823/components/172_expose_IX_section_12_1_through_12_2.tex

% BEGIN INCLUDED SGA7I_English_current_buildable_source_20260823/components/173_expose_IX_section_12_3.tex
\medskip
\noindent
\SGARefTarget{sga7.sga7i.eix.n12-3}{12.3}. Let $C$ be a proper separable curve over a field $k$, which for
simplicity we suppose separably closed (or even algebraically closed, if
the reader prefers). If $\widetilde C$ is the normalization of $C$, it is
known that $\underline{\operatorname{Pic}}^o_{\widetilde C/k}$ has
``reductive rank'' zero (that is, its maximal torus is zero). Consequently
the maximal torus of $\underline{\operatorname{Pic}}^o_{C/k}$ is contained
in

\begin{equation*}
 N=\ker\bigl(\underline{\operatorname{Pic}}_{C/k}
       \longrightarrow\underline{\operatorname{Pic}}_{\widetilde C/k}\bigr),
\end{equation*}

\noindent
and is therefore canonically isomorphic to the maximal torus of this last
group. The structure of $N$ is easily made explicit by EGA IV, 21.8.5:
$N$ is canonically an extension of a torus $T$, which we shall describe,
by a connected unipotent group $U$ (which itself has, over $k$, a
composition series whose quotients are isomorphic to the additive group
$\mathbb G_a$). Moreover, if $T'$ is the maximal torus of $N$, then
$T'\cap U$ is both unipotent and of multiplicative type, and is therefore
zero (SGA 3, XVII, 2.4), while the image of $T'$ in $T$ is all of $T$
(SGA 3, XII, 7.1 e)). Thus $T'\to T$ is an isomorphism. Hence $T$ is
canonically isomorphic to the maximal torus of $N$, and therefore to that
of $\underline{\operatorname{Pic}}^o_{C/k}$.
\footnote{The printed source has $\underline{\operatorname{Pic}}^o_{X/k}$
at the end of this argument. The curve throughout \SGARefLink{sga7.sga7i.eix.n12-3}{12.3} is $C$, and the
preceding and following occurrences both read
$\underline{\operatorname{Pic}}^o_{C/k}$; the English text restores $C$.}

To make the structure of $N$ explicit, it is convenient to suppose that
$k$ is separably closed. The irreducible components $C_i$ of $C$ are then
geometrically irreducible, and the nonnormal points of $C$ are radicial
over $k$. Let $J$ be the set of irreducible components of $\widetilde C$,
and for every $x\in C$ let $J(x)$ be the set of ``branches'' of $C$ through
$x$ (that is, the points of $\widetilde C$ above $x$). Finally put

\begin{equation}\tag{\SGARefTarget{sga7.sga7i.eix.n12-3-1}{12.3.1}}
 R(x)=\mathbb Z^{J(x)}/\operatorname{Im}\mathbb Z,
\end{equation}

\noindent
where $\mathbb Z\to\mathbb Z^{J(x)}$ is the diagonal map, and

\begin{equation}\tag{\SGARefTarget{sga7.sga7i.eix.n12-3-2}{12.3.2}}
 R=\left(\bigoplus_x R(x)\right)/\operatorname{Im}\mathbb Z^J.
\end{equation}

\noindent
Clearly it is enough to take the sum over the finite set of points of $C$
which belong to at least two branches of $C$. Thus $R$ is a finitely
generated $\mathbb Z$-module. There is then a canonical isomorphism

\begin{equation}\tag{\SGARefTarget{sga7.sga7i.eix.n12-3-3}{12.3.3}}
 T\ \bigl(\simeq\text{ the maximal torus of }
       \underline{\operatorname{Pic}}^o_{C/k}\bigr)
 \simeq R\otimes\mathbb G_{m,k}.
\end{equation}

\noindent
In other words, the character group $M$ of $T$ is

\begin{equation}\tag{\SGARefTarget{sga7.sga7i.eix.n12-3-4}{12.3.4}}
 M\mathrel{\mathop{=}^{\mathrm{def}}}D(T)=\check R
 \simeq\ker\left(\bigoplus_x\check R(x)\longrightarrow\mathbb Z^J\right),
\end{equation}

\noindent
where $\check R(x)$ may also be interpreted as the set of
integer-valued functions on $J(x)$ whose values have sum zero.
\footnote{The printed source numbers these two displays (12.2.3) and
(12.2.4). They belong to section \SGARefLink{sga7.sga7i.eix.n12-3}{12.3} and are cited below as (\SGARefLink{sga7.sga7i.eix.n12-3-3}{12.3.3}) and
(\SGARefLink{sga7.sga7i.eix.n12-3-4}{12.3.4}); those coherent numbers are used here.}

\medskip
\noindent
\SGARefTarget{sga7.sga7i.eix.n12-3-5}{12.3.5}. We shall need only the case in which at most two branches of $C$
pass through any point $x$ (as happens, in particular, for ordinary double
points). If $\operatorname{card}J(x)=2$, then $R(x)$ is isomorphic to
$\mathbb Z$, but the isomorphism depends on an ordering of the two branches
$(C',C'')$ through $x$, or equivalently on the choice of one branch $C'$:
to $C'$ one associates the element of $R(x)$ which is the canonical image
of $0_{C'}+1_{C''}$. One then obtains an isomorphism

\begin{equation}\tag{\SGARefTarget{sga7.sga7i.eix.n12-3-6}{12.3.6}}
 R=\mathbb Z^I/\operatorname{Im}\mathbb Z^J,
 \qquad
 T\simeq\mathbb G_m^I/\operatorname{Im}\mathbb G_m^J,
 \qquad
 M=\ker\bigl(\mathbb Z^I\longrightarrow\mathbb Z^J\bigr),
\end{equation}

\noindent
where $I$ is the set of points $x\in C$ which belong to two branches of
$C$, after a choice of signs as just explained.

\medskip
\noindent
\underline{Remarks} \SGARefTarget{sga7.sga7i.eix.n12-3-7}{12.3.7}. A sometimes convenient way to describe the
character group $M$ of the maximal torus $T$ of
$\underline{\operatorname{Pic}}^o_{C/k}$ is the following. Consider the
``unoriented one-dimensional graph'' $\Gamma(C)$ (N. Bourbaki,
\emph{Groupes et alg\`ebres de Lie}, Chapter IV, Appendix), whose vertices
are the irreducible components of $\widetilde C$. Two vertices
$\widetilde C_i$ and $\widetilde C_j$ are joined by a set of edges in
bijective correspondence with the pairs of \underline{distinct} branches
of $C$, contained respectively in $\widetilde C_i$ and
$\widetilde C_j$, which pass through the same point of $C$. This set is
$\widetilde C_i\times_C\widetilde C_j$ when $i\neq j$; when $i=j$, it is

\begin{equation*}
 \bigl(\widetilde C_i\times_C\widetilde C_i-
       \Delta(\widetilde C_i)\bigr)/(\text{symmetry}),
\end{equation*}

\noindent
where $\Delta:\widetilde C\to\widetilde C\times_C\widetilde C$ is the
diagonal. (It is therefore empty if $C_i$ is geometrically unibranch.) One
then has canonical isomorphisms

\begin{equation}\tag{\SGARefTarget{sga7.sga7i.eix.n12-3-7.occ-2}{12.3.7}}
 M=D(T)\xleftarrow{\ \sim\ }H_1(\Gamma(C),\mathbb Z),
 \qquad
 \check M=R\simeq H^1(\Gamma(C),\mathbb Z).
\end{equation}

We shall not need this combinatorial interpretation of $M$, and shall
confine ourselves to defining the canonical homomorphism

\begin{equation}\tag{\SGARefTarget{sga7.sga7i.eix.n12-3-8}{12.3.8}}
 H_1(\Gamma(C),\mathbb Z)\longrightarrow M,
\end{equation}

\noindent
leaving it to the reader to prove that this is an isomorphism. Let $I$ be
the set of points of $C$ which belong to at least two branches of $C$,
let $J$, as above, be the set of irreducible components of $\widetilde C$
(or of $C$, which amounts to the same thing), and let $K$ be the set of
``branches'' of $C$ at the points of $I$, that is, the inverse image of
$I$ in $\widetilde C$. There is a canonical map

\begin{equation}\tag{\SGARefTarget{sga7.sga7i.eix.n12-3-9}{12.3.9}}
 h:K\longrightarrow I\times J,
 \qquad h=(f,g),
\end{equation}

\noindent
and the description (\SGARefLink{sga7.sga7i.eix.n12-3-4}{12.3.4}) of $M$ is expressed in terms of $h$ as

\begin{equation}\tag{\SGARefTarget{sga7.sga7i.eix.n12-3-10}{12.3.10}}
 M\simeq\ker\left(
  \mathbb Z^{(h)}:\mathbb Z^{(K)}\longrightarrow
  \mathbb Z^{(I\times J)}\right).
\end{equation}

Indeed, the expression (\SGARefLink{sga7.sga7i.eix.n12-3-3}{12.3.3}) for $R=\check M$ may plainly be
interpreted as the cokernel of the homomorphism
$\mathbb Z^{I\times J}\to\mathbb Z^K$ induced by $h$, which is the
transpose of $\mathbb Z^{(h)}$. The graph $\Gamma(C)$ is likewise described
in terms of $h$ by

\begin{equation*}
 \Gamma_0=J,
 \qquad
 \Gamma_1=(K\times_I K-\Delta(K))/(\text{symmetry}).
\end{equation*}

Choose an orientation on each edge of $\Gamma(C)$, and denote by
$\Gamma_1$ the set of these oriented edges $(x,C',C'')$. Then

\begin{equation}\tag{\SGARefTarget{sga7.sga7i.eix.n12-3-11}{12.3.11}}
 H_1(\Gamma(C),\mathbb Z)=
 \ker\left(d:\mathbb Z^{(\Gamma_1)}\longrightarrow\mathbb Z^{(J)}\right),
\end{equation}

\noindent
where the differential is

\begin{equation}\tag{\SGARefTarget{sga7.sga7i.eix.n12-3-12}{12.3.12}}
 d(x,C',C'')=g(C'')-g(C').
\end{equation}

The homomorphism (\SGARefLink{sga7.sga7i.eix.n12-3-8}{12.3.8}) is induced by the following homomorphism $v$ of
chain complexes:\footnote{The printed display defining (\SGARefLink{sga7.sga7i.eix.n12-3-8}{12.3.8}) carries
no number even though the following text cites it, and (\SGARefLink{sga7.sga7i.eix.n12-3-13}{12.3.13}) writes
$\mathbb Z^{(\Gamma)}$ where the edge chain group is
$\mathbb Z^{(\Gamma_1)}$. The English text supplies the cited label and
the required subscript. The working transcription's
$\widetilde C_i\times_C\widetilde C_1$ in the loop-edge description is
also restored from the scan as
$\widetilde C_i\times_C\widetilde C_i$.}

\begin{equation}\tag{\SGARefTarget{sga7.sga7i.eix.n12-3-13}{12.3.13}}
 \begin{split}
 v=(v_0,v_1):C_\bullet(\Gamma(C))
   &=\left[\mathbb Z^{(\Gamma_1)}\longrightarrow\mathbb Z^{(J)}\right]\\
   &\longrightarrow
     \left[\mathbb Z^{(K)}\longrightarrow\mathbb Z^{(I\times J)}\right],
 \end{split}
\end{equation}

\noindent
given by

\begin{equation}\tag{\SGARefTarget{sga7.sga7i.eix.n12-3-14}{12.3.14}}
 v_1(x,C',C'')=(x,C'')-(x,C'),
 \qquad
 v_0(\widetilde C_j)=(0,\widetilde C_j).
\end{equation}

% END INCLUDED SGA7I_English_current_buildable_source_20260823/components/173_expose_IX_section_12_3.tex

% BEGIN INCLUDED SGA7I_English_current_buildable_source_20260823/components/174_expose_IX_section_12_4_and_theorem_12_5.tex
\medskip
\noindent
\SGARefTarget{sga7.sga7i.eix.n12-4}{12.4}. We now place ourselves in the case considered at the end of \SGARefLink{sga7.sga7i.eix.n12-2}{12.2}.
By (\SGARefLink{sga7.sga7i.eix.n12-1-12}{12.1.12}) and (\SGARefLink{sga7.sga7i.eix.n12-3-6}{12.3.6}), applied with $C=X_o$, the character group
$M$ of the maximal torus $T_o$ of the special fibre $A_o$ of the
N\'eron model $A$ of
$A=\underline{\operatorname{Pic}}^o_{X_\eta/\eta}$ is given by

\begin{equation}\tag{\SGARefTarget{sga7.sga7i.eix.n12-4-1}{12.4.1}}
 M=\ker\left(\mathbb Z^I\xrightarrow{\ d\ }\mathbb Z^J\right),
\end{equation}

\noindent
where $I$ is the set of singular points of $X_o$, and $J$ is the set of
its irreducible components, the homomorphism
$\mathbb Z^I\to\mathbb Z^J$ being given by

\begin{equation}\tag{\SGARefTarget{sga7.sga7i.eix.n12-4-2}{12.4.2}}
 d(x_i)=C'_i-C''_i.
\end{equation}

Here $C'_i$ and $C''_i$ are the two irreducible components of
$\widetilde C$ corresponding to the branches through the double point
$x_i$. The operator $d$ depends on the choice of an order
$(C'_i,C''_i)$ on these two branches, and the isomorphism (\SGARefLink{sga7.sga7i.eix.n12-4-1}{12.4.1})
plainly depends on this system of choices. Nevertheless, if, for
$i\in I$, we denote by

\begin{equation}\tag{\SGARefTarget{sga7.sga7i.eix.n12-4-3}{12.4.3}}
 \varphi_{x_i}:M\longrightarrow\mathbb Z
\end{equation}

\noindent
the composite

\begin{equation*}
 M\longrightarrow\mathbb Z^I
   \xrightarrow{\ \operatorname{pr}_i\ }\mathbb Z,
\end{equation*}

\noindent
where the first homomorphism comes from (\SGARefLink{sga7.sga7i.eix.n12-4-1}{12.4.1}), then, on returning to
the expression (\SGARefLink{sga7.sga7i.eix.n12-3-4}{12.3.4}), one sees at once that $\varphi_{x_i}$ is
well-defined up to sign. More precisely, it depends only on the order
chosen between the two branches $C'_i,C''_i$ through $x_i$, and changes
sign when that order is reversed. It follows in particular that the
bilinear form

\begin{equation}\tag{\SGARefTarget{sga7.sga7i.eix.n12-4-4}{12.4.4}}
 \varphi_{x_i}\otimes\varphi_{x_i}
 \in \check M\otimes_{\mathbb Z}\check M
 \simeq\operatorname{Hom}(M\otimes M,\mathbb Z)
\end{equation}

\noindent
is determined canonically by the singular point $x_i$, without any sign
ambiguity. We may therefore put

\begin{equation}\tag{\SGARefTarget{sga7.sga7i.eix.n12-4-5}{12.4.5}}
 u=\sum_i\varphi_{x_i}\otimes\varphi_{x_i}:
 M\otimes_{\mathbb Z}M\longrightarrow\mathbb Z.
\end{equation}

This gives a well-defined pairing, independent of every choice of
signs. It is plainly \underline{positive definite}, since it is the
restriction to $M\hookrightarrow\mathbb Z^I$ of the standard quadratic
form $\sum_i X_i^2$ on $\mathbb Z^I$.

The notation $M$ for the character group of $T_o$ is not, a priori,
consistent with that of the preceding sections, where this group was
written $M'$, while $M$ denoted the analogous group associated with the
dual abelian scheme $A'_K$ of $A_K$. But since $A_K$ is a Jacobian, it
has a canonical polarization $\xi$ [5]\footnote{The printed citation
following the canonical polarization is ``[5]'' followed by a blank
pinpoint. The English text preserves the reference without inventing a
missing locator.}, which induces a canonical
isomorphism

\begin{equation}\tag{\SGARefTarget{sga7.sga7i.eix.n12-4-6}{12.4.6}}
 \psi_\xi:A_K\xrightarrow{\ \sim\ }A'_K.
\end{equation}

\noindent
By means of this isomorphism we shall identify $A_K$ and $A'_K$, so that
there is no longer any need to distinguish $M$ from $M'$. In particular,
for every prime number $\ell$, the monodromy pairing (\SGARefLink{sga7.sga7i.eix.n9-1-2}{9.1.2}) may be
regarded as a pairing

\begin{equation}\tag{\SGARefTarget{sga7.sga7i.eix.n12-4-7}{12.4.7}}
 u_\ell:M_\ell\otimes_{\mathbb Z_\ell}M_\ell
 \longrightarrow\mathbb Z_\ell.
\end{equation}

With this understood, we can state the principal result of the present
section:

\medskip
\noindent
\underline{Theorem} \SGARefTarget{sga7.sga7i.eix.n12-5}{12.5} (\underline{Picard--Lefschetz formula}).
\underline{Let $S$ be a strictly local trait, and let $X$ be a regular
scheme, projective and flat over $S$, with one-dimensional fibres, whose
geometric generic fibre is smooth and connected and whose geometric
special fibre has no singularities other than ordinary quadratic
singularities. With the notation introduced in \SGARefLink{sga7.sga7i.eix.n12-4}{12.4}, for every prime
number $\ell$, the monodromy pairing (\SGARefLink{sga7.sga7i.eix.n12-4-7}{12.4.7}) is equal to the pairing
obtained from (\SGARefLink{sga7.sga7i.eix.n12-4-5}{12.4.5}) by extension of scalars from $\mathbb Z$ to
$\mathbb Z_\ell$.}

% END INCLUDED SGA7I_English_current_buildable_source_20260823/components/174_expose_IX_section_12_4_and_theorem_12_5.tex

% BEGIN INCLUDED SGA7I_English_current_buildable_source_20260823/components/175_expose_IX_sections_12_6_and_12_7.tex
\medskip
\noindent
\SGARefTarget{sga7.sga7i.eix.n12-6}{12.6}. Notice that this statement proves \SGARefLink{sga7.sga7i.eix.n10-4}{10.4} in the particular case
under consideration when, in \SGARefLink{sga7.sga7i.eix.n10-4}{10.4} b), one takes for $\xi$ the canonical
polarization considered there. As we showed in \SGARefLink{sga7.sga7i.eix.n10-5}{10.5}, this suffices to
imply \SGARefLink{sga7.sga7i.eix.n10-4}{10.4} in the general case.

\medskip
\noindent
\SGARefTarget{sga7.sga7i.eix.n12-7}{12.7}. \underline{Proof of \SGARefLink{sga7.sga7i.eix.n12-5}{12.5}. The case $\ell\ne p$.} It relies
essentially on the \underline{cohomological Picard--Lefschetz formula}
in the situation considered in \SGARefLink{sga7.sga7i.eix.n12-5}{12.5}, a formula giving the action of the
monodromy group $I$ on the following group:\footnote{The working
transcription adds a dual superscript to $T_\ell(A_{\bar\eta})$ in
(\SGARefLink{sga7.sga7i.eix.n12-7-1}{12.7.1}). A direct 5,000-dpi inspection of the printed formula shows no
such superscript, in agreement with the standard identification
$H^1(A_{\bar\eta},\mathbb Z_\ell)=
\operatorname{Hom}(T_\ell(A_{\bar\eta}),\mathbb Z_\ell)$.}

\begin{equation}\tag{\SGARefTarget{sga7.sga7i.eix.n12-7-1}{12.7.1}}
 H^1(X_{\bar\eta},\mathbb Z_\ell)
 \simeq H^1(A_{\bar\eta},\mathbb Z_\ell)
 \simeq\operatorname{Hom}
   \bigl(T_\ell(A_{\bar\eta}),\mathbb Z_\ell\bigr)
\end{equation}

\noindent
when $\ell\ne p$. This formula, which will be proved in a later expos\'e
(Expos\'e XV), and indeed in a substantially more general case, is
essentially identical to assertion \SGARefLink{sga7.sga7i.eix.n12-5}{12.5} for the prime $\ell$ in
question, in view of the definition in \SGARefLink{sga7.sga7i.eix.n9-2}{9.2} of $u_\ell$ for $\ell\ne p$.
It therefore remains to treat the case $\ell=p$. We shall do this by a
method which reduces it to characteristic zero (and, moreover, to the
case where $X_o$ has only one singular point), which we now explain.

\medskip
\noindent
\underline{Remarks \SGARefTarget{sga7.sga7i.eix.n12-7-2}{12.7.2}.} Notice also that, once the announced
reduction has been made, the usual arguments would allow us to reduce to
the case in which the base field is the field $\mathbb C$ of complex
numbers, and then to invoke the Picard--Lefschetz formula in its classical
form---except that there does not seem to be a satisfactory reference
for that formula. The proof of the Picard--Lefschetz formula given in our
Seminar will proceed similarly, by reduction to the transcendental case,
which is amenable to properly transcendental methods. One may therefore
say that, in any event, the proof of \SGARefLink{sga7.sga7i.eix.n12-5}{12.5} (and hence of \SGARefLink{sga7.sga7i.eix.n10-4}{10.4}) proposed
here is transcendental in nature. Let us also point out that the only
other known proof of \SGARefLink{sga7.sga7i.eix.n10-4}{10.4} comes from rigid-analytic theory [24], and must
for that reason likewise be regarded as essentially ``transcendental.''
It would be interesting to find purely algebraic proofs of these results
(although it is by no means evident that this is possible!).

% END INCLUDED SGA7I_English_current_buildable_source_20260823/components/175_expose_IX_sections_12_6_and_12_7.tex

% BEGIN INCLUDED SGA7I_English_current_buildable_source_20260823/components/176_expose_IX_section_12_8_opening_through_before_12_8_9.tex
\medskip
\noindent
\SGARefTarget{sga7.sga7i.eix.n12-8}{12.8}. \underline{Proof of \SGARefLink{sga7.sga7i.eix.n12-5}{12.5}: reduction to the case in which $S$ has
characteristic zero.} We shall use an argument based on the same principle
as the proof of the ``theorem on the essentially tame action of monodromy
on $\pi_1$'' in Expos\'e V.\footnote{The printed source leaves the internal
pinpoint following ``V'' blank; no section target is supplied.} Notice that
we may assume $S$ complete, by \SGARefLink{sga7.sga7i.eix.n10-3}{10.3}. Let $W$ be a Cohen $p$-ring with
residue field $k$, and choose a homomorphism $W\longrightarrow V$ inducing
the identity on residue fields (EGA $0_{\mathrm{IV}}$, 19.8.6). Let
$(x_\lambda)_{\lambda\in L}$ be the set of singular points of $X_o$, and
consider the formal moduli scheme $\underline S$ of $X_o$ over $W$, together
with the ``universal curve''
\[
 \underline f:\underline X\longrightarrow\underline S
\]
whose special fibre is $X_o$ (Expos\'e VI).\footnote{The printed source
likewise leaves the internal pinpoint following ``VI'' blank.} We saw in
the cited place that $\underline S$ is isomorphic to the formal spectrum of
a power-series ring $W[[T_1,\ldots,T_N]]$, and that in $\underline S$ one
can find a normal-crossings divisor

\begin{equation}\tag{\SGARefTarget{sga7.sga7i.eix.n12-8-1}{12.8.1}}
 D=\sum_{\lambda\in L}D_\lambda,
\end{equation}

\noindent
whose irreducible components $D_\lambda$ are indexed by $L$, each
$D_\lambda$ being described by an equation of the form $T_{i_\lambda}=0$.
(If desired, one may suppose $i_\lambda=\lambda$ after indexing $L$ by a
set of integers $[1,r]$.) The component $D_\lambda$ is the image under
$\underline f:\underline X\longrightarrow\underline S$ of the component
$D'_\lambda$ through the singular point $x_\lambda$ of $X_o$ in the
non-smoothness locus
\[
 D'=\bigcup_{\lambda\in L}D'_\lambda
\]
of $\underline f$. By the construction itself,
$f:X\longrightarrow S$ is isomorphic to the pullback of
$\underline f:\underline X\longrightarrow\underline S$ by a suitable
$W$-morphism $g:S\longrightarrow\underline S$:

\begin{equation}\tag{\SGARefTarget{sga7.sga7i.eix.n12-8-2}{12.8.2}}
\begin{tikzcd}
 \underline X \arrow[d,"\underline f"'] &
 X \arrow[l] \arrow[d,"f"] \\
 \underline S & S \arrow[l,"g"']
\end{tikzcd}
\qquad .
\end{equation}

\noindent
Set

\begin{equation}\tag{\SGARefTarget{sga7.sga7i.eix.n12-8-3}{12.8.3}}
 P=\underline{\operatorname{Pic}}^{o}_{\underline X/\underline S}.
\end{equation}

\noindent
Since the special fibre of $\underline f$ is geometrically reduced of
dimension $1$, it follows that $\underline f$ is cohomologically flat in
dimension zero,\footnote{The working transcription reads $\underline F$
here. Direct inspection of the printed page shows $\underline f$, as also
required by the sentence's reference to the special fibre of the universal
curve.} by EGA III, 7.8.1, and that $P$ is formally smooth over
$\underline S$ (the latter fact follows from
$H^2(X_o,\underline{\mathcal O}_{X_o})=0$, since $X_o$ has dimension $1$).
For simplicity, we shall assume that $P$ \underline{is representable by a
group scheme over $\underline S$}; it will necessarily be a smooth group
scheme with connected fibres. This is a special case of an unpublished
result of Mumford cited in [22] (and considerably generalized by Raynaud in
the note cited there, when the base is a trait). Since no proof of either
Mumford's or Raynaud's result has been published, the reader may prefer to
use the general results of M. Artin [1]. They give the representability of
$P$---under the sole assumption that $\underline f$ is proper, flat, and
cohomologically flat in dimension zero---by an ``algebraic space'' in
Artin's sense, which would suffice for our purpose. Put

\begin{equation}\tag{\SGARefTarget{sga7.sga7i.eix.n12-8-4}{12.8.4}}
 U=\underline S-\operatorname{supp}D.
\end{equation}

\noindent
Thus $U$ is the largest open subset of $\underline S$ over which
$\underline X$ is smooth, and hence
\[
 P|_U=\underline{\operatorname{Pic}}^{o}_{\underline X_U/U}
\]
is an abelian scheme over $U$. Its pullback to $\eta$ by
$g:\eta\longrightarrow U$ is the abelian scheme $A_\eta$ over $\eta$ that
we wish to study. More precisely, there is a canonical isomorphism over
$S$ (\SGARefLink{sga7.sga7i.eix.n12-1-10}{12.1.10}):

\begin{equation}\tag{\SGARefTarget{sga7.sga7i.eix.n12-8-5}{12.8.5}}
 P\times_{\underline S}S\simeq A^{o}.
\end{equation}

\noindent
Since the special fibre $P_o=A_o^{o}$ is $\ell$-divisible for every
$\ell$, the same holds for $P$ itself; all its fibres are therefore
extensions of an abelian variety by a torus. Consider the pro-$\ell$ group
$T_\ell(P)$ over $\underline S$, whose pullback by $g$ is precisely
$T_\ell(A^{o})$. Applying the constructions of \SGARefLink{sga7.sga7i.eix.n7-3}{7.3}, one defines a
filtration of this pro-group by a ``fixed part'' and a ``toric part'':

\begin{equation}\tag{\SGARefTarget{sga7.sga7i.eix.n12-8-6}{12.8.6}}
 T_\ell(P)\supset T_\ell(P)^f=T_\ell(\widehat P)
 \supset T_\ell(P)^t=T_\ell(T)\supset0.
\end{equation}

\noindent
Its pullback by $g$ is the analogous filtration of $T_\ell(A^{o})$
studied in sections 2 and 5.

Consider the restriction of this filtration over the open subset $U$ of
$\underline S$. There we have the canonical self-duality of $P_U$, as the
Jacobian of a smooth curve scheme over $U$, and hence a corresponding
pairing

\begin{equation}\tag{\SGARefTarget{sga7.sga7i.eix.n12-8-7}{12.8.7}}
 \varphi:T_\ell(P_U)\otimes T_\ell(P_U)
 \longrightarrow T_\ell(\mathbb G_{mU}).
\end{equation}

\noindent
This is a \underline{self-duality}, in the sense of Barsotti--Tate groups,
of $T_\ell(P_U)$.\footnote{The printed source reads $T_\ell(B_U)$ here,
although no $B_U$ has been introduced and the pairing and every adjacent
term involve $P_U$. The English text restores the internally forced
$P_U$.} We saw in \SGARefLink{sga7.sga7i.eix.n7-4-8}{7.4.8} that, for this pairing,
$T_\ell(P_U)^f$ \underline{and} $T_\ell(P_U)^t$ \underline{are each
other's orthogonal complements}. It follows that
\[
 E_U=T_\ell(P_U)/T_\ell(P_U)^t
\]
is canonically isomorphic to the Barsotti--Tate pro-group dual to
$T_\ell(P_U)^f=F_U$. It therefore extends canonically to a Barsotti--Tate
group $E$ over $\underline S$, namely

\begin{equation}\tag{\SGARefTarget{sga7.sga7i.eix.n12-8-8}{12.8.8}}
 E=D(F),\qquad\text{where}\qquad F=T_\ell(P)^f.
\end{equation}

\noindent
We are now in a typical situation of mixed extensions (\SGARefLink{sga7.sga7i.eix.n9-3}{9.3}), and can repeat
the constructions of \SGARefLink{sga7.sga7i.eix.n9-4}{9.4}--\SGARefLink{sga7.sga7i.eix.n9-6}{9.6} (for whose having been written in the stale
and piss-poor setting of a base scheme reduced to a trait we apologize!):
$T_\ell(P_U)$ is a mixed extension of $E_U$ by $F_U$, where $E$ and $F$
are Barsotti--Tate groups over $\underline S$ which are extensions, dual to
one another:

% END INCLUDED SGA7I_English_current_buildable_source_20260823/components/176_expose_IX_section_12_8_opening_through_before_12_8_9.tex

% BEGIN INCLUDED SGA7I_English_current_buildable_source_20260823/components/177_expose_IX_section_12_8_9_through_12_10.tex
\begin{equation}\tag{\SGARefTarget{sga7.sga7i.eix.n12-8-9}{12.8.9}}
\begin{tikzcd}[row sep=small,column sep=small]
0 \arrow[r] & P \arrow[rr] \arrow[d,dash,"\wr"]
 & & F \arrow[r] \arrow[d,dash,"\wr"] & R \arrow[r] & 0 & \\
& T_\ell(P)^t\simeq
  \underline{\operatorname{Hom}}
    \bigl(M,T_\ell(\mathbb G_{m\underline S})\bigr)
& & T_\ell(P)^f & & & \\
& & & & & & \\
0 \arrow[r] & R \arrow[rr] \arrow[d,dash,"\wr"]
& & E \arrow[r] & Q \arrow[r] \arrow[d,dash,"\wr"] & 0 & \\
& D(R) & & & M\otimes\mathbb Z_{\ell\,\underline S} & & .
\end{tikzcd}
\end{equation}

\noindent
The printed upper identification in (\SGARefLink{sga7.sga7i.eix.n12-8-9}{12.8.9}) omits the outer closing
parenthesis after
$\underline{\operatorname{Hom}}(M,T_\ell(\mathbb G_{m\underline S}))$;
the balanced expression has been supplied here. Notice that $R$ is
canonically self-dual, as follows from the self-duality (\SGARefLink{sga7.sga7i.eix.n12-8-7}{12.8.7}) and the
orthogonality statement noted above. The only difference from the cited
construction is that the pairing obtained on $M_\ell\times M_\ell$ takes
values in the group
\[
 \bigl(\mathbb G_m(U)/\mathbb G_m(\underline S)\bigr)
 \otimes_{\mathbb Z}\mathbb Z_\ell,
\]
which is no longer canonically isomorphic to $\mathbb Z_\ell$, but rather
to $\mathbb Z_\ell^{\,L}$, by means of the isomorphism

\begin{equation}\tag{\SGARefTarget{sga7.sga7i.eix.n12-8-10}{12.8.10}}
 \mathbb Z^L\simeq
 \mathbb G_m(U)/\mathbb G_m(\underline S).
\end{equation}

\noindent
It maps the element indexed by $\lambda$ in the canonical basis of
$\mathbb Z^L$ to the class of $T_{i_\lambda}$ in the right-hand side,
where $T_{i_\lambda}$ is an equation for the divisor $D_\lambda$. We thus
obtain a canonical pairing

\begin{equation}\tag{\SGARefTarget{sga7.sga7i.eix.n12-8-11}{12.8.11}}
 \underline{u}_\ell:M_\ell\otimes M_\ell
 \longrightarrow\mathbb Z_\ell^{\,L}\qquad .
\end{equation}

\noindent
It is essentially immediate that the monodromy pairing $u_\ell$ (\SGARefLink{sga7.sga7i.eix.n12-4-7}{12.4.7})
over $S$ which we propose to study is the composite of
$\underline{u}_\ell$ with the homomorphism
$\mathbb Z_\ell^{\,L}\longrightarrow\mathbb Z_\ell$ obtained, by
tensoring with $\mathbb Z_\ell$ over $\mathbb Z$, from the homomorphism

\begin{center}
\begin{tikzcd}
 \mathbb G_m(U)/\mathbb G_m(\underline S)
   \arrow[r] \arrow[d,dash,"\wr"'] &
 \mathbb G_m(\eta)/\mathbb G_m(S)
   \arrow[d,dash,"\wr"'] \\
 \mathbb Z^L & \mathbb Z
\end{tikzcd}
\end{center}

\noindent
induced by $g:S\longrightarrow\underline S$ (which sends $\eta$ into
$U$). This homomorphism is given by an element

\begin{equation}\tag{\SGARefTarget{sga7.sga7i.eix.n12-8-12}{12.8.12}}
 (d_\lambda)_{\lambda\in L}\in\mathbb Z^L,\qquad
 d_\lambda\geqslant1\quad\text{for every }\lambda\in L,
\end{equation}

\noindent
where, for every $\lambda\in L$, $d_\lambda$ is the multiplicity of the
inverse image in $S$ of the divisor $D_\lambda$ on $\underline S$. It in
fact depends only on the completion
$\widehat{\underline{\mathcal O}}_{X,x_\lambda}$ as a $V$-algebra
(Expos\'e VI).\footnote{Both printed citations to Expos\'e VI in this
sentence leave their internal section pinpoints blank.} Moreover, it is
known (Expos\'e VI) that, for a given $\lambda$, one has $d_\lambda=1$ if
and only if $X$ is \underline{regular} at $x_\lambda$.\footnote{The
printed source has the lowercase letter $x$ as the subject of ``is
regular.'' The following sentence and the mathematical statement require
the scheme $X$.} Since in \SGARefLink{sga7.sga7i.eix.n12-5}{12.5} we assumed $X$ regular, that is, all the
$d_\lambda$ equal to $1$, assertion \SGARefLink{sga7.sga7i.eix.n12-5}{12.5} will therefore follow from the
following lemma, in which $X,S,g$ no longer occur.

\medskip
\noindent
\underline{Lemma} \SGARefTarget{sga7.sga7i.eix.n12-8-13}{12.8.13}. \underline{If $\lambda\in L$, the composite of
(\SGARefLink{sga7.sga7i.eix.n12-8-11}{12.8.11}) with $\operatorname{pr}_\lambda$ is equal to the form
$\varphi_{x_\lambda}\otimes\varphi_{x_\lambda}$ of (\SGARefLink{sga7.sga7i.eix.n12-4-4}{12.4.4}).}

\medskip
To prove this last assertion, consider the completed strict localization
$\underline S^\lambda$ of $\underline S$ at the generic point of
$D_\lambda$,\footnote{In this reduction the printed source drops the
subscript $\lambda$ from $D_\lambda$ and prints an un-underlined $X$ as
the right-hand factor of the following fibre product. Both marks are
restored: the localization is at the component $D_\lambda$, and the
fibre product is formed from the universal curve $\underline X$ over
$\underline S$.} and put
\[
 \underline X^\lambda=
 \underline X\times_{\underline S}\underline S^\lambda.
\]
The special fibre $\underline X_o^\lambda$ then has only one singular
point $y_\lambda$, which is again an ordinary quadratic singularity. In
the corresponding monodromy pairing

\begin{equation}\tag{\SGARefTarget{sga7.sga7i.eix.n12-8-14}{12.8.14}}
 u_\ell^\lambda:M_\ell^\lambda\otimes M_\ell^\lambda
 \longrightarrow\mathbb Z_\ell\qquad ,
\end{equation}

\noindent
$M^\lambda$ consequently has rank $1$ or $0$. We leave to the reader the
verification of the following easy facts, which make explicit, in the
special case in which we need them, the behaviour under localization of
the preceding constructions.

\medskip
\noindent
a) There is a unique subtorus $T_\lambda$ of
$P\times_{\underline S}D_\lambda$ whose restriction to the generic point
$s_\lambda$ of $D_\lambda$ is the maximal torus of $P_{s_\lambda}$ (which
has dimension $0$ or $1$).

\medskip
\noindent
b) Consider the transpose homomorphism $M\longrightarrow M^\lambda$ on
character groups of the inclusion
$T_{\lambda o}\hookrightarrow T_o$, where $T_o$ is the maximal torus of
$P_o$, and the form
$\varphi_{y_\lambda}:M^\lambda\longrightarrow\mathbb Z$, defined up to
sign as in \SGARefLink{sga7.sga7i.eix.n12-4-3}{12.4.3} (with $X$ replaced by $\underline X^\lambda$). Then,
up to sign, the following diagram is commutative:

\begin{equation}\tag{\SGARefTarget{sga7.sga7i.eix.n12-8-15}{12.8.15}}
\begin{tikzcd}[column sep=huge,row sep=large]
 M \arrow[rr] \arrow[dr,"\varphi_{x_\lambda}"']
 & & M^\lambda \arrow[dl,"\varphi_{y_\lambda}"] \\
 & \mathbb Z & \qquad .
\end{tikzcd}
\end{equation}

\noindent
The working transcription misreads the right diagonal label in (\SGARefLink{sga7.sga7i.eix.n12-8-15}{12.8.15})
as $\varphi_{\lambda_\lambda}$. The printed diagram, the immediately
preceding definition, and the formula below all give
$\varphi_{y_\lambda}$. (In fact, choosing an ordering of the two branches
of $X_o$ through $x_\lambda$ determines an analogous choice at the point
$y_\lambda$ of $\underline X_o^\lambda$, and for choices corresponding in
this way the preceding diagram is commutative.)

\medskip
\noindent
c) The following diagram is commutative:

\begin{center}
\begin{tikzcd}[column sep=huge,row sep=large]
 M_\ell\otimes M_\ell
   \arrow[rr] \arrow[dr,"\operatorname{pr}_\lambda\underline{u}_\ell"']
 & & M_\ell^\lambda\otimes M_\ell^\lambda
   \arrow[dl,"u_\ell^\lambda"] \\
 & \mathbb Z_\ell & \qquad ,
\end{tikzcd}
\end{center}

\noindent
where the notation is that of \SGARefLink{sga7.sga7i.eix.n12-8-13}{12.8.13} and (\SGARefLink{sga7.sga7i.eix.n12-8-14}{12.8.14}). These relations show
at once that, to prove \SGARefLink{sga7.sga7i.eix.n12-8-13}{12.8.13}, it is enough to prove the analogous
relation for $\underline X^\lambda$:
\[
 u_\ell^\lambda=
 (\varphi_{y_\lambda}\otimes\varphi_{y_\lambda})
 \otimes_{\mathbb Z}\mathbb Z_\ell,
\]
which is nothing other than \SGARefLink{sga7.sga7i.eix.n12-5}{12.5} for $\underline X^\lambda$ over
$\underline S^\lambda$. Now $\underline S^\lambda$ is a trait of
\underline{characteristic zero}.\footnote{At this final step the printed
source writes $\underline S$ rather than the localized trait
$\underline S^\lambda$ used throughout the reduction. The English text
restores the mathematically required superscript.} This completes the
announced reduction.

\medskip
\noindent
\underline{Remarks \SGARefTarget{sga7.sga7i.eix.n12-9}{12.9}.} Notice that this reduction simultaneously
reduces us to the case in which $X_o$ has only one singular point, the case
classically considered in the Picard--Lefschetz formula. Notice also that,
in this case, $X_o$ has one or two irreducible components, corresponding
respectively to $\operatorname{rank}M=1$ and $\operatorname{rank}M=0$ (the
latter being the case in which the Jacobian of the generic fibre has good
reduction).

\medskip
\noindent
\SGARefTarget{sga7.sga7i.eix.n12-10}{12.10}. The proof given in \SGARefLink{sga7.sga7i.eix.n12-8}{12.8} yields a result more general than \SGARefLink{sga7.sga7i.eix.n12-5}{12.5},
valid without assuming that $X$ is regular, by introducing the integers
$d_\lambda\geqslant1$ of (\SGARefLink{sga7.sga7i.eix.n12-8-12}{12.8.12}): all the monodromy pairings are then
deduced from the pairing

\begin{equation}\tag{\SGARefTarget{sga7.sga7i.eix.n12-10-1}{12.10.1}}
 u=\sum_{\lambda\in L}d_\lambda\,
   \varphi_{x_\lambda}\otimes\varphi_{x_\lambda}:
   M\otimes M\longrightarrow\mathbb Z\qquad .
\end{equation}

% END INCLUDED SGA7I_English_current_buildable_source_20260823/components/177_expose_IX_section_12_8_9_through_12_10.tex

% BEGIN INCLUDED SGA7I_English_current_buildable_source_20260823/components/178_expose_IX_sections_13_1_and_13_2.tex
\medskip
\noindent
13. \underline{Links with transcendental theory: the complex analytic
case}.

For the general notions concerning analytic spaces, we refer to [2].

\medskip
\noindent
\SGARefTarget{sga7.sga7i.eix.n13-1}{13.1}. \underline{Review of analytic groups over $S$.} Let $S$ be an
analytic space, and let $A$ be an analytic space over $S$, endowed with
the structure of a commutative group over $S$ (that is, $A$ is a
commutative group in the category of analytic spaces over $S$). We assume
that $A$ is \underline{smooth} over $S$, and denote by

\begin{equation}\tag{\SGARefTarget{sga7.sga7i.eix.n13-1-1}{13.1.1}}
 \underline E=\underline{\operatorname{Lie}}(A/S)
\end{equation}

\noindent
the locally free module on $S$ obtained by pulling the relative tangent
sheaf of $A$ over $S$ back along the zero section. We shall admit here the
definition of the exponential homomorphism

\begin{equation}\tag{\SGARefTarget{sga7.sga7i.eix.n13-1-2}{13.1.2}}
 \exp:E\longrightarrow A\qquad ,
\end{equation}

\noindent
where $E$ is the analytic vector bundle on $S$ whose module of
holomorphic sections is $\underline E$. In a neighbourhood of a point
$s\in S$, this homomorphism is first defined near the zero section of $E$
by the familiar methods of the theory of formal groups in characteristic
zero, and is then extended in the evident way, by multiplicativity, to
all of $E$. It is a homomorphism of relative analytic groups over $S$;
it is \underline{\'etale}, and its image is the open subspace $A^o$ of
$A$ which is the union of the identity components of the fibres of $A$
over $S$. If $R$ is the analytic subspace which is the kernel of $\exp$,
we therefore have an exact sequence of relative analytic groups

\begin{equation}\tag{\SGARefTarget{sga7.sga7i.eix.n13-1-3}{13.1.3}}
 0\longrightarrow R\longrightarrow E\longrightarrow A^o
 \longrightarrow0\qquad .
\end{equation}

\noindent
Here the morphism $E\longrightarrow A^o$ is surjective and \'etale, so
$R$ is a relative analytic group \underline{\'etale} over $S$. It is an
analytic subspace of $E$, and is closed in $E$ if and only if $A^o$ is
separated over $S$, that is, if and only if the zero section of $A^o$ is
closed in $A^o$.

\medskip
\noindent
\SGARefTarget{sga7.sga7i.eix.n13-1-4}{13.1.4}. Of course, the exact sequence (\SGARefLink{sga7.sga7i.eix.n13-1-3}{13.1.3}) is
\underline{functorial in $A$}, and commutes with every base change
$S'\longrightarrow S$. In this way one obtains an equivalence of
categories, compatible with base changes, between the category of
relative analytic groups over $S$ which are smooth and have connected
fibres, and the category of pairs $(E,R)$ consisting of a locally free
vector bundle $E$ over $S$ and an analytic subgroup $R$ of $E$ which is
\'etale over $S$.

\medskip
\noindent
\SGARefTarget{sga7.sga7i.eix.n13-1-5}{13.1.5}. When $A$ is an ``\underline{abeloid}'' group over $S$, by which
we mean that it is smooth, proper, and separated over $S$, with connected
fibres---so that in particular $A=A^o$---then $R$ is a locally constant
group over $S$, whose rank at each point $s$ (as a $\mathbb Z$-module) is
equal to $2\operatorname{rank}_s(E)$, and conversely. In this way one
obtains the familiar interpretation of an abeloid space over $S$ in
terms of a ``\underline{local system}'' $R$ over $S$ and a
\underline{complex} analytic vector-bundle structure on the real
analytic bundle $R\otimes_{\mathbb Z}\mathbb R$ defined by $R$.

\medskip
\noindent
\SGARefTarget{sga7.sga7i.eix.n13-2}{13.2}. \underline{Formal completion along a fibre.} We now fix a point
$s\in S$, and shall work locally in a neighbourhood of $s$. For every
integer $n\geqslant0$, let $S_n$ denote the $n$-th infinitesimal
neighbourhood of $s$ in $S$, and, for every analytic space $X$ over $S$,
put likewise $X_n=X\times_S S_n$. As $n$ varies, the system of the $X_n$
then forms an inductive system of analytic spaces over $S$, which we
denote by $\widehat X$ and call the \underline{formal completion of $X$
along $X_s$}. It clearly depends functorially on $X$, and its formation
commutes with finite projective limits; hence $X\longmapsto\widehat X$
takes groups to groups. Notice that $\widehat X$ in fact depends only on
the restriction of $X$ to an arbitrary open neighbourhood $U$ of $s$ in
$S$.

\medskip
\noindent
\underline{Lemma \SGARefTarget{sga7.sga7i.eix.n13-2-1}{13.2.1}.} \underline{The data are those of \SGARefLink{sga7.sga7i.eix.n13-1}{13.1}, and
suppose that the following conditions hold:} a) $A=A^o$;
b) $R$ \underline{is constant in a neighbourhood of $s$}; and
c) $R_s$ \underline{spans the complex vector space $E_s$}.
\underline{Let $A'$ be a second smooth analytic group over $S$, and
consider the map}

\begin{equation*}
 u\longmapsto\widehat u:
 \operatorname{Hom}_{S\text{-}\mathrm{gr}}(A,A')
 \longrightarrow
 \operatorname{Hom}_{\widehat S\text{-}\mathrm{gr}}
   (\widehat A,\widehat {A}')\quad ,
\end{equation*}

\noindent
\underline{where the notation is that of \SGARefLink{sga7.sga7i.eix.n13-2}{13.2}, and the analogous map
obtained by restricting over a variable open neighbourhood $U$ of $s$
and passing to the inductive limit:}

\begin{equation}\tag{\SGARefTarget{sga7.sga7i.eix.n13-2-1-1}{13.2.1.1}}
 u\longmapsto\widehat u:
 \operatorname{Hom}_{(S,s)\text{-}\mathrm{gr}}(A,A')
 \longrightarrow
 \operatorname{Hom}_{\widehat S\text{-}\mathrm{gr}}
   (\widehat A,\widehat {A}')\quad .
\end{equation}

\noindent
(\underline{Here $(S,s)$ denotes the germ of the analytic space $S$ at
$s$.}) \underline{Then the latter map is bijective.}

Indeed, consider the exact sequence

\begin{equation*}
 0\longrightarrow R'\longrightarrow E'\longrightarrow {A'}^o
 \longrightarrow C
\end{equation*}

\noindent
analogous to (\SGARefLink{sga7.sga7i.eix.n13-1-3}{13.1.3}) and defined by $A'$. Since the homomorphisms
$A\longrightarrow A'$ necessarily factor through ${A'}^o$, and the same
holds after the base changes $U\longrightarrow S$ and
$S_n\longrightarrow S$, we may assume that $A'={A'}^o$. By the
dictionary of \SGARefLink{sga7.sga7i.eix.n13-1-4}{13.1.4}, it remains to prove that giving a germ of a vector
bundle homomorphism $E\longrightarrow E'$ which takes $R$ into $R'$ is
equivalent to giving a coherent system of homomorphisms
$E_n\longrightarrow E'_n$ which take $R_n$ into $R'_n$.

If $\underline{\mathcal O}_s$ denotes the local ring of $S$ at $s$,
giving a coherent system of homomorphisms
$E_n\longrightarrow E'_n$ is equivalent to giving a homomorphism of
$\widehat{\underline{\mathcal O}}_s$-modules

\begin{equation*}
 \widehat F:\widehat{\underline E}_s
 \longrightarrow\widehat{\underline E}'_s,
\end{equation*}

\noindent
where $\underline E_s$ and $\underline E'_s$ are the free
$\underline{\mathcal O}_s$-modules obtained at $s$ from the
$\underline{\mathcal O}_s$-modules $\underline E$ and $\underline E'$,
corresponding to $E$ and $E'$. On the other hand, a germ of a vector
bundle homomorphism $E\longrightarrow E'$ is equivalent to giving a
homomorphism of $\underline{\mathcal O}_s$-modules

\begin{equation*}
 F:\underline E_s\longrightarrow\underline E'_s\quad .
\end{equation*}

\noindent
This already shows that (\SGARefLink{sga7.sga7i.eix.n13-2-1-1}{13.2.1.1}) is injective. For surjectivity, it is
enough to prove that if $\widehat F$ is a homomorphism as above such that,
for every $n$, the induced homomorphism
$F_n:E_n\longrightarrow E'_n$ takes $R_n$ into $R'_n$, then
$\widehat F$ comes from an $F$---that is, it takes $\underline E_s$ into
$\underline E'_s$---and the corresponding germ of homomorphisms
$E\longrightarrow E'$ takes $R$ into $R'$.

By hypothesis c), $R_s=R_o$ spans $E_s$; it follows from Nakayama's
lemma that the group $R_o$ (identified either with the group of sections
of $\widehat R$ over $\widehat S$, or with the group of germs of sections
of $R$ over $S$) generates $\widehat E_s$. Thus, to prove that
$\widehat F(\underline E_s)\subset\underline E'_s$, it is enough to prove
that $\widehat F(R_o)\subset\underline E'_s$. But
$\widehat F(R_o)\subset R'_o$, where $R'_o$ is itself interpreted as the
group of germs of sections of $R'\subset E'$ in a neighbourhood of $s$.
This gives the asserted inclusion, and hence the existence of $F$.
Moreover, we see that $F(R_o)\subset R'_o$, where $R_o$ and $R'_o$ are
regarded as groups of germs of sections of $E$ and $E'$, respectively.
By b), $R_o$ generates $R$ in a neighbourhood of $s$; it follows that
$F(R)\subset R'$ in a neighbourhood of $s$, which completes the proof.

\medskip
\noindent
\underline{Corollary \SGARefTarget{sga7.sga7i.eix.n13-2-2}{13.2.2}.} \underline{If $A$ satisfies conditions
a), b), and c) of \SGARefLink{sga7.sga7i.eix.n13-2-1}{13.2.1}, it is determined up to unique isomorphism in a
neighbourhood of $s$ by the knowledge of the $\widehat S$-group
$\widehat A$.}\footnote{The printed source refers here to the conditions
of \SGARefLink{sga7.sga7i.eix.n13-2-2}{13.2.2} itself; the three conditions are introduced in Lemma \SGARefLink{sga7.sga7i.eix.n13-2-1}{13.2.1}.}

% END INCLUDED SGA7I_English_current_buildable_source_20260823/components/178_expose_IX_sections_13_1_and_13_2.tex

% BEGIN INCLUDED SGA7I_English_current_buildable_source_20260823/components/179_expose_IX_section_13_3.tex
\medskip
\noindent
\SGARefTarget{sga7.sga7i.eix.n13-3}{13.3}. \underline{The analytic Raynaud extension}.

We are now in a position to paraphrase, in analytic geometry, the
construction of the Raynaud extension given in 7. Suppose that $A$ is
as in \SGARefLink{sga7.sga7i.eix.n13-1}{13.1}, with $A=A^o$,
\footnote{The printed source has $A=A_o$ here. The required hypothesis,
and the one used below, is $A=A^o$; compare \SGARefLink{sga7.sga7i.eix.n13-1}{13.1} and Lemma \SGARefLink{sga7.sga7i.eix.n13-2-1}{13.2.1}.}
and consider $R_o=R_s$, which is a free $\mathbb Z$-module of finite
type. After restricting $S$ to a sufficiently small open neighbourhood
of $s$, and interpreting $R_o$ as the group of germs of sections of $R$
near $s$, one obtains a homomorphism

\begin{equation}\tag{\SGARefTarget{sga7.sga7i.eix.n13-3-1}{13.3.1}}
 R_{oS}\longrightarrow R\subset E,
\end{equation}

\noindent
where $R_{oS}$ denotes the constant analytic group over $S$ with value
$R_o$. On choosing the neighbourhood of $s$ small enough that every
section of $\underline{\mathcal O}_S$ whose germ at $s$ is zero is itself
zero, one sees that the preceding homomorphism is injective. Its image is
therefore a constant analytic subgroup of $R$; it is plainly the largest
constant subgroup of $R$, and is called the \underline{fixed part} of
$R$ and denoted by $R^f$:

\begin{equation}\tag{\SGARefTarget{sga7.sga7i.eix.n13-3-2}{13.3.2}}
 R_{oS}\simeq R^f\hookrightarrow R,
 \qquad\text{inducing an isomorphism }R_s^f\simeq R_s.
\end{equation}

\noindent
We may then consider the quotient group

\begin{equation}\tag{\SGARefTarget{sga7.sga7i.eix.n13-3-3}{13.3.3}}
 A^\natural=E/R^f,
\end{equation}

\noindent
and, in view of (\SGARefLink{sga7.sga7i.eix.n13-1-3}{13.1.3}), obtain a canonical exact sequence

\begin{equation}\tag{\SGARefTarget{sga7.sga7i.eix.n13-3-4}{13.3.4}}
 0\longrightarrow R/R^f\stackrel{i}{\longrightarrow}A^\natural
 \longrightarrow A\longrightarrow0.
\end{equation}

\noindent
Notice that, after restricting $S$ to a neighbourhood of $s$, $R^f$ is
always closed in $E$, and hence $A^\natural$ is separated over $S$.
Notice also that the homomorphism $A^\natural\longrightarrow A$ is
\'etale, and therefore induces an isomorphism

\begin{equation}\tag{\SGARefTarget{sga7.sga7i.eix.n13-3-5}{13.3.5}}
 \underline{\operatorname{Lie}}(A^\natural)
 \stackrel{\sim}{\longrightarrow}
 \underline{\operatorname{Lie}}(A)=\underline E.
\end{equation}

\noindent
Moreover, for every $n\geqslant0$, the canonical homomorphism
$A^\natural\longrightarrow A$ induces an isomorphism
$A_n^\natural\simeq A_n$ (since $R_n^f\stackrel{\sim}{\longrightarrow}
R_n$), and consequently an isomorphism

\begin{equation}\tag{\SGARefTarget{sga7.sga7i.eix.n13-3-6}{13.3.6}}
 \widehat{A^\natural}\stackrel{\sim}{\longrightarrow}\widehat A.
\end{equation}

\noindent
This last fact shows that $A^\natural$ does indeed play the role of the
group, likewise denoted by $A^\natural$, that we considered in 7. Here
it has been constructed without any hypothesis on the fibre
$A_s=A_o$ of the form ``an extension of an abeloid group by a torus''.
Such a hypothesis will nevertheless be necessary if we wish to recover
on $A^\natural$ the important extension structure (\SGARefLink{sga7.sga7i.eix.n7-2-3}{7.2.3}). By contrast,
the homomorphism $A^\natural\longrightarrow A$ in (\SGARefLink{sga7.sga7i.eix.n13-3-4}{13.3.4}) is essentially
transcendental in nature and has no counterpart in the purely algebraic
theory of \S~7. In \S~14 we shall return to the question of a purely
algebraic interpretation of the homomorphism
$i:R/R^f\longrightarrow A^\natural$ in (\SGARefLink{sga7.sga7i.eix.n13-3-4}{13.3.4}).

Suppose first that we are given a complex subtorus

\begin{equation}\tag{\SGARefTarget{sga7.sga7i.eix.n13-3-7}{13.3.7}}
 T_o\subset A_o,
\end{equation}

\noindent
where by a ``\underline{complex torus}'' we mean an analytic group
isomorphic to a Cartesian power of $\mathbb C^*$. Let

\begin{equation}\tag{\SGARefTarget{sga7.sga7i.eix.n13-3-8}{13.3.8}}
 T=T_o\times S
\end{equation}

\noindent
be the constant torus over $S$ with value $T_o$. Proceeding as in \SGARefLink{sga7.sga7i.eix.n5-1}{5.1},
one obtains a unique homomorphism (indeed, an immersion) of inductive
systems of analytic groups

\begin{equation*}
 \widehat T\longrightarrow\widehat{A^\natural},
\end{equation*}

\noindent
which over $S_o$ reduces to (\SGARefLink{sga7.sga7i.eix.n13-3-7}{13.3.7}). By \SGARefLink{sga7.sga7i.eix.n13-2-2}{13.2.2}, this homomorphism comes
from a unique germ of homomorphisms

\begin{equation}\tag{\SGARefTarget{sga7.sga7i.eix.n13-3-9}{13.3.9}}
 T\longrightarrow A^\natural,
\end{equation}

\noindent
and, after taking $S$ small enough around $s$, we may suppose that it is
defined over all of $S$. The corresponding homomorphism

\begin{equation*}
 \underline{\operatorname{Lie}}(T)\longrightarrow
 \underline{\operatorname{Lie}}(A^\natural)\simeq\underline E
\end{equation*}

\noindent
is injective over each $S_n$, hence over $\widehat S$, and therefore over
$\underline{\mathcal O}_s$ by flat descent. Thus there is a neighbourhood
of $s$ in $S$, which we may take to be all of $S$, on which
$\underline{\operatorname{Lie}}(T)$ identifies with a locally direct
summand $\underline E^t$ of $\underline E$; denote the corresponding
vector subbundle of $E$ by $E^t$:

\begin{equation}\tag{\SGARefTarget{sga7.sga7i.eix.n13-3-10}{13.3.10}}
 \underline{\operatorname{Lie}}(T)\xrightarrow{\ \sim\ }
 \underline E^t\subset\underline E\simeq
 \underline{\operatorname{Lie}}(A^\natural)\simeq
 \underline{\operatorname{Lie}}(A),
 \qquad E^t\subset E.
\end{equation}

\noindent
Likewise, let

\begin{equation}\tag{\SGARefTarget{sga7.sga7i.eix.n13-3-11}{13.3.11}}
 R_o^t\subset R_o,
 \qquad R^t=R_{oS}^t\subset R_{oS}\simeq R^f
\end{equation}

\noindent
denote respectively the sublattice of $R_o$ defined by the subtorus
$T_o$ of $A_o$, and the corresponding constant analytic group over $S$.
Then the homomorphism (\SGARefLink{sga7.sga7i.eix.n13-3-9}{13.3.9}) corresponds simply to the homomorphism

\begin{equation*}
 E^t/R_{oS}^t=E^t/R^t\longrightarrow E/R_{oS}=E/R^f
\end{equation*}

\noindent
induced by the inclusions

\begin{equation*}
 E^t\hookrightarrow E,
 \qquad R^t\hookrightarrow R^f.
\end{equation*}

\noindent
Using the fact that (\SGARefLink{sga7.sga7i.eix.n13-3-9}{13.3.9}) induces an immersion on the fibre over $s$,
we conclude that the homomorphism

\begin{equation*}
 R_o/R_o^t\longrightarrow E_o/E_o^t=(E/E^t)_o
\end{equation*}

\noindent
is injective. It follows at once that

\begin{equation*}
 R_o/R_o^t\longrightarrow(E/E^t)_{s'}
\end{equation*}

\noindent
is still injective for $s'$ near $s$. Thus, after shrinking $S$ around
$s$, we may suppose that (\SGARefLink{sga7.sga7i.eix.n13-3-9}{13.3.9}) \underline{is a closed immersion}.

We are therefore led to introduce

\begin{equation}\tag{\SGARefTarget{sga7.sga7i.eix.n13-3-12}{13.3.12}}
 B=A^\natural/T\simeq(E/E^t)/(R^f/R^t),
\end{equation}

\noindent
so that $A^\natural$ acquires an extension structure

\begin{equation}\tag{\SGARefTarget{sga7.sga7i.eix.n13-3-13}{13.3.13}}
 0\longrightarrow T\longrightarrow A^\natural
 \longrightarrow B\longrightarrow0,
\end{equation}

\noindent
where $B$ is a relative analytic group over $S$, smooth and separated
over $S$, with connected fibres. By construction we therefore have

\begin{equation}\tag{\SGARefTarget{sga7.sga7i.eix.n13-3-14}{13.3.14}}
 \underline{\operatorname{Lie}}(B)\simeq
 \underline E/\underline E^t,
\end{equation}

\noindent
and the exponential exact sequence for $B$ is

\begin{equation}\tag{\SGARefTarget{sga7.sga7i.eix.n13-3-15}{13.3.15}}
 0\longrightarrow R^f/R^t\longrightarrow E/E^t
 \longrightarrow B\longrightarrow0.
\end{equation}

\noindent
Notice that

\begin{equation}\tag{\SGARefTarget{sga7.sga7i.eix.n13-3-16}{13.3.16}}
 R^f/R^t=(R_o/R_o^t)_S
\end{equation}

\noindent
is a constant analytic group over $S$. It follows from \SGARefLink{sga7.sga7i.eix.n13-1-5}{13.1.5} that $B$
is an abeloid group over $S$ in a neighbourhood of $s$ if and only if
its fibre at $s$,

\begin{equation*}
 B_o=A_o/T_o,
\end{equation*}

\noindent
is an abeloid analytic group; in that case $A_o$ appears as an extension
of such a group $B_o$ by a torus $T_o$. \underline{Every way of writing
$A_o$ as an extension of an abeloid group by a complex torus $(*)$
therefore gives rise on $A$, in a neighbourhood of $s$, to a unique
extension structure} (\SGARefLink{sga7.sga7i.eix.n13-3-13}{13.3.13}) \underline{of an abeloid group $B$ by a
torus $T$, inducing the given extension structure at $s$}.

% END INCLUDED SGA7I_English_current_buildable_source_20260823/components/179_expose_IX_section_13_3.tex

% BEGIN INCLUDED SGA7I_English_current_buildable_source_20260823/components/180_expose_IX_section_13_4.tex
\medskip
\noindent
\SGARefTarget{sga7.sga7i.eix.n13-4}{13.4}. \underline{Interpretation in terms of Deligne's ``mixed Hodge
structures''.}

Let us first return to the dictionary of \SGARefLink{sga7.sga7i.eix.n13-1-4}{13.1.4}, which we shall
reformulate in a slightly different way, better suited to introducing
the ``Hodge structures'' relevant to the present situation.  For this
purpose, let $\underline R$ be the sheaf of sections of $R$ over $S$, and
introduce the $\underline{\mathcal O}_S$-module

\begin{equation}\tag{\SGARefTarget{sga7.sga7i.eix.n13-4-1}{13.4.1}}
 \mathcal E(A)=\mathcal E=
 \underline R\otimes_{\mathbb Z}\underline{\mathcal O}_S\qquad .
\end{equation}

\noindent
It is locally free of finite type when $R$ is locally constant.  (Note
that $\mathcal E$ need not be coherent when $R$ is not locally
constant.)  The homomorphism (\SGARefLink{sga7.sga7i.eix.n13-1-3}{13.1.3})
$\underline R\longrightarrow\underline E$ then defines a canonical
homomorphism

\begin{equation}\tag{\SGARefTarget{sga7.sga7i.eix.n13-4-2}{13.4.2}}
 \mathcal E=\underline R\otimes_{\mathbb Z}
 \underline{\mathcal O}_S\longrightarrow\underline E\qquad .
\end{equation}

\noindent
We denote its kernel by $\underline F(A)$, or simply by $\underline F$;
thus there is an exact sequence

\begin{equation}\tag{\SGARefTarget{sga7.sga7i.eix.n13-4-3}{13.4.3}}
 0\longrightarrow\underline F\longrightarrow\mathcal E
 \longrightarrow\underline E\qquad .
\end{equation}

\noindent
The most interesting case is that in which (\SGARefLink{sga7.sga7i.eix.n13-4-2}{13.4.2}) is an epimorphism,
that is, in which $R_s$ spans the complex vector space $E_s$ for every
$s\in S$.  We then have an exact sequence

\begin{equation}\tag{\SGARefTarget{sga7.sga7i.eix.n13-4-4}{13.4.4}}
 0\longrightarrow\underline F\longrightarrow\mathcal E
 \longrightarrow\underline E\longrightarrow0\qquad .
\end{equation}

\noindent
If $R$ is locally constant, this is a short exact sequence of locally
free modules over $S$.  Knowledge of the pair $(E,R)$ is then equivalent
to knowledge of the pair $(R,\underline F)$, where $R$ is a locally
constant analytic group over $S$ whose fibres are free $\mathbb Z$-modules
of finite type, and $\underline F$ is a locally direct-summand submodule
of $\mathcal E=\underline R\otimes_{\mathbb Z}
\underline{\mathcal O}_S$, subject in addition to

\begin{equation}\tag{\SGARefTarget{sga7.sga7i.eix.n13-4-5}{13.4.5}}
 \underline F\cap\underline R=0\qquad .
\end{equation}

\noindent
From the standpoint of Hodge structures, $\mathcal E$ should be regarded
as filtered by the ``\underline{Hodge filtration}''

\begin{equation}\tag{\SGARefTarget{sga7.sga7i.eix.n13-4-6}{13.4.6}}
 \operatorname{Fil}^{1}(\mathcal E)=\mathcal E,
 \qquad \operatorname{Fil}^{0}(\mathcal E)=\underline F,
 \qquad \operatorname{Fil}^{-1}(\mathcal E)=0\qquad .
\end{equation}

\noindent
Notice that the exact sequence (\SGARefLink{sga7.sga7i.eix.n13-4-3}{13.4.3}) still depends functorially on the
analytic group $A$ over $S$, and commutes with arbitrary base change.

Returning now to the conditions of \SGARefLink{sga7.sga7i.eix.n13-3}{13.3}, put

\begin{equation}\tag{\SGARefTarget{sga7.sga7i.eix.n13-4-7}{13.4.7}}
 \mathcal E^f=\underline R^f\otimes_{\mathbb Z}
 \underline{\mathcal O}_S,
 \qquad
 \mathcal E^t=\underline R^t\otimes_{\mathbb Z}
 \underline{\mathcal O}_S\qquad .
\end{equation}

\noindent
These modules are defined on a suitable neighbourhood of $s$, which we
may suppose to be all of $S$, and define a filtration of $\mathcal E$,
called the \underline{weight filtration},

\begin{equation}\tag{\SGARefTarget{sga7.sga7i.eix.n13-4-8}{13.4.8}}
 \mathcal E\supset\mathcal E^f\supset\mathcal E^t\supset0\qquad .
\end{equation}

\noindent
Here $\mathcal E^t$, $\mathcal E^f$, and $\mathcal E$ are to be regarded
as having weights $-2$, $-1$, and $0$, respectively; thus this is an
increasing filtration.  Like $R^t$, the module $\mathcal E^t$ depends on
the choice of a subtorus $T_o$ of $A_o=A_s$.  Since $R^f$ and $R^t$ are
locally constant, \underline{$\mathcal E^f$ and $\mathcal E^t$ are
locally free}; they may also be interpreted as

\begin{equation}\tag{\SGARefTarget{sga7.sga7i.eix.n13-4-9}{13.4.9}}
 \mathcal E^f=\mathcal E(A^\natural),
 \qquad \mathcal E^t=\mathcal E(T)\qquad .
\end{equation}

\noindent
Applying the functoriality of (\SGARefLink{sga7.sga7i.eix.n13-4-3}{13.4.3}), as $A$ varies, to the
homomorphisms

\begin{equation*}
 T\longrightarrow A^\natural\longrightarrow A
\end{equation*}

\noindent
shows at once that $\underline F(A^\natural)$ and $\underline F(T)$
identify respectively with

\begin{equation}\tag{\SGARefTarget{sga7.sga7i.eix.n13-4-10}{13.4.10}}
 \underline F(A^\natural)\simeq
 \underline F^f\stackrel{\mathrm{def}}{=}
 \underline F\cap\mathcal E^f,
 \qquad
 \underline F(T)=\underline F^t=
 \underline F\cap\mathcal E^t=0\qquad .
\end{equation}

\noindent
The last equality follows immediately from the fact that
$\underline F(T)=0$ for a torus $T$: in this case the exact sequence of
type (\SGARefLink{sga7.sga7i.eix.n13-4-3}{13.4.3}) is
$0\longrightarrow0\longrightarrow\mathcal E(T)
\longrightarrow\underline{\operatorname{Lie}}(T)\longrightarrow0$.
At the same time we see that the composite homomorphism

\begin{equation*}
 \mathcal E^t\longrightarrow\mathcal E\longrightarrow\underline E
\end{equation*}

\noindent
induces an isomorphism

\begin{equation}\tag{\SGARefTarget{sga7.sga7i.eix.n13-4-11}{13.4.11}}
 \mathcal E^t\stackrel{\sim}{\longrightarrow}\underline E^t
 \simeq\underline{\operatorname{Lie}}(T)\qquad .
\end{equation}

\noindent
\underline{Suppose now that $R_s$ spans the complex vector space $E_s$},
as happens in particular in the favourable case \SGARefLink{sga7.sga7i.eix.n13-3-16}{13.3.16}, where $A_s$ is
an extension of an abeloid group by a complex torus $T_o$.  After
restricting $S$ to a suitable neighbourhood of $s$, the canonical
homomorphism of type (\SGARefLink{sga7.sga7i.eix.n13-4-2}{13.4.2}), with $A$ replaced by $A^\natural$,

\begin{equation*}
 \mathcal E^f=\mathcal E(A^\natural)
 \longrightarrow\underline{\operatorname{Lie}}(A^\natural)
 \simeq\underline E
\end{equation*}

\noindent
is an epimorphism; a fortiori, so is (\SGARefLink{sga7.sga7i.eix.n13-4-2}{13.4.2}).  We therefore have the
exact sequences (\SGARefLink{sga7.sga7i.eix.n13-4-4}{13.4.4}) and

\begin{equation}\tag{\SGARefTarget{sga7.sga7i.eix.n13-4-12}{13.4.12}}
 0\longrightarrow\underline F^f\longrightarrow\mathcal E^f
 \longrightarrow\underline E\longrightarrow0\qquad .
\end{equation}

\noindent
In particular,

\begin{equation}\tag{\SGARefTarget{sga7.sga7i.eix.n13-4-13}{13.4.13}}
 \underline F+\mathcal E^f=\mathcal E\qquad ,
\end{equation}

\noindent
since $\mathcal E/\underline F\simeq\underline E$.  Moreover, since
$B=A^\natural/T$ is a quotient of $A^\natural$,
\footnote{The printed source says here that $B$ is a quotient of $A$;
the displayed definition $B=A^\natural/T$ gives the intended group.}
it satisfies the analogous surjectivity condition, and hence gives an
exact sequence of type (\SGARefLink{sga7.sga7i.eix.n13-4-4}{13.4.4}):

\begin{equation}\tag{\SGARefTarget{sga7.sga7i.eix.n13-4-14}{13.4.14}}
 0\longrightarrow\underline F(B)\longrightarrow\mathcal E(B)
 \longrightarrow\underline{\operatorname{Lie}}(B)
 \longrightarrow0\qquad .
\end{equation}

\noindent
Its interpretation in terms of the modules already introduced is
immediate.  Indeed, by (\SGARefLink{sga7.sga7i.eix.n13-3-12}{13.3.12}),

\begin{equation}\tag{\SGARefTarget{sga7.sga7i.eix.n13-4-15}{13.4.15}}
 \left\{
 \begin{array}{l}
  \mathcal E(B)\simeq
  (\underline R^f/\underline R^t)\otimes_{\mathbb Z}
  \underline{\mathcal O}_S
  \simeq\mathcal E^f/\mathcal E^t,\\[2ex]
  \underline{\operatorname{Lie}}(B)
  \simeq\underline E/\underline E^t.
 \end{array}
 \right.
\end{equation}

\noindent
It follows from (\SGARefLink{sga7.sga7i.eix.n13-4-11}{13.4.11}) and the relation
$\underline F\cap\mathcal E^t=0$ in (\SGARefLink{sga7.sga7i.eix.n13-4-10}{13.4.10}) that
$\underline F(B)=\ker(\mathcal E^f/\mathcal E^t
\longrightarrow\underline E/\underline E^t)$ is isomorphic to
$\underline F^f$ through the homomorphism

\begin{equation}\tag{\SGARefTarget{sga7.sga7i.eix.n13-4-16}{13.4.16}}
 \underline F^f\simeq\underline F(A^\natural)
 \stackrel{\sim}{\longrightarrow}\underline F(B)\qquad .
\end{equation}

\noindent
Let us summarize the properties just obtained for the weight filtration
(\SGARefLink{sga7.sga7i.eix.n13-4-8}{13.4.8}), and its relation with the Hodge filtration (\SGARefLink{sga7.sga7i.eix.n13-4-6}{13.4.6}) of
$\mathcal E$ defined by the submodule $\underline F$, in the favourable
case \SGARefLink{sga7.sga7i.eix.n13-3-16}{13.3.16} where $A_s$ is an extension of an abeloid analytic group
$B_o$ by a torus $T_o$.  We use this torus to define $R^t$, so that $B$
is an abeloid analytic group over $S$.

\medskip
\noindent
\underline{Proposition \SGARefTarget{sga7.sga7i.eix.n13-4-17}{13.4.17}.} \underline{Under the preceding
conditions, the Hodge and weight filtrations on}
$\mathcal E=\underline R\otimes_{\mathbb Z}\underline{\mathcal O}_S$,

\begin{equation}\tag{\SGARefTarget{sga7.sga7i.eix.n13-4-17-1}{13.4.17.1}}
 \mathcal E\supset\underline F,
 \qquad
 \mathcal E\supset\mathcal E^f\supset\mathcal E^t\supset0,
\end{equation}

\noindent
\underline{satisfy the following conditions:}

\smallskip
\noindent
a) \underline{The weight filtration is obtained, by tensoring with
$\underline{\mathcal O}_S$ over $\mathbb Z$, from the analogous
filtration on $\underline R$}

\begin{equation*}
 \underline R\supset\underline R^f\supset\underline R^t\supset0,
\end{equation*}

\noindent
\underline{where $\underline R^f$ and $\underline R^t$ are locally
constant analytic groups over $S$, with values free $\mathbb Z$-modules
of finite type}, namely
$R_o^f=R_s=\pi_1(A_s)$ and $R_o^t=\pi_1(T_o)$.

\smallskip
\noindent
b) \underline{One has the relations}

\begin{equation}\tag{*}
 \underline F\cap\mathcal E^t=0,
 \qquad \underline F+\mathcal E^f=\mathcal E\qquad .
\end{equation}

\noindent
\underline{In other words, the filtration on $\mathcal E^t$ induced by
the Hodge filtration}, with the degree conventions of (\SGARefLink{sga7.sga7i.eix.n13-4-6}{13.4.6}),
\underline{is concentrated in degree $-1$}, whereas the filtration on
$\mathcal E/\mathcal E^f$ induced by the Hodge filtration is
\underline{concentrated in degree $0$}.

\smallskip
\noindent
c) \underline{Consider on}

\begin{equation*}
 \mathcal E^{\mathrm{ab}}
 \stackrel{\mathrm{def}}{=}\mathcal E^f/\mathcal E^t
 \simeq\underline R^{\mathrm{ab}}\otimes_{\mathbb Z}
 \underline{\mathcal O}_S,
 \qquad
 \underline R^{\mathrm{ab}}=\underline R^f/\underline R^t
 \quad\text{(cf. a))},
\end{equation*}

\noindent
\underline{the filtration induced by the Hodge filtration on
$\mathcal E$}.  It is therefore defined by the submodule
$\underline F^{\mathrm{ab}}$ of $\mathcal E^{\mathrm{ab}}$ which is the
image of

\begin{equation*}
 \underline F^f\stackrel{\mathrm{def}}{=}
 \underline F\cap\mathcal E^f\longrightarrow\mathcal E^{\mathrm{ab}}.
\end{equation*}

\noindent
The arrow is a monomorphism by the first relation $(*)$.
\underline{Then the pair}

\begin{equation*}
 \bigl(\underline R^{\mathrm{ab}},
 \underline F^{\mathrm{ab}}\subset
 \underline R^{\mathrm{ab}}\otimes_{\mathbb Z}
 \underline{\mathcal O}_S=\mathcal E^{\mathrm{ab}}\bigr)
\end{equation*}

\noindent
\underline{is isomorphic to the pair}
$(\underline R(B),\underline F(B))$
\underline{associated with an abeloid group over $S$}, as explained
after (\SGARefLink{sga7.sga7i.eix.n13-4-4}{13.4.4}).  Thus

\begin{equation*}
 B\simeq\mathcal E^{\mathrm{ab}}/
 (F^{\mathrm{ab}}+R^{\mathrm{ab}}),
\end{equation*}

\noindent
where the symbols without underlining denote the corresponding analytic
groups over $S$.

\medskip
\noindent
\underline{Remark \SGARefTarget{sga7.sga7i.eix.n13-4-18}{13.4.18}.} In Deligne's terminology [4], the conclusion
of the preceding proposition, after restriction to an open subset $U$
of $S$ on which $R$ is locally constant, says that $R$, endowed with the
two filtrations (\SGARefLink{sga7.sga7i.eix.n13-4-17-1}{13.4.17.1}) on
$\mathcal E=\underline R\otimes_{\mathbb Z}\underline{\mathcal O}_S$,
is an ``analytic family of mixed Hodge structures over $U$''.  The
situation just described served as the model for Deligne's conjectures
on N\'eron theory in higher dimensions, stated in [9, 9.8--9.13].

There is, however, an important feature of the algebraic theory studied
in \S~7 which is missing from the description in \SGARefLink{sga7.sga7i.eix.n13-4-17}{13.4.17}.  Suppose that
$A$ is an abeloid group over an open subset $U$ complementary to a
nowhere-dense analytic subset $Z$ of $S$---say $Z=\{s\}$ when
$\dim S=1$---so that, by (\SGARefLink{sga7.sga7i.eix.n13-1-5}{13.1.5}), $R|U$ is a local system over $U$.
Then the local system $R/R^t$ over $U$ ought to be \underline{constant};
more precisely, it ought to be dual to the local system ${R'}^f$, where
$R'$ is the analytic group, \`etale over $S$, which would correspond
(as $R$ corresponds to $A$) to a suitable analytic group $A'$ over $S$
such that $A'|U$ is the abeloid group dual to $A|U$.  In the notation of
\S~5, the algebraic analogue is the duality between
$T_\ell(A_K)/T_\ell(A_K)^t$ and $T_\ell(A'_K)^f$ in (\SGARefLink{sga7.sga7i.eix.n5-5-9}{5.5.9}), which
follows from the duality theorem \SGARefLink{sga7.sga7i.eix.n5-2}{5.2}.

I do not know whether such a statement holds under only the hypotheses
of \SGARefLink{sga7.sga7i.eix.n13-4-17}{13.4.17}, or even whether the action of $\pi_1(U,t)$ on $R_t$ is
unipotent for $t\in U$.  Nor do the naive methods used in the present
paragraph seem capable of proving such a result, even under the
additional hypothesis that $A|U$ is endowed with a ``polarization'',
that is, with an invertible sheaf ample on every fibre.  This is,
however, something that the ``semi-simple group people'' apparently know
how to prove, at least when $S$ is nonsingular of dimension $1$; see the
report [9].  In the next subsection we shall show that everything works
under relative algebraicity hypotheses on $A$, and at the same time
relate the result to the construction of \S~7.

% END INCLUDED SGA7I_English_current_buildable_source_20260823/components/180_expose_IX_section_13_4.tex

% BEGIN INCLUDED SGA7I_English_current_buildable_source_20260823/components/181_expose_IX_section_13_5.tex
\medskip
\noindent
\SGARefTarget{sga7.sga7i.eix.n13-5}{13.5}. \underline{The case of a relatively algebraic situation.}
We now suppose that $A$ comes from a relative group scheme over $S$
which is smooth and separated, and whose fibres are extensions of
abelian schemes by tori.  We denote this relative scheme by the same
letter $A$.  (We allow ourselves this abuse of notation because the
functor $G\longmapsto G^{\mathrm{an}}$, from the category of relative
group schemes of the type under consideration to the category of
analytic groups over $S$, is fully faithful, as would be easy to prove
using results of relative GAGA type [10].)  For everything concerning
the notion of a relative scheme and the relations between relative
schemes over $S$ and analytic spaces over $S$, we refer to [10] (see
also [2, Exp. 15]).

If $\underline O$ denotes the local ring of $S$ at $s$, then the
relative scheme $A$ is determined near $s$ by the ordinary scheme
$A_{\underline O}$ over $\operatorname{Spec}(\underline O)$ obtained
from $A$ by localization at $s$.  It is then clear that the formal
completion of $A_{\underline O}$ along its special fibre (isomorphic to
$A_s=A_o$) is isomorphic to the formal completion of the analytic group
$A$ along $A_s$ (\SGARefLink{sga7.sga7i.eix.n13-2}{13.2}), where, by the abuse of notation introduced
above, the $A_n=A\times_S S_n$ are identified with algebraic schemes
over $S$.  In the favourable case in which $\hat A$, viewed as a formal
group scheme over $\hat{\underline O}$, is algebraizable (\SGARefLink{sga7.sga7i.eix.n7-2-1}{7.2.1}), as is
the case in particular when $S$ is normal at $s$ and $A$ is
quasi-projective over $S$ (\SGARefLink{sga7.sga7i.eix.n7-1-5}{7.1.5}), we therefore obtain a group scheme
$A^{\natural}_{\hat{\underline O}}$ over
$\operatorname{Spec}(\hat{\underline O})$, which is an extension of an
abelian scheme $B_{\hat{\underline O}}$ by a torus.  For simplicity, we
suppose that $B_{\hat{\underline O}}$ is projective over
$\hat{\underline O}$ (as is indeed the case, we know, if
$\underline O$, and hence $\hat{\underline O}$, is normal).  We shall
then prove:

\medskip
\noindent
\underline{Proposition \SGARefTarget{sga7.sga7i.eix.n13-5-1}{13.5.1}.} \underline{Under the preceding
conditions, the analytic group $A^{\natural}$ of (\SGARefLink{sga7.sga7i.eix.n13-3-3}{13.3.3}) comes, in a
neighbourhood of $s$, from a relative group scheme determined up to a
unique isomorphism, which is an extension of a projective relative
abelian scheme over $S$ by a torus.  If
$A^{\natural}_{\underline O}$ denotes the group scheme it defines over
$\underline O=\underline O_{S,s}$, then the group scheme}

\begin{equation*}
 A^{\natural}_{\underline O}\otimes_{\underline O}
 \hat{\underline O}
\end{equation*}

\noindent
\underline{is canonically isomorphic, with its extension structure, to
the Raynaud extension $A^{\natural}_{\hat{\underline O}}$ introduced
above.}

The latter assertion follows immediately from the first, in view of
\SGARefLink{sga7.sga7i.eix.n7-2-1}{7.2.1}, since, by (\SGARefLink{sga7.sga7i.eix.n13-3-6}{13.3.6}), the formal completion of
$A^{\natural}_{\underline O}\otimes_{\underline O}\hat{\underline O}$
along its special fibre is $\hat A_{\underline O}$.  To prove the
algebraicity of $A^{\natural}$ near $s$, proceeding as in \SGARefLink{sga7.sga7i.eix.n7-2-1}{7.2.1}, we are
immediately reduced to proving that of the abelian quotient $B$ of
$A^{\natural}$, which follows from the next lemma.

\medskip
\noindent
\underline{Lemma \SGARefTarget{sga7.sga7i.eix.n13-5-2}{13.5.2}.} \underline{Let $B$ be an abeloid analytic
group over the analytic space $S$, and let $s$ be a point of $S$ such
that the completion $\hat B$ along the fibre $B_s$ is algebraizable;
that is, the $B_n=B\times_S S_n$ are algebraizable and $\hat B$, now
viewed as a formal scheme over $\hat{\underline O}_{S,s}$, is
algebraizable, or equivalently comes from an abelian scheme
$B_{\hat{\underline O}}$ over $\hat{\underline O}$.  Suppose moreover
that $B_{\hat{\underline O}}$ is projective over
$\hat{\underline O}$.  Then $B$ is projective over $S$ in a
neighbourhood of $s$.}

The hypothesis means that there is a polarization of the formal
analytic space $\hat B$, that is, a sequence of polarizations $p_n$ of
the $B_n$ which ``glue together''.  The conclusion means that there is
a polarization $p$ of $B$ over a neighbourhood of $s$.  Let $P$ be the
analytic space over $S$ of relative polarizations of $B$ over $S$,
which represents the functor

\begin{equation*}
 S'\longmapsto
 \{\text{polarizations of $B_{S'}$ over $S'$}\}
\end{equation*}

\noindent
on the category of analytic spaces $S'$ over $S$.  This is an open
subfunctor of the N\'eron--Severi functor $\underline{\mathrm{NS}}_{B/S}$
over $S$.  The simplest standard arguments, which we omit here, show
that $\underline{\mathrm{NS}}_{B/S}$, and hence also $P$, is represented
by an \underline{unramified} analytic space over $S$.  The hypothesis
then says that there is a \underline{formal} section $p$ of $P$ over
$S$ at $s$.  Since $P$ is unramified and therefore locally quasi-finite
over $S$, it follows from [2, Exp. 18, no.~1] that every formal section
of $P$ over $S$ at $s$ comes from an analytic section near $s$, that is,
from a polarization of $B$ over a neighbourhood of $s$.  This proves
the assertion.

\medskip
\noindent
\underline{Remarks \SGARefTarget{sga7.sga7i.eix.n13-5-3}{13.5.3}.}

\smallskip
\noindent
a) If in \SGARefLink{sga7.sga7i.eix.n13-5-2}{13.5.2} we omit the hypothesis that
$B_{\hat{\underline O}}$ is projective over $\hat{\underline O}$,
\footnote{The printed source says ``projective over $S$'' here; the
hypothesis of Lemma \SGARefLink{sga7.sga7i.eix.n13-5-2}{13.5.2} and the ensuing normality argument require
``projective over $\hat{\underline O}$''.}
we can still conclude that $B$ is algebraizable near $s$.  For this one
uses an unpublished result (due independently to M.~Artin and
P.~Deligne, using the theory of the sheaves $\underline T^i$ of VI 2)
which asserts that, if $X$ is proper over $S$ and $X_{\mathrm{red}}$ is
algebraizable, then $X$ is algebraizable.  This reduces us to the case
where $S$ is reduced.  On the other hand, when $S$ is normal, so that
$\hat{\underline O}$ is normal, every abelian scheme over
$\hat{\underline O}$ is known to be projective over
$\hat{\underline O}$ (see, for example, [23]), and 13.5.2 applies.  To
obtain the case of an arbitrary reduced $S$, one uses non-flat descent
from the normalization $S'$ of $S$, by means of the following
(delicate) result, which we shall not prove here: a proper morphism of
noetherian schemes $f:S'\longrightarrow S$ which is an epimorphism
(that is, such that
$\underline O_S\longrightarrow f_*(\underline O_{S'})$ is injective)
is an \underline{effective descent morphism} for the fibred category of
relative abelian schemes.

\smallskip
\noindent
b) It appears that the proof of \SGARefLink{sga7.sga7i.eix.n13-5-2}{13.5.2} should also prove the analogous
result for every analytic space $B$ proper over $S$.  Probably more
delicate is the question whether an analytic space $X$ proper over $S$
whose formal completion $\hat X$ along $X_s$ is algebraizable is itself
algebraizable over $S$ near $s$.  We merely point out that here the
notion of relative algebraizability of $X$ over $S$ should be taken in
a broader sense than in [10]: near $s$, $X$ should be required to come
from an \underline{algebraic space} over $\underline O_{S,s}$ in the
sense of M.~Artin [1].

\medskip
\noindent
\SGARefTarget{sga7.sga7i.eix.n13-5-4}{13.5.4}. Let us place ourselves under the conditions of \SGARefLink{sga7.sga7i.eix.n13-5-1}{13.5.1}.  For
every integer $n>0$, choose also the standard isomorphism of groups over
$\mathbb C$

\begin{equation*}
 \mu_{n\mathbb C}\simeq(\mathbb Z/n\mathbb Z)_{\mathbb C},
 \qquad m\longmapsto\exp(2\pi i m/n)\qquad .
\end{equation*}

\noindent
By passing to the inverse limit over the powers of a fixed prime
$\ell$, this gives a canonical isomorphism of $\ell$-adic sheaves over
$\mathbb C$,

\begin{equation*}
 \mathbb Z_\ell(1)_{\mathbb C}\simeq\mathbb Z_{\ell\mathbb C}
 \qquad ,
\end{equation*}

\noindent
which we shall use to identify the two sides.

Applying the exact sequence of
$\underline{\operatorname{Ext}}^i_{\mathbb Z}
(\mathbb Z/n\mathbb Z,-)$ to the exact sequence (\SGARefLink{sga7.sga7i.eix.n13-1-3}{13.1.3}), we obtain a
canonical isomorphism

\begin{equation}\tag{\SGARefTarget{sga7.sga7i.eix.n13-5-5}{13.5.5}}
 {}_nA^{\mathrm o}\simeq
 \underline R\otimes_{\mathbb Z}\mathbb Z/n\mathbb Z
 \qquad ,
\end{equation}

\noindent
which, as $n$ runs through the powers of a given prime $\ell$, gives an
isomorphism of inverse systems of groups over $S$,

\begin{equation}\tag{\SGARefTarget{sga7.sga7i.eix.n13-5-6}{13.5.6}}
 T_\ell(A^{\mathrm o})\simeq
 \underline R\otimes_{\mathbb Z}\mathbb Z_\ell
 \qquad ,
\end{equation}

\noindent
where $\mathbb Z_\ell$ is regarded as the inverse system of the
$\mathbb Z/\ell^{n+1}\mathbb Z$.  In the situation considered here, in
which $A$ comes from a separated group scheme over $S$, it is then
immediately clear that the
$\underline R\otimes\mathbb Z/n\mathbb Z$ may be regarded as relative
schemes.  This gives a meaning to the symbol
$(\underline R\otimes_{\mathbb Z}\mathbb Z_\ell)_{\underline O}$
(whereas $R_{\underline O}$ would in general have none!), and yields a
canonical isomorphism

\begin{equation}\tag{\SGARefTarget{sga7.sga7i.eix.n13-5-7}{13.5.7}}
 T_\ell(A_{\underline O})\simeq
 (\underline R\otimes_{\mathbb Z}\mathbb Z_\ell)_{\underline O}
 \qquad .
\end{equation}

\noindent
It follows immediately from the constructions of paragraph 2 and \SGARefLink{sga7.sga7i.eix.n13-3}{13.3}
that the preceding isomorphism carries the fixed part to the fixed part
and the toric part to the toric part:

\begin{equation}\tag{\SGARefTarget{sga7.sga7i.eix.n13-5-8}{13.5.8}}
 T_\ell(A_{\underline O})^f\simeq
 (\underline R^f\otimes_{\mathbb Z}\mathbb Z_\ell)_{\underline O},
 \qquad
 T_\ell(A_{\underline O})^t\simeq
 (\underline R^t\otimes_{\mathbb Z}\mathbb Z_\ell)_{\underline O}
 \qquad .
\end{equation}

\noindent
If $A'$ is a relative analytic group over $S$ satisfying the same
hypotheses as $A$, and if $W_{\underline O}$ is a biextension of
$(A_{\underline O},A'_{\underline O})$ by
$\mathbb G_{m\underline O}$, hence defines over a neighbourhood of $s$
a biextension of $(A,A')$ by $\mathbb G_{mS}$ (now using the notion in
the complex-analytic sense, obtained by working in the site of analytic
spaces over $S$ for a suitable topology), then, by VIII \SGARefLink{sga7.sga7i.eviii.n2-2}{2.2}, $W$
defines pairings

\begin{equation}\tag{\SGARefTarget{sga7.sga7i.eix.n13-5-9}{13.5.9}}
 (\underline R\otimes\mathbb Z_\ell)\times
 (\underline R'\otimes_{\mathbb Z}\mathbb Z_\ell)
 \simeq T_\ell(A)\times T_\ell(A')
 \longrightarrow\mathbb Z_\ell(1)_S
 \simeq\mathbb Z_{\ell S}\qquad ,
\end{equation}

\noindent
and it is again clear that the corresponding pairing

\begin{equation}\tag{\SGARefTarget{sga7.sga7i.eix.n13-5-10}{13.5.10}}
 T_\ell(A_{\underline O})\times T_\ell(A'_{\underline O})
 \longrightarrow\mathbb Z_\ell(1)_S
\end{equation}

\noindent
is precisely the one associated, by loc. cit., with the biextension
$W_{\underline O}$ of $(A_{\underline O},A'_{\underline O})$.  It is
also not difficult to see that, after restricting to a neighbourhood
of $s$, the pairings (\SGARefLink{sga7.sga7i.eix.n13-5-9}{13.5.9}) arise from a unique pairing (no longer
depending on $\ell$)

\begin{equation}\tag{\SGARefTarget{sga7.sga7i.eix.n13-5-11}{13.5.11}}
 \underline R\otimes\underline R'\longrightarrow\mathbb Z
 \qquad .
\end{equation}

\noindent
With these preliminaries, the algebraic results concerning the
$T_\ell(A_{\underline O})$ can be used to draw conclusions about the
sheaf $\underline R$.  Thus the pairing (\SGARefLink{sga7.sga7i.eix.n13-5-11}{13.5.11}) is zero on
$\underline R^t\times\underline R'{}^f$ and on
$\underline R^f\times\underline R'{}^t$ (\SGARefLink{sga7.sga7i.eix.n7-4-8}{7.4.8}).  Moreover, if $S$ is
nonsingular of dimension $1$ and, over $U=S-\{s\}$, $A$ and $A'$ are
relative abelian schemes dual to one another through $W$, then
$\underline R^t_U$ and $\underline R'{}^f_U$, as well as
$\underline R^f_U$ and $\underline R'{}^t_U$, are exact annihilators of
one another over a sufficiently small neighbourhood of $s$ (\SGARefLink{sga7.sga7i.eix.n7-4-8}{7.4.8}).
Since (\SGARefLink{sga7.sga7i.eix.n13-5-11}{13.5.11}) is a perfect duality over $U$, it follows that there is
a duality over $U$ between $\underline R/\underline R^t$ and
$\underline R'{}^f$.  In particular,
$\underline R/\underline R^t$ is constant, and a fortiori the action of
$\pi_1(U,t)$ on $\underline R_t$ is unipotent of level $2$.  Similarly,
over $U$ there is a perfect duality between
$\underline R/\underline R^f$ and $\underline R'{}^t$.  Since the
latter is also isomorphic to the dual torus of the lattice $M_S$ over
$S$, where $M$ is the group defined in (\SGARefLink{sga7.sga7i.eix.n5-5-4}{5.5.4}), we obtain a canonical
isomorphism

\begin{equation}\tag{\SGARefTarget{sga7.sga7i.eix.n13-5-12}{13.5.12}}
 \underline R_U/\underline R^f_U\simeq M_U
 \qquad ,
\end{equation}

\noindent
and the homomorphism $i$ of (\SGARefLink{sga7.sga7i.eix.n13-3-4}{13.3.4}) may be written as a homomorphism

\begin{equation}\tag{\SGARefTarget{sga7.sga7i.eix.n13-5-13}{13.5.13}}
 i:M_U\longrightarrow A^\natural
 \qquad ,
\end{equation}

\noindent
which we shall encounter again in a slightly different context in the
next paragraph.

% END INCLUDED SGA7I_English_current_buildable_source_20260823/components/181_expose_IX_section_13_5.tex

% BEGIN INCLUDED SGA7I_English_current_buildable_source_20260823/components/182_expose_IX_section_14_1.tex
\medskip
\noindent
14. \underline{Links with transcendental theory: the rigid-analytic
case.  The canonical homomorphism}
$\underline M_K\longrightarrow A_K^{\natural\mathrm o}$.

\medskip
\noindent
\SGARefTarget{sga7.sga7i.eix.n14-1}{14.1}. We assume here the foundations of the theory of rigid-analytic
spaces, for which one may consult [12], [13], and [34].

Let $S$ be a complete trait, and let $A_K$ be an abelian scheme over its
field of fractions.  In \SGARefLink{sga7.sga7i.eix.n8-1}{8.1}, by descent from the Raynaud extension in
the semistable case, we constructed an extension
$G_K=A_K^{\natural\mathrm o}$ (\SGARefLink{sga7.sga7i.eix.n8-1-1}{8.1.1})

\begin{equation}\tag{\SGARefTarget{sga7.sga7i.eix.n14-1-1}{14.1.1}}
 0\longrightarrow T_K\longrightarrow G_K\longrightarrow B_K
 \longrightarrow0
\end{equation}

\noindent
of an abelian scheme $B_K$ by a torus $T_K$.  Considering the dual
abelian scheme $A'_K$, we likewise have an extension
$G'_K=A'{}^{\natural\mathrm o}_K$

\begin{equation}\tag{\SGARefTarget{sga7.sga7i.eix.n14-1-2}{14.1.2}}
 0\longrightarrow T'_K\longrightarrow G'_K\longrightarrow B'_K
 \longrightarrow0\qquad ,
\end{equation}

\noindent
where $B'_K$ is moreover identified with the abelian scheme dual to
$B_K$ (\SGARefLink{sga7.sga7i.eix.n8-2}{8.2}).  As in (\SGARefLink{sga7.sga7i.eix.n8-1-8}{8.1.8}), let $\underline M_K$ and
$\underline M'_K$ denote the lattices over $K$ dual to $T'_K$ and
$T_K$, respectively:

\begin{equation}\tag{\SGARefTarget{sga7.sga7i.eix.n14-1-3}{14.1.3}}
 \underline M_K=D(T'_K),\qquad \underline M'_K=D(T_K)\qquad .
\end{equation}

\noindent
Thus the extension structures (\SGARefLink{sga7.sga7i.eix.n14-1-2}{14.1.2}) and (\SGARefLink{sga7.sga7i.eix.n14-1-1}{14.1.1}) may respectively be
interpreted (VIII \SGARefLink{sga7.sga7i.eviii.n3-7}{3.7}) as equivalent to the data of homomorphisms

\begin{equation}\tag{\SGARefTarget{sga7.sga7i.eix.n14-1-4}{14.1.4}}
 \underline M_K\xrightarrow{\ p_K\ }B_K,\qquad
 \underline M'_K\xrightarrow{\ p'_K\ }B'_K\qquad .
\end{equation}

Let $W_K^{\natural}$ be the biextension of $(B_K,B'_K)$ by
$\mathbb G_{mK}$ which defines the canonical duality between these
abelian schemes (\SGARefLink{sga7.sga7i.eix.n8-2}{8.2}), and let $E_K$ be the biextension of
$(\underline M_K,\underline M'_K)$ by $\mathbb G_{mK}$ obtained by
pulling $W_K^{\natural}$ back along the homomorphisms (\SGARefLink{sga7.sga7i.eix.n14-1-4}{14.1.4}).  One
checks immediately that trivializations of $E_K$ correspond
canonically and bijectively to homomorphisms

\begin{equation}\tag{\SGARefTarget{sga7.sga7i.eix.n14-1-5}{14.1.5}}
 i:\underline M_K\longrightarrow G_K
\end{equation}

\noindent
lifting $p_K:\underline M_K\longrightarrow B_K$, or, symmetrically, to
homomorphisms

\begin{equation}\tag{\SGARefTarget{sga7.sga7i.eix.n14-1-6}{14.1.6}}
 i':\underline M'_K\longrightarrow G'_K
\end{equation}

\noindent
lifting $p'_K:\underline M'_K\longrightarrow B'_K$.

With this understood, write a superscript ${}^{\mathrm{an}}$ for the
rigid-analytic groups over $K$ associated with the algebraic groups
under consideration [13].  In [24] one defines a canonical exact
sequence

\begin{equation}\tag{\SGARefTarget{sga7.sga7i.eix.n14-1-7}{14.1.7}}
 0\longrightarrow\underline M_K^{\mathrm{an}}
 \xrightarrow{\ i^{\mathrm{an}}\ }G_K^{\mathrm{an}}
 \xrightarrow{\ e\ }A_K^{\mathrm{an}}\longrightarrow0\qquad .
\end{equation}

\noindent
Here $i^{\mathrm{an}}$, which may also be interpreted as a homomorphism
$i:\underline M_K\longrightarrow G_K$ of ordinary group schemes,
\underline{lifts the homomorphism} $p$ in (\SGARefLink{sga7.sga7i.eix.n14-1-4}{14.1.4}), that is, it is of
the form considered in (\SGARefLink{sga7.sga7i.eix.n14-1-5}{14.1.5}).  This exact sequence is the
rigid-analytic analogue of the exact sequence (\SGARefLink{sga7.sga7i.eix.n13-3-4}{13.3.4}).  It depends
functorially on $A_K$, and its formation is compatible with every
change of traits $S'\longrightarrow S$.

The latter fact also shows that, in order to define and study (\SGARefLink{sga7.sga7i.eix.n14-1-7}{14.1.7}),
finite descent reduces us to the case where $A_K$ has semistable
reduction.  In this case the group schemes
$T_K,G_K,B_K,\underline M_K,\ldots$ are the generic fibres of group
schemes $T,G,B,M,\ldots$ over $S$, with $G$ an extension of the abelian
scheme $B$ by the torus $T$ (paragraph 7).  We may then consider the
formal completion $\hat G$ of $G$ along its special fibre, and by
construction there is a canonical isomorphism (\SGARefLink{sga7.sga7i.eix.n7-2-6}{7.2.6})

\begin{equation}\tag{\SGARefTarget{sga7.sga7i.eix.n14-1-8}{14.1.8}}
 \hat G\xrightarrow{\ \sim\ }\widehat{A^{\mathrm o}}
 \qquad .
\end{equation}

\noindent
Passing to the ``generic fibres'', which are rigid-analytic groups over
$K$, gives a canonical isomorphism

\begin{equation}\tag{\SGARefTarget{sga7.sga7i.eix.n14-1-9}{14.1.9}}
 \hat G^{\mathrm{an}}\xrightarrow{\ \sim\ }
 \widehat{A^{\mathrm o}}\,{}^{\mathrm{an}}
 \qquad .
\end{equation}

\noindent
On the other hand, there are canonical homomorphisms (rigid-analytic
open immersions, but not in general isomorphisms)

\begin{equation}\tag{\SGARefTarget{sga7.sga7i.eix.n14-1-10}{14.1.10}}
 \hat G^{\mathrm{an}}\lhook\joinrel\longrightarrow G_K^{\mathrm{an}},
 \qquad
 \widehat{A^{\mathrm o}}\,{}^{\mathrm{an}}
 \lhook\joinrel\longrightarrow A_K^{\mathrm{an}}
 \qquad .
\end{equation}

\noindent
The homomorphism $e$ in (\SGARefLink{sga7.sga7i.eix.n14-1-7}{14.1.7}) is uniquely characterized by the
commutativity of the diagram

\begin{equation}\tag{\SGARefTarget{sga7.sga7i.eix.n14-1-11}{14.1.11}}
\begin{tikzcd}
 \hat G^{\mathrm{an}} \arrow[r] \arrow[d]
   & \widehat{A^{\mathrm o}}\,{}^{\mathrm{an}} \arrow[d] \\
 G_K^{\mathrm{an}} \arrow[r] & A_K^{\mathrm{an}}
\end{tikzcd}
\end{equation}

\noindent
where the unnamed arrows are those of (\SGARefLink{sga7.sga7i.eix.n14-1-9}{14.1.9}) and (\SGARefLink{sga7.sga7i.eix.n14-1-10}{14.1.10}).

\medskip
\noindent
\underline{Remarks \SGARefTarget{sga7.sga7i.eix.n14-1-12}{14.1.12}.}

\smallskip
\noindent
a) Notice that the uniqueness of $e$ making the preceding diagram
commute is immediate, since the rigid-analytic space
$G_K^{\mathrm{an}}$ is connected.  The existence of $e$, on the other
hand, is by no means trivial, and on this point we merely refer to the
cited note of M.~Raynaud.  Moreover, the definition of
$i^{\mathrm{an}}$ given in loc. cit. is likewise transcendental in
nature, although it may be interpreted as an object of an essentially
algebraic nature (\SGARefLink{sga7.sga7i.eix.n14-1-5}{14.1.5}), once the construction of $G$ itself, and
hence of $G_K$, has been obtained through the constructions of
\underline{formal geometry} in paragraph 7.  I know of no purely
algebro-geometric description of it, but below (\SGARefLink{sga7.sga7i.eix.n14-4}{14.4}) I shall propose
an ``arithmetic'' characterization, obtained by reducing to a situation
of finite type over the ring $\mathbb Z$.  Thus, from the standpoint of
algebraic geometry, the homomorphism (\SGARefLink{sga7.sga7i.eix.n14-1-5}{14.1.5}) seems to be a more hidden
piece of structure than those encountered in the preceding paragraphs.

\smallskip
\noindent
b) As explained above, the homomorphism $i$ is canonically associated
with a well-determined trivialization of the biextension $E_K$ of
$(\underline M_K,\underline M'_K)$, or equivalently with a
well-determined homomorphism (\SGARefLink{sga7.sga7i.eix.n14-1-6}{14.1.6}) lifting $p'_K$.  One then proves
that the latter is precisely the homomorphism occurring in the exact
sequence

\begin{equation}\tag{\SGARefTarget{sga7.sga7i.eix.n14-1-13}{14.1.13}}
 0\longrightarrow\underline M'{}^{\mathrm{an}}_K
 \xrightarrow{\ i'{}^{\mathrm{an}}\ }G'{}^{\mathrm{an}}_K
 \longrightarrow A'{}^{\mathrm{an}}_K\longrightarrow0
\end{equation}

\noindent
analogous to (\SGARefLink{sga7.sga7i.eix.n14-1-7}{14.1.7}), associated with the abelian scheme $A'_K$ dual
to $A_K$ (unless it is the opposite???).

% END INCLUDED SGA7I_English_current_buildable_source_20260823/components/182_expose_IX_section_14_1.tex

% BEGIN INCLUDED SGA7I_English_current_buildable_source_20260823/components/183_expose_IX_section_14_2.tex
\par\medskip
\noindent
\SGARefTarget{sga7.sga7i.eix.n14-2}{14.2}. \underline{Relations of the homomorphism}
$i:\underline M_K\longrightarrow G_K$ \underline{with monodromy.}

Let $n$ be an integer $>0$.  From the exact sequence (\SGARefLink{sga7.sga7i.eix.n14-1-7}{14.1.7}), by
means of the exact sequence of the
$\underline{\operatorname{Ext}}^i(\mathbb Z/n\mathbb Z,-)$ on the
site of rigid-analytic spaces over $K$, one then obtains a canonical
exact sequence (taking account of the fact that
$\underline M_K^{\mathrm{an}}$ is torsion-free and
$G_K^{\mathrm{an}}$ is $n$-divisible):

\begin{equation*}
 0\longrightarrow{}_nG_K^{\mathrm{an}}
 \longrightarrow{}_nA_K^{\mathrm{an}}
 \longrightarrow\underline M_{K,n}^{\mathrm{an}}
 \longrightarrow0\qquad ,
\end{equation*}

\noindent
where
$\underline M_{K,n}\overset{\mathrm{dfn}}{=}
\underline M_K/n\underline M_K$.  Since all the terms occurring in
this exact sequence come from finite group schemes over $K$, it may be
rewritten as arising from an exact sequence

\begin{equation}\tag{\SGARefTarget{sga7.sga7i.eix.n14-2-1}{14.2.1}}
 0\longrightarrow{}_nG_K\longrightarrow{}_nA_K
 \longrightarrow\underline M_{K,n}\longrightarrow0\qquad .
\end{equation}

\noindent
As $n$ varies, these exact sequences behave in the expected way, and,
on passing to the inverse limit over the powers of a prime number
$\ell$, one obtains the analogous exact sequence

\begin{equation}\tag{\SGARefTarget{sga7.sga7i.eix.n14-2-2}{14.2.2}}
 0\longrightarrow T_\ell(G_K)\longrightarrow T_\ell(A_K)
 \longrightarrow\underline M_K\otimes_{\mathbb Z}\mathbb Z_\ell
 \longrightarrow0\qquad .
\end{equation}

\noindent
But one already has such an exact sequence by purely algebraic means,
thanks to the first formula (\SGARefLink{sga7.sga7i.eix.n8-1-2}{8.1.2}) and to (\SGARefLink{sga7.sga7i.eix.n8-1-7}{8.1.7}), which give

\begin{equation}\tag{\SGARefTarget{sga7.sga7i.eix.n14-2-3}{14.2.3}}
 T_\ell(G_K)\xrightarrow{\ \sim\ }T_\ell(A_K)^{\mathrm{ef}}
 \subset T_\ell(A_K),\qquad
 T_\ell(A_K)/T_\ell(A_K)^{\mathrm{ef}}
 \simeq\underline M_K\otimes_{\mathbb Z}\mathbb Z_\ell\qquad .
\end{equation}

\noindent
With this understood, I claim that \underline{the homomorphisms}
(\SGARefLink{sga7.sga7i.eix.n14-2-3}{14.2.3}) \underline{arising from the algebraic theory are precisely
those deduced from the exact sequence} (\SGARefLink{sga7.sga7i.eix.n14-2-2}{14.2.2}) \underline{arising
from the transcendental theory}.  For the first homomorphism, which
corresponds to the first homomorphism of (\SGARefLink{sga7.sga7i.eix.n14-2-2}{14.2.2}), our assertion follows
at once, after reduction to the semistable situation, from the
commutativity of (\SGARefLink{sga7.sga7i.eix.n14-1-11}{14.1.11}), since the vertical arrows there induce
isomorphisms on the $T_\ell$, while applying $T_\ell$ to the horizontal
arrows gives, respectively, the arrow under consideration in (\SGARefLink{sga7.sga7i.eix.n14-2-3}{14.2.3})
and that in (\SGARefLink{sga7.sga7i.eix.n14-2-2}{14.2.2}).  We shall admit the asserted compatibility for
the two isomorphisms obtained between
$T_\ell(A_K)/T_\ell(A_K)^{\mathrm{ef}}$ and
$\underline M_K\otimes_{\mathbb Z}\mathbb Z_\ell$.  Notice, moreover,
that \underline{this compatibility} (for a given $\ell$)
\underline{already suffices to characterize uniquely, by
transcendental means, the homomorphism} (\SGARefLink{sga7.sga7i.eix.n14-1-5}{14.1.5}) \underline{occurring
in} (\SGARefLink{sga7.sga7i.eix.n14-1-7}{14.1.7}), interpreted as an isomorphism from
$\underline M_K^{\mathrm{an}}$ onto $\operatorname{Ker}e$.

One may also say that the homomorphism $e$ of (\SGARefLink{sga7.sga7i.eix.n14-1-7}{14.1.7}) identifies the
$T_\ell(A_K)/T_\ell(A_K)^{\mathrm{ef}}$ with the $\ell$-adic
completions $\underline R_K\otimes\mathbb Z_\ell$ of a single lattice
$\underline R_K=\operatorname{Ker}e$ over $K$, and that this
``integral structure'' on the preceding system of pro-$\ell$ groups is
compatible, in the evident sense, with the integral structure arising
from formulas (\SGARefLink{sga7.sga7i.eix.n8-1-7}{8.1.7}) (orthogonality theorem \SGARefLink{sga7.sga7i.eix.n5-2}{5.2}).  This then defines
canonically an isomorphism

\begin{equation*}
 i:\underline M_K\xrightarrow{\ \sim\ }\operatorname{Ker}e
 \qquad .
\end{equation*}

Since the homomorphism $i$ (\SGARefLink{sga7.sga7i.eix.n14-1-5}{14.1.5}) enables one to reconstruct the
usual extension structure on $T_\ell(A_K)$, namely (\SGARefLink{sga7.sga7i.eix.n14-2-2}{14.2.2}), and hence
also its structure as a \underline{blended extension}, described in
\SGARefLink{sga7.sga7i.eix.n9-6}{9.6}, by means of the exact sequence deduced from (\SGARefLink{sga7.sga7i.eix.n14-1-1}{14.1.1}),

\begin{equation*}
\begin{array}{ccccccccc}
0&\longrightarrow&T_\ell(T_K)&\longrightarrow&T_\ell(G_K)
 &\longrightarrow&T_\ell(B_K)&\longrightarrow&0\\[5ex]
 &&\underline M'_K\otimes\mathbb Z_\ell(1)&&&&&&
\end{array}
\end{equation*}

\noindent
it also enables one, in the semistable case, to reconstruct the
\underline{monodromy pairings} (\SGARefLink{sga7.sga7i.eix.n9-1-2}{9.1.2})

\begin{equation}\tag{\SGARefTarget{sga7.sga7i.eix.n14-2-4}{14.2.4}}
 u_\ell:\underline M\otimes\underline M'
 \longrightarrow\mathbb Z_\ell\qquad ,
\end{equation}

\noindent
these having been defined precisely in terms of blended extensions in
\SGARefLink{sga7.sga7i.eix.n9-6}{9.6}.  Let us merely indicate the result of translating the construction
of \SGARefLink{sga7.sga7i.eix.n9-6}{9.6}.  For this purpose, observe that in the semistable case $p_K$
and $p'_K$ arise from homomorphisms

\begin{equation*}
 p:\underline M\longrightarrow B,\qquad
 p':\underline M'\longrightarrow B'\qquad ,
\end{equation*}

\noindent
where $B,B'$ are abelian schemes over $S$, dual to one another.  Pulling
back by $(p,p')$ therefore gives a biextension $E$ of
$(\underline M,\underline M')$ by $\mathbb G_{mS}$ whose generic fibre
is the biextension $E_K$ already considered in \SGARefLink{sga7.sga7i.eix.n4-1}{4.1}.

With this said, suppose for simplicity that $\underline M$ and
$\underline M'$ are constant, with values $M$ and $M'$.  For every
pair $(m,m')\in M\times M'$, $E_{m,m'}$ is a torsor under
$\mathbb G_{mS}$, and hence corresponds to an invertible module
$L_{m,m'}$ over the ring $V$ of $S$.  It therefore admits a basis
$u_0$, determined up to multiplication by a unit of $V$.  Consequently
one obtains a canonical ``valuation'' map

\begin{equation*}
 \underline v:E_{m,m'}(K)\longrightarrow\mathbb Z
\end{equation*}

\noindent
which associates with every basis $u$ of
$L_{m,m'}\otimes_V K$, written $u=\lambda u_0$ ($\lambda\in K^*$),
the element

\begin{equation*}
 \underline v(u)=\underline v(\lambda)\in\mathbb Z\qquad ,
\end{equation*}

\noindent
where $\underline v$ on the right denotes the valuation associated
with $V$ (normalized in the usual way, so that its value group is
$\mathbb Z$); the resulting integer is plainly independent of the
choice of $u_0$.  If $\varphi$ is now a trivialization of $E$, one may
consider the corresponding map

\begin{equation}\tag{\SGARefTarget{sga7.sga7i.eix.n14-2-5}{14.2.5}}
 (m,m')\longmapsto\underline v\bigl(\varphi(m,m')\bigr):
 M\times M'\longrightarrow\mathbb Z\qquad .
\end{equation}

\noindent
With this understood, one finds that \underline{the monodromy
homomorphism} (\SGARefLink{sga7.sga7i.eix.n14-2-4}{14.2.4}) \underline{takes its values in} $\mathbb Z$
\underline{and is precisely the homomorphism} (\SGARefLink{sga7.sga7i.eix.n14-2-5}{14.2.5}), \underline{in
which} $\varphi$ \underline{is the splitting of} $E_K$
\underline{associated with the homomorphism} $i$ (\SGARefLink{sga7.sga7i.eix.n14-1-5}{14.1.5}).

\medskip
\noindent
\underline{Remarks \SGARefTarget{sga7.sga7i.eix.n14-2-6}{14.2.6}.}  In particular, one recovers (by means of
the transcendental rigid-analytic theory) the fact that the monodromy
pairing is integral and independent of $\ell$ (\SGARefLink{sga7.sga7i.eix.n10-4}{10.4} a)), proved by a
different (and likewise transcendental) method in \SGARefLink{sga7.sga7i.eix.n12-5}{12.5}.  According to
M.~Raynaud, rigid-analytic theory also recovers the positivity assertion
\SGARefLink{sga7.sga7i.eix.n10-4}{10.4} b) (which is likewise contained in \SGARefLink{sga7.sga7i.eix.n12-5}{12.5}).

\medskip
\noindent
\underline{Remarks \SGARefTarget{sga7.sga7i.eix.n14-2-7}{14.2.7}.}  Regard
$\underline M_K\otimes_{\mathbb Z}\mathbb Q$ as an extension of
$\underline M_K\otimes_{\mathbb Z}\mathbb Q/\mathbb Z$ by
$\underline M_K$.  Using the homomorphism (\SGARefLink{sga7.sga7i.eix.n14-1-5}{14.1.5})
$i:\underline M_K\longrightarrow G_K$, one deduces an extension
$A_K^b$

\begin{equation}\tag{\SGARefTarget{sga7.sga7i.eix.n14-2-8}{14.2.8}}
 0\longrightarrow G_K\longrightarrow A_K^b
 \longrightarrow\underline M_K\otimes\mathbb Q/\mathbb Z
 \longrightarrow0\qquad .
\end{equation}

\noindent
Observe moreover that, since $G_K$ is a divisible group, the datum of
such an extension is equivalent to that of a family of extensions of
the form (\SGARefLink{sga7.sga7i.eix.n14-2-2}{14.2.2}), namely those deduced from (\SGARefLink{sga7.sga7i.eix.n14-2-8}{14.2.8}) by applying the
$T_\ell$.  Since $i$ is a monomorphism, the extensions in question
which define (\SGARefLink{sga7.sga7i.eix.n14-2-8}{14.2.8}) are precisely the extensions (\SGARefLink{sga7.sga7i.eix.n14-2-2}{14.2.2}) arising
from the exact sequence (\SGARefLink{sga7.sga7i.eix.n14-2-1}{14.2.1}).  As these latter extensions are
defined by purely algebraic means, the same is true of the extension
(\SGARefLink{sga7.sga7i.eix.n14-2-8}{14.2.8}).  In fact, the characterization just given shows at once that
this extension is canonically isomorphic to the group, likewise denoted
$A_K^b$, introduced in (8.2.3), in view of the isomorphism \SGARefLink{sga7.sga7i.eix.n11-11}{11.11} from
the group of connected components of the latter to
$\underline M_K\otimes\mathbb Q/\mathbb Z$.

% END INCLUDED SGA7I_English_current_buildable_source_20260823/components/183_expose_IX_section_14_2.tex

% BEGIN INCLUDED SGA7I_English_current_buildable_source_20260823/components/184_expose_IX_section_14_3.tex
\par\medskip
\noindent
\SGARefTarget{sga7.sga7i.eix.n14-3}{14.3}. \underline{Application to an arithmetic characterization of}
$i:\underline M_K\longrightarrow G_K$.

The characterization we have in mind rests on the observation that,
for every prime number $\ell$, the extension (\SGARefLink{sga7.sga7i.eix.n14-2-2}{14.2.2}) of
$\underline M_K\otimes\mathbb Z_\ell$ by $T_\ell(G_K)$ deduced from
$i$ is known by purely algebraic means.  In particular, the image of
$i$ under the canonical homomorphism

\begin{equation}\tag{\SGARefTarget{sga7.sga7i.eix.n14-3-1}{14.3.1}}
 \operatorname{Hom}(\underline M_K,G_K)\longrightarrow
 \prod_\ell\operatorname{Ext}^1
 \bigl(\underline M_K\otimes\mathbb Z_\ell,T_\ell(G_K)\bigr)
 \qquad .
\end{equation}

\noindent
is known.  In the case (to which one may always reduce by descent) in
which $\underline M_K$ is constant, with value $M$, the preceding
homomorphism may be interpreted as

\begin{equation}\tag{\SGARefTarget{sga7.sga7i.eix.n14-3-2}{14.3.2}}
 \operatorname{Hom}\bigl(M,G_K(K)\bigr)\longrightarrow
 \prod_\ell\operatorname{Hom}
 \bigl(M,H^1(K,T_\ell(G_K))\bigr),
\end{equation}

\noindent
deduced from the canonical homomorphism

\begin{equation}\tag{\SGARefTarget{sga7.sga7i.eix.n14-3-3}{14.3.3}}
 G_K(K)\longrightarrow\prod_\ell H^1(K,T_\ell(G_K))
 \qquad ,
\end{equation}

\noindent
which comes from the connecting homomorphisms in the Kummer exact
sequences

\begin{equation*}
 0\longrightarrow{}_nG_K\longrightarrow G_K
 \xrightarrow{\ n\ }G_K\longrightarrow0\qquad .
\end{equation*}

(The homomorphism (\SGARefLink{sga7.sga7i.eix.n14-3-1}{14.3.1}) could in fact always be interpreted as the
connecting homomorphism analogous to (\SGARefLink{sga7.sga7i.eix.n14-3-3}{14.3.3}) for the twisted group
$\check{\underline M}_K\otimes G_K=
\underline{\operatorname{Hom}}(\underline M_K,G_K)$.)  The kernel of
(\SGARefLink{sga7.sga7i.eix.n14-3-3}{14.3.3}) plainly consists of the elements of $G_K(K)$ which are
\underline{divisible by} $n$ \underline{for every integer}
$n\geq1$, or, as we shall say, which are \underline{infinitely
divisible}.  Thus this already determines $i$ modulo an element of
$\operatorname{Hom}(M,g)$, where $g$ is the subgroup of $G_K(K)$
formed by its infinitely divisible elements:

\begin{equation*}
 g=\bigcap_{n\geq1}nG_K(K)\qquad .
\end{equation*}

We now have the following result.

\medskip
\noindent
\underline{Lemma \SGARefTarget{sga7.sga7i.eix.n14-3-4}{14.3.4}.}  \underline{Let} $V$ \underline{be a
discrete valuation ring whose residue field} $k$ \underline{has
characteristic} $p>0$ \underline{and is of finite type over its prime
field, and let} $G_K$ \underline{be a group scheme of finite type over}
$K$, \underline{of unipotent rank zero if} $\operatorname{char}K=0$.
\underline{Then the group} $g$ \underline{of infinitely divisible
elements of} $G_K(K)$ \underline{is zero; that is,} (\underline{if}
$G_K$ \underline{is divisible}) \underline{the homomorphism} (\SGARefLink{sga7.sga7i.eix.n14-3-3}{14.3.3})
\underline{is injective}.

After making a finite extension $K'/K$ and replacing $V$ by its
normalization in $K'$, we may suppose that $G_K$ admits a composition
series whose successive quotients are of one of the following types:
a constant finite group, an abelian variety, $\mathbb G_{mK}$, or
$\mathbb G_{aK}$.  We are therefore immediately reduced to the case
where $G_K$ itself is of one of these types.  The case of a constant
finite group, or of $\mathbb G_{aK}$, is trivial.  In the abelian case,
$G_K$ may be regarded as the generic fibre of its N\'{e}ron model $G$,
so we are reduced to showing that, if $G$ is a scheme of finite type
over $S$, the subgroup of $G(S)$ formed by the infinitely divisible
points is zero.

This is true of the group $G(k)$, as follows at once from the usual
d\'{e}vissage of $G_o$ and from the \underline{Mordell--Weil theorem}
[16] (which says that $G_o(k)$ is a finitely generated $\mathbb Z$-module
if $G_o$ is an abelian scheme over the field of finite type $k$).
Consequently $g$ is contained in the kernel $G(S)^{(1)}$ of
$G(S)\longrightarrow G(k)$.  The kernel of
$G(S)^{(1)}\longrightarrow G(V/\underline m^{,n+1})$ has a filtration
whose quotients are isomorphic to vector groups over $k$ (after
supposing, as we may, that $S$ is complete), and these are killed by
$p$.  It follows that $g$ is contained in the kernel of the preceding
homomorphism for every $n$, and hence that it is zero.  The case
$G_K=\mathbb G_{mK}$ reduces to the preceding one, since
$G_K(K)=K^*$ is an extension of $\mathbb Z$ (which has no nonzero
infinitely divisible elements) by
$V^*=\mathbb G_m(S)$, and the latter satisfies the same condition by
what we have just proved.

\medskip
\noindent
\underline{Remarks \SGARefTarget{sga7.sga7i.eix.n14-3-4-1}{14.3.4.1}.}  Notice that in \SGARefLink{sga7.sga7i.eix.n14-3-4}{14.3.4} it is essential
to require divisibility by \underline{all} integers $\geq1$, and not
only by the powers of a fixed prime $\ell$; equivalently, one must
consider the kernel of (\SGARefLink{sga7.sga7i.eix.n14-3-3}{14.3.3}), rather than restricting the right-hand
side to a single factor $H^1(K,T_\ell(G_K))$.  Indeed, if $G_K$ comes
from a smooth group scheme $G$ over $S$ and $S$ is henselian, every
element of the kernel of $G(S)\longrightarrow G(k)$ is infinitely
$\ell$-divisible for every $\ell\neq p$!  The argument proving \SGARefLink{sga7.sga7i.eix.n14-3-4}{14.3.4}
shows that the subgroup of $G_K(K)$ formed by the infinitely
$p$-divisible elements (that is, the kernel of
$G_K(K)\longrightarrow H^1(K,T_p(G_K))$ when $G_K$ is $p$-divisible)
is a \underline{finite} group.

\medskip
\noindent
\SGARefTarget{sga7.sga7i.eix.n14-3-5}{14.3.5}. The preceding considerations show that, when the residue
field $k$ has characteristic $p>0$ and is of finite type in the
absolute sense, the homomorphism $i$ (\SGARefLink{sga7.sga7i.eix.n14-1-5}{14.1.5}) is completely determined
by its image under (\SGARefLink{sga7.sga7i.eix.n14-3-1}{14.3.1}), that is, by the extensions (\SGARefLink{sga7.sga7i.eix.n14-2-2}{14.2.2}).  To
deduce from this a principle characterizing $i$ in the general case,
with no restriction on $k$, observe that in paragraph 7 we were able to
construct the scheme $G$ over $S$ under conditions considerably more
general than those in which $S$ is a complete trait and one has an
abelian variety with semistable reduction over its field of fractions.

More generally, let $S$ denote the spectrum of a normal noetherian ring
$V$, separated and complete for an $\underline m$-adic topology, and
consider a smooth quasi-projective group scheme $A$ over $S$, with
connected fibres.  Suppose that
$A_o=A\times_S S_o$, where
$S_o=\operatorname{Spec}(V/\underline m)$, is an extension of an
abelian scheme $B_o$ by an isotrivial torus $T_o$, and that $A_U$ is an
abelian scheme, where $U\subset S-S_o$ is a nonempty open subset.  In
7.22 we then constructed a scheme $G$ over $S$, an extension of an
abelian scheme $B$ by a torus $T$, characterized up to unique
isomorphism by the isomorphism (\SGARefLink{sga7.sga7i.eix.n14-1-8}{14.1.8})

\begin{equation}\tag{\SGARefTarget{sga7.sga7i.eix.n14-3-5-1}{14.3.5.1}}
 \hat G\xrightarrow{\ \sim\ }\hat A
\end{equation}

\noindent
on the formal completions along $G_o$ and $A_o$.  For every prime
number $\ell$ one deduces an isomorphism (\SGARefLink{sga7.sga7i.eix.n7-3-5}{7.3.5})

\begin{equation}\tag{\SGARefTarget{sga7.sga7i.eix.n14-3-5-2}{14.3.5.2}}
 T_\ell(G)\xrightarrow{\ \sim\ }T_\ell(A)^f\qquad ,
\end{equation}

\noindent
and hence, on restricting to $U$, the structure of an extension of
pro-$\ell$ Barsotti--Tate groups

\begin{equation}\tag{\SGARefTarget{sga7.sga7i.eix.n14-3-5-3}{14.3.5.3}}
\begin{tikzcd}[row sep=large,column sep=large]
0 \arrow[r] & T_\ell(G_U) \arrow[r] \arrow[d,equal]
 & T_\ell(A_U) \arrow[r]
 & T_\ell(A_U)/T_\ell(A_U)^f \arrow[r] & 0\quad .\\
& T_\ell(A_U)^f
\end{tikzcd}
\end{equation}

Suppose now that a second group scheme $A'$ over $S$ is given,
satisfying the same hypotheses as $A$, together with a biextension $W$
of $(A,A')$ by $\mathbb G_m$ which over $U$ induces a duality between
the abelian schemes $A_U$ and $A'_U$.  (In fact, by VIII \SGARefLink{sga7.sga7i.eviii.n7-1}{7.1}, giving
$W$ is equivalent to giving such a duality between $A_U$ and $A'_U$.)
Then $W$ defines pairings

\begin{equation}\tag{\SGARefTarget{sga7.sga7i.eix.n14-3-5-4}{14.3.5.4}}
 T_\ell(A_U)\times T_\ell(A'_U)\longrightarrow\mathbb Z_\ell(1)
 \qquad ,
\end{equation}

\noindent
which give rise to the orthogonality theorem (\SGARefLink{sga7.sga7i.eix.n7-4-8}{7.4.8}):
$T_\ell(A_U)^f$ is the orthogonal of
$T_\ell(A'_U)^t=T_\ell(T'_U)$, where $T'$ is the toric part of the
Raynaud extension $G'={A'}^{\natural}$ associated with $A'$.  Thus one
again obtains an isomorphism

\begin{equation}\tag{\SGARefTarget{sga7.sga7i.eix.n14-3-5-5}{14.3.5.5}}
 T_\ell(A_U)/T_\ell(A_U)^f
 \simeq\underline M_U\otimes\mathbb Z_\ell\qquad ,
\end{equation}

\noindent
where $\underline M=D(T')$ is the lattice over $S$ dual to the torus
$T'$.  Consequently (\SGARefLink{sga7.sga7i.eix.n14-3-5-3}{14.3.5.3}) and (\SGARefLink{sga7.sga7i.eix.n14-3-5-5}{14.3.5.5}) give an extension

\begin{equation}\tag{\SGARefTarget{sga7.sga7i.eix.n14-3-5-6}{14.3.5.6}}
 0\longrightarrow T_\ell(G_U)\longrightarrow T_\ell(A_U)
 \longrightarrow\underline M_U\otimes\mathbb Z_\ell
 \longrightarrow0\qquad .
\end{equation}

With this understood, we may state the following conjecture.

% END INCLUDED SGA7I_English_current_buildable_source_20260823/components/184_expose_IX_section_14_3.tex

% BEGIN INCLUDED SGA7I_English_current_buildable_source_20260823/components/185_expose_IX_sections_14_4_and_14_5.tex
\par\medskip
\noindent
\underline{Conjecture \SGARefTarget{sga7.sga7i.eix.n14-4}{14.4}.}  The hypotheses being those just set out in
\SGARefLink{sga7.sga7i.eix.n14-3-5}{14.3.5}, there exists a homomorphism

\begin{equation}\tag{\SGARefTarget{sga7.sga7i.eix.n14-4-1}{14.4.1}}
 i:\underline M_U\longrightarrow G_U
\end{equation}

\noindent
such that, for every prime number $\ell$, the extension (\SGARefLink{sga7.sga7i.eix.n14-3-5-6}{14.3.5.6}) is
isomorphic to the extension deduced from $i$ as explained for (\SGARefLink{sga7.sga7i.eix.n14-2-2}{14.2.2}).
More precisely, there is one and only one way of associating with every
situation of the type considered in \SGARefLink{sga7.sga7i.eix.n14-3-5}{14.3.5} a homomorphism (\SGARefLink{sga7.sga7i.eix.n14-4-1}{14.4.1})
satisfying the preceding condition, in such a way that the formation of
$i$ is compatible with every base change $S'\longrightarrow S$
associated with a continuous homomorphism $V\longrightarrow V'$ of
noetherian adic rings.

\medskip
\noindent
\SGARefTarget{sga7.sga7i.eix.n14-4-2}{14.4.2}. It seems that the existence of $i$, compatibly with base change,
follows from M. Raynaud's rigid-analytic theory, in view of the fact that
he is developing it in a framework more general than that stated in
\SGARefLink{sga7.sga7i.eix.n14-1}{14.1}, working, as here, over general noetherian adic rings.  Moreover,
to prove the existence and uniqueness asserted in \SGARefLink{sga7.sga7i.eix.n14-4}{14.4}, one is
immediately reduced to considering only those $S$ for which $S_o$ is a
scheme of finite type over $\mathbb Z$.  Indeed, $V$ is the filtered
inductive limit of subrings $V_i$ of finite type over $\mathbb Z$; these
may be supposed normal by Nagata's theorem (EGA IV 7.8), and, for $i$
sufficiently large, the situation $(A,A',W)$ comes from an analogous
situation over $V_i$, or alternatively over the separated completion of
$V_i$ for the topology defined by $m_i=\underline m V_i$.  This gives
the asserted reduction.

\medskip
\noindent
\SGARefTarget{sga7.sga7i.eix.n14-4-3}{14.4.3}. Observe now that $S-S_o$ is a Jacobson scheme (EGA IV 10.5.8).
Thus, to see that two solutions $i,i'$ of the problem, for a given
situation $(A,A',W)$, are equal, it is enough to prove that these
homomorphisms are equal at every \underline{closed} point $\eta$ of
$U$.  Let $S'$ be the normalization of the closure of $\eta$ in $S$.
Then $S'$ is a complete trait, and the uniqueness result of \SGARefLink{sga7.sga7i.eix.n4-3}{4.3} may be
applied to the pullback to $S'$ of the situation over $S$, since the
residue field of $S'$ is \underline{finite}.  Once uniqueness is thus
established, it implies compatibility of the formation of $i$ with base
change, subject to existence.

\medskip
\noindent
\underline{Remarks \SGARefTarget{sga7.sga7i.eix.n14-5}{14.5}.}  It would remain to formulate with the desired
generality, and to establish, a compatibility between the homomorphism
$i$ introduced in \SGARefLink{sga7.sga7i.eix.n14-4}{14.4} by a rigid-analytic method and the analogous
homomorphism (\SGARefLink{sga7.sga7i.eix.n13-5-13}{13.5.13}) of complex-analytic theory---or, better still, to
find a common generalization of them!


% END INCLUDED SGA7I_English_current_buildable_source_20260823/components/185_expose_IX_sections_14_4_and_14_5.tex

% BEGIN INCLUDED SGA7I_English_current_buildable_source_20260823/components/186_expose_IX_bibliography.tex
\bigskip
\noindent\underline{BIBLIOGRAPHY}

\begin{enumerate}
\item[{[1]}] M. Artin, \emph{Algebraization of formal moduli I}
  (forthcoming).
\item[{[2]}] H. Cartan, \emph{Familles d'espaces complexes et fondements
  de la g\'{e}om\'{e}trie analytique}, S\'{e}minaire ENS 1960/61.
\item[{[3]}] P. Deligne, D. Mumford, \emph{The irreducibility of the space
  of curves of given genus}, Publ. Math. no.~36 (1969).
\item[{[4]}] P. Deligne, work in preparation on Hodge theory (forthcoming
  in Publ. Math., 1970).
\item[{[5]}] M. Demazure, J. Giraud, M. Raynaud, \emph{Sch\'{e}mas
  ab\'{e}liens}, S\'{e}minaire Orsay 1967/68, forthcoming.
\item[{[6]}] A. Grothendieck, \emph{Le Groupe de Brauer III}, in
  \emph{Dix Expos\'{e}s sur la cohomologie des sch\'{e}mas}, North-Holland
  Publishing Co.
\item[{[7]}] A. Grothendieck, \emph{Crystals and the De Rham cohomology of
  schemes} (notes by I. Coates and O. Jussila), in \emph{Dix expos\'{e}s
  sur la cohomologie des sch\'{e}mas}, North-Holland Publishing Co.
\item[{[8]}] A. Grothendieck, \emph{Un th\'{e}or\`{e}me sur les
  homomorphismes de sch\'{e}mas ab\'{e}liens}, Invent. Math. \textbf{2}
  (1966), pp.~59--78.
\item[{[9]}] P. A. Griffiths, \emph{Report on variation of Hodge
  structures} (forthcoming).
\item[{[10]}] M. Hakim, \emph{Topos annel\'{e}s et sch\'{e}mas relatifs},
  forthcoming in Grundlehren der Mathematik, Springer (1970).
\item[{[11]}] J. I. Igusa, \emph{Abstract vanishing cycle theory},
  Proceedings of the Japan Academy \textbf{34}, no.~9 (1958),
  pp.~589--593.
\item[{[12]}] R. Kiehl, \emph{Der Endlichkeitssatz f\"{u}r eigentliche
  Abbildungen}, Invent. Math. \textbf{2} (1967), pp.~191--214.
\item[{[13]}] R. Kiehl.
\item[{[14]}] S. L. Kleiman, \emph{Algebraic cycles and the Weil
  conjectures}, in \emph{Dix expos\'{e}s sur la cohomologie des sch\'{e}mas},
  North-Holland Publishing Co.
\item[{[15]}] S. Lang, \emph{Abelian varieties}, Interscience Publishers,
  no.~7 (1959).
\item[{[16]}] S. Lang and A. N\'{e}ron, \emph{Rational points of abelian
  varieties over function fields}, Amer. J. Math. \textbf{81}, no.~1
  (1959), pp.~95--118.
\item[{[16bis]}] U. I. Manin, \emph{Th\'{e}orie des groupes formels
  commutatifs}, Uspekhi Mat. Nauk \textbf{XVIII} (1963), pp.~4--90.
\item[{[17]}] D. Mumford, \emph{Geometric invariant theory}, Ergebnisse
  der Mathematik 34, Springer (1965).
\item[{[18]}] D. Mumford, \emph{Abstract Theta functions}, Advanced
  Science Seminar in Algebraic Geometry, Bowdoin College, 1967 (notes by
  H. Pittie).
\item[{[19]}] A. N\'{e}ron, \emph{Mod\`{e}les minimaux des vari\'{e}t\'{e}s
  ab\'{e}liennes sur les corps locaux et globaux}, Publ. Math. 21 (1964),
  pp.~5--128.
\item[{[20]}] M. Raynaud, \emph{Caract\'{e}ristique d'Euler--Poincar\'{e}
  d'un faisceau et cohomologie des vari\'{e}t\'{e}s ab\'{e}liennes}, in
  \emph{Dix expos\'{e}s sur la cohomologie des sch\'{e}mas}, North-Holland
  Publishing Co.
\item[{[21]}] M. Raynaud, \emph{Mod\`{e}les de N\'{e}ron}, C. R. Acad.
  Sci. Paris \textbf{262} (February 1966), pp.~413--416.
\item[{[22]}] M. Raynaud, \emph{Sp\'{e}cialisation du foncteur de Picard},
  C. R. Acad. Sci. \textbf{264}, pp.~941--943 and 1001--1004.
\item[{[22bis]}] M. Raynaud, \emph{Sp\'{e}cialisation du foncteur de
  Picard}, forthcoming in Publ. Math.
\item[{[23]}] M. Raynaud, \emph{Faisceaux amples sur les sch\'{e}mas en
  groupes et les espaces homog\`{e}nes}, Paris thesis, 1968 (forthcoming
  in Lecture Notes, Springer, 1970).
\item[{[24]}] M. Raynaud.
\item[{[25]}] M. Rosenlicht, \emph{Some basic theorems on algebraic
  groups}, Amer. J. Math. \textbf{78}, no.~2 (1956), pp.~401--443.
\item[{[25bis]}] J.-P. Serre, \emph{Quelques propri\'{e}t\'{e}s des
  vari\'{e}t\'{e}s ab\'{e}liennes en caract\'{e}ristique $p$}, Amer. J.
  Math. \textbf{80} (1958), pp.~715--739.
\item[{[26]}] J.-P. Serre, \emph{Corps locaux}, Act. Sci. Ind. 1296
  (1962), Hermann, Paris.
\item[{[27]}] J.-P. Serre, \emph{Sur les corps locaux \`{a} corps
  r\'{e}siduel alg\'{e}briquement clos}, Bull. Soc. Math. France
  \textbf{89} (1961), pp.~105--154.
\item[{[28]}] J.-P. Serre, \emph{Abelian $\ell$-adic representations and
  elliptic curves}, Benjamin Inc. (1968).
\item[{[29]}] J.-P. Serre and J. Tate, \emph{Good reduction of abelian
  varieties}, Ann. of Math. \textbf{88}, no.~3 (1968), pp.~492--517.
\item[{[30]}] J.-P. Serre and J. Tate, \emph{Elliptic curves and formal
  groups} (notes of the Summer Institute on Algebraic Geometry, Woods
  Hole, 1964).
\item[{[31]}] J.-P. Serre, \emph{Morphismes universels et vari\'{e}t\'{e}
  d'Albanese}, S\'{e}minaire Chevalley 1958/59, Expos\'{e} 10.
\item[{[32]}] J. Tate, \emph{WC-groups over $p$-adic fields}, S\'{e}minaire
  Bourbaki no.~156 (1957).
\item[{[33]}] J. Tate, \emph{$p$-divisible groups}, Proceedings of a
  Conference on Local Fields, Driebergen (1966), Springer.
\item[{[34]}] J. Tate, \emph{Rigid analytic spaces} (mimeographed by
  IHES).
\end{enumerate}


% END INCLUDED SGA7I_English_current_buildable_source_20260823/components/186_expose_IX_bibliography.tex

% BEGIN INCLUDED SGA7I_English_current_buildable_source_20260823/components/187_expose_IX_terminal_publisher_matter.tex
\clearpage
\begingroup
\small
\sloppy
Lecture Notes in Mathematics


Comprehensive leaflet on request

Vol. 111: K. H. Mayer, Relationen zwischen charakteristischen Zahlen. III, 99 Seiten. 1969. DM 16,-

Vol. 112: Colloquium on Methods of Optimization. Edited by N. N. Moiseev. IV, 293 pages. 1970. DM 18,-

Vol. 113: R. Wille, Kongruenzklassengeometrien. III, 99 Seiten. 1970. DM 16,-

Vol. 114: H. Jacquet and R. P. Langlands, Automorphic Forms on GL (2). VII, 548 pages. 1970. DM 24,-

Vol. 115: K. H. Roggenkamp and V. Huber-Dyson, Lattices over Orders I. XIX, 290 pages. 1970. DM 18,-

Vol. 116: Sèminaire Pierre Lelong (Analyse) Année 1969. IV, 195 pages. 1970. DM 16,-

Vol. 117 : Y. Meyer, Nombres de Pisot, Nombres de Salem et Analyse Harmonique. 63 pages. 1970. DM 16,-

Vol. 118: Proceedings of the 15th Scandinavian Congress, Oslo 1968. Edited by K. E. Aubert and W. Ljunggren. IV, 162 pages. 1970. DM 16,-

Vol. 119: M. Raynaud, Faisceaux amples sur les schémas en groupes et les espaces homogènes. III, 219 pages. 1970. DM 16,-

Vol. 120: D. Siefkes, Büchi's Monadic Second Order Successor Arithmetic. XII, 130 Seiten. 1970. DM 16,-

Vol. 121: H. S. Bear, Lectures on Gleason Parts. III, 47 pages. 1970. DM 16,-

Vol. 122: H. Zieschang, E. Vogt und H.-D. Coldewey, Flächen und ebene diskontinuierliche Gruppen. VIII, 203 Seiten. 1970. DM 16,-

Vol. 123: A. V. Jategaonkar, Left Principal Ideal Rings.VI, 145 pages. 1970. DM 16,-

Vol. 124: Séminare de Probabilités IV. Edited by P. A. Meyer. IV, 282 pages. 1970. DM 20,-

Vol. 125: Symposium on Automatic Demonstration. V, 310 pages.1970. DM 20,-

Vol. 126: P. Schapira, Théorie des Hyperfonctions. XI,157 pages.1970. DM 16,-

Vol. 127: I. Stewart, Lie Algebras. IV, 97 pages. 1970. DM 16,-

Vol. 128: M. Takesaki, Tomita's Theory of Modular Hilbert Algebras and its Applications. II, 123 pages. 1970. DM 16,-

Vol. 129: K. H. Hofmann, The Duality of Compact Semigroups and C$^*$- Bigebras. XII, 142 pages. 1970. DM 16,-

Vol. 130: F. Lorenz, Quadratische Formen über Körpern. II, 77 Seiten. 1970. DM 16,-

Vol. 131: A Borel et al., Seminar on Algebraic Groups and Related Finite Groups. VII, 321 pages. 1970. DM 22,-

Vol. 132: Symposium on Optimization. III, 348 pages. 1970. DM 22,-

Vol. 133: F. Topsøe, Topology and Measure. XIV, 79 pages. 1970. DM 16,-

Vol. 134: L. Smith, Lectures on the Eilenberg-Moore Spectral Sequence. VII, 142 pages. 1970. DM 16,-

Vol. 135: W. Stoll, Value Distribution of Holomorphic Maps into Compact Complex Manifolds. II, 267 pages. 1970. DM 18,-

Vol. 136 : M. Karoubi et al., Séminaire Heidelberg-Saarbrücken-Strasbuorg sur la K-Théorie. IV, 264 pages. 1970. DM 18,-

Vol. 137 : Reports of the Midwest Category Seminar IV. Edited by S. MacLane. III, 139 pages. 1970. DM 16,-

Vol. 138: D. Foata et M. Schützenberger, Théorie Géométrique des Polynômes Eulériens. V, 94 pages. 1970. DM 16,-

Vol. 139: A. Badrikian, Séminaire sur les Fonctions Aléatoires Linéaires et les Mesures Cylindriques. VII, 221 pages. 1970. DM 18,-

Vol. 140: Lectures in Modern Analysis and Applications II. Edited by C. T. Taam. VI, 119 pages. 1970. DM 16,-

Vol. 141: G. Jameson, Ordered Linear Spaces. XV, 194 pages. 1970. DM 16,-

Vol. 142: K. W. Roggenkamp, Lattices over Orders II. V, 388 pages. 1970. DM 22,-

Vol. 143: K. W. Gruenberg, Cohomological Topics in Group Theory. XIV, 275 pages. 1970. DM 20,-

Vol. 144: Seminar on Differential Equations and Dynamical Systems, II. Edited by J. A. Yorke. VIII, 268 pages. 1970. DM 20,-

Vol. 145: E. J. Dubuc, Kan Extensions in Enriched Category Theory. XVI, 173 pages. 1970. DM 16,-

Vol. 146: A. B. Altman and S. Kleiman, Introduction to Grothendieck Duality Theory. II, 192 pages. 1970. DM 18,-

Vol. 147: D. E. Dobbs, Cech Cohomological Dimensions for Commutative Rings. VI, 176 pages. 1970. DM 16,-

Vol. 148: R. Azencott, Espaces de Poisson des Groupes Localement Compacts. IX, 141 pages. 1970. DM 16,-

Vol. 149: R. G. Swan and E. G. Evans, K-Theory of Finite Groups and Orders. IV, 237 pages. 1970. DM 20,-

Vol. 150: Heyer, Dualität lokalkompakter Gruppen. XIII, 372 Seiten. 1970. DM 20,-

Vol. 151: M. Demazure et A. Grothendieck, Schémas en Groupes I. (SGA 3). XV, 562 pages. 1970. DM 24,-

Vol. 152: M. Demazure et A. Grothendieck, Schémas en Groupes II. (SGA 3). IX, 654 pages. 1970. DM 24,-

Vol. 153: M. Demazure et A. Grothendieck, Schémas en Groupes III. (SGA 3). VIII, 529 pages. 1970. DM 24,-

Vol. 154: A. Lascoux et M. Berger, Variétés Kähleriennes Compactes. VII, 83 pages. 1970. DM 16,-

Vol. 155: Several Complex Variables I, Maryland 1970. Edited by J. Horváth. IV, 214 pages. 1970. DM 18,-

Vol. 156: R. Hartshorne, Ample Subvarieties of Algebraic Varieties. XIV, 256 pages. 1970. DM 20,-

Vol. 157: T. tom Dieck, K. H. Kamps und D. Puppe, Homotopietheorie. VI, 265 Seiten. 1970. DM 20,-

Vol. 158: T. G. Ostrom, Finite Translation Planes. IV. 112 pages. 1970. DM 16,-

Vol. 159: R. Ansorge und R. Hass. Konvergenz von Differenzenverfahren für lineare und nichtlineare Anfangswertaufgaben. VIII, 145 Seiten. 1970. DM 16,-

Vol. 160: L. Sucheston, Constributions to Ergodic Theory and Probability. VII, 277 pages. 1970. DM 20,-

Vol. 161: J. Stasheff, H-Spaces from a Homotopy Point of View. VI, 95 pages. 1970. DM 16,-

Vol. 162: Harish-Chandra and van Dijk, Harmonic Analysis on Reductive p-adic Groups. IV, 125 pages. 1970. DM 16,-

Vol. 163: P. Deligne, Equations Différentielles à Points Singuliers Reguliers. III, 133 pages. 1970. DM 16,-

Vol. 164: J. P. Ferrier, Seminaire sur les Algebres Complètes. II, 69 pages. 1970. DM 16,-

Vol. 165: J. M. Cohen, Stable Homotopy. V, 194 pages. 1970. DM 16,-

Vol. 166: A. J. Silberger, PGL$_2$ over the p-adics: its Representations, Spherical Functions, and Fourier Analysis. VII, 202 pages. 1970. DM 18,-

Vol. 167: Lavrentiev, Romanov and Vasiliev, Multidimensional Inverse Problems for Differential Equations. V, 59 pages. 1970. DM 16,-

Vol. 168: F. P. Peterson, The Steenrod Algebra and its Applications: A conference to Celebrate N. E. Steenrod's Sixtieth Birthday. VII, 317 pages. 1970. DM 22,-

Vol. 169: M. Raynaud, Anneaux Locaux Henséliens. V, 129 pages. 1970. DM 16,-

Vol. 170: Lectures in Modern Analysis and Applications III. Edited by C. T. Taam. VI, 213 pages. 1970. DM 18,-

Vol. 171: Set-Valued Mappings, Selections and Topological Properties of $2^X$. Edited by W. M. Fleischman. X, 110 pages. 1970. DM 16,-

Vol. 172: Y.-T. Siu and G. Trautmann, Gap-Sheaves and Extension of Coherent Analytic Subsheaves. V, 172 pages. 1971. DM 16,-

Vol. 173: J. N. Mordeson and B. Vinograde, Structure of Arbitrary Purely Inseparable Extension Fields. IV, 138 pages. 1970. DM 16,-

Vol. 174: B. Iversen, Linear Determinants with Applications to the Picard Scheme of a Family of Algebraic Curves. VI, 69 pages. 1970. DM 16,-

Vol. 175: M. Brelot, On Topologies and Boundaries in Potential Theory. VI, 176 pages. 1971. DM 18,-

Vol. 176: H. Popp, Fundamentalgruppen algebraischer Mannigfaltigkeiten. IV, 154 Seiten. 1970. DM 16,-

Vol. 177: J. Lambek, Torsion Theories, Additive Semantics and Rings of Quotients. VI, 94 pages. 1971. DM 16,-

Please turn over


%%%%%%%%%%  scan idx 537  |  folio 526  %%%%%%%%%%

Vol. 178: Th. Bröcker und T. tom Dieck, Kobordismentheorie. XVI, 191 Seiten. 1970. DM 18,-

Vol. 179: Seminaire Bourbaki -- vol. 1968/69. Exposés 347-363. IV. 295 pages. 1971. DM 22,-

Vol. 180: Séminaire Bourbaki -- vol. 1969/70. Exposés 364-381. IV, 310 pages. 1971. DM 22,-

Vol. 181: F. DeMeyer and E. Ingraham, Separable Algebras over Commutative Rings. V, 157 pages. 1971. DM 16.-

Vol. 182: L. D. Baumert. Cyclic Difference Sets. VI, 166 pages. 1971. DM 16,-

Vol. 183: Analytic Theory of Differential Equations. Edited by P. F. Hsieh and A. W. J. Stoddart. VI, 225 pages. 1971. DM 20,-

Vol. 184: Symposium on Several Complex Variables, Park City, Utah, 1970. Edited by R. M. Brooks. V, 234 pages. 1971. DM 20,-

Vol. 185: Several Complex Variables II, Maryland 1970. Edited by J. Horváth. III, 287 pages. 1971. DM 24,-

Vol. 186: Recent Trends in Graph Theory. Edited by M. Capobianco/J. B. Frechen/M. Krolik. VI, 219 pages. 1971. DM 18.-

Vol. 187: H. S. Shapiro, Topics in Approximation Theory. VIII, 275 pages. 1971. DM 22,-

Vol. 188: Symposium on Semantics of Algorithmic Languages. Edited by E. Engeler. VI, 372 pages. 1971. DM 26,-

Vol. 189: A. Weil, Dirichlet Series and Automorphic Forms. V, 164 pages. 1971. DM 16,-

Vol. 190: Martingales. A Report on a Meeting at Oberwolfach, May 17-23, 1970. Edited by H. Dinges. V, 75 pages. 1971. DM 16,-

Vol. 191: Séminaire de Probabilités V. Edited by P. A. Meyer. IV, 372 pages. 1971. DM 26,-

Vol. 192: Proceedings of Liverpool Singularities -- Symposium I. Edited by C. T. C. Wall. V, 319 pages. 1971. DM 24,-

Vol. 193: Symposium on the Theory of Numerical Analysis. Edited by J. Ll. Morris. VI, 152 pages. 1971. DM 16,-

Vol. 194: M. Berger, P. Gauduchon et E. Mazet. Le Spectre d'une Variété Riemannienne. VII, 251 pages. 1971. DM 22,-

Vol. 195: Reports of the Midwest Category Seminar V. Edited by J.W. Gray and S. Mac Lane.III, 255 pages. 1971. DM 22,-

Vol. 196: H-spaces -- Neuchâtel (Suisse)- Août 1970. Edited by F. Sigrist, V, 156 pages. 1971. DM 16,-

Vol. 197: Manifolds -- Amsterdam 1970. Edited by N. H. Kuiper. V, 231 pages. 1971. DM 20,-

Vol. 198: M. Hervé, Analytic and Plurisubharmonic Functions in Finite and Infinite Dimensional Spaces. VI, 90 pages. 1971. DM 16.-

Vol. 199: Ch. J. Mozzochi, On the Pointwise Convergence of Fourier Series. VII, 87 pages. 1971. DM 16,-

Vol. 2 00: U. Neri, Singular Integrals. VII, 272 pages. 1971. DM 22,-

Vol. 201 : J. H. van Lint, Coding Theory. VII, 136 pages. 1971. DM 16,-

Vol. 202: J. Benedetto, Harmonic Analysis on Totally Disconnected Sets. VIII, 261 pages. 1971. DM 22,-

Vol. 203: D. Knutson, Algebraic Spaces. VI, 261 pages. 1971. DM 22,-

Vol. 204: A. Zygmund, Intégrales Singulières. IV, 53 pages. 1971. DM 16,-

Vol. 205: Séminaire Pierre Lelong (Analyse) Année 1970. VI, 243 pages. 1971. DM 20,-

Vol. 206: Symposium on Differential Equations and Dynamical Systems. Edited by D. Chillingworth. XI, 173 pages. 1971. DM 16,-

Vol. 207: L. Bernstein, The Jacobi-Perron Algorithm -- Its Theory and Application. IV, 161 pages. 1971. DM 16,-

Vol. 208: A. Grothendieck and J. P. Murre, The Tame Fundamental Group of a Formal Neighbourhood of a Divisor with Normal Crossings on a Scheme. VIII, 133 pages. 1971. DM 16,-

Vol. 209: Proceedings of Liverpool Singularities \quad Symposium II. Edited by C. T. C. Wall. V, 280 pages. 1971. DM 22,-

Vol. 210: M. Eichler, Projective Varieties and Modular Forms. III, 118 pages. 1971. DM 16,-

Vol. 211: Théorie des Matroïdes. Edité par C. P. Bruter. III, 108 pages. 1971. DM 16,-

Vol. 212: B. Scarpellini, Proof Theory and Intuitionistic Systems. VII, 291 pages. 1971. DM 24,-

Vol. 213: H. Hogbe-Nlend, Théorie des Bornologies et Applications. V, 168 pages. 1971. DM 18,-

Vol. 214: M. Smorodinsky, Ergodic Theory, Entropy. V, 64 pages. 1971. DM 16,-

Vol. 215: P. Antonelli, D. Burghelea and P. J. Kahn, The Concordance-Homotopy Groups of Geometric Automorphism Groups. X, 140 pages. 1971. DM 16,-

Vol. 216: H. Maaß, Siegel's Modular Forms and Dirichlet Series. VII, 328 pages. 1971. DM 20,-

Vol. 217: T. J. Jech, Lectures in Set Theory with Particular Emphasis on the Method of Forcing. V, 137 pages. 1971. DM 16,-

Vol. 218: C. P. Schnorr, Zufälligkeit und Wahrscheinlichkeit. IV, 212.Seiten 1971. DM 20,-

Vol. 219: N. L. Alling and N. Greenleaf, Foundations of the Theory of Klein Surfaces. IX, 117 pages. 1971. DM 16,-

Vol. 220: W. A. Coppel, Disconjugacy. V, 148 pages. 1971. DM 16,-

Vol. 221: P. Gabriel und F. Ulmer, Lokal präsentierbare Kategorien. V, 200 Seiten. 1971. DM 18,-

Vol. 222: C. Meghea, Compactification des Espaces Harmoniques. III, 108 pages. 1971. DM 16,-

Vol. 223: U. Felgner, Models of ZF-Set Theory. VI, 173 pages. 1971. DM 16,-

Vol. 224: Revètements Etales et Groupe Fondamental. (SGA 1). Dirigé par A. Grothendieck XXII, 447 pages. 1971. DM 30,-

Vol. 225: Théorie des Intersections et Théorème de Riemann-Roch. (SGA 6). Dirigé par P. Berthelot, A. Grothendieck et L. Illusie. XII, 700 pages. 1971. DM 40,-

Vol. 226: Seminar on Potential Theory, II. Edited by H. Bauer. IV, 170 pages. 1971. DM 18,-

Vol. 227: H. L. Montgomery, Topics in Multiplicative Number Theory. IX, 178 pages. 1971. DM 18,-

Vol. 228: Conference on Applications of Numerical Analysis. Edited by J. Ll. Morris. X, 358 pages. 1971. DM 26,-

Vol. 229: J. Väisälä, Lectures on n-Dimensional Quasiconformal Mappings. XIV, 144 pages. 1971. DM 16,-

Vol. 230: L. Waelbroeck, Topological Vector Spaces and Algebras. VII, 158 pages. 1971. DM 16,-

Vol. 231: H. Reiter, L$^1$-Algebras and Segal Algebras. XI, 113 pages. 1971. DM 16,-

Vol. 232: T. H. Ganelius, Tauberian Remainder Theorems. VI, 75 pages. 1971. DM 16,-

Vol. 233: C. P. Tsokos and W. J. Padgett. Random Integral Equations with Applications to Stochastic Systems. VII, 174 pages. 1971. DM 18,-

Vol. 234: A. Andreotti and W. Stoll. Analytic and Algebraic Dependence of Meromorphic Functions. III, 390 pages. 1971. DM 26,-

Vol. 235: Global Differentiable Dynamics. Edited by O. Hájek, A. J. Lohwater, and R. McCann. X, 140 pages. 1971. DM 16,-

Vol. 236 : M. Barr, P. A. Grillet, and D. H. van Osdol. Exact Categories and Categories of Sheaves. VII, 239 pages. 1971, DM 20,-

Vol. 237: B. Stenström. Rings and Modules of Quotients. VII, 136 pages. 1971. DM 16,-

Vol. 238: Der kanonische Modul eines Cohen-Macaulay-Rings. Herausgegeben von Jürgen Herzog und Ernst Kunz. VI, 103 Seiten. 1971. DM 16,-

Vol. 239: L. Illusie, Complexe Cotangent et Déformations I. XV, 355 pages. 1971. DM 26,-

Vol. 240: A. Kerber, Representations of Permutation Groups I. VII, 192 pages. 1971. DM 18,-

Vol. 241: S. Kaneyuki, Homogeneous Bounded Domains and Siegel Domains. V, 89 pages. 1971. DM 16,-

Vol. 242: R. R. Coifman et G. Weiss, Analyse Harmonique Non-Commutative sur Certains Espaces. V, 160 pages. 1971. DM 16,-

Vol. 243: Japan-United States Seminar on Ordinary Differential and Functional Equations. Edited by M. Urabe. VIII, 332 pages. 1971. DM 26,-

Vol. 244: Séminaire Bourbaki -- vol. 1970/71. Exposés 382-399. IV, 356 pages. 1971. DM 26,-

Vol. 245: D. E. Cohen, Groups of Cohomological Dimension One. V, 99 pages. 1972. DM 16,-


%%%%%%%%%%  scan idx 538  |  folio 527  %%%%%%%%%%

\noindent\rule{\linewidth}{0.4pt}

Vol. 246: Lectures on Rings and Modules. Tulane University Ring and Operator Theory Year, 1970--1971. Volume I. X, 661 pages. 1972. DM 40,--

Vol. 247: Lectures on Operator Algebras. Tulane University Ring and Operator Theory Year, 1970--1971. Volume II. XI, 786 pages. 1972. DM 40,--

Vol. 248: Lectures on the Applications of Sheaves to Ring Theory. Tulane University Ring and Operator Theory Year, 1970--1971. Volume III. VIII, 315 pages. 1971. DM 26,--

Vol. 249: Symposium on Algebraic Topology. Edited by P. J. Hilton. VII, 111 pages. 1971. DM 16,--

Vol. 250: B. Jónsson, Topics in Universal Algebra. VI, 220 pages. 1972. DM 20,--

Vol. 251: The Theory of Arithmetic Functions. Edited by A. A. Gioia and D. L. Goldsmith VI, 287 pages. 1972. DM 24,--

Vol. 252: D. A. Stone, Stratified Polyhedra. IX, 193 pages. 1972. DM 18,--

Vol. 253: V. Komkov, Optimal Control Theory for the Damping of Vibrations of Simple Elastic Systems. V, 240 pages. 1972. DM 20,--

Vol. 254: C. U. Jensen, Les Foncteurs Dérivés de $\varprojlim$ et leurs Applications en Théorie des Modules. V, 103 pages. 1972. DM 16,--

Vol. 255: Conference in Mathematical Logic -- London '70. Edited by W. Hodges. VIII, 351 pages. 1972. DM 26,--

Vol. 256: C. A. Berenstein and M. A. Dostal, Analytically Uniform Spaces and their Applications to Convolution Equations. VII, 130 pages. 1972. DM 16,--

Vol. 257: R. B. Holmes, A Course on Optimization and Best Approximation. VIII, 233 pages. 1972. DM 20,--

Vol. 258: Séminaire de Probabilités VI. Edited by P. A. Meyer. VI, 253 pages. 1972. DM 22,--

Vol. 259: N. Moulis, Structures de Fredholm sur les Variétés Hilbertiennes. V, 123 pages. 1972. DM 16,--

Vol. 260: R. Godement and H. Jacquet, Zeta Functions of Simple Algebras. IX, 188 pages. 1972. DM 18,--

Vol. 261: A. Guichardet, Symmetric Hilbert Spaces and Related Topics. V, 197 pages. 1972. DM 18,--

Vol. 262: H. G. Zimmer, Computational Problems, Methods, and Results in Algebraic Number Theory. V, 103 pages. 1972. DM 16,--

Vol. 263: T. Parthasarathy, Selection Theorems and their Applications. VII, 101 pages. 1972. DM 16,--

Vol. 264: W. Messing, The Crystals Associated to Barsotti-Tate Groups: with Applications to Abelian Schemes. III, 190 pages. 1972. DM 18,--

Vol. 265: N. Saavedra Rivano, Catégories Tannakiennes. II, 418 pages. 1972. DM 26,--

Vol. 266: Conference on Harmonic Analysis. Edited by D. Gulick and R. L. Lipsman. VI, 323 pages. 1972. DM 24,--

Vol. 267: Numerische Lösung nichtlinearer partieller Differential- und Integro-Differentialgleichungen. Herausgegeben von R. Ansorge und W. Törnig, VI, 339 Seiten. 1972. DM 26,--

Vol. 268: C. G. Simader, On Dirichlet's Boundary Value Problem. IV, 238 pages. 1972. DM 20,--

Vol. 269: Théorie des Topos et Cohomologie Etale des Schémas. (SGA 4). Dirigé par M. Artin, A. Grothendieck et J. L. Verdier. XIX, 525 pages. 1972. DM 50,--

Vol. 270: Théorie des Topos et Cohomologie Etale des Schémas. Tome 2. (SGA 4). Dirigé par M. Artin, A. Grothendieck et J. L. Verdier. V, 314 pages. 1972. DM 50,--

Vol. 271: J. P. May, The Geometry of Iterated Loop Spaces. IX, 175 pages. 1972. DM 18,--

Vol. 272: K. R. Parthasarathy and K. Schmidt, Positive Definite Kernels, Continuous Tensor Products, and Central Limit Theorems of Probability Theory. VI, 107 pages. 1972. DM 16,--

Vol. 273: U. Seip, Kompakt erzeugte Vektorräume und Analysis. IX, 119 Seiten. 1972. DM 16,--

Vol. 274: Toposes, Algebraic Geometry and Logic. Edited by. F. W. Lawvere. VI, 189 pages. 1972. DM 18,--

Vol. 275: Séminaire Pierre Lelong (Analyse) Année 1970--1971. VI, 181 pages. 1972. DM 18,--

Vol. 276: A. Borel, Représentations de Groupes Localement Compacts. V, 98 pages. 1972. DM 16,--

Vol. 277: Séminaire Banach. Edité par C. Houzel. VII, 229 pages. 1972. DM 20,--

Vol. 278: H. Jacquet, Automorphic Forms on GL(2). Part II. XIII, 142 pages. 1972. DM 16,--

Vol. 279: R. Bott, S. Gitler and I. M. James, Lectures on Algebraic and Differential Topology. V, 174 pages. 1972. DM 18,--

Vol. 280: Conference on the Theory of Ordinary and Partial Differential Equations. Edited by W. N. Everitt and B. D. Sleeman. XV, 367 pages. 1972. DM 26,--

Vol. 281: Coherence in Categories. Edited by S. Mac Lane. VII, 235 pages. 1972. DM 20,--

Vol. 282: W. Klingenberg und P. Flaschel, Riemannsche Hilbertmannigfaltigkeiten. Periodische Geodätische. VII, 211 Seiten. 1972. DM 20,--

Vol. 283: L Illusie, Complexe Cotangent et Déformations II. VII, 304 pages. 1972. DM 24,--

Vol. 284: P. A. Meyer, Martingales and Stochastic Integrals I. VI, 89 pages. 1972. DM 16,--

Vol. 285: P. de la Harpe, Classical Banach-Lie Algebras and Banach-Lie Groups of Operators in Hilbert Space. III, 160 pages. 1972. DM 16,--

Vol. 286: S. Murakami, On Automorphisms of Siegel Domains. V, 95 pages. 1972. DM 16,--

Vol. 288: Groupes de Monodromie en Géométrie Algébrique. (SGA 7 I). Dirigé par A. Grothendieck. IX, 523 pages. 1972. DM 50,--

\noindent\rule{\linewidth}{0.4pt}


%%%%%%%%%%  scan idx 539  |  folio 528  %%%%%%%%%%

ISBN 3-540-05987-3

ISBN 0-387-05987-3
\endgroup


% END INCLUDED SGA7I_English_current_buildable_source_20260823/components/187_expose_IX_terminal_publisher_matter.tex

\end{document}

% SGA-CUMULATIVE-END 08_SGA7I_EN_EXPOSES_I_II_VI_IX_READER.tex

% SGA-CUMULATIVE-BEGIN 09_SGA7II_EN_COMPLETE_READER.tex
% Complete editable cumulative English source. All mathematical text is embedded below.
% Compile in a fresh directory. The filecontents blocks materialize the included
% source files, preserving native TeX file-boundary semantics. No external text
% inputs are required. Standard TeX distribution packages remain dependencies.
% The companion ZIP preserves the original modular source and build instructions.
% Original input: components/001_expose_X_opening_through_1_9.tex
\begin{filecontents*}[overwrite,noheader]{sga7ii-cumulative-input-001.tex}
%%%%%%%%%% source scan idx 8 | folio 1 %%%%%%%%%%

\begin{center}
\underline{INTERSECTIONS ON REGULAR SURFACES}

\medskip
\underline{by P. Deligne}
\end{center}

In the first part of this exposé, after some preliminaries from intersection theory
(contained in $[5]\cup[7]$), we give Artin's proof of a special case of a theorem of
Raynaud [9] (1.13), and an a priori proof of a theorem which Néron [8] had verified
case by case (1.15).  The theorems of [9] used in [2], and hence in IX, follow fairly
easily from (1.13).

In the second part, continuing in the same vein, I explain how the ``Italians'' [1],
[3] calculated intersection numbers on surfaces.  The main results, 2.9 and 2.11,
may also be found in Northcott [12] and [13].  To conclude, in 2.27, 2.30, and 2.34
I give a conjectural relation between the dimension of the versal formal moduli
scheme of a singularity of a reduced curve, the invariant $\delta$ of the singularity,
and an invariant measuring its symmetry.

This exposé will not be used in the remainder of the seminar.

\medskip
1. \underline{Intersection numbers on arithmetic surfaces}

1.0. Throughout this section, let $X$ be a regular scheme of dimension 2, let $Y$ be
a reduced divisor on $X$, and let $k$ be a subfield of
$H^0(Y,\underline{O}_Y)$.  We assume that $Y$ is a curve (EGA II 7.4.1) proper over
$k$.  Here are three examples of this situation.

\underline{Example} 1.1. Let $(S,\eta,s)$ be a trait and let
$f:X\longrightarrow S$ be a regular, proper, and flat $S$-scheme whose fibres are
purely one-dimensional.  We take $Y=(f^{-1}(s))_{\mathrm{red}}$ and $k=k(s)$.  From
1.10 onward we shall restrict ourselves to this case.

%%%%%%%%%% source scan idx 9 | folio 2 %%%%%%%%%%

\underline{Example} 1.2. Let $X_o$ be the spectrum of a normal local ring of
dimension 2, let $s$ be the closed point of $X_o$, and let
$f:X\longrightarrow X_o$ be a regular connected $X_o$-scheme, proper over $X_o$,
distinct from $X_o$ but agreeing with $X_o$ over $X_o\setminus\{s\}$.  We take
$Y=(f^{-1}(s))_{\mathrm{red}}$ and $k=k(s)$.

\underline{Example} 1.3. We take $X$ to be a complete regular surface over a field
$k$, and $Y$ to be a reduced curve on $X$.

Recall the following lemma (SGA 6 IV, Ex.~1.9).

\underline{Lemma} 1.4. \underline{Let} $Z_o$ \underline{be a closed subscheme of a
Noetherian scheme} $Z$. \underline{Then the inclusion of} $Z_o$ \underline{in} $Z$
\underline{induces an isomorphism from the Grothendieck group of the category of
coherent} $\underline{O}_{Z_o}$\underline{-modules to the Grothendieck group of the
category of coherent sheaves on} $Z$ \underline{with support in} $Z_o$.

This lemma allows us to define the Euler--Poincaré characteristic $\chi(\mathcal F)$
of a coherent sheaf $\mathcal F$ on $X$ supported on $Y$ so that:

(i) If $\mathcal F$ is a coherent sheaf of $\underline{O}_Y$-modules, then
\begin{equation}\tag{1.4.1}
\chi(\mathcal F)=\dim_k H^0(\mathcal F)-\dim_k H^1(\mathcal F).
\end{equation}

(ii) $\chi$ is additive with respect to short exact sequences.

If the cohomology sheaves of a complex $K$ of $\underline{O}_X$-modules are coherent
sheaves supported on $Y$, we set
\begin{equation}\tag{1.4.2}
\chi(K)=\sum_i(-1)^i\chi(\underline{H}^i(K)).
\end{equation}
If $\mathcal F$ is supported at a point $y$ of $Y$, then
\begin{equation}\tag{1.4.3}
\chi(\mathcal F)=\operatorname{length}_{\underline{O}_y}(\mathcal F_y)[k(y):k].
\end{equation}

%%%%%%%%%% source scan idx 10 | folio 3 %%%%%%%%%%

Under the hypotheses of (1.1), (1.2), or (1.3), one has
\begin{equation}\tag{1.4.4}
\chi(\mathcal F)=\operatorname{length}f_*\mathcal F-
\operatorname{length}R^1f_*\mathcal F.
\end{equation}
Let $D$ be a divisor $\geq0$ supported on $Y$, and let $\mathcal L$ be an invertible
sheaf on $D$.  There is then a unique way to define the degree
$\deg_D(\mathcal L)$ such that
\begin{equation}\tag{1.4.5}
\left\{
\begin{array}{l}
\deg_D(\mathcal L'\otimes\mathcal L'')=
\deg_D(\mathcal L')+\deg_D(\mathcal L''),\\[3ex]
\text{for }\mathcal L\subset\underline O,
\text{ one has }\deg_D(\mathcal L)=-\chi(\underline O/\mathcal L).
\end{array}
\right.
\end{equation}
The Riemann--Roch formula on $D$ is then
\begin{equation}\tag{1.4.6}
\chi(\mathcal L)=\chi(\underline{O}_D)+\deg_D(\mathcal L).
\end{equation}
If $D=\sum n_iD_i$ is the decomposition of $D$ into reduced irreducible divisors,
then
\begin{equation}\tag{1.4.7}
\deg_D(\mathcal L)=\sum n_i\deg_{D_i}(\mathcal L).
\end{equation}

\underline{Definition} 1.5. \underline{Let} $D$ \underline{and} $E$
\underline{be two divisors} $\geq0$ \underline{on} $X$, \underline{and suppose that}
$D$ \underline{is supported on} $Y$. \underline{The intersection number of} $D$
\underline{and} $E$ \underline{(relative to the field} $k$ (1.0)) \underline{is
defined by}
\begin{equation*}
(D,E)=\chi(\underline{O}_D\overset{\mathbb L}{\otimes}\underline{O}_E)
=\chi(\underline{O}_D\otimes\underline{O}_E)
-\chi(\operatorname{Tor}^1(\underline{O}_D,\underline{O}_E)).
\end{equation*}
Indeed, under the hypotheses made here,
$\operatorname{Tor}^2(\underline{O}_D,\underline{O}_E)=0$.

\underline{Proposition} 1.6. \underline{Under the hypotheses of} 1.5,
\underline{one has}:

(i) \quad \underline{The intersection number} $(D,E)$ \underline{is additive in the
divisors} $D$ \underline{and} $E$; \underline{this allows us, by linearity, to define
the intersection number of two divisors which are not necessarily positive.}

%%%%%%%%%% source scan idx 11 | folio 4 %%%%%%%%%%

(ii) \quad $(D,E)=\deg_D(\underline O(E))$.

(iii) \quad \underline{If} $D$ \underline{and} $E$ \underline{meet at only finitely
many points, then} $(D,E)=\chi(\underline{O}_D\otimes\underline{O}_E)\geq0$,
\underline{and equality can occur only if} $D$ \underline{and} $E$
\underline{are disjoint.}

(iv) \quad \underline{If} $D$ \underline{and} $E$ \underline{are supported on} $Y$,
\underline{then} $(D,E)=(E,D)$.

(v) \quad \underline{If} $N$ \underline{is the conormal sheaf of} $D$,
\underline{then}
\begin{equation*}
D^2=(D,D)=-\deg_D(N).
\end{equation*}

(vi) \quad \underline{If} $E$, \underline{not necessarily positive, is the divisor
of a meromorphic function on} $X$, \underline{then} $(D,E)=0$.

\underline{Proof}. The exact sequence
\begin{equation*}
0\longrightarrow\underline O(-E)\longrightarrow\underline O_X
\longrightarrow\underline O_E\longrightarrow0
\end{equation*}
gives, by (1.4.6),
\begin{equation*}
\begin{array}{rl}
(D,E)\overset{\mathrm{def}}{=}&
\chi(\underline{O}_D\overset{\mathbb L}{\otimes}\underline{O}_E)\\
=&\chi(\underline{O}_D\overset{\mathbb L}{\otimes}\underline{O}_X)
-\chi(\underline{O}_D\overset{\mathbb L}{\otimes}\underline O(-E))\\[1ex]
=&\chi(\underline{O}_D)-\chi(\underline O(-E)|D)
=-\deg_D(\underline O(-E))=\deg_D(\underline O(E)).
\end{array}
\end{equation*}
This proves (ii), (vi), and biadditivity, by (1.4.5) and (1.4.7).  Under the
hypotheses of (iii), one has
$\operatorname{Tor}^1(\underline{O}_D,\underline{O}_E)=0$, which gives the
conclusion.  Under the hypotheses of (v), one has
\begin{equation*}
\underline{O}_D\otimes\underline{O}_D=\underline{O}_D
\quad\text{and}\quad
\operatorname{Tor}^1(\underline{O}_D,\underline{O}_D)\simeq N
\quad(\text{SGA 6 VII 2.5}),
\end{equation*}
hence
\begin{equation*}
(D,D)=\chi(\underline{O}_D)-\chi(N)=-\deg_D(N).
\end{equation*}
Finally, (iv) follows from the symmetry of the tensor product.

Notice that (ii), (iv), and (vi) remain true without assuming the divisors positive.

%%%%%%%%%% source scan idx 12 | folio 5 %%%%%%%%%%

\underline{Proposition} 1.7. \underline{Let} $(Y_i)_{1\leq i\leq k}$
\underline{be the irreducible components of} $Y$, $f$ \underline{a meromorphic
function on} $X$, \underline{and}
\begin{equation*}
\operatorname{div}(f)=\sum_i n_iY_i+Z
\end{equation*}
\underline{its divisor, with} $|Z|\cap Y$ \underline{finite. Then, putting}
$D_i=n_iY_i$, \underline{one has, for} $(a_i)\in\mathbb Q^k$:
\begin{equation*}
\left(\sum_i a_iD_i\right)^2
=-\sum_{i<j}(a_i-a_j)^2(D_i,D_j)-\sum_i a_i^2(D_i,Z).
\end{equation*}

\underline{Proof} (after [7]). By 1.6 (v), (vi),
\begin{equation*}
\begin{aligned}
\left(\sum_i a_iD_i\right)^2
&=\sum_i a_i\left(D_i,\sum_j a_jD_j\right)
=\sum_i a_i\left(D_i,\sum_j a_jD_j-a_i\operatorname{div}f\right)\\
&=\sum_i a_i\left(D_i,\sum_{j\ne i}(a_j-a_i)D_j-a_iZ\right)\\
&=\sum_{i\ne j}a_i(a_j-a_i)(D_i,D_j)-\sum_i a_i^2(D_i,Z)\\
&=-\sum_{i<j}(a_i-a_j)^2(D_i,D_j)-\sum_i a_i^2(D_i,Z).
\end{aligned}
\end{equation*}

\underline{Corollary} 1.8. \underline{Under the hypotheses of Example} 1.1,
\underline{take for} $f$ \underline{a uniformizer} $\pi$ \underline{of} $S$.
\underline{Here} $Z=0$ \underline{and}
\begin{equation*}
\left(\sum_i k_iY_i\right)^2
=-\sum_{i<j}\frac{1}{n_in_j}(k_in_j-k_jn_i)^2(Y_i,Y_j);
\end{equation*}
\underline{thus the intersection form is negative, and if} $Y$
\underline{is connected, its kernel consists precisely of the multiples of}
$\operatorname{div}(\pi)$.

\underline{Corollary} 1.9 (Mumford [7]). \underline{Under the hypotheses of Example}
1.2, \underline{take for} $f$ \underline{a nonzero element of the maximal ideal of
the local ring} $\underline O_s$. \underline{Here} $Z>0$, \underline{so the
intersection form is negative definite.}
\end{filecontents*}
% Original input: components/002_expose_X_1_10_through_section_1.tex
\begin{filecontents*}[overwrite,noheader]{sga7ii-cumulative-input-002.tex}
%%%%%%%%%% source scan idx 13 | folio 6 %%%%%%%%%%

1.10. \underline{For the remainder of this section we restrict ourselves to the
case considered in Example} 1.1.

Lichtenbaum [5] proved that, in this case, $X$ is automatically projective over
$S$, so that there is no difficulty in defining the complex
$Rf^{!}\underline{O}_S$ ([4] III 8.7, p.~190).  Since $f$ has relative dimension
1 and $X$ and $S$ are regular (Gorenstein would suffice), there exists ([4] V 8.3
and V 9.1) an invertible sheaf $\omega_{X/S}$ on $X$, unique up to unique
isomorphism, such that

\begin{equation}\tag{1.10.1}
Rf^{!}\underline{O}_S \;\simeq\; \omega_{X/S}[1] \quad .
\end{equation}
We denote by $K$ any divisor (effective or not) such that
$\omega_{X/S}\simeq\underline{O}(K)$.

Let $D$ be a divisor $\geq0$ supported on $Y$, and let $N$ be its conormal sheaf:

\begin{center}
\begin{tikzcd}[column sep=large, row sep=large]
D \arrow[rr, "i"] \arrow[dr, swap, "g"] & {} & X \arrow[dl, "f"] \\
{} & S & {}
\end{tikzcd}
\end{center}
One has

\begin{equation}\tag{1.10.2}
Rg^{!}\underline{O}_S
=Ri^{!}Rf^{!}\underline{O}_S
=Ri^{!}\omega_{X/S}[1]
=\omega_{X/S}\big|_D\otimes N^{\vee}[0].
\end{equation}
By the duality theorem ([4] V 11.1, p.~210),

\begin{equation}\tag{1.10.3}
Rg_*Rg^{!}\underline{O}_S
\;\simeq\;
\operatorname{RHom}(Rg_*\underline{O}_D,\underline{O}_S) \quad .
\end{equation}
Moreover, for every $\underline{O}_S$-module $M$ supported at the closed point,

\begin{equation}\tag{1.10.4}
\chi\bigl(\operatorname{RHom}(M,\underline{O}_S)\bigr)=-\chi(M),
\end{equation}
so that here, using (1.4.4) and (1.10.3), we obtain

\begin{equation*}
\chi(Rg^{!}\underline{O}_S)=-\chi(\underline{O}_D) \quad .
\end{equation*}
Since, by (1.4.6),

\begin{equation*}
\chi(Rg^{!}\underline{O}_S)
=\chi(\underline{O}_D)+\deg_D(Rg^{!}\underline{O}_S),
\end{equation*}

%%%%%%%%%% source scan idx 14 | folio 7 %%%%%%%%%%

we conclude from (1.10.2) that

\begin{equation*}
-2\chi(\underline{O}_D)
=\deg_D(Rg^{!}\underline{O}_S)
=\deg_D(\omega_{X/S}|_D\otimes N^{\vee})
=\deg_D(\omega_{X/S})-\deg_D(N),
\end{equation*}
which, by 1.6(ii) and (v), may be rewritten as follows.

\underline{Proposition} 1.11. \underline{Under the hypotheses of} (1.1),
\underline{for every divisor} $D\geq0$ \underline{supported on} $Y$,
\underline{one has}

\begin{equation*}
\chi(\underline{O}_D)=-\frac12(D,D+K) \quad .
\end{equation*}

1.12. We denote by $d$ the greatest common divisor of the geometric
multiplicities of the irreducible components of $X_s=f^{-1}(s)$; that is, the
greatest common divisor of the lengths of the local rings at the generic points
of the geometric special fibre of $f$.  If $X_\eta$ has a rational point, or
equivalently if $f:X\longrightarrow S$ has a section $x$, then $d=1$; more
precisely, $f$ is smooth at $x(s)$.  Indeed, $x$ is a closed immersion of regular
schemes, hence a regular immersion, and EGA IV 17.12.1 applies.

Here is Artin's proof of a special case of a theorem of Raynaud [9], which only
assumes that $d$ is prime to the characteristic $p$ of $k$.

\underline{Theorem} 1.13. \underline{With the preceding notation, if} $Y$
\underline{is geometrically connected and} $d=1$, \underline{then} $X$
\underline{is cohomologically flat over} $S$ \underline{in dimension} $0$
(EGA III 7.8.1); \underline{more precisely} (EGA III 7.6.9(ii)),
\underline{one has}

\begin{equation*}
k\xrightarrow{\;\sim\;}H^0(X_s,\underline{O}_{X_s}).
\end{equation*}
It is easy to reduce to the case in which $S$ is strictly local.  Since $Y$ is
geometrically connected, $H^0(Y,\underline{O}_Y)$ is a purely inseparable
extension of $k$; and since $(d,p)=1$, at least one irreducible component of $Y$
is

%%%%%%%%%% source scan idx 15 | folio 8 %%%%%%%%%%

geometrically reduced.  It follows readily that
$H^0(Y,\underline{O}_Y)=k$.

Let $D$ be a divisor on $X$ such that

\begin{equation*}
Y\leq D<X_s \quad .
\end{equation*}
Since $d=1$, $D$ is not a rational multiple of $X_s$; hence, by 1.8,

\begin{equation*}
(D,X_s-D)=-(D,D)>0 \quad .
\end{equation*}
The divisor $Y$ therefore has an irreducible component $E$ such that

\begin{equation*}
0<E\leq X_s-D \quad\text{and}\quad (D,E)>0 \quad .
\end{equation*}
Consider the diagram

\begin{center}
\begin{tikzcd}[column sep=large, row sep=large]
0 \arrow[r] & \underline{O}(-D-E) \arrow[r] \arrow[d, "\varphi"] &
\underline{O} \arrow[r] \arrow[d, equal] &
\underline{O}_{D+E} \arrow[r] \arrow[d, "\psi"] & 0 \\
0 \arrow[r] & \underline{O}(-D) \arrow[r] &
\underline{O} \arrow[r] & \underline{O}_D \arrow[r] & 0
\end{tikzcd}
\end{center}
The snake lemma gives an isomorphism between $\operatorname{Ker}(\psi)$ and
$\operatorname{Coker}(\varphi)$.  This cokernel is precisely the restriction of
$\underline{O}(-D)$ to $E$, whence an exact sequence

\begin{equation*}
0\longrightarrow\underline{O}(-D)|_E
\longrightarrow\underline{O}_{D+E}
\longrightarrow\underline{O}_D\longrightarrow0 \quad .
\end{equation*}
By construction,
$\deg_E(\underline{O}(-D)|_E)=-(D,E)<0$, so that
$H^0(\underline{O}(-D)|_E)=0$ and the restriction map

\begin{equation*}
0\longrightarrow H^0(\underline{O}_{D+E})
\longrightarrow H^0(\underline{O}_D)
\end{equation*}
is injective.  This allows us, by induction, to construct a sequence of divisors

\begin{equation*}
Y=D_0\leq D_1\leq\cdots\leq D_n=X_s
\end{equation*}
such that $H^0(\underline{O}_{D_i})$ injects into
$H^0(\underline{O}_{D_{i-1}})$.  Thus
$H^0(\underline{O}_{X_s})=H^0(\underline{O}_Y)=k$, which completes the proof of
1.13.

%%%%%%%%%% source scan idx 16 | folio 9 %%%%%%%%%%

The following fact is proved in [5] or [10].

\underline{Theorem} 1.14 (Castelnuovo, Enriques, \ldots).
\underline{Under the hypotheses of} (1.1), \underline{let} $D$
\underline{be a reduced irreducible divisor contained in} $Y$.
\underline{Suppose that} $H^1(D,\underline{O}_D)=0$ \underline{and that, on putting}
$k'=H^0(D,\underline{O}_D)$, \underline{one has}

\begin{equation*}
(D,D)=-[k':k] \quad .
\end{equation*}
\underline{Then} $D$ \underline{is a projective line over} $k'$, \underline{and
there exist a regular} $S$\underline{-scheme} $X_1$, \underline{a closed point}
$x_1$ \underline{of} $X_1$ \underline{with residue field} $k'$, \underline{and an
isomorphism from} $X$ \underline{to the} $S$\underline{-scheme obtained from}
$X_1$ \underline{by blowing up} $x_1$, \underline{under which} $D$
\underline{is identified with the fibre over} $x_1$.

A divisor $D$ having these properties is called an \underline{exceptional curve
of the first kind}.

I cannot resist the pleasure of proving the following proposition of Néron [8],
Proposition 1$'$, p.~90.

\underline{Proposition} 1.15. \underline{Suppose that} $X_\eta$
\underline{is an elliptic curve (in particular, that it has a section), and that}
$X$ \underline{contains no exceptional curve of the first kind.  Then, with the
notation of} (1.10.1), \underline{the canonical map}

\begin{equation*}
f^*f_*\omega_{X/S}\xrightarrow{\;\sim\;}\omega_{X/S}
\end{equation*}
\underline{is universally an isomorphism}.

Since $d=1$ (1.12), by (1.13) it is enough to prove that $\omega_{X/S}$ is
isomorphic to $\underline{O}_X$.  Since $\omega_{X/S}$ is trivial on the elliptic
curve $X_\eta$, we know a priori that

\begin{equation*}
\omega_{X/S}\simeq\underline{O}\Bigl(\sum_i n_iY_i\Bigr),
\end{equation*}

%%%%%%%%%% source scan idx 17 | folio 10 %%%%%%%%%%

where the $Y_i$ are the irreducible components of $Y$; and since
$\operatorname{div}(\pi)\sim0$, we may assume the $n_i$ positive.  If
$\sum_i n_iY_i$ is proportional to $\operatorname{div}(\pi)$, it is an integral
multiple of $\operatorname{div}(\pi)$ and we are done.  Otherwise,

\begin{equation*}
\Bigl(\sum_i n_iY_i\Bigr)^2<0,
\end{equation*}
there is an $i$ such that

\begin{equation*}
\Bigl(Y_i,\sum_j n_jY_j\Bigr)<0,
\end{equation*}
and the following lemma applies.

\underline{Lemma} 1.16. \underline{Under the hypotheses of} 1.1,
\underline{suppose that} $Y$ \underline{is connected and that either} $Y$
\underline{is reducible or} $\chi(\underline{O}_Y)\leq0$.
\underline{Then an irreducible component} $Y_i$ \underline{of} $Y$
\underline{is an exceptional curve of the first kind if and only if}
$(Y_i,K)<0$.

Let $k'=H^0(Y_i,\underline{O}_{Y_i})$.  The Riemann--Roch formula 1.11 gives

\begin{equation*}
\chi(\underline{O}_{Y_i})
=[k':k]\bigl(1-\dim_{k'}H^1(Y_i,\underline{O}_{Y_i})\bigr)
=-\frac12\bigl((Y_i,Y_i)+(Y_i,K)\bigr) \quad .
\end{equation*}
If $Y_i$ is exceptional of the first kind, it follows that
$(Y_i,K)=-[k':k]<0$.  Conversely, if $(Y_i,K)<0$, then
$\chi(\underline{O}_{Y_i})>0$, which implies
$\chi(\underline{O}_{Y_i})=[k':k]$ and that $Y$ is reducible.  Hence, by 1.8,
$(Y_i,Y_i)<0$.  The negative integers $(Y_i,Y_i)$ and $(Y_i,K)$ are divisible by
$[k':k]$, and therefore necessarily

\begin{equation*}
(Y_i,Y_i)=(Y_i,K)=-[k':k] \quad .
\end{equation*}

\underline{Exercise}. (a) Under the hypotheses of (1.1), suppose that $X_\eta$ is
a principal homogeneous space under an elliptic curve and that $X$ contains no
exceptional curve of the first kind.  Show that on $X$ the Riemann--Roch formula
(1.11) takes the form

\begin{equation*}
\chi(\underline{O}_D)=-\frac12(D,D).
\end{equation*}

%%%%%%%%%% source scan idx 18 | folio 11 %%%%%%%%%%

(b) Suppose in addition that the residue field $k(s)$ is algebraically closed.
Show that the diagram of irreducible components of $X_s$ is a multiple of one of
the diagrams considered by Néron [8], pp.~124--125, fourth column (cases a through
c8).  The diagram determines the genera of the reduced irreducible components,
their multiplicities in $X_s$, and the nature and disposition of the singularities
of $(X_s)_{\mathrm{red}}$.

First treat the case in which $X_s$ is irreducible.  Once this case has been
excluded, prove that every reduced irreducible component is a projective line of
self-intersection $-2$.  If $X_s=\sum_i n_iD_i$ is the decomposition of $X_s$
into reduced irreducible divisors, rewrite this condition as

\begin{equation}\tag{*}
n_i=\frac12\sum_{j\ne i}n_j(D_i,D_j) \quad .
\end{equation}
On the other hand, $X_s$ is connected, so that:

(**) For any $i$ and $j$, there exists a chain
$i=i_1,i_2,\ldots,i_n=j$ with
$(D_{i_k},D_{i_{k+1}})\ne0$.

Now find all solutions in integers of $(*)+(**)$.  First show that, except in
Néron's cases b2 and c2, the intersection numbers $(D_i,D_j)$ are either zero or
one.

\end{filecontents*}
% Original input: components/003_expose_X_section_2_opening_through_2_9_statement.tex
\begin{filecontents*}[overwrite,noheader]{sga7ii-cumulative-input-003.tex}
2. \underline{Italian intersections}

Let $S'$ and $S$ be two smooth projective surfaces over a field $k$, and let
$f:S'\longrightarrow S$ be a birational morphism.  If $D$ and $E$ are two
divisors on $S$, then

\begin{equation*}
(D,E)=(f^*D,f^*E),
\end{equation*}

%%%%%%%%%% source scan idx 19 | folio 12 %%%%%%%%%%

as one sees, for example, by moving $D$ and $E$.

It is therefore natural to seek a ``birational'' expression for the intersection
number.  I propose to explain how the Italians proceeded.  My principal source
of information is [1], which contains a long bibliography.  See also Manin [6]
and Northcott [12], [13].

(2.0) Let $S$ be the spectrum of a regular local ring of dimension 2, let $s$ be
the closed point of $S$, and let $k$ be a subfield of $k(s)$ over which $k(s)$ is
finite.  Intersection numbers and Euler--Poincaré characteristics will be
computed over $k$.

In this section, an $S$-scheme $p:S'\longrightarrow S$ will be called a
\underline{birational transform of} $S$ if it is regular and connected, and if
$p$ is proper and induces an isomorphism
$S'\setminus p^{-1}(s)\longrightarrow S\setminus\{s\}$.  Everything said about
$S$ will also apply to the localizations of its birational transforms at a
closed point.

(2.1) Let $f:(S'',p'')\longrightarrow(S',p')$ be a morphism of $S$-schemes
between birational transforms of $S$.  It is known [5] that $f$ has a unique
factorization

\begin{equation*}
f=f_n\cdots f_1,\qquad f_i:S_{i+1}\longrightarrow S_i,
\end{equation*}
such that $S_{i+1}$ is obtained from $S_i$ by blowing up a finite set $X_i$ of
closed points, and $f_i(X_{i+1})\subset X_i$.  The set $X_i$ is the set of points
of $S_i$ at which the projection of $S''$ onto $S_i$ is not an isomorphism.

In particular, $f$ induces an isomorphism over the complement of a finite subset
of $S'$ containing the generic points of the irreducible components of
$p'^{-1}(s)$.  Hence $f^{-1}$ defines an injection

%%%%%%%%%% source scan idx 20 | folio 13 %%%%%%%%%%

from the set $\operatorname{Irr}(p'^{-1}(s))$ of irreducible components of
$p'^{-1}(s)$ into $\operatorname{Irr}(p''^{-1}(s))$.

The birational transforms of $S$ have no $S$-automorphisms.  We write $I$ for the
set of $S$-isomorphism classes of birational transforms of $S$, ordered by
domination.  The set $I$ is filtered to the right and has least element $S$.

\underline{Definition} 2.2. \underline{The set} $S^*$ \underline{of points of}
$S$ \underline{infinitely near to} $s$ \underline{is the direct-limit set}

\begin{equation*}
S^*=\varinjlim_I\operatorname{Irr}(p^{-1}(s)) \quad .
\end{equation*}
If an element $t\in S^*$ is represented on a birational transform $S'$ of $S$ by
a complete reduced irreducible curve $T$, the field
$H^0(T,\underline{O}_T)$ depends only on $t$, not on the choice of $S'$.  It is
called the \underline{residue field} $k(t)$ of $t$.

Let $S_1$ be the blow-up of $S$ along $s$.  The projective line $s^*$, the inverse
image of $\{s\}$ in $S_1$, is called the \underline{curvetta} of $s$ and defines
an element $s^*$ of $S^*$ with residue field $k(s)$.

Let $S'$ be a birational transform of $S$, let $t$ be a closed point of $S'$, and
let $S'_t$ be the spectrum of the local ring of $S'$ at $t$.  There is an evident
injection $(S'_t)^*\hookrightarrow S^*$; in particular, $t$ defines an element
$t^*$ of $S^*$ with the same residue field.  We write $(S')^*$ for the union in
$S^*$ of the $(S'_t)^*$ as $t$ ranges over the closed points of $S'$.

The following proposition follows from the description 2.1 of the birational
transforms of $S$.

%%%%%%%%%% source scan idx 21 | folio 14 %%%%%%%%%%

\underline{Proposition} 2.3. \underline{For every} $t\in S^*$,
\underline{there is a unique sequence of birational transforms}
$(S_i)_{0\leq i\leq n+1}$ \underline{of} $S$ \underline{and closed points}
$t_i\in S_i$ $(0\leq i\leq n)$ \underline{such that}

(i)\quad $S_0=S$, \underline{and} $S_{i+1}$ \underline{is the blow-up of} $S_i$
\underline{along} $\{t_i\}$;

(ii)\quad $t_{i+1}\in t_i^*$ $(0\leq i<n)$, \underline{and} $t=t_n^*$.

\underline{Moreover, for every birational transform} $(S',p')$ \underline{of}
$S$, \underline{the element} $t$ \underline{belongs to the image of}
$\operatorname{Irr}(p'^{-1}(s))$ \underline{in} $S^*$ \underline{if and only if}
$S'$ \underline{dominates} $S_{n+1}$.

The integer $n$ is called the order of $t$.  We write $S(t)$ (respectively
$S_+(t)$) for the birational transform $S_n$ (respectively $S_{n+1}$) of $S$.
By abuse of notation, we also write $t$ for the closed point $t_n$ of $S_n$, and
$t^*$ for its inverse image $t_n^*$ in $S_{n+1}$.

We equip $S^*$ with the following order:

\begin{equation*}
t_1\prec t_2
\quad\Longleftrightarrow\quad
S_+(t_2)\ \text{dominates}\ S_+(t_1) \quad .
\end{equation*}
If $t_1\prec t_2$, then $k(t_2)$ is an extension of $k(t_1)$.  The ordered set
$S^*$ has least element $s^*$, and every bounded subset of $S^*$ is finite and
totally ordered ($S^*$ is a ``tree rooted at $s^*$'').  Moreover, the order $n$ of
$t\in S^*$ is the distance in the tree $S^*$ from $s^*$ to $t$.

The following proposition makes the description 2.1 of the birational
transforms of $S$ explicit.

\underline{Proposition} 2.4. \underline{The function assigning to a birational
transform} $p:S'\to S$ \underline{the subset}
$\operatorname{Irr}(p^{-1}(s))$ \underline{of} $S^*$ \underline{induces a
bijection between} $I$ (2.1) \underline{and the finite subsets of} $S^*$
\underline{which, together with each element,}

%%%%%%%%%% source scan idx 22 | folio 15 %%%%%%%%%%

\underline{contain all its predecessors}.

Let $D$ be a divisor $\geq0$ on $S$, and let $t\in S^*$.  We say that $t$
\underline{belongs to} $D$ if the closed point $t$ of $S(t)$ (2.3) belongs to the
strict transform (the schematic closure of $D-\{s\}$) of $D$ on $S(t)$.

\underline{Definition} 2.5. \underline{Let} $t,t'\in S^*$.
\underline{We say that} $t'$ \underline{is proximate to} $t$ \underline{if}
$t'\in S_+(t)^*$ (2.2) \underline{and belongs to the curvetta} $t^*$ \underline{of}
$t$ (2.3).

If $t'$ is an immediate successor of $t$, then $t'$ is proximate to $t$.  If in a
sequence $(t_i)_{0\leq i\leq n}$ every $t_{i+1}$ is an immediate successor of
$t_i$, and if $t_n$ is proximate to $t_0$, then every $t_i$ with $0<i\leq n$ is
proximate to $t_0$, while $t_n$ is not proximate to $t_i$ for
$0<i\leq n-1$.

(2.6) Let $A$ be a Noetherian local ring of dimension $n$, let $T$ be its
spectrum, let $m$ be its maximal ideal, let $k=A/m$ be its residue field, and let
$P=\sum_{0\leq i<n}a_iX^i$ be its Hilbert--Samuel polynomial:

\begin{equation*}
P(r)=\dim_k(m^r/m^{r+1})\qquad\text{for sufficiently large }r.
\end{equation*}
Recall that the multiplicity $\mu_s$ of the closed point $s$ in $T$ is defined by

\begin{equation*}
a_{n-1}=\frac{\mu_s}{(n-1)!} \quad .
\end{equation*}
When $n=1$, this simply says

\begin{equation*}
\mu_s=\dim_k(m^r/m^{r+1})\qquad\text{for sufficiently large }r.
\end{equation*}
Geometrically, $\mu_s$ may be defined as the degree, in the projective space
$\mathbf P_k(m/m^2)$, of the special fibre of the blow-up of $T$ along the closed
subscheme $\{s\}$.

%%%%%%%%%% source scan idx 23 | folio 16 %%%%%%%%%%

Recall that if $Y$ and $Z$ are two closed subschemes of a scheme $X$, then the
strict transform of $Y$ in the blow-up $\widetilde X$ of $X$ along $Z$---the
schematic closure of $Y\setminus Z$ in $\widetilde X$---is precisely the blow-up
of $Y$ along $Y\cap Z$.

Suppose that $T$ is a closed subscheme of a regular scheme $U$.  Let
$p:U_1\longrightarrow U$ be the blow-up of $U$ along $\{s\}$, let $s^*$ be the
inverse image of $s$ in $U_1$, and let $T'$ be the strict transform of $T$ in
$U_1$.  The preceding discussion shows that the multiplicity of $T$ at $s$ is
the degree, in the projective space $s^*$, of the subscheme $s^*\cap T'$.

Suppose further that $T$ is a divisor on $U$.  Let $T^*$ be its total transform
in $U_1$ (the inverse-image divisor), and define $\nu$ by
$T^*=\nu s^*+T'$.  If $X$ is a line in $s^*$, and Euler--Poincaré
characteristics are calculated as alternating sums of dimensions over $k$, then

\begin{equation}\tag{2.6.1}
(X,s^*)=
\chi(\underline{O}_X\overset{\mathbb L}{\otimes}\underline{O}_{s^*})
=\deg_X\underline{O}(s^*)=-1,
\end{equation}
and we have seen that

\begin{equation}\tag{2.6.2}
(X,T')=
\chi(\underline{O}_X\overset{\mathbb L}{\otimes}\underline{O}_{T'})
=\deg_{s^*}(s^*\cap T')=\mu_s \quad .
\end{equation}
The projection formula

\begin{equation*}
Rp_*(\underline{O}_X\overset{\mathbb L}{\otimes}\underline{O}_{T^*})
=Rp_*(\underline{O}_X\overset{\mathbb L}{\otimes}Lp^*\underline{O}_T)
=Rp_*\underline{O}_X\overset{\mathbb L}{\otimes}\underline{O}_T
\end{equation*}
and its corollary

\begin{equation*}
(X,T^*)=\nu(X,s^*)+(X,T')=0
\end{equation*}
give, by (2.6.1) and (2.6.2), the identity

\begin{equation}\tag{2.6.3}
\mu_s=\nu \quad .
\end{equation}

%%%%%%%%%% source scan idx 24 | folio 17 %%%%%%%%%%

\underline{Proposition} 2.7. \underline{Let} $T$ \underline{be a divisor}
$\geq0$ \underline{on a regular scheme} $U$, $s$ \underline{a closed point of}
$U$, $U_1$ \underline{the blow-up of} $U$ \underline{along} $\{s\}$, $s^*$
\underline{the inverse image of} $s$ \underline{in} $U_1$, $T'$ \underline{the
strict transform of} $T$ \underline{on} $U_1$, \underline{and} $T^*$
\underline{its total transform}.

\underline{The following three quantities are equal}:

(i)\quad \underline{the multiplicity, in Samuel's sense, of the point} $s$
\underline{on} $T$;

(ii)\quad \underline{the degree, in the projective space} $s^*$,
\underline{of} $s^*\cap T'$;

(iii)\quad \underline{the multiplicity of} $s^*$ \underline{in the divisor}
$T^*$.

2.8. Return to the hypotheses of 2.0.  Let $D$ be a divisor $\geq0$ on $S$ and
let $t\in S^*$.  The \underline{multiplicity} $\mu_t(D)$ \underline{of} $D$
\underline{at} $t$ is, by definition, the multiplicity at $t$ of the strict
transform of $D$ on $S(t)$ (2.3).  The support of the function
$t\mapsto\mu_t(D)$ is the set of points infinitely near to $s$ that belong to
$D$ (2.4).  When $D$ is regular, the numbers
$\mu_t(D)[k(t):k(s)]$ are zero or one: indeed, the strict transforms $D'$ of $D$
are isomorphic to $D$, and one applies the definition (2.6) of multiplicity.

Intersection numbers can be calculated in terms of multiplicities.

\underline{Theorem} 2.9. \underline{If the divisors} $D$ \underline{and} $E$
\underline{meet only at} $s$, \underline{then}

\begin{equation*}
(D,E)=\sum_{t\in S^*}\mu_t(D)\mu_t(E)[k(t):k],
\end{equation*}
\underline{where the sum is over the points infinitely near to} $s$
\underline{which belong to both} $D$ \underline{and} $E$.
\end{filecontents*}
% Original input: components/004_expose_X_theorem_2_9_proof_through_2_11.tex
\begin{filecontents*}[overwrite,noheader]{sga7ii-cumulative-input-004.tex}
\underline{Proof}. Let $p:S'\longrightarrow S$ be a birational transform of $S$.
We denote the total transforms of $D$ and $E$ by $D^*$ and $E^*$, and their
strict transforms by $D'$ and $E'$.

%%%%%%%%%% source scan idx 25 | folio 18 %%%%%%%%%%

\underline{Lemma} 2.9.1. (i) \underline{The total transform} $D^*$
\underline{is characterized by}:

(a)\quad $D^*-D'$ \underline{is supported on} $p^{-1}(s)$;

(b)\quad \underline{if} $X$ \underline{is supported on} $p^{-1}(s)$,
\underline{then} $(D^*,X)=0$.

(ii)\quad $(D,E)=(D^*,E^*)=(D^*,E')=(D',E^*)$.

Assertion (i)(a) is immediate, and (i)(b) follows from 1.6(vi).  That (a) and (b)
determine $D^*$ follows from 1.9.  By construction, the cokernel of the map
$\underline{O}_E\longrightarrow p_*\underline{O}_{E'}$ is finite; hence

\begin{equation*}
\begin{aligned}
(D^*,E')
&=\chi(Lp^*\underline{O}_D\otimes\underline{O}_{E'})\\
&=\chi(\underline{O}_D\otimes Rp_*\underline{O}_{E'})
=\chi(\underline{O}_D\otimes\underline{O}_E)
=(D,E),
\end{aligned}
\end{equation*}
which, by (i), implies (ii).

Take for $S'$ the blow-up of $S$ at $s$.  By (2.9.1), (2.6.3), and 2.7,

\begin{equation}\tag{2.9.2}
\begin{aligned}
(D,E)
&=(D^*,E')
=(D'+\mu_s(D)s^*,E')\\
&=(D',E')+\mu_s(D)(s^*,E')\\
&=(D',E')+\mu_s(D)\mu_s(E)[k(s):k].
\end{aligned}
\end{equation}
Theorem 2.9 now follows by induction on $(D,E)$, through repeated application of
(2.9.2).

(2.10) Let $A$ be a reduced Noetherian local ring of dimension one, let $s$ be the
closed point of $\operatorname{Spec}(A)$, let $\widetilde A$ be the normalization
of $A$, and let $k$ be a field over which $k(s)$ is finite.  Suppose that
$\widetilde A$ is finite over $A$ (equivalently, that the completion $\widehat A$
is reduced).  The \underline{$\delta$-invariant of the singularity of}
$\operatorname{Spec}(A)$ \underline{at} $s$ is the length over $k$ of the
$A$-module $\widetilde A/A$.

%%%%%%%%%% source scan idx 26 | folio 19 %%%%%%%%%%

\underline{Theorem} 2.11. \underline{Let} $D$ \underline{be a reduced curve on}
$S$. \underline{Then}:

(i)\quad \underline{for every} $t\in S^*$, \underline{with the notation of} 2.8,
\underline{one has}

\begin{equation*}
\mu_t(D)
=\sum_{\substack{t'\ \mathrm{proximate}\\\mathrm{to}\ t}}
\mu_{t'}(D)[k(t'):k(t)] \quad .
\end{equation*}

(ii)\quad \underline{Suppose that the normalization} $\widetilde D$
\underline{of} $D$ \underline{is finite over} $D$, \underline{and let}
$\delta=\chi(\underline{O}_{\widetilde D}/\underline{O}_D)$
\underline{be the invariant of the singularity of} $D$ \underline{at} $s$.
\underline{For all but finitely many} $t\in S^*$, $\mu_t(D)$ \underline{is zero
or one, and}

\begin{equation*}
\delta
=\sum_{t\in S^*}\frac12\mu_t(D)\bigl(\mu_t(D)-1\bigr)[k(t):k] \quad .
\end{equation*}

Let $p:S_1\longrightarrow S$ be the blow-up of $S$ along $\{s\}$, let $D^*$ be
the total transform of $D$, and let $D'$ be its strict transform.  By (2.7),

\begin{equation}\tag{2.11.1}
D^*=D'+\mu_s(D)s^*
\end{equation}
and
\begin{equation}\tag{2.11.2}
\mu_s(D)=(s^*,D')\,[k(s):k]^{-1} \quad .
\end{equation}
The curvetta $s^*$ is nonsingular.  Thus, by 2.8 and Definition 2.5, formula 2.9,
applied to the divisors $s^*$ and $D'$ on $S_1$, becomes

\begin{equation}\tag{2.11.3}
(s^*,D')
=\sum_{\substack{t'\ \mathrm{proximate}\\\mathrm{to}\ s}}
\mu_{t'}(D)[k(t'):k] \quad .
\end{equation}
Assertion (i) for $s=t$ follows from (2.11.2) and (2.11.3); the general case is
obtained by applying this result to the localization of $S(t)$ at $t$ (2.3).

Let $m$ be the maximal ideal defining $\{s\}$, and let $n\geq0$.  It is known
(SGA VII 3.5) that

\begin{equation}\tag{2.11.4}
p_*\underline{O}_{S_1}(-ns^*)=m^n
\quad\text{and}\quad
R^ip_*\underline{O}_{S_1}(-ns^*)=0\quad(i>0).
\end{equation}

%%%%%%%%%% source scan idx 27 | folio 20 %%%%%%%%%%

It follows from this and from the exact sequence

\begin{equation*}
0\longrightarrow\underline{O}_{S_1}(-ns^*)
\longrightarrow\underline{O}_{S_1}
\longrightarrow\underline{O}_{ns^*}
\longrightarrow0
\end{equation*}
that

\begin{equation}\tag{2.11.5}
p_*\underline{O}_{ns^*}=\underline{O}_S/m^n
\quad\text{and}\quad
R^ip_*\underline{O}_{ns^*}=0\quad(i>0),
\end{equation}
and that

\begin{equation}\tag{2.11.6}
\chi(\underline{O}_{ns^*})
=\chi(\underline{O}_S/m^n)
=\frac{n(n+1)}2[k(s):k] \quad .
\end{equation}
Put $n=0$ in (2.11.4).  For every complex $K$ of coherent sheaves on $S$, the
projection formula gives

\begin{equation*}
K\xrightarrow{\;\sim\;}Rp_*Lp^*K
\qquad
\bigl(\simeq K\overset{\mathbb L}{\otimes}Rp_*\underline{O}_{S_1}\bigr).
\end{equation*}
In particular, since
$Lp^*\underline{O}_D=p^*\underline{O}_D=\underline{O}_{D^*}$, one has

\begin{equation}\tag{2.11.7}
p_*\underline{O}_{D^*}=\underline{O}_D
\quad\text{and}\quad
R^ip_*\underline{O}_{D^*}=0\quad(i>0).
\end{equation}
Put $E=\mu_s(D)s^*$.  By (2.11.1), (2.9.1)(i)(b), and (2.6.1),

\begin{equation}\tag{2.11.8}
\begin{aligned}
\chi(\underline{O}_{D'\cap E})
&=(D',E)=(D^*-E,E)=-(E,E)\\
&=-\mu_s(D)^2(s^*,s^*)
=\mu_s(D)^2[k(s):k].
\end{aligned}
\end{equation}
The exact sequence

\begin{equation*}
0\longrightarrow\underline{O}_{D^*}
\longrightarrow\underline{O}_{D'}\oplus\underline{O}_E
\longrightarrow\underline{O}_{D'\cap E}
\longrightarrow0
\end{equation*}
gives, by direct image and (2.11.7), the exact sequence

\begin{equation*}
0\longrightarrow\underline{O}_D
\longrightarrow p_*\underline{O}_{D'}\oplus p_*\underline{O}_E
\longrightarrow p_*\underline{O}_{D'\cap E}
\longrightarrow0,
\end{equation*}
and hence

\begin{equation*}
0\longrightarrow p_*\underline{O}_E
\longrightarrow p_*\underline{O}_{D'\cap E}
\longrightarrow p_*\underline{O}_{D'}/\underline{O}_D
\longrightarrow0 \quad .
\end{equation*}

%%%%%%%%%% source scan idx 28 | folio 21 %%%%%%%%%%

Consequently, by (2.11.6) and (2.11.8),

\begin{equation}\tag{2.11.9}
\begin{aligned}
\chi(p_*\underline{O}_{D'}/\underline{O}_D)
&=\mu_s(D)^2[k(s):k]
-\frac12\mu_s(D)\bigl(\mu_s(D)+1\bigr)[k(s):k]\\
&=\frac12\mu_s(D)\bigl(\mu_s(D)-1\bigr)[k(s):k] \quad .
\end{aligned}
\end{equation}
The schemes $D$ and $D'$ are reduced, differ only over $s$, and $D'$ is finite
over $D$.  They therefore have the same normalization, and the $\delta$-invariant
of the singularity of $D$ at $s$ is the sum of
$\chi(p_*\underline{O}_{D'}/\underline{O}_D)$ and the invariants $\delta$ of the
singularities of $D'$ at the points of $D'\cap s^*$.  We prove 2.11(ii) by
induction on $\delta$.  If $\mu_s(D)>1$, the induction hypothesis applies to
$D'$, and 2.11(ii) follows from (2.11.9).  If $\mu_s(D)=1$, then
$(D',s^*)=1$ and $D'\cap s^*$ is a regular divisor on $D'$; by EGA IV 19.1.1,
$D'$ is regular.  Formula (2.11.9) then shows that $D$ is itself regular, and
2.11(ii) follows from 2.8.

\end{filecontents*}
% Original input: components/005_expose_X_2_12_through_2_14.tex
\begin{filecontents*}[overwrite,noheader]{sga7ii-cumulative-input-005.tex}
2.12. Let $t\in S^*$, and let $S'$ be a birational transform of $S$ which
dominates $S_+(t)$ (2.3).  The \underline{total transform} of $t$ on $S'$,
again denoted by $t^*$, is the total transform on $S'$ of the divisor $t^*$ on
$S_+(t)$.

Let $\nu:S^*\longrightarrow\mathbf N$ be a function of finite support, and
suppose that, for every $t\in S^*$,

\begin{equation}\tag{2.12.1}
\nu(t)\geq
\sum_{\substack{t'\ \mathrm{proximate}\\\mathrm{to}\ t}}
\nu(t')[k(t'):k(t)]
\end{equation}
(the proximate-point inequality).  Let $f:S'\longrightarrow S$ be a birational
transform of $S$ which dominates $S_+(t)$ for each $t$ in the support of $\nu$,
and let
\[
\nu^*=\sum_{t\in S^*}\nu(t)t^*
\]
be the corresponding divisor on $S'$.

%%%%%%%%%% source scan idx 29 | folio 22 %%%%%%%%%%

We write $I(\nu)$ for the ideal of $\underline{O}_S$ obtained by direct image
under $f$ of the ideal of $\underline{O}_{S'}$ defining $\nu^*$:

\begin{equation*}
I(\nu)=f_*\underline{O}_{S'}(-\nu^*) \quad .
\end{equation*}
By (2.11.4), this ideal does not depend on the choice of $S'$.  The ideals of
$S$ obtained in this way are those which Zariski [11] calls ``complete.''

\underline{Theorem} 2.13. \underline{Under the preceding hypotheses, one has}

\begin{equation*}
\chi(\underline{O}_S/I(\nu))
=\sum_{t\in S^*}\frac{\nu(t)(\nu(t)+1)}2[k(t):k],
\end{equation*}
\underline{and}

\begin{equation*}
R^if_*\underline{O}_{S'}(-\nu^*)=0
\quad\text{for }i>0.
\end{equation*}

In old French, this is expressed briefly by saying: ``requiring a plane curve to
pass $n$ times through a point $t$, whether proper or infinitely near, imposes
$n(n+1)/2$ conditions, generally independent, on the coefficients of its
equation.''

We first prove the following lemma.

\underline{Lemma} 2.14. \underline{Let} $p:S_1\longrightarrow S$
\underline{be the blow-up of} $S$ \underline{along} $\{s\}$; \underline{let}
$T$ \underline{be a closed subscheme of} $S_1$ \underline{with ideal} $q$;
\underline{let} $n$ \underline{be an integer}; \underline{let}
$J=\underline{O}_{S_1}(-s^*)$ \underline{be the ideal defining the divisor}
$s^*$; \underline{put}
$Y=\operatorname{Spec}(\underline{O}_{S_1}/J^nq)$; \underline{and let} $D$
\underline{be a regular divisor on} $S$, \underline{with strict transform} $D'$
\underline{on} $S_1$. \underline{Suppose that} $T\cap s^*$ \underline{and}
$T\cap D'$ \underline{are finite}.

(i)\quad \underline{If}
$\operatorname{length}_{k(s)}(T\cap s^*)\leq n+1$, \underline{then}

\begin{equation*}
R^ip_*(qJ^n)=0\quad\text{for }i>0,
\end{equation*}
\underline{and, if} $T$ \underline{is Artinian,}

\begin{equation*}
\operatorname{length}_{k(s)}
\bigl(\underline{O}_S/p_*(qJ^n)\bigr)
=\frac{n(n+1)}2+\operatorname{length}_{k(s)}(T).
\end{equation*}

%%%%%%%%%% source scan idx 30 | folio 23 %%%%%%%%%%

(ii)\quad \underline{If}
$\operatorname{length}_{k(s)}(T\cap s^*)\leq n$, \underline{then}

\begin{equation*}
\operatorname{length}_{k(s)}
\Bigl(D\cap
\operatorname{Spec}\bigl(\underline{O}_S/p_*(qJ^n)\bigr)\Bigr)
=n+\operatorname{length}_{k(s)}(D'\cap T) \quad .
\end{equation*}

Consider the exact sequence on $S_1$

\begin{equation}\tag{2.14.1}
0\longrightarrow J^nq
\longrightarrow J^n
\xrightarrow{\,u\,}J^n\otimes\underline{O}_T
\longrightarrow0 \quad .
\end{equation}
To prove that $R^ip_*(J^nq)=0$ for $i>0$, it is enough, by (2.11.4), to show
that $p_*(u)$ is surjective.  By the theorem on formal functions, it suffices to
prove that the maps

\begin{equation*}
u_k:p_*(J^n\otimes\underline{O}/J^k)
\longrightarrow
p_*(J^n\otimes\underline{O}_T\otimes\underline{O}/J^k)
\end{equation*}
are surjective.

Write $\operatorname{Gr}_J$ for the associated graded object of the $J$-adic
filtration.  The module $J^m/J^{m+1}$ is isomorphic to
$\underline{O}_{s^*}(m)$; hence $R^1p_*(J^m/J^{m+1})=0$, and the map
\[
p_*(J^n/J^{n+k+1})\longrightarrow p_*(J^n/J^{n+k})
\]
is surjective.  Therefore, to prove by induction on $k$ that the maps $u_k$ are
surjective, it suffices to prove the surjectivity of

\begin{equation*}
v_k:p_*\bigl(J^n\otimes\operatorname{Gr}_J^k(\underline{O})\bigr)
\longrightarrow
p_*\bigl(J^n\otimes\operatorname{Gr}_J^k(\underline{O}_T)\bigr)
\qquad(k\geq0).
\end{equation*}
The module $\operatorname{Gr}_J^k(\underline{O}_T)$ is a quotient of
$J^k/J^{k+1}\otimes\underline{O}_T$, and thus has length at most $n+1$ over
$k(s)$.  The kernel $N$ of the surjective map

\begin{equation*}
v'_k:J^n\otimes\operatorname{Gr}_J^k(\underline{O})
\simeq\underline{O}_{s^*}(n+k)
\longrightarrow
J^n\otimes\operatorname{Gr}_J^k(\underline{O}_T)
\end{equation*}
is therefore an invertible sheaf on the projective line $s^*$ of degree
$\geq-1$.  Hence $R^1p_*N=0$, and $v_k=p_*(v'_k)$ is surjective.

If $T$ is Artinian, (2.14.1) now gives

\begin{equation*}
\operatorname{length}_{k(s)}
\bigl(p_*J^n/p_*(J^nq)\bigr)
=\operatorname{length}_{k(s)}(\underline{O}_T),
\end{equation*}

%%%%%%%%%% source scan idx 31 | folio 24 %%%%%%%%%%

and 2.14(i) follows from (2.11.4) and (2.11.6).

We now prove (ii).  Since $D$ is regular, its strict transform $D'$ is
isomorphic to $D$ and meets $s^*$ with length one.  The assumption
$\operatorname{length}_{k(s)}(T\cap s^*)\leq n$ therefore implies
\[
\operatorname{length}_{k(s)}\bigl((D'\cup T)\cap s^*\bigr)\leq n+1,
\]
so that $D'\cup T$ satisfies the hypothesis of (i).

The projection $p$ induces an isomorphism between $D$ and $D'$.  Under this
isomorphism, the ideal of $\underline{O}_D$ defining the subscheme
\[
D\cap\operatorname{Spec}\bigl(\underline{O}_S/p_*(qJ^n)\bigr)
\]
corresponds to the ideal of $D'$ generated by the global sections of $qJ^n$ on
$S_1$.  This ideal is contained in the ideal of $D'$ defining $D'\cap Y$.
It is in fact equal to it, because every section of
$\underline{O}_{D'\cup Y}$ extends to a section of $\underline{O}_{S_1}$:
if $I_{D'\cup Y}$ is the ideal defining $D'\cup Y$, the obstruction to
extension lies in $R^1p_*(I_{D'\cup Y})$.  One has

\begin{equation*}
\begin{aligned}
I_{D'\cup Y}
&=\underline{O}(-D')\cap qJ^n\\
&=(\underline{O}(-D')\cap J^n)\cap qJ^n\\
&=\underline{O}(-D')J^n\cap qJ^n\\
&=(\underline{O}(-D')\cap q)J^n,
\end{aligned}
\end{equation*}
and by (i), applied to $D'\cup T$, there is no obstruction.

Thus

\begin{equation}\tag{2.14.2}
\operatorname{length}_{k(s)}
\Bigl(D\cap
\operatorname{Spec}\bigl(\underline{O}_S/p_*(qJ^n)\bigr)\Bigr)
=\operatorname{length}_{k(s)}(D'\cap Y) \quad .
\end{equation}
The isomorphism

\begin{equation*}
J^n\otimes
\bigl(\underline{O}_{D'}/q\underline{O}_{D'}\bigr)
\xrightarrow{\;\sim\;}
\underline{O}_{D'}J^n/
\underline{O}_{D'}qJ^n
\end{equation*}
gives the exact sequence

\begin{equation*}
0\longrightarrow
J^n\otimes\underline{O}_{D'\cap T}
\longrightarrow\underline{O}_{D'\cap Y}
\longrightarrow\underline{O}_{D'\cap ns^*}
\longrightarrow0,
\end{equation*}
and hence

%%%%%%%%%% source scan idx 32 | folio 25 %%%%%%%%%%

\begin{equation}\tag{2.14.3}
\operatorname{length}_{k(s)}(D'\cap Y)
=n(D',s^*)+\operatorname{length}_{k(s)}(D'\cap T).
\end{equation}
Since $(D',s^*)=1$, assertion 2.14(ii) follows from (2.14.2) and (2.14.3).

We now prove, by induction on the number of elements in the support of $\nu$,
the assertions of 2.13 and the following assertion:

\begin{equation}\tag{2.14.4}
\text{if $D$ is a regular divisor on $S$, then}\qquad
\chi\Bigl(D\cap\operatorname{Spec}(\underline{O}_S/I(\nu))\Bigr)
=\sum_{\substack{t\in S^*\\t\ \mathrm{on}\ D}}
\nu(t)[k(t):k].
\end{equation}
If $\nu=0$, then $\nu^*=0$, $I(\nu)=\underline{O}_S$, and 2.13 and (2.14.4)
are immediate (the second assertion of 2.13 follows by repeated application of
(2.11.4) with $n=0$).  Otherwise $\nu(s)>0$, and there is an $S$-morphism
$g:S'\longrightarrow S_1$.

For each closed point $a$ of $S_1$, let $S_{1a}$ be the localization of $S_1$
at $a$, let $\nu_a$ be the restriction of $\nu$ to $S_{1a}^*$, let
$S'_a=S'\times_{S_1}S_{1a}$, let
$I(\nu_a)=g_*\underline{O}_{S'_a}(-\nu_a^*)$, and put
\[
T_a=\operatorname{Spec}(\underline{O}_{S_1}/I(\nu_a)),\qquad
T=\coprod_aT_a,
\]
where the union is finite and disjoint.  Let $q$ be the ideal of $T$.  With the
notation of 2.14, one has

\begin{equation}\tag{2.14.5}
g_*\underline{O}_{S'}(-\nu^*)=qJ^{\nu(s)}.
\end{equation}
The induction hypothesis applies to $(S_{1a},S'_a,\nu_a)$.  Hence

\begin{equation}\tag{2.14.6}
R^ig_*\underline{O}_{S'}(-\nu^*)=0\quad\text{for }i>0,
\end{equation}
\begin{equation}\tag{2.14.7}
\chi(T)=
\sum_{t\in S^*\setminus\{s\}}
\frac{\nu(t)(\nu(t)+1)}2[k(t):k],
\end{equation}
and, by (2.14.4) applied to the divisor $s^*$ on $S_1$ and by (2.12.1),

\begin{equation*}
\operatorname{length}_{k(s)}(T\cap s^*)
=\sum_{\substack{t\ \mathrm{proximate}\\\mathrm{to}\ s}}
\nu(t)[k(t):k(s)]
\leq\nu(s).
\end{equation*}

%%%%%%%%%% source scan idx 33 | folio 26 %%%%%%%%%%

This inequality allows us to apply 2.14.  From (2.14.5) we obtain

\begin{equation*}
R^ip_*g_*\underline{O}_{S'}(-\nu^*)=0
\quad\text{for }i>0
\end{equation*}
and
\begin{equation*}
\operatorname{length}_{k(s)}
\bigl(\underline{O}_S/p_*g_*\underline{O}_{S'}(-\nu^*)\bigr)
=\frac{\nu(s)(\nu(s)+1)}2+\operatorname{length}_{k(s)}(T).
\end{equation*}
The assertions of 2.13 therefore follow from (2.14.6) and (2.14.7).  Finally,
2.14(ii) gives

\begin{equation*}
\operatorname{length}_{k(s)}
\Bigl(D\cap\operatorname{Spec}(\underline{O}_S/I(\nu))\Bigr)
=\nu(s)+\operatorname{length}_{k(s)}(D'\cap T).
\end{equation*}
Applying (2.14.4) to $D'$ gives

\begin{equation*}
\begin{aligned}
\chi\Bigl(D\cap\operatorname{Spec}(\underline{O}_S/I(\nu))\Bigr)
&=\nu(s)[k(s):k]
+\sum_{\substack{t\ \mathrm{on}\ D\\t\ne s}}
\nu(t)[k(t):k]\\
&=\sum_{\substack{t\in S^*\\t\ \mathrm{on}\ D}}
\nu(t)[k(t):k].
\end{aligned}
\end{equation*}
This completes the proof of 2.13.
\end{filecontents*}
% Original input: components/006_expose_X_2_15_through_2_24.tex
\begin{filecontents*}[overwrite,noheader]{sga7ii-cumulative-input-006.tex}
(2.15) Fragments of the preceding discussion generalize to higher dimension.  Let
$S$ be the spectrum of a regular local ring of dimension $n$, with closed point
$s$, where $k(s)$ is a finite extension of $k$.  We may again consider the set
$I$ of schemes $f:S'\longrightarrow S$ obtained from $S$ by successive blow-ups
of closed points, and put, as in 2.1,

\begin{equation*}
S^*=\varinjlim_{S'\in I}\operatorname{Irr}\bigl(f^{-1}(s)\bigr).
\end{equation*}
The set $S^*$ will again be called the set of \underline{points of $S$
infinitely near to $s$}.

If $X$ is a closed subscheme of $S$ and $t\in S^*$, the multiplicity of $X$ at
$t$ is defined as in 2.8.

Let $f:S'\longrightarrow S$ be a composite of blow-up morphisms at closed
points, let $D$ be a divisor on $S$, and let $D^*=f^*D$ be its total transform
on $S'$.

%%%%%%%%%% source scan idx 34 | folio 27 %%%%%%%%%%

Let $E$ be a curve on $S$, let $E'$ be its strict transform on $S'$, and suppose
that $D\cap E$ is supported at $s$.  The projection formula again gives

\begin{equation}\tag{2.16}
(D,E)=(D^*,E').
\end{equation}
Proceeding as in 2.9, one proves with the aid of (2.16) that

\begin{equation}\tag{2.17}
(D,E)=\sum_{t\in S^*}\mu_t(D)\mu_t(E)[k(t):k] \quad .
\end{equation}
Define proximate points as in 2.5.  Proceeding as in 2.11, one verifies that

\begin{equation}\tag{2.18}
\mu_t(E)=
\sum_{\substack{t'\ \mathrm{proximate}\\\mathrm{to}\ t}}
\mu_{t'}(E)[k(t'):k(t)] \quad .
\end{equation}
On the other hand, formula 2.11(ii) is false in general (see the examples in
2.25).

2.19. Let $S$ be a regular projective surface over a field $k$.  Write $S^*$
for the disjoint sum, taken over the closed points of $S$,

\begin{equation*}
S^*=\coprod_{s\in S}\operatorname{Spec}(\underline{O}_{s,S})^*,
\end{equation*}
and call $S^*$ the set of \underline{points infinitely near to the points of}
$S$.

Let $I$ be the set of surfaces obtained from $S$ by successive blow-ups of
closed points.  The set $I$ is cofinal in the set of regular models of the
function field of $S$, ordered by domination.  For $i\in I$, let $S_i$ denote
the corresponding surface, equipped with $p_i:S_i\longrightarrow S$.  We write
$\operatorname{Div}^{\pm}(S_i)$ for the group of divisors, effective or not, on
$S_i$.

\underline{Definition} 2.20. (i) \underline{The group of ``curves'' on $S$ is
the inductive limit}

\begin{equation*}
\varinjlim_{i\in I}\operatorname{Div}^{\pm}(S_i).
\end{equation*}

(ii) \underline{The category of ``invertible sheaves'' on $S$ is the inductive
limit}

%%%%%%%%%% source scan idx 35 | folio 28 %%%%%%%%%%

(SGA 4 VI 6) \underline{of the categories of invertible sheaves on the}
$(S_i)_{i\in I}$.

These objects are birational invariants of $S$.

\underline{Example} 2.21. (i) Every ``curve'' $C$ on $S$ defines an
``invertible sheaf'' $\underline{O}(C)$ on $S$.

(ii) Every divisor $C$ on $S$ defines a ``curve'' on $S$, denoted by $C^*$.

(iii) If $t\in S^*$, define the surface $S_+(t)$ and the divisor $t^*$ on
$S_+(t)$ as in 2.3; the ``curve'' defined by this divisor will again be denoted
by $t^*$.

(2.22) One verifies easily that a ``curve'' $C$ (respectively an ``invertible
sheaf'' $\mathcal L$) can be written in one and only one way in the form

\begin{equation*}
C=C_0^*-\sum_{t\in S^*}i_t t^*
\qquad\text{respectively}\qquad
\mathcal L=\mathcal L_0\otimes\underline{O}
\left(-\sum_{t\in S^*}i_t t^*\right),
\end{equation*}
where $C_0$ (respectively $\mathcal L_0$) is a divisor (respectively an
invertible sheaf) on $S$, and where the integers $i_t$ are almost all zero.  We
say that $C$ (respectively $\mathcal L$) is the curve $C_0$ (respectively the
invertible sheaf $\mathcal L_0$) endowed with the virtual multiplicities $i_t$.

The completion
$\operatorname{Div}^{\pm}(S)\times\mathbb Z^{S^*}$ of the group
$\operatorname{Div}^{\pm}(S)\times\mathbb Z^{(S^*)}$ of ``curves'' on $S$ is
again a birational invariant of $S$; it is called the group of ``generalized
curves'' on $S$.  On $S$, a ``generalized curve'' $C$ is represented by a
divisor $C_0$ endowed with a possibly infinite collection of virtual
multiplicities.

``Generalized invertible sheaves'' on $S$ are defined in the same way.

%%%%%%%%%% source scan idx 36 | folio 29 %%%%%%%%%%

\underline{Example} 2.23. Suppose that $k$ is perfect.  Then the invertible
sheaf $\Omega^2_{S/k}$, endowed with virtual multiplicity $-1$ at every point
$t\in S^*$, is a birational invariant of $S$; denote it by
$\Omega^2_{S/k}{}^*$.

\underline{Example} 2.24. Let $f:S_1\longrightarrow S_2$ be a dominant
morphism of regular projective surfaces over $k$.  The inverse image of a
``curve'' or an ``invertible sheaf'' on $S_2$ is defined without difficulty by
passing to a model of $S_2$ on which this ``curve'' or ``invertible sheaf'' is
an ordinary divisor or invertible sheaf.  By passage to the limit (I have not
checked whether this passage to the limit is possible), one similarly defines
the inverse image of a ``generalized curve'' or a ``generalized invertible
sheaf.''  Suppose that $k$ is perfect and that $f$ is generically \'{e}tale.
The \underline{different} $\mathfrak d(f)$ of $f$ is then the effective divisor
on $S_1$ characterized by

\begin{equation*}
f^*\Omega^2_{S_2/k}\bigl(\mathfrak d(f)\bigr)
\overset{\sim}{\longrightarrow}\Omega^2_{S_1/k} \quad .
\end{equation*}
This different is not a birational invariant of $f$; to obtain a birational
invariant one must consider the generalized coincidence curve

\begin{equation*}
\mathfrak d^*(f)=\mathfrak d(f)+\sum_{t\in S_1^*}t^*
-f^*\left(\sum_{t\in S_2^*}t^*\right),
\end{equation*}
or, briefly,

\begin{equation*}
\mathfrak d^*(f)=\Omega^2_{S_1/k}{}^*
-f^*\Omega^2_{S_2/k}{}^* \quad .
\end{equation*}
This ``generalized curve'' is supported on the closed subset of $S_1$ where
$f$ is not \'{e}tale.
\end{filecontents*}
% Original input: components/007_expose_X_example_2_25.tex
\begin{filecontents*}[overwrite,noheader]{sga7ii-cumulative-input-007.tex}
(2.25) \underline{Example}.

Let $k$ be an algebraically closed field and let $E_n$ be affine $n$-space
over $k$.  Let $C$ be a reduced curve in $E_n$ having a singularity at the
origin.  With this singularity we associate its \underline{Enriques diagram}
(cf. [3]): its vertices are the points of $C$ infinitely near the origin
(2.4 and 2.15), and each vertex is joined by a line to the vertices proximate
to it (2.5 and 2.15).  This diagram determines the multiplicities of these
points on $C$ (2.11 and 2.18) and, when $n=2$, the invariant $\delta$ (2.10
and 2.11) of the singularity of $C$ at the origin.

For the simplest plane singularities one obtains the following table.  A
number next to a vertex is its multiplicity.

\begin{center}
\small
\setlength{\tabcolsep}{5pt}
\renewcommand{\arraystretch}{1.35}
\begin{tabular}{c c c c}
\textit{Sketch} & \textit{Equation} & $\delta$ & \textit{Enriques diagram} \\

\begin{tikzpicture}[baseline=-.5ex,x=.55cm,y=.55cm]
  \draw[line width=.45pt] (-.8,0)--(.8,0);
\end{tikzpicture}
& $y=0$ & $0$ &
\begin{tikzpicture}[baseline=-.5ex,x=.56cm,y=.56cm]
  \foreach \x in {0,...,6}{
    \fill (\x,0) circle (1.05pt);
    \node[above=1pt,scale=.72] at (\x,0) {$1$};
  }
  \draw (0,0)--(6,0);
  \node[right=2pt,scale=.85] at (6,0) {$\cdots$};
\end{tikzpicture}
\\

\begin{tikzpicture}[baseline=-.5ex,x=.55cm,y=.55cm]
  \draw[line width=.45pt] (-.75,-.42)--(.75,.42);
  \draw[line width=.45pt] (-.75,.42)--(.75,-.42);
\end{tikzpicture}
& $xy=0$ & $1$ &
\begin{tikzpicture}[baseline=-.5ex,x=.72cm,y=.62cm]
  \fill (0,0) circle (1.05pt);
  \node[left=1pt,scale=.72] at (0,0) {$2$};
  \foreach \x in {1,...,5}{
    \fill (\x,.48) circle (1.05pt);
    \fill (\x,-.48) circle (1.05pt);
    \node[above=1pt,scale=.72] at (\x,.48) {$1$};
    \node[above=1pt,scale=.72] at (\x,-.48) {$1$};
  }
  \draw (0,0)--(1,.48)--(5,.48);
  \draw (0,0)--(1,-.48)--(5,-.48);
  \node[right=2pt,scale=.85] at (5,.48) {$\cdots$};
  \node[right=2pt,scale=.85] at (5,-.48) {$\cdots$};
\end{tikzpicture}
\\

\begin{tikzpicture}[baseline=-.5ex,x=.55cm,y=.55cm]
  \draw[line width=.45pt] (-.85,0)--(.85,0);
  \draw[line width=.45pt,domain=-.72:.72,samples=40]
    plot (\x,{.72*\x*\x});
\end{tikzpicture}
& $y(y-x^2)=0$ & $2$ &
\begin{tikzpicture}[baseline=-.5ex,x=.68cm,y=.62cm]
  \fill (0,0) circle (1.05pt);
  \fill (.9,0) circle (1.05pt);
  \node[above=1pt,scale=.72] at (0,0) {$2$};
  \node[above=1pt,scale=.72] at (.9,0) {$2$};
  \foreach \x/\lab in {1.8/1,2.7/1,3.6/1,4.5/0}{
    \fill (\x,.48) circle (1.05pt);
    \fill (\x,-.48) circle (1.05pt);
    \ifnum\lab=1
      \node[above=1pt,scale=.72] at (\x,.48) {$1$};
      \node[below=1pt,scale=.72] at (\x,-.48) {$1$};
    \fi
  }
  \draw (0,0)--(.9,0)--(1.8,.48)--(4.5,.48);
  \draw (.9,0)--(1.8,-.48)--(4.5,-.48);
  \node[right=2pt,scale=.72] at (4.5,-.48) {$\cdot\ \cdot$};
\end{tikzpicture}
\\

\begin{tikzpicture}[baseline=-.5ex,x=.55cm,y=.55cm]
  \draw[line width=.45pt]
    (0,0) .. controls (.28,.05) and (.55,.35) .. (.72,.78);
  \draw[line width=.45pt]
    (0,0) .. controls (.28,-.05) and (.55,-.35) .. (.72,-.78);
\end{tikzpicture}
& $y^2=x^3$ & $1$ &
\begin{tikzpicture}[baseline=-.5ex,x=.72cm,y=.65cm]
  \coordinate (A) at (0,.42);
  \coordinate (B) at (.9,.42);
  \coordinate (C) at (1.7,-.25);
  \coordinate (D) at (2.6,-.25);
  \coordinate (E) at (3.5,-.25);
  \foreach \p in {A,B,C,D,E}{\fill (\p) circle (1.05pt);}
  \node[above=1pt,scale=.72] at (A) {$2$};
  \foreach \p in {B,C,D,E}{\node[above=1pt,scale=.72] at (\p) {$1$};}
  \draw (A)--(B)--(C)--(D)--(E);
  \draw (A)--(C);
  \node[right=2pt,scale=.85] at (E) {$\cdots$};
\end{tikzpicture}
\end{tabular}
\end{center}

If $n>2$, the invariant $\delta$ is no longer determined by this diagram.
For example, the union of three concurrent coplanar lines (respectively,
noncoplanar lines) has invariant $\delta=3$ (respectively, $\delta=2$),
although the Enriques diagram is the same.  I do not know whether this
phenomenon also occurs for unibranch curves.

Here are two examples of diagrams in the case $n=3$ which cannot occur for
plane curves (cf. 2.5).  We give, in order, the equations of the curve, the
equations of the pure transform after blowing up the origin, parametric
equations of the curve, the invariant $\delta$, and the Enriques diagram.

%%%%%%%%%% source scan idx 38 | folio 31 %%%%%%%%%%

\begin{equation}\tag{2.25.1}
\left\{
\begin{gathered}
  \left\{\begin{aligned}y^2&=xz^4,\\ z^3&=x^2,
  \end{aligned}\right.\qquad
  \left\{\begin{aligned}(y/z)^2&=(x/z)^3,\\ z&=(x/z)^2,
  \end{aligned}\right.\\[3pt]
  \left\{\begin{aligned}x&=u^6,\\y&=u^7,\\z&=u^4,
  \end{aligned}\right.\qquad \delta=5
\end{gathered}
\right.
\end{equation}

\begin{center}
\begin{tikzpicture}[x=.78cm,y=.72cm]
  \coordinate (v1) at (0,0);
  \coordinate (v2) at (1,0);
  \coordinate (v3) at (1.75,-.85);
  \coordinate (v4) at (1.75,-1.70);
  \coordinate (v5) at (2.65,-2.35);
  \coordinate (v6) at (3.65,-2.62);
  \foreach \p in {v1,v2,v3,v4,v5,v6}{\fill (\p) circle (1.15pt);}
  \node[above left=1pt] at (v1) {$4$};
  \node[above=1pt] at (v2) {$2$};
  \node[right=2pt] at (v3) {$1$};
  \node[right=2pt] at (v4) {$1$};
  \node[above right=1pt] at (v5) {$1$};
  \node[above right=1pt] at (v6) {$1$};
  \draw (v1)--(v2) (v1)--(v3) (v1)--(v4);
  \draw (v2)--(v3) (v2)--(v4) (v3)--(v4);
  \draw (v4)--(v5)--(v6);
  \fill (4.08,-2.86) circle (.7pt);
  \fill (4.38,-3.05) circle (.7pt);
  \fill (4.68,-3.24) circle (.7pt);
\end{tikzpicture}
\end{center}

\begin{equation}\tag{2.25.2}
\left\{
\begin{gathered}
  \left\{\begin{aligned}y^2z&=x^3,\\ z^3&=y^2,
  \end{aligned}\right.\qquad
  \left\{\begin{aligned}(y/z)^2&=(x/z)^3,\\ z&=(y/z)^2,
  \end{aligned}\right.\\[3pt]
  \left\{\begin{aligned}x&=u^8,\\y&=u^9,\\z&=u^6,
  \end{aligned}\right.\qquad \delta=10
\end{gathered}
\right.
\end{equation}

\begin{center}
\begin{tikzpicture}[x=.72cm,y=.65cm]
  \coordinate (w1) at (0,0);
  \coordinate (w2) at (1,0);
  \coordinate (w3) at (1.75,-.72);
  \coordinate (w4) at (1.75,-1.44);
  \coordinate (w5) at (1.75,-2.16);
  \coordinate (w6) at (1.75,-2.88);
  \coordinate (w7) at (2.55,-3.70);
  \coordinate (w8) at (3.35,-4.52);
  \foreach \p in {w1,w2,w3,w4,w5,w6,w7,w8}{\fill (\p) circle (1.15pt);}
  \node[above left=1pt] at (w1) {$6$};
  \node[above=1pt] at (w2) {$2$};
  \foreach \p in {w3,w4,w5,w6}{\node[right=2pt] at (\p) {$1$};}
  \node[above right=1pt] at (w7) {$1$};
  \node[above right=1pt] at (w8) {$1$};
  \draw (w1)--(w2);
  \foreach \p in {w3,w4,w5,w6}{\draw (w1)--(\p);}
  \draw (w2)--(w3) (w2)--(w4);
  \draw (w3)--(w4)--(w5)--(w6)--(w7)--(w8);
  \fill (3.72,-4.86) circle (.7pt);
  \fill (4.02,-5.10) circle (.7pt);
  \fill (4.32,-5.34) circle (.7pt);
\end{tikzpicture}
\end{center}
\end{filecontents*}
% Original input: components/008_expose_X_2_26_through_2_30.tex
\begin{filecontents*}[overwrite,noheader]{sga7ii-cumulative-input-008.tex}
2.26. Theorem 2.27 below was stated as a false conjecture in the first version
of this exposé (circulated by the IHES).  It owes much to conversations with
D. S. Rim and includes some of his results (D. S. Rim, ``Torsion
Differentials and Deformations,'' Theorem 2.7(i)).

%%%%%%%%%% source scan idx 39 | folio 32 %%%%%%%%%%

Let $C$ be a reduced projective curve over an algebraically closed field $k$,
and let $s\in C(k)$.  We shall write

\begin{description}[leftmargin=2.7cm,style=nextline]
\item[$A_s$] for the local ring of $C$ at $s$ (or its completion, which would
make no difference below);
\item[$\widetilde A_s$] for the normalization of $A_s$;
\item[$\delta(s)$] for $\dim_k(\widetilde A_s/A_s)$;
\item[$m_0(s)$] for the dimension of the formal versal scheme $T$ of
deformations of $\operatorname{Spec}(A_s)$.
\end{description}

For an irreducible component $E$ of $T$, put likewise
$m_0^E(s)=\dim(E)$.

The module $\widetilde\tau$ of derivations of $\widetilde A_s$ and the module
$\tau$ of derivations of $A_s$ are torsion-free.  If $A_s$ is integral, they
embed in the module of derivations of the fraction field of $A_s$; in general,
they embed in the module of derivations of the total quotient ring of $A_s$
(a product of fields).  Moreover, through these embeddings, $\tau$ and
$\widetilde\tau$ are commensurable: both
$\tau/(\tau\cap\widetilde\tau)$ and
$\widetilde\tau/(\tau\cap\widetilde\tau)$ are finite-dimensional over $k$.
Set

\begin{equation*}
m_1(s)=[\widetilde\tau:\tau]
\mathrel{\mathop=\limits_{\mathrm{def}}}
\dim_k\frac{\widetilde\tau}{\tau\cap\widetilde\tau}
-\dim_k\frac{\tau}{\tau\cap\widetilde\tau}.
\end{equation*}

\underline{Theorem} 2.27. \underline{Let $E$ be an irreducible component of
the formal versal scheme $T$ of deformations of $A_s$. If the deformation of
$A_s$ over the generic point $e$ of $E$ is smooth---that is, if the image of
the formal non-smoothness subscheme (defined by a Jacobian ideal: VI 5), which
is finite over $T$, does not contain $e$---then}

\begin{equation*}
3\delta(s)=m_0^E(s)+m_1(s).
\end{equation*}

%%%%%%%%%% source scan idx 40 | folio 33 %%%%%%%%%%

The validity of 2.27 depends only on the completion of $A_s$.  Replacing
$(C,s)$ by a pair $(C',s)$ formally isomorphic to it at $s$, we may suppose
that the following conditions hold:

\begin{enumerate}[label=(\alph*),leftmargin=1.2cm]
\item $s$ is the only singular point of $C$;
\item $C$ has no infinitesimal automorphisms.
\end{enumerate}

Condition (b) ensures that the deformation functor of $C$ is pro-representable.
It is even effectively pro-representable, since there is no obstruction to
lifting a prescribed ample invertible sheaf on $C$.  Let
$f:X_0\longrightarrow T_0$ be the universal deformation of $C$, and let $t_0$
be the closed point of $T_0$, so that $C\overset\sim\longrightarrow X_{0t_0}$.

\underline{Lemma} 2.28. \underline{There exist a scheme $T_1$ of finite type
over $k$, a point $t_1\in T_1(k)$, and a proper flat morphism
$f:X_1\longrightarrow T_1$, whose fiber at $t_1$ is $C$, such that for every
$u\in T_1(k)$ the completion of $T_1$ at $u$ is a formal moduli scheme for the
fiber $X_{1u}$.}

Here is a direct proof of this corollary of general results of Artin.  Let
$\mathcal L$ be a very ample invertible sheaf on $C$ such that
$H^1(C,\mathcal L)=0$, and put
$\mathbb P=\mathbb P(H^0(C,\mathcal L))$.  We have an embedding
$v:C\hookrightarrow\mathbb P$.  Let $T_2$ be the Hilbert scheme parametrizing
the subvarieties of $\mathbb P$ having the same Hilbert polynomial as
$v(C)\subset\mathbb P$, and let $t_2\in T_2$ correspond to $v(C)$.  The
completion $\widehat T_2$ of $T_2$ at $t_2$ is smooth over $T_0$.  The vertical
tangent bundle of $\widehat T_2/T_0$ is induced by the subbundle $L$ of the
tangent bundle of $T_2$ which is the image of
$\operatorname{Lie}(\underline{\operatorname{Aut}}(\mathbb P))$ under the
action of $\underline{\operatorname{Aut}}(\mathbb P)$ on $T_2$.

In a neighborhood of $t_2$, let $f_1,\ldots,f_n$, with
$n=\dim\underline{\operatorname{Aut}}(\mathbb P)$, be equations defining a
subscheme $T'_1$ of $T_2$ which passes through $t_2$ and whose completion at
$t_2$ is étale over $T_0$.  Take for $T_1$ a neighborhood of $t_2$ in $T'_1$
on which the tangent bundle is transverse to $L$.

%%%%%%%%%% source scan idx 41 | folio 34 %%%%%%%%%%

2.29. The morphism of functors

\begin{equation*}
(\text{deformations of }C)\longrightarrow(\text{deformations of }A_s)
\end{equation*}

corresponds to a morphism $p:T_0\longrightarrow T$.  By VI 2.10, $p$ is
smooth.  Let $t$ be the closed point of $T$.  The Zariski tangent space of
$p^{-1}(t)$ at $t_0$ classifies the locally trivial deformations of $C$ over
$k[\varepsilon]$, where $\varepsilon^2=0$.  Putting $E_0=p^{-1}(E)$, we obtain

\begin{equation}\tag{2.29.1}
\dim(E_0)=\dim(E)+\dim H^1(C,T^1_{C/k}).
\end{equation}

Since $X_0/T_0$ is flat, all its fibers have the same arithmetic genus $p$,
where $1-p=\chi(\underline O)$.  By hypothesis, the fiber at the generic point
of $E_0$ is smooth.  The formal moduli scheme of a smooth proper curve being
smooth of dimension $3p-3$, Lemma 2.28 gives

\begin{equation}\tag{2.29.2}
\dim(E_0)=3p-3.
\end{equation}

Let $\widetilde C$ be the normalization of $C$, of arithmetic genus
$\widetilde p=p-\delta(s)$.  We have

\begin{equation}\tag{2.29.3}
\begin{aligned}
\dim H^1(C,T^1_{C/k})
&=-\chi(C,T^1_{C/k})\\
&=-\chi(\widetilde C,T^1_{\widetilde C/k})+m_1(s)\\
&=[3(p-\delta(s))-3]+m_1(s).
\end{aligned}
\end{equation}

Formulas (2.29.1), (2.29.2), and (2.29.3) prove 2.27.

\underline{Conjecture} 2.30. \underline{Every reduced proper curve has a
smooth deformation.}

This conjecture implies, by 2.27 and 2.28, that the formal versal scheme of
deformations of $A_s$ is always equidimensional, of dimension
$3\delta(s)-m_1(s)$.
\end{filecontents*}
% Original input: components/009_expose_X_2_31_through_bibliography.tex
\begin{filecontents*}[overwrite,noheader]{sga7ii-cumulative-input-009.tex}
%%%%%%%%%% source scan idx 42 | folio 35 %%%%%%%%%%

2.31. Let $C$ be a reduced proper curve over $k$, of arithmetic genus $p$.
We shall write

\begin{description}[leftmargin=2.7cm,style=nextline]
\item[$M_0$] for the dimension of the formal versal scheme $T$ of deformations
of $C$;
\item[$M_1$] for the dimension of the Lie algebra of
$\underline{\operatorname{Aut}}(C)$.
\end{description}

For an irreducible component $E$ of $T$, write $M_0^E$ for the dimension of
$E$.  The global analogue of 2.27 is the following.

\underline{Proposition} 2.32. \underline{If the deformation of $C$ over the
generic point of $E$ is smooth, then}

\begin{equation*}
3p-3=M_0^E-M_1.
\end{equation*}

Let $X\subset C(k)$ be a finite set of smooth points of $C$ such that every
vector field on $C$ which vanishes at every point $x\in X$ is zero.  The
deformation functor of $(C,X)$ is pro-representable and satisfies a variant of
2.28.  The formal moduli scheme of $(C,X)$ is smooth over $T$, of relative
dimension $\#X-M_1$, and the analogue of 2.28 gives

\begin{equation*}
M_0^E+(\#X-M_1)=3p-3+\#X,
\end{equation*}

which proves 2.32.

Suppose now that $k$ has \underline{characteristic zero}.  Then
$M_1=\dim\underline{\operatorname{Aut}}(C)$.  Here is an analogous
interpretation of $m_1(s)$.

%%%%%%%%%% source scan idx 43 | folio 36 %%%%%%%%%%

\underline{Lemma} 2.33. \underline{If $k$ has characteristic zero, every
derivation of $A_s$ extends to a derivation of $\widetilde A_s$.}

If $D$ is a derivation of $A_s$, then

\begin{equation*}
\exp(Dt)=\sum_{n\geq0}\frac{D^nt^n}{n!}
\end{equation*}

is an automorphism of the completion at $(s,0)$ of
$A_s\otimes_k k[t]$.  By transport of structure, this automorphism acts on the
normalization $(\widetilde A_s\otimes_k k[t])^{\wedge}$ of that ring.  Reducing
modulo $t^2$, we obtain an automorphism $1+Dt$ of
$\widetilde A_s^{\wedge}\otimes_k k[t]/(t^2)$, and $D$ thereby extends to a
derivation of $\widetilde A_s^{\wedge}$.  Since it also extends to the fraction
field of $A_s$ (to the total quotient ring when $A_s$ is not integral), it
extends to $\widetilde A_s$.

Suppose that $s$ is singular.  A similar argument then shows that the vector
field on $\operatorname{Spec}(\widetilde A_s)$ defined by $D$ vanishes at every
point above $s$.

Let $c$ be an ideal of $\widetilde A_s$ contained in the conductor of
$\widetilde A_s/A_s$.  Put

\begin{equation*}
\widetilde G_c=\underline{\operatorname{Aut}}(\widetilde A_s/c),
\end{equation*}

and let $G_c$ be the subgroup of $\widetilde G_c$ which stabilizes the
subalgebra $A_s/c$.  If $d\subset c$, then $\widetilde G_d$ stabilizes $c/d$
and maps onto $\widetilde G_c$.  Notice that
$g\in\widetilde G_c$ is determined by $g(t)$ for a uniformizer $t$, and that
$g(t)$ may be any uniformizer.  Consequently the pointed homogeneous space
$\widetilde G_c/G_c$ is independent of $c$; denote it by $\widetilde G/G$.

Let $\widetilde\tau_0$ be the $\widetilde A_s$-module of vector fields on
$\operatorname{Spec}(\widetilde A_s)$ which vanish at the points above $s$.
When $s$ is singular, one has

\begin{equation*}
\dim_k(\widetilde\tau/\widetilde\tau_0)
=\text{the number of branches},
\end{equation*}

and

\begin{equation*}
\widetilde\tau_0/\tau\overset\sim\longrightarrow
\text{the tangent space at the origin of }\widetilde G/G.
\end{equation*}

%%%%%%%%%% source scan idx 44 | folio 37 %%%%%%%%%%

\underline{Proposition} 2.34. \underline{In characteristic zero, when $s$ is
singular, the integer $m_1(s)$ is the sum of the number of branches and the
dimension of $\widetilde G/G$.}

\bigskip
\begin{center}
\underline{\textsc{Bibliography}}
\end{center}

\begin{description}[leftmargin=1.1cm,labelwidth=.8cm,labelsep=.3cm]
\item[{[1]}] P. Defrise, \textit{Étude locale des correspondances rationnelles
entre surfaces algébriques}, thèse d'agrégation, Brussels, 1949.

\item[{[2]}] P. Deligne and D. Mumford, ``The irreducibility of the space of
curves of a given genus,'' \textit{Publ. Math. IHES} \underline{36} (1969),
75--110.

\item[{[3]}] Enriques--Chisini, \textit{Lezioni sulla teoria geometrica delle
equazioni e delle funzioni algebriche}, Zanichelli, Bologna, 1924.

\item[{[4]}] R. Hartshorne, \textit{Residues and Duality}, Lecture Notes in
Mathematics \underline{20}, Springer, 1966.

\item[{[5]}] S. Lichtenbaum, ``Curves over discrete valuation rings,''
\textit{Amer. J. Math.} 90, no. 2 (April 1968), 380--405.

\item[{[6]}] I. Manin, ``Rational surfaces over perfect fields,''
\textit{Publ. Math. IHES}, no. 30, 55--114.

\item[{[7]}] D. Mumford, ``The topology of normal singularities of an
algebraic surface and a criterion for simplicity,'' \textit{Publ. Math. IHES},
no. 9, 5--22.

\item[{[8]}] A. Néron, \textit{Modèles minimaux des variétés abéliennes sur
les corps locaux et globaux}, Publ. Math., no. 21.

%%%%%%%%%% source scan idx 45 | folio 38 %%%%%%%%%%

\item[{[9]}] M. Raynaud, ``Spécialisation du foncteur de Picard. Critères
numériques de représentabilité,'' \textit{C. R. Acad. Sci. Paris} 264
(June 1967), 1001--1004.

\item[{[10]}] Saint-Donat, \textit{Sur un théorème de G. Castelnuovo et
F. Enriques}, thèse de troisième cycle, Lyon.

\item[{[11]}] O. Zariski, ``On polynomial ideals defined by infinitely near
base points,'' \textit{Amer. J. Math.} 60 (1938), 151.

\item[{[12]}] D. G. Northcott, ``On the genus formula for plane curves,''
\textit{J. London Math. Soc.} 30 (1955), 376--382.

\item[{[13]}] D. G. Northcott, ``Abstract dilatations and infinitely near
points,'' \textit{Proc. Cambridge Philos. Soc.} 52 (1956), 178--197.

\item[] The notion of an ``infinitely near point'' is discussed in Lipman's
appendix to Chapter II of the following work.

\item[{[14]}] O. Zariski, \textit{Algebraic Surfaces}, Ergebnisse
\underline{61}, Springer-Verlag.
\end{description}
\end{filecontents*}
% Original input: components/010_expose_XI_opening_through_1_1_5.tex
\begin{filecontents*}[overwrite,noheader]{sga7ii-cumulative-input-010.tex}
\clearpage
\begin{center}
{\large\textbf{SGA 7}}\\[1.5em]
{\Large\underline{EXPOSÉ XI}}\\[2.5em]
{\Large\underline{COHOMOLOGY OF COMPLETE INTERSECTIONS}}\\[1.5em]
{\large\underline{by P. DELIGNE}}
\end{center}

\bigskip
\noindent\underline{1. Cohomology of projective space and complete
intersections}

Let $X$ be a nonsingular complex projective variety, purely of dimension $n$,
and let $\underline O(1)$ be an ample invertible sheaf on $X$.  By Akizuki and
Nakano [1], Kodaira's vanishing theorem [4] implies that

\begin{equation}\tag{*}
H^i(X,\Omega_X^j(-1))=0\qquad\text{for }i+j<n.
\end{equation}

Again by [1], this result implies the ``weak Lefschetz theorem'' in complex
cohomology.  We propose to prove, by brute force, that if $X$ is a projective
space over a field of arbitrary characteristic, then (*) remains true.  This
will allow us, following [1], to compare the Hodge and de Rham cohomology of
complete intersections over a field of arbitrary characteristic with that of
projective space.

We shall then spell out some consequences of the results obtained.  In §2 we
give qualitative results deduced from Hirzebruch's numerical formula (2.2).  I
wish to thank N. Katz for his help in the preparation of that section.

1.0. Let $\underline E$ be a locally free module of rank $r+1$ on a scheme
$S$.  \underline{We assume that $r>0$.}  Recall that
$\mathbb P(\underline E)$ denotes the relative projective space over $S$ of
hyperplanes of $\underline E$.  Let $p:\mathbb P(\underline E)\to S$ be the
projection, and let

\begin{equation*}
\eta\in H^0\bigl(S,R^1p_*\Omega^1_{\mathbb P(\underline E)/S}\bigr)
\end{equation*}

be the first Chern class, in the Hodge sense, of the invertible sheaf
$\underline O(1)$.

%%%%%%%%%% source scan idx 47 | folio 40 %%%%%%%%%%

Recall that $\eta$ is the image of the class of $\underline O(1)$ in
$H^0(S,R^1p_*\underline O^*)$ under the logarithmic differential

\begin{equation*}
df/f:\underline O^*\longrightarrow
\Omega^1_{\mathbb P(\underline E)/S}.
\end{equation*}

For every sheaf $\underline F$ of $\underline O$-modules on
$\mathbb P(\underline E)$, put

\begin{equation*}
\underline F(n)=\underline F\otimes\underline O(1)^{\otimes n}
\qquad(n\in\mathbb Z).
\end{equation*}

\underline{Theorem} 1.1. (i) \underline{The sheaves
$R^ip_*(\Omega^j_{\mathbb P(\underline E)/S}(n))$ are locally free, and their
formation is compatible with arbitrary base change.}

(ii) \underline{For $0\leq i\leq r$, the section $\eta^i$ of
$R^ip_*(\Omega^i_{\mathbb P(\underline E)/S})$ defines an isomorphism from
$\underline O_S$ onto this sheaf. Moreover,}

\begin{equation*}
R^ip_*\Omega^j_{\mathbb P(\underline E)/S}=0
\qquad\underline{\text{if }}i\ne j\ \underline{\text{ or }}\ i>r.
\end{equation*}

(iii) \underline{If $n\ne0$, then, outside the following two cases,
$R^ip_*(\Omega^j_{\mathbb P(\underline E)/S}(n))$ vanishes:}

\begin{align}
i&=0\quad\underline{\text{and}}\quad n>j,
  \tag{iii$_a$}\\
i&=r\quad\underline{\text{and}}\quad n<j-r.
  \tag{iii$_b$}
\end{align}

Let

\begin{equation*}
\mathbb V(\underline E)=\operatorname{Spec}
\bigl(\operatorname{Sym}^*\underline E\bigr)
\end{equation*}

be the affine space defined by $\underline E$, let $e$ be its zero section,
and put
$\mathbb V(\underline E)^*=\mathbb V(\underline E)\setminus e(S)$.  Let $a$
and $a'$ be the projections of $\mathbb V(\underline E)$ and
$\mathbb V(\underline E)^*$ onto $S$, and let $q$ be the projection of
$\mathbb V(\underline E)^*$ onto $\mathbb P(\underline E)$:

\begin{center}
\begin{tikzcd}[row sep=large,column sep=large]
\mathbb V(\underline E)
  \arrow[dr,"a"]
& \mathbb V(\underline E)^*
  \arrow[l,hook']
  \arrow[r,"q"]
  \arrow[d,"a'"']
& \mathbb P(\underline E)
  \arrow[dl,"p"]
\\
{} & S\arrow[ul,bend right,"e"'] & {}
\end{tikzcd}
\qquad .
\end{center}

The punctured affine space $\mathbb V(\underline O(1))^*$ over
$\mathbb P(\underline E)$ is canonically isomorphic to
$\mathbb V(\underline E)^*$.  Consequently, for every quasi-coherent sheaf
$\underline F$ of $\underline O$-modules on $\mathbb P(\underline E)$, there
is a canonical isomorphism

%%%%%%%%%% source scan idx 48 | folio 41 %%%%%%%%%%

\begin{equation}\tag{1.1.1}
q_*q^*\underline F\simeq
\bigoplus_{n\in\mathbb Z}\underline F(n).
\end{equation}

This grading of $q_*q^*\underline F$ can also be obtained from the action on
this sheaf of the group $\mathbb G_m$ of homotheties of
$\mathbb V(\underline E)$.

Relative vector fields on $\mathbb V(\underline E)$ are identified with
$S$-morphisms from $\mathbb V(\underline E)$ to itself, the target being
identified with the tangent space at the origin.  The vector field $X$ defined
by the identity map of $\mathbb V(\underline E)$ is invariant under
homotheties.  It is tangent to the fibers of $q$ and defines a trivialization

\begin{equation}\tag{1.1.2}
X:\Omega^1_{\mathbb V(\underline E)^*/\mathbb P(\underline E)}
\longrightarrow\underline O_{\mathbb V(\underline E)^*}.
\end{equation}

If the $x_i$ form a basis of $\underline E$, then

\begin{equation*}
X=\sum_i x_i\,dx^i.
\end{equation*}

Since $\Omega^1_{\mathbb V(\underline E)^*/\mathbb P(\underline E)}$ has
rank one, the exact sequence

\begin{equation*}
0\longrightarrow q^*\Omega^1_{\mathbb P(\underline E)/S}
\longrightarrow\Omega^1_{\mathbb V(\underline E)^*/S}
\longrightarrow
\Omega^1_{\mathbb V(\underline E)^*/\mathbb P(\underline E)}
\longrightarrow0
\end{equation*}

induces exact sequences

\begin{equation*}
0\longrightarrow q^*\Omega^i_{\mathbb P(\underline E)/S}
\longrightarrow\Omega^i_{\mathbb V(\underline E)^*/S}
\longrightarrow
\Omega^1_{\mathbb V(\underline E)^*/\mathbb P(\underline E)}
\otimes q^*\Omega^{i-1}_{\mathbb P(\underline E)/S}
\longrightarrow0.
\end{equation*}

Using (1.1.2), these become

\begin{equation}\tag{1.1.3}
0\longrightarrow q^*\Omega^i_{\mathbb P(\underline E)/S}
\longrightarrow\Omega^i_{\mathbb V(\underline E)^*/S}
\xrightarrow{\ \iota_X\ }
q^*\Omega^{i-1}_{\mathbb P(\underline E)/S}
\longrightarrow0.
\end{equation}

Canonically,

\begin{equation}\tag{1.1.4}
\Omega^i_{\mathbb V(\underline E)/S}
=a^*\bigwedge^i\underline E.
\end{equation}

The short exact sequences (1.1.3) are therefore obtained by splitting the
following long exact sequence on $\mathbb V(\underline E)^*$:

\begin{equation}\tag{1.1.5}
\cdots\longrightarrow {a'}^*\bigwedge^i\underline E
\xrightarrow{\ \iota_X\ }
{a'}^*\bigwedge^{i-1}\underline E
\xrightarrow{\ \iota_X\ }
{a'}^*\bigwedge^{i-2}\underline E\longrightarrow\cdots.
\end{equation}

\end{filecontents*}
% Original input: components/011_expose_XI_1_1_6_through_1_2.tex
\begin{filecontents*}[overwrite,noheader]{sga7ii-cumulative-input-011.tex}
This long exact sequence is simply the restriction to
$\mathbb V(\underline E)^*$ of the Koszul resolution of
$\underline O_{e(S)}$ on $\mathbb V(\underline E)$.

Push the long exact sequence (1.1.5) forward to $\mathbb P(\underline E)$ by
the affine morphism $q$.  By (1.1.1) and the formula
$q_*{a'}^*=q_*q^*p^*$, one obtains a long exact sequence

\begin{equation}\tag{1.1.6}
\cdots\longrightarrow
\bigoplus_{n\in\mathbb Z}p^*\bigwedge^i\underline E(n)
\xrightarrow{\ \iota_X\ }
\bigoplus_{n\in\mathbb Z}p^*\bigwedge^{i-1}\underline E(n)
\longrightarrow\cdots .
\end{equation}

It splits into short exact sequences, the images of (1.1.3),

\begin{equation}\tag{1.1.7}
0\longrightarrow
\bigoplus_{n\in\mathbb Z}\Omega^i_{\mathbb P(\underline E)/S}(n)
\longrightarrow
\bigoplus_{n\in\mathbb Z}p^*\bigwedge^i\underline E(n)
\xrightarrow{\ \iota_X\ }
\bigoplus_{n\in\mathbb Z}\Omega^{i-1}_{\mathbb P(\underline E)/S}(n)
\longrightarrow0.
\end{equation}

The exact sequences (1.1.7) are direct sums of the exact sequences

\begin{equation}\tag{1.1.8}
0\longrightarrow\Omega^i_{\mathbb P(\underline E)/S}(n)
\longrightarrow p^*\bigwedge^i\underline E(n-i)
\xrightarrow{\ \iota_X\ }
\Omega^{i-1}_{\mathbb P(\underline E)/S}(n)
\longrightarrow0.
\end{equation}

Recall (EGA III, 2.1.15, or FAC) that the sheaves
$R^ip_*\underline O(k)$ are locally free and their formation is compatible
with arbitrary base change, and that

\begin{equation}\tag{1.1.9}
R^ip_*\underline O(k)=0
\quad\text{if }i\ne0,r,
\quad\text{or if }i=0, k<0,
\quad\text{or if }i=r, k>-r-1,
\end{equation}

and hence that $R^ip_*p^*\bigwedge^j\underline E(k)$ vanishes under the
same conditions.

The exact sequences (1.1.8) give a resolution

\begin{equation}\tag{1.1.10}
0\longrightarrow\Omega^i_{\mathbb P(\underline E)/S}(n)
\xrightarrow{\ E\ }p^*\bigwedge^i\underline E(n-i)
\xrightarrow{\ \iota_X\ }\cdots\longrightarrow
p^*\bigwedge^0\underline E(n)\longrightarrow0
\end{equation}

of $\Omega^i_{\mathbb P(\underline E)/S}(n)$.  The spectral sequence
defined by this resolution degenerates ($E_1=E_\infty$) for degree reasons
by (1.1.9); for $0<j<r$, it gives

\begin{equation}\tag{1.1.11}
R^jp_*\Omega^i_{\mathbb P(\underline E)/S}(n)
=\mathcal H\left(\cdots\xrightarrow{\ \iota_X\ }
p_*\bigwedge^{i-j}\underline E(n-i+j)
\xrightarrow{\ \iota_X\ }\cdots\right).
\end{equation}

Since $r\geq1$, a depth argument shows that

\begin{equation}\tag{1.1.12}
\bigoplus_{n\in\mathbb Z}p_*p^*\bigwedge^i\underline E(n)
=p_*q_*q^*p^*\bigwedge^i\underline E
={a'}_*{a'}^*\bigwedge^i\underline E
=a_*a^*\bigwedge^i\underline E.
\end{equation}

%%%%%%%%%% source scan idx 50 | folio 43 %%%%%%%%%%

Thus $p_*p^*\bigwedge^i\underline E(n)$ is the homogeneous part of weight
$n+i$ in $a_*a^*\bigwedge^i\underline E$, namely
$\operatorname{Sym}^n(\underline E)\otimes\bigwedge^i\underline E$.
Remembering that (1.1.5) comes from the Koszul resolution, (1.1.11) shows
that, for $0<j<r$, the sheaf
$R^jp_*\Omega^i_{\mathbb P(\underline E)/S}(n)$ can be nonzero only when
$i-j=0$ and $n-i+j=0$, that is, only when $n=0$ and $i=j$.  In that case,
moreover, $R^ip_*\Omega^i_{\mathbb P(\underline E)/S}$ is isomorphic to
$\underline O_S$.

For $j=0$, (1.1.10) also gives

\begin{equation}\tag{1.1.13}
p_*\Omega^i_{\mathbb P(\underline E)/S}(n)
=\ker\!\left(\iota_X:
p_*\bigwedge^i\underline E(n-i)\longrightarrow
p_*\bigwedge^{i-1}\underline E(n-i+1)\right),
\end{equation}

and the same argument shows that $\iota_X$ is injective when $n-i\leq0$.

By (1.1.8), with $n=0$ and $i=r+1$, one has

\begin{equation*}
\Omega^r_{\mathbb P(\underline E)/S}\simeq
p^*\bigwedge^{r+1}\underline E(-r-1).
\end{equation*}

When $S$ is the spectrum of a field, the vanishing assertions in 1.1(ii)
and (iii) now follow from the preceding argument for $j\ne r$, and follow
from it by duality ([2], V, 11.2, p.~211) for $j=r$.  In particular, for
any $j$ and $n$ there is at most one $i$ for which
$R^ip_*\Omega^j_{\mathbb P(\underline E)/S}(n)\ne0$.  This proves (i),
since $\Omega^j_{\mathbb P(\underline E)/S}$ is flat over $S$ (the hardest
part is locating the reference in EGA III; we suggest 6.10.5, 7.4.1, and
7.5.5, followed by a passage to the limit).

Knowing that $R^ip_*\Omega^i_{\mathbb P(\underline E)/S}$ is an invertible
sheaf for $0\leq i\leq r$, it remains only to prove that, when $S$ is the
spectrum of a field, $\eta^i\ne0$.  Indeed, $\eta^r$ is the cohomology
class of a rational point (the intersection of $r$ hyperplanes), hence a
basis of $H^r(\mathbb P^r,\Omega^r)$ (FGA, Exposé 149, no.~4).

1.2. The Euler--Poincaré characteristics of the $\Omega^i(n)$ may be
calculated directly from (1.1.10) or (1.1.13).  The rather unappetizing
result is

\begin{equation*}
\chi\bigl(\mathbb P^r_k,\Omega^j_{\mathbb P^r_k}(n)\bigr)
=\sum_{0\leq i\leq j}(-1)^{i-j}
\binom{r+1}{i}\binom{r+n-i}{r}.
\end{equation*}

\end{filecontents*}
% Original input: components/012_expose_XI_proposition_1_3.tex
\begin{filecontents*}[overwrite,noheader]{sga7ii-cumulative-input-012.tex}
%%%%%%%%%% source scan idx 51 | folio 44 %%%%%%%%%%

\underline{Proposition} 1.3. \underline{Let $f:X\to S$ be a proper smooth
morphism, let $\underline W$ be a locally free sheaf of rank $c$ on $X$,
let $\underline F$ be a coherent sheaf on $X$, let $d>0$ be an integer,
and let $s:\underline W\to\underline O_X$ be a section of
$\underline W^\vee$.  Assume that $S$ is noetherian and that:}

\begin{enumerate}[label=(\alph*),leftmargin=2.4em]
\item \underline{The closed subscheme $H$ of $X$ defined by the vanishing
of $s$ is smooth over $S$.}

\item \underline{Locally on $X$, if the $s_i$ are the coordinates of $s$
in a basis of $\underline W^\vee$, the $s_i$ form both an
$\underline O_X$-regular sequence and an $\underline F$-regular sequence.}

\item \underline{For every sequence of integers $k_\ell$, not all zero,
one has}
\end{enumerate}

\begin{equation*}
R^if_*\left(\bigotimes_\ell\bigwedge^{k_\ell}\underline W
\otimes\Omega^j_{X/S}\otimes\underline F\right)=0
\qquad\underline{\text{for }}i+j<d.
\end{equation*}

\underline{Then:}

\begin{align*}
R^if_*\left(\Omega^j_{X/S}\otimes\underline F\right)
&\xrightarrow{\ \sim\ }
R^if_*\left(\Omega^j_{H/S}\otimes\underline F\right)
&&\underline{\text{for }}i+j<d-c, \tag{i}\\
R^if_*\left(\Omega^j_{X/S}\otimes\underline F\right)
&\lhook\joinrel\longrightarrow
R^if_*\left(\Omega^j_{H/S}\otimes\underline F\right)
&&\underline{\text{for }}i+j=d-c. \tag{ii}
\end{align*}

The conormal sheaf $N_{H/X}$ of $H$ in $X$ is isomorphic to the restriction
of $\underline W$ to $H$, so the exact sequence

\begin{equation}\tag{1.3.1}
0\longrightarrow N_{H/X}\longrightarrow
\Omega^1_{X/S}\otimes\underline O_H\longrightarrow
\Omega^1_{H/S}\longrightarrow0
\end{equation}

defines a decreasing filtration on
$\Omega^j_{X/S}\otimes\underline O_H$ whose successive quotients are

\begin{equation}\tag{1.3.2}
\operatorname{Gr}^k\bigl(\Omega^j_{X/S}\otimes\underline O_H\bigr)
\simeq\bigwedge^k\underline W\otimes\Omega^{j-k}_{H/S}.
\end{equation}

It follows from (b) that one also has

\begin{equation}\tag{1.3.3}
\operatorname{Gr}^k\bigl(\Omega^j_{X/S}\otimes\underline O_H
\otimes\underline F\bigr)
=\bigwedge^k\underline W\otimes\Omega^{j-k}_{H/S}
\otimes\underline F.
\end{equation}

We prove (i) and (ii) by induction on $j$, simultaneously for all
$\underline F$ satisfying (b) and (c).  The assertions are empty for
$j<0$.

%%%%%%%%%% source scan idx 52 | folio 45 %%%%%%%%%%

Suppose, then, that (i) and (ii) have been proved for $j'<j$.  The
induction hypothesis applies to the sheaves
$\bigwedge^k\underline W\otimes\underline F$ and, by (c), gives

\begin{equation}\tag{1.3.4}
R^if_*\left(\bigwedge^k\underline W\otimes
\Omega^{j-k}_{H/S}\otimes\underline F\right)=0
\quad\text{for }k>0\text{ and }i+j<d-c+k.
\end{equation}

The filtration of
$\Omega^j_{X/S}\otimes\underline O_H\otimes\underline F$ considered
above defines a spectral sequence

\begin{equation*}
E_1^{p,i-p}=R^if_*\left(\bigwedge^p\underline W\otimes
\Omega^{j-p}_{H/S}\otimes\underline F\right)
\Longrightarrow
R^if_*\left(\Omega^j_{X/S}\otimes\underline O_H
\otimes\underline F\right).
\end{equation*}

By (1.3.4), $E_1^{p,i-p}=0$ when $p>0$ and $i+j\leq d-c$; hence

\begin{equation}\tag{1.3.5}
\left\{
\begin{aligned}
R^if_*\left(\Omega^j_{X/S}\otimes\underline O_H\otimes\underline F\right)
&\xrightarrow{\ \sim\ }
R^if_*\left(\Omega^j_{H/S}\otimes\underline F\right)
&&\text{if }i+j<d-c,\\
R^if_*\left(\Omega^j_{X/S}\otimes\underline O_H\otimes\underline F\right)
&\lhook\joinrel\longrightarrow
R^if_*\left(\Omega^j_{H/S}\otimes\underline F\right)
&&\text{if }i+j=d-c.
\end{aligned}
\right.
\end{equation}

By hypothesis (b), with a suitable differential, the sheaves
$\bigwedge^k\underline W\otimes\Omega^j_{X/S}\otimes\underline F$ form a
left Koszul resolution of
$\Omega^j_{X/S}\otimes\underline O_H\otimes\underline F$.  This
filtration defines a spectral sequence

\begin{equation}\tag{1.3.6}
E_1^{-p,q}=R^qf_*\left(\bigwedge^p\underline W\otimes
\Omega^j_{X/S}\otimes\underline F\right)
\Longrightarrow
R^{q-p}f_*\left(\Omega^j_{X/S}\otimes\underline O_H
\otimes\underline F\right).
\end{equation}

By hypothesis (c), $E_1^{-p,q}=0$ for $p>0$ and $q+j<d$, and trivially
for $p\notin[0,c]$.  Consequently,

\begin{equation}\tag{1.3.7}
\left\{
\begin{aligned}
R^if_*\left(\Omega^j_{X/S}\otimes\underline F\right)
&\xrightarrow{\ \sim\ }
R^if_*\left(\Omega^j_{H/S}\otimes\underline F\right)
&&\text{if }i+j<d-c,\\
R^if_*\left(\Omega^j_{X/S}\otimes\underline F\right)
&\lhook\joinrel\longrightarrow
R^if_*\left(\Omega^j_{H/S}\otimes\underline F\right)
&&\text{if }i+j=d-c.
\end{aligned}
\right.
\end{equation}

Assertions (i) and (ii) now follow from (1.3.5) and (1.3.7).

\end{filecontents*}
% Original input: components/013_expose_XI_1_4_through_1_6.tex
\begin{filecontents*}[overwrite,noheader]{sga7ii-cumulative-input-013.tex}
%%%%%%%%%% source scan idx 53 | folio 46 %%%%%%%%%%

1.4. Let $\underline E$ be a locally free sheaf of rank $r+1$ on a scheme
$S$, let $\mathbb P(\underline E)$ be the corresponding projective bundle,
let $d\leq r$ be an integer, and let
$\underline a=(a_1,\ldots,a_d)$ be a sequence of integers.  A relative
\underline{complete intersection} of multidegree $\underline a$ in
$\mathbb P(\underline E)$ is a closed subscheme which has codimension $d$
in every fiber and which, locally on $S$ for the étale topology, is the
intersection of $d$ hypersurfaces of respective degrees
$a_1,\ldots,a_d$.

By 1.1, hypotheses (b) and (c) of 1.3 hold for
$X/S=\mathbb P(\underline E)/S$, $\underline F=\underline O_X$,

\begin{equation*}
\underline W=\bigoplus_i\underline O(-a_i),
\end{equation*}

and a section $s$ of $\underline W^\vee$ defining a relative complete
intersection $Y\subset\mathbb P(\underline E)$ of multidegree
$\underline a$.

\underline{Theorem} 1.5. \underline{Under the hypotheses of 1.4, let
$n=r-d$ be the relative dimension of $Y$, and let
$\eta\in H^0(S,R^1f_*\Omega^1_{Y/S})$ be the first Chern class of
$\underline O(1)$ in the diagram}

\begin{center}
\begin{tikzcd}[column sep=large,row sep=large]
Y\arrow[rr,hook]\arrow[dr,dash,"f"']&&
\mathbb P(\underline E)\arrow[dl,"g"]\\
&S&{}
\end{tikzcd}
\qquad .
\end{center}

\underline{If $Y$ is smooth over $S$, then:}

\begin{enumerate}[label=(\roman*),leftmargin=2.6em]
\item \underline{The sheaves $R^if_*\Omega^j_{Y/S}$ are locally free,
and their formation is compatible with arbitrary base change.}

\item \underline{The spectral sequence}
\begin{equation*}
E_1^{p,q}=R^qf_*\Omega^p_{Y/S}\Longrightarrow
R^{p+q}f_*(\Omega^*_{Y/S}),
\end{equation*}
\underline{relating Hodge and de Rham cohomology degenerates
($E_1=E_\infty$); its abutment is locally free, and its formation is
compatible with arbitrary base change.}
\end{enumerate}

\noindent(iii)$_a$ \underline{One has}

\begin{equation*}
R^if_*\Omega^j_{Y/S}=0
\qquad\underline{\text{if }}i\ne j\ \underline{\text{ and }}\ i+j\ne n.
\end{equation*}

%%%%%%%%%% source scan idx 54 | folio 47 %%%%%%%%%%

\noindent(iii)$_b$ \underline{$\eta^i$ is a basis of
$R^if_*\Omega^i_{Y/S}$ if $2i<n$.}

\noindent(iii)$_c$ \underline{If $n<2i\leq2n$, and if $\alpha_i$ is the
basis of $R^if_*\Omega^i_{Y/S}$ dual to the basis $\eta^{n-i}$ of
$R^{n-i}f_*\Omega^{n-i}_{Y/S}$, then}

\begin{equation*}
\eta^i=\left(\prod_{h=1}^d a_h\right)\alpha_i.
\end{equation*}

\noindent(iii)$_d$ \underline{If $n=2m$, then $\eta^m\ne0$ (provided
$S\ne\varnothing$), and the quotient
$R^mf_*\Omega^m_{Y/S}/\underline O_S\eta^m$ is locally free.}

It is enough to prove 1.5 in a ``universal'' case, which permits us to
assume that $S$ is smooth over $\operatorname{Spec}(\mathbb Z)$.  In this
case (ii) follows from (i), since it holds over $\mathbb C$, and therefore
over every field of characteristic zero, by Hodge theory.  Moreover, since
$S$ is reduced, (i) follows from (iii), assumed proved when $S$ is the
spectrum of a field.  Indeed, by (iii), for each $j$ the functions

\begin{equation*}
s\longmapsto\dim_{k(s)}H^i(Y_s,\Omega^j_{Y_s})
\end{equation*}

are constant on $S$, except perhaps when $i=j$; hence they are all constant,
because their alternating sum is constant (EGA III, 7.9.2).  The conclusion
then follows from EGA III, 6.10.5, 7.6.9, and 7.7.5.

By 1.4 and 1.3, one has

\begin{equation}\tag{1.5.1}
\left\{
\begin{aligned}
R^ig_*\Omega^j_{\mathbb P(\underline E)/S}
&\xrightarrow{\ \sim\ }R^if_*\Omega^j_{Y/S}
&&\text{for }i+j<n,\\
R^ig_*\Omega^j_{\mathbb P(\underline E)/S}
&\lhook\joinrel\longrightarrow R^if_*\Omega^j_{Y/S}
&&\text{for }i+j=n.
\end{aligned}
\right.
\end{equation}

When $S$ is the spectrum of a field, assertions (iii)$_a$, (iii)$_b$, and
(iii)$_d$ follow from (1.5.1) and Serre duality ([2], V, 11.2, p.~211).
The cohomology class

%%%%%%%%%% source scan idx 55 | folio 48 %%%%%%%%%%

$\eta^n$ is the class of a zero-cycle of degree
$\prod_{h=1}^d a_h$ (the degree of the complete intersection), so

\begin{equation*}
\eta^n=\left(\prod_{h=1}^d a_h\right)\alpha_n.
\end{equation*}

By definition, $\eta^{n-i}\alpha_i=\alpha_n$, and (iii)$_c$ follows.

Thus (i) and (ii) hold; in the universal case, the assertions in (iii)
follow from (i) and from (iii) in the case already proved, where $S$ is the
spectrum of a field.

\underline{Theorem} 1.6. \underline{Keep the hypotheses and notation of
1.4 and 1.5, assume that $Y$ is smooth, let $\ell$ be a prime invertible in
all residue fields of $S$, and let
$\eta_\ell\in H^0(S,R^2f_*\mathbb Z_\ell(1))$ denote the $\ell$-adic first
Chern class of $\underline O(1)$ (SGA 5, VII).  Then:}

\begin{enumerate}[label=(\roman*),leftmargin=2.6em]
\item \underline{$R^if_*\mathbb Z_\ell=0$ if $i$ is odd and $i\ne n$.}

\item \underline{If $2i<n$, then $R^{2i}f_*\mathbb Z_\ell(i)$ is
isomorphic to $\mathbb Z_\ell$ and has $\eta_\ell^i$ as a basis.}

\item \underline{If $n<2i\leq2n$, then
$R^{2i}f_*\mathbb Z_\ell(i)$ is isomorphic to $\mathbb Z_\ell$ and has}
\begin{equation*}
\eta_\ell^i\bigg/\prod_{h=1}^d a_h
\end{equation*}
\underline{as a basis.}

\item \underline{The $\mathbb Z_\ell$-sheaf $R^nf_*\mathbb Z_\ell$ is
constructible, twisted constant, and torsion-free.  If $n$ is even, then
$\eta_\ell^{n/2}\ne0$ (provided $S\ne\varnothing$), and
$R^nf_*\mathbb Z_\ell(n/2)/\mathbb Z_\ell\eta_\ell^{n/2}$ is torsion-free.}
\end{enumerate}

The $\mathbb Z_\ell$-sheaves $R^if_*\mathbb Z_\ell$ are constructible and
twisted constant, and their formation is compatible with arbitrary base
change, by SGA 4, XVI, 2.2.  It therefore suffices to prove 1.6 when $S$ is
the spectrum of an algebraically closed field.  Repeated application of
SGA 4, XIV, 3.3 then shows, for every $m$, that the canonical map

\begin{equation*}
H_Y^q\bigl(\mathbb P(\underline E),\mathbb Z/\ell^m\bigr)
\longrightarrow H^q(Y,\mathbb Z/\ell^m)
\end{equation*}

%%%%%%%%%% source scan idx 56 | folio 49 %%%%%%%%%%

is bijective for $q>r+d$ and surjective for $q=r+d$.  From this and
Poincaré duality (SGA 4, XVIII), one deduces (SGA 5, VII) that the
restriction map

\begin{equation*}
H^q\bigl(\mathbb P(\underline E),\mathbb Z_\ell\bigr)
\longrightarrow H^q(Y,\mathbb Z_\ell)
\end{equation*}

is bijective for $q<n$, and injective with torsion-free cokernel for $q=n$.
This proves (i) for $i<n$, as well as (ii) and (iv).  As in 1.5, one checks
that $\eta_\ell^n$ is $\prod_{h=1}^d a_h$ times the canonical class; (i)
and (iii) then follow by Poincaré duality (SGA 4, XVIII).

\end{filecontents*}
% Original input: components/014_expose_XI_1_7_through_1_9.tex
\begin{filecontents*}[overwrite,noheader]{sga7ii-cumulative-input-014.tex}
1.7. If $Y$ is a complete intersection of dimension at least three in a
projective space $\mathbb P^r(k)$ over a field $k$, then
$\operatorname{Pic}(Y)$ is known to be isomorphic to $\mathbb Z$ and to be
generated by $\underline O(1)$ (SGA 2, XII, 3.7).  In dimension two, when
$Y$ is smooth, one can still prove the following.

\underline{Theorem} 1.8. \underline{Let $Y$ be a smooth complete-intersection
surface in the projective space $\mathbb P^r(k)$ over an algebraically
closed field $k$.  Then $\underline{\operatorname{Pic}}^\tau(Y)=0$, and the
class of $\underline O(1)$ is not divisible in the Néron--Severi group
$\operatorname{NS}(Y)$.}

Since $H^1(Y,\underline O)=0$ by 1.5(iii)$_a$, one knows that
$\underline{\operatorname{Pic}}^0(Y)=0$.  The vanishing of
$\underline{\operatorname{Pic}}^\tau(Y)$ then follows from the familiar
lemma below, the fact that $H^2(Y,\mathbb Z_\ell(1))$ is torsion-free by
1.6(iv), and the fact that $H^0(Y,\Omega_Y^1)=0$ by 1.5(iii)$_a$.

\underline{Lemma} 1.9. \underline{Let $Y$ be a complete algebraic variety
over an algebraically closed field $k$.}

\begin{enumerate}[label=(\roman*),leftmargin=2.6em]
\item \underline{If $\ell$ is prime to the characteristic of $k$, then the
$\ell$-torsion of $\operatorname{NS}(Y)$ identifies with the $\ell$-torsion
of $H^2(Y,\mathbb Z_\ell(1))$.}

%%%%%%%%%% source scan idx 57 | folio 50 %%%%%%%%%%

\item \underline{If $k$ has nonzero characteristic $p$ and $Y$ is smooth,
then the kernel of multiplication by $p$ in $\operatorname{Pic}(Y)$
identifies with the subgroup of $H^0(Y,\Omega_Y^1)$ consisting of locally
logarithmic differentials.}
\end{enumerate}

The Kummer exact sequences

\begin{equation*}
0\longrightarrow\mu_{\ell^n}\longrightarrow\mathbb G_m
\xrightarrow{\ x\mapsto x^{\ell^n}\ }\mathbb G_m\longrightarrow0
\end{equation*}

give exact sequences

\begin{equation*}
0\longrightarrow
\operatorname{Pic}(Y)/\ell^n\operatorname{Pic}(Y)
\longrightarrow H^2(Y,\mu_{\ell^n})
\longrightarrow H^2(Y,\mathbb G_m)_{\ell^n}
\longrightarrow0,
\end{equation*}

where

\begin{equation*}
\operatorname{Pic}(Y)/\ell^n\operatorname{Pic}(Y)
=\operatorname{NS}(Y)\otimes\mathbb Z/\ell^n.
\end{equation*}

Passing to the inverse limit gives

\begin{equation*}
0\longrightarrow\operatorname{NS}(Y)\otimes\mathbb Z_\ell
\longrightarrow H^2(Y,\mathbb Z_\ell(1))
\longrightarrow
\operatorname{Hom}\bigl(\mathbb Q_\ell/\mathbb Z_\ell,
H^2(Y,\mathbb G_m)\bigr)
\longrightarrow0.
\end{equation*}

The cokernel is torsion-free because $\mathbb Q_\ell/\mathbb Z_\ell$ is
divisible.  This proves (i).

To prove (ii), it is enough to observe that the exact sequence

\begin{equation*}
0\longrightarrow\underline O_Y^*
\xrightarrow{\ x\mapsto x^p\ }\underline O_Y^*
\xrightarrow{\ df/f\ }\Omega_Y^1
\end{equation*}

defines a short exact sequence

\begin{equation*}
0\longrightarrow\underline O_Y^*\longrightarrow\underline O_Y^*
\longrightarrow\Omega^1_{Y,\mathrm{loc}\,\log}\longrightarrow0,
\end{equation*}

and then to take the corresponding long exact cohomology sequence.  This
argument is due to Cartier, who proves moreover that, if $C$ denotes the
Cartier operator on $H^0(Y,\Omega_Y^1)$, then the locally logarithmic
differentials form the kernel of $C-\operatorname{Id}$.

Suppose, for a contradiction, that the class of $\underline O(1)$ in
$\operatorname{NS}(Y)$ has the form $nu$, with $n>1$.  If $\ell$ divides
$n$ and is prime to the characteristic, then in $\ell$-adic cohomology one
again has

%%%%%%%%%% source scan idx 58 | folio 51 %%%%%%%%%%

\begin{equation*}
\eta_\ell=n\,c_1(u),
\end{equation*}

contradicting 1.6(iv).

Finally, if the characteristic $p$ of $k$ is nonzero and $p$ divides $n$,
then in Hodge cohomology one again has

\begin{equation*}
\eta=n\,c_1(u)=0,
\end{equation*}

contradicting 1.5(iii)$_d$ and completing the proof.

\end{filecontents*}
% Original input: components/015_expose_XI_2_1_through_2_4.tex
\begin{filecontents*}[overwrite,noheader]{sga7ii-cumulative-input-015.tex}
2. \underline{Numerical results}

2.1. Let $\underline a=(a_i)_{1\leq i\leq d}$ be a sequence of integers
$\geq1$, and let $V_n(\underline a)$ denote a smooth complete intersection
of dimension $n$ and multidegree $\underline a$ in a projective space
$\mathbb P^{n+d}$ over a field $k$.  Put

\begin{equation*}
h^{p,q}(\underline a)
=\dim_k H^q\bigl(V_{p+q}(\underline a),\Omega^p\bigr)
\end{equation*}

and

\begin{equation*}
h_0^{p,q}(\underline a)=h^{p,q}(\underline a)-\delta_{p,q}.
\end{equation*}

Thus the $h_0^{p,q}(\underline a)$ are the Hodge numbers of the
middle-dimensional primitive cohomology of a complete intersection of
multidegree $\underline a$.  They depend only on $\underline a$, $p$, and
$q$, by 1.5.  Set

\begin{equation}\tag{2.1.1}
H(\underline a)=\sum_{p,q\geq0}h_0^{p,q}(\underline a)y^pz^q
\in\mathbb Z[[y,z]].
\end{equation}

For $d=1$ and $\underline a=(a)$, put
$h^{p,q}(a)=h^{p,q}(\underline a)$ and $H(a)=H(\underline a)$.  For $d=0$,
put $\sup(a_i)=1$.

The following formula is the special case $k=0$ of formula (2), p.~160,
in Hirzebruch [3]:

%%%%%%%%%% source scan idx 59 | folio 52 %%%%%%%%%%

\begin{equation}\tag{2.2}
\sum_{n=0}^{\infty}\sum_p
\chi\bigl(V_n(\underline a),\Omega^p\bigr)y^pz^{n+d}
=\frac{1}{(1+yz)(1-z)}
\prod_{i=1}^d
\frac{(1+yz)^{a_i}-(1-z)^{a_i}}
{(1+yz)^{a_i}z+y(1-z)^{a_i}}.
\end{equation}

We shall transform this into the following formula.

\underline{Theorem} 2.3 (Hirzebruch). \underline{With the notation of 2.1,
one has}

\begin{equation*}
H(\underline a)=\frac{1}{(1+y)(1+z)}\left[
\prod_{i=1}^d
\frac{(1+y)^{a_i}-(1+z)^{a_i}}
{-(1+y)^{a_i}z+(1+z)^{a_i}y}-1\right].
\end{equation*}

By 1.5,

\begin{equation*}
\chi\bigl(V_n(\underline a),\Omega^p\bigr)
=(-1)^p+(-1)^{n-p}h_0^{p,n-p}(\underline a).
\end{equation*}

Consequently, the left-hand side $I$ of (2.2) can be rewritten as

\begin{align*}
I
&=\sum_{n=0}^{\infty}\sum_{p+q=n}
\left((-1)^p+(-1)^q h_0^{p,q}(\underline a)\right)y^pz^{p+q+d}\\
&=\sum_{p,q\geq0}(-1)^p(yz)^pz^qz^d
+\sum_{p,q\geq0}(-1)^q h_0^{p,q}(\underline a)(yz)^pz^qz^d\\
&=z^d\left[\frac{1}{(1+yz)(1-z)}
+H(\underline a)(yz,-z)\right].
\end{align*}

The right-hand side $II$ of (2.2) can be rewritten as

\begin{equation*}
II=\frac{z^d}{(1+yz)(1-z)}
\prod_{i=1}^d
\frac{(1+yz)^{a_i}-(1-z)^{a_i}}
{(1+yz)^{a_i}z+yz(1-z)^{a_i}}.
\end{equation*}

After cancelling $z^d$ in $I=II$, one obtains

\begin{equation*}
H(\underline a)(yz,-z)
=\frac{1}{(1+yz)(1-z)}\left[
\prod_{i=1}^d
\frac{(1+yz)^{a_i}-(1-z)^{a_i}}
{(1+yz)^{a_i}z+(1-z)^{a_i}yz}-1\right],
\end{equation*}

and 2.3 follows by a change of variables.

%%%%%%%%%% source scan idx 60 | folio 53 %%%%%%%%%%

\underline{Corollary} 2.4. \underline{With the notation of 2.1:}

\begin{enumerate}[label=(\roman*),leftmargin=2.6em]
\item \underline{For $d=1$, one has}
\begin{equation*}
\begin{aligned}
H(a)
&=-\frac{(1+y)^{a-1}-(1+z)^{a-1}}
{(1+y)^az-(1+z)^ay}\\[1ex]
&=\frac{\displaystyle\sum_{i,j\geq0}
\binom{a-1}{i+j+1}y^iz^j}
{\displaystyle1-\sum_{i,j\geq1}\binom{a}{i+j}y^iz^j}.
\end{aligned}
\end{equation*}

\item \underline{Writing $|P|$ for the number of elements of a subset
$P$ of $[1,d]$, one has}
\begin{equation*}
H(\underline a)=
\sum_{\substack{P\subset[1,d]\\P\ne\varnothing}}
(1+y)^{|P|-1}(1+z)^{|P|-1}
\prod_{i\in P}H(a_i).
\end{equation*}
\end{enumerate}

By 2.3,

\begin{align*}
H(a)
&=-\frac{1}{(1+y)(1+z)}\left[
\frac{(1+y)^a-(1+z)^a}{(1+y)^az-(1+z)^ay}+1\right]\\
&=-\frac{1}{(1+y)(1+z)}\left[
\frac{(1+y)^a(1+z)-(1+z)^a(1+y)}
{(1+y)^az-(1+z)^ay}\right]\\
&=-\frac{(1+y)^{a-1}-(1+z)^{a-1}}
{(1+y)^az-(1+z)^ay}.
\end{align*}

Let $N$ and $D$ be the numerator and denominator in this last expression,
so that $H(a)=-N/D$.

%%%%%%%%%% source scan idx 61 | folio 54 %%%%%%%%%%

One has

\begin{equation*}
N=\sum_k\binom{a-1}{k}(y^k-z^k)
=(y-z)\sum_{i,j\geq0}\binom{a-1}{i+j+1}y^iz^j
\end{equation*}

and

\begin{align*}
D
&=\sum_k\binom ak(y^kz-yz^k)
=(z-y)+yz\sum_{k\geq1}\binom ak(y^{k-1}-z^{k-1})\\
&=(y-z)\left(-1+yz\sum_{i,j\geq0}
\binom a{i+j+2}y^iz^j\right)\\
&=(y-z)\left(-1+\sum_{i,j\geq1}
\binom a{i+j}y^iz^j\right).
\end{align*}

This proves (i).

Now invert the formula in 2.3 when $d=1$.  One finds

\begin{equation*}
-\frac{(1+y)^a-(1+z)^a}{(1+y)^az-(1+z)^ay}
=1+(1+y)(1+z)H(a).
\end{equation*}

Substituting this identity into 2.3 gives

\begin{align*}
H(\underline a)
&=\frac{1}{(1+y)(1+z)}\left[
\prod_{i=1}^d\left(1+(1+y)(1+z)H(a_i)\right)-1\right]\\
&=\frac{1}{(1+y)(1+z)}\left[
\sum_{\substack{P\subset[1,d]\\P\ne\varnothing}}
(1+y)^{|P|}(1+z)^{|P|}\prod_{i\in P}H(a_i)\right],
\end{align*}

which proves (ii).
\end{filecontents*}
% Original input: components/016_expose_XI_theorem_2_5_through_2_6.tex
\begin{filecontents*}[overwrite,noheader]{sga7ii-cumulative-input-016.tex}
\underline{Theorem} 2.5. \underline{With the notation of 2.1:}

\begin{enumerate}[label=(\roman*),leftmargin=2.6em]
\item \underline{If $p+q=p'+q'$ and
$p\leq p'\leq q'\leq q$, then}
\begin{equation*}
h_0^{p,q}(\underline a)\leq h_0^{p',q'}(\underline a).
\end{equation*}

\item \underline{If $p\leq q$, then $h_0^{p,q}(\underline a)\ne0$ if
and only if}
\begin{equation*}
p>\left\lfloor
\frac{p+q+d-\sum_i a_i}{\sup_i(a_i)}
\right\rfloor.
\end{equation*}
\underline{Here $\lfloor\ \rfloor$ denotes the integer part of the
quotient.}
\end{enumerate}

%%%%%%%%%% source scan idx 62 | folio 55 %%%%%%%%%%

Let (*) be the following property of a formal power series
$P=\sum a_{p,q}y^pz^q$:

\begin{quote}
(*) the $a_{p,q}$ are nonnegative, $a_{p,q}=a_{q,p}$, and whenever
$p+q=p'+q'$ and $p\leq p'\leq q'\leq q$, one has
$a_{p,q}\leq a_{p',q'}$.
\end{quote}

Assertion (i) follows at once from the formulas in 2.4, the identity
$(1-x)^{-1}=\sum_{n\geq0}x^n$, and the following lemma, whose proof is left
to the reader.

\underline{Lemma} 2.5.1. \underline{The set of formal power series
satisfying (*) is stable under addition and multiplication.}

We first prove (ii) in the special case of hypersurfaces ($d=1$).  For
$p\leq q$, the desired formula can then be rewritten as

\begin{equation}\tag{2.5.2}
h_0^{p,q}(a)\ne0
\quad\Longleftrightarrow\quad
p\geq\left\lfloor\frac{p+q+1}{a}\right\rfloor
\end{equation}

or, equivalently, $(p+1)a\geq p+q+2$.  From the second formula in 2.4(i),
one sees that $h_0^{p,q}(a)$ is nonzero if and only if there exist pairs
$(e,f)$ and $(c_i,d_i)$ such that

\begin{equation}\tag{2.5.3}
\left\{
\begin{aligned}
(e,f)&\geq(0,0),& e+f&\leq a-2,\\
(c_i,d_i)&\geq(1,1),& c_i+d_i&\leq a,\\
(p,q)&=(e,f)+\sum_i(c_i,d_i).&&
\end{aligned}
\right.
\end{equation}

%%%%%%%%%% source scan idx 63 | folio 56 %%%%%%%%%%

For fixed $p+q$, the smallest possible $p$ is obtained by taking $e=0$
and every $c_i=1$; it is the number of indices $i$.  Thus it is the lower
bound of those $p$ for which

\begin{equation*}
p\geq\frac{(p+q)-(a-2)}{a},
\qquad\text{i.e.}\qquad (p+1)a\geq p+q+2.
\end{equation*}

We now pass to the general case.  If $\underline a'$ is obtained from
$\underline a$ by deleting the $a_i$ equal to one, then
$H(\underline a)=H(\underline a')$.  In view of the form of 2.5(ii), we may
therefore assume that every $a_i>1$.  If $d=0$, then
$h_0^{p,q}=0$ and the assertion is correct; we lose no generality by
assuming $d>0$ and $a_d=\sup_i(a_i)$.

By 2.4(ii) and (2.5.2), the condition
$h_0^{p,q}(\underline a)\ne0$ is equivalent to the existence of a nonempty
subset $P$ of $[1,d]$, a pair $(e,f)$, and pairs
$(c_i,d_i)_{i\in P}$ such that

\begin{equation}\tag{2.5.4}
\left\{
\begin{aligned}
(|P|-1,|P|-1)&\geq(e,f)\geq(0,0),\\
c_i,d_i&\geq
\left\lfloor\frac{c_i+d_i+1}{a_i}\right\rfloor,
& (c_i,d_i)&\geq(0,0),\\
(p,q)&=(e,f)+\sum_{i\in P}(c_i,d_i).&&
\end{aligned}
\right.
\end{equation}

For fixed $n=p+q$, seek the smallest possible value of $p$.  Put
$n_i=c_i+d_i$.  For such a $p$, conditions (2.5.4) become

%%%%%%%%%% source scan idx 64 | folio 57 %%%%%%%%%%

\begin{equation}\tag{2.5.5}
\left\{
\begin{aligned}
a_i=2&\Longrightarrow n_i\text{ is even}
&&\text{(otherwise (2.5.4) is impossible)},\\
c_i&\geq\left\lfloor\frac{n_i+1}{a_i}\right\rfloor,
&c_i&\leq n_i,\\
(|P|-1,|P|-1)&\geq(e,f)\geq(0,0),&&\\
n&=e+f+\sum_{i\in P}n_i,&&\\
p&=e+\sum_{i\in P}c_i.&&
\end{aligned}
\right.
\end{equation}

This description first shows that the smallest $p$ is obtained for
$P=[1,d]$, to which case we henceforth restrict.  If

\begin{equation}\tag{2.5.6}
\sum_i(a_i-2)+(d-1)\geq n,
\end{equation}

then a solution exists with $p=c_i=e=0$.  Otherwise, if $a_d\ne2$, the
best choice is $n_i=a_i-2$ for $i\ne d$, so that $c_i=0$ for $i\ne d$,
$e=0$, $f=d-1$, and

\begin{equation*}
n_d=n-(d-1)-\sum_{i\ne d}(a_i-2).
\end{equation*}

Then

\begin{equation}\tag{2.5.7}
p=\left\lfloor
\frac{n-(d-1)-\sum_{i\ne d}(a_i-2)+1}{a_d}
\right\rfloor,
\end{equation}

and hence, for $p\leq q$, one has
$h_0^{p,q}(\underline a)\ne0$ if and only if

\begin{equation}\tag{2.5.8}
p\geq\left\lfloor
\frac{n-\sum_{i\ne d}(a_i-1)+1}{a_d}
\right\rfloor.
\end{equation}

Formula (2.5.8) reduces to (2.5.2) for $d=1$ and is trivially true when
(2.5.6) holds.  If $d\geq2$, (2.5.6) fails, and all the $a_i$ equal two,
the best choice is instead

%%%%%%%%%% source scan idx 65 | folio 58 %%%%%%%%%%

$n_i=0$ for $i\ne d$, $e=0$, and $f=d-1$ or $d-2$, chosen congruent to
$n$ modulo two, with $n_d=n-f$.  If $n\equiv d-1\pmod2$, this proves
(2.5.7) as above.  Otherwise, observe that

\begin{equation*}
\left\lfloor\frac{n-(d-2)+1}{2}\right\rfloor
=\left\lfloor\frac{n-(d-1)+1}{2}\right\rfloor,
\end{equation*}

so (2.5.7), and hence (2.5.8), remains valid when $a_d=2$.

It remains to verify that (2.5.8) is equivalent to 2.5(ii).  Indeed,

\begin{equation*}
\left\lfloor\frac{n-\sum_{i\ne d}(a_i-1)+1}{a_d}\right\rfloor
=\left\lfloor\frac{n+d-\sum_{i\ne d}a_i}{a_d}\right\rfloor
=\left\lfloor\frac{p+q+d-\sum_i a_i}{a_d}\right\rfloor+1.
\end{equation*}

2.6. If $X$ is a complete intersection of dimension $n$ in
$\mathbb P^r$, it follows at once from (2.5.8) that the primitive part of
$H^n(X)$ vanishes only when $X$ is an odd-dimensional quadric.

\end{filecontents*}
% Original input: components/017_expose_XI_2_7_through_bibliography.tex
\begin{filecontents*}[overwrite,noheader]{sga7ii-cumulative-input-017.tex}
\underline{Definition} 2.7. \underline{If $X$ is a proper smooth algebraic
variety over $\mathbb C$, the Hodge level of $H^n(X)$ is the supremum of
the numbers $q-p$ for which $p+q=n$ and $h^{p,q}(X)\ne0$.}

The Hodge level $\ell$ is therefore congruent to $n$ modulo two, and is at
most
\[
\min\bigl(n,2\dim(X)-n\bigr).
\]
It is nonnegative unless $H^n(X)=0$, in which case it is $-\infty$.
\footnote{The printed source reads
$\min(n,\dim(X)-n)$.  The factor $2$ is required by the Hodge ranges
$0\leq p,q\leq\dim(X)$; the English text corrects this evident source
typo.}

\underline{Corollary} 2.8. \underline{Retain the notation of 2.1, put
$k=\mathbb C$, and suppose that $a_d=\sup_i(a_i)$.  Let $\ell$ be an
integer with $0\leq\ell<n$, congruent to $n$ modulo two.  The Hodge level
of $H^n(V_n(\underline a))$ is at most $\ell$ if and only if}

%%%%%%%%%% source scan idx 66 | folio 59 %%%%%%%%%%

\begin{equation*}
\sum_{i\ne d}(a_i-1)
\leq (\ell+1)-(a_d-2)\frac{n-\ell}{2}.
\end{equation*}

By (2.5.8), if the Hodge level is nonnegative, it is the maximum of zero
and

\begin{equation*}
n-2\left\lfloor
\frac{n-\sum_{i\ne d}(a_i-1)+1}{a_d}
\right\rfloor.
\end{equation*}

The inequality

\begin{equation*}
n-2\left\lfloor
\frac{n-\sum_{i\ne d}(a_i-1)+1}{a_d}
\right\rfloor\leq\ell
\end{equation*}

may be rewritten as

\begin{equation*}
\frac{n-\ell}{2}
\leq
\left\lfloor
\frac{n-\sum_{i\ne d}(a_i-1)+1}{a_d}
\right\rfloor,
\end{equation*}

or equivalently as

\begin{equation*}
\sum_{i\ne d}(a_i-1)
\leq 1+n-a_d\frac{n-\ell}{2}
=\ell+1-(a_d-2)\frac{n-\ell}{2}.
\end{equation*}

2.9. Using formula 2.8, it is easy, for small $\ell$, to enumerate the
multidegrees of complete intersections in $\mathbb P^r$ of dimension
$n>1$ whose $H^n$ has level $\ell$.  The table given in the first version
of this exposé was, however, incorrect; the table below is extracted from
the more complete one given by M. Rapoport in his article
``Complément à l'article de P. Deligne \emph{La conjecture de Weil pour
les surfaces K3}'', \emph{Inventiones Math.} \underline{15}, 227--236
(1972).  For $\ell\leq1$, the complete intersections
$V_n(\underline a)$ of dimension $n$ and multidegree
$\underline a=(a_1,\ldots,a_d)$, with $n>\ell$, $d\geq1$, and $a_i\geq2$,
having Hodge level $\ell$, are the following:

%%%%%%%%%% source scan idx 67 | folio 60 %%%%%%%%%%

\begin{equation*}
\begin{array}{ll}
\ell=-\infty & V_{2k+1}(2) \qquad (k\geq0),\\[2ex]
\ell=0 & V_{2k}(2)\ (k\geq1),\quad
          V_{2k}(2,2)\ (k\geq1),\quad V_2(3),\\[2ex]
\ell=1 & V_{2k+1}(2,2)\ (k\geq1),\quad
          V_{2k+1}(2,2,2)\ (k\geq1),\\
& V_3(3),\quad V_5(3),\quad V_3(3,2),\quad V_3(4).
\end{array}
\end{equation*}

2.10. For small $\underline a$, the generating functions of the
$h_0^{p,q}(\underline a)$ are, by 2.4,

\begin{equation*}
\begin{array}{rcl}
H(2)&=&\dfrac{1}{1-yz},\\[3ex]
H(3)&=&\dfrac{2+y+z}{1-3yz-y^2z-yz^2},\\[3ex]
H(2,2)&=&\dfrac{y+z}{(1-yz)^2}
          +\dfrac{2}{(1-yz)^2}+\dfrac{1}{1-yz},
\end{array}
\end{equation*}

so that

\begin{equation*}
\begin{array}{rcll}
h_0^{p,p}(2)&=&1&,\\[2ex]
h_0^{1,1}(3)&=&6&,\\[2ex]
h_0^{p,p}(2,2)&=&2p+3&.
\end{array}
\end{equation*}

%%%%%%%%%% source scan idx 68 | folio 61 %%%%%%%%%%

\bigskip
\begin{center}
\underline{BIBLIOGRAPHY}
\end{center}

\begin{enumerate}[label={[\arabic*]},leftmargin=3em]
\item Y. Akizuki and S. Nakano, ``Note on Kodaira--Spencer's proof of
Lefschetz theorems'', \emph{Proc. Jap. Acad.} 30, 266--272 (1954).

\item R. Hartshorne, \emph{Residues and Duality}, Lecture Notes in
Mathematics 20, Springer, 1966.

\item F. Hirzebruch, \emph{Topological Methods in Algebraic Geometry},
3rd ed., Springer-Verlag, 1966.

\item K. Kodaira, ``On a differential geometric method in the theory of
analytic stacks'', \emph{Proc. Nat. Acad. Sci. USA} 39, 1268--1273
(1953).
\end{enumerate}
\end{filecontents*}
% Original input: components/018_expose_XII_opening_through_1_3.tex
\begin{filecontents*}[overwrite,noheader]{sga7ii-cumulative-input-018.tex}
\clearpage
\begin{center}
S G A 7

\bigskip
\underline{E X P O S É XII}

\bigskip
\underline{QUADRICS}

\bigskip
\underline{by P. DELIGNE}
\end{center}

\bigskip
0.0. The present exposé collects several well-known results about
quadrics which we shall need in order to prove the Picard--Lefschetz
formula.

0.1. Let $m$ be a section of an $\mathcal O_S$-module $M$ on a scheme
$S$.  We shall say that $m$ is \underline{nowhere zero} on $S$ if, for
every $s\in S$, the image of $m$ in the $k(s)$-vector space $M_s$ obtained
from $M$ by the base change $\{s\}\longrightarrow S$ is nonzero.  We use
the same terminology when $S=\Spec(A)$ and $m$ is an element of an
$A$-module $M$.

\medskip
1. \underline{Quadratic forms}

1.1. Let $S=\Spec(A)$ be an affine scheme, let $V$ be a locally free
$A$-module of rank $n$, and let $Q$ be a quadratic form on $V$
(Bourbaki, \emph{Algèbre}, Ch.~9, \S3, no.~4).  Recall that $Q$ is called
\underline{nondegenerate} if the associated bilinear form

\begin{equation*}
\Phi(x,y)=Q(x+y)-Q(x)-Q(y)
\end{equation*}

is nondegenerate, that is, if it induces an isomorphism from $V$ to its
dual $V^*$.  If $2$ is not invertible in $A$, this can occur only when
$n$ is even (loc. cit.).

%%%%%%%%%% source scan idx 70 | folio 63 %%%%%%%%%%

The quadratic forms on $V$ form a projective $A$-module; a quadratic form
$Q$ is ``nowhere zero'' on $S$ (0.1) if and only if the elements $Q(v)$
generate the unit ideal of $A$.

The notation $\mathbb P(V)$ is used in the sense of EGA; recall that if
$A$ is a field, then $\mathbb P(V)$ is the projective space whose set of
rational points is the quotient $(V^*\setminus\{0\})/A^*$.

The quadratic form $Q$ identifies with a section of the invertible sheaf
$\mathcal O(2)$ on the projective space $\mathbb P(V^*)$ over $S$.  When
$Q$ is nowhere zero, the (``homogeneous'') equation $Q=0$ defines a
subscheme of $\mathbb P(V^*)$, flat and purely of relative dimension
$n-2$ over $S$.  This scheme is called the \underline{quadric defined by
$Q$}.  The form $Q$ is called \underline{ordinary} if it is nowhere zero
on $S$ and the quadric that it defines is smooth over $S$.

The form $Q$ is ordinary if and only if every form obtained from it by
extending the scalars from $A$ to a field is ordinary.  If $A$ is a
field, then:

\begin{enumerate}[label=\alph*),leftmargin=2.5em]
\item if $n$ is even or $\operatorname{char}(A)\ne2$, then $Q$ is ordinary
if and only if $Q$ is nondegenerate;

\item if $n$ is odd and $\operatorname{char}(A)=2$, then $Q$ is ordinary
if and only if the kernel $N$ of the associated (alternating) bilinear
form $\Phi$ has dimension one and $Q$ does not vanish on $N$.
\end{enumerate}

\underline{Proposition} 1.2. \underline{If $Q$ is ordinary and
$n=2m$ (respectively, $n=2m+1$), then, locally for the étale topology on
$S$, $V$ has a basis $e_1,\ldots,e_n$ such that}

%%%%%%%%%% source scan idx 71 | folio 64 %%%%%%%%%%

\begin{equation*}
Q\!\left(\sum_{i=1}^{n}x_i e_i\right)
=\sum_{i=1}^{m}x_i x_{i+m}
\end{equation*}

\noindent\underline{(respectively,}

\begin{equation*}
Q\!\left(\sum_{i=1}^{n}x_i e_i\right)
=\sum_{i=1}^{m}x_i x_{i+m}+\lambda x_{2m+1}^{2},
\qquad (\lambda\text{ invertible}).
\end{equation*}

\footnote{In both displayed sums the printed source has upper limit
$m-1$.  That leaves a two-dimensional radical and contradicts the
statement; the standard ordinary normal form, and the induction that
follows, require the upper limit $m$.}

We prove 1.2 by induction on $m$.  If $m=0$, the assertion is clear.  If
$m>0$, the quadric with equation $Q=0$ is smooth over $S$ and has
nonempty geometric fibers; it therefore has sections locally for the
étale topology.  Locally, such a section is defined by an element
$e\in V$, nowhere zero on $\Spec(A)$, with $Q(e)=0$.  At no point of $S$
does $e$ belong to the kernel of $\Phi$; hence, locally, one can find
$f'$ such that $\Phi(e,f')=1$.  If

\begin{equation*}
f=-Q(f')e+f',
\end{equation*}

then

\begin{equation*}
Q(e)=Q(f)=0,\qquad \Phi(e,f)=1.
\end{equation*}

\footnote{The final term in the printed definition of $f$ is $f$; it is
the evident source typo $f'$, as is confirmed by the displayed
conclusion.}

If $V_2$ is the orthogonal complement of the submodule $V_1$ generated
by $e$ and $f$, then $V=V_1\oplus V_2$, and the conclusion follows by
applying the induction hypothesis to $V_2$.

1.3. For the definition and principal properties of the Clifford algebra
of $Q$, we refer to Bourbaki, \emph{Algèbre}, Ch.~9, \S9.  We denote by
$\rho$ the canonical map from $V$ to $C(Q)$.

\underline{Lemma} 1.3.1. \underline{Under the hypotheses of 1.1, for
every $\lambda\in A$ there exists a unique homomorphism}

\begin{equation*}
[\lambda]\colon C^+(\lambda Q)\longrightarrow C^+(Q)
\end{equation*}

\underline{such that, for $x,y\in V$,}

\begin{equation*}
[\lambda]\bigl(\rho(x)\rho(y)\bigr)
=\lambda\rho(x)\rho(y).
\end{equation*}

%%%%%%%%%% source scan idx 72 | folio 65 %%%%%%%%%%

This homomorphism is induced on passing to the quotient by the
endomorphism $[\lambda]$ of the even part $T^+(V)$ of the tensor algebra
of $V$, which acts as $\lambda^r$ on $T^{2r}(V)$.  Indeed,

\begin{equation*}
[\lambda]\bigl(x\otimes x-\lambda Q(x)\bigr)
=\lambda\bigl(x\otimes x-Q(x)\bigr).
\end{equation*}

Suppose that $Q$ is nowhere zero on $S$.  Lemma 1.3.1 then defines a
transitive system of isomorphisms between the algebras $C^+(\lambda Q)$
for $\lambda\in A^*$.

Consequently, the algebra $C^+(Q)$ depends (up to unique isomorphism)
only on $V$ and on $Q$ up to multiplication by a factor; in other words,
$C^+(Q)$ depends only on $V$ and on the subscheme of $\mathbb P(V^*)$
with equation $Q=0$.  Every homothety with invertible ratio is an
automorphism of the pair $(V,X\subset\mathbb P(V^*))$; one checks that
these automorphisms act trivially on $C^+(V,X)$.  It follows that the
algebra $C^+(V,X)$ depends only on the pair of schemes
$X\subset\mathbb P(V^*)$ over $A$.  We shall denote it by
$C^+(X,\mathbb P(V^*))$.
\end{filecontents*}
% Original input: components/019_expose_XII_1_4_through_1_9.tex
\begin{filecontents*}[overwrite,noheader]{sga7ii-cumulative-input-019.tex}
1.4. Let $V$ be a locally free $A$-module of rank $2m>0$, and let $Q$
be a nondegenerate quadratic form on $V$.  Suppose given a decomposition
$V=W_1\oplus W_2$ into two totally singular subspaces.  The form
$Q(w_1+w_2)$ then places $W_1$ and $W_2$ in perfect duality.  Recall that,
for $f\in W_2$ (identified with $W_1^*$), contraction by $f$ is the
endomorphism of degree $-1$ of $\bigwedge W_1$ defined by

\begin{equation*}
f\mathbin{\lrcorner}(e_1\wedge\cdots\wedge e_p)
=\sum_{i=1}^{p}(-1)^{i-1}\langle e_i,f\rangle
e_1\wedge\cdots\wedge\widehat e_i\wedge\cdots\wedge e_p.
\end{equation*}

\footnote{The printed formula has $(-1)^i$ and begins the remaining wedge
with an undefined $e_0$.  The English restores the standard contraction
formula $(-1)^{i-1}e_1\wedge\cdots\widehat e_i\cdots\wedge e_p$, required
by the Clifford action displayed immediately below.  The reading was
checked in a targeted high-detail enlargement of folio 65.}

Endow the $A$-module $\bigwedge W_1$ with the $\mathbb Z/2$-grading for
which

%%%%%%%%%% source scan idx 73 | folio 66 %%%%%%%%%%

\begin{equation*}
(\bigwedge W_1)^+=\bigoplus_i\bigwedge^{2i}W_1,
\qquad
(\bigwedge W_1)^-=\bigoplus_i\bigwedge^{2i+1}W_1.
\end{equation*}

The algebra $\operatorname{End}(\bigwedge W_1)$ is then
$\mathbb Z/2$-graded: its even (respectively, odd) endomorphisms are those
that preserve the grading (respectively, have degree one modulo two).

We refer to Bourbaki, \emph{Algèbre}, Ch.~9, \S9, no.~4, for the proof of
the existence of a unique isomorphism $s$ of $\mathbb Z/2$-graded algebras
from $C(Q)$ to $\operatorname{End}(\bigwedge W_1)$ such that

\begin{equation*}
s\bigl(\rho((e,0))\bigr)=e\wedge(-),
\qquad
s\bigl(\rho((0,f))\bigr)=f\mathbin{\lrcorner}(-).
\end{equation*}

Temporarily denote by $e(W_1,W_2)$ the element of $C^+(Q)$ for which
$s(e(W_1,W_2))$ is the identity on $(\bigwedge W_1)^+$ and zero on
$(\bigwedge W_1)^-$.  The center $Z(Q)$ of $C^+(Q)$ is isomorphic to
$A\oplus A$, this decomposition being defined by the idempotents
$e(W_1,W_2)$ and $1-e(W_2,W_1)$.  Moreover, $C^+(Q)$ is a matrix algebra
over $Z(Q)$.

From this and 1.2 one deduces, as in loc. cit.:

\underline{Proposition} 1.5. \underline{Under the hypotheses of 1.1, if
$n=2m>0$ and $Q$ is nondegenerate, then the center $Z(Q)$ of $C^+(Q)$ is
an étale algebra, locally free of rank two over $A$, and $C^+(Q)$ is an
Azumaya algebra over $Z(Q)$.}

1.6. Let $V$ and $Q$ be as in 1.4, and let $W$ be a totally singular
direct summand of $V$, locally free of rank $m$.  Let $W'_1$ be a
complement of $W$, so that $V=W\oplus W'_1$.  Every complement $W'_2$ of
$W$ identifies with the graph of a map $u:W'_1\to W$.  The form $\Phi$
places $W$ and $W'_1$ in duality, so $u$ is uniquely determined by the
bilinear form on $W'_1$

%%%%%%%%%% source scan idx 74 | folio 67 %%%%%%%%%%

\begin{equation*}
B(x,y)=\Phi(x,u(y)),
\end{equation*}

and $W'_2$ is totally singular if and only if

\begin{equation*}
B(x,x)=-Q(x)
\end{equation*}

on $W'_1$.  Since such forms $B$ always exist, $W$ has totally singular
complements.  Having chosen one of them, say $W'$, the set of totally
singular complements $W'_0$ of $W$ identifies with the set of alternating
forms on $W'$, or equivalently with the set of sections of the affine
space $\mathbb E(\bigwedge^2W')$ over $S$.

For each section $x$ of $\mathbb E(\bigwedge^2W')$, let $W'_0(x)$ be the
corresponding complement of $W$ and let
$e(x)=e(W,W'_0(x))$.  The function $x\mapsto e(x)$ is compatible with
every extension of scalars and hence defines a morphism of $S$-schemes

\begin{equation*}
e:\mathbb E(\bigwedge^2W')\longrightarrow\Spec(Z(Q)).
\end{equation*}

The $S$-scheme $\mathbb E(\bigwedge^2W')$ is flat with geometrically
connected fibers, whereas $\Spec(Z(Q))$ is étale over $S$.  The morphism
$e$ therefore factors through $S$, hence through a section of
$\Spec(Z(Q))$ over $S$; equivalently, $e(x)$ is independent of $x$.
\footnote{The printed sentence says that $e$ factors through a section of
the affine space $\mathbb E(\bigwedge^2W')$.  The connected-to-étale
argument instead gives the factorization through $S$ and a section of
$\Spec(Z(Q))$, exactly as the following conclusion requires.}
This justifies the following notation.

\underline{Notation} 1.7. \underline{For $W$ as above and any totally
singular complement $W'$ of $W$, denote by $e(W)$ the idempotent
$e(W,W')\in Z(Q)$.}

1.8. We saw in 1.3 that, for invertible $\lambda$, $C^+(Q)$ identifies
with $C^+(\lambda Q)$.  A continuity argument analogous to the preceding
one shows that this identification carries

\begin{equation*}
e(W)\in Z(Q)\qquad\text{to}\qquad e(W)\in Z(\lambda Q).
\end{equation*}

\footnote{The printed cross-reference is 1.6; the scalar-change
identification is the construction of 1.3.1.}

%%%%%%%%%% source scan idx 75 | folio 68 %%%%%%%%%%

1.9. Recall that if $B$ and $C$ are two $\mathbb Z/2$-graded algebras,
then multiplication in their graded tensor product $B\otimes' C$ is
defined by the Koszul rule

\begin{equation*}
(b_1\otimes c_1)(b_2\otimes c_2)
=(-1)^{\deg(c_1)\deg(b_2)}b_1b_2\otimes c_1c_2
\end{equation*}

when the $b_i$ and $c_i$ are homogeneous.

If $(V_1,Q_1)$ and $(V_2,Q_2)$ are two quadratic modules, then

\begin{equation*}
C(Q_1\oplus Q_2)=C(Q_1)\otimes' C(Q_2)
\end{equation*}

(Bourbaki, loc. cit.).

If $P$ and $Q$ are two finite-type, locally free, $\mathbb Z/2$-graded
$A$-modules, we identify the $\mathbb Z/2$-graded algebras
$\operatorname{End}(P)\otimes'\operatorname{End}(Q)$ and
$\operatorname{End}(P\otimes Q)$ by setting, for homogeneous
$u,v,p,q$,

\begin{equation*}
(u\otimes v)(p\otimes q)
=(-1)^{\deg(v)\deg(p)}u(p)\otimes v(q).
\end{equation*}
\end{filecontents*}
% Original input: components/020_expose_XII_1_10_through_2_2.tex
\begin{filecontents*}[overwrite,noheader]{sga7ii-cumulative-input-020.tex}
\underline{Lemma} 1.10. \underline{Let $(V,Q)$ be the quadratic module
which is the sum of two quadratic modules $(V_1,Q_1)$ and $(V_2,Q_2)$,
where $V_i$ is locally free of rank $2m_i$ and $Q_i$ is nondegenerate.
Let $V_i=W_i\oplus W'_i$ be a decomposition of $V_i$ into two totally
singular subspaces $(i=1,2)$.  Then the diagram}

\begin{center}
\begin{tikzcd}[column sep=huge,row sep=large]
C(Q) \arrow[rr,dash,"\sim"] \arrow[d,dash,"\wr"]
  & {} & C(Q_1)\otimes' C(Q_2) \arrow[d,equal] \\
\operatorname{End}(\bigwedge(W_1\oplus W_2))
  \arrow[r,phantom,"\simeq"]
  & \operatorname{End}(\bigwedge W_1\otimes\bigwedge W_2)
  \arrow[r,phantom,"\simeq"]
  & \operatorname{End}(\bigwedge W_1)\otimes'
    \operatorname{End}(\bigwedge W_2)
\end{tikzcd}
\end{center}

\underline{is commutative.}

%%%%%%%%%% source scan idx 76 | folio 69 %%%%%%%%%%

The verification is left to the reader.  One deduces at once from this
lemma that

\begin{equation}\tag{1.10.1}
e(W_1\oplus W_2)
=e(W_1)e(W_2)+(1-e(W_1))(1-e(W_2)).
\end{equation}

1.11. The following construction is borrowed from Micali and Villamayor.
If $p:Z\longrightarrow S$ is a double étale covering of $S$, then $Z$ may
be regarded as a $\mathbb Z/2$-torsor over $S$.  Denote the addition of
torsors by $*$; for a section $s$ of $Z$, let $e(s)$ be the idempotent of
$p_*\mathcal O_Z$ that equals one on $s$ and zero elsewhere.  With the
notation of 1.10, there is a unique isomorphism

\begin{equation}\tag{1.11.1}
\Spec(Z(Q_1))*\Spec(Z(Q_2))\longrightarrow\Spec(Z(Q)),
\end{equation}

whose underlying addition morphism is

\begin{equation*}
+:\Spec(Z(Q_1))\times_S\Spec(Z(Q_2))
\longrightarrow\Spec(Z(Q)),
\end{equation*}

and which, after every base change, satisfies

\begin{equation*}
e(s+t)=e(s)e(t)+(1-e(s))(1-e(t)).
\end{equation*}

If $s(W)$ is the section of $\Spec(Z(Q))$ such that $e(s(W))=e(W)$, then
formula (1.10.1) becomes

\begin{equation*}
s(W_1\oplus W_2)=s(W_1)+s(W_2).
\end{equation*}

\underline{Proposition} 1.12. \underline{Let $A$ be a field, let $V$ be
a vector space of dimension $2m>0$ over $A$, let $Q$ be a nondegenerate
quadratic form on $V$, and let $W_1,W_2$ be two totally isotropic
subspaces of $V$ of dimension $m$.  Then $e(W_1)=e(W_2)$ if and only if
$\dim(W_1/(W_1\cap W_2))$ is even.}

\footnote{The printed proposition says that $V$ has dimension $2m$ ``over
$V$''; the base field introduced in the same sentence is $A$.}

%%%%%%%%%% source scan idx 77 | folio 70 %%%%%%%%%%

(Cf. Bourbaki, \emph{Algèbre}, Ch.~9, \S6, Exercise~28.)  Let
$e_1,\ldots,e_m,f_1,\ldots,f_m$ be a basis of $V$ such that

\begin{enumerate}[label=\alph*),leftmargin=2.5em]
\item $Q(\sum_i x_i e_i+\sum_i y_i f_i)=\sum_i x_i y_i$;
\item $W_1$ is generated by the $e_i$;
\item $W_2$ is generated by $(e_i)_{1\leq i\leq s}$ and
$(f_i)_{s<i\leq m}$.
\end{enumerate}

Then $V$ is the sum of the $V_i=Ae_i+Af_i$, and

\begin{equation*}
W_\alpha=\bigoplus_i(W_\alpha\cap V_i),
\qquad \alpha=1,2.
\end{equation*}

The addition formula 1.11 therefore reduces the proof to the case
$\dim V=2$.  If $W_1=W_2$, the assertion is clear.  Otherwise, if
$W_1=Ae$ and $W_2=Af$ with $\Phi(e,f)=1$, then
$e(W_1)=fe$ and $e(W_2)=ef=1-fe$.

\medskip
2. \underline{Quadrics}

2.1. Let $k$ be an algebraically closed field and $V$ a vector space of
dimension $n+2$ over $k$.  A smooth quadric of dimension $n$ over $k$ is
a $k$-scheme $X$ isomorphic to the subscheme of $\mathbb P(V^*)$ defined
by the equation $Q=0$, for some ordinary quadratic form $Q$ over $k$.

\underline{Remark} 2.2. For $n=0,1,2$, $X$ is a smooth quadric of
dimension $n$ over $k$ if and only if

\begin{equation*}
\begin{array}{rcl}
n=0 &:& X\simeq\Spec(k)\amalg\Spec(k),\\[1ex]
n=1 &:& X\simeq\mathbb P^1_k,\\[1ex]
n=2 &:& X\simeq\mathbb P^1_k\times\mathbb P^1_k.
\end{array}
\end{equation*}
\end{filecontents*}
% Original input: components/021_expose_XII_2_3_through_2_8.tex
\begin{filecontents*}[overwrite,noheader]{sga7ii-cumulative-input-021.tex}
%%%%%%%%%% source scan idx 78 | folio 71 %%%%%%%%%%

\underline{Lemma} 2.3. \underline{With the notation of 2.1, suppose that
$X$ is defined by $Q=0$.}

\begin{enumerate}[label=(\roman*),leftmargin=2.6em]
\item $\Omega_X^n\simeq\mathcal O_X(-n)$.

\item
\begin{equation*}
\operatorname{Pic}(X)\simeq
\begin{cases}
0, & n=0,\\
\mathbb Z, & n=1;\ \mathcal O_X(1)\text{ is the square of a generator},\\
\mathbb Z\oplus\mathbb Z, & n=2,\\
\mathbb Z, & n\geq 3;\ \mathcal O_X(1)\text{ is a generator}.
\end{cases}
\end{equation*}

\item $H^i(X,\mathcal O_X)=0$ for $i>0$, and
$H^1(X,\mathcal O_X(1))=0$.

\item The natural map
\begin{equation*}
H^0\bigl(\mathbb P(V^*),\mathcal O_{\mathbb P(V^*)}(1)\bigr)
=V^*\xrightarrow{\sim}H^0(X,\mathcal O_X(1))
\end{equation*}
is an isomorphism; and, for $n>0$,
$H^0(X,\mathcal O_X)=k$.
\end{enumerate}

Assertion (i) follows from the fact that the conormal sheaf of $X$ in
$\mathbb P(V^*)$ is $\mathcal O_X(-2)$ and that the canonical sheaf of
$\mathbb P^r$ is $\mathcal O(-r-1)$.  Assertion (ii) follows from 2.2
when $n\leq 2$, and from SGA~2, XII, 3.7 when $n\geq 3$.  Finally, (iii)
and (iv) follow from FAC, no.~78.

\underline{Definition} 2.4. \underline{A smooth quadric of dimension
$n$ over a scheme $S$ is an $S$-scheme $f:X\longrightarrow S$, proper
and smooth over $S$, whose geometric fibres are smooth quadrics.}

For $n=0$, a smooth quadric of dimension zero over $S$ is simply a
double \'{e}tale covering of $S$.  For $n=1$, it is simply a
Severi--Brauer variety of dimension one over $S$.

2.5. Let $\rho:X\longrightarrow S$ be a smooth quadric of dimension $n$
over $S$.  By 2.3(i), the inverse of the invertible sheaf
$\Omega^n_{X/S}$ is ample.

%%%%%%%%%% source scan idx 79 | folio 72 %%%%%%%%%%

By 2.3(iii), the Picard scheme $\operatorname{Pic}_{X/S}$ is smooth over
$S$ (because $H^2(X,\mathcal O_X)=0$), unramified over $S$ (because
$H^1(X,\mathcal O_X)=0$), and essentially proper over $S$ (because
$X/S$ is smooth).  Thus, locally on $S$ for the \'{e}tale topology, it
is a constant discrete group having the structure described in
2.3(ii).\footnote{The print cites 2.3(ii) for the cohomological
vanishing and 2.3(iii) for the Picard-group computation.  The two
references are interchanged here.}

In particular, if $X$ has a section, there is an invertible sheaf $L$ on
$X$ such that $L^{\otimes(-n)}$ and $\Omega^n_{X/S}$ have the same image
in $\operatorname{Pic}_{X/S}$; any two such sheaves are locally
isomorphic on $S$.  By 2.3(iii) and [1], 0.5, p.~19,
$\rho_*L$ is locally free and its formation commutes with every base
change.  Hence $\rho_*L$ is locally free of rank $n+2$, the map
$\rho^*\rho_*L\longrightarrow L$ is surjective, and the canonical map
\begin{equation*}
X\longrightarrow\mathbb P(\rho_*L)
\end{equation*}
is a closed immersion.  It identifies $X$ with a relative divisor of
degree two in $\mathbb P(\rho_*L)$; locally on $S$, $X$ is the quadric in
$\mathbb P(\rho_*L)$ defined by an ordinary quadratic form on
$(\rho_*L)^*$, unique up to an invertible factor.

\footnote{After introducing the projection as $\rho$, the print changes
its symbol to $p$ in the direct-image formulas.  The notation is
regularized here; the reference to 2.3(iii) is retained.}

2.6. Suppose that $n>0$.  If $L_1$ and $L_2$ are two invertible sheaves
of the type considered above, then locally on $S$ there is an
isomorphism $u:L_1\xrightarrow{\sim}L_2$.  It is unique up to an
invertible factor, since $\rho_*\mathcal O_X=\mathcal O_S$ by 2.3(iv).
Consequently all the isomorphisms $u:L_1\xrightarrow{\sim}L_2$ induce
the same isomorphism
\begin{equation*}
\mathbb P(u):\mathbb P(\rho_*L_1)\xrightarrow{\sim}
\mathbb P(\rho_*L_2).
\end{equation*}
Thus $\mathbb P(\rho_*L_1)$ is canonically isomorphic to
$\mathbb P(\rho_*L_2)$ and depends only on $X/S$, not on the choice of
$L$.  We denote it by $P(X)$.

%%%%%%%%%% source scan idx 80 | folio 73 %%%%%%%%%%

The projective space $P(X)$ is defined locally on $S$ for the \'{e}tale
topology.  By descent it defines a Severi--Brauer variety over $S$,
again denoted $P(X)$, in which $X$ is a relative divisor of degree two.

Locally on $S$, the pair $X\hookrightarrow P(X)$ defines a Clifford
algebra $C^+(X,P(X))$ (1.3).  Descent yields an algebra over $S$, which
we denote simply by $C^+(X)$.

2.7. Suppose that $n=2m>0$.  The centre $Z$ of $C^+(X)$ is then an
$\mathcal O_S$-algebra defining a double \'{e}tale covering $Z(X)$ of
$S$.

A \underline{generator} of $X$ is a subscheme $D$ of $X$ which is a
linear subspace of dimension $m$ in $P(X)$.

Locally on $S$, when $P(X)\simeq\mathbb P(V^*)$ and $X$ is defined by
$Q=0$, a generator $D$ of $X$ is, by definition, identified with a
locally direct-summand submodule $W(D)$ of $V$, totally singular and of
rank $m+1$.  By 1.8, the idempotent $e(W(D))$ in the centre of $C^+(X)$
depends only on $D\subset X$.  Let $e(D)$ denote the corresponding
section of $Z(X)$, namely the one on which $e(W(D))$ has value one.

The generators of $X/S$ are the $S$-points of a scheme
$\operatorname{Gen}(X)$ projective over $S$, and $e$ defines a morphism
of $S$-schemes

\begin{equation}\tag{2.7.1}
e:\operatorname{Gen}(X)\longrightarrow Z(X).
\end{equation}

When $n=0$, we set $Z(X)=X$, call a section of $X$ a generator of $X$,
and define $e$ to be the identity.

%%%%%%%%%% source scan idx 81 | folio 74 %%%%%%%%%%

\underline{Proposition} 2.8. \underline{Under the hypotheses of 2.7,
$p:\operatorname{Gen}(X)\longrightarrow S$ is projective and smooth,
and (2.7.1) is its Stein factorization.}

The question is local for the \'{e}tale topology on $S$.  Locally, if
$D$ is a generator, then by 1.6 the generators $D'$ disjoint from $D$
are identified with the $S$-points of a suitable affine space over $S$.
In this way $\operatorname{Gen}(X)$ is covered by open subschemes smooth
over $S$.

To prove that the geometric fibres of $e$ are geometrically connected,
it is enough to consider the case in which $S$ is the spectrum of an
algebraically closed field $k$ and $X$ is defined by $Q=0$ in
$\mathbb P(V^*)$.

Let $W_1$ and $W_2$ be two totally isotropic subspaces of rank $m+1$ in
$V$ such that $e(W_1)=e(W_2)$, that is, such that
$W_1/(W_1\cap W_2)$ has even dimension (1.12).  Proceeding as in 1.12,
one obtains an orthogonal decomposition
\begin{equation*}
V=V_0\oplus\bigoplus_i V_i
\end{equation*}
such that
\begin{enumerate}[label=\alph*),leftmargin=2.5em]
\item $W_\alpha=(W_\alpha\cap V_0)\oplus
  \displaystyle\bigoplus_i(W_\alpha\cap V_i)$ for $\alpha=1,2$;
\item $W_1\cap V_0=W_2\cap V_0$;
\item $\dim V_i=4$ and
  $V_i=(W_1\cap V_i)\oplus(W_2\cap V_i)$.
\end{enumerate}
It remains to prove that $W_1\cap V_i$ and $W_2\cap V_i$ belong to the
same connected family of maximal totally singular subspaces of $V_i$.
In other words, the assertion reduces to the case of quadrics of
dimension two.  Such a quadric is isomorphic to
$\mathbb P^1\times\mathbb P^1$, and its generators are the subschemes
$\{t\}\times\mathbb P^1$ and $\mathbb P^1\times\{t\}$; the proposition
is then evident.
\end{filecontents*}
% Original input: components/022_expose_XII_3_1_through_3_5.tex
\begin{filecontents*}[overwrite,noheader]{sga7ii-cumulative-input-022.tex}
\medskip
3. \underline{Cohomology of quadrics}

3.1. Let $p:X\longrightarrow S$ be a smooth quadric of dimension $n$
over $S$, and let $\ell$ be a prime number invertible on $S$.  We write
\begin{equation*}
\eta\in H^0\bigl(S,R^2p_*\mathbb Z_\ell(1)\bigr)
\end{equation*}
for the ``cohomology class of a hyperplane section.''  Locally on $S$
for the \'{e}tale topology, $P(X)\simeq\mathbb P(V^*)$; by definition,
$\eta$ is then the image in
$H^0(S,R^2p_*\mathbb Z_\ell(1))$ of the first Chern class
\begin{equation*}
c_1(\mathcal O(1))\in H^2(X,\mathbb Z_\ell(1)).
\end{equation*}
\footnote{The displayed definition in the print has $R^1p_*$, whereas
the following sentence and the degree of $c_1(\mathcal O(1))$ both give
$R^2p_*$.}

3.2. Suppose that $n=2m$.  A generator $D$ of $X$ then defines a
cohomology class
$c\ell'(D)\in H^n(X,\mathbb Z_\ell(m))$, whose image in
$H^0(S,R^np_*\mathbb Z_\ell(m))$ is denoted $c\ell(D)$.  By 2.8, the
latter depends only on the section $e(D)$ of $Z(X)$.  We therefore obtain
a morphism of \'{e}tale sheaves on $S$
\begin{equation}\tag{3.2.1}
c\ell:Z(X)\longrightarrow R^np_*\mathbb Z_\ell(m).
\end{equation}

\underline{Theorem} 3.3. \underline{Let $p:X\longrightarrow S$ be a
smooth quadric of dimension $n$ over $S$, and let $\ell$ be a prime
number invertible on $S$.}

\begin{enumerate}[label=(\roman*),leftmargin=2.6em]
\item $R^{2i+1}p_*\mathbb Z_\ell=0$.

%%%%%%%%%% source scan idx 83 | folio 76 %%%%%%%%%%

\item If $0\leq 2i<n$ (respectively, if $n<2i\leq 2n$), then
$R^{2i}p_*\mathbb Z_\ell(i)$ is canonically isomorphic to the constant
sheaf $\mathbb Z_\ell$, and is generated by $\eta^i$ (respectively, by
$\eta^i/2$).

\item If $n=2m$, the map deduced from (3.2.1),
\begin{equation}\tag{3.3.1}
\mathbb Z_\ell^{Z(X)}\longrightarrow R^np_*\mathbb Z_\ell(m),
\end{equation}
is an isomorphism.  If $\alpha$ and $\beta$ are two disjoint sections
of $Z(X)$ (assumed isomorphic to $S\amalg S$, as is locally the case),
then
\begin{align*}
\text{(a)}\quad &\eta^m=c\ell(\alpha)+c\ell(\beta),\\[1ex]
\text{(b)}\quad &
\begin{cases}
\operatorname{Tr}(c\ell(\alpha)^2)
=\operatorname{Tr}(c\ell(\beta)^2)=1,
& c\ell(\alpha)c\ell(\beta)=0,\quad m\text{ even},\\[1ex]
c\ell(\alpha)^2=c\ell(\beta)^2=0,
& \operatorname{Tr}(c\ell(\alpha)c\ell(\beta))=1,
  \quad m\text{ odd},
\end{cases}\\[1ex]
\text{(c)}\quad &\eta\bigl(c\ell(\alpha)-c\ell(\beta)\bigr)=0.
\end{align*}
\end{enumerate}

\footnote{The print repeats the tag (3.2.1) on the map in part (iii),
but the proof itself refers to this map as (3.3.1).}

Assertions (i) and (ii) are recalled only for reference (XI, 1.6,
supplemented by XI, 2.6 when $2i+1=n$ in (i)).

It is enough to verify (iii) when $S$ is the spectrum of an
algebraically closed field $k$ (one could even take
$S=\Spec(\mathbb C)$).

\underline{Formula (a).}  We may suppose that $X$ is the quadric in
$\mathbb P^{2m+1}(k)$ with equation
\begin{equation*}
\sum_{i=0}^{m}x_i x_{i+m+1}=0.
\end{equation*}
The class $\eta^m$ is the cohomology class of the plane section defined
by
\begin{equation*}
x_i=0\qquad(0\leq i<m),
\end{equation*}
and this cycle is the sum of the generators $D_1,D_2$ defined by

%%%%%%%%%% source scan idx 84 | folio 77 %%%%%%%%%%

\begin{align*}
D_1:&\quad x_i=0 &&(0\leq i\leq m),\\
D_2:&\quad x_i=0 &&(0\leq i<m),\qquad x_{2m+1}=0.
\end{align*}
By 1.12, $e(D_1)\neq e(D_2)$, and hence
\begin{equation*}
\eta^m=c\ell(D_1)+c\ell(D_2)
=c\ell(\alpha)+c\ell(\beta).
\end{equation*}

\footnote{The printed ambient projective space is
$\mathbb P^{m+1}$, although the displayed nondegenerate equation has
$2m+2$ coordinates.  The required ambient space is
$\mathbb P^{2m+1}$.}

\underline{Formula (b).}  If generators $D_1,D_2$ are disjoint
(respectively, meet transversally in one point), then
$c\ell(D_1)c\ell(D_2)=0$ (respectively,
$\operatorname{Tr}(c\ell(D_1)c\ell(D_2))=1$).  By 1.12, when $m$ is
even we have $e(D_1)\neq e(D_2)$ (respectively,
$e(D_1)=e(D_2)$); when $m$ is odd we have $e(D_1)=e(D_2)$
(respectively, $e(D_1)\neq e(D_2)$).  Formula (b) follows.

\underline{Formula (c).}  This is clear when $n=0$.  Otherwise, by
(ii), it is enough to prove
\begin{equation*}
\eta^m\bigl(c\ell(\alpha)-c\ell(\beta)\bigr)=0.
\end{equation*}
Indeed,
\begin{equation*}
\operatorname{Tr}\!\left(
\eta^m(c\ell(\alpha)-c\ell(\beta))\right)
=\operatorname{Tr}\!\left(c\ell(\alpha)^2-c\ell(\beta)^2\right)=0.
\end{equation*}

\underline{Verification.}  When $m$ is even,
\begin{equation*}
2=\operatorname{Tr}(\eta^{2m})
=\operatorname{Tr}\!\left(c\ell(\alpha)^2
+2c\ell(\alpha)c\ell(\beta)+c\ell(\beta)^2\right)=1+1;
\end{equation*}
when $m$ is odd,
\begin{equation*}
2=\operatorname{Tr}(\eta^{2m})
=\operatorname{Tr}\!\left(c\ell(\alpha)^2
+2c\ell(\alpha)c\ell(\beta)+c\ell(\beta)^2\right)=2.
\end{equation*}
\footnote{In these two top-degree intersection computations the print
has $\eta^2$ in place of $\eta^{2m}$, and once prints
$c\ell(\alpha^2)$ in place of $c\ell(\alpha)^2$.}

That (3.3.1) is an isomorphism now follows from the fact that the two
bilinear forms in two variables
\begin{equation*}
B(x,y)=x_1y_2+x_2y_1,
\qquad
B(x,y)=x_1y_1+x_2y_2
\end{equation*}
have discriminant $\pm1$.

%%%%%%%%%% source scan idx 85 | folio 78 %%%%%%%%%%

\underline{Verification} 3.4. Let $X$ be a smooth quadric of dimension
$n$ over a finite field $\mathbb F_q$.  In view of 3.3, the Lefschetz
trace formula gives the classical formula
\begin{equation*}
\#X(\mathbb F_q)=
\begin{cases}
\displaystyle\sum_{i=0}^{n}q^i,
  & n\text{ odd},\\[2ex]
\displaystyle\sum_{i=0}^{n}q^i+\varepsilon q^m,
  & n=2m,\quad \varepsilon=\pm1.
\end{cases}
\end{equation*}
When $n$ is even, $\varepsilon=1$ if $X$ is split, that is, if it has a
rational generator, and $\varepsilon=-1$ otherwise.

3.5. Let $f:X\longrightarrow S$ be a smooth quadric of dimension
$n=2m$.  The \underline{primitive part} of
$R^nf_*\mathbb Z_\ell(m)$ is the orthogonal complement of $\eta^m$;
with the notation of 3.3, it is generated by
$c\ell(\alpha)-c\ell(\beta)$.  The \underline{primitive quotient} of
$R^nf_*\mathbb Z_\ell(m)$ is the quotient by
$\mathbb Z_\ell\eta^m$.  The primitive part and primitive quotient of
$R^nf_*\mathbb Z_\ell$ are obtained by tensoring with
$\mathbb Z_\ell(-m)$.

\end{filecontents*}
% Original input: components/023_expose_XII_3_6_through_bibliography.tex
\begin{filecontents*}[overwrite,noheader]{sga7ii-cumulative-input-023.tex}
3.6. Let $V$ be a locally free vector bundle of rank $n+2$ on $S$, and
let $X$ be the smooth quadric over $S$ defined by an ordinary quadratic
form $Q$ on $V$.  Let $H$ moreover be a hyperplane in
$\mathbb P(V^*)$ meeting $X$ transversally, and put $X^o=X-H$.  The
schemes $X$ and $Y=X\cap H$ are smooth quadrics over $S$, of dimensions
$n$ and $n-1$.  We may therefore use 3.3 to compute the cohomology of
$X^o=X-Y$.

%%%%%%%%%% source scan idx 86 | folio 79 %%%%%%%%%%

\begin{equation}\tag{3.6.1}
\begin{tikzcd}[column sep=huge,row sep=huge]
X^o \arrow[r,hook] \arrow[rd,"f"']
  & X \arrow[d,"p"]
  & Y \arrow[l,hook] \arrow[ld,"q"]\\
{} & S & \mathord{.}
\end{tikzcd}
\end{equation}

If $n=0$, then $Y=\varnothing$ and $X=X^o$.  Suppose that $n>0$.  We
have two long exact sequences dual to one another:
\begin{equation}\tag{3.6.2}
\cdots\xrightarrow{\partial}R^if_!\mathbb Z_\ell
\longrightarrow R^ip_*\mathbb Z_\ell
\xrightarrow{r_i}R^iq_*\mathbb Z_\ell\longrightarrow\cdots
\end{equation}
and
\begin{equation}\tag{3.6.3}
\cdots\xrightarrow{\partial}R^{i-2}q_*\mathbb Z_\ell(-1)
\longrightarrow R^ip_*\mathbb Z_\ell
\longrightarrow R^if_*\mathbb Z_\ell\longrightarrow\cdots.
\end{equation}

For $n$ even, the restriction homomorphism $r_i$ is an isomorphism for
$i\neq n,2n$; for $n$ odd, it is an isomorphism for $i\neq n-1$.  If
$n=2m$ and $\alpha$ is the cohomology class of a generator, then
\begin{equation*}
r_n(\alpha)=\tfrac12\eta^m,
\end{equation*}
so $r_n$ is surjective and its kernel is the primitive part of
$R^np_*\mathbb Z_\ell$.  If $n=2m+1$, then
$r_{2m}(\eta^m)=\eta^m$, so $r_{2m}$ is injective and its cokernel is the
primitive quotient of $R^{2m}q_*\mathbb Z_\ell$.  It follows that
$R^if_!\mathbb Z_\ell=0$ for $i\neq n,2n$, and that
$R^if_!\mathbb Z_\ell$ is a twisted constant free
$\mathbb Z_\ell$-sheaf of rank one for $i=n,2n$.  By duality, or from
(3.6.3), one likewise shows that $R^if_*\mathbb Z_\ell=0$ for
$i\neq0,n$, that $f_*\mathbb Z_\ell=\mathbb Z_\ell$, and that
$R^nf_*\mathbb Z_\ell$ has rank one.

\footnote{This paragraph prints $r=2m+1$ instead of $n=2m+1$ and cites
3.5.3 instead of the displayed sequence (3.6.3).}

Suppose that $n=2m$ and retain the notation of 3.3.

%%%%%%%%%% source scan idx 87 | folio 80 %%%%%%%%%%

\begin{center}
\begin{tikzcd}[column sep=large,row sep=large]
0 \arrow[r]
 & R^nf_!\mathbb Z_\ell(m) \arrow[r] \arrow[d,"\varphi"]
 & R^np_*\mathbb Z_\ell(m) \arrow[r] \arrow[d,equal]
 & R^nq_*\mathbb Z_\ell(m) \arrow[r]
 & 0\\
0
 & R^nf_*\mathbb Z_\ell(m) \arrow[l]
 & R^np_*\mathbb Z_\ell(m) \arrow[l]
 & R^{n-2}q_*\mathbb Z_\ell(m-1) \arrow[l]
 & 0\mathord{.} \arrow[l]
\end{tikzcd}
\end{center}

The sheaf $R^nf_!\mathbb Z_\ell(m)$ then has two natural opposite
generators $\pm\delta$, whose images in
$R^np_*\mathbb Z_\ell(m)$ are
$\pm(c\ell(\alpha)-c\ell(\beta))$.  By 3.3(iii)(b),
\begin{equation*}
\operatorname{Tr}(\delta^2)
=(c\ell(\alpha)-c\ell(\beta))^2=(-1)^m\cdot2,
\end{equation*}
and $\pm\varphi(\delta)$ is twice the natural dual generators
$\pm\delta'$ of $R^nf_*\mathbb Z_\ell(m)$, namely the opposite images
of $c\ell(\alpha)$ and $c\ell(\beta)$.

Suppose now that $n=2m+1$ and apply the notation of 3.3 to $Y$.

\begin{center}
\begin{tikzcd}[column sep=large,row sep=large]
0 \arrow[r]
 & R^{2m}p_*\mathbb Z_\ell(m) \arrow[r]
 & R^{2m}q_*\mathbb Z_\ell(m) \arrow[r,"\partial"]
 & R^nf_!\mathbb Z_\ell(m) \arrow[r] \arrow[d,"\varphi"]
 & 0\\
{}
 & R^{2m+2}p_*\mathbb Z_\ell(m+1)
 & R^{2m}q_*\mathbb Z_\ell(m) \arrow[l]
 & R^nf_*\mathbb Z_\ell(m+1) \arrow[l]
 & 0\mathord{.} \arrow[l]
\end{tikzcd}
\end{center}

\footnote{The leftmost term in the printed lower row has twist $(m)$;
duality and the other three terms require the displayed twist $(m+1)$.}

The sheaf $R^nf_!\mathbb Z_\ell(m)$ then has two natural opposite
generators $\pm\delta$, namely
$\partial c\ell(\alpha)$ and $\partial c\ell(\beta)$.  We have
\begin{equation*}
\delta\,\varphi(\delta)=\delta^2
=(\partial c\ell(\alpha))^2
=\partial\bigl(c\ell(\alpha)\,\partial c\ell(\alpha)\bigr)=0,
\end{equation*}
and therefore $\varphi(\delta)=0$.  On the other hand,
$R^nf_*\mathbb Z_\ell(m+1)$ has two natural opposite generators
$\pm\delta'$, whose images in $R^{2m}q_*\mathbb Z_\ell(m)$ are
$\pm(c\ell(\alpha)-c\ell(\beta))$, and
$\operatorname{Tr}(\delta\delta')=\pm1$.

%%%%%%%%%% source scan idx 88 | folio 81 %%%%%%%%%%

These results are assembled in the following table.

\underline{Table} 3.7. \underline{Cohomology of affine quadrics of
dimension $n$, $X^o=X-Y$.}

\medskip
\noindent\textbf{A.} The cohomology is torsion-free.  The nonzero Betti
numbers are
\begin{align*}
\text{ordinary cohomology:}\quad &b_0=b_n=1,\\
\text{cohomology with compact supports:}\quad &b_{2n}=b_n=1,
\end{align*}
except that $b_0=2$ when $n=0$.

\medskip
\noindent\textbf{B.} If $n=2m>0$, then
\begin{align*}
R^nf_!\mathbb Z_\ell(m)
  &=\text{the primitive part of }R^{2m}p_*\mathbb Z_\ell(m),\\
R^nf_*\mathbb Z_\ell(m)
  &=\text{the primitive quotient of }R^{2m}p_*\mathbb Z_\ell(m).
\end{align*}
If $n=2m+1$, then
\begin{align*}
R^nf_!\mathbb Z_\ell(m)
  &=\text{the primitive quotient of }R^{2m}q_*\mathbb Z_\ell(m),\\
R^nf_*\mathbb Z_\ell(m+1)
  &=\text{the primitive part of }R^{2m}q_*\mathbb Z_\ell(m).
\end{align*}
Locally, these groups have natural generators defined up to sign;
$\delta$ denotes those for cohomology with compact supports, and
$\delta'$ those for ordinary cohomology.

\medskip
\noindent\textbf{C.} The map
\begin{equation*}
\varphi:R^nf_!\mathbb Z_\ell\longrightarrow R^nf_*\mathbb Z_\ell
\end{equation*}
is zero when $n$ is odd, and sends $\pm\delta$ to $\pm2\delta'$ when
$n>0$ is even.  Moreover,
\begin{equation*}
\operatorname{Tr}(\delta\delta')=\pm1,
\qquad
\operatorname{Tr}(\delta^2)=
\begin{cases}
0, & n=2m+1,\\
(-1)^m\cdot2, & n=2m.
\end{cases}
\end{equation*}

\underline{Verification} 3.8. The complex affine quadric $S$ in
$\mathbb C^n$ with equation $\sum_i z_i^2=1$ is diffeomorphic to the
tangent bundle of the real sphere $S\cap\mathbb R^n$.  This agrees with
A and C.

\medskip
\underline{BIBLIOGRAPHY}

\begin{thebibliography}{9}
\bibitem{MumfordGIT}
D. Mumford, \emph{Geometric Invariant Theory}, Ergebnisse \textbf{34},
Springer-Verlag, 1965.
\end{thebibliography}
\end{filecontents*}
% Original input: components/024_expose_XIII_title_introduction_through_0_2.tex
\begin{filecontents*}[overwrite,noheader]{sga7ii-cumulative-input-024.tex}
\clearpage
\begin{center}
{\large SGA 7 --- EXPOS\'E XIII}\par
\vspace{2em}
{\Large\underline{The Formalism of Vanishing Cycles}}\par
\vspace{1.5em}
{\large by P. Deligne}
\end{center}

\medskip
\noindent\textbf{Source contents}
\begin{center}
\begin{tabular}{@{}llr@{}}
Introduction && 2\\
\S1. & Sheaves of sets & 6\\
1.1 & Galois sheaves & 6\\
1.2 & The topos $Y\times_s S$ & 8\\
1.3 & The functor $\Psi$ & 13\\
1.4 & Variation & 16\\
\S2. & Vanishing cycles & 17\\
2.1 & The functor $R\Psi$ & 17\\
2.2 & A compatibility (trace morphisms) & 26\\
2.3 & Finiteness theorem (equal characteristic zero) & 29\\
2.4 & Isolated singularities & 30
\end{tabular}
\end{center}

%%%%%%%%%% source scan idx 90 | folio 83 %%%%%%%%%%

\medskip
\underline{Introduction}

In this expos\'e we take up again the formalism introduced in no.~2 of
Grothendieck's Expos\'e I.  The only new point is the introduction of the
algebraic analogue of \underline{variation} in the classical theory.

The topological picture I have in mind is as follows.  Let $X$ be an
analytic space, let $D\subset\mathbb C$ be a disc centred at $0$, put
$D^*=D-\{0\}$, and let $f:X\longrightarrow D$ be a proper morphism (the
nonproper case will be considered later).  Replace $D$ by a sufficiently
small disc centred at $0$, replace $X$ by the inverse image of this small
disc, and continue to denote the result by $f:X\longrightarrow D$.
According to the complex-analytic variant of Thom's isotopy theorem,
$f^{-1}(D^*)$ is a topological fibre bundle over $D^*$.  I do not know
whether one can always construct a pair $(\Gamma,r)$ of the kind
described below, but it seems plausible that one can.

\begin{enumerate}[label=(\alph*),leftmargin=2.6em]
\item $\Gamma$ is a locally coherent system of local trivializations of
the bundle $f^{-1}(D^*)$ over $D^*$.  In other words, $\Gamma$ defines a
$\pi_1(D^*)$-equivariant trivialization of the inverse image of $X$ on a
universal covering $\widetilde D^*$ of $D^*$.

\item $r$ is a retraction of $X$ onto the special fibre
$X_0=f^{-1}(0)$, compatible with $\Gamma$: if, on a connected open
subset $U$ of $D^*$, $\Gamma$ identifies
$X_U=f^{-1}(U)$ with $X_1\times U$, then $r=r(x_1,u)$ is independent of
$u$.
\end{enumerate}

Let $(X,f,\Gamma,r)$ be given.  Let $t\in D$ and
$X_t=f^{-1}(t)$ be the ``general fibre.''  The bundle $\varphi^*X$,
obtained from $X/D$ by pullback along
\begin{equation*}
\varphi:[0,1]\longrightarrow D^*,
\qquad x\longmapsto \exp(2\pi i x)\,t,
\end{equation*}
is trivialized by $\Gamma$.  The \underline{monodromy}
$T:X_t\longrightarrow X_t$ is the homeomorphism
\begin{equation*}
X_t=(\varphi^*X)_0\xrightarrow[\Gamma]{\sim}
(\varphi^*X)_1=X_t.
\end{equation*}
Let $r_t:X_t\longrightarrow X_0$ be the restriction of $r$ to $X_t$.
Then $r_tT=r_t$.  One reconstructs $(X,f)$ from
$(X_t,X_0,T,r_t)$ as follows:
\begin{enumerate}[label=(\alph*),leftmargin=2.6em]
\item identify the ends $X_t\times\{0\}$ and $X_t\times\{1\}$ of
$X_t\times[0,1]$ using $T$, obtaining $f':X'\longrightarrow S^1$;
\item $r_t$ defines $r':X'\longrightarrow X_0$, and $(X,f)$ is
identified with the map of mapping cones
\begin{equation*}
\operatorname{cone}(f'):
\operatorname{cone}(r':X'\longrightarrow X_0)
\longrightarrow D=\operatorname{cone}(S^1\longrightarrow\mathrm{pt}).
\end{equation*}
\end{enumerate}

%%%%%%%%%% source scan idx 91 | folio 84 %%%%%%%%%%

Put $\Psi=(r,f):X\longrightarrow X_0\times D$.  The higher Leray direct
images $R^i\Psi_*\mathbb Z$ satisfy the following condition:

\begin{quote}
(*) The space $D^*$ is covered by open subsets $U$ such that the
restriction of the sheaf $F$ to $X_0\times U$ is pulled back from a sheaf
on $X_0$.
\end{quote}

\underline{Construction} 0.1. \underline{The sheaves $F$ on
$X_0\times D$ satisfying (*) are identified with systems
$(F_0,F_t,T,\alpha)$ in which}
\begin{enumerate}[label=\alph*),leftmargin=2.5em]
\item \underline{$F_0$ is a sheaf on $X_0$};
\item \underline{$F_t$ is a sheaf on $X_0$, endowed with an action of
the generator $T$ of $\pi_1(D^*,t)$};
\item \underline{$\alpha$ is a morphism from $F_0$ to the sheaf of
$T$-invariants of $F_t$}.
\end{enumerate}

Let $G$ be a sheaf on $X_0\times D^*$ satisfying (*), and let
$\widetilde D^*$ be the universal covering of $D^*$, with base point
$t$.  The pullback of $G$ to $X_0\times\widetilde D^*$ is, in one and
only one way, pulled back from a sheaf $G_t$ on $X_0$.  The sheaf $G_t$
is endowed with an action of $\pi_1(D^*,t)$, and $G$ is obtained from
$\operatorname{pr}_1^*G_t$ on $X_0\times\widetilde D^*$ by taking the
quotient by $\pi_1(D^*,t)$.

A sheaf $F$ on $X_0\times D$ is determined by
\begin{enumerate}[label=\arabic*),leftmargin=2.5em]
\item its restriction $F_0$ to $X_0$;
\item its restriction $G$ to $X_0\times D^*$;
\item a morphism $\alpha:F_0\longrightarrow i^*j_*G$, where
$i:X_0\hookrightarrow X_0\times D$ and
$j:X_0\times D^*\hookrightarrow X_0\times D$.
\end{enumerate}
If $F$ satisfies (*), then $G$ is described by $(G_t,T)$ on $X_0$, and
the asserted description of $F$ follows.

For the sheaves $R^i\Psi_*\mathbb Z$ we have
\begin{equation*}
(R^i\Psi_*\mathbb Z)_0=
\begin{cases}
\mathbb Z, & i=0,\\
0, & i>0,
\end{cases}
\end{equation*}
whereas, for $i>0$, the $(R^i\Psi_*\mathbb Z)_t$ deserve the name
\underline{sheaves of vanishing cycles} (cf. the discussion at the
beginning of paragraph 2 of Expos\'e I).

%%%%%%%%%% source scan idx 92 | folio 85 %%%%%%%%%%

This discussion is useful only after it is refined by passing to the
derived category: the $R^i\Psi_*\mathbb Z$ are obtained from an object
of $D^+(X_0\times D)$, and indeed from an object
$R\Psi_*\mathbb Z$ of the derived category of the category of abelian
sheaves on $X_0\times D$ satisfying (*).  This object determines not
only the sheaves of vanishing cycles and their monodromy, but also the
classical ``variation'' morphisms (cf. 1.4 and XIV, 3.1).

We shall define its analogue in algebraic geometry.  If $S$ is a
strictly henselian trait and $f:X\longrightarrow S$ is a morphism of
schemes, we shall define an object $R\Psi(\mathbb Z/\ell^n)$ of the
derived category of the category of abelian sheaves on the product
topos $(X_0)_{\mathrm{et}}\times S_{\mathrm{et}}$.  A sheaf on this
topos may be interpreted as a system $(F_0,F_t,T,\alpha)$ in which
\begin{enumerate}[label=\alph*),leftmargin=2.5em]
\item[$\alpha$)] $F_0$ is a sheaf on $X_0$;
\item[$\beta$)] $F_t$ is a sheaf on $X_0$, endowed with a continuous action $T$
of the inertia group $I$;
\item[$\gamma$)] $\alpha$ is a morphism from $F_0$ to the sheaf of $I$-invariants
of $F_t$.
\end{enumerate}

We shall have to work in a slightly more general framework than that
described above:
\begin{enumerate}[label=\alph*),leftmargin=2.5em]
\item The construction of $R\Psi(\mathbb Z/\ell^n)$ does not require
$f$ to be proper.  For a topological interpretation, one would repeat
the preceding construction after replacing $X$ by a neighbourhood of
the special fibre $X_0$.

\item To construct $R\Psi(\mathbb Z/\ell^n)$ and prove its properties,
one must consider not only the constant sheaf $\mathbb Z/\ell^n$, but
also general sheaves and even complexes of sheaves.

\item When the residue field $k$ of the henselian trait $S$ is not
assumed separably closed, the product topos
$(X_0)_{\mathrm{et}}\times S_{\mathrm{et}}$ must be replaced by the
fibred product topos
\begin{equation*}
(X_0)_{\mathrm{et}}
\times_{\Spec(k)_{\mathrm{et}}}S_{\mathrm{et}}.
\end{equation*}
The reader who dislikes $2$-fibred products of topoi may take the Galois
description of this topos in 1.2.4 as its definition.
\end{enumerate}

%%%%%%%%%% source scan idx 93 | folio 86 %%%%%%%%%%

0.2. \underline{Notation and terminology}.

\begin{description}
\item[(0.2.1) \underline{trait}:]
the spectrum of a discrete valuation ring.

\item[(0.2.2) \underline{henselian}:]
for a scheme, the spectrum of a henselian local ring.

\item[(0.2.3) \underline{strictly henselian}:]
for a ring, henselian with separably closed residue field; for a scheme,
the spectrum of such a ring.

\item[(0.2.4) \underline{geometric point} of $S$:]
a morphism $\bar x:\Spec(\bar k)\longrightarrow S$ with image $x\in S$,
where $k(\bar x)=_{\mathrm{def}}\bar k$ is a separable closure of
$\kappa(x)$.  We say that $\bar x$ is localized at $x$.  By abuse of
language, we shall sometimes also call a morphism
$\bar x:\Spec(\bar k)\longrightarrow S$ a geometric point when $\bar k$
is a separably closed extension of $\kappa(x)$; such a morphism factors
uniquely through a geometric point of $S$ in the proper sense.

\item[(0.2.5) $s,\eta,\bar s,\bar\eta,\bar S$:]
If $S=\Spec(V)$ is a henselian trait, we shall normally denote by $s$
the closed point of $S$, by $\eta$ its generic point, by $\bar\eta$ a
generic geometric point of $S$ (localized at $\eta$), and by $\bar s$
the corresponding geometric point localized at $s$ (I, 0.0.3).  Let
$\bar S$ denote the spectrum of the normalization of $V$ in
$k(\bar\eta)$; its residue field is an inseparable extension of
$k(\bar s)$.  Write $\operatorname{Gal}(\bar\eta/\eta)$ for the Galois
group of $k(\bar\eta)$ over $k(\eta)$, and similarly for $s$.  There is
an exact sequence
\begin{equation*}
0\longrightarrow I\longrightarrow
\operatorname{Gal}(\bar\eta/\eta)\longrightarrow
\operatorname{Gal}(\bar s/s)\longrightarrow0,
\end{equation*}
where $I$ is the inertia group.

\item[(0.2.6) \underline{sheaf}:]
on a scheme $X$, always means a sheaf on the \'{e}tale site
$X_{\mathrm{et}}$ of $X$.

\item[(0.2.7) \underline{prime to}:]
a torsion sheaf $F$ on a scheme $S$ (or a torsion abelian group $F$) is
prime to the residue characteristics of $S$ if, for every $s\in S$,
multiplication by the characteristic exponent of $k(s)$ is an
automorphism of $F$.
\end{description}
\end{filecontents*}
% Original input: components/025_expose_XIII_section_1_opening_through_1_1.tex
\begin{filecontents*}[overwrite,noheader]{sga7ii-cumulative-input-025.tex}
%%%%%%%%%% source scan idx 94 | folio 87 %%%%%%%%%%

\medskip
1. \underline{Sheaves of sets}

The sheaves of vanishing cycles will be defined as values of derived
functors of the left exact functor $\Psi$ defined in no.~3 for every
sheaf of sets.

\medskip
1.1. \underline{Galois sheaves.}

1.1.1. Let $Y$ be a scheme over a field $k$ with separable closure
$\bar{k}$, and put $\bar{Y}=Y\otimes_k\bar{k}$.  Let $G$ be a profinite
group and let
\begin{equation*}
u:G\longrightarrow\operatorname{Gal}(\bar{k}/k)
\end{equation*}
be a continuous homomorphism from $G$ to the Galois group of $\bar{k}$
over $k$.  The Galois group $\operatorname{Gal}(\bar{k}/k)$ acts on the
left on $\bar{Y}$ by transport of structure.  The group $G$ acts on
$\bar{Y}$ through $u$.

Let $\mathcal F$ be a sheaf of sets on (the \'{e}tale site of) $\bar Y$.
An action of $G$ on $\mathcal F$ \underline{compatible} with the action
of $G$ on $\bar Y$ is an action of $G$ by automorphisms of the pair
$(\bar Y,\mathcal F)$ which induces the given action on $\bar Y$.  In
other words, it is a system of isomorphisms
\begin{equation*}
\sigma(g):u(g)_*\mathcal F\longrightarrow\mathcal F
\end{equation*}
satisfying $\sigma(gh)=\sigma(g)\,\sigma(h)$.  If $G$ acts on
$\mathcal F$, then for every $U$ \'{e}tale over $Y$, with
$\bar U=U\otimes_k\bar k$ \'{e}tale over $\bar Y$, the group $G$ acts on
$\mathcal F(\bar U)$.

\underline{Definition 1.1.2.} \underline{We say that $G$ acts
continuously on $\mathcal F$ if, for every quasi-compact $U$ \'{e}tale
over $Y$, the action of $G$ on the discrete set $\mathcal F(\bar U)$ is
continuous.}

It would amount to the same thing to consider only affine $U$.

If $\mathcal F$ is a sheaf of sets on $Y$, with inverse image
$\bar{\mathcal F}$ on $\bar Y$, then
$\operatorname{Gal}(\bar k/k)$ acts on $\bar{\mathcal F}$ by transport
of structure, compatibly with its action on $\bar Y$.

\underline{Reminder 1.1.3.}
\begin{enumerate}[label=(\roman*),leftmargin=2.7em]
\item \underline{The action of $\operatorname{Gal}(\bar k/k)$ on
$\bar{\mathcal F}$ is continuous.}

%%%%%%%%%% source scan idx 95 | folio 88 %%%%%%%%%%

\item \underline{The functor}
\begin{equation*}
\mathcal F\longmapsto\bar{\mathcal F}
\end{equation*}
\underline{where $\bar{\mathcal F}$ is endowed with the action of
$\operatorname{Gal}(\bar k/k)$, is an equivalence from the category of sheaves of sets on
$Y$ to the category of sheaves of sets on $\bar Y$ endowed with a
continuous action of $\operatorname{Gal}(\bar k/k)$ compatible with the
action of $\operatorname{Gal}(\bar k/k)$ on $\bar Y$.}
\end{enumerate}

Write $\bar k$ as the direct limit of finite Galois extensions $k_i$.
Let
\begin{equation*}
\varphi_i:Y_i=Y\otimes_k k_i\longrightarrow Y,
\qquad
\bar\varphi:\bar Y\longrightarrow Y
\end{equation*}
be the projections.  For $U$ \'{e}tale over $Y$, assumed affine, put
$U_i=U\otimes_k k_i$.  Then
\begin{equation*}
\bar{\mathcal F}(\bar U)
=\varinjlim_i(\varphi_i^*\mathcal F)(U_i).
\end{equation*}
On $(\varphi_i^*\mathcal F)(U_i)$, the group
$\operatorname{Gal}(\bar k/k)$ acts through
$\operatorname{Gal}(k_i/k)$, which proves (i).

The functor in (ii) has as right adjoint the functor
\begin{equation*}
\mathcal G\longmapsto
(\bar\varphi_*\mathcal G)^{\operatorname{Gal}(\bar k/k)}.
\end{equation*}
For every $\mathcal F$ on $Y$, one has
\begin{equation*}
\bar\varphi_*\bar\varphi^*\mathcal F
=\varinjlim_i\varphi_{i*}\varphi_i^*\mathcal F.
\end{equation*}
Thus, to prove that
\begin{equation*}
\mathcal F\xrightarrow{\sim}
(\bar\varphi_*\bar\varphi^*\mathcal F)^{
  \operatorname{Gal}(\bar k/k)},
\end{equation*}
it is enough to prove, for every $i$, that
\begin{equation*}
\mathcal F\xrightarrow{\sim}
(\varphi_{i*}\varphi_i^*\mathcal F)^{
  \operatorname{Gal}(k_i/k)}.
\end{equation*}
It is even enough to prove that this arrow becomes an isomorphism after
extension of scalars from $k$ to $k_i$; then $\operatorname{Spec}(k_i)$
becomes a sum of copies of $\operatorname{Spec}(k)$ permuted simply
transitively by $\operatorname{Gal}(k_i/k)$, and the assertion becomes
trivial.

The exact functor in (ii) is therefore fully faithful.  To conclude, it
remains to verify that the arrow
\begin{equation*}
\bar\varphi^*(\bar\varphi_*\mathcal G)^{
  \operatorname{Gal}(\bar k/k)}\longrightarrow\mathcal G
\end{equation*}
is an epimorphism.  Let $\bar\varphi_i:\bar Y\longrightarrow Y_i$.
The hypothesis implies that
\begin{equation*}
\varinjlim_i
\bar\varphi_i^*(\bar\varphi_{i*}\mathcal G)^{
  \operatorname{Gal}(\bar k/k_i)}
\longrightarrow\mathcal G
\end{equation*}
%%%%%%%%%% source scan idx 96 | folio 89 %%%%%%%%%%

is an epimorphism.

We have
\begin{equation*}
\bar\varphi^*(\bar\varphi_*\mathcal G)^{
  \operatorname{Gal}(\bar k/k)}
=
\bar\varphi_i^*\!\left(
  \varphi_i^*\varphi_{i*}
  \left((\bar\varphi_{i*}\mathcal G)^{
    \operatorname{Gal}(\bar k/k_i)}\right)^{
    \operatorname{Gal}(k_i/k)}
\right),
\end{equation*}
and the conclusion follows by observing, by the same Galois-descent
argument as above, that for $\mathcal G'$ on $Y_i$, endowed with an
action of $\operatorname{Gal}(k_i/k)$ compatible with its action on
$Y_i$, one has
\begin{equation*}
\varphi_i^*(\varphi_{i*}\mathcal G')^{
  \operatorname{Gal}(k_i/k)}
\xrightarrow{\sim}\mathcal G'.
\end{equation*}
\end{filecontents*}
% Original input: components/026_expose_XIII_section_1_2_opening_through_1_2_4.tex
\begin{filecontents*}[overwrite,noheader]{sga7ii-cumulative-input-026.tex}
\subsection*{1.2. The topos $Y\times_s S$}

\subsubsection*{1.2.1.}
Let $S$ be a henselian local scheme with closed point
$i:s\hookrightarrow S$.  The functor
\begin{equation*}
S'\longmapsto S'\times_Ss=i^*S'
\end{equation*}
is an equivalence from the category of finite \'{e}tale $S$-schemes
$S'$ to the category of finite \'{e}tale $s$-schemes $s'$ (spectra of
\'{e}tale $k(s)$-algebras).  The inverse functor $\mathrm{sp}^*$, from
the \'{e}tale site of $s$ to the \'{e}tale site of $S$, is a morphism
of sites
\begin{equation*}
\mathrm{sp}:S\longrightarrow s,
\end{equation*}
the \underline{specialization morphism} (cf. SGA~4, VIII, 7).  If
$s'$ is \'{e}tale over $s$, one says that $\mathrm{sp}^*(s')$ is
obtained from $S$ by \underline{extension of the residue field}.  For
a sheaf $F$ on $S$, there is a functorial isomorphism
\begin{equation}\tag{1.2.1.1}
\mathrm{sp}_*F=i^*F.
\end{equation}
For every open subscheme of $S$, $j:U\hookrightarrow S$, or for every
point of $S$, $j:\eta\hookrightarrow S$, we shall again write
$\mathrm{sp}$ for the composite morphism $\mathrm{sp}\circ j$.  For a
sheaf $F$ on $U$ (respectively on $\eta$), one has
\begin{equation}\tag{1.2.1.2}
\mathrm{sp}_*F=i^*j_*F.
\end{equation}

\subsubsection*{1.2.2.}
Let $S$ be a henselian trait, and let
$s,\eta,\bar s,\bar\eta,I$ be as in 0.2.5.
\begin{center}
\begin{tikzcd}
s \arrow[r, hook, "i"] & S & \eta \arrow[l, hook, "j"']
\end{tikzcd}
.\end{center}
Here is the Galois-theoretic interpretation of sheaves on $s$, $\eta$,
and $S$, and of the specialization morphism.

\begin{enumerate}[label=(\alph*)]
\item Sheaves $F$ on $s$ (respectively on $\eta$) are identified, by
$F\mapsto F_{\bar s}$ (respectively $F\mapsto F_{\bar\eta}$), with
sets endowed with a continuous action of $\operatorname{Gal}(\bar s/s)$
(respectively $\operatorname{Gal}(\bar\eta/\eta)$) (SGA~4, VIII, 2.1,
or 1.1.3).

\item A sheaf $F$ on $S$ defines sheaves $F_s=i^*F$ and
$F_\eta=j^*F$ on $s$ and $\eta$, and an adjunction arrow
$F\longrightarrow j_*j^*F$ inducing
\begin{equation*}
\varphi:F_s=i^*F\longrightarrow i^*j_*j^*F
=\mathrm{sp}_*F_\eta.
\end{equation*}
It is known (SGA~4, IV, 9.5.4) that the functor
\begin{equation*}
F\longmapsto(F_s,F_\eta,\varphi)
\end{equation*}
is an equivalence from the category of sheaves on $S$ to the category
of triples $(F_s,F_\eta,\varphi)$, where $F_s$ is a sheaf on $s$,
$F_\eta$ is a sheaf on $\eta$, and
$\varphi:F_s\longrightarrow i^*j_*F_\eta$.

Under the equivalence in (a), the functor $i^*j_*$ is identified with
the functor of ``invariants under $I$''.  Thus sheaves on $S$ are
identified with triples consisting of a set $F_{\bar s}$ endowed with
an action of $\operatorname{Gal}(\bar s/s)$, a set $F_{\bar\eta}$
endowed with an action of $\operatorname{Gal}(\bar\eta/\eta)$, and an
equivariant morphism
$\varphi:F_{\bar s}\longrightarrow F_{\bar\eta}$.

\item Under (a) and (b), the morphisms $\mathrm{sp}$, $j$, $i$, and
$\mathrm{sp}=\mathrm{sp}\circ j$ are expressed as follows:
\begin{equation*}
\begin{array}{llcl}
\mathrm{sp}:S\longrightarrow s:\;&
 \mathrm{sp}^*(F_{\bar s})&=&(F_{\bar s},F_{\bar s},\mathrm{Id}),\\[2pt]
&\mathrm{sp}_*(F_{\bar s},F_{\bar\eta},\varphi)&=&F_{\bar s},\\[3pt]
j:\eta\longrightarrow S:\;&
 j^*(F_{\bar s},F_{\bar\eta},\varphi)&=&F_{\bar\eta},\\[2pt]
&j_*(F_{\bar\eta})&=&(F_{\bar\eta}^{I},F_{\bar\eta},
 \text{inclusion}),\\[3pt]
i:s\longrightarrow S:\;&
 i^*(F_{\bar s},F_{\bar\eta},\varphi)&=&F_{\bar s},\\[2pt]
&i_*(F_{\bar s})&=&(F_{\bar s},\{0\},\text{projection}),\\[3pt]
\mathrm{sp}=\mathrm{sp}\circ j:\eta\longrightarrow s:\;&
 \mathrm{sp}^*(F_{\bar s})&=&F_{\bar s},\\[2pt]
&\mathrm{sp}_*(F_{\bar\eta})&=&F_{\bar\eta}^{I}.
\end{array}
\end{equation*}
\end{enumerate}

\subsubsection*{1.2.3.}
Let $S$ be a henselian trait as in 1.2.2, and let $Y$ be a scheme over
$s$.  We shall have to work with the 2-fiber product $Y\times_sS$ of
the \'{e}tale topoi of $Y$ and $S$ over the \'{e}tale topos of $s$.
By Giraud~[1], this 2-fiber product exists, and it is readily seen to
admit the Galois-theoretic description given below.  Those who dislike
2-fiber products may take this description as the definition.

\subsubsection*{Construction 1.2.4.}
\underline{Sheaves $F$ on $Y\times_sS$ are identified with triples
$(F_s,F_\eta,\varphi)$ in which:}
\begin{enumerate}[label=(\alph*)]
\item $F_s$ is a sheaf on $Y$, or equivalently a sheaf $F_{\bar s}$
on $\bar Y=Y\otimes_s\bar s$, endowed with a continuous action of
$\operatorname{Gal}(\bar s/s)$ compatible with the action of
$\operatorname{Gal}(\bar s/s)$ on $\bar Y$;

\item $F_\eta$ is a sheaf $F_{\bar\eta}$ on $\bar Y$, endowed with a
continuous action of $\operatorname{Gal}(\bar\eta/\eta)$ compatible
with the action (via 1.1.1) of $\operatorname{Gal}(\bar\eta/\eta)$ on
$\bar Y$;

\item $\varphi$ is an equivariant morphism from $F_{\bar s}$ to
$F_{\bar\eta}$.
\end{enumerate}
Similarly, a \underline{sheaf $F_\eta$ on $Y\times_s\eta$} is an
object as in 1.2.4(b).  There is a diagram of topoi
\begin{equation}\tag{1.2.4.1}
\begin{tikzcd}[column sep=large, row sep=small]
Y \arrow[dd] & {} & Y \arrow[r, hook, "i"] \arrow[dd]
  & Y\times_sS \arrow[dd]
  & Y\times_s\eta \arrow[l, hook', "j"'] \arrow[dd] & {} \\
{} & {} \arrow[l, "\mathrm{sp}"'] & {} & {} & {} & {} \\
s & {} & s \arrow[r, hook, "i"] & S
  & \eta \arrow[l, hook', "j"'] & {.}
\end{tikzcd}
\end{equation}
The topos $Y\times_sS$ is the union of the open subtopos
$Y\times_s\eta$ and its closed complement $Y$.  The formulas in
1.2.2(c) remain valid for
$\mathrm{sp}:Y\times_sS\longrightarrow Y$ (also denoted
$\mathrm{pr}_1$), $j:Y\times_s\eta\hookrightarrow Y\times_sS$,
$i:Y\hookrightarrow Y\times_sS$, and
$\mathrm{sp}:Y\times_s\eta\longrightarrow Y$ (also denoted
$\mathrm{pr}_1$).

\paragraph{1.2.4.2.}
If 1.2.4 is taken as the definition of $Y\times_sS$, one must show by
a direct construction that, up to unique isomorphism, $Y\times_sS$
does not depend on the choice of a separable closure $k(\bar\eta)$ of
$k(\eta)$.  This is done by interpreting the $F_\eta$ of 1.2.4(b) as
a functor which, to each separable closure $k(\bar\eta)$ of $k(\eta)$
defining a separable closure $k(\bar s)$ of $k(s)$, associates a sheaf
on $Y\otimes_{k(s)}k(\bar s)$; the $F_\eta$ of 1.2.4(b) is the
restriction of this functor to $\{\bar\eta\}$.

\end{filecontents*}
% Original input: components/027_expose_XIII_section_1_2_5_through_1_2_9.tex
\begin{filecontents*}[overwrite,noheader]{sga7ii-cumulative-input-027.tex}
\subsubsection*{1.2.5.}
Every geometric point $\bar x$ of $Y_{\bar s}$ defines a geometric
point $(\bar x,\bar\eta)$ of $Y\times_s\eta$, whose corresponding
fiber functor is
\begin{equation*}
F_\eta\longmapsto(F_{\bar\eta})_{\bar x}.
\end{equation*}
It defines geometric points $(\bar x,\bar\eta)$ and
$(\bar x,\bar s)$ of $Y\times_sS$ by
\begin{equation*}
F_{(\bar x,\bar\eta)}=(F_{\bar\eta})_{\bar x}
\qquad\text{and}\qquad
F_{(\bar x,\bar s)}=(F_{\bar s})_{\bar x}.
\end{equation*}
One verifies that all the points of the topoi $Y\times_s\eta$ and
$Y\times_sS$ are obtained in this way.  We shall not use this.  One
also verifies at once that these sets of points are conservative.

\subsubsection*{${}^{*}$Remark 1.2.6.}
The construction in 1.2.5 defines an equivalence from the category
$\operatorname{Point}(Y\times_s\eta)$ (respectively
$\operatorname{Point}(Y\times_sS)$) of points of $Y\times_s\eta$
(respectively of $Y\times_sS$) to the 2-fiber product
\begin{equation*}
\operatorname{Point}(Y)\times_{\operatorname{Point}(s)}
\operatorname{Point}(\eta).
\end{equation*}
This equivalence also follows from the fact that $Y\times_s\eta$
(respectively $Y\times_sS$) is a 2-fiber product of topoi.${}^{*}$

\subsubsection*{1.2.7.}
By the functoriality of 2-fiber products, diagram (1.2.4.1) is
functorial in $Y$ and $S$.  In particular:
\begin{enumerate}[label=\alph*)]
\item Every morphism $f:Y\longrightarrow Y'$ of schemes over $s$
induces a morphism of topoi from $Y\times_sS$ to $Y'\times_sS$.  One
has
\begin{equation}\tag{1.2.7.1}
f^*(F_{\bar s},F_{\bar\eta},\varphi)
=(f^*F_{\bar s},f^*F_{\bar\eta},f^*(\varphi)),
\end{equation}
and, when $f$ is quasi-compact,
\begin{equation}\tag{1.2.7.2}
f_*(F_{\bar s},F_{\bar\eta},\varphi)
=(f_*F_{\bar s},f_*F_{\bar\eta},f_*(\varphi)).
\end{equation}
Analogous formulas hold for
$f:Y\times_s\eta\longrightarrow Y'\times_s\eta$.

\item Every surjective morphism of henselian traits
$f:S'\longrightarrow S$ induces a morphism of topoi from
$Y_{s'}\times_{s'}S'$ to $Y\times_sS$.  Again writing $f$ for the
natural morphism
$Y_{s'}=Y\otimes_{k(s)}k(s')\longrightarrow Y$, one has
\begin{equation}\tag{1.2.7.3}
f^*(F_{\bar s},F_{\bar\eta},\varphi)
=(f^*F_{\bar s},f^*F_{\bar\eta},f^*(\varphi)).
\end{equation}
When $f$ is finite, $f_*$ is expressed as an induced representation.
\end{enumerate}

\subsubsection*{1.2.8.}
Let $f:Y\longrightarrow Y'$ be a separated morphism, locally of finite
type, between $s$-schemes.  Define a ``direct image with proper support''
functor
\begin{equation*}
f_!: (\text{sheaves of pointed sets on }Y\times_sS)
\longrightarrow
(\text{sheaves of pointed sets on }Y'\times_sS)
\end{equation*}
by the formula
\begin{equation}\tag{1.2.8.1}
f_!(F_{\bar s},F_{\bar\eta},\varphi)
=(f_!F_{\bar s},f_!F_{\bar\eta},f_!(\varphi)).
\end{equation}
When $f$ is a locally closed immersion, $f_!$ is called the
extension-by-zero functor.

An analogous functor is defined for $Y\times_s\eta$.

\subsubsection*{1.2.9.}
Let $f:Y\hookrightarrow Y'$ be a locally closed immersion of
$s$-schemes.  The functor $f_!$ has a right adjoint $f^!$, ``sections
with support in $Y$'', satisfying
\begin{equation}\tag{1.2.9.1}
f^!(F_{\bar s},F_{\bar\eta},\varphi)
=(f^!F_{\bar s},f^!F_{\bar\eta},f^!(\varphi)).
\end{equation}
Suppose only that $f$ is quasi-finite.  The functor $f_!$ from
abelian sheaves on $Y\times_sS$ to abelian sheaves on
$Y'\times_sS$ still has a right adjoint $f^!$, given by (1.2.9.1)
(cf. SGA~4, XVIII, 3.1.8).  If $f$ is \'{e}tale, then $f^!=f^*$.

Analogous functors are defined for $Y\times_s\eta$.
\end{filecontents*}
% Original input: components/028_expose_XIII_section_1_3.tex
\begin{filecontents*}[overwrite,noheader]{sga7ii-cumulative-input-028.tex}
\subsection*{1.3. The functor $\Psi$}

\subsubsection*{1.3.1.}
Let $S$ be a henselian trait as in 0.2.5, and let
$p:X\longrightarrow S$.  We propose to define a left exact functor
$\Psi$ from sheaves on $X$ to sheaves on the topos $X_s\times_sS$
constructed in 1.2.  In general, this functor is not the direct-image
functor of a morphism of topoi.

We shall also define
\begin{equation*}
\Psi_\eta:(\text{sheaves on }X_\eta)
\longrightarrow(\text{sheaves on }X_s\times_s\eta).
\end{equation*}

\subsubsection*{1.3.2.}
Put $\bar X=X\times_S\bar S$.  There is a Cartesian diagram
\begin{equation}\tag{1.3.2.1}
\begin{tikzcd}
X_{\bar s} \arrow[r, hook, "\bar i"] \arrow[d]
  & \bar X \arrow[d]
  & X_{\bar\eta} \arrow[l, hook', "\bar j"'] \arrow[d] \\
X_s \arrow[r, hook] & X & X_\eta \arrow[l, hook']
\end{tikzcd}
\end{equation}
over the diagram
\begin{center}
\begin{tikzcd}
\bar s \arrow[r, hook] \arrow[d]
  & \bar S \arrow[d]
  & \bar\eta \arrow[l, hook'] \arrow[d] \\
s \arrow[r, hook] & S & \eta \arrow[l, hook']
\end{tikzcd}
.\end{center}
Let $F$ be a sheaf on $X_\eta$, with inverse image $F_{\bar\eta}$
on $X_{\bar\eta}$.  Put
\begin{equation}\tag{1.3.2.2}
\Psi_\eta(F)=\bar i^*\bar j_*F_{\bar\eta}.
\end{equation}
By transport of structure, this sheaf is endowed with a continuous
action of $\operatorname{Gal}(\bar\eta/\eta)$, compatible with the
action of $\operatorname{Gal}(\bar\eta/\eta)$ on $X_{\bar s}$.  The
functor
\begin{equation}\tag{1.3.2.3}
\Psi_\eta:(X_{\eta,\mathrm{et}})^{\sim}
\longrightarrow(X_s\times_s\eta)^{\sim}
\end{equation}
is left exact.

\subsubsection*{1.3.3.}
Let $F$ be a sheaf on $X$, and let $F_\eta$ and $F_s$ be its
restrictions to $X_\eta$ and $X_s$.  Put
\begin{equation}\tag{1.3.3.1}
(\Psi(F))_\eta=\Psi_\eta(F_\eta)
\end{equation}
and
\begin{equation}\tag{1.3.3.2}
(\Psi(F))_s=F_s.
\end{equation}
Let $\bar F$ be the inverse image of $F$ on $\bar X$, and let
$\varphi'$ be the adjunction morphism
\begin{equation*}
\varphi':\bar F\longrightarrow\bar j_*\bar j^*\bar F.
\end{equation*}
This morphism induces
\begin{equation*}
\varphi:F_{\bar s}=\bar i^*\bar F
\longrightarrow
\bar i^*\bar j_*\bar j^*\bar F
=(\Psi_\eta(F_\eta))_{\bar\eta}.
\end{equation*}
The triple
\begin{equation*}
\Psi(F)=((\Psi(F))_s,(\Psi(F))_\eta,\varphi)
\end{equation*}
is a sheaf on $X_s\times_sS$.  The functor
\begin{equation}\tag{1.3.3.3}
\Psi:(X_{\mathrm{et}})^{\sim}
\longrightarrow(X_s\times_sS)^{\sim}
\end{equation}
is left exact.

\subsubsection*{1.3.4.}
When there is no risk of confusion, we shall write
\begin{equation*}
\Psi_\eta(F)=(\Psi(F))_\eta,
\qquad
\Psi_{\bar\eta}(F)=(\Psi(F))_{\bar\eta},
\qquad\text{and}\qquad
\Psi_{\bar\eta}(F)=(\Psi_\eta(F))_{\bar\eta}.
\end{equation*}

\subsubsection*{1.3.5.}
For the remainder of this section, we list functoriality properties of
$\Psi$ (cf. 1.2.7--1.2.9).  They will be useful chiefly in their
derived form (2.1).

Let $f:X\longrightarrow X'$ be a morphism of $S$-schemes, giving a
diagram
\begin{equation}\tag{1.3.5.1}
\begin{tikzcd}
X_{\bar s} \arrow[r, hook, "\bar i"] \arrow[d, "f"]
  & \bar X \arrow[d, "f"]
  & X_{\bar\eta} \arrow[l, hook', "\bar j"'] \arrow[d, "f"] \\
X'_{\bar s} \arrow[r, hook, "\bar i"]
  & \bar X' & X'_{\bar\eta} \arrow[l, hook', "\bar j"']
\end{tikzcd}
\end{equation}
(the printed diagram inadvertently omits the bar from its top-middle
$\bar X$).  The ambiguity in the notation for the arrows simplifies
the formulas and causes no confusion: only one interpretation makes
sense of the composites below.  Likewise, we write $\Psi$ for the
functor (1.3.3.3) for both $X$ and $X'$.

The morphisms of functors described below for $\Psi$ have variants for
$\Psi_\eta$, which we generally omit to state explicitly.

\subsubsection*{1.3.6.}
The base-change morphism
$\bar i^*f_*\longrightarrow f_*\bar i^*$ induces a morphism of
functors
\begin{equation}\tag{1.3.6.1}
\Psi f_*\longrightarrow f_*\Psi,
\end{equation}
which is an isomorphism when $f$ is proper (SGA~4, XII, 5.1(i)).

The essential case is $X'=S$.  One obtains a morphism
\begin{equation}\tag{1.3.6.2}
f_*\longrightarrow f_*\Psi
\end{equation}
whose essential component is
\begin{equation}\tag{1.3.6.3}
\Gamma(X_{\bar\eta},F)
\longrightarrow
\Gamma(X_{\bar s},\Psi_{\bar\eta}(F));
\end{equation}
this morphism is an isomorphism when $f$ is proper.

\subsubsection*{1.3.7.}
The base-change morphism
$f^*\bar j_*\longrightarrow\bar j_*f^*$ induces a morphism of
functors
\begin{equation}\tag{1.3.7.1}
f^*\Psi\longrightarrow\Psi f^*.
\end{equation}

\subsubsection*{1.3.8.}
Suppose that $f$ is quasi-finite.  The functors $f_!$ commute with base
change; the morphism
$f_!\bar j_*\longrightarrow\bar j_*f_!$ induces a morphism of
functors
\begin{equation}\tag{1.3.8.1}
f_!\Psi\longrightarrow\Psi f_!.
\end{equation}
When $f$ is finite, this is the inverse of (1.3.6.1).

\subsubsection*{1.3.9.}
The situation is dual for $f^!$ (we suppose that $f$ is quasi-finite
when working with abelian sheaves, or that it is an immersion when
working with sheaves of pointed sets).  This functor commutes with
direct images, and the morphism
$\bar i^*f^!\longrightarrow f^!\bar i^*$ induces a morphism of
functors
\begin{equation}\tag{1.3.9.1}
\Psi f^!\longrightarrow f^!\Psi.
\end{equation}
When $f$ is \'{e}tale, this is the inverse of (1.3.7.1).

\subsubsection*{1.3.10.}
Finally, a base change $S'\xrightarrow{f}S$ gives a base-change
morphism
\begin{equation}\tag{1.3.10.1}
f^*\Psi\longrightarrow\Psi f^*.
\end{equation}

\end{filecontents*}
% Original input: components/029_expose_XIII_section_1_4.tex
\begin{filecontents*}[overwrite,noheader]{sga7ii-cumulative-input-029.tex}
\subsection*{1.4. The variation}

\subsubsection*{1.4.1.}
Let $S$ be a henselian trait as in 0.2.5, and let $Y$ be a scheme over
$s$.  Let $\Lambda$ be a ring.  More generally, one could take for
$\Lambda$ a sheaf of rings on $Y$.  For every object $K$ of the
derived category $D(Y\times_sS,\Lambda)$ of the category of sheaves of
$\Lambda$-modules on $Y\times_sS$, we propose to define a
\underline{variation} morphism.

\subsubsection*{1.4.2.}
A complex of sheaves of $\Lambda$-modules $K$ on $Y\times_sS$ is
identified with a triple $(K_s,K_\eta,\varphi)$ in which
\begin{enumerate}[label=\alph*)]
\item $K_s$ (respectively $K_\eta$) is a complex of sheaves
$K_{\bar s}$ (respectively $K_{\bar\eta}$) on $\bar Y$, endowed with a
continuous action of $\operatorname{Gal}(\bar s/s)$ (respectively
$\operatorname{Gal}(\bar\eta/\eta)$);

\item $\varphi$ is an equivariant morphism
$\varphi:K_{\bar s}\longrightarrow K_{\bar\eta}$.
\end{enumerate}
Here and below we retain the scheme $Y$ fixed in 1.4.1; the printed
opening sentence of 1.4.2 inadvertently writes $X\times_sS$.

Such a complex $K$ is always homotopic to a complex
$K'=(K'_s,K'_\eta,\varphi')$ such that $\varphi'$ is injective and the
exact sequence
\begin{equation}\tag{1.4.2.1}
0\longrightarrow K'_{\bar s}
\longrightarrow K'_{\bar\eta}
\longrightarrow\operatorname{coker}(\varphi')
\longrightarrow0
\end{equation}
is split degree by degree.  It is enough to take $K'_s=K_s$ and to
take for $K'_\eta$ the direct sum of $K_\eta$ and the cone of
$\mathrm{sp}^*K_s$.  Put
\begin{equation*}
\Phi(K)=\operatorname{coker}(\varphi').
\end{equation*}
The exact sequence (1.4.2.1) defines a distinguished triangle
\begin{equation}\tag{1.4.2.2}
\begin{tikzcd}[column sep=large, row sep=large]
{} & \Phi(K) \arrow[dl] & {} \\
\mathrm{sp}^*K_s \arrow[rr] & {} & K_\eta \arrow[ul]
\end{tikzcd}
\end{equation}
in the category $K(Y\times_s\eta,\Lambda)$ of complexes of sheaves of
$\Lambda$-modules on $Y\times_s\eta$, up to homotopy.  This triangle
depends only on $K\in\operatorname{Ob}K(Y\times_sS,\Lambda)$.

The construction passes to the derived category and associates with
every object $K$ of $D(Y\times_sS,\Lambda)$ a distinguished triangle
of $D(Y\times_s\eta,\Lambda)$.

Triangle (1.4.2.2) is useful chiefly through the long exact
hypercohomology sequence that it induces:
\begin{equation*}
\cdots\longrightarrow H^i(\bar Y,K_{\bar s})
\longrightarrow H^i(\bar Y,K_{\bar\eta})
\longrightarrow H^i(\bar Y,\Phi(K)_{\bar\eta})
\longrightarrow\cdots.
\end{equation*}

\subsubsection*{1.4.3.}
The inertia group $I$ acts trivially on $K'_{\bar s}$.  Hence, for
every $\sigma\in I$, the endomorphism $\sigma-1$ of
$K'_{\bar\eta}$ (notation (1.4.2.1)) factors through a morphism
\begin{equation*}
\operatorname{Var}(\sigma):
\operatorname{coker}(\varphi')\longrightarrow K'_{\bar\eta}.
\end{equation*}
This construction passes to the derived category: for
$K\in\operatorname{Ob}D(Y\times_sS,\Lambda)$ and $\sigma\in I$, it
defines the \underline{variation}
\begin{equation}\tag{1.4.3.1}
\operatorname{Var}(\sigma):
\Phi(K)_{\bar\eta}\longrightarrow K_{\bar\eta}.
\end{equation}
Writing $q$ for the arrow from $K_\eta$ to $\Phi(K)$, one has
\begin{align}\tag{1.4.3.2}
\sigma&=1+\operatorname{Var}(\sigma)q
&&\text{on }K_{\bar\eta},\\
\tag{1.4.3.3}
\sigma&=1+q\operatorname{Var}(\sigma)
&&\text{on }\Phi(K)_{\bar\eta},
\end{align}
and
\begin{equation}\tag{1.4.3.4}
\operatorname{Var}(\sigma\tau)
=\operatorname{Var}(\sigma)q\operatorname{Var}(\tau)
 +\operatorname{Var}(\sigma)+\operatorname{Var}(\tau).
\end{equation}

\end{filecontents*}
% Original input: components/030_expose_XIII_section_2_1_1_through_2_1_6.tex
\begin{filecontents*}[overwrite,noheader]{sga7ii-cumulative-input-030.tex}
\section*{2. Vanishing cycles}

\subsection*{2.1. The functor $R\Psi$}

\subsubsection*{2.1.1.}
Let $S$ be a henselian trait as in 0.2.5, and let $\Lambda$ be a
torsion ring whose torsion is prime to the residual characteristic
$p$ of $S$ (0.2.7).  More generally, one could take for $\Lambda$ a
torsion sheaf whose torsion is prime to the residual characteristics;
the most important case is
$\Lambda=\mathbb Z/\ell^n$ with $\ell$ an invertible prime in
$\mathcal O_S$.

Let $X$ be a scheme over $S$.  The functor $\Psi$ (respectively
$\Psi_\eta$) carries sheaves of $\Lambda$-modules on $X$
(respectively on $X_\eta$) to sheaves of $\Lambda$-modules on
$X_s\times_sS$ (respectively on $X_s\times_s\eta$).  Let $R\Psi$ and
$R\Psi_\eta$ be the derived functors
\begin{equation*}
R\Psi:D^+(X,\Lambda)\longrightarrow D^+(X_s\times_sS,\Lambda),
\end{equation*}
\begin{equation*}
R\Psi_\eta:D^+(X_\eta,\Lambda)
\longrightarrow D^+(X_s\times_s\eta,\Lambda),
\end{equation*}
and let $R\Phi$ be the composite functor $\Phi\circ R\Psi$:
\begin{equation*}
R\Phi:D^+(X,\Lambda)
\longrightarrow D^+(X_s\times_s\eta,\Lambda).
\end{equation*}

\subsubsection*{2.1.2.}
Consider the commutative diagram (1.3.2.1)
\begin{equation}\tag{1.3.2.1}
\begin{tikzcd}
X_{\bar s} \arrow[r, hook, "\bar i"] \arrow[d]
  & \bar X \arrow[d]
  & X_{\bar\eta} \arrow[l, hook', "\bar j"'] \arrow[d] \\
X_s \arrow[r, hook, "i"] & X & X_\eta \arrow[l, hook', "j"']
\end{tikzcd}
\end{equation}
If $F$ is an injective sheaf of modules on $X$, its restriction to the
open subscheme $X_\eta$ is still injective.  If $F_\eta$ is injective
on $X_\eta$, its inverse image $F_{\bar\eta}$ is acyclic.  Thus the
derived forms of (1.3.3.1) and (1.3.3.2) are
\begin{equation}\tag{2.1.2.1}
R\Psi(K)_s=i^*K
\qquad(K\in\operatorname{Ob}D^+(X,\Lambda)),
\end{equation}
\begin{equation}\tag{2.1.2.2}
R\Psi(K)_\eta=R\Psi_\eta(K_\eta)
\qquad(K\in\operatorname{Ob}D^+(X,\Lambda)),
\end{equation}
and
\begin{equation}\tag{2.1.2.3}
R\Psi_\eta(K)_{\bar\eta}
=\bar i^*R\bar j_*K_{\bar\eta}
\qquad(K\in\operatorname{Ob}D^+(X_\eta,\Lambda)).
\end{equation}
This allows us to write, without conflict,
\begin{equation*}
R\Psi_\eta(K)=(R\Psi(K))_\eta,
\qquad
R\Psi_{\bar\eta}(K)=(R\Psi(K))_{\bar\eta},
\qquad\text{and}\qquad
R\Psi_{\bar\eta}(K)=(R\Psi_\eta(K))_{\bar\eta}.
\end{equation*}
With this notation, the distinguished triangle (1.4.2.2) becomes
\begin{equation}\tag{2.1.2.4}
\begin{tikzcd}[row sep=large, column sep=huge]
{} & R\Phi(K) \arrow[dl] & {} \\
\mathrm{sp}^*i^*K \arrow[rr] & {}
  & R\Psi_\eta(K_\eta) \arrow[ul]
\end{tikzcd}
.\end{equation}
The printed bottom-right term has $K_s$; the defining identity
(2.1.2.2) requires $K_\eta$, as written here.

For $\sigma\in I$, there is a \underline{variation} morphism
\begin{equation}\tag{2.1.2.5}
\operatorname{Var}(\sigma):R\Phi(K)_{\bar\eta}
\longrightarrow R\Psi_{\bar\eta}(K).
\end{equation}
The sheaves $R^i\Phi(\Lambda)$, or more generally $R^i\Phi(K)$, are
the \underline{sheaves of vanishing cycles}.

\subsubsection*{2.1.3.}
Let $(\bar x,\bar\eta)$ be the geometric point of
$X_s\times_s\eta$ defined by a geometric point $\bar x$ of
$X_{\bar s}$ (1.2.5).  We compute the fiber of
$R\Psi_\eta(K_\eta)$ at $(\bar x,\bar\eta)$ using (2.1.2.3)
(cf. I, 2.3).  The strict henselization $X_{(\bar x)}$ of $X$ at
$\bar x$ is a scheme over the strict henselization $S_{\mathrm{nr}}$
of $S$, with geometric generic point $\eta_{\mathrm{nr}}$.

\subsubsection*{Proposition 2.1.4.}
\underline{One has}
\begin{equation*}
(R\Psi(K_\eta))_{(\bar x,\bar\eta)}
=R\Gamma\!\left(
X_{(\bar x)}\times_{\eta_{\mathrm{nr}}}\bar\eta,K
\right);
\end{equation*}
\underline{in particular,}
\begin{equation*}
R^i\Psi(K_\eta)_{(\bar x,\bar\eta)}
=\mathbb H^i\!\left(
X_{(\bar x)}\times_{\eta_{\mathrm{nr}}}\bar\eta,K
\right).
\end{equation*}

Apply SGA~4, XV, 2.1 (the local-acyclicity theorem for smooth
morphisms).  This gives the following restatement.

\subsubsection*{Restatement 2.1.5.}
\underline{If $f$ is smooth and $F$ is a locally constant sheaf, then}
\begin{equation*}
R\Phi(F)=0.
\end{equation*}
More generally, if $K\in\operatorname{Ob}D^+(X,\Lambda)$, then
$R\Phi(K)=0$ wherever $f$ is smooth and the sheaves
$\mathcal H^i(K)$ are locally constant.  See also 2.1.11.

\subsubsection*{2.1.6.}
Let $f:Y\longrightarrow Y'$ be a morphism of schemes over $s$.
The functors $f^*$, $f_*$, $f^!$, and $f_!$ (1.2.7--1.2.9) have
natural analogues in the derived category.
\begin{enumerate}[label=\alph*)]
\item The functor $f^*$ is exact and therefore derives trivially.

\item The functor $f_*$ has a right derived functor $Rf_*$.  When $f$
is quasi-compact, one has
\begin{equation}\tag{2.1.6.1}
(Rf_*K)_{\bar s}=Rf_{\bar s*}(K_{\bar s})
\end{equation}
and
\begin{equation}\tag{2.1.6.2}
(Rf_*K)_{\bar\eta}=Rf_{\bar s*}(K_{\bar\eta}).
\end{equation}

\item When $f$ is quasi-finite, the functor $f_!$ is exact and
therefore derives trivially.  In general, the derived functor of $f_!$
is pathological, and one proceeds as in SGA~4, XVII.  Suppose that $f$
is separated and of finite type and that $Y'$ is noetherian.  There is
then [2] a factorization $f=\widetilde f\,j$, where $j$ is an open
immersion and $\widetilde f$ is proper, and one defines
\begin{equation*}
Rf_!=R\widetilde f_*\circ j_!.
\end{equation*}
The arguments of SGA~4, XVII show that $Rf_!$ is independent of the
choice of the factorization $\widetilde f\,j$.  The analogues of
(2.1.6.1) and (2.1.6.2) hold.

\item We shall use $Rf^!$ only for a quasi-finite morphism, in which
case it is the derived functor of $f^!$.  Again, the analogues of
(2.1.6.1) and (2.1.6.2) hold.
\end{enumerate}

\end{filecontents*}
% Original input: components/031_expose_XIII_section_2_1_7_through_2_1_8.tex
\begin{filecontents*}[overwrite,noheader]{sga7ii-cumulative-input-031.tex}
\subsubsection*{2.1.7.}
We now review the functoriality properties of $R\Psi$, that is, we
derive 1.3.5--1.3.10.  Let $f:X\longrightarrow X'$ be a morphism of
$S$-schemes, and retain the notation of 1.3.5.

The morphism (1.3.6.1) defines
\begin{equation}\tag{2.1.7.1}
R\Psi\,Rf_*\longrightarrow Rf_*\,R\Psi.
\end{equation}
By the proper base-change theorem (SGA~4, XII, 5.1),
\underline{this morphism is an isomorphism if $f$ is proper.}

The morphism (1.3.7.1) defines
\begin{equation}\tag{2.1.7.2}
f^*R\Psi\longrightarrow R\Psi f^*.
\end{equation}
By the smooth base-change theorem (SGA~4, XVI, 1.2),
\underline{this morphism is an isomorphism if $f$ is smooth.}

The isomorphism (2.1.7.1) for $f$ proper and the morphism (1.3.8.1)
define a morphism
\begin{equation}\tag{2.1.7.3}
Rf_!R\Psi\longrightarrow R\Psi Rf_!,
\end{equation}
which, when $f$ is proper, is the inverse of (2.1.7.1).  The print
mistakenly refers here to the later special case (2.1.8.1).

When $f$ is quasi-finite, (1.3.9.1) derives to
\begin{equation}\tag{2.1.7.4}
R\Psi Rf^!\longrightarrow Rf^!R\Psi,
\end{equation}
which, when $f$ is \'{e}tale, is the inverse of (2.1.7.2).  The printed
reference to (2.1.7.1) is incompatible with the two functors involved.

Let $f:S'\longrightarrow S$ be a change of trait, and put
$X'=X\times_SS'$:
\begin{center}
\begin{tikzcd}
X' \arrow[r, "f"] \arrow[d, "p"'] & X \arrow[d, "p"] \\
S' \arrow[r, "f"] & S
\end{tikzcd}
\end{center}
(for the ambiguity in the notation, see 1.3.5).  Suppose that $f$ is
surjective.  The base-change arrow (1.3.10.1) derives to
\begin{equation}\tag{2.1.7.5}
f^*R\Psi\longrightarrow R\Psi f^*.
\end{equation}
The same statements hold for $R\Psi_\eta$.

\subsubsection*{2.1.8.}
The essential case of (2.1.7.1) and (2.1.7.3) is $X'=S$.  For
$K\in\operatorname{Ob}D^+(X_{\bar\eta},\Lambda)$ one then obtains
morphisms
\begin{equation}\tag{2.1.8.1}
Rf_*\longrightarrow Rf_*R\Psi
\end{equation}
and
\begin{equation}\tag{2.1.8.2}
Rf_!R\Psi\longrightarrow Rf_!,
\end{equation}
hence
\begin{equation}\tag{2.1.8.3}
\mathbb H^i(X_{\bar\eta},K)
\longrightarrow
\mathbb H^i(X_{\bar s},R\Psi_{\bar\eta}(K))
\end{equation}
and
\begin{equation}\tag{2.1.8.4}
\mathbb H_c^i(X_{\bar s},R\Psi_{\bar\eta}(K))
\longrightarrow
\mathbb H_c^i(X_{\bar\eta},K),
\end{equation}
and the commutativity of the diagrams
\begin{equation}\tag{2.1.8.5}
\begin{tikzcd}[column sep=large]
(R^if_*K)_{\bar s} \arrow[r, "\mathrm{sp}"] \arrow[d]
  & \mathbb H^i(X_{\bar\eta},K) \arrow[d] \\
\mathbb H^i(X_{\bar s},K) \arrow[r]
  & \mathbb H^i(X_{\bar s},R\Psi_{\bar\eta}(K))
\end{tikzcd}
\end{equation}
and
\begin{equation}\tag{2.1.8.6}
\begin{tikzcd}[column sep=large]
\mathbb H_c^i(X_{\bar s},K) \arrow[rr, "\mathrm{sp}"] \arrow[dr]
  & {} & \mathbb H_c^i(X_{\bar\eta},K) \\
{} & \mathbb H_c^i(X_{\bar s},R\Psi_{\bar\eta}(K)) \arrow[ur] & {}
\end{tikzcd}
.\end{equation}

For $K\in\operatorname{Ob}D^+(X_\eta,\Lambda)$ and $\sigma\in I$,
the following diagrams are commutative:
\begin{equation}\tag{2.1.8.7}
\begin{tikzcd}[column sep=small, cells={nodes={font=\small}}]
\mathbb H^i(X_{\bar\eta},K)
  \arrow[rr, dash, "\sigma-1"] \arrow[d]
  & {} & \mathbb H^i(X_{\bar\eta},K) \arrow[d] \\
\mathbb H^i(X_{\bar s},R\Psi_{\bar\eta}(K)) \arrow[r]
  & \mathbb H^i(X_{\bar s},R\Phi_{\bar\eta}(K))
      \arrow[r, "\operatorname{Var}(\sigma)"]
  & \mathbb H^i(X_{\bar s},R\Psi_{\bar\eta}(K))
\end{tikzcd}
\end{equation}
and
\begin{equation}\tag{2.1.8.8}
\begin{tikzcd}[column sep=small, cells={nodes={font=\small}}]
\mathbb H_c^i(X_{\bar\eta},K)
  \arrow[rr, "\sigma-1"]
  & {} & \mathbb H_c^i(X_{\bar\eta},K) \\
\mathbb H_c^i(X_{\bar s},R\Psi_{\bar\eta}(K)) \arrow[r] \arrow[u]
  & \mathbb H_c^i(X_{\bar s},R\Phi_{\bar\eta}(K))
      \arrow[r, "\operatorname{Var}(\sigma)"]
  & \mathbb H_c^i(X_{\bar s},R\Psi_{\bar\eta}(K)) \arrow[u]
\end{tikzcd}
.\end{equation}
In (2.1.8.8) the English restores the bar omitted from the printed
bottom-left $X_{\bar s}$ and the missing closing parenthesis in the
bottom-middle term.  In (2.1.8.7), however, the upper connector really
is printed as a rule without an arrowhead, whereas (2.1.8.8) has an
arrow; this source-visible distinction is retained.

When $f$ is proper, (2.1.8.1) and (2.1.8.2), as well as (2.1.8.3) and
(2.1.8.4), are mutually inverse isomorphisms.  The distinguished
triangle (2.1.2.4) gives a long exact sequence
\begin{equation}\tag{2.1.8.9}
\cdots\longrightarrow\mathbb H^i(X_{\bar s},K)
\longrightarrow\mathbb H^i(X_{\bar\eta},K)
\longrightarrow\mathbb H^i(X_{\bar s},R\Phi_{\bar\eta}(K))
\longrightarrow\cdots.
\end{equation}
Thus the vanishing complex $R\Phi(K)$ expresses the difference between
the cohomologies of the geometric fibers $X_{\bar s}$ and
$X_{\bar\eta}$, with coefficients in the inverse image of $K$.
By (2.1.8.7), the variation $\operatorname{Var}(\sigma)$ determines
the action of the inertia group $I$ on
$\mathbb H^i(X_{\bar\eta},K)$.  These facts underlie all applications
of the method of vanishing cycles.

\end{filecontents*}
% Original input: components/032_expose_XIII_section_2_1_9_through_2_1_13.tex
\begin{filecontents*}[overwrite,noheader]{sga7ii-cumulative-input-032.tex}
The properness hypothesis imposed above on $f$ may be replaced by
regularity-at-infinity conditions on $f$ and $K$.  For example, one has
the following result (for
$K\in\operatorname{Ob}D(X_{\bar\eta},\Lambda)$).

\subsubsection*{Proposition 2.1.9.}
\underline{Suppose that $f$ admits a factorization $f=f_1k$}
\begin{center}
\begin{tikzcd}[column sep=huge]
X \arrow[r, hook, "k"] & X_1 \arrow[r, "f_1"] & S,
\end{tikzcd}
\end{center}
\underline{where $f_1$ is proper and $k$ identifies $X$ with the
complement in $X_1$ of a relative normal-crossings divisor $D$}
(thus, by hypothesis, $f_1$ is smooth on a neighborhood of $D$).
\underline{If, on a neighborhood $U$ of $D$, the sheaves
$\mathcal H^i(K)$ are locally constant on $U-(D\cup X_0)$ and tamely
ramified along the normal-crossings divisor $D\cup X_0$, then
(2.1.8.1) and (2.1.8.2) are isomorphisms.}

The hypercohomology spectral sequence for $K$ shows that it is enough
to treat the case in which $K$ consists of a sheaf $F_\eta$ in degree
zero.  After replacing $S$ by its normalization in a finite extension
(which is permissible, as one sees directly from the definitions) and
applying Abhyankar's lemma, we may suppose that $F_\eta$ is the
restriction to $X_\eta$ of a sheaf $F$ on $X$, locally constant on
$U\cap X$ and tamely ramified along $D$.  With the notation of 2.1.2,
let $\overline F$ be the inverse image of $F$ on $\overline X$, and
write $\bar f:\overline X\longrightarrow\overline S$.

To prove that (2.1.8.1) is an isomorphism, it is enough to prove that
\[
R\bar f_*\bigl(R\bar j_*\bar j^*\overline F\bigr)
\]
commutes with base change by $\bar s\longrightarrow\overline S$.  In
the distinguished triangle
\[
\longrightarrow\overline F\longrightarrow
R\bar j_*\bar j^*\overline F\longrightarrow\Delta\longrightarrow,
\]
the sheaves $\mathcal H^i(\Delta)$ are supported on $X_{\bar s}$; for
$\Delta$, the base-change statement in question is tautological.  It
therefore remains to prove that $R\bar f_*(\overline F)$ commutes with
base change by $\bar s\longrightarrow\overline S$, or, in view of the
base-change theorem for the proper morphism $f_1$, that
$Rk_*(\overline F)$ commutes with base change
$X_{1\bar s}\longrightarrow X_1$.  This is Lemma 2.1.10 below.

To prove that (2.1.8.2) is an isomorphism, it is enough to prove, for
the sheaf $F$, that extension by zero $k_!$ commutes with the functor
$R\bar j_*$; that is, on a neighborhood of $D$,
\[
R\Phi(k_!F)=0\qquad\hbox{on }X_{1\bar s}.
\]
This is Lemma 2.1.11 below.

\subsubsection*{Lemma 2.1.10.}
\underline{Let} $f:X\longrightarrow S$ \underline{be smooth, let}
$D\subset X$ \underline{be a relative normal-crossings divisor, put}
$U=X-D\mathrel{\mathop{\hookrightarrow}^{j}}X$, \underline{and let}
$F$ \underline{be a sheaf of $\Lambda$-modules on $U$ which is locally
constant and tamely ramified along} $D$.  \underline{Formation of}
$Rj_*F$ \underline{commutes with arbitrary base change}
$T\longrightarrow S$.

The question is local on $X$ for the \'{e}tale topology.  We may
therefore suppose that there is a morphism $g:X'\longrightarrow X$,
where $X'$ is the normalization of $X$ in a finite \'{e}tale
surjective covering $g':U'\longrightarrow U$, tamely ramified along
$D$, such that $g'^*F$ is constant on $U'$.  By Abhyankar's lemma, we
may further suppose that $X'$ is smooth over $S$ and that $U'$ is the
complement in $X'$ of a relative normal-crossings divisor $D'$.

The sheaf $F$ injects into $g'_*g'^*F$, and the cokernel again
satisfies the hypotheses of the lemma.  Iterating this construction
(the same $g$ continues to work) gives a resolution of $F$ by sheaves
of the form $g'_*G$, where $G$ is constant and $g'$ is as above.  It
therefore suffices to prove the assertion for such a sheaf, or,
equivalently, for $(X'/S,D',G)$ with $G$ constant on $U'$.  This is a
special case of SGA~5, II, 4.4.

\subsubsection*{Lemma 2.1.11.}
\underline{Let $S$ be a henselian trait, let $f:X\longrightarrow S$
be smooth, let $D\subset X$ be a relative normal-crossings divisor,
put $U=X-D\mathrel{\mathop{\hookrightarrow}^{j}}X$, and let $F$ be a
sheaf of $\Lambda$-modules on $U$ which is locally constant and tamely
ramified along $D$.  Then}
\[
\underline{R\Phi(j_!F)=0.}
\]

The d\'{e}vissage just used reduces us to the case in which $F$ is
constant, with value $G$.  We may suppose that $D$ is the union of
smooth divisors $(D_i)_{i\in I}$.  For $J\subset I$, let $D_J$ be the
intersection of the $D_i$ for $i\in J$.  The sheaf $F$ has a
resolution $F^\bullet$ whose terms are sums of sheaves $G_J$, constant
on $D_J$ and extended by zero to $X$.  It is therefore enough to prove
the lemma for one such sheaf, and this follows from (2.1.5), applied to
$D_J/S$.

I conjecture that (2.1.7.5) is an isomorphism when $S$ and $S'$ are
excellent.  Here is a partial result (see also 2.4.2).

\subsubsection*{Proposition 2.1.12.}
\underline{The morphism (2.1.7.5) is an isomorphism in the following
cases:}

\medskip
\noindent(a) $S$ is excellent and $S'$ is its completion;

\noindent(b) $S$ and $S'$ have equal characteristic zero.

In case (b), we reduce to the situation in which $S$ and $S'$ are
strictly local and a uniformizer of $S$ is still a uniformizer of
$S'$ (extend the scalars without changing
$\bar S'\longrightarrow\bar S$).  In both cases (a) and (b), writing
an overbar for normalization in a separable closure of the fraction
field, one then has
\[
\bar S'\simeq S'\times_S\bar S.
\]
Returning to the definitions, it is enough to prove that
$S'\longrightarrow S$ is universally locally acyclic (apply SGA~4,
XVI, 1.1).  This follows from I, 0.5.1 (based on Artin's version of
N\'{e}ron desingularization) and SGA~4, XV, 2.1 (the local-acyclicity
theorem for smooth morphisms).

\subsubsection*{2.1.13.}
Suppose that $f:X\longrightarrow S$ is of finite type.  The
cohomological dimension (SGA~4, XVII, 1.2.9) of $R\Psi$ is then at most
$\dim(X_\eta)$, as follows from 2.1.4 and SGA~4, XIV, 3.2 (cf. I,
4.2(i)).  Thus $R\Psi$ extends to the entire derived category,
\[
R\Psi:D(X,\Lambda)\longrightarrow D(X_s\times_sS,\Lambda),
\]
and it carries $D^-$ into $D^-$.  Let
$L\in\operatorname{Ob}D^-(\Lambda)$ be a bounded-above complex of
right $\Lambda$-modules, and let
$K\in\operatorname{Ob}D^-(X,\Lambda)$.  There is then an isomorphism
\begin{equation}\tag{2.1.13.1}
L\overset{\mathbf L}{\otimes}_\Lambda R\Psi(K)
\xrightarrow{\ \sim\ }
R\Psi\bigl(L\overset{\mathbf L}{\otimes}_\Lambda K\bigr).
\end{equation}
To verify (2.1.13.1), first replace $L$ by a complex of free
$\Lambda$-modules.  Since $R\Psi$ is triangulated and has finite
cohomological dimension, it is enough to verify the assertion after
replacing $L$ by one of its terms.  Passage to the filtered colimit
reduces us to the case in which $L$ is a finite free module, and then
to the tautological case $L=\Lambda$.

It follows from (2.1.13.1), by the usual argument (cf. SGA~4, XVII,
5.2.11), that if $K\in\operatorname{Ob}D(X_\eta,\Lambda)$ has finite
Tor-dimension at most $d$, where $d\in\mathbb Z$ (SGA~4, XVII, 4.1.9),
then $R\Psi_\eta(K)$ likewise has Tor-dimension at most $d$.
\end{filecontents*}
% Original input: components/033_expose_XIII_section_2_2.tex
\begin{filecontents*}[overwrite,noheader]{sga7ii-cumulative-input-033.tex}
\subsection*{2.2. A compatibility (trace morphisms)}

\subsubsection*{2.2.1.}
Let $S$ be a henselian trait, and let
$u:X_s\longrightarrow Y_s$ be a separated, quasi-finite, flat
morphism of finite presentation between $s$-schemes.  For every
abelian sheaf $F$ on $Y_s$, there is a trace morphism
\[
\operatorname{Tr}_u:u_!u^*F\longrightarrow F
\]
(SGA~4, XVII, 6.2.3).  If $F$ is an abelian sheaf on
$Y_s\times_sS$ or on $Y_s\times_s\eta$, this gives an analogous
morphism $\operatorname{Tr}_u:u_!u^*F\longrightarrow F$.

Now let $u:X\longrightarrow Y$ be a separated, quasi-finite, flat
morphism of finite presentation between $S$-schemes.  We again write
$u$ for the induced morphism
\[
X_s\times_sS\longrightarrow Y_s\times_sS,
\]
and we use $\Psi$ for the functors of 1.3.3, both for $X$ and for $Y$.

\subsubsection*{Proposition 2.2.2.}
\underline{For every abelian sheaf $F$ on $Y$, the following diagram
is commutative:}
\begin{center}
\begin{tikzcd}[column sep=large, row sep=large]
u_!u^*\Psi F
  \arrow[r, "{(1.3.7.1)}"] \arrow[d, "\operatorname{Tr}_u"']
  & u_!\Psi u^*F \arrow[r, "{(1.3.8.1)}"]
  & \Psi u_!u^*F \arrow[d, "\operatorname{Tr}_u"] \\
\Psi F \arrow[rr, equal] & {} & \Psi F.
\end{tikzcd}
\end{center}
The verification is left to the reader.  There is a variant for
$\Psi_\eta$ and for $F$ on $Y_\eta$.

\subsubsection*{2.2.3.}
Let $\Lambda$ be as in 2.1.1, and let $u:X\longrightarrow Y$ be a
quasi-finite flat morphism between separated, finite-type, flat
$S$-schemes:
\begin{center}
\begin{tikzcd}[column sep=large, row sep=large]
X \arrow[rr, "u"] \arrow[dr, "f"'] & {} & Y \arrow[dl, "g"] \\
{} & S & {}
\end{tikzcd}
\end{center}
The print refers here to 2.2.1, but the coefficient ring $\Lambda$ is
fixed in 2.1.1.  We again write $f$ and $g$ for the morphisms
\[
X_s\times_sS\longrightarrow S,
\qquad
Y_s\times_sS\longrightarrow S,
\]
and suppose that $\dim Y_\eta\leq n$.  For every
$K\in\operatorname{Ob}D(S,\Lambda)$, there is a trace morphism
(SGA~4, XVIII, 2.9)
\[
\operatorname{Tr}_f:Rf_!f^*K(n)[2n]\longrightarrow K.
\]

\subsubsection*{Corollary 2.2.4.}
\underline{The following diagram is commutative:}
\begin{center}
\resizebox{\textwidth}{!}{%
\begin{tikzcd}[ampersand replacement=\&, column sep=large, row sep=large]
Rg_!u_!u^*R\Psi\,g^*K(n)[2n]
  \arrow[r, "{(1.3.7.1)}"] \arrow[d, "\operatorname{Tr}_u"']
  \& Rf_!R\Psi\,f^*K(n)[2n]
      \arrow[r, "{(2.1.10.1)}"]
  \& Rf_!f^*K(n)[2n] \arrow[d, "\operatorname{Tr}_f"] \\
Rg_!R\Psi\,g^*K(n)[2n]
  \arrow[r, "{(2.1.10.1)}"]
  \& Rg_!g^*K(n)[2n] \arrow[r, "\operatorname{Tr}_g"]
  \& K
\end{tikzcd}}
\end{center}
\underline{If $u$ is \'{e}tale, (1.3.7.1) is an isomorphism.}

With the arguments suppressed, the diagram (2.2.4) expands as
follows:
\begin{center}
\resizebox{\textwidth}{!}{%
\begin{tikzcd}[ampersand replacement=\&, column sep=large, row sep=huge]
Rg_!u_!u^*R\Psi\,g^*
  \arrow[r] \arrow[d, "\operatorname{Tr}_u"']
  \arrow[rd, phantom, "{(1)}" description]
  \& Rg_!R\Psi\,u_!u^*g^*
      \arrow[r] \arrow[d, "\operatorname{Tr}_u"']
      \arrow[rd, phantom, "{(2)}" description]
  \& Rg_!u_!u^*g^*
      \arrow[r, "\sim"] \arrow[d, "\operatorname{Tr}_u"']
      \arrow[rd, phantom, "{(3)}" description]
  \& Rf_!f^* \arrow[d, "\operatorname{Tr}_f"] \\
Rg_!R\Psi\,g^* \arrow[r, equal]
  \& Rg_!R\Psi\,g^* \arrow[r]
  \& Rg_!g^* \arrow[r, "\operatorname{Tr}_g"]
  \& K.
\end{tikzcd}}
\end{center}
The commutativity of (1) follows from 2.2.2; that of (2) is the
functoriality of (2.1.10.2); and that of (3) is the transitivity of the
trace (SGA~4, XVIII, 2.9, Variant 3).

\subsubsection*{2.2.4.1.}
There is an analogous result for $\Psi_\eta$ and
$K\in\operatorname{Ob}D(\eta,\Lambda)$.

\subsubsection*{2.2.5.}
Let $\Phi$ be a subset of $Y_s$ which is proper over $s$, and write
$Rg_\Phi$ for the ``direct image with support in $\Phi$'' functor.
Let $\operatorname{Tr}_g$ denote the composite
\begin{equation}\tag{2.2.5.1}
\begin{split}
\operatorname{Tr}_g:
Rg_\Phi R\Psi g^*K(n)[2n]
&\longrightarrow Rg_!R\Psi g^*K(n)[2n] \\
&\longrightarrow Rg_!g^*K(n)[2n]
\longrightarrow K,
\end{split}
\end{equation}
and similarly for $\Psi_\eta$ and $K\in D(\eta,\Lambda)$.  The print
has $Rf_\Phi$ as the first functor in (2.2.5.1); the immediately
preceding definition and the remainder of the composite require
$Rg_\Phi$, which is used here.

If $\Phi$ is a finite set of closed points, then
\[
Rg_\Phi\circ u_!=Rf_{u^{-1}\Phi},
\]
and hence there is a trace morphism
\[
\operatorname{Tr}_u:
Rf_{u^{-1}\Phi}u^*=Rg_\Phi u_!u^*\longrightarrow Rg_\Phi.
\]
By 2.2.4, the following diagram is commutative:
\begin{center}
\resizebox{\textwidth}{!}{%
\begin{tikzcd}[ampersand replacement=\&, column sep=large, row sep=large]
Rf_{u^{-1}\Phi}u^*R\Psi\,g^*K(n)[2n]
  \arrow[r] \arrow[d, "\operatorname{Tr}_u"']
  \& Rf_{u^{-1}\Phi}R\Psi\,f^*K(n)[2n]
      \arrow[r, "\operatorname{Tr}_f"]
  \& K \arrow[d, equal] \\
Rg_\Phi R\Psi\,g^*K(n)[2n]
  \arrow[rr, "\operatorname{Tr}_g"']
  \& {} \& K.
\end{tikzcd}}
\end{center}

Suppose that $u$ is \'{e}tale, let $x\in X_s$, and suppose that $x$
and $y=u(x)$ have the same residue field.  Then
\[
Rf_{\{x\}}\xrightarrow{\ \sim\ }Rg_{\{y\}}.
\]
The print has $\{x\}$ also on the right, but the support for $g$ lies
in $Y_s$ and is $\{y\}$, as the following diagram confirms.  This
isomorphism, given by $\operatorname{Tr}_u$, expresses the \'{e}tale
locality of cohomology with support at $x$.  The diagram
\begin{equation}\tag{2.2.5.2}
\begin{tikzcd}[column sep=huge, row sep=large]
Rf_{\{x\}}R\Psi\,f^*K(n)[2n]
  \arrow[r, "\operatorname{Tr}_f"] \arrow[d, equal]
  & K \arrow[d, equal] \\
Rg_{\{y\}}R\Psi\,g^*K(n)[2n]
  \arrow[r, "\operatorname{Tr}_g"]
  & K
\end{tikzcd}
\end{equation}
is commutative.

\subsubsection*{Scholium 2.2.6.}
The compatibility (2.2.5.2) says that the trace morphism (2.2.5.1),
for $\Phi=\{y\}$, is a local invariant: it depends only on the
henselization of the $S$-scheme $X$ at $x$.  The same holds for
$R\Psi_\eta$.  The print's references to (2.3.5.2) and (2.3.5.1) are
the immediately preceding (2.2.5.2) and (2.2.5.1); likewise,
$\Phi\subset Y_s$ makes $\{y\}$, rather than the printed $\{x\}$,
the relevant support.
\end{filecontents*}
% Original input: components/034_expose_XIII_section_2_3.tex
\begin{filecontents*}[overwrite,noheader]{sga7ii-cumulative-input-034.tex}
\subsection*{2.3. Finiteness theorem (in equal characteristic zero)}

\subsubsection*{Theorem 2.3.1.}
\underline{Let $S$ be a henselian trait as in 0.2.5, let $\Lambda$ be
a noetherian torsion ring, and let $f:X\longrightarrow S$ be a
morphism of finite type.  Suppose that $S$ has equal characteristic
zero.  Then, for every constructible sheaf of $\Lambda$-modules $F$
on $X_\eta$, the sheaves $R^i\Psi_{\bar\eta}(F)$ on $X_{\bar s}$ are
constructible.}

We begin by repeating the opening of the proof of SGA~4, XIV, 1.10.
Apply SGA~4, IX, 2.14(ii).  There is a resolution $F^\bullet$ of $F$
\begin{center}
\begin{tikzcd}[column sep=large, row sep=large]
{} & F \arrow[d] & {} & {} \\
0 \arrow[r] & F^0 \arrow[r] & F^1 \arrow[r] & \cdots
\end{tikzcd}
\end{center}
by sheaves of the form $\bigoplus_i p'_{i*}G_i$, where
$p'_i:Y_{i\eta}\longrightarrow X_\eta$ is finite and $G_i$ is a
constructible constant sheaf.  We may suppose that $Y_{i\eta}$ is
reduced; let $p_i:Y_i\longrightarrow X$ be the normalization of $X$
in $Y_{i\eta}$.  Using the spectral sequence of the resolution
$F^\bullet$ and applying (2.1.7.1) in the elementary case of the
finite morphism $p_i$, we see that it suffices to prove the theorem
for each pair $(Y_i,G_i)$.

We next prove the theorem when $F$ is the constant sheaf defined by a
$\Lambda$-module $M$.  Recall that $\Lambda$ is a
$\mathbb Z/n$-algebra for a suitable $n$.  The module $M$ may
therefore be d\'{e}viss\'{e} into modules $M_i$ over
$\Lambda/\ell_i\Lambda$, where each $\ell_i\mid n$ is a suitable
prime.  This reduces us to the case in which $\Lambda$ is a
$\mathbb Z/\ell$-algebra, with $\ell$ prime.  Formula (2.1.13.1) then
gives
\[
M\otimes_{\mathbb Z/\ell}
R^i\Psi_{\bar\eta}(\mathbb Z/\ell)
\xrightarrow{\ \sim\ }
R^i\Psi_{\bar\eta}(M).
\]
It remains to treat the constant sheaf $F=\mathbb Z/\ell$.

Choose a descending sequence of closed subsets
\[
X_\eta=X'_0\supset\cdots\supset X'_n
\]
such that $(X'_i-X'_{i+1})_{\mathrm{red}}$ is smooth over $\eta$, and
let $X_i=\overline{X'_i}$ be the closure in $X$.  The constant sheaf
$\mathbb Z/\ell$ is d\'{e}viss\'{e} into the constant sheaves
$\mathbb Z/\ell$ on the strata $X'_i-X'_{i+1}$, extended by zero.
Thus:

\medskip
\noindent\underline{Reduction 1.}
It is enough to prove Theorem 2.3.1 for a sheaf
$j_!\mathbb Z/\ell$, where
$j:U\hookrightarrow X_\eta$ is the inclusion of a smooth open subset.

By resolution of singularities there is a morphism
$g:X_1\longrightarrow X$ such that

\medskip
\noindent(a) $g$ is proper and induces an isomorphism
$g^{-1}(U)\xrightarrow{\sim}U$;

\noindent(b) $X_1$ is regular and $X_1-g^{-1}(U)$ is a
normal-crossings divisor in $X_1$, which is a union of regular
divisors $D_\alpha$.

Let $j_1:g^{-1}(U)\hookrightarrow X_1$.  Then
\[
Rg_*j_{1!}\mathbb Z/\ell=j_!\mathbb Z/\ell.
\]
Applying (2.1.7.1) and the finiteness theorem SGA~4, XIV, 1.1, we see
that it suffices to prove Theorem 2.3.1 for
$(X_1,j_{1!}\mathbb Z/\ell)$.

There is a resolution
\[
0\longrightarrow j_{1!}\mathbb Z/\ell
\longrightarrow\mathbb Z/\ell
\longrightarrow\bigoplus_\alpha(\mathbb Z/\ell)_{D_\alpha}
\longrightarrow
\bigoplus_{\alpha,\beta}
(\mathbb Z/\ell)_{D_\alpha\cap D_\beta}
\longrightarrow\cdots,
\]
and hence:

\medskip
\noindent\underline{Reduction 2.}
It is enough to prove Theorem 2.3.1 for the constant sheaf
$\mathbb Z/\ell$ on a regular $S$-scheme $X$ such that $X_s$ is a
normal-crossings divisor in $X$.

In this final case one uses the explicit computation I, 3.3, which in
turn uses the purity theorem SGA~4, XIX, 3.2 and the characteristic-zero
hypothesis.
\end{filecontents*}
% Original input: components/035_expose_XIII_section_2_4_1_through_2_4_4.tex
\begin{filecontents*}[overwrite,noheader]{sga7ii-cumulative-input-035.tex}
\subsection*{2.4. Isolated singularities}

\subsubsection*{2.4.1.}
Throughout this section, for simplicity, we suppose that $S$ is a
strictly henselian trait.  Let $\Lambda$ be a noetherian torsion ring
whose torsion is prime to the residual characteristic of $S$, let
$f:X\longrightarrow S$ be a morphism of finite type, and let $F$ be a
constructible sheaf of $\Lambda$-modules on $X$.  A point $x\in X$ is
called a smooth point of $(X,F)$ if $f$ is smooth at $x$ and $F$ is
locally constant in a neighborhood of $x$.  Let $\Sigma$ denote the
set of points where $(X,F)$ is not smooth, and let $\Sigma_s$ be its
special fiber.  We know that $R\Phi(F)=0$ outside $\Sigma_s$.

\subsubsection*{Theorem 2.4.2.}
\underline{At every isolated point $x$ of $\Sigma_s$, the modules
$R^i\Phi_{\bar\eta}(F)_x$ are $\Lambda$-modules of finite type and
are invariant under change of trait (via 2.1.7.1).}

In equal characteristic zero, these assertions were proved in
2.1.12 and 2.3.1 without assuming that $x$ is isolated.

Let $x$ be an isolated point of $\Sigma_s$.  The theorem is local in
a neighborhood of $x$, so we may suppose that $f$ is proper.  We then
have the long exact sequence (2.1.8.9)
\[
\cdots\longrightarrow H^i(X_s,F)
\longrightarrow H^i(X_{\bar\eta},F)
\longrightarrow
\mathbb H^i\bigl(X_s,R\Phi_{\bar\eta}(F)\bigr)
\longrightarrow\cdots .
\]
In this sequence, $H^i(X_s,F)$ and $H^i(X_{\bar\eta},F)$ are
invariant under change of trait and are $\Lambda$-modules of finite
type (SGA~4, XIV, 1.1).  The modules
$\mathbb H^i(X_s,R\Phi_{\bar\eta}(F))$ therefore have the same
properties.  Recall the following lemma.

\subsubsection*{Lemma 2.4.3.}
\underline{Let $(T,\Lambda)$ be a ringed topos, let $Z$ be a closed
subtopos of $T$, and let $i:Z\hookrightarrow T$ be the inclusion.}
\underline{Then the functor}
\[
i_*:D(Z,\Lambda)\longrightarrow D(T,\Lambda)
\]
\underline{is an equivalence from $D(Z,\Lambda)$ onto the full
subcategory of $D(T,\Lambda)$ consisting of the complexes $K$ for
which $\mathcal H^i(K)$ is supported on $Z$.}

Indeed, for such a complex $K$, the morphism
$K\longrightarrow i_*i^*K$ is a quasi-isomorphism.

We may consequently regard $R\Phi_{\bar\eta}(F)$ as an object of
\[
D^b(\Sigma_s,\Lambda)
=D^b(\Sigma_s-\{x\},\Lambda)\oplus D^b(\{x\},\Lambda).
\]
Passing to a direct summand, we find that
\[
\mathbb H^i\bigl(\{x\},R\Phi_{\bar\eta}(F)\bigr)
=R^i\Phi_{\bar\eta}(F)_x
\]
has the properties stated in 2.4.2.  This completes the proof.

\subsubsection*{Remark 2.4.4.}
The same proof shows that, if $K$ is constructible in
$D^+(X,\Lambda)$ and $x$ is isolated in the support of
$R\Phi_{\bar\eta}(K)$, then the cohomology modules of the stalk
$R\Phi_{\bar\eta}(K)_x$ are of finite type over $\Lambda$.  If, after
a change of trait, $x$ remains an isolated point of the support, these
modules are invariant under that change of trait.
\end{filecontents*}
% Original input: components/036_expose_XIII_section_2_4_5_through_2_4_6_6.tex
\begin{filecontents*}[overwrite,noheader]{sga7ii-cumulative-input-036.tex}
\subsubsection*{2.4.5.}
Let $x\in X_s$ be an isolated point where $(X,F)$ is not smooth.
More generally, let $K\in\operatorname{Ob}D^b(X,\Lambda)$, and let
$x$ be an isolated point of the support of $R\Phi(K)$.  Let
$i:\{x\}\hookrightarrow X_s$ be the inclusion.  The complex
$i_*i^*R\Phi(K)$, which is isomorphic to $R\Phi(K)$ in a neighborhood
of $x$, is a direct summand of $R\Phi(K)$ by 2.4.3.  Thus, for
$\sigma\in I$, the variation morphism induces
\[
\operatorname{Var}(\sigma):
i_*i^*R\Phi_{\bar\eta}(K)\longrightarrow R\Psi_{\bar\eta}(K),
\]
that is,
\begin{equation}\tag{2.4.5.1}
\operatorname{Var}(\sigma):
\bigl(R\Phi_{\bar\eta}(K)\bigr)_x
\longrightarrow
R\Gamma_{\{x\}}\bigl(R\Psi_{\bar\eta}(K)\bigr).
\end{equation}
The printed $R\Phi(F)$ in the direct-summand sentence is
$R\Phi(K)$ in this general formulation.

\subsubsection*{2.4.6.}
Suppose that the support $\Sigma$ of $R\Phi(K)$ consists of finitely
many isolated points.  For every $\sigma\in I$, the endomorphism
$\sigma-1$ of $R\Psi_{\bar\eta}(K)$ is then the composite displayed
in the following diagram:
\begin{equation}\tag{2.4.6.1}
\resizebox{\textwidth}{!}{%
\begin{tikzcd}[ampersand replacement=\&, column sep=huge, row sep=huge]
R\Psi_{\bar\eta}(K)
  \arrow[r] \arrow[d, "\sigma-1"']
  \& R\Phi_{\bar\eta}(K)
     =\displaystyle\bigoplus_{x\in\Sigma}
       i_{x*}\bigl(R\Phi_{\bar\eta}(K)\bigr)_x
     \arrow[d, "{\bigoplus i_{x*}\operatorname{Var}(\sigma)}"'] \\
R\Psi_{\bar\eta}(K)
  \& \displaystyle\bigoplus_{x\in\Sigma}
       i_{x*}R\Gamma_{\{x\}}R\Psi_{\bar\eta}(K)
     \arrow[l]
\end{tikzcd}}
\end{equation}

If $f$ is proper, applying $H^i$ to (2.4.6.1) shows that the
endomorphism $\sigma-1$ of $\mathbb H^i(X_{\bar\eta},K)$ is the
composite
\begin{equation}\tag{2.4.6.2}
\resizebox{\textwidth}{!}{%
\begin{tikzcd}[ampersand replacement=\&, column sep=large, row sep=huge]
\mathbb H^i(X_{\bar\eta},K)
  \arrow[r, "\sim"] \arrow[d, "\sigma-1"']
  \& \mathbb H^i\bigl(X_s,R\Psi_{\bar\eta}(K)\bigr)
     \arrow[r]
  \& \mathbb H^i\bigl(X_s,R\Phi_{\bar\eta}(K)\bigr)
     =\displaystyle\bigoplus_{x\in\Sigma}R^i\Phi_{\bar\eta}(K)_x
     \arrow[d, "{\bigoplus\operatorname{Var}(\sigma)}"'] \\
\mathbb H^i(X_{\bar\eta},K)
  \& \mathbb H^i\bigl(X_s,R\Psi_{\bar\eta}(K)\bigr)
     \arrow[l, "\sim"']
  \& \displaystyle\bigoplus_{x\in\Sigma}
     \mathbb H^i_{\{x\}}\bigl(X_s,R\Psi_{\bar\eta}(K)\bigr)
     \arrow[l]
\end{tikzcd}}
\end{equation}
This describes the action of $I$ on $\mathbb H^i(X_{\bar\eta},K)$;
its essential ingredients are the variation morphisms (2.4.5.1),
which are local on $X$.  We also have the long exact sequence
\begin{equation}\tag{2.4.6.3}
\longrightarrow \mathbb H^i(X_s,K)
\longrightarrow \mathbb H^i(X_{\bar\eta},K)
\longrightarrow \bigoplus_{x\in\Sigma}R^i\Phi_{\bar\eta}(K)_x
\longrightarrow\cdots .
\end{equation}

The same argument applies when $f$ is not necessarily proper,
provided that $(X,K)$ satisfies 2.1.9.  The diagram (2.4.6.5) below
defines a morphism
\begin{equation}\tag{2.4.6.4}
\operatorname{Var}(\sigma):
R\Gamma(X_{\bar\eta},K)\longrightarrow R\Gamma_c(X_{\bar\eta},K)
\end{equation}
such that the endomorphism $\sigma-1$ of
$R\Gamma(X_{\bar\eta},K)$ (respectively,
$R\Gamma_c(X_{\bar\eta},K)$) is the composite
$\alpha\operatorname{Var}(\sigma)$ (respectively,
$\operatorname{Var}(\sigma)\alpha$), where
\[
\alpha:R\Gamma_c(X_{\bar\eta},K)\longrightarrow
R\Gamma(X_{\bar\eta},K).
\]
\begin{equation}\tag{2.4.6.5}
\resizebox{\textwidth}{!}{%
\begin{tikzcd}[ampersand replacement=\&, column sep=huge, row sep=huge]
R\Gamma(X_{\bar\eta},K)
  \arrow[r, "\sim"] \arrow[d, "\operatorname{Var}(\sigma)"']
  \& R\Gamma\bigl(X_s,R\Psi_{\bar\eta}(K)\bigr)
     \arrow[r]
  \& \displaystyle\bigoplus_{x\in\Sigma}
     \bigl(R\Phi_{\bar\eta}(K)\bigr)_x
     \arrow[d, "\bigoplus"']
  \& {} \\
R\Gamma_c(X_{\bar\eta},K)
  \arrow[r, "\sim"']
  \& R\Gamma_c\bigl(X_s,R\Psi_{\bar\eta}(K)\bigr)
  \& \displaystyle\bigoplus_{x\in\Sigma}
     R\Gamma_{\{x\}}\bigl(X_s,R\Psi_{\bar\eta}(K)\bigr)
     \arrow[l]
  \& .
\end{tikzcd}}
\end{equation}

In this case we also have the long exact sequences
\begin{equation}\tag{2.4.6.6}
\resizebox{\textwidth}{!}{%
\begin{tikzcd}[ampersand replacement=\&, column sep=large, row sep=large]
\cdots \arrow[r]
  \& \mathbb H^i_c(X_s,K) \arrow[r] \arrow[d]
  \& \mathbb H^i_c(X_{\bar\eta},K) \arrow[r] \arrow[d]
  \& \displaystyle\bigoplus_{x\in\Sigma}
     \bigl(R^i\Phi_{\bar\eta}(K)\bigr)_x
     \arrow[r] \arrow[d, equal]
  \& \cdots \\
\cdots \arrow[r]
  \& \mathbb H^i(X_s,K) \arrow[r]
  \& \mathbb H^i(X_{\bar\eta},K) \arrow[r]
  \& \displaystyle\bigoplus_{x\in\Sigma}
     R^i\Phi_{\bar\eta}(K)_x \arrow[r]
  \& \cdots
\end{tikzcd}}
\end{equation}

The printed reference (2.4.4.1) in the discussion following
(2.4.6.2) is the variation morphism (2.4.5.1).
\end{filecontents*}
% Original input: components/037_expose_XIII_second_2_4_6_and_bibliography.tex
\begin{filecontents*}[overwrite,noheader]{sga7ii-cumulative-input-037.tex}
\subsubsection*{2.4.6. (second occurrence in the source)}
Let $X_{s(x)}$ be the henselization of $X_s$ at $x$, and put
\[
X^*_{s(x)}=X_{s(x)}-\{x\}.
\]
The long hypercohomology exact sequences of
$(X_{s(x)},\{x\})$, with coefficients in $K_s$,
$R\Psi_{\bar\eta}(K)$, and $R\Phi_{\bar\eta}(K)$, give the following
commutative diagram of long exact sequences:
\begin{equation}\tag{2.4.6.1}
\resizebox{\textwidth}{!}{%
\begin{tikzcd}[ampersand replacement=\&, column sep=large, row sep=large]
{}
  \& 0 \arrow[r] \arrow[d]
  \& R^{n-1}\Phi_{\bar\eta}(K)_x \arrow[r] \arrow[d]
  \& R^{n-1}\Phi_{\bar\eta}(K)_x \arrow[r] \arrow[d]
  \& 0 \arrow[d]
  \& {} \\
{} \arrow[r]
  \& \mathbb H^{n-1}(X^*_{s(x)},K)
     \arrow[r, "j"] \arrow[d, equal]
  \& \mathbb H^n_{\{x\}}(X_s,K) \arrow[r] \arrow[d]
  \& \underline{H}^{n}(K)_x \arrow[r] \arrow[d]
  \& \mathbb H^n(X^*_{s(x)},K)
     \arrow[r] \arrow[d, equal]
  \& {} \\
{} \arrow[r]
  \& \mathbb H^{n-1}(X^*_{s(x)},K)
     \arrow[r, "j", "{\mbox{\textcircled{1}}}"'] \arrow[d]
  \& \mathbb H^n_{\{x\}}\bigl(X_s,R\Psi_{\bar\eta}(K)\bigr)
     \arrow[r] \arrow[d]
  \& R^n\Psi_{\bar\eta}(K)_x \arrow[r] \arrow[d]
  \& \mathbb H^n(X_{s(x)},K) \arrow[r] \arrow[d]
  \& {} \\
{}
  \& 0 \arrow[r]
  \& R^n\Phi_{\bar\eta}(K)_x \arrow[r, equal]
  \& R^n\Phi_{\bar\eta}(K)_x \arrow[r]
  \& 0
  \& {}
\end{tikzcd}}
\end{equation}
The image of the arrow \textcircled{1} is contained in the inertia
invariants.

The source prints both this paragraph and the preceding paragraph as
2.4.6, and prints both principal diagrams as (2.4.6.1).  The repeated
numbers are retained here; this second occurrence is the one cited in
Expos\'e XV, 2.2.6.

\renewcommand{\refname}{Bibliography}
\begin{thebibliography}{9}
\bibitem{GiraudClassifyingTopos}
J.~Giraud, ``Classifying topos,'' pp.~43--56 in
\emph{Toposes, Algebraic Geometry and Logic},
Lecture Notes in Mathematics 274, Springer-Verlag, 1971.

\bibitem{NagataImbedding}
M.~Nagata, ``A generalization of the imbedding problem,''
\emph{J. Math. Kyoto} \textbf{3}, no.~1 (1963), pp.~89--102.
\end{thebibliography}
\end{filecontents*}
% Original input: components/038_expose_XIV_title_introduction_through_1_1_5.tex
\begin{filecontents*}[overwrite,noheader]{sga7ii-cumulative-input-038.tex}
\clearpage
\begin{center}
{\large SGA 7 --- EXPOS\'E XIV}\par
\vspace{2em}
{\Large\underline{Comparison with Transcendental Theory}}\par
\vspace{1.5em}
{\large by P. Deligne}
\end{center}

\medskip
\noindent\textbf{Source contents}
\begin{center}
\begin{tabular}{@{}llr@{}}
Introduction && 1\\
\S1. & Transcendental formalism of vanishing cycles & 1\\
1.1 & Fundamental group of the punctured disc: sign conventions & 2\\
1.2 & Sheaves of sets & 5\\
1.3 & The functor $R\Psi$ & 10\\
1.4 & A special case & 12\\
\S2. & The comparison theorem & 16\\
\S3. & Isolated singularities & 20\\
3.1 & The vanishing-cycle variety & 20\\
3.2 & The Picard--Lefschetz formula (transcendental case) & 31\\
\S4. & de Rham cohomology & 38
\end{tabular}
\end{center}

\medskip
\underline{Introduction}

To compare the theory of vanishing cycles in Expos\'e XIII with its
transcendental analogue over $\mathbb C$, we proceed in two stages.

\begin{enumerate}[label=(\alph*),leftmargin=2.6em]
\item We give a transcendental formalism modeled on that of
Expos\'e XIII and prove a comparison theorem by the usual methods
(already used to prove the finiteness theorem in XIII, 2.3).

\item In the case of isolated singularities, we explain in more
concrete terms the meaning of the transcendental formalism introduced
here.
\end{enumerate}

Besides the comparison theorem, this expos\'e contains a proof of the
Picard--Lefschetz formula in the transcendental case.

In \S4 we give a result preliminary to the Hodge theory of vanishing
cycles.  That section will not be used later in the seminar.

\section*{1. Transcendental formalism of vanishing cycles}

\subsection*{1.1. The fundamental group of the punctured disc:
sign conventions}

We shall have to consider punctured discs $D^*$ and their fundamental
group $\mathbb Z$.  Since this group is abelian, sign ambiguities may
arise.  We eliminate them by forgetting that it is abelian.

\subsubsection*{1.1.1.}
Let $X$ be a nice connected topological space and let $x\in X$.
The elements of the fundamental group are homotopy classes of loops
based at $x$.  If $\alpha,\beta\in\pi_1(X,x)$ are represented by
loops $a,b$, then $\alpha\beta$ is represented by the loop obtained
by traversing $b$ and then $a$.

\subsubsection*{1.1.2.}
Let $F$ be a locally constant sheaf on $X$.  For every path
$a:[0,1]\longrightarrow X$, the inverse image $a^*F$ on $[0,1]$ is
locally constant, hence constant, and therefore gives a bijection
\[
a:F_{a(0)}
\xleftarrow{\ \sim\ }H^0([0,1],a^*F)
\xrightarrow{\ \sim\ }F_{a(1)}.
\]
This bijection depends only on the homotopy class of $a$.  Through
this construction, $\pi_1(X,x)$ acts on the left on $F_x$.  The
functor $F\longmapsto F_x$ is an equivalence from the category of
locally constant sheaves on $X$ to the category of sets equipped with
an action of $\pi_1(X,x)$.

\subsubsection*{1.1.3.}
As $x$ varies, the groups $\pi_1(X,x)$ are the fibers of a locally
constant sheaf of groups $\Pi_1(X)$, which acts on every locally
constant sheaf $F$.  If $\pi_1(X,x)$ is abelian, then $\Pi_1(X)$ is
constant; its value is denoted by $\pi_1(X)$.

\subsubsection*{1.1.4.}
The group $\pi_1(X,x)$ acts on the universal covering
\[
p:(\widetilde X,\widetilde x)\longrightarrow(X,x).
\]
Regarded as an \'etal\'e space over $X$, and hence as a sheaf (the
sheaf of its local sections), this covering corresponds under 1.1.2
to the set $\pi_1(X,x)$ with distinguished element $e$, on which
$\pi_1(X,x)$ acts by left translation.  The group also acts on this
$\pi_1$-set by automorphisms of $\pi_1$-sets: an element
$\sigma\in\pi_1(X,x)$ acts by
\[
g\longmapsto g\sigma^{-1}.
\]
This construction identifies
\[
\pi_1(X,x)=\operatorname{Aut}_X(\widetilde X).
\]

If $F$ is locally constant on $X$, the bijection
\begin{equation}\tag{1.1.4.1}
H^0(\widetilde X,p^*F)
\longrightarrow(p^*F)_{\widetilde x}=F_x
\end{equation}
commutes with the action of $\pi_1(X,x)$: on the left the action is
transport of structure through its action on $\widetilde X$, while on
the right it is the action of 1.1.2.  The equivalence of 1.1.2 may
therefore be reformulated by saying that
\[
F\longmapsto
\bigl(\widetilde X\longmapsto
H^0(\widetilde X,p^*F)\bigr)
\]
is an equivalence from the category of locally constant sheaves on
$X$ to the category of functors
\[
(\text{universal coverings of }X)\longrightarrow(\mathrm{Sets}).
\]

\subsubsection*{1.1.5.}
If $\pi_1(X,x)$ is abelian and hence independent of $x$, the action
of 1.1.4 of $\pi_1(X)$ on $\widetilde X$ is independent of the base
point: $\pi_1(X)$ acts on every universal covering $\widetilde X$ of
$X$.

One must bear in mind that this action, which is the only one we shall
use, corresponds to the inverse of the action in 1.1.3.
\end{filecontents*}
% Original input: components/039_expose_XIV_section_1_1_6_through_1_1_11.tex
\begin{filecontents*}[overwrite,noheader]{sga7ii-cumulative-input-039.tex}
\subsubsection*{1.1.6.}
Let $D$ be a disc, that is, a topological space homeomorphic to the
standard disc
\[
\{z\in\mathbb C:|z|<1\}.
\]
Let $0\in D$ and put $D^*=D-\{0\}$.  Suppose that $D$ is equipped
with the structure of an analytic manifold, and hence is oriented.
The fundamental group $\pi_1(D^*)$ is then canonically isomorphic to
$\mathbb Z$.  For the standard disc, with center $0$ and
$t\in D^*$, the positive generator of $\pi_1(D^*,t)$ is represented
by the loop
\[
x\longmapsto \exp(2\pi i x)t.
\]

\subsubsection*{1.1.7.}
Let $z:D\longrightarrow\mathbb C$ be a uniformizing parameter, so
that $(dz)_0\ne0$, and let $\widetilde D^*$ be a universal covering
of $D^*$.  Write $\log z$ for any function on $\widetilde D^*$ such
that
\[
\exp(\log z)=z.
\]
The pair $(\widetilde D^*,\log z)$ is unique up to unique
isomorphism.

For $\alpha\in\mathbb C$, put
\[
z^\alpha=\exp(\alpha\log z).
\]
Let $T$ be the positive generator of $\pi_1(D^*)$, acting on
$\widetilde D^*$ as in 1.1.5 and on functions on
$\widetilde D^*$ by transport of structure.  Then
\begin{equation}\tag{1.1.7.1}
T\log z=\log z+2\pi i,
\end{equation}
and
\begin{equation}\tag{1.1.7.2}
Tz^\alpha=\exp(2\pi i\alpha)z^\alpha.
\end{equation}

\subsubsection*{1.1.8.}
Let $D$ be a disc, let $0\in D$, and put $D^*=D-\{0\}$.  If
$D'\subset D$ is a smaller disc containing $0$, restriction to
$D'^{*}$ is an equivalence from the category of locally constant
sheaves on $D^*$ to the category of locally constant sheaves on
$D'^{*}$.  Indeed, for $t\in D'^{*}$,
\[
\pi_1(D'^{*},t)\xrightarrow{\ \sim\ }\pi_1(D^*,t).
\]

Suppose now that $D$ is equipped with a differentiable-manifold
structure.  Let $T$ be the tangent space to $D$ at $0$, and put
$T^*=T-\{0\}$.  For $D'\subset D$ as above and an embedding
$f:D'\hookrightarrow T$ satisfying $f(0)=0$ and $(df)_0=\mathrm{id}$,
the composite
\[
(\text{restriction to }D')^{-1}\circ f^*
\]
is an equivalence $[f]$ from the category of locally constant sheaves
on $T^*$ to the category of locally constant sheaves on $D^*$.
Up to canonical isomorphism it depends only on the germ of $f$ at
$0$, subject to $(df)_0=\mathrm{id}$.  A path from one such germ to
another defines an isomorphism of equivalences, and homotopic paths
define the same isomorphism.  Since the space of such germs is
contractible---by the retraction to a fixed germ $f_0$ given by
$(1-\lambda)f+\lambda f_0$---we obtain an equivalence, well defined
up to unique isomorphism, from locally constant sheaves on $T^*$ to
locally constant sheaves on $D^*$.

\subsubsection*{1.1.9.}
Let $D$ be a disc, let $0\in D$, put $D^*=D-\{0\}$, and let
$p:\widetilde D^*\longrightarrow D^*$ be a universal covering.  We
write
\[
p:\widetilde D\longrightarrow D
\]
for the topological space over $D$ obtained from
$\widetilde D^*$ by adjoining a point $0$, in such a way that:
\begin{enumerate}[label=(\alph*),leftmargin=2.6em]
\item $\widetilde D^*$ is open in $\widetilde D$;
\item $p:\widetilde D\longrightarrow D$ extends
  $p:\widetilde D^*\longrightarrow D^*$ and $p(0)=0$;
\item the sets $p^{-1}(U)$, as $U$ ranges over the neighborhoods of
  $0$ in $D$, form a neighborhood basis of $0$ in $\widetilde D$.
\end{enumerate}

\subsubsection*{1.1.10.}
Let $X$ be a nonsingular complex algebraic curve, let $0\in X$, and
let $D$ be a disc neighborhood of $0$ in $X$.  The ring of germs at
$0$ of holomorphic functions on $D$ which are algebraic over the
field $K(X)$ of rational functions on $X$ is the henselization
$\mathcal O_0^h$ of the local ring $\mathcal O_0$ of $X$ at $0$.
Let $\widetilde D^*$ be a universal covering of
$D^*=D-\{0\}$.  The field of germs at $0\in\widetilde D$ of
meromorphic functions on $\widetilde D^*$ which are algebraic over
$K(X)$ is an algebraic closure $\overline K$ of $K(X)$, and hence
also of the fraction field $K_0^h$ of $\mathcal O_0^h$.

The group $\pi_1(D^*)$ of $D$-automorphisms of
$\widetilde D^*$ (1.1.5) acts by transport of structure on
$\overline K$, giving
\begin{equation}\tag{1.1.10.1}
\begin{tikzcd}[column sep=huge, row sep=large]
\pi_1(D^*) \arrow[r] \arrow[d, equal]
  & \operatorname{Gal}(\overline K/K_0^h) \arrow[d, equal]
  & {} \\
\mathbb Z
  & \widehat{\mathbb Z}(1)
  & .
\end{tikzcd}
\end{equation}
It follows from (1.1.7.2) that this diagram is compatible with the
natural inclusion, over $\mathbb C$, of $\mathbb Z$ into
$\widehat{\mathbb Z}(1)$.  This inclusion is the composite of
$\mathbb Z\hookrightarrow\widehat{\mathbb Z}$ and the isomorphism
\[
\widehat{\mathbb Z}\xrightarrow{\ \sim\ }
\widehat{\mathbb Z}(1),
\]
which is the inverse limit of the isomorphisms
\[
\mathbb Z/n\longrightarrow\mu_n,
\qquad
k\longmapsto\exp\left(\frac{2\pi i}{n}k\right).
\]

\subsubsection*{1.1.11.}
Every sheaf $F$ on $\Spec(K_0^h)$ defines a locally constant sheaf on
$D^*$.  We use the equivalence of 1.1.4 and, to each universal
covering $\widetilde D^*$ of $D^*$, associate $F_{\overline K}$ for
the field $\overline K$ constructed in 1.1.10.
\end{filecontents*}
% Original input: components/040_expose_XIV_section_1_2_1_through_1_2_4.tex
\begin{filecontents*}[overwrite,noheader]{sga7ii-cumulative-input-040.tex}
\subsection*{1.2. Sheaves of sets}

\subsubsection*{1.2.1.}
Let $D$ be a disc (1.1.6), let $0\in D$, and put
$D^*=D-\{0\}$.  We denote by $\mathcal D^*$ the topos of locally
constant sheaves of sets on $D^*$.  By 1.1.8 this topos does not
``change'' if $D$ is replaced by a smaller disc $D'$ containing $0$,
nor, when $D$ is endowed with the structure of a differentiable
manifold, if $D$ is replaced by its tangent space at $0$.

For $t\in D^*$, the functor $F\longmapsto F_t$ identifies this topos
with the topos of sets equipped with an action of $\pi_1(D^*)$
(1.1.2).  By 1.1.4, the functor
\[
F\longmapsto
\bigl(\widetilde D^*\longmapsto
H^0(\widetilde D^*,p^*F)\bigr)
\]
identifies it with the topos of functors
\[
(\text{universal coverings of }D^*)
\longrightarrow(\mathrm{Sets}),
\]
and these identifications are related by (1.1.4.1).

\subsubsection*{1.2.2.}
Let $F$ be a sheaf on $D$ whose restriction to $D^*$ is locally
constant.  There is a unique morphism of sheaves
\[
\underline{F_0}\xrightarrow{\ \alpha\ }F
\]
whose map on the stalk at $0$ is the identity.  As explained in
XIII, 0.1 (with $X_0=(P^t)$), the functor
\[
F\longmapsto(F_0,F|_{D^*},\alpha)
\]
identifies these sheaves with triples consisting of a set $F_0$, a
locally constant sheaf $F_\eta$ on $D^*$, and a morphism
\[
\alpha:
(\text{constant sheaf with value }F_0\text{ on }D^*)
\longrightarrow F_\eta.
\]

\subsubsection*{1.2.3.}
We denote by $\mathcal D$ the topos of sheaves of sets on $D$ whose
restriction to $D^*$ is locally constant.  The inclusion of
$\mathcal D$ into the topos of sheaves on $D$ is the inverse-image
functor of a morphism of topoi $D\longrightarrow\mathcal D$, which
fits into the commutative diagram
\[
\begin{tikzcd}
\{0\} \arrow[r, hook] \arrow[d, equal]
  & D \arrow[d, "\beta"]
  & D^* \arrow[l, hook'] \arrow[d, "\beta"] & {} \\
\{0\} \arrow[r, hook]
  & \mathcal D
  & \mathcal D^* \arrow[l, hook'] & {}.
\end{tikzcd}
\]
As in 1.1.1, one checks that this topos does not change when $D$ is
replaced by a smaller disc or by the tangent space to $D$ at $0$.

\subsubsection*{1.2.4.}
Let $Y$ be a topological space.  The product topos
$Y\times\mathcal D$ has the following three descriptions (for their
equivalence, see XIII, 0.1).

\medskip
\noindent\textit{(a)} $Y\times\mathcal D$ is the topos of sheaves
$F$ on $Y\times D$ satisfying

\smallskip
\noindent $(*)$ $D^*$ is covered by open subsets $U$ such that
$F|_{Y\times U}$ is the inverse image of a sheaf on $Y$.

\medskip
\noindent\textit{(b)} Let $t\in D^*$.  The topos
$Y\times\mathcal D$ is the topos of triples $(F_0,F_t,\alpha)$ in
which

\smallskip
\noindent $\alpha)$ $F_0$ is a sheaf on $Y$;

\noindent $\beta)$ $F_t$ is a sheaf on $Y$ equipped with an action
of $\pi_1(D^*)$;

\noindent $\gamma)$ $\alpha:F_0\longrightarrow F_t$ is a morphism
whose image is contained in the invariants.

\medskip
\noindent\textit{(c)} The topos $Y\times\mathcal D$ is the topos of
triples $(F_0,F_\eta,\alpha)$ in which

\smallskip
\noindent $\alpha')$ $F_0$ is a sheaf on $Y$;

\noindent $\beta')$ $F_\eta$ is a functor, with values denoted by
$F_{\bar\eta}$,
\[
(\text{universal coverings of }D^*)
\longrightarrow(\text{sheaves on }Y);
\]

\noindent $\gamma')$ $\alpha$ is a morphism from $F_0$ (regarded as
the constant functor with value $F_0$ from the universal coverings
of $D^*$ to the sheaves on $Y$) to $F_\eta$.

It is immaterial here whether this topos is literally the product
topos of $Y$ and $\mathcal D$.  Likewise, we define
$Y\times\mathcal D^*$, at our choice, as

\medskip
\noindent\textit{(a$^*$)} the topos of sheaves on $Y\times D^*$
satisfying $(*)$;

\noindent\textit{(b$^*$)} for $t\in D^*$, the topos of sheaves $F_t$
on $Y$ equipped with an action of $\pi_1(D^*)$;

\noindent\textit{(c$^*$)} the topos of functors
\[
(\text{universal coverings of }D^*)
\longrightarrow(\text{sheaves on }Y).
\]

\end{filecontents*}
% Original input: components/041_expose_XIV_section_1_2_5_through_1_2_6.tex
\begin{filecontents*}[overwrite,noheader]{sga7ii-cumulative-input-041.tex}
\subsubsection*{1.2.5.}
The topos $Y\times\mathcal D$ has the same functoriality properties
in $Y$ as the topoi $Y\times_S S$ of Expos\'e XIII, and we shall not
list them.  To fix notation, we merely record the diagram
\[
\begin{tikzcd}[row sep=huge, column sep=huge]
Y \arrow[d, ""{name=A}]
  & Y \arrow[d, ""{name=B}] \arrow[r, hook, "i"]
  & Y\times\mathcal D \arrow[d]
  & Y\times\mathcal D^* \arrow[d] \arrow[l, hook', "j"'] \\
(P^t)
  & \{0\} \arrow[r, hook, "i"]
  & \mathcal D
  & \mathcal D^* \arrow[l, hook', "j"']
    \arrow[from=B, to=A, "\mathrm{sp}"']
\end{tikzcd}
\]

\subsubsection*{1.2.6.}
Let $K$ be a complex of sheaves of abelian groups on
$Y\times\mathcal D$.  As in XIII, 1.4.2, one obtains from $K$ a
distinguished triangle in $K(Y\times\mathcal D^*)$:
\[
\begin{tikzcd}[row sep=large, column sep=large]
{} & \Phi(K) \arrow[dl] & {} \\
\mathrm{sp}^*K_0 \arrow[rr] & {} & K_\eta \arrow[ul]
\end{tikzcd}
\]
This construction passes to the derived category and defines a
functor
\[
D(Y\times\mathcal D)
\longrightarrow
(\text{distinguished triangles of }D(Y\times\mathcal D^*)).
\]

Suppose that $D$ is oriented (for example, complex analytic), and
let $T$ be the positive generator of $\pi_1(D^*)$ (1.1.6).  For
$K\in\operatorname{Ob}D(Y\times\mathcal D)$, we define, as in XIII,
1.4.3, the \underline{variation} $V=\operatorname{Var}(T)$,
\[
V:\Phi(K)\longrightarrow K_\eta.
\]
The endomorphism $T-\mathrm{Id}$ of $K_\eta$ is the composite
\[
T-\mathrm{Id}:K_\eta\longrightarrow\Phi(K)
\longrightarrow K_\eta.
\]

\end{filecontents*}
% Original input: components/042_expose_XIV_section_1_2_7_through_1_2_8.tex
\begin{filecontents*}[overwrite,noheader]{sga7ii-cumulative-input-042.tex}
\subsubsection*{1.2.7.}
Let $D$ be a disc, let $0\in D$, put $D^*=D-\{0\}$, let
$f:X\longrightarrow D$ be a continuous map, and put
$X_0=f^{-1}(0)$.  We shall define a functor
\[
\Psi:(\text{sheaves on }X)
\longrightarrow(\text{sheaves on }X_0\times\mathcal D).
\]
Let $\widetilde D^*$ be a universal covering of $D^*$, let
$\widetilde D$ be as in 1.1.9, and put
\[
\overline X=X\times_D\widetilde D,
\qquad
X^*=X-X_0=X\times_DD^*,
\qquad
\overline X^*=\overline X-X_0=X\times_D\widetilde D^*.
\]
There is a commutative diagram
\begin{equation}\tag{1.2.7.1}
\begin{tikzcd}
X_0 \arrow[r, hook, "\overline i"] \arrow[d, equal]
  & \overline X \arrow[d, "p"]
  & \overline X^* \arrow[l, hook, "\overline j"']
      \arrow[d, "p'"] \\
X_0 \arrow[r, hook, "i"]
  & X
  & X^* \arrow[l, hook, "j"']
\end{tikzcd}
\end{equation}
cartesian over the diagram
\begin{equation}\tag{1.2.7.2}
\begin{tikzcd}
\{0\} \arrow[r, hook, "\overline i"] \arrow[d, equal]
  & \widetilde D \arrow[d, "p"]
  & \widetilde D^* \arrow[l, hook, "\overline j"']
      \arrow[d, "p'"] \\
\{0\} \arrow[r, hook, "i"]
  & D
  & D^* \arrow[l, hook, "j"']
\end{tikzcd}
\end{equation}
\noindent\emph{Editorial note.}  In (1.2.7.1) the printed top-right
node lacks the overline.  The preceding definition and the map $p'$
force that node to be $\overline X^*$, which is the notation used
here.

For a sheaf $F$ on $X$, we set (with the notation of 1.2.4(c))
\[
\Psi(F)=
\bigl(i^*F,\ \overline i^*\overline j_*((jp')^*F),\ \alpha\bigr),
\]
where
\[
\begin{split}
\alpha:i^*F=\overline i^*p^*F
&\longrightarrow
\overline i^*\overline j_*((jp')^*F)\\
&=\overline i^*\overline j_*\overline j^*p^*F
\end{split}
\]
is obtained, after applying $\overline i^*$, from the adjunction map
$\mathrm{Id}\longrightarrow\overline j_*\overline j^*$.

We also define a functor
\[
\Psi_\eta:(\text{sheaves on }X^*)
\longrightarrow(\text{sheaves on }X_0\times\mathcal D^*)
\]
by
\[
\Psi_\eta(F)=\overline i^*\overline j_*({p'}^*F).
\]

\subsubsection*{1.2.8.}
The functoriality properties of $\Psi$ are the same as those of its
analogue studied in XIII, 1.3.5--1.3.9.

\end{filecontents*}
% Original input: components/043_expose_XIV_section_1_3_1_through_1_3_3.tex
\begin{filecontents*}[overwrite,noheader]{sga7ii-cumulative-input-043.tex}
\subsection*{1.3. The functor $R\Psi$}

\subsubsection*{1.3.1.}
Let $D$ be a disc, let $0\in D$, put $D^*=D-\{0\}$, let
$f:X\longrightarrow D$ be a continuous map, and put
$X_0=f^{-1}(0)$.  The functors $R\Psi$ and $R\Psi_\eta$ are the
derived functors of $\Psi$ and $\Psi_\eta$:
\[
R\Psi:D^+(X)\longrightarrow D^+(X_0\times\mathcal D),
\]
\[
R\Psi_\eta:D^+(X^*)\longrightarrow D^+(X_0\times\mathcal D^*).
\]
We denote by $R\Phi$ the composite functor $\Phi\circ R\Psi$:
\[
R\Phi:D^+(X)\longrightarrow D^+(X_0\times\mathcal D^*).
\]
The properties of these functors are exactly analogous to those of
the functors with the same names in XIII, 2.1.  Their proofs are much
more elementary, since in the topological setting the base-change
theorems for a proper morphism or for a smooth morphism (that is,
one locally isomorphic to $\mathbb R^n\times S\longrightarrow S$)
are easy to prove.  The minor changes required are as follows.

\subsubsection*{1.3.2.}
The functor $Rf_!$ is defined only for a separated morphism between
locally compact spaces, and is the derived functor of the direct-image
functor with proper supports $f_!$.

\subsubsection*{1.3.3.}
For a proper (and separated) morphism $f:X\longrightarrow D$, the
fundamental isomorphism (XIII, (2.1.8.3)) must be rewritten as follows.
Let $p:\widetilde D^*\longrightarrow D^*$ be a universal covering of
$D^*$.  For every neighborhood $U$ of $0$ in $D$, put
\[
\widetilde U^*=p^{-1}(U),
\qquad
X_{\widetilde U^*}=X\times_D\widetilde U^*.
\]
For $K\in\operatorname{Ob}D^+(f^{-1}(D^*))$, one has
\begin{equation}\tag{1.3.3.1}
\varinjlim_U\mathbb H^i(X_{\widetilde U^*},K)
\xrightarrow{\ \sim\ }
\mathbb H^i(X_0,R\Psi_{\bar\eta}(K)).
\end{equation}
This isomorphism is especially useful when
$X^*=f^{-1}(D^*)$ is a topological fiber bundle over $D^*$.  As a
disc $U\subset D$ shrinks to $0$, the Leray spectral sequence for
$X_{\widetilde U^*}\longrightarrow\widetilde U^*$ then shows that,
for every locally constant sheaf $\Lambda$ on $X^*$ and every
$t\in\widetilde U^*$,
\[
H^i(X_{\widetilde U^*},\Lambda)
\xrightarrow{\ \sim\ }H^i(X_t,\Lambda).
\]
Thus, for $t\in D^*$,
\begin{equation}\tag{1.3.3.2}
H^i(X_t,\Lambda)
\xrightarrow{\ \sim\ }
H^i(X_0,R\Psi_t(\Lambda)),
\end{equation}
provided that $X/D$ is proper, $X^*$ is a fiber bundle, and
$\Lambda$ is locally constant.

In practice, to verify that $X^*$ is a fiber bundle over $D^*$ one
uses Thom's isotopy theorem, one form of which is stated below.  For
the proof of this theorem, none of whose consequences will be used in
the subsequent expos\'es, we refer to Whitney and Mather.

\end{filecontents*}
% Original input: components/044_expose_XIV_section_1_3_4_through_1_3_6.tex
\begin{filecontents*}[overwrite,noheader]{sga7ii-cumulative-input-044.tex}
\subsubsection*{1.3.4. Preliminary.}
Let $f:X\longrightarrow Y$ be a smooth morphism (by definition, one
with no critical point) of differentiable manifolds.  A
\underline{connection} $\nabla$ on $X/Y$ is a vector subbundle of the
tangent bundle of $X$ complementary to $\ker(df)$.  For a vector
field $u$ on $Y$, there is a unique vector field $\nabla(u)$ on $X$
which lifts $u$ and lies in this subbundle.  The connection is called
integrable if
\[
\nabla([u,v])=[\nabla(u),\nabla(v)].
\]
When $f$ is proper, an integrable connection gives local
trivializations over $Y$ of $X/Y$,
\[
\alpha:f^{-1}(U)\xrightarrow{\ \sim\ }X_0\times U,
\]
compatible with projection to $U$ and satisfying
$d\alpha(\nabla(u))=(0,u)$.  When $f$ is not proper, the same result
holds under suitable conditions at infinity; for example, it is
enough that $\nabla(u)g$ vanish for a suitable exhaustion function
$g$.

\subsubsection*{1.3.5. Thom's isotopy theorem.}
\textit{Let
\[
D=\{z\in\mathbb C:|z|<1\},
\qquad D^*=D-\{0\},
\]
and let $f:X\longrightarrow D$ be a proper morphism of analytic
spaces.  After replacing $D$ by a smaller disc and $X$ by its inverse
image, there exists a stratification $\Sigma$ of
$X^*=f^{-1}(D^*)$ having the following properties:}

\begin{enumerate}[label=(\alph*),leftmargin=2.7em]
\item \textit{$\Sigma$ is a filtration
\[
X^*=F_0\supset F_1\supset\cdots\supset F_n
\supset F_{n+1}=\varnothing
\]
of $X^*$ by closed subsets such that the closure $\overline F_i$ of
$F_i$ in $X$ is an analytic subspace of $X$.}

\item \textit{The open strata $F_i-F_{i+1}$ are smooth over $D^*$.}

\item \textit{$\Sigma$ satisfies Whitney's conditions (a) and (b).}
\end{enumerate}

\textit{For such a stratification $\Sigma$, there are integrable
connections $\nabla_i$ on the $F_i-F_{i+1}$ such that}

\begin{enumerate}[label=(\roman*),leftmargin=2.7em]
\item \textit{locally on $D^*$, $\nabla_i$ defines a trivialization
of $F_i-F_{i+1}$;}

\item \textit{these trivializations glue to give local
trivializations of $X^*/D^*$.}
\end{enumerate}

\textit{Finally, if $(Y_\alpha)$ is a finite family of closed
analytic subspaces of $X$, then, after shrinking $D$ further, one can
choose $\Sigma$ so that every $Y_\alpha$ is a union of connected
components of open strata $(F_i-F_{i+1})$.  A connection $\nabla$ as
above then locally trivializes $(X,(Y_\alpha))$.}

\subsubsection*{1.3.6.}
The isomorphism (1.3.3.2) remains valid for hypercohomology with
values in a complex of sheaves $K$ when $X$ is proper over $D$ and
the following condition holds:

\medskip
\noindent $(*)$ For every $i\in\mathbb Z$, the space $D^*$ is
covered by open subsets $U$ for which there exist an isomorphism
compatible with projection to $U$,
\[
\varphi:f^{-1}(U)\xrightarrow{\ \sim\ }Y\times U,
\]
and a sheaf $G$ on $Y$ such that, on $f^{-1}(U)$,
$\underline H^i(K)$ is isomorphic to $\varphi^*\operatorname{pr}_1^*G$.

Under this hypothesis, the triangle of 1.2.6 gives a long exact
sequence
\begin{equation}\tag{1.3.6.1}
\cdots\longrightarrow
\mathbb H^i(X_0,K)\longrightarrow
\mathbb H^i(X_t,K)\longrightarrow
\mathbb H^i(X_0,R\Phi_t(K))\longrightarrow\cdots.
\end{equation}
\noindent\emph{Editorial note.}  The print cites the triangle
``1.2.1'' here; the distinguished triangle used is the one constructed
in 1.2.6.

By Thom's theorem, condition $(*)$ holds (after shrinking $D$) whenever
$X$ admits a complex analytic stratification such that all the
$\underline H^i(K)$ are locally constant on every open stratum.

\end{filecontents*}
% Original input: components/045_expose_XIV_section_1_4_1_through_1_4_2.tex
\begin{filecontents*}[overwrite,noheader]{sga7ii-cumulative-input-045.tex}
\subsection*{1.4. A special case}

This section will not be used later in the seminar.  It formalizes a
construction sketched in the Introduction to Expos\'e XIII.

\subsubsection*{1.4.1.}
Let $D$ be an oriented disc, let $0\in D$, and put
$D^*=D-\{0\}$.  A \underline{topological space} $\mathcal X^*$ over
$\mathcal D^*$ is a functor
\[
(\text{universal coverings of }D^*)
\longrightarrow(\text{topological spaces}).
\]
After a universal covering $\widetilde D^*$ has been chosen, this is
a topological space $X_\eta$ equipped with an action of
$\pi_1(D^*)$; when $\widetilde D^*$ is the universal covering of
$(D^*,t)$, we denote it by $\mathcal X_t^*$.

A \underline{sheaf} on $\mathcal X^*$ is a functor
\[
(\text{universal coverings of }D^*)
\longrightarrow
(\text{topological spaces equipped with a sheaf}),
\]
or, equivalently, a $\pi_1(D^*)$-equivariant sheaf on $X_\eta$.

To $\mathcal X^*$ we associate the topological space
\[
\beta^*\mathcal X^*
=X_\eta\times\widetilde D^*/\pi_1(D^*)
\]
over $D^*$.  Every sheaf $F$ on $\mathcal X^*$ descends to a sheaf
$\beta^*F$ on $\beta^*\mathcal X^*$, giving a commutative diagram of
topoi
\[
\begin{tikzcd}
\beta^*\mathcal X^* \arrow[r] \arrow[d]
  & \mathcal X^* \arrow[d] \\
D^* \arrow[r, "\beta"]
  & \mathcal D^*.
\end{tikzcd}
\]

\subsubsection*{1.4.2.}
A \underline{topological space} $\mathcal X$ over $\mathcal D$
consists of
\begin{enumerate}[label=(\alph*),leftmargin=2.7em]
\item a topological space $\mathcal X^*$ over $\mathcal D^*$;
\item a topological space $X_0$;
\item a $\pi_1(D^*)$-invariant morphism
\[
\mathrm{sp}:X_\eta\longrightarrow X_0,
\]
where $\pi_1(D^*)$ acts trivially on $X_0$.
\end{enumerate}

A \underline{sheaf} on $\mathcal X$ consists of a sheaf $F_0$ on
$X_0$, a sheaf on $\mathcal X^*$ given by $F_\eta$, and an equivariant
morphism
\[
\mathrm{sp}:\mathrm{sp}^*F_0\longrightarrow F_\eta
\qquad
(\text{or }F_0\longrightarrow\mathrm{sp}_*F_\eta).
\]
There is a commutative diagram
\[
\begin{tikzcd}
X_0 \arrow[r, hook] \arrow[d]
  & \mathcal X \arrow[d]
  & \mathcal X^* \arrow[l, hook] \arrow[d] \\
\{0\} \arrow[r, hook]
  & \mathcal D
  & \mathcal D^* \arrow[l, hook]
\end{tikzcd}
\]
and, in addition, a morphism of topoi
\[
\Psi:\mathcal X\longrightarrow X_0\times\mathcal D.
\]
It is obtained from the functoriality of the preceding construction,
applied to
\[
(X_\eta,X_0,\mathrm{sp})
\longrightarrow(X_0,X_0,\mathrm{Id}),
\]
and satisfies
\[
\Psi(F_0,F_\eta,\mathrm{sp})
=(F_0,\mathrm{sp}_*F_\eta,\mathrm{sp}).
\]
This morphism induces
\[
\Psi_\eta:\mathcal X^*\longrightarrow X_0\times\mathcal D^*.
\]

\end{filecontents*}
% Original input: components/046_expose_XIV_section_1_4_3_through_1_4_5.tex
\begin{filecontents*}[overwrite,noheader]{sga7ii-cumulative-input-046.tex}
\subsubsection*{1.4.3.}
To a topological space $\mathcal X$ over $\mathcal D$ we associate a
topological space $X=\beta^*\mathcal X$ over $D$ as follows.  To
$X^*=\beta^*\mathcal X^*$ we adjoin the fiber $X_0$ over $0$.  A
fundamental system of neighborhoods of $x_0\in X_0$ is formed by the
sets
\[
V\cup
\bigl(\mathrm{sp}^{-1}(V)\times\widetilde U^*/\pi_1(D^*)\bigr),
\]
where $V$ is a neighborhood of $x_0$ in $X_0$, $U$ is a neighborhood
of $0$ in $D$, and
$\widetilde U^*=U\times_D\widetilde D^*$.  Let $j$ (respectively $i$)
denote the inclusion of $X^*$ (respectively $X_0$) into $X$.  For a
sheaf $F$ on $\mathcal X^*$, one has
\[
i^*j_*\beta^*F
=(\mathrm{sp}_*F_\eta)^{\pi_1(D^*)}.
\]
It follows that there is a commutative diagram of topoi
\[
\begin{tikzcd}
X \arrow[r, "\beta"] \arrow[d]
  & \mathcal X \arrow[d] \\
D \arrow[r, "\beta"]
  & \mathcal D
\end{tikzcd}
\]
such that, for every sheaf $F$ on $\mathcal X$, there is an
isomorphism
\begin{equation}\tag{1.4.3.1}
\Psi(F)\xrightarrow{\ \sim\ }\Psi(\beta^*F).
\end{equation}

\subsubsection*{Proposition 1.4.4.}
The map (1.4.3.1) derives to an isomorphism
\[
R\Psi\xrightarrow{\ \sim\ }R\Psi\circ\beta^*
\]
between functors from $D^+(\mathcal X)$ to
$D^+(X_0\times\mathcal D)$.

Indeed, let $\overline j$ be the inclusion of
\[
\overline X^*=X\times_D\widetilde D^*
\simeq X_\eta\times\widetilde D^*
\]
into $\overline X=X\times_D\widetilde D$.  The proposition reduces to
checking that, for every flabby sheaf of abelian groups $F$ on
$X_\eta$, the sheaf $\operatorname{pr}_1^*F$ is acyclic for
$\overline j_*$.  This in turn reduces (with $B=\widetilde D^*$) to
the following familiar assertion, proved by a homotopy argument.  If
$F$ is a sheaf of abelian groups on a space $Y$ and $B$ is a disc,
then
\[
H^i(Y,F)\xrightarrow{\ \sim\ }
H^i(Y\times B,\operatorname{pr}_1^*F).
\]

\subsubsection*{1.4.5.}
Proposition 1.4.4 gives a method for computing $R\Psi$ for sheaves on
$X$ of the form $\beta^*F$, reducing the computation of $R\Psi$
essentially to that of $R\,\mathrm{sp}_*$.

\end{filecontents*}
% Original input: components/047_expose_XIV_section_1_4_6_through_1_4_7.tex
\begin{filecontents*}[overwrite,noheader]{sga7ii-cumulative-input-047.tex}
\subsubsection*{1.4.6.}
Later we shall consider the case in which $X_\eta$ is locally compact
and $X_0$ is obtained from $X_\eta$ by contracting a compact subset
$V\subset X_\eta$ to a point $0$, with $\mathrm{sp}$ the contraction
map and the action of $\pi_1(D^*)$ trivial outside $V$.  Proper base
change for $\mathrm{sp}$ then gives, for
$K\in\operatorname{Ob}D^+(\mathcal X^*)$,
\begin{equation}\tag{1.4.6.1}
R\Psi_\eta(\beta^*K)_0=R\Gamma(V,K)
\end{equation}
and
\begin{equation}\tag{1.4.6.2}
R\Gamma_{\{0\}}R\Psi_\eta(\beta^*K)
=R\Gamma_V(X_\eta,K),
\end{equation}
the latter being cohomology with supports.

For $K$ on $\mathcal X$, suppose that
\[
\mathrm{sp}:\mathrm{sp}^*K_0\longrightarrow K_\eta
\]
is a quasi-isomorphism outside $V$.  Then $R\Phi(K)=0$ outside $0$.
In particular, for a constant sheaf $\Lambda$ on $\mathcal X$,
\[
R^i\Phi(\Lambda)=\widetilde H^i(V,\Lambda),
\qquad\text{supported at }0,
\]
where the right-hand side is reduced cohomology.

\subsubsection*{1.4.7.}
The case of interest to us is that in which $X_\eta$ is a
differentiable manifold and $V\subset X_\eta$ is a submanifold with
boundary and interior $V^\circ$.  Then $X-\{0\}$ is a differentiable
manifold obtained by gluing $\beta^*X_\eta$ to the constant bundle
over $D$ with fiber $X_\eta-V$.  The map
\begin{equation}\tag{1.4.7.1}
H_c^i(V^\circ,\Lambda)
\xrightarrow{\ \sim\ }H_V^i(X_\eta,\Lambda)
\end{equation}
is an isomorphism.  These groups are also denoted by
$H^i(V\bmod\partial V,\Lambda)$.  The isomorphism (1.4.7.1) defines
an isomorphism
\begin{equation}\tag{1.4.7.2}
R\Gamma_c(V^\circ,\Lambda)
\xrightarrow{\ \sim\ }
R\Gamma_{\{0\}}R\Psi_\eta(\Lambda),
\end{equation}
and the \underline{variation} is defined, by the procedure of XIII,
2.4.5 and through (1.4.7.2), as a morphism whose source is the reduced
cohomology of $V$:
\begin{equation}\tag{1.4.7.3}
\operatorname{Var}:
R\widetilde\Gamma(V,\Lambda)
\longrightarrow R\Gamma_c(V^\circ,\Lambda).
\end{equation}

The automorphism induced by $T$ on the pair $(V,\partial V)$ is again
denoted by $T$, and its restriction $\partial T$ to $\partial V$ is
the identity.  The variation (1.4.7.3) should coincide with the
classical variation morphism (3.1.9) attached to such an automorphism
$T$.  A statement of this kind will be proved in 3.1.10.

\end{filecontents*}
% Original input: components/048_expose_XIV_section_2_1_through_2_5.tex
\begin{filecontents*}[overwrite,noheader]{sga7ii-cumulative-input-048.tex}
\section*{2. The comparison theorem}

\subsubsection*{2.1.}
For every scheme $X$ of finite type over $\mathbb C$, we denote by
$X_{\mathrm{cl}}$ the space $X(\mathbb C)$ endowed with its usual
topology.

Let $S$ be a smooth curve over $\mathbb C$, that is, a smooth
separated scheme of finite type over $\mathbb C$, pure of dimension
$1$.  Let $0\in S(\mathbb C)$, let $f:X\longrightarrow S$ be a
separated $S$-scheme of finite type, and let $\Lambda$ be a
noetherian torsion ring (for example $\mathbb Z/n\mathbb Z$).  Put
$X_0=f^{-1}(0)$.

\subsubsection*{2.2.}
We denote by $S^h$ the henselized henselian trait of $S$ at $0$, and
by $\eta$ the generic point of $S^h$.  There is a functor
\[
\begin{aligned}
\Psi:(\text{sheaves on }X)
&\xrightarrow{\ \mathrm{restriction}\ }
(\text{sheaves on }X\times_SS^h)\\
&\xrightarrow{\ \Psi\ }
(\text{sheaves on }X_0\times S^h).
\end{aligned}
\]
Its derived functor is the composite
\[
R\Psi:D(X,\Lambda)
\longrightarrow D(X\times_SS^h,\Lambda)
\xrightarrow{\ R\Psi\ }D(X_0\times S^h,\Lambda).
\]

\subsubsection*{2.3.}
Let $D\subset S_{\mathrm{cl}}$ be a disc containing $0$.  The
composite functor $\Psi_{\mathrm{cl}}$ (or simply $\Psi$)
\[
\begin{aligned}
\Psi_{\mathrm{cl}}:(\text{sheaves on }X_{\mathrm{cl}})
&\xrightarrow{\ \mathrm{restriction}\ }
(\text{sheaves on }X_{\mathrm{cl}}\times_{S_{\mathrm{cl}}}D)\\
&\xrightarrow{\ \Psi\ }
(\text{sheaves on }(X_0)_{\mathrm{cl}}\times\mathcal D)
\end{aligned}
\]
is independent, up to unique isomorphism, of the choice of $D$;
recall that $\mathcal D$ does not change when $D$ is replaced by a
smaller disc.  Its derived functor is the composite
\[
R\Psi_{\mathrm{cl}}:D(X_{\mathrm{cl}},\Lambda)
\longrightarrow D(f_{\mathrm{cl}}^{-1}(D),\Lambda)
\xrightarrow{\ R\Psi\ }
D((X_0)_{\mathrm{cl}}\times\mathcal D,\Lambda).
\]

\subsubsection*{2.4.}
Likewise, we define
\[
\begin{aligned}
R\Psi_\eta &:D(X-X_0,\Lambda)
  \longrightarrow D(X_0\times\eta,\Lambda),\\
R\Psi_{\eta\,\mathrm{cl}} &:D((X-X_0)_{\mathrm{cl}},\Lambda)
  \longrightarrow D((X_0)_{\mathrm{cl}}\times\mathcal D^*,\Lambda).
\end{aligned}
\]

\subsubsection*{2.5.}
The construction of 1.1.11 defines a morphism of topoi from
$\mathcal D^*$ to $\eta$, which extends to
\[
\varepsilon:\mathcal D\longrightarrow S^h.
\]
Together with
$\varepsilon:(X_0)_{\mathrm{cl}}\longrightarrow X_0$ (SGA 4, XI,
4), this morphism defines
\[
\varepsilon:(X_0)_{\mathrm{cl}}\times\mathcal D
\longrightarrow X_0\times S^h.
\]
If $\widetilde D^*$ is a universal covering of $D^*$ and
$k(\bar\eta)$ is the corresponding algebraic closure (1.1.10) of
$k(\eta)$, then
\[
\varepsilon^*(F_0,F_\eta,\alpha)
=\bigl(\varepsilon^*F_0,\varepsilon^*F_{\bar\eta},
\varepsilon^*\alpha\bigr).
\]

\end{filecontents*}
% Original input: components/049_expose_XIV_section_2_6.tex
\begin{filecontents*}[overwrite,noheader]{sga7ii-cumulative-input-049.tex}
\subsubsection*{2.6.}
Let $\widetilde D^*$ be a universal covering of $D^*$, let
$k(\bar\eta)$ be the corresponding algebraic closure of $k(\eta)$,
and let $\overline S$ be the normalization of $S^h$ in
$k(\bar\eta)$.  Write $\overline S$ as an inverse limit of affine
$S$-schemes $S_i$.  For every $i$, after replacing $D$ by a smaller
disc $D_i$, there is a natural morphism of $S_{\mathrm{cl}}$-spaces
\[
a_i:\widetilde D_i\longrightarrow(S_i)_{\mathrm{cl}}.
\]
Let $j_i:U_i\hookrightarrow S_i$ be the open complement in $S_i$ of
the image $0$ of the closed point of $\overline S$, and let
$i:0\hookrightarrow S_i$.

Let $F$ be a sheaf on $X$, and put $F_{\mathrm{cl}}=\varepsilon^*F$.
We again denote by $j_i$ the inclusion of $X\times_SU_i$ into
$X\times_SS_i$, by $i$ the inclusion of $X_0$ into
$X\times_SS_i$, and by $F$ the inverse image of $F$ on
$X\times_SS_i$.  Then
\[
\Psi(F)=
\left(i^*F,\ \varinjlim_i i^*j_{i*}j_i^*F,\
\varinjlim_i\varphi_i\right).
\]

The functors $\varepsilon^*$ and inverse image commute.  The morphism
\[
\varepsilon^*j_{i*}\longrightarrow(j_{i,\mathrm{cl}})_*\varepsilon^*
\]
therefore gives
\[
c_i:
(i^*j_{i*}j_i^*F)_{\mathrm{cl}}
\longrightarrow i^*j_{i*}j_i^*(F_{\mathrm{cl}}).
\]

Let $\overline j$ (respectively $\overline j_i$) be the inclusion of
$\widetilde D^*$ (respectively $\widetilde D_i^*$) into
$\widetilde D$ (respectively $\widetilde D_i$), as well as the
induced inclusions
\[
X\times_S\widetilde D^*\hookrightarrow X\times_S\widetilde D,
\qquad
X\times_S\widetilde D_i^*\hookrightarrow X\times_S\widetilde D_i.
\]
The adjunction map $\mathrm{Id}\longrightarrow a_{i*}a_i^*$ defines
\[
i^*j_{i*}j_i^*F_{\mathrm{cl}}
\longrightarrow
i^*\overline j_{i*}\overline j_i^*\overline F_{\mathrm{cl}}
=i^*\overline j_*\overline j^*\overline F_{\mathrm{cl}}.
\]
Passing to the limit gives
\[
c_\eta:\Psi_{\bar\eta}(F)_{\mathrm{cl}}
\longrightarrow\Psi_{\mathrm{cl}}(F_{\mathrm{cl}}).
\]
This morphism defines the \underline{comparison morphism}
\begin{equation}\tag{2.6.1}
c:\varepsilon^*\Psi(F)
\longrightarrow\Psi_{\mathrm{cl}}(\varepsilon^*F).
\end{equation}
Similarly, one defines
\begin{equation}\tag{2.6.2}
c_\eta:\varepsilon^*\Psi_\eta(F)
\longrightarrow\Psi_\eta(\varepsilon^*F).
\end{equation}

\end{filecontents*}
% Original input: components/050_expose_XIV_section_2_7_through_2_8.tex
\begin{filecontents*}[overwrite,noheader]{sga7ii-cumulative-input-050.tex}
\subsubsection*{2.7.}
The morphisms (2.6.1) and (2.6.2) derive to
\[
c:\varepsilon^*R\Psi
\longrightarrow R\Psi_{\mathrm{cl}}\circ\varepsilon^*
\]
and
\[
c_\eta:\varepsilon^*R\Psi_\eta
\longrightarrow R\Psi_{\eta\,\mathrm{cl}}\circ\varepsilon^*.
\]

\subsubsection*{Theorem 2.8.}
Let $K\in\operatorname{Ob}D(X,\Lambda)$ (respectively,
$K\in\operatorname{Ob}D(X-X_0,\Lambda)$) be a complex of sheaves of
$\Lambda$-modules whose cohomology sheaves $\underline H^i(K)$ are
constructible.  Then the comparison morphism
\[
c:\varepsilon^*R\Psi(K)
\longrightarrow R\Psi_{\mathrm{cl}}(\varepsilon^*K)
\]
(respectively,
\[
c:\varepsilon^*R\Psi_\eta(K)
\longrightarrow R\Psi_{\eta\,\mathrm{cl}}(\varepsilon^*K))
\]
is an isomorphism.

It suffices to prove the assertion for $R\Psi_\eta$.  Moreover, both
sides of 2.7 are triangulated functors of finite cohomological
dimension, so it is enough to treat the case in which $K$ is a single
sheaf $F$ placed in degree $0$.  The proof is exactly analogous to
that of the finiteness theorem XIII, 2.3.1.  Since $R\Psi$ and
$R\Psi_{\mathrm{cl}}$ have the same variance properties, the
d\'evissages used there (with permission to shrink $S$) reduce the
problem to the case in which $X$ is regular, $f^{-1}(0)$ is a normal
crossings divisor in $X$, and $F=\mathbb Z/\ell\mathbb Z$.  One then
compares Expos\'e I, 3.3 with its transcendental analogue.

\end{filecontents*}
% Original input: components/051_expose_XIV_section_3_1_1_through_3_1_3.tex
\begin{filecontents*}[overwrite,noheader]{sga7ii-cumulative-input-051.tex}
\section*{3. Isolated singularities}

\subsection*{3.1. The manifold of vanishing cycles}

\subsubsection*{3.1.1.}
Let
\[
D=\{z\in\mathbb C\mid |z|<1\},
\]
let $X$ be an analytic space with $0\in X$, and let
$f:X\to D$ be a morphism such that $f(0)=0$.  We suppose that $f$ is
smooth away from $0$ (in particular, $X$ is smooth away from $0$).
We intend to study the vanishing cycles near $0$; hence there is no
loss of generality in supposing that $X$ is embedded in
$\mathbb C^n\times D\subset\mathbb C^{n+1}$ compatibly with the
projection to $D$, and in such a way that $0$ maps to $0$.

\subsubsection*{3.1.2.}
Let $\delta(x)$ be a real-analytic function on $X$ (the restriction to
$X$ of a real-analytic function on $\mathbb C^{n+1}$) such that, for
$x\in X_0=f^{-1}(0)$,

\begin{enumerate}[label=(\alph*)]
\item $\delta(x)\geq 0$;
\item $\delta(x)=0$ if and only if $x=0$.
\end{enumerate}

For example, we may take $\delta$ to be the square of the Euclidean
distance from $x$ to $0$ in $\mathbb C^{n+1}$.  The critical locus of
the restriction of $\delta$ to $X_0$ is real analytic and therefore,
in a suitable neighborhood of $0$, has only finitely many connected
components.  Replacing $X$ by an open neighborhood of $0$ and applying
Sard's theorem, we may consequently obtain the following situation:

\begin{enumerate}[label=(\alph*),start=3]
\item for a suitable $r>0$, the space
  $X_0\cap\delta^{-1}([0,r])$ is compact;
\item on $X_0$, $\delta$ has no critical value in $(0,r]$.
\end{enumerate}

After shrinking $X$ further, we then have:

\begin{enumerate}[label=(\alph*),start=5]
\item there exists a $C^\infty$ function $\tau(s)$, $0\leq s\leq r$,
  with $\tau(0)=0$ and $\tau'(s)>0$ for $s>0$, such that, for
  $0<s\leq r$ and $0\leq |t|^2\leq\tau(s)$, the fiber
  $X_t=f^{-1}(t)$ is transverse to the real hypersurface
  $\delta(x)=s$ in $X$, and the box
  \[
  B_{t,d}=\{x\in X\mid \delta(x)\leq d\ \text{and}\ |f(x)|^2<t\}
  \]
  is proper over the open disc $D_t$ of radius $t^{1/2}$.
\end{enumerate}

Let $L=X_0\cap\delta^{-1}(r)$.  By (c) and (d), there is a
homeomorphism from $X_0\cap\delta^{-1}([0,r])$ to the cone
\[
L\times[0,r]\big/\bigl(L\times\{0\}\text{ contracted to a point}\bigr)
\]
which transforms $\delta$ into the second coordinate and is
$C^\infty$ away from the origin.

\subsubsection*{3.1.3.}
Let $d\leq r$ and $t\leq\tau(d)$.  Away from $0$, $B_{t,d}$ is a
manifold with corners, and the projection $f:B_{t,d}\to D_t$ is a
fibration.  In particular:

\begin{enumerate}[label=(\alph*)]
\item over $D_t^*=D_t-\{0\}$, $B_{t,d}$ induces a $C^\infty$ bundle
  of compact manifolds with boundary;
\item the boundary $\delta(x)=d$ of this bundle extends to a
  $C^\infty$ bundle $\partial B_{t,d}$ of compact manifolds over
  $D_t$.
\end{enumerate}

Let $y\in D_t^*$.  We call the manifold with boundary
$V=f^{-1}(y)$ the \underline{manifold of vanishing cycles}.  Since the
disc is contractible, the bundle $\partial B_{t,d}$ over $D_t$ is
isomorphic to the constant bundle with fiber $\partial V$, and the
space of bundle isomorphisms
\[
\alpha:\partial B_{t,d}\xrightarrow{\ \sim\ }\partial V\times D_t
\]
that restrict to the identity on the fiber over $y$ is contractible.
Choose one, denoted by $\alpha$.  Lift the vector field
$\partial/\partial\theta$ on $S_{|y|}$, the circle of radius $|y|$,
to a vector field $u$ on $f^{-1}(S_{|y|})$ which is compatible with
$\alpha$ on the boundary, so that
\[
d\alpha(u)=(0,\partial/\partial\theta).
\]
The space of such lifts is contractible.  Integrating $u$ once around
the circle in the positive direction defines a diffeomorphism $T$ of
$V$ onto itself which is the identity on $\partial V$; this is the
\underline{monodromy transformation}.  It is well defined up to
isotopy relative to $\partial V$.

In a sense whose precise formulation is left to the reader, $(V,T)$
is independent of the choices of $t$ and $d$.  I do not know whether
it is independent of the choice of $\delta$.  This is nevertheless
the case if we restrict $\delta$ to be the restriction to $X$ of a
real-analytic function $g$ on $\mathbb C^{n+1}$ such that, on the
Zariski tangent space of $X_0$ at $0$, one has $dg=0$ and $d^2g>0$
(positive definite).  Indeed, for $\delta$ and $\delta'$ of this
kind, $r$ can be chosen as in 3.1.2(c),(d), uniformly for every
$\lambda\delta+(1-\lambda)\delta'$, $0\leq\lambda\leq1$.

\end{filecontents*}
% Original input: components/052_expose_XIV_section_3_1_4_through_3_1_5.tex
\begin{filecontents*}[overwrite,noheader]{sga7ii-cumulative-input-052.tex}
\subsubsection*{3.1.4.}
Let $(V_1,T_1)$ be obtained from $(V,T)$ by adjoining a collar:
\[
V_1-V^\circ\simeq\partial V\times[0,1],
\qquad
T_1|_{V_1-V^\circ}=\operatorname{id}.
\]
Define the following space over $D_t$ (cf. 1.4.3):
\[
f':W\longrightarrow D_t.
\]

\begin{enumerate}[label=(\alph*)]
\item Let $\widetilde D_t^*$ be a universal covering of $D_t^*$ with
  base point $y$, and let $W^*$ be the bundle over $D_t^*$ obtained as
  the quotient of $V_1\times\widetilde D_t^*$ by the equivalence
  relation generated by
  \[
  (v,y)\sim(T_1v,e^{2\pi i}y).
  \]
\item Since $T_1|_{V_1-V^\circ}$ is trivial, $W^*$ can be glued to
  the trivial bundle over $D_t$ with fiber $V_1-V$, giving a space
  $W_1$ over $D_t$.
\item Put $W=W_1\cup\{0\}$ and $f'(0)=0$.  Give $W$ the unique
  topology inducing that of $W_1$, for which $f'$ is continuous and,
  for $t'<t$, ${f'}^{-1}(\overline D_{t'})$ is the one-point
  compactification of
  $W_1\cap {f'}^{-1}(\overline D_{t'})$.
\end{enumerate}

Equip $W-\{0\}$ with the integrable connection of 1.3.4 for which
the trivialization, supplied by (a), of the inverse image of $W$ on
$\widetilde D^*$ is horizontal.

\subsubsection*{Proposition 3.1.5.}
There exists a homeomorphism
\[
\Psi:B_{t,d}\xrightarrow{\ \sim\ }W
\]
compatible with the projections to $D_t$ and carrying $0$ to $0$;
this homeomorphism may be chosen to be a diffeomorphism away from
$0$.

This result will not be used in the remainder of the seminar.

Let $d_0<d_1\leq r$, let $t_0\leq\tau(d_0)$, and let
$y\in D_{t_0}^*$.  We have
\[
B_{t_0,d_0}\subset B_{t_0,d_1},
\]
and $B_{t_0,d_1}-B_{t_0,d_0}^{\circ}$ is a bundle over $D_{t_0}$.
Trivialize it and, over the circle $S_{|y|}$, lift
$\partial/\partial\theta$ compatibly with this trivialization.  Let
$T_1$ be the corresponding endomorphism of
\[
V_1=f^{-1}(y)\cap B_{t_0,d_1},
\qquad
V_0=f^{-1}(y)\cap B_{t_0,d_0}.
\]
Then $(V_0\subset V_1,T_1)$ is isomorphic to
$(V\subset V_1,T_1)$ of 3.1.4.  Let $W$ be constructed from
$(V_0\subset V_1,T_1)$.  We have ${f'}^{-1}(y)=V_1$ canonically.
The proof will show that $\Psi$ can be chosen to induce the identity
\[
f^{-1}(y)\xrightarrow{\ \sim\ }{f'}^{-1}(y).
\]

The following drawings of $f$ and of its restriction to
$f^{-1}(\mathbb R)$ summarize the proof.

\begin{center}
\begin{tikzpicture}[x=1cm,y=1cm,line width=.45pt]
  % Total space and base disc.
  \draw (-4.2,2.6) arc[start angle=180,end angle=360,x radius=1.05,y radius=.28];
  \draw[dashed] (-2.1,2.6) arc[start angle=0,end angle=180,x radius=1.05,y radius=.28];
  \draw (-4.2,2.6)--(-4.2,.15);
  \draw (-2.1,2.6)--(-2.1,.15);
  \draw (-4.2,.15) arc[start angle=180,end angle=360,x radius=1.05,y radius=.28];
  \draw[dashed] (-2.1,.15) arc[start angle=0,end angle=180,x radius=1.05,y radius=.28];
  \draw (-3.15,2.75)--(-3.15,.05);
  \fill (-3.15,1.35) circle (1.8pt);
  \node[left] at (-3.15,1.35) {$0$};
  \draw (-3.65,-1.45) circle (.58);
  \fill (-3.65,-1.45) circle (1.2pt);
  \node[right] at (-3.08,-1.78) {$D_t$};

  % Real slice and its base.
  \draw (1.2,.1) rectangle (4.0,2.75);
  \foreach \yy in {.28,.48,2.36,2.56}{\draw (1.2,\yy)--(4.0,\yy);}
  \draw (1.2,2.15) .. controls (2.05,2.10) and (2.25,1.75) .. (2.60,1.48)
                    .. controls (2.95,1.75) and (3.18,2.10) .. (4.0,2.15);
  \draw (1.2,.72) .. controls (2.05,.76) and (2.28,1.15) .. (2.60,1.48)
                   .. controls (2.93,1.15) and (3.17,.76) .. (4.0,.72);
  \draw (1.2,1.92)--(2.60,1.48)--(4.0,1.92);
  \draw (1.2,1.05)--(2.60,1.48)--(4.0,1.05);
  \draw (1.2,1.48)--(4.0,1.48);
  \fill (2.60,1.48) circle (1.8pt);
  \draw (1.15,-1.45)--(4.05,-1.45);
  \fill (2.60,-1.45) circle (1.8pt);
\end{tikzpicture}
\end{center}

Let $B^+\subset B_{t_0,d_1}$ be the set of points $x$ such that
\[
|f(x)|^2\leq\tau(\delta(x)).
\]
Let $\varphi_n$ be a sequence of $C^\infty$ functions from
$[0,t_0]$ to $(0,d_1]$ such that, for a sequence
$t_0>r_1>r_2>\cdots$ tending to $0$,

\begin{enumerate}[label=(\alph*)]
\item for $|t|^2\geq r_i$, $\varphi_i$ is the inverse function of
  $\tau$;
\item for $|t|^2<r_i$ and $j>i$,
  $\varphi_j(t)<\varphi_i(t)$.
\end{enumerate}

\begin{center}
\begin{tikzpicture}[x=1cm,y=1cm,>=Latex,line width=.5pt]
  \draw[->] (0,0)--(4.5,0) node[right] {$t$};
  \draw[->] (0,0)--(0,3.4) node[above] {$d$};
  \draw (0,3.0)--(3.05,3.0)--(3.05,0);
  \node[left] at (0,3.0) {$d_1$};
  \node[below] at (3.05,0) {$t_0$};
  \draw (0,0) .. controls (.35,1.05) and (1.35,2.15) .. (3.05,2.45);
  \node[right] at (3.05,2.45) {$t=\tau(d)$};
  \draw (0,1.45) .. controls (.45,1.42) and (.80,1.40) .. (1.18,1.68)
        .. controls (1.45,1.88) and (1.68,2.00) .. (1.92,2.10);
  \node at (.55,1.72) {$\varphi_0$};
  \draw (0,.58) .. controls (.23,.56) and (.35,.65) .. (.50,.90)
        .. controls (.67,1.18) and (.84,1.40) .. (1.18,1.68);
\end{tikzpicture}
\end{center}

The space $B^+-\{0\}$ is the union of the bundles $B_i^+$ over
$D_t$ defined by
\[
\delta(x)\geq\varphi_i(|f(x)|^2).
\]
A standard argument, using the fact that $\delta$ has no critical
points on the fibers of $B^+-\{0\}\to D_t$, shows that there is a
diffeomorphism $\alpha_i:B_i^+\to B_{i+1}^+$, compatible with the
projection to $D_t$ and equal to the identity outside a small
neighborhood of the subspace
\[
\delta(x)=\varphi_i(|f(x)|^2),
\qquad |f(x)|^2\leq r_i,
\]
of $B_i^+$.  Composing the $\alpha_i$ and passing to the limit gives
a diffeomorphism
\[
\alpha:B_0^+-\{x\mid f(x)=0,\ \delta(x)=\varphi_0(0)\}
\xrightarrow{\ \sim\ }B^+-\{0\}.
\]
Transport a trivialization of the bundle $B_0^+$ over $D_t$ by
$\alpha$ to obtain an integrable connection $\nabla'$ on
$B^+-\{0\}$.  We may then:

\begin{enumerate}[label=(\alph*)]
\item extend the restriction of this connection to
  $f^{-1}(S_{|y|})\cap B^+$ to a connection on
  $f^{-1}(S_{|y|})/S_{|y|}$; this extension gives $T_1$;
\item extend to $B_{t_0,d_1}$ the lift, supplied by $\nabla'$ on
  $B^+-\{0\}$, of the radial vector field on $D_t$.
\end{enumerate}

There is then a unique integrable connection $\nabla$ on
$B_{t_0,d_1}$ extending $\nabla'$, (a), and (b), and a unique
horizontal diffeomorphism from $B_{t_0,d_1}-\{0\}$ to $W-\{0\}$
that is the identity on $f^{-1}(y)$.  This is the required
diffeomorphism.

\end{filecontents*}
% Original input: components/053_expose_XIV_section_3_1_6_through_3_1_8.tex
\begin{filecontents*}[overwrite,noheader]{sga7ii-cumulative-input-053.tex}
\subsubsection*{3.1.6. A generalization.}
The preceding discussion generalizes to the case in which
$f:X\to D$ is a morphism such that $0\in X$ is an isolated
non-smooth point \underline{in} $X_0$.  After shrinking $X$ and $D$,
the non-smooth locus of $f$ is then finite over $D$.  The manifold of
vanishing cycles $V$ has a finite set $\Sigma$ of critical points;
$V-\Sigma$ is a differentiable manifold with boundary (its boundary
is compact, although the manifold is not).  The monodromy
transformation $T$ is the identity on the boundary and stabilizes
$\Sigma$.  Finally, Proposition 3.1.5 remains valid: the fibered
homeomorphism $\Psi$ is a diffeomorphism away from the critical locus,
which it carries to the part of $W$ arising from $\Sigma$.

\subsubsection*{3.1.7.}
Proposition 3.1.5 allows us to use the calculations of 1.4.  We obtain:

\begin{enumerate}[label=(\Alph*)]
\item The vanishing-cycle groups $(R^i\Phi(\Lambda))_0$ are the
  reduced cohomology groups, with coefficients in $\Lambda$, of the
  manifold of vanishing cycles $V$.
\item The groups
  $H^i_{\{0\}}(X_0,R\Psi(\Lambda))$ are the cohomology groups with
  proper support of $V^\circ$, with coefficients in $\Lambda$.
\item The variation
  \[
  \operatorname{Var}:\widetilde H^i(V,\Lambda)
  \longrightarrow H_c^i(V^\circ,\Lambda)
  =H^i(V\bmod\partial V,\Lambda)
  \]
  is the classical variation defined by the automorphism $T$ of $V$,
  which is the identity on $\partial V$.
\end{enumerate}

\subsubsection*{3.1.8.}
Assertions (A), (B), and (C) are in fact more elementary than
Proposition 3.1.5.  Here is a direct proof and construction which
spares us from checking that the isomorphisms in (A) and (B) are
independent of the choices made in the construction of 3.1.5.

Let $d\leq r$ and $t\leq\tau(d)$, and let $B$ be the open box
\[
B=\{x\mid\delta(x)<d\ \text{and}\ |f(x)|^2<t\}.
\]
Let $\widetilde D^*$ be a universal covering of $D^*$, and consider
the diagram
\begin{equation}\tag{3.1.8.1}
\begin{tikzcd}[column sep=large,row sep=large]
\{0\} \arrow[r,hook,"\bar\imath"] \arrow[d,equal]
  & \widetilde D_t \arrow[d,"p"]
  & \widetilde D_t^* \arrow[l,hook',"\bar\jmath"'] \arrow[d,"p'"] \\
\{0\} \arrow[r,hook,"i"]
  & D_t
  & D_t^* \arrow[l,hook',"j"']
\end{tikzcd}
\end{equation}
obtained from (1.2.7.2) by restriction to $D_t$.  Let
\begin{equation}\tag{3.1.8.2}
\begin{tikzcd}[column sep=large,row sep=large]
B_0 \arrow[r,hook,"\bar\imath"] \arrow[d,equal]
  & \overline B \arrow[d,"p"]
  & \overline B^* \arrow[l,hook',"\bar\jmath"'] \arrow[d,"p'"] \\
B_0 \arrow[r,hook,"i"]
  & B
  & B^* \arrow[l,hook',"j"']
\end{tikzcd}
\end{equation}
be the inverse image of $f:B\to D$ over (3.1.8.1).  We again write
$f$ for the projection from (3.1.8.2) to (3.1.8.1).

\medskip
\noindent\textup{(A)}
By definition, $R^i\Psi_\eta(\Lambda)_0$ is the direct limit, as $U$
runs through the neighborhoods of $0$ in $X$, of the cohomology groups
\[
H^i(U\times_D\widetilde D^*,\Lambda).
\]
Let $U$ run through the fundamental system of neighborhoods of $0$
formed by boxes $B$ as above.  The fiber product
$U\times_D\widetilde D^*=\overline B^*$ is a bundle with fiber
$V^\circ$ over the contractible space $\widetilde D_t^*$.  Hence the
direct system of the groups $H^i(\overline B^*,\Lambda)$ is constant,
with value
\[
H^i(V^\circ,\Lambda)=H^i(V,\Lambda).
\]
The resulting isomorphisms
\begin{equation}\tag{3.1.8.3}
R^i\Psi_\eta(\Lambda)_0\simeq H^i(V,\Lambda)
\end{equation}
may be refined to the following composite isomorphism (cf. (B)
below):
\begin{equation}\tag{3.1.8.4}
\begin{tikzcd}[column sep=small,row sep=large,cells={nodes={font=\small}}]
R\Psi_\eta(\Lambda)_0
  & R\Gamma(B_0,R\Psi_\eta(\Lambda)) \arrow[l,"\sim"']
  & (Rf_*R\bar\jmath_*\Lambda)_0 \arrow[l,"\sim"']
  & R\Gamma(\widetilde D_t,Rf_*R\bar\jmath_*\Lambda)
      \arrow[l,"\sim"'] \arrow[d,equal] \\
{}
  & R\Gamma(V,\Lambda) \arrow[r,"\sim"]
  & R\Gamma(V^\circ,\Lambda)
  & R\Gamma(\overline B^*,\Lambda) \arrow[l,"\sim"']
\end{tikzcd}
\end{equation}
The complex $R\Phi(\Lambda)_0$ is defined by the distinguished
triangle
\[
\Lambda\longrightarrow(R\Psi_\eta\Lambda)_0
\longrightarrow R\Phi(\Lambda)_0\longrightarrow{}.
\]
Consequently, the groups $R^i\Phi(\Lambda)_0$ are the reduced
cohomology groups of $V$.

\medskip
\noindent\textup{(B)}
Let $\mathfrak F$ be the family of closed subsets of $\overline B^*$
whose closure in $\overline B$ is proper over $\widetilde D$.  The
arrows in the following diagram are isomorphisms:
\begin{equation}\tag{3.1.8.5}
\begin{tikzcd}[column sep=small,row sep=large,cells={nodes={font=\small}}]
R\Gamma_{\{0\}}(R\Psi_\eta(\Lambda))
  \arrow[r,"\sim","\text{\textcircled{1}}"']
  & R\Gamma_c(B_0,R\Psi_\eta(\Lambda))
  & (Rf_!R\bar\jmath_*\Lambda)_0
      \arrow[l,"\sim"',"\text{\textcircled{2}}"] \\
R\Gamma_c(V^\circ,\Lambda)
  & R\Gamma_{\mathfrak F}(\overline B^*,\Lambda)
      \arrow[l,"\sim"',"\text{\textcircled{5}}"]
      \arrow[r,equal,"\text{\textcircled{4}}"']
  & R\Gamma(\widetilde D_t,Rf_!R\bar\jmath_*\Lambda)
      \arrow[u,"\text{\textcircled{3}}","\wr"']
\end{tikzcd}
\end{equation}
Indeed, \textcircled{1} is an isomorphism because $B_0$ is a cone
(3.1.2) and $R^i\Psi_\eta(\Lambda)$ is constant away from $0$;
\textcircled{2} by the proper base-change theorem;
\textcircled{3} because $Rf_!R\bar\jmath_*\Lambda$ is constant away
from $0$; \textcircled{4} is the derived-category version of the
spectral sequence of a composite functor; and \textcircled{5} because
$\widetilde D^*$ is contractible.

The isomorphism (3.1.8.5) supplies the desired isomorphisms
\begin{equation}\tag{3.1.8.6}
R^i\Gamma_{\{0\}}(R\Psi_\eta(\Lambda))
\simeq H_c^i(V^\circ,\Lambda).
\end{equation}
We shall verify (C) in 3.1.10.

\end{filecontents*}
% Original input: components/054_expose_XIV_section_3_1_9_through_3_1_10.tex
\begin{filecontents*}[overwrite,noheader]{sga7ii-cumulative-input-054.tex}
\subsubsection*{3.1.9.}
Cf. [3].  Let $X$ be a topological space, let
$i:Y\hookrightarrow X$ be a closed subset, let $T$ be an
automorphism of $X$ inducing the identity on $Y$, and let $\Lambda$ be
an abelian group.  To define the variation $\operatorname{Var}$,
choose a $T$-equivariant resolution $\Lambda_X^*$ of $\Lambda$ on
$X$, a resolution $\Lambda_Y^*$ of the constant sheaf $\Lambda$ on
$Y$ on which, by convention, $T$ acts trivially, and an equivariant
morphism
\[
\varphi:\Lambda_X^*\longrightarrow i_*\Lambda_Y^*
\]
resolving the evident morphism
$\varphi_0:\Lambda_X\to i_*\Lambda_Y$.  Suppose that the sheaves
$\Lambda_X^i$ and $\Lambda_Y^i$ are acyclic for the functors to be
considered.  Passing to global sections gives
\[
\Gamma(\varphi):\Gamma(\Lambda_X^*)
\longrightarrow\Gamma(\Lambda_Y^*).
\]
For simplicity, suppose that every $\Gamma(\varphi)^i$ is surjective;
otherwise, make it surjective by adding the cone of
$\Gamma(\Lambda_Y^*)$ to $\Gamma(\Lambda_X^*)$.  Put
\[
R\Gamma(X\bmod Y,\Lambda)=\operatorname{Ker}(\Gamma(\varphi)).
\]
Under reasonable hypotheses, if $\mathfrak F$ is the family of closed
subsets of $X$ disjoint from $Y$, there is a quasi-isomorphism
\[
\Gamma_{\mathfrak F}(X-Y,\Lambda_X^*)
\xrightarrow{\ \sim\ }R\Gamma(X\bmod Y,\Lambda).
\]
The endomorphism $T-1$ of $\Gamma(\Lambda_X^*)$, where $T$ acts by
transport of structure, is zero on $\Lambda$ (the constants) and
factors through $\operatorname{Ker}(\Gamma(\varphi))$.  Hence, in the
derived category and under the same hypotheses, it gives a variation
morphism
\[
\operatorname{Var}:R\widetilde\Gamma(X,\Lambda)
\longrightarrow R\Gamma_{\mathfrak F}(X-Y,\Lambda).
\]
Without these hypotheses, one obtains instead
\[
\operatorname{Var}:R\widetilde\Gamma(X,\Lambda)
\longrightarrow R\Gamma(X\bmod Y,\Lambda).
\]

This description simplifies when $T$ is the identity on a
neighborhood of $Y$ and acts as the identity on $\Lambda_X^*$ in that
neighborhood.  The endomorphism $T-1$ of $\Gamma(\Lambda_X^*)$ then
factors through $\Gamma_{\mathfrak F}(\Lambda_X^*)$ and supplies the
variation.

\subsubsection*{Proposition 3.1.10.}
Under the isomorphisms of 3.1.8(A),(B), the variation
\[
(R\Phi_\eta(\Lambda))_0
\longrightarrow R\Gamma_{\{0\}}(R\Psi_\eta(\Lambda))
\]
is equal to the variation of 3.1.9 defined by the automorphism $T$ of
$V$, which is the identity on $\partial V$.

The reader is invited to draw a small picture rather than read the
proof.  Let ${}_fB$ be the closed box of 3.1.2 and let
${}_fC\subset{}_fB$ be a collar along the boundary; thus
${}_fC\simeq\partial_fB\times[0,1]$ as a bundle over $D_t$, through
an isomorphism carrying $\partial_fB\times\{1\}$ identically onto
$\partial_fB$.  Let
\[
{}_f\alpha:{}_fC\xrightarrow{\ \sim\ }
(\partial V\times[0,1])\times D_t
\]
be a trivialization of the bundle ${}_fC$ over $D_t$ whose fiber at
$y\in D_t$ induces the identity
\[
\partial V=\partial(f^{-1}(y))
\xrightarrow{\ \sim\ }\partial V\times\{1\}.
\]
Put $C={}_fC\cap B$ and let
\[
\alpha:C\xrightarrow{\ \sim\ }
(\partial V\times[0,1))\times D_t
\]
be induced by ${}_f\alpha$.  We retain the notation of (3.1.8.2).

Let $R$ be the equivalence relation on $\overline B$ defined by
$xRy$ if either $x=y$, or $p(x),p(y)\in C$ and
$\alpha(p(x))$, $\alpha(p(y))$ have the same image in
$\partial V\times[0,1)$.  Put
\[
{}_1\overline B=\overline B/R,
\qquad
{}_1\overline B^*=\overline B^*/R,
\qquad
{}_1\overline C^*=p^{-1}(C)/R.
\]
There is a commutative diagram
\[
\begin{tikzcd}[column sep=large,row sep=large]
B_0 \arrow[r,hook,"\bar\imath"] \arrow[d,equal]
  & \overline B \arrow[d,"q"]
  & \overline B^* \arrow[l,hook',"\bar\jmath"'] \arrow[d,"q'"] \\
B_0 \arrow[r,hook,"\bar\imath"]
  & {}_1\overline B
  & {}_1\overline B^* \arrow[l,hook',"{}_1\bar\jmath"']
\end{tikzcd}
\]
Let $\Lambda^*$ be a $\pi_1(D^*)$-equivariant resolution of
$\Lambda$ by acyclic sheaves on $\overline B^*$.  By definition, the
variation is computed from the morphism of complexes
\[
\Lambda[0]\longrightarrow
\bar\imath^*\bar\jmath_*(\Lambda^*).
\]
Since $\Lambda\xrightarrow{\sim}Rq_*\Lambda$ and
$\Lambda\xrightarrow{\sim}Rq'_*\Lambda$, it may equally be computed
using an equivariant resolution of $\Lambda$ on
${}_1\overline B^*$ and the morphism
\[
\Lambda[0]\longrightarrow
\bar\imath^*{}_1\bar\jmath_*(\Lambda).
\]
Under (A) and (B), the variation is thereby identified with the
composite morphism
\begin{equation}\tag{3.1.10.1}
\begin{tikzcd}[column sep=small,row sep=large,cells={nodes={font=\small}}]
R\Gamma(V,\Lambda) \arrow[r,"\sim"]
  & R\Gamma(V^\circ,\Lambda)
  & R\Gamma({}_1\overline B^*,\Lambda)
      \arrow[l,"\sim"',"\text{\textcircled{1}}"]
      \arrow[d,"\operatorname{Var}"] \\
{}
  & R\Gamma_c(V^\circ,\Lambda)
  & R\Gamma({}_1\overline B^*\bmod{}_1\overline C^*,\Lambda)
      \arrow[l,"\sim"',"\text{\textcircled{2}}"]
\end{tikzcd}
\end{equation}
where \textcircled{1} and \textcircled{2} are induced by the inclusion
of $V^\circ$ in ${}_1\overline B^*$, and where
$\operatorname{Var}$ is the variation of 3.1.9 for
${}_1\overline B^*$ and the action of the positive generator of
$\pi_1(D^*)$.  This group acts on ${}_1\overline B^*$ through the
action described in 1.1.5 and acts trivially on
${}_1\overline C^*$.

Choose a connection $\nabla$ on $B^*/D^*$ compatible with the
trivialization of $C$.  This connection trivializes $\overline B^*$
and, on passing to the quotient, defines a map
$u:{}_1\overline B^*\to V$.  For $T$ defined as in 3.1.3, the diagram
\[
\begin{tikzcd}[column sep=large,row sep=large]
{}_1\overline B^* \arrow[r,"u"] \arrow[d,"T"']
  & V \arrow[d,"T"] \\
{}_1\overline B^* \arrow[r,"u"]
  & V
\end{tikzcd}
\]
commutes.  The map $u$ induces isomorphisms in cohomology, so
Proposition 3.1.10 follows from (3.1.10.1) and the functoriality of
variation.

Proposition 3.1.10 is meant to justify an explicit calculation to be
made later in which, as here, $T$ is the identity on a neighborhood
of $\partial V$.  We shall take a cohomology class $c$ represented by
a current $C$ and use the fact that $\operatorname{Var}(c)$ is
represented by the current $TC-C$, compactly supported in $V^\circ$.

\end{filecontents*}
% Original input: components/055_expose_XIV_section_3_1_11_through_3_1_13.tex
\begin{filecontents*}[overwrite,noheader]{sga7ii-cumulative-input-055.tex}
\subsubsection*{3.1.11.}
In practice, $T$ is often calculated in two steps.

\begin{enumerate}[label=(\alph*)]
\item Fix $t$ and $d$ with $t\leq\tau(d)$, choose
  $y\in D_t^*$, and choose a connection $\nabla$ on
  $f^{-1}(S_{|y|})\cap B_{t,d}$.  This connection defines a
  diffeomorphism $T_0$ of
  \[
  V_0=f^{-1}(y)\cap B_{t,d}
  \]
  onto itself, which need not be the identity on the boundary.
\item Choose a trivialization of $\partial B_{t,d}$ over $D_t$.
  Integrating $\nabla$ gives a diffeomorphism $\varphi(\theta)$ from
  $\partial V$ to
  \[
  f^{-1}(e^{2\pi i\theta}y)\cap\partial B_{t,d},
  \]
  which is identified with $\partial V$ by the trivialization.  For
  $0\leq\theta\leq1$, $\varphi(\theta)$ is a path in
  $\operatorname{Diff}(\partial V)$ joining the identity to
  $\partial T_0$.
\end{enumerate}

The resulting pair $(T_0,\varphi)$ depends on the arbitrary choices
made.  Since the space of these choices is contractible, it is
independent of them up to homotopy.  For a transformation $T$ of the
kind considered in 3.1.3, $(T,\text{constant path})$ is a pair of the
kind considered here.

\subsubsection*{3.1.12.}
The map
\[
\operatorname{Diff}(V)\longrightarrow\operatorname{Diff}(\partial V)
\]
is a fibration.  Consequently, the assignment
$T\mapsto(T,\text{constant path})$ induces a bijection between:

\begin{enumerate}[label=\alph*)]
\item diffeomorphisms $T:V\to V$ satisfying
  $\partial T=\operatorname{id}$, up to isotopy; and
\item pairs $(T_0,\varphi)$, where $T_0:V\xrightarrow{\sim}V$ and
  $\varphi$ is a homotopy class of paths from
  $\operatorname{id}$ to $\partial T_0$, up to isotopy.
\end{enumerate}

A pair $(T_0,\varphi)$ calculated by 3.1.11(a),(b) therefore
determines $T$.  Concretely, calculate $T$ from $(T_0,\varphi)$ as
follows.

\medskip
\noindent\textup{($\alpha$)} Choose a collar embedding
  \[
  [0,1]\times\partial V\xhookrightarrow{\ \alpha\ }V,
  \qquad \alpha(1,x)=x,
  \]
  and a representative
  $\varphi:[0,1]\to\operatorname{Diff}(\partial V)$ such that its
  extension by $\operatorname{id}$ for negative parameters is
  $C^\infty$.  On the collar put
  \[
  \beta(x)=\operatorname{pr}_1\alpha^{-1}(x).
  \]
  Thus $\beta$ has no critical point, takes values in $[0,1]$, and
  equals $1$ on $\partial V$.

\medskip
\noindent\textup{($\beta$)} Take $T=T_0$ away from the collar and, on the collar, set
  \[
  T(x)=T_0\bigl(\varphi(\beta(x))^{-1}(x)\bigr).
  \]

\subsubsection*{3.1.13.}
Here is a variant of 3.1.11(b) for calculating the homotopy class of
the path $\varphi$.  On every ray $\arg(z)=\theta$, choose a
trivialization of the bundle $\partial B_{t,d}$ which varies
continuously with $\theta$.  Use the trivialization over the ray
through $y$ to identify
\[
(\partial V)_0=f^{-1}(0)\cap\partial B_{t,d}
\]
with $\partial V$.  Let
\[
\Psi_\theta:V\longrightarrow
f^{-1}(e^{i\theta}y)\cap B_{t,d}
\]
be obtained by integrating $\nabla$.  Transporting
$\partial\Psi_\theta$ along the rays gives a diffeomorphism of
$(\partial V)_0$, hence of $\partial V$, onto itself.  This is
$\varphi(\theta)$.
\end{filecontents*}
% Original input: components/056_expose_XIV_section_3_2_1_through_3_2_5.tex
\begin{filecontents*}[overwrite,noheader]{sga7ii-cumulative-input-056.tex}
\subsection*{3.2. The Picard--Lefschetz formula (transcendental case)}

\subsubsection*{3.2.1.}
We shall make the results of 3.1 explicit for the map
\[
f=\sum_{i=1}^{n+1}z_i^2:\mathbb C^{n+1}\longrightarrow\mathbb C.
\]
For $\delta$ take the function $\sum |z_i|^2$, and take the box
$B=B_{a,b}$ defined by
\[
\left|\sum z_i^2\right|\leq a^2,
\qquad
\sum |z_i|^2\leq b^2,
\]
where $0<a<b$.  It is only to simplify the notation that we use a
closed box rather than a box as in 3.1.2.

\subsubsection*{3.2.2.}
The manifold of vanishing cycles $V$ has equations
\begin{equation}\tag{3.2.2.1}
\begin{array}{rcl}
\displaystyle\sum z_i^2&=&a^2,\\
\displaystyle\sum |z_i|^2&\leq&b^2.
\end{array}
\end{equation}
Write $z=(z_i)$ in terms of its real and imaginary parts,
$z=x+iy$.  The preceding equations become
\begin{equation}\tag{3.2.2.2}
\begin{array}{rcl}
x^2-y^2&=&a^2,\\
x\mathbin{\cdot}y&=&0
  \qquad\bigl(\text{i.e. }x\perp y\text{ in }\mathbb R^{n+1}\bigr),\\
x^2+y^2&\leq&b^2.
\end{array}
\end{equation}
\footnote{The printed source has $\mathbb R^n$ in (3.2.2.2), while
$z$ has $n+1$ coordinates.  The English text corrects this evident
dimension typo to $\mathbb R^{n+1}$.}
Introduce on $V$ the new coordinate system
\[
X=\frac{x}{|x|},
\qquad
Y=y\left(\frac{b^2-a^2}{2}\right)^{-1/2}.
\]
Then $x=X(a^2+y^2)^{1/2}$, and (3.2.2.2) becomes
\begin{equation}\tag{3.2.2.3}
\begin{array}{rcl}
X^2&=&1,\\
X\mathbin{\cdot}Y&=&0,\\
Y^2&\leq&1.
\end{array}
\end{equation}
This identifies $V$ with the tangent ball bundle of the sphere $S^n$.

\subsubsection*{3.2.3.}
The map
\[
\theta\longmapsto \exp(2\pi i\theta)a
\]
identifies $\mathbb R$ with the universal covering of the circle
$S_a$ in $\mathbb C$ having radius $a$ and center $0$.  The pullback
to $\mathbb R$ of the bundle $B$ is isomorphic to $\mathbb R\times V$.
We trivialize it by means of the isomorphism
\[
v\longmapsto \exp(\pi i\theta)v
\]
from $V$ to the fiber of $B$ over $\exp(2\pi i\theta)a$.  Let $T_0$
be the map furnished by this trivialization as in 3.1.11.  In the
coordinates (3.2.2.3),
\[
T_0(X,Y)=(-X,-Y).
\]

\subsubsection*{3.2.4.}
The bundle over the disc $D_a$ of radius $a$ which forms the boundary
of the ``bundle'' $B$ has equations
\[
\sum |z_i|^2=b^2,
\qquad
f(z)\in D_a,
\]
the projection to the disc being $f$.  We shall describe a
trivialization of the inverse image of this bundle over each ray
\[
e^{2\pi i\theta}\mathbb R^+\cap D_a.
\]

Over the ray $\mathbb R^+\cap D_a$, take the isomorphism with
$(\mathbb R^+\cap D_a)\times\partial V$ which, in the coordinates
(3.2.2.3) on $\partial V$, is given by
\[
z=x+iy\longmapsto(f(z),X,Y)
  =\left(f(z),\frac{x}{|x|},\frac{y}{|y|}\right).
\]
Over $(e^{2\pi i\theta}\mathbb R^+)\cap D_a$, use the image of this
trivialization under multiplication by $e^{\pi i\theta}$.  Apply
3.1.13 and use on $\partial V$ the coordinates $(X,Y)$ satisfying
\begin{equation}\tag{3.2.4.1}
X^2=Y^2=1,
\qquad
X\mathbin{\cdot}Y=0.
\end{equation}
If $\varphi(\theta)$ takes $(X,Y)$ to $(X_\theta,Y_\theta)$, one finds
\begin{equation}\tag{3.2.4.2}
X_\theta+iY_\theta
=e^{\pi i\theta}(X+iY)
\qquad(0\leq\theta\leq1).
\end{equation}

\subsubsection*{3.2.5.}
In summary, the manifold of vanishing cycles $V$ is the tangent ball
bundle of $S^n$.  The monodromy transformation is represented by
$(T_0,\varphi)$, where $T_0$ is the differential of the antipodal map
and $\varphi$ is the following path from $\operatorname{Id}$ to
$\partial T_0$: under $\varphi(\theta)$, each unit tangent vector is
parallel-transported along the geodesic which it determines through a
distance $+\pi\theta$ (the plus sign means ``in the direction in which
the vector points'').  At $\theta=1$ one reaches the antipodal vector.

By 3.1.13 this gives the following description of $T$.  Let
$\varphi:[0,1]\to[0,1]$ be a $C^\infty$ map which is $0$ near $0$ and
$1$ near $1$.  The map $T$ moves $u$, along the geodesic which $u$
determines, through the distance
\[
-\pi\bigl(1-\varphi(|u|)\bigr).
\]
Near $|u|=1$ this is the identity; near $|u|=0$ it is the antipodal
map.
\end{filecontents*}
% Original input: components/057_expose_XIV_section_3_2_7_through_3_2_9.tex
\begin{filecontents*}[overwrite,noheader]{sga7ii-cumulative-input-057.tex}
\subsubsection*{3.2.7.}
Let $W$ be a differentiable manifold of dimension $d$, and let
\[
p:T(W)\longrightarrow W
\]
be its tangent bundle, regarded as a manifold.  If
$(x_1,\ldots,x_d)$ is a local coordinate system on $W$, then
\[
(x_1,\ldots,x_d,dx_1,\ldots,dx_d)
\]
is a local coordinate system on $T(W)$.  Let $\varepsilon'$ be the
orientation of $T(W)$ defined by this coordinate system.  It is
independent of the choice of the $x_i$, and is therefore globally
defined.

Suppose that $W$ is orientable.  The submanifold $W\subset T(W)$ (the
zero section), endowed with an orientation (respectively, with an
orientation of its normal bundle), defines an element $w$ of
$H_d(T(W))$ (respectively, of $H_c^d(T(W))$).  The self-intersection
of $W\subset T(W)$ is the Euler--Poincare characteristic:
\[
\int w\wedge w=\chi(W)
\]
for the orientation $\varepsilon'$.

\footnote{The printed sentence mixes three dimension symbols: it calls
the dimension $x$, indexes the coordinates by $n$, and places the
classes in degree $d$.  The English text consistently uses $d$ in the
general statement; the specialization below uses $d=n$.}

For the manifold of vanishing cycles $V$, its complex structure
defines an orientation $\varepsilon$.  It is related to the
orientation $\varepsilon'$ of the tangent bundle of $S^n$ by
\[
\varepsilon=(-1)^{\frac{n(n-1)}2}\varepsilon'.
\]

\subsubsection*{3.2.8.}
The \underline{vanishing cycle} $\delta$ is well defined up to sign.
For $n\ne0$, it is the element of $H_c^n(V^\circ)$ defined by the
submanifold $S^n\subset V^\circ$.  For $n=0$, the space $V$ consists
of two points, and $\delta$ is their difference.  With the complex
orientation of $V$, 3.2.7 gives
\begin{equation}\tag{3.2.8.1}
(\delta,\delta)
=(-1)^{\frac{n(n-1)}2}\bigl(1+(-1)^n\bigr).
\end{equation}

\subsubsection*{3.2.9.}
Put
\[
\widetilde H_c^i(V^\circ)
=\operatorname{Ker}\!\left(
H_c^i(V^\circ)\xrightarrow{\ \operatorname{Tr}\ }\mathbb Z
\right);
\]
this is $H_c^i(V^\circ)$ for $i\ne2n$.  Then
\[
\begin{aligned}
\widetilde H_c^i(V^\circ)&=0 &&\text{for }i\ne n,\\
\widetilde H_c^n(V^\circ)&=\mathbb Z,
  &&\text{generated by }\delta.
\end{aligned}
\]
Dually,
\[
\begin{aligned}
\widetilde H^i(V)&=0 &&\text{for }i\ne n,\\
\widetilde H^n(V)&\xrightarrow[\sim]{\ (\,\cdot\,,\delta)\ }\mathbb Z.
\end{aligned}
\]
For $n\ne0$, the image of $\delta\in H_c^n(V^\circ)$ in $H^n(V)$
is zero when $n$ is odd, and is divisible by two (and no further)
when $n$ is even.
\end{filecontents*}
% Original input: components/058_expose_XIV_section_3_2_10_through_3_2_11.tex
\begin{filecontents*}[overwrite,noheader]{sga7ii-cumulative-input-058.tex}
\subsubsection*{3.2.10.}
Let us calculate the variation
\[
\operatorname{Var}:\widetilde H^n(V)
\longrightarrow H_c^n(V^\circ).
\]
We first use the orientation $\varepsilon'$ and suppose that $n\ne0$.
Let $\alpha$ be an orientation of the tangent bundle of $S^n$, and
let $\delta\in H_c^n(V^\circ)$ be defined by
$S^n\subset V^\circ$ together with the orientation $\alpha$ of its
normal bundle.  The cocycle $\delta'\in H^n(V)$ such that
$(\delta',\delta)=1$ is represented by a fiber of the fibration of
$V^\circ$ over $S^n$.  The normal bundle of this fiber is the constant
bundle with value the tangent space at one point of $S^n$, and it is
endowed with the orientation $\alpha$.

This assertion is independent of the sign convention chosen to define
the map
\[
(\text{submanifold + orientation of its normal bundle})
\longrightarrow\text{cohomology}.
\]
Indeed, changing the convention changes $\delta$ to $-\delta$ and
$\delta'$ to $-\delta'$.  The assertion is valid when
\[
\operatorname{Tr}:H_c^{2n}(V,\mathbb R)\longrightarrow\mathbb R
\]
is defined as follows:
\begin{description}[leftmargin=2.8em,labelsep=.5em]
\item[\textup{($\alpha$)}] since $\Omega_V^*$ is a soft resolution of $\underline{\mathbb R}$,
  $\Gamma_c(\Omega_V^*)$ represents $R\Gamma_c(V,\mathbb R)$, whence
  an isomorphism
  \[
  H^i\bigl(\Gamma_c(V,\Omega_V^*)\bigr)
  \simeq H_c^i(V,\mathbb R)
  \]
  (the iterated coboundary is not used);
\item[\textup{($\beta$)}] $\operatorname{Tr}$ is integration.
\end{description}

The class $T(\delta')$ coincides with $\delta'$ in a neighborhood of
$\partial V$, so that
\(
\operatorname{Var}(\delta')=T(\delta')-\delta'
\)
has compact support in $V^\circ$.  We have
\[
T(\delta')-\delta'\sim-\delta,
\qquad\text{and hence}
\]
\[
\operatorname{Var}(\delta')
=-(\delta',\delta)\,\delta.
\]

Here is the picture for $n=1$.  We draw the tangent bundle of $S^1$
so that vectors pointing in the chosen counterclockwise orientation of
$S^1$ point outward.

\begin{center}
\begin{tikzpicture}[>=Latex,line width=.45pt,scale=.92]
  \begin{scope}[shift={(-2.8,0)}]
    \draw (0,0) circle (2.15);
    \draw (0,0) circle (1.20);
    \draw (0,0) circle (.43);
    \draw (0,.43)--(0,2.15);
    \draw[->] (-.02,1.88)--(-.48,1.88);
    \draw[->] (-.02,1.13)--(-.48,1.13);
    \node[right] at (.10,1.10) {$\delta'$};
    \draw[->] (-.85,.85)--(-1.16,1.16);
    \draw[->] (-.85,-.85)--(-1.17,-1.17);
    \node[above left] at (-1.02,-.98) {$\delta$};
    \draw[->] (.85,-.85)--(1.17,-1.17);
    \draw[->] (1.43,1.12)--(1.43,1.72)
      node[right] {$2$};
    \draw[->] (1.43,1.12)--(.83,1.12)
      node[below] {$1$};
    \node at (1.03,1.63) {$\varepsilon$};
    \node[below=5pt] at (0,-2.15)
      {$\delta,\ \delta'\ \text{and}\ \varepsilon$};
  \end{scope}
  \begin{scope}[shift={(2.8,0)}]
    \draw (0,0) circle (2.15);
    \draw (0,0) circle (.43);
    \draw (0,.43)--(0,2.15);
    \draw
      (0,2.15)
      .. controls (.02,1.68) and (1.18,1.42) .. (1.28,.30)
      .. controls (1.40,-.72) and (.62,-1.35) .. (-.12,-1.31)
      .. controls (-1.05,-1.27) and (-1.42,-.42) .. (-1.10,.48)
      .. controls (-.83,1.10) and (-.34,1.25) .. (0,1.02);
    \draw[->] (1.08,.78)--(.74,.48);
    \node[below=5pt] at (0,-2.15)
      {$\operatorname{Var}(\delta')$};
  \end{scope}
\end{tikzpicture}
\end{center}

For the orientation $\varepsilon$, we therefore have
\[
\operatorname{Var}(x)
=-\frac{\varepsilon}{\varepsilon'}(x,\delta)\,\delta.
\]

\subsubsection*{Theorem 3.2.11.}
\underline{The group $\widetilde H_c^n(V^\circ)$ is isomorphic to
$\mathbb Z$, generated by the vanishing cycle $\delta$.  The group
$\widetilde H^n(V)$ is also isomorphic to $\mathbb Z$ and is dual to
the preceding group.  All the other groups
$\widetilde H_c^i(V^\circ)$ and $\widetilde H^i(V)$ vanish.}

\medskip
\noindent\underline{The variation is given by}
\[
\operatorname{Var}(x)
=-(-1)^{\frac{n(n-1)}2}(x,\delta)\,\delta,
\qquad\underline{\text{and}}
\]
\[
(\delta,\delta)
=(-1)^{\frac{n(n-1)}2}\bigl(1+(-1)^n\bigr).
\]
According to the congruence class of $n$ modulo $4$, we have the
following table:
\[
\begin{array}{lcccc}
n\bmod4&0&1&2&3\\[2pt]
(\delta,\delta)&2&0&-2&0\\[2pt]
\operatorname{Var}(x)=\pm(x,\delta)\delta&-&-&+&+\\[2pt]
T(\delta)=(-1)^{n+1}\delta&-\delta&\delta&-\delta&\delta
\end{array}
\]
As a check, we indeed have
\[
T(\delta)=\delta+\operatorname{Var}(\delta).
\]
\end{filecontents*}
% Original input: components/059_expose_XIV_section_4_1_through_4_4.tex
\begin{filecontents*}[overwrite,noheader]{sga7ii-cumulative-input-059.tex}
\section*{4. De Rham cohomology}

\subsubsection*{4.1.}
Let $D$ be a complex analytic disc, let $0\in D$, put
$D^*=D-\{0\}$, and let $\widetilde D^*$ be a universal covering of
$D^*$.  Let $X$ be an analytic space, let
$f:X\longrightarrow D$ be a morphism, and consider the diagram of
(1.2.7.1):
\begin{equation}\tag{4.1.1}
\begin{tikzcd}[column sep=large,row sep=large]
X_0 \arrow[r,hook,"\bar i"] \arrow[d,equal]
  & \bar X \arrow[d,"p"']
  & \bar X^* \arrow[l,hook',"\bar j"'] \arrow[d,"p'"'] \\
X_0 \arrow[r,hook,"i"]
  & X
  & X^* \arrow[l,hook',"j"']
\end{tikzcd}
\end{equation}
Suppose that $X^*=X-X_0$ is nonsingular (that is, smooth).  Its
covering $\bar X^*$ is therefore smooth.  Denote by
$\Omega_{\bar X^*}^*$ the De Rham complex of sheaves of holomorphic
differential forms on $\bar X^*$.  It is a resolution of the constant
sheaf $\mathbb C$ (the holomorphic Poincare lemma), and hence gives a
morphism
\begin{equation}\tag{4.1.2}
\bar j_*\Omega_{\bar X^*}^*
\longrightarrow R\bar j_*\mathbb C
\end{equation}
in $D^+(\bar X,\mathbb C)$.

\subsubsection*{Proposition 4.2.}
\underline{The morphism (4.1.2) is an isomorphism.}

It is clear that (4.1.2) is a quasi-isomorphism at every point of
$\bar X^*\subset\bar X$.\footnote{The printed proof says (4.1.1),
which is the preceding diagram rather than the morphism under
discussion; the English restores the evident reference to (4.1.2).}
It is therefore enough to show that every point $x\in X_0$ has a
fundamental system of neighborhoods $U$ in $\bar X$ such that
\[
H^q\bigl(\bar j^{-1}(U),\Omega^p\bigr)=0
\qquad\text{for }q>0.
\]
As $V$ runs through a fundamental system of neighborhoods of $x$ in
$X$, the sets $p^{-1}(V)$ form a fundamental system of neighborhoods
of $x$ in $\bar X$.  In particular, one may take the $p^{-1}(V)$ for
$V$ a Stein neighborhood of $x$ in $X$.  If $V$ is Stein, then
$V-X_0$ is still Stein (since $X_0\cap V$ is defined by one equation
in $V$), as is its covering
\[
p'^{-1}(V-X_0)=\bar j^{-1}p^{-1}(V).
\]
Consequently, for every coherent sheaf $\mathcal F$ on $\bar X^*$,
\[
H^q\bigl(\bar j^{-1}p^{-1}(V),\mathcal F\bigr)=0
\qquad\text{for }q>0,
\]
which proves the proposition.

\subsubsection*{Corollary 4.3.}
\underline{The isomorphism (4.1.2) induces an isomorphism}
\begin{equation}\tag{4.3.1}
\bar i^*\bar j_*\Omega_{\bar X^*}^*
\xrightarrow{\ \sim\ }R\Psi_\eta(\mathbb C).
\end{equation}

\subsubsection*{Variant 4.4.}
\underline{Let $V$ be a locally constant sheaf of finite-dimensional
complex vector spaces on $X^*$ (a complex local system).  The
quasi-isomorphism}
\[
V\longrightarrow\Omega_{X^*}^*(V)
\]
\underline{of the Poincare lemma induces a quasi-isomorphism}
\begin{equation}\tag{4.4.1}
\bar i^*\bar j_*\Omega_{\bar X^*}^*(p'^*V)
\xrightarrow{\ \sim\ }R\Psi_\eta(V).
\end{equation}

We shall show that (4.3.1) remains true when
$\bar i^*\bar j_*\Omega_{\bar X^*}^*$ is replaced by a subcomplex
defined by growth and monodromy conditions.
\end{filecontents*}
% Original input: components/060_expose_XIV_section_4_5_through_4_8.tex
\begin{filecontents*}[overwrite,noheader]{sga7ii-cumulative-input-060.tex}
\subsubsection*{4.5.}
Let $f:X^*\longrightarrow D^*$ be a topological space over $D^*$,
let $F$ be a sheaf of $f^\circ\mathcal O_{D^*}$-modules on $X^*$,
and let $\bar F$ be its inverse image on
\[
\bar X^*=X^*\times_{D^*}\widetilde D^*.
\]

Let $z:D\hookrightarrow\mathbb C$ be a holomorphic embedding with
$z(0)=0$.  We shall always be free to shrink $D$ below; for a
sufficiently small disc the existence of $z$ is trivial.  Choose on
$\widetilde D^*$ a function $\log z$ such that
$\exp(\log z)=z$, and put
\[
z^\alpha=\exp(\alpha\log z)
\qquad(\alpha\in\mathbb C).
\]

The group $\pi_1(D^*)$ acts, by transport of structure, on
$H^0(\bar X^*,\bar F)$.  An element
$f\in H^0(\bar X^*,\bar F)$ is of \underline{finite determination}
if its transforms under $\pi_1(D^*)$ span a finite-dimensional
complex vector space $V(f)$.  Let $T$ be the positive generator of
$\pi_1(D^*)$.  The section $f$ is of finite determination if and only
if there is a nonzero polynomial $P\in\mathbb C[X]$ such that
$P(T)(f)=0$.  The polynomial of least degree with this property is
then both the minimal polynomial and the characteristic polynomial of
$T$ acting on $V(f)$.

We say that $f$ is of \underline{unipotent finite determination} if
there is an $A$ such that
\[
(T-1)^A f=0.
\]
On the space of these sections define
\[
N=\frac{\log T}{2\pi i}
\]
by the finite sum
\[
N(f)=-\frac1{2\pi i}
\sum_{n>0}\frac{(1-T)^n}{n}\,f.
\]
If $f$ (respectively, $g$) is a section of unipotent finite
determination in $H^0(D^*,\mathcal O)$ (respectively, in
$H^0(\bar X^*,\bar F)$), then
\[
N(fg)=N(f)g+fN(g),
\]
since the logarithm of an automorphism is a derivation.

We say that $f$ is of
\underline{quasi-unipotent finite determination} if there are
$A$ and $B>0$ such that
\[
(T^B-1)^A f=0.
\]

\subsubsection*{Examples 4.6.}
For $X=D^*$ and $F=\mathcal O$, one has
\[
\bigl(T-\exp(2\pi i\alpha)\bigr)^{k+1}
  \bigl(z^\alpha(\log z)^k\bigr)=0,
\]
\[
N(\log z)=1,
\]
and
\[
N\left(\frac{(\log z)^k}{k!}\right)
=\frac{(\log z)^{k-1}}{(k-1)!},
\qquad
(0\text{ for }k=0).
\]

\subsubsection*{Proposition 4.7.}
\underline{A section $f\in H^0(\bar X^*,\bar F)$ is of finite
determination if and only if it can be written}
\begin{equation}\tag{4.7.1}
f=\sum_{\alpha,k}f_{\alpha,k}\,
z^\alpha\frac{(\log z)^k}{k!},
\end{equation}
\underline{where $0\leq\operatorname{Re}(\alpha)<1$,
$k\in\mathbb N$, and the
$f_{\alpha,k}\in H^0(X^*,F)$ are almost all zero.  For $f$ of finite
determination, the decomposition (4.7.1) is unique.}

By 4.6, every section of the form (4.7.1) is of finite determination.
Conversely, let $f$ be of finite determination, with $P(T)(f)=0$.
Write
\[
P(X)=\prod_{\beta\in\mathbb C}(X-\beta)^{n'(\beta)}
=\prod_{0\leq\operatorname{Re}(\alpha)<1}
 \bigl(X-\exp(2\pi i\alpha)\bigr)^{n(\alpha)},
\]
where
$n(\alpha)=n'(\exp(2\pi i\alpha))$ is almost always zero.
Decomposing the $\mathbb C[X]$-module $V(f)$ into its isotypic
components, one finds a unique decomposition
\[
f=\sum_\alpha f'_\alpha
\]
such that
\[
\bigl(T-\exp(2\pi i\alpha)\bigr)^{n(\alpha)}f'_\alpha=0.
\]

Choose a polynomial $Q_\alpha\in\mathbb C[X]$ congruent to $1$
modulo
$\bigl(X-\exp(2\pi i\alpha)\bigr)^{n(\alpha)}$ and to $0$ modulo
$\bigl(X-\exp(2\pi i\alpha')\bigr)^{n(\alpha')}$ for every
$\alpha'\ne\alpha$.  Then
$f'_\alpha=Q_\alpha(T)(f)$.  Put
\begin{equation}\tag{4.7.2}
f_\alpha=z^{-\alpha}f'_\alpha
=z^{-\alpha}Q_\alpha(T)(f).
\end{equation}
Since
\[
T(z^{-\alpha}f'_\alpha)
=\exp(-2\pi i\alpha)z^{-\alpha}T(f'_\alpha),
\]
we have
\[
(T-1)^{n(\alpha)}f_\alpha=0.
\]

If
\[
f_\alpha=\sum_k f_{\alpha,k}\frac{(\log z)^k}{k!},
\qquad f_{\alpha,k}\in H^0(X^*,F),
\]
then 4.6 and the identity
$\sum_j(-1)^j\binom{k}{j}=\delta_{0,k}$ give
\begin{equation}\tag{4.7.3}
f_{\alpha,k}
=\sum_{j\geq0}(-1)^j\frac{(\log z)^j}{j!}
  N^{j+k}f_\alpha,
\end{equation}
which proves uniqueness.

Conversely, if $f_\alpha$ is a section of unipotent finite
determination, the $f_{\alpha,k}$ defined by (4.7.3) vanish for large
$k$ (because $N^k f_\alpha=0$ for large $k$) and are invariant under
$T$; hence they identify with sections of $F$ on $X^*$.  Equivalently,
we must show that $Nf_{\alpha,k}=0$.  Apply to $N^k f_\alpha$ the
identity
\[
\begin{aligned}
N\sum_{j\geq0}(-1)^j\frac{(\log z)^j}{j!}N^j
&=\sum_{j>0}(-1)^j\frac{(\log z)^{j-1}}{(j-1)!}N^j\\
&\quad+\sum_{j\geq0}(-1)^j\frac{(\log z)^j}{j!}N^{j+1}=0.
\end{aligned}
\]
Finally,
\[
\sum_{k\geq0}f_{\alpha,k}\frac{(\log z)^k}{k!}
=\sum_{j,k\geq0}(-1)^j
 \frac{(\log z)^{j+k}}{j!k!}N^{j+k}f_\alpha
\]
\[
=\sum_{n\geq0}
 \left(\sum_j(-1)^j\binom nj\right)
 \frac{(\log z)^n}{n!}N^n f_\alpha
=f_\alpha,
\]
which proves 4.7.

\subsubsection*{Remark 4.8.1.}
The condition $\operatorname{Re}(\alpha)\in[0,1)$ in 4.7 may be
replaced by $\alpha\in C$, where $C$ is any fundamental domain for
$\mathbb C\longrightarrow\mathbb C/\mathbb Z$.

\subsubsection*{Remark 4.8.2.}
Let $f$ be of finite determination, with decomposition (4.7.1).  For
$0\leq\operatorname{Re}(\alpha)<1$ and for $\alpha$ such that some
$f_{\alpha,k}\ne0$, put
\[
m(\alpha)=1+\sup\{k\mid f_{\alpha,k}\ne0\}.
\]
If all the $f_{\alpha,k}$ vanish, put $m(\alpha)=0$.  The
characteristic polynomial of $T$ acting on $V(f)$ is
\[
\operatorname{Pc}(T,V(f))
=\prod_\alpha\bigl(X-\exp(2\pi i\alpha)\bigr)^{m(\alpha)}.
\]
In particular, $f$ is of quasi-unipotent finite determination if and
only if the $\alpha$ in (4.7.1) may all be taken in $\mathbb Q$.
\end{filecontents*}
% Original input: components/061_expose_XIV_section_4_9_through_4_17.tex
\begin{filecontents*}[overwrite,noheader]{sga7ii-cumulative-input-061.tex}
\subsubsection*{4.9.}
Return to the more restrictive hypotheses of 4.1, and let $F$ be a
coherent sheaf on $X$ whose restriction to $X^*$ is locally free.

\subsubsection*{Proposition 4.10.}
\underline{Let $f$ be a section of finite determination of the inverse
image $\bar F$ of $F$ on $\bar X^*$.  Then $f$ has moderate growth
along $X_0$ if and only if the $f_{\alpha,k}$ in the decomposition
(4.7.1) are meromorphic along $X_0$.}

For the definition of ``moderate growth,'' we refer to [1], 2.10.  It
is clear from (4.7.1), (4.7.2), and (4.7.3) that $f$ has moderate
growth if and only if the $f_{\alpha,k}$ do.  For sections of $F$ on
$X^*$, moderate growth along $X_0$ is known to be equivalent to being
meromorphic along $X_0$.

For $X=D$ and $F=\mathcal O$, 4.10 has the following elementary
corollary.

\subsubsection*{Corollary 4.11.}
\underline{Let $K$ be the field of germs at $0$ of meromorphic
functions on $D$, and let $K^m$ (respectively, $K^{mu}$ and $K^{mqu}$)
be the $K$-algebra of germs at $0$ of holomorphic functions on
$\widetilde D^*$ which have moderate growth at $0$ and finite
determination (respectively, unipotent finite determination and
quasi-unipotent finite determination).  Then the following are bases
over $K$:}
\[
\begin{array}{lll}
K^m:&z^\alpha(\log z)^k,
  &0\leq\operatorname{Re}(\alpha)<1,\ k\in\mathbb N,\\[2pt]
K^{mu}:&(\log z)^k,&k\in\mathbb N,\\[2pt]
K^{mqu}:&z^\alpha(\log z)^k,
  &\alpha\in\mathbb Q,\ 0\leq\alpha<1,\ k\in\mathbb N.
\end{array}
\]

Under the hypotheses of 4.9, denote by $\Psi_\eta^m(F)$
(respectively, $\Psi_\eta^{mu}(F)$ and $\Psi_\eta^{mqu}(F)$) the
subsheaf of $\bar i^*\bar j_*\bar F$ consisting of sections of
$\bar F$, near points of $X_0$, which have moderate growth along
$X_0$ and finite determination (respectively, unipotent finite
determination and quasi-unipotent finite determination).  Let
\[
F^m\subset j_*j^*F
\]
be the sheaf on $X$ of sections of $j_*j^*F$ which are meromorphic
along $X_0$.  Proposition 4.10 has the following corollary, independent
of the arbitrary choices of $z$ and $\log z$.

\subsubsection*{Corollary 4.12.}
\underline{The natural morphisms}
\[
\textup{(i)}\qquad
K^m\otimes_K i^*F^m\longrightarrow\Psi_\eta^m(F),
\]
\[
\textup{(ii)}\qquad
K^{mu}\otimes_K i^*F^m\longrightarrow\Psi_\eta^{mu}(F),
\]
\[
\textup{(iii)}\qquad
K^{mqu}\otimes_K i^*F^m\longrightarrow\Psi_\eta^{mqu}(F)
\]
\underline{are isomorphisms.}

The main result of this section is the following.

\subsubsection*{Theorem 4.13.}
\underline{The natural morphism induced by (4.3.1),}
\begin{equation}\tag{4.13.1}
\Psi_\eta^{mqu}(\Omega_X^*)
\longrightarrow R\Psi_\eta(\mathbb C),
\end{equation}
\underline{is an isomorphism in $D^+(X_0,\mathbb C)$.}

We shall in fact prove the more general results 4.15 and 4.17.  Their
statements require a few preliminaries which a reader interested only
in 4.13 may omit.

\subsubsection*{4.14.}
Let $V$ be a locally constant sheaf of finite-dimensional complex
vector spaces on $X^*$ (a local system).  Let
\[
\mathcal V=V\otimes_{\mathbb C}\mathcal O
\]
be the corresponding locally free $\mathcal O$-module, and let
$\nabla$ be the integrable connection on $\mathcal V$ for which the
local sections of $V$ are horizontal.  It is known that $\mathcal V$
admits an extension $\widetilde{\mathcal V}$ to a coherent sheaf on
$X$ such that $\nabla$ is \underline{regular} along $X_0$ relative to
this extension.  It is also known that any two such extensions are
meromorphically equivalent ([1], II, 5.7).  Thus the sheaf
$\Psi_\eta^m(\widetilde{\mathcal V})$, and its variants
$\Psi^{mu}$ and $\Psi^{mqu}$, does not depend on the chosen extension;
we denote it by $\Psi_\eta^m(\mathcal V)$.  Likewise define
\[
\Psi_\eta^m\bigl(\Omega_{X^*}^p(V)\bigr)
=\Psi_\eta^m(\Omega_X^p\otimes\mathcal V),
\]
and its variants $\Psi^{mu}$ and $\Psi^{mqu}$.

\subsubsection*{Proposition 4.15.}
\underline{The natural morphism induced by (4.4.1),}
\begin{equation}\tag{4.15.1}
\Psi_\eta^m\bigl(\Omega_{X^*}^*(V)\bigr)
\longrightarrow R\Psi_\eta(V),
\end{equation}
\underline{is an isomorphism in $D^+(X_0,\mathbb C)$.}

\subsubsection*{4.16.}
A local system $V$ as in 4.14 is said to be \underline{unipotent}
(respectively, \underline{quasi-unipotent}) along $X_0$ if, for every
morphism $u$ from the standard disc $D_0$ to $X$ such that
$u^{-1}(X_0)=\{0\}$, the monodromy transformation of the local system
$u^*V$ on $D_0^*$ is unipotent (respectively, quasi-unipotent).

\subsubsection*{Proposition 4.17.}
\underline{If $V$ is quasi-unipotent, the natural morphism induced by
(4.4.1),}
\begin{equation}\tag{4.17.1}
\Psi_\eta^{mqu}\bigl(\Omega_{X^*}^*(V)\bigr)
\longrightarrow R\Psi_\eta(V),
\end{equation}
\underline{is an isomorphism in $D^+(X_0,\mathbb C)$.}

Theorem 4.13 is the special case of 4.17 in which
$V=\underline{\mathbb C}$.
\end{filecontents*}
% Original input: components/062_expose_XIV_section_4_18.tex
\begin{filecontents*}[overwrite,noheader]{sga7ii-cumulative-input-062.tex}
\subsubsection*{4.18.}
We shall first treat the case in which $X$ is smooth and
$(X_0)_{\mathrm{red}}$ is a normal-crossings divisor in $X$.  Since
the theorems are local, we may suppose that $f$ is the map
\[
f(\underline z)=\prod_{i=1}^d z_i^{n_i},
\qquad 0<d\leq e,\quad n_i>0,
\]
from $X=D^e$ to $D$, where this time $D$ denotes the standard disc
$D=\{z\in\mathbb C\mid |z|<1\}$.  It is enough to prove that the
fiber at $0\in D^e$ of (4.15.1) (respectively, (4.17.1)) is a
quasi-isomorphism.

Let $U$ be a neighborhood of $0$ in $D^e$ of the form
$|z_i|<a$, where $0<a\leq1$.  Define $U_0$, $\bar U$, $U^*$, and
$\bar U^*$ by taking the inverse image of $U$ in (4.1.1).  The
cohomology of $U^*$ or $\bar U^*$ is independent of $a$, and hence
\[
H^*(\bar U^*,V)\xrightarrow{\ \sim\ }
  \bigl(R^*\Psi_\eta(V)\bigr)_0.
\]
The fundamental group $\mathbb Z^d$ of
$X^*=(D^*)^d\times D^{e-d}$ (which is also that of $U^*$) is
abelian.  Therefore $V$ has a finite filtration whose successive
quotients are rank-one local systems.  By devissage, we may thus
suppose that $V$ has rank one.

For $\widetilde{\mathcal V}$ as in 4.14, write
$j_*^m\Omega_{X^*}^*(V)$ for the subcomplex of
$j_*\Omega_{X^*}^*(V)$ whose degree-$p$ term is
$(\Omega_X^p\otimes\widetilde{\mathcal V})^m$.  By [1], II 6.4.1
and 6.8, the natural map induced by the quasi-isomorphism
$V\xrightarrow{\ \approx\ }\Omega_{X^*}^*(V)$ (the holomorphic
Poincare lemma),
\begin{equation}\tag{4.18.1}
j_*^m\Omega_{X^*}^*(V)\longrightarrow Rj_*V,
\end{equation}
is a quasi-isomorphism.  Thus
\begin{equation}\tag{4.18.2}
H^*\bigl((j_*^m\Omega_{X^*}^*(V))_0\bigr)
 \xrightarrow{\ \sim\ }(R^*j_*V)_0
 \xleftarrow{\ \sim\ }H^*(U^*,V).
\end{equation}
\footnote{The printed right-hand term in (4.18.2) is $H^*(U,V)$,
although $V$ is defined on $U^*$; the English restores the evident
$U^*$.}
This will be the only analytic ingredient in this part 4.18 of the
proof.

\subsubsection*{Lemma 4.18.3.}
\underline{For a rank-one $V$, one has}

\noindent\textup{(i)}\qquad
\uline{if $V$ is nontrivial:}\quad $H^*(U^*,V)=0$,

\medskip
\noindent\textup{(ii)}\qquad
\uline{if $V=\underline{\mathbb C}$:}\quad
$H^*(U^*,V)=\bigwedge\nolimits^*\mathbb C^d$.

Indeed, $U^*$ has the homotopy type of a torus.

Let $\eta$ be the generator of $H^1(D^*,\mathbb Z)$, and also its
inverse image in $H^1(U^*,\mathbb Z)$.

\subsubsection*{Lemma 4.18.4.}
\underline{For a rank-one $V$, the natural map}
\[
H^*(U^*,V)
 \xrightarrow[(4.18.1)]{\ \sim\ }
H^*\bigl((j_*^m\Omega_{X^*}^*(V))_0\bigr)
 \longrightarrow
H^*\bigl(\Psi_\eta^{mu}(\Omega_{X^*}^*(V))_0\bigr)
\]
\underline{identifies
$H^*(\Psi_\eta^{mu}(\Omega_{X^*}^*(V))_0)$ with
$H^*(U^*,V)/\eta H^*(U^*,V)$.}
\footnote{The printed left-hand term is $H^*(U^*,\mathbb C)$,
while both the hypothesis and the conclusion concern the general
rank-one system $V$; the English restores $V$.}

Use 4.12(ii) and the identity
$d\log f=\sum_i n_i\,dz_i/z_i$.  The complex
$\Psi_\eta^{mu}(\Omega_{X^*}^*(V))_0$ identifies with the simple
complex $\mathrm{s}K$ associated (using direct sums) with the double
complex $K$ whose terms are
\[
K^{-p,n+p}=\frac{(\log f)^p}{p!}\cdot
  (j_*^m\Omega_{X^*}^n(V))_0.
\]
This complex lies in a bad quadrant
$(p\leq0,q\geq0,p+q\geq0)$.  Nevertheless, $\mathrm{s}K$ is the
inductive limit of the subcomplexes
$\mathrm{s}\,\sigma_{p\geq p_0}(K)$ whose terms are
\[
(\mathrm{s}\,\sigma_{p\geq p_0}K)^n
=\bigoplus_{\substack{p+q=n\\p\geq p_0}}K^{pq},
\]
and this suffices to justify the argument below.

The spectral sequence for the filtration by the first degree is
\[
E_1^{-p,n+p}=H^n(U^*,V)\qquad(p\geq0),
\]
with
\[
d_1=2\pi i\,\eta\wedge:
H^n(U^*,V)\longrightarrow H^{n+1}(U^*,V).
\]
If $V$ is nontrivial, then $E_1=0$ by 4.18.3, and 4.18.4 reduces
to $0=0$.  Suppose that $V=\underline{\mathbb C}$.  The complexes
$(E_1,d_1)$ are then Koszul complexes, possibly truncated, so that,
by acyclicity of the Koszul complex,
\[
E_2^{pq}=0\qquad\text{for }p\ne0,
\]
and
\[
E_2^{0n}=H^n(U^*,V)/\eta H^{n-1}(U^*,V),
\]
which proves 4.18.4.

For $\beta\in\mathbb C^*$, let $V_\beta$ denote the rank-one local
system on $D^*$ whose inverse image on $\widetilde D^*$ is
trivialized and whose monodromy is $\beta$.  We use the same notation
for its inverse image on $U^*$.  On $\bar U^*$, the inverse image of
$V\otimes V_\beta$ is canonically isomorphic to that of $V$.

\subsubsection*{Lemma 4.18.5.}
\underline{For a rank-one $V$, the map}
\[
\bigoplus_{\beta\in\mathbb C^*}
H^*(U^*,V\otimes V_\beta)/
 \eta H^*(U^*,V\otimes V_\beta)
\longrightarrow H^*(\bar U^*,V)
\]
\underline{is bijective.}

Let $k$ be the greatest common divisor of the $n_i$.  The space
$\bar U^*$ has $k$ connected components, indexed by the $k$th roots
of unity.  On the component $\bar U_\xi^*$ indexed by $\xi$, the
inverse image of the function $z^{1/k}$ on $\widetilde D^*$
coincides with the inverse image of the function
\[
\xi\prod_{i=1}^d z_i^{n_i/k}
\]
on $U^*$.  The component $\bar U_\xi^*$ has the homotopy type of a
torus, and the sequence
\[
0\longrightarrow\pi_1(\bar U_\xi^*)
 \longrightarrow\pi_1(U^*)\longrightarrow\pi_1(D^*)
\]
is exact.  Write $\varphi:\mathbb Z^d\to\mathbb Z$ for the map
$\pi_1(U^*)\to\pi_1(D^*)$.  If $\ker(\varphi)$ acts nontrivially on
$V$, both sides of 4.18.5 vanish.\footnote{The print refers here to
4.18.4; the two sides under discussion are those of 4.18.5.}
Otherwise, after replacing $V$ by a suitable $V\otimes V_\beta$, we
may suppose that $V$ is trivial, that is, $V=\mathbb C$.  In complex
cohomology we then have
\[
H^0(\bar U^*)\otimes
 H^*(U^*)/\eta H^*(U^*)
 \xrightarrow{\ \sim\ }H^*(\bar U^*).
\]
The result follows on observing that $V_\beta$ is trivial on $U^*$
when $\beta^k=1$, that
\[
\bigoplus_{\beta^k=1}H^0(U^*,V_\beta)
=\bigoplus_{\beta^k=1}\mathbb C
 \xrightarrow{\ \sim\ }H^0(\bar U^*,\mathbb C),
\]
and that $H^*(U^*,V_\beta)=0$ when $\beta^k\ne1$.

\subsubsection*{Remark 4.18.6.}
If $V$ is quasi-unipotent, the sum over $\mathbb C^*$ in 4.18.5 may
be replaced by a sum over the roots of unity; the proof is the same.

Applying 4.12 gives a commutative diagram
\begin{center}
\begin{tikzcd}[
  column sep=2.2em,
  row sep=large,
  cells={nodes={font=\scriptsize}}
]
H^*\Psi_\eta^m(\Omega_{X^*}^*(V))_0
  \arrow[d]
& \displaystyle\bigoplus_\beta
  H^*\Psi_\eta^{mu}(\Omega_{X^*}^*(V\otimes V_\beta))
  \arrow[l,"\sim"',"4.12"]
& \displaystyle\bigoplus_\beta
  H^*(U^*,V\otimes V_\beta)/
    \eta H^*(U^*,V\otimes V_\beta)
  \arrow[l,"\sim"',"4.18.4"]
  \arrow[d,"\wr"',"4.18.5"]
\\
H^*R\Psi(V)_0
& {}
& H^*(\bar U^*,V)\arrow[ll,"\sim"']
\end{tikzcd}
\end{center}
from which 4.15 follows.  Proposition 4.17 is obtained by replacing
$\Psi^m$ with $\Psi^{mqu}$, the sum over $\beta$ with a sum over the
roots of unity, and 4.18.5 with 4.18.6.
\end{filecontents*}
% Original input: components/063_expose_XIV_section_4_19_through_bibliography.tex
\begin{filecontents*}[overwrite,noheader]{sga7ii-cumulative-input-063.tex}
\subsubsection*{4.19.}
Let us deduce 4.15 and 4.17 from the special case 4.18.  Let
$u:Y\longrightarrow X$ be a resolution of $(X,X_0)$, that is, a
proper morphism inducing an isomorphism
$Y^*\xrightarrow{\ \sim\ }X^*$, such that $Y$ is smooth and
$(Y_0)_{\mathrm{red}}$ is a normal-crossings divisor in $Y$.  If $V$
is quasi-unipotent along $X_0$, its inverse image $u^*V$ is
quasi-unipotent along $Y_0$.

For every coherent sheaf $F$ on $X$, denote by $F^m$, or by
$j_*^m j^*F$, the subsheaf of $j_*j^*F$ formed by the meromorphic
sections.  If $f$ is an equation for $X_0$, then $j_*^m j^*F$ is the
inductive limit of the system
\begin{equation}\tag{4.19.1}
F\xrightarrow{f}F\xrightarrow{f}F\longrightarrow\cdots.
\end{equation}
The following lemma is proved in [2].

\subsubsection*{Lemma 4.19.2.}
\underline{For every coherent sheaf $F$ on $X$, one has}
\[
F^m\xrightarrow{\ \sim\ }u_*(u^*F)^m
\]
\underline{and $R^iu_*(u^*F)^m=0$ for $i>0$.}

By Grauert's theorem, the $R^iu_*u^*F$ are coherent sheaves.  The
kernels and cokernels of $F\longrightarrow u_*u^*F$, as well as the
$R^iu_*u^*F$ for $i>0$, vanish outside $X_0$ and are therefore
locally annihilated by a power of $f$.  Since $u$ is proper,
$R^iu_*$ commutes with filtered inductive limits, in particular with
the inductive limit (4.19.1), and (4.19.2) follows.

From 4.19.2 and 4.12 one deduces that, under the hypotheses just
stated and still denoting by $u$ the map $Y_0\to X_0$ induced by
$u$, one has the following.

\subsubsection*{Lemma 4.19.3.}
\underline{The morphisms}
\[
\Psi_\eta^m(F)\longrightarrow
u_*\Psi_\eta^m(u^*F)\longrightarrow
Ru_*\Psi_\eta^m(u^*F)
\]
\underline{are isomorphisms.}

That the second arrow is an isomorphism means that
\[
R^iu_*\Psi_\eta^m(u^*F)=0\qquad(i>0).
\]
This follows from 4.12 and from
$R^iu_*(i^*u^*F^m)=0$ for $i>0$, which in turn follows from
$R^iu_*(u^*F^m)=0$ for $i>0$ and the proper base-change theorem for
$u$.

To prove 4.15, it remains to consider the commutative diagram
\begin{center}
\begin{tikzcd}[
  column sep=large,
  row sep=large,
  cells={nodes={font=\small}}
]
\Psi^m(\Omega_{X^*}^*(V))
  \arrow[d]
  \arrow[r,equal,"(4.19.3)"']
& u_*\Psi^m(\Omega_{Y^*}^*(u^*V))
  \arrow[r,"\sim","(4.19.3)"']
& Ru_*\Psi^m(\Omega_{Y^*}^*(u^*V))
  \arrow[d,"\wr"',"(4.18)"]
\\
R\Psi(V)
  \arrow[rr,"(1.3.1;\ \mathrm{XIII}\ 2.1.7.1)"]
& {}
& Ru_*R\Psi(u^*V)
\end{tikzcd}
\end{center}
The proof of 4.17 is the same, with $\Psi^m$ replaced everywhere by
$\Psi^{mqu}$.

\subsubsection*{Remark 4.20.}
Theorem 4.13 gives an incarnation of the monodromy theorem at the
level of complexes.

\subsection*{Bibliography}

\begin{enumerate}[label={[\arabic*]},leftmargin=3em]
\item P. Deligne, \emph{Equations diff\'erentielles \`a points
  singuliers r\'eguliers}, Lecture Notes in Mathematics 163,
  Springer-Verlag.
\item A. Grothendieck, \emph{On the De Rham cohomology of algebraic
  varieties}, Publ. Math. IHES \underline{29}, 95--103 (1966).
\item K. Lamotke, \emph{Isolated critical points and monodromy}.
\end{enumerate}
\end{filecontents*}
% Original input: components/064_expose_XV_title_through_1_1_4.tex
\begin{filecontents*}[overwrite,noheader]{sga7ii-cumulative-input-064.tex}
\clearpage
\begin{center}
{\large SGA 7 --- EXPOS\'E XV}\par
\vspace{2em}
{\Large\underline{The Picard--Lefschetz Formula}}\par
\vspace{1.5em}
{\large by P. Deligne}
\end{center}

\medskip
\noindent\textbf{Source contents}
\begin{center}
\begin{tabular}{@{}llr@{}}
\S1. & Ordinary quadratic singularities (normal forms) & 2\\
1.1 & Reminders & 2\\
1.2 & Ordinary quadratic singularities & 4\\
1.3 & Moduli of quadratic singularities & 10\\
\S2. & Computation of vanishing cycles in a standard ordinary
  quadratic case & 13\\
2.1 & Preliminary: the cohomology of a cone & 13\\
2.2 & Vanishing cycles & 17\\
\S3. & The Picard--Lefschetz formula & 23\\
3.1 & Preliminary results & 23\\
3.2 & Even relative dimension & 24\\
3.3 & Odd relative dimension & 26\\
3.4 & Summary & 31
\end{tabular}
\end{center}

\section*{1. Ordinary quadratic singularities (normal forms)}

In this section we take up the results of VI 6 again, adopting a
``henselian'' rather than a ``formal'' point of view.  We give normal
forms for ordinary quadratic singularities and their deformations,
into which they may be put locally for the etale topology.

These results follow easily from VI 6, the implicit-function theorem
(in Tougeron's form), and R. Elkik's theorem on versal henselian
moduli of isolated singularities.  At the time of the oral seminar,
the latter was not yet available; a weaker theorem, covering only the
case of complete intersections (which was enough for us), had been
proved and will not be published.

\subsection*{1.1. Reminders}

\subsubsection*{1.1.1.}
Let $f:X\longrightarrow S$ be a morphism, with $X$ locally
describable by $p$ equations in relative affine space
$\mathbb A_S^{n+p}$ of dimension $n+p$ over $S$.  The Jacobian ideal
$J^n(X/S)$ was defined in VI 5.  If $X$ is the closed subscheme of
$\mathbb A_S^{n+p}$ defined by the equations
$f_i=0$ $(1\leq i\leq p)$, it is the ideal sheaf generated by the
$p\times p$ minors of the Jacobian matrix
$(\partial f_i/\partial x_j)$.  The following theorem is well known.
A special case of its formal analogue is proved in Bourbaki,
\emph{Alg\`ebre commutative}, III, \S4, no.~6, Corollary 3 to
Theorem 2; the $C^\infty$ analogue is proved in Tougeron,
\emph{Id\'eaux de fonctions diff\'erentiables} (thesis, Rennes,
May 1967); the case considered here is treated in Artin,
\emph{Algebraic approximation of structures over complete local
rings}, Publ. Math. IHES \underline{36} (Zariski's volume),
Lemma 5.10.

\subsubsection*{Theorem 1.1.2 (Tougeron--Artin).}
\underline{Let $S$ be the spectrum of a henselian local ring $A$,
let $f:X\to S$ be a morphism as in 1.1.1, let $a$ and $I$ be two
proper ideals of $A$, and let $\delta'$ be an ideal sheaf of $X$
contained in $J^n(X/S)$.  Suppose that $I$ is of finite type.  Let
$s$ be a section of $f$ over $S/I$, that is, $s\in X(A/I)$.  Denote
by $s^*\delta'^2$ the ideal of $A$ inverse to the ideal
$s^*\delta'^2$ of $A/I$.}

\underline{If $(s^*\delta'^2)\,a\supset I$, then there is a section
$s_1\in X(A)$ of $f$ congruent to $s$ modulo
$(s^*\delta')\,a$.}

\subsubsection*{1.1.3.}
Let $S$ be the spectrum of a henselian local ring, with closed point
$s$.  Let $X'_0$ be a scheme of finite type over $k(s)$, smooth away
from a closed point $x_0$, and let $X_0$ be the henselization of
$X'_0$ at $x_0$.  We are interested in systems
\begin{center}
\begin{tikzcd}[column sep=large,row sep=large]
X_0 \arrow[r,hook] \arrow[d]
  & X \arrow[d]
  & {} \\
s \arrow[r,hook]
  & T \arrow[r]
  & S
\end{tikzcd}
\end{center}
of the following kind:
\begin{enumerate}[label=(\alph*),leftmargin=2.6em]
\item $T$ is henselian local and essentially of finite presentation
over $S$ (I 0.0.5); the morphism $T\to S$ is local and induces an
isomorphism of residue fields;

\item $X$ is henselian local, flat, and essentially of finite
presentation over $T$;

\item the special fiber of $X/T$ is equipped with a
$k(s)$-isomorphism to $X_0$.
\end{enumerate}
These are the \underline{henselian $S$-deformations} of $X_0$.
A \underline{morphism} from $(X,T)$ to $(X',T')$ is a commutative
diagram of $S$-schemes
\begin{center}
\begin{tikzcd}[column sep=large,row sep=large]
X \arrow[r] \arrow[d] & X' \arrow[d] \\
T \arrow[r] & T'
\end{tikzcd}
\end{center}
which induces the identity of $X_0$ after passage to the special
fibers.  A deformation $X/T$ is called \underline{versal}
(minimally versal) if:
\begin{enumerate}[label=(\alph*),leftmargin=2.6em]
\item for every morphism $X''/T''\longrightarrow X'/T'$ for which
$T''\to T'$ is a closed immersion, the map
\[
\operatorname{Hom}(X'/T',X/T)
\longrightarrow\operatorname{Hom}(X''/T'',X/T)
\]
is surjective.  In particular, on taking $X''/T''=X_0/s$, one has
$\operatorname{Hom}(X'/T',X/T)\ne\varnothing$;

\item for
$T'=\operatorname{Spec}(k(s)[\varepsilon]/(\varepsilon^2))$, the
``inverse-image'' map
\[
\operatorname{Hom}(T',T)\longrightarrow
\{\text{isomorphism classes of deformations of $X_0$ over $T'$}\}
\]
is bijective.
\end{enumerate}

\subsubsection*{Theorem 1.1.4 (R. Elkik).}
\underline{There exists a versal henselian $S$-deformation of $X_0$
(and it is unique up to a nonunique isomorphism).  If $S$ is
noetherian, a deformation $X/T$ is such a deformation if and only if
$X$ induces over the completion $\widehat T$ of $T$ a versal formal
$S$-deformation of $X_0$ (a minimally formal versal deformation;
VI 1.9 and VI 4).}

For the proof of this delicate theorem, we refer to R. Elkik
(ENS seminar, 1971/72).
\end{filecontents*}
% Original input: components/065_expose_XV_section_1_2_1_through_1_2_6_reduction.tex
\begin{filecontents*}[overwrite,noheader]{sga7ii-cumulative-input-065.tex}
\subsection*{1.2. Ordinary quadratic singularities}

Let $y$ be a closed point of a scheme $Y$ of finite type over a field
$k$ of characteristic $p$, and let $n$ be the dimension of $Y$ at
$y$.

\subsubsection*{Definition 1.2.1 (= VI 6.6).}
\begin{enumerate}[label=(\roman*),leftmargin=2.6em]
\item \underline{If $k$ is algebraically closed, $y$ is called an
ordinary quadratic point of $Y$ if the completion
$\widehat{\mathcal O}_{Y,y}$ of the local ring of $Y$ at $y$ is
isomorphic to the quotient of $k[[x_1,\ldots,x_{n+1}]]$ by the ideal
generated by a single formal power series $f$ which can be written}
\begin{equation}\tag{1.2.1.1}
f(\underline X)=Q(\underline X)+\text{terms of order }>2,
\end{equation}
\underline{where $Q(\underline X)$ is an ordinary quadratic form.}
\footnote{The printed display is numbered (4.2.1) and has an extra
equals sign after $Q(\underline X)$; the English restores the local
numbering and the asserted expansion $f=Q+$ higher-order terms.}

\item \underline{If $\bar k$ is an algebraic closure of $k$, $y$ is
called an ordinary quadratic point of $Y$ if the points of
$\bar Y=Y\otimes_k\bar k$ lying over $y$ are ordinary quadratic
points of $\bar Y$.}
\end{enumerate}

\subsubsection*{1.2.2.}
Replacing ``ordinary quadratic form'' by ``nondegenerate quadratic
form'' in this definition gives the notion of a
\underline{nondegenerate quadratic point} (VI 6.0).  A point $y$ is
a nondegenerate quadratic point of $Y$ if and only if it is an
ordinary quadratic point and either $p\ne2$ or $n$ is odd.

\subsubsection*{Example 1.2.3.}
Let $Q(\underline X)$ be a nondegenerate quadratic form in $n+1$
variables over $k$.  Let $Y_0$ be the quadratic cone in
$\mathbb A_k^{n+1}$ with equation $Q=0$.  Then the origin, denoted by
$y_0$, is a nondegenerate quadratic point of $Y_0$.

\subsubsection*{Example 1.2.4.}
Suppose that $p=2$ and $n=2m$.  Let $a\in k$, and let $Q$ be a
polynomial
\[
Q(\underline X)=(X_0^2-a)
 +\sum_{0<i\leq j\leq2m}a_{ij}x_ix_j.
\]
Suppose that the quadratic form in $2m$ variables
$\sum_{0<i\leq j}a_{ij}x_ix_j$ is nondegenerate, and denote by $y_0$
the closed point of $\mathbb A_k^{n+1}$ with coordinates
$(\sqrt a,0,\ldots,0)$.  Then $y_0$ is an ordinary quadratic point
of the closed subscheme $Y_0\subset\mathbb A_k^{n+1}$ with equation
$Q=0$.

\subsubsection*{Example 1.2.5.}
Let $k'$ be a finite separable extension of $k$, and let $Y_0$ be a
$k'$-scheme of the kind considered in 1.2.3 (respectively, 1.2.4).
If $Y_0$ is regarded as a $k$-scheme, then $y_0$ is a nondegenerate
(respectively, ordinary) quadratic point of $Y_0$.  The extension
$k'/k$ identifies with the largest separable subextension of
$k(y_0)/k$.

\subsubsection*{Theorem 1.2.6.}
\underline{Suppose that $y$ is an ordinary quadratic point of the
$k$-scheme $Y$, and let $k'$ be the largest separable subextension
of $k(y)/k$.  There then exist a $k'$-scheme
$Y_0\subset\mathbb A_{k'}^{n+1}$ of the kind in 1.2.3 or 1.2.4 and a
$k$-isomorphism $\varphi$ between the henselization $Y_{(y)}$ of $Y$
at $y$ and the henselization $Y_{0(y_0)}$ of $Y_0$ at $y_0$.}

First observe that there is an etale neighborhood $Y'$ of $y$ in
$Y$
\begin{center}
\begin{tikzcd}[row sep=large,column sep=large]
{} & Y'\arrow[d]\\
y\arrow[ur,hook]\arrow[r,hook]&Y
\end{tikzcd}
\end{center}
which is a $k'$-scheme.  Indeed, $Y_{(y)}$ is canonically a
$k'$-scheme.  Replacing $(k,Y,y)$ by $(k',Y',y)$ reduces us to the
case $k=k'$, that is, to the case in which $k(y)/k$ is radicial.
\end{filecontents*}
% Original input: components/066_expose_XV_section_1_2_7_through_1_2_12.tex
\begin{filecontents*}[overwrite,noheader]{sga7ii-cumulative-input-066.tex}
The $k(y)$-vector space $(\Omega^1_{Y/k})_y$ has dimension $n+1$.
It is enough to check this after extending scalars from $k(y)$ to its
algebraic closure and completing at $y$, where it is immediate.  A
neighborhood of $y$ in $Y$ is therefore a closed subscheme of a
smooth $k$-scheme $Z$ of dimension $n+1$.  Moreover, after replacing
$Y$ and $Z$ by suitable neighborhoods of $y$, the ideal in
$\mathcal O_Z$ defining $Y$ is principal (the same argument as
above), say with equation $f$.

\subsubsection*{Lemma 1.2.7.}
\underline{Suppose that $y$ is a nondegenerate quadratic point of
$Y$ and that $k(y)/k$ is radicial.  Then, in a neighborhood of $y$,
$J^n_{Y/k}$ (1.1.1) is generated by one element, supported at $y$,
and of rank one over $k$.  In particular, $k(y)=k$.}

It is enough to prove the first assertion after passing to an
algebraic closure of $k$ and completing at $y$.  It is then immediate
by direct calculation: with the notation of 1.2.1, Nakayama's lemma
shows that the $\partial f/\partial x_i$ (and hence, a fortiori, $f$
and the $\partial f/\partial x_i$) generate the maximal ideal.

Let us prove 1.2.6 when $y$ is nondegenerate.  We may and shall
suppose that $k(y)=k$, that $Y$ is defined by an equation $f$ in a
smooth scheme $Z$, and that there is an etale morphism
\[
x:Z\longrightarrow\mathbb A_k^{n+1}
\]
(a ``system of local coordinates $x_i$'') taking $y$ to $0$.  There
is then a nondegenerate quadratic form $Q$ in $n+1$ variables such
that, at $y$,
\[
f=Q(x_i)+\text{terms of order }>2.
\]
Let $Y_0$ be defined by $Q$ as in 1.2.3.\footnote{The print refers to
4.3 here; the normal form just invoked is Example 1.2.3.}

\underline{On $Y$}, the system $(x_i)$ is an ``approximate solution''
of the equation $Q(x_i)=0$: if $\mathfrak m$ is the ideal defining
$y$, then
\[
Q(x_i)\in\mathfrak m^3\qquad\text{on }Y.
\]
Since $Q$ is nondegenerate, the ideal of $\mathbb A_k^{n+1}$
generated by the $\partial Q/\partial x_i$ is the maximal ideal at
$0$, and its inverse image under
$x=(x_i):Y\to\mathbb A_k^{n+1}$ is $\mathfrak m$.  By the
implicit-function theorem 1.1.2 (applied with $a=\mathfrak m$ and
$\delta'=J$), there are functions $x'_i$ on the henselization
$Y_{(y)}$, congruent to the $x_i$ modulo $\mathfrak m^2$, which
satisfy $Q(x'_i)=0$.  They give the required isomorphism between
$Y_{(y)}$ and $Y_{0(y_0)}$.

\subsubsection*{Lemma 1.2.8.}
\underline{Suppose that $y$ is a degenerate ordinary quadratic point
($p=2$ and $n$ even), and that $k(y)/k$ is radicial.  Then, in a
neighborhood of $y$, $T^1_{Y/k}$ is generated by one element,
supported at $y$, and of rank two over $k$.  In particular, either
$k(y)=k$, or $k(y)$ is isomorphic to $k(\sqrt a)$ for some
$a\in k$, $a\notin k^2$.}

Proceed as in 1.2.7.  With the notation of 1.2.1, it is enough to
check that
\[
\dim_k\left(
k[[x_1,\ldots,x_{n+1}]]\big/
  (f,\partial f/\partial x_i)
\right)=2.
\]
In suitable coordinates $x_0,x_1,\ldots,x_{2m}$, the formal power
series $f$ can be written
\[
f=x_0^2+\sum_{i=1}^m x_ix_{i+m}+R,
\qquad\operatorname{ord}(R)\geq3.
\]
Let $\mathfrak n'$ be the ideal generated by $f$ and the
$\partial f/\partial x_i$, let $\mathfrak n$ be the ideal generated
by $x_0^2$ and the $x_i$ for $i\ne0$, and let $\mathfrak q$ be the
maximal ideal.  We prove that $\mathfrak n=\mathfrak n'$.  Plainly
$\mathfrak n'\subset\mathfrak n$, and by Nakayama it is enough to
check that
\[
\mathfrak n'+\mathfrak n\mathfrak q=\mathfrak n.
\]
Modulo $\mathfrak n\mathfrak q$, one has $f\equiv x_0^2$; modulo
$\mathfrak n\mathfrak q+(x_0^2)$, one has
\[
\frac{\partial f}{\partial x_i}\equiv x_{i+m},
\qquad
\frac{\partial f}{\partial x_{i+m}}\equiv x_i
\qquad(1\leq i\leq m),
\]
which proves the assertion.

Under the hypotheses of 1.2.8, if $k=k(y)$, we shall let $a$ denote
any square in $k$.

\subsubsection*{Lemma 1.2.9.}
\underline{Let $V$ be a vector bundle on a scheme $S$, and let $T$
be a closed subscheme of $\mathbb A_S(V)$ locally free of rank two
over $S$.  Then $T$ is contained in one and only one line in
$\mathbb A_S(V)$.}

The question is local for the fppf topology on $S$, so we may suppose
that $T/S$ has a section.  By translation, we may suppose that $T$
contains the zero section.  Let $\mathfrak n$ be the ideal of $T$
defining $0$; it is locally free of rank one over $S$.  Every
homogeneous linear form $v\in V$ on $\mathbb A_S(V)$ defines an
element of $\mathfrak n$, and hence an epimorphism
$V\twoheadrightarrow\mathfrak n$.  Its kernel defines the required
line.

\subsubsection*{1.2.10.}
In particular, if $A$ is a radicial rank-two closed subscheme of
affine space $\mathbb A_k^{n+1}$ over a field $k$ of characteristic
$2$, there is an $a\in k$ such that, in suitable affine coordinates,
$A$ has equations
\[
x_0^2-a=0,
\qquad x_i=0\quad(1\leq i\leq n).
\]
It is enough to take the $x_i$ which vanish on the line defined by
$A$.

\subsubsection*{1.2.11.}
Let us prove 1.2.6 for a degenerate $y$ ($p=2$ and $n=2m$).  We may
and shall suppose that $k(y)=k(\sqrt a)$, that $Y$ is defined by an
equation $f$ in a smooth scheme $Z$, and that there is a system of
local coordinates at $y$
\[
x=(x_i)_{0\leq i\leq2m}:
Z\longrightarrow\mathbb A_k^{2m+1}
\]
taking $y$ to the closed point $(\sqrt a,0,\ldots,0)$, such that, in
a neighborhood of $y$,
\[
J^n_{Y/k}\simeq
\mathcal O_Y/(x_0^2-a,x_1,\ldots,x_{2m}).
\]
After replacing $f$ by a suitable multiple, one has
\[
f=(x_0^2-a)+
 \sum_{1\leq i\leq j\leq2m}a_{ij}x_ix_j+R,
\]
where
\begin{enumerate}[label=(\alph*),leftmargin=2.6em]
\item the quadratic form in $2m$ variables
$\sum a_{ij}x_ix_j$ is nondegenerate;
\item after extending scalars to $\bar k$, $R$ has order at least
$3$ in $x_0-\sqrt a$ and the $x_i$.
\end{enumerate}
Let $Y_0$ be defined, as in 1.2.4, by
\[
Q=x_0^2-a+\sum a_{ij}x_ix_j.
\]

\underline{On $Y$}, the system $(x_i)$ may again be regarded as an
approximate solution of $Q(X_i)=0$.  This time, the ideal $\delta$,
the inverse image under $x$ of the ideal generated by the
$\partial Q/\partial X_i$, satisfies
$\mathcal O_Y/\delta\simeq J^n_{Y/k}$; at $y$ it is
\[
\delta=(x_0^2-a,(x_i)_{1\leq i}).
\]
Let $\mathfrak q$ be the maximal ideal defining $y$.  To conclude as
in the nondegenerate case by applying the implicit-function theorem,
it remains to verify that
\begin{equation}\tag{1.2.11.1}
Q(x_i)\equiv0\pmod{\delta^2\mathfrak q}
\qquad\underline{\text{on }Y}.
\end{equation}
Let $\bar{\mathfrak q}$ be the ideal defining $y$ after scalar
extension to an algebraic closure of $k$.  One checks that
\[
\delta^2\bar{\mathfrak q}\cap\mathcal O_{Y,y}
=\delta^2\mathfrak q.
\]
It is therefore enough to verify (1.2.11.1) after extending scalars
to $\bar k$; we may suppose that $a=0$.\footnote{The print refers at
this point to (4.11.1); the displayed congruence under proof is
(1.2.11.1).}
At $y$ one then has
\[
\delta^2\mathfrak q=
(x_0^5,(x_0^3x_i)_{1\leq i},
 (x_0x_ix_j)_{1\leq i,j},
 (x_ix_jx_k)_{1\leq i,j,k}).
\]
It follows that $\delta^2\mathfrak q\supset\mathfrak q^5$ and that
\[
x_0^2f\equiv x_0^4\pmod{\delta^2\mathfrak q},
\qquad
\delta^2\mathfrak q+(x_0^2f)\supset\mathfrak q^4,
\]
\[
x_0f\equiv x_0^3
 \pmod{\delta^2\mathfrak q+\mathfrak q^4},
\qquad
\delta^2\mathfrak q+(x_0f)\supset\mathfrak q^3.
\]
Finally,
\[
Q(x_i)\equiv0\pmod{\delta^2\mathfrak q+(f)}
\qquad\underline{\text{on }Z},
\]
which is (1.2.11.1).

\subsubsection*{Remark 1.2.12.}
Let $Q$ be a not necessarily homogeneous, nonzero quadratic form in
$n+1$ variables, and suppose that the closed subscheme
$Y\subset\mathbb A_k^{n+1}$ with equation $Q$ has an ordinary
quadratic singularity at a point $y$.  A nonhomogeneous linear change
of variables then puts $Y$ into the form 1.2.3 or 1.2.4.  In the
nondegenerate case $k(y)=k$, and it is enough to put $y$ at the
origin.  In the degenerate case it is enough that the support of
$T^1_{Y/k}$ lie on the line $x_i=0$ $(1\leq i\leq2m)$
(cf. 1.2.9).
\end{filecontents*}
% Original input: components/067_expose_XV_section_1_3.tex
\begin{filecontents*}[overwrite,noheader]{sga7ii-cumulative-input-067.tex}
\subsection*{1.3. Moduli of quadratic singularities}

The following proposition follows from R. Elkik's theorem 1.1.4 and
from explicit calculations left to the reader (and essentially
contained in VI 6).

\subsubsection*{Proposition 1.3.1.}
\underline{Let $S=\operatorname{Spec}(A)$ be a henselian local scheme
with closed point $s$, and let
$Y_0\subset\mathbb A_{k(s)}^{n+1}$ be a scheme of the kind described
in 1.2.3 or 1.2.4.}

\begin{enumerate}[label=(\roman*),leftmargin=2.6em]
\item \underline{Suppose that $Y_0$ is of the kind described in
1.2.3, defined by the equation}
\[
Q(\underline X)=\sum_{i\leq j}\bar a_{ij}x_i x_j=0
\qquad(\bar a_{ij}\in k(s)),
\]
\underline{and lift the $\bar a_{ij}$ to elements $a_{ij}$ of $A$.
Then the versal henselian $S$-scheme of deformations of the
singularity of $Y_0$ at $y_0$ is}
\[
T=\operatorname{Spec}(A\{b\}),
\]
\underline{the versal deformation being the henselization at $y_0$
of the $T$-scheme $Y\subset\mathbb A_T^{n+1}$ with equation}
\[
\sum a_{ij}x_i x_j-b=0.
\]

\item \underline{Suppose that $Y_0$ is of the kind described in
1.2.4, defined by the equation}
\[
Q(\underline x)=(x_0^2-\bar a)+\sum\bar a_{ij}x_i x_j=0
\qquad(\bar a,\bar a_{ij}\in k(s)),
\]
\underline{and lift the $\bar a_{ij}$ and $\bar a$ to elements
$a_{ij}$ and $a$ of $A$.  Then the versal henselian $S$-scheme of
deformations of the singularity of $Y_0$ at $y_0$ is}
\[
T=\operatorname{Spec}(A\{b,c\}),
\]
\underline{the versal deformation being the henselization at $y_0$
of the $T$-scheme $Y\subset\mathbb A_T^{n+1}$ with equation}
\[
(x_0^2-a)+\sum a_{ij}x_i x_j+b x_0+c=0.
\]
\end{enumerate}

The following corollary follows at once from 1.2.6 and 1.3.1.

\subsubsection*{Corollary 1.3.2.}
\underline{Let $S=\operatorname{Spec}(A)$ be a henselian local
scheme with closed point $s$, and let $f:X\to S$ be a flat morphism
of finite presentation.  Suppose that, at the closed point $x$ of
$X_s$, the scheme $X_s$ has an ordinary quadratic singularity, that
$k(x)$ is radicial over $k(s)$, and that $X_s$ has dimension $n$.}

\begin{enumerate}[label=(\roman*),leftmargin=2.6em]
\item If $x$ is nondegenerate, \underline{there exist a
nondegenerate quadratic form $Q$ in $n+1$ variables with
coefficients in $A$, and an element $b$ in the maximal ideal of $A$,
such that the henselized $S$-scheme of $X$ at $x$ is isomorphic to
the henselization at the origin of the closed subscheme of
$\mathbb A_S^{n+1}$ with equation $Q-b=0$, as in 1.3.1(i).}

\item If $x$ is degenerate
($n=2m$, $\operatorname{char}(k(s))=2$), \underline{there exists a
nonhomogeneous quadratic form in $n+1$ variables with coefficients
in $A$}
\begin{equation}\tag{1.3.2.1}
Q(\underline X)=x_0^2+b x_0+c+
 \sum_{1\leq i\leq j\leq 2m}a_{ij}x_i x_j,
\end{equation}
\underline{such that $b$ lies in the maximal ideal, the quadratic
form $\sum a_{ij}x_i x_j$ in $2m$ variables is nondegenerate, and
the henselized $S$-scheme of $X$ at $x$ is isomorphic to the
henselization at $(\sqrt c,0,\ldots,0)$ of the closed subscheme of
$\mathbb A_S^{n+1}$ with equation $Q=0$.}
\footnote{The print labels this display (4.14.2) and writes $x_0$ in
its first term.  The surrounding normal forms and Remark 1.3.3
require the local tag (1.3.2.1) and the quadratic term $x_0^2$; the
English restores both.}
\end{enumerate}

\subsubsection*{Remark 1.3.3.}
In case (i), suppose that the henselization of $X_s$ at $x$ is
isomorphic, by an isomorphism $\varphi$, to the henselization at $O$
of the closed subscheme
$X_0\subset\mathbb A_k^{n+1}$ with equation
\[
\sum\bar a_{ij}x_i x_j=0,
\]
and suppose that the $a_{ij}$ lift the $\bar a_{ij}$.  One may then
take
\[
Q=\sum a_{ij}x_i x_j.
\]

In case (ii), suppose that the henselization of $X_s$ at $x$ is
isomorphic, by an isomorphism $\varphi$, to the henselization at
$(\sqrt a,0,\ldots,0)$ of the closed subscheme
$X_0\subset\mathbb A_k^{n+1}$ with equation
\[
(x_0^2+\bar a)+\sum\bar a_{ij}x_i x_j=0,
\]
and suppose that the $a_{ij}$ lift the $\bar a_{ij}$.  One may then
take
\[
Q=x_0^2+b x_0+c+\sum a_{ij}x_i x_j,
\qquad c\equiv\bar a\pmod{\mathfrak m_A}.
\]
In both cases, the isomorphism promised in 1.3.2 may be chosen to
extend $\varphi$.

\subsubsection*{Corollary 1.3.4.}
\underline{Let $f:X\to S$ be as in 1.3.2.  There is a neighborhood
$U$ of $x$ in $X$ such that the nonsmooth points of $f$ contained in
$U$ are ordinary quadratic points of their respective fibers.}
\end{filecontents*}
% Original input: components/068_expose_XV_section_2_1_1_through_lemma_2_1_3.tex
\begin{filecontents*}[overwrite,noheader]{sga7ii-cumulative-input-068.tex}
\section*{2. Computation of vanishing cycles in a standard ordinary
quadratic case}

\subsection*{2.1. Preliminary: the cohomology of a cone}

\subsubsection*{2.1.1.}
Let $Y\subset\mathbb P^r$ be a projective variety over an
algebraically closed field $k$.  We use the following notation.

The affine cone over $Y$ is denoted by
$X\subset\mathbb A^{r+1}$, and its vertex by $O$;

$X_{(O)}$ is the henselization of $X$ at $O$;

$X^*=X-\{O\}$ and $X^*_{(O)}=X_{(O)}-\{O\}$;

$X_1\subset\mathbb P^{r+1}$ is the projective cone over $Y$, so
$X=X_1-Y$;

$\widetilde X$, $\widetilde X_{(O)}$, and $\widetilde X_1$ are the
blowups of $X$, $X_{(O)}$, and $X_1$ at $O$;

$Y_0$ is the inverse image of $O$ in $\widetilde X$,
$\widetilde X_{(O)}$, or $\widetilde X_1$;

$h:\widetilde X_1\to Y$ is the natural projection;

$y_0,y_\infty:Y\to\widetilde X_1$ are the sections of $h$ whose
images are $Y_0$ and $Y$;

$F$ is an abelian torsion group of order prime to the characteristic
of $k$.  All cohomology groups considered below have coefficients in
$F$, which will not be displayed.

Our purpose in this subsection is to relate the cohomology of
$X^*_{(O)}$ to that of $Y$.

\subsubsection*{Proposition 2.1.2.}
\underline{There are isomorphisms}
\[
\text{(i)}\qquad
H^i(X)\xrightarrow{\ \sim\ }H^i(\{O\})=
\begin{cases}
F,& i=0,\\
0,& i>0;
\end{cases}
\]
\[
\text{(ii)}\qquad
H^i_{\{O\}}(X)\xrightarrow{\ \sim\ }H^i_c(X).
\]
This proposition is proved by a homotopy argument, which we now
formalize.

Let $k$ be an algebraically closed field, let $U$ and $V$ be two
$k$-schemes, let $\Lambda$ be a torsion ring of order prime to the
characteristic of $k$, and let
\[
K\in\operatorname{Ob}D^+(U,\Lambda),
\qquad L\in\operatorname{Ob}D^+(V,\Lambda).
\]
By a ``morphism'' $f$ from $(U,K)$ to $(V,L)$ we mean a pair
consisting of a morphism $f:U\to V$ and a morphism
$\varphi:f^*L\to K$.  Such a morphism induces
\[
f^*:H^*(V,L)\longrightarrow H^*(U,K).
\]

We say that two ``morphisms'' $f_0$ and $f_1$ are homotopic if there
exist a connected $k$-scheme $T$ of finite type, two points $0$ and
$1$ of $T$, and a ``morphism'' $H$ from
$(U\times T,\operatorname{pr}_1^*K)$ to $(V,L)$ whose fibers at $0$
and $1$ induce $f_0$ and $f_1$.

\subsubsection*{Lemma 2.1.3.}
\underline{If $f_0$ is homotopic to $f_1$, then
$f_0^*=f_1^*$.}

There is a sequence of points
$0=x_0,x_1,\ldots,x_n=1$ of $T$, smooth connected curves
$\Gamma_i$, and morphisms $a_i:\Gamma_i\to T$ such that $x_i$ and
$x_{i+1}$ belong to the image of $a_i$.  To construct the
$\Gamma_i$, normalize a connected one-dimensional closed subscheme
of $T$ containing $0$ and $1$.  This reduces us to the case in which
$T$ is a smooth curve.  Consider the commutative square
\begin{center}
\begin{tikzcd}[row sep=large,column sep=large]
U\times T \arrow[r,"\operatorname{pr}_1"]
  \arrow[d,"\operatorname{pr}_2"']
  & U \arrow[d,"f"] \\
T \arrow[r,"t"'] & k .
\end{tikzcd}
\end{center}
By smooth base change, applied to $t$, one has
\[
t^*Rf_*K\xrightarrow{\ \sim\ }
R\operatorname{pr}_{2*}(\operatorname{pr}_1^*K).
\]
The map $f_i^*$ $(i=0,1)$ therefore factors through
\begin{center}
\begin{tikzcd}[row sep=large,column sep=large]
H^n(V,L) \arrow[r,"H^*"]
  & H^n(U\times T,\operatorname{pr}_1^*K) \arrow[r]
  & H^0\bigl(T,R^n\operatorname{pr}_{2*}
      (\operatorname{pr}_1^*K)\bigr)
      \arrow[dd,"\text{fiber at }i"] \\
{} & H^0\bigl(T,t^*H^n(U,K)\bigr) \arrow[ur,"\wr"] & {} \\
{} & {} & H^n(U,K) \arrow[ul,"\wr"']
\end{tikzcd}
\end{center}
and is consequently independent of $i$.
\end{filecontents*}
% Original input: components/069_expose_XV_section_2_1_3_through_2_1_6.tex
\begin{filecontents*}[overwrite,noheader]{sga7ii-cumulative-input-069.tex}
\subsubsection*{2.1.3.}
Let us prove 2.1.2.  Scalar multiplication gives a homotopy between
the identity map of $X$ and the constant map from $X$ to $X$ with
value $O$; this proves (i).

Using scalar multiplications whose factors tend to infinity, one
sees that $Y$ is a deformation retract of $X_1^*$.  Assertion (ii)
then follows from the five lemma and the diagram
\begin{center}
\begin{tikzcd}[row sep=large,column sep=large]
{} \arrow[r,"\partial"]
  & H^i_{\{O\}}(X) \arrow[r] \arrow[d]
  & H^i(X_1) \arrow[r] \arrow[d,equal]
  & H^i(X_1^*) \arrow[r,"\partial"] \arrow[d,"\wr"']
  & {} \\
{} \arrow[r,"\partial"]
  & H^i_c(X) \arrow[r]
  & H^i(X_1) \arrow[r]
  & H^i(Y) \arrow[r,"\partial"]
  & {} .
\end{tikzcd}
\end{center}

\subsubsection*{Corollary 2.1.4.}
\[
H^i(X^*)\xrightarrow{\ \sim\ }H^i(X^*_{(O)}).
\]
Indeed, apply the five lemma to
\begin{center}
\begin{tikzcd}[row sep=large,column sep=large]
{} \arrow[r]
  & H^i_{\{O\}}(X) \arrow[r] \arrow[d,"\wr"]
  & H^i(X) \arrow[r]
      \arrow[d,"\wr"',"{(2.1.2(i))}"]
  & H^i(X^*) \arrow[r] \arrow[d]
  & {} \\
{} \arrow[r]
  & H^i_{\{O\}}(X_{(O)}) \arrow[r]
  & H^i(X_{(O)}) \arrow[r]
  & H^i(X^*_{(O)}) \arrow[r]
  & {} .
\end{tikzcd}
\end{center}

\subsubsection*{2.1.5.}
The Leray spectral sequences for the restrictions of $h$ to
$\widetilde X_1-Y_0$ and $\widetilde X_1-Y$ show that the restriction
maps
\[
H^i(\widetilde X_1-Y_0)\xrightarrow{\ \sim\ }H^i(Y),
\qquad
H^i(\widetilde X)\xrightarrow{\ \sim\ }H^i(Y_0)=H^i(Y)
\]
are isomorphisms.  Under these isomorphisms, the long exact
cohomology sequences of the pairs
$(\widetilde X_1-Y_0,Y)$ and $(\widetilde X_1-Y,Y_0)$ form the rows
of the diagram
\begin{equation}\tag{2.1.5.1}
\begin{tikzcd}[column sep=large]
H^{i-1}(X^*) \arrow[r] \arrow[d,equal]
  & H^{i-2}(Y)(-1) \arrow[r,"\eta"'] \arrow[d,equal]
  & H^i(Y) \arrow[r] \arrow[d,dash,"-1"]
  & H^i(X^*) \arrow[r] \arrow[d,equal]
  & {} \\
H^{i-1}(X^*) \arrow[r]
  & H^{i-2}(Y)(-1) \arrow[r,"-\eta"']
  & H^i(Y) \arrow[r]
  & H^i(X^*) \arrow[r]
  & {} .
\end{tikzcd}
\end{equation}

\subsubsection*{Lemma 2.1.6.}
\underline{The arrows denoted by $\eta$ and $-\eta$ in (2.1.5.1)
are cup products with the cohomology class of a hyperplane section.
Diagram (2.1.5.1) is commutative.}

The first assertion follows because $\eta$ (respectively $-\eta$)
is the restriction to $Y$ (respectively to $Y_0\simeq Y$) of the
cohomology class of the invertible sheaf $\mathcal O(Y)$
(respectively $\mathcal O(Y_0)$) on $\widetilde X_1-Y_0$
(respectively $\widetilde X_1-Y$).

Use the Leray spectral sequence for $h$.  The two rows of (2.1.5.1)
are obtained by applying the functor $H(Y,-)$ to the distinguished
triangles
\begin{center}
\begin{tikzcd}[column sep=large]
{} & y_{\infty*}Ry_\infty^!F \arrow[r]
  & R(h|\widetilde X_1-Y_0)_*F
      \arrow[r,"\textcircled{\scriptsize 1}"]
  & R(h|X^*)_*F \arrow[r,"\textcircled{\scriptsize 2}"]
  & {} \\
\text{and}
  & y_{0*}Ry_0^!F \arrow[r]
  & R(h|\widetilde X_1-Y)_*F \arrow[r]
  & R(h|X^*)_*F \arrow[r]
  & {} .
\end{tikzcd}
\end{center}
The nonzero cohomology sheaves of the complexes at the vertices of
these triangles occur in degrees $(2)$, $(0)$, and $(0$ and $1)$.
The arrows \textcircled{\scriptsize 1} and
\textcircled{\scriptsize 2} are therefore determined by their effect
on cohomology sheaves.  This reduces us to the case in which $Y$ is
a point, which is left to the reader.
\end{filecontents*}
% Original input: components/070_expose_XV_section_2_1_7_through_2_1_8.tex
\begin{filecontents*}[overwrite,noheader]{sga7ii-cumulative-input-070.tex}
\subsubsection*{2.1.7.}
The long exact sequence of the pair
$(\widetilde X_{(O)},Y_0)$ is a local analogue of the second row of
(2.1.5.1).  One has
\[
H^i(\widetilde X_{(O)})\xrightarrow{\ \sim\ }H^i(Y_0)
\]
by proper base change for the morphism
$\widetilde X_{(O)}\to X_{(O)}$.  The long exact sequence in question
may therefore be written
\begin{equation}\tag{2.1.7.1}
\begin{tikzcd}[column sep=large]
H^{i-1}(X^*_{(O)}) \arrow[r]
  & H^{i-2}(Y_0)(-1) \arrow[r,"-\eta"]
  & H^i(Y_0) \arrow[r]
  & H^i(X^*_{(O)}) & {} \\
H^{i-1}(X^*) \arrow[r] \arrow[u,"\wr"]
  & H^{i-2}(Y)(-1) \arrow[r,"-\eta"] \arrow[u,equal]
  & H^i(Y) \arrow[r] \arrow[u,equal]
  & H^i(X^*) \arrow[u,dash,"\wr"]
  & {} ,
\end{tikzcd}
\end{equation}
the second row of this commutative diagram being the second row of
(2.1.5.1).
\footnote{The printed tag omits the dot between 1 and 7; the English
uses the local number (2.1.7.1).}

The following technical lemma will be used in 2.2.7.

\subsubsection*{Lemma 2.1.8.}
\begin{center}
\begin{tikzcd}[column sep=large]
H^{n-1}(Y_0) \arrow[d,equal]
  & H^{n-1}(\widetilde X) \arrow[l,"\sim"'] \arrow[r]
  & H^{n-1}(X^*) \arrow[r]
  & H^n_{\{x\}}(X) \arrow[r]
  & H^n_c(X_s) \arrow[d,equal] \\
H^{n-1}(Y) \arrow[rrrr,"\partial"]
  & {} & {} & {} & H^n_c(X_s)
\end{tikzcd}
\end{center}
\underline{is anticommutative.}
\footnote{The print numbers this lemma 2.7.8.  Its position and the
forward reference in 2.2.7 require 2.1.8.}

By 2.1.7, this is equivalent to proving the commutativity of
\begin{center}
\begin{tikzcd}[column sep=large,row sep=large]
H^{n-1}(X_1-\{O\}) \arrow[r,"\partial"] \arrow[d,"\wr"]
  & H^n_{\{O\}}(X_1)=H^n_{\{O\}}(X) \arrow[d] \\
H^{n-1}(Y) \arrow[r,"\partial"]
  & H^n_c(X).
\end{tikzcd}
\end{center}
This is general nonsense.
\end{filecontents*}
% Original input: components/071_expose_XV_section_2_2_1_through_2_2_4.tex
\begin{filecontents*}[overwrite,noheader]{sga7ii-cumulative-input-071.tex}
\subsection*{2.2. Vanishing cycles}

\subsubsection*{2.2.1.}
Let $S$ be a henselian trait, and let
$s,\eta,\bar s,\bar\eta$ have the meaning of XIII, 0.2.5.  Put
$\Lambda=\mathbb Z/k$, where $k$ is prime to the residual
characteristic of $S$.  Let $\mathbb A^{n+1}_S$ be standard affine
$(n+1)$-space over $S$, and let $X$ be the closed subscheme of
$\mathbb A^{n+1}_S$ defined by a quadratic equation $Q=0$ which is
nonzero modulo a uniformizer:
\[
Q(\underline x)=\sum_{i\leq j}a_{ij}x_ix_j+\sum_i b_ix_i+c.
\]
Regard $\mathbb A^{n+1}_S$ as the complement in
$\mathbb P^{n+1}_S$ of the hyperplane at infinity $H$.  Then $X$ is
obtained from the quadric $X_1\subset\mathbb P^{n+1}_S$ with
homogeneous equation
\begin{equation}\tag{2.2.1.1}
\sum_{i\leq j}a_{ij}x_ix_j+\sum_i b_ix_iz+cz^2=0
\end{equation}
by removing $Y=X_1\cap H$.

We make the following assumptions:

\begin{enumerate}
\item[(a)] $Y$ is a smooth quadric over $S$, that is, the quadratic
form $\sum_{i\leq j}a_{ij}x_ix_j$ is ordinary;

\item[(b)] the quadric $X_s$ is singular, that is,
$X_{\bar s}$ is a quadratic cone.
\end{enumerate}
\footnote{The print has $X'_s$ in (b).  The equivalent clause and the
immediately following definition of the singular point show that
$X_s$ is intended.}

Let $x_0$ denote the singular point of $X_s$, and let
$f:X\to S$ be the projection.

\subsubsection*{Remark 2.2.2.}
If $x_0$ is a rational point, a translation puts it at the origin of
the coordinates.  If $x_0=0$, the $b_i$ and $c$ are divisible by a
uniformizer.

The subscheme $A$ of $X_s$ defined by the equations
$\partial Q/\partial x_i$ is concentrated at $x_0$.  Except in the
exceptional case where $k(s)$ has characteristic $2$ and $n+1$ is
even, this scheme has degree one over $k(s)$, and hence $x_0$ is
rational.  In the exceptional case it is ramified over $k(s)$ and
has rank $2$, so that $k(x_0)$ is either $k(s)$ or a purely
inseparable quadratic extension of $k(s)$.  The assertions about
$A$ are checked most easily after passing to the algebraic closure
$k(\bar s)$ of $k(s)$.

\subsubsection*{Proposition 2.2.3.}
\underline{Under the assumptions of 2.2.1, the morphisms}
\[
H^i(X_{\bar\eta},\Lambda)
 \xrightarrow{\ \sim\ }
\mathbb H^i\bigl(X_{\bar s},R\psi_{\bar\eta}(\Lambda)\bigr)
 \xrightarrow{\ \sim\ }
R^i\psi_{\bar\eta}(\Lambda)_{x_0}
\]
\[
H_c^i(X_{\bar\eta},\Lambda)
 \xleftarrow{\ \sim\ }
\mathbb H_c^i\bigl(X_{\bar s},R\psi_{\bar\eta}(\Lambda)\bigr)
 \xleftarrow{\ \sim\ }
\mathbb H^i_{\{x_0\}}\bigl(X_{\bar s},R\psi_{\bar\eta}(\Lambda)\bigr)
\]
\underline{are isomorphisms.}

That the left-hand arrows are isomorphisms follows from 2.1.8.6 and
2.1.10.5; the assertion for the right-hand arrows is a corollary of
2.1.2.  Indeed, there is a distinguished triangle
\[
(\Lambda\text{ on }X_{\bar s})[0]
 \longrightarrow R\psi_{\bar\eta}(\Lambda)
 \longrightarrow R\Phi_{\bar\eta}(\Lambda),
\]
and 2.1.2 applies to $(\Lambda\text{ on }X_{\bar s})$, whereas
$R\Phi_{\bar\eta}(\Lambda)$ is supported at $x_0$.

\subsubsection*{Corollary 2.2.4.}
\underline{If the geometric generic fiber $X_{\bar\eta}$ is
singular, that is, if it is still a quadratic cone, then all the
sheaves $R^i\Phi(\Lambda)$ vanish.}

By 2.1.2,
\[
\Lambda=H^0(X_{\bar s},\Lambda)
 \xrightarrow{\ \sim\ }H^0(X_{\bar\eta},\Lambda),
\qquad
0=H^i(X_{\bar s},\Lambda)=H^i(X_{\bar\eta},\Lambda)
\quad(i>0).
\]
We then apply the long exact sequence XIII, 2.1.8.9 (valid by
XIII, 2.1.9).
\footnote{The print uses the letter $o$ in the two occurrences of
$H^o$; the cohomological degree is $0$.}
\end{filecontents*}
% Original input: components/072_expose_XV_section_2_2_5_A_through_C.tex
\begin{filecontents*}[overwrite,noheader]{sga7ii-cumulative-input-072.tex}
\subsubsection*{2.2.5.}
Suppose now that
\[
(*)\qquad\text{the generic fiber $X_\eta$ is smooth.}
\]
For simplicity we shall also assume that $S$ is strictly henselian.
The sheaves of vanishing cycles are easily computed with the aid of
2.2.3 and XII, 3.7.  The cohomology with coefficients in $\Lambda$
considered here follows from the $\ell$-adic cohomology of the cited
passage by the universal-coefficient formula (cf. 2.1.13).  The
results are as follows.

\medskip
\noindent\underline{A. Cohomology without supports.}

One has $R^i\psi_{\bar\eta}(\Lambda)=0$ for $i\ne0,n$.

\begin{enumerate}
\item[1)] \underline{If $n\ne0$:}
$\psi_{\bar\eta}(\Lambda)=\Lambda$.  Moreover,
$R^n\psi_{\bar\eta}(\Lambda)$ is (noncanonically) isomorphic to the
sheaf $\Lambda$ on $x_0$, extended by zero to $X_s$.

\item[2)] \underline{For arbitrary $n$:}
even when $n=0$, one checks that
$R^i\Phi_{\bar\eta}(\Lambda)=0$ for $i\ne n$, and that
$R^n\Phi_{\bar\eta}(\Lambda)$ is (noncanonically) isomorphic to the
sheaf $\Lambda$ on $x_0$, extended by zero to $X_s$.
\end{enumerate}

\medskip
\noindent\underline{B. Cohomology with support at $x_0$.}

One has
\[
H^i_{\{x_0\}}\bigl(X_s,R\psi_{\bar\eta}(\Lambda)\bigr)=0
\qquad(i\ne n,2n).
\]
There is, moreover, a trace morphism
\begin{equation}\tag{2.2.5.1}
\operatorname{tr}:
H^{2n}_{\{x_0\}}\bigl(X_s,R\psi_{\bar\eta}(\Lambda(n))\bigr)
\longrightarrow\Lambda.
\end{equation}
\underline{If $n\ne0$:} the morphism (2.2.5.1) is an isomorphism.
Moreover,
\[
H^n_{\{x_0\}}\bigl(X_s,R\psi_{\bar\eta}(\Lambda(n))\bigr)
\]
is (noncanonically) isomorphic to $\Lambda$.

\medskip
\noindent\underline{C. Duality.}
For
\[
a\in R^n\psi_{\bar\eta}(\Lambda)_{x_0},
\qquad
b\in H^n_{\{x_0\}}\bigl(X_s,R\psi_{\bar\eta}(\Lambda(n))\bigr),
\]
put $(a,b)=\operatorname{Tr}(a\wedge b)$, where $\operatorname{Tr}$
is the morphism XIII, 2.2.6
\[
\operatorname{Tr}:
H^{2n}_{\{x_0\}}\bigl(X_s,R\psi_{\bar\eta}(\Lambda(n))\bigr)
\longrightarrow\Lambda.
\]
It follows from Poincar\'e duality on $X_{\bar\eta}$ and the
isomorphisms 2.2.3 that $(a,b)$ gives a perfect duality between the
free $\Lambda$-modules
\[
R^n\psi_{\bar\eta}(\Lambda)_{x_0}
\quad\text{and}\quad
H^n_{\{x_0\}}\bigl(X_s,R\psi_{\bar\eta}(\Lambda(n))\bigr).
\]
\footnote{The print drops the bar from $\bar\eta$ in the displayed
trace map; the surrounding formulas determine the indicated term.}
\end{filecontents*}
% Original input: components/073_expose_XV_section_2_2_5_D_through_F.tex
\begin{filecontents*}[overwrite,noheader]{sga7ii-cumulative-input-073.tex}
\medskip
\noindent\underline{D. $n=2m>0$.}

By XII, 3.7, the groups
\[
H^n_{\{x_0\}}\bigl(X_s,R\psi_{\bar\eta}(\Lambda(m))\bigr)
\quad\text{and}\quad
R^n\psi_{\bar\eta}(\Lambda(m))_{x_0}
\]
have natural generators, defined up to sign; denote them by
$\delta$ and $\delta'$, respectively.  Let $\phi$ be the natural map
from the first of these groups to the second.  The signs of $\delta$
and $\delta'$ may be normalized so that
\begin{equation}\tag{2.2.5.2}
(\delta',\delta)=1;
\end{equation}
then
\begin{equation}\tag{2.2.5.3}
\phi(\delta)=(-1)^m\,2\,\delta'.
\end{equation}

Let $Z$ be the separable quadratic extension of $k(\eta)$ which is
the center of the even part of the Clifford algebra of the quadratic
form (2.2.1.1), and let $\epsilon$ be the corresponding character,
of order $1$ or $2$, of the inertia group $I$.  For $\sigma\in I$,
one has
\begin{equation}\tag{2.2.5.4}
\sigma\delta=\epsilon(\sigma)\delta
\end{equation}
and
\begin{equation}\tag{2.2.5.5}
\sigma\delta'=\epsilon(\sigma)\delta'.
\end{equation}
Formulas (2.2.5.2), (2.2.5.3), and (2.2.5.4), in cohomology with
coefficients in $\mathbb Z/2k$, imply that the variation is
\begin{equation}\tag{2.2.5.6}
\operatorname{Var}(\sigma)(a)=
\frac{\epsilon(\sigma)-1}{2}\,(-1)^m\,(a,\delta)\,\delta.
\end{equation}

\medskip
\noindent\underline{E. $n=2m+1$---Vanishing cycles.}

By XII, 3.7, the groups
\[
H^n_{\{x_0\}}\bigl(X_s,R\psi_{\bar\eta}(\Lambda(m))\bigr)
\quad\text{and}\quad
R^n\psi_{\bar\eta}(\Lambda(m+1))_{x_0}
\]
have natural generators, defined up to sign; denote them by
$\delta$ and $\delta'$, respectively.  With the notation of D, their
signs may be normalized so that
\begin{equation}\tag{2.2.5.7}
(\delta',\delta)=1;
\end{equation}
then
\begin{equation}\tag{2.2.5.8}
\phi(\delta)=0.
\end{equation}

\medskip
\noindent\underline{F. $n=2m+1$---The variation.}

The inertia group $I$ acts trivially on the cohomology of $X_\eta$.
This follows from XII, 3.7 and from the fact that $I$ acts trivially
on the cohomology of $Y_\eta$, since $Y$ is proper and smooth over
$S$.

For $\sigma$ in the inertia group one necessarily has
\begin{equation}\tag{2.2.5.9}
D(\sigma)(\delta')=\lambda(\sigma)\delta,
\qquad \lambda(\sigma)\in\Lambda(1).
\end{equation}
The additivity formula (XIII, 1.4.3.4) becomes
\[
\lambda(\sigma\tau)=\lambda(\sigma)+\lambda(\tau).
\]
Let
\[
\epsilon:I\longrightarrow\Lambda(1)=\mu_k
\]
be the homomorphism such that
\[
\sigma\bigl(\sqrt[k]{t}\bigr)
=\epsilon(\sigma)\sqrt[k]{t}
\]
for a uniformizer $t$.  In view of the structure of the tame inertia
group, $\lambda$ is necessarily a multiple of $\epsilon$.  Formula
(2.2.5.9) may therefore be rewritten
\begin{equation}\tag{2.2.5.10}
\operatorname{Var}(\sigma)(a)
=\lambda_X\epsilon(\sigma)(a,\delta)\delta
\end{equation}
for a suitable $\lambda_X\in\Lambda$ (depending on $X/S$).  We shall
determine $\lambda_X$ later.
\footnote{In (2.2.5.6) and (2.2.5.10), the print suppresses the comma
in $(a,\delta)$; this is the pairing defined in C.}
\end{filecontents*}
% Original input: components/074_expose_XV_section_2_2_6_through_2_2_7.tex
\begin{filecontents*}[overwrite,noheader]{sga7ii-cumulative-input-074.tex}
\subsubsection*{2.2.6.}
We shall give the cycles $\delta$ and $\delta'$ of D and E a
description which is local for the \`etale topology.

The case of even $n$ is easy.  If $2\nmid k$
($\Lambda=\mathbb Z/k$), then $\pm\delta$ is characterized by
(2.2.5.3) and (2.2.5.4), that is, by
\begin{equation}\tag{2.2.6.1}
(\delta,\delta)=(-1)^m\,2.
\end{equation}
In general, $\pm\delta$ is the image of any cycle
$\widetilde\delta$, in cohomology with coefficients in
$\mathbb Z/2^ak$ for $a$ sufficiently large, which satisfies
(2.2.6.1).

Suppose that $n$ is odd.  Then $x_0$ is a rational point (2.2.2).
Let $X_{s(0)}$ be the henselization of $X_s$ at $x_0$, let
$\widetilde X_{s(0)}$ be its blowup at $x_0$, let $Y_0$ be the
inverse image of $x_0$ in $\widetilde X_{s(0)}$ (a quadric), and put
\[
X^*_{s(0)}=X_{s(0)}-\{x_0\}
=\widetilde X_{s(0)}-Y_0.
\]
Consider the composite map (cf. XIII, 2.4.6.1)
\begin{equation}\tag{2.2.6.2}
\begin{tikzcd}[column sep=huge,row sep=huge]
H^{n-1}(Y_0,\Lambda(m))
  & H^{n-1}(\widetilde X_{s(0)},\Lambda(m))
      \arrow[l,"\sim"'] \arrow[r]
  & H^{n-1}(X^*_{s(0)},\Lambda(m)) \arrow[d] \\
{}
  & H^n_{\{x_0\}}\bigl(X_s,R\Psi_{\bar\eta}(\Lambda(m))\bigr)
  & H^n_{\{x_0\}}(X_s,\Lambda(m)) \arrow[l]
\end{tikzcd}
\end{equation}

\subsubsection*{Lemma 2.2.7.}
\underline{The map (2.2.6.2) identifies
$H^n_{\{x_0\}}(X_s,R\Psi_{\bar\eta}(\Lambda(m)))$ with the
primitive quotient of the middle-dimensional cohomology of the
$2m$-dimensional quadric $Y_0$; the two classes $\pm\delta$ are the
images of the natural generators of this primitive quotient.}

The cycles $\pm\delta$ are characterized by their images in
$H_c^n(X_{\bar\eta},\Lambda(m))$.  Using (2.1.7.1), it is therefore
enough to prove that the composite map
\begin{center}
\begin{tikzcd}[column sep=huge,row sep=huge]
H^{n-1}(Y_0,\Lambda(m))
  & H^{n-1}(\widetilde X_s,\Lambda(m))
      \arrow[l,"\sim"'] \arrow[r]
  & H^{n-1}(X_s^*,\Lambda(m)) \arrow[d] \\
H_c^n(X_{\bar\eta},\Lambda(m))
  & H_c^n(X_s,\Lambda(m)) \arrow[l,"\operatorname{sp}"']
  & H^n_{\{x_0\}}(X_s,\Lambda(m)) \arrow[l]
\end{tikzcd}
\end{center}
sends the distinguished generators of the primitive quotient of the
first group to the distinguished generators of the last group.

Apply 2.1.8.  Up to sign, the map may be rewritten
\begin{center}
\begin{tikzcd}[column sep=huge,row sep=huge]
H^{n-1}(Y_s,\Lambda(m)) \arrow[r,"\partial"] \arrow[d,equal]
  & H_c^n(X_s,\Lambda(m)) \arrow[d] \\
H^{n-1}(Y_{\bar\eta},\Lambda(m)) \arrow[r]
  & H_c^n(X_{\bar\eta},\Lambda(m)),
\end{tikzcd}
\end{center}
and 2.2.7 follows.
\footnote{In this proof the print has $H_c^i$ for the group in which
the degree-$n$ cycle $\delta$ is characterized; the displayed maps
and all surrounding terms require $H_c^n$.  The print also abbreviates
the support point in (2.2.6.2) as $x$; the English uses the fixed
singular point $x_0$.}
\end{filecontents*}
% Original input: components/075_expose_XV_section_3_1.tex
\begin{filecontents*}[overwrite,noheader]{sga7ii-cumulative-input-075.tex}
\section*{3. The Picard--Lefschetz formula}

\subsection*{3.1. Preliminary results}

\subsubsection*{3.1.1.}
Let $S$ be a henselian trait and let $f:X\to S$ be a flat morphism
of finite type, purely of relative dimension $n$.  Suppose that the
special fiber $X_s$ is smooth except at a finite set of points
$\Sigma'$, at which $X_s$ has an ordinary quadratic singularity.

Let $\Sigma$ be the set of points $x\in\Sigma'$ such that the
generic fiber $X_\eta$ is smooth in a neighborhood of $x$.  Finally,
let $\Lambda$ be a torsion ring whose order is prime to the residual
characteristic of $S$.

\subsubsection*{Proposition 3.1.2.}
\begin{enumerate}
\item[(i)] \underline{The sheaves $R^i\Phi(\Lambda)$ vanish for
$i\ne n$.}

\item[(ii)] \underline{The sheaf $R^n\Phi(\Lambda)$ is zero outside
$\Sigma$.  Its restriction to $\Sigma$ is a sheaf of
$\Lambda$-modules of rank one.}

\item[(iii)] \underline{For every geometric point $\bar x$ of
$\Sigma$, the pairing $(a,b)$ of 2.2.5(C) puts the $\Lambda$-modules}
\[
\bigl(R^n\psi(\Lambda)\bigr)_{\bar x}
\quad\text{and}\quad
H^n_{\{\bar x\}}\bigl(X_{\bar s},R\psi_{\bar\eta}(\Lambda)\bigr)
\]
\underline{in duality.  They are free of rank one if $n\ne0$, and
of rank two if $n=0$.}
\end{enumerate}
\footnote{The print says that $R^n\Phi(\Lambda)$ is ``alone''
outside $\Sigma$ (\emph{seul}); the proof and the asserted support on
$\Sigma$ require ``zero'' (\emph{nul}).  In (iii) it prints
$H^n_{\{\bar x\}}(R^n\psi_{\bar\eta}(\Lambda))$; the pairing cited
from 2.2.5(C) has the second term
$H^n_{\{\bar x\}}(X_{\bar s},R\psi_{\bar\eta}(\Lambda))$.}

The assertions of 3.1.2 are local for the \`etale topology.  We may
therefore treat only the case where $S$ is strictly henselian and
$X$ is the hypersurface in relative affine space
$\mathbb A^{n+1}_S$ defined by a quadratic equation satisfying
2.2.1(a) and (b) (apply 1.3.2).

The morphism $f$ is smooth at every point of $X_s-\Sigma'$, hence
$R\Phi=0$ outside $\Sigma'$.  By 2.2.4, $R\Phi$ even vanishes outside
$\Sigma$.  We may now assume 2.2.5(*); assertions (ii) and (iii) then
follow from 2.2.5(A), (B), and (C).
\end{filecontents*}
% Original input: components/076_expose_XV_section_3_2.tex
\begin{filecontents*}[overwrite,noheader]{sga7ii-cumulative-input-076.tex}
\subsection*{3.2. Even relative dimension}

We retain the notation of 3.1 and take $n=2m$ to be even.  For
simplicity we suppose that $S$ is strictly henselian (the general
case would follow by descent).

\subsubsection*{Proposition 3.2.1.}
\begin{enumerate}
\item[(i)] \underline{For $x\in\Sigma$, the group
$H^n_{\{x\}}(X_s,R\psi_{\bar\eta}(\Lambda(m)))$ has a natural
generator $\delta$, well defined up to sign, characterized by being
natural in $\Lambda$ and by satisfying}
\[
(\delta,\delta)=(-1)^m\,2.
\]
\underline{For $n=0$, add the condition
$\operatorname{Tr}(\delta)=0$.}

\item[(ii)] \underline{For a suitable character, independent of
$\Lambda$,}
\[
\varepsilon_x:I\longrightarrow\{\pm1\}
\]
\underline{of the inertia group $I$, the variation is}
\[
\operatorname{Var}(\sigma)(a)
=(-1)^m\left(\frac{\varepsilon_x(\sigma)-1}{2}\right)
(a,\delta)\delta
\]
\underline{for $a\in R^n\Phi_{\bar\eta}(\Lambda(m))$.  Consequently,}
\[
\sigma(\delta)=\varepsilon_x(\sigma)\delta.
\]
\end{enumerate}
\footnote{The print has
$H^n_{\{x\}}(R^n\psi_{\bar\eta}(\Lambda(m)))$ in (i).  The proof
explicitly invokes 2.2.5(D), where the natural generator belongs to
$H^n_{\{x\}}(X_s,R\psi_{\bar\eta}(\Lambda(m)))$.}

Passing to the inverse limit, it is enough to treat the universal
case $\Lambda=\mathbb Z_\ell$.  Up to sign, only one $\delta$ can
then satisfy the conditions in (i).  As in 3.1.2, we reduce to the
case 2.2.5(*) and apply 2.2.5(D).

\subsubsection*{Complement 3.2.2.}
Suppose that, at $x\in\Sigma$, the henselization $X_{(x)}$ of $X$ at
$x$ is isomorphic to the henselization at $x_0$ of the projective
scheme $Q$ with equation
\[
\sum_{i\leq j}a_{ij}x_ix_j=0.
\]
Then the character $\varepsilon_x$ is defined by the separable
quadratic extension $Z(C^+(Q))$ of the fraction field.  This follows
from 2.2.5(D).

\subsubsection*{3.2.3.}
In residual characteristic different from $2$, Complement 3.2.2 can
be made more explicit.  In that case $I$ has a unique nontrivial
character $\varepsilon$ of order $2$.  For a uniformizer $t$, one
has
\[
\sigma(\sqrt t)=\varepsilon(\sigma)\sqrt t.
\]
On the other hand, $X_{(x)}$ is then isomorphic to the henselization
at $0$ of an affine scheme
\[
\sum a_{ij}x_ix_j=b,
\]
where $b$ lies in the maximal ideal and
$Q(x)=\sum a_{ij}x_ix_j$ is nondegenerate.  The ideal $(b)$ defines
the image of the subscheme of $X_{(x)}$ which is the support of
$J^n_{X_{(x)}/S}$.  The center of the Clifford algebra is isomorphic
to
\[
k(\eta)\left(\left((-1)^{m+1}2b\det(a_{ij})\right)^{1/2}\right)
\]
(Bourbaki, \emph{Alg\`ebre}, Ch.~9, \S9, no.~4, final remark), whence
\[
\varepsilon_x=\varepsilon^{v(b)}.
\]
The variation is zero if $b$ has even valuation; otherwise
$\varepsilon_x=\varepsilon$.
\end{filecontents*}
% Original input: components/077_expose_XV_section_3_3_1_through_3_3_3.tex
\begin{filecontents*}[overwrite,noheader]{sga7ii-cumulative-input-077.tex}
\subsection*{3.3. Odd relative dimension}

To prove the Picard--Lefschetz formula in odd relative dimension, we
shall have to reduce to the transcendental case treated in XIV,
3.2.1.1.  The specialization argument required for this forces us to
work over a base which is no longer one-dimensional.  We shall use
ad hoc arguments adapted to the precise case we need: I do not have
a usable theory over a base of dimension greater than one, and I have
preferred not to cover my ignorance with generalities.

\subsubsection*{3.3.1.}
Let $S=\operatorname{Spec}(A)$ be a strictly local noetherian scheme,
and let $f:X\to S$ be a flat morphism of finite type, purely of
relative dimension $n$, whose special fiber $X_s$ has no nonsmooth
points other than nondegenerate quadratic points.  Let $s$ be the
closed point of $S$, let $X_s$ be the special fiber, and let
$x_0\in X_s$ be a singular point.

By 1.3.2(i), in a neighborhood of $x_0$ the scheme $X/S$ is locally
isomorphic to an $S$-scheme defined in $\mathbb A^{n+1}_S$ by an
equation $Q-b=0$, where $b$ belongs to the maximal ideal and $Q$ is a
nondegenerate quadratic form.  The ideal $(b)$ of $A$ depends only
on $(X/S,x_0)$: in a neighborhood of $x_0$ there is a section
\[
s\in X(A/(b))
\]
of $X$ over $\operatorname{Spec}(A/(b))$ such that
\[
J^n_{X/S}\simeq
\mathcal O_{s(\operatorname{Spec}(A/(b)))}.
\]
We denote this section by $s_x$, and denote by $b(x)$ any generator
of $(b)$.

\subsubsection*{3.3.2.}
Suppose that $S$ is normal, and let $b\in A$.  If $n$ is invertible
in $A$ and $b\ne0$, write $\varepsilon_b^{(n)}$ for the following
homomorphism from the Galois group of the fraction field of $A$ to
$\mathbb Z/(n)(1)$:
\[
\varepsilon_b^{(n)}(\sigma)
= (\sigma\sqrt[n]{b})/\sqrt[n]{b}.
\]
One has
\[
\bigl(\varepsilon_b^{(nm)}\bigr)^m=\varepsilon_b^{(n)}.
\]
If $\Lambda$ is a $\mathbb Z/n$-algebra, we again denote by
$\varepsilon_b$ the composite
\[
\operatorname{Gal}\longrightarrow\mathbb Z/(n)(1)
\longrightarrow\Lambda(1).
\]
It is independent of $n$.

The character $\varepsilon_b$ factors through
\[
\pi_1\bigl(S-(\text{the subscheme }b=0)\bigr).
\]

\subsubsection*{3.3.3.}
Suppose that $S$ is a strictly henselian trait and that $n$ is odd,
$n=2m+1$.  Let $x$ be a singular point of $X_s$ in a neighborhood
of which $X_\eta$ is smooth, and let $X_1/S$ be a subscheme of
$\mathbb A^{n+1}_S$ with equation $Q-b=0$, as in 1.3.2(i).  The
assumption on $x$ means that $b$, although zero modulo the maximal
ideal, is not zero.

Let $X_{s(x)}$ be the henselization of $X_s$ at $x$, let
$\widetilde X_{s(x)}$ be the blowup of $X_s$ at $x$, let $Y$ be the
quadric which is the inverse image of $x$, and put
\[
X^*_{s(x)}=X_{s(x)}-\{x\}=\widetilde X_{s(x)}-Y.
\]
The composite morphism
\begin{center}
\begin{tikzcd}[column sep=huge,row sep=huge]
H^{n-1}(Y,\Lambda(m))
  & H^{n-1}(\widetilde X_{s(x)},\Lambda(m))
      \arrow[l,"\sim"'] \arrow[r]
  & H^{n-1}(X^*_{s(x)},\Lambda(m)) \arrow[d] \\
{}
  & H^n_{\{x\}}\bigl(X_s,R\Psi_{\bar\eta}(\Lambda(m))\bigr)
  & H^n_{\{x\}}(X_s,\Lambda(m)) \arrow[l]
\end{tikzcd}
\end{center}
identifies the latter group with the primitive quotient (XII, 3.5)
of the cohomology of the quadric $Y$.  Indeed, this follows from
2.2.7 in a locally isomorphic situation.  The images of the two
natural generators of this quotient (the cohomology classes of the
rulings of $Y$) are the two basis vectors
\[
\pm\delta\in
H^n_{\{x\}}\bigl(X_s,R\Psi_{\bar\eta}(\Lambda(m))\bigr).
\]
\footnote{At the definition of the punctured henselization, the print
places a tilde on the first $X^*_{s(x)}$.  The equality with
$X_{s(x)}-\{x\}$ and the identical construction in 2.2.6 determine
the untilded term used here.}
\end{filecontents*}
% Original input: components/078_expose_XV_section_3_3_4_through_proposition_3_3_6.tex
\begin{filecontents*}[overwrite,noheader]{sga7ii-cumulative-input-078.tex}
\subsubsection*{3.3.4.}
Return to the more general assumptions of 3.3.1, with $n=2m+1$.
For every singular point $x$ of $X_s$, we again denote by
$\pm\delta_x$ the sections of $R^nf_!\Lambda(m)$ over $S$ which are
the images of the generators of the primitive quotient of
$H^{n-1}(Y,\Lambda(m))$ (with the same notation as above) under the
composite morphism
\begin{center}
\begin{tikzcd}[column sep=huge,row sep=huge]
H^{n-1}(Y,\Lambda(m))
  & H^{n-1}(\widetilde X_{s(x)},\Lambda(m))
      \arrow[l,"\sim"'] \arrow[r]
  & H^{n-1}(X^*_{s(x)},\Lambda(m)) \arrow[d] \\
H^0\bigl(S,R^nf_!\Lambda(m)\bigr)
  & H_c^n(X_s,\Lambda(m)) \arrow[l,"\sim"']
  & H^n_{\{x\}}(X_s,\Lambda(m)) \arrow[l]
\end{tikzcd}
\end{center}
The class $\pm\delta_x$ is the \underline{vanishing cycle} at $x$.

When $S$ is a trait, the class
\[
\pm\delta\in H_c^n(X_{\bar\eta},\Lambda(m))
\]
is the image under XIII, 2.1.8.4 of the class of 3.3.3
\[
\pm\delta\in
H^n_{\{x\}}\bigl(X_s,R\Psi_{\bar\eta}(\Lambda(m))\bigr).
\]

We shall prove the following two results simultaneously.

\subsubsection*{Proposition 3.3.5.}
\underline{Let $S$ be a henselian trait and let $X/S$ be as in
3.3.3.  The variation is given by}
\[
\operatorname{Var}(\sigma)(a)
=(-1)^{m+1}\varepsilon_{b(x)}(\sigma)(a,\delta)\delta.
\]

\subsubsection*{Proposition 3.3.6.}
\underline{Let $S$ be strictly local and regular, and let
$f:X\to S$ be a proper flat morphism of relative dimension
$n=2m+1$ such that $X_s$ has exactly one nonsmooth point $x$, assumed
to be a nondegenerate quadratic point.  Suppose also that $b(x)$ is
a parameter.  Then Galois acts on the cohomology of the geometric
generic fiber by}
\[
\sigma(a)=a+(-1)^{m+1}
\varepsilon_{b(x)}(\sigma)(a,\delta_x)\delta_x.
\]
\end{filecontents*}
% Original input: components/079_expose_XV_proofs_3_3_5_3_3_6_A_through_C.tex
\begin{filecontents*}[overwrite,noheader]{sga7ii-cumulative-input-079.tex}
\medskip
\noindent\textup{(A)}
Let $k$ be an algebraically closed field of characteristic $0$, and
let $S$ be the henselization of $\operatorname{Spec}(k[T])$ at $0$.
We prove 3.3.5 when $b(x)$ is a uniformizer.

By the Lefschetz principle, we reduce to the case $k=\mathbb C$.
Since the question is local, we may suppose that $f:X\to S$ is
obtained by base change from the map
\[
\sum_i z_i^2:\mathbb C^{n+1}\longrightarrow\mathbb C.
\]
The construction of $\pm\delta$ given above has a transcendental
analogue in integral cohomology.  Via XIV, 2.1, the class
$\pm\delta$ therefore corresponds to a transcendental vanishing
cycle, the image of
\[
\pm\delta_{\mathrm{cl}}
\in H^n_{\{0\}}\bigl(X_0,R\Psi_{\mathrm{cl}\eta}(\mathbb Z)\bigr).
\]
By XIV, 2.1 and the corresponding property of the algebraic cycles
$\pm\delta$, the image of $\pm\delta_{\mathrm{cl}}$ in
\[
H^n_{\{0\}}\bigl(X_0,R\Psi_{\mathrm{cl}\eta}(\mathbb Z)\bigr)
\otimes\mathbb Z_\ell
\]
is a generator of this group.  Thus $\pm\delta_{\mathrm{cl}}$ are the
two generators of
$H^n_{\{0\}}(X_0,R\Psi_{\mathrm{cl}\eta}(\mathbb Z))$, and coincide
with the vanishing cycles of XIV, 3.2.  Formula 3.3.5 consequently
follows from XIV, 3.2.11.  Strictly speaking, one should have checked
compatibility between the isomorphism of XIV, 2.1 and cup products
and traces.

\medskip
\noindent\textup{(B)}
We prove 3.3.6 for $X/S$ as in (A), with $X/S$ proper in addition.
This follows from XIII, 2.4.6.2.

\medskip
\noindent\textup{(C)}
Let $p$ be a prime number, and let $S_0$ be the strict henselization
of $\operatorname{Spec}(\mathbb Z[T])$ at $(p,T)$.  We prove 3.3.6
for $X/S_0$ such that $b(x)=T$.

The restriction of the sheaf $R^if_*\Lambda$ to the strict
henselization of $\operatorname{Spec}(\mathbb Z[T])$ at $(T)$ is
governed by (B).  This sheaf is
locally constant on the complement of the smooth divisor $T=0$,
which has mixed characteristic.  Assertion (C) therefore follows
from (B) and Abhyankar's lemma.
\end{filecontents*}
% Original input: components/080_expose_XV_proofs_3_3_5_3_3_6_D_through_F.tex
\begin{filecontents*}[overwrite,noheader]{sga7ii-cumulative-input-080.tex}
\par\medskip
\noindent\textup{(D)}
In at least one case to which (C) applies, $\delta_x$ is not a torsion
element in $\ell$-adic cohomology.

Let $S_1$ be the strict henselization of
$\operatorname{Spec}(\mathbb Z)$ at $(p)$, and let
$V\subset\mathbb P^r/S_1$ be a smooth projective variety over $S_1$,
purely of relative dimension $n+1$.  Suppose that there is a
Lefschetz pencil of hyperplane sections for the special fiber $V_p$
of $V/S_1$ (XVII, 2.5).  This pencil is defined by a line $d'$ in the
special fiber of the projective space $\check{\mathbb P}^{,r}/S_1$
dual to $\mathbb P^r/S_1$.  Lift $d'$ to a line $d$ in
$\check{\mathbb P}^{,r}/S_1$, and let $X\to d$ be the pencil of
hyperplane sections of $V$ parametrized by $d$ (XVIII, 3.1).  The
exceptional locus $\Delta$ of $d$ is \`etale over $S_1$ (XVII, 6).
After \`etale localization at a point of $\Delta_p\subset d$, we
therefore obtain a situation of type (C).

To check that $\delta_x$ is not torsion, it is enough to do so after
extending scalars from $S_1$ to $\mathbb C$.  Take
$V=\mathbb P^{n+1}$ with the projective embedding defined by
$\mathcal O(N)$.  For $N\ge2$, XVII, 2.5 applies (it is valid for
every $N$ in this particular case).  We then use the following lemma.

\subsubsection*{Lemma 3.3.7.}
\underline{For a Lefschetz pencil of hypersurfaces of sufficiently
large degree $N$ in $\mathbb P^{n+1}(\mathbb C)$ ($N\ge3$ suffices),
a vanishing cycle in $\mathbb Q_\ell$-cohomology is nonzero.}

This lemma follows from the theory of Lefschetz pencils, for example
from XVIII, 6.6.1.  There is no circularity in using here this result
of XVIII, based on the ``Picard--Lefschetz formula,'' since (A) is
already available.  One may instead use the more direct, purely
transcendental variant XIX, no.~4, which depends only on XIV, 3.2.

\medskip
\noindent\textup{(E)}
We prove 3.3.5 in residual characteristic $p$ for
$(X/S,x\in X_s)$.

Let $X_0/S_0$ be as in (D).  By 1.3.2(i), there is a morphism
$g:S\to S_0$ such that the henselization at $x_0$ of
\[
g^*X_0=X_0\times_{S_0}S
\]
is $S$-isomorphic to the henselization of $X$ at $x$.  It is
therefore enough to prove 3.3.5 for $g^*X_0$.  By base change from
(C) for $X_0/S_0$, one obtains
\[
\sigma(a)=a-(-1)^m
\varepsilon_{b(x)}(\sigma)(a,\delta_x)\delta_x
\]
in $H^n(X_{0\bar\eta})$, for $\sigma$ in the inertia group.  Let
$\pi$ be a uniformizer and compare this formula with (2.2.5.9), where
$\varepsilon=\varepsilon_\pi$.  We find
\[
\lambda_X\varepsilon_\pi(\sigma)(a,\delta_x)\delta_x
=-(-1)^m\varepsilon_{b(x)}(\sigma)(a,\delta_x)\delta_x
\]
in $H^n(X_{0\bar\eta})$.  Since $\delta_x$ is not torsion, we may
cancel $(a,\delta_x)\delta_x$ from this identity.  Formula (2.2.5.9),
with the resulting value of $\lambda_X$, is precisely 3.3.5.

\medskip
\noindent\textup{(F) Proof of 3.3.6 in the general case.}

The general case of 3.3.6 reduces to the general case of 3.3.5 just
as (B) and (C) reduce to (A).
\end{filecontents*}
% Original input: components/081_expose_XV_section_3_4_summary.tex
\begin{filecontents*}[overwrite,noheader]{sga7ii-cumulative-input-081.tex}
\subsection*{3.4. Summary}

Propositions 3.1.2, 3.2.1, 3.2.2, 3.2.3, and 3.3.5 constitute
the ``Picard--Lefschetz formula.''  This formula is the algebraic
analogue of XIV, 3.2.11, but takes very different forms according to
the parity of the relative dimension.

In the most interesting case, where $X/S$ is proper, the results may
be summarized as follows.

\subsubsection*{Theorem 3.4 (Picard--Lefschetz formula).}
\underline{Let $S$ be a strictly henselian trait and let
$f:X\to S$ be a proper flat morphism, purely of relative dimension
$n=2m$ or $2m+1$.  Suppose that $X_\eta$ is smooth and that $X_s$
has only ordinary quadratic singularities.  Let $\Sigma$ be the
singular locus of $X_s$.}

\begin{enumerate}
\item[(i)] \underline{To every $x\in\Sigma$ is attached a vanishing
cycle}
\[
\pm\delta_x\in H_n(X_{\bar\eta})(m).
\]
\underline{The $\delta_x$ are pairwise orthogonal.  When $n$ is even,}
\[
(\delta_x,\delta_x)=(-1)^m,2.
\]

\item[(ii)] \underline{There are isomorphisms}
\[
H^i(X_s)\xrightarrow{\ \sim\ }H^i(X_{\bar\eta})
\]
\underline{for $i\ne n,n+1$, and an exact sequence}
\[
0\longrightarrow H^n(X_s)
 \longrightarrow H^n(X_{\bar\eta})
 \xrightarrow{\,((\,\cdot\,,\delta_x))_{x\in\Sigma}\,}
 \bigoplus_{x\in\Sigma}\Lambda(n-m)
 \longrightarrow H^{n+1}(X_s)
 \longrightarrow H^{n+1}(X_{\bar\eta})
 \longrightarrow0.
\]

\item[(iii)] \underline{There are defined local characters of the
inertia group $I$:}
\[
\begin{array}{lll}
n\text{ even}:&\varepsilon_x:I\longrightarrow\{\pm1\},\\[2pt]
n\text{ odd}:&\varepsilon_x:I\longrightarrow\Lambda(1).
\end{array}
\]
\underline{When $X$ is regular at $x$, these characters are
surjective; when $x$ is nondegenerate, they are respectively the
unique character of order $2$ and the canonical character.}

\underline{The action of the inertia group $I$ is given by}
\[
n=2m\qquad
\sigma(a)=a+(-1)^m\sum_{x\in\Sigma}
 \frac{\varepsilon_x(\sigma)-1}{2}
 (a,\delta_x)\delta_x,
\]
\[
n=2m+1\qquad
\sigma(a)=a-(-1)^m\sum_{x\in\Sigma}
 \varepsilon_x(\sigma)(a,\delta_x)\delta_x.
\]
\underline{In particular,}
\[
\begin{array}{lll}
n\text{ even}:&\sigma(\delta_x)=\varepsilon_x(\sigma)\delta_x,\\[2pt]
n\text{ odd}:&\sigma(\delta_x)=\delta_x.
\end{array}
\]
\end{enumerate}
\end{filecontents*}
% Original input: components/082_expose_XVI_title_through_1_5.tex
\begin{filecontents*}[overwrite,noheader]{sga7ii-cumulative-input-082.tex}
\clearpage
\begin{center}
{\large SGA 7 --- EXPOS\'E XVI}\par
\vspace{2em}
{\Large\underline{The Milnor Formula}}\par
\vspace{1.5em}
{\large by P. Deligne}
\end{center}

\medskip
In this expos\'e, Milnor's formula [3] (see also [2]) is transposed
to characteristic $p$.  In certain cases this formula makes it
possible to compute the dimensions of spaces of vanishing cycles.
The proof is global in nature and treats only the ``geometric'' case
in which the base is a henselian trait of equal characteristic with
perfect residue field.

\section*{1. Statement of the problem}

\subsubsection*{1.1.}
Let $S$ be a regular scheme, purely of dimension one, let $s$ be a
closed point of $S$ \underline{with perfect residue field}, let
$f:X\to S$ be a flat morphism of finite type, let $x_0$ be a closed
point of the fiber $X_s$, and let $\bar x_0$ be a geometric point of
$X$ localized at $x_0$.  Suppose that
\begin{enumerate}
\item[(a)] $x_0$ is an \underline{isolated point at which $f$ is not
smooth}: there are an integer $n$ and a neighborhood $U$ of $x_0$,
purely of relative dimension $n$, such that $U-\{x_0\}$ is smooth
over $S$;
\item[(b)] $X$ is \underline{regular} at $x_0$.
\end{enumerate}

Under these assumptions, $X$ is a relative complete intersection
over $S$ at $x_0$.  We refer to VI, 3.10 for the definition of
$D^1_{X/S}$.  On the neighborhood $U$ one has
\[
D^1_{X/S}
=\mathcal H^1\bigl(R\mathcal Hom(L\Omega_{X/S},\mathcal O_X)\bigr)
=\mathcal Ext^1(\Omega^1_{X/S},\mathcal O_X);
\]
on this neighborhood, $D^1_{X/S}$ has finite length and support
contained in $\{x_0\}$.

\subsubsection*{Definition 1.2.}
\underline{The Milnor number
$\mu=\mu(X/S,\bar x_0)$ of $X/S$ at $\bar x_0$ is the length of the
$\mathcal O_{X,x_0}$-module defined by $D^1_{X/S}$.}

\subsubsection*{Example 1.3.}
Take $S$ to be a smooth curve over an algebraically closed field
$k$.  After replacing $X$ by a neighborhood of $x_0$, we may suppose
that $X$ is smooth over $k$.  The cotangent complex then reduces to
the following resolution of $\Omega^1_{X/S}$:
\[
L\Omega_{X/S}
=\bigl[f^*\Omega^1_{S/k}\longrightarrow\Omega^1_{X/k}\bigr].
\]
Choose a local coordinate $t:S\to\mathbb E^1_k$ on $S$ at $s$
(with $t$ \'etale at $s$), again denote by $f$ the composite
$t\circ f$, and let $(x_i)_{1\le i\le n+1}$ be a system of local
coordinates on $X$ at $x_0$.  At $x_0$, the module
$f^*\Omega^1_{S/k}$ then has basis $df$, while $\Omega^1_{X/k}$ has
basis the $dx_i$.  Thus $D^1_{X/S}$ is identified with the quotient
of $\mathcal O_X$ by the ideal generated by the
$\partial f/\partial x_i$, and
\begin{equation}\tag{1.3.1}
\mu=\dim_k\!\left(
\mathcal O_{X,x_0}/
\left(\frac{\partial f}{\partial x_1},\ldots,
      \frac{\partial f}{\partial x_{n+1}}\right)\right).
\end{equation}
This is the definition given by Milnor.

\subsubsection*{Example 1.4.}
Suppose that $X$ is the closed subscheme of $\mathbb E^{n+1}_S$
defined by an equation $g$, and let
$i:X\hookrightarrow\mathbb E^{n+1}_S$ be the inclusion.  The
cotangent complex then reduces to
\[
\bigl[(g)/(g)^2\longrightarrow
i^*\Omega^1_{\mathbb E^{n+1}/S}\bigr],
\]
and $D^1_{X/S}$ is isomorphic to the quotient of $\mathcal O_X$ by
the ideal generated by
$(\partial g/\partial x_i)_{1\le i\le n+1}$.  Equivalently, it is
the quotient of $\mathcal O_{\mathbb E^{n+1}}$ by the ideal generated
by $g$ and the $\partial g/\partial x_i$.

\subsubsection*{Example 1.5.}
Let $S$ be the spectrum of a discrete valuation ring $V$ with
perfect residue field and fraction field $K$, and let $X$ be the
normalization of $S$ in a separable extension $K'$ of $K$.  By 1.4
and [5], III, 7.14, p.~68, $D^1_{X/S}$ is then isomorphic to the
quotient of $\mathcal O_X$ by the different $\mathfrak d_{X/S}$.  If
$\bar x_0$ is a geometric point of $X_s$, localized at
$x_0\in X_s$, its Milnor number is therefore
\begin{equation}\tag{1.5.1}
\mu=\operatorname{length}
\bigl(\mathcal O_{X,x_0}/\mathfrak d_{X/S}\bigr).
\end{equation}
Let $\bar s$ be a geometric point of $S$ localized at the closed
point $s$, and let $d_{K'/K}\in V$ generate the discriminant ideal,
with valuation $v(d_{K'/K})$.  It follows from (1.5.1) that
\begin{equation}\tag{1.5.2}
v(d_{K'/K})=\sum_{x\in X(\bar s)}\mu(X/S,x).
\end{equation}
\end{filecontents*}
% Original input: components/083_expose_XVI_sections_1_6_through_1_10.tex
\begin{filecontents*}[overwrite,noheader]{sga7ii-cumulative-input-083.tex}
\subsubsection*{1.6.}
Retain the assumptions of 1.1.  Let $\bar s$ be the geometric point
of $S$ which is the image of $\bar x_0$, let $S_{(\bar s)}$ be the
spectrum of the strict henselization of $\mathcal O_{S,s}$ at
$\bar s$, and let $X_{(\bar s)}$ be the inverse image of $X$ over
$S_{(\bar s)}$.  Let $\bar\eta$ be the spectrum of an algebraic
closure of the fraction field $k(\eta)$ of $S_{(\bar s)}$, and put
\[
I=\operatorname{Gal}(k(\bar\eta)/k(\eta)).
\]

Let $\ell$ be a prime number different from the residual
characteristic of $S$.  In a neighborhood of $\bar x_0$ in
$X_{\bar s}$, the sheaves of vanishing cycles
\[
\Phi^i_{\bar\eta}(\mathbb Z/\ell)
=R^i\Phi_{\bar\eta}(\mathbb Z/\ell)
\]
relative to the constant sheaf $\mathbb Z/\ell$ on
$X_{(\bar s)}$ are concentrated at $\bar x_0$.  Moreover, the
inertia group $I$ acts on the $\mathbb Z/\ell$-vector spaces
\[
\Phi^i_{\bar\eta}(\mathbb Z/\ell)_{\bar x_0}.
\]
These spaces are finite-dimensional (XIII, 2.4.2).

\subsubsection*{1.7.}
Let $I_1$ be a finite quotient of the inertia group.  Recall that
the \underline{Swan character} $\operatorname{Sw}$ of $I_1$ is the
difference $a-u$ between the Artin character $a$ and the character
$u$ of the augmentation representation.  Recall also that, up to
isomorphism, there is a unique projective $\mathbb Z_\ell[I_1]$-module
$S(I_1)$ with character $\operatorname{Sw}$ (see [6]).

Let $V$ be a finite-dimensional $\mathbb Z/\ell$-vector space with
an action $\rho:I\to\operatorname{GL}(V)$, and suppose that $\rho$
factors through $I_1$.  The \underline{wild dimension} of $V$ is
then
\begin{equation}\tag{1.7.1}
\dim\operatorname{Sw}(V)
=\dim_{\mathbb Z/\ell}
\operatorname{Hom}_{I_1}\!\left(
S(I_1)\otimes_{\mathbb Z_\ell}\mathbb Z/\ell,,V\right).
\end{equation}
This integer is independent of the choice of $I_1$, since for a
normal subgroup $H$ of $I_1$ one has
$S(I_1/H)\sim S(I_1)_H$.

We further put
\[
\dim\operatorname{tot}(V)
=\dim\operatorname{Sw}(V)+\dim_{\mathbb Z/\ell}(V)
\]
and call this number the \underline{total dimension} of $V$.

\subsubsection*{Definition 1.8.}
\underline{With the terminology of 1.6 and 1.7, the number of
vanishing cycles of $X/S$ at $\bar x_0$ (respectively the wild
number of vanishing cycles, respectively the total number of
vanishing cycles) is the integer}
\[
(-1)^n\dim\Phi_{\bar x_0}
\underset{\mathrm{def}}{=}
(-1)^n\sum_i(-1)^i
\dim\Phi^i_{\bar\eta}(\mathbb Z/\ell)_{\bar x_0},
\]
\[
(-1)^n\dim\operatorname{Sw}\Phi_{\bar x_0}
\underset{\mathrm{def}}{=}
(-1)^n\sum_i(-1)^i
\dim\operatorname{Sw}
\Phi^i_{\bar\eta}(\mathbb Z/\ell)_{\bar x_0},
\]
\[
\begin{aligned}
(-1)^n\dim\operatorname{tot}\Phi_{\bar x_0}
&\underset{\mathrm{def}}{=}
(-1)^n\sum_i(-1)^i\\[-2pt]
&\qquad{}\cdot\dim\operatorname{tot}
\Phi^i_{\bar\eta}(\mathbb Z/\ell)_{\bar x_0}\\
&=(-1)^n\bigl(
\dim\Phi_{\bar x_0}+\dim\operatorname{Sw}\Phi_{\bar x_0}
\bigr).
\end{aligned}
\]

\subsubsection*{Conjecture 1.9.}
\underline{Under the assumptions of 1.1, the Milnor number $\mu$ of
$X/S$ at $\bar x_0$ is equal to the total number of vanishing cycles
of $X/S$ at $\bar x_0$.}

\subsubsection*{Remark 1.10.}
Let $S$ be a henselian trait with imperfect residue field and let
$f:X\to S$ be a morphism which is smooth except at
$x_0\in X_s$.  Suppose that $X$ is regular at $x_0$.  Definition
1.2 of the Milnor number $\mu$ still makes sense under these
assumptions, but I no longer have a geometric interpretation of
$\mu$ to propose.
\end{filecontents*}
% Original input: components/084_expose_XVI_sections_1_11_through_1_12.tex
\begin{filecontents*}[overwrite,noheader]{sga7ii-cumulative-input-084.tex}
\subsubsection*{1.11.}
For the remainder of this section, suppose that
$S=\operatorname{Spec}(V)$ is a henselian trait with
\underline{algebraically closed} residue field and closed point $s$,
that $f:X\to S$ is a flat morphism of finite type, purely of relative
dimension $n$, that $X$ is regular, that $x_0$ is a closed point of
$X_s$, and that $X-\{x_0\}$ is smooth over $S$.

This entails no loss of generality.

\subsubsection*{Proposition 1.12.}
\underline{If $n=0$ (that is, if $X/S$ is finite), Conjecture 1.9
holds for $(X/S,x_0)$.}

We may suppose that $X$ is the normalization of $S$ in a finite
extension $K'$ of the fraction field $K$.  By 1.5,
$\mu=v(d_{K'/K})$.

The extension $K'$ is defined by an open subgroup $H$ of the inertia
group $I$.  Let $H'\subset H$ be an open normal subgroup of $I$, and
put $I'=I/H'$.  Let $s_{I/H}$ be the character of the permutation
representation of $I'$ on $I/H$, let $a_{I'}$ be the Artin
character, let $\operatorname{Sw}_{I'}$ be the Swan character, and
let $u_{I'}$ be the character of the augmentation representation.
By Artin's ``F\"uhrerdiskriminantenproduktformel,'' the discriminant
ideal is the Artin conductor of $s_{I/H}$ ([5], VI, 3, Corollary 1,
p.~111):
\[
\mu=(a_{I'},s_{I/H})
=|I'|^{-1}\sum_{i\in I'}a_{I'}(i)s_{I/H}(i).
\]
One has
\[
a_{I'}=u_{I'}+\operatorname{Sw}_{I'},
\qquad
(u_{I'},s_{I/H})=\#(I/H)-1,
\]
and
\begin{align*}
(\operatorname{Sw}_{I'},s_{I/H})
&=\dim_{\mathbb Z/\ell}\operatorname{Hom}_I\!\left(
S(I')\otimes_{\mathbb Z_\ell}\mathbb Z/\ell,
(\mathbb Z/\ell)^{I/H}\right)\\
&=\dim_{\mathbb Z/\ell}\operatorname{Hom}_I\!\left(
S(I')\otimes_{\mathbb Z_\ell}\mathbb Z/\ell,
\operatorname{Coker}\bigl(
\mathbb Z/\ell\longrightarrow(\mathbb Z/\ell)^{I/H}\bigr)
\right).
\end{align*}
Consequently,
\[
\mu=\#(I/H)-1+
\dim\operatorname{Sw}\operatorname{Coker}\bigl(
\mathbb Z/\ell\longrightarrow(\mathbb Z/\ell)^{I/H}\bigr).
\]
One checks immediately that $\Phi^i_{x_0}=0$ for $i\ne0$ and that
\[
\Phi^0_{x_0}=\operatorname{Coker}\bigl(
\mathbb Z/\ell\longrightarrow(\mathbb Z/\ell)^{I/H}\bigr);
\]
the formula follows.
\end{filecontents*}
% Original input: components/085_expose_XVI_section_1_13.tex
\begin{filecontents*}[overwrite,noheader]{sga7ii-cumulative-input-085.tex}
\subsubsection*{Proposition 1.13.}
\underline{If $X_s$ has an ordinary quadratic singularity at $x_0$,
Conjecture 1.9 holds for $(X/S,x_0)$.}

\medskip
\noindent\underline{Case 1: $n=\dim(X_s)=2m-1$.}

The problem is local for the \'etale topology on $X$.  We may
therefore suppose that $X$ is the closed subscheme of
$\mathbb E_S^{n+1}$ defined by an equation
\[
f=\sum_{i=1}^{m}x_i x_{i+m}+a=0,
\]
where $a$ lies in the maximal ideal of $V$ and $x_0$ is the origin
of the coordinates.  Moreover, $X$ is regular at $x_0$ if and only
if $a$ is a uniformizer.  By 1.4,
\[
\mu=\operatorname{length}\!\left(
\mathcal O_{\mathbb E_S^{2m},x_0}/
\left(f,\frac{\partial f}{\partial x_i}\right)\right)=1.
\]
By the ordinary-quadratic vanishing-cycle calculation of Expos\'e XV,
$\Phi^i_{x_0}=0$ for $i\ne n$, while $\Phi^n_{x_0}$ has dimension
one and inertia acts trivially on it.\footnote{After the printed word
``By,'' the source leaves the expected reference blank.  The stated
concentration and inertia action are the ordinary-quadratic
vanishing-cycle calculation established in Expos\'e XV.}
The wild number of vanishing cycles is therefore zero and the total
number is $1$, which verifies 1.9.

\medskip
\noindent\underline{Case 2: $n=\dim(X_s)=2m$ and
$\operatorname{char}(k(s))\ne2$.}

Under these assumptions we may take $X\subset\mathbb E_S^{n+1}$ to
be defined by
\[
f=\sum_{i=1}^{m}x_i x_{i+m}+x_{2m+1}^2+a=0,
\]
where $a$ is a uniformizer.  Again $\mu=1$,
$\Phi^i_{x_0}=0$ for $i\ne n$, and $\Phi^n_{x_0}$ has dimension
one.  This time the inertia group acts by $\pm1$ on
$\Phi^n_{x_0}$, hence tamely.  The wild number of vanishing cycles is
again zero, and 1.9 is verified.

\medskip
\noindent\underline{Case 3: $n=\dim(X_s)=2m$ and
$\operatorname{char}(k(s))=2$.}

This time we may take $X\subset\mathbb E_S^{n+1}$ to be defined by
\[
f=\sum_{i=1}^{m}x_i x_{i+m}
 +(x_{2m+1}^2+a x_{2m+1}+b)=0,
\]
where $a$ and $b$ lie in the maximal ideal.  The scheme $X$ is
regular if and only if $b$ is a uniformizer, and its generic fiber
is smooth if and only if $a^2-4b\ne0$.

The Milnor number of $X/S$ is the same as that of
$X_0\subset\mathbb E_S^1$ defined by
\begin{equation}\tag{1}
x^2+a x+b=0.
\end{equation}
Likewise, $\Phi^i_{x_0}=0$ for $i\ne n$,
$\Phi^n_{x_0}$ has dimension one, and $I$ acts on it by $\pm1$; the
index-two subgroup of $I$ acting trivially corresponds to the
extension of $K$ defined by (1).  The total number of vanishing
cycles of $X/S$ at $x_0$ therefore coincides with the one for
$X_0/S$, and 1.9 in this case follows from 1.12.  The Milnor number
is given by the discriminant formula of 1.5:
\begin{equation}\tag{1.13.1}
\mu=v(a^2-4b).
\end{equation}
\end{filecontents*}
% Original input: components/086_expose_XVI_section_2_1.tex
\begin{filecontents*}[overwrite,noheader]{sga7ii-cumulative-input-086.tex}
\section*{2. The ``geometric'' case}

We shall prove that Conjecture 1.9 holds in equal characteristic.
The key geometric result is the following.  Its possible use was
suggested to me by conversations with B. Saint-Donat.

\subsubsection*{Proposition 2.1.}
\underline{Let $S$ be a smooth connected projective curve over an
algebraically closed field $k$, let $X$ be a smooth proper variety
over $k$, purely of dimension $n+1$, and let $f:X\to S$ be a
morphism which is smooth away from a finite set $F$ of points of
$X$.  Then the Euler--Poincar\'e characteristics in \'etale
cohomology of $S$, $X$, and the geometric generic fiber
$X_{\bar\eta}$ of $f$ satisfy}
\begin{equation}\tag{2.1.1}
\chi(X)=\chi(S)\chi(X_{\bar\eta})
-(-1)^n\sum_{x\in F}\mu(X/S,x).
\end{equation}

The Euler--Poincar\'e characteristic of $X$ is the trace (or integral
over $X$) of the top Chern class of the tangent bundle:
\[
\chi(X)=\int_Xc_{n+1}(\Theta_X)
=(-1)^{n+1}\int_Xc_{n+1}(\Omega_X^1).
\]
Write
\[
\Omega_X^1=\left(\Omega_X^1\otimes
(f^*\Omega_S^1)^{\otimes-1}\right)\otimes f^*\Omega_S^1.
\]
The total Chern class of the invertible sheaf $f^*\Omega_S^1$ is
the inverse image of the total Chern class of $\Omega_S^1$ on $S$:
\[
c^*(f^*\Omega_S^1)=1+f^*c^1(S,\Omega_S^1),
\]
and the cup powers of $c^1$ vanish:
\[
\bigl(f^*c^1(S,\Omega_S^1)\bigr)^n
=f^*\bigl(c^1(S,\Omega_S^1)^n\bigr)=0
\qquad(n>1).
\]
It follows that
\begin{align*}
c^{n+1}(\Omega_X^1)
={}&c^n\!\left(\Omega_X^1\otimes
(f^*\Omega_S^1)^{\otimes-1}\right)
\smile f^*c_1(S,\Omega_S^1)\\
&+c^{n+1}\!\left(\Omega_X^1\otimes
(f^*\Omega_S^1)^{\otimes-1}\right).
\end{align*}

If $S$ has genus $g$, then $c_1(S,\Omega_S^1)$ is $(2g-2)$
times the cohomology class of any closed point $s$ of $S$.
Consequently, when $X_s$ is nonsingular,
\begin{align*}
\mathrm I
&=\int_Xc^n\!\left(\Omega_X^1\otimes
(f^*\Omega_S^1)^{\otimes-1}\right)
\smile f^*c_1(S,\Omega_S^1)\\
&=(2g-2)\int_Xc^n\!\left(\Omega_X^1\otimes
(f^*\Omega_S^1)^{\otimes-1}\right)
\smile f^*\operatorname{cl}(s)\\
&=(2g-2)\int_{X_s}c^n\!\left(
X_s,\Omega_X^1\otimes(f^*\Omega_S^1)^{\otimes-1}
\otimes\mathcal O_{X_s}\right).
\end{align*}
Moreover, on $X_s$ the sheaf $f^*\Omega_S^1$ is isomorphic to
$\mathcal O_{X_s}$, and
$\Omega_X^1\otimes_{\mathcal O_X}\mathcal O_{X_s}$ occurs in the
exact sequence
\[
0\longrightarrow f^*\Omega_S^1\otimes\mathcal O_{X_s}
\longrightarrow\Omega_X^1\otimes\mathcal O_{X_s}
\longrightarrow\Omega_{X_s}^1\longrightarrow0.
\]
Hence
\[
c^n\!\left(X_s,\Omega_X^1\otimes
(f^*\Omega_S^1)^{\otimes-1}\otimes\mathcal O_{X_s}\right)
=c^n(X_s,\Omega_{X_s}^1),
\]
and
\[
\mathrm I=(2g-2)\int_{X_s}c^n(X_s,\Omega_{X_s}^1)
=(-1)^{n+1}\chi(S)\chi(X_s).
\]

The canonical map
\[
f^*:f^*\Omega_S^1\longrightarrow\Omega_X^1
\]
is identified with a section $a$ of
$\Omega_X^1\otimes(f^*\Omega_S^1)^{\otimes-1}$.  After choosing
systems of local coordinates on $S$ and $X$, thereby identifying
this bundle with $\mathcal O_X^{n+1}$, the coordinates of $a$ are
the partial derivatives $\partial f/\partial x_i$.  Thus $a$
vanishes only on $F$, and by 1.3 the multiplicity of its zero at
$x\in F$ is the Milnor number $\mu(X/S,x)$.

Recall that if a section $a$ of a vector bundle $G$ of rank $d$ on
a smooth proper $k$-scheme $X$, purely of dimension $d$, vanishes
only at isolated points, then $\int_Xc^d(G)$ is the number of zeros
of $a$, counted with multiplicity.  Here
\[
\mathrm{II}=\int_Xc^{n+1}\!\left(
\Omega_X^1\otimes(f^*\Omega_S^1)^{\otimes-1}\right)
=\sum_{x\in F}\mu(X/S,x).
\]
Collecting terms gives
\[
\chi(X)=(-1)^{n+1}(\mathrm I+\mathrm{II})
=\chi(S)\chi(X_s)-(-1)^n\sum_{x\in F}\mu(X/S,x),
\]
which proves 2.1, since
$\chi(X_s)=\chi(X_{\bar\eta})$.
\end{filecontents*}
% Original input: components/087_expose_XVI_section_2_2.tex
\begin{filecontents*}[overwrite,noheader]{sga7ii-cumulative-input-087.tex}
\subsubsection*{Corollary 2.2.}
\underline{Under the preceding assumptions, if $n(X/S,x)$ is the
total number of vanishing cycles of $X/S$ at $x$ (1.8), then}
\[
\sum_{x\in F}\mu(X/S,x)=\sum_{x\in F}n(X/S,x).
\]

The Leray spectral sequence for $f$ gives
\begin{align*}
\chi(X)
&=\sum_{p,q}(-1)^{p+q}
\dim_{\mathbb Z/\ell}H^p(S,R^qf_*\mathbb Z/\ell)\\
&=\sum_q(-1)^q\chi(S,R^qf_*\mathbb Z/\ell).
\end{align*}
For every constructible sheaf $G$ of $\mathbb Z/\ell$-modules on
$S$, one has ([4] or SGA 5, X, 7.2)
\[
\chi(S,G)=\chi(S)\dim_{\mathbb Z/\ell}(G_{\bar\eta})
-\sum_{s\in S}\varepsilon_s(G),
\]
where, for each $s\in S$, after identifying $\bar\eta$ with the
algebraic closure of the fraction field of the henselization of $S$
at $s$,
\[
\varepsilon_s(G)=
\dim(G_{\bar\eta})-\dim(G_s)+\dim\operatorname{Sw}(G_{\bar\eta}).
\]
\footnote{The print has a minus sign before the Swan term.  The plus
sign is forced both by the standard Artin-conductor formula cited in
the preceding sentence and by the printed identity (c), which
identifies the alternating sum with the \emph{total} dimension
$\dim+\dim\operatorname{Sw}$.}

For $s\in S$, the exact sequences XIII, 2.4.6.3,
\[
\cdots\longrightarrow(R^qf_*\mathbb Z/\ell)_s
\longrightarrow(R^qf_*\mathbb Z/\ell)_{\bar\eta}
\longrightarrow\bigoplus_{x\in F(s)}
\Phi^q_{\bar\eta}(\mathbb Z/\ell)_x
\longrightarrow\cdots
\]
give
\begin{align*}
&\sum_q(-1)^q\left(
\dim(R^qf_*\mathbb Z/\ell)_{\bar\eta}
-\dim(R^qf_*\mathbb Z/\ell)_s\right)\\
&\hspace{4em}=\sum_{x\in F(s)}\sum_q(-1)^q
\dim\Phi^q_{\bar\eta}(\mathbb Z/\ell)_x,
\end{align*}
and
\[
\sum_q(-1)^q\dim\operatorname{Sw}
(R^qf_*\mathbb Z/\ell)_{\bar\eta}
=\sum_{x\in F(s)}\sum_q(-1)^q
\dim\operatorname{Sw}\Phi^q_{\bar\eta}(\mathbb Z/\ell)_x.
\]
Consequently,
\begin{equation}\tag{c}
\sum_q(-1)^q\varepsilon_s(R^qf_*\mathbb Z/\ell)
=\sum_{x\in F(s)}\dim\operatorname{tot}\Phi_x.
\end{equation}
The Euler--Poincar\'e formula therefore becomes
\begin{align*}
\chi(X)
&=\sum_q(-1)^q\chi(S,R^qf_*\mathbb Z/\ell)\\
&=\sum_q(-1)^q\left(
\chi(S)\dim(R^qf_*\mathbb Z/\ell)_{\bar\eta}
-\sum_s\varepsilon_s(R^qf_*\mathbb Z/\ell)\right)\\
&=\chi(S)\sum_q(-1)^q
\dim(R^qf_*\mathbb Z/\ell)_{\bar\eta}
-\sum_s\sum_q(-1)^q\varepsilon_s(R^qf_*\mathbb Z/\ell)\\
&=\chi(S)\chi(X_{\bar\eta})
-\sum_s\sum_{x\in F(s)}\dim\operatorname{tot}\Phi_x\\
&=\chi(S)\chi(X_{\bar\eta})
-(-1)^n\sum_{x\in F}n(X/S,x).
\end{align*}
Comparison with (2.1.1) proves 2.2.
\end{filecontents*}
% Original input: components/088_expose_XVI_sections_2_3_through_proposition_2_5.tex
\begin{filecontents*}[overwrite,noheader]{sga7ii-cumulative-input-088.tex}
\subsubsection*{Corollary 2.3.}
\underline{Under the preceding assumptions and with the preceding
notation, suppose that $x_0\in F$ and that, for every
$x\in F$ distinct from $x_0$, the fiber of $f$ through $x$ has an
ordinary quadratic singularity at $x$.  Then}
\[
\mu(X/S,x_0)=n(X/S,x_0).
\]
For $x\ne x_0$, Proposition 1.13 gives
\[
\mu(X/S,x)=n(X/S,x),
\]
and it remains to apply 2.2.

\subsubsection*{Theorem 2.4.}
\underline{Let $S$ be a henselian trait of equal characteristic with
perfect residue field, let $f:X\to S$ be a flat morphism of finite
type, and let $\bar x$ be a geometric point of $X$ whose image in
$X$ is a closed point $x$ of $X_s$ at which $X/S$ is not smooth,
and which is isolated as such.  Suppose that $X$ is regular at $x$.
Then the total number (1.8) of vanishing cycles of $X/S$ at $\bar x$
coincides with the Milnor number $\mu(X/S,x)$.}

Let $S'$ be the henselian trait obtained from $S$ by extending the
residue field $k(s)$ to its separable closure, and put
$X'=X\times_S S'$.  The assertion of 2.4 for $X/S$ is equivalent to
the assertion for $X'/S'$.  Since $k(s)$ is already assumed perfect,
we may therefore suppose that the residue field of $S$ is
algebraically closed.

For such an $S$, let $S'$ be its completion and put
$X'=X\times_S S'$.  The Milnor number at $x$ is the same for $X/S$
and $X'/S'$, and by XIII, 2.4.2 the same is true of the total number
of vanishing cycles.  We may thus suppose in addition that $S$ is
complete, hence isomorphic to $k[[t]]$, with $k$ algebraically
closed.  If $X_s$ has dimension $n$ at $x$, we then have:

\subsubsection*{Proposition 2.5.}
\underline{There exist a smooth projective variety $Y$ over $k$,
purely of dimension $n+1$, a morphism
$f:Y\to\mathbb P^1_k$ which is smooth outside a finite set $F$ of
points of $Y$, and a point $y_0\in F$ such that:}
\begin{enumerate}
\item[(a)] \underline{if $y\in F$ and $y\ne y_0$, then $y$ is an
ordinary quadratic singular point of the fiber of $f$ through $y$;}
\item[(b)] \underline{$f(y_0)=0$, and if
$a:S\to\mathbb P^1_k$ is the natural identification of $S$ with the
completion of $\mathbb P^1_k$ at $0$, there is an $S$-isomorphism
between the henselization of $X$ at $x$ and the henselization of
$a^*Y$ at $y_0$.}
\end{enumerate}

Let us show that 2.5 implies 2.4.  By 2.3 and 2.5(a), the Milnor
number of $Y/\mathbb P^1_k$ at $y_0$ coincides with the total number
of vanishing cycles.\footnote{The print says $Y/T$ here, but no base
$T$ is introduced; Proposition 2.5 defines the morphism
$Y\to\mathbb P^1_k$.}
As above, XIII, 2.4.2 gives the same result for $a^*Y/S$, hence also
for $(X/S,x)$, which by 2.5(b) is locally isomorphic to
$(a^*Y/S,y_0)$ for the \'etale topology.
\end{filecontents*}
% Original input: components/089_expose_XVI_proof_2_5_and_bibliography.tex
\begin{filecontents*}[overwrite,noheader]{sga7ii-cumulative-input-089.tex}
We prove 2.5.  As a $k[[t]]$-algebra, the completion of the local
ring $\mathcal O_{X,x}$ is isomorphic to a quotient
\[
\widehat{\mathcal O}_{X,x}
=k[[t,x_1,\ldots,x_{n+1}]]/(f).
\]
The implicit-function theorem XV, 1.1.2 (that is, [1], Lemma 5.10)
then shows that there is an integer $N$ such that condition 2.5(b)
holds whenever
\[
\widehat{\mathcal O}_{Y,y_0}
\simeq k[[t,x_1,\ldots,x_{n+1}]]/(f_1)
\]
as $k[[t]]$-algebras, with $f_1-f$ in the $N$th power $\mathfrak m^N$
of the maximal ideal of $k[[t,x_1,\ldots,x_{n+1}]]$.

Let $(t_0,t_1)$ and $(x_0,\ldots,x_{n+1})$ be the usual homogeneous
coordinate systems on $\mathbb P^1$ and $\mathbb P^{n+1}$.  We take
$Y/\mathbb P^1$ to be the subvariety of
$\mathbb P^1\times\mathbb P^{n+1}$ defined by an equation
\begin{equation}\tag{2.5.1}
g(t_0,t_1;x_0,\ldots,x_{n+1})=0,
\end{equation}
homogeneous of degree $d'$ in $(t_0,t_1)$ and of degree $d''$ in the
$x_i$.  More intrinsically, $g$ is a section of
\[
\mathcal O(d',d'')=\operatorname{pr}_1^*\mathcal O(d')
\otimes\operatorname{pr}_2^*\mathcal O(d'').
\]
Condition (b) holds if
\begin{equation}\tag{2.5.2}
g(1,t;1,x_1,\ldots,x_{n+1})\equiv f\pmod{\mathfrak m^N}.
\end{equation}

Fix $d'$ and $d''$.  The sections $g$ of type (2.5.1) are the
$k$-rational points of an affine space $A$.  For
$u\in\mathbb P^1\times\mathbb P^{n+1}$, let $B(u,r)$ denote the
affine space whose points are the $r$-jets at $u$ of sections of
$\mathcal O(d',d'')$.  Put
$u_0=(1,0;1,0,\ldots,0)$.  If $u\ne u_0$, then for sufficiently
large $d'$ and $d''$ the always-linear restriction map
\begin{equation}\tag{2.5.3}
A\longrightarrow B(u_0,N)\times B(u,2)
\end{equation}
is surjective.

Let $C(u)\subset B(u,2)$ be the set of $2$-jets which are the
$2$-jets either of the zero germ or of a germ $\widetilde h$ such
that
\begin{enumerate}
\item[(a)] $u$ lies on the variety $V$ defined by $\widetilde h=0$;
\item[(b)] $u$ is neither a smooth point nor an ordinary quadratic
point of the fiber through $u$ of $V/\mathbb P^1$.
\end{enumerate}
It is readily checked that $C(u)$ is closed of codimension $n+3$ in
$B(u,2)$.

Let $D$ be the affine subspace of $A$ consisting of the $g$ which
satisfy (2.5.2).  For $u\ne u_0$, let $D(u)$ be the subspace of $D$
consisting of the zero $g$ and of those $g$ such that $g(u)=0$ and
$u$ is neither a smooth point nor an ordinary quadratic point of the
fiber through $u$ of the variety $V$ defined by $g=0$.  If $d'$ and
$d''$ are sufficiently large for (2.5.3) to be surjective, then
$D(u)$ has codimension $n+3$ in $D$.  The union of the $D(u)$ for
$u\ne u_0$ therefore has codimension at least
\[
1=(n+3)-(n+2),
\]
and a general element $g$ of $D$ lies in none of the $D(u)$.  It
solves the problem.

\subsection*{Bibliography}

\begin{enumerate}[label={[\arabic*]},leftmargin=3em]
\item M. Artin, ``Algebraic approximation of structures over complete
local rings,'' \emph{Publ. Math. IHES} \underline{36} (1969)
(Zariski's volume).
\item L\^e D\~ung Tr\'ang, \emph{Singularit\'es isol\'ees des
hypersurfaces complexes}, Publ. Centre Math. \'Ecole Polytechnique,
May 1969.
\item J. Milnor, \emph{Singular Points of Complex Hypersurfaces},
Princeton University Press, 1968.
\item M. Raynaud, ``Caract\'eristique d'Euler--Poincar\'e d'un
faisceau et cohomologie des vari\'et\'es ab\'eliennes,'' S\'eminaire
Bourbaki 286, February 1965.
\item J.-P. Serre, \emph{Corps locaux}, Publ. Inst. Math. Univ.
Nancago VIII, Hermann, 1962.
\item J.-P. Serre, ``SGA 5, IX. Introduction \`a la th\'eorie de
Brauer,'' reproduced in \emph{Repr\'esentations lin\'eaires des
groupes}, 3rd part, Hermann, Paris, 1967.
\end{enumerate}
\end{filecontents*}
% Original input: components/090_expose_XVII_title_introduction_and_section_1.tex
\begin{filecontents*}[overwrite,noheader]{sga7ii-cumulative-input-090.tex}
\clearpage
\begin{center}
{\large SGA 7 --- EXPOS\'E XVII}\par
\vspace{2em}
{\Large\underline{Lefschetz Pencils: Existence Theorem}}\par
\vspace{1.5em}
{\large by N. M. Katz}
\end{center}

\section*{0. Introduction}

In the present expos\'e (which, in the oral seminar, came before
Expos\'es XIV and XV devoted to the Picard--Lefschetz formula), we
introduce the ``Lefschetz pencils'' of hyperplane sections of a smooth
projective scheme and give conditions for their existence.  In the
following expos\'e we make an elementary cohomological study of them,
based on the known structure of the cohomology of blown-up varieties.
The fact that certain hypotheses on Lefschetz pencils which we shall
introduce are often satisfied will follow from the Picard--Lefschetz
formula of Expos\'e XV, which is less elementary in nature than the
results of the present expos\'e and the next.

Some developments of the present expos\'e and the next will be given
elsewhere in a more general setting (general base schemes,
quasi-projectivity hypotheses in place of projectivity, and so forth):
in particular, the precohomological study of pencils of hyperplane
sections, and especially of Lefschetz pencils, in EGA V, and the
cohomological study of blown-up schemes in Jouanolou's Expos\'e VII of
SGA 5.  The reader will also find in EGA V proofs in which intrinsic
diagram chasing replaces the use of coordinates favored by the writer
of the present expos\'e.

I wish to thank P. Deligne and P. Griffiths for the help they gave me
in conceiving this expos\'e.  Throughout it, $k$ denotes an
algebraically closed field unless expressly stated otherwise.

\section*{1. A reminder on quadratic singularities (cf. VI, 6)}

\subsubsection*{1.1.}
Let $Y$ be a $k$-scheme pure of dimension $n-1$.
\footnote{The print first says ``of dimension $n$'' and then, in the
hypersurface presentation below, imposes $n=1+\dim Y$.  The latter is
the convention used by the displayed Hessian and normal form.}
A closed point $y\in Y$ is called an \underline{ordinary quadratic
singularity} of $Y$ if:
\begin{enumerate}
\item[(a)] the completion $\widehat{\mathcal O}_{Y,y}$ of the local
ring of $Y$ at $y$ is isomorphic to
\[
k[[T_1,\ldots,T_n]]/(Q(T)),
\]
the quotient of the formal power-series ring in
$n=1+\dim Y$ variables by the ideal generated by an element $Q(T)$;
\item[(b)] $Q$ begins in degree two, and the subvariety of
$\mathbb P^{n-1}_k$ defined by the vanishing of the homogeneous
degree-two part of $Q(T)$ is smooth over $k$ when $n\geq2$, while the
degree-two part of $Q(T)$ is nonzero when $n=1$.
\end{enumerate}

If the characteristic of $k$ is different from two, \underline{or} if
$n$ is even, condition (b) is equivalent to:
\begin{enumerate}
\item[(b$'$)] $Q$ begins in degree two and its Hessian matrix
\[
\left(
  \frac{\partial^2Q}{\partial T_i\partial T_j}(0)
\right)_{1\leq i,j\leq n}
\]
is invertible.
\end{enumerate}
This follows from the Jacobian criterion for smoothness.

The standard form for the homogeneous degree-two part $Q_2$ of a
series $Q$ satisfying (b), obtainable after a linear change of
variables, is
\begin{equation}\tag{1.1.1}
\left\{
\begin{aligned}
Q_2(T)&=\sum_{i=1}^{n/2}T_{2i-1}T_{2i}
  &&\text{if $n$ is even},\\[1ex]
Q_2(T)&=T_n^2+\sum_{i=1}^{(n-1)/2}T_{2i-1}T_{2i}
  &&\text{if $n$ is odd}.
\end{aligned}
\right.
\end{equation}
This follows from the theory of quadratic forms (Bourbaki,
\emph{Alg\`ebre}, Ch.~X).

\subsubsection*{1.2.}
A closed point $y\in Y$ satisfying (a) and (b) is called a
\underline{nondegenerate quadratic singularity}.  Since the
determinant of the Hessian matrix is identically zero in
characteristic two when $n$ is odd, nondegenerate quadratic
singularities presented in $n$ variables exist only when $n$ is even
or the characteristic of $k$ is different from two; in those cases
they are precisely the ordinary quadratic singularities.
\end{filecontents*}
% Original input: components/091_expose_XVII_section_2.tex
\begin{filecontents*}[overwrite,noheader]{sga7ii-cumulative-input-091.tex}
\section*{2. Lefschetz pencils: statement of the results}

\subsubsection*{2.1.}
Let $\mathbb P^r$ denote projective $r$-space over $k$,
$\check{\mathbb P}^{r}$ the \underline{dual} projective space of
hyperplanes in $\mathbb P^r$, and
$\operatorname{Gr}(1,\check{\mathbb P}^{r})$ the variety of lines in
$\check{\mathbb P}^{r}$.  A point
$D\in\operatorname{Gr}(1,\check{\mathbb P}^{r})$ is called a
\underline{pencil of hyperplanes} in $\mathbb P^r$.  Regarding $D$ as
a line in $\check{\mathbb P}^{r}$, we write, by abuse of notation,
$D=\{H_t\}_{t\in D}$, where the $H_t$ are the hyperplanes of
$\mathbb P^r$ which are the points of the line $D$.

The \underline{axis of $D$} is the codimension-two linear subvariety
of $\mathbb P^r$ which is the intersection $H_0\cap H_\infty$ of any
two distinct members of the pencil.  Regarding this axis as a point
of $\operatorname{Gr}(r-2,\mathbb P^r)$, the variety of
$(r-2)$-dimensional linear subvarieties of $\mathbb P^r$, the map
sending $D$ to its axis is the familiar orthogonality isomorphism
\[
\operatorname{Gr}(1,\check{\mathbb P}^{r})
\xrightarrow{\ \sim\ }
\operatorname{Gr}(r-2,\mathbb P^r).
\]

\subsubsection*{2.2.}
Let $X$ be a proper, smooth, irreducible $k$-scheme of dimension
$n\geq1$, equipped with an embedding $X\hookrightarrow\mathbb P^r$.
A pencil $D=\{H_t\}_{t\in D}$ of hyperplanes in $\mathbb P^r$ is
called a \underline{Lefschetz pencil} (relative to this embedding) if:
\begin{enumerate}
\item[(a)] the axis of $D$ meets $X$ transversely (EGA IV, 17.13.7);
\item[(b)] there is a dense open subset $U\subset D$ such that, for
every $t\in U$, the hyperplane $H_t$ meets $X$ transversely;
\item[(c)] for $t_0\in D-U$, the hyperplane $H_{t_0}$ meets $X$
transversely except at one point, which is an ordinary quadratic
singular point of $X\cap H_{t_0}$.
\end{enumerate}
We shall see below in 3.2.1 that the Lefschetz pencils form an
\underline{open subset} of
$\operatorname{Gr}(1,\check{\mathbb P}^{r})$.

\subsubsection*{2.3.}
The embedding $X\hookrightarrow\mathbb P^r$ is called a
\underline{Lefschetz embedding} if the Lefschetz pencils form a dense
open subset of $\operatorname{Gr}(1,\check{\mathbb P}^{r})$.
\footnote{More precisely, they form a dense open subset in the set of
closed points of $\operatorname{Gr}(1,\check{\mathbb P}^{r})$.  We
shall permit ourselves such abuses of language below.}

\subsubsection*{2.4.}
Recall the notion of a \underline{multiple of a given embedding}
$X\hookrightarrow\mathbb P^r$.  It is the composite of that embedding
with the $d$-th Veronese embedding
\footnote{The print calls this the $d$-th Segre embedding; the
displayed monomial map is the Veronese embedding.}
\[
S_d:\mathbb P^r\hookrightarrow
\mathbb P^{\binom{r+d}{d}-1},
\qquad
(X_0,\ldots,X_r)\longmapsto(\ldots,X^W,\ldots),
\]
where
\[
X^W=X_0^{W_0}\cdots X_r^{W_r}
\quad\text{and}\quad
W_0+\cdots+W_r=d.
\]
Thus a ``section of $X$ by a hypersurface of degree $d$'' for the
given embedding---in English, a hypersurface section of degree
$d$---is simply a hyperplane section for the $d$-th multiple of that
embedding.

\subsubsection*{Theorem 2.5.}
\underline{Let $X$ be irreducible, proper, and smooth over $k$, of
dimension $n\geq1$, and equipped with an embedding
$i:X\hookrightarrow\mathbb P^r$.  Then:}
\begin{enumerate}
\item[2.5.1.] \underline{every $d$-th multiple of $i$, with $d\geq2$,
is a Lefschetz embedding;}
\item[2.5.2.] \underline{if $k$ has characteristic zero, the given
embedding $i$ is a Lefschetz embedding.}
\end{enumerate}
The proof is divided into several pieces, which follow.
\end{filecontents*}
% Original input: components/092_expose_XVII_section_3_1.tex
\begin{filecontents*}[overwrite,noheader]{sga7ii-cumulative-input-092.tex}
\section*{3. The dual variety}

\subsubsection*{3.1.}
Let $I$ be the ideal in $\mathbb P^r$ defining $X$.  Write
$N=I/I^2$ for the conormal sheaf, locally free of rank $r-n$ on $X$,
and let $\mathbb P(N)$ be the projective bundle over $X$ formed by the
lines in $\check N$.  (When $X=\mathbb P^r$, the bundle
$\mathbb P(N)$ is empty.)  The bundle $\mathbb P(N)$ identifies with
the subvariety of $X\times\check{\mathbb P}^{r}$ formed by the pairs
$(x,H)$ for which $H$ is tangent to $X$ at $x$, by the map
\[
\nu:\mathbb P(N)\longrightarrow X\times\check{\mathbb P}^{r}.
\]
More explicitly,
\begin{equation}\tag{3.1.1}
\mathbb P(I/I^2)_x\ni F\longmapsto
\left(x,\text{ the hyperplane with equation }
\sum_i\frac{\partial F}{\partial X_i}(x)T_i=0\right),
\end{equation}
where $F$ denotes a local section of $I/I^2$.  Define
\begin{equation}\tag{3.1.2}
\varphi:\mathbb P(N)\longrightarrow\check{\mathbb P}^{r}
\end{equation}
to be the composite
\[
\mathbb P(N)\xrightarrow{\nu}
X\times\check{\mathbb P}^{r}
\xrightarrow{\operatorname{pr}_2}\check{\mathbb P}^{r}.
\]

\subsubsection*{3.1.3.}
Let $\check X$ denote the image of $\varphi$.  It is the set of
hyperplanes which do not meet $X$ transversely, and is called the
\underline{dual variety} of $X$ for the embedding
$i:X\hookrightarrow\mathbb P^r$.

\subsubsection*{Proposition 3.1.4.}
\underline{The dual variety $\check X$ is irreducible and proper, of
dimension at most $r-1$.}

\noindent\underline{Proof.}
It is the image under $\varphi$ of $\mathbb P(N)$, which is itself
irreducible of dimension $r-1$ if $X\ne\mathbb P^r$, and empty if
$X=\mathbb P^r$.

\subsubsection*{Remark 3.1.5.}
The interpretation of $\mathbb P(N)$ can be made more precise as
follows.  Let
$\mathcal Y\subset X\times\check{\mathbb P}^{r}$ be the incidence
subscheme whose fiber over every point
$H\in\check{\mathbb P}^{r}(k)$ is $X\cap H$.  The scheme $\mathcal Y$
is plainly smooth of dimension $n+r-1$, being a projective bundle of
relative dimension $r-1$ over $X$.  Thus
$\mathcal Y\to\check{\mathbb P}^{r}$ makes $\mathcal Y$ a relative
complete intersection of virtual relative dimension
$(n+r-1)-r=n-1$.

The \underline{Jacobian subscheme}
\[
J_{n-1}(\mathcal Y/\check{\mathbb P}^{r})
\]
(VI, 5.2) is then precisely $\mathbb P(N)$.  Formation of the
Jacobian subscheme of a relative complete intersection is compatible
with arbitrary base change.  Consequently, for every line
$D\subset\check{\mathbb P}^{r}$, the Jacobian subscheme
$J_{n-1}(\widetilde X/D)$, where
\[
\widetilde X=\mathcal Y\times_{\check{\mathbb P}^{r}}D,
\]
is precisely
$\mathbb P(N)\times_{\check{\mathbb P}^{r}}D$.
\end{filecontents*}
% Original input: components/093_expose_XVII_section_3_2_and_proposition_3_3.tex
\begin{filecontents*}[overwrite,noheader]{sga7ii-cumulative-input-093.tex}
\subsubsection*{Proposition 3.2.}
\underline{The subset $B\subset\check X$ formed by the hyperplanes
$H$ such that $X\cap H$ has exactly one singular point, and that point
is an ordinary quadratic singularity, is open.}

This is a special case of XV, 1.3.2, based on R. Elkik's theory of
``henselian modules'' (XV, 1.1.4).  The case in which one restricts to
\underline{nondegenerate} quadratic points---that is, where
$\operatorname{char}(k)\ne2$ \underline{or} $n=\dim X$ is even---is
also immediate from 3.3 below, whose proof does not use 3.2.

\subsubsection*{Corollary 3.2.1.}
\underline{Under the hypotheses of 3.2, let $F=\check X-B$, a closed
subset of $\check X$.  Then the Lefschetz pencils form an open subset
of $\operatorname{Gr}(1,\check{\mathbb P}^{r})$, and the following
conditions are equivalent:}
\begin{enumerate}
\item[3.2.2.] \underline{the given embedding
$i:X\hookrightarrow\mathbb P^r$ is a Lefschetz embedding;}
\item[3.2.3.] \underline{there exists a Lefschetz pencil of
hyperplanes for the given embedding;}
\item[3.2.4.] \underline{the codimension of $F$ in
$\check{\mathbb P}^{r}$ is at least two.}
\end{enumerate}

\noindent\underline{Proof.}
A line $D\subset\check{\mathbb P}^{r}$ is a Lefschetz pencil if and
only if:
\begin{enumerate}
\item[(3.2.5)] the axis of $D$ meets $X$ transversely;
\item[(3.2.6)] $D$ is not contained in $\check X$;
\item[(3.2.7)] $D$ does not meet $F$.
\end{enumerate}
Conditions (3.2.5) and (3.2.6) define open subsets of
$\operatorname{Gr}(1,\check{\mathbb P}^{r})$; since $F$ is closed, so
does (3.2.7).  Hence the Lefschetz pencils form an open subset.

Because $\dim X\leq r-1$, the open subset in (3.2.6) is nonempty and
therefore dense.  The open subset in (3.2.5) is likewise nonempty and
dense.  Indeed, choose a hyperplane $H$ meeting $X$ transversely; such
an $H$ exists because $\dim X\leq r-1$.  Since
$\dim(X\cap H)\leq r-1$, there is a hyperplane $G$ meeting
$X\cap H$ transversely, and the line $D$ through $H$ and $G$ then has
an axis meeting $X$ transversely.  Thus the embedding is Lefschetz if
and only if the open subset in (3.2.7) is nonempty, hence dense.  Since
$F$ is closed, this occurs if and only if $F$ has codimension at least
two in $\check{\mathbb P}^{r}$ (cf. EGA V, or Mumford,
\emph{Introduction to Algebraic Geometry}, p.~88, Cor.~3).

\subsubsection*{Corollary 3.2.8.}
\underline{Under the hypotheses of 3.2, suppose that
$i:X\hookrightarrow\mathbb P^r$ is a Lefschetz embedding.  Let $H_0$
be a hyperplane meeting $X$ transversely, and let
$\operatorname{Gr}(1,\check{\mathbb P}^{r})_{H_0}$ be the subvariety
of $\operatorname{Gr}(1,\check{\mathbb P}^{r})$ formed by the lines
through $H_0$.  Then the Lefschetz pencils through $H_0$ form a dense
open subset of $\operatorname{Gr}(1,\check{\mathbb P}^{r})_{H_0}$.}

\noindent\underline{Proof.}
By 3.2.1 these pencils form an open subset of
$\operatorname{Gr}(1,\check{\mathbb P}^{r})_{H_0}$.  The argument
just given to prove that the open subsets (3.2.5) and (3.2.6) are
nonempty also shows that their intersections with this subvariety are
nonempty.  It remains to prove that some lines through $H_0$ do not
meet $F$.

Choose a hyperplane $T\subset\check{\mathbb P}^{r}$ not containing
$H_0$.  There is an isomorphism
\begin{center}
\begin{tikzcd}[column sep=huge]
\operatorname{Gr}(1,\check{\mathbb P}^{r})_{H_0}
  & T \arrow[l,"\sim"']\\
\substack{\text{the line through}\\ H_0\text{ and }\tau}
  & \tau \arrow[l,mapsto]
\end{tikzcd}
\end{center}
so $\operatorname{Gr}(1,\check{\mathbb P}^{r})_{H_0}$ ``is'' a
projective space of dimension $r-1$.  The lines in it which pass
through a point of $F$ form a closed subset of dimension at most
$r-2$, being the image of the morphism
\begin{center}
\begin{tikzcd}[column sep=huge]
F \arrow[r] & \operatorname{Gr}(1,\check{\mathbb P}^{r})_{H_0}\\
f \arrow[r,mapsto] &
  \substack{\text{the line through }H_0\\\text{and through }f}.
\end{tikzcd}
\end{center}

\subsubsection*{Proposition 3.3.}
\underline{The morphism
$\varphi:\mathbb P(N)\to\check{\mathbb P}^{r}$ is unramified at a
closed point $(x_0,H)$ if and only if $x_0$ is a nondegenerate
quadratic singularity (1.2) of $X\cap H$.  In particular, the subset
of $\mathbb P(N)$ formed by the points $(x,H)$ for which $x$ is a
nondegenerate quadratic singularity of $X\cap H$ is open.}
\end{filecontents*}
% Original input: components/094_expose_XVII_proof_3_3_coordinates.tex
\begin{filecontents*}[overwrite,noheader]{sga7ii-cumulative-input-094.tex}
\noindent\underline{Proof of Proposition 3.3.}
We may restrict to the case $X\ne\mathbb P^r$.  The two properties
asserted to be equivalent are independent of the particular choice of
coordinates, so we may describe them using coordinates chosen in a
very special way.  Since $X$ is a smooth subvariety of $\mathbb P^r$
of dimension $n\geq1$, choose homogeneous coordinates
$X_0,\ldots,X_r$ so that:
\begin{enumerate}
\item[3.3.1.] the point $x_0\in\mathbb P^r$ is $(1,0,\ldots,0)$;
\item[3.3.2.] putting $x_i=X_i/X_0$ for $i=1,\ldots,r$, so that
$\widehat{\mathcal O}_{\mathbb P^r,x_0}$ identifies with
$k[[x_1,\ldots,x_r]]$, the ideal defining the trace of $X$ in
$\operatorname{Spec}\widehat{\mathcal O}_{\mathbb P^r,x_0}$ is
generated by elements $g_i$ $(i=1,\ldots,r-n)$ which, in the
completion, have the form
\[
g_i=x_{n+i}-h_i(x_1,\ldots,x_n),
\qquad i=1,\ldots,r-n,
\]
where the $h_i$ are power series in $x_1,\ldots,x_n$ alone and lie in
the square of the maximal ideal of $k[[x_1,\ldots,x_n]]$;
\item[3.3.3.] under (3.1.1), the hyperplane $H$ comes from
$g_{r-n}$.
\end{enumerate}

Put $V=\operatorname{Spec}(\widehat{\mathcal O}_{X,x_0})$.  It follows
from 3.3.2 that there is a trivialization
\[
(I/I^2)|_V\simeq\mathcal O_V^{r-n}
\]
given by
\begin{equation}\tag{3.3.4}
(\Lambda_1,\ldots,\Lambda_{r-n})\longmapsto
\sum_{i=1}^{r-n}\Lambda_i g_i.
\end{equation}
Together with 3.3.3 this gives a neighborhood $W$ of $(x_0,H)$ in
$\mathbb P(N)|_V$ and an isomorphism
$V\times\mathbb A^{r-n-1}\simeq W$, given by
\begin{equation}\tag{3.3.5}
\begin{aligned}
V\times\mathbb A^{r-n-1}&\longrightarrow\mathbb P(I/I^2)|_V,\\
(v,(\lambda_1,\ldots,\lambda_{r-n-1}))&\longmapsto
\left(v,\left[g_{r-n}+\sum_{i=1}^{r-n-1}\lambda_i g_i\right]\right),
\end{aligned}
\end{equation}
where $\lambda_i=\Lambda_i/\Lambda_{r-n}$.

In particular, the completion of the local ring of $\mathbb P(N)$ at
$(x_0,H)$ identifies with
\[
k[[x_1,\ldots,x_n,\lambda_1,\ldots,\lambda_{r-n-1}]].
\]
We now express $\varphi$ in these coordinates at the level of
completions:
\begin{equation}\tag{3.3.6}
\varphi(x,\lambda)
=\text{the hyperplane with equation }
\sum_{i=0}^{r}\varphi_i(x,\lambda)T_i=0.
\end{equation}
Since this image hyperplane contains the point $x$, whose homogeneous
coordinates are
\begin{equation}\tag{3.3.7}
x=(1,x_1,\ldots,x_n,h_1(x),\ldots,h_{r-n}(x)),
\end{equation}
we have
\begin{equation}\tag{3.3.8}
\varphi_0(x,\lambda)
=-\sum_{i=1}^{n}x_i\varphi_i(x,\lambda)
 -\sum_{j=1}^{r-n}h_j(x)\varphi_{n+j}(x,\lambda)
\end{equation}
in $\widehat{\mathcal O}_{\mathbb P(N),(x_0,H)}$.  Moreover, by
(3.1.1), if
\[
D_i=\frac{\partial}{\partial x_i}\in
\operatorname{Der}_k(
\widehat{\mathcal O}_{\mathbb P^r,x_0},
\widehat{\mathcal O}_{\mathbb P^r,x_0}),
\]
then for $1\leq i\leq r$,
\begin{equation}\tag{3.3.9}
\varphi_i(x,\lambda)
=D_i\!\left(g_{r-n}+
\sum_{j=1}^{r-n-1}\lambda_jg_j\right)(x)
\quad\text{in }\widehat{\mathcal O}_{\mathbb P(N),(x_0,H)}.
\end{equation}

Recalling that $g_i=x_{n+i}-h_i(x_1,\ldots,x_n)$, one obtains
\begin{equation}\tag{3.3.10}
\varphi_i(x,\lambda)
=-\frac{\partial}{\partial x_i}
\left(h_{r-n}+\sum_{j=1}^{r-n-1}\lambda_jh_j\right),
\qquad 1\leq i\leq n,
\end{equation}
\begin{equation}\tag{3.3.11}
\varphi_{n+j}(x,\lambda)=\lambda_j,
\qquad 1\leq j\leq r-n-1,
\end{equation}
and
\begin{equation}\tag{3.3.12}
\varphi_r(x,\lambda)=1.
\end{equation}
\footnote{The print gives incompatible ranges in (3.3.11).  The
coordinates in (3.3.5), the coefficient formula (3.3.9), and the
identity block of the printed Jacobian matrix force the indices shown
here.}

Let $\mathfrak m$ be the maximal ideal of
$\widehat{\mathcal O}_{\mathbb P(N),(x_0,H)}$.  Since
$h_i\equiv0\pmod{\mathfrak m^2}$ for $i=1,\ldots,r-n$, we find
\begin{equation}\tag{3.3.13}
\varphi_0(x,\lambda)\equiv0\pmod{\mathfrak m^2}
\end{equation}
and
\begin{equation}\tag{3.3.14}
\varphi_i(x,\lambda)
\equiv-\frac{\partial h_{r-n}}{\partial x_i}
\pmod{\mathfrak m^2},
\qquad i=1,\ldots,n.
\end{equation}
\end{filecontents*}
% Original input: components/095_expose_XVII_proof_3_3_matrix_and_example_3_4.tex
\begin{filecontents*}[overwrite,noheader]{sga7ii-cumulative-input-095.tex}
Putting (3.3.11)--(3.3.14) together, the Jacobian matrix of
$\varphi$ at $(x_0,H)$ has the block form
\begin{equation}\tag{3.3.15}
J_{(x_0,H)}(\varphi)=
\left(
\begin{array}{c|c}
0_{1\times n} & 0_{1\times(r-n-1)}\\ \hline
-\operatorname{Hess}(h_{r-n})(0) & 0_{n\times(r-n-1)}\\ \hline
0_{(r-n-1)\times n} & I_{r-n-1}\\ \hline
0_{1\times n} & 0_{1\times(r-n-1)}
\end{array}
\right).
\end{equation}
Thus $\varphi$ is unramified at $(x_0,H)$ if and only if the Hessian
matrix
\[
\left(
\frac{\partial^2h_{r-n}}{\partial x_i\partial x_j}(0)
\right)_{1\leq i,j\leq n}
\]
is invertible.  On the other hand, 3.3.2 and 3.3.3 show that the
completion of the local ring of $X\cap H$ at $x_0$ is
\begin{equation}\tag{3.3.16}
k[[x_1,\ldots,x_n]]/
\bigl(h_{r-n}(x_1,\ldots,x_n)\bigr),
\end{equation}
and the result follows from Definition 1.2.

\subsubsection*{Example 3.4.}
Suppose that $X$ is a hypersurface in $\mathbb P^r$, defined by a
homogeneous equation $F(X_0,\ldots,X_r)=0$.  Then
$\mathbb P(N)=X$, and in homogeneous coordinates
$\varphi:X\to\check{\mathbb P}^{r}$ is given by
\begin{equation}\tag{3.4.1}
\varphi(X_0,\ldots,X_r)=
\left(
\frac{\partial F}{\partial X_0},\ldots,
\frac{\partial F}{\partial X_r}
\right).
\end{equation}
This is the \underline{Gauss morphism}.

Here is an example in which $\varphi$ is everywhere ramified.  Let
$p$ be the characteristic of $k$, let $r\geq1$, and take the
hypersurface $X$ with equation
\begin{equation}\tag{3.4.2}
\sum_{i=0}^{r}X_i^{p^r+1}=0.
\end{equation}
Then
\begin{equation}\tag{3.4.3}
\varphi(X_0,\ldots,X_r)
=(X_0^{p^r},\ldots,X_r^{p^r}).
\end{equation}
Thus $\check X$ is reduced to a point if $p=0$, while if $p>0$ one
has $\check X\simeq X$ and $\varphi$ is the Frobenius morphism.  It
follows from 3.3 (or from 3.5.0 below) that if $p>2$, this embedding is
not a Lefschetz embedding.

We now return to the general case.
\end{filecontents*}
% Original input: components/096_expose_XVII_proposition_3_5_and_opening_proof.tex
\begin{filecontents*}[overwrite,noheader]{sga7ii-cumulative-input-096.tex}
\subsubsection*{Proposition 3.5.}
\underline{Suppose that the morphism
$\varphi:\mathbb P(N)\to\check{\mathbb P}^{r}$ is generically
unramified (that is, is not everywhere ramified).  Then $\varphi$
induces a birational morphism from $\mathbb P(N)$ onto $\check X$.
More precisely, the following three open subsets of $\check X$ are
identical:}
\begin{enumerate}
\item[(a)] \underline{the largest open subset on which $\check X$ is
smooth;}
\item[(b)] \underline{the largest open subset $V\subset\check X$ such
that $\varphi$ induces an isomorphism
$\varphi^{-1}(V)\xrightarrow{\sim}V$;}
\item[(c)] \underline{the open subset (cf. 3.2) formed by those
hyperplanes $H$ for which $X\cap H$ has exactly one singular point,
that point being a nondegenerate quadratic singularity.}
\end{enumerate}

\subsubsection*{Corollary 3.5.0.}
\underline{Suppose that $\operatorname{char}(k)\ne2$ or that $\dim X$
is even.  The given embedding is a Lefschetz embedding if and only if
one of the following two conditions holds:}
\begin{enumerate}
\item[(a)] $\varphi$ \underline{is generically unramified;}
\item[(b)] $\dim\check X\leq r-2$.
\end{enumerate}

Indeed, sufficiency follows from criterion (3.2.4).  In case (a),
Proposition 3.5 identifies $F$ with the singular locus of $\check X$,
so
\[
\dim F<\dim\check X=r-1;
\]
in case (b), the criterion is immediate.  Conversely, suppose that
$D$ is a Lefschetz pencil and $\dim\check X=r-1$.  There is a point
$H\in D\cap\check X$, and 3.3 together with Definition 2.2(c) shows
that $\varphi$ is unramified at the points of $\varphi^{-1}(H)$.
Hence it is generically unramified.

\noindent\underline{Proof of Proposition 3.5.}
The hypothesis implies
\[
\dim\check X=\dim\mathbb P(N)=r-1,
\]
and therefore 3.2 shows that $\check X$ is an irreducible hypersurface
in $\check{\mathbb P}^{r}$.

Let $R$ be the ramification locus of $\varphi$.  By hypothesis it is
a closed subset of $\mathbb P(N)$ of dimension at most $r-2$; hence
$\varphi(R)$ is a closed subset of $\check X$ of dimension at most
$r-2$.  Let $V'$ be the dense open subset
\begin{equation}\tag{3.5.1}
V'=(\check X-\varphi(R))\cap
   (\check X-\operatorname{Sing}(\check X))
\end{equation}
of $\check X$, and denote its complement by $F$:
\begin{equation}\tag{3.5.1.1}
F=\varphi(R)\cup\operatorname{Sing}(\check X).
\end{equation}

We first prove that $\varphi$ induces an isomorphism
$\varphi^{-1}(V')\simeq V'$ by constructing an inverse
$\psi:V'\to\varphi^{-1}(V')$.  Regard $\varphi^{-1}(V')$ as embedded
in $X\times V'$ by (3.1.1), so that $\varphi$ is the second
projection.  Constructing $\psi$ is then equivalent to constructing a
morphism $\psi_1:V'\to\mathbb P^r$ which makes the diagram
\begin{equation}\tag{3.5.2}
\begin{tikzcd}[column sep=small]
\mathbb P(N) \arrow[dd]
  & {}\supset\varphi^{-1}(V') \arrow[r,"\varphi"]
  & V' \arrow[dd] & \\
  & & & \psi_1 \\
X \arrow[rr,hook] & & \mathbb P^r &
\end{tikzcd}
\end{equation}
commute; one then puts
$\psi(v)=(\psi_1(v),v)\in\mathbb P^r\times V'$.

For any hypersurface in projective space, the Gauss morphism (3.4.1)
is defined at every smooth point.  Apply this to the hypersurface
$\check X\subset\check{\mathbb P}^{r}$.  Since $V'$ lies in the
smooth locus $V$ of $\check X$, the Gauss morphism gives
\begin{equation}\tag{3.5.3}
\psi_1:V'\longrightarrow
\check{\check{\mathbb P}}^{r}\simeq\mathbb P^r.
\end{equation}
It remains to check that (3.5.2) commutes.  Let $(x_0,H_0)$ be a
closed point of $\varphi^{-1}(V')$.  We must show that $x_0$, viewed
as a hyperplane in $\check{\mathbb P}^{r}$, is tangent to $\check X$
at $H_0$: in other words, tangency of $H_0$ to $X$ at $x_0$ implies
tangency of $x_0$ to $\check X$ at $H_0$.  Since $\varphi$ is
unramified at $(x_0,H_0)$ and $\check X$ is smooth at
$H_0=\varphi(x_0,H_0)$, its differential is an isomorphism
\begin{equation}\tag{3.5.4}
d\varphi:
T_{\mathbb P(N),(x_0,H_0)}\xrightarrow{\sim}T_{\check X,H_0}.
\end{equation}
The first part of the proof is therefore completed by the following
lemma.
\end{filecontents*}
% Original input: components/097_expose_XVII_lemma_3_6.tex
\begin{filecontents*}[overwrite,noheader]{sga7ii-cumulative-input-097.tex}
\subsubsection*{Lemma 3.6 (Euler).}
\underline{Let $(x_0,H_0)$ be a closed point of $\mathbb P(N)$, and
let $T_L$ denote the tangent bundle of the hyperplane
$L\subset\check{\mathbb P}^{r}$ defined by $x_0$.  Then, inside
$T_{\check{\mathbb P}^{r},H_0}$,}
\[
d\varphi\bigl(T_{\mathbb P(N),(x_0,H_0)}\bigr)
\subset T_{L,H_0}.
\]

\noindent\underline{Proof.}
Choose homogeneous coordinates $X_0,\ldots,X_r$ on $\mathbb P^r$ in
which $x_0$ is the point
$\underline A=(A_0,\ldots,A_r)$.  A point of $T_{X,x_0}$ is a point
of $\mathbb P^r$ with values in $k[\epsilon]/(\epsilon^2)$ of the
form
\begin{equation}\tag{3.6.1}
\underline A+\epsilon\underline B,
\end{equation}
such that, for every local section $H$ of $I$,
\begin{equation}\tag{3.6.2}
\sum_{i=0}^{r}B_i\frac{\partial H}{\partial X_i}(\underline A)=0.
\end{equation}
Choose also a local section $H$ of $I$ which gives rise to $H_0$ by
(3.1.1).  Since $\mathbb P(N)$ is locally trivial over $X$, a point
of $T_{\mathbb P(N),(x_0,H_0)}$ can be represented as
\begin{equation}\tag{3.6.3}
(\underline A+\epsilon\underline B,\ H+\epsilon F),
\end{equation}
where $F$ is a local section of $I$ and $\underline A,\underline B$
are as above.  Its image under $\varphi$ is the hyperplane over
$k[\epsilon]/(\epsilon^2)$ with equation
\begin{equation}\tag{3.6.4}
\sum_{i=0}^{r}
\frac{\partial(H+\epsilon F)}{\partial X_i}
(\underline A+\epsilon\underline B)T_i=0.
\end{equation}
Thus the assertion is equivalent to
\begin{equation}\tag{3.6.5}
\sum_{i=0}^{r}A_i
\frac{\partial(H+\epsilon F)}{\partial X_i}
(\underline A+\epsilon\underline B)=0.
\end{equation}
Because $\epsilon^2=0$, Taylor's formula gives
\begin{equation}\tag{3.6.6}
\frac{\partial(H+\epsilon F)}{\partial X_i}
(\underline A+\epsilon\underline B)
=\frac{\partial H}{\partial X_i}(\underline A)
+\epsilon\frac{\partial F}{\partial X_i}(\underline A)
+\epsilon\sum_{j=0}^{r}B_j
\frac{\partial^2H}{\partial X_j\partial X_i}(\underline A).
\end{equation}
It is therefore enough to verify
\begin{equation}\tag{3.6.7}
\sum_{i=0}^{r}A_i\frac{\partial H}{\partial X_i}(\underline A)=0,
\end{equation}
\begin{equation}\tag{3.6.8}
\sum_{i=0}^{r}A_i\frac{\partial F}{\partial X_i}(\underline A)=0,
\end{equation}
and
\begin{equation}\tag{3.6.9}
\sum_{i,j=0}^{r}A_iB_j
\frac{\partial^2H}{\partial X_i\partial X_j}(\underline A)=0.
\end{equation}
Here $H$ and $F$ are represented by homogeneous functions of degree
zero, so Euler's formula gives
\begin{equation}\tag{3.6.10}
\sum_iX_i\frac{\partial H}{\partial X_i}=0,
\qquad
\sum_iX_i\frac{\partial F}{\partial X_i}=0.
\end{equation}
Since $\partial H/\partial X_j$ is homogeneous of degree $-1$, the
same formula gives
\[
\sum_iX_i\frac{\partial^2H}{\partial X_i\partial X_j}
=-\frac{\partial H}{\partial X_j}.
\]
Consequently, (3.6.9) becomes
\begin{equation}\tag{3.6.11}
-\sum_jB_j\frac{\partial H}{\partial X_j}(\underline A)=0,
\end{equation}
which is precisely (3.6.2).

This proves that $\varphi$ induces an isomorphism
\begin{equation}\tag{3.6.12}
\varphi^{-1}(V')\xrightarrow{\sim}V',
\end{equation}
whose inverse is $\psi$.
\end{filecontents*}
% Original input: components/098_expose_XVII_completion_of_proof_3_5.tex
\begin{filecontents*}[overwrite,noheader]{sga7ii-cumulative-input-098.tex}
Now observe that the morphism
\begin{equation}\tag{3.6.13}
\left\{
\begin{aligned}
\psi &:V'\longrightarrow
  \mathbb P^r\times V'\subset\mathbb P^r\times\check X,\\
\psi(v)&=(\psi_1(v),v)
\end{aligned}
\right.
\end{equation}
extends to a morphism
\begin{equation}\tag{3.6.14}
\left\{
\begin{aligned}
\widetilde\psi &:V\longrightarrow\mathbb P^r\times V,\\
\widetilde\psi(v)&=(\widetilde\psi_1(v),v),
\end{aligned}
\right.
\end{equation}
where
\begin{equation}\tag{3.6.15}
\widetilde\psi_1:V\longrightarrow\mathbb P^r
\end{equation}
is the Gauss morphism, now regarded as defined on all of $V$.

We have just proved
\begin{equation}\tag{3.6.16}
\psi(V')\subset\mathbb P(N).
\end{equation}
Since $V'$ is dense in $V$ and $\mathbb P(N)$ is closed in
$\mathbb P^r\times\check X$, this implies
\begin{equation}\tag{3.6.17}
\widetilde\psi(V)\subset\mathbb P(N),
\end{equation}
and in fact plainly
\begin{equation}\tag{3.6.18}
\widetilde\psi(V)\subset\varphi^{-1}(V).
\end{equation}
Again by continuity, the identities
\[
\psi\circ\varphi=\operatorname{id}_{\varphi^{-1}(V')},
\qquad
\varphi\circ\psi=\operatorname{id}_{V'}
\]
extend to
\[
\widetilde\psi\circ\varphi=\operatorname{id}_{\varphi^{-1}(V)},
\qquad
\varphi\circ\widetilde\psi=\operatorname{id}_{V}.
\]
Thus the open subsets in 3.5(a) and 3.5(b) are identical.

We next show that the open subsets in 3.5(a) and 3.5(c) are
identical.  By 3.3, the latter is the subset $U\subset\check X$
defined by
\begin{equation}\tag{3.6.19}
U=\left\{H\in\check X\ \middle|\
\begin{array}{l}
\varphi^{-1}(H)\text{ consists of exactly one point, and}\\
\varphi\text{ is unramified at that point}
\end{array}
\right\}.
\end{equation}
The scheme $\varphi^{-1}(U)$ is normal, since it is smooth, and the
induced morphism
\[
\varphi:\varphi^{-1}(U)\longrightarrow U
\]
is birational, as already proved, radicial by the definition of $U$,
and proper.  It is therefore the normalization of $U$.  In
particular, $U$ is geometrically unibranch.  Since this morphism is
unramified, EGA IV, 18.10.1 shows that it is étale; EGA IV, 18.10.18
then shows that it is an open immersion.  It is consequently an
isomorphism, and $U$ is smooth.  This proves
$3.5(c)\subset 3.5(a)$.  Formula (3.6.19) makes the inclusion
$3.5(b)\subset 3.5(c)$ clear, and the equality
$3.5(a)=3.5(b)$ completes the proof.
\end{filecontents*}
% Original input: components/099_expose_XVII_section_3_7.tex
\begin{filecontents*}[overwrite,noheader]{sga7ii-cumulative-input-099.tex}
\subsubsection*{3.7.}
Theorem 2.5 is true when $\operatorname{char}(k)\ne2$, and also when
$X$ has even dimension.

\noindent\underline{Proof.}
In view of 3.5.0 and 3.3, it is enough to verify the following
assertion:

\medskip
\noindent(3.7.1) For every closed point $x_0\in X$ and every integer
$d\geq2$, there is a hypersurface $H$ of degree $d$, tangent to $X$
at $x_0$, such that $x_0$ is an ordinary quadratic singularity of
$X\cap H$.

Choose homogeneous coordinates $X_0,\ldots,X_r$ so that
$x_0=(1,0,\ldots,0)$ and the functions $X_i/X_0$, $i=1,\ldots,n$,
are local coordinates on $X$ near $x_0$.  One may take $H$ to be the
hypersurface with equation
\begin{equation}\tag{3.7.1.1}
H=X_0^{d-2}\sum_{i=1}^{n}X_i^2
\qquad\text{if }\operatorname{char}(k)\ne2,
\end{equation}
or
\begin{equation}\tag{3.7.1.2}
H=X_0^{d-2}\sum_{i=1}^{n/2}X_iX_{n/2+i}
\qquad\text{if }n\text{ is even}.
\end{equation}
When $n$ is odd (under the first hypothesis), the equivalent standard
choice recorded in the source is
\begin{equation}\tag{3.7.1.3}
H=X_0^{d-2}
\left(X_n^2+
\sum_{i=1}^{\lfloor n/2\rfloor}
X_iX_{\lfloor n/2\rfloor+i}\right).
\end{equation}

For assertion 2.5.2, note that either $\varphi$ is generically
unramified, in which case 3.5.0 applies, or it is everywhere ramified.
In characteristic zero the latter case implies
$\dim\check X\leq r-2$, so criterion (3.2.4) is then automatically
satisfied.\footnote{The print has $\dim X\leq r-2$ here.  The
alternative in Corollary 3.5.0 and the criterion being invoked require
$\dim\check X\leq r-2$.}
\end{filecontents*}
% Original input: components/100_expose_XVII_section_4_1_and_proposition_4_2.tex
\begin{filecontents*}[overwrite,noheader]{sga7ii-cumulative-input-100.tex}
\section*{4. The general case}

We shall sketch another method of proving 2.5, valid in every
characteristic, assuming the following statement, which is close to
3.1 and is a special case of XV, 1.1.4.

\medskip
\noindent(4.1) \underline{The subset of $\mathbb P(N)$ formed by the
points $(x,H)$ such that $x$ is an ordinary quadratic singularity of
$X\cap H$ is open.}

Observe that 4.1 follows from 3.3, except when
$\operatorname{char}(k)=2$ and $\dim X$ is odd.

\subsubsection*{4.1.1.}
Let $R_1$ be the closed complement in $\mathbb P(N)$ of the open set
4.1, and let
\[
F_1=\varphi(R_1)\subset\check X\subset\check{\mathbb P}^{r}
\]
be its image.

\subsubsection*{Lemma 4.1.2.}
\underline{If the open set 4.1 is nonempty, then
$\operatorname{codim}_{\check{\mathbb P}^{r}}(F_1)\geq2$.}

\noindent\underline{Proof.}
By hypothesis,
\[
\dim R_1<\dim\mathbb P(N)=r-1,
\]
and $F_1=\varphi(R_1)$.

\subsubsection*{4.1.3.}
Let $F_2$ be the subset of $\check X$ formed by the hyperplanes which
are tangent to $X$ at at least two points.  Notice that $F_2$ is not
closed, but that
\begin{equation}\tag{4.1.4}
\overline{F_2}\subset F_1\cup F_2.
\end{equation}
This follows from the fact that $F_1\cup F_2=F$ (cf. 3.2) is closed.
(Inclusion (4.1.4) ``says'' that a \underline{confluence} of two
singularities is not an ordinary quadratic singularity.)  In any
case, the construction in 4.2.4 will show that $F_2$ is constructible.

To prove 2.5.1 we shall apply criterion 3.2.4.  For this we need to
know that
\[
\operatorname{codim}_{\check{\mathbb P}^{r}}(F_1)\geq2
\quad\text{and}\quad
\operatorname{codim}_{\check{\mathbb P}^{r}}(F_2)\geq2.
\]
On the one hand, 3.7.1 and 4.1.2 show that, after replacing the given
embedding by its $d$-th multiple, $d\geq2$, one has
$\operatorname{codim}_{\check{\mathbb P}^{r}}(F_1)\geq2$.
\footnote{The print reads ``3.8.1'' and $F_2$ in this sentence.  The
ordinary-quadratic construction is (3.7.1), and Lemma 4.1.2 applies
to $F_1$; Proposition 4.2 supplies the required assertion for $F_2$.}
Thus 2.5.1 will follow once we know the following proposition.

\subsubsection*{Proposition 4.2.}
\underline{Let $X$ be a smooth irreducible subvariety of
$\mathbb P^r$, of dimension $n\geq1$.  For every integer $d\geq2$,
the subset $F_2$ of}
\[
\check{\mathbb P}^{N}
=\check{\mathbb P}^{\binom{r+d}{d}-1}
\]
\underline{formed by the hypersurfaces of degree $d$ in
$\mathbb P^r$ which are tangent to $X$ at at least two points has
codimension at least $2$ in $\check{\mathbb P}^{N}$.}
\footnote{The print has $\binom{n+d}{d}-1$ here and
$\binom{r+d}{d}-d$ in Proposition 4.3.  The degree-$d$ Veronese
embedding defined in 2.4 gives
$N=\binom{r+d}{d}-1$ in both places.}
\end{filecontents*}
% Original input: components/101_expose_XVII_proof_of_proposition_4_2.tex
\begin{filecontents*}[overwrite,noheader]{sga7ii-cumulative-input-101.tex}
\subsubsection*{4.2.1.}
Begin with the special case $d=2$ and $X$ a linear subvariety of
$\mathbb P^r$.  It is enough to find a line in
$\check{\mathbb P}^{N}$ which does not meet $F$ (cf. 3.2.4).  Choose
homogeneous coordinates $X_0,\ldots,X_r$ on $\mathbb P^r$ so that
$X$ is defined by
\[
X_{n+1}=\cdots=X_r=0,
\]
and homogeneous coordinates $(\lambda,\mu)$ on $\mathbb P^1$.

For $n$ even, $n=2k$, take the pencil of quadrics
\begin{equation}\tag{4.2.2}
H_{(\lambda,\mu)}
=\lambda X_0^2+(\lambda X_k+\mu X_0)X_{2k}
 +\sum_{i=1}^{k-1}(\lambda X_i+\mu X_{i+1})X_{k+i}.
\end{equation}
For $n$ odd, $n=2k+1$, take the pencil
\begin{equation}\tag{4.2.3}
H_{(\lambda,\mu)}
=(\lambda X_1+\mu X_0)X_0
 +\sum_{i=1}^{k}(\lambda X_{2i}+\mu X_i)X_{2i+1}.
\end{equation}
In both cases one checks that $X\cap H_{(\lambda,\mu)}$ is smooth for
$(\lambda,\mu)\ne(0,1)$, while $X\cap H_{(0,1)}$ has just one
singular point, which is moreover an ordinary quadratic singularity.
Thus the pencil does not meet $F$.

\subsubsection*{4.2.4.}
Let $\operatorname{Btg}^{d}(X)$ be the subvariety of
\[
\check{\mathbb P}^{N}\times
\bigl((X\times X)-\Delta(X\times X)\bigr)
\]
formed by triples $(H,x_1,x_2)$ such that $H$ is a hypersurface of
degree $d$ in $\mathbb P^r$, tangent to $X$ at $x_1$ and at $x_2$,
with $x_1\ne x_2$.  The image of $\operatorname{Btg}^{d}(X)$ under
$\operatorname{pr}_1$ is the set $F_2$ of Proposition 4.2, so it is
enough to prove
\begin{equation}\tag{4.2.5}
\dim\operatorname{Btg}^{d}(X)\leq N-2
\quad\text{if $d\geq3$, or if $d=2$ and $X$ is not linear.}
\end{equation}

For this purpose, let $\operatorname{Lin}(X\times X)$ be the
subvariety of $(X\times X)-\Delta(X\times X)$ given by
\begin{equation}\tag{4.2.6}
\begin{aligned}
\operatorname{Lin}(X\times X)
&=\bigl\{(x_1,x_2)\in X\times X\mid x_1\ne x_2,\\[-1ex]
&\hspace{8em}\text{the line spanned by $x_1,x_2$ is tangent to $X$ at $x_1$}
\bigr\}\\
&=\bigl\{(x_1,x_2)\in X\times X\mid x_1\ne x_2,\\[-1ex]
&\hspace{8em}x_2\in X\cap\operatorname{Tg}_X(x_1)\bigr\}.
\end{aligned}
\end{equation}
Clearly:

\medskip
\noindent(4.2.7) \underline{If $X$ is not a linear subvariety of
$\mathbb P^r$, then}
\[
\dim\operatorname{Lin}(X\times X)<2\dim X.
\]
Consequently, to prove (4.2.5) it is enough to show that the fibers of
the morphism
\begin{equation}\tag{4.2.8}
\operatorname{Btg}^{d}(X)
\xrightarrow{\ \operatorname{pr}_{2,3}\ }
(X\times X)-\Delta(X\times X),
\end{equation}
regarded as subvarieties of $\check{\mathbb P}^{N}$, have
\begin{equation}\tag{4.2.9}
\left\{
\begin{array}{ll}
\text{codimension at least $2\dim X+1$}
 &\text{over $\operatorname{Lin}(X\times X)$},\\[1ex]
\text{codimension at least $2\dim X+2$}
 &\text{over $(X\times X)-\Delta(X\times X)
                 -\operatorname{Lin}(X\times X)$.}
\end{array}
\right.
\end{equation}
Statement (4.2.9) follows from the next proposition by taking $L_i$
to be the tangent space to $X$ at $x_i$, for $i=1,2$.
\end{filecontents*}
% Original input: components/102_expose_XVII_proposition_4_3_and_proof.tex
\begin{filecontents*}[overwrite,noheader]{sga7ii-cumulative-input-102.tex}
\subsubsection*{Proposition 4.3.}
\underline{Let $x_1$ and $x_2$ be two distinct closed points of
$\mathbb P^r$, and let $L_1$ and $L_2$ be two linear subspaces of
$\mathbb P^r$ such that $x_1\in L_1$ and $x_2\in L_2$.  In the space}
\[
\check{\mathbb P}^{N}
=\check{\mathbb P}^{\binom{r+d}{d}-1}
\]
\underline{of hypersurfaces of degree $d$, the linear subspace formed
by the hypersurfaces tangent to $L_1$ at $x_1$ and to $L_2$ at $x_2$
has codimension}
\[
\left\{
\begin{aligned}
\dim L_1+\dim L_2+2,
 &\quad\text{if $d\geq3$, or if $d=2$ and $L_1\cap L_2$ does not}\\[-0.2ex]
 &\quad\text{contain the line $(x_1,x_2)$ through $x_1,x_2$},\\[1ex]
\dim L_1+\dim L_2+1,
 &\quad\text{if $d=2$ and $L_1\cap L_2$ contains the line $(x_1,x_2)$.}
\end{aligned}
\right.
\]

\noindent\underline{Proof.}
For $d\geq3$, one must prove that
$\dim L_1+\dim L_2+2$ linear conditions are independent.  It is
therefore enough to treat the case
$L_1=L_2=\mathbb P^r$, that is, to determine how many independent
conditions make a hypersurface singular at both $x_1$ and $x_2$.
Choose homogeneous coordinates so that
\[
x_1=(1,0,\ldots,0),\qquad x_2=(0,1,0,\ldots,0).
\]
The hypersurface with equation $\sum_W A_WX^W=0$ is singular at
$(1,0,\ldots,0)$ if and only if all the coefficients of the monomials
$X_0^{d-1}X_i$, $i=0,\ldots,r$, vanish.  It is singular at
$(0,1,0,\ldots,0)$ if and only if all the coefficients of the
monomials $X_1^{d-1}X_i$, $i=0,\ldots,r$, vanish.  For $d\geq3$ all
these monomials are distinct, which proves the assertion in this case.

For $d=2$, the same calculation gives
\[
\operatorname{codim}\geq\dim L_1+\dim L_2+1,
\]
because the monomial $X_0X_1$ is counted twice.

Suppose first that $L_1\cap L_2$ does not contain the line
$(x_1,x_2)$.  We reduce to the case $L_1=\mathbb P^r$ and $L_2$ is a
hyperplane in $\mathbb P^r$ which does not contain $(x_1,x_2)$.
Choose homogeneous coordinates $X_0,\ldots,X_r$ so that
\[
x_1=(1,0,\ldots,0),\qquad x_2=(0,1,0,\ldots,0),
\qquad L_2=(X_0=0).
\]
A quadric $\sum_{i\leq j}A_{ij}X_iX_j$ is singular at
$(1,0,\ldots,0)$ if and only if $A_{0j}=0$ for $j=0,\ldots,r$, and
is tangent to $X_0=0$ at $(0,1,\ldots,0)$ if and only if
$A_{1j}=0$ for $j=1,\ldots,r$.  Together these are
\[
2r+1=\dim L_1+\dim L_2+2
\]
independent conditions.

Finally, suppose that $L_1\cap L_2$ contains the line $(x_1,x_2)$.
We reduce to the case $L_1=L_2=(x_1,x_2)$.  In homogeneous
coordinates with
\[
x_1=(1,0,\ldots,0),\qquad x_2=(0,1,0,\ldots,0),
\]
the quadric $\sum_{i\leq j}A_{ij}X_iX_j$ passes through
$(1,0,\ldots,0)$ and is tangent there to $L_1$ if and only if
$A_{00}=A_{01}=0$.  It passes through $(0,1,0,\ldots,0)$ and is
tangent there to $L_2$ if and only if $A_{11}=A_{01}=0$.  Thus there
are altogether only
\[
3=\dim L_1+\dim L_2+1
\]
conditions.
\end{filecontents*}
% Original input: components/103_expose_XVII_section_5_1.tex
\begin{filecontents*}[overwrite,noheader]{sga7ii-cumulative-input-103.tex}
\section*{5. The degree of the dual variety by an
\lq\lq elementary\rq\rq\ method}

(Cf. XVIII, 3.2 for another method when $\dim X$ is even
\underline{or} $\operatorname{char}(k)\ne2$.)

We continue to use the notation of 2.2: thus
$X\hookrightarrow\mathbb P^r$ is a smooth irreducible subvariety of
dimension $n$.

\subsubsection*{5.1.}
We first consider the morphism (3.1.1)
\begin{equation}\tag{5.1.0}
\nu:\mathbb P(N)\longrightarrow X\times\check{\mathbb P}^{r}.
\end{equation}
By definition, the conormal sheaf $\check N$ occurs in the exact
sequence of coherent sheaves on $X$
\begin{equation}\tag{5.1.1}
0\longrightarrow\check N
\longrightarrow\Omega^1_{\mathbb P^r/k}|_X
\longrightarrow\Omega^1_X
\longrightarrow0.
\end{equation}
On the other hand, writing differential forms on $\mathbb P^r$ in
homogeneous coordinates gives an exact sequence of sheaves on
$\mathbb P^r$
\begin{equation}\tag{5.1.2}
0\longrightarrow\Omega^1_{\mathbb P^r/k}
\longrightarrow\mathcal O_{\mathbb P^r}(-1)^{r+1}
\longrightarrow\mathcal O_{\mathbb P^r}
\longrightarrow0.
\end{equation}
Its restriction to $X$ is therefore the exact sequence
\begin{equation}\tag{5.1.3}
0\longrightarrow\Omega^1_{\mathbb P^r/k}|_X
\longrightarrow\mathcal O_X(-1)^{r+1}
\longrightarrow\mathcal O_X
\longrightarrow0.
\end{equation}
\footnote{The print leaves the last term as
$\mathcal O_{\mathbb P^r}$ after restricting the sequence to $X$;
the restricted term is $\mathcal O_X$.}
Combining (5.1.1) and (5.1.3) gives an inclusion
\begin{equation}\tag{5.1.4}
\check N\lhook\joinrel\longrightarrow\mathcal O_X(-1)^{r+1},
\end{equation}
which, by transposition, gives a surjection
\begin{equation}\tag{5.1.5}
\mathcal O_X(1)^{r+1}\longrightarrow N.
\end{equation}
Hence there is a closed immersion
\begin{equation}\tag{5.1.6}
\alpha:\mathbb P(N)
\longrightarrow\mathbb P\bigl(\mathcal O_X(1)^{r+1}\bigr)
\simeq X\times\check{\mathbb P}^{r},
\end{equation}
which is precisely (5.1.0).

\end{filecontents*}
% Original input: components/104_expose_XVII_section_5_2.tex
\begin{filecontents*}[overwrite,noheader]{sga7ii-cumulative-input-104.tex}
\subsubsection*{5.2. The idea of the calculation.}

From now on suppose that
$\check X=\varphi(\mathbb P(N))$ has dimension $r-1$, so that it is a
hypersurface in $\check{\mathbb P}^{r}$.  Its degree is naturally
calculated by means of Chern classes [1].

For every smooth projective variety $T$, let $\operatorname{CH}(T)$
denote its Chow ring, graded by codimension.  For
$z\in\operatorname{CH}(T)$, write
\[
\int_T z
\]
for the degree of its zero-dimensional component.  The class in
$\operatorname{CH}^1(T)$ of an invertible sheaf $\mathcal L$ on $T$
will also be denoted by $\mathcal L$ whenever no confusion can arise
(for example, under the integral sign).

Let $L$ be the class of a line in
$\operatorname{CH}(\check{\mathbb P}^{r})$.  Then
\begin{equation}\tag{5.2.1}
\deg(\check X)
=\int_{\check{\mathbb P}^{r}}L\cdot\check X.
\end{equation}
Set-theoretically, $\check X$ is the image of $\mathbb P(N)$ under
$\varphi$ (3.1.2).  From the viewpoint of the proper pushforward
\[
\varphi_*:\operatorname{CH}(\mathbb P(N))
\longrightarrow\operatorname{CH}(\check{\mathbb P}^{r}),
\]
however, one has
\[
\varphi_*[\mathbb P(N)]=\deg(\varphi)[\check X].
\]
Consequently, (5.2.1) becomes
\begin{equation}\tag{5.2.2}
\deg(\check X)
=\frac{1}{\deg(\varphi)}
\int_{\check{\mathbb P}^{r}}L\cdot\varphi_*[\mathbb P(N)].
\end{equation}
Here, by abuse of notation, $\varphi$ in $\deg(\varphi)$ denotes the
morphism $\mathbb P(N)\to\check X$ induced by (3.1.2); its degree is a
nonnegative integer, taken to be zero precisely when
$\dim\check X<r-1=\dim\mathbb P(N)$.
\footnote{The print has $\dim X<r-1$.  The condition for the induced
map to the dual variety to be generically finite is
$\dim\check X=r-1$.}
The projection formula rewrites (5.2.2) as
\begin{equation}\tag{5.2.3}
\deg(\check X)
=\frac{1}{\deg(\varphi)}
\int_{\mathbb P(N)}\varphi^*(L).
\end{equation}

Let $\eta\in\operatorname{CH}(\check{\mathbb P}^{r})$ be the class of
a hyperplane.  Since $L=\eta^{r-1}$, equation (5.2.3) becomes
\begin{equation}\tag{5.2.4}
\deg(\check X)
=\frac{1}{\deg(\varphi)}
\int_{\mathbb P(N)}\bigl(\varphi^*(\eta)\bigr)^{r-1}.
\end{equation}
We now calculate $\varphi^*(\eta)$ concretely.

Let $\xi$ (respectively $\tau$) denote the canonical invertible sheaf
$\mathcal O_{\mathbb P(N)}(1)$ on $\mathbb P(N)$ (respectively
$\mathcal O_{\mathbb P(\mathcal O_X(1)^{r+1})}(1)$ on
$\mathbb P(\mathcal O_X(1)^{r+1})$).  Functoriality gives, by
(5.1.6),
\begin{equation}\tag{5.2.5}
\alpha^*(\tau)=\xi.
\end{equation}
Under the identification
\[
\mathbb P\bigl(\mathcal O_X(1)^{r+1}\bigr)
\simeq X\times\check{\mathbb P}^{r},
\]
and writing $H=\mathcal O_X(1)$, one plainly has
\begin{equation}\tag{5.2.6}
\tau=\operatorname{pr}_1^*(H)\otimes\operatorname{pr}_2^*(\eta).
\end{equation}
Let $\pi:\mathbb P(N)\to X$ be the structural morphism.  Then
\[
\pi=\operatorname{pr}_1\circ\alpha,
\qquad
\varphi=\operatorname{pr}_2\circ\alpha.
\]
Applying $\alpha^*$ to (5.2.6) and using (5.2.5), we obtain
\begin{equation}\tag{5.2.6$'$}
\xi=\alpha^*(\tau)=\pi^*(H)\otimes\varphi^*(\eta).
\end{equation}
\footnote{The print numbers both this display and the preceding one
(5.2.6); the prime distinguishes the second occurrence.}
Thus, in $\operatorname{CH}(\mathbb P(N))$,
\begin{equation}\tag{5.2.7}
\varphi^*(\eta)=\xi-\pi^*(H),
\end{equation}
and (5.2.3) becomes
\begin{equation}\tag{5.2.8}
\deg(\check X)
=\frac{1}{\deg(\varphi)}
\int_{\mathbb P(N)}\bigl(\xi-\pi^*(H)\bigr)^{r-1}.
\end{equation}
The right-hand side of (5.2.8) is calculated by means of the
following lemma (cf. 5.3.3).
\end{filecontents*}
% Original input: components/105_expose_XVII_lemma_5_3_and_first_proof.tex
\begin{filecontents*}[overwrite,noheader]{sga7ii-cumulative-input-105.tex}
\subsubsection*{Lemma 5.3.}
\underline{Let $X$ be a smooth connected quasi-projective variety of
dimension $n$, and let $M$ be a locally free sheaf of finite rank $a$
on $X$.  Let $\xi$ denote the invertible sheaf
$\mathcal O_{\mathbb P(M)}(1)$ on $\mathbb P(M)$, and let
$\pi:\mathbb P(M)\to X$ be the structural morphism.  Recall that
$\operatorname{CH}(\mathbb P(M))$ is a free
$\operatorname{CH}(X)$-module (via $\pi^*$), with basis
$1,\xi,\ldots,\xi^{a-1}$, where $\xi$ satisfies}
\begin{equation}\tag{5.3.1}
\xi^a-C_1\xi^{a-1}+C_2\xi^{a-2}-C_3\xi^{a-3}+\cdots=0,
\end{equation}
\underline{the $C_i$ being the Chern classes of $M$ in
$\operatorname{CH}(X)$.  Put
$C(M)=1+C_1+C_2+\cdots$, the total Chern class of $M$; it is an
invertible element of $\operatorname{CH}(X)$.  Then:}

\begin{enumerate}
\item[(i)] \underline{one has}
\begin{equation}\tag{5.3.2}
\pi_*\!\left(\frac{\xi^{a-1}}{1+\xi}\right)
=\frac{1}{C(M)};
\end{equation}
\item[(ii)] \underline{let
$F(T)\in\operatorname{CH}(X)[T]$ be a polynomial such that the
coefficient of $T^i$ is homogeneous of degree $n+a-1-i$.  Thus
$F(\xi)$ has maximal degree
$n+a-1=\dim\mathbb P(M)$ in
$\operatorname{CH}(\mathbb P(M))$.  Then}
\begin{equation}\tag{5.3.3}
\int_{\mathbb P(M)}F(\xi)
=(-1)^{a-1}\int_X\frac{F(-1)}{C(M)}.
\end{equation}
\end{enumerate}

\noindent\underline{Proof.}
We first deduce (5.3.3) from (5.3.2).  Write
\begin{equation}\tag{5.3.4}
\frac{1}{C(M)}=\sum_{j\geq0}b_j,\qquad
b_0=1,\qquad
b_j\in\operatorname{CH}^j(X),\qquad
b_j=0\ \text{for }j<0.
\end{equation}
Then (5.3.2) is equivalent to
\begin{equation}\tag{5.3.5}
\pi_*(\xi^{a-1+j})=(-1)^jb_j
\qquad\text{whenever }a-1+j\geq0,
\end{equation}
since $\pi_*(\xi^q)=0$ for $q<a-1$ for dimensional reasons.

By linearity it suffices to prove (5.3.3) when
$F(T)=uT^{a-1+j}$, with $u\in\operatorname{CH}^{n-j}(X)$
homogeneous.  By (5.3.5) and the projection formula,
\begin{equation}\tag{5.3.6}
\begin{aligned}
\int_{\mathbb P(M)}F(\xi)
&=\int_{\mathbb P(M)}\pi^*(u)\xi^{a-1+j}\\
&=\int_Xu\,\pi_*(\xi^{a-1+j})
=(-1)^j\int_Xu\,b_j.
\end{aligned}
\end{equation}
On the other hand,
\begin{equation}\tag{5.3.7}
\begin{aligned}
(-1)^{a-1}\int_X\frac{F(-1)}{C(M)}
&=(-1)^{a-1}\int_X
  \frac{(-1)^{a-1+j}u}{C(M)}\\
&=(-1)^j\int_X\sum_i u\,b_i.
\end{aligned}
\end{equation}
Because $\deg u=n-j$ and $\deg b_i=i$, one has
\[
\int_Xu\,b_i=0\qquad\text{for }i\ne j.
\]
Thus (5.3.7) reduces to
\begin{equation}\tag{5.3.8}
(-1)^{a-1}\int_X\frac{F(-1)}{C(M)}
=(-1)^j\int_Xu\,b_j,
\end{equation}
which proves the asserted implication.

It remains to prove (5.3.2), or equivalently (5.3.5).

\end{filecontents*}
% Original input: components/106_expose_XVII_completion_5_3_and_lemma_5_4.tex
\begin{filecontents*}[overwrite,noheader]{sga7ii-cumulative-input-106.tex}
For dimensional reasons,
\[
\pi_*(\xi^j)=0\qquad\text{if }j<a-1,
\]
whereas
\[
\pi_*(\xi^{a-1})=1_X
\]
for the geometric reason that $\xi^{a-1}$ meets each fiber in one
point.  The projection formula therefore gives
\begin{equation}\tag{5.3.9}
\pi_*\!\left(
\sum_{i=0}^{a-1}\pi^*(U_i)\xi^i
\right)=U_{a-1}.
\end{equation}
\footnote{The printed display mixes the indices $i,j$ and gives
$U_{j-i}$ on the right.  The two immediately preceding pushforward
identities force the coefficient $U_{a-1}$.}
By (5.3.9), assertion (5.3.5) follows from the next lemma, in which
$\xi$ is replaced by $T$ and $C_i$ by $S_i$, in view of (5.3.1).

\subsubsection*{Lemma 5.4.}
\underline{Let $a>0$ be an integer and let
$A=\mathbb Z[[X_1,\ldots,X_a]]$ be the ring of formal power series in
$a$ indeterminates.  Let $B$ be the $A$-algebra}
\[
B=A[[T]]\bigg/\left(\prod_{i=1}^{a}(T-X_i)\right).
\]
\underline{Thus $B$ is a free $A$-module with basis
$1,T,\ldots,T^{a-1}$:}
\begin{equation}\tag{5.4.1}
B=AT^{a-1}\oplus AT^{a-2}\oplus\cdots\oplus A.
\end{equation}
\underline{Let $S_i\in A$ be the $i$-th elementary symmetric
function}
\[
S_i=\sum_{1\leq\alpha_1<\cdots<\alpha_i\leq a}
X_{\alpha_1}\cdots X_{\alpha_i},
\]
\underline{so that}
\begin{equation}\tag{5.4.2}
\sum_{i=0}^{a}(-1)^iS_iT^{a-i}=0
\qquad\text{in }B.
\end{equation}
\underline{The element $1+S_1+\cdots+S_a$ is invertible in $A$;
write its inverse as}
\begin{equation}\tag{5.4.3}
\frac{1}{\sum_{i=0}^{a}S_i}
=\sum_{j\geq0}b_j,
\qquad
b_j\text{ homogeneous of degree }j\text{ in }A,
\qquad b_0=1.
\end{equation}
\underline{For every integer $j\geq0$, define a homogeneous element
$a_j\in A$ of degree $j$ by writing, using (5.4.1),}
\begin{equation}\tag{5.4.4}
T^{a-1+j}
=a_jT^{a-1}+{}?\,T^{a-2}+{}?\,T^{a-3}+\cdots+{}?,
\qquad ?\in A.
\end{equation}
\underline{Then, for every $j\geq0$,}
\begin{equation}\tag{5.4.5}
a_j=(-1)^jb_j.
\end{equation}

\end{filecontents*}
% Original input: components/107_expose_XVII_proof_5_4_through_proposition_5_7_2.tex
\begin{filecontents*}[overwrite,noheader]{sga7ii-cumulative-input-107.tex}
\noindent\underline{Proof.}
In view of the definition (5.4.3) of the $b_j$, it is enough to prove
\begin{equation}\tag{5.4.6}
\left\{
\begin{aligned}
a_0&=1,\\
\sum_{i=0}^{a}(-1)^iS_i\,a_{j-i}&=0
&&\text{if }j>0.
\end{aligned}
\right.
\end{equation}
Clearly $a_0=1$.  For a fixed $j>0$, the second relation is obtained
by taking the coefficient of $T^{a-1}$, using (5.4.1), in the
relation
\begin{equation}\tag{5.4.7}
\sum_{i=0}^{a}(-1)^iS_iT^{a-1+j-i}=0
\qquad\text{in }B.
\end{equation}
\footnote{The print omits the upper summation bound $a$ in (5.4.7);
the relation is (5.4.2) multiplied by $T^{j-1}$.}
This relation is obtained by multiplying (5.4.2) by $T^{j-1}$,
which proves the lemma.

\subsubsection*{Corollary 5.5.}
\underline{One has}
\begin{equation}\tag{5.5.1}
\deg(\check X)
=\frac{(-1)^n}{\deg(\varphi)}
\int_X\frac{(1+H)^{r-1}}{C(N)}.
\end{equation}

\noindent\underline{Proof.}
Apply (5.3.3), with $M=N$ the normal sheaf of
$X\hookrightarrow\mathbb P^r$, to (5.2.8).

For a smooth quasi-projective variety $S$, write
\begin{equation}\tag{5.5.2}
C(S)\stackrel{\mathrm{def}}{=}C(T_S)
\end{equation}
for its total Chern class, where $T_S$ denotes the tangent bundle of
$S$.

Dualizing the exact sequence (5.1.2) gives
\begin{equation}\tag{5.5.3}
C(\mathbb P^r)=(1+H)^{r+1}
\qquad\text{in }\operatorname{CH}(\mathbb P^r).
\end{equation}
Dualizing the exact sequence (5.1.1) gives
\begin{equation}\tag{5.5.4}
C(N)\,C(X)=C(\mathbb P^r)|_X
\qquad\text{in }\operatorname{CH}(X).
\end{equation}
Combining (5.5.3) and (5.5.4), formula (5.5.1) becomes:

\subsubsection*{Corollary 5.6.}
\underline{One has}
\begin{equation}\tag{5.6.1}
\deg(\check X)
=\frac{(-1)^n}{\deg(\varphi)}
\int_X\frac{C(X)}{(1+H)^2}.
\end{equation}
\footnote{The print has $(-1)^r$.  Substitution of (5.5.3) and
(5.5.4) into (5.5.1) does not alter its sign $(-1)^n$; formula
(5.7.3) immediately below confirms the same exponent.}

\subsubsection*{5.7.}
For a smooth projective variety $S$, define its
\underline{Euler--Poincaré characteristic} by
\begin{equation}\tag{5.7.1}
\chi(S)=\int_SC(S).
\end{equation}

\subsubsection*{Proposition 5.7.2.}
\underline{Let $Y$ be a smooth hyperplane section of $X$, and let
$\Delta$ be a smooth section of $X$ by a codimension-two linear
subspace of $\mathbb P^r$.  Then}
\begin{equation}\tag{5.7.3}
\deg(\check X)
=\frac{(-1)^n}{\deg(\varphi)}
\bigl[\chi(X)+\chi(\Delta)-2\chi(Y)\bigr].
\end{equation}

\end{filecontents*}
% Original input: components/108_expose_XVII_proof_5_7_2_section_6_and_references.tex
\begin{filecontents*}[overwrite,noheader]{sga7ii-cumulative-input-108.tex}
\noindent\underline{Proof.}
Let $i:Y\hookrightarrow X$ be the inclusion.  One has
\[
N_{Y/X}=H|_Y=i^*H.
\]
The exact sequence
\[
0\longrightarrow T_Y
\longrightarrow i^*T_X
\longrightarrow i^*H
\longrightarrow0
\]
therefore gives
\begin{equation}\tag{5.7.4}
C(Y)=i^*\!\left(\frac{C(X)}{1+H}\right).
\end{equation}
Hence
\[
\chi(Y)=\int_Yi^*\!\left(\frac{C(X)}{1+H}\right),
\]
and the projection formula gives
\begin{equation}\tag{5.7.5}
\chi(Y)
=\int_Xi_*i^*\!\left(\frac{C(X)}{1+H}\right)
=\int_X\frac{H\,C(X)}{1+H}.
\end{equation}
Iterating this calculation for $\Delta$, one obtains
\begin{equation}\tag{5.7.6}
\begin{aligned}
\chi(\Delta)
&=\int_Y\frac{H\,C(Y)}{1+H}\\
&=\int_Yi^*\!\left(
  \frac{H}{1+H}\cdot\frac{C(X)}{1+H}
  \right)
=\int_X\frac{H^2C(X)}{(1+H)^2}.
\end{aligned}
\end{equation}
It follows that
\[
\begin{aligned}
\chi(X)-2\chi(Y)+\chi(\Delta)
&=\int_XC(X)\left[
1-\frac{2H}{1+H}+\frac{H^2}{(1+H)^2}
\right]\\
&=\int_X\frac{C(X)}{(1+H)^2}.
\end{aligned}
\]
Together with (5.6.1), this proves the proposition.

\section*{6. The case of a general base}

\subsubsection*{6.1.}
Let $S$ be a scheme and let $\bar s$ be a geometric point of $S$,
defined by an algebraic closure of the residue field $k(s)$ for a
point $s\in S$.  Let $\mathbb P^r_S$ be projective $r$-space over
$S$, let $\check{\mathbb P}^{r}_S$ be its dual projective space, and
let
\[
I\subset\mathbb P^r_S\times_S\check{\mathbb P}^{r}_S
\]
be the incidence variety.  Let
$D\subset\check{\mathbb P}^{r}_S$ be a line, defined by an
$S$-valued point of the appropriate Grassmannian, and let $\Delta$ be
the axis of the corresponding pencil of hyperplanes (cf. 2.1).

Let $X$ be a proper smooth $S$-scheme, pure of dimension $n\geq1$,
equipped with a projective embedding $X\hookrightarrow\mathbb P^r_S$.
Put
\[
\widetilde X=(X\times_SD)\cap I.
\]
Thus $\widetilde X$ is the scheme over $D$ whose fibers are the
hyperplane sections of $X$ parametrized by $D$; let
$u:\widetilde X\to D$ be the projection.

We use a subscript $\bar s$ to denote base change by
$\bar s\to S$.  Suppose that $\Delta_{\bar s}$ is transverse to
$X_{\bar s}$.  Then, over a neighborhood $U$ of $s$ in $S$,
$\Delta$ is transverse to $X$.  In what follows we replace $S$ by
this neighborhood and write $U=S$.  The scheme $\widetilde X$ is then
smooth over $S$, and a neighborhood in $\widetilde X$ of the inverse
image of $\Delta$ is smooth over $D$.

Suppose moreover that one of the hyperplane sections of $X_{\bar s}$
belonging to the pencil $D_{\bar s}$ is smooth.  After again replacing
$S$ by a suitable neighborhood of $s$, the closed subscheme
$F\subset\widetilde X$ defined by the Jacobian ideal
$J^{n-1}(\widetilde X/D)$ is finite and flat over $S$.

\subsubsection*{6.2.}
The \underline{exceptional locus} of the pencil $D$ is the relative
Cartier divisor $E$ of $D/S$ defined by the finite flat
$\mathcal O_D$-module $u_*\mathcal O_F$.  Locally on $D$, if
\[
\mathcal O_D^m\xrightarrow{\ v\ }\mathcal O_D^m
\longrightarrow u_*\mathcal O_F\longrightarrow0
\]
is a presentation of $u_*\mathcal O_F$, then $E$ has equation
$\det(v)=0$.  Formation of $E$ commutes with arbitrary base change.

The divisor $E_{\bar s}$ on $D_{\bar s}$ can be written
\[
E_{\bar s}=\sum_P n(P)[P],
\]
where $n(P)$ is the sum of the colengths of the Jacobian ideal
$J^{n-1}(\widetilde X/D)$ at the critical points of $u$ mapping to
$P$.

\subsubsection*{Proposition 6.3.}
\underline{Suppose that $D_{\bar s}$ is a Lefschetz pencil of
hyperplane sections of $X_{\bar s}$ and that
$\operatorname{char}k(s)\ne2$ or $n$ is even.  There is a
neighborhood $U$ of $s$ in $S$ such that:}
\begin{enumerate}
\item[(a)] \underline{for every geometric point $\bar t$ of $U$,
$D_{\bar t}$ is a Lefschetz pencil of hyperplane sections of
$X_{\bar t}$;}
\item[(b)] \underline{over $U$, the exceptional locus $E$ is étale
over $S$.}
\end{enumerate}
To prove (a) and (b), it is enough to observe that, under these
hypotheses, $D_{\bar s}$ satisfies condition 2.2(c) if and only if
the divisor $E_{\bar s}$ is reduced.  Invoking 3.3, one could use
only the easier implication from left to right.

\subsubsection*{6.3 (exceptional case).}
\footnote{The print repeats the number 6.3 after Proposition 6.3.}
The proposition is false in the exceptional case just mentioned.
Nevertheless, if $S$ has characteristic two and $n$ is odd, there is
a neighborhood $U$ of $s$ over which (a) holds, and one can define a
divisor $E'$ étale over $U$ such that $E=2E'$.  We shall not use this
result.

\subsubsection*{References}

\noindent[1] A. Grothendieck, \emph{La théorie des classes de Chern},
Bull. Soc. Math. France \textbf{86} (1958), 137--154.

\end{filecontents*}
% Original input: components/109_expose_XVIII_title_introduction_and_section_1.tex
\begin{filecontents*}[overwrite,noheader]{sga7ii-cumulative-input-109.tex}
\clearpage
\begin{center}
{\large SGA 7 --- EXPOS\'E XVIII}\par
\vspace{2em}
{\Large\underline{Cohomological Study of Lefschetz Pencils}}\par
\vspace{1.5em}
{\large by N. Katz\footnote{Based on succinct notes of Grothendieck.}}
\end{center}

\section*{Contents}

\noindent 1. The cohomology of a projective bundle. \hfill 2

\noindent 2. The cohomology of blown-up varieties. \hfill 4

\noindent 3. The blow-up associated with a pencil. \hfill 13

\noindent 4. The cohomology of the blow-up associated with a pencil. \hfill 17

\noindent 5. Cohomological study of the morphism
$\rho:\widetilde X\longrightarrow\mathbb P^1$ associated with a pencil.
\hfill 24

\noindent\quad 5.1. The restriction homomorphism
$H^*(\widetilde X)\longrightarrow H^*(Y)$. \hfill 24

\noindent\quad 5.2. The ``Lefschetz vache theorem'' (LV). \hfill 27

\noindent\quad 5.3. Condition (A). \hfill 33

\noindent\quad 5.8. The Griffiths homomorphism. \hfill 47

\noindent\quad 5.9. An example in which hypothesis (A) is not satisfied.
\hfill 57

\noindent 6. Applications of the local theory of vanishing cycles
(Picard--Lefschetz). \hfill 58

\section*{Introduction}

In the present expos\'e we give the elementary cohomological structure of
Lefschetz pencils. Sections 1--4 give the ``well-known'' facts concerning the
cohomological structure of projective bundles and blow-ups. In Section 5 we
carry out the study named in the title, assuming that $X$ satisfies the
``Lefschetz vache theorem'' (5.2) and that the pencil $D$ in question satisfies
a certain condition (A), local in nature on $D$ (5.3). The principal results
are 5.6, 5.7, and 5.8.7; the last plays an essential role in the proof of
Griffiths's theorem in Expos\'e XX. In the final section, which is distinctly
less elementary, we use Mme Raynaud's Expos\'e XIII of SGA 1 on the
\underline{tame fundamental group}, and the \underline{Picard--Lefschetz
formula} (XVI), to prove:
\begin{enumerate}[label=(\arabic*)]
\item condition (A) is often satisfied (6.3, 6.4);
\item except when $\dim(X)$ is odd and the characteristic is $2$, the
monodromy group of a Lefschetz pencil acts irreducibly on the ``vanishing
cohomology,'' and transitively (up to sign) on the set of ``local vanishing
cycles'' associated, via the Picard--Lefschetz formula, with the various
``critical points'' of the pencil (6.6 and 6.7).
\end{enumerate}

\section*{1. The cohomology of a projective bundle}

\subsubsection*{Theorem 1.1 (Deligne).}
\underline{Let $f:X\longrightarrow S$ be a proper morphism, let $\ell$ be a
prime number invertible on $S$, let
$u\in H^2(X,\mathbb Z_\ell(1))$ (respectively,
$u\in H^2(X,\mathbb Q_\ell(1))$) be a cohomology class on $X$, and let
$\mathfrak F$ be a $\mathbb Z_\ell$-sheaf (respectively, a
$\mathbb Q_\ell$-sheaf) on $X$. Suppose that there is an integer $n$ such
that, for every integer $j$ and every integer $i>0$, the map}
\begin{equation}\tag{1.1.1}
R^{n-i}f_*\bigl(\mathfrak F(j)\bigr)
\xrightarrow{\ \wedge u^i\ }
R^{n+i}f_*\bigl(\mathfrak F(j+i)\bigr)
\end{equation}
\underline{is an isomorphism of $\mathbb Z_\ell$-sheaves (respectively,
$\mathbb Q_\ell$-sheaves) on $S$. Then, for every morphism
$g:S\longrightarrow T$, the Leray spectral sequence}
\[
E_2^{p,q}=R^pg_*R^qf_*(\mathfrak F)
\Longrightarrow R^{p+q}(gf)_*(\mathfrak F)
\]
\underline{degenerates.}

The proof is found in [1, Prop.~2.4]. Notice that, by the proper base-change
theorem (SGA 4 XII, 5.2), the hypothesis can be checked fibre by fibre.

\subsubsection*{Theorem 1.2 (cf. SGA 5 VII, 2.2.1).}
\underline{Let $\mathcal E\longrightarrow X$ be a vector bundle, locally
free of rank $1+r$. Let $f:\mathbb P(\mathcal E)\longrightarrow X$ denote
the projective bundle defined by $\mathcal E$, and let
$u\in H^2(\mathbb P(\mathcal E),\mathbb Z_\ell(1))$ be the first Chern
class of the canonical invertible sheaf
$\mathcal O_{\mathbb P(\mathcal E)}(1)$. For $\ell$ a prime number
invertible on $X$, one has:}
\begin{enumerate}[label=\alph*)]
\item \underline{the spectral sequence}
\begin{equation}\tag{1.2.1}
E_2^{p,q}=H^p\bigl(X,R^qf_*\mathbb Z_\ell\bigr)
\Longrightarrow
H^{p+q}\bigl(\mathbb P(\mathcal E),\mathbb Z_\ell\bigr)
\end{equation}
\underline{degenerates;}
\item \underline{the map}
\begin{equation}\tag{1.2.2}
\bigoplus_{0\leq j\leq r}
H^{q-2j}\bigl(X,\mathbb Z_\ell(-j)\bigr)
\longrightarrow
H^q\bigl(\mathbb P(\mathcal E),\mathbb Z_\ell\bigr),
\qquad
(\tau_{q-2j})_j\longmapsto
\sum_j f^*(\tau_{q-2j})\wedge u^j
\end{equation}
\underline{is an isomorphism.}
\end{enumerate}
\footnote{In (1.2.2) the print has the upper bound $n$, the degree
$q-2-j$, and one fewer closing parenthesis. The rank $1+r$, the defining
map, and the proof immediately below require $0\leq j\leq r$ and $q-2j$;
the corrected formula is printed here.}

\noindent\underline{Proof.}
The fibres of $f$ are projective spaces of dimension $r$, and one knows that
\begin{equation}\tag{1.2.3}
R^if_*\mathbb Z_\ell
\xleftarrow[\ \sim\ ]{\ \wedge^{i/2}u\ }
\begin{cases}
\mathbb Z_\ell\left(-\dfrac{i}{2}\right),
  & \text{if $i$ is even and $0\leq i\leq2r$},\\[2ex]
0, & \text{otherwise}.
\end{cases}
\end{equation}
Theorem 1.1 applies to $(f,\ell,u,\mathbb Z_\ell,r)$, which proves (a).
By (1.2.3),
\begin{equation}\tag{1.2.4}
E_2^{a,b}=
\begin{cases}
H^a\left(X,\mathbb Z_\ell\left(-\dfrac b2\right)\right),
  & \text{if $b$ is even and $0\leq b\leq2r$},\\[2ex]
0, & \text{otherwise}.
\end{cases}
\end{equation}
Thus the graded object associated with
$H^q(\mathbb P(\mathcal E),\mathbb Z_\ell)$ for the filtration coming
from the spectral sequence is
\[
\bigoplus_{0\leq j\leq r}
H^{q-2j}\bigl(X,\mathbb Z_\ell(-j)\bigr).
\]
If this direct sum is equipped with the filtration
\[
F^i=
\bigoplus_{\substack{0\leq j\leq r\\q-2j\leq i}}
H^{q-2j}\bigl(X,\mathbb Z_\ell(-j)\bigr),
\]
the map (1.2.2) is compatible with the filtrations and induces the identity on
the associated graded objects. This proves (b).

\end{filecontents*}
% Original input: components/110_expose_XVIII_section_2_1.tex
\begin{filecontents*}[overwrite,noheader]{sga7ii-cumulative-input-110.tex}
\section*{2. The cohomology of blown-up varieties (cf. SGA 5 VII, 8)}

\subsubsection*{2.1.}
Let $X$ be a smooth scheme over a field $k$, and let
$\Delta\subset X$ be an everywhere smooth subscheme of codimension $1+r$.
It is known (EGA IV, 19.4.9) that there is a \underline{smooth} scheme
$\widetilde X$, equipped with a projective morphism
\begin{equation}\tag{2.1.1}
f:\widetilde X\longrightarrow X,
\end{equation}
such that, on setting
\begin{equation}\tag{2.1.2}
\widetilde\Delta=f^{-1}(\Delta)
=\Delta\times_X\widetilde X,
\end{equation}
$f$ induces an isomorphism
\begin{equation}\tag{2.1.3}
f|_{\widetilde X-\widetilde\Delta}:
\widetilde X-\widetilde\Delta
\xrightarrow{\ \sim\ }X-\Delta,
\end{equation}
and there is a $\Delta$-isomorphism between $\widetilde\Delta$ and
$\check{\mathbb P}(\underline N_{\Delta/X})$, the projective bundle of
lines in the normal bundle $\underline N_{\Delta/X}$ of $\Delta$ in $X$.
Thus one has a commutative diagram
\begin{equation}\tag{2.1.4}
\begin{tikzcd}[row sep=huge,column sep=small]
\widetilde\Delta
  \arrow[rd,"f|_{\widetilde\Delta}"']
  & \simeq &
\check{\mathbb P}(\underline N_{\Delta|X})
  \arrow[ld,"\text{canonical projection}"]\\
{} & \Delta & {}
\end{tikzcd}.
\end{equation}

\subsubsection*{2.1.5.}
In the special case
\[
X=\operatorname{Spec}(A),\qquad
\Delta=\operatorname{Spec}(A/I),\qquad
I=(g_0,\ldots,g_r),
\]
$\widetilde X$ is obtained by the Hopf construction (cf. EGA IV, 19.4.11).
One defines a morphism from $X-\Delta$ to $\mathbb P^r$ by
\[
x\in X-\Delta\longmapsto
\bigl(g_0(x),\ldots,g_r(x)\bigr)\in\mathbb P^r
\]
and regards its graph $\Gamma$ as embedded in $X\times\mathbb P^r$. Then
$\widetilde X$ is defined to be the schematic closure of $\Gamma$ in
$X\times\mathbb P^r$, and the morphism
$f:\widetilde X\longrightarrow X$ is the composite
\[
\widetilde X\lhook\joinrel\longrightarrow
X\times\mathbb P^r\xrightarrow{\ \operatorname{pr}_1\ }X.
\]

\subsubsection*{2.1.6.}
Similarly (cf. again EGA IV, 19.4.11), if $X$ is projective and $\Delta$ is
the intersection of $X$ with $1+r$ hypersurfaces, say with equations
$G_i=0$, $i=0,\ldots,r$, one again considers the morphism
\[
X-\Delta\longrightarrow\mathbb P^r,\qquad
x\longmapsto\bigl(G_0(x),\ldots,G_r(x)\bigr),
\]
and obtains $\widetilde X$ as the closure of its graph in
$X\times\mathbb P^r$. Concretely, in terms of homogeneous coordinates
$\lambda_0,\ldots,\lambda_r$ on $\mathbb P^r$, $\widetilde X$ is the
subscheme of $X\times\mathbb P^r$ defined by the equations
\begin{equation}\tag{2.1.6.1}
\lambda_iG_j(X)-\lambda_jG_i(X),
\qquad 0\leq i<j\leq r.
\end{equation}
These equations show that the point $(x,\lambda)$ lies in
$\widetilde\Delta$ if and only if $x\in\Delta$; in other words,
$\check{\mathbb P}(\underline N_{\Delta/X})$ is the trivial projective
bundle\footnote{The print labels the right-hand projection in (2.1.6.2)
$\operatorname{pr}_i$. Its target is $\Delta$, and the same projection is
labelled $\operatorname{pr}_1$ in (2.1.6.3); the corrected label is used
here.}
\begin{equation}\tag{2.1.6.2}
\begin{tikzcd}[column sep=large,row sep=large]
\check{\mathbb P}(\underline N_{\Delta/X})
  \arrow[rr,phantom,"\simeq"]
  \arrow[dr,"\text{canonical projection}"']
  & {} &
\Delta\times\mathbb P^r
  \arrow[dl,"\operatorname{pr}_1"]\\
{} & \Delta & {}
\end{tikzcd}.
\end{equation}
One has the following commutative diagram:
\begin{equation}\tag{2.1.6.3}
\begin{tikzcd}[column sep=large,row sep=large]
\Delta\times\mathbb P^r
  \arrow[r,"\sim"]
  \arrow[dr,"\operatorname{pr}_1"]
&
\widetilde\Delta
  \arrow[r,hook]
  \arrow[d,"f|_{\widetilde\Delta}"]
&
\widetilde X
  \arrow[r,hook]
  \arrow[d,"f"]
&
X\times\mathbb P^r
  \arrow[dl,"\operatorname{pr}_1"]
\\
{} & \Delta \arrow[r,hook] & X & {}
\end{tikzcd},
\end{equation}
and the composite inclusion
$\Delta\times\mathbb P^r\hookrightarrow X\times\mathbb P^r$ is obtained
from $\Delta\hookrightarrow X$ by base change.
\end{filecontents*}
% Original input: components/111_expose_XVIII_theorem_2_2_through_part_1.tex
\begin{filecontents*}[overwrite,noheader]{sga7ii-cumulative-input-111.tex}
\subsubsection*{Theorem 2.2.}
\underline{Let $X$ be smooth over a field $k$, let $\Delta$ be an everywhere
smooth subscheme of codimension $1+r$, and let
$f:\widetilde X\longrightarrow X$ be the blow-up of $X$ along $\Delta$.
Then:}

\medskip
\noindent(2.2.1) \underline{The Leray spectral sequence}
\[
E_2^{p,q}=H^p\bigl(X,R^qf_*\mathbb Z_\ell\bigr)
\Longrightarrow
H^{p+q}\bigl(\widetilde X,\mathbb Z_\ell\bigr)
\]
\underline{degenerates.}

\medskip
\noindent(2.2.2) \underline{There is an additive isomorphism}
\[
\begin{aligned}
H^q(\widetilde X,\mathbb Z_\ell)
&\xleftarrow{\ \sim\ }
H^q(X,\mathbb Z_\ell)\oplus
\frac{H^q(\widetilde\Delta,\mathbb Z_\ell)}
     {\operatorname{Im}H^q(\Delta,\mathbb Z_\ell)}
\\
&\simeq
\left(\bigoplus_{j=1}^{r}
H^{q-2j}(\Delta,\mathbb Z_\ell(-j))\right)
\oplus H^q(X,\mathbb Z_\ell).
\end{aligned}
\]

\noindent\underline{Proof.}
Let
\begin{equation}\tag{2.2.3}
g:\widetilde\Delta\longrightarrow\Delta
\end{equation}
be the restriction of $f$ to $\widetilde\Delta$, and let
\begin{equation}\tag{2.2.4}
\left\{
\begin{aligned}
j&:\Delta\hookrightarrow X,\\
i&:\widetilde\Delta\hookrightarrow\widetilde X
\end{aligned}
\right.
\end{equation}
be the inclusions, so that one has the commutative diagram
\begin{equation}\tag{2.2.5}
\begin{tikzcd}[column sep=huge,row sep=large]
\widetilde X \arrow[d,"f"'] &
\widetilde\Delta \arrow[l,"i"'] \arrow[d,"g"]\\
X & \Delta \arrow[l,"j"']
\end{tikzcd}.
\end{equation}
It induces a homomorphism between the Leray spectral sequences for the
morphisms $f$ and $g$, with the sheaf $\mathbb Z_\ell$ on $\widetilde X$ and
$\widetilde\Delta$:
\begin{equation}\tag{2.2.6}
E_s^{p,q}(f)\xrightarrow{\ (i^*,j^*)\ }E_s^{p,q}(g).
\end{equation}
On the other hand, since $f$ induces an isomorphism from
$\widetilde X-\widetilde\Delta$ to $X-\Delta$, the proper base-change theorem
gives
\begin{equation}\tag{2.2.7}
R^qf_*\mathbb Z_\ell\simeq
\begin{cases}
\mathbb Z_\ell,&q=0,\\[1ex]
j_*R^qg_*\mathbb Z_\ell,&q>0.
\end{cases}
\end{equation}
Hence
\begin{equation}\tag{2.2.8}
E_2^{p,q}(f)
\xrightarrow[\ \sim\ ]{\ (i^*,j^*)\ }
E_2^{p,q}(g)
\qquad(q\ne0).
\end{equation}
We shall prove by induction on $s$ that
\begin{equation}\tag{2.2.9}
E_s^{p,q}(f)
\xrightarrow[\ \sim\ ]{\ (i^*,j^*)\ }
E_s^{p,q}(g)
\qquad(q\ne0,\ s\geq2)
\end{equation}
and
\begin{equation}\tag{2.2.10}
d_s^{p,q}(f)=0\qquad(s\geq2).
\end{equation}
For the moment, admit the formula
\begin{equation}\tag{2.2.11}
d_s^{p,s-1}(f)=0\qquad(s\geq2,\ \text{for every }p).
\end{equation}

Begin with $s=2$. We have a commutative diagram
\begin{equation}\tag{2.2.12}
\begin{tikzcd}[column sep=huge,row sep=large]
E_2^{p,q}(f)
  \arrow[r,"{(i^*,j^*)}"]
  \arrow[d,"{d_2^{p,q}(f)}"]
&
E_2^{p,q}(g)
  \arrow[d,"{0=d_2^{p,q}(g)}"]
\\
E_2^{p+2,q-1}(f)
  \arrow[r,"{(i^*,j^*)}"]
&
E_2^{p+2,q-1}(g)
\end{tikzcd}.
\end{equation}
For $q\leq0$, one has $E_2^{p+2,q-1}(f)=0$, hence
$d_2^{p,q}(f)=0$. For $q=1$, formula (2.2.11) gives
$d_2^{p,1}(f)=0$. For $q>1$, the horizontal arrows in (2.2.12) are
isomorphisms, while $d_2^{p,q}(g)=0$ by 1.2 and 2.1.4. This proves
(2.2.10) for $s=2$, and (2.2.9) for $s=2$ is exactly (2.2.8).

Suppose (2.2.9) and (2.2.10) are known for some $s\geq2$. Hypothesis
(2.2.10) implies
$E_s^{p,q}(f)\simeq E_{s+1}^{p,q}(f)$, and by 1.2 one has
$E_s^{p,q}(g)\simeq E_{s+1}^{p,q}(g)$, so there is a commutative diagram
\begin{equation}\tag{2.2.13}
\begin{tikzcd}[column sep=huge,row sep=large]
E_s^{p,q}(f)
  \arrow[r,"{(i^*,j^*)}"]
  \arrow[d,"\wr"']
&
E_s^{p,q}(g)
  \arrow[d,"\wr"]
\\
E_{s+1}^{p,q}(f)
  \arrow[r,"{(i^*,j^*)}"]
&
E_{s+1}^{p,q}(g)
\end{tikzcd},
\end{equation}
which proves (2.2.9) for $s+1$, in view of (2.2.9) for $s$. To deduce
(2.2.10) for $s+1$, consider the commutative diagram
\begin{equation}\tag{2.2.14}
\begin{tikzcd}[column sep=huge,row sep=large]
E_{s+1}^{p,q}(f)
  \arrow[r,"{(i^*,j^*)}"]
  \arrow[d,"{d_{s+1}^{p,q}(f)}"]
&
E_{s+1}^{p,q}(g)
  \arrow[d,"{0=d_{s+1}^{p,q}(g)}"]
\\
E_{s+1}^{p+s+1,q-s}(f)
  \arrow[r,"{(i^*,j^*)}"]
&
E_{s+1}^{p+s+1,q-s}(g)
\end{tikzcd}.
\end{equation}
For $q\leq s-1$, one has $E_{s+1}^{p+s+1,q-s}(f)=0$, hence
$d_{s+1}^{p,q}(f)=0$. For $q=s$, one has (2.2.11), and for $q>s$ the
horizontal arrows in (2.2.14) are isomorphisms and
$d_{s+1}^{p,q}(g)=0$ by 1.2. This gives the desired conclusion.
\footnote{At this point the print writes $d_s^{p,q}(f)=0$, although the
inductive step and diagram (2.2.14) concern $d_{s+1}^{p,q}(f)$.}

It remains to verify (2.2.11), which says that the composite edge
homomorphism
\begin{equation}\tag{2.2.15}
E_2^{p,0}(f)\longrightarrow E_3^{p,0}(f)
\longrightarrow E_4^{p,0}(f)\longrightarrow\cdots
\longrightarrow E_\infty^{p,0}(f)
\end{equation}
is an isomorphism; that is, that
\begin{equation}\tag{2.2.16}
f^*:H^p(X,\mathbb Z_\ell)
\longrightarrow H^p(\widetilde X,\mathbb Z_\ell)
\end{equation}
is injective. But it has a left inverse, namely the
\underline{Gysin homomorphism} (SGA 5 IV)
\begin{equation}\tag{2.2.17}
f_*:H^p(\widetilde X,\mathbb Z_\ell)
\longrightarrow H^p(X,\mathbb Z_\ell),
\end{equation}
which indeed satisfies
\begin{equation}\tag{2.2.18}
f_*f^*=\mathrm{id}.
\end{equation}
This follows from the projection formula and the fact that $f_*(1)=1$, since
$f$ is birational and therefore has degree one. This proves (2.2.1).

\end{filecontents*}
% Original input: components/112_expose_XVIII_theorem_2_2_completion_and_corollary_2_3.tex
\begin{filecontents*}[overwrite,noheader]{sga7ii-cumulative-input-112.tex}
As for (2.2.2), passing to cohomology in diagram (2.2.5) gives a
commutative diagram
\begin{equation}\tag{2.2.19}
\begin{tikzcd}[row sep=huge,column sep=huge]
H^p(\widetilde X,\mathbb Z_\ell)
  \arrow[r,"i^*"]
&
H^p(\widetilde\Delta,\mathbb Z_\ell)
\\
H^p(X,\mathbb Z_\ell)
  \arrow[r,"j^*"]
  \arrow[u,"f^*"']
&
H^p(\Delta,\mathbb Z_\ell)
  \arrow[u,"g^*"']
\end{tikzcd},
\end{equation}
hence a morphism
\begin{equation}\tag{2.2.20}
\frac{H^p(\widetilde X,\mathbb Z_\ell)}
     {f^*H^p(X,\mathbb Z_\ell)}
\xrightarrow{\ (i^*,j^*)\ }
\frac{H^p(\widetilde\Delta,\mathbb Z_\ell)}
     {g^*H^p(\Delta,\mathbb Z_\ell)}.
\end{equation}
Let $F^i$ denote the filtrations on
$H^\bullet(\widetilde X,\mathbb Z_\ell)$ and
$H^\bullet(\widetilde\Delta,\mathbb Z_\ell)$ coming from the Leray spectral
sequences. Then (2.2.20) can be rewritten as\footnote{In the right-hand
denominator of (2.2.21), the print reverses the notation to $H^pF^p$.
The filtration used in the preceding sentence and on the left is
$F^pH^p$; that notation is restored here.}
\begin{equation}\tag{2.2.21}
\frac{H^p(\widetilde X,\mathbb Z_\ell)}
     {F^pH^p(\widetilde X,\mathbb Z_\ell)}
\xrightarrow{\ (i^*,j^*)\ }
\frac{H^p(\widetilde\Delta,\mathbb Z_\ell)}
     {F^pH^p(\widetilde\Delta,\mathbb Z_\ell)}.
\end{equation}
By (2.2.9), the map induced by (2.2.21) on associated graded objects is an
isomorphism, so (2.2.20) is an isomorphism. On the other hand, (2.2.18) gives
a direct-sum decomposition
\begin{equation}\tag{2.2.22}
\begin{tikzcd}[column sep=7em]
H^\bullet(\widetilde X,\mathbb Z_\ell)
  \arrow[r,shift right=1.2ex,
    "f_*\oplus(1-f^*f_*)"']
&
H^\bullet(X,\mathbb Z_\ell)\oplus\Ker f_*
  \arrow[l,shift right=1.2ex,
    "f^*+\text{inclusion}"']
\end{tikzcd},
\end{equation}
and (1.2.2) gives the direct-sum decomposition
\begin{equation}\tag{2.2.23}
\frac{H^\bullet(\widetilde\Delta,\mathbb Z_\ell)}
     {g^*H^\bullet(\Delta,\mathbb Z_\ell)}
\xleftarrow{\ \sim\ }
\bigoplus_{1\leq j\leq r}
H^{\bullet-2j}(\Delta,\mathbb Z_\ell(-j)).
\end{equation}
Combining the isomorphism (2.2.20) and decompositions (2.2.22) and (2.2.23)
gives (2.2.2), as required.

\subsubsection*{Remark 2.2.24.}
Let us sketch a second, distinctly simpler proof of (2.2.2), valid when $X$
is \underline{proper} over $k$. More generally, without any hypothesis on
$X$, one obtains the corresponding result (2.2.26) in cohomology with proper
supports. The exact sequences of cohomology with proper supports, denoted
$H_c^\bullet$, for the pairs
$(\widetilde X,\widetilde\Delta)$ and $(X,\Delta)$ form the commutative
diagram
\begin{center}
\resizebox{\textwidth}{!}{%
\begin{tikzcd}[ampersand replacement=\&,column sep=small,row sep=large]
(2.2.25) \& {} \arrow[r] \&
H_c^i(\widetilde X-\widetilde\Delta,\mathbb Z_\ell)
  \arrow[r] \&
H_c^i(\widetilde X,\mathbb Z_\ell)
  \arrow[r] \&
H_c^i(\widetilde\Delta,\mathbb Z_\ell)
  \arrow[r] \&
H_c^{i+1}(\widetilde X-\widetilde\Delta,\mathbb Z_\ell)
\\
(2.2.26) \& {} \arrow[r] \&
H_c^i(X-\Delta,\mathbb Z_\ell)
  \arrow[r]
  \arrow[u,"\wr" left] \&
H_c^i(X,\mathbb Z_\ell)
  \arrow[r]
  \arrow[u,hook,"f^*" right] \&
H_c^i(\Delta,\mathbb Z_\ell)
  \arrow[r]
  \arrow[u,hook,"g^*" right] \&
H_c^{i+1}(X-\Delta,\mathbb Z_\ell)
  \arrow[u,"\wr" left]
\end{tikzcd}}
\end{center}
The quotient of the acyclic complex (2.2.25) by the acyclic subcomplex
(2.2.26) gives the \underline{acyclic} complex
\begin{equation}\tag{2.2.27}
\longrightarrow 0\longrightarrow
\frac{H_c^i(\widetilde X,\mathbb Z_\ell)}
     {f^*H_c^i(X,\mathbb Z_\ell)}
\longrightarrow
\frac{H_c^i(\widetilde\Delta,\mathbb Z_\ell)}
     {g^*H_c^i(\Delta,\mathbb Z_\ell)}
\longrightarrow0\longrightarrow.
\end{equation}
Thus, without considering spectral sequences, one finds that (2.2.20) is an
isomorphism. The arguments just given then apply and yield (2.2.2).

\subsubsection*{Corollary 2.3.}
\underline{Set}
\[
b_q(X)=\dim_{\mathbb Q_\ell}H^q(X,\mathbb Q_\ell),
\qquad
\chi(X)=\sum_q(-1)^qb_q(X).
\]
\underline{Then}\footnote{The print has $b_{q-j}(\Delta)$ in (2.3.1).
The decomposition (2.2.2) shifts cohomological degree by $2j$, so the
correct summand is $b_{q-2j}(\Delta)$.}
\begin{equation}\tag{2.3.1}
b_q(\widetilde X)
=b_q(X)+\sum_{j=1}^{r}b_{q-2j}(\Delta),
\end{equation}
\begin{equation}\tag{2.3.2}
\chi(\widetilde X)=\chi(X)+r\,\chi(\Delta).
\end{equation}
\end{filecontents*}
% Original input: components/113_expose_XVIII_section_3.tex
\begin{filecontents*}[overwrite,noheader]{sga7ii-cumulative-input-113.tex}
\section*{3. The blow-up associated with a pencil}

\subsection*{3.1.}
Let $X$ be projective, smooth, and irreducible over an algebraically closed
field $k$, of dimension $n$, and let
$D=\{H_\tau\}_{\tau\in\mathbb P^1}$ be a pencil of hypersurfaces whose axis
meets $X$ transversely, in $\Delta$.
We consider the blow-up of $X$ along $\Delta$, which we shall now describe
explicitly.\footnote{The print writes $u$ for the dimension in this sentence.
Section 3.2 immediately denotes the same dimension by $n$, which also
governs the signs in (3.2.3)--(3.2.4); that notation is used here.}

In terms of equations, let $(\lambda,\mu)$ be homogeneous coordinates on
$\mathbb P^1$, and let $F$ and $G$ be two linear forms such that $F=G=0$
defines the axis of the pencil.  The pencil is written
\begin{equation}\tag{3.1.1}
\lambda F-\mu G.
\end{equation}
By definition, the subscheme $\Delta$ is the intersection of $X$ with the
linear subvariety $G=F=0$.  By 2.1.6,
\begin{equation}\tag{3.1.2}
\widetilde X
=\{(x,(\lambda,\mu))\in X\times\mathbb P^1
       \mid \lambda F-\mu G=0\};
\end{equation}
it is the scheme-theoretic closure, in $X\times\mathbb P^1$, of the graph of
the map
\begin{equation}\tag{3.1.3}
\begin{tikzcd}[row sep=large,column sep=huge]
X-\Delta \arrow[r] & \mathbb P^1 \\
x \arrow[r] & (G(x),F(x)).
\end{tikzcd}
\end{equation}
The great advantage of $\widetilde X$ is that it is projective (being
embedded in $X\times\mathbb P^1$) and smooth (because $\Delta$ is smooth)
over $k$, and admits a morphism $\rho$ to $\mathbb P^1$, namely
\begin{equation}\tag{3.1.4}
\rho:\widetilde X\longrightarrow X\times\mathbb P^1
\xrightarrow{\operatorname{pr}_2}\mathbb P^1,
\end{equation}
whose fibers are the hyperplane sections of $X$ by the members of our
pencil:
\[
\rho^{-1}(t)=X_t\overset{\mathrm{def}}{=}X\cdot H_t.
\]

\subsection*{3.2. The degree of the dual variety: topological method}

Suppose in addition that the pencil of 3.1 is such that
\begin{equation}\tag{3.2.1}
X_t\text{ is smooth for every }t\text{ in a nonempty open subset of }
\mathbb P^1,
\end{equation}
and
\begin{equation}\tag{3.2.2}
\text{if $X_t$ is singular, its only singularities are nondegenerate
quadratic singularities (XVII 1.1).}
\end{equation}
These conditions hold in particular for a \underline{Lefschetz pencil}
(XVII 2.21) if $\operatorname{char}k\ne2$ or if $X$ has even dimension.

Let $\eta$ be the generic point of $\mathbb P^1$, let $\bar\eta$ be a
geometric point above it, let $n=\dim X$, and let $\chi$ be the
Euler--Poincar\'e characteristic with coefficients in $\mathbb Q_\ell$.
The Euler--Poincar\'e formula XVI 2.1 applied to $\rho$ gives
\begin{equation}\tag{3.2.3}
\chi(\widetilde X)
=\chi(\mathbb P^1)\,\chi(X_{\bar\eta})
 +(-1)^n\sum_{\substack{\text{singular points}\\
                         \text{of the fibers}}}1,
\end{equation}
which, in view of (2.2.2), can be rewritten
\begin{equation}\tag{3.2.4}
(-1)^n
\left[\begin{array}{c}
\text{total number of}\\[3pt]
\text{singular points}
\end{array}\right]
=\chi(X)+\chi(\Delta)-2\chi(X_{\bar\eta}).
\end{equation}
Notice that the quantity
$\chi(X)+\chi(\Delta)-2\chi(X_{\bar\eta})$ is independent of the particular
choice of pencil.  It depends only on the embedding of $X$ in some
$\mathbb P^r$, and indeed only on the Chern class of
$\underline{\mathcal O}_X(1)$; see XVII 5.7.5 and 5.7.6.

\subsubsection*{3.2.5.}
Let us make this explicit in a concrete case.  Take $X=\mathbb P^2$ and let
$\{H_\tau\}$ be a pencil of plane curves of degree $d$ satisfying 3.1,
(3.2.1), and (3.2.2).  Then
\begin{equation}\tag{3.2.6}
\begin{array}{ll}
\chi(X)=3 & \text{(by 1.2)},\\[4pt]
\chi(X_{\bar\eta})=-d(d-3) & \text{(see SGA 5 VII 7.3)},\\[4pt]
\chi(\Delta)=d^2,&
\end{array}
\end{equation}
whence
\begin{equation}\tag{3.2.7}
\text{the total number of singular points is }3(d-1)^2.
\end{equation}
Take the special case $d=3$ and consider the family of elliptic curves
(in characteristic different from $3$)
\begin{equation}\tag{3.2.8}
X^3+Y^3+Z^3+3\mu XYZ.
\end{equation}
The four singular fibers lie over $\mu^3=-1$ and $\mu=\infty$,
respectively,\footnote{The print says that the finite parameters are the
cube roots of unity.  With the printed plus sign in (3.2.8), the Hesse
factorization gives $\mu^3=-1$; cube roots of $+1$ correspond to the
polynomial with coefficient $-3\mu$.}
and each singular fiber is the union of three lines as follows:
\begin{equation}\tag{3.2.9}
\begin{tikzcd}[row sep=4.5em,column sep=4.5em]
{} & {} \arrow[dl,dash,shorten <=-9mm,shorten >=-9mm]
        \arrow[dr,dash,shorten <=-9mm,shorten >=-9mm] & {} & {} \\
{} \arrow[rr,dash,shorten <=-11mm,shorten >=-11mm] & {} & {} & {,}
\end{tikzcd}
\end{equation}
It has exactly three singular points, all nondegenerate quadratic
singularities.  Altogether this gives twelve singular points and verifies
(3.2.7) in this example.

\subsubsection*{Proposition 3.2.10.}
\underline{Let $X\hookrightarrow\mathbb P^r$ be a Lefschetz embedding}
(XVII 2.3), \underline{and suppose that $n=\dim X$ is even, or that
$\operatorname{char}k\ne2$.  Then every Lefschetz pencil, regarded as a
line in $\check{\mathbb P}^{r}$, meets $\check X$ in}
\begin{equation}\tag{3.2.11}
(-1)^n\bigl[\chi(X)+\chi(\Delta)-2\chi(X_{\bar\eta})\bigr]
\end{equation}
\underline{points, and meets it transversely} (EGA IV 17.13.7).

\underline{Proof.}
By hypothesis, every singularity of a fiber of the Lefschetz pencil is a
nondegenerate quadratic singularity.  Formula (3.2.4) therefore shows that
the number of singular fibers in the pencil $D$---each containing exactly
one singular point---is the number (3.2.11).  Now the fiber $X_\tau$ is
singular if and only if $\tau\in\check X\cap D$.

To see that the intersection is transverse, observe that (3.2.11) is
\underline{independent} of the particular Lefschetz pencil.  Since the
generic pencil is a Lefschetz pencil, (3.2.11) is the degree of $\check X$;
that is, its degree as a hypersurface in $\check{\mathbb P}^{r}$, taken by
convention to be $0$ if $\dim\check X\le r-2$.  If a line $D$ meets a
hypersurface of degree $d>0$ in $d$ distinct points $x_i$
($1\le i\le d$), the intersection must be transverse.  Indeed, the local
intersection multiplicity
\[
\mu_i=\operatorname{length}_{x_i}(D\cap\check X)
\]
is at least $1$, while $\sum_{i=1}^d\mu_i=d$.  Hence $\mu_i=1$ for every
$i$; equivalently, $D\cap\check X$ is smooth, which is precisely
transversality.\footnote{The print denotes the local length by
$\operatorname{rg}$; the standard local-length notation and the stated
intersection-multiplicity argument require $\operatorname{lg}$ (or,
equivalently, $\operatorname{length}$).}

\subsubsection*{Remark 3.2.12.}
For a Lefschetz pencil $D$ already assumed to meet $\check X$ transversely
(a condition satisfied in any case by sufficiently general $D$), the
result 3.2.10 may also be regarded as the conjunction of XVII 5.7.2 and
the fact that $\deg\varphi=1$ contained in XVII 3.5.  Of course, the fact
proved here that \underline{every} Lefschetz pencil $D$ meets $\check X$
transversely---by means of the Euler--Poincar\'e formula XVI 2.1---may also
be verified directly and elementarily, using the fact that the axis of
$D$ meets $X$ transversely.
\end{filecontents*}
% Original input: components/114_expose_XVIII_section_4_opening_and_proposition_4_2.tex
\begin{filecontents*}[overwrite,noheader]{sga7ii-cumulative-input-114.tex}
\section*{4. Cohomology of the blow-up associated with a pencil}

\subsection*{4.1.}
We return to the situation of 3.1, summarized by the commutative diagram
(see 2.1.6.3):
\begin{equation}\tag{4.1.1}
\resizebox{\textwidth}{!}{%
\begin{tikzcd}[ampersand replacement=\&,row sep=huge,column sep=huge]
X\times\mathbb P^1
  \& \widetilde X
      \arrow[l,hook',"f\times\rho"]
      \arrow[d,"f"]
      \arrow[dl,"\rho"']
  \& \widetilde\Delta
      \arrow[l,hook',"i"]
      \arrow[d,"g"]
      \arrow[r,"{(g,\rho_i)}","\sim"']
  \& \Delta\times\mathbb P^1 \\
\mathbb P^1 \& X \& \Delta\arrow[l,hook',"j"] \& {}
\end{tikzcd}}
\end{equation}
We begin by making (2.2.2) more explicit in this case.

\subsection*{Proposition 4.2.}
\underline{Under the hypotheses of 3.1:}

\smallskip
\noindent
\textup{(4.2.1)} \underline{The homomorphisms of $\mathbb Z_\ell$-modules}
\[
\begin{tikzcd}[column sep=huge]
H^\bullet(\widetilde X,\mathbb Z_\ell)
  \arrow[r,shift left=.7ex,"f_*\oplus(1-f^*f_*)"]
& H^\bullet(X,\mathbb Z_\ell)\oplus
   \operatorname{Ker}\!\left(f_*:H^\bullet(\widetilde X,\mathbb Z_\ell)
       \to H^\bullet(X,\mathbb Z_\ell)\right)
  \arrow[l,shift left=.7ex,"f^*+\operatorname{incl}"']
\end{tikzcd}
\]
\underline{are mutually inverse isomorphisms.}

\smallskip
\noindent
\textup{(4.2.2)} \underline{The homomorphisms of $\mathbb Z_\ell$-modules}
\[
\begin{tikzcd}[column sep=huge]
H^\bullet(\widetilde X,\mathbb Z_\ell)
  \arrow[r,shift left=.7ex,"f_*\oplus(-g_*i^*)"]
& H^\bullet(X,\mathbb Z_\ell)\oplus
   H^{\bullet-2}(\Delta,\mathbb Z_\ell(-1))
  \arrow[l,shift left=.7ex,"f^*+i_*g^*"']
\end{tikzcd}
\]
\underline{are mutually inverse isomorphisms.}\footnote{The print omits
the minus sign before $g_*i^*$ in (4.2.2), (4.2.6), (4.2.9), and
(4.2.10).  But Lemma 4.3 gives
$g_*i^*i_*g^*=-\mathrm{id}$, and the printed product calculation itself
uses $-g_*i^*$.  The sign is therefore restored consistently.}

\smallskip
\noindent
\textup{(4.2.3)} \underline{The algebra structure on}
\[
H^\bullet(X,\mathbb Z_\ell)\oplus
H^{\bullet-2}(\Delta,\mathbb Z_\ell(-1)),
\]
\underline{transported by (4.2.2) from the cup product on}
$H^\bullet(\widetilde X,\mathbb Z_\ell)$,
\underline{is expressed, for}
\[
a,b\in H^\bullet(X,\mathbb Z_\ell),
\qquad
x,y\in H^{\bullet-2}(\Delta,\mathbb Z_\ell(-1)),
\]
\underline{by the formulas}
\begin{equation}\tag{4.2.3.1}
(0\oplus x)\wedge(0\oplus y)
=-j_*(xy)\oplus2xy\,c_1(\underline{\mathcal O}_\Delta(1)),
\end{equation}
\begin{equation}\tag{4.2.3.2}
(a\oplus0)\wedge(b\oplus0)=ab\oplus0,
\end{equation}
\begin{equation}\tag{4.2.3.3}
(a\oplus0)\wedge(0\oplus y)=0\oplus j^*(a)y,
\end{equation}
\begin{equation}\tag{4.2.3.4}
(0\oplus x)\wedge(b\oplus0)=0\oplus xj^*(b).
\end{equation}
\underline{In particular, if}
\[
a\oplus x\in H^\bullet(X,\mathbb Z_\ell)
  \oplus H^{\bullet-2}(\Delta,\mathbb Z_\ell(-1))
\]
and
\[
b\oplus y\in H^{2n-\bullet}(X,\mathbb Z_\ell)
  \oplus H^{2n-2-\bullet}(\Delta,\mathbb Z_\ell(-1)),
\]
\underline{then}
\begin{equation}\tag{4.2.3.5}
(a\oplus x)\wedge(b\oplus y)=ab-j_*(xy).
\end{equation}
In other words, the total intersection matrix of
$H^\bullet(\widetilde X,\mathbb Z_\ell)$ is
\[
\left(
\begin{array}{c|c}
\begin{array}{c}
\text{total intersection}\\[-1pt]
\text{matrix on $X$}
\end{array} & 0\\ \hline
0 &
-\begin{array}{c}
\text{total intersection}\\[-1pt]
\text{matrix on $\Delta$}
\end{array}
\end{array}
\right).
\]

\underline{Proof.}
We have already proved (4.2.1); see (2.2.22).  For (4.2.2), first observe
that, since $g:\widetilde\Delta\to\Delta$ is a projective-line bundle, the
sequence
\begin{equation}\tag{4.2.4}
0\longrightarrow H^\bullet(\Delta,\mathbb Z_\ell)
\xrightarrow{g^*}H^\bullet(\widetilde\Delta,\mathbb Z_\ell)
\xrightarrow{g_*}H^{\bullet-2}(\Delta,\mathbb Z_\ell(-1))
\longrightarrow0
\end{equation}
is exact, by (1.2.2).  Together with (2.2.26), this gives the commutative
diagram of exact sequences
\begin{equation}\tag{4.2.5}
\resizebox{\textwidth}{!}{%
\begin{tikzcd}[ampersand replacement=\&,column sep=large,row sep=huge]
0\arrow[r]
  \&H^\bullet(X,\mathbb Z_\ell)
       \arrow[r,"f^*"]\arrow[d,"j^*" right]
  \&H^\bullet(\widetilde X,\mathbb Z_\ell)
       \arrow[r,"g_*i^*"]\arrow[d,"i^*" right]
  \&H^{\bullet-2}(\Delta,\mathbb Z_\ell(-1))
       \arrow[r]\arrow[d,equal]
  \&0\\
0\arrow[r]
  \&H^\bullet(\Delta,\mathbb Z_\ell)
       \arrow[r,"g^*"]
  \&H^\bullet(\widetilde\Delta,\mathbb Z_\ell)
       \arrow[r,"g_*"]
  \&H^{\bullet-2}(\Delta,\mathbb Z_\ell(-1))
       \arrow[r]
  \&0.
\end{tikzcd}}
\end{equation}
In view of the splitting (4.2.1), this yields the isomorphism
\begin{equation}\tag{4.2.6}
f_*\oplus(-g_*i^*):
H^\bullet(\widetilde X,\mathbb Z_\ell)
\xrightarrow{\sim}
H^\bullet(X,\mathbb Z_\ell)\oplus
H^{\bullet-2}(\Delta,\mathbb Z_\ell(-1)).
\end{equation}
The transpose of (4.2.5), after the degree reversal
$\bullet\mapsto2n-\bullet$, gives the commutative diagram of exact sequences
\begin{equation}\tag{4.2.7}
\resizebox{\textwidth}{!}{%
\begin{tikzcd}[ampersand replacement=\&,column sep=large,row sep=large]
0
  \&H^\bullet(X,\mathbb Z_\ell)\arrow[l]
  \&H^\bullet(\widetilde X,\mathbb Z_\ell)\arrow[l,"f_*"']
  \&H^{\bullet-2}(\Delta,\mathbb Z_\ell(-1))
       \arrow[l,"i_*g^*"']
  \&0\arrow[l]\\
0
  \&H^{\bullet-2}(\Delta,\mathbb Z_\ell(-2))
       \arrow[l]\arrow[u,"j_*"']
  \&H^{\bullet-2}(\widetilde\Delta,\mathbb Z_\ell(-1))
       \arrow[l,"g_*"']\arrow[u,"i_*"']
  \&H^{\bullet-2}(\Delta,\mathbb Z_\ell(-1))
       \arrow[l,"g_*"']\arrow[u,equal]
  \&0\arrow[l]
\end{tikzcd}}
\end{equation}
and hence the isomorphism
\begin{equation}\tag{4.2.8}
H^\bullet(X,\mathbb Z_\ell)\oplus
H^{\bullet-2}(\Delta,\mathbb Z_\ell(-1))
\xrightarrow[\sim]{\ f^*+i_*g^*\ }
H^\bullet(\widetilde X,\mathbb Z_\ell).
\end{equation}
\end{filecontents*}
% Original input: components/115_expose_XVIII_completion_of_section_4.tex
\begin{filecontents*}[overwrite,noheader]{sga7ii-cumulative-input-115.tex}
It remains to prove that the composite
\begin{equation}\tag{4.2.9}
\begin{aligned}
&H^\bullet(X,\mathbb Z_\ell)\oplus
 H^{\bullet-2}(\Delta,\mathbb Z_\ell(-1))
 \xrightarrow{\ f^*+i_*g^*\ }
 H^\bullet(\widetilde X,\mathbb Z_\ell)\\[-2pt]
&\hspace{55mm}\xrightarrow{\ f_*\oplus(-g_*i^*)\ }
 H^\bullet(X,\mathbb Z_\ell)\oplus
 H^{\bullet-2}(\Delta,\mathbb Z_\ell(-1))
\end{aligned}
\end{equation}
is the identity.  Now
\[
f_*f^*=\mathrm{id}\quad\text{on }H^\bullet(X,\mathbb Z_\ell)
\quad\text{by (2.2.18)},
\]
\[
g_*i^*f^*=0\quad\text{by (4.2.5)},
\qquad
f_*i_*g^*=0\quad\text{by (4.2.7)}.
\]
It therefore suffices to prove
\begin{equation}\tag{4.2.10}
-g_*i^*i_*g^*(x)=x
\quad\text{for }x\in H^\bullet(\Delta,\mathbb Z_\ell(-1)).
\end{equation}

To verify (4.2.10), first observe that
$i:\widetilde\Delta\hookrightarrow\widetilde X$ is the inclusion of a
smooth subvariety of codimension one.  Thus (SGA 5 VII)
\begin{equation}\tag{4.2.11}
i^*i_*(x)=c_1(\underline N_{\widetilde\Delta/\widetilde X})\,x,
\end{equation}
and
\begin{equation}\tag{4.2.12}
i_*(x)i_*(y)
=i_*\!\left(xy\,c_1(\underline N_{\widetilde\Delta/\widetilde X})\right).
\end{equation}
Write
$\xi_X=\underline{\mathcal O}_X(1)$,
$\xi_\Delta=\underline{\mathcal O}_\Delta(1)$, and
$\xi_{\mathbb P^1}=\underline{\mathcal O}_{\mathbb P^1}(1)$.
By EGA II 8.1.7, if $\widetilde{\mathfrak J}$ is the ideal of
$\widetilde\Delta$ in $\widetilde X$, then
\begin{equation}\tag{4.2.13}
\begin{aligned}
\underline{\check N}_{\widetilde\Delta/\widetilde X}
&\simeq\widetilde{\mathfrak J}/\widetilde{\mathfrak J}^{,2}
 \simeq\widetilde{\mathfrak J}
       \otimes_{\underline{\mathcal O}_{\widetilde X}}
       \underline{\mathcal O}_{\widetilde\Delta}\\
&\simeq\underline{\mathcal O}_{\widetilde X}(1)
       \otimes_{\underline{\mathcal O}_{\widetilde X}}
       \underline{\mathcal O}_{\widetilde\Delta}
 =\underline{\mathcal O}_{\widetilde\Delta}(1).
\end{aligned}
\end{equation}
Here $\widetilde\Delta$ is viewed as
$\mathbb P(\mathfrak J/\mathfrak J^2)$, where $\mathfrak J$ is the ideal of
$\Delta$ in $X$.  Since $\Delta$ is the intersection of $X$ with two given
hyperplanes, there is a canonical isomorphism
\[
\mathfrak J/\mathfrak J^2\simeq
\underline{\mathcal O}_\Delta^{,2}(-1),
\]
which induces
\[
\widetilde\Delta=\mathbb P(\mathfrak J/\mathfrak J^2)
\simeq\mathbb P(\underline{\mathcal O}_\Delta^{,2}(-1))
\xrightarrow[\sim]{\alpha}
\mathbb P(\underline{\mathcal O}_\Delta^{,2})
\simeq\Delta\times\mathbb P^1.
\]
The canonical isomorphism $\alpha$ of EGA II 4.1.4 carries the tautological
line bundle to
\[
(\rho i)^*\underline{\mathcal O}_{\mathbb P^1}(1)
\otimes g^*\underline{\mathcal O}_\Delta(-1).
\]
Consequently, (4.2.13) becomes
\begin{equation}\tag{4.2.14}
\underline{\check N}_{\widetilde\Delta/\widetilde X}
\simeq(\rho i)^*\underline{\mathcal O}_{\mathbb P^1}(1)
       \otimes g^*\underline{\mathcal O}_\Delta(-1).
\end{equation}

\subsection*{Lemma 4.3.}
\underline{One has}
\[
c_1(\underline N_{\widetilde\Delta/\widetilde X})
=g^*c_1(\xi_\Delta)-(\rho i)^*c_1(\xi_{\mathbb P^1}).
\]
Return now to the proof of (4.2.10).  By (4.2.11) and Lemma 4.3, it is enough
to prove
\[
g_*\!\left(g^*(x)\left[(\rho i)^*c_1(\xi_{\mathbb P^1})
                       -g^*c_1(\xi_\Delta)\right]\right)=x.
\]
By (4.2.4), $g_*g^*=0$, while
\[
\begin{tikzcd}[column sep=huge]
\widetilde\Delta\arrow[r,"g\times\rho_i","\sim"']
  &\Delta\times\mathbb P^1,
\end{tikzcd}
\]
so
\begin{equation}\tag{4.3.1}
g_*\!\left(g^*(x)(\rho i)^*c_1(\xi_{\mathbb P^1})\right)=x,
\end{equation}
and hence
\begin{equation}\tag{4.3.2}
g_*\!\left(g^*(x)c_1(\underline N_{\widetilde\Delta/\widetilde X})\right)
=-x.
\end{equation}

We now compute the algebra structure on
\[
H^\bullet(X,\mathbb Z_\ell)\oplus
H^{\bullet-2}(\Delta,\mathbb Z_\ell(-1))
\]
transported from the cup product by (4.2.2).  Let $a,b$ belong to
$H^\bullet(X,\mathbb Z_\ell)$ and $x,y$ to
$H^{\bullet-2}(\Delta,\mathbb Z_\ell(-1))$.  Begin with (4.2.3.1), and write
\[
(0\oplus x)\wedge(0\oplus y)=A\oplus B.
\]
By (4.2.12) and (4.3.2),
\[
\begin{aligned}
A
&=f_*\!\left((i_*g^*x)(i_*g^*y)\right)\\
&=f_*i_*\!\left(g^*(xy)c_1(\underline N_{\widetilde\Delta/\widetilde X})\right)\\
&=j_*g_*\!\left(g^*(xy)c_1(\underline N_{\widetilde\Delta/\widetilde X})\right)
=-j_*(xy),
\end{aligned}
\]
and
\[
\begin{aligned}
B
&=-g_*i^*\!\left((i_*g^*x)(i_*g^*y)\right)\\
&=-g_*i^*i_*\!\left(g^*(xy)c_1(\underline N_{\widetilde\Delta/\widetilde X})\right)\\
&=-g_*\!\left(g^*(xy)c_1(\underline N_{\widetilde\Delta/\widetilde X})^2\right).
\end{aligned}
\]
Lemma 4.3 gives
\[
\begin{aligned}
c_1(\underline N_{\widetilde\Delta/\widetilde X})^2
&=g^*\!\left(c_1(\xi_\Delta)^2\right)
 -2g^*c_1(\xi_\Delta)(\rho i)^*c_1(\xi_{\mathbb P^1}),
\end{aligned}
\]
because $c_1(\xi_{\mathbb P^1})^2=0$.  Therefore, by (4.2.4) and (4.3.1),
\[
B=0+2xy\,c_1(\xi_\Delta),
\]
which proves (4.2.3.1).  Formula (4.2.3.2) expresses the fact that $f^*$ is
an algebra homomorphism.

To prove (4.2.3.3), the projection formula shows that
$\operatorname{Ker}f_*$ is stable under multiplication by
$\operatorname{im}f^*$.  Hence
\[
(a\oplus0)\wedge(0\oplus y)=0\oplus C,
\]
where
\begin{equation}\tag{4.3.3}
C=-g_*i^*\!\left(f^*(a)i_*g^*(y)\right).
\end{equation}
Using (4.3.2),
\[
\begin{aligned}
C
&=-g_*\!\left(i^*f^*(a)g^*(y)
      c_1(\underline N_{\widetilde\Delta/\widetilde X})\right)\\
&=-g_*\!\left(g^*j^*(a)g^*(y)
      c_1(\underline N_{\widetilde\Delta/\widetilde X})\right)\\
&=j^*(a)y.
\end{aligned}
\]
This proves (4.2.3.3).  Finally, (4.2.3.4) follows from (4.2.3.3) by graded
commutativity of the cup product.
\end{filecontents*}
% Original input: components/116_expose_XVIII_section_5_1.tex
\begin{filecontents*}[overwrite,noheader]{sga7ii-cumulative-input-116.tex}
\section*{5. Cohomological study of the morphism
$\rho:\widetilde X\longrightarrow\mathbb P^1$ associated with a pencil}

\subsection*{5.1. The restriction homomorphism
$H^*(\widetilde X)\longrightarrow H^*(Y)$}

Let $Y$ denote one of the smooth fibers of
$\rho:\widetilde X\to\mathbb P^1$, say
\[
Y=\rho^{-1}(t_0)=X\mathbin{\cdot}H_{t_0}.
\]
We have the commutative diagram
\begin{center}
\resizebox{\textwidth}{!}{%
\begin{tikzcd}[ampersand replacement=\&,column sep=large,row sep=huge]
\mathbb P^1
  \& {} \& \widetilde X\arrow[ll,"\rho"']\arrow[d,"f"]
  \& {} \& \widetilde\Delta
       \arrow[ll,hook',"i"]\arrow[d,"g"]
       \arrow[r,phantom,"\simeq"]
  \& \Delta\times\mathbb P^1 \\
\mathrm{Sp}(k)=\{t_0\}\arrow[u,hook]
  \& {} \& X
  \& {} \& \Delta
       \arrow[ll,hook',"j"]
       \arrow[ur,hook,"{\ell=\mathrm{id}\times\{t_0\}}"']
       \arrow[dlll,hook,"h"]
  \& {} \\
{} \& Y\arrow[ul]\arrow[uur,"k"]\arrow[ur,"m"']
  \& {} \& {} \& {} \& {}
\end{tikzcd}}
\end{center}

\subsection*{Proposition 5.1.1.}
\underline{Under the hypotheses of (3.1), with $Y$ denoting a smooth
fiber:}

\smallskip
\noindent
\textup{(5.1.2)} \underline{The restriction homomorphism}
\[
k^*:H^\bullet(\widetilde X,\mathbb Z_\ell)
   \longrightarrow H^\bullet(Y,\mathbb Z_\ell)
\]
\underline{is expressed (via (4.2.2)) by}
\[
H^\bullet(X,\mathbb Z_\ell)\oplus
H^{\bullet-2}(\Delta,\mathbb Z_\ell(-1))
\xrightarrow{\ m^*+h_*\ }
H^\bullet(Y,\mathbb Z_\ell).
\]

\noindent
\textup{(5.1.3)} \underline{The Gysin map}
\[
k_*:H^{\bullet-2}(Y,\mathbb Z_\ell(-1))
   \longrightarrow H^\bullet(\widetilde X,\mathbb Z_\ell)
\]
\underline{is expressed by}
\[
H^{\bullet-2}(Y,\mathbb Z_\ell(-1))
\xrightarrow{\ m_*\oplus(-h^*)\ }
H^\bullet(X,\mathbb Z_\ell)\oplus
H^{\bullet-2}(\Delta,\mathbb Z_\ell(-1)).
\]

\underline{Proof.}
The smooth subschemes $\widetilde\Delta$ and $Y$ of $\widetilde X$
intersect transversely in $\Delta$; in other words, the square of smooth
schemes
\begin{equation}\tag{5.1.4}
\begin{tikzcd}[column sep=large,row sep=large]
\Delta\arrow[r,"\ell"]\arrow[d,hook,"h"]
  &\widetilde\Delta\arrow[d,hook,"i"]\\
Y\arrow[r,"k"']&\widetilde X
\end{tikzcd}
\end{equation}
is Cartesian.  It therefore implies the formula [SGA 5, IV]
\begin{equation}\tag{5.1.5}
h_*\ell^*=k^*i_* ,\qquad \ell_*h^*=i^*k_*.
\end{equation}
The composite
\[
H^\bullet(X,\mathbb Z_\ell)\oplus
H^{\bullet-2}(\Delta,\mathbb Z_\ell(-1))
\xrightarrow{\ f^*+i_*g^*\ }
H^\bullet(\widetilde X,\mathbb Z_\ell)
\xrightarrow{\ k^*\ }
H^\bullet(Y,\mathbb Z_\ell)
\]
is the sum of
$k^*f^*=m^*$ and
$k^*i_*g^*=h_*\ell^*g^*$ by (5.1.5).  Since
$g\ell=\mathrm{id}_\Delta$, this proves (5.1.2).

The composite
\[
H^{\bullet-2}(Y,\mathbb Z_\ell(-1))
\xrightarrow{\ k_*\ }
H^\bullet(\widetilde X,\mathbb Z_\ell)
\xrightarrow{\ f_*\oplus(-g_*i^*)\ }
H^\bullet(X,\mathbb Z_\ell)\oplus
H^{\bullet-2}(\Delta,\mathbb Z_\ell(-1))
\]
is the sum of
\[
f_*k_*=m_*
\qquad\text{and}\qquad
-g_*i^*k_*=-g_*\ell_*h^*=-h^*,
\]
which proves (5.1.3).

\subsection*{Corollary 5.1.6.}
\underline{Under the hypotheses of 5.1.1, for $i\leq n-1$, the images
of}
\[
m^*:H^i(X,\mathbb Z_\ell)\longrightarrow H^i(Y,\mathbb Z_\ell)
\]
\underline{and of}
\[
k^*:H^i(\widetilde X,\mathbb Z_\ell)
   \longrightarrow H^i(Y,\mathbb Z_\ell)
\]
\underline{are equal.}

\underline{Proof.}
By the weak Lefschetz theorem for hyperplane sections (SGA 5, VII 7.1),
for $i\leq n-1$ the map
\begin{equation}\tag{5.1.7}
H^{i-2}(X,\mathbb Z_\ell(-1))
\xrightarrow{\ j^*\ }
H^{i-2}(\Delta,\mathbb Z_\ell(-1))
\end{equation}
is an isomorphism.  Thus, by (5.1.2), the image of $k^*$ is the sum of
the image of $m^*$ and the image of the composite
\[
H^{i-2}(X,\mathbb Z_\ell(-1))
\xrightarrow{\ j^*\ }
H^{i-2}(\Delta,\mathbb Z_\ell(-1))
\xrightarrow{\ h_*\ }
H^i(Y,\mathbb Z_\ell).
\]
But the projection formula gives
\[
\begin{aligned}
h_*j^*(x)
 &=h_*\bigl(h^*m^*(x)\bigr)
  =m^*(x)h_*(1)
  =\xi_Y\,m^*(x)\\
 &=m^*(x\xi_X),
 \qquad\text{since }m^*(\xi_X)=\xi_Y.
\end{aligned}
\]
Consequently $\operatorname{Im}k^*\subset\operatorname{Im}m^*$, and
the reverse inclusion follows directly from (5.1.2).  This proves the
corollary.
\end{filecontents*}
% Original input: components/117_expose_XVIII_section_5_2.tex
\begin{filecontents*}[overwrite,noheader]{sga7ii-cumulative-input-117.tex}
\subsection*{5.2. The hard Lefschetz theorem (LV)}

\subsubsection*{5.2.1.}
Let $S$ be projective of dimension $n$ over the algebraically closed field
$k$, and let $S\hookrightarrow\mathbb P$ be a projective embedding.  Denote by
\[
u\in H^2(S,\mathbb Z_\ell(1))
\]
the cohomology class of a hyperplane section.  Thus $u$ is the pullback along
$S\hookrightarrow\mathbb P$ of the first Chern class of
$\underline{\mathcal O}_{\mathbb P}(1)$ in
$H^2(\mathbb P,\mathbb Z_\ell(1))$.  Write
\[
L_S:H^\bullet(S,\mathbb Q_\ell(j))
   \longrightarrow H^{\bullet+2}(S,\mathbb Q_\ell(j+1))
\]
for multiplication by $u$.

\subsubsection*{5.2.2.}
\underline{We say that $S$ (equipped with $u$) satisfies (LV), the hard
Lefschetz theorem, if, for every (or for one) integer $j$ and every integer
$i>0$, the maps}
\[
L_S^i:H^{n-i}(S,\mathbb Q_\ell(j))
\longrightarrow H^{n+i}(S,\mathbb Q_\ell(i+j))
\]
\underline{are isomorphisms.}

\subsubsection*{5.2.2.1.}
It is conjectured that this condition holds when $S$ is smooth.  It is true
when $k$ has characteristic zero [6].  The specialization theorem
[SGA 4, XVI 2.2] therefore gives the same conclusion when $S$, together with
its projective embedding, lifts to characteristic zero.  In particular, it
holds for a smooth complete intersection.  It is also known for a smooth $S$
of dimension at most two [2], and for an abelian variety, either directly
(Lieberman [4] and [2]) or by Mumford's theorem [5] that a polarized abelian
variety lifts, together with its polarization, to characteristic zero.

Recall that, for a projective variety $S$ and $i\leq n$, the primitive part of
cohomology is defined by
\begin{equation}\tag{5.2.2.2}
\operatorname{Prim}^i(S,\mathbb Q_\ell(j))
=\operatorname{Ker}\!\left(
L^{n-i+1}:H^i(S,\mathbb Q_\ell(j))
\longrightarrow
H^{2n-i+2}(S,\mathbb Q_\ell(j+n-i+1))
\right).
\end{equation}
If $S$ satisfies (LV), then for $i\leq n$ one has the decomposition
\begin{equation}\tag{5.2.2.3}
H^i(S,\mathbb Q_\ell(j))
\xleftarrow[\sim]{\ 1+L\ }
\operatorname{Prim}^i(S,\mathbb Q_\ell(j))
\oplus H^{i-2}(S,\mathbb Q_\ell(j-1)).
\end{equation}

\subsection*{Proposition 5.2.3.}
\underline{Let $X$ be smooth and projective and suppose that it satisfies
(LV).  Let $\tau:Y\hookrightarrow X$ be the inclusion of a smooth hyperplane
section.  Then}
\begin{equation}\tag{5.2.3.1}
\operatorname{Prim}^n(X,\mathbb Q_\ell(j))
=\operatorname{Ker}\!\left(
\tau^*:H^n(X,\mathbb Q_\ell(j))
\longrightarrow H^n(Y,\mathbb Q_\ell(j))
\right).
\end{equation}

\underline{Proof.}
The projection formula for $\tau:Y\to X$ gives the commutative diagram
\begin{equation}\tag{5.2.3.2}
\begin{tikzcd}[column sep=huge,row sep=huge]
H^n(X,\mathbb Q_\ell(j))
  \arrow[rr,"L"]\arrow[dr,"\tau^*"']
  & {} &
H^{n+2}(X,\mathbb Q_\ell(j+1))
  \arrow[dl,"\tau_*"]\\
{} & H^n(Y,\mathbb Q_\ell(j)) & {}
\end{tikzcd}
\end{equation}
By duality and weak Lefschetz,
\[
\tau_*:H^n(Y,\mathbb Q_\ell(j))
\longrightarrow H^{n+2}(X,\mathbb Q_\ell(j+1))
\]
is the transpose of an isomorphism and hence is itself an isomorphism.  The
assertion follows.

\subsection*{Lemma 5.2.4.}
\underline{Let $X$ be smooth and projective, and let
$\tau:Y\hookrightarrow X$ be the inclusion of a smooth hyperplane section.
Then $X$ satisfies (LV) if and only if the following two conditions hold:}

\smallskip
\noindent
\textup{5.2.4.1.} \underline{$Y$ satisfies (LV).}

\smallskip
\noindent
\textup{5.2.4.2.} \underline{The restriction of the cup product on
$H^{n-1}(Y,\mathbb Q_\ell)$ to the subspace
$\operatorname{Im}H^{n-1}(X,\mathbb Q_\ell)$ is nondegenerate.}

\underline{Proof.}
We have
\begin{equation}\tag{5.2.4.3}
\left\{
\begin{aligned}
L_X&=\tau_*\tau^*,\\[1ex]
L_Y&=\tau^*\tau_* \qquad\text{(by 4.2.11),}
\end{aligned}
\right.
\end{equation}
and consequently
\begin{equation}\tag{5.2.4.4}
L_X^{i+1}=\tau_*L_Y^i\tau^*.
\end{equation}
In the commutative diagram
\begin{equation}\tag{5.2.4.5}
\begin{tikzcd}[column sep=huge,row sep=huge]
H^{n-i-1}(X,\mathbb Q_\ell(j))
  \arrow[r,"L_X^{i+1}"]\arrow[d,"\tau^*"]
  &H^{n+i+1}(X,\mathbb Q_\ell(j+i+1))\\
H^{n-1-i}(Y,\mathbb Q_\ell(j))
  \arrow[r,"L_Y^i"]
  &H^{n-1+i}(Y,\mathbb Q_\ell(j+i))\arrow[u,"\tau_*"']
\end{tikzcd}
\end{equation}
$\tau^*$ is an isomorphism for $i\geq1$ by weak Lefschetz, and $\tau_*$ is
an isomorphism by duality.  Hence $X$ satisfies (LV) if and only if $Y$
satisfies (LV) and
\begin{equation}\tag{5.2.4.6}
L_X:H^{n-1}(X,\mathbb Q_\ell(j))
\longrightarrow H^{n+1}(X,\mathbb Q_\ell(j+1))
\end{equation}
is an isomorphism.

It remains to identify (5.2.4.6) with 5.2.4.2.  Let
$a,b\in H^{n-1}(X,\mathbb Q_\ell(j))$.  By Poincaré duality on $X$,
(5.2.4.6) is equivalent to nondegeneracy of the bilinear form
\begin{equation}\tag{5.2.4.7}
\begin{aligned}
H^{n-1}(X,\mathbb Q_\ell(j))\times
H^{n-1}(X,\mathbb Q_\ell(j))
&\longrightarrow H^{2n}(X,\mathbb Q_\ell(2j+1)),\\
(a,b)&\longmapsto a\,L_X(b).
\end{aligned}
\end{equation}
But by (5.2.4.3) and the projection formula,
\[
aL_X(b)=a\tau_*\tau^*(b)
=\tau_*\bigl(\tau^*(a)\tau^*(b)\bigr).
\]
This proves the equivalence of (5.2.4.7) and 5.2.4.2.

\subsection*{Corollary 5.2.5.}
\underline{Let $X$ be smooth and projective of dimension $n$, and let
$\tau:Y\hookrightarrow X$ be the inclusion of a smooth hyperplane section.
Suppose that $X$ satisfies (LV).  Then, for $q\geq n$, the map}
\[
\tau^*:H^q(X,\mathbb Q_\ell(j))
\longrightarrow H^q(Y,\mathbb Q_\ell(j))
\]
\underline{is surjective.}

\underline{Proof.}
By (5.2.4.3), for $i\geq0$ one has
\[
L_Y^{i+1}=\tau^*L_X^i\tau_*,
\]
and hence the commutative diagram
\begin{center}
\begin{tikzcd}[column sep=huge,row sep=huge]
H^{n-i}(X,\mathbb Q_\ell(j+1))
  \arrow[r,"L_X^i"]
  &H^{n+i}(X,\mathbb Q_\ell(j+i+1))\arrow[d,"\tau^*"]\\
H^{n-i-2}(Y,\mathbb Q_\ell(j))
  \arrow[u,"\tau_*"']\arrow[r,"L_Y^{i+1}"]
  &H^{n+i}(Y,\mathbb Q_\ell(j+i+1)).
\end{tikzcd}
\end{center}
This gives the required result.

\subsection*{Lemma 5.2.6.}
\underline{Let $f:F\to B$ be a smooth proper morphism with $B$ connected.
Let $\underline{\mathcal O}_F(1)$ be a relatively ample invertible sheaf on
$F$, and equip each geometric fiber $F_{\bar b}$ with the ample invertible
sheaf induced by $\underline{\mathcal O}_F(1)$.  If one geometric fiber
satisfies (LV), then so do all the others.}

\underline{Proof.}
Let
$\xi_F\in H^2(F,\mathbb Z_\ell(1))$ be the first Chern class of
$\underline{\mathcal O}_F(1)$.  This class and its powers
$\xi_F^i\in H^{2i}(F,\mathbb Z_\ell(i))$ induce, by cup product, the
homomorphisms
\begin{equation}\tag{5.2.6.1}
\xi_F^i:R^{n-i}f_*\mathbb Q_\ell(j)
\longrightarrow R^{n+i}f_*\mathbb Q_\ell(j+i),
\end{equation}
where $n$ is the relative dimension of the fibers.  By proper base change
[SGA 4, XII 5.2], all geometric fibers satisfy (LV) if and only if all the
maps (5.2.6.1) are isomorphisms.  By the specialization theorem
[SGA 4, XVI 2.2], the sheaves $R^if_*\mathbb Q_\ell(j)$ are twisted constant
sheaves.  Thus the maps (5.2.6.1) are isomorphisms if and only if they are so
over one given geometric point.

\subsection*{Corollary 5.2.7.}
\underline{Let $X$ be smooth and projective.  If one smooth hyperplane
section satisfies (LV), then every smooth hyperplane section does.}

\underline{Proof.}
The universal family of smooth hyperplane sections of $X$ satisfies the
hypotheses of Lemma 5.2.6.
\end{filecontents*}
% Original input: components/118_expose_XVIII_section_5_3.tex
\begin{filecontents*}[overwrite,noheader]{sga7ii-cumulative-input-118.tex}
\subsection*{5.3. Condition (A)}

\subsubsection*{5.3.1.}
Let
\[
\rho:\widetilde X\longrightarrow\mathbb P^1
\]
be the ``fibration'' associated with a pencil whose axis meets $X$
transversely and whose generic fiber is assumed to be smooth.

Let $U$ be a nonempty open subset of $\mathbb P^1$ over which $\rho$ is
smooth, and let
\begin{equation}\tag{5.3.1.1}
\nu:U\hookrightarrow\mathbb P^1
\end{equation}
be the inclusion.  Since the sheaves
$R^i\rho_*\mathbb Q_\ell(j)$ are \underline{twisted constant over $U$}, the
sheaves
\[
\nu_*\nu^*R^i\rho_*\mathbb Q_\ell(j)
\]
are seen to be essentially \underline{independent} of the choice of $U$;
the canonical ``adjunction'' homomorphisms
\begin{equation}\tag{5.3.2}
R^i\rho_*\mathbb Q_\ell(j)
\longrightarrow
\nu_*\nu^*R^i\rho_*\mathbb Q_\ell(j)
\end{equation}
are likewise independent of the choice of $U$.

\subsubsection*{5.3.3.}
Observe that, for an \underline{arbitrary} pencil of hyperplane sections,
the weak Lefschetz theorem implies that, for $i\leq n-2$, the immersion
\[
\widetilde X\overset{(f\times\rho)}{\hookrightarrow}
X\times\mathbb P^1
\]
gives an isomorphism of sheaves on $\mathbb P^1$
\[
R^i\rho_*\mathbb Q_\ell(j)
\xleftarrow{\ \sim\ }
H^i(X,\mathbb Q_\ell(j))_{\mathbb P^1},
\qquad i\leq n-2.
\]
In other words, $R^i\rho_*\mathbb Q_\ell(j)$ is constant for
$i\leq n-2$ and, a fortiori, the maps (5.3.2) are isomorphisms for
$i\leq n-2$.

\subsubsection*{5.3.4.}
The local theory of vanishing cycles, and in particular the
\underline{Picard--Lefschetz formula} [SGA 7, XV 3.4], shows that,
\underline{for a Lefschetz pencil, the homomorphisms (5.3.2) are
isomorphisms for $i\neq n$, and are surjective for $i=n$}.

\subsubsection*{5.3.5.}
We say that $\rho$, or the pencil from which it arises,
\underline{satisfies condition (A)} if \underline{all} the homomorphisms
(5.3.2) are isomorphisms (equivalently, if this holds for every $i$ with
$j=0$).

We shall see below (6.3, 6.4) that condition (A) has a strong tendency to
hold.
\end{filecontents*}
% Original input: components/119_expose_XVIII_lemmas_5_4_and_5_5.tex
\begin{filecontents*}[overwrite,noheader]{sga7ii-cumulative-input-119.tex}
\subsection*{Lemma 5.4.}
\underline{Suppose that one smooth hyperplane section of $X$ (and hence
every such section, by Corollary 5.2.7) satisfies (LV)---for example, suppose
that $X$ satisfies (LV), cf. Lemma 5.2.4---and suppose that $\rho$ satisfies
condition (A).  Then every fiber (including every singular fiber) satisfies
(LV), and the sheaves $R^i\rho_*\mathbb Q_\ell(j)$ are constant for
$i\neq n-1$.}

\underline{Proof.}
Let
\[
L:R^i\rho_*\mathbb Q_\ell(j)
\longrightarrow R^{i+2}\rho_*\mathbb Q_\ell(j+1)
\]
denote multiplication by
$f^*c_1(\mathcal O_X(1))\in
H^2(\widetilde X,\mathbb Z_\ell(1))$; fiber by fiber, this is the operation
considered in (5.2.1).  By hypothesis, over $U$ we have an isomorphism
\begin{equation}\tag{5.4.1}
\nu^*(L^i):
\nu^*R^{n-1-i}\rho_*\mathbb Q_\ell(j)
\overset{\sim}{\longrightarrow}
\nu^*R^{n-1+i}\rho_*\mathbb Q_\ell(j+i)
\end{equation}
for every $i>0$ and every $j$.  Applying the functor $\nu_*$ gives the
\underline{isomorphisms}
\begin{equation}\tag{5.4.2}
(\nu_*\nu^*)L^i:
\nu_*\nu^*R^{n-1-i}\rho_*\mathbb Q_\ell(j)
\overset{\sim}{\longrightarrow}
\nu_*\nu^*R^{n-1+i}\rho_*\mathbb Q_\ell(j+i).
\end{equation}
In view of (A), these identify with the morphisms
\begin{equation}\tag{5.4.3}
L^i:R^{n-1-i}\rho_*\mathbb Q_\ell(j)
\overset{\sim}{\longrightarrow}
R^{n-1+i}\rho_*\mathbb Q_\ell(j+i).
\end{equation}
The conclusion follows from 5.3.3.

\subsubsection*{5.4.4.}
Let $Y$ be a smooth fiber of
$\rho:\widetilde X\to\mathbb P^1$.  By weak Lefschetz, the restriction map
\[
H^{n-1}(X,\mathbb Q_\ell(j))
\overset{m^*}{\hookrightarrow}
H^{n-1}(Y,\mathbb Q_\ell(j))
\]
is injective.  We saw in 5.2.4 that, if $X$ satisfies (LV), then the cup
product on $H^{n-1}(Y,\mathbb Q_\ell(j))$ remains nondegenerate when
restricted to $m^*H^{n-1}(X,\mathbb Q_\ell(j))$.  The orthogonal complement
of $H^{n-1}(X,\mathbb Q_\ell(j))$ in
$H^{n-1}(Y,\mathbb Q_\ell(j))$ is therefore a complementary subspace to
$H^{n-1}(X,\mathbb Q_\ell(j))$ in
$H^{n-1}(Y,\mathbb Q_\ell(j))$; following the classical terminology, it is
called the \underline{vanishing-cohomology space of $Y$}.

Let $\nu:U\hookrightarrow\mathbb P^1$ be any nonempty open subset over
which $\rho$ is smooth.  The corresponding subsheaf of
$\nu_*\nu^*R^{n-1}\rho_*\mathbb Q_\ell(j)$, the
``\underline{vanishing-cohomology sheaf}''
$E^{n-1}\rho_*\mathbb Q_\ell(j)$ on $\mathbb P^1$, is defined by
\begin{equation}\tag{5.4.4.1}
E^{n-1}\rho_*\mathbb Q_\ell(j)
=\nu_*\left\{
\begin{array}{l}
\text{the orthogonal complement in }
\nu^*R^{n-1}\rho_*\mathbb Q_\ell(j)\text{ of the}\\[1ex]
\text{constant subsheaf }H^{n-1}(X,\mathbb Q_\ell(j))_U
\end{array}
\right\}.
\end{equation}
It is clear that this sheaf is independent of the choice of $U$.

\subsection*{Lemma 5.5.}
\underline{Suppose that $X$ satisfies (LV) and that the maps (5.3.2) are
isomorphisms for $i=n-1$---for example, suppose that $\rho$ arises from a
Lefschetz pencil, cf. 5.3.4.  Then there is a direct-sum decomposition}
\begin{equation}\tag{5.5.1}
R^{n-1}\rho_*\mathbb Q_\ell(j)
\simeq
H^{n-1}(X,\mathbb Q_\ell(j))_{\mathbb P^1}
\oplus E^{n-1}\rho_*\mathbb Q_\ell(j).
\end{equation}

\underline{Proof.}
By (LV), over $U$ there is a decomposition
\begin{equation}\tag{5.5.2}
\nu^*R^{n-1}\rho_*\mathbb Q_\ell(j)
\simeq
H^{n-1}(X,\mathbb Q_\ell(j))_U
\oplus \nu^*E^{n-1}\rho_*\mathbb Q_\ell(j),
\end{equation}
to which we apply $\nu_*$.  This proves the assertion.
\end{filecontents*}
% Original input: components/120_expose_XVIII_theorem_5_6_and_computation.tex
\begin{filecontents*}[overwrite,noheader]{sga7ii-cumulative-input-120.tex}
\subsection*{Theorem 5.6.}
\underline{Let $X$ be a smooth, irreducible, projective variety satisfying
(LV) (5.2.2).  Let $\{H_\tau\}$ be a pencil of hyperplanes whose generic
member and axis meet $X$ transversely, in $X_\eta$ and $\Delta$ respectively.
Suppose that the projection
$\rho:\widetilde X\to\mathbb P^1$ satisfies condition (A).  Then:}
\footnote{The print says ``condition (5.3.2)'', but (5.3.2) is the displayed
family of adjunction maps; condition (A) is defined from those maps in 5.3.5.}

\smallskip
\noindent
\textup{(5.6.1)}\quad
\underline{the Leray spectral sequence}
\[
E_2^{p,q}=H^p(\mathbb P^1,R^q\rho_*\mathbb Q_\ell)
\Longrightarrow H^{p+q}(\widetilde X,\mathbb Q_\ell)
\]
\footnote{The printed abutment has $H^{m+q}$; the bidegree and the parallel
statement in Theorem 5.6.8 require $H^{p+q}$.}
\underline{degenerates.}

\noindent
\textup{(Note.)}\quad
\underline{This remains valid if, instead of assuming that $X$ satisfies
(LV), one assumes only that a smooth hyperplane section $Y$ does.}

\subsubsection*{5.6.2.}
\underline{Let $\eta$ denote the generic point of $\mathbb P^1$,
$\bar\eta$ the geometric generic point, $X_{\bar\eta}$ the geometric generic
fiber of $\rho$, and $\pi$ the Galois group of
$k(\bar\eta)/k(\eta)$.}

\underline{For $q\neq n-1$ there are canonical isomorphisms}
\begin{align}
E_2^{0,q}
 &=H^0(\mathbb P^1,R^q\rho_*\mathbb Q_\ell)
 \simeq H^q(X_{\bar\eta},\mathbb Q_\ell),
 \tag{5.6.2.1}\\
E_2^{2,q}
 &=H^2(\mathbb P^1,R^q\rho_*\mathbb Q_\ell)
 \simeq H^q(X_{\bar\eta},\mathbb Q_\ell(-1)),
 \tag{5.6.2.2}\\
E_2^{p,q}
 &=H^p(\mathbb P^1,R^q\rho_*\mathbb Q_\ell)=0,
 \qquad p\neq0,2.
 \tag{5.6.2.3}
\end{align}
\underline{For $q=n-1$ one has}
\begin{align}
E_2^{0,n-1}
 &=H^0(\mathbb P^1,R^{n-1}\rho_*\mathbb Q_\ell)
 \simeq H^{n-1}(X_{\bar\eta},\mathbb Q_\ell)^\pi
 \simeq H^{n-1}(X,\mathbb Q_\ell),
 \tag{5.6.2.4}\\
E_2^{2,n-1}
 &=H^2(\mathbb P^1,R^{n-1}\rho_*\mathbb Q_\ell)
 \simeq H^{n-1}(X_{\bar\eta},\mathbb Q_\ell(-1))^\pi
 \simeq H^{n-1}(X,\mathbb Q_\ell(-1)).
 \tag{5.6.2.5}
\end{align}
\underline{(The missing value $E_2^{1,n-1}$ will be found in 5.7 below.)}

\underline{Proof.}
By (5.4.3), Theorem 1.1 applies and gives (5.6.1).  The first three
formulas (5.6.2.1)--(5.6.2.3) follow from the constancy of
$R^q\rho_*\mathbb Q_\ell$ for $q\neq n-1$ (Lemma 5.4); for (5.6.2.3) one
also uses, of course, the simple connectedness of $\mathbb P^1$.

The first isomorphism in (5.6.2.4) follows from
\[
R^{n-1}\rho_*\mathbb Q_\ell
\overset{\sim}{\longrightarrow}
\nu_*\nu^*R^{n-1}\rho_*\mathbb Q_\ell.
\]
Thus
\[
H^0(\mathbb P^1,R^{n-1}\rho_*\mathbb Q_\ell)
\simeq
H^0(U,\nu^*R^{n-1}\rho_*\mathbb Q_\ell),
\]
and $\nu^*R^{n-1}\rho_*\mathbb Q_\ell$ is a
\textquotedblleft twisted constant\textquotedblright{} sheaf on $U$, hence
corresponds to the $\pi$-module
$H^{n-1}(X_{\bar\eta},\mathbb Q_\ell)$.  The second isomorphism in
(5.6.2.4) follows from degeneration together with Corollary 5.1.6.

We turn to (5.6.2.5).  Let $F^\bullet$ denote the filtration on
$H^\bullet(\widetilde X,\mathbb Q_\ell)$ defined by the spectral sequence.
By degeneration,
$F^2H^{n+1}(\widetilde X,\mathbb Q_\ell)=E_2^{2,n-1}$ and
$\operatorname{gr}^1H^{n+1}(\widetilde X,\mathbb Q_\ell)=E_2^{1,n}=0$
by (5.6.2.3).  Consequently (cf. 5.1.0),
\begin{align}
E_2^{2,n-1}
 &=F^1H^{n+1}(\widetilde X,\mathbb Q_\ell) \notag\\
 &=\operatorname{Ker}\!\left(
 H^{n+1}(\widetilde X,\mathbb Q_\ell)
 \longrightarrow H^0(\mathbb P^1,R^{n+1}\rho_*\mathbb Q_\ell)
 \right) \notag\\
 &\overset{(A)}{=}
 \operatorname{Ker}\!\left(
 k^*:H^{n+1}(\widetilde X,\mathbb Q_\ell)
 \longrightarrow H^{n+1}(X_{\bar\eta},\mathbb Q_\ell)
 \right).
\tag{5.6.3}
\end{align}
By (5.1.2), $k^*$ may be written as
\begin{equation}\tag{5.6.4}
H^{n+1}(X,\mathbb Q_\ell)
\oplus H^{n-1}(\Delta,\mathbb Q_\ell(-1))
\xrightarrow{\ m^*+h_*\ }
H^{n+1}(X_{\bar\eta},\mathbb Q_\ell).
\end{equation}
Using (5.2.2.3), this becomes
\begin{equation}\tag{5.6.5}
\begin{split}
&\operatorname{Prim}^{n-1}(X,\mathbb Q_\ell(-1))
\oplus H^{n-3}(X,\mathbb Q_\ell(-2))
\oplus H^{n-1}(\Delta,\mathbb Q_\ell(-1))\\
&\hspace{35mm}\xrightarrow{\ m^*L+m^*L^2+h_*\ }
H^{n+1}(X_{\bar\eta},\mathbb Q_\ell).
\end{split}
\end{equation}
\footnote{The printed (5.6.5) omits several closing parentheses and the plus
sign before the third-summand map $h_*$.  The three displayed source
summands and (5.6.4) force the corrected three-term map shown here.}
Now the map
\[
H^{n-3}(X,\mathbb Q_\ell(-2))
\xrightarrow{\ m^*L^2\ }
H^{n+1}(X_{\bar\eta},\mathbb Q_\ell)
\]
is an isomorphism by (LV), since
$m^*L_X^2=L_{X_{\bar\eta}}^2m^*$ (cf. 5.2.4 and 5.2.5).  Hence the kernel
of (5.6.5) identifies, by $\operatorname{pr}_1\oplus\operatorname{pr}_3$,
with
\[
\operatorname{Prim}^{n-1}(X,\mathbb Q_\ell(-1))
\oplus H^{n-1}(\Delta,\mathbb Q_\ell(-1)).
\]
By (LV) and weak Lefschetz,
\[
H^{n-1}(\Delta,\mathbb Q_\ell(-1))
\xleftarrow[\sim]{\,L\,}
H^{n-3}(\Delta,\mathbb Q_\ell(-2))
\xleftarrow[\sim]{\ }
H^{n-3}(X,\mathbb Q_\ell(-2)).
\]
Together with (5.2.2.3), this gives
\begin{equation}\tag{5.6.6}
E_2^{2,n-1}\simeq H^{n-1}(X,\mathbb Q_\ell(-1)).
\end{equation}
The other isomorphism in (5.6.2.5) follows by twisting (5.6.2.4).

\subsubsection*{Remark 5.6.7.}
The assertion (5.6.1) can be proved under substantially weaker hypotheses,
as Deligne explained to me:
\end{filecontents*}
% Original input: components/121_expose_XVIII_theorem_5_6_8_first_proof.tex
\begin{filecontents*}[overwrite,noheader]{sga7ii-cumulative-input-121.tex}
\subsection*{Theorem 5.6.8.}
\underline{Let $X$ be smooth of dimension $n$ over an algebraically closed
field $k$, and let}
\[
\rho:X\longrightarrow\mathbb P^1_k
\]
\underline{be a projective morphism whose generic fiber is smooth.  Suppose
that:}

\smallskip
\noindent
\textup{(5.6.8.1)}\quad
\underline{the geometric generic fiber, equipped with the restriction of the
relatively ample invertible sheaf $\underline{\mathcal O}_X(1)$, satisfies
(LV);}

\smallskip
\noindent
\textup{(5.6.8.2)}\quad
\underline{the maps (5.3.2) are isomorphisms for $i\neq n$ and are
surjective for $i=n$ (a condition satisfied in particular when $\rho$ arises
from a Lefschetz pencil, cf. 5.3.4).  Then the Leray spectral sequence}
\[
E_2^{p,q}=H^p(\mathbb P^1,R^q\rho_*\mathbb Q_\ell)
\Longrightarrow H^{p+q}(X,\mathbb Q_\ell)
\]
\underline{degenerates.}

\smallskip
\noindent
\underline{Remark.}
The same theorem, with the same proof, holds with $\mathbb P^1_k$ replaced
by any smooth, irreducible, projective curve.

\underline{Proof.}
For the proof we identify the sheaf $\mathbb Q_\ell(1)$ with
$\mathbb Q_\ell$; this is possible because $k$ is algebraically closed.
We may plainly suppose that $X$ is connected and that $\rho$ is
\underline{surjective}, hence flat with fibers of dimension $n-1$.  Let
\[
\xi_X\in H^2(X,\mathbb Q_\ell)
\]
be the first Chern class of $\underline{\mathcal O}_X(1)$, and let
\[
\nu:U\hookrightarrow\mathbb P^1
\]
be the inclusion of an open subset over which $\rho$ is smooth.  By
(5.6.8.1) and Lemma 5.2.6, the cup-product maps
\[
\xi_X^i:
\nu^*R^{n-1-i}\rho_*\mathbb Q_\ell
\longrightarrow
\nu^*R^{n-1+i}\rho_*\mathbb Q_\ell
\]
are isomorphisms for $i\geq1$.  Applying $\nu_*$ and using (5.6.8.2) gives,
for $i\geq2$, the isomorphisms
\begin{equation}\tag{5.6.8.3}
\xi_X^i:R^{n-1-i}\rho_*\mathbb Q_\ell
\overset{\sim}{\longrightarrow}
R^{n-1+i}\rho_*\mathbb Q_\ell,
\end{equation}
and, for $i=1$, the isomorphism
\begin{equation}\tag{5.6.8.4}
\xi_X:R^{n-2}\rho_*\mathbb Q_\ell
\overset{\sim}{\longrightarrow}
\nu_*\nu^*R^n\rho_*\mathbb Q_\ell.
\end{equation}
The latter factors as
\begin{equation}\tag{5.6.8.5}
R^{n-2}\rho_*\mathbb Q_\ell
\xrightarrow{\ \xi_X\ }
R^n\rho_*\mathbb Q_\ell
\xrightarrow{\ \text{adjunction}\ }
\nu_*\nu^*R^n\rho_*\mathbb Q_\ell.
\end{equation}
Consequently there is a canonical decomposition
\begin{equation}\tag{5.6.8.6}
R^n\rho_*\mathbb Q_\ell
\xleftarrow[\sim]{\ \xi_X+\mathrm{incl}\ }
R^{n-2}\rho_*\mathbb Q_\ell\oplus
\operatorname{Local}^n\mathbb Q_\ell,
\end{equation}
where
\begin{equation}\tag{5.6.8.7}
\operatorname{Local}^n\mathbb Q_\ell
=\operatorname{Ker}\!\left(
R^n\rho_*\mathbb Q_\ell
\longrightarrow\nu_*\nu^*R^n\rho_*\mathbb Q_\ell
\right).
\end{equation}
The sheaf $\operatorname{Local}^n\mathbb Q_\ell$ is concentrated at the
points of $\mathbb P^1$ where the fiber of $\rho$ is not smooth.  Combining
this fact with (5.6.8.3), we obtain
\begin{equation}\tag{5.6.8.8}
\xi_X^{i+1}:
H^j(\mathbb P^1,R^{n-2-i}\rho_*\mathbb Q_\ell)
\overset{\sim}{\longrightarrow}
H^j(\mathbb P^1,R^{n+i}\rho_*\mathbb Q_\ell),
\qquad i\geq0,\quad j\geq1.
\end{equation}
\footnote{The printed (5.6.8.8) has an odd degree shift, labels the map only
$\xi_X$, and starts at $i=1$.  The corrected form is forced by
(5.6.8.3), (5.6.8.6), the lower arrow of the next commutative square, and
the later specialization $E_2^{1,n-2}\simeq E_2^{1,n}$.}

We now prove degeneration.  Since $E_2^{p,q}=0$ unless $p=0,1,2$, it is
enough to prove that all the differentials
\[
d_2:H^0(\mathbb P^1,R^i\rho_*\mathbb Q_\ell)
\longrightarrow H^2(\mathbb P^1,R^{i-1}\rho_*\mathbb Q_\ell)
\]
vanish.  For $q\leq n-1$, set
\begin{equation}\tag{5.6.8.9}
\operatorname{Prim}(E_2^{p,q})
=\operatorname{Ker}\!\left(
\xi_X^{n-q}:E_2^{p,q}\longrightarrow E_2^{p,2n-q}
\right).
\end{equation}
Equations (5.6.8.3) and (5.6.8.6), in particular the injectivity of
\[
\xi_X:H^p(\mathbb P^1,R^{n-2}\rho_*\mathbb Q_\ell)
\longrightarrow H^p(\mathbb P^1,R^n\rho_*\mathbb Q_\ell),
\]
give, for $q\leq n-1$, the decomposition
\begin{equation}\tag{5.6.8.10}
E_2^{p,q}
\xleftarrow[\sim]{\ 1+\xi_X+\xi_X^2+\cdots\ }
\operatorname{Prim}(E_2^{p,q})
\oplus\operatorname{Prim}(E_2^{p,q-2})
\oplus\operatorname{Prim}(E_2^{p,q-4})
\oplus\cdots.
\end{equation}
\footnote{The print displays the first three primitive summands while its
arrow label already contains an ellipsis; the final $\oplus\cdots$ makes the
same finite decomposition explicit.}
\end{filecontents*}
% Original input: components/122_expose_XVIII_theorem_5_6_8_completion_and_corollary.tex
\begin{filecontents*}[overwrite,noheader]{sga7ii-cumulative-input-122.tex}
To prove that
\[
d_2^{p,q}=0\qquad(q\leq n-1),
\]
consider the commutative diagram
\[
\begin{tikzcd}[column sep=huge,row sep=huge]
\operatorname{Prim}(E_2^{p,q})
  \arrow[r,"\xi_X^{n-q}=0"]
  \arrow[d,"d_2^{p,q}"]
&
E_2^{p,2n-q}
  \arrow[d,"d_2^{p,2n-q}"]
\\
E_2^{p+2,q-1}
  \arrow[r,"\xi_X^{n-1-(q-1)}","\sim"']
&
E_2^{p+2,2n-q-1}.
\end{tikzcd}
\]
The lower arrow is an isomorphism by (5.6.8.8).  The diagram therefore shows
that $d_2^{p,q}$ kills $\operatorname{Prim}(E_2^{p,q})$ for $q\leq n-1$,
and (5.6.8.10) then shows that $d_2$ kills $E_2^{p,q}$ for $q\leq n-1$.
By (5.6.8.3), it also kills $E_2^{p,q}$ for $q\geq n+1$.  It remains only
to prove that
\[
d_2^{0,n}:E_2^{0,n}\longrightarrow E_2^{2,n-1}
\]
is zero, or equivalently that
\begin{equation}\tag{5.6.8.12}
\dim E_3^{2,n-1}=\dim E_2^{2,n-1}.
\end{equation}
We have already proved
\begin{equation}\tag{5.6.8.13}
\dim E_2^{p,q}=\dim E_3^{p,q}
\qquad\text{for }(p,q)\neq(0,n),(2,n-1),
\end{equation}
and for degree reasons
\begin{equation}\tag{5.6.8.14}
E_3^{p,q}=E_\infty^{p,q}
\qquad\text{for every pair }(p,q).
\end{equation}
It follows that
\[
b_i(X)=\dim H^i(X,\mathbb Q_\ell)
=\sum_{p+q=i}\dim E_3^{p,q}.
\]
In view of (5.6.8.13), we therefore have
\begin{equation}\tag{5.6.8.15}
b_{n-1}(X)
=\dim E_2^{0,n-1}
 +\dim E_2^{1,n-2}
 +\dim E_2^{2,n-3}
\end{equation}
and
\begin{equation}\tag{5.6.8.16}
b_{n+1}(X)
=\dim E_2^{0,n+1}
 +\dim E_2^{1,n}
 +\dim E_3^{2,n-1}.
\end{equation}
Poincaré duality gives
\begin{equation}\tag{5.6.8.17}
b_{n-1}(X)=b_{n+1}(X).
\end{equation}
First, by (5.6.8.8),
\[
\dim E_2^{1,n}=\dim E_2^{1,n-2}.
\]
Thus (5.6.8.17) shows that the proof of (5.6.8.12) is complete once we know
\begin{align}
\dim E_2^{2,n-3}&=\dim E_2^{0,n+1},
\tag{5.6.8.18}\\
\dim E_2^{0,n-1}&=\dim E_2^{2,n-1}.
\tag{5.6.8.19}
\end{align}
These two formulas follow from (5.6.8.2), Poincaré duality on the smooth
fibers of $\rho$, and the following lemma applied to $\mathbb P^1$ and the
sheaves $R^{n-3}\rho_*\mathbb Q_\ell$ and
$R^{n-1}\rho_*\mathbb Q_\ell$.

\subsection*{Lemma 5.6.9.}
\underline{Let $D$ be a smooth, connected, proper curve over an
algebraically closed field $k$.  Let $i:\eta\hookrightarrow D$ be the
inclusion of the generic point, and let $F$ be a constructible
$\mathbb Q_\ell$-sheaf on $D$ such that
$F\overset{\sim}{\longrightarrow}i_*i^*F$.  Then}
\[
H^2(D,F)\simeq H^0(D,\check F(-1))^\vee,
\]
\underline{where $\check F\simeq i_*((i^*F)^\vee)$.}

\subsection*{Corollary 5.6.10.}
\textup{(Uses Picard--Lefschetz.)}
\underline{Let $\rho:\widetilde X\to\mathbb P^1$ be the
\textquotedblleft fibration\textquotedblright{} arising from a Lefschetz
pencil, and suppose that the smooth fibers satisfy (LV).  Then}
\begin{equation}\tag{5.6.10.1}
H^i(X_{\bar\eta},\mathbb Q_\ell)^\pi
=\operatorname{Im}\!\left(
H^i(X,\mathbb Q_\ell)\longrightarrow
H^i(X_{\bar\eta},\mathbb Q_\ell)
\right).
\end{equation}

\underline{Proof.}
By 5.3.4, which uses the Picard--Lefschetz formula, Theorem 5.6.8 applies
and gives
\[
H^i(X_{\bar\eta},\mathbb Q_\ell)^\pi
=\operatorname{Im}\!\left(
H^i(\widetilde X,\mathbb Q_\ell)\longrightarrow
H^i(X_{\bar\eta},\mathbb Q_\ell)
\right).
\]
It remains to check that $H^i(X,\mathbb Q_\ell)$ and
$H^i(\widetilde X,\mathbb Q_\ell)$ have the same image in
$H^i(X_{\bar\eta},\mathbb Q_\ell)$.  This holds for $i\leq n-1$ by
Corollary 5.1.6, and for $i\geq n$ by Corollary 5.2.5.

We now turn to the most important computation:
\end{filecontents*}
% Original input: components/123_expose_XVIII_theorem_5_7.tex
\begin{filecontents*}[overwrite,noheader]{sga7ii-cumulative-input-123.tex}
\subsection*{Theorem 5.7.}
\underline{Under the hypotheses of 5.6, set}
\[
\begin{aligned}
E^{n-2}(\Delta,\mathbb Q_\ell(j))
&=\underline{\text{the orthogonal complement in }}
H^{n-2}(\Delta,\mathbb Q_\ell(j))\\
&\qquad\underline{\text{of the image of }}
H^{n-2}(X,\mathbb Q_\ell(j)).
\end{aligned}
\]
\underline{(This is the \textquotedblleft vanishing\textquotedblright{}
part of the cohomology of $\Delta$; cf. 5.4.4.)  Then there is a
direct-sum decomposition} (cf. 4.2)
\begin{equation}\tag{5.7.1}
E_2^{1,n-1}
=H^1(\mathbb P^1,R^{n-1}\rho_*\mathbb Q_\ell)
\xleftarrow[\;f^*+i_*g^*\;]{\sim}
\operatorname{Prim}^n(X,\mathbb Q_\ell)
\oplus E^{n-2}(\Delta,\mathbb Q_\ell(-1)).
\end{equation}

\underline{Proof.}
By degeneration,
\[
E_2^{1,n-1}\simeq
F^1H^n(\widetilde X,\mathbb Q_\ell)
\big/F^2H^n(\widetilde X,\mathbb Q_\ell).
\]
By (A),
\begin{equation}\tag{5.7.2}
F^1H^n(\widetilde X,\mathbb Q_\ell)
=\ker\!\left(
k^*:H^n(\widetilde X,\mathbb Q_\ell)\longrightarrow
H^n(X_{\bar\eta},\mathbb Q_\ell)
\right),
\end{equation}
which, by (5.1.2), identifies with
\[
\ker\!\left(
H^n(X,\mathbb Q_\ell)
\oplus H^{n-2}(\Delta,\mathbb Q_\ell(-1))
\xrightarrow{\;m^*+h_*\;}
H^n(X_{\bar\eta},\mathbb Q_\ell)
\right).
\]
Under (LV), $k^*$ can also be written as
\begin{equation}\tag{5.7.3}
\operatorname{Prim}^n(X,\mathbb Q_\ell)
\oplus H^{n-2}(X,\mathbb Q_\ell(-1))
\oplus H^{n-2}(\Delta,\mathbb Q_\ell(-1))
\xrightarrow{\;m^*+m^*L+h_*\;}
H^n(X_{\bar\eta},\mathbb Q_\ell),
\end{equation}
and the middle arrow
\begin{equation}\tag{5.7.4}
H^{n-2}(X,\mathbb Q_\ell(-1))
\xrightarrow{\;m^*L\;}
H^n(X_{\bar\eta},\mathbb Q_\ell)
\end{equation}
is an isomorphism (cf. 5.2.3 and (5.2.2.3)).  Thus, denoting by
\begin{equation}\tag{5.7.5}
\wp:H^n(X,\mathbb Q_\ell)\longrightarrow
\operatorname{Prim}^n(X,\mathbb Q_\ell)
\end{equation}
the projection onto the primitive part, we obtain
\footnote{The printed label in (5.7.6) has $2^*$ after $g_*$.
The inverse blow-up decomposition (4.2.2), used here verbatim, requires
$i^*$.}
\begin{equation}\tag{5.7.6}
F^1H^n(\widetilde X,\mathbb Q_\ell)
\xrightarrow[\;\wp f_*\oplus(-g_*i^*)\;]{\sim}
\operatorname{Prim}^n(X,\mathbb Q_\ell)
\oplus H^{n-2}(\Delta,\mathbb Q_\ell(-1)).
\end{equation}
By (5.6.2.2) and degeneration,
\begin{equation}\tag{5.7.7}
F^2H^n(\widetilde X,\mathbb Q_\ell)
\simeq E_2^{2,n-2}
\simeq H^{n-2}(X_{\bar\eta},\mathbb Q_\ell(-1)).
\end{equation}

On the other hand, the inclusion of
$F^2H^n(\widetilde X,\mathbb Q_\ell)
=H^{n-2}(X_{\bar\eta},\mathbb Q_\ell(-1))$
in $H^n(\widetilde X,\mathbb Q_\ell)$ is the Gysin map $k_*$.
\footnote{The print has $F^2H^N$ in the first sentence and omits the
Tate twist $(-1)$ from the displayed source of the Gysin map.  The
surrounding identifications and the degree-two Gysin shift force the readings
shown here.}
Strictly speaking, the Gysin map is defined only as a map
\[
H^{i-2}(X_{\bar\eta},\mathbb Q_\ell(-1))
\longrightarrow
H^i(X\times_k k(\bar\eta),\mathbb Q_\ell),
\]
but, in view of the canonical isomorphism (SGA 4 XII 5.4)
\[
H^i(X,\mathbb Q_\ell)\xrightarrow{\sim}
H^i(X\otimes_kK,\mathbb Q_\ell)
\]
for every algebraically closed field $K/k$, we hope the reader will allow
us such abuses.
By (5.1.3), the Gysin map is expressed as
\begin{equation}\tag{5.7.8}
H^{n-2}(X_{\bar\eta},\mathbb Q_\ell(-1))
\xrightarrow{\;m_*\oplus(-h^*)\;}
H^n(X,\mathbb Q_\ell)
\oplus H^{n-2}(\Delta,\mathbb Q_\ell(-1)).
\end{equation}
By weak Lefschetz,
\[
m_*H^{n-2}(X_{\bar\eta},\mathbb Q_\ell(-1))
=m_*m^*H^{n-2}(X,\mathbb Q_\ell(-1))
=L_XH^{n-2}(X,\mathbb Q_\ell(-1))
\]
is the orthogonal complement of the primitive part.  Hence we have the
commutative diagram
\begin{equation}\tag{5.7.9}
\begin{tikzcd}[column sep=huge,row sep=huge]
F^2H^n(\widetilde X,\mathbb Q_\ell)
  \arrow[r,hook]
& F^1H^n(\widetilde X,\mathbb Q_\ell)
  \arrow[d,"\wp f_*\oplus(-g_*i^*)","\wr"']
\\
H^{n-2}(X_{\bar\eta},\mathbb Q_\ell(-1))
  \arrow[u,"k_*"',"\wr"]
  \arrow[r,"0\oplus h^*"]
& \operatorname{Prim}^n(X,\mathbb Q_\ell)
  \oplus H^{n-2}(\Delta,\mathbb Q_\ell(-1)).
\end{tikzcd}
\end{equation}
This proves (5.7.1), as required.

\subsection*{Remark 5.7.10.}
In 5.7, hypothesis (A) may be replaced by the weaker assumption that, for
$q=n$, the maps
\begin{equation}\tag{5.7.11}
H^q(\widetilde X,\mathbb Q_\ell)\longrightarrow
H^0(\mathbb P^1,R^q\rho_*\mathbb Q_\ell)
\end{equation}
and
\begin{equation}\tag{5.7.12}
H^q(\widetilde X,\mathbb Q_\ell)\longrightarrow
H^q(X_{\bar\eta},\mathbb Q_\ell)
\end{equation}
have the same kernel, as our proof shows.  Without this condition, the
conclusion can fail; cf. 5.9.
\end{filecontents*}
% Original input: components/124_expose_XVIII_section_5_8_1_through_definition_5_8_3.tex
\begin{filecontents*}[overwrite,noheader]{sga7ii-cumulative-input-124.tex}
\section*{5.8. The Griffiths homomorphism}

\subsection*{5.8.1.}
Let $f:T\to S$ be a morphism of schemes.  The Leray spectral sequence
defines a filtration $(F^i)$ on the groups
$H^\bullet(T,\mathbb Q_\ell(j))$.  Define
\begin{equation}\tag{5.8.1.1}
\operatorname{Prim}^\bullet(T/S,\mathbb Q_\ell(j))
=F^1H^\bullet(T,\mathbb Q_\ell(j)).
\end{equation}
Because the Leray spectral sequence lies in the first quadrant, we have
\begin{equation}\tag{5.8.1.2}
E_\infty^{0,\bullet}\lhook\joinrel\longrightarrow E_2^{0,\bullet},
\end{equation}
and hence
\begin{equation}\tag{5.8.1.3}
\operatorname{Prim}^\bullet(T/S,\mathbb Q_\ell(j))
=\ker\!\left(
H^\bullet(T,\mathbb Q_\ell(j))
\longrightarrow H^0(S,R^\bullet f_*\mathbb Q_\ell(j))
\right).
\end{equation}
For the same reason,
\begin{equation}\tag{5.8.1.4}
E_\infty^{1,\bullet}\lhook\joinrel\longrightarrow E_2^{1,\bullet}.
\end{equation}
The composite
\begin{equation}\tag{5.8.1.5}
\begin{tikzcd}[column sep=large,row sep=large]
\operatorname{Prim}^i(T/S,\mathbb Q_\ell(j))
=F^1H^i(T,\mathbb Q_\ell(j))
  \arrow[r]
&
\operatorname{gr}^1H^i(T,\mathbb Q_\ell(j))
=E_\infty^{1,i-1}
  \arrow[d,hook]
\\
{}
&
H^1(S,R^{i-1}f_*\mathbb Q_\ell(j))
=E_2^{1,i-1}
\end{tikzcd}
\end{equation}
will be denoted
\footnote{The print has a baseline dot in place of the degree index in
(5.8.1.6).  The source and target of (5.8.1.5) force $R^{i-1}f_*$.}
\begin{equation}\tag{5.8.1.6}
U^i:\operatorname{Prim}^i(T/S,\mathbb Q_\ell(j))
\longrightarrow H^1(S,R^{i-1}f_*\mathbb Q_\ell(j)).
\end{equation}

\subsection*{5.8.2.}
Now let $X$ be projective, smooth, and irreducible, and let
$D=\{H_\tau\}$ be a pencil of hyperplanes whose generic member and axis
meet $X$ transversely.  Let
$U\xrightarrow{\nu}\mathbb P^1$ be a nonempty open subset over which the
projection $\rho:\widetilde X\to\mathbb P^1$ is smooth, and put
$\widetilde X|_U=\widetilde X\times_{\mathbb P^1}U$.  We have the
commutative diagram
\begin{equation}\tag{5.8.2.1}
\begin{tikzcd}[column sep=large,row sep=large]
{}
&
\widetilde X
  \arrow[dl,"f"']
  \arrow[d,"\rho"]
&
\widetilde X|_U
  \arrow[l,"\alpha"']
  \arrow[d,"\rho|_U"]
\\
X
&
\mathbb P^1
&
U\arrow[l,"\nu"'].
\end{tikzcd}
\end{equation}
Since $\rho|_U:\widetilde X|_U\to U$ is proper and smooth, the
specialization theorem for cohomology (SGA 4 XVI 2.2) gives
\begin{equation}\tag{5.8.2.2}
\operatorname{Prim}^\bullet(\widetilde X|_U/U,\mathbb Q_\ell(j))
=\ker\!\left(
H^\bullet(\widetilde X|_U,\mathbb Q_\ell(j))
\longrightarrow H^\bullet(X_{\bar\eta},\mathbb Q_\ell(j))
\right).
\end{equation}
Consequently, by (5.2.3.1), we have:

\subsection*{Proposition 5.8.2.3.}
\underline{Under the hypotheses of 5.8.2, for
$x\in H^n(X,\mathbb Q_\ell(j))$, the following conditions are equivalent:}
\begin{equation}\tag{5.8.2.4}
x\in\operatorname{Prim}^n(X,\mathbb Q_\ell(j))
\quad\Longleftrightarrow\quad
\alpha^*f^*(x)\in
\operatorname{Prim}^n(\widetilde X|_U/U,\mathbb Q_\ell(j)).
\end{equation}
\underline{If, in addition, $\rho$ satisfies hypothesis (A) of 5.3, then}
\begin{equation}\tag{5.8.2.5}
x\in\operatorname{Prim}^n(X,\mathbb Q_\ell(j))
\quad\Longleftrightarrow\quad
f^*(x)\in
\operatorname{Prim}^n(\widetilde X/\mathbb P^1,\mathbb Q_\ell(j)).
\end{equation}

\underline{Proof.}
Formula (5.8.2.4) is proved.  For (5.8.2.5), hypothesis (A) implies that
\[
H^0(\mathbb P^1,R^n\rho_*\mathbb Q_\ell(j))
\xrightarrow[\nu^*]{\sim}
H^0(U,\nu^*R^n\rho_*\mathbb Q_\ell(j)).
\]
Thus $f^*(x)$ is primitive relative to $\mathbb P^1$ if and only if
$\alpha^*f^*(x)$ is primitive relative to $U$.
\footnote{The last printed implication has a stray $\alpha$ after
$\mathbb P^1$ and substitutes $\alpha$ for the class $x$ in
$\alpha^*f^*(x)$.  The displayed equivalence and the preceding restriction
isomorphism determine the corrected sentence.}

\subsection*{Definition 5.8.3.}
\underline{In the situation of 5.8.2, the composite homomorphism}
\footnote{The ink beside the vertical arrow is not independently legible in
the scan.  The map is nevertheless determined uniquely: (5.8.1.6) has just
defined $U^i$ with exactly this source and target, so the degree-$n$ label is
$U^n$.}
\begin{equation}\tag{5.8.3.1}
\begin{tikzcd}[column sep=huge,row sep=huge]
\operatorname{Prim}^n(X,\mathbb Q_\ell(j))
  \arrow[r,"\alpha^*f^*"]
  \arrow[dr,"\operatorname{Grif}"']
&
\operatorname{Prim}^n(\widetilde X|_U/U,\mathbb Q_\ell(j))
  \arrow[d,"U^n"]
\\
{}
&
H^1(U,\nu^*R^{n-1}\rho_*\mathbb Q_\ell(j))
\end{tikzcd}
\end{equation}
\underline{is called the Griffiths homomorphism (for $X$, relative to the
given pencil $D$).}

Adjoin to diagram (5.8.2.1) the constant family over $U$ with fiber $X$:
\begin{equation}\tag{5.8.3.2}
\begin{tikzcd}[row sep=large,column sep=large]
{}
&
\widetilde X
  \arrow[dl,"f"']
  \arrow[dd]
&
\widetilde X|_U
  \arrow[l,"\alpha"']
  \arrow[r,"f\times\rho|_U"]
  \arrow[dd,"\rho|_U"]
&
X\times U
  \arrow[ddl,"\operatorname{pr}_2"]
\\
X&{}&{}&{}
\\
{}
&
\mathbb P^1
&
U\arrow[l,"\nu"']
&
.
\end{tikzcd}
\end{equation}

\end{filecontents*}
% Original input: components/125_expose_XVIII_theorem_5_8_4_and_lemma_5_8_5.tex
\begin{filecontents*}[overwrite,noheader]{sga7ii-cumulative-input-125.tex}
\subsection*{Theorem 5.8.4.}
\underline{Suppose that $X$ satisfies (LV) (5.2.2), and that $\rho$
satisfies hypothesis (A) (5.3).  Then, for every nonempty smoothness open
$U\xrightarrow{\nu}\mathbb P^1$ of $\rho$, the Griffiths homomorphism
\textquotedblleft modulo the fixed part\textquotedblright{}, which is by
definition the composite}
\begin{equation}\tag{5.8.4.1}
\resizebox{\textwidth}{!}{%
\begin{tikzcd}[ampersand replacement=\&,row sep=huge,column sep=huge]
\operatorname{Prim}^n(X,\mathbb Q_\ell(j))
  \arrow[r,"\operatorname{Grif}"]
\&
H^1(U,\nu^*R^{n-1}\rho_*\mathbb Q_\ell(j))
  \arrow[d]
\\
{}
\&
H^1(U,\nu^*E^{n-1}\rho_*\mathbb Q_\ell(j))
\simeq
\dfrac{H^1(U,\nu^*R^{n-1}\rho_*\mathbb Q_\ell(j))}
{(f\times\rho|_U)^*
 H^1(U,H^{n-1}(X,\mathbb Q_\ell(j))_U)}
\end{tikzcd}}
\end{equation}
\underline{is injective.}
\footnote{The print omits $\nu^*$ from the lower-left sheaf in (5.8.4.1)
and leaves its coefficient parenthesis unclosed.  Formula (5.8.4.3) and the
same lower node in (5.8.4.4) give the corrected reading.}

\underline{Proof.}
By (5.5.1), over $\mathbb P^1$ we have
\begin{equation}\tag{5.8.4.2}
R^{n-1}\rho_*\mathbb Q_\ell(j)
\simeq
H^{n-1}(X,\mathbb Q_\ell(j))_{\mathbb P^1}
\oplus E^{n-1}\rho_*\mathbb Q_\ell(j),
\end{equation}
and over $U$ we have
\begin{equation}\tag{5.8.4.3}
\nu^*R^{n-1}\rho_*\mathbb Q_\ell(j)
\simeq
H^{n-1}(X,\mathbb Q_\ell(j))_U
\oplus \nu^*E^{n-1}\rho_*\mathbb Q_\ell(j).
\end{equation}
This permits us to rewrite (5.8.4.1) as indicated:
\begin{equation}\tag{5.8.4.4}
\begin{tikzcd}[row sep=huge,column sep=large]
\operatorname{Prim}^n(X,\mathbb Q_\ell(j))
  \arrow[r,"\operatorname{Grif}"]
&
H^1(U,\nu^*R^{n-1}\rho_*\mathbb Q_\ell(j))
  \arrow[d]
\\
{}
&
H^1(U,\nu^*E^{n-1}\rho_*\mathbb Q_\ell(j)).
\end{tikzcd}
\end{equation}
By (5.7.1), the map (defined only by virtue of hypothesis (A))
\begin{equation}\tag{5.8.4.5}
u:\operatorname{Prim}^n(X,\mathbb Q_\ell(j))
\longrightarrow
H^1(\mathbb P^1,R^{n-1}\rho_*\mathbb Q_\ell(j))
\end{equation}
is injective.  Since $\mathbb P^1$ is simply connected, decomposition
(5.8.4.2) gives
\footnote{The print omits the twist $(j)$ from the targets in (5.8.4.5)
and from the left-hand side of (5.8.4.6), although the domain and the
decomposition (5.8.4.2) retain it.}
\begin{equation}\tag{5.8.4.6}
H^1(\mathbb P^1,R^{n-1}\rho_*\mathbb Q_\ell(j))
\simeq
H^1(\mathbb P^1,E^{n-1}\rho_*\mathbb Q_\ell(j)).
\end{equation}
We can therefore complete (5.8.4.4) to a commutative square:
\begin{equation}\tag{5.8.4.7}
\begin{tikzcd}[column sep=huge,row sep=huge]
\operatorname{Prim}^n(X,\mathbb Q_\ell(j))
  \arrow[r,"\operatorname{Grif}"]
  \arrow[d,hook,"u"]
&
H^1(U,\nu^*R^{n-1}\rho_*\mathbb Q_\ell(j))
  \arrow[d]
\\
H^1(\mathbb P^1,E^{n-1}\rho_*\mathbb Q_\ell(j))
  \arrow[r,"\nu^*"]
&
H^1(U,\nu^*E^{n-1}\rho_*\mathbb Q_\ell(j)).
\end{tikzcd}
\end{equation}
It remains only to see that the lower arrow is injective.  This follows from
\begin{equation}\tag{5.8.4.8}
E^{n-1}\rho_*\mathbb Q_\ell(j)
\simeq
\nu_*\nu^*E^{n-1}\rho_*\mathbb Q_\ell(j)
\end{equation}
and the following well-known lemma.

\subsection*{Lemma 5.8.5.}
\underline{Let $\nu:A\to B$ be a morphism of topoi, and let $F$ be a
sheaf on $A$.  Then the canonical map}
\begin{equation}\tag{5.8.5.1}
H^1(B,\nu_*F)\longrightarrow H^1(A,F)
\end{equation}
\underline{is injective.}

\underline{Proof.}
In a first-quadrant spectral sequence,
\begin{equation}\tag{5.8.5.2}
E_2^{1,0}\simeq E_\infty^{1,0}.
\end{equation}
Thus, for the Leray spectral sequence
\[
E_2^{p,q}=H^p(B,R^q\nu_*F)
\Longrightarrow H^{p+q}(A,F),
\]
we have
\begin{equation}\tag{5.8.5.3}
E_2^{1,0}=H^1(B,\nu_*F)
\simeq E_\infty^{1,0}
=F^1H^1(A,F).
\end{equation}
Taking $F=\nu^*E$, where $E$ is a sheaf on $B$, gives
\begin{equation}\tag{5.8.5.4}
H^1(B,\nu_*\nu^*E)
\lhook\joinrel\longrightarrow
H^1(A,\nu^*E).
\end{equation}
\end{filecontents*}
% Original input: components/126_expose_XVIII_section_5_8_6_and_corollary_5_8_7.tex
\begin{filecontents*}[overwrite,noheader]{sga7ii-cumulative-input-126.tex}
\subsection*{5.8.6.}
Let $L/k(\eta)$ be a finite extension of the rational function field of
$\mathbb P^1$, and put
\begin{equation}\tag{5.8.6.1}
S=\text{the normalization of $\mathbb P^1$ in $L$}.
\end{equation}
Make the base change $S\to\mathbb P^1$ in situation (5.8.3.2):
\begin{equation}\tag{5.8.6.1 bis}
\begin{tikzcd}[column sep=huge,row sep=large]
{} & \widetilde X \arrow[dl,"f"] \arrow[dd,"\rho"]
& \widetilde X_S \arrow[l,"\alpha_S"'] \arrow[dd,"\rho_S"] \\
X&{}&{} \\
{} & \mathbb P^1 & S\arrow[l,"\pi"'] .
\end{tikzcd}
\end{equation}
For every open immersion $V\xhookrightarrow{\mu}S$ over which $\rho_S$
is smooth, base change by $V\xrightarrow{\pi\mu}\mathbb P^1$ in
(5.8.3.2) gives
\begin{equation}\tag{5.8.6.2}
\begin{tikzcd}[column sep=huge,row sep=large]
{} & \widetilde X \arrow[dl,"f"] \arrow[dd,"\rho"]
& \widetilde X_V \arrow[l,"\alpha_V"'] \arrow[dd,"\rho_V"]
  \arrow[dr,"f\times\rho_V"] &{} \\
X&{}&{}&X\times V\arrow[dl,"\operatorname{pr}_2"] \\
{} & \mathbb P^1 & V\arrow[l,"\pi\mu"'] &{} .
\end{tikzcd}
\end{equation}
We still have (cf. (5.8.2.3))
\begin{equation}\tag{5.8.6.3}
\alpha_V^*f^*\operatorname{Prim}^n(X,\mathbb Q_\ell(j))
\subset
\operatorname{Prim}^n(\widetilde X_V/V,\mathbb Q_\ell(j)),
\end{equation}
so that we may again speak of the Griffiths homomorphism
\begin{equation}\tag{5.8.6.4}
\operatorname{Prim}^n(X,\mathbb Q_\ell(j))
\xrightarrow{\operatorname{Grif}}
H^1(V,R^{n-1}\rho_{V*}\mathbb Q_\ell(j)).
\end{equation}

\subsection*{Corollary 5.8.7.}
\underline{With the preceding notation and under the hypotheses of 5.8.4,
the Griffiths homomorphism ``modulo the fixed part,'' that is, the
composite}
\begin{equation}\tag{5.8.7.1}
\resizebox{\textwidth}{!}{%
\begin{tikzcd}[ampersand replacement=\&,column sep=huge,row sep=huge]
\operatorname{Prim}^n(X,\mathbb Q_\ell(j))
  \arrow[r,"\operatorname{Grif}"]
\& H^1(V,R^{n-1}\rho_{V*}\mathbb Q_\ell(j)) \arrow[d] \\
{}
\& H^1(V,(\pi\mu)^*E^{n-1}\rho_*\mathbb Q_\ell(j))
\simeq
\dfrac{H^1(V,R^{n-1}\rho_{V*}\mathbb Q_\ell(j))}
{(f\times\rho_V)^*H^1(V,H^{n-1}(X,\mathbb Q_\ell(j))_V)}
\end{tikzcd}}
\end{equation}
\underline{is injective.}

\underline{Proof.}
By the specialization theorem for cohomology,
\[
R^{n-1}\rho_{V*}\mathbb Q_\ell(j)
=(\pi\mu)^*R^{n-1}\rho_*\mathbb Q_\ell(j).
\]
Thus decomposition (5.8.4.2) rewrites (5.8.7.1) as indicated:
\[
\begin{tikzcd}[column sep=huge,row sep=huge]
\operatorname{Prim}^n(X,\mathbb Q_\ell(j))
  \arrow[r,"\operatorname{Grif}"]
& H^1(V,(\pi\mu)^*R^{n-1}\rho_*\mathbb Q_\ell(j)) \arrow[d] \\
{}
& H^1(V,(\pi\mu)^*E^{n-1}\rho_*\mathbb Q_\ell(j)).
\end{tikzcd}
\]
It is therefore enough to prove the injectivity of
\begin{equation}\tag{5.8.7.3}
\operatorname{Prim}^n(X,\mathbb Q_\ell(j))
\longrightarrow
H^1(V,\mu^*\pi^*E^{n-1}\rho_*\mathbb Q_\ell(j))
\end{equation}
for arbitrarily small open immersions $V\xhookrightarrow{\mu}S$.  We are
reduced to the case
\[
\left\{
\begin{aligned}
U&\xhookrightarrow{\nu}\mathbb P^1
   &&\text{is an open subset over which $\rho$ is smooth},\\
V&\xrightarrow{\pi_U}U
   &&\text{is the normalization of $U$ in $L$}.
\end{aligned}
\right.
\]
From the cartesian square
\[
\begin{tikzcd}[column sep=huge,row sep=huge]
V\arrow[r,hook,"\mu"]\arrow[d,"\pi_U"']
&S\arrow[d,"\pi"]\\
U\arrow[r,hook,"\nu"]
&\mathbb P^1
\end{tikzcd}
\]
we obtain
\[
\mu^*\pi^*E^{n-1}\rho_*\mathbb Q_\ell(j)
=\pi_U^*\nu^*E^{n-1}\rho_*\mathbb Q_\ell(j).
\]
Consequently, (5.8.7.3) factors as
\[
\begin{tikzcd}[column sep=huge,row sep=huge]
\operatorname{Prim}^n(X,\mathbb Q_\ell(j)) \arrow[r]
&H^1(U,\nu^*E^{n-1}\rho_*\mathbb Q_\ell(j))
  \arrow[d,"\pi_U^*"] \\
{}
&H^1(V,\pi_U^*\nu^*E^{n-1}\rho_*\mathbb Q_\ell(j)).
\end{tikzcd}
\]
By (5.8.4), it is enough to verify that
\[
\pi_U^*:H^1(U,\nu^*E^{n-1}\rho_*\mathbb Q_\ell(j))
\longrightarrow
H^1(V,\pi_U^*\nu^*E^{n-1}\rho_*\mathbb Q_\ell(j))
\]
is injective.  Since $\pi_U:V\to U$ is finite and flat, trace theory
[SGA 4 XVII 6.2] supplies, for every $\mathbb Q_\ell$-sheaf $F$ on $U$
(for example $F=\nu^*E^{n-1}\rho_*\mathbb Q_\ell(j)$), a map
\[
\operatorname{tr}_{\pi_U}:H^1(V,\pi_U^*F)\longrightarrow H^1(U,F)
\]
such that
\[
\operatorname{tr}_{\pi_U}\circ\pi_U^*
=\deg(\pi_U)\,\operatorname{id}.
\]
It follows that $H^*(U,F)\to H^*(V,\pi_U^*F)$ is injective, completing
the proof.
\footnote{The print drops the twist $(j)$ from one base-change display and
the frozen transcription corrupts both the source and target of the trace
map.  The corrected formulas are the standard finite-flat pullback--trace
identities used by the concluding argument.}
\end{filecontents*}
% Original input: components/127_expose_XVIII_section_5_9.tex
\begin{filecontents*}[overwrite,noheader]{sga7ii-cumulative-input-127.tex}
\subsection*{5.9. \underline{An example in which hypothesis} (A)
\underline{is not satisfied}.}

Let $X$ be a smooth quadric surface in $\mathbb P^3$.  Then $X$ satisfies
(LV) (5.2.4).  Take a pencil (a Lefschetz pencil, if desired) of
\underline{hyperplanes} whose generic member and axis meet $X$
transversely.  Then (as Deligne pointed out to me) \underline{the projection}
\[
\rho:\widetilde X\longrightarrow\mathbb P^1
\]
\underline{does not satisfy} (A).  Indeed,
$R^2\rho_*\mathbb Q_\ell$ is not constant: the singular fibers are singular
conics, hence reducible, and in a Lefschetz pencil each is the union of two
lines meeting at one point (this is the definition of a Lefschetz pencil in
this simple case).  Thus
\[
\dim_{\mathbb Q_\ell}H^2(\text{singular fiber},\mathbb Q_\ell)=2,
\]
whereas
\[
\dim_{\mathbb Q_\ell}H^2(\text{smooth fiber},\mathbb Q_\ell)=1.
\]

Moreover, the conclusion of Theorem 5.8.4 is false here.  Let
$U\xhookrightarrow{\nu}\mathbb P^1$ be a nonempty open subset over which
$\rho$ is smooth.  Since the fibers of $\rho$ over $U$ are smooth curves of
genus zero, we have
\[
\nu^*R^1\rho_*\mathbb Q_\ell=0
\]
and hence
\[
H^1(U,\nu^*R^1\rho_*\mathbb Q_\ell)=0.
\]
Consequently, the Griffiths homomorphism
\[
\operatorname{Prim}^2(X,\mathbb Q_\ell)
\xrightarrow{\operatorname{Grif}}
H^1(U,\nu^*R^1\rho_*\mathbb Q_\ell)
\]
is zero.  But on a quadric surface,
\[
\dim_{\mathbb Q_\ell}\operatorname{Prim}^2(X,\mathbb Q_\ell)=1.
\]
Likewise, the conclusion of 5.7 fails here.
\end{filecontents*}
% Original input: components/128_expose_XVIII_section_6_opening_through_6_1_2_proof_opening.tex
\begin{filecontents*}[overwrite,noheader]{sga7ii-cumulative-input-128.tex}
\section*{6. \underline{Application of the local theory of vanishing cycles}\\
\underline{(Picard--Lefschetz)}}

\subsection*{Proposition 6.1.}
\underline{Let $Z$ be a hypersurface in projective space $\mathbb P^r_k$,
$r\geq2$, where $k$ is an algebraically closed field.  Let}
\[
i_o:D_o\lhook\joinrel\longrightarrow\mathbb P^r
\]
\underline{be the inclusion of a line which meets $Z$ only in its smooth
locus and meets it transversely.  Choose a geometric point $\bar a$ of
$D_o-D_o\cap Z$, and denote by}
\[
\pi_1^{\mathrm{mod}}(D_o-D_o\cap Z,\bar a)
\]
\underline{the quotient of $\pi_1(D_o-D_o\cap Z,\bar a)$ classifying the
étale coverings of $D_o-D_o\cap Z$ which are tamely ramified at every point
of $D_o\cap Z$.  Then:}

\subsubsection*{6.1.1.}
\underline{The canonical homomorphism}
\begin{equation}\tag{6.1.1.1}
i_{o*}:\pi_1^{\mathrm{mod}}(D_o-D_o\cap Z,\bar a)
\longrightarrow
\pi_1^{\mathrm{mod}}(\mathbb P^r-Z,\bar a)
\end{equation}
\underline{is surjective.}

\subsubsection*{6.1.2.}
\underline{Every choice of a geometric point $\bar x_o\in D_o\cap Z$ and
of a ``path'' in $D_o$ from $\bar x_o$ to $\bar a$ determines a
homomorphism}
\begin{equation}\tag{6.1.2.1}
\prod_{\ell\ne\operatorname{char}(k)}\mathbb Z_\ell(1)
\longrightarrow
\pi_1^{\mathrm{mod}}(D_o-D_o\cap Z,\bar a)
\end{equation}
\underline{whose image is ``the'' inertia group at $\bar x_o$.  Composing
this morphism with $i_{o*}$ gives a homomorphism}
\begin{equation}\tag{6.1.2.2}
\prod_{\ell\ne\operatorname{char}(k)}\mathbb Z_\ell(1)
\longrightarrow
\pi_1^{\mathrm{mod}}(\mathbb P^r-Z,\bar a).
\end{equation}
\underline{If $Z$ is irreducible, all these homomorphisms are conjugate to
one another.}

\underline{Proof.}
Let $T$ be the subscheme of the Grassmannian formed by the lines in
$\mathbb P^r$ which
\begin{enumerate}[label=(\roman*)]
\item pass through $\bar a$;
\item meet $Z$ only in its smooth locus, and meet it there transversely.
\end{enumerate}
Let $\eta$ be its generic point, and let
\[
i_\eta:D_\eta\longrightarrow\mathbb P^r\otimes_k k(\eta)
\]
be the corresponding line (thus a line ``joining $\bar a$ to the generic
point of $\mathbb P^r_k$'').  Choose also a geometric point $\bar\eta$ over
$\eta$, and write $D_{\bar\eta}$ for the corresponding ``geometric'' line.
Let $o\in T$ be the point corresponding to the given line $D_o$.
\end{filecontents*}
% Original input: components/129_expose_XVIII_proof_6_1_through_bertini.tex
\begin{filecontents*}[overwrite,noheader]{sga7ii-cumulative-input-129.tex}
By [SGA 1, XIII 2.8], there is a specialization homomorphism (depending on
a path in $T$ from $\bar\eta$ to $o$, and hence defined up to conjugacy)
\begin{equation}\tag{6.1.3}
\pi_1^{\mathrm{mod}}(D_{\bar\eta}-D_{\bar\eta}\cap Z,\bar a)
\xrightarrow{\ \bar\eta\longrightarrow o\ }
\pi_1^{\mathrm{mod}}(D_o-D_o\cap Z,\bar a)
\end{equation}
such that the following assertions hold.

\subsubsection*{6.1.4.}
The local-monodromy homomorphisms (6.1.2.1)
\[
\prod_{\ell\ne\operatorname{char}(k)}\mathbb Z_\ell(1)
\longrightarrow
\pi_1^{\mathrm{mod}}(D_o-D_o\cap Z,\bar a)
\]
are obtained by composing the specialization arrow (6.1.3) with the local
monodromy homomorphisms for $D_{\bar\eta}$:
\begin{equation}\tag{6.1.4.1}
\prod_{\ell\ne\operatorname{char}(k)}\mathbb Z_\ell(1)
\longrightarrow
\pi_1^{\mathrm{mod}}(D_{\bar\eta}-D_{\bar\eta}\cap Z,\bar a)
\xrightarrow{\ \bar\eta\longrightarrow o\ }
\pi_1^{\mathrm{mod}}(D_o-D_o\cap Z,\bar a).
\end{equation}

\subsubsection*{6.1.5.}
The following diagram is commutative:
\begin{equation}\tag{6.1.5}
\begin{tikzcd}[column sep=huge,row sep=huge]
{} & \pi_1^{\mathrm{mod}}(\mathbb P^r-Z,\bar a) & {} \\
\pi_1^{\mathrm{mod}}(D_\eta-D_\eta\cap Z,\bar a)
  \arrow[ur,"i_{\eta*}"]
& {}
& \pi_1^{\mathrm{mod}}(D_o-D_o\cap Z,\bar a)
  \arrow[ul,"i_{o*}"'] \\
{} & \pi_1^{\mathrm{mod}}(D_{\bar\eta}-D_{\bar\eta}\cap Z,\bar a)
  \arrow[uu,"i_{\bar\eta*}"']
  \arrow[ul]
  \arrow[ur,"\bar\eta\longrightarrow o"'] & {} .
\end{tikzcd}
\end{equation}

Assume this for the moment, and first prove 6.1.1.  By the commutativity of
6.1.5, it is enough to prove that
\[
i_{\bar\eta*}:
\pi_1^{\mathrm{mod}}(D_{\bar\eta}-D_{\bar\eta}\cap Z,\bar a)
\longrightarrow
\pi_1^{\mathrm{mod}}(\mathbb P^r-Z,\bar a)
\]
is surjective.  For this, it is enough to prove that the analogous
homomorphism of the \underline{ordinary} fundamental groups
\begin{equation}\tag{6.1.6}
i_{\bar\eta*}:
\pi_1(D_{\bar\eta}-D_{\bar\eta}\cap Z,\bar a)
\longrightarrow
\pi_1(\mathbb P^r-Z,\bar a)
\end{equation}
is surjective.  In other words, for every étale covering
$E\to\mathbb P^r-Z$ with $E$ connected (and hence irreducible, since
$\mathbb P^r-Z$ is smooth and therefore unibranch [EGA IV, 18.10.3]), its
inverse image
\[
E_{\bar\eta}
=E\times_{\mathbb P^r-Z}(D_{\bar\eta}-D_{\bar\eta}\cap Z)
\]
over $D_{\bar\eta}-D_{\bar\eta}\cap Z$ must remain geometrically
irreducible.  It is enough to use the following variant of Bertini's theorem
([EGA V, citation left blank in the source]) for the composite
$E\to\mathbb P^r-Z\lhook\joinrel\longrightarrow\mathbb P^r$ and $s=1$.

\subsubsection*{6.1.6.1. Bertini's theorem.}
Let $k$ be a field, let $\mathbb P^r=\mathbb P^r_k$, and let
$f:X\to\mathbb P^r$ be a nonconstant morphism, where $X$ is a geometrically
irreducible algebraic scheme.  For every integer
$s<\dim\overline{f(X)}$, let $L_t$ denote the \underline{generic} linear
subspace of codimension $s$ in $\mathbb P^r$.  Then
$X_t=f^{-1}(L_t)$ is geometrically irreducible.
\end{filecontents*}
% Original input: components/130_expose_XVIII_completion_of_6_1.tex
\begin{filecontents*}[overwrite,noheader]{sga7ii-cumulative-input-130.tex}
We now prove 6.1.2.  By 6.1.4 and the commutativity of (6.1.5), it is
enough to prove that all the local-monodromy homomorphisms (6.1.3)
\underline{for} $D_{\bar\eta}$,
\begin{equation}\tag{6.1.7.1}
\prod_{\ell\ne\operatorname{char}(k)}\mathbb Z_\ell(1)
\longrightarrow
\pi_1^{\mathrm{mod}}(D_{\bar\eta}-D_{\bar\eta}\cap Z,\bar a),
\end{equation}
become conjugate after composition with the homomorphism
\begin{equation}\tag{6.1.7.2}
\pi_1^{\mathrm{mod}}(D_{\bar\eta}-D_{\bar\eta}\cap Z,\bar a)
\longrightarrow
\pi_1^{\mathrm{mod}}(D_\eta-D_\eta\cap Z,\bar a).
\end{equation}
To prove this, observe that for every finite Galois extension $K$ of
$k(\eta)$, with
\[
k(\eta)\subset K\subset k(\bar\eta),
\]
the scheme
\[
D_K\stackrel{\mathrm{def}}{=}D_\eta\times_{k(\eta)}K
\]
is an étale covering of $D_\eta$.  Thus
$D_K-D_K\cap Z$ is a tame covering of $D_\eta-D_\eta\cap Z$, and hence
$\pi_1^{\mathrm{mod}}(D_\eta-D_\eta\cap Z,\bar a)$ acts on the scheme
$D_K\cap Z$.  We take $K$ sufficiently large that every point of
$D_K\cap Z$ is $K$-rational; this is possible because, by hypothesis,
$D_\eta\cap Z$ is étale over $k(\eta)$.

It remains to prove that
$\pi_1^{\mathrm{mod}}(D_\eta-D_\eta\cap Z,\bar a)$ acts
\underline{transitively} on the points of $D_K\cap Z$.  This will prove
that the composite local-inertia homomorphisms
\[
\prod_{\ell\ne\operatorname{char}(k)}\mathbb Z_\ell(1)
\longrightarrow
\pi_1^{\mathrm{mod}}(D_\eta-D_\eta\cap Z,\bar a)
\]
attached to the various points of $D_K\cap Z$ are conjugate.  Now the
action on $D_K\cap Z$ factors through the canonical quotient
\footnote{The numerator of the printed Galois group is blank.  The field
tower immediately preceding the display determines it as $k(\bar\eta)$.}
\[
\pi_1^{\mathrm{mod}}(D_\eta-D_\eta\cap Z,\bar a)
\longrightarrow
\pi_1(\eta,\bar\eta)
=\operatorname{Gal}\bigl(k(\bar\eta)/k(\eta)\bigr).
\]
The action of $\operatorname{Gal}(k(\bar\eta)/k(\eta))$ on
$D_K\cap Z$ is the one compatible with the inclusion
\[
D_K\cap Z\lhook\joinrel\longrightarrow\mathbb P^r_k\times_k K,
\]
where the Galois group acts on $\mathbb P^r_k\times_k K$ only through
the second factor.  It follows that the desired assertion is equivalent
to saying that $D_\eta\cap Z$ has only one point, that is, that it is the
spectrum of a \underline{field}.  In elementary terms, this is precisely
the following Bertini-type fact (take an affine equation $F=0$ for $Z$
and suppose that $\bar a$ is the origin).

\subsubsection*{Fact 6.1.8.}
Let $k$ be an algebraically closed field, let $r\ge2$, and let
$F(X_1,\ldots,X_r)\in k[X_1,\ldots,X_r]$ be irreducible.  Then
\[
G(T)=F(X_1T,\ldots,X_rT)\in k(X_1,\ldots,X_r)[T]
\]
is irreducible as a polynomial in $T$ with coefficients in
$k(X_1,\ldots,X_r)$.

This completes the proof of 6.1.2.

\subsubsection*{Remark 6.1.9.}
Another method of proving 6.1.2, valid under considerably more general
conditions, would be simply to establish a general compatibility for
``monodromy homomorphisms'': the composite homomorphisms in 6.1.2 are
conjugate to the monodromy homomorphism at the generic point of $Z$, and
therefore are conjugate to one another.

\subsubsection*{Corollary 6.1.10.}
Let $X$ be a smooth connected projective variety of dimension $n$ over an
algebraically closed field $k$.  Suppose that $n=\dim X$ is even, or that
$\operatorname{char}(k)\ne2$.\footnote{The print has
$\operatorname{card}(k)\ne2$; the parallel hypothesis in (6.3.2) and the
tame-monodromy argument require the characteristic.}
Let $D$ be a Lefschetz pencil for $X$, let $\eta$ be its generic point,
and let $\bar\eta$ be a geometric point over $\eta$.  Then the natural
representation of
$\pi_1(D-D\cap\check X,\bar\eta)$ on
$H^{n-1}(X_{\bar\eta},\mathbb Q_\ell)$ factors through
$\pi_1^{\mathrm{mod}}(D-D\cap\check X,\bar\eta)$.  If this
representation is nontrivial, then $D\cap\check X\ne\varnothing$, and
the inertia group at every point of $D\cap\check X$ acts nontrivially.

\subsubsection*{Proof.}
If $\dim(\check X)\le r-2$, then $D\cap\check X$ is empty and there is
nothing to prove, since $\mathbb P^1$ is simply connected.  Otherwise
$\check X$ is a hypersurface in $\mathbb P^r$ [XVII, 3.1.4], and 6.1.1
applies by 3.2.10 and XVII 3.5, with $Z=\check X$ and $D_o=D$.  The fact
that the action of $\pi_1$ on
$H^{n-1}(X_{\bar\eta},\mathbb Q_\ell)$ is tame was established in the
local theory by the Picard--Lefschetz formula [XV, 3.4], recalled below.
Finally, the nontrivial fact that
$\pi_1^{\mathrm{mod}}(D-D\cap\check X,\bar\eta)$ is generated by the
inertia groups at the various points of $D\cap\check X$ [SGA 1, XIII
2.10] gives the required conclusion.

We shall now use the preceding result to study condition (A) of (5.3.5).
First recall the following consequence of the Picard--Lefschetz formula
[XV, 3.4].
\end{filecontents*}
% Original input: components/131_expose_XVIII_proposition_6_2_through_corollary_6_4.tex
\begin{filecontents*}[overwrite,noheader]{sga7ii-cumulative-input-131.tex}
\subsubsection*{Proposition 6.2.}
Let $X$ be a smooth connected projective variety of dimension $n$ over
the algebraically closed field $k$.  Let $D$ be a Lefschetz pencil, let
$\eta$ be its generic point, and let $\bar\eta$ be a geometric point
over $\eta$.  Then the pencil $D$ satisfies condition (A) if and only if
the inertia group at every point of $D\cap\check X$ acts nontrivially on
$H^{n-1}(X_{\bar\eta},\mathbb Q_\ell)$.

\subsubsection*{Theorem 6.3.}
Under the hypotheses of 6.2, the following assertions hold.

\medskip
\noindent
\textup{(6.3.1)} If $n=\dim X$ is odd, condition (A) is satisfied by
every Lefschetz pencil for $X$.

\medskip
\noindent
\textup{(6.3.2)} Suppose that $n=\dim X$ is even, or that
$\operatorname{char}(k)\ne2$.  Then the following conditions are
equivalent:
\begin{enumerate}
\item[(i)] $D$ does not satisfy condition (A) of (5.3.5);
\item[(ii)] $\pi_1(D-D\cap\check X)$ acts trivially on
  $H^{n-1}(X_{\bar\eta},\mathbb Q_\ell)$, and
  $D\cap\check X\ne\varnothing$;
\item[(iii)] there is a point of $D\cap\check X$ whose inertia group
  acts trivially on $H^{n-1}(X_{\bar\eta},\mathbb Q_\ell)$;
\item[(iv)] the sheaf $R^{n-1}\rho_*\mathbb Q_\ell$ is constant on
  $D$, and $D\cap\check X$ is nonempty;
\item[(v)] the restriction
  $R^{n-1}\rho_*\mathbb Q_\ell\vert_U$ is constant on $U$, where $U$
  is a nonempty open subset of $D$ over which $\rho$ is smooth, and
  $\rho$ is not smooth.
\end{enumerate}

\noindent
\textup{(6.3.3)} If, in addition, $X$ satisfies condition (LV) of
(5.2.2), conditions (i)--(v) are also equivalent to:
\begin{enumerate}
\item[(vi)] the restriction map
\[
H^{n-1}(X,\mathbb Q_\ell)
\longrightarrow H^{n-1}(X_{\bar\eta},\mathbb Q_\ell)
\]
is surjective (and hence is an isomorphism by the weak Lefschetz
theorem).
\end{enumerate}

\subsubsection*{Proof.}
Assertion 6.3.1 follows from XV 3.4.  For the equivalences in 6.3.2,
(i) is equivalent to (iii) by 6.2, (ii) to (iii) by 6.1.6, (ii) to (v)
tautologically, and (iv) to (v) by 5.3.4.  Finally, if $X$ satisfies
condition (LV), 5.6.10 gives the equivalence of (ii) and (vi).

\subsubsection*{Remark 6.3.4.}
Condition (A) does not depend on the choice of the Lefschetz pencil
$D$, but only on the chosen (Lefschetz) immersion
$i:X\longrightarrow\mathbb P^r$.  This is clear from (6.3.1) when
$\dim X$ is odd.  When $\dim X$ is even, condition (A) also says that
the local system on $\mathbb P^r-\check X$ formed by the groups
$H^{n-1}(X\cap H_s,\mathbb Q_\ell)$ is constant.  This follows from
criterion (ii), from the fact that the local system in question is
tamely ramified at the generic point of $\check X$, and from the
surjectivity in (6.1.2).

\subsubsection*{Corollary 6.4.}
Under the hypotheses of 6.2, suppose that
$n=\dim X$ is even or $\operatorname{char}(k)\ne2$, and that $X$
satisfies condition (LV).\footnote{In the print, the placement of the
conjunction attaches condition (LV) only to the second alternative.
The proof invokes the equivalence (i)$\Longleftrightarrow$(vi) of
(6.3.2), so condition (LV) is required in either case.}
Then there is an integer $d_o$ such that, for $d\ge d_o$, every
Lefschetz pencil of sections of $X$ by hypersurfaces of degree $d$
satisfies condition (A).

\subsubsection*{Proof.}
By the equivalence (i)$\Longleftrightarrow$(vi) of (6.3.2), it is
enough to prove that, if $Y(d)$ denotes a smooth hypersurface section of
$X$ of degree $d$, then
\begin{equation}\tag{6.4.1}
\dim H^{n-1}(Y(d),\mathbb Q_\ell)\longrightarrow\infty
\qquad\text{as }d\longrightarrow\infty.
\end{equation}
Indeed, we have the following precise result.

\subsubsection*{Lemma 6.4.2.}
Let $X$ be a smooth connected projective variety of dimension $n\ge1$
over an algebraically closed field $k$, and let $Y(d)$ denote a
hypersurface section of degree $d$.  Then there is a polynomial
$f(T)\in\mathbb Z[T]$ such that
\begin{align}
f(T)&=\deg(X)T^n+\text{a polynomial of degree at most }n-1,
  \tag{6.4.2.1}\\
f(d)&=\dim H^{n-1}(Y(d),\mathbb Q_\ell)\qquad(d\ge1).
  \tag{6.4.2.2}
\end{align}

\subsubsection*{Proof.}
By XVII 5.7.5, if $C(X)$ denotes the total Chern class of $X$ and $H$
the first Chern class of a hyperplane, the Euler--Poincaré
characteristic of $Y(d)$ is
\begin{equation}\tag{6.4.2.3}
\chi(Y(d))
=\int_X\frac{dH\,C(X)}{1+dH}
=\sum_{i=1}^{n}(-1)^{i-1}d^i\int_X H^iC_{n-i}.
\end{equation}
On the other hand, Poincaré duality and weak Lefschetz give
\begin{equation}\tag{6.4.2.4}
\chi(Y(d))
=(-1)^{n-1}\dim H^{n-1}(Y(d),\mathbb Q_\ell)
+2\sum_{i\ge1}(-1)^{n-1-i}
  \dim H^{n-1-i}(X,\mathbb Q_\ell).
\end{equation}
Combining (6.4.2.3) and (6.4.2.4) gives the result.

\end{filecontents*}
% Original input: components/132_expose_XVIII_sections_6_5_and_6_6.tex
\begin{filecontents*}[overwrite,noheader]{sga7ii-cumulative-input-132.tex}
\subsection*{6.5. A reminder on the Picard--Lefschetz formula (XV 3.4)}

Let $S$ be the spectrum of a strictly henselian valuation ring, with
generic point $\eta$, geometric generic point $\bar\eta$, and closed
point $s$.  Let
\[
\rho:X\longrightarrow S
\]
be a proper morphism, where $X$ is a connected regular scheme of
dimension $n$.  Suppose that $\rho$ is smooth away from a single point
$x$ in the special fiber and that $x$ is an ordinary quadratic
singularity of that fiber.

Let $I$ denote the inertia group of $S$ at $s$.  By XV 3.4, $I$ acts
trivially on $H^i(X_{\bar\eta},\mathbb Q_\ell)$ for $i\ne n-1$, and its
action on $H^{n-1}(X_{\bar\eta},\mathbb Q_\ell)$ is described as follows.

\subsubsection*{6.5.1.}
If $n=2m+1$, then $I$ acts through a quotient isomorphic to
$\mathbb Z/2\mathbb Z$ (if $\operatorname{char}(k)\ne2$, there is only
one such quotient).  There is an element
\[
\delta\in H^{n-1}(X_{\bar\eta},\mathbb Q_\ell(m)),
\]
determined uniquely up to sign and called the \underline{vanishing
cycle at $s$}, such that
\begin{equation}\tag{6.5.1.1}
\delta\cdot\delta=(-1)^m2
\quad\text{in}\quad
H^{2n-2}(X_{\bar\eta},\mathbb Q_\ell(n-1))\simeq\mathbb Q_\ell.
\end{equation}
The nontrivial element $\sigma$ of $\mathbb Z/2\mathbb Z$ acts on
$H^{n-1}(X_{\bar\eta},\mathbb Q_\ell)$ by the reflection
\begin{equation}\tag{6.5.1.2}
\sigma x=x-(-1)^m(x\cdot\delta)\delta.
\end{equation}
\footnote{The print numbers the reflection formula (6.5.2.2), although
it is the second formula of 6.5.1 and is followed by (6.5.2.1).}
Here $x\cdot\delta$ belongs to
$H^{2n-2}(X_{\bar\eta},\mathbb Q_\ell(m))\simeq\mathbb Q_\ell(-m)$,
so $(x\cdot\delta)\delta$ defines an element of
$H^{n-1}(X_{\bar\eta},\mathbb Q_\ell)$.

\subsubsection*{6.5.2.}
If $n=2m+2$, then $I$ acts through its canonical quotient
$\mathbb Z_\ell(1)$.  There is an element
\[
\delta\in H^{n-1}(X_{\bar\eta},\mathbb Q_\ell(m)),
\]
canonically determined up to sign and called the \underline{vanishing
cycle at $s$}, such that $\sigma\in\mathbb Z_\ell(1)$ acts on
$x\in H^{n-1}(X_{\bar\eta},\mathbb Q_\ell)$ by the transvection
\begin{equation}\tag{6.5.2.1}
\sigma(x)=x-(-1)^m\sigma\,(x\cdot\delta)\delta.
\end{equation}
Indeed, $x\cdot\delta$ belongs to
$H^{2n-2}(X_{\bar\eta},\mathbb Q_\ell(m))\simeq\mathbb Q_\ell(-m-1)$.
Multiplication by $\sigma\in\mathbb Z_\ell(1)$ therefore gives an
element $\sigma(x\cdot\delta)\in\mathbb Q_\ell(-m)$, whose product with
$\delta\in H^{n-1}(X_{\bar\eta},\mathbb Q_\ell(m))$ lies in
$H^{n-1}(X_{\bar\eta},\mathbb Q_\ell)$.

We return to the global situation.

\subsubsection*{Theorem 6.6.}
Let $X$ be a smooth connected projective variety of dimension $n$ over
an algebraically closed field $k$.  Suppose that $n$ is even or that
$\operatorname{char}(k)\ne2$.  Let $D$ be a Lefschetz pencil for $X$,
let $\eta$ be its generic point, let $\bar\eta$ be a geometric point
over $\eta$, and let $s_1,\ldots,s_N$ be the points of
$D\cap\check X$.

For each $i$, choose an isomorphism $\ell_i$ between $\bar\eta$ and the
geometric generic point of $S_i$, the strict henselization of $D$ at
$s_i$.  Via $\ell_i$, let $I_i$ be the inertia group at $s_i$ acting on
$H^{n-1}(X_{\bar\eta},\mathbb Q_\ell(j))$, and let $\delta_i$ be the
vanishing cycle at $s_i$ defined in 6.5.\footnote{The print calls
$I_i$ the inertia group ``of $\delta_i$'' before $\delta_i$ is defined.
The standard construction is the inertia group at $s_i$, embedded by
the chosen path $\ell_i$.}
Thus
\[
\delta_i\in H^{n-1}(X_{\bar\eta},\mathbb Q_\ell(m)),
\qquad
m=\begin{cases}
(n-1)/2,& n\text{ odd},\\
n/2-1,& n\text{ even}.
\end{cases}
\]
\footnote{The print omits the generic-fiber subscript in this display;
the preceding action and the definition in 6.5 place every
$\delta_i$ in the cohomology of $X_{\bar\eta}$.}
Choose the $\ell_i$ so that the subgroups
\[
I_i\subset
\pi=\pi_1^{\mathrm{mod}}(D-D\cap\check X,\bar\eta)
\]
topologically generate $\pi$, as is possible by SGA 1, XIII 2.10.
Then:

\medskip
\noindent
\textup{6.6.1.} If $X$ satisfies condition (LV), the vanishing
cohomology space
\[
E^{n-1}(X_{\bar\eta},\mathbb Q_\ell(m))
\stackrel{\mathrm{def}}{=}
H^{n-1}(X,\mathbb Q_\ell(m))^\perp
\subset H^{n-1}(X_{\bar\eta},\mathbb Q_\ell(m))
\]
(cf. 5.4.4) is generated by the vanishing cycles
$\delta_i$, $i=1,\ldots,N$.

\medskip
\noindent
\textup{6.6.2.} The vanishing cycles $\delta_i$ are conjugate, up to
sign, under the global monodromy group
$\pi_1(D-D\cap\check X)$.

\subsubsection*{Proof.}
For 6.6.1, the Picard--Lefschetz formula 6.5 and the fact that the
$I_i$ generate $\pi=\pi_1(D-D\cap\check X)$ give
\begin{equation}\tag{6.6.3}
H^{n-1}(X_{\bar\eta},\mathbb Q_\ell(m))^\pi
=\bigl\langle\delta_1,\ldots,\delta_N\bigr\rangle^\perp
\quad\text{in }H^{n-1}(X_{\bar\eta},\mathbb Q_\ell(m)).
\end{equation}
By 5.6.10,
\begin{equation}\tag{6.6.4}
H^{n-1}(X_{\bar\eta},\mathbb Q_\ell(m))^\pi
=\operatorname{Im}\bigl(
H^{n-1}(X,\mathbb Q_\ell(m))
\longrightarrow H^{n-1}(X_{\bar\eta},\mathbb Q_\ell(m))
\bigr).
\end{equation}
Taking orthogonals in (6.6.3) proves 6.6.1.

For 6.6.2, note simply that the vanishing cycle $\delta_i$ is
determined up to sign by the representation of the inertia group at
$s_i$.  Apply 6.1.2, which gives conjugacy in
$\pi_1^{\mathrm{mod}}(\check{\mathbb P}^{r}-\check X,\bar\eta)$ of the
composite inertia homomorphisms
\begin{center}
\begin{tikzcd}[column sep=huge,row sep=huge]
\mathbb Z_\ell(1) \arrow[r]
  & \pi_1^{\mathrm{mod}}(D-D\cap\check X,\bar\eta) \arrow[d] & {}\\
{}
  & \pi_1^{\mathrm{mod}}(\check{\mathbb P}^{r}-\check X,\bar\eta) & .
\end{tikzcd}
\end{center}
Consequently, the vanishing cycles determined by these homomorphisms
are conjugate up to sign by elements of
$\pi_1^{\mathrm{mod}}(\check{\mathbb P}^{r}-\check X,\bar\eta)$.
The proof is completed by 6.1.1, which gives the surjection
\[
\pi_1^{\mathrm{mod}}(D-D\cap\check X)
\longrightarrow
\pi_1^{\mathrm{mod}}(\check{\mathbb P}^{r}-\check X).
\]

\end{filecontents*}
% Original input: components/133_expose_XVIII_corollary_6_7_and_bibliography.tex
\begin{filecontents*}[overwrite,noheader]{sga7ii-cumulative-input-133.tex}
\subsubsection*{Corollary 6.7.}
Under the hypotheses of 6.6, suppose that $X$ satisfies condition
(LV).  Then:

\medskip
\noindent
\textup{6.7.1.} The representation of
$\pi_1(D-D\cap\check X)$ on
$E^{n-1}(X_{\bar\eta},\mathbb Q_\ell)$ is absolutely irreducible.

\medskip
\noindent
\textup{6.7.2.} If $n$ is even, every finite-index subgroup of
$\pi_1(D-D\cap\check X)$ acts absolutely irreducibly on
$E^{n-1}(X_{\bar\eta},\mathbb Q_\ell)$.

\subsubsection*{Proof.}
We may replace the module in question by its twist
\[
E=E^{n-1}(X_{\bar\eta},\mathbb Q_\ell(m)).
\]
Let
$\bar E=E\otimes_{\mathbb Q_\ell}\bar{\mathbb Q}_\ell$, and choose
$0\ne\chi\in\bar E$.  To prove 6.7.1, we must show that the
$\bar{\mathbb Q}_\ell$-linear subspace of $\bar E$ generated by $\chi$
and its conjugates under $\pi$ is all of $\bar E$.  By 6.6.1 and 6.6.2,
it is enough to prove that this subspace contains one of the vanishing
cycles $\delta_i$.

Since $X$ satisfies (LV), the restriction of the cup-product form to
\[
H^{n-1}(X,\mathbb Q_\ell(m))
\lhook\joinrel\longrightarrow
H^{n-1}(X_{\bar\eta},\mathbb Q_\ell(m))
\]
is nondegenerate; hence so is its restriction to the orthogonal
$E^{n-1}(X_{\bar\eta},\mathbb Q_\ell(m))$.  By 6.6.1, there is therefore
at least one $i$ such that
\[
\chi\cdot\delta_i\ne0.
\]
For this $i$, 6.5 supplies an element $\sigma_i$ of the inertia group at
$s_i$ such that
\begin{equation}\tag{6.7.3}
\sigma_i(\chi)-\chi
=\text{a nonzero multiple of }\delta_i.
\end{equation}
This proves 6.7.1.

For 6.7.2, $n-1$ is odd, so $\delta_i\cdot\delta_i=0$.  Iterating
6.5.2 gives, for every $k$,
\begin{equation}\tag{6.7.4}
\sigma_i^k(\chi)
=\chi+k\,\sigma_i(\chi\cdot\delta_i)\delta_i.
\end{equation}
Thus a $\mathbb Q_\ell$-subspace $M$ of $E$ is stable under the action
of $\pi$ if and only if it is stable under a finite-index subgroup:
both conditions are equivalent to
\[
\chi\in M
\quad\Longrightarrow\quad
\delta_i\in M\quad\text{for every $i$ with }
\chi\cdot\delta_i\ne0.
\]
The assertion therefore reduces to 6.7.1.

\subsection*{Bibliography}

\begin{description}
\item[{[1]}] P. Deligne, \emph{Théorème de Lefschetz et critères de
dégénérescence de suites spectrales}, Publ. Math. IHÉS \textbf{35}
(1968), 259--278.

\item[{[2]}] S. Kleiman, \emph{Algebraic cycles and the Weil
conjectures}, in \emph{Dix exposés sur la cohomologie des schémas},
North-Holland, 1968.

\item[{[3]}] S. Lefschetz, \emph{L'analyse situs et la géométrie
algébrique}, Gauthier-Villars, Paris, 1924.

\item[{[4]}] D. Lieberman, \emph{Higher Picard Varieties}, Amer. J.
Math. \textbf{90} (4) (1968), 1165--1199.\footnote{The print gives the
journal volume as 15; the article appeared in volume 90.}

\item[{[5]}] D. Mumford, unpublished.

\item[{[6]}] A. Weil, \emph{Variétés kählériennes}, Hermann, Paris,
1958.
\end{description}

\end{filecontents*}
% Original input: components/134_expose_XIX_opening_through_corollary_1_5.tex
\begin{filecontents*}[overwrite,noheader]{sga7ii-cumulative-input-134.tex}
\clearpage
\begin{center}
{\large SGA 7 --- EXPOS\'E XIX}\par
\vspace{2em}
{\Large\underline{The Noether Theorem}}\par
\vspace{1.5em}
{\large by P. Deligne}
\end{center}

\section*{Contents}

\noindent 1. Statement of the theorems. \hfill 1

\noindent 2. A complement to XVIII. \hfill 4

\noindent 3. Proof of 1.3. \hfill 5

\noindent 4. A second complement to XVIII. \hfill 8

\noindent 5. The exceptional cases of the Noether theorem. \hfill 11

\section*{1. Statement of the theorems}

The present expos\'e updates Lefschetz's proof [3, p.~108] of the Noether
theorem (1.2) stated below.

\subsubsection*{1.1.}
Let $k$ be an algebraically closed field of characteristic $p$.  Let
$d\geq1$ and let
\[
\underline a=(a_i)_{1\leq i\leq d}
\]
be a sequence of integers at least $1$.  A complete intersection $V$ of
dimension $n$ and multidegree $\underline a$ in the projective space
$\mathbb P^{n+d}$ over $k$ will be called \underline{generic} if, after an
extension of $k$, it is projectively isomorphic to a complete intersection
defined by equations whose coefficients are algebraically independent over
the prime field.  In what follows, unless expressly stated otherwise, we
suppose that $d\geq1$ and $a_i\geq2$.

\subsubsection*{Theorem 1.2.}
Let $V\subset\mathbb P^{2+d}$ be a generic complete-intersection surface of
multidegree $\underline a$.  Suppose that
\[
\underline a\ne(2),(3),(2,2).
\]
Then every (reduced) curve on $V$ is the scheme-theoretic intersection of
$V$ with a hypersurface of $\mathbb P^{2+d}$.

Let us show that the theorem follows from its corollary
\begin{equation}\tag{1.2.1}
\operatorname{Pic}(\mathbb P^{2+d})=\mathbb Z
\xrightarrow{\ \sim\ }\operatorname{Pic}(V).
\end{equation}
If $C$ is a curve on $V$, (1.2.1) implies that
$\mathcal O(C)\simeq\mathcal O(n)$ for a suitable $n$.  If $H$ is a
hyperplane section, then
\[
n=\bigl(\textstyle\prod_i a_i\bigr)^{-1}(C,H)>0.
\]
It remains to prove that the linear system induced on $V$ by $|nH|$ is
complete.  Write $V$ as the intersection of hypersurfaces $H_i$ of degree
$a_i$, and put $V_i=H_1\cap\cdots\cap H_i$.  Induction on $i$ shows that the
linear system induced by $|nH|$ on $V_i$ is complete: in the exact sequence
\[
H^0(V_i,\mathcal O(n))\longrightarrow
H^0(V_{i+1},\mathcal O(n))\longrightarrow
H^1(V_i,\mathcal O(n-a_{i+1})),
\]
the last group vanishes (FAC).

\subsubsection*{Theorem 1.3.}
Let $V\subset\mathbb P^{2n+d}$ be a generic complete intersection of
nonzero even dimension $2n$ and multidegree $\underline a$.  Suppose that
one of the following conditions holds:
\begin{enumerate}[label=(\alph*)]
\item $k$ is not of characteristic $2$, and $(2n;\underline a)$ is distinct
  from $(2n;2)$, $(2n;2,2)$, and $(2;3)$;
\item the Hodge number $h^{2n,0}$ is nonzero, i.e.
  $\sum_i(a_i-1)>2n$.
\end{enumerate}
Then every algebraic cohomology class in
$H^{2n}(V,\mathbb Q_\ell(n))$ is a rational multiple of the class
$\eta^n$ of a linear section.

Let us deduce (1.2.1) from 1.3.  The exceptional cases in 1.2 are exactly
those for which $\sum_i(a_i-1)\leq2$; hence 1.3(b) applies to the surface
$V$ of 1.2.\footnote{The print writes $d_i$ in this sum, although the
multidegree has been denoted throughout by $(a_i)$.}  We know that
$\operatorname{Pic}^{\tau}(V)=0$ (XI 1.8), so
$\operatorname{Pic}(V)$ injects into $H^2(V,\mathbb Q_\ell(1))$.
By 1.3 every element of $\operatorname{Pic}(V)$ is a rational multiple of
the class $\eta$ of $\mathcal O(1)$, and (XI 1.8) gives the conclusion.

It is clear that the validity of 1.2 and 1.3 does not depend on the chosen
generic complete intersection, but only on $p$, $n$, and $\underline a$.

\subsubsection*{Theorem 1.4.}
Let $X\subset\mathbb P$ be a smooth projective variety, pure of dimension
$2n+1$ over $k$.  Let $\ell\subset\check{\mathbb P}$ be a Lefschetz pencil
of hyperplane sections of $X$.  If $k$ has characteristic $2$, suppose in
addition that $\ell$ is a generic Lefschetz pencil.

Suppose that $X$ satisfies condition (LV) (XVIII 5.2.2), and let $V$ be the
generic hyperplane section in the pencil $\ell$.  One of the following two
conditions holds:
\begin{enumerate}[label=(\alph*)]
\item no nonzero cohomology class in the vanishing part
  $\operatorname{Ev}(V,\mathbb Q_\ell(n))$ of
  $H^{2n}(V,\mathbb Q_\ell(n))$ is algebraic;
\item the algebraic cohomology classes generate
  $\operatorname{Ev}(V,\mathbb Q_\ell(n))$.
\end{enumerate}

Let us make precise the meaning of ``generic'' in 1.4.

\begin{enumerate}[label=(\alph*)]
\item Let $G$ be the Grassmannian of lines in the projective space
  $\check{\mathbb P}$ dual to $\mathbb P$.  By a ``geometric point'' of
  $G$ we mean a morphism from the spectrum of an algebraically closed field
  to $G$.  Every geometric point $\bar s\to G$ defines a pencil $\ell$ of
  hyperplane sections of the variety obtained from $X$ by extending scalars
  from $k$ to $k(\bar s)$.  We call $\ell$ generic if $\bar s$ lies over
  the generic point of $G$.
\item Let $\ell$ be a pencil of hyperplane sections (a line in
  $\check{\mathbb P}$).  Every geometric point $\bar s\to\ell$ defines a
  hyperplane section of the variety obtained from $X$ by extending scalars
  from $k$ to $k(\bar s)$.  We call this section generic if $\bar s$ lies
  over the generic point of $\ell$.
\end{enumerate}

Let $\bar s$ be as above, with image the generic point $s$ of $\ell$, and
let $S\subset\ell$ be the finite set of points giving singular hyperplane
sections.  The monodromy group $\pi_1(\ell-S,\bar s)$, a quotient of
$\operatorname{Gal}(k(\bar s)/k(s))$, acts on
$\operatorname{Ev}(V,\mathbb Q_\ell(n))$.  If $k$ is not of characteristic
$2$, this action is irreducible (XVIII 6.7); in Section 2 below we shall show
that the assertion remains true without a characteristic hypothesis for a
generic pencil.  The subspace of $\operatorname{Ev}(V,\mathbb Q_\ell(n))$
generated by algebraic cohomology classes is stable under monodromy (that is,
under Galois).  It is therefore either zero or the whole space, which proves
1.4.

\subsubsection*{Corollary 1.5.}
Let $V\subset\mathbb P^{2n+d}$ be a generic complete intersection of even
dimension $2n$ and multidegree $\underline a$.  Then either
\begin{enumerate}[label=(\alph*)]
\item every algebraic cohomology class in
  $H^{2n}(V,\mathbb Q_\ell(n))$ is a rational multiple of $\eta^n$; or
\item $H^{2n}(V,\mathbb Q_\ell(n))$ is generated by algebraic cohomology
  classes.
\end{enumerate}
If (b) holds, then for every complete intersection of dimension $2n$ and
multidegree $\underline a$ over an algebraically closed field of
characteristic $p$, the group $H^{2n}(V,\mathbb Q_\ell(n))$ is generated by
algebraic classes.

Apply 1.4 with $X$ a generic complete intersection of dimension $2n+1$ and
multidegree $(a_1,\ldots,a_{d-1})$, and with $\ell$ a generic pencil of
sections of $X$ by hypersurfaces of degree $a_d$ (cf. XVII 2.4).  One could
equally take $X$ of multidegree $(a_1,\ldots,a_d)$ and $\ell$ a generic
pencil of hyperplane sections.  The final assertion follows because an
algebraic cohomology class remains algebraic under specialization.

\end{filecontents*}
% Original input: components/135_expose_XIX_section_2.tex
\begin{filecontents*}[overwrite,noheader]{sga7ii-cumulative-input-135.tex}
\section*{2. A complement to XVIII}

\subsubsection*{Proposition 2.1.}
Let $X\subset\mathbb P$ be a smooth projective variety, pure of odd
dimension $2n+1$ over an algebraically closed field $k$, and let
$\ell\subset\check{\mathbb P}$ be a generic Lefschetz pencil of hyperplane
sections of $X$.  Let $s$ be a geometric point of $\ell$, let $H_s$ be the
corresponding hyperplane section, let $S\subset\ell$ be the set of exceptional
values of the pencil, and let $\pi_1(\ell-S,s)$ be the monodromy group.
Suppose that $X$ satisfies condition (LV).  Then the representation of
$\pi_1(\ell-S,s)$ on the vanishing cohomology
\[
\operatorname{Ev}(H_s,\mathbb Q_\ell)
\subset H^{2n}(H_s,\mathbb Q_\ell)
\]
is absolutely irreducible.\footnote{The print says that $X$ has even
dimension $2n$.  The displayed degree $H^{2n}(H_s)$, the reflection formula,
Remark 2.2, and the application to Theorem 1.4 all force odd dimension
$2n+1$.}

Recall that $\operatorname{Ev}(H_s,\mathbb Q_\ell)$ is the orthogonal
complement of the image of $H^{2n}(X,\mathbb Q_\ell)$ in
$H^{2n}(H_s,\mathbb Q_\ell)$, and that this image equals the subgroup of
monodromy invariants (XVIII 5.6.10).

\subsubsection*{Proof.}
The projective line $\mathbb P^1$ is simply connected.  Hence
$\pi_1(\ell-S,s)$ is topologically generated by the inertia subgroups at the
points of $S$ and their conjugates.  Its image in
$\operatorname{GL}(H^{2n}(H_s,\mathbb Q_\ell))$ is therefore topologically
generated by the transformations
\[
x\longmapsto x-(-1)^n(x\cdot\delta)\delta,
\]
where $\delta$ is conjugate to a vanishing cycle.\footnote{The print gives
$H^n(H_s)$ at this point, although every surrounding occurrence and the
statement of the proposition use the middle degree $H^{2n}(H_s)$.}
Consequently, the image of $H^{2n}(X,\mathbb Q_\ell)$ in
$H^{2n}(H_s,\mathbb Q_\ell)$ is the orthogonal complement of the subspace
generated by the vanishing cycles and their conjugates; thus
$\operatorname{Ev}(H_s,\mathbb Q_\ell)$ is generated by those cycles and
their conjugates.

For a \underline{generic} Lefschetz pencil, the arguments of XVIII §6 show
that the vanishing cycles attached to the different points of $S$ are
conjugate to one another.

Let $W\subset\operatorname{Ev}(H_s,\mathbb Q_\ell)$ be a nonzero invariant
subspace.  There exist a conjugate $\delta$ of a vanishing cycle and an
element $v\in W$ such that $(v,\delta)\ne0$; this is where condition (LV) is
used.  Since
\[
v-(-1)^n(v\cdot\delta)\delta\in W,
\]
we have $\delta\in W$.\footnote{The print uses the undefined exponent $m$ in
this second occurrence; the immediately preceding reflection and the
Picard--Lefschetz formula give $n$.}  The subspace $W$, containing one
vanishing cycle and all its conjugates, is therefore
$\operatorname{Ev}(H_s,\mathbb Q_\ell)$.  The argument remains valid after
extending the coefficient field $\mathbb Q_\ell$, proving absolute
irreducibility.

\subsubsection*{Remark 2.2.}
If condition (LV) is omitted, the preceding arguments still show that every
invariant subspace of $H^{2n}(H_s,\mathbb Q_\ell)$ either contains
$\operatorname{Ev}(H_s,\mathbb Q_\ell)$ or is contained in the image of
$H^{2n}(X,\mathbb Q_\ell)$.  The same assertion holds for $X$ of even
dimension (and for every Lefschetz pencil).

\end{filecontents*}
% Original input: components/136_expose_XIX_section_3.tex
\begin{filecontents*}[overwrite,noheader]{sga7ii-cumulative-input-136.tex}
\section*{3. Proof of 1.3}

\subsubsection*{3.1.}
Under the hypotheses of the theorem, it remains to prove that 1.5(b) does
not hold.  The case $h^{0,2n}\ne0$ will be treated by Katz in XXI 4.1 (the
condition $h^{0,2n}\ne0$ is equivalent to
\[
  \sum_i a_i-(2n+r+1)\geq0,
\]
this number being the integer $m$ such that
$\Omega_V^{2n}\simeq\mathcal O(m)$).
\footnote{The print writes $d_i$ in this sum, although the multidegree is
denoted by $(a_i)$.}
The results of the present expos\'e will not be used again in the remainder
of the seminar, so there is no vicious circle.

\subsubsection*{3.2.}
When $k=\mathbb C$, one proves 1.3 by observing that if 1.5(b) holds, the
cohomology of $V$ is purely of type $(n,n)$, since algebraic cohomology is
of type $(n,n)$.  This occurs only in the exceptional cases excluded in
1.3 (XI 2.9).  By the Lefschetz principle, 1.3 is therefore true in
characteristic zero.

\subsubsection*{Proposition 3.3.}
Assume the hypotheses of 1.4, and suppose that (b) holds.  Let
$\operatorname{Ev}(V,\mathbb Q)$ be the $\mathbb Q$-vector subspace of
$\operatorname{Ev}(V,\mathbb Q_\ell(n))$ generated by the algebraic
cohomology classes.

\begin{enumerate}[label=(\alph*)]
\item There is an isomorphism
  \[
  \operatorname{Ev}(V,\mathbb Q)\otimes\mathbb Q_\ell
  \xrightarrow{\ \sim\ }
  \operatorname{Ev}(V,\mathbb Q_\ell(n)).
  \]
  On $\operatorname{Ev}(V,\mathbb Q)$ the intersection form is rational
  and definite, of sign $(-1)^n$.
\item The vanishing cycles and their conjugates belong to
  $\operatorname{Ev}(V,\mathbb Q)$; there they form an irreducible root
  system of type $A$, $D$, or $E$.
\item The monodromy group (the image of $\pi_1(\ell-S,\bar s)$ in
  $\operatorname{GL}(\operatorname{Ev}(V,\mathbb Q))$) is finite; it is the
  Weyl group of the root system in (b).\footnote{The print repeatedly says
  ``groupe de Weil''; the group attached to a root system is its Weyl
  group.}
\end{enumerate}

Apart from definiteness of the intersection form, (a) follows because the
intersection form is nondegenerate on
$\operatorname{Ev}(V,\mathbb Q_\ell(n))$ (by (LV)) and rational on
$\operatorname{Ev}(V,\mathbb Q)$.

Condition 1.4(b) says equivalently that there are algebraic cycles on $V$,
modulo numerical equivalence, which are not invariant under monodromy; this
condition is therefore independent of $\ell$.  Let $x$ be such an algebraic
cycle, and let $\delta$ be the vanishing cycle associated with $s_0\in S$
and with a path $c$ from $\bar s$ (which defines $V$) to $s_0$.  For a
suitable $c$, $(x,\delta)\ne0$.  The Galois transform
\[
  \sigma(x)=x-(-1)^n(x,\delta)\delta
\]
is again algebraic; hence $(x,\delta)\delta$ is algebraic, the image of a
cycle $y$ independent of $\ell$.  The rational number
\[
  \frac{(-1)^n}{2}(y,y)=(x,\delta)^2
\]
is consequently independent of $\ell$.  It is a square in every
$\mathbb Q_\ell$, hence a square in $\mathbb Q$; thus, up to sign,
$(x,\delta)$ belongs to $\mathbb Q$ and is independent of $\ell$.  This
shows not only that $\delta\in\operatorname{Ev}(V,\mathbb Q)$, but also that
the algebraic cycle modulo numerical equivalence
\[
  \pm\delta=\frac{\sigma(x)-x}{(x,\delta)}
\]
is independent of $\ell$.

An algebraic cycle on a variety $V$ is always defined over a finite
extension of any field of definition of $V$.  It follows that
$\pi_1(\ell-S,s)$ acts on $\operatorname{Ev}(V,\mathbb Q)$ through a finite
group.  This group is generated by the reflections associated with the
vanishing cycles $\delta$ and their conjugates (XVIII 6.6.1 and Section 2).
The cycles $\delta$ and their conjugates therefore form a root system whose
Weyl group is the monodromy group.  Moreover, all roots of this root system
are conjugate under the Weyl group; it is of type $A$, $D$, or $E$.

Up to a scalar, the intersection form is the unique invariant form under
the Weyl group.  It is therefore definite.  Its sign is determined by
computing
\[
  (\delta,\delta)=2(-1)^n.
\]

\subsubsection*{Proposition 3.4.}
If $k=\mathbb C$, then under the hypotheses of 1.4 the following conditions
are equivalent:
\begin{enumerate}[label=(\roman*)]
\item the monodromy group is finite;
\item the vanishing cohomology $\operatorname{Ev}(V)$ is purely of type
  $(n,n)$.
\end{enumerate}

Let $\widetilde X\to\ell$ be the family whose fibers are the hyperplane
sections in the pencil $\ell$, and let $V_s$ be the fiber at
$s\in\ell(\mathbb C)$.  The $\operatorname{Ev}(V_s)$ form an algebraic
variation of Hodge structures on $\ell$.  By the fixed-part theorem
[2, 4.1.2], applied to a finite covering of $\ell$ (cf. [2, 4.1.3.3]), (i)
implies that the Hodge structure on $\operatorname{Ev}(V_s)$ is locally
constant.  The monodromy group therefore acts by automorphisms of the Hodge
structure.  For a Picard--Lefschetz transformation $T$,
$\operatorname{Im}(T-1)$ is a rank-one Hodge substructure, hence is of type
$(n,n)$ and is generated by a vanishing cycle $\delta$.  Since the
vanishing cycles are of type $(n,n)$, all of
$\operatorname{Ev}(V_s)$ is of that type.

Conversely, if $\operatorname{Ev}(V)$ is purely of type $(n,n)$, the
intersection form on $\operatorname{Ev}(V)$ is definite and the monodromy
group is finite (cf. [2, 4.1.3.3]).

\subsubsection*{Remark 3.5.}
If the equivalent conditions of 3.3 hold, the argument of 3.3 shows that
the vanishing cycles and their conjugates form a root system of type $A$,
$D$, or $E$, whose Weyl group is the monodromy group.

\subsubsection*{3.6.}
Let us prove 1.3 when $k$ has finite characteristic different from $2$.
Let $S=\operatorname{Spec}(\mathcal V)$ be the spectrum of a discrete
valuation ring of mixed characteristic with residue field $k$, and let
$X\subset\mathbb P_S^{2n+d}$ be a complete intersection of relative
dimension $2n+1$ and multidegree $(a_1,\ldots,a_{d-1})$ over $S$.  Let
$\ell$ be a Lefschetz pencil of sections of $X$ by hypersurfaces of degree
$a_d$ (in the special fiber, $\ell_s$ is a Lefschetz pencil for $X_s$).
Let
\[
  f:\widetilde X\longrightarrow\ell\simeq\mathbb P^1_S
\]
be the family whose fibers are the sections in the pencil.  The exceptional
locus $T\subset\ell$ is \`etale over $S$ (XVII 6; here one uses
$p\ne2$).  The sheaf $R^{2n}f_*\mathbb Q_\ell$ is locally constant on
$\ell-T$; it is therefore tamely ramified, and its monodromy is the same in
characteristic $p$ as in characteristic zero.  If 1.5(b) holds, it follows
from 3.3(c) and 3.4 that the cohomology group $H^{2n}(V)$ of complex
complete intersections of dimension $2n$ and multidegree
$\underline a$ is purely of type $(n,n)$.  This occurs only in the cases
excluded by 1.3 (XI 2.9).

\end{filecontents*}
% Original input: components/137_expose_XIX_section_4.tex
\begin{filecontents*}[overwrite,noheader]{sga7ii-cumulative-input-137.tex}
\section*{4. A second complement to XVIII}

In this section we work over $\mathbb C$ and with integral, classical
cohomology.  Using methods closer to Lefschetz's, we revisit some of the
results of XVIII.  There they were proved only after tensoring with
$\mathbb Q$ (or rather with $\mathbb Q_\ell$, which amounts to the same
thing).

\subsubsection*{4.1.}
Let $X\subset\mathbb P$ be a smooth complex projective variety, pure of
dimension $n+1$, and let $\ell$ be a Lefschetz pencil of hyperplane
sections of $X$, with axis $A$.  We also write $\ell$ for the parameter
space of the pencil, which is isomorphic to $\mathbb P^1$.  Let
$S\subset\ell$ be the set of exceptional values.  For $s\in\ell$, write
$H_s$ for the corresponding hyperplane section of $X$ (of dimension $n$).

Let $T(A)$ be a closed tubular neighborhood of $A\cap X$ in $X$, whose
boundary is transverse to every $H_s$.

Choose two points $0,\infty\in\ell-S$, and let
$\alpha:\ell\to\mathbb P^1(\mathbb C)$ be an orientation-preserving
diffeomorphism that sends $0$ to $0$, $\infty$ to $\infty$, and $S$ to the
circle of radius one centered at $0$.
\footnote{The print gives $\ell-S$ as the domain of $\alpha$, although the
same sentence evaluates $\alpha$ on $S$; the required domain is $\ell$.}
The composite $\alpha f$ of $\alpha$ with the projection to $\ell$ makes
$X-T(A)^\circ$ a singular fibration in manifolds with boundary over
$\mathbb P^1(\mathbb C)$ (one singularity above each point of
$\alpha(S)$).  Write $c_s$ for the straight segment joining $0$ to
$\alpha(s)$.

Suppressing $\alpha$ from the notation, the picture is
\begin{center}
\begin{tikzpicture}[x=1.2cm,y=1.2cm]
  \coordinate (O) at (0,0);
  \coordinate (L) at (-1.8,1.8);
  \coordinate (S2) at (0,2.8);
  \coordinate (S1) at (2.2,2.0);
  \draw (O)--(L);
  \draw (O)--(S2);
  \draw (O)--(S1);
  \node[right] at (O) {$0$};
  \node[above right] at (S2) {$s_2$};
  \node[above right] at (S1) {$s_1$};
  \node[right] at (0,1.65) {$c_2$};
  \node[below right] at (1.25,1.14) {$c_1$};
  \fill (-0.73,0.76) circle (0.035);
  \fill (-0.55,0.91) circle (0.035);
  \fill (-0.37,1.03) circle (0.035);
\end{tikzpicture}
\end{center}

Near $s\in S$ and the singular point $x_s$ of $H_s\subset X$, suitable
local coordinates put us in the situation studied in XIV 3.2.  If
$\alpha$ is chosen holomorphic near $s$, one can even choose local
coordinates on $X$ and $\ell$ in which $f$ is written
$\sum z_i^2$ and $\alpha^{-1}(c_s)$ is given by ``the $z_i$ are real and
positive.''  Above the intersection of a neighborhood of $s$ with
$\alpha^{-1}(c_s)-\{s\}$, consider the bundle of real spheres
$\Delta'_s\subset X$ defined by ``the $z_i$ are real.''  As $t\to s$, the
fiber $\Delta'_{st}$ of $\Delta'_s$ at $t$ tends to $x_s$.

Let $\Delta_s$ denote the union of $\{x_s\}$ and a sphere bundle over
$\alpha^{-1}(c_s)-\{s\}$, contained in $X-T(A)$ and extending
$\Delta'_s$.  Then $\Delta_s$ is an embedded disk in $X$ whose boundary
$H_0\cap\Delta_s$ is a sphere representing the vanishing cycle $\delta_s$
at $s$, transported to $H_0$ along the path $\alpha^{-1}(c_s)$.

Apply Morse theory to the function $|\alpha f|^2$ to study how the
homotopy type of
\[
 \{x\in X\mid x\in T(A)\ \text{or}\ |\alpha f(x)|^2\leq r\}
\]
changes as $r$ increases.  Put
\[
 X^- = \{x\in X\mid x\in T(A)\ \text{or}\ \alpha f(x)\ne\infty\}.
\]
One finds that the inclusions
\[
 H_0\cup\bigcup_{s\in S}\Delta_s
 \hookrightarrow
 T(A)\cup H_0\cup\bigcup_{s\in S}\Delta_s
 \hookrightarrow X^-
\]
are homotopy equivalences.  Morse theory also shows that $X-X^-$ can be
deformation-retracted onto a subcomplex of real dimension at most $n$.
Altogether this gives the following result.

\subsubsection*{Theorem 4.2.}
There exists a compact subcomplex $\Sigma$ of real dimension at most $n$ in
$H_\infty^{\mathrm{aff}}=H_\infty-A$ such that the inclusion
\[
 H_0\cup\bigcup_{s\in S}\Delta_s\longrightarrow X-\Sigma
\]
is a homotopy equivalence.

\subsubsection*{4.3.}
Removing $\Sigma$ has no effect on the cohomology, homology, or homotopy
groups of $X$ in degrees $i\leq n$.  Since
$H_0\cup\bigcup_{s\in S}\Delta_s$ is obtained from $H_0$ by attaching
$(n+1)$-disks along spheres representing the vanishing cycles, homology
gives\footnote{The print writes $V_0$ in this sequence and in the diagram
below, but Section 4 defines the fiber as $H_0$.}
\[
 \mathbb Z^S\longrightarrow H_n(H_0)\longrightarrow H_n(X)
 \longrightarrow0.
\]
The analogous exact sequence in cohomology, together with the sequence
deduced from it by Poincar\'e duality, forms the commutative diagram
\begin{equation}\tag{4.3.1}
\begin{tikzcd}[row sep=large,column sep=large]
{} & \mathbb Z^S & {} & {} \\
\mathbb Z^S \arrow[r,"{\sum n_s\delta_s}"] &
H^n(H_0) \arrow[u,"{(\ ,\delta_s)}"'] \arrow[r] &
H^{n+2}(X) \arrow[r] & 0 \\
{} & H^n(X) \arrow[u] \arrow[ur,"{\eta\wedge}"'] & {} & {} \\
{} & 0 \arrow[u] & {} & {}
\end{tikzcd}
\end{equation}
\end{filecontents*}
% Original input: components/138_expose_XIX_section_5_and_bibliography.tex
\begin{filecontents*}[overwrite,noheader]{sga7ii-cumulative-input-138.tex}
\section*{5. The exceptional cases of the Noether theorem}

In this section we continue to work over $\mathbb C$, with integral
cohomology.

\subsubsection*{5.1.}
Let $2n$, $d$, and $\underline a=(a_1,\ldots,a_d)$, with $d\geq1$ and
$a_i\geq2$, be such that the middle cohomology of smooth complete
intersections $V$ of dimension $2n$ and multidegree $\underline a$ is
purely of type $(n,n)$.  In other words, either $n=0$, or
$(2n;\underline a)$ is one of $(2n;2)$, $(2n;2,2)$, and $(2;3)$
(XI 2.9).

Suppose that $V$ is obtained from a complete intersection
$X\subset\mathbb P^{2n+d}$ (respectively,
$X\subset\mathbb P^{2n+d+1}$) of dimension $2n+1$ and multidegree
$(a_1,\ldots,a_{d-1})$ (respectively, $(a_1,\ldots,a_d)$) by taking a
section by a hypersurface of degree $a_d$ (respectively, a hyperplane
section) belonging to a Lefschetz pencil $\ell$.

By (4.3.1), the vanishing part $\operatorname{Ev}(V)$ of $H^{2n}(V)$,
which equals its primitive part, is generated over $\mathbb Z$ by the
vanishing cycles.  It is the orthogonal complement of the power $\eta^n$
of the cohomology class $\eta$ of a hyperplane section.

Since the intersection form on $H^{2n}(V)$ is unimodular and
$\eta^n$\footnote{The print has $\eta^{2n}$ in the primitivity assertion,
but that class lies in degree $4n$; both the surrounding discussion and the
next display require $\eta^n\in H^{2n}(V)$.} is not $b y$, with $b>1$, for
another cohomology class $y$ (XI 1.6(iv)), the
discriminant of the intersection form on $\operatorname{Ev}(V)$ is
\[
 (\eta^n,\eta^n)=\prod_i a_i.
\]

Let $R$ be the root system in $H^{2n}(V)$ formed by the vanishing cycles
and their conjugates.  We have the following information about $R$:
\begin{enumerate}[label=(\alph*)]
\item $R$ is empty or is of type $A$, $D$, or $E$;
\item the rank of $R$ is given in XI 1.10;
\item its discriminant is $\prod_i a_i$.  In the tables, this discriminant
  is the order of the center of the corresponding simply connected complex
  simple group.  The argument by which we computed the discriminant shows
  that this center, dual to the quotient of the weight lattice by the root
  lattice, is cyclic.
\end{enumerate}

In every case this information determines $R$ completely if $n\ne0$.  One
finds:

\subsubsection*{Proposition 5.2.}
\begin{enumerate}[label=(\roman*)]
\item If $V$ is an even-dimensional quadric, $R$ is of type $A_1$.
\item If $V$ is a cubic surface, $R$ is of type $E_6$.
\item If $V$ has dimension $2n$ and is the intersection of two quadrics,
  $R$ is of type $D_{2n+3}$.
\end{enumerate}

\subsubsection*{5.3.}
Suppose that $n\ne0$.  The set $R^1$ of vectors of length $\pm2$ in
$\operatorname{Ev}(V)$ is also a root system and has the same properties
(a), (b), and (c) as $R$.  Hence $R=R^1$.  Let $\pi_1$ be the fundamental
group of the parameter space of all complete intersections
$V\subset\mathbb P^{2n+d}$ of the type under consideration.  The image of
$\pi_1$ in $\operatorname{GL}(\operatorname{Ev}(V))$ is the Weyl group of
$R$.  In every case it is the automorphism group of the triple
\[
 \bigl(H^{2n}(V),\eta^n,\text{intersection form}\bigr).
\]

\subsubsection*{5.4.}
With $R$ replaced by $R^1$ from 5.3, assertion 5.2(ii) is well known [1],
and for small values of $n$, 5.2(iii) was known, at least to Manin
($n=0$ is trivial; a surface that is the intersection of two quadrics is a
del Pezzo surface).\footnote{The print refers here to 4.2(ii) and 4.2(iii),
but the assertions cited are the immediately preceding cases of
Proposition 5.2.}  A detailed study of the intersection of two quadrics was
made by M. Reid [4].  He deduces 5.2(iii) from a study of the geometry of
the $n$-dimensional linear subspaces on an intersection of two quadrics
$V\subset\mathbb P^{2n+2}$.

\section*{Bibliography}

\begin{enumerate}[label={[\arabic*]}]
\item E. Cartan, ``Quelques remarques sur les 28 bitangentes d'une quartique
  plane et les 27 droites d'une surface cubique,'' \emph{Bull. Sci. Math.}
  \textbf{70} (1946), 42--45; \emph{Oeuvres}, part I, vol. 2, p.~1353.
\item P. Deligne, ``Th\'eorie de Hodge II,'' \emph{Publ. Math. IHES}
  \textbf{40} (1971), 5--58.
\item S. Lefschetz, \emph{L'analysis situs et la g\'eom\'etrie alg\'ebrique},
  Gauthier--Villars, 1924 (reprinted in \emph{Selected Papers}, Chelsea).
\item M. Reid, Ph.D. thesis, Cambridge, 1972.
\item J.-P. Serre, ``Faisceaux alg\'ebriques coh\'erents,''
  \emph{Ann. of Math.} \textbf{61} (1955), 197--278 (cited as FAC).
\end{enumerate}
\end{filecontents*}
% Original input: components/139_expose_XX_opening_through_corollary_1_4.tex
\begin{filecontents*}[overwrite,noheader]{sga7ii-cumulative-input-139.tex}
\clearpage
\begin{center}
{\large SGA 7 --- EXPOS\'E XX}\par
\vspace{2em}
{\Large\underline{The Griffiths Theorem}}\par
\vspace{1.5em}
{\large by N. Katz}
\end{center}

\section*{0. Introduction}

We prove Griffiths's theorem~[1], which gives a systematic (though
nonconstructive) way of ``finding'' algebraic cycles that are
$\mathbb Q_\ell$-cohomologous to zero, but no multiple of which is
algebraically equivalent to zero.  The proof given here (due to
Grothendieck) is the translation into purely algebraic terms of Griffiths's
original, more or less transcendental, proof.  In both proofs, the
irreducibility theorem XVIII~6.7 is a key point (cf.~3.1.6 and~3.2).

\section*{1. The formalism of primitive classes}

Throughout what follows, a prime number $\ell$ is fixed.  All schemes under
consideration are assumed to have characteristic different from $\ell$.
Consider the situation

\begin{equation}\tag{1.0.1}
\begin{tikzcd}[column sep=huge, row sep=huge]
X \arrow[rr, "\pi"] \arrow[dr, "f"'] & {} & Y \arrow[dl, "g"] \\
{} & S & {}
\end{tikzcd}
\end{equation}
in which we assume
\begin{equation}\tag{1.0.3}
f,\ g \quad\text{smooth of relative dimensions}\quad d_1\ \text{and}\ d_2,
\end{equation}
and
\begin{equation}\tag{1.0.4}
\pi\quad\text{proper}.
\end{equation}
Since $X$ and $Y$ are smooth over $S$, the properness of $\pi$ gives a
\underline{Gysin morphism} (SGA~5, Expos\'e~IV; the further citation is left
blank in the print)
\begin{equation}\tag{1.0.5}
\pi_*:H^q(X,\mathbb Q_\ell(i))\longrightarrow
H^{q-2(d_1-d_2)}(Y,\mathbb Q_\ell(i-d_1+d_2)).
\end{equation}
Localizing on $S$ in the \`etale topology, we likewise obtain a morphism of
sheaves on $S$\footnote{In (1.0.6) the print has an apostrophe in place of
the closing parenthesis in the exponent and has $i-d_1-d_2$ in the target
twist.  Formulae (1.0.5) and (1.0.7), as well as the Gysin degree, give the
reading displayed here.}
\begin{equation}\tag{1.0.6}
\pi_*:R^qf_*\mathbb Q_\ell(i)\longrightarrow
R^{q-2(d_1-d_2)}g_*\mathbb Q_\ell(i-d_1+d_2),
\end{equation}
and, on applying the functor $H^p(S,-)$, a morphism
\begin{equation}\tag{1.0.7}
\pi_*:H^p(S,R^qf_*\mathbb Q_\ell(i))\longrightarrow
H^p(S,R^{q-2(d_1-d_2)}g_*\mathbb Q_\ell(i-d_1+d_2)).
\end{equation}

\subsubsection*{Definition 1.0.8.}
Let $(E_r^{p,q},d_r^{p,q})$ and
$({}'E_r^{p,q},{}'d_r^{p,q})$ be two spectral sequences in a given abelian
category, and let $(a,b)\in\mathbb Z\times\mathbb Z$.  The shift of
$({}'E_r^{p,q},{}'d_r^{p,q})$ by $(a,b)$ is the spectral sequence
$({}''E_r^{p,q},{}''d_r^{p,q})$ given by\footnote{The print omits the prime
on the right-hand differential in (1.0.9).  The definition of the shift of
the primed spectral sequence requires the primed differential.}
\begin{equation}\tag{1.0.9}
\left\{
\begin{aligned}
{}''E_r^{p,q}&={}'E_r^{p+a,q+b},\\
{}''d_r^{p,q}&={}'d_r^{p+a,q+b}.
\end{aligned}
\right.
\end{equation}

\subsubsection*{1.0.10.}
A morphism of spectral sequences from
$(E_r^{p,q},d_r^{p,q})$ to $({}''E_r^{p,q},{}''d_r^{p,q})$ is called a
morphism of bidegree $(a,b)$ from $(E_r^{p,q},d_r^{p,q})$ to
$({}'E_r^{p,q},{}'d_r^{p,q})$.

The following lemma is an immediate consequence of the theory of the
\underline{Gysin morphism} (SGA~5, Expos\'e~IV).

\subsubsection*{Lemma 1.1.}
There is a canonical morphism $\pi_*$ of bidegree
$(0,-2(d_1-d_2))$\footnote{The print has $(a,-2(d_1-d_2))$.  The map
(1.0.7) leaves the first Leray index $p$ unchanged, so the first component
of the bidegree is $0$.} between the Leray spectral sequences
\begin{equation}\tag{1.1.0}
R^p\alpha_*R^qf_*\mathbb Q_\ell(i)\Longrightarrow
R^{p+q}(\alpha f)_*\mathbb Q_\ell(i)
\end{equation}
and
\begin{equation}\tag{1.1.1}
R^p\alpha_*R^qg_*\mathbb Q_\ell(i-d_1+d_2)\Longrightarrow
R^{p+q}(\alpha g)_*\mathbb Q_\ell(i-d_1+d_2),
\end{equation}
which is given on the $E_2$ term by (1.0.7), and on the abutment by
(1.0.5).

\subsubsection*{1.2.}
Now consider the situation
\begin{equation}\tag{1.2.1}
\begin{tikzcd}[column sep=large, row sep=large]
{} & X\times_S Y \arrow[dl, "\mathrm{pr}_1"']
                  \arrow[dr, "\mathrm{pr}_2"] & {} \\
X \arrow[dr, "f"'] & {} & Y \arrow[dl, "g"] \\
{} & S \arrow[d, "\alpha"] & {} \\
{} & T & {}
\end{tikzcd}
\end{equation}
in which $S,X,f,Y,g$ are as in (1.0.1)--(1.0.4), but now $f$ is also
assumed proper.\footnote{The print begins this range at the nonexistent
(1.0.2); the displayed situation is (1.0.1).}  Let
\begin{equation}\tag{1.2.2}
z\in H^a(X\times_S Y,\mathbb Q_\ell(b)).
\end{equation}
Then $z$ induces a map, again denoted by $z$,
\begin{equation}\tag{1.2.3}
z:H^q(X,\mathbb Q_\ell(i))\longrightarrow
H^{q+a-2d_1}(Y,\mathbb Q_\ell(i+b-d_1))
\end{equation}
by the formula
\begin{equation}\tag{1.2.4}
z(x)=\mathrm{pr}_{2*}\bigl(z\wedge\mathrm{pr}_1^*(x)\bigr),
\end{equation}
where $\mathrm{pr}_{2*}$ is the Gysin homomorphism.  By \`etale localization
on $S$, the class $z$ likewise defines
\begin{equation}\tag{1.2.5}
z:R^qf_*\mathbb Q_\ell(i)\longrightarrow
R^{q+a-2d_1}g_*\mathbb Q_\ell(i+b-d_1),
\end{equation}
and, after applying $H^p(S,-)$, maps
\begin{equation}\tag{1.2.6}
H^p(S,R^qf_*\mathbb Q_\ell(i))\longrightarrow
H^p(S,R^{q+a-2d_1}g_*\mathbb Q_\ell(i+b-d_1)).
\end{equation}

\subsubsection*{Lemma 1.3.}
The class $z$ defines a morphism of bidegree
$(0,a-2d_1)$\footnote{The print has $(0,b-2d_1)$.  The maps (1.2.3),
(1.2.5), and (1.2.6) change cohomological degree by $a-2d_1$; $b$ controls
the Tate twist.} between the Leray spectral sequences
\begin{equation}\tag{1.2.7}
H^p(S,R^qf_*\mathbb Q_\ell(i))\Longrightarrow
H^{p+q}(X,\mathbb Q_\ell(i))
\end{equation}
and
\begin{equation}\tag{1.2.8}
H^p(S,R^qg_*\mathbb Q_\ell(i+b-d_1))\Longrightarrow
H^{p+q}(Y,\mathbb Q_\ell(i+b-d_1)),
\end{equation}
which is given on the $E_2$ terms by (1.2.6), and on the abutments by
(1.2.3).

\noindent\underline{Proof.}
In view of (1.2.4), and since ``multiplication by $z$'' defines a morphism
of bidegree $(0,a)$ between the Leray spectral sequences
\begin{equation}\tag{1.2.9}
H^p\bigl(S,R^q(f\circ\mathrm{pr}_1)_*\mathbb Q_\ell(i)\bigr)
\Longrightarrow H^{p+q}(X\times_S Y,\mathbb Q_\ell(i))
\end{equation}
\begin{center}
\begin{tikzcd}
{} \arrow[d, "z"'] \\
{}
\end{tikzcd}
\end{center}
\begin{equation}\tag{1.2.10}
H^p\bigl(S,R^{q+a}(f\circ\mathrm{pr}_1)_*\mathbb Q_\ell(i+b)\bigr)
\Longrightarrow H^{p+q+a}(X\times_S Y,\mathbb Q_\ell(i+b)),
\end{equation}
it suffices to apply Lemma~1.1 to the situation\footnote{In (1.2.9) and
(1.2.10) the print writes $(\mathrm{pr}_1 f)_*$, whereas the map to $S$ is
$f\circ\mathrm{pr}_1$, as the printed diagram (1.2.11) itself shows.}
\begin{equation}\tag{1.2.11}
\begin{tikzcd}[column sep=large, row sep=large]
X\times_S Y \arrow[rr, "\mathrm{pr}_2"]
             \arrow[dr, "f\circ\mathrm{pr}_1"]
  & {} & Y \arrow[dl, "g"] \\
{} & S & {}
\end{tikzcd}
\qquad .
\end{equation}

\subsubsection*{Corollary 1.4.}
In the situation (1.2.1)--(1.2.2), the morphism (1.2.3)
\begin{equation}\tag{1.4.1}
z:H^q(X,\mathbb Q_\ell(i))\longrightarrow
H^{q+a-2d_1}(Y,\mathbb Q_\ell(i+b-d_1))
\end{equation}
takes primitive classes (XVIII~5.8.1) to primitive classes, and the following
diagram is commutative:
\begin{equation}\tag{1.4.2}
\begin{tikzcd}[column sep=large, row sep=large]
\mathrm{Prim}^q(X/S,\mathbb Q_\ell(i))
  \arrow[r, "z"] \arrow[d, "u_f"]
& \mathrm{Prim}^{q+a-2d_1}(Y/S,\mathbb Q_\ell(i+b-d_1))
  \arrow[d, "u_g"] \\
H^1(S,R^{q-1}f_*\mathbb Q_\ell(i))
  \arrow[r, "z"]
& H^1(S,R^{q+a-2d_1-1}g_*\mathbb Q_\ell(i+b-d_1)).
\end{tikzcd}
\end{equation}
Here $u_f$ and $u_g$ are the canonical homomorphisms of XVIII~5.8.1.
\end{filecontents*}
% Original input: components/140_expose_XX_sections_1_5_through_1_7.tex
\begin{filecontents*}[overwrite,noheader]{sga7ii-cumulative-input-140.tex}
\subsubsection*{1.5.}
Let $\bar\eta$ be a geometric point of $S$ (lying over a point $\eta$ of
$S$), and let
\begin{equation}\tag{1.5.1}
z_{\bar\eta}\in
H^a\bigl(X_{\bar\eta}\times_{k(\bar\eta)}Y_{\bar\eta},
\mathbb Q_\ell(b)\bigr)
\end{equation}
be the class induced by $z$ on the corresponding geometric fiber of
\begin{equation}\tag{1.5.2}
X\times_S Y\longrightarrow S.
\end{equation}

\subsubsection*{Proposition 1.5.3.}
Suppose that $f$ and $g$ are proper, that $S$ is connected,\footnote{The
print has the nonsensical adjective ``annexe'' in place of ``connexe'', and
uses the capital $X$ for the element in the conclusion.  The proof requires
connectedness of $S$, while (1.5.4) itself uses the element $x$.} and let
$q\in\mathbb Z$.  If the class $z_{\bar\eta}$ induces the zero
homomorphism
\begin{equation}\tag{1.5.3}
0=z_{\bar\eta}:H^{q-1}(X_{\bar\eta},\mathbb Q_\ell(i))
\longrightarrow
H^{q+a-2d_1-1}(Y_{\bar\eta},\mathbb Q_\ell(i+b-d_1)),
\end{equation}
then, for every
$x\in\mathrm{Prim}^q(X/S,\mathbb Q_\ell(i))$, one has
\begin{equation}\tag{1.5.4}
u_g(z(x))=0\quad\text{in}\quad
H^1\bigl(S,R^{q+a-2d_1-1}g_*\mathbb Q_\ell(i+b-d_1)\bigr).
\end{equation}

\noindent\underline{Proof.}
By the commutativity of (1.4.2), it is enough to see that $z$ induces the
zero homomorphism
\begin{equation}\tag{1.5.5}
z:R^{q-1}f_*\mathbb Q_\ell(i)
\stackrel{0}{\longrightarrow}
R^{q+a-2d_1-1}g_*\mathbb Q_\ell(i+b-d_1).
\end{equation}
Indeed, since $f$ and $g$ are proper and smooth, the sheaves
$R^\bullet f_*\mathbb Q_\ell(\mathord\cdot)$ and
$R^\bullet g_*\mathbb Q_\ell(\mathord\cdot)$ are twisted constant sheaves
(SGA~4, XVI~2.2), hence are equivalent to their values at $\bar\eta$,
regarded as modules under $\pi_1(S,\bar\eta)$.

\subsubsection*{1.6. A construction of Manin}
Suppose that two sections of $f:X\longrightarrow S$ are given:
\begin{equation}\tag{1.6.0}
\begin{tikzcd}[column sep=huge]
S \arrow[r, bend left=45, "t_1"]
  \arrow[r, bend right=45, "t_2"']
& X \arrow[l, "f" description]
\end{tikzcd}
\qquad ,
\end{equation}
and associate with each its cohomology class
\begin{equation}\tag{1.6.1}
\mathrm{cl}(t_i)=t_{i*}(1)
\in H^{2d_1}(X,\mathbb Q_\ell(d_1)).
\end{equation}
If the fibers of $f:X\longrightarrow S$ are geometrically connected
(more generally, if, for $s\in S$, $t_1(s)$ and $t_2(s)$ belong to the same
connected component of $X_s$), it is clear that\footnote{The print says
$f(s)$ and $g(s)$ here, although the two maps under discussion are the
sections $t_1$ and $t_2$.}
\begin{equation}\tag{1.6.2}
\mathrm{cl}(t_1)-\mathrm{cl}(t_2)
\in\mathrm{Prim}^{2d_1}(X/S,\mathbb Q_\ell(d_1)).
\end{equation}
We denote by
\begin{equation}\tag{1.6.3}
u_f(t_1-t_2)\in
H^1(S,R^{2d_1-1}f_*\mathbb Q_\ell(d_1))
\end{equation}
its image under
\begin{equation}\tag{1.6.4}
u_f:\mathrm{Prim}^{2d_1}(X/S,\mathbb Q_\ell(d_1))
\longrightarrow H^1(S,R^{2d_1-1}f_*\mathbb Q_\ell(d_1)).
\end{equation}
(It is an analogue $\mu(t_1-t_2)$ of this element that Manin studies
in~[5].)

\subsubsection*{1.6.5.}
Now let $z$ be an algebraic cycle on $X\times_S Y$, of codimension at least
$b$ on each fiber of $X\times_S Y$,\footnote{If $S$ is smooth over a scheme
$T$ (for example $T=\operatorname{Spec}k$, with $k$ a field), it would be
enough to suppose that $z$ has codimension at least $b$ on each fiber over
$T$.} which determines a cohomology class (SGA~5, Expos\'e~IV; the further
citation is blank in the print)
\begin{equation}\tag{1.6.6}
\mathrm{cl}(z)\in H^{2b}(X\times_S Y,\mathbb Q_\ell(b)).
\end{equation}
This class defines a map (cf.~(1.2.3))
\begin{equation}\tag{1.6.7}
\mathrm{cl}(z)\mathbin{\cdot}H^{2d_1}(X,\mathbb Q_\ell(d_1))
\longrightarrow H^{2b}(Y,\mathbb Q_\ell(b)),
\end{equation}
and, if the cycles $z(t_1)$ and $z(t_2)$ are defined, the compatibility of
the cup product with intersections of cycles (SGA~5, Expos\'e~IV; the
further citation is blank in the print) gives
\begin{equation}\tag{1.6.8}
\mathrm{cl}(z)\bigl(\mathrm{cl}(t_1)-\mathrm{cl}(t_2)\bigr)
=\mathrm{cl}\bigl(z(t_1)-z(t_2)\bigr).
\end{equation}
The commutativity of (1.4.2) gives us
\begin{equation}\tag{1.6.9}
\mathrm{cl}(z)\bigl(u_f(t_1-t_2)\bigr)
=u_g\bigl(\mathrm{cl}(z(t_1)-z(t_2))\bigr).
\end{equation}
As a corollary of Proposition~1.5.3, we therefore have:

\subsubsection*{Proposition 1.7.}
Let $t_1$ and $t_2$ be two sections of $f:X\longrightarrow S$, with $f$
assumed to have connected fibers, and let $z$ be an algebraic cycle on
$X\times_S Y$ of codimension $b$, such that the algebraic cycle
\begin{equation}\tag{1.7.1}
z(t_1)-z(t_2)
\end{equation}
is defined.  Then its cohomology class is primitive (XVIII~5.8.1).  If the
algebraic cycle $z_{\bar\eta}$ on
$X_{\bar\eta}\times Y_{\bar\eta}$ induces the zero homomorphism
\begin{equation}\tag{1.7.2}
z_{\bar\eta}:H^{2d_1-1}(X_{\bar\eta},\mathbb Q_\ell(d_1))
\stackrel{0}{\longrightarrow}
H^{2b-1}(Y_{\bar\eta},\mathbb Q_\ell(b)),
\end{equation}
then
\begin{equation}\tag{1.7.3}
u_g\bigl(z(t_1)-z(t_2)\bigr)=0\quad\text{in}\quad
H^1(S,R^{2b-1}g_*\mathbb Q_\ell(b)),
\end{equation}
where $u_g$ is the homomorphism of XVIII~5.8.1 occurring in~1.4.
\footnote{The sentence in the print uses an upper-case $Z_{\bar\eta}$,
whereas (1.7.2) and the cycle used throughout the section are denoted by
$z_{\bar\eta}$.}
\end{filecontents*}
% Original input: components/141_expose_XX_section_2.tex
\begin{filecontents*}[overwrite,noheader]{sga7ii-cumulative-input-141.tex}
\section*{2. A reminder on level}
\noindent(cf.~[2], 9.7 and 10.1).

\subsubsection*{2.0.1.}
Let $X$ be proper and smooth over an algebraically closed field $k$.  We saw
in (1.2.3) that, for every
\begin{equation}\tag{2.0.2}
T\quad\text{proper, smooth, and connected over $k$, of dimension $d$},
\end{equation}
and every
\begin{equation}\tag{2.0.3}
\text{algebraic cycle $z$ of codimension $b$ on $T\times X$},
\end{equation}
the cohomology class
\begin{equation}\tag{2.0.4}
\mathrm{cl}(z)\in H^{2b}(T\times X,\mathbb Q_\ell(b))
\end{equation}
gives rise to maps
\begin{equation}\tag{2.0.5}
\mathrm{cl}(z):H^q(T,\mathbb Q_\ell(i))\longrightarrow
H^{q+2b-2d}(X,\mathbb Q_\ell(i+b-d)).
\end{equation}
The \underline{coniveau filtration} $(N^j)$ on
$H^q(X,\mathbb Q_\ell(i))$ is defined by
\begin{equation}\tag{2.0.6}
N^jH^q(X,\mathbb Q_\ell(i))
=\sum_{(T,z)}\operatorname{Im}\!\left(
\mathrm{cl}(z):H^{q-2j}(T,\mathbb Q_\ell(i-j))
\longrightarrow H^q(X,\mathbb Q_\ell(i))\right),
\end{equation}
where the sum extends over the pairs
\begin{equation}\tag{2.0.7}
(T,z),\qquad
\begin{array}{l}
T/k\text{ proper, smooth, and connected, of dimension }d,\\
z\text{ an algebraic cycle on }T\times X\text{ of codimension }j+d.
\end{array}
\end{equation}
It should be noted that one sometimes speaks of the increasing filtration
$F^j$ ``by level'', defined by
\begin{equation}\tag{2.0.8}
F^jH^q(X,\mathbb Q_\ell(i))
=N^{\left[\frac{q-j}{2}\right]}H^q(X,\mathbb Q_\ell(i)).
\end{equation}
Replacing $(T,z)$ by $(\mathbb P^1\times T,\{\mathrm{pt}\}\times z)$, one
sees that
\begin{equation}\tag{2.0.9}
N^jH^q(X,\mathbb Q_\ell(i))
\subset N^{j-1}H^q(X,\mathbb Q_\ell(i)).
\end{equation}
Obviously, one has:

\subsubsection*{Proposition 2.1.}
Let $X$ be proper, smooth, and geometrically connected over a field $k$, let
$\bar k$ be an algebraic closure of $k$, and put
$X_{\bar k}=X\otimes_k\bar k$.  Then $\operatorname{Gal}(\bar k/k)$,
acting on $H^q(X_{\bar k},\mathbb Q_\ell(i))$, preserves the coniveau
filtration.

Standard arguments (cf.~[2]) prove the following statements.

\subsubsection*{Proposition 2.2.}
The level filtration is invariant under an extension $k'/k$ of
algebraically closed fields.

\subsubsection*{Proposition 2.3.}
Let $S$ be a reduced and irreducible scheme, with generic point $\eta$, and
let
\begin{equation*}
\pi:X\longrightarrow S
\end{equation*}
be a proper and smooth morphism with geometrically connected fibers.  If,
for given integers $j$ and $q$, one has on the geometric generic fiber
$X_{\bar\eta}$
\begin{equation}\tag{2.3.1}
H^q(X_{\bar\eta},\mathbb Q_\ell)
=N^jH^q(X_{\bar\eta},\mathbb Q_\ell),
\end{equation}
then there exists a nonempty open subset $U\subset S$ such that, at every
point $u\in U$, one has on the geometric fiber $X_{\bar u}$
\begin{equation}\tag{2.3.2}
H^q(X_{\bar u},\mathbb Q_\ell)
=N^jH^q(X_{\bar u},\mathbb Q_\ell).
\end{equation}

\end{filecontents*}
% Original input: components/142_expose_XX_sections_3_1_and_3_2.tex
\begin{filecontents*}[overwrite,noheader]{sga7ii-cumulative-input-142.tex}
\section*{3. Application to Lefschetz pencils}

\subsubsection*{3.1.}
From now on, we fix
\begin{equation}\tag{3.1.0}
k,\quad\text{an algebraically closed field},
\end{equation}
and
\begin{equation}\tag{3.1.1}
\begin{aligned}
X\subset\mathbb P^r,\quad
&\text{a projective, smooth, connected }k\text{-scheme}\\
&\text{of even dimension }n\text{, satisfying (LV) (XVIII~5.2.2).}
\end{aligned}
\end{equation}
Take
\begin{equation}\tag{3.1.2}
D=\{H_t\}_{t\in\mathbb P^1},
\quad\text{a Lefschetz pencil of hypersurfaces of degree $d$},
\end{equation}
with generic fiber
\begin{equation}\tag{3.1.3}
X_\eta,\qquad \eta=\text{the generic point of }\mathbb P^1.
\end{equation}
Put (cf.~XVIII~5.4.4)
\begin{equation}\tag{3.1.4}
m:X_{\bar\eta}\hookrightarrow X\times_k k(\bar\eta),
\quad\text{the inclusion},
\end{equation}
and
\begin{equation}\tag{3.1.5}
E^{n-1}(X_{\bar\eta},\mathbb Q_\ell(i))
=\bigl(m^*H^{n-1}(X,\mathbb Q_\ell(i))\bigr)^\perp
\subset H^{n-1}(X_{\bar\eta},\mathbb Q_\ell(i)).
\end{equation}
We saw in XVIII~6.7.2 that the \underline{Picard--Lefschetz formula}
implies\footnote{The print has $E^{2n-1}$ in (3.1.6).  The space defined in
(3.1.5), and used throughout the proof of Proposition~3.2, is
$E^{n-1}$.}
\begin{equation}\tag{3.1.6}
\begin{gathered}
\text{If }n\text{ is even, every finite-index subgroup of }
\operatorname{Gal}(k(\bar\eta)/k(\eta))\\
\text{acts irreducibly on }E^{n-1}(X_{\bar\eta},\mathbb Q_\ell(i)).
\end{gathered}
\end{equation}

\subsubsection*{Proposition 3.2.}
Let $D=\{H_t\}$ be a Lefschetz pencil, and suppose that $n=2m$ is even.
Then either
\begin{equation}\tag{3.2.1}
H^{n-1}(X_{\bar\eta},\mathbb Q_\ell(i))
\subset m^*H^{n-1}(X,\mathbb Q_\ell(i))
+N^{m-1}H^{n-1}(X_{\bar\eta},\mathbb Q_\ell(i)),
\end{equation}
or else:

\subsubsection*{(3.2.2)}
For every curve $C$ proper and smooth over $k(\bar\eta)$ and every
algebraic cycle $z$ on $C\times X_{\bar\eta}$ of codimension $m$, the
induced homomorphism (1.2.3)
\begin{equation}\tag{3.2.3}
\mathrm{cl}(z):H^1(C,\mathbb Q_\ell(1+i-m))
\longrightarrow H^{2m-1}(X_{\bar\eta},\mathbb Q_\ell(i))
\end{equation}
has its image in $m^*H^{2m-1}(X,\mathbb Q_\ell(i))$.

\noindent\underline{Proof.}
Suppose that the image of such a homomorphism (3.2.3) is not contained in
$m^*H^{n-1}(X,\mathbb Q_\ell(i))$, or, equivalently, that its image in
$E^{n-1}(X_{\bar\eta},\mathbb Q_\ell(i))$ is nonzero.  This image is
evidently stable under the action of a finite-index subgroup of
$\operatorname{Gal}(k(\bar\eta)/k(\eta))$.  The irreducibility (3.1.6) of
$E^{n-1}(X_{\bar\eta},\mathbb Q_\ell(i))$ therefore says that
\begin{equation*}
H^{n-1}(X_{\bar\eta},\mathbb Q_\ell(i))
=m^*H^{n-1}(X,\mathbb Q_\ell(i))+\operatorname{Im}(3.2.3).
\end{equation*}
By definition,
\begin{equation*}
\operatorname{Im}(3.2.3)
=\mathrm{cl}(z)\bigl(H^1(C,\mathbb Q_\ell(1+i-m))\bigr)
\subset N^{m-1}H^{n-1}(X_{\bar\eta},\mathbb Q_\ell(i)).
\end{equation*}
Consequently,
\begin{equation*}
H^{n-1}(X_{\bar\eta},\mathbb Q_\ell(i))
=m^*H^{n-1}(X,\mathbb Q_\ell(i))
+N^{m-1}H^{n-1}(X_{\bar\eta},\mathbb Q_\ell(i)),
\end{equation*}
which is precisely the first alternative.
\end{filecontents*}
% Original input: components/143_expose_XX_sections_3_3_and_3_4.tex
\begin{filecontents*}[overwrite,noheader]{sga7ii-cumulative-input-143.tex}
\subsubsection*{3.3. Condition (B)}

\paragraph{3.3.1.}
We shall say that the inclusion
\begin{equation*}
m:X_{\bar\eta}\hookrightarrow X\times_k k(\bar\eta)
\end{equation*}
satisfies condition (B) if there exists
\begin{equation}\tag{3.3.2}
V_{\bar\eta},\quad
\text{an algebraic cycle of codimension $n-1$ on }
X_{\bar\eta}\times\bigl(X\times_k k(\bar\eta)\bigr),
\end{equation}
such that the map it defines
\begin{equation}\tag{3.3.3}
V_{\bar\eta}:H^{n-1}(X_{\bar\eta},\mathbb Q_\ell(i))
\longrightarrow H^{n-1}(X,\mathbb Q_\ell(i))
\end{equation}
is a left inverse of
\begin{equation}\tag{3.3.4}
m^*:H^{n-1}(X,\mathbb Q_\ell(i))
\longrightarrow H^{n-1}(X_{\bar\eta},\mathbb Q_\ell(i)),
\end{equation}
that is, such that
\begin{equation}\tag{3.3.5}
V_{\bar\eta}\,m^*=\mathrm{id}.
\end{equation}

\subsubsection*{3.4.}
Condition (B) is satisfied if the operation $\wedge$ of [4, 1.4.2.1]
is algebraic (cf.~[4, 2.12(ii)]).  It is conjectured that $\wedge$ is
always algebraic, and hence that condition (B) always holds; this is
known in many cases [4].  It holds trivially when $X$ is a complete
intersection of even dimension, since in that case we know (IX~1.6)
that
\begin{equation}\tag{3.4.1}
H^{n-1}(X,\mathbb Q_\ell(i))=0.
\end{equation}
\end{filecontents*}
% Original input: components/144_expose_XX_theorem_4_1_through_section_4_3.tex
\begin{filecontents*}[overwrite,noheader]{sga7ii-cumulative-input-144.tex}
\section*{4. Griffiths's theorem}

\subsubsection*{Theorem 4.1.}
Let $X$ be projective, smooth, and connected over $k$, of even dimension
$n=2m\geq4$.  Let $D$ be a Lefschetz pencil of sections by hypersurfaces
of degree $d$.  Assume that:

\medskip
\noindent\textbf{(4.1.0)} $X$ satisfies (LV) (XVIII~5.2.2);

\noindent\textbf{(4.1.1)} the pencil $D$ satisfies hypothesis (A)
(XVIII~5.3);

\noindent\textbf{(4.1.2)} the pencil $D$ satisfies alternative (3.2.2),
that is, in the setting of Proposition~3.2, the vanishing cohomology
$E^{n-1}(X_{\bar\eta},\mathbb Q_\ell)$ is not of level $\leq1$;

\noindent\textbf{(4.1.3)} the inclusion
$m:X_{\bar\eta}\hookrightarrow X\times_k k(\bar\eta)$ satisfies
hypothesis (B) (3.3).\footnote{The print writes $m^*$ before this geometric
inclusion.  Condition (B) in 3.3 is imposed on $m$; $m^*$ denotes its
induced map in cohomology.}

\medskip
Then, if an algebraic cycle $z$ of codimension $m$ on $X$ has a
restriction to $X_{\bar\eta}$ which, as an algebraic cycle, satisfies
\begin{equation}\tag{4.1.4}
z_{\bar\eta}=z|_{X_{\bar\eta}}
\quad\text{is algebraically equivalent to zero},
\end{equation}
then
\begin{equation}\tag{4.1.5}
\operatorname{cl}(z)=0
\quad\text{in }H^{2m}(X,\mathbb Q_\ell(m));
\end{equation}
in other words, $z$ is $\mathbb Q_\ell$-homologous to zero.

\subsubsection*{4.2. A remark on the case where $X$ is a complete intersection}
As already noted (cf.~XVIII~5.2.5 and (3.4)), hypotheses (4.1.0) and
(4.1.3) then hold automatically.  We saw in XVIII~6.3.4 and 6.4 that
hypothesis (4.1.1) holds for a pencil of hypersurfaces of degree $d>2n$.
Moreover, in the next exposé we shall prove condition (4.1.2) for a
``general'' pencil of hypersurfaces of degree $d>2n$ (XXI~5.2).
Consequently, a ``general'' pencil of hypersurfaces of degree $d>2n$
satisfies all the hypotheses of Theorem~4.1.

\subsubsection*{4.3.}
Returning to the case of arbitrary $X$, it seems very plausible that, for
$d$ sufficiently large, all the conditions of 4.1 hold for $D$, at least
when $D$ is sufficiently general.  We have already seen this for (4.1.1)
(XVIII~6.4), and we noted the conjecture that (4.1.0) and (4.1.3) always
hold (with (4.1.0) known in every case when $\operatorname{char}(k)=0$).
As for condition (4.1.2), Hodge theory proves it for large $d$ when
$\operatorname{char}(k)=0$,\footnote{The print has
$\operatorname{char}(k)\ne0$ in this Hodge-theoretic clause.  The following
clause treats positive characteristic, so the required reading here is
$\operatorname{char}(k)=0$.} while heuristic arguments, given in the
appendix to the next exposé, make the same conclusion plausible when
$\operatorname{char}(k)>0$.

More generally, these arguments---Hodge theory and, respectively,
crystalline cohomology---prove, and respectively make plausible, that for
a smooth projective variety $X$ of dimension $n$, and for every Lefschetz
pencil $D$ of hypersurfaces of sufficiently large degree $d$ (respectively,
every sufficiently general pencil $D$ of hypersurfaces of degree $d$), the
vanishing cohomology
$E^{n-1}(X_{\bar\eta},\mathbb Q_\ell)$ has level exactly $n-1$; that is, it
is not of level $\leq n-3$, or equivalently is not of coniveau $1$.
\end{filecontents*}
% Original input: components/145_expose_XX_proof_of_theorem_4_1.tex
\begin{filecontents*}[overwrite,noheader]{sga7ii-cumulative-input-145.tex}
\subsubsection*{4.4. Proof of Griffiths's theorem}

Let $z$ be an algebraic cycle of codimension $m$ on $X$ such that
\begin{equation}\tag{4.4.1}
z_{\bar\eta}=z|_{X_{\bar\eta}}
\quad\text{is algebraically equivalent to zero}.
\end{equation}
By definition, there exist
\begin{equation}\tag{4.4.2}
\text{a proper, smooth, connected curve $C_{\bar\eta}$ over $k(\bar\eta)$},
\end{equation}
\begin{equation}\tag{4.4.3}
t_i(\bar\eta):\Spec k(\bar\eta)\longrightarrow C_{\bar\eta},
\qquad i=1,2,
\end{equation}
and
\begin{equation}\tag{4.4.4}
\text{an algebraic cycle $P_{\bar\eta}$ of codimension $m$ on
$C_{\bar\eta}\times X_{\bar\eta}$},
\end{equation}
such that
\begin{equation}\tag{4.4.5}
z_{\bar\eta}
=P_{\bar\eta}\bigl(t_1(\bar\eta)\bigr)
-P_{\bar\eta}\bigl(t_2(\bar\eta)\bigr)
\quad\text{on }X_{\bar\eta}.
\end{equation}
The print writes target $C$ rather than $C_{\bar\eta}$ in (4.4.3), and
both above and in (4.4.4) it writes codimension $n$ rather than $m$;
the defined curve, Theorem~4.1, and the degree shift below force the
readings used here.

Regard $P_{\bar\eta}$ as an algebraic correspondence from
$C_{\bar\eta}$ to $X_{\bar\eta}$.  It defines a cohomological
correspondence, again denoted by $P_{\bar\eta}$, or equivalently a
homomorphism
\begin{equation}\tag{4.4.6}
P_{\bar\eta}:H^*(C_{\bar\eta},\mathbb Q_\ell(i))
\longrightarrow
H^{*+2(m-1)}(X_{\bar\eta},\mathbb Q_\ell(i+m-1)).
\end{equation}
The print omits the final closing parenthesis in the target of (4.4.6).

Likewise, consider the algebraic correspondence $V_{\bar\eta}$ of
(3.3.2), from $X_{\bar\eta}$ to $X\otimes_k k(\bar\eta)$, supplied by
hypothesis (B) in (4.1.3).  It defines
\begin{equation}\tag{4.4.7}
V_{\bar\eta}:H^*(X_{\bar\eta},\mathbb Q_\ell(i+m-1))
\longrightarrow
H^*(X\otimes_k k(\bar\eta),\mathbb Q_\ell(i+m-1)).
\end{equation}
The composite algebraic correspondence
$Q_{\bar\eta}=V_{\bar\eta}\circ P_{\bar\eta}$ from $C_{\bar\eta}$ to
$X\otimes_k k(\bar\eta)$ therefore defines
\begin{equation}\tag{4.4.8}
Q_{\bar\eta}=V_{\bar\eta}\circ P_{\bar\eta}:
H^*(C_{\bar\eta},\mathbb Q_\ell(i))
\longrightarrow
H^{*+2(m-1)}(X\otimes_k k(\bar\eta),\mathbb Q_\ell(i+m-1)).
\end{equation}

Consider the inclusion
\begin{equation*}
m:X_{\bar\eta}\longrightarrow X\otimes_k k(\bar\eta).
\end{equation*}
It gives the algebraic correspondence $m^*$ from
$X\otimes_k k(\bar\eta)$ to $X_{\bar\eta}$, hence the composite
correspondence $m^*Q_{\bar\eta}$ from $C_{\bar\eta}$ to
$X_{\bar\eta}$, defining
\begin{equation}\tag{4.4.9}
m^*Q_{\bar\eta}=m^*V_{\bar\eta}P_{\bar\eta}:
H^*(C_{\bar\eta},\mathbb Q_\ell(i))
\longrightarrow
H^{*+2(m-1)}(X_{\bar\eta},\mathbb Q_\ell(i+m-1)).
\end{equation}
The print repeats an extraneous $\bar\eta$ after $m^*Q_{\bar\eta}$ in
the preceding sentence.

Now consider the algebraic correspondence
\begin{equation}\tag{4.4.10}
\begin{split}
R_{\bar\eta}
&=P_{\bar\eta}-m^*Q_{\bar\eta}
=(\id-m^*V_{\bar\eta})P_{\bar\eta}:\\
&\hspace{2em}H^*(C_{\bar\eta},\mathbb Q_\ell(i))
\longrightarrow
H^{*+2(m-1)}(X_{\bar\eta},\mathbb Q_\ell(i+m-1)).
\end{split}
\end{equation}
Since $V_{\bar\eta}$ is a left inverse of $m^*$,
$\id-m^*V_{\bar\eta}$ is a projector on $H^*(X_{\bar\eta},-)$ whose
image is a complement of
\begin{equation*}
\operatorname{Im}\bigl(m^*:H^*(X\otimes_k k(\bar\eta),-)
\longrightarrow H^*(X_{\bar\eta},-)\bigr).
\end{equation*}
By hypothesis (4.1.2), the image of
\begin{equation}\tag{4.4.11}
R_{\bar\eta}:H^1(C_{\bar\eta},\mathbb Q_\ell(i))
\longrightarrow
H^{2m-1}(X_{\bar\eta},\mathbb Q_\ell(i+m-1))
\end{equation}
lies in $\operatorname{Im}m^*$.  Consequently,
\begin{equation}\tag{4.4.12}
\text{the homomorphism (4.4.11) is zero.}
\end{equation}

Choose a finite extension
\begin{equation*}
L/k(\eta)
\end{equation*}
over which the objects
\begin{equation*}
C_{\bar\eta},\quad t_1(\bar\eta),\quad t_2(\bar\eta),\quad
Q_{\bar\eta},\quad R_{\bar\eta}
\end{equation*}
are already defined, and let $U$ be a nonempty open subset of the
normalization of $\mathbb P^1$ in $L$.  After shrinking $U$, there is a
proper smooth relative curve
\begin{equation*}
C\longrightarrow U
\end{equation*}
with geometrically connected fibers and geometric generic fiber
$C_{\bar\eta}$, together with two sections
\begin{equation*}
t_i:U\longrightarrow C,\qquad i=1,2,
\end{equation*}
whose values at $\bar\eta$ are $t_i(\bar\eta)$.

Let
\begin{equation*}
\rho:\widetilde X\longrightarrow\mathbb P^1
\end{equation*}
be the projection of XVIII~3.1.4 whose fibers are the members of the
pencil, put
\begin{equation*}
\widetilde X_U=\widetilde X\times_{\mathbb P^1}U,
\end{equation*}
and let
\begin{equation*}
M:\widetilde X_U\hookrightarrow X\times U
\end{equation*}
be the canonical immersion into the constant family over $U$ with fiber
$X$.  Set
\begin{equation*}
z_U=M^*\pr_1^*(z),
\end{equation*}
the algebraic cycle on $\widetilde X_U$ induced by $z$.

After shrinking $U$ once more, there are an algebraic cycle $R$ on
$C\times_U\widetilde X_U$ whose value on
$C_{\bar\eta}\times X_{\bar\eta}$ is $R_{\bar\eta}$, and an algebraic
cycle $Q$ on $C\times X\times U$ whose value on
$C_{\bar\eta}\times(X\otimes_k k(\bar\eta))$ is $Q_{\bar\eta}$, such
that
\begin{equation}\tag{4.4.13}
z_U=R(t_1)-R(t_2)
+(\id\times M)^*\bigl[Q(t_1)-Q(t_2)\bigr].
\end{equation}

To prove the theorem it suffices, by the injectivity of the Griffiths
homomorphism ``modulo the fixed part'' (XVIII~5.8.7),
\begin{equation*}
\operatorname{Prim}^{2m}(X,\mathbb Q_\ell(m))
\xrightarrow{\ \underline{\operatorname{grif}}\ }
\frac{H^1(U,R^{2m-1}\rho_{U*}\mathbb Q_\ell(m))}
{M^*H^1(U,H^{2m-1}(X,\mathbb Q_\ell(m)))},
\end{equation*}
to show that the image of $z_U$ under
\begin{equation*}
u_0:\operatorname{Prim}^{2m}(\widetilde X_U/U,\mathbb Q_\ell(m))
\longrightarrow H^1(U,R^{2m-1}\rho_{U*}\mathbb Q_\ell(m))
\end{equation*}
lies in
$M^*H^1(U,H^{2m-1}(X,\mathbb Q_\ell(m)))$.
The print gives the twist $n$ in the domain of the Griffiths map; the
codimension and all adjacent terms require $m$.

In view of (4.4.13) and the commutativity of (1.4.2), we have the
commutative diagram
\begin{center}
\begin{tikzcd}[column sep=4.2em,row sep=4.2em]
\operatorname{Prim}^{2}(C/U,\mathbb Q_\ell(1))
  \arrow[r,"Q"]
& \operatorname{Prim}^{2m}(X\times U/U,\mathbb Q_\ell(m))
  \arrow[r,"M^*"] \arrow[d,"u_{\pr_2}"]
& \operatorname{Prim}^{2m}(\widetilde X_U/U,\mathbb Q_\ell(m))
  \arrow[d,"u_0"] \\
{}
& H^1(U,H^{2m-1}(X,\mathbb Q_\ell(m)))
  \arrow[r,"M^*"]
& H^1(U,R^{2m-1}\rho_{U*}\mathbb Q_\ell(m)),
\end{tikzcd}
\end{center}
so it is enough to prove
\begin{equation*}
u_0\bigl(R(t_1)-R(t_2)\bigr)=0
\quad\text{in}\quad
H^1(U,R^{2m-1}\rho_{U*}\mathbb Q_\ell(m)).
\end{equation*}
This follows from (4.4.12), by 1.7.  The print has $2n-1$ and twist
$n$ in the last group and cites the nonexistent (4.3.12); the coherent
readings are $2m-1$, twist $m$, and (4.4.12).
\end{filecontents*}
% Original input: components/146_expose_XX_section_5_and_references.tex
\begin{filecontents*}[overwrite,noheader]{sga7ii-cumulative-input-146.tex}
\section*{5. The Griffiths group}

\subsubsection*{5.0.}
Let $S$ be proper and smooth over an algebraically closed field $k$.
Denote by

\medskip
\noindent\textbf{(5.0.1)} $Z^i(S)$ the group of algebraic cycles of
codimension $i$ on $S$;

\noindent\textbf{(5.0.2)} $Z^i_{\mathrm{alg}}(S)$ the subgroup of
$Z^i(S)$ consisting of cycles algebraically equivalent to zero; and, for
every prime number $\ell\ne\operatorname{char}(k)$,

\noindent\textbf{(5.0.3)} $Z^i_{\ell\text{-}\mathrm{coh}}(S)$ the subgroup
of $Z^i(S)$ consisting of cycles whose cohomology class in
$H^{2i}(S,\mathbb Q_\ell(i))$ is zero.

It is conjectured in \textup{[4]} that this group is independent of $\ell$;
this is known at least when $\operatorname{char}(k)=0$.  Set
\begin{equation}\tag{5.0.4}
\mathrm{Grif}^i_\ell(S)
 =Z^i_{\ell\text{-}\mathrm{coh}}(S)/Z^i_{\mathrm{alg}}(S).
\end{equation}
If, moreover, $S$ is projective and is equipped with a projective embedding,
set

\medskip
\noindent\textbf{(5.0.5)} $Z^{i,\mathrm{prim}}(S)$ equal to the subgroup
of $Z^i(S)$ consisting of primitive cycles (that is, cycles whose
cohomology class is primitive), and

\noindent\textbf{(5.0.6)}
$\mathrm{Prim}^{2i}_{\mathrm{alg}}(S,\mathbb Q_\ell(i))$ equal to the
$\mathbb Q$-vector subspace of $H^{2i}(S,\mathbb Q_\ell(i))$ spanned by
the cohomology classes of primitive algebraic cycles.

\subsubsection*{Corollary 5.1.}
Under the hypotheses of Theorem~4.1, one has
\begin{equation}\tag{5.1.1}
\dim_{\mathbb Q}\!\left(
  \mathrm{Grif}^m_\ell(X_{\bar\eta})\otimes_{\mathbb Z}\mathbb Q
\right)
\geq
\dim_{\mathbb Q}\!\left(
  \mathrm{Prim}^{2m}_{\mathrm{alg}}(X,\mathbb Q_\ell(m))
\right).
\end{equation}
Indeed, by (4.1.5), the kernel of the composite map
\begin{equation*}
\mathbb Q\otimes_{\mathbb Z}Z^{m,\mathrm{prim}}(X)
 \xrightarrow{\ \mathrm{res.}\ \text{to }X_{\bar\eta}\ }
\mathbb Q\otimes_{\mathbb Z}Z^m_{\ell\text{-}\mathrm{coh}}
  (X_{\bar\eta})
 \longrightarrow \mathrm{Grif}^m_\ell(X_{\bar\eta})
\end{equation*}
is contained in
  $Z^m_{\ell\text{-}\mathrm{coh}}(X)\otimes_{\mathbb Z}\mathbb Q$.
Consequently,
  $\mathrm{Prim}^{2m}_{\mathrm{alg}}(X,\mathbb Q_\ell(m))$ is a quotient
of a subspace of
  $\mathrm{Grif}^m_\ell(X_{\bar\eta})\otimes_{\mathbb Z}\mathbb Q$,
which proves the assertion.

\subsubsection*{Remark 5.2.}
For $X$ a quadric hypersurface in $\mathbb P^{2m+1}$, one has
\begin{equation*}
\dim_{\mathbb Q_\ell}\!\left(
  \mathrm{Prim}^{2m}_{\mathrm{alg}}(X,\mathbb Q_\ell(m))
\right)=1.
\end{equation*}
The print has twist $n$ here and omits the outer closing parenthesis; the
middle degree $2m$ forces the twist $m$.

In particular, when $m=2$, for $X$ a quadric in $\mathbb P^5$, the section
of $X$ by a general hypersurface of degree $>5$ contains a curve which is
$\mathbb Q_\ell$-homologous to zero but of which no nonzero multiple is
algebraically equivalent to zero.\footnote{The print inserts the literal
parenthetical ``(4.2)'' at this point.  Direct inspection of the native scan
at 9000-dpi-equivalent enlargement confirms the glyphs; the numerical
specialization from $\mathbb P^{2m+1}$ to $\mathbb P^5$ is $m=2$.}

\begin{thebibliography}{9}
\bibitem{GriffithsPeriodsXX}
P. Griffiths, ``On the Periods of Certain Rational Integrals III,''
forthcoming in \emph{Annals of Mathematics}.

\bibitem{GrothendieckBrauerIIIXX}
A. Grothendieck, ``Le Groupe de Brauer III'' (especially \S\S 9--10), in
\emph{Dix Exposés sur la Cohomologie des Schémas}, North-Holland, 1968.

\bibitem{JouanolouCohomologieXX}
J.-P. Jouanolou, \emph{Cohomologie $\ell$-adique}, thesis, Paris, 1969.

\bibitem{KleimanCyclesXX}
S. Kleiman, ``Algebraic Cycles and the Weil Conjectures,'' in
\emph{Dix Exposés sur la Cohomologie des Schémas}, North-Holland, 1968.

\bibitem{ManinCurvesXX}
J. Manin, ``Rational Points of Algebraic Curves over Function Fields,''
\emph{Translations of the American Mathematical Society}.
\end{thebibliography}
\end{filecontents*}
% Original input: components/147_expose_XXI_title_introduction_and_definition_1_0.tex
\begin{filecontents*}[overwrite,noheader]{sga7ii-cumulative-input-147.tex}
\clearpage
\begin{center}
{\Large\bfseries EXPOSÉ XXI}\par
\vspace{1.5cm}
{\Large\bfseries THE LEVEL OF THE COHOMOLOGY OF COMPLETE
INTERSECTIONS}\par
\vspace{0.8cm}
{\large by N. Katz}\par
\vspace{0.5cm}
with an Appendix by P. Deligne (\S 5)
\end{center}

\vspace{1.5cm}
\section*{Contents}
\begin{description}
\item[0.] Introduction \hfill 1
\item[1.] The Hasse--Witt matrix of a complete intersection \hfill 2
\item[2.] Coniveau over a finite field of definition \hfill 14
\item[3.] Application to complete intersections \hfill 16
\item[4.] Conclusions \hfill 19
\item[5.] Integrality theorems \hfill 21
\end{description}

\clearpage
\section*{0. Introduction}

Let $X$ be a smooth projective variety of dimension $n\geq3$, let $Y(d)$
be a section by a hypersurface of degree $d$, and let
$E^{n-1}(Y(d),\mathbb Q_\ell)$ be the ``vanishing'' cohomology of $Y(d)$
(XVIII~5.4.4).  The aim of this exposé is to prove that, for $d$
sufficiently large and $Y(d)$ sufficiently general, it is not true that the
whole of $E^{n-1}(Y(d),\mathbb Q_\ell)$ ``comes'' from the groups $H^i$,
$i<n-1$, of smooth proper algebraic varieties by means of algebraic
correspondences.  In characteristic zero this follows trivially from the
Hodge bigrading on complex cohomology and from the fact that, in coherent
cohomology,
\begin{equation*}
\dim H^{n-1}(Y(d),\mathcal O_{Y(d)})\longrightarrow\infty
\qquad\text{as }d\longrightarrow\infty.
\end{equation*}
The print has the evident typographical forms \emph{et de démontrer} and
\emph{suffisament}; the readings used above are \emph{est de démontrer} and
\emph{suffisamment}.

In characteristic $p$, the method again relates the coherent group
$H^{n-1}(Y(d),\mathcal O_{Y(d)})$ to the étale group
$H^{n-1}_{\mathrm{et}}(Y(d),\mathbb Q_\ell)$.  To do so, one works over a
finite base field.  The two are then linked by the zeta function of $Y(d)$:
this function is expressed in terms of $\ell$-adic cohomology, while its
``reduction modulo $p$'' is expressed in terms of coherent cohomology
(cf.~XXII~3.9).  For complete intersections this method works without
difficulty, owing to the detailed knowledge of both kinds of cohomology
(étale and coherent) in that case.  It yields the main result, Theorem~4.2,
as a consequence of a result on Hasse--Witt operations, Theorem~1.4 (see
1.0).

There are heuristic arguments due to A. Grothendieck, based on the yoga of
crystalline cohomology, which make the general statement plausible for every
smooth projective $X$, by essentially the same method.

Even in the special case of complete intersections, the proof makes
essential use in 2.2 of Theorem~2.1 on the integrality of the eigenvalues of
Frobenius, recently proved by P. Deligne and proved here in the Appendix
(\S 5).  The print calls this ``Appendix I,'' while the title and the document
structure identify Deligne's Appendix as \S 5.

\section*{1. The Hasse--Witt matrix of a complete intersection}

\subsubsection*{Definition 1.0.}
A proper scheme over a field $k$ of characteristic exponent $p$ is called
\emph{special} if the additive $p$-linear operation
\begin{equation}\tag{1.0.1}
\mathfrak F_p:\mathcal O_X\longrightarrow\mathcal O_X,
\qquad g\longmapsto g^p,
\end{equation}
induces a nilpotent operation
\begin{equation}\tag{1.0.2}
\mathfrak F_p:H^{\dim X}(X,\mathcal O_X)
\longrightarrow H^{\dim X}(X,\mathcal O_X).
\end{equation}
Equivalently, if
$M=\dim_k H^{\dim X}(X,\mathcal O_X)$, the $p^M$-linear operation
\begin{equation}\tag{1.0.3}
\mathfrak F_p^M=\mathfrak F_{p^M}:\mathcal O_X\longrightarrow\mathcal O_X,
\qquad g\longmapsto g^{p^M},
\end{equation}
annihilates $H^{\dim X}(X,\mathcal O_X)$.

Recall from \textup{[2]} that, for every $i\in\mathbb Z$, the homomorphism
\begin{equation*}
H^i(X,\mathcal O_X)\xrightarrow{\ \mathfrak F_p\ }H^i(X,\mathcal O_X)
\end{equation*}
induced by (1.0.1) is classically called the \emph{Hasse--Witt
homomorphism} on $H^i(X,\mathcal O_X)$.  Notice that, when $p=1$, saying
that $X$ is non-special means simply that
$H^{\dim X}(X,\mathcal O_X)\ne0$.
\end{filecontents*}
% Original input: components/148_expose_XXI_theorem_1_1_and_proposition_1_1_5.tex
\begin{filecontents*}[overwrite,noheader]{sga7ii-cumulative-input-148.tex}
\subsubsection*{Theorem 1.1.}
Fix a characteristic exponent $p$, a dimension $n\geq1$, an integer
$r\geq1$, and a multidegree $(d_1,\dots,d_r)$.  If
\begin{equation}\tag{1.1.0}
\sum_{i=1}^{r}d_i\geq(n+r)+1,
\end{equation}
then, among all complete intersections of dimension $n$ in
$\mathbb P^{n+r}$ having multidegree $(d_1,\dots,d_r)$ in characteristic
$p$, a ``sufficiently general'' member is non-special in the sense of 1.0.

\medskip
\noindent\emph{Proof.}
Let
\begin{equation}\tag{1.1.1}
f:X_U\longrightarrow U
\end{equation}
denote the ``universal family'' of complete intersections in
$\mathbb P^{n+r}$ of multidegree $(d_1,\dots,d_r)$ in characteristic $p$,
together with their defining equations.  Thus $U$ is an open subset of a
product of $r$ projective spaces over $\mathbb F_p$, of impressive
dimension.

The $\mathcal O_U$-module
\begin{equation}\tag{1.1.2}
R^nf_*\mathcal O_{X_U}
\end{equation}
is locally free, and its formation commutes with every base change
(XI~1.5(i)).  By (1.1.0), its rank $M$ is nonzero (XI~2.5).  When $p=1$,
this proves the theorem for every complete intersection satisfying (1.1.0).

Suppose now that $p\geq2$.  The iterated Hasse--Witt operation
\begin{equation}\tag{1.1.3}
\mathfrak F_p^M=\mathfrak F_{p^M}:
R^nf_*\mathcal O_{X_U}\longrightarrow R^nf_*\mathcal O_{X_U}
\end{equation}
is $p^M$-linear over $\mathcal O_U$.  Since
$R^nf_*\mathcal O_{X_U}$ is locally free, (1.1.3) is represented locally on
$U$ by an $(M,M)$-matrix $H$ with entries in $\mathcal O_U$.  Consequently,
\begin{equation}\tag{1.1.4}
\{u\in U\mid X_u\text{ is non-special}\}\quad\text{is open in }U.
\end{equation}
Indeed, this is the set of $u$ for which $H(u)\ne0$.  Since $U$ is
irreducible, it is therefore enough to find one value of $u$ for which
$X_u$ is non-special; we shall do so over a finite field.

\subsubsection*{Proposition 1.1.5.}
Let $X$ be a complete intersection over $\mathbb F_q$.  If $X$ has no
$\mathbb F_q$-rational points, then $X$ is non-special.

\medskip
\noindent\emph{Proof.}
It is plainly enough to prove, for a complete intersection $X$, the
implication
\begin{equation}\tag{1.1.5.1}
X\text{ special}\quad\Longrightarrow\quad
\operatorname{card}\bigl(X(\mathbb F_q)\bigr)\equiv1\pmod p.
\end{equation}
By the trace formula, proved in the following Exposé XXII, 3.2,
\begin{equation*}
\operatorname{card}\bigl(X(\mathbb F_q)\bigr)
\equiv\sum_{i\geq0}(-1)^i
\operatorname{tr}\bigl(\mathfrak F_q\mid H^i(X,\mathcal O_X)\bigr)
\pmod p.
\end{equation*}
Now $H^i(X,\mathcal O_X)=0$ unless $i=0$ or $i=n=\dim X$, while
$H^0(X,\mathcal O_X)\simeq\mathbb F_q$ (XI~1.5(iii)); hence this becomes
\begin{equation*}
\operatorname{card}\bigl(X(\mathbb F_q)\bigr)
\equiv1+(-1)^n
\operatorname{tr}\bigl(\mathfrak F_q\mid H^n(X,\mathcal O_X)\bigr)
\pmod p.
\end{equation*}
If $X$ is special, the operation of $\mathfrak F_q$ on
$H^n(X,\mathcal O_X)$ is nilpotent by definition and therefore has trace
zero.  This proves (1.1.5.1), and hence the proposition.

\medskip
\noindent\emph{Source note.}
The diplomatic French retains the printed forms \emph{parmi toute les},
\emph{théroème}, \emph{intersection complète vérifiante}, and
\emph{n'a pas des points}; the editorial French and the translation use the
grammatically coherent readings recorded in the source-defect register.
\end{filecontents*}
% Original input: components/149_expose_XXI_proposition_1_2.tex
\begin{filecontents*}[overwrite,noheader]{sga7ii-cumulative-input-149.tex}
\subsubsection*{Proposition 1.2.}
For every power $q$ of $p$, if
\begin{equation*}
\sum_{i=1}^{r}d_i\geq n+r+1,
\end{equation*}
there exists a complete intersection in $\mathbb P^{n+r}_{\mathbb F_q}$ of
multidegree $(d_1,\dots,d_r)$ having no $\mathbb F_q$-rational point.

\medskip
\noindent\emph{Proof.}
Begin with the case $r=1$.  For every $d\geq n+2$, we must find a
homogeneous form of degree $d$,
\begin{equation}\tag{1.2.1}
H(X_1,\dots,X_{n+2})\in
\mathbb F_q[X_1,\dots,X_{n+2}],
\end{equation}
such that
\begin{equation}\tag{1.2.2}
\begin{aligned}
&(\alpha_1,\dots,\alpha_{n+2})\in(\mathbb F_q)^{n+2},\qquad
H(\alpha_1,\dots,\alpha_{n+2})=0\quad\Longrightarrow\\
&\hspace{7em}\alpha_1=\cdots=\alpha_{n+2}=0.
\end{aligned}
\end{equation}
Set $\mathbb E=\mathbb F_{q^d}$, and choose elements
$\beta_1,\dots,\beta_{n+2}$ of $\mathbb E$ that are linearly independent
over $\mathbb F_q$.  Take
\begin{equation}\tag{1.2.3}
H(X_1,\dots,X_{n+2})=
N_{\mathbb E/\mathbb F_q}
\left(\sum_{i=1}^{n+2}\beta_iX_i\right).
\end{equation}

In the general case, since
\begin{equation}\tag{1.2.4}
\sum_{i=1}^{r}d_i\geq n+r+1,
\end{equation}
we can set $a_0=0$ and choose $r$ integers
\begin{equation}\tag{1.2.5}
1\leq a_1<a_2<\cdots<a_r=n+r+1
\end{equation}
such that
\begin{equation}\tag{1.2.6}
d_i\geq a_i-a_{i-1}\qquad(1\leq i\leq r).
\end{equation}
For each $i=1,\dots,r$, choose a homogeneous form of degree $d_i$,
\begin{equation*}
H_i(X_{a_{i-1}+1},\dots,X_{a_i})
\in\mathbb F_q[X_{a_{i-1}+1},\dots,X_{a_i}],
\end{equation*}
in the $a_i-a_{i-1}$ displayed variables and having property (1.2.2).
The required complete intersection is defined by the homogeneous forms
\begin{equation*}
H_1(X_1,\dots,X_{a_1}),\quad\dots,\quad
H_r(X_{a_{r-1}+1},\dots,X_{a_r}).
\end{equation*}

\medskip
\noindent\emph{Source note.}
The diplomatic French retains the printed misplaced parenthesis in (1.2.1),
the literal bound $d_1\geq a_o$ in (1.2.6), and the missing $+1$ in the
terminal $H_r$ block.  Direct native-pixel inspection confirms these as print
features; the construction forces the coherent forms displayed above.
\end{filecontents*}
% Original input: components/150_expose_XXI_theorem_1_3_finite_field_case.tex
\begin{filecontents*}[overwrite,noheader]{sga7ii-cumulative-input-150.tex}
\subsubsection*{Theorem 1.3.}
Let $X$ be an irreducible complete intersection of dimension $n\geq1$ over
a field $k$ of characteristic exponent $p$.  For every integer $d\geq n+1$,
a sufficiently general section of $X$ by a hypersurface of degree $d$ is
non-special in the sense of 1.0.

\medskip
\noindent\emph{Proof.}
The sections $X.H_u$ of $X$ by hypersurfaces of degree $d$ are parametrized
by the lines in $H^0(X,\mathcal O_X(d))$, that is, by
\begin{equation}\tag{1.3.1}
U=\mathbb P\bigl(H^0(X,\mathcal O_X(d))\bigr).
\end{equation}
Over $U$ we have the ``universal section of $X$ by a hypersurface of degree
$d$,''
\begin{equation}\tag{1.3.2}
f:X.H_u\longrightarrow U.
\end{equation}
Let $U'$ be the open subset of those $u\in U$ for which
$\dim(X_u)=n-1$.  By XI~1.5(i) and 2.5, the restriction
\begin{equation}\tag{1.3.3}
\left.R^{n-1}f_*\mathcal O_{X.H_u}\right|_{U'}
\end{equation}
is locally free, its formation commutes with every base change, and its rank
$M$ is at least one.  If $p=1$, this proves the theorem for every section by
a hypersurface of degree at least $n+1$.

Suppose therefore that $p>0$.  The iterated Hasse--Witt operation
$\mathfrak F_{p^M}$ is again represented locally by an $M\times M$ matrix
with entries in $\mathcal O_U$.  Consequently,
\begin{equation}\tag{1.3.4}
\{u\in U'\mid X.H_u\text{ is non-special}\}
\quad\text{is open in }U'.
\end{equation}
It remains only to prove that the set in (1.3.4) is nonempty.

First consider the case
\begin{equation}\tag{1.3.5}
k=\mathbb F_q.
\end{equation}
After replacing $\mathbb F_q$ by a finite extension, we may suppose that
\begin{equation}\tag{1.3.6}
\operatorname{card}\bigl(X(\mathbb F_q)\bigr)>0
\end{equation}
and that
\begin{equation}\tag{1.3.7}
\begin{aligned}
&\text{there are homogeneous coordinates}\\
&\qquad X_1,\dots,X_{n+r+1}\text{ in the ambient space }\mathbb P_k^{n+r}\\
&\text{such that }X\text{ does not meet the linear subvariety}\\
&\qquad X_1=\cdots=X_{n+1}=0.
\end{aligned}
\end{equation}
Choose a homogeneous form of degree $d$,
\begin{equation}\tag{1.3.8}
H(X_1,\dots,X_{n+1})\in
\mathbb F_q[X_1,\dots,X_{n+1}],
\end{equation}
having property (1.2.2), which is possible because $d\geq n+1$.  Let $H$
also denote the hypersurface in $\mathbb P^{n+r}$ defined by
$H(X_1,\dots,X_{n+1})=0$.  By construction, $X.H$ has no
$\mathbb F_q$-rational point.  In view of (1.3.6), this also shows that
$X\not\subset H$; since $X$ is irreducible, $X.H$ has dimension $n-1$, so
$H\in U'$.  Proposition~1.1.5 now shows that $X.H$ is non-special.

\medskip
\noindent\emph{Source note.}
The diplomatic French retains the printed exponent $p^n$, the base sheaf
$\mathcal O_u$, and the shortened citation ``(1.5).''  The rank-$M$
nilpotence test, the base $U$, and Proposition~1.1.5 force the coherent
readings used above.  The closing parenthesis in (1.3.1), absent only from
the inherited TeX, has been restored directly from the scan.
\end{filecontents*}
% Original input: components/151_expose_XXI_theorem_1_3_general_case.tex
\begin{filecontents*}[overwrite,noheader]{sga7ii-cumulative-input-151.tex}
\par\medskip
We now pass to the general case.  Looking at the equations defining $X$, we
may suppose that $k$ is finitely generated over $\mathbb F_p$.  There is
then a subring $B\subset k$, finitely generated as an algebra over
$\mathbb F_p$ and having fraction field $k$, such that
\begin{equation}\tag{1.3.9}
X\subset\mathbb P_k^{n+r}
\end{equation}
is the generic fiber of a relative complete intersection
$\widetilde X\subset\mathbb P_B^{n+r}$ which, fiber by fiber, has the same
multidegree as $X$ and is geometrically irreducible.

Write
\begin{equation}\tag{1.3.10}
\rho:\widetilde X\longrightarrow S=\operatorname{Spec}(B)
\end{equation}
for the projection,
\begin{equation}\tag{1.3.11}
E=\rho_*\mathcal O_{\widetilde X}(d)
\end{equation}
for the locally free $\mathcal O_S$-module, and
\begin{equation}\tag{1.3.12}
\pi:T=\check{\mathbb P}(E)\longrightarrow S
\end{equation}
for the associated bundle of lines in $E$.  The generic fiber of (1.3.12)
is precisely
\[
U=\mathbb P\bigl(H^0(X,\mathcal O_X(d))\bigr).
\]
Over $T$ we have the universal section of
$\widetilde X\subset\mathbb P_B^{n+r}$ by hypersurfaces of degree $d$,
denoted
\begin{equation}\tag{1.3.13}
\alpha:\widetilde X.H_T\longrightarrow T.
\end{equation}
These constructions fit into the following diagram of cartesian squares:
\begin{equation}\tag{1.3.14}
\begin{tikzcd}[column sep=huge,row sep=large]
\widetilde X.H_T \arrow[d,"\alpha"]
  & X.H_u=(\widetilde X.H_T)\times_T U \arrow[l]\arrow[d] \\
T \arrow[d,"\pi"]
  & U=T\times_S\operatorname{Sp}(k) \arrow[l]\arrow[d] \\
S & \operatorname{Sp}(k)\arrow[l]
\end{tikzcd}
\end{equation}
The $\mathcal O_T$-module
\begin{equation}\tag{1.3.15}
R^{n-1}\alpha_*\mathcal O_{\widetilde X.H_T}
\end{equation}
is locally free over the open subset $T'\subset T$ consisting of the
$t$ for which $\dim\alpha^{-1}(t)=n-1$, and its formation there commutes
with every base change.  Since $T$ and $U$ have the same generic point, it
is enough, in order to show that (1.3.4) is nonempty, to prove that the open
subset of $T'$
\begin{equation}\tag{1.3.16}
\bigl\{t\in T'\mid\widetilde X_{\pi(t)}.H_t
   \text{ is non-special}\bigr\}
\end{equation}
is nonempty.  But the already established case (1.3.5) shows that, for
every closed point $s\in S$, the fiber $\pi^{-1}(s)$ has nonempty
intersection with the open set (1.3.16).  This proves the theorem.

\medskip
\noindent\emph{Source note.}
The diplomatic French retains the printed top-right object
$(\widetilde X.H_T)\times_U T$ and the phrase \emph{se formation}.
The cartesian square and the definitions of $T$ and its generic fiber $U$
force $(\widetilde X.H_T)\times_T U$ and \emph{sa formation}, as used above.
The plain $E$, the exponent $0$ in $H^0$, and the outer parenthesis defining
$U$ are transcription repairs made directly from the scan rather than
editorial emendations.
\end{filecontents*}
% Original input: components/152_expose_XXI_theorem_1_4.tex
\begin{filecontents*}[overwrite,noheader]{sga7ii-cumulative-input-152.tex}
\subsubsection*{Theorem 1.4.}
Let $X\subset\mathbb P^N$ be a projective Cohen--Macaulay scheme,
equidimensional of dimension $n\geq1$, over a field $k$.  There is an
integer $d_0$ such that, for every $d\geq d_0$, if
$i:X.H\longrightarrow X$ denotes the inclusion of a section by a general
hypersurface of degree $d$, then the action of
\begin{equation}\tag{1.4.0}
\mathfrak F_p\quad\text{on}\quad
\frac{H^{n-1}(X.H,\mathcal O_{X.H})}
     {i^*H^{n-1}(X,\mathcal O_X)}
\end{equation}
is non-nilpotent.

\medskip
\noindent\emph{Proof.}
As above, we reduce to the case $k=\mathbb F_q$ and to the problem of
finding just one $H$---for any fixed sufficiently large degree for which
$\dim H^{n-1}(X.H,\mathcal O_{X.H})$ is independent of $H$; see
(1.4.4)---for which (1.4.0) holds.  After enlarging $q$, we may suppose
that the irreducible components $X_i$ of $X$ are geometrically irreducible,
that
\begin{equation}\tag{1.4.1}
\operatorname{card}\bigl(X_i(\mathbb F_q)\bigr)\geq2,
\end{equation}
and that
\begin{equation}\tag{1.4.2}
\begin{aligned}
&\text{there are projective coordinates }X_0,\dots,X_N\text{ such that}\\
&\qquad X\text{ does not meet the linear subvariety }
X_0=\cdots=X_n=0.
\end{aligned}
\end{equation}
By a vanishing theorem of Serre ([1], or SGA~2 XII~1.4), there is an
integer $d_1$ such that, for every $d\geq d_1$,
\begin{equation}\tag{1.4.3}
H^i(X,\mathcal O_X(-d))=0\qquad\text{if }i\ne n=\dim X.
\end{equation}
Thus, if $H$ is a hypersurface of degree $d\geq d_1$, then (1.4.3) gives
\begin{equation}\tag{1.4.4}
\begin{aligned}
H^i(X,\mathcal O_X)&\xrightarrow{\ \sim\ }
 H^i(X.H,\mathcal O_{X.H}) &&\text{if }i\leq n-2,\\[1ex]
H^{n-1}(X,\mathcal O_X)&\lhook\joinrel\longrightarrow
 H^{n-1}(X.H,\mathcal O_{X.H}) &&\text{injectively}.
\end{aligned}
\end{equation}

We prove the theorem with $d_0=\max(N+1,d_1)$.  For every $d\geq d_0$,
we must find a hypersurface $H$ of degree $d$, defined over $\mathbb F_q$,
which contains no $X_i$ and for which
\begin{equation}\tag{1.4.5}
\operatorname{tr}\!\left(\mathfrak F_q\ \middle|\
\frac{H^{n-1}(X.H,\mathcal O_{X.H})}
     {H^{n-1}(X,\mathcal O_X)}\right)\ne0.
\end{equation}
By the trace formula (XX~3.2.1), modulo $p$ we have
\begin{equation}\tag{1.4.6}
\begin{aligned}
\operatorname{card}\bigl((X.H)(\mathbb F_q)\bigr)
&\equiv\sum_{i=0}^{n-1}(-1)^i
 \operatorname{tr}\bigl(\mathfrak F_q\mid
 H^i(X.H,\mathcal O_{X.H})\bigr)\\
&\stackrel{(1.4.4)}{\equiv}
 \sum_{i=0}^{n-1}(-1)^i
 \operatorname{tr}\bigl(\mathfrak F_q\mid H^i(X,\mathcal O_X)\bigr)\\
&\qquad+(-1)^{n-1}\operatorname{tr}\!\left(\mathfrak F_q\ \middle|\
 \frac{H^{n-1}(X.H,\mathcal O_{X.H})}
      {H^{n-1}(X,\mathcal O_X)}\right).
\end{aligned}
\end{equation}
Let $a\in\mathbb Z/p\mathbb Z$ denote the quantity
\[
a=\sum_{i=0}^{n-1}(-1)^i
\operatorname{tr}\bigl(\mathfrak F_q\mid H^i(X,\mathcal O_X)\bigr).
\]
It is therefore enough, for every $d\geq d_0$, to find a hypersurface $H$
of degree $d$, defined over $\mathbb F_q$, which does not contain $X$ and
such that
\begin{equation}\tag{1.4.7}
\operatorname{card}\bigl((X.H)(\mathbb F_q)\bigr)\not\equiv a\pmod p.
\end{equation}

There are two cases.  If $a\ne0$, it is enough to take an $H$ over
$\mathbb F_q$ for which $(X.H)(\mathbb F_q)$ is empty, just as in (1.3.8).
Such an $H$ satisfies (1.4.7) and contains no $X_i$, because
$\operatorname{card}(X_i(\mathbb F_q))\geq2$ by hypothesis.

If $a=0$, it is enough to find an $H$ such that
$\operatorname{card}((X.H)(\mathbb F_q))=1$.  Then (1.4.7) holds and
$X_i\not\subset H$, since $\operatorname{card}(X_i(\mathbb F_q))\geq2$.
Choose projective coordinates $X_0,\dots,X_N$ such that $(1,0,\dots,0)$ is
a point of $X$.  Take $H$ to have equation $H(X_1,\dots,X_N)=0$, where $H$
is a homogeneous form satisfying (cf. (1.2.2))
\begin{equation}\tag{1.4.8}
\alpha_1,\dots,\alpha_N\in\mathbb F_q,\quad
H(\alpha_1,\dots,\alpha_N)=0
\quad\Longrightarrow\quad
\alpha_1=\cdots=\alpha_N=0.
\end{equation}
The hypersurface defined by $H$ in $\mathbb P^N$ then has only the point
$(1,0,\dots,0)$ over $\mathbb F_q$; a fortiori, $X.H$ has only this point
over $\mathbb F_q$.  This proves the theorem.

\medskip
\noindent\emph{Source note.}
The diplomatic French restores the scan's \emph{coordonnées}, the two
printed equality signs in (1.4.6), and the three-dot ellipsis in (1.4.8),
all of which differ from the inherited transcription.  The print uses the
letter-like subscript $o$ for the zero indices in $d_o$ and $X_o$, says
\emph{une $H$ défini}, and uses irregular point-count parentheses.  The
corrected French and English regularize these as $d_0$, $X_0$, masculine
agreement, and ordinary evaluation notation without changing the argument.
\end{filecontents*}
% Original input: components/153_expose_XXI_section_2.tex
\begin{filecontents*}[overwrite,noheader]{sga7ii-cumulative-input-153.tex}
\section*{2. Coniveau over a finite field of definition}

\subsection*{2.0.}
Let $X$ be a separated scheme of finite type over $\mathbb F_q$.  For every
prime number $\ell\ne\operatorname{char}(\mathbb F_q)$, write
\begin{equation}\tag{2.0.1}
H_c^i(\overline X,\mathbb Q_\ell)
\end{equation}
for the compactly supported cohomology groups of
$\overline X=X\otimes_{\mathbb F_q}\overline{\mathbb F}_q$.  These are
finite-dimensional $\mathbb Q_\ell$-vector spaces on which the Galois group
$\operatorname{Gal}(\overline{\mathbb F}_q/\mathbb F_q)$ acts.  Let
\begin{equation}\tag{2.0.2}
\varphi_q\in\operatorname{Gal}(\overline{\mathbb F}_q/\mathbb F_q)
\end{equation}
denote the Frobenius automorphism.  It is known (SGA~5 XIII~3.2.1) that the
zeta function of $X$ (XXII~3.0.1) is given by
\begin{equation}\tag{2.0.3}
Z(t,X/\mathbb F_q)=
\prod_{i\geq0}
\det\!\left(1-t\varphi_q^{-1}\mid
H_c^i(\overline X,\mathbb Q_\ell)\right)^{(-1)^{i+1}}.
\end{equation}
Deligne has recently proved, generalizing a result of Washnitzer for smooth
projective $X$, the following theorem, which is proved in the appendix below
(5.5.3).

\subsubsection*{Theorem 2.1.}
The eigenvalues of $\varphi_q^{-1}$ on
$H_c^i(\overline X,\mathbb Q_\ell)$ are algebraic integers; if
$i>\dim X$, they are divisible, as algebraic integers, by
$q^{,i-\dim X}$.

In particular, tautologically, we obtain:

\medskip
\noindent\textbf{(2.1 bis).}
The eigenvalues of $\varphi_q^{-1}$ on
$H_c^i(\overline X,\mathbb Q_\ell(-j))$ are algebraic integers divisible by
$q^j$.

\subsubsection*{Corollary 2.2.}
Let $X$ be proper, smooth, and geometrically connected over $\mathbb F_q$.
Then the eigenvalues of $\varphi_q^{-1}$ on
$N^jH^q(\overline X,\mathbb Q_\ell)$ (notation of XX~2.0.6) are algebraic
integers divisible by $q^j$.

\medskip
\noindent\emph{Proof.}
The assertion is invariant under a finite extension
$\mathbb F_{q^M}/\mathbb F_q$, which amounts to replacing the eigenvalues
under consideration by their $M$-th powers.  After making such a base
change, the definition gives finitely many projective, smooth, geometrically
connected schemes $T_i$ over $\mathbb F_q$, of dimensions $d_i$, and
algebraic cycles $Z_i$ on $T_i\times_{\mathbb F_q}X$ of codimension
$j+d_i$, such that the homomorphism of
$\operatorname{Gal}(\overline{\mathbb F}_q/\mathbb F_q)$-modules
\begin{equation*}
\bigoplus_i H^{q-2j}(\overline T_i,\mathbb Q_\ell(-j))
\xrightarrow{\ \oplus_i\operatorname{Cl}(Z_i)\ }
N^jH^q(\overline X,\mathbb Q_\ell)
\end{equation*}
is surjective.  The conclusion follows from (2.1 bis).

\medskip
\noindent\emph{Source note.}
The diplomatic French preserves the printed terms
\emph{$\mathbb Q_\ell$-vectoriels}, \emph{fonction zêta ... donné}, the
overstruck phrase following the extension $\mathbb F_{q^M}/\mathbb F_q$,
the Galois group with a blank numerator, and the displayed domain without
an explicit direct-sum sign.  The corrected French and English use standard
vector-space terminology and agreement, state the $M$-th-power operation
plainly, restore $\overline{\mathbb F}_q$ in the Galois group, and display
the direct sum required by $\oplus_i\operatorname{Cl}(Z_i)$.
\end{filecontents*}
% Original input: components/154_expose_XXI_section_3_opening_through_corollary_3_1.tex
\begin{filecontents*}[overwrite,noheader]{sga7ii-cumulative-input-154.tex}
\section*{3. Application to complete intersections}

\subsection*{3.0.}
Let $X$ be a smooth complete intersection of dimension $n\geq1$ in
$\mathbb P^N$ over a field $k$.  For
$\ell\ne\operatorname{char}(k)$, weak Lefschetz and duality give
(XI~1.6)
\begin{equation}\tag{3.0.1}
H^i(\overline X,\mathbb Q_\ell)=
\begin{cases}
0 & \text{if $i$ is odd and $i\ne n$},\\[1.5ex]
\mathbb Q_\ell\!\left(-\dfrac{i}{2}\right)
  & \text{if $i$ is even, $i\leq2n$, and $i\ne n$},\\[1.5ex]
H^n(\overline X,\mathbb Q_\ell) & \text{if $i=n$}.
\end{cases}
\end{equation}
For $k=\mathbb F_q$ this gives the following formula for the zeta function:
\begin{equation}\tag{3.0.2}
Z(t,X/\mathbb F_q)\prod_{i=0}^{n}(1-q^it)
=\det\!\left(1-t\varphi_q^{-1}\mid
\operatorname{Prim}^n(\overline X,\mathbb Q_\ell)\right)^{(-1)^{n+1}}.
\end{equation}
Here we have used the primitive decomposition and (3.0.1).

\subsubsection*{Corollary 3.0.3.}
Under the hypotheses of 3.0, with $k=\mathbb F_q$, there is a congruence
modulo $p$
\[
\det\!\left(1-t\varphi_q^{-1}\mid
H^n(\overline X,\mathbb Q_\ell)\right)
\equiv
\det\!\left(1-t\mathfrak F_q\mid H^n(X,\mathcal O_X)\right).
\]

\medskip
\noindent\emph{Proof.}
By the congruence formula (XXII~3.2.1), and since
\begin{equation}\tag{3.0.4}
H^i(X,\mathcal O_X)\simeq
\begin{cases}
\mathbb F_q & i=0,\\[1.5ex]
0 & 1\leq i\leq n-1,\\[1.5ex]
0 & i\geq n+1,
\end{cases}
\end{equation}
we have
\begin{equation}\tag{3.0.5}
Z(t,X/\mathbb F_q)\equiv
(1-t)^{-1}
\det\!\left(1-t\mathfrak F_q\mid H^n(X,\mathcal O_X)\right)^{(-1)^{n+1}}
\pmod p.
\end{equation}
Consequently, (3.0.5) and (3.0.2) give
\begin{equation}\tag{3.0.6}
\det\!\left(1-t\varphi_q^{-1}\mid
\operatorname{Prim}^n(\overline X,\mathbb Q_\ell)\right)
\equiv
\det\!\left(1-t\mathfrak F_q\mid H^n(X,\mathcal O_X)\right)
\pmod p.
\end{equation}
The desired conclusion follows because
\begin{equation}\tag{3.0.7}
H^n(\overline X,\mathbb Q_\ell)\simeq
\begin{cases}
\operatorname{Prim}^n(\overline X,\mathbb Q_\ell)
  & \text{if $n$ is odd},\\[2mm]
\operatorname{Prim}^n(\overline X,\mathbb Q_\ell)
 \oplus\mathbb Q_\ell\!\left(-\dfrac n2\right)
  & \text{if $n$ is even}.
\end{cases}
\end{equation}

\subsubsection*{Corollary 3.1.}
Let $X$ be a smooth complete intersection of dimension $n\geq1$ over
$\mathbb F_q$.  If
\begin{equation}\tag{3.1.0}
H^n(\overline X,\mathbb Q_\ell)=N^1H^n(\overline X,\mathbb Q_\ell),
\end{equation}
then
\begin{equation}\tag{3.1.1}
\mathfrak F_p\quad\text{is nilpotent on}\quad H^n(X,\mathcal O_X).
\end{equation}

\medskip
\noindent\emph{Proof.}
By Corollary~2.2, (3.1.0) implies
\begin{equation}\tag{3.1.2}
\det\!\left(1-t\varphi_q^{-1}\mid
H^n(\overline X,\mathbb Q_\ell)\right)
\equiv1\pmod p
\qquad\text{(indeed, even modulo $q$)}.
\end{equation}
Together with Corollary~3.0.3, this gives
\begin{equation}\tag{3.1.3}
\det\!\left(1-t\mathfrak F_q\mid H^n(X,\mathcal O_X)\right)=1.
\end{equation}
Thus $\mathfrak F_q$ is nilpotent on $H^n(X,\mathcal O_X)$, and hence so is
$\mathfrak F_p$.  This proves the corollary.

\medskip
\noindent\emph{Source note.}
The diplomatic French preserves the printed word order
\emph{intersection complète dans $\mathbb P^N$ lisse}, the typo
\emph{da décomposition primitive}, and the omitted index in
$1\leq{}\leq n-1$ in (3.0.4).  The corrected French and English use normal
word order, \emph{la décomposition primitive}, and restore the forced index
$i$ in the cohomology range.
\end{filecontents*}
% Original input: components/155_expose_XXI_proposition_3_2.tex
\begin{filecontents*}[overwrite,noheader]{sga7ii-cumulative-input-155.tex}
\subsubsection*{Proposition 3.2.}
Let $X$ be a smooth complete intersection of dimension $n\geq1$ over a
field $k$.  If
\begin{equation}\tag{3.2.1}
H^n(\overline X,\mathbb Q_\ell)=N^1H^n(\overline X,\mathbb Q_\ell),
\end{equation}
then $X$ is special; that is,
\begin{equation}\tag{3.2.2}
\mathfrak F_p\quad\text{is nilpotent on}\quad H^n(X,\mathcal O_X).
\end{equation}

\medskip
\noindent\emph{Proof.}
We may suppose that $k$ is finitely generated over $\mathbb F_p$ (the case
$p=0$ being trivial), and that there is a finitely generated
$\mathbb F_p$-subalgebra $B$ of $k$ such that $X/k$ is the generic fiber of a
smooth relative complete intersection $\widetilde X$:
\begin{equation}\tag{3.2.3}
\begin{tikzcd}[column sep=4.5em,row sep=3em]
\widetilde X \arrow[d,"\pi"'] &
X=\widetilde X_\eta \arrow[l] \arrow[d] \\
\operatorname{Spec}(B)=S & \eta=\operatorname{Sp}(k)\ . \arrow[l]
\end{tikzcd}
\end{equation}
The $B$-module $R^n\pi_*(\mathcal O_{\widetilde X})$ is locally free, and its
formation commutes with every base change.  Consequently the set of points
\begin{equation}\tag{3.2.4}
\{\,s\in S\mid X_s\text{ is special}\,\}
\end{equation}
is closed in $S$.

By (XX~2.3), after shrinking $S$ we may suppose that
\begin{equation}\tag{3.2.5}
\forall s\in S,\qquad
H^n(X_{\bar s},\mathbb Q_\ell)=N^1H^n(X_{\bar s},\mathbb Q_\ell).
\end{equation}
By Corollary~3.1, $X_s$ is special for every closed point $s$ of $S$.
It follows that the closed subset (3.2.4) is all of $S$ and hence, in
particular, that $X=\widetilde X_\eta$ is special.

\medskip
\noindent\emph{Source note.}
In diagram (3.2.3) the diplomatic French preserves the printed
$X=\widetilde X/_{\eta}$ notation.  The corrected French and English use the
standard generic-fiber notation $X=\widetilde X_\eta$, agreeing with the
concluding sentence of the proof.
\end{filecontents*}
% Original input: components/156_expose_XXI_section_4_and_bibliography.tex
\begin{filecontents*}[overwrite,noheader]{sga7ii-cumulative-input-156.tex}
\section*{4. Conclusions}

Combining Theorem~1.1 with Proposition~3.2 gives the following result.

\subsubsection*{Theorem 4.1.}
Fix a characteristic $p$, a dimension $n$, and a multidegree
$(d_1,\ldots,d_r)$.  Let
\begin{equation}\tag{4.1.0}
X_V\longrightarrow V
\end{equation}
denote the universal smooth complete intersection of dimension $n$ and
multidegree $(d_1,\ldots,d_r)$ in characteristic $p$.  If
\begin{equation}\tag{4.1.1}
\sum_{i=1}^{r}d_i\geq n+r+1,
\end{equation}
then there is a dense open subset $U\subset V$ such that
\begin{equation}\tag{4.1.2}
\text{for every }u\in U,\qquad
H^n(X_{\bar u},\mathbb Q_\ell)\ne N^1H^n(X_{\bar u},\mathbb Q_\ell).
\end{equation}

Combining Theorem~1.3 with Proposition~3.2 likewise gives the following.

\subsubsection*{Theorem 4.2.}
Let $X$ be a smooth complete intersection of dimension $n\geq1$ over a
field $k$.  Fix an integer
\begin{equation}\tag{4.2.1}
d\geq n+1.
\end{equation}
Let
\begin{equation}\tag{4.2.2}
X\!\cdot\!H_u\longrightarrow U
\end{equation}
denote the universal family of sections of $X$ by hypersurfaces of degree
$d$ that are smooth of dimension $n-1$.  There is a dense open subset
$T\subset U$ such that
\begin{equation}\tag{4.2.3}
\text{for every }t\in T,\qquad
H^{n-1}(X\!\cdot\!H_{\bar t},\mathbb Q_\ell)
\ne N^1H^{n-1}(X\!\cdot\!H_{\bar t},\mathbb Q_\ell).
\end{equation}

\subsection*{Bibliography}

\begin{enumerate}[label={\textup{[\arabic*]}},leftmargin=2.6em]
\item J.-P. Serre, \emph{Faisceaux algébriques cohérents}, Annals of
  Mathematics \textbf{61} (1955), 197--278.
\item J.-P. Serre, \emph{Sur la topologie des variétés algébriques en
  caractéristique $p$}, in \emph{Symposium Internacional de Topología
  Algebraica}, Mexico, 1958.
\end{enumerate}

\medskip
\noindent\emph{Source note.}
The diplomatic French preserves the printed lower limit $1$ in (4.1.1) and
the unaccented phrase \emph{de meme}.  The corrected French and English make
the summation index explicit as $i=1$ and use the ordinary accented form.
\end{filecontents*}
% Original input: components/157_expose_XXI_appendix_5_0_through_lemma_5_2_1.tex
\begin{filecontents*}[overwrite,noheader]{sga7ii-cumulative-input-157.tex}
\section*{5. Appendix: Integrality Theorems}

\begin{center}
By P. Deligne
\end{center}

\subsection*{5.0.}
Let $p$ be a prime number, $q$ a power of $p$, $\mathbb F_q$ a field with
$q$ elements, $E$ the completion of a number field at a place dividing a
prime number $\ell\ne p$, and $T$ a set of prime numbers (for example,
$T=\varnothing$).

An element $\alpha$ of an algebraic closure of $E$ is called
\emph{$T$-integral} if it is algebraic over $\mathbb Q$ and integral over
$\mathbb Z[(1/r)_{r\in T}]$.  An endomorphism $F$ of a finite-dimensional
$E$-vector space is called $T$-integral if its eigenvalues are $T$-integral.

Let $k$ be a finite field with $q'$ elements, let $\bar k$ be an algebraic
closure of $k$, and let
\[
\mathfrak G=\operatorname{Gal}(\bar k/k)
\]
be the Galois group of $k$.  We denote by $f\in\mathfrak G$ the Frobenius
substitution
\[
f(x)=x^{q'},
\]
and by $F$ the ``geometric'' Frobenius, the inverse of $f$ in $\mathfrak G$.
A representation $\rho$ of $\mathfrak G$ on a finite-dimensional
$E$-vector space is called $T$-integral if $\rho(F)$ is $T$-integral.

If $X$ is a scheme of finite type over $\mathbb F_q$ and $x$ is a closed
point of $X$, we denote by $F_x$ the geometric Frobenius element of the
Galois group of $k(x)$.  If $\bar x$ is a geometric point of $X$ lying over
$x$ and $\underline G$ is an $E$-sheaf on $X$, then $F_x$ acts, by transport
of structure, on the fiber $G_{\bar x}$.

\subsubsection*{Definition 5.1.}
A constructible $E$-sheaf $\underline G$ on a scheme $X$ of finite type over
$\mathbb F_q$ is called $T$-integral if, for every geometric point $\bar x$
of $X$ lying over a closed point $x$ of $X$, $F_x$ is a $T$-integral
automorphism of $G_{\bar x}$.

If $a:X\longrightarrow\operatorname{Spec}(\mathbb F_q)$ is a scheme of
finite type over $\mathbb F_q$ and $\underline G$ is a constructible
$E$-sheaf on $X$, we write $\underline H^*(X,\underline G)$ for the
$E$-sheaves $R^*a_*\underline G$ on
$\operatorname{Spec}(\mathbb F_q)_{\acute e t}$.  If
$\overline{\mathbb F}_q$ is an algebraic closure of $\mathbb F_q$ and
$\overline X=X\times_{\mathbb F_q}\overline{\mathbb F}_q$, these sheaves
identify with the cohomology groups $H^*(\overline X,\underline G)$ regarded
as Galois modules, that is, endowed with an action of the geometric
Frobenius $F$.  We use the same convention for cohomology with proper
support.

\subsubsection*{Lemma 5.2.1.}
Let $\underline G$ be a constructible $E$-sheaf on a scheme $X$ of finite
type over $\mathbb F_q$, of dimension at most $n$.

\begin{enumerate}[label=(\roman*)]
\item There is a closed subset $Y$ of $X$, of dimension $0$, such that the
  evident map
  \[
  \underline H^0(X,\underline G)\longrightarrow
  \underline H^0(Y,\underline G)
  \]
  is injective.
\item If $X$ is separated and $U$ is an open subset of $X$ whose complement
  $Z$ has dimension less than $n$, there is a closed subset $Y$ of $U$, of
  dimension $0$, and a surjective map
  \[
  \underline H^0(Y,\underline G)(-n)\longrightarrow
  \underline H_c^{2n}(X,\underline G).
  \]
\end{enumerate}

\medskip
\noindent\emph{Proof.}
Assertion (i) is evident.  We prove (ii).  After replacing $X$ by
$X_{\mathrm{red}}$ and shrinking $U$, we may suppose that $U$ is smooth and
that $\underline G|_U$ is twisted constant.  Since $\dim Z<n$, we have
\[
0=\underline H_c^{2n-1}(Z,\underline G)
\longrightarrow \underline H_c^{2n}(U,\underline G)
\xrightarrow{\sim}\underline H_c^{2n}(X,\underline G)
\longrightarrow \underline H_c^{2n}(Z,\underline G)=0,
\]
and Poincar\'e duality (SGA~4, Expos\'e~XVIII, and [2]) gives
\[
\underline H_c^{2n}(U,\underline G)\simeq
\underline H^0(U,\underline G^\vee(n))^\vee.
\]
The conclusion follows from (i).
\end{filecontents*}
% Original input: components/158_expose_XXI_theorem_5_2_2_through_corollary_5_3_setup.tex
\begin{filecontents*}[overwrite,noheader]{sga7ii-cumulative-input-158.tex}
\subsubsection*{Theorem 5.2.2.}
Let $X$ be a separated scheme of finite type over $\mathbb F_q$, of
dimension at most $n$, and let $\underline G$ be a constructible
$T$-integral $E$-sheaf.  Then the endomorphisms $F$ and $q^{n-i}F$ of
$\underline H_c^i(X,\underline G)$ are $T$-integral.

\medskip
\noindent\emph{Proof.}
The case $\dim X\leq0$ is trivial, and the cases $i=0$ and $i=2n$ follow
from Lemma~5.2.1.  Let $Z(X,\underline G,t)$ be the formal power series
\[
Z(X,\underline G,t)=
\prod_{x\in X}\det\!\left(1-F_xt^{\deg(x)}\,;\,\underline G_x\right)^{-1}
\in E[[t]].
\]
By SGA~5, Expos\'e~XV, 3.2, one has
\begin{equation}\tag{5.2.3}
Z(X,\underline G,t)=
\prod_i\det\!\left(1-Ft\,;\,\underline H_c^i(X,\underline G)\right)^{(-1)^{i+1}}.
\end{equation}

We first prove the theorem for $n=1$.  Then
$\underline H_c^i(X,\underline G)=0$ for $i>2$ (SGA~4, Expos\'e~X, 4.3,
or Expos\'e~XVII, 5.2.8.1).  By hypothesis the coefficients of
$Z(X,\underline G,t)$ are $T$-integral.  By the cases $i=0$ and $i=2n=2$
already established, the product of $Z(X,\underline G,t)$ with
\[
\det\!\left(1-Ft\,;\,\underline H_c^0(X,\underline G)\right)
\det\!\left(1-Ft\,;\,\underline H_c^2(X,\underline G)\right)
\]
still has $T$-integral coefficients.  This product is precisely
$\det(1-Ft\,;\,\underline H_c^1(X,\underline G))$, proving the theorem for
$n=1$.

We now proceed by induction on $n=\dim X$.  Let $U$ be a Zariski-open subset
of $X$ whose complement $Z$ has dimension less than $n$, and suppose there
is a morphism $f:U\longrightarrow Y$ to a curve, with fibers of dimension
at most $n-1$:
\begin{center}
\begin{tikzcd}
U \arrow[r,hook,"j"] \arrow[d,"f"'] & X & Z \arrow[l,hook',"i"'] \\
Y
\end{tikzcd}
\qquad .
\end{center}
The exact sequence
\[
0\longrightarrow j_!j^*\underline G\longrightarrow\underline G
\longrightarrow i_!i^*\underline G\longrightarrow0
\]
induces exact sequences
\[
\underline H_c^i(U,\underline G)\longrightarrow
\underline H_c^i(X,\underline G)\longrightarrow
\underline H_c^i(Z,\underline G).
\]
By induction it is therefore enough to verify the theorem for
$\underline G$ on $U$.  Applying the induction hypothesis to the fibers of
$f$, together with proper base change, shows that the sheaves
$R^if_!\underline G$ and $R^if_!\underline G(i-(n-1))$ on $Y$ are
$T$-integral.  The Leray spectral sequence for $f$ is
\[
E_2^{pq}=\underline H_c^p(Y,R^qf_!\underline G)
\Longrightarrow\underline H_c^{p+q}(U,\underline G).
\]
By what precedes and by the theorem already proved for $Y$, both
$E_2^{pq}$ and $E_2^{pq}(p+q-n)$ are $T$-integral.  The theorem follows.

\subsubsection*{Corollary 5.3.}
Let $X$ be a scheme of finite type over $\mathbb F_q$, of dimension at most
$1$, let $U$ be an open subset of $X$, let $j:U\hookrightarrow X$ be the
inclusion, and let $\underline G$ be a constructible $E$-sheaf on $U$.

\begin{enumerate}[label=(\alph*)]
\item If $\underline G$ is $T$-integral, then $j_*\underline G$ is
  $T$-integral.
\item If $G_x^\vee$ is $T$-integral for every $x\in U$, then
  $(j_*\underline G)_x^\vee$ is $T$-integral for every $x\in X$.
\end{enumerate}

\medskip
\noindent\emph{Proof.}
We reduce first to the case where $X$ is affine, and then to the case where
$X$ is projective.  Define $\underline H$ by the exact sequence
\[
0\longrightarrow j_!\underline G\longrightarrow j_*\underline G
\longrightarrow\underline H\longrightarrow0.
\]
If $Y$ is the finite support of $\underline H$, the corresponding
cohomology sequence gives an exact sequence
\[
\underline H^0(U,\underline G)\longrightarrow
\underline H^0(Y,\underline H)\longrightarrow
\underline H_c^1(U,\underline G).
\]
Lemma~5.2.1 and Theorem~5.2.2 show that the two outer terms are
$T$-integral.  Hence $\underline H^0(Y,\underline H)$ is $T$-integral,
which proves (a).

Let $X'$ be the normalization of the pure one-dimensional part of
$X_{\mathrm{red}}$, and choose an open subset $U'$ of $U$, dense in the
smooth one-dimensional locus of $U_{\mathrm{red}}$, such that
$\underline G|_{U'}$ is twisted constant:
\begin{center}
\begin{tikzcd}
U' \arrow[r,hook,"j'"'] \arrow[d,hook] & X' \arrow[d,hook,"p"] & {} \\
U \arrow[r,hook,"j"'] & X & Y \arrow[l,hook']
\end{tikzcd}
\qquad ,
\end{center}
Then
\[
(p_*j'_*\underline G)|_Y\simeq(j_*\underline G)|_Y.
\]
It is therefore enough to treat the case in which $X$ is smooth of
dimension one and $\underline G|_U$ is twisted constant.  The sheaf
$\underline G^\vee$ is then $T$-integral, so we already know that
$(j_*\underline G^\vee)_x$ is $T$-integral for every $x\in X$.

Let $\mathfrak G_x$ denote the Galois group of the fraction field of
$\widehat{\mathcal O}_{X,x}$, and let $\mathfrak I_x$ be the inertia group.
The sheaf $\underline G$ defines a representation of $\mathfrak G_x$, again
denoted by $\underline G$.  The $\mathfrak G_x/\mathfrak I_x=
\operatorname{Gal}(\overline{k(x)}/k(x))$-module
$(j_*\underline G^\vee)_x$ (respectively,
$(j_*\underline G)_x^\vee$) identifies with the module of invariants
(respectively, coinvariants) in $\underline G^\vee$.  It remains to prove:

\medskip
\noindent\emph{Source note.}
The diplomatic French preserves the printed coefficient field $F[[t]]$ in
the zeta series and the printed restriction of $\underline G$ to $U$ in the
normalization reduction.  The corrected French and English use the
coefficient field $E$ fixed in 5.0 and the newly chosen open $U'$.  Direct
scan comparison also repairs the inherited joined words, Fraktur Galois
symbols, and doubled consonant in \emph{D\'esignons} to the printed forms.

\end{filecontents*}
% Original input: components/159_expose_XXI_lemmas_5_3_1_and_5_3_2.tex
\begin{filecontents*}[overwrite,noheader]{sga7ii-cumulative-input-159.tex}
\subsubsection*{Lemma 5.3.1.}
Let $K$ be a local field with residue field $\mathbb F_q$, let $G$ be its
Galois group, let $I$ be the inertia group, and let $\rho$ be a continuous
representation of $G$ on an $n$-dimensional $E$-vector space $V$ (cf. 5.0).
Let $V^I$ and $V_I$ denote respectively the space of invariants and the
space of coinvariants under $I$, regarded as $G/I$-modules.  Put
$r=\dim_E V^I$, and let $\alpha_1,\ldots,\alpha_r$ be the eigenvalues of
$F$ on $V^I$.  Then there are integers $m_i\geq0$ such that the eigenvalues
of $F$ on $V_I$ are
\[
\alpha_iq^{m_i}\qquad(1\leq i\leq r).
\]

\medskip
\noindent\emph{Proof.}
Let $I'$ be the largest subgroup of $I$ of order prime to $\ell$.  The group
$\rho(I')$ is finite, since it is a compact subgroup of $\operatorname{GL}(V)$
of order prime to $\ell$.  Hence
\[
V_I\simeq\bigl(V^{\rho(I')}\bigr)_I.
\]
After replacing $V$ by $V^{I'}$, we may therefore suppose that $I'$ acts
trivially on $V$.  As $G/I$-modules one has
\[
I/I'\simeq\mathbb Z_\ell(1).
\]
Let $T$ be a topological generator of $I/I'$.  For each $\alpha$, let
$V^{(\alpha)}$ be the generalized eigenspace of $T$ on $V$ with eigenvalue
$\alpha$.  Then
\[
V=\bigoplus_\alpha V^{(\alpha)},
\]
and the decomposition
\[
V=V^{(1)}\oplus\bigoplus_{\alpha\ne1}V^{(\alpha)}
\]
is stable under $G$ and compatible with duality.  We may thus replace $V$
by $V^{(1)}$ and suppose that $I$ acts unipotently.

\subsubsection*{Lemma 5.3.2.}
Let $V$ be a finite-dimensional vector space over a field $k$ of
characteristic $0$, let $I$ be a one-parameter additive group, let
$\rho:I\longrightarrow\operatorname{GL}(V)$ be a representation, and let
$N$ be an infinitesimal generator of $I$.  Set
\[
V^{I,k}=V^I\cap\operatorname{Im}(N^k)
\]
and
\[
V_{I,k}=\operatorname{Im}\!\left(\operatorname{Ker}(N^{k+1})
\longrightarrow V_I\right).
\]
Then $N^k$ induces an isomorphism
\[
N^k:V_{I,k}/V_{I,k-1}\xrightarrow{\sim}
V^{I,k}/V^{I,k+1},
\]
and defines canonical isomorphisms
\[
V_{I,k}/V_{I,k-1}\otimes I^{\otimes k}
\simeq V^{I,k}/V^{I,k+1}.
\]

\medskip
\noindent\emph{Proof.}
Via $N$, representations $V$ of $I$ identify with $k[T]$-modules killed by
a power of $T$.  It is enough to verify Lemma~5.3.2 when $V$ corresponds to
an indecomposable $k[T]$-module $k[T]/T^k$, where the assertion is immediate.

We now finish the proof of Lemma~5.3.1.  Lemma~5.3.2 supplies increasing and
decreasing filtrations $V_{I,k}$ and $V^{I,k}$ of $V_I$ and $V^I$, together
with isomorphisms of $G/I$-modules
\[
V^{I,k}/V^{I,k+1}\simeq
\bigl(V_{I,k}/V_{I,k-1}\bigr)(k).
\]
The assertion follows.

\medskip
\noindent\emph{Source note.}
The diplomatic French enumerates the eigenvalues on $V^I$ by the dimension
$n$ of $V$, cites the lemma as ``3.2'', and calls $k[T]/T^k$ irreducible.
The corrected French and English use $r=\dim_E V^I$, the exact locator
5.3.2, and the mathematically correct term \emph{indecomposable}.
\end{filecontents*}
% Original input: components/160_expose_XXI_theorem_5_4.tex
\begin{filecontents*}[overwrite,noheader]{sga7ii-cumulative-input-160.tex}
\subsubsection*{Theorem 5.4.}
Let $X$ be a separated scheme of finite type over $\mathbb F_q$, of
dimension $n$, and let $\underline G$ be a constructible $E$-sheaf on $X$.
Suppose that, for every $x\in X$, the endomorphism $F_x$ of
$G_x^\vee$ is $T$-integral.  Then the endomorphisms $q^iF^\vee$ and
$q^nF^\vee$ of $\underline H_c^i(X,\underline G)^\vee$ are $T$-integral.

\medskip
\noindent\emph{Proof.}
As in Theorem~5.2.2, we reduce to the case in which $X$ is a smooth curve
and $\underline G$ is twisted constant on $X$.  Let $j$ be the inclusion of
$X$ in a smooth projective curve $X'$ compactifying it.

By Poincar\'e duality,
$\underline H^i(X',j_*\underline G)$ is dual to
$\underline H^{2-i}(X',j_*\underline G^\vee(1))$.  By Corollary~5.3(a),
Theorem~5.2.2 applies to $j_*\underline G^\vee$:
\begin{align*}
\underline H^0(X',j_*\underline G^\vee)
  &\simeq \underline H^2(X',j_*\underline G)^\vee(-1)
  &&\text{is \(T\)-integral},\\
\underline H^1(X',j_*\underline G^\vee)
  &\simeq \underline H^1(X',j_*\underline G)^\vee(-1)
  &&\text{is \(T\)-integral},\\
\underline H^2(X',j_*\underline G^\vee)(1)
  &\simeq \underline H^0(X',j_*\underline G)^\vee
  &&\text{is \(T\)-integral}.
\end{align*}
This proves the theorem for $j_*\underline G$.

Define $\underline H$ by the exact sequence
\[
0\longrightarrow j_!\underline G\longrightarrow j_*\underline G
\longrightarrow\underline H\longrightarrow0.
\]
If $Y$ is the support of $\underline H$, the long exact cohomology sequence
provides
\begin{center}
\begin{tikzcd}[row sep=large,column sep=large]
{} & \underline H_c^0(X,\underline G)
  \arrow[r,hook] & \underline H^0(X',j_*\underline G) & ,\\
\underline H^0(Y,\underline H) \arrow[r] &
  \underline H_c^1(X,\underline G) \arrow[r] &
  \underline H^1(X',j_*\underline G) & ,\\
{} & \underline H_c^2(X,\underline G)
  \arrow[r,"\sim"] & \underline H^2(X',j_*\underline G) & ,
\end{tikzcd}
\end{center}
Thus the theorem for $(X,\underline G)$ follows from the theorem for
$(X',j_*\underline G)$.
\end{filecontents*}
% Original input: components/161_expose_XXI_section_5_5.tex
\begin{filecontents*}[overwrite,noheader]{sga7ii-cumulative-input-161.tex}
\subsection*{5.5.}
The preceding method of fibration into curves also gives Lang--Weil
estimates~[3] for the complex absolute values of the eigenvalues of
Frobenius.

Let $X$ be a scheme of finite type over $\mathbb F_q$, let $\underline G$
be a constructible $E$-sheaf on $X$, and let $\lambda\in\mathbb R$.  Suppose
that, for every point $x\in X$, with $\#k(x)=q_x$, the eigenvalues of $F_x$
acting on $G_x$ are algebraic numbers $\alpha$ all of whose complex absolute
values satisfy
\[
|\alpha|\leq q_x^\lambda.
\]
We shall then write
\[
\log|\underline G|\leq\lambda.
\]

\subsubsection*{Lemma 5.5.1.}
Let $X$ be a curve over $\mathbb F_q$, by which we mean here a separated
scheme $X$ of finite type over $\mathbb F_q$ and of dimension at most $1$.
Let $\underline G$ be a constructible $E$-sheaf on $X$ such that
$\log|\underline G|\leq\lambda$, with $\lambda\in\mathbb R$.  Then the
complex absolute values of the eigenvalues $\alpha$ of $F$ acting on
$\underline H_c^1(X,G)$ satisfy
\[
|\alpha|\leq q^{\lambda+1}.
\]

\medskip
\noindent\emph{Proof.}
Let $\sigma$ be a complex embedding of $E$, or simply a complex embedding of
the algebraic closure of $\mathbb Q$ in $E$.  It is enough to verify that,
for every $\sigma$, the zeros $\beta$ of the rational function
$\sigma Z(X,\underline G,t)$ satisfy
\[
|\beta^{-1}|\leq q^{\lambda+1}.
\]
Indeed, Lemma~5.2.1 already shows that, if $\alpha$ is an eigenvalue of
Frobenius acting on $\underline H_c^0(X,\underline G)$ or on
$\underline H_c^2(X,\underline G)$, then
$|\sigma\alpha|\leq q^{\lambda+1}$.

We verify that the series
\[
\log\sigma Z(X,\underline G,t)
=\sum_k\frac{t^k}{k}
  \sum_{x\in X(\mathbb F_{q^k})}
  \sigma\operatorname{Tr}(F,x^*\underline G)
\]
converges for $|t|<q^{-\lambda-1}$.  Representing $X$ as a ramified covering
of $\mathbb P^1$, one obtains the elementary estimate
\[
\#X(\mathbb F_{q^k})\leq A q^k.
\]
On the other hand, the hypothesis gives
\[
\left|\sigma\operatorname{Tr}(F,x^*\underline G)\right|
\leq B(q^k)^\lambda.
\]
Thus the coefficient $c_k$ of $t^k$ in
$\log\sigma Z(X,\underline G,t)$ satisfies
\[
|c_k|\leq\frac{1}{k}\,C(q^{1+\lambda})^k,
\]
and the assertion follows.

\subsubsection*{Theorem 5.5.2.}
Let $X$ be a separated scheme of finite type over $\mathbb F_q$, of
dimension at most $n$, and let $\underline G$ be a constructible $E$-sheaf
on $X$ such that $\log|\underline G|\leq\lambda$, with
$\lambda\in\mathbb R$.  Then the complex absolute values of the eigenvalues
$\alpha$ of Frobenius acting on $\underline H_c^i(X,\underline G)$ satisfy
\[
|\alpha|\leq q^{\lambda+i}
\qquad\text{and}\qquad
|\alpha|\leq q^{\lambda+n}.
\]

\medskip
\noindent\emph{Proof.}
If $n\leq0$, the theorem is trivial.  If $n=1$, its cases $i=0$ and $i=2$
follow from Lemma~5.2.1, while the case $i=1$ follows from
Lemma~5.5.1.  We then proceed by induction on $n$, as in the proof of
Theorem~5.2.2.

\subsubsection*{Corollary 5.5.3.}
Let $X$ be a separated scheme of finite type and dimension $n$ over
$\mathbb F_q$.  Then the eigenvalues $\alpha$ of the geometric Frobenius
endomorphism $F$ of $\underline H_c^i(X,\mathbb Q_\ell)$ are algebraic
integers.  Moreover:
\begin{enumerate}[label=(\roman*)]
\item for $\ell'\ne p$, $\alpha$ is an $\ell'$-adic unit;
\item for $i>n$, $\alpha$ is divisible by $q^{i-n}$;
\item for $i\leq n$ (respectively, $i\geq n$),
  $q^i\alpha^{-1}$ (respectively, $q^n\alpha^{-1}$) is an algebraic integer;
\item for $0<i$, the complex absolute values of $\alpha$ satisfy
  \[
  |\alpha|\leq q^{i-\frac12}
  \qquad\text{and}\qquad
  |\alpha|\leq q^n.
  \]
\end{enumerate}

\medskip
\noindent\emph{Proof.}
By Theorem~5.2.2, $\alpha$ and $q^{n-i}\alpha$ are algebraic integers.
Assertions (i)--(iii) then follow from Theorem~5.4.  For $i>n$, assertion
(iv) follows from Theorem~5.5.2.  To obtain (iv) when $i\leq n$, we must
combine Theorem~5.5.2 with the Riemann hypothesis for curves.

We argue by induction on $n$, the case $n=0$ being trivial.  If $U$ is an
open subset of $X$ whose complement $Z$ has dimension less than $n$, the
exact sequence
\[
\underline H_c^i(U,\mathbb Q_\ell)\longrightarrow
\underline H_c^i(X,\mathbb Q_\ell)\longrightarrow
\underline H_c^i(Z,\mathbb Q_\ell)
\]
shows that it is enough to prove the desired assertion for $U$.  After
first replacing $X$ by $X_{\mathrm{red}}$, this reduces us to the case in
which $X$ is smooth and affine and there is a smooth morphism
$f:X\longrightarrow S$ satisfying:
\begin{enumerate}[label=(\alph*)]
\item $\dim S\leq n-1$;
\item the fibers of $f$ are purely one-dimensional.
\end{enumerate}
Since the fibers of $f$ are affine curves, one has
\[
R^0f_!\mathbb Q_\ell=0.
\]
By Weil~[4],
\begin{equation}\tag{5.5.3.1}
\log\bigl|R^1f_!\mathbb Q_\ell\bigr|\leq\frac12,
\end{equation}
and Lemma~5.2.1 applied to the fibers of $f$ gives
\begin{equation}\tag{5.5.3.2}
\log\bigl|R^2f_!\mathbb Q_\ell\bigr|\leq1.
\end{equation}
The Leray spectral sequence for $f$ reduces here to a long exact sequence;
the relevant part is
\[
\underline H_c^{i-2}(S,R^2f_!\mathbb Q_\ell)\longrightarrow
\underline H_c^i(X,\mathbb Q_\ell)\longrightarrow
\underline H_c^{i-1}(S,R^1f_!\mathbb Q_\ell).
\]
By (5.5.3.1) (respectively, (5.5.3.2)) and Theorem~5.5.2, the complex
absolute values of the eigenvalues $\alpha$ of Frobenius on
$\underline H_c^{i-1}(S,R^1f_!\mathbb Q_\ell)$ (respectively, on
$\underline H_c^{i-2}(S,R^2f_!\mathbb Q_\ell)$) satisfy
\[
|\alpha|\leq q^{i-1+\frac12}=q^{i-\frac12}
\]
(respectively,
\[
|\alpha|\leq q^{i-2+1}=q^{i-1}\leq q^{i-\frac12}).
\]
The corollary follows.

\medskip
\noindent\emph{Source note.}
The diplomatic French preserves three printed slips: the coefficient
$c_k$ is said to belong to $\sigma Z$ although the displayed series and
the factor $1/k$ make it a coefficient of $\log\sigma Z$; the proof of
Theorem 5.5.2 refers to itself rather than to the parallel induction in
Theorem 5.2.2; and the proof of Corollary 5.5.3 says to prove Theorem 5.2.2
for $U$ although the argument is proving assertion (iv) of the corollary.
The corrected French and English give the coherent readings.  Direct scan
comparison also restores the notation underlines omitted by the inherited
French TeX in the statement of Theorem 5.5.2.
\end{filecontents*}
% Original input: components/162_expose_XXI_section_5_6.tex
\begin{filecontents*}[overwrite,noheader]{sga7ii-cumulative-input-162.tex}
\subsection*{5.6.}
To obtain results analogous to the preceding ones for cohomology without
supports, it seems necessary to have resolution of singularities available.
In dimension at most $1$, the problem is solved by Corollary~5.3.

\subsubsection*{Theorem 5.6.}
Let $N$ be an integer, and suppose that strong resolution of singularities
in the sense of Hironaka is available for schemes of finite type over
$\mathbb F_q$ of dimension at most $N$ (at present one may therefore take
$N=2$~[1]).  Let $f:X\longrightarrow Y$ be a morphism of schemes of finite
type over $\mathbb F_q$, and let $\underline G$ be a constructible
$T$-integral $E$-sheaf on $X$.  If $\dim X\leq N$, then the $E$-sheaves
$R^if_*\underline G$ on $Y$ are $T$-integral.

\medskip
\noindent\emph{Proof.}
We argue by induction on the dimension $n\leq N$ of $X$.  Let
$j:U\hookrightarrow X$ be an open subset of $X$, smooth and affine, such
that $\underline G|_U$ is twisted constant and the dimension of its
complement $Z$ is less than $n$.  By resolution of singularities, the
sheaves $R^ij_*(j^*\underline G)$ are constructible
(SGA~5, Expos\'e~I; [2]); for $i>0$ they are supported on $Z$, and for
$i=0$, $j_*(j^*\underline G)$ differs from $\underline G$ only on $Z$.
The induction hypothesis therefore shows that it is enough to prove the
theorem for the sheaf $j^*\underline G$ and the morphisms $j$ and
$f\circ j$.  Thus we may further suppose that $X$ is smooth and affine and
that $\underline G$ is twisted constant, that is, comes from a continuous
representation of the fundamental group of $X$ on a finite-dimensional
$E$-vector space.

Let $R$ be the ring of integers of $E$, let $\lambda$ be a uniformizer, let
$e$ be the absolute ramification index, and choose an integer
$b>e/(\ell-1)$.  Since $X$ is smooth, hence normal, and its fundamental
group is profinite, there is a twisted-constant $R$-sheaf
$\underline G^\circ$ such that
\[
\underline G\simeq\underline G^\circ\otimes_R E.
\]
Let $p:X'\longrightarrow X$ be the \'etale covering that trivializes the
locally constant sheaf
$\underline G^\circ\otimes_R R/\lambda^b$.  The sheaf $\underline G$ is a
direct summand of $p_*p^*\underline G$, so it is enough to prove the theorem
for $p^*\underline G$ and $f\circ p$.  It remains to treat the case in which
the following conditions hold:
\[
(*)\quad
\begin{minipage}{0.82\linewidth}
$X$ is smooth and affine, and there is a twisted-constant $R$-sheaf
$\underline G^\circ$ such that
$\underline G\simeq\underline G^\circ\otimes_R E$ and
$\underline G^\circ\otimes_R R/\lambda^b$ is constant.
\end{minipage}
\]

Factor $f$ as a proper morphism followed by an open immersion,
$f=\bar f\circ j$:
\begin{center}
\begin{tikzcd}
X \arrow[r,hook,"j",swap] \arrow[d,"f"'] &
X' \arrow[dl,"\bar f"]\\
Y & {}
\end{tikzcd}
\qquad .
\end{center}
The theorem has already been proved for $\bar f$ (5.2), so it is enough to
prove it for $j$.  Applying resolution of singularities to $X'$, we may
suppose that $X'$ is smooth and that $X$ is the complement in $X'$ of a
normal-crossings divisor which is a union of smooth divisors.  Hypothesis
$(*)$ implies that $\underline G$ is tamely ramified along this divisor.
We then have:

\subsubsection*{Lemma 5.6.1.}
Let $X$ be a smooth scheme of finite type over $\mathbb F_q$, let $D$ be a
normal-crossings divisor in $X$ which is a union of smooth divisors $D_i$,
let $j:U\hookrightarrow X$ be the complement of $D$ in $X$, and let
$\underline G$ be a twisted-constant $E$-sheaf on $U$, tamely ramified
along $D$.  If $\underline G$ is $T$-integral, then the sheaves
$R^ij_*\underline G$ are $T$-integral.

We shall admit the following result (cf. SGA~1, Expos\'e~XIII).

\subsubsection*{Lemma 5.6.2.}
\begin{enumerate}[label=(\roman*)]
\item Under the hypotheses of Lemma~5.6.1, let $D'$ be a component of $D$,
  let $D''$ be the union of the other components, and put
  $U'=X-D''$ and $D^{\prime\circ}=D'\cap U'$:
  \begin{center}
  \begin{tikzcd}
  {} & D^{\prime\circ}
    \arrow[r,hook,"j'_D",swap] \arrow[d,"i^\circ"] &
    D' \arrow[d,"i"]\\
  U \arrow[r,hook,"j_1",swap] &
    U' \arrow[r,hook,"j'",swap] & X
  \end{tikzcd}
  \qquad .
  \end{center}
  Then the base-change map
  \[
  i^*Rj'_*(Rj_{1*}\underline G)
  \longrightarrow
  Rj'_{D*}\,i^{\circ *}(Rj_{1*}\underline G)
  \]
  is an isomorphism, and the sheaves
  $i^{\circ *}Rj_{1*}\underline G$ are twisted constant and tamely ramified
  along $D'-D^{\prime\circ}$.

\item If $C$ is a smooth curve on $X$, transverse to $D$ at a smooth point
  $x$ of $D$,
  \begin{center}
  \begin{tikzcd}
  C\cap U \arrow[r,"j_c",swap] \arrow[d,"k'"] &
    C \arrow[d,"k"]\\
  U \arrow[r,"j",swap] & X
  \end{tikzcd}
  \qquad ,
  \end{center}
  then the base-change map
  \[
  k^*Rj_*\underline G\longrightarrow Rj_{c*}\,k^{\prime *}\underline G
  \]
  is an isomorphism at $x$.
\end{enumerate}

Let us deduce Lemma~5.6.1 from Lemma~5.6.2 by induction on the dimension,
temporarily admitting Lemma~5.6.1 for curves.  Lemma~5.6.2(ii) and this
hypothesis show that the sheaves $R^ij_*\underline G$ are $T$-integral at
every smooth point of $D$.  In particular, in the notation of
Lemma~5.6.2(i), the sheaves $R^ij_{1*}\underline G$ are $T$-integral.
Lemma~5.6.2(i), together with the induction hypothesis applied to
$i^{\circ *}Rj_{1*}\underline G$ and to $j'_D$, therefore shows that the
sheaves $i^*R^ij_*\underline G$ are $T$-integral, whichever component $D'$
of $D$ is chosen.  This proves Lemma~5.6.1.

It remains to handle curves.  Under the hypotheses of Lemma~5.6.1, let $x$
be a point of $X-U$, let $\mathfrak G$ be the Galois group of the fraction
field of $\widehat{\mathcal O}_{x,X}$, and let $\mathfrak J$ be the inertia
group.  Denote again by $\underline G$ the representation of $\mathfrak G$
defined by the sheaf, and let $\bar x$ be a geometric point of $X$ localized
at $x$.  There are isomorphisms of Galois
$\mathfrak G/\mathfrak J$-modules
\[
(j_*\underline G)_{\bar x}\simeq\underline G^{\mathfrak J},
\]
\[
(R^1j_*\underline G)_{\bar x}\simeq\underline G_{\mathfrak J}(-1),
\]
and
\[
(R^ij_*\underline G)_{\bar x}=0\qquad\text{for }i\geq2.
\]
The assertion follows from Corollary~5.3(a) and Lemma~5.3.1.

% Source: scan idx 407 | printed folio 400.
\section*{Bibliography}
\begin{enumerate}[label={\textup{[\arabic*]}},leftmargin=2.4em]
\item S.~S. Abhyankar, \emph{Resolution of Singularities of Embedded Algebraic Surfaces}, Academic Press, New York, 1966.
\item J.-P.~Jouanolou, \emph{Cohomologie $\ell$-adique des schémas}, forthcoming in \emph{Lecture Notes in Mathematics}.
\item S.~Lang and A.~Weil, \emph{Number of Points of Varieties in Finite Fields}, \emph{Amer. J. Math.} \textbf{76} (1954), 819--827.
\item A.~Weil, \emph{Sur les courbes algébriques et les variétés qui s'en déduisent}, \emph{Actualités Sci. Indust.} no.~1041, Hermann, Paris, 1948.
\end{enumerate}

\medskip
\noindent\emph{Source note.}
Direct authority comparison restores the factorization diagram's $j$ label
to the lower side of its arrow and restores the printed Fraktur
$\mathfrak G$ in the curve case.  The diplomatic French preserves the
printed singular \emph{morphisme}, the punctuation collision before
hypothesis $(*)$, and the shortened reference 6.2(i).  Corrected French and
English use plural \emph{morphismes}, ordinary sentence punctuation, and the
exact reference Lemma 5.6.2(i).
\end{filecontents*}
% Original input: components/163_expose_XXII_title_introduction_and_section_1_0.tex
\begin{filecontents*}[overwrite,noheader]{sga7ii-cumulative-input-163.tex}
\clearpage
\begin{center}
{\Large\bfseries EXPOSÉ XXII}\par
\vspace{1.5cm}
{\Large\bfseries A CONGRUENCE FORMULA FOR THE ZETA FUNCTION}\par
\vspace{0.8cm}
{\large by N. Katz}\par
\end{center}

\vspace{1.5cm}
\section*{0. Introduction}

Let $X$ be a proper scheme over $\mathbb F_q$. Its zeta function, regarded as
an element of $\mathbb Z[[T]]$, may be reduced modulo $p$; we wish to compute
it modulo $p$. We obtain an expression of the kind one might have expected
from fixed-point formulas of ``Woods Hole'' type: an alternating product of
characteristic polynomials of a suitable ``Frobenius'' acting on any one of
the following three cohomology spaces:

\begin{enumerate}[label=\arabic*)]
\item coherent cohomology $H^i(X,\mathcal O_X)$ (3.1.1);
\item de Rham cohomology
  $H^i_{\mathrm{DR}}(X)=\mathbb H^i(X,\Omega_X^\bullet)$ (3.1.1);
\item étale cohomology $H^i_{\mathrm{et}}(X,\mathbb Z/p\mathbb Z)$ (3.1.2).
\end{enumerate}

This expression should be viewed as an improvement of a special case of the
Woods Hole trace formula, namely that of the structure sheaf and the
Frobenius endomorphism. That case is indeed a corollary of the congruence
formula (cf. 3.2.1).

The method of this exposé is completely naïve. It consists in reducing to the
case of a hypersurface and then ``reducing modulo $p$'' Dwork's computation
of the actual zeta function of a hypersurface. By ``reducing'' carefully, one
obtains the improvement 5.0.6, due to Ax~[1] (and also to Dwork in the smooth
case [4, \S 7]), of the well-known Warning--Chevalley theorem~[8].

Another method, entirely cohomological and in the spirit of SGA~5, which also
applies to ``twisted coefficients,'' has just been found by Deligne and ought
to appear in the new edition of SGA~5. It rests on the trace formula of Woods
Hole type and on Deligne's ``symmetric Künneth formula''
(SGA~4, XVII~5.5.21).

\section*{1. A reminder on ``$p$-linear'' operations}

\subsection*{1.0.}
Let $V$ be a finite-dimensional vector space over a field $k$ of
characteristic $p>0$, and let $q$ be a power of $p$. A map

\begin{equation}\tag{1.0.1}
\varphi:V\longrightarrow V
\end{equation}

is called $q$-\emph{linear} if it is additive and satisfies

\begin{equation}\tag{1.0.2}
\varphi(\lambda v)=\lambda^q\varphi(v)
\qquad\text{for every }\lambda\in k,\ v\in V.
\end{equation}

\subsection*{1.0.3.}
Notice that the $n$-th iterate $\varphi^n$ of a $q$-linear endomorphism is
$q^n$-linear and that, if $k$ is not perfect, the image $\varphi(V)$ of a
$q$-linear endomorphism need not be a $k$-vector subspace.

\subsection*{1.0.4.}
For every extension $K/k$, a $q$-linear endomorphism $\varphi$ of $V$ gives
rise to a unique $q$-linear endomorphism $\varphi_K$ of
$V_K=V\otimes_kK$ satisfying

\begin{equation}\tag{1.0.5}
\varphi_K(\mu\otimes v)=\mu^q\varphi(v)
\qquad\text{for }\mu\in K,\ v\in V.
\end{equation}

\subsection*{1.0.6.}
A $q$-linear endomorphism $\varphi$ of $V$ is called \emph{semisimple} if
its image spans $V$ over $k$, and \emph{nilpotent} if some iterate is zero.
Define the following $\varphi$-stable $k$-subspaces:

\begin{equation}\tag{1.0.7}
V_{\mathrm{ss}}=
\bigcap_{n\geq1}\langle\varphi^n(V)\rangle_k,
\end{equation}

\begin{equation}\tag{1.0.8}
V_{\mathrm{nilp}}=
\bigcup_{n\geq1}\Ker(\varphi^n:V\longrightarrow V).
\end{equation}

Clearly, the restriction of $\varphi$ to $V_{\mathrm{ss}}$ is semisimple,
while its restriction to $V_{\mathrm{nilp}}$ is nilpotent. If $k$ is
perfect, then

\begin{equation}\tag{1.0.9}
V=V_{\mathrm{ss}}\oplus V_{\mathrm{nilp}}.
\end{equation}

If $k$ is not perfect, however, (1.0.9) is false, since there exist injective
$q$-linear endomorphisms which are not semisimple. For an arbitrary extension
$K/k$, one has

\begin{equation}\tag{1.0.10}
(V_K)_{\mathrm{ss}}=V_{\mathrm{ss}}\otimes_kK.
\end{equation}

\end{filecontents*}
% Original input: components/164_expose_XXII_proposition_1_1_through_proposition_1_2_opening.tex
\begin{filecontents*}[overwrite,noheader]{sga7ii-cumulative-input-164.tex}
\subsubsection*{Proposition 1.1.}
Let $\varphi$ be a $q$-linear endomorphism of a finite-dimensional vector
space $V$ over a separably closed field $k$. Suppose that $\varphi$ is
semisimple, and put

\[
V^{1-\varphi}
=\text{the $\mathbb F_q$-vector space of fixed points of $\varphi$}
=\Ker(1-\varphi).
\]

Then

\begin{equation}\tag{1.1.1}
V\xleftarrow{\ \sim\ }
V^{1-\varphi}\otimes_{\mathbb F_q}k.
\end{equation}

\subsubsection*{Corollary 1.1.2.}
There is an isomorphism

\[
V^{1-\varphi^n}\xleftarrow{\ \sim\ }
V^{1-\varphi}\otimes_{\mathbb F_q}\mathbb F_{q^n}.
\]

\medskip
\noindent\emph{Proof of Proposition 1.1.}
(Lang, cf.~[6], p.~119.)
Let $e$ be a $k$-basis of $V$, and write

\begin{equation}\tag{1.1.3}
\varphi(e)=Ae,\qquad A\in\mathrm{GL}(n,k).
\end{equation}

Let

\begin{equation}\tag{1.1.4}
F:\mathrm{GL}(n,k)\longrightarrow\mathrm{GL}(n,k)
\end{equation}

denote the map which raises every coordinate to the $q$-th power. For
$g\in\mathrm{GL}(n,k)$, the basis $ge$ of $V$ consists of fixed points of
$\varphi$ if and only if

\begin{equation}\tag{1.1.5}
F(g)Ag^{-1}=1.
\end{equation}

Let $\bar k$ be an algebraic closure of $k$. The group
$\mathrm{GL}(n,\bar k)$ acts on the variety $\mathrm{GL}(n,\bar k)$ by

\begin{equation}\tag{1.1.6}
\left\{
\begin{aligned}
\mathrm{GL}(n,\bar k)\times\mathrm{GL}(n,\bar k)
&\longrightarrow\mathrm{GL}(n,\bar k),\\
(g,A)&\longmapsto F(g)Ag^{-1}.
\end{aligned}
\right.
\end{equation}

For fixed $A$, the map

\begin{equation}\tag{1.1.7}
g\longmapsto F(g)Ag^{-1}
\end{equation}

is étale, as one sees by computing its tangent map. Its orbits are therefore
open. Since $\mathrm{GL}(n,\bar k)$ is irreducible, there is only one orbit;
in other words, the map

\begin{equation}\tag{1.1.8}
\begin{aligned}
\mathrm{GL}(n,\bar k)&\longrightarrow\mathrm{GL}(n,\bar k),\\
g&\longmapsto F(g^{-1})\,g
\end{aligned}
\end{equation}

is surjective. Since it is also étale and $k$ is separably closed, the map

\begin{equation}\tag{1.1.9}
\begin{aligned}
\mathrm{GL}(n,k)&\longrightarrow\mathrm{GL}(n,k),\\
g&\longmapsto F(g^{-1})\,g
\end{aligned}
\end{equation}

is surjective. Thus $V$ has a $k$-basis consisting of fixed points of
$\varphi$. Since $\varphi$ is $q$-linear, every fixed point of $\varphi$ lies
in the $\mathbb F_q$-vector space spanned by such a $k$-basis. This proves
Proposition~1.1.

\subsubsection*{Corollary 1.1.10.}
Let $\varphi$ be a $q$-linear endomorphism of a finite-dimensional vector
space $V$ over a separably closed field $k$. There is a canonical
isomorphism

\begin{equation}\tag{1.1.11}
V_{\mathrm{ss}}\xleftarrow{\ \sim\ }
V^{1-\varphi}\otimes_{\mathbb F_q}k.
\end{equation}

Indeed, it is enough to apply Proposition~1.1 to $V_{\mathrm{ss}}$.

\subsubsection*{Proposition 1.2.}
Let $\varphi$ be a $q$-linear endomorphism of a finite-dimensional vector
space over a separably closed field $k$. Then the additive map

\[
1-\varphi:V\longrightarrow V
\]

is surjective.

\medskip
\noindent\emph{Proof.}
There is an exact sequence of abelian groups

\[
0\longrightarrow V_{\mathrm{ss}}\longrightarrow V
\longrightarrow V/V_{\mathrm{ss}}\longrightarrow0.
\]
\end{filecontents*}
% Original input: components/165_expose_XXII_proposition_1_2_completion_through_corollary_2_1.tex
\begin{filecontents*}[overwrite,noheader]{sga7ii-cumulative-input-165.tex}
on which $1-\varphi$ acts. By the snake lemma, it is enough to treat
$V_{\mathrm{ss}}$, which is immediate from Proposition~1.1 by using a basis
consisting of fixed points, and $V/V_{\mathrm{ss}}$, on which $\varphi$ is
nilpotent and hence $1-\varphi$ is invertible.

\medskip
\noindent\emph{Source note.}
The diplomatic French preserves the print's hyphen in \emph{peut-être}, the
singular phrase \emph{coefficient tordus}, \emph{Una application}, the
omitted verb in \emph{si elle additive}, and two spacing/punctuation
collisions. The corrected French and English regularize those surface defects
without changing the mathematical content. Formula (1.0.7) was read directly
from a high-detail authority crop: its large operator is an intersection.

\section*{2. Coherent cohomology and étale cohomology}

\subsubsection*{Proposition 2.0.}
Let $\bar X$ be a proper scheme over a separably closed field $k$, and let

\begin{equation}\tag{2.0.1}
\mathfrak F_p:H^i(\bar X,\mathcal O_{\bar X})
\longrightarrow H^i(\bar X,\mathcal O_{\bar X})
\end{equation}

denote the $p$-linear map induced by

\begin{equation}\tag{2.0.2}
\mathfrak F_p:\mathcal O_{\bar X}\longrightarrow\mathcal O_{\bar X},
\qquad \mathfrak F_p(f)=f^p.
\end{equation}

Denote by

\[
H^i_{\mathrm{et}}(\bar X,\mathbb Z/p\mathbb Z)
\]

the étale cohomology groups with coefficients in the constant sheaf
$\mathbb Z/p\mathbb Z$. Then there are canonical isomorphisms

\begin{equation}\tag{2.0.3}
H^i_{\mathrm{et}}(\bar X,\mathbb Z/p\mathbb Z)
\simeq
H^i(\bar X,\mathcal O_{\bar X})^{1-\mathfrak F_p}.
\end{equation}

\medskip
\noindent\emph{Proof.}
There is an exact sequence of sheaves on $\bar X_{\mathrm{et}}$:

\begin{equation}\tag{2.0.4}
0\longrightarrow\mathbb Z/p\mathbb Z\longrightarrow\mathcal O_{\bar X}
\xrightarrow{\ 1-\mathfrak F_p\ }\mathcal O_{\bar X}\longrightarrow0,
\end{equation}

which gives an exact cohomology sequence

\begin{equation}\tag{2.0.5}
\cdots\longrightarrow H^i_{\mathrm{et}}(\bar X,\mathbb Z/p\mathbb Z)
\longrightarrow H^i_{\mathrm{et}}(\bar X,\mathcal O_{\bar X})
\xrightarrow{\ 1-\mathfrak F_p\ }
H^i_{\mathrm{et}}(\bar X,\mathcal O_{\bar X})\longrightarrow\cdots .
\end{equation}

One has
$H^i_{\mathrm{et}}(\bar X,\mathcal O_{\bar X})\simeq
H^i(\bar X,\mathcal O_{\bar X})$; this is a finite-dimensional
$k$-vector space because $\bar X$ is proper. Applying Proposition~1.2 gives
exact sequences

\begin{equation}\tag{2.0.6}
0\longrightarrow H^i_{\mathrm{et}}(\bar X,\mathbb Z/p\mathbb Z)
\longrightarrow H^i(\bar X,\mathcal O_{\bar X})
\xrightarrow{\ 1-\mathfrak F_p\ }
H^i(\bar X,\mathcal O_{\bar X})\longrightarrow0,
\end{equation}

which proves the assertion. Proposition~1.1.10 now gives the following
corollary of Proposition~2.0.

\subsubsection*{Corollary 2.1.}
For $\bar X$ proper over a separably closed field $k$, there is a canonical
isomorphism

\begin{equation}\tag{2.1.1}
H^i(\bar X,\mathcal O_{\bar X})_{\mathrm{ss}}
\xleftarrow{\ \sim\ }
H^i_{\mathrm{et}}(\bar X,\mathbb Z/p\mathbb Z)
\otimes_{\mathbb F_p}k.
\end{equation}
\end{filecontents*}
% Original input: components/166_expose_XXII_section_2_2_through_proposition_2_2_5.tex
\begin{filecontents*}[overwrite,noheader]{sga7ii-cumulative-input-166.tex}
\subsection*{2.2.}
Let now $f:X\longrightarrow\Spec(\mathbb F_q)$ be a proper scheme over
$\mathbb F_q$. Let $k$ be an algebraic closure of $\mathbb F_q$, and put

\begin{equation}\tag{2.2.1}
\bar X=X\times_{\mathbb F_q}k.
\end{equation}

The étale sheaf $R^if_*(\mathbb Z/p\mathbb Z)$ on
$\Spec(\mathbb F_q)$ identifies with
$H^i_{\mathrm{et}}(\bar X,\mathbb Z/p\mathbb Z)$, viewed as a
$\Gal(k/\mathbb F_q)$-module. Let

\begin{equation}\tag{2.2.2}
\mathfrak F_q:H^i(X,\mathcal O_X)\longrightarrow H^i(X,\mathcal O_X)
\end{equation}

denote the linear operation (it is $q$-linear, but the ground field is
$\mathbb F_q$) induced by

\begin{equation}\tag{2.2.3}
\mathfrak F_q:\mathcal O_X\longrightarrow\mathcal O_X,
\qquad \mathfrak F_q(f)=f^q.
\end{equation}

\subsubsection*{Proposition 2.2.4.}
Let $\varphi_q\in\Gal(k/\mathbb F_q)$ be the ``Frobenius'' element
$\varphi_q(\lambda)=\lambda^q$, and let $i\in\mathbb Z$. There is an
identity of characteristic polynomials

\begin{equation}\tag{2.2.4.1}
\det_{\mathbb F_p}\!\left(
1-t\varphi_q^{-1}\mid H^i_{\mathrm{et}}(\bar X,\mathbb Z/p\mathbb Z)
\right)
=
\det_{\mathbb F_q}\!\left(
1-t\mathfrak F_q\mid H^i(X,\mathcal O_X)
\right).
\end{equation}

Here the left-hand side is the characteristic polynomial of the element
$\varphi_q^{-1}$ of $\pi_1(\Spec(\mathbb F_q))$ acting on the Galois module
$H^i_{\mathrm{et}}(\bar X,\mathbb Z/p\mathbb Z)$.

Proposition~2.2.4 is a corollary of the following result.

\subsubsection*{Proposition 2.2.5.}
Let $X$ be proper over $\Spec(\mathbb F_q)$. There is a canonical
isomorphism

\begin{equation}\tag{2.2.6}
H^i(X,\mathcal O_X)_{\mathrm{ss}}\otimes_{\mathbb F_q}k
\simeq
H^i_{\mathrm{et}}(\bar X,\mathbb Z/p\mathbb Z)
\otimes_{\mathbb F_p}k
\end{equation}

which carries the automorphism induced by $\mathfrak F_q$ on the left-hand
side to the automorphism induced by $\varphi_q^{-1}$ on the right-hand side.
The latter may also be interpreted as
$(\mathrm{Fr}_q^X)^*\otimes_{\mathbb F_p}k$, where
$\mathrm{Fr}_q^X:\bar X\longrightarrow\bar X$ is the Frobenius
$k$-endomorphism (cf. SGA~5, XVIII~2.3.4).

\medskip
\noindent\emph{Proof.}
Since

\begin{equation}\tag{2.2.7}
H^i(\bar X,\mathcal O_{\bar X})\simeq
H^i(X,\mathcal O_X)\otimes_{\mathbb F_q}k,
\end{equation}

Proposition~1.0.10 gives

\begin{equation}\tag{2.2.8}
H^i(\bar X,\mathcal O_{\bar X})_{\mathrm{ss}}\simeq
H^i(X,\mathcal O_X)_{\mathrm{ss}}\otimes_{\mathbb F_q}k.
\end{equation}

Thus the isomorphism (2.2.6) is obtained by composing (2.2.8) with
(2.1.1).

In view of (2.2.7), the $q$-linear operation

\begin{equation}\tag{2.2.9}
\mathfrak F_q:H^i(X,\mathcal O_X)\longrightarrow H^i(X,\mathcal O_X)
\end{equation}

has a $q$-linear extension (cf. 1.0.5)

\begin{equation}\tag{2.2.10}
\overline{\mathfrak F}_q:H^i(\bar X,\mathcal O_{\bar X})
\longrightarrow H^i(\bar X,\mathcal O_{\bar X}),
\end{equation}

and also a $k$-linear extension, again denoted by
$\overline{\mathfrak F}_q$,

\begin{equation}\tag{2.2.11}
\overline{\mathfrak F}_q:H^i(\bar X,\mathcal O_{\bar X})
\longrightarrow H^i(\bar X,\mathcal O_{\bar X}).
\end{equation}

Let

\begin{equation}\tag{2.2.12}
\varphi_q:H^i(\bar X,\mathcal O_{\bar X})
\longrightarrow H^i(\bar X,\mathcal O_{\bar X})
\end{equation}

denote the $q$-linear operation whose fixed points are
$H^i(X,\mathcal O_X)$. Under (2.2.7), the operators
$\mathfrak F_q$, $\varphi_q$, and $\overline{\mathfrak F}_q$ may also be
written respectively as
$\mathfrak F_q\otimes\mathrm{id}_k$,
$\mathrm{id}\otimes\varphi_q$, and
$\mathfrak F_q\otimes\varphi_q$. They therefore commute with one another,
and

\begin{equation}\tag{2.2.13}
\overline{\mathfrak F}_q
=\varphi_q\mathfrak F_q
=\mathfrak F_q\circ\varphi_q.
\end{equation}

They all commute with $\mathfrak F_p$, so $\mathfrak F_q$ and $\varphi_q$
act on both
$H^i(\bar X,\mathcal O_{\bar X})^{1-\overline{\mathfrak F}_q}$ and
$H^i(\bar X,\mathcal O_{\bar X})^{1-\mathfrak F_p}$. By (2.2.13),
$\mathfrak F_q$ and $\varphi_q$ are inverse to one another on
$H^i(\bar X,\mathcal O_{\bar X})^{1-\overline{\mathfrak F}_q}$. Hence the
action of $\mathfrak F_q$ on

\[
H^i(\bar X,\mathcal O_{\bar X})_{\mathrm{ss}}
\simeq
H^i(\bar X,\mathcal O_{\bar X})^{1-\overline{\mathfrak F}_q}
\otimes_{\mathbb F_q}k
\qquad\text{(cf. 1.1.11)}
\]

is inverse to the action of $\varphi_q\otimes\mathrm{id}_k$ on

\[
H^i(\bar X,\mathcal O_{\bar X})^{1-\overline{\mathfrak F}_q}
\otimes_{\mathbb F_q}k.
\]

Writing

\begin{equation}\tag{2.2.14}
H^i(\bar X,\mathcal O_{\bar X})^{1-\overline{\mathfrak F}_q}
\simeq
H^i(\bar X,\mathcal O_{\bar X})^{1-\mathfrak F_p}
\otimes_{\mathbb F_p}\mathbb F_q,
\end{equation}

the automorphism $\varphi_q$ of the left-hand side of (2.2.14) becomes,
via (2.2.14), the automorphism induced by $\varphi_q$ on
$H^i(\bar X,\mathcal O_{\bar X})^{1-\mathfrak F_p}
\otimes_{\mathbb F_p}\mathbb F_q$. Consequently the automorphism induced by
$\mathfrak F_q$ on
$H^i(\bar X,\mathcal O_{\bar X})_{\mathrm{ss}}$ becomes, via (2.1.1), the
automorphism induced by $\varphi_q^{-1}$ on
$H^i_{\mathrm{et}}(\bar X,\mathbb Z/p\mathbb Z)
\otimes_{\mathbb F_p}k$. This proves Proposition~2.2.5.

\medskip
\noindent\emph{Source note.}
The final printed sentence of the proof says $\varphi_q$ where the proposition
and the immediately preceding inverse relation require $\varphi_q^{-1}$.
The English follows the mathematically coherent reading and the corrected
French workpass; the diplomatic French retains the printed wording.
\end{filecontents*}
% Original input: components/167_expose_XXII_section_3_opening_through_theorem_3_3.tex
\begin{filecontents*}[overwrite,noheader]{sga7ii-cumulative-input-167.tex}
\section*{3. The congruence formula: statements and elementary equivalences}

I wish to thank Grothendieck and Deligne for explaining to me the variants
(3.1.2), (3.2.1)--(3.2.2), and (3.3.1) of (3.1.1).

\subsection*{3.0.}
Let $X$ be a scheme of finite type over $\mathbb F_q$. The degree of a closed
point $y\in X$ is defined to be the degree of its residue field
$\mathbb F_q(y)/\mathbb F_q$. The zeta function of $X/\mathbb F_q$, as an
element of $\mathbb Z[[t]]$, is defined by the formula (cf. also SGA~5,
XIII~3.1)

\begin{equation}\tag{3.0.1}
Z(t,X/\mathbb F_q)=\prod_y(1-t^{\deg(y)})^{-1}.
\end{equation}

By the ``zeta function modulo $p$'' we shall always mean the image of
$Z(t,X/\mathbb F_q)$ in $\mathbb F_q[[t]]$ under the ``reduction modulo $p$''
map

\[
\mathbb Z[[t]]\longrightarrow\mathbb F_q[[t]].
\]

\subsubsection*{Theorem 3.1.}
Let $X$ be a proper scheme over $\mathbb F_q$. Then

\begin{equation}\tag{3.1.1}
Z(t,X/\mathbb F_q)
\equiv
\prod_{i\geq0}
\det\!\left(1-t\mathfrak F_q\mid H^i(X,\mathcal O_X)\right)^{(-1)^{i+1}}
\pmod p,
\end{equation}

or, equivalently by Proposition~2.2.4,

\begin{equation}\tag{3.1.2}
Z(t,X/\mathbb F_q)
\equiv
\prod_{i\geq0}
\det\!\left(1-t\varphi_q^{-1}\mid
H^i_{\mathrm{et}}(\bar X,\mathbb Z/p\mathbb Z)\right)^{(-1)^{i+1}}
\pmod p.
\end{equation}

\subsubsection*{Corollary 3.2.}
Let $X$ be a proper scheme over $\mathbb F_q$. Then

\begin{equation}\tag{3.2.1}
\sum_{i\geq0}(-1)^i
\operatorname{Tr}\!\left(\mathfrak F_q\mid H^i(X,\mathcal O_X)\right)
\equiv \#X(\mathbb F_q)\pmod p,
\end{equation}

or, equivalently by Proposition~2.2.4,

\begin{equation}\tag{3.2.2}
\sum_{i\geq0}(-1)^i
\operatorname{Tr}\!\left(\varphi_q^{-1}\mid
H^i_{\mathrm{et}}(\bar X,\mathbb Z/p\mathbb Z)\right)
\equiv \#X(\mathbb F_q)\pmod p.
\end{equation}

Indeed, take the coefficients of $t$ in the congruences (3.1.1) and (3.1.2).

\subsubsection*{Theorem 3.3.}
Let $X$ be a proper scheme over $\mathbb F_q$, and denote by
$H^i_{\mathrm{DR}}(X/\mathbb F_q)$ the $\mathbb F_q$-vector spaces
$H^i(X,\Omega^\bullet_{X/\mathbb F_q})$, on which $\mathfrak F_q$ acts. Then

\begin{equation}\tag{3.3.1}
Z(t,X/\mathbb F_q)
\equiv
\prod_{i\geq0}
\det\!\left(1-t\mathfrak F_q\mid
H^i_{\mathrm{DR}}(X/\mathbb F_q)\right)^{(-1)^{i+1}}
\pmod p.
\end{equation}

Indeed, consider the spectral sequence

\begin{equation}\tag{3.3.2}
E_1^{p,q}=H^q(X,\Omega^p_{X/\mathbb F_q})
\Longrightarrow H^{p+q}_{\mathrm{DR}}(X/\mathbb F_q).
\end{equation}

The operator $\mathfrak F_q$ acts on it through the endomorphism
$\mathfrak F_q=(\mathrm{Fr}_q^X)^*$ induced by
$\mathrm{Fr}_q^X:X\longrightarrow X$, and kills $E_1^{p,q}$ for $p\ne0$.
It follows that (3.3.1) is equivalent to (3.1.1).

Let us add a small complement. Since formation of the ``semisimple part''
(cf. 1.0.7) is an exact functor, (3.3.2) gives a spectral sequence

\begin{equation}\tag{3.3.3}
(E_1^{p,q})_{\mathrm{ss}}
\Longrightarrow H^{p+q}_{\mathrm{DR}}(X/\mathbb F_q)_{\mathrm{ss}}.
\end{equation}

Since $(E_1^{p,q})_{\mathrm{ss}}=0$ for $p\ne0$, we obtain for every $i$ an
isomorphism

\begin{equation}\tag{3.3.4}
H^i_{\mathrm{DR}}(X/\mathbb F_q)_{\mathrm{ss}}
\xrightarrow{\ \sim\ }
H^i(X,\mathcal O_X)_{\mathrm{ss}}
\end{equation}

that commutes with $\mathfrak F_q$. Hence

\begin{equation}\tag{3.3.5}
\det\!\left(1-t\mathfrak F_q\mid H^i_{\mathrm{DR}}(X/\mathbb F_q)\right)
=
\det\!\left(1-t\mathfrak F_q\mid H^i(X,\mathcal O_X)\right).
\end{equation}
\end{filecontents*}
% Original input: components/168_expose_XXII_theorem_3_4_through_lemma_3_9.tex
\begin{filecontents*}[overwrite,noheader]{sga7ii-cumulative-input-168.tex}
\subsubsection*{Theorem 3.4.}
Let $X$ be a separated scheme of finite type over $\mathbb F_q$. Denoting by
$H_c^i$ étale cohomology with proper supports (SGA~4, XVII~5.1.9(iii)), one
has

\begin{equation}\tag{3.4.1}
Z(t,X/\mathbb F_q)
\equiv
\prod_{i\geq0}
\det\!\left(1-t\varphi_q^{-1}\mid
H_c^i(\bar X,\mathbb Z/p\mathbb Z)\right)^{(-1)^{i+1}}
\pmod p.
\end{equation}

We begin by proving the equivalence of (3.4.1) and (3.1.1). For the moment,
write

\begin{equation}\tag{3.4.2}
Z(X)\quad\text{in place of }Z(t,X/\mathbb F_q)
\end{equation}

and

\begin{equation}\tag{3.4.3}
Z'(X)\quad\text{in place of }\quad
\prod_i\det\!\left(1-t\varphi_q^{-1}\mid
H_c^i(\bar X,\mathbb Z/p\mathbb Z)\right)^{(-1)^{i+1}}.
\end{equation}

\subsubsection*{Lemma 3.5.}
If $\{U_i\}$ is a finite covering of $X$ by open subschemes, then

\begin{equation}\tag{3.5.1}
Z'(X)=
\prod_{j\geq0}
\prod_{1\leq i_0<\cdots<i_j\leq n}
Z'(U_{i_0}\cap\cdots\cap U_{i_j})^{(-1)^j}.
\end{equation}

\medskip
\noindent\emph{Proof.}
This follows immediately from the localization spectral sequence (SGA~4,
XVII~6.2.10)

\[
E_1^{-a,b}=\bigoplus_{i_0<\cdots<i_b}
H_c^a(\bar U_{i_0}\cap\cdots\cap\bar U_{i_b},\mathbb Z/p\mathbb Z)
\Longrightarrow H_c^{-a+b}(\bar X,\mathbb Z/p\mathbb Z).
\]

It may also be obtained by induction on $n$, using Lemma~3.6 below.

\subsubsection*{Lemma 3.6.}
If $U$ is an open subscheme of $X$ and $Y$ is the closed subscheme $X-U$, then

\begin{equation}\tag{3.6.1}
Z'(X)=Z'(U)Z'(Y).
\end{equation}

\medskip
\noindent\emph{Proof.}
This follows from the exact cohomology sequence

\begin{equation}\tag{3.6.2}
\longrightarrow H_c^i(\bar U,\mathbb Z/p\mathbb Z)
\longrightarrow H_c^i(\bar X,\mathbb Z/p\mathbb Z)
\longrightarrow H_c^i(\bar Y,\mathbb Z/p\mathbb Z)\longrightarrow.
\end{equation}

\subsubsection*{Lemma 3.7 (M\"obius).}
Let $B$ be a finite set and let $\{A_i\}$ be a finite covering by subsets.
Then

\begin{equation}\tag{3.7.1}
\operatorname{card}(B)=
\sum_{j\geq0}(-1)^j
\sum_{i_0<\cdots<i_j}
\operatorname{card}(A_{i_0}\cap\cdots\cap A_{i_j}).
\end{equation}

\medskip
\noindent\emph{Proof.}
There is a Leray spectral sequence for the rational cohomology of the discrete
space $B$ and the open covering $\{A_i\}$,

\[
E_1^{p,q}=\bigoplus_{i_0<\cdots<i_q}
H^p(A_{i_0}\cap\cdots\cap A_{i_q},\mathbb Q)
\Longrightarrow H^{p+q}(B,\mathbb Q),
\]

and one computes the Euler--Poincar\'e characteristic of $B$.

\subsubsection*{Lemma 3.8.}
If $\{U_i\}$ is a finite covering of $X$ by subschemes, then

\begin{equation}\tag{3.8.1}
Z(X)=
\prod_{j\geq0}
\prod_{i_0<\cdots<i_j}
Z(U_{i_0}\cap\cdots\cap U_{i_j})^{(-1)^j}.
\end{equation}

\medskip
\noindent\emph{Proof.}
One has

\begin{equation}\tag{3.8.2}
\operatorname{card}(X(\mathbb F_{q^n}))
=\sum_{r\mid n}r\,\operatorname{card}\{y\in X\mid\deg(y)=r\},
\end{equation}

which gives

\begin{equation}\tag{3.8.3}
Z(X)=\exp\!\left(
\sum_{n\geq1}\operatorname{card}(X(\mathbb F_{q^n}))\frac{t^n}{n}
\right).
\end{equation}

The conclusion follows from (3.7.1).

\subsubsection*{Lemma 3.9.}
Suppose that formula (3.4.1), or equivalently (3.1.2) and hence (3.1.1), holds
when $X$ is a projective hypersurface. Then it holds for every separated
scheme $X$ of finite type.

\medskip
\noindent\emph{Proof.}
By Lemmas~3.8 and 3.6, (3.4.1) holds for the complement
$\mathbb P^n-H$ of a hypersurface. Since

\[
\bigcap_i(\mathbb P^n-H_i)=\mathbb P^n-\bigcup_iH_i
\]

is again the complement of a hypersurface, the open covering
$\{\mathbb P^n-H_i\}$ of $\mathbb P^n-\bigcap_iH_i$ allows us to apply
Lemmas~3.8 and 3.5 and deduce (3.4.1) for
$\mathbb P^n-\bigcap_iH_i$. By Lemmas~3.8 and 3.6, (3.4.1) then holds for
$\bigcap_iH_i$, and hence for every projective $X$. Another application of
Lemmas~3.8 and 3.5 gives (3.4.1) for every quasi-projective $X$, in particular
for every affine $X$. Finally, if $X$ is separated and of finite type, choose
a finite covering by affine open subschemes and apply Lemmas~3.8 and 3.5.

\medskip
\noindent\emph{Source note.}
The printed opening references three nonexistent or malformed variants of
(3.1.1); the corrected French and English cite the actual formulas (3.1.2),
(3.2.1)--(3.2.2), and (3.3.1). In the proof of Theorem~3.3 the print writes
$\mathrm{fr}_p^{X*}$ although the map displayed beside it is
$\mathrm{fr}_q^X$, and twice says that Frobenius kills spectral-sequence terms
for $q\ne0$ rather than for positive differential-form degree $p\ne0$.
Finally, the proof of Lemma~3.9 twice cites nonexistent (3.1.3) where its own
argument is proving (3.4.1). The diplomatic French retains all these printed
forms; corrected French and English use the coherent readings just stated.
\end{filecontents*}
% Original input: components/169_expose_XXII_section_4_0_through_lemma_4_2.tex
\begin{filecontents*}[overwrite,noheader]{sga7ii-cumulative-input-169.tex}
\section*{4. The calculation of the zeta function of a hypersurface by
Dwork's method}

\subsection*{4.0.}
Let

\begin{equation}\tag{4.0.1}
f(X)=f(X_1,\ldots,X_{n+2})\in\mathbb F_q[X]
\end{equation}

be homogeneous of degree $d$, and let its zero locus define the hypersurface
$V\subset\mathbb P^{n+1}$. Denote by $V^*$ the open subset of $V$ on which
$\prod_{i=1}^{n+2}X_i$ is invertible, and by $V^*_{\mathrm{aff}}$ the affine
cone over $V^*$.

Choose a nontrivial additive character

\begin{equation}\tag{4.0.2}
\chi:\mathbb F_{q^\nu}\longrightarrow\mu_p(C),
\end{equation}

where $C$ is a fixed algebraically closed field of characteristic zero (for
example $C=\mathbb C$). Thus

\begin{equation}\tag{4.0.3}
\text{for }x\in\mathbb F_{q^\nu},\qquad
\sum_{y\in\mathbb F_{q^\nu}}\chi(xy)=
\begin{cases}
q^\nu,&x=0,\\
0,&x\ne0.
\end{cases}
\end{equation}

For $x=(x_1,\ldots,x_{n+2})\in(\mathbb F_{q^\nu}^*)^{n+2}$ we therefore
have

\begin{equation}\tag{4.0.4}
1+\sum_{x_0\in\mathbb F_{q^\nu}^*}\chi(x_0f(x))
=\sum_{x_0\in\mathbb F_{q^\nu}}\chi(x_0f(x))
=
\begin{cases}
q^\nu,&f(x)=0,\\
0,&\text{otherwise}.
\end{cases}
\end{equation}

Summing (4.0.4) over $x\in(\mathbb F_{q^\nu}^*)^{n+2}$ gives

\begin{equation}\tag{4.0.5}
\operatorname{card}V^*_{\mathrm{aff}}(\mathbb F_{q^\nu})
=\frac1{q^\nu}
\sum_{x\in(\mathbb F_{q^\nu}^*)^{n+2}}
\left\{1+\sum_{x_0\in\mathbb F_{q^\nu}^*}\chi(x_0f(x))\right\}.
\end{equation}

Since

\begin{equation}\tag{4.0.6}
\operatorname{card}V^*_{\mathrm{aff}}(\mathbb F_{q^\nu})
=(q^\nu-1)\operatorname{card}V^*(\mathbb F_{q^\nu}),
\end{equation}

formula (4.0.5) becomes

\begin{equation}\tag{4.0.7}
(q^\nu-1)\operatorname{card}V^*(\mathbb F_{q^\nu})
=\frac{(q^\nu-1)^{n+2}}{q^\nu}
+\frac1{q^\nu}
\sum_{(x_0,x)\in(\mathbb F_{q^\nu}^*)^{n+3}}\chi(x_0f(x)).
\end{equation}

Write $f(X)$ explicitly as a sum of monomials,

\begin{equation}\tag{4.0.8}
f(X)=\sum_w a_wX^w,
\end{equation}

where $w=(w_1,\ldots,w_{n+2})$, $\sum_iw_i=d$, and
$X^w\mathrel{\mathop:}=X_1^{w_1}\cdots X_{n+2}^{w_{n+2}}$. Then (4.0.7)
may be rewritten as

\begin{equation}\tag{4.0.9}
(q^\nu-1)\operatorname{card}V^*(\mathbb F_{q^\nu})
=\frac{(q^\nu-1)^{n+2}}{q^\nu}
+\frac1{q^\nu}
\sum_{(x_0,x)\in(\mathbb F_{q^\nu}^*)^{n+3}}
\prod_w\chi(a_wx_0X^w).
\end{equation}

\subsection*{4.1.}
We must now make explicit our ``additive character with values in
characteristic zero.'' One would like simply to take

\begin{equation}\tag{4.1.0}
\chi:x\longmapsto
\exp\!\left(2\pi i\,\operatorname{tr}_{\mathbb F_{q^\nu}/\mathbb F_p}(x)\right)
=\xi_p^{\operatorname{tr}_{\mathbb F_{q^\nu}/\mathbb F_p}(x)}.
\end{equation}

Let

\begin{align}
\tag{4.1.1}
\Omega&=\text{an algebraic closure of }\mathbb Q_p,\\
\tag{4.1.2}
\operatorname{ord}:\Omega^*&\longrightarrow\mathbb Q
  &&\text{be the valuation normalized by }\operatorname{ord}(p)=1,\\
\tag{4.1.3}
\xi_p&\in\Omega
  &&\text{be a primitive $p$-th root of unity.}
\end{align}

One checks (cf. [2]) that there exists

\begin{equation}\tag{4.1.4}
\pi\in\mathbb Q_p(\xi_p),\qquad
\operatorname{ord}(\pi)=\frac1{p-1},
\end{equation}

such that

\begin{equation}\tag{4.1.5}
\sum_{n\geq0}\frac{\pi^{p^n}}{p^n}=0.
\end{equation}

(Notice also that the left-hand side converges in $\mathbb Q_p(\xi_p)$.)

Define the Artin--Hasse exponential

\begin{equation}\tag{4.1.6}
E(X)=\exp\!\left(\sum_{n\geq0}\frac{X^{p^n}}{p^n}\right)
\in\mathbb Q_p[[X]],
\end{equation}

which, miraculously, lies in $\mathbb Z_p[[X]]$. Define

\begin{equation}\tag{4.1.7}
\Theta(X)=E(\pi X).
\end{equation}

Its fundamental properties are summarized in the following lemma.

\subsubsection*{Lemma 4.2.}

\begin{equation}\tag{4.2.1}
\Theta(X)=\sum_n B_n X^n,\qquad B_0=1,\quad B_1=\pi,\quad
\operatorname{ord}(B_n)\geq\frac n{p-1}.
\end{equation}

\begin{equation}\tag{4.2.2}
\Theta(1)=\xi_p,\qquad\text{a primitive $p$-th root of unity.}
\end{equation}

Let $t\in\Omega$, and suppose that, for a given integer $a>0$,

\begin{equation}\tag{4.2.3}
t=t^{p^a};
\end{equation}

that is, $t$ is a Teichmüller representative in
$W(\mathbb F_{p^a})\subset\Omega$. Then
$\sum_{j=0}^{a-1}t^{p^j}\in\mathbb Z_p$, so the formal series

\[
(\Theta(X))^{\sum_{j=0}^{a-1}t^{p^j}}
\]

is defined, and one has the identity of formal power series

\begin{equation}\tag{4.2.3.1}
\prod_{j=0}^{a-1}\Theta(t^{p^j}X)
=(\Theta(X))^{\sum_{j=0}^{a-1}t^{p^j}}.
\end{equation}
\end{filecontents*}
% Original input: components/170_expose_XXII_proof_of_lemma_4_2_and_corollary_4_3.tex
\begin{filecontents*}[overwrite,noheader]{sga7ii-cumulative-input-170.tex}
\medskip
\noindent\emph{Proof.}
Since $E(X)\in\mathbb Z_p[[X]]$, (4.2.1) follows from the definition. For
(4.2.2), we have

\begin{equation}\tag{4.2.4}
\log(\Theta(1))=0
\qquad\text{by the choice of $\pi$ in (4.1.5),}
\end{equation}

and

\begin{equation}\tag{4.2.5}
\operatorname{ord}(\Theta(1)-1)=\frac1{p-1}
\qquad\text{by (4.2.1).}
\end{equation}

By (4.2.4), there is an integer $n\geq0$ such that
$\Theta(1)^{p^n}=1$ (this is the characterization of the kernel of the
logarithm), and (4.2.5) implies that $n=1$.

As for (4.2.3.1), after taking logarithms it follows by an immediate
calculation, once one knows that
$\sum_{j=0}^{a-1}t^{p^j}\in\mathbb Z_p$. But the $t^{p^j}$,
$0\leq j\leq a-1$, are precisely the conjugates over $\mathbb Z_p$ of
$t\in W(\mathbb F_{p^a})$; hence

\[
\sum_{j=0}^{a-1}t^{p^j}
=\operatorname{Tr}_{W(\mathbb F_{p^a})/\mathbb Z_p}(t)\in\mathbb Z_p.
\]

Evaluating the two sides of (4.2.3.1) at $X=1$ gives the following result.

\subsubsection*{Corollary 4.3.}
Let $t\in\Omega$ satisfy $t^{p^a}=t$, and denote its residue class by
$\bar t\in\mathbb F_{p^a}$. Then

\begin{equation}\tag{4.3.1}
\prod_{j=0}^{a-1}\Theta(t^{p^j})
=\xi_p^{\operatorname{tr}_{\mathbb F_{p^a}/\mathbb F_p}(\bar t)}.
\end{equation}

\medskip
\noindent\emph{Source note.}
The print has three mutually exposed notation slips in Lemma~4.2: it writes
$A_1=\pi$ in the coefficient expansion in (4.2.1), places the Teichmüller
representative in $W(\mathbb F_{q^a})$ although both $t^{p^a}=t$ and
Corollary~4.3 use $\mathbb F_{p^a}$, and prints $\sum t^j$ in the exponent
of (4.2.3.1) although the product and the following proof use
$\sum t^{p^j}$. The diplomatic French preserves those forms; corrected
French and English use $B_1$, $\mathbb F_{p^a}$, and $t^{p^j}$.
\end{filecontents*}
% Original input: components/171_expose_XXII_section_4_4.tex
\begin{filecontents*}[overwrite,noheader]{sga7ii-cumulative-input-171.tex}
\subsection*{4.4.}
Recall that the residue field of the ring of integers of
$\mathbb Q_p(\xi_p)$ is $\mathbb F_p$. Let

\begin{equation}\tag{4.4.0}
K=\text{the unramified extension of $\mathbb Q_p(\xi_p)$ with residue
field $\mathbb F_q$},
\end{equation}

and let

\begin{equation}\tag{4.4.1}
\tau\in\operatorname{Gal}(K/\mathbb Q_p(\xi_p))
\end{equation}

be the ``Frobenius'', that is, the automorphism reducing to the Frobenius
automorphism of $\mathbb F_q/\mathbb F_p$, $x\longmapsto x^p$.

Let

\begin{equation}\tag{4.4.2}
F(X)=\sum_w A_wX^w\in K[X_1,\ldots,X_{n+2}]
\end{equation}

be the lifting of $f(X)$ whose coefficients satisfy $A_w^q=A_w$.
Putting (4.3.1) and (4.0.9) together, and taking into account the
definition (4.1.0) of $\chi$, one obtains

\begin{equation}\tag{4.4.3}
\begin{aligned}
&(q^\nu-1)\operatorname{card}V^*(\mathbb F_{q^\nu})
-\frac{(q^\nu-1)^{n+2}}{q^\nu}\\
&\quad=\frac1{q^\nu}
\sum_{(X_0,X)\in(\mu_{q^\nu-1}(\Omega))^{n+3}}
\prod_w\prod_{j=0}^{\operatorname{ord}(q^\nu)-1}
\Theta\!\left(A_w^{p^j}X_0^{p^j}X^{p^jw}\right).
\end{aligned}
\end{equation}

Define

\[
H(X)=H(X_0,\ldots,X_{n+2})\in K[[X_0,\ldots,X_{n+2}]]
\]

by

\begin{equation}\tag{4.4.4}
H(X)=\prod_w
\Theta\!\left(A_wX_0X_1^{w_1}\cdots X_{n+2}^{w_{n+2}}\right).
\end{equation}

We let $\operatorname{Gal}(K/\mathbb Q_p(\xi_p))$ act on
$K[[X_0,\ldots,X_{n+2}]]$ through its action on the coefficients; the
componentwise action of
$\sigma\in\operatorname{Gal}(K/\mathbb Q_p(\xi_p))$ will be denoted

\begin{equation}\tag{4.4.5}
G(X)\longmapsto G^\sigma(X).
\end{equation}

In particular,

\[
H^{\tau^j}(X)=\prod_w
\Theta\!\left(A_w^{\tau^j}X_0X_1^{w_1}\cdots
X_{n+2}^{w_{n+2}}\right).
\]

Since $A_w^{\tau^j}=A_w^{p^j}$, formula (4.4.3) may therefore be rewritten
as

\begin{equation}\tag{4.4.6}
\begin{aligned}
(q^\nu-1)\operatorname{card}V^*(\mathbb F_{q^\nu})
&=\frac{(q^\nu-1)^{n+2}}{q^\nu}\\
&\quad+\frac1{q^\nu}
\sum_{(X_0,X)\in(\mu_{q^\nu-1}(\Omega))^{n+3}}
\prod_{j=0}^{\operatorname{ord}(q^\nu)-1}H^{\tau^j}(X^{p^j}).
\end{aligned}
\end{equation}
\end{filecontents*}
% Original input: components/172_expose_XXII_section_4_5_setup_and_operators.tex
\begin{filecontents*}[overwrite,noheader]{sga7ii-cumulative-input-172.tex}
\subsection*{4.5.}
We shall now express the sum in (4.4.6) as a trace.

For each $b\in\mathbb Q$, define a $p$-adic Banach $K$-algebra

\begin{equation}\tag{4.5.1}
L(b)=\text{the subring of $K[[X_0,\ldots,X_{n+2}]]$ consisting of the
series}
\end{equation}

which can be written

\begin{equation}\tag{4.5.2}
\sum_{dU_0=\sum_{i\geq1}U_i}
B_UX_0^{U_0}X_1^{U_1}\cdots X_{n+2}^{U_{n+2}},
\qquad U=(U_0,\ldots,U_{n+2}),
\end{equation}

and which satisfy

\begin{equation}\tag{4.5.3}
\inf_U\bigl(\operatorname{ord}(B_U)-bU_0\bigr)>-\infty.
\end{equation}

It is a $p$-adic Banach algebra for the order function (the negative of the
logarithm of the norm)

\begin{equation}\tag{4.5.4}
\inf_U\bigl(\operatorname{ord}(B_U)-bU_0\bigr)
=\operatorname{ord}\!\left(\sum_U B_UX^U\right).
\end{equation}

It follows from (4.4.4) and (4.2.1) that

\[
H(X)\in L\!\left(\frac1{p-1}\right).
\]

Let us introduce some operators. For every integer $k>0$, define a
continuous operator between Banach spaces over $K$ by

\begin{equation}\tag{4.5.5}
\left\{
\begin{aligned}
\psi_k:L(b)&\longrightarrow L(kb),\\
\psi_k\!\left(\sum_U B_UX^U\right)&=\sum_U B_{kU}X^U.
\end{aligned}
\right.
\end{equation}

For every pair $(b,b')$ with $b'<b$, there is a ``restriction'' operator,
which is completely continuous (that is, a uniform limit of finite-rank
operators):

\begin{equation}\tag{4.5.6}
\left\{
\begin{aligned}
L(b)&\xrightarrow{\ \text{restriction}\ }L(b'),\\
\sum_U B_UX^U&\longmapsto\sum_U B_UX^U.
\end{aligned}
\right.
\end{equation}

Each element $G$ of the algebra $L(b)$ gives the continuous endomorphism
of multiplication by $G$,

\begin{equation}\tag{4.5.7}
\left\{
\begin{aligned}
G:L(b)&\longrightarrow L(b),\\
\eta&\longmapsto\eta G.
\end{aligned}
\right.
\end{equation}

These operators are related by the following lemma, whose proof is found in
[2] (see also [4]).
\end{filecontents*}
% Original input: components/173_expose_XXII_dwork_trace_lemma_through_formula_4_5_12.tex
\begin{filecontents*}[overwrite,noheader]{sga7ii-cumulative-input-173.tex}
\subsubsection*{Lemma 4.5.8 (Dwork trace lemma).}
Let $b\in\mathbb Q$, $b>0$, let $k\in\mathbb Z$, $k\geq2$, and let
$G\in L(b/k)$. Then the composite

\[
L(b)\xrightarrow{\ \text{restriction}\ }L(b/k)
\xrightarrow{\ \times G\ }L(b/k)
\xrightarrow{\ \psi_k\ }L(b),
\]

denoted by $\psi_k\circ G$, is completely continuous, and its trace
(cf. [4]) is given by

\begin{equation}\tag{4.5.9}
\operatorname{tr}(\psi_k\circ G)
=\frac1{(k-1)^{n+2}}
\sum_{X\in(\mu_{k-1}(\Omega))^{n+3}}G(X).
\end{equation}

Apply this lemma with

\[
\left\{
\begin{aligned}
k&=q^\nu,\\
G(X)&=\prod_{j=0}^{\operatorname{ord}(q^\nu)-1}
H^{\tau^j}(X^{p^j})
\in L\!\left(\frac{p}{(p-1)q^\nu}\right),\\
b&=\frac1{p-1}.
\end{aligned}
\right.
\]

Formula (4.4.6) becomes

\begin{equation}\tag{4.5.10}
\begin{aligned}
(q^\nu-1)\operatorname{card}V^*(\mathbb F_{q^\nu})
&=\frac{(q^\nu-1)^{n+2}}{q^\nu}\\
&\quad+\frac{(q^\nu-1)^{n+3}}{q^\nu}
\operatorname{tr}\!\left(
\left.\psi_{q^\nu}\circ
\prod_{j=0}^{\operatorname{ord}(q^\nu)-1}H^{\tau^j}(X^{p^j})
\right|L\!\left(\frac1{p-1}\right)\right).
\end{aligned}
\end{equation}

Let $L^0(b)$ be the ideal of the elements of $L(b)$ which ``vanish at
zero'', that is, which have no constant term. It is stable under
$\psi_k\circ G$, and evidently

\begin{equation}\tag{4.5.11}
\operatorname{tr}(\psi_k\circ G)
=G(0)+\operatorname{tr}\!\left(\left.\psi_k\circ G\right|L^0(b)\right).
\end{equation}

Since $H(0)=1$, formula (4.5.10) may be rewritten as

\begin{equation}\tag{4.5.12}
\begin{aligned}
\operatorname{card}V^*(\mathbb F_{q^\nu})
&=(q^\nu-1)^{n+1}\\
&\quad+\frac{(q^\nu-1)^{n+2}}{q^\nu}
\operatorname{tr}\!\left(
\left.\psi_{q^\nu}\circ
\prod_{j=0}^{\operatorname{ord}(q^\nu)-1}H^{\tau^j}(X^{p^j})
\right|L^0\!\left(\frac1{p-1}\right)\right).
\end{aligned}
\end{equation}
\end{filecontents*}
% Original input: components/174_expose_XXII_alpha_and_support_decomposition.tex
\begin{filecontents*}[overwrite,noheader]{sga7ii-cumulative-input-174.tex}
Consider now the $\tau^{-1}$-linear operator

\begin{equation}\tag{4.5.13}
\tau^{-1}\circ\psi_p\circ H(X):L(b)\longrightarrow L(b).
\end{equation}

Its iterates are

\begin{equation}\tag{4.5.14}
(\tau^{-1}\circ\psi_p\circ H(X))^a
=\tau^{-a}\circ\psi_{p^a}\circ
H(X)H^\tau(X^p)\cdots H^{\tau^{a-1}}(X^{p^{a-1}}),
\end{equation}

and, if $\operatorname{ord}(q)$ divides $a$, then $\tau^{-a}=\operatorname{id}$.
Define the $\tau^{-1}$-linear operation

\begin{equation}\tag{4.5.15}
\left\{
\begin{aligned}
\alpha:L^0\!\left(\frac1{p-1}\right)
&\longrightarrow L^0\!\left(\frac1{p-1}\right),\\
\alpha&=\frac1p\,\tau^{-1}\circ\psi_p\circ H(X).
\end{aligned}
\right.
\end{equation}

Then (4.5.12) takes the form

\begin{equation}\tag{4.5.16}
\operatorname{card}V^*(\mathbb F_{q^\nu})
=(q^\nu-1)^{n+1}
+(q^\nu-1)^{n+2}
\operatorname{tr}\!\left(
\left.\alpha^{\operatorname{ord}(q^\nu)}
\right|L^0\!\left(\frac1{p-1}\right)\right).
\end{equation}

Put

\begin{equation}\tag{4.5.17}
S=\{1,\ldots,n+2\}.
\end{equation}

For every nonempty subset $A\subset S$, define the subspace

\begin{equation}\tag{4.5.18}
\begin{aligned}
L^A(b)=\{\,&\text{series $\sum_U B_UX^U\in L(b)$ such that, for every
monomial $X^U$ with $B_U\ne0$,}\\
&\text{the indices $i$ for which $U_i>0$ are exactly $0$ and the
$j\in A$}\,\}.
\end{aligned}
\end{equation}

Then

\begin{equation}\tag{4.5.19}
L^0(b)=\bigoplus_{\substack{A\subset S\\A\ne\varnothing}}L^A(b).
\end{equation}

\medskip
\noindent\emph{Source note.}
The print drops $\nu$ from the factor $q^\nu-1$ in (4.5.10) and from the
denominator $q^\nu$ in (4.5.12), although both formulas are direct rewrites
of (4.4.6). It also transposes the formula number (4.5.14) as (5.4.14).
The diplomatic French preserves all three printed forms; corrected French
and English use the algebraically forced readings displayed above.
\end{filecontents*}
% Original input: components/175_expose_XXII_support_quotients_and_point_counts.tex
\begin{filecontents*}[overwrite,noheader]{sga7ii-cumulative-input-175.tex}
\par\medskip
And, for every nonempty subset $B\subset S$, one has

\begin{equation}\tag{4.5.20}
\bigoplus_{A\supset B}L^A(b)
=L(b)\cap\left(\prod_{i\in B}X_i\right)K[[X]].
\end{equation}

Clearly every subspace $\bigoplus_{A\supset B}L^A(b)$ is stable under
$\alpha$. It follows that the same is true of
$\bigoplus_{A\not\supset B}L^A(b)$, which is a sum of subspaces of the
preceding type. Define

\begin{equation}\tag{4.5.21}
\alpha_B=\text{the $\tau^{-1}$-linear endomorphism induced by $\alpha$ on}
\quad
\frac{\displaystyle\bigoplus_{A\supset B}
L^A\!\left(\frac1{p-1}\right)}
{\displaystyle\bigoplus_{\substack{A\supset B\\A\ne B}}
L^A\!\left(\frac1{p-1}\right)}
\simeq L^B\!\left(\frac1{p-1}\right).
\end{equation}

Explicitly, $\alpha_B$ is given by

\begin{equation}\tag{4.5.22}
\alpha_B=\frac{\tau^{-1}}p\,\psi_p\,H_B(X),
\end{equation}

an operation under which $L^B$ is stable, where $H_B(X)$ is obtained from
$H(X)$ by substituting $X_i=0$ for $i\notin\{0\}\cup B$.

Rewriting (4.5.16) once more gives

\begin{equation}\tag{4.5.23}
\operatorname{card}V^*(\mathbb F_{q^\nu})
=(q^\nu-1)^{n+1}
+(q^\nu-1)^{n+2}
\sum_{\substack{A\subset S\\A\ne\varnothing}}
\operatorname{tr}\!\left((\alpha_A)^{\operatorname{ord}(q^\nu)}\right).
\end{equation}

For $A\subset S$, $A\ne\varnothing$, define the hypersurface

\begin{equation}\tag{4.5.24}
V_A\subset\mathbb P^{\operatorname{card}(A)-1}
\simeq
\mathbb P^{n+1}\cap
\{X_i=0\text{ for every }i\notin A\}
\end{equation}

by $f_A(X)$, the polynomial obtained from $f(X)$ by substituting $X_i=0$
for $i\notin A$.

By compatibility of the formation of the series $H$ from $f(X)$ with
substitution of $0$ for some of the variables, (4.5.23) gives, for every
$A\subset S$, $A\ne\varnothing$,

\begin{equation}\tag{4.5.25}
\operatorname{card}V_A^*(\mathbb F_{q^\nu})
=(q^\nu-1)^{\operatorname{card}(A)-1}
+(q^\nu-1)^{\operatorname{card}(A)}
\sum_{B\subset A}
\operatorname{tr}\!\left((\alpha_B)^{\operatorname{ord}(q^\nu)}\right).
\end{equation}

On the other hand,

\begin{equation}\tag{4.5.26}
\operatorname{card}V(\mathbb F_{q^\nu})
=\sum_{\substack{A\subset S\\A\ne\varnothing}}
\operatorname{card}V_A^*(\mathbb F_{q^\nu}).
\end{equation}

After an easy summation, (4.5.25) therefore gives

\begin{equation}\tag{4.5.27}
\begin{aligned}
\operatorname{card}V(\mathbb F_{q^\nu})
={}&1+q^\nu+q^{2\nu}+\cdots
+q^{\nu(\operatorname{card}(S)-1)}\\
&+\sum_{\substack{B\subset S\\B\ne\varnothing}}
(q^\nu-1)^{\operatorname{card}(B)}
q^{\nu\operatorname{card}(S-B)}
\operatorname{tr}\!\left((\alpha_B)^{\operatorname{ord}(q^\nu)}\right).
\end{aligned}
\end{equation}

Since $\alpha$, and hence the $\alpha_B$, are $\tau^{-1}$-linear, the
operators $\alpha_B^{\operatorname{ord}(q)}$ are $K$-linear. Define the
characteristic series

\[
\Delta_B(t)\in1+tK[[t]]
\]

by

\begin{equation}\tag{4.5.28}
\Delta_B(t)=\det\!\left(
1-t\alpha_B^{\operatorname{ord}(q)}
\mathrel{\Big|}L^B\!\left(\frac1{p-1}\right)\right).
\end{equation}

Introduce the following endomorphisms of $1+tK[[t]]$:

\begin{equation}\tag{4.5.29}
\varphi:g(t)\longmapsto g^\varphi(t)=g(qt)
\end{equation}

and

\begin{equation}\tag{4.5.30}
\delta=1-\varphi:g(t)\longmapsto
g^\delta(t)=\frac{g(t)}{g(qt)}.
\end{equation}
\end{filecontents*}
% Original input: components/176_expose_XXII_zeta_formula_and_integrality_lemma.tex
\begin{filecontents*}[overwrite,noheader]{sga7ii-cumulative-input-176.tex}
Taking account of the identities

\begin{equation}\tag{4.5.31}
Z(t,V/\mathbb F_q)
=\exp\!\left(
\sum_{\nu\geq1}
\operatorname{card}V(\mathbb F_{q^\nu})\frac{t^\nu}{\nu}
\right)
\end{equation}

and

\begin{equation}\tag{4.5.32}
\det\!\left(1-t\alpha_B^{\operatorname{ord}(q)}\right)
=\exp\!\left(
-\sum_{\nu\geq1}
\operatorname{tr}\!\left((\alpha_B)^{\nu\operatorname{ord}(q)}\right)
\frac{t^\nu}{\nu}
\right),
\end{equation}

formula (4.5.27) can be rewritten as

\begin{equation}\tag{4.5.33}
Z(t,V/\mathbb F_q)\prod_{i=0}^{n+1}(1-q^it)
=\prod_{j=1}^{n+2}
\left(
\prod_{\substack{B\subset S\\\operatorname{card}(B)=j}}
\Delta_B(t)
\right)^{-(-\delta)^j\varphi^{n+2-j}}.
\end{equation}

\medskip
\noindent\emph{Source note.}
The print numbers (4.5.23) as (4.5.2.3), abbreviates
$\operatorname{card}(A)$ as $\operatorname{car}(A)$ in the first exponent
of (4.5.25), and reuses the number (4.5.23) for the final zeta identity.
The corrected French and English layers use the sequential numbers and the
full cardinality notation displayed above; the diplomatic French preserves
the printed forms.

\subsection*{4.6.}
We now study the integrality of the series $\Delta_B$. Put

\begin{equation}\tag{4.6.1}
\mathcal O=\text{the ring of integers of $K$},
\end{equation}

and define a filtration on $1+t\mathcal O[[t]]$ by

\begin{equation}\tag{4.6.2}
F^i(1+t\mathcal O[[t]])
=1+q^it\mathcal O[[q^it]]
=\varphi^i(1+t\mathcal O[[t]]).
\end{equation}

One has

\begin{align}
\varphi(F^i)&\subset F^{i+1}, \tag{4.6.3}\\
\delta(F^i)&\subset F^i, \tag{4.6.4}\\
\delta&\text{ induces the identity on }F^i/F^{i+1}. \tag{4.6.5}
\end{align}

\subsubsection*{Lemma 4.6.6.}
One has $\Delta_B(t)\in1+t\mathcal O[[t]]$; moreover, if an integer $j$
satisfies

\begin{equation}\tag{4.6.7}
jd<\operatorname{card}(B),
\end{equation}

then

\begin{equation}\tag{4.6.8}
\Delta_B(t)\in F^j(1+t\mathcal O[[t]]).
\end{equation}

\end{filecontents*}
% Original input: components/177_expose_XXII_proof_of_lemma_4_6_6.tex
\begin{filecontents*}[overwrite,noheader]{sga7ii-cumulative-input-177.tex}
\noindent\emph{Proof.}
The Banach spaces
$L^B\!\left(\frac1{p-1}\right)$ are all of type $B(I)$ in the sense of
Serre~[7]; that is, they admit orthonormal bases, namely the monomials

\begin{equation}\tag{4.6.9}
\pi^{U_0}X_0^{U_0}X_1^{U_1}\cdots X_{n+2}^{U_{n+2}}
\end{equation}

\noindent
which belong to them.  To say that this is an orthonormal basis means, by
definition, that

\begin{equation}\tag{4.6.10}
\operatorname{ord}\!\left(
\sum_U C_U\pi^{U_0}X_0^{U_0}\cdots X_{n+2}^{U_{n+2}}
\right)
=\min_U\operatorname{ord}(C_U).
\end{equation}

Let

\begin{equation}\tag{4.6.11}
R_B=\text{the $\mathcal O$-submodule of
$L^B\!\left(\frac1{p-1}\right)$ formed by the elements of order $\geq0$}.
\end{equation}

To prove the lemma, it is enough to show

\begin{equation}\tag{4.6.12}
jd<\operatorname{card}(B)
\quad\Longrightarrow\quad
(\alpha_B)^{\operatorname{ord}(q)}R_B\subset q^jR_B,
\end{equation}

or, a fortiori, to show

\begin{equation}\tag{4.6.13}
jd<\operatorname{card}(B)
\quad\Longrightarrow\quad
\alpha_B(R_B)\subset p^jR_B.
\end{equation}

By (4.5.22),

\begin{equation}\tag{4.6.14}
\alpha_B=\frac1p\,\tau^{-1}\circ\psi_p\circ H_B(X).
\end{equation}

Since

\begin{equation}\tag{4.6.15}
H(X)\in L\!\left(\frac1{p-1}\right),
\qquad \operatorname{ord}(H(X))\geq0,
\end{equation}

we have

\begin{equation}\tag{4.6.16}
H_B(X)\,R_B\subset R_B
\end{equation}

and, plainly,

\begin{equation}\tag{4.6.17}
\tau^{-1}(R_B)=R_B.
\end{equation}

Consequently, (4.6.13) follows from

\begin{equation}\tag{4.6.18}
jd<\operatorname{card}(B)
\quad\Longrightarrow\quad
\psi_p(R_B)\subset p^{j+1}R_B.
\end{equation}

We verify (4.6.18) monomial by monomial.  We have

\begin{equation}\tag{4.6.19}
\begin{cases}
\psi_p\!\left(\pi^{U_0}X_0^{U_0}\cdots X_{n+2}^{U_{n+2}}\right)=0,
&\text{unless }p\mid U_i\text{ for every }i,\\[3pt]
\psi_p\!\left(
\pi^{pU_0}X_0^{pU_0}\cdots X_{n+2}^{pU_{n+2}}
\right)
=\pi^{(p-1)U_0}
\underbrace{\pi^{U_0}X_0^{U_0}\cdots X_{n+2}^{U_{n+2}}}_{\in R_B}.
\end{cases}
\end{equation}

The conditions for
$\pi^{U_0}X_0^{U_0}\cdots X_{n+2}^{U_{n+2}}$ to belong to
$L^B\!\left(\frac1{p-1}\right)$ are

\begin{equation}\tag{4.6.20}
\begin{cases}
dU_0=\displaystyle\sum_{i\geq1}U_i,\\
U_i>0\Longleftrightarrow i\in B,\qquad i\geq1.
\end{cases}
\end{equation}

In particular,

\begin{equation}\tag{4.6.21}
dU_0\geq\operatorname{card}(B).
\end{equation}

Hence

\begin{equation}\tag{4.6.22}
\operatorname{ord}\!\left(\pi^{(p-1)U_0}\right)
=U_0\geq\frac{\operatorname{card}(B)}d>j,
\end{equation}

and, since $U_0\in\mathbb Z$,

\begin{equation}\tag{4.6.23}
\operatorname{ord}\!\left(\pi^{(p-1)U_0}\right)
=U_0\geq j+1.
\end{equation}

Together with (4.6.19), this proves (4.6.18).

\end{filecontents*}
% Original input: components/178_expose_XXII_ax_theorem_and_proof.tex
\begin{filecontents*}[overwrite,noheader]{sga7ii-cumulative-input-178.tex}
\section*{5. A theorem of Ax [1] (see also [4], Section 7)}

\subsubsection*{Theorem 5.0.}
Let $V$ be a hypersurface in $\mathbb P^{n+1}$ of degree $d>0$, defined
over $\mathbb F_q$, and let $j$ be an integer.  If

\begin{equation}\tag{5.0.1}
jd<n+2,
\end{equation}

then

\begin{equation}\tag{5.0.2}
Z(t,V/\mathbb F_q)\prod_{i=0}^{n+1}(1-q^it)
\in1+q^jt\mathbb Z[[q^jt]],
\end{equation}

and the reciprocal of every zero or pole of

\begin{equation}\tag{5.0.3}
Z(t,V/\mathbb F_q)\prod_{i=0}^{n+1}(1-q^it)
\end{equation}

\noindent
is of the form $q^j$ times an algebraic integer.

\subsubsection*{Theorem 5.0 bis.}
Let $V_{\mathrm{aff}}$ be a hypersurface in $\mathbb A^{n+2}$ defined by a
homogeneous form of degree $d$, over $\mathbb F_q$.  If

\begin{equation}\tag{5.0.4}
jd<n+2,
\end{equation}

then

\begin{equation}\tag{5.0.5}
Z(t,V_{\mathrm{aff}}/\mathbb F_q)
\in1+q^jt\mathbb Z[[q^jt]].
\end{equation}

In particular,

\begin{equation}\tag{5.0.6}
\operatorname{card}\bigl(V_{\mathrm{aff}}(\mathbb F_q)\bigr)
\equiv0\pmod{q^j},
\end{equation}

and consequently

\begin{equation}\tag{5.0.6 bis}
\operatorname{card}\bigl(V(\mathbb F_q)\bigr)
\geq\frac{q^j-1}{q-1}>0,
\end{equation}

\noindent
which improves the Chevalley--Warning theorem~[8].

\medskip
\noindent\emph{Source note.}
The print numbers the hypothesis in Theorem~5.0 bis as (5.0.7), duplicating
the number subsequently used for the point-count identity.  The corrected
French and English layers use the sequential number (5.0.4); the diplomatic
French preserves the printed number.

\medskip
\noindent\emph{Proof.}
The implication (5.0.5) $\Rightarrow$ (5.0.6) is clear on taking the
coefficient of $t$ in the zeta function.  Assertions (5.0.4) and (5.0.2)
are equivalent, since

\begin{equation}\tag{5.0.7}
\operatorname{card}V_{\mathrm{aff}}(\mathbb F_{q^\nu})
=1+(q^\nu-1)\operatorname{card}V(\mathbb F_{q^\nu}),
\end{equation}

whence

\begin{equation}\tag{5.0.8}
Z(t,V_{\mathrm{aff}}/\mathbb F_q)
=\frac{Z(t,V/\mathbb F_q)}
{Z(qt,V/\mathbb F_q)(1-t)}
\end{equation}

and

\begin{equation}\tag{5.0.9}
Z(t,V/\mathbb F_q)
=\prod_{i\geq0}
\left[(1-q^it)Z(q^it,V_{\mathrm{aff}}/\mathbb F_q)\right].
\end{equation}

\noindent
Assertions (5.0.2) and (5.0.3) are likewise equivalent, by the theorem of
Fatou~[3] and the rationality of the zeta function.  Since

\begin{equation}\tag{5.0.10}
Z(t,V/\mathbb F_q)\in\mathbb Z[[t]],
\end{equation}

it is enough to prove

\begin{equation}\tag{5.0.11}
Z(t,V/\mathbb F_q)\prod_{i=0}^{n+1}(1-q^it)
\in F^j(1+t\mathcal O[[t]]).
\end{equation}

In view of (4.5.33) and (4.6.4)--(4.6.5), it is enough to show, for
$\varnothing\ne B\subset\{1,\ldots,n+2\}$, that

\begin{equation}\tag{5.0.12}
\Delta_B(t)\in F^{j+\operatorname{card}(B)-(n+2)}.
\end{equation}

We may finish the proof by applying Lemma~4.6.6 once we have proved

\begin{equation}\tag{5.0.13}
jd<n+2
\quad\Longrightarrow\quad
\bigl(j+\operatorname{card}(B)-(n+2)\bigr)d
<\operatorname{card}(B).
\end{equation}

This may be rewritten as

\begin{equation}\tag{5.0.14}
jd<n+2
\quad\Longrightarrow\quad
(j-(n+2))d<\operatorname{card}(B)(1-d).
\end{equation}

But

\[
\inf_{\varnothing\ne B\subset\{1,\ldots,n+2\}}
\operatorname{card}(B)(1-d)
=(n+2)(1-d),
\]

so it is enough to prove

\begin{equation}\tag{5.0.15}
jd<n+2
\quad\Longrightarrow\quad
(j-(n+2))d<(n+2)(1-d),
\end{equation}

which is true.

\medskip
\noindent\emph{Source note.}
In (5.0.9) the print has the impossible factor $(1-q^if)$ rather than
$(1-q^it)$, and the last paragraph cites Lemma~4.4.6 rather than the
integrality lemma 4.6.6 used there.  Both are corrected above; the diplomatic
French retains the printed forms.

\end{filecontents*}
% Original input: components/179_expose_XXII_hodge_remark_conjecture_and_final_section.tex
\begin{filecontents*}[overwrite,noheader]{sga7ii-cumulative-input-179.tex}
\subsubsection*{Remark 5.1.}
Condition

\begin{equation}\tag{5.1.1}
jd<n+2
\end{equation}

admits a cohomological interpretation.  Let $V$ be a smooth hypersurface of
degree $d$ and dimension $n$.  As in XI, 2.1, put, for $a+b=n$,

\begin{equation}\tag{5.1.2}
h_0^{a,b}
=\dim H^b(V,\Omega_V^a)-\delta_{a,b}
=\dim E^{a,b}(V),
\end{equation}

where $E^{a,b}(V)$ is the bidegree-$(a,b)$ component of the vanishing Hodge
cohomology (compare XVII, 5.4.4)

\begin{equation}\tag{5.1.3}
E^{a,b}(V)
=\operatorname{Coker}\!\left(
H^b(\mathbb P^{n+1},\Omega_{\mathbb P^{n+1}}^a)
\longrightarrow H^b(V,\Omega_V^a)\right).
\end{equation}

Then, by XI, 2.8,

\begin{equation}\tag{5.1.4}
jd<n+2
\quad\Longleftrightarrow\quad
h_0^{a,n-a}=0\quad\text{for }a<j\text{ and }a>n-j.
\end{equation}

In other words, the condition says that the Hodge level (XI, 2.7) of the
vanishing de Rham cohomology of $V$,

\[
E_{\mathrm{DR}}^n(V)
=\operatorname{Coker}\!\left(
H_{\mathrm{Hdg}}^n(\mathbb P^{n+1})
\longrightarrow H_{\mathrm{Hdg}}^n(V)\right),
\]

is at most $n-2j$, or equivalently that its coniveau is at least $j$.
On the other hand, the conclusion also says that the eigenvalues of the
Frobenius $\mathfrak F_q$ acting on the vanishing $\ell$-adic cohomology
($\ell\ne p$) are algebraic integers divisible by $q^j$ (XXI, A.2).
Thus Ax's result 5.0 supports the following conjecture, suggested by the
philosophy of the level of motives and relating the classical Hodge and
Tate conjectures~[5].

\subsubsection*{Conjecture 5.2.}
Let $f:X\to S$ be a proper smooth morphism, with $S$ integral, of finite
type over $\operatorname{Spec}(\mathbb Z)$, and dominant over
$\operatorname{Spec}(\mathbb Z)$; let $i,j$ be integers.  Suppose that the
Hodge coniveau of $H^i$ of the generic fiber of $f$ is at least $j$.  Then,
for every point $s\in S$, the coniveau (XX, 2.0.6) of
$H^i(X_s,\mathbb Q_\ell)$ is at least $j$.  Hence (XX, 2.2), if $s$ is a
closed point of $S$, the eigenvalues of Frobenius acting on
$H^i(X_s,\mathbb Q_\ell)$ are algebraic integers divisible by $q^j$, where
$q=\operatorname{card}(k(s))$.

\section*{6. End of the proof of the congruence formula}

Taking account of (4.6.6) and (4.6.3)--(4.6.4), formula (4.5.33) gives

\begin{equation}\tag{6.0.1}
Z(t,V/\mathbb F_q)(1-t)
\equiv\Delta_S(t)^{(-1)^{n+1}}
\pmod{1+qt\mathcal O[[t]]}.
\end{equation}

We now calculate $\Delta_S(t)$ modulo $\pi\mathcal O[[t]]$.  Define, as in
(4.6.11),

\begin{equation}\tag{6.0.2}
W_S=R_S\otimes_{\mathcal O}\mathbb F_q.
\end{equation}

\medskip
\noindent\emph{Source note.}
The print cites (4.5.24), but the congruence follows from the terminal zeta
identity (4.5.33).  It also defines $W_S$ using $R_B$, although the entire
section and every following formula concern $R_S$.  The corrected French
and English layers use (4.5.33) and $R_S$; the diplomatic French retains
the printed forms.
\end{filecontents*}
% Original input: components/180_expose_XXII_reduction_mod_p.tex
\begin{filecontents*}[overwrite,noheader]{sga7ii-cumulative-input-180.tex}
\begin{equation}\tag{6.0.3}
\bar\alpha_S:W_S\longrightarrow W_S
\end{equation}

\noindent
denotes the $\tau^{-1}$-linear operation induced by $\alpha_S$;

\begin{equation}\tag{6.0.4}
R_S(i)
=\text{the submodule of $R_S$ formed by the elements of degree $i$ in $X_0$},
\end{equation}

and

\begin{equation}\tag{6.0.5}
W_S(i)=R_S(i)\otimes_{\mathcal O}\mathbb F_q.
\end{equation}

\subsubsection*{Proposition 6.0.6.}
One has

\[
\alpha_S(R_S)\subset R_S(1)+pR_S.
\]

\noindent\emph{Proof.}
We have the factorization (4.6.14)

\begin{equation}\tag{6.0.7}
\alpha_S=\frac1p\,\tau^{-1}\circ\psi_p\circ H(X).
\end{equation}

We know (see (4.6.15)) that

\begin{equation}\tag{6.0.8}
H(X)\,R_S\subset R_S,
\end{equation}

and plainly

\begin{equation}\tag{6.0.9}
\tau^{-1}\bigl(R_S(1)+pR_S\bigr)=R_S(1)+pR_S.
\end{equation}

It is therefore enough to verify

\begin{equation}\tag{6.0.10}
\frac1p\,\psi_p(R_S)\subset R_S(1)+pR_S,
\end{equation}

which is clear from (4.6.19).

\subsubsection*{Corollary 6.0.11.}
One has
\[
\bar\alpha_S(W_S)\subset W_S(1).
\]

\subsubsection*{Corollary 6.0.12.}
One has the congruence

\[
\Delta_S(t)\equiv
\det\!\left(
1-t(\bar\alpha_S)^{\operatorname{ord}(q)}
\mathrel{\Big|}W_S(1)\right)
\pmod{\pi\mathcal O[[t]]}.
\]
\end{filecontents*}
% Original input: components/181_expose_XXII_reduced_frobenius_operator.tex
\begin{filecontents*}[overwrite,noheader]{sga7ii-cumulative-input-181.tex}
\subsubsection*{Proposition 6.1.}
The $\tau^{-1}$-linear operation
$\bar\alpha_S:W_S(1)\to W_S(1)$ is equal to

\begin{equation}\tag{6.1.1}
\frac1p\,\tau^{-1}\circ\psi_p\circ(X_0f(X))^{p-1}.
\end{equation}

\noindent\emph{Proof.}
Write

\begin{equation}\tag{6.1.2}
H(X)=1+\sum_{i\geq1}H_i(X),
\end{equation}

where

\begin{equation}\tag{6.1.3}
H_i(X)\text{ has degree $i$ in }X_0.
\end{equation}

Again by (4.6.19),

\begin{equation}\tag{6.1.4}
\frac1p\,\psi_pR_S(i)\subset pR_S
\qquad\text{if }i\ne p.
\end{equation}

Since

\begin{equation}\tag{6.1.5}
H_i(X)\,R_S(j)\subset R_S(i+j),
\end{equation}

we have

\begin{equation}\tag{6.1.6}
\left(
\alpha_S-\frac1p\,\tau^{-1}\circ\psi_p\circ H_{p-1}(X)
\right)(R_S(1))
\subset pR_S.
\end{equation}

Consequently,

\begin{equation}\tag{6.1.7}
\left[
(\alpha_S)^{\operatorname{ord}(q)}
-
\left(
\frac1p\,\tau^{-1}\circ\psi_p\circ H_{p-1}(X)
\right)^{\operatorname{ord}(q)}
\right]R_S(1)
\subset pR_S.
\end{equation}

It remains to calculate $H_{p-1}(X)$.  From (4.4.4),

\begin{equation}\tag{6.1.8}
\begin{aligned}
H(X)
&=\prod_W\Theta(\pi A_WX_0X^W)\\
&=\prod_W\exp\!\left(
\sum_{r\geq0}\frac{(\pi A_WX_0X^W)^{p^r}}{p^r}
\right)\\
&=\prod_W\prod_{r\geq0}
\exp\!\left(\frac{(\pi A_WX_0X^W)^{p^r}}{p^r}\right)\\
&\equiv\prod_W\exp(\pi A_WX_0X^W)
\pmod{\text{terms of degree }\geq p}\\
&=\exp\!\left(\pi\sum_WA_WX_0X^W\right)
\pmod{\text{terms of degree }\geq p}\\
&=\exp(\pi X_0f(X))
\pmod{\text{terms of degree }\geq p}.
\end{aligned}
\end{equation}

It follows that

\begin{equation}\tag{6.1.9}
H_{p-1}(X)
=\frac{\pi^{p-1}X_0^{p-1}}{(p-1)!}\,f(X)^{p-1}.
\end{equation}

One easily checks that

\begin{equation}\tag{6.1.10}
\frac{\pi^{p-1}}p\equiv-1\pmod p,
\end{equation}

and hence

\begin{equation}\tag{6.1.11}
\frac1p\,H_{p-1}(X)
\equiv(X_0f(X))^{p-1}\pmod p.
\end{equation}

This completes the proof of Proposition~6.1.

\subsubsection*{Corollary 6.1.12.}
The linear operator
\[
(\bar\alpha_S)^{\operatorname{ord}(q)}:W_S(1)\longrightarrow W_S(1)
\]
is precisely
\[
\psi_q\circ(X_0f(X))^{q-1}.
\]

\noindent\emph{Proof.}
By (6.1.1),

\begin{equation*}
\begin{aligned}
(\bar\alpha_S)^{\operatorname{ord}(q)}
&=
\left(
\frac1p\,\tau^{-1}\circ\psi_p\circ(X_0f(X))^{p-1}
\right)^{\operatorname{ord}(q)}\\
&=
\psi_q\circ
\prod_{j=0}^{\operatorname{ord}(q)-1}
\left(X_0^{p^j}f^{\tau^j}(X^{p^j})\right)^{p-1}\\
&=
\psi_q\circ
\prod_{j=0}^{\operatorname{ord}(q)-1}
(X_0f(X))^{p^j(p-1)}\\
&=\psi_q\circ(X_0f(X))^{q-1}.
\end{aligned}
\end{equation*}

\medskip
\noindent\emph{Source note.}
The first product in the printed proof has the undefined upper bound
$\operatorname{ord}(a)-1$; the next line and the definition of $q$ require
$\operatorname{ord}(q)-1$.  The corrected French and English layers use
the latter; the diplomatic French preserves the print.

\subsubsection*{Corollary 6.1.13.}
Modulo $p$, one has the congruence

\begin{equation*}
Z(t,V/\mathbb F_q)
\equiv
\frac{
\det\!\left(
1-t\,\psi_q\circ(X_0f(X))^{q-1}
\mathrel{\Big|}W_S(1)
\right)^{(-1)^{n+1}}
}{1-t}.
\end{equation*}

\end{filecontents*}
% Original input: components/182_expose_XXII_cohomology_of_the_structure_sheaf.tex
\begin{filecontents*}[overwrite,noheader]{sga7ii-cumulative-input-182.tex}
\subsection*{6.2.}
It remains to calculate the cohomology groups of $\mathcal O_V$.
Put

\begin{equation}\tag{6.2.1}
\mathfrak F_q
=\text{the operation of raising to the $q$-th power in $\mathcal O_V$}.
\end{equation}

There is an exact sequence of sheaves on
$\mathbb P^{n+1}=\mathbb P$,

\begin{equation}\tag{6.2.2}
0\longrightarrow\mathcal O_{\mathbb P}(-d)
\xrightarrow{\,f\,}\mathcal O_{\mathbb P}
\longrightarrow\mathcal O_V\longrightarrow0,
\end{equation}

and the following diagram is commutative:

\begin{equation}\tag{6.2.3}
\begin{tikzcd}[column sep=large,row sep=large]
0 \arrow[r] &
\mathcal O_{\mathbb P}(-d) \arrow[r,"f"] \arrow[d,"f^{q-1}\mathfrak F_q"'] &
\mathcal O_{\mathbb P} \arrow[r] \arrow[d,"\mathfrak F_q"'] &
\mathcal O_V \arrow[r] \arrow[d,"\mathfrak F_q"'] &
0\\
0 \arrow[r] &
\mathcal O_{\mathbb P}(-d) \arrow[r,"f"] &
\mathcal O_{\mathbb P} \arrow[r] &
\mathcal O_V \arrow[r] &
0 .
\end{tikzcd}
\end{equation}

Taking account of

\begin{equation}\tag{6.2.4}
\begin{cases}
H^i(\mathbb P^{n+1},\mathcal O_{\mathbb P})=0,
&i\ne0,\\
H^i(\mathbb P^{n+1},\mathcal O_{\mathbb P}(-d))=0,
&i\ne0,n+1,
\end{cases}
\end{equation}

the long exact cohomology sequence associated with (6.2.2) gives

\begin{equation}\tag{6.2.5}
\begin{cases}
H^0(V,\mathcal O_V)\simeq\mathbb F_q,\\
H^i(V,\mathcal O_V)=0,&i\ne0,n,\\
H^n(V,\mathcal O_V)
\simeq H^{n+1}(\mathbb P^{n+1},\mathcal O_{\mathbb P}(-d)).
\end{cases}
\end{equation}

The commutative diagram (6.2.3) therefore gives

\begin{equation}\tag{6.2.6}
\det\!\left(1-t\mathfrak F_q
\mathrel{\Big|}H^0(V,\mathcal O_V)\right)=1-t
\end{equation}

and

\begin{equation}\tag{6.2.7}
\det\!\left(1-t\mathfrak F_q
\mathrel{\Big|}H^n(V,\mathcal O_V)\right)
=
\det\!\left(
1-tf^{q-1}\mathfrak F_q
\mathrel{\Big|}
H^{n+1}(\mathbb P^{n+1},\mathcal O_{\mathbb P}(-d))
\right).
\end{equation}

It remains only to prove

\begin{equation}\tag{6.2.8}
\begin{aligned}
&\det\!\left(
1-tf^{q-1}\mathfrak F_q
\mathrel{\Big|}
H^{n+1}(\mathbb P^{n+1},\mathcal O_{\mathbb P}(-d))
\right)\\
&\qquad=
\det\!\left(
1-t\,\psi_q\circ(X_0f(X))^{q-1}
\mathrel{\Big|}W_S(1)
\right),
\end{aligned}
\end{equation}

which is a consequence of the following proposition.

\end{filecontents*}
% Original input: components/183_expose_XXII_proposition_6_3_and_bibliography.tex
\begin{filecontents*}[overwrite,noheader]{sga7ii-cumulative-input-183.tex}
\subsubsection*{Proposition 6.3.}
There is an isomorphism

\begin{equation}\tag{6.3.1}
W_S(1)\simeq
\operatorname{Hom}_{\mathbb F_q}\!\left(
H^{n+1}(\mathbb P^{n+1},\mathcal O_{\mathbb P}(-d)),
\mathbb F_q\right)
\end{equation}

with respect to which the operator

\begin{equation}\tag{6.3.2}
\psi_q\circ(X_0f(X))^{q-1}:W_S(1)\longrightarrow W_S(1)
\end{equation}

is the transpose of

\begin{equation}\tag{6.3.3}
f(X)^{q-1}\mathfrak F_q:
H^{n+1}(\mathbb P^{n+1},\mathcal O_{\mathbb P}(-d))
\longrightarrow
H^{n+1}(\mathbb P^{n+1},\mathcal O_{\mathbb P}(-d)).
\end{equation}

\noindent\emph{Proof.}
Computing
$H^{n+1}(\mathbb P^{n+1},\mathcal O_{\mathbb P}(-d))$
by means of the standard covering of $\mathbb P^{n+1}$ gives

\begin{equation}\tag{6.3.4}
\begin{aligned}
H^{n+1}(\mathbb P^{n+1},\mathcal O_{\mathbb P}(-d))
&=\text{the quotient of the $\mathbb F_q$-vector space with basis}\\
&\quad X_1^{W_1}\cdots X_{n+2}^{W_{n+2}},
\qquad W_i\in\mathbb Z,\quad\sum_iW_i=-d,\\
&\quad\text{by the $\mathbb F_q$-subspace with basis those monomials}\\
&\quad X_1^{W_1}\cdots X_{n+2}^{W_{n+2}}
\text{ for which }\exists i\text{ with }W_i\geq0.
\end{aligned}
\end{equation}

On the other hand, $W_S(1)=R_S(1)\otimes_{\mathcal O}\mathbb F_q$, and

\begin{equation}\tag{6.3.5}
\begin{aligned}
R_S(1)
&=\text{the free $\mathcal O$-module with basis the monomials}\\
&\quad
(\pi X_0)X_1^{U_1}\cdots X_{n+2}^{U_{n+2}},
\qquad
\sum_iU_i=d,\quad U_i>0\ \forall i.
\end{aligned}
\end{equation}

The pairing

\begin{equation}\tag{6.3.6}
(\ ,\ ):
W_S(1)\times
H^{n+1}(\mathbb P^{n+1},\mathcal O_{\mathbb P}(-d))
\longrightarrow\mathbb F_q
\end{equation}

defined by

\begin{equation}\tag{6.3.7}
\left(
\operatorname{class}\!\left(
\pi X_0X_1^{U_1}\cdots X_{n+2}^{U_{n+2}}
\right),
X_1^{W_1}\cdots X_{n+2}^{W_{n+2}}
\right)
=
\begin{cases}
1,&U_i=-W_i\text{ for every }i,\\
0,&\text{otherwise},
\end{cases}
\end{equation}

is plainly a perfect duality of $\mathbb F_q$-vector spaces.  One checks
immediately that
$\psi_q\circ(X_0f(X))^{q-1}$ is the transpose of
$f(X)^{q-1}\mathfrak F_q$, which proves Proposition~6.3.

\medskip
\noindent\emph{Source note.}
In (6.3.4) the print has $H^{i+1}$, while Proposition~6.3, (6.3.3), and the
sentence introducing the computation all require $H^{n+1}$.  The corrected
French and English layers use $n+1$; the diplomatic French preserves the
printed exponent.

\section*{Bibliography}

\begin{enumerate}[label={[\arabic*]},leftmargin=2.4em]
\item J.~Ax, \emph{Zeroes of Polynomials Over Finite Fields},
  \emph{Amer. J. Math.} \textbf{86} (1964), 255--261.
\item B.~Dwork, \emph{On the Rationality of the Zeta Function of an Algebraic
  Variety}, \emph{Amer. J. Math.} \textbf{82} (1960), 631--648.
\item B.~Dwork, \emph{An Algebraic Variety}, in \emph{Proceedings of the
  Woods Hole Summer Institute 1964} (mimeographed).
\item B.~Dwork, \emph{On the Zeta Function of a Hypersurface II},
  \emph{Ann. of Math.} (2) \textbf{80} (1964), 227--299.
\item A.~Grothendieck, \emph{Le Groupe de Brauer III}, in
  \emph{Dix Exposés sur la Cohomologie des Schémas},
  North-Holland Publishing Co., 1968.
\item J.-P.~Serre, \emph{Groupes algébriques et corps de classes},
  Hermann, Paris, 1959.
\item J.-P.~Serre, \emph{Endomorphismes complètement continus des espaces de
  Banach $p$-adiques}, \emph{Publ. Math. I.H.E.S.} no.~12 (1962).
\item C.~Chevalley, \emph{Démonstration d'une hypothèse de M.~Artin},
  \emph{Abh. Math. Sem. Univ. Hamburg} \textbf{11} (1936), 73--75.
\end{enumerate}
\end{filecontents*}
% SGA 7 II — English source-first working reader
% Controlling French input: diplomatic and corrected French layers through
% scan index 445 and volume EOF.
% Ultimate authority: Number12.pdf, SHA-256
% FA679DEBFC8ADA3232D7E752A1837FC6CE474488E20A44D7641CF296876E1297.
\documentclass[11pt,a4paper,leqno]{article}

\usepackage[T1]{fontenc}
\usepackage[utf8]{inputenc}
\usepackage[french,english]{babel}
\usepackage{amsmath,amssymb,amsfonts,mathrsfs}
\usepackage{mathtools}
\usepackage{tikz}
\usetikzlibrary{cd,arrows.meta}
\usepackage{geometry}
\geometry{margin=2.4cm}
\usepackage{enumitem}
\usepackage[normalem]{ulem}
\usepackage{hyperref}
\hypersetup{
  colorlinks=true,
  linkcolor=blue,
  citecolor=blue,
  urlcolor=blue,
  hypertexnames=false,
  pdftitle={SGA 7-II: Monodromy Groups in Algebraic Geometry --- English Reader}
}
\let\mathunderline\underline
\renewcommand{\underline}[1]{\ifmmode\mathunderline{#1}\else\uline{#1}\fi}
\emergencystretch=3em

\providecommand{\Gal}{\operatorname{Gal}}
\providecommand{\Spec}{\operatorname{Spec}}
\providecommand{\Hom}{\operatorname{Hom}}
\providecommand{\Ext}{\operatorname{Ext}}
\providecommand{\Tor}{\operatorname{Tor}}
\providecommand{\Ker}{\operatorname{Ker}}
\providecommand{\Coker}{\operatorname{Coker}}
\providecommand{\Im}{\operatorname{Im}}
\providecommand{\codim}{\operatorname{codim}}
\providecommand{\pr}{\operatorname{pr}}
\providecommand{\id}{\mathrm{id}}

\begin{document}
\selectlanguage{english}

\begin{titlepage}
\centering
\vspace*{2.2in}
{\LARGE\textbf{SGA 7-II}}\\[0.6em]
{\Large Monodromy Groups in Algebraic Geometry}\\[1.5em]
{\large Complete English reader}\\[1.2em]
{\normalsize Exposés X--XXII}
\vfill
\end{titlepage}
\tableofcontents
\clearpage

\phantomsection\addcontentsline{toc}{section}{Exposé X}
\input{sga7ii-cumulative-input-001.tex}
\input{sga7ii-cumulative-input-002.tex}
\input{sga7ii-cumulative-input-003.tex}
\input{sga7ii-cumulative-input-004.tex}
\input{sga7ii-cumulative-input-005.tex}
\input{sga7ii-cumulative-input-006.tex}
\input{sga7ii-cumulative-input-007.tex}
\input{sga7ii-cumulative-input-008.tex}
\input{sga7ii-cumulative-input-009.tex}
\phantomsection\addcontentsline{toc}{section}{Exposé XI}
\input{sga7ii-cumulative-input-010.tex}
\input{sga7ii-cumulative-input-011.tex}
\input{sga7ii-cumulative-input-012.tex}
\input{sga7ii-cumulative-input-013.tex}
\input{sga7ii-cumulative-input-014.tex}
\input{sga7ii-cumulative-input-015.tex}
\input{sga7ii-cumulative-input-016.tex}
\input{sga7ii-cumulative-input-017.tex}
\phantomsection\addcontentsline{toc}{section}{Exposé XII}
\input{sga7ii-cumulative-input-018.tex}
\input{sga7ii-cumulative-input-019.tex}
\input{sga7ii-cumulative-input-020.tex}
\input{sga7ii-cumulative-input-021.tex}
\input{sga7ii-cumulative-input-022.tex}
\input{sga7ii-cumulative-input-023.tex}
\phantomsection\addcontentsline{toc}{section}{Exposé XIII}
\input{sga7ii-cumulative-input-024.tex}
\input{sga7ii-cumulative-input-025.tex}
\input{sga7ii-cumulative-input-026.tex}
\input{sga7ii-cumulative-input-027.tex}
\input{sga7ii-cumulative-input-028.tex}
\input{sga7ii-cumulative-input-029.tex}
\input{sga7ii-cumulative-input-030.tex}
\input{sga7ii-cumulative-input-031.tex}
\input{sga7ii-cumulative-input-032.tex}
\input{sga7ii-cumulative-input-033.tex}
\input{sga7ii-cumulative-input-034.tex}
\input{sga7ii-cumulative-input-035.tex}
\input{sga7ii-cumulative-input-036.tex}
\input{sga7ii-cumulative-input-037.tex}
\phantomsection\addcontentsline{toc}{section}{Exposé XIV}
\input{sga7ii-cumulative-input-038.tex}
\input{sga7ii-cumulative-input-039.tex}
\input{sga7ii-cumulative-input-040.tex}
\input{sga7ii-cumulative-input-041.tex}
\input{sga7ii-cumulative-input-042.tex}
\input{sga7ii-cumulative-input-043.tex}
\input{sga7ii-cumulative-input-044.tex}
\input{sga7ii-cumulative-input-045.tex}
\input{sga7ii-cumulative-input-046.tex}
\input{sga7ii-cumulative-input-047.tex}
\input{sga7ii-cumulative-input-048.tex}
\input{sga7ii-cumulative-input-049.tex}
\input{sga7ii-cumulative-input-050.tex}
\input{sga7ii-cumulative-input-051.tex}
\input{sga7ii-cumulative-input-052.tex}
\input{sga7ii-cumulative-input-053.tex}
\input{sga7ii-cumulative-input-054.tex}
\input{sga7ii-cumulative-input-055.tex}
\input{sga7ii-cumulative-input-056.tex}
\input{sga7ii-cumulative-input-057.tex}
\input{sga7ii-cumulative-input-058.tex}
\input{sga7ii-cumulative-input-059.tex}
\input{sga7ii-cumulative-input-060.tex}
\input{sga7ii-cumulative-input-061.tex}
\input{sga7ii-cumulative-input-062.tex}
\input{sga7ii-cumulative-input-063.tex}
\phantomsection\addcontentsline{toc}{section}{Exposé XV}
\input{sga7ii-cumulative-input-064.tex}
\input{sga7ii-cumulative-input-065.tex}
\input{sga7ii-cumulative-input-066.tex}
\input{sga7ii-cumulative-input-067.tex}
\input{sga7ii-cumulative-input-068.tex}
\input{sga7ii-cumulative-input-069.tex}
\input{sga7ii-cumulative-input-070.tex}
\input{sga7ii-cumulative-input-071.tex}
\input{sga7ii-cumulative-input-072.tex}
\input{sga7ii-cumulative-input-073.tex}
\input{sga7ii-cumulative-input-074.tex}
\input{sga7ii-cumulative-input-075.tex}
\input{sga7ii-cumulative-input-076.tex}
\input{sga7ii-cumulative-input-077.tex}
\input{sga7ii-cumulative-input-078.tex}
\input{sga7ii-cumulative-input-079.tex}
\input{sga7ii-cumulative-input-080.tex}
\input{sga7ii-cumulative-input-081.tex}
\phantomsection\addcontentsline{toc}{section}{Exposé XVI}
\input{sga7ii-cumulative-input-082.tex}
\input{sga7ii-cumulative-input-083.tex}
\input{sga7ii-cumulative-input-084.tex}
\input{sga7ii-cumulative-input-085.tex}
\input{sga7ii-cumulative-input-086.tex}
\input{sga7ii-cumulative-input-087.tex}
\input{sga7ii-cumulative-input-088.tex}
\input{sga7ii-cumulative-input-089.tex}
\phantomsection\addcontentsline{toc}{section}{Exposé XVII}
\input{sga7ii-cumulative-input-090.tex}
\input{sga7ii-cumulative-input-091.tex}
\input{sga7ii-cumulative-input-092.tex}
\input{sga7ii-cumulative-input-093.tex}
\input{sga7ii-cumulative-input-094.tex}
\input{sga7ii-cumulative-input-095.tex}
\input{sga7ii-cumulative-input-096.tex}
\input{sga7ii-cumulative-input-097.tex}
\input{sga7ii-cumulative-input-098.tex}
\input{sga7ii-cumulative-input-099.tex}
\input{sga7ii-cumulative-input-100.tex}
\input{sga7ii-cumulative-input-101.tex}
\input{sga7ii-cumulative-input-102.tex}
\input{sga7ii-cumulative-input-103.tex}
\input{sga7ii-cumulative-input-104.tex}
\input{sga7ii-cumulative-input-105.tex}
\input{sga7ii-cumulative-input-106.tex}
\input{sga7ii-cumulative-input-107.tex}
\input{sga7ii-cumulative-input-108.tex}
\phantomsection\addcontentsline{toc}{section}{Exposé XVIII}
\input{sga7ii-cumulative-input-109.tex}
\input{sga7ii-cumulative-input-110.tex}
\input{sga7ii-cumulative-input-111.tex}
\input{sga7ii-cumulative-input-112.tex}
\input{sga7ii-cumulative-input-113.tex}
\input{sga7ii-cumulative-input-114.tex}
\input{sga7ii-cumulative-input-115.tex}
\input{sga7ii-cumulative-input-116.tex}
\input{sga7ii-cumulative-input-117.tex}
\input{sga7ii-cumulative-input-118.tex}
\input{sga7ii-cumulative-input-119.tex}
\input{sga7ii-cumulative-input-120.tex}
\input{sga7ii-cumulative-input-121.tex}
\input{sga7ii-cumulative-input-122.tex}
\input{sga7ii-cumulative-input-123.tex}
\input{sga7ii-cumulative-input-124.tex}
\input{sga7ii-cumulative-input-125.tex}
\input{sga7ii-cumulative-input-126.tex}
\input{sga7ii-cumulative-input-127.tex}
\input{sga7ii-cumulative-input-128.tex}
\input{sga7ii-cumulative-input-129.tex}
\input{sga7ii-cumulative-input-130.tex}
\input{sga7ii-cumulative-input-131.tex}
\input{sga7ii-cumulative-input-132.tex}
\input{sga7ii-cumulative-input-133.tex}
\phantomsection\addcontentsline{toc}{section}{Exposé XIX}
\input{sga7ii-cumulative-input-134.tex}
\input{sga7ii-cumulative-input-135.tex}
\input{sga7ii-cumulative-input-136.tex}
\input{sga7ii-cumulative-input-137.tex}
\input{sga7ii-cumulative-input-138.tex}
\phantomsection\addcontentsline{toc}{section}{Exposé XX}
\input{sga7ii-cumulative-input-139.tex}
\input{sga7ii-cumulative-input-140.tex}
\input{sga7ii-cumulative-input-141.tex}
\input{sga7ii-cumulative-input-142.tex}
\input{sga7ii-cumulative-input-143.tex}
\input{sga7ii-cumulative-input-144.tex}
\input{sga7ii-cumulative-input-145.tex}
\input{sga7ii-cumulative-input-146.tex}
\phantomsection\addcontentsline{toc}{section}{Exposé XXI}
\input{sga7ii-cumulative-input-147.tex}
\input{sga7ii-cumulative-input-148.tex}
\input{sga7ii-cumulative-input-149.tex}
\input{sga7ii-cumulative-input-150.tex}
\input{sga7ii-cumulative-input-151.tex}
\input{sga7ii-cumulative-input-152.tex}
\input{sga7ii-cumulative-input-153.tex}
\input{sga7ii-cumulative-input-154.tex}
\input{sga7ii-cumulative-input-155.tex}
\input{sga7ii-cumulative-input-156.tex}
\input{sga7ii-cumulative-input-157.tex}
\input{sga7ii-cumulative-input-158.tex}
\input{sga7ii-cumulative-input-159.tex}
\input{sga7ii-cumulative-input-160.tex}
\input{sga7ii-cumulative-input-161.tex}
\input{sga7ii-cumulative-input-162.tex}
\phantomsection\addcontentsline{toc}{section}{Exposé XXII}
\input{sga7ii-cumulative-input-163.tex}
\input{sga7ii-cumulative-input-164.tex}
\input{sga7ii-cumulative-input-165.tex}
\input{sga7ii-cumulative-input-166.tex}
\input{sga7ii-cumulative-input-167.tex}
\input{sga7ii-cumulative-input-168.tex}
\input{sga7ii-cumulative-input-169.tex}
\input{sga7ii-cumulative-input-170.tex}
\input{sga7ii-cumulative-input-171.tex}
\input{sga7ii-cumulative-input-172.tex}
\input{sga7ii-cumulative-input-173.tex}
\input{sga7ii-cumulative-input-174.tex}
\input{sga7ii-cumulative-input-175.tex}
\input{sga7ii-cumulative-input-176.tex}
\input{sga7ii-cumulative-input-177.tex}
\input{sga7ii-cumulative-input-178.tex}
\input{sga7ii-cumulative-input-179.tex}
\input{sga7ii-cumulative-input-180.tex}
\input{sga7ii-cumulative-input-181.tex}
\input{sga7ii-cumulative-input-182.tex}
\input{sga7ii-cumulative-input-183.tex}

\end{document}

% SGA-CUMULATIVE-END 09_SGA7II_EN_COMPLETE_READER.tex
